% Options for packages loaded elsewhere
% Options for packages loaded elsewhere
\PassOptionsToPackage{unicode}{hyperref}
\PassOptionsToPackage{hyphens}{url}
%
\documentclass[
  english,
  letterpaper,
]{scrbook}
\usepackage{xcolor}
\usepackage[margin=1in]{geometry}
\usepackage{amsmath,amssymb}
\setcounter{secnumdepth}{-\maxdimen} % remove section numbering
\usepackage{iftex}
\ifPDFTeX
  \usepackage[T1]{fontenc}
  \usepackage[utf8]{inputenc}
  \usepackage{textcomp} % provide euro and other symbols
\else % if luatex or xetex
  \usepackage{unicode-math} % this also loads fontspec
  \defaultfontfeatures{Scale=MatchLowercase}
  \defaultfontfeatures[\rmfamily]{Ligatures=TeX,Scale=1}
\fi
\usepackage{lmodern}
\ifPDFTeX\else
  % xetex/luatex font selection
\fi
%% Support for zero-width non-joiner characters.
\makeatletter
\def\zerowidthnonjoiner{%
  % Prevent ligatures and adjust kerning, but still support hyphenating.
  \texorpdfstring{%
    \TextOrMath{\nobreak\discretionary{-}{}{\kern.03em}%
      \ifvmode\else\nobreak\hskip\z@skip\fi}{}%
  }{}%
}
\makeatother
\ifPDFTeX
  \DeclareUnicodeCharacter{200C}{\zerowidthnonjoiner}
\else
  \catcode`^^^^200c=\active
  \protected\def ^^^^200c{\zerowidthnonjoiner}
\fi
%% End of ZWNJ support
% Use upquote if available, for straight quotes in verbatim environments
\IfFileExists{upquote.sty}{\usepackage{upquote}}{}
\IfFileExists{microtype.sty}{% use microtype if available
  \usepackage[]{microtype}
  \UseMicrotypeSet[protrusion]{basicmath} % disable protrusion for tt fonts
}{}
\makeatletter
\@ifundefined{KOMAClassName}{% if non-KOMA class
  \IfFileExists{parskip.sty}{%
    \usepackage{parskip}
  }{% else
    \setlength{\parindent}{0pt}
    \setlength{\parskip}{6pt plus 2pt minus 1pt}}
}{% if KOMA class
  \KOMAoptions{parskip=half}}
\makeatother
% Make \paragraph and \subparagraph free-standing
\makeatletter
\ifx\paragraph\undefined\else
  \let\oldparagraph\paragraph
  \renewcommand{\paragraph}{
    \@ifstar
      \xxxParagraphStar
      \xxxParagraphNoStar
  }
  \newcommand{\xxxParagraphStar}[1]{\oldparagraph*{#1}\mbox{}}
  \newcommand{\xxxParagraphNoStar}[1]{\oldparagraph{#1}\mbox{}}
\fi
\ifx\subparagraph\undefined\else
  \let\oldsubparagraph\subparagraph
  \renewcommand{\subparagraph}{
    \@ifstar
      \xxxSubParagraphStar
      \xxxSubParagraphNoStar
  }
  \newcommand{\xxxSubParagraphStar}[1]{\oldsubparagraph*{#1}\mbox{}}
  \newcommand{\xxxSubParagraphNoStar}[1]{\oldsubparagraph{#1}\mbox{}}
\fi
\makeatother


\usepackage{longtable,booktabs,array}
\newcounter{none} % for unnumbered tables
\usepackage{calc} % for calculating minipage widths
% Correct order of tables after \paragraph or \subparagraph
\usepackage{etoolbox}
\makeatletter
\patchcmd\longtable{\par}{\if@noskipsec\mbox{}\fi\par}{}{}
\makeatother
% Allow footnotes in longtable head/foot
\IfFileExists{footnotehyper.sty}{\usepackage{footnotehyper}}{\usepackage{footnote}}
\makesavenoteenv{longtable}
\usepackage{graphicx}
\makeatletter
\newsavebox\pandoc@box
\newcommand*\pandocbounded[1]{% scales image to fit in text height/width
  \sbox\pandoc@box{#1}%
  \Gscale@div\@tempa{\textheight}{\dimexpr\ht\pandoc@box+\dp\pandoc@box\relax}%
  \Gscale@div\@tempb{\linewidth}{\wd\pandoc@box}%
  \ifdim\@tempb\p@<\@tempa\p@\let\@tempa\@tempb\fi% select the smaller of both
  \ifdim\@tempa\p@<\p@\scalebox{\@tempa}{\usebox\pandoc@box}%
  \else\usebox{\pandoc@box}%
  \fi%
}
% Set default figure placement to htbp
\def\fps@figure{htbp}
\makeatother



\ifLuaTeX
\usepackage[bidi=basic,shorthands=off,bidi=bidi]{babel}
\else
\usepackage[bidi=default,shorthands=off,bidi=bidi]{babel}
\fi
\ifLuaTeX
  \usepackage{selnolig} % disable illegal ligatures
\fi


\setlength{\emergencystretch}{3em} % prevent overfull lines

\providecommand{\tightlist}{%
  \setlength{\itemsep}{0pt}\setlength{\parskip}{0pt}}



 


\babelprovide[import]{arabic}
\babelprovide[import]{persian}
\babelfont[arabic]{rm}{Amiri}
\babelfont[persian]{rm}{Amiri}
% Avoid the KOMA-Script/bidi title table conflict established in BH11924.
\renewcommand{\maketitle}{}
\makeatletter
\@ifpackageloaded{bookmark}{}{\usepackage{bookmark}}
\makeatother
\makeatletter
\@ifpackageloaded{caption}{}{\usepackage{caption}}
\AtBeginDocument{%
\ifdefined\contentsname
  \renewcommand*\contentsname{Table of contents}
\else
  \newcommand\contentsname{Table of contents}
\fi
\ifdefined\listfigurename
  \renewcommand*\listfigurename{List of Figures}
\else
  \newcommand\listfigurename{List of Figures}
\fi
\ifdefined\listtablename
  \renewcommand*\listtablename{List of Tables}
\else
  \newcommand\listtablename{List of Tables}
\fi
\ifdefined\figurename
  \renewcommand*\figurename{Figure}
\else
  \newcommand\figurename{Figure}
\fi
\ifdefined\tablename
  \renewcommand*\tablename{Table}
\else
  \newcommand\tablename{Table}
\fi
}
\@ifpackageloaded{float}{}{\usepackage{float}}
\floatstyle{ruled}
\@ifundefined{c@chapter}{\newfloat{codelisting}{h}{lop}}{\newfloat{codelisting}{h}{lop}[chapter]}
\floatname{codelisting}{Listing}
\newcommand*\listoflistings{\listof{codelisting}{List of Listings}}
\makeatother
\makeatletter
\makeatother
\makeatletter
\@ifpackageloaded{caption}{}{\usepackage{caption}}
\@ifpackageloaded{subcaption}{}{\usepackage{subcaption}}
\makeatother
\usepackage{bookmark}
\IfFileExists{xurl.sty}{\usepackage{xurl}}{} % add URL line breaks if available
\urlstyle{same}
% fallback for those not using the hyperref driver hyperxmp:
\makeatletter
\@ifundefined{xmpquote}{\newcommand{\xmpquote}[1]{#1}}{}
\makeatother
\hypersetup{
  pdftitle={Translations of Bahá'u'lláh},
  pdflang={en},
  hidelinks,
  pdfcreator={LaTeX via pandoc}}


\title{Translations of Bahá'u'lláh}
\author{}
\date{}
\begin{document}
\frontmatter
\maketitle

\renewcommand*\contentsname{Table of contents}
{
\setcounter{tocdepth}{2}
\tableofcontents
}

\mainmatter
\bookmarksetup{startatroot}

\chapter{Introduction}\label{sec-introduction}

This collection brings together the currently available publication
editions of Bahá'u'lláh's writings. The tablets are arranged in
numerical order by their BH identifiers. The order is provisional and
does not imply chronology.

The texts preserve the original-language passage, transliteration, and
contemporary English translation from the publication editions.

\textbf{Tablets in this edition:} 16.

\bookmarksetup{startatroot}

\chapter{The Kitab-i-Aqdas (The Most Holy
Book)}\label{the-kitab-i-aqdas-the-most-holy-book}

\section{BH00001}\label{bh00001}

\emph{Translated on 9 Qudrat, 183 BE using GPT 6.1 Sol}

\subsection{1}\label{section}

\begin{otherlanguage}{arabic}

بسمه الحاکم علی ما کان و ما یکون

\end{otherlanguage}

Bismihi al-Ḥākim ʿalā mā kāna wa-mā yakūn.

In His name, the Ruler over all that has been and all that will be.

\subsection{2}\label{section-1}

\begin{otherlanguage}{arabic}

انّ اوّل ما کتب اللّه علی العباد عرفان مشرق وحیه و مطلع امره الّذی کان مقام
نفسه فی عالم الأمر و الخلق من فاز به قد فاز بکلّ الخیر و الّذی منع انّه من
اهل الضّلال ولو یأتی بکلّ الأعمال اذا فزتم بهذا المقام الأسنی و الأفق
الأعلی ینبغی لکلّ نفسٍ ان یتّبع ما امر به من لدی المقصود لأنّهما معا لا یقبل
احدهما دون الآخر هذا ما حکم به مطلع الألهام

\end{otherlanguage}

Inna awwala mā kataba Allāhu ʿalā al-ʿibād ʿirfānu mashriqi waḥyihi
wa-maṭlaʿi amrihi alladhī kāna maqāma nafsihi fī ʿālami al-amri
wa-al-khalq. Man fāza bihi qad fāza bi-kulli al-khayr wa-alladhī muniʿa
innahu min ahli al-ḍalāl wa-law yaʾtī bi-kulli al-aʿmāl. Idhā fuztum
bi-hādhā al-maqāmi al-asnā wa-al-ufuqi al-aʿlā yanbaghī li-kulli nafsin
an yattabiʿa mā umira bihi min ladā al-Maqṣūd li-annahumā maʿan lā
yuqbalu aḥaduhumā dūna al-ākhar. Hādhā mā ḥakama bihi maṭlaʿu al-ilhām.

The first duty God has prescribed for His servants is to recognize the
Dayspring of His revelation and the Dawning Place of His Cause, who
represents His own Self in the world of command and creation. Whoever
attains this has attained all good, and whoever is deprived of it is
among those who have gone astray, even if he performs every deed. Once
you have attained this most exalted station and highest horizon, every
soul must follow what has been commanded by the Desired One, for these
two belong together: neither is accepted without the other. This is what
the Dawning Place of inspiration has decreed.

\subsection{3}\label{section-2}

\begin{otherlanguage}{arabic}

انّ الّذین اوتوا بصآئر من اللّه یرون حدود اللّه السّبب الأعظم لنظم العالم و
حفظ الأمم و الّذی غفل انّه من همج رعاع انّا امرناکم بکسر حدودات النّفس و
الهوی لا ما رقم من القلم الأعلی انّه لروح الحیوان لمن فی الأمکان قد ماجت
بحور الحکمة و البیان بما هاجت نسمة الرّحمن اغتنموا یا اولی الألباب انّ
الّذین نکثوا عهد اللّه فی اوامره و نکصوا علی اعقابهم اولٓئک من اهل الضّلال
لدی الغنیّ المتعال

\end{otherlanguage}

Inna alladhīna ūtū baṣāʾira min Allāh yarawna ḥudūda Allāhi al-sababa
al-aʿẓama li-naẓmi al-ʿālami wa-ḥifẓi al-umam wa-alladhī ghafala innahu
min hamajin raʿāʿ. Innā amarnākum bi-kasri ḥudūdāti al-nafsi wa-al-hawā
lā mā ruqima mina al-qalami al-aʿlā. Innahu la-rūḥu al-ḥayawān li-man fī
al-imkān. Qad mājat buḥūru al-ḥikmati wa-al-bayān bi-mā hājat nasmatu
al-Raḥmān. Ightanimū yā ulī al-albāb. Inna alladhīna nakathū ʿahda
Allāhi fī awāmirihi wa-nakaṣū ʿalā aʿqābihim ulāʾika min ahli al-ḍalāl
ladā al-Ghaniyyi al-Mutaʿāl.

Those endowed by God with insight see God's laws as the greatest means
of ordering the world and safeguarding its peoples. Whoever is heedless
belongs among the ignorant rabble. We have commanded you to break the
bounds of self and desire, not what has been inscribed by the Most
Exalted Pen. It is indeed the spirit of life for all who dwell in the
realm of existence. The oceans of wisdom and utterance have surged as
the breeze of the All-Merciful has stirred. Seize this opportunity,
people of understanding. Those who have broken God's covenant concerning
His commandments and turned back on their heels are among those who have
gone astray in the sight of the Self-Sufficient, the Most Exalted.

\subsection{4}\label{section-3}

\begin{otherlanguage}{arabic}

یا ملاء الأرض اعلموا انّ اوامری سرج عنایتی بین عبادی و مفاتیح رحمتی
لبریتی کذلک نزّل الأمر من سمآء مشیة ربّکم مالک الأدیان لو یجد احد حلاوة
البیان الّذی ظهر من فم مشیة الرّحمن لینفق ما عنده ولو یکون خزائن الأرض
کلّها لیثبت امراً من اوامره المشرقة من افق العنایة و الألطاف

\end{otherlanguage}

Yā malaʾa al-arḍ iʿlamū anna awāmirī suruju ʿināyatī bayna ʿibādī
wa-mafātīḥu raḥmatī li-bariyyatī. Kadhālika nuzzila al-amru min samāʾi
mashīʾati rabbikum Māliki al-adyān. Law yajidu aḥadun ḥalāwata al-bayāni
alladhī ẓahara min fami mashīʾati al-Raḥmān la-yunfiqu mā ʿindahu wa-law
yakūnu khazāʾina al-arḍi kullahā li-yuthbita amran min awāmirihi
al-mushriqati min ufuqi al-ʿināyati wa-al-alṭāf.

People of the earth, know that My commandments are the lamps of My
loving care among My servants and the keys to My mercy for My creatures.
Thus has the command been sent down from the heaven of the will of your
Lord, the Master of religions. If anyone were to taste the sweetness of
the utterance that has issued from the mouth of the will of the
All-Merciful, he would expend all he possesses, even if it were all the
treasures of the earth, to establish a single one of His commandments
that shine from the horizon of loving care and favor.

\subsection{5}\label{section-4}

\begin{otherlanguage}{arabic}

قل من حدودی یمرّ عرف قمیصی و بها تنصب اعلام النّصر علی القنن و الأتلال قد
تکلّم لسان قدرتی فی جبروت عظمتی مخاطباً لبریتی ان اعملوا حدودی حبّاً لجمالی
طوبی لحبیب وجد عرف المحبوب من هذه الکلمة الّتی فاحت منها نفحات الفضل علی
شأنٍ لا توصف بالأذکار لعمری من شرب رحیق الأنصاف من ایادی الألطاف انّه یطوف
حول اوامری المشرقة من افق الأبداع

\end{otherlanguage}

Qul min ḥudūdī yamurru ʿarfu qamīṣī wa-bihā tunṣabu aʿlāmu al-naṣri ʿalā
al-qunani wa-al-atlāl. Qad takallama lisānu qudratī fī jabarūti ʿaẓamatī
mukhāṭiban li-bariyyatī an iʿmalū ḥudūdī ḥubban li-jamālī. Ṭūbā
li-ḥabībin wajada ʿarfa al-Maḥbūbi min hādhihi al-kalimati allatī fāḥat
minhā nafaḥātu al-faḍli ʿalā shaʾnin lā tūṣafu bi-al-adhkār. La-ʿamrī
man shariba raḥīqa al-inṣāfi min ayādī al-alṭāf innahu yaṭūfu ḥawla
awāmirī al-mushriqati min ufuqi al-ibdāʿ.

Say: From My laws comes the fragrance of My garment, and through them
the banners of victory are raised upon peaks and hills. The Tongue of My
power has spoken in the dominion of My grandeur, addressing My
creatures: ``Observe My laws out of love for My beauty.'' Blessed is the
lover who has caught the fragrance of the Beloved from this utterance,
from which the breezes of grace have wafted in a manner beyond
description. By My life, whoever drinks the choice wine of fairness from
the hands of loving favor will circle around My commandments that shine
from the horizon of creation.

\subsection{6}\label{section-5}

\begin{otherlanguage}{arabic}

لا تحسبنّ انّا نزّلنا لکم الأحکام بل فتحنا ختم الرّحیق المختوم باصابع القدرة
و الأقتدار یشهد بذلک ما نزّل من قلم الوحی تفکّروا یا اولی الأفکار

\end{otherlanguage}

Lā taḥsabanna annā nazzalnā lakumu al-aḥkām bal fataḥnā khatma al-raḥīqi
al-makhtūmi bi-aṣābiʿi al-qudrati wa-al-iqtidār. Yashhadu bi-dhālika mā
nuzzila min qalami al-waḥy. Tafakkarū yā ulī al-afkār.

Do not think that We have revealed mere laws to you. Rather, with the
fingers of power and might, We have broken the seal of the sealed choice
wine. What has been sent down by the Pen of revelation bears witness to
this. Reflect, people of thought.

\subsection{7}\label{section-6}

\begin{otherlanguage}{arabic}

قد کتب علیکم الصّلوة تسع رکعات للّه منزل الآیات حین الزّوال و فی البکور و
الآصال و عفونا عدّة اخری امراً فی کتاب اللّه انّه لهو الأمر المقتدر المختار
و اذا اردتم الصّلوة ولّوا وجوهکم شطری الأقدس المقام المقدّس الّذی جعله اللّه
مطاف الملاء الأعلی و مقبل اهل مدائن البقآء و مصدر الأمر لمن فی الأرضین و
السّموات و عند غروب شمس الحقیقة و التّبیان المقرّ الّذی قدّرناه لکم انّه لهو
العزیز العلاّم

\end{otherlanguage}

Qad kutiba ʿalaykumu al-ṣalātu tisʿa rakaʿātin li-Allāhi Munzili al-āyāt
ḥīna al-zawāli wa-fī al-bukūri wa-al-āṣāl. Wa-ʿafawnā ʿiddatan ukhrā
amran fī kitābi Allāh. Innahu la-huwa al-Āmiru al-Muqtadiru al-Mukhtār.
Wa-idhā aradtumu al-ṣalāta wallū wujūhakum shaṭrī al-aqdasa al-maqāma
al-muqaddasa alladhī jaʿalahu Allāhu maṭāfa al-malaʾi al-aʿlā wa-muqbila
ahli madāʾini al-baqāʾ wa-maṣdara al-amri li-man fī al-araḍīna
wa-al-samāwāt. Wa-ʿinda ghurūbi shamsi al-ḥaqīqati wa-al-tibyān
al-maqarra alladhī qaddarnāhu lakum. Innahu la-huwa al-ʿAzīzu al-ʿAllām.

Obligatory prayer has been prescribed for you, consisting of nine units,
offered to God, the Revealer of verses, at noon, in the morning, and in
the late afternoon. We have exempted you from another number, as a
decree in the Book of God. He is indeed the Ordainer, the All-Powerful,
who freely chooses. When you wish to perform obligatory prayer, turn
your faces toward My Most Holy presence, the sacred Station that God has
made the center around which the Concourse on high circles, the point
toward which the inhabitants of the cities of eternity turn, and the
source of command for all in the earth and the heavens. At the setting
of the Sun of truth and explanation, turn toward the place We have
appointed for you. He is indeed the Almighty, the All-Knowing.

\subsection{8}\label{section-7}

\begin{otherlanguage}{arabic}

کلّ شئ تحقّق بامره المبرم اذا اشرقت من افق البیان شمس الأحکام لکلٍ ان
یتّبعوها ولو بامرٍ تنفطر عنه سموات افئدة الأدیان انّه یفعل ما یشآء و لا
یسئل عمّا شآء و ما حکم به المحبوب انّه لمحبوب و مالک الاختراع انّ الّذی وجد
عرف الرّحمن و عرف مطلع هذا البیان انّه یستقبل بعینه السّهام لاثبات الأحکام
بین الأنام طوبی لمن اقبل و فاز بفصل الخطاب

\end{otherlanguage}

Kullu shayʾin taḥaqqaqa bi-amrihi al-mubram. Idhā ashraqat min ufuqi
al-bayāni shamsu al-aḥkām li-kullin an yattabiʿūhā wa-law bi-amrin
tanfaṭiru ʿanhu samāwātu afʾidati al-adyān. Innahu yafʿalu mā yashāʾ
wa-lā yusʾalu ʿammā shāʾ wa-mā ḥakama bihi al-Maḥbūbu innahu la-maḥbūb
wa-Māliki al-ikhtirāʿ. Inna alladhī wajada ʿarfa al-Raḥmān wa-ʿarafa
maṭlaʿa hādhā al-bayān innahu yastaqbilu bi-ʿaynihi al-sihāma li-ithbāti
al-aḥkāmi bayna al-anām. Ṭūbā li-man aqbala wa-fāza bi-faṣli al-khiṭāb.

Everything has come into being through His irrevocable command. When the
Sun of laws shines from the horizon of utterance, everyone must follow
them, even if a command should cause the heavens of the hearts of
religions to split apart. He does whatever He wills and is not
questioned about what He wills. Whatever the Beloved decrees is indeed
beloved, by the Lord of creation! Whoever has caught the fragrance of
the All-Merciful and recognized the Dawning Place of this utterance will
face arrows with his own eyes to establish the laws among humanity.
Blessed is the one who has turned toward Him and attained the decisive
word.

\subsection{9}\label{section-8}

\begin{otherlanguage}{arabic}

یقرء ما نزّل قبلها والاّعفا اللّه عنه انّه له قد فصّلنا الصّلوة فی ورقةٍ اخری
طوبی لمن عمل بما امر به من لدن مالک الرّقاب قد نزّلت فی صلوة المیت ستّ
تکبیرات من اللّه منزل الایات و الّذی عنده علم القرآئة له ان و العزیز
الغفّار

\end{otherlanguage}

Yaqraʾu mā nuzzila qablahā wa-illā ʿafā Allāhu ʿanhu {[}corrupted
text{]} qad faṣṣalnā al-ṣalāta fī waraqatin ukhrā ṭūbā li-man ʿamila
bi-mā umira bihi min ladun Māliki al-riqāb. Qad nuzzilat fī ṣalāti
al-mayyiti sittu takbīrātin min Allāhi Munzili al-āyāt wa-alladhī
ʿindahu ʿilmu al-qirāʾati {[}corrupted text{]} al-ʿAzīzu al-Ghaffār.

He reads what was revealed before it; otherwise, God has pardoned him.
{[}corrupted original text{]} We have set out the obligatory prayer in
detail on another sheet. Blessed is the one who acts according to what
he has been commanded by the Master of all people. Six repetitions of
``God is Most Great'' have been revealed for the prayer for the dead by
God, the Revealer of verses, and the one who knows how to read
{[}corrupted original text{]} the Almighty, the All-Forgiving.

\subsection{10}\label{section-9}

\begin{otherlanguage}{arabic}

لا یبطل الشّعر صلوتکم و لا ما منع عن الرّوح مثل العظام و غیرها البسوا
السّمور کما تلبسون الخزّ و السّنجاب و ما دونهما انّه ما نهی فی الفرقان ولکن
اشتبه علی العلمآء انّه لهو العزیز العلاّم

\end{otherlanguage}

Lā yubṭilu al-shaʿru ṣalātakum wa-lā mā muniʿa ʿani al-rūḥi mithla
al-ʿiẓāmi wa-ghayrihā. Ilbasū al-sammūra kamā talbasūna al-khazza
wa-al-sinjāba wa-mā dūnahumā. Innahu mā nahā fī al-Furqān wa-lākin
ishtabaha ʿalā al-ʿulamāʾ. Innahu la-huwa al-ʿAzīzu al-ʿAllām.

Hair does not invalidate your prayer, nor does anything devoid of life,
such as bones and the like. Wear sable as you wear silk, squirrel fur,
and other materials. He did not forbid this in the Qur'an, but the
scholars misunderstood. He is indeed the Almighty, the All-Knowing.

\subsection{11}\label{section-10}

\begin{otherlanguage}{arabic}

قد فرض علیکم الصّلوة و الصّوم من اوّل البلوغ امراً من لدی اللّه ربّکم و ربّ
ابائکم الأولین من کان فی نفسه ضعف من المرض او الهرم عفا اللّه عنه فضلاً من
عنده انّه لهو الغفور الکریم قد اذن اللّه لکم السّجود علی کل شئً طاهرٍ و رفعنا
عنه حکم الحدّ فی الکتاب انّ اللّه یعلم و انتم لا تعلمون من لم یجد المآء
یذکر خمس مرّاتٍ بسم اللّه الأطهر الأطهر ثمّ یشرع فی العمل هذا ما حکم به مولی
العالمین و البلدان الّتی طالت فیها اللّیالی و الایام فلیصلّوا بالسّاعات و
المشاخص الّتی منها تحدّدت الأوقات انّه لهو المبین الحکیم

\end{otherlanguage}

Qad furiḍa ʿalaykumu al-ṣalātu wa-al-ṣawmu min awwali al-bulūgh amran
min ladā Allāhi rabbikum wa-rabbi ābāʾikumu al-awwalīn. Man kāna fī
nafsihi ḍaʿfun mina al-maraḍi aw al-haram ʿafā Allāhu ʿanhu faḍlan min
ʿindihi innahu la-huwa al-Ghafūru al-Karīm. Qad adhina Allāhu lakumu
al-sujūda ʿalā kulli shayʾin ṭāhirin wa-rafaʿnā ʿanhu ḥukma al-ḥaddi fī
al-kitāb. Inna Allāha yaʿlamu wa-antum lā taʿlamūn. Man lam yajidi
al-māʾa yadhkuru khamsa marrātin bismi Allāhi al-Aṭhari al-Aṭhar thumma
yashraʿu fī al-ʿamal. Hādhā mā ḥakama bihi Mawlā al-ʿālamīn.
Wa-al-buldānu allatī ṭālat fīhā al-layālī wa-al-ayyām fa-li-yuṣallū
bi-al-sāʿāti wa-al-mashākhiṣi allatī minhā taḥaddadati al-awqāt. Innahu
la-huwa al-Mubīnu al-Ḥakīm.

Prayer and fasting have been made obligatory for you from the onset of
maturity, as a command from God, your Lord and the Lord of your
forefathers. Whoever is weakened by illness or old age has been exempted
by God as a grace from Him. He is indeed the All-Forgiving, the
Generous. God has permitted you to prostrate yourselves on any clean
surface, and We have removed from the Book the restriction that applied
to it. God knows, and you do not know. Whoever cannot find water is to
repeat five times, ``In the name of God, the Most Pure, the Most Pure,''
and then proceed with the act. This is what the Lord of the worlds has
decreed. In lands where nights and days are long, they are to pray
according to clocks and instruments by which the times are determined.
He is indeed the Expounder, the All-Wise.

\subsection{12}\label{section-11}

\begin{otherlanguage}{arabic}

قد عفونا عنکم صلوة الآیات اذا ظهرت اذکروا اللّه بالعظمة و الأقتدار انّه هو
السّمیع البصیر قولوا العظمة للّه ربّ ما یری و ما لا یری ربّ العالمین

\end{otherlanguage}

Qad ʿafawnā ʿankum ṣalāta al-āyāt. Idhā ẓaharat udhkurū Allāha
bi-al-ʿaẓamati wa-al-iqtidār. Innahu huwa al-Samīʿu al-Baṣīr. Qūlū
al-ʿaẓamatu li-Allāhi rabbi mā yurā wa-mā lā yurā rabbi al-ʿālamīn.

We have exempted you from the prayer of the signs. When they appear,
remember God with grandeur and power. He is the All-Hearing, the
All-Seeing. Say: ``Grandeur belongs to God, Lord of all that is seen and
unseen, Lord of the worlds.''

\subsection{13}\label{section-12}

\begin{otherlanguage}{arabic}

کتب علیکم الصّلوة فرادی قد رفع حکم الجماعة الاّ فی صلوة المیت انّه لهو
الأمر الحکیم

\end{otherlanguage}

Kutiba ʿalaykumu al-ṣalātu furādā qad rufiʿa ḥukmu al-jamāʿati illā fī
ṣalāti al-mayyit. Innahu la-huwa al-Āmiru al-Ḥakīm.

You have been required to perform obligatory prayer individually. The
ordinance of congregational prayer has been abolished, except for the
prayer for the dead. He is indeed the Ordainer, the All-Wise.

\subsection{14}\label{section-13}

\begin{otherlanguage}{arabic}

قد عفا اللّه عن النّسآء حین ما یجدن الدّم الصّوم و الصّلوة و لهنّ ان یتوضّأن و
یسبّحن خمساً و تسعین مرّة من زوالٍ الی زوالٍ سبحان اللّه ذی الطّلعة و الجمال
هذا ما قدّر فی الکتاب ان انتم من العالمین

\end{otherlanguage}

Qad ʿafā Allāhu ʿani al-nisāʾi ḥīna mā yajidna al-dama al-ṣawma
wa-al-ṣalāta wa-lahunna an yatawaḍḍaʾna wa-yusabbiḥna khamsan wa-tisʿīna
marratan min zawālin ilā zawālin subḥāna Allāhi dhī al-ṭalʿati
wa-al-jamāl. Hādhā mā quddira fī al-kitābi in antum mina al-ʿālimīn.

God has exempted women from fasting and obligatory prayer during
menstruation. They are to perform ablutions and repeat ninety-five
times, from one noon to the next, ``Glory be to God, the Possessor of
radiance and beauty.'' This is what has been ordained in the Book, if
you are among those who know.

\subsection{15}\label{section-14}

\begin{otherlanguage}{persian}

و لکم و لهنّ فی الاسفار اذ انزلتم و استرحتم المقام الأمن مکان کلّ صلوة
سجدة واحدة و اذکروا فیها سبحان اللّه ذی العظمة و الاجلال و الموهبة و
الافضال و الّذی عجز یقول سبحان اللّه انّه یکفیه بالحق انّه لهو الکافی الباقی
الغفور الرّحیم و بعد اتمام السّجود لکم و لهنّ ان تقعدوا علی هیکل التّوحید و
تقولوا ثمانی عشرة مرّة سبحان اللّه ذی الملک و الملکوت کذلک یبین اللّه سبل
الحقّ و الهدی و انّها انتهت الی سبیل واحد و هو هذا الصّراط المستقیم اشکروا
اللّه بهذا الفضل العظیم احمدوا اللّه بهذه الموهبة الّتی احاطت السّموات و
الأرضین اذکروا اللّه بهذه الرّحمة الّتی سبقت العالمین

\end{otherlanguage}

Wa-lakum wa-lahunna fī al-asfāri idhā nazaltum wa-istaraḥtumu al-maqāma
al-amna makāna kulli ṣalātin sajdatun wāḥidatun wa-udhkurū fīhā subḥāna
Allāhi dhī al-ʿaẓamati wa-al-ijlāli wa-al-mawhibati wa-al-ifḍāl.
Wa-alladhī ʿajaza yaqūlu subḥāna Allāh innahu yakfīhi bi-al-ḥaqq innahu
la-huwa al-Kāfī al-Bāqī al-Ghafūru al-Raḥīm. Wa-baʿda itmāmi al-sujūdi
lakum wa-lahunna an taqʿudū ʿalā haykali al-tawḥīdi wa-taqūlū thamāniya
ʿashrata marratan subḥāna Allāhi dhī al-mulki wa-al-malakūt. Kadhālika
yubayyinu Allāhu subula al-ḥaqqi wa-al-hudā wa-innahā intahat ilā
sabīlin wāḥidin wa-huwa hādhā al-ṣirāṭu al-mustaqīm. Ushkurū Allāha
bi-hādhā al-faḍli al-ʿaẓīm. Iḥmadū Allāha bi-hādhihi al-mawhibati allatī
aḥāṭati al-samāwāti wa-al-araḍīn. Udhkurū Allāha bi-hādhihi al-raḥmati
allatī sabaqati al-ʿālamīn.

For you and for women, when traveling, upon stopping and resting in a
safe place, a single prostration is to replace each obligatory prayer.
During it say, ``Glory be to God, the Possessor of grandeur and majesty,
bounty and grace.'' Anyone unable to do this says, ``Glory be to God'';
this will indeed suffice. He is indeed the All-Sufficing, the
Everlasting, the All-Forgiving, the Merciful. After completing the
prostrations, you and women are to sit in the posture of unity and say
eighteen times, ``Glory be to God, the Possessor of the earthly and
heavenly kingdoms.'' Thus God makes clear the ways of truth and
guidance, and they have converged in a single way: this straight path.
Thank God for this great grace. Praise God for this bounty that has
encompassed the heavens and the earth. Remember God for this mercy that
has preceded all the worlds.

\subsection{16}\label{section-15}

\begin{otherlanguage}{arabic}

قل قد جعل اللّه مفتاح الکنز حبّی المکنون لو انتم تعرفون لولا المفتاح لکان
مکنوناً فی ازل الازال لو انتم توقنون قل هذا لمطلع الوحی و مشرق الأشراق
الّذی به اشرقت الآفاق لو انتم تعلمون انّ هذا لهو القضآء المثبت و به ثبت کلّ
قضآءٍ محتوم

\end{otherlanguage}

Qul qad jaʿala Allāhu miftāḥa al-kanzi ḥubbī al-maknūna law antum
taʿrifūn. Lawlā al-miftāḥu la-kāna maknūnan fī azali al-āzāl law antum
tūqinūn. Qul hādhā la-maṭlaʿu al-waḥyi wa-mashriqu al-ishrāqi alladhī
bihi ashraqati al-āfāqu law antum taʿlamūn. Inna hādhā la-huwa al-qaḍāʾu
al-muthbatu wa-bihi thabata kullu qaḍāʾin maḥtūm.

Say: God has made My hidden love the key to the treasure, if only you
would recognize it. Without the key, it would have remained concealed
for all eternity, if only you were certain. Say: This is indeed the
Dawning Place of revelation and the Dayspring of radiance through which
the horizons have been illumined, if only you knew. This is indeed the
established decree, and through it every irrevocable decree has been
established.

\subsection{17}\label{section-16}

\begin{otherlanguage}{arabic}

یا قلم الأعلی قل یا ملاء الأنشآء قد کتبنا علیکم الصّیام ایاماً معدودات و
جعلنا النّیروز عیداً لکم بعد اکمالها کذلک اضآئت شمس البیان من افق الکتاب
من لدن مالک المبدء و المأب و اجعلوا الأیام الزّآئدة عن الشّهور قبل شهر
الصّیام انّا جعلناها مظاهر الهآء بین اللّیالی و الأیّام لذا ما تحدّدت بحدود
السّنة و الشّهور ینبغی لأهل البهآء ان یطعموا فیها انفسهم و ذوی القربی ثمّ
الفقرآء و المساکین و یهلّلنّ و یکبّرنّ و یسبّحنّ و یمجّدنّ ربّهم بالفرح و
الأنبساط و اذا تمّت ایام الاعطآء قبل الامساک فلیدخلنّ فی الصّیام کذلک حکم
مولی الأنام لیس علی المسافر و المریض و الحامل و المرضع من حرجٍ عفا اللّه
عنهم فضلاً من عنده انّه لهو العزیز الوهّاب

\end{otherlanguage}

Yā qalama al-aʿlā qul yā malaʾa al-inshāʾ qad katabnā ʿalaykumu
al-ṣiyāma ayyāman maʿdūdātin wa-jaʿalnā al-nayrūza ʿīdan lakum baʿda
ikmālihā. Kadhālika aḍāʾat shamsu al-bayāni min ufuqi al-kitābi min
ladun Māliki al-mabdaʾi wa-al-maʾāb. Wa-ijʿalū al-ayyāma al-zāʾidata
ʿani al-shuhūri qabla shahri al-ṣiyām. Innā jaʿalnāhā maẓāhira al-hāʾi
bayna al-layālī wa-al-ayyām li-dhā mā taḥaddadat bi-ḥudūdi al-sanati
wa-al-shuhūr. Yanbaghī li-ahli al-Bahāʾi an yuṭʿimū fīhā anfusahum
wa-dhawī al-qurbā thumma al-fuqarāʾa wa-al-masākīna wa-yuhallilanna
wa-yukabbiranna wa-yusabbiḥanna wa-yumajjidanna rabbahum bi-al-faraḥi
wa-al-inbisāṭ. Wa-idhā tammat ayyāmu al-iʿṭāʾi qabla al-imsāk
fa-li-yadkhulanna fī al-ṣiyām. Kadhālika ḥakama Mawlā al-anām. Laysa
ʿalā al-musāfiri wa-al-marīḍi wa-al-ḥāmili wa-al-murḍiʿi min ḥaraj. ʿAfā
Allāhu ʿanhum faḍlan min ʿindihi innahu la-huwa al-ʿAzīzu al-Wahhāb.

Most Exalted Pen, say: People of creation, We have prescribed fasting
for you for a specified number of days and have made Naw-Rúz a festival
for you upon their completion. Thus has the Sun of utterance shone from
the horizon of the Book, from the Lord of the beginning and the return.
Place the days in excess of the months before the month of fasting. We
have made them manifestations of the letter Há among the nights and
days; thus they have not been bounded by the limits of the year and the
months. During them, the people of Bahá should provide food for
themselves and their relatives, as well as the poor and needy, and with
joy and gladness proclaim the oneness and greatness of their Lord,
praise Him, and glorify Him. When the days of giving that precede
abstinence have ended, let them enter the fast. Thus has the Lord of
humanity decreed. Travelers, the sick, pregnant women, and nursing
mothers are under no obligation. God has exempted them as a grace from
Him. He is indeed the Almighty, the Bestower.

\subsection{18}\label{section-17}

\begin{otherlanguage}{arabic}

هذه حدود اللّه الّتی رقمت من القلم الأعلی فی الزّبر و الألواح تمسّکوا باوامر
اللّه و احکامه و لا تکونوا من الّذین اخذوا اصول انفسهم و نبذوا اصول اللّه
ورآئهم بما اتّبعوا الظّنون و الأوهام کفّوا انفسکم عن الأکل و الشّرب من
الطّلوع الی الأفول ایّاکم ان یمنعکم الهوی عن هذا الفضل الّذی قدّر فی الکتاب

\end{otherlanguage}

Hādhihi ḥudūdu Allāhi allatī ruqimat mina al-qalami al-aʿlā fī al-zuburi
wa-al-alwāḥ. Tamassakū bi-awāmiri Allāhi wa-aḥkāmihi wa-lā takūnū mina
alladhīna akhadhū uṣūla anfusihim wa-nabadhū uṣūla Allāhi warāʾahum
bi-mā ittabaʿū al-ẓunūna wa-al-awhām. Kuffū anfusakum ʿani al-akli
wa-al-shurbi mina al-ṭulūʿi ilā al-ufūl. Iyyākum an yamnaʿakumu al-hawā
ʿan hādhā al-faḍli alladhī quddira fī al-kitāb.

These are the laws of God inscribed by the Most Exalted Pen in the
scriptures and Tablets. Hold fast to God's commandments and ordinances,
and do not be among those who have adopted principles of their own and
cast God's principles behind them, following conjectures and vain
imaginings. Refrain from eating and drinking from sunrise to sunset.
Beware lest desire deprive you of this grace ordained in the Book.

\subsection{19}\label{section-18}

\begin{otherlanguage}{persian}

قد کتب لمن دان باللّه الدّیان ان یغسل فی کلّ یومٍ یدیه ثمّ وجهه و یقعد مقبلاً
الی اللّه و یذکر خمساً و تسعین مرّةً اللّه ابهی کذلک حکم فاطر السّمآء اذ استوی
علی اعراش الأسماء بالعظمة و الأقتدار کذلک توضّأوا للصّلوة امراً من اللّه
الواحد المختار

\end{otherlanguage}

Qad kutiba li-man dāna bi-Allāhi al-Dayyāni an yaghsila fī kulli yawmin
yadayhi thumma wajhahu wa-yaqʿuda muqbilan ilā Allāhi wa-yadhkura
khamsan wa-tisʿīna marratan Allāhu Abhā. Kadhālika ḥakama Fāṭiru
al-samāʾi idh istawā ʿalā aʿrāshi al-asmāʾi bi-al-ʿaẓamati
wa-al-iqtidār. Kadhālika tawaḍḍaʾū li-al-ṣalāti amran min Allāhi
al-Wāḥidi al-Mukhtār.

Everyone who believes in God, the Judge, has been required each day to
wash his hands and then his face, sit facing God, and repeat ninety-five
times, ``God is Most Glorious.'' Thus has the Creator of heaven decreed
as He sat upon the thrones of names in grandeur and power. Likewise,
perform ablutions for obligatory prayer, as a command from God, the One,
who freely chooses.

\subsection{20}\label{section-19}

\begin{otherlanguage}{arabic}

قد حرّم علیکم القتل و الزّنا ثمّ الغیبة و الافترآء اجتنبوا عمّا نهیتم عنه فی
الصّحآئف و الألواح

\end{otherlanguage}

Qad ḥurrima ʿalaykumu al-qatlu wa-al-zinā thumma al-ghībatu
wa-al-iftirāʾ. Ijtanibū ʿammā nuhītum ʿanhu fī al-ṣaḥāʾifi wa-al-alwāḥ.

Murder, adultery, backbiting, and slander have been forbidden to you.
Avoid whatever you have been forbidden in the scriptures and Tablets.

\subsection{21}\label{section-20}

\begin{otherlanguage}{arabic}

قد قسمنا المواریث علی عدد الزّآء منها قدّر لذرّیاتکم من کتاب الطّآء علی عدد
المقت و للأزواج من کتاب الحآء علی عدد التّآء و الفآء و للابآء من کتاب
الزّآء علی عدد التّآء و الکاف و للأمّهات من کتاب الواو علی عدد الرّفیع و
للأخوان من کتاب الهآء عدد الشّین و للأخوات من کتاب الدّال عدد الرّآء و
المیم و للمعلّمین من کتاب الجیم عدد القاف و الفآء کذلک حکم مبشّری الّذی
یذکرنی فی اللّیالی و الأسحار انّا لمّا سمعنا ضجیج الذّرّیات فی الأصلاب زدنا
ضعف مالهم و نقصنا عن الاخری انّه لهو المقتدر علی ما یشآء یفعل بسلطانه کیف
اراد

\end{otherlanguage}

Qad qasamnā al-mawārītha ʿalā ʿadadi al-zāʾ. Minhā quddira
li-dhurriyyātikum min kitābi al-ṭāʾ ʿalā ʿadadi al-maqt, wa-li-l-azwāji
min kitābi al-ḥāʾ ʿalā ʿadadi al-tāʾ wa-l-fāʾ, wa-li-l-ābāʾi min kitābi
al-zāʾ ʿalā ʿadadi al-tāʾ wa-l-kāf, wa-li-l-ummahāti min kitābi al-wāw
ʿalā ʿadadi al-rafīʿ, wa-li-l-ikhwāni min kitābi al-hāʾ ʿadada al-shīn,
wa-li-l-akhawāti min kitābi al-dāl ʿadada al-rāʾ wa-l-mīm,
wa-li-l-muʿallimīna min kitābi al-jīm ʿadada al-qāf wa-l-fāʾ. Kadhālika
ḥakama mubashshirī alladhī yadhkurunī fī al-layālī wa-l-asḥār. Innā
lammā samiʿnā ḍajīja al-dhurriyyāti fī al-aṣlāb zidnā ḍiʿfa mā lahum
wa-naqaṣnā ʿani al-ukhrā. Innahu la-huwa al-muqtadiru ʿalā mā yashāʾ,
yafʿalu bi-sulṭānihi kayfa arād.

We have divided inheritances according to the number of Zāʾ. Of these,
your offspring have been allotted from the Book of Ṭāʾ according to the
number of maqt; spouses, from the Book of Ḥāʾ according to the number of
Tāʾ and Fāʾ; fathers, from the Book of Zāʾ according to the number of
Tāʾ and Kāf; mothers, from the Book of Wāw according to the number of
rafīʿ; brothers, from the Book of Hāʾ, the number of Shīn; sisters, from
the Book of Dāl, the number of Rāʾ and Mīm; and teachers, from the Book
of Jīm, the number of Qāf and Fāʾ. Thus decreed My Herald, who remembers
Me in the night and at dawn. When We heard the clamor of the offspring
still in their fathers' loins, We doubled their share and reduced those
of the others. He indeed has power over whatever He wills; through His
sovereignty He does as He wishes.

\subsection{22}\label{section-21}

\begin{otherlanguage}{arabic}

من مات و لم یکن له ذرّیة ترجع حقوقهم الی بیت العدل لیصرفوها امنآء الرّحمن
فی الأیتام و الارامل و ما ینتفع به جمهور النّاس لیشکروا ربّهم العزیز
الغفّار

\end{otherlanguage}

Man māta wa-lam yakun lahu dhurriyyatun tarjiʿu ḥuqūquhum ilā bayti
al-ʿadl li-yaṣrifūhā umanāʾu al-Raḥmān fī al-aytāmi wa-l-arāmili wa-mā
yantafiʿu bihi jumhūru al-nās li-yashkurū rabbahumu al-ʿazīza
al-ghaffār.

If someone dies without offspring, their shares revert to the House of
Justice, so that the trustees of the All-Merciful may spend them on
orphans, widows, and whatever benefits the general public, that they may
give thanks to their Lord, the Almighty, the Ever-Forgiving.

\subsection{23}\label{section-22}

\begin{otherlanguage}{arabic}

و الّذی له ذرّیة و لم یکن مادونها عمّا حدّد فی الکتاب یرجع الثّلثان ممّا ترکه
الی الذّرّیة و الثّلث الی بیت العدل کذلک حکم الغنی المتّعال بالعظمة و
الأجلال

\end{otherlanguage}

Wa-alladhī lahu dhurriyyatun wa-lam yakun mā dūnahā ʿammā ḥuddida fī
al-kitāb yarjiʿu al-thuluthāni mimmā tarakahu ilā al-dhurriyyati
wa-l-thuluthu ilā bayti al-ʿadl. Kadhālika ḥakama al-ghaniyyu al-mutaʿāl
bi-l-ʿaẓamati wa-l-ijlāl.

If someone has offspring but none of the other heirs specified in the
Book, two-thirds of what he leaves revert to the offspring and one-third
to the House of Justice. Thus has the Self-Sufficient, the Most Exalted,
decreed in majesty and glory.

\subsection{24}\label{section-23}

\begin{otherlanguage}{arabic}

و الّذی لم یکن له من یرثه و کان له ذو القربٰی من ابنآء الأخ و الأخت و
بناتهما فلهم الثّلثان والاّ للأعمام و الأخوال و العمّات و الخالات و من
بعدهم و بعدهنّ لابنآئهم و ابنآئهنّ و بناتهم و بناتهنّ و الثلث یرجع الی مقرّ
العدل امراً فی الکتاب من لدی اللّه مالک الرّقاب

\end{otherlanguage}

Wa-alladhī lam yakun lahu man yarithuhu wa-kāna lahu dhawū al-qurbā min
abnāʾi al-akhi wa-l-ukhti wa-banātihimā fa-lahumu al-thuluthān, wa-illā
fa-li-l-aʿmāmi wa-l-akhwāli wa-l-ʿammāti wa-l-khālāt, wa-min baʿdihim
wa-baʿdihinna li-abnāʾihim wa-abnāʾihinna wa-banātihim wa-banātihinna,
wa-l-thuluthu yarjiʿu ilā maqarri al-ʿadl, amran fī al-kitābi min ladā
Allāhi māliki al-riqāb.

If someone has none to inherit from him but has relatives among the sons
and daughters of his brother or sister, two-thirds go to them;
otherwise, to his paternal and maternal uncles and aunts, and after them
to their sons and daughters. One-third reverts to the Seat of Justice,
as an ordinance in the Book from God, the Master of all people.

\subsection{25}\label{section-24}

\begin{otherlanguage}{arabic}

من مات و لم یکن له احد من الّذین نزّلت اسمائهم من القلم الأعلی ترجع
الأموال کلّها الی المقرّ المذکور لتصرف فیما امر اللّه به انّه لهو المقتدر
الأمّار

\end{otherlanguage}

Man māta wa-lam yakun lahu aḥadun mina alladhīna nuzzilat asmāʾuhum mina
al-qalami al-aʿlā tarjiʿu al-amwālu kulluhā ilā al-maqarri al-madhkūr
li-tuṣrafa fīmā amara Allāhu bihi. Innahu la-huwa al-muqtadiru al-ammār.

If someone dies leaving none of those whose names have been sent down by
the Most Exalted Pen, all the property reverts to the aforementioned
Seat, to be spent on what God has ordained. He indeed is the
All-Powerful, the Supreme Ordainer.

\subsection{26}\label{section-25}

\begin{otherlanguage}{arabic}

و جعلنا الدّار المسکونة و الالبسة المخصوصه للذّرّیة من الذّکران دون الأناث و
الورّاث انّه لهو المعطی الفیاض

\end{otherlanguage}

Wa-jaʿalnā al-dāra al-maskūnata wa-l-albisata al-makhṣūṣata
li-l-dhurriyyati mina al-dhukrāni dūna al-ināthi wa-l-wurrāth. Innahu
la-huwa al-muʿṭī al-fayyāḍ.

We have assigned the residence and personal clothing to the male
offspring, to the exclusion of the female offspring and other heirs. He
indeed is the Giver, the Most Bountiful.

\subsection{27}\label{section-26}

\begin{otherlanguage}{arabic}

انّ الّذی مات فی ایام والده و له ذرّیة اولٓئک یرثون ما لأبیهم فی کتاب اللّه
اقسموا بینهم بالعدل الخالص کذلک ماج بحر الکلام و قذف لئالی الأحکام من
لدن مالک الأنام

\end{otherlanguage}

Inna alladhī māta fī ayyāmi wālidihi wa-lahu dhurriyyatun ulāʾika
yarithūna mā li-abīhim fī kitābi Allāh. Iqsimū baynahum bi-l-ʿadli
al-khāliṣ. Kadhālika māja baḥru al-kalāmi wa-qadhafa laʾāliʾa al-aḥkāmi
min ladun māliki al-anām.

If someone dies during his father's lifetime and leaves offspring, they
inherit their father's share as prescribed in the Book of God. Divide it
among them with pure justice. Thus has the ocean of utterance surged and
cast forth the pearls of ordinances from the Lord of all people.

\subsection{28}\label{section-27}

\begin{otherlanguage}{persian}

و الّذی ترک ذرّیةً ضعافاً سلّموا مالهم الی امین لیتجرّ لهم الی ان یبلغوا رشدهم
او الی محلّ الشّراکة ثمّ عینوا للأمین حقّاً ممّا حصل من التّجارة و الأقتراف

\end{otherlanguage}

Wa-alladhī taraka dhurriyyatan ḍiʿāfan sallimū mālahum ilā amīnin
li-yattajira lahum ilā an yablughū rushdahum aw ilā maḥalli al-shirāka.
Thumma ʿayyinū li-l-amīni ḥaqqan mimmā ḥaṣala mina al-tijārati
wa-l-iqtirāf.

If someone leaves offspring who are minors, entrust their property to a
trustworthy person to trade on their behalf until they reach maturity or
the age of partnership. Then assign the trustee a share of the proceeds
from trade and earnings.

\subsection{29}\label{section-28}

\begin{otherlanguage}{arabic}

کلّ ذلک بعد ادآء حقّ اللّه و الدّیون لو تکون علیه و تجهیز الأسباب للکفن و
الدّفن و حمل المیت بالعزّة و الاءعتزاز کذلک حکم مالک المبدء و المأب

\end{otherlanguage}

Kullu dhālika baʿda adāʾi ḥaqqi Allāhi wa-l-duyūni law takūnu ʿalayhi
wa-tajhīzi al-asbābi li-l-kafani wa-l-dafni wa-ḥamli al-mayyiti
bi-l-ʿizzati wa-l-iʿtizāz. Kadhālika ḥakama māliku al-mabdaʾi
wa-l-maʾāb.

All this follows payment of the Right of God and any debts owed by the
deceased, provision for shrouding and burial, and the conveyance of the
deceased with dignity and honor. Thus has the Lord of the beginning and
the return decreed.

\subsection{30}\label{section-29}

\begin{otherlanguage}{arabic}

قل هذا لهو العلم المکنون الّذی لن یتغیر لانّه بدء بالطّآء المدلّة علی الأسم
المخزون الظّاهر الممتنع المنیع و ما خصّصناه للذّرّیات هذا من فضل اللّه علیهم
لیشکروا ربّهم الرّحمن الرّحیم تلک حدود اللّه لا تعتدوها باهوآء انفسکم اتّبعوا
ما امرتم به من مطلع البیان و المخلصون یرون حدود اللّه ماء الحیوان لاهل
الأدیان و مصباح الحکمة و الفلاح لمن فی الأرضین و السّموات

\end{otherlanguage}

Qul hādhā la-huwa al-ʿilmu al-maknūnu alladhī lan yataghayyara li-annahu
budiʾa bi-l-ṭāʾi al-mudillati ʿalā al-ismi al-makhzūni al-ẓāhiri
al-mumtaniʿi al-manīʿ. Wa-mā khaṣṣaṣnāhu li-l-dhurriyyāti hādhā min
faḍli Allāhi ʿalayhim li-yashkurū rabbahumu al-Raḥmāna al-raḥīm. Tilka
ḥudūdu Allāhi lā taʿtadūhā bi-ahwāʾi anfusikum. Ittabiʿū mā umirtum bihi
min maṭlaʿi al-bayān. Wa-l-mukhliṣūna yarawna ḥudūda Allāhi māʾa
al-ḥayawāni li-ahli al-adyāni wa-miṣbāḥa al-ḥikmati wa-l-falāḥi li-man
fī al-arḍīna wa-l-samāwāt.

Say: This is indeed the hidden knowledge that will never change, for it
began with Ṭāʾ, which points to the treasured Name, the Manifest, the
Inaccessible, the Unassailable. What We have set aside for the offspring
is a grace bestowed on them by God, that they may thank their Lord, the
All-Merciful, the Compassionate. These are God's limits; do not
transgress them at the prompting of your own desires. Follow what you
have been commanded by the Dawning-Place of utterance. The sincere see
God's ordinances as the water of life for the followers of religions and
the lamp of wisdom and salvation for all in the earth and the heavens.

\subsection{31}\label{section-30}

\begin{otherlanguage}{arabic}

قد کتب اللّه علی کلّ مدینةٍ ان یجعلوا فیها بیت العدل و یجتمع فیه النّفوس علی
عدد البهآء و ان ازداد لا بأس و یرون کانّهم یدخلون محضر اللّه العلی الأعلی
و یرون من لا یری و ینبغی لهم ان یکونوا امنآء الرّحمن بین الاءمکان و وکلاء
اللّه لمن علی الأرض کلّها و یشاوروا فی مصالح العباد لوجه اللّه کما یشاورون
فی امورهم و یختاروا ما هو المختار کذلک حکم ربّکم العزیز الغفّار ایاکم ان
تدعوا ما هو المنصوص فی اللّوح اتّقوا اللّه یا اولی الانظار

\end{otherlanguage}

Qad kataba Allāhu ʿalā kulli madīnatin an yajʿalū fīhā bayta al-ʿadli
wa-yajtamiʿa fīhi al-nufūsu ʿalā ʿadadi al-Bahāʾ, wa-in izdāda lā baʾs.
Wa-yarawna ka-annahum yadkhulūna maḥḍara Allāhi al-ʿaliyyi al-aʿlā
wa-yarawna man lā yurā. Wa-yanbaghī lahum an yakūnū umanāʾa al-Raḥmāni
bayna al-imkāni wa-wukalāʾa Allāhi li-man ʿalā al-arḍi kullihā
wa-yushāwirū fī maṣāliḥi al-ʿibādi li-wajhi Allāhi kamā yushāwirūna fī
umūrihim wa-yakhtārū mā huwa al-mukhtār. Kadhālika ḥakama rabbukumu
al-ʿazīzu al-ghaffār. Iyyākum an tadaʿū mā huwa al-manṣūṣu fī al-lawḥ.
Ittaqū Allāha yā ulī al-anẓār.

God has prescribed that a House of Justice be established in every city,
in which souls numbering the numerical value of Bahāʾ shall gather; if
the number is greater, there is no objection. They should regard
themselves as entering the presence of God, the Most High, the Most
Exalted, and as seeing Him who cannot be seen. They should be the
trustees of the All-Merciful among created beings and God's
representatives for all who dwell on earth. They should consult, for
God's sake, on the interests of His servants as they consult on their
own affairs, and choose what is best. Thus has your Lord, the Almighty,
the Ever-Forgiving, decreed. Beware of abandoning what is expressly
stated in the Tablet. Be mindful of God, you who possess insight.

\subsection{32}\label{section-31}

\begin{otherlanguage}{arabic}

یا ملاء الأنشآء عمّروا بیوتاً باکمل ما یمکن فی الامکان باسم مالک الأدیان
فی البلدان و زینوها بما ینبغی لها لا بالصّور و الامثال ثمّ اذکروا فیها
ربّکم الرّحمن بالرّوح و الرّیحان الا بذکره تستنیر الصّدور و تقرّ الأبصار

\end{otherlanguage}

Yā malaʾa al-inshāʾ, ʿammirū buyūtan bi-akmali mā yumkinu fī al-imkāni
bi-smi māliki al-adyāni fī al-buldāni wa-zayyinūhā bi-mā yanbaghī lahā
lā bi-l-ṣuwari wa-l-amthāl. Thumma udhkurū fīhā rabbakumu al-Raḥmāna
bi-l-rūḥi wa-l-rayḥān. Alā bi-dhikrihi tastanīru al-ṣudūru wa-taqarru
al-abṣār.

People of creation! Build houses throughout the lands in the name of the
Lord of religions, as perfectly as is possible in the world of being,
and adorn them with what befits them, not with images and statues. Then
remember within them your Lord, the All-Merciful, with joy and gladness.
Truly, through remembrance of Him hearts are illumined and eyes are
cheered.

\subsection{33}\label{section-32}

\begin{otherlanguage}{persian}

قد حکم اللّه لمن استطاع منکم حجّ البیت دون النّسآء عفا اللّه عنهنّ رحمة من
عنده انّه لهو المعطی الوهّاب

\end{otherlanguage}

Qad ḥakama Allāhu li-man istaṭāʿa minkum ḥajja al-bayti dūna al-nisāʾ.
ʿAfā Allāhu ʿanhunna raḥmatan min ʿindihi. Innahu la-huwa al-muʿṭī
al-wahhāb.

God has ordained pilgrimage to the House for those of you who are able,
except women. God has exempted them as a mercy from Him. He indeed is
the Giver, the All-Bountiful.

\subsection{34}\label{section-33}

\begin{otherlanguage}{arabic}

یا اهل البهآء قد وجب علی کلّ واحدٍ منکم الاشتغال بامر من الامور من الصّنآیع
و الاقتراف و امثالها و جعلنا اشتغالکم بها نفس العبادة للّه الحق تفکّروا یا
قوم فی رحمة اللّه و الطافه ثمّ اشکروه فی العشی و الأشراق لا تضیعوا اوقاتکم
بالبطالة و الکسالة و اشتغلوا بما ینتفع به انفسکم و انفس غیرکم کذلک قضی
الأمر فی هذا اللّوح الّذی لاحت من افقه شمس الحکمة و التّبیان ابغض النّاس عند
اللّه من یقعد و یطلب تمسّکوا بحبل الأسباب متوکّلین علی اللّه مسبّب الأسباب

\end{otherlanguage}

Yā ahla al-Bahāʾ, qad wajaba ʿalā kulli wāḥidin minkumu al-ishtighālu
bi-amrin mina al-umūri mina al-ṣanāʾiʿi wa-l-iqtirāfi wa-amthālihā.
Wa-jaʿalnā ishtighālakum bihā nafsa al-ʿibādati li-llāhi al-ḥaqq.
Tafakkarū yā qawmi fī raḥmati Allāhi wa-alṭāfihi thumma ushkurūhu fī
al-ʿashiyyi wa-l-ishrāq. Lā tuḍayyiʿū awqātakum bi-l-baṭālati
wa-l-kasālati wa-ishtaghilū bi-mā yantafiʿu bihi anfusukum wa-anfusu
ghayrikum. Kadhālika quḍiya al-amru fī hādhā al-lawḥi alladhī lāḥat min
ufuqihi shamsu al-ḥikmati wa-l-tibyān. Abghaḍu al-nāsi ʿinda Allāhi man
yaqʿudu wa-yaṭlub. Tamassakū bi-ḥabli al-asbābi mutawakkilīna ʿalā
Allāhi musabbibi al-asbāb.

People of Bahāʾ! Every one of you is required to engage in an
occupation, such as a craft, a trade, or the like. We have made your
engagement in such work the very worship of God, the True One. Reflect,
My people, on God's mercy and favors, then thank Him in the evening and
at daybreak. Do not waste your time in idleness and laziness, but occupy
yourselves with what benefits you and others. Thus has the matter been
decreed in this Tablet, from whose horizon the sun of wisdom and
exposition has shone. The most detestable of people in God's sight is
one who sits and begs. Hold fast to the cord of means, placing your
trust in God, the Originator of all means.

\subsection{35}\label{section-34}

\begin{otherlanguage}{arabic}

قد حرّم علیکم تقبیل الأیادی فی الکتاب هذا ما نهیتم عنه من لدن ربّکم العزیز
الحکّام لیس لأحدٍ ان یستغفر عند احد توبُوا الی اللّه تلقآء انفسکم انّه لهو
الغافر المعطی العزیز التّوّاب

\end{otherlanguage}

Qad ḥurrima ʿalaykum taqbīlu al-ayādī fī al-kitāb. Hādhā mā nuhītum
ʿanhu min ladun rabbikumu al-ʿazīzi al-ḥakkām. Laysa li-aḥadin an
yastaghfira ʿinda aḥad. Tūbū ilā Allāhi tilqāʾa anfusikum. Innahu
la-huwa al-ghāfiru al-muʿṭī al-ʿazīzu al-tawwāb.

Kissing hands has been forbidden to you in the Book. This is what your
Lord, the Almighty, the Supreme Judge, has prohibited you from doing. No
one may seek forgiveness in the presence of another. Turn to God in
repentance on your own. He indeed is the Forgiver, the Giver, the
Almighty, the Ever-Relenting.

\subsection{36}\label{section-35}

\begin{otherlanguage}{arabic}

یا عباد الرّحمن قوموا علی خدمة الأمر علی شأنٍ لا تأخذکم الأحزان من الّذین
کفروا بمطلع الآیات لمّا جآء الوعد و ظهر الموعود اختلف النّاس و تمسّک کلّ حزبٍ
بما عنده من الظّنون و الأوهام

\end{otherlanguage}

Yā ʿibāda al-Raḥmān, qūmū ʿalā khidmati al-amri ʿalā shaʾnin lā
taʾkhudhukumu al-aḥzānu mina alladhīna kafarū bi-maṭlaʿi al-āyāt. Lammā
jāʾa al-waʿdu wa-ẓahara al-mawʿūdu ikhtalafa al-nāsu wa-tamassaka kullu
ḥizbin bi-mā ʿindahu mina al-ẓunūni wa-l-awhām.

Servants of the All-Merciful! Arise to serve the Cause in such a way
that those who disbelieve in the Dawning-Place of the verses cannot
afflict you with sorrow. When the promise was fulfilled and the Promised
One appeared, the people fell into disagreement, and every party clung
to its own conjectures and imaginings.

\subsection{37}\label{section-36}

\begin{otherlanguage}{arabic}

من النّاس من یقعد صفّ النّعال طلباً لصدر الجلال قل من انت یا ایها الغافل
الغرّار و منهم من یدّعی الباطن و باطن الباطن قل یا ایها الکذّاب تاللّه ما
عندک انّه من القشور ترکناها لکم کما تترک العظام للکلاب تاللّه الحقّ لو یغسل
احد ارجل العالم و یعبد اللّه علی الادغال و الشّواجن و الجبال و القنان و
الشناخیب و عند کلّ حجر و شجر و مدر و لا یتضوّع منه عرف رضآئی لن یقبل ابداً
هذا ما حکم به مولی الانام کم من عبد اعتزل فی جزائر الهند و منع عن نفسه
ما احلّه اللّه له و حمل الرّیاضات و المشقّات و لم یذکر عند اللّه منزل الآیات
لا تجعلوا الأعمال شرک الآمال و لا تحرموا انفسکم عن هذا المأل الّذی کان
امل المقرّبین فی ازل الازال قل روح الاعمال هو رضائی و علّق کلّ شئ بقبولی
اقرؤا الالواح لتعرفوا ما هو المقصود فی کتب اللّه العزیز الوهّاب من فاز
بحبّی حقّ له ان یقعد علی سریر العقیان فی صدر الامکان و الّذی منع عنه لو
یقعد علی التّراب انّه یستعیذ منه الی اللّه مالک الأدیان

\end{otherlanguage}

Mina al-nāsi man yaqʿudu ṣaffa al-niʿāli ṭalaban li-ṣadri al-jalāl. Qul
man anta yā ayyuhā al-ghāfilu al-gharrār. Wa-minhum man yaddaʿī
al-bāṭina wa-bāṭina al-bāṭin. Qul yā ayyuhā al-kadhdhāb, ta-Allāhi mā
ʿindaka innahu mina al-qushūri taraknāhā lakum kamā tutraku al-ʿiẓāmu
li-l-kilāb. Ta-Allāhi al-ḥaqq, law yaghsilu aḥadun arjula al-ʿālami
wa-yaʿbudu Allāha ʿalā al-adghāli wa-l-shawājini wa-l-jibāli wa-l-qināni
wa-l-shanākhībi wa-ʿinda kulli ḥajarin wa-shajarin wa-madarin wa-lā
yataḍawwaʿu minhu ʿarfu riḍāʾī lan yuqbala abadan. Hādhā mā ḥakama bihi
mawlā al-anām. Kam min ʿabdin iʿtazala fī jazāʾiri al-Hindi wa-manaʿa
ʿan nafsihi mā aḥallahu Allāhu lahu wa-ḥamala al-riyāḍāti
wa-l-mashaqqāti wa-lam yudhkar ʿinda Allāhi munzili al-āyāt. Lā tajʿalū
al-aʿmāla sharaka al-āmāli wa-lā taḥrimū anfusakum ʿan hādhā al-maʾāli
alladhī kāna amala al-muqarrabīna fī azali al-āzāl. Qul rūḥu al-aʿmāli
huwa riḍāʾī wa-ʿulliqa kullu shayʾin bi-qabūlī. Iqraʾū al-alwāḥa
li-taʿrifū mā huwa al-maqṣūdu fī kutubi Allāhi al-ʿazīzi al-wahhāb. Man
fāza bi-ḥubbī ḥaqqa lahu an yaqʿuda ʿalā sarīri al-ʿiqyāni fī ṣadri
al-imkān, wa-alladhī muniʿa ʿanhu law yaqʿudu ʿalā al-turābi innahu
yastaʿīdhu minhu ilā Allāhi māliki al-adyān.

Among the people is one who sits among the shoes in pursuit of the seat
of honor. Say: Who are you, heedless deceiver? Among them is one who
claims knowledge of the inner meaning and its innermost meaning. Say:
Liar! By God, what you possess is only husks that We have left to you,
as bones are left to dogs. By God, the True One, even if someone were to
wash the feet of the whole world and worship God in thickets, ravines,
mountains, peaks, and lofty summits, and beside every stone, tree, and
clod of earth, yet the fragrance of My good pleasure did not waft from
him, he would never be accepted. This is what the Lord of all people has
decreed. How many a servant has withdrawn to the islands of India,
denied himself what God has made lawful for him, and borne austerities
and hardships, yet has not been remembered before God, the Revealer of
the verses. Do not make your deeds a snare for your hopes, and do not
deprive yourselves of this ultimate goal, which has been the hope of
those brought near to God from all eternity. Say: The spirit of deeds is
My good pleasure, and everything depends upon My acceptance. Read the
Tablets, that you may know what is intended in the Books of God, the
Almighty, the All-Bountiful. Whoever has attained My love is entitled to
sit on a throne of pure gold in the seat of honor in the world of being.
But if one who has been deprived of it should sit upon the dust, the
dust would seek refuge from him with God, the Lord of religions.

\subsection{38}\label{section-37}

\begin{otherlanguage}{persian}

من یدّعی امراً قبل اتمام الف سنةٍ کاملةٍ انّه کذّاب مفترٍ نسئل اللّه بان یؤیّده
علی الرّجوع ان تاب انّه هو التّوّاب و ان اصرّ علی ما قال یبعث علیه من لا
یرحمه انّه شدید العقاب من یأوّل هذه الآیة او یفسّرها بغیر ما نزّل فی الظّاهر
انّه محروم من روح اللّه و رحمته الّتی سبقت العالمین خافوا اللّه و لا تتّبعوا
ما عندکم من الاوهام اتّبعوا ما یأمرکم به ربّکم العزیز الحکیم سوف یرتفع
النّعاق من اکثر البلدان اجتنبوا یا قوم و لا تتّبعوا کلّ فاجرٍ لئیمٍ هذا ما
اخبرناکم به اذ کنّا فی العراق و فی ارض السّرّ و فی هذا المنظر المنیر

\end{otherlanguage}

Man yaddaʿī amran qabla itmāmi alfi sanatin kāmilatin innahu kadhdhābun
muftarin. Nasʾalu Allāha bi-an yuʾayyidahu ʿalā al-rujūʿ. In tāba innahu
huwa al-tawwāb, wa-in aṣarra ʿalā mā qāla yabʿathu ʿalayhi man lā
yarḥamuhu. Innahu shadīdu al-ʿiqāb. Man yuʾawwilu hādhihi al-āyata aw
yufassiruhā bi-ghayri mā nuzzila fī al-ẓāhiri innahu maḥrūmun min rūḥi
Allāhi wa-raḥmatihi allatī sabaqati al-ʿālamīn. Khāfū Allāha wa-lā
tattabiʿū mā ʿindakum mina al-awhām. Ittabiʿū mā yaʾmurukum bihi
rabbukumu al-ʿazīzu al-ḥakīm. Sawfa yartafiʿu al-nuʿāqu min akthari
al-buldān. Ijtanibū yā qawmi wa-lā tattabiʿū kulla fājirin laʾīm. Hādhā
mā akhbarnākum bihi idh kunnā fī al-ʿIrāqi wa-fī arḍi al-sirri wa-fī
hādhā al-manẓari al-munīr.

Whoever claims a revelation before the completion of a full thousand
years is a lying impostor. We ask God to help him turn back. If he
repents, God is indeed the Ever-Relenting; but if he persists in what he
has said, He will raise up against him one who will show him no mercy.
He is indeed severe in punishment. Whoever interprets this verse or
explains it contrary to its revealed, outward meaning is deprived of the
Spirit of God and His mercy, which has encompassed all worlds. Fear God
and do not follow your own imaginings. Follow what your Lord, the
Almighty, the All-Wise, commands you. Clamor will soon arise from most
lands. Shun it, My people, and do not follow any wicked, base person.
This is what We told you when We were in Iraq, in the Land of Mystery,
and in this radiant place.

\subsection{39}\label{section-38}

\begin{otherlanguage}{arabic}

یا اهل الأرض اذا غربت شمس جمالی و سترت سمآء هیکلی لا تضطربوا قوموا علی
نصرة امری و ارتفاع کلمتی بین العالمین انّا معکم فی کلّ الأحوال و ننصرکم
بالحقّ انّا کنّا قادرین من عرفنی یقوم علی خدمتی بقیامٍ لا تقعده جنود السّموات
و الأرضین

\end{otherlanguage}

Yā ahla al-arḍ, idhā gharabat shamsu jamālī wa-sutirat samāʾu haykalī lā
taḍṭaribū. Qūmū ʿalā nuṣrati amrī wa-irtifāʿi kalimatī bayna al-ʿālamīn.
Innā maʿakum fī kulli al-aḥwāli wa-nanṣurukum bi-l-ḥaqq. Innā kunnā
qādirīn. Man ʿarafanī yaqūmu ʿalā khidmatī bi-qiyāmin lā tuqʿiduhu
junūdu al-samāwāti wa-l-arḍīn.

People of the earth! When the sun of My beauty has set and the heaven of
My person has been concealed, do not be troubled. Arise to champion My
Cause and exalt My Word among the peoples of the world. We are with you
in all circumstances and will assuredly aid you. We are indeed powerful.
Whoever has recognized Me will arise to serve Me with such resolve that
the hosts of the heavens and the earth cannot hold him back.

\subsection{40}\label{section-39}

\begin{otherlanguage}{arabic}

انّ الناس نیام لو انتبهوا سرعوا بالقلوب الی اللّه العلیم الحکیم و نبذوا ما
عندهم ولو کان کنوز الدّنیا کلّها لیذکرهم مولیهم بکلمةٍ من عنده کذلک ینبّئکم
من عنده علم الغیب فی لوحٍ ما ظهر فی الامکان و ما اطّلع به الاّ نفسه
المهیمنة علی العالمین قد اخذهم سکر الهوی علی شأنٍ لا یرون مولی الوری الّذی
ارتفع ندآئه من کلّ الجهات لا اله الاّ انا العزیز الحکیم

\end{otherlanguage}

Inna al-nāsa niyāmun, law intabahū saraʿū bi-l-qulūbi ilā Allāhi
al-ʿalīmi al-ḥakīmi wa-nabadhū mā ʿindahum wa-law kāna kunūza al-dunyā
kullahā li-yadhkurahum mawlāhum bi-kalimatin min ʿindihi. Kadhālika
yunabbiʾukum man ʿindahu ʿilmu al-ghaybi fī lawḥin mā ẓahara fī
al-imkāni wa-mā iṭṭalaʿa bihi illā nafsuhu al-muhayminatu ʿalā
al-ʿālamīn. Qad akhadhahum sukru al-hawā ʿalā shaʾnin lā yarawna mawlā
al-warā alladhī irtafaʿa nidāʾuhu min kulli al-jihāt: lā ilāha illā anā
al-ʿazīzu al-ḥakīm.

People are asleep. Were they to awaken, they would hasten in their
hearts to God, the All-Knowing, the All-Wise, and cast away all they
possess, even if it were all the treasures of the world, so that their
Lord might remember them with a word from Him. Thus does He who
possesses knowledge of the unseen inform you in a Tablet that has not
appeared in the world of being and of which none but His own Self, the
Guardian over all worlds, has knowledge. The intoxication of desire has
so seized them that they do not see the Lord of creation, whose call has
risen from every direction: ``There is no God but Me, the Almighty, the
All-Wise.''

\subsection{41}\label{section-40}

\begin{otherlanguage}{arabic}

قل لا تفرحوا بما ملکتموه فی العشی و فی الاءشراق یملکه غیرکم کذلک یخبرکم
العلیم الخبیر قل هل رایتم لما عندکم من قرارٍ او وفاءٍ لا و نفسی الرّحمن لو
انتم من المنصفین تمرّ ایام حیوتکم کما تمرّ الاریاح و یطوی بساط عزّکم کما
طوی بساط الأوّلین تفکّروا یا قوم این ایّامکم الماضیة و این اعصارکم الخالیة
طوبی لأیّام مضت بذکر اللّه و لاوقات صرفت فی ذکره الحکیم لعمری لا تبقی عزّة
الاعزّآء و لا زخارف الاغنیآء و لا شوکة الاشقیآء سیفنی الکلّ بکلمة من عنده
انّه لهو المقتدر العزیز القدیر لا ینفع النّاس ما عندهم من الاثاث و ما
ینفعهم غفلوا عنه سوف ینتبهون و لا یجدون ما فات عنهم فی ایّام ربّهم العزیز
الحمید لو یعرفون ینفقون ما عندهم لتذکر اسمائهم لدی العرش الا انّهم من
المیتین

\end{otherlanguage}

Qul lā tafraḥū bi-mā malaktumūhu fī al-ʿashī wa-fī al-ishrāq yamlikuhu
ghayrukum. Kadhālika yukhbirukum al-ʿAlīm al-Khabīr. Qul hal raʾaytum
li-mā ʿindakum min qarārin aw wafāʾin? Lā wa-nafsī al-Raḥmān law antum
min al-munṣifīn. Tamurru ayyāmu ḥayātikum ka-mā tamurru al-aryāḥ
wa-yuṭwā bisāṭu ʿizzikum ka-mā ṭuwiya bisāṭu al-awwalīn. Tafakkarū yā
qawm ayna ayyāmukum al-māḍiya wa-ayna aʿṣārukum al-khāliya. Ṭūbā
li-ayyāmin maḍat bi-dhikri Allāh wa-li-awqātin ṣurifat fī dhikrihi
al-Ḥakīm. La-ʿamrī lā tabqā ʿizzatu al-aʿizzāʾ wa-lā zakhārifu
al-aghniyāʾ wa-lā shawkatu al-ashqiyāʾ. Sa-yafnā al-kullu bi-kalimatin
min ʿindihi innahu la-huwa al-Muqtadir al-ʿAzīz al-Qadīr. Lā yanfaʿu
al-nāsa mā ʿindahum min al-athāth wa-mā yanfaʿuhum ghafalū ʿanhu. Sawfa
yantabihūna wa-lā yajidūna mā fāta ʿanhum fī ayyāmi Rabbihim al-ʿAzīz
al-Ḥamīd. Law yaʿrifūna yunfiqūna mā ʿindahum li-tudhkara asmāʾuhum ladā
al-ʿarsh. Alā innahum min al-mayyitīn.

Say: Do not rejoice in what you possess in the evening, for at dawn
someone else will possess it. Thus the All-Knowing, the All-Informed
informs you. Say: Have you found any permanence or faithfulness in what
you possess? No, by My own Self, the All-Merciful, if you judge fairly.
The days of your lives pass like the winds, and the carpet of your glory
will be rolled up as the carpet of those before you was rolled up.
Reflect, people: where are your past days, and where are your bygone
ages? Blessed are the days spent remembering God and the hours devoted
to remembering Him, the All-Wise. By My life, neither the glory of the
mighty, nor the finery of the rich, nor the power of the wicked will
endure. All will perish at a word from Him. He is indeed the
All-Powerful, the Almighty, the Omnipotent. People's possessions do not
benefit them, while they neglect what does benefit them. Soon they will
awaken and will not find what they missed in the days of their Lord, the
Almighty, the Praiseworthy. If they understood, they would spend all
they possess so that their names might be mentioned before the Throne.
Yet they are indeed among the dead.

\subsection{42}\label{section-41}

\begin{otherlanguage}{arabic}

من النّاس من غرّته العلوم و بها من اسمی القیّوم و اذا سمع صوت النّعال عن
خلفه یری نفسه اکبر من نمرود قل این هو یا ایها المردود تاللّه انّه لفی اسفل
الجحیم قل یا معشر العلمآء اما تسمعون صریر قلمی الأعلی و اما ترون هذه
الشّمس المشرقة من الأفق الأبهی الی م اعتکفتم علی اصنام اهوائکم دعوا
الأوهام و توجّهوا الی اللّه مولیکم القدیم

\end{otherlanguage}

Min al-nāsi man gharrathu al-ʿulūm wa-bihā {[}corrupted text{]} ismī
al-Qayyūm wa-idhā samiʿa ṣawta al-niʿāl ʿan khalfihi yarā nafsahu akbara
min Nimrūd. Qul ayna huwa yā ayyuhā al-mardūd? Ta-Allāh innahu la-fī
asfali al-jaḥīm. Qul yā maʿshara al-ʿulamāʾ a-mā tasmaʿūna ṣarīra qalamī
al-aʿlā wa-a-mā tarawna hādhihi al-shamsa al-mushriqata min al-ufuq
al-abhā? Ilā ma iʿtakaftum ʿalā aṣnāmi ahwāʾikum? Daʿū al-awhām
wa-tawajjahū ilā Allāh mawlākum al-Qadīm.

Among the people is one whom learning has deluded, and through it
{[}corrupted original text{]} My name, the Self-Subsisting. When he
hears the sound of sandals behind him, he considers himself greater than
Nimrod. Say: Where is he, you who have been rejected? By God, he is in
the lowest depths of hell. Say: Scholars! Do you not hear the scratching
of My Most Exalted Pen, and do you not see this Sun shining from the
most glorious horizon? Why have you remained devoted to the idols of
your desires? Abandon your imaginings and turn toward God, your ancient
Lord.

\subsection{43}\label{section-42}

\begin{otherlanguage}{arabic}

قد رجعت الاوقاف المختصّة للخیرات الی اللّه مظهر الآیات لیس لأحد ان یتصرّف
فیها الاّ بعد اذن مطلع الوحی و من بعده یرجع الحکم الی الأغصان و من بعدهم
الی بیت العدل ان تحقّق امره فی البلاد لیصرفوها فی البقاع المرتفعة فی هذا
الأمر و فیما امروا به من لدن مقتدر قدیر والاّ ترجع الی اهل البهآء الّذین
لا یتکلّمون الاّ بعد اذنه و لا یحکمون الاّ بما حکم اللّه فی هذا اللّوح اولٓئک
اولیآء النّصر بین السّموات و الأرضین لیصرفوها فیما حدّد فی الکتاب من لدن
عزیز کریم

\end{otherlanguage}

Qad rajaʿat al-awqāfu al-mukhtaṣṣatu li-l-khayrāt ilā Allāh muẓhir
al-āyāt. Laysa li-aḥadin an yataṣarrafa fīhā illā baʿda idhni maṭlaʿ
al-waḥy wa-min baʿdihi yarjiʿu al-ḥukmu ilā al-Aghṣān wa-min baʿdihim
ilā Bayt al-ʿAdl in taḥaqqaqa amruhu fī al-bilād li-yaṣrifūhā fī
al-biqāʿ al-murtafiʿa fī hādhā al-amr wa-fī-mā umirū bihi min ladun
Muqtadir Qadīr. Wa-illā tarjiʿu ilā ahl al-Bahāʾ alladhīna lā
yatakallamūna illā baʿda idhnihi wa-lā yaḥkumūna illā bi-mā ḥakama Allāh
fī hādhā al-lawḥ. Ulāʾika awliyāʾu al-naṣr bayna al-samāwāt wa-l-araḍīn
li-yaṣrifūhā fī-mā ḥuddida fī al-kitāb min ladun ʿAzīz Karīm.

Endowments dedicated to charitable purposes have reverted to God, the
Revealer of signs. No one may dispose of them without permission from
the Dawning Place of Revelation. After Him, authority passes to the
Branches, and after them to the House of Justice, if it has been
established in the lands, so that they may spend them on the places
exalted in this Cause and on whatever they have been commanded by One
who is All-Powerful and Omnipotent. Otherwise, they revert to the people
of Bahá who speak only with His permission and judge only according to
what God has decreed in this Tablet. These are the champions of victory
between the heavens and the earth. They are to spend them as prescribed
in the Book by One who is Almighty and Generous.

\subsection{44}\label{section-43}

\begin{otherlanguage}{arabic}

لا تجزعوا فی المصائب و لا تفرحوا ابتغوا امراً بین الأمرین هو التّذکّر فی
تلک الحالة و التّنبّه علی ما یرد علیکم فی العاقبة کذلک ینبئکم العلیم
الخبیر

\end{otherlanguage}

Lā tajzaʿū fī al-maṣāʾib wa-lā tafraḥū. Ibtaghū amran bayna al-amrayn
huwa al-tadhakkur fī tilka al-ḥāla wa-l-tanabbuh ʿalā mā yaridu ʿalaykum
fī al-ʿāqiba. Kadhālika yunabbiʾukum al-ʿAlīm al-Khabīr.

Do not despair in adversity, nor rejoice. Seek a course between the two:
remembrance in that condition and awareness of what will come to you in
the end. Thus the All-Knowing, the All-Informed instructs you.

\subsection{45}\label{section-44}

\begin{otherlanguage}{arabic}

لا تحلقوا رؤسکم قد زینها اللّه بالشّعر و فی ذلک لآیات لمن ینظر الی مقتضیات
الطّبیعة من لدن مالک البریه انّه لهو العزیز الحکیم و لا ینبغی ان یتجاوز عن
حدّ الآذان هذا ما حکم به مولی العالمین

\end{otherlanguage}

Lā taḥliqū ruʾūsakum qad zayyanahā Allāh bi-l-shaʿr wa-fī dhālika
la-āyātun li-man yanẓuru ilā muqtaḍayāt al-ṭabīʿa min ladun Mālik
al-bariyya. Innahu la-huwa al-ʿAzīz al-Ḥakīm. Wa-lā yanbaghī an yata
jāwaza ʿan ḥadd al-ādhān. Hādhā mā ḥakama bihi Mawlā al-ʿālamīn.

Do not shave your heads. God has adorned them with hair, and in this are
signs for anyone who considers the requirements of nature ordained by
the Sovereign of creation. He is indeed the Almighty, the All-Wise. Hair
should not extend beyond the ears. This is what the Lord of the worlds
has decreed.

\subsection{46}\label{section-45}

\begin{otherlanguage}{arabic}

قد کتب علی السّارق النّفی و الحبس و فی الثّالث فاجعلوا فی جبینه علامةً یعرف
بها لئلاّ تقبله مدن اللّه و دیاره ایّاکم ان تأخذکم الرّأفة فی دین اللّه
اعملوا ما امرتم به من لدن مشفق رحیم انّا ربّیناکم بسیاط الحکمة و الأحکام
حفظاً لانفسکم و ارتفاعاً لمقاماتکم کما یربّی الابآء ابنائهم لعمری لو تعرفون
ما اردناه لکم من اوامرنا المقدّسة لتفدون ارواحکم لهذا الأمر المقدّس العزیز
المنیع

\end{otherlanguage}

Qad kutiba ʿalā al-sāriq al-nafy wa-l-ḥabs wa-fī al-thālith fa-jʿalū fī
jabīnihi ʿalāmatan yuʿrafu bihā li-allā taqbalahu mudunu Allāh
wa-diyāruhu. Iyyākum an taʾkhudhakum al-raʾfatu fī dīn Allāh. Iʿmalū mā
umirtum bihi min ladun Mushfiq Raḥīm. Innā rabbaynākum bi-siyāṭ al-ḥikma
wa-l-aḥkām ḥifẓan li-anfusikum wa-irtifāʿan li-maqāmātikum ka-mā yurabbī
al-ābāʾ abnāʾahum. La-ʿamrī law taʿrifūna mā aradnāhu lakum min
awāmirinā al-muqaddasa la-tafdūna arwāḥakum li-hādhā al-amr al-muqaddas
al-ʿAzīz al-Manīʿ.

Exile and imprisonment have been prescribed for the thief. On the third
offense, place a mark on his forehead by which he may be recognized, so
that the cities and lands of God will not admit him. Do not let pity
restrain you in carrying out God's religion. Do what you have been
commanded by One who is Compassionate and Merciful. We have trained you
with the lashes of wisdom and laws to protect you and elevate your
stations, just as fathers bring up their children. By My life, if you
understood what We have intended for you through Our sacred
commandments, you would sacrifice your lives for this holy, mighty,
unassailable Cause.

\subsection{47}\label{section-46}

\begin{otherlanguage}{arabic}

من اراد ان یستعمل اوانی الذّهب و الفضّة لا بأس علیه ایّاکم ان تنغمس ایادیکم
فی الصّحاف و الصّحان خذوا ما یکون اقرب الی اللّطافة انّه اراد ان یریکم علی
آداب اهل الرّضوان فی ملکوته الممتنع المنیع تمسّکوا باللّطافة فی کلّ الأحوال
لئلاّ تقع العیون علی ما تکرهه انفسکم و اهل الفردوس و الّذی تجاوز عنها یحبط
عمله فی الحین و ان کان له عذر یعفو اللّه عنه انّه لهو العزیز الکریم

\end{otherlanguage}

Man arāda an yastaʿmila awāniya al-dhahab wa-l-fiḍḍa lā baʾsa ʿalayhi.
Iyyākum an tanghamisa ayādīkum fī al-ṣiḥāf wa-l-ṣiḥān. Khudhū mā yakūnu
aqraba ilā al-laṭāfa. Innahu arāda an yuriyakum ʿalā ādāb ahl al-Riḍwān
fī malakūtihi al-mumtaniʿ al-manīʿ. Tamassakū bi-l-laṭāfa fī kull
al-aḥwāl li-allā taqaʿa al-ʿuyūnu ʿalā mā takrahuhu anfusukum wa-ahlu
al-firdaws. Wa-alladhī tajāwaza ʿanhā yaḥbaṭu ʿamaluhu fī al-ḥīn wa-in
kāna lahu ʿudhrun yaʿfū Allāh ʿanhu. Innahu la-huwa al-ʿAzīz al-Karīm.

Anyone who wishes to use vessels of gold and silver may do so. Take care
not to plunge your hands into bowls and platters. Adopt what is most in
keeping with refinement. He wishes to see you observe the manners of the
people of Paradise in His inaccessible, unassailable Kingdom. Adhere to
refinement in all circumstances, so that no eyes fall upon anything
repugnant to you or to the people of Paradise. Whoever departs from it
will have his deeds immediately brought to nothing; but if he has an
excuse, God will pardon him. He is indeed the Almighty, the Generous.

\subsection{48}\label{section-47}

\begin{otherlanguage}{arabic}

لیس لمطلع الأمر شریک فی العصمة الکبری انّه لمظهر یفعل ما یشآء فی ملکوت
الانشآء قد خصّ اللّه هذا المقام لنفسه و ما قدّر لأحد نصیب من هذا الشّأن
العظیم المنیع هذا امر اللّه قد کان مستوراً فی حجب الغیب اظهرناه فی هذا
الظّهور و به خرقنا حجاب الّذین ما عرفوا حکم الکتاب و کانوا من الغافلین

\end{otherlanguage}

Laysa li-maṭlaʿ al-amr sharīkun fī al-ʿiṣma al-kubrā. Innahu la-maẓharu
yafʿalu mā yashāʾ fī malakūt al-inshāʾ. Qad khaṣṣa Allāh hādhā al-maqāma
li-nafsihi wa-mā qaddara li-aḥadin naṣīban min hādhā al-shaʾn al-ʿaẓīm
al-manīʿ. Hādhā amru Allāh qad kāna mastūran fī ḥujub al-ghayb aẓharnāhu
fī hādhā al-ẓuhūr wa-bihi kharaqnā ḥijāba alladhīna mā ʿarafū ḥukma
al-kitāb wa-kānū min al-ghāfilīn.

The Dawning Place of the Cause has no partner in the Most Great
Infallibility. He is indeed the manifestation of ``He does whatever He
wills'' in the kingdom of creation. God has reserved this station for
Himself and has allotted no one a share in this exalted, unassailable
rank. This is God's decree, formerly concealed behind the veils of the
unseen. We have disclosed it in this Revelation, and through it We have
torn apart the veil of those who did not understand the ruling of the
Book and were among the heedless.

\subsection{49}\label{section-48}

\begin{otherlanguage}{arabic}

کتب علی کلّ اب تربیة ابنه و بنته بالعلم و الخطّ و دونهما عمّا حدّد فی اللّوح
و الّذی ترک ما امر به فللأمنآء ان یأخذوا منه ما یکون لازماً لتربیتهما ان
کان غنیاً والّا یرجع الی بیت العدل انّا جعلناه مأوی الفقرآء و المساکین انّ
الّذی ربّی ابنه او ابناً من الأبنآء کأنّه ربّی احد ابنائی علیه بهائی و عنایتی
و رحمتی الّتی سبقت العالمین

\end{otherlanguage}

Kutiba ʿalā kull ab tarbiyatu ibnihi wa-bintihi bi-l-ʿilm wa-l-khaṭṭ
wa-dūnihimā ʿammā ḥuddida fī al-lawḥ. Wa-alladhī taraka mā umira bihi
fa-li-l-umanāʾ an yaʾkhudhū minhu mā yakūnu lāziman li-tarbiyatihimā in
kāna ghaniyyan wa-illā yarjiʿu ilā Bayt al-ʿAdl. Innā jaʿalnāhu maʾwā
al-fuqarāʾ wa-l-masākīn. Inna alladhī rabbā ibnahu aw ibnan min al-abnāʾ
ka-annahu rabbā aḥada abnāʾī. ʿAlayhi bahāʾī wa-ʿināyatī wa-raḥmatī
allatī sabaqat al-ʿālamīn.

Every father is required to educate his son and daughter in knowledge
and writing, and in the other subjects prescribed in the Tablet. If
anyone neglects what he has been commanded, the trustees may take from
him what is needed for their education if he is wealthy; otherwise, the
matter devolves upon the House of Justice. We have made it a refuge for
the poor and needy. Whoever educates his own child or another child is
as though he had educated one of My children. Upon him rest My glory, My
loving care, and My mercy, which has encompassed all worlds.

\subsection{50}\label{section-49}

\begin{otherlanguage}{arabic}

قد حکم اللّه لکلّ زان و زانیة دیة مسلّمة الی بیت العدل و هی تسعة مثاقیل من
الذّهب و ان عادا مرّة اخری عودوا بضعف الجزآء هذا ما حکم به مالک الأسمآء فی
الأولی و فی الأخری قدّر لهما عذاب مهین من ابتلی بمعصیة فله ان یتوب و یرجع
الی اللّه انّه یغفر لمن یشآء و لا یسئل عمّا شآء انّه لهو التوّاب العزیز
الحمید

\end{otherlanguage}

Qad ḥakama Allāh li-kull zānin wa-zāniyatin diyatan musallamatan ilā
Bayt al-ʿAdl wa-hiya tisʿatu mathāqīla min al-dhahab wa-in ʿādā marratan
ukhrā ʿūdū bi-ḍiʿf al-jazāʾ. Hādhā mā ḥakama bihi Mālik al-asmāʾ fī
al-ūlā wa-fī al-ukhrā qaddara lahumā ʿadhābun muhīn. Man ibtuliya
bi-maʿṣiyatin fa-lahu an yatūba wa-yarjiʿa ilā Allāh. Innahu yaghfiru
li-man yashāʾ wa-lā yusʾalu ʿammā shāʾ. Innahu la-huwa al-Tawwāb
al-ʿAzīz al-Ḥamīd.

God has decreed that every adulterer and adulteress pay a fine to the
House of Justice of nine mithqals of gold. If they repeat the offense,
double the penalty. This is what the Sovereign of Names has decreed in
this world; in the world to come He has ordained a humiliating
punishment for them both. Anyone who commits a sin may repent and return
to God. He forgives whomever He wills and is not questioned about what
He wills. He is indeed the Ever-Forgiving, the Almighty, the
Praiseworthy.

\subsection{51}\label{section-50}

\begin{otherlanguage}{arabic}

ایاکم ان تمنعکم سبحات الجلال عن زلال هذا السّلسال خذوا اقداح الفلاح فی
هذا الصّباح باسم فالق الاصباح ثمّ اشربوا بذکره العزیز البدیع

\end{otherlanguage}

Iyyākum an tamnaʿakum subuḥātu al-jalāl ʿan zulāl hādhā al-salsāl.
Khudhū aqdāḥa al-falāḥ fī hādhā al-ṣabāḥ bi-smi Fāliq al-iṣbāḥ thumma
ishrabū bi-dhikrihi al-ʿAzīz al-Badīʿ.

Do not let the veils of glory keep you from the pure water of this
flowing spring. Take up the cups of salvation this morning in the name
of Him who causes the dawn to break, then drink in remembrance of Him,
the Almighty, the Incomparable.

\subsection{52}\label{section-51}

\begin{otherlanguage}{arabic}

انّا حلّلنا لکم اصغآء الأصوات و النّغمات ایّاکم ان یخرجکم الاصغآء عن شأن
الأدب و الوقار افرحوا بفرح اسمی الأعظم الّذی به تولّهت الافئدة و انجذبت
عقول المقرّبین انّا جعلناه مرقاةً لعروج الأرواح الی الأفق الأعلی لا تجعلوه
جناح النّفس و الهوی انّی اعوذ ان تکونوا من الجاهلین

\end{otherlanguage}

Innā ḥallalnā lakum iṣghāʾa al-aṣwāt wa-l-naghamāt. Iyyākum an
yukhrijakum al-iṣghāʾ ʿan shaʾn al-adab wa-l-waqār. Ifraḥū bi-faraḥ ismī
al-aʿẓam alladhī bihi tawallahat al-afʾida wa-injadhabat ʿuqūlu
al-muqarrabīn. Innā jaʿalnāhu mirqātan li-ʿurūj al-arwāḥ ilā al-ufuq
al-aʿlā. Lā tajʿalūhu janāḥa al-nafs wa-l-hawā. Innī aʿūdhu an takūnū
min al-jāhilīn.

We have made it lawful for you to listen to voices and melodies. Take
care that listening does not lead you beyond the bounds of propriety and
dignity. Rejoice in the joy of My Most Great Name, through which hearts
have been enraptured and the minds of those brought near have been drawn
in attraction. We have made it a ladder for souls to ascend to the
highest horizon. Do not make it a wing for self and desire. I seek
refuge from your being among the ignorant.

\subsection{53}\label{section-52}

\begin{otherlanguage}{arabic}

قد ارجعنا ثلث الدّیات کلّها الی مقرّ العدل و نوصی رجاله بالعدل الخالص
لیصرفوا ما اجتمع عندهم فیما امروا به من لدن علیم حکیم یا رجال العدل
کونوا رعاة اغنام اللّه فی مملکته و احفظوهم عن الذّئاب الّذین ظهروا بالأثواب
کما تحفظون ابنائکم کذلک ینصحکم النّاصح الأمین

\end{otherlanguage}

Qad arjaʿnā thultha al-diyāt kullihā ilā maqarr al-ʿadl wa-nūṣī rijālahu
bi-l-ʿadl al-khāliṣ li-yaṣrifū mā ijtamaʿa ʿindahum fī-mā umirū bihi min
ladun ʿAlīm Ḥakīm. Yā rijāl al-ʿadl kūnū ruʿāta aghnām Allāh fī
mamlakatihi wa-iḥfaẓūhum ʿan al-dhiʾāb alladhīna ẓaharū bi-l-athwāb
ka-mā taḥfaẓūna abnāʾakum. Kadhālika yanṣaḥukum al-Nāṣiḥ al-Amīn.

We have assigned a third of all fines to the seat of justice, and We
counsel its members to practice pure justice, so that they may spend
what has accumulated with them as they have been commanded by One who is
All-Knowing and All-Wise. Members of the House of Justice! Be shepherds
of God's sheep in His kingdom, and protect them from the wolves that
have appeared in disguise, just as you protect your own children. Thus
the faithful Counselor advises you.

\subsection{54}\label{section-53}

\begin{otherlanguage}{arabic}

اذا اختلفتم فی امر فارجعوه الی اللّه ما دامت الشّمس مشرقة من افق هذه
السّماء و اذا غربت ارجعوا الی ما نزّل من عنده انّه لیکفی العالمین قل یا قوم
لا یأخذکم الاضطراب اذا غاب ملکوت ظهوری و سکنت امواج بحر بیانی انّ فی
ظهوری لحکمة و فی غیبتی حکمة اخری ما اطّلع بها الاّ اللّه الفرد الخبیر و
نراکم من افقی الأبهی و ننصر من قام علی نصرة امری بجنود من الملاء الأعلی
و قبیل من الملائکة المقرّبین

\end{otherlanguage}

Idhā ikhtalaftum fī amrin fa-arjiʿūhu ilā Allāh mā dāmat al-shamsu
mushriqatan min ufuq hādhihi al-samāʾ wa-idhā gharabat irjiʿū ilā mā
nuzzila min ʿindihi. Innahu la-yakfī al-ʿālamīn. Qul yā qawm lā
yaʾkhudhkum al-iḍṭirāb idhā ghāba malakūtu ẓuhūrī wa-sakanat amwāju baḥr
bayānī. Inna fī ẓuhūrī la-ḥikmatan wa-fī ghaybatī ḥikmatun ukhrā mā
iṭṭalaʿa bihā illā Allāh al-Fard al-Khabīr. Wa-narākum min ufuqī al-abhā
wa-nanṣuru man qāma ʿalā nuṣrati amrī bi-junūdin min al-malaʾ al-aʿlā
wa-qabīlin min al-malāʾika al-muqarrabīn.

If you disagree over a matter, refer it to God as long as the Sun shines
from the horizon of this heaven. When it has set, turn to what has been
sent down from Him. It is indeed sufficient for all worlds. Say: People!
Do not be troubled when the kingdom of My revelation has disappeared and
the waves of the ocean of My utterance have grown still. In My
appearance there is a wisdom, and in My absence there is another wisdom,
known to none but God, the One, the All-Informed. We see you from My
most glorious horizon and will aid anyone who arises to champion My
Cause with hosts from the Concourse on high and a company of angels
brought near.

\subsection{55}\label{section-54}

\begin{otherlanguage}{arabic}

یا ملاء الأرض تاللّه الحقّ قد انفجرت من الأحجار الأنهار العذبة السّائغة بما
اخذتها حلاوة بیان ربّکم المختار و انتم من الغافلین دعوا ما عندکم ثمّ طیروا
بقوادم الانقطاع فوق الابداع کذلک یأمرکم مالک الاختراع الّذی بحرکة قلمه
قلّب العالمین

\end{otherlanguage}

Yā malaʾ al-arḍ ta-Allāh al-Ḥaqq qad infajarat min al-aḥjār al-anhāru
al-ʿadhbatu al-sāʾigha bi-mā akhadhathā ḥalāwatu bayāni Rabbikum
al-Mukhtār wa-antum min al-ghāfilīn. Daʿū mā ʿindakum thumma ṭīrū
bi-qawādim al-inqiṭāʿ fawqa al-ibdāʿ. Kadhālika yaʾmurukum Mālik
al-ikhtirāʿ alladhī bi-ḥarakati qalamihi qallaba al-ʿālamīn.

People of the earth! By God, the Truth, sweet, refreshing rivers have
burst from the rocks, so seized have they been by the sweetness of the
utterance of your Lord, the Unconstrained; yet you remain heedless.
Abandon what you possess, then soar on the wings of detachment above
creation. Thus commands you the Sovereign of creation, who has
transformed the worlds with a movement of His Pen.

\subsection{56}\label{section-55}

\begin{otherlanguage}{arabic}

هل تعرفون من ایّ افقٍ ینادیکم ربّکم الأبهی و هل علمتم من ایّ قلم یأمرکم ربّکم
مالک الأسمآء لا و عمری لو عرفتم لترکتم مقبلین بالقلوب الی شطر المحبوب و
اخذکم اهتزاز الکلمة علی شأن یهتزّ منه العالم الأکبر و کیف هذا العالم
الصّغیر کذلک هطلت من سماء عنایتی امطار مکرمتی فضلاً من عندی لتکونوا من
الشّاکرین

\end{otherlanguage}

Hal taʿrifūna min ayyi ufuqin yunādīkum Rabbukum al-Abhā wa-hal ʿalimtum
min ayyi qalamin yaʾmurukum Rabbukum Mālik al-asmāʾ? Lā wa-ʿamrī law
ʿaraftum la-taraktum muqbilīna bi-l-qulūb ilā shaṭr al-maḥbūb
wa-akhadhakum ihtizāzu al-kalima ʿalā shaʾnin yahtazzu minhu al-ʿālam
al-akbar wa-kayfa hādhā al-ʿālam al-ṣaghīr. Kadhālika haṭalat min samāʾ
ʿināyatī amṭāru makrumatī faḍlan min ʿindī li-takūnū min al-shākirīn.

Do you know from what horizon your all-glorious Lord calls you, and have
you understood from what Pen your Lord, the Sovereign of Names, commands
you? No, by My life! If you knew, you would let go, turning your hearts
toward the Beloved, and the stirring of the Word would seize you with
such force that the greater world would tremble---how much more this
lesser world. Thus the rains of My generosity have poured down from the
heaven of My loving care, as a grace from Me, so that you may be
thankful.

\subsection{57}\label{section-56}

\begin{otherlanguage}{arabic}

و امّا الشّجاج و الضّرب تختلف احکامهما باختلاف مقادیرهما و حکم الدّیان لکلّ
مقدار دیة معینة انّه لهو الحاکم العزیز المنیع لو نشآء نفصّلها بالحقّ وعداً
من عندنا انّه لهو الموفّی العلیم

\end{otherlanguage}

Wa-ammā al-shijāj wa-l-ḍarb takhtalifu aḥkāmuhumā bi-ikhtilāf
maqādīrihimā wa-ḥakama al-Dayyān li-kull miqdārin diyatan muʿayyana.
Innahu la-huwa al-Ḥākim al-ʿAzīz al-Manīʿ. Law nashāʾ nufaṣṣiluhā
bi-l-ḥaqq waʿdan min ʿindinā. Innahu la-huwa al-Muwaffī al-ʿAlīm.

As for wounds and blows, their rulings vary according to their severity,
and the Judge has decreed a specified compensation for each degree. He
is indeed the Ruler, the Almighty, the Unassailable. Should We wish, We
will set them out in detail in truth---a promise from Us. He is indeed
the One who fulfills His promise, the All-Knowing.

\subsection{58}\label{section-57}

\begin{otherlanguage}{arabic}

قد رقم علیکم الضّیافة فی کلّ شهر مرّةً واحدةً ولو بالمآء انّ اللّه اراد ان یؤلّف
بین القلوب ولو باسباب السّموات و الأرضین

\end{otherlanguage}

Qad ruqima ʿalaykum al-ḍiyāfatu fī kull shahr marratan wāḥidatan wa-law
bi-l-māʾ. Inna Allāh arāda an yuʾallifa bayna al-qulūb wa-law bi-asbāb
al-samāwāt wa-l-araḍīn.

You have been enjoined to offer hospitality once each month, even if
only with water. God wishes to unite hearts, even through the means of
the heavens and the earth.

\subsection{59}\label{section-58}

\begin{otherlanguage}{arabic}

ایّاکم ان تفرّقکم شئونات النّفس و الهوی کونوا کالأصابع فی الید و الأرکان
للبدن کذلک یعظکم قلم الوحی ان انتم من الموقنین

\end{otherlanguage}

Iyyākum an tufarriqakum shuʾūnātu al-nafs wa-l-hawā. Kūnū ka-l-aṣābiʿ fī
al-yad wa-l-arkān li-l-badan. Kadhālika yaʿiẓukum Qalam al-waḥy in antum
min al-mūqinīn.

Do not let the promptings of self and desire divide you. Be like the
fingers of one hand and the limbs of one body. Thus the Pen of
Revelation counsels you, if you are among those who are certain.

\subsection{60}\label{section-59}

\begin{otherlanguage}{arabic}

فانظروا فی رحمة اللّه و الطافه انّه یأمرکم بما ینفعکم بعد اذ کان غنیاً عن
العالمین لن تضرّنا سیئاتکم کما لا تنفعنا حسناتکم انّما ندعوکم لوجه اللّه
یشهد بذلک کلّ عالم بصیر

\end{otherlanguage}

Fa-nẓurū fī raḥmat Allāh wa-alṭāfihi innahu yaʾmurukum bi-mā yanfaʿukum
baʿda idh kāna ghaniyyan ʿan al-ʿālamīn. Lan taḍurranā sayyiʾātukum
ka-mā lā tanfaʿunā ḥasanātukum. Innamā nadʿūkum li-wajh Allāh. Yashhadu
bi-dhālika kullu ʿālim baṣīr.

Consider God's mercy and kindness: He commands you to do what benefits
you, although He has no need of any of the worlds. Your evil deeds will
never harm Us, just as your good deeds do not benefit Us. We call you
only for the sake of God. Every discerning person of knowledge bears
witness to this.

\subsection{61}\label{section-60}

\begin{otherlanguage}{arabic}

اذا ارسلتم الجوارح الی الصّید اذکروا اللّه اذاً یحلّ ما امسکن لکم ولو تجدونه
میتاً انّه لهو العلیم الخبیر ایّاکم ان تسرفوا فی ذلک کونوا علی صراط العدل و
الانصاف فی کلّ الأمور کذلک یأمرکم مطلع الظّهور ان انتم من العارفین

\end{otherlanguage}

Idhā arsaltum al-jawāriḥ ilā al-ṣayd udhkurū Allāh idhan yaḥill mā
amsakna lakum wa-law tajidūnahu maytan innahu la-huwa al-ʿalīm
al-khabīr. Iyyākum an tusrifū fī dhālik kūnū ʿalā ṣirāṭ al-ʿadl
wa-al-inṣāf fī kull al-umūr. Kadhālika yaʾmurukum maṭlaʿ al-ẓuhūr in
antum min al-ʿārifīn.

When you send hunting animals after game, mention God; then whatever
they catch for you is lawful, even if you find it dead. He is indeed the
All-Knowing, the All-Informed. Beware of excess in this matter. Follow
the path of justice and fairness in all things. Thus the Dawning-Place
of Revelation commands you, if you are among those who understand.

\subsection{62}\label{section-61}

\begin{otherlanguage}{arabic}

انّ اللّه قد امرکم بالمودّة فی ذوی القربی و ما قدّر لهم حقّاً فی اموال النّاس
انّه لهو الغنی عن العالمین

\end{otherlanguage}

Inna Allāh qad amarakum bi-al-mawadda fī dhawī al-qurbā wa-mā qaddara
lahum ḥaqqan fī amwāl al-nās innahu la-huwa al-ghaniyy ʿan al-ʿālamīn.

God has commanded you to show affection toward His kindred, but has
assigned them no right to the property of others. He is indeed
independent of all the worlds.

\subsection{63}\label{section-62}

\begin{otherlanguage}{arabic}

من احرق بیتاً متعمّداً فاحرقوه و من قتل نفساً عامداً فاقتلوه خذوا سنن اللّه
بأیادی القدرة و الاقتدار ثمّ اترکوا سنَنَ الجاهلین و ان تحکموا لهما حبساً
ابدیّاً لا بأس علیکم فی الکتاب انّه لهو الحاکم علی ما یرید

\end{otherlanguage}

Man aḥraqa baytan mutaʿammidan fa-aḥriqūhu wa-man qatala nafsan ʿāmidan
fa-uqtulūhu. Khudhū sunan Allāh bi-ayādī al-qudra wa-al-iqtidār thumma
utrukū sunan al-jāhilīn. Wa-in taḥkumū lahumā ḥabsan abadiyyan lā baʾs
ʿalaykum fī al-kitāb innahu la-huwa al-ḥākim ʿalā mā yurīd.

Whoever deliberately burns down a house, burn him; whoever intentionally
kills a person, put him to death. Take hold of God's laws with the hands
of power and might, and abandon the ways of the ignorant. If you
sentence them both to lifelong imprisonment, no blame attaches to you in
the Book. He indeed rules as He wills.

\subsection{64}\label{section-63}

\begin{otherlanguage}{persian}

بواحدة من الامآء استراحت نفسه و نفسها و من اتّخذ بکراً لخدمته قد کتب اللّه
علیکم النّکاح ایّاکم ان تجاوزوا عن الاثنتین و الّذی اقتنع لا بأس علیه کذلک
کان الأمر من قلم الوحی بالحقّ مرقوماً تزوّجوا یا قوم لیظهر منکم من یذکرنی
بین عبادی هذا من امری علیکم اتخذوه لانفسکم معیناً

\end{otherlanguage}

Bi-wāḥida min al-imāʾ istarāḥat nafsuhu wa-nafsuhā wa-man ittakhadha
bikran li-khidmatihi qad kataba Allāh ʿalaykum al-nikāḥ iyyākum an
tujāwizū ʿan al-ithnatayn wa-alladhī iqtanaʿa lā baʾs ʿalayhi kadhālika
kāna al-amr min qalam al-waḥy bi-al-ḥaqq marqūman. Tazawwajū yā qawm
li-yaẓhara minkum man yadhkurunī bayna ʿibādī hādhā min amrī ʿalaykum
ittakhidhūhu li-anfusikum muʿīnan.

With one of the maidservants, his soul and hers find rest. As for
whoever takes a virgin into his service, God has prescribed marriage for
you. Beware of exceeding two, and whoever is content incurs no blame.
Thus has the decree been inscribed in truth by the Pen of Revelation.
Marry, O people, so that from you may appear someone who will remember
Me among My servants. This is My command to you; take it as an aid to
yourselves.

\subsection{65}\label{section-64}

\begin{otherlanguage}{arabic}

یا ملاء الانشآء لا تتّبعوا انفسکم انّها لأمّارة بالبغی و الفحشآء اتبعوا
مالک الأشیآء الّذی یأمرکم بالبرّ و التّقوی انّه کان عن العالمین غنیاً ایّاکم
ان تفسدوا فی الأرض بعد اصلاحها و من افسد انّه لیس منّا و نحن برٴآء منه
کذلک کان الأمر من سمآء الوحی بالحقّ مشهوداً

\end{otherlanguage}

Yā malaʾ al-inshāʾ lā tattabiʿū anfusakum innahā la-ammāra bi-al-baghy
wa-al-faḥshāʾ. Ittabiʿū mālik al-ashyāʾ alladhī yaʾmurukum bi-al-birr
wa-al-taqwā innahu kāna ʿan al-ʿālamīn ghaniyyan. Iyyākum an tufsidū fī
al-arḍ baʿda iṣlāḥihā wa-man afsada innahu laysa minnā wa-naḥnu buraʾāʾ
minhu. Kadhālika kāna al-amr min samāʾ al-waḥy bi-al-ḥaqq mashhūdan.

O people of creation, do not follow your own selves, for they constantly
urge wrongdoing and indecency. Follow the Lord of all things, who
commands you to practice righteousness and reverence for God. He is
independent of all the worlds. Beware of spreading corruption on earth
after it has been set right. Whoever spreads corruption is not one of
Us, and We disavow him. Thus has the decree been made manifest in truth
from the heaven of revelation.

\subsection{66}\label{section-65}

\begin{otherlanguage}{arabic}

انّه قد حدّد فی البیان برضآء الطّرفین انّا لمّا اردنا المحبّة و الوداد و اتّحاد
العباد لذا علّقناه باذن الأبوین بعدهما لئلاّ تقع بینهم الضّغینة و البغضآء و
لنا فیه مأرب اخری و کذلک کان الأمر مقضیّا

\end{otherlanguage}

Innahu qad ḥuddida fī al-bayān bi-riḍā al-ṭarafayn innā lammā aradnā
al-maḥabba wa-al-widād wa-ittiḥād al-ʿibād lidhā ʿallaqnāhu bi-idhn
al-abawayn baʿdahumā li-allā taqaʿa baynahum al-ḍaghīna wa-al-baghḍāʾ
wa-lanā fīhi maʾrib ukhrā wa-kadhālika kāna al-amr maqḍiyyan.

In the Bayan, it was made conditional on the consent of both parties.
Since We desire love, affection, and unity among Our servants, We have
also made it dependent on the permission of the parents after that of
the couple, so that resentment and hatred may not arise among them. We
have other purposes in this as well. Thus has the matter been decreed.

\subsection{67}\label{section-66}

\begin{otherlanguage}{arabic}

لا یحقّق الصّهار الاّ بالامهار قد قدّر للمدن تسعة عشر مثقالاً من الذّهب
الابریز و للقری من الفضّة و من اراد الزّیادة حرّم علیه ان یتجاوز عن خمسة و
تسعین مثقالا کذلک کان الأمر بالعزّ مسطورا و الّذی اقتنع بالدّرجة الأولی خیر
له فی الکتاب انّه یغنی من یشآء باسباب السّموات و الأرض و کان اللّه علی کلّ
شئً قدیراً

\end{otherlanguage}

Lā yuḥaqqaq al-ṣihār illā bi-al-imhār qad quddira li-al-mudun tisʿata
ʿashara mithqālan min al-dhahab al-ibrīz wa-li-al-qurā min al-fiḍḍa
wa-man arāda al-ziyāda ḥurrima ʿalayhi an yatajāwaza ʿan khamsa
wa-tisʿīn mithqālan kadhālika kāna al-amr bi-al-ʿizz masṭūran.
Wa-alladhī iqtanaʿa bi-al-daraja al-ūlā khayr lahu fī al-kitāb innahu
yughnī man yashāʾ bi-asbāb al-samāwāt wa-al-arḍ wa-kāna Allāh ʿalā kull
shayʾ qadīran.

Marriage is not established without a dowry. For towns, nineteen
mithqals of pure gold have been prescribed, and for villages, the same
amount of silver. Whoever wishes to increase this is forbidden to exceed
ninety-five mithqals. Thus has the decree been inscribed in majesty. For
whoever is content with the lowest amount, this is better for him
according to the Book. He enriches whomever He wills through the means
of the heavens and the earth. God has power over all things.

\subsection{68}\label{section-67}

\begin{otherlanguage}{arabic}

قد کتب اللّه لکلّ عبد اراد الخروج من وطنه ان یجعل میقاتا لصاحبته فی ایة
مدّة اراد ان اتی و وفی بالوعد انّه اتّبع امر مولاه و کان من المحسنین من قلم
الأمر مکتوبا والاّ ان اعتذر بعذر حقیقی فله ان یخبر قرینته و یکون فی غایة
الجهد للرّجوع الیها و ان فات الأمران فلها تربّص تسعة اشهر معدودات و بعد
اکمالها لا بأس علیها فی اختیار الزّوج و ان صبرت انّه یحبّ الصّابرات و
الصّابرین اعملوا اوامری و لا تتّبعوا کلّ مشرک کان فی اللّوح اثیماً و ان اتی
الخبر حین تربّصها لها ان تأخذ المعروف انّه اراد الاصلاح بین العباد و
الاماء ایّاکم ان ترتکبوا ما یحدث به العناد بینکم کذلک قضی الأمر و کان
الوعد مأتیاً و ان اتاها خبر الموت او القتل و ثبت بالشّیاع او بالعدلین لها
ان تلبث فی البیت اذا مضت اشهر معدودات لها الاختیار فیما تختار هذا ما حکم
به من کان علی الأمر قویا

\end{otherlanguage}

Qad kataba Allāh li-kull ʿabd arāda al-khurūj min waṭanihi an yajʿala
mīqātan li-ṣāḥibatihi fī ayyat mudda arād. In atā wa-wafā bi-al-waʿd
innahu ittabaʿa amr mawlāhu wa-kāna min al-muḥsinīn min qalam al-amr
maktūban wa-illā in iʿtadhara bi-ʿudhr ḥaqīqī fa-lahu an yukhbira
qarīnatahu wa-yakūna fī ghāyat al-juhd li-al-rujūʿ ilayhā wa-in fāta
al-amrān fa-lahā tarabbuṣ tisʿat ashhur maʿdūdāt wa-baʿda ikmālihā lā
baʾs ʿalayhā fī ikhtiyār al-zawj wa-in ṣabarat innahu yuḥibb al-ṣābirāt
wa-al-ṣābirīn. Iʿmalū awāmirī wa-lā tattabiʿū kull mushrik kāna fī
al-lawḥ athīman. Wa-in atā al-khabar ḥīna tarabbuṣihā lahā an taʾkhudha
al-maʿrūf innahu arāda al-iṣlāḥ bayna al-ʿibād wa-al-imāʾ iyyākum an
tartakibū mā yaḥduthu bihi al-ʿinād baynakum kadhālika quḍiya al-amr
wa-kāna al-waʿd maʾtiyyan. Wa-in atāhā khabar al-mawt aw al-qatl
wa-thabata bi-al-shiyāʿ aw bi-al-ʿadlayn lahā an talbatha fī al-bayt
idhā maḍat ashhur maʿdūdāt lahā al-ikhtiyār fīmā takhtār hādhā mā ḥakama
bihi man kāna ʿalā al-amr qawiyyan.

God has prescribed that every man who wishes to leave his homeland set a
time for his wife, of whatever duration he wishes. If he returns and
fulfills his promise, he has obeyed his Lord's command and is recorded
by the Pen of Command among those who do good. Otherwise, if he has a
genuine excuse, he must inform his wife and make every effort to return
to her. If neither occurs, she is to wait nine full months; after
completing them, she incurs no blame in choosing a husband. If she waits
patiently, He loves women and men who are patient. Follow My
commandments, and do not follow any who associates partners with God and
is recorded as sinful in the Tablet. If news arrives during her waiting
period, she may do what is right. He desires reconciliation between His
male and female servants. Beware of doing anything that creates
hostility among you. Thus has the matter been decreed, and the promise
is fulfilled. If news of his death or killing reaches her and is
established by widespread report or by two just witnesses, she is to
remain in the house. When the appointed months have passed, she is free
to choose as she wishes. This is what He who is mighty in command has
decreed.

\subsection{69}\label{section-68}

\begin{otherlanguage}{persian}

و ان حدث بینهما کدورة او کره لیس له ان یطلّقها و له ان یصبر سنة کاملة لعلّ
تسطع بینهما رائحة المحبّة و ان کملت و ما فاحت فلا بأس فی الطّلاق انّه کان
علی کلّ شئ حکیماً قد نهاکم اللّه عمّا عملتم بعد طلقات ثلاث فضلاً من عنده
لتکونوا من الشّاکرین فی لوحٍ کان من قلم الأمر مسطوراً و الّذی طلّق له
الاختیار فی الرّجوع بعد انقضاء کلّ شهر بالمودّة و الرضآء ما لم تستحصن و اذا
استحصنت تحقّق الفصل بوصل آخر و قضی الأمر الاّ بعد امر مبین کذلک کان الأمر
من مطلع الجمال فی لوح الجلال بالاجلال مرقوماً

\end{otherlanguage}

Wa-in ḥadatha baynahumā kudūra aw kurh laysa lahu an yuṭalliqahā wa-lahu
an yaṣbira sana kāmila laʿalla tasṭaʿu baynahumā rāʾiḥat al-maḥabba
wa-in kamulat wa-mā fāḥat fa-lā baʾs fī al-ṭalāq innahu kāna ʿalā kull
shayʾ ḥakīman. Qad nahākum Allāh ʿammā ʿamiltum baʿda ṭalaqāt thalāth
faḍlan min ʿindihi li-takūnū min al-shākirīn fī lawḥ kāna min qalam
al-amr masṭūran. Wa-alladhī ṭallaqa lahu al-ikhtiyār fī al-rujūʿ baʿda
inqiḍāʾ kull shahr bi-al-mawadda wa-al-riḍā mā lam tastaḥṣin wa-idhā
istaḥṣanat taḥaqqaqa al-faṣl bi-waṣl ākhar wa-quḍiya al-amr illā baʿda
amr mubīn kadhālika kāna al-amr min maṭlaʿ al-jamāl fī lawḥ al-jalāl
bi-al-ijlāl marqūman.

If estrangement or aversion arises between them, he may not divorce her,
but must wait a full year, so that perhaps the fragrance of love may
spread between them. If the year is completed and that fragrance has not
spread, there is no objection to divorce. He is wise in all things. In a
Tablet inscribed by the Pen of Command, God has forbidden you what you
formerly practiced after three divorces, as a favor from Him, so that
you may be thankful. A man who has divorced his wife may choose to
return to her at the end of any month, with affection and mutual
consent, as long as she has not remarried. If she has remarried, their
separation is confirmed by another union, and the matter is settled,
unless a manifest change occurs. Thus has the decree been inscribed with
majesty by the Dawning-Place of Beauty in the Tablet of Glory.

\subsection{70}\label{section-69}

\begin{otherlanguage}{arabic}

و الّذی سافر و سافرت معه ثمّ حدث بینهما الاختلاف فله ان یؤتیها نفقة سنة
کاملة و یرجعها الی المقرّ الّذی خرجت عنه او یسلّمها بید امین و ما تحتاج به
فی السّبیل لیبلّغها الی محلّها انّ ربّک یحکم کیف یشآء بسلطان کان علی العالمین
محیطا

\end{otherlanguage}

Wa-alladhī sāfara wa-sāfarat maʿahu thumma ḥadatha baynahumā al-ikhtilāf
fa-lahu an yuʾtiyahā nafaqat sana kāmila wa-yurjiʿahā ilā al-maqarr
alladhī kharajat ʿanhu aw yusallimahā bi-yad amīn wa-mā taḥtāju bihi fī
al-sabīl li-yuballighahā ilā maḥallihā inna rabbaka yaḥkum kayfa yashāʾ
bi-sulṭān kāna ʿalā al-ʿālamīn muḥīṭan.

If a man travels and his wife travels with him, and disagreement then
arises between them, he must provide her with maintenance for a full
year and return her to the place from which she departed, or entrust her
to a trustworthy person, together with whatever she needs for the
journey, to convey her to her home. Your Lord decrees as He wills, with
an authority encompassing all the worlds.

\subsection{71}\label{section-70}

\begin{otherlanguage}{persian}

و الّتی طلّقت بما ثبت علیها منکر لا نفقة لها ایّام تربّصها کذلک کان نیّر
الأمر من افق العدل مشهودا انّ اللّه احبّ الوصل و الوفاق و ابغض الفصل و
الطّلاق عاشروا یا قوم بالرّوح و الرّیحان لعمری سیفنی من فی الامکان و ما
یبقی هو العمل الطّیّب و کان اللّه علی ما اقول شهیدا یا عبادی اصلحوا ذات
بینکم ثمّ استمعوا ما ینصحکم به القلم الأعلی و لا تتّبعوا جبّاراً شقیا

\end{otherlanguage}

Wa-allatī ṭulliqat bi-mā thabata ʿalayhā munkar lā nafaqa lahā ayyām
tarabbuṣihā kadhālika kāna nayyir al-amr min ufuq al-ʿadl mashhūdan.
Inna Allāh aḥabba al-waṣl wa-al-wifāq wa-abghaḍa al-faṣl wa-al-ṭalāq.
ʿĀshirū yā qawm bi-al-rūḥ wa-al-rayḥān la-ʿamrī sa-yafnā man fī al-imkān
wa-mā yabqā huwa al-ʿamal al-ṭayyib wa-kāna Allāh ʿalā mā aqūl shahīdan.
Yā ʿibādī aṣliḥū dhāta baynikum thumma istamiʿū mā yanṣaḥukum bihi
al-qalam al-aʿlā wa-lā tattabiʿū jabbāran shaqiyyan.

A woman divorced because of proven misconduct has no right to
maintenance during her waiting period. Thus has the Sun of the Cause
appeared above the horizon of justice. God loves union and harmony and
detests separation and divorce. Associate with one another, O people, in
joy and gladness. By My life, all who exist will pass away; what remains
is the good deed. God bears witness to what I say. O My servants, set
right your relations with one another, then listen to the counsel of the
Most Exalted Pen, and do not follow any wretched tyrant.

\subsection{72}\label{section-71}

\begin{otherlanguage}{arabic}

ایّاکم ان تغرّنکم الدّنیا کما غرّت قوماً قبلکم اتّبعوا حدود اللّه و سننه ثمّ
اسلکوا هذا الصّراط الّذی کان بالحقّ ممدوداً انّ الّذین نبذوا البغی و الغوی و
اتّخذوا التّقوی اولئک من خیرة الخلق لدی الحقّ یذکرهم الملأ الأعلی و اهل هذا
المقام الّذی کان باسم اللّه مرفوعاً

\end{otherlanguage}

Iyyākum an taghurrannakum al-dunyā kamā gharrat qawman qablakum.
Ittabiʿū ḥudūd Allāh wa-sunanahu thumma uslukū hādhā al-ṣirāṭ alladhī
kāna bi-al-ḥaqq mamdūdan. Inna alladhīna nabadhū al-baghy wa-al-ghawā
wa-ittakhadhū al-taqwā ulāʾika min khīrat al-khalq ladā al-ḥaqq
yadhkuruhum al-malaʾ al-aʿlā wa-ahl hādhā al-maqām alladhī kāna bi-ism
Allāh marfūʿan.

Beware lest the world deceive you as it deceived people before you.
Follow God's ordinances and laws, and walk this path that has been laid
out in truth. Those who cast aside wrongdoing and waywardness and
embrace reverence for God are among the finest of His creatures in the
sight of the True One. The Concourse on high remembers them, as do the
inhabitants of this station that has been exalted in God's name.

\subsection{73}\label{section-72}

\begin{otherlanguage}{arabic}

قد حرّم علیکم بیع الامآء و الغلمان لیس لعبد ان یشتری عبداً نهیاً فی لوح
اللّه کذلک کان الأمر من قلم العدل بالفضل مسطورا و لیس لأحد ان یفتخر علی
احد کلّ ارقّاء له و ادلاّء علی انّه لا اله الاّ هو انّه کان علی کلّ شئ حکیما

\end{otherlanguage}

Qad ḥurrima ʿalaykum bayʿ al-imāʾ wa-al-ghilmān laysa li-ʿabd an
yashtarī ʿabdan nahyan fī lawḥ Allāh kadhālika kāna al-amr min qalam
al-ʿadl bi-al-faḍl masṭūran. Wa-laysa li-aḥad an yaftakhira ʿalā aḥad
kull ariqqāʾ lahu wa-adillāʾ ʿalā annahu lā ilāha illā huwa innahu kāna
ʿalā kull shayʾ ḥakīman.

You are forbidden to sell female or male slaves. No person may buy
another person: this is prohibited in the Tablet of God. Thus has the
decree been graciously inscribed by the Pen of Justice. No one may exalt
himself above another. All are His bondservants and bear witness that
there is no God but Him. He is wise in all things.

\subsection{74}\label{section-73}

\begin{otherlanguage}{arabic}

زینوا انفسکم بطراز الأعمال و الّذی فاز بالعمل فی رضاه انّه من اهل البهآء
قد کان لدی العرش مذکوراً انصروا مالک البریة بالأعمال الحسنة ثمّ بالحکمة و
البیان کذلک امرتم فی اکثر الألواح من لدی الرّحمن انّه کان علی ما اقول
علیماً لا یعترض احد علی احد و لا یقتل نفس نفساً هذا ما نهیتم عنه فی کتاب
کان فی سرادق العزّ مستورا اتقتلون من احیاه اللّه بروح من عنده انّ هذا خطأ
قد کان لدی العرش کبیرا اتّقوا اللّه و لا تخربوا ما بناه اللّه بأیادی الظّلم
و الطّغیان ثمّ اتّخذوا الی الحقّ سبیلا لما ظهرت جنود العرفان برایات البیان
انهزمت قبائل الادیان الاّ من اراد ان یشرب کوثر الحیوان فی رضوانٍ کان من
نفس السّبحان موجودا

\end{otherlanguage}

Zayyinū anfusakum bi-ṭirāz al-aʿmāl wa-alladhī fāza bi-al-ʿamal fī
riḍāhu innahu min ahl al-bahāʾ qad kāna ladā al-ʿarsh madhkūran. Unṣurū
mālik al-bariyya bi-al-aʿmāl al-ḥasana thumma bi-al-ḥikma wa-al-bayān
kadhālika umirtum fī akthar al-alwāḥ min ladā al-raḥmān innahu kāna ʿalā
mā aqūl ʿalīman. Lā yaʿtariḍ aḥad ʿalā aḥad wa-lā yaqtul nafs nafsan
hādhā mā nuhītum ʿanhu fī kitāb kāna fī surādiq al-ʿizz mastūran.
A-taqtulūna man aḥyāhu Allāh bi-rūḥ min ʿindihi inna hādhā khaṭaʾ qad
kāna ladā al-ʿarsh kabīran. Ittaqū Allāh wa-lā tukhribū mā banāhu Allāh
bi-ayādī al-ẓulm wa-al-ṭughyān thumma ittakhidhū ilā al-ḥaqq sabīlan.
Lammā ẓaharat junūd al-ʿirfān bi-rāyāt al-bayān inhazamat qabāʾil
al-adyān illā man arāda an yashraba kawthar al-ḥayawān fī riḍwān kāna
min nafs al-subḥān mawjūdan.

Adorn yourselves with the ornament of deeds. Whoever succeeds in acting
according to His good pleasure is numbered among the people of Baha
before the Throne. Aid the Lord of creation through good deeds, and
through wisdom and utterance. Thus have you been commanded in most of
the Tablets by the All-Merciful. He knows what I say. Let no one contend
with another, and let no person kill another. This is what you have been
forbidden in a Book concealed within the pavilion of majesty. Would you
kill someone whom God has brought to life through a spirit from Him?
This is a grievous wrong before the Throne. Fear God, and do not destroy
with the hands of oppression and tyranny what God has built. Take a path
toward the True One. When the hosts of understanding appeared with the
banners of utterance, the tribes of the religions were put to flight,
except those who wished to drink from the river of life in a paradise
brought into being by the All-Glorious Himself.

\subsection{75}\label{section-74}

\begin{otherlanguage}{arabic}

قد حکم اللّه بالطّهارة علی مآء النّطفة رحمةً من عنده علی البریة اشکروه
بالرّوح و الرّیحان و لا تتّبعوا من کان عن مطلع القرب بعیدا قوموا علی خدمة
الأمر فی کلّ الأحوال انّه یؤیّدکم بسلطان کان علی العالمین محیطا تمسّکوا بحبل
اللّطافة علی شأن لا یری من ثیابکم اآار الأوساخ هذا ما حکم به من کان الطف
من کلّ لطیف و الّذی له عذر لا بأس علیه انّه لهو الغفور الرّحیم طهّروا کلّ
مکروه بالمآء الّذی لم یتغیر بالثلٰث ایّاکم ان تستعملوا المآء الّذی تغیر
بالهوآء او بشئ آخر کونوا عنصر اللّطافة بین البریّة هذا ما اراد لکم مولاکم
العزیز الحکیم

\end{otherlanguage}

Qad ḥakama Allāh bi-al-ṭahāra ʿalā māʾ al-nuṭfa raḥmatan min ʿindihi
ʿalā al-bariyya ushkurūhu bi-al-rūḥ wa-al-rayḥān wa-lā tattabiʿū man
kāna ʿan maṭlaʿ al-qurb baʿīdan. Qūmū ʿalā khidmat al-amr fī kull
al-aḥwāl innahu yuʾayyidukum bi-sulṭān kāna ʿalā al-ʿālamīn muḥīṭan.
Tamassakū bi-ḥabl al-laṭāfa ʿalā shaʾn lā yurā min thiyābikum
{[}corrupted text{]} al-awsākh hādhā mā ḥakama bihi man kāna alṭaf min
kull laṭīf wa-alladhī lahu ʿudhr lā baʾs ʿalayhi innahu la-huwa
al-ghafūr al-raḥīm. Ṭahhirū kull makrūh bi-al-māʾ alladhī lam
yataghayyar bi-al-thalāth iyyākum an tastaʿmilū al-māʾ alladhī
taghayyara bi-al-hawāʾ aw bi-shayʾ ākhar kūnū ʿunṣur al-laṭāfa bayna
al-bariyya hādhā mā arāda lakum mawlākum al-ʿazīz al-ḥakīm.

God has decreed that semen is pure, as a mercy from Him to His
creatures. Thank Him with joy and gladness, and do not follow those who
are far from the Dawning-Place of Nearness. Arise to serve the Cause
under all circumstances. He will support you with an authority
encompassing all the worlds. Hold fast to the cord of cleanliness, so
that no {[}corrupted original text{]} of dirt may be seen on your
clothes. This is what He who is more refined than all who possess
refinement has decreed. Whoever has an excuse incurs no blame. He is
indeed the Ever-Forgiving, the Merciful. Cleanse every unclean thing
with water that has not undergone any of the three changes. Beware of
using water that has been altered by the air or by anything else. Be the
embodiment of cleanliness among creation. This is what your Lord, the
Almighty, the All-Wise, desires for you.

\subsection{76}\label{section-75}

\begin{otherlanguage}{arabic}

و کذلک رفع اللّه حکم دون الطّهارة عن کلّ الأشیاء و عن مللٍ اخری موهبة من
اللّه انّه لهو الغفور الکریم قد انغمست الأشیآء فی بحر الطّهارة فی اوّل
الرّضوان اذ تجلّینا علی من فی الامکان باسمآئنا الحسنی و صفاتنا العلیا هذا
من فضلی الّذی احاط العالمین لتعاشروا مع الأدیان و تبلّغوا امر ربّکم الرّحمن
هذا لاکلیل الاعمال لو انتم من العارفین

\end{otherlanguage}

Wa-kadhālika rafaʿa Allāh ḥukm dūn al-ṭahāra ʿan kull al-ashyāʾ wa-ʿan
milal ukhrā mawhiba min Allāh innahu la-huwa al-ghafūr al-karīm. Qad
inghamasat al-ashyāʾ fī baḥr al-ṭahāra fī awwal al-riḍwān idh tajallaynā
ʿalā man fī al-imkān bi-asmāʾinā al-ḥusnā wa-ṣifātinā al-ʿulyā hādhā min
faḍlī alladhī aḥāṭa al-ʿālamīn li-tuʿāshirū maʿa al-adyān wa-tuballighū
amr rabbikum al-raḥmān hādhā la-iklīl al-aʿmāl law antum min al-ʿārifīn.

Likewise, God has removed the designation of impurity from all things
and from other religious communities, as a gift from God. He is indeed
the Ever-Forgiving, the Generous. All things were immersed in the sea of
purity on the first day of Ridvan, when We revealed Our splendor to all
who exist through Our most beautiful names and Our most exalted
attributes. This is from My grace, which encompasses all the worlds, so
that you may associate with the followers of other religions and convey
the Cause of your Lord, the All-Merciful. This is indeed the crown of
deeds, if you are among those who understand.

\subsection{77}\label{section-76}

\begin{otherlanguage}{persian}

و حکم باللّطافة الکبری و تغسیل ما تغبّر من الغبار و کیف الأوساخ المنجمدة و
دونها اتّقوا اللّه و کونوا من المطهّرین و الّذی یری فی کسائه وسخ انّه لا یصعد
دعائه الی اللّه و یجتنب عنه ملأ عالون استعملوا مآء الورد ثمّ العطر الخالص
هذا ما احبّه اللّه من الأوّل الّذی لا اوّل له لیتضوّع منکم ما اراد ربّکم العزیز
الحکیم

\end{otherlanguage}

Wa-ḥakama bi-al-laṭāfa al-kubrā wa-taghsīl mā taghabbara min al-ghubār
wa-kayfa al-awsākh al-munjamida wa-dūnihā ittaqū Allāh wa-kūnū min
al-muṭahhirīn. Wa-alladhī yarā fī kisāʾihi wasakh innahu lā yaṣʿadu
duʿāʾuhu ilā Allāh wa-yajtanibu ʿanhu malaʾ ʿālūn. Istaʿmilū māʾ al-ward
thumma al-ʿiṭr al-khāliṣ hādhā mā aḥabbahu Allāh min al-awwal alladhī lā
awwala lahu li-yataḍawwaʿa minkum mā arāda rabbukum al-ʿazīz al-ḥakīm.

He has decreed the utmost cleanliness and the washing of whatever has
gathered dust---how much more, then, encrusted dirt and other
impurities. Fear God and be among those who keep themselves pure.
Whoever has visible dirt on his clothing, his prayer does not ascend to
God, and the exalted Concourse shuns him. Use rosewater and pure
perfume. This is what God has loved from the beginning that has no
beginning, so that the fragrance your Lord, the Almighty, the All-Wise,
desires may spread from you.

\subsection{78}\label{section-77}

\begin{otherlanguage}{arabic}

قد عفا اللّه عنکم ما نزّل فی البیان من محو الکتب و اذّناکم بان تقرؤوا من
العلوم ما ینفعکم لا ما ینتهی الی المجادلة فی الکلام هذا خیر لکم ان انتم
من العارفین

\end{otherlanguage}

Qad ʿafā Allāh ʿankum mā nuzzila fī al-bayān min maḥw al-kutub
wa-adhinnākum bi-an taqraʾū min al-ʿulūm mā yanfaʿukum lā mā yantahī ilā
al-mujādala fī al-kalām hādhā khayr lakum in antum min al-ʿārifīn.

God has released you from the command revealed in the Bayan to destroy
books. We have permitted you to study those branches of knowledge that
benefit you, not those that end in verbal disputation. This is better
for you, if you are among those who understand.

\subsection{79}\label{section-78}

\begin{otherlanguage}{arabic}

یا معشر الملوک قد اتی المالک و الملک للّه المهیمن القیّوم الاّ تعبدوا الاّ
اللّه و توجّهوا بقلوب نوراء الی وجه ربّکم مالک الأسمآء هذا امر لا یعادله ما
عندکم لو انتم تعرفون

\end{otherlanguage}

Yā maʿshar al-mulūk, qad atā al-Mālik wa-l-mulk li-llāh al-Muhaymin
al-Qayyūm. Allā taʿbudū illā Allāh wa-tawajjahū bi-qulūb nūrāʾ ilā wajh
rabbikum Mālik al-Asmāʾ. Hādhā amr lā yuʿādiluhu mā ʿindakum law antum
taʿrifūn.

O kings, the Sovereign has come, and sovereignty belongs to God, the
All-Protecting, the Self-Subsisting. Worship none but God, and turn with
radiant hearts toward the face of your Lord, the Lord of Names. This is
a Cause that nothing you possess can equal, if only you understood.

\subsection{80}\label{section-79}

\begin{otherlanguage}{arabic}

انّا نریکم تفرحون بما جمعتموه لغیرکم و تمنعون انفسکم عن العوالم الّتی لم
یحصها الاّ لوحی المحفوظ قد شغلتکم الأموال عن المأل هذا لا ینبغی لکم لو
انتم تعلمون طهّروا قلوبکم عن ذفر الدّنیا مسرعین الی ملکوت ربّکم فاطر الأرض
و السّمآء الّذی به ظهرت الزّلازل و ناحت القبائل الاّ من نبذ الوری و اخذ ما
امر به فی لوح مکنون

\end{otherlanguage}

Innā narākum tafraḥūn bi-mā jamaʿtumūhu li-ghayrikum wa-tamnaʿūn
anfusakum ʿan al-ʿawālim allatī lam yuḥṣihā illā lawḥī al-maḥfūẓ. Qad
shaghalatkum al-amwāl ʿan al-maʾāl. Hādhā lā yanbaghī lakum law antum
taʿlamūn. Ṭahhirū qulūbakum ʿan dhafar al-dunyā musriʿīn ilā malakūt
rabbikum Fāṭir al-arḍ wa-l-samāʾ alladhī bihi ẓaharat al-zalāzil
wa-nāḥat al-qabāʾil illā man nabadha al-warā wa-akhadha mā umira bihi fī
lawḥ maknūn.

We see you rejoicing in what you have amassed for others, while
depriving yourselves of worlds that none but My preserved Tablet has
numbered. Your wealth has distracted you from your final destination.
This does not befit you, if only you knew. Cleanse your hearts of the
stench of this world, hastening toward the kingdom of your Lord, the
Creator of earth and heaven, through whom earthquakes have appeared and
the peoples have lamented, except those who have cast aside all creation
and embraced what they were commanded in a hidden Tablet.

\subsection{81}\label{section-80}

\begin{otherlanguage}{arabic}

هذا یوم فیه فاز الکلیم بانوار القدیم و شرب زلال الوصال من هذا القدح الّذی
به سجّرت البحور قل تاللّه الحقّ انّ الطّور یطوف حول مطلع الظّهور و الرّوح ینادی
من الملکوت هلّموا و تعالوا یا ابنآء الغرور هذا یوم فیه سرع کوم اللّه شوقا
للقائه و صاح الصّهیون قد اتی الوعد و ظهر ما هو المکتوب فی الواح اللّه
المتعالی العزیز المحبوب

\end{otherlanguage}

Hādhā yawm fīhi fāza al-Kalīm bi-anwār al-Qadīm wa-shariba zulāl
al-wiṣāl min hādhā al-qadaḥ alladhī bihi sujjirat al-buḥūr. Qul ta-llāh
al-Ḥaqq inna al-Ṭūr yaṭūf ḥawla maṭlaʿ al-ẓuhūr wa-l-Rūḥ yunādī min
al-malakūt halummū wa-taʿālaw yā abnāʾ al-ghurūr. Hādhā yawm fīhi saraʿa
kawm Allāh shawqan li-liqāʾihi wa-ṣāḥa al-Ṣahyūn qad atā al-waʿd
wa-ẓahara mā huwa al-maktūb fī alwāḥ Allāh al-Mutaʿālī al-ʿAzīz
al-Maḥbūb.

This is the day when the One Who Conversed with God attained the lights
of the Ancient One and drank the pure water of reunion from this cup
through which the seas were set ablaze. Say: By God, the Truth! Sinai
circles around the Dayspring of Revelation, and the Spirit calls from
the kingdom: ``Come, draw near, you children of delusion!'' This is the
day when the hill of God has hastened in longing to meet Him, and Zion
has cried out: ``The promise has come, and what was written in the
Tablets of God, the Most Exalted, the Almighty, the Beloved, has
appeared.''

\subsection{82}\label{section-81}

\begin{otherlanguage}{arabic}

یا معشر الملوک قد نزّل النّاموس الأکبر فی المنظر الأنور و ظهر کل امر مستتر
من لدن مالک القدر الّذی به اتت السّاعة و انشقّ القمر و فصّل کلّ امر محتوم

\end{otherlanguage}

Yā maʿshar al-mulūk, qad nuzzila al-nāmūs al-akbar fī al-manẓar al-anwar
wa-ẓahara kull amr mustatir min ladun Mālik al-qadar alladhī bihi atat
al-sāʿa wa-nshaqqa al-qamar wa-fuṣṣila kull amr maḥtūm.

O kings, the Most Great Law has been sent down in the most radiant
Vision, and every hidden matter has been revealed by the Lord of
Destiny, through whom the Hour has come, the moon has split, and every
irrevocable decree has been set forth.

\subsection{83}\label{section-82}

\begin{otherlanguage}{arabic}

یا معشر الملوک انتم الممالیک قد ظهر المالک بأحسن الطّراز و یدعوکم الی
نفسه المهیمن القیّوم ایّاکم ان یمنعکم الغرور عن مشرق الظّهور او تحجبکم
الدّنیا عن فاطر السّمآء قوموا علی خدمة المقصود الّذی خلقکم بکلمة من عنده و
جعلکم مظاهر القدرة لما کان و ما یکون

\end{otherlanguage}

Yā maʿshar al-mulūk, antum al-mamālīk. Qad ẓahara al-Mālik bi-aḥsan
al-ṭirāz wa-yadʿūkum ilā nafsihi al-Muhaymin al-Qayyūm. Iyyākum an
yamnaʿakum al-ghurūr ʿan mashriq al-ẓuhūr aw taḥjubakum al-dunyā ʿan
Fāṭir al-samāʾ. Qūmū ʿalā khidmat al-Maqṣūd alladhī khalaqakum bi-kalima
min ʿindihi wa-jaʿalakum maẓāhir al-qudra li-mā kāna wa-mā yakūn.

O kings, you are but servants. The Sovereign has appeared in the finest
adornment and calls you to Himself, the All-Protecting, the
Self-Subsisting. Beware lest pride keep you from the Dawning-Place of
Revelation, or the world veil you from the Creator of heaven. Arise to
serve the Desired One, who created you through a word from Him and made
you manifestations of power over all that has been and all that will be.

\subsection{84}\label{section-83}

\begin{otherlanguage}{arabic}

تاللّه لا نرید ان نتصرّف فی ممالککم بل جئنا لتصرّف القلوب انّها لمنظر البهآء
یشهد بذلک ملکوت الأسمآء لو انتم تفقهون و الّذی اتّبع مولاه انّه اعرض عن
الدّنیا کلّها و کیف هذا المقام المحمود دعوا البیوت ثمّ اقبلوا الی الملکوت
هذا ما ینفعکم فی الآخرة و الأولی یشهد بذلک مالک الجبروت لو انتم تعلمون

\end{otherlanguage}

Ta-llāh lā nurīd an nataṣarraf fī mamālikikum bal jiʾnā li-taṣarruf
al-qulūb. Innahā la-manẓar al-Bahāʾ yashhadu bi-dhālika malakūt al-asmāʾ
law antum tafqahūn. Wa-lladhī ittabaʿa mawlāhu innahu aʿraḍa ʿan
al-dunyā kullihā wa-kayfa hādhā al-maqām al-maḥmūd. Daʿū al-buyūt thumma
aqbilū ilā al-malakūt. Hādhā mā yanfaʿukum fī al-ākhira wa-l-ūlā
yashhadu bi-dhālika Mālik al-jabarūt law antum taʿlamūn.

By God! We do not wish to take possession of your kingdoms; rather, We
have come to take possession of hearts. They are indeed the place where
Bahá is revealed; the kingdom of names bears witness to this, if only
you understood. Whoever has followed his Lord has turned away from the
whole world---how much more so in this praiseworthy station! Leave your
houses and turn toward the kingdom. This is what will benefit you in the
next life and in this one. The Lord of the realm of power bears witness
to this, if only you knew.

\subsection{85}\label{section-84}

\begin{otherlanguage}{arabic}

طوبی لملک قام علی نصرة امری فی مملکتی و انقطع عن سوائی انّه من اصحاب
السّفینة الحمرآء التی جعلها اللّه لأهل البهآء ینبغی لکلٍ ان یعزّزوه و یوقّروه
و ینصروه لیفتح المدن بمفاتیح اسمی المهیمن علی من فی ممالک الغیب و الشّهود
انّه بمنزلة البَصَر للبشر و الغرّة الغرآء لجبین الانشآء و رأس الکرم لجسد
العالم انصروه یا اهل البهآء بالأموال و النّفوس

\end{otherlanguage}

Ṭūbā li-malik qāma ʿalā nuṣrat amrī fī mamlakatī wa-nqaṭaʿa ʿan siwāʾī.
Innahu min aṣḥāb al-safīna al-ḥamrāʾ allatī jaʿalahā Allāh li-ahl
al-Bahāʾ. Yanbaghī li-kull an yuʿazzizūhu wa-yuwaqqirūhu wa-yanṣurūhu
li-yaftaḥ al-mudun bi-mafātīḥ ismī al-Muhaymin ʿalā man fī mamālik
al-ghayb wa-l-shuhūd. Innahu bi-manzilat al-baṣar li-l-bashar
wa-l-ghurra al-gharrāʾ li-jabīn al-inshāʾ wa-raʾs al-karam li-jasad
al-ʿālam. Unṣurūhu yā ahl al-Bahāʾ bi-l-amwāl wa-l-nufūs.

Blessed is the king who arises to aid My Cause in My kingdom and
detaches himself from all but Me. He is among the companions of the
Crimson Ark that God has ordained for the people of Bahá. Everyone
should honor, revere, and support him, so that he may open the cities
with the keys of My name, the All-Protecting over all who dwell in the
kingdoms of the unseen and the seen. He is as sight to humanity, a
radiant adornment on the brow of creation, and the head of generosity to
the body of the world. Support him, O people of Bahá, with your wealth
and your lives.

\subsection{86}\label{section-85}

\begin{otherlanguage}{arabic}

یا ملک النّمسة کان مطلع نور الأحدیّة فی سجن عکّآء اذ قصدت المسجد الأقصی
مررت و ما سئلت عنه بعد اذ رفع به کلّ بیت و فتح کلّ باب منیف قد جعلناه مقبل
العالم لذکری و انت نبذت المذکور اذ ظهر بملکوت اللّه ربّک و ربّ العالمین کنّا
معک فی کلّ الأحوال و وجدناک متمسّکا بالفرع غافلا عن الأصل انّ ربّک علی ما
اقول شهید قد اخذتنا الأحزان بما رأیناک تدور لاسمنا و لا تعرفنا امام وجهک
افتح البصر لتنظر هذا المنظر الکریم و تعرف من تدعوه فی اللّیالی و الأیّام و
تری النّور المشرق من هذا الأفق اللّمیع

\end{otherlanguage}

Yā malik al-Nimsa, kāna maṭlaʿ nūr al-aḥadiyya fī sijn ʿAkkā idh qaṣadta
al-masjid al-aqṣā. Mararta wa-mā saʾalta ʿanhu baʿda idh rufiʿa bihi
kull bayt wa-futiḥa kull bāb munīf. Qad jaʿalnāhu muqbil al-ʿālam
li-dhikrī wa-anta nabadhta al-madhkūr idh ẓahara bi-malakūt Allāh
rabbika wa-rabb al-ʿālamīn. Kunnā maʿaka fī kull al-aḥwāl wa-wajadnāka
mutamassikan bi-l-farʿ ghāfilan ʿan al-aṣl. Inna rabbaka ʿalā mā aqūl
shahīd. Qad akhadhatnā al-aḥzān bi-mā raʾaynāka tadūr li-sminā wa-lā
taʿrifunā amāma wajhika. Iftaḥ al-baṣar li-tanẓur hādhā al-manẓar
al-karīm wa-taʿrif man tadʿūhu fī al-layālī wa-l-ayyām wa-tarā al-nūr
al-mushriq min hādhā al-ufuq al-lamīʿ.

O king of Austria, the Dayspring of the light of oneness was in the
prison of Acre when you set out for the Farthest Mosque. You passed by
without asking about Him, through whom every house has been exalted and
every lofty gate opened. We made it the place toward which the world
turns to remember Me, yet you cast aside the One remembered when He
appeared with the kingdom of God, your Lord and the Lord of all worlds.
We were with you in every circumstance and found you clinging to the
branch, heedless of the root. Your Lord bears witness to what I say.
Sorrow seized Us when We saw you journeying for the sake of Our name,
yet failing to recognize Us before your very eyes. Open your eyes to
behold this noble Vision, recognize the One you invoke by night and day,
and see the light shining from this radiant horizon.

\subsection{87}\label{section-86}

\begin{otherlanguage}{arabic}

قل یا ملک برلین اسمع النّدآء من هذا الهیکل المبین انّه لا اله الاّ انا
الباقی الفرد القدیم ایّاک ان یمنعک الغرور عن مطلع الظّهور او یحجبک الهوی
عن مالک العرش و الثّری کذلک ینصحک القلم الأعلی انّه لهو الفضّال الکریم اذکر
من کان اعظم منک شأنا و اکبر منک مقاماً این هو و ما عنده انتبه و لا تکن من
الرّاقدین انّه نبذ لوح اللّه ورائه اذ اخبرناه بما ورد علینا من جنود
الظّالمین لذا اخذته الذّلّة من کلّ الجهات الی ان رجع الی التّراب خسران عظیم
یا ملک تفکّر فیه و فی امثالک الّذین سخّروا البلاد و حکموا علی العباد قد
انزلهم الرّحمن من القصور الی القبور اعتبر و کن من المتذکّرین

\end{otherlanguage}

Qul yā malik Birlīn, ismaʿ al-nidāʾ min hādhā al-haykal al-mubīn. Innahu
lā ilāha illā anā al-Bāqī al-Fard al-Qadīm. Iyyāka an yamnaʿaka
al-ghurūr ʿan maṭlaʿ al-ẓuhūr aw yaḥjubaka al-hawā ʿan Mālik al-ʿarsh
wa-l-tharā. Kadhālika yanṣaḥuka al-Qalam al-Aʿlā innahu la-huwa
al-Faḍḍāl al-Karīm. Udhkur man kāna aʿẓam minka shaʾnan wa-akbar minka
maqāman. Ayna huwa wa-mā ʿindahu? Intabih wa-lā takun min al-rāqidīn.
Innahu nabadha lawḥ Allāh warāʾahu idh akhbarnāhu bi-mā warada ʿalaynā
min junūd al-ẓālimīn. Li-dhā akhadhat-hu al-dhilla min kull al-jihāt ilā
an rajaʿa ilā al-turāb. Khusrān ʿaẓīm. Yā malik tafakkar fīhi wa-fī
amthālika alladhīna sakhkharū al-bilād wa-ḥakamū ʿalā al-ʿibād. Qad
anzalahum al-Raḥmān min al-quṣūr ilā al-qubūr. Iʿtabir wa-kun min
al-mutadhakkirīn.

Say: O king of Berlin, hear the call from this manifest Temple: ``There
is no god but Me, the Everlasting, the One, the Ancient.'' Beware lest
pride keep you from the Dayspring of Revelation, or desire veil you from
the Lord of the throne and the earth. Thus the Most Exalted Pen counsels
you. He is indeed the Most Bountiful, the Generous. Remember the one who
was greater than you in power and higher than you in station. Where is
he, and where is all he possessed? Awaken, and do not be among those who
sleep. He cast the Tablet of God behind him when We told him what had
befallen Us at the hands of the armies of the oppressors. Therefore
humiliation overtook him from every side, until he returned to dust.
What a tremendous loss! O king, reflect on him and on those like you who
subdued lands and ruled over their peoples. The All-Merciful brought
them down from palaces to graves. Take heed, and be among those who
remember.

\subsection{88}\label{section-87}

\begin{otherlanguage}{arabic}

انّا ما اردنا منکم شیئا انّما ننصحکم لوجه اللّه و نصبر کما صبرنا بما
وردعلینا منکم یا معشر السّلاطین

\end{otherlanguage}

Innā mā aradnā minkum shayʾan innamā nanṣaḥukum li-wajh Allāh wa-naṣbir
kamā ṣabarnā bi-mā warada ʿalaynā minkum yā maʿshar al-salāṭīn.

We have desired nothing from you. We counsel you only for the sake of
God, and We endure patiently, as We have endured what you have inflicted
on Us, O sovereigns.

\subsection{89}\label{section-88}

\begin{otherlanguage}{arabic}

یا ملوک امریقا و رؤسآء الجمهور فیها اسمعوا ما تغنّ به الورقآء علی غصن
البقآء انّه لا اله الاّ انا الباقی الغفور الکریم زیّنوا هیکل الملک بطراز
العدل و التّقی و رأسه باکلیل ذکر ربّکم فاطر السّمآء کذلک یأمرکم مطلع
الأسمآء من لدن علیمٍ حکیمٍ قد ظهر الموعود فی هذا المقام المحمود الّذی به
ابتسم ثغر الوجود من الغیب و الشّهود اغتنموا یوم اللّه انّ لقائه خیر لکم عمّا
تطلع الشّمس علیها ان انتم من العارفین یا معشر الأمرآء اسمعوا ما ارتفع من
مطلع الکبریآء انّه لا اله الاّ انا النّاطق العلیم اجبروا الکسیر بایادی
العدل و کسّروا الصّحیح الظّالم بسیاط اوامر ربّکم الأمر الحکیم

\end{otherlanguage}

Yā mulūk Amrīqā wa-ruʾasāʾ al-jumhūr fīhā, ismaʿū mā taghann bihi
al-warqāʾ ʿalā ghuṣn al-baqāʾ. Innahu lā ilāha illā anā al-Bāqī
al-Ghafūr al-Karīm. Zayyinū haykal al-mulk bi-ṭirāz al-ʿadl wa-l-tuqā
wa-raʾsahu bi-iklīl dhikr rabbikum Fāṭir al-samāʾ. Kadhālika yaʾmurukum
maṭlaʿ al-asmāʾ min ladun ʿAlīm Ḥakīm. Qad ẓahara al-Mawʿūd fī hādhā
al-maqām al-maḥmūd alladhī bihi ibtasama thaghr al-wujūd min al-ghayb
wa-l-shuhūd. Ightanimū yawm Allāh inna liqāʾahu khayr lakum ʿammā taṭluʿ
al-shams ʿalayhā in antum min al-ʿārifīn. Yā maʿshar al-umarāʾ, ismaʿū
mā irtafaʿa min maṭlaʿ al-kibriyāʾ. Innahu lā ilāha illā anā al-Nāṭiq
al-ʿAlīm. Ijburū al-kasīr bi-ayādī al-ʿadl wa-kassirū al-ṣaḥīḥ al-ẓālim
bi-siyāṭ awāmir rabbikum al-Āmir al-Ḥakīm.

O kings of America and presidents of its republics, listen to what the
Dove sings upon the branch of eternity: ``There is no god but Me, the
Everlasting, the All-Forgiving, the Generous.'' Adorn the body of
sovereignty with justice and piety, and its head with the crown of
remembrance of your Lord, the Creator of heaven. Thus the Dayspring of
Names commands you on behalf of One who is All-Knowing, All-Wise. The
Promised One has appeared in this praiseworthy station, through which
all existence, unseen and seen, has smiled. Seize the opportunity of the
Day of God. Meeting Him is better for you than all that the sun shines
upon, if you are among those who understand. O rulers, listen to what
has risen from the Dayspring of Grandeur: ``There is no god but Me, the
Speaker, the All-Knowing.'' Bind up the broken with the hands of
justice, and break the oppressor who is unbroken with the scourges of
the commands of your Lord, the Ordainer, the All-Wise.

\subsection{90}\label{section-89}

\begin{otherlanguage}{persian}

یا معشر الرّوم نسمع بینکم صوت البوم ء اخذکم سکر الهوی ام کنتم من الغافلین
یا ایتها النّقطة الواقعة فی شاطئ البحرین قد استقرّ علیک کرسیّ الظّلم و
اشتعلت فیک نار البغضآء علی شأن ناح بها الملأ الأعلی و الّذین یطوفون حول
کرسی رفیع نری فیک الجاهل یحکم علی العاقل و الظّلام یفتخر علی النّور و انّک
فی غرور مبین اغرّتک زینتک الظّاهرة سوف تفنی و ربّ البریّة و تنوح البنات و
الأرامل و ما فیک من القبائل کذلک ینبئک العلیم الخبیر

\end{otherlanguage}

Yā maʿshar al-Rūm, nasmaʿ baynakum ṣawt al-būm. A-akhadhakum sukr
al-hawā am kuntum min al-ghāfilīn? Yā ayyatuhā al-nuqṭa al-wāqiʿa fī
shāṭiʾ al-baḥrayn, qad istaqarra ʿalayki kursī al-ẓulm wa-shtaʿalat fīki
nār al-baghḍāʾ ʿalā shaʾn nāḥa bihā al-malaʾ al-aʿlā wa-lladhīna yaṭūfūn
ḥawla kursī rafīʿ. Narā fīki al-jāhil yaḥkum ʿalā al-ʿāqil wa-l-ẓalām
yaftakhir ʿalā al-nūr wa-innaki fī ghurūr mubīn. A-gharratki zīnatuki
al-ẓāhira? Sawfa tafnā wa-rabb al-bariyya wa-tanūḥ al-banāt wa-l-arāmil
wa-mā fīki min al-qabāʾil. Kadhālika yunabbiʾuki al-ʿAlīm al-Khabīr.

O people of Rūm, We hear the owl's cry among you. Has the intoxication
of desire seized you, or are you among the heedless? O spot situated on
the shore of the two seas, the throne of oppression has been established
upon you, and the fire of hatred has blazed within you so fiercely that
the exalted Concourse and those who circle a lofty throne have lamented.
We see the ignorant ruling over the wise within you, and darkness
boasting over light. You are indeed in manifest delusion. Has your
outward splendor deceived you? By the Lord of creation, it will perish,
and the daughters, the widows, and all the peoples within you will
lament. Thus the All-Knowing, the All-Aware informs you.

\subsection{91}\label{section-90}

\begin{otherlanguage}{arabic}

یا شواطی نهر الرّین قد رأیناک مغطّاة بالدّماء بما سلّ علیک سیوف الجزآء و لک
مرّة اخری و نسمع حنین البرلین ولو انّها الیوم علی عزّ مبین

\end{otherlanguage}

Yā shawāṭiʾ nahr al-Rayn, qad raʾaynāki mughaṭṭāt bi-l-dimāʾ bi-mā sulla
ʿalayki suyūf al-jazāʾ wa-laki marra ukhrā wa-nasmaʿ ḥanīn al-Birlīn
wa-law annahā al-yawm ʿalā ʿizz mubīn.

O banks of the Rhine, We have seen you covered in blood, for the swords
of retribution have been drawn against you; and you will have another
turn. We hear the lament of Berlin, though today she enjoys manifest
glory.

\subsection{92}\label{section-91}

\begin{otherlanguage}{arabic}

یا ارض الطّآء لا تحزنی من شئ قد جعلک اللّه مطلع فرح العالمین لو یشآء یبارک
سریرک بالّذی یحکم بالعدل و یجمع اغنام اللّه الّتی تفرّقت من الذّئاب انّه یواجه
اهل البهآء بالفرح و الانبساط الا انّه من جوهر الخلق لدی الحقّ علیه بهآء
اللّه و بهآء من فی ملکوت الأمر فی کلّ حین

\end{otherlanguage}

Yā arḍ al-Ṭāʾ, lā taḥzanī min shayʾ. Qad jaʿalaki Allāh maṭlaʿ faraḥ
al-ʿālamīn. Law yashāʾ yubārik sarīraki bi-lladhī yaḥkum bi-l-ʿadl
wa-yajmaʿ aghnām Allāh allatī tafarraqat min al-dhiʾāb. Innahu yuwājih
ahl al-Bahāʾ bi-l-faraḥ wa-l-inbisāṭ. Alā innahu min jawhar al-khalq
ladā al-Ḥaqq. ʿAlayhi bahāʾ Allāh wa-bahāʾ man fī malakūt al-amr fī kull
ḥīn.

O land of Ṭā, grieve over nothing. God has made you the dayspring of joy
for all worlds. If He wills, He will bless your throne with one who
rules with justice and gathers the sheep of God that have been scattered
by the wolves. He will meet the people of Bahá with joy and warmth.
Indeed, in the sight of the True One, he is among the choicest of
creation. Upon him be the glory of God and the glory of all who dwell in
the kingdom of His Cause at every moment.

\subsection{93}\label{section-92}

\begin{otherlanguage}{arabic}

افرحی بما جعلک اللّه افق النّور بما ولد فیک مطلع الظّهور و سمّیت بهذا الاسم
الّذی به لاح نیّر الفضل و اشرقت السّموات و الأرضون

\end{otherlanguage}

Ifraḥī bi-mā jaʿalaki Allāh ufuq al-nūr bi-mā wulida fīki maṭlaʿ
al-ẓuhūr wa-summīti bi-hādhā al-ism alladhī bihi lāḥa nayyir al-faḍl
wa-ashraqat al-samāwāt wa-l-araḍūn.

Rejoice that God has made you the horizon of light, for the Dayspring of
Revelation was born within you, and you were called by this name through
which the sun of grace has shone and the heavens and the earth have been
illumined.

\subsection{94}\label{section-93}

\begin{otherlanguage}{arabic}

سوف تنقلب فیک الأمور و یحکم علیک جمهور النّاس انّ ربّک لهو العلیم المحیط
اطمئنّی بفضل ربّک انّه لا تنقطع عنک لحظات الألطاف سوف یأخذک الاطمینان بعد
الاضطراب کذلک قضی الأمر فی کتاب بدیع

\end{otherlanguage}

Sawfa tanqalib fīki al-umūr wa-yaḥkum ʿalayki jumhūr al-nās. Inna
rabbaki la-huwa al-ʿAlīm al-Muḥīṭ. Iṭmaʾinnī bi-faḍl rabbiki innahu lā
tanqaṭiʿ ʿanki laḥaẓāt al-alṭāf. Sawfa yaʾkhudhuki al-iṭmiʾnān baʿda
al-iḍṭirāb. Kadhālika quḍiya al-amr fī kitāb badīʿ.

Affairs within you will change, and the people will rule over you. Your
Lord is indeed the All-Knowing, the All-Encompassing. Be assured of your
Lord's grace: His glances of favor toward you will not cease.
Tranquillity will embrace you after turmoil. Thus has the matter been
decreed in a wondrous Book.

\subsection{95}\label{section-94}

\begin{otherlanguage}{arabic}

یا ارض الخآء نسمع فیک صوت الرّجال فی ذکر ربّک الغنی المتّعال طوبی لیوم فیه
تنصب رایات الأسمآء فی ملکوت الانشآء باسمی الأبهی یومئذ یفرح المخلصون
بنصر اللّه و ینوح المشرکون

\end{otherlanguage}

Yā arḍ al-Khāʾ, nasmaʿ fīki ṣawt al-rijāl fī dhikr rabbiki al-Ghanī
al-Mutaʿāl. Ṭūbā li-yawm fīhi tunṣab rāyāt al-asmāʾ fī malakūt al-inshāʾ
bi-smī al-Abhā. Yawmaʾidh yafraḥ al-mukhliṣūn bi-naṣr Allāh wa-yanūḥ
al-mushrikūn.

O land of Khā, We hear within you the voices of men remembering your
Lord, the Self-Sufficient, the Most Exalted. Blessed is the day when the
banners of the names will be raised in the kingdom of creation through
My name, the All-Glorious. On that day the sincere will rejoice in God's
victory, and those who associate partners with Him will lament.

\subsection{96}\label{section-95}

\begin{otherlanguage}{arabic}

لیس لأحد ان یعترض علی الّذین یحکمون علی العباد دعوا لهم ما عندهم و توجّهوا
الی القلوب

\end{otherlanguage}

Laysa li-aḥad an yaʿtariḍ ʿalā alladhīna yaḥkumūn ʿalā al-ʿibād. Daʿū
lahum mā ʿindahum wa-tawajjahū ilā al-qulūb.

No one may oppose those who rule over the people. Leave to them what
they possess, and turn your attention to hearts.

\subsection{97}\label{section-96}

\begin{otherlanguage}{arabic}

یا بحر الأعظم رشّ علی الأمم ما امرت به من لدن مالک القدم و زین هیاکل
الأنام بطراز الأحکام الّتی بها تفرح القلوب و تقرّ العیون

\end{otherlanguage}

Yā baḥr al-aʿẓam, rushsha ʿalā al-umam mā umirta bihi min ladun Mālik
al-qidam wa-zayyin hayākil al-anām bi-ṭirāz al-aḥkām allatī bihā tafraḥ
al-qulūb wa-taqarr al-ʿuyūn.

O Most Great Ocean, sprinkle upon the nations what you have been
commanded by the Lord of Eternity, and adorn the bodies of humanity with
the laws through which hearts rejoice and eyes find delight.

\subsection{98}\label{section-97}

\begin{otherlanguage}{arabic}

و الّذی تملّک مأة مثقال من الذّهب فتسعة عشر مثقالا للّه فاطر الأرض و السّمآء
ایّاکم یا قوم ان تمنعوا انفسکم عن هذا الفضل العظیم قد امرناکم بهذا بعد اذ
کنّا غنیا عنکم و عن کلّ من فی السّموات و الأرضین انّ فی ذلک لحکم و مصالح لم
یحط بها علم احد الاّ اللّه العالم الخبیر قل بذلک اراد تطهیر اموالکم و
تقرّبکم الی مقامات لا یدرکها الاّ من شاء اللّه انّه لهو الفضّال العزیز الکریم
یا قوم لا تخونوا فی حقوق اللّه و لا تصرّفوا فیها الاّ بعد اذنه کذلک قضی
الأمر فی الألواح و فی هذا اللّوح المنیع من خان اللّه یخان بالعدل و الّذی
عمل بما امر ینزل علیه البرکة من سمآء عطآء ربّه الفیاض المعطی الباذل
القدیم انّه اراد لکم ما لا تعرفونه الیوم سوف یعرفه القوم اذا طارت الأرواح
و طویت زرابی الافراح کذلک یذکرکم من عنده لوح حفیظ

\end{otherlanguage}

Wa-alladhī tamallaka miʾata mithqālin mina al-dhahabi fa-tisʿata ʿashara
mithqālan li-llāhi fāṭiri al-arḍi wa-al-samāʾi. Iyyākum yā qawmi an
tamnaʿū anfusakum ʿan hādhā al-faḍli al-ʿaẓīmi. Qad amarnākum bi-hādhā
baʿda idh kunnā ghaniyyan ʿankum wa-ʿan kulli man fī al-samāwāti
wa-al-araḍīna. Inna fī dhālika la-ḥikaman wa-maṣāliḥa lam yuḥiṭ bihā
ʿilmu aḥadin illā Allāhu al-ʿālimu al-khabīru. Qul bi-dhālika arāda
taṭhīra amwālikum wa-taqarrubakum ilā maqāmātin lā yudrikuhā illā man
shāʾa Allāhu. Innahu la-huwa al-faḍḍālu al-ʿazīzu al-karīmu. Yā qawmi lā
takhūnū fī ḥuqūqi Allāhi wa-lā taṣarrafū fīhā illā baʿda idhnihi.
Kadhālika quḍiya al-amru fī al-alwāḥi wa-fī hādhā al-lawḥi al-manīʿi.
Man khāna Allāha yukhānu bi-al-ʿadli wa-alladhī ʿamila bi-mā umira
yanzilu ʿalayhi al-barakatu min samāʾi ʿaṭāʾi rabbihi al-fayyāḍi
al-muʿṭī al-bādhili al-qadīmi. Innahu arāda lakum mā lā taʿrifūnahu
al-yawma. Sawfa yaʿrifuhu al-qawmu idhā ṭārati al-arwāḥu wa-ṭuwiyat
zarābiyyu al-afrāḥi. Kadhālika yudhakkirukum man ʿindahu lawḥun ḥafīẓun.

Whoever possesses a hundred mithqals of gold, nineteen mithqals of it
belong to God, the Creator of earth and heaven. Take care, O people, not
to deprive yourselves of this great bounty. We have commanded you to
observe this, although We have always been independent of you and of all
who are in the heavens and on earth. In this lie wisdom and benefits
beyond the knowledge of anyone except God, the All-Knowing, the
All-Informed. Say: By this He has intended to purify your possessions
and bring you near to stations that none can comprehend except those
whom God wills. He is indeed the Bountiful, the Almighty, the Generous.
O people, do not act dishonestly concerning the Right of God, nor
dispose of it without His permission. Thus has the command been decreed
in the Tablets and in this exalted Tablet. Whoever betrays God will, in
justice, be betrayed; and whoever acts as commanded will receive
blessings from the heaven of the bounty of his Lord, the All-Bountiful,
the Bestower, the Munificent, the Ancient. He has intended for you what
you do not understand today. The people will understand it when their
spirits take flight and the carpets of joy are rolled up. Thus does He
who possesses a preserving Tablet admonish you.

\subsection{99}\label{section-98}

\begin{otherlanguage}{arabic}

قد حضرت لدی العرش عرائض شتّی من الّذین امنوا و سئلوا فیها اللّه ربّ ما یری و
ما لا یری ربّ العالمین لذا نزّلنا اللّوح و زینّاه بطراز الأمر لعلّ النّاس
باحکام ربّهم یعملون و کذلک سئلنا من قبل فی سنین متوالیات و امسکنا القلم
حکمة من لدنّا الی ان حضرت کتب من انفسٍ معدودات فی تلک الأیّام لذا اجبناهم
بالحقّ بما تحیی به القلوب

\end{otherlanguage}

Qad ḥaḍarat ladā al-ʿarshi ʿarāʾiḍu shattā mina alladhīna āmanū
wa-saʾalū fīhā Allāha rabba mā yurā wa-mā lā yurā rabba al-ʿālamīna.
Lidhā nazzalnā al-lawḥa wa-zayyannāhu bi-ṭirāzi al-amri laʿalla al-nāsa
bi-aḥkāmi rabbihim yaʿmalūna. Wa-kadhālika suʾilnā min qablu fī sinīna
mutawāliyātin wa-amsaknā al-qalama ḥikmatan min ladunnā ilā an ḥaḍarat
kutubun min anfusin maʿdūdātin fī tilka al-ayyāmi. Lidhā ajabnāhum
bi-al-ḥaqqi bi-mā taḥyā bihi al-qulūbu.

Various petitions have reached the Throne from those who have believed,
in which they have addressed questions to God, the Lord of the seen and
the unseen, the Lord of all worlds. We have therefore sent down this
Tablet and adorned it with the ornament of command, so that people may
act according to the laws of their Lord. Likewise, We were questioned in
previous years, one after another, but in Our wisdom We withheld the Pen
until letters from a few individuals arrived in these days. We have
therefore answered them in truth with that which brings hearts to life.

\subsection{100}\label{section-99}

\begin{otherlanguage}{arabic}

قل یا معشر العلمآء لا تزنوا کتاب اللّه بما عندکم من القواعد و العلوم انّه
لقسطاس الحقّ بین الخلق قد یوزن ما عند الأمم بهذا القسطاس الأعظم و انّه
بنفسه لو انتم تعلمون

\end{otherlanguage}

Qul yā maʿshara al-ʿulamāʾi lā tazinū kitāba Allāhi bi-mā ʿindakum mina
al-qawāʿidi wa-al-ʿulūmi. Innahu la-qisṭāsu al-ḥaqqi bayna al-khalqi.
Qad yūzanu mā ʿinda al-umami bi-hādhā al-qisṭāsi al-aʿẓami wa-innahu
bi-nafsihi law antum taʿlamūna.

Say: O assembly of scholars, do not weigh the Book of God by the rules
and sciences you possess. It is the true balance among creation.
Whatever the nations possess is weighed in this Most Great Balance,
while it stands by itself, if only you knew.

\subsection{101}\label{section-100}

\begin{otherlanguage}{arabic}

تبکی علیکم عین عنایتی لانّکم ما عرفتم الّذی دعوتموه فی العشی و الاشراق و
فی کلّ اصیلٍ و بکور توجّهوا یا قوم بوجوه بیضآء و قلوب نورآء الی البقعة
المبارکة الحمرآء الّتی فیها تنادی سدرة المنتهی انّه لا اله الاّ انا المهیمن
القیّوم

\end{otherlanguage}

Tabkī ʿalaykum ʿaynu ʿināyatī li-annakum mā ʿaraftumu alladhī
daʿawtumūhu fī al-ʿashiyyi wa-al-ishrāqi wa-fī kulli aṣīlin wa-bukūrin.
Tawajjahū yā qawmi bi-wujūhin bayḍāʾa wa-qulūbin nawrāʾa ilā al-buqʿati
al-mubārakati al-ḥamrāʾi allatī fīhā tunādī sidratu al-muntahā innahu lā
ilāha illā anā al-muhayminu al-qayyūmu.

The eye of My loving care weeps over you, because you have not
recognized the One you have called upon at evening and sunrise, and at
every dusk and dawn. Turn, O people, with radiant faces and luminous
hearts toward the blessed crimson Spot, where the Lote Tree of the
Utmost Boundary calls out: ``There is no God but Me, the All-Protecting,
the Self-Subsisting.''

\subsection{102}\label{section-101}

\begin{otherlanguage}{arabic}

یا معشر العلمآء هل یقدر احد منکم ان یستنّ معی فی میدان المکاشفة و العرفان
او یجول فی مضمار الحکمة و التّبیان لا و ربّی الرّحمن کلّ من علیها فان و هذا
وجه ربّکم العزیز المحبوب

\end{otherlanguage}

Yā maʿshara al-ʿulamāʾi hal yaqdiru aḥadun minkum an yastanna maʿī fī
maydāni al-mukāshafati wa-al-ʿirfāni aw yajūla fī miḍmāri al-ḥikmati
wa-al-tibyāni. Lā wa-rabbī al-raḥmāni. Kullu man ʿalayhā fānin wa-hādhā
wajhu rabbikumu al-ʿazīzi al-maḥbūbi.

O assembly of scholars, can any of you race with Me in the field of
unveiling and recognition, or range over the arena of wisdom and
exposition? No, by My Lord, the All-Merciful! Everyone on earth passes
away, and this is the face of your Lord, the Almighty, the Beloved.

\subsection{103}\label{section-102}

\begin{otherlanguage}{arabic}

یا قوم انّا قدّرنا العلوم لعرفان المعلوم و انتم احتجبتم بها عن مشرقها الّذی
به ظهر کلّ امر مکنون لو عرفتم الأفق الّذی منه اشرقت شمس الکلام لنبذتم
الأنام و ما عندهم و اقبلتم الی المقام المحمود

\end{otherlanguage}

Yā qawmi innā qaddarnā al-ʿulūma li-ʿirfāni al-maʿlūmi wa-antum
iḥtajabtum bihā ʿan mashriqihā alladhī bihi ẓahara kullu amrin maknūnin.
Law ʿaraftumu al-ufuqa alladhī minhu ashraqat shamsu al-kalāmi
la-nabadhtumu al-anāma wa-mā ʿindahum wa-aqbaltum ilā al-maqāmi
al-maḥmūdi.

O people, We ordained the sciences for the recognition of the One who is
their object, yet you have allowed them to veil you from their dawning
place, through whom every hidden thing has been revealed. Had you
recognized the horizon from which the sun of utterance has risen, you
would have cast aside humanity and all it possesses and turned toward
the praiseworthy station.

\subsection{104}\label{section-103}

\begin{otherlanguage}{arabic}

قل هذه لسمآء فیها کنز امّ الکتاب لو انتم تعقلون هذا لهو الّذی به صاحت
الصّخرة و نادت السّدرة علی الطّور المرتفع علی الأرض المبارکة الملک للّه
الملک العزیز الودود

\end{otherlanguage}

Qul hādhihi la-samāʾun fīhā kanzu ummi al-kitābi law antum taʿqilūna.
Hādhā la-huwa alladhī bihi ṣāḥati al-ṣakhratu wa-nādati al-sidratu ʿalā
al-ṭūri al-murtafiʿi ʿalā al-arḍi al-mubārakati al-mulku li-llāhi
al-maliki al-ʿazīzi al-wadūdi.

Say: This is indeed a heaven in which lies the treasure of the Mother
Book, if only you understood. This is the One through whom the Rock
cried out and the Lote Tree proclaimed upon the Mount rising above the
blessed land: ``Sovereignty belongs to God, the King, the Almighty, the
All-Loving.''

\subsection{105}\label{section-104}

\begin{otherlanguage}{arabic}

انّا ما دخلنا المدارس و ما طالعنا المباحث اسمعوا ما یدعوکم به هذا الأمی
الی اللّه الأبدی انّه خیر لکم عمّا کنز فی الأرض لو انتم تفقهون

\end{otherlanguage}

Innā mā dakhalnā al-madārisa wa-mā ṭālaʿnā al-mabāḥitha. Ismaʿū mā
yadʿūkum bihi hādhā al-ummiyyu ilā Allāhi al-abadiyyi. Innahu khayrun
lakum ʿammā kuniza fī al-arḍi law antum tafqahūna.

We have not entered schools or studied learned treatises. Listen to that
with which this unlettered One calls you to God, the Eternal. It is
better for you than all the treasures stored in the earth, if only you
understood.

\subsection{106}\label{section-105}

\begin{otherlanguage}{arabic}

انّ الّذی یاوّل ما نزّل من سمآء الوحی و یخرجه عن الظّاهر انّه ممّن حرّف کلمة
اللّه العلیا و کان من الأخسرین فی کتاب مبین

\end{otherlanguage}

Inna alladhī yuʾawwilu mā nuzzila min samāʾi al-waḥyi wa-yukhrijuhu ʿani
al-ẓāhiri innahu mimman ḥarrafa kalimata Allāhi al-ʿulyā wa-kāna mina
al-akhsarīna fī kitābin mubīnin.

Whoever interprets what has been sent down from the heaven of revelation
and takes it away from its evident meaning is among those who have
distorted the exalted Word of God, and is numbered among the greatest
losers in a clear Book.

\subsection{107}\label{section-106}

\begin{otherlanguage}{persian}

قد کتب علیکم تقلیم الاظفار و الدّخول فی مآء یحیط هیاکلکم فی کلّ اسبوع و
تنظیف ابدانکم بما استعملتموه من قبل ایّاکم ان تمنعکم الغفلة عمّا امرتم به
من لدن عزیز عظیم ادخلوا مآءَ بکرا و المستعمل منه لا یجوز الدّخول فیه ایّاکم
ان تقربوا خزائن حمامات العجم من قصدها وجد رائحتها المنتنة قبل وروده فیها
تجنّبوا یا قوم و لا تکوننّ من الصّاغرین انّه یشبه بالصّدید و الغسلین ان انتم
من العارفین و کذلک حیاضهم المنتنة اترکوها و کونوا من المقدّسین انّا اردنا
ان نراکم مظاهر الفردوس فی الأرض لیتضوّع منکم ما تفرح به افئدة المقرّبین و
الّذی یصبّ علیه المآء و یغسل به بدنه خیر له و یکفیه عن الدّخول انّه اراد ان
یسهّل علیکم الأمور فضلا من عنده لتکونوا من الشّاکرین

\end{otherlanguage}

Qad kutiba ʿalaykum taqlīmu al-aẓfāri wa-al-dukhūlu fī māʾin yuḥīṭu
hayākilakum fī kulli usbūʿin wa-tanẓīfu abdānikum bi-mā istaʿmaltumūhu
min qablu. Iyyākum an tamnaʿakumu al-ghaflatu ʿammā umirtum bihi min
ladun ʿazīzin ʿaẓīmin. Udkhulū māʾan bikran wa-al-mustaʿmalu minhu lā
yajūzu al-dukhūlu fīhi. Iyyākum an taqrabū khazāʾina ḥammāmāti
al-ʿajami. Man qaṣadahā wajada rāʾiḥatahā al-muntinata qabla wurūdihi
fīhā. Tajannabū yā qawmi wa-lā takūnanna mina al-ṣāghirīna. Innahu
yushbihu bi-al-ṣadīdi wa-al-ghislīni in antum mina al-ʿārifīna.
Wa-kadhālika ḥiyāḍuhumu al-muntinatu utrukūhā wa-kūnū mina
al-muqaddasīna. Innā aradnā an narākum maẓāhira al-firdawsi fī al-arḍi
li-yataḍawwaʿa minkum mā tafraḥu bihi afʾidatu al-muqarrabīna.
Wa-alladhī yuṣabbu ʿalayhi al-māʾu wa-yaghsilu bihi badanahu khayrun
lahu wa-yakfīhi ʿani al-dukhūli. Innahu arāda an yusahhila ʿalaykumu
al-umūra faḍlan min ʿindihi li-takūnū mina al-shākirīna.

You are required to trim your nails, immerse yourselves each week in
water that surrounds your bodies, and clean your bodies with what you
have used before. Take care that negligence does not prevent you from
doing what you have been commanded by One who is Almighty, Most Great.
Enter fresh water; immersion in water that has been used is not
permitted. Take care not to approach the reservoirs of the Persian
baths. Anyone making their way to them encounters their foul smell
before entering. Avoid them, O people, and do not be among the abased.
Such water resembles pus and foul discharge, if you understand.
Likewise, abandon their foul-smelling pools and be among the pure. We
have desired to see you as manifestations of paradise on earth, so that
a fragrance may spread from you that gladdens the hearts of those
brought near to God. For anyone who has water poured over him and washes
his body with it, this is better and suffices in place of immersion. He
has desired to make matters easy for you, as a bounty from Him, so that
you may be thankful.

\subsection{108}\label{section-107}

\begin{otherlanguage}{arabic}

قد حرّمت علیکم ازواج آبائکم انّا نستحیی ان نذکر حکم الغلمان اتّقوا الرّحمن
یا ملاء الامکان و لا ترتکبوا ما نهیتم عنه فی اللّوح و لا تکونوا فی هیماء
الشّهوات من الهائمین

\end{otherlanguage}

Qad ḥurrimat ʿalaykum azwāju ābāʾikum. Innā nastaḥyī an nadhkura ḥukma
al-ghilmāni. Ittaqū al-raḥmāna yā malaʾa al-imkāni wa-lā tartakibū mā
nuhītum ʿanhu fī al-lawḥi wa-lā takūnū fī haymāʾi al-shahawāti mina
al-hāʾimīna.

Marriage to your fathers' wives has been forbidden to you. We are
ashamed to mention the ruling concerning boys. Fear the All-Merciful, O
people of the world of existence, and do not commit what you have been
forbidden in the Tablet. Do not be among those who wander in the
wilderness of desires.

\subsection{109}\label{section-108}

\begin{otherlanguage}{arabic}

لیس لأحد ان یحرّک لسانه امام النّاس اذ یمشی فی الطّرق و الأسواق بل ینبغی
لمن اراد الذّکر ان یذکر فی مقام بنی لذکر اللّه او فی بیته هذا اقرب بالخلوص
و التّقوی کذلک اشرقت شمس الحکم من افق البیان طوبی للعاملین

\end{otherlanguage}

Laysa li-aḥadin an yuḥarrika lisānahu amāma al-nāsi idh yamshī fī
al-ṭuruqi wa-al-aswāqi. Bal yanbaghī li-man arāda al-dhikra an yadhkura
fī maqāmin buniya li-dhikri Allāhi aw fī baytihi. Hādhā aqrabu
bi-al-khulūṣi wa-al-taqwā. Kadhālika ashraqat shamsu al-ḥukmi min ufuqi
al-bayāni. Ṭūbā li-al-ʿāmilīna.

No one may move his tongue in remembrance before others while walking in
the streets and marketplaces. Rather, whoever wishes to engage in
remembrance should do so in a place built for the remembrance of God or
in his own home. This is closer to sincerity and piety. Thus has the sun
of the ordinance shone from the horizon of utterance. Blessed are those
who act accordingly.

\subsection{110}\label{section-109}

\begin{otherlanguage}{arabic}

قد فرض لکلّ نفس کتاب الوصیة و له ان یزین رأسه بالاسم الأعظم و یعترف فیه
بوحدانیّة اللّه فی مظهر ظهوره و یذکر فیه ما اراد من المعروف لیشهد له فی
عوالم الأمر و الخلق و یکون له کنزاً عند ربّه الحافظ الأمین

\end{otherlanguage}

Qad furiḍa li-kulli nafsin kitābu al-waṣiyyati wa-lahu an yuzayyina
raʾsahu bi-al-ismi al-aʿẓami wa-yaʿtarifa fīhi bi-waḥdāniyyati Allāhi fī
maẓhari ẓuhūrihi wa-yadhkura fīhi mā arāda mina al-maʿrūfi li-yashhada
lahu fī ʿawālimi al-amri wa-al-khalqi wa-yakūna lahu kanzan ʿinda
rabbihi al-ḥāfiẓi al-amīni.

Everyone is required to write a will. He should adorn its heading with
the Most Great Name, acknowledge in it the oneness of God in the
Manifestation of His revelation, and set down in it whatever good he
wishes, so that it may bear witness for him in the worlds of command and
creation and be a treasure for him with his Lord, the Guardian, the
Faithful.

\subsection{111}\label{section-110}

\begin{otherlanguage}{arabic}

قد انتهت الأعیاد الی العیدین الأعظمین امّا الأوّل ایّام فیها تجلّی الرّحمن
علی من فی الامکان باسمائه الحسنی و صفاته العلیا و الآخر یوم فیه بعثنا من
بشّر النّاس بهذا الاسم الّذی به قامت الأموات و حشر من فی السّموات و الأرضین
و الآخرین فی یومین کذلک قضی الأمر من لدن آمرٍ علیم

\end{otherlanguage}

Qad intahati al-aʿyādu ilā al-ʿīdayni al-aʿẓamayni. Ammā al-awwalu
ayyāmun fīhā tajallā al-raḥmānu ʿalā man fī al-imkāni bi-asmāʾihi
al-ḥusnā wa-ṣifātihi al-ʿulyā wa-al-ākharu yawmun fīhi baʿathnā man
bashshara al-nāsa bi-hādhā al-ismi alladhī bihi qāmati al-amwātu
wa-ḥushira man fī al-samāwāti wa-al-araḍīna wa-al-ākharayni fī yawmayni.
Kadhālika quḍiya al-amru min ladun āmirin ʿalīmin.

The festivals have culminated in the two Most Great Festivals. The first
comprises the days when the All-Merciful revealed Himself to all in the
world of existence through His most beautiful names and most exalted
attributes. The other is the day on which We raised up the One who
announced to humanity this Name, through which the dead have risen and
all who are in the heavens and on earth have been gathered together. The
other two festivals fall on two days. Thus has the matter been decreed
by an all-knowing Ordainer.

\subsection{112}\label{section-111}

\begin{otherlanguage}{arabic}

طوبی لمن فاز بالیوم الأوّل من شهر البهآء الّذی جعله اللّه لهذا الاسم العظیم
طوبی لمن یظهر فیه نعمة اللّه علی نفسه انّه ممّن اظهر شکر اللّه بفعله المدلّ
علی فضله الّذی احاط العالمین قل انّه لصدر الشّهور و مبدئها و فیه تمرّ نفحة
الحیات علی الممکنات طوبی لمن ادرکه بالرّوح و الرّیحان نشهد انّه من الفائزین

\end{otherlanguage}

Ṭūbā li-man fāza bi-al-yawmi al-awwali min shahri al-bahāʾi alladhī
jaʿalahu Allāhu li-hādhā al-ismi al-ʿaẓīmi. Ṭūbā li-man yuẓhiru fīhi
niʿmata Allāhi ʿalā nafsihi. Innahu mimman aẓhara shukra Allāhi
bi-fiʿlihi al-dālli ʿalā faḍlihi alladhī aḥāṭa al-ʿālamīna. Qul innahu
la-ṣadru al-shuhūri wa-mabdaʾuhā wa-fīhi tamurru nafḥatu al-ḥayāti ʿalā
al-mumkināti. Ṭūbā li-man adrakahu bi-al-rūḥi wa-al-rayḥāni. Nashhadu
innahu mina al-fāʾizīna.

Blessed is whoever reaches the first day of the month of Bahāʾ, which
God has appointed for this great Name. Blessed is whoever displays on
that day God's bounty to him. He is among those who have shown gratitude
to God through an act that bears witness to His grace, which encompasses
all worlds. Say: It is the foremost of the months and their beginning,
and in it the breath of life passes over all created things. Blessed is
whoever reaches it with joy and gladness. We testify that he is among
those who have attained.

\subsection{113}\label{section-112}

\begin{otherlanguage}{arabic}

قل انّ العید الاعظم لسلطان الأعیاد اذکروا یا قوم نعمة اللّه علیکم اذ کنتم
رقدآء ایقظکم من نسمات الوحی و عرّفکم سبیله الواضح المستقیم

\end{otherlanguage}

Qul inna al-ʿīda al-aʿẓama la-sulṭānu al-aʿyādi. Udhkurū yā qawmi
niʿmata Allāhi ʿalaykum idh kuntum ruqadāʾa ayqaẓakum min nasamāti
al-waḥyi wa-ʿarrafakum sabīlahu al-wāḍiḥa al-mustaqīma.

Say: The Most Great Festival is indeed the sovereign of festivals.
Remember, O people, God's bounty to you: when you were asleep, He
awakened you through the breezes of revelation and made known to you His
clear and straight path.

\subsection{114}\label{section-113}

\begin{otherlanguage}{arabic}

اذا مرضتم ارجعوا الی الحذّاق من الأطبّآء انّا ما رفعنا الأسباب بل اثبتناها
من هذا القلم الّذی جعله اللّه مطلع امره المشرق المنیر

\end{otherlanguage}

Idhā mariḍtum irjiʿū ilā al-ḥudhdhāqi mina al-aṭibbāʾi. Innā mā rafaʿnā
al-asbāba bal athbatnāhā min hādhā al-qalami alladhī jaʿalahu Allāhu
maṭlaʿa amrihi al-mushriqa al-munīra.

When you are ill, consult skilled physicians. We have not set aside the
use of means; rather, We have affirmed it through this Pen, which God
has made the radiant, luminous dawning place of His Cause.

\subsection{115}\label{section-114}

\begin{otherlanguage}{arabic}

قد کتب اللّه علی کلّ نفس ان یحضر لدی العرشبما عنده ممّا لا عدل له انّا عفونا
عن ذلک فضلا من لدنّا انّه هو المعطی الکریم

\end{otherlanguage}

Qad kataba Allāhu ʿalā kulli nafsin an yuḥḍira ladā al-ʿarshi bi-mā
ʿindahu mimmā lā ʿidla lahu. Innā ʿafawnā ʿan dhālika faḍlan min
ladunnā. Innahu huwa al-muʿṭī al-karīmu.

God had required every person to bring before the Throne whatever he
possessed that was without equal. We have waived this as a bounty from
Us. He is indeed the Bestower, the Generous.

\subsection{116}\label{section-115}

\begin{otherlanguage}{arabic}

طوبی لمن توجّه الی مشرق الأذکار فی الأسحار ذاکرا متذکّرا مستغفرا و اذا دخل
یقعد صامتا لاصغآء آیات اللّه الملک العزیز الحمید قل مشرق الأذکار انّه کلّ
بیت بنی لذکری فی المدن و القری کذلک سمّی لدی العرش ان انتم من العارفین

\end{otherlanguage}

Ṭūbā li-man tawajjaha ilā mashriqi al-adhkāri fī al-asḥāri dhākiran
mutadhakkiran mustaghfiran wa-idhā dakhala yaqʿudu ṣāmitan li-iṣghāʾi
āyāti Allāhi al-maliki al-ʿazīzi al-ḥamīdi. Qul mashriqu al-adhkāri
innahu kullu baytin buniya li-dhikrī fī al-muduni wa-al-qurā. Kadhālika
summiya ladā al-ʿarshi in antum mina al-ʿārifīna.

Blessed is whoever turns toward the Dawning Place of Remembrance at
dawn, remembering, reflecting, and seeking forgiveness, and who, upon
entering, sits silently to listen to the verses of God, the King, the
Almighty, the All-Praised. Say: The Dawning Place of Remembrance is
every house built for My remembrance in cities and villages. Thus has it
been named before the Throne, if you understand.

\subsection{117}\label{section-116}

\begin{otherlanguage}{arabic}

و الّذین یتلون آیات الرّحمن بأحسن الألحان اولئک یدرکون منها ما لا یعادله
ملکوت ملک السّموات و الأرضین و بها یجدون عرف عوالمی الّتی لا یعرفها الیوم
الاّ من اوتی البصر من هذا المنظر الکریم قل انّها تجذب القلوب الصّافیة الی
العوالم الرّوحانیة الّتی لا تعبّر بالعبارة و لا تشار بالاشارة طوبی للسّامعین

\end{otherlanguage}

Wa-alladhīna yatlūna āyāti al-raḥmāni bi-aḥsani al-alḥāni ulāʾika
yudrikūna minhā mā lā yuʿādiluhu malakūtu mulki al-samāwāti
wa-al-araḍīna wa-bihā yajidūna ʿarfa ʿawālimī allatī lā yaʿrifuhā
al-yawma illā man ūtiya al-baṣara min hādhā al-manẓari al-karīmi. Qul
innahā tajdhibu al-qulūba al-ṣāfiyata ilā al-ʿawālimi al-rūḥāniyyati
allatī lā tuʿabbaru bi-al-ʿibārati wa-lā tushāru bi-al-ishārati. Ṭūbā
li-al-sāmiʿīna.

Those who recite the verses of the All-Merciful in the most beautiful
melodies perceive in them what the entire dominion of the heavens and
earth cannot equal. Through them they sense the fragrance of My worlds,
which no one knows today except those who have been granted sight by
this gracious Vision. Say: These verses draw pure hearts toward
spiritual worlds that cannot be expressed in words or indicated by any
allusion. Blessed are those who listen.

\subsection{118}\label{section-117}

\begin{otherlanguage}{arabic}

انصروا یا قوم اصفیائی الّذین قاموا علی ذکری بین خلقی و ارتفاع کلمتی فی
مملکتی اولئک انجم سمآء عنایتی و مصابیح هدایتی للخلائق اجمعین و الّذی
یتکلّم بغیر ما نزّل فی الواحی انّه لیس منّی ایّاکم ان تتّبعوا کلّ مدّعٍ اثیم قد
زینت الألواح بطراز ختم فالق الأصباح الّذی ینطق بین السّموات و الأرضین
تمسّکوا بالعروة الوثقی و حبل امری المحکم المتین

\end{otherlanguage}

Unṣurū yā qawmi aṣfiyāʾī alladhīna qāmū ʿalā dhikrī bayna khalqī
wa-irtifāʿi kalimatī fī mamlakatī. Ulāʾika anjumu samāʾi ʿināyatī
wa-maṣābīḥu hidāyatī li-al-khalāʾiqi ajmaʿīna. Wa-alladhī yatakallamu
bi-ghayri mā nuzzila fī alwāḥī innahu laysa minnī. Iyyākum an tattabiʿū
kulla muddaʿin athīmin. Qad zuyyinati al-alwāḥu bi-ṭirāzi khatmi fāliqi
al-aṣbāḥi alladhī yanṭiqu bayna al-samāwāti wa-al-araḍīna. Tamassakū
bi-al-ʿurwati al-wuthqā wa-ḥabli amrī al-muḥkami al-matīni.

Support, O people, My chosen ones who have arisen to remember Me among
My creation and exalt My Word in My kingdom. They are the stars of the
heaven of My loving care and the lamps of My guidance to all created
beings. Whoever speaks contrary to what has been sent down in My Tablets
is not of Me. Take care not to follow any sinful pretender. The Tablets
have been adorned with the seal of the One who causes the dawn to break,
who speaks amid the heavens and earth. Hold fast to the surest handle
and to the firm, strong cord of My Cause.

\subsection{119}\label{section-118}

\begin{otherlanguage}{arabic}

قد اذن اللّه لمن اراد ان یتعلّم الألسن المختلفة لیبلّغ امر اللّه شرق الأرض و
غربها و یذکره بین الدّول و الملل علی شأن تنجذب به الافئدة و یحیی به کلّ
عظم رمیم

\end{otherlanguage}

Qad adhina Allāhu li-man arāda an yataʿallama al-alsuna al-mukhtalifata
li-yuballigha amra Allāhi sharqa al-arḍi wa-gharbahā wa-yadhkurahu bayna
al-duwali wa-al-milali ʿalā shaʾnin tanjadhibu bihi al-afʾidatu wa-yaḥyā
bihi kullu ʿaẓmin ramīmin.

God has permitted whoever wishes to learn different languages to do so,
so that he may convey the Cause of God to the east and west of the earth
and proclaim it among nations and peoples in such a way that hearts are
attracted and every decayed bone is brought to life.

\subsection{120}\label{section-119}

\begin{otherlanguage}{arabic}

لیس للعاقل ان یشرب ما یذهب به العقل و له ان یعمل ما ینبغی للانسان لا ما
یرتکبه کلّ غافلٍ مریب

\end{otherlanguage}

Laysa li-al-ʿāqili an yashraba mā yadhhabu bihi al-ʿaqlu wa-lahu an
yaʿmala mā yanbaghī li-al-insāni lā mā yartakibuhu kullu ghāfilin
murībin.

A person of understanding must not drink what takes away reason, but
should do what befits a human being, not what every heedless, dubious
person commits.

\subsection{121}\label{section-120}

\begin{otherlanguage}{arabic}

زینوا رؤسکم باکلیل الأمانة و الوفآء و قلوبکم بردآء التّقوی و السنکم
بالصّدق الخالص و هیاکلکم بطراز الآداب کلّ ذلک من سجیة الانسان لو انتم من
المتبصّرین یا اهل البهآء تمسّکوا بحبل العبودیة للّه الحقّ بها تظهر مقاماتکم
و تثبت اسمائکم و ترتفع مراتبکم و اذکارکم فی لوح حفیظ ایّاکم ان یمنعکم من
علی الأرض عن هذا المقام العزیز الرّفیع قد وصّیناکم بها فی اکثر الألواح و
فی هذا اللّوح الّذی لاح من افقه نیّر احکام ربّکم المقتدر الحکیم

\end{otherlanguage}

Zayyinū ruʾūsakum bi-iklīl al-amāna wa-l-wafāʾ wa-qulūbakum bi-ridāʾ
al-taqwā wa-alsinatakum bi-l-ṣidq al-khāliṣ wa-hayākilakum bi-ṭirāz
al-ādāb. Kull dhālika min sajiyyat al-insān law antum min
al-mutabaṣṣirīn. Yā ahl al-Bahāʾ, tamassakū bi-ḥabl al-ʿubūdiyya li-llāh
al-Ḥaqq. Bihā taẓhar maqāmātukum wa-tathbut asmāʾukum wa-tartafiʿ
marātibukum wa-adhkārukum fī lawḥ ḥafīẓ. Iyyākum an yamnaʿakum man ʿalā
al-arḍ ʿan hādhā al-maqām al-ʿazīz al-rafīʿ. Qad waṣṣaynākum bihā fī
akthar al-alwāḥ wa-fī hādhā al-lawḥ alladhī lāḥa min ufuqihi nayyir
aḥkām rabbikum al-Muqtadir al-Ḥakīm.

Adorn your heads with the crown of trustworthiness and fidelity, your
hearts with the robe of reverence for God, your tongues with pure
truthfulness, and your bodies with the ornament of good manners. All
these belong to human nature, if you have insight. People of Bahá, hold
fast to the cord of servitude to God, the Truth. Through it your
stations become manifest, your names are established, and your ranks and
remembrance are exalted in a preserving Tablet. Beware lest anyone on
earth keep you from this mighty, lofty station. We have enjoined it upon
you in most of the Tablets, and in this Tablet, from whose horizon the
sun of the ordinances of your Lord, the All-Powerful, the All-Wise, has
shone.

\subsection{122}\label{section-121}

\begin{otherlanguage}{arabic}

اذا غیض بحر الوصال و قضی کتاب المبدء فی المال توجّهوا الی من اراده اللّه
الّذی انشعب من هذا الأصل القدیم

\end{otherlanguage}

Idhā ghīḍa baḥr al-wiṣāl wa-quḍiya kitāb al-mabdaʾ fī al-maʾāl, tawajj
ahū ilā man arādahu Allāh alladhī inshaʿaba min hādhā al-aṣl al-qadīm.

When the ocean of reunion has ebbed and the Book of the Beginning has
reached its end, turn to the one whom God has intended, who has branched
from this ancient Root.

\subsection{123}\label{section-122}

\begin{otherlanguage}{arabic}

فانظروا فی النّاس و قلّة عقولهم یطلبون ما یضرّهم و یترکون ما ینفعهم الا
انّهم من الهائمین انّا نری بعض النّاس ارادوا الحرّیة و یفتخرون بها اولیک فی
جهلٍ مبین

\end{otherlanguage}

Fa-nẓurū fī al-nās wa-qillat ʿuqūlihim. Yaṭlubūna mā yaḍurruhum
wa-yatrukūna mā yanfaʿuhum. Alā innahum min al-hāʾimīn. Innā narā baʿḍ
al-nās arādū al-ḥurriyya wa-yaftakhirūna bihā. Ulāʾika fī jahl mubīn.

Consider people and how little understanding they have. They seek what
harms them and abandon what benefits them. Truly, they wander aimlessly.
We see that some people desire liberty and take pride in it. They are in
manifest ignorance.

\subsection{124}\label{section-123}

\begin{otherlanguage}{arabic}

انّ الحرّیة تنتهی عواقبها الی الفتنة الّتی لا تخمد نارها کذلک یخبرکم المحصی
العلیم فاعلموا انّ مطالع الحرّیة و مظاهرها هی الحیوان و للانسان ینبغی ان
یکون تحت سننٍ تحفظه عن جهل نفسه و ضرّ الماکرین انّ الحرّیة تخرج الانسان عن
شئون الأدب و الوقار و تجعله من الارذلین

\end{otherlanguage}

Inna al-ḥurriyya tantahī ʿawāqibuhā ilā al-fitna allatī lā tukh m ad
nāruhā. Kadhālika yukhbirukum al-Muḥṣī al-ʿAlīm. Fa-ʿlamū anna maṭāliʿ
al-ḥurriyya wa-maẓāhirahā hiya al-ḥayawān. Wa-li-l-insān yanbaghī an
yakūna taḥta sunan taḥfaẓuhu ʿan jahl nafsihi wa-ḍarr al-mākirīn. Inna
al-ḥurriyya tukhrij al-insān ʿan shuʾūn al-adab wa-l-waqār wa-tajʿaluhu
min al-ardhalīn.

Liberty ultimately leads to turmoil whose fire cannot be extinguished.
Thus the All-Reckoning, the All-Knowing informs you. Know that animals
are the sources and manifestations of liberty. A human being should be
subject to laws that protect him from his own ignorance and the harm of
schemers. Liberty takes a person beyond the bounds of courtesy and
dignity and makes him one of the most abject.

\subsection{125}\label{section-124}

\begin{otherlanguage}{arabic}

فانظروا الخلق کالأغنام لابدّ لها من راعٍ لیحفظها انّ هذا لحقٌّ یقین انّا
نصدّقها فی بعض المقامات دون الآخر انّا کنّا عالمین

\end{otherlanguage}

Fa-nẓurū al-khalq ka-l-aghnām. Lā budda lahā min rāʿ li-yaḥfaẓahā. Inna
hādhā la-ḥaqq yaqīn. Innā nuṣaddiquhā fī baʿḍ al-maqāmāt dūna al-ākhar.
Innā kunnā ʿālimīn.

Look upon humanity as sheep that must have a shepherd to protect them.
This is indeed a certain truth. We approve liberty in some
circumstances, but not in others. We have always been all-knowing.

\subsection{126}\label{section-125}

\begin{otherlanguage}{arabic}

قل الحرّیة فی اتّباع اوامری لو انتم من العارفین لو اتّبع النّاس ما نزّلناه
لهم من سمآء الوحی لیجدنّ انفسهم فی حرّیة بحتة طوبی لمن عرف مراد اللّه فیما
نزّل من سمآء مشیته المهیمنة علی العالمین قل الحرّیة الّتی تنفعکم انّها فی
العبودیة للّه الحق و الّذی وجد حلاوتها لا یبدّلها بملکوت ملک السّموات و
الأرضین

\end{otherlanguage}

Qul al-ḥurriyya fī ittibāʿ awāmirī law antum min al-ʿārifīn. Law
ittabaʿa al-nās mā nazzalnāhu lahum min samāʾ al-waḥy la-yajidanna
anfusahum fī ḥurriyya baḥta. Ṭūbā li-man ʿarafa murād Allāh fīmā nazzala
min samāʾ mashīʾatihi al-muhaymina ʿalā al-ʿālamīn. Qul al-ḥurriyya
allatī tanfaʿukum innahā fī al-ʿubūdiyya li-llāh al-Ḥaqq. Wa-alladhī
wajada ḥalāwatahā lā yubaddiluhā bi-malakūt mulk al-samāwāt wa-l-araḍīn.

Say: Liberty lies in following My commandments, if you understand. If
people followed what We have sent down to them from the heaven of
revelation, they would surely find themselves in pure liberty. Blessed
is the one who recognizes God's purpose in what He has sent down from
the heaven of His will, which reigns over all worlds. Say: The liberty
that benefits you lies in servitude to God, the Truth. Whoever has
tasted its sweetness would not exchange it for sovereignty over the
dominion of the heavens and the earth.

\subsection{127}\label{section-126}

\begin{otherlanguage}{arabic}

حرّم علیکم السّؤال فی البیان عفا اللّه عن ذلک لتسئلواما تحتاج به انفسکم لا
ما تکلّم به رجال قبلکم اتّقوا اللّه و کونوا من المتّقین اسئلوا ما ینفعکم فی
امر اللّه و سلطانه قد فتح باب الفضل علی من فی السّموات و الأرضین

\end{otherlanguage}

Ḥurrima ʿalaykum al-suʾāl fī al-Bayān. ʿAfā Allāh ʿan dhālika li-tasʾalū
mā taḥtāju bihi anfusukum lā mā takallama bihi rijāl qablakum. Ittaqū
Allāh wa-kūnū min al-muttaqīn. Isʾalū mā yanfaʿukum fī amr Allāh
wa-sulṭānihi. Qad futiḥa bāb al-faḍl ʿalā man fī al-samāwāt wa-l-araḍīn.

Asking questions was forbidden to you in the Bayán. God has waived this
prohibition so that you may ask about what you yourselves need, not what
people before you discussed. Be mindful of God and be among those who
revere Him. Ask about what benefits you in the Cause of God and His
sovereignty. The door of grace has been opened to all who are in the
heavens and the earth.

\subsection{128}\label{section-127}

\begin{otherlanguage}{arabic}

انّ عدّة الشّهور تسعة عشر شهراً فی کتاب اللّه قد زین اوّلها بهذا الاسم المهیمن
علی العالمین

\end{otherlanguage}

Inna ʿiddat al-shuhūr tisʿata ʿashara shahran fī kitāb Allāh. Qad
zuyyina awwaluhā bi-hādhā al-ism al-muhaymin ʿalā al-ʿālamīn.

The number of months in the Book of God is nineteen. The first of them
has been adorned with this Name, which reigns over all worlds.

\subsection{129}\label{section-128}

\begin{otherlanguage}{arabic}

قد حکم اللّه دفن الأموات فی البلّور و الأحجار الممتنعة او الاخشاب الصّلبة
اللّطیفة و وضع الخواتیم المنقوشة فی اصابعهم انّه لهو المقتدر العلیم

\end{otherlanguage}

Qad ḥakama Allāh dafn al-amwāt fī al-billawr wa-l-aḥjār al-mumtaniʿa aw
al-akhshāb al-ṣalba al-laṭīfa wa-waḍʿ al-khawātīm al-manqūsha fī
aṣābiʿihim. Innahu la-huwa al-Muqtadir al-ʿAlīm.

God has decreed that the dead be buried in crystal, durable stone, or
fine hardwood, and that engraved rings be placed on their fingers. He is
indeed the All-Powerful, the All-Knowing.

\subsection{130}\label{section-129}

\begin{otherlanguage}{persian}

یکتب للرّجال و للّه ما فی السّموات و الأرض و ما بینهما و کان اللّه بکلّ شئ
علیماً و للورقات و للّه ملک السّموات و الأرض و ما بینهما و کان اللّه علی کلّ
شئ قدیرا هذا ما نزّل من قبل و ینادی نقطة البیان و یقول یا محبوب الامکان
انطق فی هذا المقام بما تتضوّع به نفحات الطافک بین العالمین انّا اخبرنا
الکلّ بأن لا یعادل بکلمة منک ما نزّل فی البیان انّک انت المقتدر علی ما تشآء
لا تمنع عبادک عن فیوضات بحر رحمتک انّک انت ذو الفضل العظیم قد استجبنا ما
اراد انّه لهو المحبوب المجیب لو ینقش علیها ما نزّل فی الحین من لدی اللّه
انّه خیر لهم و لهنّ انّا کنّا حاکمین قد بدئت من اللّه و رجعت الیه منقطعا عمّا
سواه و متمسّکا باسمه الرّحمن الرّحیم کذلک یختصّ اللّه من یشآء بفضل من عنده
انّه لهو المقتدر القدیر

\end{otherlanguage}

Yuktab li-l-rijāl: Wa-li-llāh mā fī al-samāwāt wa-l-arḍ wa-mā baynahumā
wa-kāna Allāh bi-kull shayʾ ʿalīman. Wa-li-l-waraqāt: Wa-li-llāh mulk
al-samāwāt wa-l-arḍ wa-mā baynahumā wa-kāna Allāh ʿalā kull shayʾ
qadīran. Hādhā mā nuzzila min qabl. Wa-yunādī Nuqṭat al-Bayān wa-yaqūl:
Yā maḥbūb al-imkān, unṭuq fī hādhā al-maqām bi-mā tataḍawwaʿu bihi
nafaḥāt alṭāfika bayna al-ʿālamīn. Innā akhbarnā al-kull bi-an lā
yuʿādalu bi-kalima minka mā nuzzila fī al-Bayān. Innaka anta al-muqtadir
ʿalā mā tashāʾ. Lā tamnaʿ ʿibādaka ʿan fuyūḍāt baḥr raḥmatika. Innaka
anta dhū al-faḍl al-ʿaẓīm. Qad istajabnā mā arāda. Innahu la-huwa
al-Maḥbūb al-Mujīb. Law yunqashu ʿalayhā mā nuzzila fī al-ḥīn min ladā
Allāh innahu khayr lahum wa-lahunna. Innā kunnā ḥākimīn. Qad badaʾtu min
Allāh wa-rajaʿtu ilayhi munqaṭiʿan ʿammā siwāhu wa-mutamassikan
bi-ismihi al-Raḥmān al-Raḥīm. Kadhālika yakhtaṣṣu Allāh man yashāʾ
bi-faḍl min ʿindihi. Innahu la-huwa al-Muqtadir al-Qadīr.

For men, the inscription is: ``To God belongs all that is in the heavens
and the earth and whatever is between them, and God knows all things.''
For the leaves: ``To God belongs the dominion of the heavens and the
earth and whatever is between them, and God has power over all things.''
This is what was sent down before. The Point of the Bayán calls out and
says: ``Beloved of all existence, speak on this matter words through
which the fragrances of Your loving-kindness may spread among all
worlds. We have informed everyone that nothing revealed in the Bayán can
equal a single word from You. You have power to do whatever You will. Do
not withhold from Your servants the outpourings of the ocean of Your
mercy. You are the Possessor of immense grace.'' We have answered His
wish. He is indeed the Beloved, the Responsive. If what has now been
sent down from God is engraved upon them, it is better for both men and
women. We have indeed decreed: ``I came from God and have returned to
Him, detached from all but Him and holding fast to His name, the
All-Merciful, the Compassionate.'' Thus God singles out whomever He
wills for grace from His presence. He is indeed the All-Powerful, the
Almighty.

\subsection{131}\label{section-130}

\begin{otherlanguage}{arabic}

و ان تکفنوه فی خمسة اثواب من الحریر او القطن من لم یستطع یکتفی بواحدة
منهما کذلک قضی الأمر من لدن علیم خبیر حرّم علیکم نقل المیت ازید من مسافة
ساعة من المدینة ادفنوه بالرّوح و الرّیحان فی مکان قریب

\end{otherlanguage}

Wa-an tukaffinūhu fī khamsat athwāb min al-ḥarīr aw al-quṭn. Man lam
yastaṭiʿ yaktafī bi-wāḥida minhumā. Kadhālika quḍiya al-amr min ladun
ʿAlīm Khabīr. Ḥurrima ʿalaykum naql al-mayyit azyad min masāfat sāʿa min
al-madīna. Idfinūhu bi-l-rūḥ wa-l-rayḥān fī makān qarīb.

And wrap the deceased in five lengths of silk or cotton. Whoever cannot
afford this may make do with one of either. Thus has the ordinance been
decreed by One who is All-Knowing, All-Aware. You are forbidden to
transport the deceased more than an hour's journey from the city. Bury
him with joy and fragrance in a nearby place.

\subsection{132}\label{section-131}

\begin{otherlanguage}{arabic}

قد رفع اللّه ما حکم به البیان فی تحدید الأسفا انّه لهو المختار یفعل ما
یشآء و یحکم ما یرید

\end{otherlanguage}

Qad rafaʿa Allāh mā ḥakama bihi al-Bayān fī taḥdīd {[}corrupted text{]}.
Innahu la-huwa al-Mukhtār yafʿalu mā yashāʾ wa-yaḥkumu mā yurīd.

God has lifted what the Bayán ordained concerning the restriction of
{[}corrupted original text{]}. He is indeed free to choose: He does
whatever He wills and ordains whatever He desires.

\subsection{133}\label{section-132}

\begin{otherlanguage}{arabic}

یا ملاء الانشآء اسمعوا ندآء مالک الأسمآء انّه ینادیکم من شطر سجنه الأعظم
انّه لا اله الاّ انا المقتدر المتکبّر المتسخّر المتعالی العلیم الحکیم انّه لا
اله الاّ هو المقتدر علی العالمین لو یشآء یاخذ العالم بکلمة من عنده ایّاکم
ان تتوقّفوا فی هذا الأمر الّذی خضع له الملأ الأعلی و اهل مدائن الأسمآء
اتّقوا اللّه و لا تکوننّ من المحتجبین احرقوا الحجبات بنار حبّی و السّبحات
بهذا الاسم الّذی به سخّرنا العالمین

\end{otherlanguage}

Yā malaʾ al-inshāʾ, ismaʿū nidāʾ Mālik al-Asmāʾ. Innahu yunādīkum min
shaṭr sijnihi al-aʿẓam: Innahu lā ilāha illā anā al-Muqtadir
al-Mutakabbir al-Mutasakhkhir al-Mutaʿālī al-ʿAlīm al-Ḥakīm. Innahu lā
ilāha illā huwa al-muqtadir ʿalā al-ʿālamīn. Law yashāʾ yaʾkhudh
al-ʿālam bi-kalima min ʿindihi. Iyyākum an tatawaqqafū fī hādhā al-amr
alladhī khaḍaʿa lahu al-malaʾ al-aʿlā wa-ahl madāʾin al-asmāʾ. Ittaqū
Allāh wa-lā takūnanna min al-muḥtajibīn. Aḥriqū al-ḥujubāt bi-nār ḥubbī
wa-l-subuḥāt bi-hādhā al-ism alladhī bihi sakhkharnā al-ʿālamīn.

People of creation, hear the call of the Lord of Names. He calls to you
from His Most Great Prison: ``There is no God but Me, the All-Powerful,
the Supremely Great, the One in service, the Most Exalted, the
All-Knowing, the All-Wise.'' There is no God but Him, who has power over
all worlds. If He wills, He can seize the world with a word from
Himself. Beware of hesitating in this Cause, before which the Concourse
on high and the inhabitants of the cities of names have bowed. Be
mindful of God, and do not be among those who are veiled. Burn away the
veils with the fire of My love, and the veils of glory with this Name
through which We have subdued all worlds.

\subsection{134}\label{section-133}

\begin{otherlanguage}{persian}

و ارفعنّ البیتین فی المقامین و المقامات الّتی فیها استقرّ عرش ربّکم الرّحمن
کذلک یأمرکم مولی العارفین

\end{otherlanguage}

Wa-irfaʿanna al-baytayn fī al-maqāmayn wa-l-maqāmāt allatī fīhā
istaqarra ʿarsh rabbikum al-Raḥmān. Kadhālika yaʾmurukum Mawlā
al-ʿārifīn.

Raise up the two Houses in the two places, and the sites where the
throne of your Lord, the All-Merciful, was established. Thus commands
you the Lord of those who know.

\subsection{135}\label{section-134}

\begin{otherlanguage}{persian}

ایّاکم ان تمنعکم شئونات الأرض عمّا امرتم به من لدن قوی امین کونوا مظاهر
الاستقامة بین البریّة علی شأن لا تمنعکم شبهات الّذین کفروا باللّه اذ ظهر
بسلطان عظیم ایّاکم ان یمنعکم ما نزّل فی الکتاب عن هذا الکتاب الّذی ینطق
بالحقّ انّه لا اله الاّ انا العزیز الحمید انظروا بعین الانصاف الی من اتی من
سمآء المشیّة و الاقتدار و لا تکوننّ من الظّالمین

\end{otherlanguage}

Iyyākum an tamnaʿakum shuʾūnāt al-arḍ ʿammā umirtum bihi min ladun
Qawiyy Amīn. Kūnū maẓāhir al-istiqāma bayna al-bariyya ʿalā shaʾn lā
tamnaʿukum shubuhāt alladhīna kafarū bi-llāh idh ẓahara bi-sulṭān ʿaẓīm.
Iyyākum an yamnaʿakum mā nuzzila fī al-kitāb ʿan hādhā al-kitāb alladhī
yanṭiqu bi-l-ḥaqq: Innahu lā ilāha illā anā al-ʿAzīz al-Ḥamīd. Unẓurū
bi-ʿayn al-inṣāf ilā man atā min samāʾ al-mashīʾa wa-l-iqtidār wa-lā
takūnanna min al-ẓālimīn.

Beware lest earthly affairs keep you from what you have been commanded
by One who is Strong and Trustworthy. Be manifestations of steadfastness
among humanity, so that the doubts of those who disbelieved in God when
He appeared with mighty sovereignty cannot deter you. Beware lest what
was revealed in the Book keep you from this Book, which declares the
truth: ``There is no God but Me, the Almighty, the Praiseworthy.'' Look
with fairness upon the One who has come from the heaven of will and
power, and do not be among the unjust.

\subsection{136}\label{section-135}

\begin{otherlanguage}{persian}

ثمّ اذکروا ما جری من قلم مبشّری فی ذکر هذا الظّهور و ما ارتکبه اولو الطّغیان
فی ایّامه الا انّهم من الأخسرین قال ان ادرکتم ما نظهره انتم من فضل اللّه
تسئلون لیمنّ علیکم باستوائه علی سرائرکم فانّ ذلک عزّ ممتنع منیع ان یشرب کأس
مآء عندکم اعظم من ان تشربنّ کلّ نفس مآء وجوده بل کلّ شئ ان یا عبادی تدرکون

\end{otherlanguage}

Thumma udhkurū mā jarā min qalam mubashshirī fī dhikr hādhā al-ẓuhūr
wa-mā irtakabahu ulū al-ṭughyān fī ayyāmihi. Alā innahum min
al-akhsarīn. Qāla: In adraktum mā nuẓhiruhu antum min faḍl Allāh
tasʾalūn li-yamunna ʿalaykum bi-istiwāʾihi ʿalā sarāʾirikum fa-inna
dhālika ʿizz mumtaniʿ manīʿ. An yashraba kaʾs māʾ ʿindakum aʿẓam min an
tashrabanna kull nafs māʾ wujūdihi bal kull shayʾ, in yā ʿibādī
tudrikūn.

Then remember what flowed from the pen of My Herald concerning this
Revelation, and what the tyrannical committed in His days. Truly, they
are among the greatest losers. He said: ``If you attain the presence of
the One We shall manifest, ask of God's grace that He may favor you by
taking His seat in your innermost hearts, for that is an inaccessible,
exalted honor. For Him to drink a cup of water in your presence is
greater than for every soul---indeed, every thing---to drink the water
of its own existence, if you understand, My servants.''

\subsection{137}\label{section-136}

\begin{otherlanguage}{arabic}

هذا ما نزّل من عنده ذکرا لنفسی لو انتم تعلمون و الّذی تفکّر فی هذه الآیات و
اطّلع بما ستر فیهنّ من اللآلیٴ المخزونة تاللّه انّه یجد عرف الرّحمن من شطر
السّجن و یسرع بقلبه الیه باشتیاق لا تمنعه جنود السّموات و الأرضین قل هذا
لظهور تطوف حوله الحجّة و البرهان کذلک انزله الرّحمن ان انتم من المنصفین قل
هذا روح الکتب قد نفخ به فی القلم الأعلی و انصعق من فی الانشآء الاّ من
اخذته نفحات رحمتی و فوحات الطافی المهیمنة علی العالمین

\end{otherlanguage}

Hādhā mā nuzzila min ʿindihi dhikran li-nafsī law antum taʿlamūn.
Wa-alladhī tafakkara fī hādhihi al-āyāt wa-iṭṭalaʿa bi-mā sutira fīhinna
min al-laʾāliʾ al-makhzūna, ta-llāh innahu yajidu ʿarf al-Raḥmān min
shaṭr al-sijn wa-yusriʿu bi-qalbihi ilayhi bi-ishtiyāq lā tamnaʿuhu
junūd al-samāwāt wa-l-araḍīn. Qul hādhā la-ẓuhūr taṭūfu ḥawlahu al-ḥujja
wa-l-burhān. Kadhālika anzalahu al-Raḥmān in antum min al-munṣifīn. Qul
hādhā rūḥ al-kutub qad nufikha bihi fī al-qalam al-aʿlā wa-inṣaʿiqa man
fī al-inshāʾ illā man akhadhathu nafaḥāt raḥmatī wa-fawḥāt alṭāfī
al-muhaymina ʿalā al-ʿālamīn.

This is what was sent down from Him in remembrance of Me, if you only
knew. Whoever reflects on these verses and discovers the treasured
pearls concealed within them will, by God, perceive the fragrance of the
All-Merciful from the direction of the prison and hasten toward Him in
his heart with a longing that the armies of heaven and earth cannot
restrain. Say: This is indeed a Revelation around which proof and
testimony circle. Thus has the All-Merciful sent it down, if you judge
fairly. Say: This is the spirit of the scriptures, breathed into the
Most Exalted Pen, at which all creation has fallen thunderstruck, except
those whom the breezes of My mercy and the fragrances of My
loving-kindness, which reigns over all worlds, have embraced.

\subsection{138}\label{section-137}

\begin{otherlanguage}{arabic}

یا ملاء البیان اتّقوا الرّحمن ثمّ انظروا ما انزله فی مقام آخر قال انّما
القبلة من یظهره اللّه متی ینقلب تنقلب الی ان یستقرّ کذلک نزّل من لدن مالک
القدر اذ اراد ذکر هذا المنظر الأکبر تفکّروا یا قوم و لا تکوننّ من الهآئمین
لو تنکرونه بأهوائکم الی آیة قبلة تتوجّهون یا معشر الغافلین تفکّروا فی هذه
الآیة ثمّ انصفوا باللّه لعلّ تجدون لآلیٴ الأسرار من البحر الّذی تموّج باسمی
العزیز المنیع

\end{otherlanguage}

Yā malaʾ al-Bayān, ittaqū al-Raḥmān thumma unẓurū mā anzalahu fī maqām
ākhar. Qāla: Innamā al-qibla man yuẓhiruhu Allāh matā yanqalib tanqalib
ilā an yastaqirr. Kadhālika nuzzila min ladun Mālik al-Qadar idh arāda
dhikr hādhā al-manẓar al-akbar. Tafakkarū yā qawm wa-lā takūnanna min
al-hāʾimīn. Law tunkirūnahu bi-ahwāʾikum ilā ayyat qibla tataw ajjahūn
yā maʿshar al-ghāfilīn? Tafakkarū fī hādhihi al-āya thumma anṣifū
bi-llāh laʿalla tajidūna laʾāliʾ al-asrār min al-baḥr alladhī tamawwaja
bi-ismī al-ʿAzīz al-Manīʿ.

People of the Bayán, be mindful of the All-Merciful, and consider what
He revealed elsewhere. He said: ``The Qibla is none other than the One
whom God will manifest. Whenever He moves, it moves, until He comes to
rest.'' Thus was it sent down by the Lord of Destiny when He wished to
speak of this Most Great Vision. Reflect, people, and do not be among
those who wander aimlessly. If you deny Him because of your desires,
toward what Qibla will you turn, heedless ones? Reflect on this verse,
and then, by God, judge fairly, so that you may find the pearls of
mysteries in the ocean that has surged through My name, the Almighty,
the Inaccessible.

\subsection{139}\label{section-138}

\begin{otherlanguage}{arabic}

لیس لأحد ان یتمسّک الیوم الاّ بما ظهر فی هذا الظّهور هذا حکم اللّه من قبل و
من بعد و به زیّن صحف الأوّلین هذا ذکر اللّه من قبل و من بعد قد طرّز به دیباج
کتاب الوجود ان انتم من الشّاعرین هذا امر اللّه من قبل و من بعد ایّاکم ان
تکونوا من الصّاغرین لا یغنیکم الیوم شئ و لیس لأحد مهرب الاّ اللّه العلیم
الحکیم من عرفنی فقد عرف المقصود من توجّه الی قد توجّه الی المعبود کذلک فصّل
فی الکتاب و قضی الأمر من لدی اللّه ربّ العالمین من یقرء آیة من آیاتی لخیر
له من ان یقرء کتب الأوّلین و الأخرین هذا بیان الرّحمن ان انتم من السّامعین
قل هذا حقّ العلم لو انتم من العارفین

\end{otherlanguage}

Laysa li-aḥad an yatamassaka al-yawm illā bi-mā ẓahara fī hādhā
al-ẓuhūr. Hādhā ḥukm Allāh min qabl wa-min baʿd wa-bihi zuyyina ṣuḥuf
al-awwalīn. Hādhā dhikr Allāh min qabl wa-min baʿd qad ṭurriza bihi dībā
j kitāb al-wujūd in antum min al-shāʿirīn. Hādhā amr Allāh min qabl
wa-min baʿd. Iyyākum an takūnū min al-ṣāghirīn. Lā yughnīkum al-yawm
shayʾ wa-laysa li-aḥad mahrab illā Allāh al-ʿAlīm al-Ḥakīm. Man ʿarafanī
fa-qad ʿarafa al-maqṣūd. Man tawajjaha ilayya qad tawajjaha ilā
al-maʿbūd. Kadhālika fuṣṣila fī al-kitāb wa-quḍiya al-amr min ladā Allāh
rabb al-ʿālamīn. Man yaqraʾ āya min āyātī la-khayr lahu min an yaqraʾ
kutub al-awwalīn wa-l-ākhirīn. Hādhā bayān al-Raḥmān in antum min
al-sāmiʿīn. Qul hādhā ḥaqq al-ʿilm law antum min al-ʿārifīn.

Today no one may hold fast to anything but what has appeared in this
Revelation. This is God's decree, before and hereafter, and with it the
scriptures of the ancients were adorned. This is God's remembrance,
before and hereafter, with which the brocade of the Book of Existence
has been embroidered, if you perceive. This is God's command, before and
hereafter. Beware of being among the abased. Nothing will avail you
today, and no one has any refuge except God, the All-Knowing, the
All-Wise. Whoever has recognized Me has recognized the One who is
sought; whoever has turned toward Me has turned toward the One who is
worshipped. Thus has it been set out in the Book, and the matter decreed
by God, the Lord of all worlds. For anyone to read a single one of My
verses is better than reading the books of the ancients and those who
came after them. This is the utterance of the All-Merciful, if you
listen. Say: This is true knowledge, if you understand.

\subsection{140}\label{section-139}

\begin{otherlanguage}{arabic}

ثمّ انظروا ما نزّل فی مقام اخر لعلّ تدعون ما عندکم مقبلین الی اللّه ربّ
العالمین قال لا یحلّ الاقتران ان لم یکن فی البیان و ان یدخل من احد یحرم
علی الآخر ما یملک من عنده الاّ و ان یرجع ذلک بعد ان یرفع امر من نظهره
بالحقّ او ما قد ظهر بالعدل و قبل ذلک فلتقربنّ لعلّکم بذلک امر اللّه ترفعون
کذلک تغرّدت الورقآء علی الافنان فی ذکر ربّها الرّحمن طوبی للسّامعین

\end{otherlanguage}

Thumma unẓurū mā nuzzila fī maqām ākhar laʿalla tadaʿūna mā ʿindakum
muqbilīn ilā Allāh rabb al-ʿālamīn. Qāla: Lā yaḥillu al-iqtirān in lam
yakun fī al-Bayān wa-in yadkhul min aḥad yaḥrumu ʿalā al-ākhar mā
yamliku min ʿindihi illā wa-an yarjiʿa dhālika baʿda an yurfaʿa amr man
nuẓhiruhu bi-l-ḥaqq aw mā qad ẓahara bi-l-ʿadl wa-qabla dhālika
fa-la-taqrubanna laʿallakum bi-dhālika amr Allāh tarfaʿūn. Kadhālika
tagharradat al-warqāʾ ʿalā al-afnān fī dhikr rabbihā al-Raḥmān. Ṭūbā
li-l-sāmiʿīn.

Then consider what was revealed elsewhere, so that you may relinquish
what you possess and turn toward God, the Lord of all worlds. He said:
``Marriage is not lawful unless it is within the Bayán. If one of the
two enters it, whatever that one possesses becomes unlawful to the
other, unless the other returns, after the Cause of the One We shall
manifest in truth, or what has appeared in justice, has been raised up.
Before that, you must draw close, so that through this you may exalt the
Cause of God.'' Thus has the Dove sung upon the branches in remembrance
of her Lord, the All-Merciful. Blessed are those who listen.

\subsection{141}\label{section-140}

\begin{otherlanguage}{arabic}

یا ملاء البیان اقسمکم بربّکم الرّحمن بأن تنظروا فیما نزّل بالحقّ بعین
الانصاف و لا تکوننّ من الّذین یرون برهان اللّه و ینکرونه الا انّهم من
الهالکین قد صرّح نقطة البیان فی هذه الآیة بارتفاع امری قبل امره یشهد بذلک
کلّ منصف علیمٍ کما ترونه الیوم انّه ارتفع علی شأن لا ینکره الا الّذین سکرت
ابصارهم فی الأولی و فی الأخری لهم عذاب مهین

\end{otherlanguage}

Yā malaʾa al-Bayān, uqsimukum bi-rabbikumu al-Raḥmān bi-an tanẓurū fīmā
nuzzila bi-al-ḥaqq bi-ʿayni al-inṣāf wa-lā takūnunna mina alladhīna
yarawna burhāna Allāh wa-yunkirūnahu. Alā innahum mina al-hālikīn. Qad
ṣarraḥa Nuqṭatu al-Bayān fī hādhihi al-āyah bi-irtifāʿi amrī qabla
amrihi, yashhadu bi-dhālika kullu munṣifin ʿalīm, kamā tarawnahu
al-yawma innahu irtafaʿa ʿalā shaʾnin lā yunkiruhu illā alladhīna
sakirat abṣāruhum fī al-ūlā wa-fī al-ukhrā lahum ʿadhābun muhīn.

O people of the Bayán! I implore you by your Lord, the All-Merciful, to
look with fairness upon what has been revealed in truth, and not to be
among those who see the proof of God yet deny it. They are indeed among
those who perish. In this verse, the Point of the Bayán explicitly
declared that My Cause would be exalted before His own. Every
fair-minded, knowledgeable person bears witness to this, just as you see
today that it has been exalted to a degree that none deny except those
whose sight is bewildered in this world; in the next, a humiliating
punishment awaits them.

\subsection{142}\label{section-141}

\begin{otherlanguage}{arabic}

قل تاللّه انّی لمحبوبه و الان یسمع ما ینزل من سمآء الوحی و ینوح بما
ارتکبتم فی ایّامه خافوا اللّه و لا تکوننّ من المعتدین قل یا قوم ان لن
تؤمنوا به لا تعترضوا علیه تاللّه یکفی ما اجتمع علیه من جنود الظّالمین

\end{otherlanguage}

Qul ta-Allāhi innī la-maḥbūbuhu wa-al-āna yasmaʿu mā yanzilu min samāʾi
al-waḥy wa-yanūḥu bi-mā irtakabtum fī ayyāmihi. Khāfū Allāh wa-lā
takūnunna mina al-muʿtadīn. Qul yā qawmi in lan tuʾminū bihi lā
taʿtariḍū ʿalayhi. Ta-Allāhi yakfī mā ijtamaʿa ʿalayhi min junūdi
al-ẓālimīn.

Say: By God, I am indeed His Beloved, and He now hears what descends
from the heaven of revelation and laments what you have committed in His
days. Fear God, and do not be among the transgressors. Say: O people! If
you will not believe in Him, do not oppose Him. By God, the hosts of
oppressors already arrayed against Him are enough.

\subsection{143}\label{section-142}

\begin{otherlanguage}{arabic}

انّه قد انزل بعض الأحکام لئلّا یتحرّک القلم الأعلی فی هذا الظّهور الاّ علی
ذکر مقاماته العلیا و منظره الأسنی و انّا لمّا اردنا الفضل فصّلناها بالحقّ و
خفّفنا ما اردناه لکم انّه لهو الفضّال الکریم

\end{otherlanguage}

Innahu qad anzala baʿḍa al-aḥkām li-allā yataḥarraka al-Qalamu al-Aʿlā
fī hādhā al-ẓuhūr illā ʿalā dhikri maqāmātihi al-ʿulyā wa-manẓarihi
al-asnā. Wa-innā lammā aradnā al-faḍla faṣṣalnāhā bi-al-ḥaqq
wa-khaffafnā mā aradnāhu lakum. Innahu la-huwa al-Faḍḍālu al-Karīm.

He revealed certain laws so that, in this Revelation, the Most Exalted
Pen would move only to mention His highest stations and His most radiant
appearance. Yet when We desired to bestow Our bounty, We set them out in
truth and lightened what We had intended for you. He is indeed the
Bountiful, the Generous.

\subsection{144}\label{section-143}

\begin{otherlanguage}{arabic}

قد اخبرکم من قبل بما ینطق به هذا الذّکر الحکیم قال و قوله الحقّ انّه ینطق
فی کلّ شأن انّه لا اله الاّ انا الفرد الواحد العلیم الخبیر هذا مقام خصّه
اللّه لهذا الظّهور الممتنع البدیع هذا من فضل اللّه ان انتم من العارفین هذا
من امره المبرم و اسمه الأعظم و کلمته العلیا و مطلع اسمائه الحسنی لو انتم
من العالمین بل به تظهر المطالع و المشارق تفکروا یا قومِ فیما نزّل بالحقّ و
تدبّروا فیه و لا تکوننّ من المعتدین

\end{otherlanguage}

Qad akhbarakum min qablu bi-mā yanṭiqu bihi hādhā al-Dhikru al-Ḥakīm.
Qāla wa-qawluhu al-ḥaqq: innahu yanṭiqu fī kulli shaʾnin innahu lā ilāha
illā anā al-Fardu al-Wāḥidu al-ʿAlīmu al-Khabīr. Hādhā maqāmun khaṣṣahu
Allāhu li-hādhā al-ẓuhūri al-mumtaniʿi al-badīʿ. Hādhā min faḍli Allāh
in antum mina al-ʿārifīn. Hādhā min amrihi al-mubram wa-ismihi al-aʿẓam
wa-kalimatihi al-ʿulyā wa-maṭlaʿi asmāʾihi al-ḥusnā law antum mina
al-ʿālimīn. Bal bihi taẓharu al-maṭāliʿu wa-al-mashāriq. Tafakkarū yā
qawmi fīmā nuzzila bi-al-ḥaqq wa-tadabbarū fīhi wa-lā takūnunna mina
al-muʿtadīn.

He previously informed you of what this wise Remembrance would proclaim.
He said---and His word is the truth---``He proclaims in every
circumstance: `There is no God but Me, the Single, the One, the
All-Knowing, the All-Informed.'\,'' This is a station God has reserved
for this inaccessible, wondrous Revelation. This is from the bounty of
God, if you are among those who recognize. This proceeds from His
irrevocable command, His Greatest Name, His most exalted Word, and the
Dawning-Place of His most beautiful names, if you are among those who
know. Indeed, through Him the dawning-places and daysprings appear.
Reflect, O people, on what has been revealed in truth, ponder it, and do
not be among the transgressors.

\subsection{145}\label{section-144}

\begin{otherlanguage}{arabic}

عاشروا مع الأدیان بالرّوح و الرّیحان لیجدوا منکم عرف الرّحمن ایّاکم ان
تأخذکم حمیّة الجاهلیّة بین البریّة کلّ بدء من اللّه و یعود الیه انّه لمبدء
الخلق و مرجع العالمین

\end{otherlanguage}

ʿĀshirū maʿa al-adyān bi-al-rūḥi wa-al-rayḥān li-yajidū minkum ʿarfa
al-Raḥmān. Iyyākum an taʾkhudhakum ḥamiyyatu al-jāhiliyyah bayna
al-bariyyah. Kullun badaʾa mina Allāh wa-yaʿūdu ilayhi. Innahu
la-mabdaʾu al-khalqi wa-marjiʿu al-ʿālamīn.

Associate with the followers of all religions in joy and fellowship, so
that they may perceive from you the fragrance of the All-Merciful.
Beware lest the zealotry of ignorance seize you among the people. All
have originated from God and return to Him. He is indeed the source of
creation and the One to whom all worlds return.

\subsection{146}\label{section-145}

\begin{otherlanguage}{arabic}

ایّاکم ان تدخلوا بیتاً عند فقدان صاحبه الاّ بعد اذنه تمسّکوا بالمعروف فی کلّ
الأحوال و لا تکوننّ من الغافلین

\end{otherlanguage}

Iyyākum an tadkhulū baytan ʿinda fuqdāni ṣāḥibihi illā baʿda idhnihi.
Tamassakū bi-al-maʿrūfi fī kulli al-aḥwāl wa-lā takūnunna mina
al-ghāfilīn.

Beware of entering a house in its owner's absence without first
obtaining permission. Adhere to what is right in all circumstances, and
do not be among the heedless.

\subsection{147}\label{section-146}

\begin{otherlanguage}{arabic}

قد کتب علیکم تزکیة الاقوات و ما دونها بالزّکوة هذا ما حکم به منزل الآیات
فی هذا الرّقّ المنیع سوف نفصّل لکم نصابها اذا شآء اللّه و اراد انّه یفصّل ما
یشآء بعلم من عنده انّه لهو العلاّم الحکیم

\end{otherlanguage}

Qad kutiba ʿalaykum tazkiyatu al-aqwāti wa-mā dūnahā bi-al-zakāh. Hādhā
mā ḥakama bihi Munzilu al-Āyāt fī hādhā al-raqqi al-manīʿ. Sawfa
nufaṣṣilu lakum niṣābahā idhā shāʾa Allāhu wa-arād. Innahu yufaṣṣilu mā
yashāʾu bi-ʿilmin min ʿindihi. Innahu la-huwa al-ʿAllāmu al-Ḥakīm.

You are required to purify your provisions and other possessions through
the payment of zakát. This is what the Revealer of the Verses has
decreed in this exalted parchment. We will set out its minimum
assessment for you when God wishes and wills. He sets out whatever He
wishes through knowledge from His own presence. He is indeed the
All-Knowing, the All-Wise.

\subsection{148}\label{section-147}

\begin{otherlanguage}{arabic}

لا یحلّ السّؤال و من سئل حرّم علیه العطآء قد کتب علی الکّل ان یکسب و الّذی
عجز فللوکلآء و الأغنیآء ان یعینوا له ما یکفیه اعملوا حدود اللّه و سننه ثمّ
احفظوها کما تحفظون اعینکم و لا تکوننّ من الخاسرین

\end{otherlanguage}

Lā yaḥillu al-suʾālu wa-man suʾila ḥurrima ʿalayhi al-ʿaṭāʾ. Qad kutiba
ʿalā al-kulli an yaksiba wa-alladhī ʿajaza fa-li-al-wukalāʾi
wa-al-aghniyāʾi an yuʿīnū lahu mā yakfīhi. Iʿmalū ḥudūda Allāhi
wa-sunanahu thumma iḥfaẓūhā kamā taḥfaẓūna aʿyunakum wa-lā takūnunna
mina al-khāsirīn.

Begging is not permitted, and giving to one who begs is forbidden.
Everyone is required to earn a living. For anyone unable to do so, it is
for the trustees and the wealthy to provide what is sufficient. Observe
God's laws and ordinances, and guard them as you guard your own eyes. Do
not be among those who suffer loss.

\subsection{149}\label{section-148}

\begin{otherlanguage}{arabic}

قد منعتم فی الکتاب عن الجدال و النّزاع و الضّرب و امثالها عمّا تحزن به
الافئدة و القلوب من یحزن احداً فله ان ینفق تسعة عشر مثقالاً من الذّهب هذا
ما حکم به مولی العالمین انّه قد عفا ذلک عنکم فی هذا الظّهور و یوصیکم بالبرّ
و التّقوی امراً من عنده فی هذا اللّوح المنیر لا ترضوا لأحدٍ ما لا ترضونه
لانفسکم اتّقوا اللّه و لا تکوننّ من المتکبّرین کلّکم خلقتم من المآء و ترجعون
الی التّراب تفکّروا فی عواقبکم و لا تکوننّ من الظّالمین اسمعوا ما تتلو
السّدرة علیکم من آیات اللّه انّها لقسطاس الهدی من اللّه ربّ الآخرة و الأولی و
بها تطیر النّفوس الی مطلع الوحی و تستضئی افئدة المقبلین تلک حدود اللّه قد
فرضت علیکم و تلک اوامر اللّه قد امرتم بها فی اللّوح اعملوا بالرّوح و
الرّیحان هذا خیر لکم ان انتم من العارفین

\end{otherlanguage}

Qad muniʿtum fī al-kitābi ʿani al-jidāli wa-al-nizāʿi wa-al-ḍarbi
wa-amthālihā ʿammā taḥzanu bihi al-afʾidatu wa-al-qulūb. Man yuḥzin
aḥadan fa-lahu an yunfiqa tisʿata ʿashara mithqālan mina al-dhahab.
Hādhā mā ḥakama bihi Mawlā al-ʿālamīn. Innahu qad ʿafā dhālika ʿankum fī
hādhā al-ẓuhūr wa-yūṣīkum bi-al-birri wa-al-taqwā amran min ʿindihi fī
hādhā al-lawḥi al-munīr. Lā tarḍaw li-aḥadin mā lā tarḍawnahu
li-anfusikum. Ittaqū Allāh wa-lā takūnunna mina al-mutakabbirīn.
Kullukum khuliqtum mina al-māʾi wa-tarjiʿūna ilā al-turāb. Tafakkarū fī
ʿawāqibikum wa-lā takūnunna mina al-ẓālimīn. Ismaʿū mā tatlū al-Sidratu
ʿalaykum min āyāti Allāh. Innahā la-qisṭāsu al-hudā mina Allāhi Rabbi
al-ākhirati wa-al-ūlā wa-bihā taṭīru al-nufūsu ilā maṭlaʿi al-waḥyi
wa-tastaḍīʾu afʾidatu al-muqbilīn. Tilka ḥudūdu Allāhi qad furiḍat
ʿalaykum wa-tilka awāmiru Allāhi qad umirtum bihā fī al-lawḥ. Iʿmalū
bi-al-rūḥi wa-al-rayḥān. Hādhā khayrun lakum in antum mina al-ʿārifīn.

In the Book you have been forbidden to engage in argument, conflict,
violence, and similar acts that bring sorrow to hearts and souls.
Whoever causes another sorrow is to expend nineteen mithqáls of gold.
This is what the Lord of all worlds has decreed. In this Revelation He
has relieved you of that obligation and enjoins goodness and piety upon
you as a command from Him in this luminous Tablet. Do not wish for
anyone what you would not wish for yourselves. Fear God, and do not be
among the arrogant. You were all created from water and will return to
dust. Reflect on your ultimate fate, and do not be among the wrongdoers.
Listen to the verses of God that the Lote Tree recites to you. They are
indeed the standard of guidance from God, the Lord of the world to come
and of this world. Through them souls soar to the Dawning-Place of
revelation, and the hearts of those who turn toward Him are illumined.
These are the laws of God that have been made binding upon you, and
these are the commandments of God that you have been enjoined to follow
in the Tablet. Observe them with joy and gladness. This is better for
you, if you are among those who recognize.

\subsection{150}\label{section-149}

\begin{otherlanguage}{arabic}

اتلوا آیات اللّه فی کلّ صباح و مسآء انّ الّذی لم یتل لم یوف بعهد اللّه و
میثاقه و الّذی اعرض عنها الیوم انّه ممّن اعرض عن اللّه فی ازل الآزال اتّقنّ
اللّه یا عبادی کلّکم اجمعون لا تغرّنکم کثرة القرائة و الاعمال فی اللیل و
النّهار لو یقرء احدٌ آیةٌ من الآیات بالرّوح و الرّیحان خیر له من ان یتلو
بالکسالة صحف اللّه المهیمن القیوم اتلوا آیات اللّه علی قدر لا تأخذکم
الکسالة و الأحزان لا تحملوا علی الأرواح ما یکسلها و یثقلها بل ما یخفّفها
لتطیر باجنحة الآیات الی مطلع البینات هذا اقرب الی اللّه لو انتم تعقلون

\end{otherlanguage}

Utlū āyāti Allāhi fī kulli ṣabāḥin wa-masāʾ. Inna alladhī lam yatlu lam
yūfi bi-ʿahdi Allāhi wa-mīthāqihi wa-alladhī aʿraḍa ʿanhā al-yawma
innahu mimman aʿraḍa ʿani Allāhi fī azali al-āzāl. Ittaqinna Allāha yā
ʿibādī kullukum ajmaʿūn. Lā taghurrannakum kathratu al-qirāʾati
wa-al-aʿmāli fī al-layli wa-al-nahār. Law yaqraʾu aḥadun āyatun mina
al-āyāti bi-al-rūḥi wa-al-rayḥān khayrun lahu min an yatluwa
bi-al-kasālati ṣuḥufa Allāhi al-Muhaymini al-Qayyūm. Utlū āyāti Allāhi
ʿalā qadrin lā taʾkhudhukumu al-kasālatu wa-al-aḥzān. Lā taḥmilū ʿalā
al-arwāḥi mā yuksiluhā wa-yuthqiluhā bal mā yukhaffifuhā li-taṭīra
bi-ajniḥati al-āyāti ilā maṭlaʿi al-bayyināt. Hādhā aqrabu ilā Allāhi
law antum taʿqilūn.

Recite the verses of God every morning and evening. Whoever does not
recite them has not fulfilled God's covenant and pledge, and whoever
turns away from them today is among those who have turned away from God
from all eternity. Fear God, My servants, each and every one of you. Do
not let abundant recitation and deeds by night and day delude you. For
someone to read a single verse with joy and gladness is better than to
read with weariness the scriptures of God, the All-Protecting, the
Self-Subsisting. Recite the verses of God to such an extent that
weariness and sorrow do not overtake you. Do not burden souls with what
wearies and weighs them down, but rather with what lightens them, so
that they may soar on the wings of the verses to the Dawning-Place of
clear proofs. This is nearer to God, if you understand.

\subsection{151}\label{section-150}

\begin{otherlanguage}{arabic}

علّموا ذرّیاتکم ما نزّل من سمآء العظمة و الاقتدار لیقرٴوا الواح الرّحمن بحسن
الألحان فی الغرف المبنیة فی مشارق الأذکار انّ الّذی اخذه جذب محبّة اسمی
الرّحمن انّه یقرء آیات اللّه علی شأن تنجذب به افئدة الرّاقدین هنیئاً لمن شرب
رحیق الحیوان من بیان ربّه الرّحمن بهذا الاسم الّذی به نسف کلّ جبل باذخ رفیع

\end{otherlanguage}

ʿAllimū dhurriyyātikum mā nuzzila min samāʾi al-ʿaẓamati wa-al-iqtidār
li-yaqraʾū alwāḥa al-Raḥmāni bi-ḥusni al-alḥāni fī al-ghurafi
al-mabniyyati fī mashāriqi al-adhkār. Inna alladhī akhadhahu jadhbu
maḥabbati ismī al-Raḥmān innahu yaqraʾu āyāti Allāhi ʿalā shaʾnin
tanjadhibu bihi afʾidatu al-rāqidīn. Hanīʾan li-man shariba raḥīqa
al-ḥayawāni min bayāni rabbihi al-Raḥmān bi-hādhā al-ismi alladhī bihi
nusifa kullu jabalin bādhikhin rafīʿ.

Teach your children what has been revealed from the heaven of grandeur
and power, so that they may recite the Tablets of the All-Merciful in
beautiful melodies in the chambers built within the Dawning-Places of
Remembrance. Whoever is seized by the rapture of love for My name, the
All-Merciful, recites the verses of God in such a manner that the hearts
of those who sleep are enraptured. Blessed is the one who has drunk the
choice wine of life from the utterance of his Lord, the All-Merciful,
through this Name by which every lofty, towering mountain has been
scattered.

\subsection{152}\label{section-151}

\begin{otherlanguage}{arabic}

کتب علیکم تجدید اسباب البیت بعد انقضآء تسع عشرة سنة کذلک قضی الأمر من
لدن علیم خبیر انّه اراد تلطیفکم و ما عندکم اتّقوا اللّه و لا تکوننّ من
الغافلین و الّذی لم یستطع عفا اللّه عنه انّه لهو الغفور الکریم

\end{otherlanguage}

Kutiba ʿalaykum tajdīdu asbābi al-bayti baʿda inqiḍāʾi tisʿa ʿashrata
sanah. Kadhālika quḍiya al-amru min ladun ʿalīmin khabīr. Innahu arāda
talṭīfakum wa-mā ʿindakum. Ittaqū Allāha wa-lā takūnunna mina
al-ghāfilīn. Wa-alladhī lam yastaṭiʿ ʿafā Allāhu ʿanhu. Innahu la-huwa
al-Ghafūru al-Karīm.

You are required to renew your household furnishings after nineteen
years have elapsed. Thus has the matter been decreed by One who is
All-Knowing, All-Informed. He desires refinement for you and for what
you possess. Fear God, and do not be among the heedless. Whoever is
unable to do so is excused by God. He is indeed the All-Forgiving, the
Generous.

\subsection{153}\label{section-152}

\begin{otherlanguage}{arabic}

اغسلوا ارجلکم کلّ یوم فی الصّیف و فی الشّتآء کلّ ثلثة ایّام مرّة واحدة

\end{otherlanguage}

Ighsilū arjulakum kulla yawmin fī al-ṣayfi wa-fī al-shitāʾi kulla
thalāthati ayyāmin marratan wāḥidah.

Wash your feet every day in summer, and once every three days in winter.

\subsection{154}\label{section-153}

\begin{otherlanguage}{arabic}

و من اغتاظ علیکم قابلوه بالرّفق و الّذی زجرکم لا تزجروه دعوه بنفسه و
توکّلوا علی اللّه المنتقم العادل القدیر

\end{otherlanguage}

Wa-man ightāẓa ʿalaykum qābilūhu bi-al-rifqi wa-alladhī zajarakum lā
tazjurūhu. Daʿūhu bi-nafsihi wa-tawakkalū ʿalā Allāhi al-Muntaqimi
al-ʿĀdili al-Qadīr.

Meet anyone who is angry with you with gentleness, and do not rebuke
anyone who rebukes you. Leave that person alone, and put your trust in
God, the Avenger, the Just, the All-Powerful.

\subsection{155}\label{section-154}

\begin{otherlanguage}{arabic}

قد منعتم عن الارتقآء الی المنابر من اراد ان یتلو علیکم الآیات ربّه فلیقعد
علی الکرسی الموضوع علی السّریر و یذکر اللّه ربّه و ربّ العالمین قد احبّ اللّه
جلوسکم علی السّرر و الکراسی لعزّ ما عندکم من حبّ اللّه و مطلع امره المشرق
المنیر

\end{otherlanguage}

Qad muniʿtum ʿani al-irtiqāʾi ilā al-manābir. Man arāda an yatluwa
ʿalaykumu al-āyāti rabbihi fa-l-yaqʿud ʿalā al-kursiyyi al-mawḍūʿi ʿalā
al-sarīri wa-yadhkuri Allāha rabbahu wa-rabba al-ʿālamīn. Qad aḥabba
Allāhu julūsakum ʿalā al-sururi wa-al-karāsī li-ʿizzi mā ʿindakum min
ḥubbi Allāhi wa-maṭlaʿi amrihi al-mushriqi al-munīr.

You have been forbidden to ascend pulpits. Anyone who wishes to recite
the verses of his Lord to you should sit on a chair placed upon a
platform and remember God, his Lord and the Lord of all worlds. God has
wished you to sit on platforms and chairs, in honor of the love you bear
for God and for the radiant, luminous Dawning-Place of His Cause.

\subsection{156}\label{section-155}

\begin{otherlanguage}{arabic}

حرّم علیکم المیسر و الأفیون اجتنبوا یا معشر الخلق و لا تکوننّ من
المتجاوزین ایّاکم ان تستعملوا ما تکسل به هیاکلکم و یضرّ ابدانکم انّا ما
اردنا لکم الا ما ینفعکم یشهد بذلک کلّ الأشیآء لو انتم تسمعون

\end{otherlanguage}

Ḥurrima ʿalaykumu al-maysiru wa-al-afyūn. Ijtanibū yā maʿshara al-khalqi
wa-lā takūnunna mina al-mutajāwizīn. Iyyākum an tastaʿmilū mā taksalu
bihi hayākilukum wa-yaḍurru abdānakum. Innā mā aradnā lakum illā mā
yanfaʿukum. Yashhadu bi-dhālika kullu al-ashyāʾi law antum tasmaʿūn.

Gambling and opium have been forbidden to you. Avoid them, O people of
creation, and do not be among those who overstep the bounds. Beware of
using anything that makes your bodies sluggish and harms them. We have
desired for you only what benefits you. All things bear witness to this,
if you listen.

\subsection{157}\label{section-156}

\begin{otherlanguage}{arabic}

اذا دعیتم الی الولآئم و العزائم اجیبوا بالفرح و الانبساط و الّذی وفی
بالوعد انّه امن من الوعید هذا یوم فیه فصّل کلّ امر حکیم

\end{otherlanguage}

Idhā duʿītum ilā al-walāʾimi wa-al-ʿazāʾimi ajībū bi-al-faraḥi
wa-al-inbisāṭ. Wa-alladhī wafā bi-al-waʿdi innahu amina mina al-waʿīd.
Hādhā yawmun fīhi fuṣṣila kullu amrin ḥakīm.

When you are invited to feasts and festive gatherings, respond with joy
and delight. Whoever keeps a promise is safe from the threat of
punishment. This is a day on which every wise decree has been set out.

\subsection{158}\label{section-157}

\begin{otherlanguage}{arabic}

قد ظهر سرّ التّنکیس لرمز الرّئیس طوبی لمن ایده اللّه علی الاقرار بالسّتّة الّتی
ارتفعت بهذا الألف القائمة الا انّه من المخلصین کم من ناسک اعرض و کم من
تارک اقبل و قال لک الحمد یا مقصود العالمین انّ الأمر بید اللّه یعطی من
یشآء ما یشآء و یمنع عمّن یشآء ما اراد یعلم خافیة القلوب و ما یتحرّک به
اعین اللّامزین کم من غافلٍ اقبل بالخلوص اقعدناه علی سریر القبول و کم من
عاقل رجعناه الی النّار عدلا من عندنا انّا کنّا حاکمین انّه لمظهر یفعل اللّه
ما یشآء و المستقرّ علی عرش یحکم ما یرید

\end{otherlanguage}

Qad ẓahara sirru al-tankīsi li-ramzi al-raʾīs. Ṭūbā li-man ayyadahu
Allāhu ʿalā al-iqrāri bi-al-sittati allatī irtafaʿat bi-hādhihi al-alifi
al-qāʾimah. Alā innahu mina al-mukhliṣīn. Kam min nāsikin aʿraḍa wa-kam
min tārikٍ aqbala wa-qāla laka al-ḥamdu yā Maqṣūda al-ʿālamīn. Inna
al-amra bi-yadi Allāh yuʿṭī man yashāʾu mā yashāʾu wa-yamnaʿu ʿamman
yashāʾu mā arād. Yaʿlamu khāfiyata al-qulūbi wa-mā tataḥarraku bihi
aʿyunu al-lāmizīn. Kam min ghāfilin aqbala bi-al-khulūṣi aqʿadnāhu ʿalā
sarīri al-qabūli wa-kam min ʿāqilin rajaʿnāhu ilā al-nāri ʿadlan min
ʿindinā. Innā kunnā ḥākimīn. Innahu la-maẓharu yafʿalu Allāhu mā yashāʾu
wa-al-mustaqirru ʿalā ʿarshi yaḥkumu mā yurīd.

The mystery of the inversion in the symbol of the Chief has been
revealed. Blessed is the one whom God has enabled to acknowledge the six
raised up by this upright alif. He is indeed among the sincere. How many
a devout worshipper has turned away, and how many a neglectful one has
turned toward Him and said, ``Praise be to You, O Desire of all
worlds!'' The matter is in God's hands: He gives to whomever He wishes
whatever He wishes, and withholds from whomever He wishes whatever He
wills. He knows what hearts conceal and what moves the eyes of those who
find fault. How many a heedless person who turned toward Us with
sincerity have We seated upon the throne of acceptance, and how many a
learned person have We consigned to the fire, in justice on Our part. We
have indeed been the judge. He is indeed the manifestation of ``God does
whatever He wishes'' and the One established upon the throne of ``He
decrees whatever He wills.''

\subsection{159}\label{section-158}

\begin{otherlanguage}{arabic}

طوبی لمن وجد عرف المعانی من اثر هذا القلم الّذی اذا تحرّک فاحت نسمة اللّه
فیما سواه و اذا توقّف ظهرت کینونة الاطمینان فی الامکان تعالی الرّحمن مظهر
هذا الفضل العظیم قل بما حمل الظّلم ظهر العدل فیما سواه و بما قبل الذّلّة
لاح عزّ اللّه بین العالمین

\end{otherlanguage}

Ṭūbā li-man wajada ʿarfa al-maʿānī min athari hādhā al-qalami alladhī
idhā taḥarraka fāḥat nasmatu Allāhi fīmā siwāhu wa-idhā tawaqqafa
ẓaharat kaynūnatu al-iṭmiʾnāni fī al-imkān. Taʿālā al-Raḥmānu muẓhiru
hādhā al-faḍli al-ʿaẓīm. Qul bi-mā ḥamala al-ẓulma ẓahara al-ʿadlu fīmā
siwāhu wa-bi-mā qabila al-dhillata lāḥa ʿizzu Allāhi bayna al-ʿālamīn.

Blessed is the one who has perceived the fragrance of inner meanings
from the traces of this Pen, which, when it moves, spreads the breeze of
God through all that is other than Him, and, when it rests, manifests
the very essence of tranquillity throughout the world of being. Exalted
is the All-Merciful, who manifests this great bounty. Say: Because He
bore oppression, justice appeared throughout all that is other than Him,
and because He accepted humiliation, the glory of God shone among all
worlds.

\subsection{160}\label{section-159}

\begin{otherlanguage}{arabic}

حرّم علیکم حمل آلات الحرب الاّ حین الضّرورة و احلّ لکم لبس الحریر قد رفع
اللّه عنکم حکم الحدّ فی اللّباس و اللّحی فضلاً من عنده انّه لهو الآمر العلیم
اعملوا ما لا تنکره العقول المستقیمة و لا تجعلوا انفسکم ملعب الجاهلین
طوبی لمن تزیّن بطراز الآداب و الأخلاق انّه ممّن نصر ربّه بالعمل الواضح
المبین

\end{otherlanguage}

Ḥurrima ʿalaykum ḥamlu ālāti al-ḥarbi illā ḥīna al-ḍarūrati wa-uḥilla
lakum lubsu al-ḥarīr. Qad rafaʿa Allāhu ʿankum ḥukma al-ḥaddi fī
al-libāsi wa-al-liḥā faḍlan min ʿindihi. Innahu la-huwa al-Āmiru
al-ʿAlīm. Iʿmalū mā lā tunkiruhu al-ʿuqūlu al-mustaqīmatu wa-lā tajʿalū
anfusakum malʿaba al-jāhilīn. Ṭūbā li-man tazayyana bi-ṭirāzi al-ādābi
wa-al-akhlāq. Innahu mimman naṣara rabbahu bi-al-ʿamali al-wāḍiḥi
al-mubīn.

You are forbidden to carry weapons except when necessary, and you are
permitted to wear silk. As a bounty from Him, God has removed the
restrictions on clothing and beards. He is indeed the Ordainer, the
All-Knowing. Do what sound minds do not disapprove of, and do not make
yourselves the playthings of the ignorant. Blessed is the one adorned
with good manners and character. Such a person is among those who have
aided their Lord through clear and manifest deeds.

\subsection{161}\label{section-160}

\begin{otherlanguage}{arabic}

عمّروا دیار اللّه و بلاده ثمّ اذکروه فیها بترنّمات المقرّبین انّما تعمّر القلوب
باللّسان کما تعمّر البیوت و الدّیار بالید و اسباب آخر قد قدّرنا لکلّ شئ سبباً
من عندنا تمسّکوا به و توکّلوا علی الحکیم الخبیر

\end{otherlanguage}

ʿAmmirū diyāra Allāhi wa-bilādahu thumma udhkurūhu fīhā bi-tarannumāṭi
al-muqarrabīn. Innamā tuʿammaru al-qulūbu bi-al-lisāni kamā tuʿammaru
al-buyūtu wa-al-diyāru bi-al-yadi wa-asbābin ukhar. Qad qaddarnā
li-kulli shayʾin sababan min ʿindinā. Tamassakū bihi wa-tawakkalū ʿalā
al-Ḥakīmi al-Khabīr.

Develop God's cities and lands, and remember Him there in the melodies
of those brought near to Him. Hearts are built up through the tongue,
just as houses and cities are built through the hand and other means. We
have appointed a means for everything from Our own presence. Hold fast
to it, and put your trust in the All-Wise, the All-Informed.

\subsection{162}\label{section-161}

\begin{otherlanguage}{arabic}

طوبی لمن اقرّ باللّه و آیاته و اعترف بانّه لا یسئل عمّا یفعل هذه کلمة قد
جعلها اللّه طراز العقاید و اصلها و بها یقبل عمل العاملین اجعلوا هذه
الکلمة نصب عیونکم لئلّا تزلّکم اشارات المعرضین

\end{otherlanguage}

Ṭūbā li-man aqarra bi-Allāh wa-āyātihi wa-iʿtarafa bi-annahu lā yusʾalu
ʿammā yafʿal. Hādhihi kalima qad jaʿalahā Allāh ṭirāz al-ʿaqāʾid
wa-aṣlahā wa-bihā yuqbalu ʿamal al-ʿāmilīn. Ijʿalū hādhihi al-kalima
nuṣba ʿuyūnikum li-allā tuzillakum ishārāt al-muʿriḍīn.

Blessed is the one who acknowledges God and His verses and recognizes
that He is not to be questioned about what He does. This is a statement
that God has made the adornment and foundation of all beliefs, and
through it the deeds of those who act are accepted. Keep this statement
before your eyes, so that the insinuations of those who turn away do not
cause you to stumble.

\subsection{163}\label{section-162}

\begin{otherlanguage}{arabic}

لو یحلّ ما حرّم فی ازل الآزال او بالعکس لیس لأحد ان یعترض علیه و الّذی توقّف
فی اقلّ من انٍ انّه من المعتدین

\end{otherlanguage}

Law yuḥillu mā ḥarrama fī azal al-āzāl aw bi-l-ʿaks laysa li-aḥad an
yaʿtariḍa ʿalayhi wa-alladhī tawaqqafa fī aqalla min ānin innahu min
al-muʿtadīn.

Were He to make lawful what He had forbidden from all eternity, or the
reverse, no one would have the right to object to Him. Whoever hesitates
for less than an instant is among the transgressors.

\subsection{164}\label{section-163}

\begin{otherlanguage}{persian}

و الّذی ما فاز بهذا الأصل الأسنی و المقام الأعلی تحرّکه اریاح الشّبهات و
تقلّبه مقالات المشرکین من فاز بهذا الأصل قد فاز بالاستقامة الکبری حبّذا
هذا المقام الأبهی الّذی بذکره زین کلّ لوح منیع کذلک یعلّمکم اللّه ما یخلّصکم
عن الرّیب و الحیرة و ینجّیکم فی الدّنیا و الآخرة انّه هو الغفور الکریم هو
الّذی ارسل الرّسل و انزل الکتب علی انّه لا اله الّا انا العزیز الحکیم

\end{otherlanguage}

Wa-alladhī mā fāza bi-hādhā al-aṣl al-asnā wa-l-maqām al-aʿlā
tuḥarrikuhu aryāḥ al-shubuhāt wa-tuqallibuhu maqālāt al-mushrikīn. Man
fāza bi-hādhā al-aṣl qad fāza bi-l-istiqāma al-kubrā. Ḥabbadhā hādhā
al-maqām al-abhā alladhī bi-dhikrihi zuyyina kull lawḥ manīʿ. Kadhālika
yuʿallimukum Allāh mā yukhalliṣukum ʿan al-rayb wa-l-ḥayra wa-yunajjīkum
fī al-dunyā wa-l-ākhira. Innahu huwa al-ghafūr al-karīm. Huwa alladhī
arsala al-rusul wa-anzala al-kutub ʿalā annahu lā ilāha illā anā
al-ʿazīz al-ḥakīm.

Whoever has not attained this most exalted principle and this highest
station is shaken by the winds of doubt and tossed about by the
assertions of those who associate partners with God. Whoever has
attained this principle has attained the greatest steadfastness. How
excellent is this most glorious station, with whose mention every
exalted Tablet is adorned! Thus God teaches you what will free you from
doubt and bewilderment and save you in this world and the next. He is
indeed the All-Forgiving, the All-Bountiful. He is the One who sent the
Messengers and revealed the Books, declaring: ``There is no god but Me,
the Almighty, the All-Wise.''

\subsection{165}\label{section-164}

\begin{otherlanguage}{persian}

یا ارض الکاف و الرّآء انّا نراک علی ما لا یحبّه اللّه و نری منک ما لا اطّلع
به احد الّا اللّه العلیم الخبیر و نجد ما یمرّ منک فی سرّ السّرّ عندنا علم کلّ
شئ فی لوح مبین لا تحزنی بذلک سوف یظهر اللّه فیک اولی باس شدید یذکروننی
باستقامة لا تمنعهم اشارات العلمآء و لا تحجبهم شبهات المریبین اولئک
ینظرون اللّه باعینهم و ینصرونه بأنفسهم الا انّهم من الرّاسخین

\end{otherlanguage}

Yā arḍ al-kāf wa-l-rāʾ innā narāki ʿalā mā lā yuḥibbuhu Allāh wa-narā
minki mā lā iṭṭalaʿa bihi aḥad illā Allāh al-ʿalīm al-khabīr wa-najidu
mā yamurru minki fī sirr al-sirr. ʿIndanā ʿilm kull shayʾ fī lawḥ mubīn.
Lā taḥzanī bi-dhālika sawfa yuẓhiru Allāh fīki ulī baʾs shadīd
yadhkurūnanī bi-istiqāma lā tamnaʿuhum ishārāt al-ʿulamāʾ wa-lā
taḥjubuhum shubuhāt al-murībīn. Ulāʾika yanẓurūna Allāh bi-aʿyunihim
wa-yanṣurūnahu bi-anfusihim alā innahum min al-rāsikhīn.

O land of Kāf and Rāʾ! We see you in a condition that God does not love,
and We see from you what no one knows but God, the All-Knowing, the
All-Informed. We perceive what passes within you in the deepest secrecy.
With Us is the knowledge of all things in a clear Tablet. Do not grieve
over this. God will soon raise up within you people of mighty strength
who will remember Me with a steadfastness that neither the insinuations
of the learned can hinder nor the doubts of those who sow doubt can
veil. They will see God with their own eyes and aid Him with their own
selves. They are indeed among those firmly established.

\subsection{166}\label{section-165}

\begin{otherlanguage}{arabic}

یا معشر العلمآء لمّا نزّلت الآیات و ظهرت البینات رأیناکم خلف الحجبات ان
هذا الّا شئ عجاب قد افتخرتم باسمی و غفلتم عن نفسی اذ اتی الرّحمن بالحجّة و
البرهان انّا خرقنا الأحجاب ایّاکم ان تحجبوا النّاس بحجاب اخر کسّروا سلاسل
الأوهام باسم مالک الأنام و لا تکوننّ من الخادعین اذا اقبلتم الی اللّه و
دخلتم هذا الأمر لا تفسدوا فیه و لا تقیسوا کتاب اللّه بأهوائکم هذا نصح
اللّه من قبل و من بعد یشهد بذلک شهداء اللّه و اصفیائه انّا کلّ له شاهدون

\end{otherlanguage}

Yā maʿshar al-ʿulamāʾ lammā nuzzilat al-āyāt wa-ẓaharat al-bayyināt
raʾaynākum khalf al-ḥujubāt. In hādhā illā shayʾ ʿujāb. Qad iftakhartum
bi-ismī wa-ghafaltum ʿan nafsī idh atā al-Raḥmān bi-l-ḥujja wa-l-burhān.
Innā kharaqnā al-aḥjāb iyyākum an taḥjubū al-nās bi-ḥijāb ākhar. Kassirū
salāsil al-awhām bi-ism mālik al-anām wa-lā takūnanna min al-khādiʿīn.
Idhā aqbaltum ilā Allāh wa-dakhaltum hādhā al-amr lā tufsidū fīhi wa-lā
taqīsū kitāb Allāh bi-ahwāʾikum. Hādhā nuṣḥ Allāh min qabl wa-min baʿd
yashhadu bi-dhālika shuhadāʾ Allāh wa-aṣfiyāʾuhu innā kull lahu
shāhidūn.

O company of the learned! When the verses were revealed and the clear
proofs appeared, We saw you behind veils. This is truly an astonishing
thing. You took pride in My name but neglected My own Self when the
All-Merciful came with testimony and proof. We have torn through the
veils; beware lest you veil the people with another veil. Break the
chains of vain imaginings in the name of the Lord of all people, and do
not be among the deceivers. When you turn to God and enter this Cause,
do not spread corruption within it or measure the Book of God against
your own desires. This is God's counsel, before and hereafter. God's
witnesses and His chosen ones bear witness to this; We all bear witness
to it.

\subsection{167}\label{section-166}

\begin{otherlanguage}{arabic}

اذکروا الشّیخ الّذی سمّی بمحمّد قبل حسن و کان من اعلم العلمآء فی عصره لمّا
ظهر الحقّ اعرض عنه هو و امثاله و اقبل الی اللّه من ینقّی القمح و الشّعیر و
کان یکتب علی زعمه احکام اللّه فی اللّیل و النّهار و لمّا اتی المختار ما نفعه
حرف منها لو نفعه لم یعرض عن وجه به انارت وجوه المقرّبین لو آمنتم باللّه
حین ظهوره ما اعرض عنه النّاس و ما ورد علینا ما ترونه الیوم اتّقوا اللّه و
لا تکوننّ من الغافلین

\end{otherlanguage}

Udhkurū al-shaykh alladhī summiya bi-Muḥammad qabla Ḥasan wa-kāna min
aʿlam al-ʿulamāʾ fī ʿaṣrihi lammā ẓahara al-ḥaqq aʿraḍa ʿanhu huwa
wa-amthāluhu wa-aqbala ilā Allāh man yunaqqī al-qamḥ wa-l-shaʿīr.
Wa-kāna yaktubu ʿalā zaʿmihi aḥkām Allāh fī al-layl wa-l-nahār wa-lammā
atā al-mukhtār mā nafaʿahu ḥarf minhā. Law nafaʿahu lam yuʿriḍ ʿan wajh
bihi anārat wujūh al-muqarrabīn. Law āmantum bi-Allāh ḥīna ẓuhūrihi mā
aʿraḍa ʿanhu al-nās wa-mā warada ʿalaynā mā tarawnahu al-yawm. Ittaqū
Allāh wa-lā takūnanna min al-ghāfilīn.

Remember the shaykh whose name was Muḥammad followed by Ḥasan, who was
among the most learned scholars of his age. When the Truth appeared, he
and others like him turned away from Him, while a sifter of wheat and
barley turned to God. Night and day the shaykh wrote what he imagined
were the laws of God, yet when the Chosen One came, not a single letter
of them benefited him. Had they benefited him, he would not have turned
away from the Face through which the faces of those brought near to God
were illumined. Had you believed in God when He appeared, the people
would not have turned away from Him, nor would We have suffered what you
see today. Fear God, and do not be among the heedless.

\subsection{168}\label{section-167}

\begin{otherlanguage}{persian}

ایّاکم ان تمنعکم الأسماء عن مالکها او یحجبکم ذکر عن هذا الذّکر الحکیم
استعیذوا باللّه یا معشر العلمآء و لا تجعلوا انفسکم حجاباً بینی و بین خلقی
کذلک یعظکم اللّه و یأمرکم بالعدل لئلّا تحبط اعمالکم و انتم غافلون انّ الّذی
اعرض عن هذا الأمر هل یقدر ان یثبت حقاً فی الابداع لا و مالک الاختراع ولکنّ
النّاس فی حجاب مبین قل به اشرقت شمس الحجّة و لاح نیر البرهان لمن فی
الامکان اتّقوا اللّه یا اولی الابصار و لا تنکرون ایّاکم ان یمنعکم ذکر النّبی
عن هذا النّبأ الأعظم او الولایة عن ولایة اللّه المهیمنة علی العالمین قد
خلق کلّ اسم بقوله و علّق کلّ امر بأمره المبرم العزیز البدیع قل هذا یوم اللّه
لا یذکر فیه الاّ نفسه المهیمنة علی العالمین هذا امر اضطرب منه ما عندکم من
الأوهام و التّماثیل

\end{otherlanguage}

Iyyākum an tamnaʿakum al-asmāʾ ʿan mālikihā aw yaḥjubakum dhikr ʿan
hādhā al-dhikr al-ḥakīm. Istaʿīdhū bi-Allāh yā maʿshar al-ʿulamāʾ wa-lā
tajʿalū anfusakum ḥijāban baynī wa-bayna khalqī. Kadhālika yaʿiẓukum
Allāh wa-yaʾmurukum bi-l-ʿadl li-allā taḥbaṭa aʿmālukum wa-antum
ghāfilūn. Inna alladhī aʿraḍa ʿan hādhā al-amr hal yaqdiru an yuthbita
ḥaqqan fī al-ibdāʿ lā wa-mālik al-ikhtirāʿ wa-lākinna al-nās fī ḥijāb
mubīn. Qul bihi ashraqat shams al-ḥujja wa-lāḥa nayyir al-burhān li-man
fī al-imkān. Ittaqū Allāh yā ulī al-abṣār wa-lā tunkirūn. Iyyākum an
yamnaʿakum dhikr al-nabī ʿan hādhā al-nabaʾ al-aʿẓam aw al-wilāya ʿan
wilāyat Allāh al-muhaymina ʿalā al-ʿālamīn. Qad khalaqa kull ism
bi-qawlihi wa-ʿallaqa kull amr bi-amrihi al-mubram al-ʿazīz al-badīʿ.
Qul hādhā yawm Allāh lā yudhkaru fīhi illā nafsuhu al-muhaymina ʿalā
al-ʿālamīn. Hādhā amr iḍṭaraba minhu mā ʿindakum min al-awhām
wa-l-tamāthīl.

Beware lest names keep you from their Lord, or any remembrance veil you
from this wise Remembrance. Seek refuge in God, O company of the
learned, and do not make yourselves a veil between Me and My creation.
Thus God admonishes you and commands you to act justly, lest your deeds
come to nothing while you remain unaware. Can one who has turned away
from this Cause establish any truth in creation? No, by the Lord of all
that has been brought into being! Yet the people are plainly veiled.
Say: Through Him the sun of testimony has shone and the luminary of
proof has dawned for all who dwell in the realm of existence. Fear God,
you who have insight, and do not deny. Beware lest mention of the
Prophet keep you from this Most Great Announcement, or guardianship from
the guardianship of God, which encompasses all the worlds. He has
created every name through His word and made every matter dependent on
His irrevocable, mighty, wondrous command. Say: This is the Day of God,
in which none is to be mentioned but His own Self, who holds sway over
all the worlds. This is a Cause that has shaken the imaginings and idols
you possess.

\subsection{169}\label{section-168}

\begin{otherlanguage}{persian}

قد نری منکم من یأخذ الکتاب و یستدلّ به علی اللّه کما استدلّت کلّ ملّة بکتابها
علی اللّه المهیمن القیوم قل تاللّه الحقّ لا تغنیکم الیوم کتب العالم و لا ما
فیه من الصّحف الاّ بهذا الکتاب الّذی ینطق فی قطب الابداع انّه لا اله الاّ انا
العلیم الحکیم

\end{otherlanguage}

Qad narā minkum man yaʾkhudhu al-kitāb wa-yastadillu bihi ʿalā Allāh
kamā istadallat kull milla bi-kitābihā ʿalā Allāh al-muhaymin al-qayyūm.
Qul ta-Allāh al-ḥaqq lā tughnīkum al-yawm kutub al-ʿālam wa-lā mā fīhi
min al-ṣuḥuf illā bi-hādhā al-kitāb alladhī yanṭiqu fī quṭb al-ibdāʿ
innahu lā ilāha illā anā al-ʿalīm al-ḥakīm.

We see among you those who take the Book and use it to argue against
God, just as every religious community has used its own Book to argue
against God, the All-Protecting, the Self-Subsisting. Say: By God, the
Truth! Today neither the books of the world nor any of its scriptures
can avail you except through this Book, which proclaims at the heart of
creation: ``There is no god but Me, the All-Knowing, the All-Wise.''

\subsection{170}\label{section-169}

\begin{otherlanguage}{arabic}

یا معشر العلمآء ایّاکم ان تکونوا سبب الاختلاف فی الأطراف کما کنتم علّة
الاعراض فی اوّل الأمر اجمعوا النّاس علی هذه الکلمة الّتی بها صاحت الحصاة
الملک للّه مطلع الآیات کذلک یعظکم اللّه فضلاً من عنده انّه لهو الغفور الکریم

\end{otherlanguage}

Yā maʿshar al-ʿulamāʾ iyyākum an takūnū sabab al-ikhtilāf fī al-aṭrāf
kamā kuntum ʿillat al-iʿrāḍ fī awwal al-amr. Ijmaʿū al-nās ʿalā hādhihi
al-kalima allatī bihā ṣāḥat al-ḥaṣāh al-mulk li-Allāh maṭlaʿ al-āyāt.
Kadhālika yaʿiẓukum Allāh faḍlan min ʿindihi innahu la-huwa al-ghafūr
al-karīm.

O company of the learned! Beware lest you become a cause of discord
throughout the lands, just as you were the cause of people turning away
at the beginning of the Cause. Gather the people around this word,
through which the pebbles have cried out: ``Sovereignty belongs to God,
the Dawning-Place of the verses.'' Thus God admonishes you as a bounty
from His presence. He is indeed the All-Forgiving, the All-Bountiful.

\subsection{171}\label{section-170}

\begin{otherlanguage}{persian}

اذکروا الکریم اذ دعوناه الی اللّه انّه استکبر بما اتّبع هواه بعد اذ ارسلنا
الیه ما قرّت به عین البرهان فی الامکان و تمّت حجّة اللّه علی من فی السّموات و
الأرضین انّا امرناه بالاقبال فضلا من الغنی المتعال انّه ولی مدبراً الی ان
اخذته زبانیة العذاب عدلا من اللّه انّا کنّا شاهدین

\end{otherlanguage}

Udhkurū al-Karīm idh daʿawnāhu ilā Allāh innahu istakbara bi-mā ittabaʿa
hawāhu baʿda idh arsalnā ilayhi mā qarrat bihi ʿayn al-burhān fī
al-imkān wa-tammat ḥujjat Allāh ʿalā man fī al-samāwāt wa-l-araḍīn. Innā
amarnāhu bi-l-iqbāl faḍlan min al-ghaniyy al-mutaʿāl. Innahu wallā
mudbiran ilā an akhadhat-hu zabāniyat al-ʿadhāb ʿadlan min Allāh innā
kunnā shāhidīn.

Remember Karīm when We called him to God. He grew arrogant, following
his own desires, after We had sent him that which gladdened the eye of
proof in the realm of existence and completed God's testimony to all who
dwell in the heavens and the earth. We commanded him to turn toward Us,
as a bounty from the All-Sufficient, the Most Exalted. He turned his
back until the wardens of torment seized him, as an act of justice from
God. We were indeed witnesses.

\subsection{172}\label{section-171}

\begin{otherlanguage}{arabic}

اخرقنّ الأحجاب علی شأن یسمع اهل الملکوت صوت خرقها هذا امر اللّه من قبل و
من بعد طوبی لمن عمل بما امر ویلٌ للتّارکین

\end{otherlanguage}

Ikhriqanna al-aḥjāb ʿalā shaʾn yasmaʿu ahl al-malakūt ṣawt kharqihā.
Hādhā amr Allāh min qabl wa-min baʿd. Ṭūbā li-man ʿamila bi-mā umira
waylun li-l-tārikīn.

Tear through the veils so that the inhabitants of the Kingdom hear the
sound of their tearing. This is God's command, before and hereafter.
Blessed is the one who acts as commanded; woe to those who neglect it.

\subsection{173}\label{section-172}

\begin{otherlanguage}{arabic}

انّا ما اردنا فی الملک الاّ ظهور اللّه و سلطانه و کفی باللّه علی شهیداً انّا
ما اردنا فی الملکوت الاّ علوّ امر اللّه و ثنائه و کفی باللّه علی وکیلا انّا
ما اردنا فی الجبروت الاّ ذکر اللّه و ما نزّل من عنده و کفی باللّه معینا

\end{otherlanguage}

Innā mā aradnā fī al-mulk illā ẓuhūr Allāh wa-sulṭānahu wa-kafā bi-Allāh
ʿalayya shahīdan. Innā mā aradnā fī al-malakūt illā ʿuluww amr Allāh
wa-thanāʾihi wa-kafā bi-Allāh ʿalayya wakīlan. Innā mā aradnā fī
al-jabarūt illā dhikr Allāh wa-mā nuzzila min ʿindihi wa-kafā bi-Allāh
muʿīnan.

In the earthly realm, We have desired only the manifestation of God and
His sovereignty, and God is sufficient as a witness over Me. In the
Kingdom, We have desired only the exaltation of God's Cause and His
praise, and God is sufficient as My guardian. In the realm of power, We
have desired only the remembrance of God and what has been sent down
from His presence, and God is sufficient as a helper.

\subsection{174}\label{section-173}

\begin{otherlanguage}{persian}

طوبی لکم یا معشر العلمآء فی البهآء تاللّه انتم امواج البحر الأعظم و انجم
سمآء الفضل و الویة النّصر بین السّموات و الأرضین انتم مطالع الاستقامة بین
البریة و مشارق البیان لمن فی الامکان طوبی لمن اقبل الیکم ویل للمعرضین
ینبغی الیوم لمن شرب رحیق الحیوان من ید الطاف ربّه الرّحمن ان یکون نبّاضاً
کالشّریان فی جسد الامکان لیتحرّک به العالم و کلّ عظم رمیم

\end{otherlanguage}

Ṭūbā lakum yā maʿshar al-ʿulamāʾ fī al-Bahāʾ. Ta-Allāh antum amwāj
al-baḥr al-aʿẓam wa-anjum samāʾ al-faḍl wa-alwiyat al-naṣr bayna
al-samāwāt wa-l-araḍīn. Antum maṭāliʿ al-istiqāma bayna al-bariyya
wa-mashāriq al-bayān li-man fī al-imkān. Ṭūbā li-man aqbala ilaykum
waylun li-l-muʿriḍīn. Yanbaghī al-yawm li-man shariba raḥīq al-ḥayawān
min yad alṭāf rabbihi al-Raḥmān an yakūna nabbāḍan ka-l-shiryān fī jasad
al-imkān li-yataḥarraka bihi al-ʿālam wa-kull ʿaẓm ramīm.

Blessed are you, O company of the learned in Bahā! By God, you are the
waves of the Most Great Ocean, the stars of the heaven of bounty, and
the banners of victory between the heavens and the earth. You are the
dawning-places of steadfastness among creation and the daysprings of
utterance for all who dwell in the realm of existence. Blessed is the
one who turns toward you; woe to those who turn away. Today, whoever has
drunk the choice wine of life from the hand of the loving kindness of
his Lord, the All-Merciful, should pulse like an artery in the body of
existence, so that through him the world and every crumbling bone may be
stirred to life.

\subsection{175}\label{section-174}

\begin{otherlanguage}{arabic}

یا اهل الانشآء اذا طارت الورقاء عن ایک الثّنآء و قصدت المقصد الأقصی
الأخفی ارجعوا ما لا عرفتموه من الکتاب الی الفرع المنشعب من هذا الأصل
القویم

\end{otherlanguage}

Yā ahl al-inshāʾ idhā ṭārat al-warqāʾ ʿan ayk al-thanāʾ wa-qaṣadat
al-maqṣad al-aqṣā al-akhfā irjiʿū mā lā ʿaraftumūhu min al-kitāb ilā
al-farʿ al-munshaʿib min hādhā al-aṣl al-qawīm.

O people of creation! When the Dove has flown from the grove of praise
and made its way to the farthest, most hidden destination, refer
whatever you have not understood in the Book to the Branch that has
sprung from this firmly established Root.

\subsection{176}\label{section-175}

\begin{otherlanguage}{arabic}

یا قلم الأعلی تحرّک علی اللّوح باذن ربّک فاطر السّمآء ثمّ اذکر اذ اراد مطلع
التّوحید مکتب التّجرید لعلّ الاحرار یطّلعنّ علی قدر سمّ الابرة بما هو خلف
الأستار من اسرار ربّک العزیز العلاّم قل انّا دخلنا مکتب المعانی و التّبیان
حین غفلة من فی الامکان و شاهدنا ما انزله الرّحمن و قبلنا ما اهداه لی من
آیات اللّه المهیمن القیوم و سمعنا ما شهد به فی اللّوح انّا کنّا شاهدین و
اجبناه بأمر من عندنا انّا کنّا آمرین

\end{otherlanguage}

Yā qalam al-aʿlā taḥarrak ʿalā al-lawḥ bi-idhn rabbika fāṭir al-samāʾ
thumma udhkur idh arāda maṭlaʿ al-tawḥīd maktab al-tajrīd laʿalla
al-aḥrār yaṭṭaliʿunna ʿalā qadr samm al-ibra bi-mā huwa khalf al-astār
min asrār rabbika al-ʿazīz al-ʿallām. Qul innā dakhalnā maktab al-maʿānī
wa-l-tibyān ḥīna ghaflat man fī al-imkān wa-shāhadnā mā anzalahu
al-Raḥmān wa-qabilnā mā ahdāhu lī min āyāt Allāh al-muhaymin al-qayyūm
wa-samiʿnā mā shahida bihi fī al-lawḥ innā kunnā shāhidīn wa-ajabnāhu
bi-amr min ʿindinā innā kunnā āmirīn.

O Most Exalted Pen! Move across the Tablet by permission of your Lord,
the Creator of heaven. Then recount when the Dawning-Place of Divine
Unity set out for the school of detachment, so that the free may
glimpse, to the extent of a needle's eye, the mysteries of your Lord,
the Almighty, the All-Knowing, that lie behind the veils. Say: We
entered the school of inner meanings and exposition while those in the
realm of existence were heedless. We beheld what the All-Merciful had
sent down and accepted what He offered Me as a gift of the verses of
God, the All-Protecting, the Self-Subsisting. We heard what He testified
to in the Tablet; We were indeed witnesses. And We answered Him with a
command from Our presence; We were indeed the One who commanded.

\subsection{177}\label{section-176}

\begin{otherlanguage}{arabic}

یا ملاء البیان انّا دخلنا مکتب اللّه اذ انتم راقدون و لاحظنا اللّوح اذ انتم
نائمون تاللّه الحقّ قد قرئناه قبل نزوله و انتم غافلون قد احطنا الکتاب اذ
کنتم فی الأصلاب هذا ذکری علی قدرکم لا علی قدر اللّه یشهد بذلک ما فی علم
اللّه لو انتم تعرفون و یشهد بذلک لسان اللّه لو انتم تفقهون تاللّه لو نکشف
الحجاب انتم تنصعقون

\end{otherlanguage}

Yā malaʾ al-bayān innā dakhalnā maktab Allāh idh antum rāqidūn
wa-lāḥaẓnā al-lawḥ idh antum nāʾimūn. Ta-Allāh al-ḥaqq qad qaraʾnāhu
qabla nuzūlihi wa-antum ghāfilūn. Qad aḥaṭnā al-kitāb idh kuntum fī
al-aṣlāb. Hādhā dhikrī ʿalā qadrikum lā ʿalā qadr Allāh yashhadu
bi-dhālika mā fī ʿilm Allāh law antum taʿrifūn wa-yashhadu bi-dhālika
lisān Allāh law antum tafqahūn. Ta-Allāh law nakshif al-ḥijāb antum
tanṣaʿiqūn.

O people of the Bayān! We entered the school of God while you lay
sleeping, and We examined the Tablet while you slumbered. By God, the
Truth! We read it before it was revealed, while you were heedless. We
encompassed the Book while you were still in the loins of your fathers.
This account of Mine is according to your capacity, not according to
God's measure. What is in the knowledge of God bears witness to this, if
only you knew, and the Tongue of God bears witness to it, if only you
understood. By God! Were We to lift the veil, you would be struck
senseless.

\subsection{178}\label{section-177}

\begin{otherlanguage}{arabic}

ایّاکم ان تجادلوا فی اللّه و امره انّه ظهر علی شأن احاط ما کان و ما یکون لو
نتکلّم فی هذا المقام بلسان اهل الملکوت لنقول قد خلق اللّه ذلک المکتب قبل
خلق السّموات و الأرض و دخلنا فیه قبل ان یقترن الکاف برکنها النّون هذا لسان
عبادی فی ملکوتی تفکّروا فیما ینطق به لسان اهل جبروتی بما علّمناهم علما من
لدنّا و ما کان مستورا فی علم اللّه و ما ینطق به لسان العظمة و الاقتدار فی
مقامه المحمود

\end{otherlanguage}

Iyyākum an tujādilū fī Allāh wa-amrihi innahu ẓahara ʿalā shaʾn aḥāṭa mā
kāna wa-mā yakūn. Law natakallam fī hādhā al-maqām bi-lisān ahl
al-malakūt la-naqūl qad khalaqa Allāh dhālika al-maktab qabla khalq
al-samāwāt wa-l-arḍ wa-dakhalnā fīhi qabla an yaqtarina al-kāf
bi-ruknihā al-nūn. Hādhā lisān ʿibādī fī malakūtī tafakkarū fīmā yanṭiqu
bihi lisān ahl jabarūtī bi-mā ʿallamnāhum ʿilman min ladunnā wa-mā kāna
mastūran fī ʿilm Allāh wa-mā yanṭiqu bihi lisān al-ʿaẓama wa-l-iqtidār
fī maqāmihi al-maḥmūd.

Beware of disputing about God and His Cause. He has appeared in such a
manner as to encompass all that has been and all that will be. Were We
to speak of this matter in the language of the inhabitants of the
Kingdom, We would say: God created that school before creating the
heavens and the earth, and We entered it before the Kāf was joined to
its counterpart, the Nūn. This is the language of My servants in My
Kingdom. Reflect on what the tongue of the inhabitants of My realm of
power utters through the knowledge We have taught them from Our
presence, on what has been concealed in God's knowledge, and on what the
Tongue of grandeur and power utters in its praiseworthy station.

\subsection{179}\label{section-178}

\begin{otherlanguage}{arabic}

لیس هذا امر تلعبون به بأوهامکم و لیس هذا مقام یدخل فیه کلّ جبان موهوم
تاللّه هذا مضمار المکاشفة و الانقطاع و میدان المشاهدة و الارتفاع لا یجول
فیه الاّ فوارس الرّحمن الّذین نبذوا الامکان اولئک انصار اللّه فی الأرض و
مشارق الاقتدار بین العالمین

\end{otherlanguage}

Laysa hādhā amr talʿabūna bihi bi-awhāmikum wa-laysa hādhā maqām
yadkhulu fīhi kull jabān mawhūm. Ta-Allāh hādhā miḍmār al-mukāshafa
wa-l-inqiṭāʿ wa-maydān al-mushāhada wa-l-irtifāʿ lā yajūlu fīhi illā
fawāris al-Raḥmān alladhīna nabadhū al-imkān. Ulāʾika anṣār Allāh fī
al-arḍ wa-mashāriq al-iqtidār bayna al-ʿālamīn.

This is not a Cause with which you may toy through your imaginings, nor
is this a station that every deluded coward may enter. By God! This is
the course of unveiling and detachment, the arena of vision and ascent.
None may ride within it but the horsemen of the All-Merciful who have
cast aside the realm of existence. They are God's helpers on earth and
the daysprings of power among the worlds.

\subsection{180}\label{section-179}

\begin{otherlanguage}{arabic}

ایّاکم ان یمنعکم ما فی البیان عن ربّکم الرّحمن تاللّه انّه قد نزّل لذکری لو
انتم تعرفون لا یجد منه المخلصون الاّ عرف حبّی و اسمی المهیمن علی کلّ شاهد و
مشهود قل یا قوم توجّهوا الی ما نزّل من قلمی الأعلی ان وجدتم منه عرف اللّه
لا تعترضوا علیه و لا تمنعوا انفسکم عن فضل اللّه و الطافه کذلک ینصحکم اللّه
انّه لهو النّاصح العلیم

\end{otherlanguage}

Iyyākum an yamnaʿakum mā fī al-bayān ʿan rabbikum al-Raḥmān. Ta-Allāh
innahu qad nuzzila li-dhikrī law antum taʿrifūn. Lā yajidu minhu
al-mukhliṣūn illā ʿarf ḥubbī wa-ismī al-muhaymin ʿalā kull shāhid
wa-mashhūd. Qul yā qawm tawajjahū ilā mā nuzzila min qalamī al-aʿlā in
wajadtum minhu ʿarf Allāh lā taʿtariḍū ʿalayhi wa-lā tamnaʿū anfusakum
ʿan faḍl Allāh wa-alṭāfihi. Kadhālika yanṣaḥukum Allāh innahu la-huwa
al-nāṣiḥ al-ʿalīm.

Beware lest what is in the Bayān keep you from your Lord, the
All-Merciful. By God! It was sent down to make mention of Me, if only
you knew. The sincere find in it only the fragrance of My love and of My
name, which holds sway over all who witness and all that is witnessed.
Say: O people! Turn toward what has been sent down from My Most Exalted
Pen. If you find in it the fragrance of God, do not object to it or
deprive yourselves of God's bounty and loving kindness. Thus God
counsels you. He is indeed the All-Knowing Counsellor.

\subsection{181}\label{section-180}

\begin{otherlanguage}{arabic}

ما لا عرفتموه من البیان فاسئلوا اللّه ربّکم و ربّ ابائکم الأوّلین انّه لو
یشآء یبین لکم ما نزّل فیه و ما ستر فی بحر کلماته من لآلیٴ العلم و الحکمة
انّه لهو المهیمن علی الأسمآء لا اله الاّ هو المهیمن القیوم

\end{otherlanguage}

Mā lā ʿaraftumūhu mina al-Bayāni fa-sʾalū Allāha rabbakum wa-rabba
ābāʾikumu al-awwalīna. Innahu law yashāʾu yubayyinu lakum mā nuzzila
fīhi wa-mā sutira fī baḥri kalimātihi min laʾāliʾi al-ʿilmi
wa-al-ḥikmati. Innahu la-huwa al-muhayminu ʿalā al-asmāʾi lā ilāha illā
huwa al-muhayminu al-qayyūmu.

Whatever you have not understood in the Bayān, ask God, your Lord and
the Lord of your forefathers. If He wishes, He will explain to you what
has been sent down in it and the pearls of knowledge and wisdom
concealed in the ocean of His words. He is indeed the Guardian over all
names. There is no god but Him, the All-Protecting, the Self-Subsisting.

\subsection{182}\label{section-181}

\begin{otherlanguage}{arabic}

قد اضطرب النّظم من هذا النّظم الأعظم و اختلف التّرتیب بهذا البدیع الّذی ما
شهدت عین الابداع شبهه

\end{otherlanguage}

Qadi iḍṭaraba al-naẓmu min hādhā al-naẓmi al-aʿẓami wa-ikhtalafa
al-tartību bi-hādhā al-badīʿi alladhī mā shahidat ʿaynu al-ibdāʿi
shibhahu.

The established order has been shaken by this Most Great Order, and its
arrangement has been transformed by this wondrous new order, the like of
which the eye of creation has never witnessed.

\subsection{183}\label{section-182}

\begin{otherlanguage}{arabic}

اغتمسوا فی بحر بیانی لعلّ تطّلعون بما فیه من لآلیٴ الحکمة و الأسرار ایّاکم
ان توقّفوا فی هذا الأمر الّذی به ظهرت سلطنة اللّه و اقتداره اسرعوا الیه
بوجوه بیضآء هذا دین اللّه من قبل و من بعد من اراد فلیقبل و من لم یرد فانّ
اللّه لغنی عن العالمین

\end{otherlanguage}

Ightamisū fī baḥri bayānī laʿalla taṭṭaliʿūna bi-mā fīhi min laʾāliʾi
al-ḥikmati wa-al-asrāri. Iyyākum an tawaqqafū fī hādhā al-amri alladhī
bihi ẓaharat salṭanatu Allāhi wa-iqtidāruhu. Asriʿū ilayhi bi-wujūhin
bayḍāʾa. Hādhā dīnu Allāhi min qablu wa-min baʿdu. Man arāda fa-līuqbil
wa-man lam yurid fa-inna Allāha la-ghaniyyun ʿani al-ʿālamīna.

Immerse yourselves in the ocean of My utterance, so that you may
discover the pearls of wisdom and mysteries within it. Beware of
hesitating in this Cause, through which God's sovereignty and might have
been revealed. Hasten to Him with radiant faces. This is the religion of
God, before and hereafter. Whoever wishes, let him turn toward it; and
whoever does not wish---God is indeed independent of all the worlds.

\subsection{184}\label{section-183}

\begin{otherlanguage}{persian}

قل هذا لقسطاس الهدی لمن فی السّموات و الأرض و البرهان الأعظم لو انتم
تعرفون قل به ثبت کلّ حجّة فی الأعصار لو انتم توقنون قل به استغنی کلّ فقیر و
تعلّم کلّ عالمٍ و عرج من اراد الصّعود الی اللّه ایّاکم ان تختلفوا فیه کونوا
کالجبال الرّواسخ فی امر ربّکم العزیز الودود

\end{otherlanguage}

Qul hādhā la-qisṭāsu al-hudā li-man fī al-samāwāti wa-al-arḍi
wa-al-burhānu al-aʿẓamu law antum taʿrifūna. Qul bihi thabata kullu
ḥujjatin fī al-aʿṣāri law antum tūqinūna. Qul bihi istaghnā kullu
faqīrin wa-taʿallama kullu ʿālimin wa-ʿaraja man arāda al-ṣuʿūda ilā
Allāhi. Iyyākum an takhtalifū fīhi. Kūnū ka-al-jibāli al-rawāsikhi fī
amri rabbikumu al-ʿazīzi al-wadūdi.

Say: This is indeed the balance of guidance for all who are in the
heavens and on earth, and the greatest proof, if only you knew. Say:
Through it every proof throughout the ages has been established, if only
you were certain. Say: Through it every poor person has been enriched,
every learned person has gained knowledge, and whoever has desired to
ascend to God has risen. Beware of differing over it. Be like firmly
rooted mountains in the Cause of your Lord, the Almighty, the
All-Loving.

\subsection{185}\label{section-184}

\begin{otherlanguage}{arabic}

قل یا مطلع الاعراض دع الاغماض ثمّ انطق بالحقّ بین الخلق تاللّه قد جرت دموعی
علی خدودی بما اراک مقبلاً الی هواک و معرضاً عمّن خلقک و سوّاک اذکر فضل مولاک
اذ ربّیناک فی اللّیالی و الأیّام لخدمة الأمر اتّق اللّه و کن من التّائبین هبنی
اشتبه علی النّاس امرک هل یشتبه علی نفسک خف عن اللّه ثمّ اذکر اذ کنت قائما
لدی العرش و کتبت ما القیناک من ایات اللّه المهیمن المقتدر القدیر ایّاک ان
تمنعک الحمیّة عن شطر الأحدیة توجّه الیه و لا تخف من اعمالک انّه یغفر من
یشآء بفضل من عنده لا اله الاّ هو الغفور الکریم انّما ننصحک لوجه اللّه ان
اقبلت فلنفسک و ان اعرضت انّ ربّک غنی عنک و عن الّذین اتّبعوک بوهم مبین قد
اخذ اللّه من اغواک فارجع الیه خاضعاً خاشعاً متذلّلاً انّه یکفّر عنک سیئاتک انّ
ربّک لهو التّوّاب العزیز الرّحیم

\end{otherlanguage}

Qul yā maṭlaʿa al-iʿrāḍi daʿi al-ighmāḍa thumma inṭiq bi-al-ḥaqqi bayna
al-khalqi. Ta-Allāhi qad jarat dumūʿī ʿalā khudūdī bi-mā arāka muqbilan
ilā hawāka wa-muʿriḍan ʿamman khalaqaka wa-sawwāka. Udhkur faḍla mawlāka
idh rabbaynāka fī al-layālī wa-al-ayyāmi li-khidmati al-amri. Ittaqi
Allāha wa-kun mina al-tāʾibīna. Habnī ishtabaha ʿalā al-nāsi amruka hal
yashtabihu ʿalā nafsika. Khaf ʿani Allāhi thumma udhkur idh kunta
qāʾiman ladā al-ʿarshi wa-katabta mā alqaynāka min āyāti Allāhi
al-muhaymini al-muqtadiri al-qadīri. Iyyāka an tamnaʿaka al-ḥamiyyatu
ʿan shaṭri al-aḥadiyyati. Tawajjah ilayhi wa-lā takhaf min aʿmālika.
Innahu yaghfiru man yashāʾu bi-faḍlin min ʿindihi lā ilāha illā huwa
al-ghafūru al-karīmu. Innamā nanṣaḥuka li-wajhi Allāhi. In aqbalta
fa-li-nafsika wa-in aʿraḍta inna rabbaka ghaniyyun ʿanka wa-ʿani
alladhīna ittabaʿūka bi-wahmin mubīnin. Qad akhadha Allāhu man aghwāka
fa-irjiʿ ilayhi khāḍiʿan khāshiʿan mutadhallilan. Innahu yukaffiru ʿanka
sayyiʾātika inna rabbaka la-huwa al-tawwābu al-ʿazīzu al-raḥīmu.

Say: O source of rejection! Give up your willful blindness, then speak
the truth among the people. By God! My tears have flowed down My cheeks
because I see you turning toward your own desires and away from Him who
created and fashioned you. Remember the bounty of your Lord, when We
nurtured you through the nights and days for service to the Cause. Fear
God and be among those who repent. Suppose your true condition is
unclear to others: is it unclear to yourself? Fear God, then remember
when you stood before the Throne and wrote down the verses of God, the
All-Protecting, the All-Powerful, the Almighty, that We imparted to you.
Beware lest pride keep you from turning toward the Divine Unity. Turn to
Him and do not fear because of your deeds. He forgives whomever He
wishes through His grace. There is no god but Him, the Ever-Forgiving,
the Generous. We counsel you only for the sake of God. If you turn
toward Him, it is for your own good; and if you turn away, your Lord has
no need of you or of those who have followed you in manifest delusion.
God has seized the one who led you astray, so return to Him in
submission, humility, and lowliness. He will indeed efface your
misdeeds. Your Lord is indeed the Ever-Relenting, the Almighty, the
Merciful.

\subsection{186}\label{section-185}

\begin{otherlanguage}{arabic}

هذا نصح اللّه لو انت من السّامعین هذا فضل اللّه لو انت من المقبلین هذا ذکر
اللّه لو انت من الشّاعرین هذا کنز اللّه لو انت من العارفین

\end{otherlanguage}

Hādhā nuṣḥu Allāhi law anta mina al-sāmiʿīna. Hādhā faḍlu Allāhi law
anta mina al-muqbilīna. Hādhā dhikru Allāhi law anta mina al-shāʿirīna.
Hādhā kanzu Allāhi law anta mina al-ʿārifīna.

This is God's counsel, if only you would listen. This is God's grace, if
only you would turn toward Him. This is God's remembrance, if only you
would perceive. This is God's treasure, if only you would recognize it.

\subsection{187}\label{section-186}

\begin{otherlanguage}{arabic}

هذا کتاب اصبح مصباح القدم للعالم و صراطه الأقوم بین العالمین قل انّه
لمطلع علم اللّه لو انتم تعلمون و مشرق اوامر اللّه لو انتم تعرفون

\end{otherlanguage}

Hādhā kitābun aṣbaḥa miṣbāḥa al-qidami li-al-ʿālami wa-ṣirāṭahu
al-aqwama bayna al-ʿālamīna. Qul innahu la-maṭlaʿu ʿilmi Allāhi law
antum taʿlamūna wa-mashriqu awāmiri Allāhi law antum taʿrifūna.

This is a Book that has become the lamp of eternity for the world and
His straightest path among all peoples. Say: It is indeed the dawning
place of God's knowledge, if only you knew, and the dayspring of God's
commandments, if only you recognized it.

\subsection{188}\label{section-187}

\begin{otherlanguage}{arabic}

لا تحملوا علی الحیوان ما یعجز عن حمله انّا نهیناکم عن ذلک نهیاً عظیماً فی
الکتاب کونوا مظاهر العدل و الانصاف بین السّموات و الأرضین

\end{otherlanguage}

Lā taḥmilū ʿalā al-ḥayawāni mā yaʿjizu ʿan ḥamlihi. Innā nahaynākum ʿan
dhālika nahyan ʿaẓīman fī al-kitābi. Kūnū maẓāhira al-ʿadli wa-al-inṣāfi
bayna al-samāwāti wa-al-arḍīna.

Do not load an animal with more than it can carry. We have most
emphatically forbidden you to do so in the Book. Be manifestations of
justice and fairness throughout the heavens and the earth.

\subsection{189}\label{section-188}

\begin{otherlanguage}{arabic}

من قتل نفساً خطأً فله دیة مسلّمة الی اهلها و هی مائة مثقال من الذّهب اعملوا
بما امرتم به فی اللّوح و لا تکوننّ من المتجاوزین

\end{otherlanguage}

Man qatala nafsan khaṭaʾan fa-lahu diyatun musallamatun ilā ahlihā
wa-hiya miʾatu mithqālin mina al-dhahabi. Iʿmalū bi-mā umirtum bihi fī
al-lawḥi wa-lā takūnanna mina al-mutajāwizīna.

Whoever kills a person by mistake must pay compensation to that person's
family, amounting to one hundred mithqāls of gold. Act according to what
you have been commanded in the Tablet, and do not be among those who
transgress.

\subsection{190}\label{section-189}

\begin{otherlanguage}{arabic}

یا اهل المجالس فی البلاد اختاروا لغة من اللّغات لیتکلّم بها من علی الأرض و
کذلک من الخطوط انّ اللّه یبین لکم ما ینفعکم و یغنیکم عن دونکم انّه لهو
الفضّال العلیم الخبیر هذا سبب الاتّحاد لو انتم تعلمون و العلة الکبری
للاتّفاق و التّمدّن لو انتم تشعرون انّا جعلنا الأمرین علامتین لبلوغ العالم
الأوّل و هو الاسّ الأعظم نزّلناه فی الواح اخری و الثّانی نزل فی هذا اللّوح
البدیع

\end{otherlanguage}

Yā ahla al-majālisi fī al-bilādi ikhtārū lughatan mina al-lughāti
li-yatakallama bihā man ʿalā al-arḍi wa-kadhālika mina al-khuṭūṭi. Inna
Allāha yubayyinu lakum mā yanfaʿukum wa-yughnīkum ʿan dūnikum. Innahu
la-huwa al-faḍḍālu al-ʿalīmu al-khabīru. Hādhā sababu al-ittiḥādi law
antum taʿlamūna wa-al-ʿillatu al-kubrā li-al-ittifāqi wa-al-tamadduni
law antum tashʿurūna. Innā jaʿalnā al-amrayni ʿalāmatayni li-bulūghi
al-ʿālami. Al-awwalu wa-huwa al-ussu al-aʿẓamu nazzalnāhu fī alwāḥin
ukhrā wa-al-thānī nazala fī hādhā al-lawḥi al-badīʿi.

O members of assemblies in the countries of the world! Choose one
language from among the languages for everyone on earth to speak, and
likewise one script from among the scripts. God makes clear to you what
will benefit you and make you independent of others. He is indeed the
Most Bountiful, the All-Knowing, the All-Informed. This is a means of
unity, if only you knew, and the greatest means of harmony and
civilization, if only you perceived. We have made these two matters
signs of the world's maturity. The first, which is the greatest
foundation, We have sent down in other Tablets, and the second has been
sent down in this wondrous Tablet.

\subsection{191}\label{section-190}

\begin{otherlanguage}{arabic}

قد حرّم علیکم شرب الأفیون انّا نهیناکم عن ذلک نهیاً عظیماً فی الکتاب و الّذی
شرب انّه لیس منّی اتّقوا اللّه یا اولی الألباب

\end{otherlanguage}

Qad ḥurrima ʿalaykum shurbu al-afyūni. Innā nahaynākum ʿan dhālika
nahyan ʿaẓīman fī al-kitābi wa-alladhī shariba innahu laysa minnī.
Ittaqū Allāha yā ulī al-albābi.

Consuming opium has been forbidden to you. We have most emphatically
forbidden you to do so in the Book, and whoever consumes it is not of
Me. Fear God, O people of understanding.

\bookmarksetup{startatroot}

\chapter{The Kitab-i-Iqan (The Book of
Certitude)}\label{the-kitab-i-iqan-the-book-of-certitude}

\section{BH00002}\label{bh00002}

\emph{Translated on 9 Qudrat, 183 BE}

\subsection{1}\label{section-191}

\begin{otherlanguage}{arabic}

بسم ربّنا العلیّ الأعلی

\end{otherlanguage}

Bismi Rabbinā al-ʿAliyyi al-Aʿlā

In the name of our Lord, the Exalted, the Most Exalted.

\subsection{2}\label{section-192}

\begin{otherlanguage}{arabic}

الباب المذکور فی بیان انّ العباد لن یصلوا الی شاطئ بحر العرفان الّا
بالانقطاع الصّرف عن کلّ من فی السّموات و الأرض قدّسوا انفسکم یا اهل الأرض
لعلّ تصلنّ الی المقام الّذی قدّر اللّه لکم و تدخلنّ فی سرادق جعله اللّه فی سمآء
البیان مرفوعاً

\end{otherlanguage}

Al-bāb al-madhkūr fī bayān anna al-ʿibād lan yaṣilū ilā shāṭiʾ baḥr
al-ʿirfān illā bi-l-inqiṭāʿ al-ṣirf ʿan kull man fī al-samāwāt wa-l-arḍ.
Qaddisū anfusakum yā ahl al-arḍ laʿalla taṣilunna ilā al-maqām alladhī
qaddara Allāh lakum wa-tadkhulunna fī surādiq jaʿalahu Allāh fī samāʾ
al-bayān marfūʿan.

This chapter explains that people will never reach the shore of the
ocean of recognition except through utter detachment from all who are in
the heavens and on earth. Sanctify yourselves, people of the earth, so
that you may reach the station God has ordained for you and enter a
pavilion that God has raised in the heaven of utterance.

\subsection{3}\label{section-193}

\begin{otherlanguage}{persian}

جوهر این باب آنکه سالکین سبیل ایمان و طالبین کؤوس ایقان باید نفوس خود را
از جمیع شئونات عرضیّه پاک و مقدّس نمایند یعنی گوش را از استماع اقوال و قلب
را از ظنونات متعلّقه بسبحات جلال و روح را از تعلّق باسباب ظاهره و چشم را
از ملاحظهٔ کلمات فانیه و متوکّلین علی اللّه و متوسّلین الیه سالک شوند تا
آنکه قابل تجلّیات اشراقات شموس علم و عرفان الهی و محلّ ظهورات فیوضات غیب
نامتناهی گردند زیرا اگر عبد بخواهد اقوال و اعمال و افعال عباد را از عالم
و جاهل میزان معرفت حقّ و اولیای او قرار دهد هرگز برضوان معرفت ربّ العزّة
داخل نشود و بعیون علم و حکمت سلطان احدیّت فائز نگردد و هرگز بسرمنزل بقا
نرسد و از جام قرب و رضا مرزوق نگردد

\end{otherlanguage}

Jawhar-i īn bāb ānkih sālikīn-i sabīl-i īmān va ṭālibīn-i kuʾūs-i īqān
bāyad nufūs-i khud rā az jamīʿ-i shuʾūnāt-i ʿaraḍiyyih pāk va muqaddas
namāyand, yaʿnī gūsh rā az istimāʿ-i aqvāl va qalb rā az ẓunūnāt-i
mutaʿalliqih bi-subuḥāt-i jalāl va rūḥ rā az taʿalluq bi-asbāb-i ẓāhirih
va chashm rā az mulāḥiẓiy-i kalimāt-i fāniyih, va mutavakkilīn ʿalā
Allāh va mutavassilīn ilayh sālik shavand tā ānkih qābil-i tajalliyāt-i
ishrāqāt-i shumūs-i ʿilm va ʿirfān-i ilāhī va maḥall-i ẓuhūrāt-i
fuyūḍāt-i ghayb-i nāmutanāhī gardand. Zīrā agar ʿabd bikhvāhad aqvāl va
aʿmāl va afʿāl-i ʿibād rā az ʿālim va jāhil mīzān-i maʿrifat-i Ḥaqq va
awliyā-yi ū qarār dihad, hargiz bi-Riḍvān-i maʿrifat-i Rabb al-ʿIzzat
dākhil nashavad va bi-ʿuyūn-i ʿilm va ḥikmat-i Sulṭān-i Aḥadiyyat fāʾiz
nagardad va hargiz bi-sar-manzil-i baqā narisad va az jām-i qurb va riḍā
marzūq nagardad.

The essence of this chapter is that those who tread the path of faith
and seek the cups of certitude must cleanse and sanctify themselves from
all contingent concerns: their ears from listening to sayings, their
hearts from conjectures concerning the veils of glory, their spirits
from attachment to outward means, and their eyes from attending to
perishable words. They must pursue their path trusting in God and
turning to Him, so that they may become receptive to the radiant
revelations of the suns of divine knowledge and recognition, and a place
where the outpourings of the infinite Unseen may appear. For if a person
wishes to make the sayings, deeds, and actions of others, whether
learned or ignorant, the measure of recognition of God and His chosen
ones, that person will never enter the paradise of recognition of the
Lord of Might, nor attain the wellsprings of the knowledge and wisdom of
the Sovereign of Oneness, nor ever reach the abiding home, nor receive
sustenance from the cup of nearness and good pleasure.

\subsection{4}\label{section-194}

\begin{otherlanguage}{persian}

ناظر بایّام قبل شوید که چه ‌قدر مردم از اعالی و ادانی همیشه منتظر ظهورات
احدیّه در هیاکل قدسیّه بوده‌اند بقسمی که در جمیع اوقات و اوان مترصّد و منتظر
و دعاها و تضرّعها مینمودند که شاید نسیم رحمت الهیّه بوزیدن آید و جمال
موعود از سرادق غیب بعرصهٔ ظهور قدم گذارد و چون ابواب عنایت مفتوح میگردید
و غمام مکرمت مرتفع و شمس غیب از افق قدرت ظاهر میشد جمیع تکذیب مینمودند و
از لقاء او که عین لقآء اللّه است احتراز میجستند چنانچه تفصیل آن در جمیع
کتب سماویّه مذکور و مسطور است

\end{otherlanguage}

Nāẓir bi-ayyām-i qabl shavīd kih chih qadr mardum az aʿālī va adānī
hamīshih muntaẓir-i ẓuhūrāt-i aḥadiyyih dar hayākil-i qudsiyyih
būdih-and, bi-qismī kih dar jamīʿ-i awqāt va avān mutaraṣṣid va muntaẓir
va duʿāhā va taḍarruʿhā mīnamūdand kih shāyad nasīm-i raḥmat-i ilāhiyyih
bi-vazīdan āyad va jamāl-i mawʿūd az surādiq-i ghayb bi-ʿarṣiy-i ẓuhūr
qadam gudhārad. Va chūn abvāb-i ʿināyat maftūḥ mīgardīd va ghamām-i
mukramat murtafiʿ va shams-i ghayb az ufuq-i qudrat ẓāhir mīshud, jamīʿ
takdhīb mīnamūdand va az liqā-yi ū kih ʿayn-i liqāʾ Allāh ast iḥtirāz
mījustand, chinānchih tafṣīl-i ān dar jamīʿ-i kutub-i samāviyyih madhkūr
va masṭūr ast.

Consider past ages: how many people, both high and low, were always
awaiting the manifestations of divine Oneness in holy human temples. At
every moment and in every season they watched and waited, offering
prayers and supplications that perhaps the breeze of divine mercy would
begin to blow and the promised Beauty would step from the pavilion of
the unseen into the realm of manifestation. Yet when the gates of grace
were opened, the cloud of bounty was raised, and the Sun of the unseen
appeared above the horizon of power, they all denied Him and shunned
meeting Him, although meeting Him is the very meeting with God. The
details of this are recorded and set down in all the heavenly books.

\subsection{5}\label{section-195}

\begin{otherlanguage}{persian}

حال قدری تأمّل نمائید که سبب اعتراض ناس بعد از طلب و آمال ایشان چه بود و
بقسمی هم اعتراض مینمودند که زبان و بیان و تقریر و تحریر همه از ذکر آن
عاجز و قاصر است و احدی از مظاهر قدسیّه و مطالع احدیّه ظاهر نشد مگر آنکه
باعتراض و انکار و احتجاج ناس مبتلا گشت چنانچه میفرماید یا حسرةً علی
العباد ما یأتیهم من رسول الّا کانوا به یستهزئون و در مقام دیگر میفرماید و
همّت کلّ امّة برسولهم لیأخذوه و جادلوا بالباطل لیدحضوا به الحقّ

\end{otherlanguage}

Ḥāl qadrī taʾammul namāyīd kih sabab-i iʿtirāḍ-i nās baʿd az ṭalab va
āmāl-i īshān chih būd, va bi-qismī ham iʿtirāḍ mīnamūdand kih zabān va
bayān va taqrīr va taḥrīr hamih az dhikr-i ān ʿājiz va qāṣir ast. Va
aḥadī az maẓāhir-i qudsiyyih va maṭāliʿ-i aḥadiyyih ẓāhir nashud magar
ānkih bi-iʿtirāḍ va inkār va iḥtijāj-i nās mubtalā gasht. Chinānchih
mīfarmāyad: Yā ḥasratan ʿalā al-ʿibād mā yaʾtīhim min rasūl illā kānū
bihi yastahziʾūn. Va dar maqām-i dīgar mīfarmāyad: Wa-hammat kull ummat
bi-rasūlihim li-yaʾkhudhūhu wa-jādalū bi-l-bāṭil li-yudḥiḍū bihi
al-ḥaqq.

Now reflect a little: what caused people to object after all their
seeking and hopes? They objected so fiercely that speech and expression,
spoken account and written record, are all powerless and inadequate to
describe it. Not one of the holy Manifestations or Dawning-Places of
Oneness appeared without being afflicted by the objections, denials, and
disputations of the people. As He says: ``Alas for the servants! No
messenger comes to them without their mocking him.'' And elsewhere He
says: ``Every nation plotted against its messenger to seize him, and
they disputed with falsehood to refute the truth with it.''

\subsection{6}\label{section-196}

\begin{otherlanguage}{persian}

و همچنین کلمات منزله که از غمام قدرت صمدانیّه و سماء عزّت ربّانیّه نازل شده
زیاده از حدّ احصا و احاطهٔ عباد است و اولو الأفئدة و صاحبان بصر را سورهٔ
هود کفایت میکند قدری در آن سورهٔ مبارکه تأمّل فرمائید و بفطرت اصلیّه تدبّر
نمائید تا قدری بر بدایع امور انبیا و رد و تکذیب کلمات نفی اطّلاع یابید
شاید ناس را از موطن غفلت نفسانیّه بآشیان وحدت و معرفت الهیّه پرواز دهید و
از زلال حکمت لایزال و اثمار شجرهٔ علم ذی الجلال بیاشامید و مرزوق گردید
اینست نصیب انفس مجرّده از مائدهٔ منزلهٔ قدسیّهٔ باقیه

\end{otherlanguage}

Va hamchunīn kalimāt-i munazzalih kih az ghamām-i qudrat-i ṣamadāniyyih
va samā-yi ʿizzat-i rabbāniyyih nāzil shudih ziyādih az ḥadd-i iḥṣā va
iḥāṭiy-i ʿibād ast, va ulū al-afʾidih va ṣāḥibān-i baṣar rā sūriy-i Hūd
kifāyat mīkunad. Qadrī dar ān sūriy-i mubārakih taʾammul farmāyīd va
bi-fiṭrat-i aṣliyyih tadabbur namāyīd tā qadrī bar badāyiʿ-i umūr-i
anbiyā va radd va takdhīb-i kalimāt-i nafy iṭṭilāʿ yābīd. Shāyad nās rā
az mawṭin-i ghaflat-i nafsāniyyih bi-āshiyān-i vaḥdat va maʿrifat-i
ilāhiyyih parvāz dihīd va az zulāl-i ḥikmat-i lā-yazāl va athmār-i
shajariy-i ʿilm-i Dhī al-Jalāl biyāshāmīd va marzūq gardīd. Īnast
naṣīb-i anfus-i mujarradih az māʾidiy-i munazzaliy-i qudsiyyiy-i
bāqiyih.

Likewise, the revealed words that have descended from the cloud of
eternal power and the heaven of divine might are beyond people's
capacity to count or encompass. For those endowed with hearts and
vision, the Surah of Hud suffices. Reflect a little on that blessed
surah and ponder it with your original, unspoiled nature, so that you
may gain some awareness of the wondrous affairs of the Prophets and the
rejection and denial voiced by the words of negation. Perhaps you may
enable people to soar from the abode of heedlessness of the self to the
nest of unity and divine recognition, and may drink the pure water of
everlasting wisdom and be nourished by the fruits of the tree of
knowledge of the Lord of Glory. This is the portion of detached souls
from the holy, enduring table sent down from heaven.

\subsection{7}\label{section-197}

\begin{otherlanguage}{persian}

اگر بر ابتلای انبیا و علّت و سبب اعتراضات عباد بر آن شموس هویّه آگاه شوید
بر اکثری از امور اطّلاع یابید و دیگر هر چه اعتراضات مردم را بر مشارق شموس
صفات احدیّه بیشتر ملاحظه کنید در دین خود و امر اللّه محکمتر و راسختر شوید
لهذا بعضی از حکایات انبیا مجملاً در این الواح ذکر میشود تا معلوم و مبرهن
آید که در جمیع اعصار و اقران بر مظاهر قدرت و مطالع عزّت وارد میآوردند
آنچه را که قلم از ذکرش خجل و منفعل است شاید این اذکار سبب شود که بعضی از
ناس از اعراض و اعتراض علما و جهّال ارض مضطرب نشوند و بلکه بر ایقان و
اطمینانشان بیفزاید

\end{otherlanguage}

Agar bar ibtilā-yi anbiyā va ʿillat va sabab-i iʿtirāḍāt-i ʿibād bar ān
shumūs-i huviyyih āgāh shavīd, bar aktharī az umūr iṭṭilāʿ yābīd. Va
dīgar har chih iʿtirāḍāt-i mardum rā bar mashāriq-i shumūs-i ṣifāt-i
aḥadiyyih bīshtar mulāḥiẓih kunīd, dar dīn-i khud va Amr Allāh muḥkamtar
va rāsikhtar shavīd. Lihādhā baʿḍī az ḥikāyāt-i anbiyā mujmalan dar īn
alvāḥ dhikr mīshavad tā maʿlūm va mubarhan āyad kih dar jamīʿ-i aʿṣār va
aqrān bar maẓāhir-i qudrat va maṭāliʿ-i ʿizzat vārid mīāvardand ānchih
rā kih qalam az dhikrash khajil va munfaʿil ast. Shāyad īn adhkār sabab
shavad kih baʿḍī az nās az iʿrāḍ va iʿtirāḍ-i ʿulamā va juhhāl-i arḍ
muḍṭarib nashavand va balkih bar īqān va iṭmīnānishān biyafzāyad.

If you become aware of the trials of the Prophets and the reasons and
causes of people's objections to those Suns of divine Being, you will
gain insight into many matters. Moreover, the more you consider people's
objections to the Dawning-Places of the suns of the attributes of
Oneness, the firmer and more steadfast you will become in your religion
and in the Cause of God. For this reason, some accounts of the Prophets
are briefly mentioned in these Tablets, so that it may become clear and
evident that in every age and era people inflicted upon the
Manifestations of power and Dawning-Places of might things that the pen
is ashamed and abashed to recount. Perhaps these accounts will keep some
people from being troubled by the rejection and objections of the
learned and ignorant of the earth, and will instead increase their
certitude and assurance.

\subsection{8}\label{section-198}

\begin{otherlanguage}{persian}

و از جملهٔ انبیا نوح بود که نهصد و پنجاه سال نوحه نمود و عباد را بوادی
ایمن روح دعوت فرمود و احدی او را اجابت ننمود و در هر یوم بقدری ایذا و
اذیّت بر آن وجود مبارک وارد میآوردند که یقین بر هلاکت او مینمودند و چه
مراتب سخریّه و استهزاء و کنایه که بر آن حضرت وارد شد چنانچه میفرماید و
کلّما مرّ علیه ملأ من قومه سخروا منه قال ان تسخروا منّا فانّا نسخر منکم کما
تسخرون فسوف تعلمون و بعد از مدّتها چند مرتبه وعدهٔ انزال نصر باصحاب خود
فرمودند بوعدۀ معیّن و در هر مرتبه بدا شد و بعضی از آن اصحاب معدوده بعلّت
ظهور بدا اعراض مینمودند چنانچه تفصیل آن در اکثر کتب مشهوره ثبت شده و
البتّه بنظر عالی رسیده یا میرسد تا آنکه باقی نماند از برای آن حضرت مگر
چهل نفس و یا هفتاد و دو نفس چنانچه در کتب و اخبار مذکور است تا آنکه
بالأخره نداء ربّ لا تذر علی الأرض من الکافرین دیّاراً از جان برکشید

\end{otherlanguage}

Va az jumliy-i anbiyā Nūḥ būd kih nuhṣad va panjāh sāl nawḥih namūd va
ʿibād rā bi-vādī-yi ayman-i rūḥ daʿvat farmūd va aḥadī ū rā ijābat
nanamūd. Va dar har yawm bi-qadrī īdhā va adhiyyat bar ān vujūd-i
mubārak vārid mīāvardand kih yaqīn bar halāk-i ū mīnamūdand. Va chih
marātib-i sukhriyyih va istihzā va kināyih kih bar ān Ḥaḍrat vārid shud.
Chinānchih mīfarmāyad: Wa-kullamā marra ʿalayhi malaʾ min qawmihi
sakhirū minhu qāla in taskharū minnā fa-innā naskhar minkum kamā
taskharūn fa-sawfa taʿlamūn. Va baʿd az muddathā chand martabih vaʿdiy-i
inzāl-i naṣr bi-aṣḥāb-i khud farmūdand bi-vaʿdiy-i muʿayyan, va dar har
martabih badā shud va baʿḍī az ān aṣḥāb-i maʿdūdih bi-ʿillat-i ẓuhūr-i
badā iʿrāḍ mīnamūdand. Chinānchih tafṣīl-i ān dar akthar-i kutub-i
mashhūrih thabt shudih va al-battih bi-naẓar-i ʿālī rasīdih yā mīrasad,
tā ānkih bāqī namānd az barā-yi ān Ḥaḍrat magar chihil nafs va yā haftād
va du nafs, chinānchih dar kutub va akhbār madhkūr ast. Tā ānkih
bi-l-ākharah nidā-yi Rabbi lā tadhar ʿalā al-arḍ min al-kāfirīn dayyāran
az jān bar-kashīd.

Among the Prophets was Noah, who lamented for nine hundred and fifty
years and called people to the blessed valley of the spirit, yet no one
answered Him. Each day they inflicted such harm and suffering upon that
blessed Being that they felt certain He would perish. What mockery,
ridicule, and taunts were directed against Him! As He says: ``Whenever a
company of His people passed by Him, they mocked Him. He said: `If you
mock us, we will indeed mock you as you mock us. Soon you will
know.'\,'' After long periods, He several times promised His companions
that help would be sent down at an appointed time. Each time a change in
the decree occurred, and some of those few companions turned away
because of the appearance of that change. The details are recorded in
most of the well-known books and have surely come, or will come, to your
attention. At last, only forty souls, or seventy-two souls, remained
with Him, as is mentioned in the books and traditions. Finally, from the
depths of His soul He raised the cry: ``My Lord, do not leave a single
unbeliever dwelling on the earth.''

\subsection{9}\label{section-199}

\begin{otherlanguage}{persian}

حال قدری تأمّل باید که سبب چه بود در این مدّت آن عباد باین قسم اعتراض
نمودند و احتراز جستند و از قمیص نفی بخلع اثبات مفتخر و فائز نشدند و دیگر
چرا در وعده‌های الهی بدا شد که سبب ادبار بعضی مقبلین شود بسیار تأمّل باید
تا بر اسرار امور غیبی واقف شوید و از طیب معنوی گلستان حقیقی بوئی برید و
تصدیق نمائید که امتحانات الهیّه همیشه در مابین عباد او بوده و خواهد بود
تا نور از ظلمت و صدق از کذب و حقّ از باطل و هدایت از ضلالت و سعادت از
شقاوت و خار از گل ممتاز و معلوم شود چنانچه فرموده الم أ حسب النّاس ان
یترکوا ان یقولوا آمنّا و هم لا یفتنون

\end{otherlanguage}

Ḥāl qadrī taʾammul bāyad kih sabab chih būd dar īn muddat ān ʿibād bi-īn
qism iʿtirāḍ namūdand va iḥtirāz justand va az qamīṣ-i nafy bi-khilaʿ-i
ithbāt muftakhar va fāʾiz nashudand. Va dīgar chirā dar vaʿdih-hā-yi
ilāhī badā shud kih sabab-i idbār-i baʿḍī muqbilīn shavad. Bisyār
taʾammul bāyad tā bar asrār-i umūr-i ghaybī vāqif shavīd va az ṭīb-i
maʿnavī-yi gulistān-i ḥaqīqī būʾī barīd va taṣdīq namāyīd kih
imtiḥānāt-i ilāhiyyih hamīshih dar mābayn-i ʿibād-i ū būdih va khvāhad
būd tā nūr az ẓulmat va ṣidq az kidhb va ḥaqq az bāṭil va hidāyat az
ḍalālat va saʿādat az shaqāvat va khār az gul mumtāz va maʿlūm shavad.
Chinānchih farmūdih: Alif Lām Mīm. A-ḥasiba al-nās an yutrakū an yaqūlū
āmannā wa-hum lā yuftanūn.

Now some reflection is needed: why did those people object and keep
their distance in this way throughout that period, and why did they not
exchange the garment of negation for the honor of attaining the robes of
affirmation? And why did a change occur in the divine promises, causing
some who had turned toward Him to turn away? Much reflection is needed,
so that you may become aware of the mysteries of unseen affairs, catch
the spiritual fragrance of the garden of truth, and acknowledge that
divine tests have always existed among His servants and always will, so
that light may be distinguished from darkness, truthfulness from
falsehood, truth from error, guidance from misguidance, felicity from
misery, and the thorn from the flower. As He has said: ``Alif. Lam. Mim.
Do people suppose they will be left to say, `We believe,' without being
tested?''

\subsection{10}\label{section-200}

\begin{otherlanguage}{persian}

و بعد از نوح جمال هود از مشرق ابداع مشرق شد و قریب هفتصد سنه او ازید
باختلاف اقوال مردم را برضوان قرب ذی الجلال دعوت نمود و چه مقدار بلایا که
بمثل غیث هاطل بر آن حضرت بارید تا آنکه کثرت دعوت سبب کثرت اعراض شد و شدّت
اهتمام علّت شدّت اغماض گردید و لا یزید الکافرین کفرهم الّا خسارا

\end{otherlanguage}

Va baʿd az Nūḥ jamāl-i Hūd az mashriq-i ibdāʿ mushriq shud va qarīb-i
haftṣad sanih aw azyad bi-ikhtilāf-i aqvāl mardum rā bi-Riḍvān-i qurb-i
Dhī al-Jalāl daʿvat namūd. Va chih miqdār balāyā kih bi-mithl-i ghayth-i
hāṭil bar ān Ḥaḍrat bārīd, tā ānkih kathrat-i daʿvat sabab-i kathrat-i
iʿrāḍ shud va shiddat-i ihtimām ʿillat-i shiddat-i ighmāḍ gardīd. Wa-lā
yazīdu al-kāfirīn kufruhum illā khasāran.

After Noah, the beauty of Hud shone from the dawning-place of creation.
For nearly seven hundred years or more, according to differing accounts,
He called people to the paradise of nearness to the Lord of Glory. How
many afflictions rained upon Him like a downpour! His repeated calls
only increased their rejection, and the intensity of His efforts only
intensified their disregard. ``Their unbelief increases the unbelievers
in nothing but loss.''

\subsection{11}\label{section-201}

\begin{otherlanguage}{persian}

و بعد هیکل صالحی از رضوان غیبی معنوی قدم بیرون نهاد و عباد را بشریعهٔ قرب
باقیه دعوت نمود و صد سنه او ازید امر باوامر الهی و نهی از مناهی میفرمود
ثمری نبخشید و اثری ظاهر نیامد و چند مرتبه غیبت اختیار فرمود با آنکه آن
جمال ازلی ناس را جز بمدینهٔ احدیّه دعوت نمینمود چنانچه میفرماید و الی ثمود
اخاهم صالحاً قال یا قوم اعبدوا اللّه ما لکم من اله غیره الی آخر القول
قالوا یا صالح قد کنت فینا مرجوّاً قبل هذا أ تنهانا ان نعبد ما یعبد آبائنا
و انّنا لفی شکّ ممّا تدعونا الیه مریب و هیچ فائده نبخشید تا آنکه بصیحه‌ئی
جمیع بنار راجع شدند

\end{otherlanguage}

Va baʿd haykal-i Ṣāliḥī az Riḍvān-i ghaybī-yi maʿnavī qadam bīrūn nihād
va ʿibād rā bi-sharīʿiy-i qurb-i bāqiyih daʿvat namūd va ṣad sanih aw
azyad amr bi-avāmir-i ilāhī va nahy az manāhī mīfarmūd. Thamarī
nabakhshīd va atharī ẓāhir nayāmad va chand martabih ghaybat ikhtiyār
farmūd, bā ānkih ān jamāl-i azalī nās rā juz bi-madīniy-i aḥadiyyih
daʿvat namīnamūd. Chinānchih mīfarmāyad: Wa-ilā Thamūd akhāhum Ṣāliḥan
qāla yā qawm iʿbudū Allāh mā lakum min ilāh ghayruhu, ilā ākhir al-qawl.
Qālū yā Ṣāliḥ qad kunta fīnā marjuwwan qabla hādhā a-tanhānā an naʿbuda
mā yaʿbudu ābāʾunā wa-innanā la-fī shakkin mimmā tadʿūnā ilayhi murīb.
Va hīch fāʾidih nabakhshīd tā ānkih bi-ṣayḥihʾī jamīʿ bi-nār rājiʿ
shudand.

After Him, the person of Salih stepped forth from the unseen, spiritual
paradise and called people to the waters of everlasting nearness. For a
hundred years or more, He enjoined the divine commandments and forbade
what was prohibited. It bore no fruit, and no effect appeared. Several
times He withdrew from sight, even though that eternal Beauty called
people to nothing but the city of Oneness. As He says: ``And to Thamud
We sent their brother Salih. He said: `My people, worship God; you have
no god other than Him'\,''---and so on to the end of the passage. ``They
said: `Salih, before this you were someone in whom we placed our hopes.
Do you forbid us to worship what our fathers worship? We are indeed in
troubling doubt about what you call us to.'\,'' Nothing availed, until,
at a single cry, they all returned to the fire.

\subsection{12}\label{section-202}

\begin{otherlanguage}{persian}

و بعد جمال خلیل کشف نقاب نمود و علم هدی مرتفع شد و اهل ارض را بنور تقی
دعوت فرموده هر چه مبالغه در نصیحت فرمود جز حسد ثمری نیاورد و غیر غفلت
حاصلی نبخشید الّا الّذینهم انقطعوا بکلّهم الی اللّه و عرجوا بجناحی الایقان
الی مقام جعله اللّه عن الادراک مرفوعاً و تفصیل آن حضرت مشهور است که چه
مقدار اعدا احاطه نمودند تا آنکه نار حسد و اعراض افروخته شد و بعد از
حکایت نار آن سراج الهی را از بلد اخراج نمودند چنانچه در همهٔ رسائل و کتب
مذکور است

\end{otherlanguage}

Va baʿd jamāl-i Khalīl kashf-i niqāb namūd va ʿalam-i hudā murtafiʿ shud
va ahl-i arḍ rā bi-nūr-i tuqā daʿvat farmūdih. Har chih mubālaghih dar
naṣīḥat farmūd juz ḥasad thamarī nayāvard va ghayr-i ghaflat ḥāṣilī
nabakhshīd, illā alladhīna hum inqaṭaʿū bi-kullihim ilā Allāh wa-ʿarajū
bi-janāḥay al-īqān ilā maqām jaʿalahu Allāh ʿan al-idrāk marfūʿan. Va
tafṣīl-i ān Ḥaḍrat mashhūr ast kih chih miqdār aʿdā iḥāṭih namūdand tā
ānkih nār-i ḥasad va iʿrāḍ afrūkhtih shud. Va baʿd az ḥikāyat-i nār ān
sirāj-i ilāhī rā az balad ikhrāj namūdand, chinānchih dar hamiy-i
rasāʾil va kutub madhkūr ast.

After Him, the beauty of the Friend unveiled Himself, the banner of
guidance was raised, and He called the people of the earth to the light
of piety. However earnestly He counseled them, it bore no fruit but envy
and yielded nothing but heedlessness, except among those who detached
themselves wholly for God and ascended on the wings of certitude to a
station God has raised beyond comprehension. The details of His story
are well known: how many enemies surrounded Him, until the fire of envy
and rejection was kindled. After the episode of the fire, they expelled
that divine Lamp from the city, as is mentioned in all the epistles and
books.

\subsection{13}\label{section-203}

\begin{otherlanguage}{persian}

و بعد زمان او منقضی شد تا نوبت به موسی رسید و آن حضرت بعصای امر و بیضای
معرفت از فاران محبّت الهیّه با ثعبان قدرت و شوکت صمدانیّه از سینای نور
بعرصهٔ ظهور ظاهر شد و جمیع من فی الملک را بملکوت بقا و اثمار شجرهٔ وفا
دعوت نمود و شنیده شد که فرعون و ملأ او چه اعتراضها که نمودند و چه مقدار
احجار ظنونات از انفس مشرکه بر آن شجرهٔ طیّبه وارد آمد تا بحدّی که فرعون و
ملأ او همّت گماشتند که آن نار سدرهٔ ربّانیّه را از ماء تکذیب و اعراض افسرده
و مخمود نمایند و غافل از اینکه نار حکمت الهیّه از آب عنصری افسرده نشود و
سراج قدرت ربّانیّه از بادهای مخالف خاموشی نپذیرد بلکه در این مقام ماء سبب
اشتعال شود و باد علّت حفظ لو انتم بالبصر الحدید تنظرون و فی رضی اللّه
تسلکون و چه بیانی خوش فرمود مؤمن آل فرعون چنانچه حکایت او را ربّ العزّة
برای حبیب خود میفرماید و قال رجل مؤمن من آل فرعون یکتم ایمانه أ تقتلون
رجلاً ان یقول ربّی اللّه و قد جآءکم بالبیّنات من ربّکم و ان یک کاذباً فعلیه
کذبه و ان یک صادقاً یصبکم بعض الّذی یعدکم انّ اللّه لا یهدی من هو مسرف کذّاب
و بالأخره امر بجائی کشید که همین مؤمن را بنهایت عذاب شهید نمودند الا
لعنة اللّه علی القوم الظّالمین

\end{otherlanguage}

Va baʿd zamān-i ū munqaḍī shud tā nawbat bi Mūsā rasīd va ān Ḥaḍrat
bi-ʿaṣā-yi amr va bayḍā-yi maʿrifat az Fārān-i maḥabbat-i ilāhiyya bā
thuʿbān-i qudrat va shawkat-i ṣamadāniyya az Sīnā-yi nūr bi-ʿarṣa-yi
ẓuhūr ẓāhir shud va jamīʿ-i man fī al-mulk rā bi-malakūt-i baqā va
athmār-i shajara-yi vafā daʿvat namūd va shinīda shud ki Firʿawn va
malaʾ-i ū chi iʿtirāḍ-hā ki namūdand va chi miqdār aḥjār-i ẓunūnāt az
anfus-i mushrika bar ān shajara-yi ṭayyiba vārid āmad tā bi-ḥaddī ki
Firʿawn va malaʾ-i ū himmat gumāshtand ki ān nār-i sidra-yi rabbāniyya
rā az māʾ-i takdhīb va iʿrāḍ afsurda va makhmūd namāyand va ghāfil az
īnki nār-i ḥikmat-i ilāhiyya az āb-i ʿunṣurī afsurda nashavad va sirāj-i
qudrat-i rabbāniyya az bād-hā-yi mukhālif khāmūshī napadhīrad balki dar
īn maqām māʾ sabab-i ishtiʿāl shavad va bād ʿillat-i ḥifẓ law antum
bi-l-baṣari al-ḥadīdi tanẓurūn wa fī riḍā Allāhi taslukūn va chi bayānī
khush farmūd muʾmin-i āl-i Firʿawn chunānchi ḥikāyat-i ū rā Rabb
al-ʿIzza barā-yi Ḥabīb-i khud mīfarmāyad wa qāla rajulun muʾminun min
āli Firʿawna yaktumu īmānahu a-taqtulūna rajulan an yaqūla rabbī Allāhu
wa qad jāʾakum bi-l-bayyināti min rabbikum wa in yaku kādhiban
fa-ʿalayhi kadhibuhu wa in yaku ṣādiqan yuṣibkum baʿḍu alladhī yaʿidukum
inna Allāha lā yahdī man huwa musrifun kadhdhāb va bi-l-ākhara amr
bi-jāʾī kashīd ki hamīn muʾmin rā bi-nihāyat-i ʿadhāb shahīd namūdand
alā laʿnatu Allāhi ʿalā al-qawmi al-ẓālimīn.

Afterward, His time came to an end, until the turn of Moses arrived.
With the staff of command and the white hand of knowledge, He came forth
from the Paran of divine love; with the serpent of power and eternal
majesty, He appeared from the Sinai of light in the arena of revelation.
He called all who dwell in the realm of creation to the kingdom of
everlasting life and the fruits of the tree of faithfulness. You have
heard what objections Pharaoh and his chiefs raised, and how many stones
of suspicion were hurled by idolatrous souls at that goodly Tree. It
reached the point where Pharaoh and his chiefs set themselves to quench
and extinguish the fire of that divine Lote-Tree with the water of
denial and rejection, unaware that the fire of divine wisdom cannot be
quenched by material water, nor can the lamp of divine power be
extinguished by contrary winds. Rather, in this station water causes it
to blaze and wind preserves it, if you look with penetrating sight and
walk in the way of God's good pleasure. How fine were the words spoken
by the believer of Pharaoh's household, whose story the Lord of Glory
recounts to His Beloved: ``A believing man from Pharaoh's household, who
concealed his faith, said: `Would you kill a man for saying, ``My Lord
is God,'' when he has brought you clear proofs from your Lord? If he is
lying, his lie will be held against him; but if he is truthful, some of
what he threatens you with will befall you. God does not guide anyone
who is a transgressor and a liar.'\,'' In the end, matters reached such
a point that they put this very believer to death with the utmost
torment. Indeed, God's curse rests upon the unjust.

\subsection{14}\label{section-204}

\begin{otherlanguage}{persian}

حال قدری در این امورات تأمّل فرمائید که چه سبب اینگونه اختلافات بوده که
هر ظهور حقّی که در امکان از افق لامکان ظاهر میشد اینگونه فساد و اغتشاش و
ظلم و انقلاب در اطراف عالم ظاهر و هویدا میگشت با اینکه جمیع انبیا در حین
ظهور خود مردم را بشارت میدادند به نبیّ بعد و علامتی از برای ظهور بعد ذکر
میفرمودند چنانچه در همهٔ کتب مسطور است با وجود طلب و انتظار ناس بمظاهر
قدسیّه و ذکر علامات در کتب چرا باید اینگونه امور در عالم رو دهد که جمیع
انبیا و اصفیا را در هر عهد و عصر اینگونه ظلم و جبر و تعدّی نمایند چنانچه
میفرماید أ فکلّما جآءکم رسول بما لا تهوی انفسکم استکبرتم ففریقاً کذّبتم و
فریقاً تقتلون میفرماید هر زمان و عهد که آمد بسوی شما رسولی از جانب
پروردگار بغیر هوای نفس شما تکبّر نمودید و موقن نشدید و گروهی از آن انبیا
را تکذیب نمودید و گروهی را میکشتید

\end{otherlanguage}

Ḥāl qadrī dar īn umūrāt taʾammul farmāʾīd ki chi sabab īn-gūna
ikhtilāfāt būda ki har ẓuhūr-i Ḥaqqī ki dar imkān az ufuq-i lā-makān
ẓāhir mīshud īn-gūna fasād va ightishāsh va ẓulm va inqilāb dar aṭrāf-i
ʿālam ẓāhir va huvaydā mīgasht bā īnki jamīʿ-i anbiyā dar ḥīn-i ẓuhūr-i
khud mardum rā bishārat mīdādand bi-nabī-yi baʿd va ʿalāmatī az barā-yi
ẓuhūr-i baʿd dhikr mīfarmūdand chunānchi dar hama-yi kutub masṭūr ast bā
vujūd-i ṭalab va intiẓār-i nās bi-maẓāhir-i qudsiyya va dhikr-i ʿalāmāt
dar kutub chirā bāyad īn-gūna umūr dar ʿālam rūy dihad ki jamīʿ-i anbiyā
va aṣfiyā rā dar har ʿahd va ʿaṣr īn-gūna ẓulm va jabr va taʿaddī
namāyand chunānchi mīfarmāyad a-fa-kullamā jāʾakum rasūlun bi-mā lā
tahwā anfusukumu istakbartum fa-farīqan kadhdhabtum wa farīqan taqtulūn
mīfarmāyad har zamān va ʿahd ki āmad bi-sū-yi shumā rasūlī az jānib-i
Parvardigār bi-ghayr-i havā-yi nafs-i shumā takabbur namūdīd va mūqin
nashudīd va gurūhī az ān anbiyā rā takdhīb namūdīd va gurūhī rā
mīkushtīd.

Now reflect a little on these matters: what caused such differences that
whenever a Manifestation of the True One appeared in the contingent
world from the horizon of the placeless, such corruption, turmoil,
oppression, and upheaval became manifest throughout the world? Yet all
the Prophets, at the time of Their appearance, gave people glad tidings
of the Prophet to come and mentioned a sign of the next revelation, as
is recorded in all the books. Given that people were seeking and
awaiting the holy Manifestations, and that the signs were mentioned in
the books, why should such things happen in the world---that in every
age and era all the Prophets and chosen Ones should suffer such
oppression, coercion, and aggression? As He says: ``Whenever a Messenger
came to you with what your souls did not desire, did you grow arrogant,
calling some impostors and killing others?'' He says that whenever, in
any age or era, a Messenger came to you from the Lord with something
contrary to your own desires, you grew arrogant and did not attain
certainty; some of those Prophets you called impostors, and others you
killed.

\subsection{15}\label{section-205}

\begin{otherlanguage}{persian}

آخر تأمّل فرمائید که سبب این افعال چه بود که باین قسم با طلعات جمال ذی
الجلال سلوک مینمودند و هر چه که در آن ازمنه سبب اعراض و اغماض آن عباد
بود حال هم سبب اغفال این عباد شده و اگر بگوئیم حجج الهیّه کامل و تمام
نبود لهذا سبب اعتراض عباد شد این کفریست صراح لأجل آنکه این بغایت از فیض
فیّاض دور است و از رحمت منبسطه بعید که نفسی را از میان جمیع عباد برگزیند
برای هدایت خلق خود و باو حجّت کافیهٔ وافیه عطا نفرماید و مع‌ذلک خلق را از
عدم اقبال باو معذّب فرماید بلکه لم‌یزل جود سلطان وجود بر همهٔ ممکنات بظهور
مظاهر نفس خود احاطه فرموده و آنی نیست که فیض او منقطع شود و یا آنکه
امطار رحمت از غمام عنایت او ممنوع گردد پس نیست این امورات محدثه مگر از
انفس محدوده که در وادی کبر و غرور حرکت مینمایند و در صحراهای بعد سیر
مینمایند و بظنونات خود و هر چه از علمای خود شنیده‌اند همان را تأسّی
مینمایند لهذا غیر از اعراض امری ندارند و جز اغماض حاصلی نخواهند و این
معلوم است نزد هر ذی بصری که اگر این عباد در ظهور هر یک از مظاهر شمس
حقیقت چشم و گوش و قلب را از آنچه دیده و شنیده و ادراک نموده پاک و مقدّس
مینمودند البتّه از جمال الهی محروم نمیماندند و از حرم قرب و وصال مطالع
قدسیّه ممنوع نمیگشتند و چون در هر زمان حجّت را بمعرفت خود که از علمای خود
شنیده بودند میزان مینمودند و بعقول ضعیفهٔ آنها موافق نمیآمد لهذا از
اینگونه امور غیر مرضیّه از ایشان در عالم ظهور بظهور میآمد

\end{otherlanguage}

Ākhir taʾammul farmāʾīd ki sabab-i īn afʿāl chi būd ki bi-īn qism bā
ṭalʿāt-i jamāl-i Dhī al-Jalāl sulūk mīnamūdand va har chi ki dar ān
azmina sabab-i iʿrāḍ va ighmāḍ-i ān ʿibād būd ḥāl ham sabab-i ighfāl-i
īn ʿibād shuda va agar bigūʾīm ḥujaj-i ilāhiyya kāmil va tamām nabūd
li-hādhā sabab-i iʿtirāḍ-i ʿibād shud īn kufrīst ṣurāḥ li-ajl-i ānki īn
bi-ghāyat az fayḍ-i Fayyāḍ dūr ast va az raḥmat-i munbasiṭa baʿīd ki
nafsī rā az miyān-i jamīʿ-i ʿibād barguzīnad barā-yi hidāyat-i khalq-i
khud va bi-ū ḥujjat-i kāfiya-yi vāfiya ʿaṭā nafarmāyad va maʿa-dhālik
khalq rā az ʿadam-i iqbāl bi-ū muʿadhdhab farmāyad balki lam-yazal jūd-i
Sulṭān-i vujūd bar hama-yi mumkināt bi-ẓuhūr-i maẓāhir-i nafs-i khud
iḥāṭa farmūda va ānī nīst ki fayḍ-i ū munqaṭiʿ shavad va yā ānki amṭār-i
raḥmat az ghamām-i ʿināyat-i ū mamnūʿ gardad pas nīst īn umūrāt-i
muḥdatha magar az anfus-i maḥdūda ki dar vādī-yi kibr va ghurūr ḥarakat
mīnamāyand va dar ṣaḥrā-hā-yi buʿd sayr mīnamāyand va bi-ẓunūnāt-i khud
va har chi az ʿulamā-yi khud shinīda-and hamān rā taʾassī mīnamāyand
li-hādhā ghayr az iʿrāḍ amrī nadārand va juz ighmāḍ ḥāṣilī nakhvāhand va
īn maʿlūm ast nazd-i har dhī baṣarī ki agar īn ʿibād dar ẓuhūr-i har yak
az maẓāhir-i Shams-i Ḥaqīqat chashm va gūsh va qalb rā az ānchi dīda va
shinīda va idrāk namūda pāk va muqaddas mīnamūdand albatta az jamāl-i
ilāhī maḥrūm namīmāndand va az ḥaram-i qurb va viṣāl-i maṭāliʿ-i
qudsiyya mamnūʿ namīgashtand va chūn dar har zamān ḥujjat rā
bi-maʿrifat-i khud ki az ʿulamā-yi khud shinīda būdand mīzān mīnamūdand
va bi-ʿuqūl-i ḍaʿīfa-yi ānhā muvāfiq namīāmad li-hādhā az īn-gūna umūr-i
ghayr-i marḍiyya az īshān dar ʿālam-i ẓuhūr bi-ẓuhūr mīāmad.

Consider, then, what caused these actions, that they should treat the
Countenances of the beauty of the All-Glorious in this way. Whatever
caused those people to turn away and close their eyes in those times has
now caused these people to be deceived as well. If we say that the
divine proofs were not complete and perfect, and therefore gave rise to
people's objections, this is blatant unbelief. For it is utterly
incompatible with the bounty of the All-Bountiful and far removed from
His all-embracing mercy that He should choose a soul from among all His
servants to guide His creation, yet fail to bestow on Him sufficient and
adequate proof, and then punish His creatures for not turning toward
Him. Rather, the generosity of the Sovereign of existence has always
encompassed all contingent beings through the appearance of the
Manifestations of His own Self. There is not a moment when His grace is
cut off or the rains of mercy are withheld from the cloud of His loving
care. These occurrences arise, then, only from narrow-minded souls who
walk in the valley of pride and arrogance and roam the deserts of
remoteness, following their own suppositions and whatever they have
heard from their religious scholars. Thus they do nothing but turn away
and seek no outcome but to close their eyes. It is clear to everyone
with discernment that if, at the appearance of each of the
Manifestations of the Sun of Truth, these people had cleansed and
sanctified their eyes, ears, and hearts from whatever they had seen,
heard, and understood, they would certainly not have been deprived of
the divine Beauty or barred from the sanctuary of nearness to and
reunion with the holy Daysprings. But in every age they weighed the
proof against their own knowledge, which they had heard from their
religious scholars; and because it did not accord with their feeble
minds, such unworthy acts appeared from them in the world of
manifestation.

\subsection{16}\label{section-206}

\begin{otherlanguage}{persian}

و در همهٔ اوقات سبب صدّ عباد و منع ایشان از شاطی بحر احدیّه علمای عصر
بوده‌اند که زمام آن مردم در کفّ کفایت ایشان بود و ایشان هم بعضی نظر بحبّ
ریاست و بعضی از عدم علم و معرفت ناس را منع مینمودند چنانچه همهٔ انبیا
باذن و اجازهٔ علمای عصر سلسبیل شهادت را نوشیدند و باعلی افق عزّت پرواز
نمودند چه ظلمها که از رؤسای عهد و علمای عصر بر سلاطین وجود و جواهر مقصود
وارد شد و باین ایّام محدودهٔ فانیه قانع شدند و از ملک لایفنی بازماندند
چنانچه چشم را از مشاهدهٔ انوار جمال محبوب بی‌نصیب نمودند و گوش را از بدایع
نغمات ورقاء مقصود محروم ساختند اینست که در جمیع کتب سماویّه ذکر احوال
علمای هر عصر شده چنانچه میفرماید یا اهل الکتاب لم تکفرون بآیات اللّه و
انتم تشهدون و همچنین میفرماید یا اهل الکتاب لم تلبسون الحقّ بالباطل و
تکتمون الحقّ و انتم تعلمون و در مقام دیگر میفرماید قل یا اهل الکتاب لم
تصدّون عن سبیل اللّه و این معلوم است که اهل کتابی که صدّ نمودند مردم را از
صراط مستقیم علمای آن عهد بوده‌اند چنانچه اسم و رسم جمیع در کتب مذکور است
و از اکثر آیات و اخبار مستفاد میشود لو انتم بطرف اللّه تنظرون

\end{otherlanguage}

Va dar hama-yi awqāt sabab-i ṣadd-i ʿibād va manʿ-i īshān az shāṭī-yi
baḥr-i aḥadiyya ʿulamā-yi ʿaṣr būda-and ki zimām-i ān mardum dar kaff-i
kifāyat-i īshān būd va īshān ham baʿḍī naẓar bi-ḥubb-i riyāsat va baʿḍī
az ʿadam-i ʿilm va maʿrifat nās rā manʿ mīnamūdand chunānchi hama-yi
anbiyā bi-idhn va ijāza-yi ʿulamā-yi ʿaṣr salsabīl-i shahādat rā
nūshīdand va bi-aʿlā ufuq-i ʿizzat parvāz namūdand chi ẓulm-hā ki az
ruʾasā-yi ʿahd va ʿulamā-yi ʿaṣr bar salāṭīn-i vujūd va javāhir-i maqṣūd
vārid shud va bi-īn ayyām-i maḥdūda-yi fāniya qāniʿ shudand va az mulk-i
lā-yafnā bāzmāndand chunānchi chashm rā az mushāhada-yi anvār-i jamāl-i
Maḥbūb bī-naṣīb namūdand va gūsh rā az badāyiʿ-i naghamāt-i varqā-yi
Maqṣūd maḥrūm sākhtand īnast ki dar jamīʿ-i kutub-i samāviyya dhikr-i
aḥvāl-i ʿulamā-yi har ʿaṣr shuda chunānchi mīfarmāyad yā ahla al-kitābi
lima takfurūna bi-āyāti Allāhi wa antum tashhadūn va hamchunīn
mīfarmāyad yā ahla al-kitābi lima talbisūna al-ḥaqqa bi-l-bāṭili wa
taktumūna al-ḥaqqa wa antum taʿlamūn va dar maqām-i dīgar mīfarmāyad qul
yā ahla al-kitābi lima taṣuddūna ʿan sabīli Allāh va īn maʿlūm ast ki
ahl-i kitābī ki ṣadd namūdand mardum rā az ṣirāṭ-i mustaqīm ʿulamā-yi ān
ʿahd būda-and chunānchi ism va rasm-i jamīʿ dar kutub madhkūr ast va az
akthar-i āyāt va akhbār mustafād mīshavad law antum bi-ṭarfi Allāhi
tanẓurūn.

At all times, it has been the religious scholars of the age, in whose
capable hands lay the reins of the people, who have turned them away and
kept them from the shore of the ocean of oneness. Some did so out of
love of leadership, and others through lack of knowledge and
understanding. Thus all the Prophets drank the life-giving waters of
martyrdom with the sanction and permission of the religious scholars of
Their age, and soared to the highest horizon of glory. What oppression
the leaders and religious scholars of each age inflicted upon the
Sovereigns of existence and the Gems of all desire! They contented
themselves with these few fleeting days and were deprived of the
imperishable kingdom. They denied their eyes the sight of the radiance
of the Beloved's beauty and deprived their ears of the wondrous melodies
of the Dove of the Desired One. For this reason, the circumstances of
the religious scholars of every age are mentioned in all the heavenly
books. As He says: ``People of the Book, why do you disbelieve in God's
signs when you yourselves bear witness?'' And likewise He says: ``People
of the Book, why do you clothe truth in falsehood and conceal the truth
when you know?'' Elsewhere He says: ``Say: People of the Book, why do
you turn others away from the path of God?'' It is clear that the people
of the Book who turned people away from the straight path were the
religious scholars of that age. The names and particulars of them all
are recorded in the books, and this can be understood from most of the
verses and traditions, if you look with the eye of God.

\subsection{17}\label{section-207}

\begin{otherlanguage}{persian}

پس قدری بدیدهٔ بصیرت الهیّه در آفاق علم ربّانی و انفس کلمات تامّهٔ صمدانیّه
تعقّل فرمائید تا جمیع اسرار حکمت روحانیّه بی‌سبحات جلال از خلف سرادق فضل و
افضال ظاهر و هویدا شود و کلّیّهٔ اعتراض مردم و احتجاجات ایشان از عدم ادراک
و عرفان حاصل شده مثلاً بیاناتی که طلعات جمال حقّ در علامات ظهور بعد
فرمودند آن بیانات را ادراک ننمودند و بحقیقت آن واصل نشدند لهذا علم فساد
برافراختند و رایات فتنه برپا نمودند و این معلومست که تأویل کلمات حمامات
ازلیّه را جز هیاکل ازلیّه ادراک ننمایند و نغمات ورقاء معنویّه را جز سامعهٔ
اهل بقا نشنود و هرگز قبطی ظلم از شراب سبطی عدل نصیب ندارد و فرعون کفر از
بیضای موسی اطّلاع نیابد چنانچه میفرماید و ما یعلم تأویله الّا اللّه و
الرّاسخون فی العلم مع‌ذلک تأویل کتاب را از اهل حجاب مستفسر شدند و علم را
از منبع او اخذ ننمودند

\end{otherlanguage}

Pas qadrī bi-dīda-yi baṣīrat-i ilāhiyya dar āfāq-i ʿilm-i rabbānī va
anfus-i kalimāt-i tāmma-yi ṣamadāniyya taʿaqqul farmāʾīd tā jamīʿ-i
asrār-i ḥikmat-i rūḥāniyya bī-subuḥāt-i jalāl az khalf-i surādiq-i faḍl
va ifḍāl ẓāhir va huvaydā shavad va kulliyya-yi iʿtirāḍ-i mardum va
iḥtijājāt-i īshān az ʿadam-i idrāk va ʿirfān ḥāṣil shuda mathalan
bayānātī ki ṭalʿāt-i jamāl-i Ḥaqq dar ʿalāmāt-i ẓuhūr-i baʿd farmūdand
ān bayānāt rā idrāk nanamūdand va bi-ḥaqīqat-i ān vāṣil nashudand
li-hādhā ʿalam-i fasād barafrākhtand va rāyāt-i fitna bar-pā namūdand va
īn maʿlūmast ki taʾvīl-i kalimāt-i ḥamāmāt-i azaliyya rā juz hayākil-i
azaliyya idrāk nanamāyand va naghamāt-i varqā-yi maʿnaviyya rā juz
sāmiʿa-yi ahl-i baqā nashnavad va hargiz qibṭī-yi ẓulm az sharāb-i
sibṭī-yi ʿadl naṣīb nadārad va Firʿawn-i kufr az bayḍā-yi Mūsā iṭṭilāʿ
nayābad chunānchi mīfarmāyad wa mā yaʿlamu taʾwīlahu illā Allāhu wa
al-rāsikhūna fī al-ʿilm maʿa-dhālik taʾvīl-i kitāb rā az ahl-i ḥijāb
mustafsir shudand va ʿilm rā az manbaʿ-i ū akhdh nanamūdand.

Reflect, then, with the eye of divine insight on the horizons of
heavenly knowledge and the inner realities of the perfect, eternal
Words, so that all the mysteries of spiritual wisdom may appear from
behind the pavilion of grace and bounty, free of the veils of glory. All
the people's objections and disputes have arisen from a lack of
understanding and recognition. For example, they did not understand the
statements that the Countenances of the beauty of the True One made
concerning the signs of the next revelation, nor did they grasp their
true meaning. Thus they raised the standard of corruption and unfurled
the banners of discord. It is clear that none but the eternal Temples
can understand the interpretation of the words of the eternal Doves, and
none but the ears of the people of everlasting life can hear the
melodies of the spiritual Dove. The Copt of oppression can never have a
share of the wine of the Israelite of justice, nor can the Pharaoh of
unbelief gain any knowledge of the white hand of Moses. As He says:
``None knows its interpretation except God and those firmly grounded in
knowledge.'' Nevertheless, they sought the interpretation of the Book
from those veiled from the truth and did not take knowledge from its
source.

\subsection{18}\label{section-208}

\begin{otherlanguage}{persian}

مثلاً چون ایّام موسی گذشت و انوار عیسی از فجر روح عالم را احاطه نمود جمیع
یهود اعتراض نمودند که آن نفس که در تورات موعود است باید مروّج و مکمّل
شرایع تورات باشد و این جوان ناصری که خود را مسیح اللّه مینامد حکم طلاق و
سبت را که از حکمهای اعظم موسی است نسخ نموده و دیگر آنکه علائم ظهور هنوز
ظاهر نشده چنانچه یهود هنوز منتظر آن ظهورند که در تورات مذکور است چه ‌قدر
از مظاهر قدس احدیّه و مطالع نور ازلیّه که بعد از موسی در ابداع ظاهر شده و
هنوز یهود بحجبات نفسیّهٔ شیطانیّه و ظنونات افکیّهٔ نفسانیّه محتجب بوده و هستند
و منتظرند که هیکل مجعول با علامات مذکوره که خود ادراک نموده‌اند کی ظاهر
خواهد شد کذلک اخذهم اللّه بذنبهم و اخذ عنهم روح الایمان و عذّبهم بنار کانت
فی هاویة الجحیم و این نبود مگر از عدم عرفان یهود عبارات مسطورۀ در تورات
را که در علائم ظهور بعد نوشته شده چون بحقیقت آن پی‌نبردند و بظاهر هم چنین
امور واقع نشد لهذا از جمال عیسوی محروم شدند و بلقآء اللّه فائز نگشتند و
کانوا من المنتظرین و لم‌یزل و لایزال جمیع امم بهمین جعلیّات افکار نالایقه
تمسّک جسته و از عیونهای لطیفهٔ رقیقهٔ جاریه خود را بی‌بهره و بی‌نصیب نمودند

\end{otherlanguage}

Mathalan chūn ayyām-i Mūsā gudhasht va anvār-i ʿĪsā az fajr-i rūḥ ʿālam
rā iḥāṭa namūd jamīʿ-i Yahūd iʿtirāḍ namūdand ki ān nafs ki dar Tawrāt
mawʿūd ast bāyad muravvij va mukammil-i sharāyiʿ-i Tawrāt bāshad va īn
javān-i Nāṣirī ki khud rā Masīḥ Allāh mīnāmad ḥukm-i ṭalāq va sabt rā ki
az ḥukm-hā-yi aʿẓam-i Mūsā ast naskh namūda va dīgar ānki ʿalāʾim-i
ẓuhūr hanūz ẓāhir nashuda chunānchi Yahūd hanūz muntaẓir-i ān ẓuhūrand
ki dar Tawrāt madhkūr ast chi qadr az maẓāhir-i quds-i aḥadiyya va
maṭāliʿ-i nūr-i azaliyya ki baʿd az Mūsā dar ibdāʿ ẓāhir shuda va hanūz
Yahūd bi-ḥujubāt-i nafsiyya-yi shayṭāniyya va ẓunūnāt-i ifkiyya-yi
nafsāniyya muḥtajib būda va hastand va muntaẓirand ki haykal-i majʿūl bā
ʿalāmāt-i madhkūra ki khud idrāk namūda-and kay ẓāhir khvāhad shud
kadhālika akhadhahumu Allāhu bi-dhanbihim wa akhadha ʿanhum rūḥa
al-īmāni wa ʿadhdhabahum bi-nārin kānat fī hāwiyati al-jaḥīm va īn nabūd
magar az ʿadam-i ʿirfān-i Yahūd ʿibārāt-i masṭūra dar Tawrāt rā ki dar
ʿalāʾim-i ẓuhūr-i baʿd nivishta shuda chūn bi-ḥaqīqat-i ān pay-naburdand
va bi-ẓāhir ham chunīn umūr vāqiʿ nashud li-hādhā az jamāl-i ʿĪsavī
maḥrūm shudand va bi-liqāʾ Allāh fāʾiz nagashtand wa kānū mina
al-muntaẓirīn va lam-yazal va lā-yazāl jamīʿ-i umam bi-hamīn jaʿliyyāt-i
afkār-i nālāʾiqa tamassuk justa va az ʿuyūn-hā-yi laṭīfa-yi raqīqa-yi
jāriya khud rā bī-bahra va bī-naṣīb namūdand.

For example, when the days of Moses had passed and the lights of Jesus
shone from the dawn of the spirit and encompassed the world, all the
Jews objected: ``The person promised in the Torah must promote and
complete the laws of the Torah. Yet this young man of Nazareth, who
calls himself the Messiah of God, has abrogated the laws of divorce and
the Sabbath, which are among the greatest laws of Moses. Moreover, the
signs of the revelation have not yet appeared.'' Thus the Jews are still
awaiting the revelation mentioned in the Torah. How many Manifestations
of the holiness of oneness and Daysprings of eternal light have appeared
in creation since Moses! Yet the Jews have remained, and still remain,
concealed behind satanic veils of self and false, selfish suppositions,
waiting to see when their imagined figure will appear with the stated
signs as they themselves have understood them. Thus God seized them for
their sin, took from them the spirit of faith, and tormented them with a
fire in the abyss of hell. This was due only to the Jews' failure to
understand the passages recorded in the Torah concerning the signs of
the next revelation. Since they did not grasp their true meaning, and
such things did not occur outwardly either, they were deprived of the
beauty of Jesus and did not attain the presence of God, and they
remained among those who wait. All peoples have always clung, and
continue to cling, to these same fabrications of unworthy thoughts,
depriving themselves of any share in the pure and gentle flowing
springs.

\subsection{19}\label{section-209}

\begin{otherlanguage}{persian}

و در کشف این اسرار بعضی از عبارات انبیا ببدایع نغمات حجازی در الواح
مسطورهٔ قبل که برای یکی از احباب نوشته شده بود مذکور گشت و حال هم بتغنّیات
خوش عراقی نظر بخواهش آن جناب در این اوراق مجدّداً ذکر مینمائیم که شاید
تشنگان صحراهای بعد را ببحر قرب دلالت نماید و گمگشتگان بیابانهای هجر و
فراق را بخیام قرب و وصال رساند تا غمام ضلالت مرتفع شود و آفتاب جهانتاب
هدایت از افق جان طالع گردد و علی اللّه اتّکل و به استعین لعلّ یجری من هذا
القلم ما یحیی به افئدة النّاس لیقومنّ الکلّ عن مراقد غفلتهم و یسمعنّ اطوار
ورقات الفردوس من شجر کان فی الرّوضة الأحدیّة من ایدی القدرة باذن اللّه
مغروساً

\end{otherlanguage}

Va dar kashf-i īn asrār baʿḍī az ʿibārāt-i anbiyā bi-badāyiʿ-i
naghamāt-i Ḥijāzī dar alvāḥ-i masṭūra-yi qabl ki barā-yi yakī az aḥbāb
nivishta shuda būd madhkūr gasht va ḥāl ham bi-taghanniyāt-i khush-i
ʿIrāqī naẓar bi-khvāhish-i ān janāb dar īn awrāq mujaddadan dhikr
mīnamāyīm ki shāyad tishnagān-i ṣaḥrā-hā-yi buʿd rā bi-baḥr-i qurb
dalālat namāyad va gumgashtagān-i biyābān-hā-yi hajr va firāq rā
bi-khiyām-i qurb va viṣāl rasānad tā ghamām-i ḍalālat murtafiʿ shavad va
āftāb-i jahāntāb-i hidāyat az ufuq-i jān ṭāliʿ gardad wa ʿalā Allāhi
attakilu wa bihi astaʿīnu laʿalla yajrī min hādhā al-qalami mā yuḥyī
bihi afʾidata al-nāsi li-yaqūmanna al-kullu ʿan marāqidi ghaflatihim wa
yasmaʿanna aṭwāra warqāti al-firdawsi min shajarin kāna fī al-rawḍati
al-aḥadiyyati min aydī al-qudrati bi-idhni Allāhi maghrūsan.

To unveil these mysteries, some of the sayings of the Prophets were
mentioned in the wondrous melodies of Hijaz in tablets previously
written for one of the friends. Now, at your request, We mention them
again in these pages in the sweet strains of Iraq, that they may perhaps
guide those thirsting in the deserts of remoteness to the ocean of
nearness, and bring those lost in the wilderness of separation and
estrangement to the tents of nearness and reunion. May the cloud of
error be lifted and the world-illumining sun of guidance rise from the
horizon of the soul. In God I place My trust, and from Him I seek help,
that perhaps there may flow from this Pen what will bring the hearts of
people to life, so that all may rise from their beds of heedlessness and
hear the melodies of the doves of paradise from a tree planted by the
hands of power, with God's permission, in the garden of oneness.

\subsection{20}\label{section-210}

\begin{otherlanguage}{persian}

بر اولی العلم معلوم و واضح بوده که چون نار محبّت عیسوی حجبات حدود یهود را
سوخت و حکم آن حضرت فی‌الجمله جریان بر حسب ظاهر یافت روزی آن جمال غیبی
ببعضی از اصحاب روحانی ذکر فراق فرمودند و نار اشتیاق افروختند و فرمودند
که من میروم و بعد میآیم و در مقام دیگر فرمودند من میروم و میآید دیگری تا
بگوید آنچه من نگفته‌ام و تمام نماید آنچه را که گفته‌ام و این دو عبارت
فی‌الحقیقه یکی است لو انتم فی مظاهر التّوحید بعین اللّه تشهدون

\end{otherlanguage}

Bar ulī al-ʿilm maʿlūm va vāḍiḥ būde ke chūn nār-e maḥabbat-e ʿĪsavī
ḥujubāt-e ḥudūd-e Yahūd rā sūkht va ḥukm-e ān Ḥaḍrat fī al-jumle jarayān
bar ḥasab-e ẓāhir yāft, rūzī ān Jamāl-e Ghaybī be-baʿḍī az aṣḥāb-e
rūḥānī dhekr-e firāq farmūdand va nār-e ishtiyāq afrūkhtand va farmūdand
ke man mīravam va baʿd mīāyam va dar maqām-e dīgar farmūdand man mīravam
va mīāyad dīgarī tā begūyad ānche man nagofte-am va tamām namāyad ānche
rā ke gofte-am va īn do ʿebārat fī al-ḥaqīqe yekī ast, law antum fī
maẓāhir al-tawḥīd bi-ʿayn Allāh tashhadūn.

It has been known and clear to those endowed with knowledge that, when
the fire of the love of Jesus had burned away the veils of Jewish
restrictions and His authority had, to some extent, outwardly prevailed,
that unseen Beauty one day spoke of separation to some of His spiritual
companions, kindled the fire of longing, and said, ``I am going away and
will come again.'' Elsewhere He said, ``I am going away, and another
will come to say what I have not said and complete what I have said.''
These two statements are, in truth, one, if you behold the
Manifestations of oneness with the eye of God.

\subsection{21}\label{section-211}

\begin{otherlanguage}{persian}

و اگر بدیدهٔ بصیرت معنوی مشاهده شود فی‌الحقیقه در عهد خاتم هم کتاب عیسی و
امر او ثابت شد در مقام اسم که خود حضرت فرمود منم عیسی و آثار و اخبار و
کتاب عیسی را هم تصدیق فرمود که من عند اللّه بوده در این مقام نه در خودشان
فرقی مشهود و نه در کتابشان غیریّتی ملحوظ زیرا که هر دو قائم بامر اللّه
بودند و هم ناطق بذکر اللّه و کتاب هر دو هم مشعر بر اوامر اللّه بود از این
جهت است که خود عیسی فرمود من میروم و مراجعت میکنم بمثل شمس که اگر شمس
الیوم بگوید من شمس یوم قبلم صادق است و اگر بگوید در حدود یومی که غیر آنم
صادقست و همچنین در ایّام ملاحظه نمائید که اگر گفته شود که کل یک شیء‌اند
صحیح و صادقست و اگر گفته شود که بحدود اسمی و رسمی غیر همند آنهم صادقست
چنانچه می‌بینی با اینکه یک شیء‌اند با وجود این در هر کدام اسمی دیگر و
خواصّی دیگر و رسمی دیگر ملحوظ میشود که در غیر آن نمیشود و بهمین بیان و
قاعده مقامات تفصیل و فرق و اتّحاد مظاهر قدسی را ادراک فرمائید تا تلویحات
کلمات آن مبدع اسماء و صفات را در مقامات جمع و فرق عارف شوی و واقف گردی و
جواب مسئلهٔ خود را در موسوم نمودن آن جمال ازلی در هر مقام خود را به اسمی
و رسمی بتمامه بیابی

\end{otherlanguage}

Va agar be-dīde-ye baṣīrat-e maʿnavī moshāhede shavad, fī al-ḥaqīqe dar
ʿahd-e Khātam ham ketāb-e ʿĪsā va amr-e ū thābet shod. Dar maqām-e esm
ke khod-e Ḥaḍrat farmūd manam ʿĪsā va āthār va akhbār va ketāb-e ʿĪsā rā
ham taṣdīq farmūd ke min ʿind Allāh būde. Dar īn maqām na dar khodeshān
farqī mashhūd va na dar ketābeshān ghayriyyatī malḥūẓ, zīrā ke har do
qāʾem be-amr Allāh būdand va ham nāṭeq be-dhekr Allāh va ketāb-e har do
ham moshʿer bar avāmer Allāh būd. Az īn jahat ast ke khod-e ʿĪsā farmūd
man mīravam va morājaʿat mīkonam. Be-methl-e shams ke agar shams-e
al-yawm begūyad man shams-e yawm-e qablam ṣādeq ast va agar begūyad dar
ḥudūd-e yawmī ke ghayr-e ānam ṣādeq ast. Va hamchenīn dar ayyām molāḥeẓe
namāʾīd ke agar gofte shavad ke koll yek shayʾ-and ṣaḥīḥ va ṣādeq ast va
agar gofte shavad ke be-ḥudūd-e esmī va rasmī ghayr-e hamand ān ham
ṣādeq ast. Chenānche mībīnī bā īnke yek shayʾ-and bā vojūd-e īn dar har
kodām esmī dīgar va khavāṣṣī dīgar va rasmī dīgar malḥūẓ mīshavad ke dar
ghayr-e ān nemīshavad. Va be-hamīn bayān va qāʿede maqāmāt-e tafṣīl va
farq va ittiḥād-e maẓāher-e qodsī rā edrāk farmāʾīd tā talvīḥāt-e
kalemāt-e ān Mobdeʿ-e asmā va ṣefāt rā dar maqāmāt-e jamʿ va farq ʿāref
shavī va vāqef gardī va javāb-e masʾale-ye khod rā dar mawsūm namūdan-e
ān Jamāl-e azalī dar har maqām khod rā be-esmī va rasmī be-tamāme
beyābī.

If viewed with the eye of spiritual insight, the Book of Jesus and His
Cause were, in truth, confirmed in the age of the Seal as well. With
regard to His name, He Himself declared, ``I am Jesus,'' and He also
affirmed that the works, accounts, and Book of Jesus were from God. From
this standpoint, no difference is seen between their persons, nor is any
distinction discerned between their Books, for both arose through the
command of God, both spoke in remembrance of God, and the Books of both
conveyed the commandments of God. This is why Jesus Himself said, ``I am
going away and will return.'' It is like the sun: if the sun of today
says, ``I am the sun of the previous day,'' it speaks the truth; and if,
with regard to the limitations of each day, it says, ``I am other than
that sun,'' it also speaks the truth. Likewise, consider the days: if it
is said that they are all one thing, this is correct and true; and if it
is said that, within the limits of their names and distinctive features,
they are different from one another, this too is true. As you see,
although they are one thing, a different name, different properties, and
different features are nevertheless discerned in each that are not
discerned in another. Through this same explanation and principle,
comprehend the stations of distinction, difference, and unity of the
holy Manifestations, so that you may recognize and understand the
allusions in the words of that Originator of names and attributes
concerning the stations of unity and distinction, and fully find the
answer to your question about that eternal Beauty designating Himself in
each station by a particular name and distinctive description.

\subsection{22}\label{section-212}

\begin{otherlanguage}{persian}

و بعد اصحاب و تلامیذ آن حضرت استدعا نمودند که علامت رجعت و ظهور چیست و
چه وقت این ظاهر خواهد شد و در چند مقام این سؤال را از آن طلعت بیمثال
نمودند و آن حضرت در هر مقام علامتی ذکر فرمودند چنانچه در اناجیل اربعه
مسطور است

\end{otherlanguage}

Va baʿd aṣḥāb va talāmīdh-e ān Ḥaḍrat estedʿā namūdand ke ʿalāmat-e
rajʿat va ẓohūr chīst va che vaqt īn ẓāher khvāhad shod va dar chand
maqām īn soʾāl rā az ān Ṭalʿat-e bīmethāl namūdand va ān Ḥaḍrat dar har
maqām ʿalāmatī dhekr farmūdand chenānche dar anājīl-e arbaʿe masṭūr ast.

Afterward, His companions and disciples asked what the sign of His
return and appearance would be, and when this would take place. On
several occasions they put this question to that incomparable
Countenance, and on each occasion He mentioned a sign, as recorded in
the four Gospels.

\subsection{23}\label{section-213}

\begin{otherlanguage}{persian}

و این مظلوم یک فقرهٔ آن را ذکر مینمایم و نعمتهای مکنونهٔ سدرهٔ مخزونه را
لوجه اللّه بر عباد اللّه مبذول میدارم تا هیاکل فانیه از اثمار باقیه محروم
نمانند که شاید برشحی از انهار بیزوال حضرت ذی الجلال که در دارالسّلام
بغداد جاری شده فائز شوند بی آنکه اجر و مزدی طلب نمایم انّما نطعمکم لوجه
اللّه لا نرید منکم جزآءً و لا شکوراً و این طعامیست که ارواح و افئدهٔ منیره
باو حیات باقیه یابند و این همان مائده‌ایست که میفرماید ربّنا انزل علینا
مائدةً من السّمآء و این مائده هرگز از اهلش مقطوع نشود و نفاد نجوید و در کلّ
حین از شجرهٔ فضل میروید و از سموات رحمت و عدل نازل میشود چنانچه فرموده
است مثلاً کلمةً طیّبةً کشجرة طیّبة اصلها ثابت و فرعها فی السّمآء تؤتی اکلها کلّ
حین

\end{otherlanguage}

Va īn maẓlūm yek feqre-ye ān rā dhekr mīnamāyam va neʿmathā-ye
maknūne-ye Sedre-ye makhzūne rā li-wajh Allāh bar ʿebād Allāh mabdhūl
mīdāram tā hayākel-e fāniye az athmār-e bāqiye maḥrūm namānand ke shāyad
be-rashḥī az anhār-e bīzavāl-e Ḥaḍrat-e Dhi al-Jalāl ke dar Dār
al-Salām-e Baghdād jārī shode fāʾez shavand bī ānke ajr va mozdī ṭalab
namāyam. Innamā nuṭʿimukum li-wajh Allāh lā nurīdu minkum jazāʾan wa-lā
shukūran. Va īn ṭaʿāmīst ke arvāḥ va afʾede-ye monīre be-ū ḥayāt-e
bāqiye yāband va īn hamān māʾede-īst ke mīfarmāyad rabbanā anzil ʿalaynā
māʾidatan min al-samāʾ. Va īn māʾede hargiz az ahlash maqṭūʿ nashavad va
nafād najūyad va dar koll-e ḥīn az shajare-ye faḍl mīrūyad va az
samavāt-e raḥmat va ʿadl nāzel mīshavad chenānche farmūde ast mathalan
kalimatan ṭayyibatan ka-shajaratin ṭayyibatin aṣluhā thābitun wa-farʿuhā
fī al-samāʾ tuʾtī ukulahā kulla ḥīn.

I, this wronged one, will mention one passage and, for the sake of God,
bestow upon the servants of God the hidden blessings of the treasured
Lote Tree, so that mortal beings may not be deprived of its everlasting
fruits. Perhaps they may attain a sprinkling from the unfailing rivers
of the Lord of Glory that have flowed in Baghdad, the Abode of Peace,
without My asking any recompense or reward. ``We feed you only for the
sake of God; We desire from you neither recompense nor thanks.'' This is
the food through which spirits and radiant hearts attain everlasting
life, and this is that very heavenly table of which He says, ``Our Lord,
send down to us a table from heaven.'' This table will never be withheld
from its people, nor will it be exhausted. At every moment it grows from
the tree of grace and descends from the heavens of mercy and justice, as
He has said: ``As a parable: a good word is like a good tree, its root
firm and its branches in heaven; it yields its fruit at all times.''

\subsection{24}\label{section-214}

\begin{otherlanguage}{persian}

حیف است که انسان از این عطیّهٔ لطیفه خود را منع نماید و از این نعمت باقیه
و حیات دائمه خود را محروم سازد پس قدر این مائدهٔ معنوی را دانسته که بلکه
از الطاف بدیعهٔ آن شمس حقیقی اجسادهای مرده حیات تازه یابند و ارواح پژمرده
بروح بی‌اندازه فائز شوند ای برادر من جهدی باید تا ایّام باقیست از اکواب
باقی چشیم همیشه نسیم جان از مصر جانان نوزد و همیشه نهرهای تبیان در جریان
نه و مدام ابواب رضوان مفتوح نماند آید وقتی که عندلیبان جنان از گلستان
قدسی بآشیانهای الهی پرواز نمایند دیگر نه نغمهٔ بلبل شنوی و نه جمال گل
بینی پس تا حمامهٔ ازلی در شور و تغنّی است و بهار الهی در جلوه و تزیین
غنیمت شمرده گوش قلب را از سروش او بی‌بهره مکن اینست نصیحت این عبد آن جناب
و احبّای خدا را فمن شآء فلیقبل و من شآء فلیعرض انّ اللّه کان غنیّاً عنه و عمّا
یشاهد و یری

\end{otherlanguage}

Ḥayf ast ke ensān az īn ʿaṭiyye-ye laṭīfe khod rā manʿ namāyad va az īn
neʿmat-e bāqiye va ḥayāt-e dāʾeme khod rā maḥrūm sāzad. Pas qadr-e īn
māʾede-ye maʿnavī rā dāneste ke balke az alṭāf-e badīʿe-ye ān Shams-e
ḥaqīqī ajsādhā-ye morde ḥayāt-e tāze yāband va arvāḥ-e pazhmorde
be-rawḥ-e bī-andāze fāʾez shavand. Ey barādar-e man jahdī bāyad tā ayyām
bāqīst az akvāb-e bāqī chashīm. Hamīshe nasīm-e jān az Meṣr-e Jānān
navazad va hamīshe nahrhā-ye tebyān dar jarayān na va modām abvāb-e
Reḍvān maftūḥ namānad. Āyad vaqtī ke ʿandalībān-e jenān az golestān-e
qodsī be-āshiyānhā-ye elāhī parvāz namāyand dīgar na naghme-ye bolbol
shenavī va na jamāl-e gol bīnī. Pas tā Ḥamāme-ye azalī dar shūr va
taghannī ast va bahār-e elāhī dar jelve va tazyīn ghanīmat shemorde
gūsh-e qalb rā az sorūsh-e ū bī-bahre makon. Īnast naṣīḥat-e īn ʿabd ān
janāb va aḥebbā-ye Khodā rā. Fa-man shāʾa fa-l-yaqbal wa-man shāʾa
fa-l-yuʿriḍ inna Allāha kāna ghaniyyan ʿanhu wa-ʿammā yushāhadu wa-yurā.

It is a pity for a person to deny himself this gracious gift and deprive
himself of this lasting blessing and enduring life. Appreciate, then,
the value of this spiritual table, so that through the wondrous favors
of that true Sun, dead bodies may perhaps find fresh life and withered
spirits attain boundless joy. My brother, we must make an effort, while
the days remain, to taste from the cups of eternity. The life-giving
breeze will not always blow from the Egypt of the Beloved; the rivers of
exposition will not always flow; nor will the gates of paradise remain
forever open. A time will come when the nightingales of paradise will
fly from the holy rose garden to their divine nests, and you will no
longer hear the nightingale's song or behold the rose's beauty. So,
while the eternal Dove is singing in rapture and the divine springtime
displays its splendor and adornment, seize the opportunity and do not
deprive the ear of your heart of its call. This is this servant's
counsel to you and to the loved ones of God. Whoever wishes, let him
accept it, and whoever wishes, let him turn away. God is indeed
independent of him and of all that is witnessed and seen.

\subsection{25}\label{section-215}

\begin{otherlanguage}{persian}

و اینست نغمات عیسی بن مریم که در رضوان انجیل بالحان جلیل در علائم ظهور
بعد فرموده در سفر اوّل که منسوب به متّی است در وقتی که سؤال نمودند از
علامات ظهور بعد جواب فرمود و للوقت من بعد ضیق تلک الأیّام تظلم الشّمس و
القمر لا یعطی ضوئه و الکواکب تتساقط من السّمآء و قوّات الأرض ترتجّ حینئذ
یظهر علامات ابن الانسان فی السّمآء و ینوح کلّ قبائل الأرض و یرون ابن
الانسان آتیاً علی سحاب السّمآء مع قوّات و مجد کبیر و یرسل ملائکته مع صوت
السّافور العظیم انتهی ترجمهٔ آن بفارسی اینست که بعد از تنگی و ابتلا که همهٔ
مردم را احاطه مینماید شمس از افاضه ممنوع میشود یعنی تاریک میگردد و قمر
از اعطای نور بازمیماند و ستاره‌های سماء بر ارض نازل میشوند و ارکان ارض
متزلزل میشود در این وقت ظاهر میگردد نشانه‌های پسر انسان در آسمان یعنی
جمال موعود و ساذج وجود بعد از ظهور این علامات از عرصهٔ غیب بعالم شهود
میآید و میفرماید در آن حین جمیع قبیله‌ها که در ارض ساکنند نوحه و ندبه
مینمایند و میبینند خلایق آن جمال احدیّه را که میآید از آسمان در حالتی که
سوار بر ابر است با قوّت و بزرگی و بخششی بزرگ و میفرستد ملائکه‌های خود را
با صدای سافور عظیم انتهی و در اسفار ثلاثهٔ دیگر که منسوب به لوقا و مرقس و
یوحنّاست همین عبارات مذکور است و چون در الواح عربیّه بتفصیل مذکور شد دیگر
در این اوراق متعرّض ذکر آنها نشدیم و اکتفا بیکی از آنها نمودیم

\end{otherlanguage}

Va īnast naghamāt-e ʿĪsā ibn-e Maryam ke dar Reḍvān-e Enjīl be-alḥān-e
jalīl dar ʿalāʾem-e ẓohūr-e baʿd farmūde. Dar sefr-e avval ke mansūb
be-Mattā ast dar vaqtī ke soʾāl namūdand az ʿalāmāt-e ẓohūr-e baʿd javāb
farmūd: Wa-li-l-waqt min baʿdi ḍīqi tilka al-ayyām tuẓlimu al-shams
wa-al-qamar lā yuʿṭī ḍawʾahu wa-al-kawākib tatasāqaṭu min al-samāʾ
wa-quwwāt al-arḍ tartajj. Ḥīnaʾidh yaẓharu ʿalāmāt ibn al-insān fī
al-samāʾ wa-yanūḥu kull qabāʾil al-arḍ wa-yarawna ibn al-insān ātiyan
ʿalā saḥāb al-samāʾ maʿa quwwāt wa-majd kabīr wa-yursilu malāʾikatahu
maʿa ṣawt al-sāfūr al-ʿaẓīm. Intahā. Tarjome-ye ān be-Fārsī īnast ke
baʿd az tangī va ebtelā ke hame-ye mardom rā eḥāṭe mīnamāyad shams az
efāḍe mamnūʿ mīshavad yaʿnī tārīk mīgardad va qamar az eʿṭā-ye nūr bāz
mīmānad va setārehā-ye samā bar arḍ nāzel mīshavand va arkān-e arḍ
motazalzel mīshavad. Dar īn vaqt ẓāher mīgardad neshānehā-ye pesar-e
ensān dar āsmān yaʿnī Jamāl-e mawʿūd va Sādhej-e vojūd baʿd az ẓohūr-e
īn ʿalāmāt az ʿarṣe-ye ghayb be-ʿālam-e shohūd mīāyad va mīfarmāyad dar
ān ḥīn jamīʿ-e qabīlehā ke dar arḍ sākenand nawḥe va nodbe mīnamāyand va
mībīnand khalāyeq ān Jamāl-e aḥadiyye rā ke mīāyad az āsmān dar ḥālatī
ke savār bar abr ast bā qovvat va bozorgī va bakhsheshī bozorg va
mīferestad malāʾekehā-ye khod rā bā ṣedā-ye sāfūr-e ʿaẓīm. Entahā. Va
dar asfār-e thalāthe-ye dīgar ke mansūb be-Lūqā va Marqos va Yūḥannāst
hamīn ʿebārāt madhkūr ast va chūn dar alvāḥ-e ʿarabiyye be-tafṣīl
madhkūr shod dīgar dar īn awrāq motaʿarreḍ-e dhekr-e ānhā nashodīm va
ektefā be-yekī az ānhā namūdīm.

These are the melodies that Jesus, son of Mary, uttered in majestic
tones in the paradise of the Gospel concerning the signs of the next
revelation. In the first book, attributed to Matthew, when asked about
the signs of the next revelation, He answered: ``Immediately after the
distress of those days, the sun will be darkened, and the moon will not
give its light, and the stars will fall from heaven, and the powers of
the earth will be shaken. Then the signs of the Son of Man will appear
in heaven, and all the tribes of the earth will mourn, and they will see
the Son of Man coming upon the clouds of heaven with power and great
glory. And He will send His angels with the sound of the great
trumpet.'' Here ends the quotation. Its translation into Persian is
this: After the hardship and tribulation that encompass all people, the
sun will be prevented from shedding its radiance---that is, it will grow
dark---and the moon will cease to give light, and the stars of heaven
will descend to earth, and the foundations of the earth will be shaken.
At that time the signs of the Son of Man will appear in heaven; that is,
the promised Beauty and pure Essence of being will come from the realm
of the unseen into the visible world after the appearance of these
signs. He also says that at that time all the tribes dwelling on earth
will lament and mourn, and people will see that Beauty of oneness coming
from heaven, riding upon a cloud, with power, grandeur, and great
bounty, and He will send His angels with the sound of the great trumpet.
Here ends the quotation. These same statements are mentioned in the
other three books, attributed to Luke, Mark, and John. Since they have
been discussed in detail in the Arabic Tablets, We have not undertaken
to mention them in these pages and have confined Ourselves to one of
them.

\subsection{26}\label{section-216}

\begin{otherlanguage}{persian}

و علمای انجیل چون عارف بمعانی این بیانات و مقصود مودعهٔ در این کلمات
نشدند و بظاهر آن متمسّک شدند لهذا از شریعهٔ فیض محمّدیّه و از سحاب فضل
احمدیّه ممنوع گشتند و جهّال آن طائفه هم تمسّک بعلمای خود جسته از زیارت جمال
سلطان جلال محروم ماندند زیرا که در ظهور شمس احمدیّه چنین علامات که مذکور
شد بظهور نیامد اینست که قرنها گذشت و عهدها بآخر رسید و آن جوهر روح بمقرّ
بقای سلطنت خود راجع شد و نفخهٔ دیگر از نفس روحانی در صور الهی دمیده شد و
نفسهای مرده از قبور غفلت و ضلالت بارض هدایت و محلّ عنایت محشور شدند و
هنوز آن گروه در انتظار که کی این علامات ظاهر شود و آن هیکل معهود بوجود
آید تا نصرت نمایند و مالها در راهش انفاق کنند و جانها در سبیلش ایثار
چنانچه امم دیگر هم بهمین ظنونات از کوثر معانی رحمت نامتناهی حضرت باری
دور مانده‌اند و بخیال خود مشغولند

\end{otherlanguage}

Va ʿolamā-ye Enjīl chūn ʿāref be-maʿānī-ye īn bayānāt va maqṣūd-e
mūdaʿe-ye dar īn kalemāt nashodand va be-ẓāher-e ān motamassek shodand
le-hādhā az sharīʿe-ye fayḍ-e Moḥammadiyye va az saḥāb-e faḍl-e
Aḥmadiyye mamnūʿ gashtand va johhāl-e ān ṭāʾefe ham tamassok
be-ʿolamā-ye khod joste az ziyārat-e Jamāl-e Solṭān-e Jalāl maḥrūm
māndand zīrā ke dar ẓohūr-e Shams-e Aḥmadiyye chenīn ʿalāmāt ke madhkūr
shod be-ẓohūr nayāmad. Īnast ke qarnhā godhasht va ʿahdhā be-ākher rasīd
va ān Jawhar-e Rūḥ be-maqarr-e baqā-ye salṭanat-e khod rājeʿ shod va
nafkhe-ye dīgar az nafas-e rūḥānī dar ṣūr-e elāhī damīde shod va
nafs-hā-ye morde az qobūr-e ghaflat va ḍalālat be-arḍ-e hedāyat va
maḥall-e ʿenāyat maḥshūr shodand va hanūz ān gorūh dar enteẓār ke key īn
ʿalāmāt ẓāher shavad va ān Haykal-e maʿhūd be-vojūd āyad tā noṣrat
namāyand va mālhā dar rāhash enfāq konand va jānhā dar sabīlash īthār.
Chenānche omam-e dīgar ham be-hamīn ẓonūnāt az Kawthar-e maʿānī-ye
raḥmat-e nāmotanāhī-ye Ḥaḍrat-e Bārī dūr mānde-and va be-khiyāl-e khod
mashghūland.

Because the scholars of the Gospel did not understand the meanings of
these statements or the purpose enshrined in these words, and clung to
their outward sense, they were barred from the stream of Muḥammad's
outpouring grace and the cloud of Aḥmad's bounty. The ignorant among
that community likewise clung to their own scholars and remained
deprived of beholding the beauty of the King of Glory, because such
signs as have been mentioned did not appear at the dawning of the Sun of
Aḥmad. Thus centuries passed and ages came to an end, and that Essence
of the Spirit returned to the eternal seat of His sovereignty. Another
blast was sounded in the divine trumpet by the spiritual breath, and
dead souls were raised from the graves of heedlessness and error to the
land of guidance and the abode of favor. Yet that people still wait for
these signs to appear and that promised Figure to come into being, so
that they may support Him, spend their wealth in His path, and sacrifice
their lives for His sake. Other religious communities, likewise, have
through these same suppositions remained far from the Kawthar of the
meanings of the Creator's boundless mercy, and are occupied with their
own imaginings.

\subsection{27}\label{section-217}

\begin{otherlanguage}{persian}

و از این عبارت گذشته بیان دیگر در انجیل هست که میفرماید السّمآء و الأرض
تزولان ولکن کلامی لا یزول که معنی آن بفارسی اینست که آسمان و زمین ممکن
است که زایل و معدوم شوند امّا کلام من هرگز زایل نمیشود و همیشه باقی و
ثابت میانهٔ ناس خواهد بود و از این راه است که اهل انجیل میگویند که حکم
انجیل هرگز منسوخ نمیشود و هر وقت و زمان که طلعت موعود با همهٔ علامتها
ظاهر شود باید شریعت مرتفعهٔ در انجیل را محکم و ثابت نماید تا در همۀ عالم
دینی باقی نماند مگر این دین و این فقره از مطالب محقّقهٔ مسلّمه است نزد
ایشان و چنان اعتقاد کرده‌اند که اگر نفسی هم مبعوث شود بجمیع علامات موعوده
و بر خلاف حکم ظاهر در انجیل حکم نماید البتّه اذعان نکنند و قبول ننمایند
بلکه تکفیر نمایند و استهزاء کنند چنانچه در ظهور شمس محمّدیّه مشهود شد حال
اگر معانی این کلمات منزلهٔ در کتب را که جمیع ناس از عدم بلوغ بآن از غایت
قصوی و سدرهٔ منتهی محجوب شده‌اند از ظهورات احدیّه در هر ظهور بتمام خضوع
سؤال مینمودند البتّه بانوار شمس هدایت مهتدی میشدند و باسرار علم و حکمت
واقف میگشتند

\end{otherlanguage}

Va az īn ʿebārat godhashte bayān-e dīgar dar Enjīl hast ke mīfarmāyad
al-samāʾ wa-al-arḍ tazūlān wa-lākin kalāmī lā yazūl, ke maʿnī-ye ān
be-Fārsī īnast ke āsmān va zamīn momken ast ke zāʾel va maʿdūm shavand
ammā kalām-e man hargiz zāʾel nemīshavad va hamīshe bāqī va thābet
miyāne-ye nās khvāhad būd. Va az īn rāh ast ke ahl-e Enjīl mīgūyand ke
ḥokm-e Enjīl hargiz mansūkh nemīshavad va har vaqt va zamān ke Ṭalʿat-e
mawʿūd bā hame-ye ʿalāmathā ẓāher shavad bāyad sharīʿat-e mortafeʿe-ye
dar Enjīl rā moḥkam va thābet namāyad tā dar hame-ye ʿālam dīnī bāqī
namānad magar īn dīn va īn feqre az maṭāleb-e moḥaqqaqe-ye mosallame ast
nazd-e īshān. Va chenān eʿteqād karde-and ke agar nafsī ham mabʿūth
shavad be-jamīʿ-e ʿalāmāt-e mawʿūde va bar khelāf-e ḥokm-e ẓāher dar
Enjīl ḥokm namāyad albatte edhʿān nakonand va qabūl nanamāyand balke
takfīr namāyand va estehzā konand chenānche dar ẓohūr-e Shams-e
Moḥammadiyye mashhūd shod. Ḥāl agar maʿānī-ye īn kalemāt-e monazzale-ye
dar kotob rā ke jamīʿ-e nās az ʿadam-e bolūgh be-ān az ghāyat-e qoṣvā va
Sedre-ye montahā maḥjūb shode-and az ẓohūrāt-e aḥadiyye dar har ẓohūr
be-tamām-e khoḍūʿ soʾāl mīnamūdand albatte be-anvār-e Shams-e hedāyat
mohtadī mīshodand va be-asrār-e ʿelm va ḥekmat vāqef mīgashtand.

Beyond this passage, there is another statement in the Gospel, where He
says, ``Heaven and earth will pass away, but My words will not pass
away.'' Its meaning in Persian is this: Heaven and earth may pass away
and cease to exist, but My words will never pass away and will always
remain enduring and established among people. This is why the people of
the Gospel say that the law of the Gospel will never be abrogated, and
that whenever the promised Countenance appears with all the signs, He
must confirm and establish the law exalted in the Gospel, so that no
religion but this religion remains in the whole world. They regard this
as an established and indisputable matter. So firmly do they believe it
that, even if someone were raised up with all the promised signs but
pronounced a ruling contrary to the law apparent in the Gospel, they
would certainly neither acknowledge nor accept Him, but would instead
declare Him an unbeliever and ridicule Him, as was witnessed at the
dawning of the Sun of Muḥammad. Now, if in every revelation they had
asked the Manifestations of oneness, with complete humility, about the
meanings of these words revealed in the scriptures---words whose meaning
all people have failed to attain and have thus remained veiled from the
ultimate goal and the Lote Tree of the farthest limit---they would
certainly have been guided by the rays of the Sun of guidance and become
acquainted with the mysteries of knowledge and wisdom.

\subsection{28}\label{section-218}

\begin{otherlanguage}{arabic}

حال این بنده رشحی از معانی این کلمات را ذکر مینمایم تا اصحاب بصیرت و
فطرت از معنی آن بجمیع تلویحات کلمات الهی و اشارات بیانات مظاهر قدسی واقف
شوند تا از هیمنهٔ کلمات از بحر اسماء و صفات ممنوع نشوند و از مصباح احدیّه
که محلّ تجلّی ذاتست محجوب نگردند

\end{otherlanguage}

Ḥāl īn bande rashḥī az maʿānī-ye īn kalemāt rā dhekr mīnamāyam tā
aṣḥāb-e baṣīrat va feṭrat az maʿnī-ye ān be-jamīʿ-e talvīḥāt-e kalemāt-e
elāhī va eshārāt-e bayānāt-e maẓāher-e qodsī vāqef shavand tā az
haymane-ye kalemāt az baḥr-e asmā va ṣefāt mamnūʿ nashavand va az
Meṣbāḥ-e aḥadiyye ke maḥall-e tajallī-ye Dhāt ast maḥjūb nagardand.

Now I, this servant, will mention a sprinkling of the meanings of these
words, so that those endowed with insight and a sound nature may,
through their meaning, become acquainted with all the allusions in the
divine words and the intimations in the utterances of the holy
Manifestations. Thus may the overwhelming majesty of the words not bar
them from the ocean of names and attributes, nor veil them from the Lamp
of oneness, which is the seat of the revelation of the divine Essence.

\subsection{29}\label{section-219}

\begin{otherlanguage}{persian}

قوله من بعد ضیق تلک الأیّام یعنی وقتی که ناس در سختی و تنگی مبتلا شوند و
این در وقتی است که آثار شمس حقیقت و اثمار سدرهٔ علم و حکمت از میان مردم
زایل شود و زمام ناس بدست جهّال افتد و ابواب توحید و معرفت که مقصود اصلی
از خلق انسانیست مسدود شود و علم بظنّ تبدیل گردد و هدایت بشقاوت راجع شود
چنانچه الیوم مشاهده میشود که زمام هر گروهی بدست جاهلی افتاده و بهر نحو
که اراده کنند حرکت میدهند و در میان ایشان از معبود جز اسمی و از مقصود جز
حرفی نمانده و بقسمی بادهای هوی و نفس غالب شده که سراجهای عقل و فؤاد را
در قلوب خاموش نموده با اینکه ابواب علم الهی بمفاتیح قدرت ربّانی مفتوح
گشته و جواهر وجود ممکنات بنور علمی و فیوضات قدسی منوّر و مهتدی گشتند
بقسمی که در هر شیء بابی از علم باز گشته و در هر ذرّه‌ آثاری از شمس مشهود
شده و با همهٔ این ظهورات علمی که عالم را احاطه نموده هنوز باب علم را
مسدود دانسته‌اند و امطار رحمت را مقطوع گرفته‌اند بظنّ تمسّک جسته از عروة
الوثقای محکم علم دور مانده‌اند و آنچه از ایشان مفهوم میشود گویا بعلم و
باب آن بالفطرة رغبتی ندارند و در خیال ظهور آن هم نیستند زیرا که در ظنّ و
گمان ابوابی برای نان یافته‌اند و در ظهور مظهر علم جز انفاق جان چیزی
نیافته‌اند لهذا البتّه از این گریزانند و بآن متمسّک و با اینکه حکم الهی را
یک میدانند از هر گوشه‌ حکمی صادر میشود و از هر محلّی امری ظاهر دو نفس بر
یک حکم ملاحظه نمیشود زیرا جز هوی الهی نجویند و بغیر از خطا سبیلی نخواهند
ریاست را نهایت وصول بمطلوب دانسته‌اند و کبر و غرور را غایت بلوغ بمحبوب
شمرده‌اند تزویرات نفسانی را مقدّم بر تقدیرات ربّانی دانند از تسلیم و رضا
گذشته‌اند و بتدبیر و ریا اشتغال نموده‌اند و بتمام قوّت و قدرت حفظ این مراتب
را مینمایند که مبادا نقصی در شوکت راه یابد و یا خللی در عزّت بهم رسد و
اگر چشمی از کحل معارف الهی روشن شود ملاحظه میکند سبعی چند را که بر
مردارهای نفوس عباد افتاده‌اند

\end{otherlanguage}

Qawluhu min baʿdi ḍīqi tilka al-ayyām. Yaʿnī vaqtī ke nās dar sakhtī va
tangī mubtalā shavand va īn dar vaqtī ast ke āthār-i shams-i ḥaqīqat va
athmār-i sidri-yi ʿilm va ḥikmat az miyān-i mardum zāyil shavad va
zimām-i nās bi-dast-i juhhāl uftad va abvāb-i tawḥīd va maʿrifat ke
maqṣūd-i aṣlī az khalq-i insānīst masdūd shavad va ʿilm bi-ẓann tabdīl
gardad va hidāyat bi-shaqāvat rājiʿ shavad. Chunānchi al-yawm mushāhidi
mī-shavad ke zimām-i har gurūhī bi-dast-i jāhilī uftādi va bi-har naḥv
ke irādi kunand ḥarakat mī-dihand va dar miyān-i īshān az maʿbūd juz
ismī va az maqṣūd juz ḥarfī na-māndi va bi-qismī bādhā-yi havā va nafs
ghālib shudi ke sirājhā-yi ʿaql va fuʾād rā dar qulūb khāmūsh namūdi. Bā
īnki abvāb-i ʿilm-i ilāhī bi-mafātīḥ-i qudrat-i rabbānī maftūḥ gashti va
javāhir-i vujūd-i mumkināt bi-nūr-i ʿilmī va fuyūḍāt-i qudsī munavvar va
muhtadī gashtand bi-qismī ke dar har shayʾ bābī az ʿilm bāz gashti va
dar har dharri āthārī az shams mashhūd shudi va bā hami-yi īn ẓuhūrāt-i
ʿilmī ke ʿālam rā iḥāṭi namūdi hanūz bāb-i ʿilm rā masdūd dānisti-and va
amṭār-i raḥmat rā maqṭūʿ girifti-and. Bi-ẓann tamassuk justi az
ʿurvatu'l-vuthqā-yi muḥkam-i ʿilm dūr māndi-and va ānchi az īshān mafhūm
mī-shavad gūyā bi-ʿilm va bāb-i ān bi'l-fiṭri raghbatī na-dārand va dar
khiyāl-i ẓuhūr-i ān ham nīstand zīrā ke dar ẓann va gumān abvābī barā-yi
nān yāfti-and va dar ẓuhūr-i maẓhar-i ʿilm juz infāq-i jān chīzī
na-yāfti-and. Li-hādhā al-batti az īn gurīzānand va bi-ān mutamassik va
bā īnki ḥukm-i ilāhī rā yak mī-dānand az har gūshi ḥukmī ṣādir mī-shavad
va az har maḥallī amrī ẓāhir. Du nafs bar yak ḥukm mulāḥiẓi na-mī-shavad
zīrā juz havā ilāhī na-jūyand va bi-ghayr az khaṭā sabīlī na-khvāhand.
Riyāsat rā nihāyat-i vuṣūl bi-maṭlūb dānisti-and va kibr va ghurūr rā
ghāyat-i bulūgh bi-maḥbūb shumurdi-and. Tazvīrāt-i nafsānī rā muqaddam
bar taqdīrāt-i rabbānī dānand az taslīm va riḍā gudhashti-and va
bi-tadbīr va riyā ishtighāl namūdi-and va bi-tamām-i quvvat va qudrat
ḥifẓ-i īn marātib rā mī-namāyand ke mabādā naqṣī dar shawkat rāh yābad
va yā khalalī dar ʿizzat bi-ham rasad va agar chashmī az kuḥl-i
maʿārif-i ilāhī rawshan shavad mulāḥiẓi mī-kunad sibuʿī chand rā ke bar
murdārhā-yi nufūs-i ʿibād uftādi-and.

His words, ``After the distress of those days,'' mean the time when
people are afflicted with hardship and distress. This is when the signs
of the Sun of Truth and the fruits of the Tree of knowledge and wisdom
disappear from among the people, the reins of humanity fall into the
hands of the ignorant, the doors of divine unity and recognition---which
are the primary purpose of human creation---are closed, knowledge turns
into conjecture, and guidance gives way to wretchedness. Today, as can
be seen, the reins of every group have fallen into the hands of an
ignorant person, and they lead people in whatever direction they wish.
Among them, nothing remains of the One they worship but a name, and
nothing of the One they seek but a word. The winds of desire and self
have so prevailed that they have extinguished the lamps of reason and
heart within them. Although the doors of divine knowledge have been
opened by the keys of heavenly power, and the gems of the existence of
created things have been illumined and guided by the light of knowledge
and the outpourings of holiness, so that in every thing a door of
knowledge has opened and in every atom signs of the Sun have become
visible, yet, despite all these revelations of knowledge encompassing
the world, they still consider the door of knowledge closed and the
rains of mercy cut off. They have clung to conjecture and remained far
from the firm, surest handle of knowledge. From what can be understood
of them, it seems that they have no natural inclination toward knowledge
or its door, nor do they even contemplate its appearance, for in
conjecture and supposition they have found ways to earn their bread,
whereas in the appearance of the Manifestation of knowledge they have
found nothing but the sacrifice of life. Thus they inevitably flee from
the latter and cling to the former. Although they regard the divine
decree as one, a ruling issues from every corner and a command appears
in every place. No two people are found to agree on one ruling, for they
seek no god but desire and wish for no path but error. They have
regarded leadership as the ultimate attainment of their goal, and
counted pride and arrogance as the highest degree of reaching the
Beloved. They give precedence to selfish deceits over divine decrees;
they have abandoned submission and contentment and occupied themselves
with scheming and hypocrisy. With all their strength and power, they
preserve these positions, lest their grandeur suffer any diminution or
their prestige any harm. If an eye is illumined by the kohl of divine
knowledge, it sees a number of beasts preying upon the dead bodies of
the souls of God's servants.

\subsection{30}\label{section-220}

\begin{otherlanguage}{persian}

حال کدام ضیق و تنگی است که ازید از مراتب مذکوره باشد که اگر نفسی طلب حقّی
و یا معرفتی بخواهد نماید نمیداند نزد کدام رود و از که جویا شود از غایت
اینکه رأیها مختلف و سبیلها متعدّد شده و این تنگی و ضیق از شرایط هر ظهور
است که تا واقع نشود ظهور شمس حقیقت نشود زیرا که صبح ظهور هدایت بعد از
لیل ضلالت طالع میشود اینست که در روایات و احادیث جمیع این مضامین هست که
کفر عالم را فرو میگیرد و ظلمت احاطه مینماید و امثال اینها چنانچه مذکور
شد و این عبد بواسطهٔ شهرت این احادیث و اختصار دیگر متعرّض ذکر عبارات حدیث
نشده‌ام

\end{otherlanguage}

Ḥāl kudām ḍīq va tangī ast ke azyad az marātib-i madhkūri bāshad ke agar
nafsī ṭalab-i ḥaqqī va yā maʿrifatī bi-khvāhad namāyad na-mī-dānad
nazd-i kudām ravad va az ke jūyā shavad az ghāyat-i īnki raʾyhā
mukhtalif va sabīlhā mutaʿaddid shudi. Va īn tangī va ḍīq az sharāyiṭ-i
har ẓuhūr ast ke tā vāqiʿ na-shavad ẓuhūr-i shams-i ḥaqīqat na-shavad
zīrā ke ṣubḥ-i ẓuhūr-i hidāyat baʿd az layl-i ḍalālat ṭāliʿ mī-shavad.
Īnast ke dar rivāyāt va aḥādīth jamīʿ-i īn maḍāmīn hast ke kufr ʿālam rā
farū mī-gīrad va ẓulmat iḥāṭi mī-namāyad va amthāl-i īnhā chunānchi
madhkūr shud va īn ʿabd bi-vāsiṭi-yi shuhrat-i īn aḥādīth va ikhtiṣār
dīgar mutaʿarriḍ-i dhikr-i ʿibārāt-i ḥadīth na-shudi-am.

Now, what distress and hardship could be greater than the conditions
described, when anyone wishing to seek truth or knowledge does not know
whom to approach or whom to ask, so greatly have opinions diverged and
paths multiplied? This hardship and distress are among the conditions of
every revelation: until they occur, the Sun of Truth does not appear,
for the dawn of the revelation of guidance rises after the night of
error. Thus the traditions and reported sayings contain all these
themes: that unbelief will overtake the world, that darkness will
encompass it, and similar statements, as already mentioned. Because
these traditions are well known, and for brevity, I, this servant, have
refrained from quoting their wording.

\subsection{31}\label{section-221}

\begin{otherlanguage}{persian}

حال اگر مقصود از این ضیق را همچو ادراک نمایند که عالم ضیق بهم رساند و یا
امورات دیگر که بخیال خود توهّم نمایند هرگز مشهود نگردد و البتّه گویند که
این شرط ظهور نیافته چنانچه گفته‌اند و میگویند باری مقصود از ضیق ضیق از
معارف الهیّه و ادراک کلمات ربّانیّه است که در ایّام غروب شمس و مرایای او
عباد در تنگی و سختی افتند و ندانند به که توجّه نمایند چنانچه مذکور شد
کذلک نعلّمک من تأویل الأحادیث و نلقی علیک من اسرار الحکمة لتطّلع بما هو
المقصود و تکون من الّذینهم شربوا من کأس العلم و العرفان

\end{otherlanguage}

Ḥāl agar maqṣūd az īn ḍīq rā hamchu idrāk namāyand ke ʿālam ḍīq bi-ham
rasānad va yā umūrāt-i dīgar ke bi-khiyāl-i khud tavahhum namāyand
hargiz mashhūd na-gardad va al-batti gūyand ke īn sharṭ ẓuhūr na-yāfti
chunānchi gufti-and va mī-gūyand. Bārī maqṣūd az ḍīq ḍīq az maʿārif-i
ilāhiyyi va idrāk-i kalimāt-i rabbāniyyi ast ke dar ayyām-i ghurūb-i
shams va marāyā-yi ū ʿibād dar tangī va sakhtī uftand va na-dānand bi-ke
tavajjuh namāyand chunānchi madhkūr shud. Kadhālika nuʿallimuka min
taʾwīli al-aḥādīthi wa-nulqī ʿalayka min asrāri al-ḥikmati li-taṭṭaliʿa
bi-mā huwa al-maqṣūdu wa-takūna mina alladhīna hum sharibū min kaʾsi
al-ʿilmi wa'l-ʿirfān.

Now, if they understand this distress to mean that the world will become
constricted, or imagine other things of their own devising, these will
never be witnessed, and they will certainly say that this condition has
not been fulfilled, as they have said and continue to say. Rather,
distress means being deprived of divine knowledge and an understanding
of the heavenly words: in the days when the Sun and its mirrors have
set, people fall into hardship and difficulty and do not know to whom
they should turn, as already mentioned. Thus We teach you something of
the interpretation of the traditions and impart to you some of the
mysteries of wisdom, so that you may discover what is intended and be
among those who have drunk from the cup of knowledge and recognition.

\subsection{32}\label{section-222}

\begin{otherlanguage}{persian}

و قوله تظلم الشّمس و القمر لا یعطی ضوئه و الکواکب تتساقط من السّمآء مقصود
از شمس و قمر که در کلمات انبیا مذکور است منحصر باین شمس و قمر ظاهری نیست
که ملاحظه میشود بلکه از شمس و قمر معانی بسیار اراده فرموده‌اند که در هر
مقام بمناسبت آن مقام معنیی اراده میفرمایند مثلاً یک معنی از شمس شمسهای
حقیقتند که از مشرق قدم طالع میشوند و بر جمیع ممکنات ابلاغ فیض میفرمایند
و این شموس حقیقت مظاهر کلّیّهٔ الهی هستند در عوالم صفات و اسمای او و همچنان
که شمس ظاهری تربیت اشیای ظاهره از اثمار و اشجار و الوان و فواکه و معادن
و دون ذلک از آنچه در عالم ملک مشهود است بامر معبود حقیقی باعانت اوست و
همچنین اشجار توحید و اثمار تفرید و اوراق تجرید و گلهای علم و ایقان و
ریاحین حکمت و بیان از عنایت و تربیت شمسهای معنوی ظاهر میشود اینست که در
حین اشراق این شموس عالم جدید میشود و انهار حیوان جاری میگردد و ابحر
احسان بموج میآید و سحاب فضل مرتفع میشود و نسمات جود بر هیاکل موجودات
میوزد و از حرارت این شمسهای الهی و نارهای معنویست که حرارت محبّت الهی در
ارکان عالم احداث میشود و از عنایت این ارواح مجرّده است که روح حیوان باقیه
بر اجساد مردگان فانیه مبذول میگردد و فی‌الحقیقه این شمس ظاهری یک آیه از
تجلّی آن شمس معنوی است و آن شمسی است که از برای او مقابلی و شبهی و مثلی و
ندّی ملاحظه نمیشود و کل بوجود او قائمند و از فیض او ظاهر و باو راجع منها
ظهرت الأشیآء و الی خزائن امرها رجعت و منها بدئت الممکنات و الی کنائز
حکمها عادت

\end{otherlanguage}

Va qawluhu taẓlumu al-shamsu wa'l-qamaru lā yuʿṭī ḍawʾahu wa'l-kawākibu
tatasāqaṭu mina al-samāʾ. Maqṣūd az shams va qamar ke dar kalimāt-i
anbiyā madhkūr ast munḥaṣir bi-īn shams va qamar-i ẓāhirī nīst ke
mulāḥiẓi mī-shavad balki az shams va qamar maʿānī-yi bisyār irādi
farmūdi-and ke dar har maqām bi-munāsibat-i ān maqām maʿnīyī irādi
mī-farmāyand. Mathalan yak maʿnī az shams shamshā-yi ḥaqīqatand ke az
mashriq-i qidam ṭāliʿ mī-shavand va bar jamīʿ-i mumkināt iblāgh-i fayḍ
mī-farmāyand va īn shumūs-i ḥaqīqat maẓāhir-i kulliyyi-yi ilāhī hastand
dar ʿavālim-i ṣifāt va asmā-yi ū. Va hamchunān ke shams-i ẓāhirī
tarbiyat-i ashyā-yi ẓāhiri az athmār va ashjār va alvān va favākih va
maʿādin va dūn-i dhālik az ānchi dar ʿālam-i mulk mashhūd ast bi-amr-i
maʿbūd-i ḥaqīqī bi-iʿānat-i ūst va hamchunīn ashjār-i tawḥīd va athmār-i
tafrīd va awrāq-i tajrīd va gulhā-yi ʿilm va īqān va rayāḥīn-i ḥikmat va
bayān az ʿināyat va tarbiyat-i shamshā-yi maʿnavī ẓāhir mī-shavad. Īnast
ke dar ḥīn-i ishrāq-i īn shumūs ʿālam jadīd mī-shavad va anhār-i ḥayavān
jārī mī-gardad va abḥur-i iḥsān bi-mawj mī-āyad va saḥāb-i faḍl murtafiʿ
mī-shavad va nasamāt-i jūd bar hayākil-i mawjūdāt mī-vazad va az
ḥarārat-i īn shamshā-yi ilāhī va nārhā-yi maʿnavīst ke ḥarārat-i
maḥabbat-i ilāhī dar arkān-i ʿālam iḥdāth mī-shavad va az ʿināyat-i īn
arvāḥ-i mujarradi ast ke rūḥ-i ḥayavān-i bāqiyi bar ajsād-i murdigān-i
fāniyi mabdhūl mī-gardad. Va fi'l-ḥaqīqi īn shams-i ẓāhirī yak āyi az
tajallī-yi ān shams-i maʿnavī ast va ān shamsī ast ke az barā-yi ū
muqābilī va shibhī va mithlī va niddī mulāḥiẓi na-mī-shavad va kull
bi-vujūd-i ū qāʾimand va az fayḍ-i ū ẓāhir va bi-ū rājiʿ. Minhā ẓaharati
al-ashyāʾu wa-ilā khazāʾini amrihā rajaʿat wa-minhā budiʾati
al-mumkinātu wa-ilā kanāʾizi ḥukmihā ʿādat.

And His words, ``The sun will be darkened, and the moon will not give
its light, and the stars will fall from heaven.'' The sun and moon
mentioned in the words of the Prophets do not refer solely to this
visible sun and moon that we observe. Rather, they have intended many
meanings by sun and moon, each meaning suited to its particular context.
For example, one meaning of the sun is the Suns of Truth, who rise from
the dawning-place of eternity and impart grace to all created things.
These Suns of Truth are the universal Manifestations of God in the
realms of His attributes and names. Just as, by the command of the true
Object of worship, the visible sun nurtures visible things---fruits,
trees, colors, produce, minerals, and all else that is seen in the
material world---so the trees of divine unity, the fruits of singleness,
the leaves of detachment, the flowers of knowledge and certitude, and
the fragrant plants of wisdom and utterance appear through the care and
nurture of the spiritual Suns. Thus, when these Suns shine, the world is
renewed, the rivers of life flow, the seas of beneficence surge, the
clouds of grace rise, and the breezes of bounty blow over the forms of
created beings. It is through the heat of these divine Suns and
spiritual fires that the warmth of divine love is kindled within the
world's foundations, and through the grace of these sanctified Spirits
that the spirit of everlasting life is bestowed upon the bodies of the
perishable dead. In truth, this visible sun is but a sign of the
radiance of that spiritual Sun, a Sun for which no counterpart,
likeness, equal, or rival can be found. All things subsist through its
existence, appear through its grace, and return to it. From it all
things have appeared, and to the treasuries of its command they have
returned; from it all created things have begun, and to the repositories
of its decree they have gone back.

\subsection{33}\label{section-223}

\begin{otherlanguage}{arabic}

و اینکه در مقام بیان و ذکر تخصیص داده میشوند ببعضی از اسماء و صفات
چنانچه شنیده‌اید و میشنوید نیست مگر برای ادراک عقول ناقصهٔ ضعیفه والّا
لم‌یزل و لایزال مقدّس بوده‌اند از هر اسمی و منزّه خواهند بود از هر وصفی
جواهر اسماء را بساحت قدسشان راهی نه و لطائف صفات را در ملکوت عزّشان سبیلی
نه فسبحان اللّه من ان یعرف اصفیائه بغیر ذواتهم او یوصف اولیائه بغیر
انفسهم فتعالی عمّا یذکر العباد فی وصفهم و تعالی عمّا هم یعرفون

\end{otherlanguage}

Va īnki dar maqām-i bayān va dhikr takhṣīṣ dādi mī-shavand bi-baʿḍī az
asmā va ṣifāt chunānchi shinīdi-īd va mī-shinavīd nīst magar barā-yi
idrāk-i ʿuqūl-i nāqiṣi-yi ḍaʿīfi va illā lam-yazal va lā-yazāl muqaddas
būdi-and az har ismī va munazzah khvāhand būd az har vaṣfī. Javāhir-i
asmā rā bi-sāḥat-i qudsishān rāhī na va laṭāʾif-i ṣifāt rā dar malakūt-i
ʿizzishān sabīlī na. Fa-subḥāna Allāhi min an yuʿrafa aṣfiyāʾuhu
bi-ghayri dhawātihim aw yūṣafa awliyāʾuhu bi-ghayri anfusihim fa-taʿālā
ʿammā yadhkuru al-ʿibādu fī waṣfihim wa-taʿālā ʿammā hum yaʿrifūn.

When, in explanation and discourse, particular names and attributes are
assigned to them, as you have heard and continue to hear, this is only
for the understanding of weak and limited minds. Otherwise, they have
always been sanctified above every name and will forever remain exalted
beyond every description. The essences of names have no access to their
court of holiness, and the subtleties of attributes have no path into
the kingdom of their glory. Glory be to God, that His chosen Ones should
be known through anything other than their own beings, or His holy Ones
described by anything other than their own selves! Exalted is He above
all that His servants say in describing them, and exalted is He above
all that they understand.

\subsection{34}\label{section-224}

\begin{otherlanguage}{persian}

و اطلاق شموس بر آن انوار مجرّده در کلمات اهل عصمت بسیار شده از آن جمله در
دعای ندبه میفرماید این الشّموس الطّالعة این الأقمار المنیرة این الأنجم
الزّاهرة پس معلوم شد که مقصود از شمس و قمر و نجوم در مقام اوّلیّه انبیا و
اولیا و اصحاب ایشانند که از انوار معارفشان عوالم غیب و شهود روشن و منوّر
است

\end{otherlanguage}

Va iṭlāq-i shumūs bar ān anvār-i mujarradi dar kalimāt-i ahl-i ʿiṣmat
bisyār shudi az ān jumli dar duʿā-yi nudbi mī-farmāyad ayna al-shumūsu
al-ṭāliʿatu ayna al-aqmāru al-munīratu ayna al-anjumu al-zāhirah. Pas
maʿlūm shud ke maqṣūd az shams va qamar va nujūm dar maqām-i avvaliyyi
anbiyā va awliyā va aṣḥāb-i īshānand ke az anvār-i maʿārifishān
ʿavālim-i ghayb va shuhūd rawshan va munavvar ast.

These sanctified Lights have often been called suns in the words of the
infallible Ones. Among these is the Prayer of Lamentation, which says,
``Where are the rising suns? Where are the shining moons? Where are the
radiant stars?'' It is thus clear that, in their primary meaning, the
sun, moon, and stars refer to the Prophets, the holy Ones, and their
companions, through the light of whose knowledge the worlds of the
unseen and the seen are bright and illumined.

\subsection{35}\label{section-225}

\begin{otherlanguage}{persian}

و در مقام دیگر مقصود از شمس و قمر و نجوم علمای ظهور قبلند که در زمان
ظهور بعد موجودند و زمام دین مردم در دست ایشانست و اگر در ظهور شمس اخری
بضیای او منوّر گشتند لهذا مقبول و منیر و روشن خواهند بود والّا حکم ظلمت در
حقّ آنها جاریست اگرچه بظاهر هادی باشند زیرا که جمیع این مراتب از کفر و
ایمان و هدایت و ضلالت و سعادت و شقاوت و نور و ظلمت منوط بتصدیق آن شمس
معنوی الهی است بر هر نفسی از علما حکم ایمان از مبدء عرفان در یوم تغابن و
احسان جاری شد حکم علم و رضا و نور و ایمان در بارهٔ او صادقست والّا حکم جهل
و نفی و کفر و ظلم در حقّ او جریان یابد

\end{otherlanguage}

Va dar maqām-i dīgar maqṣūd az shams va qamar va nujūm ʿulamā-yi ẓuhūr-i
qabland ke dar zamān-i ẓuhūr-i baʿd mawjūdand va zimām-i dīn-i mardum
dar dast-i īshānast va agar dar ẓuhūr-i shams-i ukhrā bi-ḍiyā-yi ū
munavvar gashtand li-hādhā maqbūl va munīr va rawshan khvāhand būd va
illā ḥukm-i ẓulmat dar ḥaqq-i ānhā jārīst agarchi bi-ẓāhir hādī bāshand
zīrā ke jamīʿ-i īn marātib az kufr va īmān va hidāyat va ḍalālat va
saʿādat va shaqāvat va nūr va ẓulmat manūṭ bi-taṣdīq-i ān shams-i
maʿnavī-yi ilāhī ast. Bar har nafsī az ʿulamā ḥukm-i īmān az mabdaʾ-i
ʿirfān dar yawm-i taghābun va iḥsān jārī shud ḥukm-i ʿilm va riḍā va nūr
va īmān dar bāri-yi ū ṣādiq ast va illā ḥukm-i jahl va nafy va kufr va
ẓulm dar ḥaqq-i ū jarayān yābad.

In another context, the sun, moon, and stars refer to the religious
scholars of the previous revelation who are present at the time of the
next revelation and hold in their hands the reins of the people's
religion. If, when the new Sun appears, they become illumined by its
radiance, they will be accepted, radiant, and bright. Otherwise, the
judgment of darkness applies to them, even though they may outwardly be
guides. For all these conditions---unbelief and faith, guidance and
error, blessedness and wretchedness, light and darkness---depend upon
acknowledging that divine, spiritual Sun. Any religious scholar upon
whom the Source of recognition pronounces the judgment of faith on the
day of loss and beneficence is truly entitled to the designations of
knowledge, divine approval, light, and faith. Otherwise, the judgment of
ignorance, negation, unbelief, and oppression applies to him.

\subsection{36}\label{section-226}

\begin{otherlanguage}{persian}

و این بر هر ذی بصری مشهود است که همچنان که نور ستاره محو میشود نزد اشراق
شمس ظاهره همین قسم شمس علم و حکمت و عرفان ظاهره نزد طلوع شمس حقیقت و
آفتاب معنوی محو و تاریک میشود

\end{otherlanguage}

Va īn bar har dhī baṣarī mashhūd ast ke hamchunān ke nūr-i sitāri maḥv
mī-shavad nazd-i ishrāq-i shams-i ẓāhiri hamīn qism shams-i ʿilm va
ḥikmat va ʿirfān-i ẓāhiri nazd-i ṭulūʿ-i shams-i ḥaqīqat va āftāb-i
maʿnavī maḥv va tārīk mī-shavad.

It is evident to everyone with insight that, just as the light of a star
fades before the shining of the visible sun, so the sun of outward
knowledge, wisdom, and understanding fades and grows dark before the
rising of the Sun of Truth, the spiritual Sun.

\subsection{37}\label{section-227}

\begin{otherlanguage}{persian}

و اطلاق شمس بر آن علما بمناسبت علوّ و شهرت و معروفیّت است مثل علمای مسلّم
عصر که مشهور بلاد و مسلّمند بین عباد و اگر حاکی از شمس الهی باشند از شموس
عالیه محسوبند والّا از شموس سجّین چنانچه میفرماید الشّمس و القمر بحسبان و
معنی شمس و قمر هم که در آیهٔ مذکوره هست البتّه شنیده‌اید احتیاج بذکر نیست و
هر نفسی هم که از عنصر این شمس و قمر باشد یعنی در اقبال بباطل و اعراض از
حقّ البتّه از حسبان ظاهر و بحسبان راجع خواهد شد

\end{otherlanguage}

Va iṭlāq-i shams bar ān ʿulamā bi-munāsibat-i ʿuluvv va shuhrat va
maʿrūfiyyat ast mithl-i ʿulamā-yi musallam-i ʿaṣr ke mashhūr-i bilād va
musallamand bayn-i ʿibād va agar ḥākī az shams-i ilāhī bāshand az
shumūs-i ʿāliyi maḥsūband va illā az shumūs-i sijjīn chunānchi
mī-farmāyad al-shamsu wa'l-qamaru bi-ḥusbān. Va maʿnī-yi shams va qamar
ham ke dar āyi-yi madhkūri hast al-batti shinīdi-īd iḥtiyāj bi-dhikr
nīst va har nafsī ham ke az ʿunṣur-i īn shams va qamar bāshad yaʿnī dar
iqbāl bi-bāṭil va iʿrāḍ az ḥaqq al-batti az ḥusbān ẓāhir va bi-ḥusbān
rājiʿ khvāhad shud.

These scholars are called suns because of their high standing, fame, and
renown, like the acknowledged religious scholars of the age, who are
renowned throughout the lands and accepted among the people. If they
reflect the divine Sun, they are counted among the exalted suns;
otherwise, they belong to the suns of Sijjin, as He says, ``The sun and
the moon are subject to a reckoning.'' You have surely heard the meaning
of the sun and moon in the verse just mentioned; there is no need to
repeat it. Anyone who is of the same nature as this sun and moon---that
is, in turning toward falsehood and away from truth---will certainly
have come forth from that reckoning and will return to it.

\subsection{38}\label{section-228}

\begin{otherlanguage}{persian}

پس ای سائل باید بعروة الوثقی متمسّک شویم که شاید از شام ضلالت بنور هدایت
راجع گردیم و از ظلّ نفی فرار نموده در ظلّ اثبات درآئیم و از نار حسبان آزاد
شده بنور جمال حضرت منّان منوّر گردیم و السّلام کذلک نعطیکم من اثمار شجرة
العلم لتکوننّ فی رضوان حکمة اللّه لمن المحبرین

\end{otherlanguage}

Pas ay sāʾil bāyad bi-ʿurvatu'l-vuthqā mutamassik shavīm ke shāyad az
shām-i ḍalālat bi-nūr-i hidāyat rājiʿ gardīm va az ẓill-i nafy firār
namūdi dar ẓill-i ithbāt dar-āyīm va az nār-i ḥusbān āzād shudi bi-nūr-i
jamāl-i Ḥaḍrat-i Mannān munavvar gardīm. Wa'l-salām. Kadhālika nuʿṭīkum
min athmāri shajarati al-ʿilmi li-takūnanna fī riḍwāni ḥikmati Allāhi
la-mina al-muḥbarīn.

So, seeker, we must cling to the surest handle, that perhaps we may turn
from the night of error to the light of guidance, flee the shadow of
negation and enter the shelter of affirmation, and be freed from the
fire of reckoning and illumined by the light of the beauty of the
All-Bountiful. Peace. Thus We give you some of the fruits of the Tree of
knowledge, so that you may surely be among those who rejoice in the
paradise of God's wisdom.

\subsection{39}\label{section-229}

\begin{otherlanguage}{persian}

و در مقامی هم مقصود از اطلاقات شمس و قمر و نجوم علوم و احکام مرتفعهٔ در
هر شریعت است مثل صلوة و صوم که در شریعت فرقان بعد از اخفای جمال محمّدی از
جمیع احکام محکمتر و اعظمتر است چنانچه احادیث و اخبار مشعر بر آنست و بعلّت
شهرت احتیاج ذکر نیست بلکه در هر عصری حکم صلوة محکم و مجری بوده چنانچه از
انوار مشرقه از شمس محمّدیّه

\end{otherlanguage}

Va dar maqāmī ham maqṣūd az eṭlāqāt-e shams va qamar va nojūm ʿolūm va
aḥkām-e mortafeʿe-ye dar har sharīʿat ast, mes̱l-e ṣalāt va ṣawm ke dar
sharīʿat-e Forqān baʿd az ekhfā-ye jamāl-e Moḥammadī az jamīʿ-e aḥkām
moḥkamtar va aʿẓamtar ast, chenānche aḥādīs̱ va akhbār moshʿer bar ān ast
va be-ʿellat-e shohrat eḥtīāj-e dhekr nīst. Balke dar har ʿaṣrī ḥokm-e
ṣalāt moḥkam va mojrā būde, chenānche az anvār-e moshreqe az shams-e
Moḥammadīye

In another sense, the terms sun, moon, and stars refer to the teachings
and laws established in every religious dispensation, such as prayer and
fasting, which, in the dispensation of the Quran, are the most firmly
established and greatest of all laws after the concealment of the beauty
of Muhammad. The traditions and reports attest to this, and because it
is so well known, there is no need to cite them. Indeed, in every age
the law of prayer has been firmly established and observed, as, from the
lights shining from the Sun of Muhammad,

\subsection{40}\label{section-230}

\begin{otherlanguage}{persian}

مأثور است که بر جمیع انبیا در هر عهدی حکم صلوة نازل شده نهایت آنکه در هر
عصر باقتضای وقت بقسمی و آدابی جدید مخصوص گشته و چون در هر ظهور بعد آداب
و عادات و علوم مرتفعهٔ محکمهٔ مشرقهٔ واضحهٔ ثابتۀ در ظهور قبل منسوخ میشود
لهذا تلویحاً باسم شمس و قمر ذکر نموده‌اند لیبلوکم ایّکم احسن عملاً

\end{otherlanguage}

Maʾs̱ūr ast ke bar jamīʿ-e anbiyā dar har ʿahdī ḥokm-e ṣalāt nāzel shode,
nehāyat ānke dar har ʿaṣr be-eqteżā-ye vaqt be-qesmī va ādābī jadīd
makhṣūṣ gashte. Va chon dar har ẓohūr-e baʿd ādāb va ʿādāt va ʿolūm-e
mortafeʿe-ye moḥkame-ye moshreqe-ye vāżeḥe-ye s̱ābete dar ẓohūr-e qabl
mansūkh mīshavad, lehādhā talvīḥan be-esm-e shams va qamar dhekr
namūde-and: li-yabluwakum ayyukum aḥsanu ʿamalan.

it is related that the law of prayer was revealed to all the prophets in
every age, except that in each age, according to the needs of the time,
it was given a particular form and new observances. And since in every
subsequent revelation the observances, customs, and teachings that were
exalted, firmly established, radiant, clear, and enduring in the
previous revelation are abrogated, they have therefore been referred to
figuratively as the sun and moon: ``so that He may test you as to which
of you is best in deeds.''

\subsection{41}\label{section-231}

\begin{otherlanguage}{persian}

و در حدیث هم اطلاق شمس و قمر بر صوم و صلوة شده چنانچه میفرماید الصّوم
ضیآء و الصّلوة نور ولکن روزی در محلّی نشسته بودم شخصی از علمای معروف وارد
شد و بمناسبتی این حدیث را ذکر نمود و فرمود چون صوم حرارت در مزاج احداث
مینماید لهذا بضیاء که شمس باشد تعبیر یافته و صلوة لیل چون برودت میطلبد
لهذا بنور که قمر باشد معبّر گشته ملاحظه نمودم که آن فقیر بقطره‌ئی از بحر
معانی موفّق نشده و بجذوه‌ئی از نار سدرهٔ حکمت ربّانی فائز نگشته بعد از مدّتی
در نهایت ادب اظهار داشتم که جناب آنچه فرمودید در معنی حدیث در السن و
افواه ناس مذکور است ولیکن گویا مقصود دیگر هم از حدیث مستفاد میشود بیان
آن را طلب نمود ذکر شد که خاتم انبیا و سیّد اصفیا دین مرتفع در فرقان را
تشبیه بسماء فرموده‌اند بعلّت علوّ و رفعت و عظمت و احاطهٔ آن بر جمیع ادیان و
چون در سماء ظاهره دو رکن اعظم اقوم مقرّر شده است که نیّرین باشد و بشمس و
قمر نامیده و همچنین در سماء دین هم دو نیّر مقدّر گشته که صوم و صلوة باشد
الاسلام سمآء و الصّوم شمسها و الصّلوة قمرها

\end{otherlanguage}

Va dar ḥadīs̱ ham eṭlāq-e shams va qamar bar ṣawm va ṣalāt shode,
chenānche mīfarmāyad: al-ṣawmu ḍiyāʾun wa-l-ṣalātu nūrun. Valāken rūzī
dar maḥallī neshaste būdam, shakhṣī az ʿolamā-ye maʿrūf vāred shod va
be-monāsebatī īn ḥadīs̱ rā dhekr namūd va farmūd: chon ṣawm ḥarārat dar
mezāj eḥdās̱ mīnamāyad, lehādhā be-ḍiyāʾ ke shams bāshad taʿbīr yāfte, va
ṣalāt-e leyl chon borūdat mīṭalabad, lehādhā be-nūr ke qamar bāshad
moʿabbar gashte. Molāḥeẓe namūdam ke ān faqīr be-qaṭreʾī az baḥr-e
maʿānī movaffaq nashode va be-jadhveʾī az nār-e sedre-ye ḥekmat-e
rabbānī fāʾez nagashte. Baʿd az moddatī dar nehāyat-e adab eẓhār dāshtam
ke: janāb, ānche farmūdīd dar maʿnī-ye ḥadīs̱ dar alson va afvāh-e nās
madhkūr ast, valīkan gūyā maqṣūd-e dīgar ham az ḥadīs̱ mostafād mīshavad.
Bayān-e ān rā ṭalab namūd. Dhekr shod ke Khātam-e Anbiyā va Sayyed-e
Aṣfiyā dīn-e mortafeʿ dar Forqān rā tashbīh be-samā farmūde-and,
be-ʿellat-e ʿolovv va refʿat va ʿaẓamat va eḥāṭe-ye ān bar jamīʿ-e
adyān. Va chon dar samā-ye ẓāhere do rokn-e aʿẓam-e aqvam moqarrar shode
ast ke nayyerayn bāshad va be-shams va qamar nāmīde, va hamchenīn dar
samā-ye dīn ham do nayyer moqaddar gashte ke ṣawm va ṣalāt bāshad:
al-Islāmu samāʾun wa-l-ṣawmu shamsuhā wa-l-ṣalātu qamaruhā.

In a tradition, too, the terms sun and moon have been applied to fasting
and prayer, as He says: ``Fasting is radiance, and prayer is light.''
One day, however, I was sitting in a place when a well-known religious
scholar arrived and, in the course of conversation, cited this
tradition. He said, ``Since fasting produces heat in the body, it is
described as radiance, which is the sun; and since night prayer calls
for coolness, it is described as light, which is the moon.'' I saw that
this poor man had not attained a single drop from the ocean of meanings
or a single ember from the fire of the Tree of divine wisdom. After a
while, with the utmost courtesy, I remarked, ``Sir, what you have said
about the meaning of this tradition is commonly repeated on people's
tongues and lips. Yet it seems that another meaning may also be
understood from it.'' He asked for an explanation. I said that the Seal
of the Prophets and Lord of the Chosen Ones had likened the religion
established in the Quran to heaven because of its loftiness, elevation,
grandeur, and encompassing of all religions. Just as in the visible
heaven two supremely great and steadfast pillars have been
established---the two luminaries called the sun and the moon---so in the
heaven of religion two luminaries have been ordained, namely fasting and
prayer: ``Islam is a heaven; fasting is its sun, and prayer its moon.''

\subsection{42}\label{section-232}

\begin{otherlanguage}{persian}

باری این است مقصود از تلویحات کلمات مظاهر الهی پس اطلاق شمس و قمر در این
مراتب بر این مقامات مذکوره بآیات نازله و اخبار وارده محقّق و ثابت شد
اینست که مقصود از ذکر تاریکی شمس و قمر و سقوط انجم ضلالت علما و نسخ شدن
احکام مرتفعۀ در شریعت است که مظهر آن ظهور باین تلویحات اخبار میدهد و جز
ابرار را از این کأس نصیبی نیست و جز اخیار را قسمتی نه انّ الأبرار یشربون
من کأس کان مزاجها کافوراً

\end{otherlanguage}

Bārī īn ast maqṣūd az talvīḥāt-e kalemāt-e maẓāher-e elāhī. Pas eṭlāq-e
shams va qamar dar īn marāteb bar īn maqāmāt-e madhkūre be-āyāt-e nāzele
va akhbār-e vārede moḥaqqaq va s̱ābet shod. Īn ast ke maqṣūd az dhekr-e
tārīkī-ye shams va qamar va soqūṭ-e anjom żalālat-e ʿolamā va naskh
shodan-e aḥkām-e mortafeʿe dar sharīʿat ast ke maẓhar-e ān ẓohūr be-īn
talvīḥāt ekhbār mīdahad. Va joz abrār rā az īn kaʾs naṣībī nīst va joz
akhyār rā qesmatī na: inna l-abrāra yashrabūna min kaʾsin kāna mizājuhā
kāfūran.

This, then, is what is meant by the allusions in the words of the
Manifestations of God. The application of the terms sun and moon, in
these instances, to the stations mentioned has thus been verified and
established through revealed verses and transmitted reports. It follows
that the darkening of the sun and moon and the falling of the stars
refer to the waywardness of the religious scholars and the abrogation of
the laws established in a religious dispensation, which the
Manifestation of that revelation foretells through these allusions. None
but the righteous have a share of this cup, and none but the good have a
portion of it: ``The righteous will drink from a cup whose mixture is
camphor.''

\subsection{43}\label{section-233}

\begin{otherlanguage}{persian}

و این مسلّم است که در هر ظهور بعد شمس علوم و احکام و اوامر و نواهی که در
ظهور قبل مرتفع شده و اهل آن عصر در ظلّ آن شمس و قمر معارف و اوامر منوّر و
مهتدی میشدند تاریک میشود یعنی حکمش و اثرش تمام میگردد و حال ملاحظه
فرمائید که اگر امّت انجیل مقصود از شمس و قمر را ادراک مینمودند و یا از
مظهر علم الهی مستفسر میشدند بدون اعتراض و لجاج البتّه معانی آن واضح میگشت
و اینگونه در ظلمت نفس و هوی مبتلا و گرفتار نمیشدند بلی چون علم را از
مبدء و معدنش اخذ ننمودند لهذا در وادی مهلک کفر و ضلالت بهلاکت رسیده‌اند و
هنوز مشعر نشدند که علامات کل ظاهر شد و شمس موعود از افق ظهور اشراق نمود
و شمس و قمر علوم و احکام و معارف قبل تاریک شد و غروب نمود

\end{otherlanguage}

Va īn mosallam ast ke dar har ẓohūr-e baʿd shams-e ʿolūm va aḥkām va
avāmer va navāhī ke dar ẓohūr-e qabl mortafeʿ shode va ahl-e ān ʿaṣr dar
ẓell-e ān shams va qamar-e maʿāref va avāmer monavvar va mohtadī
mīshodand tārīk mīshavad, yaʿnī ḥokmash va as̱arash tamām mīgardad. Va
ḥāl molāḥeẓe farmāyīd ke agar ommat-e Enjīl maqṣūd az shams va qamar rā
edrāk mīnamūdand va yā az maẓhar-e ʿelm-e elāhī mostafser mīshodand
bedūn-e eʿterāż va lejāj, albatte maʿānī-ye ān vāżeḥ mīgasht va īngūne
dar ẓolmat-e nafs va havā mobtalā va gereftār namīshodand. Balī chon
ʿelm rā az mabdaʾ va maʿdanash akhdh nanamūdand, lehādhā dar vādī-ye
mohleke-ye kofr va żalālat be-halākat resīde-and va hanūz moshʿer
nashodand ke ʿalāmāt koll ẓāher shod va shams-e mawʿūd az ofoq-e ẓohūr
eshrāq namūd va shams va qamar-e ʿolūm va aḥkām va maʿāref-e qabl tārīk
shod va ghorūb namūd.

It is certain that in every subsequent revelation the sun of the
teachings, laws, commands, and prohibitions established in the previous
revelation grows dark---that sun and moon of knowledge and commandments
beneath whose shelter the people of that age were illumined and guided.
That is, its authority and influence come to an end. Now consider: had
the followers of the Gospel understood what was meant by the sun and
moon, or had they asked the Manifestation of divine knowledge without
objection or obstinacy, their meanings would certainly have become
clear, and they would not have become so afflicted and entangled in the
darkness of self and desire. Indeed, because they did not acquire
knowledge from its source and mine, they have perished in the deadly
valley of unbelief and error. They still have not realized that all the
signs have appeared, that the promised Sun has shone from the horizon of
revelation, and that the sun and moon of the former teachings, laws, and
knowledge have grown dark and set.

\subsection{44}\label{section-234}

\begin{otherlanguage}{persian}

حال بچشم علم الیقین و جناحی عین الیقین بصراط حقّ الیقین قدم گذار قل اللّه
ثمّ ذرهم فی خوضهم یلعبون تا از اصحابی محسوب شوی که میفرماید انّ الّذین
قالوا ربّنا اللّه ثمّ استقاموا تتنزّل علیهم الملائکة تا جمیع این اسرار را
ببصر خود مشاهده فرمائی

\end{otherlanguage}

Ḥāl be-chashm-e ʿelm al-yaqīn va janāḥay-e ʿayn al-yaqīn be-ṣerāṭ-e ḥaqq
al-yaqīn qadam godhār: qul Allāhu thumma dharhum fī khawḍihim yalʿabūna.
Tā az aṣḥābī maḥsūb shavī ke mīfarmāyad: inna alladhīna qālū rabbunā
Allāhu thumma istaqāmū tatanazzalu ʿalayhimu l-malāʾikatu. Tā jamīʿ-e īn
asrār rā be-baṣar-e khod moshāhede farmāʾī.

Now, with the eye of certain knowledge and the two wings of certain
vision, set foot on the path of certain truth: ``Say: God! Then leave
them to play in their idle discourse.'' Thus you may be counted among
those companions of whom He says, ``Those who say, `Our Lord is God,'
and then remain steadfast---the angels descend upon them,'' and may
behold all these mysteries with your own eyes.

\subsection{45}\label{section-235}

\begin{otherlanguage}{arabic}

ای برادر من قدم روح بردار تا بادیه‌های بعیدهٔ بعد و هجر را بآنی طیّ فرمائی
و در رضوان قرب و وصل درآئی و در نفسی بانفس الهیّه فائز شوی و بقدم جسد
هرگز این مراحل طیّ نشود و مقصود حاصل نیاید و السّلام علی من اتّبع الحقّ
بالحقّ و کان علی صراط الأمر فی شاطئ العرفان باسم اللّه موقوفاً

\end{otherlanguage}

Ey barādar-e man, qadam-e rūḥ bardār tā bādiyehā-ye baʿīde-ye boʿd va
hajr rā be-ānī ṭayy farmāʾī va dar Reżvān-e qorb va vaṣl darāʾī va dar
nafasī be-anfos-e elāhīye fāʾez shavī. Va be-qadam-e jasad hargaz īn
marāḥel ṭayy nashavad va maqṣūd ḥāṣel nayāyad. Wa-l-salāmu ʿalā man
ittabaʿa l-Ḥaqqa bi-l-ḥaqqi wa-kāna ʿalā ṣirāṭi l-amri fī shāṭiʾi
l-ʿirfāni bi-smi Allāhi mawqūfan.

My brother, take the step of the spirit, so that in an instant you may
cross the far-stretching deserts of remoteness and separation, enter the
paradise of nearness and reunion, and in a single breath attain the
presence of the divine Souls. With bodily steps these stages can never
be traversed, nor can the goal be attained. Peace be upon the one who
follows the Truth through truth and stands, in the name of God, on the
path of the Cause upon the shore of recognition.

\subsection{46}\label{section-236}

\begin{otherlanguage}{persian}

اینست معنی آیهٔ مبارکه که میفرماید فلا اقسم بربّ المشارق و المغارب زیرا که
از برای هر شمسی از این شموس مذکوره محلّ اشراق و غروبست و چون علمای تفسیر
بر حقیقت این شمسهای مذکوره اطّلاع نیافتند لهذا در تفسیر این آیهٔ مبارکه
معطّل شدند و بعضی ذکر نمودند که چون آفتاب در هر روز از نقطه‌ئی طلوع
مینماید غیر از نقطهٔ یوم قبل لهذا بلفظ جمع ذکر فرموده و بعضی دیگر
نوشته‌اند که مقصود فصول اربعه است که در هر فصلی چون شمس از محلّی طالع
میشود و بمحلّی غروب مینماید لهذا مشارق و مغارب ذکر شده اینست مراتب علم
عباد و با وجود این بجواهر علم و لطائف حکمت چه جهلها و عیوبها که نسبت
میدهند

\end{otherlanguage}

Īn ast maʿnī-ye āye-ye mobārake ke mīfarmāyad: fa-lā uqsimu bi-rabbi
l-mashāriqi wa-l-maghāribi. Zīrā ke az barā-ye har shamsī az īn shomūs-e
madhkūre maḥall-e eshrāq va ghorūb ast. Va chon ʿolamā-ye tafsīr bar
ḥaqīqat-e īn shamshā-ye madhkūre eṭṭelāʿ nayāftand, lehādhā dar tafsīr-e
īn āye-ye mobārake moʿaṭṭal shodand. Va baʿżī dhekr namūdand ke chon
āftāb dar har rūz az noqṭeʾī ṭolūʿ mīnamāyad gheyr az noqṭe-ye yawm-e
qabl, lehādhā be-lafẓ-e jamʿ dhekr farmūde. Va baʿżī-ye dīgar
neveshte-and ke maqṣūd foṣūl-e arbaʿe ast ke dar har faṣlī chon shams az
maḥallī ṭāleʿ mīshavad va be-maḥallī ghorūb mīnamāyad, lehādhā mashāreq
va maghāreb dhekr shode. Īn ast marāteb-e ʿelm-e ʿebād, va bā vojūd-e īn
be-javāher-e ʿelm va laṭāʾef-e ḥekmat che jahlhā va ʿeybhā ke nesbat
mīdahand.

This is the meaning of the blessed verse in which He says, ``No! I swear
by the Lord of the places of rising and setting.'' For each of these
suns has a place of rising and setting. Since the commentators did not
discover the reality of these suns, they were at a loss to interpret
this blessed verse. Some said that because the sun rises each day from a
point different from that of the previous day, the plural form was used.
Others wrote that the four seasons are intended: because in each season
the sun rises in one place and sets in another, the expressions ``places
of rising'' and ``places of setting'' were used. Such is the extent of
people's knowledge; yet what ignorance and faults they attribute to the
Gems of knowledge and the Subtleties of wisdom!

\subsection{47}\label{section-237}

\begin{otherlanguage}{persian}

و همچنین از این بیانات واضحهٔ محکمهٔ متقنئ غیر متشابهه تفطّر سماء را که از
علائم ساعت و قیامت است ادراک نما اینست که میفرماید اذا السّمآء انفطرت
مقصود سماء ادیانست که در هر ظهور مرتفع میشود و بظهور بعد شکافته میگردد
یعنی باطل و منسوخ میشود قسم بخدا که اگر درست ملاحظه شود تفطّر این سماء
اعظم است از تفطّر سماء ظاهری قدری تأمّل فرمائید دینی که سالها مرتفع شده
باشد و جمیع در ظلّ آن نشو و نما نموده باشند و باحکام مشرقئ آن مدّتها تربیت
یافته و از آباء و اجداد جز ذکر آن را نشنیده بقسمی که چشمها جز نفوذ امرش
را ادراک نکرده و گوشها جز احکامش را استماع ننموده بعد نفسی ظاهر شود و
جمیع اینها را بقوّت و قدرت الهی تفریق نماید و فصل کند بلکه همه را نفی
فرماید حال فکر نما که این اعظم است یا آنچه این همج رعاع گمان نموده‌اند از
تفطّر سماء

\end{otherlanguage}

Va hamchenīn az īn bayānāt-e vāżeḥe-ye moḥkame-ye motqane-ye gheyr-e
motashābehe tafaṭṭor-e samā rā ke az ʿalāʾem-e sāʿat va qiyāmat ast
edrāk namā. Īn ast ke mīfarmāyad: idhā l-samāʾu infaṭarat. Maqṣūd
samā-ye adyān ast ke dar har ẓohūr mortafeʿ mīshavad va be-ẓohūr-e baʿd
shekāfte mīg ar dad, yaʿnī bāṭel va mansūkh mīshavad. Qasam be-Khodā ke
agar dorost molāḥeẓe shavad tafaṭṭor-e īn samā aʿẓam ast az tafaṭṭor-e
samā-ye ẓāherī. Qadrī taʾammol farmāyīd: dīnī ke sālhā mortafeʿ shode
bāshad va jamīʿ dar ẓell-e ān nashv va namā namūde bāshand va be-aḥkām-e
moshreqe-ye ān moddathā tarbiyat yāfte va az ābā va ajdād joz dhekr-e ān
rā nashnīde, be-qesmī ke chashmhā joz nofūdh-e amrash rā edrāk nakarde
va gūshhā joz aḥkāmash rā estemāʿ nanamūde, baʿd nafsī ẓāher shavad va
jamīʿ-e īnhā rā be-qovvat va qodrat-e elāhī tafrīq namāyad va faṣl
konad, balke hame rā nafy farmāyad. Ḥāl fekr namā ke īn aʿẓam ast yā
ānche īn hamaj-e raʿāʿ gomān namūde-and az tafaṭṭor-e samā.

Likewise, from these clear, firm, conclusive, and unambiguous
statements, understand the splitting of heaven, which is one of the
signs of the Hour and the Resurrection. Thus He says, ``When heaven is
split apart.'' What is meant is the heaven of religions, which is raised
in every revelation and split apart by the next---that is, it is
annulled and abrogated. By God, if considered rightly, the splitting of
this heaven is greater than the splitting of the visible sky. Reflect a
little: a religion has been established for years; everyone has grown
and developed under its shelter and has long been trained by its radiant
laws. From their fathers and forefathers they have heard of nothing
else, so that their eyes have perceived nothing but the sway of its
authority, and their ears have heard nothing but its laws. Then a Soul
appears and, through divine strength and power, breaks apart and severs
all these things---indeed, negates them all. Now consider whether this
is greater, or what this ignorant rabble has imagined the splitting of
heaven to be.

\subsection{48}\label{section-238}

\begin{otherlanguage}{arabic}

و دیگر زحمت و مرارت آن طلعات را ملاحظه نما که بی ‌ناصر و معین ظاهری در
مقابل جمیع اهل ارض اقامئ حدود اللّه میفرمایند با آن همه ایذا که بر آن
وجودهای مبارکهٔ لطیفئ رقیقه وارد میشود و با کمال قدرت صبر میفرمایند و با
نهایت غلبه تحمّل مینمایند

\end{otherlanguage}

Va dīgar zaḥmat va marārat-e ān ṭalʿathā rā molāḥeẓe namā ke bī-nāṣer va
moʿīn-e ẓāherī dar moqābel-e jamīʿ-e ahl-e arż eqāme-ye ḥodūd Allāh
mīfarmāyand, bā ān hame īdhā ke bar ān vojūdhā-ye mobārake-ye laṭīfe-ye
raqīqe vāred mīshavad, va bā kamāl-e qodrat ṣabr mīfarmāyand va bā
nehāyat-e ghalabe taḥammol mīnamāyand.

Consider also the hardship and bitterness endured by those radiant
Figures who, with no outward helper or supporter, establish the laws of
God in the face of all the peoples of the earth. Despite all the
suffering inflicted on those blessed, tender, and sensitive Beings, they
remain patient with consummate strength and endure with supreme
ascendancy.

\subsection{49}\label{section-239}

\begin{otherlanguage}{persian}

و همچنین معنی تبدیل ارض را ادراک نما که غمام رحمت آن سماء بر قلوبی که
نیسان مکرمت مبذول داشت تبدیل شد اراضی آن قلوب بارض معرفت و حکمت و چه
ریاحین توحید که در ریاض قلوبشان انبات شده و چه شقایقهای حقایق علم و حکمت
که از صدور منیرشان روئیده و اگر ارض قلوبشان تبدیل نمیشد چگونه رجالی که
حرفی تعلیم نگرفته‌اند و معلّم را ندیده‌اند و بهیچ دبستانی قدم نگذاشته‌اند
بکلمات و معارفی تکلّم مینمایند که احدی ادراک نتواند نمود گویا از تراب علم
سرمدی سرشته شده‌اند و از آب حکمت لدنّی عجین گشته‌اند اینست که میفرماید
العلم نور یقذفه اللّه فی قلب من یشآء و این نحو از علم است که ممدوح بوده و
هست نه علوم محدوده که از افکار محجوبهٔ کدره احداث شده و آن را گاهی از هم
سرقت مینمایند و بر دیگران افتخار میکنند

\end{otherlanguage}

Va hamchenīn maʿnī-ye tabdīl-e arż rā edrāk namā ke ghamām-e raḥmat-e ān
samā bar qolūbī ke naysān-e makramat mabdhūl dāsht, tabdīl shod arāżī-ye
ān qolūb be-arż-e maʿrefat va ḥekmat. Va che rayāḥīn-e tawḥīd ke dar
rīyāż-e qolūbeshān enbāt shode va che shaqāyeqhā-ye ḥaqāyeq-e ʿelm va
ḥekmat ke az ṣodūr-e monīreshān rūʾīde. Va agar arż-e qolūbeshān tabdīl
namīshod chegūne rejālī ke ḥarfī taʿlīm nagerefte-and va moʿallem rā
nadīde-and va be-hīch dabestānī qadam nagodhāshte-and be-kalemāt va
maʿārefī takallom mīnamāyand ke aḥadī edrāk natavānad namūd? Gūyā az
torāb-e ʿelm-e sarmadī sereshte shode-and va az āb-e ḥekmat-e ladonnī
ʿajīn gashte-and. Īn ast ke mīfarmāyad: al-ʿilmu nūrun yaqdhifuhu Allāhu
fī qalbi man yashāʾu. Va īn naḥv az ʿelm ast ke mamdūḥ būde va hast, na
ʿolūm-e maḥdūde ke az afkār-e maḥjūbe-ye kadere eḥdās̱ shode va ān rā
gāhī az ham serqat mīnamāyand va bar dīgarān eftekhār mīkonand.

Likewise, understand the meaning of the transformation of the earth:
upon whatever hearts the cloud of mercy from that heaven bestowed its
spring rain of bounty, the soil of those hearts was transformed into the
earth of knowledge and wisdom. What fragrant herbs of divine unity have
sprung up in the gardens of their hearts, and what anemones of the
realities of knowledge and wisdom have grown from their illumined
breasts! If the earth of their hearts had not been transformed, how
could men who have not learned a single letter, seen a teacher, or set
foot in any school utter words and knowledge that no one can comprehend?
It is as though they have been fashioned from the dust of eternal
knowledge and kneaded with the water of divinely bestowed wisdom. Thus
He says, ``Knowledge is a light that God casts into the heart of
whomever He wills.'' This kind of knowledge has always been, and
remains, praiseworthy---not the limited learning produced by veiled and
clouded minds, which people sometimes steal from one another and use to
vaunt themselves over others.

\subsection{50}\label{section-240}

\begin{otherlanguage}{persian}

ای کاش صدرهای عباد از نقوش این تحدیدات و کلمات مظلمه پاک و مقدّس میشد که
لعلّ بتجلّی انوار شمس علم و معانی و جواهر اسرار حکمت لدنّی فائز میگشت حال
ملاحظه نما اگر این اراضی جرزهٔ وجود تبدیل نمیشد چگونه محلّ ظهور اسرار
احدیّه و بروز جواهر هویّه میشد اینست که میفرماید یوم تبدّل الأرض غیر الأرض

\end{otherlanguage}

Ey kāsh ṣadrhā-ye ʿebād az noqūsh-e īn taḥdīdāt va kalemāt-e moẓleme pāk
va moqaddas mīshod ke laʿall be-tajallī-ye anvār-e shams-e ʿelm va
maʿānī va javāher-e asrār-e ḥekmat-e ladonnī fāʾez mīgasht. Ḥāl molāḥeẓe
namā agar īn arāżī-ye joroze-ye vojūd tabdīl namīshod chegūne maḥall-e
ẓohūr-e asrār-e aḥadīye va borūz-e javāher-e hovīye mīshod? Īn ast ke
mīfarmāyad: yawma tubaddalu l-arḍu ghayra l-arḍi.

If only people's hearts were cleansed and sanctified from the imprints
of these limitations and dark words, so that they might receive the
shining rays of the sun of knowledge and meanings, and the gems of the
mysteries of divinely bestowed wisdom. Now consider: if these barren
lands of existence were not transformed, how could they become the place
where the mysteries of divine oneness appear and the gems of the divine
Essence emerge? Thus He says, ``On the day when the earth is changed
into another earth.''

\subsection{51}\label{section-241}

\begin{otherlanguage}{arabic}

و از نسمات جود آن سلطان وجود ارض ظاهره هم تبدیل یافته لو انتم فی اسرار
الظّهور تتفکّرون

\end{otherlanguage}

Va az nasamāt-e jūd-e ān Solṭān-e Vojūd arż-e ẓāhere ham tabdīl yāfte,
law antum fī asrāri l-ẓuhūri tatafakkarūna.

Through the breezes of the bounty of that Sovereign of existence, the
outward earth, too, has been transformed, if you would but reflect on
the mysteries of revelation.

\subsection{52}\label{section-242}

\begin{otherlanguage}{persian}

و دیگر معنی این آیه را ادراک نما که میفرماید و الأرض جمیعاً قبضته یوم
القیامة و السّموات مطویّات بیمینه سبحانه و تعالی عمّا یشرکون مضمون آن اینست
که همهٔ زمین اخذ شده در دست اوست روز قیامت و آسمان پیچیده شده در دست راست
اوست حال قدری انصاف میخواهد که اگر مقصود اینست که مردم ادراک نموده‌اند چه
حسن بر آن مرتّب میشود وانگهی این مسلّم است که حقّ منیع دستی که مرئی شود
ببصر ظاهر و مرتکب این امورات شود منسوب بذات نیست بلکه کفری است محض و
افکی است صرف اقرار بر چنین امری و اگر بگوئی مظاهر امر او هستند که در
قیامت باین امر مأمور میشوند این هم بغایت بعید است و بی‌فائده بلکه مقصود
از ارض ارض معرفت و علم است و از سموات سموات ادیان حال ملاحظه فرما که
چگونه ارض علم و معرفت که از قبل مبسوط شده بود بقبضهٔ قدرت و اقتدار قبض
نمود و ارض منیعهٔ تازه در قلوب عباد مبسوط فرمود و ریاحین جدیده و گلهای
بدیعه و اشجار منیعه از صدور منیره انبات نمود

\end{otherlanguage}

Va dīgar maʿnī-yi īn āya rā idrāk namā ke mī-farmāyad: Wa-l-arḍu jamīʿan
qabḍatuhu yawma l-qiyāmati wa-l-samāwātu maṭwiyyātun bi-yamīnihi
subḥānahu wa-taʿālā ʿammā yushrikūn. Maḍmūn-i ān īn ast ke hama-yi zamīn
akhdh shuda dar dast-i ūst rūz-i qiyāmat va āsmān pīchīda shuda dar
dast-i rāst-i ūst. Ḥāl qadrī inṣāf mī-khvāhad ke agar maqṣūd īn ast ke
mardum idrāk namūda-and chi ḥusn bar ān mutarattib mī-shavad. Vāngahī īn
musallam ast ke Ḥaqq-i manīʿ dastī ke marʾī shavad bi-baṣar-i ẓāhir va
murtakib-i īn umūrāt shavad mansūb bi-dhāt nīst, balke kufrī-st maḥḍ va
ifkī-st ṣirf iqrār bar chunīn amrī. Va agar bi-gūʾī maẓāhir-i amr-i ū
hastand ke dar qiyāmat bi-īn amr maʾmūr mī-shavand, īn ham bi-ghāyat
baʿīd ast va bī-fāʾida. Balke maqṣūd az arḍ arḍ-i maʿrifat va ʿilm ast
va az samāvāt samāvāt-i adyān. Ḥāl mulāḥaẓa farmā ke chigūna arḍ-i ʿilm
va maʿrifat ke az qabl mabsūṭ shuda būd bi-qabḍa-yi qudrat va iqtidār
qabḍ namūd va arḍ-i manīʿa-yi tāza dar qulūb-i ʿibād mabsūṭ farmūd va
rayāḥīn-i jadīda va gulhā-yi badīʿa va ashjār-i manīʿa az ṣudūr-i munīra
inbāt namūd.

Further, grasp the meaning of this verse, in which He says: ``The whole
earth will be in His grasp on the Day of Resurrection, and the heavens
will be rolled up in His right hand. Glory be to Him, and exalted is He
above the partners they associate with Him.'' Its meaning is that on the
Day of Resurrection the whole earth is held in His hand, and heaven is
rolled up in His right hand. Now, a little fairness is needed: if the
intended meaning is what people have understood, what merit would follow
from it? Moreover, it is certain that no hand visible to the outward eye
and performing such acts can be attributed to the Essence of the exalted
God. Indeed, to affirm such a thing is sheer unbelief and utter
falsehood. And if you say that it is the Manifestations of His Cause who
are charged with this task at the Resurrection, this too is most
implausible and pointless. Rather, by the earth is meant the earth of
understanding and knowledge, and by the heavens, the heavens of
religions. Now consider how He gathered into the grasp of His power and
might the earth of knowledge and understanding that had previously been
spread out, spread a new and exalted earth within the hearts of His
servants, and caused fresh herbs, wondrous flowers, and lofty trees to
spring from radiant hearts.

\subsection{53}\label{section-243}

\begin{otherlanguage}{persian}

و همچنین ملاحظه کن که سموات ادیان مرتفعۀ در قبل چگونه در یمین قدرت
پیچیده شد و سماء بیان بامر اللّه مرتفع گشت و بشمس و قمر و نجوم اوامر
بدیعهٔ جدیده تزیین یافت اینست اسرار کلمات که بی‌حجاب کشف و ظاهر گشته تا
ادراک صبح معانی فرمائی و سراجهای ظنون و وهم و شک و ریب را بقوّت توکّل و
انقطاع خاموش نمائی و مصباح جدید علم و یقین در مشکاة قلب و دل برافروزی

\end{otherlanguage}

Va hamchunīn mulāḥaẓa kun ke samāvāt-i adyān-i murtafiʿa dar qabl
chigūna dar yamīn-i qudrat pīchīda shud va samā-yi Bayān bi-amr-i Allāh
murtafiʿ gasht va bi-shams va qamar va nujūm-i avāmir-i badīʿa-yi jadīda
tazyīn yāft. Īn ast asrār-i kalimāt ke bī-ḥijāb kashf va ẓāhir gashta tā
idrāk-i ṣubḥ-i maʿānī farmāʾī va sirājhā-yi ẓunūn va vahm va shakk va
rayb rā bi-quvvat-i tavakkul va inqiṭāʿ khāmūsh namāʾī va miṣbāḥ-i
jadīd-i ʿilm va yaqīn dar mishkāt-i qalb va dil bar-afrūzī.

Likewise, consider how the heavens of the religions raised up in the
past were rolled up in the right hand of power, how the heaven of the
Bayán was raised by the command of God, and how it was adorned with the
sun, moon, and stars of new and wondrous ordinances. These are the
mysteries of the words, unveiled and made manifest, so that you may
perceive the dawn of meanings, extinguish the lamps of conjecture,
fancy, doubt, and misgiving through the power of trust in God and
detachment, and kindle the new lamp of knowledge and certainty in the
niche of your inmost heart.

\subsection{54}\label{section-244}

\begin{otherlanguage}{persian}

و از جمیع این کلمات مرموزه و اشارات ملغزه که از مصادر امریّه ظاهر میشود
مقصود امتحان عباد است چنانچه مذکور شد تا معلوم شود اراضی قلوب جیّدهٔ منیره
از اراضی جرزهٔ فانیه و همیشه این از سنّت الهی در میان عباد بوده چنانچه در
کتب مسطور است

\end{otherlanguage}

Va az jamīʿ-i īn kalimāt-i marmūza va ishārāt-i mulghaza ke az maṣādir-i
amriyya ẓāhir mī-shavad maqṣūd imtiḥān-i ʿibād ast, chunānchi madhkūr
shud, tā maʿlūm shavad arāḍī-yi qulūb-i jayyida-yi munīra az arāḍī-yi
juruza-yi fāniya. Va hamīsha īn az sunnat-i ilāhī dar miyān-i ʿibād
būda, chunānchi dar kutub masṭūr ast.

The purpose of all these symbolic words and enigmatic allusions that
issue from the sources of divine authority is, as has been mentioned, to
test His servants, so that the fertile and radiant soil of hearts may be
distinguished from barren and perishable ground. This has always been
God's way among His servants, as is recorded in the Books.

\subsection{55}\label{section-245}

\begin{otherlanguage}{persian}

و همچنین آیهٔ قبله را ملاحظه فرمائید که بعد از هجرت شمس نبوّت محمّدی از
مشرق بطحا به یثرب رو به بیت المقدّس توجّه میفرمودند در وقت صلوة تا آنکه
یهود بعضی سخنهای ناشایسته بر زبان راندند که ذکرش شایستۀ این مقام نیست و
سبب تطویل کلام میشود باری آن حضرت بسیار مکدّر شدند و بلحاظ تفکّر و تحیّر در
سماء نظر میفرمودند بعد جبرئیل نازل شد و این آیه تلاوت نمود قد نری تقلّب
وجهک فی السّمآء فلنولّینّک قبلةً ترضاها تا آنکه در یومی آن حضرت با جمعی
اصحاب بفریضهٔ ظهر مشغول شدند و دو رکعت از نماز بجا آورده بودند که جبرئیل
نزول نمود و عرض کرد فولّ وجهک شطر المسجد الحرام در اثنای نماز حضرت از بیت
المقدّس انحراف جسته بکعبه مقابل شدند فی‌الحین تزلزل و اضطراب در میان اصحاب
افتاد بقسمی که جمعی نماز را بر هم زده اعراض نمودند این فتنه نبود مگر
برای امتحان عباد والّا آن سلطان حقیقی قادر بود که هیچ قبله را تغییر ندهد
و در آن عصر هم بیت المقدّس را قرار فرماید و این خلعت قبول را از وی سلب
ننماید

\end{otherlanguage}

Va hamchunīn āya-yi qibla rā mulāḥaẓa farmāʾīd ke baʿd az hijrat-i
shams-i nubuwwat-i Muḥammadī az mashriq-i Baṭḥā bi-Yathrib rū bi-Bayt
al-Muqaddas tavajjuh mī-farmūdand dar vaqt-i ṣalāt, tā ānke Yahūd baʿḍī
sukhanhā-yi nāshāyista bar zabān rāndand ke dhikrash shāyista-yi īn
maqām nīst va sabab-i taṭvīl-i kalām mī-shavad. Bārī ān Ḥaḍrat bisyār
mukaddar shudand va bi-liḥāẓ-i tafakkur va taḥayyur dar samā naẓar
mī-farmūdand. Baʿd Jabraʾīl nāzil shud va īn āya tilāvat namūd: Qad narā
taqalluba wajhika fī l-samāʾi fa-la-nuwalliyannaka qiblatan tarḍāhā. Tā
ānke dar yawmī ān Ḥaḍrat bā jamʿī aṣḥāb bi-farīḍa-yi ẓuhr mashghūl
shudand va du rakʿat az namāz bi-jā āvarda būdand ke Jabraʾīl nuzūl
namūd va ʿarḍ kard: Fa-walli wajhaka shaṭra l-masjidi l-ḥarām. Dar
athnā-yi namāz Ḥaḍrat az Bayt al-Muqaddas inḥirāf justa bi-Kaʿba muqābil
shudand. Fī l-ḥīn tazalzul va iḍṭirāb dar miyān-i aṣḥāb uftād bi-qismī
ke jamʿī namāz rā bar ham zada iʿrāḍ namūdand. Īn fitna nabūd magar
barā-yi imtiḥān-i ʿibād, va-illā ān Sulṭān-i ḥaqīqī qādir būd ke hīch
qibla rā taghyīr na-dahad va dar ān ʿaṣr ham Bayt al-Muqaddas rā qarār
farmāyad va īn khilʿat-i qabūl rā az vay salb na-namāyad.

Likewise, consider the verse concerning the direction of prayer. After
the Sun of Muḥammad's prophethood had migrated from its dawning-place in
Baṭḥá to Yathrib, He would turn toward Jerusalem at the time of prayer,
until the Jews uttered certain unseemly remarks that are not fitting to
mention here and would prolong the discussion. He became deeply
distressed and gazed toward heaven in thought and perplexity. Then
Gabriel descended and recited this verse: ``We have indeed seen your
face turning toward heaven; We will surely turn you toward a direction
of prayer that will please you.'' Then one day, while He and a group of
His companions were performing the obligatory noon prayer and had
completed two units of prayer, Gabriel descended and said: ``Turn your
face toward the Sacred Mosque.'' In the midst of prayer, He turned away
from Jerusalem and faced the Kaaba. Immediately, turmoil and agitation
arose among His companions, to such a degree that some broke off their
prayers and turned away. This trial was only to test His servants;
otherwise, that true Sovereign could have left the direction of prayer
entirely unchanged, retained Jerusalem in that age as well, and not
stripped it of this robe of acceptance.

\subsection{56}\label{section-246}

\begin{otherlanguage}{persian}

چنانچه در عهد اکثری انبیا که بعد از موسی مبعوث برسالت شدند مثل داود و
عیسی و دون آنها از انبیای اعظم که مابین این دو نبی آمدند هیچ حکم قبله
تغییر داده نشد و همهٔ این مرسلین از جانب ربّ العالمین مردم را بتوجّه همان
جهت امر میفرمودند و نسبت همهٔ اراضی هم بآن سلطان حقیقی یکیست مگر هر ارضی
را که در ظهور مظاهر خود تخصیص بامری دهد چنانچه میفرماید و للّه المشرق و
المغرب فأینما تولّوا فثمّ وجه اللّه با وجود تحقّق این امور چرا تبدیل شد که
سبب جزع و فزع عباد شود و علّت تزلزل و اضطراب اصحاب گردد بلی اینگونه امور
که سبب وحشت جمیع نفوس است واقع نمیشود مگر برای آنکه کل بمحکّ امتحان اللّه
درآیند تا صادق و کاذب از هم تمیز و تفصیل یابد اینست که بعد از اختلاف ناس
میفرماید و ما جعلنا القبلة الّتی کنت علیها الّا لنعلم من یتّبع الرّسول ممّن
ینقلب علی عقبیه که مضمون آن اینست ما نگردانیدیم و بر هم نزدیم قبله را که
آن بیت المقدّس باشد مگر آنکه بدانیم که متابعت تو مینماید و که راجع بر
عقبیه میشود یعنی اعراض مینماید و اطاعت نمیکند و صلوة را باطل نموده فرار
مینماید حمر مستنفرة فرّت من قسورة

\end{otherlanguage}

Chunānchi dar ʿahd-i aktharī anbiyā ke baʿd az Mūsā mabʿūth bi-risālat
shudand, mithl-i Dāvūd va ʿĪsā va dūn-i ānhā az anbiyā-yi aʿẓam ke
mābayn-i īn du nabī āmadand, hīch ḥukm-i qibla taghyīr dāda na-shud va
hama-yi īn mursalīn az jānib-i Rabb al-ʿālamīn mardum rā bi-tavajjuh-i
hamān jihat amr mī-farmūdand. Va nisbat-i hama-yi arāḍī ham bi-ān
Sulṭān-i ḥaqīqī yakī-st magar har arḍī rā ke dar ẓuhūr-i maẓāhir-i khud
takhṣīṣ bi-amrī dahad, chunānchi mī-farmāyad: Wa-li-llāhi l-mashriqu
wa-l-maghribu fa-aynamā tuwallū fa-thamma wajhu llāh. Bā vujūd-i
taḥaqquq-i īn umūr chirā tabdīl shud ke sabab-i jazaʿ va fazaʿ-i ʿibād
shavad va ʿillat-i tazalzul va iḍṭirāb-i aṣḥāb gardad? Balī īn-gūna umūr
ke sabab-i vaḥshat-i jamīʿ-i nufūs ast vāqiʿ na-mī-shavad magar barā-yi
ānke kull bi-miḥakk-i imtiḥān-i Allāh dar-āyand tā ṣādiq va kādhib az
ham tamyīz va tafṣīl yābad. Īn ast ke baʿd az ikhtilāf-i nās
mī-farmāyad: Wa-mā jaʿalnā l-qiblata llatī kunta ʿalayhā illā li-naʿlama
man yattabiʿu l-rasūla mimman yanqalibu ʿalā ʿaqibayh. Ke maḍmūn-i ān īn
ast: Mā na-gardānīdīm va bar ham na-zadīm qibla rā ke ān Bayt
al-Muqaddas bāshad magar ānke bi-dānīm ke mutābaʿat-i tu mī-namāyad va
ke rājiʿ bar ʿaqibayh mī-shavad, yaʿnī iʿrāḍ mī-namāyad va iṭāʿat
na-mī-kunad va ṣalāt rā bāṭil namūda firār mī-namāyad. Ḥumurun
mustanfiratun farrat min qaswara.

Thus, in the time of most of the prophets raised up as Messengers after
Moses, such as David and Jesus, and other great prophets who came
between these two, the ordinance concerning the direction of prayer was
never changed. All these Messengers, on behalf of the Lord of the
worlds, commanded people to turn in that same direction. All lands stand
in the same relation to that true Sovereign, except any land He singles
out for a particular purpose when His Manifestations appear. As He says:
``To God belong the East and the West; wherever you turn, there is the
face of God.'' These things being established, why was it changed,
causing His servants distress and alarm and His companions turmoil and
agitation? Indeed, such things, which fill every soul with dismay, occur
only so that all may be put to the touchstone of God's testing, and the
truthful may be distinguished and separated from the false. Thus, after
the disagreement among the people, He says: ``We appointed the direction
of prayer that you formerly observed only to know who follows the
Messenger and who turns back on his heels.'' Its meaning is this: We did
not change and set aside the direction of prayer, namely Jerusalem,
except to know who follows you and who turns back on his heels---that
is, who turns away, refuses to obey, invalidates his prayer, and flees.
``Like startled wild asses fleeing from a lion.''

\subsection{57}\label{section-247}

\begin{otherlanguage}{persian}

اگر قدری تأمّل رود در همین مطلب و بیان ابوابهای معانی و تبیان مفتوح بینید
و جمیع علم و اسرار آن را بی‌حجاب مشاهده فرمائید و نیست این امور مگر برای
تربیت و خلاصی نفوس از قفس نفس و هوی والّا آن سلطان حقیقی لم‌یزل بذات خود
غنی بوده از معرفت موجودات و لایزال بکینونت خود مستغنی خواهد بود از عبادت
ممکنات یک نسیم از غنای او جمیع عالم را بخلع غنا مفتخر نماید و یک قطره از
بحر جود او همهٔ هستی را بحیات باقیه مشرّف فرماید ولیکن چون مقصود امتیاز حقّ
از باطل و شمس از ظلّ است اینست که در کلّ حین امتحانهای منزله از جانب ربّ
العزّة چون غیث هاطل جاریست

\end{otherlanguage}

Agar qadrī taʾammul ravad dar hamīn maṭlab va bayān, abvābhā-yi maʿānī
va tibyān maftūḥ bīnīd va jamīʿ-i ʿilm va asrār-i ān rā bī-ḥijāb
mushāhada farmāʾīd. Va nīst īn umūr magar barā-yi tarbiyat va khalāṣī-yi
nufūs az qafas-i nafs va havā, va-illā ān Sulṭān-i ḥaqīqī lam-yazal
bi-dhāt-i khud ghanī būda az maʿrifat-i mawjūdāt va lā-yazāl
bi-kaynūnat-i khud mustaghnī khvāhad būd az ʿibādat-i mumkināt. Yak
nasīm az ghinā-yi ū jamīʿ-i ʿālam rā bi-khilaʿ-i ghinā muftakhir namāyad
va yak qaṭra az baḥr-i jūd-i ū hama-yi hastī rā bi-ḥayāt-i bāqiya
musharraf farmāyad. Va-līkin chūn maqṣūd imtiyāz-i ḥaqq az bāṭil va
shams az ẓill ast, īn ast ke dar kull-i ḥīn imtiḥānhā-yi munazzala az
jānib-i Rabb al-ʿIzza chūn ghayth-i hāṭil jārī-st.

If you reflect a little on this very subject and explanation, you will
see the gates of meanings and elucidation opened and behold all
knowledge and its mysteries unveiled. These things are only for the
education of souls and their deliverance from the cage of self and
desire. Otherwise, that true Sovereign has from eternity been, in His
Essence, independent of the recognition of created beings, and will
forever remain, in His own Being, independent of the worship of
contingent beings. A single breeze from His self-sufficiency would honor
the entire world with robes of sufficiency, and a single drop from the
ocean of His generosity would confer everlasting life upon all
existence. But since the purpose is to distinguish truth from falsehood
and the sun from the shadow, the tests sent down by the Lord of Might
pour forth at every moment like abundant rain.

\subsection{58}\label{section-248}

\begin{otherlanguage}{persian}

اگر قدری در انبیای قبل و ظهور ایشان تعقّل رود امر بسیار بر اهل دیار سهل
شود بقسمی که از افعال و اقوالی که مخالف نفس و هوی است محتجب نمیمانند و
همهٔ حجبات را بنار سدرهٔ عرفان محترق نمایند و بر عرش سکون و اطمینان مستریح
شوند مثلاً موسی بن عمران که یکی از انبیای معظّم و صاحب کتاب بود در اوّل امر
قبل از بعثت روزی در سوق میگذشت دو نفر با یکدیگر معارضه مینمودند یکی از
آن دو نفس از موسی استمداد جست آن حضرت او را اعانت نموده مدّعی را بقتل
رسانید چنانچه در کتاب مسطور است و ذکر تفصیل مایهٔ تعویق و تعطیل مقصود
میشود و این خبر در مدینه اشتهار یافت و آن حضرت را خوف غالب شد چنانچه نصّ
کتابست تا آنکه بخبر انّ الملأ یأتمرون بک لیقتلوک مخبر شد و از مدینه بیرون
تشریف بردند و در مدین در خدمت شعیب اقامه فرمودند و در مراجعت در وادی
مبارکه که برّیّهٔ سینا باشد وارد شد و تجلّی سلطان احدیّه را از شجرهٔ لاشرقیّه و
لاغربیّه مشاهده نمود و ندای جانفزای روحانی را از نار موقدهٔ ربّانی استماع
فرمود و مأمور بهدایت انفس فرعونی گشت تا مردم را از وادی نفس و هوی نجات
داده بصحراهای دلفزای روح و هدی وارد نماید و از سلسبیل انقطاع جمیع من فی
الابداع را از حیرت بعد بدار السّلام قرب رساند و چون در منزل فرعون وارد شد
و تبلیغ نمود بآنچه مأمور بود فرعون زبان به بی‌ادبی گشود و گفت آیا تو
نبودی که قتل نفس نمودی و از کافران شدی مثل اینکه ربّ العظمة خبر داد از
لسان فرعون که به موسی عرض نمود و فعلت فعلتک الّتی فعلت و انت من الکافرین
قال فعلتها اذاً و انا من الضّالّین ففررت منکم لمّا خفتکم فوهب لی ربّی حکماً و
جعلنی من المرسلین

\end{otherlanguage}

Agar qadrī dar anbiyā-yi qabl va ẓuhūr-i īshān taʿaqqul ravad, amr
bisyār bar ahl-i diyār sahl shavad, bi-qismī ke az afʿāl va aqvālī ke
mukhālif-i nafs va havā-st muḥtajib na-mī-mānand va hama-yi ḥujubāt rā
bi-nār-i Sidra-yi ʿirfān muḥtariq namāyand va bar ʿarsh-i sukūn va
iṭmīnān mustarīḥ shavand. Mathalan Mūsā ibn-i ʿImrān ke yakī az
anbiyā-yi muʿaẓẓam va ṣāḥib-i kitāb būd, dar avval-i amr qabl az biʿthat
rūzī dar sūq mī-gudhasht. Du nafar bā yakdīgar muʿāraḍa mī-namūdand.
Yakī az ān du nafs az Mūsā istimdād just. Ān Ḥaḍrat ū rā iʿānat namūda
muddaʿī rā bi-qatl rasānīd, chunānchi dar Kitāb masṭūr ast, va dhikr-i
tafṣīl māya-yi taʿvīq va taʿṭīl-i maqṣūd mī-shavad. Va īn khabar dar
madīna ishtihār yāft va ān Ḥaḍrat rā khawf ghālib shud, chunānchi naṣṣ-i
Kitāb ast, tā ānke bi-khabar-i Inna l-malaʾa yaʾtamirūna bika
li-yaqtulūka mukhbar shud va az madīna bīrūn tashrīf burdand va dar
Madyan dar khidmat-i Shuʿayb iqāma farmūdand. Va dar murājaʿat dar
vādī-yi mubāraka ke barriyya-yi Sīnā bāshad vārid shud va tajallī-yi
Sulṭān-i aḥadiyya rā az shajara-yi lā-sharqiyya va lā-gharbiyya
mushāhada namūd va nidā-yi jān-fazā-yi rūḥānī rā az nār-i mūqada-yi
rabbānī istimāʿ farmūd va maʾmūr bi-hidāyat-i anfus-i Firʿawnī gasht tā
mardum rā az vādī-yi nafs va havā najāt dāda bi-ṣaḥrāhā-yi dil-fazā-yi
rūḥ va hudā vārid namāyad va az salsabīl-i inqiṭāʿ jamīʿ-i man fī
l-ibdāʿ rā az ḥayrat-i buʿd bi-Dār al-Salām-i qurb rasānad. Va chūn dar
manzil-i Firʿawn vārid shud va tablīgh namūd bi-ānchi maʾmūr būd,
Firʿawn zabān bi-bī-adabī gushūd va guft: Āyā tu na-būdī ke qatl-i nafs
namūdī va az kāfirān shudī? Mithl-i īnke Rabb al-ʿAẓama khabar dād az
lisān-i Firʿawn ke bi-Mūsā ʿarḍ namūd: Wa-faʿalta faʿlataka llatī
faʿalta wa-anta mina l-kāfirīn. Qāla faʿaltuhā idhan wa-anā mina
l-ḍāllīn. Fa-farartu minkum lammā khiftukum fa-wahaba lī rabbī ḥukman
wa-jaʿalanī mina l-mursalīn.

If people were to reflect a little on the prophets of the past and their
appearance, the matter would become very easy for the inhabitants of
these lands. They would no longer remain veiled by deeds and words
contrary to self and desire, but would burn away every veil with the
fire of the Tree of recognition and rest upon the throne of tranquillity
and assurance. For example, Moses, son of ʿImrán, who was one of the
exalted prophets and the bearer of a Book, was passing through the
marketplace one day early in His life, before His prophetic mission. Two
men were disputing with one another. One of them sought Moses' help, and
He assisted him and killed his opponent, as is recorded in the Book. To
recount the details would delay and interrupt the purpose of this
discussion. News of this spread through the city, and fear overcame Him,
as the text of the Book states. Then He was informed that ``the chiefs
are plotting to kill you,'' and He left the city and took up residence
in Midian in the service of Shuʿayb. On His return, He entered the
blessed valley, the wilderness of Sinai, and beheld the revelation of
the Sovereign of Oneness from the Tree that is neither of the East nor
of the West. He heard the life-giving spiritual call from the divinely
kindled fire and was charged with guiding the people of Pharaoh, so that
He might deliver them from the valley of self and desire and lead them
into the heart-cheering expanses of spirit and guidance, and, through
the flowing waters of detachment, bring all who dwell in creation from
the bewilderment of remoteness to the Abode of Peace of nearness. When
He entered Pharaoh's dwelling and proclaimed what He had been charged to
convey, Pharaoh spoke insolently and said: ``Was it not you who killed a
man and became one of the unbelievers?'' Thus the Lord of Grandeur has
reported Pharaoh's words to Moses: ``You did that deed of yours, and you
are one of the unbelievers.'' He said: ``I did it then, when I was among
those who had gone astray. So I fled from you when I feared you; then my
Lord bestowed wisdom upon me and made me one of the Messengers.''

\subsection{59}\label{section-249}

\begin{otherlanguage}{persian}

حال تفکّر در فتنه‌های الهی و بدایع امتحانهای او کن که نفسی که معروفست بقتل
نفس و خود هم اقرار بر ظلم مینماید چنانچه در آیه مذکور است و سی سنه او
اقلّ هم بر حسب ظاهر در بیت فرعون تربیت یافته و از طعام و غذای او بزرگ شده
یکمرتبه او را از مابین عباد برگزیده و بامر هدایت کبری مأمور فرمود و حال
آنکه آن سلطان مقتدر قادر بر آن بود که موسی را از قتل ممنوع فرماید تا
باین اسم در بین عباد معروف نباشد که سبب وحشت قلوب شود و علّت احتراز نفوس
گردد

\end{otherlanguage}

Ḥāl tafakkur dar fitnahā-yi ilāhī va badāyiʿ-i imtiḥānhā-yi ū kun ke
nafsī ke maʿrūf ast bi-qatl-i nafs va khud ham iqrār bar ẓulm
mī-namāyad, chunānchi dar āya madhkūr ast, va sī sana aw aqall ham bar
ḥasab-i ẓāhir dar bayt-i Firʿawn tarbiyat yāfta va az ṭaʿām va ghadhā-yi
ū buzurg shuda, yak-martaba ū rā az mābayn-i ʿibād bar-guzīda va
bi-amr-i hidāyat-i kubrā maʾmūr farmūd. Va ḥāl ānke ān Sulṭān-i muqtadir
qādir bar ān būd ke Mūsā rā az qatl mamnūʿ farmāyad tā bi-īn ism dar
bayn-i ʿibād maʿrūf na-bāshad ke sabab-i vaḥshat-i qulūb shavad va
ʿillat-i iḥtirāz-i nufūs gardad.

Now reflect on the divine trials and the wonders of His tests: a person
known to have killed a man, who Himself acknowledges His wrongdoing, as
is mentioned in the verse, and who, to outward appearances, had been
brought up for thirty years or less in Pharaoh's house and had grown up
on his food and sustenance---suddenly God chose Him from among His
servants and charged Him with the task of supreme guidance. Yet that
mighty Sovereign could have prevented Moses from killing, so that He
would not be known among His servants by this reputation, causing dismay
in their hearts and leading souls to shun Him.

\subsection{60}\label{section-250}

\begin{otherlanguage}{persian}

و همچنین در حالت مریم مشاهده نما که آن طلعت کبری از عظمت امر و تحیّر
آرزوی عدم فرمود چنانچه مستفاد از آیهٔ مبارکه میشود که بعد از تولّد عیسی
مریم ناله نمود و باین کلمه زبان گشود یا لیتنی متّ قبل هذا و کنت نسیاً
منسیّا که ترجمهٔ آن این است ای کاش مرده بودم قبل از ظهور این امر و بودم از
فراموش‌شدگان قسم بخدا که کبدها از استماع این سخن میگدازد و روانها میریزد
و این اضطراب و حزن نبود مگر از شماتت اعدا و اعتراض اهل کفر و شقا آخر
تفکّر نمائید که مریم چه جواب با مردم میگفت طفلی که پدر او معیّن نباشد
چگونه میتوان بمردم معیّن نمود که این از روح القدس است این بود که آن مخدّرهٔ
بقا آن طفل را برداشته بمنزل مراجعت فرمود تا چشم قوم بر او افتاد گفتند یا
اخت هارون ما کان ابوک امرء سوء و ما کانت امّک بغیّاً مضمون آن این است که ای
خواهر هارون نبود پدر تو مرد بدی و نبود مادر تو بدکار

\end{otherlanguage}

Va hamchunīn dar ḥālat-i Maryam mushāhada namā ke ān Ṭalʿat-i kubrā az
ʿaẓamat-i amr va taḥayyur ārzū-yi ʿadam farmūd, chunānchi mustafād az
āya-yi mubāraka mī-shavad ke baʿd az tavallud-i ʿĪsā Maryam nāla namūd
va bi-īn kalima zabān gushūd: Yā laytanī mittu qabla hādhā wa-kuntu
nasyan mansiyyan. Ke tarjuma-yi ān īn ast: Ay kāsh murda būdam qabl az
ẓuhūr-i īn amr va būdam az farāmūsh-shudagān. Qasam bi-Khudā ke kabidhā
az istimāʿ-i īn sukhan mī-gudāzad va ravānhā mī-rīzad. Va īn iḍṭirāb va
ḥuzn nabūd magar az shamātat-i aʿdā va iʿtirāḍ-i ahl-i kufr va shaqā.
Ākhir tafakkur namāʾīd ke Maryam chi javāb bā mardum mī-guft? Ṭiflī ke
pidar-i ū muʿayyan na-bāshad chigūna mī-tavān bi-mardum muʿayyan namūd
ke īn az Rūḥ al-Qudus ast? Īn būd ke ān Mukhaddara-yi baqā ān ṭifl rā
bar-dāshta bi-manzil murājaʿat farmūd. Tā chashm-i qawm bar ū uftād
guftand: Yā ukhta Hārūna mā kāna abūki imraʾa sawʾin wa-mā kānat ummuki
baghiyyan. Maḍmūn-i ān īn ast ke ay khvāhar-i Hārūn, na-būd pidar-i tu
mard-i badī va na-būd mādar-i tu badkār.

Likewise, consider Mary's condition: that exalted Countenance,
overwhelmed by the gravity of the matter and by bewilderment, wished for
nonexistence. This is understood from the blessed verse stating that,
after the birth of Jesus, Mary lamented and uttered these words: ``If
only I had died before this and been utterly forgotten.'' Its
translation is: ``If only I had died before this matter came to pass and
been among the forgotten.'' By God, hearts melt and souls dissolve upon
hearing these words. This anguish and sorrow arose only from the taunts
of enemies and the objections of the people of unbelief and wickedness.
Reflect, then: what answer could Mary give the people? How could she
make it clear to them that a child whose father was unknown was of the
Holy Spirit? So that veiled Lady of eternity took the child and returned
home. As soon as her people caught sight of her, they said: ``O sister
of Aaron, your father was not an evil man, nor was your mother
unchaste.'' Its meaning is: ``O sister of Aaron, your father was not a
wicked man, and your mother was not a woman of loose morals.''

\subsection{61}\label{section-251}

\begin{otherlanguage}{arabic}

حال ناظر باین فتنهٔ کبری و امتحان اعظم شوید و از همه گذشته همان جوهر روح
که در میان قوم بنسبت بی‌پدری معروف بوده او را پیغمبری بخشید و حجّت خود
نمود بر کلّ اهل سموات و ارض

\end{otherlanguage}

Ḥāl nāẓir bi-īn fitniy-i kubrā va imtiḥān-i aʿẓam shavīd va az hamih
gudhashtih hamān jawhar-i rūḥ kih dar miyān-i qawm bi-nisbat-i bī-pidarī
maʿrūf būd ū rā payghambarī bakhshīd va ḥujjat-i khud nimūd bar kull-i
ahl-i samāvāt va arḍ.

Now consider this supreme ordeal and this greatest test. Beyond all
this, He bestowed prophethood upon that very Essence of the Spirit who
was known among the people as fatherless, and made Him His proof to all
the inhabitants of the heavens and the earth.

\subsection{62}\label{section-252}

\begin{otherlanguage}{persian}

حال مشاهده فرمائید که چه‌ قدر امور مظاهر ظهور مغایر نفس و هوای عباد از
سلطان ایجاد ظاهر میشود و چون بر این جواهر اسرار مطّلع شوی بمقصود آن نگار
اطّلاع یابی و اقوال و افعال آن ملیک بااقتدار را مثل هم ملاحظه نمائی بقسمی
که آنچه در افعال او مشاهده شود در کلمات او هم ملاحظه گردد و هر چه در
کلمات او ملاحظه گردد در افعال او بنظر آید اینست که این افعال و اقوال در
ظاهر نقمتند برای فجّار و در باطن رحمتند برای ابرار اگر بدیدهٔ قلب ملاحظه
رود کلمات منزله از سماء مشیّت با امور مظهره از ملکوت قدرت یک شیء مشاهده
شود و بر یک قسم ادراک گردد چنانچه مذکور شد

\end{otherlanguage}

Ḥāl mushāhidih farmāyīd kih chih qadr umūr-i maẓāhir-i ẓuhūr mughāyir-i
nafs va havā-yi ʿibād az Sulṭān-i ījād ẓāhir mīshavad va chūn bar īn
javāhir-i asrār muṭṭaliʿ shavī bi-maqṣūd-i ān Nigār iṭṭilāʿ yābī va
aqvāl va afʿāl-i ān Malīk-i bā-iqtidār rā mithl-i ham mulāḥiẓih namāyī
bi-qismī kih ānchih dar afʿāl-i ū mushāhidih shavad dar kalimāt-i ū ham
mulāḥiẓih gardad va har chih dar kalimāt-i ū mulāḥiẓih gardad dar
afʿāl-i ū bi-naẓar āyad īnast kih īn afʿāl va aqvāl dar ẓāhir niqmatand
barā-yi fujjār va dar bāṭin raḥmatand barā-yi abrār agar bi-dīdiy-i qalb
mulāḥiẓih ravad kalimāt-i munazzalih az samā-yi mashiyyat bā umūr-i
muẓharih az malakūt-i qudrat yak shayʾ mushāhidih shavad va bar yak qism
idrāk gardad chunānchih madhkūr shud.

Now observe how greatly the affairs of the Manifestations of Revelation,
brought forth by the Sovereign of creation, run counter to the desires
and inclinations of His servants. When you become acquainted with these
gems of mystery, you will understand the purpose of that Beloved and
regard the words and deeds of that mighty King as alike, so that
whatever is seen in His deeds may also be discerned in His words, and
whatever is discerned in His words may be seen in His deeds. Thus these
deeds and words are outwardly a punishment for the wicked, but inwardly
a mercy for the righteous. If viewed with the eye of the heart, the
words sent down from the heaven of His will and the events brought forth
from the kingdom of His power will be seen as one and understood in the
same way, as has been mentioned.

\subsection{63}\label{section-253}

\begin{otherlanguage}{persian}

حال ای برادر ملاحظه نما اگر در این عهد چنین اموری ظاهر شود و چنین حکایت
بروز نماید چه خواهند نمود قسم بمربّی وجود و منزل کلمات که در حین بی‌تکلّم
حکم بر کفر و امر بر قتل نمایند کجا گوش میدهند که گفته شود که عیسی از
نفخهٔ روح القدس ظاهر شده و یا موسی از امر مبرم مأمور گشته اگر صد‌هزار خروش
برآری بگوش احدی نرود که بی‌پدری مبعوث برسالت گشته و یا قاتلی از شجرهٔ نار
انّی انا اللّه آورده

\end{otherlanguage}

Ḥāl ay barādar mulāḥiẓih namā agar dar īn ʿahd chunīn umūrī ẓāhir shavad
va chunīn ḥikāyat burūz namāyad chih khvāhand nimūd qasam bi-Murabbī-yi
vujūd va Munzil-i kalimāt kih dar ḥīn bī-takallum ḥukm bar kufr va amr
bar qatl namāyand kujā gūsh mīdihand kih guftih shavad kih ʿĪsā az
nafkhiy-i Rūḥ al-Qudus ẓāhir shudih va yā Mūsā az amr-i mubram maʾmūr
gashtih agar ṣad-hizār khurūsh bar ārī bi-gūsh-i aḥadī naravad kih
bī-pidarī mabʿūth bi-risālat gashtih va yā qātilī az shajariy-i nār innī
anā Allāh āvardih.

Now consider, my brother: if such things were to appear in this age and
such a story were to unfold, what would people do? I swear by the One
who nurtures all existence and sends down the words: at once, without a
word, they would pronounce a verdict of unbelief and order an execution.
How would they listen to anyone saying that Jesus had appeared through
the breath of the Holy Spirit, or that Moses had been commissioned by an
irrevocable command? Though you cry out a hundred thousand times, no one
would listen to the assertion that a fatherless one had been sent as a
Messenger, or that a murderer had brought from the tree of fire the
words, ``Truly, I am God.''

\subsection{64}\label{section-254}

\begin{otherlanguage}{persian}

چشم انصاف اگر باز شود از جمیع این بیانات مشهود میگردد که مظهر همهٔ این
امور و نتیجهٔ همه الیوم ظاهر است با اینکه امثال این امور در این ظهور واقع
نشده با وجود این متمسّک بظنونات انفس مردوده شده چه نسبتها که داده‌اند و چه
بلایا که وارد آورده‌اند که در ابداع شبه آن بظهور نیامده

\end{otherlanguage}

Chashm-i inṣāf agar bāz shavad az jamīʿ-i īn bayānāt mashhūd mīgарdad
kih maẓhar-i hamiy-i īn umūr va natījiy-i hamih al-yawm ẓāhir ast bā
īnkih amthāl-i īn umūr dar īn ẓuhūr vāqiʿ nashudih bā vujūd-i īn
mutamassik bi-ẓunūnāt-i anfus-i mardūdih shudih chih nisbat-hā kih
dādih-and va chih balāyā kih vārid āvardih-and kih dar ibdāʿ shibh-i ān
bi-ẓuhūr nayāmadih.

If the eye of fairness is opened, all these statements make it evident
that the Manifestation of all these things, and the culmination of them
all, is manifest today. Although nothing like these events has occurred
in this Revelation, people have nevertheless clung to the idle
suppositions of rejected souls. What accusations they have made, and
what afflictions they have inflicted, the like of which has never
appeared in creation!

\subsection{65}\label{section-255}

\begin{otherlanguage}{persian}

اللّه اکبر بیان که باین مقام رسید رائحهٔ روحانی از صبح صمدانی مرور نمود و
صبای صبحگاهی از مدینهٔ سبای لا‌یزالی وزید و اشارتش جان را بشارت تازه بخشید
و روح را فتوحی بی‌اندازه بساط جدیدی مبسوط نمود و ارمغان بیشمار بیکران از
آن یار بی‌نشان آورد که خلعت ذکر از قدّ لطیفش بسی قاصر است و رداء بیان از
قامت منیرش بس کوتاه بی‌لفظ رمز معانی کشف مینماید و بی‌لسان اسرار تبیان
میگوید و بلبلهای شاخسار هجر و فراق را ناله و افغان میآموزد و قاعده و
رسوم عشق و عاشقی و رمز دلدادگی تعلیم مینماید و گلهای بدیع رضوان قرب و
وصال را رسم دلبری و آداب عشوه‌گری تلقین مینماید و اسرار حقایق بر شقایق
بستان عشق میبخشد و دقایق رموز و رقایق آن را در صدر عشّاق ودیعه میگذارد
بقسمی عنایت در این ساعت فرموده که روح القدس بغایت حسرت میبرد قطره را
امواج بحری داده و ذرّه را طراز خورشیدی عنایت نموده الطاف بمقامی رسیده که
جعل قصد نافهٔ مشک نموده و خفّاش در مقابل آفتاب مقرّ گزیده مردگان را بنفخۀ
حیات از قبور جسد مبعوث نموده و جاهلان را بر صدر علم منزل داده و ظالمان
را بر فراز عدل محل معین نموده

\end{otherlanguage}

Allāhu akbar bayān kih bi-īn maqām rasīd rāʾiḥiy-i rūḥānī az ṣubḥ-i
ṣamadānī murūr nimūd va ṣabā-yi ṣubḥgāhī az madīniy-i Sabā-yi lā-yazālī
vazīd va ishāratash jān rā bishārat-i tāzih bakhshīd va rūḥ rā futūḥī
bī-andāzih bisāṭ-i jadīdī mabsūṭ nimūd va armughān-i bīshumār-i bīkarān
az ān Yār-i bī-nishān āvard kih khilʿat-i dhikr az qadd-i laṭīfash basī
qāṣir ast va ridā-yi bayān az qāmat-i munīrash bas kūtāh bī-lafẓ ramz-i
maʿānī kashf mīnamāyad va bī-lisān asrār-i tibyān mīgūyad va
bulbul-hā-yi shākhsār-i hajr va firāq rā nālih va afghān mīāmūzad va
qāʿidih va rusūm-i ʿishq va ʿāshiqī va ramz-i dildādigī taʿlīm mīnamāyad
va gul-hā-yi badīʿ-i Riḍvān-i qurb va viṣāl rā rasm-i dilbarī va ādāb-i
ʿishvih-garī talqīn mīnamāyad va asrār-i ḥaqāyiq bar shaqāyiq-i būstān-i
ʿishq mībakhshad va daqāyiq-i rumūz va raqāyiq-i ān rā dar ṣadr-i
ʿushshāq vadīʿih mīgudhārad bi-qismī ʿināyat dar īn sāʿat farmūdih kih
Rūḥ al-Qudus bi-ghāyat ḥasrat mībarad qaṭrih rā amvāj-i baḥrī dādih va
dharrih rā ṭirāz-i khurshīdī ʿināyat nimūdih alṭāf bi-maqāmī rasīdih kih
juʿal qaṣd-i nāfiy-i mushk nimūdih va khuffāsh dar muqābil-i āftāb
maqarr guzīdih murdigān rā bi-nafkhiy-i ḥayāt az qubūr-i jasad mabʿūth
nimūdih va jāhilān rā bar ṣadr-i ʿilm manzil dādih va ẓālimān rā bar
farāz-i ʿadl maḥall muʿayyan nimūdih.

God is Most Great! When the discourse reached this point, a spiritual
fragrance wafted from the everlasting Dawn, and the morning breeze blew
from the city of eternal Sheba. Its intimation brought fresh glad
tidings to the soul and boundless joy to the spirit. It spread out a new
array and brought countless, limitless gifts from that traceless
Friend---gifts whose delicate form the robe of mention is far too short
to clothe, and whose luminous stature the mantle of expression cannot
encompass. Without words it discloses the mysteries of meaning, and
without a tongue it speaks the secrets of exposition. It teaches the
nightingales on the branches of separation and estrangement to lament
and cry out, and instructs them in the ways and customs of love and
loving, and the secret of surrendering the heart. It teaches the
wondrous flowers of the paradise of nearness and reunion the art of
captivating hearts and the manners of enchantment. It bestows the
secrets of realities upon the anemones in the garden of love and
deposits the subtleties and refinements of its mysteries in the breasts
of lovers. Such favor has it shown at this hour that the Holy Spirit
yearns for it intensely. It has given the drop the waves of a sea and
bestowed upon the atom the splendor of a sun. Its favors have reached
such a height that the dung beetle has sought the musk pod, and the bat
has made its home before the sun. With the breath of life it has raised
the dead from the graves of the body, given the ignorant a seat of honor
in knowledge, and assigned the unjust a place on the heights of justice.

\subsection{66}\label{section-256}

\begin{otherlanguage}{persian}

و عالم هستی بجمیع این عنایات حامله گشته تا کی اثر این عنایت غیبی در
خاکدان ترابی ظاهر شود و تشنگان از پا افتاده را بکوثر زلال محبوب رساند و
گمگشتگان صحرای بعد و نیستی را بسرادق قرب و هستی معشوق فائز گرداند و در
ارض قلوب که این حبّه‌های قدس انبات نماید و از ریاض نفوس که شقایقهای حقایق
غیبی بشکفد باری نچنان سدرهٔ عشق در سینای حبّ مشتعل شده که به آبهای بیان
افسرده گردد و یا اتمام پذیرد عطش این حوت را بحور ننشاند و این سمندر ناری
جز در نار روی یار مقرّ نگزیند پس ای برادر سراج روح را در مشکاة قلب بدهن
حکمت برافروز و بزجاج عقل حفظش نما تا نفسهای انفس مشرکه آن را خاموش نکند
و از نور بازندارد کذلک نوّرنا افق سمآء البیان من انوار شموس الحکمة و
العرفان لیطمئنّ بها قلبک و تکون من الّذین طاروا بأجنحة الایقان فی هوآء
محبّة ربّهم الرّحمن

\end{otherlanguage}

Va ʿālam-i hastī bi-jamīʿ-i īn ʿināyāt ḥāmilih gashtih tā kay athar-i īn
ʿināyat-i ghaybī dar khākdān-i turābī ẓāhir shavad va tishnigān-i az pā
uftādih rā bi-Kawthar-i zulāl-i Maḥbūb rasānad va gumgashtigān-i
ṣaḥrā-yi buʿd va nīstī rā bi-surādiq-i qurb va hastī-yi Maʿshūq fāʾiz
gardānad va dar arḍ-i qulūb-i kih īn ḥabbih-hā-yi quds inbāt namāyad va
az riyāḍ-i nufūs-i kih shaqāyiq-hā-yi ḥaqāyiq-i ghaybī bishkufad bārī
nachunān sidriy-i ʿishq dar Sīnā-yi ḥubb mushtaʿil shudih kih
bi-āb-hā-yi bayān afsurdih gardad va yā itmām padhīrad ʿaṭash-i īn ḥūt
rā buḥūr ninishānad va īn samandar-i nārī juz dar nār-i rūy-i Yār maqarr
naguzīnad pas ay barādar sirāj-i rūḥ rā dar mishkāt-i qalb bi-duhn-i
ḥikmat bar afrūz va bi-zujāj-i ʿaql ḥifẓash namā tā nafas-hā-yi anfus-i
mushrikih ān rā khāmūsh nakunad va az nūr bāz nadārad kadhālika nawwarnā
ufuqa samāʾi al-bayān min anwāri shumūsi al-ḥikmati wa-al-ʿirfān
li-yaṭmaʾinna bihā qalbuka wa-takūna mina alladhīna ṭārū bi-ajniḥati
al-īqān fī hawāʾi maḥabbati rabbihimu al-Raḥmān.

And the world of being has become pregnant with all these favors. When
will the effect of this unseen grace appear in this earthly realm, bring
those who have collapsed from thirst to the Beloved's crystal-clear
Kawthar, and admit those lost in the wilderness of remoteness and
nonexistence to the pavilion of nearness and of the Loved One's
existence? In the soil of whose hearts will these seeds of holiness
sprout, and from the gardens of whose souls will the anemones of unseen
realities blossom? In any event, the tree of love is not so ablaze upon
the Sinai of affection that the waters of words could quench or exhaust
it. Oceans cannot slake the thirst of this fish, and this fiery
salamander will make its home nowhere but in the fire of the Friend's
face. So, my brother, kindle the lamp of the spirit within the niche of
the heart with the oil of wisdom, and guard it with the glass of
understanding, so that the breaths of souls who associate partners with
God may neither extinguish it nor keep it from shining. Thus have We
illumined the horizon of the heaven of utterance with the lights of the
suns of wisdom and spiritual knowledge, so that your heart may be
reassured by them and you may be among those who have flown on the wings
of certitude in the atmosphere of the love of their Lord, the
All-Merciful.

\subsection{67}\label{section-257}

\begin{otherlanguage}{persian}

و قوله حینئذ یظهر علامة ابن الانسان فی السّمآء میفرماید بعد از کسوف شمس
معارف الهیّه و سقوط نجوم احکام مثبته و خسوف قمر علم که مربّی عباد است و
انعدام اعلام هدایت و فلاح و ظلمت صبح صدق و صلاح ظاهر میشود علامت ابن
الانسان در آسمان و مقصود از سماء سماء ظاهره است که قریب ظهور آن فلک
سموات معدلت و جریان فلک هدایت بر بحر عظمت در آسمان نجمی بر حسب ظاهر پیدا
میشود که مبشّر است خلق سموات را بظهور آن نیّر اعظم و همچنین در آسمان معنی
نجمی ظاهر میشود که مبشّر است اهل ارض را بآن فجر اقوم اکرم و این دو علامت
در سماء ظاهره و سماء باطنه قبل از ظهور هر نبی ظاهر گشته چنانچه شنیده‌اند

\end{otherlanguage}

Va qawluhu ḥīnaʾidhin yaẓharu ʿalāmatu ibni al-insāni fī al-samāʾ
mīfarmāyad baʿd az kusūf-i shams-i maʿārif-i ilāhiyyih va suqūṭ-i
nujūm-i aḥkām-i muthbitih va khusūf-i qamar-i ʿilm kih murabbī-yi ʿibād
ast va inʿidām-i aʿlām-i hidāyat va falāḥ va ẓulmat-i ṣubḥ-i ṣidq va
ṣalāḥ ẓāhir mīshavad ʿalāmat-i Ibn al-Insān dar āsimān va maqṣūd az samā
samā-yi ẓāhirih ast kih qarīb-i ẓuhūr-i ān falak-i samāvāt-i maʿdalat va
jarayān-i fulk-i hidāyat bar baḥr-i ʿaẓamat dar āsimān najmī bar ḥasab-i
ẓāhir paydā mīshavad kih mubashshir ast khalq-i samāvāt rā bi-ẓuhūr-i ān
Nayyir-i Aʿẓam va hamchunīn dar āsimān-i maʿnī najmī ẓāhir mīshavad kih
mubashshir ast ahl-i arḍ rā bi-ān fajr-i aqvam-i akram va īn du ʿalāmat
dar samā-yi ẓāhirih va samā-yi bāṭinih qabl az ẓuhūr-i har nabī ẓāhir
gashtih chunānchih shinīdih-and.

And His words, ``Then the sign of the Son of Man will appear in
heaven.'' He says that after the eclipse of the sun of divine knowledge,
the falling of the stars of established ordinances, the eclipse of the
moon of learning that nurtures His servants, the disappearance of the
banners of guidance and salvation, and the darkening of the dawn of
truth and righteousness, the sign of the Son of Man will appear in
heaven. By heaven is meant the visible heaven: as the appearance of that
Orb of the heavens of justice and the sailing of the Ark of guidance
upon the sea of grandeur draw near, a star appears in the visible sky,
announcing to the inhabitants of the heavens the appearance of that Most
Great Luminary. Likewise, a star appears in the heaven of inner meaning,
announcing to the people of the earth that most steadfast and most noble
Dawn. These two signs in the outward and inward heavens have appeared
before the advent of every Prophet, as people have heard.

\subsection{68}\label{section-258}

\begin{otherlanguage}{persian}

از جمله خلیل الرّحمن که قبل از ظهور آن حضرت نمرود خوابی دید و کهنه را
خواست اخبار دادند بر طلوع نجمی در سماء و همچنین شخصی در ارض ظاهر شد مردم
را بشارت میداد بظهور آن حضرت

\end{otherlanguage}

Az jumlih Khalīl al-Raḥmān kih qabl az ẓuhūr-i ān Ḥaḍrat Nimrūd khvābī
dīd va kahanih rā khvāst ikhbār dādand bar ṭulūʿ-i najmī dar samā va
hamchunīn shakhṣī dar arḍ ẓāhir shud mardum rā bishārat mīdād bi-ẓuhūr-i
ān Ḥaḍrat.

Among them was the Friend of the All-Merciful. Before His appearance,
Nimrod had a dream and summoned the soothsayers, who foretold the rising
of a star in heaven. Likewise, a person appeared on earth who gave the
people glad tidings of His advent.

\subsection{69}\label{section-259}

\begin{otherlanguage}{persian}

و بعد از او حکایت کلیم اللّه بود که کهنهٔ آن زمان فرعون را خبر دادند که
کوکبی در سماء طالع شده که دالّ است بر انعقاد نطفه‌ئی‌ که هلاک تو و قوم تو
بر دست اوست و همچنین عالمی پیدا شد که شبها بنی‌اسرائیل را بشارت و تسلّی
میفرمود و اطمینان میداد چنانچه در کتب مسطور است و اگر تفصیل این امور ذکر
شود این رساله کتابی میشود و دیگر آنکه دوست ندارم حکایات واقعهٔ قبل را ذکر
نمایم و خدا شاهد حال است که این بیان هم که میشود نیست مگر از کمال حبّ بآن
جناب که شاید جمعی فقرای ارض بر شاطی غنا وارد شوند و یا گروهی از جاهلان
بر بحر علم وارد گردند و یا تشنگان معرفت بر سلسبیل حکمت واصل آیند والّا
این عبد اشتغال باین مقالات را ذنبی عظیم میدانم و عصیانی کبیر میشمرم

\end{otherlanguage}

Va baʿd az ū ḥikāyat-i Kalīm Allāh būd kih kahaniy-i ān zamān Firʿawn rā
khabar dādand kih kawkabī dar samā ṭāliʿ shudih kih dāll ast bar
inʿiqād-i nuṭfih-ī kih halāk-i tu va qawm-i tu bar dast-i ūst va
hamchunīn ʿālimī paydā shud kih shab-hā Banī-Isrāʾīl rā bishārat va
tasallī mīfarmūd va iṭmīnān mīdād chunānchih dar kutub masṭūr ast va
agar tafṣīl-i īn umūr dhikr shavad īn risālih kitābī mīshavad va dīgar
ānkih dūst nadāram ḥikāyāt-i vāqiʿiy-i qabl rā dhikr namāyam va Khudā
shāhid-i ḥāl ast kih īn bayān ham kih mīshavad nīst magar az kamāl-i
ḥubb bi-ān janāb kih shāyad jamʿī fuqarā-yi arḍ bar shāṭī-yi ghinā vārid
shavand va yā gurūhī az jāhilān bar baḥr-i ʿilm vārid gardand va yā
tishnigān-i maʿrifat bar Salsabīl-i ḥikmat vāṣil āyand va-illā īn ʿabd
ishtighāl bi-īn maqālāt rā dhanbī ʿaẓīm mīdānam va ʿiṣyānī kabīr
mīshumaram.

After Him came the story of the One who conversed with God. The
soothsayers of that time informed Pharaoh that a star had risen in
heaven, indicating the conception of one at whose hands he and his
people would perish. Likewise, a learned man appeared who, by night,
brought glad tidings, comfort, and reassurance to the children of
Israel, as is recorded in the books. If these matters were recounted in
detail, this treatise would become a book. Moreover, I do not wish to
recount stories of past events. God bears witness that even these
statements are made only out of the fullness of My love for you, so that
perhaps some of the poor of the earth may reach the shore of wealth, or
a group of the ignorant may approach the ocean of knowledge, or those
thirsting for understanding may attain the Salsabil of wisdom.
Otherwise, I, this servant, consider engagement in such discussions a
grave sin and regard it as a great transgression.

\subsection{70}\label{section-260}

\begin{otherlanguage}{persian}

و همچنین نزدیک ظهور عیسی شد چند نفر از مجوس که اطّلاع یافتند بر ظهور نجم
عیسی در سماء باثر آن نجم آمدند تا داخل شدند بشهری که مقرّ سلطنت هیرودس
بود و در آن ایّام سلطنت آن ممالک در قبضهٔ تصرّف او بود

\end{otherlanguage}

Va hamchunīn nazdīk-i ẓuhūr-i ʿĪsā shud chand nafar az majūs kih iṭṭilāʿ
yāftand bar ẓuhūr-i najm-i ʿĪsā dar samā bi-athar-i ān najm āmadand tā
dākhil shudand bi-shahrī kih maqarr-i salṭanat-i Hīrūdis būd va dar ān
ayyām salṭanat-i ān mamālik dar qabḍiy-i taṣarruf-i ū būd.

Likewise, when the advent of Jesus drew near, several of the Magi,
having learned of the appearance of the star of Jesus in heaven,
followed that star until they entered the city that was the seat of
Herod's rule. In those days, sovereignty over those lands was in his
grasp.

\subsection{71}\label{section-261}

\begin{otherlanguage}{persian}

و کانوا قائلین این هو المولود ملک الیهود لأنّا قد رأینا نجمه فی المشرق و
وافینا لنسجد له و بعد از تفحّص معلوم نمودند که در بیت‌اللّحم یهود آن طفل
متولّد شد این علامت در سماء ظاهره و علامت در سماء باطنه که سماء علم و
معانی باشد ظهور یحیی بن زکریّا بود که مردم را بشارت میداد بظهور آن حضرت
چنانچه میفرماید انّ اللّه یبشّرک بیحیی مصدّقاً بکلمة من اللّه و سیّداً و حصورا
مقصود از کلمه حضرت عیسی است که یحیی مبشّر بظهور او بود و در الواح سماوی
هم مسطور است کان یوحنّا یکرز فی برّیّة یهودا قائلاً توبوا فقد اقترب ملکوت
السّموات و مقصود از یوحنّا یحیی است

\end{otherlanguage}

Wa-kānū qāʾilīna ayna huwa al-mawlūdu maliku al-Yahūd li-annā qad
raʾaynā najmahu fī al-mashriqi wa-wāfaynā li-nasjuda lahu va baʿd az
tafaḥḥuṣ maʿlūm nimūdand kih dar Bayt al-Laḥm-i Yahūd ān ṭifl mutavallid
shud īn ʿalāmat dar samā-yi ẓāhirih va ʿalāmat dar samā-yi bāṭinih kih
samā-yi ʿilm va maʿānī bāshad ẓuhūr-i Yaḥyā ibn Zakariyyā būd kih mardum
rā bishārat mīdād bi-ẓuhūr-i ān Ḥaḍrat chunānchih mīfarmāyad inna Allāha
yubashshiruka bi-Yaḥyā muṣaddiqan bi-kalimatin mina Allāhi wa-sayyidan
wa-ḥaṣūran maqṣūd az kalimih Ḥaḍrat-i ʿĪsā ast kih Yaḥyā mubashshir
bi-ẓuhūr-i ū būd va dar alvāḥ-i samāvī ham masṭūr ast kāna Yūḥannā
yukarrizu fī barriyyati Yahūdhā qāʾilan tūbū fa-qad iqtaraba malakūtu
al-samāwāt va maqṣūd az Yūḥannā Yaḥyā ast.

They said, ``Where is the newborn King of the Jews? For we have seen His
star in the east and have come to bow down before Him.'' After making
inquiries, they learned that the child had been born in Bethlehem of
Judea. This was the sign in the outward heaven. The sign in the inward
heaven, which is the heaven of knowledge and inner meanings, was the
appearance of John, son of Zechariah, who gave the people glad tidings
of His advent. As He says, ``God gives you glad tidings of John, who
will confirm a Word from God, noble and chaste.'' By the Word is meant
Jesus, whose advent John heralded. It is also recorded in the heavenly
Tablets: ``John was preaching in the wilderness of Judea, saying,
`Repent, for the kingdom of heaven has drawn near.'\,'' By Yuhanna is
meant Yahya.

\subsection{72}\label{section-262}

\begin{otherlanguage}{persian}

و همچنین قبل از ظهور جمال محمّدی آثار سماء ظاهره ظاهر شد و آثار باطنه که
مردم را در ارض بشارت میدادند بظهور آن شمس هویّه چهار نفر بودند واحداً بعد
واحد چنانچه روزبه که موسوم به سلمان شد بشرف خدمتشان مشرّف بود و زمان وفات
هر یک میرسید روزبه را نزد دیگری میفرستاد تا نوبت به چهارم رسید و او در
حین موت فرمود ای روزبه بعد از تکفین و تدفین من برو به حجاز که شمس محمّدی
اشراق مینماید و بشارت باد ترا بلقای آن حضرت

\end{otherlanguage}

Va hamchunīn qabl az ẓuhūr-i Jamāl-i Muḥammadī āthār-i samā-yi ẓāhirih
ẓāhir shud va āthār-i bāṭinih kih mardum rā dar arḍ bishārat mīdādand
bi-ẓuhūr-i ān Shams-i huviyyih chahār nafar būdand wāḥidan baʿda wāḥid
chunānchih Rūzbih kih mawsūm bi-Salmān shud bi-sharaf-i khidmatishān
musharraf būd va zamān-i vafāt-i har yak mīrasīd Rūzbih rā nazd-i dīgarī
mīfiristād tā nawbat bi-chahārum rasīd va ū dar ḥīn-i mawt farmūd ay
Rūzbih baʿd az takfīn va tadfīn-i man biraw bi-Ḥijāz kih Shams-i
Muḥammadī ishrāq mīnamāyad va bishārat bād turā bi-liqā-yi ān Ḥaḍrat.

Likewise, before the appearance of the Beauty of Muhammad, signs
appeared in the outward heaven. The inward signs, who gave the people on
earth glad tidings of the advent of that Sun of divine Being, were four
persons, one after another. Ruzbih, who came to be known as Salman, was
privileged to serve them. As the time of each one's death approached, he
would send Ruzbih to the next, until the turn of the fourth came. At the
hour of his death, he said, ``Ruzbih, after you have shrouded and buried
me, go to the Hijaz, for the Sun of Muhammad is dawning. Glad tidings to
you of meeting Him!''

\subsection{73}\label{section-263}

\begin{otherlanguage}{arabic}

تا رسید باین امر بدیع منیع و اکثر از منجّمان خبر ظهور نجم را در سماء
ظاهره داده‌اند و همچنین در ارض هم نورین نیّرین احمد و کاظم قدّس‌اللّه تربتهما

\end{otherlanguage}

Tā rasīd bi-īn Amr-i badīʿ-i manīʿ va akthar az munajjimān khabar-i
ẓuhūr-i najm rā dar samā-yi ẓāhirih dādih-and va hamchunīn dar arḍ ham
nūrayn-i nayyirayn Aḥmad va Kāẓim qaddasa Allāhu turbatahumā.

Until it came to this wondrous, exalted Cause: most of the astronomers
have reported the appearance of a star in the outward heaven. Likewise,
on earth there were the two shining lights, Ahmad and Kazim---may God
sanctify their resting places.

\subsection{74}\label{section-264}

\begin{otherlanguage}{persian}

پس از این معانی مبرهن شد که قبل از ظهور هر یک از مرایای احدیّه علامات آن
ظهور در آسمان ظاهر و آسمان باطن که محلّ شمس علم و قمر حکمت و انجم معانی و
بیانست ظاهر میشود و آن ظهور انسان کاملست قبل از هر ظهور برای تربیت و
استعداد عباد از برای لقای آن شمس هویّه و قمر احدیّه

\end{otherlanguage}

Pas az īn maʿānī mubarhan shud kih qabl az ẓuhūr-i har yak az marāyā-yi
aḥadiyyih ʿalāmāt-i ān ẓuhūr dar āsimān-i ẓāhir va āsimān-i bāṭin kih
maḥall-i shams-i ʿilm va qamar-i ḥikmat va anjum-i maʿānī va bayānast
ẓāhir mīshavad va ān ẓuhūr-i insān-i kāmil ast qabl az har ẓuhūr barā-yi
tarbiyat va istiʿdād-i ʿibād az barā-yi liqā-yi ān Shams-i huviyyih va
Qamar-i aḥadiyyih.

These considerations thus establish that, before the appearance of each
of the Mirrors of divine unity, signs of that Revelation appear in the
outward heaven and in the inward heaven, which is the seat of the sun of
knowledge, the moon of wisdom, and the stars of inner meaning and
utterance. This is the appearance, before each Revelation, of a perfect
human being to educate and prepare the servants for meeting that Sun of
divine Being and Moon of unity.

\subsection{75}\label{section-265}

\begin{otherlanguage}{persian}

و قوله و ینوح کلّ قبائل الأرض و یرون ابن الانسان آتیاً علی سحاب السّمآء مع
قوّات و مجد کبیر تلویح این بیان اینست یعنی در آن وقت نوحه میکنند عباد از
جهت فقدان شمس جمال الهی و قمر علم و انجم حکمت لدنّی و در آن اثنا مشاهده
میشود که آن طلعت موعود و جمال معبود از آسمان نازل میشود در حالتی که بر
ابر سوار است یعنی آن جمال الهی از سموات مشیّت ربّانی در هیکل بشری ظهور
میفرماید و مقصود از سماء نیست مگر جهت علوّ و سموّ که آن محلّ ظهور آن مشارق
قدسیّه و مطالع قدمیّه است و این کینونات قدیمه اگرچه بحسب ظاهر از بطن امّهات
ظاهر میشوند ولیکن فی‌الحقیقه از سموات امر نازلند و اگرچه بر ارض ساکنند
ولیکن بر رفرف معانی متّکیند و در حینی که میان عباد مشی مینمایند در هواهای
قرب طائرند بی‌حرکت رجل در ارض روح مشی نمایند و بی‌پر بمعارج احدیّه پرواز
فرمایند در هر نفسی مشرق و مغرب ابداع را طیّ فرمایند و در هر آنی ملکوت غیب
و شهاده را سیر نمایند بر عرش لا یشغله شأن عن شأن واقفند و بر کرسیّ کلّ یوم
هو فی شأن ساکن از علوّ قدرت سلطان قدم و سموّ مشیّت ملیک اعظم مبعوث میشوند
اینست که میفرماید از آسمان نازل میشود

\end{otherlanguage}

Va qawluhu: ``Wa yanūḥu kullu qabāʾili al-arḍi wa yarawna ibna al-insāni
ātiyan ʿalā saḥābi al-samāʾi maʿa quwwātin wa majdin kabīr.'' Talvīḥ-i
īn bayān īnast, yaʿnī dar ān vaqt nawḥih mīkunand ʿibād az jihat-i
fuqdān-i shams-i jamāl-i ilāhī va qamar-i ʿilm va anjum-i ḥikmat-i
ladunnī, va dar ān a s̱nā mushāhidih mīshavad kih ān ṭalʿat-i mawʿūd va
jamāl-i maʿbūd az āsmān nāzil mīshavad dar ḥālatī kih bar abr savār ast.
Yaʿnī ān jamāl-i ilāhī az samāvāt-i mashīyyat-i rabbānī dar haykal-i
basharī ẓuhūr mīfarmāyad. Va maqṣūd az samāʾ nīst magar jihat-i ʿuluvv
va sumuvv kih ān maḥall-i ẓuhūr-i ān mashāriq-i qudsīyyih va maṭāliʿ-i
qidamīyyih ast. Va īn kaynūnāt-i qadīmih agarchih bi-ḥasab-i ẓāhir az
baṭn-i ummahāt ẓāhir mīshavand, valīkin fī al-ḥaqīqih az samāvāt-i amr
nāziland. Va agarchih bar arḍ sākinand, valīkin bar rafraf-i maʿānī
muttakiʾand. Va dar ḥīnī kih miyān-i ʿibād mashy mīnamāyand, dar
havāhā-yi qurb ṭāʾirand. Bī-ḥarakat-i rijl dar arḍ-i rūḥ mashy namāyand
va bī-par bi-maʿārij-i aḥadīyyih parvāz farmāyand. Dar har nafasī
mashriq va maghrib-i ibdāʿ rā ṭayy farmāyand va dar har ānī malakūt-i
ghayb va shahādih rā sayr namāyand. Bar ʿarsh-i ``lā yashghaluhu shaʾnun
ʿan shaʾn'' vāqifand va bar kursī-yi ``kulla yawmin huwa fī shaʾn''
sākin. Az ʿuluvv-i qudrat-i sulṭān-i qidam va sumuvv-i mashīyyat-i
malīk-i aʿẓam mabʿū s̱ mīshavand. Īnast kih mīfarmāyad az āsmān nāzil
mīshavad.

And His words: ``All the tribes of the earth will mourn, and they will
see the Son of Man coming upon the clouds of heaven with power and great
glory.'' The allusion in this statement is this: at that time the
servants will mourn the loss of the sun of divine beauty, the moon of
knowledge, and the stars of divinely bestowed wisdom. Meanwhile, that
promised Countenance and adored Beauty will be seen descending from
heaven, riding upon a cloud. That is, that divine Beauty appears from
the heavens of the Lord's will in a human temple. By heaven is meant
only the state of exaltation and loftiness, which is the place of
appearance of those sacred Daysprings and eternal Dawning-places.
Although these eternal Beings outwardly emerge from their mothers'
wombs, in reality they descend from the heavens of the divine command.
Although they dwell on earth, they recline upon the celestial couch of
inner meanings. While walking among the servants, they soar in the skies
of nearness. Without moving their feet, they walk upon the earth of the
spirit, and without wings they fly to the heights of oneness. With every
breath they traverse the east and west of creation, and at every moment
they journey through the kingdom of the unseen and the seen. They stand
upon the throne of ``No concern distracts Him from another'' and dwell
upon the seat of ``Every day He is engaged in a task.'' They are sent
forth from the heights of the power of the Sovereign of eternity and the
loftiness of the will of the supreme King. This is why He says that He
descends from heaven.

\subsection{76}\label{section-266}

\begin{otherlanguage}{persian}

و لفظ سماء در بیانات شموس معانی بر مراتب کثیره اطلاق میشود مثلاً سماء امر
و سماء مشیّت و سماء اراده و سماء عرفان و سماء ایقان و سماء تبیان و سماء
ظهور و سماء بطون و امثال آن و در هر مقام از لفظ سماء معنیی اراده
میفرماید که غیر از واقفین اسرار احدیّه و شاربین کؤوس ازلیّه احدی ادراک
ننماید مثلاً میفرماید و فی السّمآء رزقکم و ما توعدون و حال آنکه رزق از ارض
انبات مینماید و همچنین الأسمآء تنزل من السّمآء با اینکه از لسان عباد
اسماء ظاهر میشود اگر قدری مرآت قلب را از غبار غرض پاک و لطیف فرمائی جمیع
تلویحات کلمات کلمۀ جامعهٔ ربوبیّه را در هر ظهوری ادراک مینمائی و بر اسرار
علم واقف میشوی ولیکن تا حجبات علمیّه را که مصطلح بین عباد است بنار انقطاع
نسوزانی بصبح نورانی علم حقیقی فائز نگردی

\end{otherlanguage}

Va lafẓ-i samāʾ dar bayānāt-i shumūs-i maʿānī bar marātib-i ka s̱īrih
iṭlāq mīshavad, ma s̱alan samāʾ-i amr va samāʾ-i mashīyyat va samāʾ-i
irādih va samāʾ-i ʿirfān va samāʾ-i īqān va samāʾ-i tibyān va samāʾ-i
ẓuhūr va samāʾ-i buṭūn va am s̱āl-i ān. Va dar har maqām az lafẓ-i samāʾ
maʿnīyī irādih mīfarmāyad kih ghayr az vāqifīn-i asrār-i aḥadīyyih va
shāribīn-i kuʾūs-i azalīyyih aḥadī idrāk nanamāyad. Ma s̱alan mīfarmāyad:
``Wa fī al-samāʾi rizqukum wa mā tūʿadūn,'' va ḥāl ānkih rizq az arḍ
inbāt mīnamāyad. Va hamchunīn: ``Al-asmāʾu tanzilu mina al-samāʾ,'' bā
īnkih az lisān-i ʿibād asmāʾ ẓāhir mīshavad. Agar qadrī mirʾāt-i qalb rā
az ghubār-i gharaḍ pāk va laṭīf farmāʾī, jamīʿ-i talvīḥāt-i kalimāt-i
kalimih-yi jāmiʿih-yi rubūbīyyih rā dar har ẓuhūrī idrāk mīnamāʾī va bar
asrār-i ʿilm vāqif mīshavī. Valīkin tā ḥujubāt-i ʿilmīyyih rā kih
muṣṭalaḥ bayn-i ʿibād ast bi-nār-i inqiṭāʿ nasūzānī, bi-ṣubḥ-i nūrānī-yi
ʿilm-i ḥaqīqī fāʾiz nagardī.

The word ``heaven'' is applied to many degrees in the utterances of the
Suns of inner meaning: for example, the heaven of the divine command,
the heaven of will, the heaven of purpose, the heaven of recognition,
the heaven of certainty, the heaven of exposition, the heaven of
manifestation, the heaven of concealment, and the like. In each
instance, He intends by the word ``heaven'' a meaning that no one can
grasp except those acquainted with the mysteries of oneness and those
who drink from the cups of eternity. For example, He says: ``In heaven
are your sustenance and what you are promised,'' whereas sustenance
grows from the earth. Likewise: ``Names descend from heaven,'' although
names issue from the tongues of the servants. If you cleanse and polish
the mirror of your heart even a little from the dust of prejudice, you
will grasp all the allusions in the utterances of the all-embracing Word
of divinity in every Revelation and become acquainted with the mysteries
of knowledge. But until you burn away, with the fire of detachment, the
veils of learning current among the servants, you will not attain the
radiant dawn of true knowledge.

\subsection{77}\label{section-267}

\begin{otherlanguage}{persian}

و علم بدو قسم منقسم است علم الهی و علم شیطانی آن از الهامات سلطان حقیقی
ظاهر و این از تخیّلات انفس ظلمانی باهر معلّم آن حضرت باری و معلّم این وساوس
نفسانی بیان آن اتّقوا اللّه یعلّمکم اللّه و بیان این العلم حجاب الأکبر اثمار
آن شجر صبر و شوق و عرفان و محبّت و اثمار این شجر کبر و غرور و نخوت و از
بیانات صاحبان بیان که در معنی علم فرموده‌اند هیچ رائحهٔ این علوم ظلمانی که
ظلمت آن همهٔ بلاد را فراگرفته استشمام نمیشود این شجر جز بغی و فحشا ثمری
نیاورد و جز غل و بغضا حاصلی نبخشد ثمرش سمّ قاتل است و ظلّش نار مهلک فنعم
ما قال تمسّک بأذیال الهوی و اخلع الحیا و خلّ سبیل النّاسکین و ان جلّوا

\end{otherlanguage}

Va ʿilm bi-du qism munqasim ast: ʿilm-i ilāhī va ʿilm-i shayṭānī. Ān az
ilhāmāt-i sulṭān-i ḥaqīqī ẓāhir va īn az takhayyulāt-i anfus-i ẓulmānī
bāhir. Muʿallim-i ān Ḥaḍrat-i Bārī va muʿallim-i īn vasāvis-i nafsānī.
Bayān-i ān: ``Ittaqū Allāha yuʿallimukumu Allāh,'' va bayān-i īn:
``Al-ʿilmu ḥijābu al-akbar.'' A s̱mār-i ān shajar ṣabr va shawq va ʿirfān
va maḥabbat, va a s̱mār-i īn shajar kibr va ghurūr va nikhvat. Va az
bayānāt-i ṣāḥibān-i bayān kih dar maʿnī-yi ʿilm farmūdih-and hīch
rāʾiḥih-yi īn ʿulūm-i ẓulmānī kih ẓulmat-i ān hamih-yi bilād rā
farāgiriftih istishmām namīshavad. Īn shajar juz baghy va faḥshā s̱amarī
nayāvarad va juz ghill va baghḍā ḥāṣilī nabakhshad. S̱amarash samm-i
qātil ast va ẓillash nār-i muhlik. Fa-niʿma mā qāl: ``Tamassak
bi-adhyāli al-hawā wa-ikhlaʿi al-ḥayā wa khalli sabīla al-nāsikīna wa in
jallū.''

Knowledge is divided into two kinds: divine knowledge and satanic
knowledge. The former springs from the inspirations of the true
Sovereign; the latter arises from the imaginings of darkened souls. The
teacher of the former is the Creator; the teacher of the latter is the
promptings of the self. The former is expressed in ``Fear God; God will
teach you,'' and the latter in ``Knowledge is the greatest veil.'' The
fruits of the former tree are patience, longing, recognition, and love;
the fruits of the latter are pride, arrogance, and conceit. In the
utterances of the masters of exposition concerning the meaning of
knowledge, not a trace can be detected of these dark sciences whose
darkness has encompassed every land. This tree bears no fruit but
oppression and indecency, and yields nothing but malice and hatred. Its
fruit is deadly poison, and its shade a consuming fire. How well it has
been said: ``Cling to the hem of desire and cast off shame; leave the
path of the ascetics, however exalted they may be.''

\subsection{78}\label{section-268}

\begin{otherlanguage}{persian}

پس باید صدر را از جمیع آنچه شنیده شده پاک نمود و قلب را از همهٔ تعلّقات
مقدّس فرمود تا محلّ ادراک الهامات غیبی شود و خزینهٔ اسرار علوم ربّانی گردد
اینست که میفرماید السّالک فی النّهج البیضآء و الرّکن الحمرآء لن یصل الی
مقام وطنه الّا بالکفّ الصّفر عمّا فی ایدی النّاس این است شرط سالک درست تفکّر و
تعقّل فرموده تا بی‌حجاب بر مقصود کتاب واقف شوی

\end{otherlanguage}

Pas bāyad ṣadr rā az jamīʿ-i ānchih shunīdih shudih pāk namūd va qalb rā
az hamih-yi taʿalluqāt muqaddas farmūd tā maḥall-i idrāk-i ilhāmāt-i
ghaybī shavad va khazīnih-yi asrār-i ʿulūm-i rabbānī gardad. Īnast kih
mīfarmāyad: ``Al-sāliku fī al-nahji al-bayḍāʾ wa al-rukni al-ḥamrāʾ lan
yaṣila ilā maqāmi waṭanihi illā bi-al-kaffi al-ṣifr ʿammā fī aydī
al-nās.'' Īn ast sharṭ-i sālik. Durust tafakkur va taʿaqqul farmūdih tā
bī-ḥijāb bar maqṣūd-i kitāb vāqif shavī.

One must therefore cleanse the breast of all that has been heard and
sanctify the heart from every attachment, so that it may become a place
for apprehending unseen inspirations and a treasury of the mysteries of
divine knowledge. This is why He says: ``The traveler on the snow-white
path and at the crimson pillar will never reach his homeland unless his
hands are empty of what is in the hands of the people.'' This is the
condition for the traveler. Reflect and consider carefully, so that you
may understand the intent of the Book without a veil.

\subsection{79}\label{section-269}

\begin{otherlanguage}{persian}

باری از مطلب دور ماندیم اگرچه همه ذکر مطلب است ولیکن قسم بخدا آنچه
میخواهم اختصار نمایم و باقلّ کفایت کنم می‌بینم زمام قلم از دست رفته و با
وجود این چه‌ قدر از لآلی بیشمار که ناسفته در صدف قلب مانده و چه مقدار
حوریّات معانی که در غرفه‌های حکمت مستور گشته که احدی مسّ آنها ننموده لم
یطمثهنّ انس قبلهم و لا جانّ و با همۀ این بیانات گویا حرفی از مقصود ذکر نشد
و رمزی از مطلوب مذکور نیامد تا کی محرمی یافت شود و احرام حرم دوست بندد و
بکعبهٔ مقصود واصل گردد و بی ‌گوش و لسان اسرار بیان بشنود و بیابد

\end{otherlanguage}

Bārī az maṭlab dūr māndīm, agarchih hamih dhikr-i maṭlab ast. Valīkin
qasam bi-Khudā ānchih mīkhvāham ikhtiṣār namāyam va bi-aqall kifāyat
kunam, mībīnam zimām-i qalam az dast raftih. Va bā vujūd-i īn chih qadr
az laʾālī-yi bīshumār kih nāsuftih dar ṣadaf-i qalb māndih va chih
miqdār ḥūrīyyāt-i maʿānī kih dar ghurfihhā-yi ḥikmat mastūr gashtih kih
aḥadī mass-i ānhā nanamūdih: ``Lam yaṭmithhunna insun qablahum wa lā
jānn.'' Va bā hamih-yi īn bayānāt gūyā ḥarfī az maqṣūd dhikr nashud va
ramzī az maṭlūb madhkūr nayāmad. Tā kay maḥramī yāft shavad va iḥrām-i
ḥaram-i Dūst bandad va bi-Kaʿbih-yi maqṣūd vāṣil gardad va bī-gūsh va
lisān asrār-i bayān bishnavad va biyābad.

We have, however, strayed from the subject, although all of this
concerns the subject. Yet, I swear by God, whenever I wish to be brief
and content Myself with the least, I find that the reins of the pen have
slipped from My hand. Even so, how many countless pearls remain
unpierced within the shell of the heart, and how many maidens of inner
meaning remain concealed in the chambers of wisdom, untouched by anyone:
``Neither human nor jinn has touched them before them.'' Despite all
these utterances, it is as though not a word of the intended meaning has
been spoken, nor a hint of the object sought given. When will a
confidant be found who will put on the pilgrim's garb for the sanctuary
of the Friend, reach the Kaaba of desire, and, without ear or tongue,
hear and apprehend the mysteries of utterance?

\subsection{80}\label{section-270}

\begin{otherlanguage}{persian}

پس از این بیانات محکمهٔ واضحهٔ لائحه مقصود از سماء در آیۀ منزله معلوم شد و
مفهوم گشت و اینکه میفرماید با ابر و غمام نازل میشود مقصود از ابر آن
اموریست که مخالف نفس و هوای ناس است چنانچه ذکر شد در آیهٔ مذکوره أ فکلّما
جآءکم رسول بما لا تهوی انفسکم استکبرتم ففریقاً کذّبتم و فریقاً تقتلون مثلاً
از قبیل تغییر احکام و تبدیل شرایع و ارتفاع قواعد و رسوم عادیّه و تقدّم
مؤمنین از عوام بر معرضین از علما و همچنین ظهور آن جمال ازلی بر حدودات
بشریّه از اکل و شرب و فقر و غنا و عزّت و ذلّت و نوم و یقظه و امثال آن از آن
چیزهائی که مردم را بشبهه میاندازد و منع مینماید همۀ این حجبات بغمام
تعبیر شده

\end{otherlanguage}

Pas az īn bayānāt-i muḥkamih-yi vāḍiḥih-yi lāʾiḥih maqṣūd az samāʾ dar
āyih-yi munazzalih maʿlūm shud va mafhūm gasht. Va īnkih mīfarmāyad bā
abr va ghamām nāzil mīshavad, maqṣūd az abr ān umūrist kih mukhālif-i
nafs va havā-yi nās ast, chunānchih dhikr shud dar āyih-yi madhkūrih:
``A-fa-kullamā jāʾakum rasūlun bi-mā lā tahwā anfusukumu istakbartum
fa-farīqan kadhdhabtum wa farīqan taqtulūn.'' Ma s̱alan az qabīl-i
taghyīr-i aḥkām va tabdīl-i sharāyiʿ va irtifāʿ-i qavāʿid va rusūm-i
ʿādīyyih va taqaddum-i muʾminīn az ʿavāmm bar muʿriḍīn az ʿulamā. Va
hamchunīn ẓuhūr-i ān jamāl-i azalī bar ḥudūdāt-i basharīyyih az akl va
shurb va faqr va ghinā va ʿizzat va dhillat va nawm va yaqẓih va am
s̱āl-i ān az ān chīzhāʾī kih mardum rā bi-shubhih mīandāzad va manʿ
mīnamāyad. Hamih-yi īn ḥujubāt bi-ghamām taʿbīr shudih.

These firm, clear, and evident explanations have thus made known and
clarified what is meant by heaven in the revealed verse. As for His
statement that He descends with clouds, by clouds are meant those things
that run counter to people's inclinations and desires, as mentioned in
the verse already cited: ``Whenever a messenger came to you with what
your souls did not desire, did you grow arrogant, calling some liars and
killing others?'' These include, for example, the changing of laws, the
replacement of religious codes, the abolition of established rules and
customary practices, and the precedence of believers among the common
people over those among the learned who turn away. Likewise, they
include the appearance of that eternal Beauty within human limitations,
such as eating and drinking, poverty and wealth, honor and abasement,
sleep and wakefulness, and similar things that cause people to doubt and
hold them back. All these veils are described as clouds.

\subsection{81}\label{section-271}

\begin{otherlanguage}{persian}

و اینست آن غمامی که سموات علم و عرفان کلّ من فی الأرض بآن میشکافد و شقّ
میگردد چنانچه میفرماید یوم تشقّق السّمآء بالغمام و همچنان که غمام ابصار
ناس را منع مینماید از مشاهدهٔ شمس ظاهری همین قسم هم این شئونات مذکوره
مردم را منع مینماید از ادراک آن شمس حقیقی چنانچه مذکور است در کتاب از
لسان کفّار و قالوا ما لهذا الرّسول یأکل الطّعام و یمشی فی الأسواق لو لا
انزل الیه ملک فیکون معه نذیراً مثل اینکه ملاحظه میشد از انبیا فقر ظاهری و
ابتلای ظاهری و همچنین ملزومات عنصری جسدی از قبیل جوع و امراض و حوادث
امکانیّه چون این مراتب از آن هیاکل قدسیّه ظاهر میشد مردم در صحراهای شک و
ریب و بیابانهای وهم و تحیّر متحیّر میماندند که چگونه میشود نفسی از جانب
خدا بیاید و اظهار غلبه نماید بر کلّ من علی الأرض و علّت خلق موجودات را
بخود نسبت دهد چنانچه فرموده لولاک لما خلقت الافلاک و مع‌ذلک باین قسمها
مبتلا بامور جزئیّه شود چنانچه شنیده‌اند از ابتلای هر نبی و اصحاب او از فقر
و امراض و ذلّت چنانچه سرهای اصحاب ایشان را در شهرها بهدیه میفرستادند و
ایشان را منع مینمودند از آنچه بآن مأمور بودند و هر کدام در دست اعدای دین
مبتلا بودند بقسمی که بر ایشان وارد میاوردند آنچه اراده مینمودند

\end{otherlanguage}

Va īnast ān ghamāmī kih samāvāt-i ʿilm va ʿirfān-i kull-i man fī al-arḍ
bi-ān mīshikāfad va shaqq mīgardad, chunānchih mīfarmāyad: ``Yawma
tashaqqaqu al-samāʾu bi-al-ghamām.'' Va hamchunān kih ghamām abṣār-i nās
rā manʿ mīnamāyad az mushāhidih-yi shams-i ẓāhirī, hamīn qism ham īn
shuʾūnāt-i madhkūrih mardum rā manʿ mīnamāyad az idrāk-i ān shams-i
ḥaqīqī, chunānchih madhkūr ast dar kitāb az lisān-i kuffār: ``Wa qālū mā
li-hādhā al-rasūli yaʾkulu al-ṭaʿāma wa yamshī fī al-aswāqi law lā
unzila ilayhi malakun fa-yakūna maʿahu nadhīrā.'' Mi s̱l-i īnkih
mulāḥiẓih mīshud az anbiyā faqr-i ẓāhirī va ibtilā-yi ẓāhirī va
hamchunīn malzūmāt-i ʿunṣurī-yi jasadī az qabīl-i jūʿ va amrāḍ va ḥavādi
s̱-i imkānīyyih. Chūn īn marātib az ān hayākil-i qudsīyyih ẓāhir mīshud,
mardum dar ṣaḥrāhā-yi shakk va rayb va biyābānhā-yi vahm va taḥayyur
mutaḥayyir mīmāndand kih chigūnih mīshavad nafsī az jānib-i Khudā
biyāyad va iẓhār-i ghalabih namāyad bar kull-i man ʿalā al-arḍ va
ʿillat-i khalq-i mawjūdāt rā bi-khud nisbat dahad, chunānchih farmūdih:
``Lawlāka la-mā khalaqtu al-aflāk,'' va maʿa dhālik bi-īn qismhā mubtalā
bi-umūr-i juzʾīyyih shavad. Chunānchih shunīdih-and az ibtilā-yi har
nabī va aṣḥāb-i ū az faqr va amrāḍ va dhillat, chunānchih sarhā-yi
aṣḥāb-i īshān rā dar shahrhā bi-hadyih mīfiristādand va īshān rā manʿ
mīnamūdand az ānchih bi-ān maʾmūr būdand. Va har kudām dar dast-i
aʿdā-yi dīn mubtalā būdand bi-qismī kih bar īshān vārid mīāvardand
ānchih irādih mīnamūdand.

This is the cloud through which the heavens of the knowledge and
understanding of all who dwell on earth are cleft and split apart, as He
says: ``The day when heaven is split apart by clouds.'' Just as clouds
prevent people's eyes from seeing the outward sun, so too do these
conditions prevent them from recognizing that true Sun, as recorded in
the Book in the words of the unbelievers: ``They said, `What is the
matter with this messenger, that he eats food and walks in the markets?
Why has an angel not been sent down to him, to be a warner alongside
him?'\,'' For example, outward poverty and outward affliction were
observed in the prophets, as were the necessities of elemental, bodily
existence, such as hunger, illnesses, and the vicissitudes of contingent
life. When these conditions appeared in those sacred Temples, people
remained bewildered in the deserts of doubt and uncertainty and the
wildernesses of delusion and perplexity. How, they wondered, could
someone come from God, proclaim his ascendancy over all who dwell on
earth, and declare himself the reason for the creation of all
beings---as He has said, ``Were it not for you, I would not have created
the heavens''---and yet be afflicted in these ways by such petty
matters? They have heard of the afflictions of every prophet and his
companions: poverty, illnesses, and abasement. The heads of their
companions were sent as gifts to the cities, and they were prevented
from carrying out what they had been commanded to do. Each suffered at
the hands of the enemies of religion, who inflicted upon them whatever
they wished.

\subsection{82}\label{section-272}

\begin{otherlanguage}{persian}

و این معلوم است که تغییرات و تبدیلات که در هر ظهور واقع میشود همان
غمامیست تیره که حایل میشود بصر عرفان عباد را از معرفت آن شمس الهی که از
مشرق هویّه اشراق فرموده زیرا که سالها عباد بر تقلید آباء و اجداد باقی
هستند و بآداب و طریقی که در آن شریعت مقرّر شده تربیت یافته‌اند یکمرتبه
بشنوند و یا ملاحظه نمایند شخصی که در میان ایشان بوده و در جمیع حدودات
بشریّه با ایشان یکسانست و مع‌ذلک جمیع آن حدودات شرعیّه که در قرنهای متواتره
بآن تربیت یافته‌اند و مخالف و منکر آن را کافر و فاسق و فاجر دانسته‌اند همه
را از میان بردارد البتّه این امور حجاب و غمام است از برای آنهائی که
قلوبشان از سلسبیل انقطاع نچشیده و از کوثر معرفت نیاشامیده و بمجرّد استماع
این امور چنان محتجب از ادراک آن شمس میمانند که دیگر بی ‌سؤال و جواب حکم
بر کفرش میکنند و فتوی بر قتلش میدهند چنانچه دیده‌اند و شنیده‌اند از قرون
اولی و این زمان نیز ملاحظه شد

\end{otherlanguage}

Va īn maʿlūm ast kih taghyīrāt va tabdīlāt kih dar har ẓuhūr vāqiʿ
mīshavad hamān ghamāmīst tīrih kih ḥāʾil mīshavad baṣar-i ʿirfān-i ʿibād
rā az maʿrifat-i ān shams-i ilāhī kih az mashriq-i huvīyyih ishrāq
farmūdih. Zīrā kih sālhā ʿibād bar taqlīd-i ābāʾ va ajdād bāqī hastand
va bi-ādāb va ṭarīqī kih dar ān sharīʿat muqarrar shudih tarbiyat
yāftih-and. Yikmartabih bishnavand va yā mulāḥiẓih namāyand shakhṣī kih
dar miyān-i īshān būdih va dar jamīʿ-i ḥudūdāt-i basharīyyih bā īshān
yiksānast va maʿa dhālik jamīʿ-i ān ḥudūdāt-i sharʿīyyih kih dar
qarnhā-yi mutavātirih bi-ān tarbiyat yāftih-and va mukhālif va munkir-i
ān rā kāfir va fāsiq va fājir dānistih-and hamih rā az miyān bardārad.
Al-battih īn umūr ḥijāb va ghamām ast az barā-yi ānhāʾī kih qulūbashān
az salsabīl-i inqiṭāʿ nachishīdih va az kaw s̱ar-i maʿrifat nayāshāmīdih.
Va bi-mujarrad-i istimāʿ-i īn umūr chunān muḥtajib az idrāk-i ān shams
mīmānand kih dīgar bī-suʾāl va javāb ḥukm bar kufrash mīkunand va fatvā
bar qatl ash mīdihand, chunānchih dīdih-and va shunīdih-and az qurūn-i
ūlā va īn zamān nīz mulāḥiẓih shud.

It is clear that the changes and alterations that occur in every
Revelation are that dark cloud which prevents the servants' eye of
understanding from recognizing the divine Sun shining from the dayspring
of divine Being. For years the servants have continued to imitate their
forefathers and ancestors and have been brought up in the observances
and ways prescribed in that religious law. Suddenly they hear or see
that a person who has lived among them, and who is like them in every
human limitation, nonetheless abolishes all those religious ordinances
in which they have been brought up through successive centuries, and
whose opponents and deniers they have regarded as unbelievers, sinners,
and evildoers. These things are certainly a veil and a cloud for those
whose hearts have not tasted the Salsabil of detachment or drunk from
the Kawthar of recognition. As soon as they hear of these things, they
become so veiled from recognizing that Sun that, without question or
discussion, they pronounce him an unbeliever and issue a decree for his
death, as people have seen and heard of past ages, and as has also been
witnessed in this time.

\subsection{83}\label{section-273}

\begin{otherlanguage}{persian}

پس باید جهدی نمود تا باعانت غیبی از این حجبات ظلمانی و غمام امتحانات
ربّانی از مشاهدهٔ آن جمال نورانی ممنوع نشویم و او را بنفس او بشناسیم و اگر
هم حجّت بخواهیم بیک حجّت و برهان اکتفا نمائیم تا بمنبع فیض نامتناهی که
جمیع فیوضات نزد او معدوم صرفست فائز گردیم نه آنکه هر روز بخیالی اعتراض
نمائیم و بهوائی تمسّک جوئیم

\end{otherlanguage}

Pas bāyad jahdī namūd tā bi-iʿānat-i ghaybī az īn ḥujubāt-i ẓulmānī va
ghamām-i imtiḥānāt-i rabbānī az mushāhidih-yi ān jamāl-i nūrānī mamnūʿ
nashavīm va ū rā bi-nafs-i ū bishināsīm. Va agar ham ḥujjat bikhvāhīm,
bi-yik ḥujjat va burhān iktifā namāʾīm tā bi-manbaʿ-i fayḍ-i nāmutanāhī
kih jamīʿ-i fuyūḍāt nazd-i ū maʿdūm-i ṣirfast fāʾiz gardīm, nah ānkih
har rūz bi-khayālī iʿtirāḍ namāʾīm va bi-havāʾī tamassuk jūʾīm.

We must therefore strive, with unseen assistance, not to let these dark
veils and clouds of divine tests prevent us from beholding that radiant
Beauty, and to recognize Him through His own Self. If we do seek
evidence, let us be content with a single proof and demonstration, so
that we may attain the source of infinite grace, beside which all other
gifts of grace are sheer nothingness, rather than raising an objection
each day on the basis of some fancy and clinging to some desire.

\subsection{84}\label{section-274}

\begin{otherlanguage}{persian}

سبحان‌اللّه با وجود اینکه از قبل این امورات را بتلویحات عجیبه و اشارات
غریبه خبر داده‌اند تا جمیع ناس اطّلاع یابند و در آن روز خود را از بحر
البحور فیوضات محروم نسازند مع‌ذلک امر چنین واقع میشود که مشهود است و این
مضمونات در قرآن هم نازل شده چنانچه میفرماید هل ینظرون الّا ان یأتیهم اللّه
فی ظلل من الغمام و علمای ظاهر بعضی این آیه را از علائم قیامت موهوم که
خود تعقّل نموده‌اند گرفته‌اند و مضمون آن اینست آیا انتظار میکشند مگر اینکه
بیاید آنها را خدا در سایه‌ئی از ابر و حال آنکه این مضمون در اکثر کتب
سماوی مذکور است و در همهٔ اماکن در ذکر علامات ظهور بعد ذکر فرموده‌اند
چنانچه از قبل ذکر شد

\end{otherlanguage}

Subḥān Allāh! Bā vujūd-e īnke az qabl īn umūrāt rā be talvīḥāt-e ʿajībe
va ishārāt-e gharībe khabar dāde-and tā jamīʿ-e nās iṭṭilāʿ yāband va
dar ān rūz khud rā az baḥr al-buḥūr-e fuyūḍāt maḥrūm nasāzand,
maʿa-dhālik amr chunīn vāqiʿ mī-shavad ke mashhūd ast. Va īn maḍmūnāt
dar Qurʾān ham nāzil shude, chunānche mī-farmāyad: hal yanẓurūna illā an
yaʾtiyahumu Allāhu fī ẓulalin mina al-ghamām. Va ʿulamā-ye ẓāhir baʿḍī
īn āye rā az ʿalāʾim-e qiyāmat-e mawhūm ke khud taʿaqqul namūde-and
gerefte-and. Va maḍmūn-e ān īnast: āyā intiẓār mī-kashand magar īnke
biyāyad ānhā rā Khudā dar sāye-ī az abr? Va ḥāl ānke īn maḍmūn dar
akthar-e kutub-e samāvī madhkūr ast va dar hame-ye amākin dar dhikr-e
ʿalāmāt-e ẓuhūr-e baʿd dhikr farmūde-and, chunānche az qabl dhikr shud.

Glory be to God! Although these matters were foretold through wondrous
allusions and extraordinary references, so that all people might become
aware and not deprive themselves on that day of the ocean of oceans of
grace, things nevertheless happen as is now witnessed. These themes have
also been revealed in the Qur'an, as He says: ``Are they waiting for
anything but that God should come to them in the shadows of the
clouds?'' Some of the outwardly learned have taken this verse to refer
to the signs of the imaginary resurrection they themselves have
conceived. Its meaning is this: ``Are they waiting for anything but that
God should come to them in the shadow of a cloud?'' Yet this theme is
mentioned in most of the heavenly Books, and wherever it appears, it is
mentioned among the signs of the subsequent Revelation, as was stated
before.

\subsection{85}\label{section-275}

\begin{otherlanguage}{persian}

و همچنین میفرماید یوم تأتی السّمآء بدخان مبین یغشی النّاس هذا عذاب الیم که
مضمون آن اینست روزی که میآید آسمان بدودی آشکار و فرو میگیرد مردم را و
اینست عذاب الیم و همین امورات را که مغایر انفس خبیثه و مخالف هوای ناس
است حضرت ربّ العزّة محک و میزان قرار داده و بآنها امتحان میفرماید عباد خود
را و تمیز میدهد سعید را از شقی و معرض را از مقبل چنانچه مذکور شد و
اختلافات و نسخ و هدم رسومات عادیّه و انعدام اعلام محدوده را بدخان در آیهٔ
مذکوره تعبیر فرموده و کدام دخان است اعظم از این دخان که فرو گرفته همهٔ
ناس را و عذابیست برای آنها که هر چه میخواهند رفع آن نمایند قادر نیستند و
بنار نفس در هر حین بعذابی جدید معذّبند زیرا که هر چه میشنوند که این امر
بدیع الهی و حکم منیع صمدانی در اطراف ارض ظاهر شده و هر روز در علوّ است
ناری جدید در قلوبشان مشتعل میشود و آنچه ملاحظه مینمایند از قدرت و انقطاع
و ثبوت این اصحاب که هر روز بعنایت الهی محکمتر و راسختر میشوند اضطراب
تازه در نفوسشان ظاهر میگردد در این ایّام که بحمد اللّه سطوت الهی چنان غلبه
فرموده که جرئت تکلّم ندارند و اگر یکی از اصحاب حقّ را که صد‌هزار جان بدل و
جان رایگان در ره دوست ایثار مینماید ملاقات نمایند از خوف اظهار ایمان
میکنند و چون خلوت میکنند بسب و لعن مشغول میشوند چنانچه میفرماید اذا
لقوکم قالوا آمنّا و اذا خلوا عضّوا علیکم الأنامل من الغیظ قل موتوا بغیظکم
انّ اللّه علیم بذات الصّدور

\end{otherlanguage}

Va hamchunīn mī-farmāyad: yawma taʾtī al-samāʾu bi-dukhānin mubīn,
yaghshā al-nāsa hādhā ʿadhābun alīm. Ke maḍmūn-e ān īnast: rūzī ke
mī-āyad āsmān be dūdī āshkār va furū mī-gīrad mardum rā va īnast
ʿadhāb-e alīm. Va hamīn umūrāt rā ke mughāyir-e anfus-e khabīthe va
mukhālif-e havā-ye nās ast Ḥaḍrat-e Rabb al-ʿIzzat maḥakk va mīzān qarār
dāde va be ānhā imtiḥān mī-farmāyad ʿibād-e khud rā va tamyīz mī-dahad
saʿīd rā az shaqī va muʿriḍ rā az muqbil, chunānche madhkūr shud. Va
ikhtilāfāt va naskh va hadm-e rusūmāt-e ʿādiyye va inʿidām-e aʿlām-e
maḥdūde rā be dukhān dar āye-ye madhkūre taʿbīr farmūde. Va kudām dukhān
ast aʿẓam az īn dukhān ke furū gerefte hame-ye nās rā va ʿadhābīst
barā-ye ānhā ke har che mī-khāhand rafʿ-e ān namāyand qādir nīstand va
be nār-e nafs dar har ḥīn be ʿadhābī jadīd muʿadhdhaband? Zīrā ke har
che mī-shenavand ke īn amr-e badīʿ-e ilāhī va ḥukm-e manīʿ-e ṣamadānī
dar aṭrāf-e arḍ ẓāhir shude va har rūz dar ʿuluvv ast, nārī jadīd dar
qulūbeshān mushtaʿil mī-shavad. Va ānche mulāḥiẓe mī-namāyand az qudrat
va inqiṭāʿ va thubūt-e īn aṣḥāb ke har rūz be ʿināyat-e ilāhī muḥkamtar
va rāsikhtar mī-shavand, iḍṭirāb-e tāze dar nufūseshān ẓāhir mī-gardad.
Dar īn ayyām ke bi-ḥamd Allāh saṭvat-e ilāhī chunān ghalabe farmūde ke
jurʾat-e takallum nadārand va agar yekī az aṣḥāb-e Ḥaqq rā ke ṣad-hazār
jān be del va jān rāygān dar rah-e Dūst īthār mī-namāyad mulāqāt
namāyand, az khawf iẓhār-e īmān mī-kunand va chun khalvat mī-kunand be
sabb va laʿn mashghūl mī-shavand. Chunānche mī-farmāyad: idhā laqūkum
qālū āmannā wa-idhā khalaw ʿaḍḍū ʿalaykumu al-anāmila mina al-ghayẓ, qul
mūtū bi-ghayẓikum inna Allāha ʿalīmun bi-dhāti al-ṣudūr.

He also says: ``The day when heaven brings forth a manifest smoke, which
envelops the people: this is a painful torment.'' Its meaning is this:
``The day when heaven comes with a visible smoke and envelops the
people, and this is a painful torment.'' The Lord of Glory has made
these very things, which are repugnant to impure souls and contrary to
people's desires, a touchstone and a balance. Through them He tests His
servants and distinguishes the blessed from the wretched, and those who
turn away from those who turn toward Him, as has been mentioned. In the
verse cited, He has referred to differences, the abrogation and
destruction of customary observances, and the disappearance of
established authorities as ``smoke.'' What smoke is greater than this
smoke that has enveloped all people and is a torment to them, which they
are powerless to dispel, however much they wish to do so? At every
moment they suffer a fresh torment in the fire of self. For whenever
they hear that this wondrous Cause of God and this exalted decree of the
Eternal have appeared throughout the earth and are gaining ascendancy
each day, a new fire is kindled in their hearts. And whenever they
observe the strength, detachment, and steadfastness of these companions,
who grow firmer and more resolute each day through divine grace, a fresh
agitation arises within them. In these days, praise be to God, divine
might has so prevailed that they do not dare to speak. If they encounter
one of the companions of the True One, who would freely and
wholeheartedly sacrifice a hundred thousand lives in the path of the
Beloved, they profess faith out of fear; but when they are alone, they
occupy themselves with abuse and curses. As He says: ``When they meet
you, they say, `We believe'; but when they are alone, they bite their
fingertips in rage at you. Say: `Die in your rage! God indeed knows what
lies within the breasts.'\,''

\subsection{86}\label{section-276}

\begin{otherlanguage}{persian}

و عنقریب است که اعلام قدرت الهی را در همۀ بلاد مرتفع بینی و آثار غلبه و
سلطنت او را در جمیع دیار ظاهر مشاهده فرمائی باری اکثر علما چون این آیات
را ادراک ننموده‌اند و از مقصود قیامت واقف نشده‌اند لهذا جمیع را بقیامت
موهوم من حیث لا یشعر تفسیر مینمایند خدای واحد شاهد است که اگر قدری بصیرت
باشد از تلویح همین دو آیه جمیع مطالب که مقصود است ادراک میشود و بصبح
منیر ایقان بعنایت رحمن واصل میگردند کذلک تغنّ علیک حمامة البقآء علی افنان
سدرة البهآء لعلّ تکوننّ فی مناهج العلم و الحکمة باذن اللّه سالکا

\end{otherlanguage}

Va ʿan-qarīb ast ke aʿlām-e qudrat-e ilāhī rā dar hame-ye bilād murtafiʿ
bīnī va āthār-e ghalabe va salṭanat-e ū rā dar jamīʿ-e diyār ẓāhir
mushāhade farmāʾī. Bārī akthar-e ʿulamā chun īn āyāt rā idrāk
nanamūde-and va az maqṣūd-e qiyāmat vāqif nashude-and, li-hādhā jamīʿ rā
be qiyāmat-e mawhūm min ḥaythu lā yashʿur tafsīr mī-namāyand. Khudā-ye
vāḥid shāhid ast ke agar qadrī baṣīrat bāshad az talvīḥ-e hamīn du āye
jamīʿ-e maṭālib ke maqṣūd ast idrāk mī-shavad va be ṣubḥ-e munīr-e īqān
be ʿināyat-e Raḥmān vāṣil mī-gardand. Kadhālika taghannu ʿalayka
ḥamāmatu al-baqāʾ ʿalā afnāni sidrati al-bahāʾ laʿallaka takūnanna fī
manāhiji al-ʿilmi wa-al-ḥikmati bi-idhni Allāhi sālikan.

Soon you will see the banners of divine power raised in every land and
behold the signs of His ascendancy and sovereignty manifest in every
region. In short, because most of the learned have not understood these
verses or grasped what is meant by resurrection, they unwittingly
interpret them all as referring to an imaginary resurrection. The one
God is witness that, with even a little insight, all the intended
meanings can be understood from the allusions in these two verses alone,
and through the grace of the All-Merciful one can attain the radiant
dawn of certitude. Thus the dove of eternity sings to you upon the
branches of the Tree of Glory, so that, by God's permission, you may
walk the paths of knowledge and wisdom.

\subsection{87}\label{section-277}

\begin{otherlanguage}{persian}

و قوله یرسل ملائکته الی آخر القول مقصود از این ملائکه آن نفوسی هستند که
بقوّهٔ روحانیّه صفات بشریّه را بنار محبّت الهی سوختند و بصفات عالین و کرّوبیّین
متّصف گشتند چنانچه حضرت صادق در وصف کرّوبیّین میفرماید قومی از شیعیان ما
هستند خلف عرش و از ذکر خلف العرش اگرچه معانی بسیار منظور بوده هم بر حسب
ظاهر و هم بر حسب باطن ولیکن در یک مقام مدلّست بر عدم وجود شیعه چنانچه در
مقام دیگر میفرماید مؤمن مثل کبریت احمر است و بعد بمستمع میفرماید آیا
کبریت احمر دیده‌ئی ملتفت شوید باین تلویح که ابلغ از تصریح است دلالت میکند
بر عدم وجود مؤمن این قول آن حضرت و حال مشاهده کن چه‌ قدر از این خلق
بی‌انصاف که رائحۀ ایمان نشنیده‌اند مع‌ذلک کسانی را که بقول ایشان ایمان محقّق
میشود نسبت بکفر میدهند

\end{otherlanguage}

Va qawluhu: yursilu malāʾikatahu, ilā ākhir al-qawl. Maqṣūd az īn
malāʾike ān nufūsī hastand ke be quvve-ye rūḥāniyye ṣifāt-e bashariyye
rā be nār-e maḥabbat-e ilāhī sūkhtand va be ṣifāt-e ʿāllīn va karūbiyyīn
muttaṣif gashtand. Chunānche Ḥaḍrat-e Ṣādiq dar vaṣf-e karūbiyyīn
mī-farmāyad: qawmī az shīʿiyān-e mā hastand khalf-e ʿarsh. Va az dhikr-e
khalf al-ʿarsh agarchi maʿānī-ye bisyār manẓūr būde, ham bar ḥasab-e
ẓāhir va ham bar ḥasab-e bāṭin, valīkin dar yek maqām mudill ast bar
ʿadam-e vujūd-e shīʿe. Chunānche dar maqām-e dīgar mī-farmāyad: muʾmin
mithl-e kibrīt-e aḥmar ast. Va baʿd be mustamiʿ mī-farmāyad: āyā
kibrīt-e aḥmar dīde-ī? Multafit shavīd be īn talvīḥ ke ablagh az taṣrīḥ
ast. Dalālat mī-kunad bar ʿadam-e vujūd-e muʾmin īn qawl-e ān Ḥaḍrat. Va
ḥāl mushāhade kun che qadr az īn khalq-e bī-inṣāf ke rāʾiḥe-ye īmān
nashenīde-and, maʿa-dhālik kasānī rā ke be qawl-e īshān īmān muḥaqqaq
mī-shavad nisbat be kufr mī-dahand.

As for His words, ``He will send His angels,'' to the end of the
passage: these angels mean those souls who, through spiritual power,
have burned away human attributes in the fire of divine love and have
become endowed with the attributes of the exalted ones and the cherubim.
Thus Ṣādiq says in describing the cherubim: ``They are a company of our
followers behind the Throne.'' Although many meanings, both outward and
inward, were intended by the expression ``behind the Throne,'' in one
sense it indicates that no true follower exists. As He says elsewhere:
``A believer is like red sulphur.'' He then asks His listener: ``Have
you ever seen red sulphur?'' Attend to this allusion, which is more
eloquent than an explicit statement. This saying of His indicates that
no believer exists. And now observe how many of these unjust people, who
have never caught the fragrance of faith, nevertheless accuse of
unbelief those through whose word faith itself is established.

\subsection{88}\label{section-278}

\begin{otherlanguage}{persian}

باری چون این وجودات قدسیّه از عوارض بشریّه پاک و مقدّس گشتند و متخلّق باخلاق
روحانیّین و متّصف باوصاف مقدّسین شدند لهذا اسم ملائکه بر آن نفوس مقدّسه
اطلاق گشته باری اینست معنی این کلمات که هر فقرهٔ آن بآیات واضحه و دلیلهای
متقنه و براهین لائحه اظهار شد

\end{otherlanguage}

Bārī chun īn vujūdāt-e qudsiyye az ʿavāriḍ-e bashariyye pāk va muqaddas
gashtand va mutakhalliq be akhlāq-e rūḥāniyyīn va muttaṣif be awṣāf-e
muqaddasīn shudand, li-hādhā ism-e malāʾike bar ān nufūs-e muqaddase
iṭlāq gashte. Bārī īnast maʿnī-ye īn kalimāt ke har faqre-ye ān be
āyāt-e vāḍiḥe va dalīlhā-ye mutqane va barāhīn-e lāʾiḥe iẓhār shud.

In short, because these holy beings have become pure and sanctified from
human limitations, have acquired the character of spiritual beings, and
have become endowed with the attributes of the sanctified, the name
``angels'' has been applied to those holy souls. Such is the meaning of
these words, every phrase of which has been explained through clear
verses, firm evidence, and manifest proofs.

\subsection{89}\label{section-279}

\begin{otherlanguage}{persian}

و چون امم عیسی باین معانی نرسیدند و این علامات بر حسب ظاهر چنانچه خود و
علمای ایشان ادراک نموده‌اند ظاهر نشد لهذا بمظاهر قدسیّه از آن یوم تا بحال
اقبال ننمودند و از جمیع فیوضات قدسیّه محروم شدند و از بدایع کلمات صمدانیّه
محجوب گشتند اینست شأن این عباد در یوم معاد و اینقدر ادراک ننمودند که اگر
در هر عصری علائم ظهور مطابق آنچه در اخبار است در عالم ظاهر ظاهر شود دیگر
که را یارای انکار و اعراض میماند و چگونه میان سعید و شقی و مجرم و متّقی
تفصیل میشود مثلاً انصاف دهید اگر این عبارات که در انجیل مسطور است بر حسب
ظاهر ظاهر شود و ملائکه با عیسی بن مریم از سماء ظاهره با ابری نازل شوند
دیگر که یارای تکذیب دارد و یا که لایق انکار و قابل استکبار باشد بلکه
فی‌الفور همۀ اهل ارض را اضطراب بقسمی احاطه میکند که قادر بر حرف و تکلّم
نیستند تا چه رسد برد و قبول و نظر بعدم ادراک این معانی بود که جمعی از
علمای نصاری بآن حضرت معارضه نمودند که اگر تو آن نبیّ موعودی چرا با تو
نیستند آن ملائکه که در کتب ما مسطور است که باید با جمال موعود بیایند تا
او را اعانت نمایند در امر او و منذر باشند برای عباد چنانچه ربّ العزّة از
لسان ایشان خبر داده لو لا انزل الیه ملک فیکون معه نذیرا که مضمون آن این
است چرا فرو فرستاده نشد با محمّد ملکی پس باشد با او بیم‌دهنده و ترساننده
مردمان را

\end{otherlanguage}

Va chun umam-e ʿĪsā be īn maʿānī naresīdand va īn ʿalāmāt bar ḥasab-e
ẓāhir chunānche khud va ʿulamā-ye īshān idrāk namūde-and ẓāhir nashud,
li-hādhā be maẓāhir-e qudsiyye az ān yawm tā be ḥāl iqbāl nanamūdand va
az jamīʿ-e fuyūḍāt-e qudsiyye maḥrūm shudand va az badāyiʿ-e kalimāt-e
ṣamadāniyye maḥjūb gashtand. Īnast shaʾn-e īn ʿibād dar yawm-e maʿād. Va
īn-qadr idrāk nanamūdand ke agar dar har ʿaṣrī ʿalāʾim-e ẓuhūr muṭābiq-e
ānche dar akhbār ast dar ʿālam-e ẓāhir ẓāhir shavad, dīgar ke rā yārā-ye
inkār va iʿrāḍ mī-mānad va chegūne miyān-e saʿīd va shaqī va mujrim va
muttaqī tafṣīl mī-shavad? Mathalan inṣāf dahīd: agar īn ʿibārāt ke dar
Injīl masṭūr ast bar ḥasab-e ẓāhir ẓāhir shavad va malāʾike bā ʿĪsā
ibn-e Maryam az samāʾ-e ẓāhere bā abrī nāzil shavand, dīgar ke yārā-ye
takdhīb dārad va yā ke lāyiq-e inkār va qābil-e istikbār bāshad? Balke
fī al-fawr hame-ye ahl-e arḍ rā iḍṭirāb be qismī iḥāṭe mī-kunad ke qādir
bar ḥarf va takallum nīstand, tā che resad be radd va qabūl. Va naẓar be
ʿadam-e idrāk-e īn maʿānī būd ke jamʿī az ʿulamā-ye Naṣārā be ān Ḥaḍrat
muʿāraḍe namūdand ke: agar tu ān nabī-ye mawʿūdī, cherā bā tu nīstand ān
malāʾike ke dar kutub-e mā masṭūr ast ke bāyad bā Jamāl-e Mawʿūd
biyāyand tā ū rā iʿānat namāyand dar amr-e ū va mundhir bāshand barā-ye
ʿibād? Chunānche Rabb al-ʿIzzat az lisān-e īshān khabar dāde: law lā
unzila ilayhi malakun fa-yakūna maʿahu nadhīran. Ke maḍmūn-e ān īn ast:
cherā furū ferestāde nashud bā Muḥammad malakī, pas bāshad bā ū
bīm-dahande va tarsānande mardumān rā?

Because the followers of Jesus did not grasp these meanings, and these
signs did not outwardly appear as they and their learned men had
understood them, they have not turned toward the holy Manifestations
from that day to the present. They have been deprived of all holy
outpourings of grace and veiled from the wonders of the words of the
Eternal. Such is the condition of these servants on the Day of Return.
They have failed to understand even this much: if, in every age, the
signs of the Revelation were to appear in the visible world exactly as
described in the traditions, who would retain the power to deny or turn
away, and how would the blessed be distinguished from the wretched, or
the sinner from the God-fearing? Consider fairly, for example: if these
statements recorded in the Gospel were fulfilled outwardly, and the
angels descended with Jesus, son of Mary, from the visible heaven upon a
cloud, who would have the power to deny Him, or be capable of rejection
and arrogance? Rather, agitation would immediately seize all the
inhabitants of the earth to such a degree that they would be unable to
utter a word, much less reject or accept. It was because they did not
understand these meanings that a number of Christian divines challenged
Him, saying: ``If you are that promised Prophet, why are those angels
not with you who, according to our Books, must come with the promised
Beauty to assist Him in His Cause and warn His servants?'' Thus the Lord
of Glory has reported their words: ``Why has an angel not been sent down
to Him, to be a warner with Him?'' Its meaning is this: ``Why has an
angel not been sent down with Muḥammad to warn and admonish the people
alongside Him?''

\subsection{90}\label{section-280}

\begin{otherlanguage}{persian}

اینست که در همۀ اعهاد و اعصار اینگونه اعتراضات و اختلافات در میان مردم
بوده و همیشۀ ایّام مشغول بزخارف قول میشدند که فلان علامت ظاهر نشد و فلان
برهان باهر نیامد و این مرضها عارض نمیشد مگر آنکه تمسّک بعلمای عصر میجستند
در تصدیق و تکذیب این جواهر مجرّده و هیاکل الهیّه و ایشان هم نظر باستغراق
در شئونات نفسیّه و اشتغال بامورات دنیّۀ فانیه این شموس باقیه را مخالف علم
و ادراک و معارض جهد و اجتهاد خود میدیدند و معانی کلمات الهیّه و احادیث و
اخبار حروفات احدیّه را هم بر سبیل ظاهر بادراک خود معنی و بیان مینمودند
لهذا خود و جمیع ناس را از نیسان فضل و رحمت ایزدی مأیوس و مهجور نمودند با
اینکه خود مذعن و مقرّند بحدیث مشهور که میفرماید حدیثنا صعب مستصعب و در
جای دیگر میفرماید انّ امرنا صعب مستصعب لا یحتمله الّا ملک مقرّب او نبیّ مرسل
او عبد امتحن اللّه قلبه للایمان و مسلّم است نزد خود ایشان که هیچیک از این
ثلاثه در حقّ ایشان صادق نیست دو قسم اوّل که واضحست و امّا ثالث هرگز از
امتحانات الهی سالم نماندند و در ظهور محکّ الهی جز غش چیزی از ایشان بظهور
نرسید

\end{otherlanguage}

Īnast ke dar hame-ye aʿhād va aʿṣār īn-gūne iʿtirāḍāt va ikhtilāfāt dar
miyān-e mardum būde va hamīshe-ye ayyām mashghūl be zakhārif-e qawl
mī-shudand ke fulān ʿalāmat ẓāhir nashud va fulān burhān bāhir nayāmad.
Va īn maraḍhā ʿāriḍ namī-shud magar ānke tamassuk be ʿulamā-ye ʿaṣr
mī-justand dar taṣdīq va takdhīb-e īn javāhir-e mujarrade va hayākil-e
ilāhiyye. Va īshān ham naẓar be istighrāq dar shuʾūnāt-e nafsiyye va
ishtighāl be umūrāt-e daniyye-ye fāniye īn shumūs-e bāqiye rā mukhālif-e
ʿilm va idrāk va muʿāriḍ-e jahd va ijtihād-e khud mī-dīdand. Va
maʿānī-ye kalimāt-e ilāhiyye va aḥādīth va akhbār-e ḥurūfāt-e aḥadiyye
rā ham bar sabīl-e ẓāhir be idrāk-e khud maʿnī va bayān mī-namūdand.
Li-hādhā khud va jamīʿ-e nās rā az naysān-e faḍl va raḥmat-e Īzadī
maʾyūs va mahjūr namūdand. Bā īnke khud mudhʿin va muqirrand be ḥadīth-e
mashhūr ke mī-farmāyad: ḥadīthunā ṣaʿbun mustaṣʿab. Va dar jā-ye dīgar
mī-farmāyad: inna amranā ṣaʿbun mustaṣʿabun lā yaḥtamiluhu illā malakun
muqarrabun aw nabiyyun mursalun aw ʿabdun imtaḥana Allāhu qalbahu
li-al-īmān. Va musallam ast nazd-e khud-e īshān ke hīch-yek az īn
thalāthe dar ḥaqq-e īshān ṣādiq nīst. Du qism-e avval ke vāḍiḥ ast va
ammā thālith hargiz az imtiḥānāt-e ilāhī sālim namāndand va dar ẓuhūr-e
maḥakk-e ilāhī juz ghishsh chīzī az īshān be ẓuhūr naresīd.

Thus, in every age and era, objections and disputes of this kind have
arisen among the people, and they have continually occupied themselves
with specious talk, saying that this sign did not appear and that proof
did not shine forth. These maladies arose only because they relied on
the learned of their age to affirm or deny these detached Essences and
divine Temples. Those learned men, immersed in selfish concerns and
occupied with base and fleeting affairs, regarded these everlasting Suns
as contrary to their own knowledge and understanding, and opposed to
their efforts and learned judgments. They also interpreted and explained
the meanings of the divine words and the traditions and accounts of the
Letters of Unity according to their own outward understanding. They
thereby deprived themselves and all people of hope in the spring rains
of divine grace and mercy and kept them far removed from them. Yet they
themselves acknowledge and affirm the well-known tradition that says:
``Our utterance is difficult, exceedingly difficult.'' Elsewhere He
says: ``Our Cause is indeed difficult, exceedingly difficult; none can
bear it except an angel brought near, a Prophet sent by God, or a
servant whose heart God has tested for faith.'' They themselves concede
that none of these three descriptions applies to them. The first two
cases are clear; as for the third, they have never emerged unscathed
from divine tests, and when the divine touchstone appeared, nothing but
impurity was revealed in them.

\subsection{91}\label{section-281}

\begin{otherlanguage}{persian}

سبحان‌اللّه با وجود اقرار باین حدیث علمائی که در مسائل شرعیّه هنوز در ظنّ و
شکّند چگونه در غوامض مسائل اصول الهیّه و جواهر اسرار کلمات قدسیّه اظهار علم
مینمایند و میگویند فلان حدیث که از علائم ظهور قائم است هنوز ظاهر نشده با
اینکه رائحهٔ معانی احادیث را ابداً ادراک ننموده‌اند و غافل از اینکه جمیع
علامات ظاهر شد و صراط امر کشیده گشت و المؤمنون کالبرق علیه یمرّون و هم
لظهور العلامة ینتظرون قل یا ملأ الجهّال فانتظروا کما کان الّذین من قبلکم
لمن المنتظرین

\end{otherlanguage}

Subḥān Allāh! Bā vujūd-e iqrār be īn ḥadīth, ʿulamāʾī ke dar masāʾil-e
sharʿiyye hanūz dar ẓann va shakkand chegūne dar ghavāmiḍ-e masāʾil-e
uṣūl-e ilāhiyye va javāhir-e asrār-e kalimāt-e qudsiyye iẓhār-e ʿilm
mī-namāyand va mī-gūyand fulān ḥadīth ke az ʿalāʾim-e ẓuhūr-e Qāʾim ast
hanūz ẓāhir nashude, bā īnke rāʾiḥe-ye maʿānī-ye aḥādīth rā abadan idrāk
nanamūde-and va ghāfil az īnke jamīʿ-e ʿalāmāt ẓāhir shud va ṣirāṭ-e amr
kashīde gasht. Wa-al-muʾminūna ka-al-barqi ʿalayhi yamurrūna wa-hum
li-ẓuhūri al-ʿalāmati yantaẓirūn. Qul yā malaʾa al-juhhāli fa-intaẓirū
kamā kāna alladhīna min qablikum la-mina al-muntaẓirīn.

Glory be to God! Despite acknowledging this tradition, how can learned
men who remain uncertain and doubtful even in questions of religious law
claim knowledge of the intricacies of divine principles and the
innermost mysteries of the holy words? They say that a certain tradition
concerning the signs of the Qāʾim's appearance has not yet been
fulfilled, although they have never caught even the fragrance of the
meanings of the traditions. They are unaware that all the signs have
appeared and the bridge of the Cause has been laid. The believers pass
over it like lightning, while they wait for the sign to appear. Say:
``You company of the ignorant! Wait, then, as those before you also
waited.''

\subsection{92}\label{section-282}

\begin{otherlanguage}{persian}

و اگر از ایشان سؤال شود از شرایط ظهور انبیای بعد که در کتب قبل است از
جملهٔ آنها علامات ظهور و اشراق شمس محمّدیست چنانچه مذکور شد و بر حسب ظاهر
هیچیک ظاهر نشد مع‌ذلک بچه دلیل و برهان نصاری و امثال آنها را رد مینمائید
و حکم بر کفر آنها نموده‌اید چون عاجز از جواب میشوند تمسّک باین نمایند که
این کتب تحریف شده و من عند اللّه نبوده و نیست و حال آنکه خود عبارات آیه
شهادت میدهد بر اینکه من عند اللّه است و مضمون همین آیه در قرآن هم موجود
است لو انتم تعرفون براستی میگویم مقصود از تحریف را در این مدّت ادراک
ننموده‌اند

\end{otherlanguage}

Va agar az īshān suʾāl shavad az sharāʾiṭ-e ẓuhūr-e anbiyā-ye baʿd ke
dar kutub-e qabl ast, az jumle-ye ānhā ʿalāmāt-e ẓuhūr va ishrāq-e
shams-e Muḥammadīst chunānche madhkūr shud, va bar ḥasab-e ẓāhir
hīch-yek ẓāhir nashud, maʿa-dhālik be che dalīl va burhān Naṣārā va
amthāl-e ānhā rā radd mī-namāyīd va ḥukm bar kufr-e ānhā namūde-īd? Chun
ʿājiz az javāb mī-shavand tamassuk be īn namāyand ke īn kutub taḥrīf
shude va min ʿind Allāh nabūde va nīst. Va ḥāl ānke khud-e ʿibārāt-e āye
shahādat mī-dahad bar īnke min ʿind Allāh ast va maḍmūn-e hamīn āye dar
Qurʾān ham mawjūd ast, law antum taʿrifūn. Be rāstī mī-gūyam maqṣūd az
taḥrīf rā dar īn muddat idrāk nanamūde-and.

If they are questioned about the conditions for the appearance of later
Prophets that are recorded in the earlier Books---among them the signs
of the appearance and dawning of the Muḥammadan Sun, as mentioned
before---and are asked, ``Since none of these was outwardly fulfilled,
by what evidence and proof do you nevertheless reject the Christians and
others like them and pronounce them unbelievers?'' they are unable to
answer. They then resort to the claim that these Books have been
corrupted and were not, and are not, from God. Yet the very wording of
the verse testifies that it is from God, and the substance of this same
verse is also found in the Qur'an, if you but knew. Truly I say:
throughout this time they have failed to understand what is meant by
corruption.

\subsection{93}\label{section-283}

\begin{otherlanguage}{persian}

بلی در آیات منزله و کلمات مرایای احمدیّه ذکر تحریف عالین و تبدیل مستکبرین
هست ولیکن در مواضع مخصوصه ذکر شده و از آن جمله حکایت ابن صوریاست در وقتی
که اهل خیبر در حکم قصاص زنای محصن و محصنه از نقطۀ فرقان سؤال نمودند و آن
حضرت فرمود حکم خدا رجم است و ایشان انکار نمودند که در تورات چنین حکمی
نیست حضرت فرمود از علمای خود که را مسلّم و کلام او را مصدّقید ابن صوریا را
قبول نمودند و حضرت او را احضار نمود و فرمود اقسمک باللّه الّذی فلق لکم
البحر و انزل علیکم المنّ و ظلّل لکم الغمام و نجّاکم من فرعون و ملإه و فضّلکم
علی النّاس بأن تذکر لنا ما حکم به موسی فی قصاص الزّانی المحصن و الزّانیة
المحصنة که مضمون آن این است که آن حضرت ابن صوریا را باین قسمهای مؤکّد قسم
دادند که در تورات حکم قصاص در زنای محصن چه نازل شده عرض نمود یا محمّد رجم
است آن حضرت فرمود پس چرا این حکم میان یهود منسوخ شده و مجری نیست عرض
نمود چون بختنصّر بیت المقدّس را بسوخت و جمیع یهود را بقتل رساند دیگر یهودی
در ارض باقی نماند الّا معدودی قلیل و علمای آن عصر نظر بقلّت یهود و کثرت
عمالقه بمشاوره جمع شدند که اگر موافق حکم تورات عمل شود آنچه از دست
بختنصّر نجات یافتند بحکم کتاب مقتول میشوند و باین مصالح حکم قتل را از
میان بالمرّه برداشتند باری در این بین جبرئیل بر قلب منیرش نازل شد و این
آیه عرض نمود یحرّفون الکلم عن مواضعه

\end{otherlanguage}

Balī, dar āyāt-i munazzala va kalimāt-i marāyā-yi Aḥmadiyya dhikr-i
taḥrīf-i ʿālīn va tabdīl-i mustakbirīn hast, valīkin dar mavāḍiʿ-i
makhṣūṣa dhikr shuda, va az ān jumla ḥikāyat-i Ibn-i Ṣūriyā-st. Dar
vaqtī ki ahl-i Khaybar dar ḥukm-i qiṣāṣ-i zinā-yi muḥṣan va muḥṣana az
Nuqṭa-yi Furqān suʾāl namūdand va ān Ḥaḍrat farmūd ḥukm-i Khudā rajm
ast, va īshān inkār namūdand ki dar Tawrāt chunīn ḥukmī nīst. Ḥaḍrat
farmūd az ʿulamā-yi khud ki rā musallam va kalām-i ū rā muṣaddiqīd?
Ibn-i Ṣūriyā rā qabūl namūdand va Ḥaḍrat ū rā iḥḍār namūd va farmūd:
Uqsimuka bi-llāhi alladhī falaqa lakumu al-baḥra wa-anzala ʿalaykumu
al-manna wa-ẓallala lakumu al-ghamāma wa-najjākum min Firʿawna
wa-malaʾihi wa-faḍḍalakum ʿalā al-nāsi bi-an tadhkura lanā mā ḥakama
bihi Mūsā fī qiṣāṣi al-zānī al-muḥṣani wa-al-zāniyati al-muḥṣanati. Ki
maḍmūn-i ān īn ast ki ān Ḥaḍrat Ibn-i Ṣūriyā rā bi-īn qasam-hā-yi
muʾakkad qasam dādand ki dar Tawrāt ḥukm-i qiṣāṣ dar zinā-yi muḥṣan chi
nāzil shuda. ʿArḍ namūd: Yā Muḥammad, rajm ast. Ān Ḥaḍrat farmūd: Pas
chirā īn ḥukm miyān-i Yahūd mansūkh shuda va mujrā nīst? ʿArḍ namūd:
Chūn Bukhtnaṣṣar Bayt al-Muqaddas rā bisūkht va jamīʿ-i Yahūd rā bi-qatl
rasānd, dīgar Yahūdī dar arḍ bāqī namānd illā maʿdūdī qalīl, va
ʿulamā-yi ān ʿaṣr naẓar bi-qillat-i Yahūd va kathrat-i ʿAmāliqa
bi-mushāvara jamʿ shudand ki agar muvāfiq-i ḥukm-i Tawrāt ʿamal shavad,
ānchi az dast-i Bukhtnaṣṣar najāt yāftand bi-ḥukm-i Kitāb maqtūl
mīshavand, va bi-īn maṣāliḥ ḥukm-i qatl rā az miyān bi-al-marra
bardāshtand. Bārī, dar īn bayn Jibraʾīl bar qalb-i munīrash nāzil shud
va īn āya ʿarḍ namūd: Yuḥarrifūna al-kalima ʿan mawāḍiʿihi.

Yes, the revealed verses and the words of the mirrors of Aḥmad mention
the distortion wrought by the haughty and the alterations made by the
arrogant, but these are mentioned in specific instances. Among them is
the account of Ibn Ṣūriyā. When the people of Khaybar questioned the
Point of the Qur'an about the prescribed punishment for adultery
committed by a married man or woman, He said that God's decree was
stoning. They denied this, saying that there was no such decree in the
Torah. He asked, ``Whom among your learned men do you acknowledge as an
authority and accept as truthful?'' They accepted Ibn Ṣūriyā, so He
summoned him and said, ``I adjure you by God, who parted the sea for
you, sent down manna upon you, shaded you with clouds, delivered you
from Pharaoh and his chiefs, and favored you above the peoples, to tell
us what Moses decreed as the punishment for a married adulterer and a
married adulteress.'' The substance of this is that He bound Ibn Ṣūriyā
by these solemn oaths to state what punishment for adultery by a married
person had been revealed in the Torah. He replied, ``Muḥammad, it is
stoning.'' He then asked, ``Why, then, has this decree been abrogated
among the Jews and ceased to be enforced?'' He replied, ``When
Nebuchadnezzar burned Jerusalem and put all the Jews to death, only a
small number of Jews remained in the land. The learned men of that age,
considering how few the Jews were and how numerous the Amalekites,
gathered to consult, reasoning that if they acted according to the
decree of the Torah, those who had escaped Nebuchadnezzar would be put
to death by the decree of the Book. For these considerations, they
abolished the death penalty altogether.'' Meanwhile, Gabriel descended
upon His radiant heart and presented this verse: ``They distort the
words from their proper places.''

\subsection{94}\label{section-284}

\begin{otherlanguage}{persian}

این یک موضع بود که ذکر شد و در این مقام مقصود از تحریف نچنانست که این
همج رعاع فهم نموده‌اند چنانچه بعضی میگویند که علمای یهود و نصاری آیاتی را
که در وصف طلعت محمّدیّه بود از کتاب محو نمودند و مخالف آن را ثبت کردند این
قول نهایت بی‌معنی و بی‌اصلست آیا میشود کسی که معتقد بکتابی گشته و من عند
اللّه دانسته آن را محو نماید و از این گذشته تورات در همهٔ روی ارض بود
منحصر به مکّه و مدینه نبود که بتوانند تغییر دهند و یا تبدیل نمایند بلکه
مقصود از تحریف همین است که الیوم جمیع علمای فرقان بآن مشغولند و آن تفسیر
و معنی نمودن کتابست بر هوی و میل خود و چون یهود در زمان آن حضرت آیات
تورات را که مدلّ بر ظهور آن حضرت بود بهوای خود تفسیر نمودند و ببیان آن
حضرت راضی نشدند لهذا حکم تحریف در بارهٔ آنها صدور یافت چنانچه الیوم مشهود
است که چگونه تحریف نمودند امّت فرقان آیات کتاب را در علامات ظهور بمیل و
هوای خود تفسیر مینمایند چنانچه مشهود است

\end{otherlanguage}

Īn yak mawḍiʿ būd ki dhikr shud va dar īn maqām maqṣūd az taḥrīf na
chunān ast ki īn hamaj-i raʿāʿ fahm namūda-and. Chunānchi baʿḍī mīgūyand
ki ʿulamā-yi Yahūd va Naṣārā āyātī rā ki dar vaṣf-i Ṭalʿat-i
Muḥammadiyya būd az Kitāb maḥv namūdand va mukhālif-i ān rā thabt
kardand. Īn qawl nihāyat bī-maʿnī va bī-aṣl ast. Āyā mīshavad kasī ki
muʿtaqid bi-kitābī gashta va min ʿindi Allāh dānista ān rā maḥv namāyad?
Va az īn gudhashta Tawrāt dar hama-yi rū-yi arḍ būd, munḥaṣir bi-Makka
va Madīna nabūd ki bitavānand taghyīr dihand va yā tabdīl namāyand.
Balki maqṣūd az taḥrīf hamīn ast ki al-yawm jamīʿ-i ʿulamā-yi Furqān
bi-ān mashghūl-and, va ān tafsīr va maʿnī namūdan-i Kitāb ast bar havā
va mayl-i khud. Va chūn Yahūd dar zamān-i ān Ḥaḍrat āyāt-i Tawrāt rā ki
mudall bar ẓuhūr-i ān Ḥaḍrat būd bi-havā-yi khud tafsīr namūdand va
bi-bayān-i ān Ḥaḍrat rāḍī nashudand, li-hādhā ḥukm-i taḥrīf dar bāra-yi
ānhā ṣudūr yāft. Chunānchi al-yawm mashhūd ast ki chigūna taḥrīf
namūdand ummat-i Furqān, āyāt-i Kitāb rā dar ʿalāmāt-i ẓuhūr bi-mayl va
havā-yi khud tafsīr mīnamāyand, chunānchi mashhūd ast.

This was one instance that has been mentioned. In this case, distortion
does not mean what these ignorant rabble have understood it to mean.
Some, for example, say that the Jewish and Christian scholars erased
from the Book the verses describing the beauty of Muḥammad and recorded
their opposite. This claim is utterly meaningless and unfounded. Could
someone who has come to believe in a book and regards it as coming from
God erase it? Moreover, the Torah existed throughout the earth; it was
not confined to Mecca and Medina, so that they could change or alter it.
Rather, distortion means precisely what all the learned men of the
Qur'an are engaged in today: interpreting and explaining the Book
according to their own desires and inclinations. Since the Jews in His
time interpreted according to their own desires those verses of the
Torah that indicated His appearance, and would not accept His
explanation, the judgment that they had distorted it was pronounced
against them. Likewise, it is evident today how the people of the Qur'an
have distorted it: they interpret the verses of the Book concerning the
signs of the Manifestation according to their own inclinations and
desires, as is plain to see.

\subsection{95}\label{section-285}

\begin{otherlanguage}{persian}

و در مقام دیگر میفرماید و قد کان فریق منهم یسمعون کلام اللّه ثمّ یحرّفونه
من بعد ما عقلوه و هم یعلمون و این آیه هم مدلّ است بر تحریف معانی کلام
الهی نه بر محو کلمات ظاهریّه چنانچه از آیه مستفاد میشود و عقول مستقیمه هم
ادراک مینماید

\end{otherlanguage}

Va dar maqām-i dīgar mīfarmāyad: Wa-qad kāna farīqun minhum yasmaʿūna
kalāma Allāhi thumma yuḥarrifūnahu min baʿdi mā ʿaqalūhu wa-hum
yaʿlamūna. Va īn āya ham mudall ast bar taḥrīf-i maʿānī-yi kalām-i
ilāhī, na bar maḥv-i kalimāt-i ẓāhiriyya, chunānchi az āya mustafād
mīshavad va ʿuqūl-i mustaqīma ham idrāk mīnamāyad.

Elsewhere He says, ``A group of them used to hear the word of God and
then distort it after they had understood it, knowingly.'' This verse,
too, indicates the distortion of the meanings of the divine word, not
the erasure of its outward wording, as can be understood from the verse
and as sound minds also perceive.

\subsection{96}\label{section-286}

\begin{otherlanguage}{persian}

و در موضع دیگر میفرماید فویل للّذین یکتبون الکتاب بأیدیهم ثمّ یقولون هذا
من عند اللّه لیشتروا به ثمناً قلیلاً الی آخر الآیة و این آیه در شأن علمای
یهود و بزرگان ایشان نازل شد که آن علما بواسطهٔ استرضای خاطر اغنیا و
استجلاب زخارف دنیا و اظهار غل و کفر الواحی چند بر ردّ حضرت نوشتند و
بدلائلی چند مستدلّ شدند که ذکر آنها جایز نه و نسبت دادند ادلّه‌های خود را
که از اسفار تورات مستفاد گشته

\end{otherlanguage}

Va dar mawḍiʿ-i dīgar mīfarmāyad: Fa-waylun li-alladhīna yaktubūna
al-kitāba bi-aydīhim thumma yaqūlūna hādhā min ʿindi Allāhi li-yashtarū
bihi thamanan qalīlan, ilā ākhiri al-āyati. Va īn āya dar shaʾn-i
ʿulamā-yi Yahūd va buzurgān-i īshān nāzil shud ki ān ʿulamā bi-vāsiṭa-yi
istirḍā-yi khāṭir-i aghniyā va istijlāb-i zakhārif-i dunyā va iẓhār-i
ghill va kufr alvāḥī chand bar radd-i Ḥaḍrat nivishtand va bi-dalāʾilī
chand mustadill shudand ki dhikr-i ānhā jāʾiz na, va nisbat dādand
adilla-hā-yi khud rā ki az asfār-i Tawrāt mustafād gashta.

Elsewhere He says, ``Woe to those who write the Book with their own
hands and then say, `This is from God,' to sell it for a paltry price,''
and so on to the end of the verse. This verse was revealed concerning
the Jewish scholars and leaders. To please the wealthy, attract worldly
riches, and express their malice and unbelief, those scholars wrote
several treatises rejecting Him. They supported their case with
arguments that it is not permissible to mention, and claimed that their
proofs had been derived from the books of the Torah.

\subsection{97}\label{section-287}

\begin{otherlanguage}{arabic}

چنانچه الیوم مشاهده میشود که چه مقدار رد بر این امر بدیع علمای جاهل عصر
نوشته‌اند و گمان نموده‌اند که این مفتریات مطابق آیات کتاب و موافق کلمات
اولی الألبابست

\end{otherlanguage}

Chunānchi al-yawm mushāhada mīshavad ki chi miqdār radd bar īn Amr-i
badīʿ ʿulamā-yi jāhil-i ʿaṣr nivishta-and va gumān namūda-and ki īn
muftariyāt muṭābiq-i āyāt-i Kitāb va muvāfiq-i kalimāt-i ulī al-albāb
ast.

Likewise, today one can see how many refutations of this wondrous Cause
the ignorant scholars of the age have written, imagining that these
fabrications accord with the verses of the Book and agree with the words
of those endowed with understanding.

\subsection{98}\label{section-288}

\begin{otherlanguage}{persian}

باری مقصود از این اذکار این بود که اگر بگویند این علائم مذکوره که از
انجیل ذکر شد تحریف یافته و رد نمایند و متمسّک بآیات و اخبار شوند مطّلع
باشید که کذب محض و افترای صرفست بلی ذکر تحریف باین معنی که ذکر شد در
اماکن مشخّصه هست چنانچه بعضی از آن را ذکر نمودیم تا معلوم و مبرهن شود بر
هر ذی بصری که احاطهٔ علوم ظاهره هم نزد بعضی از امّیّین الهی هست دیگر
معارضین باین خیال نیفتند و معارضه ننمایند که فلان آیه دلیل بر تحریفست و
این اصحاب از عدم اطّلاع ذکر این مراتب و مطالب را نموده‌اند و دیگر آنکه
اکثر آیات که مشعر بر تحریفست در بارهٔ یهود نازل شده لو انتم فی جزائر علم
الفرقان تحبرون

\end{otherlanguage}

Bārī, maqṣūd az īn adhkār īn būd ki agar bigūyand īn ʿalāʾim-i madhkūra
ki az Injīl dhikr shud taḥrīf yāfta va radd namāyand va mutamassik
bi-āyāt va akhbār shavand, muṭṭaliʿ bāshīd ki kidhb-i maḥḍ va iftirā-yi
ṣirf ast. Balī, dhikr-i taḥrīf bi-īn maʿnī ki dhikr shud dar amākin-i
mushakhkhaṣa hast, chunānchi baʿḍī az ān rā dhikr namūdīm tā maʿlūm va
mubarhan shavad bar har dhī baṣarī ki iḥāṭa-yi ʿulūm-i ẓāhira ham nazd-i
baʿḍī az ummiyyīn-i ilāhī hast. Dīgar muʿāriḍīn bi-īn khiyāl nayuftand
va muʿāraḍa nanamāyand ki fulān āya dalīl bar taḥrīf ast va īn aṣḥāb az
ʿadam-i iṭṭilāʿ dhikr-i īn marātib va maṭālib rā namūda-and. Va dīgar
ānki akthar-i āyāt ki mushʿir bar taḥrīf ast dar bāra-yi Yahūd nāzil
shuda, law antum fī jazāʾiri ʿilmi al-Furqāni taḥburūna.

The purpose of these remarks was that, should they claim that the
aforementioned signs cited from the Gospel have been distorted, reject
them, and appeal to verses and traditions, you may know that this is
pure falsehood and sheer fabrication. Yes, distortion in the sense
explained above is mentioned in specific passages. We have cited some of
these so that it may become clear and evident to everyone with insight
that some of the unlettered ones of God also possess a comprehensive
knowledge of outward learning. Let the opponents no longer imagine or
object that a particular verse proves distortion and that these
companions have discussed these matters and subjects out of ignorance.
Furthermore, most of the verses referring to distortion were revealed
concerning the Jews, if you would delight in the islands of the
knowledge of the Qur'an.

\subsection{99}\label{section-289}

\begin{otherlanguage}{persian}

اگرچه از بعضی حمقای ارض شنیده شد که انجیل سماوی در دست نصاری نیست و
بآسمان رفته دیگر غافل از اینکه از همین قول نسبت کمال ظلم و جبر برای حضرت
باری جلّ و عزّ ثابت میشود زیرا بعد از آنکه شمس جمال عیسی از میان قوم غایب
شد و بفلک چهارم ارتقا فرمود و کتاب حقّ جلّ ذکره که اعظم برهان اوست میان
خلق او آن هم غایب شود دیگر آن خلق از زمان عیسی تا زمان اشراق شمس محمّدی
بچه متمسّکند و بکدام امر مأمور و دیگر چگونه مورد انتقام منتقم حقیقی
میشوند و محلّ نزول عذاب و سیاط سلطان معنوی میگردند از همه گذشته انقطاع
فیض فیّاض و انسداد باب رحمت سلطان ایجاد لازم میآید فنعوذ باللّه عمّا یظنّ
العباد فی حقّه فتعالی عمّا هم یعرفون

\end{otherlanguage}

Agarchi az baʿḍī ḥumaqā-yi arḍ shinīda shud ki Injīl-i samāvī dar dast-i
Naṣārā nīst va bi-āsmān rafta, dīgar ghāfil az īnki az hamīn qawl
nisbat-i kamāl-i ẓulm va jabr barā-yi Ḥaḍrat-i Bārī jalla va ʿazza
thābit mīshavad. Zīrā baʿd az ānki shams-i jamāl-i ʿĪsā az miyān-i qawm
ghāʾib shud va bi-falak-i chahārum irtiqā farmūd va Kitāb-i Ḥaqq jalla
dhikruhu ki aʿẓam burhān-i ū-st miyān-i khalq-i ū ān ham ghāʾib shavad,
dīgar ān khalq az zamān-i ʿĪsā tā zamān-i ishrāq-i shams-i Muḥammadī
bi-chi mutamassik-and va bi-kudām amr maʾmūr? Va dīgar chigūna mawrid-i
intiqām-i Muntaqim-i Ḥaqīqī mīshavand va maḥall-i nuzūl-i ʿadhāb va
siyāṭ-i Sulṭān-i Maʿnavī mīgardand? Az hama gudhashta inqiṭāʿ-i fayḍ-i
Fayyāḍ va insidād-i bāb-i raḥmat-i Sulṭān-i Ījād lāzim mīāyad.
Fa-naʿūdhu bi-llāhi ʿammā yaẓunnu al-ʿibādu fī ḥaqqihi, fa-taʿālā ʿammā
hum yaʿrifūna.

Yet some of the foolish people of the earth have been heard to say that
the heavenly Gospel is not in the hands of the Christians but has
ascended to heaven, unaware that this very statement attributes the
utmost injustice and oppression to the Creator, exalted and mighty is
He. For if, after the sun of the beauty of Jesus disappeared from among
His people and ascended to the fourth heaven, the Book of God---exalted
be His mention---which is His greatest proof among His creatures, also
disappeared, what could those people have held to from the time of Jesus
until the dawning of the Muḥammadan Sun, and what command were they
required to follow? How, then, could they become subject to the
retribution of the true Avenger and the recipients of the punishment and
scourges of the spiritual Sovereign? Above all, this would necessarily
imply the cessation of the grace of the All-Bountiful and the closing of
the door of mercy of the Sovereign of creation. We take refuge in God
from what His servants imagine concerning Him! Exalted is He above all
that they comprehend.

\subsection{100}\label{section-290}

\begin{otherlanguage}{persian}

ای عزیز در این صبح ازلی که انوار اللّه نور السّموات و الأرض عالم را احاطه
نموده و سرادق عصمت و حفظ و یأبی اللّه الّا ان یتمّ نوره مرتفع گشته و ید
قدرت و بیده ملکوت کلّ شیء مبسوط و قائم شده کمر همّت را محکم باید بست که
شاید بعنایت و مکرمت الهی در مدینۀ قدسیّهٔ انّا للّه وارد شویم تا بمواقع عزّ
الیه راجعون مقرّ یابیم انشآءاللّه باید چشم دل را از اشارات آب و گل پاک
نمود تا ادراک مراتب ما لا نهایهٔ عرفان نمائید و حقّ را اظهر از آن بینید که
در اثبات وجودش بدلیلی محتاج شوید و یا بحجّتی تمسّک جوئید

\end{otherlanguage}

Ay ʿazīz, dar īn ṣubḥ-i azalī ki anvār-i Allāhu nūru al-samāwāti
wa-al-arḍi ʿālam rā iḥāṭa namūda va surādiq-i ʿiṣmat va ḥifẓ-i wa-yaʾbā
Allāhu illā an yutimma nūrahu murtafiʿ gashta va yad-i qudrat-i
wa-bi-yadihi malakūtu kulli shayʾin mabsūṭ va qāʾim shuda, kamar-i
himmat rā muḥkam bāyad bast ki shāyad bi-ʿināyat va makramat-i ilāhī dar
Madīna-yi Qudsiyya-yi innā li-llāhi vārid shavīm tā bi-mavāqiʿ-i ʿizz-i
ilayhi rājiʿūna maqarr yābīm, in shāʾa Allāh. Bāyad chashm-i dil rā az
ishārāt-i āb va gil pāk namūd tā idrāk-i marātib-i mā lā nihāya-yi
ʿirfān namāyīd va Ḥaqq rā aẓhar az ān bīnīd ki dar ithbāt-i vujūdash
bi-dalīlī muḥtāj shavīd va yā bi-ḥujjatī tamassuk jūyīd.

Dear one, in this eternal dawn, when the radiance of ``God is the light
of the heavens and the earth'' has encompassed the world, the canopy of
protection and preservation expressed in ``God will not allow anything
but the perfection of His light'' has been raised, and the hand of power
proclaimed in ``In His hand is the dominion of all things'' has been
stretched out and established, we must firmly resolve that, through
divine favor and bounty, we may enter the holy city of ``We belong to
God'' and find our abode in the realms of glory of ``To Him we return,''
God willing. The eye of the heart must be cleansed of the allusions of
water and clay, so that you may perceive the infinite degrees of
spiritual knowledge and see God as too manifest for you to need evidence
to establish His existence or to resort to any proof.

\subsection{101}\label{section-291}

\begin{otherlanguage}{persian}

ای سائل محبّ اگر در هوای روح روحانی طائری حقّ را ظاهر فوق کلّ شیء بینی
بقسمی که جز او را نیابی کان اللّه و لم یکن معه من شیء و این مقام مقدّس از
آنست که بدلیلی مدلّل شود و یا آنکه ببرهانی باهر آید و اگر در فضای قدس
حقیقت سایری کلّ اشیاء بمعروفیّت او معروفند و او بنفسه معروف بوده و خواهد
بود و اگر در ارض دلیل ساکنی کفایت کن بآنچه خود فرموده أ و لم یکفهم انّا
انزلنا علیک الکتاب اینست حجّتی که خود قرار فرموده و اعظم از این حجّت نبوده
و نیست دلیله آیاته و وجوده اثباته

\end{otherlanguage}

Ay sāʾil-i muḥibb, agar dar havā-yi rūḥ rūḥānī ṭāʾirī, Ḥaqq rā ẓāhir
fawqa kulli shayʾ bīnī, bi-qismī ki juz ū rā nayābī. Kāna Allāhu wa-lam
yakun maʿahu min shayʾin. Va īn maqām muqaddas az ān ast ki bi-dalīlī
mudallal shavad va yā ānki bi-burhānī bāhir āyad. Va agar dar faḍā-yi
quds-i ḥaqīqat sāʾirī, kull-i ashyā bi-maʿrūfiyyat-i ū maʿrūf-and va ū
bi-nafsihi maʿrūf būda va khvāhad būd. Va agar dar arḍ-i dalīl sākinī,
kifāyat kun bi-ānchi khud farmūda: A-wa-lam yakfihim annā anzalnā
ʿalayka al-kitāba? Īn ast ḥujjatī ki khud qarār farmūda va aʿẓam az īn
ḥujjat nabūda va nīst. Dalīluhu āyātuhu wa-vujūduhu ithbātuhu.

Loving seeker, if you are a spiritual bird in the air of the spirit, you
will see God manifest above all things, so that you will find none but
Him. ``God was, and there was nothing with Him.'' This station is too
holy to be established by evidence or made manifest through proof. If
you journey in the holy expanse of reality, all things are known through
the knowledge of Him, while He has always been and will always be known
through Himself. If you dwell in the land of evidence, be content with
what He Himself has said: ``Is it not enough for them that We have sent
down to you the Book?'' This is the proof that He Himself has ordained,
and no greater proof has ever existed or exists now. His verses are His
evidence, and His existence is His proof.

\subsection{102}\label{section-292}

\begin{otherlanguage}{persian}

در این وقت از اهل بیان و عرفا و حکما و علما و شهدای آن استدعا مینمایم که
وصایای الهی را که در کتاب فرموده فراموش ننمایند و همیشه ناظر باصل امر
باشند که مبادا حین ظهور آن جوهر الجواهر و حقیقة الحقائق و نور الأنوار
متمسّک ببعضی عبارات کتاب شوند و بر او وارد بیاورند آنچه را که در کور
فرقان وارد آمد چه که آن سلطان هویّه قادر است بر اینکه جمیع بیان و خلق آن
را بحرفی از بدایع کلمات خود قبض روح فرماید و یا بحرفی جمیع را حیات بدیعۀ
قدمیّه بخشد و از قبور نفس و هوی محشور و مبعوث نماید ملتفت و مراقب بوده که
جمیع منتهی بایمان باو و ادراک ایّام و لقای او میشود لیس البرّ ان تولّوا
وجوهکم قبل المشرق و المغرب ولکنّ البرّ من آمن باللّه و الیوم الآخر اسمعوا
یا اهل البیان ما وصّیناکم بالحقّ لعلّ تسکننّ فی ظلّ کان فی ایّام اللّه ممدوداً

\end{otherlanguage}

Dar īn vaqt az ahl-i Bayān va ʿurafā va ḥukamā va ʿulamā va shuhadā-yi
ān istidʿā mīnamāyam ki vaṣāyā-yi ilāhī rā ki dar Kitāb farmūda farāmūsh
nanamāyand va hamīsha nāẓir bi-aṣl-i Amr bāshand ki mabādā ḥīn-i ẓuhūr-i
ān Jawhar al-Javāhir va Ḥaqīqat al-Ḥaqāʾiq va Nūr al-Anvār mutamassik
bi-baʿḍī ʿibārāt-i Kitāb shavand va bar ū vārid biyāvarand ānchi rā ki
dar kawr-i Furqān vārid āmad. Chi ki ān Sulṭān-i Huviyya qādir ast bar
īnki jamīʿ-i Bayān va khalq-i ān rā bi-ḥarfī az badāyiʿ-i kalimāt-i khud
qabḍ-i rūḥ farmāyad va yā bi-ḥarfī jamīʿ rā ḥayāt-i badīʿa-yi qidamiyya
bakhshad va az qubūr-i nafs va havā maḥshūr va mabʿūth namāyad. Multafit
va murāqib būda ki jamīʿ muntahī bi-īmān bi-ū va idrāk-i ayyām va
liqā-yi ū mīshavad. Laysa al-birra an tuwallū wujūhakum qibala
al-mashriqi wa-al-maghribi wa-lākinna al-birra man āmana bi-llāhi
wa-al-yawmi al-ākhiri. Ismaʿū yā ahla al-Bayāni mā waṣṣaynākum
bi-al-ḥaqqi laʿalla taskununna fī ẓillin kāna fī ayyāmi Allāhi mamdūdan.

At this time I appeal to the people of the Bayān, its mystics,
philosophers, scholars, and witnesses, not to forget the divine counsels
He has given in the Book, and always to keep their gaze upon the essence
of the Cause. Otherwise, when that Essence of essences, that Reality of
realities, that Light of lights appears, they may cling to certain
expressions in the Book and inflict upon Him what was inflicted in the
cycle of the Qur'an. For that Sovereign of divine Being has the power,
through a single letter of His wondrous words, to take the life of the
entire Bayān and its people, or through a single letter to bestow upon
all a wondrous, everlasting life and raise and resurrect them from the
graves of self and desire. Be attentive and vigilant, for everything
culminates in faith in Him, attaining His days, and meeting Him.
``Righteousness is not that you turn your faces toward the east and the
west; rather, righteous is the one who believes in God and the Last
Day.'' Listen, people of the Bayān, to what We have counseled you in
truth, so that you may dwell beneath a shade stretched out in the days
of God.

\subsection{103}\label{section-293}

\begin{otherlanguage}{arabic}

الباب المذکور فی بیان انّ شمس الحقیقة و مظهر نفس اللّه لیکوننّ سلطاناً علی
من فی السّموات و الأرض و ان لن یطیعه احد من اهل الأرض و غنیّاً عن کلّ من فی
الملک و ان لم یکن عنده دینار کذلک نظهر لک من اسرار الأمر و نلقی علیک من
جواهر الحکمة لتطیرنّ بجناحی الانقطاع فی الهوآء الّذی کان عن الأبصار مستوراً

\end{otherlanguage}

Al-bābu al-madhkūru fī bayāni anna Shamsa al-Ḥaqīqati wa-Maẓhara Nafsi
Allāhi la-yakūnanna sulṭānan ʿalā man fī al-samāwāti wa-al-arḍi wa-in
lan yuṭīʿahu aḥadun min ahli al-arḍi wa-ghaniyyan ʿan kulli man fī
al-mulki wa-in lam yakun ʿindahu dīnār. Kadhālika nuẓhiru laka min
asrāri al-amri wa-nulqī ʿalayka min jawāhiri al-ḥikmati li-taṭīranna
bi-janāḥayi al-inqiṭāʿi fī al-hawāʾi alladhī kāna ʿani al-abṣāri
mastūran.

The aforementioned chapter explains that the Sun of Reality and the
Manifestation of God's own Self will assuredly be sovereign over all who
are in the heavens and on earth, even if no one on earth obeys Him, and
independent of all who are in the realm of creation, even if He
possesses not a single dinar. Thus We disclose to you some of the
mysteries of the Cause and impart to you some of the gems of wisdom, so
that you may soar on the wings of detachment in the atmosphere hidden
from sight.

\subsection{104}\label{section-294}

\begin{otherlanguage}{arabic}

لطائف و جواهر این باب آنکه بر صاحبان نفوس زکیّه و مرایای قدسیّه مبرهن و
واضح شود که شموس حقیقت و مرایای احدیّت در هر عصر و زمان که از خیام غیب
هویّه بعالم شهاده ظهور میفرمایند برای تربیت ممکنات و ابلاغ فیض بر همۀ
موجودات با سلطنتی قاهر و سطوتی غالب ظاهر میشوند چه که این جواهر مخزونه و
کنوز غیبیّهٔ مکنونه محلّ ظهور یفعل اللّه ما یشآء و یحکم ما یریدند

\end{otherlanguage}

Laṭāʾif va javāhir-i īn bāb ānkih bar ṣāḥibān-i nufūs-i zakiyyih va
marāyā-yi qudsiyyih mubarhan va vāḍiḥ shavad kih shumūs-i ḥaqīqat va
marāyā-yi aḥadiyyat dar har ʿaṣr va zamān kih az khiyām-i ghayb-i
huviyyih bi-ʿālam-i shahādih ẓuhūr mīfarmāyand barā-yi tarbiyat-i
mumkināt va iblāgh-i fayḍ bar hamih-yi mawjūdāt bā salṭanatī qāhir va
saṭvatī ghālib ẓāhir mīshavand chih kih īn javāhir-i makhzūnih va
kunūz-i ghaybiyyih-yi maknūnih maḥall-i ẓuhūr-i yafʿalu Allāhu mā yashāʾ
wa yaḥkumu mā yurīd-and.

The subtleties and gems of this chapter are that it should become clear
and demonstrable to those with pure souls and sanctified hearts that, in
every age and time, when the Suns of Reality and Mirrors of Oneness
appear from the tabernacles of the unseen divine Being in the visible
world, they appear with an all-subduing sovereignty and an irresistible
might, to educate all contingent beings and impart grace to all that
exists. For these hidden gems and concealed, unseen treasures are the
places where ``God does whatever He wills and ordains whatever He
desires'' is manifested.

\subsection{105}\label{section-295}

\begin{otherlanguage}{persian}

و بر اولی العلم و افئدهٔ منیره واضحست که غیب هویّه و ذات احدیّه مقدّس از
بروز و ظهور و صعود و نزول و دخول و خروج بوده و متعالیست از وصف هر واصفی
و ادراک هر مدرکی لم‌یزل در ذات خود غیب بوده و هست و لا‌یزال بکینونت خود
مستور از ابصار و انظار خواهد بود لا تدرکه الأبصار و هو یدرک الأبصار و هو
اللّطیف الخبیر چه میان او و ممکنات نسبت و ربط و فصل و وصل و یا قرب و بعد
و جهت و اشاره بهیچوجه ممکن نه زیرا که جمیع من فی السّموات و الأرض بکلمئ
امر او موجود شدند و بارادهٔ او که نفس مشیّت است از عدم و نیستی بحت باتّ
بعرصهٔ شهود و هستی قدم گذاشتند

\end{otherlanguage}

Va bar ulī al-ʿilm va afʾidih-yi munīrih vāḍiḥ ast kih ghayb-i huviyyih
va dhāt-i aḥadiyyih muqaddas az burūz va ẓuhūr va ṣuʿūd va nuzūl va
dukhūl va khurūj būdih va mutaʿālīst az vaṣf-i har vāṣifī va idrāk-i har
mudrikī. Lam-yazal dar dhāt-i khud ghayb būdih va hast va lā-yazāl
bi-kaynūnat-i khud mastūr az abṣār va anẓār khvāhad būd. Lā tudrikuhu
al-abṣār wa huwa yudriku al-abṣār wa huwa al-Laṭīf al-Khabīr. Chih
miyān-i ū va mumkināt nisbat va rabṭ va faṣl va vaṣl va yā qurb va buʿd
va jihat va ishārih bi-hīch-vajh mumkin nah zīrā kih jamīʿ-i man fī
al-samāwāt wa al-arḍ {[}corrupted text{]} amr-i ū mawjūd shudand va
bi-irādih-yi ū kih nafs-i mashiyyat ast az ʿadam va nīstī-yi baḥt-i bātt
bi-ʿarṣih-yi shuhūd va hastī qadam gudhāshtand.

It is clear to those endowed with knowledge and illumined hearts that
the unseen divine Being and the Essence of Oneness have always been
sanctified above emergence and manifestation, ascent and descent, entry
and departure, and are exalted beyond the description of anyone who
describes and the understanding of anyone who understands. He has ever
been, and remains, hidden in His Essence, and in His own Being He will
forever remain concealed from eyes and sight. ``Eyes do not perceive
Him, but He perceives all eyes, and He is the Subtle, the All-Aware.''
For no relationship or connection, separation or union, nearness or
remoteness, direction or reference is in any way possible between Him
and contingent beings, since all who are in the heavens and on earth
came into existence {[}corrupted original text{]} His command, and by
His will, which is volition itself, stepped from nonexistence and sheer,
absolute nothingness into the realm of visibility and existence.

\subsection{106}\label{section-296}

\begin{otherlanguage}{persian}

سبحان‌اللّه بلکه میانهٔ ممکنات و کلمهٔ او هم نسبت و ربطی نبوده و نخواهد بود
و یحذّرکم اللّه نفسه بر این مطلب برهانیست واضح و کان اللّه و لم یکن معه من
شیء دلیلی است لائح چنانچه جمیع انبیا و اوصیا و علما و عرفا و حکما بر عدم
بلوغ معرفت آن جوهر الجواهر و بر عجز از عرفان و وصول آن حقیقة الحقائق مقرّ
و مذعنند

\end{otherlanguage}

Subḥān Allāh, balkih miyānih-yi mumkināt va kalimih-yi ū ham nisbat va
rabṭī nabūdih va nakhvāhad būd. Wa yuḥadhdhirukumu Allāhu nafsahu bar īn
maṭlab burhānīst vāḍiḥ va kāna Allāhu wa lam yakun maʿahu min shayʾ
dalīlī ast lāʾiḥ chunānchih jamīʿ-i anbiyā va awṣiyā va ʿulamā va ʿurafā
va ḥukamā bar ʿadam-i bulūgh-i maʿrifat-i ān jawhar al-javāhir va bar
ʿajz az ʿirfān va vuṣūl-i ān ḥaqīqat al-ḥaqāʾiq muqirr va mudhʿinand.

Glory be to God! Rather, there has never been, nor will there ever be,
any relationship or connection even between contingent beings and His
Word. ``God cautions you concerning Himself'' is a clear proof of this,
and ``God was, and there was nothing with Him'' is an evident testimony.
Thus all the prophets, their appointed successors, the scholars,
mystics, and sages acknowledge and affirm that knowledge of that Essence
of essences is unattainable, and that they are powerless to recognize or
attain that Reality of realities.

\subsection{107}\label{section-297}

\begin{otherlanguage}{persian}

و چون ابواب عرفان ذات ازل بر وجه ممکنات مسدود شد لهذا باقتضای رحمت واسعهٔ
سبقت رحمته کلّ شیء و وسعت رحمتی کلّ شیء جواهر قدس نورانی را از عوالم روح
روحانی بهیاکل عزّ انسانی در میان خلق ظاهر فرمود تا حکایت نمایند از آن ذات
ازلیّه و ساذج قدمیّه و این مرایای قدسیّه و مطالع هویّه بتمامهم از آن شمس
وجود و جوهر مقصود حکایت مینمایند مثلاً علم ایشان از علم او و قدرت ایشان
از قدرت او و سلطنت ایشان از سلطنت او و جمال ایشان ازجمال او و ظهور ایشان
از ظهور او و ایشانند مخازن علوم ربّانی و مواقع حکمت صمدانی و مظاهر فیض
نامتناهی و مطالع شمس لایزالی چنانچه میفرماید لا فرق بینک و بینهم الّا
بأنّهم عبادک و خلقک و اینست مقام انا هو و هو انا که در حدیث مذکور است

\end{otherlanguage}

Va chūn abvāb-i ʿirfān-i dhāt-i azal bar vajh-i mumkināt masdūd shud
li-hādhā bi-iqtiḍā-yi raḥmat-i vāsiʿih-yi sabaqat raḥmatuhu kulla shayʾ
va wasiʿat raḥmatī kulla shayʾ javāhir-i quds-i nūrānī rā az ʿavālim-i
rūḥ-i rūḥānī bi-hayākil-i ʿizz-i insānī dar miyān-i khalq ẓāhir farmūd
tā ḥikāyat namāyand az ān dhāt-i azaliyyih va sādhij-i qidamiyyih va īn
marāyā-yi qudsiyyih va maṭāliʿ-i huviyyih bi-tamāmihim az ān shams-i
vujūd va jawhar-i maqṣūd ḥikāyat mīnamāyand mathalan ʿilm-i īshān az
ʿilm-i ū va qudrat-i īshān az qudrat-i ū va salṭanat-i īshān az
salṭanat-i ū va jamāl-i īshān az jamāl-i ū va ẓuhūr-i īshān az ẓuhūr-i ū
va īshānand makhāzin-i ʿulūm-i rabbānī va mavāqiʿ-i ḥikmat-i ṣamadānī va
maẓāhir-i fayḍ-i nāmutanāhī va maṭāliʿ-i shams-i lā-yazālī chunānchih
mīfarmāyad lā farqa baynaka wa baynahum illā bi-annahum ʿibāduka wa
khalquka va īnast maqām-i anā huwa wa huwa anā kih dar ḥadīth madhkūr
ast.

Since the doors to knowing the eternal Essence are closed to contingent
beings, He has, in accordance with His all-embracing mercy---``His mercy
has preceded all things'' and ``My mercy has encompassed all
things''---caused luminous Gems of holiness to appear among creation,
from the spiritual realms of the spirit, in temples of human glory, so
that they may reflect that eternal Essence and primordial Reality. These
sanctified Mirrors and Dawning-places of the divine Being all reflect
that Sun of existence and Essence of all desire. For example, their
knowledge is from His knowledge, their power from His power, their
sovereignty from His sovereignty, their beauty from His beauty, and
their manifestation from His manifestation. They are the treasuries of
divine knowledge, the repositories of the wisdom of the Self-Sufficient,
the manifestations of boundless grace, and the dawning-places of the
everlasting Sun. As He says, ``There is no distinction between You and
them, except that they are Your servants and Your creation.'' This is
the station of ``I am He, and He is I,'' mentioned in the tradition.

\subsection{108}\label{section-298}

\begin{otherlanguage}{persian}

و احادیث و اخبار مدلّهٔ بر این مطلب بسیار است و این بنده نظر باختصار متعرّض
ذکر آنها نشدم بلکه آنچه در آسمانها و زمین است محالّ بروز صفات و اسمای
الهی هستند چنانچه در هر ذرّه آثار تجلّی آن شمس حقیقی ظاهر و هویداست که
گویا بدون ظهور آن تجلّی در عالم ملکی هیچ شیء بخلعت هستی مفتخر نیاید و
بوجود مشرّف نشود چه آفتابهای معارف که در ذرّه مستور شده و چه بحرهای حکمت
که در قطره پنهان گشته خاصّه انسان که از بین موجودات باین خلع تخصیص یافته
و باین شرافت ممتاز گشته چنانچه جمیع صفات و اسمای الهی از مظاهر انسانی
بنحو اکمل و اشرف ظاهر و هویداست و کلّ این اسماء و صفات راجع باوست اینست
که فرموده الانسان سرّی و انا سرّه و آیات متواتره که مدلّ و مشعر بر این مطلب
رقیق لطیف است در جمیع کتب سماویّه و صحف الهیّه مسطور و مذکور است چنانچه
میفرماید سنریهم آیاتنا فی الآفاق و فی انفسهم و در مقام دیگر میفرماید و
فی انفسکم أ فلا تبصرون و در مقام دیگر میفرماید و لا تکونوا کالّذین نسوا
اللّه فأنساهم انفسهم چنانچه سلطان بقا روح من فی سرادق العمآء فداه
میفرماید من عرف نفسه فقد عرف ربّه

\end{otherlanguage}

Va aḥādīth va akhbār-i mudillih-yi bar īn maṭlab bisyār ast va īn bandih
naẓar bi-ikhtiṣār mutaʿarriḍ-i dhikr-i ānhā nashudam balkih ānchih dar
āsmānhā va zamīn ast maḥāll-i burūz-i ṣifāt va asmā-yi ilāhī hastand
chunānchih dar har dharrih āthār-i tajallī-yi ān shams-i ḥaqīqī ẓāhir va
huvaydāst kih gūyā bidūn-i ẓuhūr-i ān tajallī dar ʿālam-i mulkī hīch
shayʾ bi-khilʿat-i hastī muftakhir nayāyad va bi-vujūd musharraf
nashavad. Chih āftābhā-yi maʿārif kih dar dharrih mastūr shudih va chih
baḥrhā-yi ḥikmat kih dar qaṭrih pinhān gashtih khāṣṣih insān kih az
bayn-i mawjūdāt bi-īn khilaʿ takhṣīṣ yāftih va bi-īn sharāfat mumtāz
gashtih chunānchih jamīʿ-i ṣifāt va asmā-yi ilāhī az maẓāhir-i insānī
bi-naḥv-i akmal va ashraf ẓāhir va huvaydāst va kull-i īn asmā va ṣifāt
rājiʿ bi-ūst īnast kih farmūdih al-insānu sirrī wa anā sirruhu va āyāt-i
mutavātirih kih mudill va mushʿir bar īn maṭlab-i raqīq-i laṭīf ast dar
jamīʿ-i kutub-i samāviyyih va ṣuḥuf-i ilāhiyyih masṭūr va madhkūr ast
chunānchih mīfarmāyad sa-nurīhim āyātinā fī al-āfāqi wa fī anfusihim va
dar maqām-i dīgar mīfarmāyad wa fī anfusikum a-fa-lā tubṣirūn va dar
maqām-i dīgar mīfarmāyad wa lā takūnū ka-alladhīna nasū Allāha
fa-ansāhum anfusahum chunānchih sulṭān-i baqā rūḥ-i man fī surādiq
al-ʿamāʾ fidāh mīfarmāyad man ʿarafa nafsahu fa-qad ʿarafa rabbahu.

There are many traditions and reports bearing witness to this matter,
but for the sake of brevity this servant has not undertaken to mention
them. Indeed, whatever is in the heavens and on earth is a place where
the divine names and attributes appear. Thus, in every atom the traces
of the radiance of that true Sun are manifest and evident, so that,
without the appearance of that radiance in the earthly realm, it is as
though no thing could be honored with the robe of being or ennobled by
existence. How many suns of knowledge lie concealed within an atom, and
how many oceans of wisdom are hidden within a drop! This is especially
true of the human being, who, among all beings, has been singled out for
these robes and distinguished by this honor. All the divine names and
attributes are manifest and evident in human manifestations in the most
perfect and noble manner, and all these names and attributes refer to
the human being. Thus He has said, ``The human being is My mystery, and
I am his mystery.'' Numerous verses indicating and intimating this
delicate, subtle truth are recorded and mentioned in all the heavenly
Books and divine Scriptures. As He says, ``We will show them Our signs
on the horizons and within themselves.'' Elsewhere He says, ``And within
yourselves---do you not see?'' And elsewhere He says, ``Do not be like
those who forgot God, so He caused them to forget themselves.''
Likewise, the Sovereign of eternity---may the souls of all who dwell in
the tabernacle of the cloud of unknowing be sacrificed for Him---says,
``Whoever knows himself has known his Lord.''

\subsection{109}\label{section-299}

\begin{otherlanguage}{arabic}

قسم بخدا ای مخدوم من اگر قدری در این عبارات تفکّر فرمائی ابواب حکمت الهیّه
و مصاریع علم نامتناهی را بر وجه خود گشوده یابی

\end{otherlanguage}

Qasam bi-Khudā ay makhdūm-i man agar qadrī dar īn ʿibārāt tafakkur
farmāʾī abvāb-i ḥikmat-i ilāhiyyih va maṣārīʿ-i ʿilm-i nāmutanāhī rā bar
vajh-i khud gushūdih yābī.

By God, my honored sir, if you reflect a little on these words, you will
find the doors of divine wisdom and the portals of boundless knowledge
opened before you.

\subsection{110}\label{section-300}

\begin{otherlanguage}{persian}

باری از این بیانات معلوم شد که جمیع اشیاء حاکی از اسماء و صفات الهیّه
هستند هر کدام بقدر استعداد خود مدلّ و مشعرند بر معرفت الهیّه بقسمی که
احاطه کرده است ظهورات صفاتیّه و اسمائیّه همهٔ غیب و شهود را اینست که
میفرماید أ یکون لغیرک من الظّهور ما لیس لک حتّی یکون هو المظهر لک عمیت عین
لا تراک و باز سلطان بقا میفرماید ما رأیت شیئاً الّا و قد رأیت اللّه فیه او
قبله او بعده و در روایت کمیل نور اشرق من صبح الأزل فیلوح علی هیاکل
التّوحید آثاره و انسان که اشرف و اکمل مخلوقات است اشدّ دلالةً و اعظم حکایةً
است از سایر معلومات و اکمل انسان و افضل و الطف او مظاهر شمس حقیقتند بلکه
ماسوای ایشان موجودند بارادهٔ ایشان و متحرّکند بافاضهٔ ایشان لولاک لما خلقت
الأفلاک بلکه کل در ساحت قدس ایشان معدوم صرف و مفقود بحتند بلکه منزّه است
ذکر ایشان از ذکر غیر و مقدّس است وصف ایشان از وصف ماسوی و این هیاکل قدسیّه
مرایای اوّلیّۀ ازلیّه هستند که حکایت نموده‌اند از غیب الغیوب و از کلّ اسماء و
صفات او از علم و قدرت و سلطنت و عظمت و رحمت و حکمت و عزّت و جود و کرم و
جمیع این صفات از ظهور این جواهر احدیّه ظاهر و هویداست

\end{otherlanguage}

Bārī az īn bayānāt maʿlūm shud kih jamīʿ-i ashyāʾ ḥākī az asmā va
ṣifāt-i ilāhiyyih hastand har kudām bi-qadr-i istiʿdād-i khud mudill va
mushʿirand bar maʿrifat-i ilāhiyyih bi-qismī kih iḥāṭih kardih ast
ẓuhūrāt-i ṣifātiyyih va asmāʾiyyih hamih-yi ghayb va shuhūd rā īnast kih
mīfarmāyad a-yakūnu li-ghayrika mina al-ẓuhūri mā laysa laka ḥattā
yakūna huwa al-muẓhira laka ʿamiyat ʿaynun lā tarāka va bāz sulṭān-i
baqā mīfarmāyad mā raʾaytu shayʾan illā wa qad raʾaytu Allāha fīhi aw
qablahu aw baʿdahu va dar rivāyat-i Kumayl nūrun ashraqa min ṣubḥi
al-azali fa-yalūḥu ʿalā hayākili al-tawḥīdi āthāruhu va insān kih ashraf
va akmal-i makhlūqāt ast ashaddu dalālatan va aʿẓamu ḥikāyatan ast az
sāʾir-i maʿlūmāt va akmal-i insān va afḍal va alṭaf-i ū maẓāhir-i
shams-i ḥaqīqatand balkih māsivā-yi īshān mawjūdand bi-irādih-yi īshān
va mutaḥarrikand bi-ifāḍih-yi īshān lawlāka la-mā khalaqtu al-aflāka
balkih kull dar sāḥat-i quds-i īshān maʿdūm-i ṣirf va mafqūd-i baḥtand
balkih munazzah ast dhikr-i īshān az dhikr-i ghayr va muqaddas ast
vaṣf-i īshān az vaṣf-i māsivā va īn hayākil-i qudsiyyih marāyā-yi
avvaliyyih-yi azaliyyih hastand kih ḥikāyat nimūdih-and az ghayb
al-ghuyūb va az kull-i asmā va ṣifāt-i ū az ʿilm va qudrat va salṭanat
va ʿaẓamat va raḥmat va ḥikmat va ʿizzat va jūd va karam va jamīʿ-i īn
ṣifāt az ẓuhūr-i īn javāhir-i aḥadiyyih ẓāhir va huvaydāst.

In short, these statements make clear that all things reflect the divine
names and attributes. Each, according to its capacity, points and
alludes to the knowledge of God, so that the manifestations of His names
and attributes encompass everything unseen and seen. Thus He says,
``Does anything other than You possess a manifestation that You do not
possess, such that it could be what makes You manifest? Blind is the eye
that does not see You.'' Again, the Sovereign of eternity says, ``I have
seen nothing without seeing God within it, or before it, or after it.''
And in the tradition of Kumayl: ``A light shone from the dawn of
eternity, and its traces gleam upon the temples of divine unity.'' The
human being, who is the noblest and most perfect of creatures, points to
and reflects Him more powerfully than all other known things. The most
perfect, excellent, and refined among human beings are the
Manifestations of the Sun of Reality. Indeed, all others exist by their
will and move through the grace they impart: ``Were it not for you, I
would not have created the celestial spheres.'' Indeed, all are sheer
nonexistence and utter nothingness in their holy presence. Rather,
mention of them is exalted above mention of others, and their
description is sanctified above the description of anything else. These
holy Temples are the primal, eternal Mirrors that reflect the Unseen of
the unseen and all His names and attributes---knowledge, power,
sovereignty, grandeur, mercy, wisdom, glory, generosity, and bounty. All
these attributes become manifest and evident through the appearance of
these Gems of Oneness.

\subsection{111}\label{section-301}

\begin{otherlanguage}{persian}

و این صفات مختصّ ببعضی دون بعضی نبوده و نیست بلکه جمیع انبیای مقرّبین و
اصفیای مقدّسین باین صفات موصوف و باین اسماء موسومند نهایت بعضی در بعضی
مراتب اشدّ ظهوراً و اعظم نوراً ظاهر میشوند چنانچه میفرماید تلک الرّسل فضّلنا
بعضهم علی بعض پس معلوم و محقّق شد که محلّ ظهور و بروز جمیع این صفات عالیه
و اسمای غیر متناهیه انبیا و اولیای او هستند خواه بعضی از این صفات در آن
هیاکل نوریّه بر حسب ظاهر ظاهر شود و خواه نشود نه اینست که اگر صفتی بر حسب
ظاهر از آن ارواح مجرّده ظاهر نشود نفی آن صفت از آن محالّ صفات الهیّه و
معادن اسماء ربوبیّه شود لهذا بر همۀ این وجودات منیره و طلعات بدیعه حکم
جمیع صفات اللّه از سلطنت و عظمت و امثال آن جاریست اگرچه بر حسب ظاهر
بسلطنت ظاهره و غیر آن ظاهر نشوند و این فقره بر هر ذی بصری ثابت و محقّق
است دیگر احتیاج برهان نیست

\end{otherlanguage}

Va īn ṣifāt mukhtaṣṣ bi-baʿḍī dūn-i baʿḍī nabūdih va nīst balkih jamīʿ-i
anbiyā-yi muqarrabīn va aṣfiyā-yi muqaddasīn bi-īn ṣifāt mawṣūf va bi-īn
asmā mawsūmand nihāyat baʿḍī dar baʿḍī marātib ashaddu ẓuhūran va aʿẓamu
nūran ẓāhir mīshavand chunānchih mīfarmāyad tilka al-rusulu faḍḍalnā
baʿḍahum ʿalā baʿḍ. Pas maʿlūm va muḥaqqaq shud kih maḥall-i ẓuhūr va
burūz-i jamīʿ-i īn ṣifāt-i ʿāliyih va asmā-yi ghayr-i mutanāhiyih anbiyā
va awliyā-yi ū hastand khvāh baʿḍī az īn ṣifāt dar ān hayākil-i nūriyyih
bar ḥasab-i ẓāhir ẓāhir shavad va khvāh nashavad nah īnast kih agar
ṣifatī bar ḥasab-i ẓāhir az ān arvāḥ-i mujarradih ẓāhir nashavad nafy-i
ān ṣifat az ān maḥāll-i ṣifāt-i ilāhiyyih va maʿādin-i asmā-yi
rubūbiyyih shavad li-hādhā bar hamih-yi īn vujūdāt-i munīrih va ṭalʿāt-i
badīʿih ḥukm-i jamīʿ-i ṣifāt Allāh az salṭanat va ʿaẓamat va amthāl-i ān
jārīst agarchih bar ḥasab-i ẓāhir bi-salṭanat-i ẓāhirih va ghayr-i ān
ẓāhir nashavand va īn fiqrih bar har dhī baṣarī thābit va muḥaqqaq ast
dīgar iḥtiyāj-i burhān nīst.

These attributes have never been, and are not, exclusive to some rather
than others. Rather, all the prophets brought near to God and all His
sanctified chosen ones possess these attributes and bear these names. It
is only that some appear, in certain stations, with a more intense
manifestation and a greater light. As He says, ``These messengers---We
have favored some of them above others.'' It is thus clear and
established that His prophets and holy ones are the places where all
these exalted attributes and limitless names appear and are manifested,
whether or not some of these attributes are outwardly manifested in
those luminous Temples. If an attribute is not outwardly manifested by
those detached Spirits, this does not mean that it is to be denied to
those repositories of divine attributes and mines of the names of
lordship. Therefore, all the attributes of God, such as sovereignty,
grandeur, and the like, apply to all these luminous Beings and wondrous
Countenances, even if they do not outwardly appear with visible
sovereignty or other such attributes. This matter is established and
certain to everyone endowed with sight, and requires no further proof.

\subsection{112}\label{section-302}

\begin{otherlanguage}{persian}

بلی این عباد چون از عیون صافیهٔ منیرئ علوم الهیّه تفاسیر کلمات قدسیّه را
اخذ ننموده‌اند لهذا تشنه و افسرده در وادی ظنون و غفلت سایرند و از بحر عذب
فرات معرض شده در حول ملح اجاج طائفند چنانچه در وصف ایشان ورقاء هویّه بیان
فرموده و ان یروا سبیل الرّشد لا یتّخذوه سبیلاً و ان یروا سبیل الغیّ یتّخذوه
سبیلاً ذلک بأنّهم کذّبوا بآیاتنا و کانوا عنها غافلین که ترجمهٔ آن اینست اگر
ببینند راه صلاح و رستگاری را آن را اخذ نمینمایند و بآن اقبال نمیکنند و
امّا اگر راه باطل و طغیان و ضلالت را مشاهده کنند آن را برای خود راه وصول
بحقّ قرار دهند و این اقبال بباطل و اعراض از حقّ ظاهر نشد یعنی باین ضلالت و
گمراهی مبتلا نشدند مگر بجزای آنکه تکذیب کردند آیات ما را و بودند از نزول
آیات ما و ظهورات آن غفلت‌کنندگان

\end{otherlanguage}

Balī īn ʿibād chūn az ʿuyūn-i ṣāfiyih-yi {[}corrupted text{]} ʿulūm-i
ilāhiyyih tafāsīr-i kalimāt-i qudsiyyih rā akhdh nanimūdih-and li-hādhā
tishnih va afsurdih dar vādī-yi ẓunūn va ghaflat sāyirand va az baḥr-i
ʿadhb-i furāt muʿriḍ shudih dar ḥawl-i milḥ-i ujāj ṭāʾifand chunānchih
dar vaṣf-i īshān varqā-yi huviyyih bayān farmūdih wa in yaraw sabīla
al-rushdi lā yattakhidhūhu sabīlan wa in yaraw sabīla al-ghayyi
yattakhidhūhu sabīlan dhālika bi-annahum kadhdhabū bi-āyātinā wa kānū
ʿanhā ghāfilīn kih tarjumih-yi ān īnast agar bibīnand rāh-i ṣalāḥ va
rastgārī rā ān rā akhdh namīnamāyand va bi-ān iqbāl namīkunand va ammā
agar rāh-i bāṭil va ṭughyān va ḍalālat rā mushāhidih kunand ān rā
barā-yi khud rāh-i vuṣūl bi-Ḥaqq qarār dihand va īn iqbāl bi-bāṭil va
iʿrāḍ az Ḥaqq ẓāhir nashud yaʿnī bi-īn ḍalālat va gumrāhī mubtalā
nashudand magar bi-jazā-yi ānkih takdhīb kardand āyāt-i mā rā va būdand
az nuzūl-i āyāt-i mā va ẓuhūrāt-i ān ghaflat-kunandigān.

Indeed, because these servants have not drawn the interpretations of the
holy words from the pure {[}corrupted original text{]} springs of divine
knowledge, they wander, thirsty and despondent, in the valley of
conjecture and heedlessness. Turning away from the ocean of sweet,
refreshing water, they circle around bitter salt water. As the Dove of
the divine Being has declared in describing them, ``If they see the way
of right guidance, they do not take it as a way; but if they see the way
of error, they take it as a way. This is because they denied Our signs
and were heedless of them.'' Its translation is this: If they see the
path of righteousness and salvation, they do not adopt it or turn toward
it; but if they behold the path of falsehood, rebellion, and error, they
make it their path to attaining God. This turning toward falsehood and
turning away from God came about---that is, they became afflicted with
this error and waywardness---only as a recompense for their denial of
Our verses and their heedlessness of the descent of Our verses and their
manifestations.

\subsection{113}\label{section-303}

\begin{otherlanguage}{persian}

چنانچه مشاهده شد در این ظهور بدیع منیع که کرورها آیات الهیّه از سماء قدرت
و رحمت نازل شد با وجود این جمیع خلق اعراض نموده و تمسّک جسته‌اند باقوال
عبادی که یک حرف از آن را ادراک نمینمایند از این جهت است که در امثال این
مسائل واضحه شبهه نموده و خود را از رضوان علم احدیّه و ریاض حکمت صمدیّه
محروم نموده‌اند

\end{otherlanguage}

Chunānchih mushāhidih shud dar īn ẓuhūr-i badīʿ-i manīʿ kih kurūrhā
āyāt-i ilāhiyyih az samā-yi qudrat va raḥmat nāzil shud bā vujūd-i īn
jamīʿ-i khalq iʿrāḍ nimūdih va tamassuk justih-and bi-aqvāl-i ʿibādī kih
yik ḥarf az ān rā idrāk namīnamāyand az īn jihat ast kih dar amthāl-i īn
masāʾil-i vāḍiḥih shubhih nimūdih va khud rā az riḍvān-i ʿilm-i
aḥadiyyih va riyāḍ-i ḥikmat-i ṣamadiyyih maḥrūm nimūdih-and.

As has been witnessed in this wondrous, exalted Revelation, countless
millions of divine verses have descended from the heaven of power and
mercy. Despite this, all people have turned away and clung to the
sayings of servants who do not comprehend a single letter of those
verses. For this reason they have cast doubt on such clear matters and
deprived themselves of the paradise of the knowledge of Oneness and the
gardens of the wisdom of the Self-Sufficient.

\subsection{114}\label{section-304}

\begin{otherlanguage}{persian}

باری راجع بمطلب میشویم که سؤال از آن شده بود که سلطنت قائم با آنکه در
احادیث مأثوره از انجم مضیئه وارد شده با وجود این اثری از سلطنت ظاهر نشد
بلکه خلاف آن تحقّق یافت چنانچه اصحاب و اولیای او در دست ناس مبتلا و محصور
بوده و هستند و در نهایت ذلّت و عجز در ملک ظاهرند بلی سلطنتی که در کتب در
حقّ قائم مذکور است حقّ و لا ریب فیه ولیکن آن نه آن سلطنت و حکومتی است که
هر نفسی ادراک نماید و دیگر آنکه جمیع انبیای قبل که بشارت داده‌اند مردم را
بظهور بعد همئ آن مظاهر قبل ذکر سلطنت ظهور بعد را نموده چنانچه در کتب قبل
مسطور است و این تخصیص بقائم ندارد و در حقّ جمیع آن مظاهر قبل و بعد حکم
سلطنت و جمیع صفات و اسماء ثابت و محقّق است زیرا که مظاهر صفات غیبیّه و
مطالع اسرار الهیّه‌اند چنانچه مذکور شد

\end{otherlanguage}

Bārī, rājeʿ be maṭlab mī-shavīm ke soʾāl az ān shode būd ke salṭanat-e
Qāʾim, bā ānke dar aḥādīs̱-e maʾs̱ūre az anjom-e moḍīʾe vāred shode, bā
vojūd-e īn as̱arī az salṭanat ẓāher nashod balke khelāf-e ān taḥaqqoq
yāft, chenānche aṣḥāb va awlīyā-ye ū dar dast-e nās mobtalā va maḥṣūr
būde va hastand va dar nehāyat-e ḏellat va ʿajz dar molk ẓāherand. Balī,
salṭanatī ke dar kotob dar ḥaqq-e Qāʾim maḏkūr ast ḥaqq va lā rayba
fīhi, valīkan ān na ān salṭanat va ḥokūmatī ast ke har nafsī edrāk
namāyad. Va dīgar ānke jamīʿ-e anbiyā-ye qabl ke beshārat dāde-and
mardom rā be ẓohūr-e baʿd, hame-ye ān maẓāher-e qabl ḏekr-e salṭanat-e
ẓohūr-e baʿd rā namūde, chenānche dar kotob-e qabl masṭūr ast, va īn
takhṣīṣ be Qāʾim nadārad. Va dar ḥaqq-e jamīʿ-e ān maẓāher-e qabl va
baʿd ḥokm-e salṭanat va jamīʿ-e ṣefāt va asmā s̱ābet va moḥaqqaq ast,
zīrā ke maẓāher-e ṣefāt-e ghaybīye va maṭāleʿ-e asrār-e elāhīye-and,
chenānche maḏkūr shod.

Let us return, then, to the matter about which the question was asked:
although the sovereignty of the Qāʾim is mentioned in the traditions
handed down from the radiant stars, no trace of that sovereignty
appeared; instead, the opposite came to pass, for his companions and
devoted followers have been, and remain, afflicted and confined at the
hands of the people, and appear in the world in utter abasement and
helplessness. Indeed, the sovereignty mentioned in the books concerning
the Qāʾim is true, and there is no doubt about it. But it is not the
sovereignty and rule that everyone imagines. Furthermore, all the former
Prophets who gave the people glad tidings of a subsequent
Revelation---all those former Manifestations---mentioned the sovereignty
of the Revelation to come, as is recorded in the earlier books; this is
not peculiar to the Qāʾim. Sovereignty and all the attributes and names
are firmly and assuredly established for all those Manifestations, both
past and future, for they are the manifestations of the unseen
attributes and the dawning places of divine mysteries, as has been
mentioned.

\subsection{115}\label{section-305}

\begin{otherlanguage}{persian}

و دیگر آنکه مقصود از سلطنت احاطه و قدرت آن حضرتست بر همهٔ ممکنات خواه در
عالم ظاهر باستیلای ظاهری ظاهر شود و یا نشود و این بسته باراده و مشیّت خود
آن حضرتست ولیکن بر آن جناب معلوم بوده که سلطنت و غنا و حیات و موت و حشر
و نشر که در کتب قبل مذکور است مقصود این نیست که الیوم این مردم احصا و
ادراک مینمایند بلکه مراد از سلطنت سلطنتی است که در ایّام ظهور هر یک از
شموس حقیقت بنفسه لنفسه ظاهر میشود و آن احاطۀ باطنیّه است که بآن احاطه
مینماید کلّ من فی السّموات و الأرض را و بعد باستعداد کون و زمان و خلق در
عالم ظاهر بظهور میآید چنانچه سلطنت حضرت رسول حال در میان ناس ظاهر و
هویداست و در اوّل امر آن حضرت آن بود که شنیدید چه مقدار اهل کفر و ضلال که
علمای آن عصر و اصحاب ایشان باشند بر آن جوهر فطرت و ساذج طینت وارد آوردند
چه مقدار خاشاکها و خارها که بر محلّ عبور آن حضرت میریختند و این معلوم است
که آن اشخاص بظنون خبیثئ شیطانیّهٔ خود اذیّت بآن هیکل ازلی را سبب رستگاری
خود میدانستند زیرا که جمیع علمای عصر بمثل عبداللّه اُبَی و ابوعامر راهب و
کعب بن اشرف و نضر بن حارث جمیع آن حضرت را تکذیب نمودند و نسبت بجنون و
افترا دادند و نسبتهائی که نعوذ باللّه من ان یجری به المداد او یتحرّک علیه
القلم او یحمله الألواح بلی این نسبتها بود که سبب ایذای مردم نسبت بآن
حضرت شد و این معلوم و واضح است که علمای وقت اگر کسی را رد و طرد نمایند و
از اهل ایمان ندانند چه بر سر آن نفس میآید چنانچه بر سر این بنده آمد و
دیده شد

\end{otherlanguage}

Va dīgar ānke maqṣūd az salṭanat eḥāṭe va qodrat-e ān Ḥaḍrat ast bar
hame-ye momkenāt, khāh dar ʿālam-e ẓāher be estīlā-ye ẓāherī ẓāher
shavad va yā nashavad, va īn baste be erāde va mashīyyat-e khod-e ān
Ḥaḍrat ast. Valīkan bar ān janāb maʿlūm būde ke salṭanat va ghenā va
ḥayāt va mawt va ḥashr va nashr ke dar kotob-e qabl maḏkūr ast, maqṣūd
īn nīst ke al-yawm īn mardom eḥṣā va edrāk mī-namāyand. Balke morād az
salṭanat salṭanatī ast ke dar ayyām-e ẓohūr-e har yek az shomūs-e
ḥaqīqat bi-nafsihi li-nafsihi ẓāher mī-shavad, va ān eḥāṭe-ye bāṭenīye
ast ke be ān eḥāṭe mī-namāyad kull-e man fī al-samāwāt wa al-arḍ rā, va
baʿd be esteʿdād-e kawn va zamān va khalq dar ʿālam-e ẓāher be ẓohūr
mī-āyad. Chenānche salṭanat-e Ḥaḍrat-e Rasūl ḥāl dar mīyān-e nās ẓāher
va hovaydāst, va dar avval-e amr-e ān Ḥaḍrat ān būd ke shenīdīd. Che
meqdār ahl-e kofr va ḍalāl ke ʿolamā-ye ān ʿaṣr va aṣḥāb-e īshān bāshand
bar ān jawhar-e feṭrat va sādej-e ṭīnat vāred āvardand! Che meqdār
khāshākhā va khārhā ke bar maḥall-e ʿobūr-e ān Ḥaḍrat mī-rīkhtand! Va īn
maʿlūm ast ke ān ashkhāṣ be ẓonūn-e khabīs̱e-ye shayṭānīye-ye khod
aḏīyyat be ān haykal-e azalī rā sabab-e rastegārī-ye khod mī-dānestand,
zīrā ke jamīʿ-e ʿolamā-ye ʿaṣr be mes̱l-e ʿAbduʾllāh-e Ubayy va Abū
ʿĀmer-e Rāheb va Kaʿb ibn Ashraf va Naḍr ibn Ḥāres̱, jamīʿ ān Ḥaḍrat rā
takḏīb namūdand va nesbat be jonūn va efterā dādand, va nesbathāyī ke
naʿūdhu bi-Allāhi min an yajrīya bihi al-midādu aw yataḥarraka ʿalayhi
al-qalamu aw yaḥmilahu al-alwāḥ. Balī, īn nesbathā būd ke sabab-e īḏā-ye
mardom nesbat be ān Ḥaḍrat shod. Va īn maʿlūm va vāḍeḥ ast ke ʿolamā-ye
vaqt agar kasī rā radd va ṭard namāyand va az ahl-e īmān nadānand, che
bar sar-e ān nafs mī-āyad, chenānche bar sar-e īn bande āmad va dīde
shod.

Furthermore, sovereignty means his all-encompassing power over all
contingent things, whether or not it appears in the outward world as
outward ascendancy; this depends upon his own will and purpose. You
should know, however, that the sovereignty, wealth, life, death,
gathering, and resurrection mentioned in the earlier books do not mean
what these people today reckon and understand them to mean. Rather,
sovereignty means the sovereignty that appears, of itself and for
itself, in the days of the revelation of each of the Suns of Truth. It
is that inward dominion by which he encompasses all who are in the
heavens and on earth, and which later appears in the outward world
according to the capacity of existence, the time, and the people. Thus
the sovereignty of the Messenger is now manifest and evident among the
people, whereas at the beginning of his Cause his circumstances were as
you have heard. How much the people of unbelief and error---the
religious scholars of that age and their followers---inflicted upon that
Essence of innate nature and Embodiment of purity! How much rubbish and
how many thorns they scattered along his path! It is well known that
those people, in their wicked, satanic imaginings, regarded the torment
they inflicted upon that eternal Temple as a means of their own
salvation. For all the religious scholars of that age, such as
ʿAbduʾllāh Ubayy, Abū ʿĀmir the monk, Kaʿb ibn Ashraf, and Naḍr ibn
Ḥārith, all denounced him as false and accused him of madness and
fabrication, and made accusations from which we seek refuge in God lest
ink should flow with them, or the pen move to record them, or the
tablets bear them. Indeed, these accusations were what caused the people
to torment him. It is clear and evident what befalls a person when the
religious scholars of the time reject and cast him out and refuse to
count him among the faithful, as happened to this servant and was
witnessed.

\subsection{116}\label{section-306}

\begin{otherlanguage}{persian}

اینست که آن حضرت فرمود ما اوذی نبیّ بمثل ما اوذیت و در فرقان نسبتها که
دادند و اذیّتها که بآن حضرت نمودند همه مذکور است فارجعوا الیه لعلّکم
بمواقع الأمر تطّلعون حتّی قسمی بر آن حضرت سخت شد که احدی با آن حضرت و
اصحاب او چندی معاشرت نمینمود و هر نفسی که خدمت آن حضرت میرسید کمال اذیّت
را باو وارد مینمودند

\end{otherlanguage}

Īn ast ke ān Ḥaḍrat farmūd: mā ūdhiya nabiyyun bi-mithli mā ūdhītu. Va
dar Forqān nesbathā ke dādand va aḏīyyathā ke be ān Ḥaḍrat namūdand hame
maḏkūr ast. Fa-irjiʿū ilayhi laʿallakum bi-mawāqiʿi al-amri taṭṭaliʿūn.
Ḥattā qesmī bar ān Ḥaḍrat sakht shod ke aḥadī bā ān Ḥaḍrat va aṣḥāb-e ū
chandī moʿāsherat nemī-namūd, va har nafsī ke khedmat-e ān Ḥaḍrat
mī-resīd kamāl-e aḏīyyat rā be ū vāred mī-namūdand.

Thus he said, ``No Prophet has been afflicted as I have been
afflicted.'' All the accusations they made and the torments they
inflicted upon him are recorded in the Qurʾān. Refer to it, that you may
become acquainted with the circumstances of the matter. His situation
became so grievous that for a time no one would associate with him and
his companions, and anyone who entered his presence was subjected to the
utmost persecution.

\subsection{117}\label{section-307}

\begin{otherlanguage}{persian}

در این موقع یک آیه ذکر مینمایم که اگر چشم بصیرت باز کنی تا زنده هستی بر
مظلومی آن حضرت نوحه و ندبه نمائی و آن آیه در وقتی نازل شد که آن حضرت از
شدّت بلایا و اعراض ناس بغایت افسرده و دلتنگ بود جبرئیل از سدرة المنتهای
قرب نازل شد و این آیه تلاوت نمود و ان کان کبر علیک اعراضهم فان استطعت ان
تبتغی نفقاً فی الأرض او سلّماً فی السّمآء که ترجمهٔ آن اینست که اگر بزرگ است
بر تو اعراض معرضین و سخت است بر تو ادبار منافقین و ایذای ایشان پس اگر
مستطیعی و میتوانی طلب کن نقبی در زیر ارض یا نردبانی بسوی آسمان که تلویح
بیان اینست که چاره نیست و دست از تو بر نمیدارند مگر آنکه در زیر زمین
پنهان شوی و یا بآسمان فرار نمائی

\end{otherlanguage}

Dar īn mawqeʿ yek āye ḏekr mī-namāyam ke agar chashm-e baṣīrat bāz konī
tā zende hastī bar maẓlūmī-ye ān Ḥaḍrat nawḥe va nodbe namāyī. Va ān āye
dar vaqtī nāzel shod ke ān Ḥaḍrat az sheddat-e balāyā va eʿrāḍ-e nās be
ghāyat afsorde va del-tang būd. Jebraʾīl az Sedrat al-Montahā-ye qorb
nāzel shod va īn āye telāvat namūd: wa in kāna kabura ʿalayka iʿrāḍuhum
fa-in istaṭaʿta an tabtaghiya nafaqan fī al-arḍi aw sullaman fī
al-samāʾ. Ke tarjome-ye ān īn ast ke agar bozorg ast bar to eʿrāḍ-e
moʿreḍīn va sakht ast bar to edbār-e monāfeqīn va īḏā-ye īshān, pas agar
mostaṭīʿī va mī-tavānī ṭalab kon naqbī dar zīr-e arḍ yā nardebānī be
sū-ye āsmān. Ke talvīḥ-e bayān īn ast ke chāre nīst va dast az to bar
nemī-dārand magar ānke dar zīr-e zamīn penhān shavī va yā be āsmān farār
namāyī.

Here I will mention a verse which, if you open the eye of insight, will
cause you to mourn and lament the wrongs he suffered for as long as you
live. This verse was revealed at a time when he was profoundly grieved
and distressed by the severity of his afflictions and the people's
rejection. Gabriel descended from the Lote Tree of the Utmost Limit of
nearness and recited this verse: ``And if their rejection weighs heavily
upon you, then, if you can, seek a tunnel in the earth or a ladder into
heaven.'' Its translation is this: If the rejection of those who turn
away weighs heavily upon you, and the aversion of the hypocrites and
their persecution are hard for you to bear, then, if you are able and
have the power, seek a tunnel beneath the earth or a ladder toward
heaven. This utterance implies that there is no remedy, and they will
not leave you alone unless you hide beneath the earth or flee to heaven.

\subsection{118}\label{section-308}

\begin{otherlanguage}{persian}

و حال امروز مشاهده نما که چه ‌قدر از سلاطین باسم آن حضرت تعظیم مینمایند و
چه‌ قدر از بلاد و اهل آن که در ظلّ او ساکنند و بنسبت بآن حضرت افتخار دارند
چنانچه بر منابر و گلدسته‌ها این اسم مبارک را بکمال تعظیم و تکریم ذکر
مینمایند و سلاطینی هم که در ظلّ آن حضرت داخل نشده‌اند و قمیص کفر را تجدید
ننموده‌اند ایشان هم ببزرگی و عظمت آن شمس عنایت مقرّ و معترفند اینست سلطنت
ظاهره که مشاهده میکنی و این لابدّ است از برای جمیع انبیا که یا در حیات و
یا بعد از عروج ایشان بموطن حقیقی ظاهر و ثابت میشود چنانچه الیوم ملاحظه
میگردد ولیکن آن سلطنت که مقصود است لم‌یزل و لا‌یزال طائف حول ایشانست و
همیشه با ایشان است و آنی انفکاک نیابد و آن سلطنت باطنیّه است که احاطه
نموده کلّ من فی السّموات و الأرض را

\end{otherlanguage}

Va ḥāl emrūz moshāhede namā ke che qadr az salāṭīn be esm-e ān Ḥaḍrat
taʿẓīm mī-namāyand va che qadr az belād va ahl-e ān ke dar ẓell-e ū
sākenand va be nesbat be ān Ḥaḍrat eftekhār dārand, chenānche bar
manāber va goldastehā īn esm-e mobārak rā be kamāl-e taʿẓīm va takrīm
ḏekr mī-namāyand. Va salāṭīnī ham ke dar ẓell-e ān Ḥaḍrat dākhel
nashode-and va qamīṣ-e kofr rā tajdīd nanamūde-and, īshān ham be bozorgī
va ʿaẓamat-e ān shams-e ʿenāyat moqerr va moʿtarefand. Īn ast salṭanat-e
ẓāhere ke moshāhede mī-konī, va īn lā-bodd ast az barā-ye jamīʿ-e anbiyā
ke yā dar ḥayāt va yā baʿd az ʿorūj-e īshān be mawṭen-e ḥaqīqī ẓāher va
s̱ābet mī-shavad, chenānche al-yawm molāḥeẓe mī-gardad. Valīkan ān
salṭanat ke maqṣūd ast lam-yazal va lā-yazāl ṭāʾef-e ḥawl-e īshān ast va
hamīshe bā īshān ast va ānī enfekāk nayābad, va ān salṭanat-e bāṭenīye
ast ke eḥāṭe namūde kull-e man fī al-samāwāt wa al-arḍ rā.

Now observe today how many sovereigns bow in reverence to his name, and
how many cities and their inhabitants dwell beneath his shadow and take
pride in their association with him. From pulpits and minarets they
mention this blessed name with the utmost reverence and honor. Even
those sovereigns who have not entered beneath his shadow and have not
exchanged their garment of unbelief acknowledge and confess the
greatness and majesty of that Sun of loving-kindness. This is the
outward sovereignty that you behold, and it must inevitably appear and
become established for all the Prophets, either during their lifetime or
after their ascent to their true homeland, as can be observed today. Yet
the sovereignty intended here has always revolved, and will always
revolve, around them. It is always with them and is never separated from
them for an instant. It is that inward sovereignty which encompasses all
who are in the heavens and on earth.

\subsection{119}\label{section-309}

\begin{otherlanguage}{persian}

و از جملهٔ سلطنت آنست که از آن شمس احدیّه ظاهر شد آیا نشنیدی که بیک آیه
چگونه میانهٔ نور و ظلمت و سعید و شقی و مؤمن و کافر فصل فرمود و جمیع
اشارات و دلالات قیامت که شنیدی از حشر و نشر و حساب و کتاب و غیره کل
بتنزیل همان یک آیه هویدا شد و بعرصهٔ شهود آمد و همچنین آن آیهٔ منزله رحمت
بود برای ابرار یعنی انفسی که در حین استماع گفتند ربّنا سمعنا و اطعنا و
نقمت شد برای فجّار یعنی آنهائی که بعد از استماع گفتند سمعنا و عصینا و سیف
اللّه بود برای فصل مؤمن از کافر و پدر از پسر چنانچه دیده‌اید آنهائی که
اقرار نمودند با آنهائی که انکار نمودند در صدد جان و مال هم برآمدند چه
پدرها که از پسرها اعراض نمودند و چه عاشقها که از معشوقها احتراز جستند و
چنان حادّ و برنده بود این سیف بدیع که همهٔ نسبتها را از هم قطع نمود و از
یک جهت ملاحظه فرمائید چگونه وصل نمود مثل آنکه ملاحظه شد که جمعی از ناس
که سالها شیطان نفس تخم کینه و عدوان مابین ایشان کاشته بود بسبب ایمان
باین امر بدیع منیع چنان متّحد و موافق شدند که گویا از یک صلب ظاهر شده‌اند
کذلک یؤلّف اللّه بین قلوب الّذینهم انقطعوا الیه و آمنوا بآیاته و کانوا من
کوثر الفضل بأیادی العزّ من الشّاربین و دیگر آنکه چه‌ قدر از مردم مختلف
العقاید و مختلف المذهب و مختلف المزاج که از این نسیم رضوان الهی و
بهارستان قدس معنوی قمیص جدید توحید پوشیدند و از کأس تفرید نوشیدند

\end{otherlanguage}

Va az jomle-ye salṭanat ān ast ke az ān shams-e aḥadīye ẓāher shod. Āyā
nashenīdī ke be yek āye chegūne mīyāne-ye nūr va ẓolmat va saʿīd va
shaqī va moʾmen va kāfer faṣl farmūd? Va jamīʿ-e eshārāt va dalālāt-e
qīyāmat ke shenīdī az ḥashr va nashr va ḥesāb va ketāb va ghayre, koll
be tanzīl-e hamān yek āye hovaydā shod va be ʿarṣe-ye shohūd āmad. Va
hamchenīn ān āye-ye monazzale raḥmat būd barā-ye abrār, yaʿnī anfosī ke
dar ḥīn-e estemāʿ goftand: rabbanā samiʿnā wa aṭaʿnā; va neqmat shod
barā-ye fojjār, yaʿnī ānhāyī ke baʿd az estemāʿ goftand: samiʿnā wa
ʿaṣaynā. Va sayf-e Allāh būd barā-ye faṣl-e moʾmen az kāfer va pedar az
pesar, chenānche dīde-īd ānhāyī ke eqrār namūdand bā ānhāyī ke enkār
namūdand dar ṣadad-e jān va māl-e ham bar-āmadand. Che pedarhā ke az
pesarhā eʿrāḍ namūdand va che ʿāsheqhā ke az maʿshūqhā eḥterāz jostand!
Va chenān ḥādd va borrande būd īn sayf-e badīʿ ke hame-ye nesbathā rā az
ham qaṭʿ namūd. Va az yek jehat molāḥeẓe farmāyīd chegūne vaṣl namūd.
Mes̱l-e ānke molāḥeẓe shod ke jamʿī az nās ke sālhā shayṭān-e nafs
tokhm-e kīne va ʿodvān mābayn-e īshān kāshte būd, be sabab-e īmān be īn
amr-e badīʿ-e manīʿ chenān mottaḥed va movāfeq shodand ke gūyā az yek
ṣolb ẓāher shode-and. Kadhālika yuʾallifu Allāhu bayna qulūbi alladhīna
hum inqaṭaʿū ilayhi wa āmanū bi-āyātihi wa kānū min kawthari al-faḍli
bi-ayādī al-ʿizzi mina al-shāribīn. Va dīgar ānke che qadr az mardom-e
mokhtalef al-ʿaqāyed va mokhtalef al-madhhab va mokhtalef al-mezāj ke az
īn nasīm-e Reḍvān-e elāhī va bahārestān-e qods-e maʿnavī qamīṣ-e jadīd-e
tawḥīd pūshīdand va az kaʾs-e tafrīd nūshīdand.

Among the instances of sovereignty is that which appeared from that Sun
of Oneness. Have you not heard how, through a single verse, he
distinguished light from darkness, the blessed from the wretched, and
the believer from the unbeliever? All the signs and tokens of the
Resurrection of which you have heard---the gathering, the rising of the
dead, the reckoning, the record, and the rest---became evident and
entered the realm of visible existence through the revelation of that
one verse. Likewise, that revealed verse was mercy for the righteous,
meaning those souls who, upon hearing it, said, ``Our Lord, we have
heard and obeyed''; and it became retribution for the wicked, meaning
those who, after hearing it, said, ``We have heard and disobeyed.'' It
was the sword of God, separating believer from unbeliever and father
from son. As you have seen, those who acknowledged it and those who
denied it sought one another's lives and possessions. How many fathers
turned away from their sons, and how many lovers shunned their beloveds!
So keen and cutting was this wondrous sword that it severed every bond.
Yet consider, from another perspective, how it united people. For
example, it was observed that people among whom the satan of self had
for years sown the seeds of malice and enmity became, through faith in
this wondrous, exalted Cause, so united and harmonious that they seemed
to have sprung from a single father. Thus God unites the hearts of those
who have detached themselves for His sake, believed in His verses, and
drunk from the river of grace at the hands of glory. Furthermore, how
many people of differing beliefs, differing religions, and differing
temperaments have, through this breeze from the divine Paradise and this
springtime garden of spiritual holiness, put on the new garment of unity
and drunk from the cup of singleness!

\subsection{120}\label{section-310}

\begin{otherlanguage}{persian}

اینست معنی حدیث مشهور که فرموده گرگ و میش از یک محل میخورند و میآشامند و
حال نظر بعدم معرفت این جهّال فرمائید بمثل امم سابقه هنوز منتظرند که کی
این حیوانات بر یک خوان مجتمع میشوند اینست رتبهٔ ناس گویا هرگز از جام
انصاف ننوشیده‌اند و هرگز در سبیل عدل قدم نگذاشته‌اند از همه گذشته این امر
وقوعش چه حسنی در عالم احداث مینماید فنعم ما نزّل فی شأنهم لهم قلوب لا
یفقهون بها و لهم اعین لا یبصرون بها

\end{otherlanguage}

Īn ast maʿnī-ye ḥadīs̱-e mashhūr ke farmūde gorg va mīsh az yek maḥall
mī-khorand va mī-āshāmand. Va ḥāl naẓar be ʿadam-e maʿrefat-e īn johhāl
farmāyīd: be mes̱l-e omam-e sābeqe hanūz montaẓerand ke key īn ḥayvānāt
bar yek khān mojtameʿ mī-shavand. Īn ast rotbe-ye nās. Gūyā hargez az
jām-e enṣāf nanūshīde-and va hargez dar sabīl-e ʿadl qadam
nagozāshte-and. Az hame goḏashte, īn amr voqūʿash che ḥosnī dar ʿālam
eḥdās̱ mī-namāyad? Fa-niʿma mā nuzzila fī shaʾnihim: lahum qulūbun lā
yafqahūna bihā wa lahum aʿyunun lā yubṣirūna bihā.

This is the meaning of the well-known tradition which states that the
wolf and the lamb will eat and drink in the same place. Now consider the
lack of understanding of these ignorant people: like the peoples of the
past, they are still waiting for the time when these animals will gather
at one table. Such is the people's understanding! It is as though they
had never drunk from the cup of fairness and never set foot on the path
of justice. Besides all this, what good would the occurrence of such a
thing bring about in the world? How fitting is what has been revealed
concerning them: ``They have hearts with which they do not understand,
and eyes with which they do not see.''

\subsection{121}\label{section-311}

\begin{otherlanguage}{persian}

و دیگر آنکه ملاحظه فرمائید بتنزیل همین یک آیهٔ منزلهٔ از سماء مشیّت چگونه
حساب خلایق کشیده شد که هر کس اقرار نمود و اقبال جست حسنات او بر سیّئات
زیادتی نمود و جمیع خطایای او معفوّ شد و مغفور آمد کذلک یصدق فی شأنه بأنّه
سریع الحساب و کذلک یبدّل اللّه السّیّئات بالحسنات لو انتم فی آفاق العلم و
انفس الحکمة تتفرّسون و همچنین هر کس از جام حبّ نصیب برداشت از بحر فیوضات
سرمدیّه و غمام رحمت ابدیّه حیات باقیهٔ ابدیّۀ ایمانیّه یافت و هر نفسی که قبول
ننمود بموت دائمی مبتلا شد و مقصود از موت و حیات که در کتب مذکور است موت
و حیات ایمانیست و از عدم ادراک این معنی است که عامّهٔ ناس در هر ظهور
اعتراض نمودند و بشمس هدایت مهتدی نشدند و جمال ازلی را مقتدی نگشتند

\end{otherlanguage}

Va dīgar ānke molāḥeẓe farmāyīd be tanzīl-e hamīn yek āye-ye monazzale
az samā-ye mashīyyat chegūne ḥesāb-e khalāyeq keshīde shod, ke har kas
eqrār namūd va eqbāl jost ḥasanāt-e ū bar sayyeʾāt zīyādatī namūd va
jamīʿ-e khaṭāyā-ye ū maʿfovv shod va maghfūr āmad. Kadhālika yaṣduqu fī
shaʾnihi bi-annahu sarīʿu al-ḥisāb, wa kadhālika yubaddilu Allāhu
al-sayyiʾāti bi-al-ḥasanāti law antum fī āfāqi al-ʿilmi wa anfusi
al-ḥikmati tatafarrasūn. Va hamchenīn har kas az jām-e ḥobb naṣīb
bardāsht az baḥr-e foyūḍāt-e sarmadīye va ghamām-e raḥmat-e abadīye
ḥayāt-e bāqīye-ye abadīye-ye īmānīye yāft, va har nafsī ke qabūl nanamūd
be mawt-e dāʾemī mobtalā shod. Va maqṣūd az mawt va ḥayāt ke dar kotob
maḏkūr ast mawt va ḥayāt-e īmānī-st. Va az ʿadam-e edrāk-e īn maʿnī ast
ke ʿāmme-ye nās dar har ẓohūr eʿterāḍ namūdand va be shams-e hedāyat
mohtadī nashodand va jamāl-e azalī rā moqtadī nagashtand.

Furthermore, consider how, through the revelation of this single verse
sent down from the heaven of the divine Will, the reckoning of all
people was carried out: whoever acknowledged it and turned toward it
found that his good deeds outweighed his evil deeds, and all his
transgressions were pardoned and forgiven. Thus the statement that He is
swift in reckoning is true of Him; and thus God changes evil deeds into
good deeds, if you discern the horizons of knowledge and the inner
realities of wisdom. Likewise, whoever partook of the cup of love
received, from the ocean of everlasting outpourings and the cloud of
eternal mercy, the enduring, everlasting life of faith; and every soul
who refused to accept was afflicted with perpetual death. The death and
life mentioned in the books mean the death and life of faith. It is
because they have failed to grasp this meaning that the generality of
people have objected in every Revelation, have not been guided by the
Sun of guidance, and have not taken the eternal Beauty as their guide.

\subsection{122}\label{section-312}

\begin{otherlanguage}{persian}

چنانچه وقتی که سراج محمّدی در مشکاة احمدیّه مشتعل شد بر مردم حکم بعث و حشر
و حیات و موت فرمود این بود که اعلام مخالفت مرتفع شد و ابواب استهزاء
مفتوح گشت چنانچه از زبان مشرکین روح الأمین خبر داده و لئن قلت انّکم
مبعوثون من بعد الموت لیقولنّ الّذین کفروا ان هذا الّا سحر مبین مضمون آن
اینست که اگر بگوئی باین مشرکین که شما مبعوث شده‌اید بعد از مردن هرآینه
میگویند آنهائی که کافر شده‌اند بخدا و آیات او نیست این مگر سحری ظاهر و
آشکار و هویدا و در جای دیگر میفرماید و ان تعجب فعجب قولهم أ ئذا کنّا
تراباً أ ئنّا لفی خلق جدید که ترجمهٔ آن اینست میفرماید اگر عجب میداری پس
عجبست قول کافران و معرضان که میگویند آیا ما تراب بودیم و از روی استهزاء
میگفتند که آیا مائیم مبعوث‌شدگان اینست که در مقام دیگر قهراً لهم میفرماید
أ فعیینا بالخلق الأوّل بل هم فی لبس من خلق جدید مضمون آن این است که آیا
ما عاجز و مانده شدیم از خلق اوّل بلکه این مشرکین در شک و شبهه هستند از
خلق جدید

\end{otherlanguage}

Chunānche vaqtī ke sirāj-i Muḥammadī dar mishkāt-i Aḥmadiyye mushtaʿil
shod, bar mardom ḥukm-i baʿth va ḥashr va ḥayāt va mawt farmūd. Īn būd
ke aʿlām-i mukhālifat murtafiʿ shod va abvāb-i istihzā maftūḥ gasht.
Chunānche az zabān-i mushrikīn Rūḥ al-Amīn khabar dāde: Wa-laʾin qulta
innakum mabʿūthūna min baʿdi al-mawti la-yaqūlanna alladhīna kafarū in
hādhā illā siḥrun mubīn. Maḍmūn-i ān īn ast ke agar begūʾī be-īn
mushrikīn ke shomā mabʿūth shode-īd baʿd az mordan, har-āyine mīgūyand
ānhāʾī ke kāfir shode-and be-Khudā va āyāt-i ū: Nīst īn magar siḥrī
ẓāhir va āshkār va huvaydā. Va dar jā-yi dīgar mīfarmāyad: Wa-in taʿjab
fa-ʿajabun qawluhum a-idhā kunnā turāban a-innā la-fī khalqin jadīd. Ke
tarjume-yi ān īn ast, mīfarmāyad: Agar ʿajab mīdārī pas ʿajab ast qawl-i
kāfirān va muʿriḍān ke mīgūyand āyā mā turāb būdīm, va az rū-yi istihzā
mīguftand ke āyā māʾīm mabʿūth-shodegān. Īn ast ke dar maqām-i dīgar
qahran lahum mīfarmāyad: A-fa-ʿayīnā bi-al-khalqi al-awwali bal hum fī
labsin min khalqin jadīd. Maḍmūn-i ān īn ast ke āyā mā ʿājiz va mānde
shodīm az khalq-i avval? Balke īn mushrikīn dar shakk va shubhe hastand
az khalq-i jadīd.

Thus, when the lamp of Muḥammad was kindled in the niche of Aḥmad, He
pronounced upon the people the judgment of resurrection, gathering,
life, and death. Then the banners of opposition were raised and the
gates of mockery were opened. As the Faithful Spirit has reported the
words of the idolaters: ``If you say, `You will be raised after death,'
those who disbelieve will surely say, `This is nothing but manifest
sorcery.'\,'' Its meaning is this: if you tell these idolaters that they
have been raised after dying, those who have disbelieved in God and His
verses will assuredly say, ``This is nothing but plain, evident, and
manifest sorcery.'' Elsewhere He says: ``If you wonder, wondrous indeed
is their saying, `When we have become dust, shall we truly be in a new
creation?'\,'' Its translation is this: He says, ``If you wonder, then
wondrous is the saying of the unbelievers and those who turn away, who
say, `Were we dust?'\,'' And in mockery they would say, ``Are we the
ones who have been raised?'' Thus, elsewhere, to confound them, He says:
``Were We wearied by the first creation? Rather, they are in doubt about
a new creation.'' Its meaning is this: have We become powerless and
exhausted by the first creation? Rather, these idolaters are in doubt
and uncertainty about the new creation.

\subsection{123}\label{section-313}

\begin{otherlanguage}{persian}

و علمای تفسیر و اهل ظاهر چون معانی کلمات الهیّه را ادراک ننمودند و از
مقصود اصلی محتجب ماندند لهذا بقاعدۀ نحو استدلال نمودند اذا که بر سر ماضی
درآید معنی مستقبل افاده میشود و بعد در کلماتی که کلمهٔ اذا نازل نگشته
متحیّر ماندند مثل اینکه میفرماید و نفخ فی الصّور ذلک یوم الوعید و جآءت کلّ
نفس معها سائق و شهید که معنی ظاهر آن این است دمیده شد در صور و آنست یوم
وعید که بنظرها بسیار بعید بود و آمد هر نفسی برای حساب و با اوست راننده و
گواه و در مثل این مواقع یا کلمهٔ اذا را مقدّر گرفتند و یا مستدلّ شدند بر
اینکه چون قیامت محقّق الوقوعست لهذا بفعل ماضی ادا شد که گویا گذشته است
ملاحظه فرمائید چه ‌قدر بی ‌ادراک و تمیزند نفخهٔ محمّدیّه را که باین صریحی
میفرماید ادراک نمیکنند و از افاضهٔ این نقرهٔ الهی خود را محروم مینمایند و
منتظر صور اسرافیل که یکی از عباد اوست میشوند با اینکه تحقّق وجود اسرافیل
و امثال او ببیان خود آن حضرت شده قل أ تستبدلون الّذی هو خیر لکم فبئس ما
استبدلتم بغیر حقّ و کنتم قوم سوء اخسرین

\end{otherlanguage}

Va ʿulamā-yi tafsīr va ahl-i ẓāhir chūn maʿānī-yi kalimāt-i ilāhiyye rā
idrāk nanamūdand va az maqṣūd-i aṣlī muḥtajib māndand, li-hādhā
be-qāʿide-yi naḥv istidlāl namūdand: idhā ke bar sar-i māḍī dar-āyad
maʿnī-yi mustaqbal ifāde mīshavad. Va baʿd dar kalimātī ke kalime-yi
idhā nāzil nagashte mutaḥayyir māndand. Mithl-i īnke mīfarmāyad:
Wa-nufikha fī al-ṣūri dhālika yawmu al-waʿīdi wa-jāʾat kullu nafsin
maʿahā sāʾiqun wa-shahīd. Ke maʿnī-yi ẓāhir-i ān īn ast: Damīde shod dar
ṣūr va ān ast yawm-i vaʿīd ke be-naẓarhā besyār baʿīd būd, va āmad har
nafsī barā-yi ḥisāb va bā ūst rānande va guvāh. Va dar mithl-i īn
mavāqiʿ yā kalime-yi idhā rā muqaddar giriftand va yā mustadill shodand
bar īnke chūn qiyāmat muḥaqqaq al-vuqūʿ ast li-hādhā be-fiʿl-i māḍī adā
shod ke gūyā gudhashte ast. Mulāḥiẓe farmāʾīd che qadr bī-idrāk va
tamīzand. Nafkhe-yi Muḥammadiyye rā ke be-īn ṣarīḥī mīfarmāyad idrāk
namīkunand va az ifāḍe-yi īn naqre-yi ilāhī khud rā maḥrūm mīnamāyand va
muntaẓir-i ṣūr-i Isrāfīl ke yakī az ʿibād-i ūst mīshavand, bā īnke
taḥaqquq-i vujūd-i Isrāfīl va amthāl-i ū be-bayān-i khud-i ān Ḥaḍrat
shode. Qul a-tastabdilūna alladhī huwa khayrun lakum fa-biʾsa mā
istabdaltum bi-ghayri ḥaqqin wa-kuntum qawma sawʾin akhsarīn.

Because the interpreters of scripture and those who cling to outward
meanings did not understand the meanings of the divine words and
remained veiled from their essential purpose, they resorted to a rule of
grammar: when ``idhā'' precedes a past-tense verb, it conveys a future
meaning. But they were then bewildered by passages in which the word
``idhā'' had not been revealed, such as His words: ``The trumpet was
sounded: that is the Day of the Threat. And every soul came, with it a
driver and a witness.'' Its outward meaning is this: the trumpet was
sounded, and that is the Day of the Threat, which seemed so remote to
people's eyes; and every soul came for the reckoning, accompanied by a
driver and a witness. In such instances, they either assumed the word
``idhā'' to be implicit or argued that, since the Resurrection was
certain to occur, it was expressed in the past tense as though it had
already happened. Consider how devoid they are of understanding and
discernment. They do not perceive the blast of Muḥammad, which He speaks
of so explicitly, and deprive themselves of the outpouring of this
divine trumpet call. They await the trumpet of Isrāfīl, who is but one
of His servants, even though the very existence of Isrāfīl and others
like him has been established through that exalted One's own words. Say:
``Would you exchange that which is better for you? Wretched indeed is
the exchange you have made without right, and you are an evil people,
the greatest losers.''

\subsection{124}\label{section-314}

\begin{otherlanguage}{persian}

بلکه مقصود از صور صور محمّدیست که بر همۀ ممکنات دمیده شد و قیامت قیام آن
حضرت بود بر امر الهی و غافلین که در قبور اجساد مرده بودند همه را بخلعت
جدیدهٔ ایمانیّه مخلّع فرمود و بحیات تازهٔ بدیعه زنده نمود اینست وقتی که آن
جمال احدیّه اراده فرمود که رمزی از اسرار بعث و حشر و جنّت و نار و قیامت
اظهار فرماید جبرئیل وحی این آیه آورد فسینغضون الیک رؤوسهم و یقولون متی
هو قل عسی ان یکون قریباً یعنی زود است این گمراهان وادی ضلالت سرهای خود را
از روی استهزاء حرکت میدهند و میگویند چه زمان خواهد این امور ظاهر شد تو
در جواب بگو که شاید اینکه نزدیک باشد تلویح همین یک آیه مردم را کافیست
اگر بنظر دقیق ملاحظه نمایند

\end{otherlanguage}

Balke maqṣūd az ṣūr ṣūr-i Muḥammadī-st ke bar hame-yi mumkināt damīde
shod, va qiyāmat qiyām-i ān Ḥaḍrat būd bar amr-i ilāhī, va ghāfilīn ke
dar qubūr-i ajsād murde būdand hame rā be-khilʿat-i jadīde-yi īmāniyye
mukhallaʿ farmūd va be-ḥayāt-i tāze-yi badīʿe zende namūd. Īn ast vaqtī
ke ān Jamāl-i Aḥadiyye irāde farmūd ke ramzī az asrār-i baʿth va ḥashr
va jannat va nār va qiyāmat iẓhār farmāyad, Jabraʾīl vaḥy-i īn āye
āvard: Fa-sa-yunghiḍūna ilayka ruʾūsahum wa-yaqūlūna matā huwa qul ʿasā
an yakūna qarīban. Yaʿnī zūd ast īn gumrāhān-i vādī-yi ḍalālat sarhā-yi
khud rā az rū-yi istihzā ḥarakat mīdihand va mīgūyand che zamān khvāhad
īn umūr ẓāhir shod. Tu dar javāb begū ke shāyad īnke nazdīk bāshad.
Talvīḥ-i hamīn yak āye mardom rā kāfī-st agar be-naẓar-i daqīq mulāḥiẓe
namāyand.

Rather, the trumpet means the trumpet of Muḥammad, which was sounded
over all contingent beings, and the Resurrection was His arising for the
Cause of God. All the heedless, who lay dead in the graves of their
bodies, He clothed in the new robe of faith and quickened with a fresh
and wondrous life. Thus, when that Beauty of Oneness wished to disclose
a hint of the mysteries of resurrection, gathering, paradise, fire, and
the Day of Resurrection, Gabriel brought the revelation of this verse:
``They will shake their heads at you and say, `When will it be?' Say:
`Perhaps it is near.'\,'' That is, these wanderers in the valley of
error will soon shake their heads in mockery and say, ``When will these
things appear?'' In reply, say, ``Perhaps it is near.'' The allusion in
this single verse is sufficient for the people, if they examine it with
close attention.

\subsection{125}\label{section-315}

\begin{otherlanguage}{arabic}

سبحان‌اللّه چه‌ قدر آن قوم از سبل حقّ دور بودند با اینکه قیامت بقیام آن حضرت
قائم بود و علامات و انوار او همۀ ارض را احاطه نموده بود مع‌ذلک سخریّه
مینمودند و معتکف بودند بتماثیلی که علمای عصر بافکار عاطل باطل جسته‌اند و
از شمس عنایت ربّانیّه و امطار رحمت سبحانیّه غافل گشته‌اند بلی جعل از روائح
قدس ازل محرومست و خفّاش از تجلّی آفتاب جهانتاب در گریز

\end{otherlanguage}

Subḥān Allāh, che qadr ān qawm az subul-i Ḥaqq dūr būdand! Bā īnke
qiyāmat be-qiyām-i ān Ḥaḍrat qāʾim būd va ʿalāmāt va anvār-i ū hame-yi
arḍ rā iḥāṭe namūde būd, maʿa-dhālik sukhriyye mīnamūdand va muʿtakif
būdand be-tamāthīlī ke ʿulamā-yi ʿaṣr bā-afkār-i ʿāṭil-i bāṭil juste-and
va az shams-i ʿināyat-i rabbāniyye va amṭār-i raḥmat-i subḥāniyye ghāfil
gashte-and. Balī, joʿal az ravāʾiḥ-i quds-i azal maḥrūm ast va khuffāsh
az tajallī-yi āftāb-i jahāntāb dar gurīz.

Glory be to God! How far those people were from the paths of the True
One! Although the Resurrection had come about through His arising, and
His signs and lights had encompassed the whole earth, they nevertheless
mocked and remained devoted to images that the scholars of the age had
sought with their idle and vain imaginings. They had become heedless of
the sun of divine favor and the showers of heavenly mercy. Indeed, the
dung beetle is deprived of the fragrances of eternal holiness, and the
bat flees from the radiance of the world-illuminating sun.

\subsection{126}\label{section-316}

\begin{otherlanguage}{persian}

و این مطلب در همهٔ اعصار در حین ظهور مظاهر حقّ بوده چنانچه عیسی میفرماید
لا بدّ لکم بأن تولدوا مرّة اخری و در مقام دیگر میفرماید من لم یولد من
المآء و الرّوح لا یقدر ان یدخل ملکوت اللّه المولود من الجسد جسد هو و
المولود من الرّوح هو روح که ترجمهٔ آن اینست نفسی که زنده نشده است از ماء
معرفت الهی و روح قدسی عیسوی قابل ورود و دخول در ملکوت ربّانی نیست زیرا هر
چه از جسد ظاهر شد و تولّد یافت پس اوست جسد و متولّد شده از روح که نفس
عیسوی باشد پس اوست روح خلاصهٔ معنی آنکه هر عبادی که از روح و نفخهٔ مظاهر
قدسیّه در هر ظهور متولّد و زنده شدند بر آنها حکم حیات و بعث و ورود در جنّت
محبّت الهیّه میشود و من‌دون آن حکم غیر آن که موت و غفلت و ورود در نار کفر و
غضب الهی است میشود و در جمیع کتب و الواح و صحائف مردمی که از جامهای لطیف
معارف نچشیده‌اند و بفیض روح القدس وقت قلوب ایشان فائز نشده بر آنها حکم
موت و نار و عدم بصر و قلب و سمع شده چنانچه از قبل ذکر شده لهم قلوب لا
یفقهون بها

\end{otherlanguage}

Va īn maṭlab dar hame-yi aʿṣār dar ḥīn-i ẓuhūr-i maẓāhir-i Ḥaqq būde.
Chunānche ʿĪsā mīfarmāyad: Lā budda lakum bi-an tūladū marratan ukhrā.
Va dar maqām-i dīgar mīfarmāyad: Man lam yūlad min al-māʾi wa-al-rūḥi lā
yaqdiru an yadkhula malakūta Allāh. Al-mawlūdu min al-jasadi jasadun
huwa wa-al-mawlūdu min al-rūḥi huwa rūḥ. Ke tarjume-yi ān īn ast: Nafsī
ke zende nashode ast az māʾ-i maʿrifat-i ilāhī va rūḥ-i qudsī-yi ʿĪsavī
qābil-i vurūd va dukhūl dar malakūt-i rabbānī nīst, zīrā har che az
jasad ẓāhir shod va tavallud yāft pas ūst jasad, va mutavallid shode az
rūḥ ke nafas-i ʿĪsavī bāshad pas ūst rūḥ. Khulāṣe-yi maʿnī ānke har
ʿibādī ke az rūḥ va nafkhe-yi maẓāhir-i qudsiyye dar har ẓuhūr
mutavallid va zende shodand bar ānhā ḥukm-i ḥayāt va baʿth va vurūd dar
jannat-i maḥabbat-i ilāhiyye mīshavad, va min-dūn-i ān ḥukm-i ghayr-i ān
ke mawt va ghaflat va vurūd dar nār-i kufr va ghaḍab-i ilāhī ast
mīshavad. Va dar jamīʿ-i kutub va alvāḥ va ṣaḥāʾif mardumī ke az
jāmhā-yi laṭīf-i maʿārif nachashīde-and va be-fayḍ-i Rūḥ al-Qudus-i vaqt
qulūb-i īshān fāʾiz nashode bar ānhā ḥukm-i mawt va nār va ʿadam-i baṣar
va qalb va samʿ shode. Chunānche az qabl dhikr shode: Lahum qulūbun lā
yafqahūna bihā.

This has been the case in every age at the appearance of the
Manifestations of the True One. Thus Jesus says: ``You must be born
again.'' Elsewhere He says: ``Whoever is not born of water and the
Spirit cannot enter the Kingdom of God. What is born of the flesh is
flesh, and what is born of the Spirit is spirit.'' Its translation is
this: a soul that has not been quickened by the water of divine
knowledge and the holy spirit of Jesus is not fit to enter and gain
admission to the divine Kingdom, for whatever appears and is born of the
flesh is flesh, and whatever is born of the spirit, which is the breath
of Jesus, is spirit. In brief, the meaning is that all servants who, in
each Revelation, are born and quickened through the spirit and breath of
the holy Manifestations are pronounced to have attained life,
resurrection, and entrance into the paradise of divine love. Those apart
from them receive the opposite judgment: death, heedlessness, and
entrance into the fire of unbelief and divine wrath. In all the books,
tablets, and scriptures, those who have not tasted the choice cups of
knowledge, and whose hearts have not received the grace of the Holy
Spirit of the age, have been condemned to death and fire and pronounced
to be without sight, heart, or hearing. As has already been mentioned:
``They have hearts with which they do not understand.''

\subsection{127}\label{section-317}

\begin{otherlanguage}{persian}

و در مقام دیگر در انجیل مسطور است که روزی یکی از اصحاب عیسی والدش وفات
نمود و او خدمت حضرت معروض داشت و اجازه خواست که برود و او را دفن و کفن
نموده راجع شود آن جوهر انقطاع فرمود دع الموتی لیدفنوه الموتی یعنی واگذار
مرده‌ها را تا دفن کنند مرده‌ها

\end{otherlanguage}

Va dar maqām-i dīgar dar Injīl masṭūr ast ke rūzī yakī az aṣḥāb-i ʿĪsā
vālidash vafāt namūd va ū khidmat-i Ḥaḍrat maʿrūḍ dāsht va ijāze khvāst
ke biravad va ū rā dafn va kafan namūde rājiʿ shavad. Ān Jawhar-i
Inqiṭāʿ farmūd: Daʿ al-mawtā li-yadfinūhu al-mawtā. Yaʿnī vāgudhār
murdehā rā tā dafn kunand murdehā.

Elsewhere in the Gospel it is recorded that one day the father of one of
Jesus's disciples died. He informed Jesus and asked permission to go,
bury and shroud his father, and return. That Essence of Detachment said:
``Leave the dead, so that the dead may bury him.'' That is, leave the
dead to bury the dead.

\subsection{128}\label{section-318}

\begin{otherlanguage}{persian}

و همچنین دو نفر از اهل کوفه خدمت حضرت امیر آمدند یکی را بیتی بود که
ارادهٔ بیع آن داشت و دیگری مشتری بود و قرار بر آن داده بودند که باطّلاع آن
حضرت این مبایعه وقوع یابد و قباله مسطور گردد آن مظهر امر الهی بکاتب
فرمودند که بنویس قد اشتری میّت عن میّت بیتاً محدوداً بحدود اربعة حدّ الی
القبر و حدّ الی اللّحد و حدّ الی الصّراط و حدّ امّا الی الجنّة و امّا الی النّار
حال اگر این دو نفر از صور حیات علوی زنده شده بودند و از قبر غفلت بمحبّت
آن حضرت مبعوث گشته بودند البتّه اطلاق موت بر ایشان نمیشد

\end{otherlanguage}

Va hamchunīn du nafar az ahl-i Kūfe khidmat-i Ḥaḍrat-i Amīr āmadand.
Yakī rā baytī būd ke irāde-yi bayʿ-i ān dāsht va dīgarī mushtarī būd va
qarār bar ān dāde būdand ke be-iṭṭilāʿ-i ān Ḥaḍrat īn mubāyaʿe vuqūʿ
yābad va qabāle masṭūr gardad. Ān Maẓhar-i Amr-i Ilāhī be-kātib
farmūdand ke binivīs: Qad ishtarā mayyitun ʿan mayyitin baytan maḥdūdan
bi-ḥudūdin arbaʿatin: ḥaddun ilā al-qabri wa-ḥaddun ilā al-laḥdi
wa-ḥaddun ilā al-ṣirāṭi wa-ḥaddun immā ilā al-jannati wa-immā ilā
al-nār. Ḥāl agar īn du nafar az ṣūr-i ḥayāt-i ʿAlavī zende shode būdand
va az qabr-i ghaflat be-maḥabbat-i ān Ḥaḍrat mabʿūth gashte būdand,
albatte iṭlāq-i mawt bar īshān namīshod.

Likewise, two men from Kūfa came into the presence of the Commander of
the Faithful. One owned a house that he wished to sell, and the other
was the buyer. They had agreed that the sale should take place with His
knowledge and that the deed should be drawn up. That Manifestation of
the Cause of God told the scribe to write: ``A dead man has purchased
from a dead man a house bounded on four sides: one boundary leads to the
grave, another to the tomb, another to the bridge, and another either to
paradise or to the fire.'' Now, had these two men been quickened by the
trumpet of ʿAlī's life and raised from the grave of heedlessness through
love for Him, they would certainly not have been described as dead.

\subsection{129}\label{section-319}

\begin{otherlanguage}{persian}

و هرگز در هیچ عهد و عصر جز حیات و بعث و حشر حقیقی مقصود انبیا و اولیا
نبوده و نیست اگر قدری تعقّل شود در همین بیان آن حضرت کشف جمیع امور میشود
که مقصود از لحد و قبر و صراط و جنّت و نار چه بود ولیکن چه چاره که جمیع
ناس در لحد نفس محجوب و در قبر هوی مدفونند خلاصه اگر قدری از زلال معرفت
الهی مرزوق شوید میدانید که حیات حقیقی حیات قلب است نه حیات جسد زیرا که
در حیات جسد همۀ ناس و حیوانات شریکند ولیکن این حیات مخصوص است بصاحبان
افئدهٔ منیره که از بحر ایمان شاربند و از ثمرهٔ ایقان مرزوق و این حیات را
موت از عقب نباشد و این بقا را فنا از پی نیاید چنانچه فرموده‌اند المؤمن حیّ
فی الدّارین اگر مقصود حیات ظاهرهٔ جسدی باشد که مشاهده میشود موت آن را اخذ
مینماید

\end{otherlanguage}

Va hargiz dar hīch ʿahd va ʿaṣr juz ḥayāt va baʿth va ḥashr-i ḥaqīqī
maqṣūd-i anbiyā va awliyā nabūde va nīst. Agar qadrī taʿaqqul shavad dar
hamīn bayān-i ān Ḥaḍrat kashf-i jamīʿ-i umūr mīshavad ke maqṣūd az laḥd
va qabr va ṣirāṭ va jannat va nār che būd. Va-līkin che chāre ke jamīʿ-i
nās dar laḥd-i nafs maḥjūb va dar qabr-i havā madfūnand. Khulāṣe agar
qadrī az zulāl-i maʿrifat-i ilāhī marzūq shavīd mīdānīd ke ḥayāt-i
ḥaqīqī ḥayāt-i qalb ast na ḥayāt-i jasad, zīrā ke dar ḥayāt-i jasad
hame-yi nās va ḥayvānāt sharīkand, va-līkin īn ḥayāt makhṣūṣ ast
be-ṣāḥibān-i afʾide-yi munīre ke az baḥr-i īmān shāriband va az
thamare-yi īqān marzūq. Va īn ḥayāt rā mawt az ʿaqab nabāshad va īn baqā
rā fanā az pay nayāyad. Chunānche farmūde-and: Al-muʾminu ḥayyun fī
al-dārayn. Agar maqṣūd ḥayāt-i ẓāhire-yi jasadī bāshad ke mushāhede
mīshavad mawt ān rā akhdh mīnamāyad.

In no age or era have the Prophets and holy ones ever intended anything
other than true life, resurrection, and gathering, nor do they now. A
little reflection on this very statement of His will disclose the
meaning of all these matters: what was intended by the tomb, the grave,
the bridge, paradise, and the fire. Yet what can be done, when all
people are veiled in the tomb of self and buried in the grave of desire?
In short, if you receive even a little of the pure water of divine
knowledge, you will know that true life is the life of the heart, not
the life of the body. For all people and animals share in the life of
the body, whereas this life belongs only to those with illumined hearts,
who drink from the ocean of faith and are nourished by the fruit of
certitude. This life is not followed by death, nor is this enduring
existence followed by extinction. As has been said: ``The believer is
alive in both worlds.'' If outward, bodily life were intended, it is
plain that death overtakes it.

\subsection{130}\label{section-320}

\begin{otherlanguage}{persian}

و همچنین بیانات دیگر که در همۀ کتب مذکور و ثبت شده مدلّ است بر این مطلب
عالی و کلمهٔ متعالی و همچنین آیۀ مبارکه که در حقّ حمزهٔ سیّد الشّهدآء و
ابوجهل نازل شد برهانیست واضح و حجّتی است لائح که میفرماید أ و من کان میّتاً
فأحییناه و جعلنا له نوراً یمشی به فی النّاس کمن مثله فی الظّلمات لیس بخارجٍ
منها و این آیه در وقتی از سماء مشیّت نازل شد که حمزه بردای مقدّس ایمان
متردّی شده بود و ابوجهل در کفر و اعراض ثابت و راسخ بود از مصدر الوهیّت
کبری و مکمن ربوبیّت عظمی حکم حیات بعد از موت در بارهٔ حمزه شد و بر خلاف در
حقّ ابوجهل این بود که نائرۀ کفر در قلوب مشرکین مشتعل شد و هوای اعراض
بحرکت آمد چنانچه فریاد برآوردند که حمزه چه زمان مرد و کی زنده شد و چه
وقت این حیات بر او عرضه گشت و چون این بیانات شریفه را ادراک نمینمودند و
باهل ذکر هم تمسّک نجستند تا رشحی از کوثر معانی بر آنها مبذول فرمایند لهذا
این نوع فسادها در عالم جریان یافت

\end{otherlanguage}

Va hamchunīn bayānāt-i dīgar ke dar hame-yi kutub madhkūr va thabt shode
mudill ast bar īn maṭlab-i ʿālī va kalime-yi mutaʿālī. Va hamchunīn
āye-yi mubārake ke dar ḥaqq-i Ḥamze-yi Sayyid al-Shuhadāʾ va Abū Jahl
nāzil shod burhānī-st vāḍiḥ va ḥujjatī-st lāʾiḥ, ke mīfarmāyad: A-wa-man
kāna maytan fa-aḥyaynāhu wa-jaʿalnā lahu nūran yamshī bihi fī al-nāsi
ka-man mathaluhu fī al-ẓulumāti laysa bi-khārijin minhā. Va īn āye dar
vaqtī az samā-yi mashiyyat nāzil shod ke Ḥamze be-ridā-yi muqaddas-i
īmān mutaraddī shode būd va Abū Jahl dar kufr va iʿrāḍ thābit va rāsikh
būd. Az maṣdar-i ulūhiyyat-i kubrā va makman-i rubūbiyyat-i ʿuẓmā ḥukm-i
ḥayāt baʿd az mawt dar bāre-yi Ḥamze shod va bar khilāf dar ḥaqq-i Abū
Jahl. Īn būd ke nāʾire-yi kufr dar qulūb-i mushrikīn mushtaʿil shod va
havā-yi iʿrāḍ be-ḥarakat āmad. Chunānche faryād bar-āvurdand ke Ḥamze
che zamān mord va kay zende shod va che vaqt īn ḥayāt bar ū ʿarḍe gasht?
Va chūn īn bayānāt-i sharīfe rā idrāk namīnamūdand va be-ahl-i dhikr ham
tamassuk najostand tā rashḥī az Kawthar-i maʿānī bar ānhā mabdhūl
farmāyand, li-hādhā īn nawʿ fasādhā dar ʿālam jarayān yāft.

Likewise, other statements mentioned and recorded in all the books
attest to this exalted theme and sublime word. The blessed verse
revealed concerning Ḥamza, the Lord of Martyrs, and Abū Jahl is also a
clear proof and an evident testimony. He says: ``Is one who was dead,
whom We brought to life and for whom We made a light by which he walks
among the people, like one who is in darkness and cannot emerge from
it?'' This verse descended from the heaven of the divine Will when Ḥamza
had clothed himself in the sacred robe of faith, while Abū Jahl remained
firm and steadfast in unbelief and rejection. From the Source of supreme
Divinity and the Sanctuary of the most exalted Lordship, the judgment of
life after death was pronounced upon Ḥamza, and the opposite upon Abū
Jahl. This kindled the fire of unbelief in the hearts of the idolaters
and stirred the winds of rejection, so that they cried out: ``When did
Ḥamza die, and when did he come to life? At what time was this life
bestowed upon him?'' Because they did not understand these noble
statements, nor did they turn to the people of remembrance so that they
might bestow upon them a sprinkling from the Kawthar of meanings, such
disorders spread throughout the world.

\subsection{131}\label{section-321}

\begin{otherlanguage}{persian}

چنانچه الیوم می‌‌بینی که با وجود شمس معانی جمیع ناس از اعالی و ادانی تمسّک
بجعلهای ظلمانی و مظاهر شیطانی جسته‌اند و متّصل مسائل مشکلهٔ خود را از ایشان
مستفسر میشوند و ایشان نظر بعدم عرفان چیزی جواب میگویند که ضرری بر اسباب
ظاهرهٔ ایشان نرساند و این معلوم و واضح است که جعل خود قسمتی از نسیم مشک
بقا نبرده و برضوان ریاحین معنوی قدم نگذاشته با وجود این چگونه میتواند
رائحهٔ عطر بمشامّ دیگران رساند لم‌یزل شأن این عباد این بوده و خواهد بود و
لن یفوز بآثار اللّه الّا الّذینهم اقبلوا الیه و اعرضوا عن مظاهر الشّیطان و
کذلک اثبت اللّه حکم الیوم من قلم العزّة علی لوح کان خلف سرادق العزّ مکنونا
اگر ملتفت باین بیانات شوید و تفکّر در ظاهر و باطن آن بفرمائید جمیع مسائل
مشکله را که الیوم سدّی شده میان عباد و معرفت یوم التّناد عارف شوی دیگر
احتیاج بسؤال نخواهی داشت انشآءاللّه امیدواریم که از شاطی بحر الهی لب‌تشنه
و محروم برنگردید و از حرم مقصود لایزالی بی‌بهره راجع نشوید دیگر تا همّت و
مجاهدهٔ شما چه کند

\end{otherlanguage}

Chonānche al-yawm mībīnī ke bā vojūd-e shams-e maʿānī jamīʿ-e nās az
aʿālī va adānī tamassok be joʿalhā-ye ẓolmānī va maẓāher-e shayṭānī
joste-and va mottaṣel masāʾel-e moshkele-ye khod rā az īshān mostafser
mīshavand va īshān naẓar be ʿadam-e ʿerfān chīzī javāb mīgūyand ke
żararī bar asbāb-e ẓāhere-ye īshān narasānad va īn maʿlūm va vāżeḥ ast
ke joʿal khod qesmatī az nasīm-e moshk-e baqā naborde va be reżvān-e
rayāḥīn-e maʿnavī qadam nagozāshte bā vojūd-e īn chegūne mītavānad
rāʾeḥe-ye ʿeṭr be mashāmm-e dīgarān rasānad lam-yazal shaʾn-e īn ʿebād
īn būde va khwāhad būd va lan yafūza bi-āthāri Allāhi illā alladhīna-hum
aqbalū ilayhi va aʿraḍū ʿan maẓāhiri al-shayṭān va kadhālika athbata
Allāhu ḥukma al-yawmi min qalami al-ʿizzati ʿalā lawḥin kāna khalfa
surādiqi al-ʿizzi maknūnan agar moltafet be īn bayānāt shavīd va
tafakkor dar ẓāher va bāṭen-e ān befar māʾīd jamīʿ-e masāʾel-e moshkele
rā ke al-yawm saddī shode miyān-e ʿebād va maʿrefat-e yawm al-tanād
ʿāref shavī dīgar eḥtiyāj be soʾāl nakhwāhī dāsht in shāʾa Allāh
omīdvārīm ke az shāṭī-ye baḥr-e elāhī lab-teshne va maḥrūm barnagardīd
va az ḥaram-e maqṣūd-e lā-yazālī bī-bahre rājeʿ nashavīd dīgar tā hemmat
va mojāhede-ye shomā che konad

Thus today you see that, despite the presence of the Sun of meanings,
all people, high and low alike, have clung to dark dung beetles and
satanic manifestations. They continually ask them about their difficult
questions, and these, lacking understanding, give answers that will do
no harm to their worldly interests. Yet it is clear and evident that the
dung beetle itself has had no share in the breeze of the musk of
eternity and has never set foot in the paradise of spiritual flowers.
How, then, can it convey a sweet fragrance to the senses of others? Such
has always been, and will always be, the way of these servants. None
will attain the traces of God except those who have turned toward Him
and turned away from the manifestations of Satan. Thus has God inscribed
the decree of this day with the Pen of glory upon a Tablet concealed
behind the tabernacle of majesty. If you heed these words and reflect
upon their outward and inward meanings, you will understand all the
difficult questions that today have become a barrier between people and
recognition of the Day of Mutual Calling. You will no longer need to ask
questions. God willing, We hope that you will not return thirsty and
deprived from the shore of the divine ocean, nor come back empty-handed
from the sanctuary of the everlasting Object of desire. It now remains
to be seen what your resolve and striving will accomplish.

\subsection{132}\label{section-322}

\begin{otherlanguage}{persian}

باری مقصود از این بیانات واضحه اثبات سلطنت آن سلطان السّلاطین بود حال
انصاف دهید که این سلطنت که بیک حرف و بیان اینهمه تصرّف و غلبه و هیمنه
داشته باشد اکبر و اعظم است یا سلطنت این سلاطین که بعد از اعانت رعایا و
فقرا ایشان را چند صباحی مردم بحسب ظاهر تمکین مینمایند ولیکن بقلب همه
معرض و مدبرند و این سلطنت بحرفی عالم را مسخّر نموده و حیات بخشیده و وجود
افاضه فرموده ما للتّراب و ربّ الأرباب چه میتوان ذکر نسبت نمود که همۀ
نسبتها منقطع است از ساحت قدس سلطنت او و اگر خوب ملاحظه شود خدّام درگه او
سلطنت مینمایند بر همۀ مخلوقات و موجودات چنانچه ظاهر شده و میشود

\end{otherlanguage}

Bārī maqṣūd az īn bayānāt-e vāżeḥe ethbāt-e salṭanat-e ān Solṭān
al-Salāṭīn būd ḥāl enṣāf dehīd ke īn salṭanat ke be yek ḥarf va bayān
īn-hame taṣarrof va ghalabe va haymane dāshte bāshad akbar va aʿẓam ast
yā salṭanat-e īn salāṭīn ke baʿd az eʿānat-e raʿāyā va foqarā īshān rā
chand ṣabāḥī mardom be ḥasab-e ẓāher tamkīn mīnamāyand valīkan be qalb
hame moʿreż va modberand va īn salṭanat be ḥarfī ʿālam rā mosakhkhar
namūde va ḥayāt bakhshīde va vojūd efāże farmūde mā li-l-turābi va Rabbi
al-Arbāb che mītavān dhekr-e nesbat namūd ke hame-ye nesbat-hā monqaṭeʿ
ast az sāḥat-e qods-e salṭanat-e ū va agar khūb molāḥeẓe shavad
khoddām-e dargah-e ū salṭanat mīnamāyand bar hame-ye makhlūqāt va
mawjūdāt chonānche ẓāher shode va mīshavad

In short, the purpose of these clear statements was to establish the
sovereignty of that King of kings. Now judge fairly: is a sovereignty
that exercises such influence, ascendancy, and authority through a
single word and utterance greater and more exalted, or the sovereignty
of these kings, whom people outwardly obey for a few days through the
assistance of their subjects and the poor, while in their hearts they
all reject them and turn away? This sovereignty, however, has subdued
the world with a word, bestowed life, and conferred existence. What has
dust to do with the Lord of lords? What comparison can be mentioned,
when every relationship is severed from the holy court of His
sovereignty? If the matter is carefully considered, the servants of His
threshold exercise sovereignty over all creatures and all beings, as has
been and continues to be made manifest.

\subsection{133}\label{section-323}

\begin{otherlanguage}{persian}

باری اینست یک معنی از سلطنت باطنی که نظر باستعداد و قابلیّت ناس ذکر شد و
از برای آن نقطۀ وجود و طلعت محمود سلطنتهاست که این مظلوم قادر بر اظهار
آن رتبه نیست و خلق لایق ادراک آن نه فسبحان‌ اللّه عمّا یصف العباد فی سلطنته
و تعالی عمّا هم یذکرون

\end{otherlanguage}

Bārī īnast yek maʿnī az salṭanat-e bāṭenī ke naẓar be esteʿdād va
qābeliyyat-e nās dhekr shod va az barā-ye ān noqṭe-ye vojūd va ṭalʿat-e
maḥmūd salṭanat-hāst ke īn maẓlūm qāder bar eẓhār-e ān rotbe nīst va
khalq lāyeq-e edrāk-e ān na fa-subḥāna Allāhi ʿammā yaṣifu al-ʿibādu fī
salṭanatihi va taʿālā ʿammā hum yadhkurūn

This, then, is one meaning of inward sovereignty, mentioned with regard
to the capacity and receptivity of the people. That Point of existence
and praiseworthy Countenance possesses sovereignties whose degree this
wronged One is unable to express, and which creation is unfit to
comprehend. Glorified, then, is God above what His servants describe of
His sovereignty, and exalted above what they recount.

\subsection{134}\label{section-324}

\begin{otherlanguage}{persian}

سؤالی مینمائیم از آن جناب که اگر مقصود از سلطنت حکم ظاهر و غلبه و اقتدار
ظاهر ملکی باشد که همۀ ناس مقهور شوند و بظاهر مطیع و منقاد گردند تا
دوستان مستریح و معزّز و دشمنان مخذول و منکوب شوند پس در حقّ ربّ العزّة که
مسلّماً سلطنت باسم اوست و جمیع بعظمت و شوکت او معترفند این نوع از سلطنت
صادق نمیآید چنانچه مشاهده مینمائی که اکثر ارض در تصرّف دشمنان اوست و جمیع
بر خلاف رضای او حرکت مینمایند و همه کافر و معرض و مدبرند از آنچه بآن امر
فرموده و مقبل و فاعلند آنچه را نهی نموده و دوستان او همیشه در دست دشمنان
مبتلا و مقهورند چنانچه همهٔ اینها اظهر من الشّمس واضحست

\end{otherlanguage}

Soʾālī mīnamāʾīm az ān janāb ke agar maqṣūd az salṭanat ḥokm-e ẓāher va
ghalabe va eqtedār-e ẓāher-e molkī bāshad ke hame-ye nās maqhūr shavand
va be ẓāher moṭīʿ va monqād gardand tā dūstān mostarīḥ va moʿazzaz va
doshmanān makhdhūl va mankūb shavand pas dar ḥaqq-e Rabb al-ʿEzzat ke
mosallaman salṭanat be esm-e ūst va jamīʿ be ʿaẓamat va shawkat-e ū
moʿtarefand īn nawʿ az salṭanat ṣādeq namīāyad chonānche moshāhede
mīnamāʾī ke akthar-e arż dar taṣarrof-e doshmanān-e ūst va jamīʿ bar
khelāf-e reżā-ye ū ḥarekat mīnamāyand va hame kāfer va moʿreż va
modberand az ānche be ān amr farmūde va moqbel va fāʿeland ānche rā nahy
namūde va dūstān-e ū hamīshe dar dast-e doshmanān mobtalā va maqhūrand
chonānche hame-ye īnhā aẓhar min al-shams vāżeḥast

We ask you a question: if sovereignty means outward rule, ascendancy,
and worldly power, such that all people are subdued and outwardly
obedient and submissive, so that friends are at ease and honored and
enemies are forsaken and crushed, then this kind of sovereignty does not
hold true of the Lord of glory, to whom sovereignty unquestionably
belongs and whose grandeur and majesty all acknowledge. For you see that
most of the earth is under the control of His enemies, and all act
contrary to His pleasure. They all disbelieve in, reject, and turn away
from what He has commanded, while embracing and carrying out what He has
forbidden. His friends are always afflicted and subjugated at the hands
of His enemies. All these things are clearer than the sun.

\subsection{135}\label{section-325}

\begin{otherlanguage}{persian}

پس بدان ای سائل طالب که هرگز سلطنت ظاهره نزد حقّ و اولیای او معتبر نبوده
و نخواهد بود و دیگر آنکه اگر مقصود از غلبه و قدرت قدرت و غلبهٔ ظاهری باشد
کار بسیار بر آن جناب سخت میشود مثل آنکه میفرماید و انّ جندنا لهم الغالبون
و در مقام دیگرمیفرماید یریدون ان یطفئوا نور اللّه بأفواههم و یأبی اللّه
الّا ان یتمّ نوره ولو کره الکافرون و دیگر هو الغالب فوق کلّ شیء مثل اینکه
اکثری از فرقان صریح بر این مطلب است

\end{otherlanguage}

Pas bedān ey sāʾel-e ṭāleb ke hargez salṭanat-e ẓāhere nazd-e Ḥaqq va
awliyā-ye ū moʿtabar nabūde va nakhwāhad būd va dīgar ānke agar maqṣūd
az ghalabe va qodrat qodrat va ghalabe-ye ẓāherī bāshad kār besyār bar
ān janāb sakht mīshavad methl-e ānke mīfarmāyad va inna jundanā la-humu
al-ghālibūn va dar maqām-e dīgar mīfarmāyad yurīdūna an yuṭfiʾū nūra
Allāhi bi-afwāhihim va yaʾbā Allāhu illā an yutimma nūrahu va law kariha
al-kāfirūn va dīgar huwa al-ghālibu fawqa kulli shayʾ methl-e īnke
aktharī az Forqān ṣarīḥ bar īn maṭlab ast

Know, then, O questioning seeker, that outward sovereignty has never
had, and will never have, any standing before God and His holy ones.
Furthermore, if power and ascendancy mean outward power and ascendancy,
your difficulty becomes very great. For He says, ``Indeed, Our host will
surely be victorious.'' Elsewhere He says, ``They wish to extinguish the
light of God with their mouths, but God refuses to do anything but
perfect His light, even though the unbelievers detest it.'' And again:
``He is the One who prevails over all things.'' Much of the Qurʾān
explicitly affirms this matter.

\subsection{136}\label{section-326}

\begin{otherlanguage}{persian}

و اگر مقصود این باشد که این همج رعاع میگویند مفرّی برای ایشان نمیماند مگر
انکار جمیع این کلمات قدسیّه و اشارات ازلیّه را نمایند زیرا که جندی از حسین
بن علی اعلی در ارض نبوده که اقرب الی ‌اللّه باشد و آن حضرت بر روی ارض مثلی
و شبهی نداشت لولاه لم ‌یکن مثله فی الملک با وجود این شنیدید که چه واقع شد
الا لعنة اللّه علی القوم الظّالمین

\end{otherlanguage}

Va agar maqṣūd īn bāshad ke īn hamaj-e raʿāʿ mīgūyand mafarrī barā-ye
īshān namīmānad magar enkār-e jamīʿ-e īn kalemāt-e qodsiyye va eshārāt-e
azaliyye rā namāyand zīrā ke jondī az Ḥosayn ebn-e ʿAlī aʿlā dar arż
nabūde ke aqrab ilā Allāh bāshad va ān ḥażrat bar rū-ye arż methlī va
shabhī nadāsht lawlāhu lam yakun mithluhu fī al-mulk bā vojūd-e īn
shenīdīd ke che vāqeʿ shod alā laʿnatu Allāhi ʿalā al-qawmi al-ẓālimīn

If the intended meaning were what this ignorant rabble says, they would
have no escape except to deny all these sacred words and eternal
allusions. For there was no soldier on earth more exalted than Ḥusayn,
son of ʿAlī, or nearer to God. He had no equal or likeness on the face
of the earth: ``But for him, there would have been none like him in the
world.'' Yet you have heard what happened. Indeed, the curse of God
rests upon the people who do wrong.

\subsection{137}\label{section-327}

\begin{otherlanguage}{persian}

حال اگر بر حسب ظاهر تفسیر کنید این آیه هیچ در حقّ اولیای خدا و جنود او بر
حسب ظاهر صادق نمیآید چه که آن حضرت که جندیّتش مثل شمس لائح و واضحست در
نهایت مغلوبیّت و مظلومیّت در ارض طفّ کأس شهادت را نوشیدند و همچنین در آیهٔ
مبارکه که میفرماید یریدون ان یطفئوا نور اللّه بأفواههم و یأبی اللّه الّا ان
یتمّ نوره ولو کره الکافرون اگر بر ظاهر ملکی تفسیر شود هرگز موافق نیاید
زیرا که همیشه انوار الهی را بر حسب ظاهر اطفا نمودند و سراجهای صمدانی را
خاموش کردند مع‌ذلک غلبه از کجا ظاهر میشود و منع در آیهٔ شریفه که میفرماید
و یأبی اللّه الّا ان یتمّ نوره چه معنی دارد چنانچه ملاحظه شد جمیع انوار از
دست مشرکین در محلّ امنی نیاسودند و شربت راحتی نیاشامیدند و مظلومیّت این
انوار بقسمی بود که هر نفسی بر آن جواهر وجود وارد میآورد آنچه را اراده
مینمود چنانچه همه را احصا و ادراک نمودند مع‌ذلک چگونه این مردم از عهدهٔ
معانی و بیان این کلمات الهی و آیات عزّ صمدانی برمیآیند

\end{otherlanguage}

Ḥāl agar bar ḥasab-e ẓāher tafsīr konīd īn āye hīch dar ḥaqq-e awliyā-ye
Khodā va jonūd-e ū bar ḥasab-e ẓāher ṣādeq namīāyad che ke ān ḥażrat ke
jondiyyatash methl-e shams lāʾeḥ va vāżeḥast dar nehāyat-e maghlūbiyyat
va maẓlūmiyyat dar arż-e Ṭaff kaʾs-e shahādat rā nūshīdand va hamchenīn
dar āye-ye mobārake ke mīfarmāyad yurīdūna an yuṭfiʾū nūra Allāhi
bi-afwāhihim va yaʾbā Allāhu illā an yutimma nūrahu va law kariha
al-kāfirūn agar bar ẓāher-e molkī tafsīr shavad hargez movāfeq nayāyad
zīrā ke hamīshe anvār-e elāhī rā bar ḥasab-e ẓāher eṭfā namūdand va
serājhā-ye ṣamadānī rā khāmūsh kardand maʿa-dhālek ghalabe az kojā ẓāher
mīshavad va manʿ dar āye-ye sharīfe ke mīfarmāyad va yaʾbā Allāhu illā
an yutimma nūrahu che maʿnī dārad chonānche molāḥeẓe shod jamīʿ-e anvār
az dast-e moshrekīn dar maḥall-e amnī nayāsūdand va sharbat-e rāḥatī
nayāshāmīdand va maẓlūmiyyat-e īn anvār be qesmī būd ke har nafsī bar ān
javāher-e vojūd vāred mīāvard ānche rā erāde mīnamūd chonānche hame rā
eḥṣā va edrāk namūdand maʿa-dhālek chegūne īn mardom az ʿohde-ye maʿānī
va bayān-e īn kalemāt-e elāhī va āyāt-e ʿezz-e ṣamadānī barmīāyand

Now, if you interpret this verse outwardly, it does not hold true at
all, in outward terms, of the holy ones of God and His hosts. For that
holy person, whose station as a soldier is as manifest and clear as the
sun, drank the cup of martyrdom in the land of Ṭaff in utter defeat and
oppression. Likewise, if the blessed verse in which He says, ``They wish
to extinguish the light of God with their mouths, but God refuses to do
anything but perfect His light, even though the unbelievers detest it,''
is interpreted in an outward, worldly sense, it will never accord with
events. For they have always outwardly extinguished the divine lights
and put out the lamps of the Eternal. Where, then, does victory appear?
And what does the refusal mean in the noble verse in which He says,
``But God refuses to do anything but perfect His light''? As has been
observed, none of these lights found rest in a place safe from the
idolaters or drank a draught of ease. So oppressed were these lights
that anyone inflicted upon those gems of existence whatever he wished,
as people have recounted and understood all these things. How, then, can
these people account for the meanings and explain these divine words and
verses of eternal glory?

\subsection{138}\label{section-328}

\begin{otherlanguage}{persian}

باری مقصود نچنانست که ادراک نمودند بلکه مقصود از غلبه و قدرت و احاطه
مقامی دیگر و امری دیگر است مثلاً ملاحظه فرمائید غلبهٔ ترشّحات دم آن حضرت را
که بر تراب ترشّح نموده و بشرافت و غلبهٔ آن دم تراب چگونه غلبه و تصرّف در
اجساد و ارواح ناس فرموده چنانچه هر نفسی برای استشفا بذرّه‌ئی‌ از آن مرزوق
شد شفا یافت و هر وجود که برای حفظ مال قدری از آن تراب مقدّس را بیقین کامل
و معرفت ثابتهٔ راسخه در بیت نگاه داشت جمیع مالش محفوظ ماند و این مراتب
تأثیرات آن است در ظاهر و اگر تأثیرات باطنیّه را ذکر نمایم البتّه خواهند
گفت تراب را ربّ الأرباب دانسته و از دین خدا بالمرّه خارج گشته

\end{otherlanguage}

Bārī maqṣūd nachonānast ke edrāk namūdand balke maqṣūd az ghalabe va
qodrat va eḥāṭe maqāmī dīgar va amrī dīgar ast mathalan molāḥeẓe
farmāʾīd ghalabe-ye tarashshoḥāt-e dam-e ān ḥażrat rā ke bar torāb
tarashshoḥ namūde va be sharāfat va ghalabe-ye ān dam torāb chegūne
ghalabe va taṣarrof dar ajsād va arvāḥ-e nās farmūde chonānche har nafsī
barā-ye esteshfā be dharre-ī az ān marzūq shod shefā yāft va har vojūd
ke barā-ye ḥefẓ-e māl qadrī az ān torāb-e moqaddas rā be yaqīn-e kāmel
va maʿrefat-e thābete-ye rāsekhe dar bayt negāh dāsht jamīʿ-e mālash
maḥfūẓ mānd va īn marāteb taʾthīrāt-e ān ast dar ẓāher va agar
taʾthīrāt-e bāṭeniyye rā dhekr namāyam albatte khwāhand goft torāb rā
Rabb al-Arbāb dāneste va az dīn-e Khodā be-l-marre khārej gashte

In short, the intended meaning is not what they have understood. Rather,
ascendancy, power, and encompassing dominion refer to another station
and another matter. Consider, for example, the ascendancy of the drops
of His blood that fell upon the dust, and how, through the nobility and
ascendancy of that blood, the dust has exercised power and influence
over the bodies and souls of the people. Whoever was granted a particle
of it for healing was healed, and whoever kept a little of that sacred
dust in his house to safeguard his possessions, with complete certainty
and firm, deeply rooted recognition, had all his possessions preserved.
These are its outward effects. Were I to mention its inward effects,
they would certainly say, ``He has regarded the dust as the Lord of
lords and has departed entirely from the religion of God.''

\subsection{139}\label{section-329}

\begin{otherlanguage}{persian}

و همچنین ملاحظه نما با اینکه بنهایت ذلّت آن حضرت شهید شد و احدی نبود که
آن حضرت را در ظاهر نصرت نماید و یا غسل دهد و کفن نماید مع‌ذلک حال چگونه
از اطراف و اکناف بلاد چه‌ قدر از مردم که شدّ رحال مینمایند برای حضور در آن
ارض که سر بر آن آستان بمالند اینست غلبه و قدرت الهی و شوکت و عظمت ربّانی

\end{otherlanguage}

Va hamchenīn molāḥeẓe namā bā īnke be nehāyat-e dhellat ān ḥażrat shahīd
shod va aḥadī nabūd ke ān ḥażrat rā dar ẓāher noṣrat namāyad va yā ghosl
dehad va kafan namāyad maʿa-dhālek ḥāl chegūne az aṭrāf va aknāf-e belād
che qadr az mardom ke shadd-e reḥāl mīnamāyand barā-ye ḥożūr dar ān arż
ke sar bar ān āstān bemāland īnast ghalabe va qodrat-e elāhī va shawkat
va ʿaẓamat-e rabbānī

Consider also how, although He was martyred in the utmost abasement and
there was no one to give Him outward aid, or to wash and shroud Him, so
many people now set out from every corner of the lands to be present in
that land and rub their heads upon that threshold. This is divine
ascendancy and power, heavenly majesty and grandeur.

\subsection{140}\label{section-330}

\begin{otherlanguage}{persian}

و همچه تصوّر ننمائی که این امور بعد از شهادت آن حضرت واقع شده و چه ثمری
برای آن حضرت مترتّب است زیرا که آن حضرت همیشه حیّ است بحیات الهی و در رفرف
امتناع قرب و سدرهٔ ارتفاع وصل ساکن و این جواهر وجود در مقام انفاق کل
قائمند یعنی جان و مال و نفس و روح همه را در راه دوست انفاق نموده و
مینمایند و هیچ رتبه‌ئی نزدشان احبّ از این مقام نیست عاشقان جز رضای معشوق
مطلبی ندارند و جز لقای محبوب منظوری نجویند

\end{otherlanguage}

Va hamcho taṣavvor nanamāʾī ke īn omūr baʿd az shahādat-e ān ḥażrat
vāqeʿ shode va che thamarī barā-ye ān ḥażrat motaratteb ast zīrā ke ān
ḥażrat hamīshe ḥayy ast be ḥayāt-e elāhī va dar rafraf-e emtenāʿ-e qorb
va sedre-ye ertefāʿ-e vaṣl sāken va īn javāher-e vojūd dar maqām-e
enfāq-e koll qāʾemand yaʿnī jān va māl va nafs va rūḥ hame rā dar rāh-e
Dūst enfāq namūde va mīnamāyand va hīch rotbe-ī nazdashān aḥabb az īn
maqām nīst ʿāsheqān joz reżā-ye maʿshūq maṭlabī nadārand va joz leqā-ye
maḥbūb manẓūrī najūyand

Do not imagine that, since these things occurred after His martyrdom,
they can bear no fruit for Him. For He is ever alive with divine life,
dwelling in the lofty realm of unapproachable nearness and upon the Lote
Tree of exalted union. These gems of existence abide in the station of
giving up everything: that is, they have offered, and continue to offer,
life, possessions, self, and spirit---all in the path of the Friend. No
rank is dearer to them than this station. Lovers have no desire but the
pleasure of the one they love and seek no goal but reunion with the
Beloved.

\subsection{141}\label{section-331}

\begin{otherlanguage}{persian}

دیگر اگر بخواهم رشحی از اسرار شهادت و ثمرهای آن را ذکر نمایم البتّه این
الواح کفایت نکند و بانتها نرساند انشآء‌اللّه امیدواریم که نسیم رحمتی بوزد
و شجرهٔ وجود از ربیع الهی خلعت جدید پوشد تا باسرار حکمت ربّانی پی‌بریم و
بعنایت او از عرفان کلّ شیء بی‌نیاز گردیم تا حال نفسی مشهود نگشت که به این
مقام فائز آید مگر معدودی قلیل که هیچ معروف نیستند تا بعد قضای الهی چه
اقتضا نماید و از خلف سرادق امضا چه ظاهر شود کذلک نذکر لکم من بدایع امر
اللّه و نلقی علیکم من نغمات الفردوس لعلّکم بمواقع العلم تصلون و من ثمرات
العلم ترزقون پس بیقین باید دانست که این شموس عظمت اگرچه بر نقطهٔ تراب
جالس باشند بر عرش اعظم ساکنند و اگر فلسی نزدشان موجود نباشد بر رفرف غنا
طائرند و در حینی که در دست دشمنان مبتلایند بر یمین قدرت و غلبه ساکن و در
کمال ذلّت ظاهره بر عرش عزّت صمدانی جالس و متّکی و در نهایت عجز ظاهری بر
کرسیّ سلطنت و اقتدار قائم

\end{otherlanguage}

Dīgar agar bekhwāham rashḥī az asrār-e shahādat va thamarhā-ye ān rā
dhekr namāyam albatte īn alvāḥ kefāyat nakonad va be entehā narasānad in
shāʾa Allāh omīdvārīm ke nasīm-e raḥmatī bevazad va shajare-ye vojūd az
rabīʿ-e elāhī khelʿat-e jadīd pūshad tā be asrār-e ḥekmat-e rabbānī
pey-barīm va be ʿenāyat-e ū az ʿerfān-e koll-e shayʾ bī-niyāz gardīm tā
ḥāl nafsī mashhūd nagasht ke be īn maqām fāʾez āyad magar maʿdūdī qalīl
ke hīch maʿrūf nīstand tā baʿd qażā-ye elāhī che eqteżā namāyad va az
khalf-e sorādeq-e emżā che ẓāher shavad kadhālika nadhkuru lakum min
badāʾiʿi amri Allāhi va nulqī ʿalaykum min naghamāti al-Firdaws
laʿallakum bi-mawāqiʿi al-ʿilmi taṣilūn va min thamarāti al-ʿilmi
turzaqūn pas be yaqīn bāyad dānest ke īn shomūs-e ʿaẓamat agarche bar
noqṭe-ye torāb jāles bāshand bar ʿarsh-e aʿẓam sākenand va agar felsī
nazdashān mawjūd nabāshad bar rafraf-e ghenā ṭāʾerand va dar ḥīnī ke dar
dast-e doshmanān mobtalāyand bar yamīn-e qodrat va ghalabe sāken va dar
kamāl-e dhellat-e ẓāhere bar ʿarsh-e ʿezzat-e ṣamadānī jāles va mottakī
va dar nehāyat-e ʿajz-e ẓāherī bar korsī-ye salṭanat va eqtedār qāʾem

Were I to wish to mention even a sprinkling of the mysteries of
martyrdom and its fruits, these Tablets would certainly not suffice or
bring the matter to its conclusion. God willing, We hope that a breeze
of mercy will blow and that the tree of existence will put on a new robe
through the divine springtime, so that we may discover the mysteries of
divine wisdom and, through His grace, become independent of the
knowledge of all things. So far, no one has been seen to attain this
station except a very few who are entirely unknown. It remains to be
seen what the divine decree will require hereafter and what will appear
from behind the tabernacle of fulfillment. Thus do We recount to you
some of the wonders of the Cause of God and impart to you some of the
melodies of Paradise, that you may reach the wellsprings of knowledge
and be nourished by its fruits. You must therefore know with certainty
that these Suns of grandeur, though seated upon a spot of dust, dwell
upon the most exalted throne. Though they possess not a penny, they soar
in the realm of wealth. While afflicted at the hands of their enemies,
they abide at the right hand of power and ascendancy; in the utmost
outward abasement, they sit and recline upon the throne of eternal
glory; and in the utmost outward helplessness, they occupy the seat of
sovereignty and power.

\subsection{142}\label{section-332}

\begin{otherlanguage}{persian}

این است که عیسی بن مریم روزی بر کرسی جالس شدند و بنغمات روح القدس
بیاناتی فرمودند که مضمون آن اینست ای مردم غذای من از گیاه ارض است که بآن
سدّ جوع مینمایم و فراش من سطح زمین است و سراج من در شبها روشنی ماه است و
مرکوب من پاهای من است و کیست از من غنی‌‌تر بر روی زمین قسم بخدا که صد‌هزار
غنا طائف حول این فقر است و صد‌هزار ملکوت عزّت طالب این ذلّت اگر برشحی از
بحر این معانی فائز شوی از عالم ملک و هستی درگذری و چون طیر نار در حول
سراج بهّاج جان بازی

\end{otherlanguage}

Īn ast ke ʿĪsā ibn-i Maryam rūzī bar korsī jāles shodand va
be-naghamāt-e Rūḥ al-Qudus bayānātī farmūdand ke maḍmūn-e ān īn ast: ey
mardom, ghadhā-ye man az giyāh-e arḍ ast ke be-ān sadd-e jūʿ mīnamāyam
va ferāsh-e man saṭḥ-e zamīn ast va serāj-e man dar shab-hā rowshanī-ye
māh ast va markūb-e man pāhā-ye man ast va kīst az man ghanī-tar bar
rū-ye zamīn? Qasam be-Khodā ke ṣad-hezār ghenā ṭāʾef-e ḥowl-e īn faqr
ast va ṣad-hezār malakūt-e ʿezzat ṭāleb-e īn dhellat. Agar be-rashḥī az
baḥr-e īn maʿānī fāʾez shavī az ʿālam-e molk va hastī dargodharī va chūn
ṭeyr-e nār dar ḥowl-e serāj-e bahhāj jān bāzī.

Thus Jesus, son of Mary, was seated one day upon a chair and spoke in
the melodies of the Holy Spirit, words whose purport is this: ``O
people, My food is the vegetation of the earth, with which I satisfy My
hunger; My bed is the surface of the ground; My lamp at night is the
light of the moon; and My mount is My own feet. Who on the face of the
earth is richer than I?'' By God, a hundred thousand riches circle
around this poverty, and a hundred thousand kingdoms of glory seek this
abasement. If you attain a sprinkling from the ocean of these meanings,
you will pass beyond the world of possessions and existence and, like
the moth, offer up your life as you circle the radiant lamp.

\subsection{143}\label{section-333}

\begin{otherlanguage}{persian}

و مثل این از حضرت صادق ذکر شده که روزی شخصی از اصحاب در خدمت آن حضرت
شکایت از فقر نمود آن جمال لایزالی فرمودند که تو غنی هستی و از شراب غنا
آشامیده‌ئی آن فقیر از بیان طلعت منیر متحیّر شد که چگونه غنی هستم که بفلسی
محتاجم آن حضرت فرمود آیا محبّت ما را نداری عرض نمود بلی یا ابن رسول اللّه
فرمود آیا به هزار دینار این را مبایعه مینمائی عرض نمود بجمیع دنیا و آنچه
در آنست نمیدهم حضرت فرمودند آیا نفسی که چنین چیزی نزد او باشد که او را
بعالم ندهد چگونه فقیر است

\end{otherlanguage}

Va methl-e īn az Ḥaḍrat-e Ṣādiq dhekr shode ke rūzī shakhṣī az aṣḥāb dar
khedmat-e ān Ḥaḍrat shekāyat az faqr namūd. Ān jamāl-e lā-yazālī
farmūdand ke to ghanī hastī va az sharāb-e ghenā āshāmīdeʾī. Ān faqīr az
bayān-e ṭalʿat-e monīr motaḥayyer shod ke chegūne ghanī hastam ke
be-falsī moḥtājam? Ān Ḥaḍrat farmūd: āyā maḥabbat-e mā rā nadārī? ʿArḍ
namūd: balī yā ibn-i Rasūl Allāh. Farmūd: āyā be-hezār dīnār īn rā
mobāyaʿe mīnamāʾī? ʿArḍ namūd: be-jamīʿ-e donyā va ānche dar ān ast
namīdaham. Ḥaḍrat farmūdand: āyā nafsī ke chonīn chīzī nazd-e ū bāshad
ke ū rā be-ʿālam nadahad chegūne faqīr ast?

Something similar is related of Ṣādiq: one day, one of His companions
complained of poverty in His presence. That everlasting Beauty said,
``You are rich and have drunk of the wine of wealth.'' That poor man was
bewildered by the words of that radiant Countenance: ``How am I rich
when I am in need of a penny?'' He asked, ``Do you not love Us?'' He
replied, ``Yes, O son of the Messenger of God.'' He asked, ``Would you
sell this for a thousand dinars?'' He replied, ``I would not give it up
for the whole world and all that is in it.'' He said, ``How can someone
be poor who possesses something he would not exchange for the world?''

\subsection{144}\label{section-334}

\begin{otherlanguage}{persian}

و این فقر و غنا و ذلّت و عزّت و سلطنت و قدرت و مادون آن که نزد این همج
رعاع معتبر است در آن ساحت مذکور نیست چنانچه میفرماید یا ایّها النّاس انتم
الفقرآء الی اللّه و اللّه هو الغنیّ پس مقصود از غنا غنای از ماسوی است و از
فقر فقر باللّه

\end{otherlanguage}

Va īn faqr va ghenā va dhellat va ʿezzat va salṭanat va qodrat va
mādūn-e ān ke nazd-e īn hamaj-e raʿāʿ moʿtabar ast dar ān sāḥat madhkūr
nīst. Chenānche mīfarmāyad: yā ayyuhā al-nāsu antumu al-fuqarāʾu ilā
Allāhi wa-Allāhu huwa al-Ghanī. Pas maqṣūd az ghenā ghenā-ye az māsavā
ast va az faqr faqr-e bi-Allāh.

This poverty and wealth, abasement and glory, sovereignty and power, and
all the other things esteemed by this ignorant rabble have no mention in
that court. As He says, ``O people, you are in need of God, and God is
the Self-Sufficient.'' Thus, wealth means independence from all else,
and poverty means need of God.

\subsection{145}\label{section-335}

\begin{otherlanguage}{persian}

و دیگر آنکه روزی عیسی بن مریم را یهود احاطه نمودند و خواستند که آن حضرت
اقرار فرماید بر اینکه ادّعای مسیحی و پیغمبری نمودند تا حکم بر کفر آن حضرت
نمایند و حدّ قتل بر او جاری سازند تا آنکه آن خورشید سماء معانی را در مجلس
فیلاطس و قیافا که اعظم علمای آن عصر بود حاضر نمودند و جمیع علما در آن
محضر حضور همرساندند و جمع کثیری برای تماشا و استهزاء و اذیّت آن حضرت
مجتمع شدند و هر چه از آن حضرت استفسار نمودند که شاید اقرار بشنوند حضرت
سکوت فرمودند و هیچ متعرّض جواب نشدند تا آنکه ملعونی برخاست و آمد در مقابل
آن حضرت و قسم داد آن حضرت را که آیا تو نگفتی که منم مسیح اللّه و منم ملک
الملوک و منم صاحب کتاب و منم مخرّب یوم سبت آن حضرت رأس مبارک را بلند
نموده فرمودند اما تری بأنّ ابن الانسان قد جلس عن یمین القدرة و القوّة یعنی
آیا نمی‌بینی که پسر انسان جالس بر یمین قدرت و قوّت الهی است و حال آنکه بر
حسب ظاهر هیچ اسباب قدرت نزد آن حضرت موجود نبود مگر قدرت باطنیّه که احاطه
نموده بود کلّ من فی السّموات و الأرض را دیگر چه ذکر نمایم که بعد از این
قول بر آن حضرت چه وارد آمد و چگونه باو سلوک نمودند بالأخره چنان در صدد
ایذا و قتل آن حضرت افتادند که بفلک چهارم فرار نمود

\end{otherlanguage}

Va dīgar ānke rūzī ʿĪsā ibn-i Maryam rā Yahūd eḥāṭe namūdand va khāstand
ke ān Ḥaḍrat eqrār farmāyad bar īnke eddeʿā-ye Masīḥī va peyghambarī
namūdand tā ḥokm bar kofr-e ān Ḥaḍrat namāyand va ḥadd-e qatl bar ū jārī
sāzand. Tā ānke ān khorshīd-e samā-ye maʿānī rā dar majles-e Fīlāṭos va
Qayāfā ke aʿẓam-e ʿolamā-ye ān ʿaṣr būd ḥāḍer namūdand va jamīʿ-e ʿolamā
dar ān maḥḍar ḥoḍūr ham-rasāndand va jamʿ-e kathīrī barā-ye tamāshā va
estehzā va adhiyyat-e ān Ḥaḍrat mojtameʿ shodand. Va har che az ān
Ḥaḍrat estefsār namūdand ke shāyad eqrār beshnavand Ḥaḍrat sokūt
farmūdand va hīch motaʿarreḍ-e javāb nashodand tā ānke malʿūnī barkhāst
va āmad dar moqābel-e ān Ḥaḍrat va qasam dād ān Ḥaḍrat rā ke āyā to
nagoftī ke manam Masīḥ Allāh va manam Malek al-Molūk va manam ṣāḥeb-e
ketāb va manam mokharreb-e yowm-e sabt? Ān Ḥaḍrat raʾs-e mobārak rā
boland namūde farmūdand: a-mā tarā bi-anna ibna al-insāni qad jalasa ʿan
yamīni al-qudrati wa-al-quwwa? Yaʿnī āyā namībīnī ke Pesar-e Ensān jāles
bar yamīn-e qodrat va qovvat-e elāhī ast? Va ḥāl ānke bar ḥasab-e ẓāher
hīch asbāb-e qodrat nazd-e ān Ḥaḍrat mowjūd nabūd magar qodrat-e
bāṭeniyye ke eḥāṭe namūde būd kulla man fī al-samāwāti wa-al-arḍ rā.
Dīgar che dhekr namāyam ke baʿd az īn qowl bar ān Ḥaḍrat che vāred āmad
va chegūne be-ū solūk namūdand? Be-al-ākhare chenān dar ṣadad-e īdhā va
qatl-e ān Ḥaḍrat oftādand ke be-falak-e chahārom farār namūd.

Furthermore, one day the Jews surrounded Jesus, son of Mary, and sought
to have Him admit that He had claimed to be the Messiah and a Prophet,
so that they could pronounce Him an unbeliever and carry out the death
penalty upon Him. At length they brought that Sun of the heaven of
meanings before the assembly of Pilate and Caiaphas, who was the
foremost religious scholar of that age. All the religious scholars
attended that gathering, and a large crowd assembled to watch, mock, and
torment Him. However much they questioned Him in the hope of hearing an
admission, He remained silent and offered no reply. Then an accursed man
arose, came before Him, and adjured Him, saying, ``Did You not say, `I
am the Messiah of God, I am the King of kings, I am the bearer of a
Book, and I am the destroyer of the Sabbath'?'' He raised His blessed
head and said, ``Do you not see that the Son of Man is seated at the
right hand of power and might?'' That is, ``Do you not see that the Son
of Man is seated at the right hand of divine power and might?'' Yet
outwardly He possessed no means of power, save that inward power which
encompassed all who are in the heavens and on earth. What more should I
relate of what befell Him after these words and how they treated Him? In
the end, they became so intent on tormenting and killing Him that He
fled to the fourth heaven.

\subsection{146}\label{section-336}

\begin{otherlanguage}{persian}

و همچنین در انجیل لوقا مذکور است که روزی دیگر آن حضرت بر یکی از یهود
گذشت که بمرض فلج مبتلا شده بود و بر سریر افتاده چون آن حضرت را دید
بقرائن شناخت آن حضرت را و استغاثه نمود و آن حضرت فرمودند قم عن سریرک
فانّک مغفورة خطایاک چند یهود که در آن مکان حضور داشتند اعتراض نمودند که
هل یمکن لأحد ان یغفر الخطایا الّا اللّه فالتفت المسیح الیهم و قال ایّما
اسهل ان اقول له قم فاحمل سریرک ام اقول له مغفورة خطایاک لتعلموا بأنّ لابن
الانسان سلطاناً علی الأرض لمغفرة الخطایا که ترجمهٔ آن بفارسی اینست چون آن
حضرت بآن عاجز مسکین فرمودند که برخیز بدرستی که معاصی تو آمرزیده شد جمعی
از یهود اعتراض نمودند که آیا جز پروردگار غالب قادر کسی قادر بر غفران
عباد هست آن حضرت ملتفت بایشان شده فرمودند که آیا کدام اسهل است نزد شما
از اینکه بگویم باین عاجز فالج برخیز و برو و یا آنکه بگویم آمرزیده است
گناهان تو تا آنکه بدانید که از برای پسر انسان سلطانیست در ارض برای آمرزش
ذنوب مذنبان اینست سلطنت حقیقی و اقتدار اولیای الهی همهٔ این تفاصیل که
مکرّر ذکر میشود از همه مقام و همه جا مقصود اینست که بر تلویحات کلمات
اصفیای الهی مطّلع شوید که شاید از بعضی عبارات قدم نلغزد و قلب مضطرب نشود

\end{otherlanguage}

Va hamchenīn dar Enjīl-e Lūqā madhkūr ast ke rūzī dīgar ān Ḥaḍrat bar
yekī az Yahūd godhasht ke be-maraḍ-e falaj mobtalā shode būd va bar
sarīr oftāde. Chūn ān Ḥaḍrat rā dīd be-qarāʾen shenākht ān Ḥaḍrat rā va
esteghāthe namūd va ān Ḥaḍrat farmūdand: qum ʿan sarīrika fa-innaka
maghfūratun khaṭāyāka. Chand Yahūd ke dar ān makān ḥoḍūr dāshtand
eʿterāḍ namūdand ke hal yumkinu li-aḥadin an yaghfira al-khaṭāyā illā
Allāh? Fa-iltafata al-Masīḥu ilayhim wa-qāla: ayyumā ashalu an aqūla
lahu qum fa-iḥmil sarīraka am aqūla lahu maghfūratun khaṭāyāka
li-taʿlamū bi-anna li-ibni al-insāni sulṭānan ʿalā al-arḍi li-maghfirati
al-khaṭāyā? Ke tarjome-ye ān be-Fārsī īn ast: chūn ān Ḥaḍrat be-ān
ʿājez-e meskīn farmūdand ke barkhīz be-dorostī ke maʿāṣī-ye to āmorzīde
shod jamʿī az Yahūd eʿterāḍ namūdand ke āyā joz Parvardegār-e ghāleb-e
qāder kasī qāder bar ghofrān-e ʿebād hast? Ān Ḥaḍrat moltafet be-īshān
shode farmūdand ke āyā kodām ashal ast nazd-e shomā az īnke begūyam
be-īn ʿājez-e fālej barkhīz va berow va yā ānke begūyam āmorzīde ast
gonāhān-e to tā ānke bedānīd ke az barā-ye Pesar-e Ensān solṭānī-st dar
arḍ barā-ye āmorzesh-e dhonūb-e modhnebān. Īn ast salṭanat-e ḥaqīqī va
eqtedār-e owliyā-ye elāhī. Hame-ye īn tafāṣīl ke mokarrar dhekr mīshavad
az hame maqām va hame jā maqṣūd īn ast ke bar talvīḥāt-e kalemāt-e
aṣfiyā-ye elāhī moṭṭaleʿ shavīd ke shāyad az baʿḍī ʿebārāt qadam
nalaghzad va qalb moḍṭareb nashavad.

Likewise, it is related in the Gospel of Luke that on another day He
passed by a Jew who was afflicted with paralysis and lay upon a bed.
When he saw Him, he recognized Him by the signs and cried out for help.
He said, ``Rise from your bed, for your sins are forgiven.'' Some Jews
who were present there objected, ``Can anyone forgive sins except God?''
Christ turned to them and said, ``Which is easier: to say to him, `Rise
and carry your bed,' or to say to him, `Your sins are forgiven,' so that
you may know that the Son of Man has authority on earth to forgive
sins?'' The Persian translation is this: when He said to that poor,
helpless man, ``Rise, for your sins have indeed been forgiven,'' a group
of Jews objected, ``Can anyone other than the all-conquering, almighty
Lord forgive His servants?'' He turned to them and said, ``Which is
easier in your eyes: for Me to say to this helpless paralytic, `Rise and
go,' or to say, `Your sins are forgiven,' so that you may know that the
Son of Man has authority on earth to forgive the sins of sinners?'' This
is the true sovereignty and power of the holy ones of God. The purpose
of all these details, repeatedly mentioned in every connection and in
every place, is that you may become aware of the allusions in the words
of God's chosen ones, so that certain expressions may not cause your
feet to slip or your heart to be troubled.

\subsection{147}\label{section-337}

\begin{otherlanguage}{arabic}

و بقدم یقین در صراط حقّ الیقین قدم گذاریم که لعلّ نسیم رضا از ریاض قبول
الهی بوزد و این فانیان را بملکوت جاودانی رساند و عارف شوی بر معانی سلطنت
و امثال آن که در اخبار و آیات ذکر یافته و دیگر آنکه بر آن جناب محقّق و
معلوم بوده آنچه را که یهود و نصاری بآن تمسّک جسته‌اند و بر جمال محمّدی
اعتراض مینمودند بعینه در این زمان اصحاب فرقان بهمان تشبّث نموده و بر نقطهٔ
بیان روح من فی ملکوت الأمر فداه اعتراض مینمایند این بی‌خردان را مشاهده
فرما که حرف یهودان را الیوم میگویند و شاعر نیستند فنعم ما نزّل من قبل فی
شأنهم ذرهم فی خوضهم یلعبون و لعمرک انّهم لفی سکرتهم یعمهون

\end{otherlanguage}

Va be-qadam-e yaqīn dar ṣerāṭ-e ḥaqq al-yaqīn qadam godhārīm ke laʿall-e
nasīm-e reḍā az riyāḍ-e qabūl-e elāhī bevazad va īn fāniyān rā
be-malakūt-e jāvedānī rasānad va ʿāref shavī bar maʿānī-ye salṭanat va
amthāl-e ān ke dar akhbār va āyāt dhekr yāfte. Va dīgar ānke bar ān
janāb moḥaqqaq va maʿlūm būde ānche rā ke Yahūd va Naṣārā be-ān tamassok
joste-and va bar jamāl-e Moḥammadī eʿterāḍ mīnamūdand be-ʿeyne dar īn
zamān aṣḥāb-e Forqān be-hamān tashabboth namūde va bar Noqṭe-ye Bayān
rūḥu man fī malakūti al-amri fidāhu eʿterāḍ mīnamāyand. Īn bī-kheradān
rā moshāhede farmā ke ḥarf-e Yahūdān rā al-yowm mīgūyand va shāʿer
nīstand. Fa-niʿma mā nuzzila min qablu fī shaʾnihim: dharhum fī
khawḍihim yalʿabūn. Wa-la-ʿamruka innahum la-fī sakratihim yaʿmahūn.

And let us tread with certainty the path of true certainty, that perhaps
the breeze of good pleasure may blow from the gardens of divine
acceptance and bring these mortal beings to the everlasting Kingdom, and
that you may recognize the meanings of sovereignty and similar terms
mentioned in the traditions and verses. Furthermore, it has been clearly
established and known to you that the followers of the Qur'an now cling
to precisely the same things to which the Jews and Christians clung in
objecting to the Beauty of Muḥammad, and raise objections against the
Point of the Bayān---may the souls of all in the kingdom of the Cause be
sacrificed for Him. Behold these senseless people: today they repeat the
words of the Jews and are unaware of it. How fitting, then, is what was
sent down before concerning them: ``Leave them to amuse themselves in
their vain discourse,'' and ``By your life, they wander blindly in their
intoxication.''

\subsection{148}\label{section-338}

\begin{otherlanguage}{persian}

چون غیب ازلی و ساذج هویّه شمس محمّدی را از افق علم و معانی مشرق فرمود از
جمله اعتراضات علمای یهود آن بود که بعد از موسی نبی مبعوث نشود بلی طلعتی
در کتاب مذکور است که باید ظاهر شود و ترویج ملّت و مذهب او را نماید تا
شریعهٔ شریعت مذکورهٔ در تورات همهٔ ارض را احاطه نماید اینست که از لسان آن
ماندگان وادی بعد و ضلالت سلطان احدیّت میفرماید و قالت الیهود ید اللّه
مغلولة غلّت ایدیهم و لعنوا بما قالوا بل یداه مبسوطتان ترجمهٔ آن اینست که
گفتند یهودان دست خدا بسته شده بسته باد دستهای خود ایشان و ملعون شدند
بآنچه افترا بستند بلکه دستهای قدرت الهی همیشه باز و مهیمن است ید اللّه
فوق ایدیهم

\end{otherlanguage}

Chūn Ghayb-e azalī va Sādhaj-e hoviyye shams-e Moḥammadī rā az ofoq-e
ʿelm va maʿānī moshreq farmūd az jomle eʿterāḍāt-e ʿolamā-ye Yahūd ān
būd ke baʿd az Mūsā nabī mabʿūth nashavad. Balī ṭalʿatī dar ketāb
madhkūr ast ke bāyad ẓāher shavad va tarvīj-e mellat va madhhab-e ū rā
namāyad tā sharīʿe-ye sharīʿat-e madhkūre-ye dar Towrāt hame-ye arḍ rā
eḥāṭe namāyad. Īn ast ke az lesān-e ān māndegān-e vādī-ye boʿd va
ḍalālat Solṭān-e aḥadiyyat mīfarmāyad: wa-qālati al-Yahūdu yadu Allāhi
maghlūlatun ghullat aydīhim wa-luʿinū bi-mā qālū bal yadāhu mabsūṭatān.
Tarjome-ye ān īn ast ke goftand Yahūdān dast-e Khodā baste shode baste
bād dast-hā-ye khod-e īshān va malʿūn shodand be-ānche efterā bastand
balke dast-hā-ye qodrat-e elāhī hamīshe bāz va mohaymen ast. Yadu Allāhi
fawqa aydīhim.

When the eternal Unseen, the pure Essence of Being, caused the Sun of
Muḥammad to rise from the horizon of knowledge and meanings, one of the
objections of the Jewish religious scholars was that no prophet would be
raised up after Moses. They conceded that a figure was mentioned in the
Book who must appear and promote His faith and religion, so that the
stream of the law set forth in the Torah might encompass the whole
earth. Thus the Sovereign of Oneness speaks of the words of those who
remained in the valley of remoteness and error: ``The Jews said, `God's
hand is fettered.' Fettered be their own hands, and cursed were they for
what they said! Rather, both His hands are outstretched.'' Its
translation is this: the Jews said that God's hand had been bound. May
their own hands be bound, and they were cursed for the falsehood they
uttered. Rather, the hands of divine power are always open and
all-controlling. ``God's hand is above their hands.''

\subsection{149}\label{section-339}

\begin{otherlanguage}{persian}

اگرچه شرح نزول این آیه را علمای تفسیر مختلف ذکر نموده‌اند ولیکن بر مقصود
ناظر شوید که میفرماید نچنین است یهود خیال نمودند که سلطان حقیقی طلعت
موسوی را خلق نمود و خلعت پیغمبری بخشید و دیگر دستهایش مغلول و بسته شد و
قادر نیست بر ارسال رسولی بعد از موسی ملتفت این قول بی‌معنی شوید که چه ‌قدر
از شریعهٔ علم و دانش دور است و الیوم جمیع این مردم بامثال این مزخرفات
مشغولند و هزار سال بیش میگذرد که این آیه را تلاوت مینمایند و بر یهود من
حیث لا یشعر اعتراض مینمایند و ملتفت نشدند و ادراک ننمودند باینکه خود سرّاً
و جهراً میگویند آنچه را که یهود بآن معتقدند چنانچه شنیده‌اید که میگویند
جمیع ظهورات منتهی شده و ابواب رحمت الهی مسدود گشته دیگر از مشارق قدس
معنوی شمسی طالع نمیشود و از بحر قدم صمدانی امواجی ظاهر نگردد و از خیام
غیب ربّانی هیکلی مشهود نیاید اینست ادراک این همج رعاع فیض کلّیّه و رحمت
منبسطه که بهیچ عقلی و ادراکی انقطاع آن جایز نیست جایز دانسته و از اطراف
و جوانب کمر ظلم بسته و همّت گماشته‌اند که نار سدره را بماء ملح ظنون مخمود
نمایند و غافل از اینکه زجاج قدرت سراج احدیّت را در حصن حفظ خود محفوظ
میدارد و همین ذلّت کافیست این گروه را که از اصل مقصود محروم ماندند و از
لطیفه و جوهر امر محجوب گشتند لأجل آنکه منتها فیض الهی که برای عباد مقدّر
شده لقآء اللّه و عرفان اوست که کل بآن وعده داده شده‌اند و این نهایت فیض
فیّاض قدم است برای عباد او و کمال فضل مطلق است برای خلق او که هیچیک از
این عباد بآن مرزوق نشدند و باین شرافت کبری مشرّف نگشتند و با اینکه چه ‌قدر
از آیات منزله که صریح باین مطلب عظیم و امر کبیر است مع‌ذلک انکار نموده‌اند
و بهوای خود تفسیر کرده‌اند چنانچه میفرماید و الّذین کفروا بآیات ‌اللّه و
لقائه اولئک یئسوا من رحمتی واولئک لهم عذاب الیم و همچنین میفرماید الّذین
یظنّون انّهم ملاقوا ربّهم و انّهم الیه راجعون و در مقام دیگر قال الّذین یظنّون
انّهم ملاقوا اللّه کم من فئة قلیلة غلبت فئة کثیرة و در مقامی دیگر فمن کان
یرجو لقآء ربّه فلیعمل عملاً صالحاً و در مقامی دیگر یدبّر الأمر یفصّل الآیات
لعلّکم بلقآء ربّکم توقنون

\end{otherlanguage}

Agarche sharḥ-e nozūl-e īn āye rā ʿolamā-ye tafsīr mokhtalef dhekr
namūde-and valīkan bar maqṣūd nāẓer shavīd ke mīfarmāyad: na chonīn ast.
Yahūd khiyāl namūdand ke Solṭān-e ḥaqīqī ṭalʿat-e Mūsavī rā khalq namūd
va khelʿat-e peyghambarī bakhshīd va dīgar dast-hāyash maghlūl va baste
shod va qāder nīst bar ersāl-e rasūlī baʿd az Mūsā. Moltafet-e īn qowl-e
bī-maʿnī shavīd ke che qadr az sharīʿe-ye ʿelm va dānesh dūr ast va
al-yowm jamīʿ-e īn mardom be-amthāl-e īn mozakhrafāt mashghūland va
hezār sāl bīsh mīgodharad ke īn āye rā talāvat mīnamāyand va bar Yahūd
min ḥaythu lā yashʿur eʿterāḍ mīnamāyand va moltafet nashodand va edrāk
nanamūdand be-īnke khod serran va jahran mīgūyand ānche rā ke Yahūd
be-ān moʿtaqedand. Chenānche shenīde-īd ke mīgūyand jamīʿ-e ẓohūrāt
montahī shode va abvāb-e raḥmat-e elāhī masdūd gashte dīgar az
mashāreq-e qods-e maʿnavī shamsī ṭāleʿ namīshavad va az baḥr-e qedam-e
ṣamadānī amvājī ẓāher nagardad va az khiyām-e ghayb-e rabbānī hey kalī
mashhūd nayāyad. Īn ast edrāk-e īn hamaj-e raʿāʿ. Feyḍ-e kolliyye va
raḥmat-e monbaseṭe ke be-hīch ʿaqlī va edrākī enqeṭāʿ-e ān jāʾez nīst
jāʾez dāneste va az aṭrāf va javāneb kamar-e ẓolm baste va hemmat
gomāshte-and ke nār-e Sedre rā be-māʾ-e melḥ-e ẓonūn makhmūd namāyand va
ghāfel az īnke zojāj-e qodrat serāj-e aḥadiyyat rā dar ḥeṣn-e ḥefẓ-e
khod maḥfūẓ mīdārad. Va hamīn dhellat kāfī-st īn gorūh rā ke az aṣl-e
maqṣūd maḥrūm māndand va az laṭīfe va jowhar-e amr maḥjūb gashtand
li-ajl-e ānke montahā feyḍ-e elāhī ke barā-ye ʿebād moqaddar shode leqāʾ
Allāh va ʿerfān-e ū-st ke koll be-ān vaʿde dāde shode-and va īn
nehāyat-e feyḍ-e Fayyāḍ-e qedam ast barā-ye ʿebād-e ū va kamāl-e faḍl-e
moṭlaq ast barā-ye khalq-e ū ke hīch-yek az īn ʿebād be-ān marzūq
nashodand va be-īn sharāfat-e kobrā mosharraf nagashtand. Va bā īnke che
qadr az āyāt-e monzale ke ṣarīḥ be-īn maṭlab-e ʿaẓīm va amr-e kabīr ast
maʿa-dhālek enkār namūde-and va be-havā-ye khod tafsīr karde-and.
Chenānche mīfarmāyad: wa-alladhīna kafarū bi-āyāti Allāhi wa-liqāʾihi
ulāʾika yaʾisū min raḥmatī wa-ulāʾika lahum ʿadhābun alīm. Va hamchenīn
mīfarmāyad: alladhīna yaẓunnūna annahum mulāqū rabbihim wa-annahum
ilayhi rājiʿūn. Va dar maqām-e dīgar: qāla alladhīna yaẓunnūna annahum
mulāqū Allāhi kam min fiʾatin qalīlatin ghalabat fiʾatan kathīra. Va dar
maqāmī dīgar: fa-man kāna yarjū liqāʾa rabbihi fa-l-yaʿmal ʿamalan
ṣāliḥan. Va dar maqāmī dīgar: yudabbiru al-amra yufaṣṣilu al-āyāti
laʿallakum bi-liqāʾi rabbikum tūqinūn.

Although the commentators have given differing accounts of the
circumstances in which this verse was revealed, consider its intended
meaning, for He says, ``It is not so.'' The Jews imagined that the true
Sovereign had created the person of Moses and bestowed upon Him the robe
of prophethood, after which His hands became fettered and bound, and He
was powerless to send a Messenger after Moses. Consider how far this
senseless assertion lies from the stream of knowledge and understanding.
Today all these people occupy themselves with similar absurdities. For
more than a thousand years they have recited this verse and unwittingly
objected to the Jews, without noticing or understanding that they
themselves say, in private and in public, the very things the Jews
believe. As you have heard, they say that all revelations have ended and
the gates of divine mercy have been closed: no sun will ever again rise
from the dawning places of spiritual holiness, no waves will appear from
the ocean of eternal, all-sufficing being, and no figure will be seen
emerging from the tabernacles of the divine Unseen. Such is the
understanding of this ignorant rabble. They have deemed it possible for
that universal grace and all-embracing mercy to cease, though no reason
or understanding can admit of its cessation. On every side they have
girded themselves for oppression and bent their efforts toward
extinguishing the fire of the Lote Tree with the briny water of their
conjectures, unaware that the glass of power keeps the lamp of Oneness
safe within the fortress of its protection. This abasement alone is
enough for this people: they have remained deprived of the very object
of their quest and veiled from the subtle essence and substance of the
Cause. For the highest divine grace ordained for His servants is to meet
God and recognize Him, which has been promised to all. This is the
utmost grace of the eternal Bestower to His servants and the perfection
of His absolute bounty to His creation. Yet not one of these people has
received it or been honored with this supreme distinction. Although so
many revealed verses explicitly proclaim this momentous subject and
mighty matter, they have nevertheless denied it and interpreted it
according to their own desires. As He says, ``Those who disbelieve in
the signs of God and in meeting Him have despaired of My mercy, and a
painful punishment awaits them.'' Likewise He says, ``Those who are
certain that they will meet their Lord and that they will return to
Him.'' And elsewhere: ``Those who were certain that they would meet God
said, `How many a small company has overcome a great company!'\,'' And
in another place: ``Whoever hopes to meet his Lord, let him do righteous
deeds.'' And elsewhere: ``He directs the affairs and sets out the signs
in detail, so that you may be certain of meeting your Lord.''

\subsection{150}\label{section-340}

\begin{otherlanguage}{persian}

جمیع این آیات مدلّهٔ بر لقا را که حکمی محکمتر از آن در کتب سماوی ملحوظ
نگشته انکار نموده‌اند و از این رتبهٔ بلند اعلی و مرتبهٔ ارجمند ابهی خود را
محروم ساخته‌اند و بعضی ذکر نموده‌اند که مقصود از لقا تجلّی اللّه است در
قیامت و حال آنکه اگر گویند تجلّی عام مقصود است این در همۀ اشیاء موجود است
چنانچه از قبل ثابت شد که همۀ اشیاء محل و مظهر تجلّی آن سلطان حقیقی هستند
و آثار اشراق شمس مجلّی در مرایای موجودات موجود و لائحست بلکه اگر انسان را
بصر معنوی الهی مفتوح شود ملاحظه مینماید که هیچ شیء بی‌ظهور تجلّی پادشاه
حقیقی موجود نه چنانچه همهٔ ممکنات و مخلوقات را ملاحظه مینمائید که حاکیند
از ظهور و بروز آن نور معنوی و ابواب رضوان الهی را مشاهده میفرمائید که در
همهٔ اشیاء مفتوح گشته برای ورود طالبین در مداین معرفت و حکمت و دخول
واصلین در حدائق علم و قدرت و در هر حدیقه‌ئی‌‌ عروس معانی ملاحظه آید که در
غرفه‌های کلمات در نهایت تزیین و تلطیف جالسند و اکثر آیات فرقانی بر این
مطلب روحانی مدلّ و مشعر است و ان من شیء الّا یسبّح بحمده شاهدیست ناطق و کلّ
شیء احصیناه کتاباً گواهیست صادق حال اگر مقصود از لقآء اللّه لقاء این
تجلّیات باشد پس جمیع ناس بلقاء طلعت لایزال آن سلطان بیمثال مشرّفند دیگر
تخصیص بقیامت چرا

\end{otherlanguage}

Jamīʿ-e īn āyāt-e mudilla-ye bar liqā rā ke ḥukmī muḥkamtar az ān dar
kutub-e samāvī malḥūẓ nagashte enkār namūde-and va az īn rotbe-ye
boland-e aʿlā va martabe-ye arjomand-e abhā khod rā maḥrūm sākhte-and va
baʿżī dhekr namūde-and ke maqṣūd az liqā tajallī-ye Allāh ast dar
qiyāmat va ḥāl ānke agar gūyand tajallī-ye ʿāmm maqṣūd ast īn dar
hame-ye ashyāʾ mowjūd ast chenānche az qabl thābet shod ke hame-ye
ashyāʾ maḥall va maẓhar-e tajallī-ye ān solṭān-e ḥaqīqī hastand va
āthār-e eshrāq-e shams-e mujallī dar marāyā-ye mowjūdāt mowjūd va lāʾeḥ
ast balke agar ensān rā baṣar-e maʿnavī-ye elāhī maftūḥ shavad molāḥeẓe
mīnamāyad ke hīch shayʾ bī-ẓohūr-e tajallī-ye pādeshāh-e ḥaqīqī mowjūd
na chenānche hame-ye momkenāt va makhlūqāt rā molāḥeẓe mīnamāyīd ke
ḥākī-and az ẓohūr va borūz-e ān nūr-e maʿnavī va abvāb-e reżvān-e elāhī
rā moshāhede mīfarmāyīd ke dar hame-ye ashyāʾ maftūḥ gashte barā-ye
vorūd-e ṭālebīn dar madāyen-e maʿrefat va ḥekmat va dokhūl-e vāṣelīn dar
ḥadāʾeq-e ʿelm va qodrat va dar har ḥadīqe-ī ʿarūs-e maʿānī molāḥeẓe
āyad ke dar ghorfe-hā-ye kalemāt dar nehāyat-e tazyīn va talṭīf
jāles-and va akthar-e āyāt-e Forqānī bar īn maṭlab-e rūḥānī modall va
moshʿer ast wa-in min shayʾin illā yusabbiḥu bi-ḥamdihi shāhedī-st nāṭeq
va kulla shayʾin aḥṣaynāhu kitāban govāhī-st ṣādeq ḥāl agar maqṣūd az
liqāʾ-e Allāh liqāʾ-e īn tajalliyāt bāshad pas jamīʿ-e nās be-liqāʾ-e
ṭalʿat-e lā-yazāl-e ān solṭān-e bī-methāl mosharraf-and dīgar takhṣīṣ
be-qiyāmat cherā

They have denied all these verses that attest to the meeting with God,
although no more firmly established decree is to be found in the
heavenly Books, and have deprived themselves of this lofty, most exalted
rank and this precious, most glorious station. Some have stated that
what is meant by this meeting is the revelation of God on the Day of
Resurrection. Yet if they say that His universal revelation is intended,
this exists in all things, as it was established earlier that all things
are the seats and manifestations of the revelation of that true
Sovereign. The traces of the radiance of the illuminating Sun are
present and visible in the mirrors of created things. Indeed, if a
person's divine, spiritual sight is opened, they will observe that
nothing exists without manifesting the revelation of the true King. Thus
you see that all contingent beings and created things bear witness to
the manifestation and appearance of that spiritual Light, and you behold
the gates of the divine Paradise opened in all things, so that seekers
may enter the cities of knowledge and wisdom, and those who have
attained may enter the gardens of learning and power. In every garden,
the bride of inner meanings is seen seated in the chambers of words,
adorned with the utmost beauty and refinement. Most of the verses of the
Qur'an indicate and attest to this spiritual truth. ``There is nothing
that does not glorify Him with praise'' is an eloquent witness, and ``We
have recorded everything in a book'' is a faithful testimony. Now, if
meeting God means encountering these revelations, then all people are
honored with meeting the everlasting countenance of that peerless
Sovereign. Why, then, should it be restricted to the Day of
Resurrection?

\subsection{151}\label{section-341}

\begin{otherlanguage}{persian}

و اگر گویند مقصود تجلّی خاصّ است آنهم اگر در عین ذاتست در حضرت علم ازلاً
چنانچه جمعی از صوفیّه این مقام را تعبیر بفیض اقدس نموده‌اند بر فرض تصدیق
این رتبه صدق لقا برای نفسی در این مقام صادق نیاید لأجل آنکه این رتبه در
غیب ذات محقّق است و احدی بآن فائز نشود السّبیل مسدود و الطّلب مردود افئدهٔ
مقرّبین باین مقام طیران ننماید تا چه رسد بعقول محدودین و محتجبین

\end{otherlanguage}

Va agar gūyand maqṣūd tajallī-ye khāṣṣ ast ān ham agar dar ʿayn-e dhāt
ast dar ḥażrat-e ʿelm azalan chenānche jamʿī az ṣūfiyye īn maqām rā
taʿbīr be-fayż-e aqdas namūde-and bar farż-e taṣdīq-e īn rotbe ṣedq-e
liqā barā-ye nafsī dar īn maqām ṣādeq nayāyad li-ajl-e ānke īn rotbe dar
ghayb-e dhāt moḥaqqaq ast va aḥadī be-ān fāʾez nashavad al-sabīlu
masdūdun wa-l-ṭalabu mardūdun afʾede-ye moqarrabīn be-īn maqām ṭayarān
nanamāyad tā che resad be-ʿoqūl-e maḥdūdīn va moḥtajebīn

And if they say that His particular revelation is intended, then, if
this is within the Essence itself, in the eternal presence of divine
knowledge---as some Sufis have described this station as the Most Holy
Outpouring---even if this rank is acknowledged, the term ``meeting''
cannot apply to any soul at this station. For this rank exists in the
hidden depths of the Essence, and no one can attain it. ``The way is
barred, and the quest is rejected.'' The hearts of those brought near to
God cannot soar to this station, much less the minds of those who are
limited and veiled.

\subsection{152}\label{section-342}

\begin{otherlanguage}{persian}

و اگر گویند تجلّی ثانیست که معبّر بفیض مقدّس شده این مسلّماً در عالم خلق است
یعنی در عالم ظهور اوّلیّه و بروز بدعیّه و این مقام مختصّ بانبیا و اولیای
اوست چه که اعظم و اکبر از ایشان در عوالم وجود موجود نگشته چنانچه جمیع بر
این مطلب مقرّ و مذعنند و ایشانند محالّ و مظاهر جمیع صفات ازلیّه و اسماء
الهیّه و ایشانند مرایائی که تمام حکایت مینمایند و جمیع آنچه بایشان راجع
است فی‌ا‌لحقیقه بحضرت ظاهر مستور راجع و معرفت مبدء و وصول باو حاصل نمیشود
مگر بمعرفت و وصول این کینونات مشرقه از شمس حقیقت پس از لقاء این انوار
مقدّسه لقآء اللّه حاصل میشود و از علمشان علم ‌اللّه و از وجهشان وجه‌ اللّه و
از اوّلیّت و آخریّت و ظاهریّت و باطنیّت این جواهر مجرّده ثابت میشود از برای آن
شمس حقیقت بأنّه هو الأوّل و الآخر و الظّاهر و الباطن و همچنین سایر اسماء
عالیه و صفات متعالیه لهذا هر نفسی که باین انوار مضیئهٔ ممتنعه و شموس
مشرقهٔ لائحه در هر ظهور موفّق و فائز شد او بلقآء اللّه فائز است و در مدینهٔ
حیات ابدیّهٔ باقیه وارد و این لقا میسّر نشود برای احدی الّا در قیامت که قیام
نفس‌ اللّه است بمظهر کلّیّهٔ خود

\end{otherlanguage}

Va agar gūyand tajallī-ye thānī-st ke moʿabbar be-fayż-e moqaddas shode
īn mosallaman dar ʿālam-e khalq ast yaʿnī dar ʿālam-e ẓohūr-e avvaliyye
va borūz-e badʿiyye va īn maqām mokhtaṣṣ be-anbiyā va owliyā-ye ū-st che
ke aʿẓam va akbar az īshān dar ʿavālem-e vojūd mowjūd nagashte chenānche
jamīʿ bar īn maṭlab moqerr va modhʿen-and va īshān-and maḥāll va
maẓāher-e jamīʿ-e ṣefāt-e azaliyye va asmāʾ-e elāhiyye va īshān-and
marāyāʾī ke tamām ḥekāyat mīnamāyand va jamīʿ-e ānche be-īshān rājeʿ ast
fī-l-ḥaqīqe be-ḥażrat-e ẓāher-e mastūr rājeʿ va maʿrefat-e mabdaʾ va
voṣūl be-ū ḥāṣel namīshavad magar be-maʿrefat va voṣūl-e īn kaynūnāt-e
moshreqe az shams-e ḥaqīqat pas az liqāʾ-e īn anvār-e moqaddase liqāʾ-e
Allāh ḥāṣel mīshavad va az ʿelm-eshān ʿelm-e Allāh va az vajh-eshān
vajh-e Allāh va az avvaliyyat va ākhiriyyat va ẓāhiriyyat va
bāṭiniyyat-e īn javāher-e mojarrade thābet mīshavad az barā-ye ān
shams-e ḥaqīqat bi-annahu huwa al-Awwalu wa-l-Ākhiru wa-l-Ẓāhiru
wa-l-Bāṭinu va hamchenīn sāyer-e asmāʾ-e ʿāliye va ṣefāt-e motaʿāliye
le-hādhā har nafsī ke be-īn anvār-e możīʾe-ye momtaneʿe va shomūs-e
moshreqe-ye lāʾeḥe dar har ẓohūr movaffaq va fāʾez shod ū be-liqāʾ-e
Allāh fāʾez ast va dar madīne-ye ḥayāt-e abadiyye-ye bāqiye vāred va īn
liqā moyassar nashavad barā-ye aḥadī illā dar qiyāmat ke qiyām-e nafs-e
Allāh ast be-maẓhar-e kolliyye-ye khod

And if they say that it is the second revelation, which has been termed
the Holy Outpouring, this assuredly belongs to the realm of
creation---that is, to the realm of primal manifestation and original
appearance. This station is reserved for His Prophets and holy ones, for
none greater or more exalted than they has come into being in the realms
of existence, as all acknowledge and affirm. They are the seats and
manifestations of all the eternal attributes and divine names, and they
are the mirrors that reflect Him fully. Everything that pertains to them
pertains in truth to Him who is both manifest and concealed. Knowledge
of the Source and attainment to Him are possible only through knowledge
of, and attainment to, these Beings who shine forth from the Sun of
Truth. Thus, by meeting these holy Lights, one attains to the meeting
with God; through knowledge of them, to knowledge of God; and through
their faces, to the face of God. Through the firstness and lastness, the
outwardness and inwardness of these immaterial Essences, it is
established of that Sun of Truth that ``He is the First and the Last,
the Outward and the Inward,'' and likewise with the other lofty names
and transcendent attributes. Therefore, any soul who, in any revelation,
succeeds in attaining to these shining, inaccessible Lights and these
radiant, resplendent Suns has attained to the meeting with God and
entered the city of everlasting, enduring life. This meeting is possible
for no one except at the Resurrection, which is the arising of God
Himself in His universal Manifestation.

\subsection{153}\label{section-343}

\begin{otherlanguage}{persian}

و اینست معنی قیامت که در کلّ کتب مسطور و مذکور است و جمیع بشارت داده
شده‌اند بآن یوم حال ملاحظه فرمائید که آیا یومی از این یوم عزیزتر و بزرگتر
و معظّم‌تر تصوّر میشود که انسان چنین روز را از دست بگذارد و از فیوضات این
یوم که بمثابهٔ ابر نیسان از قبل رحمن در جریانست خود را محروم نماید و بعد
از آنکه بتمام دلیل مدلّل شد که یومی اعظم از این یوم و امری اعزّ از این امر
نه چگونه میشود که انسان بحرف متوهّمین و ظانّین از چنین فضل اکبر مأیوس گردد
و بعد از همۀ این دلائل محکمهٔ متقنه که هیچ عاقلی را گریزی نه و هیچ عارفی
را مفرّی نه آیا روایت مشهور را نشنیده‌اند که میفرماید اذا قام القائم قامت
القیامة و همچنین ائمّۀ هدی و انوار لاتطفی هل ینظرون الّا ان یأتیهم اللّه فی
ظلل من الغمام را که مسلّماً از امورات محدثۀ در قیامت میدانند بحضرت قائم و
ظهور او تفسیر نموده‌اند

\end{otherlanguage}

Va īn ast maʿnī-ye qiyāmat ke dar koll-e kotob masṭūr va madhkūr ast va
jamīʿ beshārat dāde shode-and be-ān yowm ḥāl molāḥeẓe farmāyīd ke āyā
yowmī az īn yowm ʿazīztar va bozorgtar va moʿaẓẓamtar taṣavvor mīshavad
ke ensān chenīn rūz rā az dast begoẕārad va az foyūżāt-e īn yowm ke
be-mathābe-ye abr-e naysān az qebal-e Raḥmān dar jarayān ast khod rā
maḥrūm namāyad va baʿd az ānke be-tamām-e dalīl modallal shod ke yowmī
aʿẓam az īn yowm va amrī aʿazz az īn amr na chegūne mīshavad ke ensān
be-ḥarf-e motavahhemīn va ẓānnīn az chenīn fażl-e akbar maʾyūs gardad va
baʿd az hame-ye īn dalāʾel-e moḥkame-ye motqane ke hīch ʿāqelī rā gorīzī
na va hīch ʿārefī rā mafarrī na āyā revāyat-e mashhūr rā nashenīde-and
ke mīfarmāyad idhā qāma al-Qāʾimu qāmati al-qiyāmatu va hamchenīn
aʾemme-ye hodā va anvār-e lā tuṭfā hal yanẓurūna illā an yaʾtiyahumu
Allāhu fī ẓulalin mina al-ghamāmi rā ke mosallaman az omūrāt-e
moḥdathe-ye dar qiyāmat mīdānand be-ḥażrat-e Qāʾem va ẓohūr-e ū tafsīr
namūde-and

This is the meaning of resurrection that is recorded and mentioned in
all the Books, and all have been given the glad tidings of that Day. Now
consider: can any day be imagined more precious, greater, or more august
than this Day, that a person should let such a day slip away and deprive
themselves of its outpourings, which stream from the All-Merciful like
April rain? And after it has been fully demonstrated that no day is
greater than this Day and no Cause more glorious than this Cause, how
can anyone, because of the words of those who indulge in fancies and
conjectures, lose hope of such surpassing grace? After all these firm
and conclusive proofs, which no rational or discerning person can evade,
have they not heard the well-known tradition that says, ``When the Qāʾim
arises, the Resurrection has come''? Likewise, the Imams of guidance and
the unquenchable Lights have interpreted ``Are they waiting for anything
but that God should come to them in canopies of clouds''---which they
assuredly regard as one of the events of the Resurrection---as referring
to the Qāʾim and His appearance.

\subsection{154}\label{section-344}

\begin{otherlanguage}{persian}

پس ای برادر معنی قیامت را ادراک نما و گوش را از حرفهای این مردم مردود
پاک فرما اگر قدری بعوالم انقطاع قدم گذاری شهادت میدهید که یومی اعظم از
این یوم و قیامتی اکبر از این قیامت متصوّر نیست و یک عمل در این یوم مقابل
است با اعمال صد‌هزار سنه بلکه استغفراللّه از این تحدید زیرا که مقدّس است
عمل این یوم از جزای محدود و این همج رعاع چون معنی قیامت و لقای الهی را
ادراک ننمودند لهذا از فیض او بالمرّه محجوب ماندند با اینکه مقصود از علم و
زحمات آن وصول و معرفت این مقامست مع‌ذلک همه مشغول بعلوم ظاهره شده‌اند
چنانچه آنی منفکّ نیستند و از جوهر علم و معلوم چشم پوشیده‌اند گویا نمی از
یمّ علم الهی ننوشیدند و بقطره‌ئی از سحاب فیض رحمانی فائز نگشتند

\end{otherlanguage}

Pas ey barādar maʿnī-ye qiyāmat rā edrāk namā va gūsh rā az ḥarf-hā-ye
īn mardom-e mardūd pāk farmā agar qadrī be-ʿavālem-e enqeṭāʿ qadam
goẕārī shahādat mīdahīd ke yowmī aʿẓam az īn yowm va qiyāmatī akbar az
īn qiyāmat motaṣavvar nīst va yek ʿamal dar īn yowm moqābel ast bā
aʿmāl-e ṣad-hezār sane balke astaghfiru Allāh az īn taḥdīd zīrā ke
moqaddas ast ʿamal-e īn yowm az jazā-ye maḥdūd va īn hamaj-e raʿāʿ chon
maʿnī-ye qiyāmat va liqā-ye elāhī rā edrāk nanamūdand le-hādhā az fayż-e
ū bel-marre maḥjūb māndand bā īnke maqṣūd az ʿelm va zaḥamāt-e ān voṣūl
va maʿrefat-e īn maqām ast maʿa-dhālek hame mashghūl be-ʿolūm-e ẓāhere
shode-and chenānche ānī monfakk nīstand va az jowhar-e ʿelm va maʿlūm
chashm pūshīde-and gūyā namī az yamm-e ʿelm-e elāhī nanūshīdand va
be-qaṭre-ī az saḥāb-e fayż-e raḥmānī fāʾez nagashtand

So, my brother, grasp the meaning of resurrection and cleanse your ears
of the words of these rejected people. If you take even a few steps into
the realms of detachment, you will testify that no day greater than this
Day, and no resurrection mightier than this Resurrection, can be
imagined. A single deed in this Day is equivalent to the deeds of a
hundred thousand years---indeed, I ask God's forgiveness for setting
such a limit, for a deed performed in this Day is sanctified beyond any
limited recompense. Because this ignorant rabble has not understood the
meaning of resurrection and meeting God, they have remained entirely
veiled from His grace. Although the purpose of learning and its labors
is to attain and recognize this station, they have nevertheless all
occupied themselves with outward learning, never leaving it for a
moment, and have shut their eyes to the essence of knowledge and its
Object. It is as though they have not tasted the slightest moisture from
the ocean of divine knowledge or received a single drop from the clouds
of the All-Merciful's grace.

\subsection{155}\label{section-345}

\begin{otherlanguage}{persian}

حال ملاحظه فرمائید اگر کسی در یوم ظهور حقّ ادراک فیض لقا و معرفت مظاهر حقّ
را ننماید آیا صدق عالم بر او میشود اگرچه هزار سنه تحصیل کرده باشد و جمیع
علوم محدودهٔ ظاهریّه را اخذ نموده باشد و این بالبدیهه معلومست که تصدیق علم
در حقّ او نمیشود ولیکن اگر نفسی حرفی از علم ندیده باشد و باین شرافت کبری
فائز شود البتّه او از علمای ربّانی محسوبست زیرا بغایت قصوای علم و نهایت و
منتهای آن فائز گشته

\end{otherlanguage}

Ḥāl molāḥeẓe farmāyīd agar kasī dar yowm-e ẓohūr-e Ḥaqq edrāk-e fayż-e
liqā va maʿrefat-e maẓāher-e Ḥaqq rā nanamāyad āyā ṣedq-e ʿālem bar ū
mīshavad agarche hezār sane taḥṣīl karde bāshad va jamīʿ-e ʿolūm-e
maḥdūde-ye ẓāhiriyye rā akhdh namūde bāshad va īn bel-badīhe maʿlūm ast
ke taṣdīq-e ʿelm dar ḥaqq-e ū namīshavad valīkan agar nafsī ḥarfī az
ʿelm nadīde bāshad va be-īn sherāfat-e kobrā fāʾez shavad albatte ū az
ʿolamā-ye rabbānī maḥsūb ast zīrā be-ghāyat-e qoṣvā-ye ʿelm va nehāyat
va montahā-ye ān fāʾez gashte

Now consider: if someone, on the day of God's manifestation, does not
attain the grace of meeting Him and recognizing the Manifestations of
God, can that person truly be called learned, even after a thousand
years of study and the acquisition of all the limited branches of
outward learning? It is self-evident that such a person cannot be
acknowledged as possessing knowledge. But if a soul has never
encountered a single letter of learning and yet attains this supreme
distinction, that soul is assuredly counted among the divinely learned,
having attained the ultimate goal of knowledge, its end and
consummation.

\subsection{156}\label{section-346}

\begin{otherlanguage}{persian}

و این رتبه هم از علائم ظهور است چنانچه میفرماید یجعل اعلاکم اسفلکم و
اسفلکم اعلاکم و همچنین در فرقان میفرماید و نرید ان نمنّ علی‌الّذین استضعفوا
فی الأرض و نجعلهم ائمّة و نجعلهم الوارثین و این مشاهده شد که الیوم چه
مقدار از علما نظر باعراض در اسفل اراضی جهل ساکن شده‌اند و اسامیشان از
دفتر عالین و علما محو شده و چه مقدار از جهّال نظر باقبال باعلی افق علم
ارتقا جستند و اسمشان در الواح علم بقلم قدرت ثبت گشته کذلک یمحو اللّه ما
یشآء و یثبت و عنده امّ الکتاب اینست که گفته‌اند طلب الدّلیل عند حصول
المدلول قبیح و الاشتغال بالعلم بعد الوصول الی المعلوم مذموم قل یا اهل
الأرض هذا فتی ناریّ یرکض فی برّیّة الرّوح و یبشّرکم بسراج‌ اللّه و یذکّرکم
بالأمر الّذی کان عن افق القدس فی شطر العراق تحت حجبات النّور بالسّتر مشهودا

\end{otherlanguage}

Va īn rotbe ham az ʿalāʾem-e ẓohūr ast chenānche mīfarmāyad yajʿalu
aʿlākum asfalakum wa-asfalakum aʿlākum va hamchenīn dar Forqān
mīfarmāyad wa-nurīdu an namunna ʿalā alladhīna stuḍʿifū fī al-arḍi
wa-najʿalahum aʾimmatan wa-najʿalahumu al-wārithīna va īn moshāhede shod
ke al-yowm che meqdār az ʿolamā naẓar be-eʿrāż dar asfal-e arāżī-ye jahl
sāken shode-and va asāmī-shān az daftar-e ʿālīn va ʿolamā maḥv shode va
che meqdār az johhāl naẓar be-eqbāl be-aʿlā ofoq-e ʿelm erteqā jostand
va esm-eshān dar alvāḥ-e ʿelm be-qalam-e qodrat thabt gashte kadhālika
yamḥū Allāhu mā yashāʾu wa-yuthbitu wa-ʿindahu ummu al-kitābi īn ast ke
gofte-and ṭalabu al-dalīli ʿinda ḥuṣūli al-madlūli qabīḥun
wa-l-ishtighālu bi-l-ʿilmi baʿda al-wuṣūli ilā al-maʿlūmi madhmūmun qul
yā ahla al-arḍi hādhā fatan nāriyyun yarkuḍu fī barriyyati al-rūḥi
wa-yubashshirukum bi-sirāji Allāhi wa-yudhakkirukum bi-l-amri alladhī
kāna ʿan ufuqi al-qudsi fī shaṭri al-ʿIrāqi taḥta ḥujubāti al-nūri
bi-l-sitri mashhūdan

This rank is also among the signs of the revelation, as He says, ``He
will make the highest among you the lowest, and the lowest among you the
highest.'' Likewise, He says in the Qur'an, ``We desire to favor those
who have been oppressed in the land, and to make them leaders, and to
make them the inheritors.'' It has been witnessed today how many of the
learned, because they turned away, have come to dwell in the lowest
depths of ignorance, their names erased from the register of the exalted
and the learned; and how many of the unlearned, because they turned
toward Him, have ascended to the highest horizon of knowledge, their
names inscribed by the Pen of power upon the tablets of knowledge. Thus,
``God effaces whatever He wills and confirms, and with Him is the Mother
Book.'' This is why it has been said, ``To seek proof when that which is
proven has been attained is unseemly, and to occupy oneself with
knowledge after reaching the Object of knowledge is blameworthy.'' Say:
People of the earth! This is a fiery Youth who races through the
wilderness of the Spirit, brings you glad tidings of the lamp of God,
and reminds you of the Cause that, from the horizon of holiness in the
region of Iraq, was beheld in concealment beneath the veils of light.

\subsection{157}\label{section-347}

\begin{otherlanguage}{arabic}

ای دوست من اگر قدری در سموات معانی فرقان طیران فرمائی و در ارض معرفت
الهی که در آن مبسوط گشته تفرّج نمائی بسیار از ابواب علوم بر وجه آن جناب
مفتوح شود و خواهید یقین نمود بر اینکه جمیع این امور که الیوم این عباد را
منع مینماید از ورود در شاطی بحر ازلی بعینها در ظهور نقطهٔ فرقان هم مردم
آن عصر را منع نموده از اقرار بآن شمس و اذعان بآن و همچنین بر اسرار رجعت
و بعثت مطّلع شوی و باعلی غرف یقین و اطمینان مقرّ یابی

\end{otherlanguage}

Ey dūst-e man agar qadrī dar samāvāt-e maʿānī-ye Forqān ṭayarān farmāʾī
va dar arż-e maʿrefat-e elāhī ke dar ān mabsūṭ gashte tafarroj namāʾī
besyār az abvāb-e ʿolūm bar vajh-e ān janāb maftūḥ shavad va khāhīd
yaqīn namūd bar īnke jamīʿ-e īn omūr ke al-yowm īn ʿebād rā manʿ
mīnamāyad az vorūd dar shāṭī-ye baḥr-e azalī be-ʿaynehā dar ẓohūr-e
Noqṭe-ye Forqān ham mardom-e ān ʿaṣr rā manʿ namūde az eqrār be-ān shams
va edhʿān be-ān va hamchenīn bar asrār-e rajʿat va baʿthat moṭṭaleʿ
shavī va be-aʿlā ghoraf-e yaqīn va eṭmīnān maqarr yābī

My friend, if you soar a little in the heavens of the Qur'an's meanings
and wander through the land of divine knowledge spread out within it,
many gates of learning will open before you. You will become certain
that all the things that today prevent these servants from reaching the
shore of the eternal Ocean likewise prevented the people of that age, at
the appearance of the Point of the Qur'an, from acknowledging that Sun
and submitting to Him. You will also become acquainted with the
mysteries of return and resurrection and find a dwelling in the loftiest
chambers of certainty and assurance.

\subsection{158}\label{section-348}

\begin{otherlanguage}{persian}

از جمله اینکه روزی جمعی از مجاحدان آن جمال بیمثال و محرومان از کعبهٔ
لایزال از روی استهزاء عرض نمودند انّ اللّه عهد الینا ان لا نؤمن لرسول حتّی
یأتینا بقربان تأکله النّار مضمون آن اینست که پروردگار عهد کرده است بما که
ایمان نیاوریم برسولی مگر آنکه معجزهٔ هابیل و قابیل را ظاهر فرماید یعنی
قربانی کند و آتشی از آسمان بیاید و آن را بسوزاند چنانچه در حکایت هابیل
شنیده‌اند و در کتب مذکور است آن حضرت در جواب فرمودند قد جآءکم رسل من قبلی
بالبیّنات و بالّذی قلتم فلم قتلتموهم ان کنتم صادقین ترجمهٔ آن اینست که آن
حضرت فرمودند آمد بسوی شما پیش از من رسولهای پروردگار با بیّنات ظاهرات و
بآنچه شما میطلبید پس چرا کشتید آن رسل پروردگار را اگر هستید راست‌گویان
حال انصاف دهید بر حسب ظاهر آن عباد که در عصر و عهد آن حضرت بوده‌اند کجا
در عهد آدم یا انبیای دیگر بودند که چند هزار سال فاصله بود از عهد آدم تا
آن زمان مع‌ذلک چرا آن جوهر صدق نسبت قتل هابیل و یا انبیای دیگر را بعباد
زمان خود فرمودند چاره نداری یا اینکه نعوذ باللّه نسبت کذب و یا کلام لغو
بآن حضرت بدهی یا بگوئی آن اشقیا همان اشقیا بودند که در هر عصری با نبیّین
و مرسلین معارضه مینمودند تا آنکه بالأخره همه را شهید نمودند

\end{otherlanguage}

Az jumleh īnkih rūzī jamʿī az mujāḥidān-i ān jamāl-i bī-mithāl va
maḥrūmān az Kaʿbih-yi lā-yazāl az rū-yi istihzā ʿarḍ namūdand: Inna
Allāha ʿahida ilaynā allā nuʾmina li-rasūlin ḥattā yaʾtiyanā bi-qurbānin
taʾkuluhu al-nār. Maḍmūn-i ān īnast kih Parvardigār ʿahd kardih ast
bi-mā kih īmān nayāvarīm bi-rasūlī magar ānkih muʿjizih-yi Hābīl va
Qābīl rā ẓāhir farmāyad, yaʿnī qurbānī kunad va ātashī az āsmān biyāyad
va ān rā bisūzānad, chunānchih dar ḥikāyat-i Hābīl shinīdih-and va dar
kutub madhkūr ast. Ān Ḥaḍrat dar javāb farmūdand: Qad jāʾakum rusulun
min qablī bi-l-bayyināti wa-bi-alladhī qultum fa-lima qataltumūhum in
kuntum ṣādiqīn. Tarjumih-yi ān īnast kih ān Ḥaḍrat farmūdand: Āmad
bi-sū-yi shumā pīsh az man rasūlhā-yi Parvardigār bā bayyināt-i ẓāhirāt
va bi-ānchih shumā mīṭalabīd, pas chirā kushtīd ān rusul-i Parvardigār
rā agar hastīd rāst-gūyān? Ḥāl inṣāf dihīd: bar ḥasab-i ẓāhir ān ʿibād
kih dar ʿaṣr va ʿahd-i ān Ḥaḍrat būdih-and kujā dar ʿahd-i Ādam yā
anbiyā-yi dīgar būdand kih chand hizār sāl fāṣilih būd az ʿahd-i Ādam tā
ān zamān? Maʿa-dhālik chirā ān jawhar-i ṣidq nisbat-i qatl-i Hābīl va yā
anbiyā-yi dīgar rā bi-ʿibād-i zamān-i khud farmūdand? Chārih nadārī yā
īnkih naʿūdhu bi-Allāh nisbat-i kidhb va yā kalām-i laghv bi-ān Ḥaḍrat
bidahī yā bigūʾī ān ashqiyā hamān ashqiyā būdand kih dar har ʿaṣrī bā
nabiyyīn va mursilīn muʿāriḍih mīnamūdand tā ānkih bi-l-ākhirih hamih rā
shahīd namūdand.

Among these instances is that one day a group of those who denied that
peerless Beauty and were deprived of the everlasting Kaaba said
mockingly, ``God has pledged to us that we should not believe in any
messenger until he brings us an offering that fire consumes.'' Its
meaning is that the Lord has pledged to us that we should not believe in
a messenger unless he manifests the miracle of Abel and Cain---that is,
makes an offering, and a fire comes down from heaven and burns it---as
they have heard in the story of Abel and as is recorded in the books. In
reply, He said, ``Messengers came to you before Me with clear proofs and
with what you have asked for. Why, then, did you kill them, if you are
truthful?'' The translation of this is that He said, ``Before Me, the
Lord's messengers came to you with manifest proofs and with what you
seek. Why, then, did you kill those messengers of the Lord, if you are
truthful?'' Now judge fairly: outwardly, how could those people who
lived in His age and time have been present in the days of Adam or the
other Prophets, when several thousand years separated the time of Adam
from that age? Nevertheless, why did that Essence of truth attribute the
murder of Abel or the other Prophets to the people of His own time? You
have no choice but either to ascribe falsehood or idle speech to
Him---God forbid---or to say that those wicked people were the very same
wicked people who in every age opposed the Prophets and Messengers until
they finally put them all to death.

\subsection{159}\label{section-349}

\begin{otherlanguage}{arabic}

درست در این بیان تفکّر فرما تا نسیم خوش عرفان از مصر رحمن بوزد و جان را
از بیان خوش جانان بحدیقهٔ عرفان رساند این بود که مردم غافل چون معانی این
بیانات بالغهٔ کامله را ادراک نمینمودند و جواب را بگمان خود مطابق سؤال
نمی‌یافتند لهذا نسبت عدم علم و جنون بآن جواهر علم و عقل میدادند

\end{otherlanguage}

Durust dar īn bayān tafakkur farmā tā nasīm-i khush-i ʿirfān az Miṣr-i
Raḥmān bivazad va jān rā az bayān-i khush-i Jānān bi-ḥadīqih-yi ʿirfān
rasānad. Īn būd kih mardum-i ghāfil chūn maʿānī-yi īn bayānāt-i
bālighih-yi kāmilih rā idrāk namīnamūdand va javāb rā bi-gumān-i khud
muṭābiq-i suʾāl namīyāftand, li-hādhā nisbat-i ʿadam-i ʿilm va junūn
bi-ān javāhir-i ʿilm va ʿaql mīdādand.

Reflect carefully on this utterance, so that the sweet breeze of
recognition may blow from the Egypt of the All-Merciful and, through the
Beloved's sweet utterance, carry the soul to the garden of knowledge.
Thus, because heedless people did not understand the meanings of these
consummate and perfect utterances, and imagined that the answer did not
correspond to the question, they attributed ignorance and madness to
those Essences of knowledge and reason.

\subsection{160}\label{section-350}

\begin{otherlanguage}{persian}

و همچنین در آیهٔ دیگر میفرماید تعرّضاً باهل زمان و کانوا من قبل یستفتحون
علی‌الّذین کفروا فلمّا جآءهم ما عرفوا کفروا به فلعنة اللّه علی الکافرین
میفرماید بودند این گروه که با کفّار مجاهده و قتال مینمودند در راه خدا و
طلب فتح مینمودند برای نصرت امر اللّه پس چون آمد ایشان را آن کسی که شناخته
بودند کافر شدند باو پس لعنت خدا بر کافران حال ملاحظه فرمائید که از آیه
چنین مستفاد میشود که مردم زمان آن حضرت همان مردمی بودند که در عهد انبیای
قبل برای ترویج آن شریعت و ابلاغ امر اللّه مجادله و محاربه مینمودند و حال
آنکه مردم عهد عیسی و موسی غیر مردم زمان آن حضرت بودند و دیگر آنکه آن کسی
را که از قبل شناخته بودند موسی بود صاحب تورات و عیسی بود صاحب انجیل
مع‌ذلک چرا آن حضرت میفرماید چون آمد بسوی ایشان آن کسی که او را شناخته
بودند که عیسی باشد یا موسی باو کافر شدند و حال آنکه آن حضرت نظر بظاهر
موسوم باسم دیگر بودند که محمّد باشد و از مدینهٔ دیگر ظاهر شدند و بلسان
دیگر و شرع دیگر آمدند مع‌ذلک چگونه حکم آیه ثابت میشود و ادراک میگردد

\end{otherlanguage}

Va hamchunīn dar āyih-yi dīgar mīfarmāyad taʿarruḍan bi-ahl-i zamān:
Wa-kānū min qablu yastaftiḥūna ʿalā alladhīna kafarū fa-lammā jāʾahum mā
ʿarafū kafarū bihi fa-laʿnatu Allāhi ʿalā al-kāfirīn. Mīfarmāyad: Būdand
īn gurūh kih bā kuffār mujāhadih va qitāl mīnamūdand dar rāh-i Khudā va
ṭalab-i fatḥ mīnamūdand barā-yi nuṣrat-i Amruʾllāh, pas chūn āmad īshān
rā ān kasī kih shinākhtih būdand kāfir shudand bi-ū, pas laʿnat-i Khudā
bar kāfirān. Ḥāl mulāḥiẓih farmāʾīd kih az āyih chunīn mustafād mīshavad
kih mardum-i zamān-i ān Ḥaḍrat hamān mardumī būdand kih dar ʿahd-i
anbiyā-yi qabl barā-yi tarvīj-i ān sharīʿat va iblāgh-i Amruʾllāh
mujādilih va muḥāribih mīnamūdand, va ḥāl ānkih mardum-i ʿahd-i ʿĪsā va
Mūsā ghayr-i mardum-i zamān-i ān Ḥaḍrat būdand. Va dīgar ānkih ān kasī
rā kih az qabl shinākhtih būdand Mūsā būd ṣāḥib-i Tawrāt va ʿĪsā būd
ṣāḥib-i Injīl. Maʿa-dhālik chirā ān Ḥaḍrat mīfarmāyad chūn āmad bi-sū-yi
īshān ān kasī kih ū rā shinākhtih būdand kih ʿĪsā bāshad yā Mūsā bi-ū
kāfir shudand, va ḥāl ānkih ān Ḥaḍrat naẓar bi-ẓāhir mawsūm bi-ism-i
dīgar būdand kih Muḥammad bāshad va az madīnih-yi dīgar ẓāhir shudand va
bi-lisān-i dīgar va sharʿ-i dīgar āmadand? Maʿa-dhālik chigūnih ḥukm-i
āyih thābit mīshavad va idrāk mīgardad?

Likewise, in another verse He says, rebuking the people of His time,
``Previously they had sought victory over those who disbelieved; but
when there came to them what they recognized, they disbelieved in it. So
God's curse is upon the unbelievers.'' He says: These people used to
strive and fight against the unbelievers in the path of God and seek
victory in support of the Cause of God. But when the One they had
recognized came to them, they disbelieved in Him. So God's curse is upon
the unbelievers. Now consider how this verse implies that the people of
His time were the very same people who, in the days of the former
Prophets, had struggled and fought to promote that religious law and
proclaim the Cause of God, whereas the people of the days of Jesus and
Moses were different from the people of His time. Moreover, the one they
had previously recognized was Moses, the bearer of the Torah, and Jesus,
the bearer of the Gospel. Nevertheless, why does He say that when the
One they had recognized---that is, Jesus or Moses---came to them, they
disbelieved in Him, even though He was outwardly known by another name,
Muḥammad, appeared from another city, and came with another language and
another religious law? How, then, is the statement of this verse
established and understood?

\subsection{161}\label{section-351}

\begin{otherlanguage}{persian}

حال حکم رجوع را ادراک فرما که بچه صریحی در خود فرقان نازل شده و احدی تا
الیوم ادراک آن ننموده حال چه میفرمائید اگر میفرمائید که آن حضرت رجعت
انبیای قبل بودند چنانچه از آیه مستفاد میشود و همچنین اصحاب او هم رجعت
اصحاب قبل خواهند بود چنانچه از آیات مذکوره هم رجعت عباد قبل واضح و
لائحست و اگر انکار کنند بر خلاف حکم کتاب که حجّت اکبر است قائل شده‌اند پس
همین قسم حکم رجع و بعث و حشر را در ایّام ظهور مظاهر هویّه ادراک نما تا
رجوع ارواح مقدّسه را در اجساد صافیهٔ منیره بعین رأس ملاحظه فرمائی و
غبارهای جهل و نفس ظلمانی را بآب رحمت علم رحمانی پاک و منزّه نمائی که شاید
بقوّت یزدانی و هدایت سبحانی و سراج نورانی سبیل صبح هدایت را از شام ضلالت
تمیز دهی و فرق گذاری

\end{otherlanguage}

Ḥāl ḥukm-i rujūʿ rā idrāk farmā kih bi-chih ṣarīḥī dar khud-i Furqān
nāzil shudih va aḥadī tā al-yawm idrāk-i ān nanamūdih. Ḥāl chih
mīfarmāʾīd? Agar mīfarmāʾīd kih ān Ḥaḍrat rajʿat-i anbiyā-yi qabl
būdand, chunānchih az āyih mustafād mīshavad, va hamchunīn aṣḥāb-i ū ham
rajʿat-i aṣḥāb-i qabl khwāhand būd, chunānchih az āyāt-i madhkūrih ham
rajʿat-i ʿibād-i qabl vāḍiḥ va lāʾiḥast. Va agar inkār kunand bar
khilāf-i ḥukm-i Kitāb kih ḥujjat-i akbar ast qāʾil shudih-and. Pas hamīn
qism ḥukm-i rajʿ va baʿth va ḥashr rā dar ayyām-i ẓuhūr-i maẓāhir-i
Huviyyih idrāk namā tā rujūʿ-i arvāḥ-i muqaddasih rā dar ajsād-i
ṣāfiyih-yi munīrih bi-ʿayn-i raʾs mulāḥiẓih farmāʾī va ghubārhā-yi jahl
va nafs-i ẓulmānī rā bi-āb-i raḥmat-i ʿilm-i Raḥmānī pāk va munazzah
namāʾī kih shāyad bi-quvvat-i Yazdānī va hidāyat-i Subḥānī va sirāj-i
nūrānī sabīl-i ṣubḥ-i hidāyat rā az shām-i ḍalālat tamyīz dihī va farq
gudhārī.

Now understand how explicitly the truth of return has been revealed in
the Qur'an itself, though no one to this day has understood it. What do
you now say? If you say that He was the return of the former Prophets,
as the verse implies, then His companions must likewise be the return of
the former companions, as the return of the people of former times is
also clear and evident from the verses cited. But if they deny this,
they have asserted something contrary to the pronouncement of the Book,
which is the greatest proof. In the same way, then, understand the truth
of return, resurrection, and gathering in the days when the
Manifestations of the divine Essence appear, so that you may behold with
your own eyes the return of sanctified spirits in pure and radiant
bodies, and wash away the dust of ignorance and the darkened self with
the water of the mercy of divine knowledge. Perhaps, through divine
power, heavenly guidance, and a radiant lamp, you may distinguish and
tell apart the path of the morning of guidance from the evening of
error.

\subsection{162}\label{section-352}

\begin{otherlanguage}{persian}

و دیگر معلوم آن جناب بوده که حاملان امانت احدیّه که در عوالم ملکیّه بحکم
جدید و امر بدیع ظاهر میشوند چون این اطیار عرش باقی از سماء مشیّت الهی
نازل میگردند و جمیع بر امر مبرم ربّانی قیام میفرمایند لهذا حکم یک نفس و
یک ذات را دارند چه جمیع از کأس محبّت الهی شاربند و از اثمار شجرهٔ توحید
مرزوق و این مظاهر حقّ را دو مقام مقرّر است یکی مقام صرف تجرید و جوهر تفرید
و در این مقام اگر کل را بیک اسم و رسم موسوم و موصوف نمائی بأسی نیست
چنانچه میفرماید لا نفرّق بین احد من رسله زیرا که جمیع مردم را بتوحید الهی
دعوت میفرمایند و بکوثر فیض و فضل نامتناهی بشارت میدهند و کل بخلع نبوّت
فائزند و برداء مکرمت مفتخر اینست که نقطهٔ فرقان میفرماید امّا النّبیّون فأنا
و همچنین میفرماید منم آدم اوّل و نوح و موسی و عیسی و همین مضمون را طلعت
علوی هم فرموده‌اند و امثال این بیانات که مشعر بر توحید آن مواقع تجرید است
از مجاری بیانات ازلیّه و مخازن لآلی علمیّه ظاهر شده و در کتب مذکور گشته و
این طلعات مواقع حکم و مطالع امرند و امر مقدّس از حجبات کثرت و عوارضات
تعدّد است اینست که میفرماید و ما امرنا الّا واحدة و چون امر واحد شد البتّه
مظاهر امر هم واحدند و همچنین ائمّهٔ دین و سراجهای یقین فرمودند اوّلنا محمّد
و آخرنا محمّد و اوسطنا محمّد

\end{otherlanguage}

Va dīgar maʿlūm-i ān janāb būdih kih ḥāmilān-i amānat-i Aḥadiyyih kih
dar ʿavālim-i mulkiyyih bi-ḥukm-i jadīd va amr-i badīʿ ẓāhir mīshavand,
chūn īn aṭyār-i ʿarsh-i bāqī az samāʾ-i mashiyyat-i Ilāhī nāzil
mīgardand va jamīʿ bar amr-i mubram-i Rabbānī qiyām mīfarmāyand,
li-hādhā ḥukm-i yik nafs va yik dhāt rā dārand, chih jamīʿ az kaʾs-i
maḥabbat-i Ilāhī shāriband va az athmār-i shajarih-yi tawḥīd marzūq. Va
īn maẓāhir-i Ḥaqq rā du maqām muqarrar ast: yikī maqām-i ṣirf-i tajrīd
va jawhar-i tafrīd, va dar īn maqām agar kull rā bi-yik ism va rasm
mawsūm va mawṣūf namāʾī baʾsī nīst, chunānchih mīfarmāyad: Lā nufarriqu
bayna aḥadin min rusulihi. Zīrā kih jamīʿ mardum rā bi-tawḥīd-i Ilāhī
daʿvat mīfarmāyand va bi-kawthar-i fayḍ va faḍl-i nāmutanāhī bishārat
mīdihand va kull bi-khilaʿ-i nubuwwat fāʾizand va bi-ridā-yi makrumat
muftakhir. Īnast kih Nuqṭih-yi Furqān mīfarmāyad: Ammā al-nabiyyūna
fa-anā. Va hamchunīn mīfarmāyad: Manam Ādam-i avval va Nūḥ va Mūsā va
ʿĪsā. Va hamīn maḍmūn rā Ṭalʿat-i ʿAlavī ham farmūdih-and. Va amthāl-i
īn bayānāt kih mushʿir bar tawḥīd-i ān mavāqiʿ-i tajrīd ast az majārī-yi
bayānāt-i azaliyyih va makhāzin-i laʾālī-yi ʿilmiyyih ẓāhir shudih va
dar kutub madhkūr gashtih. Va īn ṭalʿāt mavāqiʿ-i ḥukm va maṭāliʿ-i
amrand va amr muqaddas az ḥujubāt-i kathrat va ʿavāriḍāt-i taʿaddud ast.
Īnast kih mīfarmāyad: Wa-mā amrunā illā wāḥidah. Va chūn amr vāḥid shud
al-battih maẓāhir-i amr ham vāḥidand. Va hamchunīn aʾimmih-yi dīn va
sirājhā-yi yaqīn farmūdand: Awwalunā Muḥammad wa-ākhirunā Muḥammad
wa-awsaṭunā Muḥammad.

Furthermore, you should know that the bearers of the trust of divine
Unity, who appear in the earthly realms with a new decree and a wondrous
Cause, are regarded as one soul and one essence, since these birds of
the everlasting Throne descend from the heaven of the divine will and
all arise to uphold God's irrevocable Cause. For all drink from the cup
of divine love and are nourished by the fruits of the tree of oneness.
Two stations are assigned to these Manifestations of God. One is the
station of pure abstraction and the essence of singleness. In this
station, there is no harm in calling and describing them all by one name
and designation, as He says, ``We make no distinction between any of His
messengers.'' For they all summon people to divine unity and bring glad
tidings of the river of limitless grace and bounty. All are invested
with the robes of prophethood and honored with the mantle of glory. Thus
the Point of the Qur'an says, ``As for the Prophets, I am they.''
Likewise, He says, ``I am the first Adam, Noah, Moses, and Jesus.'' The
Countenance of `Alí has also expressed this same meaning. Utterances
such as these, which indicate the oneness of those embodiments of
abstraction, have issued from the channels of eternal utterance and the
treasuries of the pearls of knowledge, and have been recorded in the
books. These Countenances are the seats of divine decree and the dawning
places of the Cause, and the Cause is sanctified above the veils of
multiplicity and the contingencies of plurality. Thus He says, ``Our
Cause is but one.'' Since the Cause is one, its Manifestations are
necessarily also one. Likewise, the Imams of religion, those lamps of
certitude, have said, ``Our first is Muḥammad, our last is Muḥammad, and
our middle is Muḥammad.''

\subsection{163}\label{section-353}

\begin{otherlanguage}{persian}

باری معلوم و محقّق آن جناب بوده که جمیع انبیا هیاکل امر اللّه هستند که در
قمایص مختلفه ظاهر شدند و اگر بنظر لطیف ملاحظه فرمائی همه را در یک رضوان
ساکن بینی و در یک هوا طائر و بر یک بساط جالس و بر یک کلام ناطق و بر یک
امر آمر اینست اتّحاد آن جواهر وجود و شموس غیر محدود و معدود پس اگر یکی از
این مظاهر قدسیّه بفرماید من رجوع کلّ انبیا هستم صادقست و همچنین ثابت است
در هر ظهور بعد صدق رجوع ظهور قبل و چون رجوع انبیا موافق و مطابق آیات و
اخبار ثابت شد رجوع اولیا هم ثابت و محقّق است و این رجوع اظهر از آن است که
بدلیل و برهان محتاج شود مثلاً ملاحظه فرمائید از جملهٔ انبیا نوح بود که چون
مبعوث بنبوّت شد و بقیام الهی بر امر قیام فرمود هر نفسی که باو مؤمن و بامر
او مذعن شد او فی‌الحقیقه بحیات جدیده مشرّف شد و در حقّ او صادق میآمد حیات
بدیع و روح جدید زیرا که او قبل از ایمان بخدا و اذعان بمظهر نفس او کمال
علایق را باموال و اسباب متعلّقهٔ بدنیا از قبیل زن و فرزند و اطعمه و اشربه
و امثال ذلک داشته بقسمی که اوقات لیل و نهار را مصروف بر اخذ زخارف و
اسباب تعیّش داشته و همّت در تحصیل اشیای فانیه گماشته و از این مراتب گذشته
قبل از ورود در لجّهٔ ایمان بحدودات آباء و اجداد و اتّباع آداب و شرایع ایشان
چنان راسخ و محکم بود که اگر حکم بقتل او میشد شاید رضا میداد و راضی بر
تغییر حرفی از امور تقلیدیّه که در میان قوم بود نمیشد چنانچه همهٔ قوم نداء
انّا وجدنا آبائنا علی امّة و انّا علی آثارهم مقتدون برآوردند

\end{otherlanguage}

Bārī maʿlūm va muḥaqqaq-i ān janāb būdih kih jamīʿ-i anbiyā hayākil-i
Amruʾllāh hastand kih dar qamāyiṣ-i mukhtalifih ẓāhir shudand, va agar
bi-naẓar-i laṭīf mulāḥiẓih farmāʾī hamih rā dar yik Riḍvān sākin bīnī va
dar yik havā ṭāʾir va bar yik bisāṭ jālis va bar yik kalām nāṭiq va bar
yik amr āmir. Īnast ittiḥād-i ān javāhir-i vujūd va shumūs-i ghayr-i
maḥdūd va maʿdūd. Pas agar yikī az īn maẓāhir-i qudsiyyih bifarmāyad man
rujūʿ-i kull-i anbiyā hastam ṣādiqast. Va hamchunīn thābit ast dar har
ẓuhūr-i baʿd ṣidq-i rujūʿ-i ẓuhūr-i qabl. Va chūn rujūʿ-i anbiyā muvāfiq
va muṭābiq-i āyāt va akhbār thābit shud rujūʿ-i awliyā ham thābit va
muḥaqqaq ast. Va īn rujūʿ aẓhar az ān ast kih bi-dalīl va burhān muḥtāj
shavad. Mathalan mulāḥiẓih farmāʾīd az jumlih-yi anbiyā Nūḥ būd kih chūn
mabʿūth bi-nubuwwat shud va bi-qiyām-i Ilāhī bar amr qiyām farmūd har
nafsī kih bi-ū muʾmin va bi-amr-i ū mudhʿin shud ū fī al-ḥaqīqih
bi-ḥayāt-i jadīdih musharraf shud va dar ḥaqq-i ū ṣādiq mīāmad ḥayāt-i
badīʿ va rūḥ-i jadīd. Zīrā kih ū qabl az īmān bi-Khudā va idhʿān
bi-maẓhar-i nafs-i ū kamāl-i ʿalāyiq rā bi-amvāl va asbāb-i mutaʿalliqih
bi-dunyā az qabīl-i zan va farzand va aṭʿimih va ashribih va amthāl-i
dhālik dāshtih, bi-qismī kih awqāt-i layl va nahār rā maṣrūf bar akhdh-i
zakhārif va asbāb-i taʿayyush dāshtih va himmat dar taḥṣīl-i ashyā-yi
fāniyih gumāshtih. Va az īn marātib gudhashtih qabl az vurūd dar
lujjih-yi īmān bi-ḥudūdāt-i ābā va ajdād va ittibāʿ-i ādāb va sharāyiʿ-i
īshān chunān rāsikh va muḥkam būd kih agar ḥukm bi-qatl-i ū mīshud
shāyad riḍā mīdād va rāḍī bar taghyīr-i ḥarfī az umūr-i taqlīdiyyih kih
dar miyān-i qawm būd namīshud. Chunānchih hamih-yi qawm nidā-yi innā
wajadnā ābāʾanā ʿalā ummatin wa-innā ʿalā āthārihim muqtadūn bar-āvur
dand.

In short, you should know and be certain that all the Prophets are the
temples of the Cause of God, who have appeared in different garments. If
you look with discerning vision, you will see them all dwelling in one
Paradise, flying in one sky, seated upon one carpet, speaking one word,
and enjoining one Cause. Such is the unity of those Essences of being,
those Suns beyond limitation and number. Thus, if one of these holy
Manifestations should say, ``I am the return of all the Prophets,'' He
speaks the truth. Likewise, in every subsequent revelation, the truth of
the return of the previous revelation is established. Since the return
of the Prophets has been established in accordance with the verses and
traditions, the return of their chosen ones is also established and
certain. This return is too evident to need proof or demonstration.
Consider, for example, Noah, who was among the Prophets. When He was
raised to prophethood and arose with divine resolve to uphold the Cause,
every soul who believed in Him and acknowledged His Cause was truly
honored with a new life, and a wondrous life and a new spirit could
truly be attributed to that person. For before believing in God and
acknowledging the Manifestation of His own Self, that person had been
utterly attached to possessions and worldly concerns, such as wife and
children, food and drink, and the like, to such a degree that he spent
his days and nights acquiring worldly ornaments and the means of
livelihood, and devoted his efforts to obtaining perishable things.
Beyond all this, before entering the ocean of faith, he was so firmly
and steadfastly committed to the precepts of his forefathers and to
following their customs and religious laws that he might have consented
to being put to death, but would not have consented to changing a single
letter of the inherited practices current among his people. Thus all the
people raised the cry, ``We found our fathers following a way, and we
follow in their footsteps.''

\subsection{164}\label{section-354}

\begin{otherlanguage}{persian}

و همین گروه با همهٔ این حجبات محدوده و حدودات مذکوره بمجرّد اینکه صهبای
ایمان را از کأس ایقان از ایادی مظاهر سبحان مینوشیدند بالمرّه تقلیب میشدند
بقسمی که از زن و فرزند و اموال و اثقال و جان و ایمان بلکه از کلّ ماسوی
میگذشتند و بقسمی غلبات شوق الهی و جذبات ذوق صمدانی ایشان را اخذ مینمود
که دنیا را و آنچه در آن هست بپر کاهی ندانسته آیا حکم خلق جدید و رجوع در
اینها نمیشود و همچنین ملاحظه شد که این نفوس قبل از فوز بعنایت بدیع جدید
الهی جان خود را بصد‌هزار حیله و تدبیر از موارد هلاکت حفظ مینمودند بقسمی
که از خاری احتراز میجستند و از روباهی فی‌المثل فرار مینمودند و بعد از شرف
بفوز اکبر و عنایت عظمی صدهزار جان رایگان انفاق میفرمودند بلکه نفوس
مقدّسشان از قفس تن بیزار و یک نفر از این جنود در مقابل گروهی مقاتله
مینمود مع‌ذلک چگونه میشود که اگر این نفوس همان نفوس قبل باشند اینگونه
امورات که مخالف عادات بشریّه و منافی هوای جسمانیّه است از ایشان ظاهر شود

\end{otherlanguage}

Va hamīn gurūh bā hamih-yi īn ḥujubāt-i maḥdūdih va ḥudūdāt-i madhkūrih
bi-mujarrad-i īnkih ṣahbā-yi īmān rā az kaʾs-i īqān az ayādī-yi
maẓāhir-i Subḥān mīnūshīdand bi-l-marrih taqlīb mīshudand, bi-qismī kih
az zan va farzand va amvāl va athqāl va jān va īmān balkih az kull-i
mā-sivā mīgudhashtand. Va bi-qismī ghalabāt-i shawq-i Ilāhī va
jadhabāt-i dhawq-i Ṣamadānī īshān rā akhdh mīnamūd kih dunyā rā va
ānchih dar ān hast bi-par-i kāhī nadānistih. Āyā ḥukm-i khalq-i jadīd va
rujūʿ dar īnhā namīshavad? Va hamchunīn mulāḥiẓih shud kih īn nufūs qabl
az fawz bi-ʿināyat-i badīʿ-i jadīd-i Ilāhī jān-i khud rā bi-ṣad-hizār
ḥīlih va tadbīr az mavārid-i halākat ḥifẓ mīnamūdand, bi-qismī kih az
khārī iḥtirāz mījustand va az rūbāhī fī al-mathal farār mīnamūdand. Va
baʿd az sharaf bi-fawz-i akbar va ʿināyat-i ʿuẓmā ṣad-hizār jān rāygān
infāq mīfarmūdand. Balkih nufūs-i muqaddasishān az qafas-i tan bīzār va
yik nafar az īn junūd dar muqābil-i gurūhī muqātilih mīnamūd.
Maʿa-dhālik chigūnih mīshavad kih agar īn nufūs hamān nufūs-i qabl
bāshand īngūnih umūrāt kih mukhālif-i ʿādāt-i bashariyyih va munāfī-yi
havā-yi jismāniyyih ast az īshān ẓāhir shavad?

Yet these very people, despite all these restricting veils and the
limitations already mentioned, were completely transformed as soon as
they drank the wine of faith from the cup of certitude at the hands of
the Manifestations of the All-Glorious. They would relinquish wife and
children, possessions and belongings, life and faith---indeed,
everything besides God. The surging power of divine longing and the
compelling attractions of heavenly delight so seized them that they
considered the world and everything in it worth less than a wisp of
straw. Does not the designation of a new creation and return apply to
them? Likewise, it was observed that before attaining God's wondrous new
favor, these souls would protect their lives from mortal danger by a
hundred thousand stratagems and devices, to such a degree that they
would avoid a thorn and flee, for example, from a fox. But after being
honored with the greatest attainment and the supreme favor, they would
freely give away a hundred thousand lives. Indeed, their sanctified
souls were weary of the cage of the body, and a single soldier among
them would fight against a multitude. Nevertheless, if these souls were
the same as before, how could such deeds, contrary to human habits and
incompatible with bodily desires, proceed from them?

\subsection{165}\label{section-355}

\begin{otherlanguage}{persian}

باری این مطلب واضحست که بدون تغییر و تبدیل الهی محالست این قسم آثار و
افعال که بهیچوجه شباهت بآثار و افعال قبل ندارد از ایشان ظاهر شود و در
عرصهٔ کون بوجود آید چنانچه اضطرابشان باطمینان تبدیل میشد و ظنّ بیقین تغییر
مییافت و خوف بجرئت مبادله میگشت اینست شأن اکسیر الهی که در یک حین عباد
را تقلیب میفرماید

\end{otherlanguage}

Bārī īn maṭlab vāżeḥ ast ke bedūn-e taghyīr va tabdīl-e ilāhī moḥāl ast
īn qesm āthār va afʿāl ke be-hīch vajh shabāhat be-āthār va afʿāl-e qabl
nadārad az īshān ẓāher shavad va dar ʿarṣe-ye kawn be-vojūd āyad.
Chenānche eżṭerābeshān be-eṭmīnān tabdīl mīshod va ẓann be-yaqīn taghyīr
mīyāft va khawf be-jorʾat mobādele mīgasht. Īn ast shaʾn-e eksīr-e ilāhī
ke dar yek ḥīn ʿebād rā taqlīb mīfarmāyad.

In short, it is clear that without divine transformation and change, it
would be impossible for such deeds and actions, bearing no resemblance
whatsoever to their former deeds and actions, to proceed from them and
come into being in the world of existence. Thus their agitation was
transformed into assurance, their doubt into certainty, and their fear
into courage. Such is the power of the divine elixir, which transforms
people in a single moment.

\subsection{166}\label{section-356}

\begin{otherlanguage}{persian}

مثلاً در مادّهٔ نحاسی ملاحظه فرمائید که اگر در معدن خود از غلبهٔ یبوست محفوظ
بماند در مدّت هفتاد سنه بمقام ذهبی میرسد اگرچه بعضی خود نحاس را ذهب
میدانند که بواسطهٔ غلبهٔ یبوست مریض شده و بمقام خود نرسیده

\end{otherlanguage}

Mathalan dar mādde-ye noḥāsī molāḥeẓe farmāyīd ke agar dar maʿdan-e khod
az ghalabe-ye yobūsat maḥfūẓ bemānad dar moddat-e haftād sane be-maqām-e
dhahabī mīrasad. Agarche baʿżī khod-e noḥās rā dhahab mīdānand ke
be-vāseṭe-ye ghalabe-ye yobūsat marīż shode va be-maqām-e khod narasīde.

Consider, for example, the substance of copper: if it remains in its
mine, protected from an excess of dryness, it attains the state of gold
within seventy years. Some, however, regard copper itself as gold that
has become diseased through an excess of dryness and has not attained
its proper state.

\subsection{167}\label{section-357}

\begin{otherlanguage}{persian}

باری در هر حال اکسیر کامل مادّهٔ نحاسی را در آنی بمقام ذهبی میرساند و
منازل هفتادساله را بآنی طیّ نماید آیا آن ذهب را بعد میتوان گفت که نحاسست
و یا بعالم ذهبی نرسیده و حال آنکه محک در میان است و صفات ذهبی را از
نحاسی معیّن و واضح مینماید

\end{otherlanguage}

Bārī dar har ḥāl eksīr-e kāmel mādde-ye noḥāsī rā dar ānī be-maqām-e
dhahabī mīrasānad va manāzel-e haftād-sāle rā be-ānī ṭeyy namāyad. Āyā
ān dhahab rā baʿd mītavān goft ke noḥās ast va yā be-ʿālam-e dhahabī
narasīde va ḥāl ānke maḥakk dar miyān ast va ṣefāt-e dhahabī rā az
noḥāsī moʿayyan va vāżeḥ mīnamāyad?

In any case, the perfect elixir brings the substance of copper to the
state of gold in an instant, traversing in a moment the stages of
seventy years. Can that gold then be called copper, or said not to have
attained the state of gold, when the touchstone is at hand and clearly
distinguishes the qualities of gold from those of copper?

\subsection{168}\label{section-358}

\begin{otherlanguage}{persian}

همچنین این نفوس هم از اکسیر الهی در آنی عالم ترابی را طیّ نموده بعوالم
قدسی قدم گذارند و بقدمی از مکان محدود بلامکان الهی واصل شوند جهدی باید
تا باین اکسیر فائز شوی که در یک آن مغرب جهل را بمشرق علم رساند و ظلمت
لیل ظلمانی را بصبح نورانی فائز گرداند و بعید صحرای ظنّ را بچشمهٔ قرب و
یقین دلالت کند و هیاکل فانیه را برضوان باقی مشرّف فرماید حال اگر در حقّ
این ذهب حکم نحاسی صادق میآید در حقّ این عباد هم حکم عباد قبل از فوز
بایمان صادق و محقّق است

\end{otherlanguage}

Hamchenīn īn nofūs ham az eksīr-e ilāhī dar ānī ʿālam-e torābī rā ṭeyy
namūde be-ʿavālem-e qodsī qadam godhārand va be-qadamī az makān-e maḥdūd
be-lāmakān-e ilāhī vāṣel shavand. Jahdī bāyad tā be-īn eksīr fāʾez shavī
ke dar yek ān maghreb-e jahl rā be-mashreq-e ʿelm rasānad va ẓolmat-e
leyl-e ẓolmānī rā be-ṣobḥ-e nūrānī fāʾez gardānad va baʿīd-e ṣaḥrā-ye
ẓann rā be-cheshme-ye qorb va yaqīn delālat konad va hayākel-e fāniye rā
be-reżvān-e bāqī mosharraf farmāyad. Ḥāl agar dar ḥaqq-e īn dhahab
ḥokm-e noḥāsī ṣādeq mīāyad dar ḥaqq-e īn ʿebād ham ḥokm-e ʿebād-e qabl
az fawz be-īmān ṣādeq va moḥaqqaq ast.

Likewise, through the divine elixir, these souls traverse the earthly
world in an instant and enter the realms of holiness; in a single step
they pass from limited space into the divine realm beyond space. An
effort is needed for you to attain this elixir, which in a single moment
brings the west of ignorance to the east of knowledge, leads the
darkness of a gloomy night to a radiant dawn, guides those far away in
the wilderness of doubt to the spring of nearness and certainty, and
honors mortal forms with admission to the everlasting paradise. Now, if
the designation of copper holds true of this gold, then the designation
of people as they were before attaining faith also truly and certainly
holds true of these servants.

\subsection{169}\label{section-359}

\begin{otherlanguage}{persian}

ای برادر از این بیانات شافیهٔ کافیهٔ وافیه اسرار خلق جدید و رجوع و بعث بی‌
حجاب و نقاب ظاهر و هویداست انشآءاللّه بتأییدات غیبیّه جامهٔ کهنه را از جسم
و جان دور کنی و بخلع جدیدهٔ باقیه مفتخر گردی

\end{otherlanguage}

Ey barādar, az īn bayānāt-e shāfiye-ye kāfiye-ye vāfiye asrār-e khalq-e
jadīd va rojūʿ va baʿth bī-ḥejāb va neqāb ẓāher va hoveydāst. In shāʾ
Allāh be-taʾyīdāt-e gheybiyye jāme-ye kohne rā az jesm va jān dūr konī
va be-khelaʿ-e jadīde-ye bāqiye moftakhar gardī.

O brother, these healing, sufficient, and comprehensive explanations
reveal the mysteries of a new creation, return, and resurrection,
without veil or concealment. God willing, through unseen assistance, you
will strip the old garment from body and soul and be honored with new
and enduring robes.

\subsection{170}\label{section-360}

\begin{otherlanguage}{persian}

اینست که در هر ظهور بعد انفسی که سبقت یافتند بایمان از کلّ من علی الأرض و
شربت زلال معرفت را از جمال احدیّت نوشیدند و باعلی معارج ایمان و ایقان و
انقطاع ارتفاع جستند حکم رجوع انفس قبل که در ظهور قبل باین مراتب فائز
شده‌اند بر این اصحاب ظهور بعد میشود اسماً و رسماً و فعلاً و قولاً و امراً زیرا
آنچه از عباد قبل ظاهر شد از این عباد بعد بعینه ظاهر و هویدا گشت مثلاً اگر
شاخسار گلی در مشرق ارض باشد و در مغرب هم از شاخۀ دیگر آن گل ظاهر شود
اطلاق گل بر او میشود دیگر در این مقام نظر بحدودات شاخه و هیئت آن نیست
بلکه نظر برائحه و عطریست که در هر دو ظاهر است

\end{otherlanguage}

Īn ast ke dar har ẓohūr-e baʿd anfosī ke sabqat yāftand be-īmān az
koll-e man ʿalā al-arḍ va sharbat-e zolāl-e maʿrefat rā az jamāl-e
aḥadiyyat nūshīdand va be-aʿlā maʿārej-e īmān va īqān va enqeṭāʿ ertefāʿ
jostand ḥokm-e rojūʿ-e anfos-e qabl ke dar ẓohūr-e qabl be-īn marāteb
fāʾez shode-and bar īn aṣḥāb-e ẓohūr-e baʿd mīshavad, isman va rasman va
fiʿlan va qawlan va amran. Zīrā ānche az ʿebād-e qabl ẓāher shod az īn
ʿebād-e baʿd be-ʿeyne ẓāher va hoveydā gasht. Mathalan agar shākhsār-e
golī dar mashreq-e arḍ bāshad va dar maghreb ham az shākhe-ye dīgar ān
gol ẓāher shavad eṭlāq-e gol bar ū mīshavad. Dīgar dar īn maqām naẓar
be-ḥodūdāt-e shākhe va heyʾat-e ān nīst balke naẓar be-rāʾeḥe va
ʿeṭrī-st ke dar har do ẓāher ast.

Thus, in every subsequent revelation, those souls who embrace faith
before all others on earth, drink the clear draught of recognition from
the Beauty of Oneness, and ascend to the highest summits of faith,
certitude, and detachment are regarded as the return of those souls who
attained these ranks in the preceding revelation. This designation
applies to these companions of the later revelation in name, character,
deed, word, and purpose. For whatever appeared in the former servants
has appeared and become manifest, in its very likeness, in these later
servants. For example, if a rosebush stands in the east of the earth,
and in the west the same rose appears on another branch, the name of
rose applies to it. Here one looks not at the limitations of the branch
or its form, but rather at the fragrance and perfume manifest in both.

\subsection{171}\label{section-361}

\begin{otherlanguage}{persian}

پس نظر را از حدودات ظاهره طاهر و منزّه کن تا همه را بیک اسم و یک رسم و یک
ذات و یک حقیقت مشاهده نمائی و اسرار رجوع کلمات را هم در حروفات نازله
ملاحظه فرمائی قدری تفکّر در اصحاب عهد نقطهٔ فرقان نما که چگونه از جمیع
جهات بشریّه و مشتهیات نفسیّه بنفحات قدسیّهٔ آن حضرت پاک و مقدّس و منقطع گشتند
و قبل از همهٔ اهل ارض بشرف لقا که عین لقآء اللّه بود فائز شدند و از کلّ اهل
ارض منقطع گشتند چنانچه شنیده‌اید که در مقابل آن مظهر ذی الجلال چگونه جان
نثار میفرمودند و حال همان ثبوت و رسوخ و انقطاع را بعینه ملاحظه فرما در
اصحاب نقطهٔ بیان راجع شده چنانچه ملاحظه فرموده‌اید که چگونه این اصحاب از
بدایع جود ربّ الأرباب علم انقطاع بر رفرف امتناع برافراشتند باری این انوار
از یک مصباح ظاهر شده‌اند و این اثمار از یک شجره روئیده‌اند فی‌الحقیقه فرقی
ملحوظ نه و تغییری مشهود نه کلّ ذلک من فضل اللّه یؤتیه من یشآء من خلقه
انشآءاللّه از ارض نفی احتراز جوئیم و ببحر اثبات درآئیم تا عوالم جمع و فرق
و توحید و تفریق و تحدید و تجرید الهی را ببصری که مقدّس از عناصر و اضداد
است مشاهده کنیم و باعلی افق قرب و قدس حضرت معانی پرواز نمائیم

\end{otherlanguage}

Pas naẓar rā az ḥodūdāt-e ẓāhere ṭāher va monazzah kon tā hame rā be-yek
esm va yek rasm va yek dhāt va yek ḥaqīqat moshāhede namāʾī va asrār-e
rojūʿ-e kalemāt rā ham dar ḥorūfāt-e nāzele molāḥeẓe farmāʾī. Qadrī
tafakkor dar aṣḥāb-e ʿahd-e noqṭe-ye forqān namā ke chegūne az jamīʿ-e
jahāt-e bashariyye va moshtahayāt-e nafsiyye be-nafaḥāt-e qodsiyye-ye ān
ḥażrat pāk va moqaddas va monqaṭeʿ gashtand va qabl az hame-ye ahl-e arḍ
be-sharaf-e leqā ke ʿeyn-e liqāʾ Allāh būd fāʾez shodand va az koll-e
ahl-e arḍ monqaṭeʿ gashtand. Chenānche shenīde-īd ke dar moqābel-e ān
maẓhar-e dhī al-jalāl chegūne jān nethār mīfarmūdand. Va ḥāl hamān
thobūt va rosūkh va enqeṭāʿ rā be-ʿeyne molāḥeẓe farmā dar aṣḥāb-e
noqṭe-ye bayān rājeʿ shode, chenānche molāḥeẓe farmūde-īd ke chegūne īn
aṣḥāb az badāyeʿ-e jūd-e rabb al-arbāb ʿalam-e enqeṭāʿ bar rafraf-e
emtenāʿ bar afrāshtand. Bārī īn anvār az yek meṣbāḥ ẓāher shode-and va
īn athmār az yek shajare rūʾīde-and. Fī al-ḥaqīqe farqī malḥūẓ na va
taghyīrī mashhūd na. Kullu dhālika min faḍli Allāh yuʾtīhi man yashāʾu
min khalqihi. In shāʾ Allāh az arḍ-e nafy eḥterāz jūʾīm va be-baḥr-e
ethbāt dar āʾīm tā ʿavālem-e jamʿ va farq va tawḥīd va tafrīq va taḥdīd
va tajrīd-e ilāhī rā be-baṣarī ke moqaddas az ʿanāṣor va ażdād ast
moshāhede konīm va be-aʿlā ofoq-e qorb va qods-e ḥażrat-e maʿānī parvāz
namāʾīm.

Therefore, cleanse and sanctify your vision from outward limitations, so
that you may see them all as having one name, one character, one
essence, and one reality, and may also discern the mysteries of the
return of the Words in the revealed Letters. Reflect a little on the
companions in the age of the Point of the Furqan: how, through His holy
breaths, they became purified, sanctified, and detached from every human
concern and selfish desire; how, before all the people of the earth,
they attained the honor of meeting Him, which was the very meeting with
God, and detached themselves from all the people of the earth. You have
heard how they sacrificed their lives before that glorious
Manifestation. Now observe how that same steadfastness, firmness, and
detachment have returned, in their very likeness, in the companions of
the Point of the Bayan. You have seen how, through the wondrous bounty
of the Lord of Lords, these companions raised the banner of detachment
upon the heights of inaccessibility. In short, these lights have
appeared from one lamp, and these fruits have grown from one tree. In
truth, no difference can be discerned and no change observed. All this
is from the grace of God, which He bestows on whomever He wills among
His creatures. God willing, may we shun the land of negation and enter
the ocean of affirmation, so that, with a vision sanctified from the
elements and their opposites, we may behold the divine worlds of union
and distinction, unity and differentiation, limitation and
transcendence, and soar to the highest horizon of nearness to the Lord
of inner meanings and His holiness.

\subsection{172}\label{section-362}

\begin{otherlanguage}{persian}

پس از این بیانات معلوم شد که اگر در آخر لا آخر طلعتی بیاید و قیام نماید
بر امری که قیام نمود بر آن طلعت اوّل لا اوّل هرآینه صدق طلعت اوّل بر طلعت
آخر میشود زیرا که طلعت آخر لا آخر قیام نمود بهمان امر که طلعت اوّل لا اوّل
بر آن قیام نمود اینست که نقطهٔ بیان روح ما سواه فداه شموس احدیّه را بشمس
مثال زده‌اند که اگر از اوّل لا اوّل الی آخر لا آخر طلوع نماید همان شمس است
که طالع میشود حال اگر گفته شود این شمس همان شمس اوّلیّه است صحیح است و اگر
گفته شود که رجوع آن شمس است ایضاً صحیح است و همچنین از این بیان صادق
میآید ذکر ختمیّت بر طلعت بدء و بالعکس زیرا که آنچه طلعت ختم بر آن قیام
مینماید بعینه همان است که جمال بدء بر آن قیام فرموده

\end{otherlanguage}

Pas az īn bayānāt maʿlūm shod ke agar dar ākher-e lā ākher ṭalʿatī
biyāyad va qiyām namāyad bar amrī ke qiyām namūd bar ān ṭalʿat-e avval-e
lā avval har-āyene ṣedq-e ṭalʿat-e avval bar ṭalʿat-e ākher mīshavad.
Zīrā ke ṭalʿat-e ākher-e lā ākher qiyām namūd be-hamān amr ke ṭalʿat-e
avval-e lā avval bar ān qiyām namūd. Īn ast ke noqṭe-ye bayān, rūḥ-e mā
sevāh fedāh, shomūs-e aḥadiyye rā be-shams methāl zade-and ke agar az
avval-e lā avval ilā ākher-e lā ākher ṭolūʿ namāyad hamān shams ast ke
ṭāleʿ mīshavad. Ḥāl agar gofte shavad īn shams hamān shams-e avvaliyye
ast ṣaḥīḥ ast va agar gofte shavad ke rojūʿ-e ān shams ast ayżan ṣaḥīḥ
ast. Va hamchenīn az īn bayān ṣādeq mīāyad dhekr-e khatmiyyat bar
ṭalʿat-e badʾ va be-l-ʿaks zīrā ke ānche ṭalʿat-e khatm bar ān qiyām
mīnamāyad be-ʿeyne hamān ast ke jamāl-e badʾ bar ān qiyām farmūde.

These explanations make it clear that if, at the end that has no end, a
Manifestation should come and arise to uphold the same Cause as the
Manifestation at the beginning that has no beginning, the designation of
the first Manifestation would assuredly apply to the last. For the
Manifestation at the end that has no end has arisen to uphold the same
Cause as the Manifestation at the beginning that has no beginning. Thus
the Point of the Bayan---may the souls of all besides Him be sacrificed
for Him---has likened the Suns of Oneness to the sun: even if it rises
from the beginning that has no beginning until the end that has no end,
it is the same sun that rises. Now, if it is said that this sun is that
first sun, it is correct; and if it is said that it is the return of
that sun, it is likewise correct. Similarly, this explanation shows that
the designation of finality truly applies to the Manifestation of the
beginning, and vice versa, for what the Manifestation of the end arises
to uphold is precisely what the Beauty of the beginning arose to uphold.

\subsection{173}\label{section-363}

\begin{otherlanguage}{persian}

و این مطلب با اینکه چه‌ قدر واضحست نزد شاربان صهبای علم و ایقان مع‌ذلک چه
مقدار از نفوس که بسبب عدم بلوغ باین مطلب بذکر خاتم النّبیّین محتجب شده از
جمیع فیوضات محجوب و ممنوع شده‌اند با اینکه خود آن حضرت فرمود امّا النّبیّون
فأنا و همچنین فرمودند منم آدم و نوح و موسی و عیسی چنانچه ذکر شد مع‌ذلک
تفکّر نمینمایند بعد از آنکه بر آن جمال ازلی صادق میآید باینکه فرمودند منم
آدم اوّل همین قسم صادق میآید که بفرمایند منم آدم آخر و همچنان که بدء
انبیا را که آدم باشد بخود نسبت دادند همین قسم ختم انبیا هم بآن جمال الهی
نسبت داده میشود و این بسی واضحست که بعد از آنکه بدء النّبیّین بر آن حضرت
صادقست همان قسم ختم النّبیّین صادق آید

\end{otherlanguage}

Va īn maṭlab bā īnke che-qadr vāżeḥ ast nazd-e shārebān-e ṣahbā-ye ʿelm
va īqān maʿa-dhālek che meqdār az nofūs ke be-sabab-e ʿadam-e bolūgh
be-īn maṭlab be-dhekr-e khātam al-nabiyyīn moḥtajeb shode az jamīʿ-e
foyūżāt maḥjūb va mamnūʿ shode-and. Bā īnke khod-e ān ḥażrat farmūd:
ammā al-nabiyyūna fa-anā. Va hamchenīn farmūdand: manam Ādam va Nūḥ va
Mūsā va ʿĪsā, chenānche dhekr shod. Maʿa-dhālek tafakkor nemīnamāyand.
Baʿd az ānke bar ān jamāl-e azalī ṣādeq mīāyad be-īnke farmūdand manam
Ādam-e avval, hamīn qesm ṣādeq mīāyad ke befarmāyand manam Ādam-e ākher.
Va hamchenān ke badʾ-e anbiyā rā ke Ādam bāshad be-khod nesbat dādand
hamīn qesm khatm-e anbiyā ham be-ān jamāl-e ilāhī nesbat dāde mīshavad.
Va īn basī vāżeḥ ast ke baʿd az ānke badʾ al-nabiyyīn bar ān ḥażrat
ṣādeq ast hamān qesm khatm al-nabiyyīn ṣādeq āyad.

Although this matter is so clear to those who drink the wine of
knowledge and certitude, how many souls, through failing to grasp it,
have nevertheless been veiled by the expression ``Seal of the Prophets''
and have been shut out and deprived of every outpouring of grace. Yet He
Himself said, ``As for the Prophets, I am they.'' He also said, ``I am
Adam, Noah, Moses, and Jesus,'' as already mentioned. Still they do not
reflect. Since the words ``I am the first Adam'' truly apply to that
eternal Beauty, it would be equally true for Him to say, ``I am the last
Adam.'' And just as He attributed to Himself the beginning of the
Prophets, who was Adam, so too the last of the Prophets is attributed to
that divine Beauty. It is abundantly clear that, since the designation
``beginning of the Prophets'' truly applies to Him, the designation
``last of the Prophets'' likewise applies.

\subsection{174}\label{section-364}

\begin{otherlanguage}{persian}

و باین مطلب جمیع اهل ارض در این ظهور ممتحن شده‌اند چنانچه اکثری بهمین قول
تمسّک جسته از صاحب قول معرض شده‌اند و نمیدانم این قوم از اوّلیّت و آخریّت حقّ
جلّ ذکره چه ادراک نموده‌اند اگر مقصود از اوّلیّت و آخریّت اوّلیّت و آخریّت ملکی
باشد هنوز که اسباب ملکی بآخر نرسیده پس چگونه آخریّت بر آن ذات احدیّت صادق
میآید بلکه در این رتبه اوّلیّت نفس آخریّت و آخریّت نفس اوّلیّت باشد

\end{otherlanguage}

Va be-īn maṭlab jamīʿ-e ahl-e arḍ dar īn ẓohūr momtaḥan shode-and,
chenānche aktharī be-hamīn qawl tamassok joste az ṣāḥeb-e qawl moʿreż
shode-and. Va nemīdānam īn qawm az avvaliyyat va ākheriyyat-e Ḥaqq,
jalla dhikruhu, che edrāk namūde-and. Agar maqṣūd az avvaliyyat va
ākheriyyat avvaliyyat va ākheriyyat-e molkī bāshad hanūz ke asbāb-e
molkī be-ākher narasīde, pas chegūne ākheriyyat bar ān dhāt-e aḥadiyyat
ṣādeq mīāyad? Balke dar īn rotbe avvaliyyat nafs-e ākheriyyat va
ākheriyyat nafs-e avvaliyyat bāshad.

All the people of the earth have been tested by this matter in this
revelation, for most have clung to this very saying and turned away from
the One who uttered it. I do not know what these people have understood
by the firstness and lastness of God, exalted be His mention. If
firstness and lastness mean firstness and lastness in the earthly world,
the conditions of that world have not yet come to an end; how, then, can
lastness truly apply to that Essence of Oneness? Rather, in this
station, firstness is lastness itself, and lastness is firstness itself.

\subsection{175}\label{section-365}

\begin{otherlanguage}{persian}

باری همان قسمی که در اوّل لا اوّل صدق آخریّت بر آن مربّی غیب و شهود میآید
همان قسم هم بر مظاهر او صادق میآید و در حینی که اسم اوّلیّت صادقست همان
حین اسم آخریّت صادق و در حینی که بر سریر بدئیّت جالسند همان حین بر عرش
ختمیّت ساکن و اگر بصر حدید یافت شود مشاهده مینماید که مظهر اوّلیّت و آخریّت
و ظاهریّت و باطنیّت و بدئیّت و ختمیّت این ذوات مقدّسه و ارواح مجرّده و انفس
الهیّه هستند و اگر در هوای قدس کان اللّه و لم یکن معه من شیء طائر شوی جمیع
این اسماء را در آن ساحت معدوم صرف و مفقود بحت بینی و دیگر هیچ باین حجبات
و اشارات و کلمات محتجب نشوی چه لطیف و بلند است این مقام که جبرئیل بی‌دلیل
سبیل نجوید و طیر قدسی بی‌اعانت غیبی طیران نتواند

\end{otherlanguage}

Bārī hamān qesmī ke dar avval-e lā avval ṣedq-e ākheriyyat bar ān
morabbī-ye gheyb va shohūd mīāyad hamān qesm ham bar maẓāher-e ū ṣādeq
mīāyad. Va dar ḥīnī ke esm-e avvaliyyat ṣādeq ast hamān ḥīn esm-e
ākheriyyat ṣādeq va dar ḥīnī ke bar sarīr-e badʾiyyat jālesand hamān ḥīn
bar ʿarsh-e khatmiyyat sāken. Va agar baṣar-e ḥadīd yāft shavad
moshāhede mīnamāyad ke maẓhar-e avvaliyyat va ākheriyyat va ẓāheriyyat
va bāṭeniyyat va badʾiyyat va khatmiyyat īn dhavāt-e moqaddase va
arvāḥ-e mojarrade va anfos-e ilāhiyye hastand. Va agar dar havā-ye
qods-e kāna Allāhu wa-lam yakun maʿahu min shayʾ ṭāʾer shavī jamīʿ-e īn
asmā rā dar ān sāḥat maʿdūm-e ṣerf va mafqūd-e baḥt bīnī va dīgar hīch
be-īn ḥojobāt va eshārāt va kalemāt moḥtajeb nashavī. Che laṭīf va
boland ast īn maqām ke Jebraʾīl bī-dalīl sabīl najūyad va ṭeyr-e qodsī
bī-eʿānat-e gheybī ṭayarān natavānad.

In short, just as lastness truly applies, in the beginning that has no
beginning, to the Nurturer of the unseen and the seen, so it also truly
applies to His Manifestations. At the very moment when the name of
firstness applies, the name of lastness also applies; and at the very
moment when they sit upon the throne of beginning, they also dwell upon
the throne of finality. If keen vision is attained, it will perceive
that these sanctified Essences, these transcendent Spirits, and these
divine Souls are the manifestations of firstness and lastness,
outwardness and inwardness, beginning and finality. And if you soar in
the holy atmosphere of ``God was, and nothing whatsoever was with Him,''
you will see all these names in that realm as sheer nonexistence and
utter absence, and will no longer be veiled by these veils, allusions,
and words. How subtle and exalted is this station, where Gabriel cannot
find his way without a guide, and the holy bird cannot fly without
unseen assistance!

\subsection{176}\label{section-366}

\begin{otherlanguage}{persian}

حال قول حضرت امیر را ادراک نما که فرموده کشف سبحات الجلال من غیر اشارة و
از جملهٔ سبحات مجلّله علمای عصر و فقهای زمان ظهورند که جمیع نظر بعدم ادراک
و اشتغال و حبّ بریاست ظاهره تسلیم امر اللّه نمینمایند بلکه گوش نمیدهند تا
نغمهٔ الهی را بشنوند بل یجعلون اصابعهم فی آذانهم و عباد هم چون ایشان را
من دون اللّه ولیّ خود اخذ نموده‌اند منتظر رد و قبول این خشبهای مسنّده هستند
زیرا از خود بصر و سمع و قلبی ندارند که تمیز و تفصیل دهند میانهٔ حقّ و باطل

\end{otherlanguage}

Ḥāl qawl-e Ḥażrat-e Amīr rā edrāk namā ke farmūde: kashfu subuḥāti
al-jalāli min ghayri ishāra. Va az jomle-ye soboḥāt-e mojallele
ʿolamā-ye ʿaṣr va foqahā-ye zamān-e ẓohūrand ke jamīʿ naẓar be-ʿadam-e
edrāk va eshteghāl va ḥobb be-riyāsat-e ẓāhere taslīm-e amr Allāh
nemīnamāyand balke gūsh nemīdahand tā naghme-ye ilāhī rā beshnavand. Bal
yajʿalūna aṣābiʿahum fī ādhānihim. Va ʿebād ham chon īshān rā min dūn
Allāh valī-ye khod akhdh namūde-and montaẓer-e radd va qabūl-e īn
khashabhā-ye mosannade hastand zīrā az khod baṣar va samʿ va qalbī
nadārand ke tamīz va tafṣīl dehand miyāne-ye Ḥaqq va bāṭel.

Now grasp the saying of His Holiness the Commander: ``Unveiling the
splendors of glory without allusion.'' Among these veiling splendors are
the learned and the jurists of the age of revelation, who, through their
lack of understanding, their preoccupations, and their love of outward
leadership, all refuse to submit to the Cause of God. Indeed, they will
not lend an ear to hear the divine melody. ``Rather, they put their
fingers in their ears.'' Since the people have taken them as their
guardians instead of God, they await the rejection or acceptance of
these propped-up timbers, for they have no sight, hearing, or heart of
their own with which to distinguish and discriminate between truth and
falsehood.

\subsection{177}\label{section-367}

\begin{otherlanguage}{persian}

با اینکه همهٔ انبیا و اصفیا و اولیا من عند اللّه امر فرمودند که بچشم و گوش
خود بشنوند و ملاحظه نمایند مع‌ذلک معتنی بنصح انبیا نگشته تابع علمای خود
بوده و خواهند بود و اگر مسکینی و یا فقیری که عاری از لباس اهل علم باشد
بگوید یا قوم اتّبعوا المرسلین جواب گویند که این همه علما و فضلا با این
ریاست ظاهره و البسهٔ مقطّعهٔ لطیفه نفهمیده‌اند و حقّ را از باطل ادراک
ننموده‌اند و تو و امثال تو ادراک نموده‌اید و نهایت تعجّب مینمایند از چنین
قولی با اینکه امم سلف اکثر و اعظم و اکبرند و اگر کثرت و لباس علم سبب و
علّت علم و صدق باشد البتّه امم سابقه اولی و اسبقند

\end{otherlanguage}

Bā īnke hame-ye anbiyā va aṣfiyā va awliyā min ʿindi Allāh amr farmūdand
ke be-chashm va gūsh-e khod beshnavand va molāḥeẓe namāyand, maʿa-dhālek
moʿtanī be-noṣḥ-e anbiyā nagashte tābeʿ-e ʿolamā-ye khod būde va
khwāhand būd. Va agar meskīnī va yā faqīrī ke ʿārī az lebās-e ahl-e ʿelm
bāshad begūyad, yā qawmi ittabiʿū al-mursalīn, javāb gūyand ke īn hame
ʿolamā va fożalā bā īn reyāsat-e ẓāhere va albese-ye moqaṭṭaʿe-ye laṭīfe
nafahmīde-and va Ḥaqq rā az bāṭel edrāk nanamūde-and va to va amthāl-e
to edrāk namūde-īd, va nehāyat-e taʿajj ob mīnamāyand az chonīn qawlī.
Bā īnke omam-e salaf akthar va aʿẓam va akbarand, va agar kathrat va
lebās-e ʿelm sabab va ʿellat-e ʿelm va ṣedq bāshad, albatte omam-e
sābeqe awlā va asbaqand.

Although all the Prophets, chosen ones, and holy ones commanded, on
God's behalf, that people hear with their own ears and observe with
their own eyes, they have nevertheless paid no heed to the Prophets'
counsel, and have followed, and will continue to follow, their own
religious scholars. If some humble or poor person, unadorned with the
garb of learning, says, ``My people, follow the Messengers,'' they
reply, ``Have all these scholars and learned men, with their outward
authority and finely tailored robes, failed to understand and to
distinguish truth from falsehood, while you and your kind have
understood?'' They are utterly astonished at such a statement, even
though the peoples of former times were more numerous, greater, and
mightier. If numbers and the garb of learning were the grounds and
causes of knowledge and truthfulness, the peoples of former times would
certainly have the stronger and prior claim.

\subsection{178}\label{section-368}

\begin{otherlanguage}{persian}

و با اینکه این فقره هم معلوم و واضحست که در جمیع احیان ظهور مظاهر قدسیّه
علمای عهد مردم را از سبیل حقّ منع مینمودند چنانچه در جمیع کتب و صحف سماوی
مذکور و مسطور است و احدی از انبیا مبعوث نشد مگر آنکه محلّ بغض و انکار و
رد و سبّ علما گشت قاتلهم اللّه بما فعلوا من قبل و من بعد کانوا یفعلون حال
کدام سبحات جلال اعظم از این هیاکل ضلالست واللّه کشف آن اعظم امور است و
خرقش اکبر اعمال وفّقنا اللّه و ایّاکم یا معشر الرّوح لعلّکم بذلک فی زمن
المستغاث توفّقون و من لقآء اللّه فی ایّامه لا تحتجبون

\end{otherlanguage}

Va bā īnke īn feqre ham maʿlūm va vāżeḥ ast ke dar jamīʿ-e aḥyān-e
ẓohūr-e maẓāher-e qodsiyye ʿolamā-ye ʿahd mardom rā az sabīl-e Ḥaqq manʿ
mīnamūdand, chenānche dar jamīʿ-e kotob va ṣoḥof-e samāvī madhkūr va
masṭūr ast. Va aḥadī az anbiyā mabʿūth nashod magar ānke maḥall-e boghż
va enkār va radd va sabb-e ʿolamā gasht. Qātalahumu Allāhu bi-mā faʿalū
min qablu wa-min baʿdu kānū yafʿalūn. Ḥāl kodām soboḥāt-e jalāl aʿẓam az
īn hayākel-e żalāl ast? Wa-Allāh, kashf-e ān aʿẓam-e omūr ast va
kharqash akbar-e aʿmāl. Waffaqanā Allāhu wa-iyyākum yā maʿshara al-rūḥ,
laʿallakum bi-dhālika fī zamani al-Mustaghāth tuwaffaqūn wa-min liqāʾi
Allāhi fī ayyāmihi lā taḥtajibūn.

It is also known and clear that whenever the holy Manifestations
appeared, the religious scholars of the age kept people from the path of
Truth, as is recorded and written in all the heavenly Books and
Scriptures. Not a single Prophet was raised up without becoming the
object of the scholars' hatred, denial, rejection, and abuse. May God
confound them for what they did before and continued to do afterward!
Now, what splendors of glory are greater than these embodiments of
error? By God, unveiling them is the greatest of undertakings, and
tearing through them the greatest of deeds. May God grant success to us
and to you, O company of the Spirit, so that through this you may be
aided in the time of the Mustaghāth and may not be veiled from meeting
God in His days.

\subsection{179}\label{section-369}

\begin{otherlanguage}{persian}

و همچنین ذکر خاتم النّبیّین و امثال آن از سبحات مجلّله است که کشف آن از
اعظم امور است نزد این همج رعاع و جمیع باین حجبات محدوده و سبحات مجلّلهٔ
عظیمه محتجب مانده‌اند آیا نغمهٔ طیر هویّه را نشنیده‌اند که میفرماید الف
فاطمه نکاح نمودم که همه بنت محمّد بن عبداللّه خاتم النّبیّین بودند حال
ملاحظه فرما که چه‌ قدر از اسرار در سرادق علم الهی مستور است و چه‌ مقدار
جواهر علم او در خزائن عصمت مکنون تا یقین نمائی که صنع او را بدایت و
نهایتی نبوده و نخواهد بود و فضای قضای او اعظم از آنست که ببیان تحدید شود
و یا طیر افئده آن را طیّ نماید و تقدیرات قدریّهٔ او اکبر از آنست که بادراک
نفسی منتهی شود خلق او از اوّل لا اوّل بوده و آخری او را اخذ نکرده و مظاهر
جمال او الی نهایت لا نهایه خواهند بود و ابتدائی او را ندیده حال در همین
بیان ملاحظه فرما که چگونه حکم آن بر جمیع این طلعات صدق مینماید

\end{otherlanguage}

Va hamchenīn dhekr-e Khātam al-Nabiyyīn va amthāl-e ān az soboḥāt-e
mojallale ast ke kashf-e ān az aʿẓam-e omūr ast nazd-e īn hamaj-e raʿāʿ,
va jamīʿ be-īn ḥojobāt-e maḥdūde va soboḥāt-e mojallale-ye ʿaẓīme
moḥtajeb mānde-and. Āyā naghme-ye ṭayr-e hoviyye rā nashenīde-and ke
mīfarmāyad, alf Fāṭeme nekāḥ namūdam ke hame bent-e Moḥammad ebn-e ʿAbd
Allāh Khātam al-Nabiyyīn būdand? Ḥāl molāḥeẓe farmā ke che qadr az asrār
dar sorādeq-e ʿelm-e elāhī mastūr ast va che meqdār javāher-e ʿelm-e ū
dar khazāʾen-e ʿeṣmat maknūn, tā yaqīn namāʾī ke ṣonʿ-e ū rā bedāyat va
nehāyatī nabūde va nakhwāhad būd, va fażā-ye qażā-ye ū aʿẓam az ān ast
ke be-bayān taḥdīd shavad va yā ṭayr-e afʾede ān rā ṭayy namāyad, va
taqdīrāt-e qadariyye-ye ū akbar az ān ast ke be-edrāk-e nafsī montahī
shavad. Khalq-e ū az avval-e lā avval būde va ākherī ū rā akhdh nakarde,
va maẓāher-e jamāl-e ū ilā nehāyat-e lā nehāye khwāhand būd va ebtedāʾī
ū rā nadīde. Ḥāl dar hamīn bayān molāḥeẓe farmā ke chegūne ḥokm-e ān bar
jamīʿ-e īn ṭalaʿāt ṣedq mīnamāyad.

Likewise, the designation ``Seal of the Prophets'' and similar
expressions are among the majestic splendors whose unveiling is one of
the greatest undertakings in the eyes of this ignorant rabble. All have
remained veiled by these limiting veils and these mighty, majestic
splendors. Have they not heard the melody of the Bird of divine Being,
who says, ``I married a thousand Fāṭimihs, all of whom were daughters of
Muḥammad son of ʿAbd Allāh, the Seal of the Prophets''? Now consider how
many mysteries lie hidden in the tabernacle of divine knowledge, and how
many jewels of His knowledge are concealed in the treasuries of
infallibility, so that you may know with certainty that His creation has
never had, and will never have, a beginning or an end. The expanse of
His decree is too vast to be bounded by words or traversed by the bird
of the heart, and His ordained decrees are too great to be encompassed
by anyone's understanding. His creation has existed from a beginning
that has no beginning, and no end has overtaken it; the Manifestations
of His beauty will endure to an end that has no end, and no beginning
has beheld them. Now reflect on this very utterance, and see how its
meaning holds true of all these Countenances.

\subsection{180}\label{section-370}

\begin{otherlanguage}{arabic}

و همچنین نغمهٔ جمال ازلی حسین بن علی را ادراک نما که به سلمان میفرماید که
مضمون آن اینست بودم با الف آدم که فاصلهٔ هر آدم به آدم بعد خمسین الف سنه
بود و با هر یک ولایت پدرم را عرض نمودم و تفصیلی ذکر میفرماید تا آنکه
میفرماید الف مرّه جهاد نمودم در سبیل الهی که اصغر و کوچکتر از همه مثل
غزوهٔ خیبر بود که پدرم با کفّار محاربه و مجادله نمود حال اسرار ختم و رجع و
لا اوّلیّت و لا آخریّت صنع همه را از این دو روایت ادراک فرما

\end{otherlanguage}

Va hamchenīn naghme-ye Jamāl-e azalī Ḥosayn ebn-e ʿAlī rā edrāk namā ke
be-Salmān mīfarmāyad ke mażmūn-e ān īn ast: būdam bā alf Ādam ke
fāṣele-ye har Ādam be-Ādam-e baʿd khamsīn alf sana būd va bā har yek
velāyat-e pedaram rā ʿarż namūdam. Va tafṣīlī dhekr mīfarmāyad tā ānke
mīfarmāyad: alf marre jehād namūdam dar sabīl-e elāhī ke aṣghar va
kūchektar az hame methl-e ghazve-ye Khaybar būd ke pedaram bā koffār
moḥārebe va mojādele namūd. Ḥāl asrār-e khatm va rajʿ va lā avvaliyyat
va lā ākheriyyat-e ṣonʿ hame rā az īn do revāyat edrāk farmā.

Likewise, grasp the melody of the eternal Beauty, Ḥusayn son of ʿAlī,
who addresses Salmān in words whose meaning is this: ``I was with a
thousand Adams, and the interval between each Adam and the next was
fifty thousand years. To each one I presented my father's authority.''
He gives further details until He says, ``A thousand times I fought in
the path of God, and the least and smallest of them all was like the
battle of Khaybar, when my father fought and contended with the
unbelievers.'' Now grasp from these two traditions all the mysteries of
finality and return, and of creation's having neither beginning nor end.

\subsection{181}\label{section-371}

\begin{otherlanguage}{persian}

باری ای حبیب من مقدّسست نغمهٔ لاهوت که باستماع و عقول ناسوت محدود شود نملهٔ
وجود کجا تواند بعرصهٔ معبود قدم گذارد اگرچه نفوس ضعیفه از عدم ادراک این
بیانات معضله را انکار نمایند و امثال این احادیث را نفی کنند بلی لا یعرف
ذلک الّا اولو الألباب قل هو الختم الّذی لیس له ختم فی الابداع و لا بدء له
فی الاختراع اذاً یا ملأ الارض فی ظهورات البدء تجلّیات الختم تشهدون

\end{otherlanguage}

Bārī, ey ḥabīb-e man, moqaddas ast naghme-ye lāhūt ke be-estemāʿ va
ʿoqūl-e nāsūt maḥdūd shavad. Namle-ye vojūd kojā tavānad be-ʿarṣe-ye
Maʿbūd qadam godhārad? Agarche nofūs-e żaʿīfe az ʿadam-e edrāk īn
bayānāt-e moʿżele rā enkār namāyand va amthāl-e īn aḥādīth rā nafy
konand. Balī, lā yaʿrifu dhālika illā ulū al-albāb. Qul huwa al-khatmu
alladhī laysa lahu khatmun fī al-ibdāʿ wa-lā badʾa lahu fī al-ikhtirāʿ.
Idhan yā malaʾa al-arḍ fī ẓuhūrāti al-badʾ tajalliyāti al-khatm
tashhadūn.

In short, my beloved, the melody of the divine realm is too exalted to
be confined within the hearing and understanding of the human world. How
can the ant of existence set foot in the court of the One who is
worshipped? Yet feeble souls, through lack of understanding, deny these
abstruse utterances and reject such traditions. Indeed, none understands
this except those endowed with insight. Say: He is the End who has no
end in creation and no beginning in origination. Thus, people of the
earth, in the manifestations of the Beginning you witness the
revelations of the End.

\subsection{182}\label{section-372}

\begin{otherlanguage}{persian}

بسیار تعجّب است که این قوم در بعضی از مراتب که مطابق میل و هوای ایشان است
متمسّک بآیهٔ منزلهٔ در فرقان و احادیث اولی الایقان میشوند و از بعضی که
مغایر هوای ایشانست بالمرّه اعراض مینمایند قل أ تؤمنون ببعض الکتاب و
تکفرون ببعض ما لکم کیف تحکمون ما لا تشعرون مثل آنکه در کتاب مبین ربّ
العالمین بعد از ذکر ختمیّت فی قوله تعالی ولکنّه رسول اللّه و خاتم النّبیّین
جمیع ناس را بلقای خود وعده فرموده چنانچه آیات مدلّهٔ بر لقای آن ملیک بقا
در کتاب مذکور است و بعضی از قبل ذکر شده و خدای واحد شاهد مقال است که هیچ
امری اعظم از لقا و اصرح از آن در فرقان ذکر نیافته فهنیئاً لمن فاز به فی
یوم اعرض عنه اکثر النّاس کما انتم تشهدون

\end{otherlanguage}

Besyār taʿajj ob ast ke īn qawm dar baʿżī az marāteb ke moṭābeq-e meyl
va havā-ye īshān ast motamassek be-āye-ye monazzale dar Forqān va
aḥādīth-e ūlī al-īqān mīshavand va az baʿżī ke moghāyer-e havā-ye īshān
ast be-almarre eʿrāż mīnamāyand. Qul a-tuʾminūna bi-baʿḍi al-kitābi
wa-takfurūna bi-baʿḍ? Mā lakum kayfa taḥkumūna mā lā tashʿurūn? Methl-e
ānke dar Ketāb-e mobīn Rabb-e ʿālamīn baʿd az dhekr-e khātemiyyat fī
qawlihi taʿālā, wa-lākinnahu rasūlu Allāhi wa-khātama al-nabiyyīn,
jamīʿ-e nās rā be-leqā-ye khod vaʿde farmūde, chenānche āyāt-e modelle
bar leqā-ye ān Malīk-e baqā dar Ketāb madhkūr ast va baʿżī az qabl dhekr
shode. Va Khodā-ye vāḥed shāhed-e maqāl ast ke hīch amrī aʿẓam az leqā
va aṣraḥ az ān dar Forqān dhekr nayāfte. Fa-hanīʾan li-man fāza bihi fī
yawmin aʿraḍa ʿanhu aktharu al-nāsi kamā antum tashhadūn.

It is astonishing that, in matters that accord with their inclinations
and desires, these people hold fast to the verse revealed in the Qurʾān
and the traditions of those endowed with certitude, while utterly
turning away from those that run counter to their desires. Say: Do you
believe in part of the Book and disbelieve in part? What is wrong with
you? How can you judge what you do not understand? For example, in the
clear Book, after mentioning finality in His exalted words, ``But he is
the Messenger of God and the Seal of the Prophets,'' the Lord of the
worlds has promised all people a meeting with Himself. Thus, verses
testifying to a meeting with that King of eternity are recorded in the
Book, and some have already been cited. The one God bears witness to
these words: nothing greater than this meeting, or more explicitly
stated, has been mentioned in the Qurʾān. Blessed, then, is the one who
attains it on a day when most people turn away from it, as you
yourselves witness.

\subsection{183}\label{section-373}

\begin{otherlanguage}{persian}

و مع‌ذلک بحکم اوّل از امر ثانی معرض گشته‌اند با اینکه حکم لقا در یوم قیام
منصوص است در کتاب و قیامت هم بدلائل واضحه ثابت و محقّق شد که مقصود قیام
مظهر اوست بر امر او و همچنین از لقا لقای جمال اوست در هیکل ظهور او اذ
انّه لا تدرکه الأبصار و هو یدرک الأبصار و با جمیع این مطالب ثابته و
بیانات واضحه من حیث لا یشعر بذکر ختم تمسّک جسته‌اند و از موجد ختم و بدء در
یوم لقای او بالمرّه محتجب مانده‌اند و لو یؤاخذ اللّه النّاس بما کسبوا ما ترک
علی ظهرها من دابّة ولکن یؤخّرهم الی اجل مسمّی و از همهٔ این مراتب چشم پوشیده
اگر این قوم بقطره‌ئی از چشمهٔ لطیف یفعل ما یشآء و یحکم ما یرید میآشامیدند
هیچ اینگونه اعتراضات غیر مرضیّه بر محلّ امر نمینمودند امر و قول و فعل در
قبضهٔ قدرت اوست کلّ شیء فی قبضة قدرته اسیر و انّ ذلک علیه سهل یسیر فاعلست
آنچه را اراده نماید و عامل است آنچه میل فرماید من قال لم و بم فقد کفر و
اگر این عباد قدری بشعور بیایند از آنچه مرتکب شده‌اند هلاک میشوند و
بدستهای خود خود را بنار که مقرّ و مرجع ایشانست راجع مینمایند آیا
نشنیده‌اند که میفرماید لا یسأل عمّا یفعل و با این بیانات چگونه میتوان
جسارت نمود و بزخارف قول مشغول شد

\end{otherlanguage}

Va maʿa-dhālek be-ḥokm-e avval az amr-e thānī moʿreż gashte-and, bā īnke
ḥokm-e leqā dar yawm-e qiyām manṣūṣ ast dar Ketāb, va qiyāmat ham
be-dalāʾel-e vāżeḥe thābet va moḥaqqaq shod ke maqṣūd qiyām-e Maẓhar-e
ūst bar amr-e ū, va hamchenīn az leqā leqā-ye jamāl-e ūst dar haykal-e
ẓohūr-e ū, idh innahu lā tudrikuhu al-abṣāru wa-huwa yudriku al-abṣār.
Va bā jamīʿ-e īn maṭāleb-e thābete va bayānāt-e vāżeḥe min ḥaythu lā
yashʿur be-dhekr-e khatm tamassok joste-and va az Mūjed-e khatm va badʾ
dar yawm-e leqā-ye ū be-almarre moḥtajeb mānde-and. Wa-law yuʾākhidhu
Allāhu al-nāsa bi-mā kasabū mā taraka ʿalā ẓahrihā min dābbatin wa-lākin
yuʾakhkhiruhum ilā ajalin musamman. Va az hame-ye īn marāteb chashm
pūshīde, agar īn qawm be-qaṭreʾī az chashme-ye laṭīf-e yafʿalu mā yashāʾ
wa-yaḥkumu mā yurīd mī-āshāmīdand, hīch īngūne eʿterāżāt-e gheyr-e
marżiyye bar Maḥall-e amr namīnamūdand. Amr va qawl va feʿl dar qabże-ye
qodrat-e ūst. Kullu shayʾin fī qabḍati qudratihi asīr, wa-inna dhālika
ʿalayhi sahlun yasīr. Fāʿel ast ānche rā erāde namāyad va ʿāmel ast
ānche meyl farmāyad. Man qāla lima wa-bima fa-qad kafar. Va agar īn
ʿebād qadrī be-shoʿūr beyāyand, az ānche mortakeb shode-and halāk
mīshavand va be-dasthā-ye khod khod rā be-nār ke maqarr va marjaʿ-e
īshān ast rājeʿ mīnamāyand. Āyā nashenīde-and ke mīfarmāyad, lā yusʾalu
ʿammā yafʿal? Va bā īn bayānāt chegūne mītavān jesārat namūd va
be-zakhāref-e qawl mashghūl shod?

Nevertheless, on the strength of the first pronouncement, they have
turned away from the second, although the promise of meeting on the Day
of Resurrection is explicitly stated in the Book. It has also been
clearly established and demonstrated that resurrection means the rising
of His Manifestation to proclaim His Cause, and that meeting Him means
meeting His Beauty in the person of His Manifestation, for ``Eyes do not
perceive Him, but He perceives the eyes.'' Despite all these established
truths and clear statements, they have unwittingly clung to the mention
of finality and remained utterly veiled from the Originator of the end
and the beginning on the day of meeting Him. ``If God were to take
people to task for what they have earned, He would leave no living
creature on the earth's surface, but He gives them respite until an
appointed time.'' Leaving all these considerations aside, if these
people drank but a drop from the gentle spring of ``He does whatever He
wills and decrees whatever He desires,'' they would never raise such
unseemly objections against the Seat of His Cause. Command, word, and
deed lie within the grasp of His power. Everything is held captive in
the grasp of His power, and that is easy and effortless for Him. He does
whatever He wills and carries out whatever He desires. Whoever says
``Why?'' and ``For what reason?'' has indeed disbelieved. If these
servants came to even a little awareness, they would perish at what they
have committed and, with their own hands, consign themselves to the fire
that is their abode and destination. Have they not heard Him say, ``He
is not questioned about what He does''? In the face of these statements,
how can anyone presume to be so bold and indulge in specious words?

\subsection{184}\label{section-374}

\begin{otherlanguage}{persian}

سبحان‌اللّه جهل و نادانی عباد بمقام و حدّی رسیده که بعلم و ارادهٔ خود مقبل
شده از علم و ارادهٔ حقّ جلّ و عزّ معرض گشته‌اند

\end{otherlanguage}

Sobḥān Allāh, jahl va nādānī-ye ʿebād be-maqām va ḥaddī resīde ke
be-ʿelm va erāde-ye khod moqbel shode az ʿelm va erāde-ye Ḥaqq jalla
wa-ʿazza moʿreż gashte-and.

Glory be to God! The ignorance and folly of these servants have reached
such a degree that they have turned toward their own knowledge and will,
and away from the knowledge and will of God, exalted and mighty is He.

\subsection{185}\label{section-375}

\begin{otherlanguage}{arabic}

حال انصاف دهید اگر این عباد موقن باین کلمات درّیّه و اشارات قدسیّه شوند و
حقّ را یفعل ما یشآء بدانند دیگر چگونه باین مزخرفات تشبّث مینمایند و تمسّک
میجویند بلکه آنچه بفرماید بجان اقرار نمایند و مذعن شوند قسم بخدا که اگر
تقدیرات مقدّره و حکمتهای قدریّه سبقت نیافته بود ارض جمیع این عباد را معدوم
مینمود ولکن یؤخّر ذلک الی میقات یوم معلوم

\end{otherlanguage}

Ḥāl enṣāf dehīd, agar īn ʿebād mūqen be-īn kalemāt-e dorriyye va
eshārāt-e qodsiyye shavand va Ḥaqq rā yafʿalu mā yashāʾ bedānand, dīgar
chegūne be-īn mozakhrafāt tashabboth mīnamāyand va tamassok mījūyand?
Balke ānche befarmāyad be-jān eqrār namāyand va modhʿen shavand. Qasam
be-Khodā ke agar taqdīrāt-e moqaddare va ḥekmathā-ye qadariyye sabqat
nayāfte būd, arż jamīʿ-e īn ʿebād rā maʿdūm mīnamūd. Wa-lākin
yuʾakhkhiru dhālika ilā mīqāti yawmin maʿlūm.

Now judge fairly: if these servants became certain of these pearl-like
words and holy allusions, and acknowledged that God ``does whatever He
wills,'' how could they continue to cling to and rely upon such
absurdities? Rather, they would wholeheartedly acknowledge and submit to
whatever He said. I swear by God that, had the ordained decrees and the
wisdom of providence not already been established, the earth would have
annihilated all these servants. But He postpones this until the
appointed hour of a known day.

\subsection{186}\label{section-376}

\begin{otherlanguage}{persian}

باری هزار و دویست و هشتاد سنه از ظهور نقطهٔ فرقان گذشت و جمیع این همج
رعاع در هر صباح تلاوت فرقان نموده‌اند و هنوز بحرفی از مقصود فائز نشدند و
خود قرائت میکنند بعضی آیات را که صریح بر مطالب قدسیّه و مظاهر عزّ صمدیّه
است مع‌ذلک هیچ ادراک ننموده‌اند و این مدّت این قدر ادراک نشده که مقصود از
تلاوت کتب و قرائت صحف در هر عصری ادراک معانی آن و بلوغ بمعارج اسرار آن
بوده والّا تلاوت بی‌معرفت را البتّه فائدهٔ کلّی نباشد

\end{otherlanguage}

Bārī, hezār va devīst va hashtād sane az ẓohūr-e Noqṭe-ye Forqān
godhasht va jamīʿ-e īn hamaj-e raʿāʿ dar har ṣabāḥ telāvat-e Forqān
namūde-and va hanūz be-ḥarfī az maqṣūd fāʾez nashodand. Va khod qerāʾat
mīkonand baʿżī āyāt rā ke ṣarīḥ bar maṭāleb-e qodsiyye va maẓāher-e
ʿezz-e ṣamadiyye ast, maʿa-dhālek hīch edrāk nanamūde-and. Va īn moddat
īn qadr edrāk nashode ke maqṣūd az telāvat-e kotob va qerāʾat-e ṣoḥof
dar har ʿaṣrī edrāk-e maʿānī-ye ān va bolūgh be-maʿārej-e asrār-e ān
būde, va-ellā telāvat-e bī-maʿrefat rā albatte fāʾede-ye kollī nabāshad.

In short, one thousand two hundred and eighty years have passed since
the appearance of the Point of the Qurʾān, and all this ignorant rabble
have recited the Qurʾān every morning, yet have not grasped a single
letter of its purpose. They themselves read verses that explicitly
concern holy truths and the Manifestations of everlasting glory, yet
they have understood nothing. Throughout this time, they have not
understood even this much: that the purpose of reciting the Books and
reading the Scriptures in every age has been to comprehend their
meanings and ascend to the heights of their mysteries. Otherwise,
recitation without understanding can certainly yield no substantial
benefit.

\subsection{187}\label{section-377}

\begin{otherlanguage}{persian}

چنانچه شخصی در یومی نزد این فقیر بحر معانی حاضر بود و ذکری از علائم
قیامت و حشر و نشر و حساب بمیان آمد و اصرار مینمود که حساب خلایق در ظهور
بدیع چگونه شد که احدی اطّلاع نیافته بعد قدری از صور علمیّه و شئونات حکمیّه
بقدر ادراک و فهم سامع القا شد و بعد ذکر شد که این مدّت مگر تلاوت فرقان
ننموده‌ئی و آیهٔ مبارکه را که میفرماید فیومئذ لا یسأل عن ذنبه انس و لا جانّ
را ندید‌ه‌اید و بمقصود ملتفت نشده‌اید که معنی سوٴال چنان نیست که ادراک
نموده‌اید بلکه سوٴال بلسان و بیان نیست چنانچه همین آیه مشعر و مدلّ بر آنست
و بعد میفرماید یعرف المجرمون بسیماهم فیؤخذ بالنّواصی و الأقدام

\end{otherlanguage}

Chenānche shakhṣī dar yawmī nazd-e īn faqīr-e baḥr-e maʿānī ḥāżer būd va
dhekrī az ʿalāʾem-e qiyāmat va ḥashr va nashr va ḥesāb be-meyān āmad va
eṣrār mīnamūd ke ḥesāb-e khalāʾeq dar ẓohūr-e badīʿ chegūne shod ke
aḥadī eṭṭelāʿ nayāfte. Baʿd qadrī az ṣovar-e ʿelmiyye va shoʾūnāt-e
ḥekmiyye be-qadr-e edrāk va fahm-e sāmeʿ elqā shod va baʿd dhekr shod ke
īn moddat magar telāvat-e Forqān nanamūdeʾī va āye-ye mobārake rā ke
mīfarmāyad, fa-yawmaʾidhin lā yusʾalu ʿan dhanbihi insun wa-lā jānn, rā
nadīde-īd? Va be-maqṣūd moltafet nashode-īd ke maʿnī-ye soʾāl chenān
nīst ke edrāk namūde-īd, balke soʾāl be-lesān va bayān nīst, chenānche
hamīn āye moshʿer va modell bar ān ast. Va baʿd mīfarmāyad, yuʿrafu
al-mujrimūna bi-sīmāhum fa-yuʾkhadhu bi-al-nawāṣī wa-al-aqdām.

For example, one day someone was in the presence of this poor one, the
Ocean of meanings, and the conversation turned to the signs of the
Resurrection, the gathering, the raising of the dead, and the reckoning.
He kept insisting, ``How did the reckoning of all creatures take place
in this wondrous Revelation without anyone knowing of it?'' Some
conceptions of knowledge and aspects of wisdom were then imparted to him
according to his capacity to comprehend and understand. Afterward, he
was asked, ``Have you not been reciting the Qurʾān all this time? Have
you not seen the blessed verse that says, `On that day neither human nor
jinn will be questioned about his sin'?'' You have failed to grasp its
purpose: questioning does not mean what you have understood it to mean.
Rather, this questioning is not conducted through speech and words, as
this very verse indicates and attests. Then He says, ``The guilty will
be known by their countenances and seized by their forelocks and their
feet.''

\subsection{188}\label{section-378}

\begin{otherlanguage}{persian}

اینست که از وجهه حساب خلایق کشیده میشود و کفر و ایمان و عصیان جمیع ظاهر
میگردد چنانچه الیوم مشهود است که بسیما اهل ضلالت از اصحاب هدایت معلوم و
واضحند و اگر این عباد خالصاً للّه و طلباً لرضائه در آیات کتاب ملاحظه نمایند
جمیع آنچه را که میطلبند البتّه ادراک مینمایند بقسمی که جمیع امور واقعه در
این ظهور را از کلّی و جزئی در آیات او ظاهر و مکشوف ادراک مینمایند حتّی
خروج مظاهر اسماء و صفات را از اوطان و اعراض و اغماض ملّت و دولت را و سکون
و استقرار مظهر کلّیّه در ارض معلوم مخصوص ولکن لا یعرف ذلک الّا اولو الألباب

\end{otherlanguage}

Īn ast ke az vajhe ḥesāb-e khalāʾeq kashīde mīshavad va kofr va īmān va
ʿeṣyān-e jamīʿ ẓāher mīgardad, chenānche al-yawm mashhūd ast ke be-sīmā
ahl-e żalālat az aṣḥāb-e hedāyat maʿlūm va vāżeḥand. Va agar īn ʿebād
khāliṣan li-llāh va ṭalaban li-riḍāʾihi dar āyāt-e Ketāb molāḥeẓe
namāyand, jamīʿ-e ānche rā ke mīṭalaband albatte edrāk mīnamāyand,
be-qesmī ke jamīʿ-e omūr-e vāqeʿe dar īn ẓohūr rā az kollī va jozʾī dar
āyāt-e ū ẓāher va makshūf edrāk mīnamāyand, ḥattā khorūj-e maẓāher-e
asmā va ṣefāt rā az awṭān va eʿrāż va eghmāż-e mellat va dawlat rā va
sokūn va esteqrār-e Maẓhar-e kolliyye dar arż-e maʿlūm-e makhṣūṣ.
Wa-lākin lā yaʿrifu dhālika illā ulū al-albāb.

Thus it is by their countenances that all creatures are called to
account, and everyone's unbelief, faith, and disobedience become
apparent. So it is evident today that those in error are clearly
distinguished by their countenances from the companions of guidance. If
these servants were to ponder the verses of the Book solely for God's
sake and in pursuit of His pleasure, they would certainly understand all
that they seek. Indeed, they would perceive all the events of this
Revelation, both great and small, plainly disclosed in His verses---even
the departure of the manifestations of the names and attributes from
their homelands, the rejection and disregard of people and government,
and the dwelling and establishment of the universal Manifestation in a
particular, designated land. But none understands this except those
endowed with insight.

\subsection{189}\label{section-379}

\begin{otherlanguage}{arabic}

اختم القول بما نزّل علی محمّد من قبل لیکون ختامه المسک الّذی یهدی النّاس الی
رضوان قدس منیر قال و قوله الحقّ و اللّه یدعو الی دارالسّلام و یهدی من یشآء
الی صراط مستقیم لهم دارالسّلام عند ربّهم و هو ولیّهم بما کانوا یعملون لیسبق
هذا الفضل علی العالم و الحمد للّه ربّ العالمین

\end{otherlanguage}

Akhtimu al-qawla bi-mā nuzzila ʿalā Muḥammadin min qablu li-yakūna
khitāmuhu al-miska alladhī yahdī al-nāsa ilā riḍwāni qudsin munīr. Qāla
wa-qawluhu al-ḥaqq: wa-Allāhu yadʿū ilā dāri al-salāmi wa-yahdī man
yashāʾu ilā ṣirāṭin mustaqīm. Lahum dāru al-salāmi ʿinda rabbihim
wa-huwa waliyyuhum bi-mā kānū yaʿmalūn. Li-yasbiqa hādhā al-faḍlu ʿalā
al-ʿālam, wa-al-ḥamdu li-llāhi rabbi al-ʿālamīn.

I conclude this discourse with what was sent down to Muḥammad before, so
that its seal may be the musk that guides people to a radiant paradise
of holiness. He said---and His word is truth---``God calls to the abode
of peace and guides whomever He wills to a straight path.'' ``Theirs is
the abode of peace with their Lord, and He is their protector because of
what they have done.'' Thus may this grace take precedence throughout
the world. Praise be to God, the Lord of all worlds.

\subsection{190}\label{section-380}

\begin{otherlanguage}{persian}

بیان را در هر مطلب مکرّر نمودیم که شاید هر نفسی از اعالی و ادانی از این
بیانات بقدر و اندازهٔ خود قسمت و نصیب بردارد و اگر نفسی از ادراک بیانی
عاجز باشد از بیان دیگر مقصود خود را ادراک نماید لیعلم کلّ اناس مشربهم

\end{otherlanguage}

Bayān rā dar har maṭlab mokarrar namūdīm ke shāyad har nafsī az aʿālī va
adānī az īn bayānāt be-qadr va andāze-ye khod qesmat va naṣīb bardārad
va agar nafsī az edrāk-e bayānī ʿājez bāshad az bayān-e dīgar maqṣūd-e
khod rā edrāk namāyad, li-yaʿlama kullu unāsin mashrabahum.

We have repeated the explanation of every subject so that perhaps every
soul, whether high or low, may receive a share and portion of these
explanations according to their own capacity and measure. If someone is
unable to grasp one explanation, they may understand what they seek
through another, ``so that every people may know their drinking place.''

\subsection{191}\label{section-381}

\begin{otherlanguage}{persian}

قسم بخدا که این حمامهٔ ترابی را غیر این نغمات نغمه‌هاست و جز این بیانات
رموزها که هر نکته‌ئی از آن مقدّس است از آنچه بیان شد و از قلم جاری گشت تا
مشیّت الهی چه وقت قرار گیرد که عروسهای معانی بی‌حجاب از قصر روحانی قدم
ظهور بعرصهٔ قدم گذارند و ما من امر الّا بعد اذنه و ما من قدرة الّا بحوله و
قوّته و ما من اله الّا هو له الخلق و الأمر و کلّ بأمره ینطقون و من اسرار
الرّوح یتکلّمون

\end{otherlanguage}

Qasam be-Khodā ke īn ḥamāme-ye torābī rā gheyr-e īn naghamāt naghmehāst
va joz īn bayānāt romūzhā ke har nokteʾī az ān moqaddas ast az ānche
bayān shod va az qalam jārī gasht tā mashiyyat-e elāhī che vaqt qarār
gīrad ke ʿarūshā-ye maʿānī bī-ḥejāb az qaṣr-e rūḥānī qadam-e ẓohūr
be-ʿarṣe-ye qedam goẕārand. Wa-mā min amrin illā baʿda idhnihi wa-mā min
qudratin illā bi-ḥawlihi wa-quwwatihi wa-mā min ilāhin illā huwa lahu
al-khalqu wa-al-amru wa-kullun bi-amrihi yanṭiqūna wa-min asrāri al-rūḥi
yatakallamūn.

I swear by God that this earthly dove has melodies other than these
songs, and mysteries beyond these explanations, every subtle point of
which is sanctified beyond what has been explained and has flowed from
the pen. It remains to be seen when the divine will shall decree that
the brides of inner meanings step unveiled from their spiritual palace
into the arena of eternity. Nothing comes to pass except with His
permission, and there is no power except through His might and strength.
There is no God but Him. His are creation and command, and all speak at
His command and tell of the mysteries of the Spirit.

\subsection{192}\label{section-382}

\begin{otherlanguage}{persian}

از قبل دو مقام از برای شموس مشرقه از مشارق الهیّه بیان نمودیم یکی مقام
توحید و رتبهٔ تفرید چنانچه از قبل ذکر شد لا نفرّق بین احد منهم و مقام دیگر
مقام تفصیل و عالم خلق و رتبهٔ حدودات بشریّه است در این مقام هر کدام را
هیکلی معیّن و امری مقرّر و ظهوری مقدّر و حدودی مخصوص است چنانچه هر کدام
باسمی موسوم و بوصفی موصوف و بامری بدیع و شرعی جدید مأمورند چنانچه
میفرماید تلک الرّسل فضّلنا بعضهم علی بعض منهم من کلّم‌‌‌ اللّه ‌‌و رفع بعضهم
درجات و آتینا عیسی بن مریم البیّنات و ایّدناه بروح القدس

\end{otherlanguage}

Az qabl do maqām az barā-ye shomūs-e moshreqe az mashāreq-e elāhiyye
bayān namūdīm: yekī maqām-e tawḥīd va rotbe-ye tafrīd, chenānche az qabl
ẕekr shod, lā nufarriqu bayna aḥadin minhum; va maqām-e dīgar maqām-e
tafṣīl va ʿālam-e khalq va rotbe-ye ḥodūdāt-e bashariyye ast. Dar īn
maqām har kodām rā heykalī moʿayyan va amrī moqarrar va ẓohūrī moqaddar
va ḥodūdī makhṣūṣ ast, chenānche har kodām be-esmī mawsūm va be-vaṣfī
mawṣūf va be-amrī badīʿ va sharʿī jadīd maʾmūrand, chenānche mīfarmāyad:
tilka al-rusulu faḍḍalnā baʿḍahum ʿalā baʿḍin minhum man kallama Allāhu
wa-rafaʿa baʿḍahum darajātin wa-ātaynā ʿĪsā ibna Maryama al-bayyināti
wa-ayyadnāhu bi-rūḥi al-qudus.

We previously explained two stations belonging to the suns that rise
from the divine daysprings. One is the station of unity and the rank of
oneness, as was mentioned before: ``We make no distinction between any
of them.'' The other is the station of distinction, the world of
creation, and the rank of human limitations. In this station, each has a
particular form, an appointed cause, a predestined revelation, and
specific limitations. Thus each bears a name, is characterized by an
attribute, and is charged with a new cause and a new law. As He says:
``These Messengers---We have favored some above others. Among them are
those to whom God spoke, and some He raised in degree. We gave Jesus son
of Mary clear proofs and strengthened Him with the Holy Spirit.''

\subsection{193}\label{section-383}

\begin{otherlanguage}{persian}

نظر باختلاف این مراتب و مقاماتست که بیانات و کلمات مختلفه از آن ینابیع
علوم سبحانی ظاهر میشود والّا فی‌الحقیقه نزد عارفین معضلات مسائل الهیّه جمیع
در حکم یک کلمه مذکور است چون اکثر ناس اطّلاع بر مقامات مذکوره نیافته‌اند
اینست که در کلمات مختلفهٔ آن هیاکل متّحده مضطرب و متزلزل میشوند

\end{otherlanguage}

Naẓar be-ekhtelāf-e īn marāteb va maqāmāt ast ke bayānāt va kalemāt-e
mokhtalefe az ān yanābīʿ-e ʿolūm-e sobḥānī ẓāher mīshavad va-ellā fī
al-ḥaqīqe nazd-e ʿārefīn moʿḍalāt-e masāʾel-e elāhiyye jamīʿ dar ḥokm-e
yek kaleme maẕkūr ast. Chon akthar-e nās eṭṭelāʿ bar maqāmāt-e maẕkūre
nayāfte-and īn ast ke dar kalemāt-e mokhtalefe-ye ān hayākel-e mottaḥede
moẓṭareb va motazalzel mīshavand.

It is because these ranks and stations differ that diverse utterances
and words appear from those wellsprings of divine knowledge. Otherwise,
in truth, to those who understand the intricacies of divine questions,
they are all regarded as a single utterance. Since most people have not
become acquainted with these stations, they are troubled and shaken by
the differing words of those united Temples.

\subsection{194}\label{section-384}

\begin{otherlanguage}{persian}

باری معلوم بوده و خواهد بود که جمیع این اختلافات کلمات از اختلافات
مقاماتست اینست که در مقام توحید و علوّ تجرید اطلاق ربوبیّت و الوهیّت و
احدیّت صرفه و هویّهٔ بحته بر آن جواهر وجود شده و میشود زیرا که جمیع بر عرش
ظهور اللّه ساکنند و بر کرسیّ بطون ‌اللّه واقف یعنی ظهور اللّه بظهورشان ظاهر و
جمال ‌اللّه از جمالشان باهر چنانچه نغمات ربوبیّه از این هیاکل احدیّه ظاهر شد

\end{otherlanguage}

Bārī maʿlūm būde va khvāhad būd ke jamīʿ-e īn ekhtelāfāt-e kalemāt az
ekhtelāfāt-e maqāmāt ast. Īn ast ke dar maqām-e tawḥīd va ʿolovv-e
tajrīd eṭlāq-e robūbiyyat va olūhiyyat va aḥadiyyat-e ṣerfe va
hoviyye-ye baḥte bar ān javāher-e vojūd shode va mīshavad zīrā ke jamīʿ
bar ʿarsh-e ẓohūr-e Allāh sākenand va bar korsī-ye boṭūn-e Allāh vāqef,
yaʿnī ẓohūr-e Allāh be-ẓohūrashān ẓāher va jamāl-e Allāh az jamālashān
bāher, chenānche naghamāt-e robūbiyye az īn hayākel-e aḥadiyye ẓāher
shod.

In short, it has been and will remain clear that all these differences
in utterance arise from differences of station. Thus, in the station of
unity and the heights of detachment, lordship, divinity, pure oneness,
and unalloyed identity have been and are ascribed to those essences of
existence. For they all dwell upon the throne of God's revelation and
are established upon the seat of God's concealment. That is, God's
revelation is manifested through their revelation, and God's beauty
shines through their beauty. Thus the melodies of lordship have sounded
from these Temples of oneness.

\subsection{195}\label{section-385}

\begin{otherlanguage}{persian}

و در مقام ثانی که مقام تمیز و تفصیل و تحدید و اشارات و دلالات ملکیّه است
عبودیّت صرفه و فقر بحت و فنای باتّ از ایشان ظاهر است چنانچه میفرماید انّی
عبد اللّه و ما انا الّا بشر مثلکم

\end{otherlanguage}

Va dar maqām-e thānī ke maqām-e tamyīz va tafṣīl va taḥdīd va eshārāt va
dalālāt-e molkiyye ast ʿobūdiyyat-e ṣerfe va faqr-e baḥt va fanā-ye bātt
az īshān ẓāher ast, chenānche mīfarmāyad: innī ʿabdu Allāh, va mā anā
illā basharun mithlukum.

In the second station, which is the station of differentiation,
distinction, limitation, and earthly signs and indications, pure
servitude, utter poverty, and complete self-effacement are manifest in
them. As He says: ``I am the servant of God,'' and ``I am only a human
being like you.''

\subsection{196}\label{section-386}

\begin{otherlanguage}{arabic}

و از این بیانات محقّقهٔ مثبته ادراک فرما مسائل خود را که سوٴال نموده بودی
تا در دین الهی راسخ شوی و از اختلافات بیانات انبیا و اصفیا متزلزل نشوی

\end{otherlanguage}

Va az īn bayānāt-e moḥaqqaqe-ye mothbete edrāk farmā masāʾel-e khod rā
ke soʾāl namūde būdī tā dar dīn-e elāhī rāsekh shavī va az ekhtelāfāt-e
bayānāt-e anbiyā va aṣfiyā motazalzel nashavī.

From these verified and substantiated explanations, understand the
questions you had asked, so that you may become steadfast in the
religion of God and not be shaken by differences in the utterances of
the Prophets and chosen ones.

\subsection{197}\label{section-387}

\begin{otherlanguage}{persian}

و اگر شنیده شود از مظاهر جامعه انّی انا اللّه حقّ است و ریبی در آن نیست
چنانچه بکرّات مبرهن شد که بظهور و صفات و اسمای ایشان ظهور اللّه و اسم ‌اللّه
و صفة ‌اللّه در ارض ظاهر اینست که میفرماید و ما رمیت اذ رمیت ولکنّ ‌اللّه رمی
و همچنین انّ الّذین یبایعونک انّما یبایعون ‌اللّه و اگر نغمهٔ انّی رسول ‌اللّه
برآرند این نیز صحیح است و شکّی در آن نه چنانچه میفرماید ما کان محمّد ابا
احد من رجالکم ولکنّ رسول ‌اللّه و در این مقام همه مرسلند از نزد آن سلطان
حقیقی و کینونة ازلی و اگر جمیع ندای انا خاتم النّبیّین برآرند آن هم حقّ است
و شبهه را راهی نه و سبیلی نه زیرا که جمیع حکم یک ذات و یک نفس و یک روح و
یک جسد و یک امر دارند و همه مظهر بدئیّت و ختمیّت و اوّلیّت و آخریّت و ظاهریّت
و باطنیّت آن روح الأرواح حقیقی و ساذج السّواذج ازلیند و همچنین اگر
بفرمایند نحن عباد اللّه این نیز ثابت و ظاهر است چنانچه بظاهر در منتها
رتبهٔ عبودیّت ظاهر شده‌اند احدی را یارای آن نه که بآن نحو از عبودیّت در
امکان ظاهر شود اینست که از آن جواهر وجود در مقام استغراق در بحار قدس
صمدی و ارتقاء بمعارج معانی سلطان حقیقی اذکار ربوبیّه و الوهیّه ظاهر شد اگر
درست ملاحظه شود در همین رتبه منتهای نیستی و فنا در خود مشاهده نموده‌اند
در مقابل هستی مطلق و بقای صرف که گویا خود را معدوم صرف دانسته‌اند و ذکر
خود را در آن ساحت شرک شمرده‌اند زیرا که مطلق ذکر در این مقام دلیل هستی و
وجود است و این نزد واصلان بس خطا چه جای آنکه ذکر غیر شود و قلب و لسان و
دل و جان بغیر ذکر جانان مشغول گردد و یا چشم غیر جمال او ملاحظه نماید و
یا گوش غیر نغمهٔ او شنود و یا رجل در غیر سبیل او مشی نماید

\end{otherlanguage}

Va agar shenīde shavad az maẓāher-e jāmeʿe innī anā Allāh, ḥaqq ast va
raybī dar ān nīst, chenānche be-karrāt mobarhan shod ke be-ẓohūr va
ṣefāt va asmā-ye īshān ẓohūr-e Allāh va esm-e Allāh va ṣefat-e Allāh dar
arḍ ẓāher. Īn ast ke mīfarmāyad: wa-mā ramayta idh ramayta wa-lākinna
Allāha ramā, va hamchenīn: inna alladhīna yubāyiʿūnaka innamā yubāyiʿūna
Allāh. Va agar naghme-ye innī rasūlu Allāh barārand īn nīz ṣaḥīḥ ast va
shakkī dar ān na, chenānche mīfarmāyad: mā kāna Muḥammadun abā aḥadin
min rijālikum wa-lākin rasūla Allāh. Va dar īn maqām hame morsaland az
nazd-e ān solṭān-e ḥaqīqī va kaynūnat-e azalī. Va agar jamīʿ nedā-ye anā
khātamu al-nabiyyīn barārand ān ham ḥaqq ast va shobhe rā rāhī na va
sabīlī na zīrā ke jamīʿ ḥokm-e yek ẕāt va yek nafs va yek rūḥ va yek
jasad va yek amr dārand va hame maẓhar-e badʾiyyat va khatmiyyat va
avvaliyyat va ākhiriyyat va ẓāhiriyyat va bāṭiniyyat-e ān rūḥ al-arvāḥ-e
ḥaqīqī va sāẕaj al-savāẕaj-e azalī-and. Va hamchenīn agar befarmāyand
naḥnu ʿibādu Allāh īn nīz thābet va ẓāher ast chenānche be-ẓāher dar
montahā rotbe-ye ʿobūdiyyat ẓāher shode-and aḥadī rā yārā-ye ān na ke
be-ān naḥv az ʿobūdiyyat dar emkān ẓāher shavad. Īn ast ke az ān
javāher-e vojūd dar maqām-e esteghrāq dar beḥār-e qods-e ṣamadī va
erteqā be-maʿārej-e maʿānī-ye solṭān-e ḥaqīqī aẕkār-e robūbiyye va
olūhiyye ẓāher shod. Agar dorost molāḥaẓe shavad dar hamīn rotbe
montahā-ye nīstī va fanā dar khod moshāhede namūde-and dar moqābel-e
hastī-ye moṭlaq va baqā-ye ṣerf ke gūyā khod rā maʿdūm-e ṣerf
dāneste-and va ẕekr-e khod rā dar ān sāḥat sherk shemorde-and zīrā ke
moṭlaq-e ẕekr dar īn maqām dalīl-e hastī va vojūd ast va īn nazd-e
vāṣelān bas khaṭā che jā-ye ānke ẕekr-e gheyr shavad va qalb va lesān va
del va jān be-gheyr-e ẕekr-e jānān mashghūl gardad va yā chashm gheyr-e
jamāl-e ū molāḥaẓe namāyad va yā gūsh gheyr-e naghme-ye ū shenavad va yā
rejl dar gheyr-e sabīl-e ū mashy namāyad.

If the all-embracing Manifestations are heard to say, ``I am God,'' it
is true, and there is no doubt in it. For it has repeatedly been
demonstrated that through their revelation, attributes, and names, the
revelation of God, the name of God, and the attribute of God are
manifested on earth. Thus He says: ``It was not you who threw when you
threw, but God who threw,'' and likewise: ``Those who pledge allegiance
to you are only pledging allegiance to God.'' If they raise the song,
``I am the Messenger of God,'' this too is correct and beyond doubt. As
He says: ``Muḥammad is not the father of any of your men, but the
Messenger of God.'' In this station, they are all sent by that true
Sovereign and eternal Being. If they all proclaim, ``I am the Seal of
the Prophets,'' that too is true, and doubt has no way or access to it,
for they are all regarded as one essence, one self, one spirit, one
body, and one cause. They all manifest the beginning and the end, the
firstness and the lastness, the outwardness and the inwardness of that
true Spirit of spirits and eternal Essence of essences. Likewise, if
they say, ``We are the servants of God,'' this too is established and
manifest, for outwardly they have appeared in the utmost degree of
servitude. No one in the contingent world has the capacity to manifest
such servitude. Thus, while immersed in the oceans of eternal holiness
and ascending the heights of the inner meanings of the true Sovereign,
those essences of existence have voiced declarations of lordship and
divinity. If this is considered carefully, in this very rank they have
perceived in themselves the utmost nothingness and self-effacement
before absolute existence and pure permanence, as though they regarded
themselves as utter nonexistence and considered any mention of
themselves in that court to be the association of partners with God. For
in this station, any mention at all is evidence of being and existence,
and to those who have attained, this is a grave error. How much more so
if anyone else were mentioned, and heart and tongue, soul and spirit
were occupied with anything other than remembrance of the Beloved, or if
the eye beheld anything but His beauty, the ear heard anything but His
melody, or the foot walked in any path but His.

\subsection{198}\label{section-388}

\begin{otherlanguage}{arabic}

در این زمان نسمة اللّه وزیده و روح ‌اللّه احاطه نموده قلم از حرکت ممنوع و
لسان از بیان مقطوع گشته

\end{otherlanguage}

Dar īn zamān nasamat Allāh vazīde va rūḥ Allāh eḥāṭe namūde qalam az
ḥarekat mamnūʿ va lesān az bayān maqṭūʿ gashte.

At this moment the breeze of God has blown and the Spirit of God has
encompassed all. The pen has been restrained from moving, and the tongue
has been stopped from speaking.

\subsection{199}\label{section-389}

\begin{otherlanguage}{persian}

باری نظر باین مقام ذکر ربوبیّه و امثال ذلک از ایشان ظاهر شده و در مقام
رسالت اظهار رسالت فرمودند و همچنین در هر مقام باقتضای آن ذکری فرمودند و
همه را نسبت بخود داده‌اند از عالم امر الی عالم خلق و از عوالم ربوبیّه الی
عوالم ملکیّه اینست که آنچه بفرمایند و هر چه ذکر نمایند از الوهیّت و ربوبیّت
و نبوّت و رسالت و ولایت و امامت و عبودیّت همه حقّست و شبهه در آن نیست پس
باید تفکّر در این بیانات که استدلال شده نمود تا دیگر از اختلافات اقوال
مظاهر غیبیّه و مطالع قدسیّه احدی را اضطراب و تزلزل دست ندهد

\end{otherlanguage}

Bārī naẓar be-īn maqām ẕekr-e robūbiyye va amthāl-e ẕālek az īshān ẓāher
shode va dar maqām-e resālat eẓhār-e resālat farmūdand va hamchenīn dar
har maqām be-eqteżā-ye ān ẕekrī farmūdand va hame rā nesbat be-khod
dāde-and az ʿālam-e amr elā ʿālam-e khalq va az ʿavālem-e robūbiyye elā
ʿavālem-e molkiyye. Īn ast ke ānche befarmāyand va har che ẕekr namāyand
az olūhiyyat va robūbiyyat va nobovvat va resālat va velāyat va emāmat
va ʿobūdiyyat hame ḥaqq ast va shobhe dar ān nīst. Pas bāyad tafakkor
dar īn bayānāt ke estedlāl shode namūd tā dīgar az ekhtelāfāt-e aqvāl-e
maẓāher-e gheybiyye va maṭāleʿ-e qodsiyye aḥadī rā eżṭerāb va tazalzol
dast nadahad.

In short, it is with reference to this station that declarations of
lordship and the like have issued from them, while in the station of
messengership they have declared their messengership. Likewise, in every
station they have spoken as that station required and have ascribed all
these things to themselves, from the world of command to the world of
creation, and from the realms of lordship to the earthly realms. Thus,
whatever they say and whatever they declare concerning divinity,
lordship, prophethood, messengership, guardianship, leadership, and
servitude is all true and beyond doubt. One must therefore reflect on
these explanations that have been offered as proof, so that differences
in the sayings of the Manifestations of the unseen and the Daysprings of
holiness may no longer cause anyone distress or uncertainty.

\subsection{200}\label{section-390}

\begin{otherlanguage}{persian}

باری در کلمات شموس حقیقت باید تفکّر نمود و اگر ادراک نشد باید از واقفین
مخازن علم سوٴال شود تا بیان فرمایند و رفع اشکال نمایند نه آنکه بعقل ناقص
خود کلمات قدسیّه را تفسیر نمایند و چون مطابق نفس و هوای خود نیابند بنای
رد و اعتراض گذارند چنانچه الیوم علما و فقهای عصر که بر مسند علم و فضل
نشسته‌اند و جهل را علم نام گذاشته‌اند و ظلم را عدل نامیده‌اند اگر مجعولات
خاطر خود را از شمس حقیقی سوٴال نمایند و جواب موافق آنچه فهمیده و یا از
کتاب مثل خود ادراک نموده‌اند نشنوند البتّه نفی علم از آن معدن و منبع علم
نمایند چنانچه در هر زمانی این واقع شد

\end{otherlanguage}

Bārī dar kalemāt-e shomūs-e ḥaqīqat bāyad tafakkor namūd va agar edrāk
nashod bāyad az vāqefīn-e makhāzen-e ʿelm soʾāl shavad tā bayān
farmāyand va rafʿ-e eshkāl namāyand na ānke be-ʿaql-e nāqeṣ-e khod
kalemāt-e qodsiyye rā tafsīr namāyand va chon moṭābeq-e nafs va havā-ye
khod nayāband banā-ye radd va eʿterāż goẕārand. Chenānche al-yawm ʿolamā
va foqahā-ye ʿaṣr ke bar masnad-e ʿelm va fażl nesheste-and va jahl rā
ʿelm nām goẕāshte-and va ẓolm rā ʿadl nāmīde-and agar majʿūlāt-e
khāṭer-e khod rā az shams-e ḥaqīqī soʾāl namāyand va javāb movāfeq-e
ānche fahmīde va yā az ketāb-e methl-e khod edrāk namūde-and nashnavand
al-batte nafī-ye ʿelm az ān maʿdan va manbaʿ-e ʿelm namāyand chenānche
dar har zamānī īn vāqeʿ shod.

In short, one must reflect on the words of the Suns of Truth. If they
are not understood, one must ask those acquainted with the treasuries of
knowledge to explain them and resolve the difficulty, rather than
interpret the sacred words through one's own deficient reason and, on
finding them contrary to one's own desires and inclinations, proceed to
reject and challenge them. Thus today, the scholars and jurists of the
age, who occupy the seats of learning and excellence, who have named
ignorance knowledge and called oppression justice, may question the true
Sun about the fabrications of their own minds. If they do not hear an
answer that agrees with what they have understood or gathered from the
books of others like themselves, they will certainly deny knowledge to
that mine and source of knowledge, as has happened in every age.

\subsection{201}\label{section-391}

\begin{otherlanguage}{persian}

مثل اینکه مذکور شد در سوٴال از اهلّه که از سیّد وجود نمودند و آن حضرت بامر
الهی جواب فرمود که هی مواقیت للنّاس بعد از استماع نفی علم از آن حضرت
نمودند

\end{otherlanguage}

Methl-e īnke maẕkūr shod dar soʾāl az ahelle ke az sayyed-e vojūd
namūdand va ān ḥażrat be-amr-e elāhī javāb farmūd ke hiya mawāqītu
li-l-nās, baʿd az estemāʿ nafī-ye ʿelm az ān ḥażrat namūdand.

For example, as was mentioned concerning the question about the new
moons that they put to the Lord of existence, He answered at God's
command, ``They are measures of time for people.'' After hearing this,
they denied that He possessed knowledge.

\subsection{202}\label{section-392}

\begin{otherlanguage}{persian}

و همچنین در آیهٔ روح که میفرماید و یسألونک عن الرّوح قل الرّوح من امر ربّی و
چون این جواب مذکور شد کل فریاد واویلا برآوردند که جاهلی که نمیداند روح
چه چیز است خود را عالم علم لدنّی میداند و الیوم چون علمای عصر باسم آن
حضرت مفتخرند و آبای خود را هم مذعن دیده‌اند لهذا تقلیداً حکمش را قبول
دارند چنانچه اگر انصاف باشد و الیوم در جواب امثال این مسائل چنین جواب
بشنوند البتّه رد نمایند و اعتراض کنند و همان سخنهای قبل را اعاده نمایند
چنانچه نمودند با اینکه آن جواهر وجود مقدّسند از کلّ این علمهای مجعوله و
منزّهند از جمیع این کلمات محدوده و متعالیند از ادراک هر مدرکی کلّ این علوم
نزد آن علم کذب صرفست و جمیع این ادراکات افک محض بلکه هر چه از آن معادن
حکمت الهی و مخازن علم صمدانی ظاهر میشود علم همانست و العلم نقطة کثّرها
الجاهلون دلیل بر آن و العلم نور یقذفه‌ اللّه فی قلب من یشآء مثبت این بیان

\end{otherlanguage}

Va hamchenīn dar āye-ye rūḥ ke mīfarmāyad: wa-yasʾalūnaka ʿani al-rūḥi
quli al-rūḥu min amri rabbī. Va chon īn javāb maẕkūr shod koll faryād-e
vāveylā barāvardand ke jāhelī ke nemīdānad rūḥ che chīz ast khod rā
ʿālem-e ʿelm-e ladonnī mīdānad. Va al-yawm chon ʿolamā-ye ʿaṣr be-esm-e
ān ḥażrat moftakherand va ābā-ye khod rā ham moẕʿen dīde-and le-hāẕā
taqlīdan ḥokmash rā qabūl dārand. Chenānche agar enṣāf bāshad va al-yawm
dar javāb-e amthāl-e īn masāʾel chonīn javāb beshnavand al-batte radd
namāyand va eʿterāż konand va hamān sokhanhā-ye qabl rā eʿāde namāyand
chenānche namūdand. Bā īnke ān javāher-e vojūd moqaddasand az koll-e īn
ʿelmhā-ye majʿūle va monazzahand az jamīʿ-e īn kalemāt-e maḥdūde va
motaʿālī-and az edrāk-e har modrekī. Koll-e īn ʿolūm nazd-e ān ʿelm
keẕb-e ṣerf ast va jamīʿ-e īn edrākāt efk-e maḥż balke har che az ān
maʿāden-e ḥekmat-e elāhī va makhāzen-e ʿelm-e ṣamadānī ẓāher mīshavad
ʿelm hamān ast, va al-ʿilmu nuqṭatun kaththarahā al-jāhilūn dalīl bar ān
va al-ʿilmu nūrun yaqdhifuhu Allāhu fī qalbi man yashāʾ mothbet-e īn
bayān.

Likewise, in the verse concerning the Spirit, He says: ``They ask you
about the Spirit. Say: `The Spirit is of the command of my Lord.'\,''
When this answer was given, they all raised an outcry, saying, ``An
ignorant man who does not know what the Spirit is claims to possess
knowledge bestowed directly by God!'' Today, because the scholars of the
age take pride in His name and have seen their forefathers acknowledge
Him, they accept His pronouncement through imitation. Yet, to judge
fairly, if today they heard such an answer to questions of this kind,
they would certainly reject it, object to it, and repeat those same
earlier statements, as indeed they have done. Yet those essences of
existence are sanctified beyond all these fabricated sciences, free from
all these limited words, and exalted beyond the understanding of anyone
who understands. All these sciences are sheer falsehood beside that
knowledge, and all these perceptions are pure delusion. Rather,
knowledge is whatever issues from those mines of divine wisdom and
treasuries of eternal knowledge. ``Knowledge is a point that the
ignorant have multiplied'' is evidence of this, and ``Knowledge is a
light that God casts into the heart of whomever He wills'' confirms this
statement.

\subsection{203}\label{section-393}

\begin{otherlanguage}{persian}

باری چون معنی علم را ادراک ننموده‌اند و افکار مجعولهٔ خود را که ناشی از
مظاهر جهل شده اسم آن را علم گذاشته بر مبدء علوم وارد آورده‌اند آنچه
دیده‌اید و شنیده‌اید

\end{otherlanguage}

Bārī chon maʿnī-ye ʿelm rā edrāk nanamūde-and va afkār-e majʿūle-ye khod
rā ke nāshī az maẓāher-e jahl shode esm-e ān rā ʿelm goẕāshte bar
mabdaʾ-e ʿolūm vāred āvarde-and ānche dīde-īd va shenīde-īd.

In short, because they have not understood the meaning of knowledge and
have given the name of knowledge to their own fabricated ideas, which
arise from manifestations of ignorance, they have inflicted upon the
Source of all knowledge what you have seen and heard.

\subsection{204}\label{section-394}

\begin{otherlanguage}{persian}

مثلاً در کتاب یکی از عباد که مشهور بعلم و فضل است و خود را از صنادید قوم
شمرده و جمیع علمای راشدین را رد و سب نموده چنانچه در همه جای از کتاب او
تلویحاً و تصریحاً مشهود است و این بنده چون ذکر او را بسیار شنیده بودم
اراده نمودم که از رسائل او قدری ملاحظه نمایم هرچند این بنده اقبال
بملاحظهٔ کلمات غیر نداشته و ندارم ولیکن چون جمعی از احوال ایشان سوٴال
نموده و مستفسر شده بودند لهذا لازم گشت که قدری در کتب او ملاحظه رود و
جواب سائلین بعد از معرفت و بصیرت داده شود باری کتب عربیّهٔ او بدست نیفتاد
تا اینکه شخصی روزی ذکر نمود که کتابی از ایشان که مسمّی به ارشادالعوامست
در این بلد یافت میشود اگرچه از این اسم رائحهٔ کبر و غرور استشمام شد که
مردم را عوام و خود را عالم فرض نموده و جمیع مراتب او فی‌الحقیقه از همین
اسم کتاب معلوم و مبرهن شد که در سبیل نفس و هوی سالکند و در تیه جهل و عمی
ساکن گویا حدیث مشهور را فراموش نموده‌اند که میفرماید العلم تمام المعلوم و
القدرة و العزّة تمام الخلق با وجود این کتاب را طلب نموده چند روز معدود
نزد بنده بود و گویا دو مرتبه در او ملاحظه شد از قضا مرتبهٔ ثانی جائی بدست
آمد که حکایت معراج سیّد لولاک بود ملاحظه شد که قریب بیست علم او ازید شرط
معرفت معراج نوشته‌اند و همچو مستفاد شد که اگر نفسی این علوم را درست ادراک
ننموده باشد بمعرفت این امر عالی متعالی فائز نگردد و از جملهٔ علوم علم
فلسفه و علم کیمیا و علم سیمیا را مذکور نموده و ادراک این علوم فانیهٔ
مردوده را شرط ادراک علوم باقیهٔ قدسیّه شمرده

\end{otherlanguage}

Mathalan dar kitāb-e yekī az ʿibād ke mashhūr be-ʿilm va faḍl ast va
khod rā az ṣanādīd-e qawm shomorde va jamīʿ-e ʿolamā-ye rāshidīn rā radd
va sabb namūde chenānche dar hame jāy az kitāb-e ū talvīḥan va taṣrīḥan
mashhūd ast va īn bande chon dhekr-e ū rā besyār shenīde būdam erāde
namūdam ke az rasāʾel-e ū qadrī molāḥeẓe namāyam harchand īn bande eqbāl
be-molāḥeẓe-ye kalemāt-e gheyr nadāshte va nadāram valīken chon jamʿī az
aḥvāl-e īshān soʾāl namūde va mostafser shode būdand lehādhā lāzem gasht
ke qadrī dar kotob-e ū molāḥeẓe ravad va javāb-e sāʾelīn baʿd az
maʿrefat va baṣīrat dāde shavad bārī kotob-e ʿarabiyye-ye ū be-dast
nayoftād tā īnke shakhṣī rūzī dhekr namūd ke kitābī az īshān ke mosammā
be Irshād al-ʿawāmm ast dar īn balad yāft mīshavad agarche az īn esm
rāʾeḥe-ye kebr va ghorūr estesḥām shod ke mardom rā ʿavāmm va khod rā
ʿālem farḍ namūde va jamīʿ-e marāteb-e ū fī-l-ḥaqīqe az hamīn esm-e
kitāb maʿlūm va mobarhan shod ke dar sabīl-e nafs va havā sālekand va
dar tīh-e jahl va ʿamā sāken gūyā ḥadīth-e mashhūr rā farāmūsh
namūde-and ke mīfarmāyad al-ʿilmu tamāmu al-maʿlūmi wa-l-qudratu
wa-l-ʿizzatu tamāmu al-khalq bā vojūd-e īn kitāb rā ṭalab namūde chand
rūz-e maʿdūd nazd-e bande būd va gūyā do martabe dar ū molāḥeẓe shod az
qaḍā martabe-ye thānī jāʾī be-dast āmad ke ḥekāyat-e miʿrāj-e sayyed-e
lawlāka būd molāḥeẓe shod ke qarīb-e bīst ʿilm aw azyad sharṭ-e
maʿrefat-e miʿrāj neveshte-and va hamcho mostafād shod ke agar nafsī īn
ʿolūm rā dorost edrāk nanamūde bāshad be-maʿrefat-e īn amr-e ʿālī-ye
motaʿālī fāʾez nagardad va az jomle-ye ʿolūm ʿilm-e falsafe va ʿilm-e
kīmiyā va ʿilm-e sīmiyā rā madhkūr namūde va edrāk-e īn ʿolūm-e
fāniye-ye mardūde rā sharṭ-e edrāk-e ʿolūm-e bāqiye-ye qodsiyye shomorde

For example, consider the book of one of the servants who is renowned
for learning and accomplishment, who has counted himself among the
leaders of the people and has rejected and reviled all the rightly
guided scholars, as is evident throughout his book, both implicitly and
explicitly. As I, this servant, had heard much about him, I decided to
examine some of his treatises. Although this servant has never been
inclined, nor am I now, to examine the words of others, a number of
people had asked and inquired about him. It therefore became necessary
to examine his books to some extent, so that their questions might be
answered with knowledge and insight. His Arabic books, however, could
not be obtained, until one day someone mentioned that a book of his
entitled ``Guidance for the Common People'' could be found in this city.
Although the odor of pride and conceit could be detected in this title,
since he had assumed the people to be common and himself learned, and
although his entire standing was in truth made clear and manifest by
this very title---that he walks the path of self and desire and dwells
in the wilderness of ignorance and blindness, apparently having
forgotten the well-known tradition that says, ``Knowledge is the
entirety of the known, and power and might are the entirety of
creation''---I nevertheless requested the book. It remained with this
servant for a few days, and I believe I examined it twice. On the second
occasion, I happened upon a passage recounting the ascension of the Lord
of ``But for You.'' I observed that he had set down some twenty sciences
or more as prerequisites for understanding the ascension. The
implication was that unless a person had properly grasped these
sciences, he could not attain an understanding of this lofty and
transcendent matter. Among these sciences he mentioned philosophy,
alchemy, and occult science, and counted an understanding of these
perishable and rejected sciences as a prerequisite for understanding the
eternal and sacred sciences.

\subsection{205}\label{section-395}

\begin{otherlanguage}{persian}

سبحان‌اللّه با این ادراک چه اعتراضات و تهمتها که بهیاکل علم نامتناهی الهی
وارد آورده فنعم ما قال متّهم داری کسانی را که حقّ کرد امین مخزن هفتم طبق و
یک نفر از اهل بصیرت و دانش و صاحبان علوم و عقول ملتفت این مزخرفات نشده
با اینکه بر هر صاحب بصیرتی واضح و هویداست که اینگونه علمها لم‌یزل مردود
حقّ بوده و هست و چگونه علومی که مردود است نزد علمای حقیقی ادراک آن شرط
ادراک معارج معراج میشود با اینکه صاحب معراج حرفی از این علوم محدودهٔ
محجوبه حمل نفرموده و قلب منیر آن سیّد لولاک از جمیع این اشارات مقدّس و
منزّه بوده چه خوب میگوید جمله ادراکات بر خرهای لنگ حقّ سوار باد پرّان چون
خدنگ واللّه هر کس بخواهد سرّ معراج را ادراک نماید و یا قطره‌ئی از عرفان این
بحر بنوشد اگر هم این علوم نزد او باشد یعنی مرآت قلب او از نقوش این علوم
غبار گرفته باشد البتّه باید پاک و منزّه نماید تا سرّ این امر در مرآت قلب او
تجلّی نماید

\end{otherlanguage}

Subḥāna Allāh bā īn edrāk che eʿterāḍāt va tohmathā ke be-hayākel-e
ʿilm-e nāmotanāhī-ye elāhī vāred āvarde fa-niʿma mā qāl motّaham dārī
kasānī rā ke Ḥaqq kard amīn-e makhzan-e haftom ṭabaq va yek nafar az
ahl-e baṣīrat va dānesh va ṣāḥebān-e ʿolūm va ʿoqūl moltafet-e īn
mozakhrafāt nashode bā īnke bar har ṣāḥeb-e baṣīratī vāḍeḥ va hovaydāst
ke īngūne ʿelmhā lam-yazal mardūd-e Ḥaqq būde va hast va chegūne ʿolūmī
ke mardūd ast nazd-e ʿolamā-ye ḥaqīqī edrāk-e ān sharṭ-e edrāk-e
maʿārej-e miʿrāj mīshavad bā īnke ṣāḥeb-e miʿrāj ḥarfī az īn ʿolūm-e
maḥdūde-ye maḥjūbe ḥaml nafarmūde va qalb-e monīr-e ān sayyed-e lawlāka
az jamīʿ-e īn eshārāt moqaddas va monazzah būde che khūb mīgūyad jomle
edrākāt bar kharhā-ye lang Ḥaqq savār-e bād parrān chon khadang va-llāh
har kas bekhāhad serr-e miʿrāj rā edrāk namāyad va yā qaṭre-ʾī az
ʿerfān-e īn baḥr benūshad agar ham īn ʿolūm nazd-e ū bāshad yaʿnī
merʾāt-e qalb-e ū az noqūsh-e īn ʿolūm ghobār gerefte bāshad albatte
bāyad pāk va monazzah namāyad tā serr-e īn amr dar merʾāt-e qalb-e ū
tajallī namāyad

Glory be to God! With such an understanding, what objections and
accusations he has directed against the Embodiments of God's infinite
knowledge! How well it has been said: ``You accuse those whom God has
made the trusted custodians of the treasury of the seventh sphere.'' Yet
not one person of insight and learning, not one possessed of knowledge
and understanding, has taken notice of this nonsense, although it is
clear and evident to everyone with insight that such sciences have
always been, and remain, rejected by God. How can an understanding of
sciences rejected by the true scholars become a prerequisite for
understanding the heights of the ascension, when He who ascended bore
not a single letter of these limited and veiled sciences, and the
radiant heart of that Lord of ``But for You'' was sanctified and
purified of all these allusions? How well it is said: ``All
understandings ride on lame donkeys; Truth rides the wind, flying like
an arrow.'' By God, whoever wishes to grasp the mystery of the ascension
or drink a drop of the knowledge of this ocean, even if he possesses
these sciences---that is, if the mirror of his heart has gathered dust
from the impressions of these sciences---must assuredly cleanse and
purify it, so that the mystery of this matter may be revealed in the
mirror of his heart.

\subsection{206}\label{section-396}

\begin{otherlanguage}{persian}

و الیوم متغمّسان بحر علوم صمدانی و ساکنان فلک حکمت ربّانی مردم را از تحصیل
این علوم نهی میفرمایند و صدور منیرشان بحمد اللّه منزّه از این اشارات است و
مقدّس از این حجبات حجاب اکبر را که میفرماید العلم حجاب الأکبر بنار محبّت
یار سوختیم و خیمهٔ دیگر برافراختیم و باین افتخار مینمائیم که الحمد للّه
سبحات جلال را بنار جمال محبوب سوختیم و جز مقصود در قلب و دل جا ندادیم نه
بعلمی جز علم باو متمسّکیم و نه بمعلومی جز تجلّی انوار او متشبّث

\end{otherlanguage}

Va al-yawm motaghammesān-e baḥr-e ʿolūm-e ṣamadānī va sākenān-e folk-e
ḥekmat-e rabbānī mardom rā az taḥṣīl-e īn ʿolūm nahy mīfarmāyand va
ṣodūr-e monīreshān be-ḥamd-e Allāh monazzah az īn eshārāt ast va
moqaddas az īn ḥojobāt ḥejāb-e akbar rā ke mīfarmāyad al-ʿilmu ḥijābu
al-akbar be-nār-e maḥabbat-e yār sūkhtīm va kheyme-ye dīgar barafrākhtīm
va be-īn eftekhār mīnamāyīm ke al-ḥamdu li-llāh soboḥāt-e jalāl rā
be-nār-e jamāl-e maḥbūb sūkhtīm va joz maqṣūd dar qalb va del jā nadādīm
na be-ʿelmī joz ʿilm-e be-ū motamassekīm va na be-maʿlūmī joz tajallī-ye
anvār-e ū motashabbeth

Today, those immersed in the ocean of divine knowledge and those who
dwell in the ark of heavenly wisdom forbid people to acquire these
sciences. Their radiant hearts, praise be to God, are free from these
allusions and sanctified from these veils. We have burned the greatest
veil, of which it is said, ``Knowledge is the greatest veil,'' in the
fire of love for the Friend, and have raised another tent. We take pride
in this: praise be to God, we have burned the veils of glory in the fire
of the Beloved's beauty and have given none but the Desired One a place
in our inmost heart. We hold to no knowledge but knowledge of Him, and
cling to nothing known but the revelation of His lights.

\subsection{207}\label{section-397}

\begin{otherlanguage}{persian}

باری بسیار متعجّب شدم در این بیانات ندیدم مگر اینکه میخواهد بر مردم
برساند که جمیع این علوم نزد ایشانست با وجود اینکه قسم بخدا نسیمی از ریاض
علم الهی نشنیده و بر حرفی از اسرار حکمت ربّانی اطّلاع نیافته بلکه اگر معنی
علم گفته شود البتّه مضطرب شود و جبل وجود او مندک گردد با وجود این اقوال
سخیفهٔ بی‌معنی چه دعویهای زیاده از حدّ نموده

\end{otherlanguage}

Bārī besyār motaʿajjeb shodam dar īn bayānāt nadīdam magar īnke mīkhāhad
bar mardom berasānad ke jamīʿ-e īn ʿolūm nazd-e īshānst bā vojūd-e īnke
qasam be-Khodā nasīmī az riyāḍ-e ʿilm-e elāhī nashenīde va bar ḥarfī az
asrār-e ḥekmat-e rabbānī eṭṭelāʿ nayāfte balke agar maʿnī-ye ʿilm gofte
shavad albatte moḍṭareb shavad va jabal-e vojūd-e ū mondakk gardad bā
vojūd-e īn aqvāl-e sakhīfe-ye bī-maʿnī che daʿvīhā-ye ziyāde az ḥadd
namūde

In short, I was greatly astonished. In these statements I saw nothing
but a desire to convey to people that he possesses all these sciences,
even though, I swear by God, he has not sensed a breeze from the gardens
of divine knowledge or learned a single letter of the mysteries of
heavenly wisdom. Indeed, were the meaning of knowledge to be explained,
he would surely be thrown into turmoil and the mountain of his being
would crumble. Despite these foolish, meaningless statements, what
extravagant claims he has made!

\subsection{208}\label{section-398}

\begin{otherlanguage}{persian}

سبحان‌اللّه چه ‌قدر متعجّبم از مردمی که باو گرویده‌اند و تابع چنین شخصی
گشته‌اند بتراب قناعت نموده و اقبال جسته‌اند و از ربّ الأرباب معرض گشته‌اند و
از نغمهٔ بلبل و جمال گل بنعیب زاغ و جمال کلاغ قناعت نموده‌اند و دیگر چه
چیزها ملاحظه شد از کلمات مجعولهٔ این کتاب فی‌الحقیقه حیف است که قلم بتحریر
ذکر آن مطالب مشغول شود و یا اوقات مصروف بآن گردد ولیکن اگر محکّی یافت
میشد حقّ از باطل و نور از ظلمت و شمس از ظلّ معلوم میآمد

\end{otherlanguage}

Subḥāna Allāh che qadr motaʿajjebam az mardomī ke be-ū geravīde-and va
tābeʿ-e chenīn shakhṣī gashte-and be-torāb qanāʿat namūde va eqbāl
joste-and va az Rabb al-arbāb moʿreḍ gashte-and va az naghme-ye bolbol
va jamāl-e gol be-naʿīb-e zāgh va jamāl-e kalāgh qanāʿat namūde-and va
dīgar che chīzhā molāḥeẓe shod az kalemāt-e majʿūle-ye īn kitāb
fī-l-ḥaqīqe ḥeyf ast ke qalam be-taḥrīr-e dhekr-e ān maṭāleb mashghūl
shavad va yā awqāt maṣrūf be-ān gardad valīken agar meḥakkī yāft mīshod
Ḥaqq az bāṭel va nūr az ẓolmat va shams az ẓell maʿlūm mīāmad

Glory be to God! How astonished I am at the people who have put their
faith in him and become followers of such a person! They have contented
themselves with dust and turned toward it, while turning away from the
Lord of lords. In place of the nightingale's song and the beauty of the
rose, they have contented themselves with the raven's croaking and the
crow's beauty. And what other things were observed among the fabricated
statements of this book! It is truly a pity for the pen to be occupied
with recording such matters, or for time to be spent on them. Yet if a
touchstone could be found, truth would be distinguished from falsehood,
light from darkness, and the sun from the shadow.

\subsection{209}\label{section-399}

\begin{otherlanguage}{persian}

از جمله علومی که این مرد مدّعی بآن شده صنعت کیمیاست بسیار طالبم که سلطانی
و یا نفسی که مقتدر باشد ظهور این علم را از عالم لفظ بعالم شهود و از قول
بفعل از ایشان طلب نماید و این بی‌علم فانی هم که دعوی اینگونه علوم ننموده
و بلکه کون این علوم و فقدان آن را علّت علم و جهل نمیدانم با این مرد در
همین فقره قیام نمائیم تا صدق و کذب معلوم شود ولیکن چه فائده از ناس این
زمان جز زخم سنان ندیده‌ام و غیر سمّ قاتل چیزی نچشیده‌ام هنوز اثر حدید بر
گردن باقیست و هنوز علائم جفا از تمام بدن ظاهر

\end{otherlanguage}

Az jomle ʿolūmī ke īn mard moddaʿī be-ān shode ṣanʿat-e kīmiyāst besyār
ṭālebam ke solṭānī va yā nafsī ke moqtader bāshad ẓohūr-e īn ʿilm rā az
ʿālam-e lafẓ be-ʿālam-e shohūd va az qawl be-feʿl az īshān ṭalab namāyad
va īn bī-ʿilm-e fānī ham ke daʿvī-ye īngūne ʿolūm nanamūde va balke
kawn-e īn ʿolūm va foqdān-e ān rā ʿellat-e ʿilm va jahl namīdānam bā īn
mard dar hamīn faqare qiyām namāyīm tā ṣedq va kedhb maʿlūm shavad
valīken che fāʾede az nās-e īn zamān joz zakhm-e senān nadīde-am va
gheyr-e samm-e qātel chīzī nacheshīde-am hanūz athar-e ḥadīd bar gardan
bāqīst va hanūz ʿalāʾem-e jafā az tamām-e badan ẓāher

Among the sciences this man has claimed to possess is the art of
alchemy. I greatly desire that a sovereign, or someone with the power to
do so, should require him to bring this science from the realm of words
into the realm of visible reality, and from speech into action. Let this
unlearned, mortal one, who has made no claim to such sciences---and who,
indeed, does not regard their possession or absence as the basis of
knowledge or ignorance---undertake this very matter alongside this man,
so that truth and falsehood may become clear. But what use is this? From
the people of this age I have experienced nothing but spear wounds and
tasted nothing but deadly poison. The mark of the iron still remains
upon my neck, and the signs of cruelty are still visible over my whole
body.

\subsection{210}\label{section-400}

\begin{otherlanguage}{persian}

و در مراتب علم و جهل و عرفان و ایقان او در کتابی که ترک نشد از آن امری
ذکر شده اینست که میفرماید انّ شجرة الزّقّوم طعام الأثیم و بعد بیانات دیگر
میفرماید تا اینکه منتهی میشود باین ذکر ذق انّک انت العزیز الکریم ملتفت
شوید که چه واضح و صریح وصف او در کتاب محکم مذکور شده و این شخص هم خود را
در کتاب خود از بابت خفض جناح عبد اثیم ذکر نموده اثیم فی الکتاب و عزیز
بین الأنعام و کریم فی الاسم

\end{otherlanguage}

Va dar marāteb-e ʿilm va jahl va ʿerfān va īqān-e ū dar kitābī ke tark
nashod az ān amrī dhekr shode īnast ke mīfarmāyad inna shajarata
al-zaqqūmi ṭaʿāmu al-athīm va baʿd bayānāt-e dīgar mīfarmāyad tā īnke
montahī mīshavad be-īn dhekr dhuq innaka anta al-ʿazīzu al-karīm
moltafet shavīd ke che vāḍeḥ va ṣarīḥ vaṣf-e ū dar kitāb-e moḥkam
madhkūr shode va īn shakhṣ ham khod rā dar kitāb-e khod az bābat-e
khafḍ-e janāḥ ʿabd-e athīm dhekr namūde athīm fī al-kitāb va ʿazīz bayna
al-anʿām va karīm fī al-ism

His degrees of knowledge and ignorance, recognition and certitude, are
mentioned in the Book from which nothing has been left out. Thus He
says, ``Indeed, the tree of Zaqqum is the food of the sinner.'' He then
makes other statements until He concludes with these words: ``Taste! You
are indeed the mighty, the generous.'' Consider how clearly and
explicitly his description is recorded in the unerring Book. This
person, too, by way of humility, has called himself a sinful servant in
his own book: sinful in the Book, mighty among the cattle, and generous
in name.

\subsection{211}\label{section-401}

\begin{otherlanguage}{persian}

تفکّر در آیهٔ مبارکه نموده تا معنی و لا رطب و لا یابس الّا فی کتاب مبین
درست در لوح قلب ثبت شود با وجود این جمعی معتقد او شده و از موسی علم و
عدل اعراض نموده به سامریّ جهل تمسّک جسته‌اند و از شمس معانی که در سماء
لایزالی الهی مشرقست معرض گشته‌اند و کأن لم یکن انگاشته‌اند

\end{otherlanguage}

Tafakkor dar āye-ye mobārake namūde tā maʿnī-ye wa-lā raṭbin wa-lā
yābisin illā fī kitābin mubīn dorost dar lawḥ-e qalb thabt shavad bā
vojūd-e īn jamʿī moʿtaqed-e ū shode va az Mūsā-ye ʿilm va ʿadl eʿrāḍ
namūde be Sāmirī-ye jahl tamassok joste-and va az shams-e maʿānī ke dar
samā-ye lāyazālī-ye elāhī moshreq ast moʿreḍ gashte-and va ka-an lam
yakun angāshte-and

Reflect upon the blessed verse, so that the meaning of ``There is
nothing moist or dry but it is in a clear Book'' may be firmly inscribed
upon the tablet of your heart. Despite this, some have come to believe
in him and, turning away from the Moses of knowledge and justice, have
clung to the Samiri of ignorance. They have turned away from the Sun of
inner meanings, shining in the heaven of God's eternity, and have
regarded it as though it did not exist.

\subsection{212}\label{section-402}

\begin{otherlanguage}{persian}

باری ای برادر من لآلی علم ربّانی جز از معدن الهی بدست نیاید و رائحهٔ ریحان
معنوی جز از گلزار حقیقی استشمام نشود و گلهای علوم احدیّه جز از مدینهٔ قلوب
صافیه نروید و البلد الطّیّب یخرج نباته باذن ربّه و الّذی خبث لا یخرج الّا
نکدا

\end{otherlanguage}

Bārī ey barādar-e man laʾālī-ye ʿilm-e rabbānī joz az maʿdan-e elāhī
be-dast nayāyad va rāʾeḥe-ye reyḥān-e maʿnavī joz az golzār-e ḥaqīqī
estesḥām nashavad va golhā-ye ʿolūm-e aḥadiyye joz az madīne-ye qolūb-e
ṣāfiye narūyad wa-l-baladu al-ṭayyibu yakhruju nabātuhu bi-idhni rabbihi
wa-lladhī khabutha lā yakhruju illā nakidā

In short, my brother, the pearls of heavenly knowledge can be obtained
only from the divine mine, the fragrance of the spiritual herb can be
inhaled only from the garden of truth, and the flowers of the knowledge
of divine unity grow only in the city of pure hearts. ``The good land
brings forth its vegetation by its Lord's permission, while that which
is poor brings forth only scantily.''

\subsection{213}\label{section-403}

\begin{otherlanguage}{arabic}

و چون مفهوم گشت که تغنّیات ورقاء هویّه را احدی ادراک ننماید الّا اهلش لهذا
بر هر نفسی لازم و واجبست که مشکلات مسائل الهیّه و معضلات اشارات مطالع
قدسیّه را بر صاحبان افئدهٔ منیره و حاملان اسرار احدیّه عرضه دارد تا
بتأییدات ربّانی و افاضات الهی حلّ مسائل شود نه بتأییدات علوم اکتسابی
فاسألوا اهل الذّکر ان کنتم لا تعلمون

\end{otherlanguage}

Va chon mafhūm gasht ke taghanniyāt-e varqā-ye hoviyye rā aḥadī edrāk
nanamāyad illā ahlash lehādhā bar har nafsī lāzem va vājeb ast ke
moshkelāt-e masāʾel-e elāhiyye va moʿḍalāt-e eshārāt-e maṭāleʿ-e
qodsiyye rā bar ṣāḥebān-e afʾede-ye monīre va ḥāmelān-e asrār-e aḥadiyye
ʿarḍe dārad tā be-taʾyīdāt-e rabbānī va efāḍāt-e elāhī ḥall-e masāʾel
shavad na be-taʾyīdāt-e ʿolūm-e ektesābī fa-sʾalū ahla al-dhikri in
kuntum lā taʿlamūn

Since it has become clear that no one can understand the songs of the
Dove of divine Being except those qualified to do so, it is necessary
and incumbent upon every soul to submit the difficulties of divine
questions and the intricacies of the allusions of the holy Dawning
Places to those whose hearts are illumined and who bear the mysteries of
divine unity, so that these questions may be resolved through heavenly
assistance and divine outpourings, not through the support of acquired
learning. ``Ask the people of remembrance if you do not know.''

\subsection{214}\label{section-404}

\begin{otherlanguage}{persian}

ولیکن ای برادر من شخص مجاهد که اراده نمود قدم طلب و سلوک در سبیل معرفت
سلطان قدم گذارد باید در بدایت امر قلب را که محلّ ظهور و بروز تجلّی اسرار
غیبی الهی است از جمیع غبارات تیرهٔ علوم اکتسابی و اشارات مظاهر شیطانی پاک
و منزّه فرماید و صدر را که سریر ورود و جلوس محبّت محبوب ازلی است لطیف و
نظیف نماید و همچنین دل را از علاقهٔ آب و گل یعنی از جمیع نقوش شبحیّه و صور
ظلّیّه مقدّس گرداند بقسمی که آثار حبّ و بغض در قلب نماند که مبادا آن حبّ او
را بجهتی بیدلیل میل دهد و یا بغض او را از جهتی منع نماید چنانچه الیوم
اکثری باین دو وجه از وجهۀ باقی و حضرت معانی بازمانده‌اند و بی‌شبان در
صحراهای ضلالت و نسیان میچرند و باید در کلّ حین توکّل بحقّ نماید و از خلق
اعراض کند و از عالم تراب منقطع شود و بگسلد و بربّ الأرباب دربندد و نفس
خود را بر احدی ترجیح ندهد و افتخار و استکبار را از لوح قلب بشوید و بصبر
و اصطبار دل بندد و صمت را شعار خود نماید و از تکلّم بی‌فائده احتراز کند چه
زبان ناریست افسرده و کثرت بیان سمّی است هلاک‌کننده نار ظاهری اجساد را
محترق نماید و نار لسان ارواح و افئده را بگدازد اثر آن نار بساعتی فانی
شود و اثر این نار بقرنی باقی ماند

\end{otherlanguage}

Valīken ey barādar-e man shakhṣ-e mojāhed ke erāde namūd qadam-e ṭalab
va solūk dar sabīl-e maʿrefat-e Solṭān-e qidam godhārad bāyad dar
bedāyat-e amr qalb rā ke maḥall-e ẓohūr va borūz-e tajallī-ye asrār-e
gheybī-ye elāhī ast az jamīʿ-e ghobārāt-e tīre-ye ʿolūm-e ektesābī va
eshārāt-e maẓāher-e sheyṭānī pāk va monazzah farmāyad va ṣadr rā ke
sarīr-e vorūd va jolūs-e maḥabbat-e maḥbūb-e azalī ast laṭīf va naẓīf
namāyad va hamchenīn del rā az ʿalāqe-ye āb va gel yaʿnī az jamīʿ-e
noqūsh-e shabaḥiyye va ṣovar-e ẓelliyye moqaddas gardānad be-qesmī ke
āthār-e ḥobb va boghḍ dar qalb namānad ke mabādā ān ḥobb ū rā be-jahatī
bī-dalīl meyl dahad va yā boghḍ ū rā az jahatī manʿ namāyad chenānche
al-yawm aktharī be-īn do vajh az vejhe-ye bāqī va Ḥaḍrat-e maʿānī
bāzmānde-and va bī-shabān dar ṣaḥrāhā-ye ḍalālat va nesyān mīcharand va
bāyad dar koll-e ḥīn tavakkol be-Ḥaqq namāyad va az khalq eʿrāḍ konad va
az ʿālam-e torāb monqaṭeʿ shavad va bogselad va be Rabb al-arbāb
darbandad va nafs-e khod rā bar aḥadī tarjīḥ nadahad va eftekhār va
estekbār rā az lawḥ-e qalb beshūyad va be-ṣabr va eṣṭebār del bandad va
ṣamt rā sheʿār-e khod namāyad va az takallom-e bī-fāʾede eḥterāz konad
che zabān nārīst afsorde va kathrat-e bayān sammī ast halāk-konande
nār-e ẓāherī ajsād rā moḥtareq namāyad va nār-e lesān arvāḥ va afʾede rā
begodāzad athar-e ān nār be-sāʿatī fānī shavad va athar-e īn nār
be-qarnī bāqī mānad

But, my brother, when a seeker striving earnestly resolves to set foot
upon the path of search and spiritual wayfaring toward the knowledge of
the King of eternity, he must first cleanse and purify his heart---the
place where the hidden divine mysteries are revealed and made
manifest---from all the dark dust of acquired learning and the allusions
of satanic manifestations. He must refine and cleanse his breast, the
throne upon which love for the eternal Beloved enters and takes its
seat. He must likewise sanctify his heart from attachment to water and
clay---that is, from all phantom impressions and shadowy forms---so that
no trace of love or hatred remains in his heart, lest that love incline
him in some direction without reason, or hatred prevent him from taking
a direction. Today, indeed, most people have through these two things
been kept from the eternal Goal and the Presence of inner meanings, and
wander without a shepherd in the deserts of error and forgetfulness. At
every moment he must put his trust in God, turn away from creation,
detach and sever himself from the world of dust, and bind himself to the
Lord of lords. He must not prefer himself to anyone, must wash pride and
arrogance from the tablet of his heart, devote himself to patience and
steadfastness, make silence his practice, and avoid useless speech. For
the tongue is a smoldering fire, and excessive speech a deadly poison.
Outward fire burns bodies, while the fire of the tongue consumes souls
and hearts. The effects of that fire vanish in an hour, while the
effects of this fire endure for a century.

\subsection{215}\label{section-405}

\begin{otherlanguage}{persian}

و غیبت را ضلالت شمرد و بآن عرصه هرگز قدم نگذارد زیرا غیبت سراج منیر قلب
را خاموش نماید و حیات دل را بمیراند بقلیل قانع باشد و از طلب کثیر فارغ
مصاحبت منقطعین را غنیمت شمارد و عزلت از متمسّکین و متکبّرین را نعمت شمرد
در اسحار باذکار مشغول شود و بتمام همّت و اقتدار در طلب آن نگار کوشد غفلت
را بنار حبّ و ذکر بسوزاند و از ما سوی اللّه چون برق درگذرد و بر بی‌نصیبان
نصیب بخشد و از محرومان عطا و احسان دریغ ندارد رعایت حیوان را منظور نماید
تا چه رسد بانسان و اهل بیان و از جانان جان دریغ ندارد و از شماتت خلق از
حقّ احتراز نجوید و آنچه برای خود نمی‌پسندد برای غیر نپسندد و نگوید آنچه را
وفا نکند و از خاطئان در کمال استیلا درگذرد و طلب مغفرت نماید و بر عاصیان
قلم عفو درکشد و بحقارت ننگرد زیرا حسن خاتمه مجهول است ای بسا عاصی که در
حین موت بجوهر ایمان موفّق شود و خمر بقا چشد و بملأ اعلی شتابد و بسا مطیع
و مؤمن که در وقت ارتقای روح تقلیب شود و باسفل درکات نیران مقرّ یابد باری
مقصود از جمیع این بیانات متقنه و اشارات محکمه آنست که سالک و طالب باید
جز خدا را فنا داند و غیر معبود را معدوم شمرد

\end{otherlanguage}

Va ghībat rā ḍalālat shomorad va be ān ʿarṣe hargaz qadam nagodhārad
zīrā ghībat serāj-e monīr-e qalb rā khāmūsh namāyad va ḥayāt-e del rā
bemīrānad. Be qalīl qāneʿ bāshad va az ṭalab-e kathīr fāregh.
Moṣāḥebat-e monqaṭeʿīn rā ghanīmat shomārad va ʿozlat az motamassekīn va
motakabberīn rā neʿmat shomorad. Dar asḥār be adhkār mashghūl shavad va
be tamām-e hemmat va eqtedār dar ṭalab-e ān negār kūshad. Ghaflat rā be
nār-e ḥobb va dhekr besūzānad va az mā sivā Allāh chon barq dargodharad
va bar bī-naṣībān naṣīb bakhshad va az maḥrūmān ʿaṭā va eḥsān darīgh
nadārad. Reʿāyat-e ḥeyvān rā manẓūr namāyad tā che resad be ensān va
ahl-e bayān. Va az jānān jān darīgh nadārad va az shamātat-e khalq az
Ḥaqq eḥterāz najūyad va ānche barā-ye khod nemīpasandad barā-ye gheyr
napasandad va nagūyad ānche rā vafā nakonad va az khāṭeʾān dar kamāl-e
estīlā dargodharad va ṭalab-e maghferat namāyad va bar ʿāṣiyān qalam-e
ʿafv darkeshad va be ḥeqārat nanegarad zīrā ḥosn-e khāteme majhūl ast.
Ey basā ʿāṣī ke dar ḥīn-e mowt be jowhar-e īmān movaffaq shavad va
khamr-e baqā cheshad va be malaʾ-e aʿlā shetābad va basā moṭīʿ va moʾmen
ke dar vaqt-e erteqā-ye rūḥ taqlīb shavad va be asfal-e darakāt-e nīrān
maqarr yābad. Bārī maqṣūd az jamīʿ-e īn bayānāt-e motqane va eshārāt-e
moḥkame ānast ke sālek va ṭāleb bāyad joz Khodā rā fanā dānad va gheyr-e
Maʿbūd rā maʿdūm shomorad.

He must regard backbiting as error and never set foot in its domain, for
backbiting extinguishes the heart's shining lamp and destroys its life.
He must be content with little and free from the desire for more. He
must prize the company of the detached and count seclusion from the
attached and arrogant a blessing. At dawn he must occupy himself with
remembrance and strive with all his resolve and strength to seek that
Beloved. He must burn away heedlessness with the fire of love and
remembrance, and pass like lightning beyond all but God. He must give a
share to those who have none and withhold neither gifts nor kindness
from the deprived. He must show consideration for animals, how much more
for human beings and those endowed with speech. He must not withhold his
life from the Beloved or shun the Truth for fear of people's reproaches.
He must not wish for others what he does not wish for himself, nor say
what he will not fulfill. Even when fully in power, he must pardon those
who have erred and seek forgiveness for them. He must draw the pen of
pardon over sinners and not look upon them with contempt, for no one
knows whether their end will be good. How many a sinner, at the moment
of death, attains the essence of faith, tastes the wine of everlasting
life, and hastens to the company on high; and how many an obedient
believer undergoes a reversal when his spirit ascends and finds his
abode in the lowest depths of the fires of hell. In short, the purpose
of all these well-founded statements and firm allusions is that the
wayfarer and seeker must regard all but God as passing away and count
all but the One who is worshipped as nonexistent.

\subsection{216}\label{section-406}

\begin{otherlanguage}{persian}

و این شرایط از صفات عالین و سجیّهٔ روحانیّین است که در شرایط مجاهدین و مشی
سالکین در مناهج علم الیقین ذکر یافت و بعد از تحقّق این مقامات برای سالک
فارغ و طالب صادق لفظ مجاهد در بارهٔ او صادق میآید و چون بعمل و الّذین
جاهدوا فینا مؤیّد شد البتّه ببشارت لنهدینّهم سبلنا مستبشر خواهد شد

\end{otherlanguage}

Va īn sharāyeṭ az ṣefāt-e ʿālīn va sajiyye-ye rūḥāniyyīn ast ke dar
sharāyeṭ-e mojāhedīn va mashy-e sālekīn dar manāhej-e ʿelm al-yaqīn
dhekr yāft. Va baʿd az taḥaqqoq-e īn maqāmāt barā-ye sālek-e fāregh va
ṭāleb-e ṣādeq lafẓ-e mojāhed dar bāre-ye ū ṣādeq mīāyad va chon be
ʿamal-e ``wa-alladhīna jāhadū fīnā'' moʾayyad shod albatte be beshārat-e
``la-nahdiyannahum subulanā'' mostabsher khwāhad shod.

These conditions belong to the attributes of the exalted and the
character of the spiritually minded, and have been mentioned in
connection with the requirements for those who strive and the progress
of wayfarers along the paths of the knowledge of certainty. Once these
stations have been realized in the detached wayfarer and sincere seeker,
the term ``one who strives'' truly applies to him. When he is confirmed
in acting according to ``those who strive in Us,'' he will assuredly
rejoice in the glad tidings, ``We will surely guide them in Our ways.''

\subsection{217}\label{section-407}

\begin{otherlanguage}{persian}

و چون سراج طلب و مجاهده و ذوق و شوق و عشق و وله و جذب و حبّ در قلب روشن
شد و نسیم محبّت از شطر احدیّه وزید ظلمت ضلالت شک و ریب زایل شود و انوار
علم و یقین همهٔ ارکان وجود را احاطه نماید در آن حین بشیر معنوی ببشارت
روحانی از مدینهٔ الهی چون صبح صادق طالع شود و قلب و نفس و روح را بصور
معرفت از نوم غفلت بیدار نماید و عنایات و تأییدات روح القدس صمدانی حیات
تازهٔ جدید مبذول دارد بقسمی که خود را صاحب چشم جدید و گوش بدیع و قلب و
فؤاد تازه می‌بیند و رجوع بآیات واضحهٔ آفاقیّه و خفیّات مستورهٔ انفسیّه مینماید
و بعین ‌اللّه بدیعه در هر ذرّه بابی مفتوح مشاهده نماید برای وصول بمراتب عین
الیقین و حقّ الیقین و نور الیقین و در جمیع اشیاء اسرار تجلّی وحدانیّه و
آثار ظهور صمدانیّه ملاحظه کند

\end{otherlanguage}

Va chon serāj-e ṭalab va mojāhede va dhowq va showq va ʿeshq va valah va
jadhb va ḥobb dar qalb rowshan shod va nasīm-e maḥabbat az shaṭr-e
aḥadiyye vazīd ẓolmat-e ḍalālat-e shakk va reyb zāyel shavad va anvār-e
ʿelm va yaqīn hame-ye arkān-e vojūd rā eḥāṭe namāyad. Dar ān ḥīn
bashīr-e maʿnavī be beshārat-e rūḥānī az madīne-ye elāhī chon ṣobḥ-e
ṣādeq ṭāleʿ shavad va qalb va nafs va rūḥ rā be ṣūr-e maʿrefat az nowm-e
ghaflat bīdār namāyad va ʿenāyāt va taʾyīdāt-e Rūḥ al-Qodos-e ṣamadānī
ḥayāt-e tāze-ye jadīd mabdhūl dārad be qesmī ke khod rā ṣāḥeb-e chashm-e
jadīd va gūsh-e badīʿ va qalb va foʾād-e tāze mībīnad va rojūʿ be āyāt-e
vāḍeḥe-ye āfāqiyye va khafiyyāt-e mastūre-ye anfosiyye mīnamāyad va be
ʿeyn-Allāh-e badīʿe dar har dharre bābī maftūḥ moshāhede namāyad barā-ye
voṣūl be marāteb-e ʿayn al-yaqīn va ḥaqq al-yaqīn va nūr al-yaqīn va dar
jamīʿ-e ashyā asrār-e tajallī-ye vaḥdāniyye va āthār-e ẓohūr-e
ṣamadāniyye molāḥeẓe konad.

When the lamp of search and striving, delight and longing, love and
rapture, attraction and affection is lit in the heart, and the breeze of
love blows from the realm of oneness, the darkness of the error of doubt
and uncertainty will vanish, and the lights of knowledge and certainty
will encompass every part of his being. At that moment, the herald of
inner meaning will rise like the true dawn from the divine city with
spiritual glad tidings, and with the trumpet of recognition awaken the
heart, self, and spirit from the sleep of heedlessness. The favors and
confirmations of the eternal Holy Spirit will bestow a fresh, new life,
so that he will find himself possessed of a new eye, a wondrous ear, and
a renewed heart and inmost being. He will turn to the clear signs in the
world and the concealed mysteries within the soul, and with the wondrous
eye of God will see in every atom an open door through which to attain
the stations of the vision of certainty, the truth of certainty, and the
light of certainty. In all things he will behold the mysteries of the
revelation of divine unity and the traces of the manifestation of the
Eternal.

\subsection{218}\label{section-408}

\begin{otherlanguage}{persian}

قسم بخدا که اگر سالک سبیل هدی و طالب معارج تقی باین مقام بلند اعلی واصل
گردد رائحهٔ حقّ را از فرسنگهای بعیده استنشاق نماید و صبح نورانی هدایت را
از مشرق کلّ شیء ادراک کند و هر ذرّه و هر شیء او را دلالت بر محبوب و مطلوب
نماید و چنان ممیّز شود که حقّ را از باطل چون شمس از ظلّ فرق گذارد مثلاً اگر
نسیم حقّ از مشرق ابداع وزد و او در مغرب اختراع باشد البتّه استشمام کند و
همچنین جمیع آثار حقّ را از کلمات بدیعه و اعمال منیعه و افعال لمیعه از
افعال و اعمال و آثار ماسوی امتیاز دهد چنانچه اهل لؤلؤ لؤلؤ را از حجر و
انسان ربیع را از خریف و حرارت را از برودت و دماغ جان چون از زکام کون و
امکان پاک شد البتّه رائحهٔ جانان را از منازل بعیده بیابد و از اثر آن رائحه
بمصر ایقان حضرت منّان وارد شود و بدایع حکمت حضرت سبحانی را در آن شهر
روحانی مشاهده کند و جمیع علوم مکنونه را از اطوار ورقهٔ شجرهٔ آن مدینه
استماع نماید و از تراب آن مدینه تسبیح و تقدیس ربّ الأرباب بگوش ظاهر و
باطن شنود و اسرار رجوع و ایاب را بچشم سر ملاحظه فرماید چه ذکر نمایم از
آثار و علامات و ظهورات و تجلّیات که بامر سلطان اسماء و صفات در آن مدینه
مقدّر شده بی‌آب رفع عطش نماید و بی‌نار حرارت محبّة اللّه بیفزاید در هر گیاهی
حکمت بالغهٔ معنوی مستور است و بر شاخسار هر گل هزار بلبل ناطقه در جذب و
شور از لاله‌های بدیعش سرّ نار موسوی ظاهر و از نفحات قدسیّه‌اش نفخهٔ روح القدس
عیسوی باهر بی‌ذهب غنا بخشد و بی‌فنا بقا عطا فرماید در هر ورقش نعیمی مکنون
و در هر غرفه‌اش صد‌هزار حکمت مخزون

\end{otherlanguage}

Qasam be Khodā ke agar sālek-e sabīl-e hodā va ṭāleb-e maʿārej-e toqā be
īn maqām-e boland-e aʿlā vāṣel gardad rāʾeḥe-ye Ḥaqq rā az farsanghā-ye
baʿīde estenshāq namāyad va ṣobḥ-e nūrānī-ye hedāyat rā az mashreq-e
koll-e shayʾ edrāk konad va har dharre va har shayʾ ū rā delālat bar
Maḥbūb va Maṭlūb namāyad va chenān momayyez shavad ke Ḥaqq rā az bāṭel
chon shams az ẓell farq godhārad. Mathalan agar nasīm-e Ḥaqq az
mashreq-e ebdāʿ vazad va ū dar maghreb-e ekhterāʿ bāshad albatte
esteshmām konad va hamchenīn jamīʿ-e āthār-e Ḥaqq rā az kalemāt-e badīʿe
va aʿmāl-e manīʿe va afʿāl-e lamīʿe az afʿāl va aʿmāl va āthār-e māsivā
emtiyāz dahad chenānche ahl-e loʾloʾ loʾloʾ rā az ḥajar va ensān rabīʿ
rā az kharīf va ḥarārat rā az borūdat. Va demāgh-e jān chon az zokām-e
kown va emkān pāk shod albatte rāʾeḥe-ye Jānān rā az manāzel-e baʿīde
beyābad va az athar-e ān rāʾeḥe be Meṣr-e īqān-e Ḥażrat-e Mannān vāred
shavad va badāyeʿ-e ḥekmat-e Ḥażrat-e Sobḥānī rā dar ān shahr-e rūḥānī
moshāhede konad va jamīʿ-e ʿolūm-e maknūne rā az aṭvār-e varaqe-ye
shajare-ye ān madīne estemāʿ namāyad va az torāb-e ān madīne tasbīḥ va
taqdīs-e Rabb al-Arbāb be gūsh-e ẓāher va bāṭen shenavad va asrār-e
rojūʿ va eyāb rā be chashm-e sar molāḥeẓe farmāyad. Che dhekr namāyam az
āthār va ʿalāmāt va ẓohūrāt va tajalliyāt ke be amr-e Solṭān-e asmā va
ṣefāt dar ān madīne moqaddar shode. Bī-āb rafʿ-e ʿaṭash namāyad va
bī-nār ḥarārat-e maḥabbat Allāh bīfzāyad. Dar har giyāhī ḥekmat-e
bāleghe-ye maʿnavī mastūr ast va bar shākhsār-e har gol hazār bolbol-e
nāṭeqe dar jadhb va shūr. Az lālehā-ye badīʿash serr-e nār-e Mūsavī
ẓāher va az nafaḥāt-e qodsiyye-ash nafkhe-ye Rūḥ al-Qodos-e ʿĪsavī
bāher. Bī-dhahab ghenā bakhshad va bī-fanā baqā ʿaṭā farmāyad. Dar har
varaqash naʿīmī maknūn va dar har ghorfe-ash ṣad-hezār ḥekmat makhzūn.

I swear by God that if the wayfarer on the path of guidance and the
seeker of the heights of piety attains this lofty and most exalted
station, he will inhale the fragrance of the Truth from many distant
leagues and perceive the radiant dawn of guidance from the dawning place
of every thing. Every atom and every thing will direct him to the
Beloved, the One he seeks. His discernment will become such that he will
distinguish truth from falsehood as he distinguishes the sun from
shadow. For example, if the breeze of Truth blows from the east of
creation while he is in its west, he will assuredly perceive its
fragrance. Likewise, he will distinguish all the signs of the
Truth---wondrous words, exalted deeds, and radiant acts---from the acts,
deeds, and signs of all others, just as those who know pearls
distinguish a pearl from a stone, and a person distinguishes spring from
autumn and heat from cold. When the soul's sense of smell is cleared of
the congestion of contingent existence, it will assuredly detect the
fragrance of the Beloved from distant places. Through the influence of
that fragrance, he will enter the Egypt of the certainty of the
All-Bountiful, behold the wonders of the wisdom of the All-Glorious in
that spiritual city, and hear all hidden knowledge from the manifold
expressions of a leaf of that city's tree. With his outward and inward
ears he will hear from the dust of that city the praise and
glorification of the Lord of lords, and with his own eyes behold the
mysteries of return and homecoming. What can I say of the signs, tokens,
manifestations, and revelations ordained in that city by the command of
the Sovereign of names and attributes? Without water it quenches thirst,
and without fire it increases the heat of the love of God. In every
plant a consummate spiritual wisdom lies concealed, and on the branches
of every flower a thousand singing nightingales are enraptured and
astir. From its wondrous tulips the mystery of the fire of Moses is
revealed, and from its holy fragrances the breath of the Holy Spirit of
Jesus shines forth. Without gold it bestows wealth, and without passing
away it grants everlasting life. In every leaf a delight lies hidden,
and in every chamber a hundred thousand wisdoms are stored.

\subsection{219}\label{section-409}

\begin{otherlanguage}{arabic}

و مجاهدین فی‌‌اللّه بعد از انقطاع از ماسوی چنان بآن مدینه انس گیرند که آنی
از آن منفکّ نشوند دلائل قطعیّه را از سنبل آن محفل شنوند و براهین واضحه را
از جمال گل و نوای بلبل اخذ نمایند و این مدینه در رأس هزار سنه او ازید او
اقلّ تجدید شود و تزیین یابد

\end{otherlanguage}

Va mojāhedīn fī Allāh baʿd az enqeṭāʿ az māsivā chenān be ān madīne ons
gīrand ke ānī az ān monfakk nashavand. Dalāʾel-e qaṭʿiyye rā az sonbol-e
ān maḥfel shenavand va barāhīn-e vāḍeḥe rā az jamāl-e gol va navā-ye
bolbol akhdh namāyand va īn madīne dar raʾs-e hazār sane aw azyad aw
aqall tajdīd shavad va tazyīn yābad.

After detaching themselves from all else, those who strive in God become
so closely attached to that city that they do not part from it for a
moment. They hear conclusive proofs from the hyacinth of that gathering
and draw clear demonstrations from the beauty of the rose and the song
of the nightingale. This city is renewed and adorned at the end of a
thousand years, or more or less.

\subsection{220}\label{section-410}

\begin{otherlanguage}{persian}

پس ای حبیب من باید جهدی نمود تا بآن مدینه واصل شویم و بعنایت الهیّه و
تفقّدات ربّانیّه کشف سبحات جلال نمائیم تا باستقامت تمام جان پژمرده را در ره
محبوب تازه نثار نمائیم و صدهزار عجز و نیاز آریم تا بآن فوز فائز شویم و
آن مدینه کتب الهیّه است در هر عهدی مثلاً در عهد موسی تورات بود و در زمن
عیسی انجیل و در عهد محمّد رسول اللّه فرقان و در این عصر بیان و در عهد من
یبعثه اللّه کتاب او که رجوع کلّ کتب بآنست و مهیمن است بر جمیع کتب و در این
مداین ارزاق مقدّر است و نعم باقیه مقرّر غذای روحانی بخشد و نعمت قدمانی
چشاند بر اهل تجرید نعمت توحید عطا فرماید بی‌نصیبان را نصیب کرم نماید و
آوارگان صحرای جهل را کأس علم عنایت کند و هدایت و عنایت و علم و معرفت و
ایمان و ایقان کلّ من فی السّموات و الأرض در این مدائن مکنون و مخزون گشته

\end{otherlanguage}

Pas ey ḥabīb-e man bāyad jahdī namūd tā be ān madīne vāṣel shavīm va be
ʿenāyat-e elāhiyye va tafaq qodāt-e rabbāniyye kashf-e soboḥāt-e jalāl
namāʾīm tā be esteqāmat-e tamām jān-e pazhmorde rā dar rah-e Maḥbūb tāze
nethār namāʾīm va ṣad-hezār ʿajz va niyāz ārīm tā be ān fowz fāʾez
shavīm. Va ān madīne kotob-e elāhiyye ast dar har ʿahdī mathalan dar
ʿahd-e Mūsā Towrāt būd va dar zaman-e ʿĪsā Enjīl va dar ʿahd-e Muḥammad
Rasūl Allāh Forqān va dar īn ʿaṣr Bayān va dar ʿahd-e Man Yabʿathuhu
Allāh ketāb-e ū ke rojūʿ-e koll-e kotob be ānast va mohaymen ast bar
jamīʿ-e kotob. Va dar īn madāyen arzāq moqaddar ast va neʿam-e bāqiye
moqarrar. Ghedhā-ye rūḥānī bakhshad va neʿmat-e qedmānī cheshānad. Bar
ahl-e tajrīd neʿmat-e towḥīd ʿaṭā farmāyad bī-naṣībān rā naṣīb karam
namāyad va āvāregān-e ṣaḥrā-ye jahl rā kaʾs-e ʿelm ʿenāyat konad va
hedāyat va ʿenāyat va ʿelm va maʿrefat va īmān va īqān-e kulli man fī
al-samāwāti wa-al-arḍ dar īn madāʾen maknūn va makhzūn gashte.

So, my beloved, we must make an effort to reach that city and, through
divine favor and the Lord's loving care, remove the veils of glory, so
that with complete steadfastness we may offer our withered souls anew in
the path of the Beloved and bring a hundred thousand expressions of
helplessness and supplication, that we may attain that triumph. That
city is the divine scriptures in every age. For example, in the age of
Moses it was the Torah; in the time of Jesus, the Gospel; in the age of
Muhammad, the Messenger of God, the Quran; in this age, the Bayan; and
in the age of Him whom God will make manifest, His Book, to which all
books refer and which holds authority over every book. In these cities
sustenance has been ordained and lasting blessings established. They
bestow spiritual nourishment and impart the taste of ancient bounty.
They grant the blessing of divine unity to those who are detached,
graciously give a share to those who have none, and offer the cup of
knowledge to those wandering in the desert of ignorance. The guidance,
grace, knowledge, recognition, faith, and certainty of all who are in
the heavens and on earth lie concealed and stored in these cities.

\subsection{221}\label{section-411}

\begin{otherlanguage}{persian}

مثلاً فرقان از برای امّت رسول حصن محکم بوده که در زمان او هر نفسی داخل او
شد از رمی شیاطین و رمح مخالفین و ظنونات مجتثّه و اشارات شرکیّه محفوظ ماند
و همچنین مرزوق شد بفواکه طیّبهٔ احدیّه و اثمار علم شجرهٔ الهیّه و از انهار
ماء غیر آسن معرفت نوشید و خمر اسرار توحید و تفرید چشید

\end{otherlanguage}

Mathalan Forqān az barā-ye ommat-e Rasūl ḥeṣn-e moḥkam būde ke dar
zamān-e ū har nafsī dākhel-e ū shod az ramy-e shayāṭīn va romḥ-e
mokhālefīn va ẓonūnāt-e mojtath the va eshārāt-e sherk iyye maḥfūẓ mānd
va hamchenīn marzūq shod be favākeh-e ṭayyebe-ye aḥadiyye va athmār-e
ʿelm-e shajare-ye elāhiyye va az anhār-e māʾ-e ghayr-e āsen-e maʿrefat
nūshīd va khamr-e asrār-e towḥīd va tafrīd cheshīd.

For example, the Quran was a strong fortress for the community of the
Messenger. In its time, every soul who entered it remained safe from the
missiles of devils, the spears of opponents, rootless conjectures, and
idolatrous insinuations. He was also nourished by the goodly fruits of
oneness and the fruits of knowledge from the divine tree, drank from the
rivers of the untainted water of recognition, and tasted the wine of the
mysteries of divine unity and singleness.

\subsection{222}\label{section-412}

\begin{otherlanguage}{persian}

چنانچه جمیع مایحتاج آن امّت در احکام دین و شریعت سیّد المرسلین در آن رضوان
مبین موجود و معیّن گشته و آنست حجّت باقیه برای اهلش بعد از نقطهٔ فرقان زیرا
مسلّم است حکم آن و محقّق الوقوعست امر آن و جمیع مأمور باتّباع آن بوده‌اند تا
ظهور بدیع در سنهٔ ستّین و آنست که طالبان را برضوان وصال میرساند و مجاهدان
و مهاجران را بسرادق قرب فائز فرماید دلیلی است محکم و حجّتی است اعظم و غیر
آن را از روایات و کتب و احادیث این فخر نه زیرا حدیث و صاحبان حدیث وجود و
قولشان بحکم کتاب ثابت و محقّق شده و دیگر آنکه در احادیث اختلاف بسیار است
و شبهه بیشمار

\end{otherlanguage}

Chenānche jamīʿ-e mā-yaḥtāj-e ān ommat dar aḥkām-e dīn va sharīʿat-e
Sayyed al-Morsalīn dar ān reżvān-e mobīn mowjūd va moʿayyan gashte va
ānast ḥojjat-e bāqiye barā-ye ahlash baʿd az Noqṭe-ye Forqān zīrā
mosallam ast ḥokm-e ān va moḥaqqaq al-voqūʿ ast amr-e ān va jamīʿ maʾmūr
be ettebāʿ-e ān būde-and tā ẓohūr-e badīʿ dar sane-ye settīn. Va ānast
ke ṭālebān rā be reżvān-e veṣāl mīresānad va mojāhedān va mohājerān rā
be sorādeq-e qorb fāʾez farmāyad. Dalīlī ast moḥkam va ḥojjatī ast aʿẓam
va gheyr-e ān rā az revāyāt va kotob va aḥādīth īn fakhr na zīrā ḥadīth
va ṣāḥebān-e ḥadīth vojūd va qowlshān be ḥokm-e ketāb thābet va moḥaqqaq
shode va dīgar ānke dar aḥādīth ekhtelāf-e besyār ast va shobhe-ye
bīshomār.

Thus, everything that community needed concerning the ordinances of
religion and the law of the Lord of the Messengers was present and
specified in that manifest paradise. It is the enduring proof for its
people after the Point of the Quran, for its ruling is acknowledged and
its decree is assured of fulfillment. All were commanded to follow it
until the wondrous Revelation in the year sixty. It brings seekers to
the paradise of reunion and admits those who strive and those who
journey forth to the tabernacle of nearness. It is a firm demonstration
and a supreme proof. Nothing else---whether traditions, books, or
reported sayings---possesses this distinction, for the existence and
words of both the reported sayings and those who transmitted them have
been established and confirmed by the authority of the Book. Moreover,
the reported sayings contain many disagreements and countless
uncertainties.

\subsection{223}\label{section-413}

\begin{otherlanguage}{persian}

چنانچه نقطهٔ فرقان در آخر امر فرمودند که انّی تارک فیکم الثّقلین کتاب اللّه
و عترتی با اینکه احادیث بسیار از منبع رسالت و معدن هدایت نازل شده بود با
وجود این جز ذکر کتاب چیزی نفرمودند و آن را سبب اعظم و دلیل اقوم برای
طالبان مقرّر فرمودند که هادی عباد باشد تا یوم معاد

\end{otherlanguage}

Chenānche Noqṭe-ye Forqān dar ākhar-e amr farmūdand ke ``innī tārikun
fīkumu al-thaqalayn kitāba Allāhi wa ʿitratī.'' Bā īnke aḥādīth-e besyār
az manbaʿ-e resālat va maʿdan-e hedāyat nāzel shode būd bā vojūd-e īn
joz dhekr-e ketāb chīzī nafarmūdand va ān rā sabab-e aʿẓam va dalīl-e
aqvam barā-ye ṭālebān moqarrar farmūdand ke hādī-ye ʿebād bāshad tā
yowm-e maʿād.

Thus, at the end of His mission, the Point of the Quran said, ``I am
leaving among you the two weighty things: the Book of God and My
family.'' Although many reported sayings had descended from the Source
of messengership and the Mine of guidance, He nevertheless mentioned
nothing but the Book, appointing it as the greatest means and the surest
guide for seekers, to guide His servants until the Day of Return.

\subsection{224}\label{section-414}

\begin{otherlanguage}{persian}

حال بچشم انصاف و قلب طاهر و نفس زکیّه ملاحظه فرمائید که در کتاب خدا که
مسلّم بین طرفین است از عامّه و خاصّه چه را حجّت برای معرفت عباد قرار فرموده
باید بنده و شما و کلّ من علی الأرض بنور آن تمسّک جسته حقّ را از باطل و
ضلالت را از هدایت تمیز دهیم و فرق گذاریم زیرا که حجّت منحصر شد بدو یکی
کتاب و دیگر عترت عترت که از میان رفته پس منحصر شد بکتاب

\end{otherlanguage}

Ḥāl be chashm-e enṣāf va qalb-e ṭāher va nafs-e zakiyye molāḥeẓe
farmāʾīd ke dar ketāb-e Khodā ke mosallam beyn-e ṭarafeyn ast az ʿāmme
va khāṣṣe che rā ḥojjat barā-ye maʿrefat-e ʿebād qarār farmūde. Bāyad
bande va shomā va kullu man ʿalā al-arḍ be nūr-e ān tamassok joste Ḥaqq
rā az bāṭel va ḍalālat rā az hedāyat tamīz dahīm va farq godhārīm zīrā
ke ḥojjat monḥaṣer shod be do yekī ketāb va dīgar ʿetrat. ʿEtrat ke az
miyān rafte pas monḥaṣer shod be ketāb.

Now consider with an impartial eye, a pure heart, and a stainless soul
what He has designated in the Book of God---which is acknowledged by
both parties, Sunni and Shia---as the proof by which His servants may
attain recognition. You, I, and everyone on earth must hold fast to its
light and distinguish truth from falsehood and error from guidance. For
the proof was confined to two things: the Book and the family. Since the
family is no longer in our midst, it is therefore confined to the Book.

\subsection{225}\label{section-415}

\begin{otherlanguage}{persian}

و اوّل کتاب میفرماید الم ذلک الکتاب لا ریب فیه هدی للمتّقین در حروف مقطّعهٔ
فرقان اسرار هویّه مستور گشته و لآلی احدیّه در صدف این حروف مخزون شده که
این مقام مجال ذکر آن نه ولیکن بر حسب ظاهر مقصود خود آن حضرتست که باو
خطاب میفرماید یا محمّد این کتاب منزل از سماء احدیّه نیست ریبی و شکّی در آن
هدایتی است برای پرهیزکاران ملاحظه فرمائید که همین فرقان را مقرّر و مقدّر
فرموده برای هدایت کلّ من فی السّموات و الأرض و بنفسه آن ذات احدیّه و غیب
هویّه شهادت داده بر آنکه شک و شبهه در آن نیست که هادی عباد است الی یوم
معاد آیا انصاف هست ثقل اعظم را که خدا شهادت بر حقّیّت آن داده و حکم بر
حقّیّت آن فرموده این عباد در آن شک نمایند و یا شبهه کنند و یا امری را که
او سبب هدایت و وصول بمعارج معرفت قرار فرموده از آن اعراض نمایند و امر
دیگر طلب نمایند و یا بحرف مزخرف ناس تشکیک نمایند که فلان چنین گفته و
فلان امر ظاهر نشده و حال آنکه اگر امری و یا احداثی غیر کتاب الهی علّت و
دلیل برای هدایت خلق بود البتّه در آیه مذکور میشد

\end{otherlanguage}

Va avval-i kitāb mīfarmāyad: Alif Lām Mīm. Dhālika al-kitābu lā rayba
fīhi hudan lil-muttaqīn. Dar ḥurūf-i muqaṭṭaʿiy-i Furqān asrār-i huviyya
mastūr gashta va laʾālī-yi aḥadiyya dar ṣadaf-i īn ḥurūf makhzūn shuda
ki īn maqām majāl-i dhikr-i ān na. Valīkin bar ḥasab-i ẓāhir maqṣūd
khud-i ān Ḥaḍrat ast ki bi-ū khiṭāb mīfarmāyad: Yā Muḥammad, īn kitāb-i
munzal az samā-yi aḥadiyya nīst raybī va shakkī dar ān, hidāyatī ast
barāy-i parhīzkārān. Mulāḥaẓa farmāyīd ki hamīn Furqān rā muqarrar va
muqaddar farmūda barāy-i hidāyat-i kull-i man fī al-samāwāti wa al-arḍ
va bi-nafsihi ān dhāt-i aḥadiyya va ghayb-i huviyya shahādat dāda bar
ānki shakk va shubha dar ān nīst ki hādī-yi ʿibād ast ilā yawm-i maʿād.
Āyā inṣāf hast thiql-i aʿẓam rā ki Khudā shahādat bar ḥaqqiyyat-i ān
dāda va ḥukm bar ḥaqqiyyat-i ān farmūda īn ʿibād dar ān shakk namāyand
va yā shubha kunand va yā amrī rā ki ū sabab-i hidāyat va vuṣūl
bi-maʿārij-i maʿrifat qarār farmūda az ān iʿrāḍ namāyand va amr-i dīgar
ṭalab namāyand va yā bi-ḥarf-i muzakhraf-i nās tashkīk namāyand ki fulān
chunīn gufta va fulān amr ẓāhir nashuda va ḥāl ānki agar amrī va yā
aḥdāthī ghayr-i kitāb-i ilāhī ʿillat va dalīl barāy-i hidāyat-i khalq
būd albatta dar āya madhkūr mīshud.

At the beginning of the Book He says, ``Alif, Lam, Mim. This is the Book
in which there is no doubt, a guidance for the God-fearing.'' Within the
disconnected letters of the Quran the mysteries of the divine Essence
lie concealed, and within the shell of these letters the pearls of
oneness are stored, but this is not the place to discuss them. In their
outward sense, however, they refer to that very Person, whom He
addresses: ``O Muhammad, this Book, sent down from the heaven of
oneness, contains no doubt or uncertainty; it is a guidance for the
God-fearing.'' Consider how He has appointed and ordained this very
Quran to guide all who are in the heavens and on earth, and how that
Essence of oneness, that hidden divine Being, has Himself testified that
there is no doubt or uncertainty in it, and that it guides His servants
until the Day of Return. Is it fair for these servants to doubt or
question the greater of the two weighty things, whose truth God has
attested and affirmed? Or to turn away from what He has made the means
of guidance and of ascent to the heights of recognition, and seek
something else? Or to entertain doubts because of people's idle words,
saying that someone has said such and such, or that a certain event has
not occurred? If anything or any events other than the divine Book were
a means and a proof for the guidance of humanity, they would certainly
have been mentioned in this verse.

\subsection{226}\label{section-416}

\begin{otherlanguage}{persian}

باری باید از امر مبرم الهی و از تقدیر مقدّر صمدانی که در آیه ذکر یافت
تجاوز ننمائیم و کتب بدیعه را مصدّق شویم چه اگر تصدیق این کتب را ننمائیم
تصدیق این آیهٔ مبارکه نشده چنانچه این واضحست که هر کس تصدیق فرقان ننمود
فی‌الحقیقه مصدّق کتب قبل از فرقان هم نبوده و این معانی از ظاهر آیه مستفاد
میشود و اگر معانی مستورهٔ آن ذکر شود و اسرار مکنونهٔ آن بیان گردد البتّه
زمان بآخر نرساند و کون حمل ننماید و کان اللّه علی ما اقول شهیداً

\end{otherlanguage}

Bārī bāyad az amr-i mubram-i ilāhī va az taqdīr-i muqaddar-i ṣamadānī ki
dar āya dhikr yāft tajāvuz nanamāyīm va kutub-i badīʿa rā muṣaddiq
shavīm, chi agar taṣdīq-i īn kutub rā nanamāyīm taṣdīq-i īn āya-yi
mubāraka nashuda. Chunānchi īn vāḍiḥ ast ki har kas taṣdīq-i Furqān
nanamūd fī al-ḥaqīqa muṣaddiq-i kutub-i qabl az Furqān ham nabūda va īn
maʿānī az ẓāhir-i āya mustafād mīshavad va agar maʿānī-yi mastūra-yi ān
dhikr shavad va asrār-i maknūna-yi ān bayān gardad albatta zamān
bi-ākhir narasānad va kawn ḥaml nanamāyad. Wa kāna Allāhu ʿalā mā aqūlu
shahīdan.

In short, we must not go beyond the irrevocable command of God and the
decree ordained by the Eternal that is mentioned in this verse, and we
must affirm the wondrous Books. For if we do not affirm these Books, we
have not affirmed this blessed verse. It is clear that whoever has not
affirmed the Quran has not, in truth, affirmed the Books that preceded
it either. These meanings can be understood from the outward sense of
the verse. Were its hidden meanings to be mentioned and its concealed
mysteries explained, time would certainly not suffice to complete the
account, nor could creation bear it. God is witness to what I say.

\subsection{227}\label{section-417}

\begin{otherlanguage}{persian}

و همچنین در جای دیگر میفرماید و ان کنتم فی ریب ممّا نزّلنا علی عبدنا فأتوا
بسورة من مثله و ادعوا شهدائکم من دون اللّه ان کنتم صادقین که ترجمهٔ ظاهر
آن اینست اگر بودید شما در شک و شبهه در آنچه ما نازل فرمودیم بر عبد خود
محمّد پس بیارید سوره‌ئی بمثل این سور منزله و بخوانید شهدای خود را یعنی
علمای خود را تا اعانت نمایند شما را در انزال سوره اگر هستید راست‌گویان
حال ملاحظه فرمائید چه مقدار عظیم است شأن آیات و بزرگست قدر آن که حجّت
بالغه و برهان کامل و قدرت قاهره و مشیّت نافذه را بآن ختم فرموده و هیچ شیء
را آن سلطان احدیّه در اظهار حجّت خود بآن شریک نفرموده چه میانهٔ حجج و دلائل
آیات بمنزلهٔ شمس است و سوای آن بمنزلهٔ نجوم و آنست حجّت باقیه و برهان ثابت
و نور مضیء از جانب سلطان حقیقی در میان عباد هیچ فضلی بآن نرسد و هیچ امری
بر آن سبقت نگیرد کنز لآلی الهیّه است و مخزن اسرار احدیّه و آنست خیط محکم و
حبل مستحکم و عروهٔ وثقی و نور لایطفی شریعهٔ معارف الهیّه از آن در جریانست و
نار حکمت بالغهٔ صمدانیّه از آن در فوران این ناریست که در یک حین دو اثر از
آن ظاهر است در مقبلین حرارت حبّ احداث نماید و در مبغضین برودت غفلت آورد

\end{otherlanguage}

Va hamchunīn dar jāy-i dīgar mīfarmāyad: Wa in kuntum fī raybin mimmā
nazzalnā ʿalā ʿabdinā faʾtū bi-sūratin min mithlihi wa udʿū shuhadāʾakum
min dūni Allāhi in kuntum ṣādiqīn. Ki tarjuma-yi ẓāhir-i ān īn ast: Agar
būdīd shumā dar shakk va shubha dar ānchi mā nāzil farmūdīm bar ʿabd-i
khud Muḥammad, pas biyārīd sūriʾī bi-mithl-i īn suvar-i munzala va
bikhvānīd shuhadā-yi khud rā yaʿnī ʿulamā-yi khud rā tā iʿānat namāyand
shumā rā dar inzāl-i sūra agar hastīd rāst-gūyān. Ḥāl mulāḥaẓa farmāyīd
chi miqdār ʿaẓīm ast shaʾn-i āyāt va buzurg ast qadr-i ān ki ḥujjat-i
bāligha va burhān-i kāmil va qudrat-i qāhira va mashiyyat-i nāfidha rā
bi-ān khatm farmūda va hīch shayʾ rā ān sulṭān-i aḥadiyya dar iẓhār-i
ḥujjat-i khud bi-ān sharīk nafarmūda. Chi miyāna-yi ḥujaj va dalāʾil
āyāt bi-manzala-yi shams ast va sivā-yi ān bi-manzala-yi nujūm va ān ast
ḥujjat-i bāqiya va burhān-i thābit va nūr-i muḍīʾ az jānib-i sulṭān-i
ḥaqīqī dar miyān-i ʿibād. Hīch faḍlī bi-ān narasad va hīch amrī bar ān
sabqat nagīrad. Kanz-i laʾālī-yi ilāhiyya ast va makhzan-i asrār-i
aḥadiyya va ān ast khayṭ-i muḥkam va ḥabl-i mustaḥkam va ʿurva-yi vuthqā
va nūr-i lā yuṭfā. Sharīʿa-yi maʿārif-i ilāhiyya az ān dar jarayān ast
va nār-i ḥikmat-i bāligha-yi ṣamadāniyya az ān dar favarān. Īn nārī ast
ki dar yak ḥīn du athar az ān ẓāhir ast: dar muqbilīn ḥarārat-i ḥubb
iḥdāth namāyad va dar mubghiḍīn burūdat-i ghaflat āvarad.

Elsewhere He also says, ``And if you are in doubt about what We have
sent down to Our servant, then produce a surah like it, and call upon
your witnesses other than God, if you are truthful.'' Its outward
meaning is this: ``If you are in doubt and uncertainty about what We
have sent down to Our servant Muhammad, then bring a surah like these
revealed surahs, and summon your witnesses---that is, your learned
men---to assist you in sending down a surah, if you are truthful.'' Now
consider how exalted is the station of the verses and how great their
worth, for He has made them the culmination of conclusive proof, perfect
demonstration, all-subduing power, and effectual will. That Sovereign of
oneness has associated nothing else with them in establishing His proof.
For among proofs and evidences, the verses are like the sun, and all
else is like the stars. They are the enduring proof, the steadfast
demonstration, and the shining light from the true Sovereign among His
servants. No bounty can equal them, and nothing can take precedence over
them. They are the treasury of divine pearls and the repository of the
mysteries of oneness; they are the strong cord, the firm rope, the
surest handle, and the light that cannot be extinguished. From them
flows the stream of divine knowledge, and from them blazes the fire of
the Eternal's consummate wisdom. This is a fire that produces two
effects at once: in those who turn toward it, it generates the warmth of
love, and in those who hate it, it brings the chill of heedlessness.

\subsection{228}\label{section-418}

\begin{otherlanguage}{persian}

ای رفیق باید از امر الهی نگذریم و بآنچه حجّت خود قرار فرموده راضی شویم و
سر بنهیم خلاصه حجّت و برهان این آیهٔ منزله اعظم از آنست که این علیل تواند
اقامهٔ دلیل نماید و اللّه یقول الحقّ و هو یهدی السّبیل و هو القاهر فوق عباده
و هو العزیز الجمیل

\end{otherlanguage}

Ay rafīq, bāyad az amr-i ilāhī nagudharīm va bi-ānchi ḥujjat-i khud
qarār farmūda rāḍī shavīm va sar binahīm. Khulāṣa ḥujjat va burhān-i īn
āya-yi munzala aʿẓam az ān ast ki īn ʿalīl tavānad iqāma-yi dalīl
namāyad. Wa Allāhu yaqūlu al-ḥaqqa wa huwa yahdī al-sabīla wa huwa
al-qāhiru fawqa ʿibādihi wa huwa al-ʿAzīzu al-Jamīl.

O friend, we must not go beyond the command of God, but must be content
with what He has appointed as His proof and bow before it. In short, the
proof and demonstration furnished by this revealed verse are too great
for this ailing one to establish by argument. God speaks the truth and
guides the way. He is supreme over His servants, and He is the Almighty,
the All-Beautiful.

\subsection{229}\label{section-419}

\begin{otherlanguage}{persian}

و همچنین میفرماید تلک آیات‌ اللّه نتلوها علیک بالحقّ فبأیّ حدیث بعد اللّه و
آیاته یؤمنون میفرماید اینست آیات منزلهٔ از سماء هویّه میخوانیم بر شما پس
بکدام سخن بعد از ظهور حقّ و نزول آیات او ایمان میآورند اگر تلویح این آیه
را ملتفت شوی میفهمی که هرگز مظهری اکبر از انبیا نبوده و حجّتی هم اکبر و
اعظم از آیات منزله در ارض ظاهر نشده بلکه اعظم از این حجّت حجّتی ممکن نه
الّا ما شآء ربّک

\end{otherlanguage}

Va hamchunīn mīfarmāyad: Tilka āyātu Allāhi natlūhā ʿalayka bi-al-ḥaqqi
fa-bi-ayyi ḥadīthin baʿda Allāhi wa āyātihi yuʾminūn. Mīfarmāyad īn ast
āyāt-i munzala az samā-yi huviyya mīkhvānīm bar shumā pas bi-kudām
sukhan baʿd az ẓuhūr-i Ḥaqq va nuzūl-i āyāt-i ū īmān mīāvarand. Agar
talvīḥ-i īn āya rā multafit shavī mīfahmī ki hargiz maẓharī akbar az
anbiyā nabūda va ḥujjatī ham akbar va aʿẓam az āyāt-i munzala dar arḍ
ẓāhir nashuda balka aʿẓam az īn ḥujjat ḥujjatī mumkin na illā mā shāʾa
rabbuka.

He also says, ``These are the verses of God that We recite to you in
truth. In what discourse, then, after God and His verses, will they
believe?'' He says: ``These are the verses sent down from the heaven of
the divine Essence, which We recite to you. In what words, then, will
they believe after the manifestation of the True One and the revelation
of His verses?'' If you grasp the allusion in this verse, you will
understand that there has never been any manifestation greater than the
Prophets, nor has any proof greater or mightier than the revealed verses
appeared on earth. Indeed, no proof greater than this is possible,
except as your Lord wills.

\subsection{230}\label{section-420}

\begin{otherlanguage}{persian}

و در جای دیگر میفرماید ویل لکلّ افّاک اثیم یسمع آیات اللّه تتلی علیه ثمّ یصرّ
مستکبراً کأن لم یسمعها فبشّره بعذاب الیم یعنی وای بر افک‌کنندهٔ گنه‌کار که
میشنود آیات نازلهٔ از سماء مشیّت الهیّه را که خوانده میشود بر او پس استکبار
مینماید که گویا نشنیده آن را پس بشارت ده او را بعذابی دردناک اشارات همین
آیه کلّ من فی السّموات و الأرض را کفایت میکند لو کان النّاس فی آیات ربّهم
یتفرّسون چنانچه الیوم میشنوید که اگر از آیات الهیّه قرائت شود احدی اعتنا
ندارد که گویا پست‌ترین امرها نزدشان آیات الهیّه است و حال آنکه اعظم از
آیات امری نبوده و نخواهد بود بگو بایشان ای بی‌خبران میگوئید آنچه را قبل
پدران شما گفتند اگر ایشان ثمری از شجرهٔ اعراض خود دیدند شما هم خواهید دید
و عنقریب با آبای خود در نار مقرّ خواهید یافت فالنّار مثواهم فبئس مثوی
الظّالمین

\end{otherlanguage}

Va dar jāy-i dīgar mīfarmāyad: Waylun li-kulli affākin athīmin yasmaʿu
āyāti Allāhi tutlā ʿalayhi thumma yuṣirru mustakbiran kaʾan lam yasmaʿhā
fa-bashshirhu bi-ʿadhābin alīm. Yaʿnī vāy bar ifk-kunanda-yi gunah-kār
ki mīshinavad āyāt-i nāzila az samā-yi mashiyyat-i ilāhiyya rā ki
khvānda mīshavad bar ū pas istikbār mīnamāyad ki gūyā nashanīda ān rā
pas bishārat dih ū rā bi-ʿadhābī dardnāk. Ishārāt-i hamīn āya kull-i man
fī al-samāwāti wa al-arḍ rā kifāyat mīkunad law kāna al-nāsu fī āyāti
rabbihim yatafarrasūn. Chunānchi al-yawm mīshinavīd ki agar az āyāt-i
ilāhiyya qirāʾat shavad aḥadī iʿtinā nadārad ki gūyā past-tarīn-i amrhā
nazdishān āyāt-i ilāhiyya ast va ḥāl ānki aʿẓam az āyāt amrī nabūda va
nakhvāhad būd. Bigū bi-īshān ay bī-khabarān mīgūyīd ānchi rā qabl
pidarān-i shumā guftand. Agar īshān thamarī az shajara-yi iʿrāḍ-i khud
dīdand shumā ham khvāhīd dīd va ʿan-qarīb bā ābā-yi khud dar nār maqarr
khvāhīd yāft. Fa-al-nāru mathwāhum fa-biʾsa mathwā al-ẓālimīn.

Elsewhere He says, ``Woe to every sinful liar who hears the verses of
God recited to him, then persists in arrogance as though he had not
heard them. So give him tidings of a painful punishment.'' That is:
``Woe to the sinful fabricator who hears the verses sent down from the
heaven of the divine will recited to him, then acts arrogantly as though
he had not heard them. So give him tidings of a painful punishment.''
The allusions in this verse alone suffice for all who are in the heavens
and on earth, if only people would discern the verses of their Lord. As
you hear today, when the divine verses are recited, no one pays
attention, as though the verses of God were the most insignificant of
all things in their eyes, whereas nothing greater than the verses has
ever existed or ever will. Say to them: ``O heedless ones, you say what
your fathers said before you. If they saw any fruit from the tree of
their rejection, you too will see it, and before long you will find your
abode in the fire with your fathers.'' The fire is their dwelling place;
how wretched is the dwelling of the wrongdoers!

\subsection{231}\label{section-421}

\begin{otherlanguage}{persian}

و در جای دیگر میفرماید و اذا علم من آیاتنا شیئاً اتّخذها هزؤاً اولئک لهم
عذاب مهین یعنی در وقتی که عالم شود از آیات ما شیئی را اخذ مینماید آن را
از روی استهزاء از برای ایشانست عذابی خوارکننده از جملهٔ استهزاء این بود
که میگفتند معجزهٔ دیگر ظاهر نما و برهان دیگر بیاور آن یکی فأسقط علینا
کسفاً من السّمآء میگفت و دیگر ان کان هذا هو الحقّ من عندک فأمطر علینا حجارةً
من السّمآء مذکور میداشت بمثل یهودان عهد موسی که تبدیل نمودند مائدهٔ سمائیّه
را باشیای خبیثهٔ سیر و پیاز آن قوم هم طلب تبدیل مینمودند آیات منزله را
بظنونات نجسهٔ کثیفه چنانچه الیوم مشاهده میشود که مائدهٔ معنویّه از سماء
رحمت الهیّه و غمام مکرمت سبحانیّه نازلست و بحور حیوان در رضوان جنان بامر
خالق کن فکان در موج و جریان و جمیع چون کلاب بر اجساد میّته مجتمع آمده‌اند
و ببرکهٔ شور که ملح اجاجست قانع گشته‌اند سبحان‌اللّه کمال تحیّر حاصلست از
عبادی که بعد از ارتفاع اعلام مدلول طلب دلیل مینمایند و بعد از ظهور شمس
معلوم باشارات علم تمسّک جسته‌اند مثل آنست که از آفتاب در اثبات نور او حجّت
طلبند و یا از باران نیسان در اثبات فیضش برهان جویند حجّت آفتاب نور اوست
که اشراق نموده و عالم را فراگرفته و برهان نیسان جود اوست که عالم را
بردای جدید تازه فرموده بلی کور از آفتاب جز گرمی حاصلی نداند و ارض جرز از
رحمت نیسانی فضلی احصا ننماید عجب نبود که از قرآن نصیبی نیست جز نقشی که
از خورشید جز گرمی نیابد چشم نابینا

\end{otherlanguage}

Va dar jāy-i dīgar mīfarmāyad: Wa idhā ʿalima min āyātinā shayʾan
ittakhadhahā huzuwan ulāʾika lahum ʿadhābun muhīn. Yaʿnī dar vaqtī ki
ʿālim shavad az āyāt-i mā shayʾī rā akhdh mīnamāyad ān rā az rūy-i
istihzā, az barāy-i īshān ast ʿadhābī khvār-kunanda. Az jumla-yi istihzā
īn būd ki mīguftand muʿjiza-yi dīgar ẓāhir namā va burhān-i dīgar
biyāvar. Ān yakī fa-asqiṭ ʿalaynā kisafan mina al-samāʾ mīguft va dīgar
in kāna hādhā huwa al-ḥaqqa min ʿindika fa-amṭir ʿalaynā ḥijāratan mina
al-samāʾ madhkūr mīdāsht. Bi-mithl-i Yahūdān-i ʿahd-i Mūsā ki tabdīl
namūdand māʾida-yi samāʾiyya rā bi-ashyā-yi khabītha-yi sīr va piyāz, ān
qawm ham ṭalab-i tabdīl mīnamūdand āyāt-i munzala rā bi-ẓunūnāt-i
najisa-yi kathīfa. Chunānchi al-yawm mushāhada mīshavad ki māʾida-yi
maʿnaviyya az samā-yi raḥmat-i ilāhiyya va ghamām-i makramat-i
subḥāniyya nāzil ast va buḥūr-i ḥayavān dar Riḍvān-i jinān bi-amr-i
khāliq-i kun fa-kān dar mawj va jarayān va jamīʿ chun kilāb bar ajsād-i
mayyita mujtamiʿ āmada-and va bi-birka-yi shūr ki milḥ-i ujāj ast qāniʿ
gashta-and. Subḥān Allāh, kamāl-i taḥayyur ḥāṣil ast az ʿibādī ki baʿd
az irtifāʿ-i aʿlām-i madlūl ṭalab-i dalīl mīnamāyand va baʿd az ẓuhūr-i
shams-i maʿlūm bi-ishārāt-i ʿilm tamassuk justa-and. Mithl-i ān ast ki
az āftāb dar ithbāt-i nūr-i ū ḥujjat ṭalaband va yā az bārān-i naysān
dar ithbāt-i fayḍash burhān jūyand. Ḥujjat-i āftāb nūr-i ūst ki ishrāq
namūda va ʿālam rā farā girifta va burhān-i naysān jūd-i ūst ki ʿālam rā
bi-ridā-yi jadīd tāza farmūda. Balī kūr az āftāb juz garmī ḥāṣilī
nadānad va arḍ-i juruz az raḥmat-i naysānī faḍlī iḥṣā nanamāyad. ʿAjab
nabvad ki az Qurʾān naṣībī nīst juz naqshī ki az khvurshīd juz garmī
nayābad chashm-i nābīnā.

Elsewhere He says, ``When he learns something of Our verses, he takes
them in mockery. For such people there is a humiliating punishment.''
That is, when he learns something of Our verses, he treats it with
mockery; for them there is a humiliating punishment. Among their
mockeries was their saying, ``Show another miracle and bring another
proof.'' One would say, ``Then cause pieces of the sky to fall upon
us,'' while another declared, ``If this is indeed the truth from You,
then rain down stones upon us from heaven.'' Like the Jews in the time
of Moses, who exchanged the heavenly food for the vile things of garlic
and onions, these people too sought to exchange the revealed verses for
filthy, defiled imaginings. Thus today it can be seen that spiritual
food is descending from the heaven of divine mercy and the clouds of
transcendent bounty, and the seas of life are surging and flowing in the
Ridvan of paradise at the command of the Creator of ``Be, and it was,''
yet all have gathered like dogs around dead bodies and contented
themselves with a brackish pool of bitter salt water. Glory be to God!
How utterly astonishing are those servants who seek evidence after the
standards of the One to whom all evidence points have been raised, and
who cling to the intimations of knowledge after the sun of the Object of
knowledge has appeared. It is as though they were to demand from the sun
a proof of its light, or seek from the April rain a demonstration of its
bounty. The sun's proof is its light, which has shone forth and
encompassed the world, and the proof of April is its generosity, which
has renewed the world with a fresh robe. Yet the blind know no benefit
from the sun but its warmth, and barren ground can reckon no bounty from
the mercy of the April rain. ``No wonder if your portion of the Quran is
nothing but its written form, for the sightless eye gains nothing from
the sun but warmth.''

\subsection{232}\label{section-422}

\begin{otherlanguage}{persian}

و در جای دیگر میفرماید و اذا تتلی علیهم آیاتنا بیّنات ما کان حجّتهم الّا ان
قالوا ائتوا بآبائنا ان کنتم صادقین یعنی در وقتی که تلاوت کرده میشود بر
ایشان آیات ما نیست حجّت ایشان مگر آنکه بگویند بیاورید پدرهای ما را اگر
هستید راست‌گویان مشاهده نما که چه حجّتها بر آن رحمتهای کاملهٔ واسعه
میگرفتند بآیاتی که حرفی از آن اعظم است از خلق سموات و ارض و مردگان وادی
نفس و هوی را بروح ایمان زنده میفرماید استهزا مینمودند و میگفتند پدرهای
ما را از قبر بیرون آر این بود اعراض و استکبار قوم و هر کدام از این آیات
برای کلّ من علی الأرض حجّتی است محکم و برهانیست معظم که جمیع ارض را کفایت
میکند لو انتم فی آیات اللّه تتفکّرون و در همین آیهٔ مذکوره لآلی اسرار
مکنونست اگر فی‌الجمله دردی یافت شود دوا میرسد

\end{otherlanguage}

Va dar jāy-i dīgar mīfarmāyad: Wa idhā tutlā ʿalayhim āyātunā bayyinātin
mā kāna ḥujjatahum illā an qālū iʾtū bi-ābāʾinā in kuntum ṣādiqīn. Yaʿnī
dar vaqtī ki tilāvat karda mīshavad bar īshān āyāt-i mā nīst ḥujjat-i
īshān magar ānki bigūyand biyāvarīd pidarhā-yi mā rā agar hastīd
rāst-gūyān. Mushāhada namā ki chi ḥujjathā bar ān raḥmathā-yi kāmila-yi
vāsiʿa mīgiriftand. Bi-āyātī ki ḥarfī az ān aʿẓam ast az khalq-i samāvāt
va arḍ va murdigān-i vādī-yi nafs va havā rā bi-rūḥ-i īmān zinda
mīfarmāyad istihzā mīnamūdand va mīguftand pidarhā-yi mā rā az qabr
bīrūn ār. Īn būd iʿrāḍ va istikbār-i qawm va har kudām az īn āyāt
barāy-i kull-i man ʿalā al-arḍ ḥujjatī ast muḥkam va burhānī ast
muʿaẓẓam ki jamīʿ-i arḍ rā kifāyat mīkunad law antum fī āyāti Allāhi
tatafakkarūn. Va dar hamīn āya-yi madhkūra laʾālī-yi asrār maknūn ast
agar fī al-jumla dardī yāft shavad davā mīrasad.

Elsewhere He says, ``When Our clear verses are recited to them, their
only argument is to say, `Bring back our fathers, if you are
truthful.'\,'' That is, when Our verses are recited to them, their only
argument is to say, ``Bring back our fathers, if you are truthful.''
Observe what objections they raised against those perfect, all-embracing
mercies. They mocked verses a single letter of which is greater than the
creation of the heavens and the earth, and which bring those dead in the
valley of self and desire to life through the spirit of faith, and they
said, ``Bring our fathers out of their graves.'' Such was the people's
rejection and arrogance. Each of these verses is a firm proof and a
mighty demonstration for all who are on earth, sufficient for the whole
earth, if only you would reflect on the verses of God. Within this very
verse lie hidden pearls of mysteries. If even a little pain is felt, the
remedy will come.

\subsection{233}\label{section-423}

\begin{otherlanguage}{persian}

گوش بحرفهای مزخرف عباد مدهید که میگویند کتاب و آیات از برای عوام حجّت
نمیشود زیرا که نمی‌فهمند و احصا نمیکنند با اینکه این قرآن حجّت است بر مشرق
و مغرب عالم اگر قوّهٔ ادراک آن در مردم نبود چگونه حجّت بر کل واقع میشد از
این قرار بر معرفت الهیّه هم نفسی مکلّف نیست و لازم نه زیرا که عرفان او
اعظم از عرفان کتاب اوست و عوام استعداد ادراک آن ندارند

\end{otherlanguage}

Gūsh bi-ḥarfhā-yi muzakhraf-i ʿibād madihīd ki mīgūyand kitāb va āyāt az
barāy-i ʿavāmm ḥujjat namīshavad zīrā ki namīfahmand va iḥṣā namīkunand
bā īnki īn Qurʾān ḥujjat ast bar mashriq va maghrib-i ʿālam. Agar
quvva-yi idrāk-i ān dar mardum nabūd chigūna ḥujjat bar kull vāqiʿ
mīshud. Az īn qarār bar maʿrifat-i ilāhiyya ham nafsī mukallaf nīst va
lāzim na zīrā ki ʿirfān-i ū aʿẓam az ʿirfān-i kitāb-i ūst va ʿavāmm
istiʿdād-i idrāk-i ān nadārand.

Do not listen to the idle words of those who say that the Book and the
verses cannot serve as proof for ordinary people because they neither
understand nor comprehend them, even though this Quran is a proof for
the East and the West of the world. If people lacked the capacity to
understand it, how could it have become a proof for all? By this
reasoning, no one is obliged to attain knowledge of God, nor is it
necessary, since recognizing Him is greater than recognizing His Book,
and ordinary people lack the capacity to comprehend it.

\subsection{234}\label{section-424}

\begin{otherlanguage}{persian}

باری این سخن بغایت لغو و غیر مقبولست همه از روی کبر و غرور گفته میشود که
مردم را از ریاض رضای الهی دور کنند و زمام آنها را محکم حفظ نمایند با
اینکه نزد حقّ این عوام از علمای ایشان که از حقّ اعراض نموده‌اند بغایت
مقبولتر و پسندیده‌ترند و فهم کلمات الهیّه و درک بیانات حمامات معنویّه هیچ
دخلی بعلم ظاهری ندارد این منوط بصفای قلب و تزکیهٔ نفوس و فراغت روح است
چنانچه حال عبادی چند موجودند که حرفی از رسوم علم ندیده‌اند و بر رفرف علم
جالسند و از سحاب فیض الهی ریاض قلوبشان بگلهای حکمت و لاله‌های معرفت تزیین
یافته فطوبی للمخلصین من انوار یوم عظیم

\end{otherlanguage}

Bārī īn sokhan be-ghāyat laghv va gheyr-e maqbūl ast. Hameh az rū-ye
kebr va ghorūr gofteh mī-shavad keh mardom rā az riyāḍ-e reḍā-ye elāhī
dūr konand va zemām-e ānhā rā moḥkam ḥefẓ namāyand, bā īnkeh nazd-e Ḥaqq
īn ʿavāmm az ʿolamā-ye īshān keh az Ḥaqq eʿrāḍ namūdeh-and be-ghāyat
maqbūltar va pasandīdehtarand. Va fahm-e kalemāt-e elāhīyeh va dark-e
bayānāt-e ḥamāmāt-e maʿnavīyeh hīch dakhlī be-ʿelm-e ẓāherī nadārad. Īn
manūṭ be-ṣafā-ye qalb va tazkīyeh-ye nofūs va farāghat-e rūḥ ast.
Chenāncheh ḥāl ʿebādī chand mawjūdand keh ḥarfī az rosūm-e ʿelm
nadīdeh-and va bar rafraf-e ʿelm jālesand va az saḥāb-e feyḍ-e elāhī
riyāḍ-e qolūbeshān be-golhā-ye ḥekmat va lāleh-hā-ye maʿrefat tazyīn
yāfteh. Fa-ṭūbā li-l-mukhliṣīna min anwāri yawmin ʿaẓīm.

In short, this claim is utterly idle and unacceptable. It is all said
out of pride and arrogance, to keep people far from the gardens of
divine good pleasure and to hold their reins firmly, even though, in the
sight of the True One, these ordinary people are far more acceptable and
pleasing than those of their religious scholars who have turned away
from Him. Understanding the words of God and comprehending the
utterances of the spiritual doves have nothing to do with outward
learning. This depends on purity of heart, purification of the soul, and
freedom of spirit. Even now there are servants who have never learned a
single letter of conventional learning, yet are seated upon the pavilion
of knowledge, and the gardens of their hearts have been adorned, through
the clouds of divine grace, with the flowers of wisdom and the tulips of
understanding. Blessed, then, are the sincere through the lights of a
mighty Day.

\subsection{235}\label{section-425}

\begin{otherlanguage}{persian}

و همچنین میفرماید و الّذین کفروا بآیات اللّه و لقائه اولئک یئسوا من رحمتی
و اولئک لهم عذاب الیم و همچنین میفرماید و یقولون ائنّا لتارکوا آلهتنا
لشاعر مجنون مضمون این آیه واضحست مشاهده فرمائید که بعد از تنزیل آیات چه
میگفتند آیا ما ترک‌کننده‌ایم خدایان خود را برای شاعری مجنون که آن حضرت را
شاعر مینامیدند و بآیات الهیّه سخریّه مینمودند و میگفتند این کلمات اساطیر
اوّلین است یعنی کلماتیست که قبل گفته شده و محمّد آن کلمات را ترکیب نموده
میگوید از خدا است

\end{otherlanguage}

Va hamchenīn mī-farmāyad: Wa-alladhīna kafarū bi-āyāti Allāhi
wa-liqāʾihi ulāʾika yaʾisū min raḥmatī wa-ulāʾika lahum ʿadhābun alīm.
Va hamchenīn mī-farmāyad: Wa-yaqūlūna a-innā la-tārikū ālihatinā
li-shāʿirin majnūn. Maḍmūn-e īn āyeh vāḍeḥ ast. Moshāhedeh farmāʾīd keh
baʿd az tanzīl-e āyāt cheh mī-goftand: āyā mā tark-konandeh-īm
khodāyān-e khod rā barā-ye shāʿerī majnūn? Keh ān Ḥaḍrat rā shāʿer
mī-nāmīdand va be-āyāt-e elāhīyeh sokhrīyeh mī-namūdand va mī-goftand īn
kalemāt asāṭīr-e avvalīn ast, yaʿnī kalemātīst keh qabl gofteh shodeh va
Moḥammad ān kalemāt rā tarkīb namūdeh mī-gūyad az Khodā ast.

He also says, ``Those who disbelieve in the verses of God and in meeting
Him have despaired of My mercy, and for them there is a painful
punishment.'' He also says, ``And they say, `Are we really to abandon
our gods for a mad poet?'\,'' The meaning of this verse is clear.
Observe what they said after the verses were sent down: ``Are we to
abandon our gods for a mad poet?'' They called Him a poet, mocked the
divine verses, and said, ``These words are tales of the
ancients''---meaning words that had been spoken before, which Muhammad
had pieced together and claimed were from God.

\subsection{236}\label{section-426}

\begin{otherlanguage}{arabic}

چنانچه الیوم بمثل همان را شنیده‌اید که نسبت باین امر میدهند و میگویند که
این کلمات را با کلمات قبل ترکیب نموده و یا کلماتیست مغلوط قد کبر قولهم و
صغر شأنهم و حدّهم

\end{otherlanguage}

Chenāncheh al-yawm be-mesl-e hamān rā shenīdeh-īd keh nesbat be-īn Amr
mī-dahand va mī-gūyand keh īn kalemāt rā bā kalemāt-e qabl tarkīb
namūdeh va yā kalemātīst maghlūṭ. Qad kabura qawluhum wa-ṣaghura
shaʾnuhum wa-ḥadduhum.

Likewise, today you have heard them make the same allegations against
this Cause, saying that these words have been pieced together from
earlier words, or that they are full of errors. Great indeed is their
claim, and small their stature and capacity!

\subsection{237}\label{section-427}

\begin{otherlanguage}{persian}

اینست که بعد از این انکارها و اعتراضات که مذکور شد گفتند بعد از موسی و
عیسی موافق کتب نباید نبیّ مستقلّ که ناسخ شریعت باشد مبعوث شود باید شخصی
بیاید که مکمّل شریعت قبل باشد این آیهٔ مبارکه که مشعر بر جمیع مطالب الهیّه
و مدلّ بر عدم انقطاع فیوضات رحمانیّه است نازل شد و لقد جآءکم یوسف من قبل
بالبیّنات فما زلتم فی شکّ ممّا جآءکم به حتّی اذا هلک قلتم لن یبعث اللّه من
بعده رسولاً کذلک یضلّ اللّه من هو مسرف مرتاب و بتحقیق آمد شما را یوسف از
پیش با بیّنه‌ها پس پیوسته بودید در شک از آنچه آمد شما را بآن تا چون هلاک
شد گفتید مبعوث نمیگرداند خدا بعد از او رسولی را همچنین اضلال میکند خدا
کسی را که اوست اسراف‌کننده و شک‌آورنده بپروردگار خود پس از این آیه ادراک
فرمائید و یقین کنید که در هر عصر امم آن عهد بآیه‌ئی از کتاب تمسّک جسته از
اینگونه حرفهای مزخرف میگفتند که دیگر نبی نباید در ابداع بیاید مثل آنکه
آیهٔ انجیل را که مذکور شده علمای آن استدلال بآن نمودند که هرگز حکم انجیل
مرتفع نمیشود و پیغمبری مستقلّ مبعوث نگردد الّا برای اثبات شریعت انجیل و
اکثری از ملل مبتلا باین مرض روحی شده‌اند

\end{otherlanguage}

Īn ast keh baʿd az īn enkārhā va eʿterāḍāt keh madhkūr shod goftand baʿd
az Mūsā va ʿĪsā movāfeq-e kotob nabāyad nabī-ye mostaqell keh nāsekh-e
sharīʿat bāshad mabʿūth shavad; bāyad shakhṣī biyāyad keh mokammel-e
sharīʿat-e qabl bāshad. Īn āyeh-ye mobārakeh keh moshʿer bar jamīʿ-e
maṭāleb-e elāhīyeh va modall bar ʿadam-e enqeṭāʿ-e foyūḍāt-e raḥmānīyeh
ast nāzel shod: Wa-laqad jāʾakum Yūsufu min qablu bi-l-bayyināti fa-mā
ziltum fī shakkin mimmā jāʾakum bihi ḥattā idhā halaka qultum lan
yabʿatha Allāhu min baʿdihi rasūlan kadhālika yuḍillu Allāhu man huwa
musrifun murtāb. Va be-taḥqīq āmad shomā rā Yūsof az pīsh bā
bayyeneh-hā, pas peyvasteh būdīd dar shakk az āncheh āmad shomā rā
be-ān, tā chon halāk shod goftīd mabʿūth nemī-gardānad Khodā baʿd az ū
rasūlī rā. Hamchenīn eḍlāl mī-konad Khodā kasī rā keh ūst esrāf-konandeh
va shakk-āvarandeh be-Parvardegār-e khod. Pas az īn āyeh edrāk farmāʾīd
va yaqīn konīd keh dar har ʿaṣr omam-e ān ʿahd be-āyeh-ī az ketāb
tamassok josteh az īngūneh ḥarfhā-ye mozakhraf mī-goftand keh dīgar nabī
nabāyad dar ebdāʿ biyāyad. Mesl-e ānkeh āyeh-ye Enjīl rā keh madhkūr
shodeh ʿolamā-ye ān estedlāl be-ān namūdand keh hargez ḥokm-e Enjīl
mortafeʿ nemī-shavad va peyghambarī mostaqell mabʿūth nagardad ellā
barā-ye ethbāt-e sharīʿat-e Enjīl. Va aktharī az melal mobtalā be-īn
maraḍ-e rūḥī shodeh-and.

Thus, after the denials and objections already mentioned, they said
that, according to the scriptures, no independent Prophet who would
abrogate the religious law should be raised up after Moses and Jesus;
rather, someone must come who would complete the previous religious law.
This blessed verse was then revealed, encompassing all divine purposes
and indicating that the outpourings of the All-Merciful never cease:
``Joseph came to you before with clear proofs, yet you remained in doubt
about what he brought you, until, when he died, you said, `God will
never raise up a messenger after him.' Thus God leads astray whoever is
transgressive and doubting.'' Indeed, Joseph came to you before with
clear proofs, yet you remained in doubt about what he brought you,
until, when he died, you said, ``God will not raise up a messenger after
him.'' Thus God leads astray whoever is transgressive and doubtful of
his Lord. Understand from this verse, then, and be certain that in every
age the peoples of that time have clung to a verse of the Book and
spoken such nonsense, saying that no further Prophet should come into
creation. Thus the scholars of the Gospel used the verse already
mentioned to argue that the law of the Gospel would never be annulled,
and that no independent Prophet would be raised up except to confirm the
religious law of the Gospel. Most religious communities have become
afflicted with this spiritual sickness.

\subsection{238}\label{section-428}

\begin{otherlanguage}{persian}

چنانچه اهل فرقان را می‌بینی که چگونه بمثل امم قبل بذکر خاتم‌ النّبیّین محتجب
گشته‌اند با اینکه خود مقرّند بر اینکه ما یعلم تأویله الّا اللّه و الرّاسخون
فی العلم بعد که راسخ در علوم و امّها و نفسها و ذاتها و جوهرها بیان
میفرماید که قدری مخالف هوای ایشان واقع میشود اینست که میشنوی که چه
میگویند و چه میکنند و نیست اینها مگر از رؤسای ناس در دین یعنی آنهائی که
الهی بجز هوی اخذ نکرده‌اند و بغیر ذهب مذهبی نیافته‌اند و بحجبات علم محتجب
گشته‌اند و بضلالت آن گمراه شده‌اند چنانچه بتصریح تمام ربّ الأنام میفرماید أ
فرأیت من اتّخذ الهه هواه و اضلّه اللّه علی علم و ختم علی سمعه و قلبه و جعل
علی بصره غشاوةً فمن یهدیه من بعد اللّه أ فلا تذکّرون یعنی آیا دیدی آن غافل
را که گرفت خدای خود خواهشهای نفس خود را و اضلال کرد او را خدا بر علمی و
مهر نهاد بر گوش و دلش و گردانید بر چشمش پرده‌ئی پس که هدایت میکند او را
از بعد خدا آیا پند نمیگیرید

\end{otherlanguage}

Chenāncheh ahl-e Forqān rā mī-bīnī keh chegūneh be-mesl-e omam-e qabl
be-dhekr-e Khātam al-Nabiyyīn moḥtajeb gashteh-and, bā īnkeh khod
moqerrand bar īnkeh: Mā yaʿlamu taʾwīlahu illā Allāhu wa-l-rāsikhūna fī
al-ʿilm. Baʿd keh rāsekh dar ʿolūm va ummihā va nafsihā va dhātihā va
jawharihā bayān mī-farmāyad keh qadrī mokhālef-e havā-ye īshān vāqeʿ
mī-shavad, īn ast keh mī-shenavī keh cheh mī-gūyand va cheh mī-konand.
Va nīst īnhā magar az roʾasā-ye nās dar dīn, yaʿnī ānhāʾī keh elāhī
be-joz havā akhdh nakardeh-and va be-gheyr-e dhahab madhhabī
nayāfteh-and va be-ḥojobāt-e ʿelm moḥtajeb gashteh-and va be-ḍalālat-e
ān gomrāh shodeh-and. Chenāncheh be-taṣrīḥ-e tamām Rabb al-Anām
mī-farmāyad: A-fa-raʾayta mani ittakhadha ilāhahu hawāhu wa-aḍallahu
Allāhu ʿalā ʿilmin wa-khatama ʿalā samʿihi wa-qalbihi wa-jaʿala ʿalā
baṣarihi ghishāwatan fa-man yahdīhi min baʿdi Allāhi a-fa-lā
tadhakkarūn. Yaʿnī āyā dīdī ān ghāfel rā keh gereft khodā-ye khod
khwāheshhā-ye nafs-e khod rā va eḍlāl kard ū rā Khodā bar ʿelmī va mohr
nahād bar gūsh va delash va gardānīd bar chashmash pardeh-ī? Pas keh
hedāyat mī-konad ū rā az baʿd-e Khodā? Āyā pand nemī-gīrīd?

Thus you see how the people of the Quran, like the peoples before them,
have become veiled by the expression ``Seal of the Prophets,'' although
they themselves acknowledge that ``None knows its interpretation except
God and those firmly grounded in knowledge.'' Yet when He who is firmly
grounded in the sciences, who is their source, their very self, their
essence and their substance, makes a statement that runs even slightly
counter to their desires, you hear what they say and what they do. These
things come only from the people's leaders in religion---that is, those
who have taken no god but desire and have found no creed but gold, who
have become veiled by the veils of learning and have been led astray by
its error. As the Lord of all people states with complete clarity:
``Have you seen the one who has taken his desire as his god, whom God
has led astray despite his knowledge, whose hearing and heart He has
sealed, and over whose sight He has placed a covering? Who, then, will
guide him after God? Will you not take heed?'' That is, have you seen
that heedless one who has made the desires of his own self his god, whom
God has led astray despite his knowledge, whose ears and heart He has
sealed, and over whose eyes He has placed a veil? Who, then, will guide
him after God? Will you not take heed?

\subsection{239}\label{section-429}

\begin{otherlanguage}{persian}

در معنی و اضلّه اللّه علی علم اگرچه در ظاهر آنست که ذکر شد ولیکن نزد این
فانی مقصود از آیه علمای عصرند که اعراض از جمال حقّ نمودند و بعلوم خود که
از نفس و هوی ناشی گشته متمسّک شده بر نبأ الهی و امر او احتجاج مینمودند قل
هو نبأ عظیم انتم عنه معرضون و همچنین میفرماید و اذا تتلی علیهم آیاتنا
بیّنات قالوا ما هذا الّا رجل یرید ان یصدّکم عمّا کان یعبد آبائکم و قالوا ما
هذا الّا افک مفتری و الحقّ یقول و چون خوانده شود بر ایشان یعنی بر آن کفرهٔ
فجره آیات قدسیّهٔ احدیّه گویند آن مشرکان از حقّ بیخبران نیست این رسول
پروردگار مگر مردی که میخواهد منع کند شما را از آنچه که میپرستیدند آن را
پدرهای شما و دیگر گفتند نیست این مگر کذبی افترا کرده شده

\end{otherlanguage}

Dar maʿnī-ye wa-aḍallahu Allāhu ʿalā ʿilm, agarcheh dar ẓāher ān ast keh
dhekr shod, valīkan nazd-e īn fānī maqṣūd az āyeh ʿolamā-ye ʿaṣrand keh
eʿrāḍ az Jamāl-e Ḥaqq namūdand va be-ʿolūm-e khod keh az nafs va havā
nāshī gashteh motamassek shodeh bar nabaʾ-e elāhī va Amr-e ū eḥtejāj
mī-namūdand. Qul huwa nabaʾun ʿaẓīmun antum ʿanhu muʿriḍūn. Va hamchenīn
mī-farmāyad: Wa-idhā tutlā ʿalayhim āyātunā bayyinātin qālū mā hādhā
illā rajulun yurīdu an yaṣuddakum ʿammā kāna yaʿbudu ābāʾukum wa-qālū mā
hādhā illā ifkun muftaran. Wa-l-Ḥaqqu yaqūl. Va chon khwāndeh shavad bar
īshān, yaʿnī bar ān kafareh-ye fajareh, āyāt-e qodsīyeh-ye aḥadīyeh,
gūyand ān moshrekān-e az Ḥaqq bī-khabarān: nīst īn Rasūl-e Parvardegār
magar mardī keh mī-khwāhad manʿ konad shomā rā az āncheh keh
mī-parastīdand ān rā pedarhā-ye shomā. Va dīgar goftand nīst īn magar
kedhbī efterā kardeh shodeh.

As for the meaning of ``whom God has led astray despite his knowledge,''
although its outward meaning is what has been stated, in the view of
this transient one the verse refers to the religious scholars of the age
who turned away from the Beauty of the True One and, clinging to their
own learning, which arose from self and desire, disputed the divine
Announcement and His Cause. ``Say: It is a mighty announcement, from
which you turn away.'' He also says, ``When Our clear verses are recited
to them, they say, `This is only a man who wishes to turn you away from
what your fathers worshipped.' And they say, `This is nothing but a
fabricated lie.'\,'' The True One speaks. And when the holy verses of
the divine Unity are recited to them---that is, to those wicked
unbelievers---those idolaters, unaware of the True One, say, ``This
Messenger of the Lord is only a man who wishes to prevent you from
worshipping what your fathers worshipped.'' They also said, ``This is
nothing but a fabricated lie.''

\subsection{240}\label{section-430}

\begin{otherlanguage}{persian}

بشنوید ندای قدس الهی و نوای خوش صمدانی را که چگونه در تلویح انذار فرموده
مکذّبین آیات را و بیزاری جسته منکرین کلمات قدسیّه را و بعد ناس را ملاحظه
فرمائید از کوثر قرب و اعراض و استکبار آن محرومان را بر آن جمال قدس با
اینکه آن جوهر لطف و کرم هیاکل عدم را بعرصهٔ قدم هدایت میفرمود و آن فقیران
حقیقی را بشریعهٔ قدسیّهٔ غنا دلالت مینمود مع‌ذلک بعضی میگفتند این مردیست
افترا‌کننده بر پروردگار عالمیان و بعضی میگفتند این منع‌کننده است ناس را از
شریعهٔ دین و ایمان و برخی نسبت جنون میدادند و امثال ذلک

\end{otherlanguage}

Beshnavīd nedā-ye qods-e elāhī va navā-ye khosh-e ṣamadānī rā keh
chegūneh dar talvīḥ endhār farmūdeh mokadhdhebīn-e āyāt rā va bīzārī
josteh monkerīn-e kalemāt-e qodsīyeh rā. Va baʿd nās rā molāḥeẓeh
farmāʾīd az Kawthar-e qorb va eʿrāḍ va estekbār-e ān maḥrūmān rā bar ān
Jamāl-e Qods, bā īnkeh ān Jawhar-e Loṭf va Karam hayākel-e ʿadam rā
be-ʿarṣeh-ye qedam hedāyat mī-farmūd va ān faqīrān-e ḥaqīqī rā
be-sharīʿeh-ye qodsīyeh-ye ghenā delālat mī-namūd. Maʿa-dhālek baʿḍī
mī-goftand īn mardīst efterā-konandeh bar Parvardegār-e ʿālamiyān va
baʿḍī mī-goftand īn manʿ-konandeh ast nās rā az sharīʿeh-ye dīn va īmān
va barkhī nesbat-e jonūn mī-dādand va amthāl-e dhālek.

Listen to the holy call of God and the sweet melody of the Eternal, how
He has warned, by allusion, those who deny the verses, and has disavowed
those who reject the holy words. Then consider the people in relation to
the Kawthar of nearness, and the rejection and arrogance of those
deprived ones toward that Holy Beauty, even though that Essence of Grace
and Generosity was guiding those embodiments of nothingness to the realm
of eternity and directing those truly impoverished ones to the holy
waters of riches. Nevertheless, some said, ``This is a man who
fabricates lies against the Lord of the worlds,'' while others said,
``He is preventing people from reaching the waters of religion and
faith,'' and still others accused Him of madness and the like.

\subsection{241}\label{section-431}

\begin{otherlanguage}{persian}

چنانچه الیوم مشاهده میکنید چه سخنهای لغو که بآن جوهر بقا گفته‌اند و چه
نسبتها و خطاها که بآن منبع و معدن عصمت داده‌اند با اینکه در کتاب الهی و
لوح قدس صمدانی در جمیع اوراق و کلمات انذار فرموده مکذّبین و معرضین آیات
منزله را و بشارت فرموده مقبلین آن را با وجود این چه ‌قدر اعتراضات که بر
آیات منزلهٔ از سموات قدسیّهٔ بدعیّه نموده‌اند و حال آنکه چشم امکان چنین فضلی
ندیده و قوّهٔ سمع اکوان چنین عنایتی نشنیده که آیات بمثابهٔ غیث نیسانی از
غمام رحمت رحمانی جاری و نازل شود چه که انبیای اولو العزم که عظمت قدر و
رفعت مقامشان چون شمس واضح و لائحست مفتخر شدند هر کدام بکتابی که در دست
هست و مشاهده شده و آیات آن احصا گشته و از این غمام رحمت رحمانی اینقدر
نازل شده که هنوز احدی احصا ننموده چنانچه بیست مجلّد الآن بدست میآید و چه
‌مقدار که هنوز بدست نیامده و چه ‌مقدار هم که تاراج شده و بدست مشرکین
افتاده و معلوم نیست چه کرده‌اند

\end{otherlanguage}

Chenāncheh al-yawm moshāhedeh mī-konīd cheh sokhanhā-ye laghv keh be-ān
Jawhar-e Baqā gofteh-and va cheh nesbathā va khaṭāhā keh be-ān Manbaʿ va
Maʿdan-e ʿEṣmat dādeh-and, bā īnkeh dar Ketāb-e Elāhī va Lawḥ-e Qods-e
Ṣamadānī dar jamīʿ-e awrāq va kalemāt endhār farmūdeh mokadhdhebīn va
moʿreḍīn-e āyāt-e monazaleh rā va beshārat farmūdeh moqbelīn-e ān rā. Bā
vojūd-e īn cheh qadr eʿterāḍāt keh bar āyāt-e monazaleh az samāvāt-e
qodsīyeh-ye badīʿīyeh namūdeh-and, va ḥāl ānkeh chashm-e emkān chenīn
faḍlī nadīdeh va qovveh-ye samʿ-e akvān chenīn ʿenāyatī nashenīdeh keh
āyāt be-mathābeh-ye ghey th-e neysānī az ghamām-e raḥmat-e raḥmānī jārī
va nāzel shavad. Cheh keh anbiyā-ye ulu'l-ʿazm keh ʿaẓamat-e qadr va
refʿat-e maqāmeshān chon shams vāḍeḥ va lāʾeḥ ast moftakher shodand har
kodām be-ketābī keh dar dast hast va moshāhedeh shodeh va āyāt-e ān eḥṣā
gashteh. Va az īn ghamām-e raḥmat-e raḥmānī īnqadr nāzel shodeh keh
hanūz aḥadī eḥṣā nanamūdeh, chenāncheh bīst mojallad al-ān be-dast
mī-āyad va cheh meqdār keh hanūz be-dast nayāmadeh va cheh meqdār ham
keh tārāj shodeh va be-dast-e moshrekīn oftādeh va maʿlūm nīst cheh
kardeh-and.

Thus today you see what idle words they have spoken against that Essence
of Eternity, and what accusations and faults they have attributed to
that Source and Mine of Infallibility, even though in the divine Book
and the holy Tablet of the Eternal, throughout all their pages and
words, He has warned those who deny and turn away from the revealed
verses, and has given glad tidings to those who accept them. Despite
this, how many objections they have raised against the verses sent down
from the wondrous heavens of holiness! Yet the eye of contingent
existence has never seen such grace, nor has the hearing of all beings
ever heard of such bounty, that verses should stream and descend like
spring rain from the clouds of the mercy of the All-Merciful. For the
Prophets endowed with constancy, the greatness of whose worth and
loftiness of whose station are as clear and manifest as the sun, were
each honoured with a Book that is available, has been seen, and whose
verses have been counted. But so much has descended from this cloud of
the mercy of the All-Merciful that no one has yet counted it. Twenty
volumes are now available, while so much has not yet come to hand, and
so much has been plundered and fallen into the hands of the idolaters,
and it is not known what they have done with it.

\subsection{242}\label{section-432}

\begin{otherlanguage}{persian}

ای برادر باید چشم گشود و تفکّر نمود و ملتجی بمظاهر الهیّه شد که شاید از
مواعظ واضحهٔ کتاب پند گیریم و از نصایح مذکورهٔ در الواح متنبّه شویم اعتراض
بر منزل آیات نکنیم امرش را بجان تسلیم کنیم و حکمش را بتمام جان و روان
قبول نمائیم و مذعن شویم که شاید در فضای رحمت وارد شویم و در شاطی فضل
مسکن یابیم و انّه بعباده لغفور رحیم

\end{otherlanguage}

Ey barādar, bāyad chashm goshūd va tafakkor namūd va moltajī
be-maẓāher-e elāhīyeh shod keh shāyad az mavāʿeẓ-e vāḍeḥeh-ye Ketāb pand
gīrīm va az naṣāyeḥ-e madhkūreh dar alvāḥ motanabbeh shavīm, eʿterāḍ bar
Monzel-e Āyāt nakonīm, Amrash rā be-jān taslīm konīm va ḥokmash rā
be-tamām-e jān va ravān qabūl namāʾīm va modhʿen shavīm keh shāyad dar
faḍā-ye raḥmat vāred shavīm va dar shāṭī-ye faḍl maskan yābīm. Wa-innahu
bi-ʿibādihi la-ghafūrun raḥīm.

O brother, we must open our eyes, reflect, and seek refuge in the
Manifestations of God, so that perhaps we may take heed of the clear
admonitions of the Book and be awakened by the counsels set forth in the
Tablets. Let us not object to the Revealer of the verses, but submit
wholeheartedly to His Cause and accept and acknowledge His decree with
all our heart and soul, so that perhaps we may enter the expanse of
mercy and find a home upon the shore of grace. He is indeed forgiving
and merciful toward His servants.

\subsection{243}\label{section-433}

\begin{otherlanguage}{persian}

و همچنین میفرماید قل یا اهل الکتاب هل تنقمون منّا الّا ان آمنّا باللّه و ما
انزل الینا و ما انزل من قبل و انّ اکثرکم فاسقون چه‌ قدر واضح است مقصود در
این آیه و چه مبرهن است حجّیّت آیات منزله و این آیه در وقتی نازل شد که کفّار
به اسلام اذیّت مینمودند و نسبت کفر میدادند چنانچه نسبت میدادند باصحاب آن
حضرت که بخدا کافر شده‌اید و بساحری کذّاب مؤمن و موقن گشته‌اید و در صدر
اسلام که هنوز امر بر حسب ظاهر قوّت نداشت در هر مقام و مکان که دوستان آن
حضرت را ملاقات مینمودند نهایت اذیّت و زجر و رجم و سب بر آن مقبلین الی
اللّه معمول میداشتند در این وقت این آیهٔ مبارکه از سماء احدیّه نازل شد
ببرهانی واضح و دلیلی لائح و تعلیم فرمود اصحاب آن حضرت را که بگوئید
بکافران و مشرکان که آیا اذیّت میکنید ما را و ستم مینمائید و عملی از ما
صادر نشد مگر آنکه ایمان آوردیم بخدا و بآیاتی که نازل شد بر ما از لسان
محمّد و همچنین آیاتی که نازل شد بر انبیای او از قبل که مقصود این است
تقصیری نداریم مگر آنکه آیات جدیدهٔ بدیعهٔ الهیّه را که بر محمّد نازل شد و
آیات قدیمه که بر انبیای قبل نازل شد جمیع را من عند اللّه دانستیم و تصدیق
و اذعان نمودیم و این دلیلی است که سلطان احدیّه تعلیم فرموده عباد خود را

\end{otherlanguage}

Va hamchenīn mī-farmāyad: Qul yā ahla al-kitābi hal tanqimūna minnā illā
an āmannā bi-Allāhi wa-mā unzila ilaynā wa-mā unzila min qablu wa-anna
aktharakum fāsiqūn. Cheh qadr vāḍeḥ ast maqṣūd dar īn āyeh va cheh
mobarhan ast ḥojjīyat-e āyāt-e monazaleh. Va īn āyeh dar vaqtī nāzel
shod keh koffār be-Eslām adhīyat mī-namūdand va nesbat-e kofr mī-dādand,
chenāncheh nesbat mī-dādand be-aṣḥāb-e ān Ḥaḍrat keh be-Khodā kāfer
shodeh-īd va be-sāḥerī kadhdhāb moʾmen va mūqen gashteh-īd. Va dar
ṣadr-e Eslām keh hanūz Amr bar ḥasab-e ẓāher qovvat nadāsht, dar har
maqām va makān keh dūstān-e ān Ḥaḍrat rā molāqāt mī-namūdand, nehāyat-e
adhīyat va zajr va rajm va sabb bar ān moqbelīn ilā Allāh maʿmūl
mī-dāshtand. Dar īn vaqt īn āyeh-ye mobārakeh az samāʾ-e aḥadīyeh nāzel
shod be-borhānī vāḍeḥ va dalīlī lāʾeḥ va taʿlīm farmūd aṣḥāb-e ān Ḥaḍrat
rā keh be-gūʾīd be-kāferān va moshrekān keh āyā adhīyat mī-konīd mā rā
va setam mī-namāʾīd va ʿamalī az mā ṣāder nashod magar ānkeh īmān
āvardīm be-Khodā va be-āyātī keh nāzel shod bar mā az lesān-e Moḥammad
va hamchenīn āyātī keh nāzel shod bar anbiyā-ye ū az qabl, keh maqṣūd īn
ast taqṣīrī nadārīm magar ānkeh āyāt-e jadīdeh-ye badīʿeh-ye elāhīyeh rā
keh bar Moḥammad nāzel shod va āyāt-e qadīmeh keh bar anbiyā-ye qabl
nāzel shod jamīʿ rā min ʿindi Allāh dānestīm va taṣdīq va edhʿān
namūdīm. Va īn dalīlī ast keh Solṭān-e Aḥadīyeh taʿlīm farmūdeh ʿebād-e
khod rā.

He also says, ``Say: O people of the Book, do you resent us for anything
except that we believe in God, in what has been sent down to us, and in
what was sent down before, and that most of you are transgressors?'' How
clear is the intent of this verse, and how evident the authority of the
revealed verses as proof! This verse was revealed at a time when the
unbelievers were persecuting Islam and accusing its followers of
unbelief. Thus they accused His companions, saying, ``You have
disbelieved in God and have come to believe with certainty in a lying
sorcerer.'' In the early days of Islam, when the Cause had not yet
gained outward strength, wherever they encountered His friends, they
subjected those who had turned toward God to the utmost persecution,
torment, stoning, and abuse. At that time this blessed verse descended
from the heaven of divine Unity with a clear proof and a manifest
argument, instructing His companions to say to the unbelievers and
idolaters, ``Do you persecute and wrong us when we have done nothing but
believe in God and in the verses sent down to us through the tongue of
Muhammad, as well as the verses sent down to His Prophets before Him?''
The meaning is this: we are guilty of nothing except that we have
recognized as coming from God all the new and wondrous divine verses
revealed to Muhammad, together with the earlier verses revealed to the
former Prophets, and have affirmed and acknowledged them. This is a
proof that the Sovereign of divine Unity has taught His servants.

\subsection{244}\label{section-434}

\begin{otherlanguage}{persian}

مع‌ذلک آیا جایز است این آیات بدیعه که احاطه فرمود شرق و غرب را از آن معرض
شوند و خود را از اهل ایمان دانند و یا آنکه مؤمن شوند منزل آیات را باین
استدلال که خود فرموده مقرّین را از اهل ایمان محسوب نفرماید حاشا ثمّ حاشا
که مقبلین و مقرّین آیات احدیّه را از ابواب رحمت خود براند و متمسّکین بحجّت
مثبته را تهدید فرماید اذ انّه مثبت الحقّ بآیاته و محقّق الأمر بکلماته و انّه
لهو المقتدر المهیمن القدیر

\end{otherlanguage}

Maʿa-dhālik āyā jāʾiz ast īn āyāt-i badīʿih kih iḥāṭih farmūd sharq va
gharb rā az ān muʿriḍ shavand va khud rā az ahl-i īmān dānand va yā
ānkih muʾmin shavand munzil-i āyāt rā bi-īn istidlāl kih khud farmūdih
muqirrīn rā az ahl-i īmān maḥsūb nafarmāyad ḥāshā thumma ḥāshā kih
muqbilīn va muqirrīn-i āyāt-i aḥadīyyih rā az abvāb-i raḥmat-i khud
birānad va mutamassikīn bi-ḥujjat-i muthbitih rā tahdīd farmāyad idh
innahu muthbitu al-ḥaqqi bi-āyātihi wa-muḥaqqiqu al-amri bi-kalimātihi
wa-innahu la-huwa al-Muqtadiru al-Muhayminu al-Qadīr.

Nevertheless, is it permissible for them to turn away from these
wondrous verses, which have encompassed East and West, and count
themselves among the people of faith? Or, if they believe in the
Revealer of the verses on the strength of this argument that He Himself
has set forth, would He not count those who acknowledge them among the
people of faith? Far from it, far indeed, that He should drive away from
the gates of His mercy those who turn toward and acknowledge the verses
of divine Unity, or threaten those who hold fast to the conclusive
proof! For He establishes the truth through His verses and confirms the
Cause through His words. He is indeed the All-Powerful, the
All-Protecting, the Almighty.

\subsection{245}\label{section-435}

\begin{otherlanguage}{persian}

و همچنین میفرماید و لو نزّلنا علیک کتاباً فی قرطاس فلمسوه بأیدیهم لقال
الّذین کفروا ان هذا الّا سحر مبین و اکثری آیات فرقانیّه مدلّ و مشعر بر این
مطلب است و این بنده اختصار نمودم باین آیات مذکوره و حال ملاحظه فرمائید
که در جمیع کتاب جز آیات را که حجّت قرار فرموده برای معرفت مظاهر جمال خود
دیگر امری ذکر شده تا بآن متمسّک شوند و اعتراض نمایند بلکه در همهٔ موارد بر
منکرین آیات و استهزاءکنندهٔ آن وعدهٔ نار فرموده‌اند چنانچه معلوم شد

\end{otherlanguage}

Va hamchunīn mīfarmāyad wa-law nazzalnā ʿalayka kitāban fī qirṭāsin
fa-lamasūhu bi-aydīhim la-qāla alladhīna kafarū in hādhā illā siḥrun
mubīn va aktharī āyāt-i furqānīyyih mudill va mushʿir bar īn maṭlab ast
va īn bandih ikhtiṣār nimūdam bi-īn āyāt-i madhkūrih va ḥāl mulāḥiẓih
farmāʾīd kih dar jamīʿ-i kitāb juz āyāt rā kih ḥujjat qarār farmūdih
barāy-i maʿrifat-i maẓāhir-i jamāl-i khud dīgar amrī dhikr shudih tā
bi-ān mutamassik shavand va iʿtirāḍ namāyand balkih dar hamiy-i mavārid
bar munkirīn-i āyāt va istihzāʾ-kunandiy-i ān vaʿdiy-i nār farmūdih-and
chunānchih maʿlūm shud.

He also says, ``Even if We had sent down to you a book on parchment, and
they had touched it with their hands, those who disbelieve would have
said, `This is nothing but manifest sorcery.'\,'' Most of the verses of
the Qur'an indicate and convey this point, but I, this servant, have
confined myself to the verses already cited. Now consider: throughout
the entire Book, has anything other than the verses, which He has
established as proof for recognizing the Manifestations of His beauty,
been mentioned for them to seize upon and raise objections? Rather, in
every instance He has promised the Fire to those who deny the verses and
mock them, as has been made clear.

\subsection{246}\label{section-436}

\begin{otherlanguage}{persian}

حال اگر کسی بیاید بکرورها از آیات و خطب و صحائف و مناجات بی‌آنکه بتعلیم
اخذ نموده باشد آیا بچه دلیل میتوان اعتراض نمود و از این فیض اکبر محروم
شد و جواب چه خواهند گفت بعد از عروج روح از جسد ظلمانی آیا متمسّک میشوند
که بفلان حدیث تمسّک جستیم و چون معنی آن را بظاهر نیافتیم لهذا بر مظاهر
امر اعتراض نمودیم و از شرایع حقّ دور گشتیم آیا نشنیده‌اید که از جملهٔ علّت
اینکه بعضی از انبیا اولو العزم بودند نزول کتاب بود بر آنها و این مسلّم
است با وجود این چگونه جایز است که بر صاحب کتب که چندین مجلّدات از او ظاهر
شده بحرفهای فلان مرد که از روی جهل بعضی کلمات برای القای شبهه در قلوب
جمع نموده و شیطان عصر شده برای اغفال عباد و اضلال من فی البلاد پیروی
نمایند و از خورشید فیض الهی بی‌بهره گردند و از همۀ این مراتب گذشته آیا از
این نفس قدسی و نفس رحمانی احتراز جویند و ادبار نمایند نمیدانم بکه تمسّک
جویند و بکدام وجه اقبال کنند بلی و لکلّ وجهة هو مولّیها فقد هدیناک
السّبیلین فی هذین المنهجین ثمّ امش علی ما تختار لنفسک و هذا قول الحقّ و ما
بعد الحقّ الّا الضّلال

\end{otherlanguage}

Ḥāl agar kasī biyāyad bi-kurūrhā az āyāt va khuṭab va ṣaḥāʾif va munājāt
bī-ānkih bi-taʿlīm akhdh nimūdih bāshad āyā bi-chih dalīl mītavān
iʿtirāḍ nimūd va az īn fayḍ-i akbar maḥrūm shud va javāb chih khvāhand
guft baʿd az ʿurūj-i rūḥ az jasad-i ẓulmānī āyā mutamassik mīshavand kih
bi-fulān ḥadīth tamassuk justīm va chūn maʿnī-yi ān rā bi-ẓāhir nayāftīm
lihādhā bar maẓāhir-i amr iʿtirāḍ nimūdīm va az sharāyiʿ-i Ḥaqq dūr
gashtīm āyā nashunīdih-īd kih az jumliy-i ʿillat-i īnkih baʿḍī az anbiyā
ūlu al-ʿazm būdand nuzūl-i kitāb būd bar ānhā va īn musallam ast bā
vujūd-i īn chigūnih jāʾiz ast kih bar ṣāḥib-i kutub kih chandīn
mujalladāt az ū ẓāhir shudih bi-ḥarfhāy-i fulān mard kih az rūy-i jahl
baʿḍī kalimāt barāy-i ilqāy-i shubhih dar qulūb jamʿ nimūdih va
shayṭān-i ʿaṣr shudih barāy-i ighfāl-i ʿibād va iḍlāl-i man fī al-bilād
payravī namāyand va az khurshīd-i fayḍ-i ilāhī bī-bahrih gardand va az
hamiy-i īn marātib gudhashtih āyā az īn nafs-i qudsī va nafs-i raḥmānī
iḥtirāz jūyand va idbār namāyand namīdānam bi-kih tamassuk jūyand va
bi-kudām vajh iqbāl kunand balī wa-li-kullin wijhatun huwa muwallīhā
fa-qad hadaynāka al-sabīlayni fī hādhayni al-manhajayni thumma imshi
ʿalā mā takhtāru li-nafsika wa-hādhā qawlu al-ḥaqqi wa-mā baʿda al-ḥaqqi
illā al-ḍalāl.

Now, if someone should come with tens of millions of verses, sermons,
scriptures, and prayers, without having acquired them through
instruction, on what grounds could anyone object and deprive himself of
this supreme outpouring of grace? What answer will they give after the
soul has ascended from the darkened body? Will they plead, ``We clung to
such-and-such a tradition, and because we did not find its meaning
outwardly fulfilled, we therefore objected to the Manifestations of the
Cause and kept away from the laws of God''? Have you not heard that one
reason certain Prophets were endowed with constancy was that a Book was
sent down to them? This is an acknowledged fact. How, then, can it be
permissible, in opposing the Author of Books from whom so many volumes
have appeared, to follow the words of some man who, out of ignorance,
has gathered together a few statements to sow doubt in people's hearts
and has become the satan of the age, deceiving the servants and leading
astray those who dwell in the land, and thus to remain deprived of the
Sun of divine grace? Beyond all this, if they shun and turn away from
this holy Being, this merciful Soul, I do not know to whom they will
cling or toward whose face they will turn. Yes, ``Everyone has a
direction toward which he turns.'' We have shown you the two ways in
these two paths; now follow whichever you choose for yourself. This is
the word of truth, and beyond truth there is nothing but error.

\subsection{247}\label{section-437}

\begin{otherlanguage}{persian}

و از جملهٔ ادلّه بر اثبات این امر آنکه در هر عهد و عصر که غیب هویّه در هیکل
بشریّه ظاهر میشد بعضی از مردمی که معروف نبودند و علاقه بدنیا و جهتی
نداشته‌اند بضیاء شمس نبوّت مستضیء و بانوار قمر هدایت مهتدی میشدند و بلقآء
اللّه فائز میگشتند لهذا این بود که علمای عصر و اغنیای عهد استهزاء
مینمودند چنانچه از لسان آن گمراهان میفرماید فقال الملأ الّذین کفروا من
قومه ما نراک الّا بشراً مثلنا و ما نراک اتّبعک الّا الّذینهم اراذلنا بادی
الرّأی و ما نری لکم علینا من فضل بل نظنّکم کاذبین اعتراض مینمودند و بآن
مظاهر قدسیّه میگفتند که متابعت شما نکرده مگر اراذل ما که اعتنائی بشأن
آنها نیست و مقصودشان این بوده که علما و اغنیا و معارف قوم بشما ایمان
نیاوردند و باین دلیل و امثال آن استدلال بر بطلان من له الحقّ مینمودند

\end{otherlanguage}

Va az jumliy-i adillih bar ithbāt-i īn amr ānkih dar har ʿahd va ʿaṣr
kih ghayb-i huvīyyih dar haykal-i basharīyyih ẓāhir mīshud baʿḍī az
mardumī kih maʿrūf nabūdand va ʿalāqih bi-dunyā va jihatī nadāshtih-and
bi-ḍiyāy-i shams-i nubuwwat mustaḍīʾ va bi-anvār-i qamar-i hidāyat
muhtadī mīshudand va bi-liqāʾ-i Allāh fāʾiz mīgashtand lihādhā īn būd
kih ʿulamāy-i ʿaṣr va aghniyāy-i ʿahd istihzāʾ mīnimūdand chunānchih az
lisān-i ān gumrāhān mīfarmāyad fa-qāla al-malaʾu alladhīna kafarū min
qawmihi mā narāka illā basharan mithlanā wa-mā narāka ittabaʿaka illā
alladhīna hum arādhilunā bādiya al-raʾyi wa-mā narā lakum ʿalaynā min
faḍlin bal naẓunnukum kādhibīn iʿtirāḍ mīnimūdand va bi-ān maẓāhir-i
qudsīyyih mīguftand kih mutābaʿat-i shumā nakardih magar arādhil-i mā
kih iʿtināʾī bi-shaʾn-i ānhā nīst va maqṣūd-shān īn būdih kih ʿulamā va
aghniyā va maʿārif-i qawm bi-shumā īmān nayāvardand va bi-īn dalīl va
amthāl-i ān istidlāl bar buṭlān-i man lahu al-ḥaqq mīnimūdand.

Among the proofs establishing this Cause is that, in every age and era
when the unseen divine Essence appeared in a human temple, certain
people who were unknown and had no attachment to the world or to any
worldly standing were illumined by the radiance of the Sun of
prophethood, guided by the light of the Moon of guidance, and blessed
with attaining the presence of God. For this reason, the learned and
wealthy of the age mocked them. Thus He speaks through the lips of those
misguided ones: ``The leaders among His people who disbelieved said, `We
see you as nothing but a human being like ourselves, and we see that
none have followed you but the lowliest among us, rash in judgment. Nor
do we see that you have any superiority over us; rather, we think you
are liars.'\,'' They objected and said to those holy Manifestations,
``None have followed you but the lowliest among us, who are of no
account.'' They meant that the learned, the wealthy, and the notables of
the people had not believed in Them. By this argument and others like
it, they sought to prove the falsity of Him to whom the truth belongs.

\subsection{248}\label{section-438}

\begin{otherlanguage}{persian}

و امّا در این ظهور اظهر و سلطنت عظمی جمعی از علمای راشدین و فضلای کاملین
و فقهای بالغین از کأس قرب و وصال مرزوق شدند و بعنایت عظمی فائز گشتند و
از کون و امکان در سبیل جانان گذشتند بعضی از اسامی آنها ذکر میشود که شاید
سبب استقامت انفس مضطربه و نفوس غیر مطمئنّه شود

\end{otherlanguage}

Va ammā dar īn ẓuhūr-i aẓhar va salṭanat-i ʿuẓmā jamʿī az ʿulamāy-i
rāshidīn va fuḍalāy-i kāmilīn va fuqahāy-i bālighīn az kaʾs-i qurb va
viṣāl marzūq shudand va bi-ʿināyat-i ʿuẓmā fāʾiz gashtand va az kawn va
imkān dar sabīl-i jānān gudhashtand baʿḍī az asāmī-yi ānhā dhikr
mīshavad kih shāyad sabab-i istiqāmat-i anfus-i muḍṭaribih va nufūs-i
ghayr-i muṭmaʾinnih shavad.

In this most manifest Revelation and this supreme sovereignty, however,
a number of rightly guided scholars, accomplished men of learning, and
mature jurists drank from the cup of nearness and reunion, attained the
supreme favor, and renounced all existence and the contingent world in
the path of the Beloved. Some of their names are mentioned here, that
this may perhaps bring steadfastness to troubled hearts and unsettled
souls.

\subsection{249}\label{section-439}

\begin{otherlanguage}{persian}

از آن جمله جناب ملّا حسین است که محلّ اشراق شمس ظهور شدند لولاه ما استوی
اللّه علی عرش رحمانیّته و ما استقرّ علی کرسیّ صمدانیّته و جناب آقا سیّد یحیی
که وحید عصر و فرید زمان خود بودند و ملّا محمّد علی زنجانی و ملّا علی بسطامی
و ملّا سعید بارفروشی و ملّا نعمة اللّه مازندرانی و ملّا یوسف اردبیلی و ملّا
مهدی خوئی و آقا سیّد حسین ترشیزی و ملّا مهدی کندی و برادر او ملّا باقر و
ملّا عبدالخالق یزدی و ملّا علی برقانی و امثال آنها که قریب چهارصد نفر
بودند که اسامی جمیع در لوح محفوظ الهی ثبت شده

\end{otherlanguage}

Az ān jumlih Jināb-i Mullā Ḥusayn ast kih maḥall-i ishrāq-i shams-i
ẓuhūr shudand lawlāhu mā istawā Allāhu ʿalā ʿarshi raḥmāniyyatihi wa-mā
istaqarra ʿalā kursiyyi ṣamadāniyyatihi va Jināb-i Āqā Siyyid Yaḥyā kih
vaḥīd-i ʿaṣr va farīd-i zamān-i khud būdand va Mullā Muḥammad ʿAlī
Zanjānī va Mullā ʿAlī Basṭāmī va Mullā Saʿīd Bārfurūshī va Mullā
Niʿmatuʾllāh Māzandarānī va Mullā Yūsuf Ardabīlī va Mullā Mihdī Khūʾī va
Āqā Siyyid Ḥusayn Turshīzī va Mullā Mihdī Kandī va barādar-i ū Mullā
Bāqir va Mullā ʿAbduʾl-Khāliq Yazdī va Mullā ʿAlī Baraghānī va amthāl-i
ānhā kih qarīb-i chahārṣad nafar būdand kih asāmī-yi jamīʿ dar lawḥ-i
maḥfūẓ-i ilāhī thabt shudih.

Among them is the honored Mullā Ḥusayn, who became the dawning-place of
the Sun of Revelation. Without him, God would not have ascended the
throne of His all-embracing mercy or been established upon the seat of
His eternal sovereignty. There were also the honored Āqā Siyyid Yaḥyā,
who was unique in his age and peerless in his time; Mullā Muḥammad ʿAlī
Zanjānī; Mullā ʿAlī Basṭāmī; Mullā Saʿīd Bārfurūshī; Mullā Niʿmatu'llāh
Māzandarānī; Mullā Yūsuf Ardabīlī; Mullā Mihdī Khūʾī; Āqā Siyyid Ḥusayn
Turshīzī; Mullā Mihdī Kandī and his brother Mullā Bāqir; Mullā
ʿAbdu'l-Khāliq Yazdī; Mullā ʿAlī Baraghānī; and others like them,
numbering nearly four hundred, all of whose names are inscribed in God's
Preserved Tablet.

\subsection{250}\label{section-440}

\begin{otherlanguage}{persian}

همهٔ اینها مهتدی و مقرّ و مذعن گشتند برای آن شمس ظهور بقسمی که اکثری از
مال و عیال گذشتند و برضای ذی الجلال پیوستند و از سر جان برای جانان
برخاستند و انفاق نمودند بجمیع آنچه مرزوق گشته بودند بقسمی که سینه‌هاشان
محلّ تیرهای مخالفین گشت و سرهاشان زینت سنان مشرکین چنانچه ارضی نماند مگر
آنکه از دم این ارواح مجرّده آشامید و سیفی نماند مگر آنکه بگردنهاشان ممسوح
گشت و دلیل بر صدق قولشان فعلشان بس آیا شهادت این نفوس قدسیّه که باین طریق
جان در راه دوست دادند که همۀ عالم از ایثار دل و جانشان متحیّر گشتند کفایت
نمیکند برای این عبادی که هستند و انکار بعضی عباد که دین را بدرهمی دادند
و بقا را بفنا تبدیل نمودند و کوثر قرب را بچشمه‌های شور معاوضه کردند و بجز
اخذ اموال ناس مرادی نجویند چنانچه مشاهده میشود که کل بزخارف دنیا مشغول
شده‌اند و از ربّ اعلی دور مانده

\end{otherlanguage}

Hamiy-i īnhā muhtadī va muqirr va mudhʿin gashtand barāy-i ān shams-i
ẓuhūr bi-qismī kih aktharī az māl va ʿiyāl gudhashtand va bi-riḍāy-i Dhī
al-Jalāl payvastand va az sar-i jān barāy-i jānān barkhāstand va infāq
nimūdand bi-jamīʿ-i ānchih marzūq gashtih būdand bi-qismī kih
sīnih-hā-shān maḥall-i tīrhāy-i mukhālifīn gasht va sarhā-shān zīnat-i
sinān-i mushrikīn chunānchih arḍī namānd magar ānkih az dam-i īn arvāḥ-i
mujarradih āshāmīd va sayfī namānd magar ānkih bi-gardanhā-shān mamsūḥ
gasht va dalīl bar ṣidq-i qawl-shān fiʿl-shān bas āyā shahādat-i īn
nufūs-i qudsīyyih kih bi-īn ṭarīq jān dar rāh-i dūst dādand kih hamiy-i
ʿālam az īthār-i dil va jān-shān mutaḥayyir gashtand kifāyat namīkunad
barāy-i īn ʿibādī kih hastand va inkār-i baʿḍī ʿibād kih dīn rā
bi-dirhamī dādand va baqā rā bi-fanā tabdīl nimūdand va kawthar-i qurb
rā bi-chashmih-hāy-i shūr muʿāvaḍih kardand va bi-juz akhdh-i amvāl-i
nās murādī najūyand chunānchih mushāhidih mīshavad kih kull
bi-zakhārif-i dunyā mashghūl shudih-and va az Rabb-i Aʿlā dūr māndih.

All these were guided to that Sun of Revelation, acknowledged Him, and
submitted to Him, to such a degree that most renounced possessions and
family and embraced the good pleasure of the Lord of Glory. They gave up
their lives for the Beloved and spent all that had been bestowed upon
them, until their breasts became targets for the arrows of their
opponents and their heads adorned the spears of the idolaters. No land
remained that had not drunk the blood of these detached spirits, and no
sword remained that had not passed across their necks. Their deeds alone
suffice to prove the truth of their words. Is the testimony of these
holy souls, who gave their lives in the path of the Friend in such a
manner that the whole world marveled at their sacrifice of heart and
soul, not sufficient for the people of today? And what of the denial of
certain people who sold their religion for a coin, exchanged eternity
for extinction, and bartered the Kawthar of nearness for brackish
springs, seeking no aim but to seize people's property? As can be seen,
they are all engrossed in the trappings of the world and remain far from
the Most Exalted Lord.

\subsection{251}\label{section-441}

\begin{otherlanguage}{persian}

حال انصاف دهید که شهادت اینها مقبول و مسموعست که قولشان و فعلشان موافق و
ظاهرشان و باطنشان مطابق بنحوی که تاهت العقول فی افعالهم و تحیّرت النّفوس
فی اصطبارهم و بما حملت اجسادهم و یا شهادت این معرضین که بجز هوای نفس
نفسی برنیارند و از قفس ظنونات باطله نجاتی نیافته‌اند و در یوم سر از فراش
برندارند مگر چون خفّاش ظلمانی در طلب دنیای فانیه کوشند و در لیل راحت
نشوند مگر در تدبیرات امورات دانیه کوشند بتدبیر نفسانی مشغول گشته و از
تقدیر الهی غافل شده‌اند روز بجان در تلاش معاشند و شب در تزیین اسباب فراش
آیا در هیچ شرع و ملّتی جایز است که باعراض این نفوس محدوده متمسّک شوند و از
اقبال و تصدیق نفوسی که از جان و مال و اسم و رسم و ننگ و نام در رضای حقّ
گذشته‌اند اغفال نمایند

\end{otherlanguage}

Ḥāl inṣāf dihīd kih shahādat-i īnhā maqbūl va masmūʿ ast kih qawl-shān
va fiʿl-shān muvāfiq va ẓāhir-shān va bāṭin-shān muṭābiq bi-naḥvī kih
tāhati al-ʿuqūlu fī afʿālihim wa-taḥayyarati al-nufūsu fī iṣṭibārihim
wa-bi-mā ḥamalat ajsāduhum va yā shahādat-i īn muʿriḍīn kih bi-juz
havāy-i nafs nafasī barnayārand va az qafas-i ẓunūnāt-i bāṭilih najātī
nayāftih-and va dar yawm sar az firāsh barnadārand magar chūn khuffāsh-i
ẓulmānī dar ṭalab-i dunyāy-i fāniyih kūshand va dar layl rāḥat nashavand
magar dar tadbīrāt-i umūrāt-i dāniyih kūshand bi-tadbīr-i nafsānī
mashghūl gashtih va az taqdīr-i ilāhī ghāfil shudih-and rūz bi-jān dar
talāsh-i maʿāsh-and va shab dar tazyīn-i asbāb-i firāsh āyā dar hīch
sharʿ va millatī jāʾiz ast kih bi-iʿrāḍ-i īn nufūs-i maḥdūdih mutamassik
shavand va az iqbāl va taṣdīq-i nufūsī kih az jān va māl va ism va rasm
va nang va nām dar riḍāy-i Ḥaqq gudhashtih-and ighfāl namāyand.

Now judge fairly: should the testimony of these people be accepted and
heard, whose words and deeds agree and whose outward and inward lives
accord, so that minds are bewildered at their deeds and souls marvel at
their endurance and at what their bodies bore? Or should it be the
testimony of these opponents, who draw no breath except in pursuit of
selfish desire and have found no escape from the cage of vain
imaginings? By day they lift their heads from their beds only to strive,
like a bat of darkness, in pursuit of the fleeting world, and at night
they find no rest, striving instead to devise schemes for their petty
affairs. They are absorbed in selfish planning and oblivious to divine
decree. By day they struggle with all their strength for their
livelihood, and by night they adorn the furnishings of their beds. Is it
permissible in any religious law or faith to cling to the rejection of
these narrow souls and disregard the acceptance and affirmation of those
who have renounced life and possessions, name and position, shame and
fame, for the good pleasure of God?

\subsection{252}\label{section-442}

\begin{otherlanguage}{persian}

آیا نبود که از قبل امر سیّدالشّهدآء را اعظم امور و اکبر دلیل بر حقّیّت آن
حضرت میشمردند و میگفتند در عالم چنین امری اتّفاق نیفتاد و حقّی باین
استقامت و ظهور ظاهر نشد با اینکه امر آن حضرت از صبح تا ظهر بیشتر امتداد
نیافت ولیکن این انوار مقدّسه هیجده سنه میگذرد که بلایا از جمیع جهات مثل
باران بر آنها بارید و بچه عشق و حبّ و محبّت و ذوق که جان رایگان در سبیل
سبحان انفاق نمودند چنانچه بر همه واضح و مبرهن است با وجود این چگونه این
امر را سهل شمرند آیا در هیچ عصر چنین امر خطیری ظاهر شد و آیا اگر این
اصحاب مجاهد فی‌اللّه نباشند دیگر که مجاهد خواهد بود و آیا اینها طالب عزّت و
مکنت و ثروت بودند و آیا مقصودی جز رضای حقّ داشتند و اگر این همه اصحاب با
این آثار عجیبه و افعال غریبه باطل باشند دیگر که سزاوار است که دعوی حقّ
نماید قسم بخدا که همین فعلشان برای جمیع من علی الأرض حجّت کافی و دلیل
وافی است لو کان النّاس فی اسرار الأمر یتفکّرون و سیعلم الّذین ظلموا ایّ
منقلب ینقلبون

\end{otherlanguage}

Āyā nabūd kih az qabl amr-i Siyyiduʾsh-Shuhadāʾ rā aʿẓam-i umūr va akbar
dalīl bar ḥaqqīyyat-i ān ḥaḍrat mīshumurdand va mīguftand dar ʿālam
chunīn amrī ittifāq nayuftād va ḥaqqī bi-īn istiqāmat va ẓuhūr ẓāhir
nashud bā īnkih amr-i ān ḥaḍrat az ṣubḥ tā ẓuhr bīshtar imtidād nayāft
valīkin īn anvār-i muqaddasih hijdah sanih mīgudharad kih balāyā az
jamīʿ-i jihāt mithl-i bārān bar ānhā bārīd va bi-chih ʿishq va ḥubb va
maḥabbat va dhawq kih jān rāygān dar sabīl-i Subḥān infāq nimūdand
chunānchih bar hamih vāḍiḥ va mubarhan ast bā vujūd-i īn chigūnih īn amr
rā sahl shumarand āyā dar hīch ʿaṣr chunīn amr-i khaṭīrī ẓāhir shud va
āyā agar īn aṣḥāb mujāhid fī Allāh nabāshand dīgar kih mujāhid khvāhad
būd va āyā īnhā ṭālib-i ʿizzat va miknat va tharvat būdand va āyā
maqṣūdī juz riḍāy-i Ḥaqq dāshtand va agar īn hamih aṣḥāb bā īn āthār-i
ʿajībih va afʿāl-i gharībih bāṭil bāshand dīgar kih sazāvār ast kih
daʿvī-yi ḥaqq namāyad qasam bi-Khudā kih hamīn fiʿl-shān barāy-i jamīʿ-i
man ʿalā al-arḍ ḥujjat-i kāfī va dalīl-i vāfī ast law kāna al-nāsu fī
asrāri al-amri yatafakkarūn wa-sa-yaʿlamu alladhīna ẓalamū ayya
munqalabin yanqalibūn.

Did they not formerly regard the episode of the Lord of Martyrs as the
greatest of events and the strongest proof of His truth? They would say
that no such event had occurred in the world, and that no truth had
appeared with such steadfastness and manifest power, even though His
ordeal lasted no longer than from morning until noon. Yet for eighteen
years afflictions have rained down upon these holy lights from every
direction. With what love, affection, devotion, and joy they freely gave
their lives in the path of the All-Glorious, as is clear and evident to
everyone! How, despite this, can they make light of this Cause? Has such
a momentous event appeared in any age? If these companions are not those
who strive in the path of God, who else could be? Were they seeking
honor, power, and wealth? Did they have any aim other than the good
pleasure of God? If all these companions, with their wondrous
achievements and extraordinary deeds, are false, who else is worthy to
lay claim to the truth? By God, their deeds alone are a sufficient proof
and a conclusive argument for all who dwell on earth, if people would
but reflect on the mysteries of the Cause. ``Those who have done wrong
will soon know to what a place of return they will return.''

\subsection{253}\label{section-443}

\begin{otherlanguage}{persian}

و از همه گذشته علامت صدق و کذب در کتاب معلوم و مقرّر شده باید ادّعا و
دعاوی کلّ عباد باین محکّ الهی زده شود تا صادق را از کاذب تمیز دهد اینست که
میفرماید فتمنّوا الموت ان کنتم صادقین حال ملاحظه فرمائید با این شهدای
صادق که نصّ کتاب شاهد بر صدق قول ایشانست چنانچه دیده‌اید که همه جان و مال
و زن و فرزند و کلّ مایملک را انفاق نموده‌اند و باعلی غرف رضوان عروج
فرمودند شهادت این طلعات عالیه و انفس منقطعه بر تصدیق این امر عالی متعالی
مقبول نیست و شهادت این گروه که برای ذهب از مذهب گذشته‌اند و برای جلوس بر
صدر از اوّل ماصدر احتراز جسته‌اند بر بطلان این نور لائح جایز و مقبولست با
اینکه جمیع مردم ایشان را شناخته‌اند و اینقدر ادراک نموده‌اند که از ذرّه‌ئی
از اعتبار ظاهری ملکی در سبیل دین الهی نمیگذرند تا چه رسد بجان و مال و
غیره

\end{otherlanguage}

Va az hamih gudhashtih ʿalāmat-i ṣidq va kidhb dar kitāb maʿlūm va
muqarrar shudih bāyad iddiʿā va daʿāvī-yi kull-i ʿibād bi-īn miḥakk-i
ilāhī zadih shavad tā ṣādiq rā az kādhib tamyīz dahad īn ast kih
mīfarmāyad fa-tamannawū al-mawta in kuntum ṣādiqīn ḥāl mulāḥiẓih
farmāʾīd bā īn shuhadāy-i ṣādiq kih naṣṣ-i kitāb shāhid bar ṣidq-i
qawl-i īshān ast chunānchih dīdih-īd kih hamih jān va māl va zan va
farzand va kull-i mā yamlik rā infāq nimūdih-and va bi-aʿlā ghuraf-i
riḍvān ʿurūj farmūdand shahādat-i īn ṭalʿathāy-i ʿāliyih va anfus-i
munqaṭiʿih bar taṣdīq-i īn amr-i ʿālī-yi mutaʿālī maqbūl nīst va
shahādat-i īn gurūh kih barāy-i dhahab az madhhab gudhashtih-and va
barāy-i julūs bar ṣadr az avval-i mā ṣadar iḥtirāz justih-and bar
buṭlān-i īn nūr-i lāʾiḥ jāʾiz va maqbūl ast bā īnkih jamīʿ-i mardum
īshān rā shinākhtih-and va īnqadr idrāk nimūdih-and kih az dharrih-ī az
iʿtibār-i ẓāhirī-yi mulkī dar sabīl-i dīn-i ilāhī namīgudharand tā chih
risad bi-jān va māl va ghayrih.

Beyond all this, the sign of truth and falsehood has been made known and
established in the Book. The claims and assertions of all people must be
tested against this divine touchstone, so that it may distinguish the
truthful from the liar. Thus He says, ``Then wish for death, if you are
truthful.'' Now consider these truthful martyrs, whose truthfulness is
attested by the explicit text of the Book. As you have seen, they all
sacrificed life and possessions, wife and children, and everything they
owned, and ascended to the loftiest chambers of Paradise. Is the
testimony of these exalted beings and detached souls in affirmation of
this lofty and transcendent Cause unacceptable, while the testimony of
this group who have forsaken their faith for gold and shunned the First
Emanation in order to occupy the seat of honor is admissible and
acceptable as proof of the falsity of this shining Light? Yet everyone
knows them and has understood at least this much: they will not
relinquish a particle of outward, worldly prestige in the path of God's
religion, much less their lives, possessions, or anything else.

\subsection{254}\label{section-444}

\begin{otherlanguage}{persian}

حال ملاحظه فرمائید که چگونه محکّ الهی بنصّ کتاب تفصیل نمود و خالص را از غش
تمیز داده و مع‌ذلک هنوز شاعر نشده‌اند و در نوم غفلت بکسب دنیای فانیه و
ریاست ظاهریّه مشغول شده‌اند

\end{otherlanguage}

Ḥāl mulāḥiẓe farmāyīd ke chegūne miḥakk-i ilāhī bi-naṣṣ-i kitāb tafṣīl
namūd va khāliṣ rā az ghishsh tamyīz dāde va maʿa-dhālik hanūz shāʿir
nashode-and va dar nawm-i ghaflat be-kasb-i dunyā-yi fāniye va riyāsat-i
ẓāhiriyye mashghūl shode-and.

Now consider how the divine touchstone has, by the explicit text of the
Book, made the distinction and separated the pure from the adulterated.
Yet they have still not become aware, and in the sleep of heedlessness
have occupied themselves with acquiring the things of this fleeting
world and outward leadership.

\subsection{255}\label{section-445}

\begin{otherlanguage}{arabic}

یا ابن الانسان قد مضی علیک ایّام و اشتغلت فیها بما تهوی به نفسک من الظّنون
و الأوهام الی متی تکون راقداً علی بساطک فارفع رأسک عن النّوم فانّ الشّمس قد
ارتفعت فی وسط الزّوال لعلّ تشرق علیک بأنوار الجمال والسّلام

\end{otherlanguage}

Yā ibn al-insān qad maḍā ʿalayka ayyāmun wa-ishtaghalta fīhā bi-mā tahwā
bihi nafsuka min al-ẓunūn wa-l-awhām ilā matā takūnu rāqidan ʿalā
bisāṭika fa-rfaʿ raʾsaka ʿan al-nawm fa-inna al-shamsa qad irtafaʿat fī
wasaṭ al-zawāl laʿalla tushriqa ʿalayka bi-anwār al-jamāl wa-l-salām.

O son of man! Days have passed over you, and you have occupied yourself
during them with the conjectures and illusions your soul desires. How
long will you remain asleep upon your couch? Raise your head from sleep,
for the sun has risen to the height of noon, that it may shine upon you
with the lights of beauty. And peace.

\subsection{256}\label{section-446}

\begin{otherlanguage}{persian}

ولکن معلوم باشد که این علما و فقها که مذکور شد هیچیک ریاست ظاهره
نداشته‌اند چه که محالست علمای مقتدر معروف عصر که بر صدر حکم جالسند و بر
سریر امر ساکن تابع حقّ شوند الّا من شآء ربّک چنین امری در عالم ظهور ننمود
مگر قلیلی و قلیل من عبادی الشّکور چنانچه در این عهد احدی از علمای مشهور
که زمام ناس در قبضهٔ حکم ایشان بود اقبال نجستند بلکه بتمام بغض و انکار در
دفع کوشیدند بقسمی که هیچ گوشی نشنیده و هیچ چشمی ندیده

\end{otherlanguage}

Va-līkin maʿlūm bāshad ke īn ʿulamā va fuqahā ke madhkūr shod hīch-yek
riyāsat-i ẓāhire nadāshte-and che ke muḥālast ʿulamā-yi muqtadir-i
maʿrūf-i ʿaṣr ke bar ṣadr-i ḥukm jālisand va bar sarīr-i amr sākin
tābiʿ-i Ḥaqq shavand illā man shāʾa rabbuka chunīn amrī dar ʿālam ẓuhūr
nanamūd magar qalīlī wa-qalīlun min ʿibādiya al-shakūr chunānche dar īn
ʿahd aḥadī az ʿulamā-yi mashhūr ke zimām-i nās dar qabḍe-yi ḥukm-i īshān
būd iqbāl najostand balke be-tamām-i bughḍ va inkār dar dafʿ kūshīdand
be-qismī ke hīch gūshī nashnīde va hīch chashmī nadīde.

But let it be known that none of these scholars and jurists who have
been mentioned held outward leadership. For it is impossible for the
powerful and renowned scholars of the age, who occupy the seat of
judgment and dwell upon the throne of authority, to become followers of
the Truth, ``except those whom your Lord wills.'' Such a thing has never
appeared in the world, save among a few: ``And few of My servants are
thankful.'' Thus, in this age, not one of the renowned scholars who held
the reins of the people within the grasp of their authority turned
toward Him. Rather, with utter hatred and denial, they strove to repel
Him in a manner that no ear has heard and no eye has seen.

\subsection{257}\label{section-447}

\begin{otherlanguage}{persian}

و ربّ اعلی روح ما سواه فداه بخصوص بجمیع علمای هر بلدی توقیعی صادر
فرموده‌اند و مراتب اعراض و اغماض هر کدام را در توقیع او بتفصیل ذکر
فرموده‌اند فاعتبروا یا اولی الأبصار و مقصود از آن ذکر آن بود که مباد اهل
بیان در ظهور مستغاث فی القیامة الأخری اعتراض نمایند که در ظهور بیان جمعی
از علما موقن گشته‌اند و چرا در این ظهور نشد و نعوذ باللّه متمسّک باینگونه
مزخرفات شوند و از جمال الهی محروم گردند بلی این علما که مذکور شد اکثری
معروف نبودند و بفضل اللّه از ریاست ظاهره و زخارف فانیه جمیع مقدّس و منزّه
بوده‌اند ذلک من فضل اللّه یؤتیه من یشآء

\end{otherlanguage}

Va Rabb-i Aʿlā rūḥ-i mā sivāh fidāh be-khuṣūṣ be-jamīʿ-i ʿulamā-yi har
baladī tawqīʿī ṣādir farmūde-and va marātib-i iʿrāḍ va ighmāḍ-i har
kudām rā dar tawqīʿ-i ū be-tafṣīl dhikr farmūde-and fa-ʿtabirū yā ulī
al-abṣār va maqṣūd az ān dhikr ān būd ke mabād ahl-i Bayān dar ẓuhūr-i
Mustaghāth fī al-qiyāmati al-ukhrā iʿtirāḍ namāyand ke dar ẓuhūr-i Bayān
jamʿī az ʿulamā mūqin gashte-and va cherā dar īn ẓuhūr nashod va naʿūdhu
bi-llāh mutamassik be-īn-gūne muzakhrafāt shavand va az jamāl-i ilāhī
maḥrūm gardand balī īn ʿulamā ke madhkūr shod aktharī maʿrūf nabūdand va
be-faḍl-i Allāh az riyāsat-i ẓāhire va zakhārif-i fāniye jamīʿ muqaddas
va munazzah būde-and dhālika min faḍli Allāh yuʾtīhi man yashāʾ.

And the Most Exalted Lord---may the souls of all besides Him be
sacrificed for Him---issued a specific epistle to the scholars of every
city, and in the epistle addressed to each described in detail the
extent of his rejection and disregard. ``Take heed, O people of
insight!'' The purpose of mentioning this was that the people of the
Bayán should not, at the appearance of Mustagháth in the latter
Resurrection, object that a number of scholars became convinced in the
Revelation of the Bayán, and ask why this has not happened in this
Revelation; lest, God forbid, they cling to such absurdities and be
deprived of the divine Beauty. Indeed, most of the scholars who have
been mentioned were not well known, and by the grace of God all were
sanctified and detached from outward leadership and fleeting vanities.
``That is the grace of God; He bestows it upon whomever He wills.''

\subsection{258}\label{section-448}

\begin{otherlanguage}{persian}

و دلیل و برهان دیگر که چون شمس بین دلائل مشرقست استقامت آن جمال ازلیست
بر امر الهی که با اینکه در سنّ شباب بودند و امری که مخالف کلّ اهل ارض از
وضیع و شریف و غنی و فقیر و عزیز و ذلیل و سلطان و رعیّت بود با وجود این
قیام بر آن امر فرمود چنانچه کل استماع نمودند و از هیچکس و هیچ نفس خوف
ننمودند و اعتنا نفرمودند آیا میشود این بغیر امر الهی و مشیّت مثبتهٔ ربّانی
قسم بخدا که اگر کسی فکر و خیال چنین امری نماید فی‌الفور هلاک شود و اگر
قلبهای عالم را در قلبش جا دهی باز جسارت بر چنین امر مهم ننماید مگر باذن
الهی باشد و قلبش متّصل بفیوضات رحمانی و نفسش مطمئن بعنایات ربّانی آیا این
را بچه حمل میکنند آیا بجنون نسبت میدهند چنانچه بانبیای قبل دادند و یا
میگویند برای ریاست ظاهره و جمع زخارف دنیای فانیه این امور را متعرّض
شده‌اند

\end{otherlanguage}

Va dalīl va burhān-i dīgar ke chūn shams bayn-i dalāʾil mushriqast
istiqāmat-i ān jamāl-i azalīst bar amr-i ilāhī ke bā īnke dar sinn-i
shabāb būdand va amrī ke mukhālif-i kull-i ahl-i arḍ az vaḍīʿ va sharīf
va ghanī va faqīr va ʿazīz va dhalīl va sulṭān va raʿiyyat būd bā
vujūd-i īn qiyām bar ān amr farmūd chunānche kull istimāʿ namūdand va az
hīch-kas va hīch nafs khawf nanamūdand va iʿtinā nafarmūdand āyā
mīshavad īn be-ghayr-i amr-i ilāhī va mashiyyat-i muthbite-yi rabbānī
qasam be-Khudā ke agar kasī fikr va khayāl-i chunīn amrī namāyad fī
al-fawr halāk shavad va agar qalbhā-yi ʿālam rā dar qalbash jā dihī bāz
jasārat bar chunīn amr-i muhimm nanamāyad magar be-idhn-i ilāhī bāshad
va qalbash muttaṣil be-fuyūḍāt-i raḥmānī va nafsash muṭmaʾinn
be-ʿināyāt-i rabbānī āyā īn rā be-che ḥaml mīkunand āyā be-junūn nisbat
mīdihand chunānche be-anbiyā-yi qabl dādand va yā mīgūyand barā-yi
riyāsat-i ẓāhire va jamʿ-i zakhārif-i dunyā-yi fāniye īn umūr rā
mutaʿarriḍ shode-and.

Another proof and testimony, shining like the sun among the proofs, is
the steadfastness of that eternal Beauty in the Cause of God. Although
He was in His youth, and His Cause stood opposed to all the people of
the earth---lowly and noble, rich and poor, mighty and abased, sovereign
and subject---He nevertheless arose to uphold that Cause, as everyone
has heard, fearing no one and paying no heed to any soul. Could this
come about except through the command of God and His confirming will? I
swear by God that if anyone were even to conceive or imagine such a
thing, he would perish immediately. Even if you placed all the hearts of
the world within his heart, he would still not dare undertake so
momentous a task unless it were by divine permission, his heart
connected to the outpourings of the All-Merciful and his soul assured of
divine favors. How do they account for this? Do they attribute it to
madness, as they did with the Prophets of old? Or do they say that He
undertook these things for the sake of outward leadership and the
accumulation of the vanities of this fleeting world?

\subsection{259}\label{section-449}

\begin{otherlanguage}{persian}

سبحان‌اللّه در اوّل از کتب خود که آن را قیّوم اسماء نامیده و اوّل و اعظم و
اکبر جمیع کتب است اخبار از شهادت خود میدهند و در مقامی این آیه را ذکر
فرموده‌اند یا بقیّة ‌اللّه قد فدیت بکلّی لک و رضیت السّبّ فی سبیلک و ما تمنّیت
الّا القتل فی محبّتک و کفی باللّه العلیّ معتصماً قدیماً

\end{otherlanguage}

Subḥān Allāh dar avval az kutub-i khud ke ān rā Qayyūm-i Asmā nāmīde va
avval va aʿẓam va akbar-i jamīʿ-i kutub ast ikhbār az shahādat-i khud
mīdihand va dar maqāmī īn āye rā dhikr farmūde-and yā Baqiyyata Allāh
qad fadaytu bi-kullī laka wa-raḍītu al-sabba fī sabīlika wa-mā
tamannaytu illā al-qatla fī maḥabbatika wa-kafā bi-llāhi al-ʿAliyyi
muʿtaṣaman qadīman.

Glory be to God! In the first of His books, which He named the
Qayyúm-i-Asmá, the first, greatest, and most exalted of all His books,
He foretells His own martyrdom. In one passage He states this verse: ``O
Remnant of God! I have sacrificed my whole being for You and have
accepted abuse in Your path. I have desired nothing but to be slain for
love of You. Sufficient is God, the Most High, as an everlasting
refuge.''

\subsection{260}\label{section-450}

\begin{otherlanguage}{arabic}

و همچنین در تفسیر هاء تمنّای شهادت خود را نموده‌اند کأنّی سمعت منادیاً ینادی
فی سرّی افد احبّ الأشیآء الیک فی سبیل اللّه کما فدی الحسین علیه السّلام فی
سبیلی و لو لا کنت ناظراً بذلک السّرّ الواقع فوالّذی نفسی بیده لو اجتمعوا
ملوک الأرض لن یقدروا ان یأخذوا منّی حرفاً فکیف عبید الّذی لیس لهم شأن بذلک
و انّهم مطرودون الی ان قال لیعلم الکلّ مقام صبری و رضائی و فدائی فی سبیل
اللّه

\end{otherlanguage}

Va hamchunīn dar Tafsīr-i Hā tamannā-yi shahādat-i khud rā namūde-and
ka-annī samiʿtu munādiyan yunādī fī sirrī ifdi aḥabba al-ashyāʾi ilayka
fī sabīli Allāh kamā fadā al-Ḥusayn ʿalayhi al-salām fī sabīlī wa-law lā
kuntu nāẓiran bi-dhālika al-sirri al-wāqiʿ fa-wa-alladhī nafsī bi-yadihi
law ijtamaʿū mulūku al-arḍ lan yaqdirū an yaʾkhudhū minnī ḥarfan
fa-kayfa ʿabīdu alladhī laysa lahum shaʾnun bi-dhālika wa-innahum
maṭrūdūn ilā an qāl li-yaʿlama al-kullu maqāma ṣabrī wa-riḍāʾī wa-fidāʾī
fī sabīli Allāh.

Likewise, in the Commentary on Há, He expressed His desire for
martyrdom: ``It is as though I heard a caller calling within My inmost
being: `Sacrifice what is dearest to You in the path of God, just as
Ḥusayn, peace be upon Him, sacrificed in My path.' Were I not mindful of
that actual mystery, I swear by Him in whose hand is My soul, even if
the kings of the earth gathered together, they would never be able to
take a single letter from Me, much less these servants who have no
standing in this matter and are indeed outcasts.'' He continues until He
says: ``That all may know the station of My patience, My contentment,
and My sacrifice in the path of God.''

\subsection{261}\label{section-451}

\begin{otherlanguage}{persian}

آیا صاحب این بیان را میتوان نسبت داد که در غیر صراط الهی مشی مینماید و
یا بغیر رضای او امری طلب نموده در همین آیه نسیم انقطاعی مکنون شده که اگر
بوزد جمیع هیاکل وجود جان را انفاق نمایند و از روان درگذرند حال ملاحظه
نمائید که چه ‌قدر ناس نسناسند و بغایت حقّ‌ناسپاس که چشم از جمیع اینها
پوشیده‌اند و بعقب مرداری چند که از بطنشان افغان مال مسلمانان میآید میدوند
و با وجود این چه نسبتهای غیر لایقه که بمطالع قدسیّه میدهند کذلک نذکر لک
ما اکتسبت ایدی الّذینهم کفروا و اعرضوا عن لقآء اللّه فی یوم القیامة و
عذّبهم اللّه بنار شرکهم و اعدّ لهم فی الآخرة عذاباً تحترق به اجسادهم و
ارواحهم ذلک بأنّهم قالوا انّ اللّه لم یکن قادراً علی شیء و کانت یده عن الفضل
مغلولة

\end{otherlanguage}

Āyā ṣāḥib-i īn bayān rā mītavān nisbat dād ke dar ghayr-i ṣirāṭ-i ilāhī
mashy mīnamāyad va yā be-ghayr-i riḍā-yi ū amrī ṭalab namūde dar hamīn
āye nasīm-i inqiṭāʿī maknūn shode ke agar bivazad jamīʿ-i hayākil-i
vujūd jān rā infāq namāyand va az ravān dargudharand ḥāl mulāḥiẓe
namāyīd ke che qadr nās nasnāsand va be-ghāyat Ḥaqq-nāsipās ke chashm az
jamīʿ-i īnhā pūshīde-and va be-ʿaqab-i murdārī chand ke az baṭnishān
afghān-i māl-i musalmānān mīāyad mīdavand va bā vujūd-i īn che
nisbathā-yi ghayr-i lāyiqe ke be-maṭāliʿ-i qudsiyye mīdihand kadhālika
nadhkuru laka mā iktasabat aydī alladhīnahum kafarū wa-aʿraḍū ʿan liqāʾi
Allāh fī yawm al-qiyāma wa-ʿadhdhabahumu Allāhu bi-nāri shirkihim
wa-aʿadda lahum fī al-ākhirati ʿadhāban taḥtariqu bihi ajsāduhum
wa-arwāḥuhum dhālika bi-annahum qālū inna Allāha lam yakun qādiran ʿalā
shayʾ wa-kānat yaduhu ʿan al-faḍli maghlūla.

Can the author of these words be accused of walking any path other than
the path of God, or of seeking anything other than His good pleasure?
Within this very verse lies a breeze of detachment which, if it were to
blow, would cause all the embodiments of existence to offer up their
lives and relinquish their souls. Now consider how brutish people are,
and how utterly ungrateful to the Truth, that they have closed their
eyes to all these things and run after a few carrion bodies from whose
bellies comes the cry of the Muslims' wealth. Yet what unworthy
accusations they level against the holy Daysprings! Thus We recount to
you what the hands of those who disbelieved and turned away from meeting
God on the Day of Resurrection have committed. God has tormented them
with the fire of their association of partners with Him, and has
prepared for them in the hereafter a torment that will burn their bodies
and their souls. This is because they said that God had no power over
anything, and that His hand was shackled from bestowing grace.

\subsection{262}\label{section-452}

\begin{otherlanguage}{persian}

و استقامت بر امر حجّتی است بزرگ و برهانیست عظیم چنانچه خاتم انبیا فرمودند
شیّبتنی الآیتین یعنی پیر نمود مرا دو آیه که هر دو مشعر بر استقامت بر امر
الهیست چنانچه میفرماید فاستقم کما امرت

\end{otherlanguage}

Va istiqāmat bar amr ḥujjatī ast buzurg va burhānīst ʿaẓīm chunānche
Khātam-i Anbiyā farmūdand shayyabatnī al-āyatayn yaʿnī pīr namūd marā du
āye ke har du mushʿir bar istiqāmat bar amr-i ilāhīst chunānche
mīfarmāyad fa-staqim kamā umirta.

Steadfastness in the Cause is a great testimony and a mighty proof. Thus
the Seal of the Prophets said, ``The two verses have made Me grow old,''
meaning that two verses, both enjoining steadfastness in the Cause of
God, had aged Him. Thus He says: ``Be steadfast, as you have been
commanded.''

\subsection{263}\label{section-453}

\begin{otherlanguage}{persian}

حال ملاحظه فرمائید که این سدرهٔ رضوان سبحانی در اوّل جوانی چگونه تبلیغ امر
اللّه فرمود و چه ‌قدر استقامت از آن جمال احدیّت ظاهر شد که جمیع من علی
الأرض بر منعش اقدام نمودند حاصلی نبخشید آنچه ایذا بر آن سدرهٔ طوبی وارد
میآوردند شوقش بیشتر و نار حبّش مشتعلتر میشد چنانچه این فقرات واضح است و
احدی انکار ندارد تا آنکه بالأخره جان را درباخت و برفیق اعلی شتافت

\end{otherlanguage}

Ḥāl mulāḥiẓe farmāyīd ke īn sidre-yi Riḍvān-i subḥānī dar avval-i javānī
chegūne tablīgh-i amr Allāh farmūd va che qadr istiqāmat az ān jamāl-i
aḥadiyyat ẓāhir shod ke jamīʿ-i man ʿalā al-arḍ bar manʿash iqdām
namūdand ḥāṣilī nabakhshīd ānche īdhā bar ān sidre-yi Ṭūbā vārid
mīāvardand shawqash bīshtar va nār-i ḥubbash mushtaʿiltar mīshod
chunānche īn faqarāt vāḍiḥ ast va aḥadī inkār nadārad tā ānke
bi-l-ākhare jān rā darbākht va be-Rafīq-i Aʿlā shitāft.

Now consider how this Tree of the divine Paradise proclaimed the Cause
of God in the earliest days of His youth, and what steadfastness that
Beauty of Oneness manifested. All who dwell on earth arose to restrain
Him, but to no avail. The more suffering they inflicted upon that
blessed Tree, the greater His longing grew and the more intensely the
fire of His love blazed. These things are clear, and no one denies them,
until at last He gave up His life and hastened to the Most Exalted
Companion.

\subsection{264}\label{section-454}

\begin{otherlanguage}{persian}

و از جملهٔ دلائل ظهور غلبه و قدرت و احاطه‌ که بنفسه از آن مظهر وجود و مظهر
معبود در اکناف و اقطار عالم ظاهر شد چنانچه آن جمال ازلی در شیراز در سنهٔ
ستّین ظاهر شدند و کشف غطا فرمودند مع‌ذلک باندک زمانی آثار غلبه و قدرت و
سلطنت و اقتدار از آن جوهر الجواهر و بحر البحور در جمیع بلاد ظاهر شد
بقسمی که از هر بلدی آثار و اشارات و دلالات و علامات آن شمس لاهوتی هویدا
گشت و چه مقدار قلوب صافیهٔ رقیقه که از آن شمس ازلیّه حکایت نمودند و چه‌ قدر
رشحات علمی از آن بحر علم لدنّی که احاطه نمود جمیع ممکنات را با اینکه در
هر بلد و مدینه جمیع علما و اعزّه بر منع و ردّ ایشان برخاستند و کمر غل و
حسد و ظلم بر دفعشان بستند و چه نفوس قدسیّه را که جواهر عدل بودند بنسبت
ظلم کشتند و چه هیاکل روح را که صرف علم و عمل از ایشان ظاهر بود ببدترین
عذاب هلاک نمودند مع کلّ ذلک هر یک از آن وجودات تا دم مرگ بذکر ذکر اللّه
مشغول بودند و در هوای تسلیم و رضا طائر و بقسمی این وجودات را تقلیب
نمودند و تصرّف فرمودند که بجز اراده‌اش مرادی نجستند و بجز امرش امری
نگزیدند رضا برضایش دادند و دل بخیالش بستند

\end{otherlanguage}

Va az jumle-yi dalāʾil ẓuhūr-i ghalabe va qudrat va iḥāṭe ke bi-nafsih
az ān maẓhar-i vujūd va muẓhir-i maʿbūd dar aknāf va aqṭār-i ʿālam ẓāhir
shod chunānche ān jamāl-i azalī dar Shīrāz dar sane-yi sittīn ẓāhir
shodand va kashf-i ghiṭā farmūdand maʿa-dhālik be-andak zamānī āthār-i
ghalabe va qudrat va salṭanat va iqtidār az ān jawhar al-javāhir va baḥr
al-buḥūr dar jamīʿ-i bilād ẓāhir shod be-qismī ke az har baladī āthār va
ishārāt va dalālāt va ʿalāmāt-i ān shams-i lāhūtī huvaydā gasht va che
miqdār qulūb-i ṣāfiye-yi raqīqe ke az ān shams-i azaliyye ḥikāyat
namūdand va che qadr rashaḥāt-i ʿilmī az ān baḥr-i ʿilm-i ladunnī ke
iḥāṭe namūd jamīʿ-i mumkināt rā bā īnke dar har balad va madīne jamīʿ-i
ʿulamā va aʿizze bar manʿ va radd-i īshān barkhāstand va kamar-i ghill
va ḥasad va ẓulm bar dafʿishān bastand va che nufūs-i qudsiyye rā ke
javāhir-i ʿadl būdand be-nisbat-i ẓulm kushtand va che hayākil-i rūḥ rā
ke ṣirf-i ʿilm va ʿamal az īshān ẓāhir būd be-badtarīn ʿadhāb halāk
namūdand maʿa kull-i dhālik har yek az ān vujūdāt tā dam-i marg
be-dhikr-i Dhikr Allāh mashghūl būdand va dar havā-yi taslīm va riḍā
ṭāʾir va be-qismī īn vujūdāt rā taqlīb namūdand va taṣarruf farmūdand ke
be-juz irāde-ash murādī najostand va be-juz amrash amrī naguzīdand riḍā
be-riḍāyash dādand va dil be-khayālash bastand.

Among the proofs is the ascendancy, power, and all-encompassing
influence that appeared, of itself, from that Manifestation of existence
and Revealer of the Adored throughout the regions and reaches of the
world. Thus that eternal Beauty appeared in Shiraz in the year sixty and
lifted the veil. Yet within a short time the signs of ascendancy, power,
sovereignty, and authority from that Essence of essences and Ocean of
oceans appeared in every land, so that in every city the traces, tokens,
proofs, and signs of that divine Sun became evident. How many pure and
tender hearts reflected that eternal Sun, and how many outpourings of
knowledge from that Ocean of innate knowledge encompassed all created
things! This was despite the fact that in every city and town all the
scholars and eminent men arose to obstruct and reject them, and girded
themselves with malice, envy, and oppression to repel them. How many
holy souls, who were jewels of justice, they killed on the accusation of
wrongdoing, and how many embodiments of spirit, who manifested nothing
but knowledge and deeds, they put to death with the cruelest torments!
Despite all this, each of those souls remained occupied until the moment
of death with remembrance of the Remembrance of God, soaring in the air
of submission and contentment. He so transformed and held sway over
these souls that they sought no desire but His will and chose no course
but His command. They made their contentment conform to His good
pleasure and fixed their hearts upon the thought of Him.

\subsection{265}\label{section-455}

\begin{otherlanguage}{persian}

حال قدری تفکّر نمائید آیا چنین تصرّف و احاطه از احدی در امکان ظاهر شده و
جمیع این قلوب منزّهه و نفوس مقدّسه بکمال رضا در موارد قضا شتافتند و در
مواقع شکایت جز شکر از ایشان ظاهر نه و در مواطن بلا جز رضا از ایشان مشهود
نه و این رتبه هم معلوم است که کلّ اهل ارض چه مقدار غل و بغض و عداوت باین
اصحاب داشتند چنانچه اذیّت و ایذای آن طلعات قدسی معنوی را علّت فوز و
رستگاری و سبب فلاح و نجاح ابدی میدانستند آیا هرگز در هیچ تاریخی از عهد
آدم تا حال چنین غوغائی در بلاد واقع شد و آیا چنین ضوضائی در میان عباد
ظاهر گشت و با این همه ایذا و اذیّت محلّ لعن جمیع ناس شدند و محلّ ملامت جمیع
عباد و گویا صبر در عالم کون از اصطبارشان ظاهر شد و وفا در ارکان عالم از
فعلشان موجود گشت

\end{otherlanguage}

Ḥāl qadrī tafakkur namāyīd āyā chunīn taṣarruf va iḥāṭe az aḥadī dar
imkān ẓāhir shode va jamīʿ-i īn qulūb-i munazzahe va nufūs-i muqaddase
be-kamāl-i riḍā dar mavārid-i qaḍā shitāftand va dar mavāqiʿ-i shikāyat
juz shukr az īshān ẓāhir na va dar mavāṭin-i balā juz riḍā az īshān
mashhūd na va īn rutbe ham maʿlūm ast ke kull-i ahl-i arḍ che miqdār
ghill va bughḍ va ʿadāvat be-īn aṣḥāb dāshtand chunānche adhiyyat va
īdhā-yi ān ṭalʿāt-i qudsī-yi maʿnavī rā ʿillat-i fawz va rastigārī va
sabab-i falāḥ va najāḥ-i abadī mīdānistand āyā hargiz dar hīch tārīkhī
az ʿahd-i Ādam tā ḥāl chunīn ghawghāʾī dar bilād vāqiʿ shod va āyā
chunīn ḍawḍāʾī dar miyān-i ʿibād ẓāhir gasht va bā īn hame īdhā va
adhiyyat maḥall-i laʿn-i jamīʿ-i nās shodand va maḥall-i malāmat-i
jamīʿ-i ʿibād va gūyā ṣabr dar ʿālam-i kawn az iṣṭibārishān ẓāhir shod
va vafā dar arkān-i ʿālam az fiʿlishān mawjūd gasht.

Now reflect a little: has such sway and all-encompassing influence ever
appeared from anyone in the world of existence? All these purified
hearts and sanctified souls hastened with perfect contentment to meet
their destiny. Where there was cause for complaint, nothing but
gratitude appeared from them, and amid affliction nothing but
contentment was seen in them. It is also well known how much malice,
hatred, and hostility all the people of the earth bore toward these
companions, to the point that they regarded tormenting and harming those
holy, spiritual countenances as a means of attainment and salvation, and
a cause of eternal deliverance and success. Has such an uproar ever
occurred in any land, in any history from the time of Adam until now?
Has such a tumult ever arisen among people? Along with all this
suffering and torment, they became the objects of everyone's curses and
everyone's reproach. It is as though patience appeared in the world of
being through their endurance, and faithfulness came into existence
throughout the world through their deeds.

\subsection{266}\label{section-456}

\begin{otherlanguage}{persian}

باری در جمیع این وقایع حادثه و حکایات وارده تفکّر فرمائید تا بر عظمت امر
و بزرگی آن مطّلع گردید تا بعنایت رحمن روح اطمینان در وجود دمیده شود و بر
سریر ایقان مستریح و جالس شوید خدای واحد شاهد است که اگر فی‌الجمله تفکّر
نمائید علاوه بر همهٔ این مطالب مقرّره و دلائل مذکوره همین رد و سب و لعن
اهل ارض بر این فوارس میدان تسلیم و انقطاع اعظم دلیل و اکبر حجّت بر حقّیّت
ایشانست و در هر آن که تفکرّ در اعتراضات جمیع مردم از علما و فضلا و جهّال
فرمائی در این امر محکمتر و راسختر و ثابت‌تر میشوی زیرا که جمیع آنچه واقع
شده از قبل معادن علم لدنّی و مواقع احکام ازلی خبر داده‌اند

\end{otherlanguage}

Bārī dar jamīʿ-i īn vaqāyiʿ-i ḥādithe va ḥikāyāt-i vāride tafakkur
farmāyīd tā bar ʿaẓamat-i amr va buzurgī-yi ān muṭṭaliʿ gardīd tā
be-ʿināyat-i Raḥmān rūḥ-i iṭmīnān dar vujūd damīde shavad va bar sarīr-i
īqān mustarīḥ va jālis shavīd Khudā-yi vāḥid shāhid ast ke agar fī
al-jumle tafakkur namāyīd ʿalāve bar hame-yi īn maṭālib-i muqarrere va
dalāʾil-i madhkūre hamīn radd va sabb va laʿn-i ahl-i arḍ bar īn
favāris-i maydān-i taslīm va inqiṭāʿ aʿẓam dalīl va akbar ḥujjat bar
ḥaqqiyyat-i īshānast va dar har ān ke tafakkur dar iʿtirāḍāt-i jamīʿ-i
mardum az ʿulamā va fuḍalā va juhhāl farmāyī dar īn amr muḥkamtar va
rāsikhtar va thābittar mīshavī zīrā ke jamīʿ-i ānche vāqiʿ shode az qabl
maʿādin-i ʿilm-i ladunnī va mavāqiʿ-i aḥkām-i azalī khabar dāde-and.

In short, reflect upon all these events that have occurred and the
accounts that have been related, so that you may recognize the grandeur
and greatness of the Cause, that through the favor of the All-Merciful
the spirit of assurance may be breathed into your being, and you may sit
at rest upon the throne of certitude. The one God is witness that, if
you reflect even a little, beyond all these established points and
proofs already mentioned, the very rejection, abuse, and curses directed
by the people of the earth against these horsemen in the field of
submission and detachment constitute the greatest proof and the
mightiest testimony to their truth. And whenever you reflect upon the
objections of all the people---scholars, learned men, and the
ignorant---you will become firmer, more deeply rooted, and more
steadfast in this Cause, for the Mines of innate knowledge and the
Sources of eternal decrees foretold all that has occurred.

\subsection{267}\label{section-457}

\begin{otherlanguage}{persian}

اگرچه این بنده ارادهٔ ذکر احادیث قبل را نداشتم ولیکن نظر بمحبّت آن جناب
چند روایتی که مناسب این مقام است ذکر مینمایم با اینکه فی‌الحقیقه احتیاج
نیست زیرا که آنچه ذکر شده جمیع ارض و من علیها را کافی است و فی‌الحقیقه
جمیع کتب و اسرار آن در این مختصر ذکر شده بقسمی که اگر کسی قدری تأمّل
نماید جمیع اسرار کلمات الهی و امور ظاهره از آن سلطان حقیقی را از آنچه
ذکر شده ادراک مینماید ولیکن چون همهٔ ناس بر یک شأن و یک مقام نیستند لهذا
چند حدیثی ذکر مینمایم تا سبب استقامت انفس متزلزله شود و اطمینان عقول
مضطربه گردد و همچنین حجّت الهی بر اعالی و ادانی عباد تام و کامل گردد

\end{otherlanguage}

Agarche īn bande erāde-ye dhekr-e aḥādīth-e qabl rā nadāshtam valīkan
naẓar be-maḥabbat-e ān janāb chand revāyatī ke monāseb-e īn maqām ast
dhekr mīnamāyam bā īnke fī'l-ḥaqīqe eḥtīāj nīst zīrā ke ānche dhekr
shode jamīʿ-e arḍ va man ʿalayhā rā kāfī ast va fī'l-ḥaqīqe jamīʿ-e
kotob va asrār-e ān dar īn mokhtaṣar dhekr shode be-qesmī ke agar kasī
qadrī taʾammol namāyad jamīʿ-e asrār-e kalemāt-e elāhī va omūr-e ẓāhere
az ān solṭān-e ḥaqīqī rā az ānche dhekr shode edrāk mīnamāyad valīkan
chon hame-ye nās bar yek shaʾn va yek maqām nīstand lehādhā chand
ḥadīthī dhekr mīnamāyam tā sabab-e esteqāmat-e anfos-e motazalzele
shavad va eṭmīnān-e ʿoqūl-e moḍṭarebe gardad va hamchenīn ḥojjat-e elāhī
bar aʿālī va adānī-ye ʿebād tāmm va kāmel gardad.

Although I, this servant, had not intended to cite the traditions of the
past, out of affection for you I will mention a few accounts appropriate
to this subject, even though there is truly no need, for what has been
stated suffices for the whole earth and all who dwell upon it. Indeed,
all the scriptures and their mysteries have been set forth in this brief
account, so that anyone who reflects a little will perceive, from what
has been stated, all the mysteries of the divine words and the acts
proceeding from that true Sovereign. Since, however, not everyone is of
the same disposition or station, I will cite a few traditions, so that
they may bring steadfastness to wavering souls and assurance to troubled
minds, and so that the proof of God may be full and complete for the
high and low among His servants.

\subsection{268}\label{section-458}

\begin{otherlanguage}{persian}

از جملهٔ احادیث اینست که میفرماید اذا ظهرت رایة الحقّ لعنها اهل الشّرق و
الغرب حال باید قدری از صهبای انقطاع نوشید و بر رفرف امتناع مقرّ گزید و
تفکّر ساعة خیر من عبادة سبعین سنة را منظور داشت که آخر سبب این امر شنیع
چه میشود که جمیع مردم با اظهار حبّ و طلب حقّ بعد از ظهور اهل حقّ را لعن
نمایند چنانچه مستفاد از حدیث میشود و این واضحست که سبب نسخ قواعد و رسوم
و عادات و آدابست که همهٔ ناس بآن محدود گشته‌اند والّا اگر جمال رحمن بر همان
رسوم و آداب حرکت نماید و تصدیق کند مردم را در آنچه بآن مشغولند دیگر چرا
این همه اختلاف و فساد در ملک ظاهر میشود و مصدّق و مثبت این حدیث شریف قوله
تعالی یوم یدعو الدّاع الی شیء نکر

\end{otherlanguage}

Az jomle-ye aḥādīth īnast ke mīfarmāyad: idhā ẓaharat rāyat al-Ḥaqq
laʿanahā ahl al-sharq wa-l-gharb. Ḥāl bāyad qadrī az ṣahbā-ye enqeṭāʿ
nūshīd va bar rafraf-e emtenāʿ maqarr gozīd va tafakkur sāʿa khayr min
ʿibādat sabʿīn sana rā manẓūr dāsht ke ākhar sabab-e īn amr-e shanīʿ che
mīshavad ke jamīʿ-e mardom bā eẓhār-e ḥobb va ṭalab-e Ḥaqq baʿd az ẓohūr
ahl-e Ḥaqq rā laʿn namāyand chenānche mostafād az ḥadīth mīshavad va īn
vāḍeḥast ke sabab naskh-e qavāʿed va rosūm va ʿādāt va ādāb ast ke
hame-ye nās be-ān maḥdūd gashte-and va-ellā agar jamāl-e Raḥmān bar
hamān rosūm va ādāb ḥarakat namāyad va taṣdīq konad mardom rā dar ānche
be-ān mashghūland dīgar cherā īn hame ekhtelāf va fasād dar molk ẓāher
mīshavad va moṣaddeq va mothbet-e īn ḥadīth-e sharīf qawluhu taʿālā:
yawma yadʿu al-dāʿi ilā shayʾin nukur.

Among these traditions is the saying: ``When the banner of the Truth
appears, the people of the East and the West will curse it.'' Now one
must drink a little of the wine of detachment, take one's seat upon the
heights of inaccessibility, and bear in mind that ``An hour's reflection
is better than seventy years of worship.'' What, after all, causes this
shameful act---that all people, despite professing love for God and
seeking Him, should curse the people of truth after the Manifestation
appears, as this tradition indicates? It is clear that the cause is the
abrogation of the rules, customs, habits, and practices within whose
bounds all people have confined themselves. Otherwise, if the Beauty of
the All-Merciful were to follow those same customs and practices and
confirm people in their pursuits, why should all this discord and
turmoil appear in the world? Confirming and substantiating this noble
tradition are His words, exalted is He: ``On the day when the caller
calls to a thing unknown.''

\subsection{269}\label{section-459}

\begin{otherlanguage}{persian}

باری چون منادی احدیّه از ورای حجبات قدسیّه مردم را بانقطاع تمام میخواند از
آنچه که در دست دارند و این ندای الهی چون مخالف هوی است لهذا اینهمه
افتتان و امتحان رو میدهد و حال مردم را ملاحظه نما که هیچ ذکر این احادیث
محکمه را که جمیع ظاهر شده نمینمایند ولیکن آن احادیثی که صحّت و سقم آن
معلوم نیست تمسّک بآنها جسته‌اند که چرا ظاهر نشد و حال آنکه آنچه را هم که
تعقّل ننموده‌اند ظاهر شد و باهر گشت و آثار و علامات حقّ بمثل شمس در وسط
سماء لائح مع‌ذلک عباد در تیه جهل و نادانی سرگردان مانده‌اند با اینکه چه
‌قدر از آیات فرقانیّه و روایات محقّقه که جمیع دالّ است بر شرع و حکم جدید و
امر بدیع باز منتظرند که طلعت موعود بر شریعت فرقان حکم فرماید چنانچه یهود
و نصاری همین حرف را میگویند

\end{otherlanguage}

Bārī chon monādī-ye aḥadīye az varā-ye ḥojobāt-e qodsīye mardom rā
be-enqeṭāʿ-e tamām mīkhwānad az ānche ke dar dast dārand va īn nedā-ye
elāhī chon mokhālef-e havā ast lehādhā īnhame eftetān va emteḥān rū
mīdahad va ḥāl-e mardom rā molāḥeẓe namā ke hīch dhekr-e īn aḥādīth-e
moḥkame rā ke jamīʿ ẓāher shode nemīnamāyand valīkan ān aḥādīthī ke
ṣeḥḥat va saqam-e ān maʿlūm nīst tamassok be-ānhā joste-and ke cherā
ẓāher nashod va ḥāl ānke ānche rā ham ke taʿaqqol nanamūde-and ẓāher
shod va bāher gasht va āthār va ʿalāmāt-e Ḥaqq be-methl-e shams dar
vasaṭ-e samāʾ lāʾeḥ maʿa-dhālek ʿebād dar tīh-e jahl va nādānī sargardān
mānde-and bā īnke che qadr az āyāt-e forqānīye va revāyāt-e moḥaqqaqe ke
jamīʿ dāll ast bar sharʿ va ḥokm-e jadīd va amr-e badīʿ bāz montaẓerand
ke ṭalʿat-e mawʿūd bar sharīʿat-e Forqān ḥokm farmāyad chenānche Yahūd
va Naṣārā hamīn ḥarf rā mīgūyand.

In short, since the Herald of divine unity calls people from beyond the
veils of holiness to complete detachment from all they possess, and
since this divine call runs counter to their desires, all these trials
and tests occur. Consider the state of the people: they make no mention
of these firmly established traditions, all of which have been
fulfilled, but cling to traditions whose authenticity or falsity is
unknown, asking why they have not been fulfilled. Yet even things they
had never conceived have appeared and become manifest, and the signs and
tokens of the Truth shine like the sun in the midst of heaven.
Nevertheless, His servants remain bewildered in the wilderness of
ignorance and unknowing. Despite the many verses of the Qur'an and
authenticated traditions, all pointing to a new law, a new ordinance,
and a wondrous Cause, they still expect the promised Countenance to rule
according to the law of the Qur'an, just as the Jews and Christians say
the same thing.

\subsection{270}\label{section-460}

\begin{otherlanguage}{persian}

و از جملهٔ کلمات مدلّهٔ بر شرع جدید و امر بدیع فقرات دعای ندبه است که
میفرماید این المدخّر لتجدید الفرائض و السّنن و این المتخیّر لاعادة الملّة و
الشّریعة و در زیارت میفرماید السّلام علی الحقّ الجدید سأل ابوعبداللّه عن
سیرة المهدیّ کیف سیرته قال (ع) یصنع ما صنع رسول اللّه (ص) و یهدم ما کان
قبله کما هدم رسول اللّه امر الجاهلیّة

\end{otherlanguage}

Va az jomle-ye kalemāt-e modelle-ye bar sharʿ-e jadīd va amr-e badīʿ
faqarāt-e Doʿā-ye Nodbe ast ke mīfarmāyad: ayna al-muddakhar li-tajdīd
al-farāʾiḍ wa-l-sunan wa-ayna al-mutakhayyar li-iʿādat al-milla
wa-l-sharīʿa. Va dar zīyārat mīfarmāyad: al-salām ʿalā al-Ḥaqq al-jadīd.
Suʾila Abū ʿAbd Allāh ʿan sīrat al-Mahdī kayfa sīratuhu qāla (ʿ):
yaṣnaʿu mā ṣanaʿa Rasūl Allāh (ṣ) wa-yahdimu mā kāna qablahu kamā hadama
Rasūl Allāh amr al-jāhiliyya.

Among the words indicating a new law and a wondrous Cause are these
passages from the Prayer of Lamentation: ``Where is the One held in
reserve to renew the obligations and practices? And where is the One
chosen to restore the faith and the law?'' And in the visitation prayer
it is said: ``Peace be upon the new Truth.'' Abu 'Abd Allah was asked
about the conduct of the Mahdi: ``What will his conduct be?'' He
replied, peace be upon him: ``He will do as the Messenger of God did,
blessings be upon him, and will demolish what came before him, just as
the Messenger of God demolished the order of the age of ignorance.''

\subsection{271}\label{section-461}

\begin{otherlanguage}{persian}

حال ملاحظه فرمائید که با وجود امثال این روایات چه استدلالها بر عدم تغییر
احکام مینمایند با اینکه مقصود از هر ظهور ظهور تغییر و تبدیل است در ارکان
عالم سرّاً و جهراً ظاهراً و باطناً چه اگر بهیچوجه امورات ارض تغییر نیابد ظهور
مظاهر کلّیّه لغو خواهد بود و با اینکه در عوالم که از کتب مشهورهٔ معتبره است
میفرماید یظهر من بنی هاشم صبیّ ذو کتاب و احکام جدید الی ان قال و اکثر
اعدائه العلمآء و در مقامی دیگر از صادق بن محمّد ذکر مینماید که فرمودند و
لقد یظهر صبیّ من بنی هاشم و یأمر النّاس ببیعته و هو ذو کتاب جدید یبایع
النّاس بکتاب جدید علی العرب شدید فان سمعتم منه شیئاً فاسرعوا الیه خوب وصیّت
ائمّهٔ دین و سرج یقین را عمل نمودند با اینکه میفرماید اگر شنیدید که جوانی
از بنی هاشم ظاهر شد و میخواند مردم را بکتاب جدید الهی و احکام بدیع ربّانی
بشتابید بسوی او مع‌ذلک جمیع حکم کفر و خروج از ایمان بآن سیّد امکان دادند و
نرفتند بسوی آن نور هاشمی و ظهور سبحانی مگر با شمشیرهای کشیده و قلبهای
پرکینه و دیگر ملاحظهٔ عداوت علما نمائید که بچه صریحی در کتب مذکور است با
وجود همۀ این احادیث ظاهرهٔ مدلّه و اشارات واضحهٔ محقّقه جمیع ناس از جوهر
صافی معرفت و بیان معرض شده‌اند و بمظاهر ضلالت و طغیان اقبال نموده‌اند و با
این روایات وارده و کلمات نازله میگویند آنچه نفسشان بآن مایل است و اگر
جوهر حقّ بیانی بفرماید که مخالف نفس و هوای این گروه واقع شود فی‌‌الفور
تکفیر نمایند و میگویند این مخالف قول ائمّهٔ دین و انوار مبین است و در شرع
متین چنین امری و حکمی صادر نشده چنانچه الیوم امثال این سخنهای بی‌فائده از
این هیاکل فانیه ظاهر شده و میشود

\end{otherlanguage}

Ḥāl molāḥeẓe farmāʾīd ke bā vojūd-e amthāl-e īn revāyāt che estedlālhā
bar ʿadam-e taghyīr-e aḥkām mīnamāyand bā īnke maqṣūd az har ẓohūr
ẓohūr-e taghyīr va tabdīl ast dar arkān-e ʿālam serran va jahran ẓāheran
va bāṭenan che agar be-hīchvajh omūrāt-e arḍ taghyīr nayābad ẓohūr-e
maẓāher-e kollīye laghv khwāhad būd va bā īnke dar ʿAvālem ke az kotob-e
mashhūre-ye moʿtabare ast mīfarmāyad: yaẓharu min Banī Hāshim ṣabiyy dhū
kitāb wa-aḥkām jadīd ilā an qāl: wa-akthar aʿdāʾihi al-ʿulamāʾ. Va dar
maqāmī dīgar az Ṣādiq ibn Muḥammad dhekr mīnamāyad ke farmūdand:
wa-laqad yaẓharu ṣabiyy min Banī Hāshim wa-yaʾmuru al-nās bi-bayʿatihi
wa-huwa dhū kitāb jadīd yubāyiʿu al-nās bi-kitāb jadīd ʿalā al-ʿArab
shadīd fa-in samiʿtum minhu shayʾan fa-asriʿū ilayhi. Khūb vaṣīyat-e
aʾemme-ye dīn va soroj-e yaqīn rā ʿamal namūdand bā īnke mīfarmāyad agar
shenīdīd ke javānī az Banī Hāshim ẓāher shod va mīkhwānad mardom rā
be-ketāb-e jadīd-e elāhī va aḥkām-e badīʿ-e rabbānī beshetābīd be-sū-ye
ū maʿa-dhālek jamīʿ ḥokm-e kofr va khorūj az īmān be-ān seyyed-e emkān
dādand va naraftand be-sū-ye ān nūr-e Hāshimī va ẓohūr-e sobḥānī magar
bā shamshīrhā-ye kashīde va qalbhā-ye porkīne va dīgar molāḥeẓe-ye
ʿadāvat-e ʿolamā namāʾīd ke be-che ṣarīḥī dar kotob madhkūr ast bā
vojūd-e hame-ye īn aḥādīth-e ẓāhere-ye modelle va eshārāt-e vāḍeḥe-ye
moḥaqqaqe jamīʿ-e nās az jawhar-e ṣāfī-ye maʿrefat va bayān moʿreḍ
shode-and va be-maẓāher-e ḍalālat va ṭoghyān eqbāl namūde-and va bā īn
revāyāt-e vārede va kalemāt-e nāzele mīgūyand ānche nafseshān be-ān
māyel ast va agar jawhar-e Ḥaqq bayānī befarmāyad ke mokhālef-e nafs va
havā-ye īn gorūh vāqeʿ shavad fī'l-fawr takfīr namāyand va mīgūyand īn
mokhālef-e qawl-e aʾemme-ye dīn va anvār-e mobīn ast va dar sharʿ-e
matīn chenīn amrī va ḥokmī ṣāder nashode chenānche al-yawm amthāl-e īn
sokhanhā-ye bī-fāʾede az īn hayākel-e fānīye ẓāher shode va mīshavad.

Now consider how, despite such traditions, they advance all manner of
arguments against any change in the ordinances, even though the purpose
of every Revelation is to bring about change and transformation in the
foundations of the world, secretly and openly, outwardly and inwardly.
For if the affairs of the earth were in no way to change, the appearance
of the universal Manifestations would be futile. Moreover, in the
'Avalim, which is among the well-known and authoritative books, it is
said: ``A youth will appear from the Banu Hashim, possessing a Book and
new ordinances,'' until it says: ``And most of his enemies will be the
religious scholars.'' Elsewhere it records that Sadiq ibn Muhammad said:
``A youth will indeed appear from the Banu Hashim and call upon the
people to pledge allegiance to him. He possesses a new Book, and will
receive the people's allegiance on the basis of a new Book, severe upon
the Arabs. If you hear anything of him, hasten to him.'' How well they
have followed the counsel of the Imams of the Faith and the lamps of
certitude! Although it says that if you hear that a youth has appeared
from the Banu Hashim, calling people to a new divine Book and wondrous
ordinances from the Lord, you should hasten to him, they nevertheless
all pronounced that Lord of all creation an unbeliever and outside the
faith. They approached that Hashimite Light and divine Manifestation
only with drawn swords and hearts full of hatred. Consider, too, how
explicitly the enmity of the religious scholars is mentioned in the
books. Despite all these manifest and demonstrative traditions and
clear, established allusions, all people have turned away from the pure
essence of knowledge and utterance and have turned toward the
manifestations of error and rebellion. Despite these recorded traditions
and revealed words, they say whatever their own selves desire. If the
Essence of Truth makes a statement contrary to the selfish desires of
this group, they immediately declare Him an unbeliever and say: ``This
contradicts the words of the Imams of the Faith and the manifest lights,
and no such command or ordinance has been issued in the firmly
established law.'' Just so have such useless words proceeded, and
continue to proceed, from these mortal frames today.

\subsection{272}\label{section-462}

\begin{otherlanguage}{persian}

حال این روایت را ملاحظه نمائید که چگونه از قبل جمیع امورات را اخبار
فرموده‌اند در اربعین ذکر فرموده یظهر من بنی هاشم صبیّ ذو احکام جدید فیدعو
النّاس و لم یجبه احد و اکثر اعدائه العلمآء فاذا حکم بشیء لم یطیعوه
فیقولون هذا خلاف ما عندنا من ائمّة الدّین الی آخر الحدیث چنانچه الیوم جمیع
همین کلمات را اعاده مینمایند و شاعر بر این نشده که آن حضرت بر عرش یفعل
ما یشآء جالسند و بر کرسیّ یحکم ما یرید ساکن

\end{otherlanguage}

Ḥāl īn revāyat rā molāḥeẓe namāʾīd ke chegūne az qabl jamīʿ-e omūrāt rā
ekhbār farmūde-and dar Arbaʿīn dhekr farmūde: yaẓharu min Banī Hāshim
ṣabiyy dhū aḥkām jadīd fa-yadʿū al-nās wa-lam yujibhu aḥad wa-akthar
aʿdāʾihi al-ʿulamāʾ fa-idhā ḥakama bi-shayʾ lam yuṭīʿūhu fa-yaqūlūna
hādhā khilāf mā ʿindanā min aʾimmat al-dīn ilā ākhir al-ḥadīth.
Chenānche al-yawm jamīʿ hamīn kalemāt rā eʿāde mīnamāyand va shāʿer bar
īn nashode ke ān ḥaḍrat bar ʿarsh-e yafʿalu mā yashāʾ jālesand va bar
korsī-ye yaḥkumu mā yurīd sāken.

Now consider this tradition, and how all these events were foretold. In
the Arba'in it is stated: ``A youth will appear from the Banu Hashim,
possessing new ordinances. He will summon the people, but no one will
answer him, and most of his enemies will be the religious scholars. When
he decrees a thing, they will not obey him, but will say: `This is
contrary to what we have received from the Imams of the Faith.'\,'' And
so the tradition continues. Today they all repeat these very words,
unaware that He is seated upon the throne of ``He does whatever He
wills'' and established upon the seat of ``He ordains whatever He
desires.''

\subsection{273}\label{section-463}

\begin{otherlanguage}{persian}

و هیچ ادراکی سبقت نیابد بر کیفیّت ظهور او و هیچ عرفانی احاطه ننماید بر
کمّیّت امر او و جمیع قولها بتصدیق او منوط است و تمام امور بامر او محتاج و
ماسوای او بامر او مخلوقند و بحکم او موجود و اوست مظهر اسرار الهی و مبیّن
حکمتهای غیب صمدانی چنانچه در بحارالأنوار و عوالم و در ینبوع از صادق بن
محمّد وارد شده که فرمود العلم سبعة و عشرون حرفاً فجمیع ما جآءت به الرّسل
حرفان و لم یعرف النّاس حتّی الیوم غیر الحرفین فاذا قام قائمنا اخرج الخمسة
و العشرین حرفاً حال ملاحظه فرمائید که علم را بیست و هفت حرف معیّن فرموده و
جمیع انبیا از آدم الی خاتم دو حرف آن را بیان فرموده‌اند و بر این دو حرف
مبعوث شده‌اند و میفرماید قائم ظاهر میفرماید جمیع این بیست و پنج حرف را از
این بیان قدر و رتبهٔ آن حضرت را ملاحظه فرما که قدرش اعظم از کلّ انبیا و
امرش اعلی و ارفع از عرفان و ادراک کلّ اولیاست و امری را که انبیا و اولیا
و اصفیا بآن اطّلاع نیافته و یا بامر مبرم الهی اظهار نداشته این همج رعاع
بعقول و علوم و ادراک ناقص خود میزان میکنند اگر مطابق نیاید رد مینمایند
ام تحسب انّ اکثرهم یسمعون او یعقلون ان هم الّا کالأنعام بل هم اضلّ سبیلا

\end{otherlanguage}

Va hīch edrākī sabqat nayābad bar keyfīyat-e ẓohūr-e ū va hīch ʿerfānī
eḥāṭe nanamāyad bar kammīyat-e amr-e ū va jamīʿ-e qawlhā be-taṣdīq-e ū
manūṭ ast va tamām-e omūr be-amr-e ū moḥtāj va māsavā-ye ū be-amr-e ū
makhlūqand va be-ḥokm-e ū mawjūd va ūst maẓhar-e asrār-e elāhī va
mobayyen-e ḥekmathā-ye ghayb-e ṣamadānī chenānche dar Biḥār al-Anvār va
ʿAvālem va dar Yanbūʿ az Ṣādiq ibn Muḥammad vāred shode ke farmūd:
al-ʿilm sabʿa wa-ʿishrūna ḥarfan fa-jamīʿ mā jāʾat bihi al-rusul ḥarfān
wa-lam yaʿrif al-nās ḥattā al-yawm ghayr al-ḥarfayn fa-idhā qāma
qāʾimunā akhraja al-khamsa wa-l-ʿishrīna ḥarfan. Ḥāl molāḥeẓe farmāʾīd
ke ʿelm rā bīst va haft ḥarf moʿayyan farmūde va jamīʿ-e anbiyā az Ādam
ilā Khātam do ḥarf-e ān rā bayān farmūde-and va bar īn do ḥarf mabʿūth
shode-and va mīfarmāyad Qāʾem ẓāher mīfarmāyad jamīʿ-e īn bīst va panj
ḥarf rā az īn bayān qadr va rotbe-ye ān ḥaḍrat rā molāḥeẓe farmā ke
qadresh aʿẓam az koll-e anbiyā va amrash aʿlā va arfaʿ az ʿerfān va
edrāk-e koll-e awliyāst va amrī rā ke anbiyā va awliyā va aṣfiyā be-ān
eṭṭelāʿ nayāfte va yā be-amr-e mobram-e elāhī eẓhār nadāshte īn hamaj-e
raʿāʿ be-ʿoqūl va ʿolūm va edrāk-e nāqeṣ-e khod mīzān mīkonand agar
moṭābeq nayāyad radd mīnamāyand: am taḥsabu anna aktharahum yasmaʿūna aw
yaʿqilūna in hum illā ka-l-anʿām bal hum aḍallu sabīlan.

No understanding can anticipate the manner of His appearance, and no
knowledge can encompass the scope of His Cause. Every utterance depends
upon His confirmation, and all affairs stand in need of His command. All
besides Him are created by His command and exist by His decree. He is
the Revealer of the divine mysteries and the Expounder of the wisdom of
the unseen, eternal Realm. Thus, in the Bihar al-Anvar, the `Avalim, and
the Yanbu', it is recorded that Sadiq ibn Muhammad said: ``Knowledge
consists of twenty-seven letters. All that the Messengers have brought
is two letters, and to this day the people have known no more than these
two letters. When our Qa'im arises, he will bring forth the remaining
twenty-five letters.'' Now consider: he has specified that knowledge
consists of twenty-seven letters, and all the Prophets, from Adam to the
Seal, have expounded two of its letters and have been sent forth with
these two letters. He says that the Qa'im will reveal all these
twenty-five letters. From this statement, consider His worth and rank:
His worth is greater than that of all the Prophets, and His Cause is
higher and more exalted than the knowledge and understanding of all the
holy ones. Yet these ignorant rabble measure by their own deficient
reason, learning, and understanding a Cause of which the Prophets, the
holy ones, and the chosen ones had no knowledge, or which, by God's
irrevocable command, they did not disclose. If it does not conform to
their measure, they reject it. ``Or do you suppose that most of them
hear or understand? They are only like cattle---indeed, they are further
astray.''

\subsection{274}\label{section-464}

\begin{otherlanguage}{persian}

آیا این حدیث مذکور را بر چه حمل مینمایند که صریح بر ظهور مطالب غیبیّه و
امورات بدیعهٔ جدیده است در ایّام آن حضرت و این امورات بدیعه سبب اختلاف ناس
میشود بقسمی که جمیع علما و فقها حکم بر قتل آن حضرت و اصحاب او کنند و همۀ
اهل ارض بر مخالفت قیام نمایند چنانچه در کافی در حدیث جابر در لوح فاطمه
در وصف قائم میفرماید علیه کمال موسی و بهآء عیسی و صبر ایّوب فیذلّ اولیائه
فی زمانه و تتهادی رؤوسهم کما تتهادی رؤوس التّرک و الدّیلم فیقتلون و یحرقون
و یکونون خائفین مرعوبین وجلین تصبغ الأرض بدمائهم و یفشو الویل و الرّنّة فی
نسائهم اولئک اولیائی حقّا حال ملاحظه فرمائید که حرفی از این حدیث باقی
نماند مگر آنکه ظاهر شد چنانچه در اکثر اماکن دم شریفشان ریخته شد و در هر
بلدی ایشان را اسیر نموده و بولایات و شهرها گردانیدند و بعضی را سوختند و
مع‌ذلک هیچ نفسی فکر ننمود که اگر قائم موعود بشریعت و احکام قبل مبعوث و
ظاهر شود دیگر ذکر این احادیث برای چه شده و چرا این همه اختلاف ظاهر میشود
تا آنکه قتل این اصحاب را واجب دانند و اذیّت این ارواح مقدّسه را سبب وصول
بمعارج قرب شمرند

\end{otherlanguage}

Āyā īn ḥadīth-e madhkūr rā bar che ḥaml mīnamāyand ke ṣarīḥ bar ẓohūr-e
maṭāleb-e ghaybīye va omūrāt-e badīʿe-ye jadīde ast dar ayyām-e ān
ḥaḍrat va īn omūrāt-e badīʿe sabab-e ekhtelāf-e nās mīshavad be-qesmī ke
jamīʿ-e ʿolamā va foqahā ḥokm bar qatl-e ān ḥaḍrat va aṣḥāb-e ū konand
va hame-ye ahl-e arḍ bar mokhālefat qīyām namāyand chenānche dar Kāfī
dar ḥadīth-e Jāber dar Lawḥ-e Fāṭeme dar vaṣf-e Qāʾem mīfarmāyad:
ʿalayhi kamāl Mūsā wa-bahāʾ ʿĪsā wa-ṣabr Ayyūb fa-yudhallu awliyāʾihi fī
zamānihi wa-tatahādā ruʾūsuhum kamā tatahādā ruʾūs al-Turk wa-l-Daylam
fa-yuqtalūna wa-yuḥraqūna wa-yakūnūna khāʾifīna marʿūbīna wajilīn
tuṣbaghu al-arḍ bi-dimāʾihim wa-yafshū al-wayl wa-l-ranna fī nisāʾihim
ulāʾika awliyāʾī ḥaqqan. Ḥāl molāḥeẓe farmāʾīd ke ḥarfī az īn ḥadīth
bāqī namānd magar ānke ẓāher shod chenānche dar akthar-e amāken dam-e
sharīfeshān rīkhte shod va dar har baladī īshān rā asīr namūde va
be-velāyāt va shahrhā gardānīdand va baʿḍī rā sūkhtand va maʿa-dhālek
hīch nafsī fekr nanamūd ke agar Qāʾem-e mawʿūd be-sharīʿat va aḥkām-e
qabl mabʿūth va ẓāher shavad dīgar dhekr-e īn aḥādīth barā-ye che shode
va cherā īn hame ekhtelāf ẓāher mīshavad tā ānke qatl-e īn aṣḥāb rā
vājeb dānand va adhīyat-e īn arvāḥ-e moqaddase rā sabab-e voṣūl
be-maʿārej-e qorb shomarand.

How do they interpret this tradition, which explicitly foretells the
revelation of hidden matters and wondrous new things in His days? These
wondrous things cause such discord among the people that all the
religious scholars and jurists decree the death of Him and His
companions, and all the people of the earth rise in opposition. Thus, in
the Kafi, in the tradition of Jabir concerning the Tablet of Fatimih, it
is said in description of the Qa'im: ``Upon him rests the perfection of
Moses, the splendor of Jesus, and the patience of Job. His friends will
be abased in his time, and their heads will be exchanged as gifts, just
as the heads of the Turks and the Daylamites are exchanged as gifts.
They will be killed and burned, and will live in fear, terror, and
dread. The earth will be stained with their blood, and wailing and
lamentation will spread among their women. These are truly My friends.''
Now consider how not a word of this tradition has remained unfulfilled.
In most places their noble blood was shed; in every town they were taken
captive and paraded through provinces and cities, and some were burned.
Yet no one reflected that, if the promised Qa'im were to be sent forth
and appear with the former law and ordinances, why should these
traditions have been related? And why should all this discord arise, to
the point that people deem it obligatory to kill these companions and
regard tormenting these holy souls as a means of ascending to the
heights of nearness to God?

\subsection{275}\label{section-465}

\begin{otherlanguage}{persian}

و دیگر ملاحظه فرمائید چگونه جمیع این امور وارده و افعال نازله در احادیث
قبل ذکر شده چنانچه در روضهٔ کافی در بیان زوراء میفرماید و فی روضة‌الکافی
عن معاویة بن وهب عن ابی‌عبداللّه قال أ تعرف الزّورآء قلت جعلت فداک یقولون
انّها بغداد قال لا ثمّ قال دخلت الرّی قلت نعم قال اتیت سوق الدّوابّ قلت نعم
قال رأیت جبل الأسود عن یمین الطّریق تلک الزّورآء یقتل فیها ثمانون رجلاً من
ولد فلان کلّهم یصلح الخلافة قلت من یقتلهم قال یقتلهم اولاد العجم

\end{otherlanguage}

Va dīgar molāḥeẓe farmāʾīd chegūne jamīʿ-e īn omūr-e vārede va afʿāl-e
nāzele dar aḥādīth-e qabl dhekr shode chenānche dar Rawḍe-ye Kāfī dar
bayān-e Zawrā mīfarmāyad: wa-fī Rawḍat al-Kāfī ʿan Muʿāwiya ibn Wahb ʿan
Abī ʿAbd Allāh qāla a-taʿrifu al-Zawrāʾ qultu juʿiltu fidāka yaqūlūna
innahā Baghdād qāla lā thumma qāla dakhalta al-Rayy qultu naʿam qāla
atayta sūq al-dawābb qultu naʿam qāla raʾayta jabal al-aswad ʿan yamīn
al-ṭarīq tilka al-Zawrāʾ yuqtalu fīhā thamānūna rajulan min wuld fulān
kulluhum yaṣluḥu al-khilāfa qultu man yaqtuluhum qāla yaqtuluhum awlād
al-ʿAjam.

Consider, moreover, how all these events that have occurred and deeds
that have befallen them were mentioned in the traditions of the past.
Thus, in the Rawdih of the Kafi, concerning Zawra, it is said: ``And in
the Rawdat al-Kafi, Mu'awiyah ibn Wahb relates from Abu 'Abd Allah: He
said, `Do you know Zawra?' I replied, `May I be your sacrifice! They say
it is Baghdad.' He said, `No.' Then he said, `Have you entered Rayy?' I
replied, `Yes.' He said, `Have you been to the livestock market?' I
replied, `Yes.' He said, `Have you seen the black mountain to the right
of the road? That is Zawra. Eighty men from the descendants of such a
one will be killed there, each of them fit for the caliphate.' I asked,
`Who will kill them?' He said, `The sons of the Persians will kill
them.'\,''

\subsection{276}\label{section-466}

\begin{otherlanguage}{persian}

اینست حکم و امر اصحاب آن حضرت که از قبل بیان فرموده‌اند و حال ملاحظه
فرمائید که زوراء موافق این روایت ارض ری است و این اصحاب را در آن مکان
ببدترین عذاب بقتل رساندند و جمیع این وجوهات قدسی را عجم شهید نموده
چنانچه در حدیث مذکور است و شنیده‌اند و بر همهٔ عالم واضح و مبرهنست حال چرا
این خراطین ارض در این احادیث که جمیع آن بمثل شمس در وسط سماء ظاهر شد
تفکّر نمینمایند و اقبال بحقّ نمیجویند و ببعضی احادیث که معنی آن را ادراک
ننموده‌اند از ظهور حقّ و جمال اللّه اعراض جسته‌اند و بسقر مقرّ گزیده‌اند نیست
این امور مگر از اعراض فقهای عصر و علمای عهد اینست که صادق بن محمّد
میفرماید فقهآء ذلک الزّمان شرّ فقهآء تحت ظلّ السّمآء منهم خرجت الفتنة و
الیهم تعود

\end{otherlanguage}

Īnast ḥokm va amr-e aṣḥāb-e ān Ḥaḍrat ke az qabl bayān farmūde-and va
ḥāl molāḥeẓe farmāʾīd ke Zawrā movāfeq-e īn revāyat arḍ-e Ray ast va īn
aṣḥāb rā dar ān makān be badtarīn ʿadhāb be qatl resāndand va jamīʿ-e īn
vojūhāt-e qodsī rā ʿAjam shahīd namūde chenānche dar ḥadīth madhkūr ast
va shenīde-and va bar hame-ye ʿālam vāḍeḥ va mobarhan ast ḥāl cherā īn
kharāṭīn-e arḍ dar īn aḥādīth ke jamīʿ-e ān be methl-e shams dar vasaṭ-e
samāʾ ẓāher shod tafakkor nemīnamāyand va eqbāl be Ḥaqq nemījūyand va be
baʿḍī aḥādīth ke maʿnī-ye ān rā edrāk nanamūde-and az ẓohūr-e Ḥaqq va
Jamāl-e Allāh eʿrāḍ joste-and va be Saqar maqarr gozīde-and nīst īn omūr
magar az eʿrāḍ-e foqahā-ye ʿaṣr va ʿolamā-ye ʿahd īnast ke Ṣādiq bin
Moḥammad mīfarmāyad fuqahāʾu dhālika al-zamān sharru fuqahāʾa taḥta
ẓilli al-samāʾ minhum kharajati al-fitnatu wa-ilayhim taʿūd.

Such is the fate and condition of that exalted One's companions,
foretold in the past. Now consider that, according to this tradition,
Zawra is the land of Rayy. These companions were put to death there with
the cruelest torments, and all these holy souls were martyred by the
Persians, as stated in the tradition, as people have heard, and as is
clear and evident to the whole world. Why, then, do these worms of the
earth not reflect on these traditions, all of which have become manifest
like the sun in the midst of heaven, and turn toward the Truth? Instead,
clinging to certain traditions whose meaning they have not understood,
they have turned away from the manifestation of the Truth and the Beauty
of God and chosen their abode in hell. These things result only from the
aversion of the jurists of this age and the religious scholars of this
time. Thus Ṣādiq ibn Muḥammad says: ``The jurists of that time are the
worst jurists beneath the shadow of heaven. From them has arisen
discord, and to them it will return.''

\subsection{277}\label{section-467}

\begin{otherlanguage}{persian}

و از فقها و علمای بیان استدعا مینمایم که چنین مشی ننمایند و بر جوهر الهی
و نور ربّانی و صرف ازلی و مبدء و منتهای مظاهر غیبی در زمن مستغاث وارد
نیاورند آنچه در این کور وارد شد و بعقول و ادراک و علم متمسّک نشوند و بآن
مظهر علوم نامتناهی ربّانی مخاصمه ننمایند اگرچه با جمیع این وصایا دیده
میشود که شخصی اعور که از رؤسای قوم است در نهایت معارضه برخیزد و همچنین
در هر بلدی بر نفی آن جمال قدسی برخیزند و اصحاب آن سلطان وجود و جوهر
مقصود در کوهها و صحراها فرار نمایند و از دست ظالمین مستور شوند و برخی
توکّل نمایند و با کمال انقطاع جان دربازند و گویا مشاهده میشود نفسی که
بکمال زهد و تقوی موصوف و معروفست بقسمی که جمیع ناس اطاعت او را فرض شمرند
و تسلیم امرش را لازم دانند بمحاربه با آن اصل شجرهٔ الهیّه قیام نماید و
بمنتهای جهد و اجتهاد بمعارضه برخیزد اینست شأن ناس

\end{otherlanguage}

Va az foqahā va ʿolamā-ye Bayān estedʿā mīnamāyam ke chenīn mashy
nanamāyand va bar Jawhar-e Elāhī va Nūr-e Rabbānī va Ṣerf-e Azalī va
Mabdaʾ va Montahā-ye maẓāher-e ghaybī dar zaman-e Mostaghāth vāred
nayāvarand ānche dar īn kowr vāred shod va be ʿoqūl va edrāk va ʿelm
motamassek nashavand va be ān Maẓhar-e ʿolūm-e nāmotanāhī-ye rabbānī
mokhāṣame nanamāyand agarche bā jamīʿ-e īn vaṣāyā dīde mīshavad ke
shakhṣī aʿvar ke az roʾasā-ye qawm ast dar nehāyat-e moʿāreḍe barkhīzad
va hamchenīn dar har baladī bar nafy-e ān Jamāl-e Qodsī barkhīzand va
aṣḥāb-e ān Solṭān-e Vojūd va Jawhar-e Maqṣūd dar kūhhā va ṣaḥrāhā farār
namāyand va az dast-e ẓālemīn mastūr shavand va barkhī tavakkol namāyand
va bā kamāl-e enqeṭāʿ jān darbāzand va gūyā moshāhede mīshavad nafsī ke
be kamāl-e zohd va taqvā mowṣūf va maʿrūf ast be qesmī ke jamīʿ-e nās
eṭāʿat-e ū rā farḍ shomarand va taslīm-e amrash rā lāzem dānand be
moḥārebe bā ān Aṣl-e Shajare-ye Elāhiyye qiyām namāyad va be montahā-ye
jahd va ejtehād be moʿāreḍe barkhīzad īnast shaʾn-e nās.

I implore the jurists and religious scholars of the Bayān not to act in
this way, and not to inflict upon the divine Essence, the heavenly
Light, the eternal Being, the Beginning and End of the unseen
Manifestations, in the time of Mustaghāth, what was inflicted in this
cycle. Let them not cling to their reason, understanding, and learning,
or dispute with that Manifestation of infinite divine knowledge. Yet,
despite all these admonitions, it is seen that a one-eyed man, one of
the leaders of the people, will arise in the fiercest opposition.
Likewise, in every city they will arise to deny that holy Beauty, and
the companions of that Sovereign of existence and Essence of all desire
will flee into the mountains and deserts and conceal themselves from the
oppressors, while some will put their trust in God and, with complete
detachment, lay down their lives. It is as though one can already see a
man so renowned and celebrated for his perfect asceticism and piety that
all people consider obedience to him obligatory and submission to his
command necessary, arising to wage war against that Root of the divine
Tree and opposing Him with the utmost effort and exertion. Such is the
way of the people.

\subsection{278}\label{section-468}

\begin{otherlanguage}{persian}

باری امیدواریم که اهل بیان تربیت شوند و در هوای روح طیران نمایند و در
فضای روح ساکن شوند حقّ را از غیر تمیز دهند و تلبیس باطل را بدیدهٔ بصیرت
بشناسند اگرچه در این ایّام رائحهٔ حسدی وزیده که قسم بمربّی وجود از غیب و
شهود که از اوّل بنای وجود عالم با اینکه آن را اوّلی نه تا حال چنین غل و
حسد و بغضائی ظاهر نشده و نخواهد شد چنانچه جمعی که رائحهٔ انصاف را
نشنیده‌اند رایات نفاق برافراخته‌اند و بر مخالفت این عبد اتّفاق نموده‌اند و
از هر جهت رمحی آشکار و از هر سمت تیری طیّار با اینکه باحدی در امری افتخار
ننمودم و بنفسی برتری نجستم مع هر نفسی مصاحبی بودم در نهایت مهربان و
رفیقی بغایت بردبار و رایگان با فقرا مثل فقرا بودم و با علما و عظما در
کمال تسلیم و رضا مع‌ذلک فواللّه الّذی لا اله الّا هو با آن همه ابتلا و بأساء
و ضرّاء که از اعدا و اولی الکتاب وارد شد نزد آنچه از احبّا وارد شد معدوم
صرف است و مفقود بحت

\end{otherlanguage}

Bārī omīdvārīm ke ahl-e Bayān tarbiyat shavand va dar havā-ye rūḥ
ṭayarān namāyand va dar faḍā-ye rūḥ sāken shavand Ḥaqq rā az ghayr
tamyīz dehand va talbīs-e bāṭel rā be dīde-ye baṣīrat beshenāsand
agarche dar īn ayyām rāʾeḥe-ye ḥasadī vazīde ke qasam be Morabbī-ye
vojūd az ghayb va shohūd ke az avval-e banā-ye vojūd-e ʿālam bā īnke ān
rā avvalī na tā ḥāl chenīn ghell va ḥasad va baghḍāʾī ẓāher nashode va
nakhāhad shod chenānche jamʿī ke rāʾeḥe-ye enṣāf rā nashenīde-and
rāyāt-e nefāq bar afrākhte-and va bar mokhālefat-e īn ʿabd ettefāq
namūde-and va az har jahat romḥī āshkār va az har samt tīrī ṭayyār bā
īnke be aḥadī dar amrī eftekhār nanamūdam va be nafsī bartarī najostam
maʿ-e har nafsī moṣāḥebī būdam dar nehāyat mehrabān va rafīqī be ghāyat
bordbār va rāyegān bā foqarā methl-e foqarā būdam va bā ʿolamā va ʿoẓamā
dar kamāl-e taslīm va reḍā maʿa-dhālek fa-wa-Allāhi alladhī lā ilāha
illā huwa bā ān hame ebtelā va baʾsāʾ va ḍarrāʾ ke az aʿdā va ūlī
al-kitāb vāred shod nazd-e ānche az aḥebbā vāred shod maʿdūm-e ṣerf ast
va mafqūd-e baḥt.

We hope, then, that the people of the Bayān will be educated, soar in
the atmosphere of the spirit, and dwell in its expanse; that they will
distinguish the Truth from all else and recognize the disguises of
falsehood with the eye of insight. Yet in these days a breath of envy
has blown such that, I swear by the One who nurtures all existence,
unseen and seen, from the beginning of the world's existence---though it
has no beginning---until now, such malice, envy, and hatred have never
appeared, nor will they ever appear again. Thus a group who have never
caught the fragrance of fairness have raised the banners of hypocrisy
and united in opposition to this servant. From every direction a spear
is raised, and from every side an arrow flies, although I have never
boasted of anything before anyone or sought superiority over any soul.
To every person I was a most loving companion and an exceedingly
forbearing and generous friend. Among the poor I was as one of the poor,
and before the learned and the eminent I showed complete submission and
acquiescence. Nevertheless, by God, besides whom there is no god, all
the trials, hardships, and afflictions inflicted by enemies and the
people of the Book are sheer nothingness and utter nonexistence beside
what has been inflicted by the friends.

\subsection{279}\label{section-469}

\begin{otherlanguage}{persian}

باری چه اظهار نمایم که امکان را اگر انصاف باشد طاقت این بیان نه و این
عبد در اوّل ورود این ارض چون فی‌الجمله بر امورات محدثهٔ بعد اطّلاع یافتم از
قبل مهاجرت اختیار نمودم و سر در بیابانهای فراق نهادم و دو سال وحده در
صحراهای هجر بسر بردم و از عیونم عیون جاری بود و از قلبم بحور دم ظاهر چه
لیالی که قوت دست نداد و چه ایّام که جسد راحت نیافت و با این بلایای نازله
و رزایای متواتره فوالّذی نفسی بیده کمال سرور موجود بود و نهایت فرح مشهود
زیرا که از ضرر و نفع و صحّت و سقم نفسی اطّلاع نبود بخود مشغول بودم و از
ماسوی غافل و غافل از اینکه کمند قضای الهی اوسع از خیال است و تیر تقدیر
او مقدّس از تدبیر سر را از کمندش نجات نه و اراده‌اش را جز رضا چاره‌ئی نه
قسم بخدا که این مهاجرتم را خیال مراجعت نبود و مسافرتم را امید مواصلت نه
و مقصود جز این نبود که محلّ اختلاف احباب نشوم و مصدر انقلاب اصحاب نگردم و
سبب ضرّ احدی نشوم و علّت حزن قلبی نگردم غیر از آنچه ذکر شد خیالی نبود و
امری منظور نه اگرچه هر نفسی محملی بست و بهوای خود خیالی نمود باری تا
آنکه از مصدر امر حکم رجوع صادر شد و لابدّاً تسلیم نمودم و راجع شدم

\end{otherlanguage}

Bārī che eẓhār namāyam ke emkān rā agar enṣāf bāshad ṭāqat-e īn bayān na
va īn ʿabd dar avval-e vorūd-e īn arḍ chon fī al-jomle bar omūrāt-e
moḥdathe-ye baʿd eṭṭelāʿ yāftam az qabl mohājerat ekhtiyār namūdam va
sar dar bīyābānhā-ye ferāq nahādam va do sāl vaḥdah dar ṣaḥrāhā-ye hajr
be sar bordam va az ʿoyūnam ʿoyūn jārī būd va az qalbam boḥūr-e dam
ẓāher che layālī ke qūt dast nadād va che ayyām ke jasad rāḥat nayāft va
bā īn balāyā-ye nāzele va razāyā-ye motavātere fa-wa-alladhī nafsī
bi-yadihi kamāl-e sorūr mowjūd būd va nehāyat-e faraḥ mashhūd zīrā ke az
ḍarar va nafʿ va ṣeḥḥat va saqam-e nafsī eṭṭelāʿ nabūd be khod mashghūl
būdam va az mā-sevā ghāfel va ghāfel az īnke kamand-e qaḍā-ye elāhī
awsaʿ az khayāl ast va tīr-e taqdīr-e ū moqaddas az tadbīr sar rā az
kamandash nejāt na va erāde-ash rā joz reḍā chāre-ī na qasam be Khodā ke
īn mohājeratam rā khayāl-e morājaʿat nabūd va mosāferatam rā omīd-e
movāṣalat na va maqṣūd joz īn nabūd ke maḥall-e ekhtelāf-e aḥbāb
nashavam va maṣdar-e enqelāb-e aṣḥāb nagardam va sabab-e ḍorr-e aḥadī
nashavam va ʿellat-e ḥozn-e qalbī nagardam ghayr az ānche dhekr shod
khayālī nabūd va amrī manẓūr na agarche har nafsī maḥmalī bast va be
havā-ye khod khayālī namūd bārī tā ānke az Maṣdar-e Amr ḥokm-e rojūʿ
ṣāder shod va lā-boddan taslīm namūdam va rājeʿ shodam.

But what can I say? Were the world of existence fair, it could not bear
this account. When this servant first arrived in this land and gained
some knowledge of the events that would later arise, I chose to depart
before they occurred and withdrew into the wilderness of separation. For
two years I lived alone in the deserts of estrangement. Springs flowed
from My eyes, and seas of blood poured from My heart. How many nights I
had no food, and how many days My body found no rest! Yet amid these
descending afflictions and successive calamities, by Him in whose hand
is My soul, perfect joy was present and the utmost gladness evident, for
I knew nothing of anyone's harm or benefit, health or illness. I was
occupied with Myself and oblivious to all else, unaware that the snare
of divine decree is wider than imagination and the arrow of His destiny
beyond all human devising. There is no escape from His snare, and no
recourse before His will but acquiescence. I swear by God that in this
departure I had no thought of returning, and in this journey no hope of
reunion. My sole purpose was to avoid becoming a cause of disagreement
among the loved ones, a source of turmoil among the companions, a cause
of harm to anyone, or a reason for sorrow to any heart. Beyond what I
have mentioned, I had no other thought or purpose, although everyone
devised an explanation and imagined something according to their own
desires. At last, the command to return issued from the Source of
authority, and I had no choice but to submit and return.

\subsection{280}\label{section-470}

\begin{otherlanguage}{persian}

دیگر قلم عاجز است از ذکر آنچه بعد از رجوع ملاحظه شد حال دو سنه میگذرد که
اعدا در اهلاک این عبد فانی بنهایت سعی و اهتمام دارند چنانچه جمیع مطّلع
شده‌اند مع‌ذلک نفسی از احباب نصرت ننموده و بهیچوجه اعانتی منظور نداشته
بلکه از عوض نصر حزنها که متوالی و متواتر قولاً و فعلاً مثل غیث هاطل وارد
میشود و این عبد در کمال رضا جان بر کف حاضرم که شاید از عنایت الهی و فضل
سبحانی این حرف مذکور مشهور در سبیل نقطه و کلمهٔ علیا فدا شود و جان دربازد
و اگر این خیال نبود فوالّذی نطق الرّوح بأمره آنی در این بلد توقّف نمینمودم
و کفی باللّه شهیداً اختم القول بلا حول و لا قوّة الّا باللّه و انّا للّه و انّا
الیه راجعون

\end{otherlanguage}

Dīgar qalam ʿājez ast az dhekr-e ānche baʿd az rojūʿ molāḥeẓe shod ḥāl
do sane mīgodharad ke aʿdā dar ehlāk-e īn ʿabd-e fānī be nehāyat-e saʿy
va ehtemām dārand chenānche jamīʿ moṭṭaleʿ shode-and maʿa-dhālek nafsī
az aḥbāb noṣrat nanamūde va be hīch vajh eʿānatī manẓūr nadāshte balke
az ʿavaḍ-e naṣr ḥoznhā ke motavālī va motavāter qawlan va feʿlan methl-e
ghayth-e hāṭel vāred mīshavad va īn ʿabd dar kamāl-e reḍā jān bar kaf
ḥāḍeram ke shāyad az ʿenāyat-e elāhī va faḍl-e sobḥānī īn ḥarf-e
madhkūr-e mashhūr dar sabīl-e Noqṭe va Kaleme-ye ʿOlyā fadā shavad va
jān darbāzad va agar īn khayāl nabūd fa-wa-alladhī naṭaqa al-Rūḥu
bi-amrihi ānī dar īn balad tavaqqof nemīnamūdam va kafā bi-Allāhi
shahīdan akhtimu al-qawla bi-lā ḥawla wa-lā quwwata illā bi-Allāh
wa-innā li-Allāhi wa-innā ilayhi rājiʿūn.

The pen is powerless to recount what was witnessed after My return. For
two years now, the enemies have exerted the utmost effort and diligence
to destroy this mortal servant, as everyone has come to know. Yet not
one of the friends has offered support or shown any thought of
assistance. Instead of help, sorrows have rained down continually and in
succession, in word and deed, like a torrential downpour. With complete
acquiescence, this servant stands ready, life in hand, so that perhaps,
through divine loving-kindness and heavenly grace, this widely mentioned
letter may be sacrificed and lay down its life in the path of the Point
and the Most Exalted Word. Were it not for this thought, by Him at whose
command the Spirit has spoken, I would not remain in this city for a
moment. God is sufficient as a witness. I close with these words:
``There is no power or strength except through God. We belong to God,
and to Him we return.''

\subsection{281}\label{section-471}

\begin{otherlanguage}{persian}

صاحبان هوش که از صهبای حبّ نوشیده‌اند و گامی بکام نفس برنداشته‌اند دلائل و
برهان و حجّت را که جمیع مشعر بر این امر بدیع و ظهور منیع الهی است اظهر از
شمس در فلک چهارم مشاهده نمایند حال اعراض خلق را از جمال الهی و اقبالشان
را بهوای نفسانی ملاحظه فرمائید با همهٔ این آیات متقنه و اشارات محکمه که
در ثقل اکبر که ودیعهٔ ربّانیّه است در بین عباد و این احادیث واضحه که اصرح
از بیان و تبیانست از همه غافل و معرض شده‌اند و چند حدیث که بادراک خود
مطابق نیافته‌اند و معنی آن را ادراک ننموده‌اند متمسّک بظاهر آنها شده و از
سلسال خمر ذی الجلال و زلال بی‌زوال جمال لا‌یزال محروم و مأیوس مانده‌اند

\end{otherlanguage}

Ṣāḥebān-e hūsh ke az ṣahbā-ye ḥobb nūshīde-and va gāmī be kām-e nafs
barnadāshte-and dalāʾel va borhān va ḥojjat rā ke jamīʿ moshʿer bar īn
Amr-e badīʿ va Ẓohūr-e manīʿ-e elāhī ast aẓhar az shams dar falak-e
chahārom moshāhede namāyand ḥāl eʿrāḍ-e khalq rā az Jamāl-e Elāhī va
eqbāles hān rā be havā-ye nafsānī molāḥeẓe farmāʾīd bā hame-ye īn āyāt-e
motqane va eshārāt-e moḥkame ke dar Theql-e Akbar ke vadīʿe-ye
rabbāniyye ast dar bayn-e ʿebād va īn aḥādīth-e vāḍeḥe ke aṣraḥ az bayān
va tebyān ast az hame ghāfel va moʿreḍ shode-and va chand ḥadīth ke bā
edrāk-e khod moṭābeq nayāfte-and va maʿnī-ye ān rā edrāk nanamūde-and
motamassek be ẓāher-e ānhā shode va az salsāl-e khamr-e Dhi al-Jalāl va
zolāl-e bī-zavāl-e Jamāl-e lā-yazāl maḥrūm va maʾyūs mānde-and.

Those endowed with discernment, who have drunk the wine of love and have
not taken a single step to gratify the self, see the evidences, proofs,
and testimonies, all of which point to this wondrous Cause and exalted
divine Revelation, as clearer than the sun in the fourth heaven. Now
consider how people have turned away from the divine Beauty and toward
their own selfish desires. Despite all these compelling verses and firm
indications in the Greater Weight, the divine trust among His servants,
and these clear traditions, more explicit than any explanation or
exposition, they have ignored and turned away from them all. Clinging
instead to the outward sense of a few traditions that they have not
found to accord with their own understanding and whose meaning they have
not grasped, they have remained deprived of, and without hope of
attaining, the flowing wine of the Lord of Majesty and the unfailing
pure waters of the everlasting Beauty.

\subsection{282}\label{section-472}

\begin{otherlanguage}{persian}

ملاحظه فرمائید که در اخبار سنهٔ ظهور آن هویّهٔ نور را هم ذکر فرموده‌اند
مع‌ذلک شاعر نشده‌اند و در نفسی از هوای نفس منقطع نگشته‌اند فی حدیث المفضّل
سأل عن الصّادق فکیف یا مولای فی ظهوره فقال (ع) فی سنة السّتّین یظهر امره و
یعلو ذکره

\end{otherlanguage}

Molāḥeẓe farmāʾīd ke dar akhbār sane-ye ẓohūr-e ān Hoviyye-ye Nūr rā ham
dhekr farmūde-and maʿa-dhālek shāʿer nashode-and va dar nafasī az
havā-ye nafs monqaṭeʿ nagashte-and fī ḥadīthi al-Mufaḍḍal saʾala ʿan
al-Ṣādiq fa-kayfa yā mawlāya fī ẓuhūrihi fa-qāla (ʿalayhi al-salām) fī
sanati al-sittīn yaẓharu amruhu wa-yaʿlū dhikruhu.

Consider that even the year of the appearance of that Being of light has
been mentioned in the traditions. Yet they have not perceived it, nor
have they detached themselves from selfish desire for even a moment. In
the tradition of Mufaḍḍal, he asked Ṣādiq: ``What, my master, of his
appearance?'' He replied, peace be upon him: ``In the year sixty his
Cause will appear and his name will be exalted.''

\subsection{283}\label{section-473}

\begin{otherlanguage}{persian}

باری تحیّر است از این عباد که چگونه با این اشارات واضحهٔ لائحه از حقّ
احتراز نموده‌اند مثلاً ذکر حزن و سجن و ابتلا که بر آن خلاصهٔ فطرت الهی وارد
شد در اخبار قبل ذکر شده فی البحار انّ فی قائمنا اربع علامات من اربعة نبیّ
موسی و عیسی و یوسف و محمّد امّا العلامة من موسی الخوف و الانتظار و امّا
العلامة من عیسی ما قالوا فی حقّه و العلامة من یوسف السّجن و التّقیّة و
العلامة من محمّد یظهر بآثار مثل قرآن با این حدیث باین محکمی که جمیع
امورات را مطابق آنچه واقع شده ذکر فرموده‌اند مع‌ذلک احدی متنبّه نشده و گمان
ندارم که بعد هم متنبّه شوند الّا من شآء ربّک انّ اللّه مسمع من یشآء و ما انا
بمسمع من فی القبور

\end{otherlanguage}

Bārī taḥayyor ast az īn ʿebād ke chegūne bā īn eshārāt-e vāḍeḥe-ye
lāʾeḥe az Ḥaqq eḥterāz namūde-and mathalan dhekr-e ḥozn va sejn va
ebtelā ke bar ān Kholāṣe-ye Feṭrat-e Elāhī vāred shod dar akhbār-e qabl
dhekr shode fī al-Biḥār inna fī qāʾiminā arbaʿa ʿalāmātin min arbaʿati
nabiyyin Mūsā wa-ʿĪsā wa-Yūsuf wa-Muḥammad ammā al-ʿalāmatu min Mūsā
al-khawfu wa-al-intiẓār wa-ammā al-ʿalāmatu min ʿĪsā mā qālū fī ḥaqqihi
wa-al-ʿalāmatu min Yūsuf al-sijnu wa-al-taqiyyatu wa-al-ʿalāmatu min
Muḥammad yaẓharu bi-āthārin mithli Qurʾān bā īn ḥadīth be īn moḥkamī ke
jamīʿ-e omūrāt rā moṭābeq-e ānche vāqeʿ shode dhekr farmūde-and
maʿa-dhālek aḥadī motanabbeh nashode va gomān nadāram ke baʿd ham
motanabbeh shavand illā man shāʾa rabbuka inna Allāha musmiʿun man
yashāʾu wa-mā anā bi-musmiʿin man fī al-qubūr.

One can only marvel at these people: how, despite these clear and
evident indications, have they shunned the Truth? For example, the
sorrow, imprisonment, and afflictions that befell that Quintessence of
the divine creation were mentioned in the traditions of the past. In the
Biḥār it is stated: ``In our Qāʾim are four signs from four prophets:
Moses, Jesus, Joseph, and Muḥammad. The sign from Moses is fear and
expectation; the sign from Jesus is what they said concerning him; the
sign from Joseph is imprisonment and concealment of his faith; and the
sign from Muḥammad is that he will appear with writings like the
Qurʾān.'' Despite this firmly established tradition, which mentions all
these events just as they occurred, no one has taken heed, nor do I
suppose they will take heed hereafter, ``except those whom your Lord
wills.'' ``God indeed causes whomever He wills to hear, but I cannot
cause those in the graves to hear.''

\subsection{284}\label{section-474}

\begin{otherlanguage}{persian}

و بر آن جناب معلوم بوده که اطیار هویّه و حمامات ازلیّه را دو بیانست بیانی
بر حسب ظاهر بی ‌رمز و نقاب و حجاب فرموده و میفرمایند تا سراجی باشد
هدایت‌کننده و نوری راه‌نماینده تا سالکین را بمعارج قدس رساند و طالبین را
ببساط انس کشاند چنانچه مذکور شد از روایات مکشوفه و آیات واضحه و بیاناتی
با حجاب و ستر فرموده و میفرمایند تا مغلّین آنچه در قلب پنهان نموده‌اند
ظاهر شود و حقایقشان باهر گردد اینست که صادق بن محمّد میفرماید و اللّه
لیمحّصنّ و اللّه لیغربلنّ اینست میزان الهی و محکّ صمدانی که عباد خود را بآن
امتحان میفرماید و احدی پی بمعانی این بیانات نبرد مگر قلوب مطمئنّه و نفوس
مرضیّه و افئدهٔ مجرّده و مقصود در امثال اینگونه بیانات معانی ظاهریّه که مردم
ادراک مینمایند نبوده و نیست اینست که میفرماید لکلّ علم سبعون وجهاً و لیس
بین النّاس الّا واحد و اذا قام القائم یبثّ باقی الوجوه بین النّاس و ایضاً قال
نحن نتکلّم بکلمة و نرید منها احدی و سبعین وجهاً و لنا لکلّ منها المخرج

\end{otherlanguage}

Va bar ān janāb maʿlūm būde ke aṭyār-e hoviyye va ḥamāmāt-e azaliyye rā
do bayān ast bayānī bar ḥasab-e ẓāher bī ramz va neqāb va ḥejāb farmūde
va mīfarmāyand tā serājī bāshad hedāyat-konande va nūrī rāh-namāyande tā
sālekīn rā be maʿārej-e qods resānad va ṭālebīn rā be besāṭ-e ons
keshānad chenānche madhkūr shod az revāyāt-e makshūfe va āyāt-e vāḍeḥe
va bayānātī bā ḥejāb va setr farmūde va mīfarmāyand tā moghellīn ānche
dar qalb penhān namūde-and ẓāher shavad va ḥaqāyeqeshān bāher gardad
īnast ke Ṣādiq bin Moḥammad mīfarmāyad wa-Allāhi la-yumaḥḥaṣanna
wa-Allāhi la-yugharbalanna īnast mīzān-e elāhī va meḥakk-e ṣamadānī ke
ʿebād-e khod rā be ān emteḥān mīfarmāyad va aḥadī pey be maʿānī-ye īn
bayānāt nabarad magar qolūb-e moṭmaʾenne va nofūs-e marḍiyye va
afʾede-ye mojarrade va maqṣūd dar amthāl-e īngūne bayānāt maʿānī-ye
ẓāheriyye ke mardom edrāk mīnamāyand nabūde va nīst īnast ke mīfarmāyad
li-kulli ʿilmin sabʿūna wajhan wa-laysa bayna al-nāsi illā wāḥid wa-idhā
qāma al-Qāʾimu yabuththu bāqiya al-wujūhi bayna al-nās wa-ayḍan qāla
naḥnu natakallamu bi-kalimatin wa-nurīdu minhā iḥdā wa-sabʿīna wajhan
wa-lanā li-kullin minhā al-makhraj.

It is known to you that the birds of divine Being and the doves of
eternity have two kinds of utterance. One is outwardly expressed,
without allusion, veil, or concealment. They have spoken and continue to
speak in this way so that it may be a guiding lamp and a light showing
the way, leading wayfarers to the heights of holiness and drawing
seekers into the court of intimacy, as has been mentioned concerning the
explicit traditions and clear verses. They have also spoken and continue
to speak in veiled and concealed utterances, so that what the malicious
have hidden in their hearts may be revealed and their true natures made
evident. Thus Ṣādiq ibn Muḥammad says: ``By God, they will surely be
tested! By God, they will surely be sifted!'' This is the divine balance
and heavenly touchstone by which He tests His servants. None can grasp
the meanings of these utterances except assured hearts, souls pleasing
to God, and detached minds. In utterances of this kind, the intended
meanings have never been, and are not, the outward meanings that people
understand. Thus he says: ``Every knowledge has seventy aspects, and
only one is known among the people. When the Qāʾim arises, he will
spread the remaining aspects among the people.'' He also said: ``We
speak a word and intend by it seventy-one meanings, and for each of them
we have an explanation.''

\subsection{285}\label{section-475}

\begin{otherlanguage}{persian}

باری ذکر این مراتب برای آنست که از بعضی روایات و بیانات که در عالم ملک
آثار آن ظاهر نشده مضطرب نشوند و حمل بر عدم ادراک خود نمایند نه بر عدم
ظهور معانی حدیث زیرا که نزد آن عباد معلوم نیست که مقصود ائمّهٔ دین چه بود
چنانچه از حدیث مستفاد میشود پس باید عباد باینگونه عبارات خود را از
فیوضات ممنوع نسازند و از اهلش سؤال نمایند تا اسرار مستوره بلا حجاب ظاهر
و واضح شود

\end{otherlanguage}

Bārī, zekr-e īn marāteb barā-ye ān ast ke az baʿżī revāyāt va bayānāt ke
dar ʿālam-e molk āthār-e ān ẓāher nashode możṭareb nashavand va ḥaml bar
ʿadam-e edrāk-e khod namāyand, na bar ʿadam-e ẓohūr-e maʿānī-ye ḥadīth,
zīrā ke nazd-e ān ʿebād maʿlūm nīst ke maqṣūd-e aʾemme-ye dīn che būd,
chonānche az ḥadīth mostafād mīshavad. Pas bāyad ʿebād be-īngūne ʿebārāt
khod rā az foyūżāt mamnūʿ nasāzand va az ahlash soʾāl namāyand tā
asrār-e mastūre belā ḥejāb ẓāher va vāżeḥ shavad.

In any case, these matters are mentioned so that people will not be
disturbed by certain traditions and utterances whose effects have not
appeared in the earthly realm, but will attribute this to their own lack
of understanding, rather than to the failure of the meanings of the
tradition to be fulfilled. For these servants do not know what the Imams
of the Faith intended, as may be gathered from the tradition. Therefore,
people should not deprive themselves of the outpourings of grace because
of such expressions, but should ask those who possess this knowledge, so
that the hidden mysteries may become manifest and clear, without a veil.

\subsection{286}\label{section-476}

\begin{otherlanguage}{persian}

ولیکن احدی از اهل ارض مشاهده نمیشود که طالب حقّ باشد تا آنکه در مسائل
غامضه رجوع بمظاهر احدیّه نماید کل در ارض نسیان ساکن و باهل بغی و طغیان
متّبع ولکنّ اللّه یفعل بهم کما هم یعملون و ینساهم کما نسوا لقائه فی ایّامه و
کذلک قضی علی الّذین کفروا و یقضی علی الّذینهم کانوا بآیاته یجحدون

\end{otherlanguage}

Valīkan aḥadī az ahl-e arż moshāhede namīshavad ke ṭāleb-e Ḥaqq bāshad
tā ānke dar masāʾel-e ghāmeże rojūʿ be-maẓāher-e aḥadīyye namāyad. Koll
dar arż-e nesyān sāken va be-ahl-e baghy va ṭoghyān mottabeʿ. Wa-lākinna
Allāha yafʿalu bihim kamā hum yaʿmalūn wa-yansāhum kamā nasū liqāʾahu fī
ayyāmihi wa-kadhālika qaḍā ʿalā alladhīna kafarū wa-yaqḍī ʿalā
alladhīnahum kānū bi-āyātihi yajḥadūn.

Yet not a single person among the people of the earth is seen to seek
the Truth, so as to turn to the Manifestations of Oneness concerning
abstruse questions. All dwell in the land of forgetfulness and follow
the people of rebellion and defiance. But God deals with them as they
themselves act, and forgets them as they forgot the meeting with Him in
His days. Thus has He passed judgment upon those who disbelieved, and He
passes judgment upon those who denied His signs.

\subsection{287}\label{section-477}

\begin{otherlanguage}{arabic}

و اختم القول بقوله تعالی و من یعش عن ذکر الرّحمن نقیّض له شیطاناً فهو له
قرین و من اعرض عن ذکری فانّ له معیشة ضنکاً

\end{otherlanguage}

Wa-akhtimu al-qawla bi-qawlihi taʿālā: Wa-man yaʿshu ʿan dhikri
al-Raḥmāni nuqayyiḍ lahu shayṭānan fa-huwa lahu qarīn. Wa-man aʿraḍa ʿan
dhikrī fa-inna lahu maʿīshatan ḍankan.

I conclude with His words, exalted is He: ``Whoever turns away from the
remembrance of the All-Merciful, We assign him a devil to be his
companion.'' And: ``Whoever turns away from My remembrance will indeed
have a life of hardship.''

\subsection{288}\label{section-478}

\begin{otherlanguage}{arabic}

و کذلک نزّل من قبل لو انتم تعقلون

\end{otherlanguage}

Wa-kadhālika nuzzila min qablu law antum taʿqilūn.

Thus was it sent down before, if only you understood.

\subsection{289}\label{section-479}

\begin{otherlanguage}{arabic}

المنزول من البآء و الهآء

\end{otherlanguage}

Al-manzūl min al-Bāʾ wa-l-Hāʾ.

Revealed from the Bāʾ and the Hāʾ.

\subsection{290}\label{section-480}

\begin{otherlanguage}{arabic}

و السّلام علی من سمع نغمة الورقآء فی سدرة المنتهی

\end{otherlanguage}

Wa-l-salāmu ʿalā man samiʿa naghmata al-Warqāʾ fī Sidrati al-Muntahā.

Peace be upon whoever has heard the melody of the Dove in the Lote Tree
of the Utmost Boundary.

\subsection{291}\label{section-481}

\begin{otherlanguage}{arabic}

فسبحان ربّنا الأعلی

\end{otherlanguage}

Fa-subḥāna Rabbinā al-Aʿlā.

Glory, then, be to our Lord, the Most Exalted.

\bookmarksetup{startatroot}

\chapter{Kitab-i-Ahdi (The Book of My
Covenant)}\label{kitab-i-ahdi-the-book-of-my-covenant}

\section{BH00003}\label{bh00003}

\emph{Translated on 2 Qudrat, 183 BE using GPT 6.1 Sol}

\subsection{1}\label{section-482}

\begin{otherlanguage}{arabic}

01 اگر افق اعلی از زخرف دنیا خالیست ولکن در خزائن توکّل و تفویض از برای
ورّاث میراث مرغوب لا عدل له گذاشتیم گنج نگذاشتیم و بر رنج نیفزودیم ایم
اللّه در ثروت خوف مستور و خطر مکنون انظروا ثمّ اذکروا ما انزله الرّحمن فی
الفرقان ویل لکلّ همزة لمزة الّذی جمع مالاً و عدّده ثروت عالم را وفائی نه
آنچه را فنا اخذ نماید و تغییر پذیرد لایق اعتنا نبوده و نیست مگر علی قدر
معلوم

\end{otherlanguage}

01 Agar ofoq-e aʿlā az zokhrof-e donyā khālīst valīken dar khazāʾen-e
tavakkol o tafvīż az barā-ye vorrās mīrās-e marghūb lā ʿadla lah
goẕāshtīm ganj nagoẕāshtīm o bar ranj nayafzūdīm īm Allāh dar servat
khowf mastūr o khaṭar maknūn onẓorū somma oẕkorū mā anzalah al-Raḥmān fī
al-Furqān wayl li-kulli humazatin lumazatin alladhī jamaʿa mālan
wa-ʿaddadah servat-e ʿālam rā vafāʾī na ānche rā fanā akhz namāyad o
taghyīr paẕīrad lāyeq-e eʿtenā nabūde o nīst magar ʿalā qadr-e maʿlūm

01 Though the Most Exalted Horizon is devoid of the adornments of the
world, We have left in the treasuries of trust and surrender an
incomparable and cherished inheritance for the heirs. We have left no
treasure and added no burden. By God, fear lies hidden in wealth, and
danger is concealed within it. Look, then remember what the All-Merciful
revealed in the Qur'an: ``Woe to every slanderer and backbiter who
amasses wealth and counts it.'' The wealth of the world has no
constancy. That which perishes or is subject to change has not been, and
is not, worthy of attention, except to a known measure.

\subsection{2}\label{section-483}

\begin{otherlanguage}{persian}

02 مقصود این مظلوم از حمل شداید و بلایا و انزال آیات و اظهار بیّنات اخماد
نار ضغینه و بغضا بوده که شاید آفاق افئدهٴ اهل عالم بنور اتّفاق منوّر گردد
و بآسایش حقیقی فائز گردد و از افق لوح الهی نیّر این بیان لائح و مشرق باید
کل بآن ناظر باشند ای اهل عالم شما را وصیّت مینمایم بآنچه سبب ارتفاع
مقامات شما است بتقوی اللّه تمسّک نمائید و بذیل معروف تشبّث کنید براستی
میگویم لسان از برای ذکر خیر است او را بگفتار زشت میالائید عفا اللّه عمّا
سلف از بعد باید کل بما ینبغی تکلّم نمایند از لعن و طعن و ما یتکدّر به
الانسان اجتناب نمایند مقام انسان بزرگست چندی قبل این کلمهٴ علیا از مخزن
قلم ابهی ظاهر امروز روزیست بزرگ و مبارک آنچه در انسان مستور بوده امروز
ظاهر شده و میشود مقام انسان بزرگست اگر بحقّ و راستی تمسّک نماید و بر امر
ثابت و راسخ باشد انسان حقیقی بمثابهٴ آسمان لدی الرّحمن مشهود شمس و قمر
سمع و بصر و انجم او اخلاق منیرهٴ مضیئه مقامش اعلی المقام و آثارش مربّی
امکان

\end{otherlanguage}

02 Maqṣūd-e īn maẓlūm az ḥaml-e shadāyed o balāyā o enzāl-e āyāt o
eẓhār-e bayyenāt ekhmād-e nār-e żaghīne o baghżā būde ke shāyad āfāq-e
afʾede-ye ahl-e ʿālam be-nūr-e ettefāq monavvar gardad o be-āsāyesh-e
ḥaqīqī fāʾez gardad o az ofoq-e lowḥ-e elāhī nayyer-e īn bayān lāyeḥ o
moshreq bāyad koll be-ān nāẓer bāshand ey ahl-e ʿālam shomā rā vaṣīyat
mīnamāyam be-ānche sabab-e ertefāʿ-e maqāmāt-e shomā ast be-taqvā Allāh
tamassok namāyīd o be-ẕayl-e maʿrūf tashabbos konīd be-rāstī mīgūyam
lesān az barā-ye ẕekr-e kheyr ast ū rā be-goftār-e zesht mayālāyīd ʿafā
Allāh ʿammā salaf az baʿd bāyad koll bi-mā yanbaghī takallom namāyand az
laʿn o ṭaʿn o mā yatakaddar bih al-insān ejtenāb namāyand maqām-e ensān
bozorg ast chandī qabl īn kalame-ye ʿolyā az makhzan-e qalam-e abhā
ẓāher emrūz rūzīst bozorg o mobārak ānche dar ensān mastūr būde emrūz
ẓāher shode o mīshavad maqām-e ensān bozorg ast agar be-ḥaqq o rāstī
tamassok namāyad o bar amr sābet o rāsekh bāshad ensān-e ḥaqīqī
be-masābe-ye āsemān ladā al-Raḥmān mashhūd shams o qamar samʿ o baṣar o
anjom-e ū akhlāq-e monīre-ye możīʾe maqāmash aʿlā al-maqām o āsārash
morabbī-ye emkān

02 The purpose of this Wronged One in bearing hardships and afflictions,
revealing verses, and making clear proofs manifest has been to
extinguish the fire of rancor and hatred, so that the horizons of the
hearts of the people of the world might be illumined by the light of
concord and they might attain true peace. From the horizon of the divine
Tablet, the sun of this utterance shines forth: All must turn toward it.
O people of the world, I counsel you to do what will raise your
stations. Hold fast to the fear of God and cling to what is right. I
tell you truly: the tongue is meant for speaking well of others; do not
defile it with ugly speech. God has forgiven what is past. Henceforth
all must speak as is fitting and refrain from curses, reproaches, and
anything that grieves a person. The station of the human being is great.
Some time ago, these exalted words issued from the treasury of the Most
Glorious Pen: ``Today is a great and blessed day. What was hidden within
the human being has become, and is becoming, manifest. The station of
the human being is great, provided one holds fast to truth and honesty
and remains steadfast and firm in the Cause. In the sight of the
All-Merciful, a true human being appears like the heavens: hearing and
sight are their sun and moon, and their radiant, luminous qualities are
their stars. Their station is the highest of stations, and their works
nurture the realm of existence.''

\subsection{3}\label{section-484}

\begin{otherlanguage}{arabic}

03 هر مقبلی الیوم عرف قمیص را یافت و بقلب طاهر بافق اعلی توجّه نمود او از
اهل بها در صحیفهٴ حمرا مذکور خذ قدح عنایتی باسمی ثمّ اشرب منه بذکری
العزیز البدیع

\end{otherlanguage}

03 Har moqbelī al-yawm ʿarf-e qamīṣ rā yāft o be-qalb-e ṭāher be-ofoq-e
aʿlā tavajjoh namūd ū az ahl-e Bahā dar ṣaḥīfe-ye ḥamrā maẕkūr khuḏ
qadāḥa ʿināyatī bi-ismī thumma ishrab minhu bi-dhikrī al-ʿazīz al-badīʿ

03 Whoever today has turned toward God, perceived the fragrance of the
garment, and directed a pure heart toward the Most Exalted Horizon is
numbered among the people of Bahá in the Crimson Scroll. Take the cup of
My favor in My name, then drink from it in remembrance of Me, the
Mighty, the Wondrous.

\subsection{4}\label{section-485}

\begin{otherlanguage}{persian}

04 ای اهل عالم مذهب الهی از برای محبّت و اتّحاد است او را سبب عداوت و
اختلاف منمائید نزد صاحبان بصر و اهل منظر اکبر آنچه سبب حفظ و علّت راحت و
آسایش عباد است از قلم اعلی نازل شده ولکن جهّال ارض چون مربّای نفس و هوسند
از حکمتهای بالغهٴ حکیم حقیقی غافلند و بظنون و اوهام ناطق و عامل

\end{otherlanguage}

04 Ey ahl-e ʿālam maẕhab-e elāhī az barā-ye maḥabbat o etteḥād ast ū rā
sabab-e ʿadāvat o ekhtelāf manamāyīd nazd-e ṣāḥebān-e baṣar o ahl-e
manẓar-e akbar ānche sabab-e ḥefẓ o ʿellat-e rāḥat o āsāyesh-e ʿebād ast
az qalam-e aʿlā nāzel shode valīken johhāl-e arż chon morabbā-ye nafs o
havasand az ḥekmathā-ye bāleghe-ye ḥakīm-e ḥaqīqī ghāfeland o be-ẓonūn o
owhām nāṭeq o ʿāmel

04 O people of the world, divine religion exists for love and unity. Do
not make it a cause of enmity and division. To those with insight and
those who behold the Most Great Vision, it is evident that what
preserves God's servants and brings them comfort and peace has been
revealed by the Most Exalted Pen. But the ignorant of the earth,
nurtured by selfish desire, are heedless of the consummate wisdom of the
true Wise One; they speak and act according to suspicions and vain
imaginings.

\subsection{5}\label{section-486}

\begin{otherlanguage}{persian}

05 یا اولیآء اللّه و امنائه ملوک مظاهر قدرت و مطالع عزّت و ثروت حقّند در
بارهٴ ایشان دعا کنید حکومت ارض بآن نفوس عنایت شد و قلوب را از برای خود
مقرّر داشت

\end{otherlanguage}

05 Yā awliyāʾ Allāh va omanāʾeh molūk maẓāher-e qodrat o maṭāleʿ-e
ʿezzat o servat-e Ḥaqqand dar bāre-ye īshān doʿā konīd ḥokūmat-e arż
be-ān nofūs ʿenāyat shod o qolūb rā az barā-ye khod moqarrar dāsht

05 O friends of God and His trusted ones, kings are manifestations of
God's power and dawning places of His might and wealth. Pray for them.
The government of the earth has been entrusted to those souls, while He
has reserved hearts for Himself.

\subsection{6}\label{section-487}

\begin{otherlanguage}{persian}

06 نزاع و جدال را نهی فرمود نهیاً عظیما فی الکتاب هذا امر اللّه فی هذا
الظّهور الأعظم و عصمه من حکم المحو و زیّنه بطراز الاثبات انّه هو العلیم
الحکیم

\end{otherlanguage}

06 Nezāʿ o jedāl rā nahy farmūd nahyan ʿaẓīman fī al-kitāb hādhā amr
Allāh fī hādhā al-ẓuhūr al-aʿẓam wa-ʿaṣamahu min ḥukm al-maḥw
wa-zayyanahu bi-ṭirāz al-ithbāt innahu huwa al-ʿalīm al-ḥakīm

06 He has emphatically forbidden conflict and contention in the Book.
This is God's command in this Most Great Revelation. He has protected it
from abrogation and adorned it with the distinction of permanence. He
is, indeed, the All-Knowing, the All-Wise.

\subsection{7}\label{section-488}

\begin{otherlanguage}{persian}

07 مظاهر حکم و مطالع امر که بطراز عدل و انصاف مزیّنند بر کل اعانت آن نفوس
لازم طوبی للأمرآء و العلمآء فی البهآء اولئک امنائی بین عبادی و مشارق
احکامی بین خلقی علیهم بهائی و رحمتی و فضلی الّذی احاط الوجود در کتاب اقدس
در این مقام نازل شده آنچه که از آفاق کلماتش انوار بخشش الهی لامع و ساطع
و مشرقست

\end{otherlanguage}

07 Maẓāher-e ḥokm o maṭāleʿ-e amr ke be-ṭerāz-e ʿadl o enṣāf mozayyanand
bar koll eʿānat-e ān nofūs lāzem ṭūbā li-l-umarāʾ wa-l-ʿulamāʾ fī
al-Bahāʾ ulāʾika umanāʾī bayna ʿibādī wa-mashāriq aḥkāmī bayna khalqī
ʿalayhim bahāʾī wa-raḥmatī wa-faḍlī alladhī aḥāṭa al-wujūd dar Ketāb-e
Aqdas dar īn maqām nāzel shode ānche ke az āfāq-e kalamātash anvār-e
bakhshesh-e elāhī lāmeʿ o sāṭeʿ o moshreq ast

07 All must assist those who manifest authority and serve as dawning
places of command, when they are adorned with justice and fairness.
``Blessed are the rulers and the learned among the people of Bahá. They
are My trusted ones among My servants and the dawning places of My
commandments among My creation. Upon them rest My glory, My mercy, and
My grace, which encompass all existence.'' On this subject, words have
been revealed in the Most Holy Book from whose horizons the lights of
divine generosity shine and radiate.

\subsection{8}\label{section-489}

\begin{otherlanguage}{arabic}

08 یا اغصانی در وجود قوّت عظیمه و قدرت کامله مکنون و مستور باو و جهة
اتّحاد او ناظر باشید نه باختلافات ظاهره از او

\end{otherlanguage}

08 Yā aghṣānī dar vojūd qovvat-e ʿaẓīme o qodrat-e kāmele maknūn o
mastūr be-ū o jehat-e etteḥād-e ū nāẓer bāshīd na be-ekhtelāfāt-e ẓāhere
az ū

08 O My Branches, a great strength and a perfect power lie concealed
within existence. Turn your attention to it and to the unity it brings,
not to the apparent differences arising from it.

\subsection{9}\label{section-490}

\begin{otherlanguage}{persian}

09 وصیّة اللّه آنکه باید اغصان و افنان و منتسبین طرّاً بغصن اعظم ناظر باشند
انظروا ما انزلناه فی کتابی الأقدس اذا غیض بحر الوصال و قضی کتاب المبدإ
فی المآل توجّهوا الی من اراده اللّه الّذی انشعب من هذا الأصل القدیم مقصود
از این آیهٴ مبارکه غصن اعظم بوده کذلک اظهرنا الأمر فضلاً من عندنا و انا
الفضّال الکریم قد قدّر اللّه مقام الغصن الأکبر بعد مقامه انّه هو الآمر
الحکیم قد اصطفینا الأکبر بعد الأعظم امراً من لدن علیم خبیر

\end{otherlanguage}

09 Vaṣīyat Allāh ānke bāyad aghṣān o afnān o montasebīn ṭorran
be-Ghoṣn-e Aʿẓam nāẓer bāshand unẓurū mā anzalnāhu fī kitābī al-Aqdas
idhā ghīḍa baḥr al-wiṣāl wa-quḍiya kitāb al-mabdaʾ fī al-maʾāl tawajjahū
ilā man arādahu Allāh alladhī inshaʿaba min hādhā al-aṣl al-qadīm maqṣūd
az īn āye-ye mobārake Ghoṣn-e Aʿẓam būde ka-dhālika aẓharnā al-amr
faḍlan min ʿindinā wa-anā al-faḍḍāl al-karīm qad qaddara Allāh maqām
al-Ghoṣn al-Akbar baʿda maqāmihi innahu huwa al-āmir al-ḥakīm qad
iṣṭafaynā al-Akbar baʿda al-Aʿẓam amran min ladun ʿalīm khabīr

09 God's testament is that the Branches, the Afnán, and all relatives
must turn toward the Most Great Branch. Consider what We have revealed
in My Most Holy Book: ``When the ocean of My presence has ebbed and the
Book of My Revelation has reached its end, turn toward Him whom God has
intended, who has branched from this Ancient Root.'' This blessed verse
refers to the Most Great Branch. Thus have We made the matter clear as a
favor from Us, and I am the Bountiful, the Generous. God has ordained
the station of the Greater Branch to follow his station. He is, indeed,
the Ordainer, the All-Wise. We have chosen the Greater after the Most
Great, by command of One who is All-Knowing, All-Informed.

\subsection{10}\label{section-491}

\begin{otherlanguage}{arabic}

10 محبّت اغصان بر کل لازم ولکن ما قدّر اللّه لهم حقّاً فی اموال النّاس

\end{otherlanguage}

10 Maḥabbat-e aghṣān bar koll lāzem valīken mā qaddara Allāh lahum
ḥaqqan fī amwāl al-nās

10 All must love the Branches, but God has assigned them no right to the
property of others.

\subsection{11}\label{section-492}

\begin{otherlanguage}{persian}

11 یا اغصانی و افنانی و ذوی قرابتی نوصیکم بتقوی اللّه و بمعروف و بما
ینبغی و بما ترتفع به مقاماتکم براستی میگویم تقوی سردار اعظمست از برای
نصرت امر الهی و جنودی که لایق این سردارست اخلاق و اعمال طیّبهٴ طاهرهٴ
مرضیّه بوده و هست

\end{otherlanguage}

11 Yā aghṣānī va afnānī va dhawī qarābatī nūṣīkum bi-taqwā Allāh
wa-bi-maʿrūf wa-bi-mā yanbaghī wa-bi-mā tartafiʿu bihi maqāmātukum
be-rāstī mīgūyam taqvā sardār-e aʿẓam ast az barā-ye noṣrat-e amr-e
elāhī o jonūdī ke lāyeq-e īn sardār ast akhlāq o aʿmāl-e ṭayyebe-ye
ṭāhere-ye marżīye būde o hast

11 O My Branches, My Afnán, and My relatives, We counsel you to fear
God, to do what is right and fitting, and to do what will raise your
stations. I tell you truly: the fear of God is the supreme commander in
the service of His Cause, and the troops worthy of this commander have
always been, and remain, pure and goodly qualities and deeds that are
pleasing to God.

\subsection{12}\label{section-493}

\begin{otherlanguage}{arabic}

12 بگو ای عباد اسباب نظم را سبب پریشانی منمائید و علّت اتّحاد را علّت
اختلاف مسازید امید آنکه اهل بها بکلمهٴ مبارکهٴ قل کلّ من عند اللّه ناظر
باشند و این کلمهٴ علیا بمثابهٴ آبست از برای اطفاء نار ضغینه و بغضا که در
قلوب و صدور مکنون و مخزونست احزاب مختلفه از این کلمهٴ واحده بنور اتّحاد
حقیقی فائز میشوند انّه یقول الحقّ و یهدی السّبیل و هو المقتدر العزیز الجمیل

\end{otherlanguage}

12 Begū ey ʿebād asbāb-e naẓm rā sabab-e parīshānī manamāyīd o ʿellat-e
etteḥād rā ʿellat-e ekhtelāf masāzīd omīd ānke ahl-e Bahā be-kalame-ye
mobārake-ye qul kullun min ʿind Allāh nāẓer bāshand o īn kalame-ye ʿolyā
be-masābe-ye āb ast az barā-ye eṭfā-ye nār-e żaghīne o baghżā ke dar
qolūb o ṣodūr maknūn o makhzūn ast aḥzāb-e mokhtalefe az īn kalame-ye
vāḥede be-nūr-e etteḥād-e ḥaqīqī fāʾez mīshavand innahu yaqūlu al-ḥaqq
wa-yahdī al-sabīl wa-huwa al-muqtadir al-ʿazīz al-jamīl

12 Say: O servants, do not turn the means of order into a cause of
turmoil, or a source of unity into a cause of division. It is hoped that
the people of Bahá will heed the blessed words ``Say: All is from God.''
These exalted words are like water for extinguishing the fire of rancor
and hatred hidden and stored in hearts. Through these single words,
differing groups can attain the light of true unity. He speaks the truth
and guides the way; He is the All-Powerful, the Mighty, the Beautiful.

\subsection{13}\label{section-494}

\begin{otherlanguage}{persian}

13 احترام و ملاحظهٴ اغصان بر کل لازم لاعزاز امر و ارتفاع کلمه و این حکم
از قبل و بعد در کتب الهی مذکور و مسطور طوبی لمن فاز بما امر به من لدن
آمر قدیم و همچنین احترام حرم و آل اللّه و افنان و منتسبین و نوصیکم بخدمة
الأمم و اصلاح العالم از ملکوت بیان مقصود عالمیان نازل شد

\end{otherlanguage}

13 Eḥterām o molāḥeẓe-ye aghṣān bar koll lāzem li-iʿzāz al-amr
wa-irtifāʿ al-kalima o īn ḥokm az qabl o baʿd dar kotob-e elāhī maẕkūr o
masṭūr ṭūbā li-man fāza bi-mā umira bihi min ladun āmir qadīm o
hamchenīn eḥterām-e ḥaram o āl Allāh o afnān o montasebīn wa-nūṣīkum
bi-khidmat al-umam wa-iṣlāḥ al-ʿālam az malakūt-e bayān-e Maqṣūd-e
ʿālamiyān nāzel shod

13 All must show respect and consideration for the Branches, for the
honor of the Cause and the exaltation of its Word. This command has been
recorded in the divine Books, both before and since. Blessed is whoever
has attained what was commanded by the Ancient Ordainer. Likewise,
respect must be shown to the Holy Household, the family of God, the
Afnán, and the relatives. ``We counsel you to serve the nations and work
for the betterment of the world'' was revealed from the realm of the
utterance of the Desire of all peoples.

\subsection{14}\label{section-495}

\begin{otherlanguage}{arabic}

14 آنچه که سبب حیات عالم و نجات اممست نصایح قلم اعلی را بگوش حقیقی اصغا
نمائید انّها خیر لکم عمّا علی الأرض یشهد بذلک کتابی العزیز البدیع

\end{otherlanguage}

14 Ānche ke sabab-e ḥayāt-e ʿālam o najāt-e omam ast naṣāyeḥ-e qalam-e
aʿlā rā be-gūsh-e ḥaqīqī eṣghā namāyīd innahā khayrun lakum ʿammā ʿalā
al-arḍ yashhadu bi-dhālika kitābī al-ʿazīz al-badīʿ

14 Listen attentively to the counsels of the Most Exalted Pen, which
bring life to the world and salvation to the nations. They are better
for you than anything on earth. My mighty and wondrous Book bears
witness to this.

\bookmarksetup{startatroot}

\chapter{Kitab-i-Badi (The Wondrous
Book)}\label{kitab-i-badi-the-wondrous-book}

\section{BH00004}\label{bh00004}

\emph{Translated on 9 Qudrat, 183 BE using GPT 6.1 Sol}

\subsection{1}\label{section-496}

\begin{otherlanguage}{arabic}

هوالله المستعان

\end{otherlanguage}

Huwa Allāh al-Mustaʿān.

He is God, whose help is sought.

\subsection{2}\label{section-497}

\begin{otherlanguage}{persian}

مكتوبت كه مشعر بر بعض روايات غير محقّقه و اخبارات غير صادقه بود باين عبد
فانى واصل از جهتى ابواب سرور گشود و از جهتى ابواب حزن مفتوح نمود امّا
سرور چه كه مشعر بر سلامتى بود و امّا حزن نظر بآنكه مشاهده شد كه بالأخره
انفس مشركه تو را از شطر احديّه ممنوع نموده‌اند و از نفحات حبّ الهّيه محروم
داشته‌اند نسئل اللّه بأن ينصرفك عنهم و يقلبك الى وجهه المنير

\end{otherlanguage}

Maktūbat kih mushʿir bar baʿḍ rivāyāt-i ghayr-i muḥaqqaqih va ikhbārāt-i
ghayr-i ṣādiqih būd bi-īn ʿabd-i fānī vāṣil, az jihatī abvāb-i surūr
gushūd va az jihatī abvāb-i ḥuzn maftūḥ namūd. Ammā surūr chih kih
mushʿir bar salāmatī būd va ammā ḥuzn naẓar bi-ānkih mushāhidih shud kih
bi'l-ākharah anfus-i mushrikih tu rā az shaṭr-i aḥadiyyih mamnūʿ
namūdih-and va az nafaḥāt-i ḥubb-i ilāhiyyih maḥrūm dāshtih-and. Nasʾalu
Allāha bi-an yanṣarifaka ʿanhum wa yuqallibaka ilā wajhihi al-munīr.

Your letter, conveying certain unverified accounts and untrue reports,
reached this transitory servant. In one respect it opened the gates of
joy, and in another it opened the gates of sorrow. As for joy, it gave
news of your well-being; as for sorrow, it became apparent that, in the
end, souls who associate partners with God have kept you from the realm
of divine unity and deprived you of the breezes of divine love. We ask
God to turn you away from them and direct you toward His radiant face.

\subsection{3}\label{section-498}

\begin{otherlanguage}{persian}

و چون اين عبد ملاحظه نمود كه آنچه مرقوم داشته از فطرت ظاهر نشده بلكه از
كلمات كاذبه انفس مشركه بوده كه نزد تو مذكورداشته‌اند لذا خالصاً لوجه اللّه
و حبّاً لك لازم شد كه جواب مذكور دارم كه شايد بقوّه ربّ الأرباب حجاب اعظم و
سبحات جلال را بانامل قدرت و انقطاع شقّ نمائى و بمقرّى كه مقدّس از ذكر و
اشاره و كلمات و عبارات است صعود نمائى و سبّاحان قلزم كبريائى و طايران
هواء عزّ ربّانى را از طيور ظلمانى تميز دهى

\end{otherlanguage}

Va chūn īn ʿabd mulāḥiẓih namūd kih ānchih marqūm dāshtih az fiṭrat
ẓāhir nashudih balkih az kalimāt-i kādhibih-yi anfus-i mushrikih būdih
kih nazd-i tu madhkūr dāshtih-and, lidhā khāliṣan li-wajhi Allāh va
ḥubban laka lāzim shud kih javāb madhkūr dāram kih shāyad bi-quvvih-yi
Rabb al-Arbāb ḥijāb-i aʿẓam va subuḥāt-i jalāl rā bi-anāmil-i qudrat va
inqiṭāʿ shaqq namāʾī va bi-maqarrī kih muqaddas az dhikr va ishārih va
kalimāt va ʿibārāt ast ṣuʿūd namāʾī va sabbāḥān-i qulzum-i kibriyāʾī va
ṭāyirān-i havā-yi ʿizz-i rabbānī rā az ṭuyūr-i ẓulmānī tamīz dihī.

When this servant observed that what you had written had not arisen from
your own nature, but consisted of the lying words of souls who associate
partners with God, which they had recounted to you, it became necessary,
solely for the sake of God and out of love for you, to give a reply.
Perhaps, through the power of the Lord of Lords, you may rend the
greatest veil and the veils of glory with the fingers of power and
detachment, ascend to an abode sanctified beyond mention, allusion,
words, and expressions, and distinguish those who swim in the ocean of
majesty and soar in the heaven of divine glory from the birds of
darkness.

\subsection{4}\label{section-499}

\begin{otherlanguage}{persian}

اگرچه مشاهده ميشود كه حجاب اكبر بصرت را از منظر اكبر بكلّى محروم نموده
ولكن چون بر حسب ظاهر كلمات نصحيّه ذكر نموده بودى لذا بر اين عبد لازم شد
كه بآنچه از طمطام يمّ عنايت الهى و ظهورات تجلّى عزّ صمدانى فايز شده‌ام بهمان
مقدار تو را اخبار نمايم كه شايد از حرم عزّ تقديس محروم نگردى و از كعبه
قدس تجريد بی‌‌نصيب نمانى

\end{otherlanguage}

Agarchih mushāhidih mīshavad kih ḥijāb-i akbar baṣarat rā az manẓar-i
akbar bi-kullī maḥrūm namūdih va lākin chūn bar ḥasab-i ẓāhir kalimāt-i
nuṣḥiyyih dhikr namūdih būdī lidhā bar īn ʿabd lāzim shud kih bi-ānchih
az ṭamṭām-i yamm-i ʿināyat-i ilāhī va ẓuhūrāt-i tajallī-yi ʿizz-i
ṣamadānī fāʾiz shudih-am bi-hamān miqdār tu rā ikhbār namāyam kih shāyad
az ḥaram-i ʿizz-i taqdīs maḥrūm nagardī va az Kaʿbih-yi quds-i tajrīd
bī-naṣīb namānī.

Although it is apparent that the greatest veil has entirely deprived
your sight of the greatest vision, since you had offered words that
outwardly took the form of counsel, it became necessary for this servant
to inform you, to the extent of what I have received, of the surging
ocean of divine favor and the manifestations of the radiance of
everlasting glory. Perhaps you may not be deprived of the sanctuary of
sanctified glory or left without a share in the sacred Kaaba of
detachment.

\subsection{5}\label{section-500}

\begin{otherlanguage}{persian}

اگرچه كلماتى از لسان مبارك استماع شد و از مضمون آن اين عبد از اكثرى از
ناس مأيوس شده و لكن خالصاً لوجه اللّه اظهار ميشود كه لعلّ بنسيمى از رضوان
عزّ تقديس فايز شوى و بثمره از اثمار سدره لاشرقيّه و لاغربيّه مرزوق گردى و
آن كلمات اينست كه فرمودند انّ النّاس اكثرهم يسترضعون اليوم من ثدى الغفلة و
الجهل و ما جائتهم ایّام فطامهم فكيف بلوغهم لذا لا تؤثّر فيهم الكلمة لانّهم
لا يفقهون و لا يشعرون و انّك فاشهد هؤلآء الخلق كاغنام يذهبهم صبی‌ من
الصّبيان كيف يشآء كذلك نزل من قبل و لكن النّاس هم لايعرفون

\end{otherlanguage}

Agarchih kalimātī az lisān-i mubārak istimāʿ shud va az maḍmūn-i ān īn
ʿabd az aktharī az nās maʾyūs shudih va lākin khāliṣan li-wajhi Allāh
iẓhār mīshavad kih laʿall bi-nasīmī az Riḍvān-i ʿizz-i taqdīs fāʾiz
shavī va bi-thamarih az athmār-i sidrih-yi lā-sharqiyyih va
lā-gharbiyyih marzūq gardī va ān kalimāt īnast kih farmūdand: Inna
al-nāsa aktharahum yastarḍiʿūna al-yawma min thadyi al-ghaflati wa
al-jahli wa mā jāʾathum ayyāmu fiṭāmihim fa-kayfa bulūghuhum lidhā lā
tuʾaththiru fīhimu al-kalimatu li-annahum lā yafqahūna wa lā yashʿurūna
wa innaka fa-ishhad hāʾulāʾi al-khalqa ka-aghnāmin yadhhabuhum ṣabiyyun
mina al-ṣibyāni kayfa yashāʾu kadhālika nazala min qablu wa lākinna
al-nāsa hum lā yaʿrifūn.

Although words were heard from the blessed tongue whose import has
caused this servant to lose hope in most people, this is nevertheless
set forth solely for the sake of God, so that perhaps you may receive a
breeze from the Paradise of sanctified glory and be granted a fruit from
the tree that is neither of the East nor of the West. These are the
words that were spoken: ``Most people today suckle at the breast of
heedlessness and ignorance. The days of their weaning have not yet come,
much less their maturity. Therefore the Word has no effect upon them,
for they neither understand nor perceive. Behold these people as sheep
whom a child leads wherever he wishes. Thus was it revealed before, yet
the people do not recognize it.''

\subsection{6}\label{section-501}

\begin{otherlanguage}{persian}

و چون در مكتوبت مطالب عديده ذكر شده بود لذا بر اين عبد لازم شد كه هر يك
از مطالب تو را مرقوم نمايد و بعد جواب مذكور دارد كه لعلّ در اين ایّام كه
ظلمت جميع ارض را احاطه نموده بعضى از قاصدين و طالبين بانوار شمس هدايت
مهتدى شوند و از كوثر حيوان بی‌‌نصيب نمانند و انّه ليهدى من يشآء و انّه على
كلّ شیء قدير

\end{otherlanguage}

Va chūn dar maktūbat maṭālib-i ʿadīdih dhikr shudih būd lidhā bar īn
ʿabd lāzim shud kih har yak az maṭālib-i tu rā marqūm namāyad va baʿd
javāb madhkūr dārad kih laʿall dar īn ayyām kih ẓulmat jamīʿ-i arḍ rā
iḥāṭih namūdih baʿḍī az qāṣidīn va ṭālibīn bi-anvār-i shams-i hidāyat
muhtadī shavand va az kawthar-i ḥayavān bī-naṣīb namānand. Wa innahu
la-yahdī man yashāʾu wa innahu ʿalā kulli shayʾin qadīr.

Since numerous points were mentioned in your letter, it became necessary
for this servant to write down each of your points and then give a
reply, so that perhaps, in these days when darkness has encompassed the
whole earth, some of those who seek and aspire may be guided by the
light of the sun of guidance and not be left without a share of the
Kawthar of life. He indeed guides whomever He wills, and He has power
over all things.

\subsection{7}\label{section-502}

\begin{otherlanguage}{persian}

اينست اوّل مطلب تو كه مرقوم داشته اميد از خداوند متعال چنان است كه ملاقات
نصيب گرداند كه شايد رفع تعيّنات ظلمانيّه و حجبات وهميّه گردد و معلوم شود
سابحين در بحر احديّه و طايرين در فضاى رفعت از ساكنين در بئر ظلمت و واقفين
در بيداى حيرت انتهى

\end{otherlanguage}

Īnast avval maṭlab-i tu kih marqūm dāshtih: Umīd az Khudāvand-i mutaʿāl
chunān ast kih mulāqāt naṣīb gardānad kih shāyad rafʿ-i taʿayyunāt-i
ẓulmāniyyih va ḥujubāt-i vahmiyyih gardad va maʿlūm shavad sābiḥīn dar
baḥr-i aḥadiyyih va ṭāyirīn dar faḍā-yi rifʿat az sākinīn dar biʾr-i
ẓulmat va vāqifīn dar baydā-yi ḥayrat. Intahā.

This is the first point you wrote: ``My hope in God, the Most Exalted,
is that He may grant us a meeting, so that perhaps the dark limitations
and veils of illusion may be removed, and those who swim in the sea of
divine unity and soar in the heights may be distinguished from those who
dwell in the pit of darkness and remain in the desert of bewilderment.''
End of quotation.

\subsection{8}\label{section-503}

\begin{otherlanguage}{persian}

اينكه خواسته بودى كه ملاقات حاصل شود و از آن ملاقات رفع حجبات گردد هر
ملاقاتى علّت كشف حجاب نبوده و نخواهد بود در اكثرى ملاقات بر حجاب بيفزايد
چنانچه نفسيكه منبع غلّ و معدن فساد بوده آمده و تو حجبات او را خرق ننموده
كه سهل است بر حجبات تو افزوده اگر تو قادر بر خرق حجاب بودى البتّه حجبات
مانعه از نفس محجوبه را كشف مينمودى و لكن انشآءاللّه از فضل الهى اميدواريم
كه ملاقاتى بروح و ريحان دست دهد و جميع در انجمن قدس مجتمع شويم و از
اشارات لفظيّه مقدّس شده بعوالم نامتناهى معانى پى بريم چه كه اين ایّام ایّام
كشف و شهود است نه ایّام تحقيق ظنون و موهوم از خدا وصل و لقا را طالبيم و
انّه لنعم المولى و نعم المجيب معلومست كه اين كلمات از خودت نبوده از وسوسه
انفس بعيده مرقوم داشته

\end{otherlanguage}

Īnkih khvāstih būdī kih mulāqāt ḥāṣil shavad va az ān mulāqāt rafʿ-i
ḥujubāt gardad, har mulāqātī ʿillat-i kashf-i ḥijāb nabūdih va nakhvāhad
būd. Dar aktharī mulāqāt bar ḥijāb bīfzāyad chunānchih nafsīkih manbaʿ-i
ghill va maʿdan-i fasād būdih āmadih va tu ḥujubāt-i ū rā kharq
nanamūdih kih sahl ast bar ḥujubāt-i tu afzūdih. Agar tu qādir bar
kharq-i ḥijāb būdī al-battih ḥujubāt-i māniʿih az nafs-i maḥjūbih rā
kashf mīnamūdī va lākin in shāʾ Allāh az faḍl-i ilāhī umīdvārīm kih
mulāqātī bi-rūḥ va rayḥān dast dihad va jamīʿ dar anjuman-i quds
mujtamiʿ shavīm va az ishārāt-i lafẓiyyih muqaddas shudih bi-ʿavālim-i
nāmutanāhī-yi maʿānī pay barīm chih kih īn ayyām ayyām-i kashf va shuhūd
ast nah ayyām-i taḥqīq-i ẓunūn va mawhūm. Az Khudā vaṣl va liqā rā
ṭālibīm va innahu la-niʿma al-mawlā wa niʿma al-mujīb. Maʿlūmast kih īn
kalimāt az khudat nabūdih az vasvasih-yi anfus-i baʿīdih marqūm dāshtih.

As for your wish that a meeting take place and that through it the veils
be removed, not every meeting has been or will be a cause of unveiling.
In most cases, a meeting adds to the veil. Thus, a person who was a
source of malice and a mine of corruption came, and far from your
rending his veils, he added to yours. Had you been capable of rending a
veil, you would certainly have removed the obstructing veils from that
veiled soul. Nevertheless, God willing, we hope through divine grace
that a meeting may take place in joy and delight, that we may all gather
in the assembly of holiness and, freed from verbal allusions, find our
way into the limitless worlds of meaning. For these days are days of
unveiling and witnessing, not days of investigating conjectures and
imaginings. From God we seek reunion and meeting; He is indeed the best
Protector and the best to answer. It is clear that these words were not
your own: you wrote them at the prompting of souls far removed.

\subsection{9}\label{section-504}

\begin{otherlanguage}{persian}

ذات نايافته از هستى بخش كى تواند كه شود هستى بخش

\end{otherlanguage}

Dhāt nāyāftih az hastī-bakhsh kay tavānad kih shavad hastī-bakhsh.

A self that has received nothing from the Giver of being---how can it
become a giver of being?

\subsection{10}\label{section-505}

\begin{otherlanguage}{persian}

امثال آن نفوس پژمرده محرومه يمين را ازشمال تميز نداده‌اند چگونه خود را
هادى و منجى فرض گرفته‌اند نفحات انفس الهيّه از فوحات انفس مشركه بسى واضح و
معلوم است نه هر كه بصورت انسان از انسان محسوب و نه هر كه بدو كلمه ناطق
از اهل بيان مذكور انفس محدوده از عالم احديّه محروم بوده و خواهند بود و
ابصر مرموده از مشاهده انوار آفتاب معانى بی‌نصيب عجب است از تو كه جوهر وهم
را نفس يقين دانسته و نفس موهوم را معلوم ذكر نموده انّهم خلقوا من الوهم و
بدئوا منه و رجعوا اليه و اليوم طايران هواى قرب الهى و سبّاحان قلزم رحمت
رحمانى چون شمس مشرق و مضيئند و واللّه اگر نفوس محتجبه ابصر خود را از رمد
نفس و هوى مطهّر سازند فى الحين باشراق نيّر آفاق‌ فايز شوند

\end{otherlanguage}

Amthāl-i ān nufūs-i pazhmurdih-yi maḥrūmih yamīn rā az shimāl tamīz
nadādih-and chigūnih khud rā hādī va munjī farḍ giriftih-and. Nafaḥāt-i
anfus-i ilāhiyyih az fawḥāt-i anfus-i mushrikih basī vāḍiḥ va maʿlūm
ast. Nah har kih bi-ṣūrat-i insān az insān maḥsūb va nah har kih bi-du
kalimih nāṭiq az ahl-i bayān madhkūr. Anfus-i maḥdūdih az ʿālam-i
aḥadiyyih maḥrūm būdih va khvāhand būd va abṣur-i marmūdih az
mushāhidih-yi anvār-i āftāb-i maʿānī bī-naṣīb. ʿAjab ast az tu kih
jawhar-i vahm rā nafs-i yaqīn dānistih va nafs-i mawhūm rā maʿlūm dhikr
namūdih. Innahum khuliqū mina al-wahmi wa budiʾū minhu wa rajaʿū ilayhi.
Va imrūz ṭāyirān-i havā-yi qurb-i ilāhī va sabbāḥān-i qulzum-i raḥmat-i
raḥmānī chūn shams mushriq va muḍīʾand va va'llāh agar nufūs-i
muḥtajibih abṣur-i khud rā az ramad-i nafs va havā muṭahhar sāzand fī
al-ḥīn bi-ishrāq-i nayyir-i āfāq fāʾiz shavand.

The likes of those withered, deprived souls have not distinguished right
from left. How have they presumed themselves to be guides and saviors?
The distinction between the breezes of divine souls and the odors of
souls who associate partners with God is abundantly clear and evident.
Not everyone with a human form is counted as human, nor is everyone who
utters two words reckoned among the people of utterance. Limited souls
have been and will remain deprived of the realm of divine unity, and
diseased eyes have no share in beholding the rays of the sun of
meanings. It is astonishing that you have taken the essence of illusion
for certainty itself, and described what is imagined as something known.
They were created from illusion, originated from it, and returned to it.
Today, those who soar in the heaven of nearness to God and swim in the
ocean of the All-Merciful's mercy are radiant and shining like the sun.
By God, if veiled souls cleanse their eyes of the disease of self and
desire, they will at once receive the radiance of the Luminary of the
horizons.

\subsection{11}\label{section-506}

\begin{otherlanguage}{persian}

و لكن چه ذكر نمايم كه امر بسيار بر تو مشتبه شده و هذا لم يكن من عندك بل
من الّذى لم يكن له ذكر عنداللّه و كان محروماً عن نفحات ایّام اللّه المهيمن
القيّوم قولك بلى اگرچه حالتم بر حسب ظاهر خوب و در كمال صحّت در گذر است ولى
از بس فسرده‌ام يكى از محرومى فيض خدمت احباب و ديگر آنكه هر قدر جان ميكنيم
و كوشش مينمائيم و از شهر بشهر و وادى بوادى ميدويم بلكه نوعى شود خيال
آسوده شود نميشود و نميگذارند عجب است از تو و اين عبد خادم بسيار متحيّر
شده از اين عبارات كه مرقوم داشته كه شهر بشهر ميدويم كه خيال آسوده شود و
نميشود چه كه از اين كلمه همچه مفهوم شد كه نعوذ باللّه ملتفت نشده كه شهر
بشهر گشتن سبب آسود‌گى خيال نبوده و نخواهد شد كاش تو در همان محلّ توقف
مينمودى و يك نفس خالصى للّه برمياوردى و البتّه آن نفس بمكمن رحمن حاضر ميشد
چه كه اگر مقصود از حركت تو عرفان حقّ بوده بايد باينمدينه بيائى مطهّرا عن
كلّ الاشارات و الدّلالات و بعد ملاقات نموده خود ببصر خود تميز دهى و حقّ را
از دونش بشناسى نه بقول نفسيكه عنداللّه مردود بوده بشأنى محتجب شوى كه خود
را از جمال مقصود بالمرّه محروم دارى و السّالك لن يبلغ الى وطن الأمن و
الأمان الّا بكفّ الصّفر عما فى الامكان

\end{otherlanguage}

Va lākin chih dhikr namāyam kih amr bisyār bar tu mushtabih shudih. Wa
hādhā lam yakun min ʿindika bal mina alladhī lam yakun lahu dhikrun
ʿinda Allāhi wa kāna maḥrūman ʿan nafaḥāti ayyāmi Allāhi al-Muhaymini
al-Qayyūm. Qawluka: Balī agarchih ḥālatam bar ḥasab-i ẓāhir khūb va dar
kamāl-i ṣiḥḥat dar gudhar ast valī az bas fasurdih-am, yakī az
maḥrūmī-yi fayḍ-i khidmat-i aḥbāb va dīgar ānkih har qadr jān mīkanīm va
kūshish mīnamāyīm va az shahr bi-shahr va vādī bi-vādī mīdavīm balkih
nawʿī shavad khiyāl āsūdih shavad namīshavad va namīgudhārand. ʿAjab ast
az tu va īn ʿabd-i khādim bisyār mutaḥayyir shudih az īn ʿibārāt kih
marqūm dāshtih kih shahr bi-shahr mīdavīm kih khiyāl āsūdih shavad va
namīshavad chih kih az īn kalimih hamchih mafhūm shud kih naʿūdhu
bi'llāh multafit nashudih kih shahr bi-shahr gashtan sabab-i āsūdigī-yi
khiyāl nabūdih va nakhvāhad shud. Kāsh tu dar hamān maḥall tavaqquf
mīnamūdī va yak nafas-i khāliṣī li'llāh bar mīāvardī va al-battih ān
nafas bi-makman-i Raḥmān ḥāḍir mīshud chih kih agar maqṣūd az ḥarakat-i
tu ʿirfān-i Ḥaqq būdih bāyad bi-īn madīnih biyāʾī muṭahharan ʿan kulli
al-ishārāti wa al-dalālāt va baʿd mulāqāt namūdih khud bi-baṣar-i khud
tamīz dihī va Ḥaqq rā az dūnash bishināsī nah bi-qawl-i nafsīkih ʿinda
Allāh mardūd būdih bi-shaʾnī muḥtajib shavī kih khud rā az jamāl-i
Maqṣūd bi'l-marrih maḥrūm dārī. Wa al-sāliku lan yablugha ilā waṭani
al-amni wa al-amāni illā bi-kaffi al-ṣifri ʿammā fī al-imkān.

But what can I say, when the matter has become so confused for you? This
did not come from you, but from one who had no mention in God's presence
and was deprived of the breezes of the days of God, the All-Protecting,
the Self-Subsisting. Your words: ``Yes, although outwardly my condition
is good and I am passing my days in perfect health, I am deeply
dispirited: first, from being deprived of the blessing of serving the
friends, and second, because however much we struggle and strive,
running from city to city and valley to valley in the hope that somehow
we may find peace of mind, it does not happen, and they do not allow
it.'' This is astonishing on your part, and this servant is greatly
bewildered by these expressions you have written: that you run from city
to city to find peace of mind, yet do not find it. For these words
suggest that, God forbid, you have failed to realize that roaming from
city to city has never been and will never be a cause of peace of mind.
If only you had remained in that same place and drawn a single breath
purely for God! That breath would certainly have reached the hidden
abode of the All-Merciful. For if the purpose of your journey was
recognition of the Truth, you should come to this city, purified of all
allusions and indications, and then meet and discern with your own eyes,
distinguishing the Truth from all else---not become so veiled by the
words of a soul rejected in God's sight that you deprive yourself
altogether of the beauty of the Desired One. The wayfarer will never
reach the homeland of safety and security except with hands empty of all
that is in the contingent world.

\subsection{12}\label{section-507}

\begin{otherlanguage}{persian}

قوله رفته بود كه قدرى آسایش حاصل نمائيم و از شرّ فتنه ايمن باشيم فراغتى
حاصل شود وساوس شياطين بهواى نفس شرّى برانگيخته خلل در اعتقادات انداخته
فتنه نوى بر پا نموده الفتنة نائمة لعن اللّه من ايقظها انتهى تكاد ان تنفطر
السّمآء ء من هذا القول و تنشق الأرض و تندكّ الجبال و تنقطع البحار عن
الأمواج و الأشجار من الأثمار فواللّه تكاد ان تنعدم اركان الوجود من الغيب
و الشّهود و بما خرج من فمك اصفرّت وجوه المقرّبين و لو يكون ذو دراية ليوقن
بأنّ وجه نقطة الأولی حينئذ اصفرّ من لطمات هذه الكلمات و اذاً ينوح فى رفيق
الأعلی و ينوح معه ملائكة المقرّبين هيچ ملاحظه نموده‌ئی كه چه نوشته مقصود
از اين بيان چه و مقصود كه و لكن از سجيّه اهل اللّه نبوده كه فوق شأن خود
تكلّم نمايند من لم يعرف حدّه و شأنه لم يعرف ربّه

\end{otherlanguage}

Qawluhu: Raftih būd kih qadrī āsāyish ḥāṣil namāyīm va az sharr-i fitnih
īman bāshīm farāghatī ḥāṣil shavad vasāvis-i shayāṭīn bi-havā-yi nafs
sharrī bar angīkhtih khalal dar iʿtiqādāt andākhtih fitnih-yi nawī bar
pā namūdih. Al-fitnatu nāʾimatun laʿana Allāhu man ayqaẓahā. Intahā.
Takādu an tanfaṭira al-samāʾu {[}corrupted text{]} min hādhā al-qawli wa
tanshaqqa al-arḍu wa tandakka al-jibālu wa tanqaṭiʿa al-biḥāru ʿani
al-amwāji wa al-ashjāru mina al-athmār. Fa-wa'llāhi takādu an tanʿadima
arkānu al-wujūdi mina al-ghaybi wa al-shuhūd wa bi-mā kharaja min famika
iṣfarrat wujūhu al-muqarrabīn wa law yakūnu dhū dirāyatin li-yūqina
bi-anna wajha Nuqṭati al-Ūlā ḥīnaʾidhin iṣfarra min laṭamāti hādhihi
al-kalimāt wa idhan yanūḥu fī Rafīqi al-Aʿlā wa yanūḥu maʿahu malāʾikatu
al-muqarrabīn. Hīch mulāḥiẓih namūdih-ī kih chih nivishtih maqṣūd az īn
bayān chih va maqṣūd kih va lākin az sajiyyih-yi ahl-i Allāh nabūdih kih
fawq-i shaʾn-i khud takallum namāyand. Man lam yaʿrif ḥaddahu wa
shaʾnahu lam yaʿrif rabbahu.

His words: ``It had been intended that we should gain some comfort, be
safe from the evil of strife, and find some respite. The whisperings of
devils, following selfish desire, have stirred up mischief, undermined
beliefs, and raised a new disturbance. Strife is asleep; may God curse
whoever awakens it.'' End of quotation. The heaven is almost rent
{[}corrupted original text{]} by this statement, the earth split apart,
the mountains crushed, the seas deprived of their waves, and the trees
of their fruits. By God, the very pillars of existence, unseen and
visible, are almost brought to nothing. What has issued from your mouth
has turned pale the faces of those brought near to God. Anyone possessed
of understanding would surely know that the face of the Primal Point has
now grown pale from the blows of these words. He therefore laments in
the companionship of the Most High, and with Him lament the angels who
are brought near. Have you considered at all what you have written, what
is meant by this statement, and whom it concerns? Yet it is not the
character of the people of God to speak beyond their station. Whoever
has not known his own limits and station has not known his Lord.

\subsection{13}\label{section-508}

\begin{otherlanguage}{persian}

بارى اين عباد كه بأيادى تربيت الهيّه تربيت شده‌اند و در ظل سدره ادب ساكنند
دوست نداشته و ندارند كه اين كلمات را بمثل تو جواب ذكر نمايند فوّضنا
الأمور الى اللّه و انّه لنعم الوكيل بشنو نداى عبد ناصح را فضل اعظم را غضب
مدان و صلح اكبر را فتنه مشمر و رحمت كليّه الهيّه را قهر مبين ماء عذب فرات
را از ملح اجاج بصر ظاهره تميز ندهد بلكه محتاج بذائقه بوده و خواهد بود و
ذائقه هم اگر بامراض وارده تغيير نمايد انّهم از ادراك عاجز خواهد بود از
خدا ميطلبيم ذائقه قدسيّه لطيفه عنايت فرمايد كه كوثر بيان الهى را از ملح
اجاج نفسانى تميز دهيم و الأمر بيده انّه على كلّ شیء حكيم و على كلّ شیء قدير

\end{otherlanguage}

Bārī īn ʿibād kih bi-ayādī-yi tarbiyat-i ilāhiyyih tarbiyat shudih-and
va dar ẓill-i sidrih-yi adab sākinand dūst nadāshtih va nadārand kih īn
kalimāt rā bi-mithl-i tu javāb dhikr namāyand. Fawwaḍnā al-umūra ilā
Allāhi wa innahu la-niʿma al-wakīl. Bishnaw nidā-yi ʿabd-i nāṣiḥ rā
faḍl-i aʿẓam rā ghaḍab madān va ṣulḥ-i akbar rā fitnih mashmur va
raḥmat-i kulliyyih-yi ilāhiyyih rā qahr mabīn. Māʾ-i ʿadhb-i furāt rā az
milḥ-i ujāj baṣar-i ẓāhirih tamīz nadihad balkih muḥtāj bi-dhāʾiqih
būdih va khvāhad būd va dhāʾiqih ham agar bi-amrāḍ-i vāridih taghyīr
namāyad ān ham az idrāk ʿājiz khvāhad būd. Az Khudā mīṭalabīm
dhāʾiqih-yi qudsiyyih-yi laṭīfih ʿināyat farmāyad kih kawthar-i bayān-i
ilāhī rā az milḥ-i ujāj-i nafsānī tamīz dihīm. Wa al-amru bi-yadihi
innahu ʿalā kulli shayʾin ḥakīmun wa ʿalā kulli shayʾin qadīr.

In any case, these servants, who have been nurtured by the hands of
divine education and dwell in the shade of the tree of courtesy, have
never wished, nor do they now wish, to answer such words from one like
you. We have entrusted all matters to God; He is indeed the best
Trustee. Listen to the call of a servant who offers sincere counsel: do
not mistake the greatest grace for wrath, count the greatest peace as
strife, or regard God's all-embracing mercy as vengeance. The outward
eye cannot distinguish sweet, refreshing water from bitter salt water;
this has always required, and will always require, the sense of taste.
If the sense of taste is altered by diseases that afflict it, it too
will be incapable of perception. We ask God to bestow a delicate,
sanctified sense of taste, so that we may distinguish the Kawthar of
divine utterance from the bitter salt water of the self. The matter is
in His hand. He is wise in all things and has power over all things.

\subsection{14}\label{section-509}

\begin{otherlanguage}{persian}

اگر آسایش ميخواهى اين بيان حقّ را بشنو كه فرموده از آلايش نفس و هوى پاك
شو تا بآسایش لايفنى فايز شوى اين آسایش در ظلّ حقّ بوده نه در دونش چه كه جز
ظلّش اليوم برهوت عنداللّه محسوب انشآء اللّه اميدواريم كه صدر را از كلمات و
اشارات مشركين ممرّد نموده تا محلّ تجلّى سلطان احديّه شود وبالمرّه از قيل و
قال و جنگ و جدال بياسايد

\end{otherlanguage}

Agar āsāyesh mīkhwāhī īn bayān-e Ḥaqq rā beshnow keh farmūdeh: az
ālāyesh-e nafs va havā pāk show tā be-āsāyesh-e lā-yafnā fāyez shavī. Īn
āsāyesh dar ẓell-e Ḥaqq būdeh nah dar dūnash, cheh keh joz ẓellash
al-yawm barahūt ʿinda Allāh maḥsūb. In shāʾ Allāh omīdvārīm keh ṣadr rā
az kalemāt va eshārāt-e moshrekīn momarrad namūdeh tā maḥall-e
tajallī-ye Solṭān-e Aḥadīyeh shavad va bel-marreh az qīl o qāl va jang o
jedāl beyāsāyad.

If you desire rest, listen to this utterance of God: ``Purify yourself
of the defilement of self and desire, so that you may attain
imperishable rest.'' This rest lies in the shade of God, not elsewhere,
for today everything outside His shade is reckoned a barren wilderness
in God's sight. God willing, We hope that you will cleanse your breast
of the words and insinuations of those who associate partners with God,
so that it may become the seat of the revelation of the Sovereign of
Oneness and find complete rest from disputation, conflict, and
contention.

\subsection{15}\label{section-510}

\begin{otherlanguage}{persian}

مذكور نموده بودى كه ازشرّ فتنه ايمن باشيم همچه مفهوم ميشود كه تو ظهورات
حقّيّه كه در عالم شده از آدم تا خاتم جميع را فتنه دانسته و اعتقاد نموده كه
تكلّم نمودن بكلمه حقّ فتنه است و لكن زين فتنه تا آن فتنه فرقى است بسيار و
لايعقلها الّا اولو الأبصار فتنه حقّ فتنه‌ايست كه صد‌هزار شريعه امن و امان از
او جارىست فواللّه اگر ذى شمّى يافت شود از اين فتنه نفحات احديّه استشمام
نمايد و ديگر ميفرمايد قوله جلّ كبريائه الم احسب النّاس ان يتركوا ان يقولوا
آمنّا و هم لايفتنون

\end{otherlanguage}

Madhkūr namūdeh būdī keh az sharr-e fetneh īmen bāshīm. Hamcheh mafhūm
mīshavad keh to ẓohūrāt-e ḥaqqīyeh keh dar ʿālam shodeh az Ādam tā
Khātam jamīʿ rā fetneh dānesteh va eʿteqād namūdeh keh takallom namūdan
be-kalemeh-ye ḥaqq fetneh ast. Va lāken zīn fetneh tā ān fetneh farqī
ast besyār, va lā yaʿqiluhā illā ulū al-abṣār. Fetneh-ye Ḥaqq fetneh-īst
keh ṣad-hezār sharīʿeh-ye amn o amān az ū jārīst. Fa-wa-Allāh agar
dhī-shammī yāft shavad az īn fetneh nafaḥāt-e Aḥadīyeh estesḥmām
namāyad. Va dīgar mīfarmāyad, qawluhu jalla kibriyāʾuhu: Alif Lām Mīm.
A-ḥasiba al-nāsu an yutrakū an yaqūlū āmannā wa-hum lā yuftanūn.

You had written, ``May we be safe from the evil of strife.'' It would
appear that you have regarded all the Manifestations of God who have
appeared in the world, from Adam to the Seal, as strife, and have
believed that speaking the word of truth is strife. Yet there is a vast
difference between this strife and that strife, and none comprehend it
except those possessed of insight. The strife of God is a strife from
which a hundred thousand streams of peace and security flow. By God, if
anyone possessed of a sense of smell were found, he would inhale the
fragrances of Oneness from this strife. Elsewhere He says---His words,
exalted be His majesty: ``Alif. Lam. Mim. Do people suppose that they
will be left to say, `We believe,' without being put to the test?''

\subsection{16}\label{section-511}

\begin{otherlanguage}{persian}

اگر آنچه از حقّ ظاهر ميشود آن را فتنه انگاشته صد‌هزار جان فداى چنين فتنه
كه انفس مشركه فرعونيّه را از مظاهر احديّه تميز داده و تفصيل فرموده ای عبد
از خدا ميطلبم كه بصر حديد بجميع ناس عنايت فرمايد و نفس زكيّه مرحمت نمايد
و اذن واعيه مبذول دارد تا كلمات مشركه را از كلمات الهيّه فرق گذارند

\end{otherlanguage}

Agar āncheh az Ḥaqq ẓāher mīshavad ān rā fetneh engāshteh, ṣad-hezār jān
fadā-ye chonīn fetneh keh anfos-e moshrekeh-ye ferʿownīyeh rā az
maẓāher-e Aḥadīyeh tamyīz dādeh va tafṣīl farmūdeh. Ey ʿAbd, az Khodā
mīṭalabam keh baṣar-e ḥadīd be-jamīʿ-e nās ʿenāyat farmāyad va nafs-e
zakīyeh marḥamat namāyad va odhon-e vāʿīyeh mabdhūl dārad tā kalemāt-e
moshrekeh rā az kalemāt-e elāhīyeh farq godhārand.

If you have deemed what appears from God to be strife, may a hundred
thousand lives be sacrificed for such strife, which has distinguished
and separated the idolatrous, Pharaoh-like souls from the Manifestations
of Oneness. O servant, I ask God to bestow keen sight upon all people,
to grant them pure souls, and to endow them with attentive ears, so that
they may distinguish idolatrous words from divine words.

\subsection{17}\label{section-512}

\begin{otherlanguage}{persian}

هرگز راضى نميشوم كه تو چنين كلمات مرقوم نمائى ذكر شده بود كه وساوس
شياطين بهواى نفس شرّى برانگيخته خلل در اعتقادات انداخته‌اى برادر من نفحات
قدسيّه ربّ العالمين را وساوس شياطين ذكر منما فواللّه يبكى كلّ شیء و الّذى
امرك بما ظهر من قلمك يلعنه كلّ الذّرّات و لكن هو لا يستشعر فى نفسه لانّ اللّه
ضرب على قلبه غشاوة النّار يقول ما لا يفقه و تكلّم بما لا تكلّم به مظاهر
النّفى كلّها و اذا هو فى خسران عظيم تكاد الوجود ان يذوب من هذه الكلمة و
لكن ملقيك فى حجاب غليظ و امّا خلل در اعتقادات كه مذكور داشته بدانكه در
اعتقادات احدى از ظهوراللّه و امره خلل نرسيده بلكه محكم و ثابت و راسختر
شده مگر آن نفوسيكه در ظهور قبل كينوناتشان از سجّين بوده و بر حسب ظاهر
بتصديق آن ظهور از علّييّن محسوب شده‌اند و چون امتحان كبرى ظاهر آن نفوس راجع
باصل خود شده و نفوس قدسيّه منيره بوطن اصليّه خود كه اقرار و اعتراف بنفس
ظهور است واصل شده‌اند طوبی‌ للواصلين

\end{otherlanguage}

Hargez rāḍī namīshavam keh to chonīn kalemāt marqūm namāʾī. Dhekr shodeh
būd keh vasāves-e shayāṭīn be-havā-ye nafs sharrī bar angīkhteh khalal
dar eʿteqādāt andākhteh. Ey barādar-e man, nafaḥāt-e qodsīyeh-ye Rabb
al-ʿālamīn rā vasāves-e shayāṭīn dhekr manamā. Fa-wa-Allāh yabkī kullu
shayʾ wa-alladhī amaraka bi-mā ẓahara min qalamika yalʿanuhu kullu
al-dharrāt, wa-lākin huwa lā yastashʿiru fī nafsihi li-anna Allāha
ḍaraba ʿalā qalbihi ghishāwata al-nār. Yaqūlu mā lā yafqahu wa-takallama
bi-mā lā takallama bihi maẓāhiru al-nafy kulluhā wa-idhā huwa fī khusrān
ʿaẓīm. Takādu al-wujūd an yadhūba min hādhihi al-kalimah wa-lākin
mulqīka fī ḥijāb ghalīẓ. Va ammā khalal dar eʿteqādāt keh madhkūr
dāshteh, bedān keh dar eʿteqādāt-e aḥadī az ẓohūr Allāh va amrihi khalal
naresīdeh balkeh moḥkam va thābet va rāsekhtar shodeh, magar ān nofūsī
keh dar ẓohūr-e qabl kaynūnātashān az Sijjīn būdeh va bar ḥasab-e ẓāher
be-taṣdīq-e ān ẓohūr az ʿIllīyīn maḥsūb shodeh-and. Va chon emteḥān-e
kobrā ẓāher, ān nofūs rājeʿ be-aṣl-e khod shodeh va nofūs-e qodsīyeh-ye
monīreh be-vaṭan-e aṣlīyeh-ye khod keh eqrār va eʿterāf be-nafs-e ẓohūr
ast vāṣel shodeh-and. Ṭūbā li-l-wāṣilīn.

I can never consent to your writing such words. It had been stated that
``the whisperings of devils, following selfish desire, have stirred up
mischief and undermined beliefs.'' My brother, do not call the holy
breezes of the Lord of the worlds the whisperings of devils. By God,
everything weeps, and every atom curses the one who instructed you in
what has issued from your pen. Yet he is unaware within himself, for God
has laid a veil of fire over his heart. He says what he does not
understand and has spoken words that none of the manifestations of
negation have spoken; he is indeed in grievous loss. Existence itself
almost melts away at this statement, yet the one who prompted you is
wrapped in a thick veil. As for the undermining of beliefs that you
mentioned, know that no one's beliefs have been undermined by the
revelation of God and His Cause. Rather, they have become firm,
steadfast, and more deeply rooted, except in the case of those souls
whose very beings were of Sijjin in the previous revelation, but who,
through their outward acknowledgment of that revelation, were reckoned
as belonging to Illiyyin. When the greatest test appeared, those souls
returned to their origin, while the radiant, holy souls reached their
original homeland, which is acknowledgment and recognition of the
Manifestation Himself. Blessed are those who have arrived.

\subsection{18}\label{section-513}

\begin{otherlanguage}{persian}

آن اعتقاداتى كه منوط بعبادت عجل شده اصفياءاللّه و احبّائه از آن مقدّس و
منزّه و مبرّا و هر نفسيكه العياذ باللّه در اين ظهور نيّر معانى محتجب ماند
عقيده نداشته و نخواهد داشت تا خلل بهمرسد آيا نشنيده نغمه ورقاى الهى را
كه ميفرمايد اگر نفسى بدوام ملك و ملكوت بقيام و قعود و سجود حقّ را ذاكر
شود و عبادت نمايد و در حين ظهور اقلّ من آن توقّف نمايد ليحبط اعماله و بقدر
همان توقّف در نار بوده و خواهد بود تا چه رسد باعراض و استكبار نسئل اللّه
بأن يعصمنا من وساوس الشّياطين الّذين كفروا باللّه جهرة و يحسبون انّهم من
العارفين لا فو‌ربّ العالمين

\end{otherlanguage}

Ān eʿteqādātī keh manūṭ be-ʿebādat-e ʿejl shodeh, aṣfiyāʾ Allāh va
aḥibbāʾihi az ān moqaddas va monazzah va mobarrā. Va har nafsī keh
al-ʿiyādhu bi-Allāh dar īn ẓohūr-e Nayyir-e Maʿānī moḥtajeb mānd,
ʿaqīdeh nadāshteh va nakhwāhad dāsht tā khalal be-ham resad. Āyā
nashenīdeh naghmeh-ye Varqā-ye Elāhī rā keh mīfarmāyad agar nafsī
be-davām-e molk va malakūt be-qīyām va qoʿūd va sojūd Ḥaqq rā dhāker
shavad va ʿebādat namāyad va dar ḥīn-e ẓohūr aqall-e men ān tavaqqof
namāyad, la-yaḥbaṭu aʿmāluhu, va be-qadr-e hamān tavaqqof dar nār būdeh
va khwāhad būd, tā cheh resad be-eʿrāḍ va estekbār. Nasʾalu Allāha bi-an
yaʿṣimanā min wasāwis al-shayāṭīn alladhīna kafarū bi-Allāhi jahratan
wa-yaḥsabūna annahum min al-ʿārifīn. Lā fa-wa-Rabb al-ʿālamīn.

God's chosen ones and His loved ones are sanctified, purified, and free
from those beliefs that depend upon worship of the calf. Any soul who,
God forbid, remains veiled in this revelation of the Luminary of
Meanings has never possessed, nor will ever possess, a belief that could
be undermined. Have you not heard the song of the divine Dove, who says
that if a soul were to remember and worship God, standing, sitting, and
bowing down, for as long as the earthly and heavenly kingdoms endure,
yet were to hesitate for less than an instant at the time of His
revelation, his deeds would surely come to nothing? For the duration of
that hesitation he has been, and will remain, in the fire---how much
more so in the case of rejection and arrogance. We ask God to protect us
from the whisperings of devils who have openly disbelieved in God and
imagine themselves to be among those who know. No, by the Lord of the
worlds!

\subsection{19}\label{section-514}

\begin{otherlanguage}{persian}

فتنه نو ذكر نموده بودى امر جديدى بر پا نشده و ظاهر نگشته اين عباد از اوّل
يوميكه بعرفان اللّه فايز شده‌ايم نظر باين ظهور قدس صمدانى داشته‌ايم و مادون
اين ظهور كليّه الهيّه را در ساحتش مفقود ميشمرديم چنانچه بعضى مذكور
ميداشتند كه لأجل مصلحت ايشان اشاره بغير ميفرمايند و بعد كه معلوم شد آنچه
ذكر شد حقّ بود و ما بعده الّا الضّلال عجب است از تو كه چنين ذكر نموده

\end{otherlanguage}

Fetneh-ye now dhekr namūdeh būdī. Amr-e jadīdī bar pā nashodeh va ẓāher
nagashteh. Īn ʿebād az avval-e yawmī keh be-ʿerfān Allāh fāyez shodeh-īm
naẓar be-īn ẓohūr-e qods-e ṣamadānī dāshteh-īm va mā-dūn-e īn ẓohūr-e
kollīyeh-ye elāhīyeh rā dar sāḥatash mafqūd mīshomordīm, chonāncheh
baʿḍī madhkūr mīdāshtand keh li-ajl-e maṣlaḥat īshān eshāreh be-ghayr
mīfarmāyand. Va baʿd keh maʿlūm shod āncheh dhekr shod ḥaqq būd va mā
baʿdahu illā al-ḍalāl. ʿAjab ast az to keh chonīn dhekr namūdeh.

You had spoken of a new disturbance. No new cause has been established
or appeared. From the very first day that we, these servants, attained
recognition of God, we have had our eyes fixed upon this holy, eternal
Revelation, and have counted everything other than this universal divine
Revelation as nonexistent in His presence. Indeed, some used to say that
He referred to another out of expediency. Later it became clear that
what had been said was the truth, and beyond it there is nothing but
error. It is astonishing that you should have spoken in this way.

\subsection{20}\label{section-515}

\begin{otherlanguage}{persian}

اگر مقصود از فتنه فساد جديد بوده واللّه الّذى لا اله الّا هو كه همان نفسيكه
اين كلمات را بتو القاء نموده او بر پا كرده چنانچه مشهود است چهل سنه با
هم جمع بوده‌اند و خود تو موقنى بآنچه ذكر شده و در اين سفر كه آن مشرك در
حزب اللّه وارد شد عامل شد آنچه را كه نمرود و شدّاد عامل نشده مع‌ذلك
ندانستيم كه سبب چه شده كه باغواى او از بصر و سمع و فؤاد خود گذشته و آنچه
او گفته مسلّم داشته

\end{otherlanguage}

Agar maqṣūd az fetneh fasād-e jadīd būdeh, wa-Allāh alladhī lā ilāha
illā huwa keh hamān nafsī keh īn kalemāt rā be-to elqā namūdeh ū bar pā
kardeh, chonāncheh mashhūd ast. Chehel saneh bā ham jamʿ būdeh-and va
khod-e to mūqenī be-āncheh dhekr shodeh. Va dar īn safar keh ān moshrek
dar ḥezb Allāh vāred shod ʿāmel shod āncheh rā keh Nimrūd va Shaddād
ʿāmel nashodeh. Maʿa dhālek nadānestīm keh sabab cheh shodeh keh
be-eghvā-ye ū az baṣar va samʿ va foʾād-e khod godhashteh va āncheh ū
gofteh mosallam dāshteh.

If by strife you meant a new outbreak of disorder, by God, besides whom
there is no god, the very person who prompted you to write these words
has stirred it up, as is evident. They had been together for forty
years, and you yourself are certain of what has been said. On this
journey, when that idolater entered among the people of God, he
committed what Nimrod and Shaddad had not committed. Even so, We have
not understood what has caused you, through his deception, to abandon
your own sight, hearing, and heart, and to accept whatever he has said
as certain.

\subsection{21}\label{section-516}

\begin{otherlanguage}{persian}

نسئل اللّه بأن يوفّقك على امره و يرزقك عرفان مظهر نفسه و يطهّرك عن اشارات
المعرضين آنچه از حقّ ظاهر رحمت است از براى موحّدين و نقمت است از براى
مشركين نسئله بأن يرزقك كأساً من مآء الحيوان الّذى جرى فى قطب الامكان لعلّ
بها تنقطع عمّا سواه و تستقرب الى مقرّ الّذى لا يرى فيه الّا انوار وجه ربّك
العلىّ العظيم و از خدا ميطلبيم كه از عبادى محسوب نشوى كه لازال برحقّ امثال
اين اعتراضات را نموده‌اند چنانچه بحسين بن علىّ روح ما سواه فداه مشركين
اعتراضاً عليه نوشتند كه تو سبب تفريق اجتماع امّت شده و در دين جدّت فتنه
نموده و بعضى بدعتها گذاشته و همين كلمات مفصّل در كتب ثبت است ملاحظه
نمائيد ديگر احتياج بذكر اين بنده نيست از ذكر فقرۀ آخر تكاد ان تنعدم
اركان كلّ شیء و لكن صبرنا و نصبر بحول اللّه و قوّته و رجعنا الحكم الى اللّه
المقتدر المهيمن القيّوم

\end{otherlanguage}

Nasʾalu Allāha bi-an yuwaffiqaka ʿalā amrihi wa-yarzuqaka ʿirfāna
maẓhari nafsihi wa-yuṭahhiraka ʿan ishārāt al-muʿriḍīn. Āncheh az Ḥaqq
ẓāher raḥmat ast az barā-ye movaḥḥedīn va neqmat ast az barā-ye
moshrekīn. Nasʾaluhu bi-an yarzuqaka kaʾsan min māʾ al-ḥayawān alladhī
jarā fī quṭb al-imkān laʿalla bihā tanqaṭiʿu ʿammā siwāhu wa-tastaqribu
ilā maqarri alladhī lā yurā fīhi illā anwāru wajhi Rabbika al-ʿAliyy
al-ʿAẓīm. Va az Khodā mīṭalabīm keh az ʿebādī maḥsūb nashavī keh lā-zāl
bar Ḥaqq amthāl-e īn eʿterāḍāt rā namūdeh-and, chonāncheh be-Ḥosayn ibn
ʿAlī, rūḥu mā siwāhu fidāhu, moshrekīn iʿtirāḍan ʿalayhi neveshtand keh
to sabab-e tafrīq-e ejtemāʿ-e ommat shodeh va dar dīn-e jaddat fetneh
namūdeh va baʿḍī bedʿathā godhāshteh. Va hamīn kalemāt mofaṣṣal dar
kotob thabt ast. Molāḥeẓeh namāʾīd, dīgar eḥtīāj be-dhekr-e īn bandeh
nīst. Az dhekr-e feqreh-ye ākhar takādu an tanʿadima arkānu kulli shayʾ,
wa-lākin ṣabarnā wa-naṣbiru bi-ḥawl Allāhi wa-quwwatihi wa-rajaʿnā
al-ḥukma ilā Allāhi al-Muqtadir al-Muhaymin al-Qayyūm.

We ask God to aid you in His Cause, to grant you recognition of the
Manifestation of His own Self, and to purify you of the insinuations of
those who turn away. What appears from God is mercy for those who
acknowledge His oneness and chastisement for those who associate
partners with Him. We ask Him to grant you a cup of the water of life
that has flowed at the heart of contingent existence, so that through it
you may detach yourself from everything but Him and draw near to the
seat where nothing is seen but the lights of the face of your Lord, the
Most High, the Almighty. We ask God that you may not be counted among
those servants who have always raised such objections against God. Thus
the idolaters wrote in protest to Husayn ibn Ali---may the souls of all
others be sacrificed for Him---``You have caused division in the unity
of the community, stirred up strife in the religion of your grandfather,
and introduced certain innovations.'' These very words are recorded at
length in the books. Consult them; there is no further need for this
servant to speak. At the mention of the final passage, the pillars of
all things almost cease to exist. Yet We have been patient, and remain
patient through God's power and might, and We have referred the judgment
to God, the All-Powerful, the All-Protecting, the Self-Subsisting.

\subsection{22}\label{section-517}

\begin{otherlanguage}{persian}

فتنه كه تصوّر نموده لم ‌يزل شأن حقّ نبوده و اين ظهور مخصوص ايتلاف و اتّحاد و
اتّفاق اهل ارض ظاهر شده كه شايد جميع امم مختلفه بر شريعه واحده جمع شوند و
كلّ بتسبيح و تهليل پروردگار ناطق گردند معلوم شد كه مصلح و مفسد را فرق
نگذاشته انشآء اللّه اميدواريم از فضل الهى كه اگر بمظاهر علّييّن عارف نشده
اقلّاً مظاهر شياطين و مطالع سجّين و مفسدين را بشناسى

\end{otherlanguage}

Fetneh keh taṣavvor namūdeh lam-yazal shaʾn-e Ḥaqq nabūdeh va īn ẓohūr
makhṣūṣ-e eʾtelāf va etteḥād va ettefāq-e ahl-e arḍ ẓāher shodeh keh
shāyad jamīʿ-e omam-e mokhtalefeh bar sharīʿeh-ye vāḥedeh jamʿ shavand
va koll be-tasbīḥ va tahlīl-e Parvardegār nāṭeq gardand. Maʿlūm shod keh
moṣleḥ va mofsed rā farq nagodhāshteh. In shāʾ Allāh omīdvārīm az faḍl-e
elāhī keh agar be-maẓāher-e ʿIllīyīn ʿāref nashodeh aqallan maẓāher-e
shayāṭīn va maṭāleʿ-e Sijjīn va mofsedīn rā beshenāsī.

The strife you have imagined has never been God's way. This Revelation
has appeared specifically to bring about harmony, unity, and concord
among the people of the earth, so that perhaps all the diverse nations
may gather around a single faith, and all may voice the praise and
oneness of the Lord. It has become clear that you have not distinguished
the one who puts things right from the one who causes corruption. God
willing, We hope, through divine grace, that if you have not recognized
the manifestations of Illiyyin, you may at least recognize the
manifestations of devils, the dawning places of Sijjin, and the workers
of corruption.

\subsection{23}\label{section-518}

\begin{otherlanguage}{arabic}

لم ‌يزل حقّ ناس را باتّحاد و اتّفاق خوانده در كلمات منزله ملاحظه نما تا بر
تو معلوم شود هذه الكلمة ما كان من عندك لانّا ظننّا بأنّك ما ارتكبت هذا
الذّنب من عند نفسك بل بما الهمك الّذى كان كأحد منّا فلمّا ارتكب ما ارتكب
اطرده اللّه عن ساحة عزّه لذا قام على الاعراض قد اتّخذ لنفسه عجلا جسدا له
خوار و يدعو النّاس اليه اعتراضاً على مظهر نفس اللّه و اثبات الرّياسة لنفسه و
له كذلك كان الأمر ان انت من العارفين

\end{otherlanguage}

Lam-yazal Ḥaqq nās rā be-etteḥād va ettefāq khwāndeh. Dar kalemāt-e
monzaleh molāḥeẓeh namā tā bar to maʿlūm shavad. Hādhihi al-kalimah mā
kāna min ʿindika li-annā ẓanannā bi-annaka mā irtakabta hādhā al-dhanba
min ʿindi nafsika bal bi-mā alhamaka alladhī kāna ka-aḥadin minnā.
Fa-lammā irtakaba mā irtakaba aṭradahu Allāhu ʿan sāḥati ʿizzihi li-dhā
qāma ʿalā al-iʿrāḍ. Qad ittakhadha li-nafsihi ʿijlan jasadan lahu khuwār
wa-yadʿū al-nāsa ilayhi iʿtirāḍan ʿalā maẓhari nafs Allāh wa-ithbāta
al-riyāsati li-nafsihi wa-lahu. Kadhālika kāna al-amru in anta min
al-ʿārifīn.

God has always called people to unity and concord. Examine the revealed
words, so that this may become clear to you. This statement did not
originate with you, for We have supposed that you did not commit this
sin on your own initiative, but through what was suggested to you by one
who was as one of Us. When he committed what he committed, God expelled
him from the court of His glory; therefore he arose in opposition. He
has taken for himself a calf, a body that lows, and calls people to it
in protest against the Manifestation of God's own Self, and to establish
leadership for himself and for it. Such is the matter, if you are among
those who know.

\subsection{24}\label{section-519}

\begin{otherlanguage}{persian}

و اين كلمات كه كبد اصفياء از آن محترق شده از معانى و روح عارى و خالى
بوده و خواهد بود و لكن در غمد آن كلمات عوض معانى اسياف شاحذه مستور و سرّاً
در كلّ حين بجسد امر وارد اى كاش صاحب بصر و درايتى يافت ميشد تا ادراك
نمايد چه ظلمى اليوم ظاهر شده و ظالمين تا بچه شأن بر ضرّ امر الهى
ايستاده‌اند و لكن نسئل اللّه بأن ينزل من غمام امره ما يلحقهم بمقرّهم فى
اسفل السّافلين

\end{otherlanguage}

Va īn kalemāt keh kabed-e aṣfiyā az ān moḥtareq shodeh az maʿānī va rūḥ
ʿārī va khālī būdeh va khwāhad būd, va lāken dar ghemd-e ān kalemāt
ʿavaḍ-e maʿānī asyāf-e shāḥedheh mastūr va serran dar koll-e ḥīn
be-jasad-e Amr vāred. Ey kāsh ṣāḥeb-e baṣar va derāyatī yāft mīshod tā
edrāk namāyad cheh ẓolmī al-yawm ẓāher shodeh va ẓālemīn tā be-cheh
shaʾn bar ḍorr-e Amr-e Elāhī īstādeh-and. Va lāken nasʾalu Allāha bi-an
yunazzila min ghamāmi amrihi mā yulḥiquhum bi-maqarrihim fī asfal
al-sāfilīn.

These words, which have set the hearts of the chosen ones ablaze, have
been and will remain devoid of meaning and spirit. Yet concealed within
the sheath of these words, in place of meanings, are sharpened swords
that secretly strike the body of the Cause at every moment. If only
someone possessed of sight and discernment could be found to perceive
what oppression has appeared today, and to what extent the oppressors
have risen to harm the Cause of God. We ask God to send down from the
clouds of His command that which will consign them to their abode in the
lowest depths.

\subsection{25}\label{section-520}

\begin{otherlanguage}{persian}

و ديگر نوشته بودى قولك از اين جهت خيال آنكه جوابی‌ نويسم نداشتم ميخواستم
دم فرو بسته از اين مقوله سخن نگويم حالت قيل و قال نمانده اين دو روز عمر
لايق آن نيست كه همه بقال و جنگ و شرّ و تأويلات زوركى بخواهش نفس گذرد
انتهى

\end{otherlanguage}

Va dīgar neveshteh būdī, qawluka: az īn jahat khayāl-e ānkeh javābī
nevīsam nadāshtam. Mīkhwāstam dam forū basteh az īn maqūleh sokhan
nagūyam. Ḥālat-e qīl o qāl namāndeh. Īn do rūz ʿomr lāyeq-e ān nīst keh
hameh be-qāl va jang va sharr va taʾvīlāt-e zūrakī be-khwāhesh-e nafs
godharad. Intahā.

You had also written---your words: ``For this reason I had no intention
of writing a reply. I wished to remain silent and not speak of such
matters. I no longer have the disposition for disputation. These two
days of life are not worth spending entirely on argument, conflict,
mischief, and forced interpretations at the bidding of the self.'' End
of quotation.

\subsection{26}\label{section-521}

\begin{otherlanguage}{persian}

كاش جواب نمينوشتى و اگر هم مينوشتى از خود مرقوم ميداشتى نه كلمات مجعوله
انفس مشركه را البتّه از چنين اذكار كه نموده دم فروبستن صدهزار هزار مرتبه
ارجح بوده كلام خوب است در ذكر و اثبات امراللّه نه در مجادله و محاربه با
او چنانچه در وصاياى الهى مذكور

\end{otherlanguage}

Kāsh javāb namīneveshtī va agar ham mīneveshtī az khod marqūm mīdāshtī
nah kalemāt-e majʿūleh-ye anfos-e moshrekeh rā. Albatteh az chonīn
adhkār keh namūdeh dam forū bastan ṣad-hezār hezār martabeh arjaḥ būdeh.
Kalām khūb ast dar dhekr va ethbāt-e Amr Allāh nah dar mojādeleh va
moḥārebeh bā ū, chonāncheh dar vaṣāyā-ye elāhī madhkūr.

If only you had not written a reply, or, if you had, had written from
yourself rather than the fabricated words of idolatrous souls.
Certainly, remaining silent would have been a hundred thousand thousand
times preferable to making such statements as you have made. Speech is
good when it serves to proclaim and establish the Cause of God, not to
dispute and wage war with Him, as is stated in the divine counsels.

\subsection{27}\label{section-522}

\begin{otherlanguage}{persian}

لسان از براى ذكر رحمن خلق شده زينهار او را بكلمات مظاهر شيطان ميالائيد
لسان مشرق انوار صدق است او را مطلع كلمات كذب منمائيد لسان مخزن بيان و
ثناى من است او را بذكر اين و آن مشغول مسازيد چه كه او را در مقام خود آيه
فرمايد نوشته بودى اين دو روزه عمر لايق و لاينفد قرار فرموده‌ام چنانچه
آنچه نطق نمايد قادر است مادام كه وجود باقى نطق از او مقطوع نشود مگر مرضى
حايل شود مابين او و ثنآء ربّ او و از او عارف بصير ادراك مينمايد كه علم
الهى انتها نداشته و نخواهد داشت لآلىٴ بيان عالمين در او مستور و لكن بر
جثّه او نيفزوده و انهار معانى از او جارى و از او ذرّه‌ئی كم نگشته

\end{otherlanguage}

Lesān az barā-ye dhekr-e Raḥmān khalq shodeh, zīnhār ū rā be-kalemāt-e
maẓāher-e shayṭān mayālāʾīd. Lesān mashreq-e anvār-e ṣedq ast, ū rā
maṭlaʿ-e kalemāt-e kedhb manamāʾīd. Lesān makhzan-e bayān va thanā-ye
man ast, ū rā be-dhekr-e īn o ān mashghūl masāzīd, cheh keh ū rā dar
maqām-e khod āyeh farmāyad. Neveshteh būdī īn do rūzeh ʿomr lāyeq va
lā-yanfad qarār farmūdeh-am, chonāncheh āncheh noṭq namāyad qāder ast.
Mādām keh vojūd bāqī noṭq az ū maqṭūʿ nashavad magar maraḍī ḥāyel shavad
mābeyn-e ū va thanā-ye Rabb-e ū. Va az ū ʿāref-e baṣīr edrāk mīnamāyad
keh ʿelm-e elāhī entehā nadāshteh va nakhwāhad dāsht. Laʾālī-ye bayān-e
ʿālamīn dar ū mastūr va lāken bar joth-theh-ye ū nayafzūdeh va anhār-e
maʿānī az ū jārī va az ū dharreh-ī kam nagashteh.

The tongue was created to remember the All-Merciful; take care not to
defile it with the words of the manifestations of Satan. The tongue is
the dawning place of the lights of truthfulness; do not make it the
rising place of words of falsehood. The tongue is the treasury of My
utterance and praise; do not occupy it with mention of this person and
that, for in its own station He designates it a sign. You had written,
``These two days of life are worthy,'' and I have made it inexhaustible,
for it is capable of uttering whatever it may. As long as existence
endures, speech will not be cut off from it, unless an illness
intervenes between it and the praise of its Lord. From it the discerning
knower perceives that divine knowledge has never had, nor will ever
have, an end. The pearls of the utterance of the worlds are concealed
within it, yet have not increased its bulk; rivers of meanings flow from
it, yet it has not diminished by a single atom.

\subsection{28}\label{section-523}

\begin{otherlanguage}{persian}

بارى در اينمقام اذكار بسيار است و اين عبد باين مختصر اكتفا نموده كه شايد
لسان ظاهر و باطن را از مرض نفس و هوى مطهّر دارى و بثناء اللّه ناطق نمائى
چه كه اليوم ناطق آن است كه بثناء حقّ قيام نمايد و الّا ناعق بوده و خواهد
بود

\end{otherlanguage}

Bārī dar īn maqām adhkār besyār ast va īn ʿabd be-īn mokhtaṣar ektefā
namūde ke shāyad lesān-e ẓāher va bāṭen rā az maraż-e nafs va havā
moṭahhar dārī va be-thanā-ye Allāh nāṭeq namāʾī che ke al-yawm nāṭeq ān
ast ke be-thanā-ye Ḥaqq qīyām namāyad va ellā nāʿeq būde va khvāhad būd.

In any case, much could be said on this subject, but this servant has
confined himself to this brief account, that perhaps you may purify your
outward and inward tongue of the disease of self and desire and cause it
to speak in praise of God. For today, one who truly speaks is one who
arises to praise God; otherwise, he has been and will remain merely a
crier.

\subsection{29}\label{section-524}

\begin{otherlanguage}{persian}

نوشته بودى حالت قيل و قال نمانده كسى بجنگ و جدال بر نخواسته عجب است كه
خود عامليد آنچه را از او شكايت مينمائيد قيل و قال عالم در اينكلمات كه
ارسال داشته مستور و مكنون گويا در اين مكتوب تو ترك نشده ذرّه‌ئی از اقوال
منكرين و مشركين كتاب هر نفسى در دستش داده شده و لكن نظر بعدم التفات مشعر
نشده و ديار بديار ميفرستند

\end{otherlanguage}

Neveshte būdī ḥālat-e qīl va qāl namānde kasī be-jang va jedāl bar
nakhāste ʿajab ast ke khvod ʿāmelīd ānche rā az ū shekāyat mīnamāʾīd qīl
va qāl-e ʿālam dar īn kalemāt ke ersāl dāshte mastūr va maknūn gūyā dar
īn maktūb-e to tark nashode dharre-ʾī az aqvāl-e monkerīn va moshrekīn
ketāb-e har nafsī dar dastash dāde shode va lāken naẓar be-ʿadam-e
eltefāt moshʿer nashode va dīyār be-dīyār mīferestand.

You had written, ``There is no disposition left for disputation; no one
has risen to fight and quarrel.'' It is astonishing that you yourselves
practice what you complain of. The disputation of the whole world lies
concealed and hidden in these words you have sent. It seems that not a
particle of the sayings of the deniers and the idolaters has been left
out of this letter of yours. Each soul's book has been placed in his
hand, yet through inattention he has not become aware of it, and they
send it from land to land.

\subsection{30}\label{section-525}

\begin{otherlanguage}{persian}

بسيار اين عبد تحيّر دارد كه چگونه تو اين كلمات ثقليّه را حمل نموده كاش
بمبدئش راجع مينمودند خوب يا بد لايق ملقى بوده و خواهد بود از فضل اعظم
الهى ميطلبيم كه تو را از آن دلالات وهميّه مقدّس قابل آن نيست كه همه بقيل و
قال و جنگ و شرّ و تأويلات زوركى بجهت خواهش نفس گذرد عجب است كه اثبات
انقطاع بلسان و ادّعا شده چنانچه نفسيكه بجميع نواهى كه تصوّر شود عامل بوده
چنانچه حال در هر ارضى يك علامتى از آن شجره اعراض باقىیست مع‌ذلك او را
منقطع دانسته و از آن گذشته نفسيكه بكلّ فواحش ظاهر و نار حرصش بشدّتى مشتعل
كه از اينجا باسلامبول بجهت لقمه نان دويده اين بفعل از او ظاهر شده و لكن
بعضى از متوهّمين را كه ملاقات نمايد يك آه طولانى و يك تحريك رأس و دو كلمه
اظهار تبرّى از دنيا رأس المنقطعين محسوب

\end{otherlanguage}

Besyār īn ʿabd taḥayyor dārad ke chegūne to īn kalemāt-e thaqīle rā ḥaml
namūde kāsh be-mabdaʾash rājeʿ mīnamūdand khūb yā bad lāyeq-e molqī būde
va khvāhad būd az fażl-e aʿẓam-e elāhī mīṭalabīm ke to rā az ān
dalālāt-e vahmīye moqaddas qābel-e ān nīst ke hame be-qīl va qāl va jang
va sharr va taʾvīlāt-e zūrakī be-jehat-e khvāhesh-e nafs godharad ʿajab
ast ke ethbāt-e enqeṭāʿ be-lesān va eddeʿā shode chenānche nafsī ke
be-jamīʿ-e navāhī ke taṣavvor shavad ʿāmel būde chenānche ḥāl dar har
arżī yek ʿalāmatī az ān shajare-ye eʿrāż bāqīst maʿa dhālek ū rā
monqaṭeʿ dāneste va az ān godhashte nafsī ke be-koll-e favāḥesh ẓāher va
nār-e ḥerṣash be-sheddatī moshtaʿel ke az īnjā be-Eslāmbūl be-jehat-e
loqme nān davīde īn be-feʿl az ū ẓāher shode va lāken baʿżī az
motavahhemīn rā ke molāqāt namāyad yek āh-e ṭūlānī va yek taḥrīk-e raʾs
va do kaleme eẓhār-e tabarrī az donyā raʾs-e monqaṭeʿīn maḥsūb.

This servant is greatly astonished that you have borne these weighty
words. If only they would return them to their source! Good or bad, they
have befitted, and will befit, the one who uttered them. We ask the
surpassing grace of God to sanctify you from those illusory indications.
It is not worthy that all should pass in disputation, warfare, mischief,
and forced interpretations for the sake of selfish desire. It is
astonishing that detachment has come to be established by words and
claims. Thus, a person who has committed every conceivable forbidden
act---so that even now, in every land, a trace of that tree of
opposition remains---has nevertheless been regarded as detached. Beyond
this, a person who has manifested every indecency, and whose fire of
greed burns so fiercely that he has run from here to Istanbul for a
morsel of bread---this has actually been manifested by him---need only
meet some of those given to delusions, utter one long sigh, shake his
head once, and say two words professing renunciation of the world, to be
reckoned the foremost of the detached.

\subsection{31}\label{section-526}

\begin{otherlanguage}{persian}

آلت حرب و حفظ مشركين ذكر موت و انقطاع از دنيا است و بتوهّم خوب تلقين
نموده‌اند بشأنيكه در صدرت نقش بسته و لكن منقطعين را نفحات ديگر و عرف ديگر
است حنظل اگر دعوى شكرى نمايد اهل ذائقه نپذيرند هر قولى را فعل لازم لازال
نفوسيكه بر حقّ اعتراض نموده‌‌اند بامثال اين كلمات ناطق بوده‌اند چنانچه نفس
مقدّسى كه در بلايائى بوده كه واللّه الّذى لا اله الّا هو كه اگر اين عباد
گمان مينمودند كه محفوظ ماند شما نسبت بدنيا داده‌ايد و نفسيكه مادام عمر
بكمال عيش و راحت مشغول بوده او را منقطع ناميده‌ايد فويل للّذين يتكلّمون
باهوآء انفسهم و لايعرفون

\end{otherlanguage}

Ālat-e ḥarb va ḥefẓ-e moshrekīn dhekr-e mawt va enqeṭāʿ az donyā ast va
be-tavahhom khūb talqīn namūde-and be-shaʾnī ke dar ṣadrat naqsh baste
va lāken monqaṭeʿīn rā nafaḥāt-e dīgar va ʿarf-e dīgar ast ḥanẓal agar
daʿvā-ye shekarī namāyad ahl-e dhāʾeqe napadhīrand har qawlī rā feʿl
lāzem lāzāl nofūsī ke bar Ḥaqq eʿterāż namūde-and be-amthāl-e īn kalemāt
nāṭeq būde-and chenānche nafs-e moqaddasī ke dar balāyāʾī būde ke
wa-Allāh alladhī lā ilāha illā Huwa ke agar īn ʿebād gomān mīnamūdand ke
maḥfūẓ mānad shomā nesbat be-donyā dāde-īd va nafsī ke mādām-e ʿomr
be-kamāl-e ʿeysh va rāḥat mashghūl būde ū rā monqaṭeʿ nāmīde-īd
fa-waylun li-lladhīna yatakallamūna bi-ahwāʾi anfusihim wa-lā yaʿrifūn.

The idolaters' weapon of warfare and defense is talk of death and
detachment from the world. Through delusion they have instilled this so
effectively that it has become impressed upon your heart. Yet those who
are detached have another breath and another fragrance. If colocynth
claims to be sugar, those with a sense of taste will not accept it.
Every word requires a deed. Those who have objected to God have always
uttered words such as these. Thus, you have attributed worldliness to a
sanctified Soul who has suffered such afflictions that---by God, besides
whom there is no other god---these servants could scarcely have thought
He would survive; yet a person who has spent his entire life in the
utmost pleasure and ease you have called detached. Woe, then, to those
who speak according to their own desires and do not understand.

\subsection{32}\label{section-527}

\begin{otherlanguage}{persian}

و ديگر ندانستم مقصود از تأويلات چه چيز است كه ذكر نموده مأوّل آنانند كه
نديده ميگويند و نرسيده حكم مينمايند بفضل اللّه آنچه ذكر نموده‌ايم از
محكمات الهيّه بوده است و از مسلّمات مطالب شما و ملأ بيانيّه بشأنيكه نازل
نشده بيان مگر آنكه جميع را امر فرموده منصوصاً من غير تأويل كه اگر نفسى
بآيات الهى ظاهر شود احدى معترض او نشود و بشأنى مبالغه فرموده‌اند در اين
فقره كه ايمان و ايقان كل را منوط باقرار نفس ظهور نموده‌اند و تو اين كلمه
محكمه الهيّه را متشابه و تأويل زوركى ناميده اگر اين مطلب از بيان برداشته
شود ديگر بچه برهان اثبات ايمان مدّعيان ميشود فبئس ما ظنّوا فى انفسهم

\end{otherlanguage}

Va dīgar nadānestam maqṣūd az taʾvīlāt che chīz ast ke dhekr namūde
moʾavvel ānānand ke nadīde mīgūyand va naresīde ḥokm mīnamāyand
be-faẓl-e Allāh ānche dhekr namūde-īm az moḥkamāt-e elāhīye būde ast va
az mosallamāt-e maṭāleb-e shomā va malaʾ-e bayānīye be-shaʾnī ke nāzel
nashode Bayān magar ānke jamīʿ rā amr farmūde manṣūṣan min ghayri taʾwīl
ke agar nafsī be-āyāt-e elāhī ẓāher shavad aḥadī moʿtareż-e ū nashavad
va be-shaʾnī mobāleghe farmūde-and dar īn feqre ke īmān va īqān-e koll
rā manūṭ be-eqrār-e nafs-e ẓohūr namūde-and va to īn kaleme-ye
moḥkame-ye elāhīye rā motashābeh va taʾvīl-e zūrakī nāmīde agar īn
maṭlab az Bayān bardāshte shavad dīgar be-che borhān ethbāt-e īmān-e
moddaʿīyān mīshavad fa-biʾsa mā ẓannū fī anfusihim.

Furthermore, I did not understand what you meant by the interpretations
you mentioned. Those who interpret are those who speak without having
seen and pass judgment without having attained. By the grace of God,
what we have stated belongs to the unequivocal divine pronouncements and
to the accepted tenets of you and the people of the Bayán. Indeed, the
Bayán was revealed for no other purpose than to command everyone,
explicitly and without interpretation, that if a Soul should appear with
divine verses, no one should object to Him. So emphatically has this
point been stressed that the faith and certitude of all have been made
dependent upon their acknowledgment of the very Manifestation. Yet you
have called this unequivocal divine statement ambiguous and a forced
interpretation. If this principle were removed from the Bayán, by what
proof could the faith of those who claim it be established? How evil,
then, is what they have supposed within themselves.

\subsection{33}\label{section-528}

\begin{otherlanguage}{persian}

بارى چه زود تحريف نمودند كلمات الهى را بعضى از هياكل جعليّه جعليّه سبب
شده‌اند كه روايح رحمن را از اهل امكان منقطع نموده‌اند فرق اين است كه اين
عباد ديده ميگويند و شما بوهم انشاء اللّه اميدواريم كه هوى را از خدا تميز
دهى اگر منصف باشى خود شهادت ميدهى كه هميشه آنچه واقع شده بر خلاف مشتهيات
نفسانيّه بوده چه كه نفس امّاره از راحت و آسایش خود نگذشته و نخواهد گذشت
نفسى بجميع آيات الهيّه ظاهر و لازال در دست اعداء مبتلی و ناطق است بآنچه
در كلّ الواح مذكور مع‌ذلك ما ظهر من عنده را بنفس و هوى تعبير نموده صد‌هزار
جان مقرّبين فداى اين هوى كه بهبوبش نفحات مسكيّه از رضوان احديّه ساطع و از
حركت آن انفس ساكنه شوقاً للقآء اللّه مهتزّ و متحرّك و از سكونش افئده مضطربه
ساكن و مطمئن و از حرارتش نار الهى در سدره‌هاى وجود انسانى مشتعل و بندآء
انّه لا اله الّا هو ناطق

\end{otherlanguage}

Bārī che zūd taḥrīf namūdand kalemāt-e elāhī rā baʿżī az hayākel-e
jaʿlīye-ye jaʿlīye sabab shode-and ke ravāyeḥ-e Raḥmān rā az ahl-e emkān
monqaṭeʿ namūde-and farq īn ast ke īn ʿebād dīde mīgūyand va shomā
be-vahm in shāʾa Allāh omīdvārīm ke havā rā az Khodā tamīz dehī agar
monṣef bāshī khvod shahādat mīdehī ke hamīshe ānche vāqeʿ shode bar
khelāf-e moshtahayāt-e nafsānīye būde che ke nafs-e ammāre az rāḥat va
āsāyesh-e khvod nagodhashte va nakhvāhad godhasht nafsī be-jamīʿ-e
āyāt-e elāhīye ẓāher va lāzāl dar dast-e aʿdā mobtalā va nāṭeq ast
be-ānche dar koll-e alvāḥ madhkūr maʿa dhālek mā ẓahara min ʿindihi rā
be-nafs va havā taʿbīr namūde ṣad-hezār jān-e moqarrabīn fedā-ye īn havā
ke be-hobūbash nafaḥāt-e meskīye az Reżvān-e aḥadīye sāṭeʿ va az
ḥarakat-e ān anfos-e sākene shawqan li-liqāʾi Allāh mohtazz va
motaḥarrek va az sokūnash afʾede-ye moẓṭarebe sāken va moṭmaʾenn va az
ḥarāratash nār-e elāhī dar sedre-hā-ye vojūd-e ensānī moshtaʿel va
be-nedā-ye innahu lā ilāha illā Huwa nāṭeq.

How quickly they have distorted the words of God! Certain counterfeit,
counterfeit figures have caused the fragrances of the All-Merciful to be
cut off from the people of the world. The difference is that these
servants speak from what they have seen, while you speak from
imagination. God willing, we hope that you will distinguish desire from
God. If you are fair, you yourself will testify that everything that has
occurred has always run counter to selfish desires, for the self that
commands evil has never relinquished its own comfort and ease, nor will
it ever do so. A Soul has appeared with all the divine signs, has
continually suffered at the hands of enemies, and speaks what is
recorded in all the Tablets. Nevertheless, you have construed what has
appeared from Him as self and desire. May a hundred thousand lives of
those brought near be sacrificed for this desire, at whose wafting musky
fragrances spread from the Paradise of oneness; through whose movement
still souls are stirred and set in motion, longing to meet God; through
whose stillness troubled hearts become calm and assured; and through
whose heat the divine fire is kindled in the trees of human existence,
voicing the call, ``There is no god but Him.''

\subsection{34}\label{section-529}

\begin{otherlanguage}{persian}

امیدواريم كه شئونات الهيّه و ظهورات ربوبيّه را هواى نفسيّه ندانى و بشطر
انصاف ناظر گردى

\end{otherlanguage}

Omīdvārīm ke shoʾūnāt-e elāhīye va ẓohūrāt-e robūbīye rā havā-ye nafsīye
nadānī va be-shaṭr-e enṣāf nāẓer gardī.

We hope that you will not regard divine acts and the manifestations of
lordship as selfish desire, and that you will turn your gaze toward
fairness.

\subsection{35}\label{section-530}

\begin{otherlanguage}{persian}

و ديگر نوشته بودى كه بر شما مشتبه شده دلم از اين ميسوزد كه ساليان سال در
ميان كهنه عرفا بوده بقول عوام با ارانده قوم بسر برده باز امر مشتبه شود
انتهى

\end{otherlanguage}

Va dīgar neveshte būdī ke bar shomā moshtabeh shode delam az īn mīsūzad
ke sāliyān-e sāl dar mīyān-e kohne ʿorafā būde be-qawl-e ʿavāmm bā
{[}corrupted text{]} qawm be-sar borde bāz amr moshtabeh shavad intahā.

You had also written, ``You have been mistaken. My heart burns at this:
after spending year upon year among seasoned mystics and, as ordinary
people say, living with the {[}corrupted original text{]} of the people,
one should still be mistaken about the matter.'' End of quotation.

\subsection{36}\label{section-531}

\begin{otherlanguage}{persian}

آنقدر بدانيد كه در مقام يقين شبهه راه ندارد و اگر بگوئيد عرفان ظهور بديع
براين عبد مشتبه شده فاحضر لتشهد و اگر بگوئيد كه در آيات منزله اين اشتباه
واقع فاحضر لتسمع و ترى و اگر بگوئيد در عمل واقع كل ذى عمل عمله منوط
بامره و معلّق باذنه مقام حقّ‌اليقين مقدّس از ظنون و اوهام عالمين و مغلّين
معلوم است كه اينكلمات از تو نبوده چه كه باين مقدار گمان نميرود كه متابعت
ظنون و اوهام نمائى

\end{otherlanguage}

Ānqadr bedānīd ke dar maqām-e yaqīn shobhe rāh nadārad va agar begūʾīd
ʿerfān-e ẓohūr-e badīʿ bar īn ʿabd moshtabeh shode fa-ḥḍur li-tashhad va
agar begūʾīd ke dar āyāt-e monzale īn eshtebāh vāqeʿ fa-ḥḍur li-tasmaʿ
wa-tarā va agar begūʾīd dar ʿamal vāqeʿ kullu dhī ʿamalin ʿamaluhu
manūṭun bi-amrihi wa-muʿallaqun bi-idhnihi maqām-e ḥaqq al-yaqīn
moqaddas az ẓonūn va awhām-e ʿālamīn va moghellīn maʿlūm ast ke īn
kalemāt az to nabūde che ke be-īn meqdār gomān nemīravad ke motābeʿat-e
ẓonūn va awhām namāʾī.

Know at least this much: doubt has no place in the station of certitude.
If you say that this servant has been mistaken in recognizing the
wondrous Manifestation, then come and witness. If you say that this
mistake concerns the revealed verses, then come, hear, and see. If you
say that it concerns deeds, the deeds of everyone who acts are dependent
upon His command and contingent upon His permission. The station of true
certitude is sanctified above the conjectures and imaginings of the
peoples of the world and those who harbor malice. It is clear that these
words did not originate with you, for it is not supposed that you would
follow conjectures and imaginings to such a degree.

\subsection{37}\label{section-532}

\begin{otherlanguage}{persian}

و جميع اين امور كه واقع شده از قبل در الواح منزله مذكور و از آن گذشته در
سنۀ اوّل ورود اين ارض جمال ابهى در قرب مراديّه در بيتى ساكن و اين عباد هم
در بيتى ديگر روزى طرف عصر از حرم بيرون تشريف آوردند و جميع عباد مهاجرين
در خدمتشان قائم از جمله اين عبد و حاجى ميرزا احمد و سیّد محمّد و سايرين
فرمودند اليوم امر غريبی‌ مشاهده شده اين عباد منتظر كه چه واقع شده مقدار
عشرة دقايق او ازيد سكوت فرمودند و بعد باين عباد توجّه نموده فرمودند كه
اليوم حين فجر طيرى بر غصنى از اغصان شجره بيت جالس و باينكلمه ناطق محمّد
آمد و بلا آمد و كرّرت الطّير هذه الكلمة ثلثة مرّات

\end{otherlanguage}

Va jamīʿ-e īn omūr ke vāqeʿ shode az qabl dar alvāḥ-e monzale madhkūr va
az ān godhashte dar sane-ye avval-e vorūd-e īn arż Jamāl-e Abhā dar
qorb-e Morādīye dar baytī sāken va īn ʿebād ham dar baytī dīgar rūzī
ṭaraf-e ʿaṣr az ḥaram bīrūn tashrīf āvardand va jamīʿ-e ʿebād-e
mohājerīn dar khedmatashān qāʾem az jomle īn ʿabd va Ḥājī Mīrzā Aḥmad va
Seyyed Moḥammad va sāyerīn farmūdand al-yawm amr-e gharībī moshāhede
shode īn ʿebād montaẓer ke che vāqeʿ shode meqdār-e ʿasharata daqāyeq aw
azyad sokūt farmūdand va baʿd be-īn ʿebād tavajjoh namūde farmūdand ke
al-yawm ḥīn-e fajr ṭayrī bar ghoṣnī az aghṣān-e shajare-ye bayt jāles va
be-īn kaleme nāṭeq Moḥammad āmad va balā āmad wa-karrarat al-ṭayr
hādhihi al-kalima thalātha marrāt.

All these things that have occurred were mentioned beforehand in the
revealed Tablets. Beyond this, in the first year after arriving in this
land, the Most Glorious Beauty was residing in a house near the
Murádíyyih, while these servants were in another house. One afternoon He
came out of the inner quarters, and all the emigrant servants stood in
His presence, including this servant, Ḥájí Mírzá Aḥmad, Siyyid Muḥammad,
and others. He said, ``Today a strange thing was witnessed.'' These
servants waited to hear what had happened. He remained silent for ten
minutes or more, then turned to these servants and said, ``Today at
dawn, a bird was perched on one of the branches of the tree in the
house, uttering these words: `Muḥammad has come, and calamity has come.'
And the bird repeated these words three times.''

\subsection{38}\label{section-533}

\begin{otherlanguage}{persian}

هيچ يك از اين عباد گمان محمّد موجود را نمينمود خود سيّد محمّد هم چنين امرى
درباره خود محال ميدانست و اين عباد بتأويلات مشغول شديم و از جمله چند
نفوس در اطراف بودند كه باين اسم موسوم بالأخره موقن شديم كه يكى از آنها
بايد در اين ارض بيايد و فسادى از او ظاهر شود و بعد از انقضاء دو سنه
كامله ظاهر شد آنچه ظاهر شد در آن حين معلوم شد كه مقصود از محمّد مذكور كه
بوده واللّه واللّه افعال و اعمالش قابل ذكر نه و بحقّ خودشان كه اگر كلمه
بغير آنچه واقع شده ذكر نموده‌ام آنچه واقع شده بعينه همين است كه ذكر شد و
همچه ندانند كه ذكر اين فقره بجهت اظهار معجزات بوده لا فو‌ربّ العالمين چه
كه موقنيم باينكه كسانيكه در صدرشان غلّ حقّ بوده بهيچ حجّت و برهانى موقن
نشده و نخواهند شد چنانچه اخبار داده شده و ان يروا كلّ آية لن يؤمنوا بها و
لكن مقصود آنكه احوال او معلوم بوده فى لوح ما غادر فيه ذرّة من اعمال
الخلايق جميعاً

\end{otherlanguage}

Hīch yek az īn ʿebād gomān-e Moḥammad-e mawjūd rā nemīnamūd khvod-e
Seyyed Moḥammad ham chenīn amrī dar bāre-ye khvod moḥāl mīdānest va īn
ʿebād be-taʾvīlāt mashghūl shodīm va az jomle chand nofūs dar aṭrāf
būdand ke be-īn esm mawsūm be-l-ākhare mūqen shodīm ke yekī az ānhā
bāyad dar īn arż bīyāyad va fasādī az ū ẓāher shavad va baʿd az
enqeżā-ye do sane-ye kāmele ẓāher shod ānche ẓāher shod dar ān ḥīn
maʿlūm shod ke maqṣūd az Moḥammad-e madhkūr ke būde wa-Allāh wa-Allāh
afʿāl va aʿmālash qābel-e dhekr na va be-ḥaqq-e khvodashān ke agar
kaleme be-ghayr-e ānche vāqeʿ shode dhekr namūde-am ānche vāqeʿ shode
be-ʿayneh hamīn ast ke dhekr shod va hamche nadānand ke dhekr-e īn feqre
be-jehat-e eẓhār-e moʿjezāt būde lā fa-wa-Rabbi al-ʿālamīn che ke
mūqenīm be-īnke kasānī ke dar ṣadrashān ghell-e Ḥaqq būde be-hīch ḥojjat
va borhānī mūqen nashode va nakhvāhand shod chenānche akhbār dāde shode
wa-in yaraw kulla āyatin lan yuʾminū bihā va lāken maqṣūd ānke aḥvāl-e ū
maʿlūm būde fī lawḥin mā ghādara fīhi dharratan min aʿmāli al-khalāʾiq
jamīʿan.

None of these servants thought of the Muḥammad who was present. Siyyid
Muḥammad himself considered such a thing impossible concerning himself.
We servants began to offer interpretations. Among others, there were
several people in the surrounding regions who bore this name. Eventually
we became certain that one of them must come to this land and cause some
disturbance. After two full years had passed, what came to pass came to
pass. It then became clear who was meant by the Muḥammad mentioned. By
God, by God, his deeds and actions are not fit to mention! By His own
truth, if I have stated a word other than what occurred---what occurred
is precisely what has been recounted. Nor should they think that this
incident has been mentioned to display miracles. No, by the Lord of the
worlds! For we are certain that those whose hearts harbor malice toward
God have never become convinced by any proof or evidence, nor will they
ever do so, as has been foretold: ``Even if they see every sign, they
will not believe in it.'' The point, however, is that his circumstances
were known in a Tablet in which He has left out not a particle of the
deeds of all created beings.

\subsection{39}\label{section-534}

\begin{otherlanguage}{persian}

اين عبارت تو بسيار تحيّر افزود چه كه نفسيكه مردود اين طايفه بوده و قابل
سخن گفتن نبوده و نيست بر تو چنان امر را مشتبه نموده كه حال در تيه شبهات
سالكى و در باديه‌هاى ظنونات ساير چه كه نيكو است كه اين پر هاى عاريتى را
بيفكنى و بپرهاى معنوى در هواهاى عزّ قدس صمدانى پرواز كنى و بفضاهاى خوش
روحانى درائى و تو كه سالها خود را از رنود ميشمردى آخر يك هبوب قواصف
ظنونيّه و وهميّه چنان ربوده و حركت داده كه وصف آن ممكن نه ديگر حقّ عالم است
كه بكجا مقرّ دهد علمه عند ربّی‌ و ما انا الّا عبد منيب

\end{otherlanguage}

Īn ʿebārat-e to besyār taḥayyor afzūd che ke nafsī ke mardūd-e īn ṭāyefe
būde va qābel-e sokhan goftan nabūde va nīst bar to chenān amr rā
moshtabeh namūde ke ḥāl dar tīh-e shobohāt sālekī va dar bādīye-hā-ye
ẓonūnāt sāyer che ke nīkū ast ke īn par-hā-ye ʿārīyatī rā bīafkanī va
be-par-hā-ye maʿnavī dar havā-hā-ye ʿezz-e qods-e ṣamadānī parvāz konī
va be-fażā-hā-ye khvosh-e rūḥānī darāʾī va to ke sālhā khvod rā az ronūd
mīshomordī ākhar yek hobūb-e qavāṣef-e ẓonūnīye va vahmīye chenān robūde
va ḥarakat dāde ke vaṣf-e ān momken na dīgar Ḥaqq ʿālem ast ke be-kojā
maqarr dehad ʿilmuhu ʿinda rabbī wa-mā anā illā ʿabdun munīb.

This statement of yours greatly increased my astonishment, for a person
who was rejected by this community, and who was never and is not now fit
to be spoken to, has so confused the matter for you that you now wander
in the wilderness of doubts and roam the deserts of conjecture. It would
be well for you to cast off these borrowed wings, fly on spiritual wings
in the skies of the glory of divine holiness, and enter the delightful
expanses of the spirit. You, who for years counted yourself among the
worldly-wise, have at last been so swept away and shaken by a single
gust of the tempests of conjecture and imagination that it cannot be
described. God alone knows where He will finally bring you to rest.
Knowledge of it is with my Lord, and I am only a servant who turns to
Him.

\subsection{40}\label{section-535}

\begin{otherlanguage}{persian}

اميدواريم كه از حبال وهميّه ظنونيّه بحبل محكمه شديده متمسّك شوى و از بيت ظنّ
و گمان بمدينه طيّبه تمكين و يقين و اطمينان وارد گردى قسم بحقّ كه بسيار دلم
بر شما سوخت و ميسوزد كه بعد از زحمتهاى لا يحصى بگير قول بيابان وهم و هوى
افتاده‌ايد نسئل اللّه بأن يخلّصك و يشربك ما يطهّرك به عمّا سمعت و يجعلك مقدّساً
عن دونه و يستقرك فى ظلّ عنايته و انّه هو ارحم الرّاحمين

\end{otherlanguage}

Omīdvārīm ke az ḥebāl-e vahmīye-ye ẓonūnīye be-ḥabl-e moḥkame-ye shadīde
motamassek shavī va az bayt-e ẓann va gomān be-madīne-ye ṭayyebe-ye
tamkīn va yaqīn va eṭmīnān vāred gardī qasam be-Ḥaqq ke besyār delam bar
shomā sūkht va mīsūzad ke baʿd az zaḥmat-hā-ye lā yuḥṣā be-gīr-e qawl-e
bīyābān-e vahm va havā oftāde-īd nasʾalu Allāha bi-an yukhalliṣaka
wa-yushribaka mā yuṭahhiruka bihi ʿammā samiʿta wa-yajʿalaka muqaddasan
ʿan dūnihi wa-yastaqirraka fī ẓilli ʿināyatihi wa-innahu huwa arḥamu
al-rāḥimīn.

We hope that you will turn from the cords of imagination and conjecture
to grasp the firm and mighty cord, and leave the house of doubt and
supposition to enter the goodly city of steadfastness, certitude, and
assurance. I swear by God that my heart has burned and still burns
deeply for you, for after countless hardships you have become ensnared
by the speech of the wilderness of illusion and desire. We ask God to
deliver you, to give you to drink what will purify you of what you have
heard, to sanctify you from all besides Him, and to establish you
beneath the shadow of His loving care. He is indeed the most merciful of
the merciful.

\subsection{41}\label{section-536}

\begin{otherlanguage}{persian}

و ديگر نوشته بودى كه دلم ميسوزد كه بجهت بعد زمان از بعثت پيغمبر امثال
ما‌ها باديان تاويلى متمسّك بوده چندى را خيال خوش كرده خون جگر ميخورديم
بخيال آنكه از دين تاويلى زوركى فارغ شده حال بفرمايش شما بقول مردم اوّل
پياله و درد بايد باز داخل بدين تأويلى شد هنوز محكمات بيان رواج نگرفته و
ظاهر نشده بدم متشابهات بايد چسبيد انتهى

\end{otherlanguage}

Va dīgar nivishtih būdī kih dilam mīsūzad kih bi-jihat-i buʿd-i zamān az
biʿthat-i Payghambar amthāl-i māhā bi-adyān-i taʾvīlī mutamassik būdih
chandī rā khiyāl khush kardih khūn-i jigar mīkhurdīm bi-khiyāl-i ānkih
az dīn-i taʾvīlī-yi zūrakī fārigh shudih ḥāl bi-farmāyish-i shumā
bi-qawl-i mardum avval-i piyālih va durd bāyad bāz dākhil bi-dīn-i
taʾvīlī shud hanūz muḥkamāt-i Bayān ravāj nagiriftih va ẓāhir nashudih
bi-dum-i mutashābihāt bāyad chasbīd. Intahā.

You had also written: ``My heart burns because, owing to the long
interval since the Prophet's mission, people like us have clung to
religions based on interpretation. For a time we consoled ourselves with
pleasant imaginings, suffering bitter anguish, thinking we had escaped a
religion of forced interpretations. Now, according to what you say---as
people put it, `the first cup, and already the dregs'---we must enter a
religion of interpretation once again. Before the unambiguous statements
of the Bayán have gained currency or become manifest, we must cling to
the tails of its ambiguous statements.'' End of quotation.

\subsection{42}\label{section-537}

\begin{otherlanguage}{persian}

آنقدر معلوم بوده كه مأوّلين بهواى نفس و اهل ظنون و اوهام لازال بر مقرّ خود
ساكن و مستريح اگر صدهزار تصريح نازل شود ابداً برشحى از آن فايز نگردند
چنانچه مشاهده ميشود بحر تصريح موّاج و كلّ در بركهء وهم و تأویل مجتمع و لكن
از براى تأويل مراتب ما لانهايه بوده يك تأويل برضاى حقّ بوده و يك تأويل
اعتراضاً على الحقّ از براى كلمات الهيّه تأويلات لانهايه و لكن احدى بآن مطّلع
نه الّا اللّه و اليوم بنصّ نقطه بيان روح ما سواه فداه حرام است بر مستظلين
شجره بيان كه حرفى از كلمات اللّه را تأويل نمايند و يا تفسير كنند چه كه
احدى مطّلع نه مگر نفس ظهور اين تأويل كه شأن خلق نبوده

\end{otherlanguage}

Ānqadr maʿlūm būdih kih muʾavvilīn bi-havā-yi nafs va ahl-i ẓunūn va
avhām lāzāl bar maqarr-i khud sākin va mustarīḥ agar ṣad-hizār taṣrīḥ
nāzil shavad abadan bi-rashḥī az ān fāʾiz nagardand chunānchih
mushāhidih mīshavad baḥr-i taṣrīḥ mavvāj va kull dar birkih-yi vahm va
taʾvīl mujtamiʿ va lākin az barā-yi taʾvīl marātib-i mā lā-nihāyih būdih
yik taʾvīl bi-riżā-yi Ḥaqq būdih va yik taʾvīl iʿtirāḍan ʿalā al-Ḥaqq az
barā-yi kalimāt-i ilāhiyyih taʾvīlāt-i lā-nihāyih va lākin aḥadī bi-ān
muṭṭaliʿ nah illā Allāh va al-yawm bi-naṣṣ-i Nuqṭih-yi Bayān rūḥ-i mā
sivāh fidāh ḥarām ast bar mustaẓillīn-i Shajarih-yi Bayān kih ḥarfī az
kalimāt-i Allāh rā taʾvīl namāyand va yā tafsīr kunand chih kih aḥadī
muṭṭaliʿ nah magar Nafs-i Ẓuhūr īn taʾvīl kih shaʾn-i khalq nabūdih.

This much has been clear: those who interpret according to their own
desires, and those given to conjecture and delusion, have always
remained settled and at ease in their own abode. Even if a hundred
thousand explicit declarations were sent down, they would never receive
so much as a drop from them. As can be seen, the ocean of explicit
declaration is surging, while all are gathered in the pool of delusion
and interpretation. Interpretation, however, has innumerable degrees:
one interpretation accords with God's good pleasure, while another is
made in objection to God. The divine words have endless interpretations,
but no one knows them except God. Today, according to the explicit text
of the Point of the Bayán---may the souls of all besides Him be
sacrificed for Him---it is forbidden for those sheltered beneath the
Tree of the Bayán to interpret or expound a single letter of the words
of God, for no one knows them except the Manifestation Himself. Such
interpretation has never been the prerogative of created beings.

\subsection{43}\label{section-538}

\begin{otherlanguage}{persian}

و امّا آن تأويلى كه تو ذكر نموده تأويلات انفس مشركه بوده كه لازال مردود
بوده و خواهد بود چنانچه اليوم تصريح را گذاشته‌اند و از نفس ظهور محتجب و
بظنون و اوهام در تفسير و تأويل كلمات قبلش مشغولند فويل لهم و بما هم
يعملون مقصود از جميع تأويلات و تفاسير و معانى كلمات و آيات آن بوده كه
معرفت حقّ حاصل شود و بعد از آنكه از نفس حقّ محتجب شده‌اند كلّ بمقام خود راجع
و آن نفوس باسفل السّافلين مستقرّ

\end{otherlanguage}

Va ammā ān taʾvīlī kih tu dhikr namūdih taʾvīlāt-i anfus-i mushrikih
būdih kih lāzāl mardūd būdih va khvāhad būd chunānchih al-yawm taṣrīḥ rā
gudhāshtih-and va az Nafs-i Ẓuhūr muḥtajib va bi-ẓunūn va avhām dar
tafsīr va taʾvīl-i kalimāt-i qablash mashghūland fa-waylun lahum va
bi-mā hum yaʿmalūn maqṣūd az jamīʿ-i taʾvīlāt va tafāsīr va maʿānī-yi
kalimāt va āyāt ān būdih kih maʿrifat-i Ḥaqq ḥāṣil shavad va baʿd az
ānkih az Nafs-i Ḥaqq muḥtajib shudih-and kull bi-maqām-i khud rājiʿ va
ān nufūs bi-asfal al-sāfilīn mustaqarr.

As for the interpretation you mentioned, it is the interpretations of
those who associate partners with God, which have always been and will
always be rejected. Today, for example, they have set aside explicit
declarations and remain veiled from the Manifestation Himself, occupying
themselves, through conjecture and delusion, with expounding and
interpreting His earlier words. Woe to them and to what they do! The
purpose of all interpretations, expositions, and meanings of words and
verses has been to attain the knowledge of God. Now that they have
become veiled from God Himself, all these return to their own station,
while those souls are lodged in the lowest depths.

\subsection{44}\label{section-539}

\begin{otherlanguage}{persian}

و اينكه نوشته بودى كه خون جگر ميخوريم و بخيال اينكه آسوده شده از دين
تأويلى زوركى فارغ شده نه واللّه خون جگر نخورده چه كه اگر خون جگر خورده
بودى البتّه بمقرّ اطهر وارد ميشدى و منظر اكبر را از مطلع بقر تميز ميدادى
اگر چه چنان معلوم و مستفاد ميشود كه بعضى از شياطين انسيّه اراده نموده‌اند
كه ما بين تو و اين پست‌ترين بريّه فساد نمايند و لكن حقّ شاهد و گواه است كه
اين عبد از كلمات تو اگر مخصوص باين عبد بود ولو هرچه بود نميرنجيدم و لكن
احترقت اكباد العالمين من اشارات كلماتك و تصريحات الفاظك لأنّ كلّها ترجع
بظنونكم الى مقرّ الّذى فواللّه الّذى لا اله الّا هو لو يطيرنّ هؤلآء الّذين
اتّخذتهم لنفسك محبوباً من دون اللّه بدوام اللّه لن يبلغن الى هوآء الّذى فيه
تهبّ نسايم اسم من اسمائه الّتى خلقت بقوله فكيف نفسه المقتدر العلى الحكيم و
لكنّ اللّه قد جعل ذيله مقدّساً عن نباح هؤلآء و قباعهم و عن مسّ ايدى هؤلآء
المشركين

\end{otherlanguage}

Va īnkih nivishtih būdī kih khūn-i jigar mīkhorīm va bi-khiyāl-i īnkih
āsūdih shudih az dīn-i taʾvīlī-yi zūrakī fārigh shudih nah va-Allāh
khūn-i jigar nakhurdih chih kih agar khūn-i jigar khurdih būdī al-battih
bi-maqarr-i aṭhar vā rid mīshudī va Manẓar-i Akbar rā az maṭlaʿ-i baqar
tamyīz mīdādī agar chih chunān maʿlūm va mustafād mīshavad kih baʿżī az
shayāṭīn-i insiyyih irādih namūdih-and kih mā bayn-i tu va īn
past-tarīn-i bariyyih fasād namāyand va lākin Ḥaqq shāhid va guvāh ast
kih īn ʿabd az kalimāt-i tu agar makhṣūṣ bi-īn ʿabd būd valaw har-chih
būd namīranjīdam va lākin iḥtaraqat akbādu al-ʿālamīn min ishārāti
kalimātika va taṣrīḥāti alfāẓika li-anna kullahā tarjiʿu bi-ẓunūnikum
ilā maqarri alladhī fa-wa-Allāhi alladhī lā ilāha illā Huwa law
yaṭīranna hāʾulāʾi alladhīna ittakhadhtahum li-nafsika maḥbūban min dūni
Allāh bi-dawāmi Allāh lan yablughanna ilā havāʾi alladhī fīhi tahubbu
nasāʾimu ismin min asmāʾihi allatī khuliqat bi-qawlihi fa-kayfa Nafsihi
al-Muqtadiri al-ʿAliyyi al-Ḥakīm va lākinna Allāha qad jaʿala dhaylahu
muqaddasan ʿan nubāḥi hāʾulāʾi va qubāʿihim va ʿan massi aydī hāʾulāʾi
al-mushrikīn.

As for your having written, ``We suffer bitter anguish, thinking that we
had found relief and escaped a religion of forced
interpretations''---no, by God, you have not suffered such anguish. Had
you done so, you would surely have entered the most pure abode and
distinguished the Most Great Vision from the place where the ox appears.
It does seem, however, that certain human devils have intended to sow
discord between you and this lowliest of creatures. Yet God bears
witness that, had your words been directed solely at this servant,
whatever they might have been, I would not have been offended. But the
hearts of all the worlds have burned at the implications of your words
and the explicit statements of your utterances, for in your suppositions
they all refer to the seat of Him who---by God, besides whom there is no
other god---even if those whom you have taken as your beloved instead of
God were to fly for as long as God endures, they would never reach the
air in which blow the breezes of one of His names, which were created by
His word, much less His own Self, the All-Powerful, the Most High, the
All-Wise. But God has sanctified the hem of His robe from their barking
and grunting and from the touch of the hands of these associators of
partners with God.

\subsection{45}\label{section-540}

\begin{otherlanguage}{persian}

بشنو ذكر اين عبد را و احتراز كن از آن نفوس فو‌اللّه بهم يتغيّر هوآء كلّ
مدينة و انّهم يدّعون الايمان بعلىّ من قبل بما نزلت عليه آيات اللّه العزيز
الكريم فلما نزلت مرّة اخرى كفروا بها الا لعنة اللّه على الكاذبين و
الظّالمين از روايح كلمات تو جز اعراض و اعتراض چيزى استشمام نشده و نميشود
لذا بسيار بر اين عبد صعب است تحرير در اينمقام و الّا فى‌الحقيقة اگر نفسى
يافت شود و از اشارات مشركين شبهه نمايد البتّه محبوبست شبهات خود را مذكور
دارد و اين عباد در اين ارض شغلى نداريم و منتظريم اگر بر نفسى شبهه‌ئی وارد
شود جواب آنرا على قدر شأننا و استطاعتنا و شأنه مذكور داريم كه شايد
بعنايت الهى قلبش از شبهات افكيّه و ظنونات وهميّه مقدّس شود و از تو همين قدر
متوقّعيم كه ببصر مقدّس در آنچه ذكر ميشود ملاحظه نمائى

\end{otherlanguage}

Bishnav dhikr-i īn ʿabd rā va iḥtirāz kun az ān nufūs fa-wa-Allāhi bihim
yataghayyaru havāʾu kulli madīnatin va innahum yaddaʿūna al-īmāna
bi-ʿAliyyin min qablu bi-mā nazalat ʿalayhi āyātu Allāhi al-ʿAzīzi
al-Karīm fa-lammā nazalat marratan ukhrā kafarū bihā alā laʿnatu Allāhi
ʿalā al-kādhibīna va al-ẓālimīn az ravāyiḥ-i kalimāt-i tu juz iʿrāż va
iʿtirāż chīzī istishmām nashudih va namīshavad lidhā bisyār bar īn ʿabd
ṣaʿb ast taḥrīr dar īn maqām va illā fī al-ḥaqīqih agar nafsī yāft
shavad va az ishārāt-i mushrikīn shubhih namāyad al-battih maḥbūb ast
shubuhāt-i khud rā madhkūr dārad va īn ʿibād dar īn arż shughlī nadārīm
va muntaẓirīm agar bar nafsī shubhih-ī vā rid shavad javāb-i ānrā ʿalā
qadri shaʾninā va istiṭāʿatinā va shaʾnihi madhkūr dārīm kih shāyad
bi-ʿināyat-i ilāhī qalbash az shubuhāt-i ifkiyyih va ẓunūnāt-i vahmiyyih
muqaddas shavad va az tu hamīn qadr mutavaqqiʿīm kih bi-baṣar-i muqaddas
dar ānchih dhikr mīshavad mulāḥiẓih namāʾī.

Listen to this servant's words and keep away from those souls. By God,
through them the air of every city is corrupted. They profess belief in
`Alí of old because the verses of God, the Almighty, the Generous, were
sent down to Him; yet when they were sent down once again, they
disbelieved in them. Surely God's curse rests upon the liars and the
wrongdoers! Nothing but rejection and objection has been, or is,
discernible in the scent of your words. It is therefore very difficult
for this servant to write on this subject. Otherwise, if someone should
truly be found who has developed doubts through the insinuations of
those who associate partners with God, it is certainly desirable that he
state his doubts. We servants have no occupation in this land, and we
wait, should any soul be troubled by a doubt, to give an answer
according to our station and ability and his station, so that perhaps,
through divine loving-kindness, his heart may be purified of false
doubts and delusive conjectures. All we expect of you is that you
consider what is said with a purified vision.

\subsection{46}\label{section-541}

\begin{otherlanguage}{persian}

از كلمات حضرت ابهى است روح ما سواه فداه كه ميفرمايند من لم يطهّر قلبه و
فؤاده و بصره و اذنه عمّا عنده لم يقدر ان يعرف ما نزل بالحقّ فى الواح عزّ
مبين لذا از خدا آمليم كه نسيمى از نسايم خوش رضون بر اهل امكان بفرستد تا
جميع را از غبار عصيان پاك و مقدّس نموده بشطر قدس رحمن متوجّه سازد و انّه
لمن دعاه قريب مجيب

\end{otherlanguage}

Az kalimāt-i Ḥaḍrat-i Abhā ast rūḥ-i mā sivāh fidāh kih mīfarmāyand man
lam yuṭahhir qalbahu va fuʾādahu va baṣarahu va udhunahu ʿammā ʿindahu
lam yaqdir an yaʿrifa mā nazala bi-al-ḥaqqi fī alwāḥi ʿizzin mubīn lidhā
az Khudā āmilīm kih nasīmī az nasāyim-i khush-i Riżvān bar ahl-i imkān
bifiristad tā jamīʿ rā az ghubār-i ʿiṣyān pāk va muqaddas namūdih
bi-shaṭr-i quds-i Raḥmān mutavajjih sāzad va innahu li-man daʿāhu
qarībun mujīb.

Among the words of His Holiness the All-Glorious---may the souls of all
besides Him be sacrificed for Him---is this utterance: ``Whoever has not
purified his heart, his inmost being, his sight, and his hearing of what
he possesses cannot recognize what has been sent down in truth in the
tablets of manifest glory.'' We therefore hope that God will send one of
the sweet breezes of Riḍván upon the people of the contingent world, to
cleanse and sanctify them all from the dust of disobedience and turn
them toward the holy presence of the All-Merciful. He is indeed near and
responsive to whoever calls upon Him.

\subsection{47}\label{section-542}

\begin{otherlanguage}{persian}

اى برادر من اين طيور هواى حبّ مظهر تفريد از خرمنهاى قدس توحيد دانه
برداشته و از نغمات قدس معنوى استماع نموده‌اند مشكل بدام اهل دنيا درآيند و
يا از وساوس كلمات معرضين و مشركين از ربّ العالمين ممنوع شوند ومحروم گردند
قاصدين اينمطلب را جز حسرت بيفزايد و جز زحمت و كسالت نصيب نداشته و
نخواهند داشت چه كه بفضل اللّه در رياض مكاشفه سايريم و در انجمن مشاهده
حاضر از وهم و تقليد گذشته‌ايم و در مقرّ شهود و مشاهده آرميده‌ايم بدنيا و
اهل آن كارى نداشته و نداريم چنانچه فعل ما بر دين ما گواه است از براى
شهريّه و لقمهء نان بدرها ندويده‌ايم و زن بسرايه نفرستاده‌ايم و بحقّ افترا
نبسته‌‌ایم فو‌اللّه الّذى لا اله الّا هو كه اگر بهمان كذب‌هائيكه در آن ارض بشما
گفته و مفترياتيكه بحقّ نسبت داده مطّلع شوى البتّه سر بصحرا گذارى و از كلّ من
على الأرض منقطع شوى لو تكون من المنصفين

\end{otherlanguage}

Ay barādar-i man īn ṭuyūr-i havā-yi ḥubb-i Maẓhar-i Tafrīd az
kharman-hā-yi quds-i tawḥīd dānih bardāshtih va az naghamāt-i quds-i
maʿnavī istimāʿ namūdih-and mushkil bi-dām-i ahl-i dunyā dar-āyand va yā
az vasāvis-i kalimāt-i muʿriżīn va mushrikīn az Rabb al-ʿālamīn mamnūʿ
shavand va maḥrūm gardand qāṣidīn-i īn maṭlab rā juz ḥasrat biyafzāyad
va juz zaḥmat va kasālat naṣīb nadāshtih va nakhvāhand dāsht chih kih
bi-faḍl-i Allāh dar riyāż-i mukāshifih sāyirīm va dar anjuman-i
mushāhidih ḥāżir az vahm va taqlīd gudhashtih-īm va dar maqarr-i shuhūd
va mushāhidih ārāmīdih-īm bi-dunyā va ahl-i ān kārī nadāshtih va nadārīm
chunānchih fiʿl-i mā bar dīn-i mā guvāh ast az barā-yi shahrīyyih va
luqmih-yi nān bi-darhā nadavīdih-īm va zan bi-sarāyih nafiristādih-īm va
bi-Ḥaqq iftirā nabastih-īm fa-wa-Allāhi alladhī lā ilāha illā Huwa kih
agar bi-hamān kidḥ-hāʾīkih dar ān arż bi-shumā guftih va muftariyātīkih
bi-Ḥaqq nisbat dādih muṭṭaliʿ shavī al-battih sar bi-ṣaḥrā gudhārī va az
kulli man ʿalā al-arḍ munqaṭiʿ shavī law takūnu min al-munṣifīn.

O my brother! These birds of the heaven of love for the Manifestation of
Oneness have gathered grain from the sacred harvests of divine unity and
have listened to spiritual melodies of holiness. They are hardly likely
to fall into the snares of the worldly, or to be kept back and deprived
of the Lord of the worlds by the insinuating words of those who turn
away and associate partners with God. Those who seek this end will only
increase their own regret; nothing but trouble and weariness has been or
will be their portion. For, by the grace of God, we roam the gardens of
revelation and are present in the gathering of direct vision. We have
passed beyond delusion and imitation and have found repose in the abode
of witnessing and beholding. We have had, and have, no concern with the
world and its people, as our deeds bear witness to our faith. We have
not run from door to door for a monthly allowance and a morsel of bread,
nor sent a woman to the palace, nor fabricated lies against God. By God,
besides whom there is no other god, if you were to learn of the lies
they have told you in that land and the fabrications they have
attributed to God, you would surely flee to the wilderness and detach
yourself from everyone on earth, if you are fair-minded.

\subsection{48}\label{section-543}

\begin{otherlanguage}{persian}

و ديگر نوشته اوّل پياله و درد اين كلمه را ذكر نموده و لكن ندانسته در چه
محلّ بايد ذكر نمود نقطه بيان روح ما سواه فداه در جميع بيان تصريحاً
ميفرمايند بآنچه خلق شده ما بين سموات و ارض چه از كتب منزله و چه از صحف
قيّمه و چه از آيات بديعه و چه مظاهر تكوينيّه و چه مظاهر تشريعيّه از نفس
ظهور محتجب نمانيد و مع آنكه اخذ عهد ولايت اين ظهور امنع اقدس را از كلّ
فرموده بقوله عزّ ذكره سبحانك اللّهمّ فاشهد بانّى بذلك الكتاب قد اخذت عهد
ولايت من تظهرنّه عن كلّ شیء قبل عهد ولايتى و كفى بك و بمن آمن بآياتك علىّ
شهدآء و انّك انت حسبی‌ عليك توكّلت و انّك كنت على كلّ شیء حسيباً كلّ از جمالش
معرض و بقاتلش متمسّك

\end{otherlanguage}

Va dīgar nivishtih avval-i piyālih va durd īn kalimih rā dhikr namūdih
va lākin nadānistih dar chih maḥall bāyad dhikr namūd Nuqṭih-yi Bayān
rūḥ-i mā sivāh fidāh dar jamīʿ-i Bayān taṣrīḥan mīfarmāyand bi-ānchih
khalq shudih mā bayn-i samāvāt va arż chih az kutub-i munzalih va chih
az ṣuḥuf-i qayyimih va chih az āyāt-i badīʿih va chih maẓāhir-i
takvīniyyih va chih maẓāhir-i tashrīʿiyyih az Nafs-i Ẓuhūr muḥtajib
namānīd va maʿa ānkih akhdh-i ʿahd-i vilāyat-i īn Ẓuhūr-i Amnaʿ-i Aqdas
rā az kull farmūdih bi-qawlihi ʿazza dhikruhu subḥānaka Allāhumma
fa-shhad bi-annī bi-dhālika al-kitābi qad akhadhtu ʿahda vilāyati man
tuẓhirannahu ʿan kulli shayʾin qabla ʿahdi vilāyatī va kafā bika va
bi-man āmana bi-āyātika ʿalayya shuhadāʾa va innaka Anta ḥasbī ʿalayka
tavakkaltu va innaka kunta ʿalā kulli shayʾin ḥasīban kull az Jamālash
muʿriż va bi-qātilash mutamassik.

You also wrote, ``The first cup, and already the dregs.'' You have used
this expression without knowing where it should be used. Throughout the
Bayán, the Point of the Bayán---may the souls of all besides Him be
sacrificed for Him---explicitly declares: ``Do not let anything created
between the heavens and the earth---whether revealed books, upright
scriptures, wondrous verses, manifestations in creation, or
manifestations of religious law---veil you from the Manifestation
Himself.'' And although He has taken from all the covenant of allegiance
to this most unassailable, most holy Manifestation, saying---exalted be
His mention---``Glory be to You, O God! Bear witness that through that
Book I have taken from all things the covenant of allegiance to Him whom
You will manifest, before the covenant of allegiance to Myself. You and
those who believe in Your verses suffice as witnesses to Me. You,
indeed, are sufficient for Me. In You have I placed My trust, and You
have ever taken account of all things,'' all have turned away from His
Beauty and clung to His slayer.

\subsection{49}\label{section-544}

\begin{otherlanguage}{persian}

فواللّه الّذى لا اله الّا هو كه مشهود بود كه بر اين جمال مبارك چه وارد
ميآورند ملاحظه كن كه در همين فقره نقطه بيان روح ما سواه فداه چه مقدار از
نوحه و ندبه و حنين مستور است حقّ و جميع عباد را در اخذ عهد اين ظهور مبارك
شاهد و گواه گرفته كه شايد اهل بيان در ظهور نفس رحمن طغيان ننمايند و عمّا
نزل فى البيان حيا كنند مع‌‌ذلك وارد شده آنچه مشاهده ميشود و هنوز ایّامى
نگذشته و در آذان كلّ كلمات منزل بيان بوده و هنوز اثر دم مطهّر آن شهيد
مظلوم بر ارض باقى كه بنفس او و كينونت او و روح او وارد آورده‌اند آنچه را
كه هيچ ملّتى وارد نيآورده اينجا است مقام آنكه بگوئى اوّل پياله و درد حال

\end{otherlanguage}

Fa-wa-Allāhi alladhī lā ilāha illā Huwa kih mashhūd būd kih bar īn
Jamāl-i Mubārak chih vā rid mīāvarand mulāḥiẓih kun kih dar hamīn
faqarih Nuqṭih-yi Bayān rūḥ-i mā sivāh fidāh chih miqdār az nawḥih va
nudbih va ḥanīn mastūr ast Ḥaqq va jamīʿ-i ʿibād rā dar akhdh-i ʿahd-i
īn Ẓuhūr-i Mubārak shāhid va guvāh giriftih kih shāyad ahl-i Bayān dar
ẓuhūr-i Nafs-i Raḥmān ṭughyān nanamāyand va ʿammā nazala fī al-Bayān
ḥayā kunand maʿa dhālik vā rid shudih ānchih mushāhidih mīshavad va
hanūz ayyāmī nagudhashtih va dar ādhān-i kull kalimāt-i Munzil-i Bayān
būdih va hanūz athar-i dam-i muṭahhar-i ān Shahīd-i Maẓlūm bar arż bāqī
kih bi-Nafs-i Ū va Kaynūnat-i Ū va Rūḥ-i Ū vā rid āvardih-and ānchih rā
kih hīch millatī vā rid nayāvardih īnjā ast maqām-i ānkih bigūʾī avval-i
piyālih va durd ḥāl.

By God, besides whom there is no other god, it was plain what they would
inflict upon this Blessed Beauty. Consider how much lamentation,
mourning, and yearning is concealed in this very passage of the Point of
the Bayán---may the souls of all besides Him be sacrificed for Him. In
taking the covenant of this blessed Manifestation, He called upon God
and all His servants as witnesses, so that perhaps the people of the
Bayán would not rebel when the All-Merciful Himself appeared, and would
feel shame before what was sent down in the Bayán. Nevertheless, what
can now be seen has come to pass. Scarcely any time has elapsed; the
words of the Revealer of the Bayán still ring in everyone's ears, and
traces of the pure blood of that wronged Martyr still remain upon the
earth. Yet they have inflicted upon His own Self, His Being, and His
Spirit what no religious community has ever inflicted. Here is the place
for you to say, ``The first cup, and already the dregs.'' Now,

\subsection{50}\label{section-545}

\begin{otherlanguage}{persian}

ملاحظه نمائيد كه چه مقدار كلمات را مغاير و مخالف و غير موقع ذكر مينمائيد
نفسى را كه ميفرمايد اخذ ولايت او را قبل از اخذ عهد ولايت خود نموده‌ام
مع‌ذلك معرضين بيانيّه نفوس مقبله باين نفس رحمانيّه را اهل تأويل شمرده‌اند
نشكو بثنآء و حزننا الى اللّه فيما ورد عليه من عباده الّذين خلقوا بقوله كن
فيكون

\end{otherlanguage}

Mulāḥiẓih namāʾīd kih chih miqdār kalimāt rā mughāyir va mukhālif va
ghayr-i mawqiʿ dhikr mīnamāʾīd Nafsī rā kih mīfarmāyad akhdh-i vilāyat-i
Ū rā qabl az akhdh-i ʿahd-i vilāyat-i khud namūdih-am maʿa dhālik
muʿriżīn-i Bayāniyyih nufūs-i muqbilih bi-īn Nafs-i Raḥmāniyyih rā ahl-i
taʾvīl shumurdih-and nashkū bi-thanāʾin va ḥuznanā ilā Allāhi fīmā
varada ʿalayhi min ʿibādihi alladhīna khuliqū bi-qawlihi kun fa-yakūn.

consider how greatly you use expressions in contrary, conflicting, and
inappropriate ways. Of this Person He says, ``I took allegiance to Him
before taking the covenant of allegiance to Myself.'' Nevertheless,
those among the people of the Bayán who turn away have counted the souls
who turn toward this all-merciful Person among the proponents of
interpretation. With praise, we bring our complaint and sorrow to God
concerning what has befallen Him at the hands of His servants, who were
created by His word, ``Be,'' and it is.

\subsection{51}\label{section-546}

\begin{otherlanguage}{persian}

و ديگر مذكور نموده بودى كه هنوز محكمات بيان رواج نگرفته و ظاهر نشده بدم
متشابهات بايد چسبيد هيچ ذى عقل و درايتى چنين كلمه از لسانش جارى نميشود
متحيّرم چگونه راضى شده كلمات مشركين را باسم خود ذكر نمود شأن اهل حقّ نبوده
كه از مقام ادب تجاوز نمايند لذا اين عباد جواب بعضى از فقرات مذكوره‌ات را
بخدا وا ميگذاريم چه كه مظلوميّت محبوب حقّ بوده

\end{otherlanguage}

Va dīgar madhkūr namūdih būdī kih hanūz muḥkamāt-i Bayān ravāj
nagiriftih va ẓāhir nashudih bi-dum-i mutashābihāt bāyad chasbīd hīch
dhī-ʿaql va dirāyatī chunīn kalimih az lisānash jārī namīshavad
mutaḥayyiram chigūnih rāżī shudih kalimāt-i mushrikīn rā bi-ism-i khud
dhikr namūd shaʾn-i ahl-i Ḥaqq nabūdih kih az maqām-i adab tajāvuz
namāyand lidhā īn ʿibād javāb-i baʿżī az faqarāt-i madhkūrih-at rā
bi-Khudā vā mīgudhārīm chih kih maẓlūmiyyat maḥbūb-i Ḥaqq būdih.

You had also stated: ``Before the unambiguous statements of the Bayán
have gained currency or become manifest, we must cling to the tails of
its ambiguous statements.'' No one possessed of reason and discernment
would utter such words. I am astonished that you consented to state in
your own name the words of those who associate partners with God. It has
never been fitting for the people of God to transgress the bounds of
courtesy. We servants therefore leave the answer to some of the passages
you have written to God, for being wronged has ever been beloved of God.

\subsection{52}\label{section-547}

\begin{otherlanguage}{persian}

و امّا در محكم و متشابه كه ذكر نموده از اين بيان همچه مستفاد ميشود كه فرق
ميانه محكم و متشابه درست در نظرت نبوده و الّا موقن ميشدى بآنكه محكمات ما
نزل فى البيان باتمّها و اكملها ظاهر شده لذا مجملى در اينمقام ذكر ميشود
بدانكه از محكم و متشابه ما بين هر حزبی‌ ذكرى مذكور و تفصيلى موجود چنانچه
در كتب ثبت است و تفصيل آن موجب تطويل خواهد شد و لكن عند الحقّ آيات محكمات
آنكه مقصود الهى از نفس آيه مستفاد شود و در اثبات آن محتاج ببرهان و ذكر
ديگر نباشد حال انصاف دهيد در حقيقت اوّليّه آن آيه كدام است نزد هر ذى بصرى
مبرهن و واضح است كه آن آيه اعظم كه در اثبات خود محتاج بدونش نبوده و
نخواهد بود در حقيقت اوّليّه نفس ظهور بوده چه كه مادونش هر چه باشد و هر كه
باشد محتاج بامرش و لاينكر ذلك الّا كلّ جاهل مردود

\end{otherlanguage}

Va ammā dar muḥkam va mutashābih kih dhikr namūdih, az īn bayān hamchih
mustafād mīshavad kih farq-i miyānih-yi muḥkam va mutashābih durust dar
naẓarat nabūdih, va illā mūqin mīshudī bi-ānkih muḥkamāt-i mā nazala fī
al-Bayān bi-atammihā va akmalihā ẓāhir shudih. Lidhā mujmalī dar īn
maqām dhikr mīshavad. Bidānkih az muḥkam va mutashābih mā bayn-i har
ḥizbī dhikrī madhkūr va tafṣīlī mawjūd, chinānchih dar kutub thabt ast,
va tafṣīl-i ān mūjib-i taṭvīl khāhad shud. Va lākin ʿinda al-Ḥaqq āyāt-i
muḥkamāt ānkih maqṣūd-i ilāhī az nafs-i āyih mustafād shavad va dar
ithbāt-i ān muḥtāj bi-burhān va dhikr-i dīgar nabāshad. Ḥāl inṣāf dahīd,
dar ḥaqīqat-i avvalīyyih ān āyih kudām ast? Nazd-i har dhī baṣarī
mubarhan va vāḍiḥ ast kih ān āyih-yi aʿẓam kih dar ithbāt-i khud muḥtāj
bi-dūnash nabūdih va nakhāhad būd, dar ḥaqīqat-i avvalīyyih nafs-i ẓuhūr
būdih, chih kih mā dūnash har chih bāshad va har kih bāshad muḥtāj
bi-amrash, va lā yunkiru dhālika illā kull jāhil mardūd.

As for what you mentioned concerning the unambiguous and the ambiguous,
your statement suggests that the distinction between them was not clear
to you. Otherwise, you would have been certain that the unambiguous
passages of what was revealed in the Bayán have appeared in their
fullest and most perfect form. A brief account is therefore given here.
Know that every community has its own account and explanation of the
unambiguous and the ambiguous, as recorded in the books; to explain this
in detail would prolong the discussion. In the sight of the True One,
however, unambiguous verses are those in which the divine intention can
be understood from the verse itself, without requiring another proof or
statement to establish it. Now judge fairly: in the primary reality,
which is that verse? It is demonstrably clear to everyone possessed of
insight that, in the primary reality, that greatest Verse, which has
never needed and will never need anything other than itself to establish
its truth, is the Manifestation Himself. For everything and everyone
besides Him depends upon His command. None denies this but an ignorant
outcast.

\subsection{53}\label{section-548}

\begin{otherlanguage}{persian}

و اگر بگوئيد آيات محكمات اوامر الهيّه است كه در الواح نازل لاشكّ فى ذلك
انّا كلّ بها موقنون و مؤمنون و من فضله آملون و اين محكمات بنسبتها الى احكم
الحاكمين محكم بوده چنانچه مشهود است كه اگر نفسى بجميع آيات محكمات عارف
شود و در تصديق آن آيه اعظم اقلّ من آن توقّف نمايد كلّ اعمال او او را ثمرى
نبخشد و هباء بوده و خواهد بود و آن نفسيكه آيات او محكمات بوده و خواهد
بود ميفرمايد جميع بيان ورقيست از اوراق جنّت نزد آن ساذج وجود و در مقام
ديگر ميفرمايد ما نزل فى البيان منوط بامره و معلّق باذنه و مشيّته و در مقام
ديگر ميفرمايد مخاطباً لملأ البيان كه مباد ببيان و ما نزل فيه از ظهور
موعود محروم شويد و امثال اين بيانات محكمات زياده از حدّ احصا نازل و مع‌ذلك
نفسى را كه جميع محكمات باراده‌اش ظاهر و جميع انهار معارف و بيان از اصبعى
از اصابع مشيّتش جارى از او معرض و بگمان خود محكمات فهميده و ادراك
نموده‌اند

\end{otherlanguage}

Va agar bigūyīd āyāt-i muḥkamāt avāmir-i ilāhīyyih ast kih dar alvāḥ
nāzil, lā shakka fī dhālika, innā kull bihā mūqinūn va muʾminūn va min
faḍlihi āmilūn. Va īn muḥkamāt bi-nisbatihā ilā Aḥkam al-Ḥākimīn muḥkam
būdih, chinānchih mashhūd ast kih agar nafsī bi-jamīʿ-i āyāt-i muḥkamāt
ʿārif shavad va dar taṣdīq-i ān āyih-yi aʿẓam aqall min ān tavaqquf
namāyad, kull-i aʿmāl-i ū ū rā thamarī nabakhshad va habāʾ būdih va
khāhad būd. Va ān nafsī kih āyāt-i ū muḥkamāt būdih va khāhad būd
mīfarmāyad jamīʿ-i Bayān varaqīst az avrāq-i jannat nazd-i ān sādhij-i
vujūd. Va dar maqām-i dīgar mīfarmāyad mā nazala fī al-Bayān manūṭ
bi-amrihi va muʿallaq bi-idhnihi va mashīyyatihi. Va dar maqām-i dīgar
mīfarmāyad mukhāṭaban li-malaʾ al-Bayān kih mabād bi-Bayān va mā nazala
fīhi az ẓuhūr-i mawʿūd maḥrūm shavīd. Va amthāl-i īn bayānāt-i muḥkamāt
ziyādih az ḥadd-i iḥṣā nāzil, va maʿa dhālik nafsī rā kih jamīʿ-i
muḥkamāt bi-irādih-ash ẓāhir va jamīʿ-i anhār-i maʿārif va bayān az
iṣbaʿī az aṣābiʿ-i mashīyyatash jārī, az ū muʿriḍ va bi-gumān-i khud
muḥkamāt fahmīdih va idrāk namūdih-and.

If you say that the unambiguous verses are the divine commandments
revealed in the Tablets, there is no doubt of this: we are all certain
of them, believe in them, and hope for His grace. These unambiguous
passages are unambiguous by virtue of their relation to the Most Just of
Judges. It is evident that if someone were to know all the unambiguous
verses, yet hesitate for less than an instant to acknowledge that
greatest Verse, none of that person's deeds would bear fruit; they have
been and will remain mere dust. He whose verses have been and will
remain unambiguous says: ``The whole Bayán is a leaf among the leaves of
Paradise in the presence of that pure Essence of Being.'' Elsewhere He
says: ``What has been revealed in the Bayán depends upon His command and
is conditional upon His permission and His will.'' Elsewhere, addressing
the people of the Bayán, He says: ``Do not let the Bayán and what has
been revealed in it deprive you of the promised Manifestation.''
Unambiguous statements such as these have been revealed beyond all
reckoning. Nevertheless, they turn away from Him through whose will all
the unambiguous passages have appeared and from one finger of whose will
all the rivers of knowledge and utterance flow, while imagining that
they have understood and comprehended the unambiguous passages.

\subsection{54}\label{section-549}

\begin{otherlanguage}{persian}

حال چشم انصاف بگشا و اهل التباس و مظاهر خنّاس را بشناس يعنى نفوسى كه
امثال ان كلمات را بر تو القاء نموده‌اند بلى لازال نفوسى خود را باسم حقّ
ناميده و ضعفاى بيچاره را در تيه ظنون و اوهام مبتلا نموده و خود برياست
مشغول شده‌اند ليس هذا اوّل سهم ورد علينا و لا اوّل رمح ورد على محبوب
العالمين

\end{otherlanguage}

Ḥāl chashm-i inṣāf bigushā va ahl-i iltibās va maẓāhir-i khannās rā
bishinās, yaʿnī nufūsī kih amthāl-i ān kalimāt rā bar tu ilqā
namūdih-and. Balī, lā zāl nufūsī khud rā bi-ism-i Ḥaqq nāmīdih va
ḍuʿafā-yi bīchārih rā dar tīh-i ẓunūn va awhām mubtalā namūdih va khud
bi-riyāsat mashghūl shudih-and. Laysa hādhā awwal sahm warada ʿalaynā va
lā awwal rumḥ warada ʿalā Maḥbūb al-ʿĀlamīn.

Now open the eye of fairness and recognize the people of deception and
the embodiments of the insidious whisperer---that is, those who have
suggested such words to you. Indeed, there have always been people who
have called themselves by the name of the True One, cast the helpless
and weak into the wilderness of conjectures and delusions, and occupied
themselves with leadership. This is not the first arrow to strike Us,
nor the first spear to strike the Beloved of the worlds.

\subsection{55}\label{section-550}

\begin{otherlanguage}{persian}

و اگر بگوئى محكمات اوصاف منزله در بيان است كه مخصوص بعضى نفوس نازل شده
اوّلا آن اوصاف را تو عارف نبوده كه مقصود كيست و ثانى فو‌نفسه الحقّ كه كلّ
اشیآء را بمنتهى وصف وصف فرموده‌اند بشأنيكه جميع اسماء و صفات را بكلّ شیء
راجع نموده‌اند و اينمقام ما ترى فى خلق الرّحمن من تفاوت بوده ما سوايش خلق
او و كلّ در اوصاف شريك چه كه در منظر اكبر كلّ در صقع واحد بوده و خواهند
بود

\end{otherlanguage}

Va agar bigūʾī muḥkamāt awṣāf-i munazzalih dar Bayān ast kih makhṣūṣ-i
baʿḍī nufūs nāzil shudih, avvalan ān awṣāf rā tu ʿārif nabūdih kih
maqṣūd kīst, va thānī fa-wa-nafsihi al-ḥaqq kih kull-i ashyāʾ rā
bi-muntahā-yi vaṣf vaṣf farmūdih-and, bi-shaʾnī kih jamīʿ-i asmāʾ va
ṣifāt rā bi-kull-i shayʾ rājiʿ namūdih-and. Va īn maqām mā tarā fī khalq
al-Raḥmān min tafāwut būdih. Mā sivāyash khalq-i ū va kull dar awṣāf
sharīk, chih kih dar manẓar-i akbar kull dar ṣuqʿ-i vāḥid būdih va
khāhand būd.

If you say that the unambiguous passages are the descriptions revealed
in the Bayán specifically concerning certain people, first, you have not
understood whom those descriptions refer to. Second, by His own true
Self, He has described all things in the most exalted terms, to such an
extent that He has referred all names and attributes to every single
thing. This is the station of ``You see no disparity in the creation of
the All-Merciful.'' All besides Him are His creation, and all share in
these attributes, for in the Most Great Perspective all have been and
will remain on a single plane.

\subsection{56}\label{section-551}

\begin{otherlanguage}{persian}

و اينكه درباره بعضى آيات الهيّه بر حسب ظاهر نازل مقصود تشويق آن نفوس بوده
و از آن گذشته جميع بآيهء تجلّى حقّ راجع ماداميكه آن آيهء تجلّى در هر شیء
موجود بكلّ اسماء و صفات موصوف ومذكور و بعد از اخذ آن معدوم بوده تفكّر فى
ذلك لتكون من العارفين

\end{otherlanguage}

Va īnkih darbārih-yi baʿḍī āyāt-i ilāhīyyih bar ḥasab-i ẓāhir nāzil,
maqṣūd tashvīq-i ān nufūs būdih, va az ān gudhashtih jamīʿ bi-āyih-yi
tajallī-yi Ḥaqq rājiʿ. Mādāmī kih ān āyih-yi tajallī dar har shayʾ
mawjūd, bi-kull-i asmāʾ va ṣifāt mawṣūf va madhkūr, va baʿd az akhdh-i
ān maʿdūm būdih. Tafakkar fī dhālika li-takūna min al-ʿārifīn.

As for the divine verses that outwardly concern certain people, their
purpose was to encourage those souls. Beyond this, all refer to the sign
of the True One's self-revelation. So long as that sign of
self-revelation is present in anything, it is described and mentioned by
all names and attributes; once that sign is withdrawn, it is
nonexistent. Reflect on this, that you may be among those who
understand.

\subsection{57}\label{section-552}

\begin{otherlanguage}{persian}

و سلمّنا كه بعضى اوصاف مخصوص نفسى كه ذكر نموده نازل آيا نشنيده كه
ميفرمايد بحروفات بيان از نفس ظهور در حين ظهور محتجب مشويد و حروفات بيان
اوّل نفوسى هستند كه بكلّ اوصاف موصوف گشته‌اند و بكلّ اسماء موسوم و همين قول
سلطان قدم دليلى است واضح بر احتجاب بعضى از حروفات چنانچه مشاهده ميشود كه
نفسى خود را از حروف و مرايا محسوب ميدارد از ذيل تقديس منقطع و بحبل نفس و
هوى متمسّك

\end{otherlanguage}

Va sallamnā kih baʿḍī awṣāf makhṣūṣ-i nafsī kih dhikr namūdih nāzil. Āyā
nashunīdih kih mīfarmāyad bi-ḥurūfāt-i Bayān az nafs-i ẓuhūr dar ḥīn-i
ẓuhūr muḥtajib mashavīd? Va ḥurūfāt-i Bayān avval nufūsī hastand kih
bi-kull-i awṣāf mawṣūf gashtih-and va bi-kull-i asmāʾ mawsūm. Va hamīn
qawl-i Sulṭān-i Qidam dalīlī ast vāḍiḥ bar iḥtijāb-i baʿḍī az ḥurūfāt,
chinānchih mushāhadih mīshavad kih nafsī khud rā az ḥurūf va marāyā
maḥsūb mīdārad, az dhayl-i taqdīs munqaṭiʿ va bi-ḥabl-i nafs va havā
mutamassik.

Even if we concede that certain descriptions were revealed specifically
concerning the person you mentioned, have you not heard Him say: ``Do
not let the Letters of the Bayán veil you from the Manifestation Himself
at the time of His appearance''? The Letters of the Bayán are the first
souls to have been described by all attributes and designated by all
names. This very statement of the Sovereign of Eternity is clear
evidence that some of the Letters would be veiled, as can be seen: a
person counts himself among the Letters and Mirrors, yet is cut off from
the hem of holiness and clings to the cord of self and desire.

\subsection{58}\label{section-553}

\begin{otherlanguage}{persian}

و هم چنين ميفرمايد كه در ظهور نيّر اعظم در كرّه اخرى طوريّون معدوم و لاشیء
ميشوند قدرى تأمّل نمائيد كه شايد بعلوّ امر اطّلاع يابيد بعد از اينكه عظمت
امر بشأنى است كه ميفرمايند طوريّون ميّت و لاشیء ميشوند حروفات و مرايا در
آن مقرّ ذكر نداشته و ندارند الّا بعد حبّهنّ و هم‌چنين در مقام ديگر ميفرمايند
قوله جلّ و عزّ حلّ لمن يظهر ان يردّ من لم يكن فوق الأرض اعلى منه اذ ذلك خلق
فى قبضته انتهى

\end{otherlanguage}

Va hamchunīn mīfarmāyad kih dar ẓuhūr-i Nayyir-i Aʿẓam dar karrih-yi
ukhrā Ṭūriyyūn maʿdūm va lā shayʾ mīshavand. Qadrī taʾammul namāyīd kih
shāyad bi-ʿuluvv-i amr iṭṭilāʿ yābīd. Baʿd az īnkih ʿaẓamat-i amr
bi-shaʾnī ast kih mīfarmāyand Ṭūriyyūn mayyit va lā shayʾ mīshavand,
ḥurūfāt va marāyā dar ān maqarr dhikr nadāshtih va nadārand illā baʿda
ḥubbihinna. Va hamchunīn dar maqām-i dīgar mīfarmāyand, qawluhu jalla va
ʿazza: ḥalla li-man yaẓhar an yarudda man lam yakun fawq al-arḍ aʿlā
minhu, idh dhālika khalq fī qabḍatihi. Intahā.

He likewise says that at the appearance of the Most Great Luminary in
another cycle, the dwellers of Sinai become nonexistent and as nothing.
Reflect a little, that perhaps you may recognize the loftiness of the
Cause. Since the greatness of the Cause is such that He says the
dwellers of Sinai become dead and as nothing, the Letters and Mirrors
have never had, and do not have, any standing in that court except after
their love. Elsewhere He also says---exalted and mighty are His words:
``It is lawful for Him who will appear to reject the one above whom
there is none on earth, for that one is a creature in His grasp.'' End
of quotation.

\subsection{59}\label{section-554}

\begin{otherlanguage}{persian}

اين بيانات موجد اسماء و صفات و ان كلمات انفس محتجبه مغلّه تاللّه اذاً قد
اشرقت شموس محكمات و آيات واضحات و دلالات متقنات و كلمات تامّات و مع‌ذلك
مينويسيد ميرزا يحيى نميشود مردود شود مع آنكه منتهى شأن او اين بوده كه
بين يدى جمال ابهى حاضر ميشد و كلمات الهى بر او القاء ميگشت فواللّه مع
افعال او احدى متعرّض او نشده خود خود را مردود نمود باعراض از حقّ الّذى خلق
بقوله كذلك فصّل لك الأمر لتكون من العالمين

\end{otherlanguage}

Īn bayānāt mūjid-i asmāʾ va ṣifāt va ān kalimāt-i anfus-i muḥtajibih-yi
mughallih. Ta-llāh, idhan qad ashraqat shumūs muḥkamāt va āyāt wāḍiḥāt
va dalālāt mutqanāt va kalimāt tāmmāt. Va maʿa dhālik mīnivīsīd Mīrzā
Yaḥyā namīshavad mardūd shavad, maʿa ānkih muntahā-yi shaʾn-i ū īn būdih
kih bayn-i yaday Jamāl-i Abhā ḥāḍir mīshud va kalimāt-i ilāhī bar ū ilqā
mīgasht. Fa-wa-llāh maʿa afʿāl-i ū aḥadī mutaʿarriḍ-i ū nashudih, khud
khud rā mardūd namūd bi-iʿrāḍ az Ḥaqq alladhī khuliqa bi-qawlihi.
Kadhālika fuṣṣila laka al-amr li-takūna min al-ʿālimīn.

These utterances bring names and attributes into being, whereas those
are the words of veiled and shackled souls. By God, the suns of
unambiguous statements, clear verses, firmly established proofs, and
perfect words have now shone forth. Nevertheless, you write that Mírzá
Yaḥyá cannot be rejected, although the utmost extent of his station was
that he would enter the presence of the All-Glorious Beauty and divine
words would be imparted to him. By God, despite his deeds, no one has
assailed him. He made himself an outcast by turning away from the True
One, by whose word he was created. Thus has the matter been explained to
you, that you may be among those who know.

\subsection{60}\label{section-555}

\begin{otherlanguage}{persian}

حال تفكّر نمائيد كه محكمات آيات الهى كدام بوده و متشابهات كدام و تعرف
الّذين اخذوا المتشابهات و تركوا المحكمات عن ورائهم و كانوا من المشركين در
اينمقام محبوب است كه از آيات محكمات كه از سماء مشيّت ابهى جلّت عظمته نازل
شده ذكر نمايم فواللّه يكفى الخلايق اجمعين تا هر صاحب ايمانى كه ذرّه از حبّ
الهى در قلب او باشد بآن تمسّك جسته خود را از ظنونات اهل ارض و مؤتفكات
مظاهر شيطان نجات دهد و ديگر در ايمان خود ثابت ماند و تزلزل مشاهده نكند و
اين بعض آيات منزله است كه در لوح يكى از احبّا لأحد من الأشقیآء نازل شده
قوله عزّ بهائه و ان رأيت اسم الها ذكره بآيات ربّك لعلّ يكون من المقبلين الى
شطر الّذى من اقبل اليه فقد نجى و من اعرض فقد هلك كذلك نطق الرّوح حينئذ بين
السّموات و الأرضين

\end{otherlanguage}

Ḥāl tafakkur namāyīd kih muḥkamāt-i āyāt-i ilāhī kudām būdih va
mutashābihāt kudām, va taʿrif alladhīna akhadhū al-mutashābihāt va
tarakū al-muḥkamāt ʿan warāʾihim va kānū min al-mushrikīn. Dar īn maqām
maḥbūb ast kih az āyāt-i muḥkamāt kih az samāʾ-i mashīyyat-i Abhā jallat
ʿaẓamatuhu nāzil shudih dhikr namāyam, fa-wa-llāh yakfī al-khalāʾiq
ajmaʿīn, tā har ṣāḥib-i īmānī kih dharrih az ḥubb-i ilāhī dar qalb-i ū
bāshad bi-ān tamassuk justih, khud rā az ẓunūnāt-i ahl-i arḍ va
muʾtafikāt-i maẓāhir-i shayṭān najāt dahad va dīgar dar īmān-i khud
thābit mānad va tazalzul mushāhadih nakunad. Va īn baʿḍ-i āyāt-i
munazzalih ast kih dar lawḥ-i yikī az aḥibbāʾ li-aḥad min al-ashqiyāʾ
nāzil shudih, qawluhu ʿazza bahāʾuhu: va in raʾayta ism al-Hā dhakkirhu
bi-āyāt rabbika laʿalla yakūn min al-muqbilīn ilā shaṭr alladhī man
aqbala ilayhi fa-qad najā va man aʿraḍa fa-qad halaka. Kadhālika naṭaqa
al-Rūḥ ḥīnaʾidh bayn al-samāwāt va al-araḍīn.

Now reflect on which of the divine verses are unambiguous and which are
ambiguous, and recognize those who have taken hold of the ambiguous
passages, cast the unambiguous ones behind their backs, and become among
those who associate partners with God. Here I wish to mention some of
the unambiguous verses revealed from the heaven of the will of the
All-Glorious---exalted be His grandeur. By God, they suffice for all
created beings, so that every person of faith who has an atom of divine
love in their heart may cling to them, rescue themselves from the
conjectures of the people of the earth and the fabrications of the
manifestations of Satan, and thereafter remain steadfast in faith and
experience no wavering. These are some of the verses revealed in a
Tablet to one of the loved ones concerning one of the wretched. He
says---mighty be His glory: ``If you see the name Há, remind him of the
verses of your Lord, that perhaps he may turn toward Him through turning
to whom one is saved, and through turning away from whom one perishes.
Thus has the Spirit spoken at this time between the heavens and the
earths.''

\subsection{61}\label{section-556}

\begin{otherlanguage}{arabic}

قل ان يا عبد اسمع قولى تاللّه الحقّ لن تجد لنفسك ناصحاً انصح منّى لاتوقّف على
الصّراط ثمّ مرّ عنه كمرّ السّحاب كذلك امرناك من قبل و نأمرك حينئذ و ان لن
تقبل منّى فسوف نأمرك بذلك من بعد و من بعد بعد الى ان تمرّ عنه بسلطانى
الغالب المقتدر الحكيم

\end{otherlanguage}

Qul in yā ʿabd, ismaʿ qawlī. Ta-llāh al-Ḥaqq, lan tajida li-nafsika
nāṣiḥan anṣaḥ minnī. Lā tawaqqaf ʿalā al-ṣirāṭ, thumma murr ʿanhu
ka-marr al-saḥāb. Kadhālika amarnāka min qabl va naʾmuruka ḥīnaʾidh, va
in lan taqbal minnī fa-sawfa naʾmuruka bi-dhālika min baʿd va min baʿd
baʿd ilā an tamurra ʿanhu bi-sulṭānī al-ghālib al-muqtadir al-ḥakīm.

Say: ``O servant, listen to My words. By God, the True One, you will
find no counselor more sincere toward you than I. Do not linger on the
path, but pass beyond it as the clouds pass. Thus We commanded you
before and command you now. If you will not accept this from Me, We will
command you to do so hereafter, and again thereafter, until you pass
beyond it through My all-conquering, all-powerful, all-wise
sovereignty.''

\subsection{62}\label{section-557}

\begin{otherlanguage}{arabic}

شقّ حجبات الأوهام بذكرى ثمّ اسمى و دع كلّ ما يمنعك عن ورائك ثمّ اظهر بغتةً
بظهور ربّك و لاتكن من الصّابرين ایّاک ان تمنعك الرّياسة عن ذكر ربّك تاللّه
الحقّ لو يسجدك كلّ من فى السّموات و الأرض و لم تكن فى ظلّى لاينفعك و بذلك
يشهد روحك لو تكون من السّامعين اتحبّ بأن تكون من الّذين منعتهم الرّياسة فى
ایّام الّتى شقّت فيها سمآء الأمر و اتى على ظلل و كفروا باللّه خالقهم و
مبدعهم الى ان افتوا على قتله بعد الّذى جائهم بحجّة من لدى اللّه و برهان
عظيم

\end{otherlanguage}

Shuqq ḥujubāt al-awhām bi-dhikrī thumma ismī, va daʿ kull mā yamnaʿuka
ʿan warāʾika, thumma iẓhar baghtatan bi-ẓuhūr rabbika va lā takun min
al-ṣābirīn. Iyyāka an tamnaʿaka al-riyāsa ʿan dhikr rabbika. Ta-llāh
al-Ḥaqq, law yasjuduka kull man fī al-samāwāt va al-arḍ va lam takun fī
ẓillī lā yanfaʿuka, va bi-dhālika yashhadu rūḥuka law takūn min
al-sāmiʿīn. A-tuḥibbu bi-an takūna min alladhīna manaʿathum al-riyāsa fī
ayyām allatī shuqqat fīhā samāʾ al-amr va atā ʿalā ẓulal, va kafarū
bi-llāh khāliqihim va mubdiʿihim ilā an aftaw ʿalā qatlihi baʿda alladhī
jāʾahum bi-ḥujja min ladā Allāh va burhān ʿaẓīm.

Tear through the veils of delusion by My remembrance and then My name,
and leave behind you everything that holds you back. Then appear
suddenly through the appearance of your Lord, and do not be among those
who wait. Beware lest leadership keep you from the remembrance of your
Lord. By God, the True One, even if everyone in the heavens and on earth
were to bow down before you, it would not benefit you if you were not
beneath My shadow. Your own spirit bears witness to this, if you are
among those who listen. Do you wish to be among those whom leadership
held back in the days when the heaven of the Cause was split apart and
He came upon clouds, and who disbelieved in God, their Creator and
Originator, until they pronounced a sentence of death against Him after
He had come to them with a testimony from God and a mighty proof?

\subsection{63}\label{section-558}

\begin{otherlanguage}{arabic}

تاللّه يا ايّها العبد ماكان مقصودى الّا تطهيرك عن كلّ ما لا ينبغى لك و كلّما
اصمت فى ذكرك لسان اللّه ينطق على لسانى و يأمرنى بالتّبليغ عليك و ما علىّ
الّا البلاغ المبين فكّر فى نفسك اقلّ من ان هل سمعت ظهوراً اعظم من هذا او من
آيات اكبر عما نزّلت بالحقّ اذا فانطق على الصّدق الخالص و كن من الّذينهم
توجّهوا الى منظر الأكبر فى يوم الّذى انقلبت فيه وجوه الخلايق اجمعين

\end{otherlanguage}

Ta-llāh, yā ayyuhā al-ʿabd, mā kāna maqṣūdī illā taṭhīruka ʿan kull mā
lā yanbaghī laka. Va kullamā aṣmut fī dhikrika lisān Allāh yanṭiqu ʿalā
lisānī va yaʾmurunī bi-al-tablīgh ʿalayka, va mā ʿalayya illā al-balāgh
al-mubīn. Fakkir fī nafsika aqall min ān, hal samiʿta ẓuhūran aʿẓam min
hādhā aw min āyāt akbar ʿammā nuzzilat bi-al-ḥaqq? Idhan fa-nṭiq ʿalā
al-ṣidq al-khāliṣ va kun min alladhīnahum tawajjahū ilā manẓar al-akbar
fī yawm alladhī inqalabat fīhi wujūh al-khalāʾiq ajmaʿīn.

By God, O servant, My purpose has been only to purify you of everything
unworthy of you. Whenever I fall silent concerning you, the tongue of
God speaks through My tongue and commands Me to convey the message to
you. My duty is only to convey it clearly. Reflect within yourself for
less than an instant: have you heard of a revelation greater than this,
or of verses greater than those revealed in truth? Then speak with pure
truthfulness and be among those who have turned toward the Most Great
Vision on the day when the faces of all created beings have been turned
away.

\subsection{64}\label{section-559}

\begin{otherlanguage}{arabic}

و ان تخاف من ايمانك خذ هذا اللّوح ثمّ احفظه فى جيب توكّلك و اذا دخلت موقف
الحشر فى يوم الّذى فيه يبعث كلّ الممكنات تلقآء وجه ربّك من نفحاته البديع
المنيع و يسئلك اللّه بأىّ حجّة آمنت بهذا الظّهور اذا فاخرج اللّوح و قل بهذا
الكتاب المنزل المبارك القديم ثمّ اقرأ ما نزل فيه تلقآء وجه ربّك مقرّ الّذى
تشهد فيه كلّ النّبيين و المرسلين تاللّه اذاً تمدّ ايادى الكلّ اليك و يأخذنّ
اللّوح و يضعنه على عيونهم شوقاً للقائى و شغفا لحبی‌ و يجدن منه روايح قدسى
العزيز المنيع كذلك فصّلنا لك الآيات لتطمئنّ فى نفسك و تكون من المطمئنّين

\end{otherlanguage}

Va in takhāf min īmānika khudh hādhā al-lawḥ thumma iḥfaẓhu fī jayb
tawakkulika, va idhā dakhalta mawqif al-ḥashr fī yawm alladhī fīhi
yubʿath kull al-mumkināt tilqāʾ wajh rabbika min nafaḥātihi al-badīʿ
al-manīʿ, va yasʾaluka Allāh bi-ayy ḥujja āmanta bi-hādhā al-ẓuhūr,
idhan fa-akhrij al-lawḥ va qul bi-hādhā al-kitāb al-munzal al-mubārak
al-qadīm. Thumma iqraʾ mā nazala fīhi tilqāʾ wajh rabbika, maqarr
alladhī tashhadu fīhi kull al-nabiyyīn va al-mursalīn. Ta-llāh, idhan
tumadd ayādī al-kull ilayka va yaʾkhudhanna al-lawḥ va yaḍaʿnahu ʿalā
ʿuyūnihim shawqan li-liqāʾī va shaghafan li-ḥubbī, va yajidanna minhu
rawāʾiḥ qudsī al-ʿazīz al-manīʿ. Kadhālika faṣṣalnā laka al-āyāt
li-taṭmaʾinna fī nafsika va takūna min al-muṭmaʾinnīn.

If you fear on account of your faith, take this Tablet and keep it in
the breast of your trust in God. When you enter the place of gathering
on the day when all contingent beings are raised before the face of your
Lord through His wondrous, unassailable breaths, and God asks you, ``By
what proof did you believe in this Revelation?'' then bring out the
Tablet and say, ``By this revealed, blessed, ancient Book.'' Then read
what has been revealed in it before the face of your Lord, in that abode
where you will behold all the Prophets and Messengers. By God, then the
hands of all will stretch out toward you, and they will take the Tablet
and place it upon their eyes, longing to meet Me and aflame with love
for Me. They will find in it the fragrances of My mighty, unassailable
holiness. Thus have We set forth the verses for you, that you may be
assured within yourself and be among those who are at peace.

\subsection{65}\label{section-560}

\begin{otherlanguage}{arabic}

و لو يعذّبك اللّه بما آمنت بآياته فى هذا الظّهور فبأىّ حجّة يعذّب الّذينهم ما
آمنوا بعلىّ من قبل و من قبله بمحمّد رسول اللّه و من قبله بعيسى بن مريم و من
قبله بالكليم و من قبله بالخليل الى ان ينتهى الظّهورات الى البديع الأوّل
الّذى خلق بارادة ربّك القادر المريد

\end{otherlanguage}

Va law yuʿadhdhibuka Allāh bi-mā āmanta bi-āyātihi fī hādhā al-ẓuhūr,
fa-bi-ayy ḥujja yuʿadhdhib alladhīnahum mā āmanū bi-ʿAlī min qabl, va
min qablihi bi-Muḥammad Rasūl Allāh, va min qablihi bi-ʿĪsā ibn Maryam,
va min qablihi bi-al-Kalīm, va min qablihi bi-al-Khalīl, ilā an
yantahiya al-ẓuhūrāt ilā al-Badīʿ al-Awwal alladhī khuliqa bi-irādat
rabbika al-Qādir al-Murīd.

If God were to punish you for believing in His verses in this
Revelation, by what proof would He punish those who did not believe in
`Alí before this, and before Him in Muḥammad, the Messenger of God, and
before Him in Jesus son of Mary, and before Him in the One who spoke
with God, and before Him in the Friend of God, until the succession of
Revelations reaches the First Wondrous One, who was created through the
will of your Lord, the Powerful, the One who wills?

\subsection{66}\label{section-561}

\begin{otherlanguage}{arabic}

ان يا اسمى انّ الأمر اظهر من ان يخفى و ابين من ان يستر و يستضىء كالشّمس فى
قطب الزّوال و انّك لو تخلص نفسك عن الحجبات لتصل اليه اقرب من ان يرتدّ بصرك
الى نفسك و انّ هذا لحقّ يقين

\end{otherlanguage}

In yā ismī, inna al-amra aẓharu min an yakhfā wa-abyanu min an yustara
wa-yastaḍīʾu ka-al-shamsi fī quṭbi al-zawāl, wa-innaka law tukhliṣu
nafsaka ʿan al-ḥujubāti la-taṣila ilayhi aqraba min an yartadda baṣaruka
ilā nafsika, wa-inna hādhā la-ḥaqqu yaqīn.

O My Name, the Cause is too manifest to be concealed and too clear to be
veiled; it shines like the sun at the zenith of noon. If you free
yourself from the veils, you will reach it sooner than your gaze can
return to yourself. This is indeed the certain truth.

\subsection{67}\label{section-562}

\begin{otherlanguage}{arabic}

اسمع قول من ينطق بالحقّ و لاتجادل بآيات اللّه بعد انزالها و لو يأمرك بذلك
كلّ العباد و لاتمنع نفسك عن فضله و لو يمنعك عن ذلك كلّ الثّقلين فانظر امر
ربّك ببصرك ثمّ اعرفه بنفسك و روحك لأنّ عرفان غيرك لم يكن دليلاً لك و اعراض
ما سواك لم يكن حجّة عليك ایّاک ان تحتجب عن الّذى لو تحتجب عنه اقلّ من ان
لتحبط اعمالك و يصدقنى فى ذلك كلّما نزل من قبل من صحايف اللّه الملك المنزل
القدير انتهى

\end{otherlanguage}

Ismaʿ qawla man yanṭiqu bi-al-ḥaqqi wa-lā tujādil bi-āyāti Allāhi baʿda
inzālihā wa-law yaʾmuruka bi-dhālika kullu al-ʿibād, wa-lā tamnaʿ
nafsaka ʿan faḍlihi wa-law yamnaʿuka ʿan dhālika kullu al-thaqalayn.
Fa-unẓur amra rabbika bi-baṣrika thumma iʿrifhu bi-nafsika wa-rūḥika
li-anna ʿirfāna ghayrika lam yakun dalīlan laka wa-iʿrāḍa mā siwāka lam
yakun ḥujjatan ʿalayka. Iyyāka an taḥtajiba ʿan alladhī law taḥtajibu
ʿanhu aqalla min ān la-taḥbaṭa aʿmāluka, wa-yuṣaddiqunī fī dhālika kullu
mā nazala min qablu min ṣaḥāʾifi Allāhi al-maliki al-munzili al-qadīr.
Intahā.

Listen to the words of Him who speaks the truth, and do not dispute
God's verses after they have been sent down, even if all people bid you
do so. Do not deprive yourself of His grace, even if all humans and jinn
would keep you from it. Examine your Lord's Cause with your own eyes,
then recognize Him with your own self and spirit, for another's
recognition is no proof for you, nor is the rejection of anyone else an
argument against you. Beware of being veiled from Him: were you to
remain veiled from Him for less than an instant, your deeds would come
to nothing. Everything previously sent down in the scriptures of God,
the King, the Revealer, the All-Powerful, bears witness to the truth of
what I say. End of quotation.

\subsection{68}\label{section-563}

\begin{otherlanguage}{persian}

اين كلمات منزله است كه ذكر شد حال انصاف دهيد و كلمات محكمات را از دونش
بشناسيد و اگر نفسى جذوه از نار محبّت الهى در قلبش باشد بآنچه از سماء مشيّت
ربّانى نازل شده عمل مينمايد و منقطعاً عن كلّ من فى السّموات و الأرض و مستقرّاً
على بساط السّكون و موقناً بامراللّه و مستغنياً عمّا سواه خود را مشاهده
مينمايد و اللّه حجّت تمام و نعمت كامل و شمس مشرق و آيات نازل و ظهور اظهر و
مع‌ذلك خود را در تيه تحقيق و ظنون و اشارات مبتلى نموده‌ايد نسئل اللّه بأن
يوفّقنا و ایّاک و يرزقنا ما قدّر من عنده و انّه لهو العزيز الكريم

\end{otherlanguage}

Īn kalemāt-e monazale ast keh dhekr shod. Ḥāl enṣāf dehīd va kalemāt-e
moḥkamāt rā az dūnash beshenāsīd va agar nafsī jadhve az nār-e
maḥabbat-e elāhī dar qalbash bāshad be-āncheh az samā-ye mashīyyat-e
rabbānī nāzel shodeh ʿamal mīnamāyad va munqaṭiʿan ʿan kulli man fī
al-samāwāti wa-al-arḍi wa-mustaqirran ʿalā bisāṭi al-sukūni wa-mūqinan
bi-amri Allāhi wa-mustaghniyan ʿammā siwāhu khod rā moshāhedeh
mīnamāyad. Va-Allāh ḥojjat tamām va neʿmat kāmel va shams moshreq va
āyāt nāzel va ẓohūr aẓhar va maʿa-dhālek khod rā dar tīh-e taḥqīq va
ẓonūn va eshārāt mobtalā nemūdeh-īd. Nasʾalu Allāha bi-an yuwaffiqanā
wa-iyyāka wa-yarzuqanā mā qaddara min ʿindihi wa-innahu la-huwa
al-ʿazīzu al-karīm.

These are the revealed words that have been cited. Now judge fairly, and
distinguish the unambiguous words from what is otherwise. If a soul has
a spark of the fire of divine love in his heart, he acts according to
what has descended from the heaven of the divine Will, and finds himself
detached from all who are in the heavens and on earth, firmly
established upon the plane of tranquillity, assured of God's Cause, and
independent of all else besides Him. By God, the proof is complete, the
blessing perfect, the sun shining, the verses descending, and the
Revelation most manifest; yet you have entangled yourselves in the
wilderness of investigation, conjectures, and allusions. We ask God to
grant success to us and to you, and to bestow upon us what He has
ordained from His presence. He is indeed the Almighty, the Generous.

\subsection{69}\label{section-564}

\begin{otherlanguage}{persian}

و ديگر قولك آفرين عجب كارى شد هنوز مردم در اصل حرف دارند باين معنى هنوز
ندانسته‌اند سخن چه بوده و احكام چه داشته حال اوضاع تازه و دين تأويلى كه
موجب هزار گونه فتنه شود برپا شده انتهى

\end{otherlanguage}

Va dīgar qawluka: Āfarīn, ʿajab kārī shod! Hanūz mardom dar aṣl ḥarf
dārand, be-īn maʿnā hanūz nadānesteh-and sokhan cheh būdeh va aḥkām cheh
dāshteh. Ḥāl awżāʿ-e tāzeh va dīn-e taʾvīlī keh mūjeb-e hezār gūneh
fetneh shavad barpā shodeh. Intahā.

And further, your statement: ``Bravo! What a thing has happened! People
still have objections to the very foundation---that is, they have not
yet understood what the message was or what laws it contained. Now new
arrangements and a religion of interpretation have been set up that will
cause a thousand kinds of turmoil.'' End of quotation.

\subsection{70}\label{section-565}

\begin{otherlanguage}{persian}

بشما بی‌‌ادبی‌ نمينمايم و لكن اشهد باللّه كه تا حال هيچ عاقلى چنين تكلّم
نمينمايد چه كه اگر منوط حرف مردم و ما عندهم باشد هرگز بايد حقّ ظاهر نشود
چه كه لازال در هر عهد و عصر كه شمس حقيقت ربّانى از افق مشيّت طالع و مشرق
كلّ حرف داشته و باشدّ اعراض معرض و جميع آن هيكل احديّه را از اهل فتنه و
فساد ميشمردند مع آنكه نزد هر ذى بصرى مشهود است كه مقصود جز اتّحاد نبوده و
نخواهد بود

\end{otherlanguage}

Be-shomā bī-adabī namīnamāyam va lāken ashhadu bi-Allāh keh tā ḥāl hīch
ʿāqelī chonīn takallom namīnamāyad, cheh keh agar manūṭ-e ḥarf-e mardom
va mā ʿendahom bāshad hargaz bāyad Ḥaqq ẓāher nashavad, cheh keh lā-zāl
dar har ʿahd va ʿaṣr keh shams-e ḥaqīqat-e rabbānī az ofoq-e mashīyyat
ṭāleʿ va moshreq koll ḥarf dāshteh va be-ashadd-e eʿrāż moʿreż va jamīʿ
ān haykal-e aḥadīyyeh rā az ahl-e fetneh va fasād mīshomordand, maʿ-e
ānkeh nazd-e har dhī-baṣarī mashhūd ast keh maqṣūd joz etteḥād nabūdeh
va nakhāhad būd.

I intend no discourtesy to you, but I call God to witness that no
rational person has ever spoken in this way. For if it depended on
people's objections and what they possess, the True One would never be
manifested. In every age and era when the Sun of divine Reality has
risen and shone from the horizon of the divine Will, all have raised
objections, turned away in the strongest rejection, and counted that
embodiment of oneness among the people of turmoil and corruption. Yet it
is evident to every discerning person that the purpose has never been,
and will never be, anything but unity.

\subsection{71}\label{section-566}

\begin{otherlanguage}{persian}

و از اين گذشته حرف اين مردم تخصيص باين ظهور نداشته در ظهور رسول اللّه
هنوز حرف دارند اگر حرف مردم منوط باشد نعوذ باللّه بايد انكار كلّ ظهورات
نمود و از آن گذشته از اين ظهور احكام بيان ثابت و ظاهر شده نه از آن انفس
متوهّمه كه شما خيال نموده‌ايد و اللّه بحرفى از آن عامل نشده‌اند و نعوذ باللّه
اگر از اعمالشان ذكر شود يفرّ الشّيطان من اعمالهم بسيار امور برتو مشتبه شده
و مستور گشته نسئل اللّه بأن يفتح بصريك بالفضل و ينقذك عن الظّلمات النّفس و
الهوى و يقرّبك الى سدرة المنتهى و انّه لهو القادر على ما یشآء و انّه بكلّ
شیء رقيباً

\end{otherlanguage}

Va az īn godhashteh ḥarf-e īn mardom takhṣīṣ be-īn ẓohūr nadāshteh, dar
ẓohūr-e Rasūl Allāh hanūz ḥarf dārand. Agar ḥarf-e mardom manūṭ bāshad,
naʿūdhu bi-Allāh, bāyad enkār-e koll-e ẓohūrāt nemūd. Va az ān
godhashteh az īn ẓohūr aḥkām-e Bayān thābet va ẓāher shodeh, nah az ān
anfos-e motavahhemeh keh shomā khayāl nemūdeh-īd. Va-Allāh be-ḥarfī az
ān ʿāmel nashodeh-and va, naʿūdhu bi-Allāh, agar az aʿmāles hān dhekr
shavad, yafirru al-shayṭānu min aʿmālihim. Besyār omūr bar to moshtabeh
shodeh va mastūr gashteh. Nasʾalu Allāha bi-an yaftaḥa baṣrayka
bi-al-faḍli wa-yunqidhaka ʿan al-ẓulumāti al-nafsi wa-al-hawā
wa-yuqarribaka ilā sidrati al-muntahā wa-innahu la-huwa al-qādiru ʿalā
mā yashāʾu wa-innahu bi-kulli shayʾin raqīban.

Moreover, these people's objections have not been confined to this
Revelation: they still object to the Revelation of the Messenger of God.
If people's objections were to be the deciding factor, we would have to
deny all Revelations---God forbid! Furthermore, the laws of the Bayān
have been established and made manifest through this Revelation, not
through those deluded souls you have imagined. By God, they have not
acted upon a single word of it. God forbid that their deeds should be
mentioned: Satan flees from their deeds. Many matters have become
confused and hidden from you. We ask God to open your eyes through His
grace, rescue you from the darkness of self and desire, and bring you
near to the Lote Tree of the Utmost Limit. He indeed has power over
whatever He wills, and He watches over all things.

\subsection{72}\label{section-567}

\begin{otherlanguage}{persian}

حقّ ناظر بقول و اعراض و اعتراض ناس نبوده و نخواهد بود يفعل ما یشآء رغماً
لانف الّذينهم كفروا باللّه و كانوا من المعرضين اگر فى‌الجمله تعقّل نمائيد
اين اعتراضات شما جميعاً بر نقطه اولى روح ما سواه فداه وارد و هم‌چنين من
قبل بكلّ رسل اسمع قولى ثمّ استغفر عمّا ارتكبت لعلّ اللّه يغفرك و يكفّر عنك
سيّئاتك و انّه لغفور رحيم

\end{otherlanguage}

Ḥaqq nāẓer be-qawl va eʿrāż va eʿterāż-e nās nabūdeh va nakhāhad būd.
Yafʿalu mā yashāʾu raghman li-anfi alladhīna-hum kafarū bi-Allāhi
wa-kānū min al-muʿriḍīn. Agar fī-al-jomleh taʿaq qol namāyīd īn
eʿterāżāt-e shomā jamīʿan bar Noqṭeh-ye Ūlā, rūḥ-e mā sevāh fedāh, vāred
va ham-chenīn men qabl be-koll-e rosol. Ismaʿ qawlī thumma istaghfir
ʿammā irtakabta laʿalla Allāha yaghfiruka wa-yukaffiru ʿanka sayyiʾātika
wa-innahu la-ghafūrun raḥīm.

The True One has never heeded, and will never heed, people's words,
rejection, or objections. He does whatever He wills, despite those who
have disbelieved in God and turned away. If you reflect even a little,
all these objections of yours apply to the Primal Point---may the souls
of all besides Him be sacrificed for Him---and likewise to all the
Messengers before Him. Listen to My words, then seek forgiveness for
what you have done, that God may forgive you and blot out your misdeeds.
He is indeed forgiving and merciful.

\subsection{73}\label{section-568}

\begin{otherlanguage}{arabic}

قول ديگر آن مبغض ملقى

\end{otherlanguage}

Qawl-e dīgar-e ān mobgheż-e molqī.

Another statement by that hateful instigator:

\subsection{74}\label{section-569}

\begin{otherlanguage}{persian}

پس بنابراين فكر خود را گر كنى تاويل به كه كنى تاويل آن نامشتبه

\end{otherlanguage}

Pas banābarīn fekr-e khod rā gar konī taʾvīl beh keh konī taʾvīl-e ān
nāmoshtabeh.

``Therefore, it is better to reinterpret your own thinking than to
reinterpret what is unambiguous.

\subsection{75}\label{section-570}

\begin{otherlanguage}{persian}

عجب است از شما كه امر بر شما مشتبه شده در صريح امر كم جو التباس انتهى

\end{otherlanguage}

ʿAjab ast az shomā keh amr bar shomā moshtabeh shodeh. Dar ṣarīḥ-e amr
kam jū eltebās. Intahā.

It is astonishing that the matter has become confused for you. Do not
seek ambiguity in an explicit matter.'' End of quotation.

\subsection{76}\label{section-571}

\begin{otherlanguage}{persian}

اين عباد فكر خود را نموده‌اند و بفضل اللّه از آجام تاويل بصحراى با فضاى
خوش تصريح رسيده‌اند مأوّل كسانى هستند كه حقّ ظاهر منيع را ميگذارند و از
تيشه ظنون اصنام اوهام ميتراشند و بعد بمصنوع خود عاكف و ساجد و خاضع
ميشوند و مع‌ذلك خود را از عابدين حقّ ميشمرند و منقطعين الى اللّه ميدانند

\end{otherlanguage}

Īn ʿebād fekr-e khod rā nemūdeh-and va be-faḍl Allāh az ājām-e taʾvīl
be-ṣaḥrā-ye bā-faẓā-ye khosh-e taṣrīḥ resīdeh-and. Moʾavvel kasānī
hastand keh Ḥaqq-e ẓāher-e manīʿ rā mīgo dhārand va az tīsheh-ye ẓonūn
aṣnām-e awhām mītarāshand va baʿd be-maṣnūʿ-e khod ʿākef va sājed va
khāżeʿ mīshavand va maʿa-dhālek khod rā az ʿābedīn-e Ḥaqq mīshomorand va
munqaṭiʿīn ilā Allāh mīdānand.

These servants have reflected for themselves and, through God's grace,
have emerged from the thickets of interpretation into the pleasant,
spacious plain of explicit statement. The interpreters are those who
abandon the manifest, exalted True One, carve idols of fancy with the
axe of conjecture, and then devote themselves to their own handiwork,
prostrating themselves and bowing before it. Yet they count themselves
among the worshippers of the True One and regard themselves as detached
from all but God.

\subsection{77}\label{section-572}

\begin{otherlanguage}{persian}

بشنو ذكر اين عبد را و قلب را از اوهام فارغ كن و بربّ الأنام ناظر شو چه كه
دوست آن است كه اگر ملاحظه نمايد دوستى در غرقاب فنا افتاده بقدر قوّه و
استعداد خود اعانت نمايد كه شايد سبب نجاح و فلاح گردد لذا اميدواريم كه
بحبل عنايت الهى متمسّك شويد و از غرقاب كدورت و ظلمت برهيد و در ظلّ اشراق
فجر الهيّه وارد شويد

\end{otherlanguage}

Beshnow dhekr-e īn ʿabd rā va qalb rā az awhām fāregh kon va be-Rabb
al-anām nāẓer show, cheh keh dūst ān ast keh agar molāḥeẓeh namāyad
dūstī dar gharqāb-e fanā oftādeh be-qadr-e qovveh va esteʿdād-e khod
eʿānat namāyad keh shāyad sabab-e najāḥ va falāḥ gardad. Ledhā omīdvārīm
keh be-ḥabl-e ʿenāyat-e elāhī motamassek shavīd va az gharqāb-e kodūrat
va ẓolmat berahīd va dar ẓell-e eshrāq-e fajr-e elāhīyyeh vāred shavīd.

Listen to this servant's words, free your heart from vain imaginings,
and turn your gaze toward the Lord of all people. A friend is one who,
on seeing a friend fallen into the whirlpool of destruction, helps him
to the extent of his strength and ability, in the hope of bringing about
his deliverance and salvation. We therefore hope that you will hold fast
to the cord of divine providence, escape the whirlpool of impurity and
darkness, and enter beneath the shelter of the radiance of the divine
dawn.

\subsection{78}\label{section-573}

\begin{otherlanguage}{persian}

شما امر باين صريحى كه الآن موجود است و بحقّ ناطق از او گذشته و چشم
پوشيده‌ايد و باب تأويل كه لازال سقيم و عليل بوده بر وجه خود باز نموده‌ايد
بشأنى كه بخود شما كذب القاء نموده و مع‌ذلك مينويسيد و صادقش ميشمريد از
جمله كاغذى از آنجا رسيد معلوم نبود كه از كيست و كه نوشته چه كه بی‌مهر و
اسم بوده و لكن بعضى كلمات كاذبه مجعوله در آن كاغذ ثبت بود و اين واضح است
كه شيطان القاء نموده و آن جعّال كه معلوم نيست كيست بالقاى او نوشته خود
نويسنده بر اين كذب مطّلع است و شما هم مطّلع و مع‌ذلك او را حقّ دانسته‌ايد و
بظنون و اوهام ميخواهيد حقّ من له الحقّ را بدونش راجع نمائيد حال ملاحظه
كنيد كه در صريح امر كه التباس جسته نسئل اللّه بأن يفتح بصرك و يؤيّدك على
عرفان مظهر نفسه و يجعلك من الفائزين

\end{otherlanguage}

Shomā amr be-īn ṣarīḥī keh al-ān mowjūd ast va be-ḥaqq nāṭeq az ū
godhashteh va chashm pūshīdeh-īd va bāb-e taʾvīl keh lā-zāl saqīm va
ʿalīl būdeh bar vajh-e khod bāz nemūdeh-īd, be-shaʾnī keh be-khod-e
shomā kedhb elqā nemūdeh va maʿa-dhālek mīnevīsīd va ṣādeqash
mīshomorīd. Az jomleh kāghadhī az ānjā resīd, maʿlūm nabūd keh az kīst
va keh neveshteh, cheh keh bī-mohr va esm būdeh va lāken baʿżī kalemāt-e
kādhebeh-ye majʿūleh dar ān kāghadh thabt būd va īn vāżeḥ ast keh
shayṭān elqā nemūdeh va ān jaʿʿāl keh maʿlūm nīst kīst be-elqā-ye ū
neveshteh. Khod-e nevīsandeh bar īn kedhb moṭṭaleʿ ast va shomā ham
moṭṭaleʿ va maʿa-dhālek ū rā ḥaqq dānesteh-īd va be-ẓonūn va awhām
mīkhāhīd ḥaqq-e man lahu al-ḥaqq rā be-dūnash rājeʿ namāyīd. Ḥāl
molāḥeẓeh konīd keh dar ṣarīḥ-e amr keh eltebās josteh. Nasʾalu Allāha
bi-an yaftaḥa baṣaraka wa-yuʾayyidaka ʿalā ʿirfāni maẓhari nafsihi
wa-yajʿalaka min al-fāʾizīn.

You have passed over and disregarded a Cause so explicit, now present
and speaking the truth, and have opened before yourselves the door of
interpretation, which has always been unsound and infirm. It has gone so
far as to instill falsehood in you yourselves, yet you write it down and
regard it as true. Among other things, a letter arrived from there whose
sender and writer were unknown, since it bore neither seal nor name. Yet
certain false, fabricated statements were recorded in that letter. It is
clear that Satan inspired them, and that the forger, whose identity is
unknown, wrote at his prompting. The writer himself knows this
falsehood, and you know it too; nevertheless, you have regarded it as
true, and through conjectures and vain imaginings you wish to assign the
right of Him to whom the right belongs to someone else. Now consider who
has sought ambiguity in an explicit matter. We ask God to open your
eyes, assist you to recognize the Manifestation of His own Self, and
make you one of those who attain their goal.

\subsection{79}\label{section-574}

\begin{otherlanguage}{persian}

قول ديگر آن معترض بلى بنده در ميان نبوده امّا چنان ميدانم كه بهتر از آن
اشخاصيكه سالها بوده‌اند اطّلاع دارم علاوه از آنكه قواعد كلّيه كه ميزان هست
در دست هست انتهى

\end{otherlanguage}

Qawl-e dīgar-e ān moʿtareż: Balī bandeh dar mīyān nabūdeh ammā chonān
mīdānam keh behtar az ān ashkhāṣī keh sālhā būdeh-and eṭṭelāʿ dāram,
ʿalāveh az ānkeh qavāʿed-e kollīyyeh keh mīzān hast dar dast hast.
Intahā.

Another statement by that objector: ``Yes, I was not among them, but I
believe I am better informed than those who were there for years.
Moreover, the general principles that serve as the standard are at
hand.'' End of quotation.

\subsection{80}\label{section-575}

\begin{otherlanguage}{persian}

اينكه ادّعاى اطّلاع نموده‌ايد نه واللّه ابداً كسى مطّلع نيست نه شما و نه غير
شما و آنچه بشما القاء كرده‌اند واللّه الّذى لا اله الّا هو كلّ خلاف بوده و
خواهد بود كاش آن نفسى را كه الآن من دون اللّه حقّ اخذ کرده‌اند مثل يكى از
آدمهاى سوق بود عجب است كه شما اراده نموده‌ايد كه باين عبد مشتبه نمائيد
جميع كلمات شما و ظنونات و تأويلات و اشارات و دلالات و محكمات و متشابهات
و ادّعاى علميّت كه در اين امر نمود‌ه‌ايد كلّ بسيّد محمّد اصفهانى راجع و منتهى
ميشود فضيلت او نزد اين عباد خوب واضح است و از قبل هم واضح بوده عجب است
از بصيرت شما خوب آدم ميشناسيد

\end{otherlanguage}

Īnkeh eddeʿā-ye eṭṭelāʿ nemūdeh-īd, nah va-Allāh, abadan kasī moṭṭaleʿ
nīst, nah shomā va nah ghayr-e shomā, va āncheh be-shomā elqā
kardeh-and, wa-Allāhi alladhī lā ilāha illā huwa, koll khelāf būdeh va
khāhad būd. Kāsh ān nafsī rā keh al-ān men dūn Allāh Ḥaqq akhdh
kardeh-and mesl-e yekī az ādamhā-ye sūq būd. ʿAjab ast keh shomā erādeh
nemūdeh-īd keh be-īn ʿabd moshtabeh namāyīd. Jamīʿ-e kalemāt-e shomā va
ẓonūnāt va taʾvīlāt va eshārāt va delālāt va moḥkamāt va motashābehāt va
eddeʿā-ye ʿelmīyyat keh dar īn amr nemūdeh-īd koll be-Sayyed Moḥammad-e
Eṣfahānī rājeʿ va montahī mīshavad. Fażīlat-e ū nazd-e īn ʿebād khūb
vāżeḥ ast va az qabl ham vāżeḥ būdeh. ʿAjab ast az baṣīrat-e shomā, khūb
ādam mīshenāsīd!

As for your claim to be informed: no, by God, no one is informed at all,
neither you nor anyone else. By God, besides whom there is no other god,
everything they have instilled in you has been, and will remain,
contrary to the truth. Would that the person whom they have now taken as
the True One apart from God were like even one of the people in the
marketplace! It is astonishing that you have sought to confuse this
servant. All your words, conjectures, interpretations, allusions,
arguments, unambiguous and ambiguous statements, and claims to learning
in this matter refer and lead back entirely to Sayyid Muḥammad of
Iṣfahān. His excellence is perfectly clear to these servants, as it was
before. How astonishing is your insight! How well you know people!

\subsection{81}\label{section-576}

\begin{otherlanguage}{persian}

بارى اميدوارم كه حقّ آن مفسد را از ميان بردارد فسادى نموده كه از اوّل
ابداع تا حال شبه آن ديده نشده و البتّه او هم ادّعاى اطّلاع نموده آنقدر
بدانيد كه واللّه الّذى لا اله الّا هو ابداً اطّلاع نداشته آنچه گفته كذب صرف و
افتراى بحت اين چهار روزه عمر قابل نيست كه انسان بفساد و عناد و ارتكاب
اينگونه امور شنيعه مشغول شود شأن تو آن بود كه بعد از ورود آن مفسد بآن
ارض او را بموعظه حسنه و كلمات حكميّه نصيحت نمائى نه آنكه تقويت او نمائى و
مرآت او شوى و حاكى از او گردى انشآءاللّه اميدواريم كه از گندم رى نخورد

\end{otherlanguage}

Bārī omīdvāram keh Ḥaqq ān mofsed rā az mīyān bardārad. Fasādī nemūdeh
keh az avval-e ebdāʿ tā ḥāl shebh-e ān dīdeh nashodeh va albatteh ū ham
eddeʿā-ye eṭṭelāʿ nemūdeh. Ānqadr bedānīd keh, wa-Allāhi alladhī lā
ilāha illā huwa, abadan eṭṭelāʿ nadāshteh, āncheh gofteh kedhb-e ṣerf va
efterā-ye baḥt. Īn chahār rūzeh ʿomr qābel nīst keh ensān be-fasād va
ʿenād va ertekāb-e īngūneh omūr-e shanīʿeh mashghūl shavad. Shaʾn-e to
ān būd keh baʿd az vorūd-e ān mofsed be-ān arż ū rā be-mowʿeẓeh-ye
ḥasaneh va kalemāt-e ḥekmīyyeh naṣīḥat namāʾī, nah ānkeh taqvīyat-e ū
namāʾī va merʾāt-e ū shavī va ḥākī az ū gardī. In shāʾ Allāh omīdvārīm
keh az gandom-e Rey nakhorad.

In any case, I hope that the True One will remove that corrupter from
our midst. He has wrought corruption the like of which has not been seen
from the beginning of creation until now. Of course, he too has claimed
to be informed. Know at least this: by God, besides whom there is no
other god, he has never had any knowledge of the matter; everything he
has said is sheer falsehood and utter fabrication. This fleeting,
four-day life is not worth spending on corruption, hostility, and the
perpetration of such disgraceful deeds. It befitted you, after that
corrupter arrived in that land, to counsel him with goodly exhortations
and words of wisdom, not to strengthen him, become his mirror, and
repeat his words. God willing, we hope that he will not eat the wheat of
Rayy.

\subsection{82}\label{section-577}

\begin{otherlanguage}{persian}

حمد ميكنيم خدا را كه اين عباد را از شرّ و خدعه او حفظ فرمود ديگر اين عبد
چه ذكر كند اين ایّامى است كه نزد متوهّمين قول بعوضه ارض جرزه مسموع است و
كلمة اللّه غير مسموع اگر انصاف داشته باشند جميع بايد خون گريه كنند و در
صحرا‌ها ناله نمايند چه كه بكلمات نقطه اولى روح ماسواه فداه استدلال حقّیّت
از براى دونش مينمايند و بر نفسش سيف ميزنند

\end{otherlanguage}

Ḥamd mīkonīm Khodā rā keh īn ʿebād rā az sharr va khodʿeh-ye ū ḥefẓ
farmūd. Dīgar īn ʿabd cheh dhekr konad? Īn ayyāmī ast keh nazd-e
motavahhemīn qawl-e baʿūżeh-ye arż-e jorozeh masmūʿ ast va kalemat Allāh
ghayr-e masmūʿ. Agar enṣāf dāshteh bāshand jamīʿ bāyad khūn geryeh
konand va dar ṣaḥrāhā nāleh namāyand, cheh keh be-kalemāt-e Noqṭeh-ye
Ūlā, rūḥ-e mā sevāh fedāh, estedlāl-e ḥaqqīyyat az barā-ye dūnash
mīnamāyand va bar nafsash sayf mīzanand.

We praise God for having protected these servants from his malice and
deceit. What more can this servant say? These are days when the deluded
listen to the words of a gnat of the barren land, but do not listen to
the Word of God. If they had any fairness, they would all weep tears of
blood and lament in the wilderness, for they use the words of the Primal
Point---may the souls of all besides Him be sacrificed for Him---to
argue for the truth of someone other than Him, while striking His own
Self with the sword.

\subsection{83}\label{section-578}

\begin{otherlanguage}{persian}

و از اينها گذشته ميخواهيم به بينيم بنقطه اولى بچه حجّت و برهان مؤمنند
فاتوا به و بها لو انتم من الصّادقين بهمان حجّتيكه بنقطه اولى مؤمن و موقنند
ما باعظم از آن بنفس ظهور مؤمن و موقنيم فرقى كه در ميان است اينست كه چون
محك الهى و ميزان عدل ربّانى بميان آمد موحّد و مشرك از هم معلوم شد انصاف
دهيد اين ظهور و شئونات و آياتش را انكار كنيم ديگر بچه رو و بچه امر دعوى
ايمان نمائيم

\end{otherlanguage}

Va az īnhā godhashteh mīkhāhīm bebīnīm be-Noqṭeh-ye Ūlā be-cheh ḥojjat
va borhān moʾmenand. Fa-ʾtū bihi wa-bihā law antum min al-ṣādiqīn.
Be-hamān ḥojjatī keh be-Noqṭeh-ye Ūlā moʾmen va mūqenand mā be-aʿẓam az
ān be-nafs-e ẓohūr moʾmen va mūqenīm. Farqī keh dar mīyān ast īnast keh
chon meḥakk-e elāhī va mīzān-e ʿadl-e rabbānī be-mīyān āmad movaḥḥed va
moshrek az ham maʿlūm shod. Enṣāf dehīd, īn ẓohūr va shoʾūnāt va āyātash
rā enkār konīm dīgar be-cheh rū va be-cheh amr daʿvā-ye īmān namāʾīm?

Beyond all this, we wish to see by what testimony and proof they believe
in the Primal Point. Bring them forth, if you are truthful. By the same
testimony through which they believe with certitude in the Primal Point,
we believe with certitude in the Manifestation Himself, on evidence
greater still. The difference is this: when the divine touchstone and
the balance of divine justice appeared, the believer in God's oneness
was distinguished from the one who joins partners with Him. Judge
fairly: if we deny this Revelation, its manifestations, and its verses,
with what face and on what grounds could we any longer claim to believe?

\subsection{84}\label{section-579}

\begin{otherlanguage}{persian}

و اينكه نوشته بوديد علاوه از آنكه قواعد كلّيه كه ميزان است در دست هست قد
كبر هذا القول منك لانّ هذا مقام قد جعله اللّه فوق العالمين لانّ نفس الظّهور
هو ميزان اللّه و قسطاسه يتقلّب معه الحقّ كما يتقلّب الظلّ مع الشّمس و كذلك نزل
فى البيان ان انت من العارفين بعينه مثل شما مثل كسى است كه يكقطعه حجر
بردارد و پيش نفسى برد كه نزد او كنائز جواهر موجود و دعوى نمايد كه باين
حجر ميخواهم جواهر تو را بشناسم و تميز دهم

\end{otherlanguage}

Va īnkih nivishtih būdīd ʿalāvih az ānkih qavāʿid-i kullīyih kih mīzān
ast dar dast hast, qad kabura hādhā al-qawlu minka li-anna hādhā maqāmun
qad jaʿalahu Allāhu fawqa al-ʿālamīn, li-anna nafsa al-ẓuhūri huwa
mīzānu Allāhi va qisṭāsuhu, yataqallabu maʿahu al-ḥaqqu kamā yataqallabu
al-ẓillu maʿa al-shamsi, va kadhālika nazala fī al-Bayāni in anta mina
al-ʿārifīn. Bi-ʿaynih mathal-i shumā mathal-i kasī ast kih yik qiṭʿih
ḥajar bardārad va pīsh-i nafsī barad kih nazd-i ū kanāʾiz-i javāhir
mawjūd va daʿvā namāyad kih bi-īn ḥajar mīkhvāham javāhir-i tu rā
bishināsam va tamyīz daham.

As for what you wrote, ``Besides, the universal principles that
constitute the balance are at hand''---grave indeed is this statement of
yours, for this is a station that God has placed above all the worlds.
The Manifestation Himself is God's balance and His scales of justice.
Truth turns with Him as the shadow turns with the sun. Thus was it
revealed in the Bayán, if you are among those who understand. Your
likeness is precisely that of someone who picks up a piece of stone,
takes it to a person who possesses treasuries of jewels, and declares,
``With this stone I intend to recognize and distinguish your jewels.''

\subsection{85}\label{section-580}

\begin{otherlanguage}{persian}

آن ميزانيكه در دست خلق است امثال خود را شايد بآن ميزان بشناسد نفس اللّه
مقدّس از آن و ابداً باو شناخته نشده و نخواهد شد ميزان الهى لم‌يزل نفس او
بوده و خواهد بود من عرفه فقد عرف الميزان و من احتجب ضلّ سعيه و كان من
الخاسرين

\end{otherlanguage}

Ān mīzānī kih dar dast-i khalq ast amthāl-i khud rā shāyad bi-ān mīzān
bishināsad. Nafs-i Allāh muqaddas az ān va abadan bi-ū shinākhtih
nashudih va nakhvāhad shud. Mīzān-i ilāhī lam-yazal nafs-i ū būdih va
khvāhad būd. Man ʿarafahu fa-qad ʿarafa al-mīzāna va man iḥtajaba ḍalla
saʿyuhu va kāna mina al-khāsirīn.

The balance in people's hands may serve to recognize others like
themselves. God's own Self is sanctified above it and has never been,
nor ever will be, known through it. The divine balance has always been,
and will always be, His own Self. Whoever has recognized Him has
recognized the balance; whoever remains veiled has wasted his efforts
and is among the losers.

\subsection{86}\label{section-581}

\begin{otherlanguage}{persian}

يقين ميزانى كه ذكر نمود‌ه‌ايد از همان ميزانى است كه در حين ظهور نقطه اولى
روح ما سواه فداه در دست علماى عصر بود و هم‌چنين در ظهور رسول اللّه من قبل
و از قبل او ظهور عيسى بن مريم كه جميع علماء و رؤسا بميزان انفس خود بر آن
شموس حقيقت اعتراض نموده تا آنكه بالأخره فتوى بر قتل آن هياكل قدسيّه دادند
الا لعنة اللّه على الظّالمين

\end{otherlanguage}

Yaqīn mīzānī kih dhikr namūdih-īd az hamān mīzānī ast kih dar ḥīn-i
ẓuhūr-i Nuqṭiy-i Ūlā, rūḥ-i mā sivāhu fidāh, dar dast-i ʿulamāy-i ʿaṣr
būd va ham-chunīn dar ẓuhūr-i Rasūl-i Allāh min qabl va az qabl-i ū
ẓuhūr-i ʿĪsā ibn-i Maryam kih jamīʿ-i ʿulamā va ruʾasā bi-mīzān-i
anfus-i khud bar ān shumūs-i ḥaqīqat iʿtirāḍ namūdih tā ānkih
bi-al-ākharah fatvā bar qatl-i ān hayākil-i qudsīyih dādand. Alā laʿnatu
Allāhi ʿalā al-ẓālimīn.

Surely the balance you have mentioned is the very balance that was in
the hands of the learned of the age when the Primal Point appeared---may
the soul of all besides Him be sacrificed for Him---and likewise at the
earlier appearance of the Messenger of God and, before Him, the
appearance of Jesus son of Mary. All the learned and the leaders used
the balance of their own selves to raise objections against those Suns
of Truth, until at last they issued decrees for the killing of those
sacred Temples. Indeed, God's curse rests upon the wrongdoers.

\subsection{87}\label{section-582}

\begin{otherlanguage}{persian}

اى راقد غافل بدانكه ميزان الهى لم‌يزل نفس او و ما يظهر من عنده بوده و
خواهد بود كلّ ذلك فضلاً من لدنه ان انتم تعلمون و اگر فى‌الحقيقه نفسى بنقطه
بيان و وصاياى او ناظر باشد اگر كلّ عالمين باسياف شاحذه باو رو نمايند و
جميع اعضايش را قطعه قطعه سازند توقّف در اين ظهور نخواهد نمود چه كه آنچه
مرقوم فرموده‌اند و از سماء مشيّت نازل شده كلّ صريح باين ظهور اعظم بوده و
خواهد بود

\end{otherlanguage}

Ay rāqid-i ghāfil, bidānkih mīzān-i ilāhī lam-yazal nafs-i ū va mā
yaẓharu min ʿindihi būdih va khvāhad būd. Kullu dhālika faḍlan min
ladunhu in antum taʿlamūn. Va agar fī-al-ḥaqīqih nafsī bi-Nuqṭiy-i Bayān
va vaṣāyāy-i ū nāẓir bāshad agar kull-i ʿālamīn bi-asyāf-i shāḥidhih
bi-ū rū namāyand va jamīʿ-i aʿḍāyash rā qiṭʿih qiṭʿih sāzand tavaqquf
dar īn ẓuhūr nakhvāhad namūd, chih kih ānchih marqūm farmūdih-and va az
samāʾ-i mashīyyat nāzil shudih kull ṣarīḥ bi-īn ẓuhūr-i aʿẓam būdih va
khvāhad būd.

O heedless sleeper! Know that the divine balance has always been, and
will always be, His own Self and whatever appears from Him. All this is
a bounty from Him, if you only knew. If someone truly looks to the Point
of the Bayán and His counsels, even if all the peoples of the worlds
should turn upon him with sharpened swords and cut every limb of his
body to pieces, he would not hesitate over this Revelation. For
everything He wrote, everything sent down from the heaven of His will,
has been and will remain an explicit reference to this Most Great
Revelation.

\subsection{88}\label{section-583}

\begin{otherlanguage}{persian}

و اگر ناظر ببيان نيستيد و حجّت الهى را كه كلّ انبیاء و رسل بآن اتيان نموده
انكار مينمائيد و آيات اللّه را كه لم‌يزل و لايزال حجّت محكمه بوده ميزان
نميدانيد فاتوا بما عندكم لو انتم من الصّادقين البتّه بر شما لازم است كه آن
ميزان و آن حجّتى كه اليوم بآن حقّیّت خود را ثابت مينمائيد بياوريد لا واللّه
لن تستطيعن و لن تقدرن‌ّ كلّ الأمور فى قبضته و كلّ شیء فى كفّ اقتداره و ما
دونه مخلوق عنده و بذلك يشهد كلّ ذى بصر عليم

\end{otherlanguage}

Va agar nāẓir bi-Bayān nīstīd va ḥujjat-i ilāhī rā kih kull-i anbiyā va
rusul bi-ān ityān namūdih inkār mīnamāyīd va āyāt-i Allāh rā kih
lam-yazal va lā-yazāl ḥujjat-i muḥkamih būdih mīzān nimīdānīd, faʾtū
bi-mā ʿindakum law antum mina al-ṣādiqīn. Al-battih bar shumā lāzim ast
kih ān mīzān va ān ḥujjatī kih al-yawm bi-ān ḥaqqīyyat-i khud rā thābit
mīnamāyīd biyāvarīd. Lā va-Allāhi lan tastaṭīʿunna va lan taqdiranna.
Kullu al-umūri fī qabḍatihi va kullu shayʾin fī kaffi iqtidārihi va mā
dūnahu makhlūqun ʿindahu va bi-dhālika yashhadu kullu dhī baṣarin ʿalīm.

If you do not look to the Bayán, if you deny the divine proof that all
the Prophets and Messengers have brought, and if you do not regard the
verses of God---which have always been and will always remain a
conclusive proof---as the balance, then bring forward what you possess,
if you are truthful. You must certainly produce that balance and that
proof by which you establish the truth of your own claim today. No, by
God! You will never be able to do so, nor will you have the power. All
affairs are in His grasp, and all things are in the hand of His might.
Everything besides Him is a creature in His presence; to this every
discerning person of insight bears witness.

\subsection{89}\label{section-584}

\begin{otherlanguage}{persian}

اگر نفسى اليوم اين شئونات الهيّه و ظهورات ربوبيّه و آيات منزله‌ئيكه در اين
ظهور امنع اقدس ظاهر شده و بمثل غيث هاطل جارى انكار نمايد انكار حق و كلّ
رسل را نموده و انّه من المشركين قد كان فى ام الألواح من قلم الأمر مرقوماً

\end{otherlanguage}

Agar nafsī al-yawm īn shuʾūnāt-i ilāhīyih va ẓuhūrāt-i rubūbīyih va
āyāt-i munazzalih-ī kih dar īn ẓuhūr-i amnaʿ-i aqdas ẓāhir shudih va
bi-mithl-i ghayth-i hāṭil jārī inkār namāyad, inkār-i Ḥaqq va kull-i
rusul rā namūdih va innahu mina al-mushrikīn qad kāna fī Ummi al-Alwāḥi
min qalami al-amri marqūman.

If anyone today denies these divine expressions, these manifestations of
lordship, and these revealed verses that have appeared in this most
exalted, most holy Revelation and flow like pouring rain, he has denied
God and all the Messengers. He is indeed among those who associate
partners with God: this has been inscribed in the Mother of Tablets by
the Pen of the Cause.

\subsection{90}\label{section-585}

\begin{otherlanguage}{persian}

يك نغمه از نغمات ورقاء بيان در ذكر ميزان مخصوص قاصدان حرم رحمن ذكر ميشود
لعلّ بمعرفت آن فائز شده از مظاهر شيطانيّه بكلّيه منقطع شوند و حجبات اوهام
را شقّ نمايند على شأن يسمعنّ اهل ملأ الأعلی صوت خرقها و يشهدنّ انفسهم فوق
ملأ البيان كلّها من الّذينهم كفروا و اشركوا و يفرّنّ من غصن الى غصن ليدخلنّ
الشّكّ و الرّيب فى قلوب الّذين قلّبتهم نسمات الرّحمن عن شمال الظّنّ الى يمين
اليقين قوله عزّ ذكره و عزّ ثنائه و عزّ كلمته و اعظم احسانه او است ميزان و
امر اوست ميزان و احوال او است ميزان و دلالات او است ميزان ما ينسب الى
الميزان ميزان انتهى

\end{otherlanguage}

Yik naghmih az naghamāt-i Varqāy-i Bayān dar dhikr-i mīzān makhṣūṣ-i
qāṣidān-i ḥaram-i Raḥmān dhikr mīshavad laʿall bi-maʿrifat-i ān fāʾiz
shudih az maẓāhir-i shayṭānīyih bi-kullīyih munqaṭiʿ shavand va
ḥujubāt-i awhām rā shaqq namāyand, ʿalā shaʾnin yasmaʿanna ahlu Malaʾi
al-Aʿlā ṣawta kharqihā va yashhadanna anfusahum fawqa malaʾi al-Bayāni
kullihā mina alladhīnahum kafarū va ashrakū va yafirrunna min ghuṣnin
ilā ghuṣnin li-yudkhilanna al-shakka va al-rayba fī qulūbi alladhīna
qallabathum nasamātu al-Raḥmāni ʿan shimāli al-ẓanni ilā yamīni
al-yaqīn. Qawluhu ʿazza dhikruhu va ʿazza thanāʾuhu va ʿazzat kalimatuhu
va aʿẓama iḥsānahu: Ū ast mīzān va amr-i ūst mīzān va aḥvāl-i ū ast
mīzān va dalālāt-i ū ast mīzān. Mā yunsabu ilā al-mīzāni mīzān. Intahā.

One melody from the melodies of the Dove of the Bayán concerning the
balance is recounted especially for those who seek the sanctuary of the
All-Merciful, so that they may attain an understanding of it, detach
themselves entirely from satanic manifestations, and tear apart the
veils of vain imaginings in such a way that the inhabitants of the
Concourse on high will hear the sound of their rending, and they will
see themselves above the entire assembly of the Bayán, consisting of
those who have disbelieved and associated partners with God and who flee
from branch to branch to instill doubt and misgiving in the hearts of
those whom the breezes of the All-Merciful have turned from the left
hand of conjecture to the right hand of certitude. His words---exalted
be His remembrance, exalted be His praise, exalted be His word, and how
great His bounty---are: ``He is the balance, His command is the balance,
His conditions are the balance, and His indications are the balance.
Whatever is related to the balance is a balance.'' End of quotation.

\subsection{91}\label{section-586}

\begin{otherlanguage}{arabic}

بارى ميزان نفس ظهور بوده و خواهد بود و آنچه از بحر جود وكرم ظاهر فرمايد
و بآنچه اليوم عاملند و بآنچه ناطقند و بآنچه آمرند ميزان كلّ از قبل و بعد
بوده و خواهد بود و اليوم هر نفسى را از اين ميزان قدس رحمانى خارج مشاهده
نمائى باطل دان و مردود شمر و نعوذ باللّه باصل ميزان جسارت منما خف عن اللّه
الّذى خلقك من قطرة ماء مهين

\end{otherlanguage}

Bārī mīzān nafs-i ẓuhūr būdih va khvāhad būd va ānchih az baḥr-i jūd va
karam ẓāhir farmāyad va bi-ānchih al-yawm ʿāmiland va bi-ānchih nāṭiqand
va bi-ānchih āmirand mīzān-i kull az qabl va baʿd būdih va khvāhad būd
va al-yawm har nafsī rā az īn mīzān-i quds-i raḥmānī khārij mushāhidih
namāʾī bāṭil dān va mardūd shumur va naʿūdhu bi-Allāh bi-aṣl-i mīzān
jasārat manamā. Khaf ʿan Allāhi alladhī khalaqaka min qaṭrati māʾin
mahīn.

In short, the balance has been and will be the Manifestation Himself and
whatever He brings forth from the ocean of bounty and generosity. What
He does today, what He says, and what He commands have been and will be
the balance for all, both before and after. Today, regard anyone you see
standing outside this holy, all-merciful balance as false, and count him
as rejected. Do not---God forbid---presume to assail the very source of
the balance. Fear God, who created you from a drop of lowly water.

\subsection{92}\label{section-587}

\begin{otherlanguage}{persian}

قول ديگران مبغض اوّلاً آنكه در اصل مدّعى و مطلب عرض نمايم همان حضرت باب كه
شما ربّ اعلى ميگوئيد در حقّیّت خود آيات خود را مقرّر فرموده‌اند و بآيات اثبات
حقّیّت خود فرموده‌و هيچ كسى را از معتقدين انكار بر اين نيست و يقين است كه
شما هم نميتوانيد انكار نمائيد بلكه تصديق شما و ديگران حضرت باب را بآيات
بوده انتهى

\end{otherlanguage}

Qawl-i dīgarān-i mubghiḍ: Avvalan ānkih dar aṣl-i muddaʿā va maṭlab ʿarḍ
namāyam hamān Ḥaḍrat-i Bāb kih shumā Rabb-i Aʿlā mīgūʾīd dar ḥaqqīyyat-i
khud āyāt-i khud rā muqarrar farmūdih-and va bi-āyāt ithbāt-i
ḥaqqīyyat-i khud farmūdih va hīch kasī rā az muʿtaqidīn inkār bar īn
nīst va yaqīn ast kih shumā ham nimītavānīd inkār namāʾīd balkih
taṣdīq-i shumā va dīgarān Ḥaḍrat-i Bāb rā bi-āyāt būdih. Intahā.

The words of other hostile persons: ``First, let me address the central
claim and issue. That same Báb whom you call the Most Exalted Lord
established His own verses as the proof of His truth and demonstrated
His truth through those verses. None of the believers denies this, and
it is certain that you too cannot deny it. Indeed, your acceptance of
the Báb, and that of others, was based on His verses.'' End of
quotation.

\subsection{93}\label{section-588}

\begin{otherlanguage}{persian}

فو‌الّذى نفسى بيده كه هيچ ذى شمّى از اين كلمات كه مرقوم داشته رايحه حقّ
استشمام ننمايد زهى حسرت و ندامت كه در اين ربيع روحانى بطراز بديع مزيّن
نشدى و از نفحات خوش رضوان معانى استشمام ننمودى كلماتت مثل كلمات محتجبين
بل احقر فو‌اللّه لو تتفكّر فيها لتكون خجلاً فى نفسك و تمحو كلّها قدر و شأن
انسان از كلماتش ظاهر و فی‌‌الحقيقه كلمه مرآت نفس است لو انت من العارفين

\end{otherlanguage}

Fa-va-alladhī nafsī bi-yadihi kih hīch dhī shammī az īn kalimāt kih
marqūm dāshtih rāyiḥiy-i Ḥaqq istishmām nanamāyad. Zihī ḥasrat va
nadāmat kih dar īn rabīʿ-i rūḥānī bi-ṭirāz-i badīʿ muzayyan nashudī va
az nafaḥāt-i khush-i Riḍvān-i maʿānī istishmām nanamūdī. Kalimātat
mithl-i kalimāt-i muḥtajibīn bal aḥqar. Fa-va-Allāhi law tatafakkaru
fīhā la-takūnu khajilan fī nafsika va tamḥū kullahā. Qadr va shaʾn-i
insān az kalimātash ẓāhir va fī-al-ḥaqīqih kalimih mirʾāt-i nafs ast,
law anta mina al-ʿārifīn.

By Him in whose hand is my soul! No one with a sense of smell could
perceive the fragrance of truth in these words you have written. What
sorrow and regret that, in this spiritual springtime, you have not
adorned yourself with a wondrous adornment or inhaled the sweet breezes
of the Riḍván of inner meanings! Your words resemble those of the
veiled---indeed, they are baser. By God! If you reflected on them, you
would feel ashamed within yourself and erase them all. A person's worth
and station are revealed by his words. In truth, a word is the mirror of
the self, if you are among those who understand.

\subsection{94}\label{section-589}

\begin{otherlanguage}{arabic}

لم يزل باين كلمات غير متقنه هياكل ظنونيّه ناس را از شطر احديّه منع
نمود‌ه‌اند چه مقام مشاهده ميشود تقديس امر و چه مقام مشاهده ميشود انفس
محتجبه اسمع قولى ثمّ دعها عن ورائك فاقبل الى اللّه بكلّك لعلّ يجرى من قلبك
بحور المعانى و البيان فى ذكر ربّك المقتدر العزيز المنّان

\end{otherlanguage}

Lam-yazal bi-īn kalimāt-i ghayr-i mutqanih hayākil-i ẓunūnīyih nās rā az
shaṭr-i aḥadīyih manʿ namūdih-and. Chih maqām mushāhidih mīshavad
taqdīs-i amr va chih maqām mushāhidih mīshavad anfus-i muḥtajibih. Ismaʿ
qawlī thumma daʿhā ʿan warāʾika fa-aqbil ilā Allāhi bi-kullika laʿalla
yajrī min qalbika buḥūru al-maʿānī va al-bayāni fī dhikri Rabbika
al-Muqtadiri al-ʿAzīzi al-Mannān.

With such unsound words, these embodiments of conjecture have always
kept people from turning toward the divine Unity. At what a station is
the sanctity of the Cause seen, and at what a station these veiled
souls! Hear my words, then leave those things behind you and turn with
your whole being toward God, so that oceans of inner meanings and
utterance may flow from your heart in remembrance of your Lord, the
All-Powerful, the Almighty, the Bestower.

\subsection{95}\label{section-590}

\begin{otherlanguage}{persian}

اگر چه قلم خجلت ميكشد در ذكر جواب چنين كلمات و لكن نظر بآنكه جواب را
واجب دانسته‌اند مذكور ميدارم كه همين كلمات شما ردّ بر شما است و اثبات
امراللّه نوشته و ملتفت نشده كذلك يجرى اللّه من القلم ما يكون غافلاً عنه
كاتبه فويل للمحتجبين من عذاب يوم عظيم

\end{otherlanguage}

Agar chih qalam khijlat mīkishad dar dhikr-i javāb-i chunīn kalimāt va
lākin naẓar bi-ānkih javāb rā vājib dānistih-and madhkūr mīdāram kih
hamīn kalimāt-i shumā radd bar shumā ast va ithbāt-i Amru'llāh nivishtih
va multafit nashudih. Kadhālika yujrī Allāhu mina al-qalami mā yakūnu
ghāfilan ʿanhu kātibuhu fa-waylun li-l-muḥtajibīna min ʿadhābi yawmin
ʿaẓīm.

Although the pen feels ashamed to set down a reply to such words, since
a reply has been deemed necessary, I state that these very words of
yours refute you. You have written a proof of the Cause of God without
realizing it. Thus God causes to flow from the pen what its writer is
unaware of. Woe, then, to the veiled, from the punishment of a momentous
day!

\subsection{96}\label{section-591}

\begin{otherlanguage}{persian}

اى برادر مرئه عقيمه در كلّ ملل بی‌قدر بوده چه كه محروم گشته از نسل ظاهره
فواللّه نفوسى كه قلوبشان عقيم گشته از ظهور ابكار معارف الهيّه اين نفوس اشدّ
محروماً بنظر آيند و حال مشاهده ميشود اكثرى از محتجبين قلب و فؤاد و بصرشان
عقيم گشته بشأنيكه آنچه خود مينويسند بادراك آن قادر نيستند تا چه رسد
باشارات كلمات الهيّه و جواهر حكم بالغه ربّانيّه تاللّه اولئك عقيم من رحمة
اللّه و من الصّدق و الانصاف و لم يظهر منهم الّا ما يشتدّ به غضب اللّه و سخطه
عليهم و لكنّهم لا يفقهون

\end{otherlanguage}

Ay barādar, marʾih-yi ʿaqīmih dar kull-i milal bī-qadr būdih chih kih
maḥrūm gashtih az nasl-i ẓāhirih. Fa-va-Allāh nufūsī kih qulūbishān
ʿaqīm gashtih az ẓuhūr-i abkār-i maʿārif-i ilāhīyih īn nufūs ashaddu
maḥrūman bi-naẓar āyand va ḥāl mushāhidih mīshavad aktharī az muḥtajibīn
qalb va fuʾād va baṣarishān ʿaqīm gashtih bi-shaʾnī kih ānchih khud
mīnivīsand bi-idrāk-i ān qādir nīstand tā chih rasad bi-ishārāt-i
kalimāt-i ilāhīyih va javāhir-i ḥikam-i bālighih-yi rabbānīyih.
Ta-Allāhi ulāʾika ʿaqīmun min raḥmati Allāhi va mina al-ṣidqi va
al-inṣāfi va lam yaẓhar minhum illā mā yashtaddu bihi ghaḍabu Allāhi va
sakhaṭuhu ʿalayhim va lākinnahum lā yafqahūn.

O brother! A barren woman has been held in little esteem among all
peoples because she has been deprived of physical offspring. By God!
Those whose hearts have become barren of the emergence of the virginal
truths of divine knowledge appear far more deprived. Now it can be seen
that the hearts, inner understanding, and sight of most of the veiled
have become so barren that they cannot comprehend what they themselves
write, much less the allusions of the divine words and the gems of
consummate heavenly wisdom. By God! These people are barren of God's
mercy, of truthfulness, and of fairness. Nothing has appeared from them
but what intensifies God's anger and wrath against them, yet they do not
understand.

\subsection{97}\label{section-592}

\begin{otherlanguage}{persian}

اينكه نوشته همان حضرت باب كه شما ربّ اعلى ميدانيد از اينكلمه معلوم ميشود
كه شما ربّ اعلى نميدانيد و يا تقيّه نمود‌ه‌ايد مثل مرشدين شما كه در بعضى
مواضع انكار مينمايند و تبرّى ميجويند و باطراف پشته پشته كتب مجعوله در
اثبات حقّیّت خود ميفرستند شما و كلّ من فى السّموات و الأرض جميعاً بدانيد بانّا
كنّا موقناً معترفاً مذعناً ناطقاً ذاكراً قائلاً منادياً مضجّاً مصرخاً مصحاً متكلّماً
مبلّغاً معجّاً بأعلى الصّوت بأنّه هو ربّ الأعلی و سدرة المنتهى و شجرة القصوى و
ملكوت العلى و جبروت العمآء و لاهوت البقآء و روح البهآء و سرّ الأعظم و
كلمة الأتمّ و مظهر القدم و هيكل الأکرم و رمز المنمنم و ربّ الأمم و البحر
الملتطم و كلمة العليا و درّة الأولی و صحيفة المكنون و كتاب المخزون جمال
الأحديّه و مظهر الهويّة و مطلع الصمديّة لولاه ما ظهر الوجود و ما عرف
المقصود و ما برز جمال المعبود تاللّه باسمه قد خلقت السّمآء و ما فيها و
الأرض و من عليها و به موّجت البحار و جرت الأنهار و اثمرت الأشجار و به
حقّقت الأدیان و ظهر جمال الرّحمن فواللّه لو نصفه الى آخر الّذى لا آخر له لن
يسكن فؤادى من عطش حبّ ذكر اسمائه و صفاته فكيف نفسه المقدّس العزيز الجميل

\end{otherlanguage}

Īnkih nivishtih hamān Ḥaḍrat-i Bāb kih shumā Rabb-i Aʿlā mīdānīd az īn
kalimih maʿlūm mīshavad kih shumā Rabb-i Aʿlā nimīdānīd va yā taqīyyih
namūdih-īd mithl-i murshidīn-i shumā kih dar baʿḍī mavāḍiʿ inkār
mīnamāyand va tabarrī mījūyand va bi-aṭrāf pushtih pushtih kutub-i
majʿūlih dar ithbāt-i ḥaqqīyyat-i khud mīfiristand. Shumā va kull-i man
fī al-samāvāti va al-arḍ jamīʿan bidānīd bi-annā kunnā mūqinan
muʿtarifan mudhʿinan nāṭiqan dhākiran qāʾilan munādiyan muḍijjan
muṣrikhan muṣiḥḥan mutakalliman muballighan muʿijjan bi-aʿlā al-ṣawti
bi-annahu huwa Rabbu al-Aʿlā va Sidratu al-Muntahā va Shajaratu al-Quṣwā
va Malakūtu al-ʿUlā va Jabarūtu al-ʿAmāʾ va Lāhūtu al-Baqāʾ va Rūḥu
al-Bahāʾ va Sirru al-Aʿẓam va Kalimatu al-Atamm va Maẓharu al-Qidam va
Haykalu al-Akram va Ramzu al-Munamnam va Rabbu al-Umam va al-Baḥru
al-Mulṭaṭim va Kalimatu al-ʿUlyā va Durratu al-Ūlā va Ṣaḥīfatu al-Maknūn
va Kitābu al-Makhzūn, Jamāl-i Aḥadīyih va Maẓhar-i Huvīyyat va Maṭlaʿ-i
Ṣamadīyyat. Lawlāhu mā ẓahara al-wujūdu va mā ʿurifa al-Maqṣūdu va mā
baraza jamālu al-Maʿbūd. Ta-Allāhi bi-ismihi qad khuliqati al-samāʾu va
mā fīhā va al-arḍu va man ʿalayhā va bihi muwwijati al-biḥāru va jarati
al-anhāru va athmarati al-ashjāru va bihi ḥuqqiqati al-adyānu va ẓahara
jamālu al-Raḥmān. Fa-va-Allāhi law naṣifuhu ilā ākhiri alladhī lā ākhira
lahu lan yaskuna fuʾādī min ʿaṭashi ḥubbi dhikri asmāʾihi va ṣifātihi
fa-kayfa nafsihi al-muqaddasi al-ʿAzīzi al-Jamīl.

You wrote, ``That same Báb whom you regard as the Most Exalted Lord.''
From these words it is clear that you do not regard Him as the Most
Exalted Lord, or else you have concealed your belief, like your guides,
who in some places deny Him and disavow Him, while sending stacks upon
stacks of fabricated books in every direction to establish the truth of
their own claims. You, and all who are in the heavens and on earth, know
that We have been certain, acknowledging, assenting, declaring,
remembering, stating, proclaiming, crying aloud, calling out, affirming,
speaking, conveying, and raising Our voice at its highest pitch that He
is the Most Exalted Lord, the Lote-Tree of the Utmost Boundary, the
Farthest Tree, the Kingdom on high, the Dominion of the Cloud, the
Divinity of everlastingness, the Spirit of Glory, the Most Great Secret,
the Most Perfect Word, the Manifestation of eternity, the Most Bountiful
Temple, the finely adorned Mystery, the Lord of the nations, the surging
Ocean, the Most Exalted Word, the Primal Pearl, the Hidden Scroll, the
Treasured Book, the Beauty of Oneness, the Manifestation of the divine
Identity, and the Dawning-Place of eternal Self-Sufficiency. Were it not
for Him, existence would not have appeared, the Desired One would not
have been known, and the beauty of the Worshipped One would not have
emerged. By God! Through His name the heaven and all that is within it,
and the earth and all who dwell upon it, were created. Through Him the
seas were made to surge, the rivers flowed, and the trees bore fruit.
Through Him the religions were established and the beauty of the
All-Merciful appeared. By God! Were We to describe Him until the end
that has no end, my heart's thirst in its love of mentioning His names
and attributes would not be stilled---how much less its thirst for His
own sanctified, mighty, beautiful Self!

\subsection{98}\label{section-593}

\begin{otherlanguage}{persian}

و شما شاهد و گواه باشيد همين لوح را بهر نفسى از نفوس معرضه كه بخواهيد
بنمائيد مرخّصيد چنانچه مرشدت الواح الهيه بخطّ غصن ربانيّه اخذ نموده و ببعضى
نشان داده و خود تبرّى جسته هذا ما فعل فى السّر و لكن فى الجهر يخضع رأسه و
يطوّل نفسه بين ملأ البيان الّذين اتّخذوه ربا لانفسهم و يستدلّ لهم اثبات نفسه
المجعولة المكذوبة و يفرح برياسته عليهم قل فويل لك يا ايّها الخائف الخائن
الخائب الخاسر ما يغنيك اليوم شیء و لو تتمسّك بكلّ ما خلق بين السّموات و
الأرضين

\end{otherlanguage}

Va shumā shāhid va gavāh bāshīd hamīn lawḥ rā bi-har nafsī az nufūs-i
muʿriḍih kih bikhvāhīd binamāʾīd murakhkhaṣīd chunānchih murshidat
alvāḥ-i ilāhīyih bi-khaṭṭ-i ghuṣn-i rabbānīyih akhdh namūdih va bi-baʿḍī
nishān dādih va khud tabarrī justih. Hādhā mā faʿala fī al-sirri va
lākin fī al-jahri yakhḍaʿu raʾsahu va yuṭawwilu nafsahu bayna malaʾi
al-Bayāni alladhīna ittakhadhūhu rabban li-anfusihim va yastadillu lahum
ithbāta nafsihi al-majʿūlati al-makdhūbati va yafraḥu bi-riyāsatihi
ʿalayhim. Qul fa-waylun laka yā ayyuhā al-khāʾifu al-khāʾinu al-khāʾibu
al-khāsir, mā yughnīka al-yawma shayʾun va law tatamassaku bi-kulli mā
khuliqa bayna al-samāvāti va al-arḍīn.

Be a witness and bear testimony: you have permission to show this very
Tablet to any of the souls who have turned away, as you wish. Your guide
obtained divine Tablets in the handwriting of the heavenly Branch and
showed them to certain people, while himself disavowing them. This is
what he did in secret. In public, however, he bows his head and exalts
himself among the people of the Bayán, who have taken him as a lord for
themselves. He offers them arguments to establish his own fabricated and
false self, and rejoices in his leadership over them. Say: Woe to you, O
fearful, treacherous, disappointed loser! Nothing will avail you today,
even if you cling to everything created between the heavens and the
earths.

\subsection{99}\label{section-594}

\begin{otherlanguage}{persian}

و ايكاش بهمين كفايت مينمودند بلكه كلماتى شبه آن نوشته و باسم حقّ ببعضى
نموده انّ اللّه يعلم ما فى سرّه و احصى كلّ شیء فى امام مبين حال شما هم اگر
اراده نمائيد اين لوح را بهر نفسيكه اراده كنيد بنمائيد فو‌اللّه انّ ربی‌ ما
حفظ نفسه و ما يحفظ ابداً و ينتظر الصّليب من اليهود و السّنان من السّنان لو
انت من العارفين بدانكه حقّ منتظر است آنچه را كه بظهور قبلش وارد شده و
هم‌چنين آنچه را كه برسول اللّه از قبل وارد شده و هم‌چنين از قبل او بروح ثمّ
من قبله بالخليل ما خاف من احد و لن يخاف بفضل اللّه و عنايته اذا فاجتمعوا
لسفك دمه المطهّر العزيز المنيع

\end{otherlanguage}

Va ey kāsh be-hamīn kefāyat mī-nemūdand balke kalemātī shebh-e ān
neveshte va be-esm-e Ḥaqq be-baʿżī nemūde. Inna Allāha yaʿlamu mā fī
sirrihi wa-aḥṣā kulla shayʾin fī imāmin mubīn. Ḥāl shomā ham agar erāde
namāyīd īn lowḥ rā be-har nafsī ke erāde konīd benamāyīd. Fa-wa-Allāhi
inna rabbī mā ḥafiẓa nafsahu wa-mā yaḥfaẓu abadan wa-yantaẓiru al-ṣalība
min al-Yahūd wa-l-sināna min al-sinān law anta min al-ʿārifīn. Bedān ke
Ḥaqq montaẓer ast ānche rā ke be-ẓohūr-e qablash vāred shode va
ham-chenīn ānche rā ke be-Rasūl Allāh az qabl vāred shode va ham-chenīn
az qabl-e ū be-Rūḥ thumma min qablihi bi-l-Khalīl. Mā khāfa min aḥadin
wa-lan yakhāfa bi-faḍli Allāhi wa-ʿināyatihi. Idhan fa-jtamiʿū li-safki
damihi al-muṭahhari al-ʿazīzi al-manīʿ.

Would that they had been content with this! Instead, they wrote words
resembling those and showed them to certain people in the name of the
True One. God knows what is in his secret heart and has recorded all
things in a manifest register. Now you too, if you wish, may show this
Tablet to any soul you choose. By God! My Lord has not protected
Himself, nor will He ever protect Himself, and He awaits the cross from
the Jews and the spear from the spear, if you are among those who
understand. Know that the True One awaits what befell His previous
Manifestation, and likewise what befell the Messenger of God before Him,
and, before Him, the Spirit, and before Him, the Friend. He has feared
no one and will never fear, through God's grace and loving care. Gather,
then, to shed His purified, mighty, inviolable blood.

\subsection{100}\label{section-595}

\begin{otherlanguage}{persian}

و اينكه نوشته بودى در حقّیّت خود آيات را مقرّر فرموده‌اند و بآيات اثبات حقّیّت
خود فرموده همين كلمه لو تعرف در اثبات اين امر از قلم جارى شده و شما
ادراك ننموده‌‌ايد چه كه اين آياتى كه ذكر نموده شما بآن موقنيد يا نه اگر
نيستيد جميع اين كلمات كه ذكر نموده‌ايد لغو بوده و خواهد بود و اگر موقنيد
چرا باين آيات منزله كه بمثابه غيث هاطله در كلّ حين نازل است منكر و معرضيد
يقين بدان كه اهل بصر و فؤاد فرصتت نميدهند كه از كلمه اوّل تجاوز نمائى در
همين يك كلمه تو برهان الهى و حجّت او را بر تو ثابت و مسجّل مينمايند بهمان
نحو كه آن شجره احديّه اثبات امر خود فرموده و ما عارف شده‌ايم حال هم بهمان
قسم امر اللّه ظاهر و آيات اللّه نازل و اين عباد موقن شده‌اند اهل انكار اين
عباد نبوده و نخواهند بود منكر نفوسى هستند كه آمنوا ببعض الكتاب و كفروا
ببعض لا واللّه بل كفروا بكلّ الكتب چه كه اليوم اگر نفسى از حرفى از اين
آيات اعراض نمايد اعراضش بر كلّ ثابت و محقّق است انّا آمنّا بما نزل من قبل و
حينئذ و يشهد بذلك كلّ شیء ان انت من الموقنين خوب استدلال مينمائى و بمحبوب
نقطه اولى باعتراض و جدال مشغولى و مينويسى كه بآيات حقّ او ثابت شده بگمان
خود بتنزيل ناظرى و لكن از منزل بسى بعيد

\end{otherlanguage}

Va īnke neveshte būdī dar ḥaqqīyat-e khod āyāt rā moqarrar farmūde-and
va be-āyāt ethbāt-e ḥaqqīyat-e khod farmūde, hamīn kaleme, law taʿrifu,
dar ethbāt-e īn amr az qalam jārī shode va shomā edrāk na-nemūde-īd. Che
ke īn āyātī ke dhekr nemūde shomā be-ān mūqenīd yā na? Agar nīstīd
jamīʿ-e īn kalemāt ke dhekr nemūde-īd laghv būde va khāhad būd, va agar
mūqenīd cherā be-īn āyāt-e monzale ke be-mathābe-ye gheyth-e hāṭele dar
koll-e ḥīn nāzel ast monker va moʿreżīd? Yaqīn bedān ke ahl-e baṣar va
foʾād forṣatat ne-mī-dehand ke az kaleme-ye avval tajāvoz namāyī. Dar
hamīn yek kaleme-ye to borhān-e elāhī va ḥojjat-e ū rā bar to thābet va
mosajjal mī-namāyand. Be-hamān naḥv ke ān shajare-ye aḥadīye ethbāt-e
amr-e khod farmūde va mā ʿāref shode-īm, ḥāl ham be-hamān qesm amr Allāh
ẓāher va āyāt Allāh nāzel va īn ʿebād mūqen shode-and. Ahl-e enkār īn
ʿebād na-būde va na-khāhand būd. Monker nofūsī hastand ke āmanū bi-baʿḍi
al-kitāb wa-kafarū bi-baʿḍ. Lā wa-Allāh bal kafarū bi-kulli al-kutub.
Che ke al-yowm agar nafsī az ḥarfī az īn āyāt eʿrāż namāyad eʿrāżash bar
koll thābet va moḥaqqaq ast. Innā āmannā bi-mā nazala min qablu
wa-ḥīnaʾidhin wa-yashhadu bi-dhālika kullu shayʾin in anta min
al-mūqinīn. Khūb estedlāl mī-namāyī va be-maḥbūb-e Noqṭe-ye Ūlā
be-eʿterāż va jedāl mashghūlī va mī-nevīsī ke be-āyāt ḥaqq-e ū thābet
shode. Be-gomān-e khod be-tanzīl nāẓerī va lāken az monzel basī baʿīd.

As for your statement that He established verses as the criterion of His
own truth and proved His own truth through verses: these very words, if
you understood, have flowed from the Pen in proof of this Cause, but you
have not grasped them. For are you certain of these verses you have
mentioned, or not? If you are not, all these words you have written have
been and will remain futile. If you are certain, why do you deny and
turn away from these revealed verses, which descend at every moment like
pouring rain? Know for certain that those endowed with sight and
understanding will not give you the chance to proceed beyond your first
statement. Through this single statement of yours, they establish and
confirm God's proof and His testimony against you. Just as that Tree of
Oneness established His Cause and we came to recognize it, so now, in
the same way, the Cause of God is manifest, the verses of God are
revealed, and these servants have attained certainty. These servants
have not been, nor will they be, among the deniers. The deniers are
those souls who have believed in part of the Book and disbelieved in
part. No, by God! Rather, they have disbelieved in all the Books. For
today, if a soul turns away from a single letter of these verses, his
rejection of them all is established and confirmed. We have believed in
what was revealed before and what is revealed now, and all things bear
witness to this, if you are among those who are certain. How well you
reason! You are occupied with objections and disputes against the
Beloved of the Primal Point, while writing that His truth was
established through verses. You imagine that you look to the revelation,
yet you are far indeed from its Revealer.

\subsection{101}\label{section-596}

\begin{otherlanguage}{persian}

و اينكه نوشته بودى كه يقين است شما هم نميتوانيد انكار نمائيد يشهد بآياته
سرّنا و كينونتنا و لساننا و قلبنا و جوارحنا ان انت من العارفين قسم بآفتاب
عزّ توحيد كه حرفى از كلمات آن مظهر اسماء و صفات نزد جمال ابهى احبّ و اعزّ
است از كلّ آنچه در آسمانها و زمينها است ديگر تا چه رسد باين عباد كه
معدوميم نزد حرفى از او و خاضعيم نزد آنچه ظاهر شده از آن شجره حقيقت و شمس
عنايت و قمر عطوفت فرق ما بين اين عباد و شما اينست كه شما لفظاً از وراء
هزار ستر و حجاب موقنيد و از نفس و كينونت و روح و ذات او معرض و كاش
باعراض هم كفايت ميرفت لا فو‌نفسه الرّحمن الرّحيم بلكه در ليالى و ایّام در
تدبير سفك دم طيّب طاهر مشغوليد چنانچه از مكرى كه در آن ارض نموده واضح شد
و اين عباد بآيات او و كلمات او و ظهور او و نفس او و روح او بجميع موقن و
مذعن و معترفيم اذا فاعرف شأن المعرضين و شأن اللّه العلى العظيم و شأن
هؤلآء الفقرآء و المساكين و

\end{otherlanguage}

Va īnke neveshte būdī ke yaqīn ast shomā ham ne-mī-tavānīd enkār
namāyīd, yashhadu bi-āyātihi sirrunā wa-kaynūnatunā wa-lisānunā
wa-qalbunā wa-jawāriḥunā in anta min al-ʿārifīn. Qasam be-āftāb-e ʿezz-e
towḥīd ke ḥarfī az kalemāt-e ān maẓhar-e asmāʾ va ṣefāt nazd-e Jamāl-e
Abhā aḥabb va aʿazz ast az koll-e ānche dar āsemānhā va zamīnhā ast.
Dīgar tā che resad be-īn ʿebād ke maʿdūmīm nazd-e ḥarfī az ū va khāżeʿīm
nazd-e ānche ẓāher shode az ān shajare-ye ḥaqīqat va shams-e ʿenāyat va
qamar-e ʿoṭūfat. Farq-e mā-beyn-e īn ʿebād va shomā īnast ke shomā
lafẓan az varā-ye hezār setr va ḥejāb mūqenīd va az nafs va keynūnat va
rūḥ va dhāt-e ū moʿreż. Va kāsh be-eʿrāż ham kefāyat mī-raft. Lā
fa-wa-nafsihi al-Raḥmāni al-Raḥīm! Balke dar layālī va ayyām dar
tadbīr-e safk-e dam-e ṭayyeb-e ṭāher mashghūlīd chenānche az makrī ke
dar ān arż nemūde vāżeḥ shod. Va īn ʿebād be-āyāt-e ū va kalemāt-e ū va
ẓohūr-e ū va nafs-e ū va rūḥ-e ū be-jamīʿ mūqen va modhʿen va
moʿtarefīm. Idhan fa-ʿrif shaʾna al-muʿriḍīn wa-shaʾna Allāhi al-ʿAliyyi
al-ʿAẓīm wa-shaʾna hāʾulāʾi al-fuqarāʾ wa-l-masākīn wa-

As for your statement, ``It is certain that you too cannot deny'': our
inmost being, our existence, our tongue, our heart, and our limbs bear
witness to His verses, if you are among those who understand. I swear by
the Sun of the glory of divine unity that a single letter of the words
of that Manifestation of the names and attributes is dearer and more
precious to the Most Glorious Beauty than all that is in the heavens and
the earths. How much more is this true of these servants, who are as
nothing before a single letter of His and are humble before all that has
appeared from that Tree of Truth, that Sun of loving care, and that Moon
of compassion. The difference between these servants and you is this:
you profess certainty in words, from behind a thousand coverings and
veils, while turning away from His Self, His being, His spirit, and His
essence. Would that you had been content merely to turn away! No, by His
own Self, the All-Merciful, the Ever-Merciful! Rather, night and day you
are occupied with devising ways to shed His good and pure blood, as
became clear from the plot contrived in that land. These servants,
however, are certain of, assent to, and acknowledge His verses, His
words, His Manifestation, His Self, and His spirit---all of them.
Recognize, then, the condition of those who turn away, the station of
God, the Most High, the Most Great, and the condition of these poor and
lowly ones, and

\subsection{102}\label{section-597}

\begin{otherlanguage}{persian}

ديگر نوشته و خودش هم آيات را از غير آيات از ديگران بهتر تميز خواهد داد
بلكه امتياز واقعى را او بايد بدهد و البتّه چنين شخصى دروغ‌گو هم نخواهد بود
و بهواى نفس حركت نخواهد نمود انتهى

\end{otherlanguage}

Dīgar neveshte va khodash ham āyāt rā az gheyr-e āyāt az dīgarān behtar
tamīz khāhad dād balke emtīyāz-e vāqeʿī rā ū bāyad bedehad va albatte
chenīn shakhṣī dorūgh-gū ham na-khāhad būd va be-havā-ye nafs ḥarakat
na-khāhad nemūd. Intahā.

You further wrote, ``He himself will also distinguish verses from what
are not verses better than others; indeed, he is the one who must make
the true distinction. Certainly, such a person will not be a liar, nor
will he act according to selfish desire.'' End of quotation.

\subsection{103}\label{section-598}

\begin{otherlanguage}{persian}

اگر شما بظنّ و گمان نوشته‌ايد اين عباد موقناً بين سموات و ارض ناطقند بأنّه
لا اله الّا هو و انّ نقطة البيان لميزانه و ظهوره و كلّما نزل من عنده لحقّ
لاريب فيه و انّه لمميّز الحقّ عن الباطل و الشّكّ عن اليقين حال اين مميّزى كه
ذكر نموده بكلّ آيات او موقن و مذعنى يا بآنچه بنفس و هواى خود تفسير نموده
و مطابق ديده اگر بجميع ما نزل من عنده موقنى كه ميفرمايند قوله عزّ ذكره كلّ
آنچه در بيان است تحفه‌ايست از خدا از براى آنظهور عزّ صمدانى و بعد
ميفرمايند سبحانك اللّهمّ يا الهى ما اصغر ذكرى و ما ينسب الىّ اذا اريد ان
انسبه اليك فلتقبلننى و ما ينسب الىّ بفضلك و انّك انت خير الفاضلين

\end{otherlanguage}

Agar shomā be-ẓann va gomān neveshte-īd īn ʿebād mūqanan beyn-e samāvāt
va arż nāṭeqand bi-annahu lā ilāha illā huwa wa-inna Nuqṭata al-Bayān
la-mīzānuhu wa-ẓuhūruhu wa-kullamā nazala min ʿindihi la-ḥaqqun lā rayba
fīhi wa-innahu la-mumayyizu al-ḥaqq ʿan al-bāṭil wa-l-shakk ʿan
al-yaqīn. Ḥāl īn momayyezī ke dhekr nemūde be-koll-e āyāt-e ū mūqen va
modhʿenī yā be-ānche be-nafs va havā-ye khod tafsīr nemūde va moṭābeq
dīde? Agar be-jamīʿ-e mā nazala min ʿindihi mūqenī ke mī-farmāyand
qawluhu ʿazza dhikruhu koll-e ānche dar Bayān ast toḥfe-īst az Khodā az
barā-ye ān ẓohūr-e ʿezz-e ṣamadānī va baʿd mī-farmāyand subḥānaka
Allāhumma yā ilāhī mā aṣghara dhikrī wa-mā yunsabu ilayya idhā urīdu an
ansibahu ilayka fa-la-taqbalannī wa-mā yunsabu ilayya bi-faḍlika
wa-innaka anta khayru al-fāḍilīn.

If you have written on the basis of conjecture and supposition, these
servants proclaim with certainty between the heavens and the earth that
there is no God but Him, that the Point of the Bayán is indeed His
Balance and His Manifestation, that everything revealed by Him is the
truth, beyond all doubt, and that He is indeed the One who distinguishes
truth from falsehood and doubt from certainty. Now, concerning this One
who distinguishes, whom you have mentioned: are you certain of and do
you assent to all His verses, or only those you have interpreted
according to your own self and desire and found to accord with them? If
you are certain of all that has been revealed by Him, He says---exalted
be His mention---``Everything in the Bayán is a gift from God for that
Manifestation of self-subsisting glory.'' And afterward He says, ``Glory
be to You, O God, my God! How insignificant is my mention and all that
is attributed to me when I wish to ascribe it to You! Accept, then,
through Your grace, me and all that is attributed to me. You are indeed
the best of those who bestow grace.''

\subsection{104}\label{section-599}

\begin{otherlanguage}{persian}

اى غافلين از ملأ بيان مع اين خضوع نقطه اولى عند ظهوره الآخر مرّة اخرى
متنبّه نشده باين اذكار مشغول گشته‌ايد فو‌اللّه اگر فی‌‌الجمله تفكّر نمائيد از
براى انفس خود نوحه كنيد چه كه در اين مقام آن وجود مبارك خود را بكمال
خشوع و فناى محض ذكر فرموده‌اند ديگر كجا مقام مرايا و ما ينسب اليهنّ مذكور
ميشود آيا سطرى از بيان مشاهده شده كه منتهى باين ذكر ابدع بديع نشده لا
فو‌نفسه الحقّ و لكن انتم لاتعرفون و اللّه در مقامى افتاده‌ايد كه گويا بهزار
سلاسل حركت ننمائيد آخر يك نفس از مفازه اشارات بيرون خرام ودر فضاى خوش
تجريد وارد شو تا حقّ منيع را بنفس مقدّسش بشناسى و از دونش منقطع شوى

\end{otherlanguage}

Ey ghāfelīn az malaʾ-e Bayān, maʿ-e īn khożūʿ-e Noqṭe-ye Ūlā ʿinda
ẓuhūrihi al-ākhar marratan ukhrā motanabbeh na-shode be-īn adhkār
mashghūl gashte-īd. Fa-wa-Allāh agar fī-l-jomle tafakkor namāyīd az
barā-ye anfos-e khod nowḥe konīd, che ke dar īn maqām ān vojūd-e mobārak
khod rā be-kamāl-e khoshūʿ va fanā-ye maḥż dhekr farmūde-and. Dīgar kojā
maqām-e marāyā va mā yunsabu ilayhinna madhkūr mī-shavad? Āyā saṭrī az
Bayān moshāhede shode ke montahī be-īn dhekr-e abdaʿ-e badīʿ na-shode?
Lā fa-wa-nafsihi al-Ḥaqq wa-lākin antum lā taʿrifūn. Va-Allāh dar maqāmī
oftāde-īd ke gūyā be-hezār salāsel ḥarakat na-namāyīd. Ākhar yek nafas
az mafāze-ye eshārāt bīrūn kherām va dar fażā-ye khosh-e tajrīd vāred
show tā Ḥaqq-e manīʿ rā be-nafs-e moqaddasash beshenāsī va az dūnash
monqaṭeʿ shavī.

O heedless ones among the people of the Bayán! Despite this humility of
the Primal Point before His other Manifestation, you have once again
failed to take heed and have occupied yourselves with these assertions.
By God! If you reflect even a little, you will lament for yourselves,
for in this passage that blessed Being has spoken of Himself with utter
humility and complete self-effacement. What place, then, is there to
mention the station of the mirrors and what is attributed to them? Has a
single line of the Bayán been seen that does not culminate in this most
wondrous, incomparable remembrance? No, by His own Self, the True One!
Yet you do not understand. By God! You have fallen into such a state
that it is as though a thousand chains prevent you from moving. At last,
step out for a moment from the wilderness of allusions and enter the
delightful expanse of detachment, so that you may recognize the
unassailable True One through His own sanctified Self and detach
yourself from all besides Him.

\subsection{105}\label{section-600}

\begin{otherlanguage}{persian}

و اينكه آن طلعت احديّه را مميّز ذكر نموده بچه دليل اگر بآيات بوده پس چرا
از اين مميّز معرضى كه بآياتى ظاهر شده كه كلّ ذى آيه خاضع است لدى حضرته الّا
الّذينهم تحرّكهم ارياح الرّياسة كيف تشآء مرّة يضربهم على قلل الجبال و مرّة
على تيه الضّلال و مرّة يطرحهم الى غمرات الفنآء و مرّة يدعهم تحت ذياب النّفس
و الهوى قل موتوا بغيظكم تاللّه قد ظهر من ثبت بقوله حجّة الأوّلين و الآخرین
و به تثبت الى آخر الّذى لا آخر له و لاينكر ذلك الّا كلّ كفّار اثيم

\end{otherlanguage}

Va īnke ān ṭalʿat-e aḥadīye rā momayyez dhekr nemūde be-che dalīl? Agar
be-āyāt būde pas cherā az īn momayyez moʿreżī ke be-āyātī ẓāher shode ke
koll-e dhī āye khāżeʿ ast ladā ḥaḍratihi illā alladhīna hum tuḥarrikuhum
aryāḥu al-riʾāsa kayfa tashāʾ. Marratan yaḍribuhum ʿalā qulali al-jibāl
wa-marratan ʿalā tīhi al-ḍalāl wa-marratan yaṭraḥuhum ilā ghamarāti
al-fanāʾ wa-marratan yadaʿuhum taḥta dhiʾābi al-nafs wa-l-hawā. Qul mūtū
bi-ghayẓikum. Ta-Allāhi qad ẓahara man thabata bi-qawlihi ḥujjatu
al-awwalīn wa-l-ākhirīn wa-bihi tathbutu ilā ākhiri alladhī lā ākhira
lahu wa-lā yunkiru dhālika illā kullu kaffārin athīm.

On what evidence have you described that Countenance of Oneness as the
One who distinguishes? If it was on the basis of verses, why, then, do
you turn away from this One who distinguishes, who has appeared with
verses before which everyone endowed with verses bows humbly in His
presence, except those whom the winds of leadership move as they please?
At one time they dash them against the mountaintops, at another against
the wilderness of error; at one time they cast them into the depths of
extinction, and at another leave them beneath the wolves of self and
desire. Say: Die in your rage! By God! The One has appeared through
whose word the proof of the first and the last was established, and
through whom it will remain established until the end that has no end.
None denies this but every sinful, inveterate disbeliever.

\subsection{106}\label{section-601}

\begin{otherlanguage}{persian}

و اينكه نوشته دروغ‌گو هم نخواهد بود اگر چه بظاهر تصديق صدق حضرت نموده و
لكن اين تصديق شبحى بوده و اخذ شد چه كه اگر موقن بصدق نقطه اولى بودى منبع
صدق را كه باراده از او حقيقت صدق مخلوق و ممدوح تكذيب نمينمودى و تصديق
صدق حضرت در اينمقام نظر بآن بوده كه باين تصديق ظاهره كينونت و حقيقتش را
تكذيب نمائى چنانچه علماى قبل حين ظهور نقطه بيان جلّت عظمته همين قسم متكلّم
كه يقين ائمّه هدى صادقند احاديث متواتره از ايشان موجود كه هر نفسى قائل
شود كه بايد قائم متولّد شود كاذب و منافق و مفترى است و قتلش واجب و هم‌چنين
استدلال مينمودند كه خاتم النّبيّين قول خدا است و در قرآن بنصّ صريح نازل و
يقين است كه خدا صادق است و نبی‌ او صادق و بعد اعتراضاً على اللّه و مظهر
نفسه باين كلمات ناطق كه اين شخص مدّعى قولش باطل و از درجه اعتبار ساقط چه
كه مخالف كتاب اللّه و سننه و نبی‌ صادق تكلّم نموده چنانچه در كتاب او است كه
نبی‌ از اوّل لا اوّل آمده و الى آخر لا آخر خواهد آمد

\end{otherlanguage}

Va īnke neveshte dorūgh-gū ham na-khāhad būd, agar-che be-ẓāher taṣdīq-e
ṣedq-e Ḥażrat nemūde va lāken īn taṣdīq shabaḥī būde va akhdh shod, che
ke agar mūqen be-ṣedq-e Noqṭe-ye Ūlā būdī manbaʿ-e ṣedq rā ke be-erāde
az ū ḥaqīqat-e ṣedq makhlūq va mamdūḥ takdhīb ne-mī-nemūdī. Va taṣdīq-e
ṣedq-e Ḥażrat dar īn maqām naẓar be-ān būde ke be-īn taṣdīq-e ẓāhere
keynūnat va ḥaqīqatash rā takdhīb namāyī. Chenānche ʿolamā-ye qabl ḥīn-e
ẓohūr-e Noqṭe-ye Bayān jallat ʿaẓamatuhu hamīn qesm motakallem ke yaqīn
aʾemme-ye hodā ṣādeqand. Aḥādīth-e motavātere az īshān mowjūd ke har
nafsī qāʾel shavad ke bāyad Qāʾem motavalled shavad kādheb va monāfeq va
moftarī ast va qatl-ash vājeb. Va ham-chenīn estedlāl mī-nemūdand ke
Khātam al-Nabīyīn qowl-e Khodā ast va dar Qorʾān be-naṣṣ-e ṣarīḥ nāzel
va yaqīn ast ke Khodā ṣādeq ast va Nabī-ye ū ṣādeq. Va baʿd iʿtirāḍan
ʿalā Allāh wa-maẓhari nafsihi be-īn kalemāt nāṭeq ke īn shakhṣ-e moddaʿī
qowlash bāṭel va az daraje-ye eʿtebār sāqeṭ che ke mokhālef-e Kitāb
Allāh va sonan-ash va Nabī-ye ṣādeq takallom nemūde, chenānche dar
ketāb-e ū ast ke Nabī az avval-e lā avval āmade va elā ākher-e lā ākher
khāhad āmad.

As for your statement, ``He will not be a liar'': although you have
outwardly affirmed His truthfulness, this affirmation was a mere
semblance and has been taken up. For had you been certain of the
truthfulness of the Primal Point, you would not have denied the Source
of truthfulness, by an act of whose will the very reality of
truthfulness is created and praised. Your affirmation of His
truthfulness here was intended to enable you, through this outward
affirmation, to deny His being and His reality. In the same way, the
divines of the past, at the appearance of the Point of the
Bayán---glorified be His majesty---said, ``The Imams of guidance are
certainly truthful. There are widely transmitted traditions from them
stating that anyone who says the Qá'im must be born is a liar, a
hypocrite, and a fabricator, and must be put to death.'' They also
argued that ``Seal of the Prophets'' is God's word, explicitly revealed
in the Qur'an, and that God is certainly truthful and His Prophet
truthful. Then, in objection to God and the Manifestation of His own
Self, they uttered these words: ``The words of this claimant are false
and devoid of credibility, for he has spoken contrary to the Book of
God, His established ways, and His truthful Prophet. Thus, in his book
it says that Prophets have come from the beginning that has no beginning
and will come until the end that has no end.''

\subsection{107}\label{section-602}

\begin{otherlanguage}{persian}

بعينه همين قسم كه شما اعتراض نموده‌ايد اولئك اعترضوا بل انتم اشدّ اعتراضاً
لو انتم تعرفون حال ملاحظه كن كه چه قدر متوهّمند آن معرضين كه اين كلماترا
بشما القا نموده‌اند گويا بحرفى از معانى كلمات الهيّه فايز نشده‌اند و بغرفه
از قلزم عزّ صمدانيه مرزوق نگشته‌اند فويل لهم و بما امرهم نفسهم و هويهم
آنچه از بيان ذكر نمائيد و تسليم هم نمائيم كه جميع را فهميده‌ايد و حال
آنكه اين محال بوده چه كه احدى بمعانى آن مطّلع نه الّا اللّه و هم‌چنين بجميع
اعمال آن عامل در حين ظهور نفع نمى‌بخشد تو را الّا بتصديق ظهور چه كه بنص
صريح ميفرمايد بما نزل فى البيان از آن طلعت رحمن محتجب نشويد تا چه رسد
باين چهار كلمه كه شما ذكر نموده‌ايد

\end{otherlanguage}

Be-ʿeyne hamīn qesm ke shomā eʿterāż nemūde-īd ulāʾika iʿtaraḍū bal
antum ashaddu iʿtirāḍan law antum taʿrifūn. Ḥāl molāḥeẓe kon ke che qadr
motavahhemand ān moʿreżīn ke īn kalemāt rā be-shomā elqā nemūde-and.
Gūyā be-ḥarfī az maʿānī-ye kalemāt-e elāhīye fāʾez na-shode-and va
be-ghorfe az qolzom-e ʿezz-e ṣamadānīye marzūq na-gashte-and. Fa-waylun
lahum wa-bi-mā amarahum nafsuhum wa-hawāhum. Ānche az Bayān dhekr
namāyīd va taslīm ham namāyīm ke jamīʿ rā fahmīde-īd va ḥāl ānke īn
moḥāl būde che ke aḥadī be-maʿānī-ye ān moṭṭaleʿ na illā Allāh va
ham-chenīn be-jamīʿ-e aʿmāl-e ān ʿāmel, dar ḥīn-e ẓohūr nafʿ
ne-mī-bakhshad to rā illā be-taṣdīq-e ẓohūr, che ke be-naṣṣ-e ṣarīḥ
mī-farmāyad bi-mā nazala fī al-Bayān az ān ṭalʿat-e Raḥmān moḥtajeb
na-shavīd tā che resad be-īn chahār kaleme ke shomā dhekr nemūde-īd.

They raised objections in precisely the same way you have done; indeed,
your objections are even more vehement, if you only understood. Now
consider how deluded are those who have turned away and instilled these
words into you. It is as though they have not attained a single letter
of the meanings of the divine words or received a single draught from
the ocean of self-subsisting glory. Woe to them and to what their selves
and desires have commanded them! Whatever you may cite from the
Bayán---even if we concede that you have understood it all, though this
is impossible, for no one knows its meanings except God, and that you
also carry out all its prescribed deeds---it will not benefit you at the
time of the Manifestation unless you acknowledge that Manifestation. For
He explicitly says, ``Do not let what has been revealed in the Bayán
veil you from that Countenance of the All-Merciful.'' How much less,
then, can these few words you have mentioned avail you!

\subsection{108}\label{section-603}

\begin{otherlanguage}{persian}

و يا بعد از ثبوت اين كلمه طيّبه مباركه كلّ آنچه ذكر نموده‌ايد و يا از بعد
بآن مستدلّ شويد جميع راجع بانفس محدوده ذاكر و مستدلّ بوده و حقّ مقدّس از كلّ
و اين بسى واضح است كه اگر نقطه اولى را صادق ميدانستى ابداً تكذيب اين ظهور
ابدع امنع نمينمودى حال بلسان صادقش ميدانى و امّا بفعل نعوذ باللّه كاذب
فاستغفر عما ظننت و قلت و كتبت و

\end{otherlanguage}

Va yā baʿd az thobūt-e īn kaleme-ye ṭayyebe-ye mobārake koll-e ānche
dhekr nemūde-īd va yā az baʿd be-ān mostadell shavīd jamīʿ rājeʿ
be-anfos-e maḥdūde-ye dhāker va mostadell būde va Ḥaqq moqaddas az koll.
Va īn basī vāżeḥ ast ke agar Noqṭe-ye Ūlā rā ṣādeq mī-dānestī abadan
takdhīb-e īn ẓohūr-e abdaʿ-e amnaʿ ne-mī-nemūdī. Ḥāl be-lesān ṣādeqash
mī-dānī va ammā be-feʿl naʿūdhu bi-Allāh kādheb. Fa-staghfir ʿammā
ẓananta wa-qulta wa-katabta wa-

Or, after the truth of this good and blessed word has been established,
everything you have mentioned, or may afterward adduce as evidence,
refers entirely to the limited selves of those who mention it and argue
from it, while the True One is sanctified above all. It is abundantly
clear that, had you regarded the Primal Point as truthful, you would
never have denied this most wondrous, most unassailable Manifestation.
Now, with your tongue you deem Him truthful, but through your
actions---God forbid---a liar. Seek forgiveness, then, for what you have
thought and said and written, and

\subsection{109}\label{section-604}

\begin{otherlanguage}{persian}

قل سبحانك اللّهمّ يا الهى انا الّذى فرّطت فى جنبك و سلكت مسالك الظّنون و
الأوهام بما كنت محتجباً عن عرفان مظهر نفسك و استدللت بما احترق كبد علىّ ثمّ
قلب محمّد رسول اللّه كانّى قتلت كلّ الوجود من الغيب و الشّهود اذا فانظرنى يا
الهى بلحظات رأفتك و لا تدعنى بنفسى و هوائى ثمّ طهّرنى عن اشارات الّتى بها
صرت محروماً عن رياض وصلك و لقائك و انّك انت المقتدر على ما تشآء و انّك انت
المهيمن القيّوم

\end{otherlanguage}

Qul subḥānaka Allāhumma yā ilāhī, anā alladhī farraṭtu fī janbika
wa-salaktu masālika al-ẓunūni wa-al-awhām, bimā kuntu muḥtajiban ʿan
ʿirfāni maẓhari nafsika, wa-istad laltu bimā iḥtaraqa kabidu ʿAliyyin
thumma qalbu Muḥammadin Rasūli Allāh. Ka-annī qataltu kulla al-wujūdi
mina al-ghaybi wa-al-shuhūd. Idhan fa-unẓurnī yā ilāhī bi-laḥaẓāti
raʾfatika wa-lā tadaʿnī bi-nafsī wa-hawāʾī, thumma ṭahhirnī ʿan ishārāti
allatī bihā ṣirtu maḥrūman ʿan riyāḍi waṣlika wa-liqāʾika, wa-innaka
anta al-muqtadiru ʿalā mā tashāʾu wa-innaka anta al-muhayminu al-qayyūm.

Say: ``Glory be to You, O God, my God! I am the one who neglected his
duty toward You and followed the paths of conjecture and delusion, for I
was veiled from recognizing the Manifestation of Your own Self. I
adduced as evidence that which burned the liver of ʿAlí and then the
heart of Muḥammad, the Messenger of God. It is as though I had slain all
existence, unseen and seen. Look upon me, then, O my God, with the
glances of Your compassion, and do not abandon me to my own self and
desires. Then purify me of the allusions through which I was deprived of
the gardens of union with You and of meeting You. You indeed have power
over whatever You will, and You are indeed the All-Protecting, the
Self-Subsisting.''

\subsection{110}\label{section-605}

\begin{otherlanguage}{persian}

و ديگر نوشته بوديد احدى از خارج و داخل و شهرت و واقع تواند شكّى نمود كه
حضرت باب تصديق ميرزا يحيى ننموده باين معنى كه باطراف و اصحاب در حقّیّت او
بعبارات مختلفه ننوشته كه بعد از من او است و اطاعت او بر همه لازم است وصىّ
و خليفه هر چه هست بعد از من او است صريحاً و بعبارات صريحه بجهت اهل تقليد
و عوام بيان فرموده‌ اين را و امّا بجهت اولى الأبصار و الأفئده عبارات ديگر
بيان فرموده كه انّ الأمر ينتهى الى اسم الوحيد لانّ ظهوره بنفسه حجّة و لا
يحتاج الى نصىّ انتهى

\end{otherlanguage}

Va dīgar nivishtih būdīd aḥadī az khārij va dākhil va shuhrat va vāqiʿ
tavānad shakkī nimūd kih Ḥaḍrat-i Bāb taṣdīq-i Mīrzā Yaḥyā nanimūdih,
bi-īn maʿnā kih bi-aṭrāf va aṣḥāb dar ḥaqqiyyat-i ū bi-ʿibārāt-i
mukhtalifih nanivishtih kih baʿd az man ū ast va iṭāʿat-i ū bar hamih
lāzim ast, vaṣī va khalīfih har chih hast baʿd az man ū ast. Ṣarīḥan va
bi-ʿibārāt-i ṣarīḥih bi-jihat-i ahl-i taqlīd va ʿavāmm bayān farmūdih īn
rā, va ammā bi-jihat-i ulī al-abṣār va al-afʾidih ʿibārāt-i dīgar bayān
farmūdih kih inna al-amra yantahī ilā ismi al-waḥīd li-anna ẓuhūrahu
bi-nafsihi ḥujjatun wa-lā yaḥtāju ilā naṣṣī. Intahā.

You also wrote: ``Can anyone, whether an outsider or an insider, whether
on the basis of common report or actual fact, doubt that the Báb
confirmed Mírzá Yaḥyá---that is, that He wrote in various terms to the
regions and the believers affirming his rightful standing: `After Me, it
is he, and obedience to him is incumbent on everyone. Whatever the
position of legatee and successor may be, after Me it is he'? He stated
this explicitly and in plain terms for the followers of tradition and
the common people. But for those endowed with insight and understanding
He used other words: `The Cause will culminate in the name of the One,
for His appearance is in itself a proof and requires no written
designation.'\,'' End of quotation.

\subsection{111}\label{section-606}

\begin{otherlanguage}{persian}

بسيار دلم بعد از ملاحظه اين كلمات سوخت بر شما چه كه بعينه مثل طلبه‌هاى
مدرسه كه حال در مدارس اثبات مطلب خود مينمايند بهمان دلايل وهميّه و كلمات
متشابه مستدل شده‌ايد بقول خود شما كه اهل فرقان را موهوم ميدانيد وهم شما
صد‌هزار درجه بيشتر از آن طايفه مشاهده ميشود الى متى تكون راقداً على البساط
فارفع رأسك عن النّوم ثمّ افتح بصرك لتشهد الشّمس مشرقة فى قطب الزّوال بگذاريد
اين كلمات را و بمنزل آن ناظر شويد وهم را بيندازيد و بسلطان يقين تشبّث
نمائيد حيف سمع و بصر كه بمشاهده و استماع اين كلمات مشغول شود

\end{otherlanguage}

Bisyār dilam baʿd az mulāḥaẓih-yi īn kalimāt sūkht bar shumā, chih kih
bi-ʿaynih mithl-i ṭalabih-hā-yi madrisih kih ḥāl dar madāris ithbāt-i
maṭlab-i khud mīnamāyand, bi-hamān dalāʾil-i vahmiyyih va kalimāt-i
mutashābih mustadill shudih-īd. Bi-qawl-i khud-i shumā kih ahl-i Furqān
rā mawhūm mīdānīd, vahm-i shumā ṣad-hizār darajih bīshtar az ān ṭāʾifih
mushāhadih mīshavad. Ilā matā takūnu rāqidan ʿalā al-bisāṭ? Fa-irfaʿ
raʾsaka ʿan al-nawm thumma iftaḥ baṣaraka li-tashhada al-shamsa
mushriqatan fī quṭbi al-zawāl. Bigudhārīd īn kalimāt rā va bi-munzil-i
ān nāẓir shavīd. Vahm rā biyandāzīd va bi-sulṭān-i yaqīn tashabbuth
namāyīd. Ḥayf-i samʿ va baṣar kih bi-mushāhadih va istimāʿ-i īn kalimāt
mashghūl shavad.

My heart burned with sorrow for you after reading these words, for you
have relied on precisely the same fanciful arguments and ambiguous
statements by which the seminary students now establish their claims in
their schools. You yourself regard the people of the Qurʾán as deluded,
yet your own delusion is seen to exceed theirs a hundred thousandfold.
How long will you lie asleep upon the carpet? Raise your head from
sleep, then open your eyes to behold the sun shining at its noonday
height. Set these words aside and look toward the One who revealed them.
Cast away delusion and cling to the Sovereign of certitude. What a waste
of hearing and sight that they should be occupied with looking at and
listening to these words!

\subsection{112}\label{section-607}

\begin{otherlanguage}{persian}

اوّلا آنكه آنچه نوشته‌ايد بر عكس بوده و از آن گذشته ميفرمايد ا تؤمنون ببعض
الكتاب و تكفرون ببعض چگونه اين كلماتى كه اصلش معلوم نيست كه بكجا راجع
است مسلّم دانسته‌ايد و جميع بيان كه ميفرمايد در حين ظهور بهيچ شیء غير او
متمسّك نشويد اين بيان رحمن غير مسموع و حال كلمات مجعوله بايد استماع نمود

\end{otherlanguage}

Avvalan ānkih ānchih nivishtih-īd bar ʿaks būdih va az ān gudhashtih
mīfarmāyad a-tuʾminūna bi-baʿḍi al-kitābi wa-takfurūna bi-baʿḍ? Chigūnih
īn kalimātī kih aṣlash maʿlūm nīst kih bi-kujā rājiʿ ast musallam
dānistih-īd va jamīʿ-i Bayān kih mīfarmāyad dar ḥīn-i ẓuhūr bi-hīch
shayʾ ghayr-i ū mutamassik nashavīd, īn bayān-i Raḥmān ghayr-i masmūʿ va
ḥāl kalimāt-i majʿūlih bāyad istimāʿ nimūd.

First, what you have written is the reverse of the truth. Moreover, He
says: ``Do you believe in part of the Book and disbelieve in part?'' How
have you accepted as indisputable these words, whose reference is not
known at all, while the whole Bayán declares that at the time of the
Manifestation you must cling to nothing other than Him? This utterance
of the All-Merciful goes unheard, while fabricated words must now be
listened to!

\subsection{113}\label{section-608}

\begin{otherlanguage}{persian}

ابداً در بيان ذكر وصايت و خلافت نبوده و نيست بقول شما اين كلمات اهل فرقان
است آنچه در بيان ذكر شده ذكر حروفات و مرايا بوده و در جميع مواضع
ميفرمايند مرايا منحصر نبوده و نخواهد بود هر نفسى باشراق شمس الهيّه مقابل
شود او از مرايا محسوب بوده و اين شأن باقى مادامىكه از تقابل انحراف نجويد
و بعد از انحراف يرجع النّور الى مقرّه و ينتهى المخلوق الى شكله و مثله

\end{otherlanguage}

Abadan dar Bayān dhikr-i viṣāyat va khilāfat nabūdih va nīst. Bi-qawl-i
shumā īn kalimāt-i ahl-i Furqān ast. Ānchih dar Bayān dhikr shudih
dhikr-i ḥurūfāt va marāyā būdih va dar jamīʿ-i mavāḍiʿ mīfarmāyand
marāyā munḥaṣir nabūdih va nakhvāhad būd. Har nafsī bi-ishrāq-i shams-i
ilāhiyyih muqābil shavad ū az marāyā maḥsūb būdih va īn shaʾn bāqī
mādāmīkih az taqābul inḥirāf najūyad va baʿd az inḥirāf yarjiʿu al-nūru
ilā maqarrihi wa-yantahī al-makhlūqu ilā shaklihi wa-mithlihi.

There has never been, nor is there now, any mention in the Bayán of
legateeship or succession. In your own words, these are the terms of the
people of the Qurʾán. What is mentioned in the Bayán is Letters and
Mirrors, and throughout it He states that the Mirrors have not been, and
will not be, restricted to a fixed number. Any soul who faces the
radiance of the divine Sun is counted among the Mirrors, and this
station endures so long as that soul does not turn away from facing it.
After that turning away, the light returns to its seat, and the creature
reverts to its own kind and likeness.

\subsection{114}\label{section-609}

\begin{otherlanguage}{persian}

و در كلّ بيان مخصوصاً منصوصاً ميفرمايد در يوم ظهور ابدا بغير او ناظر نباشيد
چه كه آن جمال قدم ببصر غير خود شناخته نميشود چنانچه نشده و الآن خود نقطه
اولى ميفرمايد كه آنچه شما مينويسيد خطا است و استدلالات شما راجع بانفس
متوهّمين قد ظهر الحقّ بسلطانه و دليله آياته و ظهوره نفسه اذاً موتوا بغيظكم
يا ملأ المتوهّمين

\end{otherlanguage}

Va dar kull-i Bayān makhṣūṣan manṣūṣan mīfarmāyad dar yawm-i ẓuhūr
abadan bi-ghayr-i ū nāẓir nabāshīd, chih kih ān Jamāl-i Qidam bi-baṣar-i
ghayr-i khud shinākhtih namīshavad, chinānchih nashudih, va al-ān khud-i
Nuqṭih-yi Ūlā mīfarmāyad kih ānchih shumā mīnivīsīd khaṭā ast va
istidlālāt-i shumā rājiʿ bi-anfus-i mutavahhimīn. Qad ẓahara al-Ḥaqqu
bi-sulṭānihi wa-dalīlihi āyātihi wa-ẓuhūrihi nafsihi. Idhan mūtū
bi-ghayẓikum yā malaʾa al-mutawahhimīn.

Throughout the Bayán He specifically and explicitly declares: ``On the
day of the Manifestation, never look toward anyone other than Him,'' for
that Ancient Beauty cannot be recognized through the eyes of anyone
other than Himself, nor has He been. Now the Primal Point Himself
declares that what you write is erroneous and that your arguments refer
back to the selves of the deluded. The True One has appeared with His
sovereignty, His proof---His verses---and His manifestation---His own
Self. Die, then, in your rage, you company of the deluded!

\subsection{115}\label{section-610}

\begin{otherlanguage}{persian}

ذكر خارج و داخل و شهرت نموده جميع اين كلمات را بوهم صرف و ظنّ بحت نوشته
امّا خارج كه ابدا اطّلاع نداشته و امّا داخل كه بعضى برفيق اعلى ارتقاء
جسته‌اند و بعضى هم كه موجودند شما آن نفوس را تكذيب نموده‌ايد و بهواهاى خود
جعليّاتى ترتيب داده كه شايد ضعفاى بيچاره را از حقّ منيع محروم نمائيد
چنانچه از قبل نموده‌اند و كاش ادراك ميشد كه شهرت را كه داده و حكمت آن چه
بوده لا يعرفها الّا اولى الأباب

\end{otherlanguage}

Dhikr-i khārij va dākhil va shuhrat nimūdih, jamīʿ-i īn kalimāt rā
bi-vahm-i ṣirf va ẓann-i baḥt nivishtih. Ammā khārij kih abadan iṭṭilāʿ
nadāshtih va ammā dākhil kih baʿḍī bi-Rafīq-i Aʿlā irtiqā justih-and va
baʿḍī ham kih mawjūdand shumā ān nufūs rā takdhīb nimūdih-īd va
bi-havāhā-yi khud jaʿliyyātī tartīb dādih kih shāyad ḍuʿafā-yi bīchārih
rā az Ḥaqq-i Manīʿ maḥrūm namāyīd, chinānchih az qabl nimūdih-and. Va
kāsh idrāk mīshud kih shuhrat rā kih dādih va ḥikmat-i ān chih būdih. Lā
yaʿrifuhā illā ulū {[}corrupted text{]}.

You have mentioned outsiders, insiders, and common report, writing all
these words out of sheer delusion and pure conjecture. As for outsiders,
they have had no knowledge at all. As for insiders, some have ascended
to the Most Exalted Companion, while those who are still present you
have branded as liars. Following your own desires, you have put together
fabrications so that you might deprive the helpless and weak of the
unassailable True One, just as others have done before. If only it were
understood who gave rise to that common report and what the wisdom
behind it was! None knows it except those endowed with {[}corrupted
original text{]}.

\subsection{116}\label{section-611}

\begin{otherlanguage}{persian}

و اين عبد متحيّر است كه تو بچه دليل ادعاى اطّلاع نموده چه كه هيچ وقت داخل
اين امر نبوده بلى عنقريب است كه روايات غير صحيحه سقيمه از ناحيه كذب انفس
مشركه ظاهر شود و البتّه معنعن هم خواهد شد آن روايات و منتهى ميشود بسيّد
محمّد اصفهانى و يا بحسين مشهور بسوخته كه اين ایّام وارد شده و در ظلّ شيطان
واقع و يا بميرزا حسين قمى الّذى تقل على وجه المحبوب لحفظ نفسه الخبيثة او
الاعرج الّذى تبرّء عن اللّه حين الّذى كان على المنبر و خود شما هم ميدانيد كه
هيچيك از اين نفوس خبر نداشته‌اند و لكن هوى شما را بر آن داشته كه تصديق من
لا صدق له را نموده خف عن اللّه و لاتقل ما لا تعلم و لاتكن من المفترين

\end{otherlanguage}

Va īn ʿabd mutaḥayyir ast kih tu bi-chih dalīl iddiʿā-yi iṭṭilāʿ
nimūdih, chih kih hīch vaqt dākhil-i īn amr nabūdih. Balī ʿan-qarīb ast
kih rivāyāt-i ghayr-i ṣaḥīḥih-yi saqīmih az nāḥiyih-yi kidhib-i anfus-i
mushrikih ẓāhir shavad va al-battih muʿanʿan ham khvāhad shud ān rivāyāt
va muntahī mīshavad bi-Sayyid Muḥammad-i Iṣfahānī va yā bi-Ḥusayn
mashhūr bi-Sūkhtih kih īn ayyām vārid shudih va dar ẓill-i shayṭān vāqiʿ
va yā bi-Mīrzā Ḥusayn-i Qummī alladhī taqala ʿalā wajhi al-Maḥbūb
li-ḥifẓi nafsihi al-khabīthati aw al-aʿraji alladhī tabarraʾa ʿan Allāhi
ḥīna alladhī kāna ʿalā al-minbar. Va khud-i shumā ham mīdānīd kih
hīch-yik az īn nufūs khabar nadāshtih-and va lākin havā shumā rā bar ān
dāshtih kih taṣdīq-i man lā ṣidqa lahu rā nimūdih. Khaf ʿan Allāh wa-lā
taqul mā lā taʿlam wa-lā takun mina al-muftarīn.

This servant is bewildered as to the grounds on which you claim to have
knowledge, for you have never been within this Cause. Indeed, before
long unsound and false reports will emerge from the realm of the lies of
idolatrous souls. Those reports will certainly be furnished with chains
of transmission, ending with Siyyid Muḥammad-i-Iṣfahání, or Ḥusayn,
known as Súkhtih, who has arrived in these days and is under the shadow
of Satan, or Mírzá Ḥusayn-i-Qummí, who spat in the face of the Beloved
to preserve his own vile self, or the lame man who disavowed God while
on the pulpit. You yourself know that none of these souls had any
knowledge, yet your desire has led you to affirm someone in whom there
is no truthfulness. Fear God, do not say what you do not know, and do
not be among those who fabricate falsehoods.

\subsection{117}\label{section-612}

\begin{otherlanguage}{persian}

اگر كلّ ما على الأرض اليوم مراياى لطيفه شوند و شهادت دهند بر امرى كفايت
نخواهد نمود چه كه شهادت كلّ منوط است بنفس ظهور چنانچه نقطه بيان روح ما
سواه فداه در ظهور آخر ميفرمايد قوله عزّ ذكره اگر كلّ ما على الأرض شهادت
دهند بر امرى و او شهادت دهد بدون آن كه آنها شهادت داده شهادت او مثل شمس
است و شهادت آنها مثل شبح شمس است كه در تقابل واقع نشده و الّا مطابق با
شهادت او ميگرديد قسم بذات اقدس الهى كه يك سطر از كلام او بهتر است از
كلام كلّ ما على الأرض بلكه استغفار ميكنم از اين ذكر افعل التّفضيل كجا
ميتواند آثار شموس در مرايا مثل آثار شمس در سماء گردد ذلك فى حدّ اللّا شیء
و هذا فى حدّ مشىّئ الشّیئ باللّه عزّ و جلّ انتهى

\end{otherlanguage}

Agar kull-i mā ʿalā al-arḍ al-yawm marāyā-yi laṭīfih shavand va shahādat
dihand bar amrī kifāyat nakhvāhad nimūd, chih kih shahādat-i kull manūṭ
ast bi-nafs-i ẓuhūr, chinānchih Nuqṭih-yi Bayān, rūḥ-i mā sivāh fidāh,
dar ẓuhūr-i ākhar mīfarmāyad, qawluhu ʿazza dhikruhu: Agar kull-i mā
ʿalā al-arḍ shahādat dihand bar amrī va ū shahādat dihad bi-dūn-i ān kih
ānhā shahādat dādih, shahādat-i ū mithl-i shams ast va shahādat-i ānhā
mithl-i shabaḥ-i shams ast kih dar taqābul vāqiʿ nashudih va illā
muṭābiq bā shahādat-i ū mīgardīd. Qasam bi-dhāt-i aqdas-i ilāhī kih yik
saṭr az kalām-i ū bihtar ast az kalām-i kull-i mā ʿalā al-arḍ, balkih
istighfār mīkunam az īn dhikr-i afʿal al-tafḍīl. Kujā mītavānad āthār-i
shumūs dar marāyā mithl-i āthār-i shams dar samāʾ gardad? Dhālika fī
ḥaddi al-lā shayʾ wa-hādhā fī ḥaddi mushayyiʾi al-shayʾ bi-Allāhi ʿazza
wa-jall. Intahā.

Even if all who are on earth today were to become refined mirrors and
testify to a matter, this would not suffice, for the testimony of all
depends on the Manifestation Himself. Thus the Point of the Bayán---may
the souls of all besides Him be sacrificed for Him---says concerning the
other Manifestation, in His words, exalted be His mention: ``If all who
are on earth testify to a matter and He testifies to something other
than what they have testified to, His testimony is like the sun, and
theirs is like an image of the sun that has not been placed facing it;
otherwise, it would accord with His testimony. I swear by the most holy
Essence of God that one line of His words is better than the words of
all who are on earth. Indeed, I ask forgiveness for this use of the
comparative. How could the effects of suns in mirrors ever become like
the effects of the sun in heaven? The former is at the level of
nothingness, and the latter at the level of the One who brings things
into being, by God, mighty and glorious is He.'' End of quotation.

\subsection{118}\label{section-613}

\begin{otherlanguage}{persian}

حال آنچه در نزد عباد مذكور و مشهور على زعمك ابداً عنداللّه معتبر نبوده و
نخواهد بود و كلّ معلّق و منوط بتصديق اين ظهور ابدع امنع اقدس لو انتم
تعلمون و بر فرض تصديق آنچه ذكر نموده‌ايد در حين ظهور آخر ميفرمايند عزّ
ذكره كه طوريون ميّت و لاشیء ميشوند حروفات و ادلّاء نور محتجب ميمانند و
شجره اثبات از نفى محسوب ميشود و هم‌چنين ميفرمايند بسا نفسى در ظهورى از
اعلى علوّ اهل جنّت است و حين ظهور بعد از ادنى دنوّ اهل نار ميگردد

\end{otherlanguage}

Ḥāl ānchih dar nazd-i ʿibād madhkūr va mashhūr ʿalā zaʿmika abadan ʿinda
Allāh muʿtabar nabūdih va nakhvāhad būd va kull muʿallaq va manūṭ
bi-taṣdīq-i īn ẓuhūr-i abdaʿ-i amnaʿ-i aqdas, law antum taʿlamūn. Va bar
farḍ-i taṣdīq-i ānchih dhikr nimūdih-īd, dar ḥīn-i ẓuhūr-i ākhar
mīfarmāyand, ʿazza dhikruhu, kih ṭūriyyūn mayyit va lā shayʾ mīshavand,
ḥurūfāt va adillāʾ-i nūr muḥtajib mīmānand va shajari h-yi ithbāt az
nafy maḥsūb mīshavad. Va ham-chunīn mīfarmāyand basā nafsī dar ẓuhūrī az
aʿlā ʿuluvv-i ahl-i jannat ast va ḥīn-i ẓuhūr-i baʿd az adnā dunuvv-i
ahl-i nār mīgardad.

Now, what you suppose to be mentioned and widely known among people has
never had, and will never have, any standing before God. Everything is
conditional upon and dependent on the confirmation of this most
wondrous, most unassailable, most holy Manifestation, if you only knew.
Even supposing that what you have mentioned is acknowledged, He
says---exalted be His mention---that at the time of the other
Manifestation the people of Sinai become dead and as nothing, the
Letters and the guides of light remain veiled, and the tree of
affirmation is reckoned among negation. He likewise says: ``Many a soul
who in one Manifestation is among the highest of the inhabitants of
paradise becomes, at the time of the next Manifestation, among the
lowest of the inhabitants of the fire.''

\subsection{119}\label{section-614}

\begin{otherlanguage}{persian}

بگو بمعرضين و مغلّين كه اى بی‌‌بصران كه خود را از كوثر معانى و تسنيم معارف
و سلسبيل بيان محروم و ممنوع نموده‌ايد در اعلى علوّ تفكّر نمائيد شايد بقطره
از بحر اعظم مشروب گرديد و هم‌چنين در مقام ديگر ميفرمايد عزّ ذكره كه اعلى
سكان جنّت قبل در ادنى دنوّ نار ساكن ميگردند حال بصر طاهر لازم تا در اين
بيانات مظهر امر تفكّر نمايد لم‌يزل مشركين در احيان ظهور بكلمات قبلش تمسّك
جسته بر مظهر نفس الهى اعتراض مينمودند چنانچه شما هم اليوم بكلمات قبلش
تمسّك جسته بر نفسش اعتراض نموده و مينمائيد ويل لهم و لكم و لمن يتّبعكم من
عذاب يوم عظيم

\end{otherlanguage}

Bigū bi-muʿriḍīn va mughillīn kih ay bī-baṣarān kih khud rā az kawthar-i
maʿānī va tasnīm-i maʿārif va salsabīl-i bayān maḥrūm va mamnūʿ
nimūdih-īd, dar aʿlā ʿuluvv tafakkur namāyīd, shāyad bi-qaṭrih az baḥr-i
aʿẓam mashrūb gardīd. Va ham-chunīn dar maqām-i dīgar mīfarmāyad, ʿazza
dhikruhu, kih aʿlā-yi sukkān-i jannat-i qabl dar adnā dunuvv-i nār sākin
mīgardand. Ḥāl baṣar-i ṭāhir lāzim tā dar īn bayānāt-i Maẓhar-i Amr
tafakkur namāyad. Lam-yazal mushrikīn dar aḥyān-i ẓuhūr bi-kalimāt-i
qablash tamassuk justih bar Maẓhar-i Nafs-i Ilāhī iʿtirāḍ mīnamūdand,
chinānchih shumā ham al-yawm bi-kalimāt-i qablash tamassuk justih bar
nafsash iʿtirāḍ nimūdih va mīnamāyīd. Waylun lahum wa-lakum wa-li-man
yattabiʿukum min ʿadhābi yawmin ʿaẓīm.

Say to those who turn away and harbor malice: ``You who lack sight and
have deprived and debarred yourselves from the Kawthar of meanings, the
Tasním of knowledge, and the Salsabíl of utterance! Reflect upon `the
highest,' so that you may perhaps drink a drop from the Most Great
Ocean.'' Likewise, in another passage He says---exalted be His
mention---that the highest inhabitants of the former paradise come to
dwell in the lowest depths of the fire. Pure sight is now needed to
reflect upon these utterances of the Manifestation of the Cause. At the
time of each Manifestation, the idolaters have always clung to His
previous words and raised objections against the Manifestation of the
divine Self, just as you today have clung to His previous words and have
raised, and continue to raise, objections against His own Self. Woe to
them, to you, and to those who follow you, from the punishment of a
tremendous day!

\subsection{120}\label{section-615}

\begin{otherlanguage}{persian}

و بشأنى غفلت نموده‌ايد كه كلماتى كه لايق ابصر مشركين است كه بآن نظر
نمايند بمقرّ اطهر امنع ميفرستيد در اينمقام ظهور قبل روح ما سواه فداه
بيانى فرموده‌اند خالصاً لوجه اللّه ذكر ميشود كه شايد معرضين متنبّه شوند و از
شمس هدايت محجوب نمانند قوله عزّ ذكره مثل آن ایّام مثل شمس است بستارگان و
هم‌چنين مثل اهل آن ظهور بالنّسبة بظاهر مثل همين است از اين جهت است كه در
آن ظهور عارفين بآن قلم را جريان ندهند و حيا كنند از تصانيف و تأليف و
انشاءات خود و ظهورات و بروزات خود انتهى ايكاش اين يك نصيحت سلطان احديّه
را استماع مينمودند و ارتكاب نمينمودند آنچه را كه از آن نهى شده‌اند لم ادر
باىّ صراط يمشون و فى اىّ مناهج يسلكون يستدلّون بآيات اللّه ثمّ بمنزله يكفرون

\end{otherlanguage}

Va bi-shaʾnī ghaflat nimūdih-īd kih kalimātī kih lāyiq-i abṣar-i
mushrikīn ast kih bi-ān naẓar namāyand bi-maqarr-i aṭhar-i amnaʿ
mīfiristīd. Dar īn maqām ẓuhūr-i qabl, rūḥ-i mā sivāh fidāh, bayānī
farmūdih-and, khāliṣan li-wajhi Allāh dhikr mīshavad kih shāyad muʿriḍīn
mutanabbih shavand va az shams-i hidāyat maḥjūb namānand. Qawluhu ʿazza
dhikruhu: Mithl-i ān ayyām mithl-i shams ast bi-sitārigān va ham-chunīn
mithl-i ahl-i ān ẓuhūr bi-al-nisbih bi-ẓāhir mithl-i hamīn ast. Az īn
jihat ast kih dar ān ẓuhūr ʿārifīn bi-ān qalam rā jarayān nadihand va
ḥayā kunand az taṣānīf va taʾlīf va inshāʾāt-i khud va ẓuhūrāt va
burūzāt-i khud. Intahā. Ay kāsh īn yik naṣīḥat-i Sulṭān-i Aḥadiyyih rā
istimāʿ mīnamūdand va irtikāb namīnamūdand ānchih rā kih az ān nahy
shudih-and. Lam adri bi-ayyi ṣirāṭin yamshūna wa-fī ayyi manāhija
yaslukūn. Yastadillūna bi-āyāti Allāhi thumma bi-munzilihi yakfurūn.

You have become so heedless that you send to the most pure, most
unassailable Seat words fit for the most discerning of the idolaters to
examine. On this subject the previous Manifestation---may the souls of
all besides Him be sacrificed for Him---made a statement. It is quoted
solely for the sake of God, so that those who turn away may perhaps be
awakened and not remain veiled from the Sun of guidance. His words,
exalted be His mention: ``The likeness of those days is that of the sun
in relation to the stars, and the likeness of the people of that
Manifestation in relation to the One who appears is likewise the same.
For this reason, in that Manifestation those who recognize Him do not
set their pens in motion, and feel shame at their own compositions,
writings, and literary productions, and their own manifestations and
appearances.'' End of quotation. If only they would listen to this one
counsel of the Sovereign of Oneness and refrain from committing what
they have been forbidden! I do not know what path they walk or what ways
they follow. They adduce the verses of God as evidence, then disbelieve
in the One who revealed them.

\subsection{121}\label{section-616}

\begin{otherlanguage}{persian}

و اينكه نوشته كه حضرت باب تصديق ميرزا يحيى نموده اولّا تو از آن تصديق
مطّلع نيستى كه مقصود چه بوده و مراد كه چه كه حوريات معانى الهيّه در غرفات
كلمات عصمت ربانيّه لم‌يزل از انظر بريّة مستور بوده و خواهد بود لم يطمئهنّ
احد الّا نفس الظّهور و الّذينهم كانوا ادلّآء عليه فى يوم ظهوره فو‌اللّه ما صدق
الّا هذا الجمال الّذى وقع بين اظفار المشركين و تحت مخاليب المنكرين ثمّ السن
الغافلين و المغلّين و بر فرض تصديق اين قول اگر ببصر دقيق ملاحظه نمائى اين
تصديق را درباره كلّ مشاهده مينمائى چه كه ميفرمايد عزّ ذكره كه شكّى نبوده كه
جميع نفوسى كه ببيان مومن شده‌اند از انبیاء و اولیاء و اصفیاء بوده و هستند
و بعد در مقام ديگر ميفرمايند قوله عزّ ذكره كلّ از براى اين است كه در يوم
ظهور باسماء محتجب نمانى بلكه نظر كنى بما يقوم به الأسماء من كلّ شیء حتّى
ذكر النّبی‌ فانّ ذلك الاسم يخلق به بما ينزل اللّه انتهى

\end{otherlanguage}

Va īnke neveshte ke Ḥażrat-e Bāb taṣdīq-e Mīrzā Yaḥyā namūde, avvalan to
az ān taṣdīq moṭṭaleʿ nīstī ke maqṣūd che būde va morād ke che, ke
ḥūriyyāt-e maʿānī-ye elāhiyye dar ghorofāt-e kalemāt-e ʿeṣmat-e
rabbāniyye lam-yazal az {[}corrupted text{]} bariyye mastūr būde va
khāhad būd. Lam {[}corrupted text{]} aḥad illā nafs al-ẓuhūr
wa-alladhīna-hum kānū adillāʾa ʿalayhi fī yawmi ẓuhūrihi. Fa-wa-llāhi mā
ṣaddaqa illā hādhā al-jamāl alladhī waqaʿa bayna aẓfār al-mushrikīn
wa-taḥta makhālīb al-munkirīn thumma alsun al-ghāfilīn wa-al-mughillīn.
Va bar farż-e taṣdīq-e īn qawl agar be-baṣar-e daqīq molāḥeẓe namāʾī īn
taṣdīq rā darbāre-ye koll moshāhede mīnamāʾī, che ke mīfarmāyad ʿazza
dhikruhu ke shakkī nabūde ke jamīʿ-e nofūsī ke be-Bayān moʾmen shode-and
az anbiyā va awliyā va aṣfiyā būde va hastand. Va baʿd dar maqām-e dīgar
mīfarmāyand, qawluhu ʿazza dhikruhu: Koll az barā-ye īn ast ke dar
yawm-e ẓohūr be-asmāʾ moḥtajeb namānī balke naẓar konī bi-mā yaqūmu bihi
al-asmāʾ min kulli shayʾ ḥattā dhikr al-nabī fa-inna dhālika al-ism
yukhlaqu bihi bi-mā yunazzilu Allāh. Intahā.

As for your statement that the Báb endorsed Mírzá Yaḥyá: first, you do
not understand that endorsement or what its purpose and intention were.
For the heavenly maidens of divine meanings, within the chambers of
words endowed with divine infallibility, have always been and will
remain concealed from the {[}corrupted original text{]} of creation. No
one has {[}corrupted original text{]} except the Manifestation Himself
and those who pointed to Him on the day of His appearance. By God, He
endorsed none other than this Beauty, who has fallen among the claws of
the idolaters, beneath the talons of the deniers, and then beneath the
tongues of the heedless and the malicious. Even supposing this claim to
be true, if you examine it with discerning sight, you will see that this
endorsement applies to everyone. For He says, exalted be His mention,
that there has been no doubt that all souls who have believed in the
Bayán have been and are prophets, saints, and chosen ones. Then,
elsewhere, He says, exalted be His mention: ``All this is so that, on
the day of the Manifestation, you may not remain veiled by names, but
may instead look to that through which the names of all things subsist,
even the designation `Prophet'; for that name is created through Him by
what God reveals.'' End of quotation.

\subsection{122}\label{section-617}

\begin{otherlanguage}{persian}

بسيار حيرت است كه هنوز در عقبه اسماء واقفيد آفرين خوب وصيّت الهى را عمل
نموده‌ايد بشأنى غافل شده‌ايد كه باسماء عابد و عاكفيد چنانچه متّصل ذكر
مينمائيد و بما يقوم به الأسماء در كلّ حين سيف كين كشيده و ميكشيد لا
تشعرون ما تعلمون

\end{otherlanguage}

Besyār ḥeyrat ast ke hanūz dar ʿaqabe-ye asmāʾ vāqefīd. Āfarīn, khūb
vaṣiyyat-e elāhī rā ʿamal namūde-īd! Be-shaʾnī ghāfel shode-īd ke
be-asmāʾ ʿābed va ʿākefīd, chenānche mottaṣel dhekr mīnamāʾīd, va bi-mā
yaqūmu bihi al-asmāʾ dar koll-e ḥīn seyf-e kīn keshīde va mīkeshīd. Lā
tashʿurūna mā taʿlamūn.

It is astonishing that you are still stopped at the obstacle of names.
Well done! How well you have carried out the divine counsel! You have
become so heedless that you worship names and remain devoted to them, as
you continually declare, while at every moment you have drawn, and
continue to draw, the sword of malice against Him through whom the names
subsist. You do not perceive what you know.

\subsection{123}\label{section-618}

\begin{otherlanguage}{persian}

جميع استدلال و اوصاف شما كه ذكر نموده و يا نمائيد نسبت باىّ نفس كان راجع
ببيان ميشود و در كلّ بيان منصوص است كه بيان و ما فى البيان معلّق بمشيّت آن
سلطان احديّه بوده و خواهد بود و آن جمال قدم قادر است بر اينكه جميع خلق او
را بكلمه لا نفى فرمايد و باسفل النّار راجع نمايد و يا جميع را بقول بلى در
ظلّ اثبات كشاند و باعلى ذروه جنّت عرفان مقرّ دهد لايسئل عمّا يفعل و كلّ عن
كلّشیء يسئلون

\end{otherlanguage}

Jamīʿ-e estedlāl va awṣāf-e shomā ke dhekr namūde va yā namāʾīd nesbat
be-ayy-e nafs kān rājeʿ be-Bayān mīshavad, va dar koll-e Bayān manṣūṣ
ast ke Bayān va mā fī al-Bayān moʿallaq be-mashiyyat-e ān solṭān-e
aḥadiyye būde va khāhad būd. Va ān jamāl-e qedam qāder ast bar īnke
jamīʿ-e khalq-e ū rā be-kaleme-ye lā nafy farmāyad va be-asfal-e nār
rājeʿ namāyad va yā jamīʿ rā be-qawl-e balā dar ẓell-e ethbāt kashānad
va be-aʿlā dhorve-ye jannat-e ʿerfān maqarr dehad. Lā yusʾalu ʿammā
yafʿalu wa-kullun ʿan kulli shayʾin yusʾalūn.

All the arguments and descriptions you have offered or may offer
concerning any person whatsoever refer back to the Bayán. Throughout the
Bayán it is expressly stated that the Bayán and all it contains have
been and will remain dependent on the will of that Sovereign of Oneness.
That ancient Beauty has the power to negate all His creation with the
word ``No'' and consign them to the lowest depths of the fire, or to
draw them all beneath the shelter of affirmation with the word ``Yes''
and establish them on the loftiest summit of the paradise of
recognition. He is not questioned about what He does, while everyone is
questioned about everything.

\subsection{124}\label{section-619}

\begin{otherlanguage}{persian}

و كاش ادراك مينمودى كه درباره كه تصديق فرموده فواللّه الّذى لا اله الّا هو
انّه ما اصبح الّا بذكره و ما امسى الّا بثنائه و ما تحرّك الّا بحبّه و ما تنفّس
الّا فى امره و انتم اعرضتم عنه و اعترضتم عليه و تمسّكتم بالّذى كان عند اللّه
احقر منكم لو انتم تعلمون محبوب امكان را كه درباره‌اش نقطه بيان ميفرمايد
قوله عزّ ذكره لو لا كان كتابه ما نزل ذلك الكتاب و لو لا كان نفسه ما
اظهرنى اللّه از او اعراض نموده و مينويسيد كه در بيان فلان را وصف فرموده
مع آنكه نه وصف را دانسته‌ايد و نه واصف و موصوف را شناخته‌ايد فو‌اللّه لو
تعرف ما فعلت لتبكى على نفسك و لو عرفت هذا الظّهور و كان عندك ملأ السّموات
و الأرض من الذّهب و الفضّة لانفقتها لتدخل فى ظلّه و تسمع نغمة من نغماته و
لكن لما احتجبت عن ذلك صرت محروماً عن نفحات اللّه المهيمن القيّوم

\end{otherlanguage}

Va kāsh edrāk mīnamūdī ke darbāre-ye ke taṣdīq farmūde. Fa-wa-llāhi
alladhī lā ilāha illā huwa innahu mā aṣbaḥa illā bi-dhikrihi wa-mā amsā
illā bi-thanāʾihi wa-mā taḥarraka illā bi-ḥubbihi wa-mā tanaffasa illā
fī amrihi wa-antum aʿraḍtum ʿanhu wa-iʿtaraḍtum ʿalayhi wa-tamassaktum
bi-alladhī kāna ʿinda Allāhi aḥqara minkum law antum taʿlamūn. Maḥbūb-e
emkān rā ke darbāre-ash Noqṭe-ye Bayān mīfarmāyad, qawluhu ʿazza
dhikruhu: Law lā kāna kitābuhu mā nazala dhālika al-kitāb wa-law lā kāna
nafsuhu mā aẓharanī Allāh, az ū eʿrāż namūde va mīnevīsīd ke dar Bayān
folān rā vaṣf farmūde, maʿ ānke na vaṣf rā dāneste-īd va na vāṣef va
mawṣūf rā shenākhte-īd. Fa-wa-llāhi law taʿrifu mā faʿalta la-tabkī ʿalā
nafsika wa-law ʿarafta hādhā al-ẓuhūr wa-kāna ʿindaka milʾ al-samāwāt
wa-al-arḍ min al-dhahab wa-al-fiḍḍa la-anfaqtahu li-tadkhula fī ẓillihi
wa-tasmaʿa naghmatan min naghamātihi wa-lākin lammā iḥtajabta ʿan
dhālika ṣirta maḥrūman ʿan nafaḥāt Allāhi al-muhaymin al-qayyūm.

If only you understood whom He endorsed! By God, besides whom there is
no god, He passed no morning except in remembrance of Him, no evening
except in His praise; He made no movement except out of love for Him and
drew no breath except in His Cause. Yet you have turned away from Him,
opposed Him, and clung to one who, in God's sight, is more abject than
you, if you only knew. You have turned away from the Beloved of all
existence, concerning whom the Point of the Bayán says, exalted be His
mention, ``Had it not been for His Book, that Book would not have been
revealed; and had it not been for His own Self, God would not have
manifested Me.'' Yet you write that He praised a certain person in the
Bayán, although you have neither understood the praise nor recognized
the one who praised or the one who was praised. By God, if you knew what
you had done, you would weep for yourself. If you recognized this
Manifestation and possessed enough gold and silver to fill the heavens
and the earth, you would spend it to enter His shadow and hear a single
melody among His melodies. But because you have remained veiled from
this, you have been deprived of the breezes of God, the All-Protecting,
the Self-Subsisting.

\subsection{125}\label{section-620}

\begin{otherlanguage}{arabic}

ان يا اخى تاللّه الّذينهم كفروا اليوم قد كان لهم غبن بعد غبن و خسران بعد
خسران و ذلّة بعد ذلّة و نقصان بعد نقصان و لكن هم لايفقهون الّا حين الّذى
يأتيهم ملائكة العذاب و يتحرّك عيونهم عن سطوات الموت اذا يجزعون و يقولون
هل لنا من سبيل اذا تضرب على فمهم بأيادى القهر و ترجعهم الى مثوئهم فى قعر
النّار كذلك قضى الأمر من لدى اللّه المقتدر القهّار و انّك فامح ما عندك ثمّ
اثبت ما عنداللّه على لوح صدرك لتجده ممرّداً منيراً بأنوار شموس الحكمة و
المعانى كذلك يعظك هذا العبد لعلّ تكون من السّامعين

\end{otherlanguage}

In yā akhī ta-llāhi alladhīna-hum kafarū al-yawm qad kāna lahum ghabnun
baʿda ghabn wa-khusrānun baʿda khusrān wa-dhillatun baʿda dhilla
wa-nuqṣānun baʿda nuqṣān wa-lākin hum lā yafqahūna illā ḥīna alladhī
taʾtīhim malāʾikat al-ʿadhāb wa-yataḥarraku ʿuyūnuhum ʿan saṭawāt
al-mawt idhan yajzaʿūna wa-yaqūlūna hal lanā min sabīl idhan taḍribu
ʿalā famihim bi-ayādī al-qahr wa-turjiʿuhum ilā mathwāhum fī qaʿr al-nār
kadhālika quḍiya al-amru min ladā Allāhi al-muqtadir al-qahhār wa-innaka
fa-imḥu mā ʿindaka thumma uthbit mā ʿinda Allāh ʿalā lawḥi ṣadrika
li-tajidahu mumarradan munīran bi-anwār shumūs al-ḥikma wa-al-maʿānī
kadhālika yaʿiẓuka hādhā al-ʿabdu laʿalla takūna min al-sāmiʿīn.

O my brother! By God, those who have disbelieved today have suffered
deprivation upon deprivation, loss upon loss, abasement upon abasement,
and diminishment upon diminishment. Yet they will not understand until
the angels of punishment come to them and their eyes move under the
assaults of death. Then they will cry out in anguish and say, ``Is there
any way open to us?'' Then the angels will strike their mouths with the
hands of subjugation and return them to their dwelling in the depths of
the fire. Thus has the matter been decreed by God, the All-Powerful, the
All-Subduing. As for you, erase what you possess, then inscribe what is
with God upon the tablet of your breast, so that you may find it
polished and illumined by the light of the suns of wisdom and meanings.
Thus does this servant counsel you, so that you may be among those who
listen.

\subsection{126}\label{section-621}

\begin{otherlanguage}{persian}

و اينكه نوشته كه فرموده‌اند وصىّ و خليفه هر چه هست بعد از من او است فو‌الّذى
تفرّد بالعزّ و الكبريا و خلق الأشیآء بلا مثال بكلمة العليا كه اگر نفسى
نفحه از رضوان بيان استشمام نموده باشد بچنين كلمات تكلّم نمينمايد بل يفرّ
من الّذين يتكلّمون بتلك الكلمات شما هنوز حشر بديع و روح بديع و كلمات بديع
ادراك ننموده‌ايد امر بيان مقدّس از آن امورى است كه از پيش گذشت قد طوى ما
عندكم من قبل و بسط بساط آخر باذن اللّه العلى الأظهر و لكن انتم ما ادركتم
لا اوّله و لا آخره حمد محبوب را كه مغلّين را چنان محروم كرده كه الآن احقر
از ملل قبل بين يدى مذكور و مشهودند مثل آن است كه جامه از يكقطعه حديد
پوشيده‌اند كه ابدا قادر بر حركت نيستند و هم‌چنين قلب و فؤادشان قدم
برداشته‌اند تازه باثر قدم اهل فرقان گذارده‌اند

\end{otherlanguage}

Va īnke neveshte ke farmūde-and vaṣī va khalīfe har che hast baʿd az man
ū ast, fa-wa-alladhī tafarrada bi-al-ʿizz wa-al-kibriyā wa-khalaqa
al-ashyāʾ bi-lā mithāl bi-kalimati al-ʿulyā, ke agar nafsī nafḥe az
Reżvān-e Bayān esteshmām namūde bāshad be-chenīn kalemāt takallom
namīnamāyad, bal yafirru min alladhīna yatakallamūna bi-tilka
al-kalimāt. Shomā hanūz ḥashr-e badīʿ va rūḥ-e badīʿ va kalemāt-e badīʿ
edrāk nanamūde-īd. Amr-e Bayān moqaddas az ān omūrī ast ke az pīsh
godhasht. Qad ṭuwiya mā ʿindakum min qabl wa-busiṭa bisāṭun ākhar
bi-idhni Allāhi al-ʿalī al-aẓhar wa-lākin antum mā adraktum lā awwalahu
wa-lā ākhirahu. Ḥamd-e maḥbūb rā ke moghellīn rā chenān maḥrūm karde ke
al-ān aḥqar az melal-e qabl bayna yaday madhkūr va mashhūdand. Methl-e
ān ast ke jāme az yek qeṭʿe ḥadīd pūshīde-and ke abadan qāder bar
ḥarekat nīstand va ham-chenīn qalb va foʾādshān. Qadam bardāshte-and
tāze be-athar-e qadam-e ahl-e Forqān godhāshte-and.

As for your statement that He said, ``Whatever the office of executor
and successor may be, after Me it belongs to him'': by Him who alone
possesses might and majesty and who created all things without a model
through His exalted Word, if anyone had inhaled a single breeze from the
Riḍván of the Bayán, he would not utter such words. Rather, he would
flee from those who utter them. You have not yet understood the wondrous
resurrection, the wondrous spirit, or the wondrous words. The Cause of
the Bayán is sanctified above those things that came before. What you
formerly possessed has been rolled up, and another expanse has been
spread out by the permission of God, the Exalted, the Most Manifest. Yet
you have understood neither its beginning nor its end. Praise be to the
Beloved, who has so utterly deprived the malicious that they are now
mentioned and seen before Him as more abject than the religious
communities of the past. It is as though they have put on a garment made
from a single piece of iron, so that they are utterly incapable of
movement; so too are their hearts and inmost beings. They have taken a
step only to set it in the footprints of the people of the Qur'an.

\subsection{127}\label{section-622}

\begin{otherlanguage}{persian}

بشنويد و بشكنيد اين سدّه‌هاى مانعه را و بفضاى بديع وسيع مالك يوم الّدين
درآئيد تا عوالمى غير عوالم خود مشاهده كنيد آنقدر بر خود سخت مگيريد ترحّم
نمائيد بر خود و بر مستضعفين عباد تازه ببقيّه وهم ما ترك من اهل الفرقان
تمسّك جسته‌ايد گاهى ذكر خلافت مجعول و وصايت مكذوب مينمائيد و گاهى كلمات
جعل نموده بنقطه اولى روح ما سواه فداه نسبت ميدهيد بشنو نداى ربّ علىّ اعلى
را كه مخاطباً للمرايا ميفرمايند

\end{otherlanguage}

Beshnavīd va beshkanīd īn sodde-hā-ye māneʿe rā va be-fażā-ye badīʿ-e
vasīʿ-e Mālek-e Yawm-e Dīn darāʾīd tā ʿavālemī gheyr-e ʿavālem-e khod
moshāhede konīd. Ānqadr bar khod sakht magīrīd. Taraḥḥom namāʾīd bar
khod va bar mostażʿafīn-e ʿebād. Tāze be-baqiyye-ye vahm-e mā taraka min
ahl al-Furqān tamassok joste-īd. Gāhī dhekr-e khelāfat-e majʿūl va
veṣāyat-e makdhūb mīnamāʾīd va gāhī kalemāt jaʿl namūde be-Noqṭe-ye Ūlā
rūḥ-e mā sevāh fedāh nesbat mīdehīd. Beshnow nedā-ye Rabb-e ʿAlī-ye Aʿlā
rā ke mukhāṭaban li-al-marāyā mīfarmāyand.

Listen, break down these obstructing barriers, and enter the wondrous,
boundless expanse of the Lord of the Day of Judgment, so that you may
behold worlds other than your own. Do not be so hard on yourselves. Have
mercy on yourselves and on the weak among His servants. You have now
clung to the remnants of delusion left by the people of the Qur'an. At
times you speak of a fabricated successorship and a falsely claimed
executorship; at other times you invent words and attribute them to the
Primal Point---may the souls of all besides Him be sacrificed for Him.
Hear the call of the Lord, the Exalted, the Most High, who says,
addressing the Mirrors:

\subsection{128}\label{section-623}

\begin{otherlanguage}{persian}

اذا اشرقت شمس البهآء عن افق البقآء انتم فاحضر وا بين يدى العرش ایّاکم ان
تقعدوا بين يديه او تسئلوا ما لا اذن لكم اتّقوا اللّه يا ايّها المرايا كلكم
اجمعون ان اسئلوا من بدايع فضله ليظهر لكم ما شآء و اراد لانّ كلّ الفضل فى
ذلك اليوم يطوف حول عرشه و يظهر من عنده ان انتم تعلمون ان اصمتوا تلقآء
العرش لانّ الصّمت فى ذلك اليوم لخير عمّا خلق بين السّموات و الأرض و لا
تكونوا من الّذين اوتوا العلم من قبل و استكبروا به على اللّه المهيمن القيّوم
لانّه فى ذلك اليوم قد كان فوق كلّ ذى علم علّاماً عالماً عليم و فوق كلّ ذى قدرة
كان قدّاراً مقتدراً قدير و فوق كلّ ذى عظمة قد كان عظّاماً متعظّماً عظيم و فوق كلّ
ذى رفعة قد كان رفّاعاً مرتفعاً رفيع و فوق كلّ ذى عزّ قد كان عزّازا متعزّزاً عزيز
وفوق كلّ ذى حكم قد كان حكّاماً حاكماً حكيم و فوق كلّ ذى جود قد كان جوّاداً
جاوداً جويد و فوق كلّ ذى امر قد كان امّاراً آمرا امير و فوق كّل ذى فضل قد كان
فضّالاً فاضلا فضيل

\end{otherlanguage}

Idhā ashraqat shams al-Bahāʾ ʿan ufuq al-baqāʾ antum fa-uḥḍurū bayna
yaday al-ʿarsh iyyākum an taqʿudū bayna yadayhi aw tasʾalū mā lā idhna
lakum ittaqū Allāha yā ayyuhā al-marāyā kullukum ajmaʿūn in isʾalū min
badāʾiʿ faḍlihi li-yuẓhira lakum mā shāʾa wa-arāda li-anna kulla al-faḍl
fī dhālika al-yawm yaṭūfu ḥawla ʿarshihi wa-yaẓharu min ʿindihi in antum
taʿlamūn in uṣmutū tilqāʾa al-ʿarsh li-anna al-ṣamt fī dhālika al-yawm
la-khayrun ʿammā khuliqa bayna al-samāwāt wa-al-arḍ wa-lā takūnū min
alladhīna ūtū al-ʿilma min qabl wa-istakbarū bihi ʿalā Allāhi
al-muhaymin al-qayyūm li-annahu fī dhālika al-yawm qad kāna fawqa kulli
dhī ʿilmin ʿallāman ʿāliman ʿalīm wa-fawqa kulli dhī qudratin kāna
qaddāran muqtadiran qadīr wa-fawqa kulli dhī ʿaẓamatin qad kāna ʿaẓẓāman
mutaʿaẓẓiman ʿaẓīm wa-fawqa kulli dhī rifʿatin qad kāna raffāʿan
murtafiʿan rafīʿ wa-fawqa kulli dhī ʿizzin qad kāna ʿazzāzan
mutaʿazzizan ʿazīz wa-fawqa kulli dhī ḥukmin qad kāna ḥakkāman ḥākiman
ḥakīm wa-fawqa kulli dhī jūdin qad kāna jawwādan jāwidan jawīd wa-fawqa
kulli dhī amrin qad kāna ammāran āmiran amīr wa-fawqa kulli dhī faḍlin
qad kāna faḍḍālan fāḍilan faḍīl.

``When the Sun of Bahá shines from the horizon of eternity, present
yourselves before the Throne. Beware of sitting in His presence or
asking what you have not been permitted to ask. Fear God, O Mirrors, all
of you without exception. Ask for the wonders of His grace, so that He
may reveal to you whatever He wishes and desires; for all grace on that
day circles His Throne and proceeds from Him, if you only knew. Remain
silent before the Throne, for silence on that day is better than all
that has been created between the heavens and the earth. Do not be among
those who were previously given knowledge and, because of it, behaved
arrogantly toward God, the All-Protecting, the Self-Subsisting. For on
that day, above everyone possessed of knowledge, He is the All-Knowing,
the Knowing, the Most Learned; above everyone possessed of power, He is
the All-Powerful, the Omnipotent, the Mighty; above everyone possessed
of greatness, He is the Supremely Great, the Most Majestic, the Great;
above everyone possessed of elevation, He is the Supreme Exalter, the
Exalted, the Lofty; above everyone possessed of might, He is the
All-Mighty, the Supremely Glorious, the Mighty; above everyone possessed
of judgment, He is the Supreme Judge, the Ruler, the Wise; above
everyone possessed of generosity, He is the All-Bountiful, the Bestower,
the Generous; above everyone possessed of command, He is the Supreme
Ordainer, the Commander, the Sovereign; and above everyone possessed of
grace, He is the All-Gracious, the Bestower of Grace, the Gracious.

\subsection{129}\label{section-624}

\begin{otherlanguage}{arabic}

كلكم خلقتم للقائه و الحضور فى ذلك الموقف الأعظم العظيم و انّه قد ينزل من
سمآء جوده ما هو خير لكم و ما نزل من عنده انّه ليغنيكم عن العالمين لن ينفع
اليوم كلّ ذى علم علمه و لا كلّ ذى فضل فضله و لا كلّ ذى عظمة عظمته و لا كلّ
ذى قدرة قدرته و لا كلّ ذى ذكر ذكره و لا كلّ ذى عمل عمله و لا كلّ ذى ركوع
ركوعه و لا كلّ ذى سجود سجوده و لا كلّ ذى توجّه توجّهه و لا كلّ ذى شرف شرفه و
لا كلّ ذى نسب عالى نسبه و لا كلّ ذى حسب متعالى حسبه و لا كلّ ذى بيان بيانه
ولا كلّ ذى نور نوره و لا كلّ ذى اسم اسمه لانّ كلّ ذلك و كلّ ما عرفتم و ادركتم
كلّها قد خلق بقوله كن فيكون و انّه لو یشآء ليبعث الممكنات بكلمة من عنده
ليقدر و انّه قد كان فوق ذلك لقدّاراً مقتدراً قدير

\end{otherlanguage}

Kullukum khuliqtum li-liqāʾihi wa-al-ḥuḍūr fī dhālika al-mawqif al-aʿẓam
al-ʿaẓīm wa-innahu qad yunazzilu min samāʾi jūdihi mā huwa khayrun lakum
wa-mā nazala min ʿindihi innahu la-yughnīkum ʿan al-ʿālamīn lan yanfaʿa
al-yawm kulla dhī ʿilmin ʿilmuhu wa-lā kulla dhī faḍlin faḍluhu wa-lā
kulla dhī ʿaẓamatin ʿaẓamatuhu wa-lā kulla dhī qudratin qudratuhu wa-lā
kulla dhī dhikrin dhikruhu wa-lā kulla dhī ʿamalin ʿamaluhu wa-lā kulla
dhī rukūʿin rukūʿuhu wa-lā kulla dhī sujūdin sujūduhu wa-lā kulla dhī
tawajjuhin tawajjuhuhu wa-lā kulla dhī sharafin sharafuhu wa-lā kulla
dhī nasabin ʿālin nasabuhu wa-lā kulla dhī ḥasabin mutaʿālin ḥasabuhu
wa-lā kulla dhī bayānin bayānuhu wa-lā kulla dhī nūrin nūruhu wa-lā
kulla dhī ismin ismuhu li-anna kulla dhālika wa-kulla mā ʿaraftum
wa-adraktum kulluhā qad khuliqa bi-qawlihi kun fa-yakūn wa-innahu law
yashāʾu li-yabʿatha al-mumkināt bi-kalimatin min ʿindihi la-yaqdiru
wa-innahu qad kāna fawqa dhālika la-qaddāran muqtadiran qadīr.

All of you were created to meet Him and to be present at that most great
and mighty gathering. He sends down from the heaven of His bounty what
is best for you, and whatever descends from Him will surely make you
independent of all the worlds. Today, no one's knowledge will benefit
him, nor anyone's grace, greatness, power, remembrance, deeds, bowing,
prostration, turning toward God, honor, lofty lineage, exalted standing,
utterance, light, or name. For all these things, and everything you have
known and understood, were created through His saying, `Be,' and it is.
If He wished to raise up all contingent beings through a word from
Himself, He would assuredly have the power; and beyond even that, He is
indeed the All-Powerful, the Omnipotent, the Mighty.

\subsection{130}\label{section-625}

\begin{otherlanguage}{arabic}

ایّاکم يا ايّها المرايا ان تغرّنكم الأسماء فى ذلك اليوم فاعلموا انتم و ما
فوقكم و مادونكم قد خلقتم لذلك اليوم اتّقوا اللّه و لا ترتكبوا ما يحزن به
فؤاده و تكوننّ من الغافلين ربّما يظهر بالحقّ و انتم راقدون على مقاعدكم و
يأتيكم رسله بالواح عزّ مبين و انتم تستكبرون و تحكمون عليه ما لا تحكمون به
على انفسكم و تقولون ليس هذا من اللّه المهيمن القيّوم

\end{otherlanguage}

Iyyākum yā ayyuhā al-marāyā an taghurrannakum al-asmāʾ fī dhālika
al-yawm fa-iʿlamū antum wa-mā fawqakum wa-mā dūnakum qad khuliqtum
li-dhālika al-yawm ittaqū Allāha wa-lā tartakibū mā yaḥzanu bihi
fuʾāduhu wa-takūnunna min al-ghāfilīn rubbamā yaẓharu bi-al-ḥaqq
wa-antum rāqidūna ʿalā maqāʿidikum wa-yaʾtīkum rusuluhu bi-alwāḥi ʿizzin
mubīn wa-antum tastakbirūna wa-taḥkumūna ʿalayhi mā lā taḥkumūna bihi
ʿalā anfusikum wa-taqūlūna laysa hādhā min Allāhi al-muhaymin al-qayyūm.

Beware, O Mirrors, lest names deceive you on that day. Know that you,
all that is above you, and all that is below you were created for that
day. Fear God, and do not commit what would grieve His heart and make
you among the heedless. He may appear in truth while you lie asleep upon
your seats, and His messengers may come to you with Tablets of manifest
glory, while you behave arrogantly, pass judgments against Him that you
would not pass against yourselves, and say, `This is not from God, the
All-Protecting, the Self-Subsisting.'

\subsection{131}\label{section-626}

\begin{otherlanguage}{persian}

سبحانك اللّهم يا الهى انت تعلم بانّى بلغت كلمتك و ما قصرت فيما امرتنى به
اسئلك بأن تعتصمنّ فى ذلك اليوم خلق البيان بأن لا يعترضوا عليك و لا
يجادلوا بياتك و لتحفظهم يا الهى بقدرتك الّتى استطلتها على العالمين

\end{otherlanguage}

Subḥānaka Allāhumma yā ilāhī anta taʿlamu bi-annī ballaghtu kalimataka
wa-mā qaṣṣartu fī-mā amartanī bihi asʾaluka bi-an taʿtaṣimanna fī
dhālika al-yawm khalq al-Bayān bi-an lā yaʿtariḍū ʿalayka wa-lā yujādilū
{[}corrupted text{]} wa-li-taḥfaẓhum yā ilāhī bi-qudratika allatī
istaṭaltahā ʿalā al-ʿālamīn.

Glory be to You, O God, my God! You know that I have conveyed Your word
and have not fallen short in what You commanded Me. I ask You to protect
the people of the Bayán on that day, so that they may neither oppose You
nor dispute {[}corrupted original text{]}. Preserve them, O my God,
through Your power, which You have exercised over all the worlds.''

\subsection{132}\label{section-627}

\begin{otherlanguage}{persian}

در اينمقام باين كلمات قناعت شد اكثر بيانات الهى كه در اين ظهور امنع اقدس
نازل شده احدى ملاحظه ننموده بعضى آن را از مواضع خود تحريف نموده‌اند و
بعضى را ستر نموده و بعضى را بحسب ظاهر كلمه تغيير داده‌اند اگر چه
فى‌الحقيقة آنچه از بيان نزد هر نفسى موجود است كفايت ميكند او را در تذكّر
باين كلمه اتمّ اعظم و كان اللّه على ذلك شهيد و عليم

\end{otherlanguage}

Dar īn maqām be-īn kalemāt qanāʿat shod. Akthar-e bayānāt-e elāhī ke dar
īn ẓohūr-e amnaʿ-e aqdas nāzel shode aḥadī molāḥeẓe nanamūde. Baʿżī ān
rā az mavāżeʿ-e khod taḥrīf namūde-and va baʿżī rā satr namūde va baʿżī
rā be-ḥasab-e ẓāher kaleme taghyīr dāde-and. Agarche fī al-ḥaqīqe ānche
az Bayān nazd-e har nafsī mawjūd ast kefāyat mīkonad ū rā dar tadhakkor
be-īn kaleme-ye atamm-e aʿẓam. Wa-kāna Allāhu ʿalā dhālika shahīdan
wa-ʿalīm.

At this point, these words have been deemed sufficient. No one has
examined most of the divine utterances revealed concerning this most
unassailable, most holy Manifestation. Some have distorted them out of
their proper context, concealed others, and changed the outward wording
of still others. In truth, however, whatever portion of the Bayán anyone
possesses is sufficient to remind him of this most perfect, most great
Word. God is a witness to this and knows it.

\subsection{133}\label{section-628}

\begin{otherlanguage}{persian}

و اينكه نوشته كه فرموده انّ الأمر ينتهى الى اسم الوحيد لانّ ظهوره بنفسه
حجّة و لايحتاج الى نصّى فو‌اللّه حين الّذى كتبت هذه الكلمة قد بكت كلّ الأشیآء
و ناح اهل ملأ الأعلی و تزلزلت اركان الوجود من الغيب و الشّهود لانّك لاثبات
عدوّا من اعدآء اللّه احتقرت شأن اللّه و عظمته و كبريائه فواللّه انّ الّذى
اتّخذته ربّاً من دون اللّه انّه يحتاج بكلّ شیء و الّذى كلّ يحتاج اليه انّه قد ظهر
بالحقّ فى سنة السّتّين و حينئذ استوى بالحقّ فتبارك اللّه احسن الخالقين

\end{otherlanguage}

Va īnkih nivishtih kih farmūdih inna al-amra yantahī ilā ismi al-Waḥīd
li-anna ẓuhūrahu bi-nafsihi ḥujjatun wa-lā yaḥtāju ilā naṣṣī fa-wa-llāhi
ḥīna alladhī katabta hādhihi al-kalimata qad bakat kullu al-ashyāʾ
wa-nāḥa ahlu malaʾi al-aʿlā wa-tazalzalat arkānu al-wujūdi min al-ghaybi
wa-al-shuhūd li-annaka li-ithbāti ʿaduwwan min aʿdāʾi Allāhi iḥtaqarta
shaʾna Allāhi wa-ʿaẓamatahu wa-kibriyāʾahu fa-wa-llāhi inna alladhī
ittakhadhtahu rabban min dūni Allāhi innahu yaḥtāju bi-kulli shayʾ
wa-alladhī kullun yaḥtāju ilayhi innahu qad ẓahara bi-al-ḥaqqi fī sanati
al-sittīn wa-ḥīnaʾidhin istawā bi-al-ḥaqqi fa-tabāraka Allāhu aḥsanu
al-khāliqīn.

And as for your having written that He said, ``The Cause will devolve
upon the name of the Unique One, because his appearance is itself a
proof and he has no need of My explicit designation'': by God, when you
wrote these words, all things wept, the inhabitants of the Most Exalted
Assembly lamented, and the foundations of existence, both unseen and
visible, were shaken, because, to establish the claim of an enemy of
God, you belittled God's station, greatness, and majesty. By God, the
one you have taken as a lord apart from God is in need of everything,
whereas the One whom all need appeared in truth in the year sixty and
then ascended His throne in truth. Blessed, then, is God, the best of
creators.

\subsection{134}\label{section-629}

\begin{otherlanguage}{arabic}

خف عن اللّه و لاتجر قلمك فى اتباع نفسك و هوئك و لاتكفر بالّذى باشارة من
قلمه ثبت اديان العالمين انّ الّذى قمت لحبّه على النّفاق كان ان يقوم تلقآء
الوجه و يلقى عليه الآيات و كان مفتخراً بذلك اذ انّك اتّبعته من عند نفسك و
به اعترضت على اللّه الّذى خلق كلّ شیء بأمر من عنده لا اله الّا هو الغالب
القاهر المقتدر المتعالى العزيز المحبوب

\end{otherlanguage}

Khaf ʿan Allāhi wa-lā tajur qalamaka fī ittibāʿi nafsika wa-hawāka wa-lā
takfur bi-alladhī bi-ishāratin min qalamihi thabatat adyānu al-ʿālamīn.
Inna alladhī qumta li-ḥubbihi ʿalā al-nifāqi kāna an yaqūma tilqāʾa
al-wajhi wa-yulqā ʿalayhi al-āyātu wa-kāna muftakhiran bi-dhālika idh
innaka ittabaʿtahu min ʿindi nafsika wa-bihi iʿtaraḍta ʿalā Allāhi
alladhī khalaqa kulla shayʾin bi-amrin min ʿindihi lā ilāha illā huwa
al-ghālibu al-qāhiru al-muqtadiru al-mutaʿālī al-ʿazīzu al-maḥbūb.

Fear God, and do not wield your pen in pursuit of your own self and
desires. Do not disbelieve in the One through a gesture of whose Pen the
religions of the worlds were established. The one for whose love you
have risen in hypocrisy used to stand before the Presence and have the
verses recited to him, and he took pride in this. Yet you have followed
him of your own accord and, through him, opposed God, who created all
things by a command from Himself. There is no God but Him, the
All-Victorious, the All-Subduing, the All-Powerful, the Most Exalted,
the Almighty, the Beloved.

\subsection{135}\label{section-630}

\begin{otherlanguage}{persian}

نفسيكه خود شما ذكر ميكنيد كه مرآت است و بر هر ذى بصرى معلوم كه مرآت از
خود تحقّقى نداشته و نخواهد داشت وجودش از براى آن بوده كه خاكى باشد از
تجلّيات شمس در باره او قائل شده كه نقطه اولى روح ماسواه فداه فرموده كه او
بنفسه حجّت است و احتياج بنصّ من ندارد فواللّه يؤيّد السنكم هوئكم و يظهر منه
ما یشآء لاما قدّر من لدى اللّه العزيز القدير اين ظهور كه آيات و ظهوراتش كلّ
عالم را احاطه نموده انكار نموده و نصّ خواسته مع آنكه كلّ بيان صريحاً منصوصاً
در اين ظهور ابدع امنع نازل شده و مع آنكه اين امر لم‌يزل و لا‌یزال مقدّس از
دليل و اشاره و عباره و دلاله بوده و خواهد بود ظهورش بنفسه حجّت و آياتش
برهان اللّه على العالمين

\end{otherlanguage}

Nafsīkih khud-i shumā dhikr mīkunīd kih mirʾāt ast va bar har dhī-baṣarī
maʿlūm kih mirʾāt az khud taḥaqquqī nadāshtih va nakhvāhad dāsht
vujūdash az barā-yi ān būdih kih khākī bāshad az tajalliyāt-i shams dar
bārih-yi ū qāʾil shudih kih Nuqṭih-yi Ūlā rūḥ-i mā-sivāh fidāh farmūdih
kih ū bi-nafsih ḥujjat ast va iḥtiyāj bi-naṣṣ-i man nadārad. Fa-wa-llāhi
yuʾayyidu alsunakum hawākum wa-yaẓharu minhu mā yashāʾu lā mā quddira
min ladā Allāhi al-ʿazīzi al-qadīr. Īn ẓuhūr kih āyāt va ẓuhūrātash
kull-i ʿālam rā iḥāṭih nimūdih inkār nimūdih va naṣṣ khvāstih maʿa ānkih
kull-i Bayān ṣarīḥan manṣūṣan dar īn ẓuhūr-i abdaʿ-i amnaʿ nāzil shudih
va maʿa ānkih īn amr lam-yazal va lā-yazāl muqaddas az dalīl va ishārih
va ʿibārih va dilālih būdih va khvāhad būd ẓuhūrash bi-nafsih ḥujjat va
āyātash burhānu Allāhi ʿalā al-ʿālamīn.

Of the very person whom you yourselves describe as a mirror---and
everyone endowed with sight knows that a mirror has never had, and will
never have, any independent reality, its existence having been for the
purpose of being dust from the radiances of the Sun---you have asserted
that the Primal Point, may the spirit of all besides Him be sacrificed
for Him, said that he is himself a proof and has no need of My explicit
designation. By God, your desire lends strength to your tongues, and
from it emerges whatever it wishes, not what has been ordained by God,
the Almighty, the All-Powerful. You have denied this Manifestation,
whose verses and revelations have encompassed the whole world, and
demanded an explicit designation, even though the whole Bayán was
revealed explicitly and expressly concerning this most wondrous, most
unassailable Manifestation. Moreover, this Cause has always been, and
will forever remain, sanctified above evidence, allusion, expression,
and indication. His manifestation is itself a proof, and His verses are
God's demonstration to all the worlds.

\subsection{136}\label{section-631}

\begin{otherlanguage}{persian}

و ايكاش دلائل السّبعه كه از كلمات اوّليّه نقطه بيان است ملاحظه مينمودى از
قراري كه معلوم ميشود مقصودى جز اعراض از حقّ و اعتراض بر او نداشته و
نداريد لو يأتيكم اللّه بكلّ الآيات لن تؤمنوا به چنانچه نقطه بيان روح
ماسواه فداه از اين ایّام خبر داد‌ه‌اند فانظروا فى كلمات اللّه لعلّ تكوننّ من
المتنبّهين

\end{otherlanguage}

Va ay kāsh Dalāʾil-i Sabʿih kih az kalimāt-i avvalīyih-yi Nuqṭih-yi
Bayān ast mulāḥiẓih mīnamūdī. Az qarārī kih maʿlūm mīshavad maqṣūdī juz
iʿrāḍ az Ḥaqq va iʿtirāḍ bar ū nadāshtih va nadārīd. Law yaʾtīkum Allāhu
bi-kulli al-āyāti lan tuʾminū bihi. Chinānchih Nuqṭih-yi Bayān rūḥ-i
mā-sivāh fidāh az īn ayyām khabar dādih-and. Fa-unẓurū fī kalimāti
Allāhi laʿalla takūnanna min al-mutanabbihīn.

If only you had examined the Seven Proofs, which is among the early
writings of the Point of the Bayán. As is apparent, you have had, and
still have, no purpose other than turning away from the Truth and
opposing Him. Even if God were to come to you with every sign, you would
not believe in Him. Thus did the Point of the Bayán, may the spirit of
all besides Him be sacrificed for Him, foretell these days. Examine,
then, the words of God, that you may be among those who take heed.

\subsection{137}\label{section-632}

\begin{otherlanguage}{persian}

و اينكه نوشته كه فرموده‌اند انّ الأمر ينتهى الى اسم الوحيد فواللّه الّذى لا
اله الّا هو كه وحيد را نشناخته و ندانسته كه مقصود كه بوده و سلّمنا كه
مقصود از وحيد مقصود اوّل معرض باللّه باشد چرا عبارت بعد را ملاحظه ننموده
كه ميفرمايد فان يظهر اللّه فى ایّامك مثلك هذا ما يورثن الأمر من اللّه
الواحد الوحيد مع‌ذلك اين ظهورى كه بجميع شئونات الهيّه ظاهر شده بقدر يك
مرآت هم نبوده كه جميع ص باظفار غلّ و بغضا حولش جمع شده‌ايد و در هر حين
بمكر جديد در تضييع امرش مشغول گشته‌ايد فافّ على الّذينهم اعرضوا عن كلّ
الآیات و اتّخذوا من الكتاب حرفاً و استدلّوا بهآء على ما عندهم و كفروا باللّه
الّذى بامر من عنده نزلت صحايف عزّ كريم و الواح قدس منيع

\end{otherlanguage}

Va īnkih nivishtih kih farmūdih-and inna al-amra yantahī ilā ismi
al-Waḥīd fa-wa-llāhi alladhī lā ilāha illā huwa kih Vaḥīd rā
nashinākhtih va nadānistih kih maqṣūd kih būdih va sallamnā kih maqṣūd
az Vaḥīd maqṣūd-i avval-i muʿriḍ bi-Allāh bāshad chirā ʿibārat-i baʿd rā
mulāḥiẓih nanimūdih kih mīfarmāyad fa-in yuẓhiri Allāhu fī ayyāmika
mithlaka hādhā mā yūrathanna al-amra min Allāhi al-wāḥidi al-waḥīd. Maʿa
dhālik īn ẓuhūrī kih bi-jamīʿ-i shuʾūnāt-i ilāhīyih ẓāhir shudih
bi-qadr-i yik mirʾāt ham nabūdih kih jamīʿ {[}corrupted text{]}
bi-aẓfār-i ghill va baghḍā ḥawlash jamʿ shudih-īd va dar har ḥīn
bi-makr-i jadīd dar taḍyīʿ-i amrash mashghūl gashtih-īd. Fa-uffin ʿalā
alladhīna hum aʿraḍū ʿan kulli al-āyāti wa-ittakhadhū min al-kitābi
ḥarfan wa-istadallū {[}corrupted text{]} ʿalā mā ʿindahum wa-kafarū
bi-Allāhi alladhī bi-amrin min ʿindihi nazalat ṣaḥāʾifu ʿizzin karīmin
wa-alwāḥu qudsin manīʿ.

And as for your having written that He said, ``The Cause will devolve
upon the name of the Unique One'': by God, besides whom there is no God,
you have neither recognized the Unique One nor understood who was
intended. Even if we concede that by the Unique One was intended the one
meant by the foremost of those who turn away from God, why have you not
considered the passage that follows, where He says, ``If God should
bring forth in your days one like you, this is the one who shall be made
heir to the Cause by God, the One, the Unique''? Despite this, has this
Manifestation, who has appeared with every divine attribute, not been
worth even a single mirror, that you have all {[}corrupted original
text{]} gathered around Him with the claws of malice and hatred, and at
every moment busied yourselves with a new scheme to undermine His Cause?
Shame on those who have turned away from all the verses, taken a single
letter from the Book, and adduced {[}corrupted original text{]} in
support of what they possess, and disbelieved in God, at whose command
the scrolls of noble glory and the tablets of unassailable holiness were
revealed.

\subsection{138}\label{section-633}

\begin{otherlanguage}{persian}

اينست كه نقطه بيان روح ما سواه فداه ميفرمايد كه آن ظهور اعظم را اهل بيان
يكى از مؤمنين هم نخواهند شمرد چه كه اگر مؤمن هم بدانند راضى نميشوند از
براى او آنچه را كه از براى خود نمى‌پسندند صدق اللّه العلىّ العظيم

\end{otherlanguage}

Īnast kih Nuqṭih-yi Bayān rūḥ-i mā-sivāh fidāh mīfarmāyad kih ān ẓuhūr-i
aʿẓam rā ahl-i Bayān yikī az muʾminīn ham nakhvāhand shumurd chih kih
agar muʾmin ham bidānand rāḍī namīshavand az barā-yi ū ānchih rā kih az
barā-yi khud namīpasandand. Ṣadaqa Allāhu al-ʿaliyyu al-ʿaẓīm.

This is why the Point of the Bayán, may the spirit of all besides Him be
sacrificed for Him, says that the people of the Bayán will not count
that Most Great Manifestation even as one of the believers. For if they
regarded Him as even a believer, they would not consent to inflict on
Him what they would not accept for themselves. God, the Most High, the
Most Great, has spoken the truth.

\subsection{139}\label{section-634}

\begin{otherlanguage}{arabic}

از جميع اينها گذشته قول الحقّ كلمة واحدة و هى ما نزّلت من عنده ان انتم من
العارفين قال و قوله الحقّ ایّاک ایّاک ایّام ظهوره ان تحتجب بالواحد البيانيّه
فانّ ذلك الواحد خلق عنده و ایّاک ایّاک ان تحتجب بكلمات ما نزلت فى البيان
فانّها كلمات نفسه فى هيكل ظهوره من قبل انتهى و

\end{otherlanguage}

Az jamīʿ-i īnhā gudhashtih qawlu al-Ḥaqqi kalimatun wāḥidatun wa-hiya mā
nuzzilat min ʿindihi in antum min al-ʿārifīn. Qāla wa-qawluhu al-ḥaqqu
iyyāka iyyāka ayyāma ẓuhūrihi an taḥtajiba bi-al-wāḥidi al-bayānīyyati
fa-inna dhālika al-wāḥida khul iqa ʿindahu wa-iyyāka iyyāka an taḥtajiba
bi-kalimāti mā nazalat fī al-Bayāni fa-innahā kalimātu nafsihi fī
haykali ẓuhūrihi min qablu intahā va.

Beyond all this, the word of the Truth is a single word: what has been
revealed by Him, if you are among those who recognize. He said---and His
word is the truth---``Beware, beware, in the days of His manifestation,
lest you be veiled by the Váḥid of the Bayán, for that Váḥid was created
in His presence. Beware, beware lest you be veiled by the words revealed
in the Bayán, for they are His own words in the Temple of His former
manifestation.'' End of quotation. And

\subsection{140}\label{section-635}

\begin{otherlanguage}{persian}

كلّ آنچه نوشته‌اند در صورت تصديق و تسليم رشحى از طمطام يمّ بيان الهى بوده
صد‌هزار امثال اين نفوس بقولى خلق شده و ميشوند و مع‌ذلك افترإ على اللّه و
كذباً عليه و بغياً على حضرته و طغياناً على بهائه و اثباتاً لأعدائه مينويسى
كه حضرت اعلى فرموده كه احتياج بنصّ من ندارد و بهذه الكلمة تبرّء منك كلّ ما
ظهر من القلم فى ممالك اللّه المقتدر العلىّ العليم

\end{otherlanguage}

Kull-i ānchih nivishtih-and dar ṣūrat-i taṣdīq va taslīm rashḥī az
ṭamṭām-i yamm-i bayān-i ilāhī būdih. Ṣad-hizār amthāl-i īn nufūs
bi-qawlī khalq shudih va mīshavand va maʿa dhālik iftirāʾan ʿalā Allāhi
wa-kadhiban ʿalayhi wa-baghyan ʿalā ḥaḍratihi wa-ṭughyānan ʿalā bahāʾihi
wa-ithbātan li-aʿdāʾihi mīnivīsī kih Ḥaḍrat-i Aʿlā farmūdih kih iḥtiyāj
bi-naṣṣ-i man nadārad wa-bi-hādhihi al-kalimati tabarraʾa minka kullu mā
ẓahara min al-qalami fī mamāliki Allāhi al-muqtadiri al-ʿaliyyi
al-ʿalīm.

Everything they have written, if acknowledged and accepted, has been but
a spray from the surging ocean of divine utterance. A hundred thousand
souls like these have been, and are being, created by a word. Yet,
fabricating a falsehood against God, lying about Him, transgressing
against His presence, rebelling against His glory, and vindicating His
enemies, you write that His Holiness the Most Exalted said that he has
no need of My explicit designation. Through these words, everything that
has issued from the Pen throughout the realms of God, the All-Powerful,
the Most High, the All-Knowing, has disowned you.

\subsection{141}\label{section-636}

\begin{otherlanguage}{persian}

اليوم جميع اسماء حسنى بملكوت خود راجع و لايفقه ذلك الّا كلّ موقن عالم بصير
و آنچه در بيان نازل شده چه از احكام و چه از اوصاف منزله در شأن حروفات و
مرايا احدى را نفع نمىبخشد مگر بتصديق اين ظهور بذلك تنفّس نفس الرّحمن فى
نفس البيان لو انتم من العارفين و اليوم چون معرضين قابل عرفان مثل اعلى كه
حجّة اللّه و برهانه بوده نبوده و نيستند و از ادراك لطايف كلمات الهيّه عاجز
مشاهده مى شوند لذا در اينمقام مثل ادنى ذكر ميشود كه شايد نفوس مقبله
بحديقه معانى باطنيّه وارد شوند

\end{otherlanguage}

Al-yawm jamīʿ-i asmāʾ-i ḥusnā bi-malakūt-i khud rājiʿ va-lā yafqahu
dhālika illā kullu mūqinin ʿālimin baṣīr. Va ānchih dar Bayān nāzil
shudih chih az aḥkām va chih az awṣāf-i munazzalih dar shaʾn-i ḥurūfāt
va marāyā aḥadī rā nafʿ namībakhshad magar bi-taṣdīq-i īn ẓuhūr.
Bi-dhālika tanaffasa nafasu al-Raḥmāni fī nafsi al-Bayāni law antum min
al-ʿārifīn. Va al-yawm chūn muʿriḍīn qābil-i ʿirfān-i mathal-i aʿlā kih
ḥujjatu Allāh va burhān-i ū būdih nabūdih va nīstand va az idrāk-i
laṭāʾif-i kalimāt-i ilāhīyih ʿājiz mushāhidih mīshavand lidhā dar īn
maqām mathal-i adnā dhikr mīshavad kih shāyad nufūs-i muqbilih
bi-ḥadīqih-yi maʿānī-yi bāṭinīyih vārid shavand.

Today all the most beautiful names have returned to their own kingdom,
and none understands this except those possessed of certitude,
knowledge, and insight. Whatever has been revealed in the Bayán, whether
laws or descriptions revealed concerning the Letters and the Mirrors,
benefits no one unless they acknowledge this Manifestation. Thus has the
breath of the All-Merciful breathed within the Bayán itself, if you are
among those who recognize. Today, since those who turn away have not
been, and are not, capable of recognizing the most exalted likeness,
which is God's proof and demonstration, and are seen to be incapable of
grasping the subtleties of the divine words, a lesser analogy is given
here, so that receptive souls may enter the garden of inner meanings.

\subsection{142}\label{section-637}

\begin{otherlanguage}{persian}

سلطانى نفسى را حاكم و مختار مينمايد در كلّ امور چنانچه ديده ميشود كه بعضى
از ولات كه از جانب سلطان در ديار حاكمند مختارند در كلّ امور مع‌ذلك آيا
سلطان قادر است بر عزل آنها يا نه اگر تكذيب نمائى كاذب بوده چه كه مشهود
است كه قادر است و هم‌چنين بعد از جلوس سلطان ثانى آيا آن فرمان و حكمى كه
حكّام مختار بوده‌اند از قبل سلطان قبل منوط باجازه حكم سلطان ثانى بوده و يا
نه شكّى نيست كه جميع منوط و معلّق بامر و اجازه سلطان بعد است اين در شأن
ولات و سلاطين ظاهريّه ذكر شد كه شايد بآن سلطان مقتدرى كه در كلّ حين در شأن
بديع بوده و خواهد بود لم‌يزل يفعل ما یشآء و يحكم ما يريد قدرت او را از او
نفى ننمائيد و يداللّه را مغلول نشمريد اقلّا قدرتش را بقدر سلاطين ظاهريّه
بدانيد

\end{otherlanguage}

Sulṭānī nafsī rā ḥākim va mukhtār mīnamāyad dar kull-i umūr chinānchih
dīdih mīshavad kih baʿḍī az vulāt kih az jānib-i sulṭān dar diyār
ḥākimand mukhtārand dar kull-i umūr. Maʿa dhālik āyā sulṭān qādir ast
bar ʿazl-i ānhā yā nah? Agar takdhīb namāʾī kādhib būdih chih kih
mashhūd ast kih qādir ast. Va hamchinīn baʿd az julūs-i sulṭān-i thānī
āyā ān farmān va ḥukmī kih ḥukkām mukhtār būdih-and az qibal-i sulṭān-i
qabl manūṭ bi-ijāzih-yi ḥukm-i sulṭān-i thānī būdih va yā nah? Shakkī
nīst kih jamīʿ manūṭ va muʿallaq bi-amr va ijāzih-yi sulṭān-i baʿd ast.
Īn dar shaʾn-i vulāt va salāṭīn-i ẓāhirīyih dhikr shud kih shāyad bi-ān
sulṭān-i muqtadirī kih dar kull-i ḥīn dar shaʾn-i badīʿ būdih va khvāhad
būd lam yazal yafʿalu mā yashāʾu wa-yaḥkumu mā yurīdu qudrat-i ū rā az ū
nafy nanamāʾīd va yadu Allāh rā maghlūl nashmurīd. Aqallan qudratash rā
bi-qadr-i salāṭīn-i ẓāhirīyih bidānīd.

A sovereign appoints someone as governor and grants him authority in all
affairs, just as some governors who rule regions on behalf of a
sovereign are seen to have authority in all affairs. Nevertheless, can
the sovereign dismiss them or not? If you deny it, you are a liar, for
it is evident that he can. Likewise, after a second sovereign ascends
the throne, are the decrees and authority under which the governors
exercised their powers on behalf of the former sovereign subject to the
sanction of the second sovereign's rule or not? There is no doubt that
all depend upon, and are subject to, the command and permission of the
succeeding sovereign. This has been mentioned concerning outward
governors and sovereigns so that you may not deny His power to that
all-powerful Sovereign who has ever been, and will ever be, manifest in
a new state at every moment, who has always done whatever He wills and
decreed whatever He desires, and so that you may not regard the hand of
God as fettered. At least consider His power equal to that of outward
sovereigns.

\subsection{143}\label{section-638}

\begin{otherlanguage}{persian}

واللّه اگر نفسى آن سلطان قدم را قادر بداند بامثال اين كلمات بر موجد اسماء
و صفات اعتراض نمينمايد اى بی‌‌نصيبان از كوثر معارف رحمن سلّمنا بقول شما
نفسى بقميص اسمى از اسماء بين ارض و سماء مزيّن شده آيا حقّ قادر نيست بر
انتزاع آن فويل لكم تاللّه يكذّبكم كلّ الوجود من الغيب و الشّهود هرگز حقّ جلّ
ذكره در حين اعطاء بنفسى مصلحت نمود كه در حين اخذ مصلحت نمايد قلم حيا
مينمايد از ذكر آنچه اين ایّام تلقآء وجه مليك علّام بآن مشغول شده‌ايد خافوا
عن اللّه ثمّ اتركوا ما عندكم و خذوا ما عند اللّه هذا خير لكم عمّا خلق بين
السّموات و الأرضين

\end{otherlanguage}

Va-llāh agar nafsī ān sulṭān-i qidam rā qādir bidānad bi-amthāl-i īn
kalimāt bar mūjid-i asmāʾ va ṣifāt iʿtirāḍ namīnamāyad. Ay bī-naṣībān az
kawthar-i maʿārif-i Raḥmān sallamnā bi-qawl-i shumā nafsī bi-qamīṣ-i
ismī az asmāʾ bayn-i arḍ va samāʾ muzayyan shudih āyā Ḥaqq qādir nīst
bar intizāʿ-i ān? Fa-waylun lakum ta-Allāhi yukadhdhibukum kullu
al-wujūdi min al-ghaybi wa-al-shuhūd. Hargiz Ḥaqq jalla dhikruhu dar
ḥīn-i iʿṭā bi-nafsī maṣlaḥat nimūd kih dar ḥīn-i akhdh maṣlaḥat namāyad?
Qalam ḥayā mīnamāyad az dhikr-i ānchih īn ayyām tilqāʾ-i vajh-i Malīk-i
ʿAllām bi-ān mashghūl shudih-īd. Khāfū ʿan Allāhi thumma utrukū mā
ʿindakum wa-khudhū mā ʿinda Allāhi hādhā khayrun lakum ʿammā khuliqa
bayna al-samāwāti wa-al-arḍīn.

By God, if anyone regarded that Sovereign of eternity as powerful, he
would not raise objections such as these against the Creator of names
and attributes. You who have no share of the Kawthar of the
All-Merciful's knowledge! Let us concede your assertion that someone has
been adorned with the robe of one of the names between earth and heaven:
is God not able to strip it away? Woe to you! By God, all existence,
unseen and visible, gives you the lie. Did God---exalted be His
mention---ever consult anyone when bestowing a gift, that He should
consult anyone when taking it away? The Pen is ashamed to mention what
you have busied yourselves with in these days before the face of the
all-knowing King. Fear God, then abandon what you possess and take what
is with God. This is better for you than all that has been created
between the heavens and the earths.

\subsection{144}\label{section-639}

\begin{otherlanguage}{persian}

و اينكه تقسيم نموده بوديد كه نقطه بيان روح ماسواه فداه بجهت اهل تقليد و
عوام بيان فرموده اين را و بجهت اولى الأبصار و الافئده بيان ديگر فرموده
هذا بغى منك على اللّه المهيمن القيّوم اين قول بسيار لغو و بی‌ معنى است قد
ظهر بنفسه الحقّ بآيات بيّنات و انّها لحجّة من عند اللّه على كلّ نفس ان انت من
العارفين حجّت الهى لم‌يزل واحد بوده و از براى كلّ همان آيات منزلات را حجّت
قرار فرموده‌اند و اگر در خلق آيه عرفان و ادراك آيات نميبود البتّه بآيات
ظاهر نميشدند فى ذلك لحكمة لأولى الأباب و لايعرفها الّا اولو الأبصار

\end{otherlanguage}

Va īnkih taqsīm nimūdih būdīd kih Nuqṭih-yi Bayān rūḥ-i mā-sivāh fidāh
bi-jahat-i ahl-i taqlīd va ʿavāmm bayān farmūdih īn rā va bi-jahat-i ulī
al-abṣār va al-afʾidih bayān-i dīgar farmūdih hādhā baghyun minka ʿalā
Allāhi al-muhaymini al-qayyūm. Īn qawl bisyār laghv va bī-maʿnī ast. Qad
ẓahara bi-nafsihi al-ḥaqqu bi-āyātin bayyinātin wa-innahā la-ḥujjatun
min ʿindi Allāhi ʿalā kulli nafsin in anta min al-ʿārifīn. Ḥujjat-i
ilāhī lam-yazal vāḥid būdih va az barā-yi kull hamān āyāt-i munazzalāt
rā ḥujjat qarār farmūdih-and va agar dar khalq āyih-yi ʿirfān va idrāk-i
āyāt namībūd albattih bi-āyāt ẓāhir namīshudand. Fī dhālika la-ḥikmatun
li-ulī {[}corrupted text{]} wa-lā yaʿrifuhā illā ulū al-abṣār.

And as for the distinction you made, saying that the Point of the Bayán,
may the spirit of all besides Him be sacrificed for Him, made this
statement for the imitators and the common people, and made another
statement for those endowed with sight and understanding: this is a
transgression on your part against God, the All-Protecting, the
Self-Subsisting. This assertion is utterly vain and meaningless. The
Truth Himself has appeared with clear verses, and these are indeed a
proof from God for every soul, if you are among those who recognize.
God's proof has always been one, and He has appointed those same
revealed verses as a proof for everyone. Had there been no sign within
creation of the capacity to recognize and understand the verses, He
would certainly not have appeared with verses. In this there is indeed
wisdom for those endowed with {[}corrupted original text{]}, and none
recognizes it except those endowed with sight.

\subsection{145}\label{section-640}

\begin{otherlanguage}{persian}

و ديگر معلوم نيست كه را اهل تقليد و عوام دانسته‌ايد و كه را اولو الافئده
و اولو الأبصار اتّخذتم لانفسكم امراً و به ضللتم و ظلمتم و لا تعرفون بدانكه
اولو الأبصار و غيره بتصديق و تكذيب معلوم ميشوند و ممتاز ميگردند هر
نفسيكه موفّق شد بقول بلى در ظهور امنع اقدس اعلى او عنداللّه از اخيار و
اولو الأبصار بوده و من دون آن از اشرار و فجّار و البتّه اولو الأبصار نزد
شما جهلائى هستند كه ادّعاى علم و حكمت نموده‌اند بشنو بيان نقطه اوّلیّه و
مظهر كليّه را كه فرموده‌اند گندم پاك‌كن ارض صاد لباس نقابت را پوشيد با
اينكه حرفى از علم نخوانده بود و شيخ محمّد حسن كه پانصد هزار بيت در عمل
باللّه نوشت حرفى ثمر نبخشيد كه در يوم قيامت او را از نار نجات دهد

\end{otherlanguage}

Va dīgar maʿlūm nīst kih rā ahl-i taqlīd va ʿavāmm dānistih-īd va kih rā
ulū al-afʾidih va ulū al-abṣār. Ittakhadhtum li-anfusikum amran wa-bihi
ḍalaltum wa-ẓalamtum wa-lā taʿrifūn. Bidānkih ulū al-abṣār va ghayrih
bi-taṣdīq va takdhīb maʿlūm mīshavand va mumtāz mīgardand. Har nafsīkih
muwaffaq shud bi-qawl-i balā dar ẓuhūr-i amnaʿ-i aqdas-i aʿlā ū ʿinda
Allāh az akhyār va ulū al-abṣār būdih va min dūn-i ān az ashrār va
fujjār. Va albattih ulū al-abṣār nazd-i shumā juhalāʾī hastand kih
iddiʿā-yi ʿilm va ḥikmat nimūdih-and. Bishnaw bayān-i Nuqṭih-yi
Avvalīyih va Maẓhar-i Kullīyih rā kih farmūdih-and gandum pāk-kun-i
arḍ-i Ṣād libās-i niqābat rā pūshīd bā īnkih ḥarfī az ʿilm nakhvāndih
būd va Shaykh Muḥammad Ḥasan kih pānṣad hizār bayt dar ʿamal bi-Allāh
nivisht ḥarfī thamar nabakhshīd kih dar yawm-i qiyāmat ū rā az nār najāt
dahad.

Moreover, it is unclear whom you have regarded as imitators and common
people, and whom as those endowed with understanding and sight. You have
devised a matter for yourselves and through it have gone astray and
acted unjustly, yet you do not know. Know that those endowed with sight
and those who are not are identified and distinguished by their
acknowledgment or denial. Every soul enabled to say ``Yes'' in this most
unassailable, most holy, most exalted Manifestation is, in God's sight,
among the righteous and those endowed with sight; all others are among
the wicked and the ungodly. Of course, those endowed with sight in your
estimation are ignorant people who have claimed knowledge and wisdom.
Listen to the utterance of the Primal Point and Universal Manifestation,
who said, ``The wheat-sifter of the land of Ṣád put on the robe of
leadership, although he had not studied a single letter of learning,
whereas Shaykh Muḥammad Ḥasan, who wrote five hundred thousand verses on
acting in accordance with God, gained not a single letter's benefit that
would save him from the fire on the Day of Resurrection.''

\subsection{146}\label{section-641}

\begin{otherlanguage}{persian}

و هم‌چنين در مقام ديگر منصوصاً در معنى اولو الأبصار ميفرمايند قوله جلّ
ثنائه اينست كه مشتبه ميگردد در نزد هر ظهورى حقّ صرف از دون او الّا نزد
اولى الأبصار كه ايشان حقایق اشیاء را كما هى مشاهده مينمايند و هميشه ناظر
بشمس حقيقت هستند و ظهور او و حدودات مرايا ايشانرا محتجب از ظاهر در آنها
نمينمايد اولئك هم المتّقون حقّا فی كتاب اللّه من قبل و من بعد اولئك هم
المهتدون انتهى

\end{otherlanguage}

Va ham-chunīn dar maqām-i dīgar manṣūṣan dar maʿnī-yi ulū al-abṣār
mīfarmāyand qawluhu jalla thanāʾuhu īnast kih mushtabih mīgardad dar
nazd-i har ẓuhūrī Ḥaqq-i ṣirf az dūn-i ū illā nazd-i ulī al-abṣār kih
īshān ḥaqāyiq-i ashyā rā kamā hiya mushāhidih mīnamāyand va hamīshih
nāẓir bi-shams-i ḥaqīqat hastand va ẓuhūr-i ū va ḥudūdāt-i marāyā
īshānrā muḥtajib az ẓāhir dar ānhā namīnamāyad ulāʾika hum al-muttaqūn
ḥaqqan fī kitābi Allāh min qabl wa-min baʿd ulāʾika hum al-muhtadūn
intahā.

Likewise, in another passage He explicitly states concerning the meaning
of those endowed with sight---His words, exalted be His praise: ``At
every Manifestation, the pure Truth becomes indistinguishable from all
else except to those endowed with sight, who behold the realities of
things as they are and always look toward the Sun of Truth and His
manifestation. The limitations of the mirrors do not veil them from the
One manifest in them. These are truly the God-fearing in the Book of
God, before and after; these are the rightly guided.'' End of quotation.

\subsection{147}\label{section-642}

\begin{otherlanguage}{persian}

اين بيان نقطه بيان روح ما سواه فداه حال ملاحظه نما تا اولو الأبصار را
بشناسى چنانچه مشاهده مينمائى كه ابدا حدودات مرايا آن نفوس مقدّسه را از
شاطى بحر احديّه منع ننموده و صدهزار امثال آن مرايا را نزد اشراق تجلّى شمس
حقيقت معدوم صرف مشاهده مينمايند آن نفوسند كه عنداللّه و مظاهر امر او
باولو الأبصار معروفند نه نفوسيكه على زعم خود مرآتى درست نموده و بآن از
حقّ منيع محتجب مانده‌اند فويل للّذين غفلوا عمّا نزل فى الكتاب و اخذوا كلمة
منه و فسّروها بما تهوى به انفسهم و اهوائهم و اعترضوا بها على اللّه المهيمن
القيّوم و ايكاش كه اين رؤساى بيان كه الآن موجودند علم ظاهره هم ميداشتند
لافو‌الّذى نفسى بيده جز كذب و افترا و اعراض از حقّ و جعل نزد اين نفوس نبوده
و نخواهد بود كذلك بيّن اللّه كلّ امر فى كتاب مبين

\end{otherlanguage}

Īn bayān-i Nuqṭiy-i Bayān rūḥ-i mā sivāh fidāh ḥāl mulāḥiẓih namā tā ulū
al-abṣār rā bishināsī chunānchih mushāhidih mīnamāʾī kih abadan
ḥudūdāt-i marāyā ān nufūs-i muqaddasih rā az shāṭī-yi baḥr-i aḥadīyyih
manʿ nanamūdih va ṣad-hizār amthāl-i ān marāyā rā nazd-i ishrāq-i
tajallī-yi shams-i ḥaqīqat maʿdūm-i ṣirf mushāhidih mīnamāyand ān
nufūsand kih ʿinda Allāh va maẓāhir-i amr-i ū bi-ulū al-abṣār maʿrūfand
nah nufūsīkih ʿalā zaʿm-i khud mirʾātī durust namūdih va bi-ān az Ḥaqq-i
manīʿ muḥtajib māndih-and fa-waylun li-lladhīna ghafalū ʿammā nazala fī
al-kitāb wa-akhadhū kalimatan minhu wa-fassarūhā bi-mā tahwā bihi
anfusuhum wa-ahwāʾuhum wa-iʿtaraḍū bihā ʿalā Allāh al-muhaymin al-qayyūm
va ay kāsh kih īn ruʾasā-yi Bayān kih al-ān mawjūdand ʿilm-i ẓāhirih ham
mīdāshtand lā fa-wa-lladhī nafsī bi-yadihi juz kidhb va iftirā va iʿrāḍ
az Ḥaqq va jaʿl nazd-i īn nufūs nabūdih va nakhwāhad būd kadhālika
bayyana Allāh kulla amrin fī kitābin mubīn.

This is the utterance of the Point of the Bayán, may the spirit of all
besides Him be sacrificed for Him. Now reflect, so that you may
recognize those endowed with sight. As you observe, the limitations of
the mirrors have never kept those sanctified souls from the shore of the
ocean of oneness, and they regard a hundred thousand such mirrors as
utter nothingness before the radiant revelation of the Sun of Truth.
These are the souls known to God and the Manifestations of His Cause as
those endowed with sight, not those who have fashioned a mirror
according to their own fancy and through it remained veiled from the
unassailable Truth. Woe, then, to those who have neglected what was sent
down in the Book, taken a word from it, interpreted it according to the
inclinations of their own selves and desires, and used it to object to
God, the All-Protecting, the Self-Subsisting. If only these present
leaders of the Bayán possessed even outward learning! No, by Him in
whose hand is my soul! Nothing but lies, slander, turning away from the
Truth, and fabrication has ever been found among these souls, nor will
it ever be. Thus has God made every matter clear in a lucid Book.

\subsection{148}\label{section-643}

\begin{otherlanguage}{persian}

و ديگر نوشته بوديد مراد آن است كه مبالغه در اين امر فرموده از براى
اشخاصيكه تميز آيات نميدهند و تصريح فرموده و امّا از براى اشخاصيكه تميز
آيات ميدهند فرموده كه او خود صاحب آيات است و در تمام نوشته‌جات شكر و حمد
پروردگار نموده كه وصىّ مرا بآيات ظاهر فرموده انتهى

\end{otherlanguage}

Va dīgar nivishtih būdīd murād ān ast kih mubālaghih dar īn amr farmūdih
az barā-yi ashkhāṣīkih tamīz-i āyāt namīdihand va taṣrīḥ farmūdih va
ammā az barā-yi ashkhāṣīkih tamīz-i āyāt mīdihand farmūdih kih ū khud
ṣāḥib-i āyāt ast va dar tamām-i nivishtih-jāt shukr va ḥamd-i
Parvardigār namūdih kih vaṣī-yi marā bi-āyāt ẓāhir farmūdih intahā.

You also wrote: ``The intention is that He emphasized this matter and
stated it explicitly for those who cannot discern the verses. But for
those who can discern the verses, He said that he himself is the
possessor of verses, and throughout His writings He has thanked and
praised the Lord, saying that He has made My successor manifest with
verses.'' End of quotation.

\subsection{149}\label{section-644}

\begin{otherlanguage}{persian}

اوّلا اكثرى از اين كلمات از حضرت اعلى روح ماسواه فداه نبوده و نيست بسيار
از اين اقوال جعل نموده و مينمايند شما بتفصيل مطّلع نيستيد و ثانيا آنكه
آيات مرايا در رتبه مرايا مذكور مثل شبح شمس در مرآت حال آن تجلّى در مرآت
در رتبه مرآت مشهود و مذكور و لكن نسبت بشمس سماء ذكرى نداشته و نخواهد
داشت چنانچه مشاهده ميشود كه اگر مراياى مالانهايه مقابل شوند بشمس در كلّ
شبح شمس ظاهر و بعد از انحراف كلّ فانى و شمس بتمام ضيا در مقرّ خود مشرق و
مضئ و نقطه اولى روح ما سواه فداه در اينمقام بيانات لايحصى فرموده‌اند
فطوبی‌ للنّاظرين ثمّ للعارفين هر ذى بصر و ذي شعورى ملاحظه نمايد از آياتش
رتبه و مقامش را ادراك مينمايد

\end{otherlanguage}

Avvalan aktharī az īn kalimāt az Ḥaḍrat-i Aʿlā rūḥ-i mā sivāh fidāh
nabūdih va nīst bisyār az īn aqvāl jaʿl namūdih va mīnamāyand shumā
bi-tafṣīl muṭṭaliʿ nīstīd va thāniyan ānkih āyāt-i marāyā dar rutbiy-i
marāyā madhkūr mithl-i shabaḥ-i shams dar mirʾāt ḥāl ān tajallī dar
mirʾāt dar rutbiy-i mirʾāt mashhūd va madhkūr va lākin nisbat bi-shams-i
samā dhikrī nadāshtih va nakhwāhad dāsht chunānchih mushāhidih mīshavad
kih agar marāyā-yi mā lā nihāyih muqābil shavand bi-shams dar kull
shabaḥ-i shams ẓāhir va baʿd az inḥirāf kull fānī va shams bi-tamām-i
ḍiyā dar maqarr-i khud mushriq va muḍīʾ va Nuqṭiy-i Ūlā rūḥ-i mā sivāh
fidāh dar īn maqām bayānāt-i lā yuḥṣā farmūdih-and fa-ṭūbā li-l-nāẓirīn
thumma li-l-ʿārifīn har dhī baṣar va dhī shuʿūrī mulāḥiẓih namāyad az
āyātash rutbih va maqāmash rā idrāk mīnamāyad.

First, most of these words have never come from His Holiness the Most
Exalted, may the spirit of all besides Him be sacrificed for Him. They
have fabricated many such statements and continue to do so; you are not
acquainted with the details. Second, the verses of the mirrors are
spoken of at the level of the mirrors, like the image of the sun in a
mirror. That reflection in the mirror is seen and spoken of at the level
of the mirror, but in relation to the sun in the sky it has never
merited mention, nor will it ever do so. As can be observed, if
infinitely many mirrors face the sun, its image appears in them all;
when they turn away, all those images vanish, while the sun shines
resplendent in its own station with all its radiance. The Primal Point,
may the spirit of all besides Him be sacrificed for Him, has made
countless statements on this subject. Blessed, then, are those who
behold, and then those who recognize. Anyone endowed with sight and
understanding who reflects will perceive his rank and station from his
verses.

\subsection{150}\label{section-645}

\begin{otherlanguage}{persian}

و ديگر آنكه در مقامى كلّ اشیآء آيات اللّه بوده و خواهد بود لو انتم فى آيات
اللّه فى انفس الخلق تنظرون و كم من آية فى السّموات و الأرض يمرّون عليها و
هم عنها معرضون و آيات الهى در كلّ شیء مشهود و موجود ولكن مادامى كه در ظلّ
حقّ مشاهده ميشوند از آيات علّيين مذكور و من دون آن از سجّين محسوب

\end{otherlanguage}

Va dīgar ānkih dar maqāmī kull-i ashyā āyāt-i Allāh būdih va khwāhad būd
law antum fī āyāti Allāh fī anfusi al-khalq tanẓurūn wa-kam min āyatin
fī al-samāwāti wa-l-arḍ yamurrūna ʿalayhā wa-hum ʿanhā muʿriḍūn va
āyāt-i ilāhī dar kull-i shayʾ mashhūd va mawjūd va lākin mā dāmī kih dar
ẓill-i Ḥaqq mushāhidih mīshavand az āyāt-i ʿillīyyīn madhkūr va min
dūn-i ān az sijjīn maḥsūb.

Moreover, in one sense all things have always been and will ever be
signs of God, if you would look upon the signs of God within created
beings. ``How many a sign in the heavens and the earth do they pass by,
turning away from it!'' The divine signs are visible and present in
everything. Yet so long as these things are seen beneath the shadow of
the Truth, they are spoken of as signs of the highest heaven; otherwise,
they are reckoned as belonging to the nethermost prison.

\subsection{151}\label{section-646}

\begin{otherlanguage}{persian}

قسم بحقّ كه از كلمات تو خجلت بردم چه كه بكلماتى تكلّم نموده كه هيچيك از
محتجبين قبل بآن تكلّم ننموده‌اند بحر اعظم مشهود و ملاحظه نموده و مينمائى
مع‌ذلك باثبات ماء غدير كه از رايحه و طعمش انسان متأذّى است مشغول شده شمس
باشراقها و انوارها ظاهر و مشهود و تو خود را از انوار و ضيائش محروم نموده
و ادلّه براى معدومى كه نديده و ادراك ننموده ذكر مينمائى واللّه كه حرفى از
بيان ادراك ننموده و بعوالم تبيان كه مقام ظهور نفحات الهيّه و فوحات قدس
صمدانيّه است قدم نگذارده بكلمه از ذوالكلمه ممنوع شده و بحجاب اشارات از
نفحات ربّ الأرباب‌ محروم گشته

\end{otherlanguage}

Qasam bi-Ḥaqq kih az kalimāt-i tu khijlat burdam chih kih bi-kalimātī
takallum namūdih kih hīch-yik az muḥtajibīn-i qabl bi-ān takallum
nanamūdih-and baḥr-i aʿẓam mashhūd va mulāḥiẓih namūdih va mīnamāʾī maʿa
dhālik bi-ithbāt-i māʾ-i ghadīr kih az rāyiḥih va ṭaʿmash insān
mutaʾadhdhī ast mashghūl shudih shams bi-ishrāqihā va anwārihā ẓāhir va
mashhūd va tu khud rā az anvār va ḍiyāʾash maḥrūm namūdih va adillih
barā-yi maʿdūmī kih nadīdih va idrāk nanamūdih dhikr mīnamāʾī va-llāh
kih ḥarfī az Bayān idrāk nanamūdih va bi-ʿavālim-i tibyān kih maqām-i
ẓuhūr-i nafaḥāt-i ilāhīyyih va fawaḥāt-i quds-i ṣamadānīyyih ast qadam
nagudhāshtih bi-kalimih az Dhu al-Kalimih mamnūʿ shudih va bi-ḥijāb-i
ishārāt az nafaḥāt-i Rabb al-Arbāb maḥrūm gashtih.

By the Truth, I felt ashamed of your words, for you have spoken words
that none of those veiled in the past have ever uttered. The Most Great
Ocean is visible, and you have beheld it and continue to behold it;
nevertheless, you have busied yourself with establishing the worth of
water in a pool whose smell and taste cause distress. The sun is
manifest and visible in its radiance and light, yet you have deprived
yourself of its light and splendor and adduce proofs for a nonexistent
thing that you have neither seen nor perceived. By God, you have not
understood a single letter of the Bayán, nor set foot in the realms of
exposition, where the divine breezes and the holy fragrances of the
Eternal are made manifest. A word has kept you from the Possessor of the
Word, and the veil of allusions has deprived you of the breezes of the
Lord of Lords.

\subsection{152}\label{section-647}

\begin{otherlanguage}{persian}

سبحان اللّه سلطان مقتدرى كه مع كمال ظهور و بروز مشركين را از مشاهده جمالش
منع نموده و مع كمال سموّ و علوّ مغلّين را از ادراك انوار وجهش بی‌نصيب گذاشته
بآيات اثبات بيّنه مينمايند و از سلطان آيات معرض كذلك كان مقتدراً على كلّ
شیء و لكنّ النّاس لا يبصرون بشنو نداى اين عبد فانى را و بعرصه روحانى وارد
شو نقطه اولى بنفسه ظاهر و انوار و تجلّياتش عالم غيب و شهاده را منوّر
فرموده بآن ناظر نيستى و بكلماتش كه ابداً عارف بحرفى از آن نبوده اثبات
حقّیّت از براى اعدايش مينمائى آفرين بر بصر شما كه يار را اغيار دانسته و
فجّار را اخيار شمرده

\end{otherlanguage}

Subḥān Allāh Sulṭān-i muqtadirī kih maʿa kamāl-i ẓuhūr va burūz
mushrikīn rā az mushāhidiy-i jamālash manʿ namūdih va maʿa kamāl-i
sumuvv va ʿuluvv mughillīn rā az idrāk-i anvār-i vajhash bī-naṣīb
gudhāshtih bi-āyāt ithbāt-i bayyinih mīnamāyand va az Sulṭān-i āyāt
muʿriḍ kadhālika kāna muqtadiran ʿalā kulli shayʾ wa-lākinna al-nāsa lā
yubṣirūn bishinaw nidā-yi īn ʿabd-i fānī rā va bi-ʿarṣiy-i rūḥānī vārid
shaw Nuqṭiy-i Ūlā bi-nafsih ẓāhir va anvār va tajallīyātash ʿālam-i
ghayb va shahādih rā munavvar farmūdih bi-ān nāẓir nīstī va
bi-kalimātash kih abadan ʿārif bi-ḥarfī az ān nabūdih ithbāt-i ḥaqqīyyat
az barā-yi aʿdāyash mīnamāʾī āfarīn bar baṣar-i shumā kih yār rā aghyār
dānistih va fujjār rā akhyār shumurdih.

Glory be to God! What a mighty Sovereign, who, despite His complete
manifestation and appearance, has prevented the idolaters from beholding
His beauty, and despite His consummate loftiness and exaltation, has
left the malicious without a share in perceiving the lights of His face.
They use verses to establish a proof, yet turn away from the Sovereign
of verses. Thus has He ever been powerful over all things, but the
people do not see. Hear the call of this evanescent servant and enter
the spiritual realm. The Primal Point Himself is manifest, and His
lights and revelations have illumined the worlds of the unseen and the
seen. You do not look toward Him, but use His words, of which you have
never understood a single letter, to establish the legitimacy of His
enemies. Bravo for your sight, which has taken the Friend for a stranger
and reckoned the wicked among the good!

\subsection{153}\label{section-648}

\begin{otherlanguage}{persian}

فواللّه اين عبد متحيّر است بلكه عالمى متحيّر و مبهوت متّصل مينويسيد كه نقطه
اولى چنين و چنان فرموده و نفس نقطه را در كلّ حين باسياف بغضا و رماح غلّ و
فحشا مجروح مينمائى تضربون و لا تشعرون مثل شما مثل نفوسى است كه در ليالى
و ایّام منتظر حقّ بوده‌اند و بعد از ظهور بكلمات قبلش استدلال بر كفر و فسادش
نموده تا آنكه بالأخره‌ بتمام ظلم سفك دم مطهّرش نموده‌اند الا لعنة اللّه على
الّذين ظلموا و يظلمون حينئذ و لا یشعرون متوهّم شنيده بوديم و لكن واللّه
باين شأن كه تو ظاهر شده نشنيده و نديده بوديم نسئل اللّه بأن يطهّرنا و ایّاک
و يجمعنا على شريعة حبّه و رضائه و لايجعلنا محروماً عن انوار وجهه المشرق
المنير

\end{otherlanguage}

Fa-va-llāh īn ʿabd mutaḥayyir ast balkih ʿālamī mutaḥayyir va mabhūt
muttaṣil mīnivīsīd kih Nuqṭiy-i Ūlā chunīn va chunān farmūdih va nafs-i
Nuqṭih rā dar kull-i ḥīn bi-asyāf-i baghḍā va rimāḥ-i ghill va faḥshā
majrūḥ mīnamāʾī taḍribūna wa-lā tashʿurūn mathal-i shumā mathal-i nufūsī
ast kih dar layālī va ayyām muntaẓir-i Ḥaqq būdih-and va baʿd az ẓuhūr
bi-kalimāt-i qablash istidlāl bar kufr va fasādash namūdih tā ānkih
bi-l-ākharah bi-tamām-i ẓulm safk-i dam-i muṭahharash namūdih-and alā
laʿnatu Allāh ʿalā alladhīna ẓalamū wa-yaẓlimūna ḥīnaʾidhin wa-lā
yashʿurūn mutavahhim shinīdih būdīm va lākin va-llāh bi-īn shaʾn kih tu
ẓāhir shudih nashinīdih va nadīdih būdīm nasʾalu Allāh bi-an yuṭahhiranā
wa-iyyāka wa-yajmaʿanā ʿalā sharīʿati ḥubbihi wa-riḍāʾihi wa-lā
yajʿalanā maḥrūman ʿan anwāri wajhihi al-mushriq al-munīr.

By God, this servant is bewildered---indeed, a whole world is bewildered
and astonished. You continually write that the Primal Point said this
and that, while at every moment you wound the Point Himself with the
swords of hatred and the spears of malice and abuse. You strike and do
not realize it. You are like those souls who awaited the Truth by night
and day, and after His appearance used His former words to argue for His
unbelief and corruption, until at last, in utter injustice, they shed
His sacred blood. Truly, God's curse rests upon those who committed
injustice and commit it now, yet do not realize it. We had heard of
deluded people, but, by God, we had neither heard nor seen anything like
the state you have displayed. We ask God to purify us and you, to gather
us upon the path of His love and good pleasure, and not to deprive us of
the lights of His radiant, luminous face.

\subsection{154}\label{section-649}

\begin{otherlanguage}{persian}

وديگر نوشته بوديد كه صريح فرموده كه آيات او از من است و من او هستم و او
من ميباشد و او را ثمره خود خوانده حال انصاف بدهيد كه خود حضرت باب آيات
را بهتر تميز ميداد يا بنده و شما و اين تقليد نيست و اگر بفرمائيد تقليد
است خود انصاف بدهيد كه از اوّل امر الى الآن از ميرزا يحيى چه قدرها آيات
بظهور رسيده و ميرسد انتهى

\end{otherlanguage}

Va dīgar nivishtih būdīd kih ṣarīḥ farmūdih kih āyāt-i ū az man ast va
man ū hastam va ū man mībāshad va ū rā th marih-yi khud khwāndih ḥāl
inṣāf bidihīd kih khud-i Ḥaḍrat-i Bāb āyāt rā bihtar tamīz mīdād yā
bandih va shumā va īn taqlīd nīst va agar bifarmāʾīd taqlīd ast khud
inṣāf bidihīd kih az avval-i amr ilā al-ān az Mīrzā Yaḥyā chih qadr-hā
āyāt bi-ẓuhūr rasīdih va mīrasad intahā.

You also wrote: ``He explicitly stated that his verses are from Me, and
I am he and he is Me, and He called him His own fruit. Now judge fairly:
was the Báb Himself better able to discern the verses, or are you and I?
This is not imitation. And if you say it is imitation, judge fairly for
yourself how many verses have appeared and continue to appear from Mírzá
Yaḥyá from the beginning of the Cause until now.'' End of quotation.

\subsection{155}\label{section-650}

\begin{otherlanguage}{persian}

بر فرض تصديق اين فقره كلّ اشیآء را بخود نسبت داده‌اند و بكلّ اسماء خود را
ناميده‌اند چه كه در كلّ آيه تجلّى خود را مشاهده فرموده‌اند مع تقديس ذاته عن
كلّ صفة و تنزیه كينونته عن كلّ اسم چنانچه در مقامى ميفرمايند قوله عزّ ذكره
قل تنزّلت حتّى قلت انّنى انا ذرّ و دون ذرّ مثل ما انّنى انا قلت انا ربّ و مربّب
كلّ ذى ربّ الى الآخر عجب است از شما كه پى باين مطلب واضح نبرده اگر بصير
باشى موقن ميشوى كه از اوّل ظهور تا حين ارقاء غير خود را مشاهده نفرموده‌اند
بشأنيكه چشم مبارك جز بر طلعت مبارك نيفتاده و اين يك مقام از مقامات آن
طلعت قدم بوده و ما اطّلع به احد الّا نفسه العليم الحكيم

\end{otherlanguage}

Bar farḍ-i taṣdīq-i īn faqrih kull-i ashyā rā bi-khud nisbat dādih-and
va bi-kull-i asmā khud rā nāmīdih-and chih kih dar kull āyiy-i
tajallī-yi khud rā mushāhidih farmūdih-and maʿa taqdīsi dhātihi ʿan
kulli ṣifah wa-tanzīhi kaynūnatihi ʿan kulli ism chunānchih dar maqāmī
mīfarmāyand qawluhu ʿazza dhikruhu qul tanazzaltu ḥattā qultu innanī anā
dharr wa-dūna dharr mithla mā innanī anā qultu anā rabb wa-murabbib
kulli dhī rabb ilā al-ākhir ʿajab ast az shumā kih pay bi-īn maṭlab-i
vāḍiḥ naburdih agar baṣīr bāshī mūqin mīshavī kih az avval-i ẓuhūr tā
ḥīn-i irqāʾ ghayr-i khud rā mushāhidih nafarmūdih-and bi-shaʾnīkih
chashm-i mubārak juz bar ṭalʿat-i mubārak nayuftādih va īn yik maqām az
maqāmāt-i ān Ṭalʿat-i Qidam būdih wa-mā iṭṭalaʿa bihi aḥadun illā
nafsuhu al-ʿalīm al-ḥakīm.

Even assuming this statement to be authentic, He has attributed all
things to Himself and called Himself by every name, for in everything He
has beheld a sign of His own revelation, while His essence remains
sanctified above every attribute and His being exalted above every name.
Thus, in one passage He says---His words, exalted be His mention: ``Say:
I have descended until I said, `I am indeed an atom, and less than an
atom,' just as I have said, `I am a Lord and the nurturer of every
possessor of lordship,'\,'' and so forth. It is astonishing that you
have not grasped this clear matter. If you possess insight, you will
become certain that from the beginning of His manifestation until the
time of His ascension He beheld none but Himself, such that His blessed
eye fell on nothing but His blessed countenance. This was one of the
stations of that Countenance of Eternity, and no one has perceived it
except He Himself, the All-Knowing, the All-Wise.

\subsection{156}\label{section-651}

\begin{otherlanguage}{persian}

و فوق اين مقام مقامات لانهايه عنداللّه بوده لاينبغى اظهارها لم يزل اناء
علم الهى بخاتم حفظ مختوم بوده و خواهد بود كلّ اشیاء از آنچه منسوب بحقّند
بندآء انّنى انا اللّه لا اله الّا انا ناطق و اين بسى واضح است كه اين كلمه و
اين رتبه از نفس آن شیء نبوده بلكه نظر بآن است كه در ظل حقّ ساكن گشته و از
حقّ محسوب شده و بعد از قطع آن نسبت همان شیء از اهل سجّين عنداللّه مذكور
چنانچه نقطه بيان روح ما سواه فداه فرموده قوله عزّ ذكره اين حجره كه بلا
ابواب است و بلا حدود معروفه امروز اعلى غرف رضوانست كه شجره حقيقت در او
مستقرّ است كه گويا ذرّات آن همه بندآء ذكر انّنى انا اللّه لا اله الّا انا ربّ
كلّ شیء ميخوانند بر كلّ حجرات اگرچه مقاعد مرءاتيّه مطرزه بذهب باشد ولى اگر
شجره حقيقت در مثل آن حجرات ساكن گردد آن وقت آن ذرّات مرايا ميخوانند
چنانچه ذرّات مراياى اعلى مقاعد عمارت صدرى كه در ایّام ارض صاد در آنجا ساكن
بوده ميخوند و ميخوانند و حكم از براى هيچ شیء نيست الّا آنكه راجع بآن نفس
ميگردد انتهى

\end{otherlanguage}

Va fawq-i īn maqām maqāmāt-i lā nihāyih ʿinda Allāh būdih lā yanbaghī
iẓhāruhā lam yazal ināʾ-i ʿilm-i ilāhī bi-khātam-i ḥifẓ makhtūm būdih va
khwāhad būd kull-i ashyā az ānchih mansūb bi-Ḥaqqand bi-nidā-yi innanī
anā Allāh lā ilāha illā anā nāṭiq va īn basī vāḍiḥ ast kih īn kalimih va
īn rutbih az nafs-i ān shayʾ nabūdih balkih naẓar bi-ān ast kih dar
ẓill-i Ḥaqq sākin gashtih va az Ḥaqq maḥsūb shudih va baʿd az qaṭʿ-i ān
nisbat hamān shayʾ az ahl-i sijjīn ʿinda Allāh madhkūr chunānchih
Nuqṭiy-i Bayān rūḥ-i mā sivāh fidāh farmūdih qawluhu ʿazza dhikruhu īn
ḥujrih kih bi-lā abvāb ast va bi-lā ḥudūd-i maʿrūfih imrūz aʿlā ghuraf-i
Riḍvān ast kih Shajariy-i Ḥaqīqat dar ū mustaqarr ast kih gūyā dharrāt-i
ān hamih bi-nidā-yi dhikr-i innanī anā Allāh lā ilāha illā anā rabbu
kulli shayʾ mīkhwānand bar kull-i ḥujurāt agarchih maqāʿid-i mirʾātīyyih
muṭarrazih bi-dhahab bāshad valī agar Shajariy-i Ḥaqīqat dar mithl-i ān
ḥujurāt sākin gardad ān vaqt ān dharrāt-i marāyā mīkhwānand chunānchih
dharrāt-i marāyā-yi aʿlā maqāʿid-i ʿimārat-i Ṣadrī kih dar ayyām-i arḍ-i
Ṣād dar ānjā sākin būdih mīkhund va mīkhwānand va ḥukm az barā-yi hīch
shayʾ nīst illā ānkih rājiʿ bi-ān nafs mīg ar dad intahā.

Above this station there have ever been limitless stations with God that
should not be disclosed. The vessel of divine knowledge has always been,
and will ever remain, sealed with the seal of protection. All things,
insofar as they are attributed to the Truth, proclaim the call, ``I am
indeed God; there is no god but Me.'' It is abundantly clear that this
utterance and this rank do not arise from the thing itself, but are due
to its having come to dwell beneath the shadow of the Truth and to be
reckoned as belonging to Him. Once that relationship is severed, that
same thing is spoken of in God's sight as belonging to the people of the
nethermost prison. Thus the Point of the Bayán, may the spirit of all
besides Him be sacrificed for Him, has said---His words, exalted be His
mention: ``This chamber, which is without doors and without customary
boundaries, is today the loftiest of the chambers of Paradise, for the
Tree of Truth is established within it. It is as though all its
particles proclaim over every other chamber, `I am indeed God; there is
no god but Me, the Lord of all things,' even though those chambers be
mirrored seats adorned with gold. But if the Tree of Truth should dwell
in such chambers, then those particles of the mirrors would proclaim it,
just as the particles of the mirrors of the loftiest seats of the Ṣadrí
mansion, where He dwelt during His days in the land of Ṣád, proclaimed
it and still proclaim it. No judgment applies to anything except insofar
as it refers back to that Person.'' End of quotation.

\subsection{157}\label{section-652}

\begin{otherlanguage}{persian}

اين شهادت كه خود آن جمال مبارك در ذرّات تراب داده قدرى تفكّر كن شايد
بكلماتى كه صدهزار رتبه در تحت وصف تراب افتاده از نفس حقّ محتجب نمانى
اگرچه قسم بآفتاب عزّ تقديس كه محتجبين اين طايفه احجب از كلّ مشاهده ميشوند
بشأنيكه نزديك بآن رسيده كه بروايت خبيثى از صدهزار قدرت مشهوده الهى و
آيات منزله صمدانى محجوب مانند بلكه باعراض و اعتراض برخيزند و اينست از
وعدههائيكه خود صاحب بيان روح ماسواه فداه فرموده و هم‌چنين اين ظهور امنع
اقدس در اكثر الواح قبل ما يظهر من بعد را خبر داده‌اند

\end{otherlanguage}

Īn shahādat kih khud ān Jamāl-i Mubārak dar dharrāt-i turāb dādih qadrī
tafakkur kun shāyad bi-kalimātī kih ṣad-hizār rutbih dar taḥt-i vaṣf-i
turāb uftādih az nafs-i Ḥaqq muḥtajib namānī agarchih qasam bi-āftāb-i
ʿizz-i taqdīs kih muḥtajibīn-i īn ṭāʾifih aḥjab az kull mushāhidih
mīshavand bi-shaʾnīkih nazdīk bi-ān rasīdih kih bi-rivāyat-i khabīthī az
ṣad-hizār qudrat-i mashhūdiy-i ilāhī va āyāt-i munazzaliy-i ṣamadānī
maḥjūb mānand balkih bi-iʿrāḍ va iʿtirāḍ bar khīzand va īnast az
vaʿdih-hāʾīkih khud Ṣāḥib-i Bayān rūḥ-i mā sivāh fidāh farmūdih va
ham-chunīn īn Ẓuhūr-i Amnaʿ-i Aqdas dar akthar-i alvāḥ-i qabl mā yaẓharu
min baʿd rā khabar dādih-and.

Reflect a little on this testimony that the Blessed Beauty Himself has
given concerning the particles of dust, so that you may not remain
veiled from the Truth Himself by words that stand a hundred thousand
degrees below the description of dust. Yet, by the sun of the glory of
sanctity, the veiled among this people are seen to be more deeply veiled
than all others, to such a degree that they have nearly reached the
point where a malicious report would veil them from a hundred thousand
visible manifestations of divine power and divinely revealed
verses---indeed, would cause them to rise up in rejection and
opposition. This is among the promises made by the Author of the Bayán
Himself, may the spirit of all besides Him be sacrificed for Him.
Likewise, this most unassailable, most holy Manifestation has foretold
in most of His earlier Tablets what will appear hereafter.

\subsection{158}\label{section-653}

\begin{otherlanguage}{persian}

مع آنكه درباره تراب تصديق انّه لا اله الّا هو فرموده‌اند ميفرمايند حكم از
براى هيچ شیء نيست الّا آنكه راجع بآن نفس ميگردد يعنى نفس حقّ چنانچه اليوم
جميع اوصاف بيان بجبروت صفات راجع و هم‌چنين جميع اسماء حسنى كه در او نازل
بملكوت اسماء و مقصود از كلّ اين اذكار ظهور آخر در كرّه اخرى بوده چه كه او
محبوب جان و مقصود روانش بوده چنانچه هر نفسى در بيان ناظر شود تصريحاً من
غير تلويح مشاهده نمايد و موقن شود باينكه مقصود از كلّ بيان آن جوهر تبيان
بوده و خواهد بود

\end{otherlanguage}

Maʿ ānkih darbārih-yi turāb taṣdīq-i annahu lā ilāha illā huwa
farmūdih-and, mīfarmāyand ḥukm az barā-yi hīch shayʾ nīst illā ānkih
rājiʿ bi-ān nafs mīgardad, yaʿnī nafs-i Ḥaqq. Chunānchih al-yawm jamīʿ-i
awṣāf-i Bayān bi-jabarūt-i ṣifāt rājiʿ va hamchunīn jamīʿ-i asmā-yi
ḥusnā kih dar ū nāzil bi-malakūt-i asmā. Va maqṣūd az kull-i īn adhkār
ẓuhūr-i ākhar dar karrih-yi ukhrā būdih, chih kih ū maḥbūb-i jān va
maqṣūd-i ravānash būdih. Chunānchih har nafsī dar Bayān nāẓir shavad,
taṣrīḥan min ghayri talwīḥ mushāhadih namāyad va mūqin shavad bi-īnkih
maqṣūd az kull-i Bayān ān jawhar-i tibyān būdih va khwāhad būd.

Although He has affirmed concerning the dust, ``There is no god but
Him,'' He says that no judgment applies to anything except insofar as it
refers back to that Person, meaning the Truth Himself. Thus, today, all
the attributes in the Bayán refer to the dominion of attributes, and
likewise all the most beautiful names revealed in it refer to the
kingdom of names. The purpose of all these references was the subsequent
Manifestation in another cycle, for He was the beloved of His soul and
the desire of His spirit. Anyone who examines the Bayán will see this
stated explicitly, not merely implied, and will become certain that the
purpose of the entire Bayán has been, and will remain, that Essence of
exposition.

\subsection{159}\label{section-654}

\begin{otherlanguage}{persian}

مثلاً ببعضى الواح فرستاده‌اند در صدر آن نازل هذا كتاب من عنداللّه المهيمن
القيّوم الى من يظهره اللّه انّه لا اله الّا انا العزيز المحبوب و حال آنكه آن
لوح بحسب ظاهر بجهت نفس ديگر بوده و هم‌چنين در كتاب بعضى نازل هذا كتاب من
اللّه المهيمن القيّوم الى اللّه المهيمن القيّوم اينست كه ميفرمايد آنچه در
بيان نازل شده در حقيقت اوّليه مقصود ظهور آخر است و در مادون آن نظر بآن
نسبتى است كه بشجره حقيقت داشته‌اند و هم‌چنين مخصوص بعضى الواح ارسال و در
صدر آن هو انت و انت هو نازل از جمله جناب ملّا محمّد على زنجانى و ديگر در
ذكر تكميل مراتب توحيد در اعداد كلّ شیء ميفرمايند لا اله الّا هو و لا اله
الّا انا و لا اله الّا اللّه و لا اله الّا انت و لا اله الّا الّذى و لا اله
الّا ايّاى ربّ العالمين جميعاً و لا اله الّا ایّاک ربّ العالمين جميعاً

\end{otherlanguage}

Mathalan bi-baʿḍī alvāḥ firistādih-and, dar ṣadr-i ān nāzil: hādhā
kitābun min ʿindi Allāhi al-muhaymini al-qayyūmi ilā man yuẓhiruhu
Allāh, innahu lā ilāha illā anā al-ʿazīzu al-maḥbūb. Va ḥāl ānkih ān
lawḥ bi-ḥasab-i ẓāhir bi-jihat-i nafs-i dīgar būdih. Va hamchunīn dar
kitāb-i baʿḍī nāzil: hādhā kitābun min Allāhi al-muhaymini al-qayyūmi
ilā Allāhi al-muhaymini al-qayyūm. Īnast kih mīfarmāyad ānchih dar Bayān
nāzil shudih dar ḥaqīqat-i avvalīyih maqṣūd ẓuhūr-i ākhar ast va dar
mādūn-i ān naẓar bi-ān nisbatī ast kih bi-shajari h-yi ḥaqīqat
dāshtih-and. Va hamchunīn makhṣūṣ-i baʿḍī alvāḥ irsāl va dar ṣadr-i ān
huwa anta va anta huwa nāzil, az jumlih Janāb-i Mullā Muḥammad ʿAlī-yi
Zanjānī. Va dīgar dar dhikr-i takmīl-i marātib-i tawḥīd dar aʿdād-i
kull-i shayʾ mīfarmāyand: lā ilāha illā huwa va lā ilāha illā anā va lā
ilāha illā Allāh va lā ilāha illā anta va lā ilāha illā alladhī va lā
ilāha illā iyyāya rabba al-ʿālamīna jamīʿan va lā ilāha illā iyyāka
rabba al-ʿālamīna jamīʿan.

For example, He sent Tablets to certain people, at the beginning of
which was revealed: ``This is a Book from God, the All-Protecting, the
Self-Subsisting, to Him whom God will make manifest. There is no god but
Me, the Almighty, the Beloved.'' Yet, outwardly, that Tablet was
intended for another person. Likewise, in the letter to another was
revealed: ``This is a Book from God, the All-Protecting, the
Self-Subsisting, to God, the All-Protecting, the Self-Subsisting.'' This
is why He says that, in its primary reality, what has been revealed in
the Bayán is intended for the subsequent Manifestation; at a lesser
level, it concerns the relationship people have had to the Tree of
Truth. Likewise, He sent Tablets specifically to certain people, at the
beginning of which was revealed, ``He is you, and you are He,'' among
them the honored Mullá Muḥammad ʿAlí of Zanján. Furthermore, in speaking
of the completion of the degrees of divine unity in the numbers of all
things, He says: ``There is no god but Him, and there is no god but Me,
and there is no god but God, and there is no god but You, and there is
no god but He who is, and there is no god but Me, the Lord of all the
worlds, and there is no god but You, the Lord of all the worlds.''

\subsection{160}\label{section-655}

\begin{otherlanguage}{persian}

اميدواريم كه از اين بيانات مالك اسماء و صفات قدرى از حدودات نفسيّه و
حجبات و اشارات انفس مردوده طاهر شويد كه اقلّا اگر در يوم ظهور اعانت حقّ
نمينمائيد بر نفسش سيف نكشيد و فتوى بر قتلش ندهيد بارى در هر شیء از اشیآء
كه منسوب الى اللّه است چه از اعلى علوّ و چه از ادنى دنوّ و اين علوّ و دنوّ
نظر بخلق ذكر ميشود بر جميع كلّ اسماء و كلّ اوصاف اطلاق شده و ميشود
ماداميكه در ظلّ حقّ مستقرّ و الّا حكم وجود بر آنها نميشود تا چه رسد بمراتب
عاليه نسئل اللّه بأن يغسّلك من كوثر الّذى جرى من اصبع الفضل و ينزع عن هيكلك
ثوب الاشارات و يخلعك بطراز اسمه العلى العظيم

\end{otherlanguage}

Umīdvārīm kih az īn bayānāt-i Mālik-i asmā va ṣifāt qadrī az ḥudūdāt-i
nafsīyih va ḥujubāt va ishārāt-i anfus-i mardūdih ṭāhir shavīd kih
aqallan agar dar yawm-i ẓuhūr iʿānat-i Ḥaqq namīnamāyīd, bar nafsash
sayf nakashīd va fatvā bar qatlash nadahīd. Bārī dar har shayʾ az ashyāʾ
kih mansūb ilā Allāh ast, chih az aʿlā-yi ʿuluvv va chih az adnā-yi
dunuvv, va īn ʿuluvv va dunuvv naẓar bi-khalq dhikr mīshavad, bar jamīʿ
kull-i asmā va kull-i awṣāf iṭlāq shudih va mīshavad mādāmīkih dar
ẓill-i Ḥaqq mustaqarr, va illā ḥukm-i vujūd bar ānhā namīshavad tā chih
rasad bi-marātib-i ʿālīyih. Nasʾalu Allāha bi-an yughassilaka min
kawthari alladhī jarā min iṣbaʿi al-faḍli va yanziʿa ʿan haykalika
thawba al-ishārāti va yakhlaʿaka bi-ṭirāzi ismihi al-ʿaliyyi al-ʿaẓīm.

We hope that through these utterances of the Sovereign of names and
attributes you may become somewhat purified of the limitations of self
and of the veils and allusions of rejected souls, so that, at the very
least, if you do not aid the Truth on the Day of Manifestation, you will
not draw a sword against Him or issue a decree for His death. In short,
all names and all attributes have been, and are, applied to every thing
that is attributed to God, whether at the highest height or the lowest
depth---and this height and depth are spoken of with reference to
creation---so long as it remains established beneath the shadow of the
Truth. Otherwise, it is not even accounted as existing, much less as
possessing exalted stations. We ask God to wash you with the heavenly
water that has flowed from the finger of grace, to strip from your
person the garment of allusions, and to invest you with the adornment of
His name, the Most High, the Most Great.

\subsection{161}\label{section-656}

\begin{otherlanguage}{persian}

و اينكه نوشته او را ثمره خود خوانده بگو بآن معرضين كه اى غافلان از معانى
كلمات الهى و اى بی‌بهره‌گان از ثمره شجره ربّانى اوّل مذهب بيان را ادراك
نمائيد كه چيست و مقصود نقطه بيان جلّت عظمته از كلمات و آياتش چه بعد
بتحقيق مشغول شويد و مباشيد از نفوسيكه كلمات الهيّه را بهواى نفسانيه تأويل
نموده‌اند و معانى آن را تحريف کرده‌اند

\end{otherlanguage}

Va īnkih nivishtih ū rā thamarih-yi khud khwāndih, bigū bi-ān muʿriḍīn
kih ay ghāfilān az maʿānī-yi kalimāt-i ilāhī va ay bī-bahrigān az
thamarih-yi shajari h-yi rabbānī, avval madhhab-i Bayān rā idrāk namāyīd
kih chīst va maqṣūd-i Nuqṭih-yi Bayān, jallat ʿaẓamatuhu, az kalimāt va
āyātash chih, baʿd bi-taḥqīq mashghūl shavīd va mabāshīd az nufūsīkih
kalimāt-i ilāhīyih rā bi-havā-yi nafsānīyih taʾvīl namūdih-and va
maʿānī-yi ān rā taḥrīf kardih-and.

As for your statement that He called him His own fruit, say to those who
turn away: ``You who are heedless of the meanings of the divine words
and deprived of the fruit of the divine Tree! First understand what the
teaching of the Bayán is, and what the Point of the Bayán---exalted be
His grandeur---intends by His words and verses; then undertake your
investigation. Do not be among those who have interpreted the divine
words according to selfish desires and distorted their meanings.''

\subsection{162}\label{section-657}

\begin{otherlanguage}{persian}

مذهب نقطه بيان اين است كه ذكر ميشود اصل ظهور را شجره ميفرمايند و بعد
مظاهر آن ظهور كلّ را از اغصان و افنان و اوراق و اثمار شجره ظهور چنانچه
مشاهده ميشود مثلاً خود نقطه اولى روح ماسواه فداه شجره ظهور الهيّه بوده و
متديّنين بدين او جميع از اوراق و اغصان و افنان و اثمار او محسوب ميشوند و
هم‌چنين در جميع ظهورات ملاحظه كن و بهر اسم و رسم كه بخواهى ذكر نما مثلاً
نقطه را بحر و سايرين را امواج و هم‌‌چنين شمس و سايرين را مرايا در اينمقام
اگر هر يك از مظاهر آن ظهور را بفرمايد كه او من است و من او حقّ بوده و لا
ريب فيه چنانچه در حروفات حىّ كه بعضى از آن نفوس موجودند و هم‌‌چنين در اعداد
لانهايه كه منتهى بواحد بيان ميشود ميفرمايد كه در كلّ ديده نميشود مگر جمال
آن ظهور

\end{otherlanguage}

Madhhab-i Nuqṭih-yi Bayān īn ast kih dhikr mīshavad: aṣl-i ẓuhūr rā
shajarih mīfarmāyand va baʿd maẓāhir-i ān ẓuhūr, kull rā az aghṣān va
afnān va awrāq va athmār-i shajari h-yi ẓuhūr, chunānchih mushāhadih
mīshavad. Mathalan khud-i Nuqṭih-yi Ūlā, rūḥ-i mā-sivāhu fidāh, shajari
h-yi ẓuhūr-i ilāhīyih būdih va mutadayyinīn bi-dīn-i ū jamīʿ az awrāq va
aghṣān va afnān va athmār-i ū maḥsūb mīshavand. Va hamchunīn dar jamīʿ-i
ẓuhūrāt mulāḥiẓih kun va bi-har ism va rasm kih bikhwāhī dhikr namā,
mathalan Nuqṭih rā baḥr va sāyirīn rā amvāj va hamchunīn shams va
sāyirīn rā marāyā. Dar īn maqām agar har yak az maẓāhir-i ān ẓuhūr rā
bifarmāyad kih ū man ast va man ū, ḥaqq būdih va lā rayba fīh.
Chunānchih dar Ḥurūfāt-i Ḥayy kih baʿḍī az ān nufūs mawjūdand va
hamchunīn dar aʿdād-i lā-nihāyih kih muntahī bi-Vāḥid-i Bayān mīshavad
mīfarmāyad kih dar kull dīdih namīshavad magar jamāl-i ān ẓuhūr.

The teaching of the Point of the Bayán is as follows: He calls the
source of the Manifestation the Tree, and then regards all the
manifestations of that Revelation as branches, twigs, leaves, and fruits
of the Tree of Revelation, as can be seen. For example, the Primal Point
Himself---may the spirit of all besides Him be sacrificed for Him---was
the Tree of divine Revelation, and all who profess His religion are
reckoned as His leaves, branches, twigs, and fruits. Consider all
Revelations in the same way, and describe them by whatever name and
designation you wish: for example, the Point as the ocean and the others
as waves, or as the sun and the others as mirrors. In this station, if
He says of any one of the manifestations of that Revelation, ``He is Me,
and I am he,'' this is true, and there is no doubt about it. Thus,
concerning the Letters of the Living, some of whom are still alive, and
likewise the limitless numbers that culminate in the Váḥid of the Bayán,
He says that in all of them nothing is seen but the beauty of that
Manifestation.

\subsection{163}\label{section-658}

\begin{otherlanguage}{persian}

جميع اشیاء در مقامى منسوب بحقّند و از حقّند و بحقّ قائمند و از حقّ ظاهرند و
بحقّ راجعند و بحقّ ناظرند و بحقّ متحرّك و اللّه الّذى لا اله الّا هو يك مقامى
وراى آنچه مشاهده مينمائى در سرادق معانى الهى مستور است كه اگر انامل قدرت
ربّ صمدانى حجاب سرادق روحانى را حركت دهد و از آن وجه معانى مستوره اشراقى
ظاهر شود ننصعق و ينصعق كلّ من فى السّموات و الأرض و لكن كور از ظهور انوار
شمس معانى مهجور و كر از صرير قلم سلطان لايزالى ممنوع

\end{otherlanguage}

Jamīʿ-i ashyāʾ dar maqāmī mansūb bi-Ḥaqq-and va az Ḥaqq-and va bi-Ḥaqq
qāʾimand va az Ḥaqq ẓāhirand va bi-Ḥaqq rājiʿand va bi-Ḥaqq nāẓirand va
bi-Ḥaqq mutaḥarrik. Va Allāhi alladhī lā ilāha illā huwa, yak maqāmī
varā-yi ānchih mushāhadih mīnamāʾī dar surādiq-i maʿānī-yi ilāhī mastūr
ast kih agar anāmil-i qudrat-i Rabb-i ṣamadānī ḥijāb-i surādiq-i rūḥānī
rā ḥarakat dihad va az ān vajh-i maʿānī-yi mastūrih ishrāqī ẓāhir
shavad, nanṣaʿiqu va yanṣaʿiqu kullu man fī al-samāwāti va al-arḍ. Va
lākin kūr az ẓuhūr-i anvār-i shams-i maʿānī mahjūr va kar az ṣarīr-i
qalam-i Sulṭān-i lā-yazālī mamnūʿ.

In one station, all things are attributed to the Truth, are from the
Truth, subsist through the Truth, appear from the Truth, return to the
Truth, look toward the Truth, and move through the Truth. By God,
besides whom there is no god! Beyond what you observe, a station lies
concealed in the tabernacle of divine meanings. Should the fingers of
the power of the Eternal Lord stir the veil of the spiritual tabernacle,
and a radiance shine forth from that countenance of hidden meanings, We
would swoon, and everyone in the heavens and on earth would swoon. Yet
the blind are shut out from the shining lights of the sun of meanings,
and the deaf are deprived of the sound of the Pen of the everlasting
Sovereign.

\subsection{164}\label{section-659}

\begin{otherlanguage}{persian}

بارى اين فضايلى كه من حيث لا يشعر بشما القاء نموده‌اند كه بشجره اعراض
راجع نموده بر فرض تسليم بكلّ اشیآء در مقامى راجع بوده و خواهد بود از ارض
تحديد بسماء. عزّ تجريد پرواز كن تا آنچه ذكر مينمايم بچشم خود ملاحظه نمائى
و همچه ندانيد كه اين خادم فانى عمر تلف نموده و آب و علف خورده و مثل شما
بوهم و تقليد اكتفا نموده بارى خوشتر آن است كه اينمطالب را بگذاريم چه كه
ذكر اينمقام فوق ادراك همج رعاع بوده وخواهد بود

\end{otherlanguage}

Bārī īn faḍāʾilī kih min ḥaythu lā yashʿur bi-shumā ilqā namūdih-and kih
bi-shajari h-yi iʿrāḍ rājiʿ namūdih, bar farḍ-i taslīm bi-kull-i ashyāʾ
dar maqāmī rājiʿ būdih va khwāhad būd. Az arḍ-i taḥdīd bi-samā-yi ʿizz-i
tajrīd parvāz kun tā ānchih dhikr mīnamāyam bi-chashm-i khud mulāḥiẓih
namāʾī. Va hamchih nadānīd kih īn khādim-i fānī ʿumr talaf namūdih va āb
va ʿalaf khurdih va mithl-i shumā bi-vahm va taqlīd iktifā namūdih. Bārī
khushtar ān ast kih īn maṭālib rā bigudhārīm, chih kih dhikr-i īn maqām
fawq-i idrāk-i hamaj-i raʿāʿ būdih va khwāhad būd.

In short, even if one were to concede these virtues that have been
conveyed to you unwittingly and that you have attributed to the tree of
rejection, they have applied, and will apply, in one station to all
things. Soar from the earth of limitation to the heaven of the glory of
detachment, so that you may see with your own eyes what I am describing.
Do not imagine that this mortal servant has wasted his life, consumed
water and fodder, and, like you, contented himself with fancy and
imitation. In any case, it is better to leave these matters aside, for
discussion of this station has been, and will remain, beyond the
understanding of the ignorant rabble.

\subsection{165}\label{section-660}

\begin{otherlanguage}{persian}

بعالم ادراك شما راجع شويم و از مطالب آن نفوس ذكر نمائيم آنكه نقطه اولى
روح ماسواه فداه اصل شجره امر و بحر اعظم و شمس حقيقت و سماء قدم را نفس
ظهور دانسته در هر ظهور و مادونش را خلق آن ظهور كه بكلمه از حركت قلمش خلق
شده‌اند چنانچه ميفرمايد اگر بشجره از اشجار مغرسه در اين ارض ظاهره نفس
ظهور بفرمايد كه اين حجّت است بر كلّ شیء هرآينه حجّت است مثل آنچه خلق
فرموده‌است حجج را از قبل و فرقى مابين اين دو مشهود نه و اين حجيّت نظر
بتجلّى شمس كلمه است در او كه بقول او خلق شده چنانچه حجر محلّ طواف شده و
معبد عباد گشته حال ملاحظه كن كه اگر آن حجر نبود و حجر ديگر را ميفرمود او
هم بهمين شرافت مشرّف ميشد و اين شرافت اسماء است در ملكوت اسماء

\end{otherlanguage}

Bi-ʿālam-i idrāk-i shumā rājiʿ shavīm va az maṭālib-i ān nufūs dhikr
namāʾīm. Ānkih Nuqṭih-yi Ūlā, rūḥ-i mā-sivāhu fidāh, aṣl-i shajari h-yi
amr va baḥr-i aʿẓam va shams-i ḥaqīqat va samā-yi qidam rā nafs-i ẓuhūr
dānistih dar har ẓuhūr va mādūnash rā khalq-i ān ẓuhūr kih bi-kalimih az
ḥarakat-i qalamash khalq shudih-and. Chunānchih mīfarmāyad agar
bi-shajarih az ashjār-i mughar rasih dar īn arḍ-i ẓāhirih nafs-i ẓuhūr
bifarmāyad kih īn ḥujjat ast bar kull-i shayʾ, har-āyinih ḥujjat ast,
mithl-i ānchih khalq farmūdih-ast ḥujaj rā az qabl va farqī mābayn-i īn
du mashhūd nah. Va īn ḥujjīyat naẓar bi-tajallī-yi shams-i kalimih ast
dar ū kih bi-qawl-i ū khalq shudih, chunānchih ḥajar maḥall-i ṭavāf
shudih va maʿbad-i ʿibād gashtih. Ḥāl mulāḥiẓih kun kih agar ān ḥajar
nabūd va ḥajar-i dīgar rā mīfarmūd, ū ham bi-hamīn shirāfat musharraf
mīshud va īn shirāfat-i asmā ast dar malakūt-i asmā.

Let Us return to the realm of your understanding and speak of the
contentions of those people. The Primal Point---may the spirit of all
besides Him be sacrificed for Him---has regarded the root of the Tree of
the Cause, the Most Great Ocean, the Sun of Truth, and the Heaven of
Eternity, in every Revelation, as the Manifestation Himself, and all
besides Him as the creation of that Revelation, brought into being by a
word through the movement of His Pen. Thus He says that if the
Manifestation Himself should declare of a tree planted in this visible
earth, ``This is a proof for all things,'' it would assuredly be a
proof, just as the proofs He created before, and no difference would be
seen between the two. This standing as a proof is due to the radiance
within it of the sun of the Word, for it was created by His utterance.
Thus a stone has become a place of circumambulation and a sanctuary of
worship for His servants. Now consider: if it had not been that stone,
and He had pronounced this concerning another stone, it too would have
been invested with this same honor. This is the honor of names in the
kingdom of names.

\subsection{166}\label{section-661}

\begin{otherlanguage}{persian}

ايكاش نفوس بالغ مشاهده ميشدند تا رشحى از طمطام اسماء در اينمقام ذكر
ميگشت و معلوم ميشد كه جميع اشیاء در آن ساحت مقدّس در صقع واحد مشهود و
امتيازات اسمائيّه نظر بحكمت بالغه الهيّه بوده كه از انظر بريّه مستور گشته
الّا من شآء ربّه بهتر آنكه اينمقام هم از محتجبين مستور ماند چنانچه مستور
بوده و هست اين قدر بدان كه اين اسامى صادق و اين شرافت باقى مادام كه در
ظلّ حقّ بوده و بآن اسم متمسّك و بعد از تجاوز و خروج از حدّ خلعت اسمى اخذ
ميشود و مطروح و مردود عنداللّه مذكور و مشهود و بشأنى از رضوان مشاهده و
مكاشفه دور مانده‌اند كه مطالب وهميّه خود را هر وقت اراده نموده كه ذكر
نمايند بروايتى از روايات ملل قبل متمسّك ميشوند و ذكر مينمايند

\end{otherlanguage}

Ay kāsh nufūs-i bāligh mushāhadih mīshudand tā rashḥī az ṭamṭām-i asmā
dar īn maqām dhikr mīgasht va maʿlūm mīshud kih jamīʿ-i ashyāʾ dar ān
sāḥat-i muqaddas dar ṣuqʿ-i vāḥid mashhūd va imtiyāzāt-i asmāʾīyih naẓar
bi-ḥikmat-i bālighih-yi ilāhīyih būdih kih az {[}corrupted text{]}
barīyih mastūr gashtih illā man shāʾa rabbuhu. Bihtar ānkih īn maqām ham
az muḥtajibīn mastūr mānad, chunānchih mastūr būdih va hast. Īn qadr
bidān kih īn asāmī ṣādiq va īn shirāfat bāqī mādām kih dar ẓill-i Ḥaqq
būdih va bi-ān ism mutamassik, va baʿd az tajāvuz va khurūj az ḥadd
khilʿat-i ismī akhdh mīshavad va maṭrūḥ va mardūd ʿinda Allāh madhkūr va
mashhūd. Va bi-shaʾnī az riḍvān-i mushāhadih va mukāshifih dūr
māndih-and kih maṭālib-i vahmīyih-yi khud rā har vaqt irādih namūdih kih
dhikr namāyand bi-rivāyatī az rivāyāt-i milal-i qabl mutamassik
mīshavand va dhikr mīnamāyand.

If only mature souls could be found, so that a drop from the surging
ocean of names might be mentioned in this connection! It would then
become clear that all things in that sacred court are seen on a single
plane, and that distinctions among names are due to the consummate
divine wisdom, which has remained hidden from the {[}corrupted original
text{]} of creation, except those whom their Lord wills. It is better
that this station too remain hidden from the veiled, as it has been and
still is. Know only this much: these names hold true and this honor
endures so long as a thing remains beneath the shadow of the Truth and
clings to that name. Once it transgresses and passes beyond the limit,
the robe of the name is taken away, and it is spoken of and seen in
God's sight as cast out and rejected. They have remained so far from the
paradise of vision and unveiling that whenever they wish to express
their fanciful notions, they seize upon and cite a tradition from among
the traditions of earlier religious communities.

\subsection{167}\label{section-662}

\begin{otherlanguage}{persian}

بگو اى بی‌بصران قبل و بعد را از اين امر بشناسيد و عارف شويد چه كه اين امر
بدونش شناخته نشده و نخواهد شد آنچه در دست ناس بوده جميع توهّمات لايغنيه
بوده و خواهد بود و آنچه ذكر نمائيد نزد اولو البصر معلوم و واضحست كه از
عدم بلوغ بچنين اذكار تمسّك جسته‌اند حقّ ظاهر و مشهود است او را ميزان كنيد و
از ميزان اللّه تجاوز منمائيد و جميع امور را باين ميزان الهى بسنجيد هر
كدام را مطابق ملاحظه نموديد قبول نمائيد چه از قبل و چه از بعد و هر كدام
مخالف مشاهده شد دعه عن ورائك و كن فى استقامة بديع

\end{otherlanguage}

Bigū ay bī-baṣarān, qabl va baʿd rā az īn amr bishināsīd va ʿārif
shavīd, chih kih īn amr bi-dūnash shinākhtih nashudih va nakhwāhad shud.
Ānchih dar dast-i nās būdih jamīʿ tavahhumāt-i lā-yughniyih būdih va
khwāhad būd. Va ānchih dhikr namāʾīd nazd-i ulū al-baṣar maʿlūm va vāḍiḥ
ast kih az ʿadam-i bulūgh bi-chunīn adhkār tamassuk justih-and. Ḥaqq
ẓāhir va mashhūd ast, ū rā mīzān kunīd va az mīzān-i Allāh tajāvuz
manamāʾīd va jamīʿ-i umūr rā bi-īn mīzān-i ilāhī bisanjīd. Har kudām rā
muṭābiq mulāḥiẓih namūdīd qabūl namāʾīd, chih az qabl va chih az baʿd,
va har kudām mukhālif mushāhadih shud, daʿhu ʿan warāʾika va kun fī
istiqāmatin badīʿ.

Say: ``You who lack insight! Recognize and understand what came before
and what comes after through this Cause, for this Cause has not been,
and will not be, known through anything but itself. Everything in
people's possession has been, and will remain, unavailing imaginings.
Whatever you may say, it is evident and clear to those with insight that
such statements have been seized upon out of immaturity. The Truth is
manifest and visible. Make Him the standard, do not transgress God's
standard, and weigh all matters by this divine standard. Accept whatever
you find to conform to it, whether from before or after; and whatever is
found to conflict with it, cast it behind you and stand firm with
wondrous steadfastness.''

\subsection{168}\label{section-663}

\begin{otherlanguage}{persian}

ایدوست زبانى من مظاهر ظهور را ثمره شجره ظهور ملاحظه كن و اين اثمار
ماداميكه بشجره متمسّك از او محسوب و بعد از قطع از شجره هبوب ارياح نفس و
هوى چنان متغيّرش نمايد كه از طعم و لطافت و طراوت جميع ممنوع شود و محروم
گردد و اين اذكار نظر بعرفان ناس ذكر شده و الّا قسم بجمال علىّ ابهى روح
ماسواه فداه كه نفس ظهور را با دونش نسبت و ربط و مشابهت و مشاكلت بهيچ‌وجه
نبوده و نخواهد بود و اين نسبتها در عوالم اسماء مذكور و الّا ساحت ظهور
مقدّس است از آنچه ذكر شده و ميشود

\end{otherlanguage}

Ay dūst-i zabānī-yi man, maẓāhir-i ẓuhūr rā thamarih-yi shajari h-yi
ẓuhūr mulāḥiẓih kun, va īn athmār mādāmīkih bi-shajarih mutamassik az ū
maḥsūb, va baʿd az qaṭʿ az shajarih hubūb-i aryāḥ-i nafs va havā chunān
mutaghayyirash namāyad kih az ṭaʿm va laṭāfat va ṭarāvat jamīʿ mamnūʿ
shavad va maḥrūm gardad. Va īn adhkār naẓar bi-ʿirfān-i nās dhikr
shudih, va illā qasam bi-Jamāl-i ʿAlī-yi Abhā, rūḥ-i mā-sivāhu fidāh,
kih nafs-i ẓuhūr rā bā dūnash nisbat va rabṭ va mushābahat va mushākalat
bi-hīch vajh nabūdih va nakhwāhad būd. Va īn nisbat-hā dar ʿavālim-i
asmā madhkūr, va illā sāḥat-i ẓuhūr muqaddas ast az ānchih dhikr shudih
va mīshavad.

My friend in words alone! Regard the manifestations of the Revelation as
the fruit of the Tree of Revelation. These fruits are reckoned as
belonging to the Tree so long as they cling to it. Once severed from the
Tree, the blowing winds of self and desire alter them so completely that
they are deprived of all flavor, delicacy, and freshness. These
descriptions have been given in accordance with people's understanding.
Otherwise, I swear by the exalted, all-glorious Beauty---may the spirit
of all besides Him be sacrificed for Him---that the Manifestation
Himself has never had, and will never have, any relationship,
connection, resemblance, or likeness whatsoever to anyone besides Him.
These relationships are spoken of in the worlds of names; otherwise, the
court of the Manifestation is sanctified above all that has been or is
being said.

\subsection{169}\label{section-664}

\begin{otherlanguage}{persian}

فتعالى عما ظننت و عرفت و علمت و قلت و تقول معانى كلمات بيان را احدى جز
نفس رحمن مطلع نبوده و نخواهد بود اى كاش در ظلّ شجره فردانيّه مستظلّ ميشدى و
از اثمارش كه بهيئت لا اله الّا هو ظاهر شده‌اند مرزوق ميگشتى و معلوم ميشد
كه مقصود منزل بيان از بيان چه بوده و لكن قضى منك بما اكتسبت يداك

\end{otherlanguage}

Fa-taʿālā ʿammā ẓananta va ʿarafta va ʿalimta va qulta va taqūl.
Maʿānī-yi kalimāt-i Bayān rā aḥadī juz nafs-i Raḥmān muṭṭaliʿ nabūdih va
nakhwāhad būd. Ay kāsh dar ẓill-i shajari h-yi fardānīyih mustaẓill
mīshudī va az athmārash kih bi-hayʾat-i lā ilāha illā huwa ẓāhir
shudih-and marzūq mīgashtī va maʿlūm mīshud kih maqṣūd-i Munzil-i Bayān
az Bayān chih būdih, va lākin quḍiya minka bi-mā iktasabat yadāk.

Exalted is He above what you have imagined, recognized, known, said, and
now say. No one but the All-Merciful Himself has known, or will ever
know, the meanings of the words of the Bayán. If only you had sought
shelter beneath the shadow of the Tree of Oneness and been nourished by
its fruits, which have appeared in the form of ``There is no god but
Him''! Then it would have become clear what the Revealer of the Bayán
intended by the Bayán. But judgment has been passed upon you through
what your own hands have earned.

\subsection{170}\label{section-665}

\begin{otherlanguage}{persian}

و اينكه نوشته كه حال انصاف بدهيد خود حضرت باب آيات را بهتر تميز ميداد يا
بنده و شما و اين تقليد نيست الى آخر قد كبر قول من القاك و عظم كفره و عدم
كينونته ميگويند و شاعر نيستند قسم بنقطه وجود در حينيكه اين كلمات از قلم
معرضين جارى كلّ اشیاء بر عقل و درايتشان هم متحيّر و هم مضحك چه كه نقطه
اولى بكينونته ظاهر و آيات الهى بمثل غيث هاطل نازل مع‌ذلك از آن اعراض
نموده‌اند و بذكر اثبات دونش مشغول گشته‌اند

\end{otherlanguage}

Va īnke neveshte ke ḥāl enṣāf bedahīd khod-e Ḥażrat-e Bāb āyāt rā behtar
tamīz mīdād yā bande va shomā va īn taqlīd nīst ilā ākhir. Qad kabura
qawlu man alqāka wa ʿaẓuma kufruhu wa ʿadamu kaynūnatihi. Mīgūyand va
shāʿer nīstand. Qasam be Noqṭe-ye Vojūd dar ḥīnī ke īn kalemāt az
qalam-e moʿreżīn jārī, koll-e ashyāʾ bar ʿaql va derāyatashān ham
motaḥayyer va ham moḍḥek, che ke Noqṭe-ye Ūlā bi-kaynūnatihi ẓāher va
āyāt-e elāhī be methl-e ghayth-e hāṭel nāzel, maʿa dhālek az ān eʿrāż
nemūde-and va be dhekr-e ethbāt-e dūnash mashghūl gashte-and.

And as for what he has written: ``Now judge fairly: did the Báb Himself
discern the verses better, or do you and I? And this is not imitation,''
and so forth. Grave indeed is the utterance of the one who prompted you,
and great are his unbelief and his lack of true existence. They speak
without awareness. I swear by the Point of Existence: as these words
flow from the pens of those who turn away, all things are both
bewildered by their reason and discernment and moved to laughter. For
the Primal Point has appeared in His own Being, and the divine verses
descend like pouring rain; nevertheless, they have turned away from Him
and busied themselves with arguments to establish the claims of another.

\subsection{171}\label{section-666}

\begin{otherlanguage}{persian}

اگر چشم بصيرت ميداشتند ملاحظه مينمودند كه نقطه بيان از ظلم مظاهر شيطان و
مطالع حسبان فواللّه بصيحه و ندبه و نوحه مشغولست چه كه فرموده آنچه وصف در
بيان نازل شده مقصود ظهور آخر است و حال مشاهده ميشود كه طلعت اولى در كرّه
اخرى ظاهر و بكلمات ظهور قبلش بر نفسش سيف كشيده‌اند چنانچه از قبل هم امثال
اين نفوس بكلمات ظهور قبلش مستدلّ شده دم مطهّرش را ريختند قل قد فعلتم و
تفعلون حينئذ ما لافعل النّمرود و لا الشّمر و لا السّنان و لا ابن ملجم و لا
الأوّل و لا الثّانى

\end{otherlanguage}

Agar chashm-e baṣīrat mīdāshtand molāḥeẓe mīnemūdand ke Noqṭe-ye Bayān
az ẓolm-e maẓāher-e shayṭān va maṭāleʿ-e ḥosbān, fa-wa-Allāh, be ṣayḥe
va nodbe va nowḥe mashghūl ast, che ke farmūde ānche vaṣf dar Bayān
nāzel shode maqṣūd ẓohūr-e ākhar ast. Va ḥāl moshāhede mīshavad ke
Ṭalʿat-e Ūlā dar karre-ye okhrā ẓāher va be kalemāt-e ẓohūr-e qablash
bar nafsash sayf keshīde-and, chenānche az qabl ham amthāl-e īn nofūs be
kalemāt-e ẓohūr-e qablash mostadell shode dam-e moṭahharash rā rīkhtand.
Qul qad faʿaltum wa tafʿalūna ḥīnaʾidhin mā lā faʿala al-Nimrūd wa lā
al-Shimr wa lā al-Sinān wa lā Ibn Muljam wa lā al-Awwal wa lā al-Thānī.

Had they possessed the eye of insight, they would have seen that the
Point of the Bayán is, by God, crying out, lamenting, and mourning
because of the oppression of the manifestations of Satan and the sources
of vain imaginings. For He has said that all the descriptions revealed
in the Bayán refer to the later Revelation. And now it is seen that the
Primal Countenance has appeared once more, yet with the words of His
previous Revelation they have drawn a sword against His very Self, just
as, in the past, souls like these invoked the words of His previous
Revelation and shed His sacred blood. Say: You have done, and are now
doing, what neither Nimrod, nor Shimr, nor Sinán, nor Ibn Muljam, nor
the First, nor the Second ever did.

\subsection{172}\label{section-667}

\begin{otherlanguage}{persian}

كاش از بصر و فؤاد كه محروم شده‌اند اقلّا شامّه ميداشتند و اللّه از آن هم
محروم شده‌اند چه كه نكهت دفراء را از عرف طيب بقاء تميز نداده‌اند و برايحه
جعل خود را از عطر رياحين رضوان عزّ صمدانى محروم نموده‌اند و اين كلمه را
موهوم مدان چنانچه بسيّد محمّد اصفهانى كه واللّه هر نفسيكه اقلّ من ذرّ صاحب
شامّه باشد رايحه منتنه‌اش را از فرسنگهاى بعيده ميشنود و احتراز ميجويد با
او مأنوس شده از عرف گلزار قدس ربّانى بعيد و بی‌نصيب مانده‌اند

\end{otherlanguage}

Kāsh az baṣar va foʾād ke maḥrūm shode-and aqallan shāmme mīdāshtand.
Va-Allāh az ān ham maḥrūm shode-and, che ke nakhat-e dafrā rā az ʿarf-e
ṭīb-e baqā tamīz nadāde-and va be rāyeḥe-ye joʿal khod rā az ʿeṭr-e
rayāḥīn-e Reżvān-e ʿezz-e ṣamadānī maḥrūm nemūde-and. Va īn kaleme rā
mowhūm madān, chenānche be Sayyed Moḥammad-e Eṣfahānī, ke va-Allāh har
nafsī ke aqall-e men dharr ṣāḥeb-e shāmme bāshad rāyeḥe-ye montene-ash
rā az farsanghā-ye baʿīde mīshenavad va eḥterāz mījūyad, bā ū maʾnūs
shode az ʿarf-e golzār-e qods-e rabbānī baʿīd va bī-naṣīb mānde-and.

If only, deprived as they are of sight and heart, they at least
possessed a sense of smell! By God, they have been deprived of that as
well, for they have failed to distinguish a foul stench from the sweet
fragrance of eternity, and for the odor of a dung beetle they have
deprived themselves of the perfume of the flowers of the Paradise of
everlasting glory. Do not regard these words as imaginary. They have
become intimate with Siyyid Muḥammad-i-Iṣfahání---whose foul odor, by
God, anyone possessing even less than an atom of a sense of smell would
detect from many leagues away and shun---and have remained far removed
from, and deprived of, the fragrance of the rose garden of divine
holiness.

\subsection{173}\label{section-668}

\begin{otherlanguage}{persian}

متحيّرم كه با چنين انفس چه ذكر نمايم و چه القاء كنم بادلّه بر جعل ميتوان
ثابت نمود كه رايحه گلزار محبوب است لا فو‌نفس المحبوب لو انت من العارفين
آيا بكر ميتوان محقّق نمود كه نغمات الهى نيكو است لا فو‌ربّ الأسماء و الصّفات
لو انت من المستشعرين

\end{otherlanguage}

Motaḥayyeram ke bā chenīn anfos che dhekr namāyam va che elqā konam. Be
adelle bar joʿal mītavān thābet nemūd ke rāyeḥe-ye golzār maḥbūb ast? Lā
fa-wa-nafsi al-Maḥbūb law anta min al-ʿārifīn. Āyā be kar mītavān
moḥaqqaq nemūd ke naghamāt-e elāhī nīkū ast? Lā fa-wa-Rabbi al-asmāʾ wa
al-ṣifāt law anta min al-mustashʿirīn.

I am bewildered as to what I should say to such souls and what I should
convey to them. Can one prove by arguments to a dung beetle that the
fragrance of a rose garden is delightful? No, by the Self of the
Beloved, if you are among those who know! Can one demonstrate to the
deaf that the divine melodies are beautiful? No, by the Lord of names
and attributes, if you are among those who perceive!

\subsection{174}\label{section-669}

\begin{otherlanguage}{persian}

ميفرمايند ببيان و ما نزّل فيه از آن ساذج حقيقت محتجب نمانيد و شما بالقاى
آن نفس مشرك اوصاف كلمات معرضين نوشته و بمقرّ اطهر ارسال داشته‌ايد و
اعتراضا على اللّه بكلمات نقطه اولى هم مستدلّ شده‌ايد از براى نفس نقطه از
كلمات قبلش استدلال مينمايند و شاعر نيستند و اگر شاعر بودند بايد آنچه از
جمال احديّه اليوم ظاهر بآن متمسّك شوند بگو اى بی‌‌بصر چشم بگشا شمس قدم مقابل
وجهت نشسته و قلب مباركش از ظلمت مجروح باو ناظر شو و رحم كن و از دونش از
او سؤال مكن ببحر معنى رسيدى از تحريف كلمات بگذر بشمس واصل شدى از سراج
منقطع شو

\end{otherlanguage}

Mīfarmāyand be Bayān va mā nuzzila fīhi az ān Sādhej-e Ḥaqīqat moḥtajeb
namānīd, va shomā be elqā-ye ān nafs-e moshrek owṣāf-e kalemāt-e
moʿreżīn neveshte va be maqarr-e aṭhar ersāl dāshte-īd va iʿtirāḍan ʿalā
Allāh be kalemāt-e Noqṭe-ye Ūlā ham mostadell shode-īd. Az barā-ye
nafs-e Noqṭe az kalemāt-e qablash estedlāl mīnemāyand va shāʿer nīstand,
va agar shāʿer būdand bāyad ānche az Jamāl-e Aḥadiyye al-yawm ẓāher be
ān motamassek shavand. Begū ey bī-baṣar chashm bogshā, Shams-e Qedam
moqābel-e vajhat neshaste va qalb-e mobārakash az ẓolmat majrūḥ. Be ū
nāẓer show va raḥm kon va az dūnash az ū soʾāl makon. Be baḥr-e maʿnā
resīdī, az taḥrīf-e kalemāt bogdhar. Be shams vāṣel shodī, az serāj
monqaṭeʿ show.

He says: ``Do not let the Bayán and what has been revealed in it veil
you from that Essence of Truth.'' Yet, at the prompting of that
idolatrous soul, you have written descriptions of the words of those who
turn away and sent them to the most holy Seat, and, in objection to God,
you have even invoked the words of the Primal Point as evidence. They
reason about the Point Himself from His earlier words, without
awareness. If they were aware, they would have to hold fast to whatever
appears today from the Beauty of Oneness. Say: O sightless one, open
your eyes! The Sun of Eternity is seated before your face, and His
blessed heart is wounded by your oppression. Look upon Him and have
mercy, and do not question Him about anyone besides Him. You have
reached the ocean of meaning: leave off distorting words. You have
attained the sun: detach yourself from the lamp.

\subsection{175}\label{section-670}

\begin{otherlanguage}{persian}

نقطه بيان ميفرمايد جميع بيان اليوم منوط بتصديق مظهر ظهور است مع‌ذلك از
كلمات محكمه الهيّه غافل شده و بهواى نفس خود مشغول گشته‌اند فواللّه مثل آن
مشاهده ميشوند كه از براى قطره ماء ملح اجاج اوصاف بحر عذب فرات ذكر نمايند
عجب است كه شما نوشته‌ايد حال انصاف دهيد و تقليد نيست مع آنكه بسلاسل تقليد
چنان بسته شده‌ايد كه اگر جنود سموات و ارض جمع شوند از گشودنش عاجز گردند

\end{otherlanguage}

Noqṭe-ye Bayān mīfarmāyad jamīʿ-e Bayān al-yawm manūṭ be taṣdīq-e
Maẓhar-e Ẓohūr ast. Maʿa dhālek az kalemāt-e moḥkame-ye elāhiyye ghāfel
shode va be havā-ye nafs-e khod mashghūl gashte-and. Fa-wa-Allāh methl-e
ān moshāhede mīshavand ke az barā-ye qaṭre-ye māʾ-e melḥ-e ojāj owṣāf-e
baḥr-e ʿadhb-e forāt dhekr namāyand. ʿAjab ast ke shomā neveshte-īd ḥāl
enṣāf dehīd va taqlīd nīst, maʿa ānke be salāsel-e taqlīd chenān baste
shode-īd ke agar jonūd-e samāvāt va arż jamʿ shavand az goshūdanash
ʿājez gardand.

The Point of the Bayán says that the entire Bayán today depends upon the
confirmation of the Manifestation of Revelation. Nevertheless, they have
neglected the decisive words of God and busied themselves with their own
desires. By God, they are seen to be like those who ascribe the
qualities of an ocean of sweet, refreshing water to a drop of bitter
brine. It is astonishing that you have written, ``Now judge fairly,''
and, ``This is not imitation,'' when you are so tightly bound in the
chains of imitation that, were the hosts of heaven and earth to gather
together, they would be powerless to undo them.

\subsection{176}\label{section-671}

\begin{otherlanguage}{persian}

و ديگر ذكر انصاف هم نموده حال اى با انصاف انصاف ده آنچه نوشته و بجمال
قدم نسبت داده جز از سيّد محمّد شنيده بكدام شرع قول مدّعى مقبول و مسموع اگر
در يكى از شرايع الهى كه از قبل لا قبل مابين ناس بوده حكمى آوردى كه مدلّ
باشد بر اينكه كلمات مدّعى مسموع است اين عباد اذعان مينمايند قول شما را
بقول نفسيكه لايعرف اليمين عن اليسار و لا النّور عن الظّلمة و جميع ذرّات
شهادت بر فسق و فجور و عنادش داده و ميدهند بر حقّ وارد آورده آنچه را كه
هيچ ظالمى وارد نياورد

\end{otherlanguage}

Va dīgar dhekr-e enṣāf ham nemūde. Ḥāl ey bā-enṣāf enṣāf deh, ānche
neveshte va be Jamāl-e Qedam nesbat dāde joz az Sayyed Moḥammad shenīde?
Be kodām sharʿ qowl-e moddaʿī maqbūl va masmūʿ? Agar dar yekī az
sharāyeʿ-e elāhī ke az qabl-e lā qabl mābayn-e nās būde ḥokmī āvordī ke
modell bāshad bar īnke kalemāt-e moddaʿī masmūʿ ast, īn ʿebād edhʿān
mīnemāyand qowl-e shomā rā. Be qowl-e nafsī ke lā yaʿrifu al-yamīn ʿan
al-yasār wa lā al-nūr ʿan al-ẓulma va jamīʿ-e dharrāt shahādat bar fesq
va fojūr va ʿenādash dāde va mīdehand, bar Ḥaqq vāred āvorde ānche rā ke
hīch ẓālemī vāred nayāvard.

He has also spoken of fairness. Now, O fair-minded one, judge fairly:
has he heard what he has written and attributed to the Ancient Beauty
from anyone but Siyyid Muḥammad? Under which religious law is the
claimant's own statement accepted and admitted as testimony? If you
produce a ruling from any of the divine laws that have existed among
people from a beginning with no beginning, showing that the claimant's
words are admissible, these servants will acknowledge your assertion. On
the word of a soul who cannot distinguish right from left or light from
darkness, and to whose wickedness, depravity, and malice every atom has
borne and continues to bear witness, you have inflicted upon the True
One what no oppressor has ever inflicted.

\subsection{177}\label{section-672}

\begin{otherlanguage}{persian}

قريب هفتاد نفس در اين سفر بوده‌اند و قريب همين عدد هم نفوس از بعد بمقرّ
ظهور وارد شده‌اند اگر يك نفس بگويد احدى باين دو مشرك باللّه معارضه نموده و
يا نزاع و مجادله و مكالمه مابين شده آنچه بقلم كذب نوشته‌اند مسلّم ميداريم
فانصف فى نفسك يا عبد مع‌ذلك بقول سيّد محمّد ظلم بحقّ نسبت داده مع آنكه آنچه
حكم فرمايند مختارند و آنچه عامل شوند عمل حقّ بوده و خواهد بود و لكن سمع
نه تا اصغا نمايد و قلب نه تا ادراك كند از براى كر هر چه گفته شود حقّ يفعل
ما یشآء و يحكم مايريد بوده چه فايده حاصل شود

\end{otherlanguage}

Qarīb-e haftād nafs dar īn safar būde-and va qarīb-e hamīn ʿadad ham
nofūs az baʿd be maqarr-e ẓohūr vāred shode-and. Agar yek nafs begūyad
aḥadī be īn do moshrek be Allāh moʿāreże nemūde va yā nezāʿ va mojādele
va mokāleme mābayn shode, ānche be qalam-e kedhb neveshte-and mosallam
mīdārīm. Fa-anṣif fī nafsika yā ʿabd. Maʿa dhālek be qowl-e Sayyed
Moḥammad ẓolm be Ḥaqq nesbat dāde, maʿa ānke ānche ḥokm farmāyand
mokhtārand va ānche ʿāmel shavand ʿamal-e Ḥaqq būde va khvāhad būd. Va
lāken samʿ na tā eṣghā namāyad va qalb na tā edrāk konad. Az barā-ye kar
har che gofte shavad Ḥaqq yafʿalu mā yashāʾ va yaḥkumu mā yurīd būde,
che fāyde ḥāṣel shavad?

Nearly seventy souls were on this journey, and about the same number
subsequently arrived at the seat of Revelation. If a single soul should
say that anyone opposed these two who associate partners with God, or
that any quarrel, dispute, or exchange of words took place between them,
We will concede what they have written with the pen of falsehood. So
judge fairly within yourself, O servant. Nevertheless, on Siyyid
Muḥammad's word, you have attributed oppression to the True One,
although He is free to decree whatever He wishes, and whatever He does
has been and will be the action of the True One. But there is no ear to
listen and no heart to understand. However often the deaf are told that
the True One does whatever He wills and ordains whatever He desires,
what benefit can result?

\subsection{178}\label{section-673}

\begin{otherlanguage}{persian}

فواللّه اوراق شجرهء طوبی‌ از ظلم آن ظالمان فرو ريخته و اغصان سدره منتهى
درهم شكسته و وجه منير علىّ اعلى از لطمات اشارات هياكل بغضا تغيير نموده و
صدر ممرّد منيرش مجروح گشته و ذيل تقديسش بغبار مفتريات انفس مشركه آلوده
شده از حقّ ميطلبم آيه تطهيرى تا ارض طيّبه مباركه عرفانش را از انفس خبيثه
نجسّه كثيفه مشركه مطهّر نمايد

\end{otherlanguage}

Fa-wa-Allāh owrāq-e shajare-ye Ṭūbā az ẓolm-e ān ẓālemān forū rīkhte va
aghṣān-e Sedre-ye Montahā darham shekaste va vajh-e monīr-e ʿAliyy-e
Aʿlā az laṭamāt-e eshārāt-e hayākel-e baghżā taghyīr nemūde va ṣadr-e
momarrad-e monīrash majrūḥ gashte va dheyl-e taqdīsash be ghobār-e
moftarayāt-e anfos-e moshreke ālūde shode. Az Ḥaqq mīṭalabam āye-ye
taṭhīrī tā arż-e ṭayyebe-ye mobārake-ye ʿerfānash rā az anfos-e
khabīthe-ye najese-ye kathīfe-ye moshreke moṭahhar namāyad.

By God, the leaves of the Tree of Blessedness have fallen through the
oppression of those oppressors, the branches of the Tree beyond which
there is no passing have been shattered, the radiant face of the
Exalted, the Most High, has been altered by the blows of the
insinuations of those embodiments of hatred, His smooth, radiant breast
has been wounded, and the hem of His holiness has been soiled by the
dust of the slanders of idolatrous souls. I ask the True One for a verse
of purification, that He may cleanse the pure and blessed land of His
recognition of wicked, unclean, filthy, idolatrous souls.

\subsection{179}\label{section-674}

\begin{otherlanguage}{persian}

اى با‌انصاف نفسى ظاهر و بكلّ آيات قدرتيّه و ظهورات عزّ صمدانيّه مشهود چنانچه
كلّ مشاهده نموده و مينمايند اگر شما منكرى از ملل مختلفه جويا شو و مع‌ذلك
بهمان آيات كه نقطه اولى روح ماسواه فداه اثبات حقّيّت خود فرموده بل اعظم
ظاهر شده و ميفرمايد يا قوم انّى لنفس علىّ بينكم و جمال محمّد فيكم و كينونة
الرّوح بين السّموات و الأرضين يا قوم خافوا عن اللّه انّى عبد آمنت باللّه و
آياته و كنت راقداً على بساطى مرّت علىّ نسائم الرّحمن و ايقظنى بالحقّ و علّمنى
علم ما كان و ما يكون و اظهرنى بطراز نفسه و انطقنى بثنائه ان انتم تعلمون

\end{otherlanguage}

Ey bā-enṣāf nafsī ẓāher va be koll-e āyāt-e qodratiyye va ẓohūrāt-e
ʿezz-e ṣamadāniyye mashhūd, chenānche koll moshāhede nemūde va
mīnemāyand. Agar shomā monkerī az melal-e mokhtalefe jūyā show. Va maʿa
dhālek be hamān āyāt ke Noqṭe-ye Ūlā rūḥ-e mā sevāh fedāh ethbāt-e
ḥaqqiyyat-e khod farmūde, bal aʿẓam, ẓāher shode va mīfarmāyad: Yā qawm
innī la-nafsu ʿAlī baynakum wa jamālu Muḥammad fīkum wa kaynūnatu al-Rūḥ
bayna al-samāwāt wa al-arḍīn. Yā qawm khāfū ʿan Allāh innī ʿabd āmantu
bi-Allāh wa āyātihi wa kuntu rāqidan ʿalā bisāṭī marrat ʿalayya nasāʾimu
al-Raḥmān wa ayqaẓanī bi-al-ḥaqq wa ʿallamanī ʿilma mā kāna wa mā yakūn
wa aẓharanī bi-ṭirāzi nafsihi wa anṭaqanī bi-thanāʾihi in antum
taʿlamūn.

O fair-minded one! A Soul has appeared, manifest with all the signs of
power and the revelations of everlasting glory, as everyone has
witnessed and continues to witness. If you deny this, inquire among the
various peoples. Moreover, He has appeared with the very same verses by
which the Primal Point---may the spirit of all besides Him be sacrificed
for Him---established His own truth, indeed with greater ones, and He
says: ``O people! I am indeed the very Self of ʿAlí among you, the
Beauty of Muḥammad in your midst, and the Being of the Spirit between
the heavens and the earths. O people! Fear God. I am a servant who has
believed in God and His verses. I lay asleep upon My couch when the
breezes of the All-Merciful passed over Me, and He awakened Me through
the truth, taught Me the knowledge of all that has been and all that
will be, arrayed Me in the adornment of His own Self, and caused Me to
speak His praise, if you but knew.''

\subsection{180}\label{section-675}

\begin{otherlanguage}{arabic}

يا قوم ان لا تؤمنوا لاتتعرّضوا بی‌ و ل اتفعلوا ما منعتم عنه فى كلّ الألواح
اتّقوااللّه و لا تشركون انّى آمنت بما عندكم من صحف اللّه و كتبه و اصدّق ما
امرتم به من لدى اللّه المقتدر المهيمن العزيز القيّوم و من كان ذى شمّ يستنشق
منّى رايحة القدس و من كان ذى ذائقة ليذوق من اثمار هذه الشّجرة ما ذاق محمّد
من سدرة المنتهى و ماذاق علىّ من سدرة البيان ان انتم تعلمون

\end{otherlanguage}

Yā qawm in lā tuʾminū lā tataʿarraḍū bī wa lā tafʿalū mā muniʿtum ʿanhu
fī kull al-alwāḥ. Ittaqū Allāh wa lā tushrikūn. Innī āmantu bi-mā
ʿindakum min ṣuḥuf Allāh wa kutubihi wa uṣaddiqu mā umirtum bihi min
ladā Allāh al-Muqtadir al-Muhaymin al-ʿAzīz al-Qayyūm. Wa man kāna dhī
shamm yastanshiqu minnī rāʾiḥata al-quds wa man kāna dhī dhāʾiqa
li-yadhūq min athmār hādhihi al-shajara mā dhāqa Muḥammad min Sidrat
al-Muntahā wa mā dhāqa ʿAlī min Sidrat al-Bayān in antum taʿlamūn.

``O people! If you do not believe, do not assail Me, and do not do what
you have been forbidden in all the Tablets. Fear God and do not
associate partners with Him. I have believed in the scriptures and books
of God that are in your possession, and I affirm what you have been
commanded by God, the All-Powerful, the All-Protecting, the Almighty,
the Self-Subsisting. Whoever possesses a sense of smell inhales from Me
the fragrance of holiness, and whoever possesses a sense of taste, let
him taste of the fruits of this Tree what Muḥammad tasted from the Tree
beyond which there is no passing, and what ʿAlí tasted from the Tree of
the Bayán, if you but knew.''

\subsection{181}\label{section-676}

\begin{otherlanguage}{arabic}

یا قوم خافوا عن اللّه انّى كنت كأحد منكم واردت ان اصمت تلقاء نفسى و لكنّ
الرّوح اهتزّنى و حرّكنى بالحقّ و انطقنى بآيات اللّه ثمّ ببيّناته و هذا ليس منّى
بل من لدى اللّه المقتدر العزيز المحبوب فواللّه لو كان فى قلوبكم حبّ اللّه
اقلّ عمّا يحصى لم تسفكوا دم هذا المظلوم لانّه لو لم يكن كما انتم تظنّون و
لكن ينسب نفسه الى محبوبكم و بحبّكم هذا الاسم لاينبغى ان تتعرّضوا عليه و
تفعلوا ما يكبّ بوجهه على التّراب روح الأعظم ثمّ حقايق ما خلق بين الكاف و
النّون

\end{otherlanguage}

Yā qawm khāfū ʿan Allāh innī kuntu ka-aḥadin minkum wa aradtu an aṣmuta
tilqāʾa nafsī wa lākinna al-Rūḥa ihtazzanī wa ḥarrakanī bi-al-ḥaqq wa
anṭaqanī bi-āyāt Allāh thumma bi-bayyinātihi wa hādhā laysa minnī bal
min ladā Allāh al-Muqtadir al-ʿAzīz al-Maḥbūb. Fa-wa-Allāh law kāna fī
qulūbikum ḥubbu Allāh aqall ʿammā yuḥṣā lam tasfikū dama hādhā al-Maẓlūm
li-annahu law lam yakun ka-mā antum taẓunnūn wa lākin yansibu nafsahu
ilā maḥbūbikum wa bi-ḥubbikum hādhā al-ism lā yanbaghī an tataʿarraḍū
ʿalayhi wa tafʿalū mā yakubbu bi-wajhihi ʿalā al-turāb rūḥa al-Aʿẓam
thumma ḥaqāʾiqa mā khuliqa bayna al-kāf wa al-nūn.

``O people! Fear God. I was as one of you and wished, of My own accord,
to remain silent, but the Spirit stirred Me, moved Me through the truth,
and caused Me to utter the verses of God and His clear proofs. This is
not from Me, but from God, the All-Powerful, the Almighty, the Beloved.
By God, if there were in your hearts a love of God smaller than could be
reckoned, you would not shed the blood of this Wronged One. For even if
He were not what you suppose, He nevertheless ascribes Himself to your
Beloved; and, because of your love for this Name, you should not assail
Him and do what casts the Spirit of the Most Great face down upon the
dust, and with it the realities of all that has been created between the
letters káf and nún.''

\subsection{182}\label{section-677}

\begin{otherlanguage}{arabic}

و يا قوم انا نهيناكم عمّا نهيتم عنه فى كلّ الألواح و امرناكم بما امرتم به
فى الواح اللّه المهيمن العزيز القيّوم و انتم ما سمعتم قولى و تجاوزتم عن
حدودكم الى ان تصرّفتم فى حرم اللّه و ارتكبتم ما ظهرت الصّيحة بين السّموات و
الأرض و لكن انتم لا تسمعون و كذلك نهيناكم عن التّصرّف فى امآء النّاس ثمّ
نسائهم و كذلك عن سفك الدّمآء و انتم بذلك قمتم علىّ بالاعراض و فعلتم ما لا
فعله المشركون و لو اقتل فى كلّ حين بكلّ الأسياف لن اصدّق افعالكم و بذلك
امرت من لدى اللّه و هذا لم يكن من عندى بل من لدنه و انّه لهو الحقّ علّام
الغيوب

\end{otherlanguage}

Wa yā qawm innā nahaynākum ʿammā nuhītum ʿanhu fī kull al-alwāḥ wa
amarnākum bi-mā umirtum bihi fī alwāḥ Allāh al-Muhaymin al-ʿAzīz
al-Qayyūm wa antum mā samiʿtum qawlī wa tajāwaztum ʿan ḥudūdikum ilā an
taṣarraftum fī ḥaram Allāh wa irtakabtum mā ẓaharat al-ṣayḥatu bayna
al-samāwāt wa al-arḍ wa lākin antum lā tasmaʿūn. Wa kadhālika nahaynākum
ʿan al-taṣarruf fī imāʾ al-nās thumma nisāʾihim wa kadhālika ʿan safk
al-dimāʾ wa antum bi-dhālika qumtum ʿalayya bi-al-iʿrāḍ wa faʿaltum mā
lā faʿalahu al-mushrikūn. Wa law uqtal fī kull ḥīn bi-kull al-asyāf lan
uṣaddiqa afʿālakum wa bi-dhālika umirtu min ladā Allāh wa hādhā lam
yakun min ʿindī bal min ladunhu wa innahu la-huwa al-Ḥaqq ʿAllām
al-ghuyūb.

``O people! We have forbidden you what you were forbidden in all the
Tablets and commanded you what you were commanded in the Tablets of God,
the All-Protecting, the Almighty, the Self-Subsisting. Yet you did not
heed My words and overstepped your bounds, until you violated the
sanctity of God's household and committed deeds that raised a cry
between the heavens and the earth; but you do not hear. Likewise, We
forbade you to violate other people's maidservants and their wives, and
likewise to shed blood. Because of this, you rose against Me in
rejection and did what even those who associate partners with God have
not done. Were I to be slain at every moment by every sword, I would
never sanction your deeds. Thus have I been commanded by God. This is
not from Me but from Him, and He is indeed the Truth, the Knower of all
that is unseen.''

\subsection{183}\label{section-678}

\begin{otherlanguage}{persian}

حال ملاحظه كن اى باانصاف كه كى بی‌انصاف است نفوس زانيه كه در اصل حرم عزّ
صمدانيّه خيانت نموده‌اند و ديگر چه ارتكابات شنيعه كه بآن عامل گشته كه
واللّه الّذى لا اله الّا هو كه قلم از ذكرش حيا مينمايد مع‌ذلك نسبت حقّ باين
نفوس داده و بحقّ نسبت ظلم راجع كرده حال ما بی‌انصافيم و شما با‌انصاف كتاب
مستيقظ را بخوان كه در آن كتاب مرشدت ميرزا يحيى فتوى بر دم جميع نفوس
مقدّسه داده چنانچه نفسى را كه منصوصاً نقطه اولى روح ماسواه فداه حرف ثالث
من يظهره اللّه ناميده بفتواى آن ظالم شهيد شده و كتاب ديگرش كه در آن مخصوص
در اماكن متعدّده فتوى بر قتل كلّ نفوس داده مع‌ذلك جميع را بحقّ راجع نموده

\end{otherlanguage}

Ḥāl molāḥeẓe kon ey bā-enṣāf ke kī bī-enṣāf ast. Nofūs-e zāniye ke dar
aṣl-e ḥaram-e ʿezz-e ṣamadāniyye khiyānat nemūde-and va dīgar che
ertekābāt-e shanīʿe ke be ān ʿāmel gashte ke va-Allāh alladhī lā ilāha
illā huwa ke qalam az dhekrash ḥayā mīnemāyad, maʿa dhālek nesbat-e ḥaqq
be īn nofūs dāde va be Ḥaqq nesbat-e ẓolm rājeʿ karde. Ḥāl mā bī-enṣāfīm
va shomā bā-enṣāf? Ketāb-e Mostayqeẓ rā bekhvān ke dar ān ketāb
morshedat Mīrzā Yaḥyā fatvā bar dam-e jamīʿ-e nofūs-e moqaddase dāde,
chenānche nafsī rā ke manṣūṣan Noqṭe-ye Ūlā rūḥ-e mā sevāh fedāh Ḥarf-e
Thāleth-e Man Yoẓheroh Allāh nāmīde be fatvā-ye ān ẓālem shahīd shode,
va ketāb-e dīgarash ke dar ān makhṣūṣ dar amāken-e motaʿaddede fatvā bar
qatl-e koll-e nofūs dāde. Maʿa dhālek jamīʿ rā be Ḥaqq rājeʿ nemūde.

Now consider, O fair-minded one, who is unfair. Adulterous souls have
acted treacherously within the very household of everlasting glory, and
what other abominable deeds they have committed---by God, besides whom
there is no other god, the pen is ashamed to mention them! Nevertheless,
you have ascribed right to these souls and attributed oppression to the
True One. Are We, then, unfair, and you fair? Read the book Mustayqiẓ,
for in that book your guide, Mírzá Yaḥyá, has issued a decree for the
blood of all holy souls. Indeed, the soul whom the Primal Point---may
the spirit of all besides Him be sacrificed for Him---explicitly named
the Third Letter of Him Whom God Shall Make Manifest was martyred by the
decree of that oppressor. And read his other book, in which, in numerous
places, he has specifically decreed the killing of all souls.
Nevertheless, you have attributed it all to the True One.

\subsection{184}\label{section-679}

\begin{otherlanguage}{persian}

و خود آن خبيث مشرك بحضرت ابهى عريضه معروض داشته و حال بخطّ خود او موجود
است و نوشته آنچه ما بكنيم باسم شما تمام ميشود چه كه شما رئيس هستيد و در
مدينه كبير مشافهه بيكى از اصحاب گفته كه تا حال آنچه از جعليات گفته‌‌ام
نقلى نبوده من بعد جعلها خواهم كرد و فسادها خواهم برپا نمود چنانچه نموده
مخصوص گفته كه جمال ابهى آدم باطراف بجهت فساد فرستاده‌اند و خود شما
ميدانيد كه آنچه گفته كذب صرف بوده و مع‌ذلك صادقش شمرده مع آنكه در ليالى و
ایّام كذب تعليم شما مينمود كه همچه بگو و همچه بنويس و شما عامل بوده و
هستيد

\end{otherlanguage}

Va khod ān khabīs-i mushrik bi-Ḥaḍrat-i Abhā ʿarīḍa maʿrūḍ dāshta va ḥāl
bi-khaṭṭ-i khod-i ū mawjūd ast va nivishta ānchi mā bikunīm bi-ism-i
shumā tamām mīshavad chi ki shumā raʾīs hastīd va dar Madīna-yi Kabīr
mushāfahatan bi-yikī az aṣḥāb gufta ki tā ḥāl ānchi az jaʿliyyāt
gufta-am naqlī nabūda min baʿd jaʿlhā khwāham kard va fasādhā khwāham
barpā namūd chunānchi namūda makhṣūṣ gufta ki Jamāl-i Abhā ādam bi-aṭrāf
bi-jihat-i fasād firistāda-and va khod-i shumā mīdānīd ki ānchi gufta
kiẕb-i ṣirf būda va maʿa-dhālik ṣādiqash shumurda maʿa ānki dar layālī
va ayyām kiẕb taʿlīm-i shumā mīnamūd ki hamcha bigū va hamcha binivīs va
shumā ʿāmil būda va hastīd.

That wicked man who associates partners with God himself submitted a
letter to His Most Glorious Presence, which still exists in his own
handwriting, and wrote: ``Whatever we do is attributed to you, because
you are the chief.'' In the Great City he told one of the companions
face to face: ``The fabrications I have uttered so far have been nothing
of consequence. From now on I will invent falsehoods and stir up
disorders''---as indeed he has done. He specifically said that the Abhá
Beauty had sent people to various regions to stir up disorder. You
yourself know that everything he said was sheer falsehood, yet you have
considered him truthful, even though day and night he taught you to lie,
saying, ``Say this and write that,'' and you have acted, and still act,
accordingly.

\subsection{185}\label{section-680}

\begin{otherlanguage}{persian}

و آنقدر ملتفت نيستيد كه آنى تفكّر نمائى كه ادرنه خلف كوه قاف نيست للّه چند
روزى بروم و بچشم خود به‌بينم حقّ و دونش در آن ارض موجود باين موفّق نشده و
بمفتريات مشغول گشته اليوم كه لقا ممكن بود بآن فايز نشدى بارى اين نعمت بی‌
زوال هميشه بدست نيايد و اين ربيع الهى هميشه ظاهر نه عنقريب است كه كلّ
محروم شوند اين عبد از خدا ميطلبد كه در دنيا و آخرت اين عباد را از ظلّ
فناء بابش محروم نفرمايد و لازال در حول بيت طايف باشيم و بانوارش مشرق و
از فواكه عرفانش مرزوق و از كوثر رضايش مشروب دنيا را بطالبانش واگذاشته‌ايم
از دنيا و عقبی‌ جز دوست نخواسته و نخواهيم انّه لخير معطى و احسن مجيب

\end{otherlanguage}

Va ānqadr multafit nīstīd ki ānī tafakkur namāʾī ki Adrianih khalf-i
Kūh-i Qāf nīst li-llāh chand rūzī biravam va bi-chashm-i khod bibīnam
Ḥaqq va dūnash dar ān arḍ mawjūd bi-īn muvaffaq nashuda va bi-muftariyāt
mashghūl gashta al-yawm ki liqā mumkin būd bi-ān fāʾiz nashudī bārī īn
niʿmat-i bī-zavāl hamīsha bi-dast nayāyad va īn rabīʿ-i ilāhī hamīsha
ẓāhir na ʿan-qarīb ast ki kull maḥrūm shavand īn ʿabd az Khudā mīṭalabad
ki dar dunyā va ākhirat īn ʿibād rā az ẓill-i finā-yi bābash maḥrūm
nafarmāyad va lā-zāl dar ḥawl-i bayt ṭāʾif bāshīm va bi-anvārash mushriq
va az favākih-i ʿirfānash marzūq va az kawthar-i riḍāyash mashrūb dunyā
rā bi-ṭālibānash vāgudhāshta-īm az dunyā va ʿuqbā juz Dūst nakhwāsta va
nakhwāhīm innahu la-khayru muʿṭin wa-aḥsanu mujīb.

You are not attentive enough to reflect for a moment: ``Adrianople is
not beyond Mount Qáf. For God's sake, let me go for a few days and see
with my own eyes. The True One and those besides Him are present in that
land.'' You have not succeeded in doing this and have instead occupied
yourself with fabrications. Today, when meeting Him was possible, you
did not attain it. In any event, this unfailing bounty will not always
be within reach, nor will this divine springtime always be manifest.
Soon all will be deprived. This servant asks God not to deprive these
servants, in this world or the next, of the shelter of the courtyard of
His gate, and that we may always circle around the House, shine with His
lights, be nourished by the fruits of His knowledge, and drink from the
Kawthar of His good pleasure. We have left the world to those who seek
it. From this world and the next we have sought, and will seek, none but
the Friend. He is indeed the best of givers and the most gracious of
responders.

\subsection{186}\label{section-681}

\begin{otherlanguage}{persian}

از اهل انصاف نيستيد كه سهل است در كوى منصفين هم مرور ننموده و قدم
نگذاشته چه كه اگر قدمى رفته بودى نمينوشتى كه تقليد نيست و حال آنكه ابدا
از بيان اطّلاع نداشته و ندارى و آنچه ذكر نموده محض تقليد بوده كه از سيّد
محمّد خبيث كاذب مفسد شنيده تا حين كه از روايح انصاف محروم بوده و لكن از
خدا ميطلبم كه از شمال ظلم و اعتساف بيمين عدل و انصاف راجع شوى اگرچه
بسيار مشكل است چه كه نحل انگبين جنس ديگر و زنبور جنس ديگر و لكن نزد قدرت
الهى هيچ چيز محال نبوده و نيست و آن محاليّت در خود اشیاء موجود و مشهود

\end{otherlanguage}

Az ahl-i inṣāf nīstīd ki sahl ast dar kūy-i munṣifīn ham murūr nanamūda
va qadam nagudhāshta chi ki agar qadamī rafta būdī naminivishtī ki
taqlīd nīst va ḥāl ānki abadan az Bayān iṭṭilāʿ nadāshta va nadārī va
ānchi dhikr namūda maḥḍ-i taqlīd būda ki az Sayyid Muḥammad-i khabīs-i
kādhib-i mufsid shinīda tā ḥīn ki az ravāyiḥ-i inṣāf maḥrūm būda va
lākin az Khudā mīṭalabam ki az shimāl-i ẓulm va iʿtisāf bi-yamīn-i ʿadl
va inṣāf rājiʿ shavī agarchi bisyār mushkil ast chi ki naḥl-i angabīn
jins-i dīgar va zanbūr jins-i dīgar va lākin nazd-i qudrat-i ilāhī hīch
chīz muḥāl nabūda va nīst va ān muḥāliyyat dar khod-i ashyā mawjūd va
mashhūd.

Not only are you not among the fair-minded; you have never even passed
through their quarter or set foot there. For had you taken a single step
there, you would not have written that this was not blind imitation,
when you have never had, and still have no, knowledge of the Bayán.
Everything you have mentioned has been mere imitation of what you heard
from that wicked, lying mischief-maker, Siyyid Muḥammad. Until now you
have been deprived of the fragrances of fairness, but I ask God that you
may turn from the left hand of oppression and injustice to the right
hand of justice and fairness. This is very difficult, for the honeybee
is one kind and the wasp another. Yet nothing has ever been, or is,
impossible for divine power; that impossibility exists and is evident in
the things themselves.

\subsection{187}\label{section-682}

\begin{otherlanguage}{persian}

فكّر لتعرف ذكر مميّزى نقطه اولى نموده بودى اگر شما ناظر بمميّز بوديد چرا آن
قول مميّز را اطاعت ننمودى كه در جميع بيان فرموده كه كلّ در حين ظهور بايد
از كلّ شیء و عن كلّ ما عندهم منقطع شوند و بنفس ظهور ناظر گردند چه كه آنچه
از قبل بآن متمسّك بوده‌اند در آن يوم كفايت نمينمايد الّا بتصديق باو بشأنيكه
ميفرمايند سطرى از آيات آن شمس ظهور اعظم است از آنچه از قبل نازل شده و
مع‌ذلك بذكر اثبات كلمات خلق او مشغول گشته و حال آنكه اطّلاع از كلمات آن
نفس نداشته و ندارى و آنچه نزدش بوده رشحى از اين بحر اعظم بوده كه باو
عنايت شده و بعد باعراضش از حقّ اخذ شده نفوسيكه بر جميع امور مطّلعند الآن
موجود جميع اين عباد را بی‌دين مدان و بی‌ مذهب مشمر فواللّه كه للّه هجرت
نموده‌ايم و الى اللّه ساير گشته‌ايم اگر چه معلوم است كه آن مفسد خبيث بشماچه
القاء نموده و جميع اين عباد را بدنيا نسبت داده خود شما فكر نمائيد كجا
دنيا بوده و بكه دنيا رسيده

\end{otherlanguage}

Fakkir li-taʿrif dhikr-i mumayyizī-yi Nuqṭa-yi Ūlā namūda būdī agar
shumā nāẓir bi-mumayyiz būdīd chirā ān qawl-i mumayyiz rā iṭāʿat
nanamūdī ki dar jamīʿ-i Bayān farmūda ki kull dar ḥīn-i ẓuhūr bāyad az
kull-i shayʾ va ʿan kulli mā ʿindahum munqaṭiʿ shavand va bi-nafs-i
ẓuhūr nāẓir gardand chi ki ānchi az qabl bi-ān mutamassik būda-and dar
ān yawm kifāyat namīnamāyad illā bi-taṣdīq bi-ū bi-shaʾnī ki mīfarmāyand
saṭrī az āyāt-i ān Shams-i Ẓuhūr aʿẓam ast az ānchi az qabl nāzil shuda
va maʿa-dhālik bi-dhikr-i ithbāt-i kalimāt-i khalq-i ū mashghūl gashta
va ḥāl ānki iṭṭilāʿ az kalimāt-i ān nafs nadāshta va nadārī va ānchi
nazdash būda rashḥī az īn Baḥr-i Aʿẓam būda ki bi-ū ʿināyat shuda va
baʿd bi-iʿrāḍash az Ḥaqq akhdh shuda nufūsī ki bar jamīʿ-i umūr
muṭṭaliʿand al-ān mawjūd jamīʿ-i īn ʿibād rā bī-dīn madān va bī-madhhab
mashmur fa-wa-llāh ki li-llāh hijrat namūda-īm va ilā Allāh sāyir
gashta-īm agar chi maʿlūm ast ki ān mufsid-i khabīs bi-shumā chi ilqā
namūda va jamīʿ-i īn ʿibād rā bi-dunyā nisbat dāda khod-i shumā fikr
namāʾīd kujā dunyā būda va bi-ki dunyā rasīda.

Reflect, so that you may understand. You had mentioned the Primal
Point's discerning judgment. If you looked to that discerning Judge, why
did you not obey His statement, repeated throughout the Bayán, that at
the time of the Manifestation all must detach themselves from everything
and from all that they possess, and look to the Manifestation Himself?
For whatever they previously clung to will not suffice on that Day
unless they acknowledge Him. So emphatically does He speak that He says
a single line of the verses of that Sun of Revelation is greater than
everything previously revealed. Nevertheless, you have occupied yourself
with establishing the validity of the words of one of His creatures,
although you have never had, and still have no, knowledge of that
person's words. Whatever he possessed was a sprinkling from this Most
Great Ocean, bestowed upon him and then taken away when he turned from
the True One. Those who are informed of all these matters are now
present. Do not regard all these servants as irreligious or reckon them
faithless. For, by God, we have emigrated for God's sake and journeyed
toward God. It is clear, however, what that wicked mischief-maker has
insinuated to you, attributing worldliness to all these servants.
Reflect for yourself: where has this worldly gain been, and to whom has
it come?

\subsection{188}\label{section-683}

\begin{otherlanguage}{persian}

و اگر اين عباد يك كلمه صدق ظاهره از آن نفوس ميشنيدند باز ممكن بود كه
بعضى از اعمالشان را تأويل كنيم متّصلا بمكر و كذب و افترا مشغول چون ملاحظه
نمودند كه اين عباد بفضل اللّه بر حقّ ثابت و مستقيمند تدبير نموده‌اند كه
شايد بنسبت بدنيا و امورات كذبه اين عباد را تضييع نمايند و مظلوميّتى از
براى خود ثابت كنند بحقّ نقطه وجود روح من فى ملكوت الغيب و الشّهود فداه كه
ظلمى نموده‌اند كه از اوّل آفرينش ابداع تا بحال احدى ارتكاب ننموده الا لعنة
اللّه على الكاذبين و المفترين و الظّالمين

\end{otherlanguage}

Va agar īn ʿibād yik kalima-yi ṣidq-i ẓāhira az ān nufūs mīshinīdand bāz
mumkin būd ki baʿḍī az aʿmālashān rā taʾvīl kunīm muttaṣilan bi-makr va
kiẕb va iftirā mashghūl chūn mulāḥaẓa namūdand ki īn ʿibād bi-faḍli
Allāh bar Ḥaqq thābit va mustaqīmand tadbīr namūda-and ki shāyad
bi-nisbat bi-dunyā va umūrāt-i kādhiba īn ʿibād rā taḍyīʿ namāyand va
maẓlūmiyyatī az barāy-i khod thābit kunand bi-ḥaqq-i Nuqṭa-yi Vujūd rūḥu
man fī malakūti al-ghaybi wa-al-shuhūdi fidāh ki ẓulmī namūda-and ki az
avval-i āfarīnish-i ibdāʿ tā bi-ḥāl aḥadī irtikāb nanamūda alā laʿnatu
Allāhi ʿalā al-kādhibīna wa-al-muftarīna wa-al-ẓālimīn.

If these servants had heard even one manifestly truthful word from those
people, it would still have been possible for us to interpret some of
their deeds favorably. They are continually occupied with deceit, lies,
and slander. When they saw that, by God's grace, these servants remained
firm and steadfast in the True One, they devised a scheme, hoping to
discredit these servants by attributing worldliness and false matters to
them, and to establish themselves as victims of oppression. By the Point
of existence---may the spirit of all in the kingdom of the unseen and
the seen be sacrificed for Him---they have committed an injustice that
no one has committed from the beginning of creation until now. Indeed,
God's curse rests upon the liars, the slanderers, and the oppressors.

\subsection{189}\label{section-684}

\begin{otherlanguage}{persian}

و اينكه ذكر آيات ميرزا يحيى نموده بودى نقطه اولى روح ماسواه فداه
ميفرمايد كه از حين غروب شمس حقيقت ديگر احدى آيات اللّه نخواهد شنيد الّا در
ظهور بعد چنانچه حال هم خود آن جمال مبارك همين را بنفسه ميفرمايد چه كه
آيات دونش ماداميكه مدلّ بر آيات الهى است بحقّ نسبت داده ميشود و وجودش مثل
وجود مرايا است كه از خود تحقّقى نداشته و ندارند و لكن مادامى كه بشمس عزّ
صمدانى مقابل اشراق در او ظاهر و بعد از انحراف محو و اينهم از آن بيان است
كه ميفرمايند در مظاهر ظهور ديده نميشود مگر نفس ظهور لذا در جميع مظاهر
ظهور چه در كور فرقان و چه در كور انجيل و چه در كور زبور و چه در كور
تورية ديده نميشود مگر نفس ظهور و هم‌چنين در سكّان بيان ديده نميشود مگر
ظهور نقطه بيان و اين مادام كه در ظلّ بيانند و العياذ باللّه اگر نفسى خارج
شد ديده نميشود در او مگر شيطان و افعال و اعمال او

\end{otherlanguage}

Va īnki dhikr-i āyāt-i Mīrzā Yaḥyā namūda būdī Nuqṭa-yi Ūlā rūḥ-i mā
sivāh fidāh mīfarmāyad ki az ḥīn-i ghurūb-i Shams-i Ḥaqīqat dīgar aḥadī
āyāt-i Allāh nakhwāhad shinīd illā dar ẓuhūr-i baʿd chunānchi ḥāl ham
khod-i ān Jamāl-i Mubārak hamīn rā bi-nafsih mīfarmāyad chi ki āyāt-i
dūnash mādāmī ki mudill bar āyāt-i ilāhī ast bi-Ḥaqq nisbat dāda
mīshavad va vujūdash mithl-i vujūd-i marāyā ast ki az khod taḥaqquqī
nadāshta va nadārand va lākin mādāmī ki bi-Shams-i ʿIzz-i Ṣamadānī
muqābil ishrāq dar ū ẓāhir va baʿd az inḥirāf maḥv va īnham az ān bayān
ast ki mīfarmāyand dar maẓāhir-i ẓuhūr dīda namīshavad magar nafs-i
ẓuhūr lidhā dar jamīʿ-i maẓāhir-i ẓuhūr chi dar kawr-i Furqān va chi dar
kawr-i Injīl va chi dar kawr-i Zabūr va chi dar kawr-i Tawrāt dīda
namīshavad magar nafs-i ẓuhūr va hamchunīn dar sukkān-i Bayān dīda
namīshavad magar ẓuhūr-i Nuqṭa-yi Bayān va īn mādām ki dar ẓill-i
Bayānand va al-ʿiyādhu bi-llāh agar nafsī khārij shud dīda namīshavad
dar ū magar shayṭān va afʿāl va aʿmāl-i ū.

As for your mention of the verses of Mírzá Yaḥyá, the Primal Point---may
the spirit of all besides Him be sacrificed for Him---says that, from
the setting of the Sun of Truth, no one will again hear the verses of
God until the next Manifestation. Now that Blessed Beauty Himself
likewise says this in person. For the verses of anyone besides Him are
attributed to the True One only so long as they point to the divine
verses. Their existence is like that of mirrors, which have never had,
and do not have, any independent reality. As long as a mirror faces the
Sun of everlasting glory, its radiance appears within it; once it turns
away, that radiance vanishes. This, too, follows from the statement in
which He says that nothing is seen in the manifestations of Revelation
but Revelation itself. Thus, in all the manifestations of Revelation,
whether in the cycle of the Qur'an, the Gospel, the Psalms, or the
Torah, nothing is seen but Revelation itself. Likewise, in the
inhabitants of the Bayán nothing is seen but the revelation of the Point
of the Bayán, and this is so only while they remain beneath the shelter
of the Bayán. Should anyone leave it---God forbid---nothing is seen in
him but Satan and his deeds and actions.

\subsection{190}\label{section-685}

\begin{otherlanguage}{persian}

و اين مسلّم است كه هر نفسى در اين ظهور توقّف نمايد از بيان خارج و از حصن
معانى محروم چنانچه خود نقطه بيان روح ماسواه فداه باين تصريح فرموده‌اند
لذا آيات هر يك از مظاهر ظهور ذكر آيتيّت بر ايشان صادق ماداميكه در ظلّ
ذوالآيه مستقرّ و ساكنند و بعد از تجاوز از اهل نار محسوب و همان تجاوز نفس
نار او است لو انتم من النّاظرين

\end{otherlanguage}

Va īn musallam ast ki har nafsī dar īn ẓuhūr tavaqquf namāyad az Bayān
khārij va az ḥiṣn-i maʿānī maḥrūm chunānchi khod-i Nuqṭa-yi Bayān rūḥ-i
mā sivāh fidāh bi-īn taṣrīḥ farmūda-and lidhā āyāt-i har yik az
maẓāhir-i ẓuhūr dhikr-i āyatiyyat bar īshān ṣādiq mādāmī ki dar ẓill-i
Dhu al-Āya mustaqarr va sākinand va baʿd az tajāvuz az ahl-i nār maḥsūb
va hamān tajāvuz nafs-i nār-i ū ast law antum mina al-nāẓirīn.

It is certain that anyone who hesitates in this Revelation is outside
the Bayán and deprived of the stronghold of inner meanings, as the Point
of the Bayán Himself---may the spirit of all besides Him be sacrificed
for Him---has expressly stated. Thus the verses of each of the
manifestations of Revelation may truly be called signs only so long as
they remain established and abiding beneath the shelter of the Possessor
of the Sign. Once they transgress that boundary, they are counted among
the people of the fire, and that very transgression is their fire, if
you are among those who see.

\subsection{191}\label{section-686}

\begin{otherlanguage}{persian}

بارى اظلال و اشباح را شمس فرض مگيريد و بتحريف معانى كلمات الهيّه مشغول
نشويد چنانچه از قبل مشغول شده‌ايد از جمال نقطه بيان كه الآن مشرق و مشهود
است حيا كنيد بخلق او از او اعراض منمائيد كلماتش رحمتى بوده از نزد او بر
عباد او او را سيف خود قرار مدهيد و بر نفس مباركش ضربت مزنيد چنانچه اهل
فرقان بكلمات رسول اللّه استدلال نموده بر كينونتش وارد آوردند آنچه وارد
آوردند

\end{otherlanguage}

Bārī aẓlāl va ashbāḥ rā shams farḍ magīrīd va bi-taḥrīf-i maʿānī-yi
kalimāt-i ilāhiyya mashghūl nashavīd chunānchi az qabl mashghūl shuda-īd
az Jamāl-i Nuqṭa-yi Bayān ki al-ān mushriq va mashhūd ast ḥayā kunīd
bi-khalq-i ū az ū iʿrāḍ manamāʾīd kalimātash raḥmatī būda az nazd-i ū
bar ʿibād-i ū ū rā sayf-i khod qarār madihīd va bar nafs-i mubārakash
ḍarbat mazanīd chunānchi ahl-i Furqān bi-kalimāt-i Rasūl-i Allāh
istidlāl namūda bar kaynūnatash vā rid āvardand ānchi vā rid āvardand.

In short, do not mistake shadows and phantoms for the sun, and do not
occupy yourselves with distorting the meanings of the divine words, as
you have done before. Feel shame before the Beauty of the Point of the
Bayán, who is now shining and visible. Do not turn away from Him because
of His creatures. His words have been a mercy from Him to His servants.
Do not make them your sword and strike His blessed person with them, as
the people of the Qur'an, invoking the words of the Messenger of God as
their proof, inflicted upon His being what they inflicted.

\subsection{192}\label{section-687}

\begin{otherlanguage}{persian}

از جميع اين امور گذشته بگو اى غافلان مگر سمع و بصر نداريد خود آيات را
بخوانيد فواللّه الّذى لا اله الّا هو بلوحى از آيات بديع مقابله نمينمايد
آنچه استماع نموده‌ايد آيا مشاهده نشد كه بالغ كامل طفل را چه قدر مرحمت
مينمايد بشأنيكه افعال او را تعريف ميكند اين تعريف و توصيف نظر بملاطفت
كبرى است و اگر همان فعل از عاقلى ظاهر شود البتّه قبيح بوده و خواهد بود در
حقّ صبی‌ افعال او در رتبه او ممدوح و از رتبه كه تجاوز نمود قبيح و غير
مقبول درست ملتفت شويد كه چه ذكر ميشود اگر چه در ناس شعور ملاحظه نشده و
ادراك مشاهده نگشته و لكن اگر شما و امثال شما عارف نشويد عبادى خلق خواهد
شد كه كوثر معانى را از كأس كلمات الهى بياشامند و بمقصود مطّلع گردند

\end{otherlanguage}

Az jamīʿ-i īn umūr gudhashta bigū ay ghāfilān magar samʿ va baṣar
nadārīd khod-i āyāt rā bikhwānīd fa-wa-llāhi alladhī lā ilāha illā huwa
bi-lawḥī az āyāt-i badīʿ muqābila namīnamāyad ānchi istimāʿ namūda-īd
āyā mushāhada nashud ki bāligh-i kāmil ṭifl rā chi qadr marḥamat
mīnamāyad bi-shaʾnī ki afʿāl-i ū rā taʿrīf mīkunad īn taʿrīf va tawṣīf
naẓar bi-mulāṭafat-i kubrā ast va agar hamān fiʿl az ʿāqilī ẓāhir shavad
albatta qabīḥ būda va khwāhad būd dar ḥaqq-i ṣabī afʿāl-i ū dar rutba-yi
ū mamdūḥ va az rutba ki tajāvuz namūd qabīḥ va ghayr-i maqbūl durust
multafit shavīd ki chi dhikr mīshavad agar chi dar nās shuʿūr mulāḥaẓa
nashuda va idrāk mushāhada nagashta va lākin agar shumā va amthāl-i
shumā ʿārif nashavīd ʿibādī khalq khwāhad shud ki kawthar-i maʿānī rā az
kaʾs-i kalimāt-i ilāhī biyāshāmand va bi-maqṣūd muṭṭaliʿ gardand.

Apart from all these matters, say: ``O heedless ones! Have you no
hearing or sight? Read the verses yourselves. By God, besides whom there
is no other god, all that you have heard cannot compare with a single
tablet of wondrous verses. Have you not seen how kindly a fully mature
adult treats a child, even praising his actions? Such praise and
commendation arise from the utmost kindness. If the same action were
performed by a person of understanding, it would certainly be, and
remain, blameworthy. A child's actions are praiseworthy at his own
stage, but once he passes beyond that stage, they are blameworthy and
unacceptable. Pay close attention to what is being said. Although no
awareness has been discerned among the people and no understanding
witnessed, if you and those like you fail to understand, servants will
be created who will drink the Kawthar of inner meanings from the cup of
the divine words and comprehend what is intended.''

\subsection{193}\label{section-688}

\begin{otherlanguage}{persian}

اكثرى اليوم رضيع عنداللّه محسوبند چه كه از ثدى غفلت مياشامند و هنوز ایّام
فطام ايشان نرسيده چه رسد ببلوغ و اگر بالغ يافت شود از آيات عزّ صمدانى كه
در اين ظهور ظاهر شده حلاوت تسنيم معرفت ربّانى و سلسبيل قدس صمدانى را
ميچشد و ادراك مينمايد و دونش را مثل سراب بقيعه مشاهده ميكند قل خافوا عن
اللّه و ل اتذكروا دونه عنده ول اتقاسوا ما سواه بنفسه ول اتظنّوا فى حقّه ما
لا يليق و لا ينبغى لكم اتّقوا اللّه يا قوم و لا تعقّبوا الّذين كفروا باللّه
فى ایّامه و جادلوا ببرهانه و غصبوا حقّه و تجاوزوا عن عدله و ظلموا و
استظلموا اولئك يلعنهم كلّ الذّرّات ان انتم تفقهون

\end{otherlanguage}

Aktharī al-yawm raḍīʿ ʿinda Allāh maḥsūband chi ki az thady-i ghaflat
mīyāshāmand va hanūz ayyām-i fiṭām-i īshān narasīda chi rasad bi-bulūgh
va agar bāligh yāft shavad az āyāt-i ʿizz-i ṣamadānī ki dar īn ẓuhūr
ẓāhir shuda ḥalāvat-i tasnīm-i maʿrifat-i rabbānī va salsabīl-i quds-i
ṣamadānī rā mīchashad va idrāk mīnamāyad va dūnash rā mithl-i sarāb
bi-qīʿa mushāhada mīkunad qul khāfū ʿani Allāhi wa-{[}corrupted text{]}
dūnahu ʿindahu {[}corrupted text{]} mā siwāhu bi-nafsihi {[}corrupted
text{]} fī ḥaqqihi mā lā yalīqu wa-lā yanbaghī lakum ittaqū Allāha yā
qawmu wa-lā taʿaqqabū alladhīna kafarū bi-llāhi fī ayyāmihi wa-jādalū
bi-burhānihi wa-ghaṣabū ḥaqqahu wa-tajāwazū ʿan ʿadlihi wa-ẓalamū
wa-istaẓlamū ulāʾika yalʿanuhum kullu al-dharrāti in antum tafqahūn.

Most people today are accounted nursing infants in God's sight, for they
drink from the breast of heedlessness, and their time of weaning has not
yet come, much less that of maturity. If a mature person is found, he
tastes and perceives, in the verses of everlasting glory manifested in
this Revelation, the sweetness of the Tasním of divine knowledge and the
Salsabíl of everlasting holiness. He sees all besides Him as a mirage on
a plain. Say: ``Fear God, and {[}corrupted original text{]} anyone
besides Him in His presence, {[}corrupted original text{]} anything
other than Him with His own Self, {[}corrupted original text{]}
concerning Him what is unfitting and unseemly for you. Fear God, O
people, and do not follow those who disbelieved in God in His days,
disputed His proof, usurped His right, transgressed His justice, and
committed oppression while claiming to be oppressed. These are the ones
whom every atom curses, if you understand.''

\subsection{194}\label{section-689}

\begin{otherlanguage}{persian}

و ديگر نوشته بوديد آيات ميرزا يحيى در اطراف منبسط است و مدّعى مقام حقّيّت
خود ميباشد به بيّنه آيات خود قطع نظر از نصّ ربّ اعلى كه اعظم دليل است انتهى
اين دعوى از اصل باطل بوده چه كه اگر خلق مستقلّا ادّعاى آيات نمايند باطل
بوده و خواهند بود و همان ادّعا دليل بر بطلان است بعينه مثل آن است كه عكس
در مرآت بگويد كه من شمس سما هستم خود تو نوشته كه او مرآت است لو نصدّقك
شكّى نيست كه وجود مرايا لانطباع تجلّى شمس بوده و خواهد بود مع اين رتبه
ادّعاى رتبه فوق باطل بوده و خواهد بود تناقض در اينمقام است نه در آنمقام
كه ذكر نموده چه كه با وصف مرآتيّت خواسته رتبه شمس ثابت نمائى آيا تجلّى كه
در مرآت ديده ميشود اگر دعوى نمايد كه من بنفس خود قائمم صادق است لا فو‌نفس
اللّه المهيمن القيّوم

\end{otherlanguage}

Va dīgar nivishta būdīd āyāt-i Mīrzā Yaḥyā dar aṭrāf munbasiṭ ast va
muddaʿī-yi maqām-i ḥaqqiyyat-i khod mībāshad bi-bayyina-yi āyāt-i khod
qaṭʿ-i naẓar az naṣṣ-i Rabb-i Aʿlā ki aʿẓam dalīl ast intahā īn daʿvā az
aṣl bāṭil būda chi ki agar khalq mustaqillan iddiʿā-yi āyāt namāyand
bāṭil būda va khwāhand būd va hamān iddiʿā dalīl bar buṭlān ast
bi-ʿaynih mithl-i ān ast ki ʿaks dar mirʾāt bigūyad ki man Shams-i Samā
hastam khod-i tu nivishta ki ū mirʾāt ast law nuṣaddiquka shakkī nīst ki
vujūd-i marāyā li-inṭibāʿ-i tajallī-yi shams būda va khwāhad būd maʿa īn
rutba iddiʿā-yi rutba-yi fawq bāṭil būda va khwāhad būd tanāquḍ dar īn
maqām ast na dar ān maqām ki dhikr namūda chi ki bā vaṣf-i mirʾātiyyat
khwāsta rutba-yi shams thābit namāʾī āyā tajallī ki dar mirʾāt dīda
mīshavad agar daʿvā namāyad ki man bi-nafs-i khod qāʾimam ṣādiq ast lā
fa-wa-nafsi Allāhi al-muhaymini al-qayyūm.

You had also written: ``The verses of Mírzá Yaḥyá are spread abroad, and
he claims the station of truth for himself on the evidence of his own
verses, quite apart from the explicit text of the Most Exalted Lord,
which is the greatest proof.'' Here ends your statement. This claim is
fundamentally invalid, for if creatures independently lay claim to
verses, they have been, and will be, false; that very claim is proof of
their falsity. It is exactly as though the reflection in a mirror were
to say, ``I am the sun in the sky.'' You yourself have written that he
is a mirror. Even if we accept your word, there is no doubt that mirrors
have existed, and will exist, to receive the imprint of the sun's
radiance. While occupying this station, a claim to the station above it
has been, and will be, invalid. The contradiction lies here, not where
you have alleged it, for while describing him as a mirror, you have
sought to establish for him the station of the sun. If the radiance seen
in a mirror claims, ``I subsist by myself,'' is it truthful? No, by the
very Self of God, the All-Protecting, the Self-Subsisting!

\subsection{195}\label{section-690}

\begin{otherlanguage}{persian}

و ديگر عجب است كه انتشار آيات موهومى را كه نشناخته حجّت اخذ نموده و دليل
بر حقّيّت او ذكر نموده مع آنكه آن آيات در رتبه خود او بوده كلّما اظهر من
الخلق فى حدّه و مقامه و از آياتيكه مهيمن است بر كلّ من فى السّموات و الأرض
و مشرق است مثل شمس از افق ابهى محتجب مانده اگر انتشار است كه انتشار اين
آيات بشأنى شده كه فوق آن متصوّر نه خف عن اللّه الّذى خلقك و لا تذكر آيات
احد تلقآء ذكر آياته لو تقابل شمس الّتى كانت فى صور المرايا شمس الّتى كانت
مشرقة عن افق القدس لتقابل آياته بآيات اللّه المهيمن العزيز القيّوم و اين
در وقتى است كه مرآت حقيقى باشد و اگر مرآت امثال شما بوده كه ذكرى نزد حقّ
نداشته و نخواهد داشت فافتح اذنك لتسمع من كلّ حرف من آيات ربّك انّه لا اله
الّا انا ان يا كلّشیء قد خلقتم بنفسى لنفسى ایّاکم ان تنكرون و لو تطهّر بصرك
لتشهد كلّ كلمات ربّك على هيكل الرّوح و ينطقن بأنّه لا اله الّا انا المقتدر
المهيمن العزيز المحبوب قل فاستغفر عمّا ارتكبت يا ايّها الغافل لانّ ما عندك
لن يذكر عند اللّه و ما عرفته يرجع اليك كذلك كان الأمر و لكن انت لا تكون
من الشّاعرين

\end{otherlanguage}

Va dīgar ʿajab ast keh enteshār-e āyāt-e mowhūmī rā keh nashenākhteh
ḥojjat akhdh namūdeh va dalīl bar ḥaqqiyyat-e ū dhekr namūdeh maʿa ānkeh
ān āyāt dar rotbeh-ye khod-e ū būdeh kullu mā aẓhara min al-khalq fī
ḥaddihi va maqāmihi va az āyātī keh mohaymen ast bar kulli man fī
al-samāwāt va al-arḍ va moshreq ast methl-e shams az ofoq-e Abhā
moḥtajeb māndeh agar enteshār ast keh enteshār-e īn āyāt beh shaʾnī
shodeh keh fowq-e ān motaṣavvar nah khaf ʿan Allāh alladhī khalaqaka va
lā tadhkur āyāti aḥadin tilqāʾa dhikri āyātihi law tuqābilu shams allatī
kānat fī ṣuwar al-marāyā shams allatī kānat mushriqatan ʿan ufuq al-quds
la-tuqābilu āyātahu bi-āyāt Allāh al-Muhaymin al-ʿAzīz al-Qayyūm va īn
dar vaqtī ast keh merʾāt ḥaqīqī bāshad va agar merʾāt amthāl-e shomā
būdeh keh dhekrī nazd-e Ḥaqq nadāshteh va nakhāhad dāsht fa-ftaḥ
udhunaka li-tasmaʿa min kulli ḥarfin min āyāti rabbika innahu lā ilāha
illā anā in yā kulla shayʾin qad khuliqtum bi-nafsī li-nafsī iyyākum an
tunkirūn va law tuṭahhir baṣaraka li-tashhada kulla kalimāti rabbika
ʿalā haykal al-rūḥ va yanṭiqna bi-annahu lā ilāha illā anā al-Muqtadir
al-Muhaymin al-ʿAzīz al-Maḥbūb qul fa-staghfir ʿammā irtakabta yā ayyuhā
al-ghāfil li-anna mā ʿindaka lan yudhkara ʿinda Allāh va mā ʿaraftahu
yarjiʿu ilayka kadhālika kāna al-amr va lākin anta lā takūnu min
al-shāʿirīn.

It is also astonishing that you have taken the circulation of imaginary
verses that you have not recognized as proof, and cited it as evidence
of his truth, even though those verses belong to his own rank. Whatever
any creature brings forth is within its own limits and station. Yet you
have remained veiled from verses that hold sway over all who are in the
heavens and on earth and shine like the sun from the All-Glorious
horizon. If circulation is the criterion, these verses have circulated
to a degree beyond which nothing can be imagined. Fear God, who created
you, and do not mention anyone's verses in the presence of His verses.
If the sun appearing in the images of mirrors could rival the sun
shining from the horizon of holiness, then his verses could rival the
verses of God, the All-Protecting, the Almighty, the Self-Subsisting.
And this is only if it is a true mirror; if it is a mirror like you, it
has never had, and will never have, any standing before God. Open your
ear, then, to hear from every letter of your Lord's verses: ``There is
no God but Me. O all created things! You have been created by My own
Self for My own Self. Beware lest you deny Me.'' If you purify your
sight, you will behold all the words of your Lord in the form of the
Spirit, declaring: ``There is no God but Me, the All-Powerful, the
All-Protecting, the Almighty, the Beloved.'' Say: ``Ask forgiveness for
what you have done, you who are heedless, for what you possess will
never be accounted anything before God, and what you have recognized
returns to you. Such has been the matter, but you are not among those
who perceive.''

\subsection{196}\label{section-691}

\begin{otherlanguage}{persian}

نسائم خريف را بنسائم ربيع چه مناسبت و نفحات علّيين را بنفحات سجّين چه
مشابهت فو‌الّذى نفسى بيده اگر بصير بودى همان لوحى كه از سماء مشيّت الهى
اتماماً لحجّته نازل و مخصوص تو ارسال شد حرفى از آن تو را كافى بود و همان
لوح را ميزان كن با آنچه از قبل نازل شده لتشهد بانّ هذا قد نزل من لدن منزل
قديم اگر در دون كلمات اللّه تعريف و يا توصيف شود اين نظر باظهار عنايت و
تشويق نفوس است و على حدّهم و شأنهم ذكر ميشود لو يقترن المخلوق الى الخالق
و يقابل السّراج بالشّمس لتقترن آيات دونه بآياته و تقابلها

\end{otherlanguage}

Nasāʾem-e kharīf rā beh nasāʾem-e rabīʿ cheh monāsebat va nafaḥāt-e
ʿilliyyīn rā beh nafaḥāt-e sijjīn cheh moshābahat fa-wa-alladhī nafsī
bi-yadihi agar baṣīr būdī hamān lowḥī keh az samā-ye mashiyyat-e elāhī
itmāman li-ḥujjatihi nāzel va makhṣūṣ-e to ersāl shod ḥarfī az ān to rā
kāfī būd va hamān lowḥ rā mīzān kon bā āncheh az qabl nāzel shodeh
li-tashhada bi-anna hādhā qad nazala min ladun Munzilin Qadīm agar dar
dūn-e kalimāt-e Allāh taʿrīf va yā towṣīf shavad īn naẓar beh eẓhār-e
ʿenāyat va tashvīq-e nofūs ast va ʿalā ḥaddihim va shaʾnihim dhekr
mīshavad law yaqtarinu al-makhlūq ilā al-khāliq va yuqābalu al-sirāj
bi-al-shams la-taqtarinu āyātu dūnihi bi-āyātihi va tuqābiluhā.

What affinity do the breezes of autumn have with the breezes of spring,
and what resemblance do the breaths of the highest heaven bear to the
breaths of the nethermost prison? By Him in whose hand is my soul! Had
you possessed insight, a single letter of that Tablet which descended
from the heaven of the divine will to complete His proof, and was sent
expressly to you, would have sufficed you. Take that very Tablet as the
balance against which to weigh what was revealed before, so that you may
witness that this has descended from an eternal Revealer. If anything
other than the words of God is praised or commended, this is to show
favor and encourage souls, and it is spoken of within their limits and
according to their station. If the creature could be joined to the
Creator and the lamp could rival the sun, then the verses of anyone
other than Him could be joined to His verses and rival them.

\subsection{197}\label{section-692}

\begin{otherlanguage}{arabic}

اتّق اللّه و لاتفسد فى الأرض بعد اصلاحها و لاتكن من الّذينهم بذكر الورقة
احتجبوا عن الشّجر و افنانه و اغصانه و اوراقه و اثماره لو يكون ورقة حقيقيّة
چون از عين كدره نوشيده از عين عذبه سايغه سايله جاريه محروم مانده و قل
بسم اللّه و باللّه ثمّ ارفع يدك و خذ منه غرفة ثم اشربها تاللّه يستنير بها
قلبك و تطمئن بها نفسك و يبرد فؤادك و تكون من الموقنين

\end{otherlanguage}

Ittaq Allāh va lā tufsid fī al-arḍ baʿda iṣlāḥihā va lā takun min
alladhīnahum bi-dhikr al-waraqah iḥtajabū ʿan al-shajar va afnānihi va
aghṣānihi va awrāqihi va athmārihi law yakūnu waraqatan ḥaqīqiyyatan
chūn az ʿayn-e kadereh nūshīdeh az ʿayn-e ʿadhbeh-ye sāʾegheh-ye
sāʾeleh-ye jāriyeh maḥrūm māndeh va qul bi-sm Allāh va bi-Allāh thumma
irfaʿ yadaka va khudh minhu ghurfatan thumma ishrabhā ta-Allāh yastanīru
bihā qalbuka va taṭmaʾinnu bihā nafsuka va yabrudu fuʾāduka va takūnu
min al-mūqinīn.

Fear God, and do not spread corruption on earth after it has been set
right. Do not be among those who, through mention of the leaf, have
become veiled from the tree, its boughs, its branches, its leaves, and
its fruits---even if it is a true leaf. Because you have drunk from a
turbid spring, you have remained deprived of the sweet, refreshing,
flowing spring. Say, ``In the name of God and by God,'' then raise your
hand, take a handful from it, and drink it. By God! Through it your
heart will be illumined, your soul will find peace, your inner heart
will be cooled, and you will be among those who are certain.

\subsection{198}\label{section-693}

\begin{otherlanguage}{persian}

بدان اينكه نوشته مدّعى مقام حقّيّت خود است ببيّنه آيات خود هذا بغى منك و منه
على اللّه المهيمن القيّوم چه كه اين رتبه مخصوص مظاهر احديّه و اين ظهور بدع
رحمانيّه بوده و خواهد بود كه بآيات اللّه و بيّناته ظاهر شده و جز آيات منزله
و ظهورات قدرتيّه بامرى متمسّك نشده چنانچه در جميع اوراق بيان مسطور است
فافتح بصرك لتعرف و اينست مقام يا من دلّ على ذاته بذاته و هم‌‌چنين مقام
دليله آياته و وجوده اثباته و هم‌‌چنين اعرفوا اللّه باللّه چه كه بر او جايز
نيست كه بغير خود استدلال بر حقّیّت خود نمايد او بنفس خود معروف بوده و
خواهد بود و ماسواه باو معروف و بنسبتهم اليه مفتخر او است ميزان الهيّه و
صراط احديّه و شمس قدميّه و بحر اعظم و كلمه اتمّ و رمز مستور و بيت معمور و
كتاب مسطور جميع باين ميزان سنجيده ميشوند و او بذات خود كلّ بنفحه از او
مبعوث و او بنفس خود

\end{otherlanguage}

Bedān īnkeh neveshteh moddaʿī-ye maqām-e ḥaqqiyyat-e khod ast beh
bayyeneh-ye āyāt-e khod hādhā baghyun minka va minhu ʿalā Allāh
al-Muhaymin al-Qayyūm cheh keh īn rotbeh makhṣūṣ-e maẓāher-e aḥadiyyeh
va īn ẓohūr-e badʿ-e raḥmāniyyeh būdeh va khāhad būd keh beh āyāt-e
Allāh va bayyenātesh ẓāher shodeh va joz āyāt-e monzaleh va ẓohūrāt-e
qodratiyyeh beh amrī motamassek nashodeh chenāncheh dar jamīʿ-e owrāq-e
Bayān masṭūr ast fa-ftaḥ baṣaraka li-taʿrifa va īnast maqām-e yā man
dalla ʿalā dhātihi bi-dhātihi va hamchenīn maqām-e dalīluhu āyātuhu va
wujūduhu ithbātuhu va hamchenīn iʿrifū Allāh bi-Allāh cheh keh bar ū
jāyez nīst keh beh gheyr-e khod estedlāl bar ḥaqqiyyat-e khod namāyad ū
beh nafs-e khod maʿrūf būdeh va khāhad būd va māsavāh beh ū maʿrūf va
bi-nisbatihim ilayhi moftakher ū ast mīzān-e elāhiyyeh va ṣerāṭ-e
aḥadiyyeh va shams-e qedamiyyeh va baḥr-e aʿẓam va kalemeh-ye atamm va
ramz-e mastūr va bayt-e maʿmūr va ketāb-e masṭūr jamīʿ beh īn mīzān
sanjīdeh mīshavand va ū beh dhāt-e khod koll beh nafḥeh az ū mabʿūth va
ū beh nafs-e khod.

As for what you wrote, that he claims the station of truth for himself
on the evidence of his own verses, know that this is an act of rebellion
by you and by him against God, the All-Protecting, the Self-Subsisting.
For this rank has belonged, and will belong, exclusively to the
Manifestations of divine unity and this wondrous Revelation of the
All-Merciful, which has appeared with the verses of God and His clear
proofs and has appealed to nothing but revealed verses and
manifestations of power, as is recorded throughout the pages of the
Bayán. Open your eyes, then, so that you may recognize. This is the
station of ``O You who have pointed to Your own Essence through Your own
Essence,'' and likewise the station of ``His proof is His verses, and
His existence is His evidence,'' and likewise ``Know God through God.''
For it is not fitting for Him to appeal to anything other than Himself
as evidence of His truth. He has been, and will be, known through His
own Self, while all besides Him are known through Him and glory in their
relationship to Him. He is the divine Balance, the Path of unity, the
Sun of eternity, the Most Great Ocean, the Most Perfect Word, the Hidden
Mystery, the Frequented House, and the Inscribed Book. All are weighed
in this Balance, while He is weighed by His own Essence. All are raised
up by a breath from Him, while He is raised up by His own Self.

\subsection{199}\label{section-694}

\begin{otherlanguage}{persian}

قل يا ايّها البعيد تقول و لاتشعر و تذكر و لاتعرف اگر فى‌الجمله شعور ملاحظه
ميشد ذكرى از معانى مستوره كه در سرادق عرفان معنوى اين ظهور مستور است
مذكور ميامد تا بر جنّت لا ظل لها فى الجمله مطّلع ميشدى و لكن حيف است كه
لآلىٴ قلزم عزّ ربّانى را بدست سارقين داد و حور معانى مقصوره را در محضر
ظالمين مكشوف آورد چه كه بچشم و بصر و فؤاد و قلب خائن بوده و هستند مثل
آنكه عبارات حضرت ابهى روح ماسواه فداه را سرقت نموده بكلمات مجعوله خود
مركب نموده‌اند و باسم خود انتشار داده‌اند و غافل از اينكه صرّافان احديّه در
صد هزار هزار حجر يكدانه لؤلؤ را بشناسند و از او بمعدنش پى‌برند چه كه از
معدن آهك تسنيم قدس الهى جارى نشود و از منبع قطران سلسبيل حكمت ربّانى سارى
نگردد و عبادى هستند كه برايحه رضوان بمكمن رحمن پى برند چه كه رايحه طيّبه
هدايت نماينده است صاحبان شامّه اصفى را برضوان جمال علىّ ابهى البتّه در ملك
ابصر حديده ظاهر شوند و صرّافان احديّه بوجود آيند و جواهر كلماتش را بجان
خريدار شوند و از كون و امكان در سبيلش بگذرند و جان در رهش ايثار نمايند
باسمش از رسم دنيا منقطع گردند و بذكرش از ذكر عالمين بی‌ نياز و بيادش از
ياد دونش فارغ

\end{otherlanguage}

Qul yā ayyuhā al-baʿīd taqūlu va lā tashʿuru va tadhkuru va lā taʿrifu
agar fī al-jomleh shoʿūrī molāḥaẓeh mīshod dhekrī az maʿānī-ye mastūreh
keh dar sorādeq-e ʿerfān-e maʿnavī-ye īn ẓohūr mastūr ast madhkūr mīāmad
tā bar jannat-e lā ẓilla lahā fī al-jomleh moṭṭaleʿ mīshodī va lāken
ḥeyf ast keh laʾālī-ye qolzom-e ʿezz-e rabbānī rā beh dast-e sāreqīn dād
va ḥūr-e maʿānī-ye maqṣūreh rā dar maḥḍar-e ẓālemīn makshūf āvard cheh
keh beh chashm va baṣar va foʾād va qalb khāʾen būdeh va hastand methl-e
ānkeh ʿebārāt-e Ḥaḍrat-e Abhā rūḥ-e māsavāh fedāh rā serqat namūdeh beh
kalemāt-e majʿūleh-ye khod morakkab namūdeh-and va beh esm-e khod
enteshār dādeh-and va ghāfel az īnkeh ṣarrāfān-e aḥadiyyeh dar ṣad hezār
hezār ḥajar yek dāneh loʾloʾ rā beshenāsand va az ū beh maʿdanash pey
barand cheh keh az maʿdan-e āhak tasnīm-e qods-e elāhī jārī nashavad va
az manbaʿ-e qaṭrān salsabīl-e ḥekmat-e rabbānī sārī nagardad va ʿebādī
hastand keh beh rāyeḥeh-ye Reḍvān beh makman-e Raḥmān pey barand cheh
keh rāyeḥeh-ye ṭayyebeh hedāyat namāyandeh ast ṣāḥebān-e shāmmeh-ye aṣfā
rā beh Reḍvān-e Jamāl-e ʿAlī-ye Abhā albatteh dar molk abṣor-e ḥadīdeh
ẓāher shavand va ṣarrāfān-e aḥadiyyeh beh vojūd āyand va javāher-e
kalemātesh rā beh jān kharīdār shavand va az kown va emkān dar sabīlash
bogdharand va jān dar rahash īthār namāyand beh esmash az rasm-e donyā
monqaṭeʿ gardand va beh dhekrash az dhekr-e ʿālamīn bī-niyāz va beh
yādash az yād-e dūnash fāregh.

Say: ``You who are far away! You speak without perceiving and make
mention without understanding.'' Had even a little awareness been
discerned, some of the hidden meanings concealed within the tabernacle
of the spiritual knowledge of this Revelation would have been mentioned,
so that you might have gained some knowledge of the paradise that has no
shade. But it would be a pity to place the pearls of the ocean of divine
glory in the hands of thieves, or to unveil the secluded maidens of
inner meanings in the presence of oppressors, for they have been, and
remain, treacherous in eye and sight, in soul and heart. Thus they have
stolen the words of the Most Glorious---may the souls of all besides Him
be sacrificed for Him---mingled them with their own fabricated words,
and circulated them in their own name, heedless that the assayers of
divine unity recognize a single pearl among a hundred thousand thousand
stones and trace it to its source. For the Tasním of divine holiness
does not flow from a limestone quarry, nor does the Salsabíl of divine
wisdom stream from a spring of pitch. There are servants who find the
hidden sanctuary of the All-Merciful through the fragrance of Riḍván,
for that sweet fragrance guides those with the purest sense of smell to
the Riḍván of the Most Exalted, the All-Glorious Beauty. Assuredly, keen
eyes will appear in the world, and assayers of divine unity will arise.
They will purchase the jewels of His words with their lives, renounce
all existence and the contingent world in His path, and sacrifice their
lives in His way. Through His name they will become detached from the
ways of the world; through His remembrance they will have no need to
mention the worlds; and through remembering Him they will be freed from
remembering anything besides Him.

\subsection{200}\label{section-695}

\begin{otherlanguage}{persian}

بگو اى بی‌‌بصر اگر بمنظر اكبر ناظر شوى و در رضوان بيان سير نمائى جز
اينكلمه نخواهى ديد لم‌يزل بنفس خود دليل بر خود بوده و بظهورش حجّت بر جميع
اهل عالم تمام و كامل و اگر از فضل و مرحمتش اتيان بآيات فرموده اين از فضل
او است مخصوص عباد او ميفرمايد بغير ما يظهر من عنده بر او استدلال مكنيد
چه كه دليل از او ظاهر و حجّت باو ثابت و برهان از او مشهود و غيرش در ساحت
مقدّسش معدوم و مفقود

\end{otherlanguage}

Begū ey bī-baṣar agar beh manẓar-e akbar nāẓer shavī va dar Reḍvān-e
Bayān seyr namāʾī joz īn kalemeh nakhāhī dīd lam-yazal beh nafs-e khod
dalīl bar khod būdeh va beh ẓohūrash ḥojjat bar jamīʿ-e ahl-e ʿālam
tamām va kāmel va agar az faḍl va marḥamatesh etyān beh āyāt farmūdeh īn
az faḍl-e ū ast makhṣūṣ-e ʿebād-e ū mīfarmāyad beh gheyr-e mā yaẓharu
min ʿindihi bar ū estedlāl makonīd cheh keh dalīl az ū ẓāher va ḥojjat
beh ū thābet va borhān az ū mashhūd va gheyrash dar sāḥat-e moqaddasash
maʿdūm va mafqūd.

Say: ``You who have no sight! If you look toward the Most Great Vision
and wander through the Riḍván of the Bayán, you will see nothing but
this statement: He has always been His own proof, and through His
appearance the proof has been made complete and perfect for all the
people of the world. If, out of His grace and mercy, He has brought
forth verses, this is a token of His grace bestowed expressly upon His
servants. He says: `Do not adduce anything as evidence of Him except
what proceeds from Him, for evidence appears from Him, proof is
established through Him, and demonstration is witnessed from Him, while
all besides Him are nonexistent and absent in His holy presence.'\,''

\subsection{201}\label{section-696}

\begin{otherlanguage}{persian}

بگو اى كر حقيقى آخر نشنيده كه فرموده ببيان و ما نزل فى البيان از آن طلعت
رحمن محتجب نمانيد و بحروفات بيان در امرش توقّف مكنيد چه كه جميع حروفات و
مرايا بقول او خلق شده‌اند و جميع بيان از سماء امرش نازل گويا شما از اهل
جابلقاى موهوم اهل فرقانيد و تازه از جنگل وهميّه بيرون آمده‌ايد و به بيشه
ظنّيّه اخرى وارد شده‌ايد حرفى از بيان مطّلع نبوده و نيستيد از سنن عارفين
بی‌خبريد و از مذهب عاشقين بغايت دور

\end{otherlanguage}

Begū ey kar-e ḥaqīqī ākhar nashenīdeh keh farmūdeh beh Bayān va mā
nazala fī al-Bayān az ān Ṭalʿat-e Raḥmān moḥtajeb namānīd va beh
ḥorūfāt-e Bayān dar amrash tavaqqof makonīd cheh keh jamīʿ-e ḥorūfāt va
marāyā beh qowl-e ū khalq shodeh-and va jamīʿ-e Bayān az samā-ye amrash
nāzel gūyā shomā az ahl-e Jābolqā-ye mowhūm-e ahl-e Forqānīd va tāzeh az
jangal-e vahmiyyeh bīrūn āmadeh-īd va beh bīsheh-ye ẓanniyyeh-ye okhrā
vāred shodeh-īd ḥarfī az Bayān moṭṭaleʿ nabūdeh va nīstīd az sonan-e
ʿārefīn bī-khabarīd va az madhhab-e ʿāsheqīn beh ghāyat dūr.

Say: ``You who are truly deaf! Have you not heard that He has said: `Do
not let the Bayán and what was revealed in the Bayán veil you from that
Countenance of the All-Merciful, and do not let the Letters of the Bayán
cause you to hesitate in His Cause'? For all the Letters and Mirrors
have been created by His word, and the whole Bayán has descended from
the heaven of His command. It is as though you belong to the imaginary
Jábulqá of the people of the Qur'an, and have only just emerged from the
forest of delusion to enter another thicket of conjecture. You have
never known, and do not now know, a single letter of the Bayán. You are
unaware of the ways of those who possess true knowledge and are far
indeed from the path of the lovers.''

\subsection{202}\label{section-697}

\begin{otherlanguage}{persian}

بگو اى مغرور دارالغرور با اين ادراك سيف جفا بر تارك لولاك كشيده و بگمان
خود عارف شده و بحقّ رسيده كاش اين راهى كه ميروى انتهايش بتركستان باشد
چنانچه گفته‌اند اين ره كه تو ميروى بتركستان است و لكن اين سبيل منتهى
باسفل نيران بوده و خواهد بود و اين عجب است كه مدّتها است در جوار حقّ آمده
چه كه مدينهء كبيره با اين ارض طيّبه مباركه اينقدر بعيد نبوده كه قاصدين را
زحمت دهد مع‌ذلك خود نيامديد كه سهل است يكنفس مقدّسى هم بساحت اقدس از تو
نيامد و ذكر لطيفى بمقرّ الطف اطهر وارد نشد و بذلك جعلت نفسك محروماً عمّا
خلقت له و كنت من المبعدين

\end{otherlanguage}

Begū ey maghrūr-e dār al-ghorūr bā īn edrāk seyf-e jafā bar tārok-e
Lawlāk keshīdeh va beh gomān-e khod ʿāref shodeh va beh Ḥaqq resīdeh
kāsh īn rāhī keh mīravī entehāyash beh Torkestān bāshad chenāncheh
gofteh-and īn rah keh to mīravī beh Torkestān ast va lāken īn sabīl
montahī beh asfal-e nīrān būdeh va khāhad būd va īn ʿajab ast keh
moddathā ast dar javār-e Ḥaqq āmadeh cheh keh Madīneh-ye Kabīreh bā īn
arḍ-e ṭayyebeh-ye mobārakeh īnqadr baʿīd nabūdeh keh qāṣedīn rā zaḥmat
dehad maʿa dhālek khod nayāmadīd keh sahl ast yek nafs-e moqaddasī ham
beh sāḥat-e aqdas az to nayāmad va dhekr-e laṭīfī beh maqarr-e alṭaf-e
aṭhar vāred nashod va bi-dhālika jaʿalta nafsaka maḥrūman ʿammā khuliqta
lahu va kunta min al-mubʿadīn.

Say: ``You who are deluded by the abode of delusion! With such
understanding, you have drawn the sword of oppression against the head
of Him addressed as `But for You,' imagining that you have attained true
knowledge and reached God. Would that the road you travel ended in
Turkestan, as it has been said: `This road you travel leads to
Turkestan.' But this path has led, and will lead, to the lowest depths
of the fires of hell. It is astonishing that you have long been in the
vicinity of God, for the Great City has not been so far from this good
and blessed land as to make the journey burdensome for those who intend
it. Yet not only did you fail to come yourself, but not a single
sanctified soul came from you to the Most Holy Court, nor did any
gracious message reach the most gracious, most pure abode. Thus you have
deprived yourself of that for which you were created and have been among
those kept far away.''

\subsection{203}\label{section-698}

\begin{otherlanguage}{persian}

در آخر اينمطلب ذكر ميشود بدانكه اين آيات همان آياتى است كه نقطه بيان بآن
مبعوث شده و محمّدا رسول اللّه بآن ظاهر و روح بآن ناطق و كليم بآن ثابت حال
قدرى تفكّر نمائيد كه معرضين باين ظهور چه ارتكاب نموده‌اند عند متغمّسين بحر
معانى و ساكنين فلك حكمت ربّانى نفوس معرضه از اين آيات بديعه الهيّه از كلّ
انبیاء و رسل معرض بوده و خواهند بود و لن يجدن رايحة الجنّة الّتى كانت نفس
الظّهور اولئك من اخسر النّاس و لكن هم لايعلمون نسئل اللّه بأن يعطيك جناح
اليقين فى هذا الأمر المبين لتطير بها الى سمآء العزّ و التّقديس و تقول
الحمد للّه ربّ العالمين

\end{otherlanguage}

Dar ākhar īn maṭlab dhekr mīshavad bedānkeh īn āyāt hamān āyātī ast keh
Noqṭeh-ye Bayān beh ān mabʿūth shodeh va Muḥammadan Rasūl Allāh beh ān
ẓāher va Rūḥ beh ān nāṭeq va Kalīm beh ān thābet ḥāl qadrī tafakkor
namāʾīd keh moʿreḍīn beh īn ẓohūr cheh ertekāb namūdeh-and ʿend-e
motaghammesīn-e baḥr-e maʿānī va sākenīn-e folk-e ḥekmat-e rabbānī
nofūs-e moʿreḍeh az īn āyāt-e badīʿeh-ye elāhiyyeh az koll-e anbiyā va
rosol moʿreḍ būdeh va khāhand būd va lan yajidanna rāʾiḥat al-jannah
allatī kānat nafs al-ẓuhūr ulāʾika min akhsar al-nās va lākin hum lā
yaʿlamūn nasʾalu Allāh bi-an yuʿṭiyaka janāḥ al-yaqīn fī hādhā al-amr
al-mubīn li-taṭīra bihā ilā samāʾ al-ʿizz va al-taqdīs va taqūla
al-ḥamdu li-Allāh rabb al-ʿālamīn.

Finally, this point is mentioned: know that these verses are the very
verses with which the Point of the Bayán was sent forth, through which
Muḥammad, the Messenger of God, appeared, through which the Spirit
spoke, and through which He who conversed with God was established. Now
reflect a little on what those who have turned away from this Revelation
have committed. In the eyes of those immersed in the ocean of inner
meanings and those dwelling in the ark of divine wisdom, souls who turn
away from these wondrous divine verses have turned, and will turn, away
from all the Prophets and Messengers. They will never perceive the
fragrance of the paradise that is the very Person of the Manifestation.
They are among the greatest losers, but they do not know. We ask God to
grant you the wing of certainty in this manifest Cause, so that you may
soar with it to the heaven of glory and holiness and say: ``Praise be to
God, the Lord of the worlds.''

\subsection{204}\label{section-699}

\begin{otherlanguage}{persian}

و اينكه نوشته بودى قطع نظر از نصّ ربّ اعلى بشنو نداى اين عبد را و ابدا از
نصّ نقطه اولى روح ماسواه فداه سخن ذكر منما چه كه و نفسه المنّان و جماله
الرّحمن كه بهيچوجه مطّلع نيستى كه مقصود از نصّ نقطه بيان كه بوده و چيست نه
شما مطّلع و نه دون شما چه كه اين امر از ابصر كلّ مستور و احدى بتفصيل آن
مطّلع نه الّا نفسان واحد منهما استشهد فى سبيل اللّه و نفس ديگر موجود و
مشركين او را كاذب دانسته‌اند كه هذيانات خود را بين ناس انتشار دهند

\end{otherlanguage}

Va īnkeh neveshteh būdī qaṭʿ-e naẓar az naṣṣ-e Rabb-e Aʿlā beshnow
nedā-ye īn ʿabd rā va abadan az naṣṣ-e Noqṭeh-ye Ūlā rūḥ-e māsavāh fedāh
sokhan dhekr manamā cheh keh va nafsihi al-Mannān va jamālihi al-Raḥmān
keh beh hīch vajh moṭṭaleʿ nīstī keh maqṣūd az naṣṣ-e Noqṭeh-ye Bayān
keh būdeh va chīst nah shomā moṭṭaleʿ va nah dūn-e shomā cheh keh īn amr
az abṣor-e koll mastūr va aḥadī beh tafṣīl-e ān moṭṭaleʿ nah illā nafsān
wāḥidun minhumā istushhida fī sabīl Allāh va nafs-e dīgar mowjūd va
moshrekīn ū rā kādheb dānesteh-and keh hadhayānāt-e khod rā beyn-e nās
enteshār dehand.

As for your having written, ``Quite apart from the explicit text of the
Most Exalted Lord,'' listen to the call of this servant, and never speak
of the explicit text of the Primal Point---may the souls of all besides
Him be sacrificed for Him. For, by His bountiful Self and His
all-merciful Beauty, you have no knowledge whatsoever of who was
intended by the explicit text of the Point of the Bayán or what it
means. Neither you nor anyone else knows, for this matter is hidden from
the eyes of all, and no one knows its details except two souls. One of
them was martyred in the path of God, and the other is still living. The
idolaters have deemed him a liar so that they may spread their own
ravings among the people.

\subsection{205}\label{section-700}

\begin{otherlanguage}{persian}

و از اين گذشته اين نصّى كه معرض باللّه ادّعاى آن مينمايد و بشما ذكر
نموده‌اند درباره اسم اللّه الديّان ميرزا اسداللّه اعظم از آن نازل بيانات
نقطه بيان را ملاحظه كن كه جميع اسماء الهى را باو راجع فرموده‌اند و تفصيل
اين در كلماتيكه از مشرق اصبع رحمن اشراق فرموده موجود و هم‌‌چنين عبد حاضر
لدى العرش بعضى الواح در اين امور مبسوط نوشته‌اند ملاحظه كن لعلّ تتذكّر فى
نفسك و تخرق احجاب التّقليد بقدرة ربّك المقتدر العزيز الحميد

\end{otherlanguage}

Va az īn godhashteh īn naṣṣī keh moʿreḍ-e bi-Allāh eddeʿā-ye ān
mīnamāyad va beh shomā dhekr namūdeh-and dar bāreh-ye Esm Allāh
al-Dayyān Mīrzā Asad Allāh aʿẓam az ān nāzel bayānāt-e Noqṭeh-ye Bayān
rā molāḥaẓeh kon keh jamīʿ-e asmā-ye elāhī rā beh ū rājeʿ farmūdeh-and
va tafṣīl-e īn dar kalemātī keh az mashreq-e eṣbaʿ-e Raḥmān eshrāq
farmūdeh mowjūd va hamchenīn ʿabd-e ḥāḍer ladā al-ʿarsh baʿḍī alvāḥ dar
īn omūr mabsūṭ neveshteh-and molāḥaẓeh kon laʿallaka tatadhakkaru fī
nafsika va takhriqu aḥjāb al-taqlīd bi-qudrati rabbika al-Muqtadir
al-ʿAzīz al-Ḥamīd.

Moreover, as for this explicit text that the one who has turned away
from God lays claim to, and which has been mentioned to you, something
greater than it was revealed concerning the Name of God, the Judge,
Mírzá Asadu'lláh. Consider the utterances of the Point of the Bayán, in
which He referred all the divine names to him. The details of this are
found in the words that have shone forth from the dawning-place of the
finger of the All-Merciful. Likewise, the servant present before the
Throne has written certain Tablets at length on these matters. Examine
them, so that you may take heed within yourself and rend the veils of
imitation through the power of your Lord, the All-Powerful, the
Almighty, the All-Praised.

\subsection{206}\label{section-701}

\begin{otherlanguage}{persian}

و همين اسم مذكور را كه نقطه بيان ديّان فرموده كه اعظم اسماء الهى است شما
اليوم دنىّ ميخوانيد و ابوالشّرور ميناميد و مرشد شما ميرزا يحيى تصريحاً من
غير تلويح در كتاب سجّين خود كه بمستيقظ معروف است فتواى قتل آن مظلوم را
داده و هم‌‌چنين نفوس ديگر را برويد و بخوانيد تا مطّلع شويد كه بر مظاهر حقّ
چه وارد شده و هم‌‌چنين نقطه اولى باو مرقوم فرموده انّك انت ان يا حرف الثّالث
المؤمن بمن يظهره اللّه و اين مقامى است كه سبقت داشته بر كلّ مقامات و

\end{otherlanguage}

Va hamīn ism-e madhkūr rā keh Nuqṭeh-ye Bayān Dayyān farmūdeh keh
aʿẓam-e asmā-ye ilāhī ast shomā al-yawm danī mīkhvānīd va Abu'sh-Shurūr
mīnāmīd va murshid-e shomā Mīrzā Yaḥyā taṣrīḥan min ghayri talwīḥ dar
kitāb-e sijjīn-e khvod keh be-Mustayqiẓ maʿrūf ast fatvā-ye qatl-e ān
maẓlūm rā dādeh va ham-chunīn nufūs-e dīgar rā. Beravīd va bekhvānīd tā
muṭṭaliʿ shavīd keh bar maẓāhir-e Ḥaqq cheh vārid shodeh va ham-chunīn
Nuqṭeh-ye Ūlā be-ū marqūm farmūdeh: innaka anta {[}corrupted text{]} yā
ḥarfa al-thālithi al-muʾminu bi-man yuẓhiruhu Allāh. Va īn maqāmī ast
keh sabqat dāshteh bar kull-e maqāmāt va

And this same person, whom the Point of the Bayán called Dayyán,
declaring this to be the greatest of the divine names, you now call base
and name the Father of Evils. Your leader, Mírzá Yaḥyá, explicitly and
without allusion, has issued a sentence of death against that wronged
one, as well as other souls, in his own book of perdition, known as the
Mustayqiẓ. Go and read it, so that you may learn what has befallen the
Manifestations of the Truth. Likewise, the Primal Point wrote to him:
``You are indeed {[}corrupted original text{]}, O third Letter, the
believer in Him whom God will make manifest.'' This is a station that
has taken precedence over all stations, and

\subsection{207}\label{section-702}

\begin{otherlanguage}{persian}

بعد بيانى در اينمقام فرموده‌اند كه از تلويح كلمات الهى چنان مستفاد ميشود
كه بلقاء اللّه در ظهور بعد فايز ميگردد و آن هيكل قدم او را باو ميشناساند
چنانچه بين يدى حاضر شد و لسان اللّه باو متكلّم مجملى از تفصيل او اينكه
هنگامى كه ميرزا يحيى فتواى قتل او را نوشت حضرت ابهى از عراق هجرت اختيار
فرموده بودند و قريب دو سنه بود كه ابدا احدى از ايشان اطّلاع نداشت و چون
رقم قتل از نزد مرشدت ظاهر شد ميرزا محمّد نامى مخصوص رفته بآذربايجان كه آن
مظلوم را شهيد نمايد ابی‌ اللّه عمّا اراد ليتمّ القول فيما ذكر فى اللّوح و اسم
مذكور از آذربايجان آمده و جميع اصحاب بر قتلش ايستادند از عرب و عجم و اين
در وقتى بود كه حضرت ابهى مراجعت فرموده بودند و در عراق تشريف داشتند
بقسمى امر شديد شد كه يومى از ایّام جمال ابهى از اوّل طلوع فجر تا قريب بعصر
يك يك از اصحاب را كه در صدد اذيّت بيچاره بودند طلبيده و نهى بليغ فرمودند

\end{otherlanguage}

Baʿd bayānī dar īn maqām farmūdeh-and keh az talvīḥ-e kalamāt-e ilāhī
chunān mustafād mīshavad keh be-liqāʾ Allāh dar ẓuhūr-e baʿd fāʾiz
mīgardad va ān Haykal-e Qidam ū rā be-ū mīshināsānad, chunāncheh bayna
yaday ḥāḍir shod va Lisān Allāh be-ū mutakallim. Mujmalī az tafṣīl-e ū
īnkeh hangāmī keh Mīrzā Yaḥyā fatvā-ye qatl-e ū rā nivisht Ḥaḍrat-e Abhā
az ʿIrāq hijrat ikhtiyār farmūdeh būdand va qarīb-e do saneh būd keh
abadan aḥadī az īshān iṭṭilāʿ nadāsht va chūn raqam-e qatl az nazd-e
murshidat ẓāhir shod Mīrzā Muḥammad-nāmī makhṣūṣ rafteh be-Ādharbāyjān
keh ān maẓlūm rā shahīd namāyad. Abā Allāhu ʿammā arāda li-yatimma
al-qawlu fīmā dhukira fī al-lawḥ. Va ism-e madhkūr az Ādharbāyjān āmadeh
va jamīʿ-e aṣḥāb bar qatl-ash īstādand az ʿArab va ʿAjam va īn dar vaqtī
būd keh Ḥaḍrat-e Abhā murājaʿat farmūdeh būdand va dar ʿIrāq tashrīf
dāshtand. Be-qismī amr shadīd shod keh yawmī az ayyām Jamāl-e Abhā az
avval-e ṭulūʿ-e fajr tā qarīb be-ʿaṣr yek yek az aṣḥāb rā keh dar
ṣadad-e adhiyyat-e bīchāreh būdand ṭalabīdeh va nahy-e balīgh farmūdand.

He then made a statement concerning this station, from the allusions of
whose divine words it is understood that he would attain the presence of
God in the next Revelation, and that the Temple of Eternity would make
Him known to him. Indeed, he came into His presence, and the Tongue of
God spoke to him. A brief account of his story is this: when Mírzá Yaḥyá
wrote the sentence of death against him, His Holiness the All-Glorious
had departed from Iraq, and for nearly two years no one had any news of
Him whatsoever. When the decree of death issued from your leader, a
certain Mírzá Muḥammad went expressly to Azerbaijan to martyr that
wronged one. God thwarted what he intended, so that the words recorded
in the Tablet might be fulfilled. The aforementioned person came from
Azerbaijan, and all the companions, Arab and Persian alike, rose up to
kill him. This was at a time when His Holiness the All-Glorious had
returned and was residing in Iraq. Matters became so grave that, one
day, from the first light of dawn until nearly the afternoon, the
All-Glorious Beauty summoned, one by one, those companions who were
intent on harming that helpless man and emphatically forbade them.

\subsection{208}\label{section-703}

\begin{otherlanguage}{persian}

و بعد از دو يوم اسم مذكور بين يدى حاضر و آنچه از مفتريات كه نسبت باو
داده بودند از خود سلب نموده و قسم ياد نمود كه آنچه نسبت باين عبد داده‌اند
كذب صرف است و جمال ابهى كمال ملاطفت باو فرمودند و بآنچه در لوح ثبت بود
فايز شد و بمنزل خود راجع گشت بعد از چند يوم نفس مذكور او را فريب داده و
از كاظمين آورده قرب خانه مرشدت آن بيچاره را شهيد نمود فواللّه در شهادت او
غبارى در عراق مرتفع گشت و بقسمى ظلمت فرو گرفت كه روز بشب تبديل شد و در
اسواق سراج برافروختند و جميع اهل عراق خائفاً مضطرباً متزلزلاً صيحه يا اللّه
برآوردند و ندا مينمودند چه از اناث و چه از ذكور كه از سفك دم اين مظلوم
است كه اين بليّه بر عراق نازل شده و مرشد مستشعر نشد و از نوم غفلت بيدار
نگشت و اللّه بظلمى او را شهيد نمودند كه آسمانها بفزع و نوحه و ندبه
درآمدند مع‌ذلك متنبّه نشده نفسى را كه حضرت اعلى روح ماسواه فداه او را فرد
الأحد ناميده و هم‌‌چنين مظهر الأحديّه و اسم اللّه الديّان خطاب فرموده در
الواح ناريّه خود بحقارت او را مينامند و تا حال احدى از مشركين او را
نشناختند چه كه ستر رحمن حايل بود قل يا قوم اتّقوا اللّه و لاتسفكوا الدّمآء
توبوا اليه يا معشر الغافلين لعلّ يكفّر عنكم جريراتكم و يغفر سيّئاتكم الّتى
ناحت بها سكّان الملك و الملكوت ثمّ اهل ملأ العالّين

\end{otherlanguage}

Va baʿd az do yawm ism-e madhkūr bayna yaday ḥāḍir va āncheh az
muftariyāt keh nisbat be-ū dādeh būdand az khvod salb namūdeh va qasam
yād namūd keh āncheh nisbat be-īn ʿabd dādeh-and kidhb-e ṣirf ast va
Jamāl-e Abhā kamāl-e mulāṭafat be-ū farmūdand va be-āncheh dar lawḥ
thabt būd fāʾiz shod va be-manzil-e khvod rājiʿ gasht. Baʿd az chand
yawm nafs-e madhkūr ū rā farīb dādeh va az Kāẓimayn āvordeh qurb-e
khāneh-ye murshidat ān bīchāreh rā shahīd namūd. Fa-wa-Allāh dar
shahādat-e ū ghubārī dar ʿIrāq murtafiʿ gasht va be-qismī ẓulmat furū
girift keh rūz be-shab tabdīl shod va dar asvāq sirāj bar-afrūkhtand va
jamīʿ-e ahl-e ʿIrāq khāʾifan muḍṭariban mutazalzilan ṣayḥeh-ye yā Allāh
bar-āvordand va nidā mīnamūdand cheh az ināth va cheh az dhukūr keh az
safk-e dam-e īn maẓlūm ast keh īn baliyyeh bar ʿIrāq nāzil shodeh va
murshid mustashʿir nashod va az nawm-e ghaflat bīdār nagasht. Va-Allāh
be-ẓulmī ū rā shahīd namūdand keh āsmānhā be-fazaʿ va nawḥeh va nudbeh
dar-āmadand. Maʿa dhālik mutanabbih nashodeh nafsī rā keh Ḥaḍrat-e Aʿlā
rūḥ-e mā sivāh fidāh ū rā Fard al-Aḥad nāmīdeh va ham-chunīn Maẓhar
al-Aḥadiyyah va Ism Allāh al-Dayyān khiṭāb farmūdeh dar alvāḥ-e
nāriyyeh-ye khvod be-ḥaqārat ū rā mīnāmand va tā ḥāl aḥadī az mushrikīn
ū rā nashinākhtand cheh keh sitr-e Raḥmān ḥāʾil būd. Qul yā qawmi ittaqū
Allāha wa-lā tasfikū al-dimāʾ. Tūbū ilayhi yā maʿshara al-ghāfilīn
laʿalla yukaffiru ʿankum jarīrātikum wa-yaghfiru sayyiʾātikumu allatī
nāḥat bihā sukkānu al-mulki wa-l-malakūt thumma ahlu malaʾi al-ʿāllīn.

Two days later, the aforementioned person came into His presence,
repudiated the fabrications attributed to him, and swore, ``Whatever
they have attributed to this servant is sheer falsehood.'' The
All-Glorious Beauty treated him with the utmost kindness. He attained
what had been recorded in the Tablet and returned to his lodging. A few
days later, the person already mentioned deceived him, brought him from
Káẓimayn, and martyred that helpless man near your leader's house. By
God, at his martyrdom a cloud of dust rose over Iraq, and darkness
descended so thickly that day turned to night and lamps were lit in the
markets. All the people of Iraq, fearful, distraught, and shaken, cried
out, ``O God!'' Women and men alike called out, ``It is because the
blood of this wronged one has been shed that this calamity has descended
upon Iraq.'' Yet the leader did not perceive it and did not awaken from
the sleep of heedlessness. By God, they martyred him with such cruelty
that the heavens were filled with terror, lamentation, and mourning.
Even so, they have not taken heed, and in their fiery tablets they speak
contemptuously of the soul whom His Holiness the Most Exalted---may the
soul of all besides Him be sacrificed for Him---named the Peerless One
of the One, and likewise addressed as the Manifestation of Oneness and
the Name of God, the Judge. To this day none of the idolaters has
recognized him, for the veil of the All-Merciful intervened. Say: O
people, fear God and do not shed blood. Turn to Him in repentance, O
company of the heedless, that He may expiate your offenses and forgive
your misdeeds, over which the inhabitants of the earthly world and the
heavenly kingdom, and then the company of the exalted ones, have
lamented.

\subsection{209}\label{section-704}

\begin{otherlanguage}{persian}

و مرشد شما در مقام استحقار آن شهيد بی‌معين ذكر مينمود كه نقطه اولى باو
مرقوم فرموده‌اند و سفارش نموده كه كاغذ و قلم خوب از براى من تحصيل نمايد
كه من بر او بنويسم بگو سلّمنا كه اين قول مقبول از همين معلوم است كه تو در
ظلّ او بوده چه كه اينگونه تلطّف از آباء ظاهر ميشود نسبت باطفال و كاش همين
قدر مستشعر ميشد لاواللّه لن يستشعر فى نفسه و لو يأتى اللّه بكلّ آية او
بقبيل من اهل السّموات او بجنود الرّوح او بصحف الأوّلين و الآخرین فواللّه از
ظلم آن انفس مشركه هياكل احديّه در نوحه و ندبه و نقطه اولى بحنين و ناله
مشغول فيا ليت يكون من ذى سمعٍ و ذى بصر و ذى قلب ليشهد و يسمع و يفقه غشت
قلوبهم و افئدتهم شواظ من النّار و مارج من النّيران و منعتهم عن التّوجّه الى
افعالهم و اعمالهم و ما ظهر من انفسهم و ما يخرج من افواههم

\end{otherlanguage}

Va murshid-e shomā dar maqām-e istiḥqār-e ān shahīd-e bī-muʿīn dhikr
mīnamūd keh Nuqṭeh-ye Ūlā be-ū marqūm farmūdeh-and va sifārish namūdeh
keh kāghadh va qalam-e khūb az barā-ye man taḥṣīl namāyad keh man bar ū
binivīsam. Begū sallamnā keh īn qawl maqbūl, az hamīn maʿlūm ast keh to
dar ẓill-e ū būdeh cheh keh īngūneh talaṭṭuf az ābāʾ ẓāhir mīshavad
nisbat be-aṭfāl va kāsh hamīn qadr mustashʿir mīshod. Lā wa-Allāh lan
yastashʿira fī nafsihi wa-law yaʾtī Allāhu bi-kulli āyatin aw bi-qabīlin
min ahli al-samāwāt aw bi-junūdi al-rūḥ aw bi-ṣuḥufi al-awwalīn
wa-l-ākhirīn. Fa-wa-Allāh az ẓulm-e ān anfus-e mushrikeh hayākil-e
aḥadiyyeh dar nawḥeh va nudbeh va Nuqṭeh-ye Ūlā be-ḥanīn va nāleh
mashghūl. Fa-yā layta yakūnu min dhī samʿin wa-dhī baṣarin wa-dhī qalbin
li-yashhada wa-yasmaʿa wa-yafqaha. Ghashat qulūbahum wa-afʾidatahum
shuwāẓun min al-nār wa-mārijun min al-nīrān wa-manaʿathum ʿan
al-tawajjuhi ilā afʿālihim wa-aʿmālihim wa-mā ẓahara min anfusihim wa-mā
yakhruju min afwāhihim.

Your leader, seeking to belittle that martyr who had no helper, would
mention that the Primal Point had written to him and requested that he
obtain good paper and pens for Him to write with. Say: Granted that this
statement is accepted, this very thing makes it clear that you were
under his shadow, for such kindness is shown by fathers toward children.
If only he had perceived even this much! No, by God, he will never
perceive it within himself, even should God come with every sign, or
with a company of the inhabitants of the heavens, or with the hosts of
the Spirit, or with the scriptures of the former and latter generations.
By God, because of the injustice of those idolatrous souls, the Temples
of Oneness are lamenting and mourning, and the Primal Point is sighing
and weeping. If only there were someone with ears, eyes, and a heart,
that he might witness, hear, and understand! A flame of fire and a blaze
of burning heat have enveloped their hearts and their inmost selves,
preventing them from attending to their deeds and actions, to what has
proceeded from their own selves, and to what comes from their mouths.

\subsection{210}\label{section-705}

\begin{otherlanguage}{persian}

بگو اى غافل جاهل سلّمنا كه مجعولات تو قبول شد نصّ نقطه بيان است در كلّ بيان
كه در حين ظهور ببيان و ما خلق فيه از نفس ظهور محتجب نشويد حمد خدا را كه
آن مشرك و امثال او از تصريح كلمات محكمه متقنه الهيّه محروم مانده تا چه
رسد بتلويح اشارات مالك اسماء و صفات بگو از خدا شرم كنيد و ذيل دين قويم
الهى را بطين ظنون و اوهام ميالائيد و شمس مشرقه سماء احديّه را بحجبات
نفسيّه مپوشانيد قل تاللّه انّ هذه لشمس ما اخذها الكسوف و قمر لايخسف ابداً
اذا موتوا بغيظكم ثمّ بنار حسدكم يا ملأ الأشرار هنوز آن قدر ادراك
ننموده‌ايد كه در حين اشراق شمس مدلول دليل لا يسمن و لا يغنى بوده چه كه
مقصود از دليل عرفان نفس مدلول بوده اذا فافتح بصراك لتراها مشرقة عن افق
العزّ و المجد و الجلال

\end{otherlanguage}

Begū ay ghāfil-e jāhil sallamnā keh majʿūlāt-e to qabūl shod, naṣṣ-e
Nuqṭeh-ye Bayān ast dar kull-e Bayān keh dar ḥīn-e ẓuhūr be-Bayān va mā
khuliqa fīh az nafs-e ẓuhūr muḥtajib nashavīd. Ḥamd-e Khudā rā keh ān
mushrik va amthāl-e ū az taṣrīḥ-e kalamāt-e muḥkameh-ye mutqaneh-ye
ilāhiyyeh maḥrūm māndeh tā cheh rasad be-talvīḥ-e ishārāt-e Mālik-e
Asmāʾ va Ṣifāt. Begū az Khudā sharm kunīd va dhayl-e dīn-e qavīm-e ilāhī
rā be-ṭīn-e ẓunūn va awhām mayālāyīd va shams-e mushriqeh-ye samā-ye
aḥadiyyeh rā be-ḥujubāt-e nafsiyyeh mapūshānīd. Qul ta-Allāhi inna
hādhihi la-shamsun mā akhadhahā al-kusūfu wa-qamarun lā yukhsafu abadan.
Idhan mūtū bi-ghayẓikum thumma bi-nāri ḥasadikum yā malaʾa al-ashrār.
Hanūz ān qadr idrāk nanamūdeh-īd keh dar ḥīn-e ishrāq-e shams-e madlūl
dalīl lā yusminu wa-lā yughnī būdeh cheh keh maqṣūd az dalīl ʿirfān-e
nafs-e madlūl būdeh. Idhan fa-ftaḥ baṣrāka li-tarāhā mushriqatan ʿan
ufuqi al-ʿizzi wa-l-majdi wa-l-jalāl.

Say: O heedless and ignorant one, granted that your fabrications are
accepted, the explicit statement of the Point of the Bayán throughout
the Bayán is that, at the time of the Revelation, you must not allow the
Bayán and all that was created within it to veil you from the Revelation
itself. Praise be to God! That idolater and those like him have remained
deprived of the explicit meaning of God's firm and incontrovertible
words, let alone the allusions contained in the intimations of the Lord
of Names and Attributes. Say: Be ashamed before God. Do not soil the hem
of God's steadfast religion with the mire of suppositions and
imaginings, and do not conceal the Sun shining from the heaven of
Oneness behind the veils of self. Say: By God, this is indeed a Sun that
no eclipse has overtaken, and a Moon that will never be eclipsed. Die,
then, in your rage and in the fire of your envy, O company of the
wicked! You have not yet understood even this much: when the Sun of the
One to whom the proof points shines forth, the proof neither nourishes
nor avails, for the purpose of the proof is to recognize the very One to
whom it points. Open your eyes, then, that you may see it shining from
the horizon of might, glory, and majesty.

\subsection{211}\label{section-706}

\begin{otherlanguage}{persian}

قوله كبر غفلته پس اگر ديگرى برآيد مدّعى مقامى بعد از نقطه بيان شود خالى
از اين نيست يا مدّعى آن بايد باشد كه در ظلّ نقطه بيان است و دوره نقطه بيان
هنوز منقضى نگشته و تكميل اين دوره نشده بايد در اينصورت كمال عبوديّت نسبت
بربّ اعلى نمايد و غايت اطاعت بفرمايش او نمايد و ابدا بقدر جوى از فرمايش
ربّ اعلى ايد تخلّف ننمايد چونكه دوره دوره اوست و خود مدّعى در ظلّ افتاده
انتهى

\end{otherlanguage}

Qawluhu, kabura ghaflatuhu: pas agar dīgarī bar-āyad muddaʿī-ye maqāmī
baʿd az Nuqṭeh-ye Bayān shavad khālī az īn nīst, yā muddaʿī ān bāyad
bāshad keh dar ẓill-e Nuqṭeh-ye Bayān ast va dawreh-ye Nuqṭeh-ye Bayān
hanūz munqaḍī nagashteh va takmīl-e īn dawreh nashodeh bāyad dar īn
ṣūrat kamāl-e ʿubūdiyyat nisbat be-Rabb-e Aʿlā namāyad va ghāyat-e
iṭāʿat be-farmāyesh-e ū namāyad va abadan be-qadr-e javī az farmāyesh-e
Rabb-e Aʿlā {[}corrupted text{]} takhalluf nanamāyad chūnkeh dawreh
dawreh-ye ūst va khvod-e muddaʿī dar ẓill uftādeh. Intahā.

His words---great is his heedlessness: ``Then, if another should arise
and claim a station after the Point of the Bayán, it cannot be
otherwise: either he must claim that he is under the shadow of the Point
of the Bayán, while the dispensation of the Point of the Bayán has not
yet ended and this dispensation has not been completed. In that case he
must show perfect servitude toward the Most Exalted Lord and the utmost
obedience to His command, and never depart, even by a barleycorn, from
the command of the Most Exalted Lord {[}corrupted original text{]}, for
the dispensation is His dispensation, and the claimant himself stands
under the shadow.'' End of quotation.

\subsection{212}\label{section-707}

\begin{otherlanguage}{persian}

حال قدرى انصاف ده و چشم بگشا كه در چه مقام واقعى و در چه رتبه قائم چه
فرق مابين خود و ملل قبل مشاهده ميكنى چه كه آن قوم كه از زلال سلسال معارف
عزّ صمدانى و كوثر افضال سلطان لايزالى محروم مانده‌اند سبب اين شده كه
بدلايل مثبته نزد خود مستدلّ و معتكف و بان دلايل بر سلطان قبايل اعتراض
نموده و از شريعه قدس احديّه محروم گشته‌اند

\end{otherlanguage}

Ḥāl qadrī inṣāf deh va chashm bigushā keh dar cheh maqām vāqiʿī va dar
cheh rutbeh qāʾim. Cheh farq mābayn-e khvod va milal-e qabl mushāhedeh
mīkunī? Cheh keh ān qawm keh az zulāl-e salsāl-e maʿārif-e ʿizz-e
ṣamadānī va kawthar-e afḍāl-e Sulṭān-e lā-yazālī maḥrūm māndeh-and sabab
īn shodeh keh be-dalāʾil-e muthbiteh nazd-e khvod mustadill va muʿtakif
va be-ān dalāʾil bar Sulṭān-e qabāʾil iʿtirāḍ namūdeh va az sharīʿeh-ye
quds-e aḥadiyyeh maḥrūm gashteh-and.

Now be a little fair, and open your eyes to see the position you occupy
and the rank in which you stand. What difference do you perceive between
yourself and the religious communities of the past? The reason those
people have remained deprived of the clear, flowing water of the
knowledge of divine glory and the heavenly stream of the favors of the
everlasting Sovereign is that they have relied upon and clung to proofs
they themselves regard as established. With those proofs they have
raised objections against the Sovereign of all peoples and have been
deprived of the holy waters of Oneness.

\subsection{213}\label{section-708}

\begin{otherlanguage}{persian}

مثلاً اهل تورئة بآن استدلال نموده كه بعد از موسى نفسى صاحب شريعت نخواهد
آمد و هر نفسى هم بيايد بايد شريعت او را ترويج نمايد و هم‌‌چنين علاماتى
معيّن نموده‌اند كه نفس ظهور بايد بآن علامات ظاهر شود و لذا ممنوع شدند از
عرفان مظاهر ظهور چه كه نفس ظهور يظهر كيف يظهره اللّه و اراد لا بما عند
النّاس و هم‌‌چنين اهل انجيل برآنند كه خود عيسى از آسمان نازل ميشود بعد از
ظلمت شمس و تاريكى قمر و سقوط كواكب و هم‌چنين ملّت فرقان كه ابدا قائل بر
اينكه صاحب شريعتى ظاهر شود نبوده و نيستند و ميگويند نفسى هزار سنه قبل از
نرجس خاتون متولّد شده و در جابلقا و جابلسا الآن موجود است و منتظر اذن
الهى است كه ظاهر شود و او ميايد و شريعت رسول اللّه را يعنى همين شريعت كه
در دست هست ترويج ميدهد و علاماتى هم در اينمقام ذكر نموده‌اند كه اگر تفصيل
ذكر شود كتابی‌ بايد تاليف نمود

\end{otherlanguage}

Mathalan ahl-e Tawrāt be-ān istidlāl namūdeh keh baʿd az Mūsā nafsī
ṣāḥib-e sharīʿat nakhvāhad āmad va har nafsī ham biyāyad bāyad
sharīʿat-e ū rā tarvīj namāyad va ham-chunīn ʿalāmātī muʿayyan
namūdeh-and keh nafs-e ẓuhūr bāyad be-ān ʿalāmāt ẓāhir shavad va lidhā
mamnūʿ shodand az ʿirfān-e maẓāhir-e ẓuhūr cheh keh nafs-e ẓuhūr yaẓharu
kayfa yuẓhiruhu Allāhu wa-arāda lā bi-mā ʿinda al-nās. Va ham-chunīn
ahl-e Injīl bar-ānand keh khvod-e ʿĪsā az āsmān nāzil mīshavad baʿd az
ẓulmat-e shams va tārīkī-ye qamar va suqūṭ-e kavākib va ham-chunīn
millat-e Furqān keh abadan qāʾil bar īnkeh ṣāḥib-e sharīʿatī ẓāhir
shavad nabūdeh va nīstand va mīgūyand nafsī hizār saneh qabl az Narjis
Khātūn mutavallid shodeh va dar Jābulqā va Jābulsā al-ān mawjūd ast va
muntaẓir-e idhn-e ilāhī ast keh ẓāhir shavad va ū mīāyad va sharīʿat-e
Rasūl Allāh rā yaʿnī hamīn sharīʿat keh dar dast hast tarvīj mīdahad va
ʿalāmātī ham dar īn maqām dhikr namūdeh-and keh agar tafṣīl dhikr shavad
kitābī bāyad taʾlīf namūd.

For example, the people of the Torah have reasoned that no one bringing
a religious law will come after Moses, and that whoever does come must
promote His law. They have likewise specified signs with which the
Manifestation Himself must appear. Thus they were prevented from
recognizing the Manifestations of Revelation, for the Manifestation
Himself appears as God makes Him appear and as He has willed, not
according to what people possess. Likewise, the people of the Gospel
hold that Jesus Himself will descend from heaven after the sun has been
darkened, the moon has grown dim, and the stars have fallen. Likewise,
the people of the Qur'an have never accepted, and do not now accept,
that a bearer of a religious law will appear. They say that a person was
born a thousand years ago to Narjis Khátún, that he now exists in
Jábulqá and Jábulsá, and that he awaits God's permission to appear. He
will come and promote the law of the Messenger of God---that is, this
very law now in their possession. They have also mentioned signs in this
connection, which would require a book to set out in detail.

\subsection{214}\label{section-709}

\begin{otherlanguage}{persian}

جميع اينمطالب را ببراهين متقنه و نصوص محكمه و دلائل محقّقه نزد خود چنان
ثابت نموده‌اند كه منكرين آن را كافر و باطل و مشرك دانسته و ميدانند و در
حين اشراق شموس عزّ تقديس چون بر خلاف آراء اين نفوس بوده لذا كلّ محجوب و
محروم مانده‌اند مع‌ذلك تازه استدلال مينمائيد كه اين نفسيكه ظاهر است يا
بايد فوق آن باشد يا در ظلّ او اگر در ظلّ او است بايد باعمال او عامل شود و
اگر فوق او رتبه را ادّعا نمايد محال بوده چه كه هنوز دوره دوره او است و
تكميل نشده

\end{otherlanguage}

Jamīʿ-e īn maṭālib rā be-barāhīn-e mutqaneh va nuṣūṣ-e muḥkameh va
dalāʾil-e muḥaqqaqeh nazd-e khvod chunān thābit namūdeh-and keh
munkirīn-e ān rā kāfir va bāṭil va mushrik dānisteh va mīdānand va dar
ḥīn-e ishrāq-e shumūs-e ʿizz-e taqdīs chūn bar khilāf-e ārā-ye īn nufūs
būdeh lidhā kull muḥtajib va maḥrūm māndeh-and. Maʿa dhālik tāzeh
istidlāl mīnamāyīd keh īn nafsī keh ẓāhir ast yā bāyad fawq-e ān bāshad
yā dar ẓill-e ū. Agar dar ẓill-e ū ast bāyad be-aʿmāl-e ū ʿāmil shavad
va agar fawq-e ū rutbeh rā iddiʿā namāyad muḥāl būdeh cheh keh hanūz
dawreh dawreh-ye ū ast va takmīl nashodeh.

They have established all these matters to their own satisfaction with
such firm arguments, explicit texts, and verified proofs that they have
regarded, and still regard, those who deny them as unbelievers,
followers of falsehood, and idolaters. When the Suns of the glory of
holiness shone forth, their appearance ran contrary to the opinions of
these souls, and so all remained veiled and deprived. Nevertheless, you
now begin to argue that this Person who has appeared must either be
above Him or under His shadow. If He is under His shadow, He must act
according to His practices; and if He claims a rank above Him, that is
impossible, for the dispensation is still His dispensation and has not
been completed.

\subsection{215}\label{section-710}

\begin{otherlanguage}{persian}

حال انصاف ده فرق شما با اين فئه چه چيز است بلى فرق لايحصى چه كه در كتب
آنقوم باشارات معضله و آيات مقنّعه مستوره تلويحاً من غير تصريح ذكر ظهور بعد
شده و لذا محتجب مانده‌اند چه كه ادراك ننموده‌اند معانى كلمات الهيّه را و
لكن شما و امثال شما مع آنكه نقطه بيان روح ما سواه فداه تصريحاً من غير
تلويح ميفرمايند بآنچه مابين سموات و ارض خلق شده چه از كلمات و چه از
اشارات از آن سلطان وجود و مليك شهود خود را منع ننمائيد دلايل ذكر ميكنيد
كه حقّ جايز نيست ظاهر شود

\end{otherlanguage}

Ḥāl inṣāf deh farq-e shomā bā īn fiʾeh cheh chīz ast? Balī farq lā
yuḥṣā, cheh keh dar kutub-e ān qawm be-ishārāt-e muʿḍileh va āyāt-e
muqannaʿeh-ye mastūreh talvīḥan min ghayri taṣrīḥ dhikr-e ẓuhūr-e baʿd
shodeh va lidhā muḥtajib māndeh-and cheh keh idrāk nanamūdeh-and
maʿānī-ye kalamāt-e ilāhiyyeh rā. Va lākin shomā va amthāl-e shomā maʿa
ānkeh Nuqṭeh-ye Bayān rūḥ-e mā sivāh fidāh taṣrīḥan min ghayri talvīḥ
mīfarmāyand be-āncheh mābayn-e samāvāt va arḍ khalq shodeh cheh az
kalamāt va cheh az ishārāt az ān Sulṭān-e vujūd va Malīk-e shuhūd khvod
rā manʿ nanamāyīd dalāʾil dhikr mīkunīd keh Ḥaqq jāʾiz nīst ẓāhir
shavad.

Now be fair: what difference is there between you and this group?
Indeed, the difference is beyond measure, for in those people's books
the next Revelation was mentioned through abstruse allusions and veiled,
hidden verses, by implication rather than explicit statement. Thus they
remained veiled, for they did not comprehend the meanings of the divine
words. But you and those like you, even though the Point of the
Bayán---may the soul of all besides Him be sacrificed for
Him---explicitly and without allusion says that you must not allow
anything created between heaven and earth, whether words or allusions,
to keep you from that Sovereign of existence and King of the visible
world, adduce proofs that God may not manifest Himself.

\subsection{216}\label{section-711}

\begin{otherlanguage}{persian}

صدهزار آفرين بر فهم و كمال شما و مرشدين شما از اين كلمات كبريّه غروريه كه
بشما القاء نموده‌اند معلوم شد كه خود را مختار ميدانند نه حقّ را بگو اى
متوهّم هرگز حقّ در هيچ ظهورى بخيال ناس ظاهر نشده و بر خلاف آنچه در دست
جميع ناس بوده ظاهر شده چه كه مقدّس بوده ذات امنع اقدسش از جميع ظنون و
اوهام و اشارات و دلالات و احدى بر نفس ظهور و كيفيّت او مطّلع نه الّا نفسه
العليم الخبير آنقدر کرده‌اند كه روى يهود و جميع محتجبين ارض را سفيد
کرده‌اند يهود يداللّه را مغلول دانسته ديگر دستور العمل از براى حقّ ذكر
نكرده كه اگر همچه ادّعا نمايد جايز نيست و اگر فلان ادّعا نمايد بايد بشئون
محدوده معيّنه ظاهر شود

\end{otherlanguage}

Ṣad-hizār āfarīn bar fahm va kamāl-e shomā va morshedīn-e shomā! Az īn
kalemāt-e kebrīyyeh-ye ghorūrīyyeh keh be-shomā elqā namūdeh-and maʿlūm
shod keh khod rā mokhtār mīdānand nah Ḥaqq rā. Begū ey motavahhem,
hargez Ḥaqq dar hīch ẓohūrī be-khīyāl-e nās ẓāher nashodeh va bar
khelāf-e āncheh dar dast-e jamīʿ-e nās būdeh ẓāher shodeh, cheh keh
moqaddas būdeh dhāt-e amnaʿ-e aqdasash az jamīʿ-e ẓonūn va avhām va
eshārāt va dalālāt va aḥadī bar nafs-e ẓohūr va keyfīyyat-e ū moṭṭaleʿ
nah, illā nafsuhu al-ʿAlīm al-Khabīr. Ānqadr kardeh-and keh rūy-e Yahūd
va jamīʿ-e moḥtajebīn-e arḍ rā sefīd kardeh-and. Yahūd yadullāh rā
maghlūl dānesteh, dīgar dastūr al-ʿamal az barā-ye Ḥaqq dhekr nakardeh
keh agar hamcheh eddeʿā namāyad jāyez nīst va agar folān eddeʿā namāyad
bāyad be-shoʾūn-e maḥdūdeh-ye moʿayyanه ẓāher shavad.

A hundred thousand bravos for your understanding and accomplishments,
and those of your guides! These grandiose, arrogant words they have
instilled in you show that they consider themselves, not God, free to
choose. Say: O deluded one, never in any Revelation has God appeared
according to people's imaginings. Rather, He has appeared contrary to
what all people possessed, for His most inaccessible, most holy Essence
has been sanctified above all conjectures, fancies, allusions, and
indications. No one knows the Revelation itself or its manner except He
Himself, the All-Knowing, the All-Informed. They have gone so far as to
make the Jews and all the veiled peoples of the earth seem blameless by
comparison. The Jews considered the hand of God fettered, but they did
not go on to prescribe instructions for God, saying that if He makes
such a claim it is inadmissible, and if He makes some other claim He
must appear in specified, limited ways.

\subsection{217}\label{section-712}

\begin{otherlanguage}{persian}

قد كبر قولك يا غافل آنقدر ادراك ننموده كه تكميل هر ظهور قبلى بظهور بعد
شده و خواهد شد چنانچه در اين ظهور بدع قدس رحمانى در سنه تسع در سرّ سرّ
نفوس مقدّسه مطهّره زكيّه در همان حين تكميل شدند و لكن انت لا تشعر و لا تفقه
و لا تعرف چه كه شكر مصر عرفان الهى قسمت طوطى معانى بوده نه قسمت زاغ
ظلمانى و معرض باللّه چون توقّف نمود و از صراط لغزيد در همان حين هيكلش از
قميص انسانى خارج و بجلود بهائم ظاهر و مشهود گشت فسبحان من يبدّل الوجود
كيف یشآء و يقدّر الأمور كيف يريد لا رادّ لأمره و لا مردّ لمشيّته كلّ شیء
منقاد لسلطنته و مشفق من خشيته و مرتعش من سطوته و لذلك الرّبّ يليق تهليل
اهل ملأ الأعلی ثمّ تسبيح اهل مداين البقآء و لكن تو از آن نفوسى بنظر ميائى
كه اراده ادراك امرى و يا تفقّه در شيئ نداشته و ندارى مقصود اثبات رياستى
است كه از براى مرشد خود محقّق نمايند و خود هم در ظلّش معروف و موسوم باشند
تاللّه انّ ظلّه برهوت و انّ مائه ملح اجاج لو انتم تفقهون

\end{otherlanguage}

Qad kabura qawluka yā ghāfil. Ānqadr edrāk nanamūdeh keh takmīl-e har
ẓohūr-e qablī be-ẓohūr-e baʿd shodeh va khāhad shod, chenāncheh dar īn
ẓohūr-e badʿ-e qods-e raḥmānī dar saneh-ye tesʿ dar serr-e serr nofūs-e
moqaddaseh-ye moṭahhareh-ye zakīyyeh dar hamān ḥīn takmīl shodand, va
lākin anta lā tashʿuru wa-lā tafqahu wa-lā taʿrifu. Cheh keh shekar-e
Meṣr-e ʿerfān-e elāhī qesmat-e ṭūṭī-ye maʿānī būdeh nah qesmat-e zāgh-e
ẓolmānī, va moʿreḍ-e bi-Allāh chūn tavaqqof namūd va az ṣerāṭ laghzīd
dar hamān ḥīn heykalash az qamīṣ-e ensānī khārej va be-jolūd-e bahāʾem
ẓāher va mashhūd gasht. Fa-subḥāna man yubaddilu al-wujūda kayfa yashāʾu
wa-yuqaddiru al-umūra kayfa yurīdu, lā rādda li-amrihi wa-lā maradda
li-mashīʾatihi, kullu shayʾin munqādun li-salṭanatihi wa-mushfiqun min
khashyatihi wa-murtaʿishun min saṭwatihi, wa-li-dhālika al-Rabbi yalīqu
tahlīlu ahli malaʾi al-aʿlā thumma tasbīḥu ahli madāʾini al-baqāʾ. Va
lāken to az ān nofūsī be-naẓar mīāyī keh erādeh-ye edrāk-e amrī va yā
tafaqqoh dar shayʾī nadāshteh va nadārī. Maqṣūd ethbāt-e rīyāsatī ast
keh az barā-ye morshed-e khod moḥaqqaq namāyand va khod ham dar ẓellash
maʿrūf va mawsūm bāshand. Ta-Allāhi inna ẓillahu barahūtun wa-inna
māʾahu milḥun ujājun law antum tafqahūn.

Grievous indeed are your words, O heedless one! You have not understood
even this much: every former Revelation has been, and will be, brought
to completion through the subsequent Revelation. Thus, in this wondrous,
holy Revelation of the All-Merciful, in the year nine, holy, purified,
stainless souls were brought to completion at that very moment, in the
innermost secret. But you neither perceive, nor comprehend, nor
recognize. For the sugar of the Egypt of divine knowledge has been the
portion of the parrot of inner meanings, not of the raven of darkness.
When the one who turns away from God halted and slipped from the path,
at that very moment his form emerged from its human garment and
appeared, plainly visible, in the skins of beasts. Glory, then, to Him
who transforms existence as He pleases and ordains affairs as He wills.
No one can repel His command or reverse His will. All things submit to
His sovereignty, stand in awe of Him, and tremble before His might. That
Lord deserves the acclamation of His unity by the people of the company
on high, and the glorification offered by the inhabitants of the cities
of eternity. But you appear to be one of those souls who have never
intended, and do not intend, to understand any matter or comprehend
anything. Their purpose is to establish a position of leadership for
their guide, so that they themselves may become known and recognized
under his shadow. By God, his shadow is a barren wilderness, and his
water bitter salt, if only you understood.

\subsection{218}\label{section-713}

\begin{otherlanguage}{persian}

بگو اى اصمّ باطنى احدى بر علم حقّ احاطه ننموده و نخواهد نمود علمه مكنون
مخزون مصون من اطّلاعكم اگر اراده فرمايد در يك حين خلق اوّلين و آخرين را
تكميل ميفرمايد و لكن تو هنوز تكميل را نفهميده كاش استعداد ادراك هم در تو
مشاهده ميشد و اين عبد كه يكى از خدّام درگاه است تو را مستعدّ مييافت و بر
تو كلماتى القاء مينمود كه از آن ادراك نمائى بعضى امور مستوره را اقلّا
آنقدر ميشد كه اگر بر تكميل ظهور اطّلاع نمييافتى ببغض خود مطّلع ميشدى
فاستحى عن اللّه و لا تقنّن له قوانين نفسك و هوئك انّه يظهر كيف یشآء لا بما
انت تشآء

\end{otherlanguage}

Begū ey aṣamm-e bāṭenī, aḥadī bar ʿelm-e Ḥaqq eḥāṭeh nanamūdeh va
nakhāhad namūd. ʿIlmuhu maknūnun makhzūnun maṣūnun min iṭṭilāʿikum. Agar
erādeh farmāyad dar yek ḥīn khalq-e avvalīn va ākherīn rā takmīl
mīfarmāyad, va lāken to hanūz takmīl rā nafahmīdeh. Kāsh esteʿdād-e
edrāk ham dar to moshāhedeh mīshod va īn ʿabd keh yekī az khoddām-e
dargāh ast to rā mostaʿedd mīyāft va bar to kalemātī elqā mīnamūd keh az
ān edrāk namāʾī baʿḍī omūr-e mastūreh rā. Aqallan ānqadr mīshod keh agar
bar takmīl-e ẓohūr eṭṭelāʿ namīyāftī be-boghḍ-e khod moṭṭaleʿ mīshodī.
Fa-staḥi ʿani Allāhi wa-lā tuqannin lahu qawānīna nafsika wa-hawāka,
innahu yaẓharu kayfa yashāʾu lā bi-mā anta tashāʾ.

Say: O inwardly deaf one, no one has encompassed God's knowledge, nor
will anyone ever encompass it. His knowledge is hidden, treasured, and
protected from your scrutiny. If He so wills, He brings the creation of
the first and the last to completion in a single instant. Yet you have
still not understood completion. If only some capacity for understanding
could be seen in you, and this servant, who is one of the attendants at
the threshold, could find you receptive and impart words to you through
which you might understand certain hidden matters. At least, even if you
did not come to understand the completion of the Revelation, you would
become aware of your own hatred. Be ashamed before God, and do not
prescribe for Him the laws of your own self and desires. He appears as
He wills, not as you will.

\subsection{219}\label{section-714}

\begin{otherlanguage}{persian}

چه مقدار جسور شده‌اند كه از براى حقّ جلّ ذكره قواعد ذكر مينمايند و شرايط
مذكور ميدارند حال شما انصاف ده آن حقّيكه بدستور العمل مشركين ظاهر شود مثل
خود آن انفس مشركه خواهد بود كلّما عرفتم و استدركتم هو مخلوق مثلكم و مردود
اليكم و اللّه پست‌ترين خلق او و عباد او ننگ ميدارند كه بآدابيكه آن نفوس
مشركه بشما القا نموده‌اند حركت نمايند لانّهم يعبدون العجل و لا يشعرون و
يتّبعون الشّيطان و لا يعرفون و يعقّبون الجبت و لا يفقهون و يشربون فضلة ما
شربه الطّاغوت و لا يعلمون

\end{otherlanguage}

Cheh meqdār jasūr shodeh-and keh az barā-ye Ḥaqq, jalla dhikruhu,
qavāʿed dhekr mīnamāyand va sharāyeṭ madhkūr mīdārand. Ḥāl shomā enṣāf
deh, ān Ḥaqqī-keh be-dastūr al-ʿamal-e moshrekīn ẓāher shavad methl-e
khod-e ān anfos-e moshrekeh khāhad būd. Kullamā ʿaraftum wa-stadraktum
huwa makhlūqun mithlukum wa-mardūdun ilaykum. Va-Allāh past-tarīn-e
khalq-e ū va ʿebād-e ū nang mīdārand keh be-ādābī-keh ān nofūs-e
moshrekeh be-shomā elqā namūdeh-and ḥarakat namāyand, li-annahum
yaʿbudūna al-ʿijla wa-lā yashʿurūna wa-yattabiʿūna al-shayṭāna wa-lā
yaʿrifūna wa-yuʿaqqibūna al-jibta wa-lā yafqahūna wa-yashrabūna faḍlata
mā sharibahu al-ṭāghūtu wa-lā yaʿlamūn.

How audacious they have become, prescribing rules and laying down
conditions for God, exalted be His remembrance! Now judge fairly: a God
who appeared according to the instructions of those who associate
partners with Him would be like those idolatrous souls themselves.
Whatever you have known and apprehended is a creature like yourselves
and is returned to you. By God, the lowliest of His creatures and
servants would be ashamed to act according to the practices those
idolatrous souls have instilled in you. For they worship the calf
without perceiving it, follow Satan without recognizing it, pursue the
idol without understanding, and drink the dregs of what the false god
has drunk, without knowing.

\subsection{220}\label{section-715}

\begin{otherlanguage}{persian}

و آنقدر ادراك ننموده‌اند كه در هر ظهور اگر حقّ جلّ كبريائه بآداب و شرايط و
اعماليكه مابين ناس بر حسب ظاهر بوده ظاهر ميشد هرگز احدى اعراض و اعتراض
نمينمود و جميع بكوثر معانى حضرت لايزالى فايز ميشدند و بلقاء اللّه مرزوق
ميگشتند آنانكه در هواء روح طايرند و ببصر حديد الهى در اشیآء ناظر صدهزار
امثال آن نفوس را در غرقاب فنا مشاهده مينمايند كه هر يك بحشيش منبته متشبّث
و لكن از براى آن حشيش قرار و استقرارى نبوده و نخواهد بود و بالأخره‌غرق
خواهد شد اگر چه حال هم نزد جالسين فلك ابهى لاشیء محض و غريق و فانى و
معدوم مشاهده ميشوند

\end{otherlanguage}

Va ānqadr edrāk nanamūdeh-and keh dar har ẓohūr agar Ḥaqq, jalla
kibriyāʾuhu, be-ādāb va sharāyeṭ va aʿmālī-keh mābeyn-e nās bar ḥasab-e
ẓāher būdeh ẓāher mīshod hargez aḥadī eʿrāḍ va eʿterāḍ namīnamūd va
jamīʿ be-kawthar-e maʿānī-ye Ḥaḍrat-e Lāyazālī fāʾez mīshodand va
be-leqāʾ Allāh marzūq mīgashtand. Ānāṅ-keh dar havā-ye rūḥ ṭāyerand va
be-baṣar-e ḥadīd-e elāhī dar ashyāʾ nāẓer ṣad-hizār amthāl-e ān nofūs rā
dar gharqāb-e fanā moshāhedeh mīnamāyand keh har yek be-ḥashīsh-e
monbeteh motashabbeth va lāken az barā-ye ān ḥashīsh qarār va esteqrārī
nabūdeh va nakhāhad būd va be-l-ākhereh gharq khāhad shod, agarcheh ḥāl
ham nazd-e jālesīn-e folk-e Abhā lā-shayʾ-e maḥḍ va gharīq va fānī va
maʿdūm moshāhedeh mīshavand.

They have not understood even this much: if, in every Revelation, God,
exalted be His majesty, had appeared in accordance with the outward
practices, conditions, and deeds prevailing among people, no one would
ever have turned away or objected. All would have attained the Kawthar
of the meanings of the Eternal One and been granted the encounter with
God. Those who soar in the atmosphere of the spirit and regard all
things with the keen sight of God see a hundred thousand souls like
these in the depths of destruction, each clinging to a growing blade of
grass. Yet that grass has never had, and will never have, any firmness
or stability, and in the end each will drown, although even now, in the
eyes of those seated in the all-glorious Ark, they are seen as sheer
nothingness, drowned, perished, and nonexistent.

\subsection{221}\label{section-716}

\begin{otherlanguage}{persian}

سمع را از قطن تقليد پاك كن بشنو كلمه حقّ را از پست‌ترين خلق او و بدانكه حقّ
جلّ شأنه بهيچ حجابی‌ محجوب نه و بهيچ حدودى محدود نخواهد بود ظاهر ميشود بهر
قسم كه اراده فرمايد چه كه لم‌يزل مقتدر بوده و خواهد بود و هرگز قدرت
كلّيّه‌اش سلب نشده و نخواهد شد حال بوهم شما اگر حقّ بدستور العمل معرضين ظاهر
شود قدرتش را چه تاويل مينمائى چه كه از بيان شما چنان مستفاد ميشود كه
لابدّ بايد باين قسم ظاهر شود و قادر نيست بغير آنچه شما نوشته‌ايد مظهرى از
مظاهر امر خود را در ارض بفرستد فبئس ما ظننتم و تظنّون و قد حبطت اعمالكم
بما خرجت من افواهكم و لكن انتم لا تستشعرون فى انفسكم و تكونن من الغافلين

\end{otherlanguage}

Samʿ rā az qoṭn-e taqlīd pāk kon, beshnow kalemeh-ye Ḥaqq rā az
past-tarīn-e khalq-e ū va bedān-keh Ḥaqq, jalla shaʾnuhu, be-hīch ḥejābī
maḥjūb nah va be-hīch ḥodūdī maḥdūd nakhāhad būd. Ẓāher mīshavad be-har
qesm keh erādeh farmāyad, cheh keh lam-yazal moqtader būdeh va khāhad
būd va hargez qodrat-e kollīyyeh-ash salb nashodeh va nakhāhad shod. Ḥāl
be-vahm-e shomā agar Ḥaqq be-dastūr al-ʿamal-e moʿreḍīn ẓāher shavad
qodratash rā cheh taʾvīl mīnamāʾī? Cheh keh az bayān-e shomā chenān
mostafād mīshavad keh lābodd bāyad be-īn qesm ẓāher shavad va qāder nīst
be-gheyr-e āncheh shomā neveshteh-īd maẓharī az maẓāher-e amr-e khod rā
dar arḍ beferestad. Fa-biʾsa mā ẓanantum wa-taẓunnūna wa-qad ḥabiṭat
aʿmālukum bi-mā kharajat min afwāhikum wa-lākin antum lā tastashʿirūna
fī anfusikum wa-takūnanna mina al-ghāfilīn.

Clear your hearing of the cotton of imitation. Hear the word of God from
the lowliest of His creatures, and know that God, exalted be His glory,
is concealed by no veil and will be confined by no limits. He appears in
whatever manner He wills, for He has always been, and will always be,
all-powerful. His absolute power has never been, and will never be,
taken from Him. Now, if, as you imagine, God were to appear according to
the instructions of those who turn away, how would you interpret His
power? For your words imply that He must necessarily appear in this
manner and is unable to send to earth a Manifestation of His Cause in
any way other than what you have written. Evil indeed is what you have
imagined and continue to imagine! Your deeds have come to nothing
because of what has issued from your mouths, yet you do not perceive
this within yourselves and are surely among the heedless.

\subsection{222}\label{section-717}

\begin{otherlanguage}{persian}

فواللّه روايح منتنه كثيفه از اشارات كلمات آن انفس مشركه بر عالم و اهل آن
مرور نموده از خدا ميطلبيم كه صرصر عنايتى بفرستد و اين ارياح را معدوم
نمايد خود آن مشركين خمسين الف سنه يوم قيامت را مستدلّند كه بيك ساعت منقضى
شد بگو اى بی‌‌بصران همان معنى در اينجا جارى پنجاه هزار سنه در ساعتى منقضى
شود حرفى نداريد و لكن اگر دو هزار سال بوهم شما در سنين معدوده منقضى شود
اعتراض مينمائيد فواللّه يكذّبكم كلّ الأشیآء و يضحك لقلّة عقولكم و عرفانكم كلّ
ما خلق بين السّموات و الأرضين

\end{otherlanguage}

Fa-wa-Allāh ravāyeḥ-e monteneh-ye kathīfeh az eshārāt-e kalemāt-e ān
anfos-e moshrekeh bar ʿālam va ahl-e ān morūr namūdeh. Az Khodā
mīṭalabīm keh ṣarṣar-e ʿenāyatī beferestad va īn aryāḥ rā maʿdūm
namāyad. Khod-e ān moshrekīn khamsīn alf saneh-ye yawm-e qīyāmat rā
mostadelland keh be-yek sāʿat monqaḍī shod. Begū ey bī-baṣarān, hamān
maʿnī dar īnjā jārī. Panjāh hizār saneh dar sāʿatī monqaḍī shavad ḥarfī
nadārīd va lāken agar do hizār sāl be-vahm-e shomā dar senīn-e maʿdūdeh
monqaḍī shavad eʿterāḍ mīnamāʾīd. Fa-wa-Allāhi yukadhdhibukum kullu
al-ashyāʾi wa-yaḍḥaku li-qillati ʿuqūlikum wa-ʿirfānikum kullu mā
khuliqa bayna al-samāwāti wa-l-arḍīn.

By God, foul, noxious odors from the implications of those idolatrous
souls' words have swept over the world and its inhabitants. We ask God
to send a mighty wind of grace and dispel these winds. Those idolaters
themselves argue that the fifty thousand years of the Day of
Resurrection elapsed in a single hour. Say: O sightless ones, the same
meaning applies here. You have nothing to say if fifty thousand years
elapse in an hour, but if two thousand years, as you imagine them,
elapse in a few years, you object. By God, all things give you the lie,
and everything created between the heavens and the earths laughs at how
little understanding and knowledge you possess.

\subsection{223}\label{section-718}

\begin{otherlanguage}{persian}

بشنو از اين عبد و سبحات اوهام را خرق كن و بحقّ بنفسه ناظر باش چه كه دون
او مخلوق اويند و بكلمه من عند او ظاهر و ابدا بدون خود معروف نگشته بلكه
ما سوايش باو معروف بوده و خواهد بود چشم بنفس ظهور و بما يظهر من عنده
داشته باش و شكر كن پروردگار را كه ظاهر فرمود مظهر نفس خود را بغتة چنانچه
نقطه بيان ميفرمايد و لقد قرب الزّوال و انتم راقدون شمس در قطب زوال مشرق و
مضئ و اين مقام شكر است نه شكايت رحمترا نقمت مدان و نعمت را غضب مشمر
اگرچه لم‌يزل ظهور مظاهر احديّه نعمت ابرار و نقمت فجّار بوده و خواهد بود
فجّار و اشرار لازال نعمت را نقمت انگاشته‌اند و عنايت صرفه را غضب دانسته‌اند

\end{otherlanguage}

Beshnow az īn ʿabd va soboḥāt-e avhām rā kharq kon va be-Ḥaqq be-nafseh
nāẓer bāsh, cheh keh dūn-e ū makhlūq-e ūyand va be-kalemeh-ye min ʿend-e
ū ẓāher va abadan be-dūn-e khod maʿrūf nagashteh balkeh mā-sevāyash be-ū
maʿrūf būdeh va khāhad būd. Chashm be-nafs-e ẓohūr va bi-mā yaẓharu min
ʿindihi dāshteh bāsh va shokr kon Parvardegār rā keh ẓāher farmūd
maẓhar-e nafs-e khod rā baghtatan, chenāncheh Noqṭeh-ye Bayān
mīfarmāyad: wa-laqad qaruba al-zawālu wa-antum rāqidūn. Shams dar qoṭb-e
zavāl moshreq va moḍīʾ va īn maqām-e shokr ast nah shekāyat. Raḥmat rā
neqmat madān va neʿmat rā ghaḍab mashmor, agarcheh lam-yazal ẓohūr-e
maẓāher-e aḥadīyyeh neʿmat-e abrār va neqmat-e fojjār būdeh va khāhad
būd. Fojjār va ashrār lāzāl neʿmat rā neqmat angāshteh-and va ʿenāyat-e
ṣerfeh rā ghaḍab dānesteh-and.

Listen to this servant, tear through the veils of vain imaginings, and
look to God Himself. For all besides Him are His creatures, made
manifest by a word from Him. He has never been known through anyone but
Himself; rather, all besides Him have been, and will be, known through
Him. Keep your eyes on the Revelation itself and on whatever appears
from Him, and thank the Lord for suddenly revealing the Manifestation of
His own Self, as the Point of the Bayán says: ``The noon has indeed
drawn near, while you lie asleep.'' The Sun shines and sheds its light
at the zenith of noon. This is an occasion for gratitude, not complaint.
Do not regard mercy as affliction or count blessing as wrath, although
the appearance of the Manifestations of oneness has always been, and
will always be, a blessing for the righteous and an affliction for the
wicked. The wicked and evil have always taken blessing for affliction
and regarded pure grace as wrath.

\subsection{224}\label{section-719}

\begin{otherlanguage}{persian}

در تكميل دوره ذكر نموده بودى بشنو بيان منزل بيان روح من فى حقايق الأكوان
فداه را كه بنصّ صريح ميفرمايند قوله جلّ كبريائه من اوّل ذلك الأمر الى قبل
ان يكمل تسعة كينونات الخلق لم تظهر و انّ كلّما قد رأيت من النّطفة الى ما
كسوناه لحماً ثمّ اصبر حتّى تشهد خلق الآخر اذا قل فتبارك اللّه احسن الخالقين
و اشهد انّ فرق القائم و القيّوم عدد تسعة ذلك ما تكملن الكينونيات فى
مقاعدهنّ ذلك فرق بين كلّ اعظم و عظيم و در مقام ديگر و لتراقبنّ فرق القائم و
القيّوم ثمّ فى سنة التّسع كلّ خير تدركون و در اينمقام نازل هذا ما وعدناك قبل
حين الّذى اجبناك اصبر حتّى يقضى عن البيان تسعة اذا قل فتبارك اللّه احسن
المبدعين

\end{otherlanguage}

Dar takmīl-e dowreh dhekr namūdeh būdī. Beshnow bayān-e Monzel-e Bayān,
rūḥu man fī ḥaqāʾiqi al-akwāni fidāhu, rā keh be-naṣṣ-e ṣarīḥ
mīfarmāyand, qawluhu jalla kibriyāʾuhu: min awwali dhālika al-amri ilā
qabli an yakmula tisʿatun kaynūnātu al-khalqi lam taẓhar wa-inna kullamā
qad raʾayta mina al-nuṭfati ilā mā kasawnāhu laḥman thumma iṣbir ḥattā
tashhada khalqa al-ākhar idhan qul fa-tabāraka Allāhu aḥsanu
al-khāliqīna wa-shhad anna farqa al-qāʾimi wa-l-qayyūmi ʿadadu tisʿatin
dhālika mā takmulanna al-kaynūniyyātu fī maqāʿidihinna dhālika farqun
bayna kulli aʿẓamin wa-ʿaẓīm. Va dar maqām-e dīgar: wa-la-turāqibanna
farqa al-qāʾimi wa-l-qayyūmi thumma fī sanati al-tisʿi kulla khayrin
tudrikūn. Va dar īn maqām nāzel: hādhā mā waʿadnāka qablu ḥīna alladhī
ajabnāka iṣbir ḥattā yuqḍā ʿani al-bayāni tisʿatun idhan qul fa-tabāraka
Allāhu aḥsanu al-mubdiʿīn.

You had spoken of the completion of the dispensation. Listen to the
words of the Revealer of the Bayán---may the soul of whoever dwells in
the realities of existence be sacrificed for Him---who states in an
explicit text, exalted be His majesty: ``From the beginning of that
Cause until before nine is completed, the essential realities of
creation have not appeared. Indeed, all that you have seen is from the
seed until We clothed it in flesh. Then wait until you witness the other
creation. Then say: `Blessed be God, the best of creators!' And bear
witness that the difference between the Qá'im and the Qayyúm is the
number nine. That is when the essential realities will be completed in
their stations. That is the difference between every `most great' and
`great.'\,'' And in another passage: ``You shall surely watch for the
difference between the Qá'im and the Qayyúm. Then, in the year nine, you
will attain all good.'' And in this connection it was revealed: ``This
is what We promised you before, at the time when We answered you: `Wait
until nine has elapsed from the Bayán. Then say: Blessed be God, the
best of originators!'\,''

\subsection{225}\label{section-720}

\begin{otherlanguage}{persian}

اى صاحب بصر و انصاف دار على زعم خود چشم بگشا و از ظهور تسع تا حين خلق
آخر مشاهده كن و بتبارك اللّه احسن الخالقين و تبارك اللّه احسن المبدعين
ناطق شو ملاحظه كن كه در تكميل خلق و اعظميّت ظهور بعد در فرق قائم و قيّوم و
هم‌‌چنين اعظم و عظيم تصريحاً من غير تلويح ميفرمايند و چون مدّت منقضى شد و
ميقات اللّه بآخر رسيد ظهر جمال التّسع بسلطان مبين اذا فزع كلّ من فى السّموات
و الأرض الّا من اتى اللّه بقلب سليم اتّقوا اللّه يا قوم و لا تتّخذوا القيّوم
مهجوراً بينكم خافوا عن اللّه و كونوا من المتّقين ایّاکم ان تحرّفوا كلمات اللّه
عن مواضعها لانّ ذلك خطأ عظيم فى الواح عزّ حفيظ

\end{otherlanguage}

Ey ṣāḥeb-e baṣar va enṣāf-dār ʿalā zaʿm-e khod, chashm bogshā va az
ẓohūr-e tesʿ tā ḥīn khalq-e ākhar moshāhedeh kon va be-tabāraka Allāhu
aḥsanu al-khāliqīn va tabāraka Allāhu aḥsanu al-mubdiʿīn nāṭeq show.
Molāḥeẓeh kon keh dar takmīl-e khalq va aʿẓamīyyat-e ẓohūr-e baʿd dar
farq-e Qāʾem va Qayyūm va ham-chenīn aʿẓam va ʿaẓīm taṣrīḥan min ghayri
talwīḥ mīfarmāyand. Va chūn moddat monqaḍī shod va mīqāt Allāh be-ākhar
rasīd, ẓahara Jamālu al-Tisʿi bi-sulṭānin mubīnin idhan faziʿa kullu man
fī al-samāwāti wa-l-arḍi illā man atā Allāha bi-qalbin salīm. Ittaqū
Allāha yā qawmu wa-lā tattakhidhū al-Qayyūma mahjūran baynakum. Khāfū
ʿani Allāhi wa-kūnū mina al-muttaqīn. Iyyākum an tuḥarrifū kalimāti
Allāhi ʿan mawāḍiʿihā li-anna dhālika khaṭaʾun ʿaẓīmun fī alwāḥi ʿizzin
ḥafīẓ.

O you who, in your own estimation, possess sight and fairness! Open your
eyes and behold the other creation from the Revelation of the year nine
until now, and proclaim: ``Blessed be God, the best of creators!'' and
``Blessed be God, the best of originators!'' Observe how explicitly,
without allusion, He speaks of the completion of creation and the
greater greatness of the subsequent Revelation in the difference between
Qá'im and Qayyúm, and likewise between ``most great'' and ``great.''
When the period elapsed and God's appointed time reached its end, the
Beauty of the year Nine appeared with manifest authority. Then all who
are in the heavens and on earth were seized with terror, except those
who came to God with a sound heart. Fear God, O people, and do not leave
the Qayyúm forsaken among you. Stand in awe of God and be among the
God-fearing. Beware of distorting the words of God from their proper
places, for that is a grievous error in the preserved Tablets of glory.

\subsection{226}\label{section-721}

\begin{otherlanguage}{arabic}

و لكن نطفه مقام بلوغ را ادراك ننمايد چنانچه اليوم ادراك بيانات الهى را
احدى ننموده الّا من شآء ربّك اگرچه اين كلمات را بظنونات باطله خود تأويل
نموده تحريف خواهند نمود چنانچه نقطه اولى جلّت عظمته خبر داده

\end{otherlanguage}

Va lāken noṭfeh maqām-e bolūgh rā edrāk nanamāyad, chenāncheh al-yawm
edrāk-e bayānāt-e elāhī rā aḥadī nanamūdeh illā man shāʾa rabbuka,
agarcheh īn kalemāt rā be-ẓonūnāt-e bāṭeleh-ye khod taʾvīl namūdeh
taḥrīf khāhand namūd, chenāncheh Noqṭeh-ye Ūlā, jallat ʿaẓamatuhu,
khabar dādeh.

But the seed cannot comprehend the station of maturity, just as today no
one has understood the divine utterances except those whom your Lord has
willed. Yet they will interpret these words according to their own vain
conjectures and distort them, as the Primal Point, exalted be His
grandeur, has foretold.

\subsection{227}\label{section-722}

\begin{otherlanguage}{persian}

حال از اينها گذشته يك ذكر مينمايم لعلّ بآن مستشعر شويد و بشعور آئيد اگر
نفسى مذكور دارد كه شجره از حين الى انقضاء دوهزار سنه باوراد لطيفه حمريّه
معطّره ظاهر خواهد شد با كمال طراوت و نضارت و عطر و روايح محبوبه و حصر
نمايد كه شما ناظر بخود آن شجره و ظهور آن باشيد نه دون او و بعد صبح از
خواب برخيزيد و بغتةً ملاحظه نمائيد كه آن شجره طيّبه مباركه باوراد حمريّه
لطيفه منيعه ظاهر شده حال چه ميكنى انكار خواهى نمود آنچه ببصر مشاهده
ميكنى و بشامّه مييابی‌ يا بآنچه ديده موقن خواهى شد لا واللّه شما و امثال
شما از آن نفوسى مشاهده ميشويد كه آن اوراد مشهوده ظاهره را باحجار ظنون و
هوى در هم ريزند كه اين زود بثمر آمده بايد دوهزار سنه بعد از اين اوراد
ظاهر شود و از اين شجره مشهود آيد افّ لكم يا ملأ المتوهّمين

\end{otherlanguage}

Ḥāl az īnhā goẕashte yek ẕekr mīnamāyam laʿalla be-ān mostashʿer shavīd
va be-shoʿūr āyīd. Agar nafsī maẕkūr dārad ke shajare az ḥīn ilā
inqiḍāʾ-e do-hezār sane be-ovrād-e laṭīfe-ye ḥamriyye-ye moʿaṭṭare ẓāher
khvāhad shod bā kamāl-e ṭarāvat va naḍārat va ʿeṭr va ravāyeḥ-e maḥbūbe
va ḥaṣr namāyad ke shomā nāẓer be-khod-e ān shajare va ẓohūr-e ān bāshīd
na dūn-e ū va baʿd ṣobḥ az khvāb barkhīzīd va baghtatan molāḥeẓe namāyīd
ke ān shajare-ye ṭayyebe-ye mobārake be-ovrād-e ḥamriyye-ye laṭīfe-ye
manīʿe ẓāher shode, ḥāl che mīkonī? Enkār khvāhī namūd ānche be-baṣar
moshāhede mīkonī va be-shāmme mīyābī yā be-ānche dīde mūqen khvāhī shod?
Lā wa-Allāh! Shomā va amthāl-e shomā az ān nofūsī moshāhede mīshavīd ke
ān ovrād-e mashhūde-ye ẓāhere rā be-aḥjār-e ẓonūn va havā dar ham rīzand
ke īn zūd be-thamar āmade, bāyad do-hezār sane baʿd az īn ovrād ẓāher
shavad va az īn shajare mashhūd āyad. Uffin lakum yā malaʾ
al-mutawahhimīn!

Now, setting these matters aside, I will mention something, so that
perhaps you may become aware of it and come to your senses. Suppose
someone states that a tree will appear, sometime between now and the end
of two thousand years, adorned with delicate, fragrant crimson blossoms,
in the fullness of freshness and verdure, with perfume and pleasing
fragrances, and expressly directs you to look to that tree itself and
its appearance, and to nothing else. Then, one morning, you rise from
sleep and suddenly see that this good and blessed tree has appeared with
delicate, sublime crimson blossoms. What will you do now? Will you deny
what you see with your eyes and perceive by smell, or will you become
certain of what you have seen? No, by God! You and those like you are
seen to be among those who would batter those visible, manifest blossoms
with the stones of conjecture and desire, saying: ``This has borne fruit
too soon. These blossoms should appear two thousand years from now and
be seen upon this tree.'' Shame on you, O company of the deluded!

\subsection{228}\label{section-723}

\begin{otherlanguage}{persian}

مع آنكه بنصّ صريح منزل بيان ابدا ظهور بعد محدود بوقتى نشده چنانچه منصوصاً
ميفرمايد اگر اين حين بخواهد ظاهر شود مختار است چه كه عارف بحين ظهور احدى
الّا اللّه نبوده و نيست و در جميع بيان ميفرمايند منوط باراده و مشيّت خود آن
سلطان احديّه است يظهر نفسه كيف یشآء و يريد انّما الدّليل آياته و الوجود على
نفسه نفسه اذ الغير يعرف به و هو لا يعرف بدونه سبحان اللّه عمّا تصفون و
هم‌‌چنين ميفرمايند چه كسى عالم بظهور نيست غير اللّه هر وقت شود بايد كلّ
تصديق بنقطه حقيقت نمايند و شكر الهى بجا آورند انتهى

\end{otherlanguage}

Maʿ-e ānke be-naṣṣ-e ṣarīḥ-e Monzel-e Bayān abadan ẓohūr-e baʿd maḥdūd
be-vaqtī nashode, chenānche manṣūṣan mīfarmāyad: agar īn ḥīn bekhvāhad
ẓāher shavad mokhtār ast, che ke ʿāref be-ḥīn-e ẓohūr aḥadī illā Allāh
nabūde va nīst. Va dar jamīʿ-e Bayān mīfarmāyand manūṭ be-erāde va
mashiyyat-e khod-e ān Solṭān-e Aḥadiyye ast. Yuẓhiru nafsahu kayfa
yashāʾu wa-yurīd. Innamā al-dalīlu āyātuhu wa-al-wujūdu ʿalā nafsihi
nafsuhu, idh al-ghayru yuʿrafu bihi wa-huwa lā yuʿrafu bi-dūnihi.
Subḥāna Allāhi ʿammā taṣifūn. Va hamchenīn mīfarmāyand: che kasī ʿālem
be-ẓohūr nīst ghayr Allāh, har vaqt shavad bāyad koll taṣdīq be-Noqṭe-ye
Ḥaqīqat namāyand va shokr-e elāhī be-jā āvarand. Intahā.

Yet, according to the explicit text of the Revealer of the Bayán, the
subsequent Revelation has never been confined to any particular time, as
He expressly states: ``If He wishes to appear at this very moment, He is
free to do so, for no one but God has ever known or now knows the time
of His appearance.'' And throughout the Bayán He states that it depends
upon the will and wish of that Sovereign of Oneness Himself: ``He
manifests Himself however He wishes and desires. His verses alone are
the proof, and His existence rests upon His own Self, since all others
are known through Him, while He is not known through anything besides
Him. Exalted is God above what you describe.'' And He likewise states:
``Since no one but God knows of the Revelation, whenever it occurs, all
must acknowledge the Point of Truth and offer thanks to God.'' End of
quotation.

\subsection{229}\label{section-724}

\begin{otherlanguage}{persian}

حال ملاحظه نمائيد كه بچه صراحت ميفرمايند كه هر وقت ظاهر شود بايد تصديق
نمايند مع‌ذلك بجاى آنكه فرموده شكر الهى بجا آورند عوض شكر شكايت آغاز
نموده و بجاى تصديق سهام كين و سيف بغضا بر آن مطاف ارواح نبيّين و مرسلين
هجوم نموده‌اند بگو اى محروم ملاحظه نما كه سلطان معلوم آن ظهور محتوم مكنون
را بهيچ شرطى مشروط فرموده فأت بها لا فو‌الّذى خلق الممكنات بكلمة من عنده
ابدا اين ظهور را معلّق و مشروط بامرى و شرطى نفرموده‌اند چه كه شرط در آن
ساحت مفقود صرف است و جز نفس او بر او احدى محيط و مطّلع نه

\end{otherlanguage}

Ḥāl molāḥeẓe namāyīd ke be-che ṣarāḥat mīfarmāyand ke har vaqt ẓāher
shavad bāyad taṣdīq namāyand. Maʿa-dhālik be-jā-ye ānke farmūde shokr-e
elāhī be-jā āvarand, ʿevaḍ-e shokr shekāyat āghāz namūde va be-jā-ye
taṣdīq sehām-e kīn va seyf-e baghḍā bar ān maṭāf-e arvāḥ-e nabiyyīn va
morsalīn hojūm namūde-and. Begū ey maḥrūm, molāḥeẓe namā ke Solṭān-e
Maʿlūm ān Ẓohūr-e Maḥtūm-e Maknūn rā be-hīch sharṭī mashrūṭ farmūde?
Faʾti bihā! Lā fa-wa-alladhī khalaqa al-mumkināti bi-kalimatin min
ʿindihi! Abadan īn ẓohūr rā moʿallaq va mashrūṭ be-amrī va sharṭī
nafarmūde-and, che ke sharṭ dar ān sāḥat mafqūd-e ṣerf ast va joz nafs-e
ū bar ū aḥadī moḥīṭ va moṭṭaleʿ na.

Now consider how explicitly He states that whenever He appears, they
must acknowledge Him. Nevertheless, instead of offering thanks to God as
He instructed, they have begun to complain rather than give thanks, and
instead of acknowledging Him, they have assailed with arrows of malice
and the sword of hatred the One around whom circle the spirits of the
Prophets and Messengers. Say: O deprived one! Consider whether the known
Sovereign has made that hidden, inevitable Revelation subject to any
condition. Produce it! No, by Him who created all contingent things
through a word from Himself! He has never made this Revelation dependent
upon or subject to any matter or condition, for in that court conditions
are utterly nonexistent, and no one but His own Self encompasses or
knows Him.

\subsection{230}\label{section-725}

\begin{otherlanguage}{persian}

ميفرمايد آياتش دليل بر اوست و ظهور و وجودش مثبت نفسش مع‌ذلك بكلمات موهومه
و اشارات نالايقه حقّ منيع را محدود نموده‌اند فويل لهم و بما اكتسبت ايديهم
كلمات الهى را بهواهاى خود تحريف منمائيد و تفسير مكنيد چه كه احدى بمقصود
آن فايز نشود چنانچه در احكام ظاهريّه رؤساى بيان را امر فرموده‌اند كه از
جناب آقا سيّد حسين عليه بهآء اللّه كه تلقاء عرش بود سؤال نمايند ميفرمايند
چه كه شما عارف بمراد اللّه نيستيد مع‌ذلك شما معرض باللّه را نفس اللّه
دانسته‌ايد

\end{otherlanguage}

Mīfarmāyad āyātash dalīl bar ūst va ẓohūr va vojūdash mothbet-e nafsash.
Maʿa-dhālik be-kalemāt-e mowhūme va eshārāt-e nālāyeqe Ḥaqq-e Manīʿ rā
maḥdūd namūde-and. Fa-waylun lahum wa-bi-mā iktasabat aydīhim. Kalemāt-e
elāhī rā be-havāhā-ye khod taḥrīf manamāyīd va tafsīr makonīd, che ke
aḥadī be-maqṣūd-e ān fāʾez nashavad. Chenānche dar aḥkām-e ẓāheriyye
roʾasā-ye Bayān rā amr farmūde-and ke az Janāb-e Āqā Seyyed Ḥoseyn,
ʿalayhi bahāʾu Allāh, ke telqāʾ-e ʿarsh būd soʾāl namāyand. Mīfarmāyand:
che ke shomā ʿāref be-morād-e Allāh nīstīd. Maʿa-dhālik shomā moʿreḍ-e
bi-Allāh rā nafs-e Allāh dāneste-īd.

He states that His verses are the proof of Him, and His appearance and
existence establish His own Self. Nevertheless, they have confined the
unassailable Truth within fanciful words and unworthy allusions. Woe to
them and for what their hands have earned! Do not distort or interpret
the words of God according to your own desires, for no one can attain
their intended meaning. Thus, concerning outward ordinances, He
instructed the leaders of the Bayán to ask Áqá Siyyid Ḥusayn, upon him
be the glory of God, who was in the presence of the Throne. He states:
``For you do not know what God intends.'' Nevertheless, you have
regarded one who turns away from God as the very Self of God.

\subsection{231}\label{section-726}

\begin{otherlanguage}{persian}

از اين گذشته خاتم النّبيّين از محكمات آيات كتاب ربّ العالمين است اگر هزار و
دويست و هفتاد سنه احدى بمعنى آن مطّلع شد حال هم معانى كلمات منزله بيان را
نفسى دون اللّه ميتواند ادراك نمايد بسيار تفكّر نمائيد در آنچه ذكر ميشود و
از قلم باذن اللّه جارى ميگردد كه شايد از صرير قلم اين ضعيف بجمال قدم فايز
شوى انّه ليهدى من یشآء الى صراط مستقيم

\end{otherlanguage}

Az īn goẕashte, Khātam al-Nabiyyīn az moḥkamāt-e āyāt-e Ketāb-e Rabb
al-ʿĀlamīn ast. Agar hezār va devīst va haftād sane aḥadī be-maʿnī-ye ān
moṭṭaleʿ shod, ḥāl ham maʿānī-ye kalemāt-e monzale-ye Bayān rā nafsī dūn
Allāh mītavānad edrāk namāyad. Besyār tafakkor namāyīd dar ānche ẕekr
mīshavad va az qalam be-eẕn-e Allāh jārī mīgardad ke shāyad az ṣarīr-e
qalam-e īn ḍaʿīf be-Jamāl-e Qedam fāʾez shavī. Innahu la-yahdī man
yashāʾu ilā ṣirāṭin mustaqīm.

Furthermore, ``the Seal of the Prophets'' is among the unequivocal
verses of the Book of the Lord of the worlds. If anyone understood its
meaning over the course of one thousand two hundred and seventy years,
then now, too, someone other than God can comprehend the meanings of the
revealed words of the Bayán. Reflect deeply on what is being stated and
flows from the Pen by God's permission, so that perhaps, through the
scratching of this feeble one's pen, you may attain the Ancient Beauty.
He indeed guides whomever He wills to a straight path.

\subsection{232}\label{section-727}

\begin{otherlanguage}{persian}

و ديگر معلوم نشد مقصود شما از تكميل چه بوده اگر مقصود تكميل آيات الهيّه
است كه آن در حين نزول كامل و ابدا در آن ساحت نقص راه نداشته و نخواهد
داشت ما نزل من عنده كان كاملاً من قبل و يكون كاملاًمن بعد و فى‌الحين و اگر
مقصود تكميل نفوس است كه نفوسى مشاهده شده هرچه مدّت طول كشيد ناقص‌تر ظاهر
شده‌اند چنانچه در منتهى دوره فرقان بعقيده شما منتهى كمال رؤساى آن عصر آن
بود كه جميع الواح اللّه را طرد نموده و برسل الهى وارد آوردند آنچه
شنيده‌ايد و يا از تكميل تكميل اعمال ظاهريه مقصود است اگر اينست كه مشاهده
شده اكثرى را همان اعمال سبب احتجاب گشته چنانچه ملاحظه شد كه عاملين فرقان
چون ظهور نيّر اعظم را مخالف اعمال خود مشاهده نمودند فتوى بر قتلش دادند و
سفك دم مطهّرش نمودند و هم‌‌چنين اهل انجيل و هم‌‌چنين اهل تورئة تا منتهى شود
ظهورات ببديع اوّل

\end{otherlanguage}

Va dīgar maʿlūm nashod maqṣūd-e shomā az takmīl che būde. Agar maqṣūd
takmīl-e āyāt-e elāhiyye ast ke ān dar ḥīn-e nozūl kāmel va abadan dar
ān sāḥat naqṣ rāh nadāshte va nakhvāhad dāsht. Mā nazala min ʿindihi
kāna kāmilan min qablu wa-yakūnu kāmilan min baʿdu wa-fī al-ḥīn. Va agar
maqṣūd takmīl-e nofūs ast ke nofūsī moshāhede shode harche moddat ṭūl
keshīd nāqeṣtar ẓāher shode-and. Chenānche dar montahā-ye dowre-ye
Forqān be-ʿaqīde-ye shomā montahā-ye kamāl-e roʾasā-ye ān ʿaṣr ān būd ke
jamīʿ-e alvāḥ-e Allāh rā ṭard namūde va be-rosol-e elāhī vāred āvardand
ānche shenīde-īd. Va yā az takmīl takmīl-e aʿmāl-e ẓāheriyye maqṣūd ast?
Agar īnast ke moshāhede shode aktharī rā hamān aʿmāl sabab-e eḥtejāb
gashte. Chenānche molāḥeẓe shod ke ʿāmelīn-e Forqān chūn ẓohūr-e
Nayyِer-e Aʿẓam rā mokhālef-e aʿmāl-e khod moshāhede namūdand fatvā bar
qatlash dādand va safk-e dam-e moṭahharash namūdand, va hamchenīn ahl-e
Enjīl va hamchenīn ahl-e Towrāt tā montahī shavad ẓohūrāt be-Badīʿ-e
Avval.

Moreover, it is not clear what you mean by completion. If you mean the
completion of the divine verses, they are complete at the moment of
revelation, and imperfection has never entered and will never enter that
realm. Whatever has descended from Him was complete before and will be
complete afterward and at this very moment. If you mean the perfecting
of souls, souls have been seen to become increasingly imperfect the more
time passed. Thus, at the end of the dispensation of the Qur'an, which
according to your belief was the height of perfection, the utmost
perfection of the leaders of that age consisted in rejecting all the
Tablets of God and inflicting upon the divine Messengers what you have
heard. Or does completion mean the perfecting of outward deeds? If so,
those very deeds have been seen to become a veil for most people. Thus
it was observed that when those who practiced the ordinances of the
Qur'an saw that the appearance of the Most Great Luminary was contrary
to their practices, they issued a decree for His death and shed His pure
blood. Likewise did the people of the Gospel, and likewise the people of
the Torah, and so on through the Revelations, back to the First
Originator.

\subsection{233}\label{section-728}

\begin{otherlanguage}{persian}

بدانكه تكميل كل‌‌شیء منوط بقول و تصديق نقطه حقيقت بوده بشنو كه نقطه بيان
روح ماسواه فداه چه ميفرمايد كه شايد از كوثر بيانات رحمن از غبار تيره
اشارات مظاهر شيطان مطهّر شوى اگرچه يقين است كه كلمات اللّه در تو اثر
ننموده و نخواهد نمود و لكن لحاظ اللّه در اكثر بيانات متوجّه افئده مخلصين
از عباد خود بوده چه كه آن نفوس مقدّسه منقطعه روايح مسك معانى را از
فرسنگهاى بعيده استشمام نمايند و بكلمات الهى از كلمات عالمين منقطع شوند
قوله تبارك ذكره فانّه لو يجعل ما على الأرض نبيّاً ليكوننّ ابنياء عند اللّه از
اين بيان رحمن حقيقت تكميل را ادراك كن كه شايد متنبّه شوى و خزف اشارات
ببحر لؤلؤ معانى الهى هديه نفرستى و موقن شوى باينكه كلّ امور از حقّ ظاهر و
بحقّ راجع و بتصديق او ثابت و محقّق

\end{otherlanguage}

Bedān ke takmīl-e koll-e shayʾ manūṭ be-qowl va taṣdīq-e Noqṭe-ye
Ḥaqīqat būde. Beshnow ke Noqṭe-ye Bayān, rūḥ-e mā-sevāh fedāh, che
mīfarmāyad, ke shāyad az kawthar-e bayānāt-e Raḥmān az ghobār-e tīre-ye
eshārāt-e maẓāher-e shayṭān moṭahhar shavī. Agarche yaqīn ast ke
kalemāt-e Allāh dar to athar nanamūde va nakhvāhad namūd, va lāken
leḥāẓ-e Allāh dar akthar-e bayānāt motavajjeh-e afʾede-ye mokhleṣīn az
ʿebād-e khod būde, che ke ān nofūs-e moqaddase-ye monqaṭeʿe ravāyeḥ-e
meshk-e maʿānī rā az farsanghā-ye baʿīde esteshmām namāyand va
be-kalemāt-e elāhī az kalemāt-e ʿālamīn monqaṭeʿ shavand. Qawluhu
tabāraka dhikruhu: fa-innahu law yajʿalu mā ʿalā al-arḍi nabiyyan
la-yakūnanna {[}corrupted text{]} ʿinda Allāh. Az īn bayān-e Raḥmān
ḥaqīqat-e takmīl rā edrāk kon ke shāyad motanabbeh shavī va khazaf-e
eshārāt be-baḥr-e loʾloʾ-e maʿānī-ye elāhī hadiyye naferestī va mūqen
shavī be-īnke koll-e omūr az Ḥaqq ẓāher va be-Ḥaqq rājeʿ va be-taṣdīq-e
ū thābet va moḥaqqaq.

Know that the perfection of all things has depended upon the word and
confirmation of the Point of Truth. Listen to what the Point of the
Bayán---may the spirit of all besides Him be His sacrifice---says, so
that perhaps, through the heavenly stream of the All-Merciful's
utterances, you may be cleansed of the dark dust of the allusions of the
manifestations of Satan. Although it is certain that the words of God
have had and will have no effect upon you, in most of His utterances
God's gaze has been directed toward the hearts of the sincere among His
servants, for those sanctified and detached souls inhale the musk-laden
fragrances of meaning from many leagues away, and through the words of
God detach themselves from the words of all the worlds. His words,
blessed be His mention, are: ``For if He were to make all that is on
earth a prophet, they would assuredly be {[}corrupted original text{]}
in the sight of God.'' From this utterance of the All-Merciful,
comprehend the reality of perfection, so that perhaps you may be
awakened, may not offer the potsherds of allusions as a gift to the
ocean of the pearls of divine meanings, and may become certain that all
things proceed from God, return to God, and are established and realized
through His confirmation.

\subsection{234}\label{section-729}

\begin{otherlanguage}{persian}

قدرى تفكّر نما مع آنكه منتهى رتبه تكميل هياكل بشريّه برتبه نبوّت است و
اكثرى اين رتبه بلند اعلى را محال دانسته‌اند كه نفسى بانمقام فايز شود
مع‌ذلك ميفرمايند اگر آن ظهور اعظم اراده فرمايد جميع ما على الأرض را باين
رتبه كه اكمل مراتب بوده نزد كلّ فايز ميفرمايد فو‌عمره لمثل هذا الرّبّ ينبغى
ان نفدى انفسنا و ارواحنا ثمّ ذواتنا و حقايقنا و هذا قليل عند ظهور عطاياه
و بروز مواهبه الّتى احاطت الثّقلين

\end{otherlanguage}

Qadrī tafakkor namā. Maʿ-e ānke montahā-ye rotbe-ye takmīl-e hayākel-e
bashariyye be-rotbe-ye nobovvat ast va aktharī īn rotbe-ye boland-e aʿlā
rā moḥāl dāneste-and ke nafsī be-ān maqām fāʾez shavad, maʿa-dhālik
mīfarmāyand agar ān Ẓohūr-e Aʿẓam erāde farmāyad jamīʿ-e mā ʿalā al-arḍ
rā be-īn rotbe ke akmal-e marāteb būde nazd-e koll fāʾez mīfarmāyad.
Fa-wa-ʿumrihi la-mithli hādhā al-Rabbi yanbaghī an nafdī anfusanā
wa-arwāḥanā thumma dhawātanā wa-ḥaqāʾiqanā, wa-hādhā qalīlun ʿinda
ẓuhūri ʿaṭāyāhu wa-burūzi mawāhibihi allatī aḥāṭati al-thaqalayn.

Reflect a little. Although the highest degree of perfection for human
beings is the rank of prophethood, and most have considered it
impossible for any soul to attain this lofty and exalted station, He
nevertheless states that if that Most Great Revelation should will it,
He would cause all who are on earth to attain this rank, which everyone
regards as the most perfect of stations. By His life! For such a Lord we
should sacrifice our selves and our souls, and then our very beings and
realities. Yet this is little in the face of the manifestation of His
gifts and the emergence of His bestowals, which have encompassed both
humankind and jinn.

\subsection{235}\label{section-730}

\begin{otherlanguage}{persian}

اگر نفسى بدوام ملك و ملكوت عمر نمايد و بكلّ اعمال عامل شود و در ظهور اقلّ
من آن توقّف نمايد جميع اعمال او كه در زمن لايحصى عامل بوده هباء ميشود
بلكه وجود و كينونتش لاشیء ميشود تا چه رسد بتكميل او و هم‌‌چنين اگر در
اينمدّت نعوذ باللّه تارك كلّ اعمال بوده و در حين ظهور فايز بقول بلى شود
جميع اعمال در حقّ او ثبت و عند اللّه كامل بوده و خواهد بود

\end{otherlanguage}

Agar nafsī be-davām-e molk va malakūt ʿomr namāyad va be-koll-e aʿmāl
ʿāmel shavad va dar ẓohūr aqall min ān tavaqqof namāyad, jamīʿ-e aʿmāl-e
ū ke dar zaman-e lā yuḥṣā ʿāmel būde habāʾ mīshavad, balke vojūd va
kaynūnatash lā shayʾ mīshavad, tā che resad be-takmīl-e ū. Va hamchenīn
agar dar īn moddat, naʿūdhu bi-Allāh, tārek-e koll-e aʿmāl būde va dar
ḥīn-e ẓohūr fāʾez be-qowl-e balā shavad, jamīʿ-e aʿmāl dar ḥaqq-e ū
thabt va ʿinda Allāh kāmel būde va khvāhad būd.

If a soul should live as long as the earthly dominion and the heavenly
kingdom endure, and perform every prescribed deed, yet hesitate for less
than an instant at the Revelation, all the deeds he has performed over
that immeasurable time become dust. Indeed, his very existence and being
become nothing, let alone his perfection. Likewise, if throughout this
period he has---God forbid---neglected every deed, but at the moment of
Revelation attains the declaration ``Yes,'' all deeds are recorded to
his credit, and he has been and will be perfect in the sight of God.

\subsection{236}\label{section-731}

\begin{otherlanguage}{persian}

بر تكميل دوره احدى الّا اللّه مطّلع نه بسا كه در اقرب من آن تكميل ظاهر و
بسا كه در الف سنه و فوق آن تكميل نميشود و اين بسى واضح است كه تكميل هر
شیء بتصديق اللّه بوده نه بطول مدّت و ظهور اوهام ناس نفس ظهور بعد بنفسه
دليل بر تكميل ما ظهر فى ظهور القبل بوده بحقّ ناظر باشيد نه بطول مدّت و
زمان چه كه حقّ قادر است در آنى كلّ اوان را منقضى فرمايد و يا آنكه يك آن را
در صدهزار سنه امتداد دهد له الخلق و الأمر كلّ بأمره يعملون

\end{otherlanguage}

Bar takmīl-e dowre aḥadī illā Allāh moṭṭaleʿ na. Basā ke dar aqrab min
ān takmīl ẓāher va basā ke dar alf sane va fowq-e ān takmīl namīshavad.
Va īn basī vāḍeḥ ast ke takmīl-e har shayʾ be-taṣdīq-e Allāh būde na
be-ṭūl-e moddat va ẓohūr-e owhām-e nās. Nafs-e ẓohūr-e baʿd be-nafsehi
dalīl bar takmīl-e mā ẓahara fī ẓuhūr al-qabl būde. Be-Ḥaqq nāẓer bāshīd
na be-ṭūl-e moddat va zamān, che ke Ḥaqq qāder ast dar ānī koll-e avān
rā monqaḍī farmāyad va yā ānke yek ān rā dar ṣad-hezār sane emtedād
dahad. Lahu al-khalqu wa-al-amru, kullun bi-amrihi yaʿmalūn.

No one but God knows when a dispensation is complete. Completion may
appear in less than an instant, or it may not occur in a thousand years
or more. It is abundantly clear that the completion of every thing
depends upon God's confirmation, not upon the length of time or the
emergence of people's imaginings. The subsequent Revelation is itself
the proof of the completion of what appeared in the preceding
Revelation. Look to God, not to the length of the period or the passage
of time, for God is able to bring all ages to an end in a single
instant, or to extend one instant over a hundred thousand years. His are
creation and command; all act at His command.

\subsection{237}\label{section-732}

\begin{otherlanguage}{arabic}

در اينمقام چند فردى از لسان ابدع امنع در عراق نازل لاجل تنبّه ذكر ميشود

\end{otherlanguage}

Dar īn maqām chand fardī az Lesān-e Abdaʿ-e Amnaʿ dar ʿErāq nāzel,
li-ajl-e tanabboh ẕekr mīshavad.

At this point, a few couplets revealed in Iraq by the most wondrous,
most exalted Tongue are cited for your awakening.

\subsection{238}\label{section-733}

\begin{otherlanguage}{persian}

سائلى مر عارفى را گفت كى تو بر اسرار الهى بردۀ پى وى تو از خمر عنايت
گشته مست هيچ يادت آيد از روز الست گفت ياد آيد مرا آن صوت و گفت كو بدى
بود و نباشد اين شگفت هست در گوشم همى آواى او وان صداى روح و جان افزاى او
عارفى ديگر كه بر تر رفته بود درّ اسرار الهى سفته بود گفت آن روز خدا آخر
نشد ما در آن روزيم آن قاصر نشد

\end{otherlanguage}

Sāʾelī mar ʿārefī rā goft: key to bar asrār-e elāhī borde-ī pey? Vey to
az khamr-e ʿenāyat gashte mast, hīch yādat āyad az rūz-e alast? Goft:
yād āyad marā ān ṣowt va goft, k-ū {[}corrupted text{]} būd va nabāshad
īn shegeft. Hast dar gūsham hamī āvā-ye ū, v-ān ṣadā-ye rūḥ va
jān-afzā-ye ū. ʿĀrefī dīgar ke bartar rafte būd, dorr-e asrār-e elāhī
softe būd, goft: ān rūz-e Khodā ākhar nashod, mā dar ān rūzīm, ān qāṣer
nashod.

A questioner said to a mystic: ``When did you come to know the divine
mysteries? O you who are drunk with the wine of grace, do you remember
anything of the day of `Am I not?'\,'' He said: ``I remember that voice
and speech, which was {[}corrupted original text{]}, and this is no
wonder. His voice is ever in my ear, that spirit- and soul-nourishing
call of His.'' Another mystic, who had ascended higher and had pierced
the pearl of the divine mysteries, said: ``That day of God has not
ended. We are in that day; it has not been cut short.''

\subsection{239}\label{section-734}

\begin{otherlanguage}{persian}

اين كلمات مخصوص آن ذكر شد كه معلوم شود ادراك اين امور منوط بعرفان عارفين
است هر نفسى اين مراتب را ادراك ننمايد و اطّلاع نيابد عجب است كه تو خود را
عارف ميدانى و از كلمات عرفا هم اطّلاع ندارى و الّا در اينمقام اعتراض
نمينمودى چه كه عرفا برآنند كه اسماء الهى در كلّ حين مؤثّر چه كه تعطيل جايز
نه و لذا گفته‌اند اسم مميت در كلّ حين كلّ موجودات را از وجود اضافى بعدم
اضافى راجع مينمايد و هم‌‌چنين اسم محيى و مبعث در كلّ ان كلّ اكوان را از عدم
اضافى بوجود اضافى جلوه ميدهد چنانچه آيه مباركه بل هم فى لبس من خلق جديد
را در اين مقام تفسير نموده‌اند

\end{otherlanguage}

Īn kalemāt makhṣūṣ-e ān ẕekr shod ke maʿlūm shavad edrāk-e īn omūr manūṭ
be-ʿerfān-e ʿārefīn ast. Har nafsī īn marāteb rā edrāk nanamāyad va
eṭṭelāʿ nayābad. ʿAjab ast ke to khod rā ʿāref mīdānī va az kalemāt-e
ʿorafā ham eṭṭelāʿ nadārī, va illā dar īn maqām eʿterāḍ namīnamūdī. Che
ke ʿorafā bar ānand ke asmāʾ-e elāhī dar koll-e ḥīn moʾathther, che ke
taʿṭīl jāyez na. Va leẕā gofte-and esm-e Momīt dar koll-e ḥīn koll-e
mowjūdāt rā az vojūd-e eḍāfī be-ʿadam-e eḍāfī rājeʿ mīnamāyad va
hamchenīn esm-e Moḥyī va Mobʿeth dar koll-e ān koll-e akvān rā az
ʿadam-e eḍāfī be-vojūd-e eḍāfī jelve mīdahad, chenānche āye-ye
mobārake-ye bal hum fī labsin min khalqin jadīd rā dar īn maqām tafsīr
namūde-and.

These words were cited specifically to make it clear that comprehending
these matters depends upon the insight of the mystics. Not every soul
can comprehend or become aware of these degrees. It is strange that you
consider yourself a mystic, yet do not even know the words of the
mystics; otherwise, you would not have raised an objection here. For the
mystics hold that the divine names are effective at every moment, since
their suspension is inadmissible. Thus they have said that the name
``the Giver of Death'' returns all beings at every moment from relative
existence to relative nonexistence, and likewise the names ``the Giver
of Life'' and ``the Raiser'' bring all created things at every instant
from relative nonexistence into relative existence. It is in this sense
that they have interpreted the blessed verse: ``Rather, they are in
confusion concerning a new creation.''

\subsection{240}\label{section-735}

\begin{otherlanguage}{persian}

در اينصورت آن خلق بعد در حين قبل موجود نه و هم‌‌چنين در حين قبل آن خلق
موجود حين بعد معلوم نه پس چگونه تصوّر تكميل شیء از اشیاء نموده‌اند چه كه
آن شیء كه در اين حين موجود است در حين بعد نيست چه كه تجدّد و حدوث است از
براى اين شیء چگونه تكميل دست ميدهد اگر بگوئى در همان حين محدثه در حين
ميشود اين طول مدّت لازم نداشت و ندارد و اگر ناس قابل ادراك معانى كلمات
الهيّه ميبودند هرآينه در اينمقام تفصيلى ذكر ميشد تا آنكه از جميع عرفان
قبل و بعد منقطع شده و بشريعه عرفان بديع كه اليوم از يمين عرش رحمن جارى
است توجّه نمايند

\end{otherlanguage}

Dar īn ṣūrat ān khalq-i baʿd dar ḥīn-i qabl mawjūd na va ham-chunīn dar
ḥīn-i qabl ān khalq-i mawjūd-i ḥīn-i baʿd maʿlūm na. Pas chigūna
taṣavvur-i takmīl-i shayʾ az ashyāʾ namūda-and, chi ki ān shayʾ ki dar
īn ḥīn mawjūd ast dar ḥīn-i baʿd nīst, chi ki tajaddud va ḥudūth ast az
barāy-i īn shayʾ. Chigūna takmīl dast mīdahad? Agar bigūʾī dar hamān
ḥīn-i muḥdatha dar ḥīn mīshavad, īn ṭūl-i muddat lāzim nadāsht va
nadārad. Va agar nās qābil-i idrāk-i maʿānī-yi kalimāt-i ilāhiyya
mībūdand, har-āyina dar īn maqām tafṣīlī dhikr mīshud tā ānki az jamīʿ-i
ʿirfān-i qabl va baʿd munqaṭiʿ shuda va bi-sharīʿa-yi ʿirfān-i badīʿ ki
al-yawm az yamīn-i ʿarsh-i Raḥmān jārī ast tavajjuh namāyand.

In this case, the subsequent creation did not exist at the preceding
moment, and likewise, at the preceding moment, the creation existing at
the subsequent moment was not known. How, then, have they conceived of
the perfection of any single thing? For the thing that exists at this
moment does not exist at the next, since this thing undergoes renewal
and origination. How can perfection be attained? If you say that it
takes place at once, in that very moment of origination, this long
duration was not and is not necessary. If people were capable of
comprehending the meanings of the divine words, a detailed explanation
would certainly be given here, so that they might detach themselves from
all knowledge, past and future, and turn toward the waters of wondrous
new knowledge that today flow from the right hand of the throne of the
All-Merciful.

\subsection{241}\label{section-736}

\begin{otherlanguage}{persian}

بشنو ذكر اين فانى را و امر را بر خود مشتبه منما كلمات الهى را احدى ادراك
ننموده و نخواهد نمود چه كه مفتاح ابواب او بيد مقتدر قدير بوده و تا ختم
اناء معارف الهيّه را يد قدرت نگشايد احدى بر آن مطّلع نشده و نخواهد شد بلكه
عرف آن را هم نيافته و نخواهند يافت چنانچه تا يد قدرت ظهور قبل ختم كنائز
علوم فرقان را نگشود احدى ادراك ننمود بارى تكميل در دست قدرت حقّ بوده و
خواهد بود

\end{otherlanguage}

Bishnaw dhikr-i īn fānī rā va amr rā bar khud mushtabih manamā.
Kalimāt-i ilāhī rā aḥadī idrāk nanamūda va nakhwāhad namūd, chi ki
miftāḥ-i abvāb-i ū bi-yad-i Muqtadir-i Qadīr būda va tā khatm-i ināʾ-i
maʿārif-i ilāhiyya rā yad-i qudrat nagushāyad aḥadī bar ān muṭṭaliʿ
nashuda va nakhwāhad shud, balki ʿarf-i ān rā ham nayāfta va nakhwāhand
yāft. Chunānchi tā yad-i qudrat-i ẓuhūr-i qabl khatm-i kanāʾiz-i ʿulūm-i
Furqān rā nagushūd aḥadī idrāk nanamūd. Bārī takmīl dar dast-i qudrat-i
Ḥaqq būda va khwāhad būd.

Listen to the words of this transient one, and do not confuse the matter
for yourself. No one has comprehended or will comprehend the divine
words, for the key to their gates has been in the hand of the Almighty,
the All-Powerful. Until the hand of power breaks the seal on the vessel
of divine knowledge, no one has gained or will gain access to it;
indeed, they have not even caught, nor will they catch, its fragrance.
Likewise, until the hand of power of the preceding Revelation broke the
seal on the treasuries of the knowledge of the Qur'an, no one
comprehended them. In short, perfection has been and will remain in the
hand of God's power.

\subsection{242}\label{section-737}

\begin{otherlanguage}{persian}

تكميل هر دوره بظهور بعد بوده و در حين ظهور تكميل حقيقى واقع و لكن
استطاعت ادراك آن را امثال آن نفوس ضعيفه نداشته و ندارند چنانچه ادراك آن
انقضاى خمسين الف سنه در ساعت واحده نشده و اگر نفسى هم ادّعا نموده بر حسب
تقليد بوده الّا من علّمه اللّه من جواهر علمه بسا از علوم ربّانى كه در سماء
علم لدنّى الهى وديعه مانده ديگر تا چه وقت و زمان از افق امكان مظاهر قدسيّه
و مطالع منيعه ظاهر شوند و مستعد ادراك گردند و از سماء مشيّت على قدر مقدور
نازل شود الأمر بيده يفعل ما یشآء و لا يسئل عمّا شآء

\end{otherlanguage}

Takmīl-i har dawra bi-ẓuhūr-i baʿd būda va dar ḥīn-i ẓuhūr takmīl-i
ḥaqīqī vāqiʿ, va lākin istiṭāʿat-i idrāk-i ān rā amthāl-i ān nufūs-i
ḍaʿīfa nadāshta va nadārand, chunānchi idrāk-i ān inqiḍā-yi khamsīn alf
sana dar sāʿat-i vāḥida nashuda va agar nafsī ham iddiʿā namūda bar
ḥasab-i taqlīd būda, illā man ʿallamahu Allāhu min jawāhiri ʿilmihi.
Basā az ʿulūm-i rabbānī ki dar samā-yi ʿilm-i ladunnī-yi ilāhī vadīʿa
mānda. Dīgar tā chi vaqt va zamān az ufuq-i imkān maẓāhir-i qudsiyya va
maṭāliʿ-i manīʿa ẓāhir shavand va mustaʿidd-i idrāk gardand va az
samā-yi mashiyyat ʿalā qadr-i maqdūr nāzil shavad. Al-amru bi-yadihi
yafʿalu mā yashāʾu wa-lā yusʾalu ʿammā shāʾa.

The perfection of each cycle has come through the subsequent Revelation,
and true perfection occurs at the moment of revelation. Yet the likes of
those feeble souls have not had and do not have the capacity to
comprehend it, just as that passage of fifty thousand years in a single
hour has not been comprehended. If anyone has claimed to do so, it has
been by imitation, except those whom God has taught from the gems of His
knowledge. Much divine learning remains deposited in the heaven of God's
innate knowledge. It remains to be seen at what time and season holy
Manifestations and exalted Daysprings will appear from the horizon of
contingent existence, and receptivity to understanding will arise, and
it will descend from the heaven of the divine will in the measure
ordained. The command is in His hand. He does whatever He wills and is
not questioned about what He has willed.

\subsection{243}\label{section-738}

\begin{otherlanguage}{persian}

و اينكه نوشته بايد از فرمايش ربّ اعلى بقدر جوى تخلّف ننمايد معلوم نشد كه
اين حكم حكم شما است بايد عمل شود يا حكم نقطه بيان نقطه بيان که ميفرمايند
عمل او ميزان كلّ بوده و بهر چه آن وجود مبارك عامل شود بايد كلّ عامل شوند
گويا اصلا بيان را نديده و اطّلاع هم ندارى بتعليم شيطان ناطقى و بتحريك او
متحرّك آخر

\end{otherlanguage}

Va īnki nivishta bāyad az farmāyish-i Rabb-i Aʿlā bi-qadr-i jawī
takhalluf nanamāyad, maʿlūm nashud ki īn ḥukm ḥukm-i shumā ast bāyad
ʿamal shavad yā ḥukm-i Nuqṭa-yi Bayān. Nuqṭa-yi Bayān ki mīfarmāyand
ʿamal-i ū mīzān-i kull būda va bi-har chi ān vujūd-i mubārak ʿāmil
shavad bāyad kull ʿāmil shavand. Gūyā aṣlan Bayān rā nadīda va iṭṭilāʿ
ham nadārī. Bi-taʿlīm-i shayṭān nāṭiqī va bi-taḥrīk-i ū mutaḥarrik,
ākhar.

As for your statement that he must not deviate from the words of the
Exalted Lord by so much as a barleycorn, it is unclear whether this is
your decree that must be obeyed or the decree of the Point of the Bayán.
The Point of the Bayán states that His action is the standard for all,
and that whatever that blessed Being does, everyone must do. It seems
you have never seen the Bayán at all and have no knowledge of it. After
all, you speak at Satan's instruction and move at his prompting.

\subsection{244}\label{section-739}

\begin{otherlanguage}{persian}

اى نفسيكه بعرفان موسومى و خود را عارف دانسته از خود هم يك اظهار وجودى
بنما و خود را مثل اغنام مكن كه هر راعى بهر جهت اراده نمايد شما را ببرد و
اىكاش راعى انسان ميشد راعى ذئب پيدا نموده‌ايد و جميع عرفان و ايقان و
ثمرات توحيدت را در ساعت اوّل ملاقات بلع نموده و حال جسدى ميّت و لاشیء
مشهودى و بقسمى هم ميّت و لاشیء ملاحظه ميشوى كه بصدهزار صور هم مهتزّ نشوى و
بحركت نيائى بحقّ معبود يكتا كه دلم بر تو و امثال تو ميسوزد چه كه كينه و
بغضا در اين قلوب نبوده و نيست و آنچه ذكر شده و ميشود خالصاً لوجه اللّه
بوده و خواهد بود

\end{otherlanguage}

Ay nafsī ki bi-ʿirfān mawsūmī va khud rā ʿārif dānista, az khud ham yak
iẓhār-i vujūdī binamā va khud rā mithl-i aghnām makun ki har rāʿī bi-har
jahat irāda namāyad shumā rā bibarad. Va ay kāsh rāʿī insān mīshud. Rāʿī
dhiʾb paydā namūda-īd va jamīʿ-i ʿirfān va īqān va thamarāt-i tawḥīdat
rā dar sāʿat-i avval-i mulāqāt balʿ namūda va ḥāl jasadī mayyit va
lā-shayʾ mashhūdī. Va bi-qismī ham mayyit va lā-shayʾ mulāḥiẓa mīshavī
ki bi-ṣad-hizār ṣūr ham muhtazz nashavī va bi-ḥarakat nayāʾī. Bi-Ḥaqq-i
Maʿbūd-i Yaktā ki dilam bar tu va amthāl-i tu mīsūzad, chi ki kīna va
baghḍā dar īn qulūb nabūda va nīst va ānchi dhikr shuda va mīshavad
khāliṣan li-wajhi Allāh būda va khwāhad būd.

O you who are known for spiritual knowledge and consider yourself a
mystic! Show some sign of an independent existence, and do not make
yourself like sheep, to be led by any shepherd in whatever direction he
wishes. If only the shepherd were human! You have found a wolf for a
shepherd, and at the very first meeting he swallowed all your spiritual
knowledge and certitude and the fruits of your belief in divine unity.
Now you appear as a dead body, a thing of nothing. Indeed, you are seen
to be so dead and insubstantial that even a hundred thousand trumpet
blasts would not stir you or bring you to move. By the one true God who
is worshipped, my heart burns for you and those like you, for malice and
hatred have never been and are not in these hearts. Whatever has been
and is being said has been and will be solely for the sake of God.

\subsection{245}\label{section-740}

\begin{otherlanguage}{persian}

و ديگر طلعت ابهى روح ما سواه له الفدآء بجميع اعمال عاملند لو انت من
المتبصّرين چن دسنه ميشود كه اين عبد ليلاً و نهاراً معاشر بوده فواللّه الّذى
لا اله الّا هو كه امرى كه فى‌الجمله مكروه باشد از اين جمال مبارك ديده نشده
جميع اشیاء بر تقديس ذات مقدّسشان شهادت داده و ميدهند در ليالى و ایّام جز
ذكر اللّه و ارتفاع كلمته از آن بيت مطهّر مقدّس استماع نشده خود شما اگر قدرى
بانصاف بيائيد شهادت ميدهيد كه مقدّسند از اوهام ما سواه و منزّهند از ادراك
ما دونه

\end{otherlanguage}

Va dīgar Ṭalʿat-i Abhā, rūḥu mā siwāhu lahu al-fidāʾ, bi-jamīʿ-i aʿmāl
ʿāmiland, law anta mina al-mutabaṣṣirīn. {[}corrupted text{]} mīshavad
ki īn ʿabd laylan va nahāran muʿāshir būda. Fa-wa-Allāhi alladhī lā
ilāha illā huwa ki amrī ki fī al-jumla makrūh bāshad az īn Jamāl-i
Mubārak dīda nashuda. Jamīʿ-i ashyāʾ bar taqdīs-i dhāt-i muqaddasishān
shahādat dāda va mīdihand. Dar layālī va ayyām juz dhikr-i Allāh va
irtifāʿ-i kalimatihi az ān bayt-i muṭahhar-i muqaddas istimāʿ nashuda.
Khud-i shumā agar qadrī bi-inṣāf biyāʾīd shahādat mīdihīd ki muqaddasand
az awhām-i mā siwāh va munazzahand az idrāk-i mā dūnih.

Furthermore, the All-Glorious Countenance---may the soul of all else be
sacrificed for Him---observes every prescribed deed, if you are among
those who see clearly. It has been {[}corrupted original text{]} that
this servant has been in His company night and day. By God, besides whom
there is no god, nothing even slightly objectionable has been seen from
this Blessed Beauty. All things have testified and continue to testify
to the sanctity of His holy Being. Throughout the nights and days,
nothing has been heard from that pure and holy house except the
remembrance of God and the exaltation of His Word. You yourself, if you
show a little fairness, will testify that He is sanctified beyond the
imaginings of all others and exalted beyond the comprehension of all
besides Him.

\subsection{246}\label{section-741}

\begin{otherlanguage}{persian}

كدام عمل فوق انفاق جان است در سبيل رحمن آيا نشنيده‌ايد كه بيست سنه در
سبيل الهى بر ايشان چه وارد شده اگر چه آن مفترى خبيث اعمال خود و مرشدش را
بجمال اطهر اقدس نسبت داده و تو هم من غير دليل و برهان از او پذيرفته و
لكن نه واللّه بقول او و امثال او ذيل مطهّر نيالايد

\end{otherlanguage}

Kudām ʿamal fawq-i infāq-i jān ast dar sabīl-i Raḥmān? Āyā nashanīda-īd
ki bīst sana dar sabīl-i ilāhī bar īshān chi vārid shuda? Agar chi ān
muftarī-yi khabīth aʿmāl-i khud va murshidash rā bi-Jamāl-i Aṭhar-i
Aqdas nisbat dāda va tu ham min ghayr-i dalīl va burhān az ū padhīrufta,
va lākin na, va-Allāh, bi-qawl-i ū va amthāl-i ū dhayl-i muṭahhar
nayālāyad.

What deed is greater than offering one's life in the path of the
All-Merciful? Have you not heard what He has suffered for twenty years
in the path of God? That wicked slanderer has indeed attributed his own
deeds and those of his guide to the most pure, most holy Beauty, and you
have accepted them from him without evidence or proof. But no, by God!
The words of that man and his like cannot stain His pure robe.

\subsection{247}\label{section-742}

\begin{otherlanguage}{persian}

و اينكه نوشته مدّعى در ظلّ افتاده اگر چه از جمال اقدس اطهر جز عبوديّت صرفه
چيزى استماع نشده و لكن اين كلمه ملقى بسيار لغو بوده و خواهد بود از
اينقرار حضرت اعلى روح ما سواه فداه در ظلّ علماء عصر بوده چه كه در آن عصر
ظاهر و اگر از علماء سلب نمايند در ظلّ رسول اللّه بوده ما صغر شأنهم و قلّ
عقلهم و صغر درايتهم و كبر قولهم و لكن شهادت ميدهم كه قبل از ظهور در ظلّ
نبی‌ بوده و حين ظهور فوق كلّ ذى رسالة

\end{otherlanguage}

Va īnki nivishta muddaʿī dar ẓill uftāda, agar chi az Jamāl-i Aqdas-i
Aṭhar juz ʿubūdiyyat-i ṣirfa chīzī istimāʿ nashuda, va lākin īn
kalima-yi mulqā bisyār laghv būda va khwāhad būd. Az īn qarār Ḥaḍrat-i
Aʿlā, rūḥu mā siwāhu fidāh, dar ẓill-i ʿulamā-yi ʿaṣr būda, chi ki dar
ān ʿaṣr ẓāhir, va agar az ʿulamā salb namāyand dar ẓill-i Rasūl-i Allāh
būda. Mā ṣaghura shaʾnuhum wa-qalla ʿaqluhum wa-ṣaghura dirāyatuhum
wa-kabura qawluhum. Va lākin shahādat mīdaham ki qabl az ẓuhūr dar
ẓill-i nabī būda va ḥīn-i ẓuhūr fawq-i kull-i dhī risāla.

As for your statement that the claimant has fallen under the shadow:
although nothing but pure servitude has been heard from the most holy,
most pure Beauty, this utterance has been and will remain utterly idle.
By this reasoning, His Holiness the Exalted One---may the soul of all
else be sacrificed for Him---was under the shadow of the religious
scholars of His age, since He appeared in that age. If they deny that He
was under the scholars' shadow, then He was under the shadow of the
Messenger of God. How low their station, how little their reason, how
meager their understanding, and how great their claim! Yet I testify
that before His Revelation He was under the shadow of a Prophet, and at
the moment of His Revelation He was above every bearer of a divine
message.

\subsection{248}\label{section-743}

\begin{otherlanguage}{persian}

قوله عظم ظلّه پس چنين شخصى اگر خود و اصحابش در تقرير و تحرير مدّعى آن شوند
كه برترى از ربّ اعلى جسته و هست زهى تناقض و خلاف صريح با مدّعاى اوّل است
ادّعاى دو امر متناقض از شخصى بزرگ بلكه از هر ذى شعورى قبيح است هيچ عاقلى
ادّعاى دو امر متناقض صريح نمينمايد خصوص اگر ادّعاى آن نمايد كه ربّ اعلى از
يك كلمه من ايجاد شده پناه بر خدا كه يك فى‌الجمله تجلّى اشخاص بزرگ ببعضى
بجهت امتحان چه ادّعاها بروز مينمايد و جمعى از عوام هم باور مينمايند و
بلاتصوّر تصديق ميكنند انتهى

\end{otherlanguage}

Qawluhu, ʿaẓuma ẓilluhu: Pas chunīn shakhṣī agar khud va aṣḥābash dar
taqrīr va taḥrīr muddaʿī-yi ān shavand ki bartarī az Rabb-i Aʿlā justa
va hast, zihī tanāquḍ va khilāf-i ṣarīḥ bā muddaʿā-yi avval ast.
Iddiʿā-yi du amr-i mutanāqiḍ az shakhṣī buzurg, balki az har
dhī-shuʿūrī, qabīḥ ast. Hīch ʿāqilī iddiʿā-yi du amr-i mutanāqiḍ-i ṣarīḥ
namīnamāyad, khuṣūṣ agar iddiʿā-yi ān namāyad ki Rabb-i Aʿlā az yak
kalima-yi man ījād shuda. Panāh bar Khudā ki yak fī al-jumla tajallī-yi
ashkhāṣ-i buzurg bi-baʿḍī bi-jahat-i imtiḥān chi iddiʿāhā burūz
mīnamāyad va jamʿī az ʿavāmm ham bāvar mīnamāyand va bi-lā taṣavvur
taṣdīq mīkunand. Intahā.

His words---may his shadow grow: ``If such a person and his companions
then claim, in speech and writing, that he has attained superiority over
the Exalted Lord and is superior to Him, what a contradiction and
explicit conflict with the original claim! For a great person, or indeed
any person possessed of understanding, to make two contradictory claims
is reprehensible. No rational person makes two explicitly contradictory
claims, especially if he claims that `the Exalted Lord was created by a
single word of mine.' God save us! What claims a mere partial revelation
of great personages to certain people as a test brings forth! And a
number of ordinary people believe them and assent without reflection.''
End of quotation.

\subsection{249}\label{section-744}

\begin{otherlanguage}{persian}

قل للمشركين اذا لو تسمعون منّى فابكوا على انفسكم ثمّ نوحوا لقلّة عقولكم و
درايتكم ذكر امثال آن معرضين در كتاب قبل و بعد نازل قوله تعالى كمثل
الحمار يحمل اسفاراً بگو تناقض وجود خود شما است او را از ميان بردار تا جز
توحيد صرف و تفريد بحت مشاهده ننمائى اگر تو بخواهى اين ادّعا را بفهمى و
آنچه فرموده‌ ادراك نمائى كه شايد بمنبع معين قدس لايزال وارد شوى اينست
قوله عزّ اجلاله يا قوم انّى لجمال علىّ بينكم و نفس محمّد فيكم و كينونة الرّوح
بين السّموات و الأرض ان انتم تعرفون و اين همان مطلبی‌ است كه بعينه نقطه
بيان فرموده آن اراده نفس اين اراده بوده و خواهد بود و آن مشيّت بعينه همين
مشيّت بوده و خواهد بود و اين مشيّت و اراده بعينه همان مشيّت و اراده بوده و
خواهد بود ذكر تناقض من غير شعور نموده چه كه فى‌الحقيقة در حقّ جميع انبیاء
اين تناقض را ثابت نموده زيرا كه هر نبی‌ بعدى در ظلّ نبی‌ قبل بوده و با ثبوت
ظلّيّت ادّعاى برترى كه نموده نعوذ باللّه بقول مشركين تناقضست و از درجه
اعتبار ساقط حاشا ثمّ حاشا

\end{otherlanguage}

Qul li-l-mushrikīn idhan law tasmaʿūna minnī fa-bkū ʿalā anfusikum
thumma nūḥū li-qillati ʿuqūlikum wa-dirāyatikum. Dhikr-i amthāl-i ān
muʿriḍīn dar kitāb-i qabl va baʿd nāzil. Qawluhu taʿālā: Ka-mathali
al-ḥimāri yaḥmilu asfāran. Bigū tanāquḍ vujūd-i khud-i shumā ast. Ū rā
az miyān bardār tā juz tawḥīd-i ṣirf va tafrīd-i baḥt mushāhada
nanamāʾī. Agar tu bikhwāhī īn iddiʿā rā bifahmī va ānchi farmūda idrāk
namāʾī ki shāyad bi-manbaʿ-i maʿīn-i quds-i lā-yazāl vārid shavī, īnast
qawluhu ʿazza ijlāluhu: Yā qawmi innī la-Jamālu ʿAliyyin baynakum
wa-nafsu Muḥammadin fīkum wa-kaynūnatu al-Rūḥi bayna al-samāwāti
wa-al-arḍi in antum taʿrifūn. Va īn hamān maṭlabī ast ki bi-ʿaynihi
Nuqṭa-yi Bayān farmūda. Ān irāda nafs-i īn irāda būda va khwāhad būd va
ān mashiyyat bi-ʿaynihi hamīn mashiyyat būda va khwāhad būd va īn
mashiyyat va irāda bi-ʿaynihi hamān mashiyyat va irāda būda va khwāhad
būd. Dhikr-i tanāquḍ min ghayr-i shuʿūr namūda, chi ki fī al-ḥaqīqa dar
ḥaqq-i jamīʿ-i anbiyā īn tanāquḍ rā thābit namūda, zīrā ki har nabī-yi
baʿdī dar ẓill-i nabī-yi qabl būda va bā thubūt-i ẓilliyyat iddiʿā-yi
bartarī ki namūda, naʿūdhu bi-Allāh, bi-qawl-i mushrikīn tanāquḍ ast va
az daraja-yi iʿtibār sāqiṭ. Ḥāshā thumma ḥāshā.

Say to those who associate partners with God: ``If you would listen to
Me, then weep for yourselves and lament the paucity of your reason and
understanding.'' Descriptions of the likes of those who turn away have
been revealed in the former and latter Books. His words, exalted is He:
``Like a donkey carrying books.'' Say: Your own existence is the
contradiction. Remove it, so that you may behold nothing but pure unity
and absolute oneness. If you wish to understand this claim and
comprehend what He has said, so that you may perhaps reach the source of
the ever-flowing waters of everlasting holiness, these are His words,
exalted be His majesty: ``O people! I am indeed the Beauty of `Alí among
you, the very Self of Muhammad in your midst, and the Being of the
Spirit between the heavens and the earth, if you only knew.'' This is
precisely the same truth that the Point of the Bayán expressed. That
intention has been and will remain this very intention, and that will
has been and will remain this very will; and this will and intention
have been and will remain precisely that same will and intention. He has
spoken of contradiction without understanding, for in truth he has
established this contradiction with respect to all the Prophets. Every
subsequent Prophet was under the shadow of the preceding Prophet; thus,
with this relationship of being under the shadow established, His claim
to superiority is---God forbid!---according to the words of those who
associate partners with God, a contradiction and devoid of validity. Far
be it, far be it!

\subsection{250}\label{section-745}

\begin{otherlanguage}{persian}

قل يا ايّها النّائم على فراش الشّهوات و المسترضع عن ثدى الغفلة و الجهل و
الهوى تناقض را در ساحت اقدس مظاهر مقدّسه الهيّه راه نبوده و نخواهد بود اين
تناقض در نفس خلق است نه در امر حقّ بپذير قول اين عبد را و نفس خود را از
ثقل حمل اين اشارات و كلمات نجات ده تا خفيف شده بهواى قدس روحانى پرواز
نمائى و بمقرّ مقدّس از تحديدات بشريّه وارد شوى تا آن ساحت را مقدّس از ذكر
تناقض و اضداد و انقلاب و اختلاف و محاربه و مكالمه و مجادله و اشاره و
دلاله يابی‌ چه كه آن مقرّ اطهر از اسطقسّات اضداديّه خلق نشده و از عناصر
اربعه مختلفه بوجود نيامده از عنصر روح ظاهر شده هوائش عين نار و نارش نفس
هواء و مائش صرف تراب كلّ متّحدند من غير تغيير و اختلاف و رطبش عين يبوست و
يبوستش عين برودت و برودتش حقيقت حرارت و از اين حرارت حركت ما فى الوجود
خلق لو لا الحرّ لم تكن الحركة و الحرارة هى الفاعل و انّما الحركة هى الفعل
و لكن

\end{otherlanguage}

Qul yā ayyuhā al-nāʾimu ʿalā firāshi al-shahawāti wa-al-mustarḍiʿu ʿan
thadyi al-ghaflati wa-al-jahli wa-al-hawā. Tanāquḍ rā dar sāḥat-i
aqdas-i maẓāhir-i muqaddasa-yi ilāhiyya rāh nabūda va nakhwāhad būd. Īn
tanāquḍ dar nafs-i khalq ast na dar amr-i Ḥaqq. Bipadhīr qawl-i īn ʿabd
rā va nafs-i khud rā az thiql-i ḥaml-i īn ishārāt va kalimāt najāt dih
tā khafīf shuda bi-havā-yi quds-i rūḥānī parvāz namāʾī va bi-maqarr-i
muqaddas az taḥdīdāt-i bashariyya vārid shavī tā ān sāḥat rā muqaddas az
dhikr-i tanāquḍ va aḍdād va inqilāb va ikhtilāf va muḥāraba va mukālama
va mujādala va ishāra va dalāla yābī, chi ki ān maqarr-i aṭhar az
usṭuqussāt-i aḍdādiyya khalq nashuda va az ʿanāṣir-i arbaʿa-yi
mukhtalifa bi-vujūd nayāmada. Az ʿunṣur-i rūḥ ẓāhir shuda. Havāyash
ʿayn-i nār va nārash nafs-i havā va māʾash ṣirf-i turāb. Kull
muttaḥidand min ghayr-i taghyīr va ikhtilāf va raṭbash ʿayn-i yubūsat va
yubūsatash ʿayn-i burūdat va burūdatash ḥaqīqat-i ḥarārat va az īn
ḥarārat ḥarakat-i mā fī al-vujūd khalq. Law lā al-ḥarru lam takuni
al-ḥarakatu wa-al-ḥarāratu hiya al-fāʿilu wa-innamā al-ḥarakatu hiya
al-fiʿlu wa-lākin

Say: ``O you who sleep on the bed of desires and suckle at the breast of
heedlessness, ignorance, and passion!'' Contradiction has never had and
will never have access to the most holy court of the sacred
Manifestations of God. This contradiction lies within creation itself,
not in the Cause of God. Accept the words of this servant and free
yourself from the burden of carrying these allusions and words, so that,
unburdened, you may fly in the atmosphere of spiritual holiness and
enter the abode sanctified from human limitations. Then you will find
that court sanctified beyond any mention of contradiction, opposites,
transformation, difference, conflict, conversation, disputation,
allusion, or indication. For that most pure abode was not created from
opposing constituents, nor did it come into being from the four
differing elements. It has appeared from the element of spirit. Its air
is fire itself, its fire is air itself, and its water is pure earth. All
are united, without change or difference. Its moisture is dryness
itself, its dryness is cold itself, and its cold is the very reality of
heat. From this heat the motion of all that exists was created. Were it
not for heat, there would be no motion. Heat is the agent, while motion
is the action, but

\subsection{251}\label{section-746}

\begin{otherlanguage}{persian}

امثال آن نفوس كه از طين حرص و نار كبر و هواء نفس و ماء غفلت خلق شده‌اند
البتّه بآن مقرّ اقدس اطهر كه مقدّس از عرفان كلّ من فى الملك بوده و خواهد بود
نتوانند تقرّب جست لو يخطر التقرّب ببالهم ليحترقون فى الحين آنچه در مظاهر
قبل ذكر شده در اين ظهور ابدع امنع هو بهو و كلمه بكلمه و حرف بحرف ذكر
ميشود يرجع التّناقض و الاضداد الى اشكالهم و امثالهم و يبقى الاتّحاد و
الايتلاف لاحبّاء اللّه فى المبدأ و المآب اى منصف بقول خود ظلّيّت فرض مگير و
مايظهر من جمال الأبهی را تناقض مبين

\end{otherlanguage}

Amthāl-i ān nufūs ki az ṭīn-i ḥirṣ va nār-i kibr va havā-yi nafs va
māʾ-i ghaflat khalq shuda-and albatta bi-ān maqarr-i aqdas-i aṭhar ki
muqaddas az ʿirfān-i kull-i man fī al-mulk būda va khwāhad būd
natavānand taqarrub just. Law yakhṭuru al-taqarrubu bi-bālihim
la-yaḥtariqūna fī al-ḥīn. Ānchi dar maẓāhir-i qabl dhikr shuda dar īn
ẓuhūr-i abdaʿ-i amnaʿ hū bi-hū va kalima bi-kalima va ḥarf bi-ḥarf dhikr
mīshavad. Yarjiʿu al-tanāquḍu wa-al-aḍdādu ilā ashkālihim wa-amthālihim
wa-yabqā al-ittiḥādu wa-al-iʾtilāfu li-aḥibbāʾi Allāhi fī al-mabdaʾi
wa-al-maʾāb. Ay munṣif, bi-qawl-i khud ẓilliyyat farḍ magīr va mā
yaẓharu min Jamāli al-Abhā rā tanāquḍ mabīn.

The likes of those souls who have been created from the clay of greed,
the fire of pride, the air of self, and the water of heedlessness
certainly cannot draw near to that most holy, most pure abode, which has
been and will remain sanctified beyond the knowledge of all who dwell in
the realm of creation. Should the thought of drawing near enter their
minds, they would burn up at once. Whatever was said in the former
Manifestations is stated in this most wondrous, most exalted Revelation,
identically, word for word and letter for letter. Contradiction and
opposites revert to those of their kind and their like, while unity and
harmony remain for the loved ones of God, in the beginning and at the
return. O fair-minded one! Do not assume, on your own assertion, a
relationship of being under the shadow, and do not regard what appears
from the All-Glorious Beauty as a contradiction.

\subsection{252}\label{section-747}

\begin{otherlanguage}{persian}

و ديگر نوشته بودند خصوص اگر ادّعاى آن نمايد كه ربّ اعلى از يكحرف من ايجاد
شده از اين قسم كلمات جعليّه شركيّه كه آن نفس مشرك بشما القا نموده از انفس
خبيثه مشركه در جميع اعصار ظهور مظاهر الهيّه ظاهر شده تخصيص باين ظهور بديع
نداشته و ندارد چنانچه در ظهور نقطه بيان روح ما سواه فداه نسبت كذب و
افترا و شرك و كفر بآن شجره طيّبه الهيّه داده‌اند و كلّ استماع نموده و
ميدانند لذا اگر از امثال آن معرضين اينكلمات ظاهر شود عجب نيست چه كه
لازال شأن آن مردود و امثال او اين بوده و خواهد بود فواللّه هر ذى بصرى از
امثال اين كلمات ايقان و تصديقش بحقّ زياد شده و خواهد شد چه كه اين امور از
سنن امر بوده و خواهد بود چنانچه در كلّ كتب از قبل و بعد مسطور است و همّت
كلّ امّة برسولهم ليأخذوه و جادلوا به بالباطل ليدحضوا به الحقّ فاخذتهم فكيف
كان عقاب و در مقام ديگر و ما يأتيهم من رسول الّا كانوا به يستهزئون لذا هر
چه از اين اعتراضات ذكر شود صدق كلمات الهى كه دركتب او ثبت شده نزد موقنين
و متبصّرين ظاهر شده و ميشود

\end{otherlanguage}

Va dīgar nivishteh būdand khuṣūṣ agar iddiʿā-yi ān namāyad kih Rabb-i
Aʿlā az yik ḥarf-i man ījād shudih. Az īn qism kalimāt-i jaʿliyyih-yi
shirkiyyih kih ān nafs-i mushrik bi-shumā ilqā namūdih az anfus-i
khabīthih-yi mushrikih dar jamīʿ-i aʿṣār-i ẓuhūr-i maẓāhir-i ilāhiyyih
ẓāhir shudih, takhṣīṣ bi-īn ẓuhūr-i badīʿ nadāshtih va nadārad.
Chunānchih dar ẓuhūr-i Nuqṭih-yi Bayān, rūḥ-i mā sivāh fidāh, nisbat-i
kidhb va iftirā va shirk va kufr bi-ān shajari h-yi ṭayyibih-yi
ilāhiyyih dādih-and va kull istimāʿ namūdih va mīdānand. Lidhā agar az
amthāl-i ān muʿriḍīn īn kalimāt ẓāhir shavad ʿajab nīst, chih kih lāzāl
shaʾn-i ān mardūd va amthāl-i ū īn būdih va khvāhad būd. Fa-va'llāh har
dhī-baṣarī az amthāl-i īn kalimāt īqān va taṣdīqash bi-Ḥaqq ziyād shudih
va khvāhad shud, chih kih īn umūr az sunan-i Amr būdih va khvāhad būd,
chunānchih dar kull-i kutub az qabl va baʿd masṭūr ast: wa-hammat kullu
ummatin bi-rasūlihim li-yaʾkhudhūhu wa-jādalū bihi bi'l-bāṭili
li-yudḥiḍū bihi al-ḥaqqa fa-akhadhtuhum fa-kayfa kāna ʿiqāb. Va dar
maqām-i dīgar: wa-mā yaʾtīhim min rasūlin illā kānū bihi yastahziʾūn.
Lidhā har chih az īn iʿtirāḍāt dhikr shavad ṣidq-i kalimāt-i ilāhī kih
dar kutub-i ū thabt shudih nazd-i mūqinīn va mutabaṣṣirīn ẓāhir shudih
va mīshavad.

They had also written: ``Especially if he claims that the Most Exalted
Lord was created from a single letter of mine.'' Fabricated, idolatrous
words of this kind, which that idolatrous soul has instilled in you,
have issued from wicked, idolatrous souls in every age when the
Manifestations of God have appeared. They have never been peculiar to
this wondrous Revelation. Thus, at the appearance of the Point of the
Bayān---may the souls of all besides Him be sacrificed for Him---they
accused that goodly divine Tree of falsehood, fabrication, idolatry, and
unbelief, as everyone has heard and knows. It is therefore no wonder if
such words issue from the likes of those who turn away, for this has
always been, and will remain, the way of that rejected one and his like.
By God, words such as these have increased, and will increase, the
certainty of everyone with insight and their acknowledgment of the True
One, for such things have always been, and will remain, among the
established ways of the Cause, as recorded in all the Books, both
earlier and later: ``Every nation plotted against its messenger to seize
him, and they disputed with falsehood to refute the truth thereby. So I
seized them, and how was My punishment!'' And elsewhere: ``No messenger
comes to them without their mocking him.'' Thus, whatever objections of
this kind are voiced, the truth of the divine words recorded in His
Books has become, and continues to become, apparent to those who are
certain and discerning.

\subsection{253}\label{section-748}

\begin{otherlanguage}{persian}

بگو اى بی‌ بصر آنچه از بيانات بديعه منيعه نازل شده حقّ نازل فرموده آيا
نشنيده كه از سماء احديّه نازل قوله جلّ كبريائه و عزّ اعزازه يا قوم انّى قد
كنت راقداً فى البيت و صامتاً عن الذّكر هبّت علىّ نسمات اللّه و احيتنى بالحقّ و
انطقنى بثنآء نفسه و جعلنى هدى و ذكرى للعالمين و كلّما اريد ان اصمت روح
القدس ينطقنى بالحقّ و روح الأعظم يهتزّنى و روح البقآء يحرك قلم البهآء ان
انتم من العارفين

\end{otherlanguage}

Bigū ay bī-baṣar, ānchih az bayānāt-i badīʿih-yi manīʿih nāzil shudih
Ḥaqq nāzil farmūdih. Āyā nashnīdih kih az samā-yi aḥadiyyih nāzil,
qawluhu jalla kibriyāʾuhu wa-ʿazza iʿzāzuhu: yā qawmi innī qad kuntu
rāqidan fī al-bayti wa-ṣāmitan ʿani al-dhikri habbat ʿalayya nasamātu
Allāhi wa-aḥyatnī bi'l-ḥaqqi wa-anṭaqanī bi-thanāʾi nafsihi wa-jaʿalanī
hudan wa-dhikrā li'l-ʿālamīn. Wa-kullamā urīdu an aṣmuta Rūḥu al-Qudusi
yunṭiqunī bi'l-ḥaqqi wa-Rūḥu al-Aʿẓami yahuzzunī wa-Rūḥu al-Baqāʾi
yuḥarriku Qalama al-Bahāʾi in antum mina al-ʿārifīn.

Say: O you who lack insight! Whatever wondrous and exalted utterances
have been sent down, the True One has sent them down. Have you not heard
what has descended from the heaven of oneness? His words---glorious is
His majesty and mighty His exaltation---are: ``O people! I was asleep in
the house and silent in remembrance when the breezes of God blew over Me
and quickened Me through the truth. He caused Me to speak in praise of
Himself and made Me a guidance and a reminder to the worlds. Whenever I
wish to remain silent, the Holy Spirit causes Me to speak the truth, the
Most Great Spirit stirs Me, and the Spirit of Eternity moves the Pen of
Bahāʾ, if you are among those who understand.''

\subsection{254}\label{section-749}

\begin{otherlanguage}{persian}

يا قوم خافوا عن اللّه ثمّ استحيوا عن جماله و لا تتكلّموا بما يلعنكم به كلّ
الذّرّات و عن ورائها لسان اللّه الملك الصّادق الأمين و ليس هذا من عندى بل من
عنده لو انتم من الشّاعرين فواللّه لو كان الأمر بيدى لسترت نفسى عن ابصاركم
و ما القيت كبدى تحت مخاليب ذياب الأرض و كان اللّه على ما اقول شهيد

\end{otherlanguage}

Yā qawmi khāfū ʿani Allāhi thumma istaḥyū ʿan jamālihi wa-lā tatakallamū
bi-mā yalʿanukum bihi kullu al-dharrāti wa-ʿan warāʾihā lisānu Allāhi
al-Maliki al-Ṣādiqi al-Amīn. Wa-laysa hādhā min ʿindī bal min ʿindihi
law antum mina al-shāʿirīn. Fa-wa'llāhi law kāna al-amru bi-yadī
la-satartu nafsī ʿan abṣārikum wa-mā alqaytu kabidī taḥta makhālībi
dhiyābi al-arḍi wa-kāna Allāhu ʿalā mā aqūlu shahīd.

O people! Fear God, and be ashamed before His Beauty. Do not speak words
for which every atom curses you, and beyond them the Tongue of God, the
King, the Truthful, the Trustworthy. This is not from Me but from Him,
if only you perceived. By God, if the matter were in My hands, I would
have concealed Myself from your eyes and would not have cast My heart
beneath the claws of the wolves of the earth. God is witness to what I
say.

\subsection{255}\label{section-750}

\begin{otherlanguage}{arabic}

ایّاکم يا قوم ان تمنعكم النّفس و الهوى عن الصّعود الى مقرّ الأقصى فانظروا
بطرف الانصاف الى حجج النّبيّين و المرسلين تاللّه انّ هذا لغلام الرّحمن قد
اظهره اللّه بين ملأ الاكوان و اصطفاه من بين بريّته و اظهره بطراز نفسه بين
العالمين و يقول فى كلّ حين مخاطباً الى ملأ اليهود ان يا ملأ العنود تاللّه
قد جائكم الموعود و انّ هذا لهو الرّوح ان تريدوا ان تصلبوه فافعلوا ما شئتم
و لا تكوننّ من الصّابرين ثمّ يخاطب ملأ الانجيل و يقول ان تريدوا ان تجادلوا
محمّداً رسول اللّه انّ هذا لمحمّد بينكم فافعلوا ما اردتم لانّه انفق روحه فى
سبيل اللّه المهيمن العزيز القدير

\end{otherlanguage}

Iyyākum yā qawmi an tamnaʿakumu al-nafsu wa'l-hawā ʿani al-ṣuʿūdi ilā
maqarri al-aqṣā fa'nẓurū bi-ṭarfi al-inṣāfi ilā ḥujaji al-nabiyyīna
wa'l-mursalīn. Ta'llāhi inna hādhā la-ghulāmu al-Raḥmāni qad aẓharahu
Allāhu bayna malaʾi al-akwāni wa'ṣṭafāhu min bayni bariyyatihi
wa-aẓharahu bi-ṭirāzi nafsihi bayna al-ʿālamīn. Wa-yaqūlu fī kulli ḥīnin
mukhāṭiban ilā malaʾi al-Yahūdi an yā malaʾa al-ʿunūdi ta'llāhi qad
jāʾakumu al-Mawʿūdu wa-inna hādhā la-huwa al-Rūḥu in turīdū an taṣlibūhu
fa'fʿalū mā shiʾtum wa-lā takūnanna mina al-ṣābirīn. Thumma yukhāṭibu
malaʾa al-Injīli wa-yaqūlu in turīdū an tujādilū Muḥammadan Rasūla
Allāhi inna hādhā la-Muḥammadun baynakum fa'fʿalū mā aradtum li-annahu
anfaqa rūḥahu fī sabīli Allāhi al-Muhaymini al-ʿAzīzi al-Qadīr.

Beware, O people, lest self and desire prevent you from ascending to the
highest abode. Look with the eye of fairness upon the proofs of the
Prophets and Messengers. By God, this is indeed the Youth of the
All-Merciful, whom God has manifested among the concourse of created
beings, chosen from among His creatures, and revealed among the worlds
adorned with His own Self. At every moment He addresses the company of
the Jews, saying: ``O company of obstinacy! By God, the Promised One has
come to you, and this is indeed the Spirit. If you wish to crucify Him,
do whatever you please, and do not delay.'' Then He addresses the people
of the Gospel and says: ``If you wish to dispute with Muḥammad, the
Messenger of God, this is indeed Muḥammad among you. Do whatever you
desire, for He has offered up His soul in the path of God, the
All-Protecting, the Almighty, the All-Powerful.''

\subsection{256}\label{section-751}

\begin{otherlanguage}{arabic}

ثمّ يخاطب اهل الفرقان و يقول يا ملأ الطّغيان ان تشائوا ان تعلّقوا مظهر نفسى
الّذى سمّى بعلىّ فى الهواء تاللّه انّ هذا لعلىّ قد حضر تلقآء وجوهكم يا ملأ
الذّياب ثمّ يا شرذمة الخنازير فافعلوا به ما شئتم و لا تصبروا اقلّ من حين
لانّه ما اتّخذ لنفسه معيناً و لا ناصراً و علّمه اللّه بالحقّ ما ورد عليه و يرد
يا جنود الشّياطين

\end{otherlanguage}

Thumma yukhāṭibu ahla al-Furqāni wa-yaqūlu yā malaʾa al-ṭughyāni in
tashāʾū an tuʿalliqū maẓhara nafsī alladhī summiya bi-ʿAliyyin fī
al-hawāʾi ta'llāhi inna hādhā la-ʿAliyyun qad ḥaḍara tilqāʾa wujūhikum
yā malaʾa al-dhiyābi thumma yā shirdhimata al-khanāzīri fa'fʿalū bihi mā
shiʾtum wa-lā taṣbirū aqalla min ḥīnin li-annahu mā ittakhadha
li-nafsihi muʿīnan wa-lā nāṣiran wa-ʿallamahu Allāhu bi'l-ḥaqqi mā
warada ʿalayhi wa-yaridu yā junūda al-shayāṭīn.

Then He addresses the people of the Qurʾān and says: ``O company of
tyranny! If you wish to suspend in the air the Manifestation of My Self
who was named ʿAlī, by God, this is indeed ʿAlī standing before your
faces, O company of wolves, O band of swine! Do to Him whatever you
please, and do not delay for even a moment, for He has taken neither
helper nor supporter for Himself. God has truly taught Him what has
befallen Him and what will befall Him, O armies of devils!''

\subsection{257}\label{section-752}

\begin{otherlanguage}{arabic}

ثمّ يخاطب ملأ البيان و يقول يا ملأ الطّغيان و الكفران قد جائكم الموعود
الّذى وعدتم به فى كلّ الألواح فو‌جماله انّه قد ظهر بالحقّ و ما حفظ نفسه فى
سبيل اللّه الملك المتعالى العزيز الحميد و كان بين يدى الأعداء فى كلّ
اللّيالى و الأیّام و نصر امر ربّه بنفسه المتعالى العزيز الجميل و ورد عليه
ما لا يذكر بالبيان و ما اطّلع به احد الّا اللّه الّذى بعثه بالحقّ و ارسله على
العالمين انتم ان تريدوا ان تقتلوه كما قتلتموه فافعلوا ما شئتم لأنّه ينتظر
ما وعد به فى كلّ الألواح ان انتم من العارفين

\end{otherlanguage}

Thumma yukhāṭibu malaʾa al-Bayāni wa-yaqūlu yā malaʾa al-ṭughyāni
wa'l-kufrāni qad jāʾakumu al-Mawʿūdu alladhī wuʿidtum bihi fī kulli
al-alwāḥi fa-wa-jamālihi innahu qad ẓahara bi'l-ḥaqqi wa-mā ḥafiẓa
nafsahu fī sabīli Allāhi al-Maliki al-Mutaʿālī al-ʿAzīzi al-Ḥamīd.
Wa-kāna bayna yaday al-aʿdāʾi fī kulli al-layālī wa'l-ayyāmi wa-naṣara
amra rabbihi bi-nafsihi al-Mutaʿālī al-ʿAzīzi al-Jamīl. Wa-warada
ʿalayhi mā lā yudhkaru bi'l-bayāni wa-mā iṭṭalaʿa bihi aḥadun illā
Allāhu alladhī baʿathahu bi'l-ḥaqqi wa-arsalahu ʿalā al-ʿālamīn. Antum
in turīdū an taqtulūhu kamā qataltumūhu fa'fʿalū mā shiʾtum li-annahu
yantaẓiru mā wuʿida bihi fī kulli al-alwāḥi in antum mina al-ʿārifīn.

Then He addresses the people of the Bayān and says: ``O company of
tyranny and unbelief! The Promised One, who was promised to you in all
the Tablets, has come to you. By His Beauty, He has appeared in truth
and has not safeguarded Himself in the path of God, the King, the Most
Exalted, the Almighty, the All-Praised. He has been in the presence of
His enemies every night and day and has aided the Cause of His Lord
through His own Self, the Most Exalted, the Almighty, the All-Beautiful.
What has befallen Him cannot be expressed in words, and no one knows it
except God, who raised Him up in truth and sent Him to the worlds. If
you wish to kill Him as you killed Him before, do whatever you please,
for He awaits what He was promised in all the Tablets, if you are among
those who understand.''

\subsection{258}\label{section-753}

\begin{otherlanguage}{persian}

اذاً يخاطب علىّ عن يمين العرش ملأ البيان و يقول يا عباد اما بشّرناكم بهذا
الظّهور فى كلّ الألواح اما اخذت عنكم عهد نفسه فى ذرّ البيان اذا فاستحيوا عن
جمال الرّحمن و لا تفعلوا ما يرتفع به ضجيج اهل ملأ الأعلی خافوا عن اللّه و
لا تكوننّ من الظّالمين

\end{otherlanguage}

Idhan yukhāṭibu ʿAliyyun ʿan yamīni al-ʿarshi malaʾa al-Bayāni wa-yaqūlu
yā ʿibādi a-mā bashsharnākum bi-hādhā al-ẓuhūri fī kulli al-alwāḥi a-mā
akhadhtu ʿankum ʿahda nafsihi fī dharri al-Bayāni idhan fa'staḥyū ʿan
jamāli al-Raḥmāni wa-lā tafʿalū mā yartafiʿu bihi ḍajīju ahli malaʾi
al-aʿlā khāfū ʿani Allāhi wa-lā takūnanna mina al-ẓālimīn.

Then ʿAlī, from the right hand of the Throne, addresses the people of
the Bayān and says: ``O servants! Did We not give you the glad tidings
of this Revelation in all the Tablets? Did I not take from you the
covenant of His very Self in the primordial realm of the Bayān? Be
ashamed, then, before the Beauty of the All-Merciful, and do not act in
a way that causes the lamentation of the people of the Concourse on high
to rise. Fear God, and do not be among the wrongdoers.''

\subsection{259}\label{section-754}

\begin{otherlanguage}{arabic}

تاللّه انّى قد فديت نفسى شوقاً للقائه و طلباً لوصاله و اكون معه فى كلّ
الأحيان و ابكى على ما ورد عليه منكم يا معشر الظّالمين فآه آه بعد وصاياى
قد فعلتم ما لا فعل احد باحد اذاً بكى و يبكى مظاهر النّبيّين و المرسلين
فواللّه يا قوم انّى جعلت كلّ الأمور فى قبضته و ما تحركت فى ایّامى الّا بذكره
و ثنائه و ما تفوّهت الّا باسمه و ما امرت العباد الّا بالدّخول فى ظلّه العزيز
المنيع فيا ليت انتم فعلتم كما فعلت اليهود بالرّوح و عملتم بما عمل ملأ
الانجيل بمحمّد رسول اللّه و ارتكبتم ما ارتكبت به ملأ الفرقان حين ظهورى
بالحقّ بسلطان مبين

\end{otherlanguage}

Ta'llāhi innī qad fadaytu nafsī shawqan li-liqāʾihi wa-ṭalaban
li-wiṣālihi wa-akūnu maʿahu fī kulli al-aḥyāni wa-abkī ʿalā mā warada
ʿalayhi minkum yā maʿshara al-ẓālimīn. Fa-āhi āhi baʿda waṣāyāya qad
faʿaltum mā lā faʿala aḥadun bi-aḥadin idhan bakā wa-yabkī maẓāhiru
al-nabiyyīna wa'l-mursalīn. Fa-wa'llāhi yā qawmi innī jaʿaltu kulla
al-umūri fī qabḍatihi wa-mā taḥarraktu fī ayyāmī illā bi-dhikrihi
wa-thanāʾihi wa-mā tafawwahtu illā bi'smihi wa-mā amartu al-ʿibāda illā
bi'l-dukhūli fī ẓillihi al-ʿazīzi al-manīʿ. Fa-yā layta antum faʿaltum
kamā faʿalati al-Yahūdu bi'l-Rūḥi wa-ʿamiltum bi-mā ʿamila malaʾu
al-Injīli bi-Muḥammadin Rasūli Allāhi wa'rtakabtum mā irtakabat bihi
malaʾu al-Furqāni ḥīna ẓuhūrī bi'l-ḥaqqi bi-sulṭānin mubīn.

By God, I sacrificed Myself out of longing to meet Him and a desire for
union with Him. I am with Him at all times and weep over what He has
suffered at your hands, O company of wrongdoers! Alas, alas! Despite My
counsels, you have done what no one has done to another. At this the
Manifestations of the Prophets and Messengers have wept and continue to
weep. By God, O people, I placed all matters in His grasp. Throughout My
days I moved only to remember and praise Him, spoke only His name, and
commanded the servants only to enter His mighty, impregnable shelter. If
only you had acted as the Jews acted toward the Spirit, done what the
people of the Gospel did to Muḥammad, the Messenger of God, and
committed what the people of the Qurʾān committed when I appeared in
truth with manifest authority!

\subsection{260}\label{section-755}

\begin{otherlanguage}{persian}

لا فو‌نفسى ما قنعتم بما فعلوا هؤلآء و تفعلون ما يستعيذ منه حقايق كلّ الملل
لو انتم من الشّاعرين اتقتلون الّذى سفك دمى فى سبيله تاللّه بعد سفكه على
الأرض نقش باسمه المتعالى العزيز البديع يا قوم فاستحيوا عن اللّه و
لاتدحضوا الحقّ بما عندكم و لاترتكبوا ما يجعلكم محروماً عن نفحات اللّه
العزيز الحكيم اسمعوا قولى ثمّ ارجعوا اليه ثمّ توبوا و استغفروا لعلّ يغفركم
ببدايع رحمته و يطهّركم عن دنس العصيان و عن رجس انفسكم يا ملأ الغافلين

\end{otherlanguage}

Lā fa-wa-nafsī mā qaniʿtum bi-mā faʿala hāʾulāʾi wa-tafʿalūna mā
yastaʿīdhu minhu ḥaqāʾiqu kulli al-milali law antum mina al-shāʿirīn.
A-taqtulūna alladhī sufika damī fī sabīlihi ta'llāhi baʿda safkihi ʿalā
al-arḍi nuqisha bi'smihi al-Mutaʿālī al-ʿAzīzi al-Badīʿ. Yā qawmi
fa'staḥyū ʿani Allāhi wa-lā tudḥiḍū al-ḥaqqa bi-mā ʿindakum wa-lā
tartakibū mā yajʿalukum maḥrūman ʿan nafaḥāti Allāhi al-ʿAzīzi al-Ḥakīm.
Ismaʿū qawlī thumma irjiʿū ilayhi thumma tūbū wa'staghfirū laʿalla
yaghfirukum bi-badāʾiʿi raḥmatihi wa-yuṭahhirukum ʿan danasi al-ʿiṣyāni
wa-ʿan rijsi anfusikum yā malaʾa al-ghāfilīn.

No, by My own Self! You have not been content with what those people
did, and you do what the very realities of all religions seek refuge
from, if only you perceived. Will you kill the One in whose path My
blood was shed? By God, after it was shed upon the earth, it was
inscribed with His name, the Most Exalted, the Almighty, the Wondrous. O
people! Be ashamed before God. Do not refute the truth with what you
possess, and do not commit what will deprive you of the breezes of God,
the Almighty, the All-Wise. Listen to My words, then return to Him;
repent and seek forgiveness, so that He may forgive you through the
wonders of His mercy and cleanse you of the stain of disobedience and
the defilement of your own selves, O company of the heedless!

\subsection{261}\label{section-756}

\begin{otherlanguage}{arabic}

فواللّه يا قوم انّى نزلت ما يرد عليه فى الألواح فسوف يظهر ما رقم من قلم
الأعلی على الواح عزّ حفيظ اذاً ينوحنّ قبايل مداين الأعلی و انتم تفرحون فى
انفسكم و تلعبون بآيات اللّه بعد الّذى بها بشّرناكم فى صحايف عزّ منيع يا قوم
انّه روحى و انا روحه و لم يزل كان مشيّتى مشيّته لاتفرقوا بينى و بينه و كلّما
يظهر من عنده لكان محبوب قلبی‌ و فؤادى لو انتم من الموحّدين

\end{otherlanguage}

Fa-wa'llāhi yā qawmi innī nazzaltu mā yaridu ʿalayhi fī al-alwāḥi
fa-sawfa yaẓharu mā ruqima min Qalami al-Aʿlā ʿalā alwāḥi ʿizzin ḥafīẓ.
Idhan yanūḥanna qabāʾilu madāʾini al-aʿlā wa-antum tafraḥūna fī
anfusikum wa-talʿabūna bi-āyāti Allāhi baʿda alladhī bihā bashsharnākum
fī ṣaḥāʾifi ʿizzin manīʿ. Yā qawmi innahu rūḥī wa-anā rūḥuhu wa-lam
yazal kāna mashīʾatī mashīʾatahu lā tufarriqū baynī wa-baynihi
wa-kullamā yaẓharu min ʿindihi la-kāna maḥbūba qalbī wa-fuʾādī law antum
mina al-muwaḥḥidīn.

By God, O people, I revealed in the Tablets what would befall Him, and
what was inscribed by the Most Exalted Pen upon the Tablets of securely
guarded glory will soon appear. Then the tribes of the cities on high
will lament, while you rejoice within yourselves and trifle with the
verses of God, after We gave you glad tidings through them in the pages
of impregnable glory. O people! He is My spirit, and I am His spirit,
and My will has always been His will. Do not distinguish between Me and
Him. Whatever proceeds from Him is beloved of My heart and inmost being,
if you are among those who affirm the divine unity.

\subsection{262}\label{section-757}

\begin{otherlanguage}{persian}

اين بدايع كلمات الهى كه از سحاب فيض رحمت رحمانى نازلشده اين عبد ذكر
نموده كه شايد عباد غافلين بشعور آيند و قياس بنفس ننمايند و آنقدر مشعر
شوند كه معرضين از اين آيات بدع ربّانى معارض و محارب با حقّ بوده و خواهند
بود چه كه منزل آيات حقّ است حال اگر با پروردگار عالميان محاربند خود
ميدانند مفرّى از براى نفسى اليوم نه و مقرّى مشهود نه مگر آنكه يا بايد
اقرار نمايند بآنچه ظاهر شده يا آنكه بايد بجميع صحف و كتب و زبر كافر شوند
اينست قول حقّ و نيست بعد از حقّ مگر گمراهى آشكار و هويدا

\end{otherlanguage}

Īn badāyiʿ-i kalimāt-i ilāhī kih az saḥāb-i fayḍ-i raḥmat-i raḥmānī
nāzil shudih īn ʿabd dhikr namūdih kih shāyad ʿibād-i ghāfilīn bi-shuʿūr
āyand va qiyās bi-nafs nanamāyand va ān-qadr mushʿir shavand kih
muʿriḍīn az īn āyāt-i budaʿ-i rabbānī muʿāriḍ va muḥārib bā Ḥaqq būdih
va khvāhand būd, chih kih munzil-i āyāt Ḥaqq ast. Ḥāl agar bā
Parvardigār-i ʿālamiyān muḥāriband khud mīdānand. Mafarrī az barā-yi
nafsī al-yawm nah va maqarrī mashhūd nah magar ānkih yā bāyad iqrār
namāyand bi-ānchih ẓāhir shudih yā ānkih bāyad bi-jamīʿ-i ṣuḥuf va kutub
va zubur kāfir shavand. Īn ast qawl-i Ḥaqq va nīst baʿd az Ḥaqq magar
gumrāhī-yi āshkār va huvaydā.

This servant has cited these wondrous divine words, which have descended
from the cloud of the outpouring of the All-Merciful's mercy, so that
the heedless servants may perhaps become aware and not judge by the
measure of their own selves. May they at least realize that those who
turn away from these wondrous divine verses have opposed, and will
continue to oppose and wage war against, the True One, for it is the
True One who sends down the verses. Now, if they are at war with the
Lord of the worlds, that is their own affair. Today there is no escape
for any soul, nor any evident refuge: either they must acknowledge what
has appeared, or they must disbelieve in all the scriptures, books, and
sacred writings. This is the word of truth, and beyond the truth there
is nothing but plain and manifest error.

\subsection{263}\label{section-758}

\begin{otherlanguage}{persian}

مثل معرضين مثل نفوسى است كه در حين ظهور نقطه اولى بعضى اعراضا عن اللّه و
اعتراضا عليه مذكور مينمودند اين سيّد ادّعاى ولايت نموده و بعضى ادّعاى نبوّت
و بعضى ادّعاى ربوبيّت و هر كدام از مغلّين بكلمه ناطق و آنقدر ادراك ننمودند
كه بايد معين نمود كه منزل كيست و منزل هم معلوم نميشود الّا از نفس آيه و
قد خسر الّذينهم كفروا و اشركوا باللّه و بغوا على اللّه المهيمن القيّوم

\end{otherlanguage}

Mathal-i muʿriḍīn mathal-i nufūsī ast kih dar ḥīn-i ẓuhūr-i Nuqṭih-yi
Ūlā baʿḍī iʿrāḍan ʿani Allāhi wa-iʿtirāḍan ʿalayhi madhkūr mīnamūdand īn
siyyid iddiʿā-yi vilāyat namūdih va baʿḍī iddiʿā-yi nubuwwat va baʿḍī
iddiʿā-yi rubūbiyyat va har kudām az mughillīn bi-kalimih nāṭiq va
ān-qadr idrāk nanamūdand kih bāyad muʿayyan namūd kih munzil kīst va
munzil ham maʿlūm namīshavad illā az nafs-i āyih. Wa-qad khasira
alladhīnahum kafarū wa-ashrakū bi'llāhi wa-baghaw ʿalā Allāhi
al-Muhaymini al-Qayyūm.

Those who turn away are like the souls who, at the appearance of the
Primal Point, turned away from God and objected to Him. Some said,
``This Siyyid has claimed spiritual authority''; others, that He had
claimed prophethood; and others, that He had claimed lordship. Each of
those malicious people voiced a different word, yet they did not
understand even this much: that one must determine who sends down the
verses, and that the One who sends them down can be known only through
the verse itself. Lost indeed are those who have disbelieved, associated
partners with God, and rebelled against God, the All-Protecting, the
Self-Subsisting.

\subsection{264}\label{section-759}

\begin{otherlanguage}{persian}

و ديگر ذكر تجلّى نموده بودى و اللّه آن نفس ملقى و مرشدش نه تجلّى را ادراك
نموده و نه مجلّى را چه كه اگر ادراك نموده بودند بحقّ اعتراض نمينمودند
سزاوار بود كه اين كلمات را بنفس معرض باللّه راجع كنى چه كه از قميص يك
اسمى از اسماء بر حقّ استكبار نموده و كافر شده بنفسى كه باو ادّعاى ايمان
مينمايد آيات او را ميخواند و بر كينونتش سيف ميزند يفعل و لايشعر يقول و
لايعرف هو و امثاله عند اللّه فى حدّ سواء الّا انّه اشرّ و اظلم اعاذنا اللّه و
احبّائه من شرّه و شرّهم و انّه لهو الحافظ الحاكم العليم الخبير

\end{otherlanguage}

Va dīgar dhikr-i tajallī namūdih būdī. Va'llāh ān nafs-i mulqī va
murshidash nah tajallī rā idrāk namūdih va nah mujallī rā, chih kih agar
idrāk namūdih būdand bi-Ḥaqq iʿtirāḍ namīnamūdand. Sazāvār būd kih īn
kalimāt rā bi-nafs-i muʿriḍ bi'llāh rājiʿ kunī, chih kih az qamīṣ-i yik
ismī az asmā bar Ḥaqq istikbār namūdih va kāfir shudih bi-nafsī kih bi-ū
iddiʿā-yi īmān mīnamāyad. Āyāt-i ū rā mīkhvānad va bar kaynūnatash sayf
mīzanad. Yafʿalu wa-lā yashʿuru yaqūlu wa-lā yaʿrifu huwa wa-amthāluhu
ʿinda Allāhi fī ḥaddin sawāʾin illā annahu asharru wa-aẓlamu aʿādhanā
Allāhu wa-aḥibbāʾahu min sharrihi wa-sharrihim wa-innahu la-huwa
al-Ḥāfiẓu al-Ḥākimu al-ʿAlīmu al-Khabīr.

You had also mentioned revelation through divine radiance. By God,
neither that soul who instilled these ideas nor his guide has understood
either the radiance or the One who sheds it, for had they understood,
they would not have objected to the True One. It would have been fitting
for you to direct these words at the one who turns away from God, for,
clothed in a single one of the names, he has exalted himself against the
True One and disbelieved in the very One in whom he professes faith. He
reads His verses and strikes with his sword at His being. He acts
without perceiving and speaks without knowing. He and his like are on an
equal footing before God, except that he is more wicked and more unjust.
May God protect us and His loved ones from his evil and theirs. He is
indeed the Preserver, the Ruler, the All-Knowing, the All-Informed.

\subsection{265}\label{section-760}

\begin{otherlanguage}{persian}

و ديگر ذكر عوام نموده بودى و اينكه بلا تصوّر تصديق ميكنند حال تو انصاف ده
اين عبد حال شش سنه ميشود كه ليلا و نهارا در خدمت حاضر و با هر يك معاشر
بوده و هستم و آنچه بچشم خود ديده ام ميگويم و شما محض قول يك مفسد خبيث
ملحدى كه آسمان سايه نينداخته بر مفسدى مثل او حال اين عباد بلا تصوّر تصديق
حقّ نموده‌اند يا تو فو‌اللّه حمير بر عقل و درايت آن نفس مشرك كه اينكلمات را
بتو القاء نموده ميخندد چه جاى انسان مع‌ذلك اين عباد را عوام دانسته و ذكر
نموده كه بلا تصوّر تصديق حقّ نموده‌ايم

\end{otherlanguage}

Va dīgar dhikr-i ʿavāmm namūdih būdī va īnkih bilā taṣavvur taṣdīq
mīkunand. Ḥāl tu inṣāf dih: īn ʿabd ḥāl shish sanih mīshavad kih laylan
va nahāran dar khidmat ḥāḍir va bā har yak muʿāshir būdih va hastam, va
ānchih bi-chashm-i khud dīdih-am mīgūyam, va shomā maḥḍ-i qawl-i yak
mufsid-i khabīth-i mulḥidī kih āsmān sāyih nayandākhtih bar mufsidī
mithl-i ū. Ḥāl īn ʿibād bilā taṣavvur taṣdīq-i Ḥaqq namūdih-and yā tu?
Fa-va-llāh ḥamīr bar ʿaql va dirāyat-i ān nafs-i mushrik kih īn kalimāt
rā bi-tu ilqā namūdih mīkhandad, chih jāy-i insān. Maʿa-dhālik īn ʿibād
rā ʿavāmm dānistih va dhikr namūdih kih bilā taṣavvur taṣdīq-i Ḥaqq
namūdih-īm.

You had also spoken of the common people, saying that they assent
without forming a conception. Now judge fairly: for six years this
servant has been present, day and night, in His service, and has
associated, and continues to associate, with each of them. I speak of
what I have seen with my own eyes, while you rely solely on the word of
a wicked, impious troublemaker---the sky has never cast its shadow over
a troublemaker like him. Now, is it these servants who have affirmed the
True One without forming a conception, or you? By God, even donkeys
laugh at the intelligence and discernment of that idolatrous soul who
has instilled these words in you, let alone human beings. Nevertheless,
you have regarded these servants as common people and stated that we
have affirmed the True One without forming a conception.

\subsection{266}\label{section-761}

\begin{otherlanguage}{persian}

يقين تو خود را عالم ميدانى اولّا اينكه اين مطلب از كجا بر شما معلوم شد چه
كه خدمت جمال ابهى نرسيده‌ايد هيچ منصفى چنين ارتكابی‌ ننموده كه نديده و
نرسيده بمجرّد حرف اخبث ناس و افسدهم بردارد من غير شعور آنچه از قلم جارى
ميشود نوشته و بشهرها بفرستد كاش بقول طلبه‌هاى مدرسه هم عمل مينموديد چه كه
نزد آن طايفه تصديق منوط بسه تصوّر بوده اوّل موضوع و ثانى محمول و ثالث نسبت
حكميّه و بعد تصديق ثابت ميشود و جارى ميگردد آن موضوعى كه تو فهميده چه
بوده و آن محمول كدام است و فواللّه ارتكبت ما لا ارتكبه احد من الجاهلين

\end{otherlanguage}

Yaqīn tu khud rā ʿālim mīdānī. Avvalan īnkih īn maṭlab az kujā bar shomā
maʿlūm shud, chih kih khidmat-i Jamāl-i Abhā narisīdih-īd. Hīch munṣifī
chunīn irtikābī nanamūdih kih nadīdih va narisīdih bi-mujarrad-i ḥarf-i
akhbath-i nās va afsadihim bardārad min ghayr-i shuʿūr ānchih az qalam
jārī mīshavad nivishtih va bi-shahrhā bifiristad. Kāsh bi-qawl-i
ṭalabih-hā-yi madrisih ham ʿamal mīnamūdīd, chih kih nazd-i ān ṭāʾifih
taṣdīq manūṭ bi-sih taṣavvur būdih: avval mawḍūʿ va thānī maḥmūl va
thālith nisbat-i ḥukmiyyih, va baʿd taṣdīq thābit mīshavad va jārī mīg
ar dad. Ān mawḍūʿī kih tu fahmīdih chih būdih va ān maḥmūl kudām ast? Va
fa-wa-llāh irtakabta mā lā irtakabahu aḥadun min al-jāhilīn.

Surely you consider yourself learned. First, how did you come to know
this, since you have not attained the presence of the All-Glorious
Beauty? No fair-minded person has ever committed such an act: without
seeing Him or attaining His presence, merely on the word of the most
wicked and corrupt of people, heedlessly taking up the pen, writing
whatever flows from it, and sending it to the cities. If only you had
acted even according to what the students in the seminaries say! For
among them assent depends on three conceptions: first, the subject;
second, the predicate; and third, the relation of judgment. Only then is
assent established and brought into effect. What was the subject you
understood, and what was the predicate? By God, you have committed what
none of the ignorant has committed.

\subsection{267}\label{section-762}

\begin{otherlanguage}{persian}

بلى اين عباد تصديق مظاهر الهيّه از قبل نموده و از بعد هم بفضل اللّه و منّه
خواهند نمود و انشآء اللّه از عنايت و مكرمت او بهيچ حجابی‌ محتجب نشويم و
بهيچ منعى ممنوع نگرديم فواللّه آنچه از آن ملحد خبيث شنيده كذب صرف و
افتراء بحت بوده و نميدانى كه بنَفَس كه مؤانس شده تاللّه الحقّ يفرّ الجحيم من
نفسه و يستعيذ الحسبان باللّه من وجهه و لقائه بحقّ ايشان كه اگر يك كلمه صدق
اين عباد از اين دو نفس استماع نموده و نميدانى چه کرده‌اند بكلّ ظلم قيام
نمودند و بعد نظر بالقاى شبهه در قلوب ناس و ستر اعمال مردوده شنيعه خود
مفترياتى بحقّ نسبت داده‌اند كه احصاى آن ممكن نه اميدوارم كه در دنيا و آخرت
خير نه‌بينند و بعذاب دنيا قبل عذاب آخرت معذّب شوند

\end{otherlanguage}

Balī īn ʿibād taṣdīq-i maẓāhir-i ilāhiyyih az qabl namūdih va az baʿd
ham bi-faḍl-i Allāh va mannih khvāhand namūd, va in shāʾ Allāh az
ʿināyat va mukarramat-i ū bi-hīch ḥijābī muḥtajib nashavīm va bi-hīch
manʿī mamnūʿ nagardīm. Fa-va-llāh ānchih az ān mulḥid-i khabīth shinīdih
kidhib-i ṣirf va iftirā-yi baḥt būdih va namīdānī kih bi-nafas kih
muʾānis shudih. Ta-llāhi al-Ḥaqq yaf irru al-jaḥīmu min nafasihi
wa-yastaʿīdhu al-ḥusbānu bi-llāhi min wajhihi wa-liqāʾihi. Bi-ḥaqq-i
īshān kih agar yak kalimih ṣidq īn ʿibād az īn du nafs istimāʿ namūdih,
va namīdānī chih kardih-and. Bi-kull-i ẓulm qiyām namūdand va baʿd naẓar
bi-ilqā-yi shubhih dar qulūb-i nās va satr-i aʿmāl-i mardūdih-yi
shanīʿih-yi khud muftariyātī bi-Ḥaqq nisbat dādih-and kih iḥṣā-yi ān
mumkin nah. Umīdvāram kih dar dunyā va ākhirat khayr nabīnand va
bi-ʿadhāb-i dunyā qabl-i ʿadhāb-i ākhirat muʿadhdhab shavand.

Yes, these servants have affirmed the Manifestations of God in the past,
and through God's grace and favor they will do so in the future. God
willing, through His care and generosity, may no veil conceal Him from
us and no hindrance hold us back. By God, whatever you have heard from
that wicked unbeliever has been sheer falsehood and pure fabrication,
and you do not know with whose breath you have become familiar. By God,
the True One, hell flees from his breath, and the fire of punishment
seeks refuge in God from his face and his presence. By His truth, these
servants have not heard a single truthful word from these two souls, and
you do not know what they have done. They have perpetrated every
injustice and then, with the aim of casting doubts into people's hearts
and concealing their own vile and reprehensible deeds, have attributed
to the True One more fabrications than can be counted. I hope that they
see no good in this world or the next, and that they suffer punishment
in this world before the punishment of the next.

\subsection{268}\label{section-763}

\begin{otherlanguage}{persian}

بارى للّه كلمه ذكر ميشود كه شايد از تصوّر و تصديق و بديهيات و نظريات خود
طاهر و مقدّس شوى و بشريعه تصديق كه عين و حقيقت تصوّر است وارد شوى نقطه
اولى جلّت عظمته ميفرمايند قبل ان يعرفكم نفسه انتم فى نار الحبّ تدخلون
ميفرمايند قبل از آنكه آن جمال قدم بشناساند نفس خود را شما در نار حبّ او
داخل شويد و اين عباد سالهاى سال است كه ذكر و ثناى ايشان را از كلّ عباد
شنيده‌ايم تا اينكه بشريعه قرب وارد شديم و در كلّ حين ديده آنچه از احدى
نديده و نشنيده‌ايم و از كلّ اين امور گذشته آيات الهى بمثل باران نيسانى در
كلّ حين نازل مع‌ذلك چگونه توقّف جايز با آنكه نقطه بيان تصريحا در كلّ بيان
ميفرمايند منتهى جدّ و جهد خود را نموده كه بقدر قول بلى توقّف ننمائى كه
همان قدر در نار خواهى بود و جميع اعمالت هباء ميگردد

\end{otherlanguage}

Bārī li-llāh kalimih dhikr mīshavad kih shāyad az taṣavvur va taṣdīq va
badīhiyyāt va naẓariyyāt-i khud ṭāhir va muqaddas shavī va
bi-sharīʿih-yi taṣdīq kih ʿayn va ḥaqīqat-i taṣavvur ast vārid shavī.
Nuqṭih-yi Ūlā jallat ʿaẓamatuhu mīfarmāyand: qabla an yuʿarrifakum
nafsahu antum fī nāri al-ḥubbi tadkhulūn. Mīfarmāyand qabl az ānkih ān
Jamāl-i Qidam bishināsānad nafs-i khud rā shomā dar nār-i ḥubb-i ū
dākhil shavīd. Va īn ʿibād sālhā-yi sāl ast kih dhikr va thanā-yi īshān
rā az kull-i ʿibād shinīdih-īm tā īnkih bi-sharīʿih-yi qurb vārid shudīm
va dar kull-i ḥīn dīdih ānchih az aḥadī nadīdih va nashinīdih-īm. Va az
kull-i īn umūr gudhashtih āyāt-i ilāhī bi-mithl-i bārān-i naysānī dar
kull-i ḥīn nāzil. Maʿa-dhālik chigūnih tavaqquf jāʾiz, bā ānkih
Nuqṭih-yi Bayān taṣrīḥan dar kull-i Bayān mīfarmāyand muntahā-yi jidd va
jahd-i khud rā namūdih kih bi-qadr-i qawl-i balī tavaqquf nanamāʾī kih
hamān qadr dar nār khvāhī būd va jamīʿ-i aʿmālat habā mīgardad.

In any case, a word is spoken for God's sake, so that perhaps you may be
cleansed and sanctified from your conceptions and assents, your
self-evident truths and theoretical judgments, and reach the water's
edge of assent, which is the very essence and reality of conception. The
Primal Point---exalted be His grandeur---says: ``Before He makes Himself
known to you, enter the fire of love.'' He says that before that Ancient
Beauty makes Himself known, you should enter the fire of His love. For
many years these servants heard all people speak of Him and praise Him,
until we reached the water's edge of nearness and at every moment
witnessed what we had neither seen nor heard from anyone else. Apart
from all these things, the verses of God descend at every moment like
spring rain. How, then, can hesitation be permissible, when the Point of
the Bayán explicitly states throughout the Bayán that you must exert
your utmost effort not to hesitate even for the time it takes to say
``Yes,'' for during that very interval you will be in the fire, and all
your deeds will come to nothing?

\subsection{269}\label{section-764}

\begin{otherlanguage}{persian}

بسيار اين عبد شرم نمود از اين كلمه تو خود تو هم تصديق اين عبد مينمودى و
لكن حجاب حايل شد نسئل اللّه بأن يرفعه بقدرته لتشهد شمس وجه ربّك مشرقاً
مضيئاً طالعاً بين العالمين و اللّه حيرت فوق حيرت است كه چه باز داشته تو را
كه بهتك حرمت حقّ كمر بسته و بمنتهى سعى بمجادله و محاربه بر خواسته فواللّه
سيف اعراضك احدّ من سيف الحديد لانّ منه يجرح جسد الظّاهر و من سيفك قطع اكباد
الوجود من الغيب و الشّهود و لكن انت تفرح فى نفسك كانّك لا ذنب لك نسئل اللّه
بأن يرجع اليك ما عملت و فعلت و ارتكبت او يكفّر عنك خطيئاتك و يهديك الى
نفسه و يرزقك لقائه العزيز المنيع

\end{otherlanguage}

Bisyār īn ʿabd sharm namūd az īn kalimih-yi tu. Khud-i tu ham taṣdīq-i
īn ʿabd mīnamūdī va lākin ḥijāb ḥāʾil shud. Nasʾalu Allāha bi-an
yarfaʿahu bi-qudratihi li-tashhada shamsa wajhi rabbika mushriqan
muḍīʾan ṭāliʿan bayna al-ʿālamīn. Va-llāh ḥayrat fawq-i ḥayrat ast kih
chih bāz dāshtih tu rā kih bi-hatk-i ḥurmat-i Ḥaqq kamar bastih va
bi-muntahā-yi saʿy bi-mujādilih va muḥāribih bar khāstih. Fa-wa-llāh
sayfu iʿrāḍika aḥaddu min sayfi al-ḥadīd li-anna minhu yujraḥu jasadu
al-ẓāhir wa-min sayfika quṭiʿa akbādu al-wujūd min al-ghaybi
wa-al-shuhūd wa-lākin anta tafraḥu fī nafsika ka-annaka lā dhanba laka.
Nasʾalu Allāha bi-an yarjiʿa ilayka mā ʿamilta wa-faʿalta wa-irtakabta
aw yukaffira ʿanka khaṭīʾātika wa-yahdiyaka ilā nafsihi wa-yarzuqaka
liqāʾahu al-ʿazīza al-manīʿ.

This servant was deeply ashamed by these words of yours. You yourself
would have affirmed this servant, but a veil intervened. We ask God to
lift it through His power, so that you may witness the sun of your
Lord's countenance rising, radiant and luminous, among the worlds. By
God, it is bewilderment upon bewilderment: what has held you back, so
that you have set yourself to violating the sanctity of the True One and
have risen with all your might to dispute and wage war? By God, the
sword of your rejection is sharper than a sword of iron, for the latter
wounds the outward body, while your sword has cut through the very
hearts of existence, both unseen and seen. Yet you rejoice within
yourself, as though you were without sin. We ask God either to return to
you what you have done, performed, and perpetrated, or to atone for your
sins, guide you to Himself, and grant you His mighty, unassailable
presence.

\subsection{270}\label{section-765}

\begin{otherlanguage}{persian}

و اينكه اين عباد را عوام ذكر نموده قسم بمربی‌ امكان همين قول شما دليل بر
حجيّت اين امر است چه كه در اوّل ظهور ستّين اعتراضا على الحقّ بعينه همين قول
را ذكر مينمودند و از قبل هم در احيان مظاهر ظهور لسان مشركين باين كلمات
ناطق چنانچه در كتاب مبين نازل قوله تعالى فقال الملأ الّذين كفروا من قومه
ما نراك الّا بشراً مثلنا و ما نراك اتّبعك الّا الّذينهم اراذلنا بادى الرّاى و
ما نرى لكم علينا من فضل بل نظنّكم كاذبين كه جميع مؤمنين بحقّ را اراذل و
عوام ميشمردند و اللّه انعام از قول آن ملقى ملحد ننگ دارند بگو خود را
بشناس و بحدّ خود تكلّم كن قدرت از آنچه القاء نموده مشهود و معلوم شد انّا لا
نشرك بربّنا احدا و لا نتوجّه الى احد دونه بحول اللّه و قوّته

\end{otherlanguage}

Va īnkih īn ʿibād rā ʿavāmm dhikr namūdih, qasam bi-Murabbī-yi imkān
hamīn qawl-i shomā dalīl bar ḥujjiyyat-i īn Amr ast, chih kih dar
avval-i ẓuhūr-i sittīn iʿtirāḍan ʿalā al-Ḥaqq bi-ʿaynih hamīn qawl rā
dhikr mīnamūdand. Va az qabl ham dar aḥyān-i maẓāhir-i ẓuhūr lisān-i
mushrikīn bi-īn kalimāt nāṭiq, chunānchih dar Kitāb-i Mubīn nāzil,
qawluhu taʿālā: fa-qāla al-malaʾu alladhīna kafarū min qawmihi mā narāka
illā basharan mithlanā wa-mā narāka ittabaʿaka illā alladhīnhum
arādhilunā bādiya al-raʾyi wa-mā narā lakum ʿalaynā min faḍlin bal
naẓunnukum kādhibīn. Kih jamīʿ-i muʾminīn bi-Ḥaqq rā arādhil va ʿavāmm
mīshumurdand. Va-llāh anʿām az qawl-i ān mulqī-yi mulḥid nang dārand.
Bigū khud rā bishinās va bi-ḥadd-i khud takallum kun. Qadrat az ānchih
ilqā namūdih mashhūd va maʿlūm shud. Innā lā nushriku bi-rabbinā aḥadan
wa-lā natawajjahu ilā aḥadin dūnahu bi-ḥawli Allāhi wa-quwwatihi.

As for your calling these servants common people, I swear by the
Educator of all existence that this very statement of yours is proof of
the validity of this Cause. For at the beginning of the Revelation of
the year sixty, they voiced precisely the same statement as an objection
to the True One. Previously, too, in the days of the Manifestations of
Revelation, the tongues of the idolaters uttered these words, as
revealed in the manifest Book---exalted are His words: ``The chiefs
among His people who disbelieved said, `We see you only as a human being
like ourselves, and we see that none but the lowliest among us have
followed you, without reflection. Nor do we see that you have any
advantage over us; rather, we think you are liars.'\,'' Thus they
counted all who believed in the True One as lowly and common people. By
God, even livestock are ashamed of the words of that unbeliever who
instilled these ideas. Say: Know yourself, and speak within your limits.
Your worth has become manifest and known from what you have put forward.
Through God's power and might, we associate no one with our Lord and
turn to no one besides Him.

\subsection{271}\label{section-766}

\begin{otherlanguage}{persian}

در اينمقام چون مشاهده شد كه معرضين از كأس غرور شاربند و از كوب كبر مشروب
چه كه عارفين باللّه را عوام ذكر نموده‌اند لذا لازم شد كه اين عبد بعضى از
آيات نقطه اولى روح ماسواه فداه را ذكر نمايم تا اقلّا عالم از جاهل تميز
داده شود قوله عزّ ذكره قل ما خلق اللّه شيئاً اعزّ من العلم عنده كلّ به يخلقون
كلّ به يرزقون كلّ به يميتون كلّ به يحيون ذلك علمكم باللّه ثمّ بمظهر نفسه فى
كلّ ظهور ثمّ ما يظهر من عند مظهر نفسه دون ذلك لم يكن عند اللّه بعلم ان انتم
تعلمون و عند ظهور الأخرى كلّ العلم علمكم به و ما يظهر من عنده لا بما قد
ظهر من قبل و علمكم بما قدّر من قبل فانكم انتم حينئذ به لا توجرون بل عنه
تسئلون ينبغى فى كلّ ظهور ادلّاء ذلك ان يحيطون بما نزل فى ذلك الظّهور لا على
فصَل الّتى انتم تفصلون الى ان قال عزّ ذكره بل الأمر ان تتركن كلّ ذلك و تعد
من انفسكم و علمكم ان انتم باللّه و آياته مؤمنون

\end{otherlanguage}

Dar īn maqām chūn mushāhidih shud kih muʿriḍīn az kaʾs-i ghurūr
shāriband va az kūb-i kibr mashrūb, chih kih ʿārifīn bi-llāh rā ʿavāmm
dhikr namūdih-and, lidhā lāzim shud kih īn ʿabd baʿḍī az āyāt-i
Nuqṭih-yi Ūlā rūḥ-i mā sivāhu fidāh rā dhikr namāyam tā aqallan ʿālim az
jāhil tamyīz dādih shavad. Qawluhu ʿazza dhikruhu: qul mā khalaqa Allāhu
shayʾan aʿazza min al-ʿilmi ʿindahu kullun bihi yukhlaqūn kullun bihi
yurzaqūn kullun bihi yumītūn kullun bihi yuḥyawn dhālika ʿilmukum
bi-llāhi thumma bi-maẓhari nafsihi fī kulli ẓuhūrin thumma mā yaẓharu
min ʿindi maẓhari nafsihi dūna dhālika lam yakun ʿinda Allāhi bi-ʿilmin
in antum taʿlamūn wa-ʿinda ẓuhūri al-ukhrā kullu al-ʿilmi ʿilmukum bihi
wa-mā yaẓharu min ʿindihi lā bi-mā qad ẓahara min qablu wa-ʿilmukum
bi-mā quddira min qablu fa-innakum antum ḥīnaʾidhin bihi lā tuʾjarūn bal
ʿanhu tusʾalūn yanbaghī fī kulli ẓuhūrin adillāʾa dhālika an yuḥīṭūna
bi-mā nazala fī dhālika al-ẓuhūri lā ʿalā fiṣali allatī antum
tufaṣṣilūn. Ilā an qāla ʿazza dhikruhu: bal al-amru an tatrukunna kulla
dhālika wa-taʿudda min anfusikum wa-ʿilmikum in antum bi-llāhi
wa-āyātihi muʾminūn.

At this point, since it was observed that those who turn away are
drinking from the chalice of conceit and are drunk from the cup of
pride---for they have called those who know God common people---it
became necessary for this servant to quote some of the verses of the
Primal Point---may the spirit of all else be sacrificed for Him---so
that at least the learned may be distinguished from the ignorant. He
says---exalted be His remembrance: ``Say: God has created nothing more
precious in His sight than knowledge. Through it all are created;
through it all are sustained; through it all are put to death; through
it all are brought to life. This is your knowledge of God, then of the
Manifestation of His Self in every Revelation, and then of what appears
from the Manifestation of His Self. Apart from this, nothing is
knowledge in God's sight, if you but knew. At the appearance of the next
Revelation, all knowledge consists in your knowledge of Him and of what
appears from Him, not of what appeared before, nor in your knowledge of
what was ordained before. For at that time you will not be rewarded for
it, but rather questioned about it. In every Revelation, its guides
should comprehend what has been sent down in that Revelation, not
according to the divisions you make.'' He continues---exalted be His
remembrance---until He says: ``Rather, the command is that you abandon
all that and count it as belonging to yourselves and your knowledge, if
you believe in God and His verses.''

\subsection{272}\label{section-767}

\begin{otherlanguage}{persian}

حال قدرى ملاحظه در اين كلمات الهيّه لازم كه شايد علمى را كه لم يزل عند
اللّه محبوب و محمود بوده ادراك نمايند و معرضين باللّه را عالم ندانيد
فو‌الّذى ايّد الرّوح بالرّوح اليوم اگر كلّ من فى السّموات و الأرض بتمام علوم
ارض ظاهر شوند و در امر بديع توقّف نمائيد از جهلا عند اللّه مذكور مؤيّد اين
بيان نصّ بيان نقطه بيان روح ما سواه فداه قوله جلّ و عزّ مخاطباً لحرف الحىّ ان
آمنتَبه سواء علمت من شیء او لا علمت فلا ينقص من فضلك من شیء و لكن ان
احتجبت عنه ان علمت كلّ علم لا ينفعك بنصّ آن جوهر وجود اليوم اطلاق اسم عالم
بر احدى نخواهد شد مگر آن نفوسيكه بقميص ايمان امر بديع مزيّن شده‌اند

\end{otherlanguage}

Ḥāl qadrī mulāḥiẓih dar īn kalimāt-i ilāhiyyih lāzim kih shāyad ʿilmī rā
kih lam yazal ʿinda Allāh maḥbūb va maḥmūd būdih idrāk namāyand va
muʿriḍīn bi-llāh rā ʿālim nadānīd. Fa-wa-alladhī ayy ada al-Rūḥa
bi-al-rūḥ al-yawm agar kull-i man fī al-samāwāti wa-al-arḍ bi-tamām-i
ʿulūm-i arḍ ẓāhir shavand va dar Amr-i Badīʿ tavaqquf namāʾīd az juhalā
ʿinda Allāh madhkūr. Muʾayyid-i īn bayān naṣṣ-i bayān-i Nuqṭih-yi Bayān
rūḥ-i mā sivāhu fidāh, qawluhu jalla wa-ʿazza mukhāṭiban li-Ḥarfi
al-Ḥayy: in āmanta bihi sawāʾun ʿalimta min shayʾin aw lā ʿalimta fa-lā
yanquṣu min faḍlika min shayʾin wa-lākin in iḥtajabta ʿanhu in ʿalimta
kulla ʿilmin lā yanfaʿuka. Bi-naṣṣ-i ān Jawhar-i Vujūd al-yawm iṭlāq-i
ism-i ʿālim bar aḥadī nakhvāhad shud magar ān nufūsī kih bi-qamīṣ-i
īmān-i Amr-i Badīʿ muzayyan shudih-and.

Now some reflection on these divine words is needed, so that perhaps
they may comprehend the knowledge that has always been beloved and
praiseworthy in God's sight, and you may not regard those who turn away
from God as learned. By Him who strengthened the Spirit through the
spirit, today, even if all who are in the heavens and on earth were to
appear with all earthly learning, and you were to hesitate before this
wondrous Cause, you would be counted among the ignorant in God's sight.
Supporting this statement is the explicit utterance of the Point of the
Bayán---may the spirit of all else be sacrificed for Him. He
says---glorified and exalted be He---addressing the Letter of the
Living: ``If you believe in Him, whether you know anything or do not
know, nothing whatsoever will diminish your merit. But if you are veiled
from Him, even if you possess all knowledge, it will not benefit you.''
According to the explicit words of that Essence of existence, today the
designation `learned' will apply to no one except those souls adorned
with the robe of faith in this wondrous Cause.

\subsection{273}\label{section-768}

\begin{otherlanguage}{persian}

عجب است مع اين آيات واضحات مع‌ذلك مؤمنين باللّه را عوام ناميده‌ايد چه زود
نفس رحمن از آن قلوب منقطع شده و از نفس سبحان محتجب گشته‌اند چه كه اگر اقلّ
من ذر رايحه علم و حكمتش بر نفوس معرضه مرور ميفرمود هرگز باين كلماتى كه
پست‌ترين ملّت فرقان سالها بآن مستدلّ بوده‌اند مشغول نميشدند فواحسرتا عليهم
بما ارتكبوا فى جنب اللّه و كانوا من الخاسرين

\end{otherlanguage}

ʿAjab ast maʿ-i īn āyāt-i vāḍiḥāt maʿa-dhālik muʾminīn bi-llāh rā ʿavāmm
nāmīdih-īd. Chih zūd nafas-i Raḥmān az ān qulūb munqaṭiʿ shudih va az
Nafs-i Subḥān muḥtajib gashtih-and, chih kih agar aqall-i min dharr
rāyiḥih-yi ʿilm va ḥikmatash bar nufūs-i muʿriḍih murūr mīfarmūd hargiz
bi-īn kalimātī kih past-tarīn millat-i Furqān sālhā bi-ān mustadill
būdih-and mashghūl namīshudand. Fa-wā ḥasratā ʿalayhim bi-mā irtakabū fī
janbi Allāhi wa-kānū min al-khāsirīn.

It is astonishing that, despite these clear verses, you have
nevertheless called those who believe in God common people. How quickly
the breath of the All-Merciful has been cut off from those hearts, and
they have become veiled from the sanctified One! For if less than a mote
of the fragrance of His knowledge and wisdom had passed over these
wayward souls, they would never have occupied themselves with these
words, which the lowliest followers of the Qur'an have used as arguments
for years. Alas for them, for what they have perpetrated against God!
They are among the losers.

\subsection{274}\label{section-769}

\begin{otherlanguage}{arabic}

نغمه ديگر آن سلطان احديّه را بشنو مخاطباً لحرف الحىّ ميفرمايند قوله عزّ
بيانه فانّه يعرف كلشیء نفسه بنفسه فمن يقدر ان يعرفه بغيره الا انّك لو
ادركت يوم ظهوره ان عرفته باعلم علمآء البيان و ان رأيته واقفاً فى امره ثمّ
ذكرت عليه اسم الانسانيّة ما انفيت حروف النّفى لاثبات مظهر الأحديّة انتهى

\end{otherlanguage}

Naghmih-yi dīgar-i ān Sulṭān-i Aḥadiyyih rā bishinav. Mukhāṭiban
li-Ḥarfi al-Ḥayy mīfarmāyand, qawluhu ʿazza bayānuhu: fa-innahu
yuʿarrifu kulla shayʾin nafsahu bi-nafsihi fa-man yaqdiru an yuʿarrifahu
bi-ghayrihi illā innaka law adrakta yawma ẓuhūrihi in ʿaraftahu
bi-aʿlami ʿulamāʾi al-Bayāni wa-in raʾaytahu wāqifan fī amrihi thumma
dhakarta ʿalayhi isma al-insāniyyati mā anfayta ḥurūfa al-nafyi
li-ithbāti Maẓhari al-Aḥadiyyah. Intahā.

Listen to another melody of that Sovereign of Oneness. Addressing the
Letter of the Living, He says---exalted be His utterance: ``He Himself
makes Himself known to all things through Himself. Who, then, can make
Him known through anything other than Him? Yet, if you attain the day of
His Revelation, and if you recognize someone as the most learned of the
learned of the Bayán, but see him hesitating before His Cause and then
apply to him the designation `human,' you have not negated the letters
of negation in order to affirm the Manifestation of Oneness.'' End of
quotation.

\subsection{275}\label{section-770}

\begin{otherlanguage}{arabic}

فما اعزّ كلماته و احكم اشاراته و اظهر برهانه و اجلّ بيانه و لكن انتم يا
ملأ البيان ظلمتم على نفس اللّه و كينونته و انّه ما قصّر فى بيانه قد شرع لكم
كلّ مناهج و بين لكم كلّ ما كان مستوراً عنكم و فصّل لكم من كلّ شیء تفصيلاً رحمة
من عنده عليكم و هدى و ذكرى للعالمين و انتم اشتبهتم امر اللّه على انفسكم
كما اشتبهوا ملل القبل على انفسهم و بذلك ضل سعيهم فى الدّنيا و الآخرة و
كانوا قوم سوء اخسرين

\end{otherlanguage}

Fa-mā aʿazza kalimātihi wa-aḥkama ishārātihi wa-aẓhara burhānahu
wa-ajalla bayānahu wa-lākin antum yā malaʾa al-Bayān ẓalamtum ʿalā nafsi
Allāhi wa-kaynūnatihi wa-innahu mā qaṣṣara fī bayānihi qad sharaʿa lakum
kulla manāhija wa-bayyana lakum kulla mā kāna mastūran ʿankum wa-faṣṣala
lakum min kulli shayʾin tafṣīlan raḥmatan min ʿindihi ʿalaykum wa-hudan
wa-dhikrā li-al-ʿālamīn wa-antum ishtabahtum amra Allāhi ʿalā anfusikum
kamā ishtabahū milalu al-qabli ʿalā anfusihim wa-bi-dhālika ḍalla
saʿyuhum fī al-dunyā wa-al-ākhirati wa-kānū qawma sawʾin akhsarīn.

How precious are His words, how firm His allusions, how manifest His
proof, and how exalted His utterance! Yet you, O people of the Bayán,
have wronged the Self of God and His very being. He has left nothing
lacking in His exposition: He has laid out every path for you, made
clear to you all that was concealed from you, and explained everything
to you in detail, as a mercy to you from Him and as guidance and a
reminder to the worlds. Yet you have confused the Cause of God for
yourselves, just as the religious communities of the past confused it
for themselves. Thus their efforts came to nothing in this world and the
next, and they were an evil people, the greatest losers.

\subsection{276}\label{section-771}

\begin{otherlanguage}{persian}

مع آنكه آن جمال رحمن ميفرمايد كه اگر اعلم علماء بيان را در يوم ظهور و
طلوع و اشراق نيّر اعظم متوقّف ديدى و ذكر انسانيّت بر او نمودى هرآينه از
نفوسى خواهى بود كه حروف نفى را در آن يوم نفى ننموده‌اند از براى اثبات
مظهر احديّه

\end{otherlanguage}

Maʿ-i ānkih ān Jamāl-i Raḥmān mīfarmāyad kih agar aʿlam-i ʿulamā-yi
Bayān rā dar yawm-i ẓuhūr va ṭulūʿ va ishrāq-i Nayyir-i Aʿẓam mutavaqqif
dīdī va dhikr-i insāniyyat bar ū namūdī har-āyinih az nufūsī khvāhī būd
kih ḥurūf-i nafy rā dar ān yawm nafy nanamūdih-and az barā-yi ithbāt-i
Maẓhar-i Aḥadiyyih.

And yet that All-Merciful Beauty says that if, on the day of the
revelation, rising, and radiance of the Most Great Luminary, you see the
most learned of the learned of the Bayán hesitating and apply to him the
designation `human,' you will assuredly be among those souls who have
not negated the letters of negation on that day in order to affirm the
Manifestation of Oneness.

\subsection{277}\label{section-772}

\begin{otherlanguage}{persian}

اليوم هر نفسى بر احدى از معرضين من اعلاهم او من ادناهم ذكر انسانيّت نمايد
از جميع فيوضات رحمانى محروم است تا چه رسد كه بخواهد از براى آن نفوس
اثبات رتبه و مقام نمايد

\end{otherlanguage}

Al-yawm har nafsī bar aḥadī az muʿriḍīn min aʿlāhum aw min adnāhum
dhikr-i insāniyyat namāyad az jamīʿ-i fuyūḍāt-i raḥmānī maḥrūm ast tā
chih rasad kih bikhvāhad az barā-yi ān nufūs ithbāt-i rutbih va maqām
namāyad.

Today, anyone who applies the designation ``human'' to any of those who
turn away, whether the highest or the lowest among them, is deprived of
all the outpourings of the All-Merciful---how much more so if he wishes
to establish a rank or station for those souls.

\subsection{278}\label{section-773}

\begin{otherlanguage}{persian}

ملاحظه نموده كه مخاطب آيه منزله يكى از حروفات حىّ است كه اسبق از كلّ بيان
بوده و خواهد بود بنصّ خود نقطه بيان كه ميفرمايند ايشانند اوّل خلق بيان و
ميفرمايند كلّ بوجود اين حروفات بشريعه احديّه و ظهور الهيّه فايز شدند و
مع‌ذلك بآن نفوس باين قسم تحديد فرموده چنانچه ذكر شد فتفكّروا يا اولو
الأبصار فتفكّروا يا اولو الأنظار فتفكّروا يا اولو الأخيار فتفكّروا يا اولو
الأبرار فتفكّروا يا اولو العلم و الحكمة و البيان

\end{otherlanguage}

Mulāḥiẓih namūdih kih mukhāṭab-i āyih-yi munazzalih yakī az ḥurūfāt-i
Ḥayy ast kih asbaq az kull-i Bayān būdih va khvāhad būd bi-naṣṣ-i khud-i
Nuqṭih-yi Bayān kih mīfarmāyand īshānand avval-i khalq-i Bayān va
mīfarmāyand kull bi-vujūd-i īn ḥurūfāt bi-sharīʿih-yi aḥadiyyih va
ẓuhūr-i ilāhiyyih fāʾiz shudand va maʿa-dhālik bi-ān nufūs bi-īn qism
taḥdīd farmūdih chunānchih dhikr shud fa-tafakkarū yā ulū al-abṣār
fa-tafakkarū yā ulū al-anẓār fa-tafakkarū yā ulū al-akhyār fa-tafakkarū
yā ulū al-abrār fa-tafakkarū yā ulū al-ʿilm wa-al-ḥikmah wa-al-bayān.

Consider that the person addressed in the revealed verse is one of the
Letters of the Living, who has preceded, and will precede, all the
people of the Bayán, according to the explicit statement of the Point of
the Bayán Himself: ``They are the first creation of the Bayán.'' He also
says: ``Through the existence of these Letters, all attained the law of
Divine Oneness and the divine Revelation.'' Yet He issued such a warning
to those souls, as has been mentioned. Reflect, then, O people of
insight! Reflect, O people of discernment! Reflect, O people of
goodness! Reflect, O people of righteousness! Reflect, O people of
knowledge, wisdom, and utterance!

\subsection{279}\label{section-774}

\begin{otherlanguage}{persian}

قل تاللّه كلّما نزل فى البيان قد رجع الى نقطة و انّها هى حينئذ مقبوضة بكفّ
القدرة و الاقتدار و نريد ان نفصّلها مرّهً اخرى رغماً لا نفسكم يا ملأ
المشركين كما فصّلناها مرّه بعد مرّة و هذه من مرّة فزعت عنها سكّان السّموات و
الأرض ثمّ سكّان مداين الأسماء ان انتم من العارفين قل انّه قد كان نفسى و ان
حينئذ قد ظهرت نفسه بالحقّ و نطقت بين السّموات و الأرضين و يقول بالحقّ ان كلّ
الآیات قد رجعت الى آياتى و بعثت فى كلماتى و ظهرت باذنى و نزلت من سماء
مشيّتى ان انتم من الموقنين ثمّ كلّ الظّهورات انتهت بظهورى و كلّ المشارق قد
اشرقت باشراقى و كلّ المطالع قد استبهت بطلوعى و ظهورى فاين ابصار النّاظرة
فاين قلوب العارفة فاين افئدة المقدّسة و اين انفس الزّكيّة و اين صدور
المنيرة ليعرف و يفقه و ينظر و يشهد ما ظهر و اشرق من افق اللّه المقتدر
العزيز العظيم قل موتوا بغيظكم انّه قد ظهر بالحقّ و لايخاف من احد و ينتظر
ما قدّر له من لدن محبوبه المقتدر العليم الحكيم

\end{otherlanguage}

Qul ta-llāhi kullamā nazala fī al-Bayān qad rajaʿa ilā nuqṭah wa-innahā
hiya ḥīnaʾidhin maqbūḍah bi-kaff al-qudrah wa-al-iqtidār wa-nurīdu an
nufaṣṣilahā marratan ukhrā raghman li-anfusikum yā malaʾ al-mushrikīn
kamā faṣṣalnāhā marratan baʿda marrah wa-hādhihi min marrah faziʿat
ʿanhā sukkān al-samāwāt wa-al-arḍ thumma sukkān madāʾin al-asmāʾ in
antum min al-ʿārifīn qul innahu qad kāna nafsī wa-in ḥīnaʾidhin qad
ẓaharat nafsuhu bi-al-ḥaqq wa-naṭaqat bayna al-samāwāt wa-al-araḍīn
wa-yaqūlu bi-al-ḥaqq inna kull al-āyāt qad rajaʿat ilā āyātī wa-buʿithat
fī kalimātī wa-ẓaharat bi-idhnī wa-nazalat min samāʾ mashīʾatī in antum
min al-mūqinīn thumma kull al-ẓuhūrāt intahat bi-ẓuhūrī wa-kull
al-mashāriq qad ashraqat bi-ishrāqī wa-kull al-maṭāliʿ qad istabhat
bi-ṭulūʿī wa-ẓuhūrī fa-ayna abṣār al-nāẓirah fa-ayna qulūb al-ʿārifah
fa-ayna afʾidah al-muqaddasah wa-ayna anfus al-zakiyyah wa-ayna ṣudūr
al-munīrah li-yaʿrif wa-yafqah wa-yanẓur wa-yashhad mā ẓahara wa-ashraqa
min ufuq Allāh al-Muqtadir al-ʿAzīz al-ʿAẓīm qul mūtū bi-ghayẓikum
innahu qad ẓahara bi-al-ḥaqq wa-lā yakhāfu min aḥad wa-yantaẓiru mā
quddira lahu min ladun maḥbūbihi al-Muqtadir al-ʿAlīm al-Ḥakīm.

Say: By God, all that was revealed in the Bayán has returned to a single
point, and it is now held in the grasp of power and might. We intend to
unfold it once again, despite you, O company of those who associate
partners with God, just as We have unfolded it time after time. This is
a time that has struck terror into the inhabitants of the heavens and
the earth, and then the inhabitants of the cities of names, if you are
among those who understand. Say: He was My own Self, and now His Self
has appeared in truth and spoken between the heavens and the earths. He
says in truth: ``All verses have returned to My verses, have been raised
to life in My words, have appeared by My permission, and have descended
from the heaven of My will, if you are among those who are certain. All
revelations have culminated in My Revelation, all daysprings have shone
through My radiance, and all dawning places have grown resplendent
through My rising and My appearance. Where are the eyes that see? Where
are the hearts that understand? Where are the sanctified souls? Where
are the pure selves? Where are the illumined breasts, that they may
recognize, comprehend, behold, and witness what has appeared and shone
from the horizon of God, the All-Powerful, the Almighty, the Most
Great?'' Say: Die in your rage! He has appeared in truth and fears no
one. He awaits what has been ordained for Him by His Beloved, the
All-Powerful, the All-Knowing, the All-Wise.

\subsection{280}\label{section-775}

\begin{otherlanguage}{arabic}

و يقول فى كلّ حين مخاطباً لاسيافكم و رماحكم و سهامكم هلمّوا و تعالوا لانّى
اكون مشتاقاً بكم و يشهد بذلك شوقى و اشتياقى ثمّ شغفى و نغماتى لو انتم من
الشّاعرين انّ الّذى قد ظهر بالحقّ هو محبوبی‌ و انا محبوبه و هو نفسى و انا
نفسه و ما يخرج من لسانى هو بيانه الابدع الاملح الأقدس المنيع و من يكون
عنده اقلّ من الشّعر شعور و اقلّ من الذّرّ بصر ليشهد ظهور اللّه و عظمته و يجد
روايح قدسه من نفحات هذه النّغمات و يكون من الرّاسخين على هذا الأمر الأعظم
المبين

\end{otherlanguage}

Wa-yaqūlu fī kull ḥīn mukhāṭiban li-asyāfikum wa-rimāḥikum wa-sihāmikum
halummū wa-taʿālaw li-annī akūnu mushtāqan bikum wa-yashhadu bi-dhālika
shawqī wa-ishtiyāqī thumma shaghafī wa-naghamātī law antum min
al-shāʿirīn inna alladhī qad ẓahara bi-al-ḥaqq huwa maḥbūbī wa-anā
maḥbūbuhu wa-huwa nafsī wa-anā nafsuhu wa-mā yakhruju min lisānī huwa
bayānuhu al-abdaʿ al-amlaḥ al-aqdas al-manīʿ wa-man yakūnu ʿindahu aqall
min al-shaʿr shuʿūr wa-aqall min al-dharr baṣar la-yashhadu ẓuhūr Allāh
wa-ʿaẓamatahu wa-yajidu rawāʾiḥ qudsihi min nafaḥāt hādhihi al-naghamāt
wa-yakūnu min al-rāsikhīn ʿalā hādhā al-amr al-aʿẓam al-mubīn.

At every moment He says, addressing your swords, your spears, and your
arrows: ``Come, draw near, for I yearn for you!'' My longing and
yearning, My ardor and My melodies bear witness to this, if you are
among those who perceive. He who has appeared in truth is My Beloved,
and I am His beloved; He is My Self, and I am His Self. What proceeds
from My tongue is His wondrous, most gracious, most holy, and
unassailable utterance. Anyone who possesses awareness finer than a hair
and sight smaller than a speck of dust will surely witness the
Revelation of God and His grandeur, perceive the fragrances of His
holiness in the breaths of these melodies, and be among those who stand
firm in this most great and manifest Cause.

\subsection{281}\label{section-776}

\begin{otherlanguage}{arabic}

قل انّى اكون حيّاً بحيوته و اتحرّك بحبّه كما انّه كان حيّاً بحيوتى و ما تنفّس الّا
بذكرى و ثنائى بين العالمين نفسى فداه و نفسه فداى كلّى فداه و كلّه فداى
كلّما اقول انّه فداى هو نفس قوله بأنّه فدائه ایّاکم ان تشهدوا الفرق يا ملأ
المتبصّرين كذلك نزل اللّه حينئذ من آيات عظمته على طور الفؤاد و قد اظهرناها
لكم لتشكروا اللّه بارئكم يا ملأ البيان و لا تكوننّ من المحتجبين ان

\end{otherlanguage}

Qul innī akūnu ḥayyan bi-ḥayātihi wa-ataḥarraku bi-ḥubbihi kamā innahu
kāna ḥayyan bi-ḥayātī wa-mā tanaffasa illā bi-dhikrī wa-thanāʾī bayna
al-ʿālamīn nafsī fidāhu wa-nafsuhu fidāya kullī fidāhu wa-kulluhu fidāya
kullamā aqūlu innahu fidāya huwa nafs qawlihi bi-annahu fidāʾuhu iyyākum
an tashhadū al-farq yā malaʾ al-mutabaṣṣirīn kadhālika nazzala Allāh
ḥīnaʾidhin min āyāt ʿaẓamatihi ʿalā ṭūr al-fuʾād wa-qad aẓharnāhā lakum
li-tashkurū Allāh bāriʾakum yā malaʾ al-Bayān wa-lā takūnanna min
al-muḥtajibīn in.

Say: I live through His life and am moved by His love, just as He lived
through My life and breathed only in remembrance and praise of Me among
the worlds. May My soul be sacrificed for Him, and His soul for Me; may
My whole being be sacrificed for Him, and His whole being for Me.
Whenever I say that He is a sacrifice for Me, it is His very statement
that He is a sacrifice for Him. Beware of seeing any distinction, O
company of the discerning! Thus has God now sent down verses of His
grandeur upon the Sinai of the heart. We have made them manifest to you
so that you may thank God, your Maker, O people of the Bayán, and not be
among those who are veiled. If

\subsection{282}\label{section-777}

\begin{otherlanguage}{arabic}

يا ايّها المستغرق فى بحر الاشارات فاعلم بانّ اعراضك و اعتراضك لا يرجع الى
نفسى بل على اللّه ربّی‌ و ربّك و ربّ آبائنا الأوّلين لانّى عبد آمنت به و برسله
و صفوته و لا اجد لنفسى من وجود لانّه قد بعثنى بالحقّ و ارسلنى على العالمين
انّك لو تريد الاعتراض اذهب اليه و قل بأىّ جهة بعثت الّذى به فزعت كينونتى ثمّ
حقايق المشركين و لو وجدته و تكون مستطيعاً فى نفسك فاذكر لى كلمة عنده لعلّ
يخلّصنى منك و من ملأ الأشرار و يجعلنى من المنقطعين قل اى ربّ انت تعرف ملأ
البيان لم ارسلت عبدك هذا اذاً فاشهده تحت اسيافهم و مخاليبهم و اظفارهم يا
من بيدك ملكوت ملك السّموات و الأرضين ارحمه يا مالك الملوك ثمّ انصره بالحقّ
و انّك انت المقتدر المتعالى الغالب العزيز المريد

\end{otherlanguage}

Yā ayyuhā al-mustaghriq fī baḥr al-ishārāt fa-iʿlam bi-anna iʿrāḍaka
wa-iʿtirāḍaka lā yarjiʿu ilā nafsī bal ʿalā Allāh rabbī wa-rabbika
wa-rabb ābāʾinā al-awwalīn li-annī ʿabd āmantu bihi wa-bi-rusulihi
wa-ṣafwatihi wa-lā ajidu li-nafsī min wujūd li-annahu qad baʿathanī
bi-al-ḥaqq wa-arsalanī ʿalā al-ʿālamīn innaka law turīdu al-iʿtirāḍ
idhhab ilayhi wa-qul bi-ayyi jihah baʿathta alladhī bihi faziʿat
kaynūnatī thumma ḥaqāʾiq al-mushrikīn wa-law wajadtahu wa-takūnu
mustaṭīʿan fī nafsika fa-udhkur lī kalimah ʿindahu laʿalla yukhalliṣanī
minka wa-min malaʾ al-ashrār wa-yajʿalanī min al-munqaṭiʿīn qul ay rabb
anta taʿrifu malaʾ al-Bayān lima arsalta ʿabdaka hādhā idhan fa-ashhadhu
taḥta asyāfihim wa-makhālībihim wa-aẓfārihim yā man bi-yadika malakūt
mulk al-samāwāt wa-al-araḍīn irḥamhu yā Mālik al-Mulūk thumma unṣurhu
bi-al-ḥaqq wa-innaka anta al-Muqtadir al-Mutaʿālī al-Ghālib al-ʿAzīz
al-Murīd.

O you who are immersed in the sea of allusions! Know that your rejection
and objections are directed not at Me, but at God, My Lord and your Lord
and the Lord of our forefathers. For I am a servant who has believed in
Him, His Messengers, and His chosen ones. I find no existence of My own,
for He has raised Me up in truth and sent Me to the worlds. If you wish
to object, go to Him and say: ``Why did You raise up One through whom my
being, and then the realities of those who associate partners with You,
have been struck with terror?'' If you find Him and are able, speak a
word on My behalf in His presence, that He may deliver Me from you and
the company of the wicked and make Me one of those who are detached.
Say: ``O Lord, You know the people of the Bayán. Why did You send this
servant of Yours? Behold Him, then, beneath their swords, their claws,
and their talons. O You in whose hand is the dominion of the kingdom of
the heavens and the earths! Have mercy on Him, O King of kings, and aid
Him in truth. You are indeed the All-Powerful, the Most Exalted, the
All-Conquering, the Almighty, He who wills.''

\subsection{283}\label{section-778}

\begin{otherlanguage}{arabic}

فو‌الّذى ارسلنى بالحقّ انّى اكون متحيّراً من فعله فى حقّى لأنّه بعثنى فى ایّام ما
هبّت فى مثلها نفحات اللّه المقتدر العزيز العليم و بذلك ابتليت بما لا ابتلى
به احد من الأصفیاء و يشهد بذلك نفسه العليم الخبير هل تظنّ بانّى انطق من
تلقآء نفسى و اكون موجود بوجودى لا فو‌ربّ العالمين بل يحركنى ارياح مشيّته
كيف شآء و اراد و من كان له دراية ليجد من اهتزاز نفسى اهتزاز اللّه الملك
المقتدر الغالب القدير هل تقدر ان تسكن تلقآء نفسك حين الّذى يأخذك حمّى
الرّعد الّذى به يرتعش كلّ الأعضا لا فو‌الّذى خلقك بالحقّ لو انت من المستشعرين
فكما يحرّكك و لا تقدر ان تسكن كذلك تحركنى ارياح مشيّة اللّه و انّك ان تريد
ان تعترض فاعترض عليه و ما انا الّا عبد منيب و لم اجد لنفسى سكوناً و لا
قرارا و لاحركة الّا بعد امره و لا ينكر ذلك الّا كلّ مغلّ معتد اثيم

\end{otherlanguage}

Fa-wa-alladhī arsalanī bi-al-ḥaqq innī akūnu mutaḥayyiran min fiʿlihi fī
ḥaqqī li-annahu baʿathanī fī ayyām mā habbat fī mithlihā nafaḥāt Allāh
al-Muqtadir al-ʿAzīz al-ʿAlīm wa-bi-dhālika ubtulītu bi-mā lā ubtuliya
bihi aḥad min al-aṣfiyāʾ wa-yashhadu bi-dhālika nafsuhu al-ʿAlīm
al-Khabīr hal taẓunnu bi-annī anṭiqu min tilqāʾ nafsī wa-akūnu mawjūd
bi-wujūdī lā fa-wa-rabb al-ʿālamīn bal yuḥarrikunī aryāḥ mashīʾatihi
kayfa shāʾa wa-arāda wa-man kāna lahu dirāyah la-yajidu min ihtizāz
nafsī ihtizāz Allāh al-Malik al-Muqtadir al-Ghālib al-Qadīr hal taqdiru
an taskuna tilqāʾ nafsika ḥīna alladhī yaʾkhudhuka ḥummā al-raʿd alladhī
bihi yartaʿishu kull al-aʿḍā lā fa-wa-alladhī khalaqaka bi-al-ḥaqq law
anta min al-mustashʿirīn fa-kamā yuḥarrikuka wa-lā taqdiru an taskuna
kadhālika tuḥarrikunī aryāḥ mashīʾat Allāh wa-innaka in turīdu an
taʿtariḍ fa-iʿtariḍ ʿalayhi wa-mā anā illā ʿabd munīb wa-lam ajid
li-nafsī sukūnan wa-lā qarāran wa-lā ḥarakah illā baʿda amrihi wa-lā
yunkiru dhālika illā kull mughill muʿtad athīm.

By Him who has sent Me in truth, I am bewildered by His dealings with
Me, for He has raised Me up in days the like of which have never been
stirred by the breaths of God, the All-Powerful, the Almighty, the
All-Knowing. Through this I have been afflicted as none of His chosen
ones has been afflicted. His own Self, the All-Knowing, the
All-Informed, bears witness to this. Do you imagine that I speak of My
own accord and exist through an existence of My own? No, by the Lord of
the worlds! Rather, the winds of His will move Me as He wishes and
desires. Anyone possessed of understanding will perceive, in the
stirring of My being, the stirring of God, the King, the All-Powerful,
the All-Conquering, the Omnipotent. Can you keep yourself still when a
shivering fever seizes you and makes every limb tremble? No, by Him who
created you in truth, if you are among those who perceive! Just as it
moves you and you cannot keep still, so the winds of the will of God
move Me. If you wish to object, object to Him. I am only a servant who
turns to Him, and I find for Myself neither stillness, repose, nor
movement except at His command. None denies this but every treacherous,
transgressing sinner.

\subsection{284}\label{section-779}

\begin{otherlanguage}{persian}

بارى در جميع ظهورات اوّل معرض علماء عصر بوده‌‌اند چه كه اگر علماء انكار حقّ
اللّه نمينمودند احدى اعراض نمينمود و آنچه فساد در ارض شده از انكار علماء
بوده چنانچه تصريحاً ميفرمايند قوله عزّ ذكره و انّ اوّل من احتجب عن نقطة
البيان ادلّآء العلم عند انفسهم و لكنّهم عند اللّه لا يدركون من شیء و لا
يتعقلون اينست كه عند اللّه امثال آن نفوس از اهل علم محسوب نه و نفوسى كه
بايمان موفّق شده‌اند بر كرسى علم مستقرّ اگرچه بحرفى از علوم ظاهره مطّلع
نباشند

\end{otherlanguage}

Bārī dar jamīʿ-i ẓuhūrāt avval muʿriḍ ʿulamā-yi ʿaṣr būdih-and chih kih
agar ʿulamā inkār-i Ḥaqq Allāh namīnamūdand aḥadī iʿrāḍ namīnamūd va
ānchih fasād dar arḍ shudih az inkār-i ʿulamā būdih chunānchih taṣrīḥan
mīfarmāyand qawluhu ʿazza dhikruhu wa-inna awwal man iḥtajaba ʿan Nuqṭat
al-Bayān adillāʾ al-ʿilm ʿinda anfusihim wa-lākinnahum ʿinda Allāh lā
yudrikūna min shayʾ wa-lā yataʿaqqalūn īnast kih ʿinda Allāh amthāl-i ān
nufūs az ahl-i ʿilm maḥsūb nah va nufūsī kih bi-īmān muwaffaq shudih-and
bar kursī-yi ʿilm mustaqarr agarchih bi-ḥarfī az ʿulūm-i ẓāhirih
muṭṭaliʿ nabāshand.

In every Revelation, the first to turn away have been the religious
scholars of the age. For had the scholars not denied God, the Truth, no
one would have turned away, and whatever corruption has arisen on earth
has resulted from the scholars' denial. Thus He explicitly
states---exalted be His mention: ``The first to be veiled from the Point
of the Bayán are those who regard themselves as the guides of learning.
Yet, in God's sight, they perceive nothing and do not understand.''
Thus, in God's sight, such souls are not counted among the learned,
while those who have been enabled to believe are established upon the
seat of knowledge, even if they know not a single letter of outward
learning.

\subsection{285}\label{section-780}

\begin{otherlanguage}{persian}

اليوم علم و جهل و اعلى و ادنى و قرب و بعد و صادق و كاذب حىّ و ميّت بالغ و
رضيع عاقل و غافل بتصديق ظهور معلوم و ممتاز ميشوند من آمن به فهو من اعلى
العلمآء و اعزّ النّاس و ابصرهم و اتقاهم و اعقلهم و اعلاهم و اقربهم و
اصدقهم و ادقّهم و افقههم كذلك نزل بالحقّ من جبروت عزّ عظيم اين است كه نقطه
بيان در اينمقام ميفرمايند قوله عزّ ذكره فلتهدينّ اللّهمّ كلّ اولى محبّتى الى
دينك فانّ الأمر من عندك على حدّ سواء لو يومن بحجّتك و آياتك ادنى الخلق
ليكون اعلاه و ابعد الخلق ليكون اقربه

\end{otherlanguage}

Al-yawm ʿilm va jahl va aʿlā va adnā va qurb va buʿd va ṣādiq va kādhib
ḥayy va mayyit bāligh va raḍīʿ ʿāqil va ghāfil bi-taṣdīq-i ẓuhūr maʿlūm
va mumtāz mīshavand man āmana bihi fa-huwa min aʿlā al-ʿulamāʾ wa-aʿazz
al-nās wa-abṣarihim wa-atqāhum wa-aʿqalihim wa-aʿlāhum wa-aqrabihim
wa-aṣdaqihim wa-adaqqihim wa-afqahihim kadhālika nazala bi-al-ḥaqq min
jabarūt ʿizz ʿaẓīm īn ast kih Nuqṭih-yi Bayān dar īn maqām mīfarmāyand
qawluhu ʿazza dhikruhu fa-la-tahdiyanna Allāhumma kulla ulī maḥabbatī
ilā dīnika fa-inna al-amr min ʿindika ʿalā ḥadd sawāʾ law yuʾminu
bi-ḥujjatika wa-āyātika adnā al-khalq la-yakūnu aʿlāhu wa-abʿad al-khalq
la-yakūnu aqrabahu.

Today, knowledge and ignorance, high and low, nearness and remoteness,
the truthful and the false, the living and the dead, the mature and the
suckling, the discerning and the heedless are known and distinguished by
their acknowledgment of the Revelation. Whoever believes in Him is among
the highest of the learned, the most honored of people, the most
perceptive, the most God-fearing, the most discerning, the most exalted,
the nearest, the most truthful, the most exact in understanding, and the
most comprehending. Thus has it descended in truth from the dominion of
great glory. Hence the Point of the Bayán says in this
connection---exalted be His mention: ``Guide, O God, all who love Me to
Your religion, for the matter rests equally with You: if the lowest of
creation believes in Your Proof and Your verses, he will become its
highest, and the most remote of creation will become its nearest.''

\subsection{286}\label{section-781}

\begin{otherlanguage}{persian}

مقصود از اين بيانات آنكه جميع بدانند كه عند اللّه اهل علم نفوسى هستند كه
عالم باو شده‌اند و در كلمه بلى توقّف ننموده‌اند آن نفوسند اهل علم و ادلّاء
آن و اهل عصمت و معدن حكمت و منبع بيان و مطلع ذكر كلّ اسماء خير بآن نفوس
راجع ماداميكه در ظلّ اللّه مستقرّند علوم ظاهريّه در اينمقام معتبر نبوده و
نخواهد بود اينست كه آن نقطه وجود و طلعت مقصود در كلمات منزله مخاطباً باهل
علم وصاياى لايحصى فرموده‌اند و جميع را وصيّت فرموده كه در حين ظهور مغرور
بعلم و حكمت و بيان خود نشده كه شايد بعلم اللّه و ما عنده فايز گردند اينست
كه نقطه بيان ميفرمايد قوله عزّ ذكره لعلّ يوم ظهور حقّ كسى بافتخار تعلّم نزد
آن مبدأ علم متعلّم گردد و چه بسيار عظيم است اين امر زيرا كه علم او غير
بنفس او نيست و ميفرمايد عزّ ذكره هيچ لذّتى اعظم‌تر در امكان خلق نشده كه كسى
استماع نمايد آيات آن را و بفهمد مراد آن و لم و بم در حقّ كلمات آن نگويد و
مقايسه با كلام غير او نكند همين قسم كه كينونت او مظهر الوهيّت و ربوبيّت
است بر كلّ شیء همين قسم كلام او مظهر الو هيّت و ربوبيّت است بر كلّ كلامها كه
اگر آن انسانى بود متكلّم هرآينه ميگفت انّى انا اللّه لا اله الّا انا و انّ ما
دونى خلقى ان يا كلّ الحروف ايّاى فاتّقون انتهى

\end{otherlanguage}

Maqṣūd az īn bayānāt ānkih jamīʿ bidānand kih ʿinda Allāh ahl-i ʿilm
nufūsī hastand kih ʿālim bi-ū shudih-and va dar kalimih-yi balā tavaqquf
nanamūdih-and ān nufūsand ahl-i ʿilm va adillā-yi ān va ahl-i ʿiṣmat va
maʿdin-i ḥikmat va manbaʿ-i bayān va maṭlaʿ-i dhikr kull-i asmā-yi khayr
bi-ān nufūs rājiʿ mādāmīkih dar ẓill-i Allāh mustaqarrand ʿulūm-i
ẓāhiriyyih dar īn maqām muʿtabar nabūdih va nakhvāhad būd īnast kih ān
Nuqṭih-yi Vujūd va Ṭalʿat-i Maqṣūd dar kalimāt-i munazzalih mukhāṭaban
bi-ahl-i ʿilm vaṣāyā-yi lā yuḥṣā farmūdih-and va jamīʿ rā vaṣiyyat
farmūdih kih dar ḥīn-i ẓuhūr maghrūr bi-ʿilm va ḥikmat va bayān-i khud
nashudih kih shāyad bi-ʿilm-i Allāh va mā ʿindahu fāʾiz gardand īnast
kih Nuqṭih-yi Bayān mīfarmāyad qawluhu ʿazza dhikruhu laʿall yawm-i
ẓuhūr-i Ḥaqq kasī bi-iftikhār-i taʿallum nazd-i ān Mabdaʾ-i ʿIlm
mutaʿallim gardad va chih bisyār ʿaẓīm ast īn amr zīrā kih ʿilm-i ū
ghayr bi-nafs-i ū nīst va mīfarmāyad ʿazza dhikruhu hīch lidhdhatī
aʿẓam-tar dar imkān khalq nashudih kih kasī istimāʿ namāyad āyāt-i ān rā
va bifahmad murād-i ān va lim va bim dar ḥaqq-i kalimāt-i ān nagūyad va
muqāyisih bā kalām-i ghayr-i ū nakunad hamīn qism kih kaynūnat-i ū
maẓhar-i ulūhiyyat va rubūbiyyat ast bar kull-i shayʾ hamīn qism kalām-i
ū maẓhar-i ulūhiyyat va rubūbiyyat ast bar kull-i kalām-hā kih agar ān
insānī būd mutakallim har-āyinih mīguft innī anā Allāh lā ilāha illā anā
wa-inna mā dūnī khalqī an yā kull al-ḥurūf iyyāya fa-ittaqūn intahā.

The purpose of these statements is that all may know that, in God's
sight, the learned are those souls who have come to know Him and have
not hesitated to utter the word ``Yes.'' These souls are the possessors
of knowledge and its guides, the possessors of purity, the mine of
wisdom, the source of utterance, and the dawning place of remembrance.
All good names refer to these souls so long as they remain beneath the
shadow of God. Outward learning has never carried weight in this
station, nor will it ever do so. Thus that Point of Existence and
Countenance of the Desired One has given countless counsels in His
revealed words addressed to the learned. He has counseled them all not
to become conceited about their own knowledge, wisdom, and utterance at
the time of the Revelation, so that they may attain the knowledge of God
and what is with Him. Hence the Point of the Bayán says---exalted be His
mention: ``Perhaps, on the day of the Revelation of the True One,
someone may attain the honor of learning by becoming a student before
that Source of Knowledge. How exceedingly great is this matter, for His
knowledge is none other than His own Self!'' He also says---exalted be
His mention: ``No greater delight has been created in the realm of
possibility than for someone to hear His verses, understand their
meaning, ask no `why' or `wherefore' concerning His words, and not
compare them with the words of anyone else. Just as His being is the
manifestation of Divinity and Lordship over all things, so His words are
the manifestation of Divinity and Lordship over all words. Were His
words a human being endowed with speech, they would assuredly say: `I am
God; there is no god but Me, and all besides Me is My creation. O all
you letters, fear Me!'\,'' End of quotation.

\subsection{287}\label{section-782}

\begin{otherlanguage}{persian}

اى اهل بيان اگر چه بحدودى ظاهر شده‌ايد كه ابدا گمان تقديس و تنزيه نميرود
كه احدى بآن فايز شود و چنان مشاهده ميشود كه اگر صد‌هزار باب معارف الهيّه
مفتوح شود احدى گوش ندهد تا چه رسد بادراك بشنويد نصح الهى را و در كلّ ایّام
تفكّر در ظهور قبل و امورات واقعه در ایّامش نمائيد كه شايد آن امور محدثه از
نظر محو نشود و سبب استقامت گردد

\end{otherlanguage}

Ay ahl-i Bayān agar chih bi-ḥudūdī ẓāhir shudih-īd kih abadan gumān-i
taqdīs va tanzīh namīravad kih aḥadī bi-ān fāʾiz shavad va chunān
mushāhadih mīshavad kih agar ṣad-hizār bāb-i maʿārif-i ilāhiyyih maftūḥ
shavad aḥadī gūsh nadahad tā chih rasad bi-idrāk bishnavīd nuṣḥ-i ilāhī
rā va dar kull-i ayyām tafakkur dar ẓuhūr-i qabl va umūrāt-i vāqiʿih dar
ayyāmash namāʾīd kih shāyad ān umūr-i muḥdathih az naẓar maḥv nashavad
va sabab-i istiqāmat gardad.

O people of the Bayán! Although you have appeared in such a state of
limitation that there is no prospect of anyone attaining sanctity and
purity, and it seems that even if a hundred thousand doors of divine
knowledge were opened, no one would listen, much less understand, heed
the divine counsel. At all times, reflect upon the previous Revelation
and the events that took place in its days, so that those occurrences
may not fade from view and may become a cause of steadfastness.

\subsection{288}\label{section-783}

\begin{otherlanguage}{persian}

در اوّل ظهور كلّ علماء بر انكار قيام نمودند بنصّ خود نقطه بيان روح ما سواه
فداه يعنى مجتهدينى كه در كلّ اطراف بودند آن طلعت احديّه بهر يك لوحى باسمه
و رسمه نازل فرمود و در آن لوح ذكر اعراض او را بتفصيل فرموده‌اند و سواد آن
الواح الآن موجود و بعلماى نجف و كربلا واحداً بعد واحد مرقوم داشته‌اند و
هم‌‌چنين بعلماى فارس اسماء بعد اسم و هم‌‌چنين بعلماى ارض صاد و ارض قاف و
ساير بلاد بشأنى ابواب انصاف مسدود بود كه يك نفر آن قوم كه از معارف و
بعلم مشهور بآن شجره علم الهى و منبع حكمت صمدانى و مطلع بيان رحمانى موقن
نشده و فايز نگشته

\end{otherlanguage}

Dar avval-e ẓohūr koll-e ʿolamā bar enkār qiyām namūdand be-naṣṣ-e
khod-e Noqṭe-ye Bayān, rūḥ-e mā sevāh fadāh, yaʿnī mojtahedīnī ke dar
koll-e aṭrāf būdand, ān Ṭalʿat-e Aḥadiyye be-har yek lowḥī bi-smihi
wa-rasmihi nāzel farmūd va dar ān lowḥ dhekr-e eʿrāż-e ū rā be-tafṣīl
farmūde-and va savād-e ān alvāḥ al-ān mowjūd va be-ʿolamā-ye Najaf va
Karbalā wāḥidan baʿda wāḥid marqūm dāshte-and va ham-chenīn be-ʿolamā-ye
Fārs asmāʾ baʿda ism va ham-chenīn be-ʿolamā-ye arż-e Ṣād va arż-e Qāf
va sāyer-e belād. Be-shaʾnī abvāb-e enṣāf masdūd būd ke yek nafar ān
qowm ke az maʿāref va be-ʿelm mashhūr be-ān Shajare-ye ʿElm-e Elāhī va
Manbaʿ-e Ḥekmat-e Ṣamadānī va Maṭlaʿ-e Bayān-e Raḥmānī mūqen nashode va
fāyez nagashte.

At the beginning of the Revelation, all the religious scholars rose in
denial, according to the explicit statement of the Point of the Bayán
Himself---may the spirit of all besides Him be sacrificed for Him. That
is, to each of the mujtahids throughout the regions, that Countenance of
Oneness revealed a tablet bearing his name and designation, and in that
tablet described his rejection in detail. Copies of those tablets are
now extant. He wrote to the scholars of Najaf and Karbalá, one after
another, and likewise to the scholars of Fárs, name after name, and
likewise to the scholars of the Land of Ṣád, the Land of Qáf, and other
lands. The doors of fairness were so firmly closed that not a single one
of those people renowned for learning and knowledge came to believe in
that Tree of divine knowledge, that Source of everlasting wisdom, that
Dawning Place of the All-Merciful's utterance, or attained His presence.

\subsection{289}\label{section-784}

\begin{otherlanguage}{persian}

و نفسيكه قبل از كلّ نفوس بشرافت كبرى فايز و باب رحمت و باب علم و باب
هدايت و باب فضل شد از براى كلّ موجودات جناب باب بهآءاللّه عليه بوده و
البتّه تفصيل آن ذات مقدّس را هر نفسى شنيده كه مابين ناس اشتهار نداشته و
صاحب منبر و محراب هم نبوده‌اند يكى از تلامذه جناب حاجى سيّد كاظم عليه رحمة
اللّه و فيوضاته بودند و اكثرى از علماى ظاهره آن نفس مقدسّه را تكفير
مينمودند و در تلاميذ سيّد مرحوم هم بعضى از نفوس بوده كه از ايشان معروف‌تر
بوده‌اند چنانچه يومى از جناب ملّا مهدى خوئى كه يكى از اصحاب جناب سيّد بود
از علم جناب باب استفسار شده او چند نفوس متعدّده را اسم برده و جميع را در
علم و فضل ظاهره ترجيح ميداد بر باب اعظم و كلمه اكرم و

\end{otherlanguage}

Va nafsī ke qabl az koll-e nofūs be-sharāfat-e kobrā fāyez va bāb-e
raḥmat va bāb-e ʿelm va bāb-e hedāyat va bāb-e fażl shod az barā-ye
koll-e mowjūdāt Janāb-e Bāb, Bahāʾuʾllāh ʿalayh, būde va albatte
tafṣīl-e ān dhāt-e moqaddas rā har nafsī shenīde ke mābeyn-e nās
eshtehār nadāshte va ṣāḥeb-e menbar va meḥrāb ham nabūde-and. Yekī az
talāmedhe-ye Janāb-e Ḥājī Seyyed Kāẓem, ʿalayhi raḥmatuʾllāh
wa-fuyūḍātuh, būdand va aktharī az ʿolamā-ye ẓāhere ān nafs-e moqaddase
rā takfīr mīnamūdand va dar talāmīdh-e Seyyed-e marḥūm ham baʿżī az
nofūs būde ke az īshān maʿrūf-tar būde-and. Chenānche yowmī az Janāb-e
Mollā Mahdī-ye Khūʾī ke yekī az aṣḥāb-e Janāb-e Seyyed būd az ʿelm-e
Janāb-e Bāb estefsār shode, ū chand nofūs-e motaʿaddede rā esm borde va
jamīʿ rā dar ʿelm va fażl-e ẓāhere tarjīḥ mīdād bar Bāb-e Aʿẓam va
Kaleme-ye Akram va

The soul who, before all others, attained the greatest honor and became
the gate of mercy, the gate of knowledge, the gate of guidance, and the
gate of grace for all beings was the honored Báb---may the glory of God
rest upon him. Surely everyone has heard the account of that holy
person: he was not widely known among the people, nor did he occupy a
pulpit or prayer niche. He was one of the students of the honored Ḥájí
Siyyid Káẓim---may God's mercy and blessings rest upon him---and most of
the outwardly learned declared that holy soul an unbeliever. Among the
students of the departed Siyyid, too, there were some who were better
known than he. Thus, one day, when the honored Mullá Mahdí Khú'í, one of
the Siyyid's companions, was asked about the learning of the honored
Báb, he named several individuals and ranked them all above that Most
Great Gate and Most Noble Word in outward learning and accomplishments,
and

\subsection{290}\label{section-785}

\begin{otherlanguage}{persian}

نفوس مشهوره معروفه هيچ كدام موقن نشدند بلكه بر اعراض و اعتراض قيام
نمودند از جمله ملّا حسن گوهر و ميرزا محيط و ميرزا كريم خان و همچنين
سايرين كه در آذربايجان و اطراف معروف بودند بارى جميع علماى معروفين چه از
فئه شيخيّه و چه از فئه ظاهره در امراللّه توقّف نمودند بلكه معرض و منكر و
نظر باعراض سلب علم از آن نفوس غافله فرموده‌اند

\end{otherlanguage}

Nofūs-e mashhūre-ye maʿrūfe hīch kodām mūqen nashodand balke bar eʿrāż
va eʿterāż qiyām namūdand, az jomle Mollā Ḥasan-e Gowhar va Mīrzā Moḥīṭ
va Mīrzā Karīm Khān va hamchenīn sāyerīn ke dar Ādharbāyjān va aṭrāf
maʿrūf būdand. Bārī jamīʿ-e ʿolamā-ye maʿrūfīn che az feʾe-ye Sheykhiyye
va che az feʾe-ye ẓāhere dar Amruʾllāh tavaqqof namūdand balke moʿreż va
monker va naẓar be-eʿrāż salb-e ʿelm az ān nofūs-e ghāfele farmūde-and.

None of the prominent, well-known figures came to believe; instead, they
rose in rejection and opposition. Among them were Mullá Ḥasan Gawhar,
Mírzá Muḥíṭ, and Mírzá Karím Khán, as well as others who were renowned
in Azerbaijan and the surrounding regions. In short, all the renowned
scholars, whether of the Shaykhí school or of the outwardly learned,
hesitated concerning the Cause of God; indeed, they turned away and
denied it. Because of their rejection, He denied those heedless souls
the designation of knowledge.

\subsection{291}\label{section-786}

\begin{otherlanguage}{persian}

و لكن جناب باب اعظم يعنى اوّل من آمن و سبعة عشر انفس كه با ايشان بودند
چون بشرافت ايمان فايز شدند كلّ عند اللّه از ادلّاء علم مذكور چنانچه
ميفرمايد اين نفوس در جميع اسماء حسنى ابواب رحمتند بر كلّ ناس و در بيان
پارسى بتفصيل مذكور و مسطور ملاحظه نمايند منتهى جدّ و جهد را نموده كه در
يوم ظهور از هيچ امرى محتجب نمانيد و ناظر بهيچ شیء عمّا خلق و يخلق نشويد
چه كه هر نفسى در آن يوم در اقلّ من آن بشیء غير حقّ ناظر شود بقدر همان از
شريعه قرب بعيد خواهد بود

\end{otherlanguage}

Va lāken Janāb-e Bāb-e Aʿẓam yaʿnī avval-e man āmana va sabʿata ʿashara
anfus ke bā īshān būdand chūn be-sharāfat-e īmān fāyez shodand koll
ʿindaʾllāh az adellā-ye ʿelm madhkūr. Chenānche mīfarmāyad īn nofūs dar
jamīʿ-e asmā-ye ḥosnā abvāb-e raḥmatand bar koll-e nās va dar Bayān-e
Pārsī be-tafṣīl madhkūr va masṭūr. Molāḥeẓe namāyand. Montahā-ye jedd va
jahd rā namūde ke dar yowm-e ẓohūr az hīch amrī moḥtajeb namānīd va
nāẓer be-hīch shayʾ ʿammā kholiqa va yokhlaq nashavīd che ke har nafsī
dar ān yowm dar aqall-e men ān be-shayʾ-e gheyr-e Ḥaqq nāẓer shavad
be-qadr-e hamān az sharīʿe-ye qorb baʿīd khāhad būd.

But the honored Most Great Gate---that is, the first to believe---and
the seventeen souls who were with him, having attained the honor of
faith, are all accounted in God's sight as guides to knowledge. Thus He
says: ``These souls, in all the most beautiful names, are gates of mercy
to all people.'' This is mentioned and recorded in detail in the Persian
Bayán; let them consult it. He exerted the utmost effort so that, on the
day of the Revelation, you would not remain veiled by anything, nor fix
your gaze on anything that has been or will be created. For if, on that
day, any soul looks for less than an instant toward anything other than
the True One, to that same extent he will be far from the waters of
nearness.

\subsection{292}\label{section-787}

\begin{otherlanguage}{persian}

قلب را از كلّ مطهّر داريد و نظر را از علامات و اشارات و قرب و بعد كه بعقل
خود ادراك نموده‌ايد بر‌داريد چه كه احدى بر مشيّت مكنونه الهيّه و اراده
مستوره‌اش اطّلاع نداشته و ندارد و در كلّ حين بشأن بديع ظاهر شده و خواهد شد
بسا است كه بشأنى ظاهر ميشود كه ابدا احدى از قبل تعقّل ننموده پناه بخدا
برده و در ليل و نهار مسئلت نمائيد كه بامرى ظاهر شود كه ادراك آن از براى
خلق سهل باشد كه لعلّ از فيوضات آن بحر اعظم و اشراق انوار شمس معانى محروم
نمانند

\end{otherlanguage}

Qalb rā az koll moṭahhar dārīd va naẓar rā az ʿalāmāt va eshārāt va qorb
va boʿd ke be-ʿaql-e khod edrāk namūde-īd bar-dārīd che ke aḥadī bar
mashiyyat-e maknūne-ye Elāhiyye va erāde-ye mastūre-ash eṭṭelāʿ nadāshte
va nadārad va dar koll-e ḥīn be-shaʾn-e badīʿ ẓāher shode va khāhad
shod. Basā ast ke be-shaʾnī ẓāher mīshavad ke abadan aḥadī az qabl
taʿaqqol nanamūde. Panāh be-Khodā borde va dar leyl va nahār masʾalat
namāyīd ke be-amrī ẓāher shavad ke edrāk-e ān az barā-ye khalq sahl
bāshad ke laʿall az foyūżāt-e ān Baḥr-e Aʿẓam va eshrāq-e anvār-e
Shams-e Maʿānī maḥrūm namānand.

Keep your heart pure of all things, and turn your gaze away from the
signs and allusions, the nearness and remoteness, that you have
apprehended through your own reason. For no one has ever known, or now
knows, God's hidden will and His concealed purpose. At every time He has
appeared, and will appear, in a wondrous manner. He may well appear in a
manner that no one has ever previously conceived. Seek refuge in God and
beseech Him night and day to appear in a way that His creatures can
readily comprehend, so that perhaps they may not be deprived of the
outpourings of that Most Great Ocean and the radiance of the lights of
the Sun of inner meanings.

\subsection{293}\label{section-788}

\begin{otherlanguage}{persian}

ظهور قبل روحى و نفسى له الفدآء در اين امر بسيار نصيحت فرموده‌اند و
ببيانات لايحصى ذكر نموده‌اند على شأن لن يخطر ببال احد امر الّا و انّه رقم
من قلمه الأعلی مثلا اگر نفسى ببيان ناظر باشد ميفرمايد ایّاک ان تحتجب به
عن منزله و اگر كسى بحروف متمسّك باشد ميفرمايد ایّاک ان تحتجب بالحروف عن
موجدهم و خالقهم و اگر كسى بعلم مفتخر باشد ميفرمايد لم يصدق اسم العلم على
احد الّا بأن يكون عالماً بنفس الظّهور و اگر نفسى بعلامات محدود باشد
ميفرمايد كه ایّاکم ان تحتجبوا بها چه كه احدى بر كيفيّت ظهور آن ذات قدم
مطّلع نه يظهر نفسه كيف یشآء و اگر كسى بتكميل ناظر باشد ميفرمايد تصديق كلّ
شیء منوط بتصديق آن منبع صدق بوده و خواهد بود اگر بفرمايد كلّ كاملند ليس
لأحد ان يقول لم و بم و اگر بفرمايد كلّ ناقصند ليس لأحد ان يعترض عليه و
اگر نفسى بوقت محدود باشد چنانچه اين ایّام كلّ بآن محتجب ميفرمايد هر وقت
ظاهر شود بايد كلّ تصديق بشجره حقيقت نمايند و شكر الهى بجا آورند و اگر
نفسى برؤسا و مراياى عصر ناظر باشد مخاطباً ايّاهم ميفرمايد كه شما بقول من
حجّت شده‌ايد زينهار كه از نفس من در ظهور آخر محتجب شويد

\end{otherlanguage}

Ẓohūr-e qabl, rūḥī wa-nafsī lahu al-fidāʾ, dar īn amr besyār naṣīḥat
farmūde-and va be-bayānāt-e lā yoḥṣā dhekr namūde-and, ʿalā shaʾnin lan
yakhṭura bi-bāli aḥadin amrun illā wa-innahu ruqima min qalamihi
al-aʿlā. Mathalan agar nafsī be-Bayān nāẓer bāshad mīfarmāyad iyyāka an
taḥtajiba bihi ʿan munzilihi va agar kasī be-ḥorūf motamassek bāshad
mīfarmāyad iyyāka an taḥtajiba bi-l-ḥurūfi ʿan mūjidihim wa-khāliqihim
va agar kasī be-ʿelm moftakher bāshad mīfarmāyad lam yaṣduq ismu
al-ʿilmi ʿalā aḥadin illā bi-an yakūna ʿāliman bi-nafsi al-ẓuhūr va agar
nafsī be-ʿalāmāt maḥdūd bāshad mīfarmāyad ke iyyākum an taḥtajibū bihā
che ke aḥadī bar keyfiyyat-e ẓohūr-e ān Dhāt-e Qedam moṭṭaleʿ na,
yuẓhiru nafsahu kayfa yashāʾ. Va agar kasī be-takmīl nāẓer bāshad
mīfarmāyad taṣdīq-e koll-e shayʾ manūṭ be-taṣdīq-e ān Manbaʿ-e Ṣedq būde
va khāhad būd. Agar befarmāyad koll kāmeland laysa li-aḥadin an yaqūla
lima wa-bima va agar befarmāyad koll nāqeṣand laysa li-aḥadin an
yaʿtariḍa ʿalayh. Va agar nafsī be-vaqt maḥdūd bāshad chenānche īn ayyām
koll be-ān moḥtajeb mīfarmāyad har vaqt ẓāher shavad bāyad koll taṣdīq
be-Shajare-ye Ḥaqīqat namāyand va shokr-e Elāhī be-jā āvarand va agar
nafsī be-roʾasā va marāyā-ye ʿaṣr nāẓer bāshad mukhāṭiban iyyāhum
mīfarmāyad ke shomā be-qowl-e man ḥojjat shode-īd zīnhār ke az nafs-e
man dar ẓohūr-e ākhar moḥtajeb shavīd.

The previous Manifestation---may my spirit and my self be sacrificed for
Him---gave abundant counsel on this matter and spoke of it in countless
utterances, so that no matter could occur to anyone's mind without
having been recorded by His Most Exalted Pen. For example, if someone
looks to the Bayán, He says: ``Beware lest you be veiled by it from the
One who revealed it.'' If someone clings to the Letters, He says:
``Beware lest you be veiled by the Letters from their Originator and
Creator.'' If someone prides himself on knowledge, He says: ``The
designation of knowledge does not truly apply to anyone unless he knows
the Manifestation Himself.'' If a soul is confined by signs, He says:
``Beware lest you be veiled by them,'' for no one knows how that Eternal
Essence will appear: ``He manifests Himself however He wills.'' If
someone looks to the attainment of perfection, He says that the
validation of all things has always depended, and will always depend, on
the validation of that Source of truth. If He says that all are perfect,
``No one may say `why' or `wherefore.'\,'' If He says that all are
imperfect, ``No one may object to Him.'' If a soul is confined by a
fixed time, as everyone is veiled by it in these days, He says that
whenever He appears, all must acknowledge the Tree of Truth and give
thanks to God. If a soul looks to the leaders and mirrors of the age, He
addresses them, saying: ``You have become proofs through My word. Beware
lest you be veiled from My own Self in the other Revelation.''

\subsection{294}\label{section-789}

\begin{otherlanguage}{persian}

و لذلك الرّبّ يليق تسبيحى و تسبيح كلّ شیء و تقديسى و تقديس كلّ شیء و ذكرى و
ذكر كلّ شیء لأنّه ما قصّر فى نصح عباده فى هذا النّبأ الأعظم العظيم و اجاب من
قبل كلّ اعتراضات الّتى كانت فى انفس عباده لئلّا يرد فى ظهوره مرّة اخرى ما
يحزن به نفسه و يكدّر به فؤاده حال ملاحظه امراللّه نمائيد و ملاحظه شأن عباد
كه چه قدر غافلند و آن نفوسيكه بقول او موجود شده عوض آنكه شكرش نمايند و
بحمدش حامد و بثنائش مثنى و بذكرش ذاكر و بنعمتش شاكر كلّ باسياف كين ظاهر
شده‌اند نسئل اللّه بأن يحفظ هذا المظلوم من هؤلآء المغلّين و المشركين و ينصر
امره و يظهر برهانه و يعلو حجّته بين السّموات و الأرضين روحى له الفدآء و
نفسى له الفدآء لأنّه ما بقى من نصح الّا و قد نصح به عباده

\end{otherlanguage}

Wa-li-dhālika al-Rabbi yalīqu tasbīḥī wa-tasbīḥu kulli shayʾin
wa-taqdīsī wa-taqdīsu kulli shayʾin wa-dhikrī wa-dhikru kulli shayʾin
li-annahu mā qaṣṣara fī nuṣḥi ʿibādihi fī hādhā al-nabaʾi al-aʿẓami
al-ʿaẓīm wa-ajāba min qablu kulla iʿtirāḍāti allatī kānat fī anfusi
ʿibādihi li-allā yarida fī ẓuhūrihi marratan ukhrā mā yaḥzanu bihi
nafsuhu wa-yukaddaru bihi fuʾāduh. Ḥāl molāḥeẓe-ye Amruʾllāh namāyīd va
molāḥeẓe-ye shaʾn-e ʿebād ke che qadr ghāfeland va ān nofūsī ke
be-qowl-e ū mowjūd shode ʿevaż-e ānke shokrash namāyand va be-ḥamdash
ḥāmed va be-thanāyash mothnī va be-dhekrash dhāker va be-neʿmatash
shāker koll be-asyāf-e kīn ẓāher shode-and. Nasʾalu Allāha bi-an yaḥfaẓa
hādhā al-Maẓlūma min hāʾulāʾi al-mughillīna wa-l-mushrikīna wa-yanṣura
amrahu wa-yuẓhira burhānahu wa-yaʿluwa ḥujjatahu bayna al-samāwāti
wa-l-araḍīn. Rūḥī lahu al-fidāʾ wa-nafsī lahu al-fidāʾ li-annahu mā
baqiya min nuṣḥin illā wa-qad naṣaḥa bihi ʿibādah.

To that Lord belong my glorification and the glorification of all
things, my sanctification and the sanctification of all things, my
remembrance and the remembrance of all things, for He neglected nothing
in counseling His servants concerning this most great and mighty
Announcement. He answered beforehand all the objections within His
servants' souls, so that when He appeared once more, nothing would
befall Him that would grieve His soul or trouble His heart. Now consider
the Cause of God and consider the condition of His servants---how
heedless they are! Those souls who were brought into being by His word,
instead of thanking Him, praising Him, extolling Him, remembering Him,
and showing gratitude for His bounty, have all appeared with swords of
malice. We ask God to protect this Wronged One from these malicious
people and those who associate partners with God, to aid His Cause, to
make His proof manifest, and to exalt His testimony throughout the
heavens and the earths. May my spirit be sacrificed for Him, and may my
self be sacrificed for Him, for no counsel remained that He had not
given His servants.

\subsection{295}\label{section-790}

\begin{otherlanguage}{persian}

مخصوص چون از علماء در ظهور قبل غرور و استكبار مشاهده فرموده‌اند اين فئه
را بسيار نصيحت نموده و بمواعظ حسنه و بيانات شافيه ايشان را متذكّر داشته
كه مباد بسنن قبل عمل نمايند و در حين مظهر ظهور از حقّ غافل شوند چه كه
باعراض آن نفوس جميع ناس معرض الّا من شآء ربّك چنانچه ميفرمايند قوله عزّ
ذكره اينست كه در نزد هر ظهورى اهل آن ظهور بعلماى آن ظهور محتجب ميشوند و
حال آنكه از وراى امر غافلند كه بيك قول لاى او كلّ اينها غير صادق ميگردند
مراقب باش ظهور حقّ را كه در يوم ظهور بتصديق اهل بيان تصديق او نكنى كه كلّ
تصديقهاى ايشان در نزد او بيك بلى تصديق صرف ميشود و بيك لا دون آن انتهى

\end{otherlanguage}

Makhṣūṣ chūn az ʿolamā dar ẓohūr-e qabl ghorūr va estekbār moshāhede
farmūde-and īn feʾe rā besyār naṣīḥat namūde va be-mavāʿeẓ-e ḥasane va
bayānāt-e shāfiye īshān rā motadhakker dāshte ke mabād be-sonan-e qabl
ʿamal namāyand va dar ḥīn-e maẓhar-e ẓohūr az Ḥaqq ghāfel shavand che ke
be-eʿrāż-e ān nofūs jamīʿ-e nās moʿreż illā man shāʾa rabbuk. Chenānche
mīfarmāyand qawluhu ʿazza dhikruh: īnast ke dar nazd-e har ẓohūrī ahl-e
ān ẓohūr be-ʿolamā-ye ān ẓohūr moḥtajeb mīshavand va ḥāl ānke az varā-ye
amr ghāfeland ke be-yek qowl-e lā-ye ū koll-e īnhā gheyr-e ṣādeq
mīgardand. Morāqeb bāsh ẓohūr-e Ḥaqq rā ke dar yowm-e ẓohūr be-taṣdīq-e
ahl-e Bayān taṣdīq-e ū nakonī ke koll-e taṣdīq-hā-ye īshān dar nazd-e ū
be-yek balā taṣdīq-e ṣerf mīshavad va be-yek lā dūn-e ān. Intahā.

In particular, since He had witnessed conceit and arrogance among the
scholars in the previous Revelation, He counseled this group at length,
reminding them through gracious admonitions and healing utterances not
to follow the ways of the past or become heedless of the True One at the
time of the Manifestation's appearance. For when those souls turn away,
all people turn away, except those whom your Lord wills. Thus He
says---His words, exalted be His mention: ``It is thus that, at every
Revelation, the people of that Revelation are veiled by its scholars,
while they remain heedless of what lies beyond the matter: by a single
`No' from Him, all these become untrue. Watch for the Revelation of the
True One, lest on the day of the Revelation you acknowledge Him on the
strength of the acknowledgment of the people of the Bayán. For in His
presence all their affirmations become pure affirmation through a single
`Yes,' and through a single `No,' the opposite.'' End of quotation.

\subsection{296}\label{section-791}

\begin{otherlanguage}{arabic}

و هم‌‌چنين در اينمقام مخاطباً لعلماء البيان ميفرمايند قوله جلّ كبريائه ان يا
علماء البيان فلترحمنّ على انفسكم ثمّ ترحمون بأن لا تحتجبنّ عن الّذى يظهر
بالحقّ بشئون دينكم فانّ كلّ دينكم عنده بمثل يوم قد خلق اللّه ذلك الدّين انتم
الى نقطة البدء تنظرون فانكم انتم مثل كلّ اولى الظّهورات مبتلون و بما ليل
راقدون و انتم عنداللّه مؤمنون و فى نهاركم مثلكم كمثل العلمآء من قبلكم لن
يحكم اللّه عليكم بشىء فلتتراقبن انفسكم ثمّ لتراقبون

\end{otherlanguage}

Va ham-chenīn dar īn maqām mukhāṭiban li-ʿulamāʾi al-Bayān mīfarmāyand
qawluhu jalla kibriyāʾuh: in yā ʿulamāʾa al-Bayāni fa-la-tarḥamunna ʿalā
anfusikum thumma tarḥamūna bi-an lā taḥtajibanna ʿani alladhī yaẓharu
bi-l-ḥaqqi bi-shuʾūni dīnikum fa-inna kulla dīnikum ʿindahu bi-mithli
yawmin qad khalaqa Allāhu dhālika al-dīn. Antum ilā nuqṭati al-badʾi
tanẓurūna fa-innakum antum mithla kulli ulī al-ẓuhūrāti mubtalūna
wa-bi-mā laylin rāqidūna wa-antum ʿinda Allāhi muʾminūna wa-fī nahārikum
mithlukum ka-mithli al-ʿulamāʾi min qablikum lan yaḥkuma Allāhu ʿalaykum
bi-shayʾin fa-la-tatarāqabinna anfusakum thumma la-turāqibūn.

Likewise, in this connection He addresses the scholars of the Bayán,
saying---His words, glorious be His majesty: ``O scholars of the Bayán!
Have mercy on yourselves, and show mercy by not allowing the concerns of
your religion to veil you from the One who appears with the truth. For
your entire religion is, in His presence, as it was on the day God
created that religion. Look to the Point of the Beginning, for you are
being tested like the people of all Revelations. In a night you lie
asleep, while in God's sight you are believers; and in your daytime you
are like the scholars before you. God will pass no judgment in your
favor. Watch over yourselves, then remain watchful.''

\subsection{297}\label{section-792}

\begin{otherlanguage}{persian}

اى ملأ بيان رحم بر خود و امراللّه نموده سبب اختلاف مشويد و ارتكاب منمائيد
امرى را كه بآن كلمه جامعه متّحده تفريق شود در امور وارده خالصاً لوجه اللّه
تفكّر نمائيد و اللّه الّذى لا اله الّا هو كه حقّ در هر وقتى از اوقات و عصرى
از اعصار كه ظاهر شده بشأنى ظاهر كه بر متبصّرين و منصفين مجال توقّف و شبهه
و ريب نبوده مگر نفوسيكه خود برخود مشتبه نموده‌اند فلتّتقنّ اللّه يا ملأ
البيان و لتعتصمنّ باللّه يا ملأ البيان و لترحمن على انفسكم يا ملأ البيان و
لاتعترضنّ على الّذى آمنتم به فى ظهور قبله و به ثبت ايمان كلّ مؤمن و ايقان
كلّ موقن و توحيد كلّ موحّد و علم كلّ عالم و ذكر كلّ ذاكر هذا خير النّصح من لدى
العبد لكم و عليكم ان انتم من المنصفين

\end{otherlanguage}

Ey malaʾ-e Bayān raḥm bar khod va Amruʾllāh namūde sabab-e ekhtelāf
mashavīd va ertekāb nanamāyīd amrī rā ke be-ān kaleme-ye jāmeʿe-ye
mottaḥede tafrīq shavad. Dar omūr-e vārede khāliṣan li-wajhi Allāh
tafakkor namāyīd. Wa-Allāhi alladhī lā ilāha illā huwa ke Ḥaqq dar har
vaqtī az owqāt va ʿaṣrī az aʿṣār ke ẓāher shode be-shaʾnī ẓāher ke bar
motabaṣṣerīn va monṣefīn majāl-e tavaqqof va shobhe va reyb nabūde magar
nofūsī ke khod bar khod moshtabeh namūde-and. Fa-la-ttaqinna Allāha yā
malaʾa al-Bayāni wa-la-taʿtaṣimunna bi-Allāhi yā malaʾa al-Bayāni
wa-la-tarḥamunna ʿalā anfusikum yā malaʾa al-Bayāni wa-lā taʿtariḍanna
ʿalā alladhī āmantum bihi fī ẓuhūri qablihi wa-bihi thabata īmānu kulli
muʾminin wa-īqānu kulli mūqinin wa-tawḥīdu kulli muwaḥḥidin wa-ʿilmu
kulli ʿālimin wa-dhikru kulli dhākir. Hādhā khayru al-nuṣḥi min ladā
al-ʿabdi lakum wa-ʿalaykum in antum mina al-munṣifīn.

O people of the Bayán! Have mercy on yourselves and on the Cause of God,
and do not become a cause of discord. Do not commit any act that would
divide the united body gathered around one word. Reflect on what has
occurred solely for the sake of God. By God, besides whom there is no
other god, whenever and in whatever age the True One has appeared, He
has appeared in such a manner that those endowed with insight and
fairness had no room for hesitation, misgiving, or doubt, except for
souls who have deliberately confused themselves. Fear God, O people of
the Bayán! Hold fast to God, O people of the Bayán! Have mercy on
yourselves, O people of the Bayán! Do not oppose the One in whom you
believed in His previous Revelation, through whom the faith of every
believer, the certitude of every assured soul, the affirmation of divine
unity by everyone who professes it, the knowledge of every learned
person, and the remembrance of everyone who remembers God have been
established. This is the best counsel from this servant for you, and it
is incumbent upon you, if you are fair-minded.

\subsection{298}\label{section-793}

\begin{otherlanguage}{persian}

و ديگر نوشته بودى و اگر مدّعى آن است كه من خود آن نقطه بيان هستم خالى از
اين نيست يا آنكه مرادش آن است كه او هستم بعينه و بشخصه كه واضح البطلان
است چه كه تكرّر تجلّيات جايز نيست و بالبديهه باطل و اگر تجدّد امثال است كه
واضح تر است بطلانش بر علاوه بايد در اينصورت هم بعينه بهمان احكام سابق
بلاتفاوت امر نمايد پس كجا تغيير جايز است انتهى

\end{otherlanguage}

Va dīgar neveshte būdī va agar moddaʿī ān ast ke man khod ān Noqṭe-ye
Bayān hastam khālī az īn nīst yā ānke morādash ān ast ke ū hastam
bi-ʿaynihi va bi-shakhṣihi ke vāżeḥ al-boṭlān ast che ke takarror-e
tajalliyāt jāyez nīst va bi-l-badīhe bāṭel va agar tajaddod-e amthāl ast
ke vāżeḥ-tar ast boṭlānash. Bar ʿalāve bāyad dar īn ṣūrat ham bi-ʿaynihi
be-hamān aḥkām-e sābeq belā tafāvot amr namāyad pas kojā taghyīr jāyez
ast. Intahā.

You also wrote: ``If he claims, `I myself am that Point of the Bayán,'
it must mean one of two things. Either he means, `I am that very same
individual in person,' which is clearly false, since the repetition of
divine manifestations is inadmissible and self-evidently false; or he
means the renewal of like beings, whose falsity is even clearer.
Moreover, even in this case he must command precisely the same former
laws, without any difference. How, then, can change be permissible?''
End of quotation.

\subsection{299}\label{section-794}

\begin{otherlanguage}{persian}

از قراريكه معلوم ميشود هنوز معنى تجلّى را هم ادراك ننموده چه كه محلّش را
ندانسته اينجا محلّ مجلّى است نه تجلّى اگر بعين بصيرت ملاحظه نمائى بمبدأ
تجلّيات لانهايه پى‌برى و اين جواب از قبل ذكر شده كه بعينه ميفرمايند آنچه
را نقطه اولى فرموده ذكر تجلّى و عدم تكرار تجلّى را بگذار و مجلّى را ادراك
كن و از دونش منقطع شو

\end{otherlanguage}

Az qarārī ke maʿlūm mī-shavad hanūz maʿnā-ye tajallī rā ham edrāk
na-namūde, che ke maḥallash rā na-dāneste. Īnjā maḥall-e mojallī ast na
tajallī. Agar be-ʿayn-e baṣīrat molāḥeẓe namāʾī be-mabdaʾ-e tajalliyāt-e
lā-nehāye pey barī. Va īn javāb az qabl dhekr shode ke be-ʿayneh
mī-farmāyand ānche rā Noqṭe-ye Ūlā farmūde. Dhekr-e tajallī va ʿadam-e
tekrār-e tajallī rā be-godhār va mojallī rā edrāk kon va az dūnash
monqaṭeʿ show.

From what appears, you have not yet understood even the meaning of
effulgence, for you have not recognized its proper place. What is at
issue here is the Source of effulgence, not effulgence itself. If you
look with the eye of insight, you will recognize the Origin of infinite
effulgences. This answer has already been given: He says precisely what
the Primal Point said. Set aside discussion of effulgence and its
non-recurrence, recognize the Source of effulgence, and detach yourself
from all else.

\subsection{300}\label{section-795}

\begin{otherlanguage}{persian}

چه ميگوئيد در نقطه اولى روح ما سواه فداه كه ميفرمايند نقطه بيان بعينه
همان نقطه فرقان است از قبل و هم‌‌چنين در جميع بيان كه ميفرمايند من ظهر عين
من يظهر و من يظهر عين من ظهر و هم‌‌چنين ميفرمايند انّنى انا ايّاه و انّه هو
ايّاى

\end{otherlanguage}

Che mī-gūʾīd dar Noqṭe-ye Ūlā, rūḥu mā siwāhu fidāhu, ke mī-farmāyand
Noqṭe-ye Bayān be-ʿayneh hamān Noqṭe-ye Forqān ast az qabl, va
ham-chenīn dar jamīʿ-e Bayān ke mī-farmāyand man ẓahara ʿaynu man
yaẓharu va man yaẓharu ʿaynu man ẓahara, va ham-chenīn mī-farmāyand
innanī anā iyyāhu wa-innahu huwa iyyāya.

What do you say about the Primal Point---may the souls of all else be
sacrificed for Him---who says that the Point of the Bayán is precisely
that same Point of the Qur'an of former times? Likewise, throughout the
Bayán He says, ``He who has appeared is the very same as He who will
appear, and He who will appear is the very same as He who has
appeared.'' He also says, ``I am indeed He, and He is indeed I.''

\subsection{301}\label{section-796}

\begin{otherlanguage}{persian}

از اين گذشته مظاهر احديّه از جهتى متّحد و از جهتى بظهورات مختلفه ظاهر از
آنجهت كه كلّ از مشرق امر مشرق و من عنداللّه تكلّم نموده‌اند و كلّ بتبليغ ناس
مأمور گشته‌اند كلّ واحد بوده اينست كه ميفرمايد لانفرّق بين احد منهم و در
مقام ديگر كه مقام شرايع و اسم و جهات ملكيّه است هر كدام باسمى و شريعتى و
هيكل و هيئتى ظاهر شده‌اند و اگر بديده بصيرت ملاحظه نمائى فواللّه الّذى لا
اله الّا هو كه كلّ را بحقيقت واحده و ظهور واحد و حركت واحده و فعل واحد و
عمل واحد و شريعت واحده و اسم واحد و رسم واحد مشاهده نمائى عجب است كه
باين بيان نرسيده اطفال اين ظهور بر اين عارف شده‌اند چه كه در الواح پارسيّه
كه در سنين توقّف در عراق از سماء عزّ احديّه نازل جميع اين مراتب بتفصيل در
آن الواح مذكور آمده فطوبی‌ للفائزين و العارفين

\end{otherlanguage}

Az īn godhashte maẓāher-e aḥadiyye az jahatī mottaḥed va az jahatī
be-ẓohūrāt-e mokhtalefe ẓāher. Az ān jahat ke koll az mashreq-e amr
moshreq va min ʿindi Allāh takallom namūde-and va koll be-tablīgh-e nās
maʾmūr gashte-and, koll vāḥed būde. Īnast ke mī-farmāyad lā nufarriqu
bayna aḥadin minhum. Va dar maqām-e dīgar ke maqām-e sharāyeʿ va esm va
jahāt-e molkiyye ast, har kodām be-esmī va sharīʿatī va heykal va
heyʾatī ẓāher shode-and. Va agar be-dīde-ye baṣīrat molāḥeẓe namāʾī,
fa-wa-Allāhi alladhī lā ilāha illā huwa, ke koll rā be-ḥaqīqat-e vāḥede
va ẓohūr-e vāḥed va ḥarakat-e vāḥede va feʿl-e vāḥed va ʿamal-e vāḥed va
sharīʿat-e vāḥede va esm-e vāḥed va rasm-e vāḥed moshāhede namāʾī. ʿAjab
ast ke be-īn bayān na-rasīde. Aṭfāl-e īn ẓohūr bar īn ʿāref shode-and,
che ke dar alvāḥ-e pārsiyye ke dar senīn-e tavaqqof dar ʿErāq az samāʾ-e
ʿezz-e aḥadiyye nāzel, jamīʿ-e īn marāteb be-tafṣīl dar ān alvāḥ madhkūr
āmade. Fa-ṭūbā li-l-fāʾizīna wa-l-ʿārifīn.

Beyond this, the Manifestations of divine unity are united in one
respect and appear in different revelations in another. Insofar as all
shine from the dawning-place of the Cause, have spoken from God, and
have been commissioned to convey the message to humanity, they are all
one. Thus He says, ``We make no distinction between any of them.'' In
another station, that of laws, names, and earthly conditions, each has
appeared with a name, a law, a human form, and an outward appearance.
Yet if you look with the eye of insight---by God, besides whom there is
no other god---you will behold them all as one reality, one revelation,
one movement, one deed, one action, one law, one name, and one practice.
It is astonishing that you have not grasped this explanation. The
children of this Revelation have come to understand it, for all these
degrees have been described in detail in the Persian Tablets that
descended from the heaven of the glory of divine unity during the years
of residence in Iraq. Blessed, then, are those who have attained and
those who understand.

\subsection{302}\label{section-797}

\begin{otherlanguage}{persian}

و تو هنوز مطّلع بر كيفيت ظهور نشده و قدرت الهيّه را مثل قدرت خود فرض گرفته
و امرش را منوط بامر خود فبئس ما ظننت انّه يظهر كيف یشآء كما ظهر بالحقّ
رغماً لأنف المشركين

\end{otherlanguage}

Va to hanūz moṭṭaleʿ bar keyfiyyat-e ẓohūr na-shode va qodrat-e elāhiyye
rā mesl-e qodrat-e khod farż gerefte va amrash rā manūṭ be-amr-e khod.
Fa-biʾsa mā ẓananta. Innahu yaẓharu kayfa yashāʾu kamā ẓahara bi-l-ḥaqqi
raghman li-anfi al-mushrikīn.

You have not yet understood the nature of the Revelation. You have
supposed divine power to be like your own power and made His Cause
dependent on your own affairs. How evil is what you have imagined! He
appears however He wills, just as He has appeared in truth, despite
those who associate partners with God.

\subsection{303}\label{section-798}

\begin{otherlanguage}{persian}

همان قسم كه نقطه اولى روح ما سواه فداه فرموده است اين ظهور بعينه همان
قسم ميفرمايد تفكّر لتعرف و اينست همان ظهور موعودى كه كلّ بآن وعده داده
شده‌اند در بيان و اليوم عمل كلّ من اعرض عن هذا الظّهور عنداللّه مردود و غير
مذكور انّه يقول حينئذ انّنى انا اللّه لا اله الّا انا كما قال النّقطة من قبل
و بعينه يقول من يأتى من بعد آيا نديده كه ميفرمايد اگر صد‌هزار مرتبه شمس
طالع شود يك شمس بوده و خواهد بود و ميفرمايد همين قسم در مظاهر احديّه
ملاحظه نمائيد

\end{otherlanguage}

Hamān qesm ke Noqṭe-ye Ūlā, rūḥu mā siwāhu fidāhu, farmūde ast, īn ẓohūr
be-ʿayneh hamān qesm mī-farmāyad. Tafakkar li-taʿrif. Va īnast hamān
ẓohūr-e mowʿūdī ke koll be-ān vaʿde dāde shode-and dar Bayān. Va al-yawm
ʿamal-e kulli man aʿraḍa ʿan hādhā al-ẓuhūr ʿinda Allāh mardūd va
gheyr-e madhkūr. Innahu yaqūlu ḥīnaʾidhin innanī anā Allāhu lā ilāha
illā anā kamā qāla al-Nuqṭatu min qablu wa-bi-ʿaynihi yaqūlu man yaʾtī
min baʿd. Āyā na-dīde ke mī-farmāyad agar ṣad-hezār martabe shams ṭāleʿ
shavad yek shams būde va khwāhad būd, va mī-farmāyad hamīn qesm dar
maẓāher-e aḥadiyye molāḥeẓe namāʾīd.

This Revelation speaks in precisely the same manner as the Primal
Point---may the souls of all else be sacrificed for Him---has spoken.
Reflect, so that you may understand. This is that very promised
Revelation which all were promised in the Bayán. Today the deeds of
everyone who turns away from this Revelation are rejected before God and
receive no mention. He now says, ``I am indeed God; there is no god but
Me,'' just as the Point said before, and precisely the same will be said
by whoever comes afterward. Have you not seen that He says, ``If the sun
rises a hundred thousand times, it has been and will be one sun''? And
He says, ``Regard the Manifestations of divine unity in this same way.''

\subsection{304}\label{section-799}

\begin{otherlanguage}{persian}

بدانكه در تجلّى دو نظر ملحوظ بنظرى مكرّر و بنظرى غير مكرّر و اين دو نظر
منسوب بتجلّيات ظاهريّه ملكيّه است و نفس تجلّى در مقرّ خود مقدّس از اين دو نظر
بوده و خواهد بود و لكن ادراك اينمقام قدرى مشكل هر نفسى بخواهد على ما هو
عليه ادراك نمايد بايد بين يدى الوجه حاضر شود ليعرف ما لاعرفه من قبل و
يطلع بما هو المستور

\end{otherlanguage}

Be-dān ke dar tajallī do naẓar malḥūẓ: be-naẓarī mokarrar va be-naẓarī
gheyr-e mokarrar. Va īn do naẓar mansūb be-tajalliyāt-e ẓāheriyye-ye
molkiyye ast, va nafs-e tajallī dar maqarr-e khod moqaddas az īn do
naẓar būde va khwāhad būd. Va lāken edrāk-e īn maqām qadrī moshkel. Har
nafsī be-khwāhad ʿalā mā huwa ʿalayhi edrāk namāyad bāyad bayna yaday
al-wajh ḥāżer shavad li-yaʿrifa mā lā ʿarafahu min qablu wa-yaṭṭaliʿa
bi-mā huwa al-mastūr.

Know that effulgence is considered from two perspectives: from one it
recurs, and from the other it does not. These two perspectives pertain
to outward, earthly effulgences, while effulgence itself, in its own
station, has been and will remain sanctified above both perspectives.
Understanding this station, however, is somewhat difficult. Anyone who
wishes to understand it as it truly is must come into the presence of
the Countenance, so that he may understand what he did not understand
before and become acquainted with what is concealed.

\subsection{305}\label{section-800}

\begin{otherlanguage}{persian}

و ديگر تجدّد امثال را محال دانسته اگر تجدّد امثال مظاهر الهيّه محال باشد
تجدّد امثال شيطانيّه كه موجود و مشهود است و مجال انكار نه چنانچه ديده
ميشود كه در هر شهرى مثالى از هياكل ظلمانيّه بعينه موجود و مشهود فاعتبروا
يا اولو الألباب

\end{otherlanguage}

Va dīgar tajaddod-e amthāl rā moḥāl dāneste. Agar tajaddod-e amthāl-e
maẓāher-e elāhiyye moḥāl bāshad, tajaddod-e amthāl-e shayṭāniyye ke
mowjūd va mashhūd ast va majāl-e enkār na, chenānche dīde mī-shavad ke
dar har shahrī methālī az hayākel-e ẓolmāniyye be-ʿayneh mowjūd va
mashhūd. Fa-ʿtabirū yā ulī al-albāb.

Furthermore, you have held the renewal of like beings to be impossible.
If the renewal of beings like the divine Manifestations is impossible,
the renewal of satanic counterparts is nonetheless a visible reality
that leaves no room for denial. As can be seen, in every city an exact
likeness of these embodiments of darkness exists and is visible. Take
heed, then, O people of understanding!

\subsection{306}\label{section-801}

\begin{otherlanguage}{persian}

بارى تجلّى از مجلّى من غير انقطاع بوده و خواهد بود در صور مظاهر اختلاف
ظاهر لو انت من العارفين و اگر قدرى برتر خرامى و از كلمات مشركين قلب خود
را مطهّر نمائى خود ادراك مينمائى آنچه را كه حال من غير التفات انكار نموده
ذكر تكرار تجلّى نيست كه جايز باشد يا غير جايز

\end{otherlanguage}

Bārī tajallī az mojallī min ghayri inqiṭāʿ būde va khwāhad būd. Dar
ṣovar-e maẓāher ekhtelāf ẓāher, law anta min al-ʿārifīn. Va agar qadrī
bartar kherāmī va az kalemāt-e moshrekīn qalb-e khod rā moṭahhar namāʾī,
khod edrāk mī-namāʾī ānche rā ke ḥāl min ghayri iltifāt enkār namūde.
Dhekr-e tekrār-e tajallī nīst ke jāyez bāshad yā gheyr-e jāyez.

In short, effulgence has proceeded and will continue to proceed without
interruption from its Source. Differences appear in the forms of the
Manifestations, if you are among those who understand. If you rise a
little higher and cleanse your heart of the words of those who associate
partners with God, you will yourself understand what you have now denied
without consideration. The issue is not the recurrence of effulgence,
such that it might be admissible or inadmissible.

\subsection{307}\label{section-802}

\begin{otherlanguage}{persian}

بشنو كلمات اين عبد را و بقلب طاهر و بصر منير در كلمات اللّه و ظهوره
ملاحظه نما و بدو كلمه مجعوله كه شيطان القاء نموده افتخار منما در امم قبل
و علماى آنها مشاهده كن كه صدهزار كتب تأليف نموده‌اند و مقصودشان از جميع
آن كتب اثبات حقّیّت حقّ بوده و بعد از ظهور كلّ بهمان كتب محروم مانده مع‌ذلك
هنوز معرضين متنبّه نشده تازه دست برساله نويسى گذاشته‌اند فويل لأناملهم و
بصرهم و يدهم در حين ظهور هيچ كلمه بمقرّ خود مستقرّ نه الّا باذن مظهر ظهور

\end{otherlanguage}

Be-shnow kalemāt-e īn ʿabd rā va be-qalb-e ṭāher va baṣar-e monīr dar
kalemāt-e Allāh va ẓohūrash molāḥeẓe namā va be-do kaleme-ye majʿūle ke
shayṭān elqā namūde eftekhār ma-namā. Dar omam-e qabl va ʿolamā-ye ānhā
moshāhede kon ke ṣad-hezār kotob taʾlīf namūde-and va maqṣūdshān az
jamīʿ-e ān kotob ethbāt-e ḥaqqiyyat-e Ḥaqq būde va baʿd az ẓohūr koll
be-hamān kotob maḥrūm mānde. Maʿa dhālek hanūz moʿreżīn motanabbeh
na-shode tāze dast be-resāle-nevīsī godhāshte-and. Fa-waylun
li-anāmilihim wa-baṣarihim wa-yadihim. Dar ḥīn-e ẓohūr hīch kaleme
be-maqarr-e khod mostaqarr na illā be-edhn-e maẓhar-e ẓohūr.

Listen to the words of this servant, and with a pure heart and clear
sight consider the words of God and His Revelation. Do not take pride in
two fabricated words that Satan has suggested. Consider the peoples of
the past and their scholars, who composed a hundred thousand books.
Their purpose in all those books was to establish the truth of the True
One, yet after His appearance they were all deprived of Him by those
very books. Nevertheless, those who turn away have still not taken heed
and have now begun writing treatises. Woe to their fingers, their sight,
and their hands! At the time of Revelation, no word is established in
its proper place except by permission of the Manifestation of
Revelation.

\subsection{308}\label{section-803}

\begin{otherlanguage}{persian}

نقطه بيان جلّ و عزّ فرموده كه در يوم ظهور بكلمات او بر او احتجاج منمائيد
بالقاى شيطان بمجعولات انفس معرضه احتجاج نموده و مينمائيد و همين احتجاجات
شما و امثال شما را خبر فرموده‌اند صدق اللّه العلىّ العظيم نصّ بيان منزل بيان
است كه ميفرمايد قوله عزّ ذكره منتهاى جدّ و جهد خود را نموده كه در يوم ظهور
حقّ از كلمات او بر او احتجاج نگردد كه كلّ بيان كلام ظهور قبل او است و او
است عالم‌تر از آنچه نازل فرموده از كلّ خلق زيرا كه روح كلّ در قبضه او است و
در نزد كلّ نيست الّا شبحى اگر بر حقّ مستقرّ باشند و الّا كه لايق ذكر نيست
چنانچه نفسى اگر افلاطون زمان باشد در هر علمى اگر تصديق نكند حقّ را آيا
ثمر خواهد بخشيد او را علم او سبحان اللّه عن ذلك بلكه از براى او علم اون
است كه علم بخدا و رسول و مظاهر امر و مظاهر حكم آن باشد و مادون آن را
اطلاق علم اولوالأفئده ننمايند انتهى

\end{otherlanguage}

Noqṭe-ye Bayān, jalla wa-ʿazza, farmūde ke dar yawm-e ẓohūr be-kalemāt-e
ū bar ū eḥtejāj ma-namāʾīd. Be-elqā-ye shayṭān be-majʿūlāt-e anfos-e
moʿreże eḥtejāj namūde va mī-namāʾīd, va hamīn eḥtejājāt-e shomā va
amthāl-e shomā rā khabar farmūde-and. Ṣadaqa Allāhu al-ʿAliyyu al-ʿAẓīm.
Naṣṣ-e bayān-e Monazzel-e Bayān ast ke mī-farmāyad, qawluhu ʿazza
dhikruhu: montahā-ye jedd va jahd-e khod rā namūde ke dar yawm-e ẓohūr-e
Ḥaqq az kalemāt-e ū bar ū eḥtejāj na-gardad, ke koll-e Bayān kalām-e
ẓohūr-e qabl-e ū ast va ū ast ʿālem-tar az ānche nāzel farmūde az koll-e
khalq, zīrā ke rūḥ-e koll dar qabże-ye ū ast va dar nazd-e koll nīst
illā shabaḥī agar bar Ḥaqq mostaqarr bāshand, va illā ke lāyeq-e dhekr
nīst. Chenānche nafsī agar Aflāṭūn-e zamān bāshad dar har ʿelmī, agar
taṣdīq na-konad Ḥaqq rā āyā thamar khwāhad bakhshīd ū rā ʿelm-e ū?
Subḥāna Allāhi ʿan dhālik. Balke az barā-ye ū ʿelm ūn ast ke ʿelm
be-Khodā va rasūl va maẓāher-e amr va maẓāher-e ḥokm-e ān bāshad, va
mā-dūn-e ān rā eṭlāq-e ʿelm ulū al-afʾede na-namāyand. Intahā.

The Point of the Bayán---glorious and exalted is He---has said, ``On the
day of Revelation, do not use His words as arguments against Him.'' At
Satan's prompting, you have used, and continue to use, the fabrications
of souls who turn away as arguments against Him. He foretold these very
arguments of yours and of others like you. God, the Most High, the Most
Great, has spoken the truth. This is the explicit statement of the
Revealer of the Bayán, who says---exalted be His mention---``He has
exerted His utmost effort so that, on the day of the appearance of the
True One, His words should not be used as arguments against Him, for the
whole Bayán is the utterance of His previous Revelation, and He knows
what He has revealed better than all creation, since the spirit of all
is in His grasp. With all others there is only a shadow, if they are
established in the truth; otherwise, it is not worthy of mention. Thus,
even if someone were the Plato of his age in every science, would his
knowledge benefit him if he did not acknowledge the True One? Far be it
from God! Rather, for him knowledge is that which is knowledge of God,
His Messenger, the Manifestations of His Cause, and the Manifestations
of His command. Those endowed with understanding do not apply the term
knowledge to anything else.'' End of quotation.

\subsection{309}\label{section-804}

\begin{otherlanguage}{persian}

قل فافتح بصرك يا ايّها الغافل ثمّ اسمع ما تغرّد الورقآء على افنان دوحة
البيان باللّه المقتدر العزيز المنّان ملاحظه نما كه امر تا بچه مقام بلند و
مرتفع است كه آن هيكل احديّه و ذات صمديّه كلمات خود را از آن ساحت اقدس اطهر
اظهر انور منع فرموده‌اند و بنصّ صريح ميفرمايند كه در حين ظهور بكلمات بيان
بر آن هيكل معانى و مظهر عزّ صمدانى احتجاج منمائيد و مع‌ذلك اهل بيان شرم
ننموده و ايكاش بكلمات او احتجاج مينمودند بلكه بكلمات مجعوله مردوده كه از
انفس خبيثه ظاهر شده بر جواهر احديّه و ساذج قدميّه احتجاج نموده و مينمايند
زهى حسرت بر اين نفوس كه بدو كلمه الفاظ مجعوله از شاطى بحر علم الهى محروم
مانده‌اند و برايحه منتنه از نفحات قدس احديّه ممنوع گشته‌اند و ايكاش كه
باحتجاج كفايت ميشد لا فو‌الّذى دلع لسان الرّوح بثنآء نفسه كه باحتجاج قانع
نشده و نخواهند شد بلكه باسياف كين و رماح بغضا بشطر علىّ اعلى توجّه نموده و
مينمايند

\end{otherlanguage}

Qul fa-ftaḥ baṣaraka yā ayyuhā al-ghāfilu thumma ismaʿ mā tugharridu
al-Warqāʾu ʿalā afnāni dawḥati al-bayāni bi-Allāhi al-Muqtadiri
al-ʿAzīzi al-Mannān. Molāḥeẓe namā ke amr tā be-che maqām boland va
mortafeʿ ast ke ān heykal-e aḥadiyye va dhāt-e ṣamadiyye kalemāt-e khod
rā az ān sāḥat-e aqdas-e aṭhar-e aẓhar-e anvar manʿ farmūde-and va
be-naṣṣ-e ṣarīḥ mī-farmāyand ke dar ḥīn-e ẓohūr be-kalemāt-e Bayān bar
ān heykal-e maʿānī va maẓhar-e ʿezz-e ṣamadānī eḥtejāj ma-namāʾīd. Va
maʿa dhālek ahl-e Bayān sharm na-namūde, va ey kāsh be-kalemāt-e ū
eḥtejāj mī-namūdand, balke be-kalemāt-e majʿūle-ye mardūde ke az anfos-e
khabīthe ẓāher shode bar javāher-e aḥadiyye va sādhaj-e qedamiyye
eḥtejāj namūde va mī-namāyand. Zehī ḥasrat bar īn nofūs ke be-do kaleme
alfāẓ-e majʿūle az shāṭī-ye baḥr-e ʿelm-e elāhī maḥrūm mānde-and va
be-rāyeḥe-ye montene az nafaḥāt-e qods-e aḥadiyye mamnūʿ gashte-and. Va
ey kāsh ke be-eḥtejāj kefāyat mī-shod. Lā fa-wa-alladhī dalaʿa lisāna
al-Rūḥi bi-thanāʾi nafsihi ke be-eḥtejāj qāneʿ na-shode va na-khwāhand
shod, balke be-asyāf-e kīn va remāḥ-e baghżā be-shaṭr-e ʿAliyy-e Aʿlā
tavajjoh namūde va mī-namāyand.

Say: Open your eyes, O heedless one, and then listen to what the Dove
sings upon the branches of the tree of utterance concerning God, the
All-Powerful, the Almighty, the All-Bountiful. Consider how lofty and
exalted the Cause is, that that Temple of divine unity and
self-subsisting Essence has barred His own words from that most holy,
most pure, most manifest, most luminous court, expressly saying that at
the time of Revelation you must not use the words of the Bayán as
arguments against that Temple of inner meanings and Manifestation of
self-subsisting glory. Nevertheless, the people of the Bayán have felt
no shame. If only they had used His words as arguments! Instead, they
have used, and continue to use, rejected, fabricated words issuing from
malicious souls as arguments against the Gems of divine unity and the
Essence of eternity. Alas for these souls, who have been deprived of the
shore of the ocean of divine knowledge by two fabricated expressions,
and have been kept from the holy breezes of divine unity by a foul odor!
If only they had been content with argument. No, by Him who has set the
tongue of the Spirit speaking in His own praise, they have not been and
will not be content with argument. Rather, with swords of malice and
spears of hatred, they have advanced and continue to advance toward the
Most High, the Most Exalted.

\subsection{310}\label{section-805}

\begin{otherlanguage}{persian}

و هم‌‌چنين ميفرمايند بيانه عزّ بيانه احدى بر بيان مطّلع نه و بعلم او محيط
نبوده و نخواهد بود مگر آن هيكل الهيّه و معدن و منبع علوم ربّانيّه مع‌ذلك
نفوسيكه بمحاربه برخواسته ادّعاى علم نموده و مينمايند مع آنكه فو‌اللّه الّذى
لا اله الّا هو كه در دبستان علم حقيقى با طفل ابجد خوان برابرى نكرده و
نخواهند كرد بلكه در آن محضر اگر وارد شوند خود را كليل مشاهده نمايند
ميفرمايد روح كلّ در قبضه اقتدار آن ظهور امنع اقدس بوده و خواهد بود و در
نزد كلّ نيست مگر شبحى از او و اينهم اگر بر حقّ ثابت باشند و الّا ميفرمايد
كه قابل ذكر نيستند

\end{otherlanguage}

Va ham-chenīn mī-farmāyand, bayānuhu ʿazza bayānuhu: aḥadī bar Bayān
moṭṭaleʿ na va be-ʿelm-e ū moḥīṭ na-būde va na-khwāhad būd magar ān
heykal-e elāhiyye va maʿdan va manbaʿ-e ʿolūm-e rabbāniyye. Maʿa dhālek
nofūsī ke be-moḥārebe bar-khāste eddeʿā-ye ʿelm namūde va mī-namāyand,
maʿ ānke fa-wa-Allāhi alladhī lā ilāha illā huwa ke dar dabestān-e
ʿelm-e ḥaqīqī bā ṭefl-e abjad-khwān barābarī na-karde va na-khwāhand
kard, balke dar ān maḥżar agar vāred shavand khod rā kalīl moshāhede
namāyand. Mī-farmāyad rūḥ-e koll dar qabże-ye eqtedār-e ān ẓohūr-e
amnaʿ-e aqdas būde va khwāhad būd va dar nazd-e koll nīst magar shabaḥī
az ū, va īn ham agar bar Ḥaqq thābet bāshand, va illā mī-farmāyad ke
qābel-e dhekr nīstand.

He also says---His utterance, exalted be His utterance---``No one is
acquainted with the Bayán, nor has anyone encompassed or will anyone
encompass its knowledge, except that divine Temple, the mine and
wellspring of divine knowledge.'' Nevertheless, those who have risen to
wage war have claimed and continue to claim knowledge, although---by
God, besides whom there is no other god---in the school of true
knowledge they have never equaled and will never equal a child learning
the alphabet. Indeed, if they entered that presence, they would find
themselves incapable. He says that the spirit of all has been and will
remain in the grasp of the power of that most unapproachable, most holy
Revelation, and that with all others there is only a shadow of it---and
even this only if they remain steadfast in the truth. Otherwise, He
says, they are not worthy of mention.

\subsection{311}\label{section-806}

\begin{otherlanguage}{persian}

حال ملاحظه كن اين نفوس را كه بچه مقام افتاده‌ند فواللّه فى اسفل الهاوية لو
انتم تشعرون اليوم كلّ علوم ارض عنداللّه معدوم و غير مذكور بلكه علم الهى كه
مقصود بوده و در كلّ كتب ربّانى و زبر صمدانى ممدوح علم بمظهر نفس او بوده در
هر ظهورى حين ظهور اليوم هر نفسى بآن فايز شد ولو حرفى عالم نباشد اعلم ناس
عنداللّه مذكور و اگر العياذ باللّه از اين فضل اعظم محروم ماند و حافظ كلّ
كتب باشد اجهل از كلّ ناس عنداللّه مشهود

\end{otherlanguage}

Ḥāl molāḥeẓe kon īn nofūs rā ke be-che maqām oftāde-and. Fa-wa-Allāhi fī
asfali al-hāwiyati law antum tashʿurūn. Al-yawm koll-e ʿolūm-e arż ʿinda
Allāh maʿdūm va gheyr-e madhkūr. Balke ʿelm-e elāhī ke maqṣūd būde va
dar koll-e kotob-e rabbānī va zobor-e ṣamadānī mamdūḥ, ʿelm be-maẓhar-e
nafs-e ū būde dar har ẓohūrī ḥīn-e ẓohūr. Al-yawm har nafsī be-ān fāyez
shod va law ḥarfī ʿālem na-bāshad aʿlam-e nās ʿinda Allāh madhkūr, va
agar al-ʿiyādhu bi-Allāh az īn fażl-e aʿẓam maḥrūm mānad va ḥāfeẓ-e
koll-e kotob bāshad ajhal az koll-e nās ʿinda Allāh mashhūd.

Now consider to what a station these souls have fallen. By God, they are
in the lowest depths of the abyss, if only you perceived it. Today all
the sciences of the earth are as nothing before God and receive no
mention. Rather, the divine knowledge that has been intended and has
been praised in all the divine Books and everlasting Scriptures is
knowledge of the Manifestation of His own Self in every Revelation, at
the time of His appearance. Today, whoever attains this, even if he does
not know a single letter, is accounted before God the most learned of
all people. But if---God forbid---he remains deprived of this supreme
bounty, even if he knows all the books by heart, he is regarded before
God as more ignorant than anyone else.

\subsection{312}\label{section-807}

\begin{otherlanguage}{persian}

بشنو بيان ديگر منزل بيان روح ماسواه فداه را قوله عزّ اجلاله و همچنين تصوّر
كن در نزد ظهور آن جمال قدم كه اگر كلّ اهل بيان در جوهر علم مثل او گردند
كه ثمر نمى‌بخشد الّا بتصديق باو فلتعتبرن ان يا اولى العلم ثمّ ايّاه تتّقون
انتهى واللّه حال خجلت ميبرد اين عبد كه چه ذكر نمايد اين بيان نقطه اولى
روح ماسواه فداه و آن كلمات شما كه عالم را بكدورت و غبار تيره نموده فارفع
رأسك عن النّوم ثمّ اشهد انوار تجلّى ربّك لعلّ تطّلع بما كنت غافلاً عنه و تكون
من التّائبين الى اللّه الّذى خلقنى و ایّاک و جعلنا من الّذينهم ادركوا ایّام
لقائه و كانوا من الفائزين از خدا ميطلبم كه نسيم عنايتى بوزد كه شايد آن
مرده را بعرصه حيات كشاند و از فضل و عنايت خود قطره از ماء حيوان و كوثر
جاريه از يمين عرش ربّك الرّحمن مشروب فرمايد و او است بر كلّ شیء قادر و
توانا

\end{otherlanguage}

Beshnow bayān-i dīgar-i Munzil-i Bayān rūḥ-i mā-sivāhu fidāh rā, qawluhu
ʿazza ijlāluhu: Va hamchenīn taṣavvor kon dar nazd-i ẓuhūr-i ān Jamāl-i
Qidam ke agar koll-i ahl-i Bayān dar jawhar-i ʿelm mes̱l-i ū gardand ke
s̱amar nemī-bakhshad illā be-taṣdīq be-ū. Fa-ltaʿtabirunna an yā ulī
al-ʿilm thumma iyyāhu tattaqūn. Intahā. Va'llāh ḥāl khejlat mībarad īn
ʿabd ke che dhekr namāyad, īn bayān-i Nuqṭe-yi Ūlā rūḥ-i mā-sivāhu fidāh
va ān kalemāt-i shomā ke ʿālam rā be-kodūrat va ghobār tīre namūde.
Fa-rfaʿ raʾsaka ʿan al-nawm thumma ishhad anwār-i tajallī rabbika
laʿalla taṭṭaliʿu bi-mā kunta ghāfilan ʿanhu wa-takūnu min al-tāʾibīn
ilā Allāh alladhī khalaqanī wa-iyyāka wa-jaʿalanā min alladhīna hum
adrakū ayyāma liqāʾihi wa-kānū min al-fāʾizīn. Az Khodā mīṭalabam ke
nasīm-i ʿenāyatī bevazad ke shāyad ān morde rā be-ʿarṣe-yi ḥayāt
keshānad va az faḍl va ʿenāyat-i khod qaṭre az māʾ-i ḥayavān va
kawthar-i jāriye az yamīn-i ʿarsh-i rabbika al-Raḥmān mashrūb farmāyad
va ū ast bar koll-i shayʾ qāder va tavānā.

Listen to another utterance of the Revealer of the Bayán---may the
spirit of all else be sacrificed for Him. His words, exalted be His
majesty: ``And likewise consider that, at the appearance of that Ancient
Beauty, even if all the people of the Bayán should become like Him in
the essence of knowledge, this yields no fruit except through affirming
Him. Take heed, then, O people of knowledge, and fear Him.'' End of
quotation. By God, this servant is now ashamed to speak: here is this
utterance of the Primal Point---may the spirit of all else be sacrificed
for Him---and there are those words of yours that have darkened the
world with gloom and dust. Raise your head from sleep, then behold the
lights of your Lord's revelation, so that you may become aware of what
you have neglected and be among those who turn in repentance to God, who
created me and you and made us among those who attained the days of
meeting Him and were among the victorious. I ask God that a breeze of
loving care may blow, that perhaps it may draw that dead one into the
realm of life, and that, through His grace and loving care, He may give
him a drop of the water of life and of the heavenly stream flowing from
the right hand of the throne of your Lord, the All-Merciful. He is
powerful and mighty over all things.

\subsection{313}\label{section-808}

\begin{otherlanguage}{persian}

از اينمراتب گذشته در يحيى بن ذكريّا و عيسى بن مريم چه ميگوئيد كه يحيى بن
ذكريّا مبعوث شد قبل از عيسى و ادّعاى نبوّت هم نمود چنانچه جميع فرق اسلام
بنبوّت او قائل و مذعنند و باحكام و شرايع هم ظاهر شدند چنانچه هر نفسى كه
بين يديه حاضر ميشد و برسالتش معترف بآدابيكه در آنوقت از سماء مشيّت نازل
باب غسلش ميداد و تلقين كلمات الهيّه ميفرمود و همچنين بشارت ميداد ناس را
بظهور كلمه اعظم يعنى روح اللّه

\end{otherlanguage}

Az īn marāteb godhashte dar Yaḥyā ibn-i Zakariyyā va ʿĪsā ibn-i Maryam
che mīgūʾīd ke Yaḥyā ibn-i Zakariyyā mabʿūs̱ shod qabl az ʿĪsā va
eddeʿā-yi nobovvat ham namūd chenānche jamīʿ-i feraq-i Eslām
be-nobovvat-i ū qāʾel va modhʿenand va be-aḥkām va sharāyeʿ ham ẓāher
shodand chenānche har nafsī ke bayn-i yadayh ḥāḍer mīshod va
be-resālatash moʿtaref be-ādābī ke dar ān vaqt az samāʾ-i mashiyyat
nāzel ba-āb ghoslash mīdād va talqīn-i kalemāt-i elāhiyye mīfarmūd va
hamchenīn beshārat mīdād nās rā be-ẓuhūr-i Kaleme-yi Aʿẓam yaʿnī Rūḥ
Allāh.

Leaving these matters aside, what do you say about John son of Zechariah
and Jesus son of Mary? John son of Zechariah was sent before Jesus and
also claimed prophethood, as all the sects of Islam acknowledge and
affirm his prophethood. He also appeared with laws and ordinances. Thus,
whenever anyone entered his presence and acknowledged his mission, he
baptized him with water according to the rites sent down at that time
from the heaven of the divine will, and instructed him in the divine
words. He also gave the people glad tidings of the appearance of the
Most Great Word, that is, the Spirit of God.

\subsection{314}\label{section-809}

\begin{otherlanguage}{persian}

بعد از آنكه روح ظاهر شد و كلمه الهيّه از خلف حجبات ستريّه باشراقات احديّه
طالع و مشرق و منير گشت ناس اعتراض نمودند چه كه آنحضرت بحسب ظاهر با
خاطئين مجالس و مؤانس بودند چنانچه در خود انجيل مسطور و مكتوب كه با
عشارين معاشرت ميفرمود چنانچه در مواضع عديده بر آن سلطان احديّه اعتراض
نمودند مخصوص در يومى از ایّام بعضى از غافلين بين يديه حاضر و اعتراض نمود
يا روح اللّه اصحاب يحيى بن ذكريّا بسيار زاهد و عابد بودند بشأنيكه در ایّام
صائم و ليالى بذكر اللّه قائم لكن آنچه از شما و اصحاب شما مشاهده ميشود
بعبادت و ذكر مشغول نيستيد بلكه اجتناب از مال ناس نمينمائيد فتفكّر الرّوح
فى نفسه ثمّ رفع رأسه و قال لو كان يحيى على ما انتم تقولون لم قتلتموه يا
ملأ المشركين ميفرمايد بقول شما من خاطى يحيى بن ذكريّا كه تقصيرى نكرده و
جميع هم مقرّيد بزهد و تقواى آن جمال مبين پس او را چرا كشتيد فبهت الّذى كفر
و خائباً خاسراً راجع گشت

\end{otherlanguage}

Baʿd az ānke Rūḥ ẓāher shod va Kaleme-yi Elāhiyye az khalf-i ḥojobāt-i
setriyye be-eshrāqāt-i aḥadiyye ṭāleʿ va moshreq va monīr gasht nās
eʿterāḍ namūdand che ke ān Ḥaḍrat be-ḥasab-i ẓāher bā khāṭeʾīn mojāles
va moʾānes būdand chenānche dar khod-i Enjīl masṭūr va maktūb ke bā
ʿashshārīn moʿāsherat mīfarmūd chenānche dar mavāḍeʿ-i ʿadīde bar ān
Solṭān-i Aḥadiyye eʿterāḍ namūdand makhṣūṣ dar yawmī az ayyām baʿḍī az
ghāfelīn bayn-i yadayh ḥāḍer va eʿterāḍ namūd yā Rūḥ Allāh aṣḥāb-i Yaḥyā
ibn-i Zakariyyā besyār zāhed va ʿābed būdand be-shaʾnī ke dar ayyām
ṣāʾem va layālī be-dhekr-i Allāh qāʾem lāken ānche az shomā va aṣḥāb-i
shomā moshāhede mīshavad be-ʿebādat va dhekr mashghūl nīstīd balke
ejtenāb az māl-i nās nemīnamāʾīd. Fa-tafakkara al-Rūḥ fī nafsihi thumma
rafaʿa raʾsahu wa-qāla law kāna Yaḥyā ʿalā mā antum taqūlūna li-ma
qataltumūhu yā malaʾ al-mushrikīn. Mīfarmāyad be-qawl-i shomā man khāṭī,
Yaḥyā ibn-i Zakariyyā ke taqṣīrī nakarde va jamīʿ ham moqerrīd be-zohd
va taqvā-yi ān Jamāl-i Mobīn pas ū rā cherā koshtīd. Fa-buhita alladhī
kafara va khāʾiban khāsiran rājeʿ gasht.

After the Spirit appeared and the divine Word rose from behind the
concealing veils, shining and resplendent with the radiance of divine
oneness, the people raised objections, because He outwardly associated
and consorted with sinners. It is recorded and written in the Gospel
itself that He associated with tax collectors. Thus, on numerous
occasions they objected to that Sovereign of Oneness. In particular, one
day some of the heedless entered His presence and objected: ``O Spirit
of God, the companions of John son of Zechariah were very ascetic and
devout, so much so that they fasted by day and stood in remembrance of
God by night. But from what we observe of you and your companions, you
do not engage in worship and remembrance; rather, you do not refrain
from taking people's property.'' The Spirit reflected within Himself,
then raised His head and said, ``If John was as you say, why did you
kill him, O company of those who associate partners with God?'' He says:
``According to you, I am a sinner. But John son of Zechariah had done no
wrong, and you all acknowledge the asceticism and piety of that manifest
Beauty. Why, then, did you kill him?'' The one who disbelieved was
confounded and went back disappointed and defeated.

\subsection{315}\label{section-810}

\begin{otherlanguage}{persian}

و هم‌‌چنين اعتراض مينمودند كه تو شريعت يحيى را تغيير داده مع آنكه هنوز امر
او مابين عباد ثابت نشده و تكميل دوره او نگشته چه كه ایّامى نگذشت مابين آن
ظهور و اين ظهور و او عباد را باب غسل ميداد و توقائل كه بروح القدس غسل
ميدهم چنانچه اصحابت داده‌اند و آنها صائم بودند و اصحاب تو آكل و او بتمام
زهد ظاهر چنانچه ابدا در اموال ناس بقدر خردلى تصرّف ننموده و اصحاب شما هر
يك كه از مزارع عبور مينمايند اجتناب از اخذ و اكل نمينمايند و ديگر بهيچ
عملى عامل نيستند فرمودند اين ایّام ایّام شادى و عيش است وقتى آيد كه بكلّ
اعمال عامل شوند و مقصودشان اين بود كه چون ایّام لقاء اللّه است فى‌الحقيقة
مقام جنّت احديّه است و در آنمقام اگر در بعضى از امور ظاهره اهمال رود حقّ
تعالى شأنه عفو خواهد فرمود و بعد از ارتقاى روح برفيق اعلى آنچه باو مأمور
شده‌اند البتّه عامل خواهند شد

\end{otherlanguage}

Va hamchenīn eʿterāḍ mīnamūdand ke to sharīʿat-i Yaḥyā rā taghyīr dāde
maʿa ānke hanūz amr-i ū mābayn-i ʿebād s̱ābet nashode va takmīl-i
dawre-yi ū nagashte che ke ayyāmī nagodhasht mābayn-i ān ẓuhūr va īn
ẓuhūr va ū ʿebād rā ba-āb ghosl mīdād va to qāʾel ke be-Rūḥ al-Qodos
ghosl mīdaham chenānche aṣḥābat dāde-and va ānhā ṣāʾem būdand va aṣḥāb-i
to ākel va ū be-tamām-i zohd ẓāher chenānche abadan dar amvāl-i nās
be-qadr-i khardalī taṣarrof nanamūde va aṣḥāb-i shomā har yek ke az
mazāreʿ ʿobūr mīnamāyand ejtenāb az akhdh va akl nemīnamāyand va dīgar
be-hīch ʿamalī ʿāmel nīstand. Farmūdand īn ayyām ayyām-i shādī va ʿeysh
ast vaqtī āyad ke be-koll-i aʿmāl ʿāmel shavand va maqṣūd-shān īn būd ke
chon ayyām-i liqāʾ-i Allāh ast fī al-ḥaqīqe maqām-i jannat-i aḥadiyye
ast va dar ān maqām agar dar baʿḍī az omūr-i ẓāhere ehmāl ravad Ḥaqq
taʿālā shaʾnuhu ʿafv khāhad farmūd va baʿd az erteqā-yi Rūḥ be-Rafīq-i
Aʿlā ānche be-ū maʾmūr shode-and albatte ʿāmel khāhand shod.

They also objected: ``You have changed the law of John, even though his
Cause has not yet been established among the people and his dispensation
has not run its course, for only a few days have passed between that
revelation and this one. He baptized people with water, while you say,
`I baptize with the Holy Spirit,' as your companions have done. They
fasted, while your companions eat. He appeared in perfect asceticism,
never taking so much as a mustard seed's worth of people's property,
whereas whenever any of your companions pass through the fields, they do
not refrain from taking and eating, nor do they observe any other
practice.'' He said: ``These are days of joy and delight. A time will
come when they will observe all the practices.'' His meaning was that,
since these are the days of meeting God, this is in truth the station of
the paradise of oneness, and in that station, should there be negligence
in certain outward matters, God---exalted be His glory---will forgive
it. After the Spirit's ascent to the Most Exalted Companion, they will
certainly carry out what they have been commanded.

\subsection{316}\label{section-811}

\begin{otherlanguage}{persian}

فواللّه ظهور نقطه اولى و اين ظهور ابدع ابهى بعينه ظهور يحيى بن ذكريا و
روح اللّه است و جميع مطابق واقع شده ‌همان قسم كه يحيى نبی‌ و رسول بود من
عنداللّه و هم‌‌چنين مبشّر بظهور بعد چنانچه ميفرمود يا قوم انّى ابشّركم بملكوت
اللّه و انّه قد اقترب و در مقام ديگر و قد اقتربت ملكوت اللّه و هم‌‌چنين صاحب
احكام و شريعت بوده و هم‌‌چنين در ایّام ظهور او روح ظاهر شده نقطه اولى روح
ما سواه فداه بعد از آنكه اخذ عهد از كلّ نموده و بشارت داده بظهور بعد
ميفرمايد و لقد قرب الزّوال و انتم راقدون كه بعينه همان مضمون است كه يحيى
بن ذكريّا بآن تكلّم نموده و بشارت داده

\end{otherlanguage}

Fa-va'llāh ẓuhūr-i Nuqṭe-yi Ūlā va īn ẓuhūr-i abdaʿ-i abhā be-ʿayneh
ẓuhūr-i Yaḥyā ibn-i Zakariyyā va Rūḥ Allāh ast va jamīʿ moṭābeq vāqeʿ
shode hamān qesm ke Yaḥyā nabī va rasūl būd min ʿindi Allāh va hamchenīn
mobashsher be-ẓuhūr-i baʿd chenānche mīfarmūd yā qawmi innī
ubashshirukum bi-malakūt Allāh wa-innahu qad iqtaraba va dar maqām-i
dīgar wa-qad iqtarabat malakūt Allāh va hamchenīn ṣāḥeb-i aḥkām va
sharīʿat būde va hamchenīn dar ayyām-i ẓuhūr-i ū Rūḥ ẓāher shode
Nuqṭe-yi Ūlā rūḥ-i mā-sivāhu fidāh baʿd az ānke akhdh-i ʿahd az koll
namūde va beshārat dāde be-ẓuhūr-i baʿd mīfarmāyad wa-laqad qaruba
al-zawāl wa-antum rāqidūn ke be-ʿayneh hamān maḍmūn ast ke Yaḥyā ibn-i
Zakariyyā be-ān takallom namūde va beshārat dāde.

By God, the revelation of the Primal Point and this most wondrous, most
glorious Revelation are exactly like the revelations of John son of
Zechariah and the Spirit of God, and everything has occurred in
correspondence with them. John was a prophet and messenger from God, and
also the herald of the revelation to follow, as he said: ``O my people,
I bring you glad tidings of the Kingdom of God, and it has drawn near,''
and elsewhere, ``And the Kingdom of God has drawn near.'' He also
brought laws and a religious code, and the Spirit appeared during the
days of his revelation. Likewise, the Primal Point---may the spirit of
all else be sacrificed for Him---after taking a covenant from everyone
and giving glad tidings of the revelation to follow, says: ``The decline
has indeed drawn near, while you lie asleep.'' This conveys exactly the
same meaning that John son of Zechariah expressed and announced as glad
tidings.

\subsection{317}\label{section-812}

\begin{otherlanguage}{persian}

و در اين ظهور هم آنچه بعيسى بن مريم اعتراض نمودند ملأ بيان هم همان
اعتراضات را نموده چنانچه شما مينويسيد كه تكميل دوره نشده و همچنين
مينويسيد كه بايد احكام تغيير داده نشود و تلويحاً ميخواهيد ذكر نمائيد كه
احكام را تغيير داده‌اند و همين را كذب محض بشما تلقين نموده‌اند چه كه مخصوص
فرموده‌اند كه ببيان پارسى عمل نمايند و

\end{otherlanguage}

Va dar īn ẓuhūr ham ānche be-ʿĪsā ibn-i Maryam eʿterāḍ namūdand malaʾ-i
Bayān ham hamān eʿterāḍāt rā namūde chenānche shomā mīnevīsīd ke
takmīl-i dawre nashode va hamchenīn mīnevīsīd ke bāyad aḥkām taghyīr
dāde nashavad va talvīḥan mīkhāhīd dhekr namāʾīd ke aḥkām rā taghyīr
dāde-and va hamīn rā kedhb-i maḥḍ be-shomā talqīn namūde-and che ke
makhṣūṣ farmūde-and ke be-Bayān-i Pārsī ʿamal namāyand va

In this Revelation too, the people of the Bayán have raised the same
objections that were raised against Jesus son of Mary. Thus you write
that the dispensation has not run its course, and you also write that
the laws must not be changed, wishing to imply that He has changed the
laws. This very claim, which is an outright lie, has been instilled in
you, for He has expressly instructed them to act according to the
Persian Bayán, and

\subsection{318}\label{section-813}

\begin{otherlanguage}{persian}

ديگر تغيير احكام دخلى برتبه ندارد كه هر ظهورى ادّعاى برترى از ظهور قبل
نمايد بايد احكام تغيير دهد و ناسخ اوامر قبل باشد چنانچه خود نقطه بيان
فرموده كه اگر اعتراضات اهل فرقان نميبود شريعت فرقان را نسخ نمينمودم و از
آن گذشته باعتقاد خود شما و جميع حضرات داود يكى از انبياى اعظم است و صاحب
كتاب مع‌ذلك ناس را بشريعت تورية امر فرمود و هم‌‌چنين از كلمه عيسى بن مريم
است كه ميفرمايد ما جآء ابن الانسان ليدين العالم بل جآء ليحيى العالم از
اين بيان همچه معلوم ميشود كه در اوّل ظهور اراده تغيير احكام قبل نداشته‌اند
چنانچه يحيى بن ذكريا هم تغيير كلّيه نداده‌اند الّا در بعض امور و آنچه مسلّم
است حكم سبت را حضرت روح تغيير دادند و يوم احد قرار فرمودند و هم‌‌چنين حكم
سيف را مرتفع فرمودند و طلاق را هم منع نمودند

\end{otherlanguage}

Dīgar taghyīr-i aḥkām dakhlī be-rotbe nadārad ke har ẓuhūrī eddeʿā-yi
bartarī az ẓuhūr-i qabl namāyad bāyad aḥkām taghyīr dahad va nāsekh-i
avāmer-i qabl bāshad chenānche khod-i Nuqṭe-yi Bayān farmūde ke agar
eʿterāḍāt-i ahl-i Forqān nemībūd sharīʿat-i Forqān rā naskh nemīnamūdam
va az ān godhashte be-eʿteqād-i khod-i shomā va jamīʿ-i ḥaḍarāt Dāvūd
yekī az anbiyā-yi aʿẓam ast va ṣāḥeb-i ketāb maʿa dhālek nās rā
be-sharīʿat-i Tawrāt amr farmūd va hamchenīn az kaleme-yi ʿĪsā ibn-i
Maryam ast ke mīfarmāyad mā jāʾa ibn al-insān li-yadīna al-ʿālam bal
jāʾa li-yuḥyiya al-ʿālam az īn bayān hamche maʿlūm mīshavad ke dar
avval-i ẓuhūr erāde-yi taghyīr-i aḥkām-i qabl nadāshte-and chenānche
Yaḥyā ibn-i Zakariyyā ham taghyīr-i kolliyye nadāde-and illā dar baʿḍ-i
omūr va ānche mosallam ast ḥokm-i sabt rā Ḥaḍrat-i Rūḥ taghyīr dāde-and
va yawm-i aḥad qarār farmūde-and va hamchenīn ḥokm-i sayf rā mortafeʿ
farmūde-and va ṭalāq rā ham manʿ namūde-and.

Furthermore, changing the laws has no bearing on rank, such that every
revelation claiming superiority to the preceding revelation must change
the laws and abrogate the previous commandments. The Point of the Bayán
Himself has said: ``Were it not for the objections of the people of the
Qur'an, I would not have abrogated the law of the Qur'an.'' Moreover,
according to your own belief and that of all your associates, David is
one of the greatest prophets and the bearer of a Book; nevertheless, he
commanded the people to follow the law of the Torah. Likewise, among the
words of Jesus son of Mary is this saying: ``The Son of Man did not come
to judge the world, but came to give life to the world.'' From this
utterance it appears that, at the beginning of His revelation, He did
not intend to change the previous laws, just as John son of Zechariah
made no wholesale change, except in certain matters. What is certain is
that the Spirit changed the law of the Sabbath and appointed Sunday, and
likewise abolished the law of the sword and also prohibited divorce.

\subsection{319}\label{section-814}

\begin{otherlanguage}{persian}

و لكن در اين ظهور اگر ملاحظه نمائيد با اينكه در اوّل امر جز عبوديّت صرفه
امرى ظاهر نشد مع‌ذلك ذياب ارض از ملأ بيان باعراض و اعتراض برخواستند و از
ايكات نفاق بقصد نيّر آفاق‌ بيرون آمدند با آنكه حكم محكم بيان است و در كلّ
الواح اللّه مسطور كه اگر نفسى بآيات ظاهر شود باو معارضه ننمائيد و مجادله
نكنيد و احتجاج ننمائيد و بشأنى در اينحكم مبالغه فرموده‌اند كه واللّه الّذى
لا اله الّا هو حكمى از اين محكم تر و ثابت تر در بيان نازل نشده گويا جميع
بيان در اين امر ابدع امنع نازل شده مع‌ذلك اوّل معرض و امثال او باعراض
برخواسته و وارد آورده‌ا‌ند آنچه را كه هيچ ظالمى عامل نشده شما اين ارض
نبوده‌ايد كه مطّلع باشيد بحقّ نقطه اولى و جميع نبيّين و مرسلين كه بحقّ و
راستى تكلّم مينمايم چون ملاحظه نمودند كه آيات الهى مثل غيث هاطل نازل و
بحر اعظم ربّانى در امواج و شجره عزّ صمدانيّه باثمار بديعه منيعه مزيّن شده و
ارياح فضليّه از يمين عرش الهيّه در هبوب و سماء مرتفعه قدميّه بانجم علوم و
معارف جديده مطرّز بنفاق برخواست و درصدد قطع شجره ربّانيّه افتاد و چون از آن
عاجز شد مكر صدرش ظاهر گشت و جميع مهاجرين مضطرب و اهل حرم در نوحه و ندبه
مشغول و جمال احديّه از بين بريّه خارج شدند بعد بافترا و كذب برخواستند و
اراده نمودند كه بمفتريات جعليّه و ظنونات افكيّه خبيثه ستر نمايند آن فعل
شنيع را و بعد عامل شدند آنچه را كه عامل شدند چنانچه خود شما هم ميدانيد
كه آن مفسد مفترى كه بآن ارض آمده چه گفته و چه كرده و بآن اصل هاويه نفسى
متعرّض نشده و سخن نگفته مع‌ذلك خود را مظلوم ناميده و بامثال خود القاء
نموده آنچه را كه شيطان باولياى خود القاء نموده و مينمايد فويل لهم و
للّذين اتّبعوهم تاللّه انّهم لفى خسران عظيم

\end{otherlanguage}

Va lāken dar īn ẓuhūr agar molāḥeẓe namāʾīd bā īnke dar avval-i amr joz
ʿobūdiyyat-i ṣerfe amrī ẓāher nashod maʿa dhālek dheʾāb-i arḍ az malaʾ-i
Bayān be-eʿrāḍ va eʿterāḍ barkhāstand va az aykāt-i nefāq be-qaṣd-i
Nayyir-i Āfāq bīrūn āmadand bā ānke ḥokm-i moḥkam-i Bayān ast va dar
koll-i alvāḥ-i Allāh masṭūr ke agar nafsī be-āyāt ẓāher shavad be-ū
moʿāreḍe nanamāʾīd va mojādele nakonīd va eḥtejāj nanamāʾīd va be-shaʾnī
dar īn ḥokm mobāleghe farmūde-and ke va'llāh alladhī lā ilāha illā huwa
ḥokmī az īn moḥkam-tar va s̱ābet-tar dar Bayān nāzel nashode gūyā jamīʿ-i
Bayān dar īn amr-i abdaʿ-i amnaʿ nāzel shode maʿa dhālek avval-i moʿreḍ
va ams̱āl-i ū be-eʿrāḍ barkhāste va vāred āvarde-and ānche rā ke hīch
ẓālemī ʿāmel nashode shomā īn arḍ nabūde-īd ke moṭṭaleʿ bāshīd be-ḥaqq-i
Nuqṭe-yi Ūlā va jamīʿ-i nabiyyīn va morsalīn ke be-ḥaqq va rāstī
takallom mīnamāyam chon molāḥeẓe namūdand ke āyāt-i elāhī mes̱l-i
ghayth-i hāṭel nāzel va baḥr-i aʿẓam-i rabbānī dar amvāj va shajare-yi
ʿezz-i ṣamadāniyye be-as̱mār-i badīʿe-yi manīʿe mozayyan shode va aryāḥ-i
faḍliyye az yamīn-i ʿarsh-i elāhiyye dar hobūb va samāʾ-i mortafeʿe-yi
qidamiyye be-anjom-i ʿolūm va maʿāref-i jadīde moṭarraz be-nefāq
barkhāst va dar ṣadad-i qaṭʿ-i shajare-yi rabbāniyye oftād va chon az ān
ʿājez shod makr-i ṣadrash ẓāher gasht va jamīʿ-i mohājerīn moḍṭareb va
ahl-i ḥaram dar nawḥe va nodbe mashghūl va Jamāl-i Aḥadiyye az bayn-i
bariyye khārej shodand baʿd be-efterā va kedhb barkhāstand va erāde
namūdand ke be-moftarayāt-i jaʿliyye va ẓonūnāt-i efkiyye-yi khabīs̱e
satr namāyand ān feʿl-i shanīʿ rā va baʿd ʿāmel shodand ānche rā ke
ʿāmel shodand chenānche khod-i shomā ham mīdānīd ke ān mofsed-i moftarī
ke be-ān arḍ āmade che gofte va che karde va be-ān aṣl-i hāviye nafsī
motaʿarreḍ nashode va sokhan nagofte maʿa dhālek khod rā maẓlūm nāmīde
va be-ams̱āl-i khod elqā namūde ānche rā ke shayṭān be-awliyā-yi khod
elqā namūde va mīnamāyad fa-waylun lahum wa-li-alladhīna ittabaʿūhum
ta'llāh innahum la-fī khusrān ʿaẓīm.

Yet consider this Revelation: although at the outset nothing appeared
but pure servitude, the wolves of the earth among the people of the
Bayán nevertheless rose in rejection and opposition, emerging from the
thickets of hypocrisy to assail the Luminary of the horizons. This was
despite the firm command of the Bayán, recorded in all the Tablets of
God, that if anyone should appear with verses, you must not oppose him,
dispute with him, or argue against him. He has emphasized this command
to such a degree that, by God, besides whom there is no other god, no
command firmer or more established than this has been revealed in the
Bayán. It is as though the entire Bayán had been revealed concerning
this most wondrous, most unassailable Cause. Nevertheless, the first to
turn away and others like him have risen in rejection and inflicted what
no oppressor has ever perpetrated. You were not in this land to know of
it. By the Primal Point and all the prophets and messengers, I speak in
truth and sincerity. When he saw the divine verses descending like
pouring rain, the most great divine Ocean surging, the Tree of
everlasting glory adorned with wondrous, lofty fruits, the winds of
grace blowing from the right hand of the divine throne, and the lofty
heaven of eternity embroidered with the stars of new sciences and
knowledge, he rose in hypocrisy and sought to cut down the divine Tree.
When he proved incapable of this, the deceit within his breast became
manifest. All the emigrants were thrown into turmoil, the members of the
holy household were occupied with lamentation and mourning, and the
Beauty of Oneness withdrew from among the people. Then they resorted to
slander and lies, intending to conceal that heinous deed with fabricated
calumnies and foul, lying suppositions. Afterward, they did what they
did, as you yourself know what that mischief-making slanderer who came
to that land has said and done. No one has assailed or spoken against
that source of the abyss. Nevertheless, he has called himself wronged
and instilled in others like himself what Satan has instilled and
continues to instill in his allies. Woe to them and to those who have
followed them! By God, they are indeed in grievous loss.

\subsection{320}\label{section-815}

\begin{otherlanguage}{persian}

قول ديگر آنكه بايد اكمل و اتمّ باشد چنانچه ظهور هر رسولى و نبيّى اعلى از
ظهور سابق بوده و ربّ اعلى هم‌چنين ادّعا نموده بالنّسبة بسابقين اينست كه هر
لاحقى نسخ سابق نموده انتهى

\end{otherlanguage}

Qawl-i dīgar ānke bāyad akmal va atamm bāshad chenānche ẓuhūr-i har
rasūlī va nabiyyī aʿlā az ẓuhūr-i sābeq būde va Rabb-i Aʿlā hamchenīn
eddeʿā namūde bi'l-nesbe be-sābeqīn īn ast ke har lāḥeqī naskh-i sābeq
namūde. Intahā.

Another statement: ``He must be more perfect and more complete, just as
the revelation of every messenger and prophet has been higher than the
preceding revelation, and the Most Exalted Lord has made the same claim
in relation to His predecessors. This is why each successor has
abrogated what came before.'' End of quotation.

\subsection{321}\label{section-816}

\begin{otherlanguage}{persian}

اوّلا آنكه اگر بصر انصاف مفتوح مشاهده مينمايد كه شبه اين ظهور در ابداع
نيامده و يشهد بذلك ما نزل من جبروت اللّه المهيمن العزيز القدير فافتح
بصراك لتشهد بأنّ جمال الظّهور قد كان حينئذ مستقرّاً على عرش العظمة و
الاستقلال و عن يمينه نقطة البيان بسلطان العزّة و الاجلال و عن يساره محمّد
رسول اللّه بانوار اللّه عزيز المتعال و فى مقابلة الوجه قد قام الرّوح بقبيل
من ملأ الأعلی و نزل بالحقّ ان انتم تفقهون ثمّ عن خلفه صفوفاً من ملائكة
السّمآء باباريق من كوثر البقآء و اكواب من التّسنيم ان انتم تعلمون و كلّهم
ينوحن و يبكينّ و يصيحن على ما ورد على جمال اللّه المهيمن العزيز القيّوم

\end{otherlanguage}

Avvalan ānke agar baṣar-e enṣāf maftūḥ, moshāhede mīnamāyad ke shebh-e
īn ẓohūr dar ebdāʿ nayāmade, va yashhad bi-dhālika mā nazala min jabarūt
Allāh al-Muhaymin al-ʿAzīz al-Qadīr. Fa-ftaḥ baṣarāka li-tashhad bi-anna
Jamāl al-Ẓuhūr qad kāna ḥīnaʾidhin mustaqirran ʿalā ʿarsh al-ʿaẓama
wa-l-istiqlāl, wa-ʿan yamīnihi Nuqṭat al-Bayān bi-sulṭān al-ʿizza
wa-l-ijlāl, wa-ʿan yasārihi Muḥammad Rasūl Allāh bi-anwār Allāh ʿAzīz
al-Mutaʿāl, wa-fī muqābalat al-wajh qad qāma al-Rūḥ bi-qabīl min Malaʾ
al-Aʿlā wa-nazala bi-l-ḥaqq, in antum tafqahūn. Thumma ʿan khalfihi
ṣufūfan min malāʾikat al-samāʾ bi-abārīq min Kawthar al-baqāʾ wa-akwāb
min al-Tasnīm, in antum taʿlamūn. Wa-kulluhum yanūḥna wa-yabkīna
wa-yaṣīḥna ʿalā mā warada ʿalā Jamāl Allāh al-Muhaymin al-ʿAzīz
al-Qayyūm.

First, if the eye of fairness is open, it sees that nothing like this
Revelation has appeared in creation. What has descended from the
dominion of God, the All-Compelling, the Almighty, the All-Powerful,
bears witness to this. Open your eyes, then, to witness that the Beauty
of Revelation is now established upon the throne of grandeur and
independence. At His right is the Point of the Bayán, with the
sovereignty of might and majesty; at His left is Muhammad, the Messenger
of God, with the lights of God, the Mighty, the Most Exalted. Before His
face the Spirit stands with a company of the heavenly concourse and has
descended in truth, if you understand. Behind Him are ranks of the
angels of heaven, with pitchers from the Kawthar of everlasting life and
cups of Tasnim, if you know. They all lament, weep, and cry out over
what has befallen the Beauty of God, the All-Compelling, the Almighty,
the Self-Subsisting.

\subsection{322}\label{section-817}

\begin{otherlanguage}{persian}

اگر صاحب بصر بودى و بمنظر اكبر راجع ميشدى بالمشاهدة و المشافهة انوار
الهى را از مظهر عزّ سلطان لايزالى مشاهده مينمودى و لكن رمد اوهام بشأنى
ابصار را عليل نموده كه اگر بمقرّ اطهر هم وارد شوند گمان نميبرم كه بلقا
فايز گردند چنانچه جعل اگر قرب رضوان مسكن نمايد از عرف اوراد محروم بوده و
خواهد بود و هم‌چنين اعمى اگر در مقابل شمس حاضر شود از انوار او ممنوع بوده
و خواهد بود ليس لهم نصيب من هذا البحر الأعظم

\end{otherlanguage}

Agar ṣāḥeb-e baṣar būdī va be-manẓar-e akbar rājeʿ mīshodī,
bi-l-mushāhada wa-l-mushāfaha anvār-e elāhī rā az maẓhar-e ʿezz-e
solṭān-e lāyazālī moshāhede mīnamūdī, va lāken ramad-e avhām be-shaʾnī
abṣār rā ʿalīl namūde ke agar be-maqarr-e aṭhar ham vāred shavand gomān
nemībaram ke be-leqā fāyez gardand. Chenānche joʿal agar qorb-e Reżvān
maskan namāyad az ʿarf-e avrād maḥrūm būde va khvāhad būd, va ham-chenīn
aʿmā agar dar moqābel-e shams ḥāżer shavad az anvār-e ū mamnūʿ būde va
khvāhad būd. Laysa lahum naṣīb min hādhā al-baḥr al-aʿẓam.

If you had sight and turned toward the Most Great Vision, you would
behold the divine lights directly, face to face, in the Manifestation of
the glory of the everlasting Sovereign. But the inflammation of vain
imaginings has so afflicted their eyes that, even if they entered the
most holy abode, I do not suppose they would attain His presence. Just
as a dung beetle, even if it dwells near Ridvan, has always been and
will remain deprived of the fragrance of its roses, so a blind person,
even if he stands before the sun, has always been and will remain
deprived of its light. They have no share in this most great Ocean.

\subsection{323}\label{section-818}

\begin{otherlanguage}{persian}

اينكه ذكر نمودى هر لاحقى بايد مقدم و اعلى بر سابق باشد اگر مضطرب و
متزلزل مشاهده نميشدى در اينمقام كلمه كه از حقّ استماع نموده‌ام مذكور
ميداشتم و اين مطلبی‌ است كه غير اللّه هر كه در آن تكلّم نمايد از فضولى
محسوب است اللّه يعلم مراده فيما نزل چه كه لئالى معانى مكنونه در صدف ابحر
كلمات الهيّه را جز نفس ظهور مطّلع نبوده و نخواهد بود كلّ ذلك علمه عند ربّی‌ و
ما اطّلع به احد الّا نفسه العالم العليم چون ملأ بيان احجب از متوهّمين قبل
مشاهده ميشوند لذا بهتر آنكه مدارا نمائيم و آنچه را ادراك نموده‌ايم مستور
داريم كه مباد فزع ديگر مرتفع شود قدرت محيطه و رحمت منبسطه الهيّه هرگز
محدود بحدودى نبوده و نخواهد بود يفعل ما یشآء و يحكم ما يريد

\end{otherlanguage}

Īnke dhekr namūdī har lāḥeqī bāyad moqaddam va aʿlā bar sābeq bāshad,
agar możṭareb va motazalzel moshāhede nemīshodī dar īn maqām kalame ke
az Ḥaqq estemāʿ namūde-am madhkūr mīdāshtam, va īn maṭlabī ast ke gheyr
Allāh har ke dar ān takallom namāyad az fożūlī maḥsūb ast. Allāh yaʿlam
murādahu fīmā nazala. Che ke laʾālī-ye maʿānī-ye maknūne dar ṣadaf-e
abḥor-e kalemāt-e elāhiyye rā joz nafs-e ẓohūr moṭṭaleʿ nabūde va
nakhvāhad būd. Kullu dhālika ʿilmuhu ʿinda rabbī wa-mā iṭṭalaʿa bihi
aḥad illā nafsuhu al-ʿĀlim al-ʿAlīm. Chon malaʾ-e Bayān aḥjab az
motavahhemīn-e qabl moshāhede mīshavand, ledhā behtar ānke modārā
namāyīm va ānche rā edrāk namūde-īm mastūr dārīm ke mabād fazaʿ-e dīgar
mortafeʿ shavad. Qodrat-e moḥīṭe va raḥmat-e monbaseṭe-ye elāhiyye
hargez maḥdūd be-ḥodūdī nabūde va nakhvāhad būd. Yafʿal mā yashāʾ
wa-yaḥkum mā yurīd.

As for your statement that every successor must take precedence over and
be higher than his predecessor: if you had not appeared troubled and
shaken, I would have mentioned here a word I heard from the True One.
This is a matter in which anyone other than God who speaks is reckoned
to be meddling. God knows His intention in what He has revealed, for
none but the Manifestation Himself has ever known or will ever know the
pearls of meaning concealed within the shells of the oceans of divine
words. Knowledge of all this is with my Lord, and no one has become
acquainted with it except His own Self, the All-Knowing, the Omniscient.
Since the people of the Bayán appear more veiled than those who clung to
vain imaginings before them, it is better that we show forbearance and
conceal what we have understood, lest another outcry arise. God's
all-encompassing power and all-embracing mercy have never been and will
never be confined within any limits. He does what He wills and decrees
what He desires.

\subsection{324}\label{section-819}

\begin{otherlanguage}{persian}

و در نسخ احكام سابق نقطه بيان روح ما سواه فداه ميفرمايد كلّ منوط باراده
نفس ظهور است اگر بخواهد اوامر و نواهى قبل را امضاء ميفرمايد و اگر بخواهد
نسخ ميفرمايد بيده الأمر ليس لأحد ان يعترض عليه او يقول لم و بم فى كلّ ما
يفعل و يأمر و من قال لم و بم فقد كفر باللّه و حارب بنفسه و جادل بسلطانه و
انكر امره و كان من المشركين

\end{otherlanguage}

Va dar naskh-e aḥkām-e sābeq Noqṭe-ye Bayān, rūḥ-e mā sevāh fedāh,
mīfarmāyad koll manūṭ be-erāde-ye nafs-e ẓohūr ast. Agar bekhvāhad
avāmer va navāhī-ye qabl rā emżā mīfarmāyad va agar bekhvāhad naskh
mīfarmāyad. Bi-yadihi al-amr. Laysa li-aḥad an yaʿtariḍ ʿalayhi aw yaqūl
lima wa-bima fī kulli mā yafʿal wa-yaʾmur. Wa-man qāla lima wa-bima
fa-qad kafara bi-Allāh wa-ḥāraba bi-nafsihi wa-jādala bi-sulṭānihi
wa-ankara amrahu wa-kāna min al-mushrikīn.

Concerning the abrogation of previous laws, the Point of the Bayán---may
the souls of all besides Him be sacrificed for Him---says that
everything depends upon the will of the Manifestation Himself. If He
wishes, He confirms the former commandments and prohibitions, and if He
wishes, He abrogates them. Authority is in His hand. No one has the
right to object to Him or to say, ``Why?'' or ``For what reason?''
concerning anything He does or commands. Whoever says, ``Why?'' or ``For
what reason?'' has disbelieved in God, fought against His own Self,
disputed His sovereignty, denied His Cause, and become one of those who
associate partners with God.

\subsection{325}\label{section-820}

\begin{otherlanguage}{persian}

قول ديگر اگر بطريق استعكاس است مثل حضرات ائمّه اطهار بالنسبة بحضرت محمّد
كه ادّعاى مثليّت بآنحضرت مينمودند و حال آنكه در ظلّ او بودند پس مطلب واضح
است چه چنين شخصى را نميسزد كه ادّعاى برترى در چنين صورتى نمايد انتهى

\end{otherlanguage}

Qowl-e dīgar: Agar be-ṭarīq-e esteʿkās ast, mesl-e ḥażarāt-e aʾemme-ye
aṭhār bi-l-nisba be-Ḥażrat-e Moḥammad ke eddeʿā-ye methliyyat be-ān
Ḥażrat mīnamūdand va ḥāl ānke dar ẓell-e ū būdand, pas maṭlab vāżeḥ ast,
che chenīn shakhṣī rā nemīsazad ke eddeʿā-ye bartarī dar chenīn ṣūratī
namāyad. Intahā.

Another statement: ``If it is by way of reflection, as with the holy and
pure Imams in relation to Muhammad, who claimed likeness to Him while
remaining beneath His shadow, then the matter is clear, for it does not
befit such a person to claim superiority under such circumstances.'' End
of quotation.

\subsection{326}\label{section-821}

\begin{otherlanguage}{persian}

در بيان قبل ذكر اينمقام شده ديگر اين كلمات آن مشرك ملقى كه احقر از كلمات
صبيان است بی‌جواب بوده و خواهد بود اى مرد عكس چه عاكس چه اين ساحت مقدّس از
ذكر و بيان است و منزّه از ادراك من فى الامكان

\end{otherlanguage}

Dar bayān-e qabl dhekr-e īn maqām shode. Dīgar īn kalemāt-e ān moshrek-e
molqī ke aḥqar az kalemāt-e ṣebyān ast bī-javāb būde va khvāhad būd. Ey
mard, ʿaks che, ʿākes che? Īn sāḥat moqaddas az dhekr va bayān ast va
monazzah az edrāk-e man fī al-imkān.

This matter has been addressed in the preceding explanation. These
further words of that prompter who associates partners with God, which
are more contemptible than the words of children, have been and will
remain unanswered. O man, what reflection? What reflector? This court is
sanctified above mention and description and exalted beyond the
understanding of all who dwell in the contingent world.

\subsection{327}\label{section-822}

\begin{otherlanguage}{persian}

در اينمقام حكايتى بنظر آمد ایّاميكه نقطه اولى جلّ و عزّ در جبل ماكو مسجون
بودند چنانچه حال هم در اين ارض مسجونند و لكن اين سجن را بآن سجن قياس مكن
چه كه حال هم مسجون دولتند و هم مسجون اهل بيان و سجن ظهور قبل مخصوص بدولت
بوده و اگر خوب ملاحظه كنى در اين سجن كلّ دول هم شريك

\end{otherlanguage}

Dar īn maqām ḥekāyatī be-naẓar āmad. Ayyāmī ke Noqṭe-ye Ūlā, jalla
wa-ʿazza, dar Jabal-e Mākū masjūn būdand, chenānche ḥāl ham dar īn arż
masjūnand, va lāken īn sejn rā be-ān sejn qiyās makon, che ke ḥāl ham
masjūn-e dowlatand va ham masjūn-e ahl-e Bayān, va sejn-e ẓohūr-e qabl
makhṣūṣ be-dowlat būde, va agar khūb molāḥeẓe konī dar īn sejn koll-e
doval ham sharīk.

At this point a story came to mind. In the days when the Primal Point,
glorified and exalted be He, was imprisoned in the mountain of Maku---as
He is now imprisoned in this land, though do not compare this
imprisonment with that one, for now He is imprisoned both by the
government and by the people of the Bayán, whereas the imprisonment of
the former Manifestation was solely at the hands of the government; and
if you consider carefully, all governments also share in this
imprisonment.

\subsection{328}\label{section-823}

\begin{otherlanguage}{persian}

بارى در آن ایّام جناب وحيد يعنى آقا سيّد يحيى عليه بهآء اللّه بنصرت امر علم
برافراخته بودند چنانچه شنيده‌ايد و در كلّ ديار علما بر اعراض و انكار
برخواسته و بر كلّ منابر نعوذاً باللّه بسبّ و لعن نقطه اولى و اصحابش مشغول
امام جمعه طهران اگر چه فى‌الحقيقة او مبغض نبود و انصافش هم از جميع برترى
داشته‌ روزى بر منبر رفته مذكور نمود كه از جمله شنيده‌ام حضرت باب هوا را
مطهّر دانسته آخر بگوئيد اى مرد در كتاب كدام فقها و علماء چنين كلمه مشاهده
نموده كه هوا مطهّر است و آنقدر ملتفت نشده بود كه مقام اجتهاد نيست مقامى
است ماوراى اين اذكار و كلمات و خود بنفسه مجدّد شرايع بوده و همچه ادراك
نموده بود كه يكى از فقها هستند و اين مسئله فقهيّه را مخالفاً للعلمآء ذكر
نموده‌اند و جايز دانسته‌اند

\end{otherlanguage}

Bārī dar ān ayyām Janāb-e Vaḥīd, yaʿnī Āqā Seyyed Yaḥyā, ʿalayhi Bahāʾ
Allāh, be-noṣrat-e Amr ʿalam bar afrākhte būdand, chenānche shenīde-īd,
va dar koll-e deyār ʿolamā bar eʿrāż va enkār bar khāste va bar koll-e
manāber, naʿūdhan bi-Allāh, be-sabb va laʿn-e Noqṭe-ye Ūlā va aṣḥābash
mashghūl. Emām-e Jomʿe-ye Ṭehrān, agarche fī al-ḥaqīqa ū mobgheż nabūd
va enṣāfash ham az jamīʿ bartarī dāshte, rūzī bar menbar rafte madhkūr
namūd ke az jomle shenīde-am Ḥażrat-e Bāb havā rā moṭahher dāneste.
Ākhar begūʾīd ey mard, dar ketāb-e kodām foqahā va ʿolamā chenīn kalame
moshāhede namūde ke havā moṭahher ast? Va ānqadr moltafet nashode būd ke
maqām-e ejtehād nīst, maqāmī ast māvarā-ye īn adhkār va kalemāt, va khod
be-nafsehe mojadded-e sharāyeʿ būde, va hamche edrāk namūde būd ke yekī
az foqahā hastand va īn masʾale-ye feqhiyye rā mukhālifan li-l-ʿulamāʾ
dhekr namūde-and va jāyez dāneste-and.

In those days, then, the honored Vahid, that is, Aqa Siyyid Yahya---upon
him be the glory of God---had raised the banner in support of the Cause,
as you have heard. In every land the religious scholars had risen in
rejection and denial, and from every pulpit they were occupied---God
forbid---with reviling and cursing the Primal Point and His companions.
The Imam-Jumih of Tehran, although in truth he bore no hatred and was
fairer than all the others, one day ascended the pulpit and said,
``Among other things, I have heard that the Báb has declared air to be a
purifier. Now tell me, O man, in the book of which jurist or scholar has
He seen such a statement, that air is a purifier?'' He had not realized
even this much: that this was not the station of independent legal
reasoning, but a station beyond such discussions and words, and that He
Himself was the renewer of religious laws. Instead, he had understood
Him to be one of the jurists who had stated this point of religious law
and deemed it permissible in opposition to the scholars.

\subsection{329}\label{section-824}

\begin{otherlanguage}{persian}

و حال شما هم نميدانيد چه نقل است و چه امر از افق اراده الهى ظاهر شده
مع‌ذلك باين كلمات كه لايق انفس ملقيه بوده و هست و شما هم بالقاى آنها
نوشته‌ايد افتخار مينمائيد و البتّه باغواى شيطان باطراف هم فرستاده‌ايد
چنانچه مخصوص اخبار فرموده‌اند و انّه لهو الصّادق العليم فسوف يبطل اللّه ما
عندهم بقوله الحقّ و يظهر مكر نفوسهم لعباده المقرّبين

\end{otherlanguage}

Va ḥāl shomā ham nemīdānīd che naql ast va che amr az ofoq-e erāde-ye
elāhī ẓāher shode. Maʿa dhālika be-īn kalemāt ke lāyeq-e anfos-e molqiye
būde va hast va shomā ham be-elqā-ye ānhā neveshte-īd eftekhār
mīnamāyīd, va albatte be-eghvā-ye sheyṭān be-aṭrāf ham ferestāde-īd,
chenānche makhṣūṣ ekhbār farmūde-and. Wa-innahu la-huwa al-Ṣādiq
al-ʿAlīm. Fa-sawfa yubṭil Allāh mā ʿindahum bi-qawlihi al-ḥaqq wa-yuẓhir
makr nufūsihim li-ʿibādihi al-muqarrabīn.

And now you too do not know what is happening or what Cause has appeared
from the horizon of the divine will. Nevertheless, you take pride in
these words, which have always befitted those who prompted them and
which you too have written at their prompting. Surely, at Satan's
instigation, you have also sent them far and wide, as He has
specifically informed us. He is indeed the Truthful, the All-Knowing.
Soon God will nullify what they possess through His true word and expose
the scheming of their souls to His servants who are near to Him.

\subsection{330}\label{section-825}

\begin{otherlanguage}{persian}

شأن شما آن بود كه اوّلا سؤال نمائيد از نفس ظهور كه ادّعا چيست و حجّت چه تا
بر آن مطّلع شويد اگر ميفرمود كه اين همان ظهورى است كه نقطه بيان بآن بشارت
داده و اخذ عهد او را از كلّ نموده ديگر اين اعتراضات جميع بی‌ثمر بوده و
خواهد بود و اين همه زحمت بخود نميداديد كه حروفات جعليّه آن انفس ملقيه را
تركيب نمائيد اتّق اللّه و لاتقل ما لا تعلم و لا تتبّع كلّ ناعق و لو تريد ان
تطّلع بالحقّ فاسئل اللّه ربّك و ربّ العالمين

\end{otherlanguage}

Shaʾn-e shomā ān būd ke avvalan soʾāl namāyīd az nafs-e ẓohūr ke eddeʿā
chīst va ḥojjat che tā bar ān moṭṭaleʿ shavīd. Agar mīfarmūd ke īn hamān
ẓohūrī ast ke Noqṭe-ye Bayān be-ān beshārat dāde va akhdh-e ʿahd-e ū rā
az koll namūde, dīgar īn eʿterāżāt jamīʿ bī-thamar būde va khvāhad būd,
va īn hame zaḥmat be-khod nemīdādīd ke ḥorūfāt-e jaʿliyye-ye ān anfos-e
molqiye rā tarkīb namāyīd. Ittaqi Allāh wa-lā taqul mā lā taʿlam wa-lā
tattabiʿ kulla nāʿiq. Wa-law turīd an taṭṭaliʿ bi-l-ḥaqq fa-sʾal Allāh
rabbaka wa-rabb al-ʿālamīn.

What befitted you was first to ask the Manifestation Himself what His
claim was and what His proof was, so that you might become acquainted
with it. If He said that this was the very Revelation whose advent the
Point of the Bayán had announced and for whom He had taken a covenant
from everyone, then all these objections were and would remain
fruitless. You would not have put yourself to so much trouble to piece
together the fabricated words of those who prompted you. Fear God, do
not say what you do not know, and do not follow everyone who cries out.
If you wish to become acquainted with the truth, ask God, your Lord and
the Lord of the worlds.

\subsection{331}\label{section-826}

\begin{otherlanguage}{persian}

قول ديگر كبر غفلة من القاه غايت امر اين است كه يكى از ادلّاء و مراياى او
ميباشند و مرايا را مراتب است و حاكى بالتّمام از نسل يكى خواهد بود رشته
مرآت تام واحدا بعد واحد خواهد بود باقى ديگر حاكى بالتّمام نخواهد بود و آن
را هم حضرت ربّ اعلى حصر در ميرزا يحيى نموده انتهى

\end{otherlanguage}

Qowl-e dīgar, kabura ghaflat man alqāhu: Ghāyat-e amr īn ast ke yekī az
adellāʾ va marāyā-ye ū mībāshand, va marāyā rā marāteb ast, va ḥākī
be-l-tamām az nasl-e yekī khvāhad būd. Reshte-ye merʾāt-e tāmm vāḥedan
baʿd-e vāḥed khvāhad būd. Bāqī dīgar ḥākī be-l-tamām nakhvāhad būd, va
ān rā ham Ḥażrat-e Rabb-e Aʿlā ḥaṣr dar Mīrzā Yaḥyā namūde. Intahā.

Another statement---how great the heedlessness of the one who prompted
it: ``At most, He is one of His guides and mirrors. The mirrors have
ranks, and the one who fully reflects will be from the lineage of one.
The succession of the perfect mirror will be one after another. The
others will not fully reflect, and the Most Exalted Lord has confined
this too to Mirza Yahya.'' End of quotation.

\subsection{332}\label{section-827}

\begin{otherlanguage}{persian}

قد اخطأ من القاك و عصى فيما ارتكب و كفر باللّه الّذى خلقه من تراب و افترى
بالّذى اليه يرجع فى منقلبه و مثوئه و هو من الّذين جعل الكتاب عضين چه كه
تازه مرآت تام و غير تام ذكر نموده واللّه قلم متحيّر است كه در جواب امثال
آن غافلين چه ذكر نمايد خود نقطه اولى در مقام مناجات ميفرمايد الهى در هر
سنه بفرست مرآتى تا آنكه مدلّ باشد بر تو و حاكى باشد از تو و بعد ميفرمايد
در هر شهرى بفرست و اگر تسليم نمائيم قول شما را مخالف است با قول نقطه
بيان منتهى رتبه مرايا اين رتبه است كه ميفرمايد قوله جلّ اجلاله قل ان يا
شموس المرايا انتم الى شمس الحقيقة تنظرون فانّ قيامكم بها لو انتم تتبصّرون
كلكم كحيتان بالمآء فى البحر تتحرّكون و تحتجبون عن الماء و تسئلون عما انتم
به قائمون

\end{otherlanguage}

Qad akhṭaʾa man alqāka wa-ʿaṣā fīmā irtakaba wa-kafara bi-Allāh alladhī
khalaqahu min turāb wa-ftarā bi-alladhī ilayhi yarjiʿ fī munqalabihi
wa-mathwāhu wa-huwa min alladhīna jaʿala al-kitāb ʿiḍīn. Che ke tāze
merʾāt-e tāmm va gheyr-e tāmm dhekr namūde. Wa-Allāh qalam motaḥayyer
ast ke dar javāb-e amthāl-e ān ghāfelīn che dhekr namāyad. Khod-e
Noqṭe-ye Ūlā dar maqām-e monājāt mīfarmāyad: Elāhī dar har sane beferest
merʾātī tā ānke modell bāshad bar to va ḥākī bāshad az to. Va baʿd
mīfarmāyad dar har shahrī beferest. Va agar taslīm namāyīm qowl-e shomā
rā mokhālef ast bā qowl-e Noqṭe-ye Bayān. Montahā rotbe-ye marāyā īn
rotbe ast ke mīfarmāyad, qawluhu jalla ijlāluhu: Qul in yā shumūs
al-marāyā antum ilā Shams al-Ḥaqīqa tanẓurūn, fa-inna qiyāmakum bihā law
antum tatabaṣṣarūn. Kullukum ka-ḥītān bi-l-māʾ fī al-baḥr tataḥarrakūn
wa-taḥtajibūn ʿan al-māʾ wa-tasʾalūn ʿammā antum bihi qāʾimūn.

The one who prompted you has erred, rebelled in what he has perpetrated,
disbelieved in God who created him from dust, and slandered the One to
whom he returns in his final destination and abode. He is among those
who have divided the Book into fragments, for he has now introduced the
notion of a perfect and an imperfect mirror. By God, the Pen is
bewildered as to what it should say in answer to such heedless people.
The Primal Point Himself says in prayer, ``My God, send in every year a
mirror, that it may point to You and reflect You.'' Then He says, ``Send
one in every city.'' Even if we were to accept your statement, it
contradicts the words of the Point of the Bayán. The highest rank of the
mirrors is this rank of which He speaks---His words, exalted be His
majesty: ``Say: O suns of the mirrors, you look toward the Sun of Truth,
for your existence depends upon it, if you perceive. You are all like
fish moving through the water in the sea, yet you are veiled from the
water and ask about that by which you subsist.''

\subsection{333}\label{section-828}

\begin{otherlanguage}{persian}

بگو اى غافل ملاحظه نما كه بشموس مرايا كه منتهى رتبه آنها است ميفرمايد كه
شما بشمس حقيقت ناظر باشيد و بعد ميفرمايند مثل شما مثل حيتان است كه در
بحر حركت مينمائيد و محتجبيد از ماء و سؤال كرده ميشويد از آنچه باو قائميد

\end{otherlanguage}

Begū ey ghāfel, molāḥeẓe namā ke be-shomūs-e marāyā ke montahā rotbe-ye
ānhā ast mīfarmāyad ke shomā be-Shams-e Ḥaqīqat nāẓer bāshīd, va baʿd
mīfarmāyand mathal-e shomā mathal-e ḥītān ast ke dar baḥr ḥarakat
mīnamāyīd va moḥtajebīd az māʾ va soʾāl karde mīshavīd az ānche be-ū
qāʾemīd.

Say: O heedless one, consider how He tells the suns of the mirrors,
which is their highest rank, to look toward the Sun of Truth. Then He
says that you are like fish moving in the sea, veiled from the water,
and questioned about that by which you subsist.

\subsection{334}\label{section-829}

\begin{otherlanguage}{persian}

فواللّه اگر اهل بيان انصاف دهند همين بيان منزل بيان جميع را كافيست از
كلمات مشركين منقطع ميشوند و معنى اين كلمات الهيّه را اليوم مشاهده
مينمايند چنانچه اين نفس معرض كه خود را مرآت ميداند و بعضى هم من غير شعور
او را مرآت فرض گرفته بر فرض تسليم از شمس حقيقت محروم مانده چنانچه مشاهده
ميشود كه بنفس خود ناظر و بموهوم متشبّث و از شمس حقيقت بعيد و ممنوع در ماء
حركت ميكند چنانچه در ظلّ اشراقات انوار الهيّه مستظلّ و لكن از او مهجور و
محجوب و محروم و هم‌‌چنين اليوم از كلّ مرايا سؤال كرده ميشود چنانچه كلّ در
موقف سؤال حاضر هر كدام از شريعه عرفان اللّه خارج شد از حجر محسوب بلكه
قابل ذكر نبوده و هر كدام عارف شد و در ظلّ عنايتش مأوى گرفت از مراياى
اوّلیّه عنداللّه مذكور كذلك فصل الأمر لعلّ النّاس هم يفقهون

\end{otherlanguage}

Fa-wa-Allāh agar ahl-e Bayān enṣāf dehand hamīn bayān-e Monzel-e Bayān
jamīʿ rā kāfīst. Az kalemāt-e moshrekīn monqaṭeʿ mīshavand va maʿnā-ye
īn kalemāt-e elāhiyye rā al-yawm moshāhede mīnamāyand. Chenānche īn
nafs-e moʿreż ke khod rā merʾāt mīdānad va baʿżī ham min ghayri shuʿūr ū
rā merʾāt farż gerefte, bar farż-e taslīm, az Shams-e Ḥaqīqat maḥrūm
mānde. Chenānche moshāhede mīshavad ke be-nafs-e khod nāẓer va be-mowhūm
motashabbeth va az Shams-e Ḥaqīqat baʿīd va mamnūʿ. Dar māʾ ḥarakat
mīkonad, chenānche dar ẓell-e eshrāqāt-e anvār-e elāhiyye mostaẓell va
lāken az ū mahjūr va maḥjūb va maḥrūm. Va ham-chenīn al-yawm az koll-e
marāyā soʾāl karde mīshavad, chenānche koll dar mowqef-e soʾāl ḥāżer.
Har kodām az sharīʿe-ye ʿerfān Allāh khārej shod az ḥajar maḥsūb, balke
qābel-e dhekr nabūde, va har kodām ʿāref shod va dar ẓell-e ʿenāyatash
maʾvā gereft az marāyā-ye avvaliyye ʿend Allāh madhkūr. Kadhālika
fuṣṣila al-amr laʿalla al-nās hum yafqahūn.

By God, if the people of the Bayán were fair, this utterance of the
Revealer of the Bayán alone would suffice them all. They would detach
themselves from the words of those who associate partners with God and
would witness today the meaning of these divine words. Thus this soul
who has turned away, who considers himself a mirror and whom some,
without understanding, have assumed to be a mirror---even granting this
assumption---has remained deprived of the Sun of Truth. As can be seen,
he looks to his own self, clings to a vain imagining, and remains
distant from and denied access to the Sun of Truth. He moves in the
water, inasmuch as he shelters beneath the radiance of the divine
lights, yet he is estranged, veiled, and deprived of Him. Likewise,
today all the mirrors are questioned, for all stand at the place of
questioning. Whichever departs from the water's edge of the recognition
of God is reckoned as stone---indeed, is unworthy of mention---and
whichever recognizes Him and takes refuge beneath the shadow of His
loving care is numbered among the foremost mirrors in the sight of God.
Thus has the matter been set forth in detail, so that people may
understand.

\subsection{335}\label{section-830}

\begin{otherlanguage}{persian}

حال قدرى در اين بيان تفكّر نما كه شايد بكوثر احديّه كه از معين كلمات الهيّه
جارى است فايز گردى بگو اى ظالم بی‌‌خبر اين تحديدات در كجاى از بيان مسطور
است كه القاء نموده افترا بر حقّ بسته كه مرآت حاكى بالتمام از نسل يكى
خواهد بود و آن را هم حضرت ربّ اعلى در فلان حصر نموده قل افتريت بربّ العرش
و الكرسىّ الرّفيع از اين كلمات حقيقت عرفان من القى اليك و اوحى اليه معلوم
شده آفرين خوب يد منبسطه الهيّه را مغلول نموده‌ايد يتكلّمون بمثل اليهود و لا
يشعرون خاك بر فم امثال آن قائلين و طين بر رأس نابالغين كه تنزيه امر و
تقديس كلمات الهيّه را بغبار جهل و تحديد آلوده نموده‌اند لم يزل امراللّه
مقدّس از اين حدود بوده

\end{otherlanguage}

Ḥāl qadrī dar īn bayān tafakkur namā keh shāyad be-kawthar-i aḥadiyyeh
keh az maʿīn-i kalimāt-i ilāhiyyeh jārī ast fāʾiz gardī. Begū ey ẓālim-i
bī-khabar īn taḥdīdāt dar kojā-yi az Bayān masṭūr ast keh ilqā namūdeh
iftirā bar Ḥaqq basteh keh mirʾāt-i ḥākī bi-t-tamām az nasl-i yekī
khwāhad būd va ān rā ham Ḥaḍrat-i Rabb-i Aʿlā dar fulān ḥaṣr namūdeh.
Qul iftarayta bi-Rabbi l-ʿarshi wa-l-kursiyyi r-rafīʿ. Az īn kalimāt
ḥaqīqat-i ʿirfān-i man alqā ilayka wa-awḥā ilayhi maʿlūm shodeh. Āfarīn
khūb yad-i munbasiṭeh-yi ilāhiyyeh rā maghlūl namūdeh-īd. Yatakallamūna
bi-mithli l-Yahūdi wa-lā yashʿurūn. Khāk bar fam-i amthāl-i ān qāʾilīn
va ṭīn bar raʾs-i nābālighīn keh tanzīh-i amr va taqdīs-i kalimāt-i
ilāhiyyeh rā be-ghubār-i jahl va taḥdīd ālūdeh namūdeh-and. Lam yazal
Amruʾllāh muqaddas az īn ḥudūd būdeh.

Now reflect a little on this utterance, so that you may perhaps attain
the Kawthar of oneness flowing from the wellspring of the divine words.
Say: O uninformed oppressor, where in the Bayán are these restrictions
recorded, which you have put forward, slandering God by claiming that
the mirror that fully reflects will come from the descendants of one
person, and that the Most Exalted Lord has further confined this to a
certain individual? Say: You have fabricated a lie against the Lord of
the throne and the exalted seat. These words reveal the true measure of
the understanding of the one who prompted you and to whom inspiration
was given. Well done! How thoroughly you have fettered the outstretched
hand of God! They speak like the Jews and do not perceive it. Dust upon
the mouths of such speakers, and clay upon the heads of the immature,
who have sullied the transcendence of the Cause and the sanctity of the
divine words with the dust of ignorance and restriction. The Cause of
God has always been sanctified above these limitations.

\subsection{336}\label{section-831}

\begin{otherlanguage}{persian}

واللّه الّذى لا اله الّا هو كه از پست‌ترين اهل ملل قبل محتجب‌تر مشاهده ميشوند
و لكن مرشد بسيار از اين كلمه مسرور شده و البتّه شما را هم نظر باين خدمت
يكى از شهداى نار قرار داده و يا خواهد داد چه كه هميشه اراده داشته كه
امراللّه را مخصوص بخود و ما يخرج من صلبه نمايد و شما پيش از اراده او امضا
داشته‌ايد مريد يعنى شما صدهزار آفرين بر شما كه از براى يك مخلوق خبيثى كه
واللّه الّذى لا اله الّا هو كه شيطان از افعالش احتراز جسته و ميجويد از خالق
گذشته افّ له و لمن اتّبعه فواللّه بافترائه اهتزّت اركان العرش و تزلزلت قوائم
كرسىّ الرّفيع

\end{otherlanguage}

Wa-llāhi lladhī lā ilāha illā Huwa keh az past-tarīn-i ahl-i milal-i
qabl muḥtajib-tar mushāhedeh mī-shavand va lākin murshid bisyār az īn
kalimeh masrūr shodeh va al-batteh shomā rā ham naẓar be-īn khidmat yekī
az shuhadā-yi nār qarār dādeh va yā khwāhad dād cheh keh hamīsheh irādeh
dāshteh keh Amruʾllāh rā makhṣūṣ be-khod va mā yakhruju min ṣulbihi
namāyad va shomā pīsh az irādeh-yi ū imḍā dāshteh-īd. Murīd yaʿnī shomā
ṣad-hizār āfarīn bar shomā keh az barā-yi yek makhlūq-i khabīthī keh
wa-llāhi lladhī lā ilāha illā Huwa keh shayṭān az afʿālash iḥtirāz
jasteh va mī-jūyad az Khāliq gudhashteh. Uffin lahu wa-li-man ittabaʿahu
fa-wa-llāhi bi-iftirāʾihi ihtazzat arkānu l-ʿarshi wa-tazalzalat
qawāʾimu kursiyyi r-rafīʿ.

By God, besides whom there is no other god, they are seen to be more
veiled than the lowliest followers of earlier religions. Yet the master
has been greatly pleased by this statement, and in consideration of this
service he has certainly appointed you, or will appoint you, as one of
the witnesses of the fire. For he has always intended to confine the
Cause of God to himself and what proceeds from his loins, and you have
endorsed this even before he intended it. A disciple---that is you! A
hundred thousand congratulations to you for forsaking the Creator for
the sake of a vile creature whose deeds, by God, besides whom there is
no other god, Satan has shunned and continues to shun. Shame on him and
on whoever follows him! By God, at his fabrication the pillars of the
throne shook and the supports of the exalted seat trembled.

\subsection{337}\label{section-832}

\begin{otherlanguage}{persian}

بشنو نداء اين عبد را و اين دفتر نفس و هوا را كه بالقاى شيطان بكلمات
ظنونيّه كاذبه مفتريه نقش نموده بشوى و از مفازه تنگ بفضاى وسيع منيع وارد
شو نديده مگو و نفهميده منويس عنقريب تو و او و ما همه بخاك راجعيم اى با
انصاف باين شدّت چرا در تضييع امراللّه كمر بسته و براى رضاى مخلوقى بی‌قدر و
بی‌‌شأن اين همه مفتريات بحقّ نسبت داده قد اخذتكم نفحات القهر و اى قهر اعظم
من ذلك تعبدون البقر و لا تعرفون ثمّ تدعون اللّه بأن يخرج لكم من صلبه عجلاً
لتعبدوه و تكوننّ من العابدين ثم من نسل هذا العجل عجلاً آخر و لو ينقطع
النّسل لتعزون فى انفسكم و تنوحون و تبكون كما تبكون فى فقدان آبائكم و
امّهاتكم بل اعظم لو تكوننّ من الشّاعرين

\end{otherlanguage}

Bishnav nidā-yi īn ʿabd rā va īn daftar-i nafs va havā rā keh be-ilqā-yi
shayṭān be-kalimāt-i ẓunūniyyeh-yi kādhibeh-yi muftariyeh naqsh namūdeh
bishūy va az mafāzeh-yi tang be-faḍā-yi vasīʿ-i manīʿ vārid shaw.
Nadīdeh magū va nafahmīdeh manavīs. ʿAn-qarīb to va ū va mā hameh
be-khāk rājiʿīm. Ey bā-inṣāf be-īn shiddat cherā dar taḍyīʿ-i Amruʾllāh
kamar basteh va barā-yi riḍā-yi makhlūqī bī-qadr va bī-shaʾn īn hameh
muftariyāt be-Ḥaqq nisbat dādeh. Qad akhadhatkum nafaḥātu l-qahri
wa-ayyu qahrin aʿẓamu min dhālika taʿbudūna l-baqara wa-lā taʿrifūna
thumma tadʿūna Allāha bi-an yukhrija lakum min ṣulbihi ʿijlan
li-taʿbudūhu wa-takūnanna mina l-ʿābidīn. Thumma min nasli hādhā l-ʿijli
ʿijlan ākhara wa-law yanqaṭiʿu n-naslu la-taʿzawna fī anfusikum
wa-tanūḥūna wa-tabkūna kamā tabkūna fī fuqdāni ābāʾikum wa-ummahātikum
bal aʿẓama law takūnanna mina sh-shāʿirīn.

Hear the call of this servant, and wash clean this book of self and
desire that, at Satan's prompting, you have inscribed with conjectural,
lying, slanderous words. Leave this narrow wasteland and enter a vast,
lofty expanse. Do not speak without seeing, and do not write without
understanding. Soon you and he and we will all return to dust. O
fair-minded one, why have you set yourself so fiercely to undermining
the Cause of God, and attributed all these fabrications to God to please
a worthless, insignificant creature? The blasts of wrath have seized
you. What wrath could be greater than this: you worship the cow and do
not know it, then implore God to bring forth a calf for you from his
loins, so that you may worship it and be among its worshippers; then,
from the descendants of this calf, another calf. Should the line of
descent be cut off, you would console yourselves, lament, and weep as
you weep at the loss of your fathers and mothers---indeed, even more, if
you were among those who perceive.

\subsection{338}\label{section-833}

\begin{otherlanguage}{persian}

ملاحظه نما كه چه قدر غافلند اين نفوس كه شمس معانى با منتهى علوّ عظمت و
جلال از افق حقيقت طالع و مشرق شده او را بقدر ما يظهر من صلب البقر هم
ندانسته‌اند بگو لعنت بر شما و بر حياى شما قد شقّ هيكل الأحديّه رداء العظمة
من ظلمكم يا ملأ العنود

\end{otherlanguage}

Mulāḥiẓeh namā keh cheh qadr ghāfil-and īn nufūs keh Shams-i Maʿānī bā
muntahā-yi ʿuluvv-i ʿaẓamat va jalāl az ufuq-i ḥaqīqat ṭāliʿ va mushriq
shodeh ū rā be-qadr-i mā yaẓharu min ṣulbi l-baqari ham nadānisteh-and.
Begū laʿnat bar shomā va bar ḥayā-yi shomā. Qad shaqqa Haykalu
l-Aḥadiyyati ridāʾa l-ʿaẓamati min ẓulmikum yā malaʾa l-ʿanūd.

Consider how heedless these souls are: the Sun of meanings has risen and
shone from the horizon of reality in the utmost sublimity of grandeur
and glory, yet they have not esteemed Him even as much as what comes
forth from the loins of the cow. Say: A curse upon you and upon your
sense of shame! The Temple of Oneness has torn the robe of grandeur
because of your oppression, O company of the obstinate.

\subsection{339}\label{section-834}

\begin{otherlanguage}{persian}

حيرت از اين است كه سيّد محمّد اصفهانى كه ابداً از اين امر اطّلاع نداشته و با
او نبوده شما او را صادق فرض گرفته‌ايد و جميع نفوسيكه مطّلع بوده‌اند كاذب
دانسته‌ايد و باغواى او شمشير كين بر جمال ربّ العالمين آخته بريز اين وساوس
نفسانيّه و خدع نسناسيّه را قسم بآفتاب معانى كه طائفين حول ننگ دارند از ذكر
اين مراياى كاذبه شما و آنقدر شما ادراك ننموده‌ايد كه وجود مرايا بنفسه
لنفسه نبوده بلكه در حين تقابل اشراقى از شمس در آنها ظاهر و بعد از انحراف
معدوم بوده و خواهند بود صدهزار مرايا بكلمه خلق شده و ميشوند وتا در ظلّ
شجره اثبات و شمس حقيقت باقيند از مرايا مذكور و الّا از احجار سجّين و اهل
نار مذكور

\end{otherlanguage}

Ḥayrat az īn ast keh Siyyid Muḥammad-i Iṣfahānī keh abadan az īn amr
iṭṭilāʿ nadāshteh va bā ū nabūdeh shomā ū rā ṣādiq farḍ girifteh-īd va
jamīʿ-i nufūsī-keh muṭṭaliʿ būdeh-and kādhib dānisteh-īd va be-ighvā-yi
ū shamshīr-i kīn bar Jamāl-i Rabbi l-ʿĀlamīn ākhteh. Birīz īn vasāvis-i
nafsāniyyeh va khudaʿ-i nasnāsiyyeh rā. Qasam be-Āftāb-i Maʿānī keh
ṭāʾifīn-i ḥawl nang dārand az dhikr-i īn marāyā-yi kādhibeh-yi shomā va
ān-qadr shomā idrāk nanamūdeh-īd keh vujūd-i marāyā bi-nafsihi
li-nafsihi nabūdeh balkeh dar ḥīn-i taqābul ishrāqī az shams dar ānhā
ẓāhir va baʿd az inḥirāf maʿdūm būdeh va khwāhand būd. Ṣad-hizār marāyā
be-kalimeh khalq shodeh va mī-shavand va tā dar ẓill-i Shajareh-yi
Ithbāt va Shams-i Ḥaqīqat bāqī-and az marāyā madhkūr va illā az aḥjār-i
sijjīn va ahl-i nār madhkūr.

What is astonishing is that you have assumed Siyyid Muḥammad of Isfahan
to be truthful, though he had no knowledge whatsoever of this matter and
was not with Him, while you have deemed all the souls who were informed
to be liars. At his instigation you have drawn the sword of hatred
against the Beauty of the Lord of the worlds. Cast away these selfish
promptings and subhuman deceits. By the Sun of meanings, those who
circle around Him are ashamed to mention these false mirrors of yours.
You have not understood even this much: the existence of the mirrors has
not been of itself and for itself. Rather, when they face the sun, a
radiance from it appears in them; after they turn away, it has vanished
and will vanish. A hundred thousand mirrors have been and are created by
a word. So long as they remain beneath the shadow of the Tree of
Affirmation and the Sun of Truth, they are reckoned as mirrors;
otherwise, they are reckoned among the stones of Sijjin and the people
of the fire.

\subsection{340}\label{section-835}

\begin{otherlanguage}{persian}

قول ديگر مسلّم ميداريم كه بر ربّ اعلى برترى دارد و قيامت او است و دوره او
مرتفع شده باينمعنى كه من همان من يظهره اللّه هستم كه وعده خود حضرت ربّ
اعلى است كه دوره بيان باو مرتفع ميشود و ناسخ خواهد بود بطلان اين سخن
اظهر من الشّمس است اوّلا خود حضرت باب بنوعى كه احدى خلاف نكند بيان يكى از
آن بيانات مجملاً اين عبارت است لو لم يكمل خلق ذلك الظّهور لم يظهر اللّه
مظهر الآخر و معنى تكميل را نيز بيان فرموده اينست كه آنچه فرموده و امر
نموده چونكه من اللّه بايد تخلّف نكند و خلقى در ظلّ او خلق شود و كلّ فرمايشات
او چه از مقامات عاليه و چه از فروعات دينيّه بظهور آيد و الّا ظهور چنين
شخصى لغو خواهد بود و اخبارى كه از حقّ است كلّ كذب خواهد بود انتهى

\end{otherlanguage}

Qawl-i dīgar: musallam mī-dārīm keh bar Rabb-i Aʿlā bartarī dārad va
qiyāmat-i ū ast va dawreh-yi ū murtafiʿ shodeh be-īn maʿnī keh man hamān
Man Yuẓhiruhuʾllāh hastam keh vaʿdeh-yi khod-i Ḥaḍrat-i Rabb-i Aʿlā ast
keh dawreh-yi Bayān be-ū murtafiʿ mī-shavad va nāsikh khwāhad būd.
Buṭlān-i īn sukhan aẓharu mina sh-shams ast. Avvalan khod-i Ḥaḍrat-i Bāb
be-nawʿī keh aḥadī khilāf nakonad bayān; yekī az ān bayānāt mujmalan īn
ʿibārat ast: law lam yukmal khalqu dhālika ẓ-ẓuhūri lam yuẓhiri Allāhu
maẓhara l-ākhar. Va maʿnī-yi takmīl rā nīz bayān farmūdeh īn ast keh
āncheh farmūdeh va amr namūdeh chūn-keh min Allāh bāyad takhalluf
nakonad va khalqī dar ẓill-i ū khalq shavad va kull-i farmāyishāt-i ū
cheh az maqāmāt-i ʿāliyeh va cheh az furūʿāt-i dīniyyeh be-ẓuhūr āyad va
illā ẓuhūr-i chunīn shakhsī laghv khwāhad būd va akhbārī keh az Ḥaqq ast
kull kidhb khwāhad būd. Intahā.

Another statement: ``Let us concede that he is superior to the Most
Exalted Lord, that he is His Resurrection, and that His dispensation has
been superseded---in the sense that `I am that very One whom God will
make manifest, promised by the Most Exalted Lord Himself, through whom
the dispensation of the Bayán will be superseded, and who will abrogate
it.' The falsity of this claim is clearer than the sun. First, the Báb
Himself has stated this in a manner that no one can dispute. One of
those statements, in brief, is this: `If the creation of that Revelation
is not completed, God will not manifest the other Manifestation.' He has
also explained the meaning of completion: since whatever He has said and
commanded is from God, it must not fail to be fulfilled. A creation must
be created beneath His shadow, and all His pronouncements, whether
concerning exalted stations or subsidiary religious ordinances, must be
realized. Otherwise, the appearance of such a person would be futile,
and all the tidings that come from God would be false.'' End of
quotation.

\subsection{341}\label{section-836}

\begin{otherlanguage}{persian}

قل يا ايّها الحمير حقّ آنچه بفرمايد حقّ است و بكلمات مشركين باطل نشود بلكه
آنچه اليوم بطلانش اظهر من الشّمس است انفس خود معرضين بوده و خواهد بود قل
خذ زمامك يا ايّها المكّار و يا ايّها المتوهّم الّذى ما شهدت عين الدّهر مثلك
نميدانم بكدام نفس آن نفوس را شبيه نمايم چه كه در اعراض شبه و كفو و مانند
ندارند

\end{otherlanguage}

Qul yā ayyuhā l-ḥamīr Ḥaqq āncheh be-farmāyad ḥaqq ast va be-kalimāt-i
mushrikīn bāṭil nashavad balkeh āncheh al-yawm buṭlānash aẓharu mina
sh-shams ast anfus-i khod-i muʿriḍīn būdeh va khwāhad būd. Qul khudh
zimāmaka yā ayyuhā l-makkāru wa-yā ayyuhā l-mutawahhimu lladhī mā
shahidat ʿaynu d-dahri mithlaka. Namī-dānam be-kodām nafs ān nufūs rā
shabīh namāyam cheh keh dar iʿrāḍ shibh va kufv va mānand nadārand.

Say: O asses! Whatever God says is true and is not invalidated by the
words of those who associate partners with Him. Rather, what is and will
be clearer than the sun today is the falsity of those very souls who
turn away. Say: Rein yourself in, O deceiver, O deluded one, whose like
the eye of time has never seen! I do not know to whom I can compare
those souls, for in turning away they have no likeness, equal, or
counterpart.

\subsection{342}\label{section-837}

\begin{otherlanguage}{persian}

بگو بآن مشرك ملقى كه اى غافل حقّ را باطل مدان و كلمات حقّ را باطل مشمر و
اللّه يا ايّها المنكر رايحه كلمات الهى از دونش بسى ممتاز بحيث لو يكون فى
شرق الابداع و يتكلّم بكلمة ليجد رايحة القدس منها من كان فى الغرب ان يكون
مطهّراً عن ارياح المشركين آنچه از حقّ جلّ جلاله و عظم كبريائه ظاهر شود فعلاً
و قولاً مطاع بوده و خواهد بود نزد اولو الافئده و بر احدى اعراض و اعتراض
جايز نه و يشهد بذلك كلّ ذى عقل و كلّ ذى دراية و شعور و لكن انّهم لما خلقوا
من الجهل لن يستشعروا و يكوننّ من الغافلين بايد اوّل عارف بحقّ شد كه چه
ميفرمايد و حجّتش چه بعد از ثبوت و اظهار بيّنه و برهان كلّ اين كلمات و
قائلين و متكلّمين آن راجع بنار بوده و خواهند بود و انّ النّار مثويهم فبئس
مثوى المشركين

\end{otherlanguage}

Begū be-ān mushrik-i mulqī keh ey ghāfil Ḥaqq rā bāṭil madān va
kalimāt-i Ḥaqq rā bāṭil mashmur. Wa-llāh yā ayyuhā l-munkir rāyiḥeh-yi
kalimāt-i ilāhī az dūnash basī mumtāz bi-ḥaythu law yakūnu fī sharqi
l-ibdāʿi wa-yatakallamu bi-kalimatin la-yajida rāʾiḥata l-qudsi minhā
man kāna fī l-gharbi in yakūnu muṭahharan ʿan aryāḥi l-mushrikīn. Āncheh
az Ḥaqq jalla jalāluhu wa-ʿaẓuma kibriyāʾuhu ẓāhir shavad fiʿlan va
qawlan muṭāʿ būdeh va khwāhad būd nazd-i ulu l-afʾideh va bar aḥadī
iʿrāḍ va iʿtirāḍ jāʾiz nah. Wa-yashhadu bi-dhālika kullu dhī ʿaqlin
wa-kullu dhī dirāyatin wa-shuʿūrin wa-lākin innahum lammā khuliqū mina
l-jahli lan yastashʿirū wa-yakūnanna mina l-ghāfilīn. Bāyad avval ʿārif
be-Ḥaqq shud keh cheh mī-farmāyad va ḥujjatash cheh. Baʿd az thubūt va
iẓhār-i bayyineh va burhān kull-i īn kalimāt va qāʾilīn va
mutakallimīn-i ān rājiʿ be-nār būdeh va khwāhand būd. Wa-inna n-nāra
mathwāhum fa-biʾsa mathwā l-mushrikīn.

Say to that prompter who associates partners with God: O heedless one,
do not deem God false, and do not reckon His words false. By God, O
denier, the fragrance of the divine words is so distinct from that of
all others that, were He in the east of creation and to speak a single
word, one in the west would perceive the fragrance of holiness from it,
provided he were purified from the winds of those who associate partners
with God. Whatever proceeds from God---glorious is His glory and great
His majesty---in deed or word has been and will be obeyed by those
endowed with hearts. No one is permitted to turn away or object. To this
every person of reason, understanding, and perception bears witness.
Yet, since they were created from ignorance, they will not perceive and
will remain among the heedless. One must first recognize God, what He
says, and what His proof is. Once this is established and clear evidence
and proof have been presented, all these words, and those who utter and
speak them, have returned and will return to the fire. The fire is
indeed their abode, and wretched is the abode of those who associate
partners with God.

\subsection{343}\label{section-838}

\begin{otherlanguage}{persian}

و اينكه نوشته حضرت اعلى فرموده لو لم يكمل خلق ذلك الظّهور الخ امر تكميل
در قبضه قدرت الهى بوده شما و امثال شما ادراك ننموده و نخواهيد نمود چه كه
ميشود هزار سنه ناس بشريعتى عامل باشند و مع‌ذلك لفظ تكميل بر او صادق نيايد
و ميشود در يومى كلّ تكميل شوند و احدى بر تكميل و غير آن جز حقّ مطّلع نه چه
كه تكميل هر امرى منوط باراده او است چنانچه از قبل ذكر شد اگر بخواهد جميع
ناس را بيك نفس راجع و از همان نفس بنفس او راجع و محشور و ظاهر ميفرمايد و
هم‌چنين بعث جميع نفوس و تكميل ايشانرا باقرار نفسى بصاحب ظهور در همان حين
حكم تكميل ظهور قبل در نفس او ميشود كلّ ذلك تقدير من لدن مقتدر قدير و لكن
محتجبين و متوهّمين باينمقام واصل و بعرفانش فايز نگشته

\end{otherlanguage}

Va īn-keh nivishteh Ḥaḍrat-i Aʿlā farmūdeh law lam yukmal khalqu dhālika
ẓ-ẓuhūri ilā ākhir, amr-i takmīl dar qabḍeh-yi qudrat-i ilāhī būdeh
shomā va amthāl-i shomā idrāk nanamūdeh va nakhwāhīd namūd cheh keh
mī-shavad hizār saneh nās be-sharīʿatī ʿāmil bāshand va maʿa-dhālik
lafẓ-i takmīl bar ū ṣādiq nayāyad va mī-shavad dar yawmī kull takmīl
shavand va aḥadī bar takmīl va ghayr-i ān juz Ḥaqq muṭṭaliʿ nah cheh keh
takmīl-i har amrī manūṭ be-irādeh-yi ū ast. Chunāncheh az qabl dhikr
shud agar bi-khwāhad jamīʿ-i nās rā be-yek nafs rājiʿ va az hamān nafs
be-nafs-i ū rājiʿ va maḥshūr va ẓāhir mī-farmāyad va ham-chunīn baʿth-i
jamīʿ-i nufūs va takmīl-i īshān-rā be-iqrār-i nafsī be-Ṣāḥib-i Ẓuhūr dar
hamān ḥīn ḥukm-i takmīl-i ẓuhūr-i qabl dar nafs-i ū mī-shavad. Kullu
dhālika taqdīrun min ladun muqtadirin qadīr. Va lākin muḥtajibīn va
mutawahhimīn be-īn maqām vāṣil va be-ʿirfānash fāʾiz nagashteh.

As for his having written that the Most Exalted One said, ``If the
creation of that Revelation is not completed,'' and so forth: the matter
of completion lies within the grasp of divine power. You and those like
you have not understood it and will not understand it. People may
practice a religious law for a thousand years, yet the term
``completion'' may not apply to it; or all may be brought to completion
in a single day. No one but God knows whether completion has occurred or
not, for the completion of every matter depends on His will. As
mentioned before, if He wishes, He causes all people to return to a
single soul, and from that same soul causes them to return to His own
Self, to be gathered and made manifest. Likewise, the resurrection of
all souls and their completion: upon a soul's acknowledgment of the
Revealer, at that very moment the completion of the preceding Revelation
is decreed in that soul. All this is an ordainment from One who is
all-powerful and mighty. Yet those who are veiled and deluded have not
attained this station or come to understand it.

\subsection{344}\label{section-839}

\begin{otherlanguage}{persian}

بسا ميشود كه شمس حقيقت ظاهر و مشرق و كلّ ناس محتجب و جميع مظاهر اثباتيّه
در نفس ظهور ظاهر و ثابت و مستقرّ و حكم نفى بر كلّ جارى و صادق و هذا سرّ خفىّ
لم يعرفه احد الّا من فتح اللّه بصره و نوّره بنور الايقان چنانچه نقطه بيان
روح ما سواه فداه ميفرمايد قوله عزّ كبريائه تو را محتجب نكند كثرت اين خلق
و انجماد آنها بر امر خود كه اگر فرض تصوّر شود در يك ظهورى شجره حقيقت ناطق
گردد و كلّ تامّل نمايند در حقّ او كلّ در لا اله منفى بوده و هستند اگر چه در
روى ارض باقى هستند ولى بقاى ايشان بقاى در نار بوده و هست و كلّ اثبات آن
نفس ممتنعه مرتفعه بوده انتهى

\end{otherlanguage}

Basā mī-shavad keh Shams-i Ḥaqīqat ẓāhir va mushriq va kull-i nās
muḥtajib va jamīʿ-i maẓāhir-i ithbātiyyeh dar nafs-i Ẓuhūr ẓāhir va
thābit va mustaqarr va ḥukm-i nafy bar kull jārī va ṣādiq. Wa-hādhā
sirrun khafiyyun lam yaʿrifhu aḥadun illā man fataḥa Allāhu baṣarahu
wa-nawwarahu bi-nūri l-īqān. Chunāncheh Nuqṭeh-yi Bayān rūḥ-i mā sivāh
fidāh mī-farmāyad qawluhu ʿazza kibriyāʾuhu: to rā muḥtajib nakonad
kathrat-i īn khalq va injimād-i ānhā bar amr-i khod keh agar farḍ
taṣavvur shavad dar yek ẓuhūrī Shajareh-yi Ḥaqīqat nāṭiq gardad va kull
taʾammul namāyand dar ḥaqq-i ū kull dar lā ilāh manfī būdeh va hastand
agar-cheh dar rū-yi arḍ bāqī hastand valī baqā-yi īshān baqā-yi dar nār
būdeh va hast va kull-i ithbāt ān nafs-i mumtaniʿeh-yi murtafiʿeh būdeh.
Intahā.

It may well be that the Sun of Truth appears and shines while all people
remain veiled, and all the manifestations of affirmation appear, are
established, and abide in the person of the Revelation Himself, while
the decree of negation applies to and holds true of all others. This is
a hidden mystery that no one has known except one whose sight God has
opened and illumined with the light of certitude. Thus the Point of the
Bayán---may the soul of all else be sacrificed for Him---says, exalted
is His majesty: ``Do not let the multitude of these people and their
rigid adherence to their own cause veil you. For if it be supposed that
in a Revelation the Tree of Truth were to speak and all were to hesitate
concerning Him, all would have been and would be negated within `there
is no god.' Though they remain on the face of the earth, their continued
existence has been and is an existence in the fire, and the whole of
affirmation would be that inaccessible, exalted Soul.'' End of
quotation.

\subsection{345}\label{section-840}

\begin{otherlanguage}{persian}

حال در اين بيان نقطه بيان چه ميگوئى و تكميل دوره آنظهور را چگونه تعقّل
مينمائى مع آنكه احدى بعرفانش فايز نشده پس عارف شو كه تكميل دوره بنفس خود
ظهور بوده و خواهد بود و لكن لايعقله الّا من شرب كوثر المعانى من بحر
الأعظم الّذى ظهر عن يمين العرش بامواج ظاهر مبين

\end{otherlanguage}

Ḥāl dar īn bayān-i Nuqṭeh-yi Bayān cheh mī-gūʾī va takmīl-i dawreh-yi ān
ẓuhūr rā chigūneh taʿaqqul mī-namāʾī maʿa ān-keh aḥadī be-ʿirfānash
fāʾiz nashodeh. Pas ʿārif shaw keh takmīl-i dawreh be-nafs-i khod-i
Ẓuhūr būdeh va khwāhad būd va lākin lā yaʿqiluhu illā man shariba
kawthara l-maʿānī min Baḥri l-Aʿẓami lladhī ẓahara ʿan yamīni l-ʿarshi
bi-amwājin ẓāhirin mubīn.

Now, what do you say about this utterance of the Point of the Bayán? How
do you conceive of the completion of the cycle of that Revelation, when
no one has attained recognition of Him? Understand, then, that the
completion of the cycle has been and will be through the person of the
Revelation Himself. Yet no one understands this except one who has drunk
the Kawthar of meanings from the Most Great Ocean, which has appeared to
the right of the throne with manifest, unmistakable waves.

\subsection{346}\label{section-841}

\begin{otherlanguage}{persian}

و هم‌‌چنين مضمون بيان نقطه بيان روح ما سواه فداه است كه ميفرمايند اينكه ما
بين ظهورين بطول ميانجامد سبب عدم استعداد ناس است و شمس حقيقت بعد از غروب
در افق ابهى ناظر بعباد خود بوده هر حين كه ملاحظه فرمايد نفسى ظاهر شده كه
تواند اصغاى كلمة اللّه نمايد در همان حين نفس خود را باو بشناساند و اقلّ از
تسع تسع عشر عشر دقيقه صبر ننمايد و در اينمقام مثل بنقطه فرقان جلّت عظمته
زده‌اند و فرموده‌اند كه مثلا بعد از غروب شمس محمّديّه هميشه ناظر بخلق خود
بوده اگر ملاحظه ميفرمود كه نفسى در ارض موجود شده كه حمل امانة اللّه نمايد
ابدا در ظهور تأخير نميفرمود و اينكه عدد غريس طول انجاميده باينجهت بوده
چنانچه حين استعداد اوّل من آمن و قبول او كلمه اعظم را اقلّ از آن در اظهار
امر تأخير نرفت و همان حين القاى كلمة اللّه باو شد و اگر يك آن قبل از آن
بالغ ميشد و بين يدى حاضر همان يك آن هم تأخير در القاء نميشد

\end{otherlanguage}

Va ham-chunīn maḍmūn-i bayān-i Nuqṭiy-i Bayān, rūḥu mā siwāhu fidāhu,
ast ki mīfarmāyand: Īnki mā bayn-i ẓuhūrayn bi-ṭūl mīanjāmad sabab-i
ʿadam-i istiʿdād-i nās ast va Shams-i Ḥaqīqat baʿd az ghurūb dar ufuq-i
Abhā nāẓir bi-ʿibād-i khud būdi, har ḥīn ki mulāḥiẓi farmāyad nafsī
ẓāhir shudi ki tavānad iṣghā-yi Kalimat Allāh namāyad, dar hamān ḥīn
nafs-i khud rā bi-ū bishināsānad va aqall az tusʿ-i tusʿ-i ʿushr-i
ʿushr-i daqīqi ṣabr nanamāyad. Va dar īn maqām mathal bi-Nuqṭiy-i
Furqān, jallat ʿaẓamatuhu, zadi-and va farmūdi-and ki mathalan baʿd az
ghurūb-i Shams-i Muḥammadiyyi hamīshi nāẓir bi-khalq-i khud būdi, agar
mulāḥiẓi mīfarmūd ki nafsī dar arḍ mawjūd shudi ki ḥaml-i Amānat Allāh
namāyad, abadan dar ẓuhūr taʾkhīr namīfarmūd. Va īnki ʿadad-i Gharīs ṭūl
anjāmīdi bi-īn jihat būdi, chunānchi ḥīn-i istiʿdād-i avval man āman va
qabūl-i ū Kalimiy-i Aʿẓam rā, aqall az ān dar iẓhār-i amr taʾkhīr naraft
va hamān ḥīn ilqā-yi Kalimat Allāh bi-ū shud. Va agar yak ān qabl az ān
bāligh mīshud va bayn-i yaday ḥāḍir, hamān yak ān ham taʾkhīr dar ilqā
namīshud.

Likewise, the substance of the utterance of the Point of the Bayán---may
the soul of all else be sacrificed for Him---is that He says: The
interval between two Revelations is prolonged because people lack
readiness. After setting in the All-Glorious horizon, the Sun of Truth
watches over His servants. Whenever He observes that a soul has appeared
who can hear the Word of God, He makes Himself known to that soul at
that very moment, and does not wait even for less than a ninth of a
ninth of a tenth of a tenth of a minute. In this connection He has given
the Point of the Qur'an---exalted is His grandeur---as an example,
saying that after the setting of the Muhammadan Sun, He continually
watched over His creation. Had He observed that a soul had come into
existence on earth who could bear the Trust of God, He would never have
delayed His Revelation. The interval extended to the number of Gharís
for this reason. Thus, when the first believer became ready and accepted
the Most Great Word, there was not the slightest delay in making the
Cause manifest, and at that very moment the Word of God was imparted to
him. Had he attained maturity and entered His presence one instant
earlier, its imparting would not have been delayed even for that one
instant.

\subsection{347}\label{section-842}

\begin{otherlanguage}{persian}

اى محتجبين در كلمات و بيان الهى تفكّر نمائيد كه شايد برشحى از طمطام بحر
معانى فايز شويد و اگر با بصر باشيد نفس اين ظهور بديع را سبب تكميل بيان
مشاهده مينمائيد كه همان حين ظهور تكميل دوره بيان شده و از آن گذشته جميع
را باحكام بيان امر فرموده‌اند و خود بنفسه عامل بوده‌اند و اگر جميع حجبات
خرق ميشد چه ميكرديد بدانيد كه در حين ظهور هر نفسيكه بين يدى حاضر شد جميع
اسماء و صفات و شرايع و احكام و اوامر كه از قبل بوده باقبال او بحقّ در نفس
او تكميل شده و ميشود اذا تفكّر لتكون من العارفين

\end{otherlanguage}

Ay muḥtajibīn, dar kalimāt va bayān-i ilāhī tafakkur namāyīd ki shāyad
bi-rashḥī az ṭamṭām-i baḥr-i maʿānī fāʾiz shavīd. Va agar bā baṣar
bāshīd, nafs-i īn ẓuhūr-i badīʿ rā sabab-i takmīl-i Bayān mushāhidi
mīnamāyīd ki hamān ḥīn-i ẓuhūr takmīl-i dawriy-i Bayān shudi, va az ān
gudhashti jamīʿ rā bi-aḥkām-i Bayān amr farmūdi-and va khud bi-nafsihi
ʿāmil būdi-and. Va agar jamīʿ-i ḥujubāt kharq mīshud chi mīkardīd?
Bidānīd ki dar ḥīn-i ẓuhūr har nafsī ki bayn-i yaday ḥāḍir shud, jamīʿ-i
asmā va ṣifāt va sharāyiʿ va aḥkām va avāmir ki az qabl būdi bi-iqbāl-i
ū bi-Ḥaqq dar nafs-i ū takmīl shudi va mīshavad. Idhan tafakkar
li-takūna mina al-ʿārifīn.

You who are veiled, reflect on the words and utterances of God, that you
may perhaps attain a sprinkling from the surging ocean of meanings. If
you have eyes to see, you will recognize the very person of this
wondrous Revelation as the means of completing the Bayán, for at the
very moment of His Revelation the cycle of the Bayán was completed.
Beyond this, He has enjoined everyone to observe the laws of the Bayán
and has Himself acted in accordance with them. Had all the veils been
torn apart, what would you have done? Know that, at the time of the
Revelation, in every soul who entered His presence, all the names,
attributes, religious laws, ordinances, and commandments that had
existed before were, and are, brought to completion through that soul's
turning toward the Truth. Reflect, then, that you may be among those who
recognize.

\subsection{348}\label{section-843}

\begin{otherlanguage}{persian}

و اينكه نوشته بودى بايد كلّ فرمايشات او چه از مقامات عاليه و چه از فروعات
دينيّه بظهور آيد و الّا ظهور چنين شخص لغو خواهد بود و كلّ اخبار حقّ كذب
بدانكه كلّ ما فى البيان بعالم ظهور آمده و بين يدى اللّه محشور گشته و لكن
آن وقت امثال آن نفوس ميّت و لاشیء بوده‌اند و لذا مستشعر نشده‌اند بعد از
ظهور اين اقوال احدى را غنى ننموده و نخواهد نمود چه كه تكميل كلّ بتصديق او
بوده و عزّ كلّ بامر او و اعمال كلّ بقبول او منوط بوده

\end{otherlanguage}

Va īnki nivishti būdī bāyad kull-i farmāyishāt-i ū chi az maqāmāt-i
ʿāliyi va chi az furūʿāt-i dīniyyi bi-ẓuhūr āyad va illā ẓuhūr-i chunīn
shakhs laghv khvāhad būd va kull-i akhbār-i Ḥaqq kidhb, bidānki kull-i
mā fī al-Bayān bi-ʿālam-i ẓuhūr āmadi va bayn-i yaday Allāh maḥshūr
gashti va lākin ān vaqt amthāl-i ān nufūs mayyit va lā-shayʾ būdi-and va
lidhā mustashʿir nashudi-and. Baʿd az ẓuhūr īn aqvāl aḥadī rā ghanī
nanamūdi va nakhvāhad namūd, chi ki takmīl-i kull bi-taṣdīq-i ū būdi va
ʿizz-i kull bi-amr-i ū va aʿmāl-i kull bi-qabūl-i ū manūṭ būdi.

As for your statement that all His pronouncements, whether concerning
exalted stations or subsidiary religious matters, must be fulfilled,
otherwise the appearance of such a person would be futile and all God's
tidings false---know that everything in the Bayán has come into the
realm of manifestation and has been gathered into the presence of God.
At that time, however, such souls were dead and as nothing, and
therefore remained unaware. After the Revelation, these assertions have
availed no one and never will, for the completion of all things has
depended on His confirmation, the glory of all on His command, and the
deeds of all on His acceptance.

\subsection{349}\label{section-844}

\begin{otherlanguage}{persian}

مثلا اگر بخواهد كلّ آنچه را در ظهور قبل عنايت فرمود اخذ نمايد احدى را
نرسد كه بر او اعتراض نمايد چه كه او قادر بوده و خواهد بود يك بيان از
نقطه بيان روح ما سواه فداه ذكر ميشود لعلّ تستشعر فى نفسك و تكون من
المستشعرين

\end{otherlanguage}

Mathalan agar bikhvāhad kull-i ānchi rā dar ẓuhūr-i qabl ʿināyat farmūd
akhdh namāyad, aḥadī rā narisad ki bar ū iʿtirāḍ namāyad, chi ki ū qādir
būdi va khvāhad būd. Yak bayān az Nuqṭiy-i Bayān, rūḥu mā siwāhu fidāhu,
dhikr mīshavad, laʿallaka tastashʿiru fī nafsika va takūna mina
al-mustashʿirīn.

For example, should He wish to take back everything He bestowed in the
preceding Revelation, no one has the right to object to Him, for He has
been and will remain powerful. An utterance of the Point of the
Bayán---may the soul of all else be sacrificed for Him---is quoted, that
you may perhaps become aware within yourself and be among those who
perceive.

\subsection{350}\label{section-845}

\begin{otherlanguage}{persian}

قوله جلّ كبريائه مثل كلّ بيان مثل جوهرى است كه نزد كسى امانت گذارد و در
يوم ظهور مظهر حقيقت اگر از كلّ خواهد اخذ نمايد آنچه بايشان داده قدر شیء
صبر ننمايند نه آنكه يكى بمسئله فروع آن محتجب گردد و يكى باصول آن و يكى
بشئون حكم و يكى بشئون عزّ بلكه كلّ از اوست و باو راجع ميگردد و او را
شناخته بآيات او و احتياط در عرفان او نكرده كه بقدر همان در نار محتجب
خواهيد بود و اگر در مابين خود و خدا توجّه ميكنيد مثال او است كه در افئده
شماست باو از او محتجب نگرديد بشناسيد كسى را كه از براى عرفان او خلق
شده‌ايد

\end{otherlanguage}

Qawluhu jalla kibriyāʾuhu: Mathal-i kull-i Bayān mathal-i jawharī ast ki
nazd-i kasī amānat gudhārad va dar yawm-i ẓuhūr-i Maẓhar-i Ḥaqīqat agar
az kull khvāhad akhdh namāyad ānchi bi-īshān dādi, qadr-i shayʾ ṣabr
nanamāyand, na ānki yakī bi-masʾaliy-i furūʿ-i ān muḥtajib gardad va
yakī bi-uṣūl-i ān va yakī bi-shuʾūn-i ḥukm va yakī bi-shuʾūn-i ʿizz,
balki kull az ūst va bi-ū rājiʿ mīgardad. Va ū rā shinākhti bi-āyāt-i ū
va iḥtiyāṭ dar ʿirfān-i ū nakardi ki bi-qadr-i hamān dar nār muḥtajib
khvāhīd būd. Va agar dar mābayn-i khud va Khudā tavajjuh mīkunīd,
mithāl-i ū ast ki dar afʾidiy-i shumāst. Bi-ū az ū muḥtajib nagardīd.
Bishināsīd kasī rā ki az barā-yi ʿirfān-i ū khalq shudi-īd.

His words---exalted is His majesty: ``The whole Bayán is like a jewel
entrusted to someone. On the day of the appearance of the Manifestation
of Truth, if He wishes to take back from everyone what He has given
them, they must not hesitate in the slightest---not that one should be
veiled by a question concerning its subsidiary matters, another by its
principles, another by the aspects of authority, and another by the
aspects of glory. Rather, everything comes from Him and returns to Him.
Recognize Him by His verses and exercise no caution in recognizing Him,
for to that same extent you will remain veiled in the fire. If, in the
privacy between yourselves and God, you turn toward Him, it is His image
that is in your hearts. Do not let it veil you from Him. Recognize the
One for whose recognition you have been created.''

\subsection{351}\label{section-846}

\begin{otherlanguage}{persian}

انتهى حال ملاحظه نمائيد كه چه قدر از شاطى بحر عرفان دور مانده‌ايد و چه
مقدار بحجاب نفس و هوى از موجد و محقّق اسماء محروم گشته‌ايد فواللّه آنچه
نقطه بيان خواسته خلاف او نموده‌ايد آنچه آيات محكمات او بوده بظنونات افكيّه
تحريف و تاويل نموده‌ايد قطره از بحر تقديس نياشاميده‌ايد و بهبوبی‌ از ارياح
عزّ تنزيه فايز نگشته‌ايد مع اينكه كوثر بيان كه از معين قلم سلطان امكان
جارى شده گاهى فروع ذكر مينمائيد و گاهى اصول مذكور ميداريد و اللّه مقصودى
نداشته و نداريد جز آنكه بحبلى از حبال نفس تمسّك جسته و بآن بر سلطان احديّه
اعتراض نمائيد

\end{otherlanguage}

Intahā. Ḥāl mulāḥiẓi namāyīd ki chi qadr az shāṭī-yi baḥr-i ʿirfān dūr
māndi-īd va chi miqdār bi-ḥijāb-i nafs va havā az Mūjid va Muḥaqqiq-i
asmā maḥrūm gashti-īd. Fa-wa-Allāh ānchi Nuqṭiy-i Bayān khvāsti khilāf-i
ū namūdi-īd. Ānchi āyāt-i muḥkamāt-i ū būdi bi-ẓunūnāt-i ifkiyyi taḥrīf
va taʾvīl namūdi-īd. Qaṭri az baḥr-i taqdīs nayāshāmīdi-īd va bi-hubūbī
az aryāḥ-i ʿizz-i tanzīh fāʾiz nagashti-īd. Maʿ-i īnki Kawthar-i Bayān
ki az maʿīn-i qalam-i Sulṭān-i imkān jārī shudi, gāhī furūʿ dhikr
mīnamāyīd va gāhī uṣūl madhkūr mīdārīd. Va Allāh maqṣūdī nadāshti va
nadārīd juz ānki bi-ḥablī az ḥibāl-i nafs tamassuk justi va bi-ān bar
Sulṭān-i Aḥadiyyi iʿtirāḍ namāyīd.

End of quotation. Now consider how far you have remained from the shore
of the ocean of recognition, and how greatly the veil of self and desire
has deprived you of the Originator and Establisher of the names. By God,
you have acted contrary to everything the Point of the Bayán desired.
You have distorted and explained away all His decisive verses through
false conjectures. You have not drunk a drop from the ocean of holiness
or received a breath of the winds of the glory of sanctification. Yet
concerning the Kawthar of utterance that has flowed from the wellspring
of the pen of the Sovereign of existence, you speak sometimes of
subsidiary matters and sometimes of principles. By God, you have had,
and still have, no aim but to seize a cord among the cords of self and
use it to object to the Sovereign of Oneness.

\subsection{352}\label{section-847}

\begin{otherlanguage}{persian}

حال ملاحظه نمائيد كه كلمات الهى چگونه محقّق شده نصّ فرموده‌اند كه در آن يوم
احدى بذكر فروع و اصول و حكم و شئون عزّ محتجب نماند چه كه آن ساذج قدم خود
بنفسه مطاع بوده و هر چه در حين ظهور بفرمايد آن حكم الهى و امر او بوده ما
بين عباد مع‌ذلك تو و امثال تو بذكر اصول و فروع احتجاج بر حقّ نموده‌ايد و
مرشد بشئون عزّ و حكم اليوم كلّ در نار مشاهده ميشويد الّا من تمسّك باللّه
المقتدر العزيز المختار

\end{otherlanguage}

Ḥāl mulāḥiẓi namāyīd ki kalimāt-i ilāhī chigūni muḥaqqaq shudi. Naṣṣ
farmūdi-and ki dar ān yawm aḥadī bi-dhikr-i furūʿ va uṣūl va ḥukm va
shuʾūn-i ʿizz muḥtajib namānad, chi ki ān Sādhij-i Qidam khud bi-nafsihi
muṭāʿ būdi va har chi dar ḥīn-i ẓuhūr bifarmāyad ān ḥukm-i ilāhī va
amr-i ū būdi mā bayn-i ʿibād. Maʿa dhālik tu va amthāl-i tu bi-dhikr-i
uṣūl va furūʿ iḥtijāj bar Ḥaqq namūdi-īd va murshid bi-shuʾūn-i ʿizz va
ḥukm. Al-yawm kull dar nār mushāhidi mīshavīd illā man tamassaka
bi-Allāh al-Muqtadir al-ʿAzīz al-Mukhtār.

Now consider how the words of God have been fulfilled. He expressly
stated that on that day no one should remain veiled by reference to
subsidiary matters, principles, authority, or the aspects of glory, for
that Essence of Eternity is Himself the One to be obeyed, and whatever
He pronounces at the time of His Revelation is the decree of God and His
command among His servants. Nevertheless, you and those like you have
argued against the Truth by invoking principles and subsidiary matters,
and the spiritual guide by invoking the aspects of glory and authority.
Today you are all seen in the fire, except those who have held fast to
God, the All-Powerful, the Almighty, the Unconstrained.

\subsection{353}\label{section-848}

\begin{otherlanguage}{persian}

اى بی‌خبران بشأنى مبالغه فرموده‌اند كه ميفرمايند در آن يوم احتياط در عرفان
او مكنيد كه بقدر همان احتياط در نار خواهيد بود مع‌ذلك وارد آورده‌ايد آنچه
را كه هيچيك از مظاهر نفى و سجّين وارد نياورده‌اند يك بيان ديگر از بيان
منزل بيان لوجه اللّه ذكر ميشود شايد اثواب رثيثه بغضا را از هيكل بيندازيد
و برداء عزّ الهى فايز شويد قوله جلّ كبريائه ولى مغرور نگشته بحبّ نقطه و
حروف حىّ كه آن روز روز امتحان است اگر كسى آن نقطه و حروف را دوست داشت و
بهدايت آن مهتدى گشت دليل است كه اين نقطه و حروف را دوست داشته انتهى

\end{otherlanguage}

Ay bī-khabarān, bi-shaʾnī mubālaghi farmūdi-and ki mīfarmāyand dar ān
yawm iḥtiyāṭ dar ʿirfān-i ū makunīd ki bi-qadr-i hamān iḥtiyāṭ dar nār
khvāhīd būd. Maʿa dhālik vārd āvardi-īd ānchi rā ki hīch-yak az
maẓāhir-i nafy va sijjīn vārd nayāvardi-and. Yak bayān-i dīgar az
bayān-i Munzil-i Bayān li-wajh Allāh dhikr mīshavad, shāyad athvāb-i
rathīthiy-i baghḍā rā az haykal biyandāzīd va bi-ridā-yi ʿizz-i ilāhī
fāʾiz shavīd. Qawluhu jalla kibriyāʾuhu: Valī maghrūr nagashti bi-ḥubb-i
Nuqṭi va Ḥurūf-i Ḥayy ki ān rūz rūz-i imtiḥān ast. Agar kasī ān Nuqṭi va
Ḥurūf rā dūst dāsht va bi-hidāyat-i ān muhtadī gasht, dalīl ast ki īn
Nuqṭi va Ḥurūf rā dūst dāshti. Intahā.

You who are unaware! So emphatically has He spoken that He says: On that
day, exercise no caution in recognizing Him, for to the extent of that
caution you will be in the fire. Nevertheless, you have inflicted what
none of the manifestations of negation and Sijjín has inflicted. Another
utterance of the Revealer of the Bayán is quoted for the sake of God,
that you may perhaps cast the tattered garments of hatred from your
bodies and attain the robe of divine glory. His words---exalted is His
majesty: ``But do not become complacent through love for the Point and
the Letters of the Living, for that day is the day of testing. If
someone loves that Point and those Letters and is guided by their
guidance, this is proof that he has loved this Point and these
Letters.'' End of quotation.

\subsection{354}\label{section-849}

\begin{otherlanguage}{persian}

ملاحظه كن كه چه مقدار محكم و متقن اين امر ابدع امنع اقدس را در كتاب ذكر
فرموده‌اند حتّى آنكه ميفرمايند بحبّ من مغرور نشويد و از آن جمال قدم محجوب
نمانيد مع‌ذلك در كلّ حين باحتجاجات لايحصى احتجاج نموده و مينمائيد و اگر
بگوئيد اين ظهور آن ظهور نيست اين همان قولى است كه كلّ ملل قبل در احيان
ظهور مظاهر احديّه گفته‌اند مثلا در حين ظهور روح اللّه و كلمه او جميع بكتب
قبل استدلالات بر بطلان او نموده و جميع معارض و منكران الّا من شآء ربّك كه
معدودى قليلى بودند و هم‌‌چنين در ظهور خاتم انبیاء كه از مشرق بطحا اشراق
فرمود كلّ منكر شدند كه تو آن ظهور موعود نيستى و ادلّه‌هاى مجعوله ذكر نموده
رسايل بر ردّ حقّ نوشتند و هم‌‌چنين در ظهور نقطه بيان مشاهده نموده و مينمائيد
كه چه كرده و ميكنند چه فرق است مابين اين طايفه و آن گروه

\end{otherlanguage}

Mulāḥiẓi kun ki chi miqdār muḥkam va mutqan īn amr-i abdaʿ-i amnaʿ-i
aqdas rā dar Kitāb dhikr farmūdi-and, ḥattā ānki mīfarmāyand bi-ḥubb-i
man maghrūr nashavīd va az ān Jamāl-i Qidam maḥjūb namānīd. Maʿa dhālik
dar kull-i ḥīn bi-iḥtijājāt-i lā-yuḥṣā iḥtijāj namūdi va mīnamāyīd. Va
agar bigūʾīd īn ẓuhūr ān ẓuhūr nīst, īn hamān qawlī ast ki kull-i
milal-i qabl dar aḥyān-i ẓuhūr-i Maẓāhir-i Aḥadiyyi gufti-and. Mathalan
dar ḥīn-i ẓuhūr-i Rūḥ Allāh va Kalimiy-i ū jamīʿ bi-kutub-i qabl
istidlālāt bar buṭlān-i ū namūdi va jamīʿ muʿāriḍ va munkirān, illā man
shāʾa rabbuka ki maʿdūdī qalīlī būdand. Va ham-chunīn dar ẓuhūr-i
Khātam-i Anbiyā ki az mashriq-i Baṭḥā ishrāq farmūd, kull munkir shudand
ki tu ān ẓuhūr-i mawʿūd nīstī va adillihā-yi majʿūli dhikr namūdi,
rasāʾil bar radd-i Ḥaqq nivishtand. Va ham-chunīn dar ẓuhūr-i Nuqṭiy-i
Bayān mushāhidi namūdi va mīnamāyīd ki chi kardi va mīkunand. Chi farq
ast mābayn-i īn ṭāʾifi va ān gurūh?

Consider how firmly and conclusively He has set forth this most
wondrous, most unassailable, most holy Cause in the Book, even saying:
Do not become complacent through love for Me, and do not remain veiled
from that Ancient Beauty. Nevertheless, at every moment you have
advanced, and continue to advance, countless arguments. If you say that
this Revelation is not that Revelation, this is the very assertion made
by all previous religious communities at the times when the
Manifestations of Oneness appeared. For example, when the Spirit of God
and His Word appeared, everyone invoked the former scriptures to argue
that He was false, and all were opponents and deniers, except those whom
your Lord willed, who were but a few. Likewise, at the Revelation of the
Seal of the Prophets, who shone from the dawning-place of Bathá, all
denied Him, saying, `You are not the promised Revelation.' They put
forward fabricated proofs and wrote treatises refuting the Truth.
Likewise, in the Revelation of the Point of the Bayán, you have seen,
and still see, what they have done and are doing. What difference is
there between this community and that group?

\subsection{355}\label{section-850}

\begin{otherlanguage}{persian}

انسان بصير بايد اليوم بنفس ظهور و اصل امر ناظر باشد اگر همان بيّنه و
برهان كه بآن اثبات نقطه اوّلیّه و طلعت احديّه نموده مشاهده نمود اقلّ من آن
توقّف جايز نه چه كه بقدر همان توقّف در نار خواهد بود و الّا از براى ناظرين
منظر اكبر نفس ظهور بنفسه حجّت بوده اين عبد متحيّر است كه چه ذكر نمايد چه
كه مشاهده ميشود كه از تصريح كلمات الهى محروم مانده‌ايد چه رسد بتلويح و
اشاره و گويا اراده ادراك هم نداشته ونداريد لذا قول اين عبد و ترنّماتش مثل
نواختن بربط است از براى نفوسى كه از سمع محروم باشند اين فرد در اينمقام
مناسب است

\end{otherlanguage}

Insān-i baṣīr bāyad al-yawm bi-nafs-i ẓuhūr va aṣl-i amr nāẓir bāshad.
Agar hamān bayyini va burhān ki bi-ān ithbāt-i Nuqṭiy-i Ūlā va Ṭalʿat-i
Aḥadiyyi namūdi mushāhidi namūd, aqall min ān tavaqquf jāʾiz na, chi ki
bi-qadr-i hamān tavaqquf dar nār khvāhad būd. Va illā az barā-yi
nāẓirīn-i Manẓar-i Akbar nafs-i ẓuhūr bi-nafsihi ḥujjat būdi. Īn ʿabd
mutaḥayyir ast ki chi dhikr namāyad, chi ki mushāhidi mīshavad ki az
taṣrīḥ-i kalimāt-i ilāhī maḥrūm māndi-īd chi risad bi-talvīḥ va ishāri.
Va gūyā irādiy-i idrāk ham nadāshti va nadārīd. Lidhā qawl-i īn ʿabd va
tarannumātash mithl-i navākhtan-i barbaṭ ast az barā-yi nufūsī ki az
samʿ maḥrūm bāshand. Īn fard dar īn maqām munāsib ast.

A person of insight must today look to the Revelation Himself and the
essence of the Cause. If he sees the same clear evidence and proof by
which he established the Primal Point and the Countenance of Oneness,
hesitation for less than an instant is not permissible, for to the
extent of that hesitation he will be in the fire. Indeed, for those who
behold the Most Great Vision, the Revelation Himself is His own proof.
This servant is bewildered as to what to say, for it is evident that you
have remained deprived of the explicit statements of God's words, let
alone their allusions and intimations. It seems that you have not even
intended to understand, nor do you now. Thus this servant's words and
melodies are like playing a lute for souls deprived of hearing. This
verse is fitting here.

\subsection{356}\label{section-851}

\begin{otherlanguage}{persian}

نكته رمز سنائى نزد نادانان چنانكه پيش كر بربط سراى پيش كور آيينه دار ضرب
اللّه على آذانكم و قراً و على قلوبكم غشاوة من النّار

\end{otherlanguage}

Nukti-yi ramz-i Sanāʾī nazd-i nādānān chunānki pīsh-i kar barbaṭ sarāy,
pīsh-i kūr āyīni dār. Ḍaraba Allāhu ʿalā ādhānikum waqran va ʿalā
qulūbikum ghishāwatan mina al-nār.

Saná'í's subtle, enigmatic words among the ignorant are like playing a
lute before the deaf or holding a mirror before the blind. God has laid
deafness upon your ears and a covering of fire upon your hearts.

\subsection{357}\label{section-852}

\begin{otherlanguage}{persian}

ان يا ايّها الغافل ميفرمايند عزّ ذكره لاتحتجبوا بحجب العلم و الحكمة فانّ
هذا ليصدّنّكم عن صراط اللّه تازه شما اشارات علميّه پيش سلطان معلوم هديه
ميفرستيد و ورقه يابسه برضوان عزّ احديّه تحفه ارسال ميداريد اى كاش صاحب علم
و حكمت هم ميبوديد لا واللّه لاشیء محض و جلود يابسه مشاهده ميشويد كه در آن
احجار صغيره باشد و در حين حركت صدائى از او ظاهر و از وجود و حركت اصلى
بی‌‌نصيب و عارى عجل جسد له خوار

\end{otherlanguage}

In yā ayyuhā al-ghāfil, mīfarmāyand, ʿazza dhikruhu: Lā taḥtajibū
bi-ḥujubi al-ʿilmi wa-al-ḥikmati fa-inna hādhā la-yaṣuddannakum ʿan
ṣirāṭi Allāh. Tāzi shumā ishārāt-i ʿilmiyyi pīsh-i Sulṭān-i Maʿlūm
hadiyyi mīfiristīd va varaqiy-i yābisi bi-Riḍvān-i ʿizz-i Aḥadiyyi tuḥfi
irsāl mīdārīd. Ay kāsh ṣāḥib-i ʿilm va ḥikmat ham mībūdīd. Lā va Allāh,
lā-shayʾ-i maḥḍ va julūd-i yābisi mushāhidi mīshavīd ki dar ān aḥjār-i
ṣaghīri bāshad va dar ḥīn-i ḥarakat ṣadāʾī az ū ẓāhir va az vujūd va
ḥarakat-i aṣlī bī-naṣīb va ʿārī. ʿIjlun jasadun lahu khuwār.

O heedless one! He says---exalted is His remembrance: ``Do not veil
yourselves with the veils of knowledge and wisdom, for this will surely
bar you from the path of God.'' Yet now you send learned allusions as
gifts to the Sovereign of all that is known, and dispatch a dried leaf
as an offering to the Ridván of the glory of Oneness. If only you
possessed knowledge and wisdom! No, by God, you are seen as sheer
nothingness and dried skins containing small stones, from which a sound
emerges when they move, yet which are bereft and devoid of true
existence and movement: ``A calf, a body that lows.''

\subsection{358}\label{section-853}

\begin{otherlanguage}{persian}

اليوم هيچ ادلّه مدلول را از امرش منع نمينمايد و هيچ شیء از آنچه مابين
سموات و ارض خلق شده كافى نبوده و نخواهد بود مگر دخول در ظلّ امر او اگر كلّ
ارض مملوّ از كتب شود مقابله نمينمايد با يكحرف كه از لسان جوهر قدم اليوم
ظاهر ميگردد بريزيد اين تحقيقات لايسمنه لايغنيه را و پاك كنيد صدور را از
وساوس عبارات اولى الاشارات كه شايد در اين ربيع الهى برضوان بديع فايز
شويد از دونش بگسليد و باو متمسّك گرديد جهد نمائيد كه شايد اليوم بكلمات و
اشارات و اسماء و مظاهر آن از موجد و محقّق آن محجوب نمانيد اينست كه نقطه
بيان جلّ اعزازه ميفرمايد قوله عزّ قوله كلّ از براى اينست كه در يوم ظهور
باسماء محتجب نمانيد بلكه نظر كنيد بما يقوم به الأسماء من كلّ شیء حتّى ذكر
النّبی‌ فانّ ذلك الاسم يخلق به بما ينزل اللّه انتهى

\end{otherlanguage}

Al-yawm hīch adilla madlūl rā az amrash manʿ namīnamāyad va hīch shayʾ
az ānchi mābayn-i samāvāt va arḍ khalq shuda kāfī nabūda va nakhvāhad
būd magar dukhūl dar ẓill-i amr-i ū. Agar kull-i arḍ mamlūv az kutub
shavad muqābila namīnamāyad bā yak ḥarf ki az lisān-i Jawhar-i Qidam
al-yawm ẓāhir mīgardad. Birīzīd īn taḥqīqāt-i lā yusminuhu lā yughnīhi
rā va pāk kunīd ṣudūr rā az vasāvis-i ʿibārāt-i ūlī al-ishārāt ki shāyad
dar īn rabīʿ-i ilāhī bi-Riḍvān-i badīʿ fāyiz shavīd. Az dūnash buguslīd
va bi-ū mutamassik gardīd. Jahd namāyīd ki shāyad al-yawm bi-kalimāt va
ishārāt va asmāʾ va maẓāhir-i ān az mūjid va muḥaqqiq-i ān maḥjūb
namānīd. Īnast ki Nuqṭa-yi Bayān jalla iʿzāzuhu mīfarmāyad qawluhu ʿazza
qawluhu: Kull az barā-yi īnast ki dar yawm-i ẓuhūr bi-asmāʾ muḥtajib
namānīd balki naẓar kunīd bi-mā yaqūmu bihi al-asmāʾ min kulli shayʾ
ḥattā dhikr al-Nabī fa-inna dhālika al-ism yukhlaqu bihi bi-mā yunazzilu
Allāh. Intahā.

Today no proofs can bar the One they attest from His Cause, and nothing
created between the heavens and the earth has been or will be sufficient
except entry beneath the shadow of His Cause. Were the whole earth
filled with books, they would not equal a single letter appearing today
from the tongue of the Essence of Eternity. Cast away these
investigations that neither nourish nor enrich, and cleanse your hearts
of the misgivings arising from the expressions of those versed in
allusions, that perhaps in this divine springtime you may attain the
wondrous Paradise. Sever yourselves from all but Him and cling to Him.
Strive, that today you may not be veiled by words, allusions, names, and
their manifestations from the One who brings them into being and gives
them reality. This is why the Point of the Bayán---exalted be His
glory---says, in His exalted words: ``All is for this purpose: that on
the Day of Manifestation you may not remain veiled by names, but rather
look to that through which the names of all things subsist, even the
designation `Prophet,' for that name is created through Him by what God
sends down.'' End of quotation.

\subsection{359}\label{section-854}

\begin{otherlanguage}{persian}

اگر منصف باشيد از اين بيان جمال رحمن ادراك مينمائيد آنچه را كه اليوم از
او محتجب شده‌ايد چنانچه باسمى از اسماء كه نميدانيد كيست و چيست از خالق آن
محروم مانده‌ايد چند نفرى جمع شده بايادى وهميّه ظنونيّه هيكلى ترتيب داده‌ايد
و بعد باو عاكف شده‌ايد صنعتم بأيادى الجهل هيكل العجل ثمّ اتّخذتموه لأنفسكم
ربّاً من دون اللّه فويل لكم يا ملأ الشّرك و النّفاق

\end{otherlanguage}

Agar munṣif bāshīd az īn bayān-i Jamāl-i Raḥmān idrāk mīnamāyīd ānchi rā
ki al-yawm az ū muḥtajib shuda-īd. Chunānchi bi-ismī az asmāʾ ki
namīdānīd kīst va chīst az khāliq-i ān maḥrūm mānda-īd. Chand nafarī
jamʿ shuda bi-ayādī-yi vahmiyya-yi ẓunūniyya haykalī tartīb dāda-īd va
baʿd bi-ū ʿākif shuda-īd. Ṣanaʿtum bi-ayādī al-jahl haykal al-ʿijl
thumma ittakhadhtumūhu li-anfusikum rabban min dūni Allāh fa-waylun
lakum yā malaʾ al-shirk wa-al-nifāq.

If you are fair-minded, you will perceive from this utterance of the
Beauty of the All-Merciful that from which you are veiled today. Through
one of the names, whose identity and nature you do not know, you have
been deprived of its Creator. A few of you have gathered together and
fashioned a figure with the hands of fancy and conjecture, and then
devoted yourselves to it. With the hands of ignorance you fashioned the
figure of the calf, then took it as a lord for yourselves instead of
God. Woe to you, people of idolatry and hypocrisy!

\subsection{360}\label{section-855}

\begin{otherlanguage}{persian}

از جميع آيات و كلمات ميگذريم امرى را اليوم مجدّداً ميزان معرفة اللّه قرار
دهيد و طرفين را بسنجيد اگر چه حقّ بايد عباد را امتحان نمايد و بسنجد نه
عباد او را چه كه نفس حقّ و ما يظهر من عنده مقدّس از ميزان ناس بوده و خواهد
بود و لكن باين راضى شديم كه حقّ از باطل معلوم شود چه كه حجّت الهى را حجّت
نميدانيد و برهان او را طرد نموده‌ايد در اينصورت اتماماً لحجّة اللّه عليكم
اظهار ميشود كه هر ميزان و برهانى كه نزد آن گروه ميزان معرفت حقّ است ميزان
قرار دهند تا حقّ از باطل معلوم شود و اين كلمه بعقايد شما ذكر ميشود تا
ديگر عذرى نماند و حقّ چون شمس بر هر نفسى ظاهر و لايح و هويدا گردد و الّا
حقّ مقدّس از ميزان عباد و برهانهم و حججهم و دلايلهم بوده و خواهد بود چه كه
ميزان اللّه اليوم نفس حقّ بوده و برهان و حجّتش ما ظهر من عنده ان انتم
تفقهون

\end{otherlanguage}

Az jamīʿ-i āyāt va kalimāt mīgudharīm. Amrī rā al-yawm mujaddadan
mīzān-i maʿrifat Allāh qarār dahīd va ṭarafayn rā bisanjīd. Agar chi
Ḥaqq bāyad ʿibād rā imtiḥān namāyad va bisanjad na ʿibād ū rā, chi ki
nafs-i Ḥaqq va mā yaẓharu min ʿindihi muqaddas az mīzān-i nās būda va
khvāhad būd. Va lākin bi-īn rāḍī shudīm ki Ḥaqq az bāṭil maʿlūm shavad
chi ki ḥujjat-i ilāhī rā ḥujjat namīdānīd va burhān-i ū rā ṭard
namūda-īd. Dar īn ṣūrat itmāman li-ḥujjati Allāh ʿalaykum iẓhār mīshavad
ki har mīzān va burhānī ki nazd-i ān gurūh mīzān-i maʿrifat-i Ḥaqq ast
mīzān qarār dahand tā Ḥaqq az bāṭil maʿlūm shavad. Va īn kalima
bi-ʿaqāyid-i shumā dhikr mīshavad tā dīgar ʿudhrī namānad va Ḥaqq chun
shams bar har nafsī ẓāhir va lāyiḥ va huvaydā gardad. Va illā Ḥaqq
muqaddas az mīzān-i ʿibād va burhānihim va ḥujajihim va dalāyilihim būda
va khvāhad būd, chi ki mīzān Allāh al-yawm nafs-i Ḥaqq būda va burhān va
ḥujjatash mā ẓahara min ʿindihi in antum tafqahūn.

We set aside all the verses and words. Today, once again, establish a
criterion for the recognition of God and weigh both sides. It is God who
should test and weigh His servants, not His servants Him, for God
Himself and whatever appears from Him have always been and will always
remain sanctified above the standards of humanity. Yet We have consented
to this so that truth may be distinguished from falsehood, since you do
not regard the divine testimony as testimony and have rejected His
proof. Therefore, to make God's testimony against you complete, it is
declared that they should adopt as their standard whatever criterion and
proof that group regards as the measure for recognizing God, so that
truth may be distinguished from falsehood. This statement is made in
accordance with your beliefs, so that no excuse may remain and the Truth
may become as manifest, evident, and plain as the sun to every soul.
Otherwise, God has always been and will always remain sanctified above
the standards of His servants and their proofs, testimonies, and
arguments, for God's standard today is His own Self, and His proof and
testimony are whatever has appeared from Him, if you understand.

\subsection{361}\label{section-856}

\begin{otherlanguage}{persian}

و اينكه بلسان شرك و اعراض القاء نموده كه ظهور چنين شخصى لغو خواهد بود و
اخبار حقّ كلّ كذب قطع اللّه لسان ملقيك و كسّراللّه يده و انامله و قلمه لانّه
تجاوز عن حدّه و افترى باللّه المهيمن المقتدر العزيز القيوم و ارتكب ما ضيعت
به حرمة الأمر بين العباد اذا يتبرئنّ منه اهل السجّين و كيف عباد اللّه
المقرّبين

\end{otherlanguage}

Va īnki bi-lisān-i shirk va iʿrāḍ ilqā namūda ki ẓuhūr-i chunīn shakhṣī
laghv khvāhad būd va akhbār-i Ḥaqq kull kidhb: qaṭaʿa Allāh lisāna
mulqīka wa-kassara Allāh yadahu wa-anāmilahu wa-qalamahu li-annahu
tajāwaza ʿan ḥaddihi wa-iftarā bi-Allāh al-Muhaymin al-Muqtadir al-ʿAzīz
al-Qayyūm wa-irtakaba mā ḍuyyiʿat bihi ḥurmatu al-Amr bayna al-ʿibād.
Idhan yatabarraʾanna minhu ahlu al-sijjīn wa-kayfa ʿibādu Allāh
al-muqarrabīn.

As for what he has suggested in the language of idolatry and
rejection---that the appearance of such a person would be futile and all
God's pronouncements false---may God cut out the tongue of the one who
prompted you, and may God break his hand, his fingers, and his pen, for
he has overstepped his bounds, fabricated lies against God, the
All-Compelling, the All-Powerful, the Almighty, the Self-Subsisting, and
committed what has violated the sanctity of the Cause among His
servants. Even the inhabitants of the nethermost prison will disown him,
how much more God's servants who are near to Him!

\subsection{362}\label{section-857}

\begin{otherlanguage}{persian}

بگو اوّلاً تو آن اخبار را ندانسته و بمعنى آن ابدا مطّلع نيستى هم‌چنانكه ملل
قبل هم نظر بعدم ادراك معانى كلمات الهيّه از شاطى بحر احديّه در حين ظهور
مظاهر ربوبيّه توقّف نموده محروم گشته‌اند از جمله در انجيل در علائم ظهور بعد
ميفرمايد شمس در آن يوم تاريك ميشود و قمر از نور ممنوع و كواكب بر ارض
ميريزند و بهمين عبارات اهل انجيل الى حال اقبال بشريعه محمّديّه ننموده‌اند
حال تو چه ميگوئى اين اخبار كذب بوده و يا آن فئه ادراك معانى آن را
ننموده‌اند

\end{otherlanguage}

Bigū avvalan tu ān akhbār rā nadānista va bi-maʿnā-yi ān abadan muṭṭaliʿ
nīstī hamchunānki milal-i qabl ham naẓar bi-ʿadam-i idrāk-i maʿānī-yi
kalimāt-i ilāhiyya az shāṭī-yi baḥr-i aḥadiyya dar ḥīn-i ẓuhūr-i
maẓāhir-i rubūbiyya tavaqquf namūda maḥrūm gashta-and. Az jumla dar
Injīl dar ʿalāyim-i ẓuhūr-i baʿd mīfarmāyad shams dar ān yawm tārīk
mīshavad va qamar az nūr mamnūʿ va kavākib bar arḍ mīrīzand. Va bi-hamīn
ʿibārāt ahl-i Injīl ilā ḥāl iqbāl bi-sharīʿa-yi Muḥammadiyya
nanamūda-and. Ḥāl tu chi mīgūyī īn akhbār kidhb būda va yā ān fiʾa
idrāk-i maʿānī-yi ān rā nanamūda-and.

Say: First, you have not understood those pronouncements and are wholly
unaware of their meaning, just as the peoples of the past, through
failing to comprehend the meanings of the divine words, held back from
the shore of the ocean of oneness at the appearance of the
Manifestations of Lordship and were deprived. For example, concerning
the signs of the next Manifestation, the Gospel says: ``The sun will be
darkened on that day, the moon will be deprived of light, and the stars
will fall upon the earth.'' Because of these very expressions, the
people of the Gospel have not, to this day, turned toward the religion
of Muhammad. Now, what do you say? Were these pronouncements false, or
did those people fail to comprehend their meanings?

\subsection{363}\label{section-858}

\begin{otherlanguage}{persian}

و هم‌‌چنين ذكر خاتم النّبيين در فرقان و هم‌‌چنين در علائم ظهور بعد و علائم
قيامت از انفطار سماء و انشقاق ارض و نسف جبال و قيام اموات از قبور ظاهره
و امثال اين بيانات كه از علائم قيامت بوده و كلّ اليوم بهمين عبارات از
مالك اسماء و صفات ممنوعند حال چه ميكنى تصديق مينمائى يا تكذيب چه كه بر
حسب ظاهر آنچه مرقوم شده ظاهر نگشته يا بايد بگوئى كلّ كذب است و يا بگوئى
كه معانى آن را ادراك ننموده‌اند و در نصّ خاتم النّبيين نميدانم چه خواهى گفت
كه نقطه اولى تصريحا فرموده كه نبی‌ از اوّل لا اوّل آمد و الى آخر لا آخر
خواهد آمد مع آنكه محال دانسته‌اند كه بعد از حضرت رسول نفسى بيايد و ادّعاى
نبوّت نمايد

\end{otherlanguage}

Va hamchunīn dhikr-i Khātam al-Nabiyyīn dar Furqān va hamchunīn dar
ʿalāyim-i ẓuhūr-i baʿd va ʿalāyim-i qiyāmat az infiṭār-i samāʾ va
inshiqāq-i arḍ va nasf-i jibāl va qiyām-i amvāt az qubūr-i ẓāhira va
amthāl-i īn bayānāt ki az ʿalāyim-i qiyāmat būda va kull al-yawm
bi-hamīn ʿibārāt az Mālik-i asmāʾ va ṣifāt mamnūʿand. Ḥāl chi mīkunī
taṣdīq mīnamāyī yā takdhīb, chi ki bar ḥasab-i ẓāhir ānchi marqūm shuda
ẓāhir nagashta. Yā bāyad bigūyī kull kidhb ast va yā bigūyī ki maʿānī-yi
ān rā idrāk nanamūda-and. Va dar naṣṣ-i Khātam al-Nabiyyīn namīdānam chi
khvāhī guft ki Nuqṭa-yi Ūlā taṣrīḥan farmūda ki nabī az avval-i lā avval
āmad va ilā ākhir-i lā ākhir khvāhad āmad maʿa ānki muḥāl dānista-and ki
baʿd az Ḥaḍrat-i Rasūl nafsī biyāyad va iddiʿā-yi nubuwwat namāyad.

Likewise, there is the mention of the ``Seal of the Prophets'' in the
Qur'an, and, among the signs of the next Manifestation and the signs of
the Resurrection, the splitting of heaven, the cleaving of the earth,
the scattering of the mountains, the rising of the dead from outward
graves, and similar statements describing the signs of the Resurrection.
Today, all are barred by these very expressions from the Sovereign of
names and attributes. Now what will you do---affirm or deny them? For,
in their outward sense, what was written has not appeared. You must
either say that all is false or say that they have not comprehended its
meanings. And I do not know what you will say about the explicit
designation ``Seal of the Prophets,'' for the Primal Point has expressly
stated that prophets have come from a beginning that has no beginning
and will come until an end that has no end, although people have deemed
it impossible for anyone to come after the Messenger and claim
prophethood.

\subsection{364}\label{section-859}

\begin{otherlanguage}{persian}

بگو اى غافل بدانكه احدى بر معانى كلمات الهيّه على ما هى عليه مطّلع نبوده و
نخواهد بود مگر مظهر علم و حكمت ربّانى چنانچه نقطه بيان تصريحا ميفرمايند
كه احدى بر جواهر معانى الهيّه كه در خزاين كلمات بيان مستور است مطّلع نشده
و نخواهد شد مگر نفس ظهور بعد و اگر كلّ ملل معانى كلمات را ادراك مينمودند
هرگز از مشارق وحى و الهام محروم نميشدند چنانچه اليوم محتجبين اهل بيان
نظر بعدم ادراك ما نزل فى البيان از جمال رحمن محروم مانده‌اند واللّه اگر
فى‌الجمله بشعور آيند از عرفان بلكه از انفس خود منقطع شوند

\end{otherlanguage}

Bigū ay ghāfil bidānki aḥadī bar maʿānī-yi kalimāt-i ilāhiyya ʿalā mā
hiya ʿalayhi muṭṭaliʿ nabūda va nakhvāhad būd magar maẓhar-i ʿilm va
ḥikmat-i rabbānī. Chunānchi Nuqṭa-yi Bayān taṣrīḥan mīfarmāyand ki aḥadī
bar javāhir-i maʿānī-yi ilāhiyya ki dar khazāyin-i kalimāt-i Bayān
mastūr ast muṭṭaliʿ nashuda va nakhvāhad shud magar nafs-i ẓuhūr-i baʿd.
Va agar kull-i milal maʿānī-yi kalimāt rā idrāk mīnamūdand hargiz az
mashāriq-i vaḥy va ilhām maḥrūm namīshudand. Chunānchi al-yawm
muḥtajibīn-i ahl-i Bayān naẓar bi-ʿadam-i idrāk-i mā nazala fī al-Bayān
az Jamāl-i Raḥmān maḥrūm mānda-and. Va-Allāh agar fī al-jumla bi-shuʿūr
āyand az ʿirfān balki az anfus-i khud munqaṭiʿ shavand.

Say: O heedless one! Know that no one has ever been or will ever be
acquainted with the true meanings of the divine words except the
Manifestation of divine knowledge and wisdom. The Point of the Bayán
expressly states that no one has attained or will attain knowledge of
the gems of divine meaning concealed in the treasuries of the words of
the Bayán except the next Manifestation Himself. Had all peoples
comprehended the meanings of the words, they would never have been
deprived of the Dawning-Places of revelation and inspiration. Today,
likewise, those among the people of the Bayán who remain veiled have
been deprived of the Beauty of the All-Merciful through failing to
comprehend what was revealed in the Bayán. By God, if they attained even
a little awareness, they would renounce their knowledge---indeed, their
very selves.

\subsection{365}\label{section-860}

\begin{otherlanguage}{persian}

اى غافل از كوثر معانى كلمات الهى را كذب مدان و عبث مشمر و خود را عاجز از
ادراك آن مشاهده نما و از عرفان هر كدام كه خود را عاجز مشاهده نمائى از
منبع و معدن علم لدنّى سؤال نما كه شايد از فيوضات بحر معانى لب‌تشنه نمانى و
در بحور وساوس شيطانى غريق و فانى نگردى حرفى از آنچه نازل شده لغو نبوده و
نخواهد بود بلكه لآلىٴ معانى كلّ در هر حرفى از كلمات منزله سبحانى مستور
بوده و لايقدر ان يخرجها الّا ذو قدرة عظيم و اگر هم بخواهى نعوذ باللّه در
كلمات منزله و ظهورات قدرتيّه لغو و فطور ملاحظه نمائى انفس محتجبه اولى است
قل فانصب مرآتك فى مقابلة وجهك لتشهد جلدك و تعرف مافيه لعلّ توفّق على
الرّجوع و تكون تائباً منيباً معترفاً مقرّاً مذعناً بخطيئاتك و سيّئاتك و انّه لهو
الغفور الرّحيم

\end{otherlanguage}

Ay ghāfil az kawthar-i maʿānī kalimāt-i ilāhī rā kidhb madān va ʿabath
mashmur va khud rā ʿājiz az idrāk-i ān mushāhida namā va az ʿirfān-i har
kudām ki khud rā ʿājiz mushāhida namāyī az manbaʿ va maʿdan-i ʿilm-i
ladunnī suʾāl namā ki shāyad az fuyūḍāt-i baḥr-i maʿānī lab-tishna
namānī va dar buḥūr-i vasāvis-i shayṭānī gharīq va fānī nagardī. Ḥarfī
az ānchi nāzil shuda laghv nabūda va nakhvāhad būd balki laʾāliʾ-i
maʿānī-yi kull dar har ḥarfī az kalimāt-i munzala-yi subḥānī mastūr būda
va lā yaqdiru an yukhrijahā illā dhū qudratin ʿaẓīm. Va agar ham
bikhvāhī naʿūdhu bi-Allāh dar kalimāt-i munzala va ẓuhūrāt-i qudratiyya
laghv va fuṭūr mulāḥiẓa namāyī anfus-i muḥtajiba awlā ast. Qul fa-inṣub
mirʾātaka fī muqābalati wajhika li-tashhada jildaka wa-taʿrifa mā fīhi
laʿalla tuwaffaqu ʿalā al-rujūʿ wa-takūna tāʾiban munīban muʿtarifan
muqirran mudhʿinan bi-khaṭīʾātika wa-sayyiʾātika wa-innahu la-huwa
al-Ghafūr al-Raḥīm.

O you who are heedless of the heavenly river of meanings! Do not regard
the divine words as false or count them as futile. Rather, recognize
your own inability to comprehend them, and concerning any that you find
yourself unable to understand, ask the Source and Wellspring of innate
knowledge, so that you may not remain thirsty for the outpourings of the
ocean of meanings or drown and perish in the seas of satanic misgivings.
Not a single letter of what has been revealed has been or will ever be
futile. Rather, the pearls of the meanings of all things lie concealed
in every letter of the divinely revealed words, and none can draw them
forth except one possessed of mighty power. And if, God forbid, you wish
to discern futility and flaws in the revealed words and manifestations
of power, the souls that remain veiled are more deserving of such
scrutiny. Say: Set your mirror before your face, that you may behold
your skin and know what is in it, so that you may be enabled to return
and become repentant, turning back, acknowledging, confessing, and
admitting your transgressions and misdeeds. He is indeed the
Ever-Forgiving, the Merciful.

\subsection{366}\label{section-861}

\begin{otherlanguage}{persian}

قوله كبر غلّه هنوز هيچ از شئونات و اوامر او بظهور نرسيده و علائم ظهور
قائم را كه اهل فرقان قائلند اگر چه علائم حين ظهور ظاهر شده ومردم نفهميده
امّا علائم بعد از ظهور ظاهر نشده اگر اهل فرقان ايراد نمايند كه علائم بعد
از ظهور چه چيز است جواب چه خواهد بود از نفهميدن علائم حين ظهور و ظاهر
نشدن علائم بعد از ظهور موجب حجاب خلق از حقّ خواهد بود و اين از حقّ دور است
خود ملاحظه فرموده غرس حبّه بدون ظهور اصول و اوراق و اغصان او حبّه ديگر كه
ثمره و قيامت او است چگونه ممكن است انتهى

\end{otherlanguage}

Qawluhu kabura ghilluhu: Hanūz hīch az shuʾūnāt va avāmir-i ū bi-ẓuhūr
narasīda va ʿalāyim-i ẓuhūr-i Qāʾim rā ki ahl-i Furqān qāyiland agar chi
ʿalāyim-i ḥīn-i ẓuhūr ẓāhir shuda va mardum nafahmīda ammā ʿalāyim-i
baʿd az ẓuhūr ẓāhir nashuda. Agar ahl-i Furqān īrād namāyand ki
ʿalāyim-i baʿd az ẓuhūr chi chīz ast javāb chi khvāhad būd. Az
nafahmīdan-i ʿalāyim-i ḥīn-i ẓuhūr va ẓāhir nashudan-i ʿalāyim-i baʿd az
ẓuhūr mūjib-i ḥijāb-i khalq az Ḥaqq khvāhad būd va īn az Ḥaqq dūr ast.
Khud mulāḥiẓa farmūda ghars-i ḥabba bidūn-i ẓuhūr-i uṣūl va avrāq va
aghṣān-i ū ḥabba-yi dīgar ki thamara va qiyāmat-i ū ast chigūna mumkin
ast. Intahā.

His words---great is his malice: ``As yet, none of His functions or
commands have appeared. As for the signs of the appearance of the Qá'im
affirmed by the people of the Qur'an, although the signs accompanying
His appearance have appeared and people have not understood them, the
signs following His appearance have not appeared. If the people of the
Qur'an object and ask what the signs following His appearance are, what
will the answer be? The failure to understand the signs accompanying His
appearance, and the failure of the signs following His appearance to
appear, will cause people to be veiled from God, and this is far removed
from God. Consider for yourself: when a seed is planted, how can there
be another seed, which is its fruit and resurrection, without its roots,
leaves, and branches appearing?'' End of quotation.

\subsection{367}\label{section-862}

\begin{otherlanguage}{persian}

آنقدر شده كه متّصل بر خود ردّ مينويسيد و مستشعر نيستيد چه كه نوشته علائم
ظهور قائم در حين ظهور ظاهر شده و مردم نفهميده‌اند حال ملاحظه كن اين عدم
فهم از چه رو بوده شكّى نيست كه از عدم ظهورات توهميّه ناس بوده چنانچه از
علائم ظهور آنچه ذكر شده از آن عبارات معانى توهّم نموده و چون موافق توهّمات
خود در عالم ملكيّه مشاهده ننمودند لذا بعرفان نقطه بيان فايز نشدند چنانچه
خود تو و امثال تو هم عارف نبوده تا آنكه منزل بيان كلمات قبل را تفسير
فرمودند و بعد از بيان هيكل رحمن قليلى تصديق نموده و اكثر انكار و اين
فعلى است كه اليوم تو و امثال تو از ملأ بيان عاملند چه كه معنى كلمات
الهيّه را ادراك ننموده‌ايد و لا يشعر صيحه ميزنيد كه نبايد اين ظهور ظاهر
شود هنوز تكميل خلق بيان نشده و اگر ادراك مينموديد البتّه باين كلمات محتجب
نميشديد حقّ را مثل خود عاجز مشاهده نموده‌ايد فو‌اللّه اگر نفسى بقدرت حقّ موقن
باشد ابدا در امثال اين كلمات توقّف نمينمايد چه كه قادر است در كلمه از
كلمات خود كلّ كتب را عود فرمايد و هم‌‌چنين بدء فرمايد انّه لهو المقتدر على
ما یشآء و انّه لهو العليم الحكيم

\end{otherlanguage}

Ānqadr shuda ki muttaṣil bar khud radd mīnivīsīd va mustashʿir nīstīd
chi ki nivishta ʿalāyim-i ẓuhūr-i Qāʾim dar ḥīn-i ẓuhūr ẓāhir shuda va
mardum nafahmīda-and. Ḥāl mulāḥiẓa kun īn ʿadam-i fahm az chi rū būda.
Shakkī nīst ki az ʿadam-i ẓuhūrāt-i tavahhumiyya-yi nās būda chunānchi
az ʿalāyim-i ẓuhūr ānchi dhikr shuda az ān ʿibārāt maʿānī tavahhum
namūda va chūn muvāfiq-i tavahhumāt-i khud dar ʿālam-i mulkiyya
mushāhida nanamūdand lidhā bi-ʿirfān-i Nuqṭa-yi Bayān fāyiz nashudand.
Chunānchi khud-i tu va amthāl-i tu ham ʿārif nabūda tā ānki Munzil-i
Bayān kalimāt-i qabl rā tafsīr farmūdand va baʿd az bayān-i Haykal-i
Raḥmān qalīlī taṣdīq namūda va akthar inkār. Va īn fiʿlī ast ki al-yawm
tu va amthāl-i tu az malaʾ-i Bayān ʿāmiland chi ki maʿnā-yi kalimāt-i
ilāhiyya rā idrāk nanamūda-īd va lā yashʿur ṣayḥa mīzanīd ki nabāyad īn
ẓuhūr ẓāhir shavad hanūz takmīl-i khalq-i Bayān nashuda. Va agar idrāk
mīnamūdīd albatta bi-īn kalimāt muḥtajib namīshudīd. Ḥaqq rā mithl-i
khud ʿājiz mushāhida namūda-īd. Fa-va-Allāh agar nafsī bi-qudrat-i Ḥaqq
mūqin bāshad abadan dar amthāl-i īn kalimāt tavaqquf namīnamāyad chi ki
qādir ast dar kalima az kalimāt-i khud kull-i kutub rā ʿawd farmāyad va
hamchunīn badʾ farmāyad. Innahu la-huwa al-Muqtadir ʿalā mā yashāʾ
wa-innahu la-huwa al-ʿAlīm al-Ḥakīm.

It has reached the point where you continually write refutations of
yourselves without realizing it. For you have written that the signs of
the appearance of the Qá'im appeared at the time of His appearance, but
people did not understand them. Now consider why this failure to
understand occurred. Undoubtedly it was because the things people
imagined did not appear: they imagined meanings for the expressions
describing the signs of His appearance, and when they did not observe
anything in the material world corresponding to their imaginings, they
failed to attain recognition of the Point of the Bayán. You yourself and
others like you likewise did not understand until the Revealer of the
Bayán interpreted the earlier words. Even after the explanation of the
Temple of the All-Merciful, only a few affirmed, while most denied. This
is precisely what you and others like you among the people of the Bayán
are doing today, for you have not comprehended the meaning of the divine
words and unthinkingly cry out that this Manifestation must not appear,
since the creation of the Bayán has not yet been completed. Had you
understood, you would certainly not have been veiled by these words. You
have regarded God as powerless like yourselves. By God, if a soul were
certain of God's power, it would never hesitate over such words, for He
is able, through a single word of His, to bring all the Books back and
likewise to originate them. He is indeed powerful over whatever He
wills, and He is indeed the All-Knowing, the All-Wise.

\subsection{368}\label{section-863}

\begin{otherlanguage}{persian}

علائم بعد از ظهور كلّ ظاهر شده و لكن شما ادراك ننموده چنانچه بقول خود شما
علايم حين ظهور را فئه فرقان عارف نشده‌اند و هر علامتى كه بعد از ظهور آن
را لازم دانسته و عارف بآن نشده چنانچه نوشته از منبع علم الهى سؤال نما تا
موقن شوى بآنكه كلّ ظاهر شده و لكن امثال آن نفوس از ادراك كلّ محتجب
مانده‌اند

\end{otherlanguage}

ʿAlāyim-i baʿd az ẓuhūr kull ẓāhir shuda va lākin shumā idrāk nanamūda
chunānchi bi-qawl-i khud-i shumā ʿalāyim-i ḥīn-i ẓuhūr rā fiʾa-yi Furqān
ʿārif nashuda-and. Va har ʿalāmatī ki baʿd az ẓuhūr ān rā lāzim dānista
va ʿārif bi-ān nashuda chunānchi nivishta az manbaʿ-i ʿilm-i ilāhī suʾāl
namā tā mūqin shavī bi-ānki kull ẓāhir shuda va lākin amthāl-i ān nufūs
az idrāk-i kull muḥtajib mānda-and.

All the signs following the Manifestation have appeared, but you have
not understood them, just as, by your own account, the people of the
Qur'an did not recognize the signs accompanying His appearance.
Concerning every sign that you have deemed necessary after the
Manifestation but have not recognized, as you have written, ask the
Source of divine knowledge, so that you may become certain that all have
appeared, but that such souls have remained veiled from comprehending
them all.

\subsection{369}\label{section-864}

\begin{otherlanguage}{persian}

و اصل علامت ظهور بعد از قائم ظهور حسينى بوده شما اوّل را بقول خود مذعنيد
و ثانى را منكر در همين كلمه اگر تفكّر نمائى بر اسرار ما لانهايه مطّلع شوى
و در ظهور حسينى بعد از قائم نوشته‌اند در كتب اخبار خود شما كه كلّ انبیاء و
مرسلين در ظلّ رايت حضرت در آيند و هم‌‌چنين نوشته اند كه قائم رجعت نمايد و
او هم بحضرت ملحق شود اگر قدرى در اين عبارات تأمّل نمائيد ادراك مينمائيد
كه بيان كتب مطابق است بما نزل فى البيان آنقدر بدان كه علائم بعد از ظهور
كلّ ظاهر شده بعضى را ادراك ننموده‌ايد و بعضى را انكار نموده‌ايد كذلك امركم
هوئكم ان انتم تعلمون

\end{otherlanguage}

Va aṣl-e ʿalāmat-e ẓohūr baʿd az Qāʾim ẓohūr-e Ḥosaynī būde shomā avval
rā be-qowl-e khod modhʿenīd va thānī rā monker dar hamīn kalame agar
tafakkor namāʾī bar asrār-e mā lā-nehāye moṭṭaleʿ shavī va dar ẓohūr-e
Ḥosaynī baʿd az Qāʾim neveshte-and dar kotob-e akhbār-e khod-e shomā ke
koll-e anbiyā va morsalīn dar ẓell-e rāyat-e Ḥażrat dar āyand va
ham-chenīn neveshte-and ke Qāʾim rajʿat namāyad va ū ham be-Ḥażrat
molḥaq shavad agar qadrī dar īn ʿebārāt taʾammol namāʾīd edrāk mīnamāʾīd
ke bayān-e kotob moṭābeq ast be-mā nazala fī al-Bayān ānqadr bedān ke
ʿalāʾem baʿd az ẓohūr koll ẓāher shode baʿżī rā edrāk nanemūde-īd va
baʿżī rā enkār nemūde-īd kadhālika amarakum hawākum in antum taʿlamūn.

The principal sign of the Manifestation after the Qāʾim was the
appearance of Ḥusayn. By your own account, you acknowledge the first and
deny the second. If you reflect on this very statement, you will
discover boundless mysteries. Concerning the appearance of Ḥusayn after
the Qāʾim, it is written in your own books of traditions that all the
Prophets and Messengers will come beneath His banner. It is likewise
written that the Qāʾim will return and will also join Him. If you
reflect a little on these passages, you will understand that what the
books say accords with what has been revealed in the Bayān. Know at
least this: all the signs following the Manifestation have appeared;
some you have not understood, and others you have denied. Thus have your
desires commanded you, if you only knew.

\subsection{370}\label{section-865}

\begin{otherlanguage}{persian}

حال چه ميگوئيد در نقطه بيان روح ما سواه فداه كه ابداً علامتى ذكر نفرموده
و منصوص فرموده‌اند كه هروقت ظاهر شود بايد كلّ تصديق نمايند و شكر الهى بجا
آورند و بظهور بنفسه ناظر شوند و بما يظهر من عنده و صريح ميفرمايند او
مختار است بهر نوع كه اراده فرمايد ظاهر شود بر احدى نيست كه باو احتجاج
نمايد و وصاياى لايحصى فرموده ملأ بيان را كه مباد ببيان و مانزل فيه بر
نفس رحمن احتجاج نمايند مع‌ذلك تو بوهميّات نفس خود احتجاج نموده و ذكر تكميل
دوره كرده نفس حقّ ظاهر و بآيات نقطه ناطق و ميفرمايد بظهور من تكميل دوره
بيان شده و ميشود مع‌ذلك آنقدر محتجبی‌ كه نفس حقّ را بقدر خلق ندانسته ميگوئى
كه مخصوص بايد در خلق تكميل شود و حال آنكه امر حقّ بدونش تكميل نشده و
نخواهد شد بلكه بنفس مقدّسش تكميل شده و ميشود لو انتم تشعرون كلّما تكلّمت به
بغى منك على اللّه المقتدر المهيمن القيّوم

\end{otherlanguage}

Ḥāl che mīgūʾīd dar Noqṭe-ye Bayān rūḥ-e mā sevāh fedāh ke abadan
ʿalāmatī dhekr nafarmūde va manṣūṣ farmūde-and ke har vaqt ẓāher shavad
bāyad koll taṣdīq namāyand va shokr-e elāhī be-jā āvarand va be-ẓohūr
be-nafseh nāẓer shavand va be-mā yaẓharu min ʿindih va ṣarīḥ mīfarmāyand
ū mokhtār ast be-har nowʿ ke erāde farmāyad ẓāher shavad bar aḥadī nīst
ke bā ū eḥtejāj namāyad va vaṣāyā-ye lā-yoḥṣā farmūde malaʾ-e Bayān rā
ke mabād be-Bayān va mā nazala fīh bar nafs-e Raḥmān eḥtejāj namāyand
maʿa-dhālek to be-vahmiyyāt-e nafs-e khod eḥtejāj nemūde va dhekr-e
takmīl-e dowre karde nafs-e Ḥaqq ẓāher va be-āyāt-e Noqṭe nāṭeq va
mīfarmāyad be-ẓohūr-e man takmīl-e dowre-ye Bayān shode va mīshavad
maʿa-dhālek ānqadr moḥtajebī ke nafs-e Ḥaqq rā be-qadr-e khalq nadāneste
mīgūʾī ke makhṣūṣ bāyad dar khalq takmīl shavad va ḥāl ānke amr-e Ḥaqq
bedūnash takmīl nashode va nakhāhad shod balke be-nafs-e moqaddasash
takmīl shode va mīshavad law antum tashʿurūn kullu mā takallamta bihi
baghyun minka ʿalā Allāh al-Muqtadir al-Muhaymin al-Qayyūm.

Now what do you say concerning the Point of the Bayān---may the souls of
all besides Him be sacrificed for Him---who mentioned no sign whatsoever
and expressly stated that whenever He appears, all must affirm Him, give
thanks to God, and look to the Manifestation Himself and to whatever
appears from Him? He explicitly says that He is free to appear in
whatever manner He wills, and that no one has the right to dispute with
Him. He gave the people of the Bayān countless counsels, lest they use
the Bayān and what was revealed in it to dispute with the very Person of
the All-Merciful. Yet you have advanced your own vain imaginings as
arguments and spoken of the completion of the dispensation. God Himself
has appeared, uttering the verses of the Point, and says, ``Through My
appearance the dispensation of the Bayān has been and is being
completed.'' Yet you are so veiled that you have not considered God
Himself to be equal even to His creation, and you say that it must
specifically be completed among the creatures. But the Cause of God has
not been and will not be completed without Him; rather, it has been and
is being completed through His holy Person, if you would only perceive.
Everything you have said is an act of rebellion on your part against
God, the All-Powerful, the All-Protecting, the Self-Subsisting.

\subsection{371}\label{section-866}

\begin{otherlanguage}{persian}

نميدانم اگر نقطه بيان علائم ذكر ميفرمودند چه ميكرديد و چه وارد مياورديد
آيا در بيان كلمه مشاهده شده كه اگر نفسى بآيات الهيّه و شئونات عزّ ربانيّه
ظاهر شود بر قتلش قيام نمائيد و يا بمفتريات نفسيّه رداء عزّ احديّه را
بيالائيد بلكه در جميع بيان ميفرمايد اگر نفسى ظاهر شود و بگويد من محبوب
شما هستم من دون آنكه از او برهانى ظاهر شود بر احدى نيست كه بر او اعتراض
نمايد

\end{otherlanguage}

Namīdānam agar Noqṭe-ye Bayān ʿalāʾem dhekr mīfarmūdand che mīkardīd va
che vāred mīāvardīd āyā dar Bayān kalame moshāhede shode ke agar nafsī
be-āyāt-e elāhiyye va shoʾūnāt-e ʿezz-e rabbāniyye ẓāher shavad bar
qatl-ash qiyām namāʾīd va yā be-moftarayāt-e nafsiyye redā-ye ʿezz-e
aḥadiyye rā beyālāʾīd balke dar jamīʿ-e Bayān mīfarmāyad agar nafsī
ẓāher shavad va begūyad man maḥbūb-e shomā hastam min dūn-e ānke az ū
borhānī ẓāher shavad bar aḥadī nīst ke bar ū eʿterāż namāyad.

I do not know what you would have done or what you would have inflicted
had the Point of the Bayān mentioned signs. Has a single word been found
in the Bayān saying that if someone appears with divine verses and
manifestations of divine glory, you should rise up to kill Him or stain
the robe of the glory of divine oneness with your self-inspired
fabrications? Rather, throughout the Bayān He says that if someone
appears and says, ``I am your Beloved,'' even without producing any
proof, no one has the right to object to Him.

\subsection{372}\label{section-867}

\begin{otherlanguage}{persian}

و حال شمس معانى از افق عنايت رحمانى مشرق و بآيات الهى ناطق بشأنيكه احدى
بر احصاى آن قادر نه مع‌ذلك بعضى ميگوئيد تكميل دوره قبل نشده بعضى ميگوئيد
زود ظاهر شده

\end{otherlanguage}

Va ḥāl Shams-e Maʿānī az ofoq-e ʿenāyat-e raḥmānī moshreq va be-āyāt-e
elāhī nāṭeq be-shaʾnī ke aḥadī bar eḥṣā-ye ān qāder na maʿa-dhālek baʿżī
mīgūʾīd takmīl-e dowre-ye qabl nashode baʿżī mīgūʾīd zūd ẓāher shode.

And now the Sun of Meanings is shining from the horizon of the grace of
the All-Merciful, uttering divine verses in such abundance that no one
can count them. Yet some of you say that the previous dispensation has
not been completed, and others say that He has appeared too soon.

\subsection{373}\label{section-868}

\begin{otherlanguage}{persian}

بگو اى بی‌‌انصافان خود جمال ابهى ميفرمايد كه معنى پارسى آن اين است كه اى
قوم بودم يكى از شما بر بستر خود آرميده و نائم نفحات رحمانيّه از يمين عزّ
احديّه وزيد و مرا بيدار نمود و بثنآء نفسش ناطق فرمود مابين آسمانها و
زمينها اى قوم قسم بخدا كه افتخار بر نفسى ننمودم و منكر آيات الهى نبوده و
نيستم و مذعنم آنچه مابين دستهاى شما است از بيان و موقنم بپروردگارى كه
جميع را خلق فرمود رحم نمائيد بر من و متعرّض نشويد و اگر باين شريعه عزّ
ربّانيّه اقبال نمينمائيد اعتراض مكنيد ايقوم اگر اختيار بدست من بود البتّه
خود را از ابصر شما ستر مينمودم و راضى نميشدم بر نفس خود آنچه را كه احدى
از ممكنات حمل آن ننموده چه كه روزى نگذشته مگر آنكه بسبّ و شماتت و استهزاء
مبتلى بوده‌ام و لكن روح القدس نطق فرمود در صدرم و روح اعظم تكلّم نمود
بلسانم و اين نيست از جانب من بلكه از جانب قادرى است كه قدرتش بر همه
اشیاء محيط و فايق است

\end{otherlanguage}

Begū ey bī-enṣāfān khod-e Jamāl-e Abhā mīfarmāyad ke maʿnī-ye pārsī-ye
ān īn ast ke ey qowm būdam yekī az shomā bar bestar-e khod ārāmīde va
nāʾem nafaḥāt-e raḥmāniyye az yamīn-e ʿezz-e aḥadiyye vazīd va marā
bīdār nemūd va be-thanā-ye nafsash nāṭeq farmūd mābeyn-e āsmānhā va
zamīnhā ey qowm qasam be-Khodā ke eftekhār bar nafsī nanemūdam va
monker-e āyāt-e elāhī nabūde va nīstam va modhʿenam ānche mābeyn-e
dasthā-ye shomā ast az Bayān va mūqenam be-Parvardegārī ke jamīʿ rā
khalq farmūd raḥm namāʾīd bar man va motaʿarreż nashavīd va agar be-īn
sharīʿe-ye ʿezz-e rabbāniyye eqbāl namīnamāʾīd eʿterāż nakonīd ey qowm
agar ekhtiyār be-dast-e man būd albatte khod rā az abṣor-e shomā satr
mīnamūdam va rāżī namīshodam bar nafs-e khod ānche rā ke aḥadī az
momkenāt ḥaml-e ān nanemūde che ke rūzī nagodhashte magar ānke be-sabb
va shamātat va estehzā mobtalā būde-am va lāken Rūḥ al-Qodos noṭq farmūd
dar ṣadram va Rūḥ-e Aʿẓam takallom nemūd be-lesānam va īn nīst az
jāneb-e man balke az jāneb-e Qāderī ast ke qodratash bar hame-ye ashyā
moḥīṭ va fāʾeq ast.

Say: You who are unjust! The All-Glorious Beauty Himself speaks, and the
meaning of His words in Persian is this: ``People! I was one of you,
resting asleep on My bed. The breezes of the All-Merciful blew from the
right hand of the glory of divine oneness, awakened Me, and caused Me to
voice His praise between the heavens and the earths. People! I swear by
God that I have not vaunted Myself over anyone, nor have I denied or do
I now deny the verses of God. I acknowledge what you possess of the
Bayān, and I am certain of the Lord who created all. Have mercy on Me
and do not assail Me. If you will not turn toward this river of divine
glory, do not object. People! Had the choice been Mine, I would
certainly have concealed Myself from your sight and would not have
accepted for Myself what no contingent being has ever borne. For not a
day has passed without My suffering abuse, taunts, and ridicule. But the
Holy Spirit spoke within My breast, and the Most Great Spirit spoke
through My tongue. This is not from Me, but from One whose power
encompasses and surpasses all things.''

\subsection{374}\label{section-869}

\begin{otherlanguage}{persian}

اى قوم قرب و بعد ظهور دست من و شما هيچكدام نبوده و نيست بلكه در دست قدرت
الهى بوده و خواهد بود و در قبضه اقتدار او مقبوض ظاهر فرمود هر قسم كه
خواسته بترسيد از خدا و نباشيد از ظلم كنندگان ايقوم از كاس كلمات الهى
كوثر معانى بياشاميد اگر روايح قدس تجريد از او استشمام نموديد در آن وقت
انصاف دهيد و وارد نياوريد آنچه را كه بآن بجزع آيند اهل غرفات رضوان و
نباشيد از غفلت كنندگان

\end{otherlanguage}

Ey qowm qorb va boʿd-e ẓohūr dast-e man va shomā hīch-kodām nabūde va
nīst balke dar dast-e qodrat-e elāhī būde va khāhad būd va dar qabże-ye
eqtedār-e ū maqbūż ẓāher farmūd har qesm ke khāste betarsīd az Khodā va
nabāshīd az ẓolm-konandegān ey qowm az kās-e kalemāt-e elāhī kowthar-e
maʿānī beyāshāmīd agar ravāyeḥ-e qods-e tajrīd az ū esteshmām nemūdīd
dar ān vaqt enṣāf dehīd va vāred nayāvarīd ānche rā ke be-ān be-jazaʿ
āyand ahl-e ghorofāt-e Reżvān va nabāshīd az ghaflat-konandegān.

``People! Whether the Manifestation comes sooner or later has never been
and is not in either My hands or yours. Rather, it has been and will be
in the hand of divine power, held within the grasp of His might. He has
brought it forth however He wished. Fear God and do not be among the
wrongdoers. People! Drink the Kawthar of meanings from the cup of the
words of God. If you inhale from it the holy fragrances of detachment,
then be fair, and do not inflict that which would cause the inhabitants
of the chambers of Paradise to cry out in anguish. Do not be among the
heedless.''

\subsection{375}\label{section-870}

\begin{otherlanguage}{persian}

اى قوم براستى ميگويم آنچه از پروردگار شما مأمور شده‌ام و نفحات صدق از كذب
بسى واضح و معلوم است بترسيد از خدا و نباشيد از دورماندگان اى قوم هرگز
متعرّض نفسى نشده‌ام و با اهل بيان در كمال رحمت و رافت و خضوع و خشوع معاشرت
نموده‌ام و در ایّاميكه جميع قلوب از سطوت مضطرب بوده و جميع نفوس از سياط
قهر متزلزل بنفس خود مابين اعدا ايستاده‌ام و در كلّ حين از سهام اعدا بر
صدرم وارد شده و احدى قادر بر احصاى ضرّم نبوده مگر پروردگار عالم دانا بسا
از شبها كه جميع راحت بوده و اين عبد در سلاسل و اغلال و بسا ایّام كه شما
در بيوت خود ساكن و اين عبد در دست اعداء مبتلا و كلّ بر آنچه ذكر شده شاهد
و گواهيد اگر بصدق تكلّم نمائيد و باشيد از راست گويان

\end{otherlanguage}

Ey qowm be-rāstī mīgūyam ānche az Parvardegār-e shomā maʾmūr shode-am va
nafaḥāt-e ṣedq az kedhb basī vāżeḥ va maʿlūm ast betarsīd az Khodā va
nabāshīd az dūr-māndegān ey qowm hargez motaʿarreż-e nafsī nashode-am va
bā ahl-e Bayān dar kamāl-e raḥmat va raʾfat va khożūʿ va khoshūʿ
moʿāsherat nemūde-am va dar ayyāmī ke jamīʿ-e qolūb az saṭvat możṭareb
būde va jamīʿ-e nofūs az seyāṭ-e qahr motazalzel be-nafs-e khod mābeyn-e
aʿdā īstāde-am va dar koll-e ḥīn az sehām-e aʿdā bar ṣadram vāred shode
va aḥadī qāder bar eḥṣā-ye żorram nabūde magar Parvardegār-e ʿālam-e
dānā basā az shabhā ke jamīʿ rāḥat būde va īn ʿabd dar salāsel va aghlāl
va basā ayyām ke shomā dar boyūt-e khod sāken va īn ʿabd dar dast-e aʿdā
mobtalā va koll bar ānche dhekr shode shāhed va govāhīd agar be-ṣedq
takallom namāʾīd va bāshīd az rāst-gūyān.

``People! I speak truthfully what I have been commanded by your Lord,
and the breezes of truth are readily and clearly distinguished from
falsehood. Fear God and do not be among those who remain far away.
People! I have never assailed anyone, and I have associated with the
people of the Bayān with the utmost mercy, kindness, humility, and
lowliness. In the days when every heart was troubled by the onslaught
and every soul trembled beneath the scourges of wrath, I Myself stood
among the enemies. At every moment their arrows struck My breast, and no
one but the knowing Lord of the world could count My afflictions. How
many nights everyone was at ease while this servant was in chains and
fetters! How many days you dwelt in your homes while this servant
suffered in the hands of the enemies! You are all witnesses to what has
been related, if you speak honestly and are among the truthful.''

\subsection{376}\label{section-871}

\begin{otherlanguage}{persian}

اى قوم شما مقتدريد يا پروردگار عالميان اگر او مقتدر است ظاهر فرمود مظهر
نفس خود را هر قسم كه اراده فرمود و اگر ذنب عبد اين است اين نبوده از نزد
او بلكه از نزد خالق زمين و آسمان اى قوم بشنويد ندايم را و تسنيم حيوان را
از كأس رحمن در اين ایّام بياشاميد و نباشيد از محرومان انتهى

\end{otherlanguage}

Ey qowm shomā moqtaderīd yā Parvardegār-e ʿālamiyān agar ū moqtader ast
ẓāher farmūd Maẓhar-e nafs-e khod rā har qesm ke erāde farmūd va agar
dhanb-e ʿabd īn ast īn nabūde az nazd-e ū balke az nazd-e Khāleq-e zamīn
va āsmān ey qowm beshnavīd nedāyam rā va tasnīm-e ḥayavān rā az kaʾs-e
Raḥmān dar īn ayyām beyāshāmīd va nabāshīd az maḥrūmān intahā.

``People! Are you powerful, or is the Lord of the worlds? If He is
powerful, He has brought forth the Manifestation of His own Self in
whatever manner He willed. If this is the servant's sin, it did not
originate with Him, but with the Creator of earth and heaven. People!
Hear My call, drink in these days the Tasnīm of life from the cup of the
All-Merciful, and do not be among the deprived.'' End of quotation.

\subsection{377}\label{section-872}

\begin{otherlanguage}{persian}

اينست معنى كلمات الهيّه كه اين عبد بقدر ادراك خود ذكر نموده كه لعلّ ملتفت
شويد كه امر در دست احدى نبوده و نيست و بر كيفيّت ظهور احدى مطّلع نبوده و
نخواهد بود كلّ در قبضه قدرت و اقتدار اسيرند و او است حاكم هر چه بخواهد
ظاهر ميفرمايد و كلّ مأمور كه در ظهور بنفس ظهور ناظر باشند نه بكلمات و
اشارات

\end{otherlanguage}

Īn ast maʿnī-ye kalemāt-e elāhiyye ke īn ʿabd be-qadr-e edrāk-e khod
dhekr nemūde ke laʿall moltafet shavīd ke amr dar dast-e aḥadī nabūde va
nīst va bar keyfiyyat-e ẓohūr aḥadī moṭṭaleʿ nabūde va nakhāhad būd koll
dar qabże-ye qodrat va eqtedār asīrand va ū ast Ḥākem har che bekhāhad
ẓāher mīfarmāyad va koll maʾmūr ke dar ẓohūr be-nafs-e ẓohūr nāẓer
bāshand na be-kalemāt va eshārāt.

This is the meaning of the divine words, which this servant has related
according to his understanding, so that you may perhaps recognize that
the matter has never been and is not in anyone's hands, and that no one
has known or will know the manner of the Manifestation's appearance. All
are held captive in the grasp of His power and might, and He is the
Ruler: He brings forth whatever He wishes. All are commanded, at the
Manifestation, to look to the Manifestation Himself, not to words and
allusions.

\subsection{378}\label{section-873}

\begin{otherlanguage}{persian}

فو‌اللّه چون آيات الهى من غير تعطيل از سماء مشيّت لايزالى نازل نفس معرض
باللّه و سيّد محمّد با هم ملحق شده تدبير نمودند كه شايد ببعضى مفتريات ناس
را از اين سلسبيل محروم نمايند اينست كه يكى باستانبول دويد و يكى متّصلاً در
اين ارض بمكاتيب مجعوله مشغول شد و اللّه الّذى لا اله الّا هو كه با كمال
رافت و شفقت بآن دو نفس و كلّ سلوك فرموده‌اند عجب است با آنكه خود شما از
اوصاف جمال مبارك شنيده و اكثرى از اهل ديار بچشم خود ديده‌اند مع‌ذلك بدو
كلمه مفتريات شجره فساد جميع را از قلب محو نموده و باين شدّت بر اعراض و
اعتراض بر خواسته نسئل اللّه بأن يوفّقك لعرفان نفسه و ينقطعك عن الاشارات
لعلّ تطّلع بما كنت محتجباً عنه و تكون من الفائزين

\end{otherlanguage}

Fa-wa-Allāh chūn āyāt-e elāhī min ghayr-e taʿṭīl az samā-ye mashiyyat-e
lā-yazālī nāzel nafs-e moʿreż bi-Allāh va Seyyed Moḥammad bā ham molḥaq
shode tadbīr nemūdand ke shāyad be-baʿżī moftarayāt nās rā az īn
salsabīl maḥrūm namāyand īn ast ke yekī be-Estānbūl davīd va yekī
mottaṣelan dar īn arż be-makātīb-e majʿūle mashghūl shod wa-Allāh
alladhī lā ilāha illā huwa ke bā kamāl-e raʾfat va shafaqat be-ān do
nafs va koll solūk farmūde-and ʿajab ast bā ānke khod-e shomā az owṣāf-e
Jamāl-e Mobārak shenīde va aktharī az ahl-e deyār be-chashm-e khod
dīde-and maʿa-dhālek be-do kalame moftarayāt-e shajare-ye fasād jamīʿ rā
az qalb maḥv nemūde va be-īn sheddat bar eʿrāż va eʿterāż bar khāste
nasʾalu Allāha bi-an yuwaffiqaka li-ʿirfāni nafsihi wa-yanqaṭiʿaka ʿan
al-ishārāt laʿallaka taṭṭaliʿu bi-mā kunta muḥtajaban ʿanhu wa-takūna
min al-fāʾizīn.

By God! When divine verses were descending without interruption from the
heaven of the everlasting Will, the one who turned away from God and
Siyyid Muḥammad joined forces and devised a plan, hoping to deprive
people of this Salsabīl through certain fabrications. Thus one rushed to
Istanbul, while the other busied himself continually in this land with
fabricated letters. By God, besides whom there is no other god, He has
treated those two souls and everyone else with the utmost kindness and
compassion. It is astonishing that, although you yourself have heard of
the qualities of the Blessed Beauty and most of the people of the region
have seen them with their own eyes, you have nevertheless erased them
all from your heart because of a few fabricated words from the tree of
corruption, and have risen up with such vehemence in rejection and
opposition. We ask God to enable you to recognize His own Self and to
detach you from allusions, so that you may discover what you have been
veiled from and be among those who attain.

\subsection{379}\label{section-874}

\begin{otherlanguage}{persian}

اگر بر احوال اين عباد مطّلع شوى كه چه گذشته در اين سنين توقّف در ادرنه و
منصف باشيد نوحه و ندبه نمائيد خداى واحد شاهد است كه در هر ساعتى از آن
قيامتى بر پا شده چه كه در هر حين آن نفسيكه او را من دون اللّه باغواى
شيطان حقّ ميدانى بامرى شنيع و فعلى قبيح ظاهر شده تو را قسم ميدهم بخداوندى
كه جميع اشیآء را باراده خود خلق فرموده كه هيچ بعقلت ميرسد كه ايشان از
نفسى كه منسوب است نان منع نمايند چه فايده كه نبوده و نيستيد محض تضييع
امر از براى لقمه نان بمدينه كبيره دويدند و چه مقدار شكايت در اين ارض
نموده‌اند و اللّه كذب نميگويم و مقصودى ندارم ولى قلب از اين محترق است كه
چنين نفوس را ربّ دانسته‌ايد و خود و ناس را از حقّ منيع محروم نموده‌ايد آخر
ایّام اين عباد است اگر چه شما توهّم نموده‌ايد كه بجهت دنيا باين ارض آمده‌ايم
و لكن خدا شاهد و گواه است كه بچه قسم آمده‌ايم و مقصود چه بوده شما خود
ميدانيد كه اين عبد در مقرّ خود محتاج بنان نبوده كه شهر بشهر براى نان بدود
بلكه نفحات رحمن اين عباد را جذب نموده لو انت من العارفين

\end{otherlanguage}

Agar bar aḥvāl-e īn ʿebād moṭṭaleʿ shavī ke che godhashte dar īn senīn-e
tavaqqof dar Edrene va monṣef bāshīd nowḥe va nodbe namāʾīd Khodā-ye
vāḥed shāhed ast ke dar har sāʿatī az ān qiyāmatī bar pā shode che ke
dar har ḥīn ān nafsī ke ū rā min dūni Allāh be-eghvā-ye sheyṭān Ḥaqq
mīdānī be-amrī shanīʿ va feʿlī qabīḥ ẓāher shode to rā qasam mīdaham
be-Khodāvandī ke jamīʿ-e ashyā rā be-erāde-ye khod khalq farmūde ke hīch
be-ʿaqlat mīresad ke īshān az nafsī ke mansūb ast nān manʿ namāyand che
fāʾede ke nabūde va nīstīd maḥż-e tażyīʿ-e amr az barā-ye loqme nān
be-Madīne-ye Kabīre davīdand va che meqdār shekāyat dar īn arż
nemūde-and va-Allāh kedhb namīgūyam va maqṣūdī nadāram valī qalb az īn
moḥtareq ast ke chonīn nofūs rā Rabb dāneste-īd va khod va nās rā az
Ḥaqq-e manīʿ maḥrūm nemūde-īd ākher-e ayyām-e īn ʿebād ast agar che
shomā tavahhom nemūde-īd ke be-jehat-e donyā be-īn arż āmade-īm va lāken
Khodā shāhed va govāh ast ke be-che qesm āmade-īm va maqṣūd che būde
shomā khod mīdānīd ke īn ʿabd dar maqarr-e khod moḥtāj-e nān nabūde ke
shahr be-shahr barā-ye nān bedavad balke nafaḥāt-e Raḥmān īn ʿebād rā
jadhb nemūde law anta min al-ʿārifīn.

If you learned of the circumstances of these servants and what has
happened during these years of residence in Adrianople, and were
fair-minded, you would weep and lament. The one God is witness that in
every hour of those years an upheaval has arisen, for at every moment
that person whom, through Satan's seduction, you regard as God instead
of God has committed some abominable deed or shameful act. I adjure you
by the God who created all things through His will: can you imagine that
He would withhold bread from someone associated with Him? But what use
is it? You were not here and are not here. Solely to discredit the
Cause, they rushed to the Great City for a morsel of bread, and how many
complaints they have made in this land! By God, I do not lie, and I have
no ulterior purpose. But my heart burns because you have regarded such
souls as the Lord and have deprived yourselves and the people of the
unapproachable True One. These are the last days of these servants.
Although you have imagined that we came to this land for worldly ends,
God is witness to how we came and what our purpose was. You yourselves
know that this servant was not in need of bread in his own dwelling,
that he should run from city to city for bread. Rather, the breezes of
the All-Merciful have drawn these servants onward, if you are among
those who understand.

\subsection{380}\label{section-875}

\begin{otherlanguage}{persian}

بارى حرمت خود و اين عباد و امر همه را ضايع كردند و مع‌ذلك به بنده نصيحت
مينمائيد و موعظه مينويسيد كه چرا بچنين نفوس مؤمن نشدى و موقن نيستى
رذالتشان بمرتبه رسيده كه مجعولا سياهه داده‌اند كه چه قدر از مهاجرين
رفته‌اند و چه قدر هم در اين ارض ساكنند و آنچه ببعضى از فقرا ميرسيد قطع
نموده و باسم خود گذرانده‌اند فو‌اللّه انسان زقّوم بخورد بهتر از چنين نانى
است چنانچه تفصيل سياهه نفوس كه آنجا داده باين ارض فرستاده‌اند و احدى هم
بر آن تفصيل مطّلع نبوده جز شجره فساد و من معه

\end{otherlanguage}

Bārī ḥormat-e khod va īn ʿebād va amr hame rā żāyeʿ kardand va
maʿa-dhālek be-bande naṣīḥat mīnamāʾīd va mowʿeẓe mīnevīsīd ke cherā
be-chonīn nofūs moʾmen nashodī va mūqen nīstī radhālatashān be-martabe
resīde ke majʿūlan siyāhe dāde-and ke che qadr az mohājerīn rafte-and va
che qadr ham dar īn arż sākenand va ānche be-baʿżī az foqarā mīresīd
qaṭʿ nemūde va be-esm-e khod godharānde-and fa-wa-Allāh ensān zaqqūm
bekhorad behtar az chonīn nānī ast chenānche tafṣīl-e siyāhe-ye nofūs ke
ānjā dāde be-īn arż ferestāde-and va aḥadī ham bar ān tafṣīl moṭṭaleʿ
nabūde joz shajare-ye fasād va man maʿahu.

In short, they have ruined their own reputation, that of these servants,
and that of the Cause---all of them. Yet you counsel this servant and
write admonitions asking why he has not come to believe in such souls
and is not certain of them. Their baseness has reached the point that
they have submitted a fabricated list stating how many of the emigrants
have left and how many still reside in this land. They have cut off what
used to reach some of the poor and have entered it in their own names.
By God, it would be better to eat Zaqqūm than such bread! Indeed, they
have sent to this land the details of the list of persons submitted
there, and no one knew of those details except the tree of corruption
and those with him.

\subsection{381}\label{section-876}

\begin{otherlanguage}{persian}

و كاش اين باقى را هم باسم خود ميگذراندند فو‌اللّه كه اين عباد ممنون ميشديم
آنقدر بدان كه بحقّ حقّ اكثرى از مهاجرين لباس ندارند و مع‌ذلك باين نفوس كه
للّه حركت نموده‌اند و بناله و نوحه با جمال احديّه باين شطر توجّه نموده‌اند
نسبت دنيا داده‌اند و از براى خود انقطاع ثابت نموده‌اند فيا ليت كنت مطّلعاً
بهم و بما عندهم و بما ارتكبوا فى الحيوة الباطلة لابقآء رياساتهم و كنت من
المطّلعين

\end{otherlanguage}

Va kāsh īn bāqī rā ham be-ism-e khod mī-godharāndand. Fa-wa-llāh ke īn
ʿebād mamnūn mī-shodīm. Ān-qadr bedān ke be-ḥaqq-e Ḥaqq aktharī az
mohājerīn lebās nadārand va maʿa-dhālek be-īn nofūs ke li-llāh ḥarakat
namūde-and va be-nāle va nowḥe bā Jamāl-e Aḥadīye be-īn shaṭr tavajjoh
namūde-and nesbat-e donyā dāde-and va az barā-ye khod enqeṭāʿ thābet
namūde-and. Fa-yā layta kunta muṭṭaliʿan bihim wa-bi-mā ʿindahum
wa-bi-mā irtakabū fī al-ḥayāt al-bāṭila li-ibqāʾ riyāsātihim wa-kunta
min al-muṭṭaliʿīn.

If only they had entered this remainder in their own names as well! By
God, we servants would have been grateful. Know at least this much: by
the truth of the True One, most of the emigrants have no clothes. Yet
they have attributed worldly motives to these souls, who set out for
God's sake and turned toward this region with the Beauty of Divine
Unity, weeping and lamenting, while claiming detachment for themselves.
If only you knew them, what they possess, and what they have committed
in this vain life to preserve their positions of leadership, and were
among those who are informed!

\subsection{382}\label{section-877}

\begin{otherlanguage}{persian}

و آنچه در غرس حبّه و تفصيل آن نوشته بوديد از جميع اين كلمات چنان مفهوم
ميشود كه بحرفى از علم الهى مطّلع نشده‌ايد و مقصود از ظهور را ندانسته‌ايد كه
چه بوده و چه خواهد بود حال از مثل اعلى گذشته و بمثل ادنى اعتراضاً على
اللّه ناطق شده تسليم مينمائيم كه همين است كه شما نوشته‌ايد و مذكور ميداريد
آيا اين حبّه سقايه لازم داشته و يا ندارد ميفرمايد منم آن هبوب ارياح احديّه
كه از يمين رضوان اعرفوا اللّه باللّه در هبوب است و تربيت كلّ شیء منوط باين
هبوب عزّ رحمانى و نفحه قدس روحانى بوده و خواهد بود و ميفرمايد منم آن كوثر
حيوانى كه باو حبّه‌‌هاى معرفت و معانى و سنبلات علم و حكمت صمدانى انبات
مينمايد و من غير هبوب اين ارياح و اين ماء عذب فرات هيچ حبّه علميّه از
رضوان صدور بريّه نروئيده و نخواهد روئيد و مع‌ذلك ميگوئيد و مينويسيد و لا
تشعرون ما تقولون

\end{otherlanguage}

Va ānche dar ghars-e ḥabbe va tafṣīl-e ān neveshte būdīd, az jamīʿ-e īn
kalemāt chenān mafhūm mī-shavad ke be-ḥarfī az ʿelm-e elāhī moṭṭaleʿ
nashode-īd va maqṣūd az ẓohūr rā nadāneste-īd ke che būde va che khāhad
būd. Ḥāl az mathal-e aʿlā godhashte va be-mathal-e adnā iʿtirāḍan ʿalā
Allāh nāṭeq shode. Taslīm mī-namāyīm ke hamīn ast ke shomā neveshte-īd
va madhkūr mī-dārīd. Āyā īn ḥabbe seqāye lāzem dāshte va yā nadārad?
Mī-farmāyad: Manam ān hobūb-e aryāḥ-e aḥadīye ke az yamīn-e Reḍvān-e
iʿrifū Allāh bi-llāh dar hobūb ast va tarbīyat-e koll-e shayʾ manūṭ
be-īn hobūb-e ʿezz-e raḥmānī va nafḥe-ye qods-e rūḥānī būde va khāhad
būd. Va mī-farmāyad: Manam ān Kowthar-e ḥayavānī ke be-ū ḥabbe-hā-ye
maʿrefat va maʿānī va sonbolāt-e ʿelm va ḥekmat-e ṣamadānī enbāt
mī-namāyad. Va men gheyr-e hobūb-e īn aryāḥ va īn māʾ-e ʿadhb-e forāt
hīch ḥabbe-ye ʿelmīye az Reḍvān-e ṣodūr-e barīye narūʾīde va nakhāhad
rūʾīd. Va maʿa-dhālek mī-gūʾīd va mī-nevīsīd wa-lā tashʿurūna mā
taqūlūn.

As for what you wrote about planting the seed and your explanation of
it, all these words suggest that you have not understood a single letter
of divine knowledge and have not known what the purpose of revelation
has been or will be. Now, passing over the most exalted analogy, you
have spoken through the lowliest analogy in protest against God. Let us
concede that it is just as you have written and stated. Has this seed
needed watering, or does it not? He says: ``I am that blowing of the
winds of divine unity that comes from the right hand of the Paradise of
`Know God through God'; the nurture of all things has depended, and will
depend, upon this glorious, merciful breeze and this holy, spiritual
breath.'' And He says: ``I am that life-giving Kawthar through which the
seeds of understanding and inner meanings, and the ears of divine
knowledge and wisdom, sprout.'' Without the blowing of these winds and
this sweet, refreshing water, no seed of knowledge has sprouted, or will
sprout, in the paradise of human hearts. Yet you speak and write, and do
not perceive what you are saying.

\subsection{383}\label{section-878}

\begin{otherlanguage}{persian}

اين انظر مشركه را بگذار و ببصر توحيد ملاحظه كن تا جميع مظاهر را نفس
واحده مشاهده نمائى و شرايع كلّ را شريعت واحده بينى جميع اين اختلافات
باقتضاى وقت و مقام و استعدادات مظاهر وجود بوده و خواهد بود خدا را اراضى
طيّبه مباركه مقدّسه بوده كه حين القاى بزر ظهور سنبلات اخرى اقرب من لمح
البصر مشاهده ميشود و اين مقام عالم امر است و خلق از ادراك آن عاجز چنانچه
ميفرمايد انّما امره اذا اراد شيئاً ان يقول له كن فيكون و اين اراضى است كه
ميفرمايد يخرج نباته باذن ربّه و لكن اراضى نفوس مشركه هرگز نميرويد الّا
آنچه را حقّ جلّ ذكره خبر داده بقوله تعالى لايخرج الّا نكداً در اراضى آن نفوس
اگر صدهزار سال هم درظهور تأخير شود بقول شما هرگز سنبلات طيّبه از آن اراضى
جرزه نروئيده و نخواهد روئيد

\end{otherlanguage}

Īn {[}corrupted text{]} moshreke rā be-godhār va be-baṣar-e towḥīd
molāḥeẓe kon tā jamīʿ-e maẓāher rā nafs-e vāḥede moshāhede namāʾī va
sharāyeʿ-e koll rā sharīʿat-e vāḥede bīnī. Jamīʿ-e īn ekhtelāfāt
be-eqteḍā-ye vaqt va maqām va esteʿdādāt-e maẓāher-e vojūd būde va
khāhad būd. Khodā rā arāḍī-ye ṭayyebe-ye mobārake-ye moqaddase būde ke
ḥīn-e elqā-ye bazr ẓohūr-e sonbolāt-e okhrā aqrab min lamḥ al-baṣar
moshāhede mī-shavad va īn maqām ʿālam-e amr ast va khalq az edrāk-e ān
ʿājez. Chenānche mī-farmāyad: Innamā amruhu idhā arāda shayʾan an yaqūla
lahu kun fa-yakūn. Va īn arāḍī ast ke mī-farmāyad: Yakhruju nabātuhu
bi-idhni rabbihi. Va lāken arāḍī-ye nofūs-e moshreke hargaz namī-rūyad
ellā ānche rā Ḥaqq jalla dhikruhu khabar dāde bi-qawlihi taʿālā: Lā
yakhruju illā nakidan. Dar arāḍī-ye ān nofūs agar ṣad-hezār sāl ham dar
ẓohūr taʾkhīr shavad be-qowl-e shomā, hargaz sonbolāt-e ṭayyebe az ān
arāḍī-ye joroze narūʾīde va nakhāhad rūʾīd.

Set aside this {[}corrupted original text{]} that associates partners
with God, and look with the eye of divine unity, so that you may see all
the Manifestations as one soul and the religious laws of all as one law.
All these differences have arisen, and will arise, according to the
requirements of time and circumstance and the capacities of the
manifestations of existence. God has had good, blessed, holy soils in
which, when the seed is sown, the appearance of further ears of grain is
witnessed in less than the blink of an eye. This station is the realm of
command, which creation is powerless to comprehend. As He says: ``His
command, when He desires a thing, is only to say to it, `Be,' and it
is.'' It is of these soils that He says: ``Its vegetation comes forth by
the permission of its Lord.'' But the soils of souls that associate
partners with God never bring forth anything except what the True
One---exalted be His mention---has announced in His exalted words: ``It
brings forth only scantily.'' Even if, as you would have it, the
revelation were delayed a hundred thousand years, good ears of grain
have never grown, and will never grow, from the barren soils of those
souls.

\subsection{384}\label{section-879}

\begin{otherlanguage}{persian}

البتّه شما خواهيد گفت دوره تمام نشده و كامل نگشته فويل للّذينهم يعرفون
نعمة اللّه ثمّ ينكرونها بچشم خود ملاحظه مينمائيد كه سنبلات معانى مكنونه از
حبّه هاى كلمات الهيّه روئيده و سر سبز و خرّم در رضوان احديّه مشهودند مع‌ذلك
انكار نموده و مينمائيد و شايد كه شما نديده و نشنيده و مطّلع نيستيد
انشآءاللّه كه چنين باشد و از آن اشخاصى نباشى كه بچشم خود ديده و انكار
نموده‌اند و بگوش خود شنيده و معرض گشته‌اند

\end{otherlanguage}

Albatte shomā khāhīd goft dowre tamām nashode va kāmel nagashte.
Fa-waylun li-lladhīna hum yaʿrifūna niʿmata Allāh thumma yunkirūnahā.
Be-cheshm-e khod molāḥeẓe mī-namāyīd ke sonbolāt-e maʿānī-ye maknūne az
ḥabbe-hā-ye kalemāt-e elāhīye rūʾīde va sar-sabz va khorram dar Reḍvān-e
Aḥadīye mashhūdand. Maʿa-dhālek enkār namūde va mī-namāyīd va shāyad ke
shomā nadīde va nashenīde va moṭṭaleʿ nīstīd. In shāʾ Allāh ke chonīn
bāshad va az ān ashkhāṣī nabāshī ke be-cheshm-e khod dīde va enkār
namūde-and va be-gūsh-e khod shenīde va moʿreḍ gashte-and.

Of course you will say that the cycle has not ended or reached
completion. Woe, then, to those who recognize God's favor and then deny
it! You see with your own eyes that ears of hidden meanings have sprung
from the seeds of divine words and stand verdant and flourishing in the
Paradise of Divine Unity. Yet you have denied, and continue to deny.
Perhaps, however, you have neither seen nor heard and are unaware. God
willing, this is so, and you are not among those who have seen with
their own eyes and denied, and heard with their own ears and turned
away.

\subsection{385}\label{section-880}

\begin{otherlanguage}{persian}

از اينها گذشته ذات اقدسش لم‌يزل مقدّس از امثال بوده و خواهد بود چه كه مثل
و امثال بآن مقرّ راه نداشته و ندارد و اينكه امثال ذكر ميشود نظر بعرفان
ناس است كه شايد از آن امثله حدوديّه ادراك بعضى از امور الهيّه را على قدر
شأنهم نمايند حقّ منيع اگر اراده فرمايد در يك حين بيك هبوب ارياح فضليّه
جميع اشجار يابسه را بطراز ربيع مزيّن فرمايد و بيك كلمه جميع حبّات وجود را
بسنبلات حكمت و علم ظاهر مينمايد گويا آيينه در پيش خود گذاشته‌ايد و آنچه
در خود مشاهده ميكنيد از عجز و جهل و شئونات ضدّيّه و دلالات حمقيّه كلّ را بحقّ
راجع نموده و مينمائيد تب الى اللّه يا ايّها النّائم و كاش از نائمين محسوب
بودى چه كه نائم خود آسوده است و هم ناس از او آسوده‌اند چنانچه گفته‌اند
عملى از براى ظالم بهتر از نوم نبوده و نيست چه كه در حالت نوم خود و ناس
هر دو آسوده‌اند تفكّر فى ظلمك يا ايّها الظّالم الّذى بظلمك ناحت الأشیاء كلّها
و اضطربت نفوس المقدّسين خف عن الله و لاتتّبع هواك فالق ما عندك ثمّ اتّخذ ما
عند اللّه بقوّة من عنده لعلّ يرتدّ اليك بصر العناية و يجعلك من المنصفين

\end{otherlanguage}

Az īnhā godhashte Dhāt-e Aqdasash lam-yazal moqaddas az amthāl būde va
khāhad būd, che ke methl va amthāl be-ān maqarr rāh nadāshte va nadārad.
Va īnke amthāl dhekr mī-shavad naẓar be-ʿerfān-e nās ast ke shāyad az ān
amthele-ye ḥodūdīye edrāk-e baʿḍī az omūr-e elāhīye rā ʿalā qadri
shaʾnihim namāyand. Ḥaqq-e manīʿ agar erāde farmāyad dar yek ḥīn be-yek
hobūb-e aryāḥ-e faḍlīye jamīʿ-e ashjār-e yābese rā be-ṭerāz-e rabīʿ
mozayyan farmāyad va be-yek kaleme jamīʿ-e ḥabbāt-e vojūd rā
be-sonbolāt-e ḥekmat va ʿelm ẓāher mī-namāyad. Gūyā āyīne dar pīsh-e
khod godhāshte-īd va ānche dar khod moshāhede mī-konīd az ʿajz va jahl
va shoʾūnāt-e ḍeddīye va dalālāt-e ḥomqīye koll rā be-Ḥaqq rājeʿ namūde
va mī-namāyīd. Tub ilā Allāh yā ayyuhā al-nāʾim. Va kāsh az nāʾemīn
maḥsūb būdī, che ke nāʾem khod āsūde ast va ham nās az ū āsūde-and.
Chenānche gofte-and ʿamalī az barā-ye ẓālem behtar az nowm nabūde va
nīst, che ke dar ḥālat-e nowm khod va nās har do āsūde-and. Tafakkar fī
ẓulmika yā ayyuhā al-ẓālim alladhī bi-ẓulmika nāḥat al-ashyāʾ kulluhā
wa-iḍṭarabat nufūs al-muqaddasīn. Khaf ʿan Allāh wa-lā tattabiʿ hawāka
fa-alqi mā ʿindaka thumma ittakhidh mā ʿinda Allāh bi-quwwatin min
ʿindihi laʿalla yartadda ilayka baṣar al-ʿināya wa-yajʿalaka min
al-munṣifīn.

Apart from all this, His most holy Essence has always been, and will
remain, sanctified above analogies, for likenesses and comparisons have
never had, and do not have, access to that seat. Analogies are mentioned
only with regard to human understanding, so that through these limited
comparisons people may perhaps comprehend some divine matters according
to their capacity. If the unapproachable True One so wills, in a single
instant, with one breath of the winds of grace, He can adorn all
withered trees with the raiment of spring; with one word He brings forth
all the seeds of existence as ears of wisdom and knowledge. It is as
though you have placed a mirror before yourself and have attributed, and
continue to attribute, to the True One everything you observe in
yourself: impotence, ignorance, contrary dispositions, and signs of
foolishness. Repent to God, you who are asleep! If only you could be
counted among the sleeping, for a sleeping person is himself at rest,
and people are at rest from him. As has been said, no deed has been or
is better for a tyrant than sleep, since while he sleeps both he and the
people are at rest. Reflect on your injustice, you tyrant whose
injustice has caused all things to lament and the souls of the
sanctified to tremble. Fear God and do not follow your desires. Cast
away what you possess, then take hold of what is with God through a
power from Him, so that the gaze of providence may return to you and He
may make you one of the fair-minded.

\subsection{386}\label{section-881}

\begin{otherlanguage}{persian}

قول من ناح بقوله قبائل العالمين عرض ديگر آنكه تعريفات حضرت باب بجهت شخص
معيّن بهاء است و بجهت مسمّاى بعينه است و يا بجهت اسم بهاء اگر بجهت مسمّاى
بعينه و شخص بخصوص است اين اوّل حرف است و محتاج بدليل است بيان فرموده كه
در كجاى بيان و بچه عبارت است و اگر بجهت اسم بهاء است اين اسم از اسماء
اللّه است تعريفات اسماء اللّه اختصاص باسم بهاء ندارد اسم جلال و جمال را
نيز تعريف فرموده و هكذا كلّ اسماء اللّه را پس هر ميرزا كمال و ميرزا جمال
بايد ادّعاى اين امر نمايد علاوه ديّان و ميرزا غوغا و شيخ اسمعيل و حاجى ملاّ
هاشم هم ادّعاى اين امر نمودند و ادّعاى آنها باطل شده پس محض ادّعا كفايت
نخواهد كرد و الّا اين حضرات بايد همان موعود باشند انتهى

\end{otherlanguage}

Qawlu man nāḥa bi-qawlihi qabāʾil al-ʿālamīn: ʿArḍ-e dīgar ānke
taʿrīfāt-e Ḥaḍrat-e Bāb be-jehat-e shakhṣ-e moʿayyan-e Bahā ast va
be-jehat-e mosammā-ye be-ʿeyneh ast va yā be-jehat-e esm-e Bahā? Agar
be-jehat-e mosammā-ye be-ʿeyneh va shakhṣ-e be-khoṣūṣ ast īn avval-e
ḥarf ast va moḥtāj be-dalīl ast. Bayān farmūde ke dar kojā-ye Bayān va
be-che ʿebārat ast. Va agar be-jehat-e esm-e Bahā ast īn esm az asmāʾ-e
Allāh ast. Taʿrīfāt-e asmāʾ-e Allāh ekhteṣāṣ be-esm-e Bahā nadārad.
Esm-e Jalāl va Jamāl rā nīz taʿrīf farmūde va hākadhā koll-e asmāʾ-e
Allāh rā. Pas har Mīrzā Kamāl va Mīrzā Jamāl bāyad eddeʿā-ye īn amr
namāyad. ʿElāve Dayyān va Mīrzā Ghowghā va Shaykh Esmāʿīl va Ḥājī Mollā
Hāshem ham eddeʿā-ye īn amr namūdand va eddeʿā-ye ānhā bāṭel shode. Pas
maḥḍ-e eddeʿā kefāyat nakhāhad kard va ellā īn ḥaḍarāt bāyad hamān
mowʿūd bāshand. Intahā.

The words of him whose speech caused the tribes of the worlds to lament:
``Another submission is this: are the praises of the Báb directed toward
the particular person of Bahá, toward a specifically designated
individual, or toward the name Bahá? If they are directed toward a
specifically designated individual and a particular person, this is the
very point at issue and requires proof. Explain where in the Bayán this
appears and in what words. If they are directed toward the name Bahá,
this name is one of the names of God. Praise of the names of God is not
confined to the name Bahá. He also praised the names Jalál and Jamál,
and likewise all the names of God. Thus every Mírzá Kamál and Mírzá
Jamál should lay claim to this Cause. Moreover, Dayyán, Mírzá Ghawghá,
Shaykh Ismáʿíl, and Ḥájí Mullá Háshim also laid claim to this Cause, and
their claims have been shown to be false. A mere claim, therefore, will
not suffice; otherwise, these gentlemen must be that very Promised
One.'' End of quotation.

\subsection{387}\label{section-882}

\begin{otherlanguage}{persian}

قل فامسك قلمك يا ايّها الّذى صرت محروماً عن رداء الأدب و الانصاف و بما جرى
من قلمك قد ضيّعت حرمة اللّه بين عباده و اذا يتبرّء منك قلمك و لوح الّذى رقمت
عليه ما رقمت تاللّه لو كانت اذنك واعية لسمعت فى حين الّذى كنت مشتغلاً بهذه
الأذکار‌ حنين اهل ملأ الأعلی ثمّ ضجيج الأصفیا‌ء ثمّ نوح الأولیاء‌ ثمّ صريخ كلّ
الأشیآء عمّا خلق بقول اللّه الملك المتعالى العلىّ العظيم خذ لسانك يا عدوّ
اللّه ثمّ امسك قلمك و دع لوحك فاستحى عن اللّه الّذى قد خلقك بأمر من عنده
لهذا الظّهور الّذى ظهر بالحقّ و انفطرت عنه سموات ايقانكم يا ملأ الغافلين ا
تظنّ بأن يعزب من علمه من شیء لا فو‌ربّ العالمين

\end{otherlanguage}

Qul fa-amsik qalamaka yā ayyuhā alladhī ṣirta maḥrūman ʿan ridāʾ al-adab
wa-l-inṣāf wa-bi-mā jarā min qalamika qad ḍayyaʿta ḥurmata Allāh bayna
ʿibādihi wa-idhā yatabarraʾu minka qalamuka wa-lawḥ alladhī raqamta
ʿalayhi mā raqamta. Ta-llāh law kānat udhunuka wāʿiya la-samiʿta fī ḥīn
alladhī kunta mushtaghilan bi-hādhihi al-adhkār ḥanīna ahl malaʾ al-aʿlā
thumma ḍajīja al-aṣfiyāʾ thumma nawḥa al-awliyāʾ thumma ṣarīkha kull
al-ashyāʾ ʿammā khuliqa bi-qawli Allāh al-Malik al-Mutaʿālī al-ʿAlī
al-ʿAẓīm. Khudh lisānaka yā ʿaduwwa Allāh thumma amsik qalamaka wa-daʿ
lawḥaka fa-istaḥi ʿan Allāh alladhī qad khalaqaka bi-amrin min ʿindihi
li-hādhā al-ẓuhūr alladhī ẓahara bi-l-ḥaqq wa-infaṭarat ʿanhu samāwāt
īqānikum yā malaʾ al-ghāfilīn. A-taẓunnu bi-an yaʿzuba min ʿilmihi min
shayʾ? Lā fa-wa-rabb al-ʿālamīn.

Say: Restrain your pen, you who have been deprived of the robe of
courtesy and fairness! Through what has flowed from your pen, you have
violated the sanctity of God among His servants. Now your pen disowns
you, as does the tablet on which you inscribed what you inscribed. By
God, had your ear been attentive, you would have heard, while you were
occupied with these statements, the lament of the company on high, then
the outcry of the chosen ones, then the mourning of the holy ones, then
the cry of all things created by the word of God, the King, the
Transcendent, the Most High, the Most Great. Hold your tongue, enemy of
God, restrain your pen, and leave your tablet! Be ashamed before God,
who created you by a command from Him for this Revelation, which has
appeared in truth and through which the heavens of your certitude have
been torn apart, you heedless company! Do you imagine that anything
escapes His knowledge? No, by the Lord of the worlds!

\subsection{388}\label{section-883}

\begin{otherlanguage}{persian}

حين الّذى جرى من قلمك ما جرى لبكت عين علىّ فى سمآء القضا و كان ناظرا اليك
و مخاطباً الى نفسك يا ايّها الشّقىّ قد سوّد اللّه وجهك لانّك ارتكبت ما لا
ارتكبه احد من الأوّلین‌ و لايرتكبه نفس من الآخرین اما نصحناكم فى كلّ
الألواح بأن لا تقولوا فى حين الظّهور لم و بم و انّك يا ايّها الملحد المكّار
قد اتّخذت امر اللّه هزوا و ضربت و تضرب على جسدى فى كلّ حين اسياف الغلّ و
النّفاق و لا تشعر و تظنّ بانّك تكون من المستقيمين فى امرى و باستدلالك تريد
ان تثبت ما ظهر منّى ثمّ تقتل فى كلّ حين نفسى المحبوب و لا تستحيى عن اللّه
الّذى اليه ترجع امور العالمين و بفعلكم قد خجلت بين ملأ المقرّبين ثمّ اصفیاء
اللّه المقدّسين ثمّ انبيائه المصطفين

\end{otherlanguage}

Ḥīn alladhī jarā min qalamika mā jarā la-bakat ʿaynu ʿAlī fī samāʾ
al-qaḍā wa-kāna nāẓiran ilayka wa-mukhāṭiban ilā nafsika: Yā ayyuhā
al-shaqī qad sawwada Allāh wajhaka li-annaka irtakabta mā lā irtakabahu
aḥadun min al-awwalīn wa-lā yartakibuhu nafsun min al-ākhirīn. A-mā
naṣaḥnākum fī kull al-alwāḥ bi-an lā taqūlū fī ḥīn al-ẓuhūr lima
wa-bima? Wa-innaka yā ayyuhā al-mulḥid al-makkār qad ittakhadhta amra
Allāh huzuwan wa-ḍarabta wa-taḍribu ʿalā jasadī fī kull ḥīn asyāfa
al-ghill wa-l-nifāq wa-lā tashʿur wa-taẓunnu bi-annaka takūnu min
al-mustaqīmīn fī amrī. Wa-bi-istidlālika turīdu an tuthbita mā ẓahara
minnī thumma taqtulu fī kull ḥīn nafsī al-Maḥbūb wa-lā tastaḥyī ʿan
Allāh alladhī ilayhi tarjiʿu umūr al-ʿālamīn. Wa-bi-fiʿlikum qad
khajiltu bayna malaʾ al-muqarrabīn thumma aṣfiyāʾ Allāh al-muqaddasīn
thumma anbiyāʾihi al-muṣṭafayn.

When what flowed from your pen came forth, the eye of ʿAlí wept in the
heaven of decree. He gazed upon you and addressed you: ``Wretched one!
God has blackened your face, for you have committed what no one among
the former generations committed and no soul among the later generations
will commit. Did We not counsel you in all the Tablets not to say, at
the time of the Revelation, `Why?' and `By what means?' Yet you,
deceitful unbeliever, have made a mockery of the Cause of God. You have
struck, and continue at every moment to strike, My body with the swords
of malice and hypocrisy, without perceiving it, while imagining yourself
among those who are steadfast in My Cause. With your arguments you seek
to establish what appeared from Me, yet at every moment you slay My
beloved Self, and you feel no shame before God, to whom the affairs of
the worlds return. Through your deeds I have been put to shame before
the company of those brought near, then before God's sanctified chosen
ones, then before His elect Prophets.''

\subsection{389}\label{section-884}

\begin{otherlanguage}{persian}

فو‌اللّه بنسبة نفسك الى البيان قد بكت حروفاته و كلماته و حقايقه و معانيه و
انّك تفرح و تضحك و تلعب و تكون من المستهزئين فسوف يأخذك اللّه بقهر من عنده
و سطوة من لدنه و انّه هو اشدّ المنتقمين اما تفكّرت فى نفسك باىّ حجّة آمنت
بنفسى و باىّ برهان نسبت نفسك الى نفسى اذا تبرّء منك و يتبرّء كلّ من فى
السّموات و الأرض و كان نفسى على ما اقول شهيد انّ الّذى بعث بينكم اما ظهر
بآياتى و ظهورى و سلطنتى و كبريائى و حجّتى و برهانى و دليلى لم كفرتم به و
اقررتم بنفسى لا فو‌نفسى العليم الخبير انّكم ما آمنتم بنفسى و لابآياتى و لا
بما نزل فى ازل الآزال من جبروت اللّه الملك المقتدر العلىّ العليم

\end{otherlanguage}

Fa-wa-llāh bi-nisbati nafsika ilā al-Bayān qad bakat ḥurūfātuhu
wa-kalimātuhu wa-ḥaqāʾiquhu wa-maʿānīhi wa-innaka tafraḥu wa-taḍḥaku
wa-talʿabu wa-takūnu min al-mustahziʾīn. Fa-sawfa yaʾkhudhuka Allāh
bi-qahrin min ʿindihi wa-saṭwatin min ladunhu wa-innahu huwa ashaddu
al-muntaqimīn. A-mā tafakkarta fī nafsika bi-ayyi ḥujjatin āmanta
bi-nafsī wa-bi-ayyi burhānin nasabta nafsaka ilā nafsī? Idhan tabarraʾa
minka wa-yatabarraʾu kullu man fī al-samāwāt wa-l-arḍ wa-kāna nafsī ʿalā
mā aqūlu shahīd. Inna alladhī buʿitha baynakum a-mā ẓahara bi-āyātī
wa-ẓuhūrī wa-salṭanatī wa-kibriyāʾī wa-ḥujjatī wa-burhānī wa-dalīlī?
Lima kafartum bihi wa-aqrartum bi-nafsī? Lā fa-wa-nafsī al-ʿAlīm
al-Khabīr innakum mā āmantum bi-nafsī wa-lā bi-āyātī wa-lā bi-mā nazala
fī azal al-āzāl min jabarūt Allāh al-Malik al-Muqtadir al-ʿAlī al-ʿAlīm.

``By God, because you have associated yourself with the Bayán, its
letters, words, realities, and meanings have wept, while you rejoice,
laugh, play, and join the mockers. Soon God will seize you with His
overpowering might and His crushing force. He is indeed the most severe
of avengers. Have you not reflected within yourself on the testimony
through which you believed in Me, and the proof through which you
associated yourself with Me? Thus everyone in the heavens and on earth
has disowned, and will disown, you, and My own Self is witness to what I
say. Has not the One raised up among you appeared with My verses, My
revelation, My sovereignty, My majesty, My testimony, My proof, and My
evidence? Why have you disbelieved in Him while acknowledging Me? No, by
My own all-knowing, all-informed Self! You have not believed in Me, nor
in My verses, nor in what descended in the eternity of eternities from
the dominion of God, the King, the All-Powerful, the Most High, the
All-Knowing.''

\subsection{390}\label{section-885}

\begin{otherlanguage}{arabic}

اذا فانصف يا ايّها الغافل ثمّ خلّص نفسك فى ساعة للّه ربّك و ربّ العالمين انّ
الّذى يتلو عليكم فى كلّ حين آيات اللّه لو لم تومن به لاتتعرّض عليه خف عن
اللّه انّه قد تجاوز عن ايمانك و ايمان مثلائك لو تكوننّ من الشّاعرين قد خجل
كلّ الأشیآء بين يدى اللّه لنسبة الّتى كانت بينك و بينها لانّ الكلّ ينتهى الى
ظلّ اسمه الصّانع كذلك كان الأمر ان انت من الشّاعرين اذا يتبرّئنّ الأشیآء عن
نسبتك اليها و يستعيذنّ باللّه من نفسك يا ايّها الّذى بفعلك احترقت اكباد
النّبيّين و المرسلين

\end{otherlanguage}

Idhan fa-anṣif yā ayyuhā al-ghāfil thumma khalliṣ nafsaka fī sāʿatin
li-llāh rabbika wa-rabb al-ʿālamīn. Inna alladhī yatlū ʿalaykum fī kull
ḥīn āyāt Allāh law lam tuʾmin bihi lā tataʿarraḍ ʿalayhi. Khaf ʿan Allāh
innahu qad tajāwaza ʿan īmānika wa-īmāni mithlāʾika law takūnanna min
al-shāʿirīn. Qad khajila kull al-ashyāʾ bayna yaday Allāh li-l-nisba
allatī kānat baynaka wa-baynahā li-anna al-kull yantahī ilā ẓill ismihi
al-Ṣāniʿ. Kadhālika kāna al-amr in anta min al-shāʿirīn. Idhan
yatabarraʾanna al-ashyāʾ ʿan nisbatika ilayhā wa-yastaʿidhanna bi-llāh
min nafsika yā ayyuhā alladhī bi-fiʿlika iḥtaraqat akbād al-nabiyyīn
wa-l-mursalīn.

Be fair, then, heedless one, and for one hour devote yourself sincerely
to God, your Lord and the Lord of the worlds. If you do not believe in
the One who recites God's verses to you at every moment, do not assail
Him. Fear God! He has dispensed with your belief and the belief of those
like you, if only you could perceive it. All things have been put to
shame before God because of the relationship between you and them, for
all ultimately rest beneath the shadow of His name, the Maker. Such is
the matter, if you are among those who perceive. Thus all things will
surely disown your association with them and seek refuge in God from
you, you whose deeds have set ablaze the hearts of the Prophets and
Messengers.

\subsection{391}\label{section-886}

\begin{otherlanguage}{persian}

و ما كان مقصودك فيما جرى من قلمك الّا سخريّة و استهزاء على اللّه و انّك
تستهزئ على نفسك و رؤسائك و لاتكون من الشّاعرين كما استهزء عباد من قبلك
على سفرآء اللّه و امنائه و انّك جئت عن ورائه بامر اضطربت عنه اركان
الاطمينان ثمّ قوائم الايقان و تزلزلت نفس السّكون و اهتزّت اركان عرش عظيم

\end{otherlanguage}

Wa-mā kāna maqṣūduka fīmā jarā min qalamika illā sukhriyyatan
wa-istihzāʾan ʿalā Allāh wa-innaka tastahziʾu ʿalā nafsika wa-ruʾasāʾika
wa-lā takūnu mina al-shāʿirīn kamā istahzaʾa ʿibādun min qablika ʿalā
sufarāʾi Allāh wa-umanāʾihi wa-innaka jiʾta ʿan warāʾihi bi-amrin
iḍṭarabat ʿanhu arkānu al-iṭmiʾnān thumma qawāʾimu al-īqān wa-tazalzalat
nafsu al-sukūn wa-ihtazzat arkānu ʿarshin ʿaẓīm.

Your purpose in what flowed from your pen was nothing but mockery and
derision of God. Yet you mock yourself and your leaders without
perceiving it, just as servants before you mocked God's emissaries and
His trusted ones. You have followed in their wake with a deed that has
shaken the foundations of assurance and the pillars of certitude,
convulsed stillness itself, and caused the supports of a mighty throne
to tremble.

\subsection{392}\label{section-887}

\begin{otherlanguage}{arabic}

اسمع قولى و لاتجاوز عن حدّك فاعرف نفسك يا ايّها المشرك باللّه و آياته و
المجادل بنفسه و المحارب بكينونته ثمّ اكف بما ارتكبت تاللّه بفعلك قد هبّت
عواصف الانتقام على العالمين فسوف تأخذك و اوليائك و انّه ما من اله الّا هو
له القدرة و الاقتدار و له العظمة و الأختيار ينزل ما يشآء بأمر من عنده و
انّه لهو المقتدر القهّار

\end{otherlanguage}

Ismaʿ qawlī wa-lā tatajāwaz ʿan ḥaddika fa-iʿrif nafsaka yā ayyuhā
al-mushriku bi-Allāh wa-āyātihi wa-al-mujādilu bi-nafsihi wa-al-muḥāribu
bi-kaynūnatihi thumma ikfā bi-mā irtakabta ta-Allāh bi-fiʿlika qad
habbat ʿawāṣifu al-intiqām ʿalā al-ʿālamīn fa-sawfa taʾkhudhuka
wa-awliyāʾaka wa-innahu mā min ilāhin illā huwa lahu al-qudratu
wa-al-iqtidār wa-lahu al-ʿaẓamatu wa-al-ikhtiyār yunazzilu mā yashāʾu
bi-amrin min ʿindihi wa-innahu la-huwa al-Muqtadiru al-Qahhār.

Listen to My words and do not overstep your bounds. Know yourself, you
who associate partners with God and His verses, dispute with His own
Self, and wage war against His Being. Enough of what you have committed!
By God, your deed has unleashed the tempests of vengeance upon the
worlds. Soon they will seize you and your allies. There is no God but
Him. His are power and might; His are majesty and the freedom to choose.
He sends down whatever He wills by a command from Himself. He is indeed
the All-Powerful, the All-Subduing.

\subsection{393}\label{section-888}

\begin{otherlanguage}{persian}

كلّ اشیآء متعجّب است كه چه تو را بآن داشته كه بالقاى آن نفس مردوده باين
شأن و جسارت با جمال احديّه بمحاربه برخواستى جمال امنع اقدس اعظم اقدمى را
كه نقطه بيان روح ما سواه فداه ميفرمايد اشاره من باو راه ندارد و هيچ
اشاره از مشيرين لعلوّ قدره و عظمة جلاله و ارتفاع اجلاله باو راجع نميشود
تو استهزاءً عليه نوشته آنچه را كه تا حال هيچ نفسى بآن ارتكاب ننموده و
ننوشته و عجب است كه خجلت نميبرى و ذكر نقطه بيان مينمائى آخر بچه دليل و
برهان اينگونه جسارت را جايز دانسته

\end{otherlanguage}

Kull-i ashyāʾ mutaʿajjib ast kih chih tu rā bi-ān dāshtih kih bi-ilqā-yi
ān nafs-i mardūdih bi-īn shaʾn va jisārat bā Jamāl-i Aḥadiyyih
bi-muḥāribih bar-khāstī Jamāl-i amnaʿ-i aqdas-i aʿẓam-i aqdamī rā kih
Nuqṭiy-i Bayān rūḥ-i mā sivāh fidāh mī-farmāyad ishāriy-i man bi-ū rāh
nadārad va hīch ishārih az mushīrīn li-ʿuluwwi qadrihi wa-ʿaẓamati
jalālihi wa-irtifāʿi ijlālihi bi-ū rājiʿ nimī-shavad tu istihzāʾan
ʿalayhi nivishtih ānchih rā kih tā ḥāl hīch nafsī bi-ān irtikāb
nanimūdih va nanivishtih va ʿajab ast kih khijlat nimī-barī va dhikr-i
Nuqṭiy-i Bayān mī-namāʾī ākhar bi-chih dalīl va burhān īn-gūnih jisārat
rā jāʾiz dānistih.

All things marvel at what could have prompted you, at the instigation of
that rejected soul, to rise up with such audacity and insolence to wage
war against the Beauty of Divine Unity. Of that most unapproachable,
most holy, most great and most ancient Beauty, the Primal Point of the
Bayán---may the souls of all besides Him be sacrificed for Him---says:
``My allusion cannot reach Him, nor can any allusion made by anyone
refer to Him, so exalted is His station, so great His glory, and so
lofty His majesty.'' Yet in mockery of Him you have written what no soul
has ever dared to commit or write. It is astonishing that you feel no
shame and still make mention of the Point of the Bayán. On what evidence
and proof, after all, have you deemed such insolence permissible?

\subsection{394}\label{section-889}

\begin{otherlanguage}{persian}

تحيّر اين عبد سهل است جميع اهل ملأ اعلى متحيّرند از امثال شما چه زود كافر
شديد بنفسى كه باو ادّعاى ايمان مينموديد گرفتم اين امر را ادراك ننمودى آخر
با اين همه ظهورات عزّ صمدانى و آيات منزله ربّانى اقلّا بايد توقّف كنى اگر چه
اگر ببيان ناظر باشى بهيچوجه سندى در دست نداشته و نداريد كه بر چنين نفسى
اعتراض نمائيد و لو هر چه بفرمايد

\end{otherlanguage}

Taḥayyur-i īn ʿabd sahl ast jamīʿ-i ahl-i Malaʾ-i Aʿlā mutaḥayyirand az
amthāl-i shumā chih zūd kāfir shudīd bi-nafsī kih bi-ū iddiʿā-yi īmān
mī-namūdīd giriftam īn amr rā idrāk nanimūdī ākhar bā īn hamih ẓuhūrāt-i
ʿizz-i ṣamadānī va āyāt-i munazzalih-yi rabbānī aqallan bāyad tavaqquf
kunī agar-chih agar bi-Bayān nāẓir bāshī bi-hīch-vajh sanadī dar dast
nadāshtih va nadārīd kih bar chunīn nafsī iʿtirāḍ namāʾīd va law har
chih bi-farmāyad.

This servant's bewilderment is of little consequence: all the
inhabitants of the Concourse on high are bewildered by the likes of you.
How quickly you have disbelieved in the very One in whom you professed
faith! Granted that you have not understood this Cause, surely, in the
face of all these manifestations of divine glory and divinely revealed
verses, you should at least withhold judgment. Indeed, if you look to
the Bayán, you have never had, nor do you now have, any warrant
whatsoever to object to such a Being, whatever He may say.

\subsection{395}\label{section-890}

\begin{otherlanguage}{persian}

افّ للّذينهم نفخوا فيك و ملئوا صدرك من غلّ اللّه و بغضه فآه آه فآه الّذى لا
يذكر بالقلم فاه الّذى لن يتمّ بالبيان على ما فعلت استرضاء لمن يفرّ عنه
السجّين بكينونته و الجحيم بحقيقته و النّار باشتعالها و لهيبها

\end{otherlanguage}

Uffin li-alladhīna hum nafakhū fīka wa-malaʾū ṣadraka min ghilli Allāh
wa-bughḍihi fa-āh āh fa-āh alladhī lā yudhkaru bi-al-qalam fa-āh alladhī
lan yatimma bi-al-bayān ʿalā mā faʿalta istirḍāʾan li-man yafirru ʿanhu
al-sijjīnu bi-kaynūnatihi wa-al-jaḥīmu bi-ḥaqīqatihi wa-al-nāru
bi-ishtiʿālihā wa-lahībihā.

Shame on those who have breathed into you and filled your breast with
malice and hatred toward God! Alas, alas! A lament that no pen can
record, a lament that no utterance can fully express, for what you have
done to please one from whom Sijjín in its very being, hell in its
reality, and fire in its blazing and flaming flee.

\subsection{396}\label{section-891}

\begin{otherlanguage}{persian}

در جميع بيان من فى البيان را از تعرّض و احتجاج با يكديگر منع فرموده و
ميفرمايند اگر هم بخواهيد با يكديگر در دلايل و براهين احتجاج نمائيد بكمال
ادب و حيا دلايل و براهين خود را ذكر نمائيد و مرقوم داريد كه مباد از آن
احتجاج بر آن جوهر عزّ و شمس معانى در ایّام ظهور و اشراق حزنى وارد شود
چنانچه ميفرمايد قوله جلّ كبريائه و ان حين ما تحبّون ان تتحاجّون بدلائل و
البرهان على اكمل الحيا تكتبون دلائلكم ثمّ على منتهى الأدب لتقولون الى ان
قال عزّ ذكره لعلّكم لاتلاقون اللّه ربّكم و تكسبون عملاً يحزن به اللّه ربّكم بما
يحزن مظهر نفسه و انتم لاتلتفتون و لا تتذكّرون انتهى

\end{otherlanguage}

Dar jamīʿ-i Bayān man fī al-Bayān rā az taʿarruḍ va iḥtijāj bā yakdīgar
manʿ farmūdih va mī-farmāyand agar ham bi-khvāhīd bā yakdīgar dar
dalāʾil va barāhīn iḥtijāj namāʾīd bi-kamāl-i adab va ḥayā dalāʾil va
barāhīn-i khud rā dhikr namāʾīd va marqūm dārīd kih mabād az ān iḥtijāj
bar ān Jawhar-i ʿIzz va Shams-i Maʿānī dar ayyām-i ẓuhūr va ishrāq ḥuznī
vārid shavad chunānchih mī-farmāyad qawluhu jalla kibriyāʾuhu wa-in ḥīna
mā tuḥibbūna an tataḥājjūna bi-dalāʾili wa-al-burhān ʿalā akmali al-ḥayā
taktubūna dalāʾilakum thumma ʿalā muntahā al-adab la-taqūlūna ilā an
qāla ʿazza dhikruhu laʿallakum lā tulāqūna Allāha rabbakum wa-taksibūna
ʿamalan yaḥzanu bihi Allāhu rabbukum bi-mā yaḥzanu maẓhara nafsihi
wa-antum lā taltafitūna wa-lā tatadhakkarūna intahā.

Throughout the Bayán, He has forbidden the people of the Bayán to assail
or dispute with one another. He says that if you wish to debate with one
another using evidence and proofs, you should state and write down your
evidence and proofs with the utmost courtesy and modesty, lest that
debate bring sorrow upon that Essence of Glory and Sun of Meanings in
the days of His manifestation and radiance. Thus He says---exalted be
His majesty: ``And when you wish to debate using evidence and proof,
write down your evidence with the utmost modesty, and then speak with
the greatest courtesy.'' Further on He says---exalted be His
remembrance: ``Lest you encounter God, your Lord, and commit a deed that
grieves God, your Lord, by grieving the Manifestation of His own Self,
while you neither take heed nor remember.'' Here ends the quotation.

\subsection{397}\label{section-892}

\begin{otherlanguage}{persian}

كاش همين قدر وصيّت آن محبوب امكان را عامل ميشدى نه و اللّه حرفى از بيان را
عامل نشده‌ايد مگر آنچه را كه بهواى خود مطابق ديده‌ايد فنعوذ باللّه من هذا
الذّنب لانّ كلماته قد كانت مقدّسة من انفسكم و اهوائكم و عرفانكم و ادراككم
كاش باحتجاج هم با مظهر ربّ الأرباب‌ كفايت مينمودى نه حيا نمودى و نه مراعات
ادب منظور داشتى ارتكاب نموده امرى را كه از اين يوم الى يوم الّذى لا آخر
له هر ذى ادب و حيائى نظر نمايد بآنچه نوشته خجل و شرمسار گردد تاللّه
ارتكبت ما احترقت به اكباد المقرّبين و

\end{otherlanguage}

Kāsh hamīn qadr vaṣiyyat-i ān Maḥbūb-i imkān rā ʿāmil mī-shudī nah
va-Allāh ḥarfī az Bayān rā ʿāmil nashudih-īd magar ānchih rā kih
bi-havā-yi khud muṭābiq dīdih-īd fa-naʿūdhu bi-Allāh min hādhā al-dhanb
li-anna kalimātihi qad kānat muqaddasatan min anfusikum wa-ahwāʾikum
wa-ʿirfānikum wa-idrākikum kāsh bi-iḥtijāj ham bā Maẓhar-i Rabb al-Arbāb
kifāyat mī-namūdī nah ḥayā namūdī va nah murāʿāt-i adab manẓūr dāshtī
irtikāb nimūdih amrī rā kih az īn yawm ilā yawm alladhī lā ākhira lahu
har dhī adab va ḥayāʾī naẓar namāyad bi-ānchih nivishtih khajil va
sharmsār gardad ta-Allāh irtakabta mā iḥtaraqat bihi akbādu
al-muqarrabīn wa-

If only you had observed even this much of the counsel of that Beloved
of all creation! No, by God, you have not acted upon a single letter of
the Bayán except what you have found to accord with your own desires. We
seek refuge in God from this sin, for His words are sanctified above
your selves, your desires, your knowledge, and your understanding. If
only you had confined yourself to disputing with the Manifestation of
the Lord of Lords! You showed neither modesty nor regard for courtesy.
You have committed a deed such that, from this day until the day that
has no end, every person possessed of courtesy and modesty who looks at
what you have written will be ashamed and abashed. By God, you have
committed what has burned the hearts of those brought near, and

\subsection{398}\label{section-893}

\begin{otherlanguage}{persian}

كاش حقّ را بمثل يكى از نفوس اهل ارض ميدانستى چه كه با آن نفوس جرئت
اينگونه جسارت و تكلّم نداشته و ندارى فو‌اللّه لو تسمع نفس الحيا يتبرّء منك و
كينونة الأدب يستعيذ باللّه من فعلك ادب از سجيّه انسان است و باو از دونش
ممتاز و هر نفسيكه باو فايز نشده البتّه عدمش بر وجودش رجحان داشته و دارد و
تو بمقرّى كه اصل ادب مخصوص آن مقرّ خلق شده معمول داشته آنچه را كه هر ذى
حيائى از ذكرش شرم دارد قد خرق حجاب ادبكم و حيائكم يا معشر الفاسقين

\end{otherlanguage}

Kāsh Ḥaqq rā bi-mithl-i yakī az nufūs-i ahl-i arḍ mī-dānistī chih kih bā
ān nufūs jurʾat-i īn-gūnih jisārat va takallum nadāshtih va nadārī
fa-wa-Allāh law tasmaʿu nafsu al-ḥayā yatabarraʾu minka wa-kaynūnatu
al-adab yastaʿīdhu bi-Allāh min fiʿlika adab az sajiyyih-yi insān ast va
bi-ū az dūnash mumtāz va har nafsī kih bi-ū fāʾiz nashudih albattih
ʿadamash bar vujūdash rujḥān dāshtih va dārad va tu bi-maqarrī kih aṣl-i
adab makhṣūṣ-i ān maqarr khalq shudih maʿmūl dāshtih ānchih rā kih har
dhī ḥayāʾī az dhikrash sharm dārad qad khuriqa ḥijābu adabikum
wa-ḥayāʾikum yā maʿshara al-fāsiqīn.

If only you had regarded the True One as you would any person on earth,
for you have never dared, nor do you now dare, to show such insolence or
speak in this way to those people. By God, if you could hear, modesty
itself disowns you, and the very being of courtesy seeks refuge in God
from your deed. Courtesy is a human virtue, distinguishing a human being
from what is beneath him. For any soul who has not attained it,
nonexistence has assuredly been, and remains, preferable to existence.
Yet toward that Seat for whom the very essence of courtesy was
specifically created, you have acted in a way that everyone possessed of
modesty is ashamed even to mention. The veil of your courtesy and
modesty has been torn apart, you company of the ungodly!

\subsection{399}\label{section-894}

\begin{otherlanguage}{persian}

جوهر قدسيكه ميفرمايد يا قوم قد جئتكم عن مشرق الرّوح بنبأ اللّه المهيمن
القيّوم و ما انطق من تلقاء نفسى بل ينطق ربّكم الرّحمن بهذا اللّسان الصّادق
المتكلّم العزيز المحبوب يا قوم لاتقتلونى باسياف البغضآء و لاتحكموا علىّ ما
يضيع به حرمة اللّه المقتدر العزيز المحبوب فيا قوم انّه يناديكم بلسانى و
ينطق فى صدرى و اقامنى على امره فو‌اللّه لا اجد لى من حركة و لا سكون الّا
بأمره تفكّروا لكى تشعرون هل يجرى كوثر الحيوان من عيون الكدرة لا فو‌ربّ
العالمين ان انتم تعرفون

\end{otherlanguage}

Jawhar-i qudsī kih mī-farmāyad yā qawmi qad jiʾtukum ʿan mashriqi al-rūḥ
bi-nabaʾi Allāh al-Muhaymini al-Qayyūm wa-mā anṭiqu min tilqāʾi nafsī
bal yanṭiqu rabbukumu al-Raḥmān bi-hādhā al-lisāni al-ṣādiqi
al-mutakallimi al-ʿAzīzi al-Maḥbūb yā qawmi lā taqtulūnī bi-asyāfi
al-baghḍāʾ wa-lā taḥkumū ʿalayya mā yaḍīʿu bihi ḥurmatu Allāh
al-Muqtadiri al-ʿAzīzi al-Maḥbūb fa-yā qawmi innahu yunādīkum bi-lisānī
wa-yanṭiqu fī ṣadrī wa-aqāmanī ʿalā amrihi fa-wa-Allāh lā ajidu lī min
ḥarakatin wa-lā sukūnin illā bi-amrihi tafakkarū li-kay tashʿurūn hal
yajrī kawtharu al-ḥayawān min ʿuyūni al-kadirah lā fa-wa-Rabbi
al-ʿālamīn in antum taʿrifūn.

That holy Essence says: ``O people, I have come to you from the
dawning-place of the Spirit with the tidings of God, the All-Protecting,
the Self-Subsisting. I do not speak of My own accord; rather, your Lord,
the All-Merciful, speaks through this truthful, eloquent, mighty and
beloved tongue. O people, do not kill Me with the swords of hatred, nor
pass against Me a judgment that violates the sanctity of God, the
All-Powerful, the Almighty, the Beloved. O people, He calls to you
through My tongue, speaks within My breast, and has raised Me up for His
Cause. By God, I find in Myself neither movement nor stillness except at
His command. Reflect, so that you may perceive. Does the heavenly river
of life flow from turbid springs? No, by the Lord of the worlds, if only
you knew!

\subsection{400}\label{section-895}

\begin{otherlanguage}{arabic}

يا قوم هذا معين مشيّة اللّه قد جرت منه انهار ارادته بالحقّ كيف شآء و انّه
لهو الحاكم بالحقّ و انّه لهو الحقّ علّام الغيوب لو تجدنّ عندكم من ذائقة
فاشربوه لعلّ تجدون منه حلاوة القول ثمّ رايحة اللّه ربّكم و هذا قول الحقّ و ما
بعد الحقّ الّا الضّلال ان انتم تفقهون

\end{otherlanguage}

Yā qawmi hādhā maʿīnu mashīʾati Allāh qad jarat minhu anhāru irādatihi
bi-al-ḥaqqi kayfa shāʾa wa-innahu la-huwa al-ḥākimu bi-al-ḥaqq wa-innahu
la-huwa al-Ḥaqqu ʿAllāmu al-ghuyūb law tajidunna ʿindakum min dhāʾiqatin
fa-ishrabūhu laʿalla tajidūna minhu ḥalāwata al-qawl thumma rāʾiḥata
Allāh rabbikum wa-hādhā qawlu al-ḥaqq wa-mā baʿda al-ḥaqqi illā al-ḍalāl
in antum tafqahūn.

O people, this is the wellspring of God's Will, from which the rivers of
His purpose have flowed in truth, as He has willed. He is indeed the One
who rules in truth; He is indeed the Truth, the Knower of the unseen. If
you find within yourselves any sense of taste, drink from it, that you
may discover in it the sweetness of utterance and the fragrance of God,
your Lord. This is the word of truth, and beyond truth there is nothing
but error, if only you understood.

\subsection{401}\label{section-896}

\begin{otherlanguage}{arabic}

هل تعترضون على الّذى ارسلنى و ينطق فى صدرى اذاً وا حسرتا عليكم و على ما
انتم به تعملون انّ الّذى جائكم بآيات اللّه و يقول انّها قد نزلت من عنده هل
ينبغى لأحد ان يعترض عليه اذا فانصفوا ان انتم منصفون يا قوم انّى لم اقدر
ان اتنفّس بنفس وحده الّا بعد اذنه و انتم تجدون نفس الرّحمن من هذه النّفس
الّتى ارسلها اللّه بالحقّ ان انتم بقلوبكم فاقبلون

\end{otherlanguage}

Hal taʿtariḍūna ʿalā alladhī arsalanī wa-yanṭiqu fī ṣadrī idhan wā
ḥasratā ʿalaykum wa-ʿalā mā antum bihi taʿmalūn inna alladhī jāʾakum
bi-āyāti Allāh wa-yaqūlu innahā qad nazalat min ʿindihi hal yanbaghī
li-aḥadin an yaʿtariḍa ʿalayhi idhan fa-anṣifū in antum munṣifūn yā
qawmi innī lam aqdir an atanaffasa bi-nafasin waḥdahu illā baʿda idhnihi
wa-antum tajidūna nafasa al-Raḥmān min hādhihi al-nafsi allatī arsalahā
Allāhu bi-al-ḥaqq in antum bi-qulūbikum fa-aqbilūn.

Do you object to the One who sent Me and speaks within My breast? Alas
for you, then, and for what you do! Is it fitting for anyone to object
to the One who has come to you with God's verses and says that they have
been sent down from Him? Be fair, then, if you are fair-minded. O
people, I cannot draw a single breath except with His permission. You
will find the breath of the All-Merciful coming from this Soul whom God
has sent in truth, if you turn toward Me with your hearts.

\subsection{402}\label{section-897}

\begin{otherlanguage}{arabic}

لاتعروا انفسكم عن ردآء الأدب و الانصاف و لاتكلّموا بما نهيتم عنه فى الواح
ربّكم المهيمن القيّوم و الّذى جعل محروماً عن الأدب انّه لعرىّ و لو يلبس حرر
الأرض كلّها و هذا قد نزل بالحقّ فى الواح عزّ محفوظ و من لا ادب له لا ايمان
له و بذلك يشهد ما نزل فى البيان ان انتم تشهدون قد خلق الأدب فى البيان
لهذا الجمال و كذلك كلّ اسم خير ان انتم تعرفون

\end{otherlanguage}

Lā tuʿarrū anfusakum ʿan ridāʾi al-adabi wa-al-inṣāf wa-lā takallamū
bi-mā nuhītum ʿanhu fī alwāḥi rabbikumu al-Muhaymini al-Qayyūm
wa-alladhī juʿila maḥrūman ʿani al-adab innahu la-ʿariyyun wa-law
yalbasu ḥurara al-arḍi kullahā wa-hādhā qad nazala bi-al-ḥaqqi fī alwāḥi
ʿizzin maḥfūẓ wa-man lā adaba lahu lā īmāna lahu wa-bi-dhālika yashhadu
mā nazala fī al-Bayān in antum tashhadūn qad khuliqa al-adabu fī
al-Bayān li-hādhā al-Jamāl wa-ka-dhālika kullu ismi khayrin in antum
taʿrifūn.

Do not strip yourselves of the robe of courtesy and fairness, and do not
utter what you have been forbidden in the tablets of your Lord, the
All-Protecting, the Self-Subsisting. Whoever is deprived of courtesy is
indeed naked, even if he wears all the silks of the earth. This has been
sent down in truth in preserved tablets of glory. Whoever has no
courtesy has no faith. To this, what has been revealed in the Bayán
bears witness, if you would bear witness. Courtesy was created in the
Bayán for this Beauty, as was every good name, if only you knew.

\subsection{403}\label{section-898}

\begin{otherlanguage}{arabic}

فو‌عمر من ظهر بالحقّ انّه ما قصد فى البيان الّا نفسى و ما تنفّس الّا بذكرى و
ما قدّر كلّ خير الّا لجمالى لو لا ذكرى ما نزل البيان و ما تكلّم جمال الرّحمن
بكلمة بينكم لانّ لمثلكم لاينبغى الّا كلماتكم ان انتم تفقهون

\end{otherlanguage}

Fa-wa-ʿumri man ẓahara bi-al-ḥaqq innahu mā qaṣada fī al-Bayān illā
nafsī wa-mā tanaffasa illā bi-dhikrī wa-mā qaddara kulla khayrin illā
li-jamālī law lā dhikrī mā nazala al-Bayān wa-mā takallama Jamālu
al-Raḥmān bi-kalimatin baynakum li-anna li-mithlikum lā yanbaghī illā
kalimātukum in antum tafqahūn.

By the life of Him who has appeared in truth! In the Bayán He intended
none but My own Self; He breathed only in remembrance of Me and ordained
every good thing solely for My Beauty. Were it not for My remembrance,
the Bayán would not have been revealed, nor would the Beauty of the
All-Merciful have uttered a word among you, for nothing befits the likes
of you but your own words, if only you understood.

\subsection{404}\label{section-899}

\begin{otherlanguage}{arabic}

قد نزل البيان لنفسى و انّى حينئذ اقرئه فى كلّ الأحيان و ما نزل حينئذ انّه
لهو البيان لو انتم تعرفون ما نزل حينئذ ليكون شاهداً على ما نزل من قبل و
ما نزل من عنده ليكون دليلاً على ما ينزل حينئذ من قلم اللّه المقتدر العزيز
المحبوب

\end{otherlanguage}

Qad nazala al-Bayānu li-nafsī wa-innī ḥīnaʾidhin aqraʾuhu fī kulli
al-aḥyān wa-mā nazala ḥīnaʾidhin innahu la-huwa al-Bayān law antum
taʿrifūn mā nazala ḥīnaʾidhin li-yakūna shāhidan ʿalā mā nazala min qabl
wa-mā nazala min ʿindihi li-yakūna dalīlan ʿalā mā yanzilu ḥīnaʾidhin
min qalami Allāh al-Muqtadiri al-ʿAzīzi al-Maḥbūb.

The Bayán was revealed for My own Self, and I now recite it at all
times. What is being revealed now is indeed the Bayán, if only you knew.
What is being revealed now is to bear witness to what was revealed
before, and what was revealed from Him is to serve as proof of what now
descends from the Pen of God, the All-Powerful, the Almighty, the
Beloved.

\subsection{405}\label{section-900}

\begin{otherlanguage}{arabic}

يا قوم انّى نفسه و هو نفسى و انّه قد فدى نفسه لنفسى و انّى قد اخذت نفسى و
روحى على يدى لانفقهما فى سبيله ثمّ فى حبّه و رضائه و يشهد بذلك كلّ الذّرّات
لو انتم تسمعون خافوا عن اللّه و لاتفرقوا بينى و بينه و لو جعل اللّه هذا
فوق قدرتكم لو اجتمعن كلّ من فى السّموات و الأرض ان يفرقن بيننا لن يقدرن و
لن يستطيعنّ

\end{otherlanguage}

Yā qawmi innī nafsuhu wa-huwa nafsī wa-innahu qad fadā nafsahu li-nafsī
wa-innī qad akhadhtu nafsī wa-rūḥī ʿalā yadayya li-unfiqahumā fī
sabīlihi thumma fī ḥubbihi wa-riḍāʾihi wa-yashhadu bi-dhālika kullu
al-dharrāt law antum tasmaʿūn khāfū ʿani Allāh wa-lā tufarriqū baynī
wa-baynihi wa-law jaʿala Allāhu hādhā fawqa qudratikum law ijtamaʿanna
kullu man fī al-samāwāti wa-al-arḍ an yufarriqna baynanā lan yaqdirna
wa-lan yastaṭīʿunn.

O people, I am His own Self, and He is My own Self. He sacrificed
Himself for Me, and I have taken My life and My soul in My hands to
expend them in His path, for His love and His good pleasure. Every atom
bears witness to this, if only you could hear. Fear God and do not
separate Me from Him, although God has placed this beyond your power.
Were all who are in the heavens and on earth to unite to separate Us,
they would have neither the power nor the ability.

\subsection{406}\label{section-901}

\begin{otherlanguage}{arabic}

و انه حينئذ كان جالساً عن يمين عرشه و كان انامله بين انيابه و ينظر متفكّرا
فيكم بما ظهر منكم يا ملأ المشركين و يضجّ و يبكى بما ورد علىّ و ما يرد منكم
من بعد و حينئذ يقصّ علىّ ما هو المكنون فى صدوركم و يبكى و ببكائه ينوحنّ اهل
ملأ الأعلی و لكن انتم لا تعرفون و لا تشعرون

\end{otherlanguage}

Wa-innahu ḥīnaʾidhin kāna jālisan ʿan yamīni ʿarshihi wa-kāna anāmiluhu
bayna anyābihi wa-yanẓuru mutafakkiran fīkum bi-mā ẓahara minkum yā
malaʾa al-mushrikīn wa-yaḍijju wa-yabkī bi-mā warada ʿalayya wa-mā
yaridu minkum min baʿd wa-ḥīnaʾidhin yaquṣṣu ʿalayya mā huwa al-maknūnu
fī ṣudūrikum wa-yabkī wa-bi-bukāʾihi yanūḥanna ahlu Malaʾi al-Aʿlā
wa-lākin antum lā taʿrifūna wa-lā tashʿurūn.

He is now seated at the right of His throne, with His fingertips between
His teeth, looking upon you in reflection on what has proceeded from
you, you company of those who associate partners with God. He cries out
and weeps over what has befallen Me and what will yet come from you. He
now recounts to Me what is hidden in your breasts and weeps; at His
weeping, the inhabitants of the Concourse on high lament. Yet you
neither know nor perceive.''

\subsection{407}\label{section-902}

\begin{otherlanguage}{persian}

مع‌ذلك وارد آورده‌‌ايد آنچه در ارض شبه نداشته اگر فى‌الجمله شاعر باشيد ادراك
مينمائيد كه آنچه اعتراض نموده‌ايد بر حقّ وارد چه كه منزل حقّ بوده و خواهد
بود و دخلى باحدى نداشته و ندارد اگر چه مشعر نخواهيد شد چنانچه ملل قبل
آنچه اعتراض نمودند من غير شعور بوده و اگر مستشعر بودند و ادراك مينمودند
كه آن اعتراضات بر حقّ وارد است ابدا باعتراض مشغول نميشدند و بهمين جهت
لازال مظاهر نفى بر مظاهر اثبات وارد آورده‌اند آنچه قلم و لسان از ذكر و
بيانش حيا ميكند

\end{otherlanguage}

Maʿa dhālik vāred āvardeh-īd ānchih dar arḍ shabah nadāshtih. Agar
fi'l-jumlih shāʿir bāshīd idrāk minamāyīd kih ānchih iʿtirāḍ namūdeh-īd
bar Ḥaqq vāred, chih kih munzil Ḥaqq būdih va khvāhad būd va dakhlī
bi-aḥadī nadāshtih va nadārad. Agarchih mushʿir nakhvāhīd shud,
chunānchih milal-i qabl ānchih iʿtirāḍ namūdand min ghayr-i shuʿūr būdih
va agar mustashʿir būdand va idrāk minamūdand kih ān iʿtirāḍāt bar Ḥaqq
vāred ast abadan bi-iʿtirāḍ mashghūl namīshudand. Va bi-hamīn jihat
lāzāl maẓāhir-i nafy bar maẓāhir-i ithbāt vāred āvardeh-and ānchih qalam
va lisān az dhikr va bayānash ḥayā mikunad.

Nevertheless, you have inflicted what has had no parallel on earth. If
you had even a little awareness, you would realize that the objections
you have raised are directed against God, for it is God who has sent
down revelation and will continue to do so; no one else has had or has
any part in it. Yet you will not perceive this, just as the objections
raised by former religious communities were made without understanding.
Had they been aware and realized that those objections were directed
against God, they would never have occupied themselves with objecting.
For this very reason, the manifestations of negation have always
inflicted upon the Manifestations of affirmation what pen and tongue are
ashamed to recount and describe.

\subsection{408}\label{section-903}

\begin{otherlanguage}{persian}

و زود است كه آنچه اخبار فرموده واقع شود چه كه باعراض و اعتراض تنها كفايت
ننموده و نخواهيد نمود جمال قدم ميفرمايد اى اصمّ وجود يك نغمه از نغمات
ورقاء ظهور قبلم ذكر مينمايم كه لعلّ فى الجمله ادراك نمائى و اگر تو ادراك
ننمائى حقايق كلّ شیء اصغا نمايند و حجاب حيا را خرق ننمايند و طراز ادب را
از هيكل خود نيندازند و آن اين است كه بحرف حىّ و وحيد اكبر عليهما بهائى
ميفرمايد در حينيكه از اين ظهور امنع اقدس سؤال نموده‌اند قوله تكبّر كبريائه
فو‌الّذى تفرّد بالعزّ و الجلال و ذرء الأشیآء لا من شیء بالمثال يستحيى اليقين
ان يوقن فيه او يسجد له و يستحيى الدّليل ان يدلّ عليه لغيره ترى لاعرفنك او
على دونه تشهد لاستدلنّ به عليه

\end{otherlanguage}

Va zūd ast kih ānchih ikhbār farmūdih vāqiʿ shavad, chih kih bi-iʿrāḍ va
iʿtirāḍ tanhā kifāyat nanamūdih va nakhvāhīd namūd. Jamāl-i Qidam
mifarmāyad: Ay aṣamm-i vujūd, yik naghmih az naghamāt-i Varqā-yi Ẓuhūr-i
Qablam dhikr minamāyam kih laʿalla fi'l-jumlih idrāk namāʾī, va agar tū
idrāk nanamāʾī ḥaqāʾiq-i kull-i shayʾ iṣghā namāyand va ḥijāb-i ḥayā rā
kharq nanamāyand va ṭirāz-i adab rā az haykal-i khud nayandāzand. Va ān
īn ast kih bi-Ḥarf-i Ḥayy va Vaḥīd-i Akbar, ʿalayhimā bahāʾī, mifarmāyad
dar ḥīnī kih az īn Ẓuhūr-i Amnaʿ-i Aqdas suʾāl namūdeh-and. Qawluhu,
takabbara kibriyāʾuhu: Fa-wa'lladhī tafarrada bi'l-ʿizzi wa'l-jalāli
wa-dharaʾa al-ashyāʾa lā min shayʾin bi'l-mithāl, yastaḥyī al-yaqīnu an
yūqina fīhi aw yasjuda lahu, wa-yastaḥyī al-dalīlu an yadulla ʿalayhi
li-ghayrihi. Turā li-aʿrifannaka aw ʿalā dūnihi tashhadu li-astadillanna
bihi ʿalayhi?

Soon what He has foretold will come to pass, for you have not been
content, nor will you be, with rejection and objection alone. The
Ancient Beauty says: O deaf one of existence! I will recount one melody
from the melodies of the Dove of My previous Revelation, that perhaps
you may understand a little. Even if you do not understand, the
realities of all things may listen, refrain from tearing the veil of
modesty, and not cast the adornment of courtesy from their forms. It is
this, addressed to the Letter of the Living and the Greater Vaḥīd---My
glory rest upon them both---when they asked about this most
unassailable, most holy Revelation. His words---exalted be His
majesty---are: ``By Him who is alone in might and majesty and brought
things into being from nothing through the exemplar! Certitude is
ashamed to attain certainty concerning Him or to bow before Him, and
proof is ashamed to point Him out to another. Are You seen, that I might
know You, or do You bear witness to anything other than Him, that I
might infer Him from it?''

\subsection{409}\label{section-904}

\begin{otherlanguage}{persian}

ساذج قدمى را كه ميفرمايد يقين حيا مينمايد از ايقان در او و سجده از براى
او و دليل حيا مينمايد از آنكه دليل شود لنفس او و تو بی‌شرم بی‌حيا بتضييع
امر اللّه چنان قيام نموده كه شبه آن در اكوان ممكن نه قد شققت بأنامل الشّرك
حجبات الحياء و الحرمة نسئل اللّه بأن يشقّ حجابك و يظهر نفسك الخبيث بين
العالمين و ينزع عنك ردآء السّتر و يعريك من اثواب الحرمة و يكشف لعباده ما
كنت عليه و انّه لمن دعاه مجيب و اللّه يا ايّها الغافل بفعلك اضطرب كلّ شیء و
شقّت ستر حجاب العصمة و تزلزلت اركان ملأ العالين

\end{otherlanguage}

Sādhij-i qidamī rā kih mifarmāyad yaqīn ḥayā minamāyad az īqān dar ū va
sajdih az barā-yi ū va dalīl ḥayā minamāyad az ānkih dalīl shavad
li-nafs-i ū, va tū bī-sharm-i bī-ḥayā bi-taḍyīʿ-i Amru'llāh chunān qiyām
namūdih kih shabah-i ān dar akvān mumkin nah. Qad shaqaqta bi-anāmili
al-shirki ḥujubāti al-ḥayāʾi wa'l-ḥurmati. Nasʾalu Allāha bi-an yashuqqa
ḥijābaka wa-yuẓhira nafsaka al-khabītha bayna al-ʿālamīna wa-yanziʿa
ʿanka ridāʾa al-sitri wa-yuʿarriyaka min athwābi al-ḥurmati wa-yakshifa
li-ʿibādihi mā kunta ʿalayhi, wa-innahu li-man daʿāhu mujībun. Wa'llāhi
yā ayyuhā al-ghāfilu, bi-fiʿlika iḍṭaraba kullu shayʾin wa-shuqqat sitru
ḥijābi al-ʿiṣmati wa-tazalzalat arkānu malaʾi al-ʿālīn.

Of that Essence of eternity He says that certitude is ashamed to attain
certainty concerning Him or to bow before Him, and proof is ashamed to
serve as a proof of His very Self. Yet you, shameless and immodest, have
arisen to undermine the Cause of God in a manner whose like is
impossible in all the worlds of being. You have torn the veils of
modesty and reverence with the fingers of association with God. We ask
God to tear away your veil, expose your wicked self among the peoples of
the world, strip you of the mantle of concealment, divest you of the
garments of respectability, and disclose to His servants what you have
been doing. He indeed answers those who call upon Him. By God, heedless
one! Through your deed all things have been thrown into turmoil, the
covering of the veil of purity has been torn, and the pillars of the
exalted Concourse have been shaken.

\subsection{410}\label{section-905}

\begin{otherlanguage}{persian}

قسم بمحبوب امكان كه من‌بعد حرمت هيچ مظهرى از مظاهر الهيّه را منظور نداشته
و نداريد مگر نفوس مشركه را آن هم لأجل وساوس و ظنونى كه در آن انفس مستور
است يك بيان ديگر از بيان منزل بيان خالصاً لوجه الرّحمن القا ميشود كه شايد
نسيمى از رضوان كلمات الهى بر اهل اكوان مرور نمايد و شايد از ميامن آن
انفس قدسيّه از خلف ستر و حجاب بعرصه شهود آيند و بمراد اللّه فايز گردند و
بآنچه از قبل ذكر شده عامل شوند و اين ایّام را غنيمت شمرند و هم‌‌چنين امثال
آن نفوس آنقدر ادراك نمايند كه بچه شأن حرمت اين ظهور عند اللّه منظور بوده
لعلّ از جسارت ممنوع شوند

\end{otherlanguage}

Qasam bi-Maḥbūb-i imkān kih min-baʿd ḥurmat-i hīch maẓharī az maẓāhir-i
ilāhiyyih rā manẓūr nadāshtih va nadārīd magar nufūs-i mushrikih rā, ān
ham li-ajl-i vasāvis va ẓunūnī kih dar ān anfus mastūr ast. Yik bayān-i
dīgar az bayān-i Munzil-i Bayān khāliṣan li-wajhi al-Raḥmān ilqā
mīshavad kih shāyad nasīmī az Riḍvān-i kalimāt-i ilāhī bar ahl-i akvān
murūr namāyad va shāyad az mayāmin-i ān anfus-i qudsiyyih az khalf-i
sitr va ḥijāb bi-ʿarṣih-yi shuhūd āyand va bi-murād-i Allāh fāʾiz
gardand va bi-ānchih az qabl dhikr shudih ʿāmil shavand va īn ayyām rā
ghanīmat shumarand. Va hamchunīn amthāl-i ān nufūs ānqadr idrāk namāyand
kih bi-chih shaʾn ḥurmat-i īn Ẓuhūr ʿinda'llāh manẓūr būdih, laʿalla az
jisārat mamnūʿ shavand.

I swear by the Beloved of all contingent being that henceforth you have
shown and show regard for the sanctity of none of the Manifestations of
God, but only for souls who associate partners with Him---and even that
is because of the insinuations and conjectures concealed within those
souls. Another utterance of the Revealer of the Bayān is now presented
solely for the sake of the All-Merciful, so that perhaps a breeze from
the paradise of divine words may pass over the inhabitants of the worlds
of being. Perhaps, through its blessings, holy souls may emerge from
behind their veils and coverings into the realm of visibility, attain
what God desires, act upon what has previously been stated, and cherish
these days as a precious opportunity. Likewise, perhaps those who
resemble such souls may understand, at least to some degree, how highly
the sanctity of this Revelation has been regarded in the sight of God,
and may thus be restrained from insolence.

\subsection{411}\label{section-906}

\begin{otherlanguage}{arabic}

قوله عزّ اعزازه فقد سمعت كتابك و ان ما فيه جوهر لو لا ما فيه ما اجبتك على
ذلك القرطاس و لا حينئذ يا على ما قدّر فى الابداع فما اعظم ذكر من قد سئلت
عنه و انّ ذلك اعلى و اعزّ و اجلّ و امنع و اقدس من ان يقدر الأفئدة بعرفانها
و الأرواح بالسّجود له و الأنفس بثنائه و الاجساد بذكر بهائه فما عظمت
مسئلتك و صغرت كينونيّتك هل شمس الّتى هى فى مراياء ظهوره فى نقطة البيان
يسئل عن شمس الّتى تلك الشّموس فى يوم ظهوره سجّاد لطلعتها ان كانت شموساً
حقيقيّة و الّا لا ينبغى لعلوّ قدسها و سموّ ذكرها و لو لا كنت من واحد الأوّل
لجعلت لك من الحدّ حيث قد سئلت عن اللّه الّذى قد خلقك و رزقك و اماتك و ابعثك
فى هيكلك هذا بالنّقطة البيان فى ذلك الظّهور المنفرد بالكيان فقل اوّلاً اذا
اردت ان تخطر بعلمك ذكره سبحان الّذى ذوالملك و الملكوت تسعة عشر مرّة الى ان
يقول عزّ ذكره و قد كتبت جوهرة فى ذكره و هو انّه لا يستشار باشارتى و لا بما
ذكر فى البيان بلى و عزّته تلك الكلمة اكبر عند اللّه عن عبادة ما على الأرض
اذ جوهر كلّ العبادة ينتهى الى ذلك فعلى ما قد عرفت اللّه فاعرف مظهر نفسه
فانّه اجلّ و اعلى من ان يكون معروفاً بدونه او مستشيراً باشارة خلقه و انّنى
انا اوّل عبد قد آمنت به و بآياته و اخذت من ابكار حدائق جنّة عرفانه حدائق
كلماته بلى و عزّته هو الحقّ لا اله الّا هو كلّ بأمره قائمون

\end{otherlanguage}

Qawluhu, ʿazza iʿzāzuhu: Fa-qad samiʿtu kitābaka wa-inna mā fīhi
jawharun, law lā mā fīhi mā ajabtuka ʿalā dhālika al-qirṭāsi wa-lā
ḥīnaʾidhin yā ʿAlī mā quddira fī al-ibdāʿi. Fa-mā aʿẓama dhikra man qad
saʾalta ʿanhu, wa-inna dhālika aʿlā wa-aʿazzu wa-ajallu wa-amnaʿu
wa-aqdasu min an yaqdiru al-afʾidatu bi-ʿirfānihā wa'l-arwāḥu
bi'l-sujūdi lahu wa'l-anfusu bi-thanāʾihi wa'l-ajsādu bi-dhikri
bahāʾihi. Fa-mā ʿaẓumat masʾalatuka wa-ṣaghurat kaynūniyyatuka. Hal
shamsu allatī hiya fī marāyā ẓuhūrihi fī Nuqṭati al-Bayāni yusʾalu ʿan
shamsi allatī tilka al-shumūsu fī yawmi ẓuhūrihi sujjādun li-ṭalʿatihā
in kānat shumūsan ḥaqīqiyyatan, wa-illā lā yanbaghī li-ʿuluwwi qudsihā
wa-sumuwwi dhikrihā? Wa-law lā kunta min Wāḥidi al-Awwali la-jaʿaltu
laka mina al-ḥaddi ḥaythu qad saʾalta ʿani Allāhi alladhī qad khalaqaka
wa-razaqaka wa-amātaka wa-abʿathaka fī haykalika hādhā bi'l-Nuqṭati
al-Bayāni fī dhālika al-ẓuhūri al-munfaridi bi'l-kiyān. Fa-qul awwalan
idhā aradta an tukhṭira bi-ʿilmika dhikrahu: Subḥāna alladhī Dhū
al-Mulki wa'l-Malakūt, tisʿata ʿashara marratan. Ilā an yaqūla, ʿazza
dhikruhu: Wa-qad katabtu jawharatan fī dhikrihi wa-huwa annahu lā
yustashāru bi-ishāratī wa-lā bi-mā dhukira fī al-Bayāni. Balā
wa-ʿizzatihi tilka al-kalimatu akbaru ʿinda Allāhi ʿan ʿibādati mā ʿalā
al-arḍi, idh jawharu kulli al-ʿibādati yantahī ilā dhālika. Fa-ʿalā mā
qad ʿarafta Allāha fa-ʿrif Maẓhara Nafsihi, fa-innahu ajallu wa-aʿlā min
an yakūna maʿrūfan bi-dūnihi aw mustashīran bi-ishārati khalqihi.
Wa-innanī anā awwalu ʿabdin qad āmantu bihi wa-bi-āyātihi wa-akhadhtu
min abkāri ḥadāʾiqi jannati ʿirfānihi ḥadāʾiqa kalimātihi. Balā
wa-ʿizzatihi huwa al-Ḥaqqu lā ilāha illā huwa, kullun bi-amrihi qāʾimūn.

His words---exalted be His honor---are: ``I have heard your letter, and
what it contains is a gem. Were it not for what it contains, I would not
have answered you on that sheet, nor then, O ʿAlī, on whatever has been
ordained in creation. How momentous is the mention of Him about whom you
have asked! This is higher, mightier, more glorious, more unassailable,
and holier than hearts can attain through recognition, spirits through
bowing before Him, souls through praising Him, or bodies through
recounting His splendor. How great is your question, and how small your
own being! Is a sun among the mirrors of His revelation in the Point of
the Bayān to be asked about that Sun before whose countenance those suns
will bow on the day of His revelation, if they are true suns? Otherwise,
they are unworthy of the loftiness of its sanctity and the sublimity of
its mention. Were you not of the First Unity, I would have imposed a
penalty upon you, since you have asked about God, who created you,
sustained you, caused you to die, and raised you up in this form of
yours through the Point of the Bayān, in that revelation unique in
existence. First, then, whenever you wish to bring His mention to mind,
say nineteen times: `Glory be to Him who possesses the kingdoms of earth
and heaven.'\,'' He continues---exalted be His mention---until He says:
``I have written a gem concerning Him, namely, that He is not to be
consulted through My indication, nor through what has been mentioned in
the Bayān. Indeed, by His might, that utterance is greater in the sight
of God than the worship of all that is on earth, for the essence of all
worship culminates in this. As you have known God, then, know the
Manifestation of His Self, for He is too glorious and exalted to be
known through anyone other than Himself or to seek counsel through the
indication of His creatures. I am indeed the first servant to believe in
Him and His verses, and from the pristine gardens of the paradise of His
recognition I have taken the gardens of His words. Indeed, by His might,
He is the Truth; there is no God but Him. All subsist through His
command.''

\subsection{412}\label{section-907}

\begin{otherlanguage}{persian}

حال از اين بيانات منزل بيان جلّت عظمته شأن مرايائى كه من دون اللّه از براى
خود حقّ فرض گرفته و شأن معرضين و متعرّضين را ملاحظه نمائيد ساذج عزّى را كه
محبوب عالميان كه بحرفى از او جميع من فى البيان خلق شده باين احترام ذكر
مينمايند كه ميفرمايد اگر تو از واحد اوّل نبودى حدّ بر تو جارى مينمودم حال
قدرى تفكّر نمائيد كه مقصود از اين تحديد كبرى چه بوده بچنين نفسى اگر
فى‌الجمله صاحب ادراك باشيد از همين كلمات صيحه زنيد و فرياد واحسرتا على ما
فرطنا فى جنب اللّه برآوريد و حال آنكه سائل خلاف ادب و احترام در عريضه خود
ذكر ننموده بمجرّد سؤال اين قسم تحديد فرموده‌اند كه او و ساير اهل بيان
متنبّه شوند و از غير نفس او از او سؤال ننمايند چه كه او بنفسه دليل خود
بوده و بدون خود معروف نه چنانچه در آخر اين توقيع ميفرمايد هر نفسيكه آن
جمال قدم را بغير عين او مشاهده نمايد ابداً باو عارف نگردد

\end{otherlanguage}

Ḥāl az īn bayānāt-i Munzil-i Bayān, jallat ʿaẓamatuhu, shaʾn-i marāyāʾī
kih min dūni'llāh az barā-yi khud ḥaqq farḍ giriftih va shaʾn-i muʿriḍīn
va mutaʿarriḍīn rā mulāḥiẓih namāyīd. Sādhij-i ʿizzī rā kih Maḥbūb-i
ʿālamiyān kih bi-ḥarfī az ū jamīʿ-i man fi'l-Bayān khalq shudih bi-īn
iḥtirām dhikr minamāyand kih mifarmāyad agar tū az Vāḥid-i Avval nabūdī
ḥadd bar tū jārī minamūdam. Ḥāl qadrī tafakkur namāyīd kih maqṣūd az īn
taḥdīd-i kubrā chih būdih bi-chunīn nafsī. Agar fi'l-jumlih ṣāḥib-i
idrāk bāshīd az hamīn kalimāt ṣayḥih zanīd va faryād-i vā ḥasratā ʿalā
mā farraṭnā fī janbi Allāh bar-āvarīd. Va ḥāl ānkih sāʾil khilāf-i adab
va iḥtirām dar ʿarīḍih-yi khud dhikr nanamūdih, bi-mujarrad-i suʾāl īn
qism taḥdīd farmūdeh-and kih ū va sāʾir-i ahl-i Bayān mutanabbih shavand
va az ghayr-i nafs-i ū az ū suʾāl nanamāyand, chih kih ū bi-nafsih
dalīl-i khud būdih va bi-dūn-i khud maʿrūf nah. Chunānchih dar ākhir-i
īn tawqīʿ mifarmāyad har nafsī kih ān Jamāl-i Qidam rā bi-ghayr-i ʿayn-i
ū mushāhidih namāyad abadan bi-ū ʿārif nagardad.

Now, in light of these utterances of the Revealer of the
Bayān---glorious is His grandeur---consider the standing of those
mirrors who, apart from God, have assumed a right for themselves, and
the standing of those who turn away and raise objections. Consider that
Essence of glory whom the Beloved of the worlds, through one letter of
whom all who are in the Bayān were created, mentions with such reverence
that He says: ``Were you not of the First Unity, I would have imposed a
penalty upon you.'' Reflect a little, then, on the purpose of this most
severe warning to such a person. If you possess even a little
understanding, these very words should cause you to cry out and raise
the lament: ``Alas for what we have neglected in our duty toward God!''
Yet the questioner had written nothing contrary to courtesy and respect
in his petition. Merely for asking, He warned him in this manner, so
that he and the rest of the people of the Bayān might take heed and not
ask anyone other than His own Self about Him. For He Himself has been
His own proof and is not known through anyone other than Himself. Thus,
at the end of this epistle, He says that whoever beholds that Ancient
Beauty with any eye other than His own will never recognize Him.

\subsection{413}\label{section-908}

\begin{otherlanguage}{persian}

قوله عزّ ذكره فاشهد بعين فؤادك و لاتنظر اليه الّا بعينه فانّ من ینظر اليه
بعينه يدركه و الّا يحتجب ان اردت اللّه و لقائه فارده فانظر اليه و هم‌‌چنين
اعزازاً لنفسه و اظهاراً لشأنه در مقام ديگر ميفرمايد قوله جلّ اعزازه الا انّه
جلّ ذكره يعرف كلّ شیء نفسه و انّى استحيى ان اقول يعرف كلّ شیء نفسه بمثل ما
انّى قد عرفت كلّ شیء بآياتى اذ كلّ ما تجد من كلّ شیء خلق له و انّ اللّه اجلّ و
اعلى من ان يعرف بخلقه بل الخلق يعرف به انتهى

\end{otherlanguage}

Qawluhu, ʿazza dhikruhu: Fa-shhad bi-ʿayni fuʾādika wa-lā tanẓur ilayhi
illā bi-ʿaynihi, fa-inna man yanẓuru ilayhi bi-ʿaynihi yudrikuhu wa-illā
yaḥtajibu. In aradta Allāha wa-liqāʾahu fa-aridhu fa-nẓur ilayhi. Va
hamchunīn iʿzāzan li-nafsihi va iẓhāran li-shaʾnihi dar maqām-i dīgar
mifarmāyad, qawluhu, jalla iʿzāzuhu: Alā innahu, jalla dhikruhu,
yuʿarrifu kulla shayʾin nafsahu, wa-innī astaḥyī an aqūla yaʿrifu kullu
shayʾin nafsahu bi-mithli mā innī qad ʿarraftu kulla shayʾin bi-āyātī,
idh kullu mā tajidu min kulli shayʾin khuliqa lahu, wa-inna Allāha
ajallu wa-aʿlā min an yuʿrafa bi-khalqihi, bali al-khalqu yuʿrafu bihi.
Intahā.

His words---exalted be His mention---are: ``Behold with the eye of your
heart, and look upon Him only with His own eye, for whoever looks upon
Him with His own eye perceives Him; otherwise, he remains veiled. If you
desire God and attainment to His presence, desire Him and look upon
Him.'' Likewise, to honor His Self and manifest His station, He says
elsewhere---glorious be His honor: ``Indeed, He---glorious be His
mention---makes His Self known to all things. I am ashamed to say that
all things know His Self, in the same way that I have made all things
known through My verses, for whatever you find among all things was
created for Him. God is too glorious and exalted to be known through His
creation; rather, creation is known through Him.'' End of quotation.

\subsection{414}\label{section-909}

\begin{otherlanguage}{persian}

ملاحظه كن كه بچه ادب و اعزاز در جميع اين بياناتيكه متعلّق بظهور بعد است
ذكر ميفرمايند ميفرمايد كه او بنفسه ميشناساند نفس مقدّس خود را بكلّ اشیآء و
من حيا ميكنم اينكه بگويم عارف ميشود نفس او را كلّ شیء اگر در اين بيانات
لسان رحمن اقلّ من آن تفكّر نمائى و بفواكه طيّبه عرفان كه در او مستور است
مرزوق شوى صنم وهم را كه مصنوع ايادى ظنون است شكسته بشطر قدم توجّه نمائى و
لكن بسيار بعيد است كه باين توفيق موفّق شوى چه كه جسارتت از حدّ گذشته

\end{otherlanguage}

Mulāḥiẓih kun kih bi-chih adab va iʿzāz dar jamīʿ-i īn bayānātī kih
mutaʿalliq bi-ẓuhūr-i baʿd ast dhikr mifarmāyand. Mifarmāyad kih ū
bi-nafsih mishināsānad nafs-i muqaddas-i khud rā bi-kull-i ashyāʾ va man
ḥayā mikunam īnkih bigūyam ʿārif mīshavad nafs-i ū rā kull-i shayʾ. Agar
dar īn bayānāt-i Lisān-i Raḥmān aqall-i min ān tafakkur namāʾī va
bi-favākih-i ṭayyibih-yi ʿirfān kih dar ū mastūr ast marzūq shavī,
ṣanam-i vahm rā kih maṣnūʿ-i ayādī-yi ẓunūn ast shikastih bi-shaṭr-i
Qidam tavajjuh namāʾī. Va lākin bisyār baʿīd ast kih bi-īn tawfīq
muvaffaq shavī, chih kih jisāratat az ḥadd gudhashtih.

Observe with what courtesy and honor He speaks in all these utterances
concerning the subsequent Revelation. He says that He Himself makes His
holy Self known to all things, and: ``I am ashamed to say that all
things know His Self.'' If you were to reflect for less than a moment on
these utterances of the Tongue of the All-Merciful and be nourished by
the wholesome fruits of recognition concealed within them, you would
shatter the idol of delusion, fashioned by the hands of conjecture, and
turn toward the Eternal One. Yet it is most unlikely that you will
receive this grace, for your insolence has exceeded all bounds.

\subsection{415}\label{section-910}

\begin{otherlanguage}{persian}

بحقّ نقطه وجود و جمال معبود كه اين عبد متحيّر است از امثال شما اين كلمات
ربّ اعلى و تو ذكر ميرزا جمال و ميرزا كمال استهزاء على اللّه مينمائى و همچه
گمان نموده‌ايد كه حرمة اللّه مابين ناس باين كلمات ضايع ميشود چنانچه معرضين
باللّه بهمين كلمات مفتريات بر جمال رحمن بسته و لكن غافل از آنكه نزد هر ذى
بصرى از همين كلماتى كه مذكور داشته‌ايد احقر و اشرّ و ابعد از هر ذكر دون
خيرى بوده و خواهيد بود هر ذى شأنى از شئون امثال آن انفس مشركه احتراز
جويد و هر ذى قدرى از هتك حرمتى كه نموده‌ايد بخدا پناه برد

\end{otherlanguage}

Bi-ḥaqq-i Nuqṭih-yi vujūd va Jamāl-i Maʿbūd kih īn ʿabd mutaḥayyir ast
az amthāl-i shumā. Īn kalimāt-i Rabb-i Aʿlā va tū dhikr-i Mīrzā Jamāl va
Mīrzā Kamāl istihzāʾan ʿala'llāh minamāʾī, va hamchih gumān namūdeh-īd
kih ḥurmatu'llāh mābayn-i nās bi-īn kalimāt ḍāyiʿ mīshavad, chunānchih
muʿriḍīn bi'llāh bi-hamīn kalimāt muftariyāt bar Jamāl-i Raḥmān bastih.
Va lākin ghāfil az ānkih nazd-i har dhī baṣarī az hamīn kalimātī kih
madhkūr dāshteh-īd aḥqar va asharr va abʿad az har dhikr-i dūn-i khayrī
būdih va khvāhīd būd. Har dhī shaʾnī az shuʾūn-i amthāl-i ān anfus-i
mushrikih iḥtirāz jūyad va har dhī qadrī az hatk-i ḥurmatī kih
namūdeh-īd bi-Khudā panāh barad.

By the Point of existence and the Beauty of the Adored One, this servant
is bewildered by people such as you. Such are the words of the Most
Exalted Lord, yet you speak of Mīrzā Jamāl and Mīrzā Kamāl in mockery of
God! You have imagined that these words will undermine the reverence for
God among the people, just as those who turn away from God have used
these very words to fabricate slanders against the Beauty of the
All-Merciful. Yet you are unaware that, in the eyes of everyone with
insight, these very words you have uttered have made and will make you
more contemptible, more wicked, and more remote than any description of
evil can convey. Everyone of standing shuns the conduct of such souls
who associate partners with God, and everyone of worth seeks refuge in
God from the desecration you have committed.

\subsection{416}\label{section-911}

\begin{otherlanguage}{persian}

و هم‌‌چنين از اين بيانات مذكوره ثابت و محقّق كه دون نفس مقدّسش دليل او نبوده
و نخواهد بود و سؤال از دون او لعرفان نفس او لم‌يزل مردود بوده و خواهد بود
و هر نفسى اراده نمايد بآن ساذج غيب عارف شود بايد بنفس او بما يظهر من
عنده عارف گردد نه بكلمات قبل و بعد فو‌الّذى نفسى بيده اوراد منيعه احديّه در
روضه اين كلمات الهيّه مشهود و اثمار معارف جنيّه از سدرات كلمات مرقومه
موجود و لكن دون صاحبان بصر و شامّه و ذائقه احديرا از آن قسمتى نه و اليوم
احدى از اهل بصر و ذائقه و فؤاد مذكور نه الّا اهل البهآء الّذين نوّر اللّه
وجوههم بأنوار جماله و قدّسهم عمّا سواه و عرفهم ما اضطربت عنه القلوب و
استقرّهم فى ظلّ سدرة امره و انقطعهم عن العالمين

\end{otherlanguage}

Va hamchunīn az īn bayānāt-i madhkūrih thābit va muḥaqqaq kih dūn-i
nafs-i muqaddasash dalīl-i ū nabūdih va nakhvāhad būd va suʾāl az dūn-i
ū li-ʿirfān-i nafs-i ū lam-yazal mardūd būdih va khvāhad būd. Va har
nafsī irādih namāyad bi-ān Sādhij-i Ghayb ʿārif shavad bāyad bi-nafs-i ū
bi-mā yaẓharu min ʿindihi ʿārif gardad nah bi-kalimāt-i qabl va baʿd.
Fa-wa'lladhī nafsī bi-yadihi, avrād-i manīʿih-yi aḥadiyyih dar rawḍih-yi
īn kalimāt-i ilāhiyyih mashhūd va athmār-i maʿārif-i janiyyih az
sidrāt-i kalimāt-i marqūmih mawjūd, va lākin dūn-i ṣāḥibān-i baṣar va
shāmmih va dhāʾiqih aḥadī rā az ān qismatī nah. Va al-yawm aḥadī az
ahl-i baṣar va dhāʾiqih va fuʾād madhkūr nah illā ahl-i al-Bahāʾ,
alladhīna nawwara Allāhu wujūhahum bi-anwāri jamālihi wa-qaddasahum
ʿammā siwāhu wa-ʿarrafahum mā iḍṭarabat ʿanhu al-qulūbu wa-istaqarrahum
fī ẓilli Sidrati Amrihi wa-inqaṭaʿahum ʿani al-ʿālamīn.

Likewise, these utterances establish beyond doubt that none other than
His holy Self has been or will be His proof, and that asking anyone
other than Him in order to recognize His Self has always been and will
remain unacceptable. Whoever wishes to recognize that Essence of the
Unseen must recognize His own Self through what proceeds from Him, not
through the words of those who came before or come after. By Him in
whose hand is my soul! In the garden of these divine words the exalted
roses of oneness are visible, and from the lote-trees of these written
words the ripe fruits of spiritual knowledge are available. Yet no one
has a share in them except those endowed with sight, smell, and taste.
Today none are accounted among those endowed with sight, taste, and
heart except the people of Bahāʾ, whose faces God has illumined with the
lights of His beauty, whom He has sanctified from all besides Him, to
whom He has made known that which has thrown hearts into turmoil, whom
He has established in the shade of the Lote-Tree of His Cause, and whom
He has detached from all the worlds.

\subsection{417}\label{section-912}

\begin{otherlanguage}{persian}

و حال چون اهل بيان بغير رضآء اللّه سايرند و در تيه وهم سالك و از مقصود
اللّه و آياته محروم چنانچه نوشته كه اين اسم مبارك اعظم اقدم ابهى در كجاى
بيان و بچه عبارت است لذا اين عبد از مقرّ اقصى و سدرة منتهى كه مقام اعرفوا
اللّه باللّه است تنزّل نموده و بارض تحديد كه مقام ذكر و اشارات است ناظر كه
شايد آنچه در اينمقام اين عبد مذكور دارد بپذيرند و بحبل كلمات الهى از
غرقاب فانى نجات يافته بعرصه نورانى باقى خرامند

\end{otherlanguage}

Va ḥāl chūn ahl-e Bayān be-ghayr-e reżā-ye Allāh sāyerand va dar tīh-e
vahm sālek va az maqṣūd-e Allāh va āyātih maḥrūm, chenānche neveshte ke
īn esm-e mobārak-e aʿẓam-e aqdam-e abhā dar kojā-ye Bayān va be-che
ʿebārat ast, ledhā īn ʿabd az maqarr-e aqṣā va sedrat-e montahā ke
maqām-e iʿrifū Allāha bi-Allāh ast tanazzol namūde va be-arż-e taḥdīd ke
maqām-e dhekr va eshārāt ast nāẓer, ke shāyad ānche dar īn maqām īn ʿabd
madhkūr dārad bepadhīrand va be-ḥabl-e kalemāt-e elāhī az gharqāb-e fānī
nejāt yāfte be-ʿarṣe-ye nūrānī-ye bāqī kherāmand.

Now, since the people of the Bayān travel contrary to God's good
pleasure, walk in the wilderness of delusion, and remain deprived of
God's intended One and His verses---as he has written, asking where in
the Bayān this blessed, Most Great, Most Ancient, Most Glorious Name
appears, and in what terms---this servant has therefore descended from
the most distant abode and the Lote-Tree of the Utmost Limit, the
station of ``Recognize God through God,'' and turned toward the realm of
limitation, the station of mention and allusion. Perhaps they will
accept what this servant states in this station and, holding to the cord
of divine words, escape the perishing depths and walk into the
everlasting realm of light.

\subsection{418}\label{section-913}

\begin{otherlanguage}{persian}

اينكه نوشته تعريفات حضرت باب بجهت شخص معيّن و يا بجهت اسم بهاء اوّلا
بگوئيد كه حضرت اعلى روح ما سواه فداه كه خود را بذكر ناميده و جميع اسماء
حسنى را بنفس مبارك خود راجع فرموده بچه دليل تصديقش نموديد و بچه برهان
موقن و مذعن گشتيد بهمان دليل و برهان حال هم ناظر شويد جميع اسماء و ملكوت
آن بقول مظهر ظهور خلق شده و ميشوند و كلّ طايف حول و لكن انّك لا تعرف و
تكون من الغافلين ميفرمايند آنچه در بيان ذكر اسماء حسنى شده كلّ مقصود ظهور
بعد بوده و خواهد بود و هم‌‌چنين آنچه ذكر خير كه در بيان نازل شده در حقيقت
اوّليّه مقصود ظهور بعد بوده و در حقيقت ثانويّه اوّل مؤمن باو

\end{otherlanguage}

Īnke neveshte taʿrīfāt-e Ḥażrat-e Bāb be-jehat-e shakhṣ-e moʿayyan va yā
be-jehat-e esm-e Bahāʾ, avvalan begūyīd ke Ḥażrat-e Aʿlā, rūḥ-e mā sevāh
fedāh, ke khod rā be-Dhekr nāmīde va jamīʿ-e asmā-ye ḥosnā rā be-nafs-e
mobārak-e khod rājeʿ farmūde, be-che dalīl taṣdīqash nemūdīd va be-che
borhān mūqen va modhʿen gashtīd? Be-hamān dalīl va borhān ḥāl ham nāẓer
shavīd. Jamīʿ-e asmā va malakūt-e ān be-qowl-e maẓhar-e ẓohūr khalq
shode va mīshavand va koll ṭāyef-e ḥowl, wa-lākin innaka lā taʿrifu
wa-takūnu mina al-ghāfilīn. Mīfarmāyand ānche dar Bayān dhekr-e asmā-ye
ḥosnā shode koll maqṣūd ẓohūr-e baʿd būde va khāhad būd va hamchenīn
ānche dhekr-e khayr ke dar Bayān nāzel shode dar ḥaqīqat-e avvaliyye
maqṣūd ẓohūr-e baʿd būde va dar ḥaqīqat-e thānaviyye avval moʾmen be-ū.

As to his writing about whether the Bāb's praises refer to a particular
person or to the name Bahāʾ, first say: By what evidence did you
acknowledge the Most Exalted One---may the spirit of all besides Him be
sacrificed for Him---who called Himself the Remembrance and referred all
the most beautiful names to His own blessed Self? By what proof did you
become certain and assent? Look now to that same evidence and proof. All
names and their kingdom have been and are created by the word of the
Manifestation of Revelation, and all circle around Him, but you do not
recognize this and are among the heedless. He says that wherever the
most beautiful names are mentioned in the Bayān, the intended referent
has always been and will always be the subsequent Revelation. Likewise,
every mention of good revealed in the Bayān has, in its primary reality,
referred to the subsequent Revelation, and in its secondary reality, to
the first to believe in Him.

\subsection{419}\label{section-914}

\begin{otherlanguage}{persian}

اليوم بنصّ نقطه بيان نفوسى كه از اين امر بديع معرضند از رداء اسميّه وصفيّه
محروم و كلّ از بهائم بين يدى اللّه محشور و مذكور قل مت بغيظك يا ايّها الّذى
تنكر هذا الفضل و تموت حين استماعه انّه قد ظهر بالحقّ و يطوف فى حوله مظاهر
الأسمآء و ملكوتها ثمّ حقايقها و مسمّياتها لو انت من العارفين و انّه لن
يحتاج الى غيره و لا باسمآء الّتى كانت بينكم بل كلّ محتاج اليه و كلّ من فضله
لمن السّائلين اذا ينوحنّ كينونات الأسمآء من ظلمك و فعلك و يستعيذن باللّه
منك و من امثالك و يقولن مخاطباً اليك يا ايّها الغافل انّا خلقنا و بعثنا
بأمر من لدى اللّه المقتدر العزيز القدير و نكون خادماً لحضرته و طائفاً حول
بابه و مفتخراً بنسبتنا اليه و هو قد كان غنيّاً عنّا ان انت من المستشعرين و
انّك لما اشتعلت بنار الحسد و البغضآء ما رضيت بانّ الّذى خلق بأمره ملكوت
الأسمآء ينسب الى نفسه اسمآء منها و هذا من ظلم الّذى ما ظهر شبهه فى
الابداع و اذا ينوحنّ قبايل مداين البقآء من ظلمكم يا ملأ المغلّين

\end{otherlanguage}

Al-yowm be-naṣṣ-e Noqṭe-ye Bayān nofūsī ke az īn amr-e badīʿ moʿreżand
az redā-ye esmiyye-ye vaṣfiyye maḥrūm va koll az bahāʾem bayna yaday
Allāh maḥshūr va madhkūr. Qul mut bi-ghayẓika yā ayyuhā alladhī tunkiru
hādhā al-faḍla wa-tamūtu ḥīna istimāʿihi. Innahu qad ẓahara bi-al-ḥaqqi
wa-yaṭūfu fī ḥawlihi maẓāhiru al-asmāʾi wa-malakūtuhā thumma ḥaqāʾiquhā
wa-musammayātuhā law anta mina al-ʿārifīn. Wa-innahu lan yaḥtāja ilā
ghayrihi wa-lā bi-asmāʾi allatī kānat baynakum bal kullun muḥtājun
ilayhi wa-kullun min faḍlihi la-mina al-sāʾilīn. Idhan yanūḥanna
kaynūnātu al-asmāʾi min ẓulmika wa-fiʿlika wa-yastaʿidhna bi-Allāhi
minka wa-min amthālika wa-yaqulna mukhāṭiban ilayka yā ayyuhā al-ghāfilu
innā khuliqnā wa-buʿithnā bi-amrin min ladā Allāhi al-Muqtadiri
al-ʿAzīzi al-Qadīr. Wa-nakūnu khādiman li-ḥaḍratihi wa-ṭāʾifan ḥawla
bābihi wa-muftakhiran bi-nisbatinā ilayhi wa-huwa qad kāna ghaniyyan
ʿannā in anta mina al-mustashʿirīn. Wa-innaka lammā ishtaʿalta bi-nāri
al-ḥasadi wa-al-baghḍāʾi mā raḍīta bi-anna alladhī khuliqa bi-amrihi
malakūtu al-asmāʾi yansibu ilā nafsihi asmāʾan minhā wa-hādhā min
al-ẓulmi alladhī mā ẓahara shibhuhu fī al-ibdāʿi wa-idhan yanūḥanna
qabāʾilu madāʾini al-baqāʾi min ẓulmikum yā malaʾa al-mughillīn.

Today, according to the explicit words of the Point of the Bayān, souls
who turn away from this wondrous Cause are deprived of the robe of names
and attributes, and all are gathered and accounted as beasts before God.
Say: Die in your rage, you who deny this grace and die upon hearing of
it! He has indeed appeared in truth, and around Him circle the
manifestations of the names and their kingdom, and then their realities
and the things they designate, if you are among those who understand. He
will never need anyone else, nor the names current among you. Rather,
all stand in need of Him, and all ask for His bounty. Then the very
beings of the names lament your injustice and your deeds, seek refuge in
God from you and your kind, and address you, saying: ``Heedless one! We
were created and raised up by a command from God, the All-Powerful, the
Almighty, the Omnipotent. We serve His presence, circle around His door,
and glory in our relationship to Him, while He has ever been independent
of us, if you are among those who perceive. Because you have become
ablaze with the fire of envy and hatred, you have refused to allow Him,
through whose command the kingdom of names was created, to ascribe to
Himself any of those names. This is an injustice the like of which has
never appeared in creation.'' Then the peoples of the cities of eternity
lament your injustice, you company of the malicious!

\subsection{420}\label{section-915}

\begin{otherlanguage}{persian}

قل انّه لهو الّذى يفتخر الأسمآء بعبد من عباده لو تنسب اليه او يسمّى بها و
انّك ما استشعرت فى نفسك و كنت من المبعدين قل تاللّه انّى لعلىّ فى ملكوت
البقآء و محمّد فى جبروت الأسمآء ثمّ الرّوح فى مداين البقآء ثمّ بالحسين فى
هذا الظّهور الكبرى و لنا اسمآء اخرى فى ممالك القدم الّتى ما اطّلع بها احد
الّا اللّه الفرد العالم الخبير مت بغيظك يا ايّها الغافل انّ شرافة الّتى قدرت
للأسمآء انّها كانت لنسبتها الى نفسى العزيز العليم و ما ارتفع اسم فى الملك
الّا بتوجّهه الى شطرى المقدّس المتعالى العزيز المنيع فو‌نفسى كلّ اسم خير يرجع
الى نفسى و كلّ ذكر بديع ينتهى الى جمالى ان انت من الموقنين

\end{otherlanguage}

Qul innahu la-huwa alladhī yaftakhiru al-asmāʾu bi-ʿabdin min ʿibādihi
law tunsabu ilayhi aw yusammā bihā wa-innaka mā istashʿarta fī nafsika
wa-kunta mina al-mubʿadīn. Qul ta-Allāhi innī la-ʿAlīyun fī malakūti
al-baqāʾi wa-Muḥammadun fī jabarūti al-asmāʾi thumma al-Rūḥu fī madāʾini
al-baqāʾi thumma bi-al-Ḥusayni fī hādhā al-ẓuhūri al-kubrā wa-lanā
asmāʾun ukhrā fī mamāliki al-qidami allatī mā iṭṭalaʿa bihā aḥadun illā
Allāhu al-Fardu al-ʿĀlimu al-Khabīr. Mut bi-ghayẓika yā ayyuhā
al-ghāfilu inna sharāfata allatī quddirat li-al-asmāʾi innahā kānat
li-nisbatihā ilā nafsī al-ʿAzīzi al-ʿAlīm. Wa-mā irtafaʿa ismun fī
al-mulki illā bi-tawajjuhihi ilā shaṭrī al-muqaddasi al-mutaʿālī
al-ʿazīzi al-manīʿ. Fa-wa-nafsī kullu ismi khayrin yarjiʿu ilā nafsī
wa-kullu dhikrin badīʿin yantahī ilā jamālī in anta mina al-mūqinīn.

Say: He is the One whose servant the names glory in, should they be
attributed to that servant or should he be called by them. Yet you have
not perceived this within yourself and have been among those kept far
away. Say: By God! I am indeed ʿAlī in the kingdom of eternity, Muḥammad
in the dominion of names, then the Spirit in the cities of eternity, and
then called Ḥusayn in this greatest Revelation. We have other names in
the realms of ancient eternity of which no one is aware except God, the
Unique, the All-Knowing, the All-Informed. Die in your rage, heedless
one! The honor ordained for the names is due to their relationship to My
Self, the Almighty, the All-Knowing. No name has been exalted in the
world except by turning toward My sanctified, transcendent, mighty,
unassailable presence. By My Self! Every name of goodness refers to My
Self, and every wondrous mention leads to My beauty, if you are among
those who are certain.

\subsection{421}\label{section-916}

\begin{otherlanguage}{arabic}

و لو تسمّى دونى بكلّ الأسماء لن يصدق عليه بل انّ حقايق تلك الأسماء يلعنك
حين الّذى تخرج من فمك و يفرّن منك و يرجعنّ الى مقرّ الأقصى هذا المقام المقدّس
الممتنع الرّفيع انّك لو تسمّى الظّلمة نوراً هل يصدق ذلك الاسم عليها لا فو‌ربّ
العالمين او تسمّى الخزف بالّلؤلؤ هل ينبغى هذا الاسم له لا فو‌ربّك الرّحمن
الرّحيم ا تعترض بالّذى كلّ الأسماء سجّاد لطلعته و لن يتحرّكن الّا باذنه و لن
يستضيئن الّا ببهائه المشرق المتعالى المقدّس اللّايح المنير خف عن اللّه و لا
تتّبع هواك و لا تتعرّض بالّذى قدّسه اللّه عن كلّ اسم و رسم ثمّ عن كلّ ذكر و بيان
و جعله بنفسه مدلّاً لنفسه و بذاته معترفاً لذاته العزيز المنيع

\end{otherlanguage}

Wa-law tusammī dūnī bi-kulli al-asmāʾi lan yaṣduqa ʿalayhi bal inna
ḥaqāʾiqa tilka al-asmāʾi yalʿanuka ḥīna alladhī takhruju min famika
wa-yafirarna minka wa-yarjiʿanna ilā maqarri al-aqṣā hādhā al-maqāmi
al-muqaddasi al-mumtaniʿi al-rafīʿ. Innaka law tusammī al-ẓulmata nūran
hal yaṣduqu dhālika al-ismu ʿalayhā? Lā fa-wa-Rabbi al-ʿālamīn. Aw
tusammī al-khazafa bi-al-luʾluʾi hal yanbaghī hādhā al-ismu lahu? Lā
fa-wa-Rabbika al-Raḥmāni al-Raḥīm. A-taʿtariḍu bi-alladhī kullu
al-asmāʾi sujjādun li-ṭalʿatihi wa-lan yataḥarrakna illā bi-idhnihi
wa-lan yastaḍiʾna illā bi-bahāʾihi al-mushriqi al-mutaʿālī al-muqaddasi
al-lāʾiḥi al-munīr? Khaf ʿan Allāhi wa-lā tattabiʿ hawāka wa-lā
tataʿarraḍ bi-alladhī qaddasahu Allāhu ʿan kulli ismin wa-rasmin thumma
ʿan kulli dhikrin wa-bayānin wa-jaʿalahu bi-nafsihi mudillan li-nafsihi
wa-bi-dhātihi muʿtarifan li-dhātihi al-ʿAzīzi al-Manīʿ.

Even if you call someone other than Me by all the names, they will not
apply to him. Rather, the realities of those names curse you as they
leave your mouth, flee from you, and return to the most distant abode,
this sanctified, inaccessible, exalted station. If you call darkness
light, does that name apply to it? No, by the Lord of the worlds! Or if
you call earthenware a pearl, does this name befit it? No, by your Lord,
the All-Merciful, the Compassionate! Do you object to Him before whose
countenance all names bow down, moving only by His permission and
shining only through His radiant, transcendent, sanctified, resplendent,
luminous glory? Fear God, do not follow your desires, and do not assail
Him whom God has sanctified above every name and designation, and above
every mention and exposition, and whom He has made, through His own
Self, the guide to Himself, and through His own Essence, the
acknowledger of His Essence, the Almighty, the Unassailable.

\subsection{422}\label{section-917}

\begin{otherlanguage}{arabic}

انّ ربّك قد خلق ملكوت الأسماء كخلق كلّ شیء بقوله كن فيكون لو انت من
النّاظرين و لو یشآء ليرجعها الى عدم البحت و انّه لهو المقتدر على ما یشآء و
انّه لهو العليم الحكيم اسمع قولى و تب الى اللّه الّذى خلقك لعرفان نفسه و
لاتجادل بأمره بعد ظهوره و لاتكن من المغلّين طهّر قلبك ثمّ اذنك لتسمع نغمات
اللّه من لسانه الأطهر الأنور الصّادق الأمين هل تظنّ بانّ غيره يقدر ان يأتى
بمثل ما اتى لا فو‌نفسه المهيمن القيّوم او دونه يقدر ان يتكلّم تلقآء الوجه
لا فو‌جماله العزيز المحبوب انّه لن يشتبه بدونه و ما يظهر منه ليكون مقدّساً
عن كلّ ما ظهر فى العالمين لو انتم تعلمون

\end{otherlanguage}

Inna Rabbaka qad khalaqa malakūta al-asmāʾi ka-khalqi kulli shayʾin
bi-qawlihi kun fa-yakūnu law anta mina al-nāẓirīn. Wa-law yashāʾu
la-yurjiʿuhā ilā ʿadami al-baḥti wa-innahu la-huwa al-muqtadiru ʿalā mā
yashāʾu wa-innahu la-huwa al-ʿAlīmu al-Ḥakīm. Ismaʿ qawlī wa-tub ilā
Allāhi alladhī khalaqaka li-ʿirfāni nafsihi wa-lā tujādil bi-amrihi
baʿda ẓuhūrihi wa-lā takun mina al-mughillīn. Ṭahhir qalbaka thumma
udhunaka li-tasmaʿa naghamāti Allāhi min lisānihi al-aṭhari al-anwari
al-ṣādiqi al-amīn. Hal taẓunnu bi-anna ghayrahu yaqdiru an yaʾtiya
bi-mithli mā atā? Lā fa-wa-nafsihi al-Muhaymini al-Qayyūm. Aw dūnahu
yaqdiru an yatakallama tilqāʾa al-wajhi? Lā fa-wa-jamālihi al-ʿAzīzi
al-Maḥbūb. Innahu lan yashtabiha bi-dūnihi wa-mā yaẓharu minhu la-yakūnu
muqaddasan ʿan kulli mā ẓahara fī al-ʿālamīna law antum taʿlamūn.

Your Lord has created the kingdom of names just as He creates
everything, through His word ``Be,'' and it is, if you are among those
who see. Should He wish, He would return them to utter nothingness. He
indeed has power over whatever He wills, and He is the All-Knowing, the
All-Wise. Listen to My words and turn in repentance to God, who created
you to recognize His Self. Do not dispute His Cause after His
appearance, and do not be among the malicious. Purify your heart, and
then your ears, so that you may hear the melodies of God from His most
pure, most luminous, truthful, trustworthy tongue. Do you imagine that
anyone else can produce the like of what He has brought? No, by His
Self, the All-Protecting, the Self-Subsisting! Or can anyone else speak
before His face? No, by His beauty, the Almighty, the Beloved! He can
never be confused with anyone else, and whatever proceeds from Him is
sanctified above all that has appeared in the worlds, if only you knew.

\subsection{423}\label{section-918}

\begin{otherlanguage}{persian}

قل انّه مرّة يسمّى نفسه بكلّ الأسماء و مرّة ينزّهها عنها و عن كلّ ما خلق بين
السّموات و الأرض و عن كلّ ما كان و ما يكون قل كلّ الأمور فى قبضته يفعل ما
یشآء بأمره و من قال لم و بم فقد كفر باللّه المقتدر المهيمن القيّوم خف عن
اللّه و لا تتّبع سبيل اوّل من اعرض باللّه تجنّب منه ثمّ استغفر ربّك لعلّ يهبّ على
حقيقتك ما يذهب به روايح هؤلآء و انّه لهو الغفور العطوف ان سمعت فلنفسك و
ان انكرت و اعرضت فعليها فانّ ربّك الرّحمن لغنىّ عمّا فى الأكوان و انّه لهو
المقدّس المتعالى العزيز المحمود

\end{otherlanguage}

Qul innahu marratan yusammī nafsahu bi-kulli al-asmāʾi wa-marratan
yunazzihuhā ʿanhā wa-ʿan kulli mā khuliqa bayna al-samāwāti wa-al-arḍi
wa-ʿan kulli mā kāna wa-mā yakūn. Qul kullu al-umūri fī qabḍatihi
yafʿalu mā yashāʾu bi-amrihi wa-man qāla lima wa-bima fa-qad kafara
bi-Allāhi al-Muqtadiri al-Muhaymini al-Qayyūm. Khaf ʿan Allāhi wa-lā
tattabiʿ sabīla awwali man aʿraḍa bi-Allāhi tajannab minhu thumma
istaghfir Rabbaka laʿalla yahubbu ʿalā ḥaqīqatika mā yadhhabu bihi
rawāʾiḥu hāʾulāʾi wa-innahu la-huwa al-Ghafūru al-ʿAṭūf. In samiʿta
fa-li-nafsika wa-in ankarta wa-aʿraḍta fa-ʿalayhā fa-inna Rabbaka
al-Raḥmāna la-ghaniyyun ʿammā fī al-akwāni wa-innahu la-huwa
al-Muqaddasu al-Mutaʿālī al-ʿAzīzu al-Maḥmūd.

Say: At one time He calls Himself by all the names, and at another He
exalts Himself above them, above everything created between the heavens
and the earth, and above all that has been and will be. Say: All affairs
are in His grasp. By His command He does whatever He wills. Whoever says
``Why?'' or ``By what means?'' has disbelieved in God, the All-Powerful,
the All-Protecting, the Self-Subsisting. Fear God and do not follow the
path of the first to turn away from God. Shun him, then seek your Lord's
forgiveness, so that there may blow upon your inner reality that which
will dispel the odors of these people. He is indeed the Ever-Forgiving,
the Most Kind. If you listen, it is for your own good; if you deny and
turn away, it is to your own detriment. Your Lord, the All-Merciful, is
independent of all that exists, and He is the Sanctified, the Most
Exalted, the Almighty, the All-Praised.

\subsection{424}\label{section-919}

\begin{otherlanguage}{persian}

بدانكه آنچه لفظ بهآء در كلمات اللّه بوده بجهت هيكل معيّن مشخّص مشهود ظاهر
واضح نازل شده و كلّ رحمت اليوم از براى كسى است كه مؤمن باشد باو و جميع
نقمت از براى نفوسى است كه باو موقن نشوند و معرض گردند بلكه آن اسمائى را
كه درباره اوّل معرض شنيده و باو راجع كرده جميع مخصوص اين ظهور بوده و
معرضين من حيث لا يشعر تحريف نموده‌اند و نيست اين ظلم مگر از آن اصل شجره
اعراض و فساد

\end{otherlanguage}

Bedānke ānche lafẓ-e Bahāʾ dar kalemāt-e Allāh būde be-jehat-e haykal-e
moʿayyan-e moshakhkhaṣ-e mashhūd-e ẓāher-e vāżeḥ nāzel shode va koll-e
raḥmat al-yowm az barā-ye kasī ast ke moʾmen bāshad be-ū va jamīʿ-e
neqmat az barā-ye nofūsī ast ke be-ū mūqen nashavand va moʿreż gardand.
Balke ān asmāʾī rā ke dar bāre-ye avval moʿreż shenīde va be-ū rājeʿ
karde jamīʿ makhṣūṣ-e īn ẓohūr būde va moʿreżīn min ḥaythu lā yashʿur
taḥrīf nemūde-and va nīst īn ẓolm magar az ān aṣl-e shajare-ye eʿrāż va
fasād.

Know that every occurrence of the word Bahāʾ in the words of God was
revealed for a particular, distinct, visible, manifest, unmistakable
person. Today all mercy is for whoever believes in Him, and all
retribution is for those souls who do not become certain of Him and turn
away. Indeed, those names that you have heard concerning the first to
turn away and have referred to him were all reserved for this
Revelation. Those who turn away have distorted them without realizing
it. This injustice comes from none other than that root of the tree of
rejection and corruption.

\subsection{425}\label{section-920}

\begin{otherlanguage}{persian}

و از اينها گذشته اگر ذكرى از اين ظهور در بيان نباشد بزعم تو و مرشدينت و
لفظ بهآء هم در كتاب نازل نگشته باشد چه ضرر بظهور دارد ميفرمايد من بهمان
حجّت نقطه اولى ظاهر شده‌ام بل اعظم لو انتم تشعرون و واللّه الّذى لا اله الّا
هو كه ابداً دوست نداشته‌اند كه جز بنفس ظهور و آيات منزله احدى در اثبات
امرشان استدلال نمايد بلى اين سؤال شما در وقتى جايز كه صاحب ظهور بفرمايد
بيّنه من اسم من است فتعالى من ذلك بيّنه نفس ظهور بوده و خواهد بود و من
دونه ما ظهر من عنده و نزل من لدنه و اينست آن حجّتى كه جميع انبیاء و رسل
بآن مبعوث شده‌اند و من اعرض فقد اعرض عن اللّه بارئه و كفر بآياته و كان من
المشركين من قلم اللّه على اللّوح مذكوراً

\end{otherlanguage}

Va az īnhā godhashte agar dhekrī az īn ẓohūr dar Bayān nabāshad
be-zaʿm-e to va morshedīnat va lafẓ-e Bahāʾ ham dar Ketāb nāzel nagashte
bāshad che żarar be-ẓohūr dārad? Mīfarmāyad man be-hamān ḥojjat-e
Noqṭe-ye Ūlā ẓāher shode-am, bal aʿẓam law antum tashʿurūn. Va wa-Allāhi
alladhī lā ilāha illā huwa ke abadan dūst nadāshte-and ke joz be-nafs-e
ẓohūr va āyāt-e monzale aḥadī dar ethbāt-e amreshān estedlāl namāyad.
Balī īn soʾāl-e shomā dar vaqtī jāyez ke ṣāḥeb-e ẓohūr befarmāyad
bayyene-ye man esm-e man ast, fa-taʿālā min dhālik. Bayyene nafs-e ẓohūr
būde va khāhad būd va min dūnihi mā ẓahara min ʿindihi va nazala min
ladunhi va īnast ān ḥojjatī ke jamīʿ-e anbiyā va rosol be-ān mabʿūth
shode-and, wa-man aʿraḍa fa-qad aʿraḍa ʿani Allāhi bāriʾihi wa-kafara
bi-āyātihi wa-kāna mina al-mushrikīna min qalami Allāhi ʿalā al-lawḥi
madhkūran.

Beyond all this, even if, as you and your guides suppose, this
Revelation were not mentioned in the Bayān and the word Bahāʾ had not
been revealed in the Book, what harm would that do to the Revelation? He
says: ``I have appeared with the same proof as the Primal
Point---indeed, a greater one, if only you perceived.'' By God, besides
whom there is no god, He has never wished anyone to argue in support of
His Cause through anything other than the Revelation itself and the
verses revealed. Your question would be legitimate only if the Bearer of
Revelation were to say, ``My evidence is My name.'' Exalted is He above
that! The evidence has always been and will always be the Revelation
itself, and besides it, whatever has proceeded from Him and descended
from His presence. This is the proof with which all the Prophets and
Messengers have been sent. Whoever turns away has turned away from God,
his Creator, disbelieved in His verses, and been recorded by the Pen of
God upon the Tablet as one of those who associate partners with Him.

\subsection{426}\label{section-921}

\begin{otherlanguage}{persian}

و خود نقطه بيان هم روح ما سواه فداه جز بآيات استدلال نفرموده‌اند چنانچه
در مجالس عديده ايشان را حاضر نمودند و برهان طلبيدند جز بآيات از آن مكمن
اسماء و صفات ظاهر نشد و هر يك از عباد هم كه طلب حجّت و برهان نمودند جواب
صادر كه بآيات الهى ناظر باشيد چه كه حجّت لم‌يزل آيات اللّه بوده و كلّ من على
الأرض را همين حجّت كافى است و در اين ظهور اعظم اكثرى از ناس آنچه
خواسته‌اند مشاهده نموده‌اند چه از آيات منزله الهى كه احدى شبه آن نديده و
چه از بيانات علميّه و خوارق عادتيّه مع‌ذلك بحقيقت و ذات و كينونت و نفس نقطه
اولى كافر شده‌اند و شاعر نيستند تازه ذكر اسماء مذكور ميداريد لازال ناس از
سلطان معلوم معرض بوده‌اند و بموهوم انفس خود مقبل

\end{otherlanguage}

Va khod-e Noqṭe-ye Bayān ham, rūḥ-e mā sevāh fedāh, joz be-āyāt estedlāl
nafarmūde-and, chenānche dar majāles-e ʿadīde īshān rā ḥāżer nemūdand va
borhān ṭalabīdand joz be-āyāt az ān makman-e asmā va ṣefāt ẓāher nashod
va har yek az ʿebād ham ke ṭalab-e ḥojjat va borhān nemūdand javāb ṣāder
ke be-āyāt-e elāhī nāẓer bāshīd, che ke ḥojjat lam-yazal āyāt-e Allāh
būde va koll-e man ʿalā al-arż rā hamīn ḥojjat kāfī ast. Va dar īn
ẓohūr-e aʿẓam aktharī az nās ānche khāste-and moshāhede nemūde-and, che
az āyāt-e monzale-ye elāhī ke aḥadī shebh-e ān nadīde va che az
bayānāt-e ʿelmiyye va khavāreq-e ʿādatiyye, maʿa-dhālek be-ḥaqīqat va
dhāt va kaynūnat va nafs-e Noqṭe-ye Ūlā kāfer shode-and va shāʿer
nīstand. Tāze dhekr-e asmā madhkūr mīdārīd. Lā-zāl nās az Solṭān-e
maʿlūm moʿreż būde-and va be-mowhūm-e anfos-e khod moqbel.

The Point of the Bayān Himself---may the spirit of all besides Him be
sacrificed for Him---likewise offered no proof other than verses. Thus,
when they brought Him before numerous assemblies and demanded proof,
nothing but verses appeared from that sanctuary of names and attributes.
Whenever any of His servants requested evidence and proof, the answer
was given: ``Look to the divine verses,'' for the proof has always been
the verses of God, and this proof alone suffices for everyone on earth.
In this Most Great Revelation, most people have witnessed whatever they
desired, whether divinely revealed verses the like of which no one has
seen, learned expositions, or extraordinary feats. Nevertheless, they
have disbelieved in the reality, essence, being, and very Self of the
Primal Point, and do not realize it. Now you bring up the names! People
have always turned away from the known Sovereign and toward the
imaginings of their own selves.

\subsection{427}\label{section-922}

\begin{otherlanguage}{persian}

اينكه نوشته اگر بجهت مسمّاى بعينه و شخص بخصوص است در كجاى بيان و بچه
عبارت است لذا خالصاً لوجه اللّه اين عبد بعضى آيات منزله كه از سماء مشيّت
علىّ اعلى در اين اسم مبارك اعزّ اقدم و ظهور امنع اقدس نازل شده ذكر مينمايم
اگر چه واللّه الّذى لا اله الّا هو هرگز جمال ابهى جز آيات منزله كه از سماء
احديّت نازل است حجّت و ميزان امر نفرموده‌اند اگر چه همين كلمات را اعتراضاً
على اللّه ذكر نموده و ابداً مقصودت فهميدن نبوده لذا اين عبد هم اتماماً لحجّة
اللّه بعضى كلمات منزل بيان را ذكر مينمايم و انّك لو لم تهتدى بها كما ما
اهتديت من قبل لعلّ يهدى بها عباد آخرون اگر چه حجّت نفس او بوده و لكن
ميخواهم همين قدر بر خود تو معلوم شود كه از آن اشخاصى هستى كه ميفرمايد و
ان يروا كلّ آية لم يؤمنوا بها و در مقام ديگر قوله تعالى انّ الّذين حقّت
عليهم كلمة ربّك لايؤمنون و لو جائتهم كلّ آية حتّى يروا العذاب الأليم

\end{otherlanguage}

Īnkih nivishtih agar bi-jihat-i musammā-yi bi-ʿaynih va shakhṣ-i
bi-khuṣūṣ ast dar kujā-yi Bayān va bi-chih ʿibārat ast, lidhā khāliṣan
li-wajhi Allāh īn ʿabd baʿżī āyāt-i munazzalih kih az samā-yi
mashiyyat-i ʿAlī-yi Aʿlā dar īn ism-i mubārak-i aʿazz-i aqdam va ẓuhūr-i
amnaʿ-i aqdas nāzil shudih dhikr mīnamāyam. Agarchih wa-Allāh alladhī lā
ilāha illā huwa hargiz Jamāl-i Abhā juz āyāt-i munazzalih kih az samā-yi
aḥadiyyat nāzil ast ḥujjat va mīzān-i amr nafarmūdih-and. Agarchih hamīn
kalimāt rā iʿtirāḍan ʿalā Allāh dhikr namūdih va abadan maqṣūdat
fahmīdan nabūdih, lidhā īn ʿabd ham itmāman li-ḥujjati Allāh baʿżī
kalimāt-i Munazzil-i Bayān rā dhikr mīnamāyam. Wa-innaka law lam tahtadī
bihā kamā mā ihtadayta min qabl laʿalla yahdī bihā ʿibād ākharūn.
Agarchih ḥujjat nafs-i ū būdih va lākin mīkhvāham hamīn qadr bar khud-i
tu maʿlūm shavad kih az ān ashkhāṣī hastī kih mīfarmāyad: wa-in yaraw
kulla āyah lam yuʾminū bihā. Va dar maqām-i dīgar qawluhu taʿālā: inna
alladhīna ḥaqqat ʿalayhim kalimatu rabbika lā yuʾminūn wa-law jāʾathum
kullu āyah ḥattā yaraw al-ʿadhāb al-alīm.

As for what he has written---``If it concerns a specifically named
individual and a particular person, where in the Bayān is this, and in
what words?''---solely for the sake of God, this servant will therefore
mention some of the revealed verses that have descended from the heaven
of the will of the Most Exalted One concerning this blessed, most
mighty, most ancient Name and this most unassailable, most holy
Revelation. Yet, by God, besides whom there is no god, the All-Glorious
Beauty has never designated anything as the proof and standard of His
Cause except the revealed verses that descend from the heaven of divine
unity. Although you have raised these very words as an objection to God
and have never intended to understand, this servant will nevertheless
mention some of the words of the Revealer of the Bayān, to make God's
proof complete. Even if you are not guided by them, just as you were not
guided before, perhaps other servants will be guided by them. Although
He Himself is the proof, I wish only that you yourself should realize
that you are among those of whom He says: ``Even if they see every sign,
they will not believe in it.'' And elsewhere, His words---exalted is He:
``Those against whom your Lord's word has come into effect will not
believe, even if every sign comes to them, until they see the painful
punishment.''

\subsection{428}\label{section-923}

\begin{otherlanguage}{persian}

چه كه شمس مشرق و بصدهزار اشراق و انوار مضئ و نقطه بيان هم در جميع الواح
ميفرمايد او را بخود او بشناسيد و بآيات او و فى‌الحقيقة بر كلّ نفوس حرام
است كه بغير نفس اقدس امنع و ما يظهر من عنده باو استدلال نمايند و لكن اين
عبد اين ارتكاب را حمل مينمايد و از سلطان قدم چشم عفو داشته و دارم و اگر
اين عبد بخواهد آنچه در اين فقره كه خواسته ذكر نمايد الواح كفايت نكند چه
كه لايحصى از جبروت بقاء نازل

\end{otherlanguage}

Chih kih shams mushriq va bi-ṣad-hizār ishrāq va anvār muḍīʾ va Nuqṭiy-i
Bayān ham dar jamīʿ-i alvāḥ mīfarmāyad ū rā bi-khud-i ū bishināsīd va
bi-āyāt-i ū, va fī al-ḥaqīqih bar kull-i nufūs ḥarām ast kih bi-ghayr-i
nafs-i aqdas-i amnaʿ va mā yaẓharu min ʿindihi bi-ū istidlāl namāyand.
Va lākin īn ʿabd īn irtikāb rā ḥaml mīnamāyad va az Sulṭān-i Qidam
chashm-i ʿafv dāshtih va dāram. Va agar īn ʿabd bikhvāhad ānchih dar īn
fiqrih kih khvāstih dhikr namāyad alvāḥ kifāyat nakunad, chih kih lā
yuḥṣā az jabarūt-i baqā nāzil.

For the Sun has risen and shines with a hundred thousand radiances and
lights, and the Point of the Bayān likewise says in all His Tablets:
``Recognize Him through Himself and through His verses.'' Indeed, it is
forbidden for any soul to adduce proof of Him through anything other
than His most holy, most unassailable Self and what appears from Him.
Yet this servant takes upon himself this transgression and has looked,
and still looks, to the Sovereign of Eternity for pardon. If this
servant wished to mention everything concerning the matter he has asked
about, tablets would not suffice, for countless utterances have
descended from the dominion of everlasting life.

\subsection{429}\label{section-924}

\begin{otherlanguage}{arabic}

لذا بچند فقره اكتفا ميرود و آن اينست قوله عزّ كبريائه و انّ بمثل ذلك نزلنا
القرآن من قبل و لكنّكم كنتم عن مرادى محتجبون ذلك ما طاف اللّيل و النّهار
عليه ثمانية واحد و انتم به فى العبادة تتوحّدون و كنتم عن سرّه بعد ما قضى
لمحتجبون ذلك ميزان الهدى فى البيان انتم به مؤمنون الى حين ما يشرق شمس
البهآء ذلك ظهور اللّه ان تعملنّ به لمؤمنون و انتم فى الرّضوان خالدون و الّا
انتم فانيون انتهى

\end{otherlanguage}

Lidhā bi-chand fiqrih iktifā mīravad va ān īnast, qawluhu ʿazza
kibriyāʾuhu: wa-inna bi-mithli dhālika nazzalnā al-Qurʾān min qabl
wa-lākinnakum kuntum ʿan murādī muḥtajibūn. Dhālika mā ṭāfa al-layl
wa-al-nahār ʿalayhi thamāniyat wāḥid wa-antum bihi fī al-ʿibādah
tatawaḥḥadūn wa-kuntum ʿan sirrihi baʿda mā quḍiya la-muḥtajibūn.
Dhālika mīzān al-hudā fī al-Bayān antum bihi muʾminūn ilā ḥīn mā
yushriqu shams al-Bahāʾ. Dhālika ẓuhūr Allāh in taʿmalunna bihi
la-muʾminūn wa-antum fī al-Riḍwān khālidūn wa-illā antum fāniyūn.
Intahā.

A few passages will therefore suffice. They are as follows, in His
words---exalted be His majesty: ``Indeed, in like manner We sent down
the Qurʾān before, but you were veiled from My intent. That is what
night and day revolved around for eight vāḥids, and through it you
profess unity in worship; yet you remained veiled from its mystery after
it had been fulfilled. That is the standard of guidance in the Bayān, in
which you believe until the time when the Sun of Bahāʾ rises. That is
the Revelation of God. If you act according to it, you are indeed
believers and will abide forever in Riḍwān; otherwise, you will
perish.'' End of quotation.

\subsection{430}\label{section-925}

\begin{otherlanguage}{persian}

و اليوم ظهور اللّه ظاهر و شمس بهآء از افق اعلى مشرق و مضئ و معرضين فانى و
معدوم صرف چنانچه مشاهده مينمائيد لو تكوننّ من المنصفين كذلك كان حيّاً فى
الأفق الأبهی لو انتم تقدرنّ ان تكوننّ من السّامعين او تكوننّ من النّاظرين
اقسمك باللّه الّذى انطق كلّشیء بثنآء نفسه كه بعد از اين كلمه محكمه و آيهء
متقنه و اين ظهور عزّ ربانيّه ديگر از براى نفسى مجال اعراض و يا توقّف مانده
لا فو‌نفسه الحقّ لو انتم تعلمون در مقام ديگر ميفرمايد قوله جلّ اجلاله و عزّ
اعزازه من اوّل ما تطلع شمس البهآء الى ان تغرب خير فى كتاب اللّه عن كلّ
اللّيل ان انتم تدركون ما خلق اللّه من شیء الّا ليؤمئذ اذ كلّ للقآء اللّه ثمّ
رضائه يعملون انتهى

\end{otherlanguage}

Va al-yawm Ẓuhūr-i Allāh ẓāhir va Shams-i Bahāʾ az ufuq-i aʿlā mushriq
va muḍīʾ va muʿriḍīn fānī va maʿdūm-i ṣirf, chunānchih mushāhidih
mīnamāyīd, law takūnunna min al-munṣifīn. Kadhālika kāna ḥayyan fī
al-ufuq al-Abhā law antum taqdirunna an takūnunna min al-sāmiʿīn aw
takūnunna min al-nāẓirīn. Uqsimuka bi-Allāh alladhī anṭaqa kulla shayʾ
bi-thanāʾi nafsihi kih baʿd az īn kalimiy-i muḥkamih va āyiy-i mutqanih
va īn ẓuhūr-i ʿizz-i rabbāniyyih dīgar az barā-yi nafsī majāl-i iʿrāḍ va
yā tavaqquf māndih? Lā fa-wa-nafsihi al-Ḥaqq law antum taʿlamūn. Dar
maqām-i dīgar mīfarmāyad, qawluhu jalla ijlāluhu va ʿazza iʿzāzuhu: min
awwal mā taṭluʿu shams al-Bahāʾ ilā an taghrub khayr fī kitāb Allāh ʿan
kull al-layl in antum tudrikūn. Mā khalaqa Allāh min shayʾ illā
{[}corrupted text{]} idh kull li-liqāʾ Allāh thumma riḍāʾihi yaʿmalūn.
Intahā.

Today the Revelation of God is manifest, and the Sun of Bahāʾ has risen
from the highest horizon and shines, while those who turn away are
perishing and utterly nonexistent, as you can see, if you are
fair-minded. Thus He was living in the All-Glorious Horizon, if you can
be among those who hear or those who see. I adjure you by God, who has
caused all things to speak in praise of Himself: after this decisive
utterance, this firmly established verse, and this revelation of divine
glory, is any room left for a soul to turn away or hesitate? No, by His
own true Self, if only you knew! Elsewhere He says---His words, exalted
be His glory and mighty His exaltation: ``From the first rising of the
Sun of Bahāʾ until its setting is better, in the Book of God, than all
the night, if you perceive. God has created nothing except {[}corrupted
original text{]}, for all labor to attain the presence of God and then
His good pleasure.'' End of quotation.

\subsection{431}\label{section-926}

\begin{otherlanguage}{persian}

بصر اطهر بايد تا بمنظر اكبر ناظر شود و معنى كلمات الهى را ادراك نمايد
بشنو ندآء سلطان بيان را كه بكمال تصريح ظهور بعد را ذكر فرموده‌اند كه لعلّ
آن هياكل اضلال بعد از ظهور توقّف ننمايند و شبه و مثل و نظير براى مظهر
توحيد نه‌بينند چه كه توحيد ذات محققّ نشود مگر بمظهر ظهور و تنزيه الهى از
اشباح و امثال مبرهن نگردد مگر بمظهر ظهور و اگر از براى مظهر ظهور شبه و
ندّ و مثلى ملحوظ گردد تنزيه ذات قدم از مثليّت معيّن نگردد و تقديس نفسش از
اثنينيّت محقّق نيايد فاستحى عن اللّه و لاتجعل له شريكاً فى الملك انّه كان
واحداً فى ذاته و كان اللّه على ما اقول شهيد و در اينمقام ميفرمايد قوله جلّ
كبريائه و لقد قرب الزّوال و انكم انتم ذلك اليوم لا تعرفون و من يكن لقائه
ذات لقائى لاترضين له ما لا ترضى نفس لنفس انتهى

\end{otherlanguage}

Baṣar-i aṭhar bāyad tā bi-manẓar-i akbar nāẓir shavad va maʿnī-yi
kalimāt-i ilāhī rā idrāk namāyad. Bishinaw nidā-yi Sulṭān-i Bayān rā kih
bi-kamāl-i taṣrīḥ ẓuhūr-i baʿd rā dhikr farmūdih-and kih laʿall ān
hayākil-i iḍlāl baʿd az ẓuhūr tavaqquf nanamāyand va shabah va mithl va
naẓīr barā-yi Maẓhar-i Tawḥīd nabīnand. Chih kih tawḥīd-i dhāt muḥaqqaq
nashavad magar bi-Maẓhar-i Ẓuhūr va tanzīh-i ilāhī az ashbāḥ va amthāl
mubarhan nagardad magar bi-Maẓhar-i Ẓuhūr. Va agar az barā-yi Maẓhar-i
Ẓuhūr shabah va nidd va mithlī malḥūẓ gardad tanzīh-i Dhāt-i Qidam az
mithliyyat muʿayyan nagardad va taqdīs-i nafsash az ithnayniyyat
muḥaqqaq nayāyad. Fa-istaḥyi ʿan Allāh wa-lā tajʿal lahu sharīkan fī
al-mulk innahu kāna wāḥidan fī dhātihi wa-kāna Allāh ʿalā mā aqūlu
shahīd. Va dar īn maqām mīfarmāyad, qawluhu jalla kibriyāʾuhu: wa-laqad
qaruba al-zawāl wa-innakum antum dhālika al-yawm lā taʿrifūn. Wa-man
yakun liqāʾuhu dhāta liqāʾī lā tarḍayanna lahu mā lā tarḍā nafs li-nafs.
Intahā.

The purest sight is needed to behold the Most Great Vision and apprehend
the meaning of the divine words. Listen to the call of the Sovereign of
the Bayān, who has spoken of the next Revelation with complete
explicitness, so that perhaps those embodiments of misguidance may not
hesitate after its appearance or see any likeness, peer, or equal to the
Manifestation of divine unity. For the oneness of the Essence is
established only through the Manifestation of Revelation, and God's
transcendence above likenesses and equals is demonstrated only through
the Manifestation of Revelation. If any likeness, rival, or equal is
envisaged for the Manifestation of Revelation, the transcendence of the
eternal Essence above all likeness is not established, nor is His own
sanctity above duality demonstrated. Feel shame before God, and assign
Him no partner in sovereignty. He is one in His Essence, and God is
witness to what I say. In this connection He says---His words, exalted
be His majesty: ``Noon has indeed drawn near, yet you do not recognize
that Day. As for Him whose presence is My very presence, do not consent
for Him to what no soul would consent to for itself.'' End of quotation.

\subsection{432}\label{section-927}

\begin{otherlanguage}{persian}

مع‌ذلك كلّ از حدود خود تجاوز نموده و بوصاياى الهيّه اعتنا ننمودند و راضى
شدند آنچه را كه يهود از براى روح راضى نشد و هم‌‌چنين امّت فرقان از براى
نقطه بيان راضى نشده و اين نيست مگر از جهل و عدم بلوغ و اگر تا حال نشنيدى
حال بشنو و از حقّ شرم نما از شمال ضلالت بيمين هدايت راجع شو عنقريب تو و
امثال تو و آنچه بآن تكلّم نموده‌ايد فانى خواهيد شد چنانچه اليوم فانى هستيد
چه كه هر نفسيكه بين يدى اللّه مذكور نيست كلّ اشیآء بر فناى او شهادت ميدهند
طهّر نفسك عن الدّنيا و لا تقل ما لا علمته و لا تذكر ما لا عرفته فاعرف حدّك
و مقدارك و لا تجاوز عن شأنك

\end{otherlanguage}

Maʿa dhālik kull az ḥudūd-i khud tajāvuz namūdih va bi-vaṣāyā-yi
ilāhiyyih iʿtinā nanamūdand va rāḍī shudand ānchih rā kih Yahūd az
barā-yi Rūḥ rāḍī nashud va hamchunīn ummat-i Furqān az barā-yi Nuqṭiy-i
Bayān rāḍī nashudih. Va īn nīst magar az jahl va ʿadam-i bulūgh. Va agar
tā ḥāl nashunīdī ḥāl bishinaw va az Ḥaqq sharm namā. Az shimāl-i ḍalālat
bi-yamīn-i hidāyat rājiʿ shaw. ʿAn-qarīb tu va amthāl-i tu va ānchih
bi-ān takallum namūdih-īd fānī khvāhīd shud, chunānchih al-yawm fānī
hastīd, chih kih har nafsī kih bayna yaday Allāh madhkūr nīst kull-i
ashyāʾ bar fanā-yi ū shahādat mīdihand. Ṭahhir nafsaka ʿan al-dunyā
wa-lā taqul mā lā ʿalimtahu wa-lā tadhkur mā lā ʿaraftahu fa-iʿraf
ḥaddaka wa-miqdāraka wa-lā tajāwaz ʿan shaʾnika.

Despite this, all have overstepped their bounds and paid no heed to the
divine counsels. They have consented to what the Jews did not consent to
for the Spirit, nor the people of the Qurʾān for the Point of the Bayān.
This comes only from ignorance and immaturity. If you have not heard
until now, hear now, and feel shame before God. Turn from the left hand
of error to the right hand of guidance. Before long, you and those like
you, and everything you have uttered, will perish, just as you are
perished today, for all things testify to the extinction of any soul who
is not remembered in the presence of God. Purify yourself from the
world, do not say what you have not learned, and do not mention what you
have not understood. Know your limits and your measure, and do not
exceed your station.

\subsection{433}\label{section-928}

\begin{otherlanguage}{persian}

تصريح ميفرمايند كه ظهور شمس در قطب زوال نزديك است و شما آن روز را
نميشناسيد و ادراك نمينمائيد چنانچه ادراك ننموده‌ايد و عارف نگشته‌ايد و
ابدا آن جمال قدم در ذكر ظهور بعد جز تصريح بيانى نفرموده‌اند بر خلاف كتب
قبل كه جميع باشارات معضله ذكر فرموده‌اند انامل قدرت در ظهور ستّين جميع آن
حجبات را خرق نمود و ذكر ظهور بعد را در جميع بيان تصريحا من غير تلويح
فرمود

\end{otherlanguage}

Taṣrīḥ mīfarmāyand kih ẓuhūr-i shams dar quṭb-i zavāl nazdīk ast va
shumā ān rūz rā namīshināsīd va idrāk namīnamāyīd, chunānchih idrāk
nanamūdih-īd va ʿārif nagashtih-īd. Va abadan ān Jamāl-i Qidam dar
dhikr-i ẓuhūr-i baʿd juz taṣrīḥ bayānī nafarmūdih-and, bar khilāf-i
kutub-i qabl kih jamīʿ bi-ishārāt-i muʿḍilih dhikr farmūdih-and.
Anāmil-i qudrat dar ẓuhūr-i sittīn jamīʿ-i ān ḥujubāt rā kharq namūd va
dhikr-i ẓuhūr-i baʿd rā dar jamīʿ-i Bayān taṣrīḥan min ghayr-i talvīḥ
farmūd.

He states explicitly that the appearance of the Sun at its noonday
zenith is near, and that you do not recognize or apprehend that Day,
just as you have neither apprehended nor recognized it. In speaking of
the next Revelation, that eternal Beauty has never spoken otherwise than
explicitly, unlike the earlier Books, which all spoke in abstruse
allusions. In the Revelation of the year sixty, the fingers of power
tore through all those veils and spoke of the next Revelation throughout
the Bayān explicitly, without allusion.

\subsection{434}\label{section-929}

\begin{otherlanguage}{persian}

امرى را كه ببصر عباد و درايتهم گذارده اينست كه در قرب و بعد آن مختلف ذكر
فرموده‌اند اگر چه در باطن جميع كلمات مختلفه بنقطه واحده راجع و ابدا
اختلافى ملحوظ نه و لكن همين قدر كه بر حسب ظاهر فى‌الجمله مختلف ذكر شده
جهت امتياز بوده چنانچه در مقامى ذكر مستغاث فرموده‌اند و در مقامى
ميفرمايند كه بمستغاث نميرسد و اميد است كه از قبل كلمة اللّه مرتفع گردد و
در مقامى ميفرمايد چه كسى عالم بظهور نيست غير اللّه هر وقت شود بايد كلّ
تصديق بنقطه حقيقت نمايند و در مقامى ميفرمايند كه اين حين اگر بخواهد ظاهر
شود مختار بوده و ليس لأحد ان يقول لم و بم و بيكى از حروفات كه سؤال از آن
ظهور اعظم نموده ميفرمايد شايد تو لقاى او را ادراك نمائى و هم‌‌چنين لوحى از
سماء مشيّت نازل و فرموده‌اند در مكتب‌خانه ظهور بعد ملاحظه فرمايد

\end{otherlanguage}

Amrī rā kih bi-baṣar-i ʿibād va dirāyatihim gudhāshtih īnast kih dar
qurb va buʿd-i ān mukhtalif dhikr farmūdih-and. Agarchih dar bāṭin
jamīʿ-i kalimāt-i mukhtalifih bi-nuqṭiy-i vāḥidih rājiʿ va abadan
ikhtilāfī malḥūẓ nah, va lākin hamīn qadr kih bar ḥasab-i ẓāhir fī
al-jumlih mukhtalif dhikr shudih jihat-i imtiyāz būdih. Chunānchih dar
maqāmī dhikr-i Mustaghāth farmūdih-and va dar maqāmī mīfarmāyand kih
bi-Mustaghāth namīrasad va umīd ast kih az qabl Kalimat-i Allāh murtafiʿ
gardad. Va dar maqāmī mīfarmāyad chih kasī ʿālim bi-ẓuhūr nīst ghayr
Allāh, har vaqt shavad bāyad kull taṣdīq bi-Nuqṭiy-i Ḥaqīqat namāyand.
Va dar maqāmī mīfarmāyand kih īn ḥīn agar bikhvāhad ẓāhir shavad mukhtār
būdih, wa-laysa li-aḥad an yaqūla lima wa-bima. Va bi-yikī az ḥurūfāt
kih suʾāl az ān ẓuhūr-i aʿẓam namūdih mīfarmāyad shāyad tu liqā-yi ū rā
idrāk namāʾī. Va hamchunīn lawḥī az samā-yi mashiyyat nāzil va
farmūdih-and dar maktab-khāniy-i ẓuhūr-i baʿd mulāḥiẓih farmāyad.

The matter He has left to the sight and discernment of His servants is
this: He has spoken variously of its nearness and remoteness. Although
inwardly all these differing words refer to a single point and no
difference whatsoever is envisaged, the measure of outward difference in
His statements has been for the sake of distinction. Thus, in one
passage He has mentioned Mustaghāth, while in another He says that it
will not reach Mustaghāth, and expresses the hope that the Word of God
will be raised up before then. Elsewhere He says that no one knows of
the Revelation except God, and that whenever it occurs, all must
acknowledge the Point of Truth. In another passage He says that if He
wishes to appear at this very moment, He is free to do so, and no one
has the right to say why or wherefore. To one of the Letters who asked
about that Most Great Revelation, He says: ``Perhaps you will attain His
presence.'' Likewise, a Tablet descended from the heaven of His will,
and He directed that it be perused in the school of the next Revelation.

\subsection{435}\label{section-930}

\begin{otherlanguage}{persian}

و اين اختلافات ظاهريّه قوليّه كلّ بكلمه محكمه رفع ميشود و تحذير بليغ
ميفرمايند كه در حين ظهور آن جوهر مقصود بما نزل فى البيان ناظر نباشيد
بلكه بنفس ظهور و ما يظهر من عنده ناظر باشيد و بر صاحبان بصر معلوم بوده
كه قبل و بعد و قرب و بعد در ساحت حقّ نبوده و نخواهد بود چنانچه زمن اوّلين
و آخرين را اگر بخواهد در يك آن مبعوث فرمايد قادر بوده و خواهد بود چه كه
اين حدودات در خلق مذكور چنانچه خمسين الف سنه قيامت در ساعتى منقضى شد

\end{otherlanguage}

Va īn ikhtilāfāt-i ẓāhiriyyiy-i qawliyyih kull bi-kalimiy-i muḥkamih
rafʿ mīshavad va taḥdhīr-i balīgh mīfarmāyand kih dar ḥīn-i ẓuhūr-i ān
Jawhar-i Maqṣūd bi-mā nuzila fī al-Bayān nāẓir nabāshīd, balkih
bi-nafs-i ẓuhūr va mā yaẓharu min ʿindihi nāẓir bāshīd. Va bar ṣāḥibān-i
baṣar maʿlūm būdih kih qabl va baʿd va qurb va buʿd dar sāḥat-i Ḥaqq
nabūdih va nakhvāhad būd. Chunānchih zaman-i avvalīn va ākhirīn rā agar
bikhvāhad dar yik ān mabʿūth farmāyad qādir būdih va khvāhad būd, chih
kih īn ḥudūdāt dar khalq madhkūr, chunānchih khamsīn alf sanih qiyāmat
dar sāʿatī munqaḍī shud.

All these apparent differences of wording are resolved by a decisive
utterance. He gives a most emphatic warning that, when that desired
Essence appears, you should not look to what has been revealed in the
Bayān, but rather to the Revelation itself and what appears from Him.
Those endowed with sight know that before and after, nearness and
remoteness, have never existed and will never exist in the court of God.
Thus, if He wishes to bring forth the time of the first and the last in
a single instant, He has always been and will always be able to do so,
for these limitations pertain to creation, just as the fifty thousand
years of the Resurrection elapsed in an hour.

\subsection{436}\label{section-931}

\begin{otherlanguage}{persian}

اى بيهوشان خمر جهل و غفلت تا كى در ظلمات سالكيد و بهواى نفس عامل دل از
تحديدات بشريّه برداريد و بمقرّ سلطان احديّه ناظر شويد كه شايد بمقام اعرفوا
اللّه باللّه وارد گرديد و از چون و چرا فارغ ائيد از شمس سؤال ننمائيد كه
اسمت چيست و دليل ضيائت چه

\end{otherlanguage}

Ay bīhūshān-i khamr-i jahl va ghaflat, tā kay dar ẓulumāt sālikīd va
bi-havā-yi nafs ʿāmil? Dil az taḥdīdāt-i bashariyyih bardārīd va
bi-maqarr-i Sulṭān-i Aḥadiyyih nāẓir shavīd kih shāyad bi-maqām-i iʿrifū
Allāh bi-Allāh vārid gardīd va az chūn va chirā fārigh āyīd. Az shams
suʾāl nanamāyīd kih ismat chīst va dalīl-i ḍiyāʾat chih.

You who are senseless with the wine of ignorance and heedlessness! How
long will you walk in darkness and act according to your own desires?
Detach your hearts from human limitations and look toward the seat of
the Sovereign of Oneness, so that perhaps you may enter the station of
``Know God through God'' and be free of how and why. Do not ask the Sun,
``What is your name, and what is the evidence of your light?''

\subsection{437}\label{section-932}

\begin{otherlanguage}{persian}

اى بی‌ بصر دليل اشراقش اشراقش و برهان ضيائش ضيائش بوده و خواهد بود بشنو
كلمه عليم خبير را كه تصريحا اخبار فرموده كه در اين ظهور قدس صمدانى كلّ
هالك الّا الّذينهم تمسّكوا بفلك القدم فى هذا الظّهور الأعظم و انّهم لأهل
سفينة القدس عند اللّه المقتدر العلىّ العظيم و من تمسّك بهذا الفلك فقد نجى و
من اعرض فقد غرق و انّ هذا لتنزيل من لدن عليم خبير

\end{otherlanguage}

Ay bī-baṣar, dalīl-i ishrāqash ishrāqash va burhān-i ḍiyāʾash ḍiyāʾash
būdih va khvāhad būd. Bishinaw kalimiy-i ʿAlīm-i Khabīr rā kih taṣrīḥan
ikhbār farmūdih kih dar īn ẓuhūr-i quds-i ṣamadānī kull hālik illā
alladhīnahum tamassakū bi-fulk al-qidam fī hādhā al-ẓuhūr al-aʿẓam
wa-innahum la-ahl safīnat al-quds ʿinda Allāh al-Muqtadir al-ʿAlī
al-ʿAẓīm. Wa-man tamassaka bi-hādhā al-fulk fa-qad najā wa-man aʿraḍa
fa-qad ghariqa wa-inna hādhā la-tanzīl min ladun ʿAlīm Khabīr.

You who are without sight! The evidence of its radiance has always been
and will always be its radiance, and the proof of its light, its light.
Listen to the word of the All-Knowing, the All-Informed, who has
explicitly foretold that in this holy and everlasting Revelation all
perish except those who hold fast to the Ark of Eternity in this Most
Great Revelation. They are indeed the people of the Ark of Holiness in
the sight of God, the All-Powerful, the Most High, the Most Great.
Whoever holds fast to this Ark is saved, and whoever turns away is
drowned. This is indeed a revelation from One who is All-Knowing,
All-Informed.

\subsection{438}\label{section-933}

\begin{otherlanguage}{persian}

قوله جلّ اجلاله و عظم استقلاله يا ملأ الأنوار انّا نحن تاللّه الحقّ ما ننطق
عن الهوى و ما ننزل حرفا من ذلك الكتاب الّا باذن اللّه الحقّ اتّقوا اللّه و لا
تشكّوا فى امر اللّه فانّ سرّ هذا الباب مستورا تحت عمآء السّطر و مرقوماً مافوق
حجاب السّتر بايدى اللّه ربّ السّتر و السّطر و لقد خلق اللّه فى حول ذلك الباب
بحورا من مآء الاكسير محمر بالدّهن الوجود و حيواناً بالثّمرة المقصود و قدّر
اللّه له سفنا من ياقوتة الرّطبة الحمراء و لا يركب فيها الّا اهل البهآء باذن
اللّه العلىّ و هو اللّه قد كان عزيزا و حكيماً هنالك يحمل عرش اللّه ملائكة
العمآء فى الأنفس الثّمان و قد كان الحكم فى امّ الكتاب مشهودا فسوف يعرضون
المجرمون على الرّحمن فى حول العرش على الحقّ بالحقّ رمزا انتهى

\end{otherlanguage}

Qawluhu jalla ijlāluhu va ʿaẓuma istiqlāluhu: yā malaʾ al-anwār innā
naḥnu ta-Allāh al-Ḥaqq mā nanṭiqu ʿan al-hawā wa-mā nunazzilu ḥarfan min
dhālika al-kitāb illā bi-idhni Allāh al-Ḥaqq. Ittaqū Allāh wa-lā
tashukkū fī amr Allāh fa-inna sirra hādhā al-bāb mastūran taḥta ʿamāʾ
al-saṭr wa-marqūman mā fawqa ḥijāb al-satr bi-aydī Allāh rabb al-satr
wa-al-saṭr. Wa-laqad khalaqa Allāh fī ḥawl dhālika al-bāb buḥūran min
māʾ al-iksīr muḥmarr bi-al-duhn al-wujūd wa-ḥayawānan bi-al-thamarah
al-maqṣūd wa-qaddara Allāh lahu sufunan min yāqūtat al-raṭbah al-ḥamrāʾ
wa-lā yarkabu fīhā illā ahl al-Bahāʾ bi-idhni Allāh al-ʿAlī wa-huwa
Allāh qad kāna ʿazīzan wa-ḥakīman. Hunālika yaḥmilu ʿarsh Allāh
malāʾikat al-ʿamāʾ fī al-anfus al-thamān wa-qad kāna al-ḥukm fī umm
al-kitāb mashhūdan. Fa-sawfa yuʿraḍūna al-mujrimūn ʿalā al-Raḥmān fī
ḥawl al-ʿarsh ʿalā al-ḥaqq bi-al-ḥaqq ramzan. Intahā.

His words---exalted be His glory and great His independence: ``O
concourse of lights! By God, the True One, We do not speak from desire,
nor do We send down a single letter of that Book except by the
permission of God, the True One. Fear God and do not doubt the Cause of
God, for the mystery of this Gate is concealed beneath the cloud of the
written line and inscribed above the veil of concealment by the hands of
God, the Lord of concealment and of the written line. Around that Gate
God has indeed created seas of the water of the elixir, reddened with
the oil of existence, and life through the intended fruit. God has
ordained for it ships of moist red ruby, and none but the people of
Bahāʾ embark upon them, by the permission of God, the Most High. He is
God, ever mighty and wise. There the angels of the cloud, in the eight
souls, bear the throne of God, and the decree is witnessed in the Mother
Book. The guilty will then be brought before the All-Merciful around the
throne, in truth and by the Truth, symbolically.'' End of quotation.

\subsection{439}\label{section-934}

\begin{otherlanguage}{persian}

اگر ادراك ميداشتيد بهمين آيات منزلات كه در قيّوم كتب نازل شده جميع در
سبيل اين ظهور كلّيه الهيّه جان ايثار مينموديد و لكن اخذ شعور از امثال آن
نفوس كلّ شده بما اكتسبت ايّديكم لذا از فيوضات رحمن و نيسان حكمت حضرت سبحان
در اين ربيع روحانى و سلطان فصول و ایّام الهى محروم گشته‌ايد و از رضاى مليك
علّام دور مانده‌ايد و در تيه شهوات و روضه اشارات مثل حيوانات بچريدن مشغول

\end{otherlanguage}

Agar edrāk mīdāshtīd, be-hamīn āyāt-e monazzalāt ke dar Qayyūm-e kotob
nāzel shode, jamīʿ dar sabīl-e īn ẓohūr-e kollīye-ye elāhīye jān īthār
mīnamūdīd, va lāken akhdh-e shoʿūr az amthāl-e ān nofūs koll shode bi-mā
iktasabat aydīkum. Ledhā az foyūḍāt-e Raḥmān va naysān-e ḥekmat-e
Ḥażrat-e Sobḥān dar īn rabīʿ-e rūḥānī va solṭān-e foṣūl va ayyām-e elāhī
maḥrūm gashte-īd va az reżā-ye Malīk-e ʿAllām dūr mānde-īd va dar tīh-e
shahavāt va rowże-ye eshārāt methl-e ḥayvānāt be-charīdan mashghūl.

Had you possessed understanding, these very revealed verses sent down in
the Qayyūm of Books would have led you all to sacrifice your lives in
the path of this universal divine Revelation. Yet understanding has been
taken away from all souls such as these because of what your hands have
earned. Thus, in this spiritual springtime, this sovereign of seasons
and these days of God, you have been deprived of the outpourings of the
All-Merciful and the April rain of the wisdom of the All-Glorious. You
have remained far from the good pleasure of the all-knowing King and are
occupied, like animals, with grazing in the wilderness of desires and
the garden of allusions.

\subsection{440}\label{section-935}

\begin{otherlanguage}{persian}

بارى اليوم جز راكبين اين فلك عزّ صمدانى احدى عند‌اللّه مذكور نه اين نفوسند
كه اقلّ از كبريت احمرند و زود است كه مشاهده نمائى كه اكثر از ناس من حيث
لا يشعر در غمرات بحر جهل و اعراض مستغرق شوند و از فلك الهى محروم مانند
الّا عدّة معدودات ديگر تا چه زمان ارياح فضليّه بوزيدن آيد و بر هياكل مردگان
نفس و هوى مرور نمايد و ايشانرا حيات بخشد و بشطر احديّه متوجّه سازد

\end{otherlanguage}

Bārī, al-yowm joz rākebīn-e īn folk-e ʿezz-e ṣamadānī aḥadī ʿinda Allāh
madhkūr na. Īn nofūsand ke aqall az kebrīt-e aḥmarand va zūd ast ke
moshāhede namāʾī ke akthar az nās min ḥaythu lā yashʿur dar ghamarāt-e
baḥr-e jahl va eʿrāż mostaghraq shavand va az folk-e elāhī maḥrūm mānand
ellā ʿedde-ye maʿdūdāt. Dīgar tā che zamān aryāḥ-e fażlīye be-vazīdan
āyad va bar hayākel-e mordegān-e nafs va havā morūr namāyad va īshānrā
ḥayāt bakhshad va be-shaṭr-e aḥadīye motavajjeh sāzad.

In short, today no one is accounted worthy of mention before God except
those who ride in this Ark of eternal glory. These souls are rarer than
red sulfur. Soon you will see most people, without realizing it,
submerged in the depths of the ocean of ignorance and rejection, and
deprived of the divine Ark, except for a numbered few. It remains to be
seen when the winds of grace will blow, pass over the bodies of those
dead in self and desire, give them life, and turn them toward the realm
of oneness.

\subsection{441}\label{section-936}

\begin{otherlanguage}{persian}

و بعد كه ميفرمايد حمل مينمايد عرش پروردگار را ملائكه عما در انفس ثمانيّه
همان برهانى است واضح كه اهل بيان موقن باين ظهور نشوند الّا قليل چنانچه
بثمانيه ذكر فرمود چه كه در اوّل هر ظهور چنين بوده تخصيص باين ظهور نداشته
و ندارد و مقصود از عرش در اينمقام امر اللّه بوده و خواهد بود و كلمة اللّه
بوده و خواهد بود و اوّل ما يتفوّه به لسان اللّه بوده و خواهد بود و آن
كلمه‌يست كه بفزع آيند از او كلّ من فى السّموات و الأرض و منصعق شوند كلّ من
فى جبروت الأمر و الخلق و آن كلمه‌ايست كه بثقل اكبر ناميده شده در الواح
اللّه و هر نفسى قابل حمل آن نه و آن كلمه‌ايست كه بعد از ظهور آن كلّ قبايل
بنوحه آيند و اركان كلّ شیء متزلزل شود و احدى قابل حمل اين عرش نه مگر
ملائكه عما كه باعانة اللّه اين امانت اكبر را حمل نمايند و بقبول آن هياكل
امر اللّه ظاهر شود و كلمه ابدع اعظم تنطّق فرمايد

\end{otherlanguage}

Va baʿd ke mīfarmāyad ḥaml mīnamāyad ʿarsh-e Parvardegār rā malāʾeke-ye
ʿamā dar anfos-e thamānīye, hamān borhānī ast vāżeḥ ke ahl-e Bayān mūqen
be-īn ẓohūr nashavand ellā qalīl, chenānche be-thamānīye dhekr farmūd.
Che ke dar avval-e har ẓohūr chenīn būde, takhṣīṣ be-īn ẓohūr nadāshte
va nadārad. Va maqṣūd az ʿarsh dar īn maqām Amr Allāh būde va khvāhad
būd va Kalemat Allāh būde va khvāhad būd va avval mā yatafawwahu bihi
Lisān Allāh būde va khvāhad būd. Va ān kalame-īst ke be-fazaʿ āyand az ū
kullu man fī al-samāwāt wa-al-arḍ va monṣaʿeq shavand kullu man fī
jabarūt al-amr wa-al-khalq. Va ān kalame-īst ke be-theql-e akbar nāmīde
shode dar alvāḥ-e Allāh va har nafsī qābel-e ḥaml-e ān na. Va ān
kalame-īst ke baʿd az ẓohūr-e ān koll-e qabāʾel be-nowḥe āyand va
arkān-e koll-e shayʾ motazalzel shavad va aḥadī qābel-e ḥaml-e īn ʿarsh
na magar malāʾeke-ye ʿamā ke be-eʿānat-e Allāh īn amānat-e akbar rā ḥaml
namāyand va be-qabūl-e ān hayākel-e Amr Allāh ẓāher shavad va kalame-ye
abdaʿ-e aʿẓam tanaṭṭoq farmāyad.

Then, where he says that the angels of the cloud, in eight souls, bear
the throne of the Lord, this is itself clear proof that only a few of
the people of the Bayān will become certain of this Revelation, as he
indicated by mentioning eight. For such has been the case at the
beginning of every Revelation; it has not been, nor is it, peculiar to
this Revelation. The throne in this passage has meant, and will mean,
the Cause of God, the Word of God, and the first utterance of the Tongue
of God. It is that Word which strikes terror into all who are in the
heavens and on earth, and causes all who are in the dominion of command
and creation to fall thunderstruck. It is that Word which has been named
the Most Great Weight in the Tablets of God, and not every soul is
capable of bearing it. It is that Word at whose appearance all kindreds
lament and the foundations of all things are shaken. No one is capable
of bearing this throne except the angels of the cloud, who, with God's
assistance, bear this most great Trust. Through accepting it, the
embodiments of the Cause of God appear, and the most wondrous, most
great Word speaks.

\subsection{442}\label{section-937}

\begin{otherlanguage}{persian}

و در يكمقام عرش نفس ظهور است كه ما بين ناس مشهود چه كه آن آيه توحيد و
جوهر تفريد در مقرّ خود واحد بوده و احدى با او نه كان اللّه و لم يكن معه من
شیء در آنمقام عرش نفس ظهور است و ملائكه ثمانيه مدلّين بر اين اسم مبارك
ابدع امنع چنانچه عددش با عدد بهآء من غير همزه مطابق و در حقيقت اوليّه
حامل آن عرش جز نفس مقدّسش نه

\end{otherlanguage}

Va dar yek maqām ʿarsh nafs-e ẓohūr ast ke mā beyn-e nās mashhūd, che ke
ān āye-ye towḥīd va jowhar-e tafrīd dar maqarr-e khod vāḥed būde va
aḥadī bā ū na: kāna Allāhu wa-lam yakun maʿahu min shayʾ. Dar ān maqām
ʿarsh nafs-e ẓohūr ast va malāʾeke-ye thamānīye modellīn bar īn esm-e
mobārak-e abdaʿ-e amnaʿ, chenānche ʿadadash bā ʿadad-e Bahāʾ min ghayr-e
hamze moṭābeq. Va dar ḥaqīqat-e avvalīye ḥāmel-e ān ʿarsh joz nafs-e
moqaddasash na.

In one station, the throne is the Manifestation himself, who is visible
among the people, for that sign of divine unity and essence of
singleness has ever been alone in his own seat, with no one beside him:
``God was, and there was nothing with him.'' In that station, the throne
is the Manifestation himself, and the eight angels point to this
blessed, most wondrous and most inaccessible Name, for its numerical
value corresponds to that of Bahāʾ without the hamza. In the primal
reality, none bears that throne except his own sanctified Self.

\subsection{443}\label{section-938}

\begin{otherlanguage}{arabic}

عرش او را اين زمان جز ذات او مى‌نيارد حمل اين از من شنو

\end{otherlanguage}

ʿArsh-e ū rā īn zamān joz dhāt-e ū mīnayārad ḥaml, īn az man shenow.

At this time none but his own Self can bear his throne---hear this from
me.

\subsection{444}\label{section-939}

\begin{otherlanguage}{persian}

و اين فرد از اشعارى است كه از لسان ابدع احلى در عراق جارى

\end{otherlanguage}

Va īn fard az ashʿārī ast ke az lesān-e abdaʿ-e aḥlā dar ʿErāq jārī.

This couplet is among the verses that flowed from the most wondrous,
most sweet Tongue in Iraq.

\subsection{445}\label{section-940}

\begin{otherlanguage}{persian}

و در مقام ديگر عرش اوّل من آمن است و اوّل من حضر و اوّل من بعث و اوّل من حمل
او است اوّل خلق و اوّل ناطق و اوّل بصير و اوّل سميع و اوّل عارف و اوّل موقن و
اوّل خبير و سمع او باب سمع الهى است از براى كلّ وجود و هر سمعى بآن باب
داخل نشد از اهل سمع نيست و عند اللّه اصمّ مذكور و هم‌‌چنين بصر او باب بصر
الهى است از براى كلّ نفوس و هم‌‌چنين علم او باب علم الهى است در ما بين بريّه
او و هم‌‌چنين قدرت او باب قدرت صمدانى است و لسان او باب لسان الهى و قلب او
باب قلب الهى و عرفان او باب عرفان الهى و بقول او ظاهر ميشود آنچه مستور
بوده و بطلب اوكوثر حيوان افاضه ميشود و بتوجّه او سحاب رحمت رحمانى اعطاى
فيض نمايد و باستماع او كلمه الهيّه از مخزن علم احديّه ظاهر شود اگر قدرى
بشعور آئى و بهمين فضل و عنايتى كه بذكر اوّل و كلمه اوّليّه و درّۀ اولى و ركن
قضا و مستسرّ عما عنايت شده بر عنايت و رحمت و كرم و فضل حقّ فى‌الجمله مطّلع
ميشوى

\end{otherlanguage}

Va dar maqām-e dīgar ʿarsh avval man āmana ast va avval man ḥaḍara va
avval man buʿitha va avval man ḥamala. Ū ast avval-e khalq va avval-e
nāṭeq va avval-e baṣīr va avval-e samīʿ va avval-e ʿāref va avval-e
mūqen va avval-e khabīr. Va samʿ-e ū bāb-e samʿ-e elāhī ast az barā-ye
koll-e vojūd va har samʿī be-ān bāb dākhel nashod az ahl-e samʿ nīst va
ʿinda Allāh aṣamm madhkūr. Va hamchenīn baṣar-e ū bāb-e baṣar-e elāhī
ast az barā-ye koll-e nofūs va hamchenīn ʿelm-e ū bāb-e ʿelm-e elāhī ast
dar mā beyn-e barīye-ye ū va hamchenīn qodrat-e ū bāb-e qodrat-e
ṣamadānī ast va lesān-e ū bāb-e lesān-e elāhī va qalb-e ū bāb-e qalb-e
elāhī va ʿerfān-e ū bāb-e ʿerfān-e elāhī. Va be-qowl-e ū ẓāher mīshavad
ānche mastūr būde va be-ṭalab-e ū Kowthar-e ḥayavān efāże mīshavad va
be-tavajjoh-e ū saḥāb-e raḥmat-e raḥmānī eʿṭā-ye feyż namāyad va
be-estemāʿ-e ū kalame-ye elāhīye az makhzan-e ʿelm-e aḥadīye ẓāher
shavad. Agar qadrī be-shoʿūr āʾī va be-hamīn fażl va ʿenāyatī ke
be-Dhekr-e Avval va Kalame-ye Avvalīye va Dorre-ye Ūlā va Rokn-e Qażā va
Mostaserr-e ʿAmā ʿenāyat shode, bar ʿenāyat va raḥmat va karam va fażl-e
Ḥaqq fī al-jomle moṭṭaleʿ mīshavī.

In another station, the throne is the first to believe, the first to
enter the divine presence, the first to be raised up, and the first to
bear the burden. He is the first of creation, the first to speak, the
first to see, the first to hear, the first to recognize, the first to be
certain, and the first to be informed. His hearing is the gate of divine
hearing for all existence, and any hearing that has not entered that
gate is not numbered among those that hear, but is accounted deaf before
God. Likewise, his sight is the gate of divine sight for all souls; his
knowledge is the gate of divine knowledge among God's creatures; his
power is the gate of eternal power; his tongue is the gate of the divine
tongue; his heart is the gate of the divine heart; and his recognition
is the gate of divine recognition. Through his speech, what was
concealed is revealed; at his request, the Kawthar of life is poured
forth; through his turning toward God, the cloud of the All-Merciful's
mercy bestows its grace; and through his listening, the divine Word
emerges from the treasury of the knowledge of oneness. If you attain a
measure of understanding, through this very grace and favor bestowed
upon the First Remembrance, the Primal Word, the First Pearl, the Pillar
of Decree, and the Hidden Mystery of the Cloud, you will gain some
awareness of the favor, mercy, generosity, and grace of the True One.

\subsection{446}\label{section-941}

\begin{otherlanguage}{persian}

بارى كور را با بصير مشابهتى نبوده و كر را با سميع مناسبتى نه در اين آيه
منزله معانى لانهايه مستور و لكن تا تشنه سلسبيل معانى يافت نشود ساقى
احديّه كاؤس معنويّه باقيه بذل ننمايد و ملاحظه كن كه بچه تصريح احوال معرضين
و مجرمين در كتاب الهى نازل شده و معنى آيه منزله مباركه مشهود گشته و آن
مجرمين از اهل بيان بوده و خواهند بود و آنچه ضرّ بر جمال رحمن وارد شود از
اهل بيان وارد خواهد شد چنانچه بنقطه بيان وارد نشد مگر از اهل فرقان و لكن
فو‌اللّه الّذى لا اله الّاهو كه بر اين ساذج قدم در كلّ حين وارد شده آنچه بر
احدى وارد نشده و اگر اهل بيان بصير بودند همان آيه مباركه اننّى انا حىّ فى
الأفق الأبهی كلّ را كفايت مينمود

\end{otherlanguage}

Bārī, kūr rā bā baṣīr moshābehatī nabūde va kar rā bā samīʿ monāsebatī
na. Dar īn āye-ye monazzale maʿānī-ye lā-nehāye mastūr va lāken tā
teshne-ye Salsabīl-e maʿānī yāft nashavad sāqī-ye aḥadīye koʾūs-e
maʿnavīye-ye bāqīye badhl nanamāyad. Va molāḥeẓe kon ke be-che taṣrīḥ
aḥvāl-e moʿreżīn va mojremīn dar Ketāb-e elāhī nāzel shode va maʿnī-ye
āye-ye monazzale-ye mobārake mashhūd gashte. Va ān mojremīn az ahl-e
Bayān būde va khvāhand būd va ānche żorr bar Jamāl-e Raḥmān vāred shavad
az ahl-e Bayān vāred khvāhad shod, chenānche be-Noqṭe-ye Bayān vāred
nashod magar az ahl-e Forqān. Va lāken fa-wa-Allāhi alladhī lā ilāha
illā Hu ke bar īn Sādhej-e Qedam dar koll-e ḥīn vāred shode ānche bar
aḥadī vāred nashode. Va agar ahl-e Bayān baṣīr būdand hamān āye-ye
mobārake innanī anā ḥayyun fī al-ufuq al-abhā koll rā kefāyat mīnamūd.

In short, the blind bear no resemblance to those who see, nor do the
deaf have any affinity with those who hear. Endless meanings lie
concealed in this revealed verse, but until someone is found who thirsts
for the Salsabīl of meanings, the Cupbearer of oneness will not offer
the everlasting cups of spiritual meaning. Consider how explicitly the
condition of those who turn away and those who are guilty has been
revealed in the Book of God, and how the meaning of the blessed revealed
verse has become manifest. Those guilty ones have been, and will be,
among the people of the Bayān, and whatever harm befalls the Beauty of
the All-Merciful will come from the people of the Bayān, just as what
befell the Point of the Bayān came only from the people of the Qurʾān.
Yet, by God, besides whom there is no other god, at every moment this
Essence of Eternity has suffered what no one else has suffered. Had the
people of the Bayān possessed sight, that blessed verse alone, ``I am
indeed alive in the All-Glorious Horizon,'' would have sufficed them
all.

\subsection{447}\label{section-942}

\begin{otherlanguage}{arabic}

و هم‌‌چنين كلّ من فى السّموات و الأرض را روحى لنفسه الفدآء ما قصّر فى نصحه
اهل البيان و تبليغه ايّاهم در هيچ عصر و عهدى ظهور قبل باين تفصيل و تصريح
ظهور بعد را ذكر ننموده چنانچه جميع شاهد و گواهيد لو انتم تنصفون بشنو ذكر
اين عبد را و از ظلمات فوق ظلمات خود را نجات ده و از براى دو روزه فانيه
خود را مستحقّ عذاب لانهايه منما و در دين اللّه و امره بتحريف مشغول مشو

\end{otherlanguage}

Va hamchenīn koll-e man fī al-samāwāt wa-al-arḍ rā, rūḥī li-nafsihi
al-fidāʾ, mā qaṣṣara fī nuṣḥihi ahla al-Bayān wa-tablīghihi iyyāhum. Dar
hīch ʿaṣr va ʿahdī ẓohūr-e qabl be-īn tafṣīl va taṣrīḥ ẓohūr-e baʿd rā
dhekr nanamūde, chenānche jamīʿ shāhed va govāhīd, law antum tunṣifūn.
Beshnow dhekr-e īn ʿabd rā va az ẓolomāt-e fowq-e ẓolomāt khod rā nejāt
deh va az barā-ye do rūze-ye fānīye khod rā mostaḥaqq-e ʿadhāb-e
lā-nehāye manamā va dar dīn Allāh va Amreh be-taḥrīf mashghūl mashow.

Likewise, may my soul and all who are in the heavens and on earth be
sacrificed for him: he left nothing undone in counseling the people of
the Bayān and conveying his message to them. In no age or era has the
preceding Manifestation spoken of the subsequent Manifestation with such
detail and explicitness, as you all witness and attest, if you judge
fairly. Listen to the words of this servant and rescue yourself from
darkness upon darkness. Do not, for the sake of two fleeting days, make
yourself deserving of endless torment, and do not occupy yourself with
distorting the religion of God and his Cause.

\subsection{448}\label{section-943}

\begin{otherlanguage}{persian}

نيكو است اين فرد در اينمقام گرد دين بهر صلاح تن به بيدينى متن تخم دنيا
در فضاى دين بمكّارى مكار

\end{otherlanguage}

Nīkū ast īn fard dar īn maqām: gerd-e dīn bahr-e ṣalāḥ-e tan be-bīdīnī
matan; tokhm-e donyā dar fażā-ye dīn be-makkārī makār.

This couplet is fitting here: ``Do not, for the body's welfare, weave
irreligion around religion; do not deceitfully sow the seed of the world
in the field of faith.''

\subsection{449}\label{section-944}

\begin{otherlanguage}{persian}

جميع آيات منزله را منتهى باين كلمات ابدع امنع اقدس اعلى كه از سماء عزّ
نقطه اولى نازل شده مينمايم كه شايد نفحات قدسش قاصدان حرم رحمن و عاشقان
جمال سبحان را معطّر نمايد اگر چه امثال آن نفوس نفحاتش را نيابند ناظراً الى
نفوس المشتاقين ذكر ميشود قسم بمليك آفاق‌ كه هر سمعى استماع نمايد بحيات
بديعه فايز شود و از اصغاى كلمات ما سوى اللّه منقطع گردد و هر قلبی‌ ادراك
نمايد البتّه منبع علوم نامتناهيه الهيّه شود و اصرح از آن ممكن نه تا ذكر
گردد و ابين از آن از ملكوت لسان تا حال ظاهر نشده

\end{otherlanguage}

Jamīʿ-e āyāt-e monazzale rā montahī be-īn kalemāt-e abdaʿ-e amnaʿ-e
aqdas-e aʿlā ke az samā-ye ʿezz-e Noqṭe-ye Ūlā nāzel shode mīnamāyam, ke
shāyad nafaḥāt-e qodsash qāṣedān-e ḥaram-e Raḥmān va ʿāsheqān-e Jamāl-e
Sobḥān rā moʿaṭṭar namāyad. Agarche amthāl-e ān nofūs nafaḥātash rā
nayāband, nāẓiran ilā nufūs al-mushtāqīn dhekr mīshavad. Qasam
be-Malīk-e āfāq ke har samʿī estemāʿ namāyad be-ḥayāt-e badīʿe fāʾez
shavad va az eṣghā-ye kalemāt-e mā sevā Allāh monqaṭeʿ gardad va har
qalbī edrāk namāyad albatte manbaʿ-e ʿolūm-e nāmotanāhīye-ye elāhīye
shavad. Va aṣraḥ az ān momken na tā dhekr gardad va abyan az ān az
malakūt-e lesān tā ḥāl ẓāher nashode.

I bring all the revealed verses to a close with these most wondrous,
most inaccessible, most holy and most exalted words, which have
descended from the heaven of the glory of the Primal Point, so that
their holy fragrances may perfume those who seek the sanctuary of the
All-Merciful and those who love the Beauty of the All-Glorious. Although
souls such as these may not perceive their fragrances, they are
mentioned with regard for the souls who yearn. I swear by the King of
all horizons that every ear that hears them will attain a wondrous life
and cease listening to the words of anyone but God, and every heart that
understands them will assuredly become a wellspring of God's boundless
knowledge. Nothing more explicit is possible to mention, and nothing
clearer has yet appeared from the kingdom of utterance.

\subsection{450}\label{section-945}

\begin{otherlanguage}{persian}

و چه قدر مشابه است اين ذكر و ثناى نقطه اولى با حكايتى كه از شيخ احمد
مرفوع عليه بهآء اللّه مذكور است چنانچه در سفرى كه بحجّ تشريف برده بحاجى
محمّد اسمعيل نامى ساكن در ارض صاد فرموده بودند سلام مرا بحضرت باب برسان و
حضرت را باسم مبارك ذكر فرموده بودند و مقصود جناب شيخ جز هدايت آن نفس و
نفوس نبوده آن شخص بعد از كلام حضرت شيخ متحيّر شده كه مقصود كيست و بچه جهت
اين كلمه طيّبه از لسان مبارك ظاهر شده تا آنكه از سفر حجّ مراجعت نموده و
سنين معدوده منقضى شده و حضرت اعلى بارض صاد وارد شدند و ذكرشان مابين ناس
مرتفع كه شخصى پيدا شده و ادّعاى بابيّت نموده آن وقت آن شخص متنبّه شده كه
مقصود شيخ چه بوده

\end{otherlanguage}

Va che qadr moshābeh ast īn dhekr va thanā-ye Noqṭe-ye Ūlā bā ḥekāyatī
ke az Shaykh Aḥmad-e marfūʿ, ʿalayhi Bahāʾ Allāh, madhkūr ast. Chenānche
dar safarī ke be-ḥajj tashrīf borde be-Ḥājī Moḥammad Esmāʿīl nāmī sāken
dar arż-e Ṣād farmūde būdand salām-e marā be-Ḥażrat-e Bāb berasān va
Ḥażrat rā be-esm-e mobārak dhekr farmūde būdand. Va maqṣūd-e Janāb-e
Shaykh joz hedāyat-e ān nafs va nofūs nabūde. Ān shakhs baʿd az kalām-e
Ḥażrat-e Shaykh motaḥayyer shode ke maqṣūd kīst va be-che jahat īn
kalame-ye ṭayyebe az lesān-e mobārak ẓāher shode, tā ānke az safar-e
ḥajj morājaʿat namūde va senīn-e maʿdūde monqażī shode va Ḥażrat-e Aʿlā
be-arż-e Ṣād vāred shodand va dhekrashān mā beyn-e nās mortafeʿ ke
shakhsī peydā shode va eddeʿā-ye bābīyat namūde. Ān vaqt ān shakhs
motanabbeh shode ke maqṣūd-e Shaykh che būde.

How closely this mention and praise by the Primal Point resembles the
account related of the departed Shaykh Aḥmad---upon him be the glory of
God. On a journey to perform the pilgrimage, he said to a certain Ḥājī
Muḥammad Ismāʿīl, who lived in the land of Ṣād, ``Convey my greetings to
the Bāb,'' mentioning him by his blessed name. The Shaykh's purpose was
nothing but to guide that soul and other souls. After hearing the
Shaykh's words, that man was bewildered as to who was intended and why
these good words had issued from his blessed tongue. He returned from
the pilgrimage, and some years passed. When the Most Exalted One arrived
in the land of Ṣād and talk of him spread among the people---that
someone had appeared and claimed to be the Bāb---the man then realized
what the Shaykh had meant.

\subsection{451}\label{section-946}

\begin{otherlanguage}{persian}

و لكن امثال شما از اين بيانات واضحه لايحه مشرقه مضيئه متنبّه نشده و
نخواهيد شد قوله عزّ اعزازه انتم قدام طائفة تظهر فيها النّقطة لا تقدمون ان
هم كانوا مؤمنين قل اولئك خير من على الأرض و لو علم اللّه خيرا منهم فى
الايمان ليظهره منهم انتم الى ابيه و امّه و ما كان معه و من آمن به من اولى
قرابته من اللّه تسلمون ان انتم تحسنن بكلّ نفس لعلّكم تدركون هذا قبل ان يظهر
و بعد ذلك انتم ستدركون و تعلمون عليك ان يا بهآء اللّه ثمّ اولى قرابتك ذكر
اللّه و ثنآء كلّ شیء فى كلّ حين و قبل حين و بعد حين انتهى

\end{otherlanguage}

Va lāken amthāl-e shomā az īn bayānāt-e vāżeḥe-ye lāʾeḥe-ye moshreqe-ye
możīʾe motanabbeh nashode va nakhvāhīd shod. Qawluhu ʿazza iʿzāzuhu:
antum quddāma ṭāʾifatin taẓharu fīhā al-Nuqṭatu lā taqaddamūn in hum
kānū muʾminīn. Qul ulāʾika khayru man ʿalā al-arḍ wa-law ʿalima Allāhu
khayran minhum fī al-īmān li-yuẓhirahu minhum. Antum ilā abīhi wa-ummihi
wa-mā kāna maʿahu wa-man āmana bihi min ulī qarābatihi min Allāhi
tusallimūn. In antum tuḥsinunna bi-kulli nafsin laʿallakum tudrikūna
hādhā qabla an yaẓhara wa-baʿda dhālika antum satudrikūna wa-taʿlamūn.
ʿAlayka an yā Bahāʾ Allāh thumma ulī qarābatika dhikru Allāhi wa-thanāʾu
kulli shayʾin fī kulli ḥīnin wa-qabla ḥīnin wa-baʿda ḥīn. Intahā.

Yet people such as you have not been, and will not be, awakened by these
clear, manifest, radiant and luminous statements. His words---exalted be
his glory: ``Do not put yourselves ahead of the people among whom the
Point will appear, if they are believers. Say: They are the best of
those on earth, and had God known any better than them in faith, he
would have caused him to appear from among them. Convey greetings from
God to his father and his mother, to whatever was with him, and to those
of his kindred who believe in him. If you show kindness to every soul,
perhaps you will grasp this before he appears; afterward you will indeed
grasp it and know. Upon you, O Bahāʾuʾllāh, and then upon your kindred,
be the remembrance of God and the praise of all things at every moment,
before every moment and after every moment.'' End of quotation.

\subsection{452}\label{section-947}

\begin{otherlanguage}{persian}

و مقصود نقطه بيان از اين آيات آنكه اراده فرموده‌اند كه اصل ظهور را معيّن
فرمايند حال اى بی‌انصاف ملاحظه كن كه بچه ادب امر فرموده‌اند باهل بيان
ميفرمايند شما قدام طايفه كه آن نقطه اوليّه از آنها ظاهر ميشود تقدّم نجوئيد
ميفرمايد اگر از مؤمنين باشند و باين بيان آخر حبل انتساب منتسبين اين ظهور
كه اليوم از حقّ معرضند منقطع چنانچه تصريحاً فرموده و ميفرمايند آن طايفه
بهترين طوايف روى ارضند چه كه اگر طايفه در ارض از آن طايفه بهتر بود
هرآينه آن طلعت احديّه و كينونت الهيّه از آن طايفه ظاهر ميشد و امر
ميفرمايند كه اهل بيان باب و امّ آن نيّر اعظم و هم‌‌چنين ذوى قرابه او از
كسانيكه مؤمنند باو من قبل اللّه سلام برسانند

\end{otherlanguage}

Va maqṣūd-e Noqṭe-ye Bayān az īn āyāt ānke erāde farmūde-and ke aṣl-e
ẓohūr rā moʿayyan farmāyand. Ḥāl ey bī-enṣāf, molāḥeẓe kon ke be-che
adab amr farmūde-and. Be-ahl-e Bayān mīfarmāyand shomā qoddām-e ṭāyefe
ke ān Noqṭe-ye Avvalīye az ānhā ẓāher mīshavad taqaddom najūʾīd.
Mīfarmāyad agar az moʾmenīn bāshand. Va be-īn bayān-e ākhar ḥabl-e
entesāb-e montasebīn-e īn ẓohūr ke al-yowm az Ḥaqq moʿreżand monqaṭeʿ,
chenānche taṣrīḥan farmūde. Va mīfarmāyand ān ṭāyefe behtarīn-e
ṭavāyef-e rū-ye arżand, che ke agar ṭāyefe dar arż az ān ṭāyefe behtar
būd har-āyene ān Ṭalʿat-e Aḥadīye va Keynūnat-e Elāhīye az ān ṭāyefe
ẓāher mīshod. Va amr mīfarmāyand ke ahl-e Bayān be-ab va omm-e ān
Nayyer-e Aʿẓam va hamchenīn dhavī-ye qarābe-ye ū az kasānī ke moʾmenand
be-ū min qibal Allāh salām berasānand.

The Point of the Bayān's purpose in these verses was to identify the
origin of the Manifestation. Now, you who are without fairness, consider
what courtesy he enjoined. He tells the people of the Bayān not to put
themselves ahead of the people from whom that Primal Point will appear.
He says, ``If they are believers.'' By this final statement, the bond of
kinship of those related to this Manifestation who today turn away from
the True One is severed, as he has explicitly stated. He says that those
people are the best of the peoples on earth, for had there been a better
people on earth, that Countenance of oneness and divine Being would
certainly have appeared from among them. He commands the people of the
Bayān to convey greetings from God to the father and mother of that Most
Great Luminary, and likewise to those of his kindred who believe in him.

\subsection{453}\label{section-948}

\begin{otherlanguage}{persian}

از احمقى ملأ بيان بعيد نيست كه بگويند اب و امّى كه حال موجود نه سلام الهى
بچه محلّ راجع ميشود اى بی‌‌بصران همان حين كه از لسان ربّ العالمين سلام نازل
بمقرّش وارد و اين تاكيد در بيان نظر بآن بوده كه امرى در ظهور واقع نشود كه
آن ساذج قدم را محزون نمايد و بعد از اين وصاياى الهيّه و آداب مذكوره و
كلمات ممتنعه منيعه نفس اقدس اطهر متوجّهاً الى مقرّ الظّهور والى ذوى قرابته
ذكر و ثنا فرموده‌اند فو‌اللّه بملاحتى اين كوثر بيان از لسان رحمن جارى شده
كه اگر بقطره آن كلّ من فى السّموات و الأرض خالصاً لوجه اللّه فايز شوند جميع
از آن خمر الهيّه سرمست شده بشطر قدم توجّه نمايند و ما سوايش را معدوم و
مفقود و لاشیء محض مشاهده كنند

\end{otherlanguage}

Az aḥmaqī-yi malaʾ-i Bayān baʿīd nīst kih bigūyand ab va ummī kih ḥāl
mawjūd nah, salām-i ilāhī bih chih maḥall rājiʿ mīshavad. Ay bī-baṣarān,
hamān ḥīn kih az lisān-i Rabb al-ʿālamīn salām nāzil, bih maqarraš
vārid. Va īn taʾkīd dar Bayān naẓar bih ān būdih kih amrī dar ẓuhūr
vāqiʿ nashavad kih ān sādhij-i qidam rā maḥzūn namāyad. Va baʿd az īn
vaṣāyā-yi ilāhiyyih va ādāb-i madhkūrih va kalimāt-i mumtaniʿih-yi
manīʿih, nafs-i aqdas-i aṭhar mutawajjihan ilā maqarri al-ẓuhūr wa-ilā
dhawī qarābatihi dhikr va thanā farmūdih-and. Fa-wa-llāh, bih malāḥatī
īn kawthar-i bayān az lisān-i Raḥmān jārī shudih kih agar bih qaṭrih-yi
ān kullu man fī al-samāwāti wa-al-arḍ khāliṣan li-wajhi Allāh fāʾiz
shavand, jamīʿ az ān khamr-i ilāhiyyih sarmast shudih, bih shaṭr-i qidam
tavajjuh namāyand va mā sivāyash rā maʿdūm va mafqūd va lā-shayʾ-i maḥḍ
mushāhidih kunand.

It would be no surprise, given the foolishness of the people of the
Bayān, if they were to say, ``To whom do the divine greetings return,
when the father and mother are no longer living?'' O sightless ones! The
very moment the greeting descended from the tongue of the Lord of the
worlds, it reached its destination. This emphasis in the Bayān was
intended to ensure that nothing should occur during the Manifestation to
sadden that Essence of Eternity. After these divine counsels, the
prescribed courtesies, and these exalted, unassailable words, that most
holy and pure Being offered remembrance and praise, turning toward the
seat of the Manifestation and his kindred. By God, this heavenly stream
of utterance has flowed from the tongue of the All-Merciful with such
sweetness that, if all who are in the heavens and on earth were to
partake of a drop of it purely for the sake of God, they would all
become intoxicated with that divine wine, turn toward the Eternal One,
and see everything other than him as nonexistent, lost, and utter
nothingness.

\subsection{454}\label{section-949}

\begin{otherlanguage}{persian}

بنصّ طلعت اعلى روح ما سواه فداه كلّ مأمور شده‌اند بثناء و سلام در حين ظهور
و لكن احدى بآن موفّق نشده و نصح الهى را لغو شمرده‌اند و وصيّتش را از قلب
محو نموده بلى اگر بثناء و سلام در بين يدى مليك علّام حاضر نشده‌اند و لكن
بسنان و سهام از كلّ اطراف هجوم آورده‌اند و للّه الأمر من قبل و من بعد و
عليه يثنى نفسه وذاته لو انتم من العارفين و چون احدى باين ثناء موفّق نشد
كه من قبل الرّحمن ابلاغ دارد لذا كينونت سبحانى در اين ظهور بدع رحمانى خود
من قبل اللّه بر نفس خود ابلاغ ذكر و ثناء فرموده‌اند قوله جلّ ثنائه ان يا
محبوب البهآء بلغت ذكرك و ثنائك ثمّ ثنآء كلّ شیء من قبلك على نفسى ليثبت
قولك فى العالمين و يظهر آياتك للعارفين و اذا اقول عليك يا محبوب البهآء
ذكر اللّه و ثنائه ثمّ ثنآء اهل ملأ الأعلی و ثنآء اهل مداين البقآء ثمّ ثنآء
اهل لجج الكبريا ثمّ ثنآء العالمين فى كلّ حين و قبل حين و بعد حين و حين حين

\end{otherlanguage}

Bih naṣṣ-i Ṭalʿat-i Aʿlā, rūḥ-i mā sivāh fidāh, kull maʾmūr shudih-and
bih thanā va salām dar ḥīn-i ẓuhūr, va lākin aḥadī bih ān muwaffaq
nashudih va nuṣḥ-i ilāhī rā laghv shumurdih-and va vaṣiyyatash rā az
qalb maḥv namūdih. Balī, agar bih thanā va salām dar bayn-i yaday-i
Malīk-i ʿAllām ḥāḍir nashudih-and, va lākin bih sinān va sihām az kull-i
aṭrāf hujūm āvardih-and. Wa-li-llāhi al-amru min qablu wa-min baʿdu
wa-ʿalayhi yuthnī nafsuhu wa-dhātuhu law antum mina al-ʿārifīn. Va chūn
aḥadī bih īn thanā muwaffaq nashud kih min qibali al-Raḥmān iblāgh
dārad, lidhā kaynūnat-i subḥānī dar īn ẓuhūr-i badʿ-i raḥmānī khud min
qibali Allāh bar nafs-i khud iblāgh-i dhikr va thanā farmūdih-and.
Qawluhu jalla thanāʾuhu: An yā maḥbūba al-Bahāʾ, ballaghtu dhikraka
wa-thanāʾaka thumma thanāʾa kulli shayʾin min qibalika ʿalā nafsī
li-yathbuta qawluka fī al-ʿālamīn wa-yaẓhara āyātuka li-l-ʿārifīn.
Wa-idhan aqūlu ʿalayka yā maḥbūba al-Bahāʾ dhikru Allāhi wa-thanāʾuhu
thumma thanāʾu ahli malaʾi al-aʿlā wa-thanāʾu ahli madāʾini al-baqāʾ
thumma thanāʾu ahli lujaji al-kibriyāʾ thumma thanāʾu al-ʿālamīn fī
kulli ḥīnin wa-qabla ḥīnin wa-baʿda ḥīnin wa-ḥīna ḥīn.

By the explicit text of the Most Exalted Countenance---may the soul of
all else be sacrificed for him---all were commanded to offer praise and
greetings at the time of the Manifestation. Yet no one has succeeded in
doing so; they have regarded the divine counsel as futile and erased his
testament from their hearts. Indeed, although they have not entered the
presence of the all-knowing Sovereign with praise and greetings, they
have assailed him from every side with spears and arrows. Authority
belongs to God, before and after, and upon him his own Self and Essence
bestow praise, if you are among those who understand. Since no one
succeeded in conveying this praise on behalf of the All-Merciful, the
divine Being has himself, in this wondrous, merciful Manifestation,
conveyed remembrance and praise to his own Self on behalf of God. His
words---exalted be his praise: ``O Beloved of Bahāʾ! I have conveyed
your remembrance and praise, and then the praise of all things, from you
to myself, so that your word may be established among the worlds and
your signs made manifest to those who understand. And now I say: Upon
you, O Beloved of Bahāʾ, be the remembrance of God and his praise, then
the praise of the dwellers in the highest Assembly, and the praise of
the dwellers in the cities of eternity, then the praise of the dwellers
in the oceans of grandeur, then the praise of all the worlds, at every
moment, before every moment, after every moment, and at the very moment
of each moment.''

\subsection{455}\label{section-950}

\begin{otherlanguage}{arabic}

در اينمقام از جبروت ابهى اين كلمات ابدع احلى فى المناجات مع اللّه العلىّ
الأعلی نازل عليك يا بهآء اللّه محبوب البهآء ذكر اللّه و بهآئه ثمّ بهآء اهل
ملأ الأعلی ثمّ بهآء اهل مداين البقآء ثمّ بهآء كلّ الاشیآء ثمّ بهآء نفسك
لنفسك بنفسك و بهآء هذا البهآء الّذى ظهر لنصرك بين العالمين فيا محبوب
البهآء فو‌عزّتك و جلالك انّك ما قصّرت فى تبليغك عبادك و تدبيرك بريّتك و ما
اردت فى كلّ ذلك الّا خضوعهم بين يدى سلطان احديّتك و الخشوع عند ظهورات انوار
وجهتك فو‌عزّتك يا محبوبی‌ انّى اجد نفسى خجلاً عمّا بلغته لظهورى بحيث ما تنفّست
الّا بذكرى و ما تكلّمت الّا لاثبات امرى و ما جرى من قلمك الّا ما كان فيه
مقصود نفسك ذكرى و ثنائى و فى كلّ شأن ظهر منك ما يكون مدلاً لنفسى و صريحاً
فى ظهورى و حاكيا عن جمالى

\end{otherlanguage}

Dar īn maqām az jabarūt-i Abhā īn kalimāt-i abdaʿ-i aḥlā fī al-munājāt
maʿa Allāhi al-ʿAliyyi al-Aʿlā nāzil: ʿAlayka yā Bahāʾa Allāh, maḥbūba
al-Bahāʾ, dhikru Allāhi wa-bahāʾuhu thumma bahāʾu ahli malaʾi al-aʿlā
thumma bahāʾu ahli madāʾini al-baqāʾ thumma bahāʾu kulli al-ashyāʾ
thumma bahāʾu nafsika li-nafsika bi-nafsika wa-bahāʾu hādhā al-Bahāʾ
alladhī ẓahara li-naṣrika bayna al-ʿālamīn. Fa-yā maḥbūba al-Bahāʾ,
fa-wa-ʿizzatika wa-jalālika innaka mā qaṣṣarta fī tablīghika ʿibādaka
wa-tadbīrika bariyyataka wa-mā aradta fī kulli dhālika illā khuḍūʿahum
bayna yaday sulṭāni aḥadiyyatika wa-al-khushūʿa ʿinda ẓuhūrāti anwāri
wajhatika. Fa-wa-ʿizzatika yā maḥbūbī, innī ajidu nafsī khajilan ʿammā
ballaghtahu li-ẓuhūrī bi-ḥaythu mā tanaffasta illā bi-dhikrī wa-mā
takallamta illā li-ithbāti amrī wa-mā jarā min qalamika illā mā kāna
fīhi maqṣūdu nafsika dhikrī wa-thanāʾī wa-fī kulli shaʾnin ẓahara minka
mā yakūnu mudillan li-nafsī wa-ṣarīḥan fī ẓuhūrī wa-ḥākiyan ʿan jamālī.

At this point, these wondrous and most sweet words descended from the
all-glorious Dominion in prayer to God, the Exalted, the Most High:
``Upon you, O Glory of God, Beloved of Bahāʾ, be the remembrance of God
and his glory, then the glory of the dwellers in the highest Assembly,
then the glory of the dwellers in the cities of eternity, then the glory
of all things, then the glory of your own Self, for yourself, through
yourself, and the glory of this Bahāʾ who has appeared to aid you among
the worlds. O Beloved of Bahāʾ! By your might and majesty, you have
fallen short in neither conveying your message to your servants nor
ordering the affairs of your creation. In all this you have desired only
their submission before the sovereignty of your oneness and their
humility at the appearances of the lights of your countenance. By your
might, O my Beloved, I find myself abashed before all that you have
conveyed for my manifestation, for you have breathed only in remembrance
of me, spoken only to establish my Cause, and nothing has flowed from
your pen in which your purpose was not my remembrance and praise. In
every circumstance, there has appeared from you that which points to me,
speaks explicitly of my manifestation, and tells of my beauty.''

\subsection{456}\label{section-951}

\begin{otherlanguage}{arabic}

و مع‌ذلك كيف اذكرك يا محبوبی‌ بعد الّذى فو‌عزّتك لترادف القضايا و تتابع
البلايا لن اجد من فرصة لابكى على نفسى و كيف ثنآء نفسك العلىّ العظيم كلّما
اريد ان اثنيك بثنآء او اتقرّب اليك ببدايع ذكرك يمنعنى غلّ اعدائك و اعراض
طغاة بريّتك بحيث لو التفت الى اليمين اجد كتاب السّجّين من احد من خلقك و فيه
ما اضطربت عنه اركان كلّ شیء ثمّ اركان ملكوت اسمائك الحسنى و كلّما اتوجّه الى
اليسار اجد كتاباً من الفجّار و فيه ما صاحت و ضجّت و ناحت حقايق اصفيائك و
افئدة انبيائك بعد الّذى انّك امرت ملأ البيان الّذين يدّعون الايمان بنفسك و
الاقرار بفردانيّتك و الاذعان عند ظهورات انوار عزّ وحدانيّتك بأن يكتب كلّ نفس
فى كلّ واحد كتاباً فى اثبات امرك فى هذا الظّهور الّذى به اشرقت شمس العزّة و
الاقتدار عن افق سمآء عزّ رحمانيّتك ليتذكرن به العباد و ينتظرن هذا النّبأ
الّذى بشّرتهم فى كلّ الواحك و زبرك و صحايف مجدك بقولك الحقّ مخاطباً لملأ
البيان و انتم فى كلّ واحد كتاب اثبات لمن نظهره بعضكم الى بعض تكتبون لعلّكم
يوم ظهوره بما تكتبون لتعملون

\end{otherlanguage}

Wa-maʿa dhālika kayfa adhkuruka yā maḥbūbī baʿda alladhī fa-wa-ʿizzatika
li-tarādufi al-qaḍāyā wa-tatābuʿi al-balāyā lan ajida min furṣatin
li-abkiya ʿalā nafsī wa-kayfa thanāʾu nafsika al-ʿAliyyi al-ʿAẓīm.
Kullamā urīdu an uthniyaka bi-thanāʾin aw ataqarraba ilayka bi-badāʾiʿi
dhikrika yamnaʿunī ghillu aʿdāʾika wa-iʿrāḍu ṭughāti bariyyatika
bi-ḥaythu law iltafattu ilā al-yamīn ajidu kitāba al-sijjīn min aḥadin
min khalqika wa-fīhi mā iḍṭarabat ʿanhu arkānu kulli shayʾin thumma
arkānu malakūti asmāʾika al-ḥusnā. Wa-kullamā atawajjahu ilā al-yasār
ajidu kitāban mina al-fujjār wa-fīhi mā ṣāḥat wa-ḍajjat wa-nāḥat
ḥaqāʾiqu aṣfiyāʾika wa-afʾidatu anbiyāʾika baʿda alladhī innaka amarta
malaʾa al-Bayān alladhīna yaddaʿūna al-īmāna bi-nafsika wa-al-iqrāra
bi-fardāniyyatika wa-al-idhʿāna ʿinda ẓuhūrāti anwāri ʿizzi
waḥdāniyyatika bi-an yaktuba kullu nafsin fī kulli wāḥidin kitāban fī
ithbāti amrika fī hādhā al-ẓuhūr alladhī bihi ashraqat shamsu al-ʿizzati
wa-al-iqtidār ʿan ufuqi samāʾi ʿizzi raḥmāniyyatika li-yatadhakkaranna
bihi al-ʿibād wa-yantaẓiranna hādhā al-nabaʾ alladhī bashshartahum fī
kulli alwāḥika wa-zuburika wa-ṣaḥāʾifi majdika bi-qawlika al-ḥaqq
mukhāṭiban li-malaʾi al-Bayān: Wa-antum fī kulli wāḥidin kitāba ithbātin
li-man nuẓhiruhu baʿḍukum ilā baʿḍin taktubūn laʿallakum yawma ẓuhūrihi
bi-mā taktubūna la-taʿmalūn.

``Yet how can I remember you, O my Beloved, when, by your might, the
succession of adversities and the unceasing calamities leave me no
opportunity to weep for myself, much less to praise your exalted and
mighty Self? Whenever I wish to praise you or draw near to you through
the wonders of your remembrance, the malice of your enemies and the
aversion of the tyrants among your creation prevent me. If I turn to the
right, I find the book of Sijjīn from one of your creatures, containing
what has shaken the foundations of all things and the foundations of the
kingdom of your most beautiful names. Whenever I turn to the left, I
find a book from the wicked, containing what has caused the inmost
realities of your chosen ones and the hearts of your prophets to cry
out, lament, and mourn. This, even though you commanded the people of
the Bayān---who profess faith in you, acknowledge your singleness, and
claim submission at the appearances of the lights of your glorious
unity---that every soul should write, in every Vāḥid, a book
establishing your Cause in this Manifestation, through which the sun of
might and power has shone from the horizon of the heaven of your
glorious mercy, so that your servants might be reminded by it and await
this announcement of which you gave them glad tidings in all your
Tablets, your scriptures, and the scrolls of your glory, through your
true words addressed to the people of the Bayān: `In every Vāḥid, you
are to write to one another a book establishing the truth of the One
whom we shall manifest, so that perhaps, on the day of his appearance,
you may act according to what you write.'\,''

\subsection{457}\label{section-952}

\begin{otherlanguage}{arabic}

و انّهم يا الهى نبذوا احكامك عن ورائهم و نسوا ما نصحتهم به فى الواحك فيا
ليت كانوا قانعاً بذلك بل كتب كلّ واحد منهم كتاباً احترقت عنه نفسك و حقايق
كلّ شیء و اكباد المقرّبين من احبّائك و المقدّسين من اصفيائك و انتشروه فى
البلاد ليضيع به حرمة اللّه بين بريّته و عزّة اللّه بين خلقه و احترامه بين
عباده فو‌عزّتك يا محبوبی‌ ما ادركت صباحاً الّا و قد ارتدّ فيه بصرى الى ما
احترق به كبدى و اضطرب به نفسى و حزن به قلبی‌ و بذلك بكيت بعيون سرّى و بكت
ببكائى عيون الّذين كرّمتهم من بريّتك و اصطفيتهم من بين خلقك و جعلتهم مهابط
وحيك و مخازن علمك و مظاهر امرك و مطالع قدرتك و مكامن الهامك و مشارق شموس
احسانك

\end{otherlanguage}

Wa-innahum yā ilāhī nabadhū aḥkāmaka ʿan warāʾihim wa-nasū mā naṣaḥtahum
bihi fī alwāḥika. Fa-yā layta kānū qāniʿan bi-dhālika bal kataba kullu
wāḥidin minhum kitāban iḥtaraqat ʿanhu nafsuka wa-ḥaqāʾiqu kulli shayʾin
wa-akbādu al-muqarrabīna min aḥibbāʾika wa-al-muqaddasīna min aṣfiyāʾika
wa-intasharūhu fī al-bilād li-yaḍīʿa bihi ḥurmatu Allāhi bayna
bariyyatihi wa-ʿizzatu Allāhi bayna khalqihi wa-iḥtirāmuhu bayna
ʿibādihi. Fa-wa-ʿizzatika yā maḥbūbī mā adraktu ṣabāḥan illā wa-qad
irtadda fīhi baṣarī ilā mā iḥtaraqa bihi kabidī wa-iḍṭaraba bihi nafsī
wa-ḥazina bihi qalbī wa-bi-dhālika bakaytu bi-ʿuyūni sirrī wa-bakat
bi-bukāʾī ʿuyūnu alladhīna karramtahum min bariyyatika wa-iṣṭafaytahum
min bayni khalqika wa-jaʿaltahum mahābiṭa waḥyika wa-makhāzina ʿilmika
wa-maẓāhira amrika wa-maṭāliʿa qudratika wa-makāmina ilhāmika
wa-mashāriqa shumūsi iḥsānika.

``O my God, they have cast your ordinances behind them and forgotten the
counsels you gave them in your Tablets. If only they had been content
with that! Instead, each of them has written a book that has set ablaze
your own Self, the realities of all things, and the inmost hearts of
those among your loved ones who are near to you and those among your
chosen ones who are sanctified. They have spread it throughout the lands
to destroy the sanctity of God among his creatures, the majesty of God
among his creation, and reverence for him among his servants. By your
might, O my Beloved, no morning has come upon me without my gaze falling
upon something that has burned my inmost heart, agitated my soul, and
saddened my heart. Thus I have wept with the eyes of my innermost being,
and with my weeping have wept the eyes of those whom you have honored
among your creatures, chosen from among your creation, and made the
recipients of your revelation, the treasuries of your knowledge, the
manifestations of your Cause, the dawning places of your power, the
repositories of your inspiration, and the daysprings of the suns of your
bounty.''

\subsection{458}\label{section-953}

\begin{otherlanguage}{arabic}

و ما ادركت من مساء الّا و فیه قد ورد علىّ ما منعت به نسمات رحمانيّتك عن
الممكنات و غلقت به ابواب فضلك على وجوه الكاينات فوعزّتك يا محبوبی‌ صرت
متحيّراً فى امرى و اذا تشاهدنى كالحوت المتبلبل على التّراب و تسمع صريخ قلبی‌
يا من بيدك جبروت الآیات و كلّ ذلك ورد علىّ بعد الّذى دعوتهم الى شطر مواهبك
و الطافك و عرفتهم مناهج امرك و رضائك و امرتهم بالخضوع لدى باب رحمانيّتك و
الورود على فنآء فردانيّتك و كلّما ناديتهم يا الهى بما الهمتنى من بدايع
كلماتك و جواهر آياتك قاموا على الاعراض بشأن لايقدر احد ان يحصيه و انّك
انت احصيته بسلطانك و علمك

\end{otherlanguage}

Wa-mā adraktu min masāʾin illā wa-fīhi qad warada ʿalayya mā muniʿat
bihi nasamātu raḥmāniyyatika ʿani al-mumkināt wa-ghuliqat bihi abwābu
faḍlika ʿalā wujūhi al-kāʾināt. Fa-wa-ʿizzatika yā maḥbūbī ṣirtu
mutaḥayyiran fī amrī wa-idhan tushāhidunī ka-al-ḥūti al-mutabalbili ʿalā
al-turāb wa-tasmaʿu ṣarīkha qalbī yā man bi-yadika jabarūtu al-āyāt.
Wa-kullu dhālika warada ʿalayya baʿda alladhī daʿawtuhum ilā shaṭri
mawāhibika wa-alṭāfika wa-ʿarraftuhum manāhija amrika wa-riḍāʾika
wa-amartuhum bi-al-khuḍūʿi ladā bābi raḥmāniyyatika wa-al-wurūdi ʿalā
fināʾi fardāniyyatika. Wa-kullamā nādaytuhum yā ilāhī bi-mā alhamtanī
min badāʾiʿi kalimātika wa-jawāhiri āyātika qāmū ʿalā al-iʿrāḍi
bi-shaʾnin lā yaqdiru aḥadun an yuḥṣiyahu wa-innaka anta aḥṣaytahu
bi-sulṭānika wa-ʿilmika.

``No evening has come upon me without there befalling me something that
has withheld the breezes of your mercy from all contingent things and
closed the doors of your grace upon all created beings. By your might, O
my Beloved, I am bewildered by my plight. You see me now like a fish
thrashing upon the dust, and you hear the cry of my heart, O you in
whose hand is the dominion of signs. All this has befallen me after I
called them toward your gifts and favors, acquainted them with the paths
of your Cause and your good pleasure, and commanded them to bow humbly
at the door of your mercy and enter the court of your singleness.
Whenever I called to them, O my God, with the wondrous words and the
gems of your verses that you inspired in me, they rose up in opposition
to a degree that no one can reckon. You alone have reckoned it through
your sovereignty and knowledge.''

\subsection{459}\label{section-954}

\begin{otherlanguage}{arabic}

فاذا يا محبوب البهآء و محيى البهآء و مشوّق البهآء و ذكر البهآء و حبيب
البهآء فابك على نفس البهآء تاللّه انّها قد بقت فريداً بين خلقك و وحيداً بين
عبادك و يفعلون به ما يريدون و ليس لى من ناصر ليمنعهم عن فعلهم او يطردهم
عن حول حرم قدسك و سرادق عزّك و اجلالك فلك الحمد فى كلّ ذلك و فى كلّ ما ورد
على نفسى فى سبيلك ولو انّ عبادك نزعوا عن هيكلى ثوب السّرور و ردآء البهجة و
العزّة و الحرمة و لكن اعطيتنى بفضلك ما لايقدر ان يتصرّف فيه احد ولو يجتمع
علىّ خلق السّموات و الأرضين فلك الحمد على ما اعطيتنى بجودك يا محبوب قلبی‌ و
مقصود قلوب العارفين و انا الّذى بحبّك لن اجزع من شیء ولو يمطر علىّ سحاب
القضا سهام البلا فوضت امرى اليك و توكّلت عليك و انت حسبی‌ و معينى و ناصرى
و بك اكتفيت عن الخلايق اجمعين و الحمد لك اذ انّك انت معبودى و معبود من فى
العالمين انتهى

\end{otherlanguage}

Fa-idhan yā maḥbūba al-Bahāʾ wa-muḥyī al-Bahāʾ wa-mushawwiqa al-Bahāʾ
wa-dhikra al-Bahāʾ wa-ḥabība al-Bahāʾ fa-bki ʿalā nafsi al-Bahāʾ.
Ta-llāhi innahā qad baqat farīdan bayna khalqika wa-waḥīdan bayna
ʿibādika wa-yafʿalūna bihi mā yurīdūn wa-laysa lī min nāṣirin
li-yamnaʿahum ʿan fiʿlihim aw yaṭrudahum ʿan ḥawli ḥarami qudsika
wa-surādiqi ʿizzika wa-ijlālika. Fa-laka al-ḥamdu fī kulli dhālika wa-fī
kulli mā warada ʿalā nafsī fī sabīlika wa-law anna ʿibādaka nazaʿū ʿan
haykalī thawba al-surūr wa-ridāʾa al-bahjati wa-al-ʿizzati wa-al-ḥurmati
wa-lākin aʿṭaytanī bi-faḍlika mā lā yaqdiru an yataṣarrafa fīhi aḥadun
wa-law yajtamiʿu ʿalayya khalqu al-samāwāti wa-al-arḍīn. Fa-laka
al-ḥamdu ʿalā mā aʿṭaytanī bi-jūdika yā maḥbūba qalbī wa-maqṣūda qulūbi
al-ʿārifīn. Wa-anā alladhī bi-ḥubbika lan ajzaʿa min shayʾin wa-law
yumṭiru ʿalayya saḥābu al-qaḍā sahāma al-balā. Fawwaḍtu amrī ilayka
wa-tawakkaltu ʿalayka wa-anta ḥasbī wa-muʿīnī wa-nāṣirī wa-bika
iktafaytu ʿani al-khalāʾiqi ajmaʿīn. Wa-al-ḥamdu laka idh innaka anta
maʿbūdī wa-maʿbūdu man fī al-ʿālamīn. Intahā.

``Now, O Beloved of Bahāʾ, giver of life to Bahāʾ, awakener of Bahāʾ's
longing, remembrance of Bahāʾ, and dear one of Bahāʾ, weep for the
person of Bahāʾ! By God, he has remained alone among your creatures and
solitary among your servants. They do to him whatever they wish, and I
have no helper to restrain them from their deeds or drive them away from
your sacred sanctuary and the pavilion of your might and majesty. Praise
be to you in all this and in all that has befallen me in your path.
Although your servants have stripped from my body the garment of joy and
the robe of gladness, majesty, and honor, you have nevertheless bestowed
on me through your grace that which no one can lay hands upon, even
should all the creatures of the heavens and the earth unite against me.
Praise be to you for what you have bestowed on me through your
generosity, O Beloved of my heart and the desire of the hearts of those
who know you. Through my love for you, I will not be dismayed by
anything, even should the clouds of destiny rain down upon me the arrows
of affliction. I have committed my affairs to you and placed my trust in
you. You are sufficient for me, my helper and my supporter, and through
you I have become independent of all creatures. Praise be to you, for
you are the one I worship and the one worshipped by all who are in the
worlds.'' End of quotation.

\subsection{460}\label{section-955}

\begin{otherlanguage}{persian}

و ايكاش مهلت ميدادند كه جمال ربّ العالمين بثنآء اللّه و ذكره مشغول شود
فو‌اللّه فرصت آنكه بر نفس خود نوحه نمايد از ظلم مشركين نيافته تا چه رسد
بذكر و ثنا مع آنكه نقطه بيان كلّ اهل بيان را امر فرموده كه در هر واحد
بعضى ببعضى كتابی‌ مرقوم دارند و در آن كتاب يكديگر را ببدايع ذكر اين ظهور
اعظم متذكّر دارند كه ذكرى باشد آن كتاب لاثبات امر اللّه و تذكّر آن نفوس تا
آنكه جميع مستعدّ شوند از براى تصديق و تسليم در حين ظهور مشيّت اوليّه و كلمه
جامعه و استواء كينونت قدميّه بر عرش رحمانيّه چنانچه فرموده قوله عزّ ذكره
انتم فى كلّ واحد كتاب اثبات لمن نظهره بعضكم الى بعض تكتبون لانّكم يوم
ظهوره بما تكتبون لتعملون و لكن آنچه اهل بيان معمول داشته‌اند اين است كه
مشاهده ميشود در رأس ميقات و وعده الهيّه كلّ الواح ردّيّه بر جمال احديّه
نوشته‌اند و بعضى ببعضى ارسال داشته‌اند و بعضى هم بساحت عرش فرستاده‌اند و از
جمله كتاب سجّين آن نفس غافل است كه بين يدى حاضر شده فافّ لوفائكم فافّ
لحيائكم فافّ لاَدَبكم فافّ لما اكتسبت ايديكم يا ملأ الغافلين

\end{otherlanguage}

Va ay kāsh muhlat mīdādand kih Jamāl-i Rabb al-ʿālamīn bih thanāʾi Allāh
wa-dhikrihi mashghūl shavad. Fa-wa-llāh, furṣat-i ānkih bar nafs-i khud
nawḥih namāyad az ẓulm-i mushrikīn nayāftih, tā chih rasad bih dhikr va
thanā. Maʿ ānkih Nuqṭih-yi Bayān kull-i ahl-i Bayān rā amr farmūdih kih
dar har Vāḥid baʿḍī bih baʿḍī kitābī marqūm dārand va dar ān kitāb
yakdīgar rā bih badāyiʿ-i dhikr-i īn ẓuhūr-i aʿẓam mutadhakkir dārand
kih dhikrī bāshad ān kitāb li-ithbāti amri Allāh va tadhakkur-i ān nufūs
tā ānkih jamīʿ mustaʿidd shavand az barā-yi taṣdīq va taslīm dar ḥīn-i
ẓuhūr-i mashiyyat-i avvaliyyih va kalimih-yi jāmiʿih va istivā-yi
kaynūnat-i qidamiyyih bar ʿarsh-i raḥmāniyyih. Chinānchih farmūdih,
qawluhu ʿazza dhikruhu: Antum fī kulli wāḥidin kitāba ithbātin li-man
nuẓhiruhu baʿḍukum ilā baʿḍin taktubūn li-annakum yawma ẓuhūrihi bi-mā
taktubūna la-taʿmalūn. Va lākin ānchih ahl-i Bayān maʿmūl dāshtih-and īn
ast kih mushāhidih mīshavad. Dar raʾs-i mīqāt va vaʿdih-yi ilāhiyyih
kull alvāḥ-i raddiyyih bar Jamāl-i Aḥadiyyih nivishtih-and va baʿḍī bih
baʿḍī irsāl dāshtih-and va baʿḍī ham bih sāḥat-i ʿarsh firistādih-and.
Va az jumlih kitāb-i sijjīn-i ān nafs-i ghāfil ast kih bayn-i yaday
ḥāḍir shudih. Fa-uffin li-wafāʾikum fa-uffin li-ḥayāʾikum fa-uffin
li-adabikum fa-uffin li-mā iktasabat aydīkum yā malaʾa al-ghāfilīn.

If only they would allow the Beauty of the Lord of the worlds time to
occupy himself with the praise and remembrance of God! By God, because
of the oppression of those who associate partners with him, he has found
no opportunity even to lament his own plight, much less to offer
remembrance and praise. Yet the Point of the Bayān commanded all the
people of the Bayān to write to one another in every Vāḥid, and in that
book to remind one another of this Most Great Manifestation through
wondrous expressions of remembrance. That book was to be a reminder
establishing the Cause of God and awakening those souls, so that all
might be prepared to acknowledge and submit at the appearance of the
Primal Will and the all-embracing Word, and the establishment of the
Eternal Being upon the throne of mercy. Thus he said---exalted be his
remembrance: ``In every Vāḥid, you are to write to one another a book
establishing the truth of the One whom we shall manifest, for on the day
of his appearance you will surely act according to what you write.'' But
what the people of the Bayān have actually done is what is now
witnessed. At the appointed time of the divine promise, they have all
written tablets refuting the Beauty of Oneness and sent them to one
another; some have also sent them to the court of the throne. Among
these is the book of Sijjīn by that heedless soul, which has been
brought into this presence. Shame on your fidelity! Shame on your
modesty! Shame on your courtesy! Shame on what your hands have wrought,
O company of the heedless!

\subsection{461}\label{section-956}

\begin{otherlanguage}{persian}

ملاحظه نمائيد اى اهل بيان كه آن سلطان سرير تقدير چه مقدار تدبير فرموده
بعينه مشابه آنكه انسان طيور را نطق تعليم ميدهند و واللّه اعظم از آن اهل
بيان را تعليم فرموده‌اند و اوقات مبارك را مصروف داشته‌اند كه مباد در حين
ظهور احدى وارد آورد امرى را كه سبب حزن آن قلب رقيق لطيف منير شود چنانچه
در مقامى بعد از نصايح مشفقه متقنه محكمه ميفرمايند قوله جلّ اجلاله فانّ
مثله جلّ ذكره كمثل الشّمس لو يقابلنه الى ما لا نهاية مراياء كلّهنّ ليستعكسن
عن تجلّى الشّمس فى حدّهم و ان لم يقابلها من احد فيطلع الشّمس و يغرب و الحجاب
للمرايا و انّى ما قصّرت عن نصحى لذلك الخلق و تدبيرى لاقبالهم الى اللّه ربّهم
و ايمانهم باللّه بارئهم انتهى

\end{otherlanguage}

Mulāḥiẓih namāyīd ay ahl-i Bayān kih ān sulṭān-i sarīr-i taqdīr chih
miqdār tadbīr farmūdih, bi-ʿaynih mushābih-i ānkih insān ṭuyūr rā nuṭq
taʿlīm mīdihand, va va-llāh aʿẓam az ān ahl-i Bayān rā taʿlīm
farmūdih-and va awqāt-i mubārak rā maṣrūf dāshtih-and kih mabād dar
ḥīn-i ẓuhūr aḥadī vārid āvarad amrī rā kih sabab-i ḥuzn-i ān qalb-i
raqīq-i laṭīf-i munīr shavad. Chinānchih dar maqāmī baʿd az naṣāyiḥ-i
mushfiqih-yi mutqanih-yi muḥkamih mīfarmāyand, qawluhu jalla ijlāluhu:
fa-inna mithlahu jalla dhikruhu ka-mathali al-shams law yuqābilnahu ilā
mā lā nihāyata marāyā kulluhunna la-yastaʿkisna ʿan tajallī al-shams fī
ḥaddihim wa-in lam yuqābilhā min aḥad fa-yaṭluʿu al-shams wa-yaghrub
wa-al-ḥijābu lil-marāyā wa-innī mā qaṣṣartu ʿan nuṣḥī li-dhālika
al-khalq wa-tadbīrī li-iqbālihim ilā Allāh rabbihim wa-īmānihim bi-Allāh
bāriʾihim. Intahā.

Consider, O people of the Bayān, how much that Sovereign of the throne
of destiny has done to prepare the way. It is just as people teach birds
to speak; indeed, by God, he has taught the people of the Bayān far more
than this, devoting his blessed time to ensuring that, at the time of
the Manifestation, no one should do anything that would grieve that
tender, gentle, radiant heart. Thus, in one passage, after giving
compassionate, firm, and well-founded counsels, he says---exalted be his
glory: ``His likeness---exalted be his remembrance---is that of the sun.
If infinitely many mirrors face him, they will all reflect the sun's
radiance according to their capacity. And if none faces it, the sun will
still rise and set, while the veiling belongs to the mirrors. I have
spared no effort in counseling these people and preparing them to turn
toward God, their Lord, and to believe in God, their Maker.'' End of
quotation.

\subsection{462}\label{section-957}

\begin{otherlanguage}{persian}

حال ملاحظه نمائيد در آن بيان كه فرموده‌اند در رأس هر واحد كتابی‌ در اثبات
ظهور بعد مرقوم دارند چه تدبير فرموده‌اند و هم‌‌چنين در اين بيان كه
ميفرمايند اگر مراياى ما لانهايه مقابل شوند بشمس هرآينه تجلّى در حدّ آن
مرايا ظاهر و اگر كلّ محتجب مانند شمس در كمال ضيا و اشراق و انوار طالع و
مشرق چنانچه مشهود است بافاضه‌اش بر كلّ مرايا و انطباع كلّ از او بر قدرش
نيفزايد و هم‌‌چنين اگر كلّ از او محتجب مانند نقصى بر او وارد نگردد و تدبير
اين بيان آن جمال رحمن اينكه اهل بيان در يوم ظهور و اشراق شمس معانى
بمرايا متمسّك نشوند چه كه وجود و عدم كلّ مرايا نزد آن شمس ممتنعه منيعه فى
حدّ سواء بوده

\end{otherlanguage}

Ḥāl mulāḥiẓih namāyīd dar ān bayān kih farmūdih-and dar raʾs-i har vāḥid
kitābī dar ithbāt-i ẓuhūr-i baʿd marqūm dārand chih tadbīr farmūdih-and,
va ham-chinīn dar īn bayān kih mīfarmāyand agar marāyā-yi mā lā-nihāyih
muqābil shavand bi-shams har-āyinih tajallī dar ḥadd-i ān marāyā ẓāhir
va agar kull muḥtajib mānand shams dar kamāl-i ḍiyā va ishrāq va anvār
ṭāliʿ va mushriq. Chinānchih mashhūd ast bi-ifāḍih-ash bar kull-i marāyā
va inṭibāʿ-i kull az ū bar qadrash nayafzāyad va ham-chinīn agar kull az
ū muḥtajib mānand naqṣī bar ū vārid nagardad. Va tadbīr-i īn bayān-i ān
Jamāl-i Raḥmān īnkih ahl-i Bayān dar yawm-i ẓuhūr va ishrāq-i shams-i
maʿānī bi-marāyā mutamassik nashavand, chih kih vujūd va ʿadam-i kull-i
marāyā nazd-i ān shams-i mumtaniʿih-yi manīʿih fī ḥadd-i savāʾ būdih.

Now consider what provision he made in that statement in which he
instructed them to write, at the beginning of every Vāḥid, a book
establishing the truth of the next Manifestation, and likewise in this
statement in which he says that, if infinitely many mirrors face the
sun, its radiance will assuredly appear according to the capacity of
those mirrors; and if all remain veiled, the sun will still rise and
shine in the fullness of its light, splendor, and radiance. As is
evident, its outpouring upon all the mirrors, and the reflection of all
of them from it, add nothing to its worth; likewise, if all remain
veiled from it, no deficiency befalls it. The purpose of this statement
by that Beauty of the All-Merciful is that the people of the Bayān
should not cling to the mirrors on the day when the Sun of inner
meanings appears and shines forth, for the existence and nonexistence of
all the mirrors are alike before that inaccessible, exalted Sun.

\subsection{463}\label{section-958}

\begin{otherlanguage}{persian}

فيا ليت كان نقطة الأولی حينئذ و يشهد بانّكم اتخذتم لانفسكم مراياء و
جعلتموها ميزاناً لمعرفة الّذى لولاه ما خلقت حقايقهن و كينوناتهنّ و اسرفتم
فى امر اللّه على شأن علّقتم تصديقكم نفسه بتصديق مراياء المحتجبين فويل لكم
يا معشر الظّالمين و الغافلين مع اين بيانات واضحه لايحه و اين تدبيرات
محكمه متقنه آن شمس عزّ احديّه بعد از اتمام واحد اوّل و ظهور نيّر آفاق‌ كلّ
ميثاق اللّه را شكسته بنفاق برخواستند و تدابير و نصايح الهى را فراموش
نمودند كانّهم اتّخذوها هزوا فويل لهم بما ارتكبوا تلقآء وجه اللّه العلىّ
العظيم

\end{otherlanguage}

Fa-yā layta kāna Nuqṭatu al-Ūlā ḥīnaʾidhin wa-yashhadu bi-annakum
ittakhadhtum li-anfusikum marāyā wa-jaʿaltumūhā mīzānan li-maʿrifati
alladhī lawlāhu mā khuliqat ḥaqāʾiquhunna wa-kaynūnātuhunna wa-asraftum
fī amri Allāh ʿalā shaʾnin ʿallaqtum taṣdīqakum nafsahu bi-taṣdīqi
marāyā al-muḥtajibīn fa-waylun lakum yā maʿshara al-ẓālimīn
wa-al-ghāfilīn. Maʿ-i īn bayānāt-i vāḍiḥih-yi lāyiḥih va īn tadbīrāt-i
muḥkamih-yi mutqanih-yi ān shams-i ʿizz-i aḥadiyyih, baʿd az itmām-i
vāḥid-i avval va ẓuhūr-i nayyir-i āfāq kull mīthāq-i Allāh rā shikastih
bi-nifāq barkhāstand va tadābīr va naṣāyiḥ-i ilāhī rā farāmūsh namūdand.
Ka-annahum ittakhadhūhā huzuwan fa-waylun lahum bi-mā irtakabū tilqāʾa
wajhi Allāh al-ʿAliyy al-ʿAẓīm.

If only the Primal Point were here now to witness that you have taken
mirrors for yourselves and made them the standard for recognizing the
One without whom their realities and their very beings would not have
been created! You have gone to such excess in the Cause of God that you
have made your acknowledgment of him dependent on the acknowledgment of
veiled mirrors. Woe to you, O company of the unjust and the heedless!
Despite these clear and manifest statements and these firm and
well-founded provisions made by that Sun of the glory of oneness, after
the completion of the first Vāḥid and the appearance of the Luminary of
the horizons, all broke the Covenant of God, rose up in hypocrisy, and
forgot the divine provisions and counsels. It is as though they had made
a mockery of them. Woe to them for what they have committed before the
face of God, the Most High, the Most Great!

\subsection{464}\label{section-959}

\begin{otherlanguage}{persian}

و اينكه نوشته بودى كه اگر اوصاف نازله بجهت اسم بهآء است اين اسم از اسماء
اللّه و تعريفات اسماء اللّه اختصاص باين اسم مبارك نداشته و استهزاء على
اللّه ذكر ميرزا كمال و ميرزا جمال نموده بودى و بذلك صرت من اوّل من استهزء
على فطرة البديع ثمّ على الخليل ثمّ على الكليم تاللّه انّك انت اوّل من استهزء
على الرّوح ثمّ استهزء بمحمّد رسول اللّه ثمّ استهزء بعلىّ حين الّذى فلق فجر
المعانى و اشرقت شمس البيان بسلطنة اللّه المهيمن القيّوم و كنت اوّل من
استهزء على هذا الظّهور الّذى به ثبت كلّ ما نزل فى ازل الآزال و به استعرج كلّ
نفس الى معراج الحقايق و دخل كلّ موقن فى رضوان اللّه المقتدر العزيز المحبوب
خذ قلمك يا ايّها المستهزء باللّه تاللّه به ناح قلم الأعلی ثمّ ما كان و ما
يكون و كلّما ارتكبت يرجع و ينتهى الى قلم اوّل من اعرض باللّه الواحد الفرد
الملك المهيمن القيّوم

\end{otherlanguage}

Va īnkih nivishtih būdī kih agar awṣāf-i nāzilih bi-jahat-i ism-i Bahāʾ
ast īn ism az asmāʾ-i Allāh va taʿrīfāt-i asmāʾ-i Allāh ikhtiṣāṣ bi-īn
ism-i mubārak nadāshtih va istihzāʾan ʿalā Allāh dhikr-i Mīrzā Kamāl va
Mīrzā Jamāl namūdih būdī. Wa-bi-dhālika ṣirta min awwali man istahzaʾa
ʿalā fiṭrati al-Badīʿ thumma ʿalā al-Khalīl thumma ʿalā al-Kalīm.
Ta-Allāh innaka anta awwalu man istahzaʾa ʿalā al-Rūḥ thumma istahzaʾa
bi-Muḥammad Rasūl Allāh thumma istahzaʾa bi-ʿAlī ḥīna alladhī falaqa
fajra al-maʿānī wa-ashraqat shamsu al-bayān bi-salṭanati Allāh
al-Muhaymin al-Qayyūm. Wa-kunta awwala man istahzaʾa ʿalā hādhā al-ẓuhūr
alladhī bihi thabata kullu mā nazala fī azal al-āzāl wa-bihi istaʿraja
kullu nafsin ilā miʿrāj al-ḥaqāʾiq wa-dakhala kullu mūqinin fī Riḍwān
Allāh al-Muqtadir al-ʿAzīz al-Maḥbūb. Khudh qalamaka yā ayyuhā
al-mustahziʾ bi-Allāh. Ta-Allāh bihi nāḥa Qalamu al-Aʿlā thumma mā kāna
wa-mā yakūn wa-kullamā irtakabta yarjiʿu wa-yantahī ilā qalami awwali
man aʿraḍa bi-Allāh al-Wāḥid al-Fard al-Malik al-Muhaymin al-Qayyūm.

As for your having written that, if the revealed descriptions concern
the name Bahāʾ, this name is one of the names of God, and that the
praises of the names of God are not exclusive to this blessed name, you
also mentioned Mīrzā Kamāl and Mīrzā Jamāl in mockery of God. Through
this you have become one of the first to mock the nature of the Wondrous
Creator, then the Friend, then the One who conversed with God. By God,
you are the first to have mocked the Spirit, then to have mocked
Muḥammad, the Messenger of God, then to have mocked ʿAlī when he caused
the dawn of inner meanings to break and the sun of utterance to shine
through the sovereignty of God, the All-Protecting, the Self-Subsisting.
You were the first to mock this Manifestation, through whom everything
sent down from all eternity was established, through whom every soul
ascended to the heights of realities, and every assured believer entered
the Paradise of God, the All-Powerful, the Almighty, the Beloved. Take
up your pen, you who mock God! By God, because of it the Most Exalted
Pen has lamented, and with it all that has been and all that will be.
Everything you have committed goes back and ultimately leads to the pen
of the first to turn away from God, the One, the Unique, the King, the
All-Protecting, the Self-Subsisting.

\subsection{465}\label{section-960}

\begin{otherlanguage}{persian}

و ان حقيقة قلمه قد فرّ عنه و حضر بين يدى اللّه و ضجّ بضجيج بكى عليه كلّ ما
وقع عليه اسم شیء و ناح بين يدى ربّه ثمّ اشتكى من صاحبه ان انت من الّذينهم
يعرفون حقيقت قلم معرض باللّه در يومى از ایّام بين يدى العرش حاضر و نوحه
نمود بنوحه كه اهل ملأ اعلى بر حال او گريستند و عجب است كه شما ناله و
حنين كلّ شیء را در آن حين نشنيده‌ايد و اصغا ننموده‌ايد و از آن يوم الى حين
جميع مقدّسين و مقرّبين ثمّ حقايق كلّ شیء بنوحه و ندبه و جزع مشغولند در
اينمقام لازم شد كه مناجات او را كه بين يدى العرش نموده ذكر نمايم لعلّ
نفوس مستعدّه از ناله آنمظلوم بمظلوميّت محبوب امكان پى‌برند كه بر آن جمال
اقدس چه وارد شده از ظلم ظالمان فئه بيان و هذا ما ناجى به القلم ربّه و ربّ
كلّ شیء و ربّ العالمين

\end{otherlanguage}

Wa-inna ḥaqīqata qalamihi qad farra ʿanhu wa-ḥaḍara bayna yaday Allāh
wa-ḍajja bi-ḍajījin bakā ʿalayhi kullu mā waqaʿa ʿalayhi ismu shayʾ
wa-nāḥa bayna yaday rabbihi thumma ishtakā min ṣāḥibihi in anta min
alladhīna hum yaʿrifūn. Ḥaqīqat-i qalam-i muʿriḍ bi-Allāh dar yawmī az
ayyām bayna yaday al-ʿarsh ḥāḍir va nawḥih namūd bi-nawḥih kih ahl-i
malaʾ-i aʿlā bar ḥāl-i ū girīstand va ʿajab ast kih shumā nālih va
ḥanīn-i kull-i shayʾ rā dar ān ḥīn nashinīdih-īd va iṣghā nanamūdih-īd
va az ān yawm ilā ḥīn jamīʿ-i muqaddasīn va muqarrabīn thumma ḥaqāyiq-i
kull-i shayʾ bi-nawḥih va nudbih va jazaʿ mashghūland. Dar īn maqām
lāzim shud kih munājāt-i ū rā kih bayna yaday al-ʿarsh namūdih dhikr
namāyam laʿalla nufūs-i mustaʿiddih az nālih-yi ān maẓlūm
bi-maẓlūmiyyat-i Maḥbūb-i imkān pay barand kih bar ān Jamāl-i Aqdas chih
vārid shudih az ẓulm-i ẓālimān-i fiʾih-yi Bayān. Wa-hādhā mā nājā bihi
al-qalamu rabbahu wa-rabba kulli shayʾin wa-rabba al-ʿālamīn.

Indeed, the reality of his pen has fled from him and appeared before
God, crying out with such a cry that everything bearing the name of a
thing wept over it. It lamented before its Lord and then complained of
its owner, if you are among those who understand. One day, the reality
of the pen of the one who has turned away from God appeared before the
Throne and lamented with such a lament that the inhabitants of the
celestial assembly wept over its plight. It is astonishing that you did
not hear or heed the lamentation and plaintive cries of all things at
that moment. From that day until now, all the sanctified ones and those
brought near, and indeed the realities of all things, have been occupied
with lamentation, mourning, and anguish. Here it has become necessary
for me to recount the prayer it offered before the Throne, so that
receptive souls may, from the lament of that wronged one, perceive the
wrongs suffered by the Beloved of all creation---what has befallen that
Most Holy Beauty through the oppression of the oppressors among the
people of the Bayān. This is how the Pen addressed its Lord, the Lord of
all things and the Lord of the worlds:

\subsection{466}\label{section-961}

\begin{otherlanguage}{persian}

سبحانك اللّهمّ يا الهى و محبوبی‌ انت الّذى بأمرك ارتفع صرير قلم الأعلی بين
الأرض و السّمآء و به قضيت على الألواح ما قضيت بعلمك المحفوظ و بأمرك
المحتوم و بما اثبتت به على اللّوح تظهر فى الابداع ما قدّر بمشيّتك و قضى
بارادتك اذاً يا الهى اسمع ندآء هذا القلم الّذى بعد عن شاطى قربك و لقائك و
ابتلى بين انامل احد من اشقيآء خلقك فو‌عزّتك يا محبوبی‌ اكون خائفاً من سطوة
قهرك و شئونات انتقامك اذاً تشهد اضطراب نفسى و تزلزل فؤادى و تبلبل جسدى
بما جعلتنى مقهورا بأنامل هذا المشرك الّذى ما استنشق روايح الانصاف من
رضوان عدلك و مواهبك و ما وجد نفحات الطّيب من قميص فضلك و الطافك

\end{otherlanguage}

Subḥānaka Allāhumma yā ilāhī wa-maḥbūbī anta alladhī bi-amrika irtafaʿa
ṣarīru Qalami al-Aʿlā bayna al-arḍ wa-al-samāʾ wa-bihi qaḍayta ʿalā
al-alwāḥ mā qaḍayta bi-ʿilmika al-maḥfūẓ wa-bi-amrika al-maḥtūm wa-bi-mā
athbatta bihi ʿalā al-lawḥ tuẓhiru fī al-ibdāʿ mā quddira bi-mashīʾatika
wa-quḍiya bi-irādatika. Idhan yā ilāhī ismaʿ nidāʾa hādhā al-qalam
alladhī baʿuda ʿan shāṭiʾi qurbika wa-liqāʾika wa-ibtuliya bayna anāmili
aḥadin min ashqiyāʾi khalqika. Fa-wa-ʿizzatika yā maḥbūbī akūnu khāʾifan
min saṭwati qahrika wa-shuʾūni intiqāmika. Idhan tashhadu iḍṭirāba nafsī
wa-tazalzula fuʾādī wa-tabalbula jasadī bi-mā jaʿaltanī maqhūran
bi-anāmili hādhā al-mushrik alladhī mā istanshaqa rawāʾiḥa al-inṣāf min
Riḍwāni ʿadlika wa-mawāhibika wa-mā wajada nafaḥāti al-ṭīb min qamīṣi
faḍlika wa-alṭāfika.

``Glory be to you, O God, my God and my Beloved! You are the One at
whose command the sound of the Most Exalted Pen rose between earth and
heaven. Through it you decreed upon the Tablets what you decreed through
your guarded knowledge and your irrevocable command; and through what
you inscribed with it upon the Tablet, you bring forth in creation what
was ordained by your will and decreed by your purpose. Now, my God, hear
the cry of this pen, which has become far removed from the shore of
nearness to you and meeting with you, and has been afflicted between the
fingers of one of the most wretched of your creatures. By your might, my
Beloved, I am afraid of the force of your wrath and the manifestations
of your vengeance. You now witness the agitation of my soul, the
trembling of my heart, and the turmoil of my body, because you have
placed me in subjection to the fingers of this one who associates
partners with you, who has not inhaled the fragrances of fairness from
the paradise of your justice and your gifts, nor sensed the sweet
breezes from the garment of your grace and your loving-kindness.

\subsection{467}\label{section-962}

\begin{otherlanguage}{persian}

اذاً يا الهى فو‌عزّتك لو تأمرنى لاشقّ بطنه من قدرة الّتى قدرت فى سرّى لعلّ بذلك
تستريح نفسى و يسكن فؤادى لانّه استعان منّى فى هتك حجاب عزّتك و حرمتك و
ارتكب ما احترقت به افئدة اصفيائك فى خبآء مجدك و قلوب امنائك فى غرفات عزّ
تقديسك و خيام قدس تجريدك اسئلك يا محبوبی‌ بسلطانك ثمّ باسمك و ظهورك و
كبريائك بانّك ان لا تأمرنى بما دعوتك به خلّصنى عن هذا الخبيث و انامله

\end{otherlanguage}

Idhan yā ilāhī fa-wa-ʿizzatika law taʾmurunī la-ashuqqanna baṭnahu min
qudrati allatī quddirat fī sirrī laʿalla bi-dhālika tastarīḥu nafsī
wa-yaskunu fuʾādī li-annahu istaʿāna minnī fī hatki ḥijābi ʿizzatika
wa-ḥurmatika wa-irtakaba mā iḥtaraqat bihi afʾidatu aṣfiyāʾika fī
khibāʾi majdika wa-qulūbu umanāʾika fī ghurufāti ʿizzi taqdīsika
wa-khiyāmi qudsi tajrīdika. Asʾaluka yā maḥbūbī bi-sulṭānika thumma
bi-ismika wa-ẓuhūrika wa-kibriyāʾika bi-annaka in lā taʾmuranī bi-mā
daʿawtuka bihi khalliṣnī ʿan hādhā al-khabīth wa-anāmilihi.

``Now, my God, by your might, if you command me, I will assuredly tear
open his belly through the power ordained within my inmost being, so
that my soul may perhaps find rest and my heart be still. For he has
enlisted my aid in rending the veil of your majesty and sanctity, and
has committed that which has burned the hearts of your chosen ones
within the pavilion of your glory, and the hearts of your trusted ones
in the chambers of your sanctified might and the tents of your holy
detachment. I ask you, my Beloved, by your sovereignty, and by your
name, your manifestation, and your grandeur: if you do not command me to
do what I have asked of you, deliver me from this wicked one and his
fingers.

\subsection{468}\label{section-963}

\begin{otherlanguage}{arabic}

فآه آه يا محبوبی‌ انت خلقتنى لذكرك و ثنآء نفسك و انّى صرت محروماً عمّا خلقت
له و ظهر منّى ما شقّ به ستر العظمة و الكبرياء و تغيّرت وجوه اهل ملأ الأعلی
و انعدمت حقايق من فى ملكوت الأسماء و تزلزلت اركان مداين البقآء فكيف يا
الهى ارفع رأسى بين الاقلام لخجلة الّتى به ناكس رأسى فى محضر المقرّبين عند
ظهورات انوار جمالك و مقعد الموحّدين على فنآء باب عظمتك فوعزّتك لو تلبسنى
رداء الغفران و تهبّ علىّ من شطر عنايتك نفحات اسمك الرّحمن و تغمسنى فى ابحر
عطوفتك و الطافك و تغفرنى عمّا ارتكبت فى ایّامك فو‌عزّتك لن يسكن فؤادى و لن
تطمئن نفسى هبنى يا محبوبی‌ بأنّك بفضلك وعنايتك طهّرتنى عن العصيان فكيف يذهب
منّى روايح خجلة الّتى تهبّ من نفسى على نفسى فى كلّ الأحيان و جعلنى محروماً عن
نفحات رضوان الّذى جعلته فوق الجنان

\end{otherlanguage}

Fa-āh āh yā maḥbūbī anta khalaqtanī li-dhikrika wa-thanāʾi nafsika
wa-innī ṣirtu maḥrūman ʿammā khuliqtu lahu wa-ẓahara minnī mā shuqqa
bihi sitru al-ʿaẓama wa-al-kibriyāʾ wa-taghayyarat wujūhu ahli malaʾ
al-aʿlā wa-inʿadamat ḥaqāʾiqu man fī malakūt al-asmāʾ wa-tazalzalat
arkānu madāʾini al-baqāʾ. Fa-kayfa yā ilāhī arfaʿu raʾsī bayna al-aqlām
li-khajlati allatī bihi nākisu raʾsī fī maḥḍari al-muqarrabīn ʿinda
ẓuhūrāti anwāri jamālika wa-maqʿadi al-muwaḥḥidīn ʿalā fināʾi bābi
ʿaẓamatika. Fa-wa-ʿizzatika law tulbisunī ridāʾa al-ghufrān wa-tahubbu
ʿalayya min shaṭri ʿināyatika nafaḥātu ismika al-Raḥmān wa-taghmisunī fī
abḥuri ʿuṭūfatika wa-alṭāfika wa-taghfirunī ʿammā irtakabtu fī ayyāmika
fa-wa-ʿizzatika lan yaskuna fuʾādī wa-lan taṭmaʾinna nafsī. Habnī yā
maḥbūbī bi-annaka bi-faḍlika wa-ʿināyatika ṭahhartanī ʿan al-ʿiṣyān
fa-kayfa yadhhabu minnī rawāʾiḥu khajlati allatī tahubbu min nafsī ʿalā
nafsī fī kulli al-aḥyān wa-jaʿalanī maḥrūman ʿan nafaḥāti Riḍwān alladhī
jaʿaltahu fawqa al-jinān.

``Alas, alas, my Beloved! You created me to remember you and praise your
very self, yet I have been deprived of that for which I was created.
From me has proceeded that which has torn the veil of majesty and
grandeur, changed the faces of the inhabitants of the celestial
assembly, reduced to nothing the realities of those in the kingdom of
names, and shaken the foundations of the cities of eternity. How, then,
my God, can I raise my head among the pens, because of the shame that
makes me bow my head in the presence of those brought near, at the
manifestations of the lights of your beauty, and in the gathering of
those who affirm your oneness in the court of the gate of your majesty?
By your might, even if you clothe me in the robe of forgiveness, the
breezes of your name, the All-Merciful, blow upon me from the direction
of your loving care, you immerse me in the oceans of your tenderness and
loving-kindness, and you forgive me for what I have committed in your
days---by your might, my heart will not be still and my soul will not
find peace. Suppose, my Beloved, that through your grace and loving care
you have cleansed me of disobedience: how can the odors of shame leave
me, when they blow from my own self upon my own self at every moment and
have deprived me of the breezes of the Paradise you have placed above
all gardens?

\subsection{469}\label{section-964}

\begin{otherlanguage}{arabic}

و انّك يا الهى و مالكى لو تجاوزت عنّى بحلمك و غفرت لى خطيئتى بفضلك و كرمك
كيف ارفع رأسى تلقآء مدين عزّك و لقائك فو‌عزّتك قد بلغت فى الذّلّة الى مقام لو
تنظر الىّ بلحظات رأفتك لتبكى بنفسك على هذا المسكين الّذى صار مأيوساً عن
نفسه و عن كلّ الجهات بما جعل محروماً من بوارق انوار جمالك الّتى اشرقت عن
افق سماء عزّك و مشيّتك فآه آه عن يأسى فى هذا اليوم الّذى فتحت فيه ابواب
وصلك لمن فى ارضك و سمائك و دعوت الكلّ الى مقرّ قربك و لقائك فآه آه عمّا قدّر
لى فى الواح قضائك و محوت به حظّى و لذيذ مناجاتى عند مشاهدة انوار وجهك فيا
ليت كنت محروماً عن كلّ ما قدّرته لخير بريّتك و ما عاشرت مع الّذين ما عرفوك و
ما سجدوك و ظهر منهم باعانتى ما احترقت عنه قلوب امنائك و اصفيائك

\end{otherlanguage}

Wa-innaka yā ilāhī wa-mālikī law tajāwazta ʿannī bi-ḥilmika wa-ghafarta
lī khaṭīʾatī bi-faḍlika wa-karamika kayfa arfaʿu raʾsī tilqāʾa madīni
ʿizzika wa-liqāʾika. Fa-wa-ʿizzatika qad balaghtu fī al-dhilla ilā
maqāmin law tanẓuru ilayya bi-laḥaẓāti raʾfatika la-tabkī bi-nafsika
ʿalā hādhā al-miskīn alladhī ṣāra maʾyūsan ʿan nafsihi wa-ʿan kulli
al-jihāt bi-mā juʿila maḥrūman min bawāriqi anwāri jamālika allatī
ashraqat ʿan ufuqi samāʾi ʿizzika wa-mashīʾatika. Fa-āh āh ʿan yaʾsī fī
hādhā al-yawm alladhī fataḥta fīhi abwāba waṣlika li-man fī arḍika
wa-samāʾika wa-daʿawta al-kulla ilā maqarri qurbika wa-liqāʾika. Fa-āh
āh ʿammā quddira lī fī alwāḥi qaḍāʾika wa-maḥawta bihi ḥaẓẓī wa-ladhīdha
munājātī ʿinda mushāhadati anwāri wajhika. Fa-yā layta kuntu maḥrūman
ʿan kulli mā qaddartahu li-khayri bariyyatika wa-mā ʿāshartu maʿa
alladhīna mā ʿarafūka wa-mā sajadūka wa-ẓahara minhum bi-iʿānatī mā
iḥtaraqat ʿanhu qulūbu umanāʾika wa-aṣfiyāʾika.

``Even if you, my God and my Master, pardon me through your forbearance
and forgive my sin through your grace and generosity, how can I raise my
head before the city of your might and your presence? By your might, I
have reached such a degree of abasement that, if you looked upon me with
the glances of your compassion, you yourself would weep over this
wretched one, who has lost hope in himself and in every direction
because he has been deprived of the flashing lights of your beauty that
shone from the horizon of the heaven of your might and your will. Alas,
alas for my despair on this day when you have opened the gates of union
with you to those in your earth and your heaven, and called everyone to
the seat of nearness to you and meeting with you! Alas, alas for what
has been ordained for me in the tablets of your decree, through which
you have effaced my portion and the sweetness of my communion with you
while beholding the lights of your face! If only I had been deprived of
all that you ordained for the good of your creatures and had not
associated with those who neither knew you nor bowed down before you,
and from whom, with my assistance, proceeded that which burned the
hearts of your trusted and chosen ones!

\subsection{470}\label{section-965}

\begin{otherlanguage}{persian}

و لو اقول هذا ورد علىّ بما كنت مستحقّاً به فو‌عزّتك يا محبوبی‌ لم اجد لنفسى
استقراراً عند هبوب ارياح قضائك و لا لكينونتى وجودا تلقآء مدين سلطان قدرتك
و اقتدارك و لا لى ذكر عند ظهورات عزّ وحدانيّتك و الّذى لم يكن له وجود تلقآء
ظهورك و قدرتك كيف يقدر ان يتحرّك بغير اذنك و ارادتك

\end{otherlanguage}

Wa-law aqūlu hādhā warada ʿalayya bi-mā kuntu mustaḥiqqan bihi
fa-wa-ʿizzatika yā maḥbūbī lam ajid li-nafsī istiqrāran ʿinda hubūbi
aryāḥi qaḍāʾika wa-lā li-kaynūnatī wujūdan tilqāʾa madīni sulṭāni
qudratika wa-iqtidārika wa-lā lī dhikran ʿinda ẓuhūrāti ʿizzi
waḥdāniyyatika wa-alladhī lam yakun lahu wujūdun tilqāʾa ẓuhūrika
wa-qudratika kayfa yaqdiru an yataḥarraka bi-ghayri idhnika
wa-irādatika.

``If I say that this has befallen me because I deserved it---by your
might, my Beloved, I have found no stability for myself when the winds
of your decree blow, no existence for my being before the city of the
sovereignty of your power and might, and no mention of myself at the
manifestations of the glory of your oneness. How can one who has no
existence before your manifestation and your power move without your
permission and your will?

\subsection{471}\label{section-966}

\begin{otherlanguage}{persian}

و لو اقول انّك ابتليتنى بذلك من غير استحقاقى هذا بغى منّى عليك لانّى ما
اطّلعت باسرار امرك و شئونات حكمتك و قضائك و بذلك صرت متحيّراً فى امرى و لن
يرتقى عرفانى الى سموات حكمتك الّتى جعلتها مستوراً عن انظر خلقك و بريّتك و
لن يطير طير علمى الى هوآء علمك الّذى جعلته مكنوناً فى كنائز عصمتك و بذلك
يزيد اضطرابی‌ و يشتدّ حزنى و بأسى و ابتلائى

\end{otherlanguage}

Wa-law aqūlu innaka ibtalaytanī bi-dhālika min ghayri istiḥqāqī hādhā
baghyun minnī ʿalayka li-annī mā iṭṭalaʿtu bi-asrāri amrika wa-shuʾūni
ḥikmatika wa-qaḍāʾika wa-bi-dhālika ṣirtu mutaḥayyiran fī amrī wa-lan
yartaqiya ʿirfānī ilā samāwāti ḥikmatika allatī jaʿaltahā mastūran ʿan
anẓar khalqika wa-bariyyatika wa-lan yaṭīra ṭayru ʿilmī ilā hawāʾi
ʿilmika alladhī jaʿaltahu maknūnan fī kanāʾizi ʿiṣmatika wa-bi-dhālika
yazīdu iḍṭirābī wa-yashtaddu ḥuznī wa-baʾsī wa-ibtilāʾī.

``And if I say that you have afflicted me with this without my deserving
it, this would be an act of rebellion on my part against you, for I have
not learned the mysteries of your Cause or the workings of your wisdom
and decree. Thus I have become bewildered about my plight. My
understanding will never ascend to the heavens of your wisdom, which you
have concealed from the sight of your creatures and your creation, nor
will the bird of my knowledge fly into the atmosphere of your knowledge,
which you have hidden in the treasuries of your infallibility. Thus my
agitation increases, and my sorrow, distress, and affliction grow more
severe.

\subsection{472}\label{section-967}

\begin{otherlanguage}{persian}

فآه آه يا محبوبی‌ كيف اذكر ما اشاهد من ظهورات فعلك و شئونات امرك مرّة
اشاهد بانّك اخذت قلم الّذى كان مثلى بأنامل رحمتك و جعلته معاشر نفسك و
مؤانس جمالك و تأخذه باصابع عزّك و كبريائك و تجرى منه بحور الحيوان الّتى
بقطرة منها بعثت حقايق الامكان و افئدة اهل الاكوان و من صريره استجذبت
قلوب المقرّبين و افئدة المخلصين و مرّة اشاهد بأنّك ابتليتنى بأنامل المشركين
و جعلتنى مقهورا تحت اصابعهم و ظهر منّى ما اضطربت منه افئدة ملأ الأعلی ثمّ
سكّان مداين البقآء و تشبّكت اكباد الّذين كرّمت وجوههم عن التّوجّه الى غيرك و
قدّستهم عن اشارات اهل ارضك و استقربتهم فى ظلّ عنايتك و افضالك فو‌عزّتك اخاف
بأنّ من عصيانى يتوقّف قلم امرك و يعطّل الواح قضائك و صحايف تقديرك فيا ليت
ما كنت موجودا و ما كنت مذكورا

\end{otherlanguage}

Fa-āh āh yā maḥbūbī, kayfa adhkuru mā ushāhidu min ẓuhūrāti fiʿlika
wa-shuʾūnāti amrika? Marratan ushāhidu bi-annaka akhadhta qalama alladhī
kāna mithlī bi-anāmili raḥmatika wa-jaʿaltahu muʿāshira nafsika
wa-muʾānisa jamālika wa-taʾkhudhuhu bi-aṣābiʿi ʿizzika wa-kibriyāʾika
wa-tujrī minhu buḥūra al-ḥayawāni allatī bi-qaṭratin minhā baʿathta
ḥaqāʾiqa al-imkāni wa-afʾidata ahli al-akwāni wa-min ṣarīrihi
istajdhabta qulūba al-muqarrabīna wa-afʾidata al-mukhliṣīna, wa-marratan
ushāhidu bi-annaka ibtalaytanī bi-anāmili al-mushrikīna wa-jaʿaltanī
maqhūran taḥta aṣābiʿihim wa-ẓahara minnī mā iḍṭarabat minhu afʾidatu
malaʾi al-aʿlā thumma sukkāni madāʾini al-baqāʾi wa-tashabbakat akbādu
alladhīna karramta wujūhahum ʿani al-tawajjuhi ilā ghayrika
wa-qaddastahum ʿan ishārāti ahli arḍika wa-istaqrabtahum fī ẓilli
ʿināyatika wa-ifḍālika. Fa-wa-ʿizzatika akhāfu bi-anna min ʿiṣyānī
yatawaqqafu qalamu amrika wa-yuʿaṭṭalu alwāḥu qaḍāʾika wa-ṣaḥāʾifu
taqdīrika. Fa-yā layta mā kuntu mawjūdan wa-mā kuntu madhkūran.

Ah, ah, my Beloved! How can I recount what I witness of the
manifestations of your action and the workings of your Cause? At one
time I see that you have taken a pen like me in the fingers of your
mercy and made it a companion to your own self and an intimate of your
beauty. You take it in the fingers of your might and majesty and cause
oceans of life to flow from it, with a drop of which you have quickened
the realities of contingent existence and the hearts of the inhabitants
of all worlds. Through its scratching you have drawn to yourself the
hearts of those brought near and the souls of the sincere. At another
time I see that you have afflicted me with the fingers of those who
associate partners with you and made me powerless beneath their fingers.
From me has appeared that which has troubled the hearts of the Concourse
on high and the inhabitants of the cities of eternity, and has entangled
the inmost hearts of those whose faces you have honored by turning them
away from all but you, whom you have sanctified from the allusions of
the people of your earth, and whom you have drawn near beneath the
shadow of your care and bounty. By your might, I fear that my
disobedience will bring the Pen of your Cause to a halt and leave the
tablets of your decree and the pages of your ordainment idle. If only I
had not existed and had never been mentioned!

\subsection{473}\label{section-968}

\begin{otherlanguage}{arabic}

فو‌عزّتك يا الهى لن اجد لنفسى وجودا بعد ذلك و ان لا ترجعنى الى العدم فاحكم
بينى و بين هذا الظّالم الّذى من فعله تغيّر وجهى بين المخلصين من بريّتك و
المقرّبين من احبّتك و انّى فى تلك السّاعة يا الهى قد هربت منه اليك و جئتك
بضجيج الفاقدين و صيحة العاصين و نوح الخاطئين اذا فاحكم بينى و بينه
بسلطان قضائك و مليك عدلك و اقتدارك هل من حاكم يا محبوبی‌ دونك لالتجآء به
و افرّ اليه لا فو‌عزّتك لا مهرب الّا انت و لا مفرّ الّا اليك

\end{otherlanguage}

Fa-wa-ʿizzatika yā ilāhī lan ajida li-nafsī wujūdan baʿda dhālika wa-in
lā turjiʿnī ilā al-ʿadami fa-uḥkum baynī wa-bayna hādhā al-ẓālimi
alladhī min fiʿlihi taghayyara wajhī bayna al-mukhliṣīna min bariyyatika
wa-al-muqarrabīna min aḥibbatika. Wa-innī fī tilka al-sāʿati yā ilāhī
qad harabtu minhu ilayka wa-jiʾtuka bi-ḍajīji al-fāqidīna wa-ṣayḥati
al-ʿāṣīna wa-nawḥi al-khāṭiʾīna. Idhan fa-uḥkum baynī wa-baynahu
bi-sulṭāni qaḍāʾika wa-malīki ʿadlika wa-iqtidārika. Hal min ḥākimin yā
maḥbūbī dūnaka li-altajiʾa bihi wa-afirra ilayhi? Lā fa-wa-ʿizzatika lā
mahraba illā anta wa-lā mafarra illā ilayka.

By your might, my God, I will find no existence for myself after this.
If you do not return me to nonexistence, then judge between me and this
oppressor, whose action has altered my face among the sincere among your
creatures and those of your loved ones who have been brought near. At
that very hour, my God, I fled from him to you and came to you with the
clamor of the bereaved, the cry of the disobedient, and the lamentation
of sinners. Judge, then, between me and him by the sovereignty of your
decree and the dominion of your justice and power. Is there any judge
besides you, my Beloved, with whom I may seek refuge and to whom I may
flee? No, by your might! There is no refuge but you and nowhere to flee
except to you.

\subsection{474}\label{section-969}

\begin{otherlanguage}{persian}

و انّ هذا مظلوم قد وقف تلقآء مدين عدلك و رحمتك و انّ هذا فقير قد تشبّث بذيل
غنائك و انّ هذا محروم قد استظلّ فى ظلّ حرم قدسك و الطافك اسئلك باسمك الّذى
به كسّرت اصنام الوهم و الهوى و زيّنت المقرّبين بخلع عرفانك و غفرانك و نزعت
عن المشركين رداء اسمائك و الطافك بأن تنزل من سحاب قدرتك و سمآء فضلك ما
يسكن به قلبی‌ و تطمئن به نفسى

\end{otherlanguage}

Wa-inna hādhā maẓlūmun qad waqafa tilqāʾa madīni ʿadlika wa-raḥmatika
wa-inna hādhā faqīrun qad tashabbatha bi-dhayli ghināʾika wa-inna hādhā
maḥrūmun qad istaẓalla fī ẓilli ḥarami qudsika wa-alṭāfika. Asʾaluka
bi-ismika alladhī bihi kassarta aṣnāma al-wahmi wa-al-hawā wa-zayyanta
al-muqarrabīna bi-khilaʿi ʿirfānika wa-ghufrānika wa-nazaʿta ʿani
al-mushrikīna ridāʾa asmāʾika wa-alṭāfika bi-an tunzila min saḥābi
qudratika wa-samāʾi faḍlika mā yaskunu bihi qalbī wa-taṭmaʾinnu bihi
nafsī.

This is a wronged one who has stood before the city of your justice and
mercy; this is a poor one who has clung to the hem of your wealth; this
is a deprived one who has sought shelter in the shadow of the sanctuary
of your holiness and kindness. I ask you by your name through which you
shattered the idols of illusion and desire, adorned those brought near
with the robes of your recognition and forgiveness, and stripped from
those who associate partners with you the mantle of your names and
kindness, to send down from the clouds of your power and the heaven of
your grace that which will calm my heart and bring peace to my soul.

\subsection{475}\label{section-970}

\begin{otherlanguage}{persian}

فآه آه يا محبوبی‌ اذا قد اخذنى الاضطراب و الاضطرار عند استقرار عرش اسمك
المختار لأنّك لو تعذّب ظالم الّذى ظلمنى بدوام ملكك و ملكوتك لن تفرح بذلك
نفسى و لن يذهب وحشتى و لن يسكن اضطرابی‌ و اضطرارى لانّ وجهى يا الهى اصفرّ
من خجلتى بما ظهر منّى و هذه من رايحة الّتى لن يذهب منّى و لو يهبّ علىّ عن
يمين عنايتك ارياح غفرانك بدوام عزّ احديّتك اذا هل ترى يا محبوبی‌ خاسراً اخسر
منّى او ذليلاً اذلّ عنّى

\end{otherlanguage}

Fa-āh āh yā maḥbūbī, idhan qad akhadhanī al-iḍṭirābu wa-al-iḍṭirāru
ʿinda istiqrāri ʿarshi ismika al-Mukhtāri li-annaka law tuʿadhdhibu
ẓālima alladhī ẓalamanī bi-dawāmi mulkika wa-malakūtika lan tafraḥa
bi-dhālika nafsī wa-lan yadhhaba waḥshatī wa-lan yaskuna iḍṭirābī
wa-iḍṭirārī li-anna wajhī yā ilāhī iṣfarra min khajlatī bi-mā ẓahara
minnī wa-hādhihi min rāʾiḥati allatī lan yadhhaba minnī wa-law yahubbu
ʿalayya ʿan yamīni ʿināyatika aryāḥu ghufrānika bi-dawāmi ʿizzi
aḥadiyyatika. Idhan hal tarā yā maḥbūbī khāsiran akhsara minnī aw
dhalīlan adhalla ʿannī?

Ah, ah, my Beloved! Agitation and distress have seized me while the
throne of your name, the Unconstrained, stands firm. Even if you were to
torment the oppressor who wronged me for as long as your sovereignty and
kingdom endure, my soul would not rejoice at it, my desolation would not
depart, and my agitation and distress would not be stilled. For my face,
my God, has grown pale with shame at what has proceeded from me. This is
a scent that will never leave me, even if the winds of your forgiveness
blow upon me from the right hand of your care for as long as the glory
of your oneness endures. Do you see, then, my Beloved, anyone more
bereft than I, or anyone more abased?

\subsection{476}\label{section-971}

\begin{otherlanguage}{arabic}

و انّى قد كنت يا الهى فى كلّ الأیّام مشتاقاً لوصالك بحيث نمت فى اللّيالى رجآء
للقائك و ما ارفعت رأسى فى الاصباح الّا شوقاً لحضرتك و كنت فى تلك الحالة
الى ان حركت ارياح قضائك عن يمين ارادتك و ظهرت ظهورات تقديرك عن افق قضائك
و انصرفنى عن شطرك الى شطر اعدائك فاه آه من هذا الهبوب فآه آه من هذا
المرور فآه آه من هذه المرسلات الّتى اخذتنى بقدرتك و اودعتنى فى محضر
المشركين بنفسك و المعاندين بجمالك فيا ليت كانوا قانعاً بما ارتكبوا فى
ایّامك و وردوا على نفسك لا فو‌عزّتك اولئك لن يستريحوا الّا بان يسفكوا دمك
بين بريّتك و يأكلوا لحمك بانياب البغضآء فى ملكوت الانشآء هؤلآء الأشقیآء
الّذين يفرّ الفرعون من كبرهم و غرورهم و يهرب النّمرود من طغيانهم و بغيهم و
يستعيذ الشّيطان بك من شرّهم و غلّهم

\end{otherlanguage}

Wa-innī qad kuntu yā ilāhī fī kulli al-ayyāmi mushtāqan li-wiṣālika
bi-ḥaythu nimtu fī al-layālī rajāʾan li-liqāʾika wa-mā arfaʿtu raʾsī fī
al-iṣbāḥi illā shawqan li-ḥaḍratika wa-kuntu fī tilka al-ḥālati ilā an
ḥarakat aryāḥu qaḍāʾika ʿan yamīni irādatika wa-ẓaharat ẓuhūrātu
taqdīrika ʿan ufuqi qaḍāʾika wa-inṣarafanī ʿan shaṭrika ilā shaṭri
aʿdāʾika. Fa-āh āh min hādhā al-hubūbi fa-āh āh min hādhā al-murūri
fa-āh āh min hādhihi al-mursalāti allatī akhadhatnī bi-qudratika
wa-awdaʿatnī fī maḥḍari al-mushrikīna bi-nafsika wa-al-muʿānidīna
bi-jamālika. Fa-yā layta kānū qāniʿan bi-mā irtakabū fī ayyāmika
wa-waradū ʿalā nafsika. Lā fa-wa-ʿizzatika ulāʾika lan yastarīḥū illā
bi-an yasfikū damaka bayna bariyyatika wa-yaʾkulū laḥmaka bi-anyābi
al-baghḍāʾi fī malakūti al-inshāʾi. Hāʾulāʾi al-ashqiyāʾu alladhīna
yafirru al-firʿawnu min kibrihim wa-ghurūrihim wa-yahrubu al-nimrūdu min
ṭughyānihim wa-baghyihim wa-yastaʿīdhu al-shayṭānu bika min sharrihim
wa-ghillihim.

Throughout all these days, my God, I longed for reunion with you, so
that I slept at night in the hope of meeting you and raised my head in
the morning only out of longing for your presence. I remained in that
state until the winds of your decree stirred from the right hand of your
will, the manifestations of your ordainment appeared from the horizon of
your decree, and I was turned away from you toward your enemies. Ah, ah,
for this blowing! Ah, ah, for this passing! Ah, ah, for these winds sent
forth, which seized me through your power and left me in the presence of
those who associate partners with your own self and oppose your beauty!
If only they had been content with what they committed in your days and
inflicted upon your own self! No, by your might! They will find no rest
until they shed your blood among your creatures and devour your flesh
with the fangs of hatred in the kingdom of creation. These are the
wretched ones from whose pride and arrogance Pharaoh flees, from whose
tyranny and oppression Nimrod runs away, and from whose evil and malice
Satan seeks refuge in you.

\subsection{477}\label{section-972}

\begin{otherlanguage}{arabic}

فآه آه يا محبوبی‌ دعوتك حين الّذى نبتت فى شاطى قدس احديّتك باسمك الرّئوف ثمّ
باسمك الرّحمن ثمّ باسمك الرّحيم ثمّ باسمك الغفور ثمّ باسمك العطوف ثمّ باسمك
الودود ثمّ باسمك الكافى ثمّ باسمك المعطى و كلّما وجدت نفسى محزوناً بشّرتها
بقربك و لقائك و كلما اضطربت سكنتها بذكر ایّام وصالك فلمّا كمل خلقى اذا
قلعنى عن مكانى احد من عبادك و نقلت من يد الى يد و من سوق الى سوق الى ان
وردت فى سوق الّذى انت عالم به بعلمك الّذى احاط كلّ شیء اذا اشترانى احد من
خلقك و بريّتك و لكن يا الهى و محبوبی‌ فو‌عزّتك لما اخذنى بيده رأيت بأنّه غافل
عن ذلك و كنت فى ذلك متحيّراً فى نفسى لان هذا الضّجيج ظهر منّى من دون اختيار

\end{otherlanguage}

Fa-āh āh yā maḥbūbī, daʿawtuka ḥīna alladhī nabattu fī shāṭiʾi qudsi
aḥadiyyatika bi-ismika al-Raʾūfi thumma bi-ismika al-Raḥmāni thumma
bi-ismika al-Raḥīmi thumma bi-ismika al-Ghafūri thumma bi-ismika
al-ʿAṭūfi thumma bi-ismika al-Wadūdi thumma bi-ismika al-Kāfī thumma
bi-ismika al-Muʿṭī wa-kullamā wajadtu nafsī maḥzūnan bashshartuhā
bi-qurbika wa-liqāʾika wa-kullamā iḍṭarabat sakkantuhā bi-dhikri ayyāmi
wiṣālika. Fa-lammā kamula khalqī idhā qalaʿanī ʿan makānī aḥadun min
ʿibādika wa-nuqiltu min yadin ilā yadin wa-min sūqin ilā sūqin ilā an
waradtu fī sūqi alladhī anta ʿālimun bihi bi-ʿilmika alladhī aḥāṭa kulla
shayʾin. Idhan ishtarānī aḥadun min khalqika wa-bariyyatika wa-lākin yā
ilāhī wa-maḥbūbī fa-wa-ʿizzatika lammā akhadhanī bi-yadihi raʾaytu
bi-annahu ghāfilun ʿan dhālika wa-kuntu fī dhālika mutaḥayyiran fī nafsī
li-anna hādhā al-ḍajīja ẓahara minnī min dūni ikhtiyārin.

Ah, ah, my Beloved! When I sprouted on the shore of the holiness of your
oneness, I called upon you by your name, the Compassionate, then by your
name, the All-Merciful, then by your name, the Merciful, then by your
name, the All-Forgiving, then by your name, the Tenderhearted, then by
your name, the All-Loving, then by your name, the All-Sufficing, then by
your name, the Bestower. Whenever I found my soul sorrowful, I cheered
it with the glad news of nearness to you and of meeting you, and
whenever it grew agitated, I calmed it by recalling the days of reunion
with you. When my growth was complete, one of your servants uprooted me
from my place, and I was passed from hand to hand and from market to
market until I reached the market known to you through your knowledge,
which encompasses all things. Then one of your creatures and created
beings bought me. But, my God and my Beloved, by your might, when he
took me in his hand, I saw that he was unaware of this. I was bewildered
within myself at this, for this clamor arose from me involuntarily.

\subsection{478}\label{section-973}

\begin{otherlanguage}{persian}

و كنت فى سرّ السّرّ دعوتك بندآء المشتاقين و صريخ الأملين بأن تكشف لى سرّ ذلك
و ما سترته عنّى بحجبات تقديرك و سبحات قضائك الى ان حملنى الى ارض الّتى
استويت فيها على عرش مظلوميّتك وجدت رايحة القميص من ترابها و نفحات التّقديس
من هوائها و اخذنى الابتهاج و السّرور على شأن رأيت نفسى على معارج العزّة و
الاجلال كانى صرت مالكا على اعلى ملوك ممالك البقآء و سلطانا على من فى
ملكوت الانشآء

\end{otherlanguage}

Wa-kuntu fī sirri al-sirri daʿawtuka bi-nidāʾi al-mushtāqīna wa-ṣarīkhi
al-āmilīna bi-an takshifa lī sirra dhālika wa-mā satartahu ʿannī
bi-ḥujubāti taqdīrika wa-subuḥāti qaḍāʾika ilā an ḥamalanī ilā arḍi
allatī istawayta fīhā ʿalā ʿarshi maẓlūmiyyatika. Wajadtu rāʾiḥata
al-qamīṣi min turābihā wa-nafaḥāti al-taqdīsi min hawāʾihā wa-akhadhanī
al-ibtihāju wa-al-surūru ʿalā shaʾnin raʾaytu nafsī ʿalā maʿāriji
al-ʿizzati wa-al-ijlāli ka-annī ṣirtu mālikan ʿalā aʿlā mulūki mamāliki
al-baqāʾi wa-sulṭānan ʿalā man fī malakūti al-inshāʾi.

In the innermost secrecy I called upon you with the call of those who
yearn and the cry of those who hope, asking you to disclose to me the
secret of this and what you had concealed from me behind the veils of
your ordainment and the splendors of your decree, until he carried me to
the land where you were established upon the throne of your wronged
condition. I perceived the fragrance of the shirt from its dust and the
breaths of holiness from its air. Joy and gladness so seized me that I
saw myself upon the heights of might and majesty, as though I had become
ruler over the highest kings of the realms of eternity and sovereign
over all who dwell in the kingdom of creation.

\subsection{479}\label{section-974}

\begin{otherlanguage}{arabic}

و كنت فى تلك الحالة الى ان بلغ حاملى الى سبيل انشعب منه السّبيلان سبيل
الى اليمين و سبيل الى اليسار و انّه انحرف من اليمين الى الشّمال فآ آه اذاً
وجدت فى نفسى الاضطراب و فزعت و فزع كلّ‌شیء من فزعى الى ان بلغنى الى باب
لما فتح وجدت رايحة الجحيم و نفحات السجّين الّتى لو تهبّ رايحة منها على
الممكنات كلّها ترجعن الى عدم البحت و اودعنى بيد هذا الظّالم الّذى قام عليك
بالاعراض و ورد منه عليك ما بدلت الأسمآء و ملكوتها و انقطعت الصّفات عن
مقاعدها فلمّا اطّلعت به و بما ارتكبت باعانتى ما ارتكب فررت عنه و عن لقائه
الى ساحة عزّ احديّتك و مقرّ عرش عظمتك

\end{otherlanguage}

Wa-kuntu fī tilka al-ḥālati ilā an balagha ḥāmilī ilā sabīlin inshaʿaba
minhu al-sabīlāni sabīlun ilā al-yamīni wa-sabīlun ilā al-yasāri
wa-innahu inḥarafa min al-yamīni ilā al-shimāli. Fa-ā āh idhan wajadtu
fī nafsī al-iḍṭirāba wa-faziʿtu wa-faziʿa kullu shayʾin min fazaʿī ilā
an ballaghanī ilā bābin lammā futiḥa wajadtu rāʾiḥata al-jaḥīmi
wa-nafaḥāti al-sijjīni allatī law tahubbu rāʾiḥatun minhā ʿalā
al-mumkināti kullihā tarjʿanna ilā ʿadami al-baḥti wa-awdaʿanī bi-yadi
hādhā al-ẓālimi alladhī qāma ʿalayka bi-al-iʿrāḍi wa-warada minhu
ʿalayka mā baddala al-asmāʾa wa-malakūtahā wa-inqaṭaʿati al-ṣifātu ʿan
maqāʿidihā. Fa-lammā iṭṭalaʿtu bihi wa-bi-mā irtakabtu bi-iʿānatī mā
irtakaba farartu ʿanhu wa-ʿan liqāʾihi ilā sāḥati ʿizzi aḥadiyyatika
wa-maqarri ʿarshi ʿaẓamatika.

I remained in that state until the one carrying me reached a road that
branched into two paths, one to the right and one to the left. He turned
from the right toward the left. Ah, ah! Then I felt agitation within me
and was terrified, and everything was terrified at my terror, until he
brought me to a door. When it opened, I perceived the stench of hell and
the breaths of the infernal prison, a single breath of which, were it to
blow upon all contingent things, would return them to utter
nonexistence. He placed me in the hand of this oppressor who rose
against you in rejection and inflicted upon you that which changed the
names and their kingdom and severed the attributes from their seats.
When I became aware of him and of what I had committed by assisting in
what he committed, I fled from him and from meeting him to the court of
the glory of your oneness and the seat of the throne of your grandeur.

\subsection{480}\label{section-975}

\begin{otherlanguage}{arabic}

اذا فانظر يا الهى الىّ بلحظات عنايتك ثمّ ارتد بصر فضلك الى هذا المظلوم
الّذى صار من عمله قانطاً من روحك و عنايتك و مأيوساً عن بدايع فضلك و اكرامك
فآهء آه من عظم بلائى و كثرة حيرتى و احتراقى لم ادر ما اطلب من بدايع فضلك
لانّ كلّما يبلغ اليه اعلى مشعرى اشاهد بأنّه لايسكن فؤادى و لن يستريح به
قلبی‌ اذاً لمّا اجد نفسى فى هذا الحال و تلك الأحوال احبّ بأن افوّض امرى بيدك
و فى قبضتك لتقدّر ما هو خير عندك لنفسى و كينونتى و حقيقتى اذا اسئلك يا
محبوبی‌ بمظاهر امرك فى تلك الأیّام و مطالع وحيك و مخازن علمك بأن تنزل علىّ
ما يستضئ به وجهى بين السّموات و الأرض و انّك انت المقتدر على ما تشآء و انّك
انت المقتدر المهيمن العزيز القيّوم

\end{otherlanguage}

Idhan fa-unẓur yā ilāhī ilayya bi-laḥaẓāti ʿināyatika thumma irtadda
baṣara faḍlika ilā hādhā al-maẓlūmi alladhī ṣāra min ʿamalihi qāniṭan
min rawḥika wa-ʿināyatika wa-maʾyūsan ʿan badāʾiʿi faḍlika wa-ikrāmika.
Fa-āh āh min ʿiẓami balāʾī wa-kathrati ḥayratī wa-iḥtirāqī. Lam adri mā
aṭlubu min badāʾiʿi faḍlika li-anna kullamā yablughu ilayhi aʿlā
mashʿarī ushāhidu bi-annahu lā yusakkinu fuʾādī wa-lan yastarīḥa bihi
qalbī. Idhan lammā ajidu nafsī fī hādhā al-ḥāli wa-tilka al-aḥwāli
uḥibbu bi-an ufawwiḍa amrī bi-yadika wa-fī qabḍatika li-tuqaddira mā
huwa khayrun ʿindaka li-nafsī wa-kaynūnatī wa-ḥaqīqatī. Idhan asʾaluka
yā maḥbūbī bi-maẓāhiri amrika fī tilka al-ayyāmi wa-maṭāliʿi waḥyika
wa-makhāzini ʿilmika bi-an tunzila ʿalayya mā yastaḍīʾu bihi wajhī bayna
al-samāwāti wa-al-arḍi wa-innaka anta al-muqtadiru ʿalā mā tashāʾu
wa-innaka anta al-Muqtadiru al-Muhayminu al-ʿAzīzu al-Qayyūmu.

Look upon me, then, my God, with the glances of your care, and turn the
gaze of your grace toward this wronged one who, through his own deed,
has despaired of your solace and care and lost hope in the wonders of
your grace and generosity. Ah, ah, for the greatness of my affliction,
the depth of my bewilderment, and my burning anguish! I do not know what
to ask of the wonders of your grace, for I see that whatever my highest
perception can reach will not calm my soul or bring my heart to rest.
Since I find myself in this state and these conditions, I wish to
entrust my affair to your hand and into your grasp, so that you may
ordain what is best in your sight for my soul, my being, and my reality.
I ask you, then, my Beloved, by the manifestations of your Cause in
these days, the dawning places of your revelation, and the treasuries of
your knowledge, to send down upon me that which will make my face shine
amid the heavens and the earth. You have power over whatever you will,
and you are the All-Powerful, the All-Subduing, the Almighty, the
Self-Subsisting.

\subsection{481}\label{section-976}

\begin{otherlanguage}{arabic}

اذا يا محبوبی‌ اخذنى الاضطراب مرّةً اخرى من خطيئتى الكبرى يا من بيدك ملكوت
الامضا و جبروت القضا و كلما اسكن نفسى و احدثها من رحمتك الّتى سبقت
الممكنات و اذكرها بعنايتك الّتى احاطت من فى الأرض و السّمآء و اقول لها
اطمئنّى و لاتجزعى انّ محبوبی‌ رحيم و سلطانى كريم و مالكى عطوف و خالقى غفور
تظهر منها نار الحسرة و الخجلة و يحترق منها صبرى و صبرها و اصطبارى و
اصطبارها و سكونى و سكونها لذا لن ينقطع صريخى بين يديك و لاينتهى ضجيجى
تلقآء وجهك فو‌عزّتك اخاف بأن يحزن من حزنى سكّان جبروت سرورك و قبائل ملكوت
ابتهاجك اسئلهم بك بأن لايمنعونى عن صريخى و حنينى تلقآء مدين وحدانيّتك لانّ
الّذينهم يطوفنّ فى حول عرشك و شربوا كوثر العزّ و الأمال و وجدوا حلاوة القرب
و الوصال ينبغى لهم بأن ينظرن الّذى بعد عن ساحة قربك و جعل محروماً من بدايع
نعمائك

\end{otherlanguage}

Idhan yā maḥbūbī akhadhanī al-iḍṭirābu marratan ukhrā min khaṭīʾatī
al-kubrā yā man bi-yadika malakūtu al-imḍāʾi wa-jabarūtu al-qaḍāʾi
wa-kullamā usakkinu nafsī wa-uḥaddithuhā min raḥmatika allatī sabaqati
al-mumkināti wa-udhakkiruhā bi-ʿināyatika allatī aḥāṭat man fī al-arḍi
wa-al-samāʾi wa-aqūlu lahā iṭmaʾinnī wa-lā tajzaʿī inna maḥbūbī raḥīmun
wa-sulṭānī karīmun wa-mālikī ʿaṭūfun wa-khāliqī ghafūrun taẓharu minhā
nāru al-ḥasrati wa-al-khajlati wa-yaḥtariqu minhā ṣabrī wa-ṣabruhā
wa-iṣṭibārī wa-iṣṭibāruhā wa-sukūnī wa-sukūnuhā. Li-dhā lan yanqaṭiʿa
ṣarīkhī bayna yadayka wa-lā yantahī ḍajījī tilqāʾa wajhika.
Fa-wa-ʿizzatika akhāfu bi-an yaḥzana min ḥuznī sukkānu jabarūti surūrika
wa-qabāʾilu malakūti ibtihājika. Asʾaluhum bika bi-an lā yamnaʿūnī ʿan
ṣarīkhī wa-ḥanīnī tilqāʾa madīni waḥdāniyyatika li-anna alladhīnahum
yaṭūfanna fī ḥawli ʿarshika wa-sharibū kawthara al-ʿizzi wa-al-āmāli
wa-wajadū ḥalāwata al-qurbi wa-al-wiṣāli yanbaghī lahum bi-an yanẓurna
alladhī baʿuda ʿan sāḥati qurbika wa-juʿila maḥrūman min badāʾiʿi
naʿmāʾika.

Now, my Beloved, agitation has seized me once again because of my most
grievous sin, you in whose hand are the kingdom of fulfillment and the
dominion of decree. Whenever I calm my soul, speak to it of your mercy
that preceded all contingent things, remind it of your care that
encompasses everyone in earth and heaven, and say to it, ``Be at peace
and do not despair; my Beloved is merciful, my Sovereign is generous, my
Master is tenderhearted, and my Creator is forgiving,'' the fire of
regret and shame rises from it, consuming my patience and its patience,
my endurance and its endurance, my tranquility and its tranquility. Thus
my cry will never cease before you, nor will my clamor end before your
face. By your might, I fear that my sorrow will grieve the inhabitants
of the dominion of your joy and the tribes of the kingdom of your
gladness. I ask them, by you, not to prevent me from crying out and
lamenting before the city of your unity, for those who circle your
throne, have drunk the heavenly waters of glory and hope, and have
tasted the sweetness of nearness and reunion should look upon the one
who is far from the court of your presence and has been deprived of your
wondrous blessings.

\subsection{482}\label{section-977}

\begin{otherlanguage}{arabic}

فآهء آه يا محبوبی‌ فآهء آه يا مقصودى فآهء آه يا منائى فآهء آه يا رجائى
فآهء آه كيف ارفع رأسى عند قلم علىّ الّذى جعلته مشرق وحيك و مطّلع الهامك
فكيف انظر الى قلم الّذى ينسب الى محمّد رسولك و به رقم اسرار قضائك و كتب
الواح امرك فو‌عزّتك يا محبوبی‌ تكاد ان ارجع الى العدم من خجلتى و ما ورد علىّ
فى ایّامى هل اقدر ان اشاهد قلم الّذى جرى ما نزلته على الرّوح من مقادير امرك
و اسرار سننك و قضائك و ظهر منه ما ثبت به تنزيه ذاتك عن الامثال و تقديس
نفسك عن الاشباح او اسمع صرير قلم الكليم الّذى به كتبَ احكامك و حكمك و
جواهر توحيدك و آثار تفريدك

\end{otherlanguage}

Fa-āh āh yā maḥbūbī fa-āh āh yā maqṣūdī fa-āh āh yā munāʾī fa-āh āh yā
rajāʾī fa-āh āh kayfa arfaʿu raʾsī ʿinda qalami ʿAliyyin alladhī
jaʿaltahu mashriqa waḥyika wa-muṭṭalaʿa ilhāmika fa-kayfa anẓuru ilā
qalami alladhī yunsabu ilā Muḥammadin rasūlika wa-bihi ruqima asrāru
qaḍāʾika wa-kutiba alwāḥu amrika? Fa-wa-ʿizzatika yā maḥbūbī takādu an
arjiʿa ilā al-ʿadami min khajlatī wa-mā warada ʿalayya fī ayyāmī. Hal
aqdiru an ushāhida qalama alladhī jarā mā nazzaltahu ʿalā al-Rūḥi min
maqādīri amrika wa-asrāri sunanika wa-qaḍāʾika wa-ẓahara minhu mā
thabata bihi tanzīhu dhātika ʿani al-amthāli wa-taqdīsu nafsika ʿani
al-ashbāḥi aw asmaʿa ṣarīra qalami al-Kalīmi alladhī bihi kataba
aḥkāmaka wa-ḥikamaka wa-jawāhira tawḥīdika wa-āthāra tafrīdika?

Ah, ah, my Beloved! Ah, ah, my Goal! Ah, ah, my Desire! Ah, ah, my Hope!
Ah, ah! How can I raise my head before the pen of ʿAlī, which you made
the dawning place of your revelation and the rising place of your
inspiration? How can I look upon the pen associated with Muḥammad, your
Messenger, through which the mysteries of your decree were inscribed and
the tablets of your Cause were written? By your might, my Beloved, I am
on the verge of returning to nonexistence from my shame and from what
has befallen me in my days. Can I look upon the pen that set down what
you revealed to the Spirit of the ordinances of your Cause and the
mysteries of your ways and decree, and from which appeared that which
established that your essence is exalted above all likenesses and your
self sanctified above all forms? Or can I hear the scratching of the pen
of the one who spoke with you, through which he wrote your laws and your
wisdom, the gems of your unity and the signs of your uniqueness?

\subsection{483}\label{section-978}

\begin{otherlanguage}{persian}

اى محبوب كيف اشاهد قلم الخليل الّذى بعثته بأمرك و جعلته سراج وحيك بين
عبادك و كلمة امرك بين بريّتك فآهء آه من خجلتى الّتى لن تذهب منّى بدوام عزّك
و بقآء كينونتك هبنى يا الهى غمّستنى فى بحور الغفران فكيف تذهب منّى روايح
العصيان الّتى ما اتى به دونى فى ایّامك فيا الهى بحلمك الّذى ما اخذت به
اعدائك بل مددتهم بسلطان قضائك بحيث كلّ من اراد ان يجرّ على وجهك سيف
البغضآء انّك شحذتهُ لحكمة الّتى كانت مستورة عن انظر بريّتك و بعنايتك الّتى ما
امسكها شئونات خلقك و عمل المشركين فى بلادك بأن ترجعنى الى العدم ثمّ
ابعثنى فى ایّامك بسلطان مشيّتك مرّة اخرى لعلّ يمحى عن قلبی‌ ما ظهر منّى فى كرّة
الأولی و اكون غافلاً عما ارتكبت لتجعل هذه الغفلة نعمة على نفسى و رحمة على
كينونتى و عناية على ذاتى و انّك انت المقتدر المتعالى العزيز الكريم انتهى

\end{otherlanguage}

Ay maḥbūb, kayfa ushāhidu qalama al-Khalīl alladhī baʿathtahu bi-amrika
wa-jaʿaltahu sirāja waḥyika bayna ʿibādika wa-kalimata amrika bayna
bariyyatika? Fa-āh āh min khajlatī allatī lan tadhhaba minnī bi-dawāmi
ʿizzika wa-baqāʾi kaynūnatika. Habnī yā ilāhī ghammastanī fī buḥūri
al-ghufrān, fa-kayfa tadhhabu minnī rawāʾiḥu al-ʿiṣyān allatī mā atā
bihi dūnī fī ayyāmika? Fa-yā ilāhī, bi-ḥilmika alladhī mā akhadhta bihi
aʿdāʾaka bal madadtahum bi-sulṭāni qaḍāʾika bi-ḥaythu kullu man arāda an
yajurra ʿalā wajhika sayfa al-baghḍāʾ innaka shaḥadhtahu li-ḥikma allatī
kānat mastūra ʿan anẓur bariyyatika, wa-bi-ʿināyatika allatī mā amsakahā
shuʾūnātu khalqika wa-ʿamalu al-mushrikīn fī bilādika, bi-an turjiʿanī
ilā al-ʿadam thumma ibʿathnī fī ayyāmika bi-sulṭāni mashīʾatika marratan
ukhrā laʿalla yumḥā ʿan qalbī mā ẓahara minnī fī karrati al-ūlā wa-akūna
ghāfilan ʿammā irtakabtu li-tajʿala hādhihi al-ghaflata niʿmatan ʿalā
nafsī wa-raḥmatan ʿalā kaynūnatī wa-ʿināyatan ʿalā dhātī wa-innaka anta
al-Muqtadir al-Mutaʿālī al-ʿAzīz al-Karīm. Intahā.

O Beloved, how can I look upon the pen of the Friend whom you sent at
your command and made the lamp of your revelation among your servants
and the word of your command among your creatures? Ah, ah, for my shame,
which will never leave me as long as your might endures and your being
remains! Suppose, my God, that you immersed me in the oceans of
forgiveness: how could the odors of disobedience leave me, when no one
but me has committed such disobedience in your days? My God, by your
forbearance, through which you have not seized your enemies but have
instead sustained them through the sovereignty of your decree, so that
whenever anyone sought to draw the sword of hatred against your face,
you sharpened it for a wisdom hidden from the eyes of your creatures;
and by your care, which neither the affairs of your creation nor the
deeds of the idolaters in your lands have withheld, I ask you to return
me to nonexistence. Then raise me up once more in your days through the
sovereignty of your will, so that what proceeded from me in the first
existence may be erased from my heart and I may be unaware of what I
committed. Make this unawareness a blessing upon my soul, a mercy upon
my being, and a favor to my self. You are indeed the All-Powerful, the
Most Exalted, the Almighty, the Generous. End of quotation.

\subsection{484}\label{section-979}

\begin{otherlanguage}{persian}

فو‌الّذى نفسى بيده اگر كلّ من فى الابداع بقميص انصاف مزيّن شوند از ضجيج قلم
و صرير آن كه بين يدى الرّحمن نموده جميع از كلّ من فى السّموات و الأرض منقطع
شده بكوى دوست توجّه نمايند و از دنيا و عما قدّر فيها منقطع گردند و لكن اين
ایّام كه چنين نفوسى مشاهده نميشوند الّا قليل قليل ديگر تا چه وقت و زمان
وجوه مقدّسه منيره از جيب احديّة بيرون خرامند و بنفحات اللّه و كلماته فايز
شوند و از غيرش بی‌‌نياز آيند انّهم من خيرة الخلق عند اللّه الملك المهيمن
القيّوم

\end{otherlanguage}

Fa-wa-alladhī nafsī bi-yadihi, agar kullu man fī al-ibdāʿ bi-qamīṣ-i
inṣāf muzayyan shavand, az ḍajīj-i qalam va ṣarīr-i ān kih bayna yaday
al-Raḥmān nimūdih, jamīʿ az kullu man fī al-samāwāt wa-al-arḍ munqaṭiʿ
shudih bi-kūy-i dūst tavajjuh namāyand va az dunyā va ʿammā quddira fīhā
munqaṭiʿ gardand. Va lākin īn ayyām kih chunīn nufūsī mushāhidih
nimīshavand illā qalīl qalīl. Dīgar tā chih vaqt va zamān vujūh-i
muqaddasih-yi munīrih az jayb-i aḥadiyyat bīrūn khirāmand va
bi-nafaḥāt-i Allāh va kalimātash fāʾiz shavand va az ghayrash bī-niyāz
āyand. Innahum min khīrati al-khalq ʿinda Allāh al-Malik al-Muhaymin
al-Qayyūm.

By him in whose hand is my soul, if all who are in creation adorned
themselves with the garment of fairness, the lament of the pen and the
sound of its scratching before the All-Merciful would cause them all to
detach themselves from everyone in the heavens and on earth and turn
toward the dwelling of the Friend, and to detach themselves from the
world and whatever has been ordained within it. Yet in these days such
souls are nowhere to be seen, except for very, very few. At what time
and hour will radiant, sanctified faces emerge from the bosom of
oneness, receive the breaths of God and his words, and become
independent of all but him? They are among the choicest of creation in
the sight of God, the King, the All-Protecting, the Self-Subsisting.

\subsection{485}\label{section-980}

\begin{otherlanguage}{persian}

اگر چه اين عبد هر مطلبی‌ مذكور ميدارد بعد خود را نادم مشاهده مينمايد چه
كه معلوم نيست حرارت حبّ الهى در آن قلوب احداث شود چنانچه در ظهور ستّين
بيانات نقطه بيان روح ما سواه فداه را كه پيش هر نفسى ميبردند در كمال
بی‌‌اعتنائى بيك دست اخذ نموده مقدار عشره دقايق او اقلّ ملاحظه نموده
ميانداخت و قالوا ان هذه الكلمات مفتريات و اگر فى‌الجمله منصف بودند حرفى
از آن را بجميع من فى السّموات و الأرض مبادله نمينمودند بشنو از اين عبد و
قلب و سمع را از كلّ آنچه شنيده و ادراك نموده طاهر نما و بعد در آيات منزله
قبل كه بر نبيّين و مرسلين نازل شده ملاحظه كن و اين آيات ابدع امنع را
مشاهده نما و خود انصاف ده و البتّه اگر بآنچه ذكر شد عمل نمائى بحقّ فايز
شوى و از وسواس شياطين منزّه گردى

\end{otherlanguage}

Agar chih īn ʿabd har maṭlabī madhkūr mīdārad, baʿd khud rā nādim
mushāhidih mīnamāyad, chih kih maʿlūm nīst ḥarārat-i ḥubb-i ilāhī dar ān
qulūb iḥdāth shavad. Chunānchih dar ẓuhūr-i sittīn bayānāt-i Nuqṭiy-i
Bayān, rūḥ-i mā sivāhu fidāh, rā kih pīsh-i har nafsī mīburdand, dar
kamāl-i bī-iʿtināʾī bi-yik dast akhdh nimūdih miqdār-i ʿasharih daqāyiq
aw aqall mulāḥiẓih nimūdih mīandākht, wa-qālū inna hādhihi al-kalimāt
muftarayāt. Va agar fī al-jumlih munṣif būdand, ḥarfī az ān rā
bi-jamīʿ-i man fī al-samāwāt wa-al-arḍ mubādilih nimīnamūdand. Bishnaw
az īn ʿabd va qalb va samʿ rā az kull-i ānchih shinīdih va idrāk nimūdih
ṭāhir namā va baʿd dar āyāt-i munzalih-yi qabl kih bar nabiyyīn va
mursalīn nāzil shudih mulāḥiẓih kun va īn āyāt-i abdaʿ-i amnaʿ rā
mushāhidih namā va khud inṣāf dih. Va al-battih agar bi-ānchih dhikr
shud ʿamal namāʾī, bi-Ḥaqq fāʾiz shavī va az vasvās-i shayāṭīn munazzah
gardī.

Whatever this servant mentions, he afterward finds himself regretting
it, for it is uncertain whether the warmth of divine love will be
kindled in those hearts. Thus, in the Revelation of the year sixty,
whenever the utterances of the Point of the Bayān---may the souls of all
besides him be sacrificed for him---were brought before anyone, he would
take them in one hand with utter indifference, look at them for ten
minutes or less, and cast them aside. They said, ``These words are
fabrications.'' Had they been even somewhat fair, they would not have
exchanged a single letter of them for all who are in the heavens and on
earth. Listen to this servant, and purify your heart and hearing of
everything you have heard and understood. Then examine the verses
revealed before, which descended upon the prophets and messengers, and
behold these most wondrous, most exalted verses, and judge fairly for
yourself. If you act upon what has been stated, you will assuredly
attain the Truth and be freed from the whisperings of devils.

\subsection{486}\label{section-981}

\begin{otherlanguage}{persian}

اراده چنان بود كه جميع آيات منزله كه از سماء مشيّت ربّ علىّ اعلى در اين امر
ابدع امنع نازل شده ذكر شود نهى فرمودند و فرمودند چندى تأمّل نمائيد تا
معلوم شود كه اهل بيان در چه رتبه و مقام واقف و قائمند و جوهر آن اخذ شود
يعنى آن نفوسى كه بغير اللّه ناظر نبوده و نيستند و بعد اگر مقتضى شد آيات
مكنونه بيان و معانى مكنونه آن از سماء مشيّت رحمانى نازل تا جمعى ادراك
نمايند كه قلم اللّه حركت ننموده مگر بوصف اين ظهور اعظم و لسان اللّه ناطق
نشده مگر بذكر اين جمال اطهر الأمر بيد اللّه و انّه كان على كلّ شیء قدير

\end{otherlanguage}

Irādih chunān būd kih jamīʿ-i āyāt-i munzalih kih az samā-yi mashiyyat-i
Rabb-i ʿAlī-yi Aʿlā dar īn amr-i abdaʿ-i amnaʿ nāzil shudih dhikr
shavad. Nahy farmūdand va farmūdand chandī taʾammul namāyīd tā maʿlūm
shavad kih ahl-i Bayān dar chih rutbih va maqām vāqif va qāʾimand va
jawhar-i ān akhdh shavad, yaʿnī ān nufūsī kih bi-ghayr-i Allāh nāẓir
nabūdih va nīstand. Va baʿd agar muqtaḍī shud, āyāt-i maknūnih-yi Bayān
va maʿānī-yi maknūnih-yi ān az samā-yi mashiyyat-i raḥmānī nāzil tā
jamʿī idrāk namāyand kih Qalam-i Allāh ḥarakat nanimūdih magar bi-vaṣf-i
īn ẓuhūr-i aʿẓam va Lisān-i Allāh nāṭiq nashudih magar bi-dhikr-i īn
jamāl-i aṭhar. Al-amru bi-yadi Allāh wa-innahu kāna ʿalā kulli shayʾin
qadīr.

The intention was to cite all the revealed verses that have descended
from the heaven of the will of the Most High, the Most Exalted Lord
concerning this most wondrous, most exalted Cause. He forbade it,
saying, ``Wait a while, so that it may become clear at what rank and
station the people of the Bayān stand, and its essence may be
gathered---that is, those souls who have not looked, and do not look, to
anyone but God. Then, if it is fitting, the hidden verses of the Bayān
and its concealed meanings will descend from the heaven of the will of
the All-Merciful, so that some may understand that the Pen of God has
moved only to describe this Most Great Revelation, and the Tongue of God
has spoken only to mention this Most Pure Beauty.'' The matter is in
God's hand, and he has power over all things.

\subsection{487}\label{section-982}

\begin{otherlanguage}{persian}

لذا اين عبد فانى بآنچه از آيات مشرقه كه از مشرق فم سلطان بيان اشراق
نموده كفايت نمود و لكن بر هر ذى بصرى اليوم لازم كه قلب را از كلّ اشارات و
دلالات و اسماء مقدّس نمايد و بنفس ظهور و ما يظهر من عنده ناظر باشد و كلّ
باين مأمور شده‌اند در بيان چه كه اسماء و اذكار كلّ خلق او بوده و خواهند
بود اليوم اكثرى از اهل بيان در عقبه اسماء واقف و عند اللّه از عبده آن
مشهود و مذكور اهل منظر اكبر عند اللّه آن نفوسى هستند كه اثمار شجره الهيّه
را در اين ربيع عزّ فردانيّه بدست خود اخذ نموده و تناول نموده‌اند و متمسّكين
بغير اين شجره عند اللّه مردود و از اهل نفى و نار محسوب

\end{otherlanguage}

Lidhā īn ʿabd-i fānī bi-ānchih az āyāt-i mushriqih kih az mashriq-i
fam-i Sulṭān-i Bayān ishrāq nimūdih kifāyat nimūd. Va lākin bar har dhī
baṣarī al-yawm lāzim kih qalb rā az kull-i ishārāt va dalālāt va asmā
muqaddas namāyad va bi-nafs-i ẓuhūr va mā yaẓharu min ʿindihi nāẓir
bāshad. Va kull bi-īn maʾmūr shudih-and dar Bayān, chih kih asmā va
adhkār kull khalq-i ū būdih va khvāhand būd. Al-yawm aktharī az ahl-i
Bayān dar ʿaqabiy-i asmā vāqif va ʿinda Allāh az ʿabadi-yi ān mashhūd va
madhkūr. Ahl-i manẓar-i akbar ʿinda Allāh ān nufūsī hastand kih athmār-i
shajariy-i ilāhiyyih rā dar īn rabīʿ-i ʿizz-i fardāniyyih bi-dast-i khud
akhdh nimūdih va tanāvul nimūdih-and va mutamassikīn bi-ghayr-i īn
shajarih ʿinda Allāh mardūd va az ahl-i nafy va nār maḥsūb.

This mortal servant has therefore contented himself with those radiant
verses that have shone from the dawning place of the mouth of the
Sovereign of Utterance. Yet today everyone endowed with sight must
sanctify the heart from all allusions, indications, and names, and look
to the Revelation itself and whatever proceeds from him. All have been
commanded to do this in the Bayān, for names and invocations have all
been, and will remain, his creation. Today most of the people of the
Bayān have halted at the obstacle of names and are seen and counted
before God among their worshippers. In God's sight, the people of the
Most Great Vista are those souls who, in this springtime of the glory of
singleness, have gathered with their own hands the fruits of the divine
Tree and eaten them. Those who cling to anything other than this Tree
are rejected before God and numbered among the people of negation and
fire.

\subsection{488}\label{section-983}

\begin{otherlanguage}{persian}

در اين وقت روح ندا فرمود كه بنويس بعباد ظلمى فوق آن مشاهده نميشود كه از
ظهور حقّ نفسى دليل و برهان من غير ما ظهر من عنده و نزل من لدنه طلب نمايد
بحقّ ناظر شويد اگر آن حجّتى كه بآن اثبات حجّيت نقطه اوّليه شده مشاهده شد و
من لدى اللّه ظاهر گشت ديگر مجال توقّف نه در اين صورت هيچ عذرى از احدى
مسموع نبوده و نخواهد بود و لو بكلمات اوّلين و آخرين متمسّك شود و مستدل
گردد كذلك يبيّن اللّه امره ان انتم تفقهون

\end{otherlanguage}

Dar īn vaqt Rūḥ nidā farmūd kih binivīs bi-ʿibād: ẓulmī fawq-i ān
mushāhidih namīshavad kih az ẓuhūr-i Ḥaqq nafsī dalīl va burhān min
ghayri mā ẓahara min ʿindihi wa-nazala min ladunhi ṭalab namāyad.
Bi-Ḥaqq nāẓir shavīd. Agar ān ḥujjatī kih bi-ān ithbāt-i ḥujjiyyat-i
Nuqṭiy-i Avvaliyyih shudih mushāhidih shud va min ladā Allāh ẓāhir
gasht, dīgar majāl-i tavaqquf nah. Dar īn ṣūrat hīch ʿudhrī az aḥadī
masmūʿ nabūdih va nakhvāhad būd, va law bi-kalimāt-i avvalīn va ākhirīn
mutamassik shavad va mustadill gardad. Kadhālika yubayyinu Allāhu amrahu
in antum tafqahūn.

At this moment the Spirit called out, ``Write to the servants: No
injustice greater than this can be seen---that someone should ask the
Manifestation of the Truth for evidence and proof other than what has
appeared from him and descended from his presence. Look to the Truth. If
that proof by which the authority of the Primal Point was established is
seen and has appeared from God, there is no longer any room for
hesitation. In that case, no excuse from anyone has been or will be
accepted, even if he clings to the words of the first and the last and
adduces them as evidence. Thus does God explain his Cause, if you
understand.''

\subsection{489}\label{section-984}

\begin{otherlanguage}{persian}

و اينكه نوشته بودى ديّان و ميرزا غوغا و شيخ اسمعيل و حاجى ملّا هاشم ادّعا
نموده و باطل شده‌اند الخ مثل اين اعتراضات را در اوّل ظهور بيان مشركين
نموده‌اند و چنين مذكور ميداشتند كه انفس معدودات قبل از سيّد باب اين ادّعا
را نموده و باطل بودند لذا نعوذ باللّه اين قول هم مثل اقوال آن نفوس بوده و
باين اقوال اعراض از حقّ نمودند لازال نزد هر ظهورى مشركين آن عصر بامثال
اين اعتراضات مشغول شده‌اند و لكن عارف بصير و منقطع خبير را امثال اين
اعتراضات از صراط اللّه منع ننموده و ننمايد بلكه سبب رسوخ و ثبوت انفس
مستقيمه خواهد شد بسى واضح بوده كه هر چه ظهور اعظم‌تر و حقّ ظاهر‌تر اعراض
احرف نفى شديدتر بوده و خواهد بود چنانچه مشاهده ميشود و لكن اين نفوسى كه
ذكر نموده اولًا از كجا بطلان اين نفوس بر شما واضح شده

\end{otherlanguage}

Va īnkih nivishtih būdī Dayyān va Mīrzā Ghawghā va Shaykh Ismāʿīl va
Ḥājī Mullā Hāshim iddiʿā nimūdih va bāṭil shudih-and ilā ākhir, mithl-i
īn iʿtirāḍāt rā dar avval-i ẓuhūr-i Bayān mushrikīn nimūdih-and va
chunīn madhkūr mīdāshtand kih anfus-i maʿdūdāt qabl az Sayyid-i Bāb īn
iddiʿā rā nimūdih va bāṭil būdand, lidhā naʿūdhu bi-Allāh īn qawl ham
mithl-i aqvāl-i ān nufūs būdih, va bi-īn aqvāl iʿrāḍ az Ḥaqq nimūdand.
Lā yazāl nazd-i har ẓuhūrī mushrikīn-i ān ʿaṣr bi-amthāl-i īn iʿtirāḍāt
mashghūl shudih-and va lākin ʿārif-i baṣīr va munqaṭiʿ-i khabīr rā
amthāl-i īn iʿtirāḍāt az ṣirāṭ-i Allāh manʿ nanimūdih va nanamāyad,
balkih sabab-i rusūkh va thubūt-i anfus-i mustaqīmih khvāhad shud. Basī
vāḍiḥ būdih kih har chih ẓuhūr aʿẓamtar va Ḥaqq ẓāhirtar, iʿrāḍ-i
aḥruf-i nafy shadīdtar būdih va khvāhad būd, chunānchih mushāhidih
mīshavad. Va lākin īn nufūsī kih dhikr nimūdih, avvalan az kujā buṭlān-i
īn nufūs bar shumā vāḍiḥ shudih?

As for your statement that Dayyān, Mīrzā Ghawghā, Shaykh Ismāʿīl, and
Ḥājī Mullā Hāshim made claims and proved false, and so forth: the
idolaters raised objections like these at the beginning of the
Revelation of the Bayān. They used to say that a few people had made
this claim before the Sayyid, the Bāb, and had been false;
therefore---God forbid---this assertion too was like the assertions of
those people. With such words they turned away from the Truth. At every
Revelation, the idolaters of that age have always occupied themselves
with objections of this kind. Yet such objections have never kept, and
will never keep, an insightful knower or a discerning, detached soul
from the path of God. Rather, they will strengthen the firmness and
steadfastness of upright souls. It has been abundantly clear that the
greater the Revelation and the more manifest the Truth, the fiercer has
been, and will be, the rejection of the letters of negation, as can be
seen. But as for these souls you have mentioned: first, how has their
falsity become clear to you?

\subsection{490}\label{section-985}

\begin{otherlanguage}{persian}

امّا الدّيّان الّذى سمّى فى كتاب اللّه بكلّ اسمائه الحسنى ابدا ادّعائى ننموده كه
مخالف باشد در سنه اوّل امرش چند مناجاتى از او در ساحت اقدس ارسال داشتند و
در آن الواح جز اظهار خضوع و خشوع و ايمان باللّه و مظاهر امر او مشاهده نشد
و لكن نفس معروف چون مشاهده نمود كه نفسى ظاهر شده و مناجات نوشته نار حسدش
بشأنى مشتعل كه فتوى بر قتل او و محبّينش داده چنانچه آن مظهر احديّه را بنص
نقطه بيان بفتواى او شهيد نمودند و بعد در كتاب خود مفترياتى باو نسبت داده
كه ابدا صدق نبوده و فوق آنچه ادّعا نموده بود از منزل بيان تصريحا درباره
او نازل فانظر فى البيان لتكون من العارفين

\end{otherlanguage}

Ammā al-Dayyān alladhī summiya fī kitābi Allāh bi-kulli asmāʾihi
al-ḥusnā, abadan iddiʿāʾī nanimūdih kih mukhālif bāshad. Dar saniy-i
avval-i amrash chand munājātī az ū dar sāḥat-i aqdas irsāl dāshtand va
dar ān alvāḥ juz iẓhār-i khuḍūʿ va khushūʿ va īmān bi-Allāh va maẓāhir-i
amr-i ū mushāhidih nashud. Va lākin nafs-i maʿrūf chūn mushāhidih nimūd
kih nafsī ẓāhir shudih va munājāt nivishtih, nār-i ḥasadash bi-shaʾnī
mushtaʿil kih fatvā bar qatl-i ū va muḥibbīnash dādih, chunānchih ān
maẓhar-i aḥadiyyih rā bi-naṣṣ-i Nuqṭiy-i Bayān bi-fatvā-yi ū shahīd
nimūdand. Va baʿd dar kitāb-i khud muftarayātī bi-ū nisbat dādih kih
abadan ṣidq nabūdih va fawq-i ānchih iddiʿā nimūdih būd az Munzil-i
Bayān taṣrīḥan darbāriy-i ū nāzil. Fa-unẓur fī al-Bayān li-takūna min
al-ʿārifīn.

As for Dayyān, who was named in the Book of God by all his most
beautiful names, he never made any claim that was contrary. In the first
year of his cause, several prayers of his were sent to the Most Holy
Court, and in those tablets nothing was seen but expressions of
humility, reverence, and faith in God and the Manifestations of his
Cause. Yet when the well-known individual saw that someone had appeared
and written prayers, the fire of his envy blazed so fiercely that he
issued a ruling for the killing of him and his loved ones. Thus, through
his ruling, they martyred one who, according to the explicit text of the
Point of the Bayān, was a manifestation of oneness. Afterward, in his
own book, he attributed fabrications to him that had no truth
whatsoever, although statements explicitly concerning him, surpassing
what he had claimed, had been revealed by the Revealer of the Bayān.
Look, then, into the Bayān, so that you may be among those who know.

\subsection{491}\label{section-986}

\begin{otherlanguage}{persian}

و جمعى از اين نفوس بوده‌اند كه نفس معروف آنها را ردّ نموده و حكم بر كفر
جميع كرده و عند اللّه آن نفوس مردود نبوده و نخواهند بود بلكه شعرٌ منهم خير
منه لو انت من المتبصّرين از جمله جناب آقا سيّد ابراهيم عليه بهآء اللّه
الأبهی كه در حقّ او آياتى نازل شده و از جمله آن آيات اينست كه ذكر ميشود
شايد كه ادراك نمائى و مطّلع شوى كه احدى بنقطه اولى و ما نزل من عنده ناظر
نبوده بلكه حبّاً لرياسته و حفظاً لعزّه مرتكب شده آنچه را كه احدى از اوّل
ابداع تا حين مرتكب نشده قوله عزّ اعزازه ان اشهد ان يا ابراهيم انّك كنت فى
يوم عرش ظهور ربّك و انّا كنّا من قبل ثمّ من بعد لظاهرين انظر قد خلقناك و
رزقناك و امتناك و احييناك الى حينئذ و انّ الّذين اوتوا الصّحف هم الى حينئذ
محتجبون فلما نزلت على اللّه ربّك ربّ ما يرى و ما لايرى ربّ العالمين قد سمعت
صوت ما يتبعن امرك و هم يحسبون انّهم فى حبّك يتعاليون قل كلّا ثمّ کلّا انّنى
انا قد حشرت و من اتّبعنى على اللّه ربّی‌ فى يوم الّذى كنت بموسى عرش ظهور اللّه
من المؤمنين و انّ هؤلآء لا يتّبعون و ان اتّبعونى لأمنوا بموسى قبل عيسى ثمّ
بمحّمد بعد عيسى ثمّ بنقطة البيان يوم القيمة

\end{otherlanguage}

Va jamʿī az īn nufūs būdih-and kih nafs-i maʿrūf ānhā rā radd nimūdih va
ḥukm bar kufr-i jamīʿ kardih va ʿinda Allāh ān nufūs mardūd nabūdih va
nakhvāhand būd, balkih shaʿrun minhum khayrun minhu law anta min
al-mutabaṣṣirīn. Az jumlih Jināb-i Āqā Sayyid Ibrāhīm, ʿalayhi bahāʾu
Allāh al-Abhā, kih dar ḥaqq-i ū āyātī nāzil shudih va az jumliy-i ān
āyāt īnast kih dhikr mīshavad, shāyad kih idrāk namāʾī va muṭṭaliʿ shavī
kih aḥadī bi-Nuqṭiy-i Ūlā va mā nazala min ʿindihi nāẓir nabūdih, balkih
ḥubban li-riyāsatihi wa-ḥifẓan li-ʿizzihi murtakib shudih ānchih rā kih
aḥadī az avval-i ibdāʿ tā ḥīn murtakib nashudih. Qawluhu ʿazza iʿzāzuhu:
An ishhad an yā Ibrāhīm innaka kunta fī yawmin ʿarsha ẓuhūri rabbika
wa-innā kunnā min qablu thumma min baʿdu la-ẓāhirīn. Unẓur qad
khalaqnāka wa-razaqnāka wa-amatnāka wa-aḥyaynāka ilā ḥīniʾidhin wa-inna
alladhīna ūtū al-ṣuḥufa hum ilā ḥīniʾidhin muḥtajibūn. Fa-lammā nazalta
ʿalā Allāh rabbika rabbi mā yurā wa-mā lā yurā rabbi al-ʿālamīn qad
samiʿta ṣawta mā yatbaʿna amraka wa-hum yaḥsabūna annahum fī ḥubbika
yataʿālayūn. Qul kallā thumma kallā innanī anā qad ḥushirtu wa-man
ittabaʿanī ʿalā Allāh rabbī fī yawmi alladhī kuntu bi-Mūsā ʿarsha ẓuhūri
Allāh min al-muʾminīn. Wa-inna hāʾulāʾi lā yattabiʿūn wa-in ittabaʿūnī
la-āmanū bi-Mūsā qabla ʿĪsā thumma bi-Muḥammad baʿda ʿĪsā thumma
bi-Nuqṭati al-Bayān yawma al-qiyāma.

There have been a number of these souls whom the well-known individual
rejected, pronouncing them all unbelievers. Before God, those souls have
not been, and will not be, rejected. Indeed, a single hair of theirs is
better than him, if you are among those who see clearly. Among them is
the honored Āqā Sayyid Ibrāhīm---upon him be the glory of God, the
All-Glorious---concerning whom verses were revealed. Among those verses
is the passage cited here, so that perhaps you may understand and
realize that no one was looking to the Primal Point and what descended
from him. Rather, out of love for his leadership and to preserve his
standing, he committed what no one had committed from the beginning of
creation until now. His words---exalted be his glory---are: ``Bear
witness, O Ibrāhīm, that you were, on a day, the throne of the
revelation of your Lord, and we were indeed manifest before and
afterward. Behold, we have created you, sustained you, caused you to
die, and brought you to life up to this moment, while those who were
given the scriptures remain veiled to this moment. When you came before
God, your Lord, the Lord of what is seen and what is unseen, the Lord of
the worlds, you heard the voice of those who follow your cause, while
they suppose that they rise to exalted heights in their love for you.
Say: No, indeed, no! I and those who followed me have been gathered
before God, my Lord, on the day when, as Moses, I was the throne of
God's revelation, among the believers. These people do not follow; had
they followed me, they would have believed in Moses before Jesus, then
in Muḥammad after Jesus, then in the Point of the Bayān on the Day of
Resurrection.''

\subsection{492}\label{section-987}

\begin{otherlanguage}{persian}

حال ملاحظه نما كه او را از عرش ظهور ناميده ميفرمايند مشاهده نما خلق
نموديم و رزق داديم و ميرانديم و زنده نموديم تو را در اين هنگام و آن
نفوسيكه بتو ايمان آورده بودند در صحف تا حين محتجبند و شنيديم ندآء
نفوسىیكه متابعت مينمودند امر تو را و گمان مينمودند كه در محبّت تو بمقامات
عاليه فايز شده‌اند بگو كلّا ثمّ كلّا من محشور شدم بين يدى اللّه با نفوسيكه
متابعت من مينمودند در حينيكه ظاهر شد عرش ظهور باسم موسى و آن نفوس محتجبه
متابعت من نكردند و اگر متابعت من نموده بودند هرآينه بمظاهر بعد و بتصديق
و ايمان بآن انفس قدسيّه فايز ميشدند

\end{otherlanguage}

Ḥāl mulāḥiẓih namā kih ū rā az ʿarsh-i ẓuhūr nāmīdih mīfarmāyand:
Mushāhidih namā, khalq nimūdīm va rizq dādīm va mīrāndīm va zindih
nimūdīm tu rā dar īn hangām va ān nufūsī kih bi-tu īmān āvardih būdand
dar ṣuḥuf tā ḥīn muḥtajiband va shinīdīm nidā-yi nufūsī kih mutābiʿat
mīnimūdand amr-i tu rā va gumān mīnimūdand kih dar maḥabbat-i tu
bi-maqāmāt-i ʿāliyih fāʾiz shudih-and. Bigū kallā thumma kallā, man
maḥshūr shudam bayna yaday Allāh bā nufūsī kih mutābiʿat-i man
mīnimūdand dar ḥīnī kih ẓāhir shud ʿarsh-i ẓuhūr bi-ism-i Mūsā va ān
nufūs-i muḥtajibih mutābiʿat-i man nakardand va agar mutābiʿat-i man
nimūdih būdand har-āyinih bi-maẓāhir-i baʿd va bi-taṣdīq va īmān bi-ān
anfus-i qudsiyyih fāʾiz mīshudand.

Now consider how he has called him a throne of revelation, saying,
``Behold, we have created you, sustained you, caused you to die, and
brought you to life at this time. Those souls who had believed in you in
the scriptures remain veiled until now. We heard the call of the souls
who followed your cause and supposed that through their love for you
they had attained exalted stations. Say: No, indeed, no! I was gathered
before God, together with the souls who followed me, when the throne of
revelation appeared in the name of Moses. Those veiled souls did not
follow me. Had they followed me, they would assuredly have attained the
subsequent Manifestations and would have acknowledged and believed in
those holy souls.''

\subsection{493}\label{section-988}

\begin{otherlanguage}{persian}

بگو اى اصمّ روزگار درست متوجّه شو و التفات نما كلمات الهى را كه تصريحاً من
غير تلويح او را عرش ظهور ناميده و بعد از اين شأن محكم متعالى كلّ اسماء
طايف او بوده لو يخاطبه بانت اللّه او هو اللّه او من اللّه او الى اللّه او
ذات اللّه و كينونة اللّه جميع اينمراتب در رتبه او ثابت و محقّق و همين نفس
را اليوم اهل بيان بفتواى نفس معروف ابو الدّواهى مينامند و قسمى باو اعتراض
نموده‌اند و حكم كفرش در اطراف منتشر كه حال مدّتها است از كلّ كناره نموده
الا لعنة اللّه على الظّالمين

\end{otherlanguage}

Begū ey aṣamm-e rūzgār, dorost motavajjeh shov va eltefāt namā kalemāt-e
elāhī rā keh taṣrīḥan min ghayr-e talvīḥ ū rā ʿarsh-e ẓohūr nāmīdeh va
baʿd az īn shaʾn-e moḥkam-e motaʿālī koll-e asmāʾ ṭāyef-e ū būdeh. Law
yukhāṭibuhu bi-anta Allāh aw huwa Allāh aw min Allāh aw ilā Allāh aw
dhāt Allāh wa kaynūnat Allāh, jamīʿ-e īn marāteb dar rotbeh-ye ū thābet
va moḥaqqaq va hamīn nafs rā al-yawm ahl-e Bayān be-fatvā-ye nafs-e
maʿrūf Abū al-Davāhī mīnāmand va qesmī be-ū eʿterāż nemūdeh-and va
ḥokm-e kofrash dar aṭrāf montasher keh ḥāl moddat-hā ast az koll kenāreh
nemūdeh. A-lā laʿnatu Allāhi ʿalā al-ẓālimīn.

Say: O deaf one of this age! Pay close attention and give heed to the
words of God, which explicitly, without any allusion, have called him
the throne of revelation. Beyond this firm and exalted station, all
names circle around him. Whether he addresses him as ``You are God,''
``He is God,'' ``from God,'' ``to God,'' ``the Essence of God,'' or
``the Being of God,'' all these stations are established and confirmed
in his rank. Yet today the people of the Bayān call this very soul Abū
al-Davāhī, at the decree of the well-known individual. They have
assailed him so severely, and the verdict declaring him an unbeliever
has been so widely circulated, that he has now withdrawn from everyone
for a long time. Indeed, God's curse rests upon the oppressors.

\subsection{494}\label{section-989}

\begin{otherlanguage}{persian}

حال ملاحظه كن چه مقدار در امراللّه اسراف نموده‌اند و بغيا على اللّه چه
اعمال را مرتكب شده‌اند و بعد ميفرمايند مخاطباً لذلك الاسم الجليل قوله عزّ
اعزازه ان يا خليلى فى الصّحف لم يكن لاعراش ظهور اللّه من حدّ لا من قبل و لا
من بعد ولكنّ النّاس عن سرّ الأمر محتجبون و ان يا ذكرى فى الكتب من بعد الصّحف
لم يكن فى الاعراش الّا ما يدلّنّ على اللّه ربّهم قل كلّ من اللّه الى اللّه
يرجعون ان يا اسمى فى البيان انظر كيف ترقبن ادلّائى فى كلّ ظهور و الى حينئذ
ما فتحت باب الاسم فى ظهور من قبل هذا من فضل اللّه لمن فى البيان و لكن
النّاس لايعلمون

\end{otherlanguage}

Ḥāl molāḥeẓeh kon cheh meqdār dar Amr Allāh esrāf nemūdeh-and va baghyan
ʿalā Allāh cheh aʿmāl rā mortakeb shodeh-and va baʿd mīfarmāyand
mukhāṭiban li-dhālika al-ism al-jalīl, qawluhu ʿazza iʿzāzuhu: In yā
khalīlī fī al-ṣuḥuf, lam yakun li-aʿrāsh ẓuhūr Allāh min ḥadd, lā min
qabl wa lā min baʿd, wa-lākinna al-nāsa ʿan sirr al-amr muḥtajibūn.
Wa-in yā dhikrī fī al-kutub min baʿd al-ṣuḥuf, lam yakun fī al-aʿrāsh
illā mā yadullanna ʿalā Allāh rabbihim. Qul kullun min Allāh ilā Allāh
yarjiʿūn. In yā ismī fī al-Bayān, unẓur kayfa tarqubanna adillāʾī fī
kull ẓuhūr, wa-ilā ḥīniʾidhin mā fataḥtu bāb al-ism fī ẓuhūr min qabl.
Hādhā min faḍl Allāh li-man fī al-Bayān, wa-lākinna al-nāsa lā yaʿlamūn.

Now consider how far they have exceeded all bounds in the Cause of God,
and what deeds they have committed in rebellion against God. He then
says, addressing that glorious name---these are his words, exalted be
his glory: ``O my friend in the Scriptures! There has been no limit to
the thrones of God's revelation, either before or afterward, but the
people are veiled from the mystery of the Cause. And O my remembrance in
the Books after the Scriptures! There has been nothing in the thrones
but what points to God, their Lord. Say: All are from God and return to
God. O my name in the Bayān! Consider how you will watch for my guides
in every revelation. Until that time, I had not opened the gate of the
name in any previous revelation. This is a grace from God to those in
the Bayān, but the people do not know.''

\subsection{495}\label{section-990}

\begin{otherlanguage}{persian}

نفسى را كه تصريح ميفرمايد تو خليل منى در صحف و ذكر منى در كلّ كتب سماوى و
اسم من در بيان ابوالدّواهى مينامند و شاعر نيستند قل قد اخذ اللّه سمعكم و
ابصاركم و افئدتكم و انتم لا تسمعون و لا تبصرون و لا تفقهون كتاب مستيقظ
كه از تأليفات معرض باللّه است كه او را ربّ اخذ نموده بخوان و ببين باين
هياكل مقدّسه چه نسبتها داده و امر بقسمى صعب شد كه اين اسم وقتى در عراق
آمده بود جميع در صدد قتل او هم برآمدند و طلعت ابهى ایّام و ليالى ناس را
منع نمودند تا آنكه از آن ارض صحيحاً و سالماً بمحلّ خود مراجعت نمود كذلك كان
الأمر و كان اللّه على ما اقول شهيد حمد خدا را كه كتاب مستيقظش نزد هر نفسى
هست و الّا فو‌الّذى خلق البريّة لا من شیء كه انكار مينمود چنانچه اكثر اعمال
را انكار نموده

\end{otherlanguage}

Nafsī rā keh taṣrīḥ mīfarmāyad to khalīl-e manī dar ṣoḥof va dhekr-e
manī dar koll-e kotob-e samāvī va esm-e man dar Bayān, Abū al-Davāhī
mīnāmand va shāʿer nīstand. Qul qad akhadha Allāhu samʿakum wa abṣārakum
wa afʾidatakum wa antum lā tasmaʿūn wa lā tubṣirūn wa lā tafqahūn.
Ketāb-e Mostayqeẓ keh az taʾlīfāt-e moʿreż be-Allāh ast keh ū rā rabb
akhdh nemūdeh bekhān va bebīn be-īn hayākel-e moqaddaseh cheh nesbat-hā
dādeh va amr be-qesmī ṣaʿb shod keh īn esm vaqtī dar ʿErāq āmadeh būd
jamīʿ dar ṣadad-e qatl-e ū ham bar-āmadand va Ṭalʿat-e Abhā ayyām va
layālī nās rā manʿ nemūdand tā ānkeh az ān arż ṣaḥīḥan va sāleman
be-maḥall-e khod morājaʿat nemūd. Kadhālika kāna al-amr wa kāna Allāhu
ʿalā mā aqūlu shahīd. Ḥamd-e Khodā rā keh ketāb-e Mostayqeẓash nazd-e
har nafsī hast va ellā fa-wa-alladhī khalaqa al-bariyya lā min shayʾ keh
enkār mīnemūd chenāncheh akthar-e aʿmāl rā enkār nemūdeh.

The very soul to whom he explicitly says, ``You are my friend in the
Scriptures, my remembrance in all the heavenly Books, and my name in the
Bayān,'' they call Abū al-Davāhī, and they are unaware. Say: God has
taken away your hearing, your sight, and your hearts, and you neither
hear, nor see, nor understand. Read the book Mostayqeẓ, among the
writings of the one who has turned away from God and whom he has taken
as his lord, and see what accusations he has made against these holy
temples. Matters became so difficult that, when this name once came to
Iraq, they all sought to kill him as well. The All-Glorious Beauty
restrained the people day and night until he returned from that land to
his own home safe and sound. Such was the matter, and God is witness to
what I say. Praise be to God that everyone has his book Mostayqeẓ;
otherwise, by the One who created all beings out of nothing, he would
have denied it, just as he has denied most of his deeds.

\subsection{496}\label{section-991}

\begin{otherlanguage}{persian}

و اى كاش بشأن يكى از اين نفوسيكه مردود شمرده عند اللّه مذكور بود قد جعلت
الوهم لنفسك ربّاً من دون اللّه و لا تشعر ما تقول و لكن اللّه يشهد ما ارتكبتم
فى الحيوة الباطلة فسوف يأخذكم بقهر من عنده و انّه اشدّ المنتقمين مع‌ذلك
چنين نفوس تاركه را عامل بيان دانسته و نفسى را كه بيان بظهور او ظاهر شده
او را تارك شمرده فو‌اللّه تكلّم ننموده مگر بهوى و سالك نيستى الّا در تيه وهم
و غرور و تظنّ فى نفسك بأنّك من المهتدين

\end{otherlanguage}

Va ey kāsh be-shaʾn-e yekī az īn nofūsī keh mardūd shomordeh ʿinda Allāh
madhkūr būd. Qad jaʿalta al-wahma li-nafsika rabban min dūni Allāh wa lā
tashʿuru mā taqūl, wa-lākinna Allāha yashhadu mā irtakabtum fī al-ḥayāt
al-bāṭila fa-sawfa yaʾkhudhukum bi-qahrin min ʿindihi wa innahu ashaddu
al-muntaqimīn. Maʿa dhālek chenīn nofūs-e tārekeh rā ʿāmel-e Bayān
dānesteh va nafsī rā keh Bayān be-ẓohūr-e ū ẓāher shodeh ū rā tārek
shomordeh. Fa-wa-Allāh takallom nanemūdeh magar be-havā va sālek nīstī
ellā dar tīh-e vahm va ghorūr, wa taẓunnu fī nafsika bi-annaka min
al-muhtadīn.

Would that he were mentioned before God as being of the rank of even one
of these souls whom he has counted as rejected! You have made delusion a
lord for yourself instead of God, and you do not realize what you are
saying. But God witnesses what you have committed in this vain life, and
he will soon seize you with a might from himself. He is indeed the most
severe of avengers. Yet you have regarded such neglectful souls as
followers of the Bayān, while counting as one who has abandoned it the
very soul through whose revelation the Bayān appeared. By God, you have
spoken only out of desire, and you walk only in the wilderness of
delusion and conceit, while imagining yourself to be among the rightly
guided.

\subsection{497}\label{section-992}

\begin{otherlanguage}{persian}

و از اين گذشته اين نفوس بآيات الهى ظاهر نشده‌اند و اگر تو ميگوئى كه صاحب
آيات بودند و مردود شمرده هذا بغى منك على اللّه المهيمن القيّوم چه كه اگر
صاحب آيات باطل شود از براى تو و من على الأرض ايمان باقى نخواهد ماند چه
كه ايمان كلّ محقّق است بمظاهر امريّه الهى كه آن مظاهر امريّه بآيات ظاهر
شده‌اند و اگر نعوذ باللّه صاحب آيات باطل شود بچه برهان اثبات حقّ نقطه بيان
مينمائى كبرت كلمة تخرج من افواهكم يا ملأ المعرضين معلوم است كه اين نفوس
را محض تضييع امر اللّه مذكور داشته و اين ظهور را هم قياس بآن نفوس نموده‌اى
غافل بدان كه آن نفوس عند‌اللّه مردود نبوده و نخواهند بود بل المردود هو
الّذى انكرهم و اعترض عليهم و حكم بقتلهم

\end{otherlanguage}

Va az īn godhashteh īn nofūs be-āyāt-e elāhī ẓāher nashodeh-and va agar
to mīgūʾī keh ṣāḥeb-e āyāt būdand va mardūd shomordeh, hādhā baghyun
minka ʿalā Allāh al-Muhaymin al-Qayyūm. Cheh keh agar ṣāḥeb-e āyāt bāṭel
shavad az barā-ye to va man ʿalā al-arḍ īmān bāqī nakhāhad mānd, cheh
keh īmān-e koll moḥaqqaq ast be-maẓāher-e amrīyeh-ye elāhī keh ān
maẓāher-e amrīyeh be-āyāt ẓāher shodeh-and va agar naʿūdhu bi-Allāh
ṣāḥeb-e āyāt bāṭel shavad be-cheh borhān ethbāt-e ḥaqq-e Noqṭeh-ye Bayān
mīnamāʾī? Kaburat kalimatan takhruju min afwāhikum yā malaʾ al-muʿriḍīn.
Maʿlūm ast keh īn nofūs rā maḥż-e tażyīʿ-e Amr Allāh madhkūr dāshteh va
īn ẓohūr rā ham qiyās be-ān nofūs nemūdeh-ī. Ghāfel, bedān keh ān nofūs
ʿinda Allāh mardūd nabūdeh va nakhāhand būd, bal al-mardūd huwa alladhī
ankarahum wa iʿtaraḍa ʿalayhim wa ḥakama bi-qatlihim.

Moreover, these souls did not appear with divine verses. If you say that
they possessed verses and were counted as rejected, this is rebellion on
your part against God, the All-Protecting, the Self-Subsisting. For if
one who possesses verses is false, no faith will remain for you or
anyone on earth. The faith of all is established through the
Manifestations of God's Cause, and those Manifestations of the Cause
appeared with verses. If, God forbid, one who possesses verses is false,
by what proof will you establish the truth of the Point of the Bayān?
Grave is the word that comes from your mouths, O company of those who
turn away! It is clear that you have mentioned these souls solely to
disparage the Cause of God, and have also compared this revelation to
those souls. Heedless one! Know that those souls have not been, and will
not be, rejected before God. Rather, the rejected one is he who denied
them, assailed them, and decreed their death.

\subsection{498}\label{section-993}

\begin{otherlanguage}{persian}

و از جمله آن نفوس كه نوشته يكى جناب حاجى ملّا هاشم است و عرايضی كه از او
بمقرّ عرش آمده الآن موجود و جز خضوع كبرى از او امرى ظاهر نشده‌ اى مدّعى
عرفان همه عرفا را تضييع نمودى شنيده شد كه يكى از حكما وقتى در ارض صاد
بود و علماى آن ارض آن بيچاره فقير را تكفير نموده امر بخروج نمودند و جمعى
بر آن فقير هجوم آورده از جمله آن نفوس ملّا على حكيم نورى مشهور بود آن
مظلوم در آن حين بهيچيك از علماء سخنى نگفت و لكن بحكيم مذكور توجّه نموده و
گفت كلمه كه ذكرش محبوب نه و مقصودش آنكه تو خود را حكيم ميشمرى چرا بر
حكيم اعتراض مينمائى حال شما با اين اعتراضات ادّعاى حكمت و عرفان هم دارى
تب الى اللّه الّذى خلقك و سوّاك و لاتعترض على الّذينهم آمنوا باللّه و آياته و
ورد عليهم فى سبيله ما لايحصيه احد من العالمين تصديق و تكذيب نفوس ابدا
بتقليد كفايت نشده و نخواهد شد ربّ لاعن يرجع اللّعن الى نفسه و ربّ رادّ يرجع
الرّد اليه و هو لا يستشعر فى نفسه و يكون من الغافلين

\end{otherlanguage}

Va az jomleh-ye ān nofūs keh neveshteh yekī Janāb-e Ḥājī Mollā Hāshem
ast va ʿarāyeżī keh az ū be-maqarr-e ʿarsh āmadeh al-ān mawjūd va joz
khożūʿ-e kobrā az ū amrī ẓāher nashodeh. Ey moddaʿī-ye ʿerfān, hameh-ye
ʿorafā rā tażyīʿ nemūdī. Shenīdeh shod keh yekī az ḥokamā vaqtī dar
arż-e Ṣād būd va ʿolamā-ye ān arż ān bīchāreh-ye faqīr rā takfīr nemūdeh
amr be-khorūj nemūdand va jamʿī bar ān faqīr hojūm āvardeh, az jomleh-ye
ān nofūs Mollā ʿAlī Ḥakīm-e Nūrī-ye mashhūr būd. Ān maẓlūm dar ān ḥīn
be-hīch-yek az ʿolamā sokhanī nagoft va lāken be-ḥakīm-e madhkūr
tavajjoh nemūdeh va goft kalameh keh dhekrash maḥbūb nah va maqṣūdash
ānkeh to khod rā ḥakīm mīshomorī cherā bar ḥakīm eʿterāż mīnamāʾī. Ḥāl
shomā bā īn eʿterāżāt eddeʿā-ye ḥekmat va ʿerfān ham dārī. Tub ilā Allāh
alladhī khalaqaka wa sawwāka wa lā taʿtariḍ ʿalā alladhīna hum āmanū
bi-Allāh wa āyātihi wa warada ʿalayhim fī sabīlihi mā lā yuḥṣīhi aḥadun
min al-ʿālamīn. Taṣdīq va takdhīb-e nofūs abadan be-taqlīd kefāyat
nashodeh va nakhāhad shod. Rubba lāʿinin yarjiʿu al-laʿnu ilā nafsihi wa
rubba rāddin yarjiʿu al-raddu ilayhi wa huwa lā yastashʿiru fī nafsihi
wa yakūnu min al-ghāfilīn.

Among the souls he has mentioned is the honored Ḥājī Mollā Hāshem. The
petitions he sent to the seat of the throne still exist, and nothing but
the utmost humility has appeared from him. O claimant to spiritual
knowledge! You have disparaged all the mystics. It has been heard that
one of the philosophers was once in the land of Ṣād, and the religious
scholars of that land pronounced that poor, helpless man an unbeliever
and ordered him to leave. A group attacked that poor man; among them was
the well-known Mollā ʿAlī Ḥakīm of Nūr. At that moment the wronged man
said nothing to any of the religious scholars, but turned to the
philosopher just mentioned and said a word that it is not pleasing to
repeat. His meaning was: You consider yourself a philosopher; why do you
assail a philosopher? Now, with all these objections, you also claim
wisdom and spiritual knowledge! Turn in repentance to God, who created
and fashioned you, and do not assail those who have believed in God and
his verses, and who have suffered in his path what no one in all the
worlds can reckon. Affirming or denying souls on the basis of imitation
has never sufficed, nor will it ever suffice. Many a curse returns upon
the one who utters it, and many a rejection returns upon the one who
rejects, while he remains unaware of it within himself and is among the
heedless.

\subsection{499}\label{section-994}

\begin{otherlanguage}{persian}

قل موتوا بغيظكم قد كسر صنم الوهم و الهوى و اللّه الّذى لا اله الّا هو كه
بوهمى متشبّث شده كه در ابداع شبه نداشته نميدانم بچه دليل و برهان نفسى كه
آن نفوس را ردّ نموده ترجيح داده بشنو نصيحت اين عبد را و ابدا نفسى را ردّ
مكن چه كه مخصوص نقطه بيان جميع را وصيّت فرموده كه احدى بر احدى تعرّض
ننمايد چنانچه نفس معروف و اللّه الّذى لا اله الّا هو با آنكه با كمال حيله و
تدبير در صدد قطع سدره ربانيّه بر آمد و كشف شد آنچه در صدرش0پنهان بود
مع‌ذلك ايشان باو تعرّض نفرمودند تا آنكه اوراقى از او بين يدى حاضر شد و
جميع در ردّ اللّه و آياته بوده و الان بخطّ خود او موجود لذا آيات الهى در ردّ
من ردّ على اللّه از جبروت امر نازل و هذا لم يكن من عنده بل من لدى اللّه
المقتدر المهيمن العزيز القيّوم

\end{otherlanguage}

Qul mūtū bi-ghayẓikum qad kusira ṣanam al-wahm wa al-hawā. Wa-Allāh
alladhī lā ilāha illā huwa keh be-vahmī motashabbeth shodeh keh dar
ebdāʿ shebh nadāshteh. Namīdānam be-cheh dalīl va borhān nafsī keh ān
nofūs rā radd nemūdeh tarjīḥ dādeh. Beshnow naṣīḥat-e īn ʿabd rā va
abadan nafsī rā radd makon, cheh keh makhṣūṣ Noqṭeh-ye Bayān jamīʿ rā
vaṣīyat farmūdeh keh aḥadī bar aḥadī taʿarroż nanamāyad. Chenāncheh
nafs-e maʿrūf, wa-Allāh alladhī lā ilāha illā huwa, bā ānkeh bā kamāl-e
ḥīleh va tadbīr dar ṣadad-e qaṭʿ-e Sedreh-ye Rabbānīyeh bar-āmad va
kashf shod āncheh dar ṣadrash {[}corrupted text{]} penhān būd, maʿa
dhālek īshān be-ū taʿarroż nafarmūdand tā ānkeh awrāqī az ū bayna yaday
ḥāżer shod va jamīʿ dar radd Allāh va āyātih būdeh va al-ān be-khaṭṭ-e
khod-e ū mawjūd. Ledhā āyāt-e elāhī dar radd-e man radda ʿalā Allāh az
jabarūt-e amr nāzel, wa hādhā lam yakun min ʿindihi bal min ladā Allāh
al-Muqtadir al-Muhaymin al-ʿAzīz al-Qayyūm.

Say: Die in your rage! The idol of delusion and desire has been
shattered. By God, besides whom there is no god, he has clung to a
delusion without parallel in creation. I do not know by what evidence or
proof you have preferred the person who rejected those souls. Hear the
counsel of this servant and never reject anyone, for the Point of the
Bayān specifically enjoined everyone not to assail anyone. Thus,
although the well-known individual---by God, besides whom there is no
god---sought with the utmost cunning and calculation to cut down the
divine Lote Tree, and what lay {[}corrupted original text{]} hidden in
his heart was disclosed, he nevertheless did not assail him until some
papers from him were brought into his presence. They were all in
rejection of God and his verses, and they still exist in his own
handwriting. Thus divine verses descended from the dominion of the Cause
in rejection of the one who rejected God. This was not from himself, but
from God, the All-Powerful, the All-Protecting, the Almighty, the
Self-Subsisting.

\subsection{500}\label{section-995}

\begin{otherlanguage}{persian}

و آيات مختص است بصاحب ظهور چنانچه نقطه بيان تصريح فرموده كه بعد از غروب
شمس حقيقت احدى استماع آيات اللّه نخواهد نمود الّا حين ظهور بعد و همين نفس
معروف ليالى و اسحار بين يدى قائم و آيات الهى بر او القاء ميشد و بعد
بمجرّد دو كلمه چنان مغرور شده كه ارتكاب نمود آنچه را كه احدى از مظاهر نفى
ارتكاب ننموده و هر نفسيكه فى‌الجمله صاحب شعور باشد ميداند كه او لازال
خدمت حضرت ابهى بوده و اگر هم چيزى ادراك نموده رشحى از طمطام اين بحر موّاج
باو افاضه شده مع‌ذلك تو نوشته آنچه را كه هيچ عاقلى ننوشته فو‌اللّه اگر
صد‌هزار امثال او بين يدى قادر بر تكلم باشند چنانچه اين عباد بچشم خود
مشاهده نموده‌اند و تو بوهم صرف ميگوئى بارى لعنت كند خدا آن شيطانى كه او و
تو هر دو را اغوا نموده و براى رياست ظاهره چه مقدار مفتريات بحقّ نسبت داده
انشآءالله اميدواريم كه بگندم رى مرزوق نشود

\end{otherlanguage}

Va āyāt mokhtaṣṣ ast be-ṣāḥeb-e ẓohūr chenāncheh Noqṭeh-ye Bayān taṣrīḥ
farmūdeh keh baʿd az ghorūb-e Shams-e Ḥaqīqat aḥadī estemāʿ-e āyāt Allāh
nakhāhad nemūd ellā ḥīn-e ẓohūr-e baʿd va hamīn nafs-e maʿrūf layālī va
asḥār bayna yaday qāʾem va āyāt-e elāhī bar ū elqā mīshod va baʿd
be-mojarrad-e do kalameh chenān maghrūr shodeh keh ertekāb nemūd āncheh
rā keh aḥadī az maẓāher-e nafy ertekāb nanemūdeh va har nafsī keh fī
al-jomleh ṣāḥeb-e shoʿūr bāshad mīdānad keh ū lā-zāl khedmat-e Ḥażrat-e
Abhā būdeh va agar ham chīzī edrāk nemūdeh rashḥī az ṭamṭām-e īn baḥr-e
mavvāj be-ū efāżeh shodeh. Maʿa dhālek to neveshteh āncheh rā keh hīch
ʿāqelī naneveshteh. Fa-wa-Allāh agar ṣad hezār amthāl-e ū bayna yaday
qāder bar takallom bāshand, chenāncheh īn ʿebād be-chashm-e khod
moshāhedeh nemūdeh-and va to be-vahm-e ṣerf mīgūʾī. Bārī laʿnat konad
Khodā ān shayṭānī keh ū va to har do rā eghvā nemūdeh va barā-ye
rīyāsat-e ẓāhereh cheh meqdār moftarayāt be-Ḥaqq nesbat dādeh. In shāʾ
Allāh omīdvārīm keh be-gandom-e Ray marzūq nashavad.

Verses belong exclusively to the bearer of a revelation, as the Point of
the Bayān has explicitly stated: after the setting of the Sun of Truth,
no one will hear the verses of God until the next revelation. This same
well-known individual would stand in his presence through the nights and
at dawn, and divine verses were imparted to him. Afterward, on the
strength of merely two words, he became so conceited that he committed
what none of the manifestations of negation had committed. Anyone with
even a little understanding knows that he was always in the service of
the All-Glorious One, and if he understood anything at all, it was a
sprinkling bestowed upon him from the depths of this surging ocean. Yet
you have written what no rational person has written. By God, if a
hundred thousand like him could speak in his presence! These servants
have witnessed it with their own eyes, whereas you speak from sheer
delusion. In any case, may God curse that satan who has led both him and
you astray, and who has attributed so many fabrications to the True One
for the sake of outward leadership. God willing, we hope that he will
not be granted the wheat of Ray.

\subsection{501}\label{section-996}

\begin{otherlanguage}{persian}

و ديگر نميدانم كه ميرزا غوغا چه ارتكاب نموده كه باين قسم بغضش در قلوب
امثال او جا گرفته بشأنى كه در رسائل در ردّ او مينويسند بسا از نفوسند كه
در غلبات ذوق و شوق كلماتى ميگويند ايشانهم سخنى ذكر نموده‌اند ديگر نبايد
باين شدّت در صدد هتك حرمت عباد باشند بترسيد از خدا و قلوب عباد او را
ميازاريد هيچ امرى از امثال شما ظاهر نشده كه فى‌الجمله رايحه حقّ از او
استشمام شود نسئل اللّه بأن يجعل صدوركم خالياً عن غلّ احبّائه و انّه كان على
كلّ شیء قدير

\end{otherlanguage}

Va dīgar namīdānam keh Mīrzā Ghawghā cheh ertekāb nemūdeh keh be-īn qesm
boghżash dar qolūb-e amthāl-e ū jā gerefteh be-shaʾnī keh dar rasāʾel
dar radd-e ū mīnevīsand. Basā az nofūsand keh dar ghalabāt-e dhowq va
showq kalemātī mīgūyand, īshān ham sokhanī dhekr nemūdeh-and. Dīgar
nabāyad be-īn sheddat dar ṣadad-e hatk-e ḥormat-e ʿebād bāshand.
Betarsīd az Khodā va qolūb-e ʿebād-e ū rā mayāzārīd. Hīch amrī az
amthāl-e shomā ẓāher nashodeh keh fī al-jomleh rāyeḥeh-ye Ḥaqq az ū
esteshmām shavad. Nasʾalu Allāha bi-an yajʿala ṣudūrakum khāliyan ʿan
ghilli aḥibbāʾihi wa innahu kāna ʿalā kull shayʾ qadīr.

Furthermore, I do not know what Mīrzā Ghawghā has done that such hatred
of him has taken root in the hearts of people like him, to the point
that they write treatises against him. Many souls utter words when
overcome by spiritual delight and longing; he too has said something.
They should not seek so fiercely to dishonor God's servants. Fear God
and do not wound the hearts of his servants. Nothing has appeared from
people like you from which even a trace of the fragrance of the True One
can be perceived. We ask God to free your hearts from rancor toward his
loved ones. He is indeed powerful over all things.

\subsection{502}\label{section-997}

\begin{otherlanguage}{persian}

قول الملقى كبر اعراضه و عظم نفاقه و اگر مدّعى اينمقام حجّتى غير از آيات
دارند بفرمائيد و اگر آيات است ميرزا يحيى علاوه از نصّ صريح و براهين قاطعه
در مقابل ايستاده ادّعاى بالاتر و بيشتر و بهتر دارد و مجلّدات كثيره در دست
دارد انتهى

\end{otherlanguage}

Qawl al-mulqī, kabura iʿrāḍuhu wa ʿaẓuma nifāquhu: Va agar moddaʿī-ye īn
maqām ḥojjatī ghayr az āyāt dārand befarmāʾīd va agar āyāt ast Mīrzā
Yaḥyā ʿelāveh az naṣṣ-e ṣarīḥ va barāhīn-e qāṭeʿeh dar moqābel īstādeh
eddeʿā-ye bālātar va bīshtar va behtar dārad va mojalladāt-e kathīreh
dar dast dārad. Intahā.

The words of the instigator---great is his aversion and immense his
hypocrisy: ``If the claimant to this station has a proof other than
verses, please state it. If it is verses, Mīrzā Yaḥyā, in addition to an
explicit text and decisive proofs, stands opposed with a higher,
greater, and better claim, and has many volumes in hand.'' End of
quotation.

\subsection{503}\label{section-998}

\begin{otherlanguage}{persian}

فو‌اللّه بكلماتى تكلّم نموده كه انسان شرم ميدارد كه ملاحظه نمايد چه زود
ميزان اللّه را تغيير داده‌ايد و حرمة اللّه را ضايع نموده‌ايد و چه زود بتحريف
كلماتش مشغول گشته‌ايد فو‌اللّه هر ذى بصرى از آنچه نوشته شأنك و شأن ملقيك را
ادراك نموده و مينمايد فويل له و للّذين اتّبعوه اين عبد متحيّر است كه كدام
يك از اين مزخرفاترا استماع نمايد و جواب ذكر كند آخر اين چهار روز عمر چه
تو را بآن داشته كه باين هذيانات وهميّه نفسيّه بالقاى شيطان مشغول شده و از
شاطى بحر احديّه محروم مانده گيرم جميع نفوس او را من دون اللّه ربّ اخذ
نمودند چه نفعى از براى تو و او حاصل لا فو‌الّذى نفسى بيده بل خسران دنيا و
آخرت

\end{otherlanguage}

Fa-wa-Allāh be-kalemātī takallom nemūdeh keh ensān sharm mīdārad keh
molāḥeẓeh namāyad. Cheh zūd mīzān Allāh rā taghyīr dādeh-īd va ḥormat
Allāh rā żāyeʿ nemūdeh-īd va cheh zūd be-taḥrīf-e kalemātash mashghūl
gashteh-īd. Fa-wa-Allāh har dhī-baṣarī az āncheh neveshteh shaʾnaka wa
shaʾna mulqīka rā edrāk nemūdeh va mīnamāyad. Fa-waylun lahu wa
li-alladhīna ittabaʿūhu. Īn ʿabd motaḥayyer ast keh kodām yek az īn
mozakhrafāt rā estemāʿ namāyad va javāb dhekr konad. Ākhar īn chahār rūz
ʿomr cheh to rā be-ān dāshteh keh be-īn hadhayānāt-e vahmīyeh-ye
nafsīyeh be-elqā-ye shayṭān mashghūl shodeh va az shāṭī-ye baḥr-e
aḥadīyeh maḥrūm māndeh? Gīram jamīʿ-e nofūs ū rā min dūn Allāh rabb
akhdh nemūdand, cheh nafʿī az barā-ye to va ū ḥāṣel? Lā fa-wa-alladhī
nafsī bi-yadihi, bal khosrān-e donyā va ākherat.

By God, he has spoken words that one is ashamed to read. How quickly you
have changed God's standard and violated his sanctity, and how quickly
you have set about distorting his words! By God, everyone with
discernment has perceived, and continues to perceive, your standing and
that of your instigator from what you have written. Woe to him and to
those who have followed him! This servant is bewildered as to which of
these absurdities to hear and answer. In the end, during these few brief
days of life, what has driven you to occupy yourself with these
fanciful, self-centered ravings at Satan's prompting, and to remain
deprived of the shore of the ocean of divine unity? Suppose that all
souls took him as a lord instead of God: what benefit would come to you
and him? None, by the One in whose hand is my soul! Rather, it would be
loss in this world and the next.

\subsection{504}\label{section-999}

\begin{otherlanguage}{persian}

نوشته كه صاحب اين ظهور غير از آيات اگر چيزى دارد بياورد و اين تصريحا
مخالف است بآنچه نقطه بيان روح ما سواه فداه در كلّ بيان نازل فرموده و من
تكلّم بهذه الكلمة او يتكلّم لعن و يلعنه كلّ الذّرّات و من فى ملكوت الأمر و
الخلق بلى و اللّه ناس را احمق دانسته‌ايد چه كه اگر ناس احمق نبودند تو و
نفس معرض قادر نبوديد بر اينكه جهرا بغير ما تكلّم به اللّه تكلّم نمائيد نقطه
بيان در كلّ بيان تصريحا فرموده كه حجّت ظهور بعد غير آيات نبوده و نخواهد
بود و تو تصريحا نوشته كه اگر حجّتى غير از آيات دارد اتيان نمايد و تقول ما
لا تشعران

\end{otherlanguage}

Neveshte ke ṣāḥeb-e īn ẓohūr gheyr az āyāt agar chīzī dārad biyāvarad va
īn taṣrīḥan mokhālef ast bā ānche Noqṭe-ye Bayān rūḥ-e mā sevāh fedāh
dar koll-e Bayān nāzel farmūde va man takallama bi-hādhihi al-kalima aw
yatakallam luʿina wa yalʿanuhu kull al-dharrāt wa man fī malakūt al-amr
wa al-khalq. Balā va Allāh nās rā aḥmaq dāneste-īd che ke agar nās aḥmaq
nabūdand to va nafs-e moʿreż qāder nabūdīd bar īnke jahran be-gheyr-e mā
takallama bihi Allāh takallom namāyīd. Noqṭe-ye Bayān dar koll-e Bayān
taṣrīḥan farmūde ke ḥojjat-e ẓohūr-e baʿd gheyr-e āyāt nabūde va
nakhāhad būd va to taṣrīḥan neveshte ke agar ḥojjatī gheyr az āyāt dārad
etyān namāyad wa taqūlu mā lā tashʿurān.

You have written that if the bearer of this revelation has anything
other than verses, he should produce it. This explicitly contradicts
what the Point of the Bayān---may the spirit of all else be sacrificed
for him---has revealed throughout the Bayān. Whoever has spoken or
speaks these words has been cursed and is cursed by every atom and
everyone in the kingdoms of command and creation. Yes, by God, you have
taken people for fools, for if people were not fools, you and that soul
who has turned away would not have been able to speak openly in terms
other than those God has spoken. The Point of the Bayān has explicitly
stated throughout the Bayān that the proof of the next revelation has
been and will be nothing other than verses, yet you have explicitly
written that if he has a proof other than verses, he should produce it.
You say what neither of you understands.

\subsection{505}\label{section-1000}

\begin{otherlanguage}{persian}

مقامى را كه اراده نموده كه از براى معرض باللّه ثابت نمائى آن مقام منوط
بقول نقطه اولى بوده و نقطه اولى بآيات الهى ثابت شده حال تفكّر كن كه چه
قدر از سبيل حقّ دور مانده‌ايد و در تيه ضلالت افتاده‌ايد تفكّر فيما ارتكبت و
تكلّمت به لعلّ تتنبّه من فعلك و تكون من المستغفرين الى اللّه الّذى خلق كلّ شیء
بأمره و اظهر مظهر نفسه بسلطانه و جعل ذيله مقدّساً منك و من ظنونك و من
اوهامك ثمّ اوهام الأصنام الّذين اتّخذتموها ربّاً من دون اللّه فويل لهم من عذاب
يوم عقيم

\end{otherlanguage}

Maqāmī rā ke erāde namūde ke az barā-ye moʿreż be-Allāh thābet namāyī ān
maqām manūṭ be-qowl-e Noqṭe-ye Ūlā būde va Noqṭe-ye Ūlā be-āyāt-e elāhī
thābet shode. Ḥāl tafakkor kon ke che qadr az sabīl-e Ḥaqq dūr mānde-īd
va dar tīh-e żalālat oftāde-īd. Tafakkar fīmā irtakabta wa takallamta
bihi laʿalla tatanabbahu min fiʿlika wa takūnu min al-mustaghfirīn ilā
Allāh alladhī khalaqa kulla shayʾ bi-amrihi wa aẓhara maẓhara nafsihi
bi-sulṭānihi wa jaʿala dhaylahu muqaddasan minka wa min ẓunūnika wa min
awhāmika thumma awhām al-aṣnām alladhīna ittakhadhtumūhā rabban min dūni
Allāh fa-waylun lahum min ʿadhābi yawmin ʿaqīm.

The station you have sought to establish for the one who has turned away
from God depended on the words of the Primal Point, and the Primal Point
was established by divine verses. Now reflect on how far you have
strayed from the path of the True One and fallen into the wilderness of
error. Reflect on what you have done and said, so that you may become
aware of your deed and be among those who seek forgiveness from God, who
created all things by his command, brought forth the Manifestation of
his own Self through his sovereignty, and kept the hem of his robe
sanctified from you, your suppositions, and your imaginings, and from
the imaginings of the idols you have taken as a lord instead of God. Woe
to them for the punishment of a barren day!

\subsection{506}\label{section-1001}

\begin{otherlanguage}{persian}

و اللّه از حمقاى بيان مطمئنّيد و حقّ داريد حال آيات را كنار ميگذاريم آنچه
نفس معروف ظاهر نمايد اگر فوق آن از منبع قدرت الهيّه ظاهر نشد تسليم
مجعولات تو و آن نفس مشرك باللّه را مينمائيم فو‌اللّه در اين امر هم ثابت
نيستيد و نخواهيد بود چه كه نفس معروف را ميشناسيم در دعوى و ادّعا اكبر من
كلّ كبير و در ظهور و فعل اصغر من كلّ صغير اى فقير بی‌بضاعت كاش اقلّا يك مثل
خودى را ربّ اخذ نموده بودى چه فايده كه بسيار غافل مشاهده ميشوى و الّا بعضى
مطالب مذكور ميداشتم تا وهم را خوب ادراك نمائى و بر مكرش مطّلع گردى

\end{otherlanguage}

Va Allāh az ḥomaqā-ye Bayān moṭmaʾennīd va ḥaqq dārīd. Ḥāl āyāt rā kenār
mīgoẕārīm ānche nafs-e maʿrūf ẓāher namāyad agar fowq-e ān az manbaʿ-e
qodrat-e elāhiyye ẓāher nashod taslīm-e majʿūlāt-e to va ān nafs-e
moshrek be-Allāh rā mīnamāyīm. Fa-wa-Allāh dar īn amr ham thābet nīstīd
va nakhāhīd būd che ke nafs-e maʿrūf rā mīshenāsīm dar daʿvī va eddeʿā
akbar min kull kabīr va dar ẓohūr va feʿl aṣghar min kull ṣaghīr. Ey
faqīr-e bī-beżāʿat kāsh aqallan yek methl-e khodī rā rabb akhẕ namūde
būdī. Che fāyede ke besyār ghāfel moshāhede mīshavī va ellā baʿżī
maṭāleb maẕkūr mīdāshtam tā vahm rā khūb edrāk namāyī va bar makrash
moṭṭaleʿ gardī.

By God, you feel secure in the fools of the Bayān, and you have reason
to. Now we will set verses aside: whatever that well-known soul may
produce, if something surpassing it does not appear from the source of
divine power, we will submit to your fabrications and those of that soul
who associates partners with God. Yet, by God, you do not stand firm in
this matter either, nor will you, for we know that well-known soul: in
claims and pretensions, greater than anyone great, and in manifestation
and action, smaller than anyone small. O destitute pauper, if only you
had at least taken someone like yourself as a lord! But what use is it?
You appear utterly heedless. Otherwise, I would have mentioned certain
matters so that you might fully comprehend this illusion and become
aware of his deceit.

\subsection{507}\label{section-1002}

\begin{otherlanguage}{persian}

و اينكه نوشته اگر آيات است ميرزا يحيى در مقابل ايستاده و ادّعاى بالاتر و
بيشتر و بهتر دارد و مجلّدات كثيره در دست دارد ما صغر شأنه و كبر قوله اين
ادّعا از آن منبع غرور و نفس و هوى مثل آن است كه نمله وادى جرزه ادّعاى
ربوبيّت نمايد بلكه اعظم عند كلّ ذى بصر منير چه كه كينونت او بحرفى از اين
آيات عزّ صمدانى خلق شده و آنچه كلمات از او ظاهر شده در رتبه او بوده فافتح
بصرك لتعرف و ما عنده لن يذكر عند اللّه ابداً الّا بأن يرجع اليه و يتوب و
يكون من التّائبين و المستغفرين فو‌اللّه ارتكب ما لا ارتكبه احد قبله

\end{otherlanguage}

Va īnke neveshte agar āyāt ast Mīrzā Yaḥyā dar moqābel īstāde va
eddeʿā-ye bālātar va bīshtar va behtar dārad va mojalladāt-e kathīre dar
dast dārad mā aṣghara shaʾnahu wa akbara qawlahu. Īn eddeʿā az ān
manbaʿ-e ghorūr va nafs va havā methl-e ān ast ke namle-ye vādī-ye
joroze eddeʿā-ye robūbiyyat namāyad balke aʿẓam ʿinda kull dhī baṣar
munīr che ke keynūnat-e ū be-ḥarfī az īn āyāt-e ʿezz-e ṣamadānī khalq
shode va ānche kalemāt az ū ẓāher shode dar rotbe-ye ū būde fa-ftaḥ
baṣaraka li-taʿrif wa mā ʿindahu lan yudhkara ʿinda Allāh abadan illā
bi-an yarjiʿa ilayhi wa yatūba wa yakūna min al-tāʾibīn wa
al-mustaghfirīn fa-wa-Allāh irtakaba mā lā irtakabahu aḥad qablahu.

As for your statement, ``If it is verses, Mīrzā Yaḥyā stands opposed
with a higher, greater, and better claim, and has many volumes in
hand''---how low his standing, and how grand his words! This claim from
that source of pride, self, and desire is like an ant in a barren valley
claiming lordship; indeed, to everyone with illumined sight it is even
greater, for his very being was created through a single letter of these
verses of eternal glory, and whatever words have issued from him have
been at his own level. Open your eyes, then, so that you may recognize
this. What he possesses will never be mentioned before God unless he
returns to him, repents, and becomes one of those who repent and seek
forgiveness. By God, he has committed what no one before him has
committed.

\subsection{508}\label{section-1003}

\begin{otherlanguage}{persian}

حال خود تو انصاف ده كه اين كلمه از چه رو ظاهر شده كه نوشته او در مقابل
ايستاده و ادّعاى بالاتر و بيشتر و بهتر دارد هر جاهلى از اين قول ننگ داشته
تا چه رسد بعاقل اى بی‌ بصر همينكه نوشته كه در مقابل ايستاده همين دليلى
است واضح بر بطلان او صاحبان شامّه مقدّسه از نفس همين كلمات روايح غرور و
استكبار و اعراض از حقّ استشمام نمايند گويا ابداً تلاوت فرقانهم ننموده كه
ميفرمايد قوله تعالى و من اظلم ممن افترى على اللّه كذباً او قال اوحى الىّ و
لم يوح اليه شیء و من قال سانزل بمثل ما انزل اللّه انتهى قدرى تفكّر در اين
آيه منزله نما كه شايد موقن شوى باينكه هر نفسى بگويد نازل ميكنم مثل آنچه
حقّ نازل فرموده همين كلمه مبطل و مكذب او بوده و خواهد بود

\end{otherlanguage}

Ḥāl khod-e to enṣāf deh ke īn kaleme az che rū ẓāher shode ke neveshte ū
dar moqābel īstāde va eddeʿā-ye bālātar va bīshtar va behtar dārad. Har
jāhelī az īn qowl nang dāshte tā che resad be-ʿāqel. Ey bī-baṣar hamīnke
neveshte ke dar moqābel īstāde hamīn dalīlī ast vāżeḥ bar boṭlān-e ū.
Ṣāḥebān-e shāmme-ye moqaddase az nafs-e hamīn kalemāt ravāyeḥ-e ghorūr
va estekbār va eʿrāż az Ḥaqq esteshmām namāyand. Gūyā abadan telāvat-e
Forqān ham nanamūde ke mīfarmāyad qawluhu taʿālā wa man aẓlamu mimmani
iftarā ʿalā Allāhi kadhiban aw qāla ūḥiya ilayya wa lam yūḥa ilayhi
shayʾun wa man qāla sa-unzilu mithla mā anzala Allāh intahā. Qadrī
tafakkor dar īn āye-ye monzale namā ke shāyad mūqen shavī be-īnke har
nafsī begūyad nāzel mīkonam methl-e ānche Ḥaqq nāzel farmūde hamīn
kaleme mobṭel va mokaẕẕeb-e ū būde va khāhad būd.

Now judge fairly for yourself: what prompted these words you have
written, that he stands opposed with a higher, greater, and better
claim? Any ignorant person would be ashamed of this statement, let alone
a person of understanding. O sightless one, the very fact that you have
written that he stands opposed is itself a clear proof of his falsity.
Those with a sanctified sense of smell perceive in these very words the
odors of pride, arrogance, and turning away from the True One. It seems
you have never even read the Qurʾān, which declares---exalted are his
words: ``Who is more unjust than one who invents a lie against God, or
says, `It has been revealed to me,' when nothing has been revealed to
him, or one who says, `I will send down the like of what God has sent
down'?'' End of quotation. Reflect a little on this revealed verse, so
that you may become certain that whenever anyone says, ``I will send
down the like of what the True One has sent down,'' this very statement
has been and will be what invalidates his claim and proves him false.

\subsection{509}\label{section-1004}

\begin{otherlanguage}{persian}

اى اصمّ حقيقى بشنو نغمه ورقاى الهى را كه بر اعلى سدره بيان تغنّى فرموده
ميفرمايد كه اگر تو در اين حين ظاهر شوى من از اوّل ساجدين و اوّل مؤمنين بتو
بوده و خواهم بود و تو حيا و شرم از اين بيان لسان رحمن ننموده و تصريحا
نوشته كه در مقابل ايستاده و بالاتر و بيشتر و بهتر ادّعا دارد و اللّه همين
كلمات برهانيست آشكار بر اعراض و كفر و نفاق امثال آن نفوس و هم‌چنين
ميفرمايد قوله عزّ ذكره اگر كلّ اهل بيان در جوهر علم مثل او گردند كه ثمر
نمى‌بخشد الّا بتصديق باو و اين بسى واضح است كه احدى در جوهر علم مثل او
ممكن نه الّا نفس نقطه آن شمس الهى درباره خود چنين ميفرمايد و تو از براى
نفسى كه ميخواهى بقول آن شمس حقيقت امرش را مجعولاً ثابت نمائى بذكر اينگونه
كلمات نالايقه مشغول شده فوا حسرتا لك و له و للّذين اتّبعو كما تاللّه ان هم
الّا على ضلال مبين

\end{otherlanguage}

Ey aṣamm-e ḥaqīqī beshnow naghme-ye Varqā-ye Elāhī rā ke bar aʿlā-ye
Sedre-ye Bayān taghannī farmūde mīfarmāyad ke agar to dar īn ḥīn ẓāher
shavī man az avval-e sājedīn va avval-e moʾmenīn be-to būde va khāham
būd. Va to ḥayā va sharm az īn bayān-e Lesān-e Raḥmān nanamūde va
taṣrīḥan neveshte ke dar moqābel īstāde va bālātar va bīshtar va behtar
eddeʿā dārad. Va Allāh hamīn kalemāt borhānīst āshkār bar eʿrāż va kofr
va nefāq-e amthāl-e ān nofūs. Va hamchenīn mīfarmāyad qawluhu ʿazza
dhikruhu agar koll-e ahl-e Bayān dar jowhar-e ʿelm methl-e ū gardand ke
thamar namībakhshad ellā be-taṣdīq be-ū va īn basī vāżeḥ ast ke aḥadī
dar jowhar-e ʿelm methl-e ū momken na ellā nafs-e Noqṭe. Ān Shams-e
Elāhī darbāre-ye khod chenīn mīfarmāyad va to az barā-ye nafsī ke
mīkhāhī be-qowl-e ān Shams-e Ḥaqīqat amrash rā majʿūlan thābet namāyī
be-ẕekr-e īngūne kalemāt-e nālāyeqe mashghūl shode fa-wā ḥasratā laka wa
lahu wa li-lladhīna ittabaʿūkumā ta-Allāh in hum illā ʿalā ḍalālin
mubīn.

O truly deaf one, hear the song of the divine Dove, who has sung upon
the highest branches of the Tree of the Bayān, saying: ``If you were to
appear at this moment, I have been and will be the first to bow down and
the first to believe in you.'' Yet you have felt no shame or
embarrassment before this utterance of the Tongue of the All-Merciful,
and have explicitly written that he stands opposed with a higher,
greater, and better claim. By God, these very words are a manifest proof
of the aversion, unbelief, and hypocrisy of such souls. He also
says---exalted be his mention: ``If all the people of the Bayān were to
become like him in the essence of knowledge, it would bear no fruit
except through acknowledgment of him. And it is perfectly clear that no
one can be like him in the essence of knowledge except the Point
himself.'' That divine Sun speaks thus of himself, while you have
occupied yourself with such unworthy statements on behalf of a soul
whose cause you wish falsely to establish through the words of that Sun
of Truth. Alas for you, for him, and for those who have followed you
both! By God, they are in nothing but manifest error.

\subsection{510}\label{section-1005}

\begin{otherlanguage}{persian}

و از اين گذشته مينويسى كه مجلّدات كثيره در دست دارد بعينه اين قول مثل
اقوال معرضين نزد هر ظهور است كه افتخار بتاليفات و تصنيفات خود نموده و بر
سلطان اسماء و صفات اعتراض نموده‌اند اگر كتب نفسى را غنى مينمود هرآينه كتب
روى ارض در حين ظهور نقطه بيان روح ما سواه فداه ناس را غنى و مستغنى
مينمود لا و اللّه كلّ فقراء لدى باب غنائه و كلّ جهلاء لدى باب علمه و كلّ
عجزاء عند ظهور قدرته و انّه لهو الغنىّ بنفسه عمّا سواه و انّه لهو المقتدر
على العالمين

\end{otherlanguage}

Va az īn goẕashte mīnevīsī ke mojalladāt-e kathīre dar dast dārad.
Be-ʿeyneh īn qowl methl-e aqvāl-e moʿreżīn nazd-e har ẓohūr ast ke
eftekhār be-taʾlīfāt va taṣnīfāt-e khod namūde va bar Solṭān-e Asmā va
Ṣefāt eʿterāż namūde-and. Agar kotob nafsī rā ghanī mīnamūd harāyene
kotob-e rū-ye arż dar ḥīn-e ẓohūr-e Noqṭe-ye Bayān rūḥ-e mā sevāh fedāh
nās rā ghanī va mostaghnī mīnamūd. Lā wa-Allāh kull fuqarāʾ ladā bābi
ghināʾihi wa kull juhalāʾ ladā bābi ʿilmihi wa kull ʿujazāʾ ʿinda ẓuhūri
qudratihi wa innahu la-huwa al-ghaniyy bi-nafsihi ʿammā siwāhu wa innahu
la-huwa al-muqtadir ʿalā al-ʿālamīn.

Besides this, you write that he has many volumes in hand. This statement
is exactly like the statements of those who turn away at every
revelation, who have taken pride in their own writings and compositions
and objected to the Sovereign of names and attributes. If books could
make a soul self-sufficient, then surely the books on the face of the
earth at the time of the revelation of the Point of the Bayān---may the
spirit of all else be sacrificed for him---would have made people
self-sufficient and free of need. No, by God! All are poor at the gate
of his wealth, all are ignorant at the gate of his knowledge, and all
are powerless before the manifestation of his power. He, in his own
Self, is indeed independent of all else, and he is indeed the
All-Powerful over all worlds.

\subsection{511}\label{section-1006}

\begin{otherlanguage}{persian}

از خود تو انصاف ميطلبم اگر نزد نفسى يكقطعه از ياقوت باشد و نزد نفسى
صد‌هزار خروار حصاة حال كدام را غنىّ ميدانى و بی‌نياز ميشمرى اگر صدهزار كتب
ظاهر شود كه ثمرى نبخشيده و نخواهد بخشيد الّا بتصديق بظهور چه كه كلمات
بمنزله مشكوة و حبّ و ذكر الهى بمنزله سراج اگر آن مشكوة از سراج محروم ماند
آلتى خواهد بود معطّل كلمات خوب است در وقتيكه نفحات حبّ الهى از او استشمام
شود و بعد از فقدان اين رتبه مثل ادنى كلمات اهل ارض ملاحظه ميشود خف عن
اللّه و لا تقابل معه احد من عباده و لا تذكر غيره فى ساحة قدسه فاخرق
الأحجاب ثمّ انطق بثنآء ربّك بين العالمين

\end{otherlanguage}

Az khod-e to enṣāf mīṭalabam agar nazd-e nafsī yek qeṭʿe az yāqūt bāshad
va nazd-e nafsī ṣad hezār kharvār ḥaṣāt ḥāl kodām rā ghanī mīdānī va
bī-niyāz mīshomorī. Agar ṣad hezār kotob ẓāher shavad ke thamarī
nabakhshīde va nakhāhad bakhshīd ellā be-taṣdīq be-ẓohūr che ke kalemāt
be-manzale-ye meshkāt va ḥobb va ẕekr-e elāhī be-manzale-ye serāj. Agar
ān meshkāt az serāj maḥrūm mānad ālatī khāhad būd moʿaṭṭal. Kalemāt khūb
ast dar vaqtīke nafaḥāt-e ḥobb-e elāhī az ū esteshmām shavad va baʿd az
foqdān-e īn rotbe methl-e adnā kalemāt-e ahl-e arż molāḥeẓe mīshavad.
Khaf ʿan Allāh wa lā tuqābil maʿahu aḥadan min ʿibādihi wa lā tadhkur
ghayrahu fī sāḥati qudsihi fa-khriq al-aḥjāb thumma inṭiq bi-thanāʾi
rabbika bayna al-ʿālamīn.

I ask you yourself to judge fairly: if one person possesses a single
ruby and another a hundred thousand donkey-loads of pebbles, which do
you consider rich and free of need? Even if a hundred thousand books
were to appear, they have borne and will bear no fruit except through
acknowledgment of the revelation, for words are like a niche, and the
love and remembrance of God are like a lamp. If that niche remains
without the lamp, it will be a useless object. Words are good when the
fragrances of divine love can be perceived in them; when they lack this
quality, they are regarded as the lowliest words of the people of the
earth. Fear God, set none of his servants against him, and mention no
one other than him in the court of his holiness. Tear through the veils,
then speak the praise of your Lord among all peoples.

\subsection{512}\label{section-1007}

\begin{otherlanguage}{persian}

در اينمقام بيان منزل بيان را استماع نما كه بيكى از حروفات حىّ وصيّت فرموده
ميفرمايد قوله عزّ ذكره اگر يك آيه از آيات آن جوهر قدم را تلاوت كنى يا ثبت
كنى اعزّ تر خواهد بود عند اللّه از آنكه كلّ بيان را ثبت كنى زيرا كه آن روز
آن يك آيه تو را نجات ميدهد ولى كلّ بيان نميدهد مثل آنكه اگر كسى در زمان
رسول خدا كلّ كتب سماويّه و آنچه در ظلّ آنها نوشته شده ثبت مينمود حكم لاشیء
در حقّ او ميشد ولى اگر يك آيه از قرآن ثبت مينمود حكم ايمان در حقّ او ميشد
اينست جوهر علم و عمل انتهى

\end{otherlanguage}

Dar īn maqām bayān-e Monzel-e Bayān rā estemāʿ namā ke be-yekī az
Ḥorūfāt-e Ḥayy vaṣiyyat farmūde mīfarmāyad qawluhu ʿazza dhikruhu agar
yek āye az āyāt-e ān Jowhar-e Qedam rā telāvat konī yā thabt konī
aʿazz-tar khāhad būd ʿinda Allāh az ānke koll-e Bayān rā thabt konī zīrā
ke ān rūz ān yek āye to rā nejāt mīdahad valī koll-e Bayān namīdahad
methl-e ānke agar kasī dar zamān-e Rasūl-e Khodā koll-e kotob-e
samāviyye va ānche dar ẓell-e ānhā neveshte shode thabt mīnamūd ḥokm-e
lā-shayʾ dar ḥaqq-e ū mīshod valī agar yek āye az Qorʾān thabt mīnamūd
ḥokm-e īmān dar ḥaqq-e ū mīshod īnast jowhar-e ʿelm va ʿamal intahā.

Here listen to the words of the Revealer of the Bayān, who gave this
counsel to one of the Letters of the Living. He says---exalted be his
mention: ``If you recite or write down a single verse from the verses of
that Essence of Eternity, it will be more precious in the sight of God
than if you were to write down the whole Bayān, for on that day that
single verse will save you, while the whole Bayān will not. Similarly,
if someone in the time of the Messenger of God had written down all the
heavenly books and everything written under their shadow, he would have
been judged as nothing; but if he had written down a single verse of the
Qurʾān, he would have been judged a believer. This is the essence of
knowledge and action.'' End of quotation.

\subsection{513}\label{section-1008}

\begin{otherlanguage}{persian}

حال ملاحظه كن كه عظمت امر بچه مقدار بلند و عظيم است و حقارت تو و كلماتت
و كلمات مشركين در چه رتبه و مقام شمس مشرقه باين حجبات محجوب نماند و نفس
ظهور بآلايش اين كلمات نيالايد انّه قد كان مقدّساً منكم و من كلماتكم و انّه
بنفسه حجّة اللّه بين بريّته و ينبغى لمن على الأرض بأن يتّبعوا ما يظهر من
عنده لا بما كان عندهم و انّ هذا لهو الحقّ و ما بعد الحقّ الّا الضّلال ان انتم
من العارفين

\end{otherlanguage}

Ḥāl molāḥeẓe kon ke ʿaẓamat-e amr be-che meqdār boland va ʿaẓīm ast va
ḥaqārat-e to va kalemātat va kalemāt-e moshrekīn dar che rotbe va maqām.
Shams-e moshreqe be-īn ḥojobāt maḥjūb namānad va nafs-e ẓohūr
be-ālāyesh-e īn kalemāt nayālāyad. Innahu qad kāna muqaddasan minkum wa
min kalimātikum wa innahu bi-nafsihi ḥujjat Allāh bayna bariyyatihi wa
yanbaghī li-man ʿalā al-arḍ bi-an yattabiʿū mā yaẓharu min ʿindihi lā
bi-mā kāna ʿindahum wa inna hādhā la-huwa al-ḥaqq wa mā baʿda al-ḥaqq
illā al-ḍalāl in antum min al-ʿārifīn.

Now consider how lofty and immense is the greatness of the Cause, and
how low the station and standing of you, your words, and the words of
those who associate partners with God. The shining Sun will not remain
concealed by these veils, nor will the Manifestation himself be defiled
by the stain of these words. He has ever been sanctified from you and
your words, and he himself is God's proof among his creatures. Those on
earth should follow what appears from him, not what they themselves
possess. This is indeed the truth, and beyond the truth there is nothing
but error, if you are among those who understand.

\subsection{514}\label{section-1009}

\begin{otherlanguage}{persian}

بسا از حمير كه مجلّدات زياد حمل نموده و لكن خود بحرفى از آن شاعر نه
چنانچه اين مثل در كتب قبل ذكر شده و در كتاب مبين تصريحاً نازل و اللّه الّذى
لا اله الّا هو مجلّداتى كه ميگوئى از براى اين خوبست كه به آب محو شود او
يرجعها احد الى مبدئه لانّ ينبغى ان يكون بظنّه قبرا لكلماته اين عباد بسيار
از آن مجلّدات ديده‌اند ايكاش صنم وهم را ميشكستى و بديده بصيرت بكلمات
مجعوله‌اش نظر مينمودى تاللّه يضحك به و بما ظهر من عنده كلّ عاقل بصير

\end{otherlanguage}

Basā az ḥamīr ke mojalladāt-e ziyād ḥaml namūde va lāken khod be-ḥarfī
az ān shāʿer na chenānche īn mathal dar kotob-e qabl ẕekr shode va dar
Ketāb-e Mobīn taṣrīḥan nāzel. Va Allāh alladhī lā ilāha illā huwa
mojalladātī ke mīgūʾī az barā-ye īn khūb ast ke be-āb maḥv shavad aw
yurjiʿahā aḥad ilā mabdaʾihi li-anna yanbaghī an yakūna bi-ẓannihi
qabran li-kalimātihi. Īn ʿebād besyār az ān mojalladāt dīde-and. Ey kāsh
ṣanam-e vahm rā mīshekastī va be-dīde-ye baṣīrat be-kalemāt-e
majʿūle-ash naẓar mīnamūdī. Ta-Allāh yaḍḥaku bihi wa bi-mā ẓahara min
ʿindihi kull ʿāqil baṣīr.

Many a donkey has carried numerous volumes without itself understanding
a single letter of them, as this likeness has been mentioned in earlier
books and explicitly revealed in the perspicuous Book. By God, besides
whom there is no god, the volumes you speak of are fit only to be washed
away with water, or for someone to return them to their source, for in
his own estimation it ought to be a grave for his words. These servants
have seen many of those volumes. If only you would shatter the idol of
illusion and look with the eye of insight at his fabricated words! By
God, every discerning person of understanding laughs at him and at what
has issued from him.

\subsection{515}\label{section-1010}

\begin{otherlanguage}{persian}

و ديگر از آن منبع كبر و غرور القاء شده بود قوله حقر قدره سهل است ادّعاى
آن دارند كه كلّ آيات مدّعى را كه در مدّت عمر نوشته در يك روز خواهم نوشت
بهتر و خوبتر و فصيح تر كه هر طفلى تميز دهد انتهى

\end{otherlanguage}

Va dīgar az ān manbaʿ-e kebr va ghorūr elqā shode būd qawluhu ḥaqura
qadruhu sahl ast eddeʿā-ye ān dārand ke koll-e āyāt-e moddaʿī rā ke dar
moddat-e ʿomr neveshte dar yek rūz khāham nevesht behtar va khūbtar va
faṣīḥ-tar ke har ṭeflī tamīz dahad intahā.

Another statement had been put forward by that source of arrogance and
pride---his words, abased be his standing: ``That is easy. He claims, `I
will write in a single day all the verses the claimant has written
during his lifetime, better, finer, and more eloquent, so that any child
can tell the difference.'\,'' End of quotation.

\subsection{516}\label{section-1011}

\begin{otherlanguage}{persian}

جميع فضيلتهاى شما و او از اين يك كلمه معلوم شده اين عبد متحيّر است كه در
اين مقام چه ذكر نمايد چه كه بكلمه تنطّق كرده كه و اللّه فرعون هم در مقابل
كليم بچنين كلمه تكلّم ننموده بيكتا خداوند بی‌مانند كه اگر مقابل يك نفر از
طايفين حول قادر بر تكلّم باشد تا چه رسد بمنبع و مظهر و مصدر فيوضات ربانيّه
و حال همين قول را تجربه نمائيد آن معرض حاضر و خدّام بيت هم حاضر اگر
توانست بشأنى از شئون با خدّام مقابله نمايد نصدّقكم فى كلّ ما قلتم مع آنكه
زياد گفتن دليل بر حقّیّت نبوده و نخواهد بود چون تو اعتراضا على اللّه باين
دلايل وهميّه مستدلّ شده لذا اين عبد در جواب باين كلمات ذاكر لئلّا يبقى لك
من عذر و حجّة و برهان و دليل و الّا ساحت اقدس اطهر از جميع اين دلايل و
اذكار و كلمات و اشارات و مدلوليكه بآن تمسّك جسته‌ايد مقدّس بوده و خواهد بود

\end{otherlanguage}

Jamīʿ-i faḍīlat-hā-yi shumā va ū az īn yak kalameh maʿlūm shudeh. Īn
ʿabd mutaḥayyir ast keh dar īn maqām cheh dhikr namāyad, cheh keh
bi-kalameh tanaṭṭuq kardeh keh va-Allāh Firʿawn ham dar muqābil-i Kalīm
bi-chunīn kalameh takallum nanamūdeh. Bi-yiktā Khudāvand-i bī-mānand keh
agar muqābil-i yak nafar az ṭāʾifīn-i ḥawl qādir bar takallum bāshad, tā
cheh rasad bi-manbaʿ va maẓhar va maṣdar-i fuyūḍāt-i rabbānīyeh. Va ḥāl
hamīn qawl rā tajrebeh namāyīd. Ān muʿriḍ ḥāḍir va khuddām-i bayt ham
ḥāḍir. Agar tavānest bi-shaʾnī az shuʾūn bā khuddām muqābeleh namāyad,
nuṣaddiqukum fī kulli mā qultum. Maʿ-i ānkeh ziyād guftan dalīl bar
ḥaqqīyat nabūdeh va nakhvāhad būd. Chūn tu iʿtirāḍan ʿalā Allāh bi-īn
dalāʾil-i vahmīyeh mustadill shudeh, lidhā īn ʿabd dar javāb bi-īn
kalamāt dhākir, li-allā yabqā laka min ʿudhrin va ḥujjatin va burhānin
va dalīl. Va illā sāḥat-i aqdas aṭhar az jamīʿ-i īn dalāʾil va adhkār va
kalamāt va ishārāt va madlūlī keh bi-ān tamassuk justeh-īd muqaddas
būdeh va khvāhad būd.

All your virtues and his have become evident from this one statement.
This servant is bewildered as to what to say here, for he has uttered
words such that, by God, even Pharaoh never spoke their like before
Moses. By the one, incomparable God, if only he could speak before a
single one of those who circle around, let alone before the source,
manifestation, and fountainhead of divine outpourings! Now put this very
assertion to the test. That opponent is present, and the servants of the
House are also present. If he can confront the servants in any respect,
we will acknowledge you in everything you have said. Yet speaking at
length has never been, and will never be, a proof of truth. Since you
have adduced these imaginary proofs in objection to God, this servant
speaks these words in reply, so that no excuse, testimony,
demonstration, or proof may remain to you. Otherwise, the most holy,
most pure court has always been, and will always remain, sanctified
above all these proofs, references, words, allusions, and the meaning to
which you have clung.

\subsection{517}\label{section-1012}

\begin{otherlanguage}{persian}

اى بی‌خبر تا حال معادل آنچه از سماء بيان نازل شده نازل و بيك آيه معادله
نمينمايد كلّ آنچه در آسمانها و زمين است و اگر تو باين شئون ناظرى پس حاضر
شو بين يدى تا ملاحظه نمائى كه حين نزول آيات احدى از عهده تحرير آن بر
نيايد جرّب يا ايّها المجرّب بعد الّذى لا ينبغى لأحد ان يجرّب اللّه بل انّه يجرب
عباده و يمتحنهم كما امتحنك و رجعك الى اسفل الجحيم

\end{otherlanguage}

Ay bī-khabar, tā ḥāl muʿādil-i āncheh az samā-yi bayān nāzil shudeh
nāzil, va bi-yak āyeh muʿādeleh namīnamāyad kull-i āncheh dar āsmān-hā
va zamīn ast. Va agar tu bi-īn shuʾūn nāẓirī, pas ḥāḍir shaw bayna yaday
tā mulāḥaẓeh namāʾī keh ḥīn-i nuzūl-i āyāt aḥadī az ʿuhdeh-yi taḥrīr-i
ān bar nayāyad. Jarrib yā ayyuhā al-mujarrib, baʿda alladhī lā yanbaghī
li-aḥadin an yujarriba Allāh, bal innahu yujarribu ʿibādahu va
yamtaḥinuhum kamā imtaḥanaka va rajaʿaka ilā asfali al-jaḥīm.

O unaware one! Thus far, the equivalent of what has descended from the
heaven of utterance has been revealed, yet all that is in the heavens
and on earth cannot equal a single verse. If you are looking to such
matters, then come into my presence, so that you may see that, when the
verses descend, no one can manage to write them down. Put it to the
test, you who would test, although it is not fitting for anyone to test
God. Rather, He tests His servants and tries them, just as He tried you
and returned you to the lowest depths of hell.

\subsection{518}\label{section-1013}

\begin{otherlanguage}{persian}

نزد هر عاقل بالغى اين ادلّه شما مثل كلمات صبيان بوده و خواهد بود قد جعلتم
امر اللّه لهواً و لعبّا و هزوا و سخریا فسوف يأخذكم اللّه بقهر من لدنه و انّه
لقهّار مقتدر قدير روايح دليل و استدلالت رياحين رضوان دليل و برهان الهى را
پژمرده نموده و اگر خود اقلّ من آن در كلماتت بانصاف ملاحظه نمائى بر قبائحش
اقرب من آن ملتفت شوى هيچ شنيده از اوّل ابداع تا حين كه نفسى بحقّ معارضه
نمايد بآنچه تو در اين كتاب سجّين خود معارضه نموده لا فو‌ربّ العالمين ما
شهدت عين الابداع احجب منك و لا عين الاختراع اغفل منك و لكن انّك انت تفرح
بما رقمت و تكون من الفرحين نسئل اللّه بأن يبدّل فرحك بالحزن الكبرى الّا بأن
تتوب و ترجع اليه و تكون من التّائبين

\end{otherlanguage}

Nazd-i har ʿāqil-i bālighī īn adilleh-yi shumā mithl-i kalamāt-i ṣibyān
būdeh va khvāhad būd. Qad jaʿaltum amra Allāhi lahwan va laʿiban va
huzuwan va sukhriyyan, fa-sawfa yaʾkhudhukumu Allāhu bi-qahrin min
ladunhu va innahu la-qahhārun muqtadirun qadīr. Ravāyiḥ-i dalīl va
istidlālat rayāḥīn-i Riḍvān-i dalīl va burhān-i ilāhī rā pazhmurdeh
namūdeh. Va agar khud aqall-i min ān dar kalamātat bi-inṣāf mulāḥaẓeh
namāʾī, bar qabāyiḥash aqrab-i min ān multafit shavī. Hīch shanīdeh az
avval-i ibdāʿ tā ḥīn keh nafsī bi-Ḥaqq muʿāraḍeh namāyad bi-āncheh tu
dar īn kitāb-i Sijjīn-i khud muʿāraḍeh namūdeh? Lā fa-wa-rabbi
al-ʿālamīn, mā shahidat ʿaynu al-ibdāʿi aḥjaba minka va lā ʿaynu
al-ikhtirāʿi aghfala minka, va lākin innaka anta tafraḥu bi-mā raqamta
va takūnu mina al-fariḥīn. Nasʾalu Allāha bi-an yubaddila faraḥaka
bi-al-ḥuzni al-kubrā illā bi-an tatūba va tarjiʿa ilayhi va takūna mina
al-tāʾibīn.

To every mature, rational person, these proofs of yours have always
been, and will always be, like the words of children. You have made the
Cause of God a diversion, a game, a mockery, and an object of ridicule.
God will therefore seize you with a subduing power from Himself. He is
indeed the All-Subduing, the All-Powerful, the Omnipotent. The odors of
your proof and argument have withered the fragrant flowers of the
paradise of divine proof and demonstration. If you yourself were to
consider your words fairly for less than a moment, you would recognize
their ugliness in less than a moment. Have you ever heard, from the
beginning of creation until now, of anyone opposing the True One with
anything like what you have used to oppose Him in this book of Sijjín of
yours? No, by the Lord of the worlds! The eye of creation has never
witnessed anyone more veiled than you, nor the eye of origination anyone
more heedless. Yet you rejoice in what you have written and are among
the exultant. We ask God to change your joy into the greatest sorrow,
unless you repent, return to Him, and become one of the repentant.

\subsection{519}\label{section-1014}

\begin{otherlanguage}{persian}

و ديگر بالقاى آن مشرك ملقى نوشته عجب‌تر از اين خودشان ادّعاى آن دارند كه
عبارات من آيه است و رسولان ايشان ميگويند اطفال هفت ساله ايشان هم آيه
ميگويند پس كجا آيات من يظهره اللّه دليل حجّيّت او خواهد بود كه خلق او هم
آيه گويد و پسر كوچك ايشانهم آيه ميگويد و اگر گويند كه آيه هم از مقام خلق
و حروفات ظاهر گردد و هم از اصل نقطه و شجره حقيقت پس نفس آيه دليل بر حقّیّت
اصل شجره نخواهد بود شیء ديگر لازم است بيان شود كه چه چيز است انتهى

\end{otherlanguage}

Va dīgar bi-ilqā-yi ān mushrik-i mulqī nivishteh: ʿAjab-tar az īn
khudishān iddiʿā-yi ān dārand keh ʿibārāt-i man āyeh ast va rasūlān-i
īshān mīgūyand aṭfāl-i haft-sāleh-yi īshān ham āyeh mīgūyand. Pas kujā
āyāt-i Man Yuẓhiruhu Allāh dalīl-i ḥujjīyat-i ū khvāhad būd keh khalq-i
ū ham āyeh gūyad va pisar-i kūchak-i īshān ham āyeh mīgūyad? Va agar
gūyand keh āyeh ham az maqām-i khalq va ḥurūfāt ẓāhir gardad va ham az
aṣl-i Nuqṭeh va Shajareh-yi Ḥaqīqat, pas nafs-i āyeh dalīl bar
ḥaqqīyat-i aṣl-i Shajareh nakhvāhad būd. Shayʾ-i dīgar lāzim ast bayān
shavad keh cheh chīz ast. Intahā.

Again, at the prompting of that insinuating idolater, he has written:
``Stranger still, he himself claims that his expressions are verses, and
his messengers say that even his seven-year-old children utter verses.
How, then, will the verses of Him Whom God will make manifest establish
His authority, when His creatures also utter verses and his little son
also utters verses? If they say that verses can appear both from the
station of the creatures and the letters, and from the original Point
and Tree of Truth, then the verse itself will not prove the truth of the
original Tree. Something else is needed; let it be explained what that
is.'' End of quotation.

\subsection{520}\label{section-1015}

\begin{otherlanguage}{arabic}

بلى انّه لو يريد ان ينطق كلّ نفس بثنآء نفسه ليقدر و انّه لهو المقتدر القدير
و لو يريد ان يظهر نفساً بكلّ الآیات ليكون قادراً بسلطانه و ليس لأحد ان
يعترض عليه و كان اللّه على ما اقول شهيد انّه انطق صبيّاً من الصّبيان ليكون
دليلاً عليه و برهاناً لنفسه و هذا مقام الّذى انتم تشكرون اللّه بذلك و
تحمدونه و انتم فتحتم السنكم بالاعتراض بهذا الذّكر الّذى ظهر باسم الأبهی
بين السّموات و الأرضين

\end{otherlanguage}

Balā, innahu law yurīdu an yunṭiqa kulla nafsin bi-thanāʾi nafsihi
la-yaqdiru va innahu la-huwa al-Muqtadiru al-Qadīr. Va law yurīdu an
yuẓhira nafsan bi-kulli al-āyāti la-yakūnu qādiran bi-sulṭānihi va laysa
li-aḥadin an yaʿtariḍa ʿalayhi va kāna Allāhu ʿalā mā aqūlu shahīd.
Innahu anṭaqa ṣabiyyan mina al-ṣibyāni li-yakūna dalīlan ʿalayhi va
burhānan li-nafsihi va hādhā maqāmu alladhī antum tashkurūna Allāha
bi-dhālika va taḥmadūnahu va antum fataḥtum alsinakum bi-al-iʿtirāḍi
bi-hādhā al-Dhikri alladhī ẓahara bi-ismi al-Abhā bayna al-samāwāti va
al-araḍīn.

Yes, if He wishes to cause every soul to speak His praise, He is able to
do so. He is indeed the All-Powerful, the Omnipotent. If He wishes to
bring forth a soul endowed with all the verses, He is able to do so
through His sovereignty, and no one has the right to object to Him. God
is witness to what I say. He has caused a child among the children to
speak, that he might be a sign pointing to Him and a proof of His own
Self. This is a matter for which you should thank and praise God, yet
you have opened your mouths in objection to this Remembrance, who has
appeared under the name of the All-Glorious between the heavens and the
earths.

\subsection{521}\label{section-1016}

\begin{otherlanguage}{persian}

حال ملاحظه كن كه هيچ جاهلى چنين اعتراض نموده و يا مينمايد لا و اللّه مگر
تو و امثال تو كه از هواء نفس و ماء اعراض و نار حسد و تراب بغضا خلق
شده‌ايد اذا موتوا بغيظكم انّه قد ظهر بالحقّ و لا مردّ له و انّ ربّكم الرّحمن قد
كتب على نفسه بأن يظهر فى الأكوان ادلّاء عليه ليقومنّ بين السّموات و الأرض و
ينطقنّ بثنآء نفسه و يرتفعنّ اعلام النّصر بين الخلايق اجمعين اولئك عباد
الّذين جعلهم اللّه مطهّراً عن دونه و منقطعاً عمّا عندكم و متمسّكاً بحبل اللّه
المحكم المتين اولئك لا يمنعهم اشارات من على الأرض و دلالات المشركين

\end{otherlanguage}

Ḥāl mulāḥaẓeh kun keh hīch jāhilī chunīn iʿtirāḍ namūdeh va yā
mīnamāyad? Lā va-Allāh, magar tu va amthāl-i tu keh az havā-yi nafs va
māʾ-i iʿrāḍ va nār-i ḥasad va turāb-i baghḍā khalq shudeh-īd. Idhan mūtū
bi-ghayẓikum, innahu qad ẓahara bi-al-ḥaqqi va lā maradda lahu va inna
rabbakumu al-Raḥmāna qad kataba ʿalā nafsihi bi-an yuẓhira fī al-akwāni
adillāʾa ʿalayhi la-yaqūmunna bayna al-samāwāti va al-arḍi va yanṭiqunna
bi-thanāʾi nafsihi va yartafiʿanna aʿlāmu al-naṣri bayna al-khalāʾiqi
ajmaʿīn. Ulāʾika ʿibādun alladhīna jaʿalahumu Allāhu muṭahharan ʿan
dūnihi va munqaṭiʿan ʿammā ʿindakum va mutamassikan bi-ḥabli Allāhi
al-muḥkami al-matīn. Ulāʾika lā yamnaʿuhum ishārātu man ʿalā al-arḍi va
dalālātu al-mushrikīn.

Now consider: has any ignorant person ever made, or does any make, such
an objection? No, by God, except you and those like you, who have been
created from the air of selfish desire, the water of rejection, the fire
of envy, and the dust of hatred. Then die in your rage! He has appeared
in truth, and there is no turning Him back. Your Lord, the All-Merciful,
has taken it upon Himself to bring forth throughout existence guides who
point to Him. They will assuredly arise between the heavens and the
earth and speak His praise, and the banners of victory will be raised
among all creatures. These are servants whom God has made pure of all
besides Him, detached from what you possess, and steadfast in holding to
the firm, mighty cord of God. Neither the allusions of those on earth
nor the arguments of the idolaters can hold them back.

\subsection{522}\label{section-1017}

\begin{otherlanguage}{persian}

بگو اى جاهل نقمت را از نعمت فرق نداده و خير محض را شرّ شمرده و جوهر رحمت
و عدل را نفس ظلم دانسته بلى طفلى از اطفال بيت در سنّ شش سالگى بآيات اللّه
تكلّم نموده و آن طفل حال حاضر است و واللّه خود همان شجره فساد بكرّات ميگفت
كه آنچه از او شده از هيچ نفسى از بديع اوّل تا حين ظاهر نشده و اكثر ليالى
و ایّام بثنائش مشغول بوده و حال امر بقسمى شده كه از حقّ منيع اعراض نموده
سهل است از اطفال هم دست برنداشته و نخواهد داشت و تا ارض را باين دماء
مطهّره صبغ ندهند دست برندارند و هذا ما نزل فى الألواح ان انتم من النّاظرين
بكرّات فرموده‌اند كه او را در طفوليّت بآيات ناطق فرموده‌ايم تا دليل باشد بر
اين ظهور رحمانى و طلوع قدس صمدانى و تو اين رحمت كبرى را نقمت شمرده عوض
آنكه شكر نمائى حقّ منيع را شكايت آغاز كرده

\end{otherlanguage}

Bigū ay jāhil, niqmat rā az niʿmat farq nadādeh va khayr-i maḥḍ rā sharr
shumurdeh va jawhar-i raḥmat va ʿadl rā nafs-i ẓulm dānisteh. Balā,
ṭiflī az aṭfāl-i bayt dar sinn-i shish-sālegī bi-āyāt-i Allāh takallum
namūdeh va ān ṭifl ḥāl ḥāḍir ast. Va va-Allāh khud hamān shajareh-yi
fasād bi-karrāt mīguft keh āncheh az ū shudeh az hīch nafsī az badīʿ-i
avval tā ḥīn ẓāhir nashudeh va akthar-i layālī va ayyām bi-thanāyash
mashghūl būdeh. Va ḥāl amr bi-qismī shudeh keh az Ḥaqq-i manīʿ iʿrāḍ
namūdeh sahl ast, az aṭfāl ham dast barnadāshteh va nakhvāhad dāsht. Va
tā arḍ rā bi-īn dimāʾ-i muṭahhareh ṣibgh nadahand dast barnadārand. Va
hādhā mā nazala fī al-alwāḥi in antum mina al-nāẓirīn. Bi-karrāt
farmūdeh-and keh ū rā dar ṭufūlīyat bi-āyāt nāṭiq farmūdeh-īm tā dalīl
bāshad bar īn ẓuhūr-i raḥmānī va ṭulūʿ-i quds-i ṣamadānī. Va tu īn
raḥmat-i kubrā rā niqmat shumurdeh, ʿivaḍ-i ānkeh shukr namāʾī Ḥaqq-i
manīʿ rā shikāyat āghāz kardeh.

Say: O ignorant one! You have failed to distinguish affliction from
blessing, have counted pure good as evil, and have regarded the essence
of mercy and justice as injustice itself. Yes, a child among the
children of the House spoke the verses of God at the age of six, and
that child is now present. By God, that very tree of corruption himself
repeatedly said that what had come from him had never appeared from
anyone, from the beginning of creation until now, and he spent most
nights and days praising him. Now matters have reached such a point
that, not content with turning away from the exalted True One, he has
not left even the children alone, nor will he do so. They will not
desist until they have dyed the earth with this pure blood. This is what
has been revealed in the Tablets, if you are among those who look. He
has repeatedly declared: ``We have caused him to speak verses in
childhood, that he might be a proof of this merciful Revelation and this
holy, eternal Dawning.'' Yet you have counted this supreme mercy as an
affliction and, instead of thanking the exalted True One, have begun to
complain.

\subsection{523}\label{section-1018}

\begin{otherlanguage}{persian}

و اينكه نوشته كه خودشان ادّعاى آن دارند كه عبارات من آيه است نزد خود
استهزاء على اللّه عبارات ذكر كرده خذ زمامك يا ايّها الجاهل هذا مقام انقطعت
عنه يداك و ايدى المشركين نقطه بيان ميفرمايد كه كلّ آنچه در ارض موجود است
آيات او بوده و خواهد بود تا چه رسد بما نزل من عنده و ظهر من لدنه و الحقّ
يقول فاعلم بانّ هذا لهو الّذى ظهر من قبل و يظهر من بعد و من فرق بينه و بين
من ظهر او يظهر انّه كفر باللّه و آياته و لو يكون ذا بصر حديد ليشهد كلّ
الظّهورات فى هذا الظّهور الّذى ظهر بالحقّ بآيات مبين

\end{otherlanguage}

Va īnkeh nivishteh keh khudishān iddiʿā-yi ān dārand keh ʿibārāt-i man
āyeh ast, nazd-i khud istihzāʾan ʿalā Allāh ʿibārāt dhikr kardeh. Khudh
zimāmaka yā ayyuhā al-jāhil, hādhā maqāmun inqaṭaʿat ʿanhu yadāka va
aydī al-mushrikīn. Nuqṭeh-yi Bayān mīfarmāyad keh kull-i āncheh dar arḍ
mawjūd ast āyāt-i ū būdeh va khvāhad būd, tā cheh rasad bi-mā nazala min
ʿindihi va ẓahara min ladunhu. Va al-Ḥaqqu yaqūlu fa-iʿlam bi-anna hādhā
la-huwa alladhī ẓahara min qablu va yaẓharu min baʿdu va man farraqa
baynahu va bayna man ẓahara aw yaẓharu innahu kafara bi-Allāhi va
āyātihi va law yakūnu dhā baṣarin ḥadīdin la-yashhadu kulla al-ẓuhūrāti
fī hādhā al-ẓuhūri alladhī ẓahara bi-al-ḥaqqi bi-āyātin mubīn.

As for his statement that ``he himself claims that his expressions are
verses,'' he has used the word ``expressions,'' imagining that he is
mocking God. Restrain yourself, O ignorant one! This is a station beyond
the reach of your hands and the hands of the idolaters. The Point of the
Bayán declares that everything existing on earth has always been, and
will always be, His signs, let alone what has descended from Him and
appeared from His presence. The True One says: Know that this is indeed
the One who appeared before and will appear hereafter. Whoever makes a
distinction between Him and anyone who has appeared or will appear has
disbelieved in God and His signs. If he possessed keen sight, he would
behold all Revelations in this Revelation, which has appeared in truth
with manifest verses.

\subsection{524}\label{section-1019}

\begin{otherlanguage}{persian}

بدانكه آنچه از اوّل لا اوّل كه حجج و براهين ظاهر شده از جميع انبیاء و
مرسلين كلّ حجّت اين ظهور بوده و آنچه در اين ظهور حجّت است لازال حجّت حجج قبل
و بعد بوده و خواهد بود و لكن اين طفل را مخصوص بآيات ناطق فرموديم تا بعضى
از نفوس كه برشحى از فيوضات ما لانهايه اين بحر اعظم فايز شده‌اند استكبار
بر منزل آيات ننمايند آن نفسى را كه باو متمسّكى لازال بين يدى بوده و برشحى
از بحر البحور محروم ماند و رجع الى مقرّه فى النّار و انّها لمقرّ الفجّار و
بئس مقرّ هؤلآء الاشرار لو انت من الموقنين

\end{otherlanguage}

Bidānkeh āncheh az avval-i lā avval keh ḥujaj va barāhīn ẓāhir shudeh az
jamīʿ-i anbiyā va mursalīn, kull ḥujjat-i īn ẓuhūr būdeh, va āncheh dar
īn ẓuhūr ḥujjat ast lā-zāl ḥujjat-i ḥujaj-i qabl va baʿd būdeh va
khvāhad būd. Va lākin īn ṭifl rā makhṣūṣ bi-āyāt nāṭiq farmūdīm tā baʿḍī
az nufūs keh bi-rashḥī az fuyūḍāt-i mā lā-nihāyeh-yi īn baḥr-i aʿẓam
fāʾiz shudeh-and istikbār bar Munzil-i āyāt nanamāyand. Ān nafsī rā keh
bi-ū mutamassikī lā-zāl bayna yaday būdeh va bi-rashḥī az Baḥr al-Buḥūr
maḥrūm mānd va rajaʿa ilā maqarrihi fī al-nāri va innahā la-maqarru
al-fujjāri va biʾsa maqarru hāʾulāʾi al-ashrāri law anta mina
al-mūqinīn.

Know that all the testimonies and proofs manifested by all the Prophets
and Messengers from the beginning that has no beginning have been proofs
of this Revelation. Whatever constitutes a proof in this Revelation has
always been, and will always be, the proof of the proofs of the past and
the future. Yet We specifically caused this child to speak verses so
that certain souls who have received a sprinkling of the boundless
outpourings of this Most Great Ocean would not exalt themselves
arrogantly above the Revealer of verses. The person to whom you cling
was continually in Our presence, yet remained deprived of a sprinkling
from the Ocean of Oceans and returned to his abode in the fire. It is
indeed the abode of the wicked, and wretched is the abode of these
evildoers, if you are among those who are certain.

\subsection{525}\label{section-1020}

\begin{otherlanguage}{persian}

بدان مابين شمس و تجلّيات آن بر جدار و احجار فرق لايحصى مشهود هر بصر ضعيفى
ادراك مينمايد تا چه رسد باَ بصر حديده و از اين گذشته مثال شمس در مرآت
ظاهره مشاهده ميشود آيا ميتوان گفت كه اين نفس شمس سماء است بلكه فرق ما
بين ارض و سماء مشهود و لاينكر هذا الّا كلّ مغلّ مبغوض چنانچه اليوم اگر سمعت
را لطيف نمائى و از آنچه شنيده از مشركين پاك و طاهر كنى كلمات معرض باللّه
را مثل كلمات صبيان بل احقر عند كلمات منزله از سماء بيان رحمن مشاهده
مينمائى قسم بخدا كه نسبت نميتوان داد بلكه دونش لايق ذكر نه و در صورتى
قابل ذكرند كه بعنايت الهى موقن باشند و بعد از انكار لاشیء محض بوده و
خواهند بود

\end{otherlanguage}

Bidān mābayn-i shams va tajalliyāt-i ān bar jidār va aḥjār farq-i lā
yuḥṣā mashhūd. Har baṣar-i ḍaʿīfī idrāk mīnamāyad, tā cheh rasad bā
baṣar-i ḥadīdeh. Va az īn gudhashteh mithāl-i shams dar mirʾāt-i ẓāhereh
mushāhedeh mīshavad. Āyā mītavān guft keh īn nafs-i shams-i samā ast?
Balkeh farq mābayn-i arḍ va samā mashhūd, va lā yunkiru hādhā illā kullu
mughillin mabghūḍ. Chunāncheh al-yawm agar samʿat rā laṭīf namāʾī va az
āncheh shanīdeh az mushrikīn pāk va ṭāhir kunī, kalamāt-i muʿriḍ-i
bi-Allāh rā mithl-i kalamāt-i ṣibyān bal aḥqar ʿinda kalamāt-i
munazzaleh az samā-yi bayān-i Raḥmān mushāhedeh mīnamāʾī. Qasam bi-Khudā
keh nisbat namītavān dād, balkeh dūnash lāyiq-i dhikr nah. Va dar ṣūratī
qābil-i dhikr-and keh bi-ʿināyat-i ilāhī mūqin bāshand va baʿd az inkār
lā-shayʾ-i maḥḍ būdeh va khvāhand būd.

Know that an immeasurable difference is evident between the sun and its
radiance upon walls and stones. Even weak sight perceives it, let alone
keen sight. Moreover, the image of the sun is seen in an outward mirror.
Can it be said that this is the very sun in heaven? Rather, a difference
as great as that between earth and heaven is evident, and none denies
this except one who is treacherous and detested. Likewise, today, if you
refine your hearing and cleanse and purify it of what you have heard
from the idolaters, you will find the words of the one who turns away
from God to be like the words of children---indeed, more
contemptible---beside the words sent down from the heaven of the
All-Merciful's utterance. I swear by God that no comparison can be made;
rather, what is other than Him is not worthy of mention. They are worthy
of mention only if they are assured of divine favor; after denial, they
have been and will remain sheer nothingness.

\subsection{526}\label{section-1021}

\begin{otherlanguage}{persian}

بدان لازال اصل حجّت آيات الهى بوده و لكن از اين ظهور اعظم من دون آيات
ظاهر شده آنچه از اتيان بمثلش كلّ عاجز بوده و خواهند بود خف عن اللّه و
لاتتّبع ما يأمرك هواك فاتّبع الحقّ الّذى جائكم من مشرق الرّوح بنبأ اللّه
المقتدر المهيمن العزيز القيّوم و انّه ليكفيكم عن كلّ شیء و من دونه لن
يكفيكم و لو تتمسّكوا بخلق السّموات و الأرضين ان انتم توقنون هذا نصح اللّه
عليك و على عباده من استنصح بنصح ربّه فله و من اعرض انّ اللّه غنى عمّا كان و
عمّا يكون

\end{otherlanguage}

Bidān lā-zāl aṣl-i ḥujjat āyāt-i ilāhī būdeh, va lākin az īn ẓuhūr-i
aʿẓam min dūn-i āyāt ẓāhir shudeh āncheh az ityān bi-mithlash kull ʿājiz
būdeh va khvāhand būd. Khaf ʿan Allāh va lā tattabiʿ mā yaʾmuruka hawāka
fa-ittabiʿ al-Ḥaqqa alladhī jāʾakum min mashriqi al-Rūḥi bi-nabaʾi
Allāhi al-Muqtadiri al-Muhaymini al-ʿAzīzi al-Qayyūm. Va innahu
la-yakfīkum ʿan kulli shayʾin va min dūnihi lan yakfīkum va law
tatamassakū bi-khalqi al-samāwāti va al-araḍīna in antum tūqinūn. Hādhā
nuṣḥu Allāhi ʿalayka va ʿalā ʿibādihi. Man istanṣaḥa bi-nuṣḥi rabbihi
fa-lahu va man aʿraḍa inna Allāha ghaniyyun ʿammā kāna va ʿammā yakūn.

Know that divine verses have always been the fundamental proof. Yet from
this Most Great Revelation there has appeared, apart from verses, that
which all have been and will remain powerless to match. Fear God and do
not follow what your desire commands you. Follow the True One, who has
come to you from the dawning-place of the Spirit with the tidings of
God, the All-Powerful, the All-Protecting, the Almighty, the
Self-Subsisting. He will indeed suffice you in place of all things,
while nothing besides Him will suffice you, even if you cling to the
creatures of the heavens and the earths, if you are certain. This is
God's counsel to you and to His servants. Whoever heeds his Lord's
counsel does so for his own benefit. As for whoever turns away, God is
independent of all that has been and all that will be.

\subsection{527}\label{section-1022}

\begin{otherlanguage}{persian}

و ديگر القاى آن معرض ملقى اگر گويند كه آنها تصديق ايشان نموده حال سخن
ميرزا يحيى هم همين سخن است كه بر فرض كه شما هم آيه بگوئيد يا هزار كس
چنانچه ربّ اعلى فرموده كه بعدد كلّ شیء اگر خواهم مبعوث بلسان آيات نمايم
باز بايد تصديق من نمايند و در ظلّ من باشند بعلّت آنكه اوّل مرآت هستم و كلّ
منعكس از من ميباشند چونكه فطرة اللّه از من ظاهر است بتصديق نقطه بيان
انتهى

\end{otherlanguage}

Va dīgar ilqā-yi ān muʿriḍ-i mulqī: Agar gūyand keh ānhā taṣdīq-i īshān
namūdeh, ḥāl sukhan-i Mīrzā Yaḥyā ham hamīn sukhan ast keh bar farḍ keh
shumā ham āyeh bigūʾīd yā hizār kas, chunāncheh Rabb-i Aʿlā farmūdeh keh
bi-ʿadad-i kull-i shayʾ agar khvāham mabʿūth bi-lisān-i āyāt namāyam,
bāz bāyad taṣdīq-i man namāyand va dar ẓill-i man bāshand, bi-ʿillat-i
ānkeh avval mirʾāt hastam va kull munʿakis az man mībāshand, chūnkeh
fiṭrat-i Allāh az man ẓāhir ast bi-taṣdīq-i Nuqṭeh-yi Bayān. Intahā.

Another insinuation from that insinuating opponent: ``If they say that
those people have acknowledged him, Mírzá Yaḥyá's statement now is
precisely this: even supposing that you too utter verses, or that a
thousand people do so---as the Most Exalted Lord has said, `If I wish, I
can raise up, to the number of all things, those who speak in the
language of verses'---they must still acknowledge me and remain beneath
my shadow, because I am the first mirror and all are reflected from me,
since the innate nature of God is manifest through me, as attested by
the Point of the Bayán.'' End of quotation.

\subsection{528}\label{section-1023}

\begin{otherlanguage}{persian}

بر فرض تصديق از براى حقّ لم يزل و لايزال مرايا بوده و اوّليّت آن را احدى
احصا ننموده و اين مخصوص بنفسى دون نفسى نبوده كلّ شیء مراياى الهيّه بوده و
خواهند بود و لكن ظهور قبل از مرايا بكلّها شكايت فرموده‌اند بقوله عزّ و جلّ
لاشكونّ اليك ان يا مرآت جودى عن كلّ المرايا كلّ بألوانهم الىّ لينظرون و
مخاطباً لذلك الاسم ميفرمايند كه توئى آن مرآت اوّلیّه كه لم‌يزل از حقّ حكايت
كرده و لايزال حكايت خواهى كرد در مقام ديگر اوّل من آمن را مرآت اوّلیّه
فرموده‌اند و اين اذكار در رتبه ملكيّه ذكر ميشود و الّا از براى مرايا نه
اوّلى بوده و نه آخرى و جميع اين مرايا بقولى خلق شده و خواهند شد ماداميكه
در ظلّ شجره فردانيّه مستقرّند انوار حقّ در آنها ظاهر بحيث لايرى فيهنّ الّا
اللّه و بعد از انحراف مظاهر شيطانيّه بوده و خواهند بود فو‌الّذى نفسى بيده كه
اين شكايتى كه از جميع مرايا فرموده‌اند برهانى است عظيم اگر باشيد از
شاعران

\end{otherlanguage}

Bar farż-i taṣdīq, az barā-yi Ḥaqq lam yazal va lā yazāl marāyā būdeh va
avvaliyyat-i ān rā aḥadī iḥṣā nanamūdeh, va īn makhṣūṣ bi-nafsī dūn-i
nafsī nabūdeh. Kull-i shayʾ marāyā-yi ilāhiyyeh būdeh va khvāhand būd,
va lākin Ẓuhūr-i qabl az marāyā bi-kullihā shikāyat farmūdeh-and,
bi-qawlihi ʿazza wa-jall: La-ashkuwanna ilayka an yā mirʾāta jūdī ʿan
kulli al-marāyā, kullun bi-alwānihim ilayya la-yanẓurūn. Va mukhāṭaban
li-dhālika al-ism mīfarmāyand keh tuʾī ān mirʾāt-i avvaliyyeh keh lam
yazal az Ḥaqq ḥikāyat kardeh va lā yazāl ḥikāyat khvāhī kard. Dar
maqām-i dīgar avval man āmana rā mirʾāt-i avvaliyyeh farmūdeh-and. Va īn
adhkār dar rutbeh-yi mulkiyyeh dhikr mīshavad, va illā az barā-yi marāyā
nah avvalī būdeh va nah ākhirī. Va jamīʿ-i īn marāyā bi-qawlī khalq
shodeh va khvāhand shod. Mādāmī keh dar ẓill-i shajareh-yi fardāniyyeh
mustaqarrand, anvār-i Ḥaqq dar ānhā ẓāhir, bi-ḥaythu lā yurā fīhinna
illā Allāh, va baʿd az inḥirāf maẓāhir-i shayṭāniyyeh būdeh va khvāhand
būd. Fa-wa-alladhī nafsī bi-yadihi keh īn shikāyatī keh az jamīʿ-i
marāyā farmūdeh-and burhānī ast ʿaẓīm, agar bāshīd az shāʿirān.

Even supposing this claim is accepted, there have always been and will
always be mirrors of the True One, and no one has reckoned their
beginning. This has not been exclusive to one soul rather than another.
All things have been and will be divine mirrors. Yet the former
Manifestation complained of all the mirrors, saying---exalted and
glorified is He: ``I will surely complain to you, O Mirror of My bounty,
of all the mirrors: each looks upon Me through its own colors.''
Addressing that Name, He says: ``You are that Primal Mirror which has
always reflected the True One and will always reflect Him.'' Elsewhere
He called the First to Believe the Primal Mirror. These designations are
mentioned at the earthly level; otherwise, the mirrors have neither a
first nor a last. All these mirrors have been and will be created by a
word. As long as they remain beneath the shadow of the Tree of Oneness,
the lights of the True One are manifest in them, so that nothing is seen
in them but God. After turning aside, however, they have been and will
be satanic manifestations. By Him in whose hand is my soul, this
complaint He made against all the mirrors is a mighty proof, if you are
among those who understand.

\subsection{529}\label{section-1024}

\begin{otherlanguage}{persian}

اين شأن مرايا كه ذكر شد بنصّ نقطه بيان و مع‌ذلك تلقاء شمس ذكر مرايا
مينمائى كه شموس مرايا عند تجلّى او معدوم صرف بوده و خواهند بود اگر نفس
ظهور بر جميع اشیاء اطلاق مرآتيّت در رتبه اوّلیّه فرمايد حقّ لاريب فيه بشنو
نداى اين عبد را و از اين اراضى جرزه خود را نجات ده و بفضاهاى قدس وارد شو
فو‌اللّه اين اذكار جز بر توهّمات و حجبات نيفزايد و هر قدر باين كلمات مشغول
شوى از زلال سلسال بی‌مثال حضرت لايزال بعيد مانى و محروم گردى بريز اين
امثال را بحقّ و ما يظهر من عنده ناظر شو چه كه كلّ باو منتهى و از او ظاهر
اين كلمات خوب است از براى امثال خود در فنون جهليّه ظاهريّه كه بزعم خود علم
ناميده‌ايد استدلال نمائيد قطره ملح اجاج ببحر عذب فرات مفرست بدانكه آنچه
لفظ وجود بر آن صادق مخلوق بوده و خواهد بود

\end{otherlanguage}

Īn shaʾn-i marāyā keh dhikr shod bi-naṣṣ-i Nuqṭeh-yi Bayān, va
maʿa-dhālik tilqāʾ-i shams dhikr-i marāyā mīnamāʾī, keh shumūs-i marāyā
ʿinda tajallī-yi ū maʿdūm-i ṣirf būdeh va khvāhand būd. Agar nafs-i
Ẓuhūr bar jamīʿ-i ashyāʾ iṭlāq-i mirʾātiyyat dar rutbeh-yi avvaliyyeh
farmāyad, ḥaqq lā rayba fīhi. Bishnaw nidā-yi īn ʿabd rā va az īn
arāżī-yi jurzeh khud rā najāt dih va bi-faḍāhā-yi quds vārid shaw.
Fa-wa-Allāh īn adhkār juz bar tavahhumāt va ḥujubāt nayafzāyad, va har
qadr bi-īn kalimāt mashghūl shavī az zulāl-i salsāl-i bī-mithāl-i
Ḥażrat-i Lā Yazāl baʿīd mānī va maḥrūm gardī. Birīz īn amthāl rā,
bi-Ḥaqq va mā yaẓharu min ʿindihi nāẓir shaw, cheh keh kull bi-ū muntahī
va az ū ẓāhir. Īn kalimāt khūb ast az barā-yi amthāl-i khud dar funūn-i
jahliyyeh-yi ẓāhiriyyeh keh bi-zaʿm-i khud ʿilm nāmīdeh-īd istidlāl
namāyīd. Qaṭreh-yi milḥ-i ujāj bi-baḥr-i ʿadhb-i furāt mafirist. Bidān
keh āncheh lafẓ-i vujūd bar ān ṣādiq makhlūq būdeh va khvāhad būd.

Such is the station of the mirrors, as stated in the explicit text of
the Point of the Bayán. Yet you speak of mirrors before the Sun, when
the very suns of the mirrors have been and will be utter nothingness at
His revelation. If the Manifestation Himself designates all things as
mirrors of the first rank, it is the truth, beyond all doubt. Hear the
call of this servant, rescue yourself from these barren lands, and enter
the expanses of holiness. By God, these references only increase vain
imaginings and veils. The more you occupy yourself with these words, the
further you remain from the peerless, pure, flowing water of the Eternal
One, and the more deprived you become. Cast aside these comparisons and
look to the True One and what appears from Him, for all things culminate
in Him and proceed from Him. Such words are suitable for presenting
arguments to people like yourself in those outward arts of ignorance
which you imagine to be knowledge. Do not send a bitter, briny drop to
the ocean of sweet, refreshing water. Know that whatever the word
``existence'' applies to has been and will be created.

\subsection{530}\label{section-1025}

\begin{otherlanguage}{persian}

ذكر فطرة اللّه نموده كلّ اشیاء بر فطرة اللّه خلق شده‌اند فطرت را ندانسته از
اين گذشته بشنو نداى فطرت را كه بهيكلى مبعوث و تلقاء باب اعظم باين كلمات
ناطق قوله عزّ بهائه و من المشركين من قال هذه الآیات ما نزلت على الفطرة
تاللّه الحقّ انّ الفطرة حينئذ قد ظهرت على هيكل خادم و قامت لدى الباب بخضوع
و اناب و يضج و يقول فويل لكم يا معشر المغلّين تاللّه انّى قد خلقت بأمر من
لدنه كذلك تشهد لنفسها و لكن لا يفقهون هؤلآء الأشرار تاللّه انها لتفتخر
بنسبتها الى نفسنا الحقّ و انّا لم يزل كنّا غنيّاً عنها قد خلقناها و كلّ شیء
بأمر من لدنا و لا ينكر ذلك الّا كلّ منكر كفّار

\end{otherlanguage}

Dhikr-i fiṭrat Allāh namūdeh. Kull-i ashyāʾ bar fiṭrat Allāh khalq
shodeh-and. Fiṭrat rā nadānesteh. Az īn gudhashteh, bishnaw nidā-yi
fiṭrat rā keh bi-haykalī mabʿūth va tilqāʾ-i Bāb-i Aʿẓam bi-īn kalimāt
nāṭiq. Qawluhu ʿazza bahāʾuhu: Wa-min al-mushrikīna man qāla hādhihi
al-āyātu mā nazalat ʿalā al-fiṭrah. Ta-Allāhi al-Ḥaqq, inna al-fiṭrata
ḥīnaʾidhin qad ẓaharat ʿalā haykali khādim wa-qāmat ladā al-bāb
bi-khuḍūʿin wa-ināb, wa-yaḍijju wa-yaqūlu fa-waylun lakum yā maʿshara
al-mughillīn. Ta-Allāhi innī qad khuliqtu bi-amrin min ladunhu.
Kadhālika tashhadu li-nafsihā wa-lākin lā yafqahūna hāʾulāʾi al-ashrār.
Ta-Allāhi innahā la-taftakhiru bi-nisbatihā ilā nafsinā al-Ḥaqq, wa-innā
lam yazal kunnā ghaniyyan ʿanhā. Qad khalaqnāhā wa-kulla shayʾin
bi-amrin min ladunnā, wa-lā yunkiru dhālika illā kullu munkirin kaffār.

You have mentioned the innate nature of God. All things have been
created according to the innate nature of God. You have not understood
innate nature. Beyond this, hear the call of innate nature, raised up in
a bodily form and speaking these words before the Most Great Gate. His
words---exalted be His glory: ``Among those who associate partners with
God are some who say, `These verses were not revealed according to
innate nature.' By God, the True One, innate nature has now appeared in
the form of a servant and has stood at the Gate in humility and
repentance, crying aloud and saying, `Woe to you, company of the
malicious! By God, I was created by a command from Him.' Thus it
testifies on its own behalf, but these wicked ones do not understand. By
God, it indeed takes pride in its relationship to Our own Self, the True
One, while We have always been independent of it. We created it and all
things by a command from Us. None denies this but every ungrateful
denier.''

\subsection{531}\label{section-1026}

\begin{otherlanguage}{persian}

و اين آياتى است كه از قبل در ذكر فطرت اصليّه الهيّه از سماء ابهى نازل و
اين از مطالب اهل فرقان است چنانچه در كتاب الهى نازل اقم للّدين حنيفاً فطرة
اللّه الّتى فطر النّاس عليها اينست كه در حديث نبوى وارد كلّ مولود يولد على
الفطره الى آخر و اين فطرت همان فطرة اللّه بوده و اين مقام صادق ماداميكه
در ظلّ ايمان مستقرّ باشد و بعد از تغيير از مشركين محسوب چنانچه باقى حديث
ميفرمايد انما ابواه يهوّدانه و يمجّسانه و ينصّرانه

\end{otherlanguage}

Va īn āyātī ast keh az qabl dar dhikr-i fiṭrat-i aṣliyyeh-yi ilāhiyyeh
az samā-yi Abhā nāzil, va īn az maṭālib-i ahl-i Furqān ast, chunāncheh
dar Kitāb-i ilāhī nāzil: Aqim li-al-dīni ḥanīfan fiṭrata Allāhi allatī
faṭara al-nāsa ʿalayhā. Īnast keh dar ḥadīth-i nabavī vārid: Kullu
mawlūdin yūladu ʿalā al-fiṭrah, ilā ākhir. Va īn fiṭrat hamān fiṭrat
Allāh būdeh, va īn maqām ṣādiq mādāmī keh dar ẓill-i īmān mustaqarr
bāshad, va baʿd az taghyīr az mushrikīn maḥsūb, chunāncheh bāqī-yi
ḥadīth mīfarmāyad: Innamā abawāhu yuhawwidānihi wa-yumajjisānihi
wa-yunaṣṣirānihi.

These are verses previously revealed from the Abhá heaven concerning the
original divine innate nature. This is among the subjects addressed by
the people of the Qur'an, as revealed in the divine Book: ``Stand
upright in devotion to religion---the innate nature of God according to
which He fashioned humanity.'' Thus it is stated in the prophetic
tradition: ``Every child is born according to innate nature,'' and so
forth. This innate nature is that same innate nature of God, and this
station holds true as long as one remains beneath the shadow of faith.
After a change, one is counted among those who associate partners with
God, as the remainder of the tradition says: ``It is only his parents
who make him a Jew, a Magian, or a Christian.''

\subsection{532}\label{section-1027}

\begin{otherlanguage}{persian}

بگو اى مهمل مبهم تو اليوم بچشم و گوش مشاهده و استماع مينمائى كه بحر آيات
موّاج و چون غيث هاطل از سماء مشيّت نازل مع‌ذلك مينويسى كه ميرزا يحيى ميگويد
كه صاحب آيات منعكس از من ميباشد اى محروم اينمقام از عاكس و معكوس مقدّس
است اگر بعدد واحد در بيان بالغ مشاهده ميشد هرآينه ارض عرفانى منبسط ميگشت
كه كلّ مذعن شوند بر آنكه حرفى از بيان ادراك ننموده‌اند و لكن تا حال كه
مشاهده نشده و بعد الأمر بيده

\end{otherlanguage}

Bigū ay muhmal-i mubham, tu al-yawm bi-chashm va gūsh mushāhadeh va
istimāʿ mīnamāʾī keh baḥr-i āyāt mavvāj va chūn ghayth-i hāṭil az
samā-yi mashiyyat nāzil. Maʿa-dhālik mīnivīsī keh Mīrzā Yaḥyā mīgūyad
keh ṣāḥib-i āyāt munʿakis az man mībāshad. Ay maḥrūm, īn maqām az ʿākis
va maʿkūs muqaddas ast. Agar bi-ʿadad-i vāḥid dar Bayān bāligh
mushāhadeh mīshod, har-āyineh arż-i ʿirfānī munbasiṭ mīgasht keh kull
mudhʿin shavand bar ānkeh ḥarfī az Bayān idrāk nanamūdeh-and. Va lākin
tā ḥāl keh mushāhadeh nashodeh, va baʿd al-amru bi-yadihi.

Say: O heedless, confused one! Today you see with your own eyes and hear
with your own ears that the ocean of verses is surging and that they
descend like pouring rain from the heaven of the divine will. Yet you
write that Mírzá Yaḥyá says, ``The revealer of verses is a reflection of
me.'' O deprived one! This station is sanctified above both the one who
casts a reflection and what is reflected. Had mature souls to the number
of Váḥid been found in the Bayán, a land of spiritual knowledge would
surely have been spread out, such that all would acknowledge that they
had not understood a single letter of the Bayán. Yet none have been seen
thus far, and hereafter the matter is in His hand.

\subsection{533}\label{section-1028}

\begin{otherlanguage}{persian}

قوله كبر غفلته قطع نظر از كلّ سخنها الآن ميرزا يحيى مدّعى هر مقاميكه فوق
آن متصوّر نشود ميباشند بدليل آيات هر كه مدّعى باشد بايد اتيان بمثل و اعلاى
از آن نمايد انتهى اين كلمات ظاهر نشده مگر از منبع غرور و معدن كذب و مطّلع
فحشا اگر تصديق مرآتيّت او شود نقطه اولى روح ما سواه فداه ميفرمايد كه اگر
مرآت ادّعا نمايد كه من شمس هستم نزد شمس واضح است كه او شبحى است از او كه
ميگويد فو‌اللّه اين بيان منزل بيان نصّ صريح است بر اينكه بعد از او بر حقّ
اعتراض نمايند و نفوسيكه دعوى مرآتيّت دارند در مقابل شمس حقيقت ادّعاى شمسى
كنند فويل له و لك تفترون على انفسكم ثمّ على اللّه و لا تشعرون

\end{otherlanguage}

Qawluhu kabura ghaflatuhu: Qaṭʿ-i naẓar az kull-i sukhanhā, al-ān Mīrzā
Yaḥyā muddaʿī-yi har maqāmī keh fawq-i ān mutaṣavvar nashavad mībāshand,
bi-dalīl-i āyāt. Har keh muddaʿī bāshad bāyad ityān bi-mithl va aʿlā-yi
az ān namāyad. Intahā. Īn kalimāt ẓāhir nashodeh magar az manbaʿ-i
ghurūr va maʿdan-i kidhb va maṭlaʿ-i faḥshā. Agar taṣdīq-i mirʾātiyyat-i
ū shavad, Nuqṭeh-yi Ūlā, rūḥ-i mā sivāhu fidāh, mīfarmāyad keh agar
mirʾāt iddiʿā namāyad keh man shams hastam, nazd-i shams vāżiḥ ast keh ū
shabaḥī ast az ū keh mīgūyad. Fa-wa-Allāh īn bayān-i Munzil-i Bayān
naṣṣ-i ṣarīḥ ast bar īnkeh baʿd az ū bar Ḥaqq iʿtirāż namāyand va nufūsī
keh daʿvā-yi mirʾātiyyat dārand dar muqābil-i Shams-i Ḥaqīqat iddiʿā-yi
shamsī kunand. Fa-waylun lahu wa-laka, taftarūna ʿalā anfusikum thumma
ʿalā Allāhi wa-lā tashʿurūn.

His words---great is his heedlessness: ``Setting all other statements
aside, Mírzá Yaḥyá now claims, on the evidence of verses, a station
above which none can be conceived. Whoever makes a claim must produce
their like or something greater.'' End of quotation. These words have
issued from nothing but the wellspring of arrogance, the mine of
falsehood, and the dawning-place of obscenity. Even if his status as a
mirror is acknowledged, the Primal Point---may the souls of all besides
Him be sacrificed for Him---says that if the mirror claims, ``I am the
sun,'' it is clear to the sun that it is an image of itself that is
speaking. By God, this statement of the Revealer of the Bayán is an
explicit declaration that after Him they would raise objections against
the True One, and that souls claiming to be mirrors would claim to be
suns before the Sun of Truth. Woe to him and to you! You invent
falsehoods about yourselves and then about God, yet you do not perceive
it.

\subsection{534}\label{section-1029}

\begin{otherlanguage}{persian}

و ديگر اينكلمه بسيار بی‌شرمى است كه نوشته اليوم هر نفسى مدّعى باشد بايد
اتيان بمثل آيات ميرزا يحيى نمايد و يا اعلى يا ايّها الاعمى انّه قد كان
كليل عند ما نزل من جبروت البيان من لدى اللّه المقتدر العزيز المنّان و اللّه
كلّ اشیاء متحيّرند از اين اقوال كه ذكر نموده وجود و كينونتش اليوم قابل ذكر
نيست تا چه رسد بكلماتش و هر ذى شعورى ميداند كه اگر هم كلمه از كلمات حقّ
نزدش يافت شود از اين مقرّ اقدس اطهر اخذ نموده بسيار بی‌‌خبرى از خدا ميطلبم
كه بصرت را بكحل عرفان موفّق فرمايد كه شايد از رمد اوهام فارغ شده بكمات
مليك علّام ناظر شود تا جميع كلمات را نزد حرفى از كلمات اللّه معدوم صرف
مشاهده نمائى كذلك نزل الأمر من جبروت الحكم فضلاً من عنده على العالمين

\end{otherlanguage}

Va dīgar īn kalameh bisyār bī-sharmī ast keh nivishteh: Al-yawm har
nafsī muddaʿī bāshad bāyad ityān bi-mithl-i āyāt-i Mīrzā Yaḥyā namāyad
va yā aʿlā. Yā ayyuhā al-aʿmā, innahu qad kāna kalīl ʿinda mā nazala min
jabarūti al-bayān min ladā Allāhi al-Muqtadiri al-ʿAzīzi al-Mannān.
Va-Allāh kull-i ashyāʾ mutaḥayyirand az īn aqvāl keh dhikr namūdeh.
Vujūd va kaynūnatash al-yawm qābil-i dhikr nīst tā cheh rasad
bi-kalimātash. Va har dhī-shuʿūrī mīdānad keh agar ham kalameh az
kalimāt-i Ḥaqq nazdash yāft shavad az īn maqarr-i aqdas-i aṭhar akhdh
namūdeh. Bisyār bī-khabarī. Az Khudā mīṭalabam keh baṣarat rā bi-kuḥl-i
ʿirfān muwaffaq farmāyad keh shāyad az ramad-i awhām fārigh shodeh
{[}corrupted text{]} Malīk-i ʿAllām nāẓir shavad tā jamīʿ-i kalimāt rā
nazd-i ḥarfī az kalimāt Allāh maʿdūm-i ṣirf mushāhadeh namāʾī. Kadhālika
nazala al-amru min jabarūti al-ḥukmi faḍlan min ʿindihi ʿalā al-ʿālamīn.

Furthermore, this statement he has written is utterly shameless: ``Today
anyone making a claim must produce verses like those of Mírzá Yaḥyá, or
greater.'' O blind one! He has indeed been powerless before what has
descended from the realm of utterance, from God, the All-Powerful, the
Almighty, the Bestower. By God, all things are bewildered by these
assertions you have made. His very existence and being are unworthy of
mention today, let alone his words. Every discerning person knows that
if even a word of the words of the True One is found with him, he has
taken it from this most holy, most pure Seat. How unaware you are! I ask
God to favor your sight with the eye-salve of spiritual knowledge, so
that it may perhaps be freed from the inflammation of vain imaginings
and look {[}corrupted original text{]} of the all-knowing King, that you
may see all words as utter nothingness beside a single letter of the
words of God. Thus has the decree descended from the dominion of
judgment, as a grace from Him to all the worlds.

\subsection{535}\label{section-1030}

\begin{otherlanguage}{persian}

قوله قلّ ادراكه حال اى برادر من و اللّه امر بر شما مشتبه شده اگر بخواهيد
واقعاً بفهميد كه غرض در كار نيست شما خود وكيل من هستيد اگر محقّق شد بر شما
بمن بنويسيد امّا با انصاف از براى شما بسيار آسان است سياست مدنيه مدخليّتى
بحقّيّت ندارد و اين حكمت عملى است كه تهذيب اخلاق و تدبير منزل و سياست
مدنيه ميگويند كه علم سلوك با خلق است و با اشخاصيكه تعلّق دارند از خدّام و
مريدان كه هر كسى را بمحبّت و خلق با آنچه موافق طبع او است حركت نمايند شما
اشتباه نموده‌ايد اين حكمت عملى را بجوهر حقيقت و اينها حجاب بر شما گشته
اگر بخواهيد تجربه نمائيد روزى را قرار بدهيد در مقام امتحان برآئيد يك
صفحه از قرآن را بخوانيد با كمال جلدى نزد كسى كه مدّعى حقّيّت است مضمون آن
صفحه را بسجع خاصّى بلا تأمّل بنويسد و دفعه ديگر بسجع و عبارت ديگر و هكذا
هر كه نوشت در مدّعاى خود صادق و الّا كاذب بخلاف ميرزا يحيى كه مدّعى اين
مقام است در هر دو خود تجربه نموده و بمن اعلام دهيد اما از روى انصاف يقين
هر كه داراى اينمقام است حقّ و كلام او حقّ و غير او كاذب انتهى

\end{otherlanguage}

Qawluhu qalla idrākuhu: Ḥāl ay barādar-i man, va-Allāh amr bar shumā
mushtabih shodeh. Agar bikhvāhīd vāqiʿan bifahmīd keh gharaż dar kār
nīst, shumā khud vakīl-i man hastīd. Agar muḥaqqaq shod bar shumā bi-man
binivīsīd, ammā bā inṣāf. Az barā-yi shumā bisyār āsān ast. Siyāsat-i
madaniyyeh mudākhilīyatī bi-ḥaqqiyyat nadārad, va īn ḥikmat-i ʿamalī ast
keh tahdhīb-i akhlāq va tadbīr-i manzil va siyāsat-i madaniyyeh
mīgūyand, keh ʿilm-i sulūk bā khalq ast va bā ashkhāṣī keh taʿalluq
dārand az khuddām va murīdān, keh har kasī rā bi-maḥabbat va khulq bā
āncheh muvāfiq-i ṭabʿ-i ū ast ḥarakat namāyand. Shumā ishtibāh
namūdeh-īd īn ḥikmat-i ʿamalī rā bi-jawhar-i ḥaqīqat, va īnhā ḥijāb bar
shumā gashteh. Agar bikhvāhīd tajrubeh namāyīd, rūzī rā qarār bidahīd,
dar maqām-i imtiḥān bar-āyīd. Yik ṣafḥeh az Qurʾān rā bikhvānīd bā
kamāl-i jaldī nazd-i kasī keh muddaʿī-yi ḥaqqiyyat ast, mażmūn-i ān
ṣafḥeh rā bi-sajʿ-i khāṣṣī bilā taʾammul binivīsad, va dafʿeh-yi dīgar
bi-sajʿ va ʿibārat-i dīgar, va hākadhā. Har keh nivisht dar muddaʿā-yi
khud ṣādiq va illā kādhib, bi-khilāf-i Mīrzā Yaḥyā keh muddaʿī-yi īn
maqām ast. Dar har du khud tajrubeh namūdeh va bi-man iʿlām dihīd, ammā
az rū-yi inṣāf. Yaqīn har keh dārā-yi īn maqām ast ḥaqq va kalām-i ū
ḥaqq va ghayr-i ū kādhib. Intahā.

His words---little is his understanding: ``Now, my brother, by God, you
have become confused about the matter. If you truly wish to understand
that no ulterior motive is involved, you yourself are my representative.
If it becomes clear to you, write to me, but be fair. It is very easy
for you. Civic governance has no bearing on the truth of a claim. This
is practical philosophy, whose branches are called the refinement of
character, household management, and civic governance: the knowledge of
how to deal with people and with those attached to one, whether servants
or disciples, treating each with affection and good character in a
manner agreeable to his disposition. You have mistaken this practical
philosophy for the essence of truth, and these things have become a veil
for you. If you wish to test this, appoint a day and undertake an
examination. Read a page of the Qur'an with the utmost speed before the
person claiming to be the true one, and let him write the substance of
that page in a particular rhymed prose without hesitation, then a second
time in a different rhyme and wording, and so on. Whoever writes it is
truthful in his claim; otherwise he is a liar, in contrast to Mírzá
Yaḥyá, who claims this station. Test both of them yourself and inform
me, but do so fairly. Certainly, whoever possesses this station is true,
and his words are true, and anyone else is a liar.'' End of quotation.

\subsection{536}\label{section-1031}

\begin{otherlanguage}{persian}

فو‌اللّه يضحك على كلماتك كلّ ذى بصر حديد گويا ابدا رايحه رضوان الهى بر تو
مرور ننموده چه كه اليوم هر نفسى كه منسوب بحقّ است لابدّ از او اثرى ظاهر و
هويداست و اللّه كلماتت پست‌تر از كلمات ملل قبل بنظر ميايد چون بغض اللّه در
قلبت بوده قلمت از لطافت ذكر و تحرير ممنوع و محروم گشته امر بر تو مشتبه
شده كه از حزب رحمن خارج شده و بحزب شيطان توجّه نموده بر هر ذى بصرى معلوم
است كه اين عباد بعد از پنج سنه معاشرت كه ليلا و نهارا آنچه بچشم خود
ديده‌ايم و بسمع خود شنيده‌ايم ذكر مينمائيم و تو بقول خبيثى كه نزد اين عبد
مردود بوده تا چه رسد نزد ساكنان اعراش معانى تمسّك جسته بالمرّه خود را از
اشراقات انوار فجر معانى كه در اين ایّام رحمانى مشرق و مضئ است ممنوع نموده

\end{otherlanguage}

Fa-wa-Allāhi yaḍḥaku ʿalā kalimātika kullu dhī baṣarin ḥadīd. Gūyā
abadan rāyiḥeh-yi Riḍvān-i ilāhī bar tu murūr nanamūdeh, cheh keh
al-yawm har nafsī keh mansūb bi-Ḥaqq ast lābudd az ū atharī ẓāhir va
huvaydāst. Va-Allāh kalimātat past-tar az kalimāt-i milal-i qabl
bi-naẓar mīāyad. Chūn bughż Allāh dar qalbat būdeh, qalamat az laṭāfat-i
dhikr va taḥrīr mamnūʿ va maḥrūm gashteh. Amr bar tu mushtabih shodeh
keh az ḥizb-i Raḥmān khārij shodeh va bi-ḥizb-i shayṭān tavajjuh
namūdeh. Bar har dhī-baṣarī maʿlūm ast keh īn ʿibād baʿd az panj saneh
muʿāsharat keh laylan va nahāran āncheh bi-chashm-i khud dīdeh-īm va
bi-samʿ-i khud shinīdeh-īm dhikr mīnamāʾīm, va tu bi-qawl-i khabīthī keh
nazd-i īn ʿabd mardūd būdeh tā cheh rasad nazd-i sākinān-i aʿrāsh-i
maʿānī tamassuk juste h, bi-al-marreh khud rā az ishrāqāt-i anvār-i
fajr-i maʿānī keh dar īn ayyām-i raḥmānī mushriq va muḍīʾ ast mamnūʿ
namūdeh.

By God, everyone with keen sight laughs at your words. It seems that the
fragrance of the divine Riḍván has never reached you, for today some
sign must necessarily be manifest and evident in every soul associated
with the True One. By God, your words appear more base than those of the
religious communities of the past. Because hatred of God has been in
your heart, your pen has been deprived of grace in expression and
writing. You have become confused about the matter, leaving the company
of the All-Merciful and turning toward the company of Satan. It is clear
to every discerning person that, after five years of association, we
servants recount what we have seen with our own eyes and heard with our
own ears, by night and day. Yet you have clung to a vile statement
rejected by this servant---let alone by the inhabitants of the thrones
of inner meanings---and have utterly deprived yourself of the radiance
of the lights of the dawn of inner meanings, which is rising and shining
in these merciful days.

\subsection{537}\label{section-1032}

\begin{otherlanguage}{persian}

حال اى با انصاف انصاف ده بر ما مشتبه شده يا بر شما و اينكه نوشته غرض در
كار نيست و اللّه نفس همين قول غرض صرف است و محض نفس و هوى چنانچه خود شما
نوشته‌ايد كه حال برو و تجربه كن اين عبد را مأمور بتجربه نموده كه بعد آنچه
ظاهر شود تو را اخبار نمايم و تو خود اينجا نبوده و تجربه ننموده لذا قبل
از تجربه چرا آنقدر مجعولات مفتريه نفسانيّه بحقّ نسبت داده و انكار صرف
نموده حال ملاحظه كن كه اهل غرض كيست تاللّه لو تنصف لتجد نفسك فى خسران
عظيم ان يا بحر الغرض خف عن اللّه الّذى خلقك بأمر من عنده و لاتتكلّم بما
يكذّبك به كلّ الذّرّات اتّق اللّه و كن من المتّقين حال ملاحظه كن كه چه قدر
افعال و اعمالت نزد حقّ رسوا و واضح است و چگونه عقلت بيد قدرت اخذ شده كه
خود غرض و بغضاى خود را بقلم مينويسى و ثابت ميكنى و مع‌ذلك شاعر نيستى اذاً
يشهد بغرضك كلّ من فى السّموات و الأرض ثمّ يشهد نفسك على غرضك و انّ هذا لحقّ
يقين

\end{otherlanguage}

Ḥāl ey bā-inṣāf, inṣāf deh: bar mā mushtabah shudeh yā bar shumā? Va
īnkih nivishteh gharaḍ dar kār nīst, va-Allāh nafs-i hamīn qawl gharaḍ-i
ṣirf ast va maḥḍ-i nafs va havā. Chinānchih khud-i shumā nivishteh-īd
kih ḥāl biraw va tajribih kun. Īn ʿabd rā maʾmūr bi-tajribih namūdeh kih
baʿd ānchih ẓāhir shavad tu rā ikhbār namāyam, va tu khud īnjā nabūdeh
va tajribih nanamūdeh. Lidhā qabl az tajribih chirā ānqadr majʿūlāt-i
muftariyih-yi nafsāniyyih bi-Ḥaqq nisbat dādih va inkār-i ṣirf namūdeh?
Ḥāl mulāḥiẓih kun kih ahl-i gharaḍ kīst. Ta-Allāh law tunṣif la-tajid
nafsaka fī khusrān ʿaẓīm. An yā baḥr al-gharaḍ, khaf ʿan Allāh alladhī
khalaqaka bi-amr min ʿindihi wa-lā tatakallam bi-mā yukadhdhibuka bihi
kull al-dharrāt. Ittaq Allāh wa-kun min al-muttaqīn. Ḥāl mulāḥiẓih kun
kih chih qadr afʿāl va aʿmālat nazd-i Ḥaqq rusvā va vāḍiḥ ast va
chigūnih ʿaqlat bi-yad-i qudrat akhdh shudeh kih khud gharaḍ va
baghḍā-yi khud rā bi-qalam mīnivīsī va thābit mīkunī va maʿa-dhālik
shāʿir nīstī. Idhan yashhad bi-gharaḍika kullu man fī al-samāwāt
wa-al-arḍ thumma tashhad nafsuka ʿalā gharaḍika wa-inna hādhā la-ḥaqq
yaqīn.

Now, you who claim to be fair, judge fairly: is it we who have become
confused, or you? As for your writing that no ulterior motive is
involved, by God, that very statement is pure ulterior motive and
nothing but self and desire. You yourself wrote, ``Now go and test it.''
You commissioned this servant to conduct a test and then inform you of
whatever became apparent, while you yourself have neither been here nor
conducted a test. Why, then, before testing, have you attributed so many
selfish fabrications and slanders to the True One and denied Him
outright? Now consider who harbors ulterior motives. By God, if you
judge fairly, you will find yourself in grievous loss. O ocean of
ulterior motives, fear God, who created you by a command from His
presence, and do not say what every atom proves false. Fear God and be
among the God-fearing. Now consider how exposed and evident your deeds
and actions are before the True One, and how your understanding has been
seized by the hand of power, so that you yourself write down and
establish your own ulterior motives and hatred, yet remain unaware. Thus
everyone in the heavens and on earth bears witness to your ulterior
motives; then your own self bears witness to your ulterior motives. This
is indeed the certain truth.

\subsection{538}\label{section-1033}

\begin{otherlanguage}{persian}

غرض مشركين بمقامى رسيده كه بعضى از امثال شما ايراد گرفته‌اند و در خلق
جمال رحمن حرف گفته‌اند و اعتراضات نموده‌اند و تو شهادت بر خلق داده و بعد
بغرض نفسانى اعتراض نموده انسان متحيّر است كه كدام يك از اقوال را تصديق
نمايد از يكطرف هيكل فساد باطراف نوشته و شكايتها نموده و اخلاق رحمانيّه را
از مظهر الوهيّه سلب نموده اين كلماتيكه بخارج نوشته‌اند از آنطرف بشما القا
نموده كه بطائفين حول بيت اين قسم نوشته اگر اخلاق صحيح است انكار آنقوم چه
صورتى دارد اگر عدمش صادق بوده اين ذكر تو چه معنى دارد و اللّه هر بصيرى
موقن است كه كلمه حقّ نگفته و نخواهند گفت باقتضاى وقت هر چه مصلحت انفس
خبيثه دانند بآن متكلّم و مستدلّند

\end{otherlanguage}

Gharaḍ-i mushrikīn bi-maqāmī rasīdeh kih baʿḍī az amthāl-i shumā īrād
giriftih-and va dar khulq-i Jamāl-i Raḥmān ḥarf guftih-and va iʿtirāḍāt
namūdih-and, va tu shahādat bar khulq dādih va baʿd bi-gharaḍ-i nafsānī
iʿtirāḍ namūdeh. Insān mutaḥayyir ast kih kudām yak az aqvāl rā taṣdīq
namāyad. Az yik ṭaraf Haykal-i Fasād bi-aṭrāf nivishteh va shikāyathā
namūdeh va akhlāq-i raḥmāniyyih rā az Maẓhar-i Ulūhiyyih salb namūdeh.
Īn kalimātīkih bi-khārij nivishtih-and az ān ṭaraf bi-shumā ilqā namūdeh
kih bi-ṭāʾifīn-i ḥawl-i Bayt īn qism nivishteh. Agar akhlāq ṣaḥīḥ ast
inkār-i ān qawm chih ṣūratī dārad? Agar ʿadamash ṣādiq būdeh īn dhikr-i
tu chih maʿnī dārad? Va-Allāh har baṣīrī mūqin ast kih kalimih-yi ḥaqq
nagu⁣fteh va nakhāhand guft. Bi-iqtiḍā-yi vaqt har chih maṣlaḥat-i
anfus-i khabīthih dānand bi-ān mutakallim va mustadill-and.

The ulterior motives of those who associate partners with God have
reached such a point that some like you have found fault with the
character of the Beauty of the All-Merciful and raised objections to it.
Yet you have testified to His character and then, out of selfish
motives, objected. One is bewildered as to which statement to believe.
On the one hand, the embodiment of corruption has written to various
places, made complaints, and denied that the Manifestation of Divinity
possesses merciful virtues. These are the statements they have written
to those outside; on the other hand, he has suggested to you that he has
written in this fashion to those who circle the House. If the good
character is genuine, what justification is there for those people's
denial? If its absence is true, what does this statement of yours mean?
By God, every discerning person is certain that they have not spoken a
word of truth and never will. According to the requirements of the
moment, they say and adduce as proof whatever they consider advantageous
to their vile selves.

\subsection{539}\label{section-1034}

\begin{otherlanguage}{arabic}

بارى اين عباد نشهد بانّه لعلى خلق عظيم و نشهد بانّه حجّة اللّه بين خلقه و
برهانه بين بريّته و مظهر امره فى بلاده و مطّلع العزّ بين عباده و كلّما ظهر
من عنده لحقّ بمثل وجوده و ما دونه معدوم فى ساحته و مفقود عند ظهورات
سلطانه و انّه لهو محبوب العارفين و مقصود العالمين من اعرض عنه فقد اعرض عن
اللّه المهيمن العزيز القدير

\end{otherlanguage}

Bārī īn ʿibād nashhadu bi-annahu la-ʿalā khuluq ʿaẓīm wa-nashhadu
bi-annahu Ḥujjat Allāh bayna khalqihi wa-burhānuhu bayna bariyyatihi
wa-maẓhar amrihi fī bilādihi wa-muṭṭalaʿ al-ʿizz bayna ʿibādihi
wa-kullamā ẓahara min ʿindihi la-ḥaqq bi-mithl wujūdihi wa-mā dūnahu
maʿdūm fī sāḥatihi wa-mafqūd ʿinda ẓuhūrāt sulṭānihi wa-innahu la-huwa
maḥbūb al-ʿārifīn wa-maqṣūd al-ʿālamīn. Man aʿraḍa ʿanhu fa-qad aʿraḍa
ʿan Allāh al-Muhaymin al-ʿAzīz al-Qadīr.

In any case, we servants testify that He is indeed of exalted character.
We testify that He is God's proof among His creation, His demonstration
among His creatures, the Manifestation of His Cause in His lands, and
the Dawning-place of glory among His servants. Whatever has appeared
from Him is indeed true, just as His own being is true; all else is
nonexistent in His court and lost before the manifestations of His
sovereignty. He is indeed the Beloved of those who know and the Desire
of all the worlds. Whoever turns away from Him has turned away from God,
the All-Protecting, the Almighty, the All-Powerful.

\subsection{540}\label{section-1035}

\begin{otherlanguage}{persian}

اى اهل دين شما را بمالك يوم الدّين قسم ميدهم كه از اين كلمات اقلّ من ذرّ
رايحه حقّ استشمام ميشود لا فو‌نفسه الحقّ اين سخنان مشابه سخنان اراذل اهل
ارض است كه در حين مخاصمه با يكديگر ذكر مينمايند كه من قدرتم از تو بيشتر
و قوّتم از تو برتر و امثال اين سخنان لغو كه ما بين جهّال متداول است و اللّه
حال نوحه بر امر لازم چه كه نقطه بيان با آنمقام و شأن چه مقدار خضوع و
خشوع و تسليم و رضا نسبت باين ظهور فرموده و مع‌ذلك تو باين كلماتيكه اراذل
از آن ننگ دارند تكلّم نموده و بآن كلمات بر بطلان حقّ استدلال كرده بلى حال
امر اللّه بشأنى تنزّل نموده كه جميع ما فى البيان منوط بقول تو و ميرزا يحيى
و سيّد محمّد و على محمّد سرّاج شده و امثال آن نفوس چنانچه در اواخر دوره
فرقان جميع امور الهى و شرع محكم ربّانى بقول شيخى از مشايخ منوط و مشروط
بود

\end{otherlanguage}

Ey ahl-i dīn, shumā rā bi-Mālik-i Yawm al-Dīn qasam mīdiham kih az īn
kalimāt aqall-i min dharr rāyiḥih-yi ḥaqq istishmām mīshavad? Lā
fa-wa-nafsihi al-Ḥaqq. Īn sukhanān mushābih-i sukhanān-i arādhil-i ahl-i
arḍ ast kih dar ḥīn-i mukhāṣimih bā yikdīgar dhikr mīnamāyand kih man
qudratam az tu bīshtar va quvvatam az tu bartar va amthāl-i īn
sukhanān-i laghv kih mā bayn-i juhhāl mutadāvil ast. Va-Allāh ḥāl nawḥih
bar Amr lāzim, chih kih Nuqṭih-yi Bayān bā ān maqām va shaʾn chih miqdār
khuḍūʿ va khushūʿ va taslīm va riḍā nisbat bi-īn Ẓuhūr farmūdeh, va
maʿa-dhālik tu bi-īn kalimātīkih arādhil az ān nang dārand takallum
namūdeh va bi-ān kalimāt bar buṭlān-i Ḥaqq istidlāl kardeh. Balī ḥāl Amr
Allāh bi-shaʾnī tanazzul namūdeh kih jamīʿ-i mā fī al-Bayān manūṭ
bi-qawl-i tu va Mīrzā Yaḥyā va Siyyid Muḥammad va ʿAlī Muḥammad-i Sarrāj
shudeh va amthāl-i ān nufūs, chinānchih dar avākhir-i dawrih-yi Furqān
jamīʿ-i umūr-i ilāhī va sharʿ-i muḥkam-i rabbānī bi-qawl-i shaykhī az
mashāyikh manūṭ va mashrūṭ būd.

O people of religion, I adjure you by the Master of the Day of Judgment:
can even the slightest trace of the fragrance of truth be perceived in
these words? No, by His own true Self! These words resemble those
uttered by the basest people on earth when quarreling with one another:
``My power is greater than yours, and my strength surpasses yours,'' and
similar idle talk current among the ignorant. By God, the Cause must now
be mourned, for the Point of the Bayān, with all His station and
dignity, has shown such humility, reverence, submission, and contentment
toward this Revelation. Yet you have uttered words of which even the
base would be ashamed, and have used those words to argue that the True
One is false. Indeed, the Cause of God has now been brought so low that
everything in the Bayān has become dependent upon your word and that of
Mīrzā Yaḥyā, Siyyid Muḥammad, ʿAlī Muḥammad the saddler, and others like
them, just as, toward the end of the dispensation of the Qur'an, all
divine affairs and the firmly established law of God depended upon, and
were conditional on, the word of one shaykh among the shaykhs.

\subsection{541}\label{section-1036}

\begin{otherlanguage}{persian}

و اين معلوم است كه از القاهاى شيطان است گويا تو از خود هيچ شعورى ندارى
بر تند باد نفس و هوى مبتلا گشته هر قسم كه ميخواهد تو را ميبرد آنقدر فكر
ننموده كه اصل اين فقره نزد مبصّرين از الحاد است و در الحاد اختلاف
نموده‌اند بعضى بر آنند كه ملحد نفوسى هستند كه خدا را باسامى مينامند كه
اذن داده نشده‌اند در كتاب الهى و بعضى برآنند كه ملحد نفوسى هستند كه از
اسماء الهى اسمائى اشتقاق نموده‌اند و بآن اسماء عاكفند من دون اللّه چنانچه
لفظ منات را از منّان و عزّى را از عزيز و لات را از اللّه اخذ نموده‌اند و بآن
اسماء ساجد و عاكفند و بعضى برآنند كه الحاد همان تحريف است و آن تحريف هم
در معانى نه در الفاظ چه كه لحد بمعنى ستر آمده و آن نفوسى هستند كه معانى
كلمات الهى را بهواهاى خود ستر نموده‌اند يعنى بهواهاى خود تفسير نموده و
مينمايند و بعضى برآنند كه ملحد نفوسى هستند كه معانى غير را بلباس الفاظ
خود جلود دهند اين معانى الحاد كه مابين ناس معروف است

\end{otherlanguage}

Va īn maʿlūm ast kih az ilqāhā-yi shayṭān ast. Gūyā tu az khud hīch
shuʿūrī nadārī. Bar tundbād-i nafs va havā mubtalā gashteh, har qism kih
mīkhāhad tu rā mībarad. Ānqadr fikr nanamūdeh kih aṣl-i īn fiqrih nazd-i
mubṣirīn az ilḥād ast. Va dar ilḥād ikhtilāf namūdih-and. Baʿḍī bar
ānand kih mulḥid nufūsī hastand kih Khudā rā bi-asāmī mīnāmand kih idhn
dādih nashudih-and dar Kitāb-i Ilāhī. Va baʿḍī bar ānand kih mulḥid
nufūsī hastand kih az asmā-yi ilāhī asmāʾī ishtiqāq namūdih-and va bi-ān
asmā ʿākif-and min dūn Allāh, chinānchih lafẓ-i Manāt rā az Mannān va
ʿUzzā rā az ʿAzīz va Lāt rā az Allāh akhdh namūdih-and va bi-ān asmā
sājid va ʿākif-and. Va baʿḍī bar ānand kih ilḥād hamān taḥrīf ast va ān
taḥrīf ham dar maʿānī nah dar alfāẓ, chih kih laḥd bi-maʿnī-yi satr
āmadeh va ān nufūsī hastand kih maʿānī-yi kalimāt-i ilāhī rā
bi-havāhā-yi khud satr namūdih-and, yaʿnī bi-havāhā-yi khud tafsīr
namūdeh va mīnamāyand. Va baʿḍī bar ānand kih mulḥid nufūsī hastand kih
maʿānī-yi ghayr rā bi-libās-i alfāẓ-i khud {[}corrupted text{]} dihand.
Īn maʿānī-yi ilḥād kih mā bayn-i nās maʿrūf ast.

It is clear that this is one of Satan's promptings. It seems you have no
understanding of your own. You have been caught in the tempest of self
and desire, which carries you wherever it wishes. You have not reflected
enough to realize that, in the eyes of the discerning, this very
proposal is a form of deviation. Opinions differ concerning deviation.
Some hold that deviators are those who call God by names not authorized
in the divine Book. Others hold that deviators are those who derive
names from the divine names and devote themselves to those names instead
of God, just as they derived the word Manāt from Mannān, ʿUzzā from
ʿAzīz, and Lāt from Allāh, and bow down and devote themselves to those
names. Others hold that deviation is distortion, and that this
distortion concerns meanings, not words, for laḥd has the meaning of
concealment: these are people who have concealed the meanings of the
divine words with their own desires---that is, they have interpreted and
continue to interpret them according to their own desires. Others hold
that deviators are those who {[}corrupted original text{]} the meanings
of another in the garb of their own words. These are the meanings of
deviation commonly known among people.

\subsection{542}\label{section-1037}

\begin{otherlanguage}{persian}

حال اين مطلبی‌ كه خواسته از اين اقسام خارج نه در هر صورت الحاد بوده و
خواهد بود ملقيت اين فقره را براى تو معجزه قرار داده و از اين گذشته اين
اسهل طرق است عجب است كه صعب بنظر شما آمده چه كه معنى موجود باشد تحديد
الفاظ بغايت سهل و آسان بوده و خواهد بود

\end{otherlanguage}

Ḥāl īn maṭlabī kih khāstih az īn aqsām khārij nah. Dar har ṣūrat ilḥād
būdeh va khāhad būd. Mulqī-at īn fiqrih rā barā-yi tu muʿjizih qarār
dādeh. Va az īn gudashteh īn ashal-i ṭuruq ast. ʿAjab ast kih ṣaʿb
bi-naẓar-i shumā āmadeh, chih kih maʿnī mawjūd bāshad taḥdīd-i alfāẓ
bi-ghāyat sahl va āsān būdeh va khāhad būd.

Now, this thing you have requested falls within these categories. In any
case, it has been and will remain deviation. The one prompting you has
presented this matter to you as a miracle. Besides, this is the easiest
of methods. It is surprising that it has seemed difficult to you, for
when the meaning is already present, determining the wording has always
been and will always be exceedingly easy and simple.

\subsection{543}\label{section-1038}

\begin{otherlanguage}{persian}

حال ملاحظه كن كه چه قدر بی‌‌شعورى و از اين مراتب هم گذشته و اللّه الّذى لا
اله الّا هو كه اعظم از آنچه بتصوّر آيد بچشم ديده و بگوش شنيده‌ايم و يقين
است كه شما موقن بقول اين عبد نبوده و نيستى و اگر ميبودى همان مكتوب كه
چندى قبل بشما ارسال شد حقّ را از باطل تميز ميدادى در اينصورت لازم است يا
خود بيائى و يا وكيل ديگر تعيين نمائى با جميع كتب سماوى بين يدى حاضر شود
اين عبد عهد مينمايد كه آيات الهى بشأنى نازل شود كه احدى قادر بر تحرير آن
نباشد حال اين قولى است كه خود شما ذكر نموده‌ايد و طلب کرده‌ايد بيائيد و
ملاحظه نمائيد

\end{otherlanguage}

Ḥāl mulāḥiẓih kun kih chih qadr bī-shuʿūrī. Va az īn marātib ham
gudashteh, va-Allāh alladhī lā ilāha illā Huwa kih aʿẓam az ānchih
bi-taṣavvur āyad bi-chashm dīdih va bi-gūsh shinīdih-īm. Va yaqīn ast
kih shumā mūqin bi-qawl-i īn ʿabd nabūdeh va nīstī, va agar mībūdī hamān
maktūb kih chandī qabl bi-shumā irsāl shud Ḥaqq rā az bāṭil tamyīz
mīdādī. Dar īn ṣūrat lāzim ast yā khud biyāʾī va yā vakīl-i dīgar taʿyīn
namāʾī bā jamīʿ-i kutub-i samāvī bayna yaday ḥāḍir shavad. Īn ʿabd ʿahd
mīnamāyad kih āyāt-i ilāhī bi-shaʾnī nāzil shavad kih aḥadī qādir bar
taḥrīr-i ān nabāshad. Ḥāl īn qawlī ast kih khud-i shumā dhikr namūdih-īd
va ṭalab kardih-īd. Biyāʾīd va mulāḥiẓih namāʾīd.

Now consider how lacking in understanding you are. Beyond all this, by
God, besides whom there is no other god, we have seen with our own eyes
and heard with our own ears things greater than can be imagined. It is
certain that you have never had, and still do not have, confidence in
this servant's word. Had you done so, you would have distinguished truth
from falsehood through that very letter sent to you some time ago. Under
these circumstances, you must either come yourself or appoint another
representative to appear in His presence with all the heavenly books.
This servant pledges that divine verses will be revealed in such a
manner that no one will be able to write them down. This, then, is a
proposal that you yourselves have mentioned and requested. Come and see.

\subsection{544}\label{section-1039}

\begin{otherlanguage}{persian}

گذشته از اينكه حقّ باين امور معروف نشده و نميشود و او است ميزان كلّ شیء و
خلق او ميزان معرفت او نبوده و نخواهند بود و اين مطلب از اصل صحيح نيست
چنانچه ذكر شد و آنقدر شعور نداريد كه ميفرمايد لاتجرّب الرّبّ چه كه بر حقّ
است كه ناس را تجربه و امتحان نمايد و خلق او را لايق و سزاوار نه كه او را
تجربه كنند ولكن چون تو اين را ميزان قرار داده و بميزان اللّه راضى نشده و
كفايت ننموده استدعا نموده‌ام كه از فضل و عنايت خود قبول فرمايند و قبول
فرموده‌اند چنانچه حينئذ ميفرمايند ان احضر بكتب السّمآء بين يدى اللّه ربّك و
ربّ الأوّلین‌ و الآخرین ثمّ اقرأ ما تشآء كيف تشآء لينزل من جبروت اللّه فى كلّ
كلمة كلمات عزّ منيع على شأن يعجز عن احصائها كلّ العالمين و باىّ عبارة تريد
و باىّ شأن تشآء ينزلها اللّه بالحقّ من سماء قدس رفيع ما يعجز عن تحريرها كلّ
الخلايق اجمعين

\end{otherlanguage}

Gudashteh az īnkih Ḥaqq bi-īn umūr maʿrūf nashudeh va namīshavad va Ū
ast mīzān-i kull-i shayʾ va khalq-i Ū mīzān-i maʿrifat-i Ū nabūdeh va
nakhāhand būd, va īn maṭlab az aṣl ṣaḥīḥ nīst chinānchih dhikr shud. Va
ānqadr shuʿūr nadārīd kih mīfarmāyad lā tujarrib al-Rabb, chih kih bar
Ḥaqq ast kih nās rā tajribih va imtiḥān namāyad va khalq-i Ū rā lāyiq va
sizāvār nah kih Ū rā tajribih kunand. Va-lākin chūn tu īn rā mīzān qarār
dādih va bi-mīzān Allāh rāḍī nashudeh va kifāyat nanamūdeh, istidʿā
namūdih-am kih az faḍl va ʿināyat-i khud qabūl farmāyand va qabūl
farmūdih-and, chinānchih ḥīnaʾidh mīfarmāyand: an iḥḍur bi-kutub
al-samāʾ bayna yaday Allāh Rabbika wa-Rabb al-awwalīn wa-al-ākhirīn
thumma iqraʾ mā tashāʾ kayfa tashāʾ li-yanzila min jabarūt Allāh fī kull
kalimah kalimāt ʿizz manīʿ ʿalā shaʾn yaʿjizu ʿan iḥṣāʾihā kull
al-ʿālamīn wa-bi-ayyi ʿibārah turīd wa-bi-ayyi shaʾn tashāʾ yunazziluhā
Allāh bi-al-ḥaqq min samāʾ quds rafīʿ mā yaʿjizu ʿan taḥrīrihā kull
al-khalāʾiq ajmaʿīn.

Quite apart from the fact that the True One has never been and will
never be known through such things, He is the measure of all things, and
His creatures have never been and will never be the measure by which He
is known. As already stated, this proposal is fundamentally unsound. You
do not even have enough understanding to recognize that He says, ``Do
not put the Lord to the test,'' for it is the prerogative of the True
One to test and try people, and it is neither fitting nor proper for His
creatures to test Him. Yet because you have made this your measure, have
not been content with God's measure, and have not found it sufficient, I
have asked Him, out of His grace and loving care, to accept, and He has
accepted. Thus He now says: ``Come with the books of heaven into the
presence of God, your Lord and the Lord of the first and the last. Then
read whatever you wish, however you wish, so that from the dominion of
God there may descend, for every word, words of unassailable glory in
such a manner that all the worlds will be unable to count them. In
whatever wording you desire and in whatever manner you wish, God will
send them down in truth from a heaven of exalted holiness---words that
all created beings together will be unable to write down.''

\subsection{545}\label{section-1040}

\begin{otherlanguage}{persian}

اين كلمات محكمات در جواب مطلب شما نازل و لكن فرموده‌اند كه يقين است بعذر
ديگر متعذّر خواهند شد چه كه كلمه حقّ لايزيد المشركين الّا خساراً نميدانم
ديگر بچه عذرى متمسّك خواهيد شد و بچه ميزانى متشبّث و اللّه يا عبد از شاطى
انصاف بسيار بعيدى و از مقرّ قرب بسيار دور نميدانم چه شما را بر آن داشته
مع آنكه بين يدى حاضر نشده و خدمت ظهور قبل هم فايز نگشته و قبل از تجربه
بقول خودت باينكلمات مفتريات مشغول گشته و كسانى را كه مطّلع بر امرند ايشان
را تكذيب نموده‌ايد و نفسى را كه مثل تو بوده و ابدا اطّلاع از جائى نداشته
او را تصديق نموده

\end{otherlanguage}

Īn kalimāt-i muḥkamāt dar javāb-i maṭlab-i shumā nāzil, va-lākin
farmūdih-and kih yaqīn ast bi-ʿudhr-i dīgar mutaʿadhdhir khāhand shud,
chih kih kalimih-yi Ḥaqq lā yazīdu al-mushrikīn illā khasāran. Namīdānam
dīgar bi-chih ʿudhrī mutamassik khāhīd shud va bi-chih mīzānī
mutashabbith. Va-Allāh yā ʿabd, az shāṭī-yi inṣāf bisyār baʿīdī va az
maqarr-i qurb bisyār dūr. Namīdānam chih shumā rā bar ān dāshteh, maʿa
ānkih bayna yaday ḥāḍir nashudeh va khidmat-i Ẓuhūr-i qabl ham fāʾiz
nagashteh va qabl az tajribih bi-qawl-i khudat bi-īn kalimāt-i
muftariyāt mashghūl gashteh va kasānī rā kih muṭṭaliʿ bar Amr-and īshān
rā takdhīb namūdih-īd va nafsī rā kih mithl-i tu būdeh va abadan iṭṭilāʿ
az jāʾī nadāshteh ū rā taṣdīq namūdeh.

These firm words have been revealed in response to your request. Yet He
has said that you will certainly resort to another excuse, for the word
of the True One ``increases those who associate partners with God in
nothing but loss.'' I do not know what further excuse you will cling to,
or what measure you will seize upon. By God, O servant, you are far
indeed from the shore of fairness and far removed from the seat of
nearness. I do not know what has prompted you to do this, when you have
neither appeared in His presence nor attained the presence of the
previous Manifestation. Before conducting a test, by your own account,
you have occupied yourself with these slanderous words, denied those who
are informed about the Cause, and believed someone who was like yourself
and had no knowledge from any quarter.

\subsection{546}\label{section-1041}

\begin{otherlanguage}{persian}

و از اينمراتب گذشته يكذّبك كلماتك بعضها بعض چنانچه در اينمقام نوشته يقين
هركه داراى اين مقام است حقّ و كلام او حقّ و غير او كاذب و بعد نوشته كه اگر
اتيان باين امر هم نمايد آنوقت هم دليل بر آن نفس موعود نيست و يكى از
اولياى دوره بيان خواهد بود و آنهم بعد از تصديق ميرزا يحيى حال خود انصاف
ده اگر يقين حقّ است ديگر چگونه منوط بتصديق دونش خواهد شد حال از اين دو
بيرون نيست يا كلمه اوّل را بايد اقرار نمائى كه لغو بوده و من غير شعور ذكر
نموده و يا بيان ثانى را و اللّه علامت قهر در امثال آن نفوس ظاهر و مشهود
شده چه كه شاعر نيستيد كه چه مينويسيد و چه ذكر مينمائيد

\end{otherlanguage}

Va az īn marātib gudashteh yukadhdhibuka kalimātuka baʿḍuhā baʿḍ,
chinānchih dar īn maqām nivishteh: yaqīn har kih dārā-yi īn maqām ast
ḥaqq va kalām-i ū ḥaqq va ghayr-i ū kādhib. Va baʿd nivishteh kih agar
ityān bi-īn amr ham namāyad ān vaqt ham dalīl bar ān nafs-i mawʿūd nīst
va yikī az awliyā-yi dawrih-yi Bayān khāhad būd va ān ham baʿd az
taṣdīq-i Mīrzā Yaḥyā. Ḥāl khud inṣāf deh: agar yaqīn Ḥaqq ast dīgar
chigūnih manūṭ bi-taṣdīq-i dūnash khāhad shud? Ḥāl az īn du bīrūn nīst:
yā kalimih-yi avval rā bāyad iqrār namāʾī kih laghv būdeh va min ghayr-i
shuʿūr dhikr namūdeh va yā bayān-i thānī rā. Va-Allāh ʿalāmat-i qahr dar
amthāl-i ān nufūs ẓāhir va mashhūd shudeh, chih kih shāʿir nīstīd kih
chih mīnivīsīd va chih dhikr mīnamāʾīd.

Beyond all this, your own words prove you false, contradicting one
another. Thus, at this point you have written: ``Certainly, whoever
possesses this station is true, and his words are true, and anyone else
is a liar.'' Then you have written that even if He accomplishes this, it
will still not be proof that He is that promised Person, and He will be
one of the holy ones of the dispensation of the Bayān---and that only
after Mīrzā Yaḥyā's confirmation. Now judge fairly for yourself: if He
is certainly the True One, how can He depend upon the confirmation of
anyone besides Him? There are only two alternatives: you must admit
either that the first statement was idle talk uttered without
understanding, or that the second statement was. By God, the sign of
wrath has become manifest and evident in such souls, for you are unaware
of what you write and what you say.

\subsection{547}\label{section-1042}

\begin{otherlanguage}{persian}

بدان اليوم آنچه بعقلت رسيده و برسد و يا بعقول فوق تو و دون تو ادراك شود
هيچيك ميزان معرفت حقّ نبوده و نخواهد بود و حقّ منيع از جميع مقدّس و منزّه و
مبرّا و اگر نفسى باعلى رتبه عرفان صعود نمايد و يا باعلى ذروه بيان ارتقاء
جويد كه نفع نمى‌بخشد او را و ثمرى نداشته و نخواهد داشت الّا بعد از اقرار
باين ظهور اعظم كذلك نزل الأمر فى البيان ثمّ فى الفرقان ثمّ فى الانجيل و
التّوريّة و الزّبور و الصّحف ان انتم توقنون

\end{otherlanguage}

Bidān al-yawm ānchih bi-ʿaqlat rasīdih va birasad va yā bi-ʿuqūl-i
fawq-i tu va dūn-i tu idrāk shavad hīch-yik mīzān-i maʿrifat-i Ḥaqq
nabūdih va nakhvāhad būd va Ḥaqq-i manīʿ az jamīʿ muqaddas va munazzah
va mubarrā va agar nafsī bi-aʿlā rutbiy-i ʿirfān ṣuʿūd namāyad va yā
bi-aʿlā dhurviy-i bayān irtiqā jūyad kih nafʿ namībakhshad ū rā va
thamarī nadāshtih va nakhvāhad dāsht illā baʿd az iqrār bi-īn ẓuhūr-i
aʿẓam. Kadhālika nazala al-amru fī al-Bayān thumma fī al-Furqān thumma
fī al-Injīl wa-al-Tawrāt wa-al-Zabūr wa-al-Ṣuḥuf in antum tūqinūn.

Know that today, whatever your mind has grasped or may grasp, or
whatever minds above or below yours may comprehend, has never been and
will never be the measure of the knowledge of the True One. The
inaccessible True One is sanctified, exalted, and free from all of
these. Even if a soul ascends to the highest rank of spiritual knowledge
or rises to the loftiest summit of utterance, this will not benefit that
soul and has never borne, nor will ever bear, any fruit except after
acknowledgment of this Most Great Revelation. Thus was the command
revealed in the Bayān, then in the Qur'an, then in the Gospel, the
Torah, the Psalms, and the Scriptures, if you are certain.

\subsection{548}\label{section-1043}

\begin{otherlanguage}{persian}

قوله كبر كذبه و جعله انصاف بدهيد كسيكه بيست سال در اين دوره پرورش يافته
بيست سال هم در دوره عرفا چهار كلمه بل هزار كلمه مطالب عرفان و صورت آيات
بنويسد كسبيّه چه مدخليّت دارد با فطرت فرق اين دو همان است كه عرض شد و اگر
هم بهمان نوع نوشتند و حال آنكه مشكل است آنوقت هم دليل بر آن شخص موعود
نيست و دليل بر انكار و تكفير ميرزا يحيى نيست بلكه يكى از اولياى دوره
بيان هستند آنهم بشرط تصديق ثمره نه بانكار و ظلم بر او انتهى

\end{otherlanguage}

Qawluhu kabura kadhibuhu va jaʿluhu: inṣāf bidahīd kasī kih bīst sāl dar
īn dawrih parvarish yāftih bīst sāl ham dar dawriy-i ʿurafā chahār
kalamih bal hazār kalamih maṭālib-i ʿirfān va ṣūrat-i āyāt binivīsad
kasbiyyih chih madkhaliyyat dārad bā fiṭrat farq-i īn du hamān ast kih
ʿarḍ shud va agar ham bi-hamān nawʿ nivishtand va ḥāl ānkih mushkil ast
ān vaqt ham dalīl bar ān shakhṣ-i mawʿūd nīst va dalīl bar inkār va
takfīr-i Mīrzā Yaḥyā nīst balkih yakī az awliyā-yi dawriy-i Bayān
hastand ān ham bi-sharṭ-i taṣdīq-i thumrih nah bi-inkār va ẓulm bar ū.
Intahā.

His words---great are his falsehood and fabrication: ``Be fair: someone
who has been nurtured for twenty years in this dispensation and twenty
years in the circle of the mystics may write four words, indeed a
thousand words, on mystical subjects and in the form of verses. What
does acquired learning have to do with innate capacity? The difference
between these two is just what was explained. Even if He were to write
in that same manner---which would be difficult---it would still not
prove that He is that promised Person, nor would it justify denying
Mīrzā Yaḥyā and declaring him an unbeliever. Rather, He would be one of
the holy ones of the dispensation of the Bayān, and that only on
condition of acknowledging the Fruit, not denying him and wronging
him.'' End of quotation.

\subsection{549}\label{section-1044}

\begin{otherlanguage}{persian}

اوّلا اى مفترى كاذب كسبيّت از كجا معلوم شد پيش كه درس خوانده‌اند و معلّم كه
بوده و ثانيا آن عرفائيكه با آن جمال اقدس معاشر بوده‌اند كه بوده يكذّبك
نفسك فيما افتريت و لكن لا تستشعر و تكوننّ من الغافلين

\end{otherlanguage}

Avvalan ay muftarī-yi kādhib kasbiyyat az kujā maʿlūm shud pīsh-i kih
dars khvāndih-and va muʿallim kih būdih va thāniyan ān ʿurafāʾī kih bā
ān Jamāl-i Aqdas muʿāshir būdih-and kih būdih. Yukadhdhibuka nafsuka
fīmā iftarayta wa-lākin lā tastashʿiru wa-takūnanna min al-ghāfilīn.

First, you lying slanderer, how was it established that His knowledge
was acquired? With whom did He study, and who was His teacher? Second,
who were those mystics who associated with that Most Holy Beauty? Your
own self gives the lie to what you have fabricated, but you do not
perceive it and are indeed among the heedless.

\subsection{550}\label{section-1045}

\begin{otherlanguage}{persian}

آنچه مبرهن و واضح اينست كه تا والدشان در حيات بود با احدى معاشر نبودند
مگر گاهى بنفوسيكه در خدمت والدشان مراوده مينمودند و جميع اهل ايران
مطّلعند كه والدشان هم بحسب ظاهر از اهل علم نبوده‌اند تفحّص لعلّ تجد الى الحقّ
سبيل

\end{otherlanguage}

Ānchih mubarhan va vāḍiḥ īnast kih tā vālidishān dar ḥayāt būd bā aḥadī
muʿāshir nabūdand magar gāhī bi-nufūsī kih dar khidmat-i vālidishān
murāvadih mīnamūdand va jamīʿ-i ahl-i Īrān muṭṭaliʿand kih vālidishān
ham bi-ḥasab-i ẓāhir az ahl-i ʿilm nabūdih-and. Tafaḥḥaṣ laʿalla tajidu
ilā al-Ḥaqq sabīl.

What is established and clear is that, while His father was alive, He
associated with no one except occasionally those who visited His father.
All the people of Iran know that His father, too, was not outwardly one
of the learned. Investigate, so that perhaps you may find a path to the
Truth.

\subsection{551}\label{section-1046}

\begin{otherlanguage}{persian}

و اگر هم بعضى از عرفا به بيت ميامدند و ملاقات واقع ميشد اين دخلى بتعليم
و تعلّم نداشته و اين كلماتى است كه ميرزا يحيى القا نموده و اگر بر فرض
هم‌چنين بوده جميع ذرّات شهادت ميدهند كه ميرزا يحيى در خدمت ايشان بوده در
هر حال او هم شريك بوده

\end{otherlanguage}

Va agar ham baʿḍī az ʿurafā bi-bayt mīāmadand va mulāqāt vāqiʿ mīshud īn
dakhlī bi-taʿlīm va taʿallum nadāshtih va īn kalimātī ast kih Mīrzā
Yaḥyā ilqā namūdih va agar bar farḍ ham-chunīn būdih jamīʿ-i dharrāt
shahādat mīdihand kih Mīrzā Yaḥyā dar khidmat-i īshān būdih dar har ḥāl
ū ham sharīk būdih.

Even if some mystics came to the house and meetings took place, this had
nothing to do with teaching and learning. These are words that Mīrzā
Yaḥyā has suggested. Even supposing it had been so, every atom bears
witness that Mīrzā Yaḥyā was in His service: in every circumstance, he
too shared in it.

\subsection{552}\label{section-1047}

\begin{otherlanguage}{persian}

حال تفكّر كن اى بی‌‌انصاف كه چه گفته و چه ميگوئى در دوره بيان با كه معاشر
بوده‌اند هر نفسى كه بين يدى حاضر ميشد كسب معارف الهيّه و شئونات حكمت
ربّانيه مينمود و ابدا در حضور امنع اقدس نفسى قادر بر تكلّم نبوده تا چه رسد
بآنكه از او چيزى اخذ فرمايند چنانچه اين عباد در عراق بوده‌ايم و بچشم خود
ديده‌ايم فاسئل من اهل العراق لعلّ تطّلع و تكون من العارفين

\end{otherlanguage}

Ḥāl tafakkur kun ay bī-inṣāf kih chih guftih va chih mīgūʾī dar dawriy-i
Bayān bā kih muʿāshir būdih-and har nafsī kih bayn-i yaday ḥāḍir mīshud
kasb-i maʿārif-i ilāhiyyih va shuʾūnāt-i ḥikmat-i rabbāniyyih mīnamūd va
abadan dar ḥuḍūr-i amnaʿ-i aqdas nafsī qādir bar takallum nabūdih tā
chih rasad bi-ānkih az ū chīzī akhdh farmāyand chunānchih īn ʿibād dar
ʿIrāq būdih-īm va bi-chashm-i khud dīdih-īm. Fa-sʾal min ahl al-ʿIrāq
laʿalla taṭṭaliʿu wa-takūnu min al-ʿārifīn.

Now reflect, you unfair one, on what you have said and what you are
saying. With whom did He associate during the dispensation of the Bayān?
Every soul who entered His presence acquired divine knowledge and the
manifold expressions of heavenly wisdom. In His most inaccessible, most
holy presence, no one was ever able to speak, much less could He have
learned anything from them. We servants were in Iraq and saw this with
our own eyes. Ask the people of Iraq, so that perhaps you may learn and
be among those who recognize the truth.

\subsection{553}\label{section-1048}

\begin{otherlanguage}{persian}

و از اين گذشته از اوّل ظهور نقطه بيان روح ما سواه فداه تا حين آسایش نبوده
چنانچه در اكثر اوقات در حبس و در دست اعدا مبتلا بوده‌اند و بعد از والد
مخصوص ميرزا يحيى ملّا مهدى كندى را آورده و معلّم او قرار فرمودند كه روزى
يكساعت ميامد و بتعليم او مشغول بود حال از ارض طاء تفحّص نمائيد اگر نفسى
يافت شد كه بگويد درس خوانده و يا كسب نموده‌اند حقّ با شما است و

\end{otherlanguage}

Va az īn gudhashtih az avval-i ẓuhūr-i Nuqṭiy-i Bayān rūḥ-i mā sivāh
fidāh tā ḥīn āsāyish nabūdih chunānchih dar akthar-i awqāt dar ḥabs va
dar dast-i aʿdā mubtalā būdih-and va baʿd az vālid makhṣūṣ-i Mīrzā Yaḥyā
Mullā Mihdī-yi Kandī rā āvardih va muʿallim-i ū qarār farmūdand kih rūzī
yak sāʿat mīāmad va bi-taʿlīm-i ū mashghūl būd ḥāl az arḍ-i Ṭāʾ tafaḥḥuṣ
namāyīd agar nafsī yāft shud kih bigūyad dars khvāndih va yā kasb
namūdih-and Ḥaqq bā shumā ast va

Beyond this, from the beginning of the revelation of the Point of the
Bayān---may the spirit of all else be sacrificed for Him---until now, He
has had no rest, having spent most of His time afflicted by imprisonment
and in the hands of enemies. After His father's death, He brought Mullā
Mihdī Kandī specifically for Mīrzā Yaḥyā and appointed him as his
teacher; he would come for an hour each day and teach him. Now make
inquiries in the land of Ṭāʾ. If anyone is found who says that He
studied or acquired learning, you are right, and

\subsection{554}\label{section-1049}

\begin{otherlanguage}{persian}

حال آنكه هيچ يك از اين امور دليل بر امرى نبوده و نخواهد بود خوب در قلبت
القاى بغضا نموده‌اند و خوب در صدرت مفتريات نقش بسته در ظهور نقطه بيان هم
بعينه همين كلمات از مشركين ظاهر شده كه با شيخيها مراوده فرموده‌اند و آنچه
نوشته از كلمات شيخ احمد اخذ نموده‌ و هم‌چنين نسبت دادند كه نزد فلان در ارض
شين تحصيل نموده‌و علم ظاهريّه كسب فرموده لعن اللّه من الهمك و القى فى صدرك
ما كان به مستحقّاً عن دونه

\end{otherlanguage}

Ḥāl ānkih hīch yak az īn umūr dalīl bar amrī nabūdih va nakhvāhad būd
khūb dar qalbat ilqā-yi baghḍā namūdih-and va khūb dar ṣadrat muftariyāt
naqsh bastih dar ẓuhūr-i Nuqṭiy-i Bayān ham bi-ʿaynih hamīn kalimāt az
mushrikīn ẓāhir shudih kih bā Shaykhīhā murāvadih farmūdih-and va ānchih
nivishtih az kalimāt-i Shaykh Aḥmad akhdh namūdih va ham-chunīn nisbat
dādand kih nazd-i fulān dar arḍ-i Shīn taḥṣīl namūdih va ʿulūm-i
ẓāhiriyyih kasb farmūdih. Laʿana Allāhu man alhamaka wa-alqā fī ṣadrika
mā kāna bihi mustaḥiqqan ʿan dūnihi.

Yet none of these things has ever been, or will ever be, proof of
anything. How thoroughly they have instilled hatred in your heart, and
how firmly slanders have been imprinted in your breast! At the
revelation of the Point of the Bayān, these very same words came from
those who associated partners with God: that He had associated with the
Shaykhīs and had taken what He wrote from the words of Shaykh Aḥmad.
They also alleged that He had studied under someone in the land of Shīn
and acquired outward learning. May God curse the one who prompted you
and cast into your breast that for which he, rather than anyone else,
deserved this curse.

\subsection{555}\label{section-1050}

\begin{otherlanguage}{persian}

و ديگر صورت آيات ذكر نموده ‌اين قولى است كه مدّتها است مشركين بآن ناطقند
گاهى قشر و لبّ ذكر نموده‌گاهى صورت و هيولى در آيات ذكر کرده‌اند كلّما خرج من
افواههم يرجع الى انفسهم و انّ ربّك الرّحمن لبرئ منهم و عما يقولون اى بيچاره
اعمى تو فرق ما بين يمين و يسار ننموده‌چگونه مميّز فطرت و غير فطرت شده حقّى
كه مميّزش تو باشى همان لايق و سزاوار تو است لو يوجد من بصر ليعرف و يوقن
بأن لايعادل بآية عمّا نزل بالحقّ كلمات العالمين جميع ناس كه از شاطى بحر
معرفت رحمانى محروم شده‌اند سبب آن بوده كه عرفان خود را ميزان عرفان حقّ
قرار داده بودند لذا از كوثر معانى محروم و ممنوع گشته‌اند چنانچه هر ملّتى
اليوم بادلّه و براهين خود منتظرند و در ظهورات مظاهر احديّه توقّف نموده‌اند و
اين حال نفوسى است كه فى‌الجمله اراده عرفان داشته‌اند و لكن امثال شما كه
ببغضاى صرف و غلّ بحت بر وجه حقّ ايستاده‌ايد ذكرى عند اللّه نداشته و نداريد

\end{otherlanguage}

Va dīgar ṣūrat-i āyāt dhikr namūdih īn qawlī ast kih muddat-hā ast
mushrikīn bi-ān nāṭiqand gāhī qishr va lubb dhikr namūdih gāhī ṣūrat va
hayūlā dar āyāt dhikr kardih-and. Kullamā kharaja min afwāhihim yarjiʿu
ilā anfusihim wa-inna rabbaka al-Raḥmāna la-barīʾun minhum wa-ʿammā
yaqūlūn. Ay bīchāriy-i aʿmā tu farq mā bayn-i yamīn va yasār nanamūdih
chigūnih mumayyiz-i fiṭrat va ghayr-i fiṭrat shudih Ḥaqqī kih
mumayyizash tu bāshī hamān lāyiq va sazāvār-i tu ast. Law yūjadu min
baṣar la-yaʿrifa wa-yūqina bi-an lā yuʿādalu bi-āyatin ʿammā nazala
bi-al-Ḥaqq kalimātu al-ʿālamīn. Jamīʿ-i nās kih az shāṭī-yi baḥr-i
maʿrifat-i Raḥmānī maḥrūm shudih-and sabab ān būdih kih ʿirfān-i khud rā
mīzān-i ʿirfān-i Ḥaqq qarār dādih būdand lidhā az kawthar-i maʿānī
maḥrūm va mamnūʿ gashtih-and chunānchih har millatī al-yawm bi-adillih
va barāhīn-i khud muntaẓirand va dar ẓuhūrāt-i maẓāhir-i aḥadiyyih
tavaqquf namūdih-and va īn ḥāl-i nufūsī ast kih fī-al-jumlih irādiy-i
ʿirfān dāshtih-and va lākin amthāl-i shumā kih bi-baghḍā-yi ṣirf va
ghill-i baḥt bar vajh-i Ḥaqq īstādih-īd dhikrī ʿinda Allāh nadāshtih va
nadārīd.

Furthermore, you have mentioned ``the form of verses.'' This is
something that those who associate partners with God have long been
saying. Sometimes they have spoken of shell and kernel, and sometimes of
form and matter, in connection with the verses. Everything that comes
from their mouths returns upon themselves. Your Lord, the All-Merciful,
is indeed free of them and of what they say. You miserable blind one!
You have not distinguished right from left; how have you become a judge
of what is innate and what is not? A true one whose judge you are is
indeed fit and worthy of you. If any sight were to be found, it would
recognize and be certain that all the words of the world cannot equal a
single verse of what has been revealed by the True One. All those who
have been deprived of the shore of the ocean of the All-Merciful's
knowledge have been deprived because they made their own understanding
the measure of the knowledge of the True One. Thus they have been barred
and excluded from the heavenly waters of meaning. Today every religious
community waits according to its own proofs and arguments and has
hesitated at the revelations of the Manifestations of Divine Unity. This
is the condition of souls who have had at least some desire for
recognition. But people like you, who have stood against the face of the
True One with sheer hatred and pure malice, have never had, and do not
have, any standing before God.

\subsection{556}\label{section-1051}

\begin{otherlanguage}{persian}

و كاش نداى اركانت را ميشنيدى كه در حين تحرير كلماتت مخاطباً اليك چه ميگفت
حتّى شعرى كه در بدنت موجود پناه بخدا برده از شرّ و مكر و خدعه تو بدانكه حقّ
منيع در كلّ يوم در شأن بديع بوده و عرفان ناس ميزان عرفان او نشده و نخواهد
شد بسا از نفوس كه در ظهور قبل از اعلى سدره رضوان محسوب و در ظهور بعد از
ادنى شجره سجّين مذكور چنانچه در اينمقام نقطه بيان روح ما سواه فداه تصريح
فرموده و از قبل ذكر شده

\end{otherlanguage}

Va kāsh nidā-yi arkānat rā mīshinīdī kih dar ḥīn-i taḥrīr-i kalimātat
mukhāṭaban ilayk chih mīguft ḥattā shaʿrī kih dar badanat mawjūd panāh
bi-Khudā burdih az sharr va makr va khudʿiy-i tu bidānkih Ḥaqq-i manīʿ
dar kull-i yawm dar shaʾn-i badīʿ būdih va ʿirfān-i nās mīzān-i ʿirfān-i
ū nashudih va nakhvāhad shud basā az nufūs kih dar ẓuhūr-i qabl az aʿlā
sidriy-i Riḍvān maḥsūb va dar ẓuhūr-i baʿd az adnā shajariy-i Sijjīn
madhkūr chunānchih dar īn maqām Nuqṭiy-i Bayān rūḥ-i mā sivāh fidāh
taṣrīḥ farmūdih va az qabl dhikr shudih.

If only you could hear the cry of your limbs, and what they were saying
to you as you wrote your words! Even a hair on your body has sought
refuge in God from your evil, your scheming, and your deceit. Know that
the inaccessible True One appears each day in a wondrous new way, and
human understanding has never been, nor will ever be, the measure of the
knowledge of Him. How many souls were counted among the loftiest trees
of Paradise in the previous revelation, but were named among the
lowliest trees of hell in the revelation that followed! The Point of the
Bayān---may the spirit of all else be sacrificed for Him---has
explicitly stated this matter, as has already been mentioned.

\subsection{557}\label{section-1052}

\begin{otherlanguage}{persian}

و اينكه از قلم جعل و كذب نوشته شده كه دليل بر انكار و تكفير ميرزا يحيى
نبوده الخ اذا يتكلّم لسان اللّه بأن لعن اللّه اوّل ظالم ظلم نقطة الأولی فى
ظهوره الأخرى و كفر بآياته و انكر بظهوره و جاحد ببرهانه و نازع بسلطانه و
حارب بنفسه و اراد سفك دمه بعد الّذى جائهم بآيات عزّ مبين يشهد كلّ الذّرّات
بانّه ما اعترض على احد و كان ارحم بكلّ نفس من نفسه على نفسه و يشهد بذلك كلّ
شیء ان انتم تنكرون

\end{otherlanguage}

Va īnkih az qalam-i jaʿl va kadhib nivishtih shudih kih dalīl bar inkār
va takfīr-i Mīrzā Yaḥyā nabūdih ilā ākhir. Idhan yatakallamu lisānu
Allāh bi-an laʿana Allāhu awwala ẓālimin ẓalama Nuqṭata al-Ūlā fī
ẓuhūrihi al-ukhrā wa-kafara bi-āyātihi wa-ankara bi-ẓuhūrihi wa-jāḥada
bi-burhānihi wa-nāzaʿa bi-sulṭānihi wa-ḥāraba bi-nafsihi wa-arāda safka
damihi baʿda alladhī jāʾahum bi-āyāti ʿizzin mubīn. Yashhadu kullu
al-dharrāt bi-annahu mā iʿtaraḍa ʿalā aḥadin wa-kāna arḥama bi-kulli
nafsin min nafsihi ʿalā nafsihi wa-yashhadu bi-dhālika kullu shayʾin in
antum tunkirūn.

As for what has been written by the pen of fabrication and falsehood,
that this would not justify denying Mīrzā Yaḥyā and declaring him an
unbeliever, and so forth: the Tongue of God now declares, ``May God
curse the first oppressor who wronged the Primal Point in His latter
revelation, disbelieved in His verses, denied His revelation, rejected
His proof, disputed His sovereignty, fought against His own Person, and
sought to shed His blood after He had come to them with verses of
manifest glory.'' Every atom bears witness that He opposed no one and
was more merciful to every soul than that soul was to itself. Everything
bears witness to this, even if you deny it.

\subsection{558}\label{section-1053}

\begin{otherlanguage}{persian}

بحقّ خودشان كه با كلّ بكمال محبّت رفتار فرموده‌اند تا آنكه بعد از تفصيل كبرى
نوشته‌جاتى از ميرزا يحيى مشاهده شد و تفصيل آن مجملا از قبل ذكر شده و در
آن نوشته‌جات جميع شرور اوّلين و آخرين را بحقّ نسبت داده بعد از مشاهده آن
الواح مفتريه مجعوله كاذبه فى ردّ من ردّ على اللّه نازل شد آنچه نازل شد و
هذا لم يكن من عنده بل من عند اللّه المنزل العزيز الكريم

\end{otherlanguage}

Bi-ḥaqq-i khudishān kih bā kull bi-kamāl-i maḥabbat raftār farmūdih-and
tā ānkih baʿd az tafṣīl-i kubrā nivishtih-jātī az Mīrzā Yaḥyā mushāhadih
shud va tafṣīl-i ān mujmalan az qabl dhikr shudih va dar ān
nivishtih-jāt jamīʿ-i shurūr-i avvalīn va ākhirīn rā bi-Ḥaqq nisbat
dādih baʿd az mushāhadiy-i ān alvāḥ-i muftariyiy-i majʿūliy-i kādhibih
fī raddi man radda ʿalā Allāh nāzil shud ānchih nāzil shud. Wa-hādhā lam
yakun min ʿindihi bal min ʿindi Allāhi al-Munzili al-ʿAzīzi al-Karīm.

By His own truth, He treated everyone with perfect love, until, after
the great affair, some writings by Mīrzā Yaḥyā were seen. An outline of
that affair has already been given. In those writings he attributed to
the True One all the evils of those who came before and those who came
after. After those slanderous, fabricated, lying writings were seen,
what was revealed was revealed in refutation of the one who had rejected
God. This was not from Him, but from God, the Revealer, the Almighty,
the Generous.

\subsection{559}\label{section-1054}

\begin{otherlanguage}{persian}

اى بی‌ انصاف و اللّه الّذى لا اله الّا هو كه اگر ذرّه‌ئی در قلبت حبّ الهى ميبود
هرآينه باو مينوشتى كه بچه دليل و برهان درصدد قطع شجره ربّانيّه افتادى و
اعتراض نمودى بكينونتى كه ميفرمايد يا قوم انّى حجّة اللّه بينكم و برهانه
فيكم و رحمة عليكم و قد جئتكم عن افق المعانى و البيان بنبأ قد كان على
الحقّ عظيماً قد بعثنى اللّه و ارسلنى بآيات الّتى ارسل بها كلّ النّبيين و
المرسلين ثمّ علىّ بالحقّ فباىّ حجّة تكفرون هذا الأمر الّذى اشرق عن افق القدس
بسلطان مبينا يا قوم هل تعترضون بالّذى به ثبت ايمانكم و ايمان كلّ مؤمن و به
رفعت اسمائكم و علت مراتبكم و ظهرت كلمة اللّه فى الآفاق و كان اللّه على ذلك
شهيدا

\end{otherlanguage}

Ay bī-inṣāf va-Allāhi alladhī lā ilāha illā huwa kih agar dharrih-ī dar
qalbat ḥubb-i ilāhī mībūd har-āyinih bi-ū mīnivishtī kih bi-chih dalīl
va burhān dar ṣadad-i qaṭʿ-i shajariy-i rabbāniyyih uftādī va iʿtirāḍ
namūdī bi-kaynūnatī kih mīfarmāyad: Yā qawmi innī ḥujjatu Allāhi
baynakum wa-burhānuhu fīkum wa-raḥmatun ʿalaykum wa-qad jiʾtukum ʿan
ufuqi al-maʿānī wa-al-bayān bi-nabaʾin qad kāna ʿalā al-Ḥaqq ʿaẓīman qad
baʿathanī Allāhu wa-arsalanī bi-āyāti allatī arsala bihā kulla
al-nabiyyīn wa-al-mursalīn thumma ʿalayya bi-al-ḥaqq fa-bi-ayyi ḥujjatin
takfurūna hādhā al-amra alladhī ashraqa ʿan ufuqi al-quds bi-sulṭānin
mubīnan yā qawmi hal taʿtariḍūna bi-alladhī bihi thabata īmānukum
wa-īmānu kulli muʾminin wa-bihi rufiʿat asmāʾukum wa-ʿalat marātibukum
wa-ẓaharat kalimatu Allāhi fī al-āfāq wa-kāna Allāhu ʿalā dhālika
shahīdan.

You unfair one! By God, besides whom there is no god, if there were an
atom of divine love in your heart, you would surely have written to him:
``By what proof and argument did you set about cutting down the divine
Tree and object to a Being who says: `My people, I am God's testimony
among you, His proof in your midst, and a mercy to you. I have come to
you from the horizon of meanings and utterance with tidings that are
momentous in the sight of the True One. God has raised Me up and sent Me
with the verses with which He sent all the Prophets and Messengers, and
then bestowed them upon Me in truth. By what argument, then, do you
reject this Cause, which has dawned from the horizon of holiness with
manifest authority? My people, do you oppose the One through whom your
faith and the faith of every believer have been established, through
whom your names have been exalted and your ranks raised, and through
whom the Word of God has appeared throughout the world? God is witness
to this.'\,''

\subsection{560}\label{section-1055}

\begin{otherlanguage}{persian}

مع‌ذلك باو كلمه نصحى للّه ذكر ننمودى و سيفى برداشته و بفعل او مشغول شده و
اين مقام اعتراض بر او بود نه بر حقّ و اللّه عامل شده آنچه را كه احدى عامل
نشده نفسى كه فتوى بر قتل كلّ نبيّين داده او را حقّ دانسته و بالّذى منه ظهر
الحقّ و به حقّق الحقّ اعتراض نموده ‌فوا حسرة عليك و بما ارتكبت فى الحيوة
الباطلة فوا اسفا عليك بما ضيّعت عمرك و ما كنت من المستشعرين

\end{otherlanguage}

Maʿa dhālik bi-ū kalimiy-i nuṣḥī li-llāh dhikr nanamūdī va sayfī
bardāshtih va bi-fiʿl-i ū mashghūl shudih va īn maqām-i iʿtirāḍ bar ū
būd nah bar Ḥaqq va-Allāh ʿāmil shudih ānchih rā kih aḥadī ʿāmil
nashudih nafsī kih fatvā bar qatl-i kull-i nabiyyīn dādih ū rā Ḥaqq
dānistih va bi-alladhī minhu ẓahara al-Ḥaqq wa-bihi ḥuqqiqa al-Ḥaqq
iʿtirāḍ namūdih. Fa-wā ḥasratan ʿalayka wa-bi-mā irtakabta fī al-ḥayāti
al-bāṭilah fa-wā asafan ʿalayka bi-mā ḍayyaʿta ʿumraka wa-mā kunta min
al-mustashʿirīn.

Nevertheless, you did not offer him a word of counsel for God's sake,
but took up a sword and occupied yourself with his deeds. This was the
occasion to object to him, not to the True One. By God, he has done what
no one has done! You have regarded as true a person who has issued a
verdict for the killing of all the Prophets, and have opposed the One
from whom the Truth has appeared and through whom the Truth has been
established. Alas for you and for what you have committed in this vain
life! Alas for you, for you have wasted your life and have not been
among those who perceive.

\subsection{561}\label{section-1056}

\begin{otherlanguage}{persian}

هر نفسيكه اقلّ من ذرّ شاعر باشد يقين مينمايد كه اگر ايشان اراده ميفرمودند
هرگز او موجود نبود بيست سنه در اين امر و بيست سنه قبل از اين امر بكمال
شفقت و مهربانى در حفظ و تعليمش جهد بليغ فرمودند ثمره آن شد كه شنيده‌ايد
قد بكت السّموات و الأرض ثمّ الهوآء و المياه و الأحجار و الأشجار و الأثمار
و الأوراق و لكن المشركين هم فى فرح عظيم

\end{otherlanguage}

Har nafsī kih aqall-i min dharr shāʿir bāshad yaqīn mīnamāyad kih agar
īshān irādih mīfarmūdand hargiz ū mawjūd nabūd bīst sanih dar īn amr va
bīst sanih qabl az īn amr bi-kamāl-i shafaqat va mihrabānī dar ḥifẓ va
taʿlīmash jahd-i balīgh farmūdand thamariy-i ān shud kih shinīdih-īd.
Qad bakat al-samāwātu wa-al-arḍ thumma al-hawāʾu wa-al-miyāhu
wa-al-aḥjāru wa-al-ashjāru wa-al-athmāru wa-al-awrāq wa-lākinna
al-mushrikīna hum fī faraḥin ʿaẓīm.

Any soul with less than an atom of awareness knows with certainty that,
had He so willed, that man would never have existed. For twenty years in
this Cause and twenty years before this Cause, He exerted the utmost
effort, with perfect compassion and kindness, to protect and educate
him. The result was what you have heard. The heavens and the earth have
wept, and likewise the air, the waters, the stones, the trees, the
fruits, and the leaves. Yet those who associate partners with God are
greatly rejoicing.

\subsection{562}\label{section-1057}

\begin{otherlanguage}{persian}

و ديگر كجا شما در اين ارض بوده كه ظلم فهميده و نسبت داده و اللّه الّذى لا
اله الّا هو كه اگر احدى بآن نفس معرض تكلّم نموده‌تا چه رسد بظلم ظالم را
مظلوم دانسته و مشرك را موحّد شمرده خود تو اقرار بر اخلاق حقّ نموده‌ و مع‌ذلك
ظلم نسبت ميدهى لا تشعر ما تقول در كلّ اوقات با نسآء متعدّده در محلّ راحت
بعيش و عشرت مشغول بوده بلى در اين ارض آن سيّد مفسد كه انشآء اللّه خير
نه‌بيند و بمقصودش نرسد باو تلقين نموده ‌آنچه تلقين نموده‌ و جميع اين امور
از تدبيرات او است اگر چه خود اصل شجره اعراض و اللّه الّذى لا اله الّا هو كه
بمكر و خدعه ظاهر شده كه در ابداع مثل و شبه نداشته و ندارد درايت شما
معلوم شد كه خوب مطلب را مى فهميد يحرّكك هوآء الشّرك و النّفاق كيف یشآء

\end{otherlanguage}

Va dīgar kojā shomā dar īn arż būdeh keh ẓolm fahmīdeh va nesbat dādeh?
Va Allāh alladhī lā ilāha illā huwa keh agar aḥadī be-ān nafs-e moʿreż
takallom namūdeh, tā cheh resad be-ẓolm! Ẓālem rā maẓlūm dānesteh va
moshrek rā movaḥḥed shomordeh. Khod-e to eqrār bar akhlāq-e Ḥaqq namūdeh
va maʿa dhālek ẓolm nesbat mīdehī. Lā tashʿuru mā taqūl. Dar koll-e
owqāt bā nesāʾ-e motaʿaddedeh dar maḥall-e rāḥat be-ʿeysh va ʿeshrat
mashghūl būdeh. Balī dar īn arż ān seyyed-e mofsed keh in shāʾ Allāh
kheyr na-bīnad va be-maqṣūdash na-resad be-ū talqīn namūdeh āncheh
talqīn namūdeh va jamīʿ-e īn omūr az tadbīrāt-e ū ast, agar cheh khod
aṣl-e shajareh-ye eʿrāż. Va Allāh alladhī lā ilāha illā huwa keh be-makr
va khodʿeh ẓāher shodeh keh dar ebdāʿ methl va shebh nadāshteh va
nadārad. Derāyat-e shomā maʿlūm shod keh khūb maṭlab rā mīfahmīd.
Yuḥarrikuka hawāʾ al-shirk wa-l-nifāq kayfa yashāʾ.

And where, moreover, were you in this land that you perceived oppression
and made this accusation? By God, besides whom there is no god, no one
has so much as spoken to that soul who turned away, let alone wronged
him! You have regarded the oppressor as the oppressed and counted the
one who associates partners with God as a believer in His oneness. You
yourself have acknowledged the character of the True One, yet you accuse
Him of oppression. You do not comprehend what you say. At all times he
was occupied with pleasure and enjoyment, with several women, in a place
of comfort. Yes, in this land that troublemaking sayyid---may God grant
that he see no good and fail to attain his purpose---instilled in him
whatever he instilled, and all these matters are of his devising,
although he himself is the root of the tree of rejection. By God,
besides whom there is no god, he has displayed such deceit and trickery
as have never had, and do not have, any equal or likeness in creation.
Your discernment has become apparent: how well you understand the
matter! The desire to associate partners with God and to practice
hypocrisy moves you however it wishes.

\subsection{563}\label{section-1058}

\begin{otherlanguage}{persian}

قوله عظم افترائه خلاصه از شما بسيار تعجّب است سابقاً هر كه از دنيا و رياست
او گذشته ممدوح بود حال اسم او را شما خوف گذاشته‌ايد و كسيرا كه حضرت اعلى
در توقيعات باو من اللّه و الى اللّه اوّل عنوان توقيع مينوشته شما بنويسيد كه
در چادر نسوان مخفى شده عجب عجب طالب حقّ بوده‌ايد و خوب زيرآب از انبیاء
سابق كشيديد عيسى كه در مغاره‌ها مأوى گزيده بود و پيغمبر كه در جوال مخفى
گشته و ذكريا كه در درخت مقرّ جسته و كلّ اولیاء از خوف و يا جهت ديگر نفس
نميكشيدند از شهر بشهر و ديار بديار فرار ميكردند

\end{otherlanguage}

Qawluhu, ʿaẓuma iftirāʾuhu: Kholāṣeh az shomā besyār taʿajj ob ast.
Sābeqan har keh az donyā va reyāsat-e ū gozasht mamdūḥ būd. Ḥāl esm-e ū
rā shomā khowf goẕāshteh-īd va kasī rā keh Ḥażrat-e Aʿlā dar towqīʿāt
be-ū min Allāh wa-ilā Allāh avval-e ʿonvān-e towqīʿ mīneveshteh shomā
benevīsīd keh dar chādor-e nesvān makhfī shodeh. ʿAjab, ʿajab! Ṭāleb-e
Ḥaqq būdeh-īd va khūb zīrāb az anbiyā-ye sābeq keshīdīd. ʿĪsā keh dar
maghāreh-hā maʾvā gozīdeh būd va Peyghambar keh dar javāl makhfī gashteh
va Zakariyyā keh dar derakht maqarr josteh va koll-e owliyā az khowf va
yā jehat-e dīgar nafas nemīkeshīdand, az shahr be-shahr va deyār
be-deyār farār mīkardand.

His words---great is his slander: ``In short, one is greatly astonished
at you. Formerly, whoever renounced the world and its leadership was
praised. Now you have called this fear, and of someone to whom the Most
Exalted One, in His epistles, would write `From God and to God' at the
beginning of the address, you write that he hid himself in women's
veils. Astonishing, astonishing! You have been seekers of the True One,
and how thoroughly you have undermined the Prophets of the past! Jesus
took refuge in caves, the Prophet hid in a sack, and Zechariah found
shelter in a tree. All the holy ones, from fear or for some other
reason, dared not breathe and fled from city to city and from land to
land.''

\subsection{564}\label{section-1059}

\begin{otherlanguage}{persian}

انتهى كذبت باللّه الّذى خلق كلّ شیء تو خبر از آنچه از نزد نقطه اولى صادر
شده نداشته و ندارى چند جزوه از سماء مشيّت نقطه بيان نازل در بعضى من اللّه
الى اللّه و در بعضى من اللّه الى من يظهره اللّه مذكور و مقصود كلّ اين ظهور
امنع اقدس بوده كه تو و امثال تو از او معرض و بكمال انكار منكر يكى مخصوص
ميرزا اسد اللّه نازل و نقطه بيان او را ديّان فرموده و بكلّ اسماء حسنى
ناميده و مرشدت او را سبّ و لعن نموده‌و فتوى بر قتل او داده و آن بيچاره را
شهيد نمودند و هم‌چنين يكى مخصوص جناب آقا سيّد ابراهيم از سماء مشيّت مقتدر
عظيم نازل و او را مرشدت ابو الدّواهى ناميده بعضى از امور وارده بر اين دو
نفر از قبل ذكر شده فاقرئه لتطّلع بحسنات مرشدك و بما ورد عليهما

\end{otherlanguage}

Intahā. Kadhabta bi-Allāh alladhī khalaqa kulla shayʾ. To khabar az
āncheh az nazd-e Noqṭeh-ye Ūlā ṣāder shodeh nadāshteh va nadārī. Chand
jozveh az samā-ye mashiyyat-e Noqṭeh-ye Bayān nāzel, dar baʿżī min Allāh
ilā Allāh va dar baʿżī min Allāh ilā man yuẓhiruhu Allāh maẕkūr, va
maqṣūd-e koll īn ẓohūr-e amnaʿ-e aqdas būdeh keh to va amthāl-e to az ū
moʿreż va be-kamāl-e enkār monker. Yekī makhṣūṣ-e Mīrzā Asad Allāh nāzel
va Noqṭeh-ye Bayān ū rā Dayyān farmūdeh va be-koll-e asmā-ye ḥosnā
nāmīdeh va morshedat ū rā sabb va laʿn namūdeh va fatvā bar qatl-e ū
dādeh va ān bīchāreh rā shahīd namūdand. Va hamchenīn yekī makhṣūṣ-e
Janāb-e Āqā Seyyed Ebrāhīm az samā-ye mashiyyat-e moqtader-e ʿaẓīm nāzel
va ū rā morshedat Abū al-Davāhī nāmīdeh. Baʿżī az omūr-e vāredeh bar īn
do nafar az qabl ẕekr shodeh. Fa-qraʾhu li-taṭṭaliʿa bi-ḥasanāti
murshidika wa-bi-mā warada ʿalayhimā.

End of quotation. You have lied, by God who created all things! You have
never known, and do not know, what issued from the Primal Point. Several
booklets descended from the heaven of the will of the Point of the
Bayān. In some, `From God to God' is mentioned, and in others, `From God
to Him whom God will make manifest.' The intended object of them all was
this most unassailable, most holy Revelation, from which you and your
like have turned away and which you deny with utter rejection. One was
revealed specifically for Mīrzā Asad Allāh, whom the Point of the Bayān
called Dayyān and named with all the most beautiful names. Your guide
reviled and cursed him and issued a decree for his death, and they
martyred that hapless man. Likewise, one was revealed specifically for
His Honor Āqā Sayyid Ibrāhīm from the heaven of the mighty, all-powerful
Will, and your guide called him the Father of Calamities. Some of the
things that befell these two have been mentioned before. Read it, then,
so that you may learn of your guide's good deeds and of what befell
them.

\subsection{565}\label{section-1060}

\begin{otherlanguage}{persian}

اين كلمات مذكور شد تا من حيث لايشعر سخن نگوئى و بهر ريحى متحرّك نشوى و از
حقّيكه اگر بخواهد جميع اسامى كه تو او را احصا نموده‌ بقبضه از تراب راجع
نمايد قادر است اعراض نكنى چه كه يك مقام از مقامات توصيف حقّ آن است كه چون
در كلّ شیء آثار تجلّيات خود مشاهده ميفرمايد لذا بادنى آيه ظاهره در ارض كلّ
اسماء و صفات راجع شده و خواهد شد چه كه مقصود از اسماء حسنى كه باو خطاب
ميشود راجع است بآن تجلّى ربانيّه كه من غير جهة در او تجلّى فرموده و تا اين
آيه تجلّى در كلّ شیء باقى بر كلّ شیء اطلاق كلّ اسماء صادق و بعد از اعراض هر
يك محروم و بمقرّ نار راجع و اگر مشاهده ميشد كه قابل عرفانى هر آينه بحر
معانى و عرفان رشحى بتو مبذول ميفرمود و لكن بهتر همانكه محروم باشى و
بهمين كلمات محدوده موتفكه محتجبه تكلّم نمائى لا تخرج من اهل السّجّين الّا
كلمات السّجّينيّة

\end{otherlanguage}

Īn kalemāt maẕkūr shod tā min ḥaythu lā yashʿur sokhan nagūʾī va be-har
rīḥī motaḥarrek nashavī va az Ḥaqqī keh agar bekhāhad jamīʿ-e asāmī keh
to ū rā eḥṣā namūdeh be-qabżeh az torāb rājeʿ namāyad qāder ast eʿrāż
nakonī. Cheh keh yek maqām az maqāmāt-e towṣīf-e Ḥaqq ān ast keh chūn
dar koll-e shayʾ āthār-e tajalliyāt-e khod moshāhedeh mīfarmāyad, leẕā
be-adnā āyeh-ye ẓāhereh dar arż koll-e asmā va ṣefāt rājeʿ shodeh va
khwāhad shod. Cheh keh maqṣūd az asmā-ye ḥosnā keh be-ū kheṭāb mīshavad
rājeʿ ast be-ān tajallī-ye rabbāniyyeh keh min ghayri jihah dar ū
tajallī farmūdeh va tā īn āyeh-ye tajallī dar koll-e shayʾ bāqī bar
koll-e shayʾ eṭlāq-e koll-e asmā ṣādeq va baʿd az eʿrāż har yek maḥrūm
va be-maqarr-e nār rājeʿ. Va agar moshāhedeh mīshod keh qābel-e ʿerfānī
har āyeneh baḥr-e maʿānī va ʿerfān rashḥī be-to mabẕūl mīfarmūd va lāken
behtar hamān keh maḥrūm bāshī va be-hamīn kalemāt-e maḥdūdeh-ye
moʾtafekeh-ye moḥtajebeh takallom namāʾī. Lā takhruju min ahl al-sijjīn
illā kalimāt al-sijjīniyyah.

These words have been mentioned so that you may not speak without
awareness, be moved by every wind, or turn away from the True One, who,
if He wishes, has the power to reduce all the names you have enumerated
to a handful of dust. For one of the stations of the True One's
description is this: since He beholds in all things the traces of His
own revelations, all names and attributes have referred, and will refer,
to the lowliest sign manifest on earth. For the most beautiful names by
which it is addressed refer to that divine revelation which He has
manifested within it without direction. As long as this sign of
revelation remains in all things, the application of all names to all
things is valid; after turning away, each is deprived and returns to the
abode of fire. Had you been seen to be capable of recognition, the Ocean
of meanings and recognition would surely have bestowed a sprinkling upon
you. But it is better that you remain deprived and speak these same
limited, perverted, veiled words. Nothing issues from the people of the
nethermost prison but words belonging to that prison.

\subsection{566}\label{section-1061}

\begin{otherlanguage}{persian}

و اينكه نوشته سابق هر كه از دنيا و رياست او گذشته ممدوح بود حال شما اسم
او را خوف گذاشته‌ايد و نوشته كه فلان در چادر نسوان مخفى شده آنچه من عند
اللّه نازل و المنزل هو الحقّ و هو صادق فيما قال و يقول خوب منقطعى يافتى
بلى منقطع از حقّ نه منقطع از ما سوى فو‌اللّه شرم دارم كه ذكر اعمالش نمايم
چه كه تو مطّلع نيستى نفوسيكه با او بوده‌اند حال موجود نفسيكه در شهوت
نفسانى بمقامى رسيد كه از حرم يزدانى نگذشت و تصرّف در حرم نقطه نمود فو‌اللّه
ستر حجاب اكبر را هتك كرد اين بی‌‌ادب بی‌‌شرم او را منقطع مينويسى

\end{otherlanguage}

Va īnkeh neveshteh sābeq har keh az donyā va reyāsat-e ū goẕashteh
mamdūḥ būd, ḥāl shomā esm-e ū rā khowf goẕāshteh-īd, va neveshteh keh
folān dar chādor-e nesvān makhfī shodeh. Āncheh min ʿind Allāh nāzel,
wa-l-munzil huwa al-Ḥaqq wa-huwa ṣādiqun fī-mā qāla wa-yaqūl. Khūb
monqaṭeʿī yāftī! Balī monqaṭeʿ az Ḥaqq, nah monqaṭeʿ az mā sevā.
Fa-wa-Allāh sharm dāram keh ẕekr-e aʿmālash namāyam cheh keh to moṭṭaleʿ
nīstī. Nofūsī keh bā ū būdeh-and ḥāl mowjūd. Nafsī keh dar shahvat-e
nafsānī be-maqāmī resīd keh az ḥaram-e Yazdānī nagoẕasht va taṣarrof dar
ḥaram-e Noqṭeh namūd. Fa-wa-Allāh setr-e ḥejāb-e akbar rā hatk kard. Īn
bī-adab-e bī-sharm, ū rā monqaṭeʿ mīnevīsī!

And as for his writing, `Formerly, whoever renounced the world and its
leadership was praised; now you have called this fear,' and his writing
that a certain person hid himself in women's veils: what was revealed
came from God, and the One who reveals it is the True One, and He is
truthful in what He has said and says. What a detached soul you have
found! Yes, detached from the True One, not detached from all else. By
God, I am ashamed to mention his deeds, for you are unaware. The people
who were with him are still alive. This is someone whose carnal desire
reached such a point that he did not spare the divine household, but
took possession of the Point's wife. By God, he tore through the
covering of the Most Great Veil. This shameless, insolent man---you
describe him as detached!

\subsection{567}\label{section-1062}

\begin{otherlanguage}{persian}

در سنين قبل هفت نفر شهداى ناريّه معيّن نموده ‌يكى از آن نفوس بعد از اطّلاع
تائب شده و رجع الى اللّه و تاب اليه و البسه اللّه ثوب غفرانه و ردآء فضله و
ارتقاه الى جنّة الأبهی فوق سدرة المنتهى الّذى سمّى بموسى فى ملكوت الأسماء
انّه لهو الّذى ارتقى الى افق الأعلی و حينئذ يكون بالمنظر الأبهی و يطوفن فى
حوله اهل مداين البقآء و آن نفس مبارك ميرزا موسى قمى عليه بهآء اللّه بود
كه شكست صنم وهم و هوى را و در فلك ابهى باذن اللّه خالق الأرض و السّمآء
مستقرّ شد و بعد الى اللّه رجوع نمود و نفس ديگر ابدا قبول اين افك ننمود كانّ
روح القدس اخبره و حفظه عن مفتريات هؤلآء المفترين

\end{otherlanguage}

Dar senīn-e qabl haft nafar shohadā-ye nāriyyeh moʿayyan namūdeh. Yekī
az ān nofūs baʿd az eṭṭelāʿ tāʾeb shodeh va rajaʿa ilā Allāh wa-tāba
ilayhi wa-albasahu Allāh thawba ghufrānihi wa-ridāʾa faḍlihi wa-artaqāhu
ilā jannat al-Abhā fawqa Sidrat al-Muntahā alladhī summiya bi-Mūsā fī
malakūt al-asmāʾ. Innahu la-huwa alladhī irtaqā ilā ufuq al-Aʿlā
wa-ḥīnaʾidhin yakūnu bi-l-manẓar al-Abhā wa-yaṭūfanna fī ḥawlihi ahlu
madāʾin al-baqāʾ. Va ān nafs-e mobārak Mīrzā Mūsā-ye Qomī ʿalayhi bahāʾ
Allāh būd keh shekast ṣanam-e vahm va havā rā va dar falak-e Abhā
be-eẕn-e Allāh khāliq al-arḍ wa-l-samāʾ mostaqarr shod va baʿd ilā Allāh
rojūʿ namūd va nafs-e dīgar abadan qabūl-e īn efk nanamūd. Ka-anna Rūḥ
al-Qudus akhbarahu wa-ḥafiẓahu ʿan muftarayāti hāʾulāʾ al-muftarīn.

In former years he designated seven people as martyrs of fire. One of
these souls, upon learning of it, repented and returned to God and
turned to Him in repentance. God clothed him in the garment of His
forgiveness and the robe of His grace, and raised him to the
All-Glorious Paradise above the Lote Tree of the Utmost Boundary---he
who was named Mūsā in the kingdom of names. He is indeed the one who
ascended to the Most Exalted Horizon, and now he dwells in the
All-Glorious Vista, while the inhabitants of the cities of eternity
circle around him. That blessed soul was Mīrzā Mūsā of Qum, upon him be
the glory of God, who shattered the idol of delusion and desire and, by
the permission of God, the Creator of earth and heaven, became
established in the All-Glorious sphere, and afterward returned to God.
Another soul never accepted this falsehood, as though the Holy Spirit
had informed him and protected him from the fabrications of these
fabricators.

\subsection{568}\label{section-1063}

\begin{otherlanguage}{persian}

بارى اگر تفصيل ذكر شود بطول ميانجامد مختصر آنكه سه نفر از آن نفوس باقى و
در آن جزوه اين نفوس را شهيد نموده‌و تصريحا نوشته كه هر كدام دختر بكرى از
براى او بفرستند مع درهم و دنانير و آن جزوه الآن در دست موجود يقين اين
امور را از انقطاع شمرده فو‌اللّه اين امور مذكوره نسبت باعمال ديگرش حسنه
بوده و اين عبد شرم مينمايم از ذكر آن تا حال ستر كبريائى عيوبش را پوشيده
بود چون خرق نمود ستر را لذا حقّ كشف فرمود العجب كلّ العجب بل العجب بنفسه
يقول العجب ثمّ العجب كه فرق ميانه خزف و لؤلؤ ننموده‌ و از اين مقرّ اقدس
گذشته مفتريات بنبيّين و مرسلين هم بسته بجهت اثبات امر موهومى زبان شتم
بسلاطين وجود گشوده

\end{otherlanguage}

Bārī agar tafṣīl ẕekr shavad be-ṭūl mīanjāmad. Mokhtaṣar ānkeh seh nafar
az ān nofūs bāqī va dar ān jozveh īn nofūs rā shahīd namūdeh va taṣrīḥan
neveshteh keh har kodām dokhtar-e bekrī az barā-ye ū beferestand maʿa
derham va danānīr va ān jozveh al-ān dar dast mowjūd. Yaqīn īn omūr rā
az enqeṭāʿ shomordeh. Fa-wa-Allāh īn omūr-e maẕkūreh nesbat be-aʿmāl-e
dīgarash ḥasaneh būdeh va īn ʿabd sharm mīnamāyam az ẕekr-e ān. Tā ḥāl
setr-e kebriyāʾī ʿoyūbash rā pūshīdeh būd. Chūn kharq namūd setr rā leẕā
Ḥaqq kashf farmūd. Al-ʿajab kull al-ʿajab bal al-ʿajab bi-nafsihi yaqūlu
al-ʿajab thumma al-ʿajab keh farq-e miyāneh-ye khazaf va loʾloʾ
nanamūdeh va az īn maqarr-e aqdas goẕashteh moftariyāt be-nabiyyīn va
morsalīn ham basteh. Be-jehat-e ethbāt-e amr-e mowhūmī zabān-e shatm
be-salāṭīn-e vojūd goshūdeh.

In any case, to give the details would take too long. In short, three of
those souls remained, and in that booklet he made these souls martyrs
and explicitly wrote that each should send him a virgin girl, together
with dirhams and dinars. That booklet is now at hand. Surely you have
counted these things as detachment. By God, these deeds, compared with
his other acts, were good deeds, and I, this servant, am ashamed to
mention those. Until now the covering of divine majesty had concealed
his faults. Since he tore the covering, the True One therefore disclosed
them. Astonishing, utterly astonishing! Indeed, astonishment itself
says, `Astonishing, and again astonishing!' For he has not distinguished
earthenware from pearls, and, going beyond this most holy Seat, has also
fabricated slanders against the Prophets and Messengers. To establish an
imaginary cause, he has opened his mouth to revile the sovereigns of
existence.

\subsection{569}\label{section-1064}

\begin{otherlanguage}{persian}

اى بی‌اطّلاع بشنو امّا عيسى بن مريم از اوّل بعثت اكثر ایّام در هيكل با رؤساى
يهود مكالمه ميفرمود و من غير تقيّه تبليغ امر مينمود چنانچه در كتب ثبت است
و تو نديده نصّ انجيل است وقتيكه آن روح مقدّس رحمانى را در مجلس يكى از
علماى يهود كه موسوم بحنّاس بود بردند و او از حضرت سؤال نمود كه قول تو
چيست و چه ميگوئى حضرت فرمودند كه چرا از من سؤال مينمائى از ملأ يهود سؤال
كن چه كه من در خلوت كلمه نگفته‌ام بلكه آنچه گفته‌ام در هيكل و در مجمع يهود
گفته‌ام بعد يكى از عسكريه كه در آن مجلس حاضر بود طپانچه بر صورت مبارك
حضرت زده كه چرا بكاهن بزرگ بگستاخى تكلّم مينمائى

\end{otherlanguage}

Ey bī-eṭṭelāʿ, beshnow. Ammā ʿĪsā ibn Maryam az avval-e beʿthat akthar-e
ayyām dar heykal bā roʾasā-ye Yahūd mokālemeh mīfarmūd va min ghayri
taqiyyah tablīgh-e amr mīnamūd chenāncheh dar kotob thabt ast va to
nadīdeh. Naṣṣ-e Enjīl ast vaqtī keh ān rūḥ-e moqaddas-e raḥmānī rā dar
majles-e yekī az ʿolamā-ye Yahūd keh mowsūm be-Ḥannās būd bordand va ū
az Ḥażrat soʾāl namūd keh qowl-e to chīst va cheh mīgūʾī, Ḥażrat
farmūdand keh cherā az man soʾāl mīnamāʾī? Az malaʾ-e Yahūd soʾāl kon
cheh keh man dar khalvat kalemeh nagoft eh-am balkeh āncheh gofteh-am
dar heykal va dar majmaʿ-e Yahūd gofteh-am. Baʿd yekī az ʿaskariyyeh keh
dar ān majles ḥāżer būd ṭapāncheh bar ṣūrat-e mobārak-e Ḥażrat zadeh keh
cherā be-kāhen-e bozorg be-gostākhī takallom mīnamāʾī?

Listen, you who are uninformed. As for Jesus son of Mary, from the
beginning of His mission He conversed with the Jewish leaders in the
Temple on most days and proclaimed the Cause without concealment, as is
recorded in the books, which you have not seen. The Gospel explicitly
states that when they brought that holy, merciful Spirit into the
assembly of one of the Jewish scholars, who was named Annas, he asked
Him, `What is Your teaching, and what do You say?' He replied, `Why do
you question Me? Question the Jewish people, for I have not spoken a
word in secret; rather, whatever I have said, I have said in the Temple
and in the gathering of the Jews.' Then one of the soldiers present in
that assembly struck His blessed face, saying, `Why do You speak
insolently to the high priest?'

\subsection{570}\label{section-1065}

\begin{otherlanguage}{persian}

بعد آن عالم حكم نمود كه حضرت را نزد قيافا كه قاضى آن سنه بود بردند يعنى
در آن سنه احكام شرعيّه ملّت يهود با او بود با اينكه نصّ كتاب است كه هر روز
در هيكل تبليغ امر اللّه جهرةً ميفرمود حال نوشته كه عيسى در مغاره‌ها بوده و
از خوف نفس نميكشيدند بسيار تعدّى نموده‌ و از شأن خود تجاوز جسته

\end{otherlanguage}

Baʿd ān ʿālem ḥokm namūd keh Ḥażrat rā nazd-e Qeyāfā keh qāżī-ye ān
saneh būd bordand, yaʿnī dar ān saneh aḥkām-e sharʿiyyeh-ye mellat-e
Yahūd bā ū būd. Bā īnkeh naṣṣ-e ketāb ast keh har rūz dar heykal
tablīgh-e amr Allāh jahratan mīfarmūd ḥāl neveshteh keh ʿĪsā dar
maghāreh-hā būdeh va az khowf nafas nemīkeshīdand. Besyār taʿaddī
namūdeh va az shaʾn-e khod tajāvoz josteh.

Then that scholar gave orders, and they took Him to Caiaphas, who was
the judge that year---that is, in that year the administration of the
Jewish community's religious law rested with him. Although the Book
explicitly states that every day He publicly proclaimed the Cause of God
in the Temple, he has now written that Jesus lived in caves and dared
not breathe from fear. He has committed a grave transgression and
overstepped his station.

\subsection{571}\label{section-1066}

\begin{otherlanguage}{persian}

در اينمقام يك بيان از روح ذكر ميشود كه شايد از كوثر بيان جمال رحمن
بحديقه سبحان وارد شوى فلمّا رأى الصّليب حمله بنفسه ثمّ اعتنقه قائلاً هلمّ يا
صليبی‌ العزيز الّذى منذ ثلثة و ثلثين سنة انا منتظرك و مفتّشاً عليك و اريد ان
اموت مسمّراً فيك حبّاً باغنامى اگر چه نفحات اين كلمه را جز روحانيّين نيابند و
لذّت اين بيان را جز اهل منظر رحمن درك ننمايند امثال آن نفوس بی‌‌نصيب بوده و
خواهند بود چنانچه ادراك اينمقام ننموده ‌اگر مينمودى باينگونه جسارت و بی‌
حرمتى بآن شموس افق تجريد تحرير نمينمودى بشنو كه چه فرموده در حينيكه جميع
يهود بر قتل آنحضرت مجتمع شده بودند و چون صليب حاضر نمودند و چشم مبارك بر
صليب افتاد فرمود بيا بيا كه سى و سه سنه است منتظر توام و مشتاق تو كجا
است اينمقام و مقاميكه تو درك نموده‌و نوشته

\end{otherlanguage}

Dar īn maqām yek bayān az Rūḥ ẕekr mīshavad keh shāyad az kawthar-e
bayān-e Jamāl-e Raḥmān be-ḥadīqeh-ye Sobḥān vāred shavī. Fa-lammā raʾā
al-ṣalīb ḥamalahu bi-nafsihi thumma iʿtanaqahu qāʾilan halumma yā ṣalībī
al-ʿazīz alladhī mundhu thalāthatin wa-thalāthīna sanah anā muntaẓiruka
wa-mufattishan ʿalayka wa-urīdu an amūta musammanran fīka ḥubban
bi-aghnāmī. Agar cheh nafaḥāt-e īn kalemeh rā joz rūḥāniyyīn nayāband va
leẕẕat-e īn bayān rā joz ahl-e manẓar-e Raḥmān dark nanamāyand. Amthāl-e
ān nofūs bī-naṣīb būdeh va khwāhand būd chenāncheh edrāk-e īn maqām
nanamūdeh. Agar mīnamūdī be-īngūneh jasārat va bī-ḥormatī be-ān shomūs-e
ofoq-e tajrīd taḥrīr nemīnamūdī. Beshnow keh cheh farmūdeh dar ḥīnī keh
jamīʿ-e Yahūd bar qatl-e ān Ḥażrat mojtameʿ shodeh būdand va chūn ṣalīb
ḥāżer namūdand va chashm-e mobārak bar ṣalīb oftād farmūd biyā biyā keh
sī va seh saneh ast montaẓer-e to-am va moshtāq-e to. Kojā ast īn maqām
va maqāmī keh to dark namūdeh va neveshteh?

Here a saying of the Spirit is mentioned, so that perhaps, through the
heavenly stream of the utterance of the Beauty of the All-Merciful, you
may enter the garden of the All-Glorious. When He saw the cross, He
carried it Himself, then embraced it, saying, `Come, My beloved cross,
for thirty-three years I have awaited you and searched for you, and I
wish to die nailed to you out of love for My sheep.' Yet none but the
spiritually minded perceive the fragrances of these words, and none but
the people of the Vista of the All-Merciful comprehend the sweetness of
this utterance. The likes of those souls have been, and will remain,
without a share, just as you have not understood this station. Had you
understood it, you would not have written with such audacity and
disrespect about those Suns of the horizon of detachment. Listen to what
He said when all the Jews had gathered to put Him to death. When they
brought the cross and His blessed eyes fell upon it, He said, `Come,
come, for thirty-three years I have awaited you and longed for you.' How
far apart are this station and the station you have understood and
written about!

\subsection{572}\label{section-1067}

\begin{otherlanguage}{persian}

و در اينمقام اختلاف شده آنچه در انجيل منصوص آن است كه آن ساذج قدم را بر
صليب آويختند و روح را تسليم فرمود و بعد از سى و شش ساعت زنده شد و بآسمان
ارتقاء جست و آنچه در اخبار اهل فرقان آن است كه قبل از آن بآسمان ارتقاء
جست و لكن حقيقت امر علمه عند ربّی‌ فى كتاب الّذى ما ترك فيه نبأ الأوّلین‌ و
الآخرین چه كه اكثرى از ناس معانى كلمات الهى را درك ننموده‌سخن گفته‌اند لذا
از اصل مقصود محجوب مانده‌اند

\end{otherlanguage}

Va dar īn maqām ekhtelāf shodeh. Āncheh dar Enjīl manṣūṣ ān ast keh ān
sāẕaj-e qedam rā bar ṣalīb āvīkhtand va rūḥ rā taslīm farmūd va baʿd az
sī va shesh sāʿat zendeh shod va be-āsmān erteqā jost. Va āncheh dar
akhbār-e ahl-e Forqān ān ast keh qabl az ān be-āsmān erteqā jost. Va
lāken ḥaqīqat-e amr ʿilmuhu ʿinda rabbī fī kitāb alladhī mā taraka fīhi
nabaʾ al-awwalīn wa-l-ākhirīn. Cheh keh aktharī az nās maʿānī-ye
kalemāt-e elāhī rā dark nanamūdeh sokhan gofteh-and leẕā az aṣl-e maqṣūd
maḥjūb māndeh-and.

There has been disagreement on this point. What is explicitly stated in
the Gospel is that they hung that Essence of Eternity upon the cross,
and He yielded up His spirit; after thirty-six hours He came to life and
ascended to heaven. What is related in the traditions of the people of
the Qurʾān is that He ascended to heaven before that. But as for the
reality of the matter, its knowledge is with my Lord in a Book in which
He has omitted none of the tidings of the first or the last. For most
people have spoken without understanding the meanings of the divine
words, and have therefore remained veiled from the underlying intent.

\subsection{573}\label{section-1068}

\begin{otherlanguage}{persian}

اين مجملى از حكايت روح كه ذكر شد و حينئذ هذا الرّوح يخاطبكم و يقول هلمّوا
هلمّوا يا ملأ المنكرين باسيافكم و رماحكم و سهامكم و انّى اشتاق كما هو كان
مشتاقاً فو‌الّذى نفس حسين بيده لاكون اشدّ اشتياقاً و اعظم انتظاراً و لكن انتم
لا تشعرون چه كه از كوثر معانى نچشيده‌ايد و از سلسبيل حبّ نياشاميده‌ايد فآهء
آه عمّا افتريت على انبیاء اللّه و ظلمت على مظاهره و مطالع وحيه و مشارق
الهامه و مكامن علمه و مخازن حكمته و اين ارتكاب را ننموده‌ مگر آنكه خواسته
عمل مرشدت را تصحيح كنى كذلك سوّلت لك نفسك امراً و جعلك من المعتدين ا فمن
كان مستوراً تحت القناع كمن فدى روحه فى سبيل اللّه

\end{otherlanguage}

Īn mojmalī az ḥekāyat-e Rūḥ ke dhekr shod va ḥīnaʾidh hādhā al-Rūḥ
yukhāṭibukum wa-yaqūlu halummū halummū yā malaʾ al-munkirīn bi-asyāfikum
wa-rimāḥikum wa-sihāmikum wa-innī ashtāqu kamā huwa kāna mushtāqan
fa-wa-alladhī nafsu Ḥusayn bi-yadihi la-akūnu ashadda ishtiyāqan
wa-aʿẓama intiẓāran wa-lākin antum lā tashʿurūn che ke az kawthar-e
maʿānī nacheshīde-īd va az salsabīl-e ḥobb nayāshāmīde-īd fa-āh āh ʿammā
iftarayta ʿalā anbiyāʾ Allāh wa-ẓalamta ʿalā maẓāhirihi wa-maṭāliʿ
waḥyihi wa-mashāriq ilhāmihi wa-makāmin ʿilmihi wa-makhāzin ḥikmatihi va
īn ertekāb rā nanemūde magar ānke khāste ʿamal-e morshedat rā taṣḥīḥ
konī kadhālika sawwalat laka nafsuka amran wa-jaʿalaka min al-muʿtadīn
a-fa-man kāna mastūran taḥt al-qināʿ ka-man fadā rūḥahu fī sabīl Allāh.

This is a brief account of the story of the Spirit that has been
related. Now this Spirit addresses you and says, ``Come, come, O company
of deniers, with your swords, your spears, and your arrows! I long for
them just as He longed for them. By Him in whose hand is the soul of
Ḥusayn, I am even more eager and more expectant, but you do not
perceive.'' For you have not tasted the heavenly abundance of meanings
or drunk from the sweet spring of love. Alas, alas for what you have
fabricated against the Prophets of God, and for the wrong you have done
to His Manifestations, the dawning-places of His revelation, the
daysprings of His inspiration, the repositories of His knowledge, and
the treasuries of His wisdom! You have committed this only because you
wished to justify your guide's conduct. Thus your own self has made a
deed seem fair to you and made you one of the transgressors. Is one who
remained concealed beneath a veil like one who sacrificed His spirit in
the path of God?

\subsection{574}\label{section-1069}

\begin{otherlanguage}{persian}

قرآن هم نخوانده كه ميفرمايد قوله تعالى فضل اللّه المجاهدين على القاعدين
اجراً عظيماً ا فمن كان قاعدا فى البيت كمن كان مجاهدا فى سبيل اللّه انفس
منقطعه را از انفس مشركه محدوده بسى فرق است قياس مكن و ذيل تقديس انبیاء
را بغبار ظنون و اوهام ميالاى اينكه نوشته كه از خوف نفس نميكشيدند همچه
نفوسى مثل تو خواهند بود فسبحانهم عن ذلك

\end{otherlanguage}

Qurʾān ham nakhānde ke mīfarmāyad qawluhu taʿālā faḍḍala Allāhu
al-mujāhidīna ʿalā al-qāʿidīna ajran ʿaẓīman a-fa-man kāna qāʿidan fī
al-bayt ka-man kāna mujāhidan fī sabīl Allāh anfos-e monqaṭeʿe rā az
anfos-e moshreke-ye maḥdūde basī farq ast qiyās makon va dheyl-e
taqdīs-e anbiyā rā be-ghobār-e ẓonūn va awhām mayālāy īnke neveshte ke
az khawf nafas nemīkeshīdand hamcho nofūsī methl-e to khāhand būd
fa-subḥānahum ʿan dhālik.

You have not even read the Qurʾān, where He says---exalted are His
words---``God has favored those who strive over those who sit idle with
a mighty reward.'' Is one who sits at home like one who strives in the
path of God? Vast is the difference between detached souls and those who
associate partners with God and are confined by limitations. Do not
compare them, and do not soil the hem of the Prophets' holiness with the
dust of conjectures and imaginings. As for your statement that they did
not dare breathe for fear, such souls would be like you. Exalted are
They above that!

\subsection{575}\label{section-1070}

\begin{otherlanguage}{persian}

اى مقلّد متوهّم چنين نفوس بچه كار خواهند آمد ما فى‌حقيقتت از اين كلمه معلوم
شد كه عقيده تو و معرضين باللّه چه بوده و چه خواهد بود جواهر وجود را مثل
خود دانسته كه شموس حقيقت را بانفس مشركه قياس كرده فو‌اللّه اضطربت من قولك
كلّ شیء آنقدر بدان كه آنچه در اعلى رتبه عرفان خود عارف شوى وهمى است از تو
كه راجع بنفس خود تو بوده و اصفياى الهى از آن مقدّس و مبرّا قسم بآفتاب عزّ
معانى كه در كلّ حين و حان بتمام اشتياق منتظر بوده‌اند كه در سبيل رحمن جان
انفاق نمايند و تو استهزاء على اللّه مينويسى كه پيغمبر در جوال رفته بود از
كلمه‌ات قدرت معلوم ديگر لفظى نبود كه مذكور نمائى كه مخصوص باين لفظ قبيح
ذكر كرده افّ لك و لحيائك فطرت مرشدت هم مثل خودت بوده و خواهد بود

\end{otherlanguage}

Ey moqalled-e motavahhem chenīn nofūs be-che kār khāhand āmad mā fī
ḥaqīqatat az īn kaleme maʿlūm shod ke ʿaqīde-ye to va moʿreżīn bi-llāh
che būde va che khāhad būd javāher-e vojūd rā methl-e khod dāneste ke
shomūs-e ḥaqīqat rā be-anfos-e moshreke qiyās karde fa-wa-Allāh
iḍṭarabat min qawlika kull shayʾ ānqadr bedān ke ānche dar aʿlā rotbe-ye
ʿerfān-e khod ʿāref shavī vahmī ast az to ke rājeʿ be-nafs-e khod-e to
būde va aṣfiyā-ye elāhī az ān moqaddas va mobarrā qasam be-āftāb-e
ʿezz-e maʿānī ke dar koll-e ḥīn va ān be-tamām-e eshtiyāq montaẓer
būde-and ke dar sabīl-e Raḥmān jān enfāq namāyand va to estehzāʾan ʿalā
Allāh mīnevīsī ke peyghambar dar javāl rafte būd az kaleme-at qadrat
maʿlūm dīgar lafẓī nabūd ke madhkūr namāʾī ke makhṣūṣ be-īn lafẓ-e qabīḥ
dhekr karde off laka wa-li-ḥayāʾika feṭrat-e morshedat ham methl-e
khodat būde va khāhad būd.

O deluded imitator! What use would such souls be? What lies within your
own reality has become evident from this statement: what your belief and
that of those who turn away from God has been and will be. You have
considered the jewels of existence to be like yourself, comparing the
Suns of Truth to souls who associate partners with God. By God,
everything has been shaken by your words! Know at least this: whatever
you may comprehend at the highest degree of your own understanding is an
imagining of yours that refers back to your own self, while God's chosen
ones are sanctified and free from it. By the Sun of the glory of
meanings, at every moment and instant They have eagerly awaited the
chance to offer up Their lives in the path of the All-Merciful. Yet you
write in mockery of God that the Prophet had gone into a sack. Your
words reveal your worth. Was there no other expression you could use,
that you chose to mention Him with this vile expression? Shame on you
and your sense of shame! Your guide's nature, too, has been and will
remain like your own.

\subsection{576}\label{section-1071}

\begin{otherlanguage}{persian}

قسم بهمان مظاهر عزّ تقديس كه اگر فى‌الجمله حبّ در تو بود اگر جميع اركانت را
باسياف برنده قطعه قطعه مينمودند چنين الفاظ ركيكه در شأن مظاهر احديّه از
لسان معتقله كليله‌ات جارى نميشد مرشدت هم از اين حرفها بسيار گفته و همين
سيّد محمّد باو استهزاء مينمود حمد محبوبی‌ را كه امثال آن وجود را از عرفان
مظاهر خود محروم نموده ‌و از حلاوت حبّ ممنوع داشته آيا نشنيدى كه خاتم
انبیاء روح ما سواه فداه در مدّت بيست سال جان مباركش در دست اعدا مبتلا و
در هيچ شب اميد صبح نبود و در هيچ صبح اميد شب نه فو‌اللّه در كلّ احيان منتظر
و آمل كه جان در ره دوست ايثار نمايد و برفيق اعلى شتابد

\end{otherlanguage}

Qasam be-hamān maẓāher-e ʿezz-e taqdīs ke agar fī al-jomle ḥobb dar to
būd agar jamīʿ-e arkānat rā be-asyāf-e borrande qeṭʿe qeṭʿe mīnemūdand
chenīn alfāẓ-e rakīke dar shaʾn-e maẓāher-e aḥadiyye az lesān-e
moʿtaqele-ye kalīle-at jārī nemīshod morshedat ham az īn ḥarfhā besyār
gofte va hamīn Seyyed Moḥammad be-ū estehzā mīnemūd ḥamd-e maḥbūbī rā ke
amthāl-e ān vojūd rā az ʿerfān-e maẓāher-e khod maḥrūm nemūde va az
ḥalāvat-e ḥobb mamnūʿ dāshte āyā nashenīdī ke Khātam-e Anbiyā rūḥ-e mā
sevāh fedāh dar moddat-e bīst sāl jān-e mobārakash dar dast-e aʿdā
mobtalā va dar hīch shab omīd-e ṣobḥ nabūd va dar hīch ṣobḥ omīd-e shab
na fa-wa-Allāh dar koll-e aḥyān montaẓer va āmel ke jān dar rah-e dūst
īthār namāyad va be-Rafīq-e Aʿlā shetābad.

By those very Manifestations of the glory of holiness, had there been
even a little love in you, such base words about the Manifestations of
oneness would not have issued from your halting, faltering tongue, even
if all your limbs had been cut to pieces with sharp swords. Your guide,
too, has said many such things, and this same Siyyid Muḥammad used to
mock Him. Praise be to the Beloved, who has deprived the likes of that
person of the recognition of His Manifestations and denied them the
sweetness of love! Have you not heard that for twenty years the blessed
life of the Seal of the Prophets---may the spirit of all else be
sacrificed for Him---was afflicted at the hands of His enemies, and that
on no night was there hope of morning, and on no morning hope of night?
By God, at all times He awaited and hoped to sacrifice His life in the
path of the Friend and hasten to the Most Exalted Companion.

\subsection{577}\label{section-1072}

\begin{otherlanguage}{persian}

اى زاغ ارض جرزه حكايت از طوطى مصر احديّه منما اى طنين ذباب از صرير قلم
اعلى ذكر مكن كه عرض خود ميبرى و زحمت ما ميارى آيا نشنيده كه حسين بن علىّ
روح من فى ملكوت الابداع فداهما بچه شوق و اشتياق جان در ره حقّ باخت و اگر
بگوئى چنانچه ملحدين گفته‌اند كه بين يدى اعدا مبتلى شد و چاره جز كشته شدن
نداشت هذا القول بغى منك و منهم على اللّه المهيمن القيّوم چه كه در حينى كه
از مدينه حركت ميفرمود بهمين نيّت و اراده حركت فرمود فو‌اللّه بشوق لقاء
محبوب حركت نمود و طلباً لوصاله قطع مراحل فرمود چنانچه در حين خروج آن جمال
احديّه از مدينه مخصوص وداع بروضه مقدّسه ساذج وجود من الغيب و الشّهود جدّ
مطهّر خود تشريف برده و باين كلمات ابدع احلى ناطق

\end{otherlanguage}

Ey zāgh-e arż-e joroze ḥekāyat az ṭūṭī-ye Meṣr-e aḥadiyye manamā ey
ṭanīn-e dhobāb az ṣarīr-e Qalam-e Aʿlā dhekr makon ke ʿerż-e khod mībarī
va zaḥmat-e mā mayārī āyā nashenīde ke Ḥosayn ebn-e ʿAlī rūḥu man fī
malakūt al-ibdāʿ fidāhumā be-che shawq va eshtiyāq jān dar rah-e Ḥaqq
bākht va agar begūʾī chenānche molḥedīn gofte-and ke bayn-e yadey-e aʿdā
mobtalā shod va chāre joz koshte shodan nadāsht hādhā al-qawlu baghyun
minka wa-minhum ʿalā Allāh al-Muhaymin al-Qayyūm che ke dar ḥīnī ke az
Madīne ḥarekat mīfarmūd be-hamīn niyyat va erāde ḥarekat farmūd
fa-wa-Allāh be-shawq-e leqā-ye maḥbūb ḥarekat nemūd va ṭalaban
li-wiṣālihi qaṭʿ-e marāḥel farmūd chenānche dar ḥīn-e khorūj ān Jamāl-e
Aḥadiyye az Madīne makhṣūṣ-e vedāʿ be-rawże-ye moqaddase-ye sādhaj-e
vojūd men al-ghayb va al-shohūd jadd-e moṭahhar-e khod tashrīf borde va
be-īn kalemāt-e abdaʿ-e aḥlā nāṭeq.

O raven of the barren land, do not speak of the parrot of the Egypt of
oneness! O buzzing of a fly, do not mention the sound of the Most
Exalted Pen, for you disgrace yourself; do not trouble Us. Have you not
heard with what ardor and longing Ḥusayn son of ʿAlī---may the spirit of
everyone in the kingdom of creation be sacrificed for them both---gave
His life in the path of the True One? If you say, as the unbelievers
have said, that He fell into the hands of His enemies and had no choice
but to be killed, this statement is an outrage by you and them against
God, the All-Protecting, the Self-Subsisting. For when He set out from
Medina, He departed with this very intention and resolve. By God, He
departed longing to meet the Beloved and traversed the stages of His
journey in quest of union with Him. Thus, when that Beauty of oneness
was leaving Medina, He went expressly to bid farewell at the sacred
shrine of His sanctified grandfather, the pure essence of existence in
the unseen and the visible, and spoke these wondrous, most sweet words:

\subsection{578}\label{section-1073}

\begin{otherlanguage}{persian}

و انت تعلم يا محبوب ما اردت لوجه اللّه معتمداً فان الصّبر منقطع منّى لحبی‌
جمال اللّه منكشفا و انت تعلم ما اراد ابن الزّنا فى دمى متعمّداً لا و حضرة
عزّك لا ابايع به لاخفية و لا جهراً اللّه قرّب يوم دمى ثمّ دمعى على التّراب
متّكئاً فيا ليت يومى يوم دمى كنت بالثّرى متعطشاً تاللّه برايحة من روايح الّتى
هبت من رضوان حبّه للّه محبوبه عطر كلّ الوجود و استجذب اهل ممالك الغيب و
الشّهود و من كلماته فى ذلك المقام تحيّرت اهل ملأ الأعلی و كلّ الوجود ناح
نوح المشتاقين و بكى بكآء العاشقين و بنار حبّه اشتعلت مشاعل الحبّ بين الأرض
و السّمآء و بتوجّهه الى ارض الطفّ لفدآء نفسه للّه قد توجّهت الممكنات الى شطر
اللّه و لكن اينمقام را محبّ خالص و عاشق صادق ادراك نموده ‌و مينمايد و قد
جعله اللّه فوق عرفانك و عرفان مرشديك

\end{otherlanguage}

Wa-anta taʿlamu yā maḥbūb mā aradtu li-wajh Allāh muʿtamidan fa-inna
al-ṣabra munqaṭiʿun minnī li-ḥubbī jamāl Allāh munkashifan wa-anta
taʿlamu mā arāda ibn al-zinā fī damī mutaʿammidan lā wa-ḥaḍrat ʿizzika
lā ubāyiʿu bihi lā khufyatan wa-lā jahran Allāh qarrib yawma damī thumma
damʿī ʿalā al-turāb muttakiʾan fa-yā layta yawmī yawma damī kuntu
bi-al-tharā mutaʿaṭṭishan ta-Allāh bi-rāʾiḥatin min rawāʾiḥ allatī
habbat min riḍwān ḥubbihi li-llāh maḥbūbihi ʿaṭṭara kull al-wujūd
wa-istajdhaba ahla mamālik al-ghayb wa-al-shuhūd wa-min kalimātihi fī
dhālika al-maqām taḥayyarat ahlu malaʾ al-aʿlā wa-kull al-wujūd nāḥa
nawḥ al-mushtāqīn wa-bakā bukāʾ al-ʿāshiqīn wa-bi-nār ḥubbihi ishtaʿalat
mashāʿil al-ḥubb bayna al-arḍ wa-al-samāʾ wa-bi-tawajjuhihi ilā arḍ
al-Ṭaff li-fidāʾ nafsihi li-llāh qad tawajjahat al-mumkināt ilā shaṭr
Allāh va lāken īn maqām rā moḥebb-e khāleṣ va ʿāsheq-e ṣādeq edrāk
nemūde va mīnemāyad wa-qad jaʿalahu Allāhu fawqa ʿirfānika wa-ʿirfān
murshidayk.

``You know, O Beloved, what I have intended for the sake of God, relying
upon Him. My patience has run out because of my love for the unveiled
beauty of God. You know what the son of adultery has deliberately
intended against my blood. No, by Your glorious presence, I will not
pledge allegiance to him, either secretly or openly. O God, bring near
the day of my bloodshed, then my tears, as I lie upon the dust. If only
my day were the day of my bloodshed, when I would lie thirsting in the
earth!'' By God, with a single fragrance among those that wafted from
the paradise of His love for God, His Beloved, He perfumed all existence
and drew to Himself the inhabitants of the kingdoms of the unseen and
the visible. At His words on that occasion, the inhabitants of the
heavenly Concourse were bewildered, and all existence lamented with the
lamentation of those who yearn and wept with the weeping of lovers.
Through the fire of His love, the torches of love were kindled between
earth and heaven, and through His turning toward the land of Ṭaff to
sacrifice Himself for God, all created things turned toward God. Yet
only a pure lover and a true devotee has comprehended, and comprehends,
this station. God has placed it beyond your understanding and that of
your two guides.

\subsection{579}\label{section-1074}

\begin{otherlanguage}{persian}

خذ زمامك يا ايّها المستهزء باللّه و احبّائه و اصفيائه و اودّائه و انّ هذا
الحسين يقول حينئذ قد جعل اللّه هذا السّجن ارض الطّفّ لنفسى و حينئذ اكون
حاضراً لسيوف القضا و مشتاقاً لسهام البلا و لن اخاف من احد و لن افرّ من نفس
و لن استر وجهى و انّه كان مستضيئاً بين السّموات و الأرضين و اقول يا ايّها
السّنان هلم بسنانك و يا ايّها الخولى فأت برمحك و يا ايّها الشّمر اسرع بخنجرك
و انّى انفقت نفسى و روحى و ذاتى فى سبيل اللّه المهيمن العزيز القيّوم در كلّ
ایّام منتظر شهادت بوده و هستم و مطّلعم آنچه تو باغواى شجره فساد ببعضى از
اركان گفته و لكن بدان كه عمل حقّ از دونش واضح و ممتاز بوده و در هيچ امرى
احدى شبيه نه او است مظهر توحيد و مطّلع تفريد مابين عباد در وقتى كه جمال
ابهى در مدينه كبير وارد شدند چند فردى فرموده‌اند بعضى از آن در اينمقام
ذكر ميشود لعلّ تتذكّر و تخشى و تمسك قلمك عن المفتريات و لعلّ تكون من
الرّاجعين و التّائبين الى اللّه الّذى خلقك و خلق كلّ شیء و كلّ اليه لمن
الرّاجعين و آن اين است

\end{otherlanguage}

Khudh zimāmaka yā ayyuhā al-mustahziʾ bi-llāh wa-aḥibbāʾihi
wa-aṣfiyāʾihi wa-awiddāʾihi wa-inna hādhā al-Ḥusayn yaqūlu ḥīnaʾidh qad
jaʿala Allāhu hādhā al-sijn arḍ al-Ṭaff li-nafsī wa-ḥīnaʾidh akūnu
ḥāḍiran li-suyūf al-qaḍāʾ wa-mushtāqan li-sihām al-balāʾ wa-lan akhāfa
min aḥad wa-lan afirra min nafs wa-lan astura wajhī wa-innahu kāna
mustaḍīʾan bayna al-samāwāt wa-al-araḍīn wa-aqūlu yā ayyuhā al-Sinān
halumma bi-sinānika wa-yā ayyuhā al-Khawlī faʾti bi-rumḥika wa-yā ayyuhā
al-Shimr asriʿ bi-khanjarika wa-innī anfaqtu nafsī wa-rūḥī wa-dhātī fī
sabīl Allāh al-Muhaymin al-ʿAzīz al-Qayyūm dar koll-e ayyām montaẓer-e
shahādat būde va hastam va moṭṭaleʿam ānche to be-eghvā-ye shajare-ye
fasād be-baʿżī az arkān gofte va lāken bedān ke ʿamal-e Ḥaqq az dūnash
vāżeḥ va momtāz būde va dar hīch amrī aḥadī shabīh na ū ast maẓhar-e
tawḥīd va moṭṭaleʿ-e tafrīd mābeyn-e ʿebād dar vaqtī ke Jamāl-e Abhā dar
Madīne-ye Kabīr vāred shodand chand fardī farmūde-and baʿżī az ān dar īn
maqām dhekr mīshavad laʿallaka tatadhakkaru wa-takhshā wa-tumsiku
qalamaka ʿan al-muftarayāt wa-laʿallaka takūnu min al-rājiʿīn
wa-al-tāʾibīn ilā Allāh alladhī khalaqaka wa-khalaqa kulla shayʾ
wa-kullun ilayhi la-min al-rājiʿīn va ān īn ast.

Rein yourself in, O you who mock God, His loved ones, His chosen ones,
and His friends! This Ḥusayn now says: God has made this prison the land
of Ṭaff for Me. Now I stand ready for the swords of destiny and yearn
for the arrows of affliction. I will fear no one, flee from no soul, nor
conceal My face, which has shone between the heavens and the earths. I
say: O Sinān, come with your spearhead! O Khawlī, bring your spear! O
Shimr, hasten with your dagger! I have offered up My self, My spirit,
and My very being in the path of God, the All-Protecting, the Almighty,
the Self-Subsisting. Every day I have awaited martyrdom, and I still do.
I am aware of what, at the prompting of the tree of corruption, you have
said to certain leading figures. But know that the action of the True
One has been clear and distinct from that of anyone else, and no one
resembles Him in any matter. He is the Manifestation of unity and is
privy to singleness among His servants. When the All-Glorious Beauty
arrived in the Great City, He uttered several verses, some of which are
mentioned here, so that perhaps you may take heed, stand in awe, and
restrain your pen from fabrications, and perhaps be among those who
return and repent to God, who created you and created all things. All
will surely return to Him. They are these:

\subsection{580}\label{section-1075}

\begin{otherlanguage}{arabic}

اى صبا از پيش جانان يكزمان خوش بران تا كوى آن زورائيان

\end{otherlanguage}

Ey ṣabā az pīsh-e jānān yek zamān khosh berān tā kūy-e ān Zawrāʾiyān.

O morning breeze, from the Beloved's presence, speed gently for a moment
to the quarter of the people of Zawrā.

\subsection{581}\label{section-1076}

\begin{otherlanguage}{persian}

پس بگويش كى مدينه كردگار چون بماندى چونكه رفت از برت يار

\end{otherlanguage}

Pas begūyash k-ey Madīne-ye Kerdgār chūn bemāndī chūnke raft az barat
yār.

Then say to her: ``O city of the Creator, how have you endured, now that
the Friend has left your side?

\subsection{582}\label{section-1077}

\begin{otherlanguage}{arabic}

يار تو در دست اعدا مبتلا چون حسين اندر زمين كربلا يكحسين

\end{otherlanguage}

Yār-e to dar dast-e aʿdā mobtalā chūn Ḥosayn andar zamīn-e Karbalā yek
Ḥosayn.

Your Friend is afflicted in the hands of enemies, like Ḥusayn on the
land of Karbalā---one Ḥusayn,

\subsection{583}\label{section-1078}

\begin{otherlanguage}{arabic}

و صدهزارانش يزيد يك حبيب و صد ابوجهل عنيد

\end{otherlanguage}

Va ṣad hazārānash Yazīd yek Ḥabīb va ṣad Abū Jahl-e ʿanīd.

And against Him a hundred thousand Yazīds; one Beloved, and a hundred
obstinate Abū Jahls.

\subsection{584}\label{section-1079}

\begin{otherlanguage}{arabic}

چون كليم اندر ميان قبطيان يا چه روح اللّه ميان جاحدان

\end{otherlanguage}

Chūn Kalīm andar miyān-e Qebṭiyān yā che Rūḥ Allāh miyān-e jāḥedān.

Like Moses among the Egyptians, or like the Spirit of God among the
deniers.

\subsection{585}\label{section-1080}

\begin{otherlanguage}{persian}

همچه يوسف اندر افتاده بچاه آن چهى كه نبودش پايان و راه

\end{otherlanguage}

Hamcho Yūsof andar oftāde be-chāh ān chahī ke nabūdash pāyān va rāh.

Like Joseph, fallen into a well---a well with neither bottom nor way
out.''

\subsection{586}\label{section-1081}

\begin{otherlanguage}{persian}

فرمودند آنچه ظاهر شده كه شده و مابقى هم ظاهر خواهد شد و ما كان مؤسّسه الّا
الّذين افتروا عليه بين العباد و انّك احد منهم و كان علم اللّه على ما فعلتم
محيط و حينئذ يقول يا ملأ البيان انّ هذا رأسى قد كان عريا بين السّموات و
الأرض و كان منتظرا لاسيافكم اذاً فاضربوه كيف شتم و لا تكوننّ من المتوقّفين
و ان هذا صدرى مشتاقاً لسهام البغضآء اذا فاضربوه كيف اردتم يا ملأ المفترين
و انّ هذا حنجرى يشتاق خنجركم ان اقطعوه لانّا انفقناه فى سبيل محبوبی‌ و
محبوب العالمين و نشكره فى كلّ ذلك و نحمده و انّه لمقصود روحى و ما ظهر و
يظهر من عنده قد كان مقصودى لو انتم من الشّاعرين

\end{otherlanguage}

Farmūdand ānche ẓāher shode ke shode va mābaqī ham ẓāher khāhad shod
wa-mā kāna muʾassisuhu illā alladhīna iftaraw ʿalayhi bayna al-ʿibād
wa-innaka aḥadun minhum wa-kāna ʿilm Allāh ʿalā mā faʿaltum muḥīṭ
wa-ḥīnaʾidh yaqūlu yā malaʾ al-bayān inna hādhā raʾsī qad kāna ʿuryan
bayna al-samāwāt wa-al-arḍ wa-kāna muntaẓiran li-asyāfikum idhan
fa-iḍribūhu kayfa {[}corrupted text{]} wa-lā takūnanna min
al-mutawaqqifīn wa-inna hādhā ṣadrī mushtāqan li-sihām al-baghḍāʾ idhan
fa-iḍribūhu kayfa aradtum yā malaʾ al-muftarīn wa-inna hādhā ḥanjarī
yashtāqu khanjarakum an iqṭaʿūhu li-annā anfaqnāhu fī sabīl maḥbūbī
wa-maḥbūb al-ʿālamīn wa-nashkuruhu fī kull dhālika wa-naḥmaduhu
wa-innahu la-maqṣūd rūḥī wa-mā ẓahara wa-yaẓharu min ʿindihi qad kāna
maqṣūdī law antum min al-shāʿirīn.

He said: What has appeared has appeared, and the rest will also appear.
None laid its foundation but those who fabricated lies against Him among
His servants, and you are one of them. God's knowledge encompasses what
you have done. Now He says: ``O people of the Bayān! Here is My head,
laid bare between the heavens and the earth, awaiting your swords.
Strike it, then, however {[}corrupted original text{]}, and do not be
among those who hesitate. Here is My breast, yearning for the arrows of
hatred. Strike it, then, however you wish, O company of fabricators!
Here is My throat, longing for your dagger. Cut it, for We have offered
it up in the path of My Beloved and the Beloved of all worlds. In all
this We thank Him and praise Him. He is indeed the desire of My soul,
and whatever has appeared and appears from Him has been My desire, if
only you could perceive.''

\subsection{587}\label{section-1082}

\begin{otherlanguage}{persian}

بدانكه شأن انبیاء بلندتر از آن است كه امثال تو بتوانند ادراك نمايند و آن
شموس احديّه لازال بشوق و اشتياق الهى بوده‌اند و آنچه در سبيل محبوب بر آن
نفوس قدسى وارد كلّ راضى و شاكر بوده‌اند و دقيقه در تبليغ امر اللّه اهمال
نفرموده‌اند آيا نشنيدى كه نوح چه مقدار در امر اللّه نوحه نمود در كلّ ایّام
مقابل قوم ميايستاد و ناس را بشريعه عزّ ربانيّه دعوت ميفرمود و در هر يوم
بعد از تبليغ بشأنى بر ايذاى آن مظهر امر مبالغه و اصرار مينمودند كه وصف
آن ممكن نه و در بعض ایّام از شدّت بلايا كه بر آنحضرت وارد ميامد گمان موت
مينمودند كلّ اين امور را حبّا للّه حمل ميفرمودند و اگر از خوف انبیاء ذكرى
در كتاب ديده چنانچه در مقدّمه كليم در كتاب اللّه مسطور امر او را نه تو و
نه مرشدينت هيچكدام ادراك ننموده‌ايد اگر بخواهى ادراك نمائى فاسئل ليبيّن
اللّه لك ما عنده و انّ عنده علم كلّ شیء و علم السّموات و الأرضين

\end{otherlanguage}

Bedānke shaʾn-e anbiyā bolandtar az ān ast ke amthāl-e to betavānand
edrāk namāyand va ān shomūs-e aḥadiyye lāzāl be-shawq va eshtiyāq-e
elāhī būde-and va ānche dar sabīl-e maḥbūb bar ān nofūs-e qodsī vāred
koll rāżī va shāker būde-and va daqīqe dar tablīgh-e amr Allāh ehmāl
nafarmūde-and āyā nashenīdī ke Nūḥ che meqdār dar amr Allāh nawḥe nemūd
dar koll-e ayyām moqābel-e qawm mī-īstād va nās rā be-sharīʿe-ye ʿezz-e
rabbāniyye daʿvat mīfarmūd va dar har yawm baʿd az tablīgh be-shaʾnī bar
īdhā-ye ān maẓhar-e amr mobāleghe va eṣrār mīnemūdand ke vaṣf-e ān
momken na va dar baʿż-e ayyām az sheddat-e balāyā ke bar ān ḥażrat vāred
mī-āmad gomān-e mawt mīnemūdand koll-e īn omūr rā ḥobban li-llāh ḥaml
mīfarmūdand va agar az khawf-e anbiyā dhekrī dar ketāb dīde chenānche
dar moqaddame-ye Kalīm dar Ketāb Allāh masṭūr amr-e ū rā na to va na
morshedīnat hīch kodām edrāk nanemūde-īd agar bekhāhī edrāk namāʾī
fa-isʾal li-yubayyina Allāhu laka mā ʿindahu wa-inna ʿindahu ʿilma kull
shayʾ wa-ʿilma al-samāwāt wa-al-araḍīn.

Know that the station of the Prophets is too exalted for the likes of
you to comprehend. Those Suns of oneness have always been filled with
divine ardor and longing. Whatever befell those holy Souls in the path
of the Beloved, They were all content and thankful, and They did not
neglect the proclamation of the Cause of God for a moment. Have you not
heard how greatly Noah lamented in the Cause of God? Every day He stood
before His people and summoned them to the waters of divine glory. Each
day, after He had delivered His message, they persisted in tormenting
that Manifestation of the Cause with an intensity beyond description. On
some days, so severe were the afflictions that befell Him that they
thought He was dead. He bore all these things out of love for God. If
you have seen mention in the Book of the Prophets' fear, as is recorded
in the Book of God in the account of Moses, neither you nor your guides
have understood His case. If you wish to understand, then ask, so that
God may explain to you what is with Him. With Him is the knowledge of
all things and the knowledge of the heavens and the earths.

\subsection{588}\label{section-1083}

\begin{otherlanguage}{persian}

و امّا حكايت ذكريّا مقدّمه او را اگر شنيده قدرى تفكّر كن كه از همان حكايت
تفصيل معلوم ميشود و اين عبد مصلحت نديده كه آنچه از منبع علم لدنّى افاضه
شده مذكور دارم آنقدر بدان كه جميع انبیاء در كمال اشتياق جان در سبيل حقّ
تسليم نموده‌اند و بعضى هم كه بشهادت ظاهريّه نرسيده‌اند در باطن بآن نفوس
مقدّسه سمّ داده‌اند و يا بتلبيسات ديگر و علمه عند ربّی‌ چنانچه مرشدت و اللّه
الّذى لا اله الّا هو كه بهمين امور مشغول شد و لكنّ اللّه حفظ مظهر نفسه
بسلطانه و انّه لهو المقتدر القدير

\end{otherlanguage}

Va ammā ḥekāyat-e Zakariyyā, moqaddame-ye ū rā agar shenīde, qadrī
tafakkor kon ke az hamān ḥekāyat tafṣīl maʿlūm mīshavad. Va īn ʿabd
maṣlaḥat nadīde ke ānche az manbaʿ-e ʿelm-e ladonnī efāże shode madhkūr
dāram. Ānqadr bedān ke jamīʿ-e anbiyā dar kamāl-e eshtiyāq jān dar
sabīl-e Ḥaqq taslīm namūde-and, va baʿżī ham ke be-shahādat-e ẓāheriyye
naresīde-and dar bāṭen be-ān nofūs-e moqaddase samm dāde-and va yā
be-talbīsāt-e dīgar, wa ʿilmuhu ʿinda rabbī. Chenānche morshedat,
wa-Allāhi alladhī lā ilāha illā huwa, ke be-hamīn omūr mashghūl shod, wa
lākinna Allāha ḥafiẓa maẓhara nafsihi bi-sulṭānihi wa innahu la-huwa
al-Muqtadiru al-Qadīr.

As for the account of Zechariah, if you have heard His story, reflect on
it a little, for the details will become clear from that very account.
This servant has not deemed it advisable to mention what has been
imparted from the source of innate knowledge. Know only this: all the
Prophets surrendered Their lives in the path of God with the utmost
longing. Even some who did not outwardly suffer martyrdom were secretly
poisoned, or subjected to other deceptions; knowledge of this rests with
my Lord. Your guide, too---I swear by God, besides whom there is no
god---occupied himself with these very things. But God protected the
Manifestation of His own Self through His sovereignty. He is indeed the
All-Powerful, the Almighty.

\subsection{589}\label{section-1084}

\begin{otherlanguage}{persian}

و امّا حكايت نقطه اولى روح ما سواه فداه وقتيكه از فارس با چاپار متعدّده آن
هيكل ربانيّه را مياوردند كه بآذربايجان ببرند بعضى از اصحاب از جمله جناب
عظيم عليه بهآء اللّه نوشتند ببعضى از احبّاى ساكنه در ارض طاء كه ما همرهان
را راضى نموده‌ايم كه قدرى زخارف بدهيم و حضرت را در شب از ما بين آن قوم
برده و مستور داريم و اين امر را در پيشگاه حضور اقدس اعلى جلّ كبريائه
معروض داشتند آن جمال احديّه تبسّم فرمودند و از تبسّم مبارك اصحاب چنان ادراك
نمودند كه راضيند بآنچه اراده نموده‌اند و بعد از اجتماع اسباب در وقتيكه در
كلين تشريف داشتند در محضر اقدس معروض داشتند كه امشب بايد از ما بين اين
قوم تشريف ببريد مجدّد تبسّمى فرمودند و فرمودند جبل آذربايجان هم قسمتى دارد
هذا من فعل الحقّ و ظهر من الحقّ و كان على الحقّ و سطع منه رايحة الحقّ و لكن
انت و امثالك لا تجدن و تكوننّ من الغافلين

\end{otherlanguage}

Va ammā ḥekāyat-e Noqṭe-ye Ūlā, rūḥ-e mā sevāh fedāh, vaqtī ke az Fārs
bā chāpār-e motaʿaddede ān haykal-e rabbāniyye rā mīāvardand ke
be-Ādharbāyjān bebarand, baʿżī az aṣḥāb, az jomle Janāb-e ʿAẓīm, ʿalayhi
bahāʾu Allāh, neveshtand be-baʿżī az aḥebbā-ye sākene dar arż-e Ṭā ke mā
hamrahān rā rāżī namūde-īm ke qadrī zakhāref bedehīm va Ḥażrat rā dar
shab az mā beyn-e ān qowm borde va mastūr dārīm. Va īn amr rā dar
pīshgāh-e ḥożūr-e aqdas-e aʿlā, jalla kibriyāʾuhu, maʿrūż dāshtand. Ān
Jamāl-e Aḥadiyye tabassom farmūdand va az tabassom-e mobārak aṣḥāb
chenān edrāk namūdand ke rāżī-and be-ānche erāde namūde-and. Va baʿd az
ejtemāʿ-e asbāb, dar vaqtī ke dar Kolayn tashrīf dāshtand, dar maḥżar-e
aqdas maʿrūż dāshtand ke emshab bāyad az mā beyn-e īn qowm tashrīf
bebarīd. Mojaddad tabassomī farmūdand va farmūdand: Jabal-e Ādharbāyjān
ham qesmatī dārad. Hādhā min fiʿli al-Ḥaqqi wa ẓahara mina al-Ḥaqqi wa
kāna ʿalā al-ḥaqqi wa saṭaʿa minhu rāʾiḥatu al-ḥaqqi wa lākin anta wa
amthāluka lā tajidunna wa takūnanna mina al-ghāfilīn.

As for the account of the Primal Point---may the spirit of all else be
sacrificed for Him---when that divine Being was being brought from Fars
with successive mounted couriers to be taken to Azerbaijan, some of His
companions, including Jináb-i-ʿAẓím, upon him be the glory of God, wrote
to some of the believers living in the Land of Ṭá: ``We have secured the
agreement of those accompanying Him to accept a little money, so that we
may take Him away from among these people at night and conceal Him.''
They submitted this matter to His most holy and exalted
presence---glorified be His majesty. That Beauty of Oneness smiled, and
from His blessed smile the companions understood that He approved of
what they intended. After the arrangements had been made, while He was
in Kulayn, they submitted in His most holy presence: ``Tonight You must
leave these people.'' Again He smiled and said: ``The mountain of
Azerbaijan also has its share.'' This was among the deeds of the True
One; it proceeded from the True One, was founded on truth, and shed the
fragrance of truth. Yet you and those like you do not perceive it and
remain among the heedless.

\subsection{590}\label{section-1085}

\begin{otherlanguage}{persian}

فرمودند كه بنويس بآن نفس كه يك سخن للّه گفته ميشود استماع نما بعداوت اين
مسجون هتك حرمت انبیاء مكن و ستر عظمت را شقّ منما و حرمة اللّه را ضايع مكن
خف عن اللّه الّذى خلقك و كلّ شیء و لا تكن من المفترين على انبیاء اللّه و
امنائه و اصفيائه فو‌الّذى نفسى بيده كلّ واحد منهم فى كلّ يوم ينادى ربّه و
يقول يا ليت لى الف روح و الف جسد و الف نفس لا فديها فى سبيلك يا محبوب
العالمين و يا مقصود المشتاقين و يا وله صدور العاشقين و اگر هم بخواهى در
اينمقامات مثالى بزنى بنفس اين مسجون راجع كن مثل افعال ديگر كه راجع نموده‌
اذيال عصمت الهى را بغبار نفس و هوى تيره مكن هذه نصيحة منّى عليك ان سمعت
فلنفسك و ان اعرضت فعليها و كان اللّه بنا و بك رقيب و شهيد

\end{otherlanguage}

Farmūdand ke benevīs be-ān nafs ke yek sokhan li-llāh gofte mīshavad,
estemāʿ namā. Be-ʿadāvat-e īn masjūn hatk-e ḥormat-e anbiyā makon va
setr-e ʿaẓamat rā shaqq manamā va ḥormat Allāh rā żāyeʿ makon. Khaf ʿani
Allāhi alladhī khalaqaka wa kulla shayʾin wa lā takun mina al-muftarīna
ʿalā anbiyāʾi Allāhi wa umanāʾihi wa aṣfiyāʾihi. Fa-wa-alladhī nafsī
bi-yadihi kullu wāḥidin minhum fī kulli yawmin yunādī rabbahu wa yaqūlu
yā layta lī alfa rūḥin wa alfa jasadin wa alfa nafsin la-afdiyahā fī
sabīlika yā Maḥbūba al-ʿālamīna wa yā Maqṣūda al-mushtāqīna wa yā walaha
ṣudūri al-ʿāshiqīn. Va agar ham bekhāhī dar īn maqāmāt methālī bezanī
be-nafs-e īn masjūn rājeʿ kon, methl-e afʿāl-e dīgar ke rājeʿ namūde.
Adhyāl-e ʿeṣmat-e elāhī rā be-ghobār-e nafs va havā tīre makon. Hādhihi
naṣīḥatun minnī ʿalayka in samiʿta fa-li-nafsika wa in aʿraḍta
fa-ʿalayhā wa kāna Allāhu binā wa bika raqībun wa shahīd.

He said: Write to that person: A word is spoken for God's sake; listen
to it. Do not, out of hostility toward this Prisoner, violate the
sanctity of the Prophets, tear apart the veil of grandeur, or dishonor
the sanctity of God. Fear God, who created you and all things, and do
not be among those who fabricate lies against God's Prophets, His
trusted ones, and His chosen ones. By Him in whose hand is my soul!
Every one of Them calls upon His Lord each day and says: ``If only I had
a thousand spirits, a thousand bodies, and a thousand souls to sacrifice
in Your path, O Beloved of the worlds, O Desire of those who yearn, O
ardent longing in the hearts of lovers!'' If you wish to offer an
example in these matters, refer it to this Prisoner Himself, just as you
have attributed other deeds to Him. Do not sully the hems of divine
purity with the dust of self and desire. This is My counsel to you: if
you listen, it is for your own good; if you turn away, it is to your own
detriment. God watches over us and you and is a witness.

\subsection{591}\label{section-1086}

\begin{otherlanguage}{persian}

بشنو نداى اين عبد را و نصيحت كن آن مفترين را بس است در تضييع كوشيدند امر
بجائى رسيده كه نزد كلّ نفوس رفته متّصلا افترا گفته و ميگويند قل خافوا عن
اللّه و لاتفتروا عليه ان اتّقوا فى انفسكم انّه حفظكم فى سنين معدودات و ستر
اعمالكم و افعالكم و انتم لما اطمئنتم فى انفسكم خرجتم عن خلف القناع و
افتريتم عليه بما لا افترى به احد من المفسدين و يشهد بذلك لسان اللّه الملك
المقتدر العزيز الكريم

\end{otherlanguage}

Beshnow nedā-ye īn ʿabd rā va naṣīḥat kon ān moftarīn rā. Bas ast dar
tażyīʿ kūshīdand. Amr be-jāʾī resīde ke nazd-e koll-e nofūs rafte
mottaṣelan efterā gofte va mīgūyand. Qul khāfū ʿani Allāhi wa lā taftarū
ʿalayhi an ittaqū fī anfusikum innahu ḥafiẓakum fī sinīnin maʿdūdātin wa
satara aʿmālakum wa afʿālakum wa antum lammā iṭmaʾnantum fī anfusikum
kharajtum ʿan khalfi al-qināʿi wa iftaraytum ʿalayhi bi-mā lā iftarā
bihi aḥadun mina al-mufsidīna wa yashhadu bi-dhālika lisānu Allāhi
al-Maliki al-Muqtadiri al-ʿAzīzi al-Karīm.

Listen to this servant's call and counsel those fabricators. Enough of
their efforts to bring dishonor! Matters have reached the point where
they have gone to everyone, continually uttering slanders, and they
continue to do so. Say: Fear God and do not fabricate lies against Him.
Be mindful of your conduct. He protected you for a number of years and
concealed your deeds and actions. Yet when you felt secure, you emerged
from behind the veil and fabricated against Him what no spreader of
corruption had ever fabricated. To this bears witness the Tongue of God,
the King, the All-Powerful, the Almighty, the Generous.

\subsection{592}\label{section-1087}

\begin{otherlanguage}{persian}

قوله عظم مكره و كبره مذكور نموديد كه جناب ميرزا بتنهائى قيام در امر حقّ
نموده‌ از كجا بايد تميز داد كه اين قيام بجهت ارتفاع امر اللّه بوده و يا
بجهت حبّ جاه و رياست و شوكت و ثروت بلى قيام در امر اللّه اشخاصى نمودند كه
در مازندران علف و چرم خوردند نه اشخاصيكه در بغداد و جاى ديگر ایّام و
ليالى در صرف چاى با سماور و اوضاع و دستگاه با رفقا و خدم و قهوه‌خانها با
اجامر و اوباش و صحبتهاى لاطايل و نهار و غذاهاى رنگين بسر رود و همّت مصروف
اطعمه و شيرينيها و مهمانيها شود و اصحاب در بدر جمع شده راحت نمايند چه
قيام در امر اللّه شده انتهى

\end{otherlanguage}

Qawluhu, ʿaẓuma makruhu wa kibruhu: Madhkūr namūdīd ke Janāb-e Mīrzā
be-tanhāʾī qiyām dar amr-e Ḥaqq namūde. Az kojā bāyad tamīz dād ke īn
qiyām be-jehat-e ertefāʿ-e amr Allāh būde va yā be-jehat-e ḥobb-e jāh va
reyāsat va showkat va thervat? Balī, qiyām dar amr Allāh ashkhāṣī
namūdand ke dar Māzandarān ʿalaf va charm khordand, na ashkhāṣī ke dar
Baghdād va jā-ye dīgar ayyām va layālī dar ṣarf-e chāy bā samāvar va
owżāʿ va dastgāh bā rofaqā va khadam va qahve-khānehā bā ajāmer va
owbāsh va ṣoḥbathā-ye lāṭāyel va nahār va ghadhāhā-ye rangīn be-sar
ravad va hemmat maṣrūf-e aṭʿeme va shīrīnihā va mehmānihā shavad va
aṣḥāb dar be-dar jamʿ shode rāḥat namāyand. Che qiyām dar amr Allāh
shode? Intahā.

His words---great are his deceit and arrogance: ``You have stated that
Jináb-i-Mírzá arose alone to serve the Cause of God. How can one
determine whether this undertaking was for the exaltation of the Cause
of God or for the love of rank, leadership, pomp, and wealth? Yes, those
who arose to serve the Cause of God were the people who ate grass and
leather in Mazandaran, not those whose days and nights in Baghdad and
elsewhere are spent drinking tea, with samovars and all their
arrangements and trappings, with friends and servants; in coffeehouses
with rabble and riffraff, in idle conversations, and over lunches and
sumptuous meals; whose efforts are devoted to food, sweets, and
banquets; and whose companions gather from door to door to take their
ease. What service to the Cause of God has been undertaken?'' End of
quotation.

\subsection{593}\label{section-1088}

\begin{otherlanguage}{persian}

يشهد بذلك كلّ الذّرّات لو نفسك الخبيثة تكون من المنكرين بعد از مقدّمه ارض
طاء نار حبّ الهى در كلّ قلوب مخمود و سراج ودّش در افئده عارفين خاموش چنانچه
در سنين اوّليّه كه در عراق وارد شدند ابدا هبوب ارياح امر از جهتى نوزيده و
نفحات قميص محبّت رحمن از شطرى استشمام نشده بشأنى امر صعب بوده كه كلّ قلوب
مضطرب و اكثرى از نفوس متزلزل و اگر هم در بعضى از ديار بعضى انفس معدوده
بوده مخمود و محتجب و مستور تو انصاف ده كه در مقابل اعدا كه ايستاد آنوقت
نه عزّتى بود و نه نفوسى تا رياست محقّق شود لذا بنفسه وحده در عراق ساكن

\end{otherlanguage}

Yashhadu bi-dhālika kullu al-dharrāti law nafsuka al-khabīthatu takūnu
mina al-munkirīn. Baʿd az moqaddame-ye arż-e Ṭā nār-e ḥobb-e elāhī dar
koll-e qolūb makhmūd va serāj-e voddash dar afʾede-ye ʿārefīn khāmūsh.
Chenānche dar senīn-e avvaliyye ke dar ʿErāq vāred shodand abadan
hobūb-e aryāḥ-e amr az jehatī navazīde va nafaḥāt-e qamīṣ-e maḥabbat-e
Raḥmān az shaṭrī esteshmām nashode. Be-shaʾnī amr ṣaʿb būde ke koll-e
qolūb możṭareb va aktharī az nofūs motazalzel va agar ham dar baʿżī az
deyār baʿżī anfos-e maʿdūde būde makhmūd va moḥtajeb va mastūr. To enṣāf
deh ke dar moqābel-e aʿdā ke īstād? Ān vaqt na ʿezzatī būd va na nofūsī
tā reyāsat moḥaqqaq shavad. Ledhā be-nafsehe vaḥdeh dar ʿErāq sāken.

Every atom bears witness to this, even if your wicked self is among
those who deny it. After the events in the Land of Ṭá, the fire of
divine love was extinguished in every heart, and the lamp of His love
had gone out in the hearts of those who knew Him. Thus, in the early
years after His arrival in Iraq, no breeze of the Cause blew from any
direction, nor was the fragrance of the garment of the All-Merciful's
love perceived from any quarter. Matters were so difficult that every
heart was troubled and most souls were shaken. Even if a few souls
remained in certain lands, they were dispirited, veiled, and concealed.
Judge fairly: who stood against the enemies? At that time there was
neither prestige nor a body of followers over whom leadership could be
exercised. Thus He dwelt in Iraq by Himself, alone.

\subsection{594}\label{section-1089}

\begin{otherlanguage}{persian}

و چون در نظر اهل عراق اين طايفه حقير و بی‌‌مقدار بودند بشأنى كه جز لعن و
سبّ و شتم از احدى نسبت باين طايفه مذكور ميشد لذا احدى مراوده نمينمود تا
از كوثر بيان معارف الهى و سلسبيل حكم نامتناهى ربّانى مرزوق شود محض مخالطه
مدّتى بقهوه كه در يم شط واقع بود تشريف بردند تا آنكه در آن محلّ با بعضى از
اهل عراق گفتگو شد و بعد بناى مراوده گذاردند و بشأنى امر اللّه معزّز شد كه
كلّ بثنآءاللّه ناطق و بشطر مقصود مايل گشتند اى بی‌ بصر اگر بقدر خفّاش بصر
ميداشتى هرآينه تصديق مينمودى كه جلوسه فى القهوه فى اقلّ من آن لخير منك و
من وجودك و عمّا عملت و تعمل ثمّ عمل مثلائك

\end{otherlanguage}

Va chon dar naẓar-e ahl-e ʿErāq īn ṭāʾefe ḥaqīr va bī-meqdār būdand
be-shaʾnī ke joz laʿn va sabb va shatm az aḥadī nesbat be-īn ṭāʾefe
madhkūr mīshod, ledhā aḥadī morāvede namīnamūd tā az kawthar-e bayān-e
maʿāref-e elāhī va salsabīl-e ḥekam-e nāmotanāhī-ye rabbānī marzūq
shavad. Maḥż-e mokhāleṭe moddatī be-qahve ke dar yam-e shaṭṭ vāqeʿ būd
tashrīf bordand tā ānke dar ān maḥall bā baʿżī az ahl-e ʿErāq goftogū
shod va baʿd banā-ye morāvede godhārdand. Va be-shaʾnī amr Allāh
moʿazzaz shod ke koll be-thanāʾ Allāh nāṭeq va be-shaṭr-e Maqṣūd māyel
gashtand. Ey bī-baṣar, agar be-qadr-e khoffāsh baṣar mīdāshtī har-āyene
taṣdīq mīnamūdī ke julūsuhu fī al-qahwati fī aqalla min ānin la-khayrun
minka wa min wujūdika wa ʿammā ʿamilta wa taʿmalu thumma ʿamalu
mithlāʾik.

Since the people of Iraq regarded this community as lowly and
insignificant, to such a degree that nothing but curses, abuse, and
insults was spoken of it, no one associated with them to receive
sustenance from the heavenly stream of the exposition of divine
knowledge and the sweet waters of boundless divine wisdom. Solely for
the sake of mingling with people, He went for a time to a coffeehouse
situated by the waters of the river. There He entered into conversation
with some of the people of Iraq, who then began to visit Him. The Cause
of God became so exalted that all spoke in praise of God and turned
toward the Desired One. O blind one! If you had even the sight of a bat,
you would assuredly acknowledge that His sitting in the coffeehouse for
less than a moment is better than you, your existence, all you have done
and will do, and the deeds of those like you.

\subsection{595}\label{section-1090}

\begin{otherlanguage}{persian}

چه قدر مشابه است اين اعتراض تو با اعتراضاتيكه يهود بعيسى بن مريم نمودند
كه اعتراضاً عليه ميگفتند كه اين ناصرى در عشارخانها با عشّار مجالس است صدق
محبوبی‌ و محبوب العارفين كه بايد آنچه از قبل در ازمنه كلّ رسل ظاهر شده در
اين ظهور بدع ربّانى ظاهر شود الحرف بالحرف و علاوه بايد نفوسى در اعراض
قيام نمايند كه اشرّ نفوس جميع اوّلين و آخرين باشند چنانچه مشاهده ميشود و
لكن چون تو بالمرّه از بصر ظاهريّه و باطنيّه محرومى لذا منكر تشعشعات انوار
شمس حقيقت و معانى بوده و خواهى بود

\end{otherlanguage}

Che qadr moshābeh ast īn eʿterāż-e to bā eʿterāżātī ke Yahūd be-ʿĪsā
ibn-e Maryam namūdand ke iʿtirāḍan ʿalayhi mīgoftand ke īn Nāṣerī dar
ʿashshār-khānehā bā ʿashshār mojāles ast. Ṣadaqa maḥbūbī va maḥbūb
al-ʿārefīn ke bāyad ānche az qabl dar azmane-ye koll-e rosol ẓāher shode
dar īn ẓohūr-e badīʿ-e rabbānī ẓāher shavad, al-ḥarfa bi-al-ḥarf. Va
ʿelāve bāyad nofūsī dar eʿrāż qiyām namāyand ke asharr-e nofūs-e jamīʿ-e
avvalīn va ākherīn bāshand, chenānche moshāhede mīshavad. Va lāken chon
to be-al-marre az baṣar-e ẓāheriyye va bāṭeniyye maḥrūmī ledhā monker-e
tashaʿshoʿāt-e anvār-e Shams-e Ḥaqīqat va Maʿānī būde va khāhī būd.

How closely this objection of yours resembles the objections the Jews
raised against Jesus, son of Mary! In protest against Him, they used to
say: ``This Nazarene sits with tax collectors in the tax houses.'' My
Beloved and the Beloved of those who know Him spoke truly: all that
appeared in the past in the ages of all the Messengers must appear in
this wondrous divine Revelation, letter for letter. Moreover, souls must
arise in opposition who are the most wicked of all people, first and
last, as can be seen. But since you are utterly deprived of both outward
and inward sight, you have denied, and will continue to deny, the
radiance of the Sun of Truth and Inner Meanings.

\subsection{596}\label{section-1091}

\begin{otherlanguage}{persian}

و بعد از علوّ امر در آن ارض از اطراف قلوب صافيه مقدّسه بشطر احديّه توجّه
نمودند و هر نفسى بين يدى حاضر چه از اهل عراق و مادونش از آنچه سؤال
نمودند بمقصود خود فايز گشتند و لكن از انكار خفّاش انوار شمس محجوب نگردد
چنانچه از انكار تو غلّا للّه و اعتراضا عليه و اعراضا منه انوار وجه مستور
نماند و در آن مدّت توقّف در عراق ساعتى آسایش نبوده و آنى راحت مشاهده نگشته
چه كه اعدا از كلّ اطراف در كمين بودند و خود تو دانسته انكار مينمائى انت
من الّذين عرفوا نعمة اللّه ثمّ انكروها قل ويل لك يا ايّها المشرك المردود

\end{otherlanguage}

Va baʿd az ʿolovv-e amr dar ān arż az aṭrāf qolūb-e ṣāfiye-ye moqaddase
be-shaṭr-e Aḥadiyye tavajjoh namūdand va har nafsī bayna yaday ḥāżer,
che az ahl-e ʿErāq va mādūnash, az ānche soʾāl namūdand be-maqṣūd-e khod
fāyez gashtand. Va lāken az enkār-e khoffāsh anvār-e shams maḥjūb
nagardad, chenānche az enkār-e to, ghillan li-llāhi wa iʿtirāḍan ʿalayhi
wa iʿrāḍan minhu, anvār-e vajh mastūr namānad. Va dar ān moddat-e
tavaqqof dar ʿErāq sāʿatī āsāyesh nabūde va ānī rāḥat moshāhede
nagashte, che ke aʿdā az koll-e aṭrāf dar kamīn būdand va khod-e to
dāneste enkār mīnamāʾī. Anta mina alladhīna ʿarafū niʿmata Allāhi thumma
ankarūhā. Qul waylun laka yā ayyuhā al-mushriku al-mardūd.

After the Cause had been exalted in that land, pure and holy hearts
turned from every direction toward the One. Every soul who entered His
presence, whether from Iraq or elsewhere, attained what they sought in
whatever they asked. Yet the bat's denial cannot veil the light of the
sun, just as your denial, born of malice toward God, protest against
Him, and rejection of Him, cannot conceal the light of His countenance.
Throughout His stay in Iraq there was not an hour of ease, nor a
moment's rest, for enemies lay in wait on every side. You yourself know
this, yet deny it. You are among those who recognized God's favor and
then denied it. Say: Woe to you, O rejected one who associates partners
with God!

\subsection{597}\label{section-1092}

\begin{otherlanguage}{persian}

و در بعضى ایّام چنان مذكور ميشد كه مأمور شده‌اند اين طايفه را گرفته بارض
طاء بفرستند مخصوص جمال ابهى امر فرمود كه كلّ از عراق بروند و اين حكم محكم
را بشأنى جارى فرمودند كه زين‌العابدين خان معروف آمده بين يدى حاضر و توسّط
از آقا سيّد حسين قنّاد نمود كه خروج او موقوف شود و اذن توقّف بفرمايند ابدا
اذن داده نشد او هم از عراق خارج شد فرمودند دوست ميدارم كه بلاياى كلّ را
وحده حامل شوم و باين عبد وارد نمايند آنچه را اراده نموده‌اند و احبّاى الهى
محفوظ مانند و تو مفترى بچنين كلمات شمس تقديس و انقطاع را بعبارات غلّيه
ميالائى آخر نشنيدى كه در دو سنه آخر در عراق چه واقع شد و چه وارد گشت
بشأنيكه اكثر ناس و جميع علماء بمعارضه برخواستند خود انصاف ده در همچه
حالتى اميد جان از براى نفسى باقى تا فكر رياست نمايد فعل حقّ بر امر حقّ
برهان و گواه و شاهد بوده و خواهد بود

\end{otherlanguage}

Va dar baʿżī ayyām chenān madhkūr mīshod ke maʾmūr shode-and īn ṭāʾefe
rā gerefte be-arż-e Ṭā beferestand. Makhṣūṣ Jamāl-e Abhā amr farmūd ke
koll az ʿErāq beravand va īn ḥokm-e moḥkam rā be-shaʾnī jārī farmūdand
ke Zayn al-ʿĀbedīn Khān-e maʿrūf āmade bayna yaday ḥāżer va tavassoṭ az
Āqā Seyyed Ḥosayn-e Qannād namūd ke khorūj-e ū mowqūf shavad va edhn-e
tavaqqof befarmāyand. Abadan edhn dāde nashod. Ū ham az ʿErāq khārej
shod. Farmūdand: Dūst mīdāram ke balāyā-ye koll rā vaḥdeh ḥāmel shavam
va be-īn ʿabd vāred namāyand ānche rā erāde namūde-and va aḥebbā-ye
elāhī maḥfūẓ mānand. Va to moftarī be-chenīn kalemāt Shams-e Taqdīs va
Enqeṭāʿ rā be-ʿebārāt-e ghelliyye mīālāʾī. Ākhar nashenīdī ke dar do
sane-ye ākhar dar ʿErāq che vāqeʿ shod va che vāred gasht, be-shaʾnī ke
akthar-e nās va jamīʿ-e ʿolamā be-moʿāreże barkhāstand? Khod enṣāf deh,
dar hamcho ḥālatī omīd-e jān az barā-ye nafsī bāqī tā fekr-e reyāsat
namāyad? Feʿl-e Ḥaqq bar amr-e Ḥaqq borhān va govāh va shāhed būde va
khāhad būd.

At certain times it was reported that orders had been issued to arrest
this community and send them to the Land of Ṭá. The Abhá Beauty
expressly commanded everyone to leave Iraq. He enforced this firm decree
so strictly that when the well-known Zaynu'l-ʿÁbidín Khán came into His
presence and interceded for Áqá Siyyid Ḥusayn the confectioner, asking
that his departure be suspended and that he be permitted to remain, no
permission whatsoever was given. He too left Iraq. He said: ``I wish to
bear everyone's afflictions alone, so that they may inflict on this
servant whatever they have intended, while the loved ones of God remain
safe.'' Yet you, a fabricator, seek with such words to sully the Sun of
Holiness and Detachment with expressions of malice. Have you not heard
what happened and what befell Him during the last two years in Iraq,
when most of the people and all the religious scholars arose in
opposition? Judge fairly for yourself: in such circumstances, would
anyone retain enough hope of surviving to think of leadership? The deeds
of the True One have been, and will remain, a proof, testimony, and
witness to His Cause.

\subsection{598}\label{section-1093}

\begin{otherlanguage}{persian}

نفسى را كه از جميع من فى الأرض منقطع و ظاهراً باهراً چون شمس ما بين اعدا
مشرق و مضئ و از كلّ اطراف سهام بغضا بر جسد مباركش وارد مع‌ذلك شرم ننموده‌و
نوشته آنچه را كه حقايق كلّ شیء نوحه نموده‌اند و از اينها گذشته بلاياى جمال
احديّه را در عراق نشنيدى در ارض طاء چه ميگوئى اى بی‌خبر و بی‌‌بصر همين ميرزا
يحيى با سليمان خان در دزاشوب بودند و بعد آمده بارض طاء در خانه سليمان
خان چند روز مستور بود و بعد كه اسباب فساد جمع شد فرار نمود تا آنكه واقع
شد آنچه واقع شد و حضرت ابهى واضحا مشهودا در مقرّ خود ساكن چنانچه آمدند و
دستگير نمودند با آنكه و اللّه الّذى لا اله الّا هو كه در همان ایّام از ارض
طفّ مراجعت فرموده بودند و ابدا اطّلاع نداشتند چهار ماه متوالى در اغلال و
سلاسل بوده و ضرّ آن ایّام احتياج ذكر ندارد كه چه وارد مينمودند و

\end{otherlanguage}

Nafsī rā ke az jamīʿ-e man fī al-arḍ monqaṭeʿ va ẓāheran bāheran chon
shams mā beyn-e aʿdā moshreq va moḍīʾ va az koll-e aṭrāf sehām-e baghḍā
bar jasad-e mobārakash vāred maʿa-dhālek sharm nanemūde va neveshte
ānche rā ke ḥaqāyeq-e koll-e shayʾ nowḥe nemūde-and va az īnhā goẕashte
balāyā-ye Jamāl-e Aḥadiyye rā dar ʿErāq nashenīdī dar arḍ-e Ṭāʾ che
mīgūʾī ey bī-khabar va bī-baṣar hamīn Mīrzā Yaḥyā bā Solaymān Khān dar
Dezāshūb būdand va baʿd āmade be-arḍ-e Ṭāʾ dar khāne-ye Solaymān Khān
chand rūz mastūr būd va baʿd ke asbāb-e fasād jamʿ shod farār nemūd tā
ānke vāqeʿ shod ānche vāqeʿ shod va Ḥaẓrat-e Abhā vāḍeḥan mashhūdan dar
maqarr-e khod sāken chonānche āmadand va dastgīr nemūdand bā ānke
wa-Allāh alladhī lā ilāha illā huwa ke dar hamān ayyām az arḍ-e Ṭaff
morājaʿat farmūde būdand va abadan eṭṭelāʿ nadāshtand chahār māh-e
motavālī dar aghlāl va salāsel būde va ḍorr-e ān ayyām eḥtiyāj-e dhekr
nadārad ke che vāred mīnemūdand va

Yet you felt no shame in writing what you wrote about One who was
detached from all who are on earth, who shone visibly and brilliantly
like the sun among His enemies, and whose blessed body was struck from
every side by the arrows of hatred---words over which the realities of
all things have lamented. Besides all this, have you not heard of the
afflictions of the Beauty of Oneness in Iraq? What do you say of the
Land of Ṭāʾ, you who are unaware and blind? This same Mīrzā Yaḥyā was
with Sulaymān Khān in Dizāshūb. He then came to the Land of Ṭāʾ and
remained concealed in Sulaymān Khān's house for a few days. Once the
means of sedition had been assembled, he fled, until what happened came
to pass. His Holiness the Most Glorious remained openly and visibly in
His own residence, so that they came and arrested Him, although, by God,
besides whom there is no other god, He had only just returned from the
Land of Ṭaff in those very days and had no knowledge of it at all. He
spent four consecutive months in chains and fetters. The harm inflicted
in those days needs no recounting---what they inflicted---and

\subsection{599}\label{section-1094}

\begin{otherlanguage}{persian}

هم‌چنين قبل از آن در مازندران كه همين ميرزا يحيى همراه بود كه ايشان را
اسير نموده‌شهر بشهر و ديار بديار گرداندند و روزى كه جميع علماء و ساير ناس
بر رجم جمال احديّه مجتمع شدند و آن جمال مبارك را در محضر ظالمين حاضر
نمودند وقع ما وقع يستحى القلم ان يذكر ما ورد على مظهر نفس اللّه و بعد
اراده نمودند كه ميرزا يحيى را تأديب نمايند فرمودند او با من بود اگر
تقصيرى شده از من بوده و باين جهت باو تعرّض ننموده ‌و بآن جمال قدم وارد
آوردند آنچه وارد آوردند و هم‌چنين توسّط از ساير همراهان فرمودند از جمله
جناب ملّا باقر حرف حىّ عليه بهآء اللّه كه الآن موجود است جويا شويد

\end{otherlanguage}

Ham-chenīn qabl az ān dar Māzandarān ke hamīn Mīrzā Yaḥyā hamrāh būd ke
īshān rā asīr nemūde shahr be-shahr va diyār be-diyār gardāndand va rūzī
ke jamīʿ-e ʿolamā va sāyer-e nās bar rajm-e Jamāl-e Aḥadiyye mojtameʿ
shodand va ān Jamāl-e Mobārak rā dar maḥẓar-e ẓālemīn ḥāẓer nemūdand
waqaʿa mā waqaʿa yastaḥyī al-qalamu an yadhkura mā warada ʿalā maẓhari
nafsi Allāh va baʿd erāde nemūdand ke Mīrzā Yaḥyā rā taʾdīb namāyand
farmūdand ū bā man būd agar taqṣīrī shode az man būde va be-īn jehat
be-ū taʿarroẓ nanemūde va be-ān Jamāl-e Qedam vāred āvardand ānche vāred
āvardand va ham-chenīn tavassoṭ az sāyer-e hamrāhān farmūdand az jomle
Janāb-e Mollā Bāqer Ḥarf-e Ḥayy ʿalayhi bahāʾu Allāh ke al-ān mowjūd ast
jūyā shavīd

Likewise, before that, in Mazandaran, when this same Mīrzā Yaḥyā was
with Him, they took Him captive and led Him from city to city and region
to region. On the day when all the religious scholars and the rest of
the people gathered to stone the Beauty of Oneness, and brought that
blessed Beauty before the oppressors, what happened came to pass. The
Pen is ashamed to recount what befell the Manifestation of God's own
Self. Afterward, they intended to punish Mīrzā Yaḥyā. He said, ``He was
with Me. If any fault has been committed, it was Mine.'' For this reason
they did not molest him, but inflicted on that Ancient Beauty what they
inflicted. He likewise interceded for His other companions, among them
the honored Mullā Bāqir, the Letter of the Living---may the glory of God
rest upon him---who is still alive. Ask him.

\subsection{600}\label{section-1095}

\begin{otherlanguage}{persian}

و از اين مراتب هم گذشته در نيالا كه دهى از دهات دار المرز است بغتةً خمسة
مأة نفس او ازيد بر آن جمال احديّه هجوم آوردند و وارد شد در آن ارض آنچه
وارد شد اگر انصاف دهى از اوّل ظهور بديع تا حين ابداً آسایش ندیده‌اند و
مع‌ذلك نوشته آنچه را كه نوشته اين امور چرا واقع شده و از براى كه واقع شده
شايد اين امور را هم انكار نمائى چه كه از آن نفوس عجيب نيست

\end{otherlanguage}

Va az īn marāteb ham goẕashte dar Niyālā ke dehī az dehāt-e Dār al-Marz
ast baghtatan khamsata miʾati nafsin aw azyad bar ān Jamāl-e Aḥadiyye
hojūm āvardand va vāred shod dar ān arḍ ānche vāred shod agar enṣāf dehī
az avval-e ẓohūr-e badīʿ tā ḥīn abadan āsāyesh nadīde-and va maʿa-dhālek
neveshte ānche rā ke neveshte īn omūr cherā vāqeʿ shode va az barā-ye ke
vāqeʿ shode shāyad īn omūr rā ham enkār namāʾī che ke az ān nofūs ʿajīb
nīst

Beyond all this, in Niyālā, a village among the villages of Dār al-Marz,
five hundred people or more suddenly attacked that Beauty of Oneness,
and what befell Him in that land befell Him. If you judge fairly, from
the beginning of this wondrous Revelation until now He has never known
ease. Yet you have written what you have written. Why did these things
happen, and for whose sake did they happen? Perhaps you will deny these
things as well, for that would be no surprise from such people.

\subsection{601}\label{section-1096}

\begin{otherlanguage}{persian}

حال از معرض باللّه مرشدت بگو كه در اين امر بر او چه وارد شده مسلّم است كه
لازال باكل و شرب و تصرّف در ابكار و نساء ناس بوده و اعماليكه و اللّه خجلت
ميكشم از ذكرش مرتكب بعد از فسادى كه در ارض طاء نمود منهزم گشته به عراق
آمده و در بيت ايشان وارد شد و بقسمى خائف كه با وجود حضرت ابهى كه جميع
ميدانند كه كلّ ناس در صدد آن وجود مبارك بودند مع‌ذلك فو‌الّذى نفسى بيده
بمجرّد يك خبر كه از ايران ميرسيد از عراق باطراف ميرفت و متّصل هم بعيش و
عشرت نفسانى مشغول مع‌ذلك آنچه مشهود است انكار نموده ‌و نوشته كه از كجا
معلوم شد كه اين قيام للّه بوده‌

\end{otherlanguage}

Ḥāl az moʿreẓ be-Allāh morshedat begū ke dar īn amr bar ū che vāred
shode mosallam ast ke lā-zāl be-akl va shorb va taṣarrof dar abkār va
nesā-ye nās būde va aʿmālī ke va-Allāh khejlat mīkasham az dhekrash
mortakeb baʿd az fasādī ke dar arḍ-e Ṭāʾ nemūd monhazem gashte be ʿErāq
āmade va dar beyt-e īshān vāred shod va be-qesmī khāʾef ke bā vojūd-e
Ḥaẓrat-e Abhā ke jamīʿ mīdānand ke koll-e nās dar ṣadad-e ān Vojūd-e
Mobārak būdand maʿa-dhālek fa-wa-alladhī nafsī bi-yadihi be-mojarrad-e
yek khabar ke az Īrān mīresīd az ʿErāq be-aṭrāf mīraft va mottaṣel ham
be-ʿeysh va ʿeshrat-e nafsānī mashghūl maʿa-dhālek ānche mashhūd ast
enkār nemūde va neveshte ke az kojā maʿlūm shod ke īn qiyām li-llāh būde

Now tell us about your guide, who has turned away from God: what has
befallen him in this Cause? It is certain that he has continually
occupied himself with eating and drinking, taking possession of other
people's virgins and wives, and committing deeds that, by God, I am
ashamed to mention. After the sedition he stirred up in the Land of Ṭāʾ,
he fled to Iraq and entered His house. He was so fearful that, despite
the presence of His Holiness the Most Glorious---whom everyone knows all
the people were seeking to harm---by Him in whose hand is my soul,
whenever a single report arrived from Iran, he would leave Iraq for the
surrounding regions. All the while he was occupied with sensual pleasure
and enjoyment. Yet you have denied what is plain to see and written,
``How was it established that this uprising was for God's sake?''

\subsection{602}\label{section-1097}

\begin{otherlanguage}{persian}

اى مغلّ معرض از كجا بر تو معلوم شد كه لغير اللّه بوده هرگز نفس مقدّس را در
آنى حفظ نفرموده‌اند در اينمقام يك فقره از مناجات كه مع اللّه نموده‌اند ذكر
ميشود لعلّ تتذكّر فى نفسك و تكون من المتذكّرين مضمون آن بلسان پارسى اينست
اى اله من هرگز نفس خود را در سبيل تو حفظ ننموده‌ام و با اعدايت مداهنه
ننمودم در كلّ اوان در دست اعدا مبتلى بوده‌ام و در احيان قضاياى وارده اراده
چنان بود كه جان انفاق نمايم و لكن تو اى پروردگار من بحكم بالغه و مصالح
مستوره از انظر بريّه عبد خود را حفظ فرمودى و لكن اين حفظ نه بجهت آسایش او
بوده بلكه لاجل قضاى آخر و توئى محمود در كلّ افعال خود و محبوب در اراده
خود حال اگر امثال تو در كلمات منزله الهيّه نظر مينمودند مطّلع ميشدند
بامورى كه احدى بر آن مطّلع نشده و لكن قضى عنك و عنهم

\end{otherlanguage}

Ey moghell-e moʿreẓ az kojā bar to maʿlūm shod ke li-ghayri Allāh būde
hargez Nafs-e Moqaddas rā dar ānī ḥefẓ nafarmūde-and dar īn maqām yek
faqre az monājāt ke maʿa Allāh nemūde-and dhekr mīshavad laʿallaka
tatadhakkaru fī nafsika wa-takūnu mina al-mutad h akkirīn maẓmūn-e ān
be-lesān-e pārsī īnast ey Elāh-e man hargez nafs-e khod rā dar sabīl-e
to ḥefẓ nanemūde-am va bā aʿdāyat modāhene nanemūdam dar koll-e avān dar
dast-e aʿdā mobtalā būde-am va dar aḥyān-e qaẓāyā-ye vārede erāde chonān
būd ke jān enfāq namāyam va lāken to ey Parvardegār-e man be-ḥekam-e
bāleghe va maṣāleḥ-e mastūre az anẓor-e bariyye ʿabd-e khod rā ḥefẓ
farmūdī va lāken īn ḥefẓ na be-jehat-e āsāyesh-e ū būde balke li-ajl-e
qaḍā-ye ākhar va toʾī maḥmūd dar koll-e afʿāl-e khod va maḥbūb dar
erāde-ye khod ḥāl agar amthāl-e to dar kalemāt-e monzale-ye elāhiyye
naẓar mīnemūdand moṭṭaleʿ mīshodand be-omūrī ke aḥadī bar ān moṭṭaleʿ
nashode va lāken quḍiya ʿanka wa-ʿanhum

You malicious denier! How did you establish that it was for anyone other
than God? He has never, for a moment, sought to safeguard His sacred
Self. Here a passage from His supplication to God is recounted, so that
you may reflect within yourself and be among those who take heed. Its
meaning in Persian is this: ``O my God, I have never sought to safeguard
Myself in Your path, nor have I compromised with Your enemies. At all
times I have suffered in the hands of enemies, and amid the events that
befell Me, My intention was to offer up My life. But You, O my Lord,
through consummate wisdom and purposes hidden from the eyes of Your
creatures, have preserved Your servant. Yet this preservation was not
for His ease, but for another decree. You are praised in all Your
actions and beloved in Your will.'' Now, if people like you had
considered the divinely revealed words, they would have become aware of
things of which no one has become aware. But it has been decreed
concerning you and them.

\subsection{603}\label{section-1098}

\begin{otherlanguage}{persian}

و اينكه نوشته بودى نصرت امر اللّه كسانى نمودند كه در مازندران چرم خوردند
فرمودند در ایّاميكه در ارض طاء محبوس بوديم چند يوم اوّل چرم هم نبود كه كسى
بخورد فوا حسرتا عليك يا ايّها الغافل عن كلّ خير و العامل بكلّ شرّ تنطق بما
لا علمته و هذا لم يكن من عندك بل بما القى الشّيطان فى صدرك نسئل اللّه بأن
يأخذه اخذ عزيز مقتدر جميع ذرّات شهادت ميدهند كه جمال ابهى وحده در مقابل
اعدا قائم بودند چنانچه بعضى از احبّا از ارض طفّ در آن ایّام نوشتند و استدعا
نمودند كه چند روزى مستور شوند چه كه امر بقسمى صعب شده بود كه هر نفسى در
كلّ يوم بين يدى وارد ميشد خبرى جديد از مكر ماكرين و اراده معرضين ذكر
مينمود و بشأنى اعدا باعراض برخواستند كه در كوچه ها و اسواق بعضى را معين
نموده‌كه در حين مرور بضرب رصاص شهيد نمايند مع‌ذلك آنى وجه مبارك را ستر
نفرموده و وحده از بيت بيرون تشريف مياوردند و ابدا اعتنا بنفسى نميفرمودند
و هميشه وا شوقا للقآء المحبوب ميفرمودند اهل ملأ اعلى از فعلشان متحيّر و
جميع بتبارك اللّه احسن المبدعين ناطق و تو جميع را انكار نموده‌چه كه بانفس
فجّار مأنوس گشته و از حقّ تبرّى جسته

\end{otherlanguage}

Va īnke neveshte būdī noṣrat-e Amr Allāh kasānī nemūdand ke dar
Māzandarān charm khordand farmūdand dar ayyāmī ke dar arḍ-e Ṭāʾ maḥbūs
būdīm chand yowm-e avval charm ham nabūd ke kasī bekhorad fa-wā ḥasratā
ʿalayka yā ayyuhā al-ghāfilu ʿan kulli khayrin wa-al-ʿāmilu bi-kulli
sharrin tanṭiqu bi-mā lā ʿalimtahu wa-hādhā lam yakun min ʿindika bal
bi-mā alqā al-shayṭānu fī ṣadrika nasʾalu Allāha bi-an yaʾkhudhahu
akhdha ʿazīzin muqtadir jamīʿ-e dharrāt shahādat mīdehand ke Jamāl-e
Abhā vaḥdahu dar moqābel-e aʿdā qāʾem būdand chonānche baʿẓī az aḥebbā
az arḍ-e Ṭaff dar ān ayyām neveshtand va estedʿā nemūdand ke chand rūzī
mastūr shavand che ke amr be-qesmī ṣaʿb shode būd ke har nafsī dar
koll-e yowm beyn-e yaday vāred mīshod khabarī jadīd az makr-e mākerīn va
erāde-ye moʿreẓīn dhekr mīnemūd va be-shaʾnī aʿdā be-eʿrāẓ barkhāstand
ke dar kūchehā va asvāq baʿẓī rā moʿayyan nemūde ke dar ḥīn-e morūr
be-ẓarb-e raṣāṣ shahīd namāyand maʿa-dhālek ānī vajh-e mobārak rā satr
nafarmūde va vaḥdahu az beyt bīrūn tashrīf mīāvardand va abadan eʿtenā
be-nafsī nemīfarmūdand va hamīshe wā shawqā li-liqāʾi al-Maḥbūb
mīfarmūdand ahl-e Malaʾ-e Aʿlā az feʿlshān motaḥayyer va jamīʿ
be-tabāraka Allāhu aḥsanu al-mubdiʿīn nāṭeq va to jamīʿ rā enkār nemūde
che ke be-anfos-e fojjār maʾnūs gashte va az Ḥaqq tabarrī joste

As for your statement that those who aided the Cause of God were the
ones who ate leather in Mazandaran, He said, ``During the days when We
were imprisoned in the Land of Ṭāʾ, for the first few days there was not
even leather for anyone to eat.'' Alas for you, you who are heedless of
every good and commit every evil! You speak of what you do not know.
This did not originate with you, but with what Satan cast into your
breast. We ask God to seize him with the grip of One who is mighty and
all-powerful. Every atom bears witness that the Most Glorious Beauty
stood alone against His enemies. Indeed, in those days some of the loved
ones wrote from the Land of Ṭaff, imploring Him to remain concealed for
a few days, for matters had become so grave that everyone who entered
His presence each day brought fresh news of the schemers' plots and the
deniers' intentions. The enemies rose in such opposition that they
stationed men in the lanes and markets to shoot Him dead as He passed.
Nevertheless, He did not conceal His blessed face for a moment, but went
out of the house alone, paying no heed to anyone. He would always say,
``Oh, how I long to meet the Beloved!'' The dwellers of the heavenly
Assembly marveled at His conduct, and all proclaimed, ``Blessed is God,
the best of originators!'' Yet you have denied it all, for you have
befriended wicked souls and disowned the True One.

\subsection{604}\label{section-1099}

\begin{otherlanguage}{persian}

بلى چيزى که فرق بوده اين بوده كه اين ظهور ابدع ناس را از قتل و نزاع و
جدال منع فرمودند چنانچه در سنين توقّف در عراق ابدا در محلّ محاربه واقع نشد
و تصريحا در آيات نازل كه نصرت و تبليغ امر بمواعظ حسنه و بيانات شافيه
بوده و خواهد بود نه بمجادله و محاربه و بشأنى اين حكم محكم نازل شده كه در
بعضى از كلمات الهى مذكور كه

\end{otherlanguage}

Balī chīzī ke farq būde īn būde ke īn Ẓohūr-e Abdaʿ nās rā az qatl va
nezāʿ va jedāl manʿ farmūdand chonānche dar senīn-e tavaqqof dar ʿErāq
abadan dar maḥall-e moḥārebe vāqeʿ nashod va taṣrīḥan dar āyāt nāzel ke
noṣrat va tablīgh-e amr be-mavāʿeẓ-e ḥasane va bayānāt-e shāfiye būde va
khvāhad būd na be-mojādele va moḥārebe va be-shaʾnī īn ḥokm-e moḥkam
nāzel shode ke dar baʿẓī az kalemāt-e elāhī madhkūr ke

Yes, the difference was that this most wondrous Revelation forbade
people to kill, dispute, or fight. Thus, throughout the years of His
stay in Iraq, no occasion for battle arose. It was expressly revealed in
the verses that the Cause has been, and will be, aided and taught
through good counsel and healing words, not through contention and
warfare. So emphatically was this binding command revealed that in some
of the divine words it is stated that

\subsection{605}\label{section-1100}

\begin{otherlanguage}{persian}

اگر نفسى بابناء و يا ذوى قرابه من ضرّى وارد آورد نزد اين عبد احسن و
محبوبتر است از اينكه متعرّض نفسى شود

\end{otherlanguage}

Agar nafsī be-abnā va yā dhavī qarābe-ye man ḍorrī vāred āvarad nazd-e
īn ʿabd aḥsan va maḥbūbtar ast az īnke motaʿarreẓ-e nafsī shavad

``If someone were to inflict harm on My children or My relatives, it
would be better and more pleasing in the sight of this servant than for
him to assail another soul.''

\subsection{606}\label{section-1101}

\begin{otherlanguage}{persian}

فانظر فيما نزل فى الألواح لعلّ تنصف فى نفسك و تكون من المنصفين و لذا در
اينمدّت امرى واقع نشد مگر اينكه چند نفر از اين طايفه خود بامر مرشدت
يكديگر را كشتند و ديگر باحدى متعرّض نشدند و در جميع الواح نازل كه مقصود
اين ظهور اتّحاد و اتّفاق بوده نه افتراق و اختلاف

\end{otherlanguage}

Fa-nẓur fī-mā nazala fī al-alwāḥi laʿallaka tunṣifu fī nafsika wa-takūnu
mina al-munṣifīn va le-dhā dar īn moddat amrī vāqeʿ nashod magar īnke
chand nafar az īn ṭāyefe khod be-amr-e morshedat yekdīgar rā koshtand va
dīgar be-aḥadī motaʿarreẓ nashodand va dar jamīʿ-e alvāḥ nāzel ke
maqṣūd-e īn ẓohūr etteḥād va ettefāq būde na efterāq va ekhtelāf

Consider, then, what has been revealed in the Tablets, so that you may
judge fairly within yourself and be among the just. Accordingly,
throughout this period nothing happened except that a few members of
this community killed one another at your guide's command. They assailed
no one else. In all the Tablets it has been revealed that the purpose of
this Revelation is unity and concord, not division and discord.

\subsection{607}\label{section-1102}

\begin{otherlanguage}{persian}

و اينكه نوشته بودى كه نه اشخاصيكه در بغداد و جاى ديگر در ایّام و ليالى در
صرف چاى با سماور و اوضاع و دستگاه الخ اى ارذل وجود و اللّه الّذى لا اله
الّا هو كه اگر اقلّ از ذرّه نجابت در تو و مرشدينت بود هرگز چنين كلمات مذكور
نمينمودى و نفس اللّه را بجهت سرور انفس خبيثه نميازردى تا حال هيچ رذلى
چنين كلمات ذكر نموده‌كه تو نموده ‌لا فو‌ربّ العالمين الّا انت و امثالك اسمع
ندآء نقطة البيان ثمّ استحى عن جماله و لاتضرب محبوبه بسيوف البغضآء و لا
مقصوده برماح الفحشا اتّق من يوم تضطرب فيه كلّ الوجود من الغيب و الشّهود و
تقشعر فيه جلود العالمين

\end{otherlanguage}

Va īnke neveshte būdī ke na ashkhāṣī ke dar Baghdād va jā-ye dīgar dar
ayyām va layālī dar ṣarf-e chāy bā samāvar va owẓāʿ va dastgāh ilā ākhir
ey ardhal-e vojūd va-Allāh alladhī lā ilāha illā huwa ke agar aqall az
dharre najābat dar to va morshedīnat būd hargez chonīn kalemāt madhkūr
nemīnemūdī va Nafs Allāh rā be-jehat-e sorūr-e anfos-e khabīthe
nemīāzordī tā ḥāl hīch radhlī chonīn kalemāt dhekr nemūde ke to nemūde
lā fa-wa-Rabbi al-ʿālamīna illā anta wa-amthāluka ismaʿ nidāʾa Nuqṭati
al-Bayān thumma istaḥyi ʿan jamālihi wa-lā taḍrib maḥbūbahu bi-suyūfi
al-baghḍāʾi wa-lā maqṣūdahu bi-rimāḥi al-faḥshā ittaqi min yawmin
taḍṭaribu fīhi kullu al-wujūdi mina al-ghaybi wa-al-shuhūdi
wa-taqshaʿirru fīhi julūdu al-ʿālamīn

As for what you wrote, ``Not those people who, in Baghdad and elsewhere,
spend their days and nights drinking tea with samovars and all their
arrangements and equipment,'' and so forth: you most abject of beings!
By God, besides whom there is no other god, if you and your guides
possessed less than an atom of decency, you would never have uttered
such words or hurt God's own Self to delight wicked souls. Has any vile
person ever uttered words like those you have uttered? No, by the Lord
of the worlds, none but you and those like you! Hear the call of the
Point of the Bayān, then feel shame before His beauty. Do not strike His
Beloved with the swords of hatred, nor the One He seeks with the spears
of indecency. Beware of a day when all existence, unseen and seen, will
tremble, and the skins of all the world's peoples will shudder.

\subsection{608}\label{section-1103}

\begin{otherlanguage}{persian}

قوله عزّ اعزازه و جلّ كبريائه ثمّ فى ليلة من الاء اللّه تسعة عشر عدّة بين
ايديكم لتحصون الى عدد المستغاث اذن لمن يقدر و لاتحزنن اذا انتم
لاتستطيعون انتهى ميفرمايد تا عدد مستغاث اگر مستطيع باشيد در يك مجلس حاضر
نمائيد اذن داده شده‌ايد و مع‌ذلك تو نفهميده اعتراضاً على اللّه ذكر چاى و
سماور و اطعمه و شيرينى و مهمانيها نموده ‌مع آنكه امر نقطه بيانست كه
ميفرمايند قوله عزّ ذكره فلتضيفن فى تسعة عشر يوماً تسعة عشر نفساً ولو انتم
ماء الواحد لتؤتون و هم‌چنين ميفرمايند كه فرض است بر هر نفسى كه عدد واحد
از بلّور دربيت خود موجود نمايد و هم‌چنين ميفرمايد قوله عزّ ذكره و انتم
اسبابكم الّتى بها فى سرّكم لتعيشون من الذّهب و الفضّة تصنعون مع‌ذلك تو پست
فطرت و ارذل اعتراض بچاى و سماور نموده‌ جميع اين بيانات را نقطه اولى ذكر
نفرموده مگر آنكه كسى تمسّك نجويد كه خبز شعير خوردم تا برتبه اعلى فايز شدم

\end{otherlanguage}

Qawluhu ʿazza iʿzāzuhu wa-jalla kibriyāʾuhu thumma fī laylatin min ālāʾi
Allāhi tisʿata ʿashara ʿiddatan bayna aydīkum li-tuḥṣūna ilā ʿadadi
al-Mustaghāth udhina li-man yaqdiru wa-lā taḥzanunna idhā antum lā
tastaṭīʿūna intahā mīfarmāyad tā ʿadad-e Mostaghāth agar mostaṭīʿ bāshīd
dar yek majles ḥāẓer namāyīd edhn dāde shode-īd va maʿa-dhālek to
nafahmīde iʿtirāḍan ʿalā Allāh dhekr-e chāy va samāvar va aṭʿeme va
shīrīnī va mehmānīhā nemūde maʿa ānke amr-e Noqṭe-ye Bayān ast ke
mīfarmāyand qawluhu ʿazza dhikruhu fa-latuḍīfunna fī tisʿata ʿashara
yawman tisʿata ʿashara nafsan wa-law antum māʾa al-wāḥidi la-tuʾtūn va
ham-chenīn mīfarmāyand ke farẓ ast bar har nafsī ke ʿadad-e vāḥed az
bollūr dar beyt-e khod mowjūd namāyad va ham-chenīn mīfarmāyad qawluhu
ʿazza dhikruhu wa-antum asbābakumu allatī bi-hā fī sirrikum la-taʿīshūna
mina al-dhahabi wa-al-fiḍḍati taṣnaʿūn maʿa-dhālek to past-feṭrat va
ardhal eʿterāẓ be-chāy va samāvar nemūde jamīʿ-e īn bayānāt rā Noqṭe-ye
Ūlā dhekr nafarmūde magar ānke kasī tamassok najūyad ke khobz-e shaʿīr
khordam tā be-rotbe-ye aʿlā fāʾez shodam

His words---exalted be His honor and great His majesty---are: ``Then, on
a night, among the favors of God, let nineteen be numbered before you,
that you may count up to the number of al-Mustaghāth. Permission is
given to whoever is able; do not grieve if you are unable.'' End of
quotation. He says that you have been given permission, if you are able,
to bring together in one gathering as many as the number of
al-Mustaghāth. Yet, without understanding, you have mentioned tea,
samovars, food, sweets, and hospitality as an objection against God,
although it is the command of the Point of the Bayān, who says---exalted
be His remembrance---``You must entertain nineteen souls every nineteen
days, even if you can offer only water to each one.'' He likewise says
that every soul is required to have in his house a number of crystal
vessels equal to one vāḥid. He also says---exalted be His
remembrance---``You shall make of gold and silver the furnishings by
which you live in your private life.'' Yet you, base-natured and most
abject, have objected to tea and samovars. The Primal Point uttered all
these statements only so that no one would claim, ``I ate barley bread
and thereby attained the highest rank.''

\subsection{609}\label{section-1104}

\begin{otherlanguage}{persian}

بلى آنقدر هست كه انسان بايد بقليل كفايت كند و شاكر باشد و بكثير از حقّ
غافل نشود اصل امر اينست و الّا رياضت و عدم آن علامت حقّ و باطل نبوده و
علامت احبّاى حقّ از جبينشان ظاهر و مشهود چنانچه اليوم بخاتم عقيق معانى بر
جبين احبّاى الهى مرقوم كه هذا من اهل اللّه بين العالمين و واللّه مقصودى از
اين بيان نداشته‌اند مگر آنكه شايد بعنايت الهى عباد از اين حدودات فارغ
شوند فرمودند اگر ممكن ميشد اليوم امر مينمودم كه جميع اهل اللّه بر اكراس
ذهب جالس شوند چه كه آنچه خلق شده در حقيقت اوّليّه مخصوص مؤمنين باللّه بوده
كجا است اين نظر و آن نظرى كه در چاى خوردن اعتراض مينمايد مع آنكه از مذهب
بيان است

\end{otherlanguage}

Balī ānqadr hast keh ensān bāyad be-qalīl kefāyat konad va shāker bāshad
va be-kathīr az Ḥaqq ghāfel nashavad. Aṣl-e amr īnast va ellā riyāḍat va
ʿadam-e ān ʿalāmat-e ḥaqq va bāṭel nabūdeh va ʿalāmat-e aḥebbā-ye Ḥaqq
az jabīneshān ẓāher va mashhūd, chenāncheh al-yawm be-khātam-e ʿaqīq-e
maʿānī bar jabīn-e aḥebbā-ye elāhī marqūm keh hādhā min ahl Allāh bayna
al-ʿālamīn. Va va-llāh maqṣūdī az īn bayān nadāshteh-and magar ānkeh
shāyad be-ʿenāyat-e elāhī ʿebād az īn ḥodūdāt fāregh shavand. Farmūdand
agar momken mīshod al-yawm amr mīnamūdam keh jamīʿ-e ahl Allāh bar
akrās-e dhahab jāles shavand, cheh keh āncheh khalq shodeh dar ḥaqīqat-e
avvalīyeh makhṣūṣ-e moʾmenīn be-Allāh būdeh. Kojā ast īn naẓar va ān
naẓarī keh dar chāy khordan eʿterāḍ mīnamāyad maʿ-e ānkeh az madhhab-e
Bayān ast.

Yes, this much is required: a person should be content with little and
be thankful, and should not let abundance make him heedless of God. This
is the essence of the matter; otherwise, neither asceticism nor its
absence has been a sign of truth or falsehood. The mark of God's loved
ones is plain and visible on their foreheads, just as today the agate
seal of inner meanings has inscribed on the foreheads of God's loved
ones: ``This is one of the people of God among the worlds.'' By God, His
only purpose in this utterance was that, through divine favor, His
servants might perhaps be freed from these restrictions. He said, ``If
it were possible, today I would command all the people of God to sit on
seats of gold,'' for everything that has been created was, in its
primary reality, intended especially for those who believe in God. How
far removed is this outlook from that of someone who objects to drinking
tea, even though he belongs to the religion of the Bayān!

\subsection{610}\label{section-1105}

\begin{otherlanguage}{persian}

اى اهل بيان تفكّر در اعتراضات معرضين نمائيد كه چه وارد نموده‌اند تفكّروا
لتكوننّ فيهم لمن المتبصّرين چه زود جميع شئونات الهيّه را تبديل نموده‌اند چه
كه مخصوص حضرت اعلى روح ماسواه فداه در كلّ يوم چاى ميل ميفرمودند و جميع
احبّا را امر فرموده‌اند بشرب چاى و ميفرمايند برياضت باين مقام نرسيده‌ام
چنانچه جمال ابهى ميفرمايند هرگز رياضت نكشيده‌ام و از اهل علم نبوده‌ام بلكه
اين عنايتى كه ظاهر شده از فضل پروردگار بوده من غير استحقاقى به و لكن
باضطرار بسيار از ليالى كه از نعم الهى بين يدى هيچ شیء نبوده حتّى نان و
هم‌چنين دو سنه متواليه بحسب ظاهر هيچ از اسباب تعيّش نبوده و لكن در كلّ حال
شاكر بوده‌ايم و راضى خواهيم بود

\end{otherlanguage}

Ey ahl-e Bayān, tafakkor dar eʿterāḍāt-e moʿreḍīn namāyīd keh cheh vāred
namūdeh-and. Tafakkarū la-takūnunna fīhim la-min al-mutabaṣṣirīn. Cheh
zūd jamīʿ-e shoʾūnāt-e elāhīyeh rā tabdīl namūdeh-and, cheh keh makhṣūṣ
Ḥaḍrat-e Aʿlā, rūḥ-e mā-sevāh fedāh, dar koll-e yawm chāy meyl
mīfarmūdand va jamīʿ-e aḥebbā rā amr farmūdeh-and be-shorb-e chāy va
mīfarmāyand be-riyāḍat be-īn maqām narasīdeh-am. Chenāncheh Jamāl-e Abhā
mīfarmāyand hargaz riyāḍat nakeshīdeh-am va az ahl-e ʿelm nabūdeh-am,
balkeh īn ʿenāyatī keh ẓāher shodeh az faḍl-e Parvardegār būdeh min
ghayr istiḥqāqī bihi. Va lāken be-eḍṭerār besyār az layālī keh az
neʿam-e elāhī bayna yadayya hīch shayʾ nabūdeh ḥattā nān, va ham-chenīn
do saneh-ye motavālīyeh be-ḥasab-e ẓāher hīch az asbāb-e taʿayyosh
nabūdeh, va lāken dar koll-e ḥāl shāker būdeh-īm va rāḍī khāhīm būd.

People of the Bayān! Reflect on the objections of those who turn away,
and on what they have brought forward. Reflect, so that you may be among
those who discern their ways. How quickly they have altered all the
divine observances! For the Most Exalted One Himself---may the spirit of
all else be sacrificed for Him---drank tea every day, commanded all His
loved ones to drink tea, and says, ``I have not attained this station
through asceticism.'' Likewise, the Abhā Beauty says, ``I have never
practiced asceticism, nor have I belonged to the learned. Rather, this
favor that has appeared has come from the grace of the Lord, without any
merit on My part. Yet, out of necessity, on many nights none of God's
bounties lay before Me, not even bread. Likewise, for two consecutive
years, to outward appearances there were no means of subsistence at all.
Yet in every condition We have been thankful and will remain content.''

\subsection{611}\label{section-1106}

\begin{otherlanguage}{persian}

و اليوم احبّاى الهى بايد بهيچ وجه بر خود سخت نگيرند در صورت امكان باغذيه
لطيفه متنعّم شوند و در كلّ احوال بحدّ اعتدال حركت نمايند و اگر غذا واحد
باشد احبّ است عند اللّه و لكن جهد نموده‌كه آن واحد از اغذيه لطيفه باشد على
قدر وسع و بايد اولو الغنا فقرا را فراموش ننمايند چه كه فقرا امانات
اللّه‌اند بين عباد ان انفقوا يا قوم ما رزقكم اللّه به و كونوا من الّذين
يؤثّرون على انفسهم و لو كان بهم خصاصة كذلك نزل فى البيان ثمّ فى الفرقان ثمّ
فى الانجيل و الزّبور و التّوريّة و صحف اللّه المهيمن العزيز القيّوم

\end{otherlanguage}

Va al-yawm aḥebbā-ye elāhī bāyad be-hīch vajh bar khod sakht nagīrand.
Dar ṣūrat-e emkān be-aghdhīyeh-ye laṭīfeh motanaʿʿem shavand va dar
koll-e aḥvāl be-ḥadd-e eʿtedāl ḥarakat namāyand. Va agar ghadhā vāḥed
bāshad aḥabb ast ʿinda Allāh, va lāken jahd namūdeh keh ān vāḥed az
aghdhīyeh-ye laṭīfeh bāshad ʿalā qadr-e vosʿ. Va bāyad ulu al-ghinā
foqarā rā farāmūsh nanamāyand, cheh keh foqarā amānāt Allāh-and bayn-e
ʿebād. An anfiqū yā qawm mā razaqakum Allāhu bihi wa-kūnū min alladhīna
yuʾthirūna ʿalā anfusihim wa-law kāna bihim khaṣāṣah. Kadhālika nazala
fī al-Bayān thumma fī al-Furqān thumma fī al-Injīl wa-al-Zabūr
wa-al-Tawrīyah wa-ṣuḥuf Allāh al-Muhaymin al-ʿAzīz al-Qayyūm.

Today God's loved ones should in no way impose hardship on themselves.
Where possible, they should enjoy fine foods and in all circumstances
conduct themselves with moderation. If there is only one dish, this is
more beloved in God's sight; yet they should endeavor, according to
their means, to make that one dish a fine one. Those who possess wealth
must not forget the poor, for the poor are God's trusts among His
servants. Spend, My people, from what God has provided for you, and be
among those who prefer others to themselves, even though they are in
need. Thus was it revealed in the Bayān, and in the Qur'an, and in the
Gospel, the Psalms, the Torah, and the scriptures of God, the
All-Protecting, the Almighty, the Self-Subsisting.

\subsection{612}\label{section-1107}

\begin{otherlanguage}{persian}

بارى امر اللّه مقدّس از جميع اين امور بوده و كلّ آنچه خلق شده ما بين سموات
و ارض مخصوص مظاهر امر خلق شده و لكن فو‌اللّه الّذى لا اله الّا هو كه آن نفوس
مشرقه از افق احديّه لازال خود بقليل كفايت فرموده‌اند و ناس را در اين
امورات بر انفس مبارك خود مقدّم داشته‌اند چنانچه بسا شده كه آنچه در بيت
حاضر بوده مخصوص احبّا فرستاده‌اند و آن ليل و يا يوم را اهل حرم گرسنه
مانده‌اند فللّه الحمد فى كلّ الأحوال و انّه اطّلع ما لا اطّلعتم به يا معشر
المفترين

\end{otherlanguage}

Bārī Amr Allāh moqaddas az jamīʿ-e īn omūr būdeh va koll-e āncheh khalq
shodeh mā bayn-e samāvāt va arḍ makhṣūṣ-e maẓāher-e Amr khalq shodeh. Va
lāken fa-wa-llāhi alladhī lā ilāha illā huwa keh ān nofūs-e moshreqeh az
ofoq-e aḥadīyeh lāzāl khod be-qalīl kefāyat farmūdeh-and va nās rā dar
īn omūrāt bar anfos-e mobārak-e khod moqaddam dāshteh-and, chenāncheh
basā shodeh keh āncheh dar bayt ḥāḍer būdeh makhṣūṣ-e aḥebbā
ferestādeh-and va ān layl va yā yawm rā ahl-e ḥaram gorosneh māndeh-and.
Fa-li-llāhi al-ḥamdu fī kull al-aḥwāl wa-innahu iṭṭalaʿa mā lā
iṭṭalaʿtum bihi yā maʿshar al-muftarīn.

In short, the Cause of God is sanctified above all these matters, and
everything created between the heavens and the earth has been created
especially for the Manifestations of the Cause. Yet, by God, besides
whom there is no other god, those Souls who shine from the horizon of
oneness have always contented themselves with little and, in these
matters, have given others precedence over Their own blessed selves. It
has often happened that They sent whatever was available in the house
expressly to the loved ones, while the members of Their household went
hungry that night or day. Praise be to God in every circumstance! He
knows what you do not know, you company of slanderers.

\subsection{613}\label{section-1108}

\begin{otherlanguage}{persian}

اگر كلّ من على الأرض و آنچه در او خلق شده از نعماى لطيفه و فواكه طيّبه در
يك مجلس حبّاً للّه از براى نفس مؤمن باللّه حاضر شود اسراف نبوده و نخواهد بود
ولكن تو و امثال تو اگر لقمهء نانى تصرّف نمائيد مسرف بوده و خواهد بود چه
كه اسراف در غفلت از حقّ مشهود است و همان لقمه نان كه بنفس آن معرض باللّه
معذّب شده پناه ميبرد بخدا از اينكه نصيب او گشته و واللّه كلّ ما يأكل ليَنوح
و يضجّ و يقول يا الهى باىّ جرم جعلتنى رزق هذا الفاجر الّذى كفر بك و بآياتك
ولكن تو شاعر نشده و نخواهى شد و از اين مراتب گذشته

\end{otherlanguage}

Agar koll-e man ʿalā al-arḍ va āncheh dar ū khalq shodeh az naʿmā-ye
laṭīfeh va favākeh-e ṭayyebeh dar yek majles ḥubban li-llāh az barā-ye
nafs-e moʾmen be-Allāh ḥāḍer shavad esrāf nabūdeh va nakhāhad būd. Va
lāken to va amthāl-e to agar loqmeh-ye nānī taṣarrof namāyīd mosref
būdeh va khāhad būd, cheh keh esrāf dar ghaflat az Ḥaqq mashhūd ast. Va
hamān loqmeh nān keh be-nafs-e ān moʿreḍ be-Allāh moʿadhdhab shodeh
panāh mībarad be-Khodā az īnkeh naṣīb-e ū gashteh. Va va-llāh kullu mā
yaʾkulu la-yanūḥu wa-yaḍijju wa-yaqūlu yā ilāhī bi-ayyi jurmin jaʿaltanī
rizqa hādhā al-fājir alladhī kafara bika wa-bi-āyātika. Va lāken to
shāʿer nashodeh va nakhāhī shod va az īn marāteb godhashteh.

If everyone on earth, together with all the fine blessings and good
fruits created upon it, were brought into a single gathering for the
sake of one believer in God, out of love for God, this would not be, nor
would it ever be, extravagance. But if you and those like you consume a
morsel of bread, you have been and will be extravagant, for extravagance
is manifest in heedlessness of God. That very morsel of bread, tormented
by that person who turns away from God, takes refuge in God from having
become his portion. By God, everything he eats laments, cries out, and
says, ``My God, for what crime have You made me the sustenance of this
wicked man who has denied You and Your signs?'' Yet you have not
perceived this, nor will you. And leaving these matters aside---

\subsection{614}\label{section-1109}

\begin{otherlanguage}{persian}

و اللّه الّذى لا اله الّا هو و بحقّ طلعت مبارك خود ايشان كه اكثر اوقات در
حرم غذاى يوم نبوده و بسا از ایّام كه قريب بعصر نان از براى اهل حرم تحصيل
ميشد آن هم بدين و بعضى احيان كه بعضى از اصحاب ضيافت مينمودند بعضى از
احباب را حبّاً للّه طلب ميفرمودند و بعد از قرائت آيات الهى و كلمات عزّ
صمدانى هر كسى بقدر خود مرزوق ميشد و اين مقصود اظهار حبّ بوده بارى در آن
ایّام خيال نميرفت كه مثل تو معرض محتجبی‌ پيدا شود كه باينگونه امور اعتراض
نمايد ولكن فو‌اللّه الّذى لا اله الّا هو كه آن نفس معرض باللّه كه جميع اين
كلمات را بتو القا نموده‌در جميع ایّام راحت بوده و با نساء متعدّده بعيش
مشغول و بحقّ خود ايشان كه لازال زخارف ملكيّه نزدش موجود و هر وقت كه
فى‌الجمله گفتگوئى در عراق ميشد عازم بصره و ديار اخرى بوده چنانچه كل
ميدانند و لكن كلمات الحقّ لا يزيدك الّا غلّا و اعراضاً

\end{otherlanguage}

Va-llāhi alladhī lā ilāha illā huwa va be-ḥaqq-e ṭalʿat-e mobārak-e
khod-e īshān keh akthar-e awqāt dar ḥaram ghadhā-ye yawm nabūdeh va basā
az ayyām keh qarīb be-ʿaṣr nān az barā-ye ahl-e ḥaram taḥṣīl mīshod ān
ham be-dayn. Va baʿḍī aḥyān keh baʿḍī az aṣḥāb ḍiyāfat mīnamūdand baʿḍī
az aḥbāb rā ḥubban li-llāh ṭalab mīfarmūdand va baʿd az qerāʾat-e āyāt-e
elāhī va kalemāt-e ʿezz-e ṣamadānī har kasī be-qadr-e khod marzūq mīshod
va īn maqṣūd eẓhār-e ḥobb būdeh. Bārī dar ān ayyām khayāl namīraft keh
methl-e to moʿreḍ-e moḥtajebī peydā shavad keh be-īngūneh omūr eʿterāḍ
namāyad. Va lāken fa-wa-llāhi alladhī lā ilāha illā huwa keh ān nafs-e
moʿreḍ be-Allāh keh jamīʿ-e īn kalemāt rā be-to elqā namūdeh dar jamīʿ-e
ayyām rāḥat būdeh va bā nesā-ye motaʿadded eh be-ʿeysh mashghūl, va
be-ḥaqq-e khod-e īshān keh lāzāl zakhāref-e molkīyeh nazdash mawjūd va
har vaqt keh fī al-jomleh goftogūʾī dar ʿErāq mīshod ʿāzem-e Baṣreh va
deyār-e okhrā būdeh chenāncheh koll mīdānand. Va lāken kalimāt al-Ḥaqq
lā yazīduka illā ghillan wa-iʿrāḍan.

By God, besides whom there is no other god, and by His own blessed
countenance, most of the time there was no food for the day in the
household. On many days, bread for the household was obtained only
toward the afternoon, and even then on credit. Sometimes, when some of
the companions offered hospitality, He would invite some of the loved
ones out of love for God. After the recitation of divine verses and
words of everlasting glory, each person would receive his share. The
purpose of this was to express love. In short, in those days it never
occurred to anyone that a veiled opponent like you would arise and
object to such things. Yet, by God, besides whom there is no other god,
that person who turns away from God and has put all these words into
your mind has lived in comfort all his days, occupied with pleasure with
numerous women. By His own truth, worldly finery has always been at that
man's disposal, and whenever there was the slightest commotion in Iraq,
he would set off for Basra or other lands, as everyone knows. Yet the
words of the True One only increase your rancor and aversion.

\subsection{615}\label{section-1110}

\begin{otherlanguage}{persian}

اى بی‌خبر بی‌بصر در حسين بن علىّ چه ميگوئى كه در حين حركت از مدينه طيّبه
مباركه با هودجهاى زرّين حركت فرمودند و با كمال اسباب ظاهره چنانچه امثال
تو ميگويند كه ساربان حضرت چون ديده بود كه بند ازار مبارك مطرّز بلآلىٴ است
لذا بعد از شهادت حضرت رفته و بعد واقع شد آنچه اين عبد شرم مينمايد از ذكر
آن البتّه آن را هم نسبت بدنيا خواهى داد

\end{otherlanguage}

Ey bī-khabar-e bī-baṣar, dar Ḥosayn ibn ʿAlī cheh mīgūʾī keh dar ḥīn-e
ḥarakat az Madīneh-ye ṭayyebeh-ye mobārakeh bā howdajhā-ye zarrīn
ḥarakat farmūdand va bā kamāl-e asbāb-e ẓāhereh. Chenāncheh amthāl-e to
mīgūyand keh sārebān-e Ḥaḍrat chon dīdeh būd keh band-e ezār-e mobārak
moṭarraz be-laʾālī ast ledhā baʿd az shahādat-e Ḥaḍrat rafteh va baʿd
vāqeʿ shod āncheh īn ʿabd sharm mīnamāyad az dhekr-e ān. Albatteh ān rā
ham nesbat be-donyā khāhī dād.

You who are unaware and blind! What do you say of Ḥusayn ibn ʿAlī, who,
when departing from the good and blessed city of Medina, set out with
golden howdahs and every outward provision? Indeed, people like you say
that, because His camel driver had seen that His blessed waistband was
embroidered with pearls, he went there after His martyrdom, and then
there occurred what this servant is ashamed to mention. No doubt you
will attribute that, too, to worldliness.

\subsection{616}\label{section-1111}

\begin{otherlanguage}{persian}

و از كلّ اين امور گذشته در اين ظهور تا بحال اسباب دنيا فراهم نيامد تا تو
يا غير تو اعتراض نمايد فى كلّ حين قد كنّا فى بلآء و شدّةعظيم لو انتم من
المنصفين و اينكه بالقاى آن معرض ملقى نوشته كه جمال احديّه با اجامر و
اوباش معاشر بوده‌اند و بصحبتهاى لاطايل مشغول بوده‌اند يكذّبك كلمات بعدك
كلمات قبلك مينويسى بيست سال با عرفا مجالس و مؤانس بوده‌اند و بعد نوشته با
اجامر و اوباش مت بغيظك يا ايّها الغافل انّك انت الّذى ذكرت فى كتاب القبل
بقوله تعالى كمثل الكلب ان تحمل عليه يلهث او تتركه يلهث در كلّ احوال بهوى
ناطقند و در مسالك نفس سالك بلى و اللّه اگر با اجامر معاشر نبودند البتّه با
سيّد محمّد و مرشدش معاشر نميشدند اجامر و او باش امثال اين نفوس بوده‌اند و
الّا آن اجامر كه بوده و اوباش چه نفوسى بوده‌اند

\end{otherlanguage}

Va az koll-e īn omūr godhashteh dar īn ẓohūr tā be-ḥāl asbāb-e donyā
farāham nayāmad tā to yā ghayr-e to eʿterāḍ namāyad. Fī kulli ḥīn qad
kunnā fī balāʾ wa-shiddah ʿaẓīm law antum min al-munṣifīn. Va īnkeh
be-elqā-ye ān moʿreḍ-e molqī neveshteh keh Jamāl-e Aḥadīyeh bā ajāmer va
awbāsh moʿāsher būdeh-and va be-ṣoḥbathā-ye lāṭāyel mashghūl būdeh-and,
yukadhdhibuka kalimātu baʿdika kalimātu qablika. Mīnevīsī bīst sāl bā
ʿorafā majāles va moʾānes būdeh-and va baʿd neveshteh bā ajāmer va
awbāsh. Mut bi-ghayẓika yā ayyuhā al-ghāfil. Innaka anta alladhī
dhukirta fī kitāb al-qabl bi-qawlihi taʿālā ka-mathal al-kalb in taḥmil
ʿalayhi yalhath aw tatrukhu yalhath. Dar koll-e aḥvāl be-havā nāṭeqand
va dar masālek-e nafs sālek. Balī va-llāh agar bā ajāmer moʿāsher
nabūdand albatteh bā Seyyed Moḥammad va morshedash moʿāsher namīshodand.
Ajāmer va awbāsh amthāl-e īn nofūs būdeh-and va ellā ān ajāmer keh būdeh
va awbāsh cheh nofūsī būdeh-and.

Apart from all these matters, in this Revelation worldly means have not
been available up to now for you or anyone else to object to. At every
moment We have been in affliction and great hardship, if you are
fair-minded. As for what you have written at the prompting of that
opponent who puts words into your mind---that the Beauty of Oneness
associated with rabble and scoundrels and engaged in idle
conversations---your later words and your earlier words give you the
lie. You write that for twenty years He sat and enjoyed companionship
with mystics, and then you write, ``with rabble and scoundrels.'' Die in
your rage, heedless one! You are the one mentioned in the former Book in
His words---exalted is He---``Like a dog: if you assail it, it pants,
and if you leave it alone, it pants.'' In every condition they speak
from desire and tread the paths of self. Yes, by God, had He not
associated with rabble, He would certainly not have associated with
Siyyid Muḥammad and his guide. People like these have been the rabble
and scoundrels. Otherwise, who were that rabble, and who were those
scoundrels?

\subsection{617}\label{section-1112}

\begin{otherlanguage}{persian}

البتّه كلمات الهى را لاطايل دانسته چه كه تكلّم نفرموده‌اند مگر بآيات الهى
برو و از اهل عراق سؤال كن بسيار كلمات صدقى بتو القا نموده‌اند يكذّبك كلّ
شیء و لكن انت لا تشعر از اين گذشته اى بی‌انصاف ظالم تو كجا بوده و كجا
ديده كه ايشان با اجامر و اوباش نشسته باشند البتّه اراذل و اوبا ش مؤمنين
باللّه را دانسته لازال در احيان هر ظهور مشركين مخصوص بهمين كلمه تكلّم
نموده‌اند بقوله تعالى و ما نراك اتّبعك الّا الّذينهم اراذلنا بادى الرّأى
چنانچه از قبل ذكر شد لايق تو آنكه بامثال اين كلمات تكلّم نمايد فو‌اللّه
همين كلمه تو دليل بر حقّیّت حقّ بوده و خواهد بود لعن اللّه اوّل ظالم ظلم نقطة
الأولی فى ظهوره الأخرى و افترى على نفسه و كفر بآياته و ملا صدره من غلّه و
بغضائه و بهذا اللّعن يلعنه كلّ من فى السّموات و الأرضين

\end{otherlanguage}

Albatteh kalemāt-e elāhī rā lāṭāyel dānesteh, cheh keh takallom
nafarmūdeh-and magar be-āyāt-e elāhī. Boro va az ahl-e ʿErāq soʾāl kon.
Besyār kalemāt-e ṣedqī be-to elqā namūdeh-and! Yukadhdhibuka kullu shayʾ
wa-lākin anta lā tashʿur. Az īn godhashteh ey bī-enṣāf-e ẓālem, to kojā
būdeh va kojā dīdeh keh īshān bā ajāmer va awbāsh neshasteh bāshand?
Albatteh arādhel va awbāsh moʾmenīn be-Allāh rā dānesteh. Lāzāl dar
aḥyān-e har ẓohūr moshrekīn makhṣūṣ be-hamīn kalemeh takallom
namūdeh-and, bi-qawlihi taʿālā wa-mā narāka ittabaʿaka illā alladhīnahum
arādhilunā bādiya al-raʾy. Chenāncheh az qabl dhekr shod lāyeq-e to
ānkeh be-amthāl-e īn kalemāt takallom namāyad. Fa-va-llāh hamīn
kalemeh-ye to dalīl bar ḥaqqīyat-e Ḥaqq būdeh va khāhad būd. Laʿana
Allāhu awwala ẓālim ẓalama Nuqṭat al-Ūlā fī ẓuhūrihi al-ukhrā wa-iftarā
ʿalā nafsihi wa-kafara bi-āyātihi wa-malaʾa ṣadrahu min ghillihi
wa-baghḍāʾihi wa-bi-hādhā al-laʿn yalʿanuhu kullu man fī al-samāwāt
wa-al-araḍīn.

Clearly, you have regarded the divine words as idle talk, for He spoke
nothing but divine verses. Go and ask the people of Iraq. What a wealth
of truthful words they have put into your mind! Everything gives you the
lie, yet you do not perceive it. Furthermore, you unjust oppressor,
where were you, and where did you see Him sitting with rabble and
scoundrels? Clearly, you have taken the believers in God to be the lowly
and the scoundrels. In the days of every revelation, the idolaters have
always uttered precisely these words, as in His saying---exalted is
He---``We see that none have followed you but those who are the lowliest
among us, at first sight.'' As mentioned before, it is fitting for you
to utter such words. By God, this very statement of yours has been and
will remain a proof of the truth of the True One. May God curse the
first oppressor who wronged the Primal Point in His latter
manifestation, slandered His person, denied His signs, and filled His
breast with his rancor and hatred. With this curse, all who are in the
heavens and the earths curse him.

\subsection{618}\label{section-1113}

\begin{otherlanguage}{persian}

و اللّه الّذى لا اله الّا هو بلايائى بر اين جمال مظلوم وارد شده كه بر احدى
از اوّل ابداع تا حين وارد نشده امرى كه مشهود عالميان است او را انكار
نموده‌ايد مثل آنكه كلّ مشاهده نموده‌اند كه در امر اللّه بنفسى مداهنه
نفرموده‌اند و بنفسى ملتجى نشده‌اند چنانچه در عراق سركار مشير خواستند
معاشرت نمايند قبول نفرمودند تا آنكه بالأخره‌بمسجد تشريف بردند چنانچه كلّ
شنيده‌اند و بعد از ورود بمدينه كبير هم نزد احدى نرفتند و بنفسى توجّه
نفرمودند و كلّ شاهد و گواهند مع‌ذلك چنين نفسى را مينويسيد كه با اجامر و
اوباش معاشر بوده و نفسى را كه براى لقمه نان بكلّ نفوس شكايت كرده كه سواد
عرايضش الآن موجود و زن مخصوص نان بسرايه فرستاده و باستانبول دويده و بهر
نفسى ملتجى شده افعال چنين نفسى مقبول و افعال حقّ غير مقبول اى و اللّه كذلك
ينبغى لكم و لا نفسكم

\end{otherlanguage}

Va-llāhi alladhī lā ilāha illā huwa, balāyāʾī bar īn Jamāl-e Maẓlūm
vāred shodeh keh bar aḥadī az avval-e ebdāʿ tā ḥīn vāred nashodeh. Amrī
keh mashhūd-e ʿālamiyān ast ū rā enkār namūdeh-īd, methl-e ānkeh koll
moshāhedeh namūdeh-and keh dar Amr Allāh be-nafsī modāheneh
nafarmūdeh-and va be-nafsī moltajī nashodeh-and. Chenāncheh dar ʿErāq
Sarkār-e Moshīr khāstand moʿāsherat namāyand, qabūl nafarmūdand tā ānkeh
be-al-ākhereh be-masjed tashrīf bordand chenāncheh koll shenīdeh-and. Va
baʿd az vorūd be-Madīneh-ye Kabīr ham nazd-e aḥadī naraftand va be-nafsī
tavajjoh nafarmūdand va koll shāhed va govāhand. Maʿa dhālek chenīn
nafsī rā mīnevīsīd keh bā ajāmer va awbāsh moʿāsher būdeh. Va nafsī rā
keh barā-ye loqmeh nān be-koll-e nofūs shekāyat kardeh keh savād-e
ʿarāyeḍash al-ān mawjūd va zan makhṣūṣ-e nān be-sarāyeh ferestādeh va
be-Estānbūl davīdeh va be-har nafsī moltajī shodeh, afʿāl-e chenīn nafsī
maqbūl va afʿāl-e Ḥaqq ghayr-e maqbūl. Ī va-llāh kadhālika yanbaghī
lakum wa-li-anfusikum.

By God, besides whom there is no other god, afflictions have befallen
this wronged Beauty such as have befallen no one from the beginning of
creation until now. You have denied something witnessed by all the
world. For example, everyone has seen that in the Cause of God He has
flattered no one and sought refuge with no one. In Iraq, when His
Excellency the Mushīr wished to meet with Him, He did not accept, until
at last He went to the mosque, as everyone has heard. After His arrival
in the Great City, too, He visited no one and turned to no one, and all
are witnesses to this. Despite this, you write of such a Person that He
associated with rabble and scoundrels. Yet there is one who complained
to everyone for a morsel of bread---the copies of his petitions still
exist---who sent a woman to the government residence specifically for
bread, ran to Istanbul, and sought refuge with every person. The deeds
of such a person are acceptable, while the deeds of the True One are
unacceptable! Yes, by God, this is what befits you and your own selves.

\subsection{619}\label{section-1114}

\begin{otherlanguage}{persian}

قدرى تفكّر كن كه همين كلمه امثال تو چه قدر عظيم است و چگونه خورشيد نورانى
را در فكر افتاده‌ايد كه بحجبات نفس و هوى مستور داريد و اعمالش را ضايع
كنيد نه و اللّه جز امثال تو احدى باين كلمات التفات ننمايد چه كه كلّ فعل حقّ
را ببصر خود دیده‌اند يومى از ایّام فرمودند كه اعظم از اين بلايا در سرادق
قضا مستور است و بايد نازل شود چنانچه خوابی‌ در عراق دیده‌ام و از آن چنين
مستفاد ميشود قال و قوله الحقّ كنت نائماً فى ليلة البلماء بعد عفراء اذاً
رأيت بأن اجتمعت فى حولى النّبيّون و المرسلون و هم قد جلسوا فى اطرافى و
كلّهم ينوحون و يبكون و يصرخون و يضجّون و انّى تحيّرت فى نفسى فسئلت عنهم اذاً
اشتدّ بكائهم و صريخهم و قالوا لنفسك يا سرّ الأعظم و يا هيكل القدم و بكوا
على شأن بكيت ببكائهم و

\end{otherlanguage}

Qadrī tafakkur kun kih hamīn kalimiy-i amthāl-i tu chih qadr ʿaẓīm ast
va chigūnih khurshīd-i nūrānī rā dar fikr uftādih-īd kih bi-ḥujubāt-i
nafs va havā mastūr dārīd va aʿmālash rā ḍāyiʿ kunīd. Nah va-Allāh, juz
amthāl-i tu aḥadī bi-īn kalimāt iltifāt nanamāyad chih kih kull fiʿl-i
Ḥaqq rā bi-baṣar-i khud dīdih-and. Yawmī az ayyām farmūdand kih aʿẓam az
īn balāyā dar surādiq-i qaḍā mastūr ast va bāyad nāzil shavad chunānchih
khwābī dar ʿIrāq dīdih-am va az ān chunīn mustafād mīshavad. Qāla
wa-qawluhu al-ḥaqq: Kuntu nāʾiman fī laylatin balmāʾ baʿda ʿafrāʾ idhan
raʾaytu bi-an ijtamaʿat fī ḥawlī al-nabiyyūna wa-l-mursalūn wa-hum qad
jalasū fī aṭrāfī wa-kulluhum yanūḥūna wa-yabkūna wa-yaṣrukhūna
wa-yaḍijjūn wa-innī taḥayyartu fī nafsī fa-saʾaltu ʿanhum idhan ishtadda
bukāʾuhum wa-ṣarīkhuhum wa-qālū li-nafsika yā sirr al-aʿẓam wa-yā haykal
al-qidam wa-bakaw ʿalā shaʾnin bakaytu bi-bukāʾihim wa-

Reflect a little on how grave this very statement by people like you is,
and how you have conceived the idea of concealing the radiant Sun behind
the veils of self and desire and rendering His deeds fruitless. No, by
God, no one but people like you would pay attention to such words, for
all have seen the deeds of the True One with their own eyes. One day He
said that afflictions greater than these lay concealed in the tabernacle
of destiny and must descend, as He had seen a dream in Iraq from which
this could be understood. He said---and His word is truth---``I was
asleep on a dark night following a bright one, when I saw that the
Prophets and Messengers had gathered around Me. They were seated on
every side of Me, and all were lamenting, weeping, crying out, and
wailing. I was bewildered within Myself and questioned them. Then their
weeping and cries intensified, and they said, `For You, O Most Great
Mystery, O Temple of Eternity!' They wept so much that I wept at their
weeping, and

\subsection{620}\label{section-1115}

\begin{otherlanguage}{arabic}

اذا سمعت بكاء اهل ملأ الأعلی و فى تلك الحالة خاطبونى و قالوا قد عظم
بلائك يا سدرة المنتهى و كبر قضائك يا سرّ الآخرة و الأولی عليك بالصّبر يا
آية الكبرى و ظهور نقطة الأولی ثمّ عليك بالصّبر يا شجرة القصوى و ظهور
القضآء فى ملكوت الامضآء فسوف ترى بعينك ما لا رءاه احد من المعشر النّبيّين
و تشهد ما لا شهده احد من العالمين و تسمع ما لا سمعه اذن الأصفیآء و
الأودّاء فصبراً صبراً يا سرّ اللّه المكنون و رمز المخزون و كلمته المحتوم و
كتابه المختوم

\end{otherlanguage}

Idhā samiʿtu bukāʾa ahl-i malaʾ-i aʿlā wa-fī tilka al-ḥālati khāṭabūnī
wa-qālū qad ʿaẓuma balāʾuka yā sidrat al-muntahā wa-kabura qaḍāʾuka yā
sirr al-ākhirati wa-l-ūlā ʿalayka bi-l-ṣabr yā āyat al-kubrā wa-ẓuhūr
nuqṭat al-ūlā thumma ʿalayka bi-l-ṣabr yā shajarat al-quṣwā wa-ẓuhūr
al-qaḍāʾ fī malakūt al-imḍāʾ fa-sawfa tarā bi-ʿaynika mā lā raʾāhu
aḥadun min al-maʿshar al-nabiyyīn wa-tashhadu mā lā shahidahu aḥadun min
al-ʿālamīn wa-tasmaʿu mā lā samiʿahu udhun al-aṣfiyāʾ wa-l-awiddāʾ
fa-ṣabran ṣabran yā sirr Allāh al-maknūn wa-ramz al-makhzūn
wa-kalimatuhu al-maḥtūm wa-kitābuhu al-makhtūm.

when I heard the weeping of the people of the Supreme Concourse, in that
state they addressed Me and said, `Great indeed is Your affliction, O
Tree beyond which there is no passing, and momentous is Your destiny, O
Mystery of the next world and this one. Be patient, O Most Great Sign
and Manifestation of the Primal Point. Again, be patient, O Uttermost
Tree and Manifestation of the Decree in the kingdom of its fulfillment.
You will see with Your own eyes what none of the company of Prophets has
seen, witness what no one in all the worlds has witnessed, and hear what
the ears of the chosen ones and the beloved have never heard. Be
patient, be patient, O hidden Mystery of God, His treasured Symbol, His
irrevocable Word, and His sealed Book.'

\subsection{621}\label{section-1116}

\begin{otherlanguage}{arabic}

و كنت معهم فى تلك الليلة خاطبتهم و خاطبونى الى ان قرب الفجر و ارفعت رأسى
عن النّوم و كنت متفكّراً فى نفسى ما بلآء الّذى ما شهده احد فى الابداع و ما
حكيته عند احد الى ان حضر تلقآء الوجه احد من علمآء العراق الّذى سمّى بعبد
السّلام افندى فلمّا حضر القيناه و اخبرناه و قصصنا له ما ارانى اللّه فى
المنام و هو حزن تلقآء الوجه بحزن عظيم كذلك كان الأمر و لكن النّاس هم فى
غفلة مبين

\end{otherlanguage}

Wa-kuntu maʿahum fī tilka al-laylati khāṭabtuhum wa-khāṭabūnī ilā an
qaruba al-fajr wa-arfaʿtu raʾsī ʿan al-nawm wa-kuntu mutafakkiran fī
nafsī mā balāʾ alladhī mā shahidahu aḥadun fī al-ibdāʿ wa-mā ḥakaytuhu
ʿinda aḥadin ilā an ḥaḍara tilqāʾa al-wajh aḥadun min ʿulamāʾ al-ʿIrāq
alladhī summiya bi-ʿAbd al-Salām Afandī fa-lammā ḥaḍara alqaynāhu
wa-akhbarnāhu wa-qaṣaṣnā lahu mā arānī Allāhu fī al-manām wa-huwa ḥazina
tilqāʾa al-wajh bi-ḥuznin ʿaẓīm kadhālika kāna al-amr wa-lākinna al-nāsa
hum fī ghaflatin mubīn.

``I was with them that night, speaking to them and hearing them speak to
Me, until dawn drew near and I raised My head from sleep. I wondered
within Myself what this affliction could be that no one in creation had
witnessed. I told no one of it until one of the scholars of Iraq, named
ʿAbd al-Salām Afandī, came into My presence. When he arrived, We related
it to him, informed him, and recounted to him what God had shown Me in
the dream. He was deeply saddened in Our presence. Such was the matter,
yet the people are in manifest heedlessness.''

\subsection{622}\label{section-1117}

\begin{otherlanguage}{persian}

اين رؤيائى است كه در عراق مشاهده فرموده و ذكر فرموده‌اند و در اين ایّام
بعضى از آن تعبير شده چنانچه شنيده‌ايم از اقوال شما آنچه را كه احدى از قبل
نشنيده و دیده‌ايم آنچه را نفسى نديده و فرمودند در كلّ احوال صابر و شاكريم
و جميع مقصود آن نفوس اينست كه چرا حقّ ظاهر شده فو‌اللّه هذا لم يكن من عندى
بل من لدن من ارسلنى بالحقّ و جعلنى رحمة على العالمين

\end{otherlanguage}

Īn ruʾyāʾī ast kih dar ʿIrāq mushāhidih farmūdih va dhikr farmūdih-and
va dar īn ayyām baʿḍī az ān taʿbīr shudih chunānchih shinīdih-īm az
aqvāl-i shumā ānchih rā kih aḥadī az qabl nashnīdih va dīdih-īm ānchih
rā nafsī nadīdih va farmūdand dar kull-i aḥvāl ṣābir va shākirīm va
jamīʿ-i maqṣūd-i ān nufūs īnast kih chirā Ḥaqq ẓāhir shudih. Fa-wa-Allāh
hādhā lam yakun min ʿindī bal min ladun man ar salanī bi-l-ḥaqq
wa-jaʿalanī raḥmatan ʿalā al-ʿālamīn.

This is a dream that He saw in Iraq and recounted, and in these days
part of it has been fulfilled: We have heard from your utterances what
no one has ever heard before and have seen what no one has seen. He
said, ``In every circumstance We are patient and thankful.'' The entire
concern of those souls is why the True One has appeared. By God, this
did not originate with Me, but with Him who sent Me with the truth and
made Me a mercy to all the worlds.

\subsection{623}\label{section-1118}

\begin{otherlanguage}{persian}

قوله كبر غلّه مذكور نموده‌ بوديد كه زن در سرايه فرستاده عرض نموده‌اند اوّلا
آنكه ظلم موجب زياده از اين خواهد بود علاوه مرحبا از شما كه با ادّعاى حقّیّت
شما و زيركى زير آب‌كشى ديگر از انبیاء و اولیاء نموديد آيا بجهت فدك حضرت
فاطمه دست حسنين را نگرفته و بدر مهاجر و انصار نبرده آفرين بر شما انتهى

\end{otherlanguage}

Qawluhu, kabura ghilluhu: Madhkūr namūdih būdīd kih zan dar sarāyih
firistādih ʿarḍ namūdih-and. Avvalan ānkih ẓulm mūjib-i ziyādih az īn
khwāhad būd. ʿAlāvih marḥabā az shumā kih bā iddiʿā-yi ḥaqqiyyat-i shumā
va zīrakī zīr-āb-kashī-yi dīgar az anbiyā va awliyā namūdīd. Āyā
bi-jahat-i Fadak Ḥaḍrat-i Fāṭimih dast-i Ḥasanayn rā nagiriftih va
bi-dar-i Muhājir va Anṣār naburdih? Āfarīn bar shumā. Intahā.

His words---great is his rancor: ``You mentioned that they had sent a
woman to the government residence to present a petition. First,
oppression will lead to more than this. Moreover, bravo to you! With
your claim to being in the right and your cleverness, you have made yet
another attempt to undermine the Prophets and holy ones. Did not Her
Holiness Fāṭimih, for the sake of Fadak, take Ḥasan and Ḥusayn by the
hand and bring them to the doors of the Emigrants and the Helpers? Well
done!'' End of quotation.

\subsection{624}\label{section-1119}

\begin{otherlanguage}{arabic}

صه لسانك يا ايّها الكافر باللّه ثمّ صه لسانك يا ايّها المشرك باللّه ثمّ صه
لسانك يا ايّها المحارب باللّه

\end{otherlanguage}

Ṣah lisānaka yā ayyuhā al-kāfir bi-Allāh thumma ṣah lisānaka yā ayyuhā
al-mushrik bi-Allāh thumma ṣah lisānaka yā ayyuhā al-muḥārib bi-Allāh.

Silence your tongue, you who disbelieve in God! Again, silence your
tongue, you who associate partners with God! Again, silence your tongue,
you who wage war against God!

\subsection{625}\label{section-1120}

\begin{otherlanguage}{arabic}

ولو اكررّ تلك الكلمة من حینئذ الى آخر الّذى لا آخر له لتكون مستحقّاً بها بل
لا تذكر عند ما ارتكبت فى جنب اللّه يا ايّها الغافل عن امر اللّه و المستكبر
على نفس اللّه

\end{otherlanguage}

Wa-law ukarrir tilka al-kalimata min ḥīniʾidhin ilā ākhir alladhī lā
ākhira lahu la-takūna mustaḥiqqan bihā bal lā tudhkaru ʿinda mā
irtakabta fī janb Allāh yā ayyuhā al-ghāfil ʿan amr Allāh wa-l-mustakbir
ʿalā nafs Allāh.

Even if I repeated those words from this moment until the end that has
no end, you would deserve them. Indeed, they are nothing beside what you
have committed against God, you who are heedless of the Cause of God and
arrogant toward God Himself.

\subsection{626}\label{section-1121}

\begin{otherlanguage}{persian}

بعد از ورود اين ارض جمال مبارك و حضرت كليم و ميرزا يحيى هر يك در بيت
عليحدّه مقرّ داشتند و اين ارض محلّى نيست كه نفسى بنفسى ظلم نمايد و ظلمى را
كه مشرك باللّه ارتكاب نمود در وقتى بود كه تفريق كبرى واقع نشده بود و بعد
از ظهور ما فى صدره تفصيل و تفريق واقع چنانچه ذكر شد و ابداً ديگر ملاقات
نشده و گفتگوى لسانى هم بميان نيامده فاسئل من اهل تلك الأرض لكى تجد الى
الصّدق سبيل

\end{otherlanguage}

Baʿd az vurūd-i īn arḍ Jamāl-i Mubārak va Ḥaḍrat-i Kalīm va Mīrzā Yaḥyā
har yak dar bayt-i ʿalāḥaddih maqarr dāshtand va īn arḍ maḥallī nīst kih
nafsī bi-nafsī ẓulm namāyad va ẓulmī rā kih mushrik bi-Allāh irtikāb
namūd dar vaqtī būd kih tafrīq-i kubrā vāqiʿ nashudih būd va baʿd az
ẓuhūr-i mā fī ṣadrihi tafṣīl va tafrīq vāqiʿ chunānchih dhikr shud va
abadan dīgar mulāqāt nashudih va guftugū-yi lisānī ham bi-miyān
nayāmadih. Fa-sʾal min ahl tilka al-arḍ li-kay tajida ilā al-ṣidq sabīl.

After arriving in this land, the Blessed Beauty, His Holiness Kalīm, and
Mīrzā Yaḥyā each resided in a separate house. This land is not a place
where one person can oppress another. The wrong committed by the one who
associates partners with God took place before the Most Great Separation
had occurred. After what lay within his breast was revealed, the
division and separation took place, as has been mentioned. Since then
there has never been another meeting, nor has any spoken exchange taken
place. Ask the people of that land, so that you may find a path to the
truth.

\subsection{627}\label{section-1122}

\begin{otherlanguage}{persian}

و ديگر آنكه مملكت مال دولت عليّه و در اين ارض جميع مسجون و غريب و بعد از
اظهار ما فى سرّه احدى معاشرت ننموده‌جز درويش صدقعلى كه بجهت خدمات او
مراوده مينمود و هم‌چنين يك نفر در خدمت جمال اقدس و يكنفر هم خدمت حضرت
كليم و باقى اصحاب كلّ در منزلهاى عليحدّه متفرّق ديگر اين ظلم چه بوده و از
چه رو وارد شده با آنكه جميع ذرّات ميدانند كه سبب ظلم آن ظالم بود كه ايشان
بافتضاح خود راضى شدند و تفريق فرمودند

\end{otherlanguage}

Va dīgar ānkih mamlakat māl-i Dawlat-i ʿAliyyih va dar īn arḍ jamīʿ
masjūn va gharīb va baʿd az iẓhār-i mā fī sirrihi aḥadī muʿāsharat
nanamūdih juz Darvīsh Ṣidqʿalī kih bi-jahat-i khadamāt-i ū murāvadih
mīnamūd va hamchunīn yak nafar dar khidmat-i Jamāl-i Aqdas va yak nafar
ham khidmat-i Ḥaḍrat-i Kalīm va bāqī aṣḥāb kull dar manzilhā-yi
ʿalāḥaddih mutafarriq. Dīgar īn ẓulm chih būdih va az chih rū vārid
shudih bā ānkih jamīʿ-i dharrāt mīdānand kih sabab-i ẓulm ān ẓālim būd
kih īshān bi-iftiḍāḥ-i khud rāḍī shudand va tafrīq farmūdand.

Moreover, the country belongs to the Ottoman government, and in this
land all are prisoners and strangers. After he disclosed what lay hidden
within him, no one associated with him except Darvīsh Ṣidqʿalī, who
visited him to attend to his needs. Likewise, one person served the Most
Holy Beauty and one served His Holiness Kalīm, while all the other
companions were dispersed in separate dwellings. What, then, was this
oppression, and how did it come about? Every atom knows that the cause
of the wrongdoing was that oppressor, who consented to his own disgrace
and brought about the separation.

\subsection{628}\label{section-1123}

\begin{otherlanguage}{persian}

و مع‌ذلك تقول ما لا تشعر اصل شجره ظلم را عدل شمرده و شجره فضل و عدل را
ظالم دانسته هذا مكر منك و منهم فى امر اللّه آنقدر درايت ندارى كه سخنى
استماع نمودى تفكّر كنى در همان ایّام كه زنش بسرايه رفته فو‌اللّه كه در اوّل
ماه مخصوص او از آن زخرفيكه در هر شهر مياوردند مع يك بقچه اسباب درويش كه
مباشر او بود الآن موجود است برده و چند نفر هم مجدّد در اين ارض وارد شده
بودند كلّ مطّلع از جمله جناب ميرزا حيدر على و سلمان و استاد عبد الكريم از
اهل صاد و مع‌ذلك مخصوصاً زن بسرايه فرستاده و بهمين لفظ گفته نسبت بجمال
ابهى كه ما را آورده و نان و لباس بما نميدهد و افتضاحى در اين ارض برپاشد
كه يكى از قنصلهاى اين ارض بسيار تعجّب نموده ‌بشخصى ذكر نمود كه امروز امر
عجبی‌ واقع شده و جميع اعجام بشماتت برخواستند كه در اين طايفه عصمت و عفّت
نيست

\end{otherlanguage}

Va maʿa dhālik taqūlu mā lā tashʿur. Aṣl-i shajariy-i ẓulm rā ʿadl
shumurdih va shajariy-i faḍl va ʿadl rā ẓālim dānistih. Hādhā makrun
minka wa-minhum fī amr Allāh. Ānqadr dirāyat nadārī kih sukhanī istimāʿ
namūdī tafakkur kunī dar hamān ayyām kih zanash bi-sarāyih raftih.
Fa-wa-Allāh kih dar avval-i māh makhṣūṣ-i ū az ān zukhrufī kih dar har
shahr mīāvardand maʿ yak buqchih asbāb Darvīsh kih mubāshir-i ū būd
al-ān mawjūd ast burdih va chand nafar ham mujaddad dar īn arḍ vārid
shudih būdand kull muṭṭaliʿ az jumlih Jināb-i Mīrzā Ḥaydar ʿAlī va
Salmān va Ustād ʿAbd al-Karīm az ahl-i Ṣād va maʿa dhālik makhṣūṣan zan
bi-sarāyih firistādih va bi-hamīn lafẓ guftih nisbat bi-Jamāl-i Abhā kih
mā rā āvardih va nān va libās bi-mā namīdihad va iftiḍāḥī dar īn arḍ
bar-pā shud kih yakī az qunṣulhā-yi īn arḍ bisyār taʿajjub namūdih
bi-shakhṣī dhikr namūd kih imrūz amr-i ʿajabī vāqiʿ shudih va jamīʿ-i
aʿjām bi-shamātat bar-khāstand kih dar īn ṭāʾifih ʿiṣmat va ʿiffat nīst.

Yet you say what you do not understand. You have accounted the very root
of the tree of oppression to be justice, and have judged the tree of
grace and justice to be oppressive. This is a plot by you and them
against the Cause of God. You do not have enough discernment to reflect
on something you have heard. In those very days when his wife went to
the government residence---by God---at the beginning of the month,
Darvīsh, who acted as his agent and is still alive, had taken his
allotted portion of the money brought in each city, together with a
bundle of belongings. Several people had also newly arrived in this
land, and all knew of it, including Jināb-i Mīrzā Ḥaydar ʿAlī, Salmān,
and Ustād ʿAbd al-Karīm of the people of Ṣād. Yet he deliberately sent
his wife to the government residence, and she said these very words
about the All-Glorious Beauty: ``He has brought us here and gives us
neither bread nor clothing.'' Such a scandal arose in this land that one
of its consuls was greatly astonished and remarked to someone, ``A
strange thing has happened today.'' All the Persians began to gloat,
saying that this community had neither chastity nor decency.

\subsection{629}\label{section-1124}

\begin{otherlanguage}{persian}

و خود نفس غافلت آنقدر دانسته كه اگر يهود بر ايشان وارد شود بقدر مقدور
زخرف فانيه از او منع نميشود مع‌ذلك نوشتى آنچه نوشتى و فدك را كثيف نمودى و
لكن مطمئنّ باش كه در هر ارضى يك سيّدة النّسائى دارى و از آن رشته مراياى تام
واحداً بعد واحد هم بسيار و در هر ارضى موجود دوازده سنه جمال ابهى در عراق
تشريف داشتند و احدى از ايشان امرى كه فى‌الجمله مغاير و مخالف باشد نديده
از جميع اهل عراق سؤال كن حال از آن رشته مراياى تام واحداً بعد واحد يكى در
عراق است بايد از اهل عراق احوال او را پرسيد باو نوشته انّا جعلناك على
العالمين نوراً منيراً و هم‌چنين هو النّور فى النّور و انّ مثل نوره كنور نور من
نور اللّه و بامّ الفواحش اين آيهء مباركه قرآن را بعينه نوشته يا مريم اقنتى
لربّك و اسجدى و اركعى مع الرّاكعين انّ اللّه طهّرك و اصطفاك على نسآء العالمين
حال حظّ كن از آن سيّدة النّسآء و از آن مرآت تام كاش آنجا هم ميرفتى و يك
سجده مينمودى لعن اللّه الّذين تحرّفوا كلمة اللّه عن موضعها و اعرضوا عن جماله
و اتّخذوا لانفسهم عجلاً من دونه فما لهؤلآء لا يكادون يفقهون ذكراً من اللّه
العلىّ العظيم

\end{otherlanguage}

Va khud nafs-i ghāfilat ānqadr dānistih kih agar Yahūd bar īshān vārid
shavad bi-qadr-i maqdūr zukhruf-i fāniyih az ū manʿ namīshavad. Maʿa
dhālik nivishtī ānchih nivishtī va Fadak rā kathīf namūdī va lākin
muṭmaʾinn bāsh kih dar har arḍī yak Sayyidat al-Nisāʾī dārī va az ān
rishtih marāyā-yi tāmm wāḥidan baʿda wāḥid ham bisyār va dar har arḍī
mawjūd. Davāzdah sanih Jamāl-i Abhā dar ʿIrāq tashrīf dāshtand va aḥadī
az īshān amrī kih fī-l-jumlih mughāyir va mukhālif bāshad nadīdih. Az
jamīʿ-i ahl-i ʿIrāq suʾāl kun. Ḥāl az ān rishtih marāyā-yi tāmm wāḥidan
baʿda wāḥid yakī dar ʿIrāq ast bāyad az ahl-i ʿIrāq aḥvāl-i ū rā pursīd.
Bi-ū nivishtih innā jaʿalnāka ʿalā al-ʿālamīn nūran munīran va hamchunīn
huwa al-nūr fī al-nūr wa-inna mathala nūrihi ka-nūri nūrin min nūr Allāh
va bi-Umm al-Favāḥish īn āyiy-i mubārakiy-i Qurʾān rā bi-ʿaynihi
nivishtih yā Maryam uqnutī li-rabbiki wa-sjudī wa-rkaʿī maʿa al-rākiʿīn
inna Allāha ṭahharaki wa-ṣṭafāki ʿalā nisāʾ al-ʿālamīn. Ḥāl ḥaẓẓ kun az
ān Sayyidat al-Nisāʾ va az ān mirʾāt-i tāmm. Kāsh ānjā ham mīraftī va
yak sajdih mīnamūdī. Laʿana Allāhu alladhīna taḥarrafū kalimat Allāh ʿan
mawḍiʿihā wa-aʿraḍū ʿan jamālihi wa-ttakhadhū li-anfusihim ʿijlan min
dūnihi fa-mā li-hāʾulāʾ lā yakādūna yafqahūna dhikran min Allāh
al-ʿAliyy al-ʿAẓīm.

Your own heedless self knows at least this much: if a Jew came to Him,
such perishable wealth as was available would not be withheld from him.
Yet you wrote what you wrote and sullied Fadak. But rest assured that in
every land you have a ``Mistress of Women,'' and many of that succession
of ``perfect mirrors, one after another,'' are to be found in every
land. The All-Glorious Beauty resided in Iraq for twelve years, and no
one saw Him do anything even remotely improper or contrary to what is
right. Ask all the people of Iraq. Now, one of that succession of
``perfect mirrors, one after another,'' is in Iraq: you should ask the
people of Iraq about him. He wrote to him, ``We have made you a shining
light to all the worlds,'' and likewise, ``He is the light within the
light, and the likeness of his light is as the light of a light from the
light of God.'' To the Mother of Obscenities he wrote this blessed verse
of the Quran, word for word: ``O Mary, be devoutly obedient to your
Lord, prostrate yourself, and bow with those who bow. God has purified
you and chosen you above the women of all the worlds.'' Now take
pleasure in that ``Mistress of Women'' and that ``perfect mirror''! If
only you had gone there too and performed a prostration! May God curse
those who have displaced the word of God from its proper place, turned
away from His Beauty, and taken a calf for themselves in His stead. What
is wrong with these people that they can scarcely understand a word from
God, the Most High, the Almighty?

\subsection{630}\label{section-1125}

\begin{otherlanguage}{persian}

و اللّه ناله مينمايم و اين كلمات مينويسم چه كه امر بجائى رسيده كه بايد
بچنين اذكار مشغول شد قد ضيّعتم حرمة اللّه و حرمة انفسكم قلب را از كثافات
اوهام مشركين طاهر نما و ذكر انبیاء و مرسلين را باذكار انفس مشركين مقترن
منما چه كه جميع آن طلعات رحمن جان در سبيل محبوب داده و براى نان بكلّ من
على الأرض اعتنا نفرموده‌اند حال انصاف ده تو از اهل قياسى با اين عباد كه
چگونه فعلى را بچه فعلى قياس نموده ‌فانصف لتنصف ثمّ ارحم امر الّذى سفك فى
سبيله دمآء المقرّبين و عباد اللّه المخلصين

\end{otherlanguage}

Va-Allāh nālih mīnamāyam va īn kalimāt mīnivīsam chih kih amr bi-jāʾī
rasīdih kih bāyad bi-chunīn adhkār mashghūl shud. Qad ḍayyaʿtum ḥurmat
Allāh wa-ḥurmat anfusikum. Qalb rā az kathāfāt-i awhām-i mushrikīn ṭāhir
namā va dhikr-i anbiyā va mursalīn rā bi-adhkār-i anfus-i mushrikīn
muqtarin manamā chih kih jamīʿ-i ān ṭalaʿāt-i Raḥmān jān dar sabīl-i
Maḥbūb dādih va barā-yi nān bi-kull-i man ʿalā al-arḍ iʿtinā
nafarmūdih-and. Ḥāl inṣāf dih tu az ahl-i qiyāsī bā īn ʿibād kih
chigūnih fiʿlī rā bi-chih fiʿlī qiyās namūdih. Fa-anṣif li-tunṣaf thumma
irḥam amr alladhī sufika fī sabīlihi dimāʾ al-muqarrabīn wa-ʿibād Allāh
al-mukhliṣīn.

By God, I lament as I write these words, for matters have reached a
point where such things must be discussed. You have dishonored God and
yourselves. Purify your heart of the filth of the delusions of those who
associate partners with God, and do not link mention of the Prophets and
Messengers with accounts of those souls who associate partners with Him.
All those Countenances of the All-Merciful gave their lives in the path
of the Beloved and paid no heed to anyone on earth for the sake of
bread. Now judge fairly: you are one who draws comparisons with these
servants---what sort of deed have you compared to what other deed? Be
fair, that you may be treated fairly, and have mercy on the Cause of Him
in whose path the blood of those brought near and of God's sincere
servants has been shed.

\subsection{631}\label{section-1126}

\begin{otherlanguage}{persian}

قوله عظم كذبه عرض ديگر آنكه در باب آمدن حاجى سيّد محمّد در اسلامبول و خودم
همه را مطّلعم و حضور داشتم شب اوّل آمده در سفارت جناب مشير الدّوله بودند و
سركار حاجى ميرزا صفى و بنده و بس و مشير الدّوله باو گفت كه چرا آمده گفت
بجهت استخلاص خود و عرض حال عالى پاشا را هم من ديدم تا آنكه رسولان ايشان
آمدند و آغاز بدگوئى نمودند و او هم بجهت اين بدگوئيها بقدر لازم گفت انتهى

\end{otherlanguage}

Qawluhu, ʿaẓuma kadhibuhu: ʿArḍ-i dīgar ānkih dar bāb-i āmadan-i Ḥājī
Siyyid Muḥammad dar Islambūl va khudam hamih rā muṭṭaliʿam va ḥuḍūr
dāshtam. Shab-i avval āmadih dar sifārat Jināb-i Mushīr al-Dawlih būdand
va Sarkār Ḥājī Mīrzā Ṣafī va bandih va bas va Mushīr al-Dawlih bi-ū guft
kih chirā āmadih guft bi-jahat-i istikhlāṣ-i khud va ʿarḍ-i ḥāl-i ʿĀlī
Pāshā rā ham man dīdam tā ānkih rasūlān-i īshān āmadand va āghāz-i
badgūʾī namūdand va ū ham bi-jahat-i īn badgūʾīhā bi-qadr-i lāzim guft.
Intahā.

His words---great is his falsehood: ``Another point I submit concerns
the arrival of Ḥājī Siyyid Muḥammad in Istanbul. I myself know all about
it and was present. On the first night he came to the embassy, Jināb-i
Mushīr al-Dawlih was there, together with His Honor Ḥājī Mīrzā Ṣafī and
myself, and no one else. Mushīr al-Dawlih asked him why he had come. He
said, `To secure my release.' I also saw the petition to ʿĀlī Pāshā.
Then their messengers came and began to speak ill of him, and in
response to this disparagement he too said as much as was necessary.''
End of quotation.

\subsection{632}\label{section-1127}

\begin{otherlanguage}{persian}

قد سوّد اللّه وجه الكاذب اگر مخصوص استخلاص آمده سياهه نفوس را كه احدى مطّلع
نبوده كه داده هشت نفر بودند بقلم داده بيست و يك نفريم و شش ماه معطّل شده
تا آنكه بهزار تدليس وصل الى مقصوده و مقصود صاحبه باعوا يوسف الأمر بدراهم
معدودة فافّ لهم و لوفائهم تاللّه بفعلهم ضيع امر اللّه بين النّاس

\end{otherlanguage}

Qad sawwada Allāhu wajh al-kādhib. Agar makhṣūṣ-i istikhlāṣ āmadih
siyāhiy-i nufūs rā kih aḥadī muṭṭaliʿ nabūdih kih dādih? Hasht nafar
būdand bi-qalam dādih bīst va yak nafarīm va shash māh muʿaṭṭal shudih
tā ānkih bi-hizār tadlīs waṣala ilā maqṣūdihi wa-maqṣūd ṣāḥibihi. Bāʿū
Yūsuf al-amr bi-darāhima maʿdūdah fa-uffin lahum wa-li-wafāʾihim
ta-Allāh bi-fiʿlihim ḍuyyiʿa amr Allāh bayna al-nās.

God has blackened the face of the liar. If he came solely to secure his
release, who supplied the list of people, of which no one had any
knowledge? There were eight of them, but he submitted in writing, ``We
are twenty-one.'' He spent six months waiting until, through a thousand
deceptions, he attained his own purpose and that of his companion. They
sold the Joseph of the Cause for a few dirhams. Shame on them and their
faithfulness! By God, through their actions the Cause of God was
dishonored among the people.

\subsection{633}\label{section-1128}

\begin{otherlanguage}{persian}

و امرنامه كه از استامبول آمده نوشته‌اند كه بعضى از همرهان ايشان آمده و
شكايت نموده‌اند كه ايشان بما نان نميدهند و الان آن امرنامه موجود مع‌ذلك
نوشته آنچه را نوشته لعن اللّه الّذين بدّلوا حرمة اللّه بزخارف الدّنيا و اذا
خلوا الى الشّياطين يتبرّئون من اللّه و امره ثمّ يكتبون الى الحمقآء من
امثالهم انّا رؤسآء الأمر و مرایاء العباد و اثمار التّوحيد و اوراق التّجريد
و لكن الّذين جعل اللّه بصرهم حديداً يشهدون و يطلعون بما ارتكبوا فى الحيوة
الباطلة و لايشتبه عليهم لانّهم يجدون عن كلّ فعل من افعالهم روايح الكذب و
النّفاق و التّزوير و الشّقاق و انّهم من جواهر الخلق بين عباد اللّه و بريّته و
عليهم ذكر اللّه و ثنائه فى كلّ حين و بعد حين و حين حين

\end{otherlanguage}

Va amr-nāme ke az Estāmbol āmade neveshte-and ke baʿżī az hamrahān-e
īshān āmade va shekāyat namūde-and ke īshān be mā nān nemīdahand va
al-ān ān amr-nāme mawjūd maʿa-dhālek neveshte ānche rā neveshte. Laʿana
Allāhu alladhīna baddalū ḥurmata Allāhi bi-zakhārifi al-dunyā wa idhā
khalaw ilā al-shayāṭīni yatabarraʾūna min Allāhi wa amrihi thumma
yaktubūna ilā al-ḥumaqāʾi min amthālihim innā ruʾasāʾu al-amri wa
marāyāʾu al-ʿibādi wa athmāru al-tawḥīdi wa awrāqu al-tajrīdi wa lākinna
alladhīna jaʿala Allāhu baṣarahum ḥadīdan yashhadūna wa yaṭṭaliʿūna
bi-mā irtakabū fī al-ḥayāti al-bāṭilati wa lā yashtabihu ʿalayhim
li-annahum yajidūna ʿan kulli fiʿlin min afʿālihim rawāʾiḥa al-kadhibi
wa al-nifāqi wa al-tazwīri wa al-shiqāqi wa innahum min jawāhiri
al-khalqi bayna ʿibādi Allāhi wa bariyyatihi wa ʿalayhim dhikru Allāhi
wa thanāʾuhu fī kulli ḥīnin wa baʿda ḥīnin wa ḥīna ḥīn.

In the official order that came from Istanbul, they wrote that some of
his companions had come and complained, ``He does not give us bread.''
That order is still in existence, yet he has written what he has
written. May God curse those who have bartered the honor of God for the
vanities of this world, and who, when alone with the devils, disavow God
and His Cause, then write to fools like themselves, ``We are the leaders
of the Cause, the mirrors of the servants, the fruits of divine unity,
and the leaves of detachment.'' But those whose sight God has made keen
witness and perceive what they have committed in this vain life. They
are not deceived, for in every one of their actions they detect the
odors of falsehood, hypocrisy, deceit, and dissension. They are among
the gems of creation amid the servants of God and His creatures. May
God's remembrance and praise rest upon them at every moment, time after
time, and at all times.

\subsection{634}\label{section-1129}

\begin{otherlanguage}{persian}

و اللّه خوب مريدى هستى و خوب بستر فحشا مشغول شده بعد از شهود و لكن بر
بصير خبير امر پوشيده نيست و مطّلعيم كه چه افتراها گفته و چگونه كمر قتل
بسته حتّى گفته‌اند كه ايشان مخصوص آدم به طهران فرستاده‌اند بجهت امر شنيع
لعنة اللّه على المفترين و الكاذبين

\end{otherlanguage}

Va Allāh khūb morīdī hastī va khūb be satr-e faḥshā mashghūl shode baʿd
az shohūd va lāken bar baṣīr-e khabīr amr pūshīde nīst va moṭṭaleʿīm ke
che efterāhā gofte va chegūne kamar-e qatl baste ḥattā gofte-and ke
īshān makhṣūṣ ādam be Ṭehrān ferestāde-and be jahat-e amr-e shanīʿ.
Laʿnatu Allāhi ʿalā al-muftarīna wa al-kādhibīn.

By God, what a fine disciple you are, and how diligently you have
occupied yourself with concealing indecency after witnessing it! But the
matter is not hidden from the All-Seeing, the All-Informed. We know what
slanders he has uttered and how he has resolved upon murder. They have
even said that He sent someone to Tehran expressly to carry out a
shameful deed. May God's curse be upon the slanderers and liars.

\subsection{635}\label{section-1130}

\begin{otherlanguage}{persian}

حال فعل حقّ را مشاهده كن كه بهمان مدينه وارد شدند و ابدا بكسى مراوده
نفرمودند و در امر اللّه با نفسى مداهنه نشد و چون آن خبيث ديده كه سفير
ايران مكدّر است از ايشان محض خدعه و تلبيس و نفاق با حقّ بآن مقرّ وارد و
التجا نموده‌و گفته آنچه گفته اگر او استخلاص ميخواست كسى باو حرف نداشت او
از حاجى ميرزا احمد معروف‌تر نبوده بهيچوجه و از سايرين كه رفته‌اند معتبرتر
نه بلكه نزد بعضى ذكر نداشته فو‌اللّه كه مخصوص شهريه آمده چنانچه اين ارض هم
عرايضى كه او و ميرزا يحيى بباب مشير نوشته‌اند موجود و

\end{otherlanguage}

Ḥāl feʿl-e Ḥaqq rā moshāhede kon ke be hamān madīne vāred shodand va
abadan be kasī morāvede nafarmūdand va dar amr-e Allāh bā nafsī modāhane
nashod va chon ān khabīth dīde ke safīr-e Īrān mokaddar ast az īshān
maḥż-e khodʿe va talbīs va nefāq bā Ḥaqq be ān maqarr vāred va eltejā
namūde va gofte ānche gofte agar ū estekhlāṣ mīkhāst kasī bā ū ḥarf
nadāsht ū az Ḥājī Mīrzā Aḥmad maʿrūf-tar nabūde be hīch vajh va az
sāyerīn ke rafte-and moʿtabar-tar na balke nazd-e baʿżī dhekr nadāshte
fa-wa-Allāh ke makhṣūṣ-e shahrīye āmade chenānche īn arż ham ʿarāyeżī ke
ū va Mīrzā Yaḥyā be bāb-e moshīr neveshte-and mawjūd va

Now observe the conduct of the True One: He entered that same city and
had no dealings with anyone, and no compromise was made with anyone
concerning the Cause of God. But when that wicked man saw that the
Persian ambassador was displeased with Him, he went to that residence
and sought refuge there, solely to deceive, dissemble, and act
hypocritically toward the True One, and said what he said. If he had
wanted his release, no one would have objected. He was by no means
better known than Ḥājī Mīrzā Aḥmad, nor more esteemed than the others
who left; indeed, some had never even heard of him. By God, he came
expressly for a stipend, as the petitions that he and Mīrzā Yaḥyā wrote
to the office of the Mushīr are also here in this land, and

\subsection{636}\label{section-1131}

\begin{otherlanguage}{persian}

از استامبول جميع امور را از اينجا استفسار نمودند كه او ميگويد هفت نفر
رفته‌اند و باقى بكسب مشغولند و من و ميرزا يحيى بيست و يكنفريم و شهريّه
نداريم سركار والى مخصوص فرستاده و استفسار نموده‌كه باب عالى همچه ذكر
نموده‌اند و تو مينويسى كه بجهت استخلاص بوده حال كه بمنتهى مراد خود رسيدند
ايكاش كه ديگر باطراف ننويسند كه نان نداريم

\end{otherlanguage}

Az Estāmbol jamīʿ-e omūr rā az īnjā estefsār namūdand ke ū mīgūyad haft
nafar rafte-and va bāqī be kasb mashghūland va man va Mīrzā Yaḥyā bīst-o
yek nafarīm va shahrīye nadārīm sarkār-e vālī makhṣūṣ ferestāde va
estefsār namūde ke bāb-e ʿālī hamche dhekr namūde-and va to mīnevīsī ke
be jahat-e estekhlāṣ būde ḥāl ke be montahā-ye morād-e khod resīdand ey
kāsh ke dīgar be aṭrāf nanevīsand ke nān nadārīm.

From Istanbul they inquired here about all these matters, saying, ``He
says that seven people have left and the rest are occupied in earning a
living, while Mīrzā Yaḥyā and I number twenty-one and have no stipend.''
His Excellency the Governor sent someone expressly to inquire, saying
that the Sublime Porte had reported this. Yet you write that it was to
secure his release. Now that they have attained the utmost of their
desire, if only they would stop writing far and wide that they have no
bread!

\subsection{637}\label{section-1132}

\begin{otherlanguage}{persian}

حال انصاف ده آن فعل حقّ و اين فعل اين گروه از كدام رايحه حقّ استشمام ميشود
و اللّه چنين رؤسا از براى امثال شما بسيار خوبند و همچه مدان كه سفير ايران
غافلند لا و اللّه يومى آيد كه خود ايشان تصديق بر فعل حقّ نمايند و منكر فعل
آن نفوس شوند و اگر هم بر تو مشتبه شود بر ايشان مشتبه نخواهد شد

\end{otherlanguage}

Ḥāl enṣāf deh ān feʿl-e Ḥaqq va īn feʿl-e īn gorūh az kodām rāyeḥe-ye
Ḥaqq esteshamām mīshavad va Allāh chonīn roʾasā az barā-ye amthāl-e
shomā besyār khūband va hamche madān ke safīr-e Īrān ghāfeland lā va
Allāh yawmī āyad ke khod-e īshān taṣdīq bar feʿl-e Ḥaqq namāyand va
monker-e feʿl-e ān nofūs shavand va agar ham bar to moshtabeh shavad bar
īshān moshtabeh nakhwāhad shod.

Now judge fairly: that is the conduct of the True One, and this is the
conduct of this group. From which can the fragrance of truth be
perceived? By God, such leaders are very well suited to people like you.
Do not imagine that the Persian ambassador is unaware. No, by God! A day
will come when he himself will affirm the conduct of the True One and
repudiate the conduct of those people. Even if you are deceived, he will
not be.

\subsection{638}\label{section-1133}

\begin{otherlanguage}{persian}

و ديگر رسولان را ندانستم كه را گفته اگر مقصود جناب عبد حاضر لدى الوجه
بوده و اللّه مخصوص امرى آمده در آن مدينه كه هيچ دخلى باو نداشته چنانچه
نزد احدى نرفته جميع شاهد و گواهند كه آن نفس خبيث در اين ارض جميع افعال
شنيعه قبيحه را بحقّ نسبت داده و ذر آن ارض هم گفته‌اند آنچه را كه احدى در
حقّ احدى نگفته عبد حاضر لدى الوجه حال ميگويد خود نفس كاذبت نيامدى و مذكور
ننمودى كه سيّد محمّد چنين و چنان ميگويد و چه مقدار كلمات نالايقه در حقّ
ايشان ذكر نمودى كه او ميگويد حتّى مذكور نمودى كه از براى شهريه خوب نيست
حرف گفته شود هر قسم هست اصلاح شود

\end{otherlanguage}

Va dīgar rasūlān rā nadānestam ke rā gofte agar maqṣūd janāb-e ʿabd-e
ḥāżer ladā al-wajh būde va Allāh makhṣūṣ-e amrī āmade dar ān madīne ke
hīch dakhlī bā ū nadāshte chenānche nazd-e aḥadī narafte jamīʿ shāhed va
govāhand ke ān nafs-e khabīth dar īn arż jamīʿ-e afʿāl-e shanīʿe-ye
qabīḥe rā be Ḥaqq nesbat dāde va {[}corrupted text{]} ān arż ham
gofte-and ānche rā ke aḥadī dar ḥaqq-e aḥadī nagofte ʿabd-e ḥāżer ladā
al-wajh ḥāl mīgūyad khod-e nafs-e kādhebat nayāmadī va madhkūr nanamūdī
ke Siyyid Moḥammad chonīn va chonān mīgūyad va che meqdār kalemāt-e
nālāyeqe dar ḥaqq-e īshān dhekr namūdī ke ū mīgūyad ḥattā madhkūr namūdī
ke az barā-ye shahrīye khūb nīst ḥarf gofte shavad har qesm hast eṣlāḥ
shavad.

As for the messengers, I do not know whom he meant. If he meant His
Honor, the servant present before the Face, by God, he came to that city
for a particular matter that had nothing to do with him, and he visited
no one. All are witnesses and can testify that in this land that wicked
man attributed every shameful and vile deed to the True One, and
{[}corrupted original text{]} that land too they have said what no one
has ever said of anyone. The servant present before the Face now says:
``Did you not come yourself, liar that you are, and report that Siyyid
Muḥammad was saying such and such? How many unseemly statements about
Him did you repeat, saying that he was uttering them! You even said that
it was not good for talk to arise over the stipend and that the matter
should be settled somehow.''

\subsection{639}\label{section-1134}

\begin{otherlanguage}{persian}

و من در جواب گفتم كه معارضه نبوده و نخواهد بود اعراض از آن نفوس لامر
اللّه واقع شده نه لاجل امور ظاهره جميع شهريه را واگذارند بآن نفوس مشركه
جنگى و جدالى نبوده و در مجالس متعدّده از تو خواستم كه در يك مقرّ با او
مجتمع شويم تا بر تو ثابت گردد كه آنچه گفته بافتراى محض نسبت داده آمدى
گفتى كه قبول نميكند

\end{otherlanguage}

Va man dar javāb goftam ke moʿāreże nabūde va nakhwāhad būd eʿrāż az ān
nofūs li-amri Allāh vāqeʿ shode na li-ajl-e omūr-e ẓāhere jamīʿ-e
shahrīye rā vāgod hārand be ān nofūs-e moshreke jangī va jedālī nabūde
va dar majāles-e motaʿaddede az to khwāstam ke dar yek maqarr bā ū
mojtameʿ shavīm tā bar to thābet gardad ke ānche gofte be efterā-ye maḥż
nesbat dāde āmadī goftī ke qabūl nemīkonad.

``I replied that there had been no dispute and would be none. Our
turning away from those people was for the sake of the Cause of God, not
for outward concerns. Let them give the entire stipend to those
idolatrous souls; there has been no quarrel or contention. On several
occasions I asked you to arrange for us to meet with him in one place,
so that it would become clear to you that everything he had said was
sheer slander. You came back and said that he would not agree.''

\subsection{640}\label{section-1135}

\begin{otherlanguage}{persian}

فو‌اللّه كلّ وجود متحيّرند از كذب شما و افتراى ما اى بی‌‌درايت سالها اين ميرزا
يحيى در خدمت ايشان بوده آيا هيچ شنيده كه نفسى از منتسبين ايشان حرف بد و
يا مكروهى نسبت باو گفته باشد هر جاهلى ميداند كه آنچه واقع شود افتضاح امر
بوده و خواهد بود و اللّه كه هميشه فرموده‌اند كه راضى بودم هزار مرتبه شهيد
شوم و اين امور بين ناس ذكر نشود كه يكى از منتسبين ايشان براى نان شكايت
نموده‌و بشهرها عريضه فرستاده لعن اللّه الّذين ارتكبوا ما ناح به كلّ شیء و
احترقت عنه اكباد المقرّبين

\end{otherlanguage}

Fa-wa-Allāh koll-e vojūd motaḥayyerand az kedhb-e shomā va efterā-ye mā
ey bī-derāyat sālhā īn Mīrzā Yaḥyā dar khedmat-e īshān būde āyā hīch
shenīde ke nafsī az montasebīn-e īshān ḥarf-e bad va yā makrūhī nesbat
be ū gofte bāshad har jāhelī mīdānad ke ānche vāqeʿ shavad efteżāḥ-e amr
būde va khwāhad būd va Allāh ke hamīshe farmūde-and ke rāżī būdam hezār
martabe shahīd shavam va īn omūr beyn-e nās dhekr nashavad ke yekī az
montasebīn-e īshān barā-ye nān shekāyat namūde va be shahrhā ʿarīże
ferestāde laʿana Allāhu alladhīna irtakabū mā nāḥa bihi kullu shayʾin wa
iḥtaraqat ʿanhu akbādu al-muqarrabīn.

By God, all existence is bewildered by your lies and our fabrications!
You who lack understanding! For years this Mīrzā Yaḥyā was in His
service. Has he ever heard anyone associated with Him say anything bad
or offensive about him? Even an ignorant person knows that whatever
happens has brought, and will bring, disgrace upon the Cause. By God, He
has always said, ``I would willingly be martyred a thousand times rather
than have these matters spoken of among the people---that one of those
associated with Him has complained about bread and sent petitions to the
cities.'' May God curse those who have committed what has caused all
things to lament and the hearts of those brought near to Him to burn.

\subsection{641}\label{section-1136}

\begin{otherlanguage}{persian}

فو‌اللّه الّذى لا اله الّا هو كه جميع اين عباد راضى بوديم كه كشته شويم و اين
حرفها بميان نيايد آنچه حال ميگوئيم مكلّفيم چه كه مشاهده ميشود كه جمعى كمر
بسته‌اند در هتك حرمت امر اللّه چه كه متديّن بدينى نبوده و نيستند و لاجل
رياست خود را بمهلكه انداخته و مياندازند قريب پنج سنه در اين ارض توقّف شده
و ابدا نفسى بباب عالى نفرستاده‌اند و مطلبی‌ اظهار نشده و آنچه وارد شده
حرمة لأمر اللّه صبر نموده‌ايم و لكن آن خبيث مرتكب شده امرى را كه نفسى از
قبل مرتكب نشده اميدواريم كه بمقصود خود نرسند چنانچه حرمت يك سلسله را
تمام نمود او كه شأنى نداشته و نخواهد داشت

\end{otherlanguage}

Fa-wa-Allāhi alladhī lā ilāha illā huwa ke jamīʿ-e īn ʿebād rāżī būdīm
ke koshte shavīm va īn ḥarfhā be miyān nayāyad ānche ḥāl mīgūyīm
mokallafīm che ke moshāhede mīshavad ke jamʿī kamar baste-and dar hatk-e
ḥormat-e amr-e Allāh che ke motadayyen be dīnī nabūde va nīstand va
li-ajl-e reyāsat khod rā be mahlake andākhte va mīandāzand qarīb-e panj
sane dar īn arż tavaqqof shode va abadan nafsī be bāb-e ʿālī
naferestāde-and va maṭlabī eẓhār nashode va ānche vāred shode ḥurmatan
li-amri Allāh ṣabr namūde-īm va lāken ān khabīth mortakeb shode amrī rā
ke nafsī az qabl mortakeb nashode omīdvārīm ke be maqṣūd-e khod narasand
chenānche ḥormat-e yek selsele rā tamām namūd ū ke shaʾnī nadāshte va
nakhwāhad dāsht.

By God, besides whom there is no god, all of us servants would willingly
have been killed rather than have such things brought up. We are
duty-bound to say what we now say, for we see that a group has resolved
to violate the honor of the Cause of God. They have never adhered to any
religion and still do not, and for the sake of leadership they have cast
themselves, and continue to cast themselves, into destruction. We have
resided in this land for nearly five years, and He has never sent anyone
to the Sublime Porte, nor has any request been made. Whatever has
befallen us, we have borne patiently out of reverence for the Cause of
God. But that wicked man has committed what no one has ever committed
before. We hope they will not attain their aim, for he has utterly
destroyed the honor of an entire community. He himself has never had,
and will never have, any standing.

\subsection{642}\label{section-1137}

\begin{otherlanguage}{persian}

خود انصاف ده اسم او در كجا مذكور بوده اولياى دولت عليه كجا او را
ميشناخته‌اند بسبب كه معروف شده خود تو ميدانى و لكن بغضآء للّه ستر مينمائى
و انت من الّذين غطّيت عيوب المشركين و ما ارتكبوا لتّضييع حرمة اللّه كذلك جرى
الأمر حینئذ من هذا القلم المتكلّم الصّادق الأمين

\end{otherlanguage}

Khod enṣāf deh esm-e ū dar kojā madhkūr būde awliyā-ye dawlat-e ʿalīye
kojā ū rā mīshenākhte-and be sabab-e ke maʿrūf shode khod-e to mīdānī va
lāken bughḍan li-llāh satr mīnamāʾī wa anta mina alladhīna ghaṭṭayta
ʿuyūba al-mushrikīna wa mā irtakabū li-taḍyīʿi ḥurmati Allāhi kadhālika
jarā al-amru ḥīnaʾidhin min hādhā al-qalami al-mutakallimi al-ṣādiqi
al-amīn.

Judge fairly yourself: where had his name ever been mentioned? How had
the officials of the exalted government ever known him? You yourself
know through whom he became known, but out of hatred for God you conceal
it. You are among those who have covered up the faults of the idolaters
and what they have committed to dishonor God. Thus has the matter now
issued from this speaking Pen, the truthful, the trustworthy.

\subsection{643}\label{section-1138}

\begin{otherlanguage}{persian}

قوله كبر حمقه خيلى خيلى عجب است از شما كه قريب بيست سال است ميرزا يحيى
را بخدائى پرستيديد كه خود مطّلعم در اصفهان در نزد خودم اقرار نموديد حال
مردود شده به آقا ميرزا محمّد على و ميرزا ابوالحسن و ديگران و ميرزا حسين
ورزنه چه جواب بگويم كه فلانيها هر روز بيكى ميچسبند و يكى را حقّ ميدانند و
بعد باطل مينمايند اين از عدم تميز آنها است بجهت آنكه حقّ باطل نميشود و
اين تناقض است دين اينها بوالهوسى است چنانچه ميرزا ابوالحسن در طهران بمن
گفت حيف از شما بچه دليل اوّل حقّ بود و خدا و حال باطل است و عبد كافر اى به
به عجب يك انصاف از شما ميپرسم شما را بخدا اين است در بعضى مقامات با شما
در وقتى حركت كرده كه در عرض هم بوديم حالا را هر يك بديگرى نميتوانيم
برترى نمائيم اگر چه هر قدر كامل شده باشيم انتهى

\end{otherlanguage}

Qawluhu kabura ḥumquhu kheylī kheylī ʿajab ast az shomā ke qarīb-e bīst
sāl ast Mīrzā Yaḥyā rā be khodāʾī parastīdīd ke khod moṭṭaleʿam dar
Eṣfahān dar nazd-e khodam eqrār namūdīd ḥāl mardūd shode be Āqā Mīrzā
Moḥammad ʿAlī va Mīrzā Abu'l-Ḥasan va dīgarān va Mīrzā Ḥoseyn-e Varzane
che javāb begūyam ke folānīhā har rūz be yekī mīchasband va yekī rā ḥaqq
mīdānand va baʿd bāṭel mīnamāyand īn az ʿadam-e tamīz-e ānhā ast be
jahat-e ānke ḥaqq bāṭel nemīshavad va īn tanāqoż ast dīn-e īnhā
bū'l-havasī ast chenānche Mīrzā Abu'l-Ḥasan dar Ṭehrān be man goft ḥeyf
az shomā be che dalīl avval ḥaqq būd va khodā va ḥāl bāṭel ast va ʿabd-e
kāfer ey bah bah ʿajab yek enṣāf az shomā mīporsam shomā rā be khodā īn
ast dar baʿżī maqāmāt bā shomā dar vaqtī ḥarakat karde ke dar ʿarż-e ham
būdīm ḥālā rā har yek be dīgarī nemītavānīm bartarī namāʾīm agar che har
qadr kāmel shode bāshīm intahā.

His words---great is his folly: ``It is very, very astonishing that for
nearly twenty years you worshipped Mīrzā Yaḥyā as God. I know this
myself, for you admitted it in my own presence in Isfahan. Now he has
been rejected. What answer should I give Āqā Mīrzā Muḥammad ʿAlī, Mīrzā
Abu'l-Ḥasan, the others, and Mīrzā Ḥusayn of Varzanih, who say that
these people attach themselves to someone new every day, consider
someone to be the truth, and then declare him false? This is due to
their lack of discernment, because truth does not become falsehood, and
this is a contradiction. Their religion is caprice. As Mīrzā Abu'l-Ḥasan
said to me in Tehran, `What a pity for you! On what grounds was he at
first the truth and God, and now false and an unbelieving servant?' Oh,
bravo! How astonishing! I ask you for just one fair judgment, for God's
sake! This is someone who, on certain occasions, associated with you at
a time when we were on an equal footing. At present neither of us can
claim superiority over the other, however much we may have advanced in
perfection.'' End of quotation.

\subsection{644}\label{section-1139}

\begin{otherlanguage}{persian}

اوّلا اين بحث بر شما و امثال شما وارد چه كه سالها نفوسى را پيشوا
ميدانستيد كه اخبث از آن نفوس در ابداع نيامده چنانچه بحكم آن نفوس مظهر
نفس الهى را بتمام ظلم شهيد نمودند و ثانياً اينكه اين عباد باو خلوص
داشته‌ايم و مقرّيم چه كه در كلمات حضرت ابهى بعضى اوصاف مشاهده ميشد و همچه
گمان مينموديم كه مرجع آن اوصاف نفس موهوم است الى ان كشف اللّه لنا ما كشف
و اطّلعنا من افعاله ما لا فعله النّمرود و لا الفرعون و لذا كسّرنا صنم الوهم
و عرفنا ما هو المقصود فى كلماته الابدع البديع ثمّ اطّلعنا به و باصل امره
ما كان مقنوعا عنك و عن مثلائك و انّك لو تعرف لتفرّ منه و تشقّ سبحات الجلال
و تتوجّه الى مشرق الجمال و تكون من الموقنين

\end{otherlanguage}

Avvalan īn baḥth bar shomā va amthāl-e shomā vāred che ke sālhā nofūsī
rā pīshvā mīdānestīd ke akhbath az ān nofūs dar ebdāʿ nayāmade chenānche
be ḥokm-e ān nofūs maẓhar-e nafs-e elāhī rā be tamām-e ẓolm shahīd
namūdand va thāniyan īnke īn ʿebād bā ū kholūṣ dāshte-īm va moqerrīm che
ke dar kalemāt-e Ḥażrat-e Abhā baʿżī awṣāf moshāhede mīshod va hamche
gomān mīnamūdīm ke marjaʿ-e ān awṣāf nafs-e mawhūm ast ilā an kashafa
Allāhu lanā mā kashafa wa iṭṭalaʿnā min afʿālihi mā lā faʿalahu
al-Nimrūdu wa lā al-Firʿawnu wa li-dhā kassar nā ṣanama al-wahmi wa
ʿarafnā mā huwa al-maqṣūdu fī kalimātihi al-abdaʿi al-badīʿi thumma
iṭṭalaʿnā bihi wa bi-aṣli amrihi mā kāna maqnūʿan ʿanka wa ʿan
muthalāʾika wa innaka law taʿrifu la-tafirru minhu wa tashuqqu subuḥāti
al-jalāli wa tatawajjahu ilā mashriqi al-jamāli wa takūnu mina
al-mūqinīn.

First, this objection applies to you and those like you, for years you
regarded as leaders people more wicked than any who have ever appeared
in creation. It was at their command that the Manifestation of the
divine Self was martyred with the utmost cruelty. Second, we admit that
we servants were sincerely devoted to him, for certain descriptions
appeared in the words of His Most Glorious Presence, and we supposed
that those descriptions referred to that imagined person, until God
revealed to us what He revealed, and we learned of deeds of his that
neither Nimrod nor Pharaoh had committed. Therefore we shattered the
idol of delusion and recognized what was intended in His most wondrous
and incomparable words. Then we learned about him and the true nature of
his affair what had been concealed from you and those like you. If you
knew, you would flee from him, tear through the veils of glory, turn
toward the Dawning-Place of Beauty, and be among those who are certain.

\subsection{645}\label{section-1140}

\begin{otherlanguage}{persian}

حمد ميكنيم خدا را كه ببازوى توكّل و انقطاع صنم وهم را شكستيم و بشريعه
يقين وارد شديم و كاش تو هم ميبودى و ببصر خود مشاهده مينمودى نحمد اللّه
بذلك و نشكره فى كلّ حين فو‌اللّه اگر ذرّه انصاف در تو ميبود هرآينه بايد ليل
و نهار باوصاف اين عباد متكلّم شوى چه كه فعل ما بر دين ما گواهى است صادق و
شاهدى است واضح و آشكار چه كه سلاسل تقليد را شكستيم و بحقّ فايز گشتيم و از
وهميكه جميع قلوب را اخذ نموده ‌بود فارغ شديم اى فقير بی‌‌بضاعت بسيار بسيار
در وهم مستغرقى فيا ليت رايته و عرفته و اطلعت بما عنده كاش ملاقات شده بود
و از عرفانيكه در مجالس ذكر مينمود مذكور ميامد و مطّلع ميشدى از آنچه اين
ایّام از او محتجبی‌

\end{otherlanguage}

Ḥamd mīkonīm Khodā rā ke be bāzū-ye tavakkol va enqeṭāʿ ṣanam-e vahm rā
shekastīm va be sharīʿe-ye yaqīn vāred shodīm va kāsh to ham mībūdī va
be baṣar-e khod moshāhede mīnamūdī naḥmadu Allāha bi-dhālika wa
nashkuruhu fī kulli ḥīn fa-wa-Allāh agar dharre enṣāf dar to mībūd
har-āyene bāyad leyl va nahār be awṣāf-e īn ʿebād motakallem shavī che
ke feʿl-e mā bar dīn-e mā govāhī ast ṣādeq va shāhedī ast vāżeḥ va
āshkār che ke salāsel-e taqlīd rā shekastīm va be Ḥaqq fāyez gashtīm va
az vahmī ke jamīʿ-e qolūb rā akhdh namūde būd fāregh shodīm ey faqīr-e
bī-beżāʿat besyār besyār dar vahm mostaghraqī fa-yā layta raʾaytahu wa
ʿaraftahu wa iṭṭalaʿta bi-mā ʿindahu kāsh molāqāt shode būd va az
ʿerfānī ke dar majāles dhekr mīnamūd madhkūr mīāmad va moṭṭaleʿ mīshodī
az ānche īn ayyām az ū moḥtajebī.

We praise God that, through the strength of trust in Him and detachment,
we shattered the idol of delusion and reached the waters of certainty.
Would that you too had been there and seen with your own eyes! We praise
God for this and thank Him at all times. By God, if you had a particle
of fairness, you would surely speak day and night of the qualities of
these servants, for our conduct is a truthful testimony to our faith and
a clear and evident witness. We broke the chains of imitation, attained
the True One, and freed ourselves from the delusion that had seized
every heart. You poor, destitute one! You are deeply, deeply immersed in
delusion. If only you had seen him, known him, and learned what he
possessed! Would that a meeting had taken place, that the spiritual
knowledge of which he spoke in gatherings had been discussed, and that
you had learned what remains veiled from you concerning him in these
days.

\subsection{646}\label{section-1141}

\begin{otherlanguage}{arabic}

و انّك ان لاتعرفنا نحن نعرفك انفسنا لتطلع و تكون من العالمين نحن عباد
الّذين جعل اللّه ابصارنا حديداً و عرفناه بنفسه و انقطعنا عن العالمين و ما
منعنا منع مانع و لا مكر ماكر و لاخدع خادع و لاريب قلوب المغلّين نحمد اللّه
بما عرفنا نفسه الّذى كان موعوداً فى الواحه المقدّس المحكم البديع و ما منعنا
عن عرفانه حجبات الّذينهم كفروا باللّه ثمّ سبحات الّذينهم اشركوا بنفسه الواحد
الفرد العزيز الحميد

\end{otherlanguage}

Wa-innaka in lā taʿrifunā naḥnu nuʿarrifuka anfusanā li-taṭṭaliʿa
wa-takūna mina al-ʿālimīn. Naḥnu ʿibād alladhīna jaʿala Allāhu abṣāranā
ḥadīdan wa-ʿarafnāhu bi-nafsihi wa-inqaṭaʿnā ʿani al-ʿālamīn. Wa-mā
manaʿanā manʿu māniʿ wa-lā makru mākir wa-lā khidʿu khādiʿ wa-lā raybu
qulūbi al-mughillīn. Naḥmadu Allāha bi-mā ʿarrafanā nafsahu alladhī kāna
mawʿūdan fī alwāḥihi al-muqaddasi al-muḥkami al-badīʿ. Wa-mā manaʿanā
ʿan ʿirfānihi ḥujubātu alladhīnahum kafarū bi-Allāh thumma subuḥātu
alladhīnahum ashrakū bi-nafsihi al-Wāḥidi al-Fardi al-ʿAzīzi al-Ḥamīd.

If you do not know us, we will make ourselves known to you, so that you
may understand and be among those who know. We are servants whose sight
God has made keen; we have recognized Him through Himself and detached
ourselves from all the worlds. No obstruction by anyone, no schemer's
plotting, no deceiver's deception, and no doubt in the hearts of the
malicious have held us back. We praise God for making known to us His
own Self, who was promised in His holy, firmly established, wondrous
Tablets. Neither the veils of those who disbelieved in God nor the veils
of glory of those who associated partners with His Self, the One, the
Unique, the Almighty, the All-Praised, prevented us from recognizing
Him.

\subsection{647}\label{section-1142}

\begin{otherlanguage}{persian}

عجب است كه باينگونه تحقيقات مشغول شده اذا ظهر الشّطّ هل يليق لأحد ان يلتفت
الى مآء الغدير اذا فانصف يا ايّها الحمير و اذا اشرقت الشّمس هل ينبغى لأحد
ان يستضئ بالسّراج و لو كان سراجاً حقيقيّاً فانصف يا ايّها الغافل الّذى بغفلتك
ناح كلّ شیء و لو انّك انت فى حجبات غليظ و اذا تموّج بحر الأعظم هل يتوجّه
العاقل الى سراب بقيعة لا فو‌ربّنا الرّحمن الرّحيم و لكن انت لن تجد ما نذكر
لانّ قلبك صار محروماً من نفحات اللّه المقتدر العلىّ الحكيم

\end{otherlanguage}

ʿAjab ast kih bi-īn-gūnih taḥqīqāt mashghūl shudih. Idhā ẓahara
al-shaṭṭu hal yalīqu li-aḥadin an yaltafita ilā māʾi al-ghadīr? Idhan
fa-anṣif yā ayyuhā al-ḥamīr. Wa-idhā ashraqati al-shamsu hal yanbaghī
li-aḥadin an yastaḍīʾa bi-al-sirāj wa-law kāna sirājan ḥaqīqiyyan?
Fa-anṣif yā ayyuhā al-ghāfil alladhī bi-ghaflatika nāḥa kullu shayʾ
wa-law innaka anta fī ḥujubātin ghalīẓ. Wa-idhā tamawwaja baḥru al-Aʿẓam
hal yatawajjahu al-ʿāqilu ilā sarābin bi-qīʿah? Lā fa-wa-Rabbinā
al-Raḥmāni al-Raḥīm. Wa-lākin anta lan tajida mā nadhkuru li-anna
qalbaka ṣāra maḥrūman min nafaḥāti Allāhi al-Muqtadiri al-ʿAliyyi
al-Ḥakīm.

It is astonishing that you have occupied yourself with investigations of
this kind. When the great river has appeared, is it fitting for anyone
to turn toward the water of a pool? Be fair, then, you ass! When the sun
has risen, should anyone seek illumination from a lamp, even if it is a
real lamp? Be fair, heedless one, through whose heedlessness all things
have lamented, though you remain within thick veils. When the Most Great
Ocean surges, would a rational person turn toward a mirage on a plain?
No, by our Lord, the All-Merciful, the Ever-Merciful! Yet you will never
grasp what We recount, for your heart has been deprived of the breezes
of God, the All-Powerful, the Most High, the All-Wise.

\subsection{648}\label{section-1143}

\begin{otherlanguage}{persian}

و اللّه نسوان ننگ دارند از ادلّه تو شيطان در هر آسمانى باسمى معروف و معلّم
ملكوت بوده و زياده از الف سنه بكلّ اسماء موسوم و بعد يكمرتبه رجيم شده اين
امور از بدايع قدرت الهى است هل تنكرها و انّك لو تنكر يصدقها عباده
الموقنين

\end{otherlanguage}

Va-Allāh nisvān nang dārand az adillih-yi tu. Shayṭān dar har āsmānī
bi-ismī maʿrūf va muʿallim-i malakūt būdih va ziyādih az alf sanih
bi-kull-i asmāʾ mawsūm va baʿd yik-martabih rajīm shudih. Īn umūr az
badāyiʿ-i qudrat-i ilāhī ast. Hal tunkiruhā wa-innaka law tunkir
yuṣaddiquhā ʿibāduhu al-mūqinīn.

By God, women are ashamed of your arguments. Satan was known by a name
in every heaven and was a teacher of the Kingdom. For more than a
thousand years he bore every name, and then, all at once, he was cast
out. These things are among the wonders of divine power. Do you deny
them? Even if you deny them, His servants who are certain affirm them.

\subsection{649}\label{section-1144}

\begin{otherlanguage}{persian}

آيا نشنيده يهوداى اسخريوطى يكى از حروفات اثنا عشريّه انجيل بوده و بعد
بشقاوتى ظاهر كه كلمة اللّه را تسليم يهود نمود و كاش رئيس شما بفعل او
ارتكاب مينمود چه كه او بعد از آن عمل شنيع نادم شده و لكن مرشدت نقطه بيان
و كلّ انبیاء و مرسلين و حرمة اللّه جميع را بزخارف ارض مبادله نموده‌ و
فروخته چنانچه ذكر شد و تو بجميع مطّلعى

\end{otherlanguage}

Āyā nashanīdih Yahūdā-yi Iskharīyūṭī yikī az ḥurūfāt-i
ithnā-ʿasharīyyih-yi Injīl būdih va baʿd bi-shiqāvatī ẓāhir kih Kalimat
Allāh rā taslīm-i Yahūd namūd? Va kāsh raʾīs-i shumā bi-fiʿl-i ū irtikāb
mīnamūd, chih kih ū baʿd az ān ʿamal-i shanīʿ nādim shudih, va lākin
murshidat Nuqṭih-yi Bayān va kull-i anbiyā va mursalīn va ḥurmat-i Allāh
jamīʿ rā bi-zakhārif-i arḍ mubādilih namūdih va furūkhtih, chunānchih
dhikr shud, va tu bi-jamīʿ muṭṭaliʿī.

Have you not heard that Judas Iscariot was one of the twelve Letters of
the Gospel, and afterward displayed such wickedness that he delivered
the Word of God into the hands of the Jews? Would that your leader had
committed his deed, for Judas repented after that heinous act. Your
guide, however, has exchanged and sold the Point of the Bayán, all the
Prophets and Messengers, and the sanctity of God for earthly vanities,
as has been recounted, and you are aware of it all.

\subsection{650}\label{section-1145}

\begin{otherlanguage}{persian}

و ديگر در اسمعيل بن حضرت صادق چه ميگوئى بعد از آنكه بجميع ولايات نوشته
كه او بعد از من امام است و پيشواى خلق بعد از مدّتى امورى از او ظاهر شد و
حضرت او را عزل فرمودند كه مخصوص زراده و بعضى از اصحاب عرض كردند يا ابن
رسول اللّه اين امر بزرگى است چگونه ميشود شخصى امام خلق باشد و حقّ و بعد
باطل شود حضرت ذكر بدا فرمود برو و روايت كلينى را ببين

\end{otherlanguage}

Va dīgar dar Ismaʿīl ibn-i Ḥaḍrat-i Ṣādiq chih mīgūʾī? Baʿd az ānkih
bi-jamīʿ-i vilāyāt nivishtih kih ū baʿd az man imām ast va pīshvā-yi
khalq, baʿd az muddatī umūrī az ū ẓāhir shud va Ḥaḍrat ū rā ʿazl
farmūdand, kih makhṣūṣ Zarādih va baʿḍī az aṣḥāb ʿarḍ kardand: Yā ibn-i
Rasūl-i Allāh, īn amr-i buzurgī ast, chigūnih mīshavad shakhṣī imām-i
khalq bāshad va ḥaqq va baʿd bāṭil shavad? Ḥaḍrat dhikr-i badā farmūd.
Biraw va rivāyat-i Kulaynī rā bibīn.

And what do you say about Ismaʿil, the son of His Holiness Ṣádiq? After
he had written to all the provinces that Ismaʿil would be the Imam and
leader of the people after him, certain things appeared from Ismaʿil
after a time, and His Holiness dismissed him. In particular, Zarádih and
some of the companions submitted: ``O son of the Messenger of God, this
is a weighty matter. How can someone be the Imam of the people and be
true, and then become false?'' His Holiness spoke of badá. Go and read
Kulayní's account.

\subsection{651}\label{section-1146}

\begin{otherlanguage}{persian}

اى فقير از هيچ امرى خبر ندارى نه از بيان و نه از فرقان و نه از تورية و
نه از انجيل و در ظهور اين امر انقلابات بسيار شده و اين عبد نظر باختصار
ذكر ننمود و در كتب خود شما مذكور است چنانچه حال هم بعضى از اطراف اسمعيل
را امام ميدانند كتب ظاهره را هم نديده و مطّلع نيستى شان شما اينست كه
بحرفهاى ّ تكلّم نمائيد و باينگونه ادلّه مستدلّ شويد بر بصير خبير مشتبه نشده
و نميشود و بسا از امور كه واقع شده و احدى مطّلع نه چه كه حقّ از قبل ستر
فرموده ايكاش كه از هوى ميگذشتى و چندى للّه در كوى دوست مقرّ ميگزيدى تا
بعضى از امور مستوره بر تو كشف ميشد

\end{otherlanguage}

Ay faqīr, az hīch amrī khabar nadārī, nah az Bayān va nah az Furqān va
nah az Tawrāt va nah az Injīl. Va dar ẓuhūr-i īn amr inqilābāt-i bisyār
shudih va īn ʿabd naẓar bi-ikhtiṣār dhikr nanamūd, va dar kutub-i khud-i
shumā madhkūr ast, chunānchih ḥāl ham baʿḍī az aṭrāf Ismaʿīl rā imām
mīdānand. Kutub-i ẓāhirih rā ham nadīdih va muṭṭaliʿ nīstī. Shaʾn-i
shumā īnast kih bi-ḥarfhā-yi {[}corrupted text{]} takallum namāyīd va
bi-īn-gūnih adillih mustadill shavīd. Bar baṣīr-i khabīr mushtabih
nashudih va namīshavad. Va basā az umūr kih vāqiʿ shudih va aḥadī
muṭṭaliʿ nah, chih kih Ḥaqq az qabl satr farmūdih. Ay kāsh kih az havā
mīgudhashtī va chandī li-llāh dar kūy-i Dūst maqarr mīguzīdī tā baʿḍī az
umūr-i mastūrih bar tu kashf mīshud.

You poor man, you know nothing of any matter---neither of the Bayán, nor
of the Qur'an, nor of the Torah, nor of the Gospel. Many reversals have
occurred in the unfolding of this matter, and this servant has not
recounted them for the sake of brevity. They are recorded in your own
books; even now, people in some regions regard Ismaʿil as the Imam. You
have not even seen the readily available books and are uninformed. Your
way is to utter words of {[}corrupted original text{]} and to argue with
proofs of this kind. The discerning and well-informed have not been
misled, nor will they be. Many things have happened of which no one is
aware, for God concealed them in the past. If only you would abandon
your desires and dwell for a time, for God's sake, in the Beloved's
neighborhood, so that some hidden matters might be disclosed to you.

\subsection{652}\label{section-1147}

\begin{otherlanguage}{persian}

و ديگر از بلعم باعور اطّلاع ندارى كه بسيار از رئيس شما مشهورتر بود ولايتش
كه مسلّم بوده مع‌ذلك بعد از ظهور شمس كليم من افق ربّك الرّحمن الرّحيم فعل ما
فعل و بذلك بدّل اللّه نوره بالنّار و انّه لعلى كلّ شیء قدير كاش از رئيس خود
اطّلاع ميداشتى فو‌اللّه الّذى لا اله الّا هو انّ يد القدرة و القوّة قد اخذ كفّا
من التّراب و نفخ فيه روحاً من ارواح الّتى خلقت بمشيّته فلمّا اهتزّ و قام البسه
بقميص الأسماء بين الأرض و السّمآء لحكمة الّتى ما اطّلع بها احد الّا نفسه
العليم الخبير زياده از اين نبوده و نخواهد بود فلمّا عاد على اللّه عاد اللّه
عليه كما عاد كذلك يجزى عباده المعتدين

\end{otherlanguage}

Va dīgar az Balʿam-i Bāʿūr iṭṭilāʿ nadārī kih bisyār az raʾīs-i shumā
mashhūrtar būd. Vilāyatash kih musallam būdih, maʿa-dhālik baʿd az
ẓuhūr-i shams-i Kalīm min ufuqi Rabbika al-Raḥmāni al-Raḥīm faʿala mā
faʿala wa-bi-dhālika baddala Allāhu nūrahu bi-al-nār wa-innahu la-ʿalā
kulli shayʾin qadīr. Kāsh az raʾīs-i khud iṭṭilāʿ mīdāshtī. Fa-wa-Allāhi
alladhī lā ilāha illā huwa inna yada al-qudrati wa-al-quwwati qad
akhadha kaffan mina al-turāb wa-nafakha fīhi rūḥan min arwāḥi allatī
khuliqat bi-mashīʾatihi fa-lammā ihtazza wa-qāma albasahu bi-qamīṣi
al-asmāʾ bayna al-arḍi wa-al-samāʾ li-ḥikmati allatī mā iṭṭalaʿa bihā
aḥad illā nafsuhu al-ʿAlīmu al-Khabīr. Ziyādih az īn nabūdih va
nakhvāhad būd. Fa-lammā ʿāda ʿalā Allāh ʿāda Allāhu ʿalayhi ka-mā ʿād.
Kadhālika yajzī ʿibādahu al-muʿtadīn.

Nor do you know about Balaam son of Beor, who was far more renowned than
your leader. His spiritual authority was acknowledged; nevertheless,
after the sun of Moses rose from the horizon of your Lord, the
All-Merciful, the Ever-Merciful, he did what he did, and through this
God changed his light into fire. He indeed has power over all things.
Would that you knew about your own leader! By God, besides whom there is
no other god, the hand of power and might took a handful of dust and
breathed into it a spirit from among the spirits created by His will.
When it stirred and stood upright, He clothed it in the garment of names
between earth and heaven, for a wisdom known to no one but His own Self,
the All-Knowing, the All-Informed. It has never been, nor will it ever
be, more than this. Then, when he turned against God, God turned against
him as he had turned. Thus does He recompense His transgressing
servants.

\subsection{653}\label{section-1148}

\begin{otherlanguage}{persian}

بشنو سخن اين عبد را سبب وهم ناس مشو رحم كن بر امر اللّه و برناس در هر عصر
نفوس متوهّمه مجتمع شده و ناس را از شاطى بحر احديّه محروم نموده‌اند قدرى
تفكّر نما كه چه مقدار توهّمات در ما بين ملأ فرقان بود از ظهور قائم و ظهور
قيامت و ظهور ساعت و بعد از ظهور نقطه اولى روح ما سواه فداه معلوم شد كه
جميع خاطى بوده‌اند و بقطره از بحر علم مشروب نه و آن جمال قدم بزحمت تمام
كسر اصنام وهم و هوى نموده ‌تازه امثال آن نفوس اراده نموده‌اند ناس را
باوهام جديد مبتلا نمايند بگو بترسيد از خداوند يكتا و مرتكب مشويد امري را
كه از چشم مقرّبين بجاى اشك خون جارى شود حقّ را باطل مدانيد و باطل را حقّ
مشمريد بلى حقّ لم‌يزل حقّ بوده و لكن تو آن حقّ را نشناخته

\end{otherlanguage}

Bishnaw sukhan-i īn ʿabd rā, sabab-i vahm-i nās mashaw. Raḥm kun bar
Amr-i Allāh va bar nās. Dar har ʿaṣr nufūs-i mutavahhimih mujtamiʿ
shudih va nās rā az shāṭī-yi baḥr-i aḥadīyyih maḥrūm namūdih-and. Qadrī
tafakkur namā kih chih miqdār tavahhumāt dar mā bayn-i malaʾ-i Furqān
būd az ẓuhūr-i Qāʾim va ẓuhūr-i qiyāmat va ẓuhūr-i sāʿat, va baʿd az
ẓuhūr-i Nuqṭih-yi Ūlā, rūḥ-i mā sivāh fidāh, maʿlūm shud kih jamīʿ khāṭī
būdih-and va bi-qaṭrih az baḥr-i ʿilm mashrūb nah. Va ān Jamāl-i Qidam
bi-zaḥmat-i tamām kasr-i aṣnām-i vahm va havā namūdih, tāzih amthāl-i ān
nufūs irādih namūdih-and nās rā bi-avhām-i jadīd mubtalā namāyand. Bigū
bitarsīd az Khudāvand-i yiktā va murtakib mashavīd amrī rā kih az
chashm-i muqarrabīn bi-jā-yi ashk khūn jārī shavad. Ḥaqq rā bāṭil
madānīd va bāṭil rā ḥaqq mashmurīd. Balī Ḥaqq lam-yazal ḥaqq būdih va
lākin tu ān Ḥaqq rā nashnākhtih.

Listen to the words of this servant, and do not cause people to fall
into delusion. Have mercy on the Cause of God and on the people. In
every age, deluded souls have gathered together and deprived people of
the shore of the ocean of oneness. Reflect a little on how many
misconceptions existed among the people of the Qur'an concerning the
appearance of the Qá'im, the advent of the Resurrection, and the coming
of the Hour. After the appearance of the Primal Point---may the souls of
all besides Him be sacrificed for Him---it became clear that they had
all been mistaken and had not drunk a single drop from the ocean of
knowledge. That Ancient Beauty shattered the idols of delusion and
desire with immense difficulty, and now others like those souls have set
out to afflict people with new delusions. Say: Fear the one God, and do
not commit a deed that would cause blood rather than tears to flow from
the eyes of those near to Him. Do not regard truth as falsehood, or
reckon falsehood as truth. Yes, the True One has always been true, but
you have not recognized that True One.

\subsection{654}\label{section-1149}

\begin{otherlanguage}{persian}

و اين مسلّم بوده كه بعد از نقطه حقّ منحصر است در حروفات حىّ چه كه كلّ خلق
بيان در ظلّ واحد اوّل خلق شده و ميشوند و حقّ هر ذى حقّى باين نفوس ثابت و
هم‌چنين كلّ صفات و اسماء حسنى و مع‌ذلك خطاب ميفرمايد بيكى از حروفات حىّ كه
از ظهور بعد سؤال نموده ‌ميفرمايد اگر بدانم كه تو در يوم اشراق نيّر معانى
بايمان باو فايز نميشوى حال حكم دون ايمان در حقّ تو جارى مينمايم و اگر
بدانم يكى از نصارى در يوم ظهور مؤمن باو ميشود حال حكم ايمان درباره او
ميفرمودم مع آنكه مؤمن نيست و فعل محبوبی‌ از او ظاهر نشده قدرى تفكر نما و
حقّ را بشناس

\end{otherlanguage}

Va īn musallam būdih kih baʿd az Nuqṭih ḥaqq munḥaṣir ast dar Ḥurūfāt-i
Ḥayy, chih kih kull-i khalq-i Bayān dar ẓill-i Vāḥid-i Avval khalq
shudih va mīshavand, va ḥaqq-i har dhī-ḥaqqī bi-īn nufūs thābit va
ham-chunīn kull-i ṣifāt va asmāʾ-i ḥusnā. Va maʿa-dhālik khiṭāb
mīfarmāyad bi-yikī az Ḥurūfāt-i Ḥayy kih az ẓuhūr-i baʿd suʾāl namūdih,
mīfarmāyad: Agar bidānam kih tu dar yawm-i ishrāq-i nayyir-i maʿānī
bi-īmān bi-ū fāʾiz namīshavī, ḥāl ḥukm-i dūn-i īmān dar ḥaqq-i tu jārī
mīnamāyam, va agar bidānam yikī az Naṣārā dar yawm-i ẓuhūr muʾmin bi-ū
mīshavad, ḥāl ḥukm-i īmān dar bārih-yi ū mīfarmūdam, maʿa ānkih muʾmin
nīst va fiʿl-i maḥbūbī az ū ẓāhir nashudih. Qadrī tafakkur namā va Ḥaqq
rā bishnās.

It has been acknowledged that, after the Point, truth is confined to the
Letters of the Living, for the entire creation of the Bayán has been and
continues to be created beneath the shadow of the First Váḥid. The right
of everyone possessing a right is established through these souls, as
are all the attributes and most beautiful names. Nevertheless,
addressing one of the Letters of the Living who had asked about the next
Revelation, He says: ``If I know that, on the day when the Luminary of
inner meanings rises, you will not attain faith in Him, I pronounce upon
you now the judgment of unbelief. And if I know that one of the
Christians will believe in Him on the day of His Revelation, I would
pronounce upon him now the judgment of faith, even though he is not a
believer and no pleasing deed has appeared from him.'' Reflect a little
and recognize the True One.

\subsection{655}\label{section-1150}

\begin{otherlanguage}{persian}

اين نفوس حقّند ماداميكه در ظلّ حقّ باشند چه كه در آنوقت ديده نميشود در آن
نفوس مگر حقّ و بعد از توقّف دام الملك فى الملك و رجع المخلوق الى شكله و
مثله آن حقّيكه لم‌يزل باطل نشده و نخواهد شد آن شجره حقيقت است فاعرف لو
تكون من العارفين كذلك فصّل الأمر من لدن منزل قديم تفكّر يا عبد لعلّ تكون من
المهتدين

\end{otherlanguage}

Īn nufūs ḥaqq-and mādāmīkih dar ẓill-i Ḥaqq bāshand, chih kih dar ān
vaqt dīdih namīshavad dar ān nufūs magar Ḥaqq. Va baʿd az tavaqquf dāma
al-mulku fī al-mulk wa-rajaʿa al-makhlūqu ilā shaklihi wa-mithlihi. Ān
Ḥaqqīkih lam-yazal bāṭil nashudih va nakhvāhad shud ān Shajarih-yi
Ḥaqīqat ast. Fa-iʿrif law takūnu mina al-ʿārifīn. Kadhālika fuṣṣila
al-amru min ladun Munzilin Qadīm. Tafakkar yā ʿabd laʿalla takūnu mina
al-muhtadīn.

These souls are true as long as they remain beneath the shadow of the
True One, for at that time nothing is seen in them but the True One.
After they hesitate, sovereignty remains within sovereignty, and the
creature returns to its own kind and likeness. That True One who has
never become false, nor ever will, is the Tree of Reality. Recognize
this, if you are among those who know. Thus has the matter been
expounded by the Ancient Revealer. Reflect, O servant, so that you may
be among the rightly guided.

\subsection{656}\label{section-1151}

\begin{otherlanguage}{persian}

اى غافل غير بالغ بدانكه در بعث مراتب ما لانهايه بوده و خواهد بود و بعث
اسماء را از دونش تميز ده بسا نفوس كه اسماء مبعوثند من غير رسم ذلك لحكمة
لن يعرفها احد و لو تريد ان تعرفها فاسئل عن اللّه ربّك بخضوع و خشوع و انابة
مبين ليعلّمك ما عنده و انّه لهو العليم الخبير در الواح قبل در اينمقام از
قلم امر فى الجمله نازل شده اگر اراده نمائى در آن الواح نظر كن شايد بهمان
قدر قناعت نمائى و مقصود از اين ذكر آنكه باسماء از منزل آن محتجب نمانى
اگر چه بشأنى محتجب مانده كه ابداً گمان نميرود كه خرق شود الّا بارادة اللّه
و مشيّته

\end{otherlanguage}

Ay ghāfil-i ghayr-i bāligh, bidānkih dar baʿth marātib-i mā lā nihāyih
būdih va khvāhad būd va baʿth-i asmāʾ rā az dūnash tamyīz dih. Basā
nufūs kih asmāʾ mabʿūthand min ghayri rasm. Dhālika li-ḥikmatin lan
yaʿrifahā aḥad. Wa-law turīdu an taʿrifahā fa-isʾal ʿan Allāhi Rabbika
bi-khuḍūʿin wa-khushūʿin wa-inābatin mubīn li-yuʿallimaka mā ʿindahu
wa-innahu la-huwa al-ʿAlīmu al-Khabīr. Dar alvāḥ-i qabl dar īn maqām az
Qalam-i Amr fī al-jumlih nāzil shudih. Agar irādih namāʾī dar ān alvāḥ
naẓar kun, shāyad bi-hamān qadr qanāʿat namāʾī. Va maqṣūd az īn dhikr
ānkih bi-asmāʾ az Munzil-i ān muḥtajib namānī, agar-chih bi-shaʾnī
muḥtajib māndih kih abadan gumān namīravad kih kharq shavad illā
bi-irādati Allāhi wa-mashīʾatihi.

O heedless, immature one! Know that resurrection has had and will have
limitless degrees, and distinguish the resurrection of names from other
kinds. Many are the souls whose names have been raised without any
outward trace. This is for a wisdom that no one will know. If you wish
to know it, ask God your Lord with humility, reverence, and manifest
repentance, so that He may teach you what is with Him. He is indeed the
All-Knowing, the All-Informed. In earlier Tablets, something on this
subject has been revealed by the Pen of the Cause. If you wish, consult
those Tablets; perhaps you will be content with that much. The purpose
of this mention is that you should not remain veiled by the names from
the One who reveals them, although you have remained so veiled that it
is inconceivable that the veil should ever be torn away except through
God's purpose and will.

\subsection{657}\label{section-1152}

\begin{otherlanguage}{persian}

نقطه بيان جلّت عظمته در تبديل كينونات هم ذكر فرموده‌اند تا چه رسد باسماء
چنانچه بعضى از آن از قبل ذكر شد و هم‌چنين ميفرمايد قوله جلّ ذكره فانّ عند
كلّ ظهور ادلّآء اثباتها يبدّلنّ بالنّفى و لكن اللّه فى ذلك الظّهور يثبت امره و
يتقن صنعه لانّ فيه لايبدّلنّ النّور بالنّار الّا فى ظهور اللّه فانّ حزنى على كلّ
من فى البيان من يؤمئذ انتهى

\end{otherlanguage}

Nuqṭih-yi Bayān, jallat ʿaẓamatuhu, dar tabdīl-i kaynūnāt ham dhikr
farmūdih-and tā chih rasad bi-asmāʾ, chunānchih baʿḍī az ān az qabl
dhikr shud. Va ham-chunīn mīfarmāyad, qawluhu jalla dhikruhu: Fa-inna
ʿinda kulli ẓuhūrin adillāʾa ithbātihā yubaddalunna bi-al-nafy
wa-lākinna Allāha fī dhālika al-ẓuhūri yuthbitu amrahu wa-yutqinu
ṣunʿahu li-anna fīhi lā yubaddalunna al-nūru bi-al-nār illā fī ẓuhūri
Allāh fa-inna ḥuznī ʿalā kulli man fī al-Bayān min {[}corrupted text{]}.
Intahā.

The Point of the Bayán---exalted be His grandeur---has spoken even of
the transformation of essential realities, let alone names, as some of
this has already been mentioned. He also says---His words, exalted be
His mention: ``At every Revelation, the exponents of its affirmation are
changed into negation. Yet in that Revelation God establishes His Cause
and perfects His handiwork, for in it light is not changed into fire
except at the Revelation of God. My sorrow is indeed for everyone in the
Bayán from {[}corrupted original text{]}.'' End of quotation.

\subsection{658}\label{section-1153}

\begin{otherlanguage}{persian}

ملاحظه نمائيد كه صريح ميفرمايند نور مبدّل بنار ميشود چنانچه نورى كه شما
آن را نور فرض گرفته‌ايد كينونتش بنار تبديل شده و لكن شاعر نيست و در كمال
غفلت مشغول بتحريف كلمات الهيّه بوده و هست و اينكه ميفرمايند الّا فى ظهور
اللّه همچه ندانيد كه قبل از ظهور تبديل نور نميشود بلى فو‌نفسه الحقّ غير از
نفس ظهور كلّ تبديل شده و ميشوند و اين مسائل موهومه نزد ناس بوده و تا اين
ظهور كشف نشده فلمّا خرقتم حجاب حرمة اللّه لذا خرق ستر حجابكم بأنامل القدرة
و القوّة

\end{otherlanguage}

Mulāḥiẓih namāyīd kih ṣarīḥ mīfarmāyand nūr mubaddal bi-nār mīshavad,
chunānchih nūrī kih shumā ān rā nūr farḍ giriftih-īd kaynūnatash bi-nār
tabdīl shudih va lākin shāʿir nīst, va dar kamāl-i ghaflat mashghūl
bi-taḥrīf-i kalimāt-i ilāhīyyih būdih va hast. Va īnkih mīfarmāyand illā
fī ẓuhūri Allāh hamchih nadānīd kih qabl az ẓuhūr tabdīl-i nūr
namīshavad. Balī fa-wa-nafsihi al-Ḥaqq ghayr az nafs-i Ẓuhūr kull tabdīl
shudih va mīshavand. Va īn masāʾil-i mawhūmih nazd-i nās būdih va tā īn
ẓuhūr kashf nashudih. Fa-lammā kharaqtum ḥijāba ḥurmati Allāh li-dhā
kharaqa sitra ḥijābikum bi-anāmili al-qudrati wa-al-quwwah.

Observe that He explicitly says that light is changed into fire. Thus
the light you have assumed to be light has had its essential reality
changed into fire, yet he is unaware, and in utter heedlessness has been
and remains occupied with distorting the divine words. Do not understand
His statement, ``except at the Revelation of God,'' to mean that light
is not transformed before the Revelation. Yes, by His own Self, the True
One, all except the Manifestation Himself have been and continue to be
transformed. These imaginary notions existed among the people, and the
matter had not been disclosed until this Revelation. But when you tore
through the veil of God's sanctity, He therefore tore away the covering
of your veil with the fingers of power and might.

\subsection{659}\label{section-1154}

\begin{otherlanguage}{persian}

فرق اينست كه در حين ارتقاء مظهر احديّه برفيق اعلى الى ظهور بعد مشيّت ظاهره
مشهود نه تا حكم نمايد اگر چه لازال انّه هو حىّ فى الأفق الأبهی و لكن چون
بر حسب ظاهر احدى بر او عارف نه و كلّ نفوس از لقا ممنوع لذا حكم تبديل
مستور تا حين ظهور چه كه در آن حين اگر نور اقلّ من آن توقّف نمايد در قول
بلى حكم نار از منبع علم لدنّى الهى در حقّ او جارى چنانچه در ظهور قبل
مشاهده شد و در اين ظهور هم ببصر ظاهر و باطن ديده ميشود كه نفوسى كه خود
را نور ميناميدند عند اللّه حكم نار در حقّ آن نفوس جارى و لكن خود من غير
شعور حركت مينمايند

\end{otherlanguage}

Farq īnast kih dar ḥīn-i irtiqā-yi Maẓhar-i Aḥadīyyih bi-Rafīq-i Aʿlā
ilā ẓuhūr-i baʿd Mashīyyat-i Ẓāhirih mashhūd nah tā ḥukm namāyad,
agar-chih lā-zāl innahu huwa ḥayyun fī al-ufuq-i al-Abhā. Va lākin chūn
bar ḥasab-i ẓāhir aḥadī bar ū ʿārif nah va kull-i nufūs az liqā mamnūʿ,
li-dhā ḥukm-i tabdīl mastūr tā ḥīn-i ẓuhūr, chih kih dar ān ḥīn agar nūr
aqall min ān tavaqquf namāyad dar qawl-i balā ḥukm-i nār az manbaʿ-i
ʿilm-i ladunnī-yi ilāhī dar ḥaqq-i ū jārī, chunānchih dar ẓuhūr-i qabl
mushāhidih shud va dar īn ẓuhūr ham bi-baṣar-i ẓāhir va bāṭin dīdih
mīshavad kih nufūsī kih khud rā nūr mīnāmīdand ʿinda Allāh ḥukm-i nār
dar ḥaqq-i ān nufūs jārī va lākin khud min ghayr-i shuʿūr ḥarakat
mīnamāyand.

The difference is that, from the time when the Manifestation of Oneness
ascends to the Companion on high until the next Revelation, the manifest
Will is not visibly present to pronounce judgment, although He is ever
alive in the All-Glorious Horizon. But since outwardly no one knows Him
and all souls are denied access to His presence, the judgment of
transformation remains concealed until the time of the Revelation. At
that time, if a light hesitates for less than an instant to say ``Yes,''
the judgment of fire is pronounced upon it from the source of innate
divine knowledge. This was witnessed in the previous Revelation, and in
this Revelation too it is seen with both the outward and the inward eye
that souls who called themselves light are judged as fire in God's
sight, yet they themselves go about unaware.

\subsection{660}\label{section-1155}

\begin{otherlanguage}{persian}

در كلّ حين پناه بخدا برده كه از نفس ظهور محتجب نمانى چه كه مقصود بالذّات
او بوده و خواهد بود او است آيه ليس كمثله شیء و او است آيه لم يلد و لم
يولد بل مظاهر لم يلد و لم يولد خلق عنده ان انتم توقنون و لكن اينمقام را
هر نفسى ادراك ننمايد الّا المنقطعون و هم فى عرفان منيع

\end{otherlanguage}

Dar koll-e ḥīn panāh be-Khodā borde ke az nafs-e ẓohūr moḥtajeb namānī,
che ke maqṣūd be-dh-dhāt ū būde va khwāhad būd. Ū ast āye-ye laysa
ka-mithlihi shayʾ, va ū ast āye-ye lam yalid wa-lam yūlad. Bal maẓāhir-i
lam yalid wa-lam yūlad khalqun ʿindahu in antum tūqinūn. Va lāken īn
maqām rā har nafsī edrāk nanamāyad illā al-munqaṭiʿūn wa-hum fī ʿirfānin
manīʿ.

At every moment, seek refuge in God, that you may not remain veiled from
the Manifestation Himself, for He has always been and will ever be the
One sought in His own right. He is the verse, ``Nothing is like Him,''
and He is the verse, ``He neither begets nor is begotten.'' Indeed, the
manifestations of ``He neither begets nor is begotten'' are created
beings in His sight, if you are certain. Yet not every soul can
comprehend this station, but only those who are detached, who dwell in
exalted understanding.

\subsection{661}\label{section-1156}

\begin{otherlanguage}{persian}

از اين جهت است كه نقطه مشيّت اوّلیّه روح ما سواه فداه در مقام ذكر حروفات و
مرايا و نور و امثالها من اعلى مراتب الأسماء و الصّفات الى ان ينتهى الى
ادنى رتبة الأشیآء ميفرمايند اينها از خود تحقّقى نداشته چون بشمس حقيقى
مقابل شده‌اند اشراق تجلّيات لانهايه بر آنها شده در اينمقام ديده نميشود در
آن مظاهر الّا اللّه ميفرمايد شما ناظر باسماء نباشيد بل بما يقوم به الأسماء
ناظر شويد و هم‌چنين ما يقوم به الصّفات چه كه كلّ باقبالهم الى اللّه و
بايمانهم باللّه باثواب بديعه منيعه اسماء مزيّن شده و ميشوند جمعى از حروفات
ناميده ميشوند و جمعى از مرايا و جمعى از ادلّاء نوريّه وهم‌چنين از مظاهر
جمال و جلال و حىّ و قديم و امثال آن از اسماء اللّه الحسنى تختلف ظهورات
التّجلّى باعتباراتهم و استعداداتهم و الّا انّ التّجلّى واحد و المجلّى واحد

\end{otherlanguage}

Az īn jahat ast ke Noqṭe-ye Mashiyyat-e Avvaliyye, rūḥ-e mā sevāh fedāh,
dar maqām-e dhekr-e ḥorūfāt va marāyā va nūr va amthāl-e ānhā, min aʿlā
marātib al-asmāʾ wa-l-ṣifāt ilā an yantahī ilā adnā rutbat al-ashyāʾ,
mīfarmāyand īnhā az khod taḥaqqoqī nadāshte, chon be-shams-e ḥaqīqī
moqābel shode-and eshrāq-e tajalliyāt-e lā-nehāye bar ānhā shode. Dar īn
maqām dīde nemīshavad dar ān maẓāher illā Allāh. Mīfarmāyad shomā nāẓer
be-asmā nabāshīd bal be-mā yaqūmu bihi al-asmāʾ nāẓer shavīd va
ham-chenīn mā yaqūmu bihi al-ṣifāt, che ke koll bi-iqbālihim ilā Allāh
va bi-īmānihim bi-Allāh be-athvāb-e badīʿe-ye manīʿe-ye asmā mozayyan
shode va mīshavand. Jamʿī az ḥorūfāt nāmīde mīshavand va jamʿī az marāyā
va jamʿī az adellā-ye nūriyye va ham-chenīn az maẓāher-e jamāl va jalāl
va ḥayy va qadīm va amthāl-e ān az asmāʾ Allāh al-ḥusnā. Takhtalifu
ẓuhūrāt al-tajallī bi-iʿtibārātihim wa-istiʿdādātihim wa-illā inna
al-tajallī wāḥid wa-l-mujallī wāḥid.

For this reason, the Point of the Primal Will---may the spirit of all
besides Him be sacrificed for Him---when speaking of the Letters, the
Mirrors, the Light, and the like, from the highest degrees of the names
and attributes down to the lowest rank of things, states that these had
no independent reality of their own. When they faced the true Sun, the
radiance of infinite revelations shone upon them. In this station,
nothing is seen in those manifestations but God. He says that you should
look not to the names, but to that through which the names subsist, and
likewise to that through which the attributes subsist. For all, through
their turning to God and their faith in God, have been and are adorned
with the wondrous, exalted robes of the names. Some are called Letters,
some Mirrors, and some luminous Guides, and likewise manifestations of
Beauty, Majesty, the Living, the Ancient, and other such most beautiful
names of God. The manifestations of the revelation differ according to
their respective conditions and capacities; otherwise, the revelation is
one, and the One who reveals it is one.

\subsection{662}\label{section-1157}

\begin{otherlanguage}{persian}

و بصير خبير ادراك مينمايد كه كلّ اين اسماء نظر بآن ايمانى است كه در آنها
بوده و هست وهم را بگذار ملاحظه در اوّل من آمن كن و هم‌چنين ثانى من آمن كه
هيچ نفسى را بيان ازيد از اين دو نفس و ساير حروفات حىّ وصف نفرموده‌اند
چنانچه در زيارتشان يا ساذج اللّه و يا كينونة اللّه و يا ذات اللّه و بامثال
اين كلمات ذكر فرموده‌اند حال اگر اوّل من آمن نفس مذكور نبود و نفس ديگر بود
البتّه اينمقام باو عنايت ميشد اگر در اين بحر متغمّس شوى ادراك مينمائى بعضى
از اسرار محجّبه مقنّعه مستوره را و فرق شمس معانى را از ظلّش تميز ميدهى و
مرتكب نميشوى آنچه را كه كلّ ذرّات بنوحه مشغولند

\end{otherlanguage}

Va baṣīr-e khabīr edrāk mīnamāyad ke koll-e īn asmā naẓar be-ān īmānī
ast ke dar ānhā būde va hast. Vahm rā begoẕār, molāḥeẓe dar avval-e man
āmana kon va ham-chenīn thānī-ye man āmana, ke hīch nafsī rā Bayān azyad
az īn do nafs va sāyer-e Ḥorūfāt-e Ḥayy vaṣf nafarmūde-and, chenānche
dar ziyārat-e shān yā sādhij Allāh va yā kaynūnat Allāh va yā dhāt Allāh
va be-amthāl-e īn kalemāt dhekr farmūde-and. Ḥāl agar avval-e man āmana
nafs-e madhkūr nabūd va nafs-e dīgar būd albatte īn maqām be-ū ʿenāyat
mīshod. Agar dar īn baḥr motaghamm es shavī edrāk mīnamāʾī baʿẓī az
asrār-e moḥajjabe-ye moqannaʿe-ye mastūre rā va farq-e shams-e maʿānī rā
az ẓell-ash tamīz mīdehī va mortakeb nemīshavī ānche rā ke koll-e
dharrāt be-nawḥe mashghūl-and.

One who is discerning and well informed understands that all these names
refer to the faith that has existed and continues to exist within them.
Set aside illusion and consider the first to believe, and likewise the
second to believe, for the Bayán has praised no souls more highly than
these two and the other Letters of the Living. Thus, in their visitation
prayers, He has addressed them as ``O pure Essence of God,'' ``O Being
of God,'' and ``O Essence of God,'' and in similar words. Now, if the
first to believe had not been the person in question, but someone else,
this station would assuredly have been bestowed upon that other person.
If you immerse yourself in this ocean, you will comprehend some of the
veiled, shrouded, hidden mysteries, distinguish the Sun of inner
meanings from its shadow, and refrain from committing that over which
every atom is lamenting.

\subsection{663}\label{section-1158}

\begin{otherlanguage}{persian}

مثلا ملاحظه كن اگر ترابی‌ بيد مبارك اخذ فرمايد و بفرمايد كه اين قبضه خاك
مسجود عالمين است حقّ لاريب فيه و هم‌چنين اگر صد قبضه خاك ديگر بردارد و
درباره كلّ همين حكم را فرمايد جارى بوده و خواهد بود چه كه در اينمقام نظر
بامر است كه از مصدر امر ظاهر شده نه بشىء من حيث هو شیء

\end{otherlanguage}

Mathalan molāḥeẓe kon agar torābī be-yad-e mobārak akhdh farmāyad va
befarmāyad ke īn qabẓe khāk masjūd-e ʿālamīn ast, ḥaqqun lā rayba fīh.
Va ham-chenīn agar ṣad qabẓe khāk-e dīgar bardārad va dar bāre-ye koll
hamīn ḥokm rā farmāyad jārī būde va khwāhad būd, che ke dar īn maqām
naẓar be-amr ast ke az maṣdar-e amr ẓāher shode na be-shayʾ min ḥaythu
huwa shayʾ.

Consider, for example: if He takes some dust in His blessed hand and
declares that this handful of earth is the object before which all
worlds bow down, it is true beyond any doubt. Likewise, if He takes a
hundred other handfuls of earth and pronounces the same decree
concerning them all, it has been and will remain in force. For in this
station, regard is given to the command that has issued from the Source
of command, not to a thing considered merely as a thing.

\subsection{664}\label{section-1159}

\begin{otherlanguage}{persian}

و اينكه امثال آن نفوس محتجبه بذكر وصايت از ما يقوم به الوصىّ و فوقه من
الولاية و النّبوّة و الرّسالة و الرّبوبيّة غافل شده‌اند اوّلا بدان فو‌اللّه الّذى
لا اله الّا هو كه ابداً نقطه اولى بچنين كلمه مخصوص نفسى در بيان تكلّم
نفرموده بلى يك توقيعى بميرزا سليمان التفات فرموده بودند در صدر آن توقيع
كلمه بود كه جميع از خواندن آن عاجز شدند خود ميرزا سليمان در ایّام مراجعت
از ماكو ميگفت اين لفظ وصايت است و بمن نوشته‌اند و او در سرّ سرّ اين ادّعاى
اين رتبه مينمود و ايكاش شما باو متشبّث ميشديد در اين امر اوعندالله
سزاوارتر است از اين نفسى كه باو تمسّك جسته‌ايد و ديگر همچه ذكرى ابداً نبوده
مگر آنكه بعضى الواح نوشته‌اند و باسم نقطه اولى شهرت داده‌اند كلّها كذب بربّ
البهآء

\end{otherlanguage}

Va īnke amthāl-e ān nofūs-e moḥtajebe be-dhekr-e veṣāyat az mā yaqūmu
bihi al-waṣī wa-fawqahu min al-wilāya wa-l-nubuwwa wa-l-risāla
wa-l-rubūbiyya ghāfel shode-and, avvalan bedān fa-wa-Allāh alladhī lā
ilāha illā huwa ke abadan Noqṭe-ye Ūlā be-chenīn kaleme makhṣūṣ-e nafsī
dar Bayān takallom nafarmūde. Balī yek tawqīʿī be-Mīrzā Solaymān eltefāt
farmūde būdand. Dar ṣadr-e ān tawqīʿ kaleme būd ke jamīʿ az khwāndan-e
ān ʿājez shodand. Khod-e Mīrzā Solaymān dar ayyām-e morājaʿat az Mākū
mīgoft īn lafẓ veṣāyat ast va be-man neveshte-and va ū dar serr-e serr
īn eddeʿā-ye īn rotbe mīnamūd. Va ey kāsh shomā be-ū motashabbeth
mīshodīd dar īn amr. Ū ʿinda Allāh sazāvārtar ast az īn nafsī ke be-ū
tamassok joste-īd. Va dīgar hamche dhekrī abadan nabūde magar ānke baʿẓī
alvāḥ neveshte-and va be-nām-e Noqṭe-ye Ūlā shohrat dāde-and. Kulluhā
kadhib bi-Rabb al-Bahāʾ.

As for such veiled souls becoming heedless, through mention of
successorship, of that through which the successor subsists, and of the
higher degrees of guardianship, prophethood, messengership, and
lordship: first, know---by God, besides whom there is no god---that the
Primal Point never spoke such a word in the Bayán specifically
concerning any individual. He did, however, graciously send a Tablet to
Mírzá Sulaymán. At the beginning of that Tablet was a word that everyone
was unable to read. Mírzá Sulaymán himself, during his return from Mákú,
would say, ``This word is `successorship,' and He has written it to
me.'' In the utmost secrecy, he laid claim to this rank. Would that you
had held fast to him in this matter! In God's sight, he is more
deserving than the person to whom you have attached yourselves. There
has never been any other such mention, except that some Tablets have
been written and circulated in the name of the Primal Point. They are
all false, by the Lord of Bahá!

\subsection{665}\label{section-1160}

\begin{otherlanguage}{persian}

و ثانيا اين رتبه رتبه نيست كه بخواهيم از نفسى سلب نمائيم يا اثبات نمائيم
چه كه آنچه فخر وصى است ان ما يقوم به الوصايه بوده و خواهد بود و بعد از
انقطاع عند اللّه ذكر نداشته و نخواهد داشت چنانچه نقطهء اولى روح ما سواه
فداه در اينمقام ميفرمايد قوله عزّ ذكره زيرا كه جوهر كلّ همان است كه آن
محتجب مانده از چيزيكه وصىّ بآن وصىّ است نه اينست كه اين شأن وصايت را هم
بشأن خاصّ نظر كنى بلكه ناظر شو بما يقوم به كه امر شجره حقيقت باشد كه اگر
فرضا رسول خدا باسم وصايت نصب نميفرمود امروز ذكر بغير وصايت ميشد كلّ اينها
از براى اين است كه در يوم ظهور باسماء محتجب نمانى بلكه نظر كنى بما يقوم
به الأسماء من كلّ شیء حتّى ذكر النّبی‌ انتهى

\end{otherlanguage}

Va thāniyan īn rotbe rotbe nīst ke bekhwāhīm az nafsī salb namāʾīm yā
ethbāt namāʾīm, che ke ānche fakhr-e vaṣī ast inna mā yaqūmu bihi
al-waṣāya būde va khwāhad būd va baʿd az enqeṭāʿ ʿinda Allāh dhekr
nadāshte va nakhwāhad dāsht. Chenānche Noqṭe-ye Ūlā rūḥ-e mā sevāh fedāh
dar īn maqām mīfarmāyad, qawluhu ʿazza dhikruhu: zīrā ke jawhar-e koll
hamān ast ke ān moḥtajeb mānde az chīzī ke vaṣī be-ān vaṣī ast. Na īn
ast ke īn shaʾn-e veṣāyat rā ham be-shaʾn-e khāṣṣ naẓar konī balke nāẓer
shaw be-mā yaqūmu bihi ke amr-e Shajare-ye Ḥaqīqat bāshad, ke agar
farẓan Rasūl-e Khodā be-nām-e veṣāyat naṣb nemīfarmūd emrūz dhekr
be-ghayr-e veṣāyat mīshod. Koll-e īnhā az barā-ye īn ast ke dar yawm-e
ẓohūr be-asmā moḥtajeb namānī balke naẓar konī be-mā yaqūmu bihi
al-asmāʾ min kulli shayʾ ḥattā dhikr al-Nabī. Intahā.

Second, this rank is not a rank that We would wish either to take away
from someone or to affirm for someone. For the glory of the successor
has always been and will ever be that through which successorship
subsists; once severed from it, he has had and will have no standing
before God. Thus the Primal Point---may the spirit of all besides Him be
sacrificed for Him---states in this connection, His words, exalted be
His mention: ``For the essence of the whole matter is precisely this: he
has remained veiled from that by which the successor is a successor. It
is not that you should regard even this station of successorship as a
special distinction; rather, look to that through which it subsists,
which is the command of the Tree of Truth. For if, hypothetically, the
Messenger of God had not appointed him under the name of successorship,
today he would be spoken of by a name other than successorship. All this
is so that, on the Day of Manifestation, you may not remain veiled by
names, but may instead look to that through which the names of all
things subsist, even the designation `Prophet.'\,'' End of quotation.

\subsection{666}\label{section-1161}

\begin{otherlanguage}{persian}

حال قدرى تفكّر نما تا عارف شوى كه كلّ بامر او معزّزند و بذكر او مفتخر و اين
شأن صادق است ماداميكه از حدّ خود تجاوز ننمايند و در حين بعث امر اوّلیّه كلّ
اين اسامى از مسمّيات خود خارج و بملكوت راجع اگر ثانيا امر مبرم نازل يصدق
عليهم كلّ الأسماء من اسماء اللّه الحسنى و الّا يرجعن الى مقرّهم قبل ان يقول
لم او بم

\end{otherlanguage}

Ḥāl qadrī tafakkor namā tā ʿāref shavī ke koll be-amr-e ū moʿazzaz-and
va be-dhekr-e ū moftakher va īn shaʾn ṣādeq ast mādāmī ke az ḥadd-e khod
tajāvoz nanamāyand. Va dar ḥīn-e baʿth-e amr-e avvaliyye koll-e īn asāmī
az mosammayāt-e khod khārej va be-malakūt rājeʿ. Agar thāniyan amr-e
mobram nāzel, yaṣduqu ʿalayhim kull al-asmāʾ min asmāʾ Allāh al-ḥusnā,
wa-illā yarjiʿna ilā maqarrihim qabla an yaqūla lima aw bima.

Now reflect a little, that you may recognize that all are exalted
through His command and honored through His mention. This condition
holds true so long as they do not overstep their bounds. At the raising
up of the primal Cause, all these names are withdrawn from those who
bear them and return to the Kingdom. If a binding command descends anew,
all the most beautiful names of God apply to them; otherwise, they
return to their abode before one can say ``Why?'' or ``By what?''

\subsection{667}\label{section-1162}

\begin{otherlanguage}{persian}

بارى نظر از اين تحديدات بردار و بآن نقطه ساريه در كلّشیء ناظر شو و آن
نقطه را موهوم مدان چنانچه از قبل كلّ عباد بتوهّمات مانده و الى حين عارف
نشده و نخواهند شد الّا من شآء ربّك و همين ميرزا يحيى كه بواسطه او از حقّ
معرض شده بين يدى مثل يكى از خدّام حاضر بوده و القاى كلمات اللّه بر او ميشد
چنانچه حال جمعى تلقاء وجه حاضرند بهمان قسم بر ايشان القاء ميشود در كلّ
ليالى و ایّام و لكن فرق است كه آن نفس مغرور چون دو كلمه از لسانش جارى شد
ارتكب ما لا ارتكب مظاهر النّفى كلّها و لكن اين نفوس حاضره تلقاء وجه بابدع
كلمات ناطق و مع‌ذلك ليس لهم من امر الّا بعد اذنه كذلك كان الأمر ان انت من
العالمين

\end{otherlanguage}

Bārī naẓar az īn taḥdīdāt bardār va be-ān noqṭe-ye sāriye dar koll-e
shayʾ nāẓer shaw va ān noqṭe rā mawhūm madān, chenānche az qabl koll-e
ʿebād be-tavahhomāt mānde va ilā ḥīn ʿāref nashode va nakhwāhand shod
illā man shāʾa Rabbuka. Va hamīn Mīrzā Yaḥyā ke be-vāseṭe-ye ū az Ḥaqq
moʿreẓ shode, bayna yaday mithl-e yekī az khoddām ḥāẓer būde va elqā-ye
kalemāt-e Allāh bar ū mīshod, chenānche ḥāl jamʿī telqāʾ-e vajh
ḥāẓer-and, be-hamān qesm bar īshān elqā mīshavad dar koll-e layālī va
ayyām. Va lāken farq ast ke ān nafs-e maghrūr chon do kaleme az
lesān-ash jārī shod irtakaba mā lā irtakaba maẓāhir al-nafy kulluhā, va
lāken īn nofūs-e ḥāẓere telqāʾ-e vajh be-abdaʿ-e kalemāt nāṭeq va maʿa
dhālik laysa lahum min amrin illā baʿda idhnihi. Kadhālika kāna al-amr
in anta min al-ʿālimīn.

In short, turn your gaze away from these limitations and look to that
Point which pervades all things. Do not regard that Point as imaginary,
as all the servants of former times remained caught in imaginings and,
to this day, have not attained recognition and never will, except those
whom your Lord wills. This very Mírzá Yaḥyá, on whose account you have
turned away from the True One, used to stand in My presence like one of
the attendants, and the words of God were imparted to him, just as a
number of people now stand before My face and receive them in the same
manner every night and day. The difference, however, is that when a mere
two words flowed from the tongue of that deluded soul, he committed what
all the manifestations of negation had never committed. These souls now
present before My face, by contrast, utter the most wondrous words, yet
undertake nothing without His permission. Such has been the matter, if
you are among those who know.

\subsection{668}\label{section-1163}

\begin{otherlanguage}{persian}

اگرچه كلمات مشركين در قلبت مقرّ گرفته و محو آن مشكل مگر آنكه از اين
سلسبيل عنايت كه در عيون كلمات بنفس كلمات جارى و سارى است بياشامى كذلك
نبّأناك و اخبرناك و قصصناك ما كنت محتجباً عنه لعلّ تخرق الأحجاب بقدرة ربّك
العزيز الوهّاب و تشهد ما لا شهدته من قبل و تعرف ما كنت غافلاً عنه و لكن لن
توفّق بذلك الّا بأن تنقطع عما خلق بين السّموات و الأرض و تتوجّه الى اللّه
الملك المقتدر العزيز المختار

\end{otherlanguage}

Agarche kalemāt-e moshrekīn dar qalb-at maqarr gerefte va maḥv-e ān
moshkel magar ānke az īn salsabīl-e ʿenāyat ke dar ʿoyūn-e kalemāt
be-nafs-e kalemāt jārī va sārī ast biyāshāmī. Kadhālika nabbaʾnāka
wa-akhbarnāka wa-qaṣaṣnāka mā kunta muḥtajiban ʿanhu laʿalla takhruqa
al-aḥjāb bi-qudrati Rabbika al-ʿAzīz al-Wahhāb wa-tashhada mā lā
shahidtahu min qablu wa-taʿrifa mā kunta ghāfilan ʿanhu. Wa-lākin lan
tuwaffaqa bi-dhālika illā bi-an tanqaṭiʿa ʿammā khuliqa bayna al-samāwāt
wa-l-arḍ wa-tatawajjaha ilā Allāh al-Malik al-Muqtadir al-ʿAzīz
al-Mukhtār.

Although the words of those who associate partners with God have lodged
in your heart, and it is difficult to efface them unless you drink from
this Salsabíl of loving-kindness that flows and courses through the
wellsprings of the words by means of the words themselves. Thus have We
told you, informed you, and recounted to you what you were veiled from,
that you may rend the veils through the power of your Lord, the
Almighty, the All-Bountiful, witness what you have never witnessed
before, and recognize what you were heedless of. Yet you will not
succeed in this unless you detach yourself from all that has been
created between the heavens and the earth and turn toward God, the King,
the All-Powerful, the Almighty, the One who chooses.

\subsection{669}\label{section-1164}

\begin{otherlanguage}{persian}

اى عباد در عرفان اين ظهور بكلمات و اشارات معرضين و تصديق و تكذيب آن نفوس
متمسّك نشويد و باصل امر ناظر شويد چه كه اكثرى از ناس لم‌يزل و لايزال محتجب
و متوهّم مشاهده شده و ميشوند و رؤساى عصر نظر باستحكام رياست خود ناس را از
شطر احديّه و عرفان آيات اللّه ممنوع و مأيوس نموده ‌و خواهند نمود مثلاً
ملاحظه در ظهور نقطه فرقان نما كه آنحضرت بآيات اللّه ظاهر شدند و در فرقان
حجّت و برهان جز آيات نازل نشده مع‌ذلك بعد از ارتقاى نقطه فرقان برفيق اعلى
و انقطاع روايح رحمن امر بدست علماء افتاد و آن نفوس بالمرّه ناس را از
عرفان آيات منع نموده‌كه شما آيات اللّه را ادراك نمينمائيد و علم آن نزد ما
است باينجهت حجّيت آيات از ميان رفت و زمام خلق بدست آن نفوس آمد و ناس را
باوهامى تربيت نمودند كه ذكر آن ممكن نه چنانچه ديده و شنيده‌ايد و آياتيكه
در ظهور بعد و قيامت و امثال آن نازل شده جميع را بهواى خود تفسير نمودند و
ناس هم متابعت آن نفوس را من حيث لايشعر نموده ‌تا آنكه امر بمقامى رسيد كه
ضرّ آن اوهام جميع بشجره ظهور بعد وارد شد و

\end{otherlanguage}

Ey ʿebād, dar ʿerfān-e īn ẓohūr be-kalemāt va eshārāt-e moʿreẓīn va
taṣdīq va takdhīb-e ān nofūs motamassek nashavīd va be-aṣl-e amr nāẓer
shavīd, che ke aktharī az nās lam-yazal va lā-yazāl moḥtajeb va
motavahhem moshāhede shode va mīshavand va roʾasā-ye ʿaṣr naẓar
be-esteḥkām-e reyāsat-e khod nās rā az shaṭr-e aḥadiyye va ʿerfān-e
āyāt-e Allāh mamnūʿ va maʾyūs namūde va khwāhand namūd. Mathalan
molāḥeẓe dar ẓohūr-e Noqṭe-ye Forqān namā ke ān Ḥaẓrat be-āyāt-e Allāh
ẓāher shodand va dar Forqān ḥojjat va borhān joz āyāt nāzel nashode.
Maʿa dhālik baʿd az erteqā-ye Noqṭe-ye Forqān be-Rafīq-e Aʿlā va
enqeṭāʿ-e ravāyeḥ-e Raḥmān amr be-dast-e ʿolamā oftād va ān nofūs
be-l-marre nās rā az ʿerfān-e āyāt manʿ namūde ke shomā āyāt-e Allāh rā
edrāk nemīnamāʾīd va ʿelm-e ān nazd-e mā ast. Be-īn jahat ḥojjiyyat-e
āyāt az miyān raft va zemām-e khalq be-dast-e ān nofūs āmad va nās rā
be-awhāmī tarbiyat namūdand ke dhekr-e ān momken na, chenānche dīde va
shenīde-īd. Va āyātī ke dar ẓohūr-e baʿd va qiyāmat va amthāl-e ān nāzel
shode jamīʿ rā be-havā-ye khod tafsīr namūdand va nās ham motābaʿat-e ān
nofūs rā min ḥaythu lā yashʿur namūde tā ānke amr be-maqāmī resīd ke
ẓorr-e ān awhām jamīʿ be-Shajare-ye Ẓohūr-e baʿd vāred shod va

O servants! In seeking to recognize this Revelation, do not cling to the
words and allusions of those who turn away, or to those souls'
affirmations and denials. Look instead to the foundation of the Cause
itself, for most people have always been and continue to be seen veiled
and deluded. The leaders of the age, intent on consolidating their own
authority, have barred people from the realm of oneness and the
recognition of God's verses, and have deprived them of hope, and will
continue to do so. Consider, for example, the Revelation of the Point of
the Qur'an: He appeared with the verses of God, and in the Qur'an no
proof or testimony was revealed other than the verses. Nevertheless,
after the Point of the Qur'an ascended to the Companion on high and the
fragrances of the All-Merciful ceased, the matter fell into the hands of
the learned. Those souls completely barred people from recognizing the
verses, saying, ``You cannot comprehend the verses of God; knowledge of
them rests with us.'' Thus the evidentiary authority of the verses
disappeared, the reins of the people passed into those souls' hands, and
they instructed people in imaginings that cannot be recounted, as you
have seen and heard. They interpreted all the verses revealed concerning
the next Manifestation, the Resurrection, and the like according to
their own desires. The people, too, followed those souls unawares, until
matters reached the point where all the harm of those imaginings befell
the Tree of the next Revelation, and

\subsection{670}\label{section-1165}

\begin{otherlanguage}{persian}

بعد از ظهور موعود در سنه ستّين مجدّدا ارياح فضل وزيد و سموات علم مرتفع شد
و ارض عرفان منبسط و آيات اللّه از سماء مشيّت نازل و لكن نظر باوهام قبل ناس
در حجيّت آيات توقّف نموده‌كه ما ادراك نمينمائيم و حال آنكه حقّ جلّ ذكره اعظم
و اكرم و ارحم از آن است كه مظهر نفس خود را بفرستد و جميع را بعرفان او كه
نفس عرفان اللّه است امر فرمايد و مع‌ذلك حجّتى بفرستد فوق ادراك ناس سبحانه
عن ذلك

\end{otherlanguage}

Baʿd az ẓohūr-e Mawʿūd dar sane-ye settīn mojaddadan aryāḥ-e faẓl vazīd
va samavāt-e ʿelm mortafeʿ shod va arẓ-e ʿerfān monbaseṭ va āyāt-e Allāh
az samā-ye mashiyyat nāzel. Va lāken naẓar be-awhām-e qabl nās dar
ḥojjiyyat-e āyāt tavaqqof namūde ke mā edrāk nemīnamāʾīm, va ḥāl ānke
Ḥaqq jalla dhikruhu aʿẓam va akram va arḥam az ān ast ke maẓhar-e nafs-e
khod rā beferestad va jamīʿ rā be-ʿerfān-e ū ke nafs-e ʿerfān-e Allāh
ast amr farmāyad va maʿa dhālik ḥojjatī beferestad fawq-e edrāk-e nās.
Subḥānahu ʿan dhālik.

After the appearance of the Promised One in the year sixty, the winds of
grace blew anew, the heavens of knowledge were raised, the earth of
recognition was spread out, and the verses of God descended from the
heaven of His will. Yet, because of their former imaginings, people
hesitated to acknowledge the evidentiary authority of the verses,
saying, ``We do not comprehend them.'' The True One---exalted be His
mention---is, however, far too great, generous, and merciful to send the
Manifestation of His own Self, command everyone to recognize Him, which
is itself the recognition of God, and yet send a proof beyond people's
comprehension. Far exalted is He above such a thing!

\subsection{671}\label{section-1166}

\begin{otherlanguage}{persian}

بالأخرة باوهام انفس خود مشغول و محجوب بشأنيكه بآيات اللّه كفايت ننموده‌مع
آنكه نصّ فرقان است او لم يكفهم انّا انزلنا عليك الكتاب و جميع مستدلّ باينكه
ادراك آيات فوق عقول و عرفان ما است امر ديگر بايد ظاهر شود قد غيّروا نعمة
اللّه و حجّته و ارادوا بغير ما اراد اللّه لذا اكثر اعراض نموده‌ بهاويه راجع
شدند و هر چه اعراض مشركين شدّت نمود آن شمس معانى محكم‌تر حجّيت آيات را نازل
فرمودند بشأنيكه كلّ حجّت و برهان را منحصر بآيات نمودند

\end{otherlanguage}

Be-l-ākhare be-awhām-e anfos-e khod mashghūl va maḥjūb be-shaʾnī ke
be-āyāt-e Allāh kefāyat nanamūde, maʿa ānke naṣṣ-e Forqān ast a-wa-lam
yakfihim annā anzalnā ʿalayka al-kitāb, va jamīʿ mostadell be-īnke
edrāk-e āyāt fawq-e ʿoqūl va ʿerfān-e mā ast, amr-e dīgar bāyad ẓāher
shavad. Qad ghayyarū niʿmata Allāh wa-ḥujjatahu wa-arādū bi-ghayri mā
arāda Allāh. Le-dhā akthar eʿrāẓ namūde be-hāviye rājeʿ shodand va har
che eʿrāẓ-e moshrekīn sheddat namūd ān Shams-e Maʿānī moḥkamtar
ḥojjiyyat-e āyāt rā nāzel farmūdand be-shaʾnī ke koll-e ḥojjat va borhān
rā monḥaṣer be-āyāt namūdand.

In the end, they were so absorbed in and veiled by their own imaginings
that they did not regard the verses of God as sufficient, even though
the explicit text of the Qur'an states, ``Is it not enough for them that
We have sent down to you the Book?'' All argued, ``Comprehension of the
verses is beyond our reason and understanding; something else must
appear.'' They altered God's blessing and His proof, and desired
something other than what God desired. Thus most turned away and
returned to the abyss. The more fiercely those who associated partners
with God turned away, the more firmly that Sun of inner meanings
revealed the evidentiary authority of the verses, to the point that He
confined all proof and testimony to the verses alone.

\subsection{672}\label{section-1167}

\begin{otherlanguage}{persian}

و رغماً لأنف المشركين اين آيه از سماء احديّة نازل لو ارادت نملة ان يتصرّف
فى القرآن و باطنه و باطن باطنه فى حكم سواد عينيها لتقدر لانّ سرّ الصّمدانيّة
و تجلّى الأحدانيّة قد تلجلجت فى كلّ شیء و اين آيه نازل كه جميع مستشعر شوند
كه ابواب رحمت و عنايت مسدود نيست هر نفسى مقبل شود و الى اللّه توجّه نمايد
حجيّت آيات اللّه را بنفس آيه ادراك مينمايد و ادراك آن مقام بعلوم ظاهره
منوط نبوده و نخواهد بود بلكه بنفوس زكيّه و قلوب طاهره مقبله منوط طوبی‌
للمقبلين

\end{otherlanguage}

Va raghman li-anfi l-mushrikīn īn āyih az samā-yi aḥadiyyat nāzil: law
arādat namlatun an yataṣarrafa fī l-Qurʾān wa-bāṭinihi wa-bāṭini
bāṭinihi fī ḥukmi sawādi ʿaynayhā la-taqdiru li-anna sirra l-ṣamadāniyya
wa-tajalliya l-aḥadāniyya qad talajlajat fī kulli shayʾ. Va īn āyih
nāzil kih jamīʿ mustashʿir shavand kih abvāb-i raḥmat va ʿināyat masdūd
nīst. Har nafsī muqbil shavad va ila Allāh tavajjuh namāyad ḥujjiyyat-i
āyāt-i Allāh rā bi-nafs-i āyih idrāk mīnamāyad va idrāk-i ān maqām
bi-ʿulūm-i ẓāhirih manūṭ nabūdih va nakhvāhad būd balkih bi-nufūs-i
zakiyyih va qulūb-i ṭāhirih-yi muqbilih manūṭ. Ṭūbā li-l-muqbilīn.

And despite those who associate partners with God, this verse descended
from the heaven of oneness: ``If an ant wished to exercise power over
the Qur'an, its inner meaning, and the inner meaning of its inner
meaning, as over the black of its own eyes, it would be able to do so,
for the mystery of self-sufficiency and the revelation of oneness have
stirred within every thing.'' This verse was revealed so that all might
realize that the doors of mercy and loving care are not closed. Whoever
turns toward God perceives the evidentiary authority of the verses of
God through the verse itself. Understanding this station has never
depended, nor will it ever depend, on outward learning; rather, it
depends on pure souls and unsullied hearts turned toward Him. Blessed
are those who turn toward Him.

\subsection{673}\label{section-1168}

\begin{otherlanguage}{persian}

و هم‌چنين آن ساذج معانى حجبات و اوهاميكه از قبل بود و ناس بآن محتجب و
متوهّم جميع را خرق نمودند و بافق انّنى انا حىّ فى الأفق الأبهی راجع شدند
حال جمعى انفس موهومه لاثبات رياست و حفظ آن اراده نموده‌اند اوّلا حجيّت آيات
را از بيان محو نمايند و در اين فقره بسيار سعى نموده‌ و مينمايند و بتمام
مكر و حيله در اين امر مشغولند و لكن ابی‌ اللّه الّا ان يثبت آياته بالحقّ و
لو تنكرها المشركون و حال در نفاق بمقامى رسيده‌اند كه اگر نفسى بحجيّت آيات
متكلّم شود باو استهزاء و سخريه مينمايند بعينه مثل ملل قبل بل هؤلآء اشدّ
جهلا و اكبر نفاقاً لو انتم تعرفون

\end{otherlanguage}

Va ham-chunīn ān sādhij-i maʿānī ḥujubāt va avhāmī kih az qabl būd va
nās bi-ān muḥtajib va mutavahhim jamīʿ rā kharq namūdand va bi-ufuq-i
innanī anā ḥayyun fī l-ufuqi l-abhā rājiʿ shudand. Ḥāl jamʿī anfus-i
mawhūmih li-ithbāt-i riyāsat va ḥifẓ-i ān irādih namūdih-and avvalan
ḥujjiyyat-i āyāt rā az Bayān maḥv namāyand va dar īn fiqrih bisyār saʿy
namūdih va mīnamāyand va bi-tamām-i makr va ḥīlih dar īn amr mashghūland
va lākin abā Allāhu illā an yuthbita āyātihi bi-l-ḥaqq wa-law tunkiruhā
l-mushrikūn. Va ḥāl dar nifāq bi-maqāmī rasīdih-and kih agar nafsī
bi-ḥujjiyyat-i āyāt mutakallim shavad bi-ū istihzā va sukhriyyih
mīnamāyand bi-ʿaynih mithl-i milal-i qabl bal hāʾulāʾi ashaddu jahlan va
akbaru nifāqan law antum taʿrifūn.

Likewise, that Quintessence of meanings tore through all the veils and
imaginings that had existed before and by which people had been veiled
and deluded, and returned to the horizon of ``Indeed, I am alive in the
all-glorious horizon.'' Now a number of deluded souls, intent on
establishing and preserving their leadership, have resolved first to
erase from the Bayán the evidentiary authority of the verses. They have
striven greatly in this matter and continue to do so, employing every
trick and stratagem. But God refuses to do anything other than establish
His verses in truth, even though those who associate partners with Him
deny them. Their hypocrisy has now reached such a degree that if anyone
speaks of the evidentiary authority of the verses, they ridicule and
mock him, exactly like the religious communities of the past. Indeed,
these people are more ignorant and greater in hypocrisy, if you only
knew.

\subsection{674}\label{section-1169}

\begin{otherlanguage}{persian}

لازال عرفان منبسطه الهيّه را نفوس محتجبه محدود نموده‌اند و باين جهت اكثر
عباد از فيوضات رحمانى و عنايات ربّانى محروم مانده‌اند چنانچه نفس ملقى يومى
از ایّام استهزاء على نقطة البيان ذكر مينمود كه فرموده‌اند كه احدى بر اتيان
يك آيه قادر نه و حال آنكه مى بينيم كه بعضى نفوس قادرند كه دو كلمه آيات
تكلّم نمايند پس از اين قرار اين نفوس هم حقّ خواهند بود فو‌اللّه من وجد رايحة
رضوان المعانى ليعرف من هذه الكلمة كفره و نفاقه و بغيه و اعتراضه على كلّ
النّبيين و المرسلين

\end{otherlanguage}

Lāzāl ʿirfān-i munbasiṭih-yi ilāhiyyih rā nufūs-i muḥtajibih maḥdūd
namūdih-and va bi-īn jihat akthar-i ʿibād az fuyūḍāt-i raḥmānī va
ʿināyāt-i rabbānī maḥrūm māndih-and. Chunānchih nafs-i mulqī yawmī az
ayyām istihzāʾan ʿalā Nuqṭati l-Bayān dhikr mīnamūd kih farmūdih-and kih
aḥadī bar ityān-i yik āyih qādir nah va ḥāl ānkih mībīnīm kih baʿḍī
nufūs qādirand kih du kalimih āyāt takallum namāyand pas az īn qarār īn
nufūs ham Ḥaqq khvāhand būd. Fa-wa-Allāhi man wajada rāʾiḥata Riḍwāni
l-maʿānī la-yaʿrifu min hādhihi l-kalimati kufrahu wa-nifāqahu
wa-baghyahu wa-iʿtirāḍahu ʿalā kulli l-nabiyyīn wa-l-mursalīn.

Veiled souls have always confined the all-embracing knowledge of God,
and for this reason most people have remained deprived of the
outpourings of the All-Merciful and the loving care of the Lord. Thus,
one day, the very person who put this notion forward said in mockery of
the Point of the Bayán: ``He has stated that no one is capable of
producing a single verse, yet we see that some people are able to utter
two words of verses. On this basis, these people too would be the True
One.'' By God, anyone who has caught the fragrance of the paradise of
meanings will recognize from this statement his unbelief, hypocrisy,
rebellion, and opposition to all the Prophets and Messengers.

\subsection{675}\label{section-1170}

\begin{otherlanguage}{persian}

مشاهده شد كه ابداً از كوثر معانى مرزوق نشده و از رضوان معنوى الهى بغايت
محروم مانده چه كه اين بسى واضح است كه مقصود حضرت از اين بيان چه بوده و
كلمه حقّ اگر چه بحرفين باشد سلطان كلمات اهل ارض بوده وخواهد بود و اللّه از
نفحات همان كلمه وحده عارفين نفحات حقّ را مييابند و برضوانش توجّه مينمايند
ملاحظه نما كه اگر از شطّ عذب فرات كاسى اخذ شود صاحبان ذايقه ادراك
مينمايند كه اين ماء از كدام منبع بوده و از عذوبت او پى‌برند بمنبع او و
حال آنكه بحسب ظاهر يك كاس ماء بوده و چون منبع علم و عرفان و عنايت و
مكرمت و فضل و قدرت و جلال و عظمت جارى و سارى است لذا از آن كلمه مخلصين و
مقرّبين ادراك مينمايند و بمنبع آن پى‌ميبرند كذلك فصّل لك لعلّ تهتدى بهدى
اللّه ربّك و لاتلتفت الى جنود الشّياطين

\end{otherlanguage}

Mushāhidih shud kih abadan az kawthar-i maʿānī marzūq nashudih va az
Riḍvān-i maʿnavī-yi ilāhī bi-ghāyat maḥrūm māndih chih kih īn basī vāḍiḥ
ast kih maqṣūd-i Ḥaḍrat az īn bayān chih būdih va kalimih-yi Ḥaqq agar
chih bi-ḥarfayn bāshad sulṭān-i kalimāt-i ahl-i arḍ būdih va khvāhad
būd. Va-Allāh az nafaḥāt-i hamān kalimih vaḥdih ʿārifīn nafaḥāt-i Ḥaqq
rā mīyāband va bi-Riḍvānash tavajjuh mīnamāyand. Mulāḥiẓih namā kih agar
az shaṭṭ-i ʿadhb-i furāt kāsī akhdh shavad ṣāḥibān-i dhāʾiqih idrāk
mīnamāyand kih īn māʾ az kudām manbaʿ būdih va az ʿudhūbat-i ū pay
barand bi-manbaʿ-i ū va ḥāl ānkih bi-ḥasab-i ẓāhir yik kās-i māʾ būdih
va chūn manbaʿ-i ʿilm va ʿirfān va ʿināyat va makramat va faḍl va qudrat
va jalāl va ʿaẓamat jārī va sārī ast lidhā az ān kalimih mukhliṣīn va
muqarrabīn idrāk mīnamāyand va bi-manbaʿ-i ān pay mībarand. Kadhālika
faṣṣala laka laʿalla tahtadī bi-hudā Allāhi rabbika wa-lā taltafit ilā
junūdi l-shayāṭīn.

It became evident that he had never been nourished by the heavenly
waters of meanings and had remained utterly deprived of the divine
paradise of spiritual understanding, for what His Holiness intended by
this statement is abundantly clear. The word of the True One, even if it
consists of only two letters, has always been and will always be
sovereign over the words of the people of the earth. By God, from the
fragrance of that word alone, those endowed with understanding perceive
the fragrances of the True One and turn toward His paradise. Consider:
if a cup is drawn from a sweet, refreshing river, those with a sense of
taste perceive the source of this water and recognize its source from
its sweetness, although outwardly it is only a cup of water. And since
the wellspring of knowledge, recognition, loving care, generosity,
grace, power, glory, and grandeur flows and pervades all things, the
sincere and those brought near perceive through that word and recognize
its source. Thus He has explained it to you, so that you may be guided
by the guidance of God, your Lord, and pay no heed to the hosts of
devils.

\subsection{676}\label{section-1171}

\begin{otherlanguage}{persian}

و اگر نفسى دو كلمه بگويد و در كلمه سيّم معطّل ماند مبرهن است كه از اين
منبع جاريه ساريه نبوده و نخواهد بود فو‌اللّه يا عبد حقّ بجميع شئونش از دونش
ممتاز بوده و خواهد بود بر خود مشتبه مكن ميخواهى تجربه كن بين يدى حاضر شو
و يك آيه از آيات اللّه سؤال نما اذا تطّلع ببحر الّذى يتموّج بذاته لذاته و لم
يكن لامواجه من نفاد و تطّلع بمكمن المعانى و العرفان و مخزن العلم و الحكمة
و تكون من الموقنين

\end{otherlanguage}

Va agar nafsī du kalimih bigūyad va dar kalimih-yi siyyum muʿaṭṭal mānad
mubarhan ast kih az īn manbaʿ-i jāriyih-yi sāriyih nabūdih va nakhvāhad
būd. Fa-wa-Allāh yā ʿAbd Ḥaqq bi-jamīʿ-i shuʾūnash az dūnash mumtāz
būdih va khvāhad būd. Bar khud mushtabih makun mīkhvāhī tajribih kun
bayna yaday ḥāḍir shaw va yik āyih az āyāt-i Allāh suʾāl namā. Idhan
taṭṭaliʿ bi-baḥri lladhī yatamawwaju bi-dhātihi li-dhātihi wa-lam yakun
li-amwājihi min nafād wa-taṭṭaliʿ bi-makmani l-maʿānī wa-l-ʿirfān
wa-makhzani l-ʿilmi wa-l-ḥikma wa-takūnu mina l-mūqinīn.

If someone utters two words but is at a loss for a third, it is evident
that these have not come, nor will they ever come, from this flowing,
all-pervading wellspring. By God, O servant, the True One has always
been and will always be distinguished from all others in every aspect of
His being. Do not confuse the matter for yourself. If you wish, put it
to the test: enter My presence and ask for a single verse of the verses
of God. Then you will behold the ocean that surges of itself and for
itself, whose waves have no end, and you will discover the sanctuary of
meanings and recognition and the treasury of knowledge and wisdom, and
be among those who are certain.

\subsection{677}\label{section-1172}

\begin{otherlanguage}{persian}

و اين ظهور اعظم را چون نقطه بيان جلّ كبريائه مخصوص در جميع بيان فرموده
بآيات ظاهر ميشود لذا بآيات ظاهر شده و الّا كلّ ما ظهر من عنده حجّة على
العالمين در اين ظهور اعظم جميع آيات قدرتيّه و آفاقيّه و انفسيّه مشاهده شده
و لاينكرها الّا كلّ مغلّ مكّار

\end{otherlanguage}

Va īn Ẓuhūr-i Aʿẓam rā chūn Nuqṭih-yi Bayān jalla kibriyāʾuhu makhṣūṣ
dar jamīʿ-i Bayān farmūdih bi-āyāt ẓāhir mīshavad lidhā bi-āyāt ẓāhir
shudih va illā kullu mā ẓahara min ʿindihi ḥujjatun ʿalā l-ʿālamīn. Dar
īn Ẓuhūr-i Aʿẓam jamīʿ-i āyāt-i qudratiyyih va āfāqiyyih va anfusiyyih
mushāhidih shudih va-lā yunkiruhā illā kullu mughillin makkār.

Since the Point of the Bayán---exalted be His majesty---specifically
stated throughout the Bayán that this Most Great Revelation would appear
through verses, it has therefore appeared through verses. Otherwise,
everything that has appeared from Him is a proof to all the worlds. In
this Most Great Revelation, all the signs of power, all the signs in the
world and within human souls, have been witnessed, and none denies them
except the treacherous and deceitful.

\subsection{678}\label{section-1173}

\begin{otherlanguage}{persian}

بر هر نفسى اليوم لازم كه از انفس خبيثه مشركه اجتناب نمايد چه كه بمكرى
ظاهر شده‌اند كه در ابداع شبه آن ممكن نه چنانچه از ارض صاد بجناب اسم اللّه
آقاسيّد مهدى عليه بهآء اللّه نوشته‌اند كه امّ الفواحش و الخطيئات مكتوبی‌ باين
ارض فرستاده مضمون آنكه در بيت ايشان لوحى مدفون بوده از خطّ حضرت اعلى و من
آن لوح را بيرون آورده‌ام و در آن لوح عباراتى مسطور و نسبت بنقطه اولى داده
كه بطلعت ابهى در سفارش ميرزا يحيى نوشته‌‌اند فو‌اللّه الّذى لا اله الّا هو كلّ
متحيّر مانديم از جعل اين نفوس و مكر اين قوم

\end{otherlanguage}

Bar har nafsī al-yawm lāzim kih az anfus-i khabīthih-yi mushrikih
ijtināb namāyad chih kih bi-makrī ẓāhir shudih-and kih dar ibdāʿ shibh-i
ān mumkin nah. Chunānchih az arḍ-i Ṣād bi-janāb-i Ismuʾllāh Āqā Siyyid
Mahdī ʿalayhi Bahāʾuʾllāh nivishtih-and kih Umm-i l-Fawāḥish
va-l-Khaṭīʾāt maktūbī bi-īn arḍ firistādih maḍmūn ānkih dar bayt-i īshān
lawḥī madfūn būdih az khaṭṭ-i Ḥaḍrat-i Aʿlā va man ān lawḥ rā bīrūn
āvardih-am va dar ān lawḥ ʿibārātī masṭūr va nisbat bi-Nuqṭih-yi Ūlā
dādih kih bi-Ṭalʿat-i Abhā dar sifārish-i Mīrzā Yaḥyā nivishtih-and.
Fa-wa-Allāhi lladhī lā ilāha illā huwa kull mutaḥayyir māndīm az jaʿl-i
īn nufūs va makr-i īn qawm.

Today everyone must avoid these wicked souls who associate partners with
God, for they have appeared with a deceit whose like is impossible in
creation. Thus, a letter has been written from the land of Ṣád to His
Honor Ismu'lláh, Áqá Siyyid Mahdí---may the glory of God rest upon
him---reporting that the Mother of Abominations and Sins has sent a
letter to that land. Its purport is that a Tablet in the handwriting of
His Holiness the Most Exalted had been buried in their house, and that
``I have brought that Tablet out.'' In that Tablet, certain passages are
inscribed and attributed to the Primal Point, purportedly written to the
All-Glorious Beauty in commendation of Mírzá Yaḥyá. By God, besides whom
there is no other god, we were all astonished at the fabrications of
these souls and the deceit of these people.

\subsection{679}\label{section-1174}

\begin{otherlanguage}{persian}

و بعد از مصدر امر حكم شده كه بروند و باو بگويند كه آن لوح را از كدام محلّ
بيت بيرون آورده و بنما لعلّ تتنبّه فى نفسها و تكون من التّائبات باين شدّت بر
افترا قيام نموده‌اند و اين امور دستور‌العملى است كه ميرزا يحيى از اينجا
مينويسد و ميفرستد لعنة اللّه على الكاذبين و المفترين اذا يتبرّء منهم نقطة
البيان فى عرش الجنان و يقول خسرانا لكم يا ملأ المشركين و نارا عليكم يا
ملأ المغرضين

\end{otherlanguage}

Va baʿd az maṣdar-i amr ḥukm shudih kih biravand va bi-ū bigūyand kih ān
lawḥ rā az kudām maḥall-i bayt bīrūn āvardih va binamā laʿalla
tatanabbahu fī nafsihā va takūnu mina l-tāʾibāt. Bi-īn shiddat bar
iftirā qiyām namūdih-and va īn umūr dastūr al-ʿamalī ast kih Mīrzā Yaḥyā
az īnjā mīnivīsad va mīfiristad. Laʿnatu Allāhi ʿalā l-kādhibīn
wa-l-muftarīn. Idhan yatabarraʾu minhum Nuqṭatu l-Bayān fī ʿarshi
l-jinān wa-yaqūlu khusrānan lakum yā malaʾa l-mushrikīn va nāran
ʿalaykum yā malaʾa l-mughriḍīn.

Afterward, the Source of authority commanded that they go and ask her
from which part of the house she had brought out that Tablet, and tell
her, ``Show it,'' so that she might come to her senses and be among
those who repent. They have risen to fabricate falsehoods with such
vehemence. These actions follow instructions that Mírzá Yaḥyá writes and
sends from here. May God's curse rest upon the liars and fabricators.
Thereupon the Point of the Bayán, on the throne of paradise, disavows
them and says: ``Ruin upon you, you who associate partners with God!
Fire upon you, you people of malice!''

\subsection{680}\label{section-1175}

\begin{otherlanguage}{persian}

اليوم بايد بصر را حديد نمود و قلب را طاهر و صدر را منزّه تا انسان خود را
از دونش بشناسد و تميز دهد و اين رتبه اگر چه بسيار بلند و عظيم است ولكن
نزد منصفين بسيار سهل و يسير اگرچه حال اهل بيان بمقامى رسيده‌اند و بامورى
بر حقّ اعتراض نموده‌اند كه از اوّل ابداع تا حين احدى در مثل آن امور بر حقّ
اعتراض ننموده ‌جميع اين عباد سالها است طايف حول بوده‌ايم و اللّه العظيم كه
هرگز نزدشان از زخرف دنيا شيئى نبوده و اگر هم بوده جميع شريك بوده‌ايم و هر
نفسى ببيت اعظم وارد گويا خود صاحب بيت است كمال عنايت درباره او مبذول و
مع‌ذلك اهل بيان اعتراضا على اللّه بكلّ ما ظهر من الحقّ اعتراض نموده‌اند حتّى
شرب چاى و همچنين در لباس و امثال آن افّ لهم و لحيائهم

\end{otherlanguage}

Al-yawm bāyad baṣar rā ḥadīd namūd va qalb rā ṭāhir va ṣadr rā munazzah
tā insān khud rā az dūnash bishināsad va tamyīz dihad va īn rutbih agar
chih bisyār buland va ʿaẓīm ast va lākin nazd-i munṣifīn bisyār sahl va
yasīr. Agarchih ḥāl ahl-i Bayān bi-maqāmī rasīdih-and va bi-umūrī bar
Ḥaqq iʿtirāḍ namūdih-and kih az avval-i ibdāʿ tā ḥīn aḥadī dar mithl-i
ān umūr bar Ḥaqq iʿtirāḍ nanamūdih. Jamīʿ-i īn ʿibād sālhā ast ṭāʾif-i
ḥawl būdih-īm va-Allāhu l-ʿaẓīm kih hargiz nazdishān az zukhruf-i dunyā
shayʾī nabūdih va agar ham būdih jamīʿ sharīk būdih-īm va har nafsī
bi-Bayt-i Aʿẓam vārid gūyā khud ṣāḥib-i bayt ast kamāl-i ʿināyat dar
bārih-yi ū mabdhūl va maʿa dhālik ahl-i Bayān iʿtirāḍan ʿalā Allāh
bi-kulli mā ẓahara mina l-Ḥaqq iʿtirāḍ namūdih-and ḥattā shurb-i chāy va
hamchunīn dar libās va amthāl-i ān. Uffin lahum va li-ḥayāʾihim.

Today one must sharpen one's sight, purify one's heart, and sanctify
one's breast, so that a person may recognize and distinguish himself
from others. Although this rank is very lofty and great, it is very easy
and simple for the fair-minded. Yet the people of the Bayán have now
reached such a state, and have objected to the True One over such
matters, that from the beginning of creation until now no one has
objected to the True One over anything like them. All of us servants
have circled around Him for years. By God the Almighty, He has never
possessed anything of this world's finery, and if there has been
anything, we have all shared in it. Anyone who enters the Most Great
House is treated as though he himself were its owner, and the utmost
loving care is lavished upon him. Nevertheless, the people of the Bayán,
in opposition to God, have objected to everything that has appeared from
the True One, even the drinking of tea, and likewise to His clothing and
similar things. Shame on them and their sense of shame!

\subsection{681}\label{section-1176}

\begin{otherlanguage}{persian}

نميدانم چه شده كه هر نفسى در عقبه واقف و مبتلى لوحى در اين ایّام از مشرق
مشيّت رحمن مخصوص يكى از احبّا مشرق و نازل در اينمقام لاجل تنبّه ذكر ميشود
شايد متنبّه شوند و بشنوند قول نقطه اولى را در كرّه اخرى از لسان ابدع ابهى
و از ضجيج و حنين و ناله آن حضرت شرم نمايند و بآنچه عامل شده كفايت كنند و
هذه صورة ما نزل من جبروت اللّه العلىّ العظيم

\end{otherlanguage}

Namīdānam chih shudih kih har nafsī dar ʿaqabih vāqif va mubtalā. Lawḥī
dar īn ayyām az mashriq-i mashiyyat-i Raḥmān makhṣūṣ-i yikī az aḥibbā
mushriq va nāzil. Dar īn maqām li-ajl-i tanabbuh dhikr mīshavad shāyad
mutanabbih shavand va bishinavand qawl-i Nuqṭih-yi Ūlā rā dar karrih-yi
ukhrā az lisān-i abdaʿ-i abhā va az ḍajīj va ḥanīn va nālih-yi ān Ḥaḍrat
sharm namāyand va bi-ānchih ʿāmil shudih kifāyat kunand. Va hādhihi
ṣūratu mā nazala min jabarūti Allāhi l-ʿAliyyi l-ʿAẓīm.

I do not know what has happened that every soul has halted at the
mountain pass and become afflicted. In these days a Tablet has dawned
and descended from the dayspring of the will of the All-Merciful,
addressed specifically to one of the loved ones. It is quoted here to
awaken them, so that they may perhaps come to their senses and hear once
again the words of the Primal Point from a wondrous, all-glorious
tongue, feel ashamed at His cries, lamentations, and weeping, and
content themselves with what they have already done. This is the text of
what descended from the dominion of God, the Most High, the Almighty.

\subsection{682}\label{section-1177}

\begin{otherlanguage}{arabic}

بسم اللّه الأقدس العلىّ الأعلی

\end{otherlanguage}

Bi-smi Allāhi l-Aqdasi l-ʿAliyyi l-Aʿlā.

In the name of God, the Most Holy, the High, the Most Exalted.

\subsection{683}\label{section-1178}

\begin{otherlanguage}{arabic}

هذا كتاب نقطة الأولی الى الّذينهم آمنوا باللّه الواحد الفرد العزيز العليم
و فيه يخاطب الّذينهم توقّفوا فى هذا الأمر من ملأ البيانييّن لعلّ يستشعرن
ببدايع كلمات اللّه و يقومن عن رقد الغفلة فى هذا الفجر المشرق المنير قل
انّا امرناكم فى الكتاب بأن لاتقدموا طايفة الّتى يظهر منها محبوب العارفين و
مقصود من فى السّموات و الأرضين و امرناكم ان ادركتم لقآء اللّه قوموا تلقآء
الوجه ثمّ انطقوا من قبلى بهذه الكلمة العزيز المنيع عليك يا بهآء اللّه و
ذوى قرابتك ذكر اللّه و ثنآء كلّ شیء فى كلّ حين و قبل حين و بعد حين و جعلنا
هذه الكلمة عزّاً لاهل البيان لعلّ بها يرتقون الى معارج القدس و يكوننّ من
الفائزين و انّهم تركوا ما امروا به بحيث ما ظهر احد تلقآء الوجه بما
امرناهم فى الواح عزّ عظيم بل رموا نحوه من كلّ الآفاق رمى النّفاق و بذلك
بكيت و بكت اهل جبروت العظمة ثمّ روح الأمين

\end{otherlanguage}

Hādhā kitābu Nuqṭati l-Ūlā ilā lladhīnahum āmanū bi-Allāhi l-Wāḥidi
l-Fardi l-ʿAzīzi l-ʿAlīm wa-fīhi yukhāṭibu lladhīnahum tawaqqafū fī
hādhā l-amri min malaʾi l-bayāniyyīn laʿalla yastashʿirna bi-badāʾiʿi
kalimāti Allāhi wa-yaqūmna ʿan raqdi l-ghaflati fī hādhā l-fajri
l-mushriqi l-munīr. Qul innā amarnākum fī l-kitābi bi-an lā taqdamū
ṭāʾifata llatī yaẓharu minhā maḥbūbu l-ʿārifīn wa-maqṣūdu man fī
l-samāwāti wa-l-arḍīn wa-amarnākum in adraktum liqāʾa Allāhi qūmū
tilqāʾa l-wajhi thumma nṭiqū min qiblī bi-hādhihi l-kalimati l-ʿazīzi
l-manīʿ: ʿalayka yā Bahāʾa Allāhi wa-dhawī qarābatika dhikru Allāhi
wa-thanāʾu kulli shayʾin fī kulli ḥīnin wa-qabla ḥīnin wa-baʿda ḥīn.
Wa-jaʿalnā hādhihi l-kalimata ʿizzan li-ahli l-Bayāni laʿalla bihā
yartaqūna ilā maʿāriji l-qudsi wa-yakūnanna mina l-fāʾizīn. Wa-innahum
tarakū mā umirū bihi bi-ḥaythu mā ẓahara aḥadun tilqāʾa l-wajhi bi-mā
amarnāhum fī alwāḥi ʿizzin ʿaẓīm bal ramaw naḥwahu min kulli l-āfāqi
ramya l-nifāqi wa-bi-dhālika bakaytu wa-bakat ahlu jabarūti l-ʿaẓamati
thumma Rūḥu l-Amīn.

This is the Book of the Primal Point to those who have believed in God,
the One, the Unique, the Almighty, the All-Knowing. In it He addresses
those among the people of the Bayán who have hesitated in this Cause, so
that they may become aware of the wonders of God's words and rise from
the sleep of heedlessness in this radiant, luminous dawn. Say: We
commanded you in the Book not to take precedence over the community from
which would appear the Beloved of those endowed with understanding and
the Desire of all who are in the heavens and on earth. We commanded you,
should you attain the presence of God, to stand before His face and then
utter on My behalf these mighty, exalted words: ``Upon You, O
Bahá'u'lláh, and upon Your kindred, be the remembrance of God and the
praise of every thing, at every moment, before every moment, and after
every moment.'' We made these words an honor for the people of the
Bayán, so that through them they might ascend to the heights of holiness
and be among those who attain their goal. Yet they abandoned what they
were commanded, so that no one appeared before His face with what We had
enjoined upon them in Tablets of great glory. Instead, from every
horizon they hurled the shafts of hypocrisy at Him. At this I wept, and
the inhabitants of the dominion of grandeur wept, and then the faithful
Spirit.

\subsection{684}\label{section-1179}

\begin{otherlanguage}{persian}

قل يا قوم فاستحيوا عن جمالى انّ الّذى قد ظهر بالحقّ انّه لبهآء العالمين لو
انتم من العارفين و انّه لبهآء اللّه و عليه ذكر اللّه و ثنائه ثم ثنآء اهل
ملأ الأعلی و ثنآء اهل جبروت البقآء و ثنآء كلّ شیء فى كلّ حين ایّاکم ان
تحتجبوا بما خلق بين الأرض و السّمآء ان اسرعوا الى رضوان رضائه و لا تكوننّ
من الرّاقدين

\end{otherlanguage}

Qul yā qawmi fa-staḥyū ʿan jamālī inna lladhī qad ẓahara bi-l-ḥaqqi
innahu la-Bahāʾu l-ʿālamīn law antum mina l-ʿārifīn wa-innahu la-Bahāʾu
Allāhi wa-ʿalayhi dhikru Allāhi wa-thanāʾuhu thumma thanāʾu ahli malaʾi
l-aʿlā wa-thanāʾu ahli jabarūti l-baqāʾ wa-thanāʾu kulli shayʾin fī
kulli ḥīn. Iyyākum an taḥtajibū bi-mā khuliqa bayna l-arḍi wa-l-samāʾ an
asriʿū ilā Riḍwāni riḍāʾihi wa-lā takūnanna mina l-rāqidīn.

Say: My people, feel shame before My beauty! He who has appeared in
truth is indeed the Glory of all the worlds, if you are among those
endowed with understanding. He is indeed Bahá'u'lláh, and upon Him be
the remembrance of God and His praise, and the praise of the inhabitants
of the highest assembly, the praise of the inhabitants of the dominion
of eternity, and the praise of every thing at every moment. Beware lest
you be veiled by anything created between earth and heaven. Hasten to
the paradise of His good pleasure, and do not be among those who sleep.

\subsection{685}\label{section-1180}

\begin{otherlanguage}{persian}

قل انّ جماله كان جمالى بالحقّ و انّ نفسه نفسى و كلّما نزّلناه فى البيان قد
نزل لأمره المحكم البديع اتّقو اللّه و لاتجادلوا بالّذى اخبرناكم به و
بشّرناكم بظهوره و اخذت عهد نفسه قبل عهد نفسى و يشهد بذلك كلّ شیء ان انتم
من المنكرين تاللّه بنغمة من نغماته ولدت حقايق كلّ شیء مرّة اخرى و بنغمة
اخرى استجذبت افئدة المقرّبين ایّاکم ان تحتجبوا بشىء عن الّذى كان لقائه ذات
لقائى و فدى نفسه فى سبيلى كما فديت فى سبيله حبّاً لجماله العزيز المنيع

\end{otherlanguage}

Qul inna jamālahu kāna jamālī bi-l-ḥaqq wa-inna nafsahu nafsī wa-kullamā
nazzalnāhu fī al-Bayān qad nazala li-amrihi al-muḥkam al-badīʿ. Ittaqū
Allāh wa-lā tujādilū bi-lladhī akhbarnākum bihi wa-bashsharnākum
bi-ẓuhūrihi wa-akhadhtu ʿahda nafsihi qabla ʿahda nafsī wa-yashhadu
bi-dhālika kull shayʾ in antum min al-munkirīn. Ta-llāh bi-naghma min
naghamātihi wulidat ḥaqāʾiq kull shayʾ marra ukhrā wa-bi-naghma ukhrā
istujdhibat afʾidat al-muqarrabīn. Iyyākum an taḥtajibū bi-shayʾ ʿan
alladhī kāna liqāʾuhu dhāta liqāʾī wa-fadā nafsahu fī sabīlī kamā
fadaytu fī sabīlihi ḥubban li-jamālihi al-ʿazīz al-manīʿ.

Say: His beauty was truly My beauty, and His own Self My Self, and
everything We sent down in the Bayán was revealed for His firm and
wondrous Cause. Fear God, and do not dispute concerning Him of whom We
informed you and whose appearance We announced to you. I took the
covenant concerning His Self before the covenant concerning My own Self,
and every thing bears witness to this, if you are among those who deny
it. By God! Through one of His melodies the realities of all things were
born anew, and through another melody the hearts of those brought near
were drawn to Him. Beware lest anything veil you from Him, meeting whom
was the very essence of meeting Me. He sacrificed Himself in My path,
just as I sacrificed Myself in His path, out of love for His mighty,
unassailable beauty.

\subsection{686}\label{section-1181}

\begin{otherlanguage}{persian}

قل لولاه ما ركب الحآء بالبآء و ما استقرّ هيكل الهآء على الواو و ما خلق ما
كان و ما يكون لو انتم من الشّاعرين و لولاه ما القيت نفسى بين يدى المشركين
و ما علّقت بين الهوا تاللّه باشتياقى اليه و شوقى الى نفسه قد حملت ما لا
حمله النّبيّين و المرسلين و رضيت كلّ ذلك على نفسى لئلّا يرد عليه ما يحزن به
فؤاده الألطف الأرقّ اللّطيف المنيع و وصّيناكم فى كلّ البيان بأن لايحزن احد
احداً لعلّ لايرد عليه من حزن و الّا مالى و ذكرى لكم و اشتغالى بكم يا ملأ
التّاركين

\end{otherlanguage}

Qul lawlāhu mā rukkiba al-Ḥāʾ bi-l-Bāʾ wa-mā istaqarra haykal al-Hāʾ
ʿalā al-Wāw wa-mā khuliqa mā kāna wa-mā yakūn law antum min al-shāʿirīn.
Wa-lawlāhu mā alqaytu nafsī bayna yaday al-mushrikīn wa-mā ʿulliqtu
bayna al-hawā. Ta-llāh bi-ishtiyāqī ilayhi wa-shawqī ilā nafsihi qad
ḥamaltu mā lā ḥamalahu al-nabiyyīn wa-l-mursalīn wa-raḍītu kulla dhālika
ʿalā nafsī li-allā yarida ʿalayhi mā yaḥzanu bihi fuʾāduhu al-alṭaf
al-araqq al-laṭīf al-manīʿ. Wa-waṣṣaynākum fī kull al-Bayān bi-an lā
yuḥzina aḥad aḥadan laʿalla lā yarida ʿalayhi min ḥuzn wa-illā mā lī
wa-dhikrī lakum wa-ishtighālī bikum yā malaʾ al-tārikīn.

Say: Were it not for Him, the Ḥāʾ would not have been joined to the Bāʾ,
nor would the form of the Hāʾ have been established upon the Wāw, nor
would anything that has been or will be have been created, if you are
among those who perceive. Were it not for Him, I would not have cast
Myself into the hands of those who associate partners with God, nor
would I have been suspended in the air. By God! In My yearning for Him
and My longing for His Self, I bore what the Prophets and Messengers did
not bear, and willingly accepted all this for Myself, lest anything
befall Him that would grieve His most tender, most delicate, gracious,
unassailable heart. Throughout the Bayán We enjoined you not to cause
anyone sorrow, that perhaps no sorrow might befall Him. Otherwise, what
have I to do with mentioning you or concerning Myself with you, O
assembly of those who have turned away?

\subsection{687}\label{section-1182}

\begin{otherlanguage}{arabic}

و انّى ما اردت فى البيان الّا نفسه و لا من الأذکار‌ الّا ذكره و لا من
الأسماء الّا اسمه المبارك الأمنع الأقدس الأبدع البديع فو‌عمرى لو ذكرت ذكر
الرّبوبيّة ما اردت الّا ربوبيّته على كلّ الأشیآء و ان جرى من قلمى ذكر
الالوهيّة ما كان مقصودى الّا اله العالمين و ان جرى من قلمى ذكر المقصود فهو
كان مقصودى و كذلك فى المحبوب انّه قد كان محبوبی‌ و محبوب العارفين و ان
ذكرت ذكر السّجود ما اردت الّا السّجود لوجهه المتعالى العزيز المنيع و ان
اثنيت نفساً ما كان مقصود قلبی‌ الّا ثنآء نفسه و ان امرت النّاس بعمل ما اردت
الّا العمل فى رضائه فى يوم ظهوره و بذلك يشهد كلّما نزل علىّ من جبروت ربّی‌
العليم الحكيم و علّقت كلّ شیء بتصديقه و رضائه و انّه لهو الّذى قد كان بنفسه
اله العالمين و مقصود القاصدين و انتم لو تدقّون الأبصار لتشهدون مظاهر يفعل
ما يشآء فى ظلّه لمن العابدين و انتم قد فعلتم بنفسه ما لا فعل امّة الفرقان
بنفسى و لا ملأ اليهود بالرّوح

\end{otherlanguage}

Wa-innī mā aradtu fī al-Bayān illā nafsahu wa-lā min al-adhkār illā
dhikrahu wa-lā min al-asmāʾ illā ismahu al-mubārak al-amnaʿ al-aqdas
al-abdaʿ al-badīʿ. Fa-wa-ʿamrī law dhakartu dhikr al-rubūbiyya mā aradtu
illā rubūbiyyatahu ʿalā kull al-ashyāʾ wa-in jarā min qalamī dhikr
al-ulūhiyya mā kāna maqṣūdī illā ilāh al-ʿālamīn wa-in jarā min qalamī
dhikr al-maqṣūd fa-huwa kāna maqṣūdī wa-kadhālika fī al-maḥbūb innahu
qad kāna maḥbūbī wa-maḥbūb al-ʿārifīn wa-in dhakartu dhikr al-sujūd mā
aradtu illā al-sujūd li-wajhihi al-mutaʿālī al-ʿazīz al-manīʿ wa-in
athnaytu nafsan mā kāna maqṣūd qalbī illā thanāʾ nafsihi wa-in amartu
al-nās bi-ʿamal mā aradtu illā al-ʿamala fī riḍāʾihi fī yawm ẓuhūrihi
wa-bi-dhālika yashhadu kullamā nazala ʿalayya min jabarūt rabbī al-ʿalīm
al-ḥakīm. Wa-ʿallaqtu kulla shayʾ bi-taṣdīqihi wa-riḍāʾihi wa-innahu
la-huwa alladhī qad kāna bi-nafsihi ilāh al-ʿālamīn wa-maqṣūd al-qāṣidīn
wa-antum law tudiqqūn al-abṣār la-tashhadūn maẓāhir yafʿalu mā yashāʾ fī
ẓillihi la-min al-ʿābidīn. Wa-antum qad faʿaltum bi-nafsihi mā lā faʿala
ummat al-Furqān bi-nafsī wa-lā malaʾ al-Yahūd bi-l-Rūḥ.

In the Bayán I intended none but His Self; by its remembrances, none but
His remembrance; and by its names, none but His blessed, most
unassailable, most holy, most wondrous, incomparable Name. By My life!
If I made mention of lordship, I intended only His lordship over all
things. If mention of divinity flowed from My pen, My purpose was none
but the God of all worlds. If mention of the Desired One flowed from My
pen, He was My Desired One; and likewise with the Beloved: He was My
Beloved and the Beloved of those who know. If I made mention of
prostration, I intended only prostration before His exalted, mighty,
unassailable countenance. If I praised anyone, My heart intended only to
praise His Self. If I commanded people to perform a deed, I intended
only that it be done for His good pleasure on the day of His appearance.
Everything sent down to Me from the dominion of My Lord, the
All-Knowing, the All-Wise, bears witness to this. I made everything
dependent on His confirmation and good pleasure. He is indeed the One
who in His own Self has been the God of all worlds and the Desire of all
who seek. If you look closely, you will surely see the Manifestations of
``He does whatever He wills'' among those who worship in His shadow. Yet
you have done to His Self what neither the people of the Qur'an did to
Me nor the assembly of the Jews to the Spirit.

\subsection{688}\label{section-1183}

\begin{otherlanguage}{arabic}

فآهء آه من حرقة قلبی‌ و حنين نفسى فيما ورد على محبوبی‌ من ملأ المشركين افّ
لكم و لوفائكم يا معشر الظّالمين انا خلقنا الوفآء و الأدب لنفسه لعلّ عند
ظهوره لا تفعلوا امراً تجزع به حقيقتى و حقايق الأشیآء و انتم تجاوزتم عما
حدّد فى كتاب اللّه الملك العلىّ العظيم و خرقتم حجاب الحيا ثمّ ستر الحرمة و
عملتم ما يستحيى من ذكره قلم الانشآء بين الأرض و السّمآء

\end{otherlanguage}

Fa-āh āh min ḥurqat qalbī wa-ḥanīn nafsī fīmā warada ʿalā maḥbūbī min
malaʾ al-mushrikīn. Uff lakum wa-li-wafāʾikum yā maʿshar al-ẓālimīn.
Innā khalaqnā al-wafāʾ wa-l-adab li-nafsihi laʿalla ʿinda ẓuhūrihi lā
tafʿalū amran tajzaʿu bihi ḥaqīqatī wa-ḥaqāʾiq al-ashyāʾ wa-antum
tajāwaztum ʿammā ḥuddida fī kitāb Allāh al-malik al-ʿalī al-ʿaẓīm
wa-kharaqtum ḥijāb al-ḥayā thumma sitr al-ḥurma wa-ʿamiltum mā yastaḥyī
min dhikrihi qalam al-inshāʾ bayna al-arḍ wa-l-samāʾ.

Alas, alas for the burning of My heart and the lamentation of My soul at
what has befallen My Beloved at the hands of those who associate
partners with God! Shame on you and your faithfulness, O company of
wrongdoers! We created faithfulness and courtesy for His sake, so that
at His appearance you might not do anything that would cause My reality
and the realities of all things to grieve. Yet you have overstepped the
bounds set in the Book of God, the King, the Most High, the Most Great.
You have torn the veil of modesty and then the covering of reverence,
and have done what the Pen of creation is ashamed to mention between
earth and heaven.

\subsection{689}\label{section-1184}

\begin{otherlanguage}{arabic}

فآهء آه بما ورد منكم على هذا المظلوم الفريد الغريب و لم ادر ما تفعلون به
من بعد لا فو‌نفسى العليم بل اعلم و عندى علم كلّ شیء فى لوح جعله اللّه
محفوظاً عن انظر المشركين و اخبرناه من قبل بما ورد عليه و يرد ولو انّه قد
كان بنفسه عالماً بما فى صدور العالمين و لن يعزب عن علمه من شیء و لا يفوت
عن قبضته ما خلق بكلمة من عنده لا اله الّا هو الفرد الباعث المحيى المميت

\end{otherlanguage}

Fa-āh āh bi-mā warada minkum ʿalā hādhā al-maẓlūm al-farīd al-gharīb
wa-lam adri mā tafʿalūna bihi min baʿd. Lā fa-wa-nafsī al-ʿalīm bal
aʿlam wa-ʿindī ʿilm kull shayʾ fī lawḥ jaʿalahu Allāh maḥfūẓan ʿan
{[}corrupted text{]} al-mushrikīn wa-akhbarnāhu min qabl bi-mā warada
ʿalayhi wa-yarid wa-law innahu qad kāna bi-nafsihi ʿāliman bi-mā fī
ṣudūr al-ʿālamīn wa-lan yaʿzuba ʿan ʿilmihi min shayʾ wa-lā yafūtu ʿan
qabḍatihi mā khuliqa bi-kalima min ʿindihi lā ilāha illā huwa al-fard
al-bāʿith al-muḥyī al-mumīt.

Alas, alas for what you have inflicted on this Wronged One, alone and a
stranger! I do not know what you will do to Him hereafter. No, by My
all-knowing Self! Rather, I do know, and with Me is the knowledge of all
things in a Tablet that God has kept safe from the {[}corrupted original
text{]} of those who associate partners with Him. We told Him beforehand
what had befallen Him and what would befall Him, although He Himself
knew what lay within the hearts of all the world. Nothing will escape
His knowledge, nor will anything created by a word from Him elude His
grasp. There is no God but Him, the Unique, the Raiser of the dead, the
Giver of life, the Bringer of death.

\subsection{690}\label{section-1185}

\begin{otherlanguage}{arabic}

قل يا قوم انّه لهو الّذى لو يريد ان يجعل كلّ من فى السّموات و الأرض حجّة
باقية من عنده ليقدر و انّ هذا عنده سهل يسير و انّه لهو الّذى قد خلق رضوان
البيان لنفسه و منه بدء كلّ شیء و يعود لو انتم من العالمين و انتم بالّذى
كان فى قبضته ملكوت الابداع ما رضيتم ان يسمّى نفسه باسم من الأسمآء بعد
الّذى انّها و ملكوتها قد خلقت بأمره العزيز المنيع فآهء آه عن غفلتكم يا ملأ
البيان فآهء آه من احتجابكم يا ملأ المشركين

\end{otherlanguage}

Qul yā qawm innahu la-huwa alladhī law yurīdu an yajʿala kulla man fī
al-samāwāt wa-l-arḍ ḥujja bāqiya min ʿindihi la-yaqdir wa-inna hādhā
ʿindahu sahl yasīr. Wa-innahu la-huwa alladhī qad khalaqa Riḍwān
al-Bayān li-nafsihi wa-minhu badʾ kull shayʾ wa-yaʿūd law antum min
al-ʿālimīn. Wa-antum bi-lladhī kāna fī qabḍatihi malakūt al-ibdāʿ mā
raḍītum an yusammiya nafsahu bi-ism min al-asmāʾ baʿda alladhī innahā
wa-malakūtahā qad khuliqat bi-amrihi al-ʿazīz al-manīʿ. Fa-āh āh ʿan
ghaflatikum yā malaʾ al-Bayān fa-āh āh min iḥtijābikum yā malaʾ
al-mushrikīn.

Say: My people, He is indeed the One who, if He wished to make everyone
in the heavens and on earth an enduring testimony from Him, would have
the power to do so. This is easy and effortless for Him. He is indeed
the One who created the Paradise of the Bayán for Himself. From Him all
things originate and to Him they return, if you are among those who
know. Yet you were unwilling that He in whose grasp was the kingdom of
creation should call Himself by any name, even though those names and
their kingdom were created by His mighty, unassailable command. Alas,
alas for your heedlessness, O people of the Bayán! Alas, alas for your
veiling, O assembly of those who associate partners with God!

\subsection{691}\label{section-1186}

\begin{otherlanguage}{persian}

و انتم لمّا اسرفتم فى انفسكم و بلغتم الى معارج العرفان بزعمكم تذكرون
الوصاية لأحد من اعدائه و تستدلّون بها على اللّه الّذى به شرعت شرايع الأدیان
فى الأوّلین‌ و الآخرین و رجعتم الى ما استدلّ به اولو الفرقان بعد الّذى
نهيناكم فى ساحته عن كلّ الأذکار‌ الّا بعد اذنه و كان اللّه على ذلك لشهيد و
خبير اذا فانظروا فى شأنكم و عرفانكم فافّ لكم و لعقولكم ثمّ درايتكم يا ملأ
الأخسرين اما علمتم بانّا طوينا ما عند النّاس و بسطنا بساطاً آخر فتبارك اللّه
الملك الباسط العزيز الكريم

\end{otherlanguage}

Wa-antum lammā asraftum fī anfusikum wa-balaghtum ilā maʿārij al-ʿirfān
bi-zaʿmikum tadhkurūna al-waṣāya li-aḥad min aʿdāʾihi wa-tastadillūna
bihā ʿalā Allāh alladhī bihi shuriʿat sharāʾiʿ al-adyān fī al-awwalīn
wa-l-ākhirīn wa-rajaʿtum ilā mā istadalla bihi ulū al-Furqān baʿda
alladhī nahaynākum fī sāḥatihi ʿan kull al-adhkār illā baʿda idhnihi
wa-kāna Allāh ʿalā dhālika la-shahīd wa-khabīr. Idhan fa-nẓurū fī
shaʾnikum wa-ʿirfānikum fa-uff lakum wa-li-ʿuqūlikum thumma dirāyatikum
yā malaʾ al-akhsarīn. A-mā ʿalimtum bi-annā ṭawaynā mā ʿinda al-nās
wa-basaṭnā bisāṭan ākhar fa-tabāraka Allāh al-malik al-bāsiṭ al-ʿazīz
al-karīm.

When you went to excess in your own estimation and, as you imagined,
attained the heights of knowledge, you invoked a succession conferred by
testament upon one of His enemies, and used it to argue against God,
through whom the laws of the religions were ordained among earlier and
later generations. You returned to the arguments used by the people of
the Qur'an, even though We had forbidden you to make any mention in His
presence except with His permission. God is indeed a witness to this and
fully informed. Consider, then, your condition and your knowledge! Shame
on you, your reason, and your understanding, O assembly of those who
have suffered the greatest loss! Did you not know that We rolled up what
people possessed and spread out another carpet? Blessed, then, is God,
the King, the One who spreads it out, the Almighty, the Generous.

\subsection{692}\label{section-1187}

\begin{otherlanguage}{arabic}

قل يا قوم لاتفتروا على نفسى و انّى ما تكلّمت الّا بذكر هذا الظّهور و ثنائه و
ما تنفّست الّا بحبّه و ما توجّهت الّا بوجهه المشرق المنير و جعلت البيان و ما
نزل فيه ورقة من اوراق حديقة الرّضوان لنفسه المهيمن العزيز القدير ایّاکم ان
تغصبوها و ترجعوها الى الّذى اراد سفك دمى مرّه اخرى بما اتّبع النّفس و الهوى
و كان من الحاربين

\end{otherlanguage}

Qul yā qawm lā taftarū ʿalā nafsī wa-innī mā takallamtu illā bi-dhikr
hādhā al-ẓuhūr wa-thanāʾihi wa-mā tanaffastu illā bi-ḥubbihi wa-mā
tawajjahtu illā bi-wajhihi al-mushriq al-munīr wa-jaʿaltu al-Bayān wa-mā
nazala fīhi waraqa min awrāq ḥadīqat al-Riḍwān li-nafsihi al-muhaymin
al-ʿazīz al-qadīr. Iyyākum an taghṣabūhā wa-turjiʿūhā ilā alladhī arāda
safka damī marra ukhrā bi-mā ittabaʿa al-nafs wa-l-hawā wa-kāna min
al-ḥāribīn.

Say: My people, do not fabricate lies about Me. I spoke only in
remembrance and praise of this Revelation; I breathed only in love for
Him; and I turned only toward His radiant, luminous countenance. I made
the Bayán and all that was revealed in it a leaf from the leaves of the
garden of Paradise for His Self, the All-Protecting, the Almighty, the
All-Powerful. Beware lest you seize it and turn it over to the one who
wished to shed My blood once again, because he followed self and desire
and was among those who waged war.

\subsection{693}\label{section-1188}

\begin{otherlanguage}{arabic}

قد فصّلنا البيان من كلمة ثمّ رجعناه اليها و امرنا الكلمة بان يحضر تلقآء
العرش ليشهد خلق قبله و يفرح به نفسه العليم الحكيم اذا فانصفوا هل ينبغى
ان يتصرّف فيها صاحبها او دونها فما لكم يا معشر المحتجبين انّا امرنا ملأ
البيان بأن يلبسن الحرير و ينظفنّ انفسهم و اثوابهم لئلّا تقع عينه على ما لا
يحبّه و كذلك فى كلّ شیء فصّلنا تفصيلاً فى كتاب مبين كلّ ذلك لنفسه لو انتم من
المنصفين

\end{otherlanguage}

Qad faṣṣalnā al-Bayān min kalima thumma rajaʿnāhu ilayhā wa-amarnā
al-kalima bi-an yaḥḍura tilqāʾ al-ʿarsh li-yashhada khalqa qablahu
wa-yafraḥa bihi nafsuhu al-ʿalīm al-ḥakīm. Idhan fa-anṣifū hal yanbaghī
an yataṣarrafa fīhā ṣāḥibuhā aw dūnahā fa-mā lakum yā maʿshar
al-muḥtajibīn. Innā amarnā malaʾ al-Bayān bi-an yalbasna al-ḥarīr
wa-yunaẓẓifunna anfusahum wa-athwābahum li-allā taqaʿa ʿaynuhu ʿalā mā
lā yuḥibbuhu wa-kadhālika fī kull shayʾ faṣṣalnā tafṣīlan fī kitāb mubīn
kull dhālika li-nafsihi law antum min al-munṣifīn.

We unfolded the Bayán from a word, then returned it to that word, and
commanded the word to appear before the Throne, that it might witness
creation before Him and that His all-knowing, all-wise Self might
rejoice in it. Judge fairly, then: should its possessor exercise
authority over it, or someone other than it? What is the matter with
you, O company of the veiled? We commanded the people of the Bayán to
wear silk and to cleanse themselves and their clothes, lest His eyes
fall upon anything He dislikes. Likewise, concerning every thing, We set
out the details in a clear Book. All this was for His sake, if you are
among those who judge fairly.

\subsection{694}\label{section-1189}

\begin{otherlanguage}{arabic}

و خلقنا السّموات و الأرض و ما قدّر بينهما لاحبّائه فكيف جماله المشرق العزيز
المنير و انتم تمسّكتم بما قدّرناه له و اعترضتم به على محبوبی‌ فما لكم يا
ملأ البغضآء و ما يغنيكم اليوم يا معشر المفسدين و انتم اعترضتم عليه و بكلّ
ما ظهر من عنده بعد ما وصّيناكم به فى الألواح بأن كلّ من يخطر بباله ذكر
اسمه الأعظم البديع يقوم عن مقرّه و يقول سبحان اللّه ذوالملك و الملكوت تسعة
عشر مرّة ثمّ سبحان اللّه ذى العزّة و الجبروت تسعة عشر مرّة الى آخر ما نزلناه
فى لوح عزّ عظيم

\end{otherlanguage}

Wa-khalaqnā al-samāwāt wa-l-arḍ wa-mā quddira baynahumā li-aḥibbāʾihi
fa-kayfa jamāluhu al-mushriq al-ʿazīz al-munīr. Wa-antum tamassaktum
bi-mā qaddarnāhu lahu wa-iʿtaraḍtum bihi ʿalā maḥbūbī fa-mā lakum yā
malaʾ al-baghḍāʾ wa-mā yughnīkum al-yawm yā maʿshar al-mufsidīn.
Wa-antum iʿtaraḍtum ʿalayhi wa-bi-kull mā ẓahara min ʿindihi baʿda mā
waṣṣaynākum bihi fī al-alwāḥ bi-anna kulla man yakhṭuru bi-bālihi dhikr
ismihi al-aʿẓam al-badīʿ yaqūmu ʿan maqarrihi wa-yaqūlu subḥāna Allāh
dhū al-mulk wa-l-malakūt tisʿata ʿashara marra thumma subḥāna Allāh dhī
al-ʿizza wa-l-jabarūt tisʿata ʿashara marra ilā ākhir mā nazzalnāhu fī
lawḥ ʿizz ʿaẓīm.

We created the heavens and the earth and all that was ordained between
them for His loved ones---how much more, then, for His radiant, mighty,
luminous Beauty! Yet you clung to what We ordained for Him and used it
to object to My Beloved. What is the matter with you, O assembly of
hatred? What will avail you today, O company of those who spread
corruption? You objected to Him and to everything that appeared from
Him, even though We had enjoined you in the Tablets that whoever calls
to mind His most great and wondrous Name should rise from his seat and
say, ``Glory be to God, the Possessor of sovereignty and dominion,''
nineteen times, and then, ``Glory be to God, the Possessor of might and
supreme power,'' nineteen times, and so on to the end of what We
revealed in a Tablet of great glory.

\subsection{695}\label{section-1190}

\begin{otherlanguage}{arabic}

و انتم كفرتم به و بآياته و ما اكتفيتم بذلك و ما لاحظتم حقوق اللّه فى حقّه
و ما راعيتم امر اللّه فى نفسه العلىّ العليم الى ان اعترضتم بكلّ افعاله
واحدا بعد واحد و كنتم من المستهزئين و منكم من قال انّه يشرب الچاى و منكم
من قال انّه يأكل الطّعام و منكم من اعترض على لباسه بعد الّذى كلّ خيط من
خيوطه يشهد بانّه لا اله الّا هو و انّه لمقصود المقرّبين

\end{otherlanguage}

Wa-antum kafartum bihi wa-bi-āyātihi wa-mā iktafaytum bi-dhālika wa-mā
lāḥaẓtum ḥuqūq Allāh fī ḥaqqihi wa-mā rāʿaytum amr Allāh fī nafsihi
al-ʿalī al-ʿalīm ilā an iʿtaraḍtum bi-kull afʿālihi wāḥidan baʿda wāḥid
wa-kuntum min al-mustahziʾīn. Wa-minkum man qāla innahu yashrabu al-chāy
wa-minkum man qāla innahu yaʾkulu al-ṭaʿām wa-minkum man iʿtaraḍa ʿalā
libāsihi baʿda alladhī kull khayṭ min khuyūṭihi yashhadu bi-annahu lā
ilāha illā huwa wa-innahu la-maqṣūd al-muqarrabīn.

You disbelieved in Him and His verses, and were not content with that.
You did not regard the rights of God in relation to Him, nor did you
observe God's command concerning His exalted, all-knowing Self, until
you objected to all His actions, one after another, and became mockers.
Some of you said, ``He drinks tea''; some said, ``He eats food''; and
some objected to His clothing, even though every thread of it bears
witness that there is no God but Him and that He is indeed the Desire of
those brought near.

\subsection{696}\label{section-1191}

\begin{otherlanguage}{arabic}

و انّى اشهد بنفسى ما كان عند حضرته فى بعض الأحيان من ثوبين ليبدّل احدهما
بالآخر كذلك يشهد لسان صدق عليم و ما كان فى بعض اللّيالى ما يسترزقن به آل
اللّه و انّه ستر امره حفظاً لأمر اللّه المحكم المتين بعد الّذى خلق كلّ شیء
لنفسه و عنده مفتاح خزاين السّموات و الأرضين افّ لحيائكم يا ملأ البيان
تاللّه خجلت من فعلكم و اذاً اتبرّء منكم يا ملأ الشّياطين فآهء آه من ابتلائه
بينكم فآهء آه عمّا ورد و يرد عليه فى كلّ حين

\end{otherlanguage}

Wa-innī ashhadu bi-nafsī mā kāna ʿinda ḥaḍratihi fī baʿḍ al-aḥyān min
thawbayn li-yubaddila aḥadahumā bi-l-ākhar kadhālika yashhadu lisān ṣidq
ʿalīm. Wa-mā kāna fī baʿḍ al-layālī mā yastarziqna bihi āl Allāh
wa-innahu satara amrahu ḥifẓan li-amr Allāh al-muḥkam al-matīn baʿda
alladhī khuliqa kull shayʾ li-nafsihi wa-ʿindahu miftāḥ khazāʾin
al-samāwāt wa-l-arḍīn. Uff li-ḥayāʾikum yā malaʾ al-Bayān. Ta-llāh
khajiltu min fiʿlikum wa-idhan atabarraʾu minkum yā malaʾ al-shayāṭīn.
Fa-āh āh min ibtilāʾihi baynakum fa-āh āh ʿammā warada wa-yaridu ʿalayhi
fī kull ḥīn.

I Myself bear witness that at times He did not possess two garments so
that He could change from one into the other. Thus bears witness a
truthful, knowing tongue. On some nights the family of God had nothing
to sustain them. Yet He concealed His circumstances to safeguard the
firm and steadfast Cause of God, even though all things were created for
His sake and He holds the key to the treasuries of the heavens and the
earth. Shame on your modesty, O people of the Bayán! By God! I am
ashamed of your conduct, and I therefore disown you, O assembly of
devils! Alas, alas for His affliction among you! Alas, alas for what has
befallen Him and befalls Him at every moment!

\subsection{697}\label{section-1192}

\begin{otherlanguage}{arabic}

يا قوم فانصفوا ثمّ تفكّروا اقلّ من ان لو انتم فى تلك الحجبات لم اظهرت نفسى
و ما ثمر ظهورى يا ملأ المنافقين قد بعثنى اللّه لخرق الأحجاب و تطهيركم
لهذا الظّهور و انتم فعلتم به ما يتذرّف به عيناى و عيون المقدّسين قد ابيضّت
وجوه ملل القبل من فعلكم لأنكم احجب منهم و اغفل من ملأ التّورية و الزّبور و
الانجيل فيا ليت ما ولدت من امّى و ما اظهرت نفسى بينكم يا ملأ الخائبين
فو‌الّذى بعثنى بالحقّ احصيت علم كلّ شیء و كلّ ما كنز فى كنائز حفظ اللّه و ما
ستر عن انظر العالمين و لكن ما احصيت نفوساً اشقى منكم و ابعد عنكم لانّا بعد
ما فصّلناه فى الألواح و ما نصحنا به انفسكم فى كلّ الأوراق ما ظننّا بأن يظهر
فى الملك احد ان يعترض على اللّه الّذى فى قبضته ملكوت ملك السّموات و الأرضين
اذا تحيّرنا من خلقكم و لم ادر باىّ كلمة خلقتم يا من تحيّر فيكم و من فعلكم
افئدة اهل ملأ العالين ثمّ افئدة المخلصين و المقرّبين كذلك قصصنا لك يا عبد
فى هذا اللّوح ما تغرّدت به حمامة البيان حينئذ لدى عرش ربّك العزيز الحميد و
انّك انت فاقرأ ما نزل فيه ثمّ احفظ لؤلؤ المعانى عن كلّ خائن سارق من ملأ
الشّياطين

\end{otherlanguage}

Yā qawm fa-anṣifū thumma tafakkarū aqall min an law antum fī tilka
al-ḥujubāt lam aẓhartu nafsī wa-mā thamar ẓuhūrī yā malaʾ al-munāfiqīn
qad baʿathanī Allāh li-kharq al-aḥjāb wa-taṭhīrikum li-hādhā al-ẓuhūr
wa-antum faʿaltum bihi mā yatadharrafa bihi ʿaynāya wa-ʿuyūn
al-muqaddasīn qad ibyaḍḍat wujūh milal al-qabl min fiʿlikum li-annakum
aḥjab minhum wa-aghfal min malaʾ al-Tawriyya wa-al-Zabūr wa-al-Injīl
fa-yā layta mā wulidtu min ummī wa-mā aẓhartu nafsī baynakum yā malaʾ
al-khāʾibīn fa-wa-alladhī baʿathanī bi-al-ḥaqq aḥṣaytu ʿilm kull shayʾ
wa-kull mā kuniza fī kanāʾiz ḥifẓ Allāh wa-mā sutira ʿan anẓar
al-ʿālamīn wa-lākin mā aḥṣaytu nufūsan ashqā minkum wa-abʿad ʿankum
li-annā baʿda mā faṣṣalnāhu fī al-alwāḥ wa-mā naṣaḥnā bihi anfusakum fī
kull al-awrāq mā ẓanannā bi-an yaẓhara fī al-mulk aḥad an yaʿtariḍa ʿalā
Allāh alladhī fī qabḍatihi malakūt mulk al-samāwāt wa-al-araḍīn idhan
taḥayyarnā min khalqikum wa-lam adri bi-ayy kalima khuliqtum yā man
taḥayyara fīkum wa-min fiʿlikum afʾidat ahl malaʾ al-ʿālīn thumma
afʾidat al-mukhliṣīn wa-al-muqarrabīn kadhālika qaṣaṣnā laka yā ʿabd fī
hādhā al-lawḥ mā tagharradat bihi ḥamāmat al-bayān ḥīnaʾidh ladā ʿarsh
rabbika al-ʿAzīz al-Ḥamīd wa-innaka anta fa-iqraʾ mā nuzzila fīhi thumma
iḥfaẓ luʾluʾ al-maʿānī ʿan kull khāʾin sāriq min malaʾ al-shayāṭīn.

O people, be fair, then reflect. At the very least, had you remained in
those veils, I would not have manifested Myself. What is the fruit of My
appearance, O assembly of hypocrites? God sent Me to tear away the veils
and purify you for this Revelation, yet you have done to Him what makes
My eyes and the eyes of the sanctified shed tears. Your conduct has
vindicated the peoples of the past, for you are more veiled than they
and more heedless than the people of the Torah, the Psalms, and the
Gospel. If only My mother had not given birth to Me, and I had not
manifested Myself among you, O assembly of the disappointed! By Him who
sent Me with the truth! I have encompassed the knowledge of all things,
everything treasured in the repositories of God's keeping, and whatever
is concealed from the most discerning of the world's people. Yet I have
not counted any souls more wretched than you or more distant from you.
For after all We set forth in detail in the Tablets, and all the counsel
We gave you in every page, We did not think that anyone would appear in
the world to object to God, in whose grasp is the dominion of the
kingdom of the heavens and the earth. Thus We were bewildered at your
creation, and I do not know by what word you were created, O you at
whom, and at whose conduct, the hearts of the inhabitants of the exalted
Assembly, and then the hearts of the sincere and those brought near, are
bewildered! Thus, O servant, We have recounted to you in this Tablet
what the Dove of Utterance sang at that time before the throne of your
Lord, the Almighty, the All-Praised. Read what has been revealed in it,
then guard the pearls of meaning from every treacherous thief among the
assembly of devils.

\subsection{698}\label{section-1193}

\begin{otherlanguage}{arabic}

و ان وجدت من ذى بصر فانشرها امام عينيه ليشهد و يكون من الفائزين لعلّ اولى
الأبصار من عبادنا الأخيار يطّلعن بما ورد على جمال المختار من هؤلآء الفجّار
الّذين اتّخذوا العجل لانفسهم ربّا من دون اللّه و يسجدونه فى العشىّ و الأبكار
و يكونن من الفرحين و انّك انت لا تحزن عمّا ورد علينا ثمّ اصبر كما صبرنا و
انّه لخير ناصر و معين ان اذكر ربّك فى اللّيالى و الأیّام ثمّ انطق بثنآء نفسه
بين عباده لعلّ بثنائه تحدث نار حبّه فى قلوب المحسنين و كلّ يقومنّ على ثنآء
اللّه ربّهم و ربّ ما يرى و ما لايرى و ربّ آبائكم الأوّلین‌ و الحمد للّه ربّ
العالمين

\end{otherlanguage}

Wa-in wajadta min dhī baṣar fa-unshurhā amāma ʿaynayhi li-yashhada
wa-yakūna min al-fāʾizīn laʿalla ulī al-abṣār min ʿibādinā al-akhyār
yaṭṭaliʿna bi-mā warada ʿalā jamāl al-Mukhtār min hāʾulāʾ al-fujjār
alladhīna ittakhadhū al-ʿijl li-anfusihim rabban min dūn Allāh
wa-yasjudūnahu fī al-ʿashiyy wa-al-ibkār wa-yakūnanna min al-fariḥīn
wa-innaka anta lā taḥzan ʿammā warada ʿalaynā thumma iṣbir kamā ṣabarnā
wa-innahu la-khayr nāṣir wa-muʿīn an udhkur rabbaka fī al-layālī
wa-al-ayyām thumma inṭiq bi-thanāʾ nafsihi bayna ʿibādihi laʿalla
bi-thanāʾihi taḥduthu nār ḥubbihi fī qulūb al-muḥsinīn wa-kull yaqūmanna
ʿalā thanāʾ Allāh rabbihim wa-rabb mā yurā wa-mā lā yurā wa-rabb
ābāʾikum al-awwalīn wa-al-ḥamdu li-llāh rabb al-ʿālamīn.

If you find someone with insight, spread them before his eyes, so that
he may witness and be among those who attain their goal. Perhaps those
endowed with insight among Our righteous servants will learn what has
befallen the Beauty of the Chosen One at the hands of these wicked
people, who have taken the calf as their lord instead of God, prostrate
themselves before it evening and morning, and rejoice. Do not grieve
over what has befallen Us, but be patient as We have been patient. He is
indeed the best helper and supporter. Remember your Lord by night and
day, then voice His praise among His servants, so that through His
praise the fire of His love may arise in the hearts of those who do
good, and all may arise to praise God, their Lord and the Lord of all
that is seen and unseen, and the Lord of your forefathers. Praise be to
God, the Lord of the worlds.

\subsection{699}\label{section-1194}

\begin{otherlanguage}{persian}

هر صاحب بصر و ادراكى از اين آيات مذكوره البتّه بر حقّ و نفس خود هر دو نوحه
نمايد امّا بر حقّ منيع چه كه در دست امثال آن ظالمين مبتلى شده بشأنيكه عين
ابداع نديده و سمع اختراع نشنيده و امّا بر نفس خود چه كه بشأنيكه بايد بر
اين امر بديع قيام نمايد ننموده‌نسئل اللّه بأن يبعث بسلطانه خلقا ليأت بما
فات من هؤلآء كه شايد آن نفوس مقدّسه بآنچه لايق و سزاوار اين ظهوراست قيام
نمايند من يطوف حول الرّوح يقول لا تستر فاخرق الأحجاب ثمّ اظهر ما كنز فى
السّرّ من علم اللّه المهيمن القيّوم و الّذين يطوفن فى حول العقل يقولن ایّاک
ایّاک ان اصبر و لاتفش ما يفزع به هياكل المغلّين و تميّز به صدور المشركين ان
الق على العباد ما ينبغى لشأنهم و مقامهم ایّاک ان تلق على الرّضيع مقامات
البلوغ او على المحروم اسرار حرم اللّه او على الأصمّ نغمات اللّه العزيز
المحبوب

\end{otherlanguage}

Har ṣāḥeb-e baṣar va edrākī az īn āyāt-e madhkūre albatte bar Ḥaqq va
nafs-e khod har do nowḥe namāyad ammā bar Ḥaqq-e manīʿ che ke dar dast-e
amthāl-e ān ẓālemīn mobtalā shode be-shaʾnī ke ʿayn-e ebdāʿ nadīde va
samʿ-e ekhterāʿ nashenīde va ammā bar nafs-e khod che ke be-shaʾnī ke
bāyad bar īn amr-e badīʿ qīyām namāyad nanamūde nasʾalu Allāh bi-an
yabʿatha bi-sulṭānihi khalqan li-yaʾta bi-mā fāta min hāʾulāʾ ke shāyad
ān nofūs-e moqaddase be-ānche lāyeq va sazāvār-e īn ẓohūr ast qīyām
namāyand man yaṭūfu ḥawla al-Rūḥ yaqūlu lā tastur fa-ikhruq al-aḥjāb
thumma aẓhir mā kuniza fī al-sirr min ʿilm Allāh al-Muhaymin al-Qayyūm
wa-alladhīna yaṭūfna fī ḥawl al-ʿaql yaqūlna iyyāka iyyāka an iṣbir
wa-lā tufshi mā yafzaʿu bihi hayākil al-mughillīn wa-tamayyazu bihi
ṣudūr al-mushrikīn an alqi ʿalā al-ʿibād mā yanbaghī li-shaʾnihim
wa-maqāmihim iyyāka an tulqi ʿalā al-raḍīʿ maqāmāt al-bulūgh aw ʿalā
al-maḥrūm asrār ḥaram Allāh aw ʿalā al-aṣamm naghamāt Allāh al-ʿAzīz
al-Maḥbūb.

Anyone endowed with sight and understanding will surely, on reading
these verses, lament both for God and for himself: for the exalted God,
because He has been afflicted at the hands of such oppressors in a
manner that the eye of creation has never seen and the ear of existence
has never heard; and for himself, because he has not arisen to serve
this wondrous Cause as he should. We ask God to raise up, through His
sovereignty, a people who will make good what these people have failed
to do, so that those sanctified souls may arise to do what is fitting
and worthy of this Revelation. The one who circles around the Spirit
says, ``Do not conceal it! Tear away the veils, then reveal what has
been secretly treasured of the knowledge of God, the All-Protecting, the
Self-Subsisting.'' And those who circle around the intellect say,
``Beware, beware! Be patient, and do not divulge what would terrify the
embodiments of malice and tear apart the breasts of those who associate
partners with God. Impart to the servants what befits their condition
and station. Beware of imparting the stations of maturity to a suckling
infant, or the mysteries of God's sanctuary to one who is deprived, or
the melodies of God, the Almighty, the Beloved, to the deaf.''

\subsection{700}\label{section-1195}

\begin{otherlanguage}{persian}

اى غير منصف غير بصير چرا از شجره اعراض سؤال ننمودى كه چرا نفسى را كه
بيست سال باسمش معروف و موسوم بودى و بكلماتش از براى خودشان ثابت نموده‌
حال فتوى بر قتلش داده و عامل شده آنچه را كه احدى عامل نشده و در كلّ حين
بمكر جديد در تضييع امرش مشغول و حال آنكه تو و امثال تو كلّ از كلمه او خلق
شده و لكن چه ذكر شود كه آنچه گفته شود ابدا در نفس مغلّين و مشركين تأثير
ننموده‌و نخواهد نمود

\end{otherlanguage}

Ey gheyr-e monṣef-e gheyr-e baṣīr cherā az shajare-ye eʿrāḍ soʾāl
nanamūdī ke cherā nafsī rā ke bīst sāl be-esmash maʿrūf va mowsūm būdī
va be-kalemātash az barā-ye khodeshān thābet namūde ḥāl fatvā bar
qatl-ash dāde va ʿāmel shode ānche rā ke aḥadī ʿāmel nashode va dar
koll-e ḥīn be-makr-e jadīd dar taḍyīʿ-e amrash mashghūl va ḥāl ānke to
va amthāl-e to koll az kaleme-ye ū khalq shode va lāken che dhekr shavad
ke ānche gofte shavad abadan dar nafs-e moghellīn va moshrekīn taʾthīr
nanamūde va nakhāhad namūd.

O unjust and undiscerning one! Why did you not ask the tree of aversion
why he has now issued a decree for the death of the One by whose name
you were known and designated for twenty years, and by whose words they
established their own standing? He has done what no one has ever done,
and at every moment he is occupied with a new scheme to undermine His
Cause, even though you and all those like you were created by His word.
But what can be said? Whatever is said has never had, and will never
have, any effect on malicious souls and those who associate partners
with God.

\subsection{701}\label{section-1196}

\begin{otherlanguage}{persian}

و اين بسى واضح است كه اگر نفسى بجذوه از نار حبّ الهى فايز ميشد ابدا باين
كلمات مشغول نميگشت و باين اعتراضات بر سلطان اسماء و صفات وارد نمياورد
آنچه شبه آن در ابداع ظاهر نشده فسوف يبعثن اللّه خلقا ينتبّهنّ من كلماته و
يعرفنّ مراده و يتوجّهنّ اليه و ينقطعنّ عن العالمين

\end{otherlanguage}

Va īn basī vāḍeḥ ast ke agar nafsī be-jadhve az nār-e ḥobb-e elāhī fāʾez
mīshod abadan be-īn kalemāt mashghūl namīgasht va be-īn eʿterāḍāt bar
Solṭān-e asmāʾ va ṣefāt vāred namīāvard ānche shebh-e ān dar ebdāʿ ẓāher
nashode fa-sawfa yabʿathanna Allāh khalqan yantabihunna min kalimātihi
wa-yaʿrifanna murādahu wa-yatawajjahunna ilayhi wa-yanqaṭiʿunna ʿan
al-ʿālamīn.

It is abundantly clear that if a soul had attained even a spark of the
fire of divine love, he would never have occupied himself with these
words or, through these objections, inflicted on the Sovereign of names
and attributes what has no precedent in creation. God will assuredly
raise up a people who will be awakened by His words, recognize His
intent, turn toward Him, and detach themselves from all the worlds.

\subsection{702}\label{section-1197}

\begin{otherlanguage}{persian}

و اينكه نوشته كه جواب ميرزا محمّد على و ميرزا ابو الحسن و ديگران را چه
بگويم شما قاضى كاينات نبوده و نيستيد قاضى نفس خود باش و فكر خود را كن كه
جواب خدا را چه بگوئى جواب خلق آسانست اى بی‌‌شرم بی‌ حيا در جواب مخلوقى چند
اظهار تحيّر نموده‌ و ابدا در فكر جواب حقّ نيفتاده حال انصاف ده تو كه از چند
نفر مخلوق اظهار شرم و تحيّر از جواب نموده‌ اين عباد در مقرّ عرش اعظم
مقاميكه جميع انبیاء و مقرّبين حاضرند جواب چه بگوئيم فاشهد نفسك فى ذلك
الموقف ثمّ ملأ البيان كلّهم اجمعين و نقطة الأولی عن يمين العرش يخاطبك و
ايّاهم بفارسى

\end{otherlanguage}

Va īnke neveshte ke javāb-e Mīrzā Moḥammad ʿAlī va Mīrzā Abu al-Ḥasan va
dīgarān rā che begūyam shomā qāḍī-ye kāʾenāt nabūde va nīstīd qāḍī-ye
nafs-e khod bāsh va fekr-e khod rā kon ke javāb-e Khodā rā che begūʾī
javāb-e khalq āsān ast ey bī-sharm-e bī-ḥayā dar javāb-e makhlūqī chand
eẓhār-e taḥayyor namūde va abadan dar fekr-e javāb-e Ḥaqq nayoftāde ḥāl
enṣāf de to ke az chand nafar makhlūq eẓhār-e sharm va taḥayyor az javāb
namūde īn ʿebād dar maqarr-e ʿarsh-e aʿẓam maqāmī ke jamīʿ-e anbīyā va
moqarrabīn ḥāḍerand javāb che begūʾīm fa-ishhad nafsaka fī dhālika
al-mawqif thumma malaʾ al-bayān kullahum ajmaʿīn wa-Nuqṭat al-Ūlā ʿan
yamīn al-ʿarsh yukhāṭibuka wa-iyyāhum be-fārsī.

As for what you wrote, ``What answer should I give Mírzá Muḥammad ʿAlí,
Mírzá Abu'l-Ḥasan, and the others?''---you have never been, and are not,
the judge of all creation. Be the judge of your own self, and consider
what answer you will give God. Answering His creatures is easy. O
shameless and immodest one! You have expressed bewilderment about
answering a few created beings, yet have never given a thought to
answering God. Now be fair: you have expressed shame before a few
created beings and bewilderment about how to answer them; what answer
are we servants to give at the seat of the Most Great Throne, where all
the Prophets and those brought near are present? Behold yourself at that
gathering, and then the entire assembly of the Bayán, all together,
while the Primal Point, at the right hand of the throne, addresses you
and them in Persian.

\subsection{703}\label{section-1198}

\begin{otherlanguage}{persian}

ذكر ميشود تا درست ادراك نمائى ميفرمايد اى ملأ بيان آيا بشما خبر ندادم در
كلّ الواح كه بعد از من ظهورى ظاهر خواهد شد البتّه كلّ ميگوئيد بلى و نفسك
الحقّ در جميع الواح بشارت فرموده و اخذ عهد آن نيّر اعظم را از كلّ فرموده
بعد ميفرمايند كه آيا نگفته‌ام كه آن جمال قدم بآيات اللّه ظاهر ميشود و آيا
نفرموده‌ام كه طوريّون در آن مقرّ منصعق و ميّت ميشوند و آيا نفرموده‌ام كه اعلى
شجره اثبات از ادنى شجره نفى ميشود و آيا نفرموده‌ام كه هر ذى نورى نزد ضياء
انوار وجهش مظلم مشاهده ميشود و آيا نفرموده‌ام كه بنفس ظهور احدى جز ذاتش
مطّلع نه و آيا نفرموده‌ام در حين ظهور بآنچه در بيان نازل شده از او محتجب
نشويد و آيا نفرموده‌ام بحروفات بيان كه اوّل خلق بيانند از آن جمال قدم
محروم نمانيد و آيا نفرموده‌ام كلّ بيان ورقه ايست از اوراق جنّت نزد آن شجره
احديّت آيا نفرموده‌ام كه او قادر است بر آنكه نفسى را كه اعلى است از كلّ خلق
ردّ فرمايد و آيا نفرموده‌ام كه كلّ بيان خلق اويند در ظهور قبل او

\end{otherlanguage}

Dhekr mīshavad tā dorost edrāk namāʾī mīfarmāyad ey malaʾ-e Bayān āyā
be-shomā khabar nadādam dar koll-e alvāḥ ke baʿd az man ẓohūrī ẓāher
khāhad shod albatte koll mīgūʾīd balā wa-nafsika al-Ḥaqq dar jamīʿ-e
alvāḥ beshārat farmūde va akhdh-e ʿahd-e ān nayyer-e aʿẓam rā az koll
farmūde baʿd mīfarmāyand ke āyā nagofte-am ke ān Jamāl-e Qedam be-āyāt-e
Allāh ẓāher mīshavad va āyā nafarmūde-am ke Ṭūrīyūn dar ān maqarr
monṣaʿeq va mayyet mīshavand va āyā nafarmūde-am ke aʿlā shajare-ye
ethbāt az adnā shajare-ye nafy mīshavad va āyā nafarmūde-am ke har dhī
nūrī nazd-e ḍīyā-ye anvār-e vajhash moẓlem moshāhede mīshavad va āyā
nafarmūde-am ke be-nafs-e ẓohūr aḥadī joz dhātash moṭṭaleʿ na va āyā
nafarmūde-am dar ḥīn-e ẓohūr be-ānche dar Bayān nāzel shode az ū
moḥtajeb nashavīd va āyā nafarmūde-am be-ḥorūfāt-e Bayān ke avval-e
khalq-e Bayān-and az ān Jamāl-e Qedam maḥrūm namānīd va āyā nafarmūde-am
koll-e Bayān varaqe-īst az owrāq-e jannat nazd-e ān shajare-ye aḥadīyyat
āyā nafarmūde-am ke ū qāder ast bar ānke nafsī rā ke aʿlā ast az koll-e
khalq radd farmāyad va āyā nafarmūde-am ke koll-e Bayān khalq-e ūyand
dar ẓohūr-e qabl-e ū.

It is recounted so that you may understand correctly. He says, ``O
assembly of the Bayán! Did I not inform you in all the Tablets that
after Me a Revelation would appear?'' You will all surely say, ``Yes, by
Your own true Self! In all the Tablets You gave these glad tidings and
received from everyone a covenant concerning that Most Great Luminary.''
He then says, ``Did I not say that that Ancient Beauty would appear with
the verses of God? Did I not say that the people of Sinai would be
struck down and die in that presence? Did I not say that the highest
tree of affirmation would become one of the lowest trees of negation?
Did I not say that every luminous being would be seen as dark before the
radiance of the lights of His face? Did I not say that no one but His
own Self knows the Revelation itself? Did I not say that, at the time of
His appearance, you should not let what has been revealed in the Bayán
veil you from Him? Did I not tell the Letters of the Bayán, who are the
first of the creation of the Bayán, not to remain deprived of that
Ancient Beauty? Did I not say that the whole Bayán is but a leaf among
the leaves of paradise before that Tree of Oneness? Did I not say that
He is able to reject a soul who is the most exalted of all creation? Did
I not say that all the people of the Bayán are His creation in His
previous Revelation?''

\subsection{704}\label{section-1199}

\begin{otherlanguage}{persian}

البتّه جميع را تصديق مينمائيد چه كه كلّ در الواح بيان منصوصا نازل و مسطور
مع اين وصاياى لايحصى و نصايح محكمه متقنه بچه دليل و برهان باينجمال مبارك
سيف جفا كشيده‌ايد و جميع بر هتك حرمتش قيام نموده‌ايد و اراده سفك دم مطهّرش
کرده‌‌ايد و از جميع اشطار باسياف نفاق هجوم آورده‌ايد مع آنكه بحجّتى ظاهر شده
كه من از قبل بآن ظاهر شده‌ام و هم‌چنين از قبل قبل محمّد رسول اللّه و از قبل
قبل قبل روح اگر بگوئيد از كجا اين آيات من عند اللّه است ميگويم بهمان دليل
كه آيات مرا ادراك نموده‌ايد كه من عند اللّه است مع آنكه جميع مقرّيد كه اين
آيات بديعه در كلّ اقطار منتشر و در كلّ حين بآياتى تكلّم فرموده كه كلّ نزد
كلمه آن معدوم بوده و خواهند بود مع‌ذلك شما نفسى را كه محارب با حقّ است او
را ربّ اخذ نموده‌ايد بعضى متمسّك شده‌ايد كه من او را ثمره خوانده‌ام مع آنكه
در كتاب من نازل كه اگر مشاهده شود بر من غصنى يا ورقى و ثمرى كه سجده
ننمايد از براى آن شمس ظهور در يوم ظهور او نيست از من و راجع نميشود بسوى
من

\end{otherlanguage}

Albatte jamīʿ rā taṣdīq mīnamāʾīd che ke koll dar alvāḥ-e Bayān manṣūṣan
nāzel va masṭūr maʿ-e īn vaṣāyā-ye lā yoḥṣā va naṣāyeḥ-e moḥkame-ye
motqane be-che dalīl va borhān be-īn Jamāl-e mobārak seyf-e jafā
keshīde-īd va jamīʿ bar hatk-e ḥormatash qīyām namūde-īd va erāde-ye
safk-e dam-e moṭahharash karde-īd va az jamīʿ-e ashṭār be-asyāf-e nefāq
hojūm āvarde-īd maʿ-e ānke be-ḥojjatī ẓāher shode ke man az qabl be-ān
ẓāher shode-am va ham-chenīn az qabl-e qabl Moḥammad Rasūl Allāh va az
qabl-e qabl-e qabl Rūḥ agar begūʾīd az kojā īn āyāt min ʿind Allāh ast
mīgūyam be-hamān dalīl ke āyāt-e marā edrāk namūde-īd ke min ʿind Allāh
ast maʿ-e ānke jamīʿ moqerrīd ke īn āyāt-e badīʿe dar koll-e aqṭār
montasher va dar koll-e ḥīn be-āyātī takallom farmūde ke koll nazd-e
kaleme-ye ān maʿdūm būde va khāhand būd maʿa dhālek shomā nafsī rā ke
moḥāreb bā Ḥaqq ast ū rā rabb akhdh namūde-īd baʿḍī motamassek shode-īd
ke man ū rā thamare khānde-am maʿ-e ānke dar ketāb-e man nāzel ke agar
moshāhede shavad bar man ghoṣnī yā varaqī va thamarī ke sajde nanamāyad
az barā-ye ān shams-e ẓohūr dar yowm-e ẓohūr-e ū nīst az man va rājeʿ
namīshavad be-sū-ye man.

``You will surely acknowledge all of these, for they are all explicitly
revealed and recorded in the Tablets of the Bayán. Despite these
countless counsels and firm, conclusive admonitions, by what evidence or
proof have you drawn the sword of oppression against this blessed
Beauty, all arisen to violate His sanctity, sought to shed His pure
blood, and attacked Him from every side with the swords of hypocrisy?
Yet He has appeared with the very proof with which I appeared before,
and with which Muḥammad, the Messenger of God, appeared before that, and
the Spirit before that. If you ask, `How do we know that these verses
are from God?' I say: by the same evidence through which you recognized
that My verses are from God. Moreover, you all acknowledge that these
wondrous verses have spread throughout every region, and that at every
moment He has uttered verses before a single word of which all have
been, and will be, as nothing. Despite this, you have taken as your lord
someone who is at war with God. Some of you cling to the fact that I
called him a fruit, although it is revealed in My Book that if a branch,
a leaf, or a fruit is seen upon Me that does not prostrate itself before
that Sun of Revelation on the day of His appearance, it is not of Me and
does not return to Me.''

\subsection{705}\label{section-1200}

\begin{otherlanguage}{persian}

و بعضى متمسّك كه من او را مرآت اخذ نموده‌ام مع آنكه صد‌هزار امثال او بنصّ من
از يك تجلّى شمس ظاهر و مشرق و حال آنكه امر آن مرآت را هم ندانسته‌ايد بر
فرض تسليم كه او آن مرآت باشد شكّى نيست كه بقول من او را مرآت دانسته‌ايد و
من منصوصاً فرموده‌ام كه از آنچه نازل شده در كتاب از آن مظهر ربّ الأرباب‌
محجوب نمانيد و هم‌چنين فرموده‌ام كه طوريّون نزد آن ظهور اعظم ميّت و لاشیء
ميشوند و صد‌هزار امثال آن مرايا نزد نفسى از آن نفوس طوريّه مشرقه از افق
احديّه معدوم صرف بوده و خواهند بود

\end{otherlanguage}

Va baʿḍī motamassek ke man ū rā merʾāt akhdh namūde-am maʿ-e ānke ṣad
hezār amthāl-e ū be-naṣṣ-e man az yek tajallī-ye shams ẓāher va moshreq
va ḥāl ānke amr-e ān merʾāt rā ham nadāneste-īd bar farḍ-e taslīm ke ū
ān merʾāt bāshad shakkī nīst ke be-qowl-e man ū rā merʾāt dāneste-īd va
man manṣūṣan farmūde-am ke az ānche nāzel shode dar ketāb az ān maẓhar-e
Rabb al-Arbāb maḥjūb namānīd va ham-chenīn farmūde-am ke Ṭūrīyūn nazd-e
ān ẓohūr-e aʿẓam mayyet va lā shayʾ mīshavand va ṣad hezār amthāl-e ān
marāyā nazd-e nafsī az ān nofūs-e Ṭūrīyye-ye moshreqe az ofoq-e aḥadīyye
maʿdūm-e ṣerf būde va khāhand būd.

``Some cling to the fact that I designated him a mirror, although,
according to My explicit text, a hundred thousand like him appear and
shine forth through a single effulgence of the Sun. Yet you have not
even understood the matter of that mirror. Even supposing that he is
that mirror, there is no doubt that you have recognized him as a mirror
because of My words. And I explicitly stated that you should not allow
what has been revealed in the Book to veil you from that Manifestation
of the Lord of Lords. I also stated that the people of Sinai would
become dead and as nothing before that Most Great Revelation. A hundred
thousand like those mirrors have been, and will be, sheer nothingness
before a single one of those Sinaitic souls shining from the horizon of
oneness.''

\subsection{706}\label{section-1201}

\begin{otherlanguage}{persian}

و اگر بگوئيد كه اين مرآت هم حال همين دعوى را نموده‌شكّى نيست كه در اين قول
كاذب بوده چه كه مخصوص نازل فرموده‌ام كه اگر مرآت دعوى شمسى نمايد نزد شمس
واضح است كه او شبحى است از او و از همين بيان من محقّق كه مرآت على زعم شما
دعوى شمسى خواهد نمود و همين نفس ادّعاى او در اين ایّام دليل بر كذب او است
چه كه بعد از ظهور اشراق شمس احديّه اعتراضاً على نفسه و بغياً على كينونته و
استكباراً على ذاته و الحاداً فى امره باين كلمات متكلّم شده و اكثرى از شما
هم متابعت او نموده‌ايد

\end{otherlanguage}

Va agar begūʾīd ke īn merʾāt ham ḥāl hamīn daʿvā rā namūde shakkī nīst
ke dar īn qowl kādheb būde che ke makhṣūṣ nāzel farmūde-am ke agar
merʾāt daʿvā-ye shamsī namāyad nazd-e shams vāḍeḥ ast ke ū shabaḥī ast
az ū va az hamīn bayān-e man moḥaqqaq ke merʾāt ʿalā zaʿm-e shomā
daʿvā-ye shamsī khāhad namūd va hamīn nafs-e eddeʿā-ye ū dar īn ayyām
dalīl bar kedhb-e ū ast che ke baʿd az ẓohūr-e eshrāq-e shams-e aḥadīyye
iʿtirāḍan ʿalā nafsihi wa-baghyan ʿalā kaynūnatihi wa-istikbāran ʿalā
dhātihi wa-ilḥādan fī amrihi be-īn kalemāt motakallem shode va aktharī
az shomā ham motābaʿat-e ū namūde-īd.

``If you say that this mirror has now made the same claim, there is no
doubt that he has lied in saying this. For I specifically revealed that
if a mirror should claim to be the sun, it is clear to the Sun that it
is only an image of Him. From this very statement of Mine, it is
established that the mirror, as you regard him, would claim to be the
sun. His very claim in these days is proof of his falsehood, for after
the dawning radiance of the Sun of Oneness appeared, he uttered these
words in opposition to His Self, in rebellion against His Being, in
arrogance before His Essence, and in impiety toward His Cause. And most
of you have followed him.''

\subsection{707}\label{section-1202}

\begin{otherlanguage}{arabic}

افّ لكم يا ملأ البيان ثمّ افّ لكم يا ملأ البيان ثمّ افّ لكم يا ملأ البيان ثمّ
افّ لكم الى آخر الّذى لا آخر له لانّ من فعلكم ذاب قلبی‌ و احترق كبدى و تشبّك
حشائى و تقطّعت اعضائى و تزلزلت اركانى و اقشعرّ جلدى و اهتزّت نفسى و جنّت
كينونتى و ضجّ سرّى و حزن ظاهرى و باطنى و اوّلى و آخرى

\end{otherlanguage}

Uffin lakum yā malaʾa al-Bayān, thumma uffin lakum yā malaʾa al-Bayān,
thumma uffin lakum yā malaʾa al-Bayān, thumma uffin lakum ilā ākhiri
alladhī lā ākhira lahu, li-anna min fiʿlikum dhāba qalbī wa-iḥtaraqa
kabidī wa-tashabbaka ḥashāʾī wa-taqaṭṭaʿat aʿḍāʾī wa-tazalzalat arkānī
wa-iqshaʿarra jildī wa-ihtazzat nafsī wa-junnat kaynūnatī wa-ḍajja sirrī
wa-ḥazina ẓāhirī wa-bāṭinī wa-awwalī wa-ākhirī.

``Shame on you, O people of the Bayān! Again, shame on you, O people of
the Bayān! Again, shame on you, O people of the Bayān! Again, shame on
you, without end! For because of your deeds My heart has melted, My
liver has burned, My entrails have become entangled, My limbs have been
torn apart, My frame has been shaken, My skin has shuddered, My soul has
trembled, My being has been driven mad, My inmost self has cried out,
and My outward and inward being, My beginning and My end, have been
overcome with sorrow.''

\subsection{708}\label{section-1203}

\begin{otherlanguage}{persian}

آنچه شما ذكر نموده‌و مينمائيد بمن نسبت ميدهيد و بقول من باو متمسّك شده‌ايد
و حجّيت نفس من ببيان ثابت و مخصوص فرموده‌ام ببيان از نفس رحمن محتجب نمانيد
لم ادر بعد ذلك باىّ عذر تتمسّكون و باىّ حبل تتشبّثون بامر من ناظر نيستيد و
نصح اللّه را نپذيرفتيد نفسى را كه حال محارب باللّه است باسمى از اسماء
موسوم نموده‌ بنفس من محاربه و مجادله مينمائيد مع آنكه مظاهر كلّ اسماء را
تصريحا فرموده‌ام كه در آن ظهور عزّ تقديس خادمند وظائف چنانچه فرموده‌ام

\end{otherlanguage}

Ānchih shumā dhikr namūdih va mīnamāyīd bi-man nisbat mīdahīd va
bi-qawl-i man bi-ū mutamassik shudih-īd va ḥujjiyyat-i nafs-i man
bi-Bayān thābit va makhṣūṣ farmūdih-am bi-Bayān az nafs-i Raḥmān
muḥtajib namānīd. Lam adri baʿda dhālika bi-ayyi ʿudhrin tatamassakūna
wa-bi-ayyi ḥablin tatashabbathūn. Bi-amr-i man nāẓir nīstīd va nuṣḥ-i
Allāh rā napadhīruftīd. Nafsī rā kih ḥāl muḥārib bi-Allāh ast bi-ismī az
asmāʾ mawsūm namūdih bi-nafs-i man muḥāribih va mujādilih mīnamāyīd,
maʿa ānkih maẓāhir-i kull-i asmāʾ rā taṣrīḥan farmūdih-am kih dar ān
ẓuhūr-i ʿizz-i taqdīs khādimand vaẓāʾif, chunānchih farmūdih-am.

``What you have mentioned and continue to mention, you attribute to Me,
and it is on the strength of My words that you have clung to him. The
authority of My own Self is established through the Bayān, and I have
specifically stated that you must not let the Bayān veil you from the
Person of the All-Merciful. I do not know, after this, what excuse you
will cling to or what cord you will grasp. You pay no heed to My command
and have not accepted God's counsel. You have designated a soul who is
now at war with God by one of the names, and you wage war and dispute
with My own Self, even though I have explicitly stated that the
manifestations of all the names are servants in their duties in that
Revelation of sanctified glory, as I have stated:''

\subsection{709}\label{section-1204}

\begin{otherlanguage}{arabic}

ان اقول انّه واحد فذلك اسم من اسمائه وان اقول انّه سبّوح فظهر ذلك الاسم اسم
من اسمائه و ان اقول انّه قدّوس فهو الّذى يتجلّى بذلك الاسم بامثاله و ان اقول
انّه عزيز فكلّ عزيز ساجد لعلوّ عزّته و ان اقول انّه محبوب فكلّ محبوب قد شغف فى
حبّ عرفانه و جلال عزّته و ان اقول انّه مهيمن فهو الّذى من يستدل يوم ظهوره قد
جعله اللّه مهيمنا على كلّ شیء بأمره و ان اقول انّه قيّوم فذلك مظهر احد من
قوّام بساط احديّته فتعالى تعالى ذكره

\end{otherlanguage}

In aqūlu innahu wāḥidun fa-dhālika ismun min asmāʾihi wa-in aqūlu innahu
subbūḥun fa-ẓahara dhālika al-ismu ismun min asmāʾihi wa-in aqūlu innahu
quddūsun fa-huwa alladhī yata jallā bi-dhālika al-ismi bi-amthālihi
wa-in aqūlu innahu ʿazīzun fa-kullu ʿazīzin sājidun li-ʿuluwwi ʿizzatihi
wa-in aqūlu innahu maḥbūbun fa-kullu maḥbūbin qad shaghifa fī ḥubbi
ʿirfānihi wa-jalāli ʿizzatihi wa-in aqūlu innahu muhayminun fa-huwa
alladhī man yastadillu yawma ẓuhūrihi qad jaʿalahu Allāhu muhayminan
ʿalā kulli shayʾin bi-amrihi wa-in aqūlu innahu qayyūmun fa-dhālika
maẓharu aḥadin min quwwāmi bisāṭi aḥadiyyatihi fa-taʿālā taʿālā
dhikruhu.

``If I say that He is One, that is one of His names; and if I say that
He is All-Glorious, that name appears as one of His names; and if I say
that He is Most Holy, He is the One who reveals Himself through that
name to those like Him. If I say that He is Mighty, every mighty one
prostrates before the loftiness of His might; and if I say that He is
Beloved, every beloved one is consumed with love for the recognition of
Him and the majesty of His might. If I say that He is All-Protecting, He
is the One who, by His command, causes God to make whoever adduces proof
on the day of His appearance a protector over all things. If I say that
He is Self-Subsisting, that is the manifestation of one of those who
stand upon the expanse of His oneness. Exalted, exalted be His
remembrance!''

\subsection{710}\label{section-1205}

\begin{otherlanguage}{persian}

ثمّ تعالى تعالى شأنه كه معنى پارسى آن اينست كه اگر بگويم او واحد است اين
واحد يكى از اسماء او بوده اگر بگويم او سبّوح است محل ظهور اين اسم اسمى از
اسماء او بوده اگر بگويم او قدّوس است او است آنچنان نفسيكه تجلّى فرموده و
ميفرمايد باين اسم بامثال او و اگر بگويم او عزيز است هر عزيزى ساجد است از
براى علوّ عزّت او و اگر بگويم او محبوب است هر محبوبی‌ شقّ نموده‌ پرده قلب خود
را در دوستى و محبّت او و اگر بگويم او مهيمن است او است آن سلطانيكه هر
نفسى استدلال نمايد بدليل او در يوم ظهور او بتحقيق ميگرداند او را خدا
مهيمن بر كلّ شیء بامر او و اگر بگويم او قيّوم است پس آن مظهر يكى از قائمين
بساط احديّت او است مع آنكه باين تصريح جميع اسماء حسنى را از خدّام درگهش
ذكر نمودم و آن جوهر مقصود و طلعت موعود را از كلّ مقدّس و منزّه و مبرّا مع‌ذلك
اسمى از كتاب من اخذ نموده‌ايد و بآن اسم تشبّث نموده‌وارد آورده‌ايد آنچه را
كه هيچ ملّتى از قبل وارد نياورده

\end{otherlanguage}

Thumma taʿālā taʿālā shaʾnuhu, kih maʿnī-yi pārsī-yi ān īn ast kih agar
bigūyam ū vāḥid ast īn vāḥid yakī az asmāʾ-i ū būdih, agar bigūyam ū
subbūḥ ast maḥall-i ẓuhūr-i īn ism ismī az asmāʾ-i ū būdih, agar bigūyam
ū quddūs ast ū ast ānchunān nafsī kih tajallī farmūdih va mīfarmāyad
bi-īn ism bi-amthāl-i ū, va agar bigūyam ū ʿazīz ast har ʿazīzī sājid
ast az barā-yi ʿuluvv-i ʿizzat-i ū, va agar bigūyam ū maḥbūb ast har
maḥbūbī shaqq namūdih pardih-yi qalb-i khud rā dar dūstī va maḥabbat-i
ū, va agar bigūyam ū muhaymin ast ū ast ān sulṭānī kih har nafsī
istidlāl namāyad bi-dalīl-i ū dar yawm-i ẓuhūr-i ū bi-taḥqīq mīgardānad
ū rā Khudā muhaymin bar kull-i shayʾ bi-amr-i ū, va agar bigūyam ū
qayyūm ast pas ān maẓhar yakī az qāʾimīn-i bisāṭ-i aḥadiyyat-i ū ast.
Maʿa ānkih bi-īn taṣrīḥ jamīʿ-i asmāʾ-i ḥusnā rā az khuddām-i dargahash
dhikr namūdam va ān jawhar-i maqṣūd va ṭalʿat-i mawʿūd rā az kull
muqaddas va munazzah va mubarrā, maʿa dhālik ismī az kitāb-i man akhdh
namūdih-īd va bi-ān ism tashabbuth namūdih vārid āvardih-īd ānchih rā
kih hīch millatī az qabl vārid nayāvardih.

``And exalted, exalted be His station!'' Its meaning in Persian is this:
``If I say that He is One, this One is one of His names. If I say that
He is All-Glorious, the locus of the manifestation of this name is one
of His names. If I say that He is Most Holy, He is that very Person who
has revealed and continues to reveal Himself through this name to those
like Him. If I say that He is Mighty, every mighty one prostrates before
the loftiness of His might. If I say that He is Beloved, every beloved
one has torn the veil of his own heart in affection and love for Him. If
I say that He is All-Protecting, He is that Sovereign by whose command
God assuredly makes every soul who adduces His proof on the day of His
appearance a protector over all things. If I say that He is
Self-Subsisting, that manifestation is one of those who stand upon the
expanse of His oneness.'' ``Although I thus explicitly described all the
most beautiful names as servants at His threshold, and that desired
Essence and promised Countenance as sanctified, exalted, and free from
them all, you have nevertheless taken a name from My Book, clung to that
name, and inflicted what no religious community has ever inflicted
before.''

\subsection{711}\label{section-1206}

\begin{otherlanguage}{persian}

كذلك نطق لسان اللّه الملك المقتدر العلىّ العظيم اين بيانات نقطه اولى روح
ما سواه فداه آيا نفسى قادر است بر جواب لا و اللّه بل جميع منفعل و خجل و
شرمساريم چه كه نفوسيكه مقبل شده‌اند على ما هو عليه بشرايط اقبال قيام
ننموده‌اند و نفوس معرضين معلوم كه چه ارتكاب نموده‌ و مينمايند مع آنكه بنصّ
نقطه اولى جميع اسماء حسنى از مهيمن و قیّوم و فوق آن و دون آن جميع نفوسى
هستند كه در امر اين ظهور بديع ببيّنه خود اين ظهور استدلال نمايند ديگر
ملاحظه كن شأن نفوسيكه ببيان و قلم در اثبات اين ظهور اعظم استدلال نموده‌و
نوشته‌اند آنچه را كه بر آن قلوب مجرّده القا شده تاللّه انّهم من اعلى الخلق و
ابهاهم بحيث ما سبقهم احد من قبل و لن يسبقهم نفس من الآخرین

\end{otherlanguage}

Kadhālika naṭaqa lisānu Allāhi al-Maliki al-Muqtadiri al-ʿAliyyi
al-ʿAẓīm. Īn bayānāt-i Nuqṭih-yi Ūlā rūḥ-i mā sivāh fidāh āyā nafsī
qādir ast bar javāb? Lā wa-Allāh, bal jamīʿ munfaʿil va khajil va
sharmsārīm, chih kih nufūsī kih muqbil shudih-and ʿalā mā huwa ʿalayhi
bi-sharāyiṭ-i iqbāl qiyām nanamūdih-and va nufūs-i muʿriḍīn maʿlūm kih
chih irtikāb namūdih va mīnamāyand. Maʿa ānkih bi-naṣṣ-i Nuqṭih-yi Ūlā
jamīʿ-i asmāʾ-i ḥusnā az Muhaymin va Qayyūm va fawq-i ān va dūn-i ān
jamīʿ nufūsī hastand kih dar amr-i īn ẓuhūr-i badīʿ bi-bayyin ih-yi
khud-i īn ẓuhūr istidlāl namāyand. Dīgar mulāḥiẓih kun shaʾn-i nufūsī
kih bi-bayān va qalam dar ithbāt-i īn ẓuhūr-i aʿẓam istidlāl namūdih va
nivishtih-and ānchih rā kih bar ān qulūb-i mujarradih ilqā shudih.
Ta-Allāhi innahum min aʿlā al-khalqi wa-abhāhum bi-ḥaythu mā sabaqahum
aḥadun min qablu wa-lan yasbiqahum nafsun mina al-ākhirīn.

Thus has spoken the Tongue of God, the King, the All-Powerful, the Most
High, the Most Great. These are the words of the Primal Point---may the
soul of all else be sacrificed for Him. Is any soul capable of answering
them? No, by God! Rather, we are all abashed, ashamed, and embarrassed,
for those souls who have turned toward Him have not fulfilled the
conditions of turning toward Him as they truly require, while it is
evident what the souls who turn away have perpetrated and continue to
perpetrate. Yet, according to the explicit text of the Primal Point, all
the most beautiful names, from the All-Protecting and the
Self-Subsisting to those above and below them, are all souls who
establish the truth of this wondrous Revelation by the proof of this
Revelation itself. Consider, then, the station of those souls who,
through speech and writing, have adduced proofs establishing this Most
Great Revelation and have written what has been cast upon those detached
hearts. By God, they are among the most exalted and most glorious of
creation, such that no one before them has surpassed them, nor will any
soul among those to come surpass them.

\subsection{712}\label{section-1207}

\begin{otherlanguage}{persian}

اى نفس امّاره چرا بذكر آن نفوس مشغول شدى و اظهار تحيّر در جواب نمودى و
تنبّه ننمودى از فعل نفوسى كه در سبيل اين امر جان دادند و فعل هر يك در
اثبات اين ظهور ابدع امنع برهانى بود واضح و حجّتى لائح بين السّموات و
الأرضين از جمله محبوب الشّهداء جناب آقا سيّد اسمعيل زواره كه بدست خود حبّاً
لهذا الظّهور سر فدا نمود و هيچ نفسى در زهد و تقوى و ورع و فضل و علم او
باللّه منكر نبوده از فعل چنين نفسى مهتدى نشدى و ذكر نفوسى كه حال معلوم
نيست كه در چه صراط قائمند بين يدى العرش فرستاده و لعلّ اللّه يؤيّدهم على
امره و ينقطعهم عنك و عن امثالك و يجعلهم من المخلصين

\end{otherlanguage}

Ay nafs-i ammārih chirā bi-dhikr-i ān nufūs mashghūl shudī va iẓhār-i
taḥayyur dar javāb namūdī va tanabbuh nanamūdī az fiʿl-i nufūsī kih dar
sabīl-i īn amr jān dādand va fiʿl-i har yak dar ithbāt-i īn ẓuhūr-i
abdaʿ-i amnaʿ burhānī būd vāḍiḥ va ḥujjatī lāʾiḥ bayna al-samāwāti
wa-al-araḍīn. Az jumlih Maḥbūb al-Shuhadā Janāb-i Āqā Siyyid Ismāʿīl-i
Zavārih kih bi-dast-i khud ḥubban li-hādhā al-ẓuhūr sar fidā namūd va
hīch nafsī dar zuhd va taqvā va varaʿ va faḍl va ʿilm-i ū bi-Allāh
munkir nabūdih. Az fiʿl-i chunīn nafsī muhtadī nashudī va dhikr-i nufūsī
kih ḥāl maʿlūm nīst kih dar chih ṣirāṭ qāʾimand bayna yaday al-ʿarsh
firistādih, va laʿalla Allāha yuʾayyiduhum ʿalā amrihi wa-yanqaṭiʿuhum
ʿanka wa-ʿan amthālika wa-yajʿalahum mina al-mukhliṣīn.

O commanding self, why did you occupy yourself with mentioning those
souls and express bewilderment at answering, yet take no heed of the
deeds of souls who gave their lives in the path of this Cause? The deed
of each was a clear proof and a shining testimony, between the heavens
and the earths, establishing this most wondrous, most unassailable
Revelation. Among them was the Beloved of Martyrs, Āqā Siyyid Ismāʿīl of
Zavārih, who sacrificed his life by his own hand out of love for this
Revelation. No soul denied his asceticism, piety, devoutness, learning,
or knowledge of God. You were not guided by the deed of such a soul, but
sent before the Throne an account of souls whose present path is
unknown. May God strengthen them in His Cause, detach them from you and
those like you, and make them among the sincere.

\subsection{713}\label{section-1208}

\begin{otherlanguage}{persian}

بارى لازم شد كه تفصيل جناب سيّد شهيد عليه بهآء اللّه و عليه عزّ اللّه و عليه
نور اللّه و عليه كبريآء اللّه و عليه سرّ اللّه و عليه عظمة اللّه و عليه ضيآء
اللّه و عليه ثنآء اللّه و ثنآء كلّ شیء و ثنآء العالمين در اينمقام ذكر شود
كه شايد مطّلع شوى بر آن نفوسى كه نقطه بيان جلّ اقتداره خبر فرموده كه قبل
از ظهور آن شمس احديّه بنار حبّش مشتعل شوند چنانچه مخاطباً لظهور بعد
ميفرمايند كه مضمون پارسى آن اين است قوله عزّ اعزازه اگر چه عباد كلّ از تو
محجوب بوده و خواهند بود و لكن از براى تو عبادى هم خواهد خلق شد كه قبل از
ظهور تو بنار محبّت تو مشتعل گردند چنانچه سيّد مذكور قبل از ظهور مشتعل شده
و البتّه امثال آن نفس هم ظاهر خواهند شد و لو انتم و امثالكم تنكرون او
تكفرون حكايت آن شهيد آنكه وارد عراق شدند و جوار بيت اعظم بيت آقا محمّد
رضا عليه بهآء اللّه منزل و مقرّ گرفتند اسم مذكور از حضرت ابهى مستدعى شده
كه وقتى از اوقات ببيت او تشريف ببرند چند يوم تأخير افتاد و در جواب توقّف
فرمودند تا آنكه يومى از ایّام قبول فرموده و بآن مقرّ تشريف بردند و جناب
مذكور بقدر استطاعت خود مجلسى آراست و چند سينى از مركّبات و ميوه و حلويات
بحضور آورد اذاً توجّه وجه اللّه الى وجه الّذى استشهد فى سبيله و خاطبه تعال و
كلّ ما حضر بين يديك من الآء اللّه و نعمائه

\end{otherlanguage}

Bārī lāzim shud kih tafṣīl-i Janāb-i Siyyid-i shahīd, ʿalayhi bahāʾu
Allāh wa-ʿalayhi ʿizzu Allāh wa-ʿalayhi nūru Allāh wa-ʿalayhi kibriyāʾu
Allāh wa-ʿalayhi sirru Allāh wa-ʿalayhi ʿaẓamatu Allāh wa-ʿalayhi ḍiyāʾu
Allāh wa-ʿalayhi thanāʾu Allāh wa-thanāʾu kulli shayʾin wa-thanāʾu
al-ʿālamīn, dar īn maqām dhikr shavad kih shāyad muṭṭaliʿ shavī bar ān
nufūsī kih Nuqṭih-yi Bayān jalla iqtidāruhu khabar farmūdih kih qabl az
ẓuhūr-i ān shams-i aḥadiyyih bi-nār-i ḥubbash mushtaʿil shavand.
Chunānchih mukhāṭiban li-ẓuhūr-i baʿd mīfarmāyand kih maḍmūn-i pārsī-yi
ān īn ast, qawluhu ʿazza iʿzāzuhu: agar chih ʿibād kull az tu maḥjūb
būdih va khvāhand būd va lākin az barā-yi tu ʿibādī ham khvāhad khalq
shud kih qabl az ẓuhūr-i tu bi-nār-i maḥabbat-i tu mushtaʿil gardand.
Chunānchih Siyyid-i madhkūr qabl az ẓuhūr mushtaʿil shudih va al-battih
amthāl-i ān nafs ham ẓāhir khvāhand shud, wa-law antum wa-amthālukum
tunkirūna aw takfurūn. Ḥikāyat-i ān shahīd ānkih vārid-i ʿIrāq shudand
va jivār-i Bayt-i Aʿẓam bayt-i Āqā Muḥammad Riḍā ʿalayhi bahāʾu Allāh
manzil va maqarr giriftand. Ism-i madhkūr az Ḥaḍrat-i Abhā mustadʿī
shudih kih vaqtī az awqāt bi-bayt-i ū tashrīf bibarand. Chand yawm
taʾkhīr uftād va dar javāb tavaqquf farmūdand tā ānkih yawmī az ayyām
qabūl farmūdih va bi-ān maqarr tashrīf burdand va Janāb-i madhkūr
bi-qadr-i istiṭāʿat-i khud majlisī ārāst va chand sīnī az murakkabāt va
mīvih va ḥalaviyyāt bi-ḥuḍūr āvard. Idhan tawajjaha wajhu Allāhi ilā
wajhi alladhī istushhida fī sabīlihi wa-khāṭabahu taʿāl wa-kul mā ḥaḍara
bayna yadayka min ālāʾi Allāhi wa-naʿmāʾihi.

It has therefore become necessary to recount here the story of the
martyred Siyyid---upon him be the glory of God, the might of God, the
light of God, the grandeur of God, the mystery of God, the greatness of
God, the radiance of God, the praise of God, the praise of all things,
and the praise of the worlds---so that perhaps you may learn of those
souls of whom the Point of the Bayān---glorious is His power---foretold
that, before the appearance of that Sun of Oneness, they would be set
ablaze with the fire of His love. Thus, addressing the subsequent
Manifestation, He says, in words whose substance in Persian is as
follows---His words, exalted be His honor: ``Although the servants have
all been, and will remain, veiled from You, yet servants will also be
created for You who, before Your appearance, will be set ablaze with the
fire of Your love.'' Thus the aforementioned Siyyid was set ablaze
before the Revelation, and others like that soul will assuredly appear,
even though you and those like you deny or disbelieve. The story of that
martyr is this: he arrived in Iraq and took up residence near the Most
Great House, in the house of Āqā Muḥammad Riḍā---upon him be the glory
of God. The aforementioned person requested that the Most Glorious One
visit his house at some time. Several days passed, and He deferred His
reply until one day He accepted and went to that dwelling. The
aforementioned Siyyid arranged a gathering according to his means and
brought into His presence several trays of citrus fruits, other fruit,
and sweets. Then the Face of God turned toward the face of him who was
martyred in His path and addressed him: ``Come and eat what has been set
before you of God's favors and bounties.''

\subsection{714}\label{section-1209}

\begin{otherlanguage}{persian}

انّه تكلّم بين يدى اللّه بخضوع و خشوع و صريخ و انابة و قال اريد من بديع
مواهبك بأن ترزقنى من نعمائك الرّوحانيّة اذا اشتعل وجه الرّوح ثمّ قال تعال يا
عبد و امر بجلوسه امام وجهه ثمّ تكلّم لسان اللّه بكلمات يترشّح منها رشحات
المعانى على كلّ ما كان و ما يكون و انّى لم اقدر ان اصفها او اذكرها و لم
ادر ما انفق عليه يد العناية من نعمائه المكنونة الرّوحانيّة بحيث استجذبت
منها نفسه و روحه و كينونته و ذاته و اخذته غلبات الشّوق على شأن غفل عن
نفسه و عن كلّ من فى السّموات و الأرضين فتوجّه بسرّه و جهره الى محبوب
العالمين الى ان انتهى المجلس و رجع الرّوح الى مقرّه

\end{otherlanguage}

Innahu takallama bayna yaday Allāhi bi-khuḍūʿin wa-khushūʿin wa-ṣarīkhin
wa-inābatin wa-qāla urīdu min badīʿi mawāhibika bi-an tarzuqanī min
naʿmāʾika al-rūḥāniyyah. Idhan ishtaʿala wajhu al-Rūḥi thumma qāla taʿāl
yā ʿabd wa-amara bi-julūsihi amāma wajhihi thumma takallama lisānu
Allāhi bi-kalimātin yatarashshaḥu minhā rashaḥātu al-maʿānī ʿalā kulli
mā kāna wa-mā yakūn wa-innī lam aqdir an aṣifahā aw adhkurahā wa-lam
adri mā anfaqa ʿalayhi yadu al-ʿināyati min naʿmāʾihi al-maknūnati
al-rūḥāniyyati bi-ḥaythu istajdh ab at minhā nafsuhu wa-rūḥuhu
wa-kaynūnatuhu wa-dhātuhu wa-akhadhat-hu ghalabātu al-shawqi ʿalā
shaʾnin ghafala ʿan nafsihi wa-ʿan kulli man fī al-samāwāti wa-al-araḍīn
fa-tawajjaha bi-sirrihi wa-jahrihi ilā Maḥbūbi al-ʿālamīn ilā an intahā
al-majlisu wa-rajaʿa al-Rūḥu ilā maqarrihi.

He spoke before God with humility, reverence, cries, and supplication,
saying, ``I desire, from Your wondrous gifts, that You grant me a share
of Your spiritual bounties.'' Then the face of the Spirit blazed, and He
said, ``Come, O servant,'' and bade him sit before His face. Then the
Tongue of God spoke words from which drops of meaning rained upon all
that has been and will be. I could neither describe nor recount them,
nor do I know what hidden spiritual bounties the hand of providence
bestowed upon him, such that his soul, his spirit, his being, and his
essence were enraptured by them. The overpowering force of longing
seized him so completely that he became oblivious to himself and to all
who are in the heavens and the earths. Inwardly and outwardly he turned
toward the Beloved of the worlds, until the gathering ended and the
Spirit returned to His dwelling.

\subsection{715}\label{section-1210}

\begin{otherlanguage}{persian}

و لكن انّه بعد استماع كلمات اللّه و ما ذاق عمّا اراد ما شهد احد فى نفسه
سكوناً و قراراً و قضت عليه ایّام معدودات و فى كلّ حين يزداد شوقه و يشتدّ شغفه
باللّه بارئه الى ان حضر فى فجر يوم من الأیّام و كنس بعمامته فنآء البيت و
رجع و اخذ سكّيناً و تجنّب عن العباد و خرج عن المدينة الى ورد شاطى الشّط قام
مقبلاً الى البيت بيد اخذ لحاه و بيد آخر قطع حنجره حبّاً للّه المقتدر المهيمن
القيّوم اذا ارتفع بين النّاس ضوضآء و ارتفعت الصّياح من كلّ النّفوس و اجتمعوا
فى حوله خلق كثير و راوا بانّ السّكين كان بيده و وضع يده على صدره فتحيّر
بذلك كلّ الوجود ثمّ اهل ملأ الأعلی ثمّ اهل مداين الكبريا ثمّ اهل ملكوت
الأسماء و جبروت البقآء و كلّهم صلّوا عليه و كبّروا على وجهه و نزلوا عليه و
طافوا فى حوله و استنشقوا روايح حبّه و انّى لو اذكر ما ظهر فى ذلك اليوم انّك
لن تقبل و لن تستطيع ان تعرف و كان اللّه على ما اقول شهيد

\end{otherlanguage}

Wa-lākin innahu baʿda istimāʿi kalimāti Allāhi wa-mā dhāqa ʿammā arāda
mā shahida aḥadun fī nafsihi sukūnan wa-qarāran wa-qaḍat ʿalayhi ayyāmun
maʿdūdātun wa-fī kulli ḥīnin yazdādu shawquhu wa-yashtaddu shaghafuhu
bi-Allāhi bāriʾihi ilā an ḥaḍara fī fajri yawmin mina al-ayyāmi
wa-kanasa bi-ʿimāmatihi fināʾa al-bayti wa-rajaʿa wa-akhadha sikkīnan
wa-tajannaba ʿani al-ʿibādi wa-kharaja ʿani al-madīnati ilā wardi
shāṭiʾi al-shaṭṭi qāma muqbilan ilā al-bayti bi-yadin akhadha liḥāhu
wa-bi-yadin ākhar qaṭaʿa ḥanjarahu ḥubban li-Allāhi al-Muqtadiri
al-Muhaymini al-Qayyūm. Idhan irtafaʿa bayna al-nāsi ḍawḍāʾun
wa-irtafaʿati al-ṣiyāḥu min kulli al-nufūsi wa-ijtamaʿū fī ḥawlihi
khalqun kathīrun wa-raʾaw bi-anna al-sikkīna kāna bi-yadihi wa-waḍaʿa
yadahu ʿalā ṣadrihi fa-taḥayyara bi-dhālika kullu al-wujūdi thumma ahlu
malaʾi al-aʿlā thumma ahlu madāyini al-kibriyā thumma ahlu malakūti
al-asmāʾi wa-jabarūti al-baqāʾi wa-kulluhum ṣallaw ʿalayhi wa-kabbarū
ʿalā wajhihi wa-nazalū ʿalayhi wa-ṭāfū fī ḥawlihi wa-istanshaqū rawāʾiḥa
ḥubbihi wa-innī law adhkuru mā ẓahara fī dhālika al-yawmi innaka lan
taqbala wa-lan tastaṭīʿa an taʿrifa wa-kāna Allāhu ʿalā mā aqūlu shahīd.

But after he had heard the words of God and tasted what he had desired,
no one saw any calm or repose in him. A few days passed, and at every
moment his longing increased and his ardent love for God, his Maker,
grew more intense, until one day he came at dawn and swept the courtyard
of the House with his turban. He returned, took a knife, withdrew from
the people, and left the city for a watering place on the riverbank. He
stood facing the House, grasped his beard with one hand, and with the
other cut his throat out of love for God, the All-Powerful, the
All-Protecting, the Self-Subsisting. Then an uproar arose among the
people, cries went up from every soul, and a great crowd gathered around
him. They saw that the knife was in his hand and that he had placed his
hand upon his breast. All existence was astonished at this, as were the
people of the Concourse on High, the inhabitants of the cities of
grandeur, and the inhabitants of the kingdom of names and the dominion
of eternity. They all offered prayers for him, proclaimed God's
greatness before his face, descended upon him, circled around him, and
inhaled the fragrances of his love. Were I to recount what appeared on
that day, you would neither accept it nor be able to comprehend it. God
is witness to what I say.

\subsection{716}\label{section-1211}

\begin{otherlanguage}{persian}

بلسان پارسى ذكر ميشود كه سيّد مذكور بعد از ورود حضرت ابهى غذاى روحانى طلب
نموده ‌فرمودند بيا و در مقابل بنشين بعد از جلوس بياناتى فرمودند فو‌اللّه
سكر خمر معانى كلمات الهيّه چنان جذبش نمود كه از خود و كونين غافل گشت و
اگر عالمين اراده نمايند كه تفصيل آنمجلس را ذكر كنند و يا بوصف آن كوثر
معانى كه از بحر بيان طلعت لايزالى جارى شد مشغول شوند البتّه خود را عاجز
مشاهده نمايند و بشأنى كلمات اللّه در قلب مقبل الى اللّه مؤثّر افتاد كه جميع
اركانش بنار محبّت الهى مشتعل و ديگر حقّ عالم است كه باو چه عنايت شد قد ذاق
ما لاعرفه احد الّا اللّه العليم الخبير

\end{otherlanguage}

Bi-lisān-i pārsī dhikr mīshavad kih Siyyid-i madhkūr baʿd az vurūd-i
Ḥaḍrat-i Abhā ghadhā-yi rūḥānī ṭalab namūdih, farmūdand biyā va dar
muqābil binishīn. Baʿd az julūs bayānātī farmūdand, fa-wa-Allāh sukr-i
khamr-i maʿānī-yi kalimāt-i ilāhiyyih chunān jadhbash namūd kih az khud
va kawnayn ghāfil gasht va agar ʿālamayn irādih namāyand kih tafṣīl-i ān
majlis rā dhikr kunand va yā bi-vaṣf-i ān kawthar-i maʿānī kih az baḥr-i
bayān-i ṭalʿat-i lā-yazālī jārī shud mashghūl shavand al-battih khud rā
ʿājiz mushāhidih namāyand. Va bi-shaʾnī kalimāt-i Allāh dar qalb-i
muqbil ilā Allāh muʾaththir uftād kih jamīʿ-i arkānash bi-nār-i
maḥabbat-i ilāhī mushtaʿil va dīgar Ḥaqq ʿālim ast kih bi-ū chih ʿināyat
shud. Qad dhāqa mā lā ʿarafahu aḥadun illā Allāhu al-ʿAlīmu al-Khabīr.

It is recounted in Persian that, after the arrival of the Most Glorious
One, the aforementioned Siyyid asked for spiritual food. He said, ``Come
and sit before Me.'' After he had seated himself, He spoke certain
words. By God, the intoxication of the wine of meanings in the divine
words so enraptured him that he became oblivious to himself and to both
worlds. If both worlds sought to recount the details of that gathering
or to describe that Kawthar of meanings which flowed from the ocean of
the utterance of the eternal Countenance, they would assuredly find
themselves powerless. The words of God so affected the heart of him who
had turned toward God that his entire frame was set ablaze with the fire
of divine love. God alone knows what was bestowed upon him. He tasted
what no one has known except God, the All-Knowing, the All-Aware.

\subsection{717}\label{section-1212}

\begin{otherlanguage}{persian}

و بعد مجلس منقضى شده جمال ابهى تشريف بردند و لكن آن سيّد در كلّ حين بشأن
بديع ظاهر و بشأنى حبّ اللّه اخذش نمود كه بالأخره‌ از اكل و شرب منقطع شد و
مدّتى باينحالت بوده تا آنكه در يومى اشعارى در مدح اللّه گفته و باصحاب داده
و الان موجود اگر نفسى ملاحظه نمايد از اشتعال كينونت آن ساذج قدس مطّلع
ميشود تا آنكه يومى از ایّام در اوّل فجر برخواسته و به بيت اعظم توجّه نمود و
بعمامه خود فناء اطهر بيت را جاروب نموده‌و بعد رفته تيغى اخذ نمود و ببيت
آقا محمّد رضا هم رفته و ديدنى نمود و بمقرّ قربانگاه دوست شتافت منقطعاً عن
كلّ من فى السّموات و الأرضين و در آن حين توجّه جميع ملائكه عالين در حولش
بوده و با او بقربانگاه توجّه نمودند تا آنكه در خارج مدينه قرب شط مقابل
بيت بدست خود حنجر مبارك خود را قطع نمود شوقاً لحبّ ربّه و شغفاً فى ودّه و
طلباً لوصاله و راضياً لقضائه و مقبلاً الى حرم وصله

\end{otherlanguage}

Va baʿd majlis munqaḍī shude Jamāl-i Abhā tashrīf burdand va lākin ān
siyyid dar kull-i ḥīn bi-shaʾn-i badīʿ ẓāhir va bi-shaʾnī ḥubb-i Allāh
akhdhash namūd kih bi'l-ākhirih az akl va shurb munqaṭiʿ shud va muddatī
bi-īn ḥālat būde tā ānkih dar yawmī ashʿārī dar madḥ-i Allāh gufte va
bi-aṣḥāb dāde va al-ān mawjūd agar nafsī mulāḥiẓih namāyad az ishtiʿāl-i
kaynūnat-i ān sādhij-i quds muṭṭaliʿ mīshavad tā ānkih yawmī az ayyām
dar avval-i fajr bar-khāste va bi-Bayt-i Aʿẓam tavajjuh namūd va
bi-ʿammāmiy-i khud fanā-yi aṭhar-i bayt rā jārūb namūde va baʿd rafte
tīghī akhdh namūd va bi-bayt-i Āqā Muḥammad Riḍā ham rafte va dīdanī
namūd va bi-maqarr-i qurbāngāh-i Dūst shitāft munqaṭiʿan ʿan kulli man
fī al-samāwāti wa'l-araḍīn va dar ān ḥīn tavajjuh-i jamīʿ-i malāʾikiy-i
ʿāllīn dar ḥawlash būde va bā ū bi-qurbāngāh tavajjuh namūdand tā ānkih
dar khārij-i madīnih qurb-i shaṭṭ muqābil-i bayt bi-dast-i khud ḥanjar-i
mubārak-i khud rā qaṭʿ namūd shawqan li-ḥubbi rabbihi wa-shaghafan fī
wuddihi wa-ṭalaban li-wiṣālihi wa-rāḍiyan li-qaḍāʾihi wa-muqbilan ilā
ḥarami waṣlihi.

Afterward, the gathering ended and the Most Glorious Beauty departed.
But that Siyyid manifested a wondrous state at every moment, and the
love of God so seized him that he eventually ceased eating and drinking.
He remained in this condition for a time, until one day he composed
verses in praise of God and gave them to the companions. They still
exist, and anyone who reads them will become aware of the fire that
consumed the very being of that embodiment of holiness. Then, one day,
he arose at the first light of dawn and turned toward the Most Great
House. With his turban he swept the most pure courtyard of the House,
and afterward went and took a blade. He also went to the house of Āqā
Muḥammad Riḍā and paid a visit, then hastened to the place of sacrifice
for the Beloved, detached from all who are in the heavens and the
earths. At that moment, all the exalted angels were gathered around him,
their attention fixed upon him, and they went with him toward the place
of sacrifice. Outside the city, near the river, facing the House, he cut
his blessed throat with his own hand, out of longing for the love of his
Lord, passionate devotion to Him, and a desire for reunion with Him,
content with His decree and turning toward the sanctuary of reunion with
Him.

\subsection{718}\label{section-1213}

\begin{otherlanguage}{persian}

و بعد در مدينه اين امر شهرت نمود بشأنيكه كلّ استماع نمودند و تفنگچى باشى
با جمعى از عرب و عجم توجّه بآن مقرّ اطهر نموده ‌مشاهده نمودند كه آن هيكل
مقدّس خوابيده و تيغ در دستش و باينحالت جان فدا نموده ‌بعد از ملاحظه كلّ
متحيّر شده و همان فعل سبب شد كه چند نفر از اعدا از افعال خود نادم شده بحبّ
الهى فايز شدند و در هيچ عصرى چنين امرى واقع نشد حال ملاحظه كن كه اسم اين
نفس مبارك هم اسمعيل بوده و لكن اين اسمعيل از قربانگاه دوست زنده برنگشت و
جان در رهش ايثار نمود فو‌اللّه اين فعل حجّت است بر كلّ من فى السّموات و الأرض
و لكن تو نيابی‌ و امثال تو نيابند تو را چه باين عوالم تو برو و ذكر وصايت
مجعوله كن و در آنمقام تحقيق كن و رساله بنويس

\end{otherlanguage}

Va baʿd dar madīnih īn amr shuhrat namūd bi-shaʾnī kih kull istimāʿ
namūdand va tufangchī-bāshī bā jamʿī az ʿarab va ʿajam tavajjuh bi-ān
maqarr-i aṭhar namūde mushāhidih namūdand kih ān haykal-i muqaddas
khābīde va tīgh dar dastash va bi-īn ḥālat jān fidā namūde baʿd az
mulāḥiẓih kull mutaḥayyir shude va hamān fiʿl sabab shud kih chand nafar
az aʿdā az afʿāl-i khud nādim shude bi-ḥubb-i ilāhī fāʾiz shudand va dar
hīch ʿaṣrī chunīn amrī vāqiʿ nashud ḥāl mulāḥiẓih kun kih ism-i īn
nafs-i mubārak ham Ismāʿīl būde va lākin īn Ismāʿīl az qurbāngāh-i Dūst
zindih bar-nagasht va jān dar rahash īthār namūd fa-wa'llāh īn fiʿl
ḥujjat ast bar kulli man fī al-samāwāti wa'l-arḍ va lākin tu nayābī va
amthāl-i tu nayāband tu rā chih bi-īn ʿavālim tu biraw va dhikr-i
viṣāyat-i majʿūlih kun va dar ān maqām taḥqīq kun va risālih binivīs.

Afterward, news of this event spread through the city until everyone had
heard it. The chief of the musketeers, with a group of Arabs and
Persians, went to that most pure spot and saw that sacred figure lying
there, the blade in his hand, having sacrificed his life in this manner.
All who saw him were astonished, and this very deed caused several of
his enemies to repent of their actions and attain the love of God. In no
age had such a thing occurred. Now consider that the name of this
blessed soul was also Ishmael, but this Ishmael did not return alive
from the place of sacrifice for the Beloved: he gave his life in His
path. By God, this deed is a proof to all who are in the heavens and on
earth, but you do not comprehend it, nor do those like you. What have
you to do with these realms? Go and speak of a fabricated successorship;
investigate that matter and write a treatise.

\subsection{719}\label{section-1214}

\begin{otherlanguage}{persian}

و هم‌چنين نفوس ديگر كه در اين ظهور جان داده‌اند و تا مقرّ قربانگاه بذكر
محبوب ناطق بوده‌اند ذكر اين نفوس را كه مثل شموس مشرقه ما‌بين ناس مشرق و
مضئ‌اند بغضاً على اللّه ننموده‌ و مينويسى كه جواب فلان و فلان را چه بگويم قد
خسر الّذينهم كفروا بآيات اللّه و اتّبعوك و امثالك فى الدّنيا و الآخرة و انّ
النّار مثويهم فبئس مثوى المنكرين

\end{otherlanguage}

Va hamchunīn nufūs-i dīgar kih dar īn ẓuhūr jān dāde-and va tā maqarr-i
qurbāngāh bi-dhikr-i Maḥbūb nāṭiq būde-and dhikr-i īn nufūs rā kih
mithl-i shumūs-i mushriqih mā-bayn-i nās mushriq va muḍīʾ-and bughḍan
ʿalā Allāh nanamūde va mīnivīsī kih javāb-i fulān va fulān rā chih
bigūyam qad khasira alladhīnahum kafarū bi-āyāti Allāh wa-ttabaʿūka
wa-amthālaka fī al-dunyā wa'l-ākhirah wa-inna al-nāra mathwāhum fa-biʾsa
mathwā al-munkirīn.

Likewise, there are other souls who have given their lives in this
Revelation and continued to speak in remembrance of the Beloved until
they reached the place of sacrifice. Out of hatred for God, you have
made no mention of these souls, who shine resplendent among the people
like rising suns, and you write, ``What answer should I give to this
person and that?'' Those who have disbelieved in the verses of God and
followed you and your like have indeed suffered loss in this world and
the next. The Fire is their dwelling place: how wretched is the dwelling
of the deniers!

\subsection{720}\label{section-1215}

\begin{otherlanguage}{persian}

و اينكه نوشته كه چه جواب گويم كه فلانيها هر روز بيكى ميچسبند و يكى را حقّ
ميدانند و باطل مينمايند اين حقّيكه ذكر نموده ‌لازال خلق بوده بقميصى از
قمايص اسماء مفتخر شد و بعد از طغيان و غرور و اعراض آن قميص از او اخذ شده
خذ زمام القلم لانّ بما جرى منه يلعنك اهل ممالك القدم و يشهد بذلك لسان
اللّه الصّادق النّاطق الأمين ميفرمايد طوريّون در كرّه اخرى عند رشح من ذلك
النّور المهيمن الحمرا ميّت و لاشیء ميشوند تو تازه ذكر مينمائى كه چگونه
ميشود مرآتى على زعم تو باطل شود و يا ثمره ساقط گردد منصوص است كه بسا
شجره اثبات كه در ظهور بعد باعراض از شجره نفى ميشود و تو حال بيك ثمره
بزعم خود چسبيده و متمسّك شده لو يكون ثمرة حقيقيّة باعراضها سقطت عن سدرة
الالهيّة و رجعت الى اوّل مقرّها فى النّار

\end{otherlanguage}

Va īnkih nivishte kih chih javāb gūyam kih fulānīhā har rūz bi-yikī
mīchasband va yakī rā ḥaqq mīdānand va bāṭil mīnamāyand īn ḥaqqī kih
dhikr namūde lāzāl khalq būde bi-qamīṣī az qamāyiṣ-i asmāʾ muftakhir
shud va baʿd az ṭughyān va ghurūr va iʿrāḍ ān qamīṣ az ū akhdh shude
khudh zimāma al-qalam li-anna bi-mā jarā minhu yalʿanuka ahlu mamāliki
al-qidam wa-yashhadu bi-dhālika lisānu Allāhi al-ṣādiqu al-nāṭiqu
al-amīn mīfarmāyad Ṭūriyyūn dar karrih-yi ukhrā ʿinda rashḥin min
dhālika al-nūri al-muhaymini al-ḥamrā mayyit va lā-shayʾ mīshavand tu
tāzih dhikr mīnamāʾī kih chigūnih mīshavad mirʾātī ʿalā zaʿm-i tu bāṭil
shavad va yā thamarah sāqiṭ gardad manṣūṣ ast kih basā shajariy-i ithbāt
kih dar ẓuhūr-i baʿd bi-iʿrāḍ az shajariy-i nafy mīshavad va tu ḥāl
bi-yik thamarah bi-zaʿm-i khud chasbīde va mutamassik shude law yakūnu
thamaratan ḥaqīqiyyatan bi-iʿrāḍihā saqaṭat ʿan sidrati al-ilāhiyyah
wa-rajaʿat ilā awwali maqarrihā fī al-nār.

As for what you have written, ``What answer should I give, when these
people attach themselves to someone new every day, regard someone as
true, and then declare him false?''---this one whom you have called true
has always been a created being. He was honored with one of the garments
of the names, and after his rebellion, pride, and turning away, that
garment was taken from him. Rein in your pen, for because of what has
flowed from it, the inhabitants of the realms of eternity curse you. To
this bears witness the truthful, speaking, trustworthy Tongue of God. He
says that, in another cycle, the people of Sinai become dead and as
nothing at a sprinkling from that all-protecting crimson Light. Yet you
now ask how a mirror, as you suppose him to be, can become false, or a
fruit can fall. It is expressly stated that many a tree of affirmation,
through turning away in the next Revelation, becomes a tree of negation.
Yet now you have attached yourself to and held fast to what you imagine
to be a fruit. Even if it were a true fruit, by turning away it would
fall from the divine Lote-Tree and return to its original abode in the
Fire.

\subsection{721}\label{section-1216}

\begin{otherlanguage}{persian}

در اينمقام يك بيان از نقطه بيان و طلعت رحمن و نفس سبحان روح من فى ممالك
الأكوان و الامكان فداه ذكر ميشود كه شايد حيا نمائى و دست تعدّى بسرادق علم
و عرفان الهى دراز ننمائى اگر چه تو و امثال تو عند اللّه از محرومين محسوب
و لكن وجه بيان بمقرّبين بوده و خواهد بود قال و قوله الحقّ و اجعل اللّهمّ تلك
الشّجرة كلّها لها لتظهرن ما فيها ثمرات ما قد خلق اللّه فيها لمن قد اراد
اللّه ان يظهر به ما اراد فانّنى انا و عزّتك ما اردت ان يكون على تلك الشّجرة
من غصن و لا ورق و لا ثمر لن يسجد له يوم ظهوره و لا يسبّحك به بما ينبغى
لعلوّ علوّ ظهوره و سموّ سموّ بطونه و ان شهدت يا الهى علىّ من غصن او ورق او
ثمر لم يسجد له يوم ظهوره فاقطعه اللّهمّ عن تلك الشّجرة فانّه لم يكن منّى و
لايرجع الىّ انتهى

\end{otherlanguage}

Dar īn maqām yik bayān az Nuqṭiy-i Bayān va Ṭalʿat-i Raḥmān va Nafs-i
Subḥān rūḥu man fī mamāliki al-akwāni wa'l-imkān fidāh dhikr mīshavad
kih shāyad ḥayā namāʾī va dast-i taʿaddī bi-surādiq-i ʿilm va ʿirfān-i
ilāhī dirāz nanamāʾī agar chih tu va amthāl-i tu ʿinda Allāh az maḥrūmīn
maḥsūb va lākin vajh-i bayān bi-muqarrabīn būde va khvāhad būd qāla
wa-qawluhu al-ḥaqq wa-jʿal allāhumma tilka al-shajarata kullahā lahā
li-taẓharanna mā fīhā thamarātu mā qad khalaqa Allāhu fīhā li-man qad
arāda Allāhu an yuẓhira bihi mā arād fa-innanī anā wa-ʿizzatika mā
aradtu an yakūna ʿalā tilka al-shajarati min ghuṣnin wa-lā waraqin wa-lā
thamarin lan yasjuda lahu yawma ẓuhūrihi wa-lā yusabbiḥaka bihi bi-mā
yanbaghī li-ʿuluwwi ʿuluwwi ẓuhūrihi wa-sumūwwi sumūwwi buṭūnihi wa-in
shahidta yā ilāhī ʿalayya min ghuṣnin aw waraqin aw thamarin lam yasjud
lahu yawma ẓuhūrihi fa-qṭaʿhu allāhumma ʿan tilka al-shajarah fa-innahu
lam yakun minnī wa-lā yarjiʿu ilayya intahā.

Here an utterance of the Point of the Bayān, the Countenance of the
All-Merciful and the Person of the All-Glorious---may the souls of all
who dwell in the realms of existence and contingency be sacrificed for
Him---is cited, so that perhaps you may feel shame and not stretch out a
hand of aggression toward the pavilion of divine knowledge and
recognition. Although you and your like are numbered among the deprived
in God's sight, the utterance has been, and will remain, addressed to
those who are near to Him. He said, and His word is truth: ``Make, O
God, that Tree wholly hers, so that the fruits of what God has created
within it may appear within it for Him through whom God has willed to
manifest what He has willed. For I, by Your might, have not desired that
there should be upon that Tree any branch, leaf, or fruit that will not
bow down before Him on the day of His appearance, or glorify You through
Him as befits the loftiest heights of His manifestation and the most
exalted sublimity of His concealment. And if, O my God, You should
behold upon Me any branch, leaf, or fruit that does not bow down before
Him on the day of His appearance, cut it off, O God, from that Tree, for
it is not of Me and does not return to Me.'' End of quotation.

\subsection{722}\label{section-1217}

\begin{otherlanguage}{persian}

نفسى فداه و روحى فداه و سرّى فداه و ظاهرى فداه و باطنى فداه و الملك و
الملكوت فداه و القدرة و اللّاهوت فداه و العزّ و الياقوت فداه و الظّهور و
النّاسوت فداه و ما خلق فى ممالك القدم فداه لانّه روحى فداه ما تحرّك قلمه
الّا فى هذا الذّكر البديع و هذا الظّهور المتباذخ المنيع و هذا العلوّ
المتعالى الرّفيع و هذه الشّمس المشرقة المنير و هذا الأمر الظّاهر الباهر
العزيز المنيع قدرى تفكّر كن كه چه ميفرمايد آن نقطه وجود و جمال معبود
ميفرمايد اى پروردگار من اراده نکرده‌ام اينكه باقى بماند بر اين شجر از غصن
و ورق و ثمرى كه ساجد نشود از براى او يوم ظهور او و تسبيح ننمايد تو را
بآن مظهر ظهور بآنچه سزاوار است از براى علوّ علوّ ظهور او و سموّ سموّ بطون او
و ميفرمايد اگر مشاهده نمودى تو اى پروردگار من بر من از غصن و ورق و ثمرى
كه ساجد نشود از براى او يوم ظهور او پس قطع كن او را اى پروردگار از اين
شجره بدرستيكه او نبوده است از من و راجع نميشود بسوى من

\end{otherlanguage}

Nafsī fidāh wa-rūḥī fidāh wa-sirrī fidāh wa-ẓāhirī fidāh wa-bāṭinī fidāh
wa'l-mulku wa'l-malakūtu fidāh wa'l-qudratu wa'l-lāhūtu fidāh wa'l-ʿizzu
wa'l-yāqūtu fidāh wa'l-ẓuhūru wa'l-nāsūtu fidāh wa-mā khuliqa fī
mamāliki al-qidami fidāh li-annahu rūḥī fidāh mā taḥarraka qalamuhu illā
fī hādhā al-dhikri al-badīʿ wa-hādhā al-ẓuhūri al-mutabādhikhi al-manīʿ
wa-hādhā al-ʿuluwwi al-mutaʿālī al-rafīʿ wa-hādhihi al-shamsi
al-mushriqati al-munīr wa-hādhā al-amri al-ẓāhiri al-bāhiri al-ʿazīzi
al-manīʿ qadrī tafakkur kun kih chih mīfarmāyad ān Nuqṭiy-i Vujūd va
Jamāl-i Maʿbūd mīfarmāyad ay Parvardigār-i man irādih nakarde-am īnkih
bāqī bimānad bar īn shajar az ghuṣn va varaq va thamarī kih sājid
nashavad az barā-yi ū yawm-i ẓuhūr-i ū va tasbīḥ nanamāyad tu rā bi-ān
maẓhar-i ẓuhūr bi-ānchih sazāvār ast az barā-yi ʿuluvv-i ʿuluvv-i
ẓuhūr-i ū va sumuvv-i sumuvv-i buṭūn-i ū va mīfarmāyad agar mushāhidih
namūdī tu ay Parvardigār-i man bar man az ghuṣn va varaq va thamarī kih
sājid nashavad az barā-yi ū yawm-i ẓuhūr-i ū pas qaṭʿ kun ū rā ay
Parvardigār az īn shajarih bi-durustī kih ū nabūde ast az man va rājiʿ
namīshavad bi-sū-yi man.

May my self be sacrificed for Him, and my soul be sacrificed for Him,
and my inmost secret be sacrificed for Him, and my outward being be
sacrificed for Him, and my inward being be sacrificed for Him; may the
earthly kingdom and the heavenly kingdom be sacrificed for Him, and
power and the divine realm be sacrificed for Him, and glory and ruby be
sacrificed for Him, and manifestation and the human realm be sacrificed
for Him, and all that has been created in the realms of eternity be
sacrificed for Him. For He---may my soul be sacrificed for Him---moved
His pen only in celebration of this wondrous Remembrance, this supremely
exalted and unassailable Revelation, this transcendent and lofty height,
this radiant, shining Sun, and this manifest, resplendent, mighty,
inviolable Cause. Reflect a little on what that Point of Existence and
Beauty of the Adored One says. He says, ``O my Lord, I have not desired
that there should remain upon this Tree any branch, leaf, or fruit that
does not bow down before Him on the day of His appearance, and does not
glorify You through that Manifestation of revelation as befits the
loftiest heights of His manifestation and the most exalted sublimity of
His concealment.'' And He says, ``If, O my Lord, You should behold upon
Me any branch, leaf, or fruit that does not bow down before Him on the
day of His appearance, then cut it off, O Lord, from this Tree. Truly,
it has not been of Me and does not return to Me.''

\subsection{723}\label{section-1218}

\begin{otherlanguage}{persian}

قدرى تفكّر نما ثمره شجره حقيقت در ظهور قبل آن است كه كلّ اهل بيان در ظلّ
اين شجره ساكن شوند اين بوده اصل مقصود از آنچه تقدير شد در شجره قبل حال
هر نفسى بآنمقام فايز شد از اثمار و افنان و اغصان و ما قدّر لها مذكور و من
دون آن ساقط و مقطوع حال قدرى نوحه نما بر ثمره كه من دون اللّه اخذ نموده ‌و
اگر ثمرهء حقيقى باشد بنصّ نقطه بيان كه ميفرمايد اگر حقّ را در يوم ظهور
تصديق ننمايد از من نبوده و نخواهد بود مع اين بيان واضح لايح آشكار تعجّب
مينمائيد كه چگونه ميشود ثمره باطل گردد و يا ساقط شود

\end{otherlanguage}

Qadrī tafakkur namā thamarah-yi shajariy-i ḥaqīqat dar ẓuhūr-i qabl ān
ast kih kull-i ahl-i Bayān dar ẓill-i īn shajarih sākin shavand īn būde
aṣl-i maqṣūd az ānchih taqdīr shud dar shajariy-i qabl ḥāl har nafsī
bi-ān maqām fāʾiz shud az athmār va afnān va aghṣān va mā quddira lahā
madhkūr va min dūn-i ān sāqiṭ va maqṭūʿ ḥāl qadrī nawḥih namā bar
thamarah kih min dūni Allāh akhdh namūde va agar thamarah-yi ḥaqīqī
bāshad bi-naṣṣ-i Nuqṭiy-i Bayān kih mīfarmāyad agar Ḥaqq rā dar yawm-i
ẓuhūr taṣdīq nanamāyad az man nabūde va nakhvāhad būd maʿ-i īn bayān-i
vāḍiḥ-i lāyiḥ-i āshkār taʿajjub mīnamāyīd kih chigūnih mīshavad thamarah
bāṭil gardad va yā sāqiṭ shavad.

Reflect a little. The fruit of the Tree of Truth in the former
Revelation is that all the people of the Bayān should dwell beneath the
shade of this Tree. This was the essential purpose of all that was
ordained in the former Tree. Now, whoever has attained that station is
numbered among its fruits, twigs, branches, and whatever has been
ordained for it; anyone else is fallen and cut off. Now lament a little
for the fruit that has taken hold of something other than God. Even if
it is a true fruit, the Point of the Bayān expressly says, ``If it does
not acknowledge the True One on the day of His appearance, it has not
been, and will not be, of Me.'' Despite this clear, evident, manifest
utterance, you wonder how a fruit can become false or fall!

\subsection{724}\label{section-1219}

\begin{otherlanguage}{persian}

حال عظمت امر اللّه را مشاهده كن و شأن خود و مرشدت را بگو اى طفل دبستان
جهل هنوز بالغ نشده چه كه اگر بالغ بودى باين كلمات مشغول نميشدى فافّ لك و
بما عندك و بما استدللت به و بما تكون مسروراً به اليوم بالغين آنانند كه
بمظهر ربّ العالمين عارف شده و بر اعراش تمكين و وقار متمكّن نزد تو و امثال
تو بسيار مشكل است خرق اين حجبات و لكن از براى حقّ جلّ ثنائه نفوسى هست كه
بنفحاتش جميع احجاب را خرق نمايند و ايشانند آن نفوسى كه كلّ نفس ايشان
مهيمن است بر كلّ اهل بيان و هم‌چنين كلمات ايشان و افعال ايشان و عرفان
ايشان چنانچه نقطه بيان ميفرمايد كه نطفه يكساله ظهور بعد اقوى است از كلّ
من فى البيان و اللّه الّذى لا اله الّا هو كه اگر در همين بيان نقطه بيان
تفكّر نمايند جميع از كلّ من فى السّموات و الأرض و از روح و نفس و ذات خود
منقطع شده بشطر احديّه توجّه نمايند

\end{otherlanguage}

Ḥāl ʿaẓamat-i Amr-i Allāh rā mushāhidih kun va shaʾn-i khud va murshidat
rā bigū ay ṭifl-i dabistān-i jahl hanūz bāligh nashude chih kih agar
bāligh būdī bi-īn kalimāt mashghūl namīshudī fa-uffin laka wa-bi-mā
ʿindaka wa-bi-mā istad\-lalta bihi wa-bi-mā takūnu masrūran bihi al-yawm
bālighīn ānānand kih bi-Maẓhar-i Rabb-i ʿālamīn ʿārif shude va bar
aʿrāsh-i tamkīn va vaqār mutamakkin nazd-i tu va amthāl-i tu bisyār
mushkil ast kharq-i īn ḥujubāt va lākin az barā-yi Ḥaqq jalla thanāʾuhu
nufūsī hast kih bi-nafaḥātash jamīʿ-i aḥjāb rā kharq namāyand va
īshānand ān nufūsī kih kull-i nafs-i īshān muhaymin ast bar kull-i ahl-i
Bayān va hamchunīn kalimāt-i īshān va afʿāl-i īshān va ʿirfān-i īshān
chinānchih Nuqṭiy-i Bayān mīfarmāyad kih nuṭfih-yi yiksālih-yi ẓuhūr-i
baʿd aqvā ast az kulli man fī al-Bayān wa'llāhi alladhī lā ilāha illā
huwa kih agar dar hamīn bayān-i Nuqṭiy-i Bayān tafakkur namāyand jamīʿ
az kulli man fī al-samāwāti wa'l-arḍ va az rūḥ va nafs va dhāt-i khud
munqaṭiʿ shude bi-shaṭr-i Aḥadiyyih tavajjuh namāyand.

Now behold the greatness of the Cause of God and the standing of
yourself and your guide. Say: O child in the school of ignorance, you
have not yet reached maturity, for had you matured, you would not have
occupied yourself with these words. Shame on you, on what you possess,
on what you have adduced as proof, and on what delights you today! The
mature are those who have recognized the Manifestation of the Lord of
the worlds and are established upon the thrones of steadfastness and
dignity. For you and your like, tearing through these veils is
exceedingly difficult. Yet God---exalted be His praise---has souls who,
through His breaths, tear through every veil. These are the souls each
one of whom holds sway over all the people of the Bayān, as do their
words, their deeds, and their recognition. Thus the Point of the Bayān
says that the year-old embryo of the next Revelation is mightier than
all who are in the Bayān. By God, besides whom there is no god, if they
reflected upon this very utterance of the Point of the Bayān, they would
all detach themselves from everyone in the heavens and on earth, and
from their own souls, selves, and beings, and turn toward the One.

\subsection{725}\label{section-1220}

\begin{otherlanguage}{persian}

و اين كلمه بسيار بزرگ است لو انتم تتفكّرون چه اگر امر عظيم نباشد چگونه
چنين بيان از منزل بيان ظاهر ميشد اذا نشهد بأن صدق اللّه العلىّ العظيم چه
كه اگر اقوى از كلّ خلق نباشد نميتواند اين سبحات اوهام را خرق نمايد و اين
اصنام كبيره را بقوّه خليل آفاق‌ بشكند كجا است مقام اين عباد و مقام عباد
قبل هؤلآء عباد يصلين عليهم اهل ملأ الأعلی ثمّ حقايق الأشیآء و لكن انتم لا
تشعرون

\end{otherlanguage}

Va īn kalimih bisyār buzurg ast law antum tatafakkarūn chih agar amr
ʿaẓīm nabāshad chigūnih chunīn bayān az Munzil-i Bayān ẓāhir mīshud
idhan nashhadu bi-anna ṣadaqa Allāhu al-ʿaliyyu al-ʿaẓīm chih kih agar
aqvā az kull-i khalq nabāshad namītavānad īn subuḥāt-i awhām rā kharq
namāyad va īn aṣnām-i kabīrih rā bi-quvvih-yi Khalīl-i āfāq bishikanad
kujā ast maqām-i īn ʿibād va maqām-i ʿibād-i qabl hāʾulāʾi ʿibādun
yuṣallīna ʿalayhim ahlu malaʾi al-aʿlā thumma ḥaqāʾiqu al-ashyāʾ
wa-lākin antum lā tashʿurūn.

This utterance is exceedingly momentous, if you would reflect. For if
the Cause were not mighty, how could such an utterance have appeared
from the Revealer of the Bayān? Thus We testify that God, the Most High,
the Most Great, has spoken the truth. For unless one is mightier than
all creation, one cannot tear through these veils of vain imaginings or
shatter these mighty idols through the power of the Friend of all the
worlds. How far apart are the station of these servants and that of the
servants of former times! These are servants upon whom the inhabitants
of the Concourse on High, and then the realities of all things, invoke
blessings, but you do not perceive.

\subsection{726}\label{section-1221}

\begin{otherlanguage}{persian}

و اينكه بقلم غلّ و افترا نوشته كه بعضى مردم نسبت عصيان بجمال ابهى ميدهند
و نوشته آنچه قلم حيا مينمايد از ذكر آن كلّما ذكرت كذب صراح جميع اين كلمات
را آن ملقى كذّاب بتو القا نموده‌ و نوشته و الّذى تكلّم بذلك افترى على اللّه
المهيمن القيّوم اذا يخاطبك ملائكة التّقديس من فردوس الأعظم و يقول يا ايّها
المشرك صه لسانك بقطرة من ابحر تقديسه خلقنا و بعثنا و جعلنا اللّه مطهّراً عن
دلالة غيره و منزّها عن ذكر ما سواه و اذا نسبّحنه و نقدسنه فى هذا المقرّ
الأعظم العظيم

\end{otherlanguage}

Va īnkih bi-qalam-i ghill va iftirā nivishte kih baʿḍī mardum nisbat-i
ʿiṣyān bi-Jamāl-i Abhā mīdihand va nivishte ānchih qalam ḥayā mīnamāyad
az dhikr-i ān kullu mā dhakarta kadhibun ṣurāḥ jamīʿ-i īn kalimāt rā ān
mulqī-yi kadhdhāb bi-tu ilqā namūde va nivishte wa'lladhī takallama
bi-dhālika iftarā ʿalā Allāhi al-muhaymini al-qayyūm idhan yukhāṭibuka
malāʾikatu al-taqdīsi min Firdawsi al-Aʿẓam wa-yaqūlu yā ayyuhā
al-mushriku ṣah lisānaka bi-qaṭratin min abḥuri taqdīsihi khuliqnā
wa-buʿithnā wa-jaʿalanā Allāhu muṭahharan ʿan dalālati ghayrihi
wa-munazzahan ʿan dhikri mā siwāhu wa-idhan nusabbiḥannahu
wa-nuqaddisannahu fī hādhā al-maqarri al-aʿẓami al-ʿaẓīm.

As for what you have written with the pen of malice and slander, that
some people accuse the Most Glorious Beauty of disobedience, and your
having written what the pen is ashamed to mention: everything you have
said is an outright lie. That lying prompter has prompted you with all
these words and written them. Whoever spoke thus has fabricated a lie
against God, the All-Protecting, the Self-Subsisting. Now the angels of
holiness address you from the Most Great Paradise, saying, ``O you who
associate partners with God, silence your tongue! From a drop of the
oceans of His holiness we were created and raised up, and God has made
us pure of the guidance of anyone other than Him and sanctified us from
the mention of all besides Him. Thus we assuredly glorify Him and
proclaim His holiness in this most great and mighty abode.''

\subsection{727}\label{section-1222}

\begin{otherlanguage}{persian}

ثمّ يخاطبك ملائكة التّنزيه من جنّة الأعلی و يقولنّ خذ قلمك يا ايّها المشرك
باللّه تاللّه بهبوب من ارياح تنزيهه قد بعثنا و خلقنا و اذا ننزهنّه و نسبّحنه
و نقدسنه فى هذا المقام الّذى جعله اللّه مطهّراً عن ذكرك و ثنائك و عرفانك و
عرفان العالمين ثمّ يخاطبك ملائكة التّجريد فى جنّة الأبهی و يقولن فاستحى يا
ايّها المعرض عن اللّه قد خلقنا اللّه من ابحر تجريده و انقطعنا عن كلّ ذكر دون
ذكره و اذا نسير فى هذه الجنّة الأعظم و نذكرنه و نحمدنه فى كلّ بكور و اصيل

\end{otherlanguage}

Thumma yukhāṭibuka malāʾikat al-tanzīh min jannat al-Aʿlā wa-yaqulna
khudh qalamaka yā ayyuhā al-mushrik bi-llāh ta-llāh bi-hubūb min aryāḥ
tanzīhihi qad buʿithnā wa-khuliqnā wa-idhan nunazzihannahu
wa-nusabbiḥannahu wa-nuqaddisannahu fī hādhā al-maqām alladhī jaʿalahu
Allāh muṭahharan ʿan dhikrika wa-thanāʾika wa-ʿirfānika wa-ʿirfān
al-ʿālamīn thumma yukhāṭibuka malāʾikat al-tajrīd fī jannat al-Abhā
wa-yaqulna fa-staḥi yā ayyuhā al-muʿriḍ ʿan Allāh qad khalaqanā Allāh
min abḥur tajrīdihi wa-nqaṭaʿnā ʿan kull dhikr dūna dhikrihi wa-idhan
nasīr fī hādhihi al-jannat al-aʿẓam wa-nadhkurannahu wa-naḥmadannahu fī
kull bukūr wa-aṣīl.

Then the angels of transcendence address you from the Most Exalted
Paradise, saying, ``Take up your pen, O you who associate partners with
God! By God, through a breath from the winds of His transcendence we
were raised up and created. Thus we proclaim His transcendence, glorify
Him, and declare His holiness in this station, which God has purified of
your mention, your praise, your recognition, and the recognition of all
the worlds.'' Then the angels of detachment address you in the
All-Glorious Paradise, saying, ``Be ashamed, O you who turn away from
God! God has created us from the oceans of His detachment, and we have
severed ourselves from every mention but His. Thus we walk in this Most
Great Paradise, remembering and praising Him every morning and
evening.''

\subsection{728}\label{section-1223}

\begin{otherlanguage}{arabic}

ثمّ يخاطبك اهل جنّة الأمر و يقولن يا ايّها الغافل عن ذكر اللّه و المحارب
بنفسه انّا خلقنا بامر الّذى جرى من قلمه العزيز البديع و انّه لهو الّذى بأمره
رقم قلم الأعلی على الألواح كلّما اراد و انّه لهو الحاكم على ما يريد انّه لا
يسئل عمّا فعل و ما سواه يسئلون لو انت من العالمين انّه لهو المختار فى نفسه
يفعل ما يشآء بسلطانه ليس لأحد ان يعترض عليه و انّما الاعتراض يرجع الى
نفسك الشّقىّ البعيد ما احصينا احجب منك من ملل القبل تاللّه بنار غلّك اشتعلت
نار السّعير خف عن اللّه و لاتنسب اليه ما يأمرك نفسك و هويك و لاتكن من
المشركين تاللّه يشهد بتقديسه كلّ الذّرّات ثمّ بتنزيهه كلّ الأشیآء و بتجريده
اهل ملأ اعلى و بتوحيد ذاته اهل مداين البقآء انّه لو يحكم لنفس الشّرك حكم
التّوحيد يقدر بسلطانه المقتدر القدير و لو يريد ان يبدّل ذنوب الممكنات
بكلمة من عنده ليكون قادراً بأمره و انّما العجز شأنك و شأن امثالك يا ايّها
الّذى جعلت نفسك محروماً عن نفحات اللّه فى هذه الأیّام الأبدع البديع

\end{otherlanguage}

Thumma yukhāṭibuka ahl jannat al-amr wa-yaqulna yā ayyuhā al-ghāfil ʿan
dhikr Allāh wa-l-muḥārib bi-nafsihi innā khuliqnā bi-amr alladhī jarā
min qalamihi al-ʿazīz al-badīʿ wa-innahu la-huwa alladhī bi-amrihi
raqama qalam al-Aʿlā ʿalā al-alwāḥ kullamā arāda wa-innahu la-huwa
al-ḥākim ʿalā mā yurīd innahu lā yusʾalu ʿammā faʿala wa-mā siwāhu
yusʾalūn law anta min al-ʿālimīn innahu la-huwa al-mukhtār fī nafsihi
yafʿalu mā yashāʾ bi-sulṭānihi laysa li-aḥad an yaʿtariḍa ʿalayhi
wa-innamā al-iʿtirāḍ yarjiʿu ilā nafsika al-shaqiyy al-baʿīd mā aḥṣaynā
aḥjab minka min milal al-qabl ta-llāh bi-nār ghillika ishtaʿalat nār
al-saʿīr khaf ʿan Allāh wa-lā tansub ilayhi mā yaʾmuruka nafsuka
wa-hawāka wa-lā takun min al-mushrikīn ta-llāh yashhadu bi-taqdīsihi
kull al-dharrāt thumma bi-tanzīhihi kull al-ashyāʾ wa-bi-tajrīdihi ahl
malaʾ aʿlā wa-bi-tawḥīd dhātihi ahl madāʾin al-baqāʾ innahu law yaḥkumu
li-nafs al-shirk ḥukm al-tawḥīd yaqdiru bi-sulṭānihi al-muqtadir
al-qadīr wa-law yurīdu an yubaddila dhunūb al-mumkināt bi-kalima min
ʿindihi la-yakūnu qādiran bi-amrihi wa-innamā al-ʿajz shaʾnuka wa-shaʾn
amthālika yā ayyuhā alladhī jaʿalta nafsaka maḥrūman ʿan nafaḥāt Allāh
fī hādhihi al-ayyām al-abdaʿ al-badīʿ.

Then the inhabitants of the Paradise of the Command address you, saying,
``O you who are heedless of the remembrance of God and wage war against
His own Self! We were created by a command that flowed from His mighty,
wondrous Pen. He is indeed the One at whose command the Most Exalted Pen
inscribed on the Tablets whatever He desired. He is indeed the Ruler
over whatever He wills. He is not questioned about what He does, while
all besides Him are questioned, if you are among those who know. He
indeed chooses freely in His own right, doing whatever He wishes through
His sovereignty. No one has the right to object to Him; any objection
only turns back upon your own wretched, wayward self. We have counted
none among the religious communities of the past more veiled than you.
By God, the fire of your malice has kindled the fire of hell! Fear God,
and do not attribute to Him what your own self and desires dictate to
you, and do not be among those who associate partners with Him. By God,
every atom bears witness to His holiness, all things to His
transcendence, the inhabitants of the Concourse on High to His
detachment, and the inhabitants of the cities of eternity to the oneness
of His Essence. If He were to pronounce belief in His oneness upon
unbelief itself, He would have the power through His all-compelling,
all-powerful sovereignty. If He wished to transform the sins of all
contingent beings with a word from Himself, He would assuredly have the
power through His command. Powerlessness belongs only to you and those
like you, O you who have deprived yourself of the breaths of God in
these most wondrous days.''

\subsection{729}\label{section-1224}

\begin{otherlanguage}{persian}

افٍّ لك و لعرفانك اما عرفت بأن الحكم يظهر منه و لايرجع اليه حكم احد و انّه
لهو الأمر فى ازل الازال و لايجرى عليه امر احد لانّ كلّ فقراء عنده و خلقوا
بقوله العزيز البديع هل ينبغى لورقة الّتى سقطت من السّدرة و اصفرّت من هوآء
نفسها يعترض على سدرة الّتى ينطق كلّ ورقة منها بانّه لا اله الّا هو العزيز
المتعالى العليم الحكيم هل ينبغى للعدم بأن يعترض على سلطان القدم او يقنّن
له قوانين نفسه لا فو‌نفسه العلىّ العظيم اذاً اسمع قولى ثمّ استغفر بدوام عمرك
لعلّ اللّه يغفرك و يطهّرك عمّا ابتليت به و ينقطعك عن نفسك و هويك و يقرّبك الى
ما هو خير لك عن ملكوت ملك السّموات و الأرضين فو‌اللّه ما اردت فيما ذكرت الّا
تضييع امر اللّه و نسئل اللّه بأن يضيّعك ثمّ الّذين ايّدوك على ما فعلت و ارتكبت
و انّه لعلى كلّ شیء قدير

\end{otherlanguage}

Uffin laka wa-li-ʿirfānika a-mā ʿarafta bi-anna al-ḥukm yaẓharu minhu
wa-lā yarjiʿu ilayhi ḥukm aḥad wa-innahu la-huwa al-amr fī azal al-āzāl
wa-lā yajrī ʿalayhi amr aḥad li-anna kull fuqarāʾ ʿindahu wa-khuliqū
bi-qawlihi al-ʿazīz al-badīʿ hal yanbaghī li-waraqa allatī saqaṭat min
al-sidra wa-ṣfarrat min hawāʾ nafsihā yaʿtariḍ ʿalā sidra allatī yanṭiqu
kull waraqa minhā bi-annahu lā ilāha illā huwa al-ʿAzīz al-Mutaʿālī
al-ʿAlīm al-Ḥakīm hal yanbaghī li-l-ʿadam bi-an yaʿtariḍa ʿalā sulṭān
al-qidam aw yuqannina lahu qawānīn nafsihi lā fa-wa-nafsihi al-ʿaliyy
al-ʿaẓīm idhan ismaʿ qawlī thumma istaghfir bi-dawām ʿumrika laʿalla
Allāh yaghfiruka wa-yuṭahhiruka ʿammā ibtulīta bihi wa-yanqaṭiʿuka ʿan
nafsika wa-hawāka wa-yuqarribuka ilā mā huwa khayr laka ʿan malakūt mulk
al-samāwāt wa-l-arḍīn fa-wa-llāh mā aradta fīmā dhakarta illā taḍyīʿ amr
Allāh wa-nasʾalu Allāh bi-an yuḍayyiʿaka thumma alladhīna ayyadūka ʿalā
mā faʿalta wa-rtakabta wa-innahu la-ʿalā kull shayʾ qadīr.

Shame on you and your understanding! Have you not recognized that
judgment proceeds from Him, while no one's judgment applies to Him? He
is indeed the Command from all eternity, and no one's command holds sway
over Him, for all are destitute before Him and were created by His
mighty, wondrous word. Is it fitting for a leaf that has fallen from the
Lote Tree and yellowed through the air of its own self to object to a
Tree whose every leaf proclaims that there is no God but Him, the
Almighty, the Most Exalted, the All-Knowing, the All-Wise? Is it fitting
for nothingness to object to the Sovereign of eternity, or to prescribe
its own laws for Him? No, by His own exalted, mighty Self! Listen, then,
to My words, and seek forgiveness throughout your life, that God may
perhaps forgive you, purify you of what has afflicted you, detach you
from your self and desires, and bring you near to what is better for you
than sovereignty over the heavens and the earth. By God, in what you
have said you intended only to bring the Cause of God to ruin. We ask
God to bring you to ruin, and those who supported you in what you have
done and committed. He indeed has power over all things.

\subsection{730}\label{section-1225}

\begin{otherlanguage}{persian}

اين عبد را ملجأ نموده‌ اگر جواب ذكر نشود بيم آنست كه بعضى از نفوس ضعيفه
باين كلمات جعليّه از حقّ محروم مانند و اگر ذكر شود و اللّه الّذى لا اله الّا
هو فطرت اصليّه راضى نميشود كه از ذكر اللّه غافل شود و بكلمات مشركين توجّه
نمايد و لكن چه چاره تكليف اينست آنچه گفته‌ايد جواب ذكر شودلعلّ آن مجعولات
ما بين عباد و سلطان يوم ميعاد حجاب نگردد

\end{otherlanguage}

Īn ʿabd rā muljaʾ namūdeh agar javāb dhekr nashavad bīm ān ast keh baʿżī
az nofūs-e żaʿīfeh be-īn kalemāt-e jaʿlīyyeh az Ḥaqq maḥrūm mānand va
agar dhekr shavad va-llāh alladhī lā ilāha illā huwa feṭrat-e aṣlīyyeh
rāżī nemīshavad keh az dhekr-e Allāh ghāfel shavad va be-kalemāt-e
moshrekīn tavajjoh namāyad va lāken cheh chāreh taklīf īn ast āncheh
gofteh-īd javāb dhekr shavad laʿall ān majʿūlāt mā beyn-e ʿebād va
Solṭān-e yowm-e mīʿād ḥejāb nagardad.

You have left this servant no choice: if no answer is given, there is a
fear that some weak souls may be deprived of the True One through these
fabricated words. Yet if an answer is given---by God, besides whom there
is no other god---one's innate nature is unwilling to become heedless of
the remembrance of God and turn its attention to the words of those who
associate partners with Him. But what can be done? The duty is to answer
what you have said, so that those fabrications may not become a veil
between His servants and the Sovereign of the Day of Promise.

\subsection{731}\label{section-1226}

\begin{otherlanguage}{persian}

و از اين كلمات كه بتو القا نموده‌مقصودى نداشته مگر آنكه بعضى نفوس را از
شطر فردوس منع نمايند و عباد را بدام كلمات صيد كنند در دوازده سنه توقّف در
عراق جميع ميدانند كه چه مقدار جدّ و جهد فرموده‌اند تا آنكه بعضى فواحش را
از ما بين اين طايفه برداشتند و اگر هم نعوذ باللّه بعضى مرتكب شده‌اند در
ستر بوده فو‌اللّه يا غافل مطّلع نيستى كه بر چه مقرّى اين اعتراضات نموده ‌لو
تعرف لتهلك نفسك فى الحين

\end{otherlanguage}

Va az īn kalemāt keh be-to elqā namūdeh maqṣūdī nadāshteh magar ānkeh
baʿżī nofūs rā az shaṭr-e ferdows manʿ namāyand va ʿebād rā be-dām-e
kalemāt ṣeyd konand dar davāzdah saneh tavaqqof dar ʿErāq jamīʿ mīdānand
keh cheh meqdār jedd o jahd farmūdeh-and tā ānkeh baʿżī favāḥesh rā az
mā beyn-e īn ṭāʾefeh bardāshtand va agar ham naʿūdhu bi-llāh baʿżī
mortakeb shodeh-and dar setr būdeh fa-wa-llāh yā ghāfel moṭṭaleʿ nīstī
keh bar cheh maqarrī īn eʿterāżāt namūdeh law taʿrif la-tahlika nafsuka
fī al-ḥīn.

In prompting you with these words, he had no purpose but to keep some
souls from turning toward Paradise and to catch His servants in the
snare of words. Everyone knows how strenuously He exerted Himself during
His twelve years' residence in Iraq, until He eliminated certain
shameful practices from among this community. And if some did commit
them---God forbid---it was in secret. By God, O heedless one, you do not
know against what station you have raised these objections! If you knew,
you would perish that very instant.

\subsection{732}\label{section-1227}

\begin{otherlanguage}{persian}

علّت و سبب اعظم كدورت جمال ابهى از ميرزا يحيى و اللّه الّذى لا اله الّا هو
اين بوده كه در حرم نقطه اولى روح ما سواه فداه تصرّف نمود با اينكه در كلّ
كتب سماوى حرام است و بی‌‌شرمى او بمقامى رسيده كه مخصوص زوجات خود را در
مكتوبات خود حرام نموده‌ مع‌ذلك دست تعدّى و خيانت بحرم مظهر مليك علام گشوده
فافّ له و لوفائه و كاش بنفس خود قناعت مينمود بلكه او را بعد از ارتكاب خود
وقف مشركين نمود و جميع اهل بيان شنيده و ميدانند سيّئات او را بايد از
معاشرين او سؤال نمود

\end{otherlanguage}

ʿEllat va sabab-e aʿẓam-e kodūrat-e Jamāl-e Abhā az Mīrzā Yaḥyā va-llāh
alladhī lā ilāha illā huwa īn būdeh keh dar ḥaram-e Noqṭeh-ye Ūlā rūḥ-e
mā sevāh fadāh taṣarrof namūd bā īnkeh dar koll-e kotob-e samāvī ḥarām
ast va bī-sharmī-ye ū be-maqāmī resīdeh keh makhṣūṣ-e zowjāt-e khod rā
dar maktūbāt-e khod ḥarām namūdeh maʿa-dhālek dast-e taʿaddī va khīyānat
be-ḥaram-e Maẓhar-e Malīk-e ʿAllām goshūdeh fa-uffin lahu wa-li-wafāʾihi
va kāsh be-nafs-e khod qanāʿat mīnamūd balkeh ū rā baʿd az ertekāb-e
khod vaqf-e moshrekīn namūd va jamīʿ-e ahl-e Bayān shenīdeh va mīdānand
sayyeʾāt-e ū rā bāyad az moʿāsherīn-e ū soʾāl namūd.

The chief reason and greatest cause of the All-Glorious Beauty's
displeasure with Mīrzā Yaḥyā---by God, besides whom there is no other
god---was that he took possession of the wife of the Primal Point, may
the souls of all besides Him be sacrificed for Him, although this is
forbidden in all the heavenly Books. His shamelessness reached such a
point that, while he specifically forbade marriage to his own wives in
his writings, he nevertheless stretched out the hand of transgression
and betrayal toward the wife of the Manifestation of the all-knowing
King. Shame on him and his faithfulness! If only he had confined himself
to his own act; instead, after committing it himself, he made her
available to those who associate partners with God. All the people of
the Bayān have heard and know of this. His associates should be asked
about his misdeeds.

\subsection{733}\label{section-1228}

\begin{otherlanguage}{persian}

و از اينها گذشته عجب در اينست كه همين فقره كه نوشته دليلى است واضح و
برهانى است قاطع بر حقّیّت اين امر و تو هنوز مستشعر نشده نفسى را كه خود تو
ميدانى كه از اهل علم و زهد نبوده و از عرفا و ارباب عمايم محسوب نه و بقول
تو با امثال مشير الدّوله معاشر بوده و بعد يكمرتبه كوثر معانى الهيّه من غير
جهات از او جارى و شموس حكمت ربانيّه از افق اراده‌اش مشرق و مضئ اى بی‌ انصاف
اين مقامى است فوق كلّ مقام لو تكون من المنصفين اين فضل اكبر را حقير شمرده
و بخيال خود اين كلمات را بجهت اغواى ناس ذكر كرده قد اخطات يا ايّها المجرم
العاصى هر عاقل بصيرى كه بشنود البتّه در امر اللّه ثابت‌تر و مستقيم‌تر شده و
خواهد شد

\end{otherlanguage}

Va az īnhā godhashteh ʿajab dar īn ast keh hamīn faqreh keh neveshteh
dalīlī ast vāżeḥ va borhānī ast qāṭeʿ bar ḥaqqīyyat-e īn Amr va to hanūz
mostashʿer nashodeh nafsī rā keh khod-e to mīdānī keh az ahl-e ʿelm va
zohd nabūdeh va az ʿorafā va arbāb-e ʿamāyem maḥsūb nah va be-qowl-e to
bā amthāl-e Moshīr al-Dowleh moʿāsher būdeh va baʿd yek-martabeh
kowthar-e maʿānī-ye elāhīyyeh min ghayr-e jahāt az ū jārī va shomūs-e
ḥekmat-e rabbānīyyeh az ofoq-e erādeh-ash moshreq va możīʾ ey bī-enṣāf
īn maqāmī ast fowq-e koll-e maqām law takūnu min al-munṣifīn īn fażl-e
akbar rā ḥaqīr shemordeh va be-khīyāl-e khod īn kalemāt rā be-jehat-e
eghvā-ye nās dhekr kardeh qad akhṭaʾta yā ayyuhā al-mujrim al-ʿāṣī har
ʿāqel-e baṣīrī keh beshnavad albatteh dar Amr-e Allāh thābet-tar va
mostaqīm-tar shodeh va khvāhad shod.

Beyond all this, what is astonishing is that the very passage he has
written is clear evidence and conclusive proof of the truth of this
Cause, yet you have still not perceived it. You yourself know that this
Person was not one of the learned or the ascetics, nor counted among the
mystics and the turbaned clergy, and that, according to you, He
associated with the likes of Moshīr al-Dowleh. Then, all at once, the
heavenly stream of divine meanings flows from Him, without dependence on
any source, and the suns of divine wisdom rise and shine from the
horizon of His will. O unfair one! This is a station above every
station, if you are among the fair-minded. You have belittled this
greatest favor and, as you imagine, spoken these words to lead people
astray. You have erred, O sinful transgressor! Every discerning person
of reason who hears them has certainly become, and will become, firmer
and more steadfast in the Cause of God.

\subsection{734}\label{section-1229}

\begin{otherlanguage}{persian}

بارى فكر ديگر براى خود نموده‌كه اين وساوس در انفس خنّاسيّه اثر نمايد نه در
انفس رحمانيّه قلوب كدره محدوده شايد باينكلمات كذبه نالايقه از ذكر سلطان
احديّه غافل شوند و لكن انفس مجرّده بجان بشطر رحمن توجّه نمايند حال از تو
ميپرسم اگر نفسى از مابين قومى كه ابدا معروف بعلم نبوده و باسم فضل و عدل
مابين ناس مشهور نه برخيزد و بجواهر كلمات اللّه ناطق شود چنانچه شنيده و
ميشنويد اين نفس اعظم است و يا نفسيكه در ميانه اهل علم تربيت شود شكّى نيست
كه اوّل ارجح بوده چه كه او بصرف قدرت الهيّه ظاهر

\end{otherlanguage}

Bārī fekr-e dīgar barā-ye khod namūdeh keh īn vasāves dar anfos-e
khannāsīyyeh athar namāyad nah dar anfos-e raḥmānīyyeh qolūb-e
kadereh-ye maḥdūdeh shāyad be-īn kalemāt-e kadhebeh-ye nālāyeqeh az
dhekr-e Solṭān-e Aḥadīyyeh ghāfel shavand va lāken anfos-e mojarradeh
be-jān be-shaṭr-e Raḥmān tavajjoh namāyand ḥāl az to mīporsam agar nafsī
az mā beyn-e qowmī keh abadan maʿrūf be-ʿelm nabūdeh va be-esm-e fażl va
ʿadl mā beyn-e nās mashhūr nah barkhīzad va be-javāher-e kalemāt-e Allāh
nāṭeq shavad chenāncheh shenīdeh va mīshenavīd īn nafs aʿẓam ast va yā
nafsī keh dar mīyāneh-ye ahl-e ʿelm tarbīyat shavad shakkī nīst keh
avval arjaḥ būdeh cheh keh ū be-ṣerf-e qodrat-e elāhīyyeh ẓāher.

In any case, devise another plan for yourself, for these insinuations
affect souls given to whispering evil, not souls that belong to the
All-Merciful. Clouded, constrained hearts may become heedless of the
remembrance of the Sovereign of Oneness through these false, unworthy
words, but detached souls turn wholeheartedly toward the All-Merciful.
Now I ask you: if someone arises from among a people who have never been
known for learning or renowned among others for excellence and justice,
and gives voice to the gems of God's words, as you have heard and
continue to hear, is this person greater, or someone educated among the
learned? There is no doubt that the former is superior, for He has
appeared solely through divine power.

\subsection{735}\label{section-1230}

\begin{otherlanguage}{persian}

خوب است كه حقّ جلّ اجلاله از لسان خود شما كلماتى فرموده كه متبصّرين را در
معرفت امثال آن نفوس كافى بوده و خواهد بود لازال حقّ محلّ اينگونه اعتراضات
بوده چنانچه بروح اللّه اعتراض نمودند كه با خاطئين و عشارين مجالس و مؤانس
است فو‌اللّه بر اين ظهور مع وصاياى نقطه بيان جلّ اعزازه بلايائى وارد شده كه
بر احدى از مظاهر قبل وارد نشده بلاياى اوّلين و آخرين معادله نمينمايد با
همين قرطاس نجسه كه بساحت اقدس ارسال داشته با آنكه از هيچ امرى مطّلع نبوده
و نيستى و محض هوى و مجرّد نفس تكلّم نموده‌ و اين از رزايائى است كه مشهود
شده و لكن بلايائى بر اين ظهور وارد كه احدى محصى آن نه الّا اللّه

\end{otherlanguage}

Khūb ast keh Ḥaqq jalla ijlāluhu az lesān-e khod-e shomā kalemātī
farmūdeh keh motabaṣṣerīn rā dar maʿrefat-e amthāl-e ān nofūs kāfī būdeh
va khvāhad būd lā yazāl Ḥaqq maḥall-e īngūneh eʿterāżāt būdeh chenāncheh
be-Rūḥ Allāh eʿterāż namūdand keh bā khāṭeʾīn va ʿashshārīn mojāles va
moʾānes ast fa-wa-llāh bar īn Ẓohūr maʿa vaṣāyā-ye Noqṭeh-ye Bayān jalla
iʿzāzuhu balāyāʾī vāred shodeh keh bar aḥadī az Maẓāher-e qabl vāred
nashodeh balāyā-ye avvalīn va ākherīn moʿādeleh nemīnamāyad bā hamīn
qerṭās-e najeseh keh be-sāḥat-e aqdas ersāl dāshteh bā ānkeh az hīch
amrī moṭṭaleʿ nabūdeh va nīstī va maḥż-e havā va mojarrad-e nafs
takallom namūdeh va īn az razāyāʾī ast keh mashhūd shodeh va lāken
balāyāʾī bar īn Ẓohūr vāred keh aḥadī moḥṣī-ye ān nah illā Allāh.

It is well that the True One---exalted be His glory---has spoken through
your own tongue words that have been, and will be, sufficient for the
discerning to recognize souls such as these. The True One has always
been the object of such objections, just as they objected to the Spirit
of God for sitting and associating with sinners and tax collectors. By
God, despite the counsels of the Point of the Bayān---exalted be His
honor---afflictions have befallen this Revelation that have befallen
none of the Manifestations of the past. The afflictions of the first and
the last do not equal this very defiled sheet that you have sent to the
Most Holy Court, although you have had, and still have, no knowledge of
any of these matters, and have spoken out of sheer desire and self
alone. This is one of the calamities that have been witnessed, yet
afflictions have befallen this Revelation that no one but God can count.

\subsection{736}\label{section-1231}

\begin{otherlanguage}{persian}

چنانچه در توقيع جناب جود عليه بهآء اللّه يك آيه نازل كه و اللّه كبد اهل
آفاق‌ را گداخته قوله جلّ كبريائه انّك رأيتنى و عاشرت معى و عرفت بحر سكونى و
جبل اصطبارى فكّر ما اقامنى على الصّيحة بين السّموات و الأرضين و لكن چه
فايده كه آيات الهى بر تو القا نموده مثل آن است كه عندليبی‌ از براى كر
تغنّى نمايد نسئل اللّه بأن يطهّر الأرض من رجسك و رجس الّذينهم كانوا معك و
انّه لمجيب دعوة السّايلين قدرت آن قدر معلوم شده كه بقول اخبث ناس و ادناهم
از حقّ معرض شده و بشجره نفس و هوى متمسّك

\end{otherlanguage}

Chenāncheh dar towqīʿ-e Janāb-e Jūd ʿalayhi bahāʾ Allāh yek āyeh nāzel
keh va-llāh kabed-e ahl-e āfāq rā godākhteh qawluhu jalla kibriyāʾuhu
innaka raʾaytanī wa-ʿāsharta maʿī wa-ʿarafta baḥr sukūnī wa-jabal
iṣṭibārī fakkir mā aqāmanī ʿalā al-ṣayḥa bayna al-samāwāt wa-l-arḍīn va
lāken cheh fāyedeh keh āyāt-e elāhī bar to elqā namūdeh methl-e ān ast
keh ʿandalībī az barā-ye kar taghannī namāyad nasʾalu Allāh bi-an
yuṭahhira al-arḍ min rijsika wa-rijs alladhīnahum kānū maʿaka wa-innahu
la-mujīb daʿwat al-sāʾilīn qadrat ān qadr maʿlūm shodeh keh be-qowl-e
akhbath-e nās va adnāhom az Ḥaqq moʿreż shodeh va be-shajareh-ye nafs va
havā motamassek.

Thus, in the Tablet to Jināb-i-Jūd---upon him be the glory of God---a
verse was revealed that, by God, has melted the hearts of the people of
every land. He says---exalted be His majesty---``You have seen Me,
associated with Me, and known the ocean of My calm and the mountain of
My patience. Consider what has made Me raise My cry between the heavens
and the earth.'' But what use is there in conveying divine verses to
you? It is like a nightingale singing to someone deaf. We ask God to
cleanse the earth of your defilement and the defilement of those who
were with you. He indeed answers the prayers of those who ask. Your
worth has become clear to this extent: at the word of the vilest and
most abject of people, you have turned away from the True One and clung
to the tree of self and desire.

\subsection{737}\label{section-1232}

\begin{otherlanguage}{persian}

و ديگر نوشته بودى كه در مجلسها صحبت ميشود كه جمال ابهى بناى فساد و فتنه
گذاشته‌اند و مشهور شده كه اوّل برادرش را قائم و پنهان كرده باسم او نانى
پخت حالا او را در كرده ميخواهند آجيل بخورند انتهى

\end{otherlanguage}

Va dīgar neveshteh būdī keh dar majleshā ṣoḥbat mīshavad keh Jamāl-e
Abhā banā-ye fasād va fetneh godhāshteh-and va mashhūr shodeh keh avval
barādarash rā qāʾem va penhān kardeh be-esm-e ū nānī pokht ḥālā ū rā dar
kardeh mīkhvāhand ājīl bokhorand intahā.

You had also written, ``It is being said in gatherings that the
All-Glorious Beauty has set about fomenting corruption and discord. It
has become common talk that He first set up His brother as the Qāʾim and
hid him away, baking bread for Himself in his name; now He has cast him
out and wants to eat nuts.'' End of quotation.

\subsection{738}\label{section-1233}

\begin{otherlanguage}{persian}

لعن اللّه من القاك و لعن اللّه من كذب و افترى ثمّ بغى و اطغى ثمّ كفر و اشقى
و اللّه بسيار ناس را احمق يافته جميع ميدانند كه ایشان حسب الأمر دولت
ايران به عراق تشريف آوردند و ميرزا يحيى بعد از مقدمه ارض طاء فرار
نموده‌به عراق آمد قل ان يا ايّها المشرك انّ الّذى انفق روحه فى سبيل اللّه و
كان بين يدى النّاس و ورد عليه فى كلّ حين سهام القضا ما لم يحصيها احد الّا
اللّه الملك العليم الخبير و قام عليه اكثر النّاس فى ایّام الّتى كان فى
العراق و فى كلّ يوم ينادى منهم بأن يأخذوه و يرسلوه الى ارض الطّاء و منهم
من قال فسوف يغرقونه فى الشّط و منهم من قال ينفيه الدّولة و كان اللّه على
ذلك لشهيد و عليم و بكت عليه عيون كلّ الأشیآء بما مسّته البأسآء و الضرآء و
انّك كتبت ما كتبت الا لعنة اللّه على الكاذبين

\end{otherlanguage}

Laʿana Allāhu man alqāka wa-laʿana Allāhu man kadhaba wa-iftarā thumma
baghā wa-aṭghā thumma kafara wa-ashqā. Va-Allāh, bisyār nās rā aḥmaq
yāftih. Jamīʿ mīdānand kih īshān ḥasb al-amr-i dawlat-i Īrān bih ʿIrāq
tashrīf āvardand va Mīrzā Yaḥyā baʿd az muqaddamih-yi Arḍ-i Ṭāʾ farār
nimūdih bih ʿIrāq āmad. Qul an yā ayyuhā al-mushrik, inna alladhī anfaqa
rūḥahu fī sabīli Allāh wa-kāna bayna yaday al-nās wa-warada ʿalayhi fī
kulli ḥīn sihāmu al-qaḍāʾ mā lam yuḥṣihā aḥad illā Allāhu al-Maliku
al-ʿAlīmu al-Khabīr, wa-qāma ʿalayhi aktharu al-nās fī ayyām allatī kāna
fī al-ʿIrāq, wa-fī kulli yawm yunādī minhum bi-an yaʾkhudhūhu
wa-yursilūhu ilā arḍ al-Ṭāʾ, wa-minhum man qāla fa-sawfa yughriqūnahu fī
al-shaṭṭ, wa-minhum man qāla yanfīhi al-dawlah, wa-kāna Allāhu ʿalā
dhālika la-shahīd wa-ʿalīm. Wa-bakat ʿalayhi ʿuyūnu kulli al-ashyāʾ
bi-mā massathu al-baʾsāʾ wa-al-ḍarrāʾ, wa-innaka katabta mā katabta. Alā
laʿnatu Allāhi ʿalā al-kādhibīn.

May God curse the one who prompted you, and may God curse the one who
lied and fabricated, then acted unjustly and led others into
transgression, then disbelieved and brought misery. By God, he has found
people very foolish. Everyone knows that He came to Iraq at the command
of the government of Iran, while Mīrzā Yaḥyā fled after the incident in
the Land of Ṭā and came to Iraq. Say: O idolater! The One who expended
His life in the path of God was in full view of the people, and at every
moment the arrows of fate rained upon Him in numbers that no one but
God, the King, the All-Knowing, the All-Informed, could count. Most of
the people rose against Him during His days in Iraq. Every day some
called for Him to be seized and sent to the Land of Ṭā; some said,
``They will drown Him in the river''; and others said, ``The government
will exile Him.'' God indeed witnessed and knew all this. The eyes of
all things wept for Him because of the suffering and adversity that
afflicted Him, yet you wrote what you wrote. Surely God's curse rests
upon the liars.

\subsection{739}\label{section-1234}

\begin{otherlanguage}{persian}

كلّ اشیآء شاهد و گواهند كه در سنين توقّف در عراق ابدا اميد جان نبود و در
كلّ حين ببلاياى جديده مبتلا و مع‌ذلك امثال آن نفوس رذليّه ذكر نان و آجيل
مينمايند و ملاحظه كن كه چه قدر از عقل ظاهره محروم شده و اللّه هر نفسى
كلماتت را ملاحظه نمايد قهر حقّ را ظاهر و آشكار در تو مشاهده مينمايد سبحان
اللّه نفسيكه عاجز است دو كلمه پارسى را بعبارات حسنه تحرير نمايد اعترض على
الّذى بارادة منه خلق قلم الأعلی و ظهر منه الصّرير الّذى به استجذبت قلوب ما
كان و ما يكون افّ لك و بما ظهر منك

\end{otherlanguage}

Kull-i ashyāʾ shāhid va guvāhand kih dar sinīn-i tavaqquf dar ʿIrāq
abadan umīd-i jān nabūd va dar kull-i ḥīn bih-balāyā-yi jadīdih mubtalā,
va maʿa dhālik amthāl-i ān nufūs-i radhīliyyih dhikr-i nān va ājīl
mīnamāyand. Va mulāḥiẓih kun kih chih qadr az ʿaql-i ẓāhirih maḥrūm
shudih. Va-Allāh, har nafsī kalimātat rā mulāḥiẓih namāyad qahr-i Ḥaqq
rā ẓāhir va āshkār dar tu mushāhidih mīnamāyad. Subḥān Allāh, nafsī kih
ʿājiz ast du kalimih-yi pārsī rā bih-ʿibārāt-i ḥasanih taḥrīr namāyad,
iʿtaraḍa ʿalā alladhī bi-irādatin minhu khuliqa Qalamu al-Aʿlā wa-ẓahara
minhu al-ṣarīru alladhī bihi istujdhibat qulūbu mā kāna wa-mā yakūn.
Uffin laka wa-bi-mā ẓahara minka.

All things bear witness and testify that during the years of residence
in Iraq there was never any hope of survival, and at every moment new
afflictions arose. Yet the likes of those vile souls speak of bread and
nuts. Consider how utterly bereft of ordinary reason you have become. By
God, anyone who reads your words sees God's wrath plainly and manifestly
in you. Glory be to God! Someone unable to write two words of Persian in
a graceful manner has objected to the One by whose will the Most Exalted
Pen was created and from whom came the sound of its writing, through
which the hearts of all that has been and will be were drawn. Shame on
you and on what has proceeded from you!

\subsection{740}\label{section-1235}

\begin{otherlanguage}{persian}

و لكن از اين مردم غافل بعيد نيست كه همين كلمات تو را بر كلمات حقّ ترجيح
دهند چنانچه سنين لانهاية لهاعجل را بر حقّ ترجيح داده و پرستيده‌اند و حال
هم بآن مشغولند بسيار حيف است كه انسان باينگونه كلمات نالايقه تكلّم نمايد
چه كه كلمات حاكيند از نفس انسانى لذا كلمات اهل منظر اكبر مزيّن است بطراز
ادب نسئل اللّه بأن لا يجعلنا عريّا من هذا الثّوب الّذى به يظهر قدر الانسان
بين ملأ الأكوان و امتيازه عن الحيوان ثمّ نسئله بأن يقرّبنا عليه و ينقطعنا
عن دونه و يطهّرنا عن روايح الوهم و التّقليد و يجعلنا من الّذين قالوا اللّه
ربّنا ثمّ استقاموا و ما منعهم استهزاء الخلايق اجمعين و لاشماتة الغافلين

\end{otherlanguage}

Va-lākin az īn mardum-i ghāfil baʿīd nīst kih hamīn kalimāt-i tu rā bar
kalimāt-i Ḥaqq tarjīḥ dihand, chunānchih sinīn lā nihāyata lahā ʿijl rā
bar Ḥaqq tarjīḥ dādih va parastīdih-and va ḥāl ham bi-ān mashghūland.
Bisyār ḥayf ast kih insān bi-īngūnih kalimāt-i nālāyiqih takallum
namāyad, chih kih kalimāt ḥākī-and az nafs-i insānī. Lidhā kalimāt-i
ahl-i Manẓar-i Akbar muzayyan ast bih-ṭirāz-i adab. Nasʾalu Allāha bi-an
lā yajʿalanā ʿariyyan min hādhā al-thawb alladhī bihi yaẓharu qadru
al-insān bayna malaʾ al-akwān wa-imtiyāzuhu ʿan al-ḥayawān. Thumma
nasʾaluhu bi-an yuqarribanā ʿalayhi wa-yanqaṭiʿanā ʿan dūnihi
wa-yuṭahhiranā ʿan rawāʾiḥ al-wahm wa-al-taqlīd wa-yajʿalanā min
alladhīna qālū Allāhu rabbunā thumma istaqāmū wa-mā manaʿahum istihzāʾu
al-khalāʾiq ajmaʿīn wa-lā shamātatu al-ghāfilīn.

Yet it would not be surprising if these heedless people preferred these
very words of yours to the words of the True One, just as for endless
years they have preferred the calf to the True One and worshipped it,
and are still doing so. It is a great pity for a human being to utter
such unworthy words, for words reflect the human soul. Thus the words of
the people of the Most Great View are adorned with the ornament of
courtesy. We ask God not to leave us stripped of this robe, through
which human worth becomes manifest among all created beings, as does the
distinction between humans and animals. Then we ask Him to draw us near
to Him, detach us from all besides Him, purify us of the odors of
delusion and imitation, and make us among those who said, ``God is our
Lord,'' and then stood firm, undeterred by the mockery of all creatures
or the gloating of the heedless.

\subsection{741}\label{section-1236}

\begin{otherlanguage}{persian}

اللّهمّ يا الهى فافتح ابصار هؤلآء ليروك ظاهراً بين خلقك و مشرقاً فى مملكتك و
انّك يا الهى لو فتحت ابصارهم ما ابتليت باستهزائهم و احجار ظنونهم و سهام
اوهامهم كما فتحت ابصار احبّائك و عرفتهم ما لاعرفته دونهم و انّك لو كشفت
الحجبات لهم كما كشفت عن وجوههم ما احتجبوا و ما اعترضوا فلمّا سترت عن
هؤلآء و كشفت لاحبّائك لذا ارتفع ضجيج الغافلين من بريّتك و صريخ المتوهّمين
من اهل مملكتك اذا اسئلك بنفسك بأن تكشف لهم حجبات الّتى منعتهم عن عرفانك و
عرفان مظهر نفسك ليجمعنّ كلّ على شاطى بحر توحيدك و مقرّ عزّ تقديسك و تفريدك و
انّك انت على ما تشآء قدير

\end{otherlanguage}

Allāhumma yā ilāhī, fa-ftaḥ abṣāra hāʾulāʾ li-yarawka ẓāhiran bayna
khalqika wa-mushriqan fī mamlakatika. Wa-innaka yā ilāhī law fataḥta
abṣārahum mā ibtulīta bi-istihzāʾihim wa-aḥjāri ẓunūnihim wa-sihāmi
awhāmihim, kamā fataḥta abṣāra aḥibbāʾika wa-ʿarraftahum mā lā ʿarafathu
dūnahum. Wa-innaka law kashafta al-ḥujubāt lahum kamā kashafta ʿan
wujūhihim mā iḥtajabū wa-mā iʿtaraḍū. Fa-lammā satarta ʿan hāʾulāʾ
wa-kashafta li-aḥibbāʾika lidhā irtafaʿa ḍajīju al-ghāfilīn min
bariyyatika wa-ṣarīkhu al-mutawahhimīn min ahli mamlakatika. Idhan
asʾaluka bi-nafsika bi-an takshifa lahum ḥujubāt allatī manaʿathum ʿan
ʿirfānika wa-ʿirfāni maẓhari nafsika li-yajmaʿanna kullun ʿalā shāṭiʾi
baḥri tawḥīdika wa-maqarri ʿizzi taqdīsika wa-tafrīdika. Wa-innaka anta
ʿalā mā tashāʾ qadīr.

O God, my God! Open the eyes of these people, that they may see You
manifest among Your creatures and shining in Your kingdom. My God, had
You opened their eyes, You would not have been afflicted by their
mockery, the stones of their suspicions, and the arrows of their
imaginings, just as You opened the eyes of Your loved ones and made
known to them what no others knew. Had You lifted the veils for these
people as You lifted them from the faces of Your loved ones, they would
not have remained veiled or raised objections. But because You concealed
Yourself from these and revealed Yourself to Your loved ones, the clamor
of the heedless among Your creatures and the cries of the deluded among
the people of Your kingdom have risen. I therefore ask You by Your own
Self to lift for them the veils that have kept them from recognizing You
and the Manifestation of Your Self, so that all may gather on the shore
of the ocean of Your oneness and at the seat of the glory of Your
holiness and singleness. You indeed have power over whatever You will.

\subsection{742}\label{section-1237}

\begin{otherlanguage}{persian}

آنانكه ذكر نموده‌اند كه ايشان نان باسم او درست نموده ‌و اللّه خجالت ميبرم
از ذكر چنين نفوس بی‌‌ادب رذلى و آن خود سيّد محمّد خبيث است اين اذكار اذكار
اولى النّجابه نبوده و نخواهد بود و هر نفسى كه چنين كلمه تكلّم نموده‌ مثل
آنست كه كسى بگويد شمس از ظلمت مستنير شده و يا عندليب از زاغ تغنّى آموخته
و يا قطعهء ياقوت از حجر كسب ضيا و لون كرده چه فايده كه آن غافل مطّلع نيست
و نفوسى كه مابين بوده جميع را تكذيب نموده‌اند كه مبادا حرف صدقى بميان آيد
و ناس بشنوند

\end{otherlanguage}

Ānān kih dhikr nimūdih-and kih īshān nān bih-ism-i ū durust
nimūdih---va-Allāh khijālat mībaram az dhikr-i chunīn nufūs-i bī-adab-i
radhilī---va ān khud-i Siyyid Muḥammad-i khabīth ast. Īn adhkār adhkār-i
ulī al-najābih nabūdih va nakhvāhad būd, va har nafsī kih chunīn kalimih
takallum nimūdih mithl-i ānast kih kasī bigūyad shams az ẓulmat mustanīr
shudih va yā ʿandalīb az zāgh taghannī āmūkhtih va yā qiṭʿih-yi yāqūt az
ḥajar kasb-i ḍiyā va lawn kardih. Chih fāʾidih kih ān ghāfil muṭṭaliʿ
nīst, va nufūsī kih mābayn būdih jamīʿ rā takdhīb nimūdih-and kih mabādā
ḥarf-i ṣidqī bih-miyān āyad va nās bishnavand.

Those who have stated that He made bread in his name---by God, I am
ashamed to mention such vile, ill-mannered souls---and it is that wicked
Siyyid Muḥammad himself. Such remarks have never been, nor will they
ever be, the remarks of noble people. Anyone who has uttered such a
statement is like someone saying that the sun has been illumined by
darkness, that the nightingale has learned to sing from the crow, or
that a ruby has acquired its radiance and color from a stone. But what
use is it? That heedless one is unaware, and they have branded as liars
all those who were present, lest a truthful word be uttered and the
people hear it.

\subsection{743}\label{section-1238}

\begin{otherlanguage}{persian}

سيّد محمّدى كه و اللّه العظيم افتخار مينمود كه در حضور مبارك بايستد حال
آمده و اينگونه مفتريات نسبت بحقّ داده حمد خدا را كه اين عباد بغير حقّ و
راستى تكلّم ننموده‌ و نخواهيم نمود و دين بدنيا مبادله ننموده‌ايم و از حقّ
چشم نپوشيده‌ايم جميع شكر مينمايند كه بايمان باللّه در اين ظهور فايز شده‌اند
و اين عبد ناطق يا ليتنى ما ولدت من امّى و لو ولدت ما رضعت و لو رضعت ما
اشدّ ظهرى فيا ليت كنت من الفاقدين در هر نفسى كه اقلّ من ذرّ حبّ اللّه باشد
راضى نشده و نخواهد شد كه چنين امور مشاهده نمايد الى اللّه فوضنا الأمر و
عليه توكّلنا و نسئل من بدايع جوده بأن يوفّقنا على الاقرار بوحدانيّته و
الاعتراف بفردانيّته و ينقطعنا عن دونه و انّه لعلى كلّ شیء قدير ثمّ نسئله
باسم الّذى به فصّل بين الحقّ و الباطل بأن يوفّقك على ذكره و عرفان مظهر نفسه
و يطهّرك عن اشارات هؤلآء الكاذبين

\end{otherlanguage}

Siyyid Muḥammadī kih va-Allāh al-ʿAẓīm iftikhār mīnamūd kih dar ḥuḍūr-i
mubārak bāyistad ḥāl āmadih va īngūnih muftariyāt nisbat bih-Ḥaqq dādih.
Ḥamd-i Khudā rā kih īn ʿibād bih-ghayr-i ḥaqq va rāstī takallum
nanimūdih va nakhvāhīm nimūd va dīn bih-dunyā mubādilih nanimūdih-īm va
az Ḥaqq chashm napūshīdih-īm. Jamīʿ shukr mīnamāyand kih bih-īmān
bi-Allāh dar īn ẓuhūr fāʾiz shudih-and va īn ʿabd nāṭiq: yā laytanī mā
wulidtu min ummī wa-law wulidtu mā raḍiʿtu wa-law raḍiʿtu mā ushidda
ẓahrī fa-yā layta kuntu min al-fāqidīn. Dar har nafsī kih aqall-i min
dharr ḥubb Allāh bāshad rāḍī nashudih va nakhvāhad shud kih chunīn umūr
mushāhidih namāyad. Ilā Allāhi fawwaḍnā al-amr wa-ʿalayhi tawakkalnā
wa-nasʾalu min badāʾiʿi jūdihi bi-an yuwaffiqanā ʿalā al-iqrār
bi-waḥdāniyyatihi wa-al-iʿtirāf bi-fardāniyyatihi wa-yanqaṭiʿanā ʿan
dūnihi wa-innahu la-ʿalā kulli shayʾ qadīr. Thumma nasʾaluhu bi-ismi
alladhī bihi fuṣṣila bayna al-ḥaqq wa-al-bāṭil bi-an yuwaffiqaka ʿalā
dhikrihi wa-ʿirfāni maẓhari nafsihi wa-yuṭahhiraka ʿan ishārāti hāʾulāʾ
al-kādhibīn.

Siyyid Muḥammad, who, by God the Almighty, used to take pride in
standing in the blessed presence, has now come forward and attributed
such fabrications to the True One. Praise be to God that we servants
have spoken nothing but truth and honesty, and never will speak
otherwise. We have not exchanged faith for worldly gain, nor have we
turned our eyes away from the True One. All give thanks that they have
attained faith in God in this Revelation, while this servant says: ``If
only I had not been born of my mother! And if I had been born, had not
been suckled; and if I had been suckled, my back had not been
strengthened. If only I had been among those lost!'' No soul possessing
even less than a mote of the love of God has ever been, or ever will be,
content to witness such things. To God we have committed the matter, and
in Him we have placed our trust. We ask, from the wonders of His bounty,
that He enable us to affirm His unity and acknowledge His singleness,
and detach us from all besides Him. He indeed has power over all things.
Then we ask Him, by the Name through which truth was distinguished from
falsehood, to enable you to remember Him and recognize the Manifestation
of His Self, and to purify you of the insinuations of these liars.

\subsection{744}\label{section-1239}

\begin{otherlanguage}{persian}

فو‌اللّه الّذى لا اله الّا هو بلايائى در اين ارض وارد كه احدى قادر بر احصاى
آن نه بكرّات فرموده‌اند آنچه در سجن و تحت سلاسل و اغلال بوده‌ام و در دست
مشركين مبتلا بشأنيكه ديار بديار برده‌اند و سر و پاى برهنه از محلّ بمحلّ
كشيده‌اند هيچكدام از بلا محسوب نه اعظم بليّه كه در ابداع شبه نداشته و
نخواهد داشت آن است كه حقّ بدست ظالمى كه باسم عدل معروف است مبتلا شود چه
كه آن بی‌انصاف اعتمادا على حمقاء الأرض لايخاف و لايبالى يفعل ما يأمره
النّفس و الهوى چنانچه مشاهده ميشود نفسى را كه آن همه تربيت و حفظ نمودم
بعد استكباراً على اللّه و اعراضاً منه و اعتراضاً عليه وارد آورده آنچه را هيچ
ظالمى وارد نياورده و علّت اين ظلم اطمينان از عباد غافلين بوده چه كه اگر
موقن بود باينكه ناس صاحب بصرند و اعمالش را تصديق نمينمايند بچنين اعمال
ارتكاب نمينمود

\end{otherlanguage}

Fa-wa-Allāh alladhī lā ilāha illā huwa, balāyāʾī dar īn arḍ vārid kih
aḥadī qādir bar iḥṣā-yi ān nah. Bih-karrāt farmūdih-and: ānchih dar sijn
va taḥt-i salāsil va aghlāl būdih-am va dar dast-i mushrikīn mubtalā,
bih-shaʾnī kih diyār bih-diyār burdih-and va sar va pāy barahnih az
maḥall bih-maḥall kashīdih-and, hīchkudām az balā maḥsūb nah. Aʿẓam
baliyyih kih dar ibdāʿ shibh nadāshtih va nakhvāhad dāsht ān ast kih
Ḥaqq bih-dast-i ẓālimī kih bih-ism-i ʿadl maʿrūf ast mubtalā shavad,
chih kih ān bī-inṣāf iʿtimādan ʿalā ḥumaqāʾ al-arḍ lā yakhāfu wa-lā
yubālī yafʿalu mā yaʾmuruhu al-nafs wa-al-hawā. Chunānchih mushāhidih
mīshavad nafsī rā kih ān hamih tarbiyat va ḥifẓ nimūdam baʿd istikbāran
ʿalā Allāh wa-iʿrāḍan minhu wa-iʿtirāḍan ʿalayhi vārid āvardih ānchih rā
hīch ẓālimī vārid nayāvardih. Va ʿillat-i īn ẓulm iṭmīnān az ʿibād-i
ghāfilīn būdih, chih kih agar mūqin būd bi-īnkih nās ṣāḥib-i baṣarand va
aʿmālash rā taṣdīq namīnamāyand bih-chunīn aʿmāl irtikāb namīnimūd.

By God, besides whom there is no other god, afflictions have befallen us
in this land that no one can count. He has repeatedly said: ``All that I
suffered in prison, in chains and fetters, and at the hands of the
idolaters, who took Me from land to land and dragged Me, bareheaded and
barefoot, from place to place---none of this counts as affliction. The
greatest affliction, which has never had and will never have its like in
creation, is for the True One to fall into the hands of an oppressor
known for justice. For that unjust man, relying on the fools of the
earth, has neither fear nor concern, and does whatever self and desire
command him.'' As can be seen, the person whom I nurtured and protected
so much afterward, out of arrogance toward God, aversion to Him, and
opposition to Him, inflicted what no oppressor had ever inflicted. The
cause of this oppression was his confidence in the heedless servants,
for had he been certain that the people possessed discernment and would
not endorse his deeds, he would not have committed such acts.

\subsection{745}\label{section-1240}

\begin{otherlanguage}{persian}

لازال باكثرى از عباد هر قدر عنايت بيشتر شده بيشتر بكفران قيام نموده‌اند
حكايت شيطان از قبل ذكر شد كه در هر يك از سموات باسمى معروف و در ملكوت
اعلى بكمال وصف موصوف و قرنها در قيام و ركوع و قعود و سجود بوده و بعد
چنان شد كه بر مراصد جالس و عباد را از شطر رحمن منع مينمود و كاش انفس
مشركه در اين ایّام باين مقدار كفايت مينمودند لا فو‌الّذى نفسى بيده لن يسكن
عطش بغضهم الّا بشرب دمى و كان الشّيطان استعاذ بين يدى الرّحمن منهم و من
اعمالهم و لكن النّاس هم لا يفقهون و لا يشعرون

\end{otherlanguage}

Lā yazāl bi-aktharī az ʿibād har qadr ʿināyat bīshtar shudih bīshtar
bih-kufrān qiyām nimūdih-and. Ḥikāyat-i shayṭān az qabl dhikr shud kih
dar har yak az samāvāt bih-ismī maʿrūf va dar malakūt-i aʿlā bih-kamāl-i
vaṣf mawṣūf va qarnhā dar qiyām va rukūʿ va quʿūd va sujūd būdih, va
baʿd chunān shud kih bar marāṣid jālis va ʿibād rā az shaṭr-i Raḥmān
manʿ mīnamūd. Va kāsh anfus-i mushrikih dar īn ayyām bi-īn miqdār
kifāyat mīnamūdand. Lā fa-wa-alladhī nafsī bi-yadihi lan yaskuna ʿaṭashu
bughḍihim illā bi-shurbi damī, wa-kāna al-shayṭānu istaʿādha bayna yaday
al-Raḥmān minhum wa-min aʿmālihim, wa-lākin al-nāsu hum lā yafqahūn
wa-lā yashʿurūn.

It has always been the case with most servants that the more favor they
have received, the more they have risen in ingratitude. The story of
Satan was mentioned earlier: he was known by a name in each of the
heavens, was described in the most perfect terms in the kingdom on high,
and for centuries stood, bowed, sat, and prostrated himself in worship.
Yet afterward he came to sit in ambush, keeping the servants from
turning toward the All-Merciful. If only the idolatrous souls in these
days had contented themselves with that much! No, by Him in whose hand
is my soul, the thirst of their hatred will never be quenched except by
drinking my blood. Satan himself has sought refuge before the
All-Merciful from them and their deeds, but the people neither
understand nor perceive.

\subsection{746}\label{section-1241}

\begin{otherlanguage}{persian}

معلوم نشد كه مقصود از آجيل كه از قلم نحست جارى شده چه بوده هرگز اشیآء
فانيه و زخارف دنيويه در ميان نبوده گويا مقصود دراهمى كه در هر شهر از باب
عالى مخصوص خدّام ميفرستاده‌اند بوده اگر اينست حال كه راجع بمشرك باللّه شده
و بمحبوب خود فايز شدند افّ لهم و لحيائهم و لوفائهم از باب اللّه منقطع شده
بهر بابی‌ توجّه نموده‌و التجا نموده‌اند و از جمال قدم شكايت کرده‌‌اند تا آنكه
بمطلوب خود رسيدند خوب حقّ يافتى محكم متمسّك شو كه شايد بمتابعت آن نفوس
باسفل النّيران وارد شوى

\end{otherlanguage}

Maʿlūm nashud kih maqṣūd az ājīl kih az qalam-i naḥsat jārī shudih chih
būdih. Hargiz ashyāʾ-i fāniyih va zakhārif-i dunyaviyyih dar miyān
nabūdih. Gūyā maqṣūd darāhimī kih dar har shahr az Bāb-i ʿĀlī makhṣūṣ-i
khuddām mīfiristādih-and būdih. Agar īnast, ḥāl kih rājiʿ bih-mushrik
bi-Allāh shudih va bih-maḥbūb-i khud fāʾiz shudand. Uffin lahum
wa-li-ḥayāʾihim wa-li-wafāʾihim. Az bāb Allāh munqaṭiʿ shudih bih-har
bābī tavajjuh nimūdih va iltijā nimūdih-and va az Jamāl-i Qidam shikāyat
kardih-and tā ānkih bih-maṭlūb-i khud rasīdand. Khūb ḥaqq yāftī, muḥkam
mutamassik shaw kih shāyad bih-mutābiʿat-i ān nufūs bih-asfal-i nīrān
vārid shavī.

It is not clear what was meant by the nuts that flowed from your
ill-omened pen. Perishable things and worldly vanities were never at
issue. Apparently, the reference was to the sums of money that the
Sublime Porte sent in each city specifically for the attendants. If that
is what was meant, these have now passed to the one who associates
partners with God, and they have attained their beloved. Shame on them,
their sense of shame, and their fidelity! Having cut themselves off from
God's door, they have turned to every door, sought refuge there, and
complained of the Ancient Beauty until they attained their object. What
a fine truth you have found! Cling firmly to it, that perhaps, by
following those souls, you may enter the lowest depths of the fires.

\subsection{747}\label{section-1242}

\begin{otherlanguage}{persian}

آىغافل نفوس رذيله كه باطّلاع نفس ملحدت مخصوص زخارف چه مقدار ذلّت نموده‌اند
و تضييع امر کرده‌اند مع‌ذلك چنين نفوس را حقّ دانسته و از حقّى كه ابدا اعتنا
بنفسى نفرموده بلكه خدّامش اعتنا نداشته‌اند معرض شده لعن اللّه من شهد برهان
اللّه و انكره و عرف تنزيهه عمّا سواه و اعرض عنه

\end{otherlanguage}

Ay ghāfil, nufūs-i radhīlih kih bi-iṭṭilāʿ-i nafs-i mulḥidat makhṣūṣ-i
zakhārif chih miqdār dhillat nimūdih-and va taḍyīʿ-i Amr kardih-and,
maʿa dhālik chunīn nufūs rā ḥaqq dānistih va az Ḥaqqī kih abadan iʿtinā
bih-nafsī nafarmūdih balkih khuddāmash iʿtinā nadāshtih-and muʿriḍ
shudih. Laʿana Allāhu man shahida burhāna Allāh wa-ankarahu wa-ʿarafa
tanzīhahu ʿammā siwāhu wa-aʿraḍa ʿanhu.

O heedless one! As your own impious self knows, those vile souls have
abased themselves to such a degree solely for worldly vanities, and have
dishonored the Cause. Yet you have regarded such souls as true and
turned away from the True One, who has never paid heed to
anyone---indeed, even His attendants have paid them no heed. May God
curse the one who witnessed God's proof and denied it, and recognized
His transcendence above all besides Him yet turned away from Him.

\subsection{748}\label{section-1243}

\begin{otherlanguage}{persian}

اى ملأ بيان اعتراضات اين مردود كه بالقاى ميرزا يحيى و سيّد محمّد نوشته
مخصوص ذكر شد تا بدانيد كه چه مقدار ضرّ وارد شده و ادراك نمائيد نوحه و
ندبه نقطه اولى را بر ظهور آخر و آن جمال مبارك چون مطّلع بوده‌اند بر نفوس
مشركه كه چه خواهند نمود لذا در جميع بيان بوصاياى لايحصى كلّ را وصيّت
فرموده‌اند و بشأنى تأكيد فرموده‌اند كه فوق آن ممكن نه چنانچه از قلم اعلى
نازل كه جميع بكمال تقديس و تنزيه خود را مقدس و منزّه نمايند و هم‌چنين لباس
و اسباب بيت را كه مباد لحاظ اللّه در حين ظهور بر شیء غير محبوب افتد و
بهمان قدر سبب حزن قلب اطهر گردد با اين وصايا شأن اهل بيان را ملاحظه
نمائيد كه مخصوص رسائل رديّه نوشته و بمقرّ عرش فرستاده‌اند اين بيان اولو
الأفئدة و اولو العزّة و اولو الوفا را كافى است

\end{otherlanguage}

Ay malaʾ-i Bayān, iʿtirāḍāt-i īn mardūd kih bih-ilqā-yi Mīrzā Yaḥyā va
Siyyid Muḥammad nivishtih makhṣūṣ dhikr shud tā bidānīd kih chih miqdār
ḍurr vārid shudih va idrāk namāyīd nawḥih va nudbih-yi Nuqṭih-yi Ūlā rā
bar Ẓuhūr-i Ākhar. Va ān Jamāl-i Mubārak chūn muṭṭaliʿ būdih-and bar
nufūs-i mushrikih kih chih khvāhand nimūd, lidhā dar jamīʿ-i Bayān
bih-vaṣāyā-yi lā yuḥṣā kull rā vaṣiyyat farmūdih-and va bih-shaʾnī
taʾkīd farmūdih-and kih fawq-i ān mumkin nah. Chunānchih az Qalam-i Aʿlā
nāzil kih jamīʿ bih-kamāl-i taqdīs va tanzīh khud rā muqaddas va
munazzah namāyand va hamchunīn libās va asbāb-i bayt rā, kih mabād liḥāẓ
Allāh dar ḥīn-i ẓuhūr bar shayʾ-i ghayr-i maḥbūb uftad va bih-hamān qadr
sabab-i ḥuzn-i qalb-i aṭhar gardad. Bā īn vaṣāyā shaʾn-i ahl-i Bayān rā
mulāḥiẓih namāyīd kih makhṣūṣ rasāʾil-i raddiyyih nivishtih va
bih-maqarr-i ʿarsh firistādih-and. Īn bayān ulū al-afʾidih va ulū
al-ʿizzih va ulū al-wafā rā kāfī ast.

O people of the Bayān! The objections that this rejected one wrote at
the prompting of Mīrzā Yaḥyā and Siyyid Muḥammad have been specifically
mentioned so that you may know how much harm has been inflicted and
understand the Primal Point's lamentation and mourning over the later
Revelation. Since that blessed Beauty knew what the idolatrous souls
would do, He admonished everyone throughout the Bayān with countless
counsels, emphasizing them to a degree beyond which nothing is possible.
Thus it was revealed by the Most Exalted Pen that all should sanctify
and purify themselves with the utmost holiness and purity, and likewise
their clothing and household furnishings, lest at the time of the
Revelation God's gaze fall upon anything displeasing and thereby cause
even that much sorrow to His most pure heart. In light of these
counsels, consider the condition of the people of the Bayān: they have
deliberately written treatises of refutation and sent them to the seat
of the throne. This statement suffices for those possessed of
understanding, honor, and fidelity.

\subsection{749}\label{section-1244}

\begin{otherlanguage}{persian}

فو‌اللّه وارد شده آنچه بر احدى از مظاهر قبل وارد نشده و با جميع بلايا و
رزايا و محن كه وارد شده شاكر و صابريم و ابدا حزنى نه و لكن از تضييع امر
اللّه قلب در احتراق است قسم بحقّ كه بر جميع لازم است نوحه و ندبه بر امر
الهى چه كه بشانى امر اللّه را تضييع نموده‌اند كه فاسقه ارض را خير النّسا
ناميده اند ملاحظه كنيد كه چگونه بطلان اين قوم را حقّ ظاهر فرمود فو‌اللّه
بهمين فعل بر كلّ حجّت تمام شد چه كه اكثرى آن فاسقه را ميشناسند و بر احوالش
مطّلعند و از آن خير النّسا خير الرّجال را عارف شويد يعنى آن نفسى كه او را
سيّده نساء نوشته فو‌اللّه لعنه كلّ الذّرّات و لكن هو ما استشعر فى نفسه و يكون
من الفرحين بی‌‌شرميش بمقامى رسيده كه زوجات خود را خير النّسا ناميده و حرم
نقطه اولى را تصرّف نموده ‌و بخشيده فويل لهم من عذاب يوم عظيم

\end{otherlanguage}

Fa-wa-llāh vārid shudih ānchih bar aḥadī az maẓāhir-i qabl vārid
nashudih va bā jamīʿ-i balāyā va razāyā va miḥan kih vārid shudih shākir
va ṣābirīm va abadan ḥuznī nah va lākin az taḍyīʿ-i Amr Allāh qalb dar
iḥtirāq ast. Qasam bi-Ḥaqq kih bar jamīʿ lāzim ast nawḥih va nudbih bar
Amr-i Ilāhī chih kih bi-shaʾnī Amr Allāh rā taḍyīʿ namūdih-and kih
fāsiqiy-i arḍ rā khayr al-nisāʾ nāmīdih-and. Mulāḥiẓih kunīd kih
chigūnih buṭlān-i īn qawm rā Ḥaqq ẓāhir farmūd. Fa-wa-llāh bi-hamīn fiʿl
bar kull ḥujjat tamām shud chih kih aktharī ān fāsiqih rā mīshināsand va
bar aḥvālash muṭṭaliʿand va az ān khayr al-nisāʾ khayr al-rijāl rā ʿārif
shavīd yaʿnī ān nafsī kih ū rā sayyidiy-i nisāʾ nivishtih. Fa-wa-llāh
laʿanahu kull al-dharrāt wa lākin huwa mā istashʿara fī nafsihi wa
yakūnu min al-fariḥīn. Bī-sharmīsh bi-maqāmī rasīdih kih zawjāt-i khud
rā khayr al-nisāʾ nāmīdih va ḥaram-i Nuqṭiy-i Ūlā rā taṣarruf namūdih va
bakhshīdih. Fa-waylun lahum min ʿadhābi yawmin ʿaẓīm.

By God, what has befallen Us has befallen none of the Manifestations of
the past. Yet despite all the afflictions, calamities, and ordeals that
have befallen Us, We remain thankful and patient, with no sorrow
whatsoever. But Our heart burns over the dishonoring of the Cause of
God. By the True One, everyone must lament and mourn for the divine
Cause, for they have so dishonored the Cause of God that they have
called a depraved woman of the earth ``the best of women.'' Consider how
the True One has exposed the falsity of these people. By God, through
this act alone the proof against all has been made complete, for most
people know that depraved woman and are acquainted with her
circumstances. And through that ``best of women,'' recognize the ``best
of men''---namely, the person who has described her as the foremost of
women. By God, every atom has cursed him, yet he remains unaware within
himself and is among those who rejoice. His shamelessness has reached
such a point that he has called his own wives ``the best of women,'' and
has taken possession of the wife of the Primal Point and given her away.
Woe to them, then, for the punishment of a tremendous day!

\subsection{750}\label{section-1245}

\begin{otherlanguage}{arabic}

در اينمقام اين آيات ابدع احلى من لسان اللّه العلى الأبهی نازل هذا ما نزل
حينئذ من جبروت عزّ عليّا يا قوم فاعلموا بانّا اصطفينا امّ نقطة الأولی و انّها
قد كانت من خيرة الاماء لدى العرش مذكوراً و حرّم اطلاق هذا الاسم على غيرها
كذلك رقم من قلم الأعلی فى لوح القضا الّذى كان فى كنائز عصمة ربّك محفوظا و
انّها لخير النّسا و بعدها تطلق على ضلع النّقطة الّتى ما خرجت عن حصن العصمة و
ما مسّتها ايدى الخائنين و كذلك كان الأمر مقضيّا و الّتى خانت انّها خرجت عن
النّقطة و قطعت نسبتها من اللّه الّذى قدّر كلّ امر فى لوح مبينا قل انّ المشركين
ظنّوا بانّا اردنا ان ننسخ ما نزل على نقطة البيان

\end{otherlanguage}

Dar īn maqām īn āyāt-i abdaʿ-i aḥlā min lisān Allāh al-ʿAlī al-Abhā
nāzil: Hādhā mā nazala ḥīnaʾidhin min jabarūti ʿizzin ʿaliyyā. Yā qawm
fa-iʿlamū bi-annā iṣṭafaynā umma Nuqṭati al-Ūlā wa innahā qad kānat min
khīrati al-imāʾ ladā al-ʿarsh madhkūrā. Wa ḥurrima iṭlāqu hādhā al-ism
ʿalā ghayrihā. Kadhālika ruqima min Qalami al-Aʿlā fī Lawḥi al-Qaḍā
alladhī kāna fī kanāʾizi ʿiṣmati rabbika maḥfūẓā. Wa innahā la-khayru
al-nisāʾ wa baʿdahā tuṭlaqu ʿalā ḍilʿi al-Nuqṭah allatī mā kharajat ʿan
ḥiṣni al-ʿiṣmah wa mā massathā aydī al-khāʾinīn. Wa kadhālika kāna
al-amru maqḍiyyā. Wa allatī khānat innahā kharajat ʿan al-Nuqṭah wa
qaṭaʿat nisbatāhā min Allāh alladhī qaddara kulla amrin fī lawḥin
mubīnā. Qul inna al-mushrikīn ẓannū bi-annā aradnā an nansakha mā nazala
ʿalā Nuqṭati al-Bayān.

At this point these wondrous, most sweet verses were revealed from the
tongue of God, the Most High, the All-Glorious: ``This is what has now
descended from the dominion of lofty glory. O people, know that We have
chosen the mother of the Primal Point, and she has been recorded before
the Throne as among the choicest of handmaidens. Applying this name to
anyone other than her has been forbidden. Thus has it been inscribed by
the Most Exalted Pen in the Tablet of Decree, preserved within the
treasuries of your Lord's unfailing protection. She is indeed the best
of women, and after her this designation applies to the wife of the
Point who has not left the stronghold of purity and whom the hands of
the treacherous have not touched. Thus has the matter been decreed. As
for the one who acted treacherously, she has departed from the Point and
severed her relationship with God, who has ordained all things in a
clear Tablet. Say: The idolaters have imagined that We intended to
abrogate what was revealed to the Point of the Bayān.''

\subsection{751}\label{section-1246}

\begin{otherlanguage}{arabic}

قل فو‌ربّی‌ الرّحمن لو نريد كما ظنّوا ليس لأحد ان يعترض على اللّه الّذى خلق كلّ
شیء بأمر من عنده و انّه كان على كلّ شیء قديرا و كلّما نريد هو ما اراد اللّه
و ما يظهر منّى هو ما ظهر منه و يشهد بذلك كلّ موحّد عليماً و لكن اللّه اراد
بهذا الظّهور ان يثبت ما نزل من عند نقطة البيان و يضع احكامه على اعناق
الفراعنة من ملأ الطّغيان و كان نفسه الحقّ على ذلك شهيدا و انّا اردنا ان
نفدى نفسنا فى سبيله كما فدى نفسه فى سبيلى اتّقوا اللّه يا قوم و لاتفتروا
على الّذى به اشرقت شمس البيان و ظهر حكمه بين الخلايق جميعاً فسوف نثبت
احكامه و نبرهن آثاره فى الأرض بقدرة و سلطان مبيناً

\end{otherlanguage}

Qul fa-wa-rabbī al-Raḥmān law nurīdu kamā ẓannū laysa li-aḥadin an
yaʿtariḍa ʿalā Allāh alladhī khalaqa kulla shayʾin bi-amrin min ʿindihi
wa innahu kāna ʿalā kulli shayʾin qadīrā. Wa kullamā nurīdu huwa mā
arāda Allāh wa mā yaẓharu minnī huwa mā ẓahara minhu wa yashhadu
bi-dhālika kullu muwaḥḥidin ʿalīmā. Wa lākin Allāh arāda bi-hādhā
al-ẓuhūr an yuthbita mā nazala min ʿindi Nuqṭati al-Bayān wa yaḍaʿa
aḥkāmahu ʿalā aʿnāqi al-farāʿinah min malaʾi al-ṭughyān wa kāna nafsuhu
al-Ḥaqq ʿalā dhālika shahīdā. Wa innā aradnā an nafdī nafsanā fī
sabīlihi kamā fadā nafsahu fī sabīlī. Ittaqū Allāh yā qawm wa lā taftarū
ʿalā alladhī bihi ashraqat shamsu al-Bayān wa ẓahara ḥukmuhu bayna
al-khalāʾiqi jamīʿā. Fa-sawfa nuthbitu aḥkāmahu wa nubarhinu āthārahu fī
al-arḍ bi-qudratin wa sulṭānin mubīnā.

Say: By my Lord, the All-Merciful, if We wished to do as they have
imagined, no one would have the right to object to God, who created all
things through a command from Himself. He is indeed powerful over all
things. Whatever We wish is what God has wished, and whatever appears
from Me is what has appeared from Him. Every knowledgeable believer in
His oneness bears witness to this. But through this Revelation God has
intended to confirm what was revealed by the Point of the Bayān and to
place His laws upon the necks of the Pharaohs among the hosts of
tyranny. His own Self, the True One, bears witness to this. We have
intended to sacrifice Ourself in His path, just as He sacrificed Himself
in My path. Fear God, O people, and do not fabricate lies against Him
through whom the Sun of the Bayān has shone and His decree has appeared
among all created beings. We will assuredly establish His laws and
demonstrate the truth of His writings on earth with power and manifest
authority.

\subsection{752}\label{section-1247}

\begin{otherlanguage}{persian}

اليوم خير النّسا امّ نقطه اولى است و حرام شده اطلاق اين اسم بر اماء اللّه و
هم‌چنين اذن داده شده اطلاق اين اسم بعدها بر ضلع نقطه كه در فارس ساكنند و
از حصن عصمت ربّانى خارج نشده‌اند و حرمة اللّه را رعايت نموده‌ و دست خائنين
بذيل عصمتش نرسيده فو‌نفسى الحقّ او است بتول عذرا و حرم اعلى و عصمت كبرى ان
اخدموها يا احبّائى بصدق مبين لانّها بقيّة اللّه بينكم ان انتم من العارفين
نسئل اللّه بأن يستقيمها على ما هى عليه و يعصمها عن مسّ الهاتكين الّذين
تحرّكهم ارياح النّفس و الهوى كيف یشآء اولئك هم فى ضلال مبين و بأن يثبتها
على هذا الأمر و انّه على كلّ شیء قدير

\end{otherlanguage}

Al-yawm khayr al-nisāʾ umm-i Nuqṭiy-i Ūlā ast va ḥarām shudih iṭlāq-i īn
ism bar imāʾ Allāh va hamchunīn idhn dādih shudih iṭlāq-i īn ism baʿdahā
bar ḍilʿ-i Nuqṭih kih dar Fārs sākinand va az ḥiṣn-i ʿiṣmat-i rabbānī
khārij nashudih-and va ḥurmat Allāh rā riʿāyat namūdih va dast-i
khāʾinīn bi-dhayl-i ʿiṣmatash narisīdih. Fa-wa-nafsī al-Ḥaqq ū ast
Batūl-i ʿAdhrā va Ḥaram-i Aʿlā va ʿIṣmat-i Kubrā. An ikhdumūhā yā
aḥibbāʾī bi-ṣidqin mubīn li-annahā baqiyyatu Allāh baynakum in antum min
al-ʿārifīn. Nasʾalu Allāh bi-an yastaqīmahā ʿalā mā hiya ʿalayhi wa
yaʿṣimahā ʿan massi al-hātikīn alladhīna tuḥarrikuhum aryāḥu al-nafs wa
al-hawā kayfa yashāʾ. Ulāʾika hum fī ḍalālin mubīn. Wa bi-an yuthbitahā
ʿalā hādhā al-amr wa innahu ʿalā kulli shayʾin qadīr.

Today the best of women is the mother of the Primal Point, and applying
this name to the handmaidens of God has been forbidden. Permission has
likewise been given to apply this name after her to the wife of the
Point who resides in Fārs, who has not left the stronghold of divine
purity, who has observed the sanctity of God, and whose robe of purity
the hands of the treacherous have not reached. By My own Self, the True
One, she is the chaste Virgin, the Most Exalted Consort, and the Most
Great Purity. Serve her, O My loved ones, with manifest sincerity, for
she is what remains of God among you, if you are among those who
understand. We ask God to keep her steadfast in her present state and to
protect her from the touch of those who violate sanctity, whom the winds
of self and desire move as they please. They are indeed in manifest
error. And We ask Him to establish her firmly in this Cause. He is
powerful over all things.

\subsection{753}\label{section-1248}

\begin{otherlanguage}{persian}

يا قوم انّى ما اردت فيما اقول نفسى بل نفس اللّه ان انتم من المنصفين و نلقى
عليكم ما اراد ربّكم العلىّ العظيم حرمت حرم اللّه بر كلّ لازم در جميع احوال
رعايتشان محبوب بوده و هم‌چنين نفوسيكه نفس الرّحمن از مابين آن نفوس ظاهر
شده عزّزوهم و وقّروهم يا قوم و لاتكوننّ من الغافلين و جميع اين فضل محقّق
مادام كه حبل نسبت منقطع نشود

\end{otherlanguage}

Yā qawm innī mā aradtu fīmā aqūlu nafsī bal nafsa Allāh in antum min
al-munṣifīn wa nulqī ʿalaykum mā arāda rabbukum al-ʿAlī al-ʿAẓīm.
Ḥurmat-i ḥaram-i Allāh bar kull lāzim dar jamīʿ-i aḥvāl riʿāyatishān
maḥbūb būdih va hamchunīn nufūsī kih Nafs-i Raḥmān az mābayn-i ān nufūs
ẓāhir shudih. ʿAzzizūhum wa waqqirūhum yā qawm wa lā takūnanna min
al-ghāfilīn. Va jamīʿ-i īn faḍl muḥaqqaq mādām kih ḥabl-i nisbat
munqaṭiʿ nashavad.

O people, in what I say I have not intended My own self, but the Self of
God, if you judge fairly. We convey to you what your Lord, the Most
High, the Almighty, has willed. Respect for the family of God is
incumbent upon all; showing them consideration in every circumstance has
been deemed desirable, as has showing consideration to those from whose
midst the Person of the All-Merciful has appeared. Honor and revere
them, O people, and do not be among the heedless. All this bounty holds
true so long as the bond of relationship is not severed.

\subsection{754}\label{section-1249}

\begin{otherlanguage}{arabic}

و انت تعلم يا الهى بانّى احبّ هؤلآء و كلّ ما نسب الى نفسك وَ اَجد منهم نفحات
قميص رحمانيّتك وفوحات ثوب عزّك و مواهبك و اشتاق لقائهم و الاجتماع معهم و
المؤانسة بهم اذا اسئلك يا الهى باسم الّذى جعلته سلطان الأسماء بين الأرض و
السّمآء و به كسّرتَ اصنام النّفس و الهوى و به غسّلت العارفين و المشتاقين و
طهّرتهم عن روايح ما سواك فى ملكوت الانشآء بأن تشتعل فى صدور هؤلآء مشاعل
حبّك ليهتدى بها اهل مملكتك الّذينهم غفلوا عن ذكرك و ثنائك و بعدوا عن شاطى
قربك و لقائك

\end{otherlanguage}

Wa anta taʿlamu yā ilāhī bi-annī uḥibbu hāʾulāʾ wa kulla mā nusiba ilā
nafsika wa ajidu minhum nafaḥāti qamīṣi raḥmāniyyatika wa fawaḥāti
thawbi ʿizzika wa mawāhibika wa ashtāqu liqāʾahum wa al-ijtimāʿa maʿahum
wa al-muʾānasata bihim. Idhan asʾaluka yā ilāhī bi-ismi alladhī
jaʿaltahu sulṭāna al-asmāʾ bayna al-arḍ wa al-samāʾ wa bihi kassarta
aṣnāma al-nafs wa al-hawā wa bihi ghassalta al-ʿārifīn wa al-mushtāqīn
wa ṭahhartahum ʿan rawāʾiḥi mā siwāka fī malakūti al-inshāʾ bi-an
tashtaʿila fī ṣudūri hāʾulāʾ mashāʿilu ḥubbika li-yahtadiya bihā ahlu
mamlakatika alladhīnahum ghafalū ʿan dhikrika wa thanāʾika wa baʿudū ʿan
shāṭiʾi qurbika wa liqāʾika.

You know, O my God, that I love these souls and all that is related to
Your own Self. From them I perceive the fragrances of the shirt of Your
all-mercifulness and the sweet scents of the robe of Your glory and Your
gifts, and I long to meet them, to gather with them, and to enjoy their
companionship. I therefore ask You, O my God, by the Name that You have
made sovereign over all names between earth and heaven, through which
You have shattered the idols of self and desire, and through which You
have washed those who know You and long for You, purifying them from the
scents of all besides You in the kingdom of creation, to kindle the
torches of Your love in the hearts of these souls, so that through them
the inhabitants of Your realm may be guided---those who have neglected
Your remembrance and praise and have strayed far from the shore of
nearness to You and meeting with You.

\subsection{755}\label{section-1250}

\begin{otherlanguage}{arabic}

فيا الهى لاتمنعهم عمّا ظهر بينهم و لاتحرمهم عمّا رشحت فى تلك الأیّام من
سحاب رحمتك و غمام فضلك ثمّ اسمعهم يا الهى فى كلّ حين نغماتك و الحانك الّتى
ارتفعت عن مقرّ عرش وحدانيّتك و كرسىّ عزّ فردانيّتك ثمّ اقبل منهم يا الهى ما
ترك عنهم فى ایّامك الّتى فيها اشرقت شمس جمالك عن افق مشيّتك و كان ان يمشى
بينهم سلطان الرّسل باقتدارك و عظمتك و انت ترى يا الهى بانّ البكا منعنى
حينئذ عن ذكرك و ثنائك و ثنآء الّذين نسبتهم الى مظهر نفسك بما خطر ببالى
ذكر محبوبی‌ و مقصودى و رجائى و ذكر ایّامه و شئوناته

\end{otherlanguage}

Fa-yā ilāhī lā tamnaʿhum ʿammā ẓahara baynahum wa lā taḥrimhum ʿammā
rashaḥta fī tilka al-ayyām min saḥābi raḥmatika wa ghamāmi faḍlika.
Thumma asmiʿhum yā ilāhī fī kulli ḥīnin naghamātika wa alḥānaka allatī
irtafaʿat ʿan maqarri ʿarshi waḥdāniyyatika wa kursiyyi ʿizzi
fardāniyyatika. Thumma iqbal minhum yā ilāhī mā turika ʿanhum fī
ayyāmika allatī fīhā ashraqat shamsu jamālika ʿan ufuqi mashīʾatika wa
kāna an yamshiya baynahum Sulṭānu al-Rusul bi-iqtidārika wa ʿaẓamatika.
Wa anta tarā yā ilāhī bi-anna al-bukāʾ manaʿanī ḥīnaʾidhin ʿan dhikrika
wa thanāʾika wa thanāʾi alladhīna nasabtahum ilā maẓhari nafsika bi-mā
khaṭara bi-bālī dhikru maḥbūbī wa maqṣūdī wa rajāʾī wa dhikru ayyāmihi
wa shuʾūnātihi.

O my God, do not keep them from what has appeared among them, and do not
deprive them of what You have showered in these days from the clouds of
Your mercy and Your grace. Let them hear, O my God, at every moment,
Your melodies and songs that have risen from the seat of the throne of
Your oneness and the seat of the glory of Your singleness. Accept from
them, O my God, what they have left behind in Your days, in which the
Sun of Your beauty has shone from the horizon of Your will and the
Sovereign of the Messengers has walked among them through Your might and
grandeur. You see, O my God, that weeping has now prevented me from
remembering and praising You and from praising those whom You have
related to the Manifestation of Your own Self, for there has come to my
mind the remembrance of my Beloved, my Desire, and my Hope, and the
remembrance of His days and His ways.

\subsection{756}\label{section-1251}

\begin{otherlanguage}{arabic}

فو‌عزّتك يا الهى بنفسى ضيع اسم الوفآء بين الأرض و السّمآء لانّ محبوبی‌ فدى
نفسه فى سبيلك و انّى اكون موجوداً فى مملكتك بين بريّتك مع انّك انت تعلم يا
الهى بانّى فى كلّ حين احبّ ان انفق روحى و نفسى فى سبيله و فى كلّ آن انتظر
هذا و لكن لم ادر باىّ جهة منعت عمّا اردت فو‌عزّتك اذاً انقطع الصّبر منّى حبّاً
للقائك و طلباً لوصالك و لاتفرح نفسى الّا حين الّذى تشاهد وجهى محمرّاً بدمى و
كلّما القيت نفسى بين يدى الأعدا انّك حفظتنى بسلطان مشيّتك و كلّما اودعت جسدى
تحت ايادى اولى البغضآء انت عصمتنى بقدرتك و قوّتك

\end{otherlanguage}

Fa-wa-ʿizzatika yā ilāhī bi-nafsī ḍuyyiʿa ismu al-wafāʾ bayna al-arḍ wa
al-samāʾ li-anna maḥbūbī fadā nafsahu fī sabīlika wa innī akūnu mawjūdan
fī mamlakatika bayna bariyyatika maʿa annaka anta taʿlamu yā ilāhī
bi-annī fī kulli ḥīnin uḥibbu an unfiqa rūḥī wa nafsī fī sabīlihi wa fī
kulli ānin antaẓiru hādhā wa lākin lam adri bi-ayyi jihatin muniʿtu
ʿammā aradtu. Fa-wa-ʿizzatika idhan inqaṭaʿa al-ṣabru minnī ḥubban
li-liqāʾika wa ṭalaban li-wiṣālika wa lā tafraḥu nafsī illā ḥīna alladhī
tushāhidu wajhī muḥmarran bi-damī. Wa kullamā alqaytu nafsī bayna yaday
al-aʿdāʾ innaka ḥafiẓtanī bi-sulṭāni mashīʾatika wa kullamā awdaʿtu
jasadī taḥta ayādī ulī al-baghḍāʾ anta ʿaṣamtanī bi-qudratika wa
quwwatika.

By Your glory, O my God, through me the very name of faithfulness has
been lost between earth and heaven, for my Beloved sacrificed Himself in
Your path, while I remain alive in Your realm among Your creatures. Yet
You know, O my God, that at every moment I long to offer up my spirit
and my life in His path, and at every instant I await this. But I do not
know why I have been prevented from attaining what I desired. By Your
glory, my patience has now come to an end through love of meeting You
and longing for union with You. My soul will rejoice only when You
behold my face reddened with my blood. Whenever I have cast myself into
the hands of my enemies, You have preserved me through the sovereignty
of Your will; and whenever I have placed my body beneath the hands of
those possessed of hatred, You have protected me through Your power and
might.

\subsection{757}\label{section-1252}

\begin{otherlanguage}{persian}

فلمّا جعلتنى محروماً عمّا اردته فى سبيلك قرّبنى الى ارض الّتى استوى عليها
محبوبی‌ و محبوب العالمين ومقصودى و مقصود العارفين لاستأنس بها و اقبّل
ترابها و اجد منها روايح قدسك و نفحات رحمتك لعلّ بذلك تسكن نفسى من
اضطرابها و روحى من احتراقها و جسدى من اهتزازه و استغفرك يا الهى عن كلّ
ذلك لانّه حینئذ يكون حيّاً فى الأفق الأبهی يشهد و يرى و بذلك تشهد نفسى و
روحى و ذاتى و قلمى و لسانى و ظاهرى و باطنى و اوّلى و آخرى لا اله الّا انت
العزيز المتعالى المقتدر المهيمن العليم الحكيم

\end{otherlanguage}

Fa-lammā jaʿaltanī maḥrūman ʿammā aradtuhu fī sabīlika qarribnī ilā arḍi
allatī istawā ʿalayhā maḥbūbī wa maḥbūbu al-ʿālamīn wa maqṣūdī wa
maqṣūdu al-ʿārifīn li-astaʾnisa bihā wa uqabbila turābahā wa ajida minhā
rawāʾiḥa qudsika wa nafaḥāti raḥmatika laʿalla bi-dhālika taskunu nafsī
min iḍṭirābihā wa rūḥī min iḥtirāqihā wa jasadī min ihtizāzihi. Wa
astaghfiruka yā ilāhī ʿan kulli dhālika li-annahu ḥīnaʾidhin yakūnu
ḥayyan fī al-Ufuqi al-Abhā yashhadu wa yarā wa bi-dhālika tashhadu nafsī
wa rūḥī wa dhātī wa qalamī wa lisānī wa ẓāhirī wa bāṭinī wa awwalī wa
ākhirī. Lā ilāha illā anta al-ʿAzīz al-Mutaʿālī al-Muqtadir al-Muhaymin
al-ʿAlīm al-Ḥakīm.

Since You have deprived me of what I desired in Your path, bring me near
to the land upon which my Beloved and the Beloved of all worlds, my
Desire and the Desire of those who know You, once dwelt, so that I may
find solace in it, kiss its dust, and perceive from it the fragrances of
Your holiness and the breezes of Your mercy. Perhaps through this my
soul may be stilled from its agitation, my spirit from its burning, and
my body from its trembling. I ask Your forgiveness, O my God, for all
this, for He is even now alive in the All-Glorious Horizon, witnessing
and seeing. To this bear witness my soul, my spirit, my being, my pen,
my tongue, my outer and inner self, my beginning and my end. There is no
God but You, the Almighty, the Most Exalted, the All-Powerful, the
All-Protecting, the All-Knowing, the All-Wise.

\subsection{758}\label{section-1253}

\begin{otherlanguage}{persian}

بدانيد ايقوم كه اين غلام رحمن در جميع احيان اراده‌اش آن بوده كه آنچه از
ظهور قبلش نقطه بيان روح من فى الأكوان فداه نازل شده ثابت نمايد حركت
ننموده‌ام مگر برضاى او چنانچه او حركت نفرمود مگر برضاى اين غلام مع آنكه
كلّ آنچه از اين مشيّت ظاهر شود نفس مشيّت و اراده او بوده و كفى بنفسه شهيداً
و لكن مشركين بيان گمان نموده‌اند و در باديه‌هاى مهلكه ظنّ مشى کرده‌اند
چنانچه نسبت داده‌اند كه اين غلام گفته كه نقطه بيان از يكحرف من خلق شده لا
و اللّه آنچه نازل شده منزل حقّ تعالى سلطنته و تعالى قدرته و تعالى عظمته و
تعالى اقتداره و تعالى عزّه و تعالى رفعته و تعالى شأنه و تعالى بهائه و
تعالى امره بوده و خواهد بود

\end{otherlanguage}

Bidānīd ay qawm kih īn Ghulām-i Raḥmān dar jamīʿ-i aḥyān irādih-ash ān
būdih kih ānchih az ẓuhūr-i qablash Nuqṭiy-i Bayān rūḥu man fī al-akwān
fidāhu nāzil shudih thābit namāyad. Ḥarakat nanamūdih-am magar
bi-riḍā-yi ū chunānchih ū ḥarakat nafarmūd magar bi-riḍā-yi īn Ghulām
maʿa ānkih kull-i ānchih az īn Mashiyyat ẓāhir shavad nafs-i mashiyyat
va irādiy-i ū būdih wa kafā bi-nafsihi shahīdā. Va lākin mushrikīn-i
Bayān gumān namūdih-and va dar bādiyih-hā-yi muhlikiy-i ẓann mashy
kardih-and chunānchih nisbat dādih-and kih īn Ghulām guftih kih Nuqṭiy-i
Bayān az yik ḥarf-i man khalq shudih. Lā wa-llāh ānchih nāzil shudih
munzil Ḥaqq taʿālā salṭanatuhu wa taʿālā qudratuhu wa taʿālā ʿaẓamatuhu
wa taʿālā iqtidāruhu wa taʿālā ʿizzuhu wa taʿālā rifʿatuhu wa taʿālā
shaʾnuhu wa taʿālā bahāʾuhu wa taʿālā amruhu būdih va khwāhad būd.

Know, O people, that this Youth of the All-Merciful has at all times
intended to establish whatever was revealed by His previous
Manifestation, the Point of the Bayān---may the souls of all in
existence be sacrificed for Him. I have acted only in accordance with
His good pleasure, just as He acted only in accordance with the good
pleasure of this Youth. Moreover, whatever appears from this Will has
been His very will and purpose, and His own Self is sufficient as a
witness. But the idolaters of the Bayān have indulged in conjecture and
wandered through the deadly wildernesses of supposition, even alleging
that this Youth has said that the Point of the Bayān was created from a
single letter of Mine. No, by God! Whatever has been revealed, its
Revealer has been and will be the True One---exalted be His sovereignty,
exalted be His power, exalted be His grandeur, exalted be His might,
exalted be His glory, exalted be His loftiness, exalted be His station,
exalted be His splendor, and exalted be His Cause.

\subsection{759}\label{section-1254}

\begin{otherlanguage}{persian}

بگوئيد اى بي‌شرمهاى ارض آنچه در اين ظهور نازل بعينه همان كلمات نقطه اوّليّه
بوده و خواهد بود و اين غلام لازال جز عبوديّت صرفه دوست نداشته فو‌الّذى نفسى
بيده عبوديّت محبوب جانم بوده بشأنيكه كلّما اشاهد موقفاًمن مواقف الأرض احبّ
ان اخرّ عليه سجّدا للّه محبوبی‌ و محبوب العارفين فيا ليت على كلّ ذرّة من
ترابها سجدت للّه ربّی‌ و ربّ العالمين و اينكه در الواح ذكر مقامات عاليه شده
ناظراً الى امر اللّه و شأنه و عزّه و اجلاله بوده چنانچه نقطهء بيان فرموده
‌انّه ينطق فى كلّ شیء بانّى انا اللّه لا اله الّا انا ان يا خلقى ايّاى فاعبدون
و اگر نظر باينمقام نبود فو‌نفسه المحبوب ما ذكرت الّا العبوديّة الصّرفة للّه
الحقّ كذلك كان الأمر و لكن النّاس هم لا يشعرون

\end{otherlanguage}

Bigūyīd ay bī-sharmhā-yi arḍ ānchih dar īn ẓuhūr nāzil bi-ʿaynihi hamān
kalimāt-i Nuqṭiy-i Avvaliyyih būdih va khwāhad būd va īn Ghulām lā yazāl
juz ʿubūdiyyat-i ṣirfih dūst nadāshtih. Fa-wa-alladhī nafsī bi-yadihi
ʿubūdiyyat maḥbūb-i jānam būdih bi-shaʾnī kih kullamā ushāhidu mawqifan
min mawāqifi al-arḍ uḥibbu an akhirra ʿalayhi sujjadan li-llāh maḥbūbī
wa maḥbūbi al-ʿārifīn. Fa-yā layta ʿalā kulli dharratin min turābihā
sajadtu li-llāh rabbī wa rabbi al-ʿālamīn. Va īnkih dar alvāḥ dhikr-i
maqāmāt-i ʿāliyih shudih nāẓiran ilā Amr Allāh wa shaʾnihi wa ʿizzihi wa
ijlālihi būdih chunānchih Nuqṭiy-i Bayān farmūdih innahu yanṭiqu fī
kulli shayʾin bi-annī anā Allāh lā ilāha illā anā an yā khalqī iyyāya
fa-uʿbudūn. Va agar naẓar bi-īn maqām nabūd fa-wa-nafsihi al-Maḥbūb mā
dhakartu illā al-ʿubūdiyyata al-ṣirfata li-llāh al-Ḥaqq. Kadhālika kāna
al-amr wa lākin al-nāsu hum lā yashʿurūn.

Say: O shameless ones of the earth! Whatever is revealed in this
Revelation has been and will be precisely the words of the Primal Point,
and this Youth has always loved nothing but pure servitude. By Him in
whose hand is my soul, servitude has been the beloved of my heart, to
such a degree that whenever I behold any spot on earth, I long to fall
prostrate upon it before God, my Beloved and the Beloved of those who
know Him. Would that upon every atom of its dust I had prostrated myself
before God, my Lord and the Lord of all worlds! The exalted stations
mentioned in the Tablets have been spoken of with regard to the Cause of
God, His station, His glory, and His majesty, just as the Point of the
Bayān has said: ``He speaks in all things, saying, `I am God; there is
no God but Me. O My creation, worship Me alone.'\,'' Were it not with
regard to this station, then by His own beloved Self, I would have
spoken of nothing but pure servitude to God, the True One. Such is the
matter, but the people do not perceive it.

\subsection{760}\label{section-1255}

\begin{otherlanguage}{persian}

در اينمقام مناجاتى كه مخصوص يكى از احبّا از سماء عزّ ابهى نازل ذكر ميشود
تا كلّ عارف شوند كه در جميع مراتب اكمل از كلّ مظاهر ظاهر شده‌اند حتّى در
مقام عبوديّت

\end{otherlanguage}

Dar īn maqām munājātī kih makhṣūṣ-i yikī az aḥibbā az samāʾ-i ʿizz-i
Abhā nāzil dhikr mīshavad tā kull ʿārif shavand kih dar jamīʿ-i marātib
akmal az kull-i maẓāhir ẓāhir shudih-and ḥattā dar maqām-i ʿubūdiyyat.

At this point a prayer revealed from the heaven of all-glorious majesty
especially for one of the loved ones is cited, so that all may recognize
that He has appeared in every degree more perfect than all the
Manifestations, even in the station of servitude.

\subsection{761}\label{section-1256}

\begin{otherlanguage}{arabic}

سبحانك اللّهمّ يا الهى تشهد السن الممكنات على سلطنتك و اقتدارك و على فقرى
و افتقارى عند ظهورات غنائك اذا يا الهى فانظر هذا العاصى الّذى لم يزل كان
ناظراً الى شطر غفرانك و قلبه متوجّهاً الى افق فضلك و مواهبك و انّى يا الهى
من اوّل يوم الّذى خلقتنى بامرك و احييتنى من نسماءت جود رحمانيّتك ما توجّهت
الى احد دونك و قمت الى مقابلة الأعدآء لسلطنتك و اقتدارك و دعوت الكلّ الى
شاطى بحر توحيدك و سماء عزّ تفريدك و ما اردت فى ایّامى حفظ نفسى من طغاة
خلقك بل اعلآء ذكرك بين بريّتك و بذلك ورد علىّ ما لا حمله احد من خلقك

\end{otherlanguage}

Subḥānaka Allāhumma yā ilāhī, tashhadu alsunu al-mumkināti ʿalā
salṭanatika wa-iqtidārika wa-ʿalā faqrī wa-iftiqārī ʿinda ẓuhūrāti
ghināʾika. Idhan yā ilāhī fa-nẓur hādhā al-ʿāṣī alladhī lam yazal kāna
nāẓiran ilā shaṭri ghufrānika wa-qalbuhu mutawajjihan ilā ufuqi faḍlika
wa-mawāhibika. Wa-innī yā ilāhī min awwali yawmin alladhī khalaqtanī
bi-amrika wa-aḥyaytanī min nasamāʾāti jūdi raḥmāniyyatika mā tawajjahtu
ilā aḥadin dūnaka wa-qumtu ilā muqābalati al-aʿdāʾi li-salṭanatika
wa-iqtidārika wa-daʿawtu al-kulla ilā shāṭiʾi baḥri tawḥīdika wa-samāʾi
ʿizzi tafrīdika. Wa-mā aradtu fī ayyāmī ḥifẓa nafsī min ṭughāti khalqika
bal iʿlāʾa dhikrika bayna bariyyatika wa-bi-dhālika warada ʿalayya mā lā
ḥamalahu aḥadun min khalqika.

Glory be to You, O God, my God! The tongues of all contingent beings
testify to Your sovereignty and power, and to my poverty and need before
the manifestations of Your wealth. Look, then, O my God, upon this
sinner, who has always gazed toward Your forgiveness, his heart turned
toward the horizon of Your grace and Your gifts. O my God, from the very
first day when You created me by Your command and brought me to life
through the breezes of Your all-merciful generosity, I have turned to no
one but You. I arose to confront the enemies for the sake of Your
sovereignty and power, and summoned everyone to the shore of the ocean
of Your oneness and the heaven of the glory of Your singleness.
Throughout my days, I have sought not to protect myself from the tyrants
among Your creatures, but to exalt Your remembrance among Your creation.
Because of this, there has befallen me what none of Your creatures has
borne.

\subsection{762}\label{section-1257}

\begin{otherlanguage}{arabic}

و كم من ایّام يا الهى كنت فريداً بين المذنبين من عبادك و كم من ليالى يا
محبوبی‌ كنت اسيراً بين الغافلين من خلقك و فى موارد البأسآء و الضّرّآء كنت
ناطقاً بثنآء نفسك بين سمائك و ارضك و ذاكراً ببدايع ذكرك فى ملكوت امرك و
خلقك و لو انّ كلّ ما ظهر منّى لاينبغى لسلطان عزّ وحدانيّتك و لايليق لشأنك و
اقتدارك فو‌عزّتك يا محبوبی‌ لم اجد لنفسى وجوداً تلقآء مدين عزّك و كلّما اريد
ان اثنى نفسك بثنآء يمنعنى فؤادى لانّ دونك لم يقدر ان يطير فى هواء ملكوت
قربك او ان يصعد الى سماء جبروت لقائك فو‌عزّتك اشاهد بانّى لو اسجد لكفّ من
التّراب الى آخر الّذى لا آخر له لنسبته الى اسمك الصّانع لا أجد نفسى بعيداً
عن التّقرّب اليه و اشاهد بأنّ عملى لا ينبغى له بل كان محدوداً بحدودات نفسى و
لو اخدم احداً من عبادك بحيث اقوم بين يديه بدوام ملكوتك و بقآء جبروتك
لنسبته الى اسمك الخالق فو‌عزّتك لأجد نفسى مقصّراً عن ادآء خدمته و محروماً عمّا
يليق له لانّ فى هذا المقام لا يرى الّا نسبتهم الى اسمائك و صفاتك

\end{otherlanguage}

Wa-kam min ayyāmin yā ilāhī kuntu farīdan bayna al-mudhnibīna min
ʿibādika wa-kam min layālin yā maḥbūbī kuntu asīran bayna al-ghāfilīna
min khalqika. Wa-fī mawāridi al-baʾsāʾi wa-l-ḍarrāʾi kuntu nāṭiqan
bi-thanāʾi nafsika bayna samāʾika wa-arḍika wa-dhākiran bi-badāʾiʿi
dhikrika fī malakūti amrika wa-khalqika. Wa-law anna kulla mā ẓahara
minnī lā yanbaghī li-sulṭāni ʿizzi waḥdāniyyatika wa-lā yalīqu
li-shaʾnika wa-iqtidārika, fa-wa-ʿizzatika yā maḥbūbī lam ajid li-nafsī
wujūdan tilqāʾa madīni ʿizzika. Wa-kullamā urīdu an uthniya nafsaka
bi-thanāʾin yamnaʿunī fuʾādī li-anna dūnaka lam yaqdir an yaṭīra fī
hawāʾi malakūti qurbika aw an yaṣʿada ilā samāʾi jabarūti liqāʾika.
Fa-wa-ʿizzatika ushāhidu bi-annī law asjudu li-kaffin mina al-turābi ilā
ākhiri alladhī lā ākhira lahu li-nisbatihi ilā ismika al-Ṣāniʿi lā ajidu
nafsī baʿīdan ʿani al-taqarrubi ilayhi wa-ushāhidu bi-anna ʿamalī lā
yanbaghī lahu bal kāna maḥdūdan bi-ḥudūdāti nafsī. Wa-law akhdumu aḥadan
min ʿibādika bi-ḥaythu aqūmu bayna yadayhi bi-dawāmi malakūtika
wa-baqāʾi jabarūtika li-nisbatihi ilā ismika al-Khāliqi fa-wa-ʿizzatika
la-ajidu nafsī muqaṣṣiran ʿan adāʾi khidmatihi wa-maḥrūman ʿammā yalīqu
lahu li-anna fī hādhā al-maqāmi lā yurā illā nisbatuhum ilā asmāʾika
wa-ṣifātika.

How many days, O my God, I was alone among the sinful among Your
servants, and how many nights, O my Beloved, I was a captive among the
heedless among Your creatures! Amid hardship and adversity, I spoke Your
praise between Your heaven and Your earth, and proclaimed the wonders of
Your remembrance in the kingdom of Your command and creation. Although
everything that has proceeded from me is unworthy of the sovereignty of
Your glorious oneness and unbefitting Your station and power, yet, by
Your glory, O my Beloved, I have found no existence of my own before the
city of Your glory. Whenever I wish to offer You praise, my heart
restrains me, for no one other than You can soar in the atmosphere of
the kingdom of Your nearness or ascend to the heaven of the dominion of
Your presence. By Your glory! I see that, were I to prostrate myself
before a handful of dust until the end that has no end, because of its
relationship to Your name, the Maker, I would not consider myself far
from drawing near to it. I see that my deed would be unworthy of it,
being confined instead within the limits of my own self. Were I to serve
one of Your servants, standing before him for as long as Your kingdom
endures and Your dominion remains, because of his relationship to Your
name, the Creator, then, by Your glory, I would find myself falling
short in serving him and deprived of what befits him, for in this
station nothing is seen but their relationship to Your names and
attributes.

\subsection{763}\label{section-1258}

\begin{otherlanguage}{arabic}

انّ الّذى كان شأنه ذلك كيف يقدر ان يذكر الّذى باشارة من اصبعه خلقت الأسمآء
و ملكوتها و الصّفات و جبروتها و باشارة اخرى ركّبت الكاف بالنّون و ظهر منها
ما عجز عن عرفانه اعلى افئدة المقرّبين من اصفيائك و ابهى مشاعر المخلصين من
اودّائك

\end{otherlanguage}

Inna alladhī kāna shaʾnuhu dhālika kayfa yaqdiru an yadhkura alladhī
bi-ishāratin min iṣbaʿihi khuliqati al-asmāʾu wa-malakūtuhā wa-l-ṣifātu
wa-jabarūtuhā wa-bi-ishāratin ukhrā rukkibati al-kāfu bi-l-nūni
wa-ẓahara minhā mā ʿajaza ʿan ʿirfānihi aʿlā afʾidati al-muqarrabīna min
aṣfiyāʾika wa-abhā mashāʿiri al-mukhliṣīna min awiddāʾika.

How can one whose condition is such make mention of Him at a gesture of
whose finger the names and their kingdom, and the attributes and their
dominion, were created, and at another gesture of whom the letter kāf
was joined to the letter nūn, and from this appeared that which the
loftiest hearts of those drawn near among Your chosen ones, and the most
radiant faculties of the sincere among Your loved ones, are powerless to
comprehend?

\subsection{764}\label{section-1259}

\begin{otherlanguage}{arabic}

فو‌عزّتك يا محبوبی‌ صرت متحيّراً فى مظاهر صنعك و مطالع قدرتك و اشاهد نفسى
عاجزاً عن عرفان ادنى آيتك و كيف عرفان نفسك اذاً اسئلك يا الهى باسمك الّذى
به طيّرت العاشقين فى هوآء ارادتك و اهتديت به المشتاقين الى رضوان قربك و
وصالك بأن تهبّ من رضوان عنايتك روايح الاطمينان على المضطرّين من احبّائك فى
تلك الأیّام الّتى احاطتهم ارياح الافتتان من كلّ الجهات بحيث اضطربت النّفوس
من سطوة قضائك و تزلزلت اركان الوجود عمّا نزل عليهم من سمآء تقديرك و بلغت
اضطرابهم الى مقام تكاد ان تخمد فى مشكوة قلوبهم سراج حبّك و ذكرك و انّك انت
المقتدر على ما تشآء و انّك انت الغفور الكريم

\end{otherlanguage}

Fa-wa-ʿizzatika yā maḥbūbī ṣirtu mutaḥayyiran fī maẓāhiri ṣunʿika
wa-maṭāliʿi qudratika wa-ushāhidu nafsī ʿājizan ʿan ʿirfāni adnā āyatika
wa-kayfa ʿirfāni nafsika. Idhan asʾaluka yā ilāhī bi-ismika alladhī bihi
ṭayyarta al-ʿāshiqīna fī hawāʾi irādatika wa-ihtadayta bihi
al-mushtāqīna ilā riḍwāni qurbika wa-wiṣālika bi-an tahubba min riḍwāni
ʿināyatika rawāʾiḥu al-iṭmiʾnāni ʿalā al-muḍṭarrīna min aḥibbāʾika fī
tilka al-ayyāmi allatī aḥāṭathum aryāḥu al-iftitāni min kulli al-jihāti
bi-ḥaythu iḍṭarabati al-nufūsu min saṭwati qaḍāʾika wa-tazalzalat arkānu
al-wujūdi ʿammā nazala ʿalayhim min samāʾi taqdīrika wa-balaghat
iḍṭirābuhum ilā maqāmin takādu an takhmuda fī mishkāti qulūbihim sirāju
ḥubbika wa-dhikrika. Wa-innaka anta al-Muqtadiru ʿalā mā tashāʾu
wa-innaka anta al-Ghafūru al-Karīmu.

By Your glory, O my Beloved! I am bewildered by the manifestations of
Your handiwork and the dawning places of Your power. I see myself
powerless to know even Your least sign; how much less can I know Your
own Self! I therefore ask You, O my God, by Your name through which You
have made the lovers soar in the atmosphere of Your will and guided
those who yearn to the paradise of Your nearness and reunion with You,
to let the fragrances of assurance waft from the paradise of Your care
over those of Your loved ones who are in distress, in these days when
the winds of trial have surrounded them from every direction. Souls have
been shaken by the force of Your decree, and the pillars of existence
have trembled at what has descended upon them from the heaven of Your
ordainment. Their agitation has reached such a degree that the lamp of
Your love and remembrance is almost extinguished in the niche of their
hearts. You indeed have power over whatever You will, and You are indeed
the Ever-Forgiving, the Generous.

\subsection{765}\label{section-1260}

\begin{otherlanguage}{persian}

فيا الهى و سيّدى تسمع ضجيج محبّيك و صريخهم من كلّ الأقطار بما ورد عليهم من
الّذين كانت قلوبهم محروماً عن نفحات حبّك و ليس لهم من معين ليعينهم و لا من
ناصر لينصرهم و كذلك ليس لاعدائهم من مانع ليمنعهم عن ضرّ هؤلآء لذا يفعلون
ما يريدون و يعملون ما يشاؤن اذا فانصر يا الهى ببديع نصرك احبّائك الّذين ما
استنصروا من غيرك و ما توجّهوا الى دونك و كانت عيونهم منتظرة لبدايع مواهبك
و الطافك ثمّ ارحمهم يا الهى ببدايع رحمتك ثمّ ادخلهم فى حصن حمايتك و عنايتك
و انّك انت الّذى يا الهى لم تزل كنت مأمن الخائفين و ملجأ المضطرّين

\end{otherlanguage}

Fa-yā ilāhī wa-sayyidī tasmaʿu ḍajīja muḥibbīka wa-ṣarīkhahum min kulli
al-aqṭāri bi-mā warada ʿalayhim mina alladhīna kānat qulūbuhum maḥrūman
ʿan nafaḥāti ḥubbika wa-laysa lahum min muʿīnin li-yuʿīnahum wa-lā min
nāṣirin li-yanṣurahum wa-kadhālika laysa li-aʿdāʾihim min māniʿin
li-yamnaʿahum ʿan ḍarri hāʾulāʾi li-dhā yafʿalūna mā yurīdūna
wa-yaʿmalūna mā yashāʾūna. Idhan fa-nṣur yā ilāhī bi-badīʿi naṣrika
aḥibbāʾaka alladhīna mā istanṣarū min ghayrika wa-mā tawajjahū ilā
dūnika wa-kānat ʿuyūnuhum muntaẓiratan li-badāʾiʿi mawāhibika
wa-alṭāfika. Thumma irḥamhum yā ilāhī bi-badāʾiʿi raḥmatika thumma
adkhilhum fī ḥiṣni ḥimāyatika wa-ʿināyatika wa-innaka anta alladhī yā
ilāhī lam tazal kunta maʾmana al-khāʾifīna wa-maljaʾa al-muḍṭarrīna.

O my God and my Lord! You hear the lamentation of those who love You and
their cries from every region over what has befallen them at the hands
of those whose hearts are deprived of the breaths of Your love. They
have no helper to assist them and no supporter to come to their aid. Nor
is there anyone to restrain their enemies from harming them; thus they
do whatever they desire and act as they please. Aid, then, O my God,
with Your wondrous assistance, Your loved ones who have sought help from
none but You and turned to no one besides You, whose eyes have awaited
the wonders of Your gifts and Your favors. Have mercy on them, O my God,
through the wonders of Your mercy, and bring them into the stronghold of
Your protection and care. You, O my God, have always been the refuge of
the fearful and the shelter of those in distress.

\subsection{766}\label{section-1261}

\begin{otherlanguage}{persian}

اسئلك بأن لاتحرم هؤلآء الضّعفآء عن بدايع جودك و افضالك و لاتدعهم بين ايدى
الّذين ما خلقت كينوناتهم الّا من نار غضبك و قهرك و ما وجدوا روايح الرّحم و
الانصاف و غرّتهم الدّنيا بغرورها الى شأن انكروا برهانك و اشركوا بنفسك و
كفروا بآياتك و سفكوا دم احبّائك و امنائك فو‌عزّتك يا محبوبی‌ ارتكبوا ما لم
يرتكبه احد من قبل و بذلك استحقّوا غضبك و سياط قهرك خذهم بسلطانك ثمّ سلّط
عليهم من لا يرحمهم الّا بأن يرجعوا اليك و يدخلوا فى ظلّ عنايتك و يتوبوا
اليك و انّك انت لم تزل كنت قادراً و لاتزال تكون مقتدراً و انّك انت المقتدر
المتعالى العادل الحكيم

\end{otherlanguage}

Asʾaluka bi-an lā taḥrima hāʾulāʾi al-ḍuʿafāʾa ʿan badāʾiʿi jūdika
wa-ifḍālika wa-lā tadaʿhum bayna aydī alladhīna mā khalaqta kaynūnātihim
illā min nāri ghaḍabika wa-qahrika wa-mā wajadū rawāʾiḥa al-raḥmi
wa-l-inṣāfi wa-gharrathumu al-dunyā bi-ghurūrihā ilā shaʾnin ankarū
burhānaka wa-ashrakū bi-nafsika wa-kafarū bi-āyātika wa-safakū dama
aḥibbāʾika wa-umanāʾika. Fa-wa-ʿizzatika yā maḥbūbī irtakabū mā lam
yartakibhu aḥadun min qablu wa-bi-dhālika istaḥaqqū ghaḍabaka wa-siyāṭa
qahrika. Khudhhum bi-sulṭānika thumma salliṭ ʿalayhim man lā yarḥamuhum
illā bi-an yarjiʿū ilayka wa-yadkhulū fī ẓilli ʿināyatika wa-yatūbū
ilayka. Wa-innaka anta lam tazal kunta qādiran wa-lā tazālu takūnu
muqtadiran wa-innaka anta al-Muqtadiru al-Mutaʿālī al-ʿĀdilu al-Ḥakīmu.

I ask You not to deprive these weak ones of the wonders of Your
generosity and bounty, and not to leave them in the hands of those whose
beings You created from nothing but the fire of Your wrath and Your
subduing might. They have never sensed the fragrances of compassion and
fairness, and the world has so beguiled them with its deceptions that
they have denied Your proof, associated partners with You, disbelieved
in Your signs, and shed the blood of Your loved ones and Your trusted
ones. By Your glory, O my Beloved! They have committed what no one
before them has committed, and have thereby deserved Your wrath and the
scourges of Your subduing might. Seize them through Your sovereignty,
and set over them one who will show them no mercy unless they return to
You, enter beneath the shadow of Your care, and turn to You in
repentance. You have always been powerful and will forever be mighty.
You are indeed the All-Powerful, the Most Exalted, the Just, the
All-Wise.

\subsection{767}\label{section-1262}

\begin{otherlanguage}{arabic}

سبحانك اللّهمّ يا الهى فانصر هذا المظلوم الّذى ابتلى بين يدى الظّالمين من
خلقك و المشركين من اعدائك بعد الّذى ما تنفّس الّا باذنك و امرك قد كنت يا
الهى راقداً على مهد السّكون مرّت علىّ ارياح فضلك و الطافك و ايقظتنى بها
بسلطانك و مواهبك و اقمتنى بين عبادك بثنآء نفسك و اعلآء كلمتك اذا اعترض
علىّ اكثر بريّتك

\end{otherlanguage}

Subḥānaka Allāhumma yā ilāhī fa-nṣur hādhā al-Maẓlūma alladhī ibtuliya
bayna yaday al-ẓālimīna min khalqika wa-l-mushrikīna min aʿdāʾika baʿda
alladhī mā tanaffasa illā bi-idhnika wa-amrika. Qad kuntu yā ilāhī
rāqidan ʿalā mahdi al-sukūni marrat ʿalayya aryāḥu faḍlika wa-alṭāfika
wa-ayqaẓtanī bihā bi-sulṭānika wa-mawāhibika wa-aqamtanī bayna ʿibādika
bi-thanāʾi nafsika wa-iʿlāʾi kalimatika idhan iʿtaraḍa ʿalayya aktharu
bariyyatika.

Glory be to You, O God, my God! Aid this Wronged One who has been
afflicted at the hands of the unjust among Your creatures and those
among Your enemies who associate partners with You, even though he has
never drawn a breath except by Your permission and command. I was
asleep, O my God, in the cradle of stillness. The winds of Your grace
and favors passed over me, and through them You awakened me by Your
sovereignty and Your gifts. You raised me up among Your servants to
praise You and exalt Your word. Then most of Your creatures opposed me.

\subsection{768}\label{section-1263}

\begin{otherlanguage}{persian}

فو‌عزّتك يا الهى ما ظننت فى حقّهم ما ظهر منهم بعد الّذى انّك بشّرتهم بهذا
الظّهور فى صحايف امرك و الواح قضائك و ما نزلت من عندك كلمة الّا و قد اخذت
بها عهد هذا الغلام من خلقك و بريّتك اذاً صرت متحيّراً يا الهى و لم ادر ما
افعل بين هؤلآء و كلّما اصمت عن بدايع ذكرك ينطقنى الرّوح بين سمائك و ارضك و
كلّما اسكن يهتزّنى ما تهبّ عن يمين مشيّتك و ارادتك و اجد نفسى كورقة الّتى
تحركها ارياح قضائك و تذهبها كيف تشآء بامرك و اذنك

\end{otherlanguage}

Fa-wa-ʿizzatika yā ilāhī mā ẓanantu fī ḥaqqihim mā ẓahara minhum baʿda
alladhī innaka bashshartahum bi-hādhā al-ẓuhūri fī ṣaḥāʾifi amrika
wa-alwāḥi qaḍāʾika wa-mā nazalat min ʿindika kalimatun illā wa-qad
akhadhta bihā ʿahda hādhā al-ghulāmi min khalqika wa-bariyyatika. Idhan
ṣirtu mutaḥayyiran yā ilāhī wa-lam adri mā afʿalu bayna hāʾulāʾi
wa-kullamā aṣmutu ʿan badāʾiʿi dhikrika yunṭiqunī al-Rūḥu bayna samāʾika
wa-arḍika wa-kullamā askunu yahtazzunī mā tahubbu ʿan yamīni mashīʾatika
wa-irādatika wa-ajidu nafsī ka-waraqatin allatī tuḥarrikuhā aryāḥu
qaḍāʾika wa-tudhhibuhā kayfa tashāʾu bi-amrika wa-idhnika.

By Your glory, O my God! I had never expected from them what they have
shown, after You had given them the glad tidings of this Revelation in
the scriptures of Your Cause and the Tablets of Your decree. Not a word
has descended from You through which You have not taken from Your
creatures and Your creation a covenant concerning this Youth. I am
therefore bewildered, O my God, and do not know what to do among these
people. Whenever I fall silent concerning the wonders of Your
remembrance, the Spirit makes me speak between Your heaven and Your
earth. Whenever I am still, what blows from the right hand of Your
purpose and will stirs me. I find myself like a leaf moved by the winds
of Your decree and carried wherever You will, by Your command and
permission.

\subsection{769}\label{section-1264}

\begin{otherlanguage}{persian}

و بما ظهر منّى يوقن كلّ بصير بانّ الأمر ليس بيدى بل بيدك و لم يكن زمام
الاختيار فى قبضتى بل فى قبضتك و اقتدارك مع‌ذلك يا الهى اجتمعوا علىّ اهل
مملكتك و ينزلن فى كلّ حين ما تفزع به حقايق اصفيائك و امنائك اذا اسئلك يا
الهى باسمك الّذى به اهتديت العاشقين الى كوثر فضلك و الطافك و استجذبت
المشتاقين الى رضوان قربك و لقائك بأن تفتح ابصار بريّتك ليشهدن فى هذا
الظّهور ظهور عزّ فردانيّتك و طلوع انوار وجهك و جمالك ثمّ طهّرهم يا الهى من
الظّنون و الأوهام ليجدن روايح التّقديس من قميص ظهورك و امرك لعلّ لايردنّ علىّ
ما تمنع به انفسهم من نفحات شئون رحمانيّتك فى ایّام ظهور مظهر نفسك و مطلع
امرك و لايرتكبن ما تجعل به ذواتهم مستحقّاً لظهورات قهرك و غضبك

\end{otherlanguage}

Wa-bi-mā ẓahara minnī yūqinu kullu baṣīrin bi-anna al-amra laysa bi-yadī
bal bi-yadika wa-lam yakun zimāmu al-ikhtiyāri fī qabḍatī bal fī
qabḍatika wa-iqtidārika. Maʿa dhālika yā ilāhī ijtamaʿū ʿalayya ahlu
mamlakatika wa-yunazzilunna fī kulli ḥīnin mā tafzaʿu bihi ḥaqāʾiqu
aṣfiyāʾika wa-umanāʾika. Idhan asʾaluka yā ilāhī bi-ismika alladhī bihi
ihtadayta al-ʿāshiqīna ilā kawthari faḍlika wa-alṭāfika wa-istajdhabta
al-mushtāqīna ilā riḍwāni qurbika wa-liqāʾika bi-an taftaḥa abṣāra
bariyyatika li-yashhadanna fī hādhā al-ẓuhūri ẓuhūra ʿizzi
fardāniyyatika wa-ṭulūʿa anwāri wajhika wa-jamālika. Thumma ṭahhirhum yā
ilāhī mina al-ẓunūni wa-l-awhāmi li-yajidanna rawāʾiḥa al-taqdīsi min
qamīṣi ẓuhūrika wa-amrika laʿalla lā yuridunna ʿalayya mā tamnaʿu bihi
anfusahum min nafaḥāti shuʾūni raḥmāniyyatika fī ayyāmi ẓuhūri maẓhari
nafsika wa-maṭlaʿi amrika wa-lā yartakibanna mā tajʿalu bihi dhawātuhum
mustaḥaqqan li-ẓuhūrāti qahrika wa-ghaḍabika.

Through what has proceeded from me, every discerning person knows with
certainty that the matter is not in my hands but in Yours, and that the
reins of choice are not in my grasp but in Your grasp and power.
Nevertheless, O my God, the people of Your realm have gathered against
me, and at every moment they bring down upon me what terrifies the
inmost beings of Your chosen ones and Your trusted ones. I therefore ask
You, O my God, by Your name through which You have guided the lovers to
the heavenly river of Your grace and favors, and drawn those who yearn
to the paradise of Your nearness and presence, to open the eyes of Your
creatures so that in this Revelation they may witness the manifestation
of the glory of Your singleness and the dawning of the lights of Your
countenance and beauty. Purify them, O my God, of conjectures and
delusions, that they may sense the fragrances of holiness from the
garment of Your Revelation and Your Cause. Perhaps they will not inflict
upon me what would deprive their own souls of the breaths of the
expressions of Your all-merciful grace in the days of the appearance of
the Manifestation of Your Self and the Dawning Place of Your Cause, nor
commit what would make their own beings deserving of the manifestations
of Your subduing might and wrath.

\subsection{770}\label{section-1265}

\begin{otherlanguage}{arabic}

و انت تعلم يا الهى بانّى كنت بين ملأ البيان كاَحد منهم و عاشرت معهم بالشّوق
و الاشتياق و دعوتهم الى نفسك فى العشىّ و الاشراق ببدايع وحيك و الهامك و
ورد علىّ منهم ما عجزت عن ذكره سكّان مداين انشائك فو‌عزّتك يا محبوبی‌ ما اصبحت
الّا و قد صرت هدفا لسهام غلّهم و ما امسيت الّا و قد ورد علىّ رماح بغضهم و مع
ما جعلتنى عالماً بما فى انفسهم و قادراً عليهم سترت و صبرت ناظراً الى ميقاتك

\end{otherlanguage}

Wa-anta taʿlamu yā ilāhī bi-annī kuntu bayna malaʾi al-bayāni ka-aḥadin
minhum wa-ʿāshartu maʿahum bi-l-shawqi wa-l-ishtiyāqi wa-daʿawtuhum ilā
nafsika fī al-ʿashiyyi wa-l-ishrāqi bi-badāʾiʿi waḥyika wa-ilhāmika
wa-warada ʿalayya minhum mā ʿajazat ʿan dhikrihi sukkānu madāʾini
inshāʾika. Fa-wa-ʿizzatika yā maḥbūbī mā aṣbaḥtu illā wa-qad ṣirtu
hadafan li-sihāmi ghillihim wa-mā amsaytu illā wa-qad warada ʿalayya
rimāḥu bughḍihim wa-maʿa mā jaʿaltanī ʿāliman bi-mā fī anfusihim
wa-qādiran ʿalayhim satartu wa-ṣabartu nāẓiran ilā mīqātika.

You know, O my God, that I lived among the people of the Bayān as one of
them, associated with them with ardent longing, and summoned them to You
at evening and at dawn through the wonders of Your revelation and
inspiration. Yet there has befallen me at their hands what the
inhabitants of the cities of Your creation are powerless to recount. By
Your glory, O my Beloved! No morning came without my becoming a target
for the arrows of their malice, and no evening came without the spears
of their hatred striking me. Although You had made me aware of what was
within them and given me power over them, I concealed it and remained
patient, looking toward Your appointed time.

\subsection{771}\label{section-1266}

\begin{otherlanguage}{arabic}

فلمّا جاء الوعد و تمّت الميقات حركت ذيل السّتر اقلّ من ان يحصى اذاً فزع من فى
جبروت الأمر و الخلق الّا الّذين خلقتهم من نار حبّك و هوآء شوقك و مآء عنايتك
و تراب فضلك اولئك يصلّين عليهم اهل ملأ الأعلی و سكّان مداين البقآء فلك
الحمد يا الهى بما عصمت الموحّدين و اهلكت المشركين و فصّلت بين الكلّ بكلمة
اخرى الّتى خرجت من فم مشيّتك و ظهرت من قلم ارادتك و بذلك اعترض علىّ عباد
الّذينهم خلقوا بكلمة امرك و بعثوا بارادتك و بلغوا فى الاعراض الى مقام
كفروا بك و بآياتك و حاربوا بنفسك

\end{otherlanguage}

Fa-lammā jāʾa al-waʿdu wa-tammat al-mīqātu ḥarraktu dhayla al-sitri
aqalla min an yuḥṣā idhan faziʿa man fī jabarūti al-amri wa-l-khalqi
illā alladhīna khalaqtahum min nāri ḥubbika wa-hawāʾi shawqika wa-māʾi
ʿināyatika wa-turābi faḍlika. Ulāʾika yuṣallīna ʿalayhim ahlu malaʾi
al-aʿlā wa-sukkānu madāʾini al-baqāʾi. Fa-laka al-ḥamdu yā ilāhī bi-mā
ʿaṣamta al-muwaḥḥidīna wa-ahlakta al-mushrikīna wa-faṣṣalta bayna
al-kulli bi-kalimatin ukhrā allatī kharajat min fami mashīʾatika
wa-ẓaharat min qalami irādatika. Wa-bi-dhālika iʿtaraḍa ʿalayya ʿibādun
alladhīnahum khuliqū bi-kalimati amrika wa-buʿithū bi-irādatika
wa-balaghū fī al-iʿrāḍi ilā maqāmin kafarū bika wa-bi-āyātika wa-ḥārabū
bi-nafsika.

When the promise came and the appointed time was fulfilled, I moved the
hem of the veil by an immeasurably small degree. Then all who dwell in
the dominion of command and creation were terrified, except those whom
You had created from the fire of Your love, the air of longing for You,
the water of Your care, and the dust of Your grace. Upon them the people
of the Concourse on high and the inhabitants of the cities of eternity
invoke blessings. Praise be to You, O my God, for protecting those who
affirm Your oneness, destroying those who associate partners with You,
and distinguishing everyone from one another through another word that
issued from the mouth of Your purpose and appeared from the Pen of Your
will. Because of this, servants who had been created by the word of Your
command and raised up by Your will opposed me. Their rejection reached
such a degree that they disbelieved in You and Your signs and waged war
against Your own Self.

\subsection{772}\label{section-1267}

\begin{otherlanguage}{persian}

فو‌عزّتك يا محبوبی‌ لن يقدر القلم ان يذكر ما ورد منهم على مظهر امرك و مطلع
وحيك و مشرق الهامك فلك الحمد فى كلّ ذلك و انّى فو‌عزّتك يا الهى قد كنت
مشتاقاً لما قدّر فى سماء قضائك و ملكوت تقديرك لانّ ما يرد علىّ فى سبيلك هو
محبوب ذاتى و مقصود نفسى و هذا لم يكن الّا بحولك و قوّتك انا الّذى يا الهى
بحبّك استغنيت عن كلّ من فى السّموات و الأرض و به لن اجزع و لو يرد علىّ ضرّ
العالمين

\end{otherlanguage}

Fa-wa-ʿizzatika yā maḥbūbī lan yaq dira al-qalamu an yadhkura mā warada
minhum ʿalā maẓhari amrika wa-maṭlaʿi waḥyika wa-mashriqi ilhāmika.
Fa-laka al-ḥamdu fī kulli dhālika wa-innī fa-wa-ʿizzatika yā ilāhī qad
kuntu mushtāqan li-mā quddira fī samāʾi qaḍāʾika wa-malakūti taqdīrika
li-anna mā yaridu ʿalayya fī sabīlika huwa maḥbūbu dhātī wa-maqṣūdu
nafsī wa-hādhā lam yakun illā bi-ḥawlika wa-quwwatika. Anā alladhī yā
ilāhī bi-ḥubbika istaghnaytu ʿan kulli man fī al-samāwāti wa-l-arḍi
wa-bihi lan ajzaʿa wa-law yaridu ʿalayya ḍurru al-ʿālamīna.

By Your glory, O my Beloved! The Pen can never recount what has befallen
the Manifestation of Your Cause, the Dawning Place of Your revelation,
and the Dayspring of Your inspiration at their hands. Praise be to You
in all of this! By Your glory, O my God, I have longed for what has been
ordained in the heaven of Your decree and the kingdom of Your
ordainment, for whatever befalls me in Your path is the beloved of my
being and the desire of my soul. This has come about only through Your
might and power. I am the one, O my God, who through Your love has
become independent of all who are in the heavens and on earth. Through
it I shall never be dismayed, even should the affliction of all the
worlds befall me.

\subsection{773}\label{section-1268}

\begin{otherlanguage}{arabic}

فيا ليت كان حينئذ حين الّذى فيه يسفك دمى على وجه الأرض بين يديك و تشهدنى
على حالة الّتى بها شهدت المقرّبين من عبادك و المصطفين من خيرة خلقك فلك
الحمد يا الهى على ما قضيت بسلطان قضائك و تقضى بتقديرك و امضائك اسئلك يا
محبوبی‌ باسمك الّذى به رفعت اعلام امرك و اشرقت انوار وجهك بأن تنزل علىّ و
على المخلصين من عبادك كلّ خير قدّرته فى الألواح ثمّ اجعل لنا مقعد صدق عندك
يا من بيدك ملكوت كلّ شیء و انّك انت المقتدر العزيز الرّحمن

\end{otherlanguage}

Fa-yā layta kāna ḥīnaʾidhin ḥīna alladhī fīhi yusfaku damī ʿalā wajhi
al-arḍi bayna yadayka wa-tashhadunī ʿalā ḥālatin allatī bihā shahidta
al-muqarrabīna min ʿibādika wa-al-muṣṭafayna min khīrati khalqika
fa-laka al-ḥamdu yā ilāhī ʿalā mā qaḍayta bi-sulṭāni qaḍāʾika wa-taqḍī
bi-taqdīrika wa-imḍāʾika asʾaluka yā maḥbūbī bi-smika alladhī bihi
rafaʿta aʿlāma amrika wa-ashraqat anwāru wajhika bi-an tunazzila ʿalayya
wa-ʿalā al-mukhliṣīna min ʿibādika kulla khayrin qaddartahu fī al-alwāḥi
thumma ijʿal lanā maqʿada ṣidqin ʿindaka yā man bi-yadika malakūtu kulli
shayʾin wa-innaka anta al-Muqtadiru al-ʿAzīzu al-Raḥmān.

If only this were the hour when my blood is shed upon the face of the
earth before You, and You behold me in the condition in which You beheld
those of Your servants who are near to You and the chosen ones among the
best of Your creation. Praise be to You, O my God, for what You have
decreed through the sovereignty of Your judgment and what You decree
through Your ordainment and sanction. I ask You, O my Beloved, by Your
name through which You have raised the banners of Your Cause and the
lights of Your countenance have shone, to send down upon me and upon the
sincere among Your servants every good thing You have ordained in the
Tablets. Then appoint for us a seat of truth in Your presence, O You in
whose hand is the dominion of all things. You are indeed the
All-Powerful, the Almighty, the All-Merciful.

\subsection{774}\label{section-1269}

\begin{otherlanguage}{persian}

اميد چنان است كه از فضل سبحانى در اين ربيع روحانى كلّ بخلع انسانى فايز
شوند و اثواب خلقهء عتيقهء رثيثه باليه را باثواب جديده الهيّه تبديل نمايند
و كلّ بر امر حقّ بمنتهاى استطاعت و قدرت خود قائم شوند و اهل بهاء جناح
يكديگر باشند از براى طيران كلّ در اين هواى روحانى و فضاى قدس نورانى و
بغير ما قدّره اللّه تكلّم ننمايند و بغير رضاى او مشى نكنند و بر امرى تكلّم
ننمايند مگر باذن اللّه و بادراك خود در معانى كلمات الهيّه اكتفا ننمايند چه
كه غير اللّه بر معانى كلمات خود على ما هى عليها مطّلع نه چنانچه بعضى من
غير اذن نسبت باغصان شجره الهيّه بعضى سخنها گفته‌اند كه عند‌اللّه مرضىّ نبوده
چه كه اينگونه بيانات سبب فساد و اختلاف در امر است مخصوص در اين ایّام كه
اكثرى از معرضين منتظرند كه فسادى در امر احداث شود

\end{otherlanguage}

Umīd chunān ast kih az faḍl-i subḥānī dar īn rabīʿ-i rūḥānī kull
bi-khilaʿ-i insānī fāʾiz shavand va athvāb-i khaliqih-yi ʿatīqih-yi
rathīthih-yi bāliyih rā bi-athvāb-i jadīdih-yi ilāhiyyih tabdīl namāyand
va kull bar amr-i Ḥaqq bi-muntahā-yi istiṭāʿat va qudrat-i khud qāʾim
shavand va ahl-i Bahāʾ janāḥ-i yakdīgar bāshand az barā-yi ṭayarān-i
kull dar īn havā-yi rūḥānī va faḍā-yi quds-i nūrānī va bi-ghayr-i mā
qaddarahu Allāh takallum nanamāyand va bi-ghayr-i riḍā-yi ū mashy
nakunand va bar amrī takallum nanamāyand magar bi-idhn-i Allāh va
bi-idrāk-i khud dar maʿānī-yi kalimāt-i ilāhiyyih iktifā nanamāyand chih
kih ghayr-i Allāh bar maʿānī-yi kalimāt-i khud ʿalā mā hiya ʿalayhā
muṭṭaliʿ nah chunānchih baʿḍī min ghayr-i idhn nisbat bi-aghṣān-i
shajariy-i ilāhiyyih baʿḍī sukhanhā guftih-and kih ʿinda Allāh marḍī
nabūdih chih kih īngūnih bayānāt sabab-i fasād va ikhtilāf dar amr ast
makhṣūṣ dar īn ayyām kih aktharī az muʿriḍīn muntaẓirand kih fasādī dar
amr iḥdāth shavad.

The hope is that, through divine grace, all may attain the robes of
humanity in this spiritual springtime, exchange their old, worn,
threadbare, and tattered garments for new, divine garments, and arise to
uphold the Cause of the True One to the utmost of their ability and
power. May the people of Bahāʾ be wings for one another, so that all may
soar in this spiritual atmosphere and this luminous expanse of holiness.
May they speak only what God has ordained, walk only in accordance with
His good pleasure, and speak of no matter except with God's permission.
They should not content themselves with their own understanding of the
meanings of the divine words, for none but God knows the meanings of His
own words as they truly are. Thus, some have, without permission, said
certain things concerning the branches of the divine Tree that have not
been pleasing to God, for such statements cause corruption and
dissension in the Cause, especially in these days when most of those who
turn away are waiting for some corruption to arise in the Cause.

\subsection{775}\label{section-1270}

\begin{otherlanguage}{persian}

اتّقوا اللّه يا قوم و لاتكوننّ من المتجاوزين ان اذكروا ما امرناكم به لا بما
امركم انفسكم و اهوائكم اتّقوا اللّه و كونوا من المتّقين انا انطقناه حين
الّذى قضت من عمره ستّة سنين متواليات ليكون دليلاً على هذا النّبأ الأعظم
العظيم حقّ او را ناطق فرمود تا دليل باشد بر امرش نفوسى كه بچنين كلمات
تكلّم نموده ‌از عظمت اين امر غافل شده‌اند لأنّه لو يريد ان ينطق كلّ شیء ليقدر
و انّه لعلى كلّ شیء قدير قبل از الف ل ايظهر غير من ظهر كذلك نزل الأمر فى
اوّل هذا الظّهور لو انتم من العارفين

\end{otherlanguage}

Ittaqū Allāha yā qawmu wa-lā takūnanna mina al-mutajāwizīn in udhkurū mā
amarnākum bihi lā bi-mā amarakum anfusukum wa-ahwāʾukum ittaqū Allāha
wa-kūnū mina al-muttaqīn innā anṭaqnāhu ḥīna alladhī qaḍat min ʿumrihi
sittatu sinīna mutawāliyātin li-yakūna dalīlan ʿalā hādhā al-nabaʾi
al-aʿẓami al-ʿaẓīm Ḥaqq ū rā nāṭiq farmūd tā dalīl bāshad bar amrash
nufūsī kih bi-chunīn kalimāt takallum namūdih az ʿaẓamat-i īn amr ghāfil
shudih-and li-annahu law yurīdu an yunṭiqa kulla shayʾin la-yaqdiru
wa-innahu la-ʿalā kulli shayʾin qadīr qabl az alf {[}corrupted text{]}
yaẓharu ghayru man ẓahara kadhālika nazala al-amru fī awwali hādhā
al-ẓuhūri law antum mina al-ʿārifīn.

Fear God, O people, and do not be among those who transgress. Remember
what We have commanded you, not what your own selves and desires have
commanded you. Fear God and be among the God-fearing. We caused him to
speak when six successive years of his life had passed, that he might be
a proof of this most great and mighty Announcement. The True One caused
him to speak so that he might be a proof of His Cause. Those who have
spoken such words have become heedless of the greatness of this Cause,
for if He wishes to cause all things to speak, He is able to do so, and
He indeed has power over all things. Before a thousand, {[}corrupted
original text{]} anyone other than the One who has appeared appears.
Thus was the command sent down at the beginning of this Revelation, if
you are among those who understand.

\subsection{776}\label{section-1271}

\begin{otherlanguage}{arabic}

و قد بعث اللّه و سوف يبعث لهذا الأمر ادلّاء ينطقن عنه و يتوجّهنّ فى كلّ حين
اليه و لايتحرّكنّ الّا بعد اذنه اولئك لا يبدلنّ نغمة من نغماتى بكلّ من فى
السّموات و الأرضين اولئك ادلّائى فى الأرض و اصفيائى فى البلاد و امنائى بين
العباد عليهم رحمتى و بهائى و ذكرى و ثنائى و هم عباد مكرمون

\end{otherlanguage}

Wa-qad baʿatha Allāhu wa-sawfa yabʿathu li-hādhā al-amri adillāʾa
yanṭiqna ʿanhu wa-yatawajjahnna fī kulli ḥīnin ilayhi wa-lā yataḥarrakna
illā baʿda idhnihi ulāʾika lā yubaddilna naghmatan min naghamātī
bi-kulli man fī al-samāwāti wa-al-araḍīn ulāʾika adillāʾī fī al-arḍi
wa-aṣfiyāʾī fī al-bilādi wa-umanāʾī bayna al-ʿibādi ʿalayhim raḥmatī
wa-bahāʾī wa-dhikrī wa-thanāʾī wa-hum ʿibādun mukramūn.

God has raised up, and will raise up, guides for this Cause who speak on
His behalf, turn toward Him at all times, and make no move without His
permission. They will not exchange a single one of My melodies for all
who are in the heavens and on earth. They are My guides on earth, My
chosen ones in the lands, and My trusted ones among the servants. Upon
them rest My mercy, My glory, My remembrance, and My praise. They are
honored servants.

\subsection{777}\label{section-1272}

\begin{otherlanguage}{persian}

آنقدر معلوم بوده كه الى حين تفصيل اين ظهور ابدع ابهى ذكر نشده لو يذكر
ليشهدنّ الكلّ بانّه لهو الّذى ما سبقه احد من قبل و لن يسبقه احد من الآخرین و
من يدّعى قبل الألف انّه افترى على اللّه لانّ بذلك يفسد الأمر و لا يستقرّ هيكل
المشيّة على عرش عزّ عظيم در اوّل اين ظهور منيع در اينمقام آياتى نازل و بعضى
از آن در بعضى از الواح مذكور ان اشهدوها لتستقرّوا على الأمر و تكوننّ من
الثّابتين

\end{otherlanguage}

Ānqadr maʿlūm būdih kih ilā ḥīn tafṣīl-i īn ẓuhūr-i abdaʿ-i abhā dhikr
nashudih law yudhkaru la-yashhadanna al-kullu bi-annahu la-huwa alladhī
mā sabaqahu aḥadun min qablu wa-lan yasbiqahu aḥadun mina al-ākhirīn
wa-man yaddaʿī qabla al-alfi innahu iftarā ʿalā Allāhi li-anna
bi-dhālika yafsudu al-amru wa-lā yastaqirru haykalu al-mashīʾati ʿalā
ʿarshi ʿizzin ʿaẓīm dar avval-i īn ẓuhūr-i manīʿ dar īn maqām āyātī
nāzil va baʿḍī az ān dar baʿḍī az alvāḥ madhkūr in ishhadūhā
li-tastaqirrū ʿalā al-amri wa-takūnanna mina al-thābitīn.

Let this much be known: until now, the details of this most wondrous,
most glorious Revelation have not been set forth. Were they set forth,
all would surely testify that He is the One whom no one before Him has
surpassed and whom no one among those to come will surpass. Whoever
makes a claim before the thousand has fabricated a lie against God, for
this corrupts the Cause, and the Temple of the Will cannot be
established upon a throne of mighty glory. At the beginning of this
exalted Revelation, verses were revealed on this subject, and some of
them are recorded in certain Tablets. Consider them, that you may stand
firm in the Cause and be among the steadfast.

\subsection{778}\label{section-1273}

\begin{otherlanguage}{persian}

و ديگر معلوم بوده كه مقصود از وادى نبيل كه در الواح ذكر شده نفس اين ظهور
اعظم بوده و خواهد بود و وادى آن ارض قدسى است كه فوق طبقات جنان خلق شده و
لايقدر ان يرد عليها الّا اهل البهآء تاللّه انّهم اهل وادى النّبيل و وادى
العظمة و وادى الاستقلال و وادى الجلال و وادى الاقتدار و وادى العزّة و
وادى البقآء و وادى التّقديس و وادى التّنزيه و وادى الكبريا و غير اين نفوس
از وادى نبيل خارج و در تيه نفس و هوى متحيّر و مبتلى و كينونة النّبيل حينئذ
يتبرّء منهم و من اعمالهم و افعالهم و ما يظهر منهم و بذلك يشهد كلّ شیء ان
انتم تعلمون

\end{otherlanguage}

Va dīgar maʿlūm būdih kih maqṣūd az vādī-yi Nabīl kih dar alvāḥ dhikr
shudih nafs-i īn ẓuhūr-i aʿẓam būdih va khvāhad būd va vādī ān arḍ-i
qudsī ast kih fawq-i ṭabaqāt-i jinān khalq shudih va lā yaqdiru an
yarida ʿalayhā illā ahlu al-Bahāʾ ta-Allāhi innahum ahlu wādī al-Nabīli
wa-wādī al-ʿaẓamati wa-wādī al-istiqlāli wa-wādī al-jalāli wa-wādī
al-iqtidāri wa-wādī al-ʿizzati wa-wādī al-baqāʾi wa-wādī al-taqdīsi
wa-wādī al-tanzīhi wa-wādī al-kibriyā va ghayr-i īn nufūs az vādī-yi
Nabīl khārij va dar tīh-i nafs va havā mutaḥayyir va mubtalā va
kaynūnatu al-Nabīli ḥīnaʾidhin yatabarraʾu minhum wa-min aʿmālihim
wa-afʿālihim wa-mā yaẓharu minhum wa-bi-dhālika yashhadu kullu shayʾin
in antum taʿlamūn.

Furthermore, let it be known that what is meant by the Valley of Nabīl
mentioned in the Tablets has been, and will be, this Most Great
Revelation itself. That valley is the holy land created above the levels
of paradise, and none but the people of Bahāʾ can enter it. By God, they
are the people of the Valley of Nabīl, the Valley of Grandeur, the
Valley of Independence, the Valley of Majesty, the Valley of Power, the
Valley of Might, the Valley of Eternity, the Valley of Sanctification,
the Valley of Transcendence, and the Valley of Glory. All others are
outside the Valley of Nabīl, bewildered and afflicted in the wilderness
of self and desire. The very being of Nabīl then disavows them, their
works, their deeds, and whatever proceeds from them. To this all things
bear witness, if you know.

\subsection{779}\label{section-1274}

\begin{otherlanguage}{persian}

ایّاکم يا قوم ان تحرفوا كلمات اللّه عن مواضعها ان اشهدوا بما شهد اللّه ثمّ
كونوا اخوانا على سرر متقابلين ان اتّحدوا على شأن لايجد احد بينكم روايح
الاختلاف و باتّحادكم يظهر امر اللّه بين بريّته و تثبت آثاره بين عباده و
تبرهن انواره بين العالمين اميد چنان است كه كلّ بكلمه جامعه مجتمع شده بذكر
و ثناى حقّ مشغول شوند و بكمال انقطاع در ارض ظاهر گردند و بشأنى بنار حبّ
الهى مشتعل شوند كه كلّ بحرارت آن مشتعل شده بشطر قدس ناظر گردند

\end{otherlanguage}

Iyyākum yā qawmu an tuḥarrifū kalimāti Allāhi ʿan mawāḍiʿihā in ishhadū
bi-mā shahida Allāhu thumma kūnū ikhwānan ʿalā sururin mutaqābilīn in
ittaḥidū ʿalā shaʾnin lā yajidu aḥadun baynakum rawāʾiḥa al-ikhtilāfi
wa-bi-ittiḥādikum yaẓharu amru Allāhi bayna bariyyatihi wa-tathbutu
āthāruhu bayna ʿibādihi wa-tubarhanu anwāruhu bayna al-ʿālamīn umīd
chunān ast kih kull bi-kalimiy-i jāmiʿih mujtamiʿ shudih bi-dhikr va
thanā-yi Ḥaqq mashghūl shavand va bi-kamāl-i inqiṭāʿ dar arḍ ẓāhir
gardand va bi-shaʾnī bi-nār-i ḥubb-i ilāhī mushtaʿil shavand kih kull
bi-ḥarārat-i ān mushtaʿil shudih bi-shaṭr-i quds nāẓir gardand.

Beware, O people, lest you distort the words of God from their proper
places. Bear witness to what God has witnessed, and be brothers seated
upon couches facing one another. Unite in such a way that no one can
detect among you the scents of discord. Through your unity, the Cause of
God will become manifest among His creatures, its traces will be
established among His servants, and its lights will be demonstrated
among the peoples of the world. The hope is that all may gather around
the all-embracing Word, occupy themselves with the remembrance and
praise of the True One, and appear on earth in perfect detachment. May
they be so ablaze with the fire of divine love that all may be kindled
by its heat and turn their gaze toward the realm of holiness.

\subsection{780}\label{section-1275}

\begin{otherlanguage}{persian}

و آنچه از آيات الهى ادراك نشود از معدن و مبدء آن سؤال نمايند مقصود از
اينكلمه آنكه مباد در معنى كلمات مابين احبّاى الهى اختلاف شود و بدانند كه
كلمات الهيّه كلّ از نقطه واحده ظاهر و باو راجع ایّاکم ان تشهدوا الاختلاف
فيهنّ اگرچه بظاهر مختلف نازل شود در باطن بكلمه واحده و نقطه واحده راجع و
منتهى چنانچه در بعضى مقام ذكر شده كه احدى بمعنى كلمات الهى مطّلع نه الّا
اللّه شكّى نيست كه اين قول حقّ است چه كه حرفى از حروفات منزله در هر عالمى
از عوالم لايتناهى اثر مخصوصى و معانى مخصوصه باقتضاى آن عالم داشته و
خواهد داشت و هم‌نين حروفات و كلمات الهيّة هر يك كنز علم لدنّى بوده وخواهد
بود و ما اطّلع بما كنز فيهنّ الّا اللّه احدى قادر بر عرفان آن على ما ينبغى
لها نبوده و نخواهد بود

\end{otherlanguage}

Va ānchih az āyāt-i ilāhī idrāk nashavad az maʿdan va mabdaʾ-i ān suʾāl
namāyand maqṣūd az īn kalimih ānkih mabād dar maʿnī-yi kalimāt mābayn-i
aḥibbā-yi ilāhī ikhtilāf shavad va bidānand kih kalimāt-i ilāhiyyih kull
az nuqṭiy-i vāḥidih ẓāhir va bi-ū rājiʿ iyyākum an tashhadū al-ikhtilāfa
fīhinna agarchih bi-ẓāhir mukhtalif nāzil shavad dar bāṭin bi-kalimiy-i
vāḥidih va nuqṭiy-i vāḥidih rājiʿ va muntahī chunānchih dar baʿḍī maqām
dhikr shudih kih aḥadī bi-maʿnī-yi kalimāt-i ilāhī muṭṭaliʿ nah illā
Allāh shakkī nīst kih īn qawl ḥaqq ast chih kih ḥarfī az ḥurūfāt-i
munazzalih dar har ʿālamī az ʿavālim-i lā-yatanāhī athar-i makhṣūṣī va
maʿānī-yi makhṣūṣih bi-iqtiḍā-yi ān ʿālam dāshtih va khvāhad dāsht va
{[}corrupted text{]} ḥurūfāt va kalimāt-i ilāhiyyih har yak kanz-i
ʿilm-i ladunnī būdih va khvāhad būd wa-mā iṭṭalaʿa bi-mā kuniza fīhinna
illā Allāh aḥadī qādir bar ʿirfān-i ān ʿalā mā yanbaghī lahā nabūdih va
nakhvāhad būd.

Whatever they do not understand of the divine verses, they should ask
about from their mine and source. The purpose of this statement is that
no disagreement over the meaning of the words should arise among the
loved ones of God, and that they should know that the divine words all
proceed from a single point and return to it. Beware lest you perceive
discord among them. Although outwardly they may be revealed in differing
forms, inwardly they return and lead to a single word and a single
point. Thus, in certain passages it has been stated that no one knows
the meaning of the divine words except God. There is no doubt that this
statement is true, for a single letter among the revealed letters has
had, and will have, a particular effect and particular meanings in each
of the infinite worlds, according to the requirements of that world. And
{[}corrupted original text{]} the divine letters and words have each
been, and will be, a treasury of knowledge bestowed by God. None but God
has known what is treasured within them. No one has been, or will be,
capable of knowing them as they deserve to be known.

\subsection{781}\label{section-1276}

\begin{otherlanguage}{persian}

و هم‌چنين در مقامى ذكر شده كه كلّ آيات الهى را درك مينمايند چه كه اگر صاحب
اينمقام نبودند حجّيت آيات مبرهن نميشد اين دو بيان مختلفند بر حسب ظاهر و
لكن نزد متبصّرين و صاحبان ابصر حديده معلومست كه مقصود از آنكه فرموده كلّ
ادراك مينمايند على قدرهم و مقدارهم بوده لا على قدر و مقدار ما نزل من لدى
اللّه المقتدر المهيمن القيّوم

\end{otherlanguage}

Va hamchunīn dar maqāmī dhikr shudih kih kull āyāt-i ilāhī rā dark
mīnamāyand chih kih agar ṣāḥib-i īn maqām nabūdand ḥujjiyyat-i āyāt
mubarhan namīshud īn du bayān mukhtalifand bar ḥasab-i ẓāhir va lākin
nazd-i mutabaṣṣirīn va ṣāḥibān-i abṣār-i ḥadīdih maʿlūm ast kih maqṣūd
az ānkih farmūdih kull idrāk mīnamāyand ʿalā qadrihim wa-miqdārihim
būdih lā ʿalā qadri wa-miqdāri mā nazala min ladā Allāhi al-Muqtadiri
al-Muhaymini al-Qayyūm.

Likewise, in another passage it has been stated that all understand the
divine verses, for if they did not possess this capacity, the authority
of the verses as proof could not be demonstrated. These two statements
differ outwardly, but to those endowed with insight and keen vision it
is clear that the statement that all understand means according to their
own capacity and measure, not according to the capacity and measure of
what has been sent down from God, the All-Powerful, the All-Protecting,
the Self-Subsisting.

\subsection{782}\label{section-1277}

\begin{otherlanguage}{persian}

اليوم نفوسيكه در هتك حرمت بيان ساعيند چنانچه ديده و شنيده‌ايد آن نفوس نظر
بحماقت مريدين خود ادّعاى امر نموده‌اند لعن اللّه على الّذين تجاوزوا عن حدود
البيان و هتكوا حرمة اللّه و نقضوا ميثاقه و كفروا بنقطة الأولى و محبوبه و
مقصوده الّذى ظهر باسمه الأبهی بين الأرض و السّمآء و منه استضآء افئدة
المقرّبين و صدور المخلصين و لكن قد اخطاؤا ما توهّموا بحول اللّه و قوّته حصن
امر بيد قدرت حفظ خواهد شد و انّه لعلى كلّ شیء قدير

\end{otherlanguage}

Al-yawm nufūsī kih dar hatk-i ḥurmat-i Bayān sāʿīand chunānchih dīdih va
shinīdih-īd ān nufūs naẓar bi-ḥamāqat-i murīdīn-i khud iddiʿā-yi amr
namūdih-and laʿnu Allāhi ʿalā alladhīna tajāwazū ʿan ḥudūdi al-Bayāni
wa-hatakū ḥurmata Allāhi wa-naqaḍū mīthāqahu wa-kafarū bi-Nuqṭati al-Ūlā
wa-maḥbūbihi wa-maqṣūdihi alladhī ẓahara bi-smihi al-Abhā bayna al-arḍi
wa-al-samāʾi wa-minhu istaḍāʾa afʾidatu al-muqarrabīna wa-ṣudūru
al-mukhliṣīn va lākin qad akhṭaʾū mā tavahhamū bi-ḥawl-i Allāh va
quvvat-i ū ḥiṣn-i amr bi-yad-i qudrat ḥifẓ khvāhad shud wa-innahu
la-ʿalā kulli shayʾin qadīr.

Today, those who strive to violate the sanctity of the Bayān, as you
have seen and heard, have put forward claims to a Cause in view of the
foolishness of their followers. God's curse be upon those who have
transgressed the bounds of the Bayān, violated the sanctity of God,
broken His covenant, and disbelieved in the Primal Point and in His
Beloved and His Desired One, who has appeared in His name, the
All-Glorious, between earth and heaven, and from whom the hearts of
those brought near and the breasts of the sincere have received light.
But they have been mistaken in what they imagined. Through God's power
and His might, the fortress of the Cause will be guarded by the hand of
power. He indeed has power over all things.

\subsection{783}\label{section-1278}

\begin{otherlanguage}{arabic}

و فى آخر القول غرّدت ورقآء الأمر على افنان سدرة عزّ رفيع و يقول يا ملأ
البيان اتّقوا اللّه و لاتفسدوا فى الأرض بعد اصلاحها و لاترتكبوا الفحشآء
طهّروا انفسكم عن كلّ ما حرّم فى كتاب اللّه العلىّ العظيم ایّاکم ان تسفكوا
الدّمآء و الّذى سفك انّه ليس منّى و كان اللّه برئ منه و هذا ما نزل فى اوّل هذا
الظّهور من ملكوت ربّكم الرّحمن الرّحيم

\end{otherlanguage}

Wa-fī ākhiri al-qawli gharradat warqāʾu al-amri ʿalā afnāni sidrati
ʿizzin rafīʿin wa-yaqūlu yā malaʾa al-Bayāni ittaqū Allāha wa-lā tufsidū
fī al-arḍi baʿda iṣlāḥihā wa-lā tartakibū al-faḥshāʾa ṭahhirū anfusakum
ʿan kulli mā ḥurrima fī kitābi Allāhi al-ʿAliyyi al-ʿAẓīm iyyākum an
tasfikū al-dimāʾa wa-alladhī safaka innahu laysa minnī wa-kāna Allāhu
barīʾan minhu wa-hādhā mā nazala fī awwali hādhā al-ẓuhūri min malakūti
rabbikumu al-Raḥmāni al-Raḥīm.

At the close of this discourse, the Dove of the Cause sang upon the
branches of the Tree of lofty glory, saying: ``O people of the Bayān,
fear God, and do not spread corruption on earth after it has been set
right, nor commit indecency. Purify yourselves of everything forbidden
in the Book of God, the Most High, the Most Great. Beware lest you shed
blood. Whoever has shed blood is not of Me, and God disavows him.'' This
is what was sent down at the beginning of this Revelation from the
kingdom of your Lord, the All-Merciful, the Compassionate.

\subsection{784}\label{section-1279}

\begin{otherlanguage}{arabic}

ایّاکم ان ترتكبوا ما يكرهه عقولكم و افئدتكم اتّقوا اللّه يا ملأ المؤمنين ان
انقطعوا عن الدّنيا و ما فيها و لاتصرّفوا فى اموال النّاس الّا بعد اذنهم كذلك
نزلنا من قبل و حينئذ من لدن عزيز جميل كونوا امنآء اللّه فى البلاد بحيث لو
تمرّنّ على مداين الذّهب لن ترتدّ ابصاركم اليها كذلك ينبغى لكم يا ملأ
الموحّدين ان انصروا ربّكم الرّحمن باعمالكم ليجدنّ منكم روايح الحقّ عباد اللّه
بين السّموات و الأرضين تخلّقوا يا قوم باخلاق اللّه و لاترتكبوا ما تضيع به
حرمة مظهر نفسه بين خلقه كذلك وصّاكم اللّه فى الواح عزّ حفيظ لا تتعرّضوا باحد
و لا تجادلوا بنفس ان وجدتم من مقبل ذكّروه بكلمة الأعظم و ان وجدتم من معرض
فاعرضوا عنه ثمّ اقبلوا الى اللّه ربّكم و انّه يكفيكم عن العالمين

\end{otherlanguage}

Iyyākum an tartakibū mā yakrahuhu ʿuqūlukum wa-afʾidatukum ittaqū Allāha
yā malaʾa al-muʾminīn in inqaṭiʿū ʿani al-dunyā wa-mā fīhā wa-lā
taṣarrafū fī amwāli al-nāsi illā baʿda idhnihim kadhālika nazzalnā min
qablu wa-ḥīnaʾidhin min ladun ʿazīzin jamīl kūnū umanāʾa Allāhi fī
al-bilādi bi-ḥaythu law tamurrunna ʿalā madāʾini al-dhahabi lan tartadda
abṣārukum ilayhā kadhālika yanbaghī lakum yā malaʾa al-muwaḥḥidīn in
unṣurū rabbakumu al-Raḥmāna bi-aʿmālikum li-yajidanna minkum rawāʾiḥa
al-Ḥaqqi ʿibādu Allāhi bayna al-samāwāti wa-al-araḍīn takhallaqū yā
qawmu bi-akhlāqi Allāhi wa-lā tartakibū mā taḍīʿu bihi ḥurmatu maẓhari
nafsihi bayna khalqihi kadhālika waṣṣākumu Allāhu fī alwāḥi ʿizzin ḥafīẓ
lā tataʿarraḍū bi-aḥadin wa-lā tujādilū bi-nafsin in wajadtum min
muqbilin dhakkirūhu bi-kalimati al-aʿẓami wa-in wajadtum min muʿriḍin
fa-aʿriḍū ʿanhu thumma aqbilū ilā Allāhi rabbikum wa-innahu yakfīkum
ʿani al-ʿālamīn.

Beware lest you commit what your minds and hearts abhor. Fear God, O
company of believers. Detach yourselves from the world and all that is
in it, and do not dispose of people's property without their permission.
Thus have We sent down before and now, from One who is Mighty and
Beautiful. Be God's trusted ones in the lands, such that, were you to
pass through cities of gold, your eyes would not turn toward them. Such
should be your conduct, O company of those who affirm His oneness. Aid
your Lord, the All-Merciful, through your deeds, so that God's servants
throughout the heavens and the earth may sense from you the fragrances
of the True One. Adorn yourselves, O people, with the attributes of God,
and do not commit what would diminish the honor of the Manifestation of
His own Self among His creatures. Thus has God counseled you in Tablets
of preserving might. Do not assail anyone or dispute with any soul. If
you find someone receptive, remind him through the word of the Most
Great; and if you find someone who turns away, turn away from him, then
turn toward God, your Lord. He will indeed suffice you against all the
peoples of the world.

\subsection{785}\label{section-1280}

\begin{otherlanguage}{persian}

ان اتّحدوا يا احبّائى ثمّ اجتمعوا على شريعة امرى و لا تكوننّ من المختلفين و
انّ اقربكم الى اللّه اخشاكم و اخشعكم لو انتم من العارفين تاللّه الحقّ لو
انتم تظلمون احبّ عندى من ان تظلموا على احد و هذا من سجيّتى و احسن خصالى لو
انتم من الموقنين ان اصبروا يا احبّائى فى البأسآء و الضّرّآء و ان يظلمكم من
ظالم فارجعوا حكمه الى اللّه الّذى بيده ملكوت كلّ شیء و انّه لهو المقتدر على
ما یشآء و انّه هو اشدّ المنتقمين

\end{otherlanguage}

In ittaḥidū yā aḥibbāʾī thumma ijtamiʿū ʿalā sharīʿati amrī wa-lā
takūnanna mina al-mukhtalifīn wa-inna aqrabakum ilā Allāhi akhshākum
wa-akhshaʿukum law antum mina al-ʿārifīn ta-Allāhi al-Ḥaqqi law antum
tuẓlamūna aḥabbu ʿindī min an taẓlimū ʿalā aḥadin wa-hādhā min sajiyyatī
wa-aḥsani khiṣālī law antum mina al-mūqinīn in iṣbirū yā aḥibbāʾī fī
al-baʾsāʾi wa-al-ḍarrāʾi wa-in yaẓlimkum min ẓālimin fa-arjiʿū ḥukmahu
ilā Allāhi alladhī bi-yadihi malakūtu kulli shayʾin wa-innahu la-huwa
al-Muqtadiru ʿalā mā yashāʾu wa-innahu huwa ashaddu al-muntaqimīn.

Unite, O My loved ones, and gather around the law of My Cause. Do not be
among those who disagree. The nearest of you to God are those who most
fear Him and are most humble before Him, if you are among those who
understand. By God, the True One, that you should suffer injustice is
dearer to Me than that you should wrong anyone. This is of My nature and
among My finest attributes, if you are among those who are certain. Be
patient, O My loved ones, in hardship and adversity. If an oppressor
wrongs you, leave judgment upon him to God, in whose hand is the
dominion of all things. He indeed has power to do whatever He wills, and
He is the most severe of avengers.

\subsection{786}\label{section-1281}

\begin{otherlanguage}{arabic}

و انّا وصّيناكم فى اكثر الألواح بأن لا تتعرّضوا بأحد و انتم فعلتم ما رجع
ضرّه الى سدرة عزّ رفيع يا قوم اتّقوا اللّه و لاتجاوزوا عمّا امرتم به فى
الألواح و انّ هذا خير لكم عمّا عندكم لو انتم من العارفين ان اجتنبوا من
الّذينهم كفروا و اشركوا باللّه و جادلوا بآياته و حاربوا بنفسه العلىّ العظيم
ثمّ اصبروا كما صبرنا و لاتحاربوا مع احد من المشركين توكّلوا فى كلّ الأمور
على اللّه و انّه ينصرنى كما نصرنى اوّل مرّة اذ كنت فى السّجن تحت السّلاسل و
الأغلال و نصرنى بالحقّ بجنود السّموات و الأرضين الى ان اخرجنى عنه بسلطان
مبين فى قبضته ملكوت كلّ شیء ينصر من یشآء كيف یشآء و انّه لهو المقتدر
القدير

\end{otherlanguage}

Wa-innā waṣṣaynākum fī akthar al-alwāḥ bi-an lā tataʿarraḍū bi-aḥad
wa-antum faʿaltum mā rajaʿa ḍarruhu ilā sidrat ʿizz rafīʿ. Yā qawm
ittaqū Allāh wa-lā tatajāwazū ʿammā umirtum bihi fī al-alwāḥ wa-inna
hādhā khayr lakum ʿammā ʿindakum law antum min al-ʿārifīn. An ijtanibū
min alladhīnahum kafarū wa-ashrakū bi-Allāh wa-jādalū bi-āyātihi
wa-ḥārabū bi-nafsihi al-ʿAlī al-ʿAẓīm. Thumma iṣbirū kamā ṣabarnā wa-lā
tuḥāribū maʿa aḥad min al-mushrikīn. Tawakkalū fī kull al-umūr ʿalā
Allāh wa-innahu yanṣurunī kamā naṣaranī awwal marra idh kuntu fī al-sijn
taḥta al-salāsil wa-al-aghlāl wa-naṣaranī bi-al-ḥaqq bi-junūd al-samāwāt
wa-al-arḍīn ilā an akhrajanī ʿanhu bi-sulṭān mubīn. Fī qabḍatihi malakūt
kull shayʾ yanṣur man yashāʾ kayfa yashāʾ wa-innahu la-huwa al-Muqtadir
al-Qadīr.

We have counseled you in most of the Tablets not to assail anyone, yet
you have done what has brought harm upon the Tree of lofty glory. O
people, fear God and do not transgress what you have been commanded in
the Tablets. This is indeed better for you than what you possess, if you
are among those who understand. Avoid those who have disbelieved,
associated partners with God, disputed His verses, and made war against
His own exalted and mighty Self. Then be patient as We have been
patient, and do not wage war against any of those who associate partners
with God. Put your trust in God in all matters. He will indeed aid Me as
He aided Me the first time, when I was in prison under chains and
fetters. He aided Me in truth with the hosts of the heavens and the
earth until He brought Me out of it with manifest authority. In His
grasp is the dominion of all things. He aids whomever He wills, however
He wills. He is indeed the All-Powerful, the Omnipotent.

\subsection{787}\label{section-1282}

\begin{otherlanguage}{arabic}

ایّاکم يا قوم ان تتمسّكوا بالأسمآء و تعترضوا على الّذى خلقها بامر من عنده
لا اله الّا هو المقتدر العالم الحكيم قل انّه خلق الأسماء كما خلق كلّ شیء ان
انتم من العالمين كلّ ذى قدرة عاجز عند سلطان قدرته و كلّ ذى علم جاهل لدى
علمه و كلّ ذى غنآء فقير لدى بابه الّذى فتح على وجه الخلايق اجمعين ایّاکم ان
يصدّنكم الأسماء عن موجدها تجنّبوا عن الّذين تمسّكوا بها و كانوا عليها لمن
العاكفين قل انّ النّبی‌ من سمع نبأى و آمن بنفسى و الرّسول من بلّغ رسالاتى و
الأمام من قام امام وجهى وفاز بأیّامى و الولىّ من دخل حصن ولايتى و انقطع عن
سوائى و الوصىّ من وصّى نفسه ثمّ العباد بحبّی‌ و ذكرى و ثنائى كذلك اشرقت شمس
البيان عن افق مشيّة ربّكم الرّحمن فى هذا الرّضوان المقدّس الممتنع المتعالى
المنيع و الحمد للّه محبوبی‌ و محبوب العارفين

\end{otherlanguage}

Iyyākum yā qawm an tatamassakū bi-al-asmāʾ wa-taʿtariḍū ʿalā alladhī
khalaqahā bi-amr min ʿindihi lā ilāha illā huwa al-Muqtadir al-ʿĀlim
al-Ḥakīm. Qul innahu khalaqa al-asmāʾ kamā khalaqa kull shayʾ in antum
min al-ʿālimīn. Kull dhī qudra ʿājiz ʿinda sulṭān qudratihi wa-kull dhī
ʿilm jāhil ladā ʿilmihi wa-kull dhī ghināʾ faqīr ladā bābihi alladhī
fataḥa ʿalā wajh al-khalāʾiq ajmaʿīn. Iyyākum an yaṣuddannakum al-asmāʾ
ʿan mūjidihā tajannabū ʿan alladhīna tamassakū bihā wa-kānū ʿalayhā
la-min al-ʿākifīn. Qul inna al-nabī man samiʿa nabaʾī wa-āmana bi-nafsī
wa-al-rasūl man ballagha risālātī wa-al-imām man qāma amāma wajhī
wa-fāza bi-ayyāmī wa-al-walī man dakhala ḥiṣn wilāyatī wa-inqaṭaʿa ʿan
siwāʾī wa-al-waṣī man waṣṣā nafsahu thumma al-ʿibād bi-ḥubbī wa-dhikrī
wa-thanāʾī. Kadhālika ashraqat shams al-bayān ʿan ufuq mashīʾat rabbikum
al-Raḥmān fī hādhā al-Riḍwān al-muqaddas al-mumtaniʿ al-mutaʿālī
al-manīʿ wa-al-ḥamd li-llāh maḥbūbī wa-maḥbūb al-ʿārifīn.

Beware, O people, lest you cling to names and object to the One who
created them by a command from Himself. There is no God but Him, the
All-Powerful, the All-Knowing, the All-Wise. Say: He created the names
just as He created all things, if you are among those who know. Everyone
endowed with power is helpless before the sovereignty of His power;
everyone endowed with knowledge is ignorant before His knowledge; and
everyone endowed with wealth is poor at His door, which He has opened to
all created beings. Beware lest the names keep you from their Creator.
Shun those who have clung to them and remained devoted to them. Say: The
prophet is the one who has heard My tidings and believed in My own Self;
the messenger is the one who has conveyed My messages; the imam is the
one who has stood before My face and attained My days; the guardian is
the one who has entered the stronghold of My guardianship and detached
himself from all but Me; and the appointed successor is the one who has
enjoined upon himself, and then upon the servants, love for Me,
remembrance of Me, and praise of Me. Thus has the sun of utterance shone
from the horizon of the will of your Lord, the All-Merciful, in this
sanctified, inaccessible, exalted, and impregnable Paradise. Praise be
to God, My Beloved and the Beloved of those who know Him.

\subsection{788}\label{section-1283}

\begin{otherlanguage}{arabic}

و اذا نخاطب الّذى افترى على اللّه و اعترض عليه و ارسل الى مقرّ القدس اوراق
النّاريّة و ذكر فيها ما احترقت عنها اكباد الأصفیا‌ء و المقرّبين لعلّ يستيقظ
عن نوم الغفلة و يتوجّه الى شطر القدس مقرّ اللّه المقتدر العزيز العليم و لو
انّه اتى بذنب اثقل عن ثقل السّموات و الأرضين و لكن لمّا تموّجت ابحر الفضل فى
تلك الأیّام لو يتوب و يرجع لعلّ ان يغفره اللّه برحمته الّتى سبقت العالمين

\end{otherlanguage}

Wa-idhā nukhāṭib alladhī iftarā ʿalā Allāh wa-iʿtaraḍa ʿalayhi wa-arsala
ilā maqarr al-quds awrāq al-nāriyya wa-dhakara fīhā mā iḥtaraqat ʿanhā
akbād al-aṣfiyāʾ wa-al-muqarrabīn laʿalla yastayqiẓ ʿan nawm al-ghafla
wa-yatawajjaha ilā shaṭr al-quds maqarr Allāh al-Muqtadir al-ʿAzīz
al-ʿAlīm. Wa-law innahu atā bi-dhanb athqal ʿan thiqal al-samāwāt
wa-al-arḍīn wa-lākin lammā tamawwajat abḥur al-faḍl fī tilka al-ayyām
law yatūb wa-yarjaʿ laʿalla an yaghfirahu Allāh bi-raḥmatihi allatī
sabaqat al-ʿālamīn.

And now We address the one who has fabricated lies against God, objected
to Him, and sent fiery pages to the seat of holiness, setting down in
them what has burned the hearts of the chosen ones and those brought
near, so that he may awaken from the sleep of heedlessness and turn
toward the holy region, the seat of God, the All-Powerful, the Almighty,
the All-Knowing. Although he has committed a sin heavier than the weight
of the heavens and the earth, yet, since the oceans of grace have surged
in these days, if he repents and returns, perhaps God will forgive him
through His mercy, which has preceded all the worlds.

\subsection{789}\label{section-1284}

\begin{otherlanguage}{arabic}

قل ان يا ايّها الغافل لو تسمع من هذا العبد فاعمل بما نذكرك به لتكون من
الفائزين طهّر نفسك عن غير اللّه ثمّ قل بسم اللّه و باللّه و غسل بالمآء و بعد
فراغك قل يا الهى و محبوبی‌ كما طهّرت ظاهرى من هذا المآء طهّر قلبی‌ و فؤادى
عن النّفس و الهوى لاعرفك بنفسك و انقطع عن العالمين ثمّ ضع وجهك على التّراب
خضّعاً للّه ربّ الأرباب‌ قل يا الهى و سيّدى اسئلك باسمك الّذى به ظهر الزّلزال فى
قبائل الأرض و السّمآء و ارتفع الضّجيج بين ملأ الانشآء و اضطربت انفس
المريبين و حقايق الغافلين بأن تغفرلى جريراتى الّتى ما شهدت عين الابداع
شبهها فو‌عزّتك يا محبوبی‌ قد جئتك بذنب لا يعادله ذنوب من خلق بقولك الابدع
البديع

\end{otherlanguage}

Qul in yā ayyuhā al-ghāfil law tasmaʿ min hādhā al-ʿabd fa-iʿmal bi-mā
nadhkuruka bihi li-takūna min al-fāʾizīn. Ṭahhir nafsaka ʿan ghayr Allāh
thumma qul bi-sm Allāh wa-bi-Allāh wa-ghassil bi-al-māʾ wa-baʿda
farāghika qul yā ilāhī wa-maḥbūbī kamā ṭahharta ẓāhirī min hādhā al-māʾ
ṭahhir qalbī wa-fuʾādī ʿan al-nafs wa-al-hawā li-aʿrifaka bi-nafsika
wa-anqaṭiʿa ʿan al-ʿālamīn. Thumma ḍaʿ wajhaka ʿalā al-turāb khuḍḍaʿan
li-llāh rabb al-arbāb qul yā ilāhī wa-sayyidī asʾaluka bi-ismika alladhī
bihi ẓahara al-zilzāl fī qabāʾil al-arḍ wa-al-samāʾ wa-irtafaʿa al-ḍajīj
bayna malaʾ al-inshāʾ wa-iḍṭarabat anfus al-murībīn wa-ḥaqāʾiq
al-ghāfilīn bi-an taghfira lī jarīrātī allatī mā shahidat ʿayn al-ibdāʿ
shibhahā. Fa-wa-ʿizzatika yā maḥbūbī qad jiʾtuka bi-dhanb lā yuʿādiluhu
dhunūb man khuliqa bi-qawlika al-abdaʿ al-badīʿ.

Say: O heedless one, if you will listen to this servant, act upon what
We tell you, so that you may be among those who attain their goal.
Purify yourself of all but God, then say, ``In the name of God and by
God,'' and wash with water. When you have finished, say, ``O my God and
my Beloved, just as You have cleansed my outward being with this water,
cleanse my heart and my inmost being of self and desire, so that I may
know You through Your own Self and detach myself from all the worlds.''
Then place your face upon the dust in humility before God, the Lord of
lords, and say, ``O my God and my Master, I ask You by Your name,
through which the earthquake has appeared among the tribes of earth and
heaven, the clamor has risen among the company of creation, and the
souls of the doubters and the realities of the heedless have been
shaken, to forgive my transgressions, the like of which the eye of
creation has never witnessed. By Your might, O my Beloved, I have come
to You with a sin that cannot be equaled by the sins of those created
through Your most wondrous, incomparable word.

\subsection{790}\label{section-1285}

\begin{otherlanguage}{persian}

اذا فارحم على هذا العاصى ثمّ البسنى ثوب الغفران و طهّرنى من تعيّنات البشر و
وجّهنى الى منظرك الأكبر بحيث لايبقى فى صدرى الّا بدايع عرفانك و لا فى قلبی‌
الّا ظهورات حبّك و لا من وجهى الّا نضرة رضوانك و انّك انت المقتدر على ما
تشآء و انّك انت المهيمن المقتدر العزيز القدير يا الهى هذا عاصى قد سرع الى
مدين غفرانك و هذا المسكين قد تمسّك باذيال ردآء غنائك و استغنائك اى ربّ
لاتحرمنى عمّا عندك ثمّ اجعلنى خالصاً لوجهك و خاضعاً عند ظهورات انوار عزّ
فردانيّتك و انّك لم تزل كنت ذوالفضل القديم و تكون على كلّ شیء قدير

\end{otherlanguage}

Idhan fa-irḥam ʿalā hādhā al-ʿāṣī thumma albisnī thawb al-ghufrān
wa-ṭahhirnī min taʿayyunāt al-bashar wa-wajjihnī ilā manẓarika al-akbar
bi-ḥaythu lā yabqā fī ṣadrī illā badāʾiʿ ʿirfānika wa-lā fī qalbī illā
ẓuhūrāt ḥubbika wa-lā min wajhī illā naḍrat riḍwānika wa-innaka anta
al-Muqtadir ʿalā mā tashāʾ wa-innaka anta al-Muhaymin al-Muqtadir
al-ʿAzīz al-Qadīr. Yā ilāhī hādhā ʿāṣī qad saraʿa ilā madīn ghufrānika
wa-hādhā al-miskīn qad tamassaka bi-adhyāl ridāʾ ghināʾika
wa-istighnāʾika. Ay rabb lā taḥrimnī ʿammā ʿindaka thumma ijʿalnī
khāliṣan li-wajhika wa-khāḍiʿan ʿinda ẓuhūrāt anwār ʿizz fardāniyyatika
wa-innaka lam tazal kunta dhū al-faḍl al-qadīm wa-takūn ʿalā kull shayʾ
qadīr.

Then have mercy on this sinner, clothe me in the robe of forgiveness,
purify me of human limitations, and turn me toward Your Most Great
Countenance, so that nothing may remain in my breast but the wonders of
the knowledge of You, nor in my heart but the manifestations of Your
love, nor shine from my face but the radiance of Your good pleasure. You
are indeed powerful to do whatever You will, and You are the
All-Protecting, the All-Powerful, the Almighty, the Omnipotent. O my
God, this is a sinner who has hastened to the city of Your forgiveness,
and this poor one has clung to the hem of the robe of Your wealth and
self-sufficiency. O Lord, do not deprive me of what is with You. Then
make me wholly devoted to You and humble before the manifestations of
the lights of the glory of Your singleness. You have ever been the
Possessor of ancient grace, and You are powerful over all things.''

\subsection{791}\label{section-1286}

\begin{otherlanguage}{arabic}

و نسئل اللّه بذلك يفتح بصراك و يعرّفك ما لا عرفته من قبل و يعلّمك ما تجعل
به غنيّا عن دونه و عنده علم السّموات و الأرض و انّه قد كان على كلّ شیء محيط

\end{otherlanguage}

Wa-nasʾal Allāh bi-dhālika yaftaḥ baṣarāka wa-yuʿarrifaka mā lā
ʿaraftahu min qabl wa-yuʿallimaka mā tujʿal bihi ghaniyyan ʿan dūnihi
wa-ʿindahu ʿilm al-samāwāt wa-al-arḍ wa-innahu qad kāna ʿalā kull shayʾ
muḥīṭ.

And We ask God, through this, to open your eyes, make known to you what
you did not know before, and teach you that through which you may be
made independent of all but Him. With Him is the knowledge of the
heavens and the earth, and He indeed encompasses all things.

\bookmarksetup{startatroot}

\chapter{Lawh-i-Ibn-i-Dhi'b (Epistle to the Son of the
Wolf)}\label{lawh-i-ibn-i-dhib-epistle-to-the-son-of-the-wolf}

\section{BH00005}\label{bh00005}

\emph{Translated on 9 Qudrat, 183 BE using GPT 6.1 Sol}

\subsection{1}\label{section-1287}

\begin{otherlanguage}{arabic}

بسم اللّه الفرد الواحد المقتدر العلیم الحکیم

\end{otherlanguage}

Bi-smi Allāh al-Fard al-Wāḥid al-Muqtadir al-ʿAlīm al-Ḥakīm.

In the name of God, the Unique, the One, the All-Powerful, the
All-Knowing, the All-Wise.

\subsection{2}\label{section-1288}

\begin{otherlanguage}{arabic}

الحمد للّه الباقی بلا فنآء و الدّائم بلا زوال و القائم بلا انتقال المهیمن
بسلطانه و الظّاهر بآیاته و الباطن بأسراره الّذى بأمره ارتفعت رایة الکلمة
العلیا فى ناسوت الانشآء و نصب علم یفعل ما یشآء بین الوری هو الّذى اظهر
امره لهدایة خلقه و انزل آیاته اظهاراً لحجّته و برهانه و زیّن دیباج کتاب
الانسان بالبیان بقوله الرّحمن علّم القرآن خلق الانسان علّمه البیان لا اله
الّا هو الفرد الواحد المقتدر العزیز المنّان

\end{otherlanguage}

Al-ḥamdu li-llāh al-Bāqī bi-lā fanāʾ wa-l-Dāʾim bi-lā zawāl wa-l-Qāʾim
bi-lā intiqāl, al-Muhaymin bi-sulṭānihi wa-l-Ẓāhir bi-āyātihi wa-l-Bāṭin
bi-asrārihi, alladhī bi-amrihi irtafaʿat rāyatu al-kalimati al-ʿulyā fī
nāsūti al-inshāʾ wa-nuṣiba ʿalamu yafʿalu mā yashāʾ bayna al-warā. Huwa
alladhī aẓhara amrahu li-hidāyati khalqihi wa-anzala āyātihi iẓhāran
li-ḥujjatihi wa-burhānihi wa-zayyana dībāja kitābi al-insān bi-l-bayān
bi-qawlihi al-Raḥmānu ʿallama al-Qurʾāna khalaqa al-insāna ʿallamahu
al-bayān. Lā ilāha illā huwa al-Fard al-Wāḥid al-Muqtadir al-ʿAzīz
al-Mannān.

Praise be to God, who remains without perishing, endures without
decline, and abides without change; who rules through His sovereignty,
is manifest through His signs, and is hidden through His mysteries; at
whose command the banner of the Most Exalted Word was raised in the
world of creation, and the standard of ``He does whatever He wills'' was
planted among humankind. He is the One who revealed His Cause to guide
His creatures, sent down His verses to make His testimony and proof
manifest, and adorned the opening page of the book of humanity with
utterance through His words: ``The All-Merciful taught the Quran,
created the human being, and taught him utterance.'' There is no God but
Him, the Unique, the One, the All-Powerful, the Almighty, the Bestower.

\subsection{3}\label{section-1289}

\begin{otherlanguage}{arabic}

النّور السّاطع من افق سمآء العطآء و الصّلوة المشرقة من مطلع ارادة اللّه مالک
ملکوت الأسمآء على الواسطة الکبری و القلم الأعلی الّذی جعله اللّه مطلع
اسمائه الحسنی و مشرق صفاته العلیا و به اشرق نور التّوحید من افق العالم و
حکم التّفرید بین الأمم الّذین اقبلوا بوجوه نورآء الی الأفق الأعلی و
اعترفوا بما نطق به لسان البیان فى ملکوت العرفان الملک و الملکوت و العظمة
و الجبروت للّه المقتدر العزیز الفیاض

\end{otherlanguage}

Al-nūru al-sāṭiʿu min ufuqi samāʾi al-ʿaṭāʾ wa-l-ṣalātu al-mushriqatu
min maṭlaʿi irādati Allāh Māliki malakūti al-asmāʾ ʿalā al-wāsiṭati
al-kubrā wa-l-qalami al-aʿlā alladhī jaʿalahu Allāhu maṭlaʿa asmāʾihi
al-ḥusnā wa-mashriqa ṣifātihi al-ʿulyā wa-bihi ashraqa nūru al-tawḥīd
min ufuqi al-ʿālam wa-ḥakama al-tafrīd bayna al-umam alladhīna aqbalū
bi-wujūhin nawrāʾ ilā al-ufuqi al-aʿlā wa-iʿtarafū bi-mā naṭaqa bihi
lisānu al-bayān fī malakūti al-ʿirfān. Al-mulku wa-l-malakūtu
wa-l-ʿaẓamatu wa-l-jabarūtu li-llāh al-Muqtadir al-ʿAzīz al-Fayyāḍ.

May the light shining from the horizon of the heaven of bounty, and the
blessing radiating from the dawning place of the will of God, the Master
of the kingdom of names, rest upon the Supreme Intermediary and the Most
Exalted Pen, which God has made the dawning place of His most beautiful
names and the dayspring of His most exalted attributes. Through it the
light of divine unity has shone from the horizon of the world, and
divine oneness has prevailed among the nations who have turned with
radiant faces toward the Most Exalted Horizon and acknowledged what the
Tongue of Utterance has proclaimed in the kingdom of knowledge:
``Sovereignty and dominion, grandeur and might belong to God, the
All-Powerful, the Almighty, the All-Bountiful.''

\subsection{4}\label{section-1290}

\begin{otherlanguage}{arabic}

یا ایّها العالم الجلیل اسمع ندآء المظلوم انّه ینصحک لوجه اللّه و یعظک بما
یقرّبک الیه فی کلّ الأحوال انّه هو الغنیّ المتعال اعلم انّ الآذان خلقت لاصغآء
النّدآء فی هذا الیوم الّذی کان مذکوراً فى الکتب و الزّبر و الألواح فی اوّل
الأمر طهّر نفسک بمآء الانقطاع و زیّن رأسک باکلیل التّقوی و هیکلک بطراز
التّوکّل علی اللّه ثمّ قم عن مقامک مقبلاً الی البیت الأعظم مطاف من فی العالم
من لدن مالک القدم و قل

\end{otherlanguage}

Yā ayyuhā al-ʿālimu al-jalīl, ismaʿ nidāʾa al-Maẓlūm. Innahu yanṣaḥuka
li-wajhi Allāh wa-yaʿiẓuka bi-mā yuqarribuka ilayhi fī kulli al-aḥwāl.
Innahu huwa al-Ghanī al-Mutaʿāl. Iʿlam anna al-ādhāna khuliqat
li-iṣghāʾi al-nidāʾ fī hādhā al-yawmi alladhī kāna madhkūran fī
al-kutubi wa-l-zuburi wa-l-alwāḥ. Fī awwali al-amr ṭahhir nafsaka
bi-māʾi al-inqiṭāʿ wa-zayyin raʾsaka bi-iklīli al-taqwā wa-haykalaka
bi-ṭirāzi al-tawakkuli ʿalā Allāh. Thumma qum ʿan maqāmika muqbilan ilā
al-bayti al-aʿẓam maṭāfi man fī al-ʿālam min ladun Māliki al-qidam
wa-qul.

Distinguished scholar, listen to the call of the Wronged One. He
counsels you for the sake of God and exhorts you to what will draw you
near to Him in every circumstance. He is indeed the Self-Sufficient, the
Most Exalted. Know that ears were created to hear the call on this Day,
which was mentioned in the Books, the Scriptures, and the Tablets.
First, cleanse yourself with the water of detachment, adorn your head
with the crown of piety, and your body with the vesture of trust in God.
Then rise from your place, turning toward the Most Great House,
appointed by the Lord of Eternity as the place around which all who are
in the world should circle, and say:

\subsection{5}\label{section-1291}

\begin{otherlanguage}{arabic}

الهی الهی و مقصودی و معبودی و سیّدی و سندی و غایة املی و رجائی ترانی
مقبلاً الیک و متمسّکاً بحبل جودک و متشبّثاً بذیل عطائک و معترفاً بتقدیس نفسک و
تنزیه ذاتک مقرّاً بوحدانیّتک و فردانیّتک اشهد انّک انت الفرد الواحد الأحد
الصّمد ما اتّخذت لنفسک شریکاً فی الملک و لا نظیراً فی الأرض قد شهدت الکائنات
بما شهد به لسان عظمتک قبلها انّک انت اللّه لا اله الّا انت لم تزل کنت مقدّساً
عن ذکر عبادک و متعالیاً عن وصف خلقک ای ربّ تری الجاهل قصد بحر علمک و
العطشان کوثر بیانک و الذّلیل خبآء عزّک و الفقیر کنز غنائک و السّائل مشرق
حکمتک و الضّعیف مطلع قدرتک و المسکین سمآء کرمک و الکلیل ملکوت ذکرک

\end{otherlanguage}

Ilāhī Ilāhī wa-maqṣūdī wa-maʿbūdī wa-sayyidī wa-sanadī wa-ghāyata amalī
wa-rajāʾī, tarānī muqbilan ilayka wa-mutamassikan bi-ḥabli jūdika
wa-mutashabbithan bi-dhayli ʿaṭāʾika wa-muʿtarifan bi-taqdīsi nafsika
wa-tanzīhi dhātika muqirran bi-waḥdāniyyatika wa-fardāniyyatik. Ashhadu
annaka anta al-Fard al-Wāḥid al-Aḥad al-Ṣamad. Mā ittakhadhta li-nafsika
sharīkan fī al-mulki wa-lā naẓīran fī al-arḍ. Qad shahidati al-kāʾinātu
bi-mā shahida bihi lisānu ʿaẓamatika qablahā annaka anta Allāhu lā ilāha
illā anta. Lam tazal kunta muqaddasan ʿan dhikri ʿibādika wa-mutaʿāliyan
ʿan waṣfi khalqik. Ay Rabb, tarā al-jāhila qaṣada baḥra ʿilmika
wa-l-ʿaṭshāna kawthara bayānika wa-l-dhalīla khibāʾa ʿizzika wa-l-faqīra
kanza ghināʾika wa-l-sāʾila mashriqa ḥikmatika wa-l-ḍaʿīfa maṭlaʿa
qudratika wa-l-miskīna samāʾa karamika wa-l-kalīla malakūta dhikrik.

My God, my God, my Desire, my Adored One, my Master, my Support, the
ultimate object of my hope and expectation! You see me turning toward
You, holding fast to the cord of Your generosity, clinging to the hem of
Your bounty, acknowledging the sanctity of Your Being and the
transcendence of Your Essence, and affirming Your oneness and
uniqueness. I testify that You are the Unique, the One, the Single, the
All-Sufficing. You have taken for Yourself no partner in sovereignty and
no equal on earth. All created things have testified to what the Tongue
of Your grandeur testified before them: that You are God; there is no
God but You. You have always been sanctified above the mention of Your
servants and exalted beyond the description of Your creatures. O Lord,
You see the ignorant one seeking the ocean of Your knowledge, the
thirsty one the river of Your utterance, the abased one the pavilion of
Your glory, the poor one the treasury of Your riches, the petitioner the
dayspring of Your wisdom, the weak one the dawning place of Your power,
the destitute one the heaven of Your generosity, and the weary one the
kingdom of Your remembrance.

\subsection{6}\label{section-1292}

\begin{otherlanguage}{arabic}

اشهد یا الهی و سلطانی بأنّک خلقتنی لذکرک و ثنائک و نصرة امرک و انّى نصرت
اعدائک الّذین نقضوا عهدک و نبذوا کتابک و کفروا بک و بآیاتک آه آه من غفلتی
و خجلتی و خطیئتی و جریرتی الّتی منعتنی عن الورود فی طمطام بحر احدیّتک و
قمقام یمّ رحمتک فآه آه ثمّ آه آه من سوء حالی و کبر عصیانی قد اظهرتنی یا
الهی لاعلآء کلمتک و اظهار امرک ولکن غفلتی منعتنی و احاطت بی بحیث قمت على
محو آثارک و سفک دمآء اولیائک و مطالع آیاتک و مشارق وحیک و مخازن اسرارک

\end{otherlanguage}

Ashhadu yā Ilāhī wa-Sulṭānī bi-annaka khalaqtanī li-dhikrika
wa-thanāʾika wa-nuṣrati amrika wa-annī naṣartu aʿdāʾaka alladhīna naqaḍū
ʿahdaka wa-nabadhū kitābaka wa-kafarū bika wa-bi-āyātik. Āh āh min
ghaflatī wa-khajlatī wa-khaṭīʾatī wa-jarīratī allatī manaʿatnī ʿan
al-wurūdi fī ṭamṭāmi baḥri aḥadiyyatika wa-qamqāmi yammi raḥmatik. Fa-āh
āh thumma āh āh min sūʾi ḥālī wa-kibari ʿiṣyānī. Qad aẓhartanī yā Ilāhī
li-iʿlāʾi kalimatika wa-iẓhāri amrika wa-lākin ghaflatī manaʿatnī
wa-aḥāṭat bī bi-ḥaythu qumtu ʿalā maḥwi āthārika wa-safki dimāʾi
awliyāʾika wa-maṭāliʿi āyātika wa-mashāriqi waḥyika wa-makhāzini
asrārik.

I testify, my God and my Sovereign, that You created me to remember and
praise You and to aid Your Cause, yet I aided Your enemies, who broke
Your covenant, cast aside Your Book, and disbelieved in You and Your
signs. Alas, alas for my heedlessness, my shame, my sin, and my
wrongdoing, which kept me from entering the vast ocean of Your oneness
and the fathomless sea of Your mercy! Alas, alas, and again alas, alas
for my wretched condition and the enormity of my rebellion! You brought
me into being, my God, to exalt Your Word and make Your Cause manifest,
but my heedlessness hindered and encompassed me, so that I arose to
efface Your signs and shed the blood of Your loved ones, the dawning
places of Your verses, the daysprings of Your revelation, and the
repositories of Your mysteries.

\subsection{7}\label{section-1293}

\begin{otherlanguage}{arabic}

ای ربّ ای ربّ ای ربّ ای ربّ ای ربّ ای ربّ ای ربّ ای ربّ ای ربّ اشهد بظلمی سقطت
اثمار سدرة عدلک و بنار عصیانی احترقت افئدة المقرّبین من خلقک و ذابت اکباد
المخلصین من عبادک فآه آه من شقوتی فآه آه من ظلمی فآه آه من بعدی و غفلتی
و جهلی و ذلّتی و اعراضی و اعتراضی کم من ایّام فیها امرت عبادک و اولیائک
علی حفظی و انّی امرتهم بضرّک و ضرّ امنائک و کم من لیال فیها ذکرتنی بفضلک و
دللتنی الی صراطک و انّی اعرضت عنک و عن آیاتک وعزّتک یا امل الموحّدین و رجآء
افئدة المنقطعین لا اجد لنفسی دونک معیناً و لا سواک سلطاناً و لا ملجأً و لا
ملاذاً فآه آه اعراضی احرق ستر عصمتی و اعتراضی شقّ حجاب حرمتی یا لیت کنت
تحت اطباق التّراب و ما ظهر سوء اعمالی بین عبادک ای ربّ تری العاصی اقبل الی
مطلع عفوک و عطائک و جبل الظّلم اراد سمآء رحمتک و غفرانک فآه آه جریراتی
العظمی منعتنی عن التّقرّب الی بساط رحمتک و خطیئاتی الکبری ابعدتنی عن ساحة
قربک انا الّذى فرّطت فی جنبک و نقضت عهدک و میثاقک و ارتکبت ما ناح به سکّان
مدائن عدلک و مطالع فضلک فی بلادک اشهد یا الهی انّی ترکت اوامرک و اخذت
اوامر نفسی و نبذت احکام کتابک و اخذت کتاب هوای فآه آه کلّما زادت شقوتی
زاد حلمک و کلّما اشتعلت نار عصیانی سترها عفوک و فضلک وعزّتک یا مقصود
العالم و محبوب الأمم صبرک غرّنی و اصطبارک شجّعنی تری یا الهی عبراتی من
خجلتی و زفراتی من غفلتی وعظمتک لا اجد لنفسی مقرّاً الّا ظلّ بساط کرمک و لا
مهرباً الّا تحت قباب رحمتک ترانی فی بحر الیأس و القنوط بعدما اسمعتنی کلمة
لا تقنطوا وعزّتک ظلمی قطع حبل املی و عصیانی سوّد وجهی امام کرسیّ عدلک ای ربّ
تری المیّت مطروحاً لدی باب عطائک و یستحیی ان یطلب کوثر عفوک من ید فضلک قد
اعطیتنی لساناً لذکرک و ثنائک و انّه نطق بما ذابت به اکباد المقرّبین من
اصفیائک و احترقت افئدة المخلصین من اهل حظائر قدسک و اعطیتنی بصراً لمشاهدة
آثارک و ملاحظة آیاتک و مظاهر صنعک و انّی نبذت ارادتک و عملت ما ناح به
المخلصون من خلقک و المنقطعون من عبادک و اعطیتنی سمعاً لأسمع به ذکرک و
ثنائک و ما انزلته من سمآء کرمک و هوآء ارادتک فآه آه انّی ترکت امرک و امرت
عبادک بسبّ امنائک و اولیائک و عملت امام کرسیّ عدلک ما ارتفعت به زفرات
الموحّدین و المخلصین من اهل مملکتک لم ادر یا الهی ایّ عصیانی اذکره تلقآء
امواج بحر جودک و ایّ خطأی انطق به عند تجلّیات انوار شموس مواهبک و الطافک

\end{otherlanguage}

Ay Rabb, ay Rabb, ay Rabb, ay Rabb, ay Rabb, ay Rabb, ay Rabb, ay Rabb,
ay Rabb! Ashhadu bi-ẓulmī saqaṭat athmāru sidrati ʿadlika wa-bi-nāri
ʿiṣyānī iḥtaraqat afʾidatu al-muqarrabīna min khalqika wa-dhābat akbādu
al-mukhliṣīna min ʿibādik. Fa-āh āh min shiqwatī, fa-āh āh min ẓulmī,
fa-āh āh min buʿdī wa-ghaflatī wa-jahlī wa-dhillatī wa-iʿrāḍī
wa-iʿtirāḍī. Kam min ayyāmin fīhā amarta ʿibādaka wa-awliyāʾaka ʿalā
ḥifẓī wa-innī amartuhum bi-ḍurrika wa-ḍurri umanāʾik. Wa-kam min layālin
fīhā dhakartanī bi-faḍlika wa-dalaltanī ilā ṣirāṭika wa-innī aʿraḍtu
ʿanka wa-ʿan āyātik. Wa-ʿizzatika yā Amala al-muwaḥḥidīn wa-Rajāʾa
afʾidati al-munqaṭiʿīn, lā ajidu li-nafsī dūnaka muʿīnan wa-lā siwāka
sulṭānan wa-lā maljaʾan wa-lā malādhan. Fa-āh āh, iʿrāḍī aḥraqa sitra
ʿiṣmatī wa-iʿtirāḍī shaqqa ḥijāba ḥurmatī. Yā layta kuntu taḥta aṭbāqi
al-turāb wa-mā ẓahara sūʾu aʿmālī bayna ʿibādik. Ay Rabb, tarā al-ʿāṣiya
aqbala ilā maṭlaʿi ʿafwika wa-ʿaṭāʾika wa-jabala al-ẓulmi arāda samāʾa
raḥmatika wa-ghufrānik. Fa-āh āh, jarīrātī al-ʿuẓmā manaʿatnī ʿan
al-taqarrubi ilā bisāṭi raḥmatika wa-khaṭīʾātī al-kubrā abʿadatnī ʿan
sāḥati qurbik. Anā alladhī farraṭtu fī janbika wa-naqaḍtu ʿahdaka
wa-mīthāqaka wa-irtakabtu mā nāḥa bihi sukkānu madāʾini ʿadlika
wa-maṭāliʿu faḍlika fī bilādik. Ashhadu yā Ilāhī annī taraktu awāmiraka
wa-akhadhtu awāmira nafsī wa-nabadhtu aḥkāma kitābika wa-akhadhtu kitāba
hawāy. Fa-āh āh, kullamā zādat shiqwatī zāda ḥilmuka wa-kullamā
ishtaʿalat nāru ʿiṣyānī satarahā ʿafwuka wa-faḍluk. Wa-ʿizzatika yā
Maqṣūda al-ʿālam wa-Maḥbūba al-umam, ṣabruka gharranī wa-iṣṭibāruka
shajjaʿanī. Tarā yā Ilāhī ʿabarātī min khajlatī wa-zafarātī min
ghaflatī. Wa-ʿaẓamatika lā ajidu li-nafsī maqarran illā ẓilla bisāṭi
karamika wa-lā mahraban illā taḥta qibābi raḥmatik. Tarānī fī baḥri
al-yaʾsi wa-l-qunūṭ baʿdamā asmaʿtanī kalimata lā taqnaṭū. Wa-ʿizzatika
ẓulmī qaṭaʿa ḥabla amalī wa-ʿiṣyānī sawwada wajhī amāma kursiyyi ʿadlik.
Ay Rabb, tarā al-mayyita maṭrūḥan ladā bābi ʿaṭāʾika wa-yastaḥyī an
yaṭluba kawthara ʿafwika min yadi faḍlik. Qad aʿṭaytanī lisānan
li-dhikrika wa-thanāʾika wa-innahu naṭaqa bi-mā dhābat bihi akbādu
al-muqarrabīna min aṣfiyāʾika wa-iḥtaraqat afʾidatu al-mukhliṣīna min
ahli ḥaẓāʾiri qudsik. Wa-aʿṭaytanī baṣaran li-mushāhadati āthārika
wa-mulāḥaẓati āyātika wa-maẓāhiri ṣunʿika wa-innī nabadhtu irādataka
wa-ʿamiltu mā nāḥa bihi al-mukhliṣūna min khalqika wa-l-munqaṭiʿūna min
ʿibādik. Wa-aʿṭaytanī samʿan li-asmaʿa bihi dhikraka wa-thanāʾaka wa-mā
anzaltahu min samāʾi karamika wa-hawāʾi irādatik. Fa-āh āh, innī taraktu
amraka wa-amartu ʿibādaka bi-sabbi umanāʾika wa-awliyāʾika wa-ʿamiltu
amāma kursiyyi ʿadlika mā irtafaʿat bihi zafarātu al-muwaḥḥidīna
wa-l-mukhliṣīna min ahli mamlakatik. Lam adri yā Ilāhī ayya ʿiṣyānī
adhkuruhu tilqāʾa amwāji baḥri jūdika wa-ayya khaṭaʾī anṭiqu bihi ʿinda
tajalliyāti anwāri shumūsi mawāhibika wa-alṭāfik.

O Lord, O Lord, O Lord, O Lord, O Lord, O Lord, O Lord, O Lord, O Lord!
I testify that through my injustice the fruits of the tree of Your
justice have fallen, and through the fire of my rebellion the hearts of
those among Your creatures who are near to You have burned, and the
innermost beings of Your sincere servants have melted. Alas, alas for my
wretchedness! Alas, alas for my injustice! Alas, alas for my remoteness,
my heedlessness, my ignorance, my abasement, my turning away, and my
opposition! How many days You commanded Your servants and Your loved
ones to protect me, while I commanded them to harm You and Your trusted
ones! How many nights You remembered me through Your grace and guided me
to Your path, while I turned away from You and Your signs! By Your
glory, Hope of those who affirm Your unity and Desire of the hearts of
the detached, I find no helper but You, and no sovereign, refuge, or
shelter other than You. Alas, alas! My turning away has burned the
covering of my protection, and my opposition has torn the veil of my
dignity. If only I had been beneath the layers of dust and the evil of
my deeds had not appeared among Your servants! O Lord, You see the
disobedient one turning toward the dawning place of Your pardon and
bounty, and the mountain of injustice seeking the heaven of Your mercy
and forgiveness. Alas, alas! My grievous wrongdoings have kept me from
approaching the court of Your mercy, and my great sins have removed me
far from the precincts of Your presence. I am the one who failed in my
duty toward You, broke Your covenant and Your pledge, and committed what
caused the inhabitants of the cities of Your justice and the dawning
places of Your grace in Your lands to lament. I testify, my God, that I
abandoned Your commands and followed the commands of my own self, cast
aside the laws of Your Book, and took up the book of my own desire.
Alas, alas! The greater my wretchedness grew, the greater Your
forbearance became, and whenever the fire of my rebellion blazed, Your
pardon and grace concealed it. By Your glory, Desire of the world and
Beloved of the nations, Your patience deluded me and Your endurance
emboldened me. You see, my God, the tears I shed from shame and the
sighs I heave over my heedlessness. By Your grandeur, I find no resting
place except the shelter of the court of Your generosity, and no escape
except beneath the domes of Your mercy. You see me in the ocean of
hopelessness and despair after You have caused me to hear the words ``Do
not despair.'' By Your glory, my injustice has severed the cord of my
hope, and my rebellion has blackened my face before the throne of Your
justice. O Lord, You see the dead one cast down at the door of Your
bounty, ashamed to ask for the river of Your pardon from the hand of
Your grace. You gave me a tongue to remember and praise You, yet it
uttered what melted the innermost beings of those among Your chosen ones
who are near to You and burned the hearts of the sincere among the
inhabitants of Your holy courts. You gave me sight to behold Your works
and contemplate Your signs and the manifestations of Your handiwork, yet
I cast aside Your will and did what caused the sincere among Your
creatures and the detached among Your servants to lament. You gave me
hearing with which to hear Your remembrance and praise and what You sent
down from the heaven of Your generosity and the atmosphere of Your will.
Alas, alas! I abandoned Your command, ordered Your servants to revile
Your trusted ones and Your loved ones, and acted before the throne of
Your justice in a way that caused the sighs of those who affirm Your
unity and the sincere among the people of Your kingdom to rise. I do not
know, my God, which act of my rebellion to recount before the waves of
the ocean of Your generosity, or which of my errors to speak of before
the effulgent lights of the suns of Your gifts and favors.

\subsection{8}\label{section-1294}

\begin{otherlanguage}{persian}

اسألک فى هذا الحین بأسرار کتابک و ما کان مکنوناً فی علمک و باللّآلئ
المستورة فی اصداف عمّان رحمتک ان تجعلنی من الّذین ذکرتهم فی کتابک و وصفتهم
فی الواحک هل قدّرت لی یا الهی بعد هذا الحزن من سرور و بعد هذا القبض من
بسط و بعد هذا العسر من یسر فآه آه قد جعلت المنابر لذکرک و ارتفاع کلمتک و
اظهار امرک و انّی ارتقیت الیها لاعلآء نقض عهدک و القیت علی العباد ما ناح
به اهل سرادق عظمتک و سکّان مدائن علمک کم من اوقات انزلت فیها مائدة بیانک
من سمآء عطائک و انّی کفرت بها و کم من احیان دعوتنی فیها الی فرات رحمتک و
انّی اعرضت عنه بما اتّبعت النّفس و الهوی وعزّتک لم ادر من ایّ ذنب استغفرک و
اتوب الیک و من ایّ ظلم ارجع الی بساط جودک و ساحة کرمک قد بلغت جریراتی و
خطیئاتی مقاماً عجز المحصون عن احصائها و المحرّرون عن تحریرها اسألک یا مبدّل
الظّلمة بالنّور و مظهر الأسرار فی الطّور ایّدنی فی کلّ الأحوال علی التّوکّل
علیک و تفویض الأمور الیک ثمّ اجعلنی یا الهی راضیاً بما رقم من قلم قضائک و
یراعة تقدیرک انّک انت المقتدر علی ما تشآء و فی قبضتک زمام من فی السّموات و
الأرضین لا اله الّا انت العلیم الحکیم

\end{otherlanguage}

Asʾaluka fī hādhā al-ḥīn bi-asrāri kitābika wa-mā kāna maknūnan fī
ʿilmika wa-bi-l-laʾāliʾi al-mastūrati fī aṣdāfi ʿummāni raḥmatika an
tajʿalanī mina alladhīna dhakartahum fī kitābika wa-waṣaftahum fī
alwāḥik. Hal qaddarta lī yā Ilāhī baʿda hādhā al-ḥuzni min surūr
wa-baʿda hādhā al-qabḍi min basṭ wa-baʿda hādhā al-ʿusri min yusr? Fa-āh
āh, qad jaʿalta al-manābira li-dhikrika wa-irtifāʿi kalimatika wa-iẓhāri
amrika wa-innī irtaqaytu ilayhā li-iʿlāʾi naqḍi ʿahdika wa-alqaytu ʿalā
al-ʿibādi mā nāḥa bihi ahlu surādiqi ʿaẓamatika wa-sukkānu madāʾini
ʿilmik. Kam min awqātin anzalta fīhā māʾidata bayānika min samāʾi
ʿaṭāʾika wa-innī kafartu bihā wa-kam min aḥyānin daʿawtanī fīhā ilā
furāti raḥmatika wa-innī aʿraḍtu ʿanhu bi-mā ittabaʿtu al-nafsa
wa-l-hawā. Wa-ʿizzatika lam adri min ayyi dhanbin astaghfiruka wa-atūbu
ilayka wa-min ayyi ẓulmin arjiʿu ilā bisāṭi jūdika wa-sāḥati karamik.
Qad balaghat jarīrātī wa-khaṭīʾātī maqāman ʿajaza al-muḥṣūna ʿan
iḥṣāʾihā wa-l-muḥarrirūna ʿan taḥrīrihā. Asʾaluka yā Mubaddila
al-ẓulmati bi-l-nūr wa-Muẓhira al-asrāri fī al-Ṭūr, ayyidnī fī kulli
al-aḥwāli ʿalā al-tawakkuli ʿalayka wa-tafwīḍi al-umūri ilayk. Thumma
ijʿalnī yā Ilāhī rāḍiyan bi-mā ruqima min qalami qaḍāʾika wa-yarāʿati
taqdīrik. Innaka anta al-Muqtadiru ʿalā mā tashāʾ wa-fī qabḍatika zimāmu
man fī al-samāwāti wa-l-araḍīn. Lā ilāha illā anta al-ʿAlīm al-Ḥakīm.

I ask You at this moment, by the mysteries of Your Book, by what lies
hidden in Your knowledge, and by the pearls concealed in the shells of
the ocean of Your mercy, to make me one of those whom You have mentioned
in Your Book and described in Your Tablets. Have You ordained for me, my
God, any joy after this sorrow, any relief after this anguish, any ease
after this hardship? Alas, alas! You established pulpits for Your
remembrance, the exaltation of Your Word, and the proclamation of Your
Cause, yet I ascended them to promote the breaking of Your covenant and
delivered to Your servants what caused the people of the pavilion of
Your grandeur and the inhabitants of the cities of Your knowledge to
lament. How often You sent down the feast of Your utterance from the
heaven of Your bounty, while I rejected it! How often You summoned me to
the sweet waters of Your mercy, while I turned away from them through
following self and desire! By Your glory, I do not know for which sin to
seek Your forgiveness and turn to You in repentance, or from which
injustice to return to the court of Your generosity and the precincts of
Your bounty. My wrongdoings and sins have reached such a degree that
those who count are unable to number them, and those who write are
unable to record them. I ask You, O You who replace darkness with light
and reveal mysteries on the Mount, to strengthen me in every
circumstance to trust in You and entrust all affairs to You. Then make
me, my God, content with what has been inscribed by the pen of Your
decree and the reed of Your ordaining. You indeed have power to do
whatever You will, and in Your grasp are the reins of all who are in the
heavens and the earth. There is no God but You, the All-Knowing, the
All-Wise.

\subsection{9}\label{section-1295}

\begin{otherlanguage}{arabic}

یا شیخ اعلم انّ مفتریات العباد و اعراضهم و اعتراضهم لا تضرّ من تمسّک بحبل
العنایة و تشبّث بأذیال رحمة مالک البریّة لعمر اللّه انّ البهآء ما نطق عن
الهوی قد انطقه الّذی انطق الأشیآء بذکره و ثنائه لا اله الّا هو الفرد
الواحد المقتدر المختار

\end{otherlanguage}

Yā Shaykh, iʿlam anna muftarayāti al-ʿibādi wa-iʿrāḍahum wa-iʿtirāḍahum
lā taḍurru man tamassaka bi-ḥabli al-ʿināyati wa-tashabbatha bi-adhyāli
raḥmati Māliki al-bariyyah. La-ʿamru Allāh, inna al-Bahāʾa mā naṭaqa ʿan
al-hawā. Qad anṭaqahu alladhī anṭaqa al-ashyāʾa bi-dhikrihi wa-thanāʾih.
Lā ilāha illā huwa al-Fard al-Wāḥid al-Muqtadir al-Mukhtār.

O Shaykh, know that people's slanders, their turning away, and their
opposition cannot harm one who holds fast to the cord of providence and
clings to the hem of the mercy of the Lord of creation. By the life of
God, Bahá has not spoken from desire. He has been made to speak by the
One who caused all things to voice His remembrance and praise. There is
no God but Him, the Unique, the One, the All-Powerful, the
Unconstrained.

\subsection{10}\label{section-1296}

\begin{otherlanguage}{persian}

صاحبان ابصار حدیده و آذان واعیه و قلوب منیره و صدور منشرحه صدق را از کذب
بشناسند و تمیز دهند این مناجات را که از لسان مظلوم جاری شده قرائت نمائید
و بقلب فارغ و سمع طاهر مقدّس در آن تفکّر فرمائید شاید نفحات انقطاع را
بیابید و بر خود و عباد رحم کنید

\end{otherlanguage}

Ṣāḥibān-i abṣār-i ḥadīde va ādhān-i vāʿiye va qulūb-i munīre va ṣudūr-i
munshariḥe ṣidq rā az kidhb bishināsand va tamyīz dihand. Īn munājāt rā
kih az lisān-i Maẓlūm jārī shudih qirāʾat namāyīd va bi-qalb-i fārigh va
samʿ-i ṭāhir-i muqaddas dar ān tafakkur farmāyīd. Shāyad nafaḥāt-i
inqiṭāʿ rā biyābīd va bar khud va ʿibād raḥm kunīd.

Those with keen sight, attentive ears, radiant hearts, and open breasts
recognize and distinguish truth from falsehood. Read this prayer that
has flowed from the tongue of the Wronged One, and reflect on it with an
unburdened heart and a pure, sanctified hearing. Perhaps you may
perceive the breezes of detachment and have mercy on yourselves and on
God's servants.

\subsection{11}\label{section-1297}

\begin{otherlanguage}{arabic}

الها معبودا مقصودا کریما رحیما جانها از تو و اقتدارها در قبضهٔ قدرت تو هر
که را بلند کنی از ملک بگذرد و بمقام و رفعناه مقاماً علیّاً رسد و هر که را
بیندازی از خاک پستتر بلکه هیچ از او بهتر پروردگارا با تباه‌کاری و
گناهکاری و عدم پرهیزکاری مقعد صدق میطلبیم و لقاء اولیائت را میجوئیم امر
امر تو و حکم آن تو و عالم قدرت زیر فرمان تو هر چه کنی عدل صرفست بل فضل
محض یک تجلّی از تجلّیات اسم رحمانت رسم عصیان را از جهان براندازد و محو
نماید و یک نسیم از نسائم یوم ظهورت عالم را بخلعت تازه مزیّن نماید ای
توانا ناتوانان را توانائی بخش و مردگان را زندگی عطا فرما شاید ترا بیابند
و بدریای آگاهیت راه یابند و بر امرت مستقیم مانند اگر از لغات مختلفهٔ عالم
عرف ثنای تو متضوّع شود همه محبوب جان و مقصود روان چه تازی چه پارسی اگر از
آن محروم ماند قابل ذکر نه چه الفاظ چه معانی ای پروردگار از تو میطلبم کل
را راه نمائی و هدایت فرمائی توئی قادر و توانا و عالم و بینا

\end{otherlanguage}

Ilāhā, maʿbūdā, maqṣūdā, karīmā, raḥīmā! Jānhā az to va eqtedārhā dar
qabḍe-ye qodrat-e to. Har ke rā boland konī az malak bogẕarad va
be-maqām-e «wa rafaʿnāhu maqāman ʿaliyyan» rasad, va har ke rā biyandāzī
az khāk pasttar balke hīch az ū behtar. Parvardegārā, bā tabāhkārī va
gonāhkārī va ʿadam-e parhīzkārī maqʿad-e ṣedq mīṭalabīm va leqā-ye
owlīyāʾat rā mījūʾīm. Amr amr-e to va ḥokm ān-e to va ʿālam-e qodrat
zīr-e farmān-e to. Har che konī ʿadl-e ṣerf ast bal faḍl-e maḥḍ. Yek
tajallī az tajallīyāt-e esm-e Raḥmānat rasm-e ʿeṣyān rā az jahān
barandāzad va maḥv namāyad va yek nasīm az nasāʾem-e yowm-e ẓohūrat
ʿālam rā be-khelʿat-e tāze mozayyan namāyad. Ey tavānā, nātavānān rā
tavānāʾī bakhsh va mordegān rā zendegī ʿaṭā farmā, shāyad to rā biyāband
va be-daryā-ye āgāhīyat rāh yāband va bar amrat mostaqīm mānand. Agar az
loghāt-e mokhtalefe-ye ʿālam ʿarf-e thanā-ye to motaḍavveʿ shavad hame
maḥbūb-e jān va maqṣūd-e ravān, che tāzī che pārsī. Agar az ān maḥrūm
mānad qābel-e dhekr na, che alfāẓ che maʿānī. Ey Parvardegār, az to
mīṭalabam koll rā rāh namāʾī va hedāyat farmāʾī. Toʾī qāder va tavānā va
ʿālem va bīnā.

O God, O Adored One, O Desired One, O Generous One, O Merciful One!
Souls come from You, and all power lies within the grasp of Your might.
Whoever You exalt rises above the angels and attains the station of ``We
raised him to a lofty station,'' and whoever You cast down is lower than
dust; indeed, nothingness is better than him. O Lord, despite our
wrongdoing, sinfulness, and lack of piety, we seek a seat of truth and
long to meet Your chosen ones. The command is Yours, the judgment is
Yours, and the realm of power is subject to Your command. Whatever You
do is pure justice---indeed, sheer grace. A single revelation of Your
name, the All-Merciful, would efface and obliterate every trace of
rebellion from the world, and a single breeze from the breezes of the
day of Your revelation would adorn the world with a new robe. O Mighty
One, grant strength to the powerless and bestow life upon the dead, so
that they may find You, gain access to the ocean of Your knowledge, and
remain steadfast in Your Cause. If the fragrance of Your praise diffuses
from the various languages of the world, all are beloved of the soul and
desired by the spirit, whether Arabic or Persian. If they are deprived
of it, neither words nor meanings are worthy of mention. O Lord, I ask
You to show everyone the way and guide them. You are the Powerful, the
Mighty, the Knowing, the Seeing.

\subsection{12}\label{section-1298}

\begin{otherlanguage}{persian}

نسأل اللّه ان یؤیّدک علی العدل و الانصاف و یعرّفک ما کان مستوراً عن العیون و
الأبصار انّه هو العزیز المختار استدعا آنکه در آنچه ظاهر شده تفکّر نمایند و
بعدل و انصاف تکلّم فرمایند شاید تجلّیات انوار آفتاب صدق و صفا پرتو افکند و
از تاریکی نادانی نجات بخشد و عالم را بنور دانائی روشن فرماید این مظلوم
مدارس نرفته مباحث ندیده لعمری انّی ما اظهرت نفسی بل اللّه اظهرنی کیف اراد
در لوح حضرت سلطان ایّده اللّه تبارک و تعالى این کلمات از لسان مظلوم جاری

\end{otherlanguage}

Nasʾalu Allāha an yuʾayyidaka ʿalā al-ʿadli wa-l-inṣāfi wa yuʿarrifaka
mā kāna mastūran ʿan al-ʿuyūni wa-l-abṣār. Innahu huwa al-ʿAzīzu
al-Mukhtār. Estedʿā ānke dar ānche ẓāher shode tafakkor namāyand va
be-ʿadl va enṣāf takallom farmāyand, shāyad tajallīyāt-e anvār-e āftāb-e
ṣedq va ṣafā partow afkanad va az tārīkī-ye nādānī najāt bakhshad va
ʿālam rā be-nūr-e dānāʾī rowshan farmāyad. Īn Maẓlūm madāres narafte,
mabāḥeth nadīde. La-ʿamrī innī mā aẓhartu nafsī bal Allāhu aẓharanī
kayfa arād. Dar lowḥ-e Ḥaḍrat-e Solṭān, ayyadahu Allāhu tabāraka wa
taʿālā, īn kalemāt az lesān-e Maẓlūm jārī:

We ask God to assist you in justice and fairness and to acquaint you
with what has been concealed from eyes and sight. He is indeed the
Almighty, the Unconstrained. The request is that you reflect on what has
appeared and speak with justice and fairness, so that the radiance of
the sun of truthfulness and purity may shine forth, deliver people from
the darkness of ignorance, and illumine the world with the light of
knowledge. This Wronged One has attended no schools and witnessed no
scholarly debates. By My life, I did not reveal Myself; rather, God
revealed Me as He willed. In the Tablet to His Majesty the Shah---may
God, blessed and exalted, assist him---these words flowed from the
tongue of the Wronged One:

\subsection{13}\label{section-1299}

\begin{otherlanguage}{persian}

یا سلطان انّی کنت کأحد من العباد و راقداً على المهاد مرّت علیّ نسائم السّبحان
و علّمنی علم ما کان لیس هذا من عندی بل من لدن عزیز علیم و امرنی بالنّدآء
بین الأرض و السّمآء بذلک ورد علیّ ما ذرفت به دموع العارفین ما قرأت ما عند
النّاس من العلوم و ما دخلت المدارس فاسأل المدینة الّتی کنت فیها لتوقن بأنّی
لست من الکاذبین هذه ورقة حرّکتها اریاح مشیّة ربّک العزیز الحمید هل لها
استقرار عند هبوب اریاح عاصفات لا ومالک الأسمآء و الصّفات بل تحرّکها کیف
ترید لیس للعدم وجود تلقآء القدم قد جآء امره المبرم و انطقنی بذکره بین
العالمین انّی لم اکن الّا کالمیّت تلقآء امره قلّبتنی ید ارادة ربّک الرّحمن
الرّحیم

\end{otherlanguage}

Yā sulṭān, innī kuntu ka-aḥadin min al-ʿibādi wa rāqidan ʿalā al-mihād.
Marrat ʿalayya nasāʾimu al-Subḥāni wa ʿallamanī ʿilma mā kān. Laysa
hādhā min ʿindī bal min ladun ʿAzīzin ʿAlīm. Wa amaranī bi-l-nidāʾi
bayna al-arḍi wa-l-samāʾ. Bi-dhālika warada ʿalayya mā dharafat bihi
dumūʿu al-ʿārifīn. Mā qaraʾtu mā ʿinda al-nāsi min al-ʿulūmi wa mā
dakhaltu al-madārisa fa-sʾal al-madīnata allatī kuntu fīhā li-tūqina
bi-annī lastu min al-kādhibīn. Hādhihi waraqatun ḥarrakathā aryāḥu
mashīʾati rabbika al-ʿAzīzi al-Ḥamīd. Hal lahā istiqrārun ʿinda hubūbi
aryāḥin ʿāṣifāt? Lā wa Māliki al-asmāʾi wa-l-ṣifāt! Bal tuḥarrikuhā
kayfa turīd. Laysa li-l-ʿadami wujūdun tilqāʾa al-qidam. Qad jāʾa amruhu
al-mubramu wa anṭaqanī bi-dhikrihi bayna al-ʿālamīn. Innī lam akun illā
ka-l-mayyiti tilqāʾa amrihi. Qallabatnī yadu irādati rabbika al-Raḥmāni
al-Raḥīm.

``O King, I was like any other servant, asleep upon My couch. The
breezes of the All-Glorious passed over Me and taught Me the knowledge
of all that has been. This is not from Me, but from One who is Almighty,
All-Knowing. He commanded Me to raise My call between earth and heaven.
Because of this, there befell Me what has caused the tears of those who
know to flow. I have not read the learning current among people, nor
have I entered the schools. Ask the city in which I lived, so that you
may be certain that I am not among those who lie. This is a leaf stirred
by the winds of the will of your Lord, the Almighty, the Praiseworthy.
Can it remain still when tempestuous winds blow? No, by the Master of
names and attributes! Rather, they move it as they wish. Nothingness has
no existence before the Eternal. His irrevocable command has come and
caused Me to speak His remembrance among all peoples. Before His command
I was only as one dead. The hand of the will of your Lord, the
All-Merciful, the Merciful, turned Me about.''

\subsection{14}\label{section-1300}

\begin{otherlanguage}{persian}

حال بهتر آنکه آن جناب خود را بماء انقطاع که از معین قلم اعلی جاری شده
طاهر نمایند و لوجه اللّه در آنچه از قبل و بعد ظاهر شده و یا نازل گشته
تفکّر کنند و بعد بحکمت و بیان در اخماد نار ضغینه و بغضا که در قلوب احزاب
عالم مکنونست بقدر مقدور ساعى و جاهد شوند مقصود از ارسال رسل و انزال کتب
معرفة اللّه و الفت و اتّحاد عباد بوده حال ملاحظه میشود شریعت الهی را سبب و
علّت بغضا و عناد نموده‌اند زهی حسرت و ندامت که اکثری به ما عندهم متمسّک و
مشغول و از ما عند اللّه غافل و محجوب

\end{otherlanguage}

Ḥāl behtar ānke ān janāb khod rā be-māʾ-e enqeṭāʿ ke az maʿīn-e Qalam-e
Aʿlā jārī shode ṭāher namāyand va li-wajhi Allāh dar ānche az qabl va
baʿd ẓāher shode va yā nāzel gashte tafakkor konand va baʿd be-ḥekmat va
bayān dar ekhmād-e nār-e ḍaghīne va baghḍā ke dar qolūb-e aḥzāb-e ʿālam
maknūn ast be-qadr-e maqdūr sāʿī va jāhed shavand. Maqṣūd az ersāl-e
rosol va enzāl-e kotob maʿrefat-e Allāh va olfat va etteḥād-e ʿebād
būde. Ḥāl molāḥeẓe mīshavad sharīʿat-e elāhī rā sabab va ʿellat-e baghḍā
va ʿenād namūde-and. Zehī ḥasrat va nadāmat ke aktharī be-mā ʿindahom
motamassek va mashghūl va az mā ʿinda Allāh ghāfel va maḥjūb.

It would now be best for you to purify yourself with the water of
detachment that has flowed from the wellspring of the Most Exalted Pen
and, for the sake of God, reflect on whatever has appeared or been
revealed before and since. Then, through wisdom and utterance, strive to
the extent of your ability to extinguish the fire of rancor and hatred
concealed in the hearts of the peoples of the world. The purpose of
sending the Messengers and revealing the Books has been the knowledge of
God and fellowship and unity among His servants. Yet now we see that
they have made the divine law a cause of hatred and hostility. What
sorrow and regret that most people cling to and occupy themselves with
what they possess, heedless of, and veiled from, what is with God.

\subsection{15}\label{section-1301}

\begin{otherlanguage}{arabic}

قل الهی الهی زیّن رأسی باکلیل العدل و هیکلی بطراز الانصاف انّک انت مالک
المواهب و الألطاف

\end{otherlanguage}

Qul: Ilāhī Ilāhī, zayyin raʾsī bi-iklīli al-ʿadli wa haykalī bi-ṭirāzi
al-inṣāf. Innaka anta Māliku al-mawāhibi wa-l-alṭāf.

Say: ``My God, my God, adorn my head with the crown of justice and my
body with the ornament of fairness. You are indeed the Master of gifts
and favors.''

\subsection{16}\label{section-1302}

\begin{otherlanguage}{persian}

عدل و انصاف دو حارسند از برای حفظ عباد و از این دو کلمات محکمهٔ مبارکه که
علّت صلاح عالم و حفظ امم است ظاهر گردد

\end{otherlanguage}

ʿAdl va enṣāf do ḥāresand az barā-ye ḥefẓ-e ʿebād va az īn do kalemāt-e
moḥkame-ye mobārake ke ʿellat-e ṣalāḥ-e ʿālam va ḥefẓ-e omam ast ẓāher
gardad.

Justice and fairness are two guardians for the protection of the people,
and from these two emerge firm and blessed words that bring about the
well-being of the world and the preservation of nations.

\subsection{17}\label{section-1303}

\begin{otherlanguage}{arabic}

در یکی از الواح از قلم مظلوم این کلمات جاری حقّ جلّ جلاله از برای ظهور
جواهر معانی از معدن انسانی آمده یعنی مشارق امر و مخازن لآلی علم او چه که
انّه تعالی غیب مکنون مستور عن الأنظار انظر ما انزله الرّحمن فی الفرقان لا
تدرکه الأبصار و هو یدرک الأبصار و هو اللّطیف الخبیر

\end{otherlanguage}

Dar yekī az alvāḥ az Qalam-e Maẓlūm īn kalemāt jārī: Ḥaqq jalla jalāluhu
az barā-ye ẓohūr-e javāher-e maʿānī az maʿdan-e ensānī āmade, yaʿnī
mashāreq-e amr va makhāzen-e laʾālī-ye ʿelm-e ū, che ke innahu taʿālā
ghaybun maknūnun mastūrun ʿan al-anẓār. Unẓur mā anzalahu al-Raḥmānu fī
al-Furqān: Lā tudrikuhu al-abṣāru wa huwa yudriku al-abṣāra wa huwa
al-Laṭīfu al-Khabīr.

In one of the Tablets these words flowed from the Pen of the Wronged
One: ``God, exalted be His glory, has come to bring forth the gems of
meaning from the mine of humanity---that is, the dawning places of His
Cause and the treasuries of the pearls of His knowledge---for He,
exalted is He, is an unseen mystery, concealed and hidden from sight.
Consider what the All-Merciful has revealed in the Qur'an: `Eyes do not
perceive Him, but He perceives all eyes, and He is the Subtle, the
All-Aware.'\,''

\subsection{18}\label{section-1304}

\begin{otherlanguage}{arabic}

الیوم دین اللّه و مذهب اللّه آنکه مذاهب مختلفه و سبل متعدّده را سبب و علّت
بغضا ننمایند این اصول و قوانین و راههای محکم متین از مطلع واحد ظاهر و از
مشرق واحد مشرق و این اختلافات نظر بمصالح وقت و زمان و قرون و اعصار بوده

\end{otherlanguage}

Al-yowm dīn-e Allāh va madhhab-e Allāh ānke madhāheb-e mokhtalefe va
sobol-e motaʿaddede rā sabab va ʿellat-e baghḍā nanamāyand. Īn oṣūl va
qavānīn va rāhhā-ye moḥkam-e matīn az maṭlaʿ-e vāḥed ẓāher va az
mashreq-e vāḥed moshreq va īn ekhtelāfāt naẓar be-maṣāleḥ-e vaqt va
zamān va qorūn va aʿṣār būde.

Today the religion of God and the faith of God require that the diverse
religions and various paths not be made a cause of hatred. These
principles, laws, and firm, well-established paths have appeared from a
single source and shone from a single dawning place. Their differences
have been due to the needs of the times, centuries, and ages.

\subsection{19}\label{section-1305}

\begin{otherlanguage}{persian}

ای اهل بها کمر همّت را محکم نمائید که شاید جدال و نزاع مذهبی از بین اهل
عالم مرتفع شود و محو گردد حبّاً للّه و لعباده بر این امر عظیم خطیر قیام
نمائید ضغینه و بغضای مذهبی ناریست عالم‌سوز و اطفاء آن بسیار صعب مگر ید
قدرت الهی ناس را از این بلاء عقیم نجات بخشد در محاربهٔ واقعهٔ بین دولتین
ملاحظه نمائید طرفین از مال و جان گذشتند چه مقدار قریه‌ها کأن لم یکن
ملاحظه شد

\end{otherlanguage}

Ey ahl-e Bahā, kamar-e hemmat rā moḥkam namāʾīd ke shāyad jedāl va
nezāʿ-e madhhabī az beyn-e ahl-e ʿālam mortafeʿ shavad va maḥv gardad.
Ḥubban li-llāhi wa li-ʿibādihi bar īn amr-e ʿaẓīm-e khaṭīr qīyām
namāʾīd. Ḍaghīne va baghḍā-ye madhhabī nārīst ʿālamsūz va eṭfā-ye ān
besyār ṣaʿb, magar yad-e qodrat-e elāhī nās rā az īn balā-ye ʿaqīm najāt
bakhshad. Dar moḥārebe-ye vāqeʿe-ye beyn-e dowlateyn molāḥeẓe namāʾīd.
Ṭarafeyn az māl va jān goẕashtand. Che meqdār qaryehā ka-an lam yakun
molāḥeẓe shod.

O people of Bahá, gird yourselves with determination, so that religious
contention and strife may be removed from among the people of the world
and effaced. Out of love for God and His servants, arise to undertake
this great and momentous task. Religious rancor and hatred are a
world-consuming fire, and extinguishing it is exceedingly difficult,
unless the hand of divine power delivers people from this devastating
affliction. Consider the war that took place between the two states.
Both sides sacrificed possessions and lives. How many villages were seen
reduced to nothing!

\subsection{20}\label{section-1306}

\begin{otherlanguage}{persian}

مشکوة بیان را این کلمه بمثابهٔ مصباح است ای اهل عالم همه بار یک دارید و
برگ یک شاخسار بکمال محبّت و اتّحاد و مودّت و اتّفاق سلوک نمائید قسم بآفتاب
حقیقت نور اتّفاق آفاق را روشن و منوّر سازد حقّ آگاه گواه این گفتار بوده و
هست

\end{otherlanguage}

Meshkāt-e bayān rā īn kaleme be-mathābe-ye meṣbāḥ ast: Ey ahl-e ʿālam,
hame bār-e yek dārīd va barg-e yek shākhsār. Be-kamāl-e maḥabbat va
etteḥād va mavaddat va ettefāq solūk namāʾīd. Qasam be-Āftāb-e Ḥaqīqat,
nūr-e ettefāq āfāq rā rowshan va monavvar sāzad. Ḥaqq-e āgāh gavāh-e īn
goftār būde va hast.

This utterance is like a lamp in the niche of speech: ``O people of the
world, you are all the fruit of one tree and the leaves of one branch.
Conduct yourselves with perfect love, unity, affection, and concord. By
the Sun of Truth, the light of concord will brighten and illumine the
world. The all-knowing God has been, and remains, witness to these
words.''

\subsection{21}\label{section-1307}

\begin{otherlanguage}{arabic}

جهد نمائید تا باین مقام بلند اعلی که مقام صیانت و حفظ عالم انسانیست فائز
شوید این قصد سلطان مقاصد و این امل ملیک آمال ولکن تا افق آفتاب عدل از
سحاب تیرهٔ ظلم فارغ نشود ظهور این مقام مشکل بنظر میآید و سحاب تیره مظاهر
ظنون و اوهامند یعنی علمای ایران گاهی بلسان شریعت و هنگامی بلسان حقیقت و
طریقت نطق نمودیم و مقصد اقصی و غایت قصوی ظهور این مقام بلند اعلى بوده و
کفی باللّه شهیداً

\end{otherlanguage}

Jahd namāʾīd tā be-īn maqām-e boland-e aʿlā ke maqām-e ṣīyānat va ḥefẓ-e
ʿālam-e ensānīst fāʾez shavīd. Īn qaṣd solṭān-e maqāṣed va īn amal
malīk-e āmāl. Va lāken tā ofoq-e āftāb-e ʿadl az saḥāb-e tīre-ye ẓolm
fāregh nashavad ẓohūr-e īn maqām moshkel be-naẓar mīāyad va saḥāb-e tīre
maẓāher-e ẓonūn va owhāmand, yaʿnī ʿolamā-ye Īrān. Gāhī be-lesān-e
sharīʿat va hengāmī be-lesān-e ḥaqīqat va ṭarīqat noṭq namūdīm va
maqṣad-e aqṣā va ghāyat-e qoṣvā ẓohūr-e īn maqām-e boland-e aʿlā būde.
Wa kafā bi-llāhi shahīdan.

Strive to attain this lofty and most exalted station, the station of
safeguarding and preserving the human world. This purpose is the
sovereign of all purposes, and this hope the king of all hopes. Yet
until the horizon of the sun of justice is freed from the dark clouds of
oppression, the realization of this station appears difficult. These
dark clouds are the embodiments of conjectures and vain
imaginings---that is, the religious scholars of Iran. At times We have
spoken in the language of religious law, and at others in the language
of mystical truth and the spiritual path. Our ultimate purpose and
highest aim have been the realization of this lofty and most exalted
station. God is sufficient as a witness.

\subsection{22}\label{section-1308}

\begin{otherlanguage}{persian}

ای اهل بها با جمیع اهل عالم بروح و ریحان معاشرت نمائید اگر نزد شما
کلمه‌ئی و یا جوهریست که دون شما از آن محروم بلسان محبّت و شفقت القا نمائید
و بنمائید اگر قبول شد و اثر نمود مقصد حاصل والّا او را باو گذارید و در
بارهٔ او دعا نمائید نه جفا لسان شفقت جذّاب قلوبست و مائدهٔ روح و بمثابهٔ
معانی است از برای الفاظ و مانند افق است از برای اشراق آفتاب حکمت و
دانائی

\end{otherlanguage}

Ey ahl-e Bahā, bā jamīʿ-e ahl-e ʿālam be-rūḥ va rayḥān moʿāsherat
namāʾīd. Agar nazd-e shomā kalemeʾī va yā jowharīst ke dūn-e shomā az ān
maḥrūm, be-lesān-e maḥabbat va shafaqat elqā namāʾīd va benamāʾīd. Agar
qabūl shod va athar namūd maqṣad ḥāṣel, va ellā ū rā be-ū goẕārīd va dar
bāre-ye ū doʿā namāʾīd na jafā. Lesān-e shafaqat jadhdhāb-e qolūb ast va
māʾede-ye rūḥ va be-mathābe-ye maʿānī ast az barā-ye alfāẓ va mānand-e
ofoq ast az barā-ye eshrāq-e āftāb-e ḥekmat va dānāʾī.

O people of Bahá, associate with all the people of the world in joy and
fellowship. If you possess a word or a gem of which others are deprived,
impart it and make it known in the language of love and compassion. If
it is accepted and takes effect, the purpose is achieved; otherwise,
leave the person to himself and pray for him, rather than mistreat him.
The language of compassion attracts hearts and is nourishment for the
spirit. It is like meaning to words and like the horizon for the shining
of the sun of wisdom and knowledge.

\subsection{23}\label{section-1309}

\begin{otherlanguage}{persian}

مقصود از علما در این موارد که ذکر شده نفوسی هستند که خود را در ظاهر
بلباس علم میآرایند و در باطن از آن محروم در ذکر این مقام در لوح حضرت
سلطان چند فقره از فقرات کلمات مکنونه که باسم صحیفهٔ فاطمیّه صلوات اللّه
علیها از قلم ابهی ظاهر ذکر میشود

\end{otherlanguage}

Maqṣūd az ʿolamā dar īn mavāred ke dhekr shode nofūsī hastand ke khod rā
dar ẓāher be-lebās-e ʿelm mīārāyand va dar bāṭen az ān maḥrūm. Dar
dhekr-e īn maqām dar lowḥ-e Ḥaḍrat-e Solṭān chand faqre az faqarāt-e
Kalemāt-e Maknūne ke be-esm-e Ṣaḥīfe-ye Fāṭemīyye, ṣalawātu Allāhi
ʿalayhā, az Qalam-e Abhā ẓāher, dhekr mīshavad.

By religious scholars in the instances mentioned here are meant those
who outwardly adorn themselves with the garb of learning but are
inwardly deprived of it. In discussing this subject in the Tablet to His
Majesty the Shah, several passages are cited from the Hidden Words,
which appeared from the Most Glorious Pen under the title of the Scroll
of Fátimih, may God's blessings rest upon her.

\subsection{24}\label{section-1310}

\begin{otherlanguage}{persian}

ای بیوفایان چرا در ظاهر دعوی شبانی کنید و در باطن ذئب اغنام من شده‌اید
مثل شما مثل ستارهٔ قبل از صبح است که در ظاهر درّی و روشن است و در باطن سبب
اضلال و هلاکت کاروانهای مدینه و دیار منست

\end{otherlanguage}

Ey bīvafāyān, cherā dar ẓāher daʿvī-ye shabānī konīd va dar bāṭen
dheʾb-e aghnām-e man shode-īd? Mathal-e shomā mathal-e setāre-ye qabl az
ṣobḥ ast ke dar ẓāher dorrī va rowshan ast va dar bāṭen sabab-e eḍlāl va
halākat-e kārvānhā-ye madīne va dīyār-e man ast.

``O faithless ones, why do you outwardly claim to be shepherds, while
inwardly you have become wolves preying on My sheep? You are like the
star before dawn, which appears radiant and bright, but in reality leads
the caravans of My city and lands astray and to destruction.''

\subsection{25}\label{section-1311}

\begin{otherlanguage}{persian}

و همچنین میفرماید ای بظاهر آراسته و بباطن کاسته مثل تو مثل آب تلخ صافیست
که کمال لطافت و صفا از آن در ظاهر مشاهده شود و چون بدست صرّاف ذائقهٔ احدیّه
افتد قطره‌ئی از آن را قبول نفرماید تجلّی آفتاب در تراب و مرآت هر دو موجود
ولکن از فرقدان تا ارض فرق دان بلکه فرق بی‌منتهى در میان

\end{otherlanguage}

Va hamchenīn mīfarmāyad: Ey be-ẓāher ārāste va be-bāṭen kāste, mathal-e
to mathal-e āb-e talkh-e ṣāfīst ke kamāl-e leṭāfat va ṣafā az ān dar
ẓāher moshāhede shavad va chon be-dast-e ṣarrāf-e dhāʾeqe-ye aḥadīyye
oftad qaṭreʾī az ān rā qabūl nafarmāyad. Tajallī-ye āftāb dar torāb va
merʾāt har do mowjūd, va lāken az Farqadān tā arḍ farq dān, balke farq-e
bīmontahā dar mīyān.

And He also says: ``O you who are outwardly adorned but inwardly
deficient, you are like clear but bitter water, which outwardly appears
perfectly delicate and pure, yet when tested by the divine Assayer's
taste, not a drop of it is accepted. The sun's radiance is present in
both dust and mirror, but know that the difference between them is as
great as that between the two Pole Stars and the earth---indeed, the
difference is infinite.''

\subsection{26}\label{section-1312}

\begin{otherlanguage}{arabic}

و همچنین میفرماید ای پسر دنیا بسا سحرگاهان تجلّی عنایت من از مشرق لامکان
بمکان تو آمد و ترا در بستر راحت بغیر مشغول دید و چون برق روحانی بمقرّ عزّ
نورانی رجوع نمود و در مکامن قرب نزد جنود قدس اظهار نداشتم و خجلت ترا
نپسندیدم

\end{otherlanguage}

Va hamchenīn mīfarmāyad: Ey pesar-e donyā, basā saḥargāhān tajallī-ye
ʿenāyat-e man az mashreq-e lāmakān be-makān-e to āmad va to rā dar
bestar-e rāḥat be-ghayr mashghūl dīd va chon barq-e rūḥānī be-maqarr-e
ʿezz-e nūrānī rojūʿ namūd va dar makāmen-e qorb nazd-e jonūd-e qods
eẓhār nadāshtam va khejlat-e to rā napasandīdam.

And He also says: ``O son of the world, many a dawn the revelation of My
loving-kindness came from the dawning place of the Placeless to your
dwelling, found you on your bed of ease occupied with another, and
returned like a spiritual flash of lightning to its seat of radiant
glory. In the retreats of nearness, before the hosts of holiness, I did
not disclose it, nor did I wish for your shame.''

\subsection{27}\label{section-1313}

\begin{otherlanguage}{arabic}

و همچنین میفرماید ای مدّعی دوستی من در سحرگاهان نسیم عنایت من بر تو مرور
نمود و ترا در فراش غفلت خفته یافت و بر حال تو گریست و بازگشت

\end{otherlanguage}

Va hamchenīn mīfarmāyad: Ey moddaʿī-ye dūstī-ye man, dar saḥargāhān
nasīm-e ʿenāyat-e man bar to morūr namūd va to rā dar ferāsh-e ghaflat
khofte yāft va bar ḥāl-e to gerīst va bāzgasht.

And He also says: ``O you who claim to love Me, at dawn the breeze of My
loving-kindness passed over you and found you asleep upon the bed of
heedlessness. It wept over your condition and returned.''

\subsection{28}\label{section-1314}

\begin{otherlanguage}{arabic}

امّا علمائی که فی‌الحقیقه بطراز علم و اخلاق مزیّنند ایشان بمثابهٔ رأسند از
برای هیکل عالم و مانند بصرند از برای امم لازال هدایت عباد بآن نفوس مقدّسه
بوده و هست نسأل اللّه ان یوفّقهم علی ما یحبّ و یرضى انّه هو مولی الوری و ربّ
الآخرة و الأولی

\end{otherlanguage}

Ammā ʿolamāʾī ke fī-l-ḥaqīqe be-ṭerāz-e ʿelm va akhlāq mozayyanand,
īshān be-mathābe-ye raʾsand az barā-ye haykal-e ʿālam va mānand-e
baṣarand az barā-ye omam. Lāzāl hedāyat-e ʿebād be-ān nofūs-e moqaddase
būde va hast. Nasʾalu Allāha an yuwaffiqahum ʿalā mā yuḥibbu wa yarḍā.
Innahu huwa Mawlā al-warā wa Rabbu al-ākhirati wa-l-ūlā.

But those scholars who are truly adorned with learning and good
character are like the head to the body of the world and like eyes to
the nations. The guidance of the people has always depended, and
continues to depend, upon those sanctified souls. We ask God to enable
them to do what He loves and approves. He is indeed the Master of
creation and the Lord of the next world and this one.

\subsection{29}\label{section-1315}

\begin{otherlanguage}{persian}

یا شیخ انّا سمعنا انّ جنابک اعرضت عنّا و اعترضت علینا حیث امرت النّاس بسبّی و
افتیت علی سفک دمآء العباد للّه درّ من قال طوعاً لقاض اتی فی حکمه عجباً افتی
بسفک دمی فی الحلّ و الحرم براستی میگویم آنچه در سبیل الهی وارد شود محبوب
جانست و مقصود روان سمّ مهلک در سبیلش شهدیست فائق و عذابش عذبیست لائق در
لوح حضرت سلطان ذکر شده ونفسه الحقّ لا اجزع من البلایا فی سبیله و لا من
الرّزایا فی حبّه قد جعل اللّه البلآء غادیة لهذه الدّسکرة الخضرآء و ذبالة
لمصباحه الّذی به اشرقت الأرض و السّمآء

\end{otherlanguage}

Yā Shaykh, innā samiʿnā anna janābaka aʿraḍta ʿannā wa-iʿtaraḍta ʿalaynā
ḥaythu amarta al-nāsa bi-sabbī wa-aftayta ʿalā safki dimāʾi al-ʿibād.
Li-llāhi darru man qāla: Ṭawʿan li-qāḍin atā fī ḥukmihi ʿajaban, aftā
bi-safki damī fī al-ḥilli wa-al-ḥaram. Be-rāstī mīgūyam ānche dar
sabīl-e ilāhī vāred shavad maḥbūb-e jān ast va maqṣūd-e ravān. Samm-e
mohlik dar sabīlash shahdīst fāʾiq va ʿadhābash ʿadhbīst lāʾiq. Dar
lawḥ-e Ḥaḍrat-e Sulṭān dhekr shode: Wa-nafsihi al-Ḥaqq, lā ajzaʿu min
al-balāyā fī sabīlihi wa-lā min al-razāyā fī ḥubbihi. Qad jaʿala Allāhu
al-balāʾa ghādiyatan li-hādhihi al-daskara al-khaḍrāʾ wa-dhubālatan
li-miṣbāḥihi alladhī bihi ashraqat al-arḍu wa-al-samāʾ.

O Shaykh, We have heard that you have turned away from Us and opposed
Us, commanding people to revile Me and issuing a ruling for the shedding
of the servants' blood. God bless the one who said: ``Willing obedience
to a judge whose judgment is a wonder: he has decreed that my blood be
shed both outside the sanctuary and within it.'' Truly I say, whatever
befalls one in the path of God is beloved of the soul and the heart's
desire. Deadly poison in His path is surpassing honey, and His torment a
fitting sweetness. In the Tablet to His Majesty the Shah it is stated:
``By His own Self, the Truth! I do not bewail afflictions in His path or
calamities endured for His love. God has made affliction a morning rain
for this verdant field and a wick for His lamp, through which earth and
heaven have been illumined.''

\subsection{30}\label{section-1316}

\begin{otherlanguage}{persian}

اقبل بقلبک الی شطر کعبة اللّه المهیمن القیّوم ثمّ ارفع یدیک باستقامة ترتفع
بها ایادی الممکنات الی سمآء فضل اللّه ربّ العالمین ثمّ وجّه الیه بتوجّه تتوجّه
به الکائنات الی افقه المشرق المنیر و قل ای ربّ ترانی مقبلاً الی سمآء جودک
و بحر عطائک و معرضاً عن دونک اسألک بتجلّیات نیّر ظهورک فی الطّور و باشراقات
شمس فضلک من افق اسمک الغفور ان تغفر لی و ارحمنی ثمّ اکتب لی من قلمک
الأعلی ما یرفعنی باسمک فی ناسوت الانشآء ای ربّ وفّقنی علی التّوجّه الیک و
اصغآء ندآء اولیائک الّذین ما اضعفتهم قوّة العالم و ما منعتهم سطوة الأمم
اقبلوا و قالوا اللّه ربّنا و ربّ من فی السّموات و الأرضین

\end{otherlanguage}

Aqbil bi-qalbika ilā shaṭri Kaʿbati Allāhi al-Muhaymini al-Qayyūm,
thumma irfaʿ yadayka bi-istiqāmatin tartafiʿu bihā ayādī al-mumkināt ilā
samāʾi faḍli Allāhi rabbi al-ʿālamīn. Thumma wajjih ilayhi bi-tawajjuhin
tatawajjahu bihi al-kāʾināt ilā ufuqihi al-mushriqi al-munīr, wa-qul: Ay
rabb, tarānī muqbilan ilā samāʾi jūdika wa-baḥri ʿaṭāʾika wa-muʿriḍan
ʿan dūnika. Asʾaluka bi-tajalliyāti nayyiri ẓuhūrika fī al-Ṭūr
wa-bi-ishrāqāti shamsi faḍlika min ufuqi ismika al-Ghafūr an taghfira lī
wa-irḥamnī, thumma uktub lī min qalamika al-aʿlā mā yarfaʿunī bi-ismika
fī nāsūti al-inshāʾ. Ay rabb, waffiqnī ʿalā al-tawajjuhi ilayka
wa-iṣghāʾi nidāʾi awliyāʾika alladhīna mā aḍʿafathum quwwatu al-ʿālam
wa-mā manaʿathum saṭwatu al-umam. Aqbalū wa-qālū: Allāhu rabbunā
wa-rabbu man fī al-samāwāti wa-al-araḍīn.

Turn your heart toward the Kaaba of God, the All-Protecting, the
Self-Subsisting. Then raise your hands with such steadfastness that the
hands of all created things may be raised through it toward the heaven
of the grace of God, the Lord of the worlds. Then turn toward Him with
such devotion that all beings may thereby turn toward His radiant,
luminous horizon, and say: ``My Lord, You see me turning toward the
heaven of Your generosity and the ocean of Your bounty, and turning away
from all besides You. I ask You, by the splendors of the Daystar of Your
revelation upon the Mount and by the rays of the sun of Your grace
shining from the horizon of Your name, the Ever-Forgiving, to forgive me
and have mercy on me. Then ordain for me, through Your Most Exalted Pen,
what will exalt me through Your name in the earthly realm of creation.
My Lord, enable me to turn toward You and to hear the call of Your loved
ones, whom the power of the world has not weakened nor the might of the
nations restrained. They turned toward You and said: `God is our Lord
and the Lord of all who are in the heavens and on earth.'\,''

\subsection{31}\label{section-1317}

\begin{otherlanguage}{persian}

یا شیخ براستی میگویم ختم رحیق مختوم باسم قیّوم برداشته شد خود را محروم
منما این مظلوم لوجه اللّه میگوید تو هم لوجه اللّه در آنچه نازل شده و ظاهر
گشته تفکّر نما شاید از فیوضات فیّاض حقیقی در این یوم مبارک نصیب برداری و
محروم نمانی لیس هذا علی اللّه بعزیز آدم خاکی از کلمهٔ الهی عرشی شد و صیّاد
ماهی دارای حکمت ربّانی گشت حضرت ابوذر راعی غنم بود سیّد امم شد

\end{otherlanguage}

Yā Shaykh, be-rāstī mīgūyam khatm-e raḥīq-e makhtūm be-ism-e Qayyūm
bardāshte shod. Khod rā maḥrūm manamā. Īn maẓlūm li-wajhi Allāh mīgūyad,
to ham li-wajhi Allāh dar ānche nāzel shode va ẓāher gashte tafakkor
namā, shāyad az foyūḍāt-e Fayyāḍ-e Ḥaqīqī dar īn yawm-e mobārak naṣīb
bardārī va maḥrūm namānī. Laysa hādhā ʿalā Allāhi bi-ʿazīz. Ādam-e khākī
az kalame-ye ilāhī ʿarshī shod va ṣayyād-e māhī dārā-ye ḥekmat-e rabbānī
gasht. Ḥaḍrat-e Abū Dharr rāʿī-ye ghanam būd, sayyed-e omam shod.

O Shaykh, truly I say, the seal of the sealed wine has been removed in
the name of the Self-Subsisting. Do not deprive yourself. This Wronged
One speaks for the sake of God; you too, for the sake of God, reflect on
what has been sent down and made manifest, so that you may receive a
share of the outpourings of the true Bestower of grace on this blessed
Day and not remain deprived. This is not difficult for God. Through the
Word of God, earthly Adam became heavenly, and a fisherman became
endowed with divine wisdom. His Honor Abū Dharr was a shepherd of sheep;
he became a leader of nations.

\subsection{32}\label{section-1318}

\begin{otherlanguage}{persian}

یا شیخ امروز روز علوم و فنون ظاهره نبوده و نیست چه که ملاحظه شد نفسی که
بیک کلمه از آن علوم آگاه نه بر کرسیّ عقیان در صدر مجلس عرفان مستوی و صاحب
علوم و دارای فنون محروم مقصود از این علوم علومىیست که از لفظ ابتدا شود و
بلفظ منتهی گردد ولکن علومی که اثر و ثمر از آن ظاهر و سبب راحت و آسایش
عباد است عنداللّه مقبول بوده و هست لو تسمع ندائی تدع ما عندک و تتوجّه الى
مقام ماج فیه بحر الحکمة و البیان و هاج عرف عنایة ربّک الرّحمن

\end{otherlanguage}

Yā Shaykh, emrūz rūz-e ʿolūm va fonūn-e ẓāhere nabūde va nīst, che ke
molāḥeẓe shod nafsī ke be-yek kalame az ān ʿolūm āgāh na, bar korsī-ye
ʿeqyān dar ṣadr-e majles-e ʿerfān mostavī va ṣāḥeb-e ʿolūm va dārā-ye
fonūn maḥrūm. Maqṣūd az īn ʿolūm ʿolūmīst ke az lafẓ ebtedā shavad va
be-lafẓ montahī gardad, va-līkan ʿolūmī ke athar va thamar az ān ẓāher
va sabab-e rāḥat va āsāyesh-e ʿebād ast ʿinda Allāh maqbūl būde va hast.
Law tasmaʿ nidāʾī tadaʿ mā ʿindaka wa-tatawajjahu ilā maqāmin māja fīhi
baḥru al-ḥikmati wa-al-bayān wa-hāja ʿarfu ʿināyati rabbika al-Raḥmān.

O Shaykh, this day has not been, and is not, the day of outward learning
and arts, for it has been seen that a person unacquainted with a single
word of such learning is seated upon a throne of pure gold at the head
of the gathering of spiritual knowledge, while one possessed of learning
and versed in the arts remains deprived. By this learning is meant
learning that begins with words and ends with words. But sciences that
yield results and bear fruit, and bring comfort and well-being to the
servants of God, have always been and remain acceptable to God. If you
hear My call, you will abandon what you possess and turn toward a
station where the ocean of wisdom and utterance surges and the fragrance
of the loving care of your Lord, the All-Merciful, wafts.

\subsection{33}\label{section-1319}

\begin{otherlanguage}{persian}

در این مقام بنظر آمد فی‌الجمله از امور گذشته ذکر شود شاید سبب ظهور عدل و
انصاف گردد این مظلوم در ایّامی که حضرت سلطان ایّده اللّه ربّه الرّحمن عزم
توجّه به اصفهان نموده اذن حاصل کرده قصد زیارت بقاع مقدّسهٔ منوّرهٔ ائمّه صلوات
اللّه علیهم نموده و بعد از رجوع نظر بگرمی هواء دارالخلافة و شدّت آن به
لواسان رفتیم و بعد از توجّه حکایت حضرت سلطان ایّده اللّه تبارک و تعالی واقع
و در آن ایّام امور منقلب و نار غضب مشتعل جمعی را اخذ نمودند از جمله این
مظلوم را لعمر اللّه ابداً داخل آن امر منکر نبودیم و در مجالس تحقیق هم عدم
تقصیر ثابت مع‌ذلک ما را اخذ نمودند و از نیاوران که در آن ایّام مقرّ سلطنت
بوده سر برهنه و پای برهنه پیاده با زنجیر بسجن طهران بردند چه که یک ظالمی
سواره همراه کلاه از سر برداشت و بسرعت تمام با جمعى از میرغضبان و فرّاشان
ما را بردند و چهار شهر در مقامی که شبه و مثل نداشت مقرّ معیّن نمودند امّا
سجن که محلّ مظلوم و مظلومان بوده فى‌الحقیقه دخمهٔ تنگ تاریک از آن افضل بوده
و چون وارد حبس شدیم بعد از ورود ما را داخل دالانی ظلمانی نمودند از آنجا
از سه پلّهٔ سراشیب گذشتیم و بمقرّی که معیّن نموده بودند رسیدیم امّا محل تاریک
و معاشر قریب صد و پنجاه نفس از سارقین اموال و قاتلین نفوس و قاطعین طرق
بوده مع این جمعیّت محل منفذ نداشت جز طریقی که وارد شدیم اقلام از وصفش
عاجز و روائح منتنه‌اش خارج از بیان و آن جمع اکثری بى لباس و فراش اللّه
یعلم ما ورد علینا فى ذاک المقام الأنتن الأظلم

\end{otherlanguage}

Dar īn maqām be-naẓar āmad fī al-jomle az omūr-e godhashte dhekr shavad,
shāyad sabab-e ẓohūr-e ʿadl va enṣāf gardad. Īn maẓlūm dar ayyāmī ke
Ḥaḍrat-e Sulṭān, ayyadahu Allāhu rabbuhu al-Raḥmān, ʿazm-e tavajjoh be
Eṣfahān namūde, edhn ḥāṣel karde qaṣd-e ziyārat-e beqāʿ-e moqaddase-ye
monavvare-ye aʾemme, ṣalawātu Allāhi ʿalayhim, namūde, va baʿd az rojūʿ
naẓar be-garmī-ye havā-ye Dār al-Khelāfe va sheddat-e ān be Lavāsān
raftīm. Va baʿd az tavajjoh ḥekāyat-e Ḥaḍrat-e Sulṭān, ayyadahu Allāhu
tabāraka wa-taʿālā, vāqeʿ va dar ān ayyām omūr monqaleb va nār-e ghaḍab
moshtaʿel. Jamʿī rā akhdh namūdand, az jomle īn maẓlūm rā. La-ʿamru
Allāh, abadan dākhel-e ān amr-e monkar nabūdīm va dar majāles-e taḥqīq
ham ʿadam-e taqṣīr thābet. Maʿa dhālek mā rā akhdh namūdand va az
Niyāvarān ke dar ān ayyām maqarr-e salṭanat būde, sar-berehne va
pā-berehne piyāde bā zanjīr be-sejn-e Ṭehrān bordand, che ke yek ẓālemī
savāre hamrāh kolāh az sar bardāsht va be-sorʿat-e tamām bā jamʿī az
mīrghaḍabān va farrāshān mā rā bordand va chahār shahr dar maqāmī ke
shabah va methl nadāsht maqarr moʿayyan namūdand. Ammā sejn ke maḥall-e
maẓlūm va maẓlūmān būde, fī al-ḥaqīqe dakhme-ye tang-e tārīk az ān afḍal
būde, va chūn vāred-e ḥabs shodīm baʿd az vorūd mā rā dākhel-e dālānī
ẓolmānī namūdand. Az ānjā az se pelle-ye sarāshīb godhashtīm va
be-maqarrī ke moʿayyan namūde būdand resīdīm. Ammā maḥall tārīk va
moʿāsher qarīb-e ṣad va panjāh nafs az sāreqīn-e amvāl va qātelīn-e
nofūs va qāṭeʿīn-e ṭoroq būde. Maʿ-e īn jamʿiyyat maḥall manfadh nadāsht
joz ṭarīqī ke vāred shodīm. Aqlām az vaṣfash ʿājez va ravāʾeḥ-e
montene-ash khārej az bayān, va ān jamʿ aktharī bī lebās va ferāsh.
Allāhu yaʿlamu mā warada ʿalaynā fī dhāka al-maqāmi al-antan al-aẓlam.

At this point it seemed fitting to mention briefly some past events, in
the hope that this might bring forth justice and fairness. In the days
when His Majesty the Shah---may God, his Lord, the All-Merciful, assist
him---had resolved to travel to Isfahan, this Wronged One obtained
permission and set out to visit the sacred, luminous shrines of the
Imams, may God's blessings be upon them. After returning, because of the
intense heat in the capital, We went to Lavasan. After Our departure,
the incident involving His Majesty the Shah---may God, blessed and
exalted be He, assist him---occurred. In those days affairs were thrown
into turmoil and the fire of anger blazed. A number of people were
arrested, among them this Wronged One. By the life of God, We had no
involvement whatever in that reprehensible affair, and Our innocence was
also established in the investigative proceedings. Nevertheless, they
arrested Us and took Us in chains from Niyavaran, then the seat of the
royal court, to the prison in Tehran, bareheaded and barefoot, on foot.
For a mounted oppressor accompanying Us snatched the hat from Our head,
and with a band of executioners and attendants they hurried Us along at
full speed. For four months they assigned Us quarters in a place without
parallel or likeness. As for the prison that housed this Wronged One and
others who had been wronged, a narrow, dark tomb would in truth have
been preferable to it. When We entered the prison, they led Us into a
dark corridor. From there We descended three steps and reached the
quarters they had assigned Us. The place was dark, and Our companions
were some hundred and fifty people: thieves, murderers, and highway
robbers. Despite this crowd, the place had no opening except the passage
by which We had entered. Pens are powerless to describe it, and its foul
odors are beyond expression. Most of those people had neither clothing
nor bedding. God knows what befell Us in that most foul-smelling, most
dark place.

\subsection{34}\label{section-1320}

\begin{otherlanguage}{persian}

و در ایّام و لیالی در سجن مذکور در اعمال و احوال و حرکات حزب بابی تفکّر
مینمودیم که مع علوّ و سموّ و ادراک آن حزب آیا چه شده که از ایشان چنین عملی
ظاهر یعنى جسارت و حرکت آن حزب نسبت بذات شاهانه و بعد این مظلوم اراده
نمود که بعد از خروج از سجن بتمام همّت در تهذیب آن نفوس قیام نماید

\end{otherlanguage}

Va dar ayyām va layālī dar sejn-e madhkūr dar aʿmāl va aḥvāl va
ḥarakāt-e ḥezb-e Bābī tafakkor mīnamūdīm ke maʿ-e ʿolovv va somovv va
edrāk-e ān ḥezb āyā che shode ke az īshān chonīn ʿamalī ẓāher, yaʿnī
jasārat va ḥarakat-e ān ḥezb nesbat be-dhāt-e shāhāne. Va baʿd īn maẓlūm
erāde namūd ke baʿd az khorūj az sejn be-tamām-e hemmat dar tahdhīb-e ān
nofūs qiyām namāyad.

During the days and nights in that prison, We reflected on the deeds,
circumstances, and conduct of the Bábí community: despite that
community's loftiness, nobility, and understanding, what had happened
that such an act should have issued from them---namely, that community's
audacious action against the person of the Shah? This Wronged One then
resolved, upon leaving prison, to arise with all His strength to refine
the character of those souls.

\subsection{35}\label{section-1321}

\begin{otherlanguage}{persian}

و در شبی از شبها در عالم رؤیا از جمیع جهات این کلمهٔ علیا اصغا شد انّا
ننصرک بک و بقلمک لا تحزن عمّا ورد علیک و لا تخف انّک من الآمنین سوف یبعث
اللّه کنوز الأرض و هم رجال ینصرونک بک و باسمک الّذی به احیا اللّه افئدة
العارفین

\end{otherlanguage}

Va dar shabī az shabhā dar ʿālam-e roʾyā az jamīʿ-e jahāt īn kalame-ye
ʿolyā eṣghā shod: Innā nanṣuruka bika wa-bi-qalamika. Lā taḥzan ʿammā
warada ʿalayka wa-lā takhaf, innaka min al-āminīn. Sawfa yabʿathu Allāhu
kunūza al-arḍ wa-hum rijālun yanṣurūnaka bika wa-bi-ismika alladhī bihi
aḥyā Allāhu afʾidata al-ʿārifīn.

One night, in a dream, these exalted words were heard from every
direction: ``We will indeed make You victorious through Yourself and
through Your Pen. Do not grieve over what has befallen You, and do not
fear: You are among those who are secure. God will raise up the
treasures of the earth: men who will aid You through Yourself and
through Your name, by which God has brought the hearts of those who
recognize Him to life.''

\subsection{36}\label{section-1322}

\begin{otherlanguage}{persian}

و چون مظلوم از سجن خارج حسب الأمر حضرت پادشاه حرسه اللّه تعالى مع غلام
دولت علیّهٔ ایران و دولت بهیّهٔ روس به عراق عرب توجّه نمودیم و بعد از ورود
باعانت الهی و فضل و رحمت ربّانی آیات بمثل غیث هاطل نازل و باطراف ارض
ارسال شد و جمیع عباد را مخصوص این حزب را بمواعظ حکیمانه و نصایح مشفقانه
نصیحت نمودیم و از فساد و نزاع و جدال و محاربه منع کردیم تا آنکه از فضل
الهی غفلت و نادانی ببر و دانائی بدل گشت و سلاح باصلاح

\end{otherlanguage}

Va chūn maẓlūm az sejn khārej, ḥasb al-amr-e Ḥaḍrat-e Pādshāh, ḥarasahu
Allāhu taʿālā, maʿ-e gholām-e Dawlat-e ʿAliyye-ye Īrān va Dawlat-e
Bahiyye-ye Rūs be ʿErāq-e ʿArab tavajjoh namūdīm. Va baʿd az vorūd
be-eʿānat-e ilāhī va faḍl va raḥmat-e rabbānī āyāt be-methl-e ghayth-e
hāṭel nāzel va be-aṭrāf-e arḍ ersāl shod, va jamīʿ-e ʿebād rā makhṣūṣ īn
ḥezb rā be-mavāʿeẓ-e ḥakīmāne va naṣāyeḥ-e moshfeqāne naṣīḥat namūdīm va
az fasād va nezāʿ va jedāl va moḥārebe manʿ kardīm, tā ānke az faḍl-e
ilāhī ghaflat va nādānī be-berr va dānāʾī badal gasht va selāḥ be-eṣlāḥ.

When this Wronged One left prison, at the command of His Majesty the
Shah---may God, exalted be He, protect him---We set out for Arab Iraq,
accompanied by an attendant of the exalted government of Iran and of the
illustrious government of Russia. After Our arrival, through divine
assistance and the grace and mercy of the Lord, verses were sent down
like pouring rain and dispatched to the regions of the earth. We
counseled all the servants of God, and especially this community, with
wise exhortations and compassionate advice, and forbade corruption,
strife, contention, and warfare, until, through God's grace,
heedlessness and ignorance were changed into righteousness and
understanding, and weapons into reform.

\subsection{37}\label{section-1323}

\begin{otherlanguage}{persian}

و در ایّام توقّف در سجن ارض طاء اگرچه نوم از زحمت سلاسل و روائح منتنه قلیل
بود ولکن بعضى از اوقات که دست میداد احساس میشد از جهت اعلای رأس چیزى بر
صدر میریخت بمثابهٔ رودخانهٔ عظیمى که از قلّهٔ جبل باذخ رفیعی بر ارض بریزد و
بآن جهت از جمیع اعضا آثار نار ظاهر و در آن حین لسان قرائت مینمود آنچه را
که بر اصغاء آن احدى قادر نه

\end{otherlanguage}

Va dar ayyām-e tavaqqof dar sejn-e Arḍ-e Ṭāʾ agarche nawm az zaḥmat-e
salāsel va ravāʾeḥ-e montene qalīl būd, va-līkan baʿḍī az awqāt ke dast
mīdād eḥsās mīshod az jehat-e aʿlā-ye raʾs chīzī bar ṣadr mīrīkht,
be-mathābe-ye rūdkhāne-ye ʿaẓīmī ke az qolle-ye jabal-e bādhekh-e rafīʿī
bar arḍ berīzad. Va be-ān jehat az jamīʿ-e aʿḍā āthār-e nār ẓāher va dar
ān ḥīn lesān qerāʾat mīnamūd ānche rā ke bar eṣghā-ye ān aḥadī qāder na.

During Our confinement in the prison of the Land of Ṭā, although sleep
was scarce because of the torment of the chains and the foul odors, on
those occasions when it did come, it felt as though something poured
from the crown of Our head onto Our breast, like a mighty river
descending to the earth from the summit of a lofty, towering mountain.
As a result, every limb showed signs of fire, and at that moment Our
tongue recited what no one could bear to hear.

\subsection{38}\label{section-1324}

\begin{otherlanguage}{arabic}

و بعضی از بیانات که در الواح مخصوص این حزب نازل ذکر میشود تا بیقین مبین
کل بدانند که این مظلوم عمل نموده آنچه را که نزد عقلا و مظاهر عدل و انصاف
محبوب و مقبول بوده

\end{otherlanguage}

Va baʿḍī az bayānāt ke dar alvāḥ makhṣūṣ-e īn ḥezb nāzel dhekr mīshavad
tā be-yaqīn-e mobīn koll bedānand ke īn maẓlūm ʿamal namūde ānche rā ke
nazd-e ʿoqalā va maẓāher-e ʿadl va enṣāf maḥbūb va maqbūl būde.

Some of the utterances revealed in Tablets addressed specifically to
this community are quoted here, so that all may know with clear
certainty that this Wronged One has acted in ways beloved and acceptable
to people of understanding and to those who embody justice and fairness.

\subsection{39}\label{section-1325}

\begin{otherlanguage}{arabic}

یا اولیآء اللّه فى بلاده و احبّائه فى دیاره یوصیکم المظلوم بالأمانة و
الدّیانة طوبی لمدینة فازت بأنوارهما بهما یرتفع مقام الانسان و یفتح باب
الاطمینان على من فى الامکان طوبی لمن تمسّک بهما و عرف شأنهما و ویل لمن
انکر مقامهما

\end{otherlanguage}

Yā awliyāʾa Allāhi fī bilādihi wa-aḥibbāʾahu fī diyārihi, yūṣīkumu
al-Maẓlūmu bi-al-amānati wa-al-diyāna. Ṭūbā li-madīnatin fāzat
bi-anwārihimā. Bihimā yartafiʿu maqāmu al-insān wa-yuftaḥu bābu
al-iṭmiʾnān ʿalā man fī al-imkān. Ṭūbā li-man tamassaka bihimā wa-ʿarafa
shaʾnahumā, wa-waylun li-man ankara maqāmahumā.

O friends of God in His lands and His loved ones in His realms! The
Wronged One exhorts you to trustworthiness and piety. Blessed is the
city that has attained their light. Through them the station of the
human being is exalted and the door of assurance is opened to all in the
contingent world. Blessed is the one who holds fast to them and
recognizes their worth, and woe to the one who denies their station.

\subsection{40}\label{section-1326}

\begin{otherlanguage}{arabic}

و در مقام دیگر این کلمات نازل انّا نأمر عباد اللّه و امائه بالعصمة و
التّقوى لیقومنّ من رقد الهوی و یتوجّهنّ الى اللّه فاطر الأرض و السّمآء کذلک
امرنا العباد حینما اشرق نیّر الآفاق من جهة العراق لیس ضرّی سجنی و بلائی و
ما ورد علیّ من طغاة العباد بل عمل الّذین ینسبون انفسهم الی نفسی و یرتکبون
ما ینوح به قلبی و قلمی انّ الّذین یفسدون فى الأرض و یتصرّفون فى اموال النّاس
و یدخلون البیوت من غیر اذن صاحبها انّى بریء منهم الّا ان یتوبوا و یرجعوا
الى اللّه الغفور الرّحیم

\end{otherlanguage}

Va dar maqām-e dīgar īn kalemāt nāzel: Innā naʾmuru ʿibāda Allāhi
wa-imāʾahu bi-al-ʿiṣmati wa-al-taqwā. Li-yaqūmanna min raqdi al-hawā
wa-yatawajjahanna ilā Allāhi fāṭiri al-arḍi wa-al-samāʾ. Kadhālika
amarnā al-ʿibāda ḥīnamā ashraqa nayyiru al-āfāq min jihati al-ʿIrāq.
Laysa ḍurrī sijnī wa-balāʾī wa-mā warada ʿalayya min ṭughāti al-ʿibād,
bal ʿamalu alladhīna yansubūna anfusahum ilā nafsī wa-yartakibūna mā
yanūḥu bihi qalbī wa-qalamī. Inna alladhīna yufsidūna fī al-arḍ
wa-yataṣarrafūna fī amwāli al-nās wa-yadkhulūna al-buyūta min ghayri
idhni ṣāḥibihā, innī barīʾun minhum illā an yatūbū wa-yarjiʿū ilā Allāhi
al-Ghafūri al-Raḥīm.

Elsewhere these words were revealed: ``We command the servants and
handmaidens of God to be chaste and God-fearing. Let them arise from the
slumber of desire and turn toward God, the Creator of earth and heaven.
Thus did We command the servants when the Daystar of the horizons shone
from the direction of Iraq. What harms Me is not My imprisonment, My
affliction, or what has been inflicted on Me by the tyrants among the
people, but rather the deeds of those who ascribe themselves to Me and
commit what causes My heart and My Pen to lament. Those who spread
corruption on earth, take possession of people's property, and enter
houses without their owners' permission---I disavow them unless they
repent and return to God, the Ever-Forgiving, the Merciful.''

\subsection{41}\label{section-1327}

\begin{otherlanguage}{arabic}

و در مقام دیگر یا ملأ الأرض سارعوا الى مرضاة اللّه و جاهدوا حقّ الجهاد فی
اظهار امره المبرم المتین قد قدّرنا الجهاد فی سبیل اللّه بجنود الحکمة و
البیان و بالأخلاق و الأعمال کذلک قضی الأمر من لدن قویّ قدیر لیس الفخر لمن
یفسد فى الأرض بعد اصلاحها اتّقوا اللّه یا قوم و لا تکونوا من الظّالمین

\end{otherlanguage}

Va dar maqām-e dīgar: Yā malaʾa al-arḍ, sāriʿū ilā marḍāti Allāh
wa-jāhidū ḥaqqa al-jihād fī iẓhāri amrihi al-mubrami al-matīn. Qad
qaddarnā al-jihāda fī sabīli Allāhi bi-junūdi al-ḥikmati wa-al-bayān
wa-bi-al-akhlāqi wa-al-aʿmāl. Kadhālika quḍiya al-amru min ladun
qawiyyin qadīr. Laysa al-fakhru li-man yufsidu fī al-arḍi baʿda
iṣlāḥihā. Ittaqū Allāha yā qawm wa-lā takūnū min al-ẓālimīn.

And elsewhere: ``O peoples of the earth! Hasten toward the good pleasure
of God, and strive with all your strength to manifest His irrevocable
and firmly established Cause. We have ordained that striving in the path
of God be undertaken with the hosts of wisdom and utterance, and with
good character and deeds. Thus has the decree been issued by One who is
Strong and All-Powerful. There is no glory for one who spreads
corruption on earth after it has been set right. Fear God, O people, and
do not be among the wrongdoers.''

\subsection{42}\label{section-1328}

\begin{otherlanguage}{arabic}

و در مقام دیگر لا تسبّوا احداً بینکم قد جئنا لاتّحاد من علی الأرض و اتّفاقهم
یشهد بذلک ما ظهر من بحر بیانی بین العباد ولکنّ القوم اکثرهم فى بعد مبین
ان یسبّکم احد او یمسّکم ضرّ فى سبیل اللّه اصبروا و توکّلوا علی السّامع البصیر
انّه یشهد و یری و یعمل ما اراد بسلطان من عنده انّه هو المقتدر القدیر قد
منعتم عن النّزاع و الجدال فی کتاب اللّه العزیز العظیم تمسّکوا بما تنتفع به
انفسکم و اهل العالم کذلک یأمرکم مالک القدم الظّاهر بالاسم الأعظم انّه هو
الآمر الحکیم

\end{otherlanguage}

Va dar maqām-e dīgar: Lā tasubbū aḥadan baynakum. Qad jiʾnā li-ittiḥādi
man ʿalā al-arḍi wa-ittifāqihim. Yashhadu bi-dhālika mā ẓahara min baḥri
bayānī bayna al-ʿibād, wa-lākinna al-qawma aktharuhum fī buʿdin mubīn.
In yasubbkum aḥadun aw yamasskum ḍurrun fī sabīli Allāh, iṣbirū
wa-tawakkalū ʿalā al-Sāmiʿi al-Baṣīr. Innahu yashhadu wa-yarā wa-yaʿmalu
mā arāda bi-sulṭānin min ʿindihi. Innahu huwa al-Muqtadiru al-Qadīr. Qad
muniʿtum ʿan al-nizāʿi wa-al-jidāli fī kitābi Allāhi al-ʿAzīzi al-ʿAẓīm.
Tamassakū bi-mā tanta fiʿu bihi anfusukum wa-ahlu al-ʿālam. Kadhālika
yaʾmurukum Māliku al-Qidam al-ẓāhiru bi-al-ismi al-Aʿẓam. Innahu huwa
al-Āmiru al-Ḥakīm.

And elsewhere: ``Do not revile anyone among you. We have come to bring
unity and concord to all who dwell on earth. What has appeared among the
servants from the ocean of My utterance bears witness to this, yet most
of the people remain manifestly far away. If anyone reviles you, or harm
touches you in the path of God, be patient and place your trust in the
All-Hearing, the All-Seeing. He indeed witnesses and sees, and does what
He wills through an authority of His own. He is indeed the All-Powerful,
the Almighty. You have been forbidden to engage in strife and contention
in the Book of God, the Mighty, the Great. Hold fast to what benefits
yourselves and the people of the world. Thus commands you the Lord of
Eternity, manifest through the Most Great Name. He is indeed the
Ordainer, the All-Wise.''

\subsection{43}\label{section-1329}

\begin{otherlanguage}{arabic}

و در مقام دیگر ایّاکم ان تسفکوا الدّمآء اخرجوا سیف اللّسان عن غمد البیان
لأنّ به تفتح مدائن القلوب انّا رفعنا حکم القتل عن بینکم انّ الرّحمة سبقت
الممکنات ان کنتم تعلمون

\end{otherlanguage}

Va dar maqām-i dīgar: Iyyākum an tasfikū al-dimāʾ. Akhrijū sayf al-lisān
ʿan ghimd al-bayān li-anna bihi tuftaḥu madāʾin al-qulūb. Innā rafaʿnā
ḥukm al-qatl ʿan baynikum. Inna al-raḥmata sabaqat al-mumkināt in kuntum
taʿlamūn.

And elsewhere: ``Beware of shedding blood. Draw the sword of the tongue
from the sheath of utterance, for through it the cities of hearts are
conquered. We have abolished the decree of killing among you. Mercy has
indeed preceded all created things, if you only knew.''

\subsection{44}\label{section-1330}

\begin{otherlanguage}{arabic}

و در مقام دیگر یا قوم لا تفسدوا فی الأرض و لا تسفکوا الدّمآء و لا تأکلوا
اموال النّاس بالباطل و لا تتّبعوا کلّ ناعق رجیم

\end{otherlanguage}

Va dar maqām-i dīgar: Yā qawm, lā tufsidū fī al-arḍ wa-lā tasfikū
al-dimāʾ wa-lā taʾkulū amwāl al-nās bi-l-bāṭil wa-lā tattabiʿū kulla
nāʿiq rajīm.

And elsewhere: ``O people! Do not spread corruption on earth, do not
shed blood, do not wrongfully consume the property of others, and do not
follow every outcast who raises a clamor.''

\subsection{45}\label{section-1331}

\begin{otherlanguage}{persian}

و در مقام دیگر آفتاب بیان الهی را غروب اخذ ننماید و افول از پی درنیاید
امروز این کلمهٔ علیا از سدرهٔ منتهى اصغا شد انّى لمن احبّنی و اخذ اوامری و
نبذ ما نهی عنه فى کتابی

\end{otherlanguage}

Va dar maqām-i dīgar: Āftāb-i bayān-i ilāhī rā ghurūb akhdh nanamāyad va
ufūl az pay darnayāyad. Imrūz īn kalimiy-i ʿulyā az Sidriy-i Muntahā
iṣghā shud: Innī li-man aḥabbanī wa-akhadha awāmirī wa-nabadha mā nuhiya
ʿanhu fī kitābī.

And elsewhere: ``The sun of divine utterance will never set, nor will
its decline follow. Today this exalted utterance was heard from the Lote
Tree beyond which there is no passing: `I belong to whoever loves Me,
holds fast to My commandments, and casts aside what has been forbidden
in My Book.'\,''

\subsection{46}\label{section-1332}

\begin{otherlanguage}{persian}

و در مقام دیگر امروز روز ذکر و ثنا و روز خدمتست خود را محروم منمائید
شمائید حروفات کلمات و کلمات کتاب و شما نهالهائی هستید که از دست عنایت در
ارض رحمت کشته شده‌اید و از امطار کرم نموّ نموده‌اید شما را از عاصفات شرک و
قاصفات کفر حفظ فرمود و بایادی شفقت تربیت نمود حال وقت اثمار و اوراقست و
اثمار سدرهٔ انسانی اعمال طیّبه و اخلاق مرضیّه بوده و هست این اثمار را از
غافلین منع منمائید اگر پذیرفتند مقصود حاصل و حیات ظاهر والّا ذروهم فى
خوضهم یلعبون یا حزب اللّه جهد نمائید شاید قلوب احزاب مختلفهٔ عالم بآب
بردباری و شفقت شما از ضغینه و بغضا پاک و پاکیزه شود و قابل و لائق تجلّیات
آفتاب حقیقت گردد

\end{otherlanguage}

Va dar maqām-i dīgar: Imrūz rūz-i dhikr va thanā va rūz-i khidmat ast.
Khud rā maḥrūm manamāyīd. Shumāyīd ḥurūfāt-i kalimāt va kalimāt-i kitāb
va shumā nihālhāʾī hastīd ke az dast-i ʿināyat dar arḍ-i raḥmat keshte
shude-īd va az amṭār-i karam numuv namūde-īd. Shumā rā az ʿāṣifāt-i
shirk va qāṣifāt-i kufr ḥifẓ farmūd va bi-ayādī-yi shafaqat tarbiyat
namūd. Ḥāl vaqt-i ithmār va awrāq ast va ithmār-i sidriy-i insānī
aʿmāl-i ṭayyibe va akhlāq-i marḍiyye būde va hast. Īn ithmār rā az
ghāfilīn manʿ manamāyīd. Agar padhīruftand maqṣūd ḥāṣil va ḥayāt ẓāhir
va-illā dharhum fī khawḍihim yalʿabūn. Yā ḥizb Allāh, jahd namāyīd
shāyad qulūb-i aḥzāb-i mukhtalifiy-i ʿālam bi-āb-i burdbārī va
shafaqat-i shumā az ḍaghīne va baghḍā pāk va pākīze shavad va qābil va
lāyiq-i tajalliyāt-i āftāb-i ḥaqīqat gardad.

And elsewhere: ``Today is the day of remembrance and praise, and the day
of service. Do not deprive yourselves. You are the letters of the words
and the words of the Book, and you are saplings planted by the hand of
loving care in the soil of mercy, which have grown through the rains of
bounty. He protected you from the tempests of associating partners with
God and the shattering winds of unbelief, and nurtured you with the
hands of compassion. Now is the time for fruit and leaves, and the
fruits of the human tree have always been, and remain, good deeds and
pleasing qualities. Do not withhold these fruits from the heedless. If
they accept them, the purpose is achieved and life is made manifest;
otherwise, leave them to amuse themselves in their idle talk. O people
of God! Strive, so that the hearts of the diverse peoples of the world
may, through the water of your forbearance and compassion, be cleansed
and purified of malice and hatred, and become fit and worthy to receive
the splendors of the Sun of Truth.''

\subsection{47}\label{section-1333}

\begin{otherlanguage}{persian}

در اشراق چهارم از اشراقات ذکر نمودیم از برای هر امری ناصری لازم و جنود
منصوره در این ظهور اعمال و اخلاق پسندیده است و قائد و سردار این جنود
تقوی اللّه بوده اوست دارای کل و حاکم بر کل

\end{otherlanguage}

Dar Ishrāq-i chahārum az Ishrāqāt dhikr namūdīm: Az barāy-i har amrī
nāṣirī lāzim va junūd-i manṣūre dar īn ẓuhūr aʿmāl va akhlāq-i pasandīde
ast va qāʾid va sardār-i īn junūd taqvā-yi Allāh būde. Ūst dārā-yi kull
va ḥākim bar kull.

In the fourth Ishráq of the Ishráqát We stated: ``Every cause requires a
helper, and the victorious hosts in this Revelation are commendable
deeds and qualities. The leader and commander of these hosts is the fear
of God. It is the sovereign of all and rules over all.''

\subsection{48}\label{section-1334}

\begin{otherlanguage}{persian}

در تجلّی سوّم از کتاب تجلّیات ذکر نمودیم سبب علوّ وجود و سموّ آن علوم و فنون
و صنائع است علم بمنزلهٔ جناحست از برای وجود و مرقاتست از برای صعود تحصیلش
بر کل لازم ولکن علومی که اهل ارض از آن منتفع شوند نه علومی که از حرف
ابتدا شود و بحرف منتهی گردد صاحبان علوم و صنائع را حقّ عظیم است بر اهل
عالم یشهد بذلک امّ البیان فى هذا المقام المبین

\end{otherlanguage}

Dar Tajallī-yi sivvum az Kitāb-i Tajalliyāt dhikr namūdīm: Sabab-i
ʿuluvv-i vujūd va sumuvv-i ān ʿulūm va funūn va ṣanāyiʿ ast. ʿIlm
bi-manziliy-i janāḥ ast az barāy-i vujūd va mirqāt ast az barāy-i ṣuʿūd.
Taḥṣīlash bar kull lāzim va-lākin ʿulūmī ke ahl-i arḍ az ān muntafiʿ
shavand na ʿulūmī ke az ḥarf ibtidā shavad va bi-ḥarf muntahī gardad.
Ṣāḥibān-i ʿulūm va ṣanāyiʿ rā ḥaqq-i ʿaẓīm ast bar ahl-i ʿālam. Yashhadu
bi-dhālika Umm al-Bayān fī hādhā al-maqām al-mubīn.

In the third Tajallí of the Book of Tajallíyát We stated: ``The
sciences, arts, and crafts are the means by which existence is exalted
and elevated. Knowledge is like wings for existence and a ladder for
ascent. Its acquisition is incumbent upon everyone, but it must be
knowledge from which the people of the earth benefit, not knowledge that
begins with words and ends with words. Those versed in the sciences and
crafts have a great claim upon the people of the world. To this the
Mother of Utterance bears witness in this manifest station.''

\subsection{49}\label{section-1335}

\begin{otherlanguage}{arabic}

فی‌الحقیقه کنز حقیقی از برای انسان علم اوست و اوست علّت عزّت و نعمت و فرح و
نشاط و بهجت و انبساط طوبی لمن تمسّک به و ویل للغافلین

\end{otherlanguage}

Fī al-ḥaqīqe kanz-i ḥaqīqī az barāy-i insān ʿilm-i ūst va ūst ʿillat-i
ʿizzat va niʿmat va faraḥ va nishāṭ va bahjat va inbisāṭ. Ṭūbā li-man
tamassaka bihi wa-wayl li-l-ghāfilīn.

In truth, a human being's real treasure is knowledge. It is the source
of honor and prosperity, joy and delight, gladness and cheer. Blessed is
whoever holds fast to it, and woe to the heedless.

\subsection{50}\label{section-1336}

\begin{otherlanguage}{persian}

آن جناب باید در جمیع احوال عباد را بآنچه سبب ظهور اخلاق روحانیّه و اعمال
طیّبه است دعوت فرمایند تا کل آگاه شوند بآنچه که سبب ارتفاع وجود است و
بهمّت کامل قصد مقام اعلى و ذروهٔ علیا نمایند و آنچه سبب اوّل است از برای
تربیت خلق خشیة اللّه بوده طوبی للفائزین

\end{otherlanguage}

Ān janāb bāyad dar jamīʿ-i aḥvāl ʿibād rā bi-ānche sabab-i ẓuhūr-i
akhlāq-i rūḥāniyye va aʿmāl-i ṭayyibe ast daʿvat farmāyand tā kull āgāh
shavand bi-ānche ke sabab-i irtifāʿ-i vujūd ast va bi-himmat-i kāmil
qaṣd-i maqām-i aʿlā va dhurviy-i ʿulyā namāyand va ānche sabab-i avval
ast az barāy-i tarbiyat-i khalq khashyat-i Allāh būde. Ṭūbā
li-l-fāʾizīn.

You should, in all circumstances, summon the servants to whatever brings
forth spiritual qualities and good deeds, so that all may become aware
of what exalts human existence and, with full resolve, seek the most
exalted station and the loftiest summit. The foremost means of educating
humanity has been the fear of God. Blessed are those who attain it.

\subsection{51}\label{section-1337}

\begin{otherlanguage}{persian}

کلمهٔ اولی که در ورق اوّل فردوس از قلم ابهی مذکور و مسطور اینست براستی
میگویم حفظ مبین و حصن متین از برای عموم اهل عالم خشیة اللّه بوده اوست سبب
اکبر از برای حفظ بشر و علّت کبری از برای صیانت وری بلی در وجود آیتی موجود
و آن انسان را از آنچه شایسته و لائق نیست منع مینماید و حراست میفرماید و
نام آن را حیا گذارده‌اند ولکن این فقره محدود است بمعدودی کل دارای این
مقام نبوده و نیستند باید سلاطین ایّام و علمای انام به دین تمسّک نمایند چه
که اوست علّت ظهور خشیة اللّه فیما سواه

\end{otherlanguage}

Kalimiy-i ūlā ke dar varaq-i avval-i Firdaws az qalam-i Abhā madhkūr va
masṭūr īn ast: Bi-rāstī mīgūyam ḥifẓ-i mubīn va ḥiṣn-i matīn az barāy-i
ʿumūm-i ahl-i ʿālam khashyat-i Allāh būde. Ūst sabab-i akbar az barāy-i
ḥifẓ-i bashar va ʿillat-i kubrā az barāy-i ṣiyānat-i varā. Balī dar
vujūd āyatī mawjūd va ān insān rā az ānche shāyiste va lāyiq nīst manʿ
mīnamāyad va ḥirāsat mīfarmāyad va nām-i ān rā ḥayā gudhāshte-and
va-lākin īn faqare maḥdūd ast bi-maʿdūdī. Kull dārā-yi īn maqām nabūde
va nīstand. Bāyad salāṭīn-i ayyām va ʿulamā-yi anām bi-dīn tamassuk
namāyand che ke ūst ʿillat-i ẓuhūr-i khashyat-i Allāh fīmā sivāh.

The first word set down and inscribed by the Most Glorious Pen on the
first leaf of Paradise is this: ``Truly I say, the sure protection and
strong fortress for all the people of the world is the fear of God. It
is the greatest means of protecting humanity and the supreme cause of
safeguarding humankind. There is indeed a sign within human beings that
restrains and guards them from what is unseemly and unworthy, and it has
been named a sense of shame. Yet this is confined to a few; not everyone
has possessed or possesses this station. The rulers of the age and the
learned among the people must hold fast to religion, for it is the means
by which the fear of God appears in all besides Him.''

\subsection{52}\label{section-1338}

\begin{otherlanguage}{persian}

کلمهٔ ثانی که در ورق ثانی از فردوس ذکر نمودیم اینست قلم بیان در این حین
مظاهر قدرت و مشارق اقتدار یعنی حضرات ملوک و سلاطین ایّدهم اللّه را نصیحت
مینماید و به دین و تمسّک بآن وصیّت میکند اوست سبب بزرگ از برای نظم جهان و
اطمینان من فی الامکان سستی ارکان دین سبب قوّت جهّال و جرئت و جسارت شده
براستى میگویم آنچه از مقام بلند دین کاست بر غفلت اشرار افزود و نتیجه
بالأخره هرج و مرج است اسمعوا یا اولی الأبصار ثمّ اعتبروا یا اولی الأنظار

\end{otherlanguage}

Kalimiy-i thānī ke dar varaq-i thānī az Firdaws dhikr namūdīm īn ast:
Qalam-i bayān dar īn ḥīn maẓāhir-i qudrat va mashāriq-i iqtidār yaʿnī
ḥaḍarāt-i mulūk va salāṭīn, ayyadahum Allāh, rā naṣīḥat mīnamāyad va
bi-dīn va tamassuk bi-ān vaṣiyyat mīkunad. Ūst sabab-i buzurg az barāy-i
naẓm-i jahān va iṭmīnān-i man fī al-imkān. Sustī-yi arkān-i dīn sabab-i
quvvat-i juhhāl va jurʾat va jisārat shude. Bi-rāstī mīgūyam ānche az
maqām-i buland-i dīn kāst bar ghaflat-i ashrār afzūd va natīje
bi-l-ākhare harj va marj ast. Ismaʿū yā ulī al-abṣār thumma iʿtabirū yā
ulī al-anẓār.

The second word We mentioned on the second leaf of Paradise is this:
``At this moment the Pen of Utterance counsels the manifestations of
power and the dawning places of authority---that is, the kings and
sovereigns, may God aid them---and exhorts them to uphold religion and
hold fast to it. It is the great means of establishing order in the
world and tranquility among all who dwell in it. The weakening of the
pillars of religion has strengthened the ignorant and made them bold and
insolent. Truly I say, whatever has diminished the lofty station of
religion has increased the heedlessness of the wicked, and the eventual
outcome is chaos and confusion. Listen, O people of insight, and take
heed, O people of discernment.''

\subsection{53}\label{section-1339}

\begin{otherlanguage}{persian}

امید آنکه آن جناب در آنچه ذکر شد باذن واعیه اصغا فرمایند شاید عباد را از
ما عندهم به ما عند اللّه آگاه نمایند از حقّ میطلبیم نور انصاف و آفتاب عدل
را از سحاب تیرهٔ غفلت نجات بخشد و ظاهر فرماید هیچ نوری بنور عدل معادله
نمینماید آنست سبب نظم عالم و راحت امم

\end{otherlanguage}

Umīd ānke ān janāb dar ānche dhikr shud bi-udhun-i vāʿiye iṣghā
farmāyand shāyad ʿibād rā az mā ʿindahum bi-mā ʿinda Allāh āgāh
namāyand. Az Ḥaqq mīṭalabīm nūr-i inṣāf va āftāb-i ʿadl rā az saḥāb-i
tīriy-i ghaflat najāt bakhshad va ẓāhir farmāyad. Hīch nūrī bi-nūr-i
ʿadl muʿādile namīnamāyad. Ān ast sabab-i naẓm-i ʿālam va rāḥat-i umam.

It is hoped that you will listen with an attentive ear to what has been
stated, so that you may turn the servants' awareness from what they
possess to what is with God. We ask God to free the light of fairness
and the sun of justice from the dark clouds of heedlessness and make
them manifest. No light can equal the light of justice. It is the means
of establishing order in the world and bringing tranquility to the
nations.

\subsection{54}\label{section-1340}

\begin{otherlanguage}{persian}

در صحیفهٔ بیان این کلمهٔ علیا مسطور و مرقوم بگو ای دوستان جهد نمائید شاید
مصیباتی که فی سبیل اللّه بر مظلوم و شما وارد شده بین ناس ضایع نشود بذیل
عفّت تمسّک نمائید و همچنین بحبل امانت و دیانت صلاح عالم را ملاحظه نمائید
نه هوای نفس را یا حزب المظلوم شمائید رعاة عالم اغنام را از ذئب نفس و هوی
مقدّس دارید و بطراز تقوی اللّه مزیّن نمائید اینست حکم محکم که از قلم قدم در
این حین جاری شده لعمر اللّه سیف الأخلاق و الآداب احدّ من سیوف الحدید انّ
الفطرة فی هذا الحین تنادی و تقول یا قوم قد اتی الیوم و اظهرنی ربّی بنور
کسفت عند اشراقه شموس البیان اتّقوا الرّحمن و لا تکونوا من الغافلین

\end{otherlanguage}

Dar Ṣaḥīfiy-i Bayān īn kalimiy-i ʿulyā masṭūr va marqūm: Bigū ay dūstān
jahd namāyīd shāyad muṣībatī ke fī sabīl-i Allāh bar Maẓlūm va shumā
vārid shude bayn-i nās ḍāyiʿ nashavad. Bi-dhayl-i ʿiffat tamassuk
namāyīd va hamchunīn bi-ḥabl-i amānat va diyānat. Ṣalāḥ-i ʿālam rā
mulāḥiẓe namāyīd na havā-yi nafs rā. Yā ḥizb al-Maẓlūm, shumāyīd ruʿāt-i
ʿālam. Aghnām rā az dhiʾb-i nafs va havā muqaddas dārīd va bi-ṭirāz-i
taqvā-yi Allāh muzayyan namāyīd. Īn ast ḥukm-i muḥkam ke az qalam-i
qidam dar īn ḥīn jārī shude. La-ʿamr Allāh, sayf al-akhlāq wa-l-ādāb
aḥaddu min suyūf al-ḥadīd. Inna al-fiṭrata fī hādhā al-ḥīn tunādī
wa-taqūl: Yā qawm, qad atā al-yawm wa-aẓharanī rabbī bi-nūr kusifat
ʿinda ishrāqihi shumūs al-bayān. Ittaqū al-Raḥmān wa-lā takūnū min
al-ghāfilīn.

In the Scroll of Utterance this exalted word is written and inscribed:
``Say: O friends! Strive, so that the afflictions that have befallen the
Wronged One and you in the path of God may not be wasted among the
people. Hold fast to the hem of chastity, and likewise to the cord of
trustworthiness and piety. Consider the welfare of the world, not your
own selfish desires. O people of the Wronged One! You are the shepherds
of the world. Keep the sheep safe from the wolf of self and desire, and
adorn them with the ornament of the fear of God. This is the firm decree
that has flowed at this moment from the Pen of Eternity. By the life of
God! The sword of good character and conduct is sharper than swords of
iron. Human nature is at this moment calling out and saying: `O people!
The Day has come, and my Lord has made me manifest with a light whose
dawning has eclipsed the suns of utterance. Fear the All-Merciful, and
do not be among the heedless.'\,''

\subsection{55}\label{section-1341}

\begin{otherlanguage}{persian}

کلمهٔ سوّم که در ورق سوّم فردوس ذکر نمودیم اینست یا ابن الانسان لو تکون
ناظراً الی الفضل ضع ما ینفعک و خذ ما ینتفع به العباد و ان تکن ناظراً الى
العدل اختر لدونک ما تختاره لنفسک انّ الانسان مرّةً یرفعه الخضوع الی سمآء
العزّة و الاقتدار و اخری ینزله الغرور الی مقام الذّلّة و الانکسار یوم عظیم
است و ندا بزرگ در لوحی از الواح این کلمهٔ علیا را ذکر نمودیم اگر عالم روح
بتمامه بقوّهٔ سامعه تبدیل شود میتوان گفت لائق اصغاء این ندا است که از افق
اعلی مرتفع والّا این آذانهای آلوده بقصص کاذبه لائق نبوده و نیست طوبی
للسّامعین و ویل للغافلین

\end{otherlanguage}

Kalimiy-i sivvum ke dar varaq-i sivvum-i Firdaws dhikr namūdīm īn ast:
Yā ibn al-insān, law takūnu nāẓiran ilā al-faḍl ḍaʿ mā yanfaʿuka
wa-khudh mā yantafiʿu bihi al-ʿibād wa-in takun nāẓiran ilā al-ʿadl
ikhtar li-dūnika mā takhtāruhu li-nafsik. Inna al-insāna marratan
yarfaʿuhu al-khuḍūʿ ilā samāʾ al-ʿizza wa-l-iqtidār wa-ukhrā yunziluhu
al-ghurūr ilā maqām al-dhilla wa-l-inkisār. Yawm ʿaẓīm ast va nidā
buzurg. Dar lawḥī az alvāḥ īn kalimiy-i ʿulyā rā dhikr namūdīm: Agar
ʿālam-i rūḥ bi-tamāmihi bi-quvviy-i sāmiʿe tabdīl shavad mītavān guft
lāyiq-i iṣghā-yi īn nidā ast ke az ufuq-i aʿlā murtafiʿ va-illā īn
ādhānhā-yi ālūde bi-qiṣaṣ-i kādhibe lāyiq nabūde va nīst. Ṭūbā
li-l-sāmiʿīn wa-wayl li-l-ghāfilīn.

The third word We mentioned on the third leaf of Paradise is this: ``O
son of man! If you look to grace, set aside what benefits you and take
up what benefits the servants; and if you look to justice, choose for
others what you choose for yourself. At one time, humility raises a
human being to the heaven of glory and power; at another, pride brings
one down to the station of abasement and humiliation.'' The Day is
momentous, and the call is mighty. In one of the Tablets We stated this
exalted word: ``If the entire world of the spirit were transformed into
the faculty of hearing, it could be said to be worthy of hearing this
call that has risen from the Most Exalted Horizon. Otherwise, these
ears, defiled by false tales, have never been and are not worthy of it.
Blessed are those who listen, and woe to the heedless.''

\subsection{56}\label{section-1342}

\begin{otherlanguage}{persian}

از حقّ جلّ جلاله سائل و آمل که مشارق ثروت و اقتدار و مطالع عزّت و اختیار
یعنی ملوک ارض ایّدهم اللّه را بر صلح اکبر تأیید فرماید اینست سبب اعظم از
برای راحت امم سلاطین آفاق وفّقهم اللّه باید باتّفاق باین امر که سبب اعظم
است از برای حفظ عالم تمسّک فرمایند امید آنکه قیام نمایند بر آنچه سبب
آسایش عباد است باید مجلس بزرگی ترتیب دهند و حضرات ملوک و یا وزرا در آن
مجلس حاضر شوند و حکم اتّحاد و اتّفاق را جاری فرمایند و از سلاح باصلاح توجّه
کنند و اگر سلطانی بر سلطانی برخیزد سلاطین دیگر بر منع او قیام نمایند در
این صورت عساکر و آلات و ادوات حرب لازم نه الّا علی قدر مقدور لحفظ بلادهم
و اگر باین خیر اعظم فائز شوند اهل مملکت کل براحت و مسرّت بامور خود مشغول
گردند و نوحه و ندبهٔ اکثری از عباد ساکن شود نسأل اللّه ان یؤیّدهم على ما
یحبّ و یرضی انّه هو ربّ العرش و الثّری و مالک الآخرة و الأولی احسن و اولی
آنکه در آن مجلس خود سلاطین عظام حاضر شوند و حکم فرمایند و هر یک از
سلاطین که بر این امر و اجرای آن قیام فرماید او سیّد سلاطین است عند اللّه
طوبی له و نعیماً له

\end{otherlanguage}

Az Ḥaqq, jalla jalāluhu, sāʾil va āmil ke mashāriq-i tharvat va iqtidār
va maṭāliʿ-i ʿizzat va ikhtiyār yaʿnī mulūk-i arḍ, ayyadahum Allāh, rā
bar ṣulḥ-i akbar taʾyīd farmāyad. Īn ast sabab-i aʿẓam az barāy-i
rāḥat-i umam. Salāṭīn-i āfāq, waffaqahum Allāh, bāyad bi-ittifāq bi-īn
amr ke sabab-i aʿẓam ast az barāy-i ḥifẓ-i ʿālam tamassuk farmāyand.
Umīd ānke qiyām namāyand bar ānche sabab-i āsāyish-i ʿibād ast. Bāyad
majlis-i buzurgī tartīb dihand va ḥaḍarāt-i mulūk va yā vuzarā dar ān
majlis ḥāḍir shavand va ḥukm-i ittiḥād va ittifāq rā jārī farmāyand va
az silāḥ bi-iṣlāḥ tavajjuh kunand va agar sulṭānī bar sulṭānī barkhīzad
salāṭīn-i dīgar bar manʿ-i ū qiyām namāyand. Dar īn ṣūrat ʿasākir va
ālāt va adavāt-i ḥarb lāzim na illā ʿalā qadr maqdūr li-ḥifẓ bilādihim
va agar bi-īn khayr-i aʿẓam fāʾiz shavand ahl-i mamlakat kull bi-rāḥat
va masarrat bi-umūr-i khud mashghūl gardand va nawḥe va nudbiy-i aktharī
az ʿibād sākin shavad. Nasʾalu Allāh an yuʾayyidahum ʿalā mā yuḥibbu
wa-yarḍā innahu huwa rabb al-ʿarsh wa-l-tharā wa-mālik al-ākhira
wa-l-ūlā. Aḥsan va awlā ānke dar ān majlis khud-i salāṭīn-i ʿiẓām ḥāḍir
shavand va ḥukm farmāyand va har yak az salāṭīn ke bar īn amr va ijrā-yi
ān qiyām farmāyad ū sayyid-i salāṭīn ast ʿinda Allāh. Ṭūbā lahu
wa-naʿīman lahu.

We ask and hope that God, exalted be His glory, will aid the dawning
places of wealth and power and the sources of honor and authority---that
is, the kings of the earth, may God aid them---to establish the Most
Great Peace. This is the greatest means of bringing tranquility to the
nations. The sovereigns of the world, may God grant them success, must
jointly uphold this cause, which is the greatest means of protecting the
world. It is hoped that they will arise to do what brings ease to the
servants. They must convene a great assembly, and the kings or their
ministers must attend it, enforce unity and concord, and turn from arms
to reconciliation. If one sovereign rises against another, the other
sovereigns must arise to restrain him. In that case, armies and the
instruments and equipment of war would be unnecessary, except for a
limited measure to protect their countries. If they attain this greatest
good, all the people of the realm will pursue their affairs in comfort
and happiness, and the lamentation and mourning of most of the servants
will be stilled. We ask God to aid them in what He loves and approves.
He is indeed the Lord of the throne and the earth, and the Master of the
next world and this one. It would be better and more fitting for the
great sovereigns themselves to attend that assembly and issue its
decrees. Whichever sovereign arises to undertake this cause and carry it
out is, in God's sight, the foremost of sovereigns. Blessed is he, and
bliss be his.

\subsection{57}\label{section-1343}

\begin{otherlanguage}{arabic}

در این اراضی هر هنگام نوبت عسکریّه میشود از هر جهتی فزع اکبر ظاهر و هر یک
از دول در هر سنه بر عساکر میافزاید چه که وزرای حرب در این فقره حریصند و
به هل من مزید ناطق از قرار مذکور دولت ایران ایّدها اللّه هم اراده نموده‌اند
بر عسکر بیفزایند نزد مظلوم حال صدهزار عسکر منظّم مرتّب کفایت مینماید امید
آنکه آن حضرت بر نور عدل بیفزایند لعمر اللّه عدل جندیست قویّ اوست در مقام
اوّل و رتبهٔ اولی فاتح افئده و قلوب و اوست مبیّن اسرار وجود و دارای رایت
محبّت و جود

\end{otherlanguage}

Dar īn arāḍī har hangām nawbat-i ʿaskariyye mīshavad az har jihatī
fazaʿ-i akbar ẓāhir va har yak az duval dar har sane bar ʿasākir
mīafzāyad che ke vuzarā-yi ḥarb dar īn faqare ḥarīṣand va bi-hal min
mazīd nāṭiq. Az qarār-i madhkūr dawlat-i Īrān, ayyadahā Allāh, ham irāde
namūde-and bar ʿaskar biyafzāyand. Nazd-i Maẓlūm ḥāl ṣad hazār ʿaskar-i
munaẓẓam-i murattab kifāyat mīnamāyad. Umīd ānke ān ḥaḍrat bar nūr-i
ʿadl biyafzāyand. La-ʿamr Allāh, ʿadl jundī-st qavī. Ūst dar maqām-i
avval va rutbiy-i ūlā fātiḥ-i afʾide va qulūb va ūst mubayyin-i asrār-i
vujūd va dārā-yi rāyat-i maḥabbat va jūd.

In these lands, whenever the time for military conscription comes, the
greatest alarm appears on every side. Each state increases its forces
every year, for the ministers of war are eager in this matter and keep
saying, ``Is there more?'' According to reports, the government of Iran,
may God aid it, also intends to increase its army. In the view of the
Wronged One, one hundred thousand well-organized and disciplined troops
would suffice at present. It is hoped that His Majesty will increase the
light of justice. By the life of God! Justice is a mighty host. It is,
first and foremost, the conqueror of hearts and souls; it reveals the
mysteries of existence and bears the banner of love and generosity.

\subsection{58}\label{section-1344}

\begin{otherlanguage}{persian}

و در خزائن علوم الهی علمی موجود که عمل بآن خوف را زائل مینماید ولکن على
قدر مقدور و آن ترتیبی است که باید از ایّام طفولیّت بآن عمل شود بسیار فرق
مینماید و آنچه از آن بکاهد بر جرئت میافزاید اگر ارادهٔ الهی مدد فرماید
شرح مبسوطى در آنچه ذکر شد از قلم بیان جاری میگردد و شاید در مراتب علوم و
فنون هم ظاهر شود آنچه که سبب تجدید عالم و امم گردد و همچنین کلمه‌ئی در
صحیفهٔ حمرا از قلم اعلی مرقوم و مسطور و آن کلمه قوّت مکنونهٔ در عباد را
بتمامها ظاهر فرماید بلکه مثل آن بر آن بیفزایدنسأل اللّه تبارک و تعالی ان
یؤیّد عباده علی ما یحبّ و یرضی

\end{otherlanguage}

Va dar khazāʾen-e ʿolūm-e elāhī ʿelmī mawjūd ke ʿamal be ān khawf rā
zāʾel mīnamāyad valāken ʿalā qadr-e maqdūr va ān tartībī ast ke bāyad az
ayyām-e ṭofūliyyat be ān ʿamal shavad besyār farq mīnamāyad va ānche az
ān bekāhad bar jorʾat mīafzāyad agar erāde-ye elāhī madad farmāyad
sharḥ-e mabsūṭī dar ānche dhekr shod az qalam-e bayān jārī mīg ardad va
shāyad dar marāteb-e ʿolūm va fonūn ham ẓāher shavad ānche ke sabab-e
tajdīd-e ʿālam va omam gardad va hamchenīn kalameʾī dar ṣaḥīfe-ye ḥamrā
az qalam-e aʿlā marqūm va masṭūr va ān kalame qovvat-e maknūne-ye dar
ʿebād rā be tamāmhā ẓāher farmāyad balke methl-e ān bar ān bīafzāyad.
Nasʾalu Allāha tabāraka wa-taʿālā an yuʾayyida ʿibādahu ʿalā mā yuḥibbu
wa-yarḍā.

And in the treasuries of divine knowledge there is a knowledge whose
application dispels fear, though only to the extent possible. It is a
method that must be practiced from childhood. It makes a great
difference, and whatever it diminishes of fear it adds to courage. If
the divine will lends its aid, a detailed explanation of what has been
mentioned will flow from the pen of utterance. Perhaps there will also
appear, in the various branches of the sciences and arts, that which
will bring about the renewal of the world and its nations. Likewise, a
word has been inscribed and recorded by the Most Exalted Pen in the
Crimson Scroll. That word will reveal in their entirety the powers
latent in human beings; indeed, it will add as much again to them. We
ask God, blessed and exalted be He, to aid His servants in what He loves
and approves.

\subsection{59}\label{section-1345}

\begin{otherlanguage}{persian}

این ایّام از جمیع جهات اعدا ظاهر و نار بغضا مشتعل ای اهل ارض لعمری و
عمرکم این مظلوم خیال ریاست نداشته و ندارد و مقصود رفع آنچه سبب اختلاف
احزاب عالم و تفریق امم است بوده و هست تا کل فارغ و آزاد شوند و بخود
پردازند استدعا آنکه دوستان این ذیل را بغبار اکاذیب نیالایند و بذکر خوارق
عادات که نزد ایشان است از شأن و مقام و تقدیس و تنزیه نکاهند

\end{otherlanguage}

Īn ayyām az jamīʿ-e jehāt aʿdā ẓāher va nār-e baghḍā moshtaʿel ey ahl-e
arḍ la-ʿamrī wa-ʿamrikum īn maẓlūm khayāl-e reyāsat nadāshte va nadārad
va maqṣūd rafʿ-e ānche sabab-e ekhtelāf-e aḥzāb-e ʿālam va tafrīq-e omam
ast būde va hast tā koll fāregh va āzād shavand va be khod pardāzand
estedʿā ānke dūstān īn dheyl rā be ghobār-e akādhīb nayālāyand va be
dhekr-e khavāreq-e ʿādāt ke nazd-e īshān ast az shaʾn va maqām va taqdīs
va tanzīh nakāhand.

These days enemies have appeared on every side, and the fire of hatred
is ablaze. People of the earth! By my life and yours, this Wronged One
has never entertained, nor does he now entertain, any ambition for
leadership. His purpose has been and remains to remove whatever causes
discord among the peoples of the world and division among its nations,
so that all may be free and at ease and attend to their own affairs. The
friends are asked not to soil this hem with the dust of falsehoods, or
to diminish his dignity, station, sanctity, and purity by recounting the
miracles they ascribe to him.

\subsection{60}\label{section-1346}

\begin{otherlanguage}{persian}

سبحان‌اللّه امروز روزیست که عقلا از این مظلوم اخذ رأی نمایند و از حقّ
بطلبند آنچه را که سبب عزّت و آسایش است ولکن کل برعکس بر اطفاء این نور
ساطع لمیع جاهد و ساعی و هر نفسی در صدد آنست که تقصیری ثابت نماید و یا
ایرادی وارد آورد امر بمقامی رسیده که از جمیع اطوار این مظلوم اظهار
نموده‌اند آنچه را که ذکر آن شایسته نیست یکی از دوستان از نفسی از نفوس
مجتمعه در مدینهٔ کبیره ذکر نموده که گفته در هر سنه مبلغ پنجاه‌هزار تومان
از وطن به عکّا میرود و بکمال افسوس این فقره را ذکر نمود معلوم نشد محاسب
که بوده و دفتردار که

\end{otherlanguage}

S obḥān Allāh emrūz rūzīst ke ʿoqalā az īn maẓlūm akhdh-e raʾy namāyand
va az Ḥaqq beṭalaband ānche rā ke sabab-e ʿezzat va āsāyesh ast valāken
koll barʿaks bar eṭfāʾ-e īn nūr-e sāṭeʿ-e lamīʿ jāhed va sāʿī va har
nafsī dar ṣadad-e ānast ke taqṣīrī thābet namāyad va yā īrādī vāred
āvarad amr be maqāmī resīde ke az jamīʿ-e aṭvār-e īn maẓlūm eẓhār
namūdeand ānche rā ke dhekr-e ān shāyeste nīst yekī az dūstān az nafsī
az nofūs-e mojtameʿe dar madīne-ye kabīre dhekr namūde ke gofte dar har
sane mablagh-e panjāh-hezār tūmān az vaṭan be ʿAkkā mīravad va be
kamāl-e afsūs īn faqare rā dhekr namūd maʿlūm nashod moḥāseb ke būde va
daftardār ke.

Glory be to God! This is the day when the wise should seek the counsel
of this Wronged One and ask God for what brings honor and well-being.
Yet all, on the contrary, are striving and laboring to extinguish this
radiant, shining light, and everyone seeks to establish some fault or
raise some objection. Matters have reached the point where they have
said things about every aspect of this Wronged One's conduct that are
unworthy of mention. One of the friends reported that a person among
those gathered in the Great City had said that every year a sum of fifty
thousand tumans goes from the homeland to Acre. He mentioned this with
the deepest regret. It is not known who the accountant was, or who the
bookkeeper.

\subsection{61}\label{section-1347}

\begin{otherlanguage}{persian}

باری این مظلوم در جمیع آنچه وارد آورده‌اند و گفته‌اند صابر و صامت چه که
اراده آنکه از عنایت حقّ جلّ جلاله و رحمت مسبوقه حکم جدال و نزاع و سفک دماء
را از عالم بقوّهٔ بیان محو نمائیم در جمیع احوال در آنچه گفته‌اند بصبر جمیل
تمسّک جسته‌ایم و بحقّ گذارده‌ایم ولکن در جواب این فقره ذکر نمودیم اگر
فی‌الحقیقه در آنچه ذکر نموده صادقست باید مالک وجود و سلطان غیب و شهود را
شکر نماید که نفسی را از ایران مبعوث فرموده که در سجن من غیر ناصر و معین
ایران را تصرّف فرموده و یک مالیات هر سنه از آن اخذ میفرماید این مقام
ستایش است نه نکوهش لو یکون من المنصفین و اگر نفسی اراده نماید بر امور
این مظلوم واقف شود جواب اینست که در بعضی از لیالی و ایّام اسرای عالم و
مظلومهای امم بی‌قوت شب را بروز آورده‌اند و روز را بشب این اذکار را دوست
نداشته و نداریم و از گوینده هم شکایت ننموده و نمینمائیم در این سجن شخص
معزّزی چندی از سنگ‌تراشی کسب معاش مینمود و در بعضى از اوقات هم بعضی بطعام
اللّه مرزوق یعنى جوع نسأل اللّه تبارک و تعالی ان یؤیّد الکلّ علی العدل و
الانصاف و یوفّقهم علی الانابة و الرّجوع انّه هو السّامع المجیب

\end{otherlanguage}

Bārī īn maẓlūm dar jamīʿ-e ānche vāred āvardeand va gofteand ṣāber va
ṣāmet che ke erāde ānke az ʿenāyat-e Ḥaqq jalla jalāluhu va raḥmat-e
masbūqe ḥokm-e jedāl va nezāʿ va safk-e demā rā az ʿālam be qovve-ye
bayān maḥv namāyīm dar jamīʿ-e aḥvāl dar ānche gofteand be ṣabr-e jamīl
tamassok josteʾīm va be Ḥaqq godhāshteʾīm valāken dar javāb-e īn faqare
dhekr namūdīm agar fī'l-ḥaqīqe dar ānche dhekr namūde ṣādeq ast bāyad
Mālek-e vojūd va Solṭān-e ghayb va shohūd rā shokr namāyad ke nafsī rā
az Īrān mabʿūth farmūde ke dar sejn min ghayr-e nāṣer va moʿīn Īrān rā
taṣarrof farmūde va yek māliyāt har sane az ān akhdh mīfarmāyad īn
maqām-e setāyesh ast na nekūhesh law yakūnu mina al-munṣifīn va agar
nafsī erāde namāyad bar omūr-e īn maẓlūm vāqef shavad javāb īnast ke dar
baʿḍī az layālī va ayyām osarā-ye ʿālam va maẓlūmhā-ye omam bī-qūt shab
rā be rūz āvardeand va rūz rā be shab īn adhkār rā dūst nadāshte va
nadārīm va az gūyande ham shekāyat nanamūde va namīnamāyīm dar īn sejn
shakhṣ-e moʿazzazī chandī az sang-tarāshī kasb-e maʿāsh mīnamūd va dar
baʿḍī az awqāt ham baʿḍī be ṭaʿām-e Allāh marzūq yaʿnī jūʿ. Nasʾalu
Allāha tabāraka wa-taʿālā an yuʾayyida al-kulla ʿalā al-ʿadli
wa-l-inṣāfi wa-yuwaffiqahum ʿalā al-inābati wa-l-rujūʿi innahu huwa
al-Sāmiʿu al-Mujīb.

In short, this Wronged One has remained patient and silent in the face
of everything they have inflicted and said, for our intention is,
through the grace of God, exalted be His glory, and His all-preceding
mercy, to efface from the world the rule of contention, conflict, and
bloodshed through the power of utterance. In all circumstances, faced
with what they have said, we have held fast to gracious patience and
left the matter to God. In reply to this statement, however, we said: If
he is indeed truthful in what he has stated, he should thank the Lord of
existence and Sovereign of the unseen and the seen for having raised up
a person from Iran who, while in prison, without helper or supporter,
has taken possession of Iran and collects a tax from it every year. This
calls for praise, not blame, if he is fair-minded. And if anyone wishes
to learn about the affairs of this Wronged One, the answer is that on
some nights and days the captives of the world and the wronged ones
among the nations have passed from night to morning and from day to
night without food. We have never liked, nor do we now like, mentioning
such things, and we have neither complained nor do we complain of the
person who made that statement. In this prison, an honored person earned
his livelihood for a time by cutting stone, and at times some have been
sustained by the food of God---that is, hunger. We ask God, blessed and
exalted be He, to aid everyone in justice and fairness and enable them
to turn to Him and return. He is indeed the All-Hearing, the Responsive.

\subsection{62}\label{section-1348}

\begin{otherlanguage}{arabic}

سبحانّک اللّهمّ یا الهی تری ما ورد علی المظلوم من الّذین ما عاشرونی و قاموا
علی ضرّی و مسکنتی بحیث عجزت الأقلام عن ذکره و الألسن عن بیانه و الألواح
عن حمله تسمع حنین قلبی و ضجیج کینونتی و ما ورد علی امنائک فی بلادک و
اصفیائک فی مملکتک من الّذین نقضوا عهدک و میثاقک ای ربّ اسألک بضجیج العشّاق
فی الآفاق و بحنینهم فی بعدهم عن ساحة قربک و بالدّمآء الّتی سفکت فی حبّک و
بالأکباد الّتی ذابت فی سبیلک ان تحفظ اولیائک من ظلم الّذین غفلوا من اسرار
اسمک المختار ای ربّ ایّدهم بقدرتک الّتی غلبت الأشیآء و وفّقهم علی الصّبر و
الاصطبار انّک انت المقتدر العزیز الوهّاب لا اله الّا انت الکریم الفیّاض

\end{otherlanguage}

Subḥānaka Allāhumma yā ilāhī tarā mā warada ʿalā al-Maẓlūmi mina
alladhīna mā ʿāsharūnī wa-qāmū ʿalā ḍurrī wa-maskanatī bi-ḥaythu
ʿajazati al-aqlāmu ʿan dhikrihi wa-l-alsunu ʿan bayānihi wa-l-alwāḥu ʿan
ḥamlihi tasmaʿu ḥanīna qalbī wa-ḍajīja kaynūnatī wa-mā warada ʿalā
umanāʾika fī bilādika wa-aṣfiyāʾika fī mamlakatika mina alladhīna naqaḍū
ʿahdaka wa-mīthāqaka ay rabbi asʾaluka bi-ḍajīji al-ʿushshāqi fī
al-āfāqi wa-bi-ḥanīnihim fī buʿdihim ʿan sāḥati qurbika wa-bi-l-dimāʾi
allatī sufikat fī ḥubbika wa-bi-l-akbādi allatī dhābat fī sabīlika an
taḥfaẓa awliyāʾaka min ẓulmi alladhīna ghafalū min asrāri ismika
al-Mukhtār ay rabbi ayyidhum bi-qudratika allatī ghalabati al-ashyāʾa
wa-waffiqhum ʿalā al-ṣabri wa-l-iṣṭibāri innaka anta al-Muqtadiru
al-ʿAzīzu al-Wahhāb lā ilāha illā anta al-Karīmu al-Fayyāḍ.

Glory be to You, O God, my God! You see what has befallen the Wronged
One at the hands of those who have never associated with me and have
risen to harm and abase me, to such a degree that pens are powerless to
record it, tongues to describe it, and tablets to bear it. You hear the
lament of my heart and the cry of my being, and what has befallen Your
trusted ones in Your lands and Your chosen ones in Your realm at the
hands of those who have broken Your covenant and pledge. My Lord, I ask
You by the cries of Your lovers throughout the world, by their
lamentation in their remoteness from the court of Your presence, by the
blood shed for love of You, and by the hearts consumed in Your path, to
protect Your loved ones from the oppression of those who are heedless of
the mysteries of Your name, the All-Choosing. My Lord, strengthen them
through Your power, which has prevailed over all things, and enable them
to be patient and steadfast. You are indeed the All-Powerful, the
Almighty, the Ever-Bestowing. There is no God but You, the Generous, the
Bountiful.

\subsection{63}\label{section-1349}

\begin{otherlanguage}{arabic}

این ایّام بعضی از عدل و انصاف گذشته‌اند و با سیف ضغینه و سنان بغضا هجوم
نموده‌اند و حال آنکه شأن منصفین آنکه مردود عالم و مطرود امم را اعانت
نمایند و ببر و تقوی تمسّک جویند اکثری الی حین ندانسته‌اند که ارادهٔ این
مظلوم چیست و بچه جهت حمل بلایای لاتحصی نموده باری لسان قلب به یا لیت
قومی یعلمون ناطق و مظلوم منقطعاً عن الکلّ باین کلمهٔ علیا متکلّم قد احاطت
الأمواج سفینة اللّه المهیمن القیّوم انّک یا ملّاح لا تضطرب عن الأریاح انّ
فالق الأصباح معک فى هذه الظّلمة الّتی بها اضطربت افئدة العباد الّا من شآء
اللّه العزیز المختار

\end{otherlanguage}

Īn ayyām baʿḍī az ʿadl va enṣāf godhashteand va bā seyf-e ḍaghīne va
senān-e baghḍā hojūm namūdeand va ḥāl ānke shaʾn-e monṣefīn ānke
mardūd-e ʿālam va maṭrūd-e omam rā eʿānat namāyand va be berr va taqvā
tamassok jūyand aktharī elā ḥīn nadānesteand ke erāde-ye īn maẓlūm chīst
va be che jehat ḥaml-e balāyā-ye lā-toḥṣā namūde bārī lesān-e qalb be yā
layta qawmī yaʿlamūn nāṭeq va maẓlūm munqaṭiʿan ʿani al-kull be īn
kalame-ye ʿolyā motakallem qad aḥāṭati al-amwāju safīnata Allāhi
al-Muhaymini al-Qayyūm innaka yā mallāḥu lā taḍṭarib ʿani al-aryāḥi inna
Fāliqa al-Aṣbāḥi maʿaka fī hādhihi al-ẓulmati allatī bihā iḍṭarabat
afʾidatu al-ʿibādi illā man shāʾa Allāhu al-ʿAzīzu al-Mukhtār.

These days some have abandoned justice and fairness and have attacked
with the sword of malice and the spear of hatred, although it befits the
fair-minded to assist one rejected by the world and cast out by its
nations, and to hold fast to righteousness and piety. Most have not yet
understood what this Wronged One intends or why he has borne countless
afflictions. In short, the tongue of the heart speaks the words, ``If
only my people knew!'' And the Wronged One, detached from all, utters
these exalted words: ``The waves have surrounded the ship of God, the
All-Protecting, the Self-Subsisting. Mariner, do not be alarmed by the
winds! The One who causes the dawn to break is with you in this
darkness, which has troubled the hearts of all people except those whom
God, the Almighty, the All-Choosing, has willed to spare.''

\subsection{64}\label{section-1350}

\begin{otherlanguage}{arabic}

یا شیخ قسم بآفتاب حقیقت که از افق سماء سجن مشرق و لائح این مظلوم جز
اصلاح قصدی نداشته یشهد بذلک کلّ عارف بصیر و کلّ عالم خبیر و فی حین البلایا
کان متمسّکاً بحبل الصّبر و الاصطبار و راضیاً بما ورد علیه من الأعدآء و هو
یقول یا الهی ترکت ارادتی لارادتک و مشیّتی لظهور مشیّتک وعزّتک لا ارید نفسی
و بقائها الّا للقیام علی خدمة امرک و لا احبّ وجودی الّا للفدآء فی سبیلک ای
ربّ تری و تعلم انّ الّذین اردنا منهم العدل و الانصاف قاموا علینا بالظّلم و
الاعتساف فی الظّاهر کانوا معی و فی الباطن نصروا اعدائی الّذین قاموا علی
هتک حرمتی الهی الهی اشهد انّک خلقت عبادک لنصرة امرک و اعلآء کلمتک و هم
نصروا اعدائک اسألک بأمرک الّذی احاط بالوجود و باسمک الّذی به سخّرت الغیب و
الشّهود ان تزیّن اهل الأرض بنور عدلک و نوّر قلوبهم بضیآء معرفتک ای ربّ انا
عبدک و ابن عبدک اشهد بوحدانیّتک و فردانیّتک و بتقدیس ذاتک و تنزیه کینونتک
ای ربّ تری امنائک بین ایادی الخائنین من خلقک و المفترین من بریّتک و تعلم
ما ورد علینا من الّذین انت اعلم بهم منّا قد ارتکبوا ما انشقّ به ستر
المقرّبین من خلقک اسألک ان تؤیّدهم علی القیام علی ما فات عنهم فی ایّام مشرق
وحیک و مطلع الهامک انّک انت المقتدر على ما تشآء و فی قبضتک زمام من فی
السّموات و الأرضین قد ارتفع ندآء الفطرة و حنینها و تنادی بأعلی النّدآء و
تقول یا قوم تاللّه انّی فزت بمظهری و منزلی هذا یوم فیه ابتسم الطّور بمکلّمه
و کرمل بمنزلها و السّدرة بمعلّمها اتّقوا اللّه و لا تکونوا من المنکرین لا
تحرموا انفسکم عمّا ظهر بالفضل خذوا کوثر البقآء باسم ربّکم مالک الأسمآء ثمّ
اشربوا منه بذکره العزیز البدیع

\end{otherlanguage}

Yā Sheykh qasam be āftāb-e ḥaqīqat ke az ofoq-e samā-ye sejn moshreq va
lāʾeḥ īn maẓlūm joz eṣlāḥ qaṣdī nadāshte yashhadu bi-dhālika kullu
ʿārifin baṣīrin wa-kullu ʿālimin khabīrin wa-fī ḥīni al-balāyā kāna
mutamassikan bi-ḥabli al-ṣabri wa-l-iṣṭibāri wa-rāḍiyan bi-mā warada
ʿalayhi mina al-aʿdāʾi wa-huwa yaqūlu yā ilāhī taraktu irādatī
li-irādatika wa-mashīʾatī li-ẓuhūri mashīʾatika wa-ʿizzatika lā urīdu
nafsī wa-baqāʾahā illā li-l-qiyāmi ʿalā khidmati amrika wa-lā uḥibbu
wujūdī illā li-l-fidāʾi fī sabīlika ay rabbi tarā wa-taʿlamu anna
alladhīna aradnā minhumu al-ʿadla wa-l-inṣāfa qāmū ʿalaynā bi-l-ẓulmi
wa-l-iʿtisāfi fī al-ẓāhiri kānū maʿī wa-fī al-bāṭini naṣarū aʿdāʾī
alladhīna qāmū ʿalā hatki ḥurmatī ilāhī ilāhī ashhadu annaka khalaqta
ʿibādaka li-nuṣrati amrika wa-iʿlāʾi kalimatika wa-hum naṣarū aʿdāʾaka
asʾaluka bi-amrika alladhī aḥāṭa bi-l-wujūdi wa-bi-ismika alladhī bihi
sakhkharta al-ghayba wa-l-shuhūda an tuzayyina ahla al-arḍi bi-nūri
ʿadlika wa-nawwir qulūbahum bi-ḍiyāʾi maʿrifatika ay rabbi anā ʿabduka
wa-ibnu ʿabdika ashhadu bi-waḥdāniyyatika wa-fardāniyyatika
wa-bi-taqdīsi dhātika wa-tanzīhi kaynūnatika ay rabbi tarā umanāʾaka
bayna ayādī al-khāʾinīna min khalqika wa-l-muftarīna min bariyyatika
wa-taʿlamu mā warada ʿalaynā mina alladhīna anta aʿlamu bihim minnā qad
irtakabū mā inshaqqa bihi sitru al-muqarrabīna min khalqika asʾaluka an
tuʾayyidahum ʿalā al-qiyāmi ʿalā mā fāta ʿanhum fī ayyāmi mashriqi
waḥyika wa-maṭlaʿi ilhāmika innaka anta al-Muqtadiru ʿalā mā tashāʾu
wa-fī qabḍatika zimāmu man fī al-samāwāti wa-l-araḍīn qad irtafaʿa
nidāʾu al-fiṭrati wa-ḥanīnuhā wa-tunādī bi-aʿlā al-nidāʾi wa-taqūlu yā
qawmi ta-Allāhi innī fuztu bi-maẓharī wa-munzilī hādhā yawmun fīhi
ibtasama al-Ṭūru bi-mukallimihi wa-Karmilu bi-munzilihā wa-l-Sidratu
bi-muʿallimihā ittaqū Allāha wa-lā takūnū mina al-munkirīn lā taḥrimū
anfusakum ʿammā ẓahara bi-l-faḍli khudhū Kawthara al-Baqāʾi bi-ismi
rabbikum Māliki al-Asmāʾi thumma ishrabū minhu bi-dhikrihi al-ʿAzīzi
al-Badīʿ.

Shaykh! By the Sun of Truth that rises and shines from the horizon of
the heaven of this prison, this Wronged One has had no purpose but
reform. Every discerning knower and every well-informed scholar bears
witness to this. In times of affliction he held fast to the cord of
patience and steadfastness and was content with what befell him at the
hands of his enemies. He says: ``My God, I have relinquished my will for
Your will, and my desire for the manifestation of Your desire. By Your
might, I desire my life and its continuance only to arise in service to
Your Cause, and I cherish my existence only to sacrifice it in Your
path. My Lord, You see and know that those from whom we sought justice
and fairness have risen against us with oppression and tyranny.
Outwardly they were with me, but inwardly they aided my enemies, who
arose to violate my honor. My God, my God! I bear witness that You
created Your servants to aid Your Cause and exalt Your Word, yet they
have aided Your enemies. I ask You, by Your Cause, which has encompassed
existence, and by Your name, through which You have subdued the unseen
and the seen, to adorn the people of the earth with the light of Your
justice, and illumine their hearts with the radiance of Your knowledge.
My Lord, I am Your servant and the son of Your servant. I bear witness
to Your oneness and uniqueness, to the holiness of Your essence and the
transcendence of Your being. My Lord, You see Your trusted ones in the
hands of the treacherous among Your creatures and the slanderers among
Your people. You know what has befallen us at the hands of those whom
You know better than we do. They have committed what has torn the veil
of those among Your creatures who are near to You. I ask You to aid them
to arise and make good what they have missed in the days of the Dawning
Place of Your revelation and the Dayspring of Your inspiration. You
indeed have power over whatever You will, and in Your grasp are the
reins of all who are in the heavens and the earth. The call of innate
nature and its lament have risen, and it cries at the top of its voice,
saying: `People! By God, I have attained the presence of the One who
manifests me and sends me down. This is the day when the Mount has
smiled at the One who converses with it, Carmel at the One who descends
upon it, and the Lote Tree at its Teacher. Fear God and do not be among
those who deny. Do not deprive yourselves of what has appeared through
grace. Take the Kawthar of everlasting life in the name of your Lord,
the Lord of Names, and drink from it in remembrance of Him, the
Almighty, the Incomparable.'\,''

\subsection{65}\label{section-1351}

\begin{otherlanguage}{arabic}

در جمیع احوال عباد را بمعروف امر نمودیم و از منکر نهی مالک وجود شاهد و
گواه که این مظلوم از حقّ برای خلق طلبیده آنچه را که سبب الفت و اتّحاد و
مودّت و اتّفاقست لعمر اللّه این مظلوم قادر بر ستر نبوده انّه اظهر ما اراد و
هو المقتدر المختار

\end{otherlanguage}

Dar jamīʿ-e aḥvāl ʿebād rā be maʿrūf amr namūdīm va az monkar nahy
Mālek-e vojūd shāhed va govāh ke īn maẓlūm az Ḥaqq barā-ye khalq
ṭalabīde ānche rā ke sabab-e olfat va etteḥād va mavaddat va ettefāq ast
la-ʿamru Allāh īn maẓlūm qāder bar satr nabūde innahu aẓhara mā arāda
wa-huwa al-Muqtadiru al-Mukhtār.

In all circumstances we have enjoined people to do good and forbidden
them to do evil. The Lord of existence bears witness that this Wronged
One has asked God, on behalf of humanity, for what fosters fellowship,
unity, affection, and concord. By the life of God, this Wronged One has
been unable to conceal it. He has manifested what He willed, and He is
the All-Powerful, the All-Choosing.

\subsection{66}\label{section-1352}

\begin{otherlanguage}{persian}

این کلمات عالیات که در لوح حضرت سلطان نازل بعضی از آن را مکرّر ذکر
مینمائیم تا آن جناب بیقین مبین بدانند که آنچه ذکر شده من عند اللّه بوده
یا سلطان انّی کنت کأحد من العباد و راقداً على المهاد مرّت علیّ نسائم السّبحان
و علّمنی علم ما کان لیس هذا من عندی بل من لدن عزیز علیم و امرنی بالنّدآء
بین الأرض و السّمآء بذلک ورد علیّ ما ذرفت به دموع العارفین ما قرأت ما عند
النّاس من العلوم و ما دخلت المدارس فاسأل المدینة الّتی کنت فیها لتوقن بأنّی
لست من الکاذبین هذه ورقة حرّکتها اریاح مشیّة ربّک العزیز الحمید هل لها
استقرار عند هبوب اریاح عاصفات لا ومالک الأسمآء و الصّفات بل تحرّکها کیف
ترید لیس للعدم وجود تلقآء القدم قد جآء امره المبرم و انطقنی بذکره بین
العالمین انّی لم اکن الّا کالمیّت تلقآء امره قلّبتنی ید ارادة ربّک الرّحمن
الرّحیم هل یقدر احد ان یتکلّم من تلقآء نفسه بما یعترض به علیه العباد من کلّ
وضیع و شریف لا والّذی علّم القلم اسرار القدم الّا من کان مؤیّداً من لدن مقتدر
قدیر

\end{otherlanguage}

Īn kalemāt-e ʿāliyāt ke dar lawḥ-e Ḥażrat-e Solṭān nāzel baʿḍī az ān rā
mokarrar dhekr mīnamāyīm tā ān janāb be yaqīn-e mobīn bedānand ke ānche
dhekr shode min ʿindi Allāh būde yā sulṭānu innī kuntu ka-aḥadin mina
al-ʿibādi wa-rāqidan ʿalā al-mihādi marrat ʿalayya nasāʾimu al-Subḥāni
wa-ʿallamanī ʿilma mā kāna laysa hādhā min ʿindī bal min ladun ʿAzīzin
ʿAlīmin wa-amaranī bi-l-nidāʾi bayna al-arḍi wa-l-samāʾi bi-dhālika
warada ʿalayya mā dharafat bihi dumūʿu al-ʿārifīn mā qaraʾtu mā ʿinda
al-nāsi mina al-ʿulūmi wa-mā dakhaltu al-madārisa fa-sʾali al-madīnata
allatī kuntu fīhā li-tūqina bi-annī lastu mina al-kādhibīn hādhihi
waraqatun ḥarrakathā aryāḥu mashīʾati rabbika al-ʿAzīzi al-Ḥamīd hal
lahā istiqrārun ʿinda hubūbi aryāḥin ʿāṣifātin lā wa-Māliki al-Asmāʾi
wa-l-Ṣifāti bal tuḥarrikuhā kayfa turīdu laysa li-l-ʿadami wujūdun
tilqāʾa al-Qidami qad jāʾa amruhu al-mubramu wa-anṭaqanī bi-dhikrihi
bayna al-ʿālamīn innī lam akun illā ka-l-mayyiti tilqāʾa amrihi
qallabatnī yadu irādati rabbika al-Raḥmāni al-Raḥīm hal yaqdiru aḥadun
an yatakallama min tilqāʾi nafsihi bi-mā yuʿtaraḍu bihi ʿalayhi
al-ʿibādu min kulli waḍīʿin wa-sharīfin lā wa-alladhī ʿallama al-Qalama
asrāra al-Qidami illā man kāna muʾayyadan min ladun Muqtadirin Qadīr.

We repeat some of these exalted words revealed in the Tablet to His
Majesty the Shah, so that you may know with clear certainty that what
has been stated came from God: ``King! I was like any other person,
asleep upon my couch. The breezes of the All-Glorious passed over me and
taught me the knowledge of all that has been. This is not from me, but
from One who is Almighty, All-Knowing. He commanded me to raise my call
between earth and heaven. Because of this, there befell me what has
caused the tears of those who recognize Him to flow. I have not studied
the sciences current among the people, nor have I entered their schools.
Ask the city in which I lived, so that you may be certain that I am not
one of those who lie. This is a leaf stirred by the winds of the will of
your Lord, the Almighty, the Praiseworthy. Can it remain still when
tempestuous winds blow? No, by the Lord of Names and Attributes! Rather,
they move it as they will. Nonexistence has no existence before the
Eternal. His irrevocable command has come and caused me to speak His
praise among the peoples of the world. Before His command I was as one
dead; the hand of the will of your Lord, the All-Merciful, the
Compassionate, turned me about. Can anyone speak of his own accord words
for which all people, both lowly and exalted, assail him? No, by the One
who taught the Pen the mysteries of eternity---unless he is aided by One
who is All-Powerful, All-Capable.''

\subsection{67}\label{section-1353}

\begin{otherlanguage}{arabic}

یا سلطان انظر بطرف العدل الی المظلوم ثمّ احکم بالحقّ فیما ورد علیه انّ اللّه
قد جعلک ظلّه بین العباد و آیة قدرته لمن فی البلاد احکم بیننا و بین الّذین
ظلمونا من دون بیّنة و لا کتاب منیر انّ الّذین حولک یحبّونک لأنفسهم و الغلام
یحبّک لنفسک و ما اراد الّا ان یقرّبک الی مقرّ الفضل و یقلّبک الی یمین العدل و
کان ربّک علی ما اقول شهیداً

\end{otherlanguage}

Yā sulṭān, unẓur bi-ṭarfi al-ʿadl ilā al-maẓlūm, thumma uḥkum bi-l-ḥaqq
fīmā warada ʿalayh. Inna Allāh qad jaʿalaka ẓillahu bayna al-ʿibād
wa-āyata qudratihi li-man fī al-bilād. Uḥkum baynanā wa-bayna alladhīna
ẓalamūnā min dūni bayyinatin wa-lā kitābin munīr. Inna alladhīna ḥawlaka
yuḥibbūnaka li-anfusihim wa-l-ghulām yuḥibbuka li-nafsika, wa-mā arāda
illā an yuqarribaka ilā maqarri al-faḍl wa-yuqallibaka ilā yamīni
al-ʿadl. Wa-kāna rabbuka ʿalā mā aqūlu shahīdan.

``King! Look upon this Wronged One with the eye of justice, then judge
truthfully what has befallen him. God has made you His shadow among His
servants and a sign of His power to those who inhabit the lands. Judge
between us and those who have wronged us without proof or an
illuminating Book. Those around you love you for their own sake, while
this Youth loves you for your own sake. He has wished only to draw you
near to the seat of grace and turn you toward the right hand of justice.
Your Lord is witness to what I say.''

\subsection{68}\label{section-1354}

\begin{otherlanguage}{arabic}

یا سلطان لو تسمع صریر القلم الأعلی و هدیر ورقآء البقآء علی افنان سدرة
المنتهی فی ذکر اللّه موجد الأسمآء و خالق الأرض و السّمآء لیبلّغک الی مقام
لا تری فى الوجود الّا تجلّی حضرة المعبود و تری الملک احقر شیء عندک تضعه
لمن اراد و تتوجّه الی افق کان بأنوار الوجه مضیئاً و لا تحمل ثقل الملک ابداً
الّا لنصرة ربّک العلیّ الأعلی اذاً یصلّی علیک الملأ الأعلی حبّذا هذا المقام
الأسنی لو ترتقی الیه بسلطان کان باسم اللّه معروفا

\end{otherlanguage}

Yā sulṭān, law tasmaʿ ṣarīra al-qalami al-aʿlā wa-hadīra warqāʾi
al-baqāʾ ʿalā afnāni sidrati al-muntahā fī dhikri Allāh mūjidi al-asmāʾ
wa-khāliqi al-arḍ wa-l-samāʾ, la-yuballighuka ilā maqāmin lā tarā fī
al-wujūd illā tajalliya ḥaḍrati al-maʿbūd, wa-tarā al-mulka aḥqara
shayʾin ʿindaka taḍaʿuhu li-man arāda wa-tatawajjahu ilā ufuqin kāna
bi-anwāri al-wajhi muḍīʾan. Wa-lā taḥmilu thiqla al-mulki abadan illā
li-nuṣrati rabbika al-ʿaliyy al-aʿlā. Idhan yuṣallī ʿalayka al-malaʾ
al-aʿlā. Ḥabbadhā hādhā al-maqāmu al-asnā law tartaqī ilayhi bi-sulṭānin
kāna bi-smi Allāh maʿrūfan.

``King! If you were to hear the sound of the Most Exalted Pen and the
cooing of the Dove of Eternity upon the branches of the Lote Tree of the
Utmost Boundary, as they remember God, the Maker of names and Creator of
earth and heaven, it would bring you to a station where you would see
nothing in existence but the revelation of the One who is worshipped.
You would regard sovereignty as the most insignificant thing,
relinquishing it to whoever desired it and turning toward a horizon
radiant with the lights of His face. You would never bear the burden of
sovereignty except to aid your Lord, the Exalted, the Most High. Then
the Concourse on high would bless you. How excellent is this most
sublime station, if you would ascend to it through a sovereignty known
by the name of God!''

\subsection{69}\label{section-1355}

\begin{otherlanguage}{persian}

آن جناب یا غیر گفته سورهٔ توحید را ترجمه نمایند تا نزد کل معلوم و مبرهن
گردد که حقّ لم‌یلد و لم‌یولد است و بابیها بربوبیّت و الوهیّت قائلند

\end{otherlanguage}

Ān janāb yā ghayr guftih sūrih-yi Tawḥīd rā tarjumih namāyand tā nazd-i
kull maʿlūm va mubarhan gardad kih Ḥaqq lam yalid wa-lam yūlad ast va
Bābīhā bi-rubūbiyyat va ulūhiyyat qāʾiland.

You, or someone else, have said that the Surah of Divine Unity should be
translated so that it may become clear and evident to everyone that God
``neither begets nor is begotten,'' whereas the Bábís profess lordship
and divinity.

\subsection{70}\label{section-1356}

\begin{otherlanguage}{persian}

یا شیخ این مقام مقام فنای از نفس و بقآء باللّه است و این کلمه اگر ذکر شود
مدلّ بر نیستى بحت باتّ است این مقام لا املک لنفسى نفعاً و لا ضرّاً و لا حیوةً
و لا نشورا است

\end{otherlanguage}

Yā Shaykh, īn maqām maqām-i fanā-yi az nafs va baqāʾ bi-Allāh ast va īn
kalimih agar dhikr shavad mudill bar nīstī-yi baḥt-i bātt ast. Īn maqām
lā amliku li-nafsī nafʿan wa-lā ḍarran wa-lā ḥayātan wa-lā nushūran ast.

Shaykh! This station is that of dying to self and abiding in God, and if
this word is mentioned, it denotes sheer and utter nothingness. This is
the station of ``I have no power to bring myself benefit or harm, life
or resurrection.''

\subsection{71}\label{section-1357}

\begin{otherlanguage}{persian}

یا شیخ علمای عصر در تجلّیات سدرهٔ بیان لابن عمران در طور عرفان چه میگویند
آن حضرت کلمه را از سدره اصغا نمود و قبول فرمود ولکن اکثری از ادراک این
مقام محروم چه که به ما عندهم مشغولند و از ما عند اللّه غافل سیّد فندرسک در
این مقام خوب گفته این سخن را درنیابد هیچ فهم ظاهری گر ابونصرستی و ور
بوعلی سیناستی آیا در این فقره که خاتم انبیا روح ما سواه فداه فرموده چه
میگویند میفرماید سترون ربّکم کما ترون البدر فى لیلة اربعة عشر و حضرت امیر
علیه السّلام در خطبهٔ طتنجیّه میفرماید فتوقّعوا ظهور مکلّم موسی من الشّجرة علی
الطّور و همچنین حسین بن علی علیه السّلام میفرماید أ یکون لغیرک من الظّهور
ما لیس لک حتّی یکون هو المظهر لک عمیت عین لا تراک

\end{otherlanguage}

Yā Shaykh, ʿulamā-yi ʿaṣr dar tajalliyāt-i sidrih-yi bayān li-ibn ʿImrān
dar Ṭūr-i ʿirfān chih mīgūyand? Ān Ḥaḍrat kalimih rā az sidrih iṣghā
namūd va qabūl farmūd, va-lākin aktharī az idrāk-i īn maqām maḥrūm, chih
kih bih mā ʿindahum mashghūland va az mā ʿinda Allāh ghāfil. Sayyid
Findirisk dar īn maqām khūb guftih: Īn sukhan rā dar nayābad hīch fahm-i
ẓāhirī, gar Abū Naṣrastī va var Bū ʿAlī Sīnāstī. Āyā dar īn fiqrih kih
Khātam-i Anbiyā, rūḥ-i mā sivāh fidāh, farmūdih chih mīgūyand?
Mīfarmāyad: Sa-tarawna rabbakum kamā tarawna al-badra fī laylati
arbaʿata ʿashar. Va Ḥaḍrat-i Amīr, ʿalayhi al-salām, dar khuṭbih-yi
Ṭatanjiyyih mīfarmāyad: Fa-tawaqqaʿū ẓuhūra mukallimi Mūsā min
al-shajarati ʿalā al-Ṭūr. Va hamchunīn Ḥusayn ibn ʿAlī, ʿalayhi
al-salām, mīfarmāyad: A-yakūnu li-ghayrika min al-ẓuhūri mā laysa laka
ḥattā yakūna huwa al-muẓhira laka? ʿAmiyat ʿaynun lā tarāk.

Shaykh! What do the scholars of this age say about the revelations of
the Lote Tree of Utterance to the son of ʿImrán on the Mount of
recognition? He heard the Word from the Tree and accepted it. Yet most
are deprived of understanding this station, for they are occupied with
what they themselves possess and heedless of what is with God. Sayyid
Findirisk has spoken well on this subject: ``No outward understanding
can grasp these words, even if you are Abú Naṣr or Abú ʿAlí Síná.'' What
do they say about this passage uttered by the Seal of the Prophets---may
the soul of all else be sacrificed for him? He says: ``You will see your
Lord as you see the full moon on the fourteenth night.'' And the
Commander of the Faithful, peace be upon him, says in the Ṭatanjiyyih
sermon: ``Await, then, the appearance of the One who spoke to Moses from
the Tree upon the Mount.'' Likewise, Ḥusayn son of ʿAlí, peace be upon
him, says: ``Can anyone other than You possess a manifestation that You
do not possess, so that he should be the one to make You manifest? Blind
is the eye that does not see You.''

\subsection{72}\label{section-1358}

\begin{otherlanguage}{arabic}

امثال این اذکار از بیانات اولیا صلوات اللّه علیهم مذکور و مشهور و در کتب
معتبره موجود طوبی لمن ینظر و یتکلّم بالصّدق الخالص نعیم از برای نفسی که
خود را بمدد کوثر بیان مقصود عالمیان از ظنون و اوهام طاهر نمود و سبحات
جلال را باسم غنیّ متعال شقّ کرد و از عالم و عالمیان گذشت و قصد سجن اعظم
نمود

\end{otherlanguage}

Amthāl-i īn adhkār az bayānāt-i awliyā, ṣalawāt Allāh ʿalayhim, madhkūr
va mashhūr va dar kutub-i muʿtabarih mawjūd. Ṭūbā li-man yanẓuru
wa-yatakallamu bi-l-ṣidqi al-khāliṣ. Naʿīm az barā-yi nafsī kih khud rā
bi-madad-i kawthar-i bayān-i Maqṣūd-i ʿālamiyān az ẓunūn va awhām ṭāhir
namūd va subuḥāt-i jalāl rā bi-ism-i Ghaniyy-i Mutaʿāl shaqq kard va az
ʿālam va ʿālamiyān gudhasht va qaṣd-i Sijn-i Aʿẓam namūd.

Such sayings from the utterances of the holy ones, God's blessings be
upon them, are well known and widely cited, and are found in
authoritative books. Blessed is the one who examines and speaks with
unalloyed truthfulness. Bliss awaits the soul who, through the heavenly
stream of the utterance of the Desire of all peoples, has purified
himself of conjectures and vain imaginings, torn through the veils of
glory in the name of the All-Sufficing, the Most Exalted, passed beyond
the world and its peoples, and set out for the Most Great Prison.

\subsection{73}\label{section-1359}

\begin{otherlanguage}{arabic}

یا شیخ نفحات وحی از دونش ممتاز و بیان الهی مابین کتب بمثابهٔ آفتاب مشرق و
لائح طوبى لمن وجد و عرف و قال لک الحمد یا مقصود العالم و لک الشّکر یا
محبوب افئدة المخلصین

\end{otherlanguage}

Yā Shaykh, nafaḥāt-i vaḥy az dūnash mumtāz va bayān-i ilāhī mābayn-i
kutub bi-mathābih-yi āftāb mushriq va lāʾiḥ. Ṭūbā li-man wajada
wa-ʿarafa wa-qāla: Laka al-ḥamdu yā Maqṣūd al-ʿālam wa-laka al-shukru yā
Maḥbūb afʾidati al-mukhliṣīn.

Shaykh! The fragrances of revelation are distinct from all else, and
divine utterance shines resplendent among books like the sun. Blessed is
the one who has found and recognized it and said: ``Praise be to You,
Desire of the world, and thanks be to You, Beloved of the hearts of the
sincere!''

\subsection{74}\label{section-1360}

\begin{otherlanguage}{persian}

مقصود از ذکر الوهیّت و ربوبیّت را عباد ملتفت نشده‌اند چه اگر بیابند از مقام
خود قیام کنند و بکلمهٔ تبنا الی اللّه ناطق گردند حضرت خاتم روح ما سواه
فداه میفرماید لنا مع اللّه حالات نحن فیها هو و هو نحن و هو هو و نحن نحن

\end{otherlanguage}

Maqṣūd az dhikr-i ulūhiyyat va rubūbiyyat rā ʿibād multafit
nashudih-and, chih agar biyāband az maqām-i khud qiyām kunand va
bi-kalimih-yi tubnā ilā Allāh nāṭiq gardand. Ḥaḍrat-i Khātam, rūḥ-i mā
sivāh fidāh, mīfarmāyad: Lanā maʿa Allāhi ḥālātun naḥnu fīhā huwa
wa-huwa naḥnu wa-huwa huwa wa-naḥnu naḥnu.

People have not understood what is meant by the mention of divinity and
lordship. If they understood, they would rise from their places and
utter the words, ``We have turned in repentance to God.'' The Seal of
the Prophets---may the soul of all else be sacrificed for him---says:
``We have states with God in which we are He and He is we, yet He is He
and we are we.''

\subsection{75}\label{section-1361}

\begin{otherlanguage}{arabic}

از این مقام هم گذشته چرا مقامات دیگر را که از قلم ابهی نازل ذکر
ننموده‌اند لسان مظلوم در اکثری از ایّام و لیالی باین کلمات عالیات ناطق
الهی الهی اشهد بوحدانیّتک و فردانیّتک و بأنّک انت اللّه لا اله الّا انت لم
تزل کنت مقدّساً عن ذکر دونک و ثنآء غیرک و لا تزال تکون بمثل ما قد کنت من
قبل و من بعد اسألک یا مالک القدم بالاسم الأعظم و بتجلّیات نیّر ظهورک فی
طور البیان و بأمواج بحر علمک فی الامکان ان تؤیّدنی علی ما یقرّبنی الیک و
یقطعنی عن دونک وعزّتک یا اله الکائنات و مقصود الممکنات احبّ ان اضع وجهی
على کلّ بقعة من بقاع ارضک لعلّه یتشرّف بمقام تشرّف بقدوم اولیائک

\end{otherlanguage}

Az īn maqām ham gudhashtih, chirā maqāmāt-i dīgar rā kih az Qalam-i Abhā
nāzil dhikr nanamūdih-and? Lisān-i Maẓlūm dar aktharī az ayyām va layālī
bi-īn kalimāt-i ʿāliyāt nāṭiq: Ilāhī ilāhī, ashhadu bi-waḥdāniyyatika
wa-fardāniyyatika wa-bi-annaka anta Allāhu lā ilāha illā anta. Lam tazal
kunta muqaddasan ʿan dhikri dūnika wa-thanāʾi ghayrika wa-lā tazālu
takūnu bi-mithli mā qad kunta min qablu wa-min baʿd. Asʾaluka yā Mālik
al-qidam bi-l-ismi al-aʿẓam wa-bi-tajalliyāti nayyiri ẓuhūrika fī Ṭūri
al-bayān wa-bi-amwāji baḥri ʿilmika fī al-imkān an tuʾayyidanī ʿalā mā
yuqarribunī ilayka wa-yaqṭaʿunī ʿan dūnika. Wa-ʿizzatika yā Ilāh
al-kāʾināt wa-Maqṣūd al-mumkināt, uḥibbu an aḍaʿa wajhī ʿalā kulli
buqʿatin min biqāʿi arḍika laʿallahu yatasharrafu bi-maqāmin tasharrafa
bi-qudūmi awliyāʾika.

Even setting this station aside, why have they not mentioned the other
stations revealed by the All-Glorious Pen? On most days and nights the
tongue of this Wronged One utters these exalted words: ``My God, my God!
I bear witness to Your oneness and singleness, and that You are God;
there is no God but You. You have always been sanctified beyond the
mention of anyone besides You and the praise of anyone other than You,
and You will forever remain as You have been, before and hereafter. I
ask You, Lord of eternity, by the Most Great Name, by the radiance of
the Sun of Your Revelation upon the Mount of Utterance, and by the waves
of the ocean of Your knowledge throughout contingent existence, to
assist me in what will draw me near to You and detach me from all
besides You. By Your might, God of all beings and Desire of all created
things! I long to lay my face upon every spot of Your earth, that it may
perhaps be honored by a place honored by the footsteps of Your holy
ones.''

\subsection{76}\label{section-1362}

\begin{otherlanguage}{persian}

لعمر اللّه اوهام عباد را از افق ایقان محروم نموده و ظنون از رحیق مختوم
منع کرده براستی میگویم و لوجه اللّه میگویم این عبد و این مظلوم شرم دارد
خود را به هستی و وجود نسبت دهد تا چه رسد بمقامات فوق آن انسان بصیر هر
هنگام بر ارض مشی مینماید خود را خجل مشاهده میکند چه که بیقین مبین میداند
که علّت نعمت و ثروت و عزّت و علوّ و سموّ و اقتدار او باذن اللّه ارضی است که
تحت جمیع اقدام عالم است و نفسی که باین مقام آگاه شد البتّه از نخوت و کبر
و غرور مقدّس و مبرّا است آنچه گفته شد من عند اللّه بوده انّه شهد و یشهد و هو
العلیم الخبیر

\end{otherlanguage}

La-ʿamru Allāh, awhām ʿibād rā az ufuq-i īqān maḥrūm namūdih va ẓunūn az
raḥīq-i makhtūm manʿ kardih. Bi-rāstī mīgūyam va li-wajhi Allāh mīgūyam,
īn ʿabd va īn Maẓlūm sharm dārad khud rā bih hastī va vujūd nisbat dihad
tā chih rasad bi-maqāmāt-i fawq-i ān. Insān-i baṣīr har hangām bar arḍ
mashy mīnamāyad khud rā khajil mushāhidih mīkunad, chih kih bi-yaqīn-i
mubīn mīdānad kih ʿillat-i niʿmat va tharvat va ʿizzat va ʿuluvv va
sumuvv va iqtidār-i ū bi-idhn Allāh arḍī ast kih taḥt-i jamīʿ-i aqdām-i
ʿālam ast. Va nafsī kih bi-īn maqām āgāh shud al-battih az nikhvat va
kibr va ghurūr muqaddas va mubarrā ast. Ānchih guftih shud min ʿindi
Allāh būdih. Innahu shahida wa-yashhadu wa-huwa al-ʿAlīm al-Khabīr.

By the life of God! Vain imaginings have deprived people of the horizon
of certitude, and conjectures have kept them from the sealed wine. I
speak truthfully, and I speak for the sake of God: this servant, this
Wronged One, is ashamed to attribute being and existence to himself, let
alone stations beyond them. Whenever an insightful person walks upon the
earth, he finds himself abashed, for he knows with clear certainty that,
by God's leave, the source of his well-being, wealth, honor, elevation,
distinction, and power is the earth beneath the feet of all the world. A
soul who becomes aware of this station is assuredly sanctified and freed
from arrogance, pride, and conceit. What has been said has come from
God. He has borne and continues to bear witness, and He is the
All-Knowing, the All-Informed.

\subsection{77}\label{section-1363}

\begin{otherlanguage}{persian}

از حقّ آذان واعیه و ابصار حدیده و صدور منشرحه و قلوب مقبله مسئلت نمائید
شاید عباد مقصود را بیابند و بشطر دوست توجّه نمایند بلایائی بر این مظلوم
وارد شد که شبه آن را عیون مشاهده نکرده و در اظهار امر بهیچوجه توقّف ننمود
و بحضرات ملوک و سلاطین ایّدهم اللّه تعالى القا نمود آنچه را که سبب آسایش و
اتّحاد و اتّفاق و عمار عالم و راحت امم است از جمله دو کلمه از ناپلیون ثالث
اصغا شد و آن سبب شد که در ادرنه لوحی باو ارسال نمودیم جواب نفرمودند و
بعد از ورود در سجن اعظم از وزیر ایشان نامه و دستخطّی رسید اوّل آن بخطّ عجمی
و آخر آن را بخطّ خودشان مرقوم داشته بودند و اظهار عنایت فرموده بودند و
مذکور داشتند که نامه را حسب الخواهش رساندم و الی حین جوابى نفرموده‌اند
ولکن بوزیر مختار خودمان در علیّه و قونسلوس آن اراضی سفارش نمودیم هر مطلبی
باشد اظهار فرمائید اجرا میگردد

\end{otherlanguage}

Az Ḥaqq ādhān-i vāʿiyih va abṣār-i ḥadīdih va ṣudūr-i munshariḥih va
qulūb-i muqbilih masʾalat namāyīd, shāyad ʿibād maqṣūd rā biyāband va
bi-shaṭr-i Dūst tavajjuh namāyand. Balāyāʾī bar īn Maẓlūm vārid shud kih
shibh-i ān rā ʿuyūn mushāhidih nakardih va dar iẓhār-i Amr bi-hīch vajh
tavaqquf nanamūd va bi-ḥaḍarāt-i mulūk va salāṭīn, ayyadahum Allāhu
taʿālā, ilqā namūd ānchih rā kih sabab-i āsāyish va ittiḥād va ittifāq
va ʿimār-i ʿālam va rāḥat-i umam ast. Az jumlih du kalimih az Nāpliyūn-i
Thālith iṣghā shud va ān sabab shud kih dar Adirnih lawḥī bi-ū irsāl
namūdīm. Javāb nafarmūdand va baʿd az vurūd dar Sijn-i Aʿẓam az vazīr-i
īshān nāmih va dastkhaṭṭī rasīd. Avval-i ān bi-khaṭṭ-i ʿajamī va ākhir-i
ān rā bi-khaṭṭ-i khudishān marqūm dāshtih būdand va iẓhār-i ʿināyat
farmūdih būdand va madhkūr dāshtand kih nāmih rā ḥasb al-khvāhish
rasāndam va ilā ḥīn javābī nafarmūdih-and, va-lākin bi-vazīr-i mukhtār-i
khudimān dar ʿAliyyih va qunsulus-i ān arāḍī sifārish namūdīm har
maṭlabī bāshad iẓhār farmāʾīd ijrā mīgardad.

Ask God for attentive ears, keen eyes, receptive breasts, and hearts
turned toward Him, so that people may perhaps grasp the purpose and turn
toward the Friend. Afflictions unlike anything eyes have witnessed
befell this Wronged One, yet he never hesitated to proclaim the Cause.
He conveyed to their majesties the kings and sovereigns---may God,
exalted be He, assist them---what would bring about the peace, unity,
concord, and prosperity of the world and the well-being of its peoples.
Among other things, two statements by Napoleon III were heard, and these
prompted us to send him a Tablet from Adrianople. He did not reply.
After our arrival in the Most Great Prison, a letter with a handwritten
note arrived from his minister. Its first part was written in Persian
script and its final part in his own hand. He expressed his kindness and
stated: ``As requested, I delivered the letter, and thus far he has
given no reply. We have, however, instructed our minister
plenipotentiary in the capital and our consul in those regions that
whatever request you may have, you should make it known, and it will be
carried out.''

\subsection{78}\label{section-1364}

\begin{otherlanguage}{persian}

از این بیان ایشان معلوم شد مقصود این عبد را اصلاح امور ظاهره دانسته‌اند
لذا آیاتی باسم ایشان در سورهٔ هیکل نازل بعضی از آن ذکر میشود تا آن جناب
بدانند که امر این مظلوم للّه و من عند اللّه بوده

\end{otherlanguage}

Az īn bayān-i īshān maʿlūm shud maqṣūd-i īn ʿabd rā iṣlāḥ-i umūr-i
ẓāhirih dānistih-and. Lidhā āyātī bi-ism-i īshān dar Sūrih-yi Haykal
nāzil. Baʿḍī az ān dhikr mīshavad tā ān janāb bidānand kih Amr-i īn
Maẓlūm li-Allāh va min ʿindi Allāh būdih.

From this statement of his, it became clear that he had understood this
servant's purpose to be the improvement of outward circumstances.
Therefore, verses addressed to him were revealed in the Surah of the
Temple. Some of them are quoted here so that you may know that the Cause
of this Wronged One has been for God and from God.

\subsection{79}\label{section-1365}

\begin{otherlanguage}{arabic}

یا ملک باریس نبّئ القسّیس ان لا یدقّ النّواقیس تاللّه الحقّ قد ظهر النّاقوس
الأفخم علی هیکل الاسم الأعظم و تدقّه اصابع مشیّة ربّک العلیّ الأعلی فی جبروت
البقآء باسمه الأبهی کذلک نزّلت آیات ربّک الکبری تارةً اخری لتقوم علی ذکر
اللّه فاطر الأرض و السّمآء فی هذه الأیّام الّتی فیها ناحت قبائل الأرض کلّها و
تزلزلت ارکان البلاد و غشّت العباد غبرة الالحاد الّا من شآء ربّک العلیم
الحکیم قل قد اتی المختار فی ظلل الأنوار لیحیی الأکوان من نفحات اسمه
الرّحمن و یتّحد العالم و یجتمعوا على هذه المائدة الّتی نزلت من السّمآء ایّاکم
ان تکفروا بنعمة اللّه بعد انزالها هذا خیر لکم ممّا عندکم لأنّه سیفنی و ما
عند اللّه یبقی انّه هو الحاکم علی ما یرید قد هبّت نسمات الغفران من شطر ربّکم
الرّحمن من اقبل الیها طهّرته عن العصیان و عن کلّ دآء و سقم طوبی لمن اقبل
الیها و ویل للمعرضین

\end{otherlanguage}

Yā malik Bārīs, nabbiʾ al-qissīs an lā yaduqqa al-nawāqīs. Ta-Allāhi
al-Ḥaqq, qad ẓahara al-nāqūsu al-afkham ʿalā haykali al-ismi al-aʿẓam
wa-taduqquhu aṣābiʿu mashīʾati rabbika al-ʿaliyy al-aʿlā fī jabarūti
al-baqāʾ bi-smihi al-Abhā. Kadhālika nuzzilat āyātu rabbika al-kubrā
tāratan ukhrā li-taqūma ʿalā dhikri Allāh fāṭiri al-arḍ wa-l-samāʾ fī
hādhihi al-ayyāmi allatī fīhā nāḥat qabāʾilu al-arḍi kulluhā
wa-tazalzalat arkānu al-bilād wa-ghashshat al-ʿibāda ghubratu al-ilḥād
illā man shāʾa rabbuka al-ʿAlīm al-Ḥakīm. Qul qad atā al-Mukhtār fī
ẓulali al-anwār li-yuḥyiya al-akwāna min nafaḥāti ismihi al-Raḥmān
wa-yattaḥida al-ʿālam wa-yajtamiʿū ʿalā hādhihi al-māʾidati allatī
nazalat min al-samāʾ. Iyyākum an takfurū bi-niʿmati Allāh baʿda
inzālihā. Hādhā khayrun lakum mimmā ʿindakum li-annahu sa-yafnā wa-mā
ʿinda Allāhi yabqā. Innahu huwa al-ḥākimu ʿalā mā yurīd. Qad habbat
nasamātu al-ghufrān min shaṭri rabbikum al-Raḥmān. Man aqbala ilayhā
ṭahharathu ʿan al-ʿiṣyān wa-ʿan kulli dāʾin wa-saqam. Ṭūbā li-man aqbala
ilayhā wa-waylun li-l-muʿriḍīn.

``King of Paris! Tell the priest to ring the bells no longer. By God,
the True One! The most majestic Bell has appeared in the temple of the
Most Great Name, and the fingers of the will of your Lord, the Exalted,
the Most High, ring it in the dominion of eternity in His All-Glorious
name. Thus have the most great verses of your Lord been sent down once
again, so that you may arise to remember God, the Creator of earth and
heaven, in these days when all the tribes of the earth have lamented,
the foundations of the countries have been shaken, and the dust of
unbelief has enveloped His servants, except those whom your Lord, the
All-Knowing, the All-Wise, has willed to spare. Say: The Chosen One has
come beneath canopies of light to quicken all beings with the fragrances
of His name, the All-Merciful, so that the world may be united and all
may gather around this table sent down from heaven. Beware of denying
God's bounty after it has been sent down. This is better for you than
what you possess, for that will perish, while what is with God will
endure. He, indeed, ordains whatever He wills. The breezes of
forgiveness have blown from the direction of your Lord, the
All-Merciful. Whoever turns toward them will be purified by them from
disobedience and from every disease and ailment. Blessed is the one who
turns toward them, and woe to those who turn away!''

\subsection{80}\label{section-1366}

\begin{otherlanguage}{arabic}

لو تتوجّه بسمع الفطرة الی الأشیآء تسمع منها قد اتی القدیم ذو المجد العظیم
یسبّح کلّ شىء بحمد ربّه منهم من عرف اللّه و یذکر و منهم من یذکر و لا یعرف
کذلک احصینا الأمر فی لوح مبین

\end{otherlanguage}

Law tatawajjahu bi-samʿi al-fiṭrati ilā al-ashyāʾ tasmaʿu minhā: Qad atā
al-Qadīm dhū al-majdi al-ʿaẓīm. Yusabbiḥu kullu shayʾin bi-ḥamdi rabbih.
Minhum man ʿarafa Allāha wa-yadhkur wa-minhum man yadhkuru wa-lā yaʿrif.
Kadhālika aḥṣaynā al-amra fī lawḥin mubīn.

``If you turn toward all things with the ear of your innate nature, you
will hear them say: `The Ancient One, possessed of great glory, has
come.' Everything glorifies its Lord with praise. Among them are those
who have recognized God and remember Him, and those who remember Him
without recognizing Him. Thus have we set forth the matter in a clear
Tablet.''

\subsection{81}\label{section-1367}

\begin{otherlanguage}{arabic}

یا ملک اسمع النّدآء من هذه النّار المشتعلة من هذه الشّجرة الخضرآء فی هذا
الطّور المرتفع علی البقعة المقدّسة البیضآء خلف قلزم البقآء انّه لا اله الّا
انا الغفور الرّحیم قد ارسلنا من ایّدناه بروح القدس لیخبرکم بهذا النّور الّذی
اشرق من افق مشیّة ربّکم العلیّ الأبهی و ظهرت فی الغرب آثاره توجّهوا الیه فی
هذا الیوم الّذى جعله اللّه غرّة الأیّام و فیه تجلّی الرّحمن علی من فی السّموات
و الأرضین قم علی خدمة اللّه و نصرة امره انّه یؤیّدک بجنود الغیب و الشّهادة و
یجعلک سلطاناً علی ما تطلع الشّمس علیه انّ ربّک هو المقتدر القدیر

\end{otherlanguage}

Yā malik, ismaʿ al-nidāʾ min hādhihi al-nāri al-mushtaʿilati min hādhihi
al-shajarati al-khaḍrāʾ fī hādhā al-Ṭūri al-murtafiʿ ʿalā al-buqʿati
al-muqaddasati al-bayḍāʾ khalfa qulzumi al-baqāʾ: Innahu lā ilāha illā
anā al-Ghafūr al-Raḥīm. Qad arsalnā man ayyadnāhu bi-rūḥi al-qudus
li-yukhbirakum bi-hādhā al-nūri alladhī ashraqa min ufuqi mashīʾati
rabbikum al-ʿaliyy al-Abhā wa-ẓaharat fī al-gharbi āthāruh. Tawajj ahū
ilayhi fī hādhā al-yawmi alladhī jaʿalahu Allāhu ghurrata al-ayyām
wa-fīhi tajallā al-Raḥmān ʿalā man fī al-samāwāti wa-l-arḍīn. Qum ʿalā
khidmati Allāh wa-nuṣrati amrih. Innahu yuʾayyiduka bi-junūdi al-ghaybi
wa-l-shahādah wa-yajʿaluka sulṭānan ʿalā mā taṭluʿu al-shamsu ʿalayh.
Inna rabbaka huwa al-Muqtadir al-Qadīr.

``King! Listen to the call from this Fire blazing in this verdant Tree
upon this lofty Mount, on the holy, snow-white Spot beyond the ocean of
eternity: `There is no God but Me, the Ever-Forgiving, the Merciful.' We
sent the one whom We strengthened with the Holy Spirit to tell you of
this Light, which has shone from the horizon of the will of your Lord,
the Exalted, the All-Glorious, and whose signs have appeared in the
West. Turn toward Him on this Day, which God has made the crown of days,
and on which the All-Merciful has revealed Himself to all who are in the
heavens and on earth. Arise to serve God and aid His Cause. He will
support you with the hosts of the unseen and the seen, and make you
sovereign over all upon which the sun rises. Your Lord is indeed the
All-Powerful, the Almighty.''

\subsection{82}\label{section-1368}

\begin{otherlanguage}{arabic}

قد فاحت نفحات الرّحمن فی الامکان طوبی لمن وجد عرفها و اقبل الیها بقلب
سلیم زیّن هیکلک بطراز اسمی و لسانک بذکری و قلبک بحبّی العزیز المنیع ما
اردنا لک الّا ما هو خیر لک ممّا عندک و من خزائن الأرض کلّها انّ ربّک هو
العلیم الخبیر قم بین العباد باسمی و قل یا ملأ الأرض اقبلوا الی من اقبل
الیکم انّه لوجه اللّه بینکم و حجّته فیکم و دلیله لکم قد جآءکم بآیات عجز
عنها العالمون انّ شجرة الطّور تنطق فی صدر العالم و روح القدس ینادی بین
الأمم قد اتی المقصود بسلطان مبین

\end{otherlanguage}

Qad fāḥat nafaḥātu al-Raḥmān fī al-imkān. Ṭūbā li-man wajada ʿarfahā
wa-aqbala ilayhā bi-qalbin salīm. Zayyin haykalaka bi-ṭirāzi ismī
wa-lisānaka bi-dhikrī wa-qalbaka bi-ḥubbī al-ʿazīz al-manīʿ. Mā aradnā
laka illā mā huwa khayrun laka mimmā ʿindaka wa-min khazāʾini al-arḍi
kullihā. Inna rabbaka huwa al-ʿalīm al-khabīr. Qum bayna al-ʿibādi
bi-ismī wa-qul yā malaʾa al-arḍi aqbilū ilā man aqbala ilaykum. Innahu
la-wajhu Allāhi baynakum wa-ḥujjatuhu fīkum wa-dalīluhu lakum. Qad
jāʾakum bi-āyātin ʿajaza ʿanhā al-ʿālamūn. Inna shajarata al-Ṭūr tanṭiqu
fī ṣadri al-ʿālam wa-rūḥu al-qudus yunādī bayna al-umam: qad atā
al-maqṣūdu bi-sulṭānin mubīn.

The breezes of the All-Merciful have wafted through creation. Blessed is
whoever has caught their fragrance and turned toward them with a pure
heart. Adorn your person with the ornament of My name, your tongue with
My remembrance, and your heart with My mighty, unassailable love. We
have desired for you only what is better for you than all you possess
and all the treasures of the earth. Your Lord is indeed the All-Knowing,
the All-Aware. Arise among the people in My name and say: ``People of
the earth! Turn toward Him who has turned toward you. He is indeed the
face of God among you, His proof in your midst, and His guide for you.
He has come to you with signs that the peoples of the world are
powerless to match. The Tree of Sinai speaks in the heart of the world,
and the Holy Spirit calls out among the nations: `The Desired One has
come with manifest authority.'\,''

\subsection{83}\label{section-1369}

\begin{otherlanguage}{arabic}

یا ملک قد سقطت انجم سمآء العلم الّذین یستدلّون بما عندهم لاثبات امری و
یذکرون اللّه باسمی فلمّا جئتهم بمجدی اعرضوا عنّی الا انّهم من السّاقطین هذا
ما اخبرکم به الرّوح اذ اتی بالحقّ و اعترض علیه علمآء الیهود الی ان ارتکبوا
ما ناح به روح القدس و ذرفت دموع المقرّبین

\end{otherlanguage}

Yā malik, qad saqaṭat anjumu samāʾi al-ʿilm alladhīna yastadillūna bi-mā
ʿindahum li-ithbāti amrī wa-yadhkurūna Allāha bi-ismī. Fa-lammā jiʾtuhum
bi-majdī aʿraḍū ʿannī. Alā innahum mina al-sāqiṭīn. Hādhā mā akhbarakum
bihi al-rūḥ idh atā bi-al-ḥaqq wa-iʿtaraḍa ʿalayhi ʿulamāʾu al-yahūd ilā
an irtakabū mā nāḥa bihi rūḥu al-qudus wa-dharafat dumūʿu al-muqarrabīn.

King! The stars of the heaven of knowledge have fallen---those who draw
on what they possess to establish the truth of My Cause and who remember
God in My name. Yet when I came to them in My glory, they turned away
from Me. They are indeed among the fallen. This is what the Spirit
foretold to you when He came with the truth and the Jewish scholars
opposed Him, until they committed what caused the Holy Spirit to lament
and the tears of those brought near to God to flow.

\subsection{84}\label{section-1370}

\begin{otherlanguage}{arabic}

قل یا ملأ الرّهبان لا تعتکفوا فی الکنائس و المعابد اخرجوا باذنی ثمّ
اشتغلوا بما تنتفع به انفسکم و انفس العباد کذلک یأمرکم مالک یوم الدّین
اعتکفوا فی حصن حبّی هذا حقّ الاعتکاف لو کنتم من العارفین من جاور البیت انّه
کالمیّت ینبغی للانسان ان یظهر منه ما ینتفع به العباد و الّذی لیس له ثمر
ینبغی للنّار کذلک یعظکم ربّکم انّه هو العزیز الکریم تزوّجوا لیقوم بعدکم احد
مقامکم انّا منعناکم عن الخیانة لا عمّا تظهر به الأمانة أ اخذتم اصول انفسکم
و نبذتم اصول اللّه ورائکم اتّقوا اللّه و لا تکونوا من الجاهلین لو لا
الانسان من یذکرنی فی ارضی و کیف تظهر صفاتی و اسمائی تفکّروا و لا تکونوا
من الّذین احتجبوا و کانوا من الرّاقدین انّ الّذى ما تزوّج انّه ما وجد مقرّاً
لیسکن فیه او یضع رأسه علیه بما اکتسبت ایدی الخائنین لیس تقدیس نفسه بما
عرفتم و عندکم من الأوهام بل بما عندنا اسألوا لتعرفوا مقامه الّذی کان
مقدّساً عن ظنون من علی الأرض کلّها طوبی للعارفین

\end{otherlanguage}

Qul yā malaʾa al-ruhbān, lā taʿtakifū fī al-kanāʾisi wa-al-maʿābid.
Ukhrujū bi-idhnī thumma ishtag hilū bi-mā tanta fiʿu bihi anfusukum
wa-anfusu al-ʿibād. Kadhālika yaʾmurukum māliku yawmi al-dīn. Iʿtakifū
fī ḥiṣni ḥubbī. Hādhā ḥaqqu al-iʿtikāf law kuntum mina al-ʿārifīn. Man
jāwara al-bayta innahu ka-al-mayyit. Yanbaghī li-al-insāni an yaẓhara
minhu mā yantafiʿu bihi al-ʿibād wa-alladhī laysa lahu thamarun yanbaghī
li-al-nār. Kadhālika yaʿiẓukum rabbukum innahu huwa al-ʿazīz al-karīm.
Tazawwajū li-yaqūma baʿdakum aḥadun maqāmakum. Innā manaʿnākum ʿani
al-khiyānati lā ʿammā taẓharu bihi al-amānah. A-akhadhtum uṣūla
anfusikum wa-nabadhtum uṣūla Allāhi warāʾakum? Ittaqū Allāha wa-lā
takūnū mina al-jāhilīn. Law lā al-insānu man yadhkurunī fī arḍī wa-kayfa
taẓharu ṣifātī wa-asmāʾī? Tafakkarū wa-lā takūnū mina alladhīna iḥtajabū
wa-kānū mina al-rāqidīn. Inna alladhī mā tazawwaja innahu mā wajada
maqarran li-yaskuna fīhi aw yaḍaʿa raʾsahu ʿalayhi bi-mā iktasabat aydī
al-khāʾinīn. Laysa taqdīsu nafsihi bi-mā ʿaraftum wa-ʿindakum mina
al-awhām bal bi-mā ʿindanā. Isʾalū li-taʿrifū maqāmahu alladhī kāna
muqaddasan ʿan ẓunūni man ʿalā al-arḍi kullihā. Ṭūbā li-al-ʿārifīn.

Say: ``Monks! Do not seclude yourselves in churches and temples. Come
out with My permission, then occupy yourselves with what benefits you
and other people. Thus the Master of the Day of Judgment commands you.
Seclude yourselves in the stronghold of My love. This is true seclusion,
if you would understand. Whoever remains secluded in the house is like
the dead. A human being should bring forth what benefits others, and
whatever bears no fruit is fit for the fire. Thus your Lord admonishes
you. He is indeed the Almighty, the All-Bountiful. Marry, so that
someone may take your place after you. We have forbidden you treachery,
not that through which faithfulness is manifested. Have you adopted
principles of your own and cast the principles of God behind you? Fear
God and do not be among the ignorant. Were it not for humanity, who
would remember Me on My earth, and how would My attributes and names be
manifested? Reflect, and do not be among those who are veiled and lie
asleep. He who did not marry found no place to dwell or to lay His head,
because of what the hands of the treacherous had done. The sanctity of
His person does not consist in what you have understood and the vain
imaginings you entertain, but in what is with Us. Ask, so that you may
recognize His station, which is sanctified above the conjectures of all
who dwell on earth. Blessed are those who understand.''

\subsection{85}\label{section-1371}

\begin{otherlanguage}{arabic}

یا ملک انّا سمعنا منک کلمة تکلّمت بها اذ سألک ملک الرّوس عمّا قضی من حکم
الغزآء انّ ربّک هو العلیم الخبیر قلت کنت راقداً فی المهاد ایقظنی ندآء
العباد الّذین ظلموا الی ان غرقوا فی البحر الأسود کذلک سمعنا و ربّک علی ما
اقول شهید نشهد انّک ما ایقظک النّدآء بل الهوی لأنّا بلوناک وجدناک فی معزل
اعرف لحن القول و کن من المتفرّسین انّا ما نحبّ ان نرجع الیک کلمة سوء حفظاً
للمقام الّذى اعطیناک فی الحیوة الظّاهرة انّا اخترنا الأدب و جعلناه سجیّة
المقرّبین انّه ثوب یوافق النّفوس من کلّ صغیر و کبیر طوبی لمن جعله طراز هیکله
ویل لمن جعل محروماً من هذا الفضل العظیم لو کنت صاحب الکلمة ما نبذت کتاب
اللّه ورآء ظهرک اذ ارسل الیک من لدن عزیز حکیم انّا بلوناک به ما وجدناک على
ما ادّعیت قم و تدارک ما فات عنک سوف تفنی الدّنیا و ما عندک و یبقی الملک
للّه ربّک و ربّ آبائک الأوّلین لا ینبغی لک ان تقتصر الأمور علی ما یهوی هواک
اتّق زفرات المظلوم احفظه من سهام الظّالمین

\end{otherlanguage}

Yā malik, innā samiʿnā minka kalimatan takallamta bihā idh saʾalaka
maliku al-Rūs ʿammā quḍiya min ḥukmi al-ghazāʾ. Inna rabbaka huwa
al-ʿalīm al-khabīr. Qulta kuntu rāqidan fī al-mihād ayqaẓanī nidāʾu
al-ʿibādi alladhīna ẓulimū ilā an ghariqū fī al-baḥri al-aswad.
Kadhālika samiʿnā wa-rabbuka ʿalā mā aqūlu shahīd. Nashhadu annaka mā
ayqaẓaka al-nidāʾ bal al-hawā li-annā balawnāka wajadnāka fī maʿzil.
Iʿrif laḥna al-qawli wa-kun mina al-mutafarrisīn. Innā mā nuḥibbu an
nurjiʿa ilayka kalimata sūʾin ḥifẓan li-al-maqāmi alladhī aʿṭaynāka fī
al-ḥayāti al-ẓāhirah. Innā ikhtarnā al-adaba wa-jaʿalnāhu sajiyyata
al-muqarrabīn. Innahu thawbun yuwāfiqu al-nufūsa min kulli ṣaghīrin
wa-kabīr. Ṭūbā li-man jaʿalahu ṭirāza haykalihi. Waylun li-man juʿila
maḥrūman min hādhā al-faḍli al-ʿaẓīm. Law kunta ṣāḥiba al-kalimati mā
nabadhta kitāba Allāhi warāʾa ẓahrika idh ursila ilayka min ladun
ʿazīzin ḥakīm. Innā balawnāka bihi mā wajadnāka ʿalā mā iddaʿayta. Qum
wa-tadārak mā fāta ʿanka. Sawfa tafnā al-dunyā wa-mā ʿindaka wa-yabqā
al-mulku li-Allāhi rabbika wa-rabbi ābāʾika al-awwalīn. Lā yanbaghī laka
an taqtaṣira al-umūra ʿalā mā yahwā hawāka. Ittaq zafarāti al-maẓlūm.
Iḥfaẓhu min sihāmi al-ẓālimīn.

King! We heard words you spoke when the king of Russia asked you about
the decision to go to war. Your Lord is indeed the All-Knowing, the
All-Aware. You said: ``I was asleep on my couch when I was awakened by
the cry of the people who were wronged until they drowned in the Black
Sea.'' Thus We heard, and your Lord is witness to what I say. We testify
that it was not the cry that awakened you, but desire, for We tested you
and found you aloof. Discern the meaning of these words and be among
those who perceive. We do not wish to return an evil word to you, out of
regard for the station We have given you in this outward life. We have
chosen courtesy and made it the disposition of those brought near to Us.
It is a garment that befits every soul, great and small. Blessed is
whoever makes it the adornment of his person. Woe to whoever is deprived
of this great bounty. Had you been true to your words, you would not
have cast the Book of God behind your back when it was sent to you by
One who is Almighty, All-Wise. We tested you through it and did not find
you to be as you claimed. Arise and make amends for what you have
missed. The world and all you possess will pass away, and sovereignty
will remain with God, your Lord and the Lord of your forefathers. You
should not confine affairs to what your own desire craves. Beware the
sighs of the Wronged One. Protect Him from the arrows of the oppressors.

\subsection{86}\label{section-1372}

\begin{otherlanguage}{arabic}

بما فعلت تختلف الأمور فی مملکتک و یخرج الملک من کفّک جزآء عملک اذاً تجد
نفسک فی خسران مبین و تأخذ الزّلازل کلّ القبائل هناک الّا بأن تقوم علی نصرة
هذا الأمر و تتّبع الرّوح فی هذا السّبیل المستقیم أ عزّک غرّک لعمری انّه لا
یدوم و سوف یزول الّا بأن تتمسّک بهذا الحبل المتین قد نری الذّلّة تسعی ورائک
و انت من الغافلین ینبغی لک اذا سمعت النّدآء من شطر الکبریآء تدع ما عندک و
تقول لبّیک یا اله من فی السّموات و الأرضین

\end{otherlanguage}

Bi-mā faʿalta takhtalifu al-umūru fī mamlakatika wa-yakhruju al-mulku
min kaffika jazāʾa ʿamalika. Idhan tajidu nafsaka fī khusrānin mubīn.
Wa-taʾkhudhu al-zalāzilu kulla al-qabāʾili hunāka illā bi-an taqūma ʿalā
nuṣrati hādhā al-amri wa-tattabiʿa al-rūḥa fī hādhā al-sabīli
al-mustaqīm. A-ʿizzuka gharraka? La-ʿamrī innahu lā yadūmu wa-sawfa
yazūlu illā bi-an tatamassaka bi-hādhā al-ḥabli al-matīn. Qad narā
al-dhillata tasʿā warāʾaka wa-anta mina al-ghāfilīn. Yanbaghī laka idhā
samiʿta al-nidāʾa min shaṭri al-kibriyāʾ tadaʿa mā ʿindaka wa-taqūla
labbayka yā ilāha man fī al-samāwāti wa-al-arḍīn.

Because of what you have done, affairs in your kingdom will fall into
disorder, and sovereignty will pass from your grasp in requital for your
actions. Then you will find yourself in manifest loss. Tremors will
seize all the peoples there unless you arise to support this Cause and
follow the Spirit on this straight path. Has your might deluded you? By
My life, it will not endure and will soon vanish unless you hold fast to
this strong cord. We see humiliation pursuing you while you remain
heedless. When you hear the call from the realm of majesty, you should
abandon what you possess and say: ``Here I am, O God of all who are in
the heavens and on earth!''

\subsection{87}\label{section-1373}

\begin{otherlanguage}{arabic}

یا ملک قد کنّا بأمّ العراق الی ان حمّ الفراق توجّهنا الی ملک الاسلام بأمره
فلمّا اتیناه ورد علینا من اولی النّفاق ما لا یتمّ بالأوراق بذلک ناح سکّان
الفردوس و اهل حظائر القدس ولکنّ القوم فی حجاب غلیظ

\end{otherlanguage}

Yā malik, qad kunnā bi-ummi al-ʿIrāq ilā an ḥumma al-firāq. Tawajjahnā
ilā maliki al-Islām bi-amrihi. Fa-lammā ataynāhu warada ʿalaynā min ulī
al-nifāqi mā lā yatimmu bi-al-awrāq. Bi-dhālika nāḥa sukkānu al-firdawsi
wa-ahlu ḥaẓāʾiri al-quds wa-lākinna al-qawma fī ḥijābin ghalīẓ.

King! We were in the capital of Iraq until the time of separation was
decreed. We journeyed to the king of Islam at his command. When We
reached him, those given to hypocrisy inflicted on Us what no pages
could fully recount. At this the inhabitants of Paradise and the people
of the holy courts lamented, but the people remain behind a thick veil.

\subsection{88}\label{section-1374}

\begin{otherlanguage}{arabic}

الی ان قلنا قد اشتدّ علینا الأمر فی کلّ یوم بل فی کلّ ساعة الی ان اخرجونا
من السّجن و ادخلونا فی السّجن الأعظم بظلم مبین اذا قیل بأیّ جرم حبسوا قالوا
انّهم ارادوا ان یجدّدوا الدّین لو کان القدیم هو المختار عندکم لم ترکتم ما
شرع فی التّوریة و الانجیل بیّنوا یا قوم لعمری لیس لکم الیوم من محیص ان کان
هذا جرمی قد سبقنی فی ذلک محمّد رسول اللّه و من قبله الرّوح و من قبله الکلیم
و ان کان ذنبی اعلآء کلمة اللّه و اظهار امره فأنا اوّل المذنبین لا ابدّل هذا
الذّنب بملکوت ملک السّموات و الأرضین

\end{otherlanguage}

Ilā an qulnā qad ishtadda ʿalaynā al-amru fī kulli yawmin bal fī kulli
sāʿatin ilā an akhrajūnā mina al-sijni wa-adkhalūnā fī al-sijni
al-aʿẓami bi-ẓulmin mubīn. Idhā qīla bi-ayyi jurmin ḥubisū qālū innahum
arādū an yujaddidū al-dīn. Law kāna al-qadīmu huwa al-mukhtāra ʿindakum
lima taraktum mā shuriʿa fī al-tawrāti wa-al-injīl? Bayyinū yā qawm.
La-ʿamrī laysa lakumu al-yawma min maḥīṣ. In kāna hādhā jurmī qad
sabaqanī fī dhālika Muḥammadun rasūlu Allāh wa-man qablihi al-rūḥ wa-man
qablihi al-kalīm. Wa-in kāna dhanbī iʿlāʾa kalimati Allāhi wa-iẓhāra
amrihi fa-anā awwalu al-mudhnibīn. Lā ubaddilu hādhā al-dhanba
bi-malakūti mulki al-samāwāti wa-al-arḍīn.

Further on We said: ``Our plight grew more severe every day, indeed
every hour, until they took Us out of the prison and cast Us into the
Most Great Prison with manifest injustice. When asked, `For what crime
were they imprisoned?' they said, `They sought to renew religion.' If
you prefer what is old, why have you abandoned what was ordained in the
Torah and the Gospel? Explain yourselves, people! By My life, today you
have no escape. If this is My crime, then Muhammad, the Messenger of
God, committed it before Me, and before Him the Spirit, and before Him
the One who spoke with God. And if My sin is to exalt the word of God
and manifest His Cause, then I am the first of sinners. I would not
exchange this sin for sovereignty over the kingdom of the heavens and
the earth.''

\subsection{89}\label{section-1375}

\begin{otherlanguage}{arabic}

الی ان قلنا کلّما ازداد البلآء زاد البهآء فی حبّ اللّه و امره بحیث ما منعنی
ما ورد علیّ من جنود الغافلین لو یستروننی فی اطباق التّراب یجدوننی راکباً
علی السّحاب و داعیاً الی اللّه المقتدر القدیر انّی فدیت بنفسی فی سبیل اللّه و
اشتاق البلایا فی حبّه و رضائه یشهد بذلک ما انا فیه من البلایا الّتی ما
حملها احد من العالمین و ینطق کلّ شعر من شعراتی بما نطق شجر الطّور و کلّ عرق
من عروقی یدعو اللّه و یقول یا لیت قطعت فى سبیلک لحیوة العالم و اتّحاد من
فیه کذلک قضی الأمر من لدن علیم خبیر

\end{otherlanguage}

Ilā an qulnā kullamā izdāda al-balāʾu zāda al-Bahāʾu fī ḥubbi Allāhi
wa-amrihi bi-ḥaythu mā manaʿanī mā warada ʿalayya min junūdi
al-ghāfilīn. Law yasturūnanī fī aṭbāqi al-turābi yajidūnanī rākiban ʿalā
al-saḥābi wa-dāʿiyan ilā Allāhi al-muqtadiri al-qadīr. Innī fadaytu
bi-nafsī fī sabīli Allāhi wa-ashtāqu al-balāyā fī ḥubbihi wa-riḍāʾihi.
Yashhadu bi-dhālika mā anā fīhi mina al-balāyā allatī mā ḥamalahā aḥadun
mina al-ʿālamīn. Wa-yanṭiqu kullu shaʿrin min shaʿarātī bi-mā naṭaqa
shajaru al-Ṭūr wa-kullu ʿirqin min ʿurūqī yadʿū Allāha wa-yaqūlu yā
layta quṭiʿtu fī sabīlika li-ḥayāti al-ʿālami wa-ittiḥādi man fīhi.
Kadhālika quḍiya al-amru min ladun ʿalīmin khabīr.

Further on We said: ``The greater the affliction, the greater Bahāʾ's
love for God and His Cause, so that nothing inflicted on Me by the hosts
of the heedless has held Me back. Were they to bury Me beneath layers of
earth, they would find Me riding upon the clouds and calling people to
God, the All-Powerful, the Almighty. I have sacrificed Myself in the
path of God and long for afflictions out of love for Him and to attain
His good pleasure. To this bear witness the afflictions I suffer, which
no one in all the world has borne. Every hair of Mine speaks what the
Tree of Sinai spoke, and every vein of Mine calls upon God and says: `If
only I could be severed in Your path, that the world might live and all
within it be united!' Thus has the matter been decreed by One who is
All-Knowing, All-Aware.''

\subsection{90}\label{section-1376}

\begin{otherlanguage}{persian}

و اعلم انّ الرّعیّة امانات اللّه بینکم احفظوهم کما تحفظون انفسکم ایّاکم ان
تجعلوا الذّئاب رعاة الأغنام و ان یمنعکم الغرور و الاستکبار عن التّوجّه الى
الفقرآء و المساکین اطلع من افق الانقطاع باسمی ثمّ اقبل الی الملکوت بأمر
ربّک المقتدر القدیر

\end{otherlanguage}

Wa-iʿlam anna al-raʿiyyata amānātu Allāhi baynakum. Iḥfaẓūhum kamā
taḥfaẓūna anfusakum. Iyyākum an tajʿalū al-dhiʾāba ruʿāta al-aghnāmi
wa-an yamnaʿakumu al-ghurūru wa-al-istikbāru ʿani al-tawajjuhi ilā
al-fuqarāʾi wa-al-masākīn. Iṭlaʿ min ufuqi al-inqiṭāʿi bi-ismī thumma
aqbil ilā al-malakūti bi-amri rabbika al-muqtadiri al-qadīr.

Know that your subjects are God's trusts among you. Protect them as you
protect yourselves. Beware of making wolves the shepherds of the sheep,
and of allowing conceit and arrogance to keep you from attending to the
poor and needy. Rise above the horizon of detachment in My name, then
turn toward the Kingdom at the command of your Lord, the All-Powerful,
the Almighty.

\subsection{91}\label{section-1377}

\begin{otherlanguage}{arabic}

الی ان قلنا زیّن جسد الملک بطراز اسمی و قم علی تبلیغ امری هذا خیر لک ممّا
عندک و یرفع اللّه به اسمک بین الملوک انّه علی کلّ شیء قدیر امش بین النّاس
باسم اللّه و سلطانه لتظهر منک آثاره بین العالمین

\end{otherlanguage}

Ilā an qulnā zayyin jasada al-mulki bi-ṭirāzi ismī wa-qum ʿalā tablīghi
amrī. Hādhā khayrun laka mimmā ʿindaka wa-yarfaʿu Allāhu bihi ismaka
bayna al-mulūk. Innahu ʿalā kulli shayʾin qadīr. Imshi bayna al-nāsi
bi-ismi Allāhi wa-sulṭānihi li-taẓhara minka āthāruhu bayna al-ʿālamīn.

Further on We said: ``Adorn the body of sovereignty with the ornament of
My name and arise to proclaim My Cause. This is better for you than all
you possess, and through it God will exalt your name among kings. He has
power over all things. Walk among the people in the name of God and by
His authority, so that His signs may appear through you among the
peoples of the world.''

\subsection{92}\label{section-1378}

\begin{otherlanguage}{arabic}

الی ان قلنا قل یا قوم هل ینبغی لکم ان تنسبوا انفسکم الى الرّحمن و ترتکبوا
ما ارتکبه الشّیطان لا وجمال السّبحان لو کنتم من العارفین قدّسوا قلوبکم عن
حبّ الدّنیا و السنکم عن الافترآء و ارکانکم عمّا یمنعکم عن التّقرّب الی اللّه
العزیز الحمید قل الدّنیا هی اعراضکم عن مطلع الوحی و اقبالکم الی ما لا
ینفعکم و ما منعکم الیوم عن شطر اللّه انّه اصل الدّنیا اجتنبوها و تقرّبوا الی
المنظر الأکبر هذا المقرّ المشرق المنیر یا قوم لا تسفکوا الدّمآء و لا
تحکموا علی نفس الّا بالحقّ کذلک امرتم من لدن علیم خبیر انّ الّذین یفسدون فی
الأرض بعد اصلاحها اولئک جاوزوا ما حدّد فی الکتاب فبئس مثوی المعتدین

\end{otherlanguage}

Ilā an qulnā qul yā qawm, hal yanbaghī lakum an tansubū anfusakum ilā
al-Raḥmāni wa-tartakibū mā irtakabahu al-shayṭān? Lā wa-jamāli al-subḥān
law kuntum mina al-ʿārifīn. Qaddisū qulūbakum ʿan ḥubbi al-dunyā
wa-alsinatakum ʿani al-iftirāʾi wa-arkānakum ʿammā yamnaʿukum ʿani
al-taqarrubi ilā Allāhi al-ʿazīzi al-ḥamīd. Qul al-dunyā hiya iʿrāḍukum
ʿan maṭlaʿi al-waḥyi wa-iqbālukum ilā mā lā yanfaʿukum. Wa-mā manaʿakumu
al-yawma ʿan shaṭri Allāhi innahu aṣlu al-dunyā. Ijtanibūhā wa-taqarrabū
ilā al-manẓari al-akbari hādhā al-maqarri al-mushriqi al-munīr. Yā qawm,
lā tasfikū al-dimāʾa wa-lā taḥkumū ʿalā nafsin illā bi-al-ḥaqq.
Kadhālika umirtum min ladun ʿalīmin khabīr. Inna alladhīna yufsidūna fī
al-arḍi baʿda iṣlāḥihā ulāʾika jāwazū mā ḥuddida fī al-kitāb. Fa-biʾsa
mathwā al-muʿtadīn.

Further on We said: ``Say: `People! Is it fitting for you to claim
allegiance to the All-Merciful while committing what Satan has
committed? No, by the beauty of the All-Glorious, if you would
understand! Purify your hearts of love for the world, your tongues of
slander, and your limbs of whatever prevents you from drawing near to
God, the Almighty, the Praiseworthy. Say: The world is your turning away
from the Dawning-Place of revelation and turning toward what does not
benefit you. Whatever keeps you today from turning toward God is the
very essence of worldliness. Shun it and draw near to the Most Great
Vision, this radiant, luminous abode. People! Do not shed blood, and do
not pass judgment on anyone except with justice. Thus have you been
commanded by One who is All-Knowing, All-Aware. Those who spread
corruption on earth after it has been set right have overstepped the
bounds prescribed in the Book. Miserable indeed is the abode of the
transgressors.'\,''

\subsection{93}\label{section-1379}

\begin{otherlanguage}{persian}

الی ان قلنا لا تخونوا فی اموال النّاس کونوا امنآء فی الأرض و لا تحرموا
الفقرآء عمّا اتاکم اللّه من فضله انّه یعطیکم ضعف ما عندکم انّه هو المعطى
الکریم یا اهل البهآء سخّروا مدائن القلوب بسیوف الحکمة و البیان انّ الّذین
یجادلون بأهوآء انفسهم اولئک فی حجاب مبین قل سیف الحکمة احرّ من الصّیف و
احدّ من سیف الحدید لو کنتم من العارفین اخرجوه باسمی و سلطانی ثمّ افتحوا به
مدائن افئدة الّذین اعتکفوا فی حصن الهوی کذلک یأمرکم قلم الأبهی اذ کان
جالساً تحت سیوف الغافلین ان اطّلعتم على خطیئة استروها لیستر اللّه عنکم انّه
هو السّتّار ذو الفضل العظیم یا ملأ الأغنیآء ان رأیتم فقیراً لا تستکبروا
علیه تفکّروا فیما خلقتم منه قد خلق کلّ من مآء مهین

\end{otherlanguage}

Ilā an qulnā lā takhūnū fī amwāli al-nās. Kūnū umanāʾa fī al-arḍi wa-lā
taḥrimū al-fuqarāʾa ʿammā ātākumu Allāhu min faḍlihi. Innahu yuʿṭīkum
ḍiʿfa mā ʿindakum. Innahu huwa al-muʿṭī al-karīm. Yā ahla al-Bahāʾ,
sakhkhirū madāʾina al-qulūbi bi-suyūfi al-ḥikmati wa-al-bayān. Inna
alladhīna yujādilūna bi-ahwāʾi anfusihim ulāʾika fī ḥijābin mubīn. Qul
sayfu al-ḥikmati aḥarru mina al-ṣayfi wa-aḥaddu min sayfi al-ḥadīd law
kuntum mina al-ʿārifīn. Akhrijūhu bi-ismī wa-sulṭānī thumma iftaḥū bihi
madāʾina afʾidati alladhīna iʿtakafū fī ḥiṣni al-hawā. Kadhālika
yaʾmurukum qalamu al-Abhā idh kāna jālisan taḥta suyūfi al-ghāfilīn. In
iṭṭalaʿtum ʿalā khaṭīʾatin usturūhā li-yastura Allāhu ʿankum. Innahu
huwa al-sattāru dhū al-faḍli al-ʿaẓīm. Yā malaʾa al-aghniyāʾ, in
raʾaytum faqīran lā tastakbirū ʿalayhi. Tafakkarū fīmā khuliqtum minhu.
Qad khuliqa kullun min māʾin mahīn.

Further on We said: ``Do not act dishonestly with the property of
others. Be trustworthy on earth, and do not deprive the poor of what God
has given you from His bounty. He will give you twice what you possess.
He is indeed the Giver, the All-Bountiful. People of Bahāʾ! Conquer the
cities of hearts with the swords of wisdom and utterance. Those who
dispute according to their own desires are indeed behind a manifest
veil. Say: The sword of wisdom is hotter than summer and sharper than a
sword of steel, if you would understand. Draw it in My name and by My
authority, then conquer with it the cities of the hearts of those who
have secluded themselves in the stronghold of desire. Thus the Pen of
the All-Glorious commands you as He sits beneath the swords of the
heedless. If you become aware of a sin, conceal it, so that God may
conceal yours. He is indeed the Concealer, possessed of immense bounty.
You who are rich! If you see a poor person, do not treat him arrogantly.
Reflect on what you were created from. Everyone was created from a lowly
fluid.''

\subsection{94}\label{section-1380}

\begin{otherlanguage}{arabic}

الی ان قلنا انظروا العالم کهیکل انسان اعترته الأمراض و برئه منوط باتّحاد
من فیه اجتمعوا على ما شرّعناه لکم و لا تتّبعوا سبل المختلفین تفکّر فی
الدّنیا و شأن اهلها انّ الّذى خلق العالم لنفسه قد حبس فی اخرب الدّیار بما
اکتسبت ایدی الغافلین و من افق السّجن یدعو النّاس الى فجر اللّه العلیّ العظیم
هل تفرح بما عندک من الزّخارف بعد اذ تعلم انّها ستفنی او تسترّ بما تحکم علی
شبر من الأرض بعد اذ کلّها لم تکن عند اهل البهآء الّا کسواد عین نملة میّتة
دعها لأهلها ثمّ اقبل الى مقصود العالمین این اهل الغرور و قصورهم انظر فی
قبورهم لتعتبر بما جعلناها عبرة للنّاظرین لو تأخذک نفحات الوحی لتفرّ من
الملک مقبلاً الی الملکوت و تنفق ما عندک للتّقرّب الی هذا المنظر الکریم

\end{otherlanguage}

Ilā an qulnā: Unẓurū al-ʿālama ka-haykali insānin iʿtarat-hu al-amrāḍu
wa-burʾuhu manūṭun bi-ittiḥādi man fīhi. Ijtamiʿū ʿalā mā sharaʿnāhu
lakum wa-lā tattabiʿū subula al-mukhtalifīn. Tafakkar fī al-dunyā
wa-shaʾni ahlihā. Inna alladhī khalaqa al-ʿālama li-nafsihi qad ḥubisa
fī akhrabi al-diyāri bi-mā iktasabat aydī al-ghāfilīn, wa-min ufuqi
al-sijni yadʿū al-nāsa ilā fajri Allāhi al-ʿAliyyi al-ʿAẓīm. Hal tafraḥu
bi-mā ʿindaka mina al-zakhārifi baʿda idh taʿlamu annahā satafnā, aw
tastarru bi-mā taḥkumu ʿalā shibrin mina al-arḍi baʿda idh kulluhā lam
takun ʿinda ahli al-Bahāʾi illā ka-sawādi ʿayni namlatin mayyitah? Daʿhā
li-ahlihā thumma aqbil ilā maqṣūdi al-ʿālamīn. Ayna ahlu al-ghurūri
wa-quṣūruhum? Unẓur fī qubūrihim li-taʿtabira bi-mā jaʿalnāhā ʿibratan
li-l-nāẓirīn. Law taʾkhudhuka nafaḥātu al-waḥyi la-tafirru mina al-mulki
muqbilan ilā al-malakūti wa-tunfiqu mā ʿindaka li-l-taqarrubi ilā hādhā
al-manẓari al-karīm.

Further on We said: ``Regard the world as a human body afflicted with
diseases, whose recovery depends on the unity of those within it. Unite
around what We have prescribed for you, and do not follow the paths of
those who differ. Reflect on the world and the condition of its people.
He who created the world for Himself has been imprisoned in the most
desolate of lands because of what the hands of the heedless have done.
From the horizon of the prison He calls the people to the dawn of God,
the Most High, the Most Great. Do you rejoice in the ornaments you
possess, knowing that they will perish, or delight in ruling over a span
of earth, when all of it is, to the people of Bahāʾ, no more than the
black of a dead ant's eye? Leave it to its people, and turn toward the
One desired by all the worlds. Where are the arrogant and their palaces?
Look at their graves and take heed, for We have made them a lesson for
those who look. If the breezes of revelation were to seize you, you
would flee from earthly dominion toward the Kingdom and spend all you
possess to draw near to this gracious Countenance.''

\subsection{95}\label{section-1381}

\begin{otherlanguage}{arabic}

و امر نمودیم یکى از ملأ حضرت روح این لوح را ارسال دارد و او ذکر نمود لوح
و ترجمهٔ آن را ارسال داشتم العلم عند اللّه العزیز العلّام

\end{otherlanguage}

Va amr namūdīm yekī az malaʾ-i Ḥaḍrat-i Rūḥ īn lawḥ rā irsāl dārad va ū
dhikr namūd: Lawḥ va tarjumeh-yi ān rā irsāl dāshtam. Al-ʿilmu ʿinda
Allāhi al-ʿAzīzi al-ʿAllām.

We instructed one of the clergy of His Holiness the Spirit to send this
Tablet, and he stated: ``I have sent the Tablet and its translation.''
Knowledge rests with God, the Almighty, the All-Knowing.

\subsection{96}\label{section-1382}

\begin{otherlanguage}{persian}

و یکى از اجزاء هیکل لوح حضرت امپراطور روس ایّده اللّه تبارک و تعالی است

\end{otherlanguage}

Va yekī az ajzā-yi haykal lawḥ-i ḥaḍrat-i impirāṭūr-i Rūs, ayyadahu
Allāhu tabāraka wa-taʿālā, ast.

One of the constituent parts of this body is the Tablet to His Majesty
the Emperor of Russia---may God, blessed and exalted be He, assist him.

\subsection{97}\label{section-1383}

\begin{otherlanguage}{persian}

یا ملک الرّوس اسمع ندآء اللّه الملک القدّوس و اقبل الى الفردوس المقرّ الّذی
فیه استقرّ من سمّی بالأسمآء الحسنی بین الملإ الأعلی و فی ملکوت الانشآء
باسم اللّه البهیّ الأبهی ایّاک ان یحجبک شیء عن التّوجّه الی ربّک الرّحمن الرّحیم
انّا سمعنا ما نادیت به مولاک فی نجواک لذا هاج عرف عنایتی و ماج بحر رحمتی
اجبناک بالحقّ انّ ربّک هو العلیم الحکیم قد نصرنی احد سفرائِک اذ کنت فی سجن
الطّآء تحت السّلاسل و الأغلال بذلک کتب اللّه لک مقاماً لم یحط به علم احد الّا
هو ایّاک ان تبدّل هذا المقام العظیم

\end{otherlanguage}

Yā malika al-Rūs, ismaʿ nidāʾa Allāhi al-Maliki al-Quddūsi wa-aqbil ilā
al-firdawsi al-maqarri alladhī fīhi istaqarra man summiya bi-l-asmāʾi
al-ḥusnā bayna al-malaʾi al-aʿlā wa-fī malakūti al-inshāʾi bi-smi Allāhi
al-Bahiyyi al-Abhā. Iyyāka an yaḥjubaka shayʾun ʿani al-tawajjuhi ilā
rabbika al-Raḥmāni al-Raḥīm. Innā samiʿnā mā nādayta bihi mawlāka fī
najwāka, li-dhā hāja ʿarfu ʿināyatī wa-māja baḥru raḥmatī. Ajabnāka
bi-l-ḥaqqi, inna rabbaka huwa al-ʿAlīmu al-Ḥakīm. Qad naṣaranī aḥadu
sufarāʾika idh kuntu fī sijni al-Ṭāʾi taḥta al-salāsili wa-l-aghlāl.
Bi-dhālika kataba Allāhu laka maqāman lam yuḥiṭ bihi ʿilmu aḥadin illā
huwa. Iyyāka an tubaddila hādhā al-maqāma al-ʿaẓīm.

``King of Russia! Hear the call of God, the King, the Most Holy, and
turn toward Paradise, the dwelling where He abides who is called by the
most beautiful names among the heavenly concourse, and in the kingdom of
creation by the name of God, the All-Glorious, the Most Glorious. Beware
lest anything keep you from turning toward your Lord, the All-Merciful,
the Merciful. We heard what you called out to your Master in your secret
prayer; therefore the fragrance of My loving care stirred and the ocean
of My mercy surged. We answered you in truth. Your Lord is indeed the
All-Knowing, the All-Wise. One of your ambassadors came to My aid when I
was in the prison of Ṭāʾ, bound in chains and fetters. For this, God has
ordained for you a station that no one's knowledge encompasses except
His own. Beware lest you exchange this great station for another.''

\subsection{98}\label{section-1384}

\begin{otherlanguage}{arabic}

الی ان قلنا قد اتی الأب و الابن فی الواد المقدّس یقول لبّیک اللّهمّ لبّیک و
الطّور یطوف حول البیت و الشّجر ینادی بأعلى النّدآء قد اتی الوهّاب راکباً علی
السّحاب طوبی لمن تقرّب الیه ویل للمبعدین

\end{otherlanguage}

Ilā an qulnā: Qad atā al-Abu wa-l-Ibnu fī al-wādi al-muqaddasi yaqūlu:
Labbayka Allāhumma labbayk. Wa-l-Ṭūru yaṭūfu ḥawla al-Bayti wa-l-shajaru
yunādī bi-aʿlā al-nidāʾi: Qad atā al-Wahhābu rākiban ʿalā al-saḥāb. Ṭūbā
li-man taqarraba ilayhi, waylun li-l-mubʿadīn.

Further on We said: ``The Father has come, and the Son in the holy
valley says: `Here I am, O God, here I am!' Sinai circles around the
House, and the Tree calls out at the top of its voice: `The
All-Bountiful has come, riding upon the clouds!' Blessed is whoever
draws near to Him; woe to those who remain far away!''

\subsection{99}\label{section-1385}

\begin{otherlanguage}{arabic}

قم بین النّاس بهذا الأمر المبرم ثمّ ادع الأمم الی اللّه العزیز العظیم لا
تکن من الّذین کانوا ان یدعوا اللّه باسم من الأسمآء فلمّا اتی المسمّى کفروا
به و اعرضوا عنه الی ان افتوا علیه بظلم مبین انظر ثمّ اذکر الأیّام الّتی
فیها اتی الرّوح و حکم علیه هیرودس قد نصر اللّه الرّوح بجنود الغیب و حفظه
بالحقّ و ارسله الی ارض اخری وعداً من عنده انّه هو الحاکم علی ما یرید انّ ربّک
یحفظ من یشآء لو یکون فى قطب البحر او فی فم الثّعبان او تحت سیوف الظّالمین

\end{otherlanguage}

Qum bayna al-nāsi bi-hādhā al-amri al-mubrami thumma udʿu al-umama ilā
Allāhi al-ʿAzīzi al-ʿAẓīm. Lā takun mina alladhīna kānū an yadʿū Allāha
bi-smin mina al-asmāʾi fa-lammā atā al-musammā kafarū bihi wa-aʿraḍū
ʿanhu ilā an aftaw ʿalayhi bi-ẓulmin mubīn. Unẓur thumma udhkur
al-ayyāma allatī fīhā atā al-Rūḥu wa-ḥakama ʿalayhi Hīrūdus. Qad naṣara
Allāhu al-Rūḥa bi-junūdi al-ghaybi wa-ḥafiẓahu bi-l-ḥaqqi wa-arsalahu
ilā arḍin ukhrā waʿdan min ʿindihi. Innahu huwa al-ḥākimu ʿalā mā yurīd.
Inna rabbaka yaḥfaẓu man yashāʾu law yakūnu fī quṭbi al-baḥri aw fī fami
al-thuʿbāni aw taḥta suyūfi al-ẓālimīn.

``Rise among the people with this irrevocable Cause, then summon the
nations to God, the Almighty, the Most Great. Do not be among those who
called upon God by one of His names, but when the One bearing that name
came, disbelieved in Him and turned away from Him, until they pronounced
judgment against Him with manifest injustice. Consider and recall the
days when the Spirit came and Herod passed judgment against Him. God
aided the Spirit with the hosts of the unseen, protected Him in truth,
and sent Him to another land, in fulfillment of a promise from Him. He
indeed decrees whatever He wills. Your Lord protects whomever He wills,
even in the midst of the sea, in the serpent's mouth, or beneath the
swords of the unjust.''

\subsection{100}\label{section-1386}

\begin{otherlanguage}{arabic}

الی ان قلنا اسمع ندائی مرّة اخری من شطر سجنی لیخبرک بما ورد علی جمالی من
مظاهر جلالی و تعرف صبری بعد قدرتی و اصطباری بعد اقتداری وعمری لو تعرف ما
نزّل من قلمی و تطّلع علی خزائن امری و لآلئ اسراری فی بحور اسمائی و اواعی
کلماتی لتفدی بنفسک فی سبیل اللّه شوقاً الی ملکوته العزیز المنیع اعلم جسمی
تحت سیوف الأعدآء و جسدی فی بلآء لا یحصی ولکنّ الرّوح فی بشارة لا یعادلها
فرح العالمین

\end{otherlanguage}

Ilā an qulnā: Ismaʿ nidāʾī marratan ukhrā min shaṭri sijnī li-yukhbiraka
bi-mā warada ʿalā jamālī min maẓāhiri jalālī wa-taʿrifa ṣabrī baʿda
qudratī wa-iṣṭibārī baʿda iqtidārī. Wa-ʿamrī, law taʿrifu mā nuzzila min
qalamī wa-taṭṭaliʿu ʿalā khazāʾini amrī wa-laʾāliʾi asrārī fī buḥūri
asmāʾī wa-awāʿī kalimātī la-tafdī bi-nafsika fī sabīli Allāhi shawqan
ilā malakūtihi al-ʿAzīzi al-Manīʿ. Iʿlam, jismī taḥta suyūfi al-aʿdāʾi
wa-jasadī fī balāʾin lā yuḥṣā, wa-lākinna al-rūḥa fī bishāratin lā
yuʿādiluhā faraḥu al-ʿālamīn.

Further on We said: ``Hear My call once more from the direction of My
prison, that it may tell you what has befallen My Beauty at the hands of
the manifestations of My majesty, and that you may know My patience
despite My power and My forbearance despite My might. By My life! If you
knew what has been sent down from My Pen and became acquainted with the
treasuries of My Cause and the pearls of My mysteries in the oceans of
My names and the vessels of My words, you would sacrifice your life in
the path of God out of longing for His mighty, inaccessible Kingdom.
Know that My body is beneath the swords of enemies and My flesh suffers
countless afflictions, yet My spirit is filled with a gladness that the
joy of all the worlds cannot equal.''

\subsection{101}\label{section-1387}

\begin{otherlanguage}{persian}

و همچنین بعضی از آیات لوح حضرت ملکه ایّدها اللّه تبارک و تعالی ذکر میشود
مقصود آنکه شاید نفحات وحی آن جناب را اخذ نماید و لوجه اللّه بر خدمت امر
قیام کنند و آنچه از الواح حضرات ملوک وفّقهم اللّه نرسیده برسانند این امر
عظیم است این خدمت عظیم در آن ارض علمای اعلام بسیارند از جمله حضرات سادات
که در بزرگی و شأن معروفند با ایشان مشورت نمائید و آنچه از قلم اعلی جاری
شده بنمائید شاید بر اصلاح عالم و تهذیب نفوس امم مؤیّد شوند و ضغینه و
بغضای مکنونهٔ مخزونهٔ در افئده را بکوثر نصایح الهی ساکن نمایند این توفیق
را از حقّ میطلبیم و لیس هذا علیه بعزیز

\end{otherlanguage}

Va hamchunīn baʿḍī az āyāt-i lawḥ-i ḥaḍrat-i malakeh, ayyadahā Allāhu
tabāraka wa-taʿālā, dhikr mī-shavad. Maqṣūd ān-ke shāyad nafaḥāt-i vaḥy
ān janāb rā akhdh namāyad va li-vajhi Allāh bar khidmat-i amr qiyām
kunand va ānchih az alvāḥ-i ḥaḍarāt-i mulūk, waffaqahumu Allāh,
narasīdeh birasānand. Īn amr ʿaẓīm ast, īn khidmat ʿaẓīm. Dar ān arḍ
ʿulamā-yi aʿlām bisyārand, az jumleh ḥaḍarāt-i sādāt keh dar buzurgī va
shaʾn maʿrūfand. Bā īshān mashvarat namāyīd va ānchih az qalam-i aʿlā
jārī shudeh binamāyīd. Shāyad bar iṣlāḥ-i ʿālam va tahdhīb-i nufūs-i
umam muʾayyad shavand va ḍaghīneh va baghḍā-yi maknūneh-yi makhzūneh dar
afʾideh rā bi-kawthar-i naṣāyiḥ-i ilāhī sākin namāyand. Īn tawfīq rā az
Ḥaqq mī-ṭalabīm, wa-laysa hādhā ʿalayhi bi-ʿazīz.

Likewise, some of the verses of the Tablet to Her Majesty the
Queen---may God, blessed and exalted be He, assist her---are cited. The
purpose is that the breezes of revelation may perhaps seize you, and
that you may arise to serve the Cause for God's sake and deliver
whichever of the Tablets to Their Majesties the kings---may God grant
them success---have not reached them. This is a great undertaking, a
great service. In that land there are many distinguished scholars, among
them the honored sayyids, who are renowned for their greatness and
standing. Consult with them and show them what has flowed from the Most
Exalted Pen. Perhaps they may be assisted in reforming the world and
refining the souls of the nations, and may quell the rancor and hatred
concealed and stored in hearts with the life-giving waters of divine
counsels. We ask God to grant this success, and this is not difficult
for Him.

\subsection{102}\label{section-1388}

\begin{otherlanguage}{arabic}

یا ایّتها الملکة فی لندن اسمعی ندآء ربّک مالک البریّة من السّدرة الالهیّة انّه
لا اله الّا انا العزیز الحکیم ضعی ما علی الأرض و زیّنی رأس الملک باکلیل
ذکر ربّک الجلیل انّه قد اتی فی العالم بمجده الأعظم و کمل ما ذکر فی الانجیل
قد تشرّف برّ الشّام بقدوم ربّه مالک الأنام و اخذ سکر خمر الوصال شطر الجنوب و
الشّمال طوبی لمن وجد عرف الرّحمن و اقبل الی مشرق الجمال فی هذا الفجر
المنیر قد اهتزّ المسجد الأقصی من نسمات ربّه الأبهی و البطحآء من ندآء اللّه
العلیّ الأعلی و کلّ حصاة منها تسبّح الرّبّ بهذا الاسم العظیم

\end{otherlanguage}

Yā ayyatuhā al-malikatu fī Landan, ismaʿī nidāʾa rabbiki Māliki
al-bariyyati mina al-sidrati al-ilāhiyyah: Innahu lā ilāha illā anā
al-ʿAzīzu al-Ḥakīm. Ḍaʿī mā ʿalā al-arḍi wa-zayyinī raʾsa al-mulki
bi-iklīli dhikri rabbiki al-Jalīl. Innahu qad atā fī al-ʿālami
bi-majdihi al-aʿẓami wa-kamula mā dhukira fī al-Injīl. Qad tasharrafa
barru al-Shāmi bi-qudūmi rabbihi Māliki al-anāmi wa-akhadha sukru khamri
al-wiṣāli shaṭra al-janūbi wa-l-shamāl. Ṭūbā li-man wajada ʿarfa
al-Raḥmāni wa-aqbala ilā mashriqi al-jamāli fī hādhā al-fajri al-munīr.
Qad ihtazza al-Masjidu al-Aqṣā min nasamāti rabbihi al-Abhā wa-l-Baṭḥāʾu
min nidāʾi Allāhi al-ʿAliyyi al-Aʿlā, wa-kullu ḥaṣātin minhā tusabbiḥu
al-Rabba bi-hādhā al-ismi al-ʿaẓīm.

``Queen in London! Hear the call of your Lord, the Sovereign of
creation, from the divine Lote-Tree: `There is no God but Me, the
Almighty, the All-Wise.' Set aside what is on earth, and adorn the head
of sovereignty with the crown of remembrance of your Lord, the
All-Glorious. He has come into the world in His most great glory, and
what was recorded in the Gospel has been fulfilled. The land of Syria
has been honored by the arrival of its Lord, the Sovereign of humankind,
and the intoxication of the wine of reunion has seized the regions of
the south and the north. Blessed is whoever has perceived the fragrance
of the All-Merciful and turned toward the Dawning-Place of Beauty in
this radiant dawn. The Aqṣā Mosque has stirred at the breezes of its
Lord, the Most Glorious, and Baṭḥāʾ at the call of God, the Most High,
the Most Exalted; every pebble within it glorifies the Lord through this
great Name.''

\subsection{103}\label{section-1389}

\begin{otherlanguage}{persian}

الی ان قلنا انّا نذکرک لوجه اللّه و نحبّ ان یعلو اسمک بذکر ربّک خالق الأرض و
السّمآء انّه علی ما اقول شهید قد بلغنا انّک منعت بیع الغلمان و الامآء هذا
ما حکم به اللّه فی هذا الظّهور البدیع قد کتب اللّه لک جزآء ذلک انّه موفّی
اجور المحسنات و المحسنین ان تتّبعی ما ارسل الیک من لدن علیم خبیر انّ الّذی
اعرض و استکبر بعدما جآءته البیّنات من لدن منزل الآیات لیحبط اللّه عمله انّه
علی کلّ شیء قدیر انّ الأعمال تقبل بعد الاقبال من اعرض عن الحقّ انّه من احجب
الخلق کذلک قدّر من لدن عزیز قدیر

\end{otherlanguage}

Ilā an qulnā: Innā nadhkuruki li-wajhi Allāhi wa-nuḥibbu an yaʿluwa
ismuki bi-dhikri rabbiki Khāliqi al-arḍi wa-l-samāʾ. Innahu ʿalā mā
aqūlu shahīd. Qad balaghanā annaki manaʿti bayʿa al-ghilmāni wa-l-imāʾ.
Hādhā mā ḥakama bihi Allāhu fī hādhā al-ẓuhūri al-badīʿ. Qad kataba
Allāhu laki jazāʾa dhālika, innahu muwaffī ujūri al-muḥsināti
wa-l-muḥsinīn, in tattabiʿī mā ursila ilayki min ladun ʿAlīmin Khabīr.
Inna alladhī aʿraḍa wa-istakbara baʿdamā jāʾat-hu al-bayyinātu min ladun
Munzili al-āyāti la-yuḥbiṭu Allāhu ʿamalahu. Innahu ʿalā kulli shayʾin
qadīr. Inna al-aʿmāla tuqbalu baʿda al-iqbāl. Man aʿraḍa ʿani al-Ḥaqqi
innahu min aḥjabi al-khalq. Kadhālika quddira min ladun ʿAzīzin Qadīr.

Further on We said: ``We make mention of you for God's sake and desire
that your name be exalted through remembrance of your Lord, the Creator
of earth and heaven. He is witness to what I say. We have learned that
you have forbidden the sale of male and female slaves. This is what God
has decreed in this wondrous Revelation. God has ordained a reward for
you for this---He indeed pays in full the rewards of women and men who
do good---if you follow what has been sent to you by One who is
All-Knowing, All-Aware. Whoever turns away and acts arrogantly after the
clear proofs have come to him from the One who sends down the verses,
God will assuredly bring his deeds to nothing. He has power over all
things. Deeds are accepted after one has turned toward Him. Whoever
turns away from the True One is among the most veiled of His creatures.
Thus has it been decreed by One who is Almighty, All-Powerful.''

\subsection{104}\label{section-1390}

\begin{otherlanguage}{arabic}

و سمعنا انّک اودعت زمام المشاورة بأیادی الجمهور نعم ما عملت لأنّ بها
تستحکم اصول ابنیة الأمور و تطمئنّ قلوب من فی ظلّک من کلّ وضیع و شریف ولکن
ینبغی لهم ان یکونوا امنآء بین العباد و یرون انفسهم وکلآء لمن علی الأرض
کلّها هذا ما وعظوا به فی اللّوح من لدن مدبّر حکیم و اذا توجّه احد الی المجمع
یحوّل طرفه الی الأفق الأعلی و یقول یا الهی اسألک باسمک الأبهی ان تؤیّدنی
علی ما تصلح به امور عبادک و تعمر به بلادک انّک انت علی کلّ شیء قدیر طوبی
لمن یدخل المجمع لوجه اللّه و یحکم بین النّاس بالعدل الخالص الا انّه من
الفائزین

\end{otherlanguage}

Wa-samiʿnā annaki awdaʿti zimāma al-mushāwarati bi-ayādī al-jumhūr.
Niʿma mā ʿamilti li-anna bihā tastaḥkimu uṣūlu abniyati al-umūri
wa-taṭmaʾinnu qulūbu man fī ẓilliki min kulli waḍīʿin wa-sharīf.
Wa-lākin yanbaghī lahum an yakūnū umanāʾa bayna al-ʿibādi wa-yarawna
anfusahum wukalāʾa li-man ʿalā al-arḍi kullihā. Hādhā mā wuʿiẓū bihi fī
al-lawḥi min ladun Mudabbirin Ḥakīm. Wa-idhā tawajjaha aḥadun ilā
al-majmaʿi yuḥawwilu ṭarfahu ilā al-ufuqi al-aʿlā wa-yaqūlu: Yā ilāhī,
asʾaluka bi-smika al-Abhā an tuʾayyidanī ʿalā mā taṣluḥu bihi umūru
ʿibādika wa-taʿmuru bihi bilāduka. Innaka anta ʿalā kulli shayʾin qadīr.
Ṭūbā li-man yadkhulu al-majmaʿa li-wajhi Allāhi wa-yaḥkumu bayna al-nāsi
bi-l-ʿadli al-khāliṣ. Alā innahu mina al-fāʾizīn.

``We have also heard that you have entrusted the reins of consultation
to the people. Well have you done, for through it the foundations of the
structure of affairs are strengthened and the hearts of all who live
under your protection, both lowly and noble, are reassured. Yet they
should be trustworthy among the people and regard themselves as
representatives of all who dwell on earth. This is the counsel given to
them in the Tablet by One who is the All-Wise Ordainer. When anyone goes
to the assembly, he should turn his gaze toward the Most Exalted Horizon
and say: `My God! I ask You by Your Most Glorious Name to assist me in
whatever will set right the affairs of Your servants and make Your lands
flourish. You indeed have power over all things.' Blessed is whoever
enters the assembly for God's sake and judges between people with pure
justice. He is indeed among those who attain success.''

\subsection{105}\label{section-1391}

\begin{otherlanguage}{arabic}

یا اصحاب المجالس هناک و فی دیار اخری تدبّروا و تکلّموا فیما یصلح به العالم
و حاله لو کنتم من المتوسّمین انظروا العالم کهیکل انسان انّه خلق صحیحاً
کاملاً اعترته الأمراض بالأسباب المختلفة المتغایرة و ما طابت نفسه یوماً بل
اشتدّ مرضه بما وقع تحت تصرّف المتطبّبین الّذین رکبوا مطیّة الهوی و کانوا من
الهائمین الّا من شآء اللّه ربّ العالمین و ان طاب عضو من اعضائه فی عصر من
الأعصار بطبیب حاذق بقیت اعضآء اخری فیما کان کذلک ینبّئکم العلیم الخبیر و
الیوم نراه تحت ایادی الّذین اخذهم سکر خمر الغرور بحیث لا یعرفون خیر
انفسهم فکیف هذا الأمر الأوعر الخطیر

\end{otherlanguage}

Yā aṣḥāba al-majālisi hunāka wa-fī diyārin ukhrā, tadabbarū wa-takallamū
fī-mā yaṣluḥu bihi al-ʿālamu wa-ḥāluhu law kuntum mina al-mutawassimīn.
Unẓurū al-ʿālama ka-haykali insānin, innahu khuliqa ṣaḥīḥan kāmilan,
iʿtarat-hu al-amrāḍu bi-l-asbābi al-mukhtalifati al-mutaghāyirah. Wa-mā
ṭābat nafsuhu yawman, bal ishtadda maraḍuhu bi-mā waqaʿa taḥta taṣarrufi
al-mutaṭabbibīna alladhīna rakibū maṭiyyata al-hawā wa-kānū mina
al-hāʾimīn, illā man shāʾa Allāhu Rabbu al-ʿālamīn. Wa-in ṭāba ʿuḍwun
min aʿḍāʾihi fī ʿaṣrin mina al-aʿṣāri bi-ṭabībin ḥādhiqin baqiyat
aʿḍāʾun ukhrā fī-mā kān. Kadhālika yunabbiʾukumu al-ʿAlīmu al-Khabīr.
Wa-l-yawma narāhu taḥta ayādī alladhīna akhadhahum sukru khamri
al-ghurūri bi-ḥaythu lā yaʿrifūna khayra anfusihim, fa-kayfa hādhā
al-amru al-awʿaru al-khaṭīr?

``Members of assemblies there and in other lands! Deliberate and speak
about what will improve the world and its condition, if you are
discerning. Regard the world as a human body: it was created sound and
whole, but diseases have afflicted it through various and differing
causes. Not for a day has it recovered; rather, its illness has grown
worse because it has fallen under the control of would-be physicians who
have mounted the steed of desire and wandered astray, except for those
whom God, the Lord of the worlds, wills otherwise. Even if a skilled
physician has healed one of its limbs in some age, other limbs have
remained as they were. Thus the All-Knowing, the All-Aware informs you.
Today We see it in the hands of those so intoxicated with the wine of
pride that they do not know their own good---how much less this most
difficult and momentous matter!''

\subsection{106}\label{section-1392}

\begin{otherlanguage}{arabic}

الی ان قلنا و ما جعله اللّه الدّریاق الأعظم و السّبب الأتمّ لصحّته هو اتّحاد
من علی الأرض علی امر واحد و شریعة واحدة هذا لا یمکن ابداً الّا بطبیب حاذق
کامل مؤیّد لعمری هذا هو الحقّ و ما بعده الّا الضّلال المبین کلّما اتی ذاک
السّبب الأعظم و اشرق ذاک النّور من مشرق القدم منعه المتطبّبون و صاروا سحاباً
بینه و بین العالم لذا ما طاب مرضه و بقی فی سقمه الی الحین انّهم لم یقدروا
على حفظه و صحّته و الّذى کان مظهر القدرة بین البریّة منع عمّا اراد بما
اکتسبت ایدی المتطبّبین

\end{otherlanguage}

Ilā an qulnā: Wa-mā jaʿalahu Allāhu al-diryāqa al-aʿẓama wa-l-sababa
al-atamma li-ṣiḥḥatihi huwa ittiḥādu man ʿalā al-arḍi ʿalā amrin wāḥidin
wa-sharīʿatin wāḥidah. Hādhā lā yumkinu abadan illā bi-ṭabībin ḥādhiqin
kāmilin muʾayyad. La-ʿamrī, hādhā huwa al-ḥaqqu wa-mā baʿdahu illā
al-ḍalālu al-mubīn. Kullamā atā dhāka al-sababu al-aʿẓamu wa-ashraqa
dhāka al-nūru min mashriqi al-qidami manaʿahu al-mutaṭabbibūna wa-ṣārū
saḥāban baynahu wa-bayna al-ʿālam. Li-dhā mā ṭāba maraḍuhu wa-baqiya fī
saqamihi ilā al-ḥīn. Innahum lam yaqdirū ʿalā ḥifẓihi wa-ṣiḥḥatihi,
wa-alladhī kāna maẓhara al-qudrati bayna al-bariyyati muniʿa ʿammā arāda
bi-mā iktasabat aydī al-mutaṭabbibīn.

Further on We said: ``What God has made the supreme remedy and the most
complete means of its recovery is the union of all who dwell on earth in
one Cause and one law. This can never be achieved except through a
skilled, perfect, divinely assisted Physician. By My life! This is the
truth, and beyond it there is nothing but manifest error. Whenever that
supreme means has appeared and that Light has shone from the
Dawning-Place of eternity, the would-be physicians have obstructed Him
and become clouds between Him and the world. Thus its illness has not
been cured, and it has remained sick to this day. They have been unable
to preserve it and restore its health, while He who was the
manifestation of power among creation was prevented from accomplishing
His purpose by what the hands of the would-be physicians had done.''

\subsection{107}\label{section-1393}

\begin{otherlanguage}{arabic}

انظروا فی هذه الأیّام الّتی اتی جمال القدم بالاسم الأعظم لحیوة العالم و
اتّحادهم انّهم قاموا علیه بأسیاف شاحذة و ارتکبوا ما فزع به الرّوح الأمین
الی ان جعلوه مسجوناً فی اخرب البلاد و انقطعت عن ذیله ایادی المقبلین اذا
قیل لهم اتی مصلح العالم قالوا قد تحقّق انّه من المفسدین مع انّهم ما عاشروه
و یرون انّه ما حفظ نفسه فی اقلّ من حین کان فی کلّ الأحیان بین ایادی اهل
الطّغیان مرّةً حبسوه و طوراً اخرجوه و تارةً اداروا به البلاد کذلک حکموا علینا
و اللّه بما اقول علیم

\end{otherlanguage}

Unẓurū fī hādhihi al-ayyāmi allatī atā Jamālu al-Qidami bi-l-ismi
al-aʿẓami li-ḥayāti al-ʿālami wa-ittiḥādihim. Innahum qāmū ʿalayhi
bi-asyāfin shāḥidhatin wa-irtakabū mā faziʿa bihi al-Rūḥu al-Amīnu ilā
an jaʿalūhu masjūnan fī akhrabi al-bilādi wa-inqaṭaʿat ʿan dhaylihi
ayādī al-muqbilīn. Idhā qīla lahum atā muṣliḥu al-ʿālami qālū qad
taḥaqqaqa annahu mina al-mufsidīn, maʿa annahum mā ʿāsharūhu wa-yarawna
annahu mā ḥafiẓa nafsahu fī aqalla min ḥīn. Kāna fī kulli al-aḥyāni
bayna ayādī ahli al-ṭughyān. Marratan ḥabasūhu wa-ṭawran akhrajūhu
wa-tāratan adārū bihi al-bilād. Kadhālika ḥakamū ʿalaynā wa-Allāhu bi-mā
aqūlu ʿalīm.

``Consider these days, in which the Ancient Beauty has come with the
Most Great Name to give life to the world and unite its people. They
have risen against Him with sharpened swords and committed what has
filled the Faithful Spirit with dread, until they imprisoned Him in the
most desolate of lands and cut off the hands of those who turn toward
Him from the hem of His robe. When they are told, `The Reformer of the
world has come,' they say, `It has been established that He is among
those who spread corruption,' although they have not associated with Him
and see that He has not protected Himself for even the briefest moment.
He has always been in the hands of the tyrannical. At one time they
imprisoned Him, at another they expelled Him, and at yet another they
drove Him from land to land. Thus have they passed judgment upon Us, and
God knows what I say.''

\subsection{108}\label{section-1394}

\begin{otherlanguage}{persian}

این نسبت فساد همان نسبت فسادیست که از قبل فراعنه بحضرت کلیم اللّه داده‌اند
اقرأ ما انزله الرّحمن فی الفرقان قوله تبارک و تعالی و لقد ارسلنا موسی
بآیاتنا و سلطان مبین الی فرعون و هامان و قارون فقالوا ساحر کذّاب فلمّا
جآءهم بالحقّ من عندنا قالوا اقتلوا ابنآء الّذین آمنوا معه و استحیوا نسائهم
و ما کید الکافرین الّا فی ضلال و قال فرعون ذرونی اقتل موسى و لیدع ربّه انّی
اخاف ان یبدّل دینکم او ان یظهر فی الأرض الفساد و قال موسی انّی عذت بربّی و
ربّکم من کلّ متکبّر لا یؤمن بیوم الحساب

\end{otherlanguage}

Īn nisbat-i fasād hamān nisbat-i fasādīst ki az qabl farāʿina
bi-Ḥaḍrat-i Kalīm Allāh dāda-and. Iqraʾ mā anzalahu al-Raḥmān fī
al-Furqān, qawluhu tabāraka wa-taʿālā: Wa-laqad arsalnā Mūsā bi-āyātinā
wa-sulṭānin mubīn ilā Firʿawn wa-Hāmān wa-Qārūn fa-qālū sāḥirun
kadhdhāb. Fa-lammā jāʾahum bi-al-ḥaqq min ʿindinā qālū uqtulū abnāʾ
alladhīna āmanū maʿahu wa-istaḥyū nisāʾahum wa-mā kaydu al-kāfirīn illā
fī ḍalāl. Wa-qāla Firʿawn dharūnī aqtul Mūsā wa-l-yadʿu rabbahu innī
akhāfu an yubaddila dīnakum aw an yuẓhira fī al-arḍ al-fasād. Wa-qāla
Mūsā innī ʿudhtu bi-rabbī wa-rabbikum min kulli mutakabbirin lā yuʾminu
bi-yawm al-ḥisāb.

This charge of spreading corruption is the very charge that the Pharaohs
of old leveled against Moses, who spoke with God. Read what the
All-Merciful revealed in the Qur'an. These are His words, blessed and
exalted is He: ``We sent Moses with Our signs and a clear authority to
Pharaoh, Haman, and Korah, but they said, `A sorcerer, a liar!' When he
brought them the truth from Us, they said, `Kill the sons of those who
believe with him, and spare their women.' Yet the scheming of the
unbelievers leads only astray. Pharaoh said, `Let me kill Moses, and let
him call upon his Lord. I fear that he will change your religion or
cause corruption to appear in the land.' Moses said, `I seek refuge in
my Lord and your Lord from every arrogant person who does not believe in
the Day of Reckoning.'\,''

\subsection{109}\label{section-1395}

\begin{otherlanguage}{arabic}

لازال مصلح عالم را مفسد دانسته‌اند و گفته‌اند آنچه را که کل شنیده‌اند هر
هنگام نیّر ظهور از افق سماء ارادهٔ الهی اشراق نمود جمعی بانکار و برخی
باعراض و حزبی بمفتریات تمسّک جستند و عباد را از شریعهٔ عنایت مالک ایجاد
محروم ساختند چنانچه حال نفوسی‌که این مظلوم را ندیده‌اند و معاشرت نکرده‌اند
گفته و میگویند آنچه را که آن جناب شنیده و میشنوند بگوئید ای قوم امروز
آفتاب بیان از افق سماء فضل مشرق و نور ظهور مکلّم طور امام ادیان ساطع و
لامع صدر و قلب و سمع و بصر را بکوثر بیان رحمن مقدّس و مطهّر سازید و بعد
توجّه نمائید لعمر اللّه از جمیع اشیاء نداء قد اتی الحقّ اصغا کنید طوبی
للمنصفین و هنیئاً للمقبلین

\end{otherlanguage}

Lā-zāl muṣliḥ-i ʿālam rā mufsid dānista-and va gufta-and ānchi rā ki
kull shinīda-and. Har hangām nayyir-i ẓuhūr az ufuq-i samā-yi irāda-yi
ilāhī ishrāq nimūd, jamʿī bi-inkār va barkhī bi-iʿrāḍ va ḥizbī
bi-muftarayāt tamassuk justand va ʿibād rā az sharīʿa-yi ʿināyat-i
Mālik-i ījād maḥrūm sākhtand. Chinānchi ḥāl nufūsī ki īn maẓlūm rā
nadīda-and va muʿāsharat nakarda-and gufta va mīgūyand ānchi rā ki ān
janāb shinīda va mīshinavand. Bigūʾīd ay qawm imrūz āftāb-i bayān az
ufuq-i samā-yi faḍl mushriq va nūr-i ẓuhūr-i Mukallim-i Ṭūr imām-i adyān
sāṭiʿ va lāmiʿ. Ṣadr va qalb va samʿ va baṣar rā bi-kawthar-i bayān-i
Raḥmān muqaddas va muṭahhar sāzīd va baʿd tavajjuh namāʾīd. La-ʿamru
Allāh az jamīʿ-i ashyā nidā-yi qad atā al-Ḥaqq iṣghā kunīd. Ṭūbā
li-al-munṣifīn wa-hanīʾan li-al-muqbilīn.

They have always regarded the Reformer of the world as a corrupter, and
have said what everyone has heard. Whenever the Luminary of Revelation
shone from the horizon of the heaven of the divine will, some resorted
to denial, others to rejection, and still others to slander, depriving
God's servants of the stream of grace of the Lord of creation. So it is
now: people who have neither seen this Wronged One nor associated with
Him have said, and continue to say, the things that you have heard and
continue to hear. Say: People! Today the sun of utterance shines from
the horizon of the heaven of grace, and the light of the revelation of
Him who spoke on Sinai shines and gleams before the religions. Sanctify
and purify your breasts, hearts, hearing, and sight with the heavenly
waters of the All-Merciful's utterance, and then turn toward Him. By the
life of God, you will hear from all things the call, ``The Truth has
come!'' Blessed are the fair-minded, and joy to those who turn toward
Him.

\subsection{110}\label{section-1396}

\begin{otherlanguage}{persian}

از جمله نسبت بسدرهٔ مبارکه گفته‌اند آنچه را که هر عالم بصیری و هر عارف
خبیری بر کذبش گواهى دهد البتّه آن جناب آیاتی که در بارهٔ حضرت کلیم نازل
شده ملاحظه و قرائت نموده‌اند قوله تبارک و تعالی قال أ لم نربّک فینا ولیداً
و لبثت فینا من عمرک سنین و فعلت فعلتک الّتی فعلت و انت من الکافرین قال
فعلتها اذاً و انا من الضّالّین ففررت منکم لمّا خفتکم فوهب لی ربّی حکماً و
جعلنی من المرسلین و در مقام دیگر قوله تبارک و تعالی و دخل المدینة علی
حین غفلة من اهلها فوجد فیها رجلین یقتتلان هذا من شیعته و هذا من عدوّه
فاستغاثه الّذی من شیعته علی الّذی من عدوّه فوکزه موسی فقضی علیه قال هذا من
عمل الشّیطان انّه عدوّ مضلّ مبین قال ربّ انّی ظلمت نفسی فاغفر لی فغفر له انّه
هو الغفور الرّحیم قال ربّ بما انعمت علیّ فلن اکون ظهیراً للمجرمین فأصبح فی
المدینة خائفاً یترقّب فاذا الّذی استنصره بالأمس یستصرخه قال له موسى انّک
لغویّ مبین فلمّا ان اراد ان یبطش بالّذی هو عدوّ لهما قال یا موسی أ ترید ان
تقتلنی کما قتلت نفساً بالأمس ان ترید الّا ان تکون جبّاراً فی الأرض و ما ترید
ان تکون من المصلحین حال سمع و بصر مقدّس و مطهّر لازم که بعدل و انصاف تمسّک
نماید حضرت کلیم هم بظلم و گمراهی اقرار نموده و همچنین بخوف و فرار و
عصیان و طلب مغفرت از حقّ جلّ جلاله کرده و بآن فائز گشت

\end{otherlanguage}

Az jumla nisbat bi-Sidra-yi mubāraka gufta-and ānchi rā ki har ʿālim-i
baṣīrī va har ʿārif-i khabīrī bar kidhibash gavāhī dihad. Albatta ān
janāb āyātī ki dar bāra-yi Ḥaḍrat-i Kalīm nāzil shuda mulāḥaẓa va
qirāʾat nimūda-and. Qawluhu tabāraka wa-taʿālā: Qāla a-lam nurabbika
fīnā walīdan wa-labithta fīnā min ʿumrika sinīn. Wa-faʿalta faʿlataka
allatī faʿalta wa-anta min al-kāfirīn. Qāla faʿaltuhā idhan wa-anā min
al-ḍāllīn. Fa-farartu minkum lammā khiftukum fa-wahaba lī rabbī ḥukman
wa-jaʿalanī min al-mursalīn. Va dar maqām-i dīgar qawluhu tabāraka
wa-taʿālā: Wa-dakhala al-madīnata ʿalā ḥīni ghaflatin min ahlihā
fa-wajada fīhā rajulayni yaqtatilān hādhā min shīʿatihi wa-hādhā min
ʿaduwwihi fa-istaghāthahu alladhī min shīʿatihi ʿalā alladhī min
ʿaduwwihi fa-wakazahu Mūsā fa-qaḍā ʿalayhi qāla hādhā min ʿamali
al-shayṭān innahu ʿaduwwun muḍillun mubīn. Qāla rabbi innī ẓalamtu nafsī
fa-ighfir lī fa-ghafara lahu innahu huwa al-Ghafūr al-Raḥīm. Qāla rabbi
bi-mā anʿamta ʿalayya fa-lan akūna ẓahīran li-al-mujrimīn. Fa-aṣbaḥa fī
al-madīnati khāʾifan yataraqqabu fa-idhā alladhī istanṣarahu bi-al-amsi
yastaṣrikhuhu qāla lahu Mūsā innaka la-ghawiyyun mubīn. Fa-lammā an
arāda an yabṭisha bi-alladhī huwa ʿaduwwun lahumā qāla yā Mūsā a-turīdu
an taqtulanī kamā qatalta nafsan bi-al-amsi in turīdu illā an takūna
jabbāran fī al-arḍ wa-mā turīdu an takūna min al-muṣliḥīn. Ḥāl samʿ va
baṣar-i muqaddas va muṭahhar lāzim ki bi-ʿadl va inṣāf tamassuk namāyad.
Ḥaḍrat-i Kalīm ham bi-ẓulm va gumrāhī iqrār nimūda va hamchunīn bi-khawf
va firār va ʿiṣyān va ṭalab-i maghfirat az Ḥaqq jalla jalāluhu karda va
bi-ān fāʾiz gasht.

Among other things, they have said of the blessed Tree what every
discerning scholar and every well-informed knower would attest to be
false. You have surely examined and read the verses revealed concerning
Moses, who spoke with God. These are His words, blessed and exalted is
He: ``He said, `Did we not raise you among us as a child, and did you
not spend years of your life among us? And you committed that deed of
yours, and you were one of the ungrateful.' He said, `I did it then,
when I was among those who had gone astray. So I fled from you when I
feared you; then my Lord granted me wisdom and made me one of the
messengers.'\,'' And elsewhere these are His words, blessed and exalted
is He: ``He entered the city at a time when its people were unaware, and
found two men fighting there, one of his own people and the other of his
enemies. The one of his own people appealed to him for help against the
one of his enemies. Moses struck him with his fist and killed him. He
said, `This is the work of Satan. He is indeed an enemy who clearly
leads astray.' He said, `My Lord, I have wronged myself, so forgive me.'
And He forgave him. He is indeed the Ever-Forgiving, the Merciful. He
said, `My Lord, because of the favor You have shown me, I will never be
a supporter of the guilty.' The next morning he was in the city, fearful
and watchful, when the man who had sought his help the day before cried
out to him for help again. Moses said to him, `You are clearly in
error.' Then, when he was about to seize the man who was an enemy to
them both, he said, `Moses! Do you intend to kill me as you killed
someone yesterday? You only want to be a tyrant in the land; you do not
want to be one of those who set things right.'\,'' Now sanctified and
purified hearing and sight are needed, that one may hold fast to justice
and fairness. Moses himself acknowledged wrongdoing and error, as well
as fear, flight, and disobedience, and sought forgiveness from God,
exalted is His glory, and received it.

\subsection{111}\label{section-1397}

\begin{otherlanguage}{persian}

یا شیخ حقّ جلّ جلاله هر حین بمظاهر نفسش ظاهر با علم یفعل ما یشآء و یحکم ما
یرید آمده لیس لأحد ان یقول لم و بم و من قال انّه اعرض عن اللّه ربّ الأرباب
در ایّام ظهور این امور ظاهر و موجود چنانچه در بارهٔ این مظلوم هم گفته‌اند
آنچه را که مقرّبین و مخلصین بر کذبش گواهی داده و میدهند لعمر اللّه این ذیل
مقدّس بوده و هست اگرچه حال جمعی اراده نموده‌اند آن را بمفتریات کذبهٔ
نالائقه بیالایند ولکنّ اللّه یعلم و هم لا یعلمون نفسی که بقدرت و قوّت الهی
مقابل جمیع احزاب عالم قیام نمود و کل را بافق اعلی دعوت فرمود او را انکار
نموده‌اند و نفوسی که لازال خلف استار و احجاب ساکن و بحفظ خود مشغول بآن
نفوس تمسّک جسته‌اند حال هم جمعی بکذب و افترا مشغول قصدی جز ادخال شبهات در
افئده و قلوب نداشته و ندارند هر نفسی از مدینهٔ کبیره باین ارض توجّه
مینماید فوراً بسیّالهٔ برقیّه خبر میدهند که وجهی سرقت نموده و به عکّا رفته
شخصی عالم و کامل و فاضل در آخر عمر بارض مقدّسه توجّه نموده و مقصودش اعتکاف
بوده در بارهٔ او نوشته‌اند آنچه را که زفرات مخلصین و مقرّبین مرتفع

\end{otherlanguage}

Yā Shaykh, Ḥaqq jalla jalāluhu har ḥīn bi-maẓāhir-i nafsash ẓāhir, bā
ʿalam-i yafʿalu mā yashāʾ va yaḥkumu mā yurīd āmada. Laysa li-aḥadin an
yaqūla lima wa-bima wa-man qāla innahu aʿraḍa ʿan Allāh Rabb al-arbāb.
Dar ayyām-i ẓuhūr īn umūr ẓāhir va mawjūd, chinānchi dar bāra-yi īn
maẓlūm ham gufta-and ānchi rā ki muqarrabīn va mukhliṣīn bar kidhibash
gavāhī dāda va mīdihand. La-ʿamru Allāh īn dhayl muqaddas būda va hast
agarchi ḥāl jamʿī irāda nimūda-and ān rā bi-muftarayāt-i kadhiba-yi
nālāʾiqa biyālāyand, wa-lākinna Allāha yaʿlamu wa-hum lā yaʿlamūn. Nafsī
ki bi-qudrat va quvvat-i ilāhī muqābil-i jamīʿ-i aḥzāb-i ʿālam qiyām
nimūd va kull rā bi-Ufuq-i Aʿlā daʿvat farmūd ū rā inkār nimūda-and va
nufūsī ki lā-zāl khalf-i astār va aḥjāb sākin va bi-ḥifẓ-i khud mashghūl
bi-ān nufūs tamassuk justa-and. Ḥāl ham jamʿī bi-kidhb va iftirā
mashghūl qaṣdī juz idkhāl-i shubahāt dar afʾida va qulūb nadāshta va
nadārand. Har nafsī az madīna-yi kabīra bi-īn arḍ tavajjuh mīnamāyad
fawran bi-sayyāla-yi barqiyya khabar mīdihand ki vajhī sirqat nimūda va
bi-ʿAkkā rafta. Shakhṣī ʿālim va kāmil va fāḍil dar ākhir-i ʿumr
bi-Arḍ-i Muqaddasa tavajjuh nimūda va maqṣūdash iʿtikāf būda dar bāra-yi
ū nivishta-and ānchi rā ki zafarāt-i mukhliṣīn va muqarrabīn murtafiʿ.

Shaykh! God, exalted is His glory, has at every time appeared through
the Manifestations of His own Self, coming with the banner of ``He does
whatever He wills'' and ``He ordains whatever He desires.'' No one has
the right to say ``Why?'' or ``For what reason?'' Whoever says this has
turned away from God, the Lord of lords. Such things appear and occur in
the days of Revelation, just as they have said of this Wronged One
things whose falsehood those who are near to God and sincere in devotion
have attested and continue to attest. By the life of God! This robe has
always been, and remains, undefiled, although some now intend to stain
it with unseemly, lying slanders. But God knows, and they do not know.
They have denied the One who, through divine might and power, stood
against all the factions of the world and summoned everyone to the Most
Exalted Horizon, and have clung instead to people who have always dwelt
behind curtains and veils, occupied with protecting themselves. Even now
some are busy with lies and slander; their sole intention has been, and
remains, to sow doubts in hearts and minds. Whenever anyone sets out
from the Great City toward this land, they immediately send word by
telegraph that he has stolen money and gone to Acre. A learned,
accomplished, and erudite man came to the Holy Land at the end of his
life, intending to live in seclusion. They wrote things about him that
have raised the sighs of the sincere and those who are near to God.

\subsection{112}\label{section-1398}

\begin{otherlanguage}{persian}

این مظلوم را مرحوم مغفور حضرت مشیرالدّوله میرزا حسین خان غفر اللّه له
شناخته بود و البتّه نزد اولیای دولت تفصیل ورود این مظلوم را در آستانه و
اقوال و اعمالش را ذکر فرموده یوم ورود میهمان‌دار دولت حاضر و ما را بمحلّی
که مأمور بود برده فی‌الحقیقه کمال محبّت و عنایت از جانب دولت نسبت باین
مظلومان ظاهر و مشهود یوم دیگر شاهزاده شجاع‌الدّوله و میرزا صفا بنیابت
مرحوم مغفور مشیرالدّوله وزیر مختار تشریف آوردند و همچنین بعضی از وزرای
دولت علیّه از جمله مرحوم کمال پاشا و بعضى دیگر و این مظلوم متوکّلاً علی
اللّه من غیر ذکر حاجت و مطلبی چهار شهر در آن ارض بوده و اعمالش نزد کل
معلوم و مشهود لا ینکرها الّا کلّ مبغض کذّاب من عرف اللّه لم یعرف دونه دوست
نداشته و نداریم امثال این امور را ذکر نمائیم

\end{otherlanguage}

Īn maẓlūm rā marḥūm-i maghfūr Ḥaḍrat-i Mushīr al-Dawla Mīrzā Ḥusayn
Khān, ghafara Allāhu lahu, shinākhta būd va albatta nazd-i awliyā-yi
dawlat tafṣīl-i vurūd-i īn maẓlūm rā dar Āstāna va aqvāl va aʿmālash rā
dhikr farmūda. Yawm-i vurūd mihmān-dār-i dawlat ḥāḍir va mā rā
bi-maḥallī ki maʾmūr būd burda. Fī-al-ḥaqīqa kamāl-i maḥabbat va ʿināyat
az jānib-i dawlat nisbat bi-īn maẓlūmān ẓāhir va mashhūd. Yawm-i dīgar
Shāhzāda Shujāʿ al-Dawla va Mīrzā Ṣafā bi-niyābat-i marḥūm-i maghfūr
Mushīr al-Dawla vazīr-i mukhtār tashrīf āvardand va hamchunīn baʿḍī az
vuzarā-yi Dawlat-i ʿAliyya az jumla marḥūm Kamāl Pāshā va baʿḍī dīgar.
Va īn maẓlūm mutawakkilan ʿalā Allāh min ghayr-i dhikr-i ḥājat va
maṭlabī chahār shahr dar ān arḍ būda va aʿmālash nazd-i kull maʿlūm va
mashhūd. Lā yunkiruhā illā kullu mubghiḍin kadhdhāb. Man ʿarafa Allāha
lam yaʿrif dūnahu. Dūst nadāshta va nadārīm amthāl-i īn umūr rā dhikr
namāʾīm.

The late Mushir al-Dawlih, Mirza Husayn Khan, may God forgive him, knew
this Wronged One, and surely related to the authorities of the state the
details of this Wronged One's arrival in Constantinople, and His words
and deeds. On the day of Our arrival, the government's host was present
and took Us to the place to which he had been directed. Indeed, the
utmost kindness and consideration on the government's part toward these
wronged ones were plain and evident. The next day Prince Shuja al-Dawlih
and Mirza Safa came on behalf of the late Mushir al-Dawlih, the minister
plenipotentiary, as did some ministers of the Ottoman government,
including the late Kamal Pasha and certain others. This Wronged One
remained in that land for four months, trusting in God, without
mentioning any need or request, and His deeds were known and evident to
all. None deny them except malicious liars. Whoever has known God has
known no one besides Him. We have never wished, nor do We now wish, to
mention such matters.

\subsection{113}\label{section-1399}

\begin{otherlanguage}{persian}

بعضی از بزرگان ایران هر هنگام وارد آن مدینه شده‌اند درب خانه‌ها لأجل شهریّه
و انعام کمال جدّ و جهد را مبذول داشته‌اند و این مظلوم اگر سبب اعزاز نبوده
علّت ذلّت هم نشده و این عمل حضرت مرحوم مغفور اعلی اللّه مقامه نظر بدوستی
این مظلوم نبوده بلکه نظر بمقتضیات حکمت و خدمتی که سرّاً در نظر داشته‌اند
بوده شهادت میدهم که در خدمت دولت امین بوده بشأنی که خیانت را در عرصه‌اش
راهی و مقامی نبوده و سبب ورود این مظلومان در سجن اعظم هم او بوده ولکن
چون در عمل خود صادق بود لائق ذکر خیر است این مظلوم لازال همّش و قصدش
ارتفاع و ارتقاء دولت و ملّت بوده نه ارتفاع مقام خود حال جمعی جمعی را جمع
نموده‌اند و بر هتک حرمت این مظلوم قیام کرده‌اند ولکنّ المظلوم یسأل اللّه
تبارک و تعالی ان یؤیّدهم على الرّجوع و یوفّقهم علی تدارک ما فات عنهم و
الانابة لدی باب عطائه انّه هو الغفور الرّحیم

\end{otherlanguage}

Baʿḍī az buzurgān-i Īrān har hangām vārid-i ān madīna shuda-and darb-i
khāna-hā li-ajl-i shahriyya va inʿām kamāl-i jidd va jahd rā mabdhūl
dāshta-and va īn maẓlūm agar sabab-i iʿzāz nabūda ʿillat-i dhillat ham
nashuda. Va īn ʿamal-i Ḥaḍrat-i marḥūm-i maghfūr, aʿlā Allāhu maqāmahu,
naẓar bi-dūstī-yi īn maẓlūm nabūda balki naẓar bi-muqtaḍayāt-i ḥikmat va
khidmatī ki sirran dar naẓar dāshta-and būda. Shahādat mīdiham ki dar
khidmat-i dawlat amīn būda bi-shaʾnī ki khiyānat rā dar ʿarṣa-ash rāhī
va maqāmī nabūda va sabab-i vurūd-i īn maẓlūmān dar Sijn-i Aʿẓam ham ū
būda va lākin chūn dar ʿamal-i khud ṣādiq būd lāyiq-i dhikr-i khayr ast.
Īn maẓlūm lā-zāl hammash va qaṣdash irtifāʿ va irtiqā-yi dawlat va
millat būda na irtifāʿ-i maqām-i khud. Ḥāl jamʿī jamʿī rā jamʿ
nimūda-and va bar hatk-i ḥurmat-i īn maẓlūm qiyām karda-and. Wa-lākinna
al-Maẓlūm yasʾalu Allāha tabāraka wa-taʿālā an yuʾayyidahum ʿalā
al-rujūʿ wa-yuwaffiqahum ʿalā tadāruki mā fāta ʿanhum wa-al-inābati ladā
bābi ʿaṭāʾihi innahu huwa al-Ghafūr al-Raḥīm.

Whenever some of the notables of Iran have entered that city, they have
made every effort at people's doors to obtain stipends and gifts. This
Wronged One, if He has not brought honor, has at least not brought
humiliation. This action of that departed gentleman, may God exalt his
station, was not prompted by friendship for this Wronged One, but by
considerations of prudence and a service he privately had in mind. I
testify that he was trustworthy in the service of the state, to such a
degree that treachery had neither access to nor a place in his conduct.
He was also the cause of these wronged ones entering the Most Great
Prison; nevertheless, because he was sincere in what he did, he deserves
to be remembered favorably. This Wronged One's concern and purpose have
always been the advancement and progress of the government and the
people, not the elevation of His own station. Now some have gathered
others around them and arisen to dishonor this Wronged One. Yet the
Wronged One asks God, blessed and exalted is He, to help them return,
enable them to make up for what they have missed, and turn in repentance
to the door of His bounty. He is indeed the Ever-Forgiving, the
Merciful.

\subsection{114}\label{section-1400}

\begin{otherlanguage}{arabic}

یا شیخ انّ قلمی ینوح لنفسی و اللّوح یبکی بما ورد علیّ من الّذی حفظناه فی
سنین متوالیات و کان ان یخدم امام وجهی فی اللّیالی و الأیّام الی ان اغواه
احد خدّامی الّذی سمّی بسیّد محمّد یشهد بذلک عباد موقنون الّذین هاجروا معی من
الزّورآء الی ان وردنا فی هذا السّجن العظیم و ورد منهما علیّ ما صاح به کلّ
عالم و ناح به کلّ عارف و ذرفت دموع المنصفین

\end{otherlanguage}

Yā Shaykh, inna qalamī yanūḥu li-nafsī wa-al-Lawḥ yabkī bi-mā warada
ʿalayya min alladhī ḥafiẓnāhu fī sinīna mutawāliyāt wa-kāna an yakhduma
imāma wajhī fī al-layālī wa-al-ayyām ilā an aghwāhu aḥadu khuddāmī
alladhī summiya bi-Sayyid Muḥammad. Yashhadu bi-dhālika ʿibādun mūqinūn
alladhīna hājarū maʿī min al-Zawrāʾ ilā an waradnā fī hādhā al-sijn
al-ʿaẓīm. Wa-warada minhumā ʿalayya mā ṣāḥa bihi kullu ʿālim wa-nāḥa
bihi kullu ʿārif wa-dharafat dumūʿu al-munṣifīn.

Shaykh! My Pen laments for Me, and the Tablet weeps over what I have
suffered at the hands of one whom We protected for successive years, who
served in My presence night and day, until one of My attendants, named
Siyyid Muhammad, led him astray. The servants endowed with certitude who
migrated with Me from Baghdad until We arrived in this mighty prison
bear witness to this. From these two I suffered what made every learned
person cry out, every knower lament, and the tears of the fair-minded
flow.

\subsection{115}\label{section-1401}

\begin{otherlanguage}{arabic}

نسأل اللّه ان یؤیّد الغافلین علی العدل و الانصاف و یعرّفهم ما غفلوا عنه انّه
هو الفضّال الکریم ای ربّ لا تمنع عبادک عن باب فضلک و لا تطردهم عن بساط
قربک ایّدهم علی کشف سبحات الجلال و خرق حجبات الأوهام و الآمال انّک انت
الغنیّ المتعال لا اله الّا انت العزیز الفضّال

\end{otherlanguage}

Nasʾalu Allāha an yuʾayyida al-ghāfilīn ʿalā al-ʿadl wa-al-inṣāf
wa-yuʿarrifahum mā ghafalū ʿanhu innahu huwa al-Faḍḍāl al-Karīm. Ay
rabb, lā tamnaʿ ʿibādaka ʿan bābi faḍlika wa-lā taṭrudhum ʿan bisāṭi
qurbika. Ayyidhum ʿalā kashfi subuḥāti al-jalāl wa-kharqi ḥujubāti
al-awhām wa-al-āmāl. Innaka anta al-Ghanī al-Mutaʿāl. Lā ilāha illā anta
al-ʿAzīz al-Faḍḍāl.

We ask God to assist the heedless to act with justice and fairness, and
to make known to them what they have overlooked. He is indeed the Most
Bountiful, the Generous. My Lord! Do not bar Your servants from the door
of Your grace, nor drive them from the court of Your presence. Help them
to remove the veils of glory and tear through the veils of imaginings
and desires. You are indeed the Self-Sufficient, the Most Exalted. There
is no God but You, the Almighty, the Most Bountiful.

\subsection{116}\label{section-1402}

\begin{otherlanguage}{arabic}

قسم بآفتاب برهان که از افق سماء ایقان اشراق نموده این مظلوم در لیالی و
ایّام بتهذیب نفوس مشغول تا آنکه نور دانائی بر تاریکی نادانی غلبه نمود

\end{otherlanguage}

Qasam bi-āftāb-i burhān ki az ufuq-i samā-yi īqān ishrāq nimūda īn
maẓlūm dar layālī va ayyām bi-tahdhīb-i nufūs mashghūl tā ānki nūr-i
dānāʾī bar tārīkī-yi nādānī ghalaba nimūd.

I swear by the sun of proof that has shone from the horizon of the
heaven of certitude: this Wronged One has been occupied night and day
with refining souls, until the light of knowledge overcame the darkness
of ignorance.

\subsection{117}\label{section-1403}

\begin{otherlanguage}{persian}

یا شیخ مکرّر ذکر شده و میشود چهل سنه بعنایت الهی و ارادهٔ قویّهٔ نافذهٔ ربّانی
حضرت سلطان ایّده اللّه را نصرت نمودیم نصرتی که نزد مظاهر عدل و انصاف ثابت
و محقّق است و لا ینکرها الّا کلّ معتد اثیم و کلّ مبغض مریب عجب آنکه وزراء
دولت و امناء ملّت الی حین باین خدمت ظاهر مبین ملتفت نشده‌اند و یا شده‌اند
نظر بحکمت ذکر نفرموده‌اند قبل از اربعین هر سنه مابین عباد مجادله و محاربه
ظاهر و قائم و بعد بجنود حکمت و بیان و نصیحت و عرفان کل بحبل متین صبر و
ذیل منیر اصطبار تمسّک جستند و تشبّث نمودند بشأنی که آنچه بر این حزب مظلوم
وارد شد تحمّل کردند و بحقّ گذاشتند مع آنکه در مازندران و رشت جمع کثیری را
ببدترین عذاب معذّب نمودند از جمله حضرت حاجی نصیر که فی‌الحقیقه نوری بود
مشرق از افق سماء تسلیم بعد از شهادت چشم او را کندند و دماغش را بریدند و
ظلم بمقامی رسید که اهل ممالک خارجه گریستند و نوحه نمودند و در باطن از
برای عیال و صغار بعضی در اطراف ممالک اعانت نمودند

\end{otherlanguage}

Yā Shaykh, mukarrar dhikr shuda va mīshavad chihil sana bi-ʿināyat-i
ilāhī va irāda-yi qaviyya-yi nāfidha-yi rabbānī Ḥaḍrat-i Sulṭān ayyadahu
Allāh rā nuṣrat nimūdīm nuṣratī ki nazd-i maẓāhir-i ʿadl va inṣāf thābit
va muḥaqqaq ast, wa-lā yunkiruhā illā kullu muʿtadin athīm wa-kullu
mubghiḍin murīb. ʿAjab ānki vuzarā-yi dawlat va umanā-yi millat ilā ḥīn
bi-īn khidmat-i ẓāhir-i mubīn multafit nashuda-and va yā shuda-and naẓar
bi-ḥikmat dhikr nafarmūda-and. Qabl az arbaʿīn har sana mābayn-i ʿibād
mujādala va muḥāraba ẓāhir va qāʾim va baʿd bi-junūd-i ḥikmat va bayān
va naṣīḥat va ʿirfān kull bi-ḥabl-i matīn-i ṣabr va dhayl-i munīr-i
iṣṭibār tamassuk justand va tashabbuth nimūdand bi-shaʾnī ki ānchi bar
īn ḥizb-i maẓlūm vārid shud taḥammul kardand va bi-Ḥaqq gudhāshtand.
Maʿa ānki dar Māzandarān va Rasht jamʿ-i kathīrī rā bi-badtarīn ʿadhāb
muʿadhdhab nimūdand az jumla Ḥaḍrat-i Ḥājī Naṣīr ki fī-al-ḥaqīqa nūrī
būd mushriq az ufuq-i samā-yi taslīm. Baʿd az shahādat chashm-i ū rā
kandand va damāghash rā burīdand va ẓulm bi-maqāmī rasīd ki ahl-i
mamālik-i khārija girīstand va nawḥa nimūdand va dar bāṭin az barā-yi
ʿiyāl va ṣighār baʿḍī dar aṭrāf-i mamālik iʿānat nimūdand.

Shaykh! It has repeatedly been stated, and is stated again, that for
forty years, through divine grace and the mighty, all-pervading will of
God, We have aided His Majesty the sovereign, may God assist him---an
aid established and confirmed in the eyes of those who embody justice
and fairness. None deny it except sinful aggressors and malicious
doubters. It is astonishing that the ministers of the state and the
trustees of the people have not yet recognized this clear and evident
service, or, if they have, have refrained from mentioning it out of
prudence. Before these forty years, strife and warfare arose among the
people every year. Afterward, through the forces of wisdom, utterance,
counsel, and spiritual knowledge, all took hold of the strong cord of
patience and clung to the luminous robe of forbearance, to such a degree
that they endured whatever befell this wronged community and left it to
God. This was so even though in Mazandaran and Rasht many were subjected
to the worst tortures, among them Haji Nasir, who was truly a light
shining from the horizon of the heaven of submission. After his
martyrdom they gouged out his eye and cut off his nose. The cruelty
reached such a degree that people in foreign countries wept and
lamented, and some in various countries secretly provided assistance for
the families and young children.

\subsection{118}\label{section-1404}

\begin{otherlanguage}{persian}

یا شیخ قلم حیا میکند از ذکر آنچه واقع شده و در ارض صاد نار ظلم مشتعل
بشأنی که هر منصفی نوحه نمود لعمرک از مدائن علم و معرفت نحیب بکا مرتفع
بقسمی که اکباد اهل بر و تقوی محترق نورین نیّرین حسنین در آن ارض رایگان
جان فدا نمودند دولت و ثروت و عزّت ایشان را منع ننمود اللّه یعلم ما ورد
علیهما و القوم اکثرهم لا یعلمون

\end{otherlanguage}

Yā Shaykh, qalam ḥayā mīkunad az dhikr-i ānchih vāqiʿ shudih va dar
arḍ-i Ṣād nār-i ẓulm mushtaʿil bi-shaʾnī kih har munṣifī nawḥih nimūd.
La-ʿamruka, az madāʾin-i ʿilm va maʿrifat naḥīb-i bukā murtafiʿ bi-qismī
kih akbād-i ahl-i birr va taqvā muḥtariq. Nūrayn-i nayyirayn-i Ḥasanayn
dar ān arḍ rāygān jān fidā nimūdand. Dawlat va tharvat va ʿizzat īshān
rā manʿ nanimūd. Allāhu yaʿlamu mā warada ʿalayhimā wa-l-qawmu
aktharuhum lā yaʿlamūn.

Shaykh! The Pen is ashamed to recount what has happened. In the land of
Sad, the fire of oppression blazed so fiercely that every fair-minded
person lamented. By your life! From the cities of knowledge and
understanding arose such cries of weeping that the hearts of the
righteous and devout were consumed. The two radiant lights, the two
Hasans, freely sacrificed their lives in that land. Position, wealth,
and honor did not hold them back. God knows what befell them, but most
of the people do not know.

\subsection{119}\label{section-1405}

\begin{otherlanguage}{arabic}

و قبل از ایشان جناب کاظم و من معه و در آخر حضرت اشرف کل شربت شهادت را
بکمال شوق و اشتیاق نوشیدند و برفیق اعلی شتافتند و همچنین در عهد سردار
عزیز خان جناب عارف باللّه میرزا مصطفی و من معه را اخذ نمودند و برفیق اعلی
و افق ابهی فرستادند باری در هر بلدی آثار ظلم ظاهر و مشهود ظلمی که شبه و
مثل نداشته مع‌ذلک نفسی بر دفاع قیام ننمود در حضرت بدیع که حامل لوح حضرت
سلطان بوده تفکّر نما که چگونه جان داد آن فارس مضمار انقطاع اکلیل ثمین
حیات را نثار دوست یکتا نمود

\end{otherlanguage}

Va qabl az īshān Jināb-i Kāẓim va man maʿahu va dar ākhir Ḥaḍrat-i
Ashraf, kull sharbat-i shahādat rā bi-kamāl-i shawq va ishtiyāq
nūshīdand va bi-Rafīq-i Aʿlā shitāftand. Va hamchunīn dar ʿahd-i Sardār
ʿAzīz Khān Jināb-i ʿārif bi-llāh Mīrzā Muṣṭafā va man maʿahu rā akhdh
nimūdand va bi-Rafīq-i Aʿlā va Ufuq-i Abhā firistādand. Bārī dar har
baladī āthār-i ẓulm ẓāhir va mashhūd, ẓulmī kih shibh va mithl
nadāshtih. Maʿa-dhālik nafsī bar difāʿ qiyām nanimūd. Dar Ḥaḍrat-i Badīʿ
kih ḥāmil-i lawḥ-i Ḥaḍrat-i Sulṭān būdih tafakkur namā kih chigūnih jān
dād. Ān fāris-i miḍmār-i inqiṭāʿ iklīl-i thamīn-i ḥayāt rā nithār-i
Dūst-i Yiktā nimūd.

Before them, Kazim and his companions, and lastly Ashraf, all drank the
cup of martyrdom with the utmost eagerness and longing, and hastened to
the Most Exalted Companion. Likewise, during the tenure of Sardar Aziz
Khan, they seized Mirza Mustafa, who knew God, and his companions, and
sent them to the Most Exalted Companion and the All-Glorious Horizon. In
short, in every city the signs of oppression were clear and
evident---oppression without parallel or equal. Yet not a soul rose to
defend himself. Consider Badi, who carried the Tablet to His Majesty the
sovereign: how he gave his life! That horseman in the arena of
detachment offered the precious crown of life to the one Beloved.

\subsection{120}\label{section-1406}

\begin{otherlanguage}{persian}

یا شیخ اگر این امور انکار شود کدام امر لائق اقرار است بیّن لوجه اللّه و لا
تکن من الصّامتین حضرت نجفعلی را اخذ نمودند و با کمال جذب و شوق قصد مقام
شهادت نمود و باین کلمه ناطق ما بها و خون‌بها را یافتیم این بگفت و جان داد
در اشراق و تجلّی نیّر انقطاع که از شطر اعلای قلب ملّا علیجان اشراق نمود نظر
نمائید نفحات کلمهٔ علیا و اقتدار قلم اعلی بشأنی جذبش کرد که میدان شهادت و
ایوان عشرت نزدش یکسان بل الأولی اولی در جناب ابابصیر و سیّد اشرف زنجانی
تفکّر نمائید امّ اشرف را حاضر نمودند که ابنش را نصیحت نماید ترغیب نمود الی
ان فاز بالشّهادة الکبری

\end{otherlanguage}

Yā Shaykh, agar īn umūr inkār shavad kudām amr lāyiq-i iqrār ast? Bayyin
li-wajhi Allāh wa-lā takun mina al-ṣāmitīn. Ḥaḍrat-i NajafʿAlī rā akhdh
nimūdand va bā kamāl-i jadhb va shawq qaṣd-i maqām-i shahādat nimūd va
bi-īn kalimih nāṭiq: mā bahā va khūn-bahā rā yāftīm. Īn biguft va jān
dād. Dar ishrāq va tajallī-yi nayyir-i inqiṭāʿ kih az shaṭr-i aʿlā-yi
qalb-i Mullā ʿAlījān ishrāq nimūd naẓar namāyīd. Nafaḥāt-i kalimih-yi
ʿulyā va iqtidār-i Qalam-i Aʿlā bi-shaʾnī jadhbash kard kih maydān-i
shahādat va ayvān-i ʿishrat nazdash yiksān, bal al-ūlā awlā. Dar Jināb-i
Abā Baṣīr va Sayyid Ashraf-i Zanjānī tafakkur namāyīd. Umm-i Ashraf rā
ḥāḍir nimūdand kih ibnash rā naṣīḥat namāyad. Targhīb nimūd ilā an fāza
bi-l-shahādati al-kubrā.

Shaykh! If these things are denied, what is worthy of acknowledgment?
Speak plainly for the sake of God, and do not be among the silent. They
seized Najaf-Ali, and with the utmost rapture and eagerness he set out
for the place of martyrdom, uttering these words: ``We have found Baha
and the price of our blood.'' He said this and gave his life. Consider
the radiance and splendor of the sun of detachment that shone from the
loftiest reaches of Mulla Ali-Jan's heart. The breezes of the Most
Exalted Word and the power of the Most Exalted Pen so enraptured him
that the field of martyrdom and the hall of pleasure were alike to
him---indeed, the former was preferable. Reflect on Aba Basir and Sayyid
Ashraf of Zanjan. They brought Ashraf's mother to counsel her son. She
encouraged him until he attained the supreme martyrdom.

\subsection{121}\label{section-1407}

\begin{otherlanguage}{arabic}

یا شیخ این حزب از خلیج اسماء گذشته‌اند و بر شاطی بحر انقطاع خرگاه
برافراشته‌اند ایشان صد‌هزار جان رایگان نثار نمایند و به ما اراده الأعدآء
تکلّم ننمایند بارادة اللّه متمسّکند و از ما عند القوم فارغ و آزاد سر دادند
و کلمهٔ نالائقه نگفتند تفکّر نمائید گویا از بحر انقطاع نوشیده‌اند زندگی
دنیا ایشان را از شهادت در سبیل الهی منع ننمود

\end{otherlanguage}

Yā Shaykh, īn ḥizb az khalīj-i asmāʾ gudhashtih-and va bar shāṭī-yi
baḥr-i inqiṭāʿ khargāh bar-afrāshtih-and. Īshān ṣad-hazār jān rāygān
nithār namāyand va bi-mā arāda al-aʿdāʾ takallum nanamāyand. Bi-irādati
Allāh mutamassikand va az mā ʿinda al-qawm fārigh va āzād. Sar dādand va
kalimih-yi nālāyiqih naguftand. Tafakkur namāyīd. Gūyā az baḥr-i inqiṭāʿ
nūshīdih-and. Zindigī-yi dunyā īshān rā az shahādat dar sabīl-i ilāhī
manʿ nanimūd.

Shaykh! This community has passed beyond the gulf of names and pitched
its tent on the shore of the ocean of detachment. They would freely
offer a hundred thousand lives rather than speak as their enemies
wished. They hold fast to the will of God and are free and detached from
what the people possess. They gave their heads but uttered no unworthy
word. Reflect: it is as though they have drunk from the ocean of
detachment. The life of this world did not hold them back from martyrdom
in the path of God.

\subsection{122}\label{section-1408}

\begin{otherlanguage}{persian}

در مازندران جمع کثیری از عباد اللّه را تمام نمودند حاکم بمفتریات جمعی را
تاراج کرد از جمله ذکر نمود حضرات اسلحه جمع کرده‌اند و بعد از تفحّص ملاحظه
شد یک لولهٔ تفنگ من غیر اسباب بوده سبحان‌اللّه این حزب بسلاح محتاج نه چه که
کمر همّت لأجل اصلاح عالم بسته‌اند جندشان اعمال طیّبه سلاحشان اخلاق مرضیّه و
سردارشان تقوی اللّه طوبی لمن انصف لعمر اللّه این حزب از صبر و سکون و تسلیم
و رضا مظاهر عدل شده‌اند و در اصطبار بمقامی رسیده‌اند که کشته شده‌اند و
نکشته‌اند مع آنکه بر مظلومهای ارض وارد شد آنچه که تاریخ عالم شبه آن را
ذکر ننموده و چشم امم مثلش را ندیده آیا سبب آنکه این بلایای عظیمه را قبول
نموده‌اند و در دفع آن دست درنیاورده‌اند چه بوده و علّت تسلیم و سکون چه سبب
منع قلم اعلی در صباح و مساء و اخذ زمام امور بقدرت و قوّت مولی الوری

\end{otherlanguage}

Dar Māzandarān jamʿ-i kathīrī az ʿibād-i Allāh rā tamām nimūdand. Ḥākim
bi-muftarayāt jamʿī rā tārāj kard. Az jumlih dhikr nimūd ḥaḍarāt asliḥih
jamʿ kardih-and, va baʿd az tafaḥḥuṣ mulāḥiẓih shud yik lūlih-yi tufang
min ghayri asbāb būdih. Subḥāna Allāh, īn ḥizb bi-silāḥ muḥtāj nah, chih
kih kamar-i himmat li-ajli iṣlāḥ-i ʿālam bastih-and. Jundishān aʿmāl-i
ṭayyibih, silāḥishān akhlāq-i marḍīyyih va sardārishān taqvā Allāh. Ṭūbā
li-man anṣafa. La-ʿamru Allāh, īn ḥizb az ṣabr va sukūn va taslīm va
riḍā maẓāhir-i ʿadl shudih-and va dar iṣṭibār bi-maqāmī rasīdih-and kih
kushtih shudih-and va nakushtih-and. Maʿa ānkih bar maẓlūmhā-yi arḍ
vārid shud ānchih kih tārīkh-i ʿālam shibh-i ān rā dhikr nanimūdih va
chashm-i umam mithlash rā nadīdih. Āyā sabab-i ānkih īn balāyā-yi
ʿaẓīmih rā qabūl nimūdih-and va dar dafʿ-i ān dast dar-nayāvardih-and
chih būdih va ʿillat-i taslīm va sukūn chih? Sabab, manʿ-i Qalam-i Aʿlā
dar ṣabāḥ va masāʾ va akhdh-i zimām-i umūr bi-qudrat va quvvat-i Mawlā
al-Varā.

In Mazandaran they put to death a great many of God's servants. The
governor plundered a number of them on the strength of false
accusations. Among other things, he stated that these people had
gathered weapons, but after a search it was found that there was only a
single gun barrel without its fittings. Glory be to God! This community
has no need of weapons, for they have resolved to reform the world.
Their armies are good deeds, their weapons commendable conduct, and
their commander the fear of God. Blessed is the one who judges fairly!
By the life of God! Through patience, calm, submission, and contentment,
this community has become an embodiment of justice. They have reached
such a degree of patient endurance that they have been killed but have
not killed, even though these wronged ones of the earth have suffered
what the history of the world has never recorded and the eyes of the
nations have never seen. Why have they accepted these tremendous
afflictions and raised no hand to ward them off? What accounts for their
submission and calm? The reason is the prohibition issued by the Most
Exalted Pen morning and evening, and the reins of affairs being held
through the power and might of the Lord of all people.

\subsection{123}\label{section-1409}

\begin{otherlanguage}{arabic}

در ابا‌بدیع تفکّر نمائید آن مظلوم را اخذ نمودند و بسب و لعن این مظلوم امر
کردند ولکن بعنایت الهی و رحمت ربّانی شهادت اختیار نمود و بآن فائز گشت ان
تعدّوا الشّهدآء فی سبیل اللّه لا تحصوهم در حضرت سیّد اسمعیل علیه سلام اللّه و
عنایته نظر نمائید که قبل از فجر باب بیت را با عمامهٔ خود جاروب نمود و بعد
در کنار شط مقبلاً الی البیت بدست خود جان نثار نمود

\end{otherlanguage}

Dar Abā-Badīʿ tafakkur namāyīd. Ān maẓlūm rā akhdh nimūdand va bi-sabb
va laʿn-i īn Maẓlūm amr kardand, va lākin bi-ʿināyat-i ilāhī va raḥmat-i
rabbānī shahādat ikhtiyār nimūd va bi-ān fāʾiz gasht. In taʿuddū
al-shuhadāʾa fī sabīli Allāh lā tuḥṣūhum. Dar Ḥaḍrat-i Sayyid Ismaʿīl
ʿalayhi salāmu Allāh wa-ʿināyatuhu naẓar namāyīd kih qabl az fajr bāb-i
bayt rā bā ʿamāmih-yi khud jārūb nimūd va baʿd dar kinār-i shaṭṭ
muqbilan ilā al-bayt bi-dast-i khud jān nithār nimūd.

Reflect on Aba-Badi. They seized that wronged man and ordered him to
revile and curse this Wronged One. Yet through divine favor and the
Lord's mercy he chose martyrdom and attained it. If you were to count
the martyrs in the path of God, you could not number them. Consider
Sayyid Ismail, upon whom be God's peace and favor, who before dawn swept
the doorway of the House with his turban, and afterward, on the
riverbank, facing the House, offered up his life with his own hand.

\subsection{124}\label{section-1410}

\begin{otherlanguage}{arabic}

در نفوذ کلمه نظر نما جمیع این نفوس را اوّل بسب و لعن امر نمودند ولکن
هیچکدام ارادهٔ خود را بر ارادة اللّه مقدّم نداشت

\end{otherlanguage}

Dar nufūdh-i kalimih naẓar namā. Jamīʿ-i īn nufūs rā avval bi-sabb va
laʿn amr nimūdand, va lākin hīchkudām irādih-yi khud rā bar irādati
Allāh muqaddam nadāsht.

Consider the penetrating power of the Word. All these souls were first
ordered to revile and curse, yet not one of them placed his own will
before the will of God.

\subsection{125}\label{section-1411}

\begin{otherlanguage}{persian}

یا شیخ از قبل حضرت ذبیح یک نفس بوده و حال این مظلوم ظاهر نمود از برای
شما آنچه را که علّت حیرت منصفین است انصف باللّه و قم على خدمة ربّک انّه
یجزیک جزآء لا تعادله کنوز الأرض و لا خزائن الملوک و السّلاطین توکّل فی کلّ
الأمور علی اللّه و فوّضها الیه انّه یعطیک اجراً کان فی الکتاب عظیماً و در این
دوروزهٔ عمر باعمالی مشغول شو که عرف رضا از آن متضوّع گردد و بطراز قبول
مزیّن شود حضرت بلال حبشی چون عملش مقبول افتاد سین او از شین عالم سبقت
گرفت امروز باید جمیع احزاب بنور اتّحاد و اتّفاق منوّر گردند باری کبر و غرور
بعضی از احزاب عالم دانائی را خراب کرده و بیت عدل و داد را ویران نموده

\end{otherlanguage}

Yā Shaykh, az qabl Ḥaḍrat-i Dhabīḥ yik nafs būdih va ḥāl īn Maẓlūm ẓāhir
nimūd az barāy-i shumā ānchih rā kih ʿillat-i ḥayrat-i munṣifīn ast.
Inṣaf bi-llāh wa-qum ʿalā khidmati rabbika innahu yajzīka jazāʾan lā
tuʿādiluhu kunūzu al-arḍi wa-lā khazāʾinu al-mulūki wa-l-salāṭīn.
Tawakkal fī kulli al-umūri ʿalā Allāh wa-fawwiḍhā ilayhi innahu yuʿṭīka
ajran kāna fī al-kitābi ʿaẓīman. Va dar īn du-rūzih-yi ʿumr bi-aʿmālī
mashghūl shaw kih ʿarf-i riḍā az ān mutaḍavviʿ gardad va bi-ṭirāz-i
qabūl muzayyan shavad. Ḥaḍrat-i Bilāl-i Ḥabashī chūn ʿamalash maqbūl
uftād sīn-i ū az shīn-i ʿālam sabqat girift. Imrūz bāyad jamīʿ-i aḥzāb
bi-nūr-i ittiḥād va ittifāq munavvar gardand. Bārī kibr va ghurūr-i
baʿḍī az aḥzāb ʿālam-i dānāʾī rā kharāb kardih va bayt-i ʿadl va dād rā
vīrān nimūdih.

Shaykh! In former times the sacrificial victim was but one soul. Now
this Wronged One has shown you what astonishes the fair-minded. Be fair,
for God's sake, and arise to serve your Lord. He will reward you with a
recompense that neither the treasures of the earth nor the treasuries of
kings and sovereigns can equal. In all affairs, place your trust in God
and commit them to Him. He will grant you a reward accounted great in
the Book. In these few days of life, occupy yourself with deeds that
give off the fragrance of His good pleasure and are adorned with the
ornament of His acceptance. Because the deeds of Bilal the Ethiopian
were accepted, his sin surpassed the shin of the world. Today all
communities must be illumined with the light of unity and concord. Yet
the pride and arrogance of some communities have laid waste the realm of
knowledge and destroyed the house of justice and equity.

\subsection{126}\label{section-1412}

\begin{otherlanguage}{persian}

یا شیخ وارد شد بر این مظلوم آنچه که شبه و مثل نداشته و کل را بکمال تسلیم
و رضا لأجل تهذیب نفوس و ارتفاع کلمة اللّه حمل نمودیم در ایّامی که در سجن
ارض میم بودیم ما را یک یوم بدست علما دادند دیگر معلومست که چه وارد شد و
اگر وقتی آن جناب در انبار حضرت سلطان وارد شوند از نائب و رئیس آن محل
بطلبند که آن دو زنجیر را که یکی به قره‌کهر و یکی به سلاسل معروفست بنمایند
قسم بنیّر عدل که چهار شهر این مظلوم در یکی از این دو معذّب و مغلول و حزنِی
ما یعقوب بثّ اقلّه و کلّ بلآء ایّوب بعض بلیّتی

\end{otherlanguage}

Yā Shaykh, vārid shud bar īn Maẓlūm ānchih kih shibh va mithl nadāshtih
va kull rā bi-kamāl-i taslīm va riḍā li-ajli tahdhīb-i nufūs va
irtifāʿ-i kalimati Allāh ḥaml nimūdīm. Dar ayyāmī kih dar sijn-i arḍ-i
Mīm būdīm mā rā yik yawm bi-dast-i ʿulamā dādand. Dīgar maʿlūmast kih
chih vārid shud. Va agar vaqtī ān jināb dar anbār-i Ḥaḍrat-i Sulṭān
vārid shavand az nāʾib va raʾīs-i ān maḥall biṭalaband kih ān du zanjīr
rā kih yikī bi-Qarah-Kahr va yikī bi-Salāsil maʿrūfast binamāyand. Qasam
bi-Nayyir-i ʿAdl kih chahār shahr īn Maẓlūm dar yikī az īn du muʿadhdhab
va maghlūl. Wa-ḥuznī mā Yaʿqūbu baththa aqallahu wa-kullu balāʾi Ayyūba
baʿḍu baliyyatī.

Shaykh! This Wronged One has suffered what has no parallel or equal, and
We have borne it all with complete submission and contentment, to refine
souls and exalt the Word of God. During the days when We were in prison
in the land of Mim, they handed Us over to the religious scholars for
one day. What then befell Us is plain enough. If you should ever enter
His Majesty the sovereign's storehouse, ask the deputy and the officer
in charge there to show you the two chains known as Qarah-Kahr and
Salasil. I swear by the Sun of Justice that for four months this Wronged
One was tormented and fettered in one of these two. The sorrow Jacob
poured out was but the smallest part of My sorrow, and all Job's
afflictions are but part of My affliction.

\subsection{127}\label{section-1413}

\begin{otherlanguage}{arabic}

و همچنین در شهادت حاجی محمّد رضا در مدینهٔ عشق تفکّر نمائید ظالمهای عالم بر
آن مظلوم وارد آوردند آنچه را که بعضی از ملل خارجه گریستند و نوحه نمودند
چه که از قرار مذکور و معلوم سی و دو زخم بر جسد مبارکش وارد مع‌ذلک احدی از
امر تجاوز ننمود و دست درنیاورد در هر حال حکم کتاب را بر ارادهٔ خود مقدّم
داشتند مع آنکه در آن مدینه جمعی از این حزب بوده و هستند

\end{otherlanguage}

Va hamchunīn dar shahādat-i Ḥājī Muḥammad Riḍā dar Madīnih-yi ʿIshq
tafakkur namāyīd. Ẓālimhā-yi ʿālam bar ān maẓlūm vārid āvardand ānchih
rā kih baʿḍī az milal-i khārijih girīstand va nawḥih nimūdand, chih kih
az qarār-i madhkūr va maʿlūm sī va du zakhm bar jasad-i mubārakash
vārid. Maʿa-dhālik aḥadī az amr tajāvuz nanimūd va dast dar-nayāvard.
Dar har ḥāl ḥukm-i kitāb rā bar irādih-yi khud muqaddam dāshtand, maʿa
ānkih dar ān madīnih jamʿī az īn ḥizb būdih va hastand.

Reflect also on the martyrdom of Haji Muhammad-Rida in the City of Love.
The oppressors of the world inflicted on that wronged man such cruelties
that some among the foreign nations wept and lamented, for, according to
what has been reported and established, his blessed body received
thirty-two wounds. Nevertheless, no one transgressed the command or
raised a hand. In every circumstance they placed the injunction of the
Book before their own will, although a number of this community lived,
and still live, in that city.

\subsection{128}\label{section-1414}

\begin{otherlanguage}{persian}

استدعا آنکه حضرت سلطان ایّده اللّه تبارک و تعالی در این امور بنفسه تفکّر
فرمایند و بعدل و انصاف حکم نمایند مع آنکه در اکثری از بلاد ایران در این
سنین اخیره این حزب کشته شدند و نکشتند مشاهده میشود بغضای افئدهٔ بعضی
زیاده از قبل ظاهر و از سیّد اعمال شفاعت مظلومها از اعدا نزد امرا بوده
البتّه بسمع بعضی رسیده که حزب مظلوم در آن مدینه نزد حاکم از قاتلها شفاعت
نمودند و تخفیف طلب کردند فاعتبروا یا اولی الأبصار

\end{otherlanguage}

Istidʿā ānkih Ḥaḍrat-i Sulṭān ayyadahu Allāhu tabāraka wa-taʿālā dar īn
umūr bi-nafsih tafakkur farmāyand va bi-ʿadl va inṣāf ḥukm namāyand.
Maʿa ānkih dar aktharī az bilād-i Īrān dar īn sinīn-i akhīrih īn ḥizb
kushtih shudand va nakushtand, mushāhadih mīshavad baghḍā-yi afʾidih-yi
baʿḍī ziyādih az qabl ẓāhir. Va az sayyid-i aʿmāl shafāʿat-i maẓlūmhā az
aʿdā nazd-i umarā būdih. Al-battih bi-samʿ-i baʿḍī rasīdih kih ḥizb-i
maẓlūm dar ān madīnih nazd-i ḥākim az qātilhā shafāʿat nimūdand va
takhfīf ṭalab kardand. Fa-ʿtabirū yā ulī al-abṣār.

It is Our request that His Majesty the sovereign, may God, blessed and
exalted, assist him, personally reflect on these matters and judge with
justice and fairness. Although in most cities of Iran in recent years
this community has been killed and has not killed, the hatred in some
hearts is seen to be more evident than before. Among the noblest deeds
has been the intercession of the wronged with the authorities on behalf
of their enemies. Some have surely heard that in that city the wronged
community interceded with the governor for the murderers and asked for
their punishment to be reduced. Take heed, then, you who have insight!

\subsection{129}\label{section-1415}

\begin{otherlanguage}{arabic}

یا شیخ این آیات محکمات در یکی از الواح از قلم ابهی نازل یا عبد اسمع ندآء
المظلوم الّذی حمل الشّدائد و البلایا فی سبیل اللّه مالک الأسمآء الی ان سجن
فى ارض الطّآء انّه دعا النّاس الی الجنّة العلیا و هم اخذوه و داروا به فی
المدن و الدّیار کم من لیل طار النّوم من عیون احبّائی حبّاً لنفسی و کم من یوم
قام علیّ الأحزاب مرّة رأیت نفسی علی اعلی الجبال و اخری فى سجن الطّآء فی
السّلاسل و الأغلال لعمر اللّه قد کنت شاکراً ناطقاً ذاکراً متوجّهاً راضیاً خاضعاً
خاشعاً فی کلّ الأحوال کذلک مضت ایّامی الی ان انتهت الی هذا السّجن الّذی به
تزلزلت الأرض و ناحت السّموات طوبی لعبد نبذ الظّنون اذ اتی المکنون برایات
الآیات انّا اخبرنا النّاس بهذا الظّهور الأعظم ولکنّ القوم فی سکر عجاب

\end{otherlanguage}

Yā Shaykh, īn āyāt-i muḥkamāt dar yikī az alvāḥ az Qalam-i Abhā nāzil:
Yā ʿabd, ismaʿ nidāʾa al-Maẓlūmi alladhī ḥamala al-shadāʾida wa-l-balāyā
fī sabīli Allāhi Māliki al-asmāʾi ilā an sujina fī arḍi al-Ṭāʾ. Innahu
daʿā al-nāsa ilā al-jannati al-ʿulyā wa-hum akhadhūhu wa-dārū bihi fī
al-muduni wa-l-diyār. Kam min laylin ṭāra al-nawmu min ʿuyūni aḥibbāʾī
ḥubban li-nafsī wa-kam min yawmin qāma ʿalayya al-aḥzāb. Marratan
raʾaytu nafsī ʿalā aʿlā al-jibāli wa-ukhrā fī sijni al-Ṭāʾi fī
al-salāsili wa-l-aghlāl. La-ʿamru Allāhi qad kuntu shākiran nāṭiqan
dhākiran mutawajjihan rāḍiyan khāḍiʿan khāshiʿan fī kulli al-aḥwāl.
Kadhālika maḍat ayyāmī ilā an intahat ilā hādhā al-sijni alladhī bihi
tazalzalati al-arḍu wa-nāḥati al-samāwāt. Ṭūbā li-ʿabdin nabadha
al-ẓunūna idh atā al-Maknūnu bi-rāyāti al-āyāt. Innā akhbarnā al-nāsa
bi-hādhā al-ẓuhūri al-aʿẓami wa-lākinna al-qawma fī sukrin ʿujāb.

Shaykh! These decisive verses were revealed by the All-Glorious Pen in
one of the Tablets: ``Servant! Hear the call of the Wronged One, who
bore hardships and afflictions in the path of God, the Lord of Names,
until He was imprisoned in the land of Ta. He called the people to the
most exalted Paradise, but they seized Him and took Him from city to
city and land to land. How many nights sleep fled from the eyes of My
loved ones out of love for Me, and how many days the factions rose
against Me! At one time I found Myself on the loftiest mountains, at
another in the prison of Ta, in chains and fetters. By the life of God!
In every circumstance I was thankful, speaking forth, remembering God,
turning toward Him, content, submissive, and humble. Thus My days passed
until they culminated in this prison, through which the earth trembled
and the heavens lamented. Blessed is the servant who cast aside
conjectures when the Hidden One came with the banners of His verses. We
have informed the people of this Most Great Revelation, but the people
are in a strange stupor.''

\subsection{130}\label{section-1416}

\begin{otherlanguage}{persian}

عند ذلک ارتفع النّدآء من شطر الحجاز نادت و قالت طوبی لک یا عکّا بما جعلک
اللّه مطلع ندائه الأحلی و مشرق آیاته الکبری و نعیماً لک بما استقرّ علیک
کرسیّ العدل و اشرق من افقک نیّر العنایة و الألطاف طوبی لمنصف انصف فی هذا
الذّکر الأعظم و ویل لکلّ غافل مرتاب

\end{otherlanguage}

ʿInda dhālika irtafaʿa al-nidāʾu min shaṭri al-Ḥijāz. Nādat wa-qālat
ṭūbā laki yā ʿAkkā bi-mā jaʿalaki Allāhu maṭliʿa nidāʾihi al-aḥlā
wa-mashriqa āyātihi al-kubrā wa-naʿīman laki bi-mā istaqarra ʿalayki
kursiyyu al-ʿadli wa-ashraqa min ufuqiki nayyiru al-ʿināyati wa-l-alṭāf.
Ṭūbā li-munṣifin anṣafa fī hādhā al-dhikri al-aʿẓami wa-waylun li-kulli
ghāfilin murtāb.

At that moment a call rose from the direction of the Hijaz. It called
out and said: ``Blessed are you, Acre, for God has made you the
dawning-place of His sweetest call and the dayspring of His greatest
signs. Bliss be yours, for the throne of justice has been established
upon you, and the sun of favor and loving-kindness has shone from your
horizon. Blessed is the fair-minded person who judges equitably
concerning this Most Great Remembrance, and woe to every heedless
doubter!''

\subsection{131}\label{section-1417}

\begin{otherlanguage}{arabic}

و بعد از شهادت بعضی از شهدا لوح برهان از سماء امر مالک ادیان نازل

\end{otherlanguage}

Va baʿd az shahādat-i baʿḍī az shuhadā Lawḥ-i Burhān az samāʾ-i amr-i
Mālik-i Adyān nāzil.

After the martyrdom of some of the martyrs, the Tablet of the Proof was
revealed from the heaven of the Cause of the Lord of Religions.

\subsection{132}\label{section-1418}

\begin{otherlanguage}{arabic}

هو المقتدر العلیم الحکیم قد احاطت اریاح البغضآء سفینة البطحآء بما اکتسبت
ایدی الظّالمین یا ایّها المشهور بالعلم قد افتیت علی الّذین ناح لهم کتب
العالم و شهد لهم دفاتر الأدیان کلّها و انّک انت یا ایّها البعید فی حجاب
غلیظ تاللّه قد حکمت علی الّذین بهم لاح افق الایمان یشهد بذلک مطالع الوحی و
مظاهر امر ربّک الرّحمن الّذین انفقوا ارواحهم و ما عندهم فی سبیله المستقیم
قد صاح من ظلمک دین اللّه فیما سواه و انّک تلعب و تکون من الفرحین لیس فی
قلبی بغضک و لا بغض احد من العباد انّ العارف یراک و امثالک فی جهل مبین انّک
لو اطّلعت علی ما فعلت لألقیت نفسک فی النّار او خرجت من البیت متوجّهاً الی
الجبال و نحت الی ان رجعت الی مقام قدّر لک من لدن مقتدر قدیر یا ایّها
الموهوم اخرق حجبات الظّنون و الأوهام لتری شمس العلم مشرقة من هذا الأفق
المنیر قد قطعت بضعة الرّسول و ظننت انّک نصرت دین اللّه کذلک سوّلت لک نفسک و
انت من الغافلین قد احترق من فعلک قلوب الملإ الأعلی و الّذین طافوا حول امر
اللّه ربّ العالمین قد ذاب کبد البتول من ظلمک و ناح اهل الفردوس فی مقام
کریم

\end{otherlanguage}

Huwa al-Muqtadir al-ʿAlīm al-Ḥakīm. Qad aḥāṭat aryāḥ al-baghḍāʾ safīnat
al-Baṭḥāʾ bi-mā iktasabat aydī al-ẓālimīn. Yā ayyuhā al-mashhūr
bi-l-ʿilm, qad aftayta ʿalā alladhīna nāḥa lahum kutub al-ʿālam
wa-shahida lahum dafātir al-adyān kullihā, wa-innaka anta yā ayyuhā
al-baʿīd fī ḥijāb ghalīẓ. Ta-Allāh, qad ḥakamta ʿalā alladhīna bihim
lāḥa ufuq al-īmān. Yashhadu bi-dhālika maṭāliʿ al-waḥy wa-maẓāhir amr
rabbika al-Raḥmān alladhīna anfaqū arwāḥahum wa-mā ʿindahum fī sabīlihi
al-mustaqīm. Qad ṣāḥa min ẓulmika dīn Allāh fī-mā siwāh, wa-innaka
talʿabu wa-takūnu min al-fariḥīn. Laysa fī qalbī bughḍuka wa-lā bughḍ
aḥad min al-ʿibād. Inna al-ʿārif yarāka wa-amthālaka fī jahl mubīn.
Innaka law iṭṭalaʿta ʿalā mā faʿalta la-alqayta nafsaka fī al-nār aw
kharajta min al-bayt mutawajjihan ilā al-jibāl wa-nuḥta ilā an rajaʿta
ilā maqām quddira laka min ladun Muqtadir Qadīr. Yā ayyuhā al-mawhūm,
ukhruq ḥujubāt al-ẓunūn wa-l-awhām li-tarā shams al-ʿilm mushriqatan min
hādhā al-ufuq al-munīr. Qad qaṭaʿta biḍʿat al-Rasūl wa-ẓananta annaka
naṣarta dīn Allāh. Kadhālika sawwalat laka nafsuka wa-anta min
al-ghāfilīn. Qad iḥtaraqa min fiʿlika qulūb al-malaʾ al-aʿlā
wa-alladhīna ṭāfū ḥawla amr Allāh Rabb al-ʿālamīn. Qad dhāba kabid
al-Batūl min ẓulmika wa-nāḥa ahl al-Firdaws fī maqām karīm.

He is the All-Powerful, the All-Knowing, the All-Wise. The winds of
hatred have surrounded the Ark of Batha because of what the hands of the
oppressors have done. You who are renowned for learning! You have issued
a ruling against those for whom the books of the world have lamented and
to whom the scriptures of every religion have borne witness. Yet you,
distant one, are behind a thick veil. By God, you have passed judgment
against those through whom the horizon of faith shone forth. To this
bear witness the Daysprings of revelation and the Manifestations of the
Cause of your Lord, the All-Merciful, who gave their lives and all they
possessed in His straight path. Because of your oppression, the religion
of God has cried out amid all else, while you amuse yourself and
rejoice. In My heart there is no hatred for you or for any other person.
One who truly knows sees you and your like in manifest ignorance. If you
realized what you have done, you would cast yourself into the fire, or
leave your home and head for the mountains, lamenting until you returned
to the station ordained for you by One who is All-Powerful, Almighty.
You who are deluded! Tear through the veils of conjecture and illusion,
so that you may see the sun of knowledge shining from this luminous
horizon. You have cut to pieces the Messenger's own flesh and imagined
that you were helping the religion of God. Thus has your self enticed
you, and you are among the heedless. Your deed has set ablaze the hearts
of the heavenly Assembly and of those who have circled around the Cause
of God, the Lord of the worlds. Your oppression has melted the heart of
the Virgin, and the inhabitants of Paradise have lamented in a noble
station.

\subsection{133}\label{section-1419}

\begin{otherlanguage}{persian}

انصف باللّه بأیّ برهان استدلّ علمآء الیهود و افتوا به علی الرّوح اذ اتی
بالحقّ و بأیّ حجّة انکر الفریسیّون و علمآء الأصنام اذ اتی محمّد رسول اللّه
بکتاب حکم بین الحقّ و الباطل بعدل اضآء بنوره ظلمات الأرض و انجذبت قلوب
العارفین و انّک استدللت الیوم بما استدلّ به علمآء الجهل فی ذاک العصر یشهد
بذلک مالک مصر الفضل فی هذا السّجن العظیم انّک اقتدیت بهم بل سبقتهم فی
الظّلم و ظننت انّک نصرت الدّین و دفعت عن شریعة اللّه العلیم الحکیم ونفسه
الحقّ ینوح من ظلمک النّاموس الأکبر و تصیح شریعة اللّه الّتی بها سرت نسمات
العدل علی من فی السّموات و الأرضین هل ظننت انّک ربحت فیما افتیت لا وسلطان
الأسمآء یشهد بخسرانک من عنده علم کلّ شیء فی لوح حفیظ

\end{otherlanguage}

Anṣif bi-Allāh, bi-ayyi burhān istadalla ʿulamāʾ al-Yahūd wa-aftaw bihi
ʿalā al-Rūḥ idh atā bi-l-ḥaqq, wa-bi-ayyi ḥujja ankara al-Farīsiyyūn
wa-ʿulamāʾ al-aṣnām idh atā Muḥammad Rasūl Allāh bi-kitāb ḥakama bayna
al-ḥaqq wa-l-bāṭil bi-ʿadl aḍāʾa bi-nūrihi ẓulumāt al-arḍ wa-injadhabat
qulūb al-ʿārifīn. Wa-innaka istadlalta al-yawm bi-mā istadalla bihi
ʿulamāʾ al-jahl fī dhāka al-ʿaṣr. Yashhadu bi-dhālika Mālik miṣr al-faḍl
fī hādhā al-sijn al-ʿaẓīm. Innaka iqtadayta bihim bal sabaqtahum fī
al-ẓulm wa-ẓananta annaka naṣarta al-dīn wa-dafaʿta ʿan sharīʿat Allāh
al-ʿAlīm al-Ḥakīm. Wa-nafsihi al-ḥaqq, yanūḥu min ẓulmika al-nāmūs
al-akbar wa-taṣīḥu sharīʿat Allāh allatī bihā sarat nasamāt al-ʿadl ʿalā
man fī al-samāwāt wa-l-arḍīn. Hal ẓananta annaka rabiḥta fī-mā aftayta?
Lā wa-Sulṭān al-asmāʾ. Yashhadu bi-khusrānika man ʿindahu ʿilm kull
shayʾ fī lawḥ ḥafīẓ.

Be fair, by God! What proof did the Jewish scholars invoke when they
issued their ruling against the Spirit as He came with the truth? And on
what evidence did the Pharisees and the scholars of idols base their
denial when Muhammad, the Messenger of God, came with a Book that judged
between truth and falsehood with a justice whose light illumined the
darkness of the earth and attracted the hearts of those who truly know?
Today you have invoked the same arguments that the scholars of ignorance
invoked in that age. The Sovereign of the city of grace bears witness to
this in this mighty prison. You have followed their example; indeed, you
have surpassed them in oppression, imagining that you were helping
religion and defending the Law of God, the All-Knowing, the All-Wise. By
His own true Self, the Most Great Law laments your oppression, and the
Law of God, through which the breezes of justice have wafted over all
who are in the heavens and on earth, cries out. Did you imagine that you
had gained anything by the ruling you issued? No, by the Sovereign of
Names! He who holds the knowledge of all things in a preserving Tablet
bears witness to your loss.

\subsection{134}\label{section-1420}

\begin{otherlanguage}{arabic}

یا ایّها الغافل انّک ما رأیتنی و ما عاشرت و ما آنست معی فی اقلّ من آن فکیف
امرت النّاس بسبّی هل اتّبعت فی ذلک هواک ام مولاک فأت بآیة ان انت من
الصّادقین نشهد انّک نبذت شریعة اللّه ورائک و اخذت شریعة نفسک انّه لا یعزب عن
علمه من شیء انّه هو الفرد الخبیر یا ایّها الغافل اسمع ما انزله الرّحمن فی
الفرقان لا تقولوا لمن القی الیکم السّلام لست مؤمناً کذلک حکم من فی قبضته
ملکوت الأمر و الخلق ان انت من السّامعین انّک نبذت حکم اللّه و اخذت حکم نفسک
فویل لک یا ایّها الغافل المریب انّک لو تنکرنی بأیّ برهان یثبت ما عندک فأت
به یا ایّها المشرک باللّه و المعرض عن سلطانه الّذی احاط العالمین

\end{otherlanguage}

Yā ayyuhā al-ghāfil, innaka mā raʾaytanī wa-mā ʿāsharta wa-mā ānasta
maʿī fī aqall min ān, fa-kayfa amarta al-nās bi-sabbī? Hal ittabaʿta fī
dhālika hawāka am Mawlāka? Faʾti bi-āya in anta min al-ṣādiqīn. Nashhadu
annaka nabadhta sharīʿat Allāh warāʾaka wa-akhadhta sharīʿat nafsika.
Innahu lā yaʿzubu ʿan ʿilmihi min shayʾ. Innahu huwa al-Fard al-Khabīr.
Yā ayyuhā al-ghāfil, ismaʿ mā anzalahu al-Raḥmān fī al-Furqān: lā taqūlū
li-man alqā ilaykum al-salām lasta muʾminan. Kadhālika ḥakama man fī
qabḍatihi malakūt al-amr wa-l-khalq, in anta min al-sāmiʿīn. Innaka
nabadhta ḥukm Allāh wa-akhadhta ḥukm nafsika, fa-wayl laka yā ayyuhā
al-ghāfil al-murīb. Innaka law tunkirunī bi-ayyi burhān yathbutu mā
ʿindaka? Faʾti bihi yā ayyuhā al-mushrik bi-Allāh wa-l-muʿriḍ ʿan
sulṭānihi alladhī aḥāṭa al-ʿālamīn.

Heedless one! You have not seen Me, associated with Me, or kept Me
company for even a moment. How, then, have you ordered people to revile
Me? In this, have you followed your own desire or your Lord? Bring forth
a sign, if you are truthful. We testify that you have cast the Law of
God behind you and adopted the law of your own self. Nothing escapes His
knowledge. He is the Unique, the All-Informed. Heedless one! Listen to
what the All-Merciful has revealed in the Qur'an: ``Do not say to one
who offers you peace, `You are not a believer.'\,'' Thus has He in whose
grasp lies the kingdom of revelation and creation decreed, if you are
among those who listen. You have cast aside God's judgment and adopted
your own. Woe to you, heedless doubter! If you deny Me, by what proof
can what you possess be established? Bring it forth, you who associate
partners with God and turn away from His sovereignty, which encompasses
all worlds.

\subsection{135}\label{section-1421}

\begin{otherlanguage}{arabic}

اعلم انّ العالم من اعترف بظهوری و شرب من بحر علمی و طار فی هوآء حبّی و نبذ
ما سوائی و اخذ ما نزّل من ملکوت بیانی البدیع انّه بمنزلة البصر للبشر و روح
الحیوان لجسد الامکان تعالی الرّحمن الّذی عرّفه و اقامه علی خدمة امره العزیز
العظیم یصلّی علیه الملأ الأعلی و اهل سرداق الکبریآء الّذین شربوا رحیقی
المختوم باسمی القویّ القدیر انّک ان تک من اهل هذا المقام الأعلی فأت بآیة
من لدی اللّه فاطر السّمآء و ان عرفت عجز نفسک خذ اعنّة هواک ثمّ ارجع الی
مولاک لعلّ یکفّر عنک سیّئاتک الّتی بها احترقت اوراق السّدرة و صاحت الصّخرة و
بکت عیون العارفین بک انشقّ ستر الرّبوبیّة و غرقت السّفینة و عقرت النّاقة و
ناح الرّوح فی مقام رفیع أ تعترض علی الّذى اتاک بما عندک و عند اهل العالم
من حجج اللّه و آیاته افتح بصرک لتری المظلوم مشرقاً من افق ارادة اللّه الملک
الحقّ المبین ثمّ افتح سمع فؤادک لتسمع ما تنطق به السّدرة الّتی ارتفعت بالحقّ
من لدی اللّه العزیز الجمیل انّ السّدرة مع ما ورد علیها من ظلمک و اعتساف
امثالک تنادی بأعلی النّدآء و تدعو الکلّ الی سدرة المنتهی و الأفق الأعلی
طوبی لنفس رأت الآیة الکبری و لأذن سمعت ندائها الأحلی و ویل لکلّ معرض اثیم

\end{otherlanguage}

Iʿlam anna al-ʿālim man iʿtarafa bi-ẓuhūrī wa-shariba min baḥr ʿilmī
wa-ṭāra fī hawāʾ ḥubbī wa-nabadha mā siwāʾī wa-akhadha mā nuzzila min
malakūt bayānī al-badīʿ. Innahu bi-manzilat al-baṣar li-l-bashar wa-rūḥ
al-ḥayawān li-jasad al-imkān. Taʿālā al-Raḥmān alladhī ʿarrafahu
wa-aqāmahu ʿalā khidmat amrihi al-ʿazīz al-ʿaẓīm. Yuṣallī ʿalayhi
al-malaʾ al-aʿlā wa-ahl surdāq al-kibriyāʾ alladhīna sharibū raḥīqī
al-makhtūm bi-smī al-Qawiyy al-Qadīr. Innaka in taku min ahl hādhā
al-maqām al-aʿlā faʾti bi-āya min ladā Allāh Fāṭir al-samāʾ, wa-in
ʿarafta ʿajz nafsika khudh aʿinnat hawāka thumma irjaʿ ilā Mawlāka
laʿalla yukaffira ʿanka sayyiʾātika allatī bihā iḥtaraqat awrāq al-sidra
wa-ṣāḥat al-ṣakhra wa-bakat ʿuyūn al-ʿārifīn. Bika inshaqqa sitr
al-rubūbiyya wa-ghariqat al-safīna wa-ʿuqirat al-nāqa wa-nāḥa al-Rūḥ fī
maqām rafīʿ. A-taʿtariḍu ʿalā alladhī atāka bi-mā ʿindaka wa-ʿinda ahl
al-ʿālam min ḥujaj Allāh wa-āyātihi? Iftaḥ baṣaraka li-tarā al-Maẓlūm
mushriqan min ufuq irādat Allāh al-Malik al-Ḥaqq al-Mubīn. Thumma iftaḥ
samʿ fuʾādika li-tasmaʿ mā tanṭiqu bihi al-sidra allatī irtafaʿat
bi-l-ḥaqq min ladā Allāh al-ʿAzīz al-Jamīl. Inna al-sidra maʿa mā warada
ʿalayhā min ẓulmika wa-iʿtisāf amthālika tunādī bi-aʿlā al-nidāʾ
wa-tadʿū al-kull ilā Sidrat al-Muntahā wa-l-ufuq al-aʿlā. Ṭūbā li-nafs
raʾat al-āya al-kubrā wa-li-udhun samiʿat nidāʾahā al-aḥlā wa-wayl
li-kull muʿriḍ athīm.

Know that a truly learned person is one who has acknowledged My
revelation, drunk from the ocean of My knowledge, soared in the
atmosphere of My love, cast aside everything but Me, and embraced what
has been sent down from the kingdom of My wondrous utterance. Such a
person is like sight to humanity and the spirit of life to the body of
contingent existence. Exalted is the All-Merciful, who has granted him
understanding and raised him up to serve His mighty and glorious Cause.
The heavenly Assembly and the inhabitants of the pavilion of grandeur,
who have drunk My choice wine sealed with My name, the Strong, the
Almighty, invoke blessings upon him. If you belong to this most exalted
station, bring forth a sign from God, the Creator of heaven. But if you
recognize your own powerlessness, rein in your desires and return to
your Lord, so that He may atone for your misdeeds, through which the
leaves of the Lote-Tree have burned, the Rock has cried out, and the
eyes of those who truly know have wept. Because of you, the veil of
lordship has been torn, the Ark has sunk, the she-camel has been
slaughtered, and the Spirit has lamented in a lofty station. Do you
object to Him who has brought you the proofs and signs of God that you
and the people of the world possess? Open your eyes, so that you may see
the Wronged One shining from the horizon of the will of God, the King,
the Truth, the Expounder. Then open the hearing of your heart, so that
you may hear what is spoken by the Lote-Tree that has been raised up in
truth by God, the Almighty, the Beautiful. Despite the oppression you
have inflicted upon it and the tyranny of those like you, the Lote-Tree
calls out with the loudest voice, summoning everyone to the Lote-Tree of
the Utmost Boundary and the most exalted horizon. Blessed is the soul
that has seen the greatest sign, and the ear that has heard its sweetest
call; and woe to every sinful one who turns away!

\subsection{136}\label{section-1422}

\begin{otherlanguage}{arabic}

یا ایّها المعرض باللّه لو تری السّدرة بعین الانصاف لتری آثار سیوفک فى
افنانها و اغصانها و اوراقها بعدما خلقک اللّه لعرفانها و خدمتها تفکّر لعلّ
تطّلع بظلمک و تکون من التّائبین أ ظننت انّا نخاف من ظلمک فاعلم ثمّ ایقن انّا
فی اوّل یوم فیه ارتفع صریر القلم الأعلی بین الأرض و السّمآء انفقنا ارواحنا
و اجسادنا و ابنائنا و اموالنا فی سبیل اللّه العلیّ العظیم و نفتخر بذلک بین
اهل الانشآء و الملإ الأعلی یشهد بذلک ما ورد علینا فی هذا الصّراط المستقیم
تاللّه قد ذابت الأکباد و صلبت الأجساد و سفکت الدّمآء و الأبصار کانت ناظرة
الی افق عنایة ربّها الشّاهد البصیر کلّما زاد البلآء زاد اهل البهآء فی حبّهم
قد شهد بصدقهم ما انزله الرّحمن فی الفرقان بقوله فتمنّوا الموت ان کنتم
صادقین هل الّذی حفظ نفسه خلف الأحجاب خیر ام الّذی انفقها فی سبیل اللّه انصف
و لا تکن فى تیه الکذب من الهائمین قد اخذهم کوثر محبّة الرّحمن علی شأن ما
منعتهم مدافع العالم و لا سیوف الأمم عن التّوجّه الى بحر عطآء ربّهم المعطی
الکریم

\end{otherlanguage}

Yā ayyuhā al-muʿriḍ, bi-Allāh, law tarā al-sidra bi-ʿayn al-inṣāf
la-tarā āthār suyūfika fī afnānihā wa-aghṣānihā wa-awrāqihā baʿdamā
khalaqaka Allāh li-ʿirfānihā wa-khidmatihā. Tafakkar laʿalla taṭṭaliʿu
bi-ẓulmika wa-takūnu min al-tāʾibīn. A-ẓananta annā nakhāfu min ẓulmika?
Fa-iʿlam thumma ayqin annā fī awwal yawm fīhi irtafaʿa ṣarīr al-Qalam
al-Aʿlā bayna al-arḍ wa-l-samāʾ anfaqnā arwāḥanā wa-ajsādanā wa-abnāʾanā
wa-amwālanā fī sabīl Allāh al-ʿAliyy al-ʿAẓīm wa-naftakhiru bi-dhālika
bayna ahl al-inshāʾ wa-l-malaʾ al-aʿlā. Yashhadu bi-dhālika mā warada
ʿalaynā fī hādhā al-ṣirāṭ al-mustaqīm. Ta-Allāh, qad dhābat al-akbād
wa-ṣulibat al-ajsād wa-sufikat al-dimāʾ wa-l-abṣār kānat nāẓiratan ilā
ufuq ʿināyat rabbihā al-Shāhid al-Baṣīr. Kullamā zāda al-balāʾ zāda ahl
al-Bahāʾ fī ḥubbihim. Qad shahida bi-ṣidqihim mā anzalahu al-Raḥmān fī
al-Furqān bi-qawlihi: fa-tamannaw al-mawt in kuntum ṣādiqīn. Hal alladhī
ḥafiẓa nafsahu khalf al-aḥjāb khayr am alladhī anfaqahā fī sabīl Allāh?
Anṣif wa-lā takun fī tīh al-kadhib min al-hāʾimīn. Qad akhadhahum
Kawthar maḥabbat al-Raḥmān ʿalā shaʾn mā manaʿathum madāfiʿ al-ʿālam
wa-lā suyūf al-umam ʿan al-tawajjuh ilā baḥr ʿaṭāʾ rabbihim al-Muʿṭī
al-Karīm.

You who turn away! By God, if you looked at the Lote-Tree with the eye
of fairness, you would see the marks of your swords on its boughs,
branches, and leaves, even though God created you to recognize and serve
it. Reflect, so that you may become aware of your oppression and be
among those who repent. Did you imagine that We fear your oppression?
Know, then, and be assured that on the very first day when the sound of
the Most Exalted Pen rose between earth and heaven, We offered Our
lives, Our bodies, Our children, and Our possessions in the path of God,
the Most High, the Most Great. In this We take pride before the people
of creation and the heavenly Assembly. What has befallen Us on this
straight path bears witness to this. By God, hearts have melted, bodies
have been crucified, and blood has been shed, while eyes remained fixed
on the horizon of the loving care of their Lord, the Witness, the
All-Seeing. The greater the affliction, the greater the love of the
people of Baha. What the All-Merciful revealed in the Qur'an bears
witness to their truthfulness in His words: ``Then wish for death, if
you are truthful.'' Is one who has kept himself safe behind veils
better, or one who has offered his life in the path of God? Be fair, and
do not be among those who wander in the wilderness of falsehood. The
river of the All-Merciful's love has so enraptured them that neither the
cannons of the world nor the swords of the nations have prevented them
from turning toward the ocean of the bounty of their Lord, the Bestower,
the All-Bountiful.

\subsection{137}\label{section-1423}

\begin{otherlanguage}{arabic}

تاللّه ما اعجزنی البلآء و ما اضعفنی اعراض العلمآء نطقت و انطق امام الوجوه
قد فتح باب الفضل و اتی مطلع العدل بآیات واضحات و حجج باهرات من لدی اللّه
المقتدر القدیر احضر بین یدی الوجه لتسمع اسرار ما سمعه ابن عمران فی طور
العرفان کذلک یأمرک مشرق ظهور ربّک الرّحمن من شطر سجنه العظیم

\end{otherlanguage}

Ta-Allāh, mā aʿjazanī al-balāʾ wa-mā aḍʿafanī iʿrāḍ al-ʿulamāʾ. Naṭaqtu
wa-anṭiqu amāma al-wujūh: qad futiḥa bāb al-faḍl wa-atā Maṭlaʿ al-ʿadl
bi-āyāt wāḍiḥāt wa-ḥujaj bāhirāt min ladā Allāh al-Muqtadir al-Qadīr.
Uḥḍur bayna yaday al-Wajh li-tasmaʿ asrār mā samiʿahu Ibn ʿImrān fī Ṭūr
al-ʿirfān. Kadhālika yaʾmuruka Mashriq ẓuhūr rabbika al-Raḥmān min shaṭr
sijnihi al-ʿaẓīm.

By God, affliction has not rendered Me powerless, nor has the rejection
of the scholars weakened Me. I have spoken, and I speak openly before
all: the gate of grace has been opened, and the Dayspring of justice has
come with clear signs and resplendent proofs from God, the All-Powerful,
the Almighty. Come into the presence of the Countenance, so that you may
hear the mysteries of what the son of Imran heard on the Sinai of
recognition. Thus does the Dayspring of the revelation of your Lord, the
All-Merciful, command you from the direction of His mighty prison.

\subsection{138}\label{section-1424}

\begin{otherlanguage}{arabic}

عند ذلک ارتفع ندآء الفطرة مرّة اخری و حنینها و تقول انّ الطّور ینادی و یقول
یا ملأ البیان اتّقوا الرّحمن انّی فزت بمکلّمی و اخذ جذب فرحی حصاة الأرض و
ترابها و السّدرة تقول یا ملأ البیان انصفوا فیما ظهر بالحقّ قد ظهرت النّار
الّتی اظهرها اللّه للکلیم یشهد بذلک کلّ ذی بصر علیم

\end{otherlanguage}

ʿInda dhālika irtafaʿa nidāʾ al-fiṭra marratan ukhrā wa-ḥanīnuhā
wa-taqūlu: inna al-Ṭūr yunādī wa-yaqūlu yā malaʾ al-Bayān, ittaqū
al-Raḥmān, innī fuztu bi-Mukallimī wa-akhadha jadhb faraḥī ḥaṣāt al-arḍ
wa-turābahā. Wa-l-sidra taqūlu: yā malaʾ al-Bayān, anṣifū fī-mā ẓahara
bi-l-ḥaqq. Qad ẓaharat al-nār allatī aẓharahā Allāh li-l-Kalīm. Yashhadu
bi-dhālika kull dhī baṣar ʿalīm.

At that moment the call of innate nature and its yearning cry rose once
again, and it said: ``Sinai calls out and says, `People of the Bayan!
Fear the All-Merciful. I have attained the presence of Him who speaks to
me, and the rapture of my joy has seized the pebbles and dust of the
earth.' And the Lote-Tree says, `People of the Bayan! Be fair concerning
what has appeared in truth. The fire that God revealed to Him who
conversed with Him has appeared. Everyone endowed with sight and
knowledge bears witness to this.'\,''

\subsection{139}\label{section-1425}

\begin{otherlanguage}{persian}

یا شیخ بعضی از شهدای این ظهور ذکر شد و همچنین بعضی از آیات که در ذکر
ایشان از ملکوت بیان نازل امید آنکه منقطعاً عن العالم در آنچه ذکر شده تفکّر
فرمایند

\end{otherlanguage}

Yā Shaykh, baʿżī az shohadā-ye īn ẓohūr dhekr shod va hamchenīn baʿżī az
āyāt ke dar dhekr-e īshān az malakūt-e bayān nāzel. Omīd ānke munqaṭiʿan
ʿan al-ʿālam dar ānche dhekr shode tafakkor farmāyand.

O Shaykh! Some of the martyrs of this revelation have been mentioned, as
have some of the verses revealed in their remembrance from the kingdom
of utterance. It is hoped that, detached from the world, you will
reflect on what has been mentioned.

\subsection{140}\label{section-1426}

\begin{otherlanguage}{persian}

حال در میرزا هادی دولت‌آبادی و صاد اصفهانی در ارض طاء تفکّر لازم اوّل بمجرّد
آنکه شنید او را بابی گفته‌اند اضطراب اخذش نمود بشأنی که وقار و سکون مفقود
گشت بر منابر ارتقا جست و نطق نمود بکلماتی که سزاوار نبود لازال گلپار‌ه‌های
عالم محض حبّ ریاست عمل نموده‌اند آنچه را که سبب و علّت گمراهی عباد گشته آن
جناب جمیع را مثل آن نفوس ندانند استقامت و ثبوت و رسوخ و اطمینان و تمکین
و وقار شهدای این ظهور ذکر شد تا آن جناب آگاه شوند و مقصود از آنچه بیان
شد از الواح ملوک و غیره آنکه آن جناب بیقین مبین بدانند که این مظلوم امر
اللّه را ستر ننمود و بآنچه مأمور بود امام وجوه عالم بافصح بیان ذکر نمود و
القا فرمود ولکن امثال آن نفوس ضعیفه مثل هادی و غیره امر اللّه را تبدیل
نمودند و نظر بزندگانی دوروزهٔ دنیا عمل کردند و گفتند آنچه را که عین عدل
گریست و قلم اعلی نوحه نمود مع آنکه از اصل امر بی‌خبر بوده و هستند و این
مظلوم لوجه اللّه اظهار نمود

\end{otherlanguage}

Ḥāl dar Mīrzā Hādī Dowlat-Ābādī va Ṣād-e Eṣfahānī dar arż-e Ṭāʾ tafakkor
lāzem. Avval be-mojarrad-e ānke shenīd ū rā Bābī gofte-and eżṭerāb
akhdhash namūd be-shaʾnī ke vaqār va sokūn mafqūd gasht. Bar manāber
erteqā jost va noṭq namūd be-kalemātī ke sazāvār nabūd. Lāzāl
gelpāreh-hā-ye ʿālam maḥż-e ḥobb-e riyāsat ʿamal namūde-and ānche rā ke
sabab va ʿellat-e gomrāhī-ye ʿebād gashte. Ān janāb jamīʿ rā methl-e ān
nofūs nadānand. Esteqāmat va thobūt va rosūkh va eṭmīnān va tamkīn va
vaqār-e shohadā-ye īn ẓohūr dhekr shod tā ān janāb āgāh shavand va
maqṣūd az ānche bayān shod az alvāḥ-e molūk va ghayreh ānke ān janāb
be-yaqīn-e mobīn bedānand ke īn Maẓlūm Amr Allāh rā satr nanamūd va
be-ānche maʾmūr būd amāma al-wujūh-e ʿālam be-afṣaḥ-e bayān dhekr namūd
va elqā farmūd. Va-līkan amthāl-e ān nofūs-e żaʿīfeh methl-e Hādī va
ghayreh Amr Allāh rā tabdīl namūdand va naẓar be-zendegānī-ye do-rūze-ye
donyā ʿamal kardand va goftand ānche rā ke ʿayn-e ʿadl gerīst va Qalam-e
Aʿlā nowḥeh namūd, maʿ-e ānke az aṣl-e amr bī-khabar būde va hastand va
īn Maẓlūm li-wajh Allāh eẓhār namūd.

Now it is necessary to reflect on Mirza Hadi Dawlat-Abadi and Sad of
Isfahan in the land of Ta. The former, as soon as he heard that he had
been called a Babi, was seized with such agitation that his dignity and
composure vanished. He mounted pulpits and uttered words that were
unworthy. The clods of the world have always, solely out of love of
leadership, done what has caused people to go astray. You should not
consider everyone to be like those souls. The steadfastness, constancy,
firmness, assurance, composure, and dignity of the martyrs of this
revelation have been mentioned so that you may become aware. The purpose
of what has been set forth from the Tablets to the kings and elsewhere
is that you may know with clear certainty that this Wronged One did not
conceal the Cause of God, but openly declared and proclaimed before the
people of the world, in the most eloquent language, what He had been
commanded to convey. Yet weak souls such as Hadi and others have altered
the Cause of God and, for the sake of this world's fleeting life, have
done and said things that have made the eye of justice weep and the Most
Exalted Pen lament, although they have been and remain unaware of the
fundamentals of the Cause. This Wronged One has made this known for the
sake of God.

\subsection{141}\label{section-1427}

\begin{otherlanguage}{persian}

یا هادی نزد اخوی رفتی و دیدی حال بساحت مظلوم توجّه نما شاید نفحات وحی و
فوحات الهام ترا تأیید نماید و بمقصود فائز گرداند هر نفسی الیوم بمشاهدهٔ
آثار فائز شود حقّ را از باطل بمثابهٔ شمس از ظلّ تمیز دهد و بمقصود آگاه گردد
حقّ شاهد و گواه که آنچه ذکر شده لوجه اللّه بوده که شاید آن جناب سبب هدایت
خلق شوند و احزاب عالم را از ظنون و اوهام نجات بخشند سبحان‌اللّه الی حین
معرضین و منکرین نمیدانند آنچه نزد مبشّر یعنی نقطه رفته از که بوده العلم
عند اللّه ربّ العالمین

\end{otherlanguage}

Yā Hādī, nazd-e akhavī raftī va dīdī. Ḥāl be-sāḥat-e Maẓlūm tavajjoh
namā shāyad nafaḥāt-e vaḥy va fawaḥāt-e elhām to rā taʾyīd namāyad va
be-maqṣūd fāʾez gardānad. Har nafsī al-yowm be-moshāhede-ye āthār fāʾez
shavad Ḥaqq rā az bāṭel be-mathābe-ye shams az ẓell tamīz dahad va
be-maqṣūd āgāh gardad. Ḥaqq shāhed va govāh ke ānche dhekr shode li-wajh
Allāh būde ke shāyad ān janāb sabab-e hedāyat-e khalq shavand va aḥzāb-e
ʿālam rā az ẓonūn va owhām najāt bakhshand. Subḥān Allāh, elā ḥīn
moʿreżīn va monkerīn nemīdānand ānche nazd-e Mobashsher yaʿnī Noqṭeh
rafte az ke būde. Al-ʿilm ʿinda Allāh Rabb al-ʿālamīn.

O Hadi! You went to My brother and saw him. Now turn toward the court of
the Wronged One, so that the breezes of revelation and the fragrances of
inspiration may strengthen you and enable you to attain your goal.
Anyone who today attains a view of the writings will distinguish truth
from falsehood as the sun is distinguished from shadow, and will become
aware of the intended purpose. God bears witness and testifies that what
has been mentioned has been for His sake, so that you may become a means
of guiding people and may deliver the peoples of the world from
conjectures and illusions. Glory be to God! To this day, those who turn
away and deny do not know who was responsible for what befell the
Herald, that is, the Point. Knowledge rests with God, the Lord of the
worlds.

\subsection{142}\label{section-1428}

\begin{otherlanguage}{persian}

یا شیخ همّت کن و بر خدمت امر قیام نما امروز رحیق مختوم امام وجوه ظاهر خذه
باسم ربّک ثمّ اشربه بذکره العزیز البدیع این مظلوم در لیالی و ایّام بتألیف
قلوب و تهذیب نفوس مشغول امورات واقعه در ایران در سنین اوّلیّه فی‌الحقیقه
سبب حزن مقرّبین و مخلصین بوده و در هر سنه قتل و غارت و تاراج و سفک دماء
موجود یک سنه در زنجان ظاهر شد آنچه که سبب فزع اکبر بود و همچنین یک سنه
در نیریز و سنهٔ دیگر در طبری تا آنکه واقعهٔ ارض طاء واقع و از آن ببعد این
مظلوم باعانت حقّ جلّ جلاله این حزب مظلوم را بما ینبغى آگاه نمود کل از ما
عندهم و ما عند القوم مقدّس و به ما عند اللّه متشبّث و ناظر

\end{otherlanguage}

Yā Shaykh, hemmat kon va bar khedmat-e Amr qiyām namā. Emrūz raḥīq-e
makhtūm amām-e vojūh ẓāher. Khudhhu bi-smi rabbika thumma ishrabhu
bi-dhikrihi al-ʿAzīz al-Badīʿ. Īn Maẓlūm dar layālī va ayyām be-taʾlīf-e
qolūb va tahdhīb-e nofūs mashghūl. Omūrāt-e vāqeʿe dar Īrān dar senīn-e
avvaliyye fī'l-ḥaqīqe sabab-e ḥozn-e moqarrabīn va mokhleṣīn būde va dar
har sane qatl va ghārat va tārāj va safk-e demāʾ mowjūd. Yek sane dar
Zanjān ẓāher shod ānche ke sabab-e fazaʿ-e akbar būd va hamchenīn yek
sane dar Nayrīz va sane-ye dīgar dar Ṭabarī, tā ānke vāqeʿe-ye arḍ-e Ṭāʾ
vāqeʿ va az ān be-baʿd īn Maẓlūm be-eʿānat-e Ḥaqq, jalla jalāluhu, īn
ḥezb-e maẓlūm rā bi-mā yanbaghī āgāh nemūd. Koll az mā ʿindahum va mā
ʿinda al-qawm moqaddas va be mā ʿinda Allāh motashabbeth va nāẓer.

O Shaykh, exert yourself and arise to serve the Cause. Today the sealed
wine is manifest before all eyes. Take it in the name of your Lord, then
drink it in remembrance of Him, the Almighty, the Incomparable. This
Wronged One is occupied day and night with uniting hearts and refining
souls. The events that occurred in Iran in the early years were truly a
cause of sorrow to those who are near to God and devoted to Him. Every
year there was killing, pillage, plunder, and bloodshed. One year,
events occurred in Zanjan that caused the greatest alarm; likewise, one
year in Nayriz, and another year in Tabari, until the incident in the
land of Ta occurred. From that time onward, this Wronged One, with the
assistance of God---exalted be His glory---made this wronged community
aware of what was fitting. All have freed themselves from what they and
the people possess, and hold fast to and look toward what is with God.

\subsection{143}\label{section-1429}

\begin{otherlanguage}{persian}

حال باید حضرت سلطان حفظه اللّه تعالی بعنایت و شفقت با این حزب رفتار
فرمایند و این مظلوم امام کعبهٔ الهی عهد مینماید از این حزب جز صداقت و
امانت امری ظاهر نشود که مغایر رأی جهان‌آرای حضرت سلطانی باشد هر ملّتی باید
مقام سلطانش را ملاحظه نماید و در آن خاضع باشد و بامرش عامل و بحکمش متمسّک
سلاطین مظاهر قدرت و رفعت و عظمت الهی بوده و هستند این مظلوم با احدی
مداهنه ننموده کل بر این فقره شاهد و گواهند ولکن ملاحظهٔ شئون سلاطین من
عند اللّه بوده و از کلمات انبیا و اولیا واضح و معلوم خدمت حضرت روح علیه
السّلام عرض نمودند یا روح اللّه أ یجوز ان تعطی جزیة لقیصر ام لا قال بلی ما
لقیصر لقیصر و ما للّه للّه منع نفرمودند و این دو کلمه یکی است نزد متبصّرین
چه که ما لقیصر اگر من عند اللّه نبود نهی میفرمودند و همچنین در آیهٔ مبارکهٔ
اطیعوا اللّه و اطیعوا الرّسول و اولی الأمر منکم مقصود از این اولو الأمر در
مقام اوّل و رتبهٔ اولی ائمّه صلوات اللّه علیهم بوده و هستند ایشانند مظاهر
قدرت و مصادر امر و مخازن علم و مطالع حکم الهی و در رتبهٔ ثانی و مقام ثانی
ملوک و سلاطین بوده‌اند یعنى ملوکی که بنور عدلشان آفاق عالم منوّر و روشن
است امید آنکه از حضرت سلطان نور عدلی اشراق نماید که جمیع احزاب امم را
احاطه کند کل باید از حقّ از برایش بطلبند آنچه را که الیوم سزاوار است

\end{otherlanguage}

Ḥāl bāyad Ḥaḍrat-e Solṭān, ḥafiẓahu Allāhu taʿālā, be-ʿenāyat va
shafaqat bā īn ḥezb raftār farmāyand va īn Maẓlūm amām-e Kaʿbe-ye elāhī
ʿahd mīnamāyad az īn ḥezb joz ṣedāqat va amānat amrī ẓāher nashavad ke
moghāyer-e raʾy-e jahān-ārā-ye Ḥaḍrat-e Solṭānī bāshad. Har mellatī
bāyad maqām-e solṭānash rā molāḥeẓe namāyad va dar ān khāḍeʿ bāshad va
be-amrash ʿāmel va be-ḥokmash motamassek. Salāṭīn maẓāher-e qodrat va
refʿat va ʿaẓamat-e elāhī būde va hastand. Īn Maẓlūm bā aḥadī modāhene
nanemūde, koll bar īn faqare shāhed va govāhand, va lāken molāḥeẓe-ye
shoʾūn-e salāṭīn min ʿindi Allāh būde va az kalemāt-e anbiyā va owliyā
vāḍeḥ va maʿlūm. Khedmat-e Ḥaḍrat-e Rūḥ, ʿalayhi al-salām, ʿarḍ
nemūdand: Yā Rūḥa Allāh, a-yajūzu an tuʿṭā jizyatun li-Qayṣar am lā?
Qāla: Balā, mā li-Qayṣar li-Qayṣar wa-mā li-llāh li-llāh. Manʿ
nafarmūdand va īn do kalame yekī ast nazd-e motabaṣṣerīn, che ke mā
li-Qayṣar agar min ʿindi Allāh nabūd nahy mīfarmūdand. Va hamchenīn dar
āye-ye mobārake-ye Aṭīʿū Allāha wa-aṭīʿū al-rasūla wa-ulī al-amri
minkum, maqṣūd az īn ulū al-amr dar maqām-e avval va rotbe-ye ūlā
aʾemme, ṣalawātu Allāhi ʿalayhim, būde va hastand. Īshānand maẓāher-e
qodrat va maṣāder-e amr va makhāzen-e ʿelm va maṭāleʿ-e ḥekam-e elāhī va
dar rotbe-ye thānī va maqām-e thānī molūk va salāṭīn būde-and, yaʿnī
molūkī ke be-nūr-e ʿadles hān āfāq-e ʿālam monavvar va rowshan ast. Omīd
ānke az Ḥaḍrat-e Solṭān nūr-e ʿadlī eshrāq namāyad ke jamīʿ-e aḥzāb-e
omam rā eḥāṭe konad. Koll bāyad az Ḥaqq az barāyash beṭalaband ānche rā
ke al-yawm sazāvār ast.

Now His Majesty the Shah---may God, exalted be He, protect him---should
treat this community with kindness and compassion. This Wronged One
pledges before the divine Kaaba that nothing but truthfulness and
trustworthiness will appear from this community, and nothing contrary to
His Majesty's world-adorning judgment. Every nation should respect the
station of its sovereign, submit to it, act according to his command,
and adhere to his decree. Kings have been and remain manifestations of
divine power, loftiness, and grandeur. This Wronged One has flattered no
one; all bear witness and testify to this. Nevertheless, regard for the
dignity of kings is ordained by God, as is clear and evident from the
words of the Prophets and holy ones. The Spirit---peace be upon
Him---was asked, ``O Spirit of God, is it permissible to pay tribute to
Caesar or not?'' He said, ``Yes. What belongs to Caesar is Caesar's, and
what belongs to God is God's.'' He did not forbid it. To those with
insight, these two statements are one, for if what belonged to Caesar
had not been from God, He would have forbidden it. Likewise, in the
blessed verse, ``Obey God, and obey the Messenger and those among you
who hold authority,'' those who hold authority have been and remain, in
the first instance and foremost rank, the Imams---may God's blessings
rest upon them. They are the manifestations of power, the sources of
command, the treasuries of knowledge, and the dawning places of divine
wisdom. In the second rank and station, the reference has been to kings
and sovereigns---that is, kings whose justice illumines and brightens
the horizons of the world. It is hoped that a light of justice will
shine from His Majesty the Shah and encompass all the peoples of the
nations. Everyone should ask God to grant him what is fitting for this
day.

\subsection{144}\label{section-1430}

\begin{otherlanguage}{arabic}

الهی الهی و سیّدی و سندی و مقصودی و محبوبی اسألک بالأسرار الّتی کانت
مکنونة فی علمک و بالآیات الّتی منها تضوّع عرف عنایتک و بأمواج بحر عطائک و
سمآء فضلک و کرمک و بالدّمآء الّتی سفکت فی سبیلک و بالأکباد الّتی ذابت فی
حبّک ان تؤیّد حضرة السّلطان بقدرتک و سلطانک لیظهر منه ما یکون باقیاً فی کتبک
و صحفک و الواحک ای ربّ خذ یده بید اقتدارک و نوّره بنور معرفتک و زیّنه بطراز
اخلاقک انّک انت المقتدر على ما تشآء و فی قبضتک زمام الأشیآء لا اله الّا
انت الغفور الکریم

\end{otherlanguage}

Ilāhī ilāhī wa-sayyidī wa-sanadī wa-maqṣūdī wa-maḥbūbī, asʾaluka
bi-l-asrāri allatī kānat maknūnatan fī ʿilmika wa-bi-l-āyāti allatī
minhā taḍawwaʿa ʿarfu ʿināyatika wa-bi-amwāji baḥri ʿaṭāʾika wa-samāʾi
faḍlika wa-karamika wa-bi-l-dimāʾi allatī sufikat fī sabīlika
wa-bi-l-akbādi allatī dhābat fī ḥubbika an tuʾayyida Ḥaḍrata al-Sulṭān
bi-qudratika wa-sulṭānika li-yaẓhara minhu mā yakūnu bāqiyan fī kutubika
wa-ṣuḥufika wa-alwāḥika. Ay Rabb, khudh yadahu bi-yadi iqtidārika
wa-nawwirhu bi-nūri maʿrifatika wa-zayyinhu bi-ṭirāzi akhlāqika. Innaka
anta al-Muqtadiru ʿalā mā tashāʾu wa-fī qabḍatika zimāmu al-ashyāʾ. Lā
ilāha illā anta al-Ghafūru al-Karīm.

My God, my God, my Lord, my support, my desire, and my beloved! I ask
You, by the mysteries concealed in Your knowledge, by the verses from
which the fragrance of Your loving care has spread, by the waves of the
ocean of Your bounty and the heaven of Your grace and generosity, by the
blood shed in Your path, and by the hearts that have melted in love for
You, to strengthen His Majesty the Shah through Your power and
sovereignty, so that there may appear from him what will endure in Your
books, Your scriptures, and Your Tablets. O Lord, take his hand with the
hand of Your might, illumine him with the light of Your knowledge, and
adorn him with the ornament of Your attributes. You have the power to do
whatever You will, and the reins of all things are in Your grasp. There
is no God but You, the Ever-Forgiving, the Generous.

\subsection{145}\label{section-1431}

\begin{otherlanguage}{arabic}

حضرت بولس قدّیس در رسالهٔ باهل رومیّه نوشته لتخضع کلّ نفس للسّلاطین العالیة
فانّه لا سلطان الّا من اللّه و السّلاطین الکائنة انّما رتّبها اللّه فمن یقاوم
السّلطان فانّه یعاند ترتیب اللّه الی ان قال لأنّه خادم اللّه المنتقم الّذی
ینفذ الغضب علی من یفعل الشّرّ میفرماید ظهور سلاطین و شوکت و اقتدارشان من
عند اللّه بوده

\end{otherlanguage}

Ḥaḍrat-e Būlos-e Qeddīs dar resāle-ye be-ahl-e Rūmiyye neveshte:
Li-takhḍaʿ kullu nafsin li-l-salāṭīni al-ʿāliyati fa-innahu lā sulṭāna
illā min Allāhi wa-l-salāṭīnu al-kāʾinatu innamā rattabahā Allāhu fa-man
yuqāwimu al-sulṭāna fa-innahu yuʿānidu tartība Allāh. Ilā an qāla:
Li-annahu khādimu Allāhi al-muntaqimu alladhī yunfidh u al-ghaḍaba ʿalā
man yafʿalu al-sharr. Mīfarmāyad ẓohūr-e salāṭīn va showkat va
eqtedāreshān min ʿindi Allāh būde.

Saint Paul wrote in his Epistle to the Romans, ``Let every soul submit
to the higher authorities, for there is no authority except from God,
and the existing authorities have been established by God. Whoever
resists authority therefore opposes God's ordering.'' He goes on to say,
``For he is God's servant, the avenger who executes wrath upon the
evildoer.'' He states that the rise of sovereigns, and their majesty and
power, are from God.

\subsection{146}\label{section-1432}

\begin{otherlanguage}{persian}

در احادیث قبل هم ذکر شده آنچه که علما دیده و شنیده‌اند نسأل اللّه تبارک و
تعالی ان یؤیّدک یا شیخ علی التّمسّک بما نزّل من سمآء عطآء اللّه ربّ العالمین
علما باید با حضرت سلطان متّحد شوند و بآنچه سبب حفظ و حراست و نعمت و ثروت
عباد است تمسّک نمایند سلطان عادل عنداللّه اقربست از کل یشهد بذلک من ینطق
فی السّجن الأعظم اللّه لا اله الّا هو الفرد الواحد المقتدر العلیم الحکیم

\end{otherlanguage}

Dar aḥādīth-e qabl ham dhekr shode ānche ke ʿolamā dīde va shenīde-and.
Nasʾalu Allāha tabāraka wa-taʿālā an yuʾayyidaka yā Shaykh ʿalā
al-tamassuki bi-mā nuzzila min samāʾi ʿaṭāʾi Allāhi Rabbi al-ʿālamīn.
ʿOlamā bāyad bā Ḥaḍrat-e Solṭān mottaḥed shavand va be-ānche sabab-e
ḥefẓ va ḥerāsat va neʿmat va thervat-e ʿebād ast tamassok namāyand.
Solṭān-e ʿādel ʿinda Allāh aqrab ast az koll. Yashhadu bi-dhālika man
yanṭiqu fī al-Sijni al-Aʿẓam. Allāhu lā ilāha illā huwa al-Fardu
al-Wāḥidu al-Muqtadiru al-ʿAlīmu al-Ḥakīm.

The traditions of the past also contain statements that the religious
scholars have seen and heard. We ask God, blessed and exalted be He, to
help you, O Shaykh, to hold fast to what has been sent down from the
heaven of the bounty of God, the Lord of the worlds. The religious
scholars should unite with His Majesty the Shah and hold fast to what
ensures the protection, security, welfare, and prosperity of the people.
A just sovereign is nearer to God than anyone else. He who speaks in the
Most Great Prison bears witness to this. God---there is no God but Him,
the Unique, the One, the All-Powerful, the All-Knowing, the All-Wise.

\subsection{147}\label{section-1433}

\begin{otherlanguage}{persian}

اگر آن جناب یک ساعت لوجه اللّه در اموری که از قبل و بعد ظاهر شده تفکّر
فرمایند از ما عنده به ما عند اللّه توجّه نمایند و سبب اعلاء کلمة اللّه
گردند آیا از اوّل ابداع الی حین از مشرق ارادهٔ الهی نوری و یا ظهوری اشراق
نموده که قبائل ارض او را قبول نموده و امرش را پذیرفته‌اند آن که بوده و
اسمش چه از حضرت خاتم روح ما سواه فداه و من قبله حضرت روح الی ان ینتهی
الی البدیع الأوّل در احیان ظهور کل مبتلا بوده‌اند بعضی را مجنون گفته‌اند و
برخی را کذّاب نامیده‌اند و عمل نموده‌اند آنچه را که قلم حیا میکند از ذکرش
لعمر اللّه ورد علیهم ما ناحت به الأشیآء کلّها ولکنّ القوم اکثرهم فی جهل
مبین نسأل اللّه ان یؤیّدهم علی الرّجوع الیه و الانابة لدی باب رحمته انّه علی
کلّ شیء قدیر

\end{otherlanguage}

Agar ān janāb yek sāʿat li-wajhi Allāh dar omūrī ke az qabl va baʿd
ẓāher shode tafakkor farmāyand, az mā ʿindahu be mā ʿinda Allāh tavajjoh
namāyand va sabab-e eʿlā-ye Kalimat Allāh gardand. Āyā az avval-e ebdāʿ
ilā ḥīn az mashreq-e erāde-ye elāhī nūrī va yā ẓohūrī eshrāq nemūde ke
qabāʾel-e arḍ ū rā qabūl nemūde va amrash rā padhīrofte-and? Ān ke būde
va esmash che? Az Ḥaḍrat-e Khātam, rūḥu mā siwāhu fidāhu, va min qablihi
Ḥaḍrat-e Rūḥ ilā an yantahiya ilā al-Badīʿ al-Awwal, dar aḥyān-e ẓohūr
koll mobtalā būde-and. Baʿḍī rā majnūn gofte-and va barkhī rā kadhdhāb
nāmīde-and va ʿamal nemūde-and ānche rā ke qalam ḥayā mīkonad az
dhekrash. La-ʿamru Allāh, warada ʿalayhim mā nāḥat bihi al-ashyāʾu
kulluhā wa-lākinna al-qawma aktharuhum fī jahlin mubīn. Nasʾalu Allāha
an yuʾayyidahum ʿalā al-rujūʿi ilayhi wa-l-inābati ladā bābi raḥmatihi.
Innahu ʿalā kulli shayʾin qadīr.

If you were to reflect for one hour, for God's sake, on the events that
have occurred in the past and more recently, you would turn from what
you possess toward what is with God and become a means of exalting the
Word of God. From the beginning of creation until now, has any light or
revelation shone from the dawning place of the divine will whose bearer
the peoples of the earth accepted and whose Cause they embraced? Who was
he, and what was his name? From the Seal---may the souls of all others
be sacrificed for Him---and, before Him, the Spirit, back to the Primal
Originator, all suffered affliction at the time of their appearance.
Some were called madmen, and others were named liars, and things were
done that the pen is ashamed to mention. By the life of God, what befell
them caused all things to lament, yet most of the people are in manifest
ignorance. We ask God to help them return to Him and turn in repentance
at the gate of His mercy. He has power over all things.

\subsection{148}\label{section-1434}

\begin{otherlanguage}{arabic}

فی هذا الحین ارتفع صریر قلمی الأعلی و قال وصّ حضرة الشّیخ بما وصّیت به احد
اغصانک لعلّ نفحات البیان تجذبه و تقرّبه الی اللّه ربّ العالمین

\end{otherlanguage}

Fī hādhā al-ḥīn irtafaʿa ṣarīru qalamī al-Aʿlā wa-qāla: Waṣṣi Ḥaḍrata
al-Shaykh bi-mā waṣṣayta bihi aḥada Aghṣānika laʿalla nafaḥāti al-bayāni
tajdhubuhu wa-tuqarribuhu ilā Allāhi Rabbi al-ʿālamīn.

At this moment the sound of My Most Exalted Pen arose, and it said,
``Counsel the honored Shaykh with the counsel You gave to one of Your
Branches, so that the breezes of utterance may attract him and draw him
near to God, the Lord of the worlds.''

\subsection{149}\label{section-1435}

\begin{otherlanguage}{arabic}

کن فی النّعمة منفقاً و فی فقدها شاکرا ً و فی الحقوق امیناً و فی الوجه طلقاً
و للفقرآء کنزاً و للأغنیآء ناصحاً و للمنادی مجیباً و فی الوعد وفیّاً و فی
الأمور منصفاً و فی الجمع صامتاً و فی القضآء عادلاً و للانسان خاضعاً و فی
الظّلمة سراجاً و للمهموم فرجاً و للظّمآن بحراً و للمکروب ملجأً و للمظلوم ناصراً
و عضداً و ظهراً و فی الأعمال متّقیاً و للغریب وطناً و للمریض شفآءً و للمستجیر
حصناً و للضّریر بصراً و لمن ضلّ صراطاً و لوجه الصّدق جمالاً و لهیکل الأمانة
طرازاً و لبیت الأخلاق عرشاً و لجسد العالم روحاً و لجند العدل رایةً و لأفق
الخیر نوراً و للأرض الطّیّبة رذاذاً و لبحر العلم فلکاً و لسمآء الکرم شمساً و
لرأس الحکمة اکلیلاً و لجبین الدّهر بیاضاً و لشجر الخضوع ثمراً نسأل اللّه ان
یحفظک من حرارة الحقد و صبارّة البرد انّه قریب مجیب کذلک نطق لسانی لأحد
اغصانی و ذکرناه لأحبّائی الّذین نبذوا الأوهام و اخذوا ما امروا به فی یوم
فیه اشرقت شمس الایقان من افق ارادة اللّه ربّ العالمین هذا یوم فیه غرّدت طیر
البیان علی الأغصان باسم ربّها الرّحمن طوبی لمن طار بأجنحة الاشتیاق الی
اللّه مالک یوم التّلاق

\end{otherlanguage}

Kun fī al-niʿmati munfiqan wa-fī faqdihā shākiran wa-fī al-ḥuqūqi amīnan
wa-fī al-wajhi ṭalqan wa-li-l-fuqarāʾi kanzan wa-li-l-aghniyāʾi nāṣiḥan
wa-li-l-munādī mujīban wa-fī al-waʿdi wafiyyan wa-fī al-umūri munṣifan
wa-fī al-jamʿi ṣāmitan wa-fī al-qaḍāʾi ʿādilan wa-li-l-insāni khāḍiʿan
wa-fī al-ẓulmati sirājan wa-li-l-mahmūmi farajan wa-li-l-ẓamʾāni baḥran
wa-li-l-makrūbi maljaʾan wa-li-l-maẓlūmi nāṣiran wa-ʿaḍudan wa-ẓahran
wa-fī al-aʿmāli muttaqiyan wa-li-l-gharībi waṭanan wa-li-l-marīḍi
shifāʾan wa-li-l-mustajīri ḥiṣnan wa-li-l-ḍarīri baṣaran wa-li-man ḍalla
ṣirāṭan wa-li-wajhi al-ṣidqi jamālan wa-li-haykali al-amānati ṭirāzan
wa-li-bayti al-akhlāqi ʿarshan wa-li-jasadi al-ʿālami rūḥan wa-li-jundi
al-ʿadli rāyatan wa-li-ufuqi al-khayri nūran wa-li-l-arḍi al-ṭayyibati
radhādhan wa-li-baḥri al-ʿilmi fulkan wa-li-samāʾi al-karami shamsan
wa-li-raʾsi al-ḥikmati iklīlan wa-li-jabīni al-dahri bayāḍan
wa-li-shajari al-khuḍūʿi thamaran. Nasʾalu Allāha an yaḥfaẓaka min
ḥarārati al-ḥiqdi wa-ṣabbārrati al-bard. Innahu qarībun mujīb. Kadhālika
naṭaqa lisānī li-aḥadi Aghṣānī wa-dhakarnāhu li-aḥibbāʾī alladhīna
nabadhū al-awhāma wa-akhadhū mā umirū bihi fī yawmin fīhi ashraqat
shamsu al-īqāni min ufuqi irādati Allāhi Rabbi al-ʿālamīn. Hādhā yawmun
fīhi gharradat ṭayru al-bayāni ʿalā al-Aghṣāni bi-smi rabbihā al-Raḥmān.
Ṭūbā li-man ṭāra bi-ajniḥati al-ishtiyāqi ilā Allāhi Māliki yawmi
al-talāq.

Be generous in prosperity and thankful when it is absent; trustworthy in
respecting others' rights, cheerful in countenance, a treasure to the
poor, a counselor to the rich, responsive to the caller, faithful to
your promises, fair in your dealings, silent in gatherings, just in
judgment, humble toward others, a lamp in darkness, a relief to the
sorrowful, an ocean to the thirsty, a refuge to the distressed, a
helper, a support, and a defender to the wronged, upright in your deeds,
a home to the stranger, a cure to the sick, a fortress to those seeking
shelter, sight to the blind, a path to those who have strayed, beauty to
the face of truthfulness, an ornament to the form of trustworthiness, a
throne to the house of virtue, a spirit to the body of the world, a
banner to the army of justice, a light to the horizon of goodness,
gentle rain to fertile soil, a ship upon the ocean of knowledge, a sun
in the heaven of generosity, a crown upon the head of wisdom, radiance
upon the brow of time, and fruit upon the tree of humility. We ask God
to protect you from the heat of malice and the bitter cold. He is near
and answers prayers. Thus My tongue spoke to one of My Branches, and We
have recounted it for My loved ones who have cast aside vain imaginings
and embraced what they were commanded on a day when the sun of certitude
shone from the horizon of the will of God, the Lord of the worlds. This
is a day when the birds of utterance have sung upon the Branches in the
name of their Lord, the All-Merciful. Blessed is the one who has flown
on the wings of longing toward God, the Lord of the Day of Meeting.

\subsection{150}\label{section-1436}

\begin{otherlanguage}{arabic}

حقّ آگاه و معشر امنا گواه که این مظلوم لازال تحت خطر عظیم بوده و لو لا
البلایا فی سبیل اللّه ما لذّ لی بقائی و ما نفعنی حیوتی و لا یخفی علی اهل
البصر و النّاظرین الی المنظر الأکبر بأنّی فی اکثر ایّامی کنت کعبد جالس تحت
سیف علّق بخیط واحد و لا یدری متی ینزل أ ینزل فی الحین او بعد حین و فی کلّ
ذلک نشکر اللّه ربّ العالمین و لسان سرّ در لیالی و ایّام باین مناجات ناطق
سبحانک یا الهى لو لا البلایا فی سبیلک من این یظهر مقام عاشقیک و لو لا
الرّزایا فی حبّک بأیّ شیء یثبت شأن مشتاقیک وعزّتک انیس محبّیک دموع عیونهم و
مؤنس مریدیک زفرات قلوبهم و غذآء قاصدیک قطعات اکبادهم و ما الذّ سمّ الرّدی
فی سبیلک و ما اعزّ سهام الأعدآء لاعلآء کلمتک یا الهی و سیّدی اشربنی فی
امرک ما اردته و انزل علیّ فی حبّک ما قدّرته وعزّتک لا ارید الّا ما ترید و لا
احبّ الّا ما انت تحبّ توکّلت علیک فی کلّ الأحوال انّک انت الغنیّ المتعال اسألک
یا الهی ان تظهر لنصرة هذا الظّهور من کان قابلاً لاسمک و سلطانک لیذکرک بین
خلقک و یرفع اعلام نصرک فی مملکتک و یزیّنهم بأخلاقک و اوامرک لا اله الّا
انت المهیمن القیّوم

\end{otherlanguage}

Ḥaqq āgāh va maʿshar-e omanā govāh ke īn Maẓlūm lāzāl taḥt-e khaṭar-e
ʿaẓīm būde. Wa-law lā al-balāyā fī sabīli Allāhi mā ladhdha lī baqāʾī
wa-mā nafaʿanī ḥayātī. Wa-lā yakhfā ʿalā ahli al-baṣari wa-l-nāẓirīna
ilā al-Manẓari al-Akbar bi-annī fī akthari ayyāmī kuntu ka-ʿabdin
jālisin taḥta sayfin ʿulliqa bi-khayṭin wāḥidin wa-lā yadrī matā
yanzilu, a-yanzilu fī al-ḥīni aw baʿda ḥīn. Wa-fī kulli dhālika nashkuru
Allāha Rabba al-ʿālamīn. Va lesān-e serr dar layālī va ayyām be-īn
monājāt nāṭeq: Subḥānaka yā ilāhī, law lā al-balāyā fī sabīlika min ayna
yaẓharu maqāmu ʿāshiqīka wa-law lā al-razāyā fī ḥubbika bi-ayyi shayʾin
yathbutu shaʾnu mushtāqīk? Wa-ʿizzatika, anīsu muḥibbīka dumūʿu
ʿuyūnihim wa-muʾnisu murīdīka zafarātu qulūbihim wa-ghidhāʾu qāṣidīka
qiṭaʿātu akbādihim. Wa-mā aladhdha samma al-radā fī sabīlika wa-mā
aʿazza sihāma al-aʿdāʾi li-iʿlāʾi kalimatik. Yā ilāhī wa-sayyidī,
ashribnī fī amrika mā aradtahu wa-anzil ʿalayya fī ḥubbika mā
qaddartahu. Wa-ʿizzatika, lā urīdu illā mā turīdu wa-lā uḥibbu illā mā
anta tuḥibb. Tawakkaltu ʿalayka fī kulli al-aḥwāl. Innaka anta
al-Ghaniyyu al-Mutaʿāl. Asʾaluka yā ilāhī an tuẓhira li-nuṣrati hādhā
al-ẓuhūri man kāna qābilan li-smika wa-sulṭānika li-yadhkuraka bayna
khalqika wa-yarfaʿa aʿlāma naṣrika fī mamlakatika wa-yuzayyinahum
bi-akhlāqika wa-awāmirika. Lā ilāha illā anta al-Muhayminu al-Qayyūm.

God knows, and the company of the trustworthy bears witness, that this
Wronged One has always been in grave danger. Were it not for afflictions
in the path of God, My continued existence would bring Me no pleasure
and My life would be of no benefit to Me. It is not hidden from those
with insight and those who gaze toward the Most Great Vision that, for
most of My days, I have been like a servant seated beneath a sword
suspended by a single thread, not knowing when it will fall---whether at
once or after a time. Through all this, We thank God, the Lord of the
worlds. And day and night My inmost tongue utters this prayer: ``Glory
be to You, my God! Were it not for afflictions in Your path, how would
the station of Your lovers be revealed? Were it not for calamities
suffered out of love for You, how would the rank of those who long for
You be established? By Your might, the companions of those who love You
are the tears of their eyes; the comfort of those devoted to You is the
sighing of their hearts; and the food of those who seek You is the
fragments of their own hearts. How sweet is the poison of death in Your
path, and how precious the arrows of enemies when borne to exalt Your
Word! My God and my Lord, give Me to drink in Your Cause whatever You
have willed, and send down upon Me, in love for You, whatever You have
ordained. By Your might, I desire only what You desire, and love only
what You love. In every circumstance I have placed My trust in You. You
are the Self-Sufficient, the Most Exalted. I ask You, my God, to raise
up in support of this Revelation someone worthy of Your name and Your
sovereignty, that he may remember You among Your creatures, raise the
banners of Your victory in Your kingdom, and adorn them with Your
attributes and Your commandments. There is no God but You, the
All-Protecting, the Self-Subsisting.''

\subsection{151}\label{section-1437}

\begin{otherlanguage}{arabic}

عند ذلک ارتفع ندآء الفطرة مرّة بعد مرّة نادت و قالت یا ملأ الأرض تاللّه انّی
لفطرة اللّه بینکم ایّاکم ان تنکرونی قد اظهرنی اللّه بنور احاط علی من فی
السّموات و الأرضین انصفوا یا قوم فی ظهوری و بروزی و اشراقی و لا تکونوا من
الظّالمین

\end{otherlanguage}

ʿInda dhālika irtafaʿa nidāʾu al-fiṭrati marratan baʿda marratin, nādat
wa-qālat: Yā malaʾa al-arḍ, ta-Allāhi innī la-fiṭratu Allāhi baynakum,
iyyākum an tunkirūnī. Qad aẓharanī Allāhu bi-nūrin aḥāṭa ʿalā man fī
al-samāwāti wa-l-araḍīn. Anṣifū yā qawmi fī ẓuhūrī wa-burūzī wa-ishrāqī
wa-lā takūnū min al-ẓālimīn.

At that moment the call of the innate nature rose again and again,
crying out and saying, ``O peoples of the earth! By God, I am God's
innate nature among you. Beware of denying Me. God has made Me manifest
with a light that has encompassed all who are in the heavens and the
earth. Judge fairly, O people, concerning My manifestation, My
emergence, and My radiance, and do not be among the unjust.''

\subsection{152}\label{section-1438}

\begin{otherlanguage}{arabic}

یا شیخ انّ المظلوم یسأل اللّه تبارک و تعالی ان یجعلک فاتح باب الانصاف و
یظهر بک امره بین العباد انّه هو المقتدر العزیز الوهّاب

\end{otherlanguage}

Yā Shaykh, inna al-Maẓlūma yasʾalu Allāha tabāraka wa-taʿālā an
yajʿalaka fātiḥa bābi al-inṣāfi wa-yuẓhira bika amrahu bayna al-ʿibād.
Innahu huwa al-Muqtadiru al-ʿAzīzu al-Wahhāb.

O Shaykh, this Wronged One asks God, blessed and exalted be He, to make
you an opener of the gate of fairness and to manifest His Cause through
you among His servants. He is the All-Powerful, the Almighty, the
Ever-Bestowing.

\subsection{153}\label{section-1439}

\begin{otherlanguage}{arabic}

یا شیخ از حقّ بطلب آذان و ابصار و افئدهٔ عالم را مقدّس فرماید و از هوای نفس
حفظ نماید چه که غرض مرضیست بزرگ انسان را از عرفان حضرت موجود محروم
مینماید و از تجلّیات انوار نیّر ایقان ممنوع میسازد از فضل و رحمت الهی سائل
و آمل که این مانع اکبر را از میان بردارد انّه هو القویّ الغالب القدیر

\end{otherlanguage}

Yā Shaykh, az Ḥaqq beṭalab ādhān va abṣār va afʾede-ye ʿālam rā moqaddas
farmāyad va az havā-ye nafs ḥefẓ namāyad, che ke gharaḍ maraḍīst bozorg,
ensān rā az ʿerfān-e Ḥażrat-e Mawjūd maḥrūm mīnamāyad va az tajallīyāt-e
anvār-e nayyer-e īqān mamnūʿ mīsāzad. Az fażl va raḥmat-e elāhī sāʾel va
āmel ke īn māneʿ-e akbar rā az miyān bardārad. Innahu huwa al-Qawī
al-Ghālib al-Qadīr.

O Shaykh, ask God to sanctify the ears, eyes, and hearts of the world
and protect them from selfish desire, for prejudice is a grave disease:
it deprives people of the recognition of the Lord of Being and bars them
from the radiance of the light of the sun of certitude. We ask and hope,
through divine grace and mercy, that He will remove this greatest
obstacle. He is indeed the Strong, the All-Subduing, the All-Powerful.

\subsection{154}\label{section-1440}

\begin{otherlanguage}{arabic}

فی هذا الحین ارتفع النّدآء عن یمین البقعة النّورآء اللّه لا اله الّا هو
الآمر الحکیم الق علی الشّیخ بقیّة لوح البرهان لیجذبه الی افق ظهور ربّه
الرّحمن لعلّه یقوم علی نصرة الأمر بآیات محکمات و براهین عالیات و ینطق بین
العباد بما نطق به لسان البرهان الملک للّه ربّ العالمین

\end{otherlanguage}

Fī hādhā al-ḥīn irtafaʿa al-nidāʾ ʿan yamīn al-buqʿa al-nawrāʾ: Allāh,
lā ilāha illā huwa al-Āmir al-Ḥakīm. Alqi ʿalā al-Shaykh baqiyyat Lawḥ
al-Burhān li-yajdhibahu ilā ufuq ẓuhūr rabbihi al-Raḥmān, laʿallahu
yaqūm ʿalā nuṣrat al-Amr bi-āyāt muḥkamāt wa-barāhīn ʿāliyāt wa-yanṭiq
bayna al-ʿibād bi-mā naṭaqa bihi lisān al-burhān: al-mulk li-llāh rabb
al-ʿālamīn.

At this moment the call rose from the right side of the luminous Spot:
``God---there is no god but Him, the Ordainer, the All-Wise. Recite to
the Shaykh the remainder of the Tablet of the Proof, that it may draw
him toward the horizon of the revelation of his Lord, the All-Merciful.
Perhaps he will arise to champion the Cause with conclusive verses and
lofty proofs, and proclaim among the servants what the tongue of proof
has proclaimed: `Sovereignty belongs to God, the Lord of all
worlds.'\,''

\subsection{155}\label{section-1441}

\begin{otherlanguage}{arabic}

اقرأ کتاب الایقان و ما انزله الرّحمن لملک باریس و امثاله لتطّلع بما قضی من
قبل و توقن بأنّا ما اردنا الفساد فی الأرض بعد اصلاحها انّما نذکّر العباد
خالصاً لوجه اللّه من شآء فلیقبل و من شآء فلیعرض انّ ربّنا الرّحمن لهو الغنیّ
الحمید یا معشر الأحزاب هذا یوم لا ینفعکم شیء من الأشیآء و لا اسم من
الأسمآء الّا بهذا الاسم الّذى جعله اللّه مظهر امره و مطلع اسمائه الحسنی لمن
فی ملکوت الانشآء نعیماً لمن وجد عرف الرّحمن و کان من الرّاسخین و لا یغنیکم
الیوم علومکم و فنونکم و لا زخارفکم و عزّکم دعوا الکلّ ورائکم مقبلین الی
الکلمة العلیا الّتی بها فصّلت الزّبر و الصّحف و هذا الکتاب المبین یا قوم
ضعوا ما الّفتموه من قلم الظّنون و الأوهام تاللّه قد اشرقت شمس العلم من افق
الیقین

\end{otherlanguage}

Iqraʾ Kitāb al-Īqān wa-mā anzalahu al-Raḥmān li-malik Bārīs wa-amthālihi
li-taṭṭaliʿa bi-mā quḍiya min qabl wa-tūqina bi-annā mā aradnā al-fasād
fī al-arḍ baʿda iṣlāḥihā. Innamā nudhakkir al-ʿibād khāliṣan li-wajh
Allāh. Man shāʾa fa-l-yuqbil wa-man shāʾa fa-l-yuʿriḍ. Inna rabbanā
al-Raḥmān la-huwa al-Ghanī al-Ḥamīd. Yā maʿshar al-aḥzāb, hādhā yawm lā
yanfaʿukum shayʾ min al-ashyāʾ wa-lā ism min al-asmāʾ illā bi-hādhā
al-ism alladhī jaʿalahu Allāh maẓhar amrihi wa-maṭlaʿ asmāʾihi al-ḥusnā
li-man fī malakūt al-inshāʾ. Naʿīman li-man wajada ʿarf al-Raḥmān
wa-kāna min al-rāsikhīn. Wa-lā yughnīkum al-yawm ʿulūmukum wa-funūnukum
wa-lā zakhārifukum wa-ʿizzukum. Daʿū al-kull warāʾakum muqbilīn ilā
al-kalima al-ʿulyā allatī bihā fuṣṣilat al-zubur wa-l-ṣuḥuf wa-hādhā
al-kitāb al-mubīn. Yā qawm, ḍaʿū mā allaftumūhu min qalam al-ẓunūn
wa-l-awhām. Ta-llāh, qad ashraqat shams al-ʿilm min ufuq al-yaqīn.

Read the Book of Certitude and what the All-Merciful revealed to the
King of Paris and others like him, so that you may learn what was
decreed before and be certain that We have not sought to spread
corruption on earth after it has been set right. We admonish the
servants solely for the sake of God. Let whoever wishes turn toward Him,
and whoever wishes turn away. Our Lord, the All-Merciful, is indeed the
Self-Sufficient, the Praiseworthy. O assemblage of contending parties,
this is a day when nothing whatsoever, and no name among all names, will
benefit you except through this Name, which God has made the
manifestation of His Cause and the dawning place of His most beautiful
names for all in the kingdom of creation. Blessed is the one who has
perceived the fragrance of the All-Merciful and remained steadfast.
Today neither your learning and arts nor your finery and prestige will
avail you. Leave everything behind, turning toward the Most Exalted
Word, through which the scriptures, the scrolls, and this clear Book
have been set forth in detail. O people, set aside what you have
composed with the pen of conjecture and vain imaginings. By God, the sun
of knowledge has risen from the horizon of certainty.

\subsection{156}\label{section-1442}

\begin{otherlanguage}{arabic}

یا ایّها الغافل ان کنت فی ریب ممّا نحن علیه انّا نشهد بما شهد اللّه قبل خلق
السّموات و الأرض انّه لا اله الّا هو العزیز الوهّاب و نشهد انّه کان واحداً فی
ذاته و واحداً فی صفاته لم یکن له شبه فی الابداع و لا شریک فی الاختراع قد
ارسل الرّسل و انزل الکتب لیبشّروا الخلق الی سوآء الصّراط

\end{otherlanguage}

Yā ayyuhā al-ghāfil, in kunta fī rayb mimmā naḥnu ʿalayhi, innā nashhadu
bi-mā shahida Allāh qabla khalq al-samāwāt wa-l-arḍ: innahu lā ilāha
illā huwa al-ʿAzīz al-Wahhāb. Wa-nashhadu annahu kāna wāḥidan fī dhātihi
wa-wāḥidan fī ṣifātihi. Lam yakun lahu shibh fī al-ibdāʿ wa-lā sharīk fī
al-ikhtirāʿ. Qad arsala al-rusul wa-anzala al-kutub li-yubashshirū
al-khalq ilā sawāʾ al-ṣirāṭ.

O heedless one, if you doubt what We profess, We testify to what God
testified before the creation of the heavens and the earth: there is no
god but Him, the Almighty, the Ever-Bestowing. And We testify that He
has ever been one in His essence and one in His attributes. He has no
likeness in creation and no partner in bringing things into being. He
sent the Messengers and revealed the Books to bring creation glad
tidings of the straight path.

\subsection{157}\label{section-1443}

\begin{otherlanguage}{arabic}

هل السّلطان اطّلع و غضّ الطّرف عن فعلک ام اخذه الرّعب بما عوت شرذمة من الذّئاب
الّذین نبذوا صراط اللّه ورائهم و اخذوا سبیلک من دون بیّنة و لا کتاب انّا
سمعنا بأنّ ممالک الایران تزیّنت بطراز العدل فلمّا تفرّسنا وجدناها مطالع
الظّلم و مشارق الاعتساف انّا نری العدل تحت مخالب الظّلم نسأل اللّه ان یخلّصه
بقوّة من عنده و سلطان من لدنه انّه هو المهیمن على من فی الأرضین و السّموات
لیس لأحد ان یعترض علی نفس فیما ورد علی امر اللّه ینبغی لکلّ من توجّه الی
الأفق الأعلی ان یتمسّک بحبل الاصطبار و یتوکّل علی اللّه المهیمن المختار یا
احبّآء اللّه اشربوا من عین الحکمة و طیروا فی هوا الحکمة و تکلّموا بالحکمة و
البیان کذلک یأمرکم ربّکم العزیز العلّام

\end{otherlanguage}

Hal al-sulṭān iṭṭalaʿa wa-ghaḍḍa al-ṭarf ʿan fiʿlika am akhadhahu
al-ruʿb bi-mā ʿawat shirdhima min al-dhiʾāb alladhīna nabadhū ṣirāṭ
Allāh warāʾahum wa-akhadhū sabīlaka min dūn bayyina wa-lā kitāb? Innā
samiʿnā bi-anna mamālik al-Īrān tazayyanat bi-ṭirāz al-ʿadl, fa-lammā
tafarrasnā wajadnāhā maṭāliʿ al-ẓulm wa-mashāriq al-iʿtisāf. Innā narā
al-ʿadl taḥta makhālib al-ẓulm. Nasʾal Allāh an yukhalliṣahu bi-quwwa
min ʿindihi wa-sulṭān min ladunhu. Innahu huwa al-Muhaymin ʿalā man fī
al-arḍīn wa-l-samāwāt. Laysa li-aḥad an yaʿtariḍa ʿalā nafs fī-mā warada
ʿalā amr Allāh. Yanbaghī li-kull man tawajjaha ilā al-ufuq al-aʿlā an
yatamassaka bi-ḥabl al-iṣṭibār wa-yatawakkala ʿalā Allāh al-Muhaymin
al-Mukhtār. Yā aḥibbāʾ Allāh, ishrabū min ʿayn al-ḥikma wa-ṭīrū fī hawā
al-ḥikma wa-takallamū bi-l-ḥikma wa-l-bayān. Kadhālika yaʾmurukum
rabbukum al-ʿAzīz al-ʿAllām.

Has the sovereign learned of your deed and turned a blind eye to it, or
has terror seized him at the howling of a band of wolves who have cast
the path of God behind them and followed your way without any evidence
or Book? We heard that the realms of Iran were adorned with the ornament
of justice, but when We examined them closely, We found them to be the
dawning places of oppression and the sources of tyranny. We see justice
beneath the claws of oppression. We ask God to deliver it through
strength from Him and authority from His presence. He is indeed the
Guardian over all who are in the earth and the heavens. No one is
entitled to object to any person over what has befallen the Cause of
God. Everyone who has turned toward the Most Exalted Horizon should hold
fast to the cord of patient endurance and put their trust in God, the
Guardian, who acts as He chooses. O loved ones of God, drink from the
spring of wisdom, soar in the air of wisdom, and speak with wisdom and
eloquence. Thus your Lord, the Almighty, the All-Knowing, commands you.

\subsection{158}\label{section-1444}

\begin{otherlanguage}{arabic}

یا غافل لا تطمئنّ بعزّک و اقتدارک مثلک کمثل بقیّة اثر الشّمس علی رؤوس الجبال
سوف یدرکها الزّوال من لدی اللّه الغنیّ المتعال قد اخذ عزّک و عزّ امثالک و هذا
ما حکم به من عنده امّ الألواح این من حارب اللّه و این من جادل بآیاته و این
من اعرض عن سلطانه و این الّذین قتلوا اصفیائه و سفکوا دمآء اولیائه تفکّر
لعلّ تجد نفحات اعمالک یا ایّها الجاهل المرتاب بکم ناح الرّسول و صاحت البتول
و خربت الدّیار و اخذت الظّلمة کلّ الأقطار یا معشر العلمآء بکم انحطّ شأن
الملّة و نکس علم الاسلام و ثلّ عرشه العظیم کلّما اراد ممیّز ان یتمسّک بما
یرتفع به شأن الاسلام ارتفعت ضوضائکم بذلک منع عمّا اراد و بقی الملک فی
خسران کبیر

\end{otherlanguage}

Yā ghāfil, lā taṭmaʾinna bi-ʿizzika wa-iqtidārika. Mathaluka ka-mathal
baqiyyat athar al-shams ʿalā ruʾūs al-jibāl, sawfa yudrikuhā al-zawāl
min ladā Allāh al-Ghanī al-Mutaʿāl. Qad ukhidha ʿizzuka wa-ʿizz
amthālika, wa-hādhā mā ḥakama bihi man ʿindahu Umm al-Alwāḥ. Ayna man
ḥāraba Allāh, wa-ayna man jādala bi-āyātihi, wa-ayna man aʿraḍa ʿan
sulṭānihi, wa-ayna alladhīna qatalū aṣfiyāʾahu wa-safakū dimāʾ
awliyāʾihi? Tafakkar laʿalla tajida nafaḥāt aʿmālika, yā ayyuhā al-jāhil
al-murtāb. Bikum nāḥa al-Rasūl wa-ṣāḥat al-Batūl wa-kharibat al-diyār
wa-akhadhat al-ẓulma kull al-aqṭār. Yā maʿshar al-ʿulamāʾ, bikum inḥaṭṭa
shaʾn al-milla wa-nukkisa ʿalam al-Islām wa-thulla ʿarshuhu al-ʿaẓīm.
Kullamā arāda mumayyiz an yatamassaka bi-mā yartafiʿu bihi shaʾn
al-Islām irtafaʿat ḍawḍāʾukum. Bi-dhālika muniʿa ʿammā arāda wa-baqiya
al-mulk fī khusrān kabīr.

O heedless one, do not rest secure in your glory and power. You are like
the last trace of sunlight on the mountaintops; it will soon fade away
at the bidding of God, the Self-Sufficient, the Most Exalted. Your glory
and the glory of those like you have been taken away. This is what He
with whom is the Mother of Tablets has decreed. Where is the one who
waged war against God? Where is the one who disputed His signs? Where is
the one who turned away from His sovereignty? And where are those who
killed His chosen ones and shed the blood of His loved ones? Reflect, so
that you may perhaps perceive the scent of your deeds, O ignorant and
doubting one. Because of you the Messenger lamented, the Chaste One
cried out, homes were laid waste, and darkness overtook every region. O
assemblage of religious scholars, because of you the standing of the
faith has declined, the banner of Islam has been lowered, and its mighty
throne has been overthrown. Whenever a discerning person sought to hold
fast to something that would exalt the standing of Islam, your clamor
arose. Thus he was prevented from carrying out his purpose, and the
realm remained in grievous loss.

\subsection{159}\label{section-1445}

\begin{otherlanguage}{persian}

یا قلمی الأعلی اذکر الرّقشآء الّتی بظلمها ناحت الأشیآء و ارتعدت فرائص
الأولیآء کذلک یأمرک مالک الأسمآء فی هذا المقام المحمود قد صاحت من ظلمک
البتول و تظنّ انّک من آل الرّسول کذلک سوّلت لک نفسک یا ایّها المعرض عن اللّه
ربّ ما کان و ما یکون انصفی یا ایّتها الرّقشآء بأیّ جرم لدغت ابنآء الرّسول و
نهبت اموالهم أ کفرت بالّذی خلقک بأمره کن فیکون قد فعلت بأبنآء الرّسول ما
لا فعلت عاد و ثمود بصالِح و هود و لا الیهود بروح اللّه مالک الوجود أ تنکر
آیات ربّک الّتی اذ نزّلت من سمآء الأمر خضعت لها کتب العالم کلّها تفکّر لتطّلع
بفعلک یا ایّها الغافل المردود سوف تأخذک نفحات العذاب کما اخذت قوماً قبلک
انتظر یا ایّها المشرک باللّه مالک الغیب و الشّهود هذا یوم اخبر به اللّه
بلسان رسوله تفکّر لتعرف ما انزله الرّحمن فی الفرقان و فی هذا اللّوح المسطور
هذا یوم فیه اتی مشرق الوحی بآیات بیّنات الّتی عجز عن احصائها المحصون هذا
یوم فیه وجد کلّ ذی شمّ عرف نسمة الرّحمن فی الامکان و سرع کلّ ذی بصر الی فرات
رحمة ربّه مالک الملوک یا ایّها الغافل تاللّه قد رجع حدیث الذّبح و الذّبیح
توجّه الی مقرّ الفدآء و ما رجع بما اکتسبت یدک یا ایّها المبغض العنود أ ظننت
بالشّهادة ینحطّ شأن الأمر لا والّذى جعله اللّه مهبط الوحی ان انت من الّذین هم
یفقهون ویل لک یا ایّها المشرک باللّه و للّذین اتّخذوک اماماً لأنفسهم من دون
بیّنة و لا کتاب مشهود کم من ظالم قام علی اطفآء نور اللّه قبلک و کم من فاجر
قتل و نهب الی ان ناحت من ظلمه الأفئدة و النّفوس قد غابت شمس العدل بما
استوی هیکل الظّلم علی اریکة البغضآء ولکنّ القوم هم لا یشعرون یا جاهل قد
قتلت ابنآء الرّسول و نهبت اموالهم قل هل الأموال کفرت باللّه ام مالکها على
زعمک انصف یا ایّها الجاهل المحجوب قد اخذت الاعتساف و نبذت الانصاف بذلک
ناحت الأشیآء و انت من الغافلین قد قتلت الکبیر و نهبت الصّغیر هل تظنّ انّک
تأکل ما جمعته بالظّلم لا ونفسی کذلک یخبرک الخبیر تاللّه لا یغنیک ما عندک و
ما جمعته بالاعتساف یشهد بذلک ربّک العلیم قد قمت على اطفآء نور الأمر سوف
تنخمد نارک امراً من عنده انّه هو المقتدر القدیر لا تعجزه شؤونات العالم و
لا سطوة الأمم یفعل ما یشآء بسلطانه و یحکم ما یرید تفکّر فی النّاقة مع انّها
من الحیوان رفعها الرّحمن الی مقام نطق السن العالم بذکرها و ثنائها انّه هو
المهیمن علی من فی السّموات و الأرض لا اله الّا هو العزیز العظیم کذلک زیّنّا
آفاق سمآء اللّوح بشموس الکلمات نعیماً لمن فاز بها و استضآء بأنوارها و ویل
للمعرضین و ویل للمنکرین و ویل للغافلین الحمد للّه ربّ العالمین

\end{otherlanguage}

Yā qalamī al-aʿlā, udhkur al-Raqshāʾ allatī bi-ẓulmihā nāḥat al-ashyāʾ
wa-irtaʿadat farāʾiṣ al-awliyāʾ. Kadhālika yaʾmuruka Mālik al-Asmāʾ fī
hādhā al-maqām al-maḥmūd. Qad ṣāḥat min ẓulmika al-Batūl wa-taẓunn
annaka min āl al-Rasūl. Kadhālika sawwalat laka nafsuka, yā ayyuhā
al-muʿriḍ ʿan Allāh rabb mā kāna wa-mā yakūn. Inṣifī yā ayyatuhā
al-Raqshāʾ, bi-ayyi jurm ladaghti abnāʾ al-Rasūl wa-nahabti amwālahum?
A-kafarti bi-alladhī khalaqaki bi-amrihi kun fa-yakūn? Qad faʿalti
bi-abnāʾ al-Rasūl mā lā faʿalat ʿĀd wa-Thamūd bi-Ṣāliḥ wa-Hūd wa-lā
al-Yahūd bi-Rūḥ Allāh Mālik al-Wujūd. A-tunkir āyāt rabbika allatī idh
nuzzilat min samāʾ al-amr khaḍaʿat lahā kutub al-ʿālam kulluhā? Tafakkar
li-taṭṭaliʿa bi-fiʿlika, yā ayyuhā al-ghāfil al-mardūd. Sawfa
taʾkhudhuka nafaḥāt al-ʿadhāb kamā akhadhat qawman qablaka. Intaẓir, yā
ayyuhā al-mushrik bi-llāh Mālik al-ghayb wa-l-shuhūd. Hādhā yawm akhbara
bihi Allāh bi-lisān rasūlihi. Tafakkar li-taʿrifa mā anzalahu al-Raḥmān
fī al-Furqān wa-fī hādhā al-lawḥ al-masṭūr. Hādhā yawm fīhi atā Mashriq
al-Waḥy bi-āyāt bayyināt allatī ʿajaza ʿan iḥṣāʾihā al-muḥṣūn. Hādhā
yawm fīhi wajada kull dhī shamm ʿarf nasamat al-Raḥmān fī al-imkān
wa-saraʿa kull dhī baṣar ilā furāt raḥmat rabbihi Mālik al-Mulūk. Yā
ayyuhā al-ghāfil, ta-llāh, qad rajaʿa ḥadīth al-dhabḥ wa-l-dhabīḥ
tawajjaha ilā maqarr al-fidāʾ wa-mā rajaʿa bi-mā iktasabat yaduka, yā
ayyuhā al-mubghiḍ al-ʿanūd. A-ẓananta bi-l-shahāda yanḥaṭṭ shaʾn al-Amr?
Lā wa-alladhī jaʿalahu Allāh mahbiṭ al-waḥy, in anta min alladhīna hum
yafqahūn. Wayl laka yā ayyuhā al-mushrik bi-llāh wa-li-alladhīna
ittakhadhūka imāman li-anfusihim min dūn bayyina wa-lā kitāb mashhūd.
Kam min ẓālim qāma ʿalā iṭfāʾ nūr Allāh qablaka wa-kam min fājir qatala
wa-nahaba ilā an nāḥat min ẓulmihi al-afʾida wa-l-nufūs. Qad ghābat
shams al-ʿadl bi-mā istawā haykal al-ẓulm ʿalā arīkat al-baghḍāʾ
wa-lākinna al-qawm hum lā yashʿurūn. Yā jāhil, qad qatalta abnāʾ
al-Rasūl wa-nahabta amwālahum. Qul hal al-amwāl kafarat bi-llāh am
mālikuhā ʿalā zaʿmika? Inṣif, yā ayyuhā al-jāhil al-maḥjūb. Qad akhadhta
al-iʿtisāf wa-nabadhta al-inṣāf. Bi-dhālika nāḥat al-ashyāʾ wa-anta min
al-ghāfilīn. Qad qatalta al-kabīr wa-nahabta al-ṣaghīr. Hal taẓunn
annaka taʾkul mā jamaʿtahu bi-l-ẓulm? Lā wa-nafsī, kadhālika yukhbiruka
al-Khabīr. Ta-llāh, lā yughnīka mā ʿindaka wa-mā jamaʿtahu bi-l-iʿtisāf.
Yashhadu bi-dhālika rabbuka al-ʿAlīm. Qad qumta ʿalā iṭfāʾ nūr al-Amr.
Sawfa tankhamid nāruka amran min ʿindihi. Innahu huwa al-Muqtadir
al-Qadīr. Lā tuʿjizuhu shuʾūnāt al-ʿālam wa-lā saṭwat al-umam. Yafʿalu
mā yashāʾ bi-sulṭānihi wa-yaḥkumu mā yurīd. Tafakkar fī al-nāqa, maʿa
annahā min al-ḥayawān rafaʿahā al-Raḥmān ilā maqām naṭaqa alsun al-ʿālam
bi-dhikrihā wa-thanāʾihā. Innahu huwa al-Muhaymin ʿalā man fī al-samāwāt
wa-l-arḍ. Lā ilāha illā huwa al-ʿAzīz al-ʿAẓīm. Kadhālika zayyannā āfāq
samāʾ al-lawḥ bi-shumūs al-kalimāt. Naʿīman li-man fāza bihā wa-istaḍāʾa
bi-anwārihā, wa-wayl li-l-muʿriḍīn wa-wayl li-l-munkirīn wa-wayl
li-l-ghāfilīn. Al-ḥamd li-llāh rabb al-ʿālamīn.

O My Most Exalted Pen, make mention of the She-Serpent, whose oppression
made all things lament and the limbs of the loved ones tremble. Thus the
Lord of Names commands you in this praiseworthy station. The Chaste One
has cried out at your oppression, yet you imagine that you belong to the
family of the Messenger. Thus has your own self deceived you, O you who
turn away from God, the Lord of all that has been and all that will be.
Be fair, O She-Serpent! For what crime did you bite the sons of the
Messenger and plunder their possessions? Have you disbelieved in Him who
created you through His command, ``Be,'' and it is? You have done to the
sons of the Messenger what neither ʿAd and Thamud did to Salih and Hud,
nor the Jews to the Spirit of God, the Lord of Being. Do you deny the
verses of your Lord, before which all the books of the world bowed down
when they were sent down from the heaven of His command? Reflect, so
that you may understand your deed, O heedless and rejected one. The
blasts of punishment will soon seize you, just as they seized a people
before you. Wait, O you who associate partners with God, the Lord of the
unseen and the seen. This is the day of which God gave tidings through
the tongue of His Messenger. Reflect, so that you may recognize what the
All-Merciful has revealed in the Quran and in this inscribed Tablet.
This is the day when the Dawning Place of Revelation has come with clear
verses that those who count have been unable to number. This is the day
when everyone endowed with the sense of smell has perceived the
fragrance of the breeze of the All-Merciful in the world of being, and
everyone endowed with sight has hastened to the sweet waters of the
mercy of his Lord, the King of Kings. O heedless one, by God, the story
of the sacrifice has returned, and the sacrificial victim has turned
toward the place of sacrifice and has not returned, because of what your
hands have done, O hateful and obstinate one. Did you imagine that
martyrdom would diminish the standing of the Cause? No, by Him whom God
has made the recipient of revelation, if you are among those who
understand. Woe to you, O you who associate partners with God, and to
those who have taken you as their leader without any evidence or
attested Book. How many oppressors before you arose to extinguish the
light of God, and how many evildoers killed and plundered until hearts
and souls lamented at their oppression! The sun of justice has set
because the embodiment of oppression has seated itself upon the throne
of hatred, but the people do not perceive it. O ignorant one, you have
killed the sons of the Messenger and plundered their possessions. Say:
Was it the possessions that disbelieved in God, or their owners, as you
allege? Be fair, O ignorant and veiled one. You have embraced tyranny
and cast fairness aside. Because of this all things have lamented, yet
you remain among the heedless. You have killed the elder and plundered
the younger. Do you imagine that you will enjoy what you have amassed
through oppression? No, by My own Self! Thus the All-Informed tells you.
By God, neither what you possess nor what you have amassed through
tyranny will avail you. Your Lord, the All-Knowing, bears witness to
this. You have arisen to extinguish the light of the Cause. Your own
fire will soon be put out by a command from Him. He is indeed the
All-Powerful, the Almighty. Neither the affairs of the world nor the
might of nations can frustrate Him. Through His sovereignty He does
whatever He pleases and decrees whatever He wills. Reflect on the
she-camel: although she was an animal, the All-Merciful exalted her to
such a station that the tongues of the world have spoken in her
remembrance and praise. He is indeed the Guardian over all who are in
the heavens and the earth. There is no god but Him, the Almighty, the
All-Glorious. Thus We have adorned the horizons of the heaven of this
Tablet with the suns of words. Blessed is the one who has attained them
and been illumined by their light. Woe to those who turn away, woe to
those who deny, and woe to the heedless. Praise be to God, the Lord of
all worlds.

\subsection{160}\label{section-1446}

\begin{otherlanguage}{arabic}

یا شیخ قد اسمعناک تغرّدات عندلیب الفردوس و اریناک الآثار الّتی انزلها اللّه
بأمره المبرم فی السّجن الأعظم لتقرّ بها عینک و تطمئنّ بها نفسک انّه هو
الفیّاض الکریم قم بقوّة البرهان علی خدمة امر اللّه ربّک الرّحمن لو تخاف من
ایمانک خذ اللّوح ثمّ احفظه فی جیب توکّلک و اذا دخلت موقف الحشر و یسألک اللّه
بأیّ حجّة آمنت بهذا الظّهور اخرج اللّوح و قل بهذا الکتاب المبارک العزیز
البدیع اذاً ترتفع الیک ایادی الکلّ و یأخذون اللّوح و یضعونه علی عیونهم و
یجدون منه عرف بیان اللّه ربّ العالمین لو یعذّبک اللّه بما آمنت بآیاته فی هذا
الظّهور فبأیّ حجّة یعذّب الّذین ما آمنوا بمحمّد رسول اللّه و من قبله بعیسی بن
مریم و من قبله بالکلیم و من قبله بالخلیل الی ان تنتهی الظّهورات الی
البدیع الأوّل الّذی خلق بارادة ربّک القادر المحیط کذلک انزلنا الآیات لأحد
قبلک و ذکرناها لک فی هذا الیوم لتعرف و تکون من الموقنین یا ایّها النّاطق
بالعلم انّ الأمر اظهر من ان یخفی و ابین من ان یستر انّه کالشّمس فی وسط
الزّوال لا ینکرها الّا کلّ مبغض مریب

\end{otherlanguage}

Yā Shaykh, qad asmaʿnāka tagharrudāt ʿandalīb al-Firdaws wa-araynāka
al-āthār allatī anzalahā Allāh bi-amrihi al-mubram fī al-Sijn al-Aʿẓam
li-taqarra bihā ʿaynuka wa-taṭmaʾinna bihā nafsuka. Innahu huwa
al-Fayyāḍ al-Karīm. Qum bi-quwwat al-burhān ʿalā khidmat amr Allāh
rabbika al-Raḥmān. Law takhāf min īmānika, khudh al-lawḥ thumma iḥfaẓhu
fī jayb tawakkulika. Wa-idhā dakhalta mawqif al-ḥashr wa-yasʾaluka Allāh
bi-ayyi ḥujja āmanta bi-hādhā al-ẓuhūr, akhrij al-lawḥ wa-qul: bi-hādhā
al-kitāb al-mubārak al-ʿazīz al-badīʿ. Idhan tartafiʿ ilayka ayādī
al-kull wa-yaʾkhudhūn al-lawḥ wa-yaḍaʿūnahu ʿalā ʿuyūnihim wa-yajidūn
minhu ʿarf bayān Allāh rabb al-ʿālamīn. Law yuʿadhdhibuka Allāh bi-mā
āmanta bi-āyātihi fī hādhā al-ẓuhūr, fa-bi-ayyi ḥujja yuʿadhdhib
alladhīna mā āmanū bi-Muḥammad Rasūl Allāh wa-min qablihi bi-ʿĪsā ibn
Maryam wa-min qablihi bi-l-Kalīm wa-min qablihi bi-l-Khalīl ilā an
tantahiya al-ẓuhūrāt ilā al-badīʿ al-awwal alladhī khuliqa bi-irādat
rabbika al-Qādir al-Muḥīṭ? Kadhālika anzalnā al-āyāt li-aḥad qablaka
wa-dhakarnāhā laka fī hādhā al-yawm li-taʿrifa wa-takūna min al-mūqinīn.
Yā ayyuhā al-nāṭiq bi-l-ʿilm, inna al-Amr aẓhar min an yakhfā wa-abyan
min an yustar. Innahu ka-l-shams fī wasaṭ al-zawāl, lā yunkiruhā illā
kull mubghiḍ murīb.

O Shaykh, We have let you hear the songs of the Nightingale of Paradise
and shown you the writings that God has sent down through His
irrevocable command in the Most Great Prison, so that your eyes may be
gladdened and your soul reassured by them. He is indeed the
All-Bountiful, the Generous. Arise, through the power of proof, to serve
the Cause of God, your Lord, the All-Merciful. If you fear the
consequences of your faith, take this Tablet and keep it in the fold of
your trust in God. When you enter the place of resurrection and God asks
you, ``On what proof did you believe in this Revelation?'' bring out the
Tablet and say, ``On this blessed, mighty, wondrous Book.'' Then the
hands of all will reach up toward you, and they will take the Tablet,
lay it upon their eyes, and perceive from it the fragrance of the
utterance of God, the Lord of all worlds. If God should punish you for
believing in His verses in this Revelation, on what grounds would He
punish those who did not believe in Muhammad, the Messenger of God, or,
before Him, in Jesus, son of Mary, or, before Him, in Him who conversed
with God, or, before Him, in the Friend of God, tracing the Revelations
back to the first creation brought into being by the will of your Lord,
the All-Powerful, the All-Encompassing? Thus We revealed the verses to
someone before you, and We have recounted them to you on this day so
that you may recognize the truth and be among those who are certain. O
you who speak of learning, the Cause is too manifest to be hidden and
too clear to be concealed. It is like the sun at high noon; none denies
it but one who is hateful and full of doubt.

\subsection{161}\label{section-1447}

\begin{otherlanguage}{persian}

در این حین سزاوار آنکه بحضرت مقصود توجّه نمائیم و باین کلمات عالیات متمسّک
شویم الهی الهی سراج امرت را بدهن حکمت برافروختی از اریاح مختلفه حفظش نما
سراج از تو زجاج از تو اسباب آسمان و زمین در قبضهٔ قدرت تو امرا را عدل
عنایت فرما و علما را انصاف توئی آن مقتدری که بحرکت قلم امر مبرمت را نصرت
فرمودی و اولیا را راه نمودی توئی مالک قدرت و ملیک اقتدار لا اله الّا انت
العزیز المختار و قل الهی الهی لک الحمد بما سقیتنی من ید عطآء اسمک القیّوم
رحیقک المختوم اسألک بأنوار فجر ظهورک و بنفوذ کلمتک العلیا و اقتدار قلمک
الأعلی الّذی بحرکته انجذبت حقائق الأشیآء ان توفّق حضرة السّلطان علی نصرة
امرک و الاقبال الی افق ظهورک و التّوجّه الی انوار وجهک ای ربّ ایّده علی ما
یقرّبه الیک ثمّ انصره بجنود السّموات و الأرض اسألک یا اله الأسمآء و فاطر
السّمآء بنور امرک و نار سدرة عنایتک ان تؤیّد حضرته علی اظهار امرک بین خلقک
ثمّ افتح علی وجهه ابواب فضلک و رحمتک و عطائک انّک انت المقتدر علی ما تشآء
بقولک کن فیکون

\end{otherlanguage}

Dar īn ḥīn sazāvār ānkih bi-Ḥaḍrat-i Maqṣūd tavajjuh namāyīm va bi-īn
kalimāt-i ʿāliyāt mutamassik shavīm: Ilāhī Ilāhī, sirāj-i amrat rā
bi-duhn-i ḥikmat bar-afrūkhtī, az aryāḥ-i mukhtalifih ḥifẓash namā.
Sirāj az tu, zujāj az tu, asbāb-i āsmān va zamīn dar qabḍiy-i qudrat-i
tu. Umarā rā ʿadl ʿināyat farmā va ʿulamā rā inṣāf. Tuʾī ān muqtadirī
kih bi-ḥarakat-i qalam amr-i mubramat rā nuṣrat farmūdī va awliyā rā rāh
namūdī. Tuʾī mālik-i qudrat va malīk-i iqtidār. Lā ilāha illā anta
al-ʿAzīz al-Mukhtār. Wa qul: Ilāhī Ilāhī, laka al-ḥamdu bi-mā saqaytanī
min yadi ʿaṭāʾi ismika al-Qayyūm raḥīqaka al-makhtūm. Asʾaluka bi-anwāri
fajri ẓuhūrika wa bi-nufūdhi kalimatika al-ʿulyā wa iqtidāri qalamika
al-aʿlā alladhī bi-ḥarakatihi injadhabat ḥaqāʾiqu al-ashyāʾ an tuwaffiqa
ḥaḍrata al-sulṭān ʿalā nuṣrati amrika wa al-iqbāli ilā ufuqi ẓuhūrika wa
al-tawajjuhi ilā anwāri wajhika. Ay Rabb, ayyidhu ʿalā mā yuqarribuhu
ilayka thumma unṣurhu bi-junūdi al-samāwāti wa al-arḍ. Asʾaluka yā Ilāha
al-asmāʾ wa Fāṭira al-samāʾ bi-nūri amrika wa nāri sidrati ʿināyatika an
tuʾayyida ḥaḍratahu ʿalā iẓhāri amrika bayna khalqika thumma iftaḥ ʿalā
wajhihi abwāba faḍlika wa raḥmatika wa ʿaṭāʾika. Innaka anta
al-muqtadiru ʿalā mā tashāʾ bi-qawlika kun fa-yakūn.

At this moment it is fitting that we turn toward Him who is the Desired
One and hold fast to these exalted words: ``My God, my God! You have
kindled the lamp of Your Cause with the oil of wisdom; protect it from
contrary winds. The lamp is Yours, the glass is Yours, and the means of
heaven and earth are in the grasp of Your power. Grant justice to rulers
and fairness to the learned. You are the All-Powerful who, through the
movement of the Pen, has aided Your irrevocable Cause and shown the way
to Your loved ones. You are the Possessor of power and the King of
might. There is no god but You, the Almighty, the Unconstrained.'' And
say: ``My God, my God! Praise be to You for giving me Your sealed wine
to drink from the hand of the bounty of Your name, the Self-Subsisting.
I ask You, by the lights of the dawn of Your Revelation, by the potency
of Your most exalted Word, and by the might of Your Most Exalted Pen,
through whose movement the realities of all things have been enraptured,
to enable His Majesty the Sultan to aid Your Cause, turn toward the
horizon of Your Revelation, and direct himself toward the lights of Your
face. O Lord, assist him in whatever will draw him near to You, and then
aid him with the hosts of the heavens and the earth. I ask You, O God of
names and Creator of heaven, by the light of Your Cause and the fire of
the Lote Tree of Your loving care, to assist His Majesty in making Your
Cause manifest among Your creatures. Then open before his face the gates
of Your grace, Your mercy, and Your bounty. You indeed have power over
whatever You will through Your word, `Be,' and it is.''

\subsection{162}\label{section-1448}

\begin{otherlanguage}{persian}

یا شیخ زمام امور بقوّت الهی و قدرت ربّانی اخذ شده بود اخذ عزیز مقتدر احدی
قادر بر فساد و فتنه نبوده حال چون قدر عنایت و الطاف را ندانستند بجزای
اعمال خود مبتلا شده و میشوند و مأمورین نظر بحرکت سرّی حبل ممدود از جمیع
جهات حزب مخالف را تحریک و تأیید نموده‌اند در مدینهٔ کبیره جمعی را بر
مخالفت این مظلوم برانگیختند و امر بمقامی رسیده که نفوس مأموره در آن ارض
تمسّک نموده‌اند بآنچه سبب ذلّت دولت و ملّت است یکی از سادات بزرگوار که نزد
اکثری از منصفین معروفش معروف و عملش مقبول و تجارتش مشهور و از اجلّهٔ تجّار
نزد کل مشهود بسمت بیروت توجّه نمود نظر بدوستی آن جناب باین مظلوم از سیّالهٔ
برقیّه بمترجم ایران خبر دادند که سیّد مذکور با آدم خود وجهی از نقد و غیره
سرقت کرده‌اند و به عکّا توجّه نموده‌اند و مقصود از این حرکت ذلّت این مظلوم
بوده ولکن هیهات که اهل این دیار باین گفتهای نالائقه از صراط مستقیم صدق و
راستی منحرف شوند باری از هر جهت هجوم نموده‌اند و طرف مقابل را تأیید
مینمایند و این مظلوم از حقّ میطلبد کل را موفّق دارد بر آنچه سزاوار ایّام
است و در لیالی و ایّام باین کلمات محکمات ناظر و ناطق الهی الهی اسألک بشمس
فضلک و بحر علمک و سمآء عدلک ان تؤیّد المنکرین على الاقرار و المعرضین علی
الاقبال و المفترین علی العدل و الانصاف ای ربّ ایّدهم علی الرّجوع الیک و
الانابة لدی باب فضلک انّک انت المقتدر علی ما تشآء و فی قبضتک زمام من فی
السّموات و الأرضین الحمد للّه ربّ العالمین

\end{otherlanguage}

Yā Shaykh, zimām-i umūr bi-quvvat-i ilāhī va qudrat-i rabbānī akhdh
shudih būd, akhdh-i ʿazīz-i muqtadir. Aḥadī qādir bar fasād va fitnih
nabūdih. Ḥāl chūn qadr-i ʿināyat va alṭāf rā nadānistand bi-jazā-yi
aʿmāl-i khud mubtalā shudih va mī-shavand va maʾmūrīn naẓar bi-ḥarakat-i
sirrī-yi ḥabl-i mamdūd az jamīʿ-i jihāt ḥizb-i mukhālif rā taḥrīk va
taʾyīd namūdih-and. Dar Madīniy-i Kabīrih jamʿī rā bar mukhālifat-i īn
maẓlūm bar-angīkhtand va amr bi-maqāmī rasīdih kih nufūs-i maʾmūrih dar
ān arḍ tamassuk namūdih-and bi-ānchih sabab-i dhillat-i dawlat va millat
ast. Yikī az sādāt-i buzurgvār kih nazd-i aktharī az munṣifīn maʿrūfash
maʿrūf va ʿamalash maqbūl va tijāratash mashhūr va az ajilliy-i tujjār
nazd-i kull mashhūd bi-samt-i Bayrūt tavajjuh namūd. Naẓar bi-dūstī-yi
ān jināb bi-īn maẓlūm az sayyāliy-i barqīyih bi-mutarjim-i Īrān khabar
dādand kih sayyid-i madhkūr bā ādam-i khud vajhī az naqd va ghayrih
sirqat kardih-and va bi-ʿAkkā tavajjuh namūdih-and va maqṣūd az īn
ḥarakat dhillat-i īn maẓlūm būdih. Va-līkin hayhāt kih ahl-i īn diyār
bi-īn guftihā-yi nālāʾiqih az ṣirāṭ-i mustaqīm-i ṣidq va rāstī munḥarif
shavand. Bārī az har jihat hujūm namūdih-and va ṭaraf-i muqābil rā
taʾyīd mī-namāyand va īn maẓlūm az Ḥaqq mī-ṭalabad kull rā muvaffaq
dārad bar ānchih sazāvār-i ayyām ast va dar layālī va ayyām bi-īn
kalimāt-i muḥkamāt nāẓir va nāṭiq: Ilāhī Ilāhī, asʾaluka bi-shamsi
faḍlika wa baḥri ʿilmika wa samāʾi ʿadlika an tuʾayyida al-munkirīna
ʿalā al-iqrār wa al-muʿriḍīna ʿalā al-iqbāl wa al-muftarīna ʿalā
al-ʿadli wa al-inṣāf. Ay Rabb, ayyidhum ʿalā al-rujūʿi ilayka wa
al-inābati ladā bābi faḍlika. Innaka anta al-muqtadiru ʿalā mā tashāʾ wa
fī qabḍatika zimāmu man fī al-samāwāti wa al-arḍīn. Al-ḥamdu li-llāhi
Rabbi al-ʿālamīn.

O Shaykh, the reins of affairs had been seized through divine strength
and heavenly power, with the grasp of One who is mighty and
all-powerful. No one could stir up corruption or sedition. Now, because
they failed to appreciate this loving care and these favors, they have
suffered, and continue to suffer, the consequences of their deeds. In
view of the secret activity of the outstretched cord, the officials have
incited and supported the opposing party on every side. In the Great
City they stirred up a group to oppose this Wronged One, and matters
have reached the point where the officials in that land have resorted to
what brings disgrace upon both state and people. One of the noble
descendants of the Prophet, whose benevolence is recognized by most
fair-minded people, whose conduct is approved, whose commerce is well
known, and whom everyone regards as one of the foremost merchants, set
out for Beirut. Because of his friendship with this Wronged One, they
sent word by telegraph to the Persian interpreter that the
aforementioned sayyid and his servant had stolen a sum in cash and other
property and had set out for Acre. The purpose of this action was to
disgrace this Wronged One. Yet far be it from the people of this region
to be turned aside by such unworthy statements from the straight path of
truthfulness and integrity. In short, they have attacked from every
direction and are supporting the opposing side. This Wronged One asks
God to enable everyone to do what befits these days, and night and day
He contemplates and utters these firmly established words: ``My God, my
God! I ask You, by the sun of Your grace, the ocean of Your knowledge,
and the heaven of Your justice, to help those who deny to acknowledge,
those who turn away to turn toward You, and those who slander to act
with justice and fairness. O Lord, assist them to return to You and turn
in repentance at the gate of Your grace. You indeed have power over
whatever You will, and in Your grasp are the reins of all who are in the
heavens and the earth. Praise be to God, the Lord of all worlds.''

\subsection{163}\label{section-1449}

\begin{otherlanguage}{persian}

عنقریب آنچه در افئده و قلوب مستور مشاهده گردد یوم یومیست که حضرت لقمان
از برای ابنش ذکر فرموده و ربّ العزّة از آن خبر داده و حبیبش را آگاه نموده
بقوله تعالی یا بنیّ انّها ان تک مثقال حبّة من خردل فتکن فی صخرة او فی
السّموات او فی الأرض یأت بها اللّه انّ اللّه لطیف خبیر امروز خائنهٔ اعین و
خافیهٔ صدور کل امام کرسیّ ظهور ظاهر و هویدا لا یعزب عن علمه من شیء یسمع و
یری و هو السّمیع البصیر بسیار عجبست که امین و خائن را از هم فرقی نه

\end{otherlanguage}

ʿAn-qarīb ānchih dar afʾidih va qulūb mastūr mushāhadih gardad. Yawm
yawmīst kih Ḥaḍrat-i Luqmān az barā-yi ibnash dhikr farmūdih va Rabb
al-ʿizzat az ān khabar dādih va Ḥabībash rā āgāh namūdih bi-qawlihi
taʿālā: Yā bunayya innahā in taku mithqāla ḥabbatin min khardalin
fa-takun fī ṣakhratin aw fī al-samāwāti aw fī al-arḍi yaʾti bihā Allāh,
inna Allāha laṭīfun khabīr. Imrūz khāʾiniy-i aʿyun va khāfiyiy-i ṣudūr
kull imām-i kursī-yi ẓuhūr ẓāhir va huvaydā. Lā yaʿzubu ʿan ʿilmihi min
shayʾ, yasmaʿu wa yarā wa huwa al-Samīʿ al-Baṣīr. Bisyār ʿajab ast kih
amīn va khāʾin rā az ham farqī nah.

Soon whatever is hidden in hearts and souls will be seen. This is the
day of which Luqman spoke to his son, of which the Lord of Glory gave
tidings, and of which He informed His Beloved through His exalted words:
``O my son, even if it is the weight of a mustard seed, and lies within
a rock or in the heavens or the earth, God will bring it forth. God is
indeed the Subtle, the All-Informed.'' Today the treachery of eyes and
the secrets of hearts are all manifest and plain before the throne of
Revelation. Nothing escapes His knowledge. He hears and sees, and He is
the All-Hearing, the All-Seeing. It is most astonishing that no
distinction is made between the trustworthy and the treacherous.

\subsection{164}\label{section-1450}

\begin{otherlanguage}{persian}

ایکاش حضرت پادشاه ممالک ایران ادام اللّه سلطانه از قناسل دولت علیّهٔ ایران
که در این جهات بوده‌اند استفسار فرمایند تا بر افعال و اعمال این مظلوم
آگاه شوند باری جمعی را از اختر و غیره برانگیختند و در انتشار مفتریات
مشغول این بسی واضح و معلوم نفسی را که مردود عباد و مطرود بلاد مشاهده
نمایند با سیف ضغینه و سهام بغضا احاطه کنند لیس هذا اوّل امر ظهر بالظّلم و
لا اوّل قارورة کسرت و لا اوّل ستر هتک فی سبیل اللّه ربّ العالمین و این مظلوم
ساکتاً صامتاً در سجن اعظم بخود مشغول و از غیر اللّه منقطع ظلم بمقامی رسیده
که اقلام عالم از تحریر آن عاجز و قاصر است

\end{otherlanguage}

Ay kāsh Ḥaḍrat-i Pādshāh-i mamālik-i Īrān, adāma Allāhu sulṭānahu, az
qanāsil-i dawlat-i ʿaliyyih-yi Īrān kih dar īn jihāt būdih-and istifsār
farmāyand tā bar afʿāl va aʿmāl-i īn maẓlūm āgāh shavand. Bārī jamʿī rā
az Akhtar va ghayrih bar-angīkhtand va dar intishār-i muftariyāt
mashghūl. Īn basī vāḍiḥ va maʿlūm, nafsī rā kih mardūd-i ʿibād va
maṭrūd-i bilād mushāhadih namāyand bā sayf-i ḍaghīnih va sihām-i baghḍā
iḥāṭih kunand. Laysa hādhā awwala amrin ẓahara bi-l-ẓulmi wa lā awwala
qārūratin kusirat wa lā awwala sitrin hutika fī sabīli Allāhi Rabbi
al-ʿālamīn. Va īn maẓlūm sākitan ṣāmitan dar Sijn-i Aʿẓam bi-khud
mashghūl va az ghayr-i Allāh munqaṭiʿ. Ẓulm bi-maqāmī rasīdih kih
aqlām-i ʿālam az taḥrīr-i ān ʿājiz va qāṣir ast.

If only His Majesty the sovereign of the realms of Persia---may God
perpetuate his sovereignty---would question the consuls of the exalted
Persian government who have been in these regions, so that he might
learn of the deeds and conduct of this Wronged One. In short, they have
stirred up a group from Akhtar and elsewhere, who are busy spreading
slanders. It is quite clear and well known that when they see someone
rejected by the people and banished from every land, they surround him
with the sword of malice and the arrows of hatred. This is not the first
matter in which oppression has appeared, nor the first vessel to be
broken, nor the first veil to be torn in the path of God, the Lord of
all worlds. This Wronged One, silent and still, attends to His own
affairs in the Most Great Prison, detached from all but God. Oppression
has reached such a degree that the pens of the world are powerless and
unable to describe it.

\subsection{165}\label{section-1451}

\begin{otherlanguage}{persian}

در این مقام لازم امری که ظاهر شد ذکر شود شاید عباد بحبل عدل و راستی تمسّک
نمایند جناب حاجی شیخ محمّد علی علیه بهآء اللّه الأبدی از تجّار معروف بوده
اکثری از اهل مدینهٔ کبیره او را میشناسند در ایاّم اخیره که سفارت ایران در
آستانه در سرّ سرّ بتحریک مشغول آن مقبل صادق را پریشان دیده‌اند تا آنکه شبی
از شبها خود را در بحر انداخت و از قضا جمعی از عباد حاضر و او را اخذ
نمودند و این عمل را هر حزبی تعبیری نمود و ذکری کرد تا آنکه مرّة اخری شبی
از شبها در جامعی رفته و خادم آن محل ذکر نمود این شخص شب را احیا داشت و
تا صبح بمناجات و دعا و عجز و ابتهال مشغول و بعد ذکرش قطع شد و این عبد
توجّه نمود مشاهده شد روح را تسلیم نموده و شیشه‌ئی خالی نزدش دیده شد مشعر
بر اینکه سم خورده باری بکمال تحیّر قوم را اطّلاع داد و دو وصیّت‌نامه از او
ظاهر شد امّا اوّل مشعر بر اقرار و اعتراف بر وحدانیّت حقّ و تقدیس ذاته تعالی
عن الأشباه و الأمثال و تنزیه کینونته عن الأوصاف و الأذکار و الأقوال و
الاقرار بظهور الأنبیآء و الأولیآء و الاعتراف بما کان مرقوماً فی کتب اللّه
مولی الوری و در ورقهٔ دیگر مناجاتی عرض کرده و در آخر ورقه ذکر نموده این
عبد و اولیا متحیّر مانده‌اند چه که در یک مقام قلم اعلی کل را از فساد و
نزاع و جدال منع فرموده و در مقامی هم از قلم اعلی این کلمهٔ علیا نازل اگر
نفسی سوء قصدی از احدی در حضور مشاهده نماید باید تعرّض نکند و بحقّ گذارد
این حکم محکم از یک جهت ظاهر و ثابت و از جهت دیگر مشاهده شد مفترین
بکلماتی نطق مینمایند که قوّهٔ بشری از حمل و اصغاء آن عاجز و قاصر است لذا
این عبد این ذنب اعظم را اختیار نمود و از بحر کرم الهی و سماء رحمت ربّانی
سائل و آملم که جریرات این عبد را از قلم فضل و عطا محو فرماید سیّئات بسیار
و خطایا بیشمار ولکن بحبل جودش متمسّکم و بذیل کرمش متشبّث حقّ شاهد و مقرّبان
درگاه آگاه که این عبد قادر بر اصغاء مقالات مغلّین نبوده لذا این عمل را
ارتکاب نمود لو یعذّبنی انّه هو محمود فی فعله و لو یغفر لی انّه مطاع فی امره

\end{otherlanguage}

Dar īn maqām lāzim amrī kih ẓāhir shud dhikr shavad shāyad ʿibād
bi-ḥabl-i ʿadl va rāstī tamassuk namāyand. Jināb-i Ḥājī Shaykh Muḥammad
ʿAlī, ʿalayhi bahāʾu Allāhi al-Abadī, az tujjār-i maʿrūf būdih, aktharī
az ahl-i Madīniy-i Kabīrih ū rā mī-shināsand. Dar ayyām-i akhīrih kih
sifārat-i Īrān dar Āstānih dar sirr-i sirr bi-taḥrīk mashghūl ān
muqbil-i ṣādiq rā parishān dīdih-and tā ānkih shabī az shabhā khud rā
dar baḥr andākht va az qaḍā jamʿī az ʿibād ḥāḍir va ū rā akhdh namūdand
va īn ʿamal rā har ḥizbī taʿbīrī namūd va dhikrī kard tā ānkih marratan
ukhrā shabī az shabhā dar jāmiʿī raftih va khādim-i ān maḥall dhikr
namūd: Īn shakhs shab rā iḥyā dāsht va tā ṣubḥ bi-munājāt va duʿā va
ʿajz va ibtihāl mashghūl va baʿd dhikrash qaṭʿ shud va īn ʿabd tavajjuh
namūd, mushāhadih shud rūḥ rā taslīm namūdih va shīshih-ʾī khālī nazdash
dīdih shud mushʿir bar īnkih samm khūrdih. Bārī bi-kamāl-i taḥayyur qawm
rā iṭṭilāʿ dād va du vaṣiyyat-nāmih az ū ẓāhir shud. Ammā avval mushʿir
bar iqrār va iʿtirāf bar vaḥdāniyyat-i Ḥaqq va taqdīs-i dhātihi taʿālā
ʿan al-ashbāhi wa al-amthāl va tanzīh-i kaynūnatihi ʿan al-awṣāfi wa
al-adhkāri wa al-aqwāl wa al-iqrāri bi-ẓuhūri al-anbiyāʾi wa al-awliyāʾ
wa al-iʿtirāfi bi-mā kāna marqūman fī kutubi Allāhi Mawlā al-warā. Va
dar varaqiy-i dīgar munājātī ʿarḍ kardih va dar ākhir-i varaqih dhikr
namūdih: Īn ʿabd va awliyā mutaḥayyir māndih-and chih kih dar yik maqām
Qalam-i Aʿlā kull rā az fasād va nizāʿ va jidāl manʿ farmūdih va dar
maqāmī ham az Qalam-i Aʿlā īn kalimiy-i ʿulyā nāzil: Agar nafsī sūʾ-i
qaṣdī az aḥadī dar ḥuḍūr mushāhadih namāyad bāyad taʿarruḍ nakunad va
bi-Ḥaqq gudhārad. Īn ḥukm-i muḥkam az yik jihat ẓāhir va thābit va az
jihat-i dīgar mushāhadih shud muftarīn bi-kalimātī nuṭq mī-namāyand kih
quvviy-i bashar az ḥaml va iṣghā-yi ān ʿājiz va qāṣir ast. Lidhā īn ʿabd
īn dhanb-i aʿẓam rā ikhtiyār namūd va az baḥr-i karam-i ilāhī va samā-yi
raḥmat-i rabbānī sāʾil va āmilam kih jarīrāt-i īn ʿabd rā az qalam-i
faḍl va ʿaṭā maḥv farmāyad. Sayyiʾāt bisyār va khaṭāyā bī-shumār
va-līkin bi-ḥabl-i jūdash mutamassikam va bi-dhayl-i karamash
mutashabbith. Ḥaqq shāhid va muqarrabān-i dargāh āgāh kih īn ʿabd qādir
bar iṣghā-yi maqālāt-i mughillīn nabūdih lidhā īn ʿamal rā irtikāb
namūd. Law yuʿadhdhibnī innahu huwa maḥmūdun fī fiʿlihi wa law yaghfir
lī innahu muṭāʿun fī amrihi.

At this point it is necessary to recount an event that occurred, so that
perhaps the people may hold fast to the cord of justice and
truthfulness. Haji Shaykh Muhammad Ali---may the glory of God, the
Eternal, rest upon him---was a well-known merchant whom most of the
people of the Great City know. In recent days, while the Persian embassy
in Constantinople was engaged in the utmost secrecy in stirring up
trouble, that sincere believer was seen to be distressed, until one
night he threw himself into the sea. As it happened, some people were
present and rescued him. Every group interpreted this deed in its own
way and spoke of it, until once again, one night, he went into a mosque.
The attendant there said: ``This man kept vigil through the night and
was occupied until morning with communion, prayer, humble supplication,
and entreaty. Then his utterances ceased, and this servant went over and
saw that he had surrendered his soul. An empty bottle was seen beside
him, indicating that he had taken poison.'' In short, utterly
bewildered, he informed the people, and two wills were found that he had
left. The first expressed acknowledgment and confession of the oneness
of God, the sanctity of His exalted Essence above likenesses and
comparisons, and the transcendence of His Being above descriptions,
mentions, and words; it acknowledged the appearance of the Prophets and
holy ones and affirmed what was recorded in the Books of God, the Lord
of humanity. In the other sheet he had offered a prayer, and at the end
of the sheet he had written: ``This servant and the loved ones have been
left bewildered, for in one passage the Most Exalted Pen has forbidden
everyone to engage in corruption, contention, and dispute, and in
another this exalted utterance has been revealed by the Most Exalted
Pen: `If anyone should witness an attempt by someone to do harm in His
presence, he must not intervene but must leave the matter to God.' On
the one hand, this firm command is manifest and established; on the
other, the slanderers have been seen to utter words that human strength
is powerless and unable to bear or hear. This servant has therefore
chosen this most grievous sin, and I ask and hope, from the ocean of
divine generosity and the heaven of heavenly mercy, that He may efface
this servant's transgressions with the pen of grace and bounty. My evil
deeds are many and my sins innumerable, yet I hold fast to the cord of
His bounty and cling to the hem of His generosity. God is witness, and
those near His threshold know, that this servant was unable to listen to
the statements of the malicious, and therefore committed this deed. If
He punishes me, He is indeed praiseworthy in His action, and if He
forgives me, He is indeed to be obeyed in His command.''

\subsection{166}\label{section-1452}

\begin{otherlanguage}{persian}

حال جناب شیخ در نفوذ کلمه تفکّر نماید شاید از شمال وهم بیمین یقین توجّه
کند این مظلوم در امر الهی با احدی مداهنه ننموده و کلمهٔ حقّ را امام وجوه
خلق بأعلی النّدآء ذکر کرده من شآء فلیقبل و من شآء فلیعرض ولکن اگر این
امور ظاهرهٔ واضحهٔ مشهوده انکار شود چه امری نزد متبصّرین بطراز قبول و اقرار
مزیّن گردد انّا نسأل اللّه تبارک و تعالی ان یغفر للمذکور و یبدّل سیّئاته
بالحسنات انّه هو المقتدر العزیز الوهّاب

\end{otherlanguage}

Ḥāl Jināb-i Shaykh dar nufūdh-i kalimih tafakkur namāyad shāyad az
shimāl-i vahm bi-yamīn-i yaqīn tavajjuh kunad. Īn maẓlūm dar amr-i ilāhī
bā aḥadī mudāhanih nanamūdih va kalimiy-i ḥaqq rā imām-i vujūh-i khalq
bi-aʿlā al-nidāʾ dhikr kardih. Man shāʾa fa-l-yuqbil wa man shāʾa
fa-l-yuʿriḍ. Va-līkin agar īn umūr-i ẓāhiriy-i vāḍiḥiy-i mashhūdih inkār
shavad chih amrī nazd-i mutabaṣṣirīn bi-ṭirāz-i qabūl va iqrār muzayyan
gardad? Innā nasʾalu Allāha tabāraka wa taʿālā an yaghfira li-l-madhkūr
wa yubaddila sayyiʾātihi bi-l-ḥasanāt. Innahu huwa al-Muqtadir al-ʿAzīz
al-Wahhāb.

Let the Shaykh now reflect upon the potency of the Word, so that perhaps
he may turn from the left hand of delusion to the right hand of
certainty. In the Cause of God this Wronged One has compromised with no
one, and has proclaimed the word of truth at the top of His voice before
the faces of all people. Let whoever wishes turn toward it, and whoever
wishes turn away. But if these manifest, clear, and witnessed events are
denied, what matter could be adorned with acceptance and acknowledgment
in the eyes of the discerning? We ask God, blessed and exalted, to
forgive the aforementioned man and to change his evil deeds into good
ones. He is indeed the All-Powerful, the Almighty, the All-Bountiful.

\subsection{167}\label{section-1453}

\begin{otherlanguage}{persian}

اموری در این ظهور ظاهر که از برای مظاهر علم و فضل و مشارق انصاف و عدل جز
تصدیق مجالی نه امروز بر شما لازم و واجب که بقوّت ملکوتی قیام نمائی و
شبهات احزاب عالم را بقدرت علم محو فرمائی تا کل مقدّس شده قصد بحر اعظم
نمایند و به ما اراده اللّه تمسّک جویند

\end{otherlanguage}

Umūrī dar īn ẓuhūr ẓāhir kih az barā-yi maẓāhir-i ʿilm va faḍl va
mashāriq-i inṣāf va ʿadl juz taṣdīq majālī nah. Imrūz bar shumā lāzim va
vājib kih bi-quvvat-i malakūtī qiyām namāʾī va shubuhāt-i aḥzāb-i ʿālam
rā bi-qudrat-i ʿilm maḥv farmāʾī tā kull muqaddas shudih qaṣd-i Baḥr-i
Aʿẓam namāyand va bi-mā arāda Allāh tamassuk jūyand.

Things have appeared in this Revelation that leave those who embody
knowledge and learning, and those who are dawning places of fairness and
justice, no room for anything but affirmation. Today it is your
necessary and binding duty to arise with heavenly strength and, through
the power of knowledge, dispel the doubts of the peoples of the world,
so that all may be purified, set their course toward the Most Great
Ocean, and hold fast to what God has willed.

\subsection{168}\label{section-1454}

\begin{otherlanguage}{persian}

باری هر معرضی بکلمه‌ئی تمسّک نموده و بر حقّ اعتراض کرده سبحان‌اللّه ذکر
الوهیّت و ربوبیّت که از اولیا و اصفیا ظاهر شده آن را اسباب اعراض و انکار
قرار داده‌اند حضرت صادق فرموده العبودیّة جوهرة کنهها الرّبوبیّة و حضرت امیر
در جواب اعرابی که از نفس سؤال نموده فرموده و ثالثها اللّاهوتیّة الملکوتیّة
و هی قوّة لاهوتیّة و جوهرة بسیطة حیّة بالذّات الی ان قال علیه السّلام فهی ذات
اللّه العلیا و شجرة طوبی و سدرة المنتهی و جنّة المأوی حضرت صادق فرموده اذا
قام قائمنا اشرقت الأرض بنور ربّها و همچنین از ابی عبداللّه علیه السّلام
حدیث طویلی ذکر شده و از جملهٔ نصوص حدیث مذکور این کلمهٔ علیا است فعند ذلک
یهبط الجبّار عزّ و جلّ من الغمام و الملائکة و فی القرآن العظیم هل ینظرون
الّا ان یأتیهم اللّه فی ظلل من الغمام و در حدیث مفضّل میفرماید یسند القائم
ظهره الی الحرم و یمدّ یده المبارکة فتری بیضآء من غیر سوء و یقول هذه ید
اللّه و یمین اللّه و عن اللّه و بأمر اللّه هر نحو این احادیث را معنی
نموده‌اند آثار قلم اعلی را هم معنی نمایند حضرت امیر فرموده انا الّذی لا
یقع علیه اسم و لا صفة و همچنین فرموده ظاهری امامة و باطنی غیب لا یدرک
قال ابو جعفر الطّوسی قلت لأبی عبداللّه انتم الصّراط فی کتاب اللّه و انتم
الزّکوة و انتم الحجّ قال یا فلان نحن الصّراط فی کتاب اللّه عزّ و جلّ و نحن
الزّکوة و نحن الصّیام و نحن الحجّ و نحن الشّهر الحرام و نحن البلد الحرام و
نحن کعبة اللّه و نحن قبلة اللّه و نحن وجه اللّه روی جابر عن ابی جعفر علیه
السّلام قال یا جابر علیک بالبیان و المعانی فقال علیه السّلام امّا البیان و
هو ان تعرف اللّه سبحانه لیس کمثله شیء فتعبده و لا تشرک به شیئاً و امّا
المعانی فنحن معانیه و نحن جنبه و یده و لسانه و امره و حکمه و علمه و حقّه
اذا شئنا شآء اللّه و یرید ما نریده ایضاً حضرت امیر علیه السّلام فرمود کیف
اعبد ربّاً لم اره و در مقام دیگر میفرماید ما رأیت شیئاً الّا و قد رأیت اللّه
قبله او بعده او معه

\end{otherlanguage}

Bārī har muʿriḍī bi-kalimih-ī tamassuk namūdih va bar Ḥaqq iʿtirāḍ
kardih. Subḥāna Allāh! Dhikr-i ulūhiyyat va rubūbiyyat kih az awliyā va
aṣfiyā ẓāhir shudih ān rā asbāb-i iʿrāḍ va inkār qarār dādih-and.
Ḥaḍrat-i Ṣādiq farmūdih: al-ʿubūdiyyatu jawharatun kunhuhā
al-rubūbiyyah. Va Ḥaḍrat-i Amīr dar javāb-i aʿrābī kih az nafs suʾāl
namūdih farmūdih: wa-thālithuhā al-lāhūtiyyatu al-malakūtiyyatu wa-hiya
quwwatun lāhūtiyyatun wa-jawharatun basīṭatun ḥayyatun bi-al-dhāt. Ilā
an qāla ʿalayhi al-salām: fa-hiya dhātu Allāhi al-ʿulyā wa-shajaratu
Ṭūbā wa-sidratu al-muntahā wa-jannatu al-maʾwā. Ḥaḍrat-i Ṣādiq farmūdih:
idhā qāma qāʾimunā ashraqati al-arḍu bi-nūri rabbihā. Va hamchunīn az
Abī ʿAbdillāh ʿalayhi al-salām ḥadīth-i ṭavīlī dhikr shudih va az
jumlih-yi nuṣūṣ-i ḥadīth-i madhkūr īn kalimih-yi ʿulyā ast: fa-ʿinda
dhālika yahbiṭu al-Jabbāru ʿazza wa-jalla mina al-ghamāmi
wa-al-malāʾikah. Wa-fī al-Qurʾāni al-ʿaẓīm: hal yanẓurūna illā an
yaʾtiyahumu Allāhu fī ẓulalin mina al-ghamām. Va dar ḥadīth-i Mufaḍḍal
mīfarmāyad: yusnidu al-Qāʾimu ẓahrahu ilā al-ḥarami wa-yamuddu yadahu
al-mubārakata fa-turā bayḍāʾa min ghayri sūʾin wa-yaqūlu hādhihi yadu
Allāhi wa-yamīnu Allāhi wa-ʿan Allāhi wa-bi-amri Allāh. Har naḥv īn
aḥādīth rā maʿnī namūdih-and āthār-i Qalam-i Aʿlā rā ham maʿnī namāyand.
Ḥaḍrat-i Amīr farmūdih: anā alladhī lā yaqaʿu ʿalayhi ismun wa-lā ṣifah.
Va hamchunīn farmūdih: ẓāhirī imāmatun wa-bāṭinī ghaybun lā yudrak. Qāla
Abū Jaʿfar al-Ṭūsī: qultu li-Abī ʿAbdillāh antumu al-ṣirāṭu fī kitābi
Allāhi wa-antumu al-zakātu wa-antumu al-ḥajj. Qāla yā fulānu naḥnu
al-ṣirāṭu fī kitābi Allāhi ʿazza wa-jalla wa-naḥnu al-zakātu wa-naḥnu
al-ṣiyāmu wa-naḥnu al-ḥajju wa-naḥnu al-shahru al-ḥarāmu wa-naḥnu
al-baladu al-ḥarāmu wa-naḥnu kaʿbatu Allāhi wa-naḥnu qiblatu Allāhi
wa-naḥnu wajhu Allāh. Rawā Jābirun ʿan Abī Jaʿfar ʿalayhi al-salām qāla:
yā Jābiru ʿalayka bi-al-bayāni wa-al-maʿānī. Fa-qāla ʿalayhi al-salām:
ammā al-bayānu wa-huwa an taʿrifa Allāha subḥānahu laysa ka-mithlihi
shayʾun fa-taʿbudahu wa-lā tushrika bihi shayʾan wa-ammā al-maʿānī
fa-naḥnu maʿānīhi wa-naḥnu janbuhu wa-yaduhu wa-lisānuhu wa-amruhu
wa-ḥukmuhu wa-ʿilmuhu wa-ḥaqquhu idhā shiʾnā shāʾa Allāhu wa-yurīdu mā
nurīdu. Ayḍan Ḥaḍrat-i Amīr ʿalayhi al-salām farmūd: kayfa aʿbudu rabban
lam arahu. Va dar maqām-i dīgar mīfarmāyad: mā raʾaytu shayʾan illā
wa-qad raʾaytu Allāha qablahu aw baʿdahu aw maʿahu.

In short, everyone who has turned away has seized upon some word and
raised objections against God. Glory be to God! They have made the
references to divinity and lordship that have issued from the saints and
chosen ones grounds for rejection and denial. His Holiness Ṣādiq has
said: ``Servitude is a jewel whose essence is lordship.'' And His
Holiness the Commander, in answer to a Bedouin who had asked about the
soul, said: ``Its third is the divine, heavenly soul, which is a divine
power and a simple substance, living in itself.'' He continued, peace be
upon him, until he said: ``It is the exalted Self of God, the Tree of
Ṭūbā, the Lote Tree of the Utmost Boundary, and the Garden of Refuge.''
His Holiness Ṣādiq has said: ``When our Qāʾim arises, the earth will
shine with the light of its Lord.'' A lengthy tradition has also been
related from Abū ʿAbdillāh, peace be upon him, and among the explicit
statements of that tradition is this exalted utterance: ``At that moment
the All-Compelling, mighty and glorious is He, will descend from the
clouds, and the angels.'' And in the mighty Qurʾān: ``Are they waiting
for anything but that God should come to them in the shadows of the
clouds?'' And in the tradition of Mufaḍḍal he says: ``The Qāʾim will
lean his back against the Sanctuary and extend his blessed hand, and it
will be seen to be white, without blemish. He will say: `This is the
hand of God, the right hand of God, from God and by the command of
God.'\,'' In whatever way they have interpreted these traditions, let
them also interpret the writings of the Most Exalted Pen. His Holiness
the Commander has said: ``I am the one to whom no name or attribute
applies.'' He has also said: ``My outward reality is the Imamate, and my
inward reality is an unseen mystery that cannot be apprehended.'' Abū
Jaʿfar al-Ṭūsī said: ``I said to Abū ʿAbdillāh: `You are the Path in the
Book of God, and you are the alms and you are the pilgrimage.' He said:
`O so-and-so! We are the Path in the Book of God, mighty and glorious is
He; we are the alms, we are fasting, we are the pilgrimage, we are the
sacred month, we are the sacred city, we are the Kaaba of God, we are
the Point of Adoration of God, and we are the Face of God.'\,'' Jābir
related that Abū Jaʿfar, peace be upon him, said: ``O Jābir! Hold fast
to the exposition and the meanings.'' He then said, peace be upon him:
``As for the exposition, it is that you should recognize God, glorified
be He---nothing is like Him---and worship Him and associate nothing with
Him. As for the meanings, we are His meanings, and we are His side, His
hand, His tongue, His command, His judgment, His knowledge, and His
truth. When we will, God wills, and He desires what we desire.'' His
Holiness the Commander, peace be upon him, also said: ``How could I
worship a Lord whom I have not seen?'' And elsewhere he says: ``I have
seen nothing without seeing God before it, or after it, or with it.''

\subsection{169}\label{section-1455}

\begin{otherlanguage}{persian}

یا شیخ در آنچه ذکر شده تفکّر فرما شاید بقوّت اسم قیّوم از رحیق مختوم
بیاشامی و بیابی آنچه را که کل از ادراکش عاجزند کمر همّت را محکم نما و قصد
ملکوت اعلی کن شاید در حین تنزیل نفحات وحی و الهام را بیابی و بآن فائز
شوی براستی میگویم از برای امر الهی شبه و مثلی نبوده و نیست حجبات اوهام
را خرق نما انّه یمدّک و یؤیّدک فضلاً من عنده و هو القویّ الغالب القدیر تا وقت
باقی و سدرهٔ مبارکه مابین بریّه بأعلی النّدآء ناطق خود را منع منما توکّل علی
اللّه و فوّض امورک الیه ثمّ احضر فى السّجن الأعظم لتسمع ما لا سمعت الآذان
شبهه و تری ما لا رأت العیون و الأبصار آیا بعد از این بیان از برای احدی
حجّتی باقیست لا ونفس اللّه القائمة علی الأمر براستی میگویم امروز کلمهٔ
مبارکهٔ ولکنّه رسول اللّه و خاتم النّبیّین به یوم یقوم النّاس لربّ العالمین
منتهی شد اشکر اللّه بهذا الفضل العظیم

\end{otherlanguage}

Yā Shaykh dar ānchih dhikr shudih tafakkur farmā shāyad bi-quvvat-i
ism-i Qayyūm az raḥīq-i makhtūm biyāshāmī va biyābī ānchih rā kih kull
az idrākash ʿājizand. Kamar-i himmat rā muḥkam namā va qaṣd-i malakūt-i
aʿlā kun shāyad dar ḥīn-i tanzīl nafaḥāt-i vaḥy va ilhām rā biyābī va
bi-ān fāʾiz shavī. Bi-rāstī mīgūyam az barā-yi Amr-i Ilāhī shibh va
mithlī nabūdih va nīst. Ḥujubāt-i awhām rā kharq namā. Innahu yumidduka
wa-yuʾayyiduka faḍlan min ʿindihi wa-huwa al-Qawiyyu al-Ghālibu
al-Qadīr. Tā vaqt bāqī va Sidrih-yi Mubārakih mābayn-i bariyyih bi-aʿlā
al-nidāʾ nāṭiq khud rā manʿ manamā. Tawakkal ʿalā Allāhi wa-fawwiḍ
umūraka ilayhi thumma iḥḍur fī al-Sijni al-Aʿẓami li-tasmaʿa mā lā
samiʿati al-ādhānu shibhahu wa-tarā mā lā raʾati al-ʿuyūnu wa-al-abṣār.
Āyā baʿd az īn bayān az barā-yi aḥadī ḥujjatī bāqīst? Lā wa-nafsi Allāhi
al-qāʾimati ʿalā al-Amr. Bi-rāstī mīgūyam imrūz kalimih-yi mubārakih-yi
wa-lākinnahu rasūlu Allāhi wa-khātamu al-nabiyyīn bi-yawma yaqūmu
al-nāsu li-rabbi al-ʿālamīn muntahī shud. Ushkur Allāha bi-hādhā
al-faḍli al-ʿaẓīm.

O Shaykh! Reflect on what has been mentioned, so that perhaps, through
the power of the name of the Self-Subsisting, you may drink of the
sealed wine and discover what all are powerless to comprehend. Gird
yourself with determination and set your course toward the highest
Kingdom, so that perhaps, at the moment of revelation, you may perceive
the breezes of revelation and inspiration and attain them. Truly I say,
the Cause of God has never had, nor does it have, any peer or likeness.
Tear through the veils of vain imaginings. He will indeed assist and
strengthen you as a favor from Him, and He is the Strong, the
All-Subduing, the All-Powerful. While time remains and the Blessed Lote
Tree speaks among the people with the loudest call, do not deprive
yourself. Put your trust in God and entrust your affairs to Him; then
come into the Most Great Prison, that you may hear what no ears have
ever heard the like of, and see what no eyes or vision have seen. After
this exposition, does anyone have any excuse left? No, by the very Self
of God, Who presides over the Cause! Truly I say, today the blessed
words, ``But he is the Messenger of God and the Seal of the Prophets,''
have culminated in ``the day when humanity will stand before the Lord of
the worlds.'' Thank God for this immense bounty.

\subsection{170}\label{section-1456}

\begin{otherlanguage}{persian}

یا شیخ نفحات وحی بدونش مشتبه نشده و نمیشود حال سدرهٔ منتهی با اثمار
لاتحصی امام وجهت حاضر خود را باوهام مثل حزب قبل میالای از نفس بیان فطرة
اللّه ظاهر و مشهود اوست گواه کل در اثبات ظهورش باحدی محتاج نبوده و نیست
حال قریب صد جلد آیات باهرات و کلمات محکمات از سماء مشیّت منزل آیات نازل و
حاضر لک ان تقصد المقصد الأقصی و الغایة القصوی و الذّروة العلیا لتسمع و
تری ما ظهر من لدی اللّه ربّ العالمین

\end{otherlanguage}

Yā Shaykh nafaḥāt-i vaḥy bi-dūnash mushtabih nashudih va namīshavad. Ḥāl
Sidrih-yi Muntahā bā athmār-i lā tuḥṣā imām-i vajhat ḥāḍir. Khud rā
bi-awhām mithl-i ḥizb-i qabl mayālāy. Az nafs-i bayān fiṭrat-i Allāh
ẓāhir va mashhūd. Ūst guvāh-i kull dar ithbāt-i ẓuhūrash bi-aḥadī muḥtāj
nabūdih va nīst. Ḥāl qarīb-i ṣad jild āyāt-i bāhirāt va kalimāt-i
muḥkamāt az samā-yi mashiyyat-i Munzil-i Āyāt nāzil va ḥāḍir. Laka an
taqṣida al-maqṣada al-aqṣā wa-al-ghāyata al-quṣwā wa-al-dhirwata
al-ʿulyā li-tasmaʿa wa-tarā mā ẓahara min ladā Allāhi rabbi al-ʿālamīn.

O Shaykh! The breezes of revelation have never been, nor will they ever
be, confused with anything else. Now the Lote Tree of the Utmost
Boundary, with its innumerable fruits, stands before your face. Do not,
like the people of the past, defile yourself with vain imaginings. The
innate nature bestowed by God is evident and manifest in the utterance
itself. He is the witness to all; to establish His Revelation, He has
never needed, nor does He need, anyone. Now nearly a hundred volumes of
resplendent verses and firmly established words have descended from the
heaven of the will of the Revealer of verses and are at hand. It is for
you to set your course toward the furthest destination, the ultimate
goal, and the loftiest summit, that you may hear and see what has
appeared from God, the Lord of the worlds.

\subsection{171}\label{section-1457}

\begin{otherlanguage}{persian}

قدری در آیات لقا که از مالک ملکوت اسماء در فرقان نازل شده تفکّر نما شاید
راه مستقیم را بیابی و سبب و علّت هدایت خلق شوی مثل شما امروز باید بر خدمت
امر قیام نماید ذلّت این مظلوم و عزّت شما هر دو بفنا راجع جهد کن شاید فائز
شوی بعملی که عرفش از عالم قطع نشود در ذکر لقا نازل شده آنچه که از برای
منکرین مجال رد و عدم قبول نمانده و نیست قوله تبارک و تعالی اللّه الّذی رفع
السّموات بغیر عمد ترونها ثمّ استوی علی العرش و سخّر الشّمس و القمر کلّ یجری
لأجل مسمّی یدبّر الأمر یفصّل الآیات لعلّکم بلقآء ربّکم توقنون و میفرماید من
کان یرجو لقآء اللّه فانّ اجل اللّه لآت و هو السّمیع العلیم الی قوله تعالی و
الّذین کفروا بآیات اللّه و لقائه اولئک یئسوا من رحمتی و اولئک لهم عذاب
الیم و همچنین میفرماید قالوا ائذا ضللنا فی الأرض ائنّا لفی خلق جدید بل هم
بلقآء ربّهم کافرون و همچنین میفرماید الا انّهم فی مریة من لقآء ربّهم الا
انّه بکلّ شیء محیط و همچنین میفرماید انّ الّذین لا یرجون لقائنا و رضوا
بالحیوة الدّنیا و اطمأنّوا بها و الّذینهم عن آیاتنا غافلون اولئک مأواهم
النّار بما کانوا یکسبون و هم‌چنین میفرماید و اذا تتلی علیهم آیاتنا بیّنات
قال الّذین لا یرجون لقائنا ائت بقرآن غیر هذا او بدّله قل ما یکون لی ان
ابدّله من تلقآء نفسی ان اتّبع الّا ما یوحی الیّ انّی اخاف ان عصیت ربّی عذاب
یوم عظیم و همچنین میفرماید ثمّ آتینا موسی الکتاب تماماً علی الّذی احسن و
تفصیلاً لکلّ شیء و هدی و رحمةً لعلّهم بلقآء ربّهم یؤمنون و همچنین میفرماید
اولئک الّذین کفروا بآیات ربّهم و لقائه فحبطت اعمالهم فلا نقیم لهم یوم
القیامة وزناً ذلک جزائهم جهنّم بما کفروا و اتّخذوا آیاتی و رسلی هزؤاً و
هم‌چنین میفرماید هل اتاک حدیث موسی اذ رأی ناراً فقال لأهله امکثوا انّی آنست
ناراً لعلّی آتیکم منها بقبس او اجد على النّار هدی فلمّا اتاها نودی یا موسی
انّی انا ربّک فاخلع نعلیک انّک بالواد المقدّس طوی و انا اخترتک فاستمع لما
یوحی انّنی انا اللّه لا اله الّا انا فاعبدنی و هم‌چنین میفرماید أ و لم
یتفکّروا فی انفسهم ما خلق اللّه السّموات و الأرض و ما بینهما الّا بالحقّ و
اجل مسمّی و انّ کثیراً من النّاس بلقآء ربّهم لکافرون و همچنین میفرماید الا
یظنّ اولئک انّهم مبعوثون لیوم عظیم یوم یقوم النّاس لربّ العالمین و همچنین
میفرماید و لقد آتینا موسی الکتاب فلا تکن فی مریة من لقائه و میفرماید کلّا
اذا دکّت الأرض دکّاً دکّاً و جآء ربّک و الملک صفّاً صفّاً و همچنین میفرماید یریدون
لیطفئوا نور اللّه بأفواههم و اللّه متمّ نوره ولو کره الکافرون و همچنین
میفرماید فلمّا قضی موسی الأجل و سار بأهله آنس من جانب الطّور ناراً قال
لأهله امکثوا انّی آنست ناراً لعلّى آتیکم منها بخبر او جذوة من النّار لعلّکم
تصطلون فلمّا اتاها نودی من شاطئ الوادی الأیمن فی البقعة المبارکة من
الشّجرة ان یا موسی انّی انا اللّه ربّ العالمین انتهی

\end{otherlanguage}

Qadrī dar āyāt-i liqā kih az Mālik-i Malakūt-i Asmā dar Furqān nāzil
shudih tafakkur namā shāyad rāh-i mustaqīm rā biyābī va sabab va
ʿillat-i hidāyat-i khalq shavī. Mithl-i shumā imrūz bāyad bar khidmat-i
Amr qiyām namāyad. Dhillat-i īn Maẓlūm va ʿizzat-i shumā har du bi-fanā
rājiʿ. Jahd kun shāyad fāʾiz shavī bi-ʿamalī kih ʿarfash az ʿālam qaṭʿ
nashavad. Dar dhikr-i liqā nāzil shudih ānchih kih az barā-yi munkirīn
majāl-i radd va ʿadam-i qabūl namāndih va nīst. Qawluhu tabāraka
wa-taʿālā: Allāhu alladhī rafaʿa al-samāwāti bi-ghayri ʿamadin tarawnahā
thumma istawā ʿalā al-ʿarshi wa-sakhkhara al-shamsa wa-al-qamara kullun
yajrī li-ajalin musamman yudabbiru al-amra yufaṣṣilu al-āyāti laʿallakum
bi-liqāʾi rabbikum tūqinūn. Va mīfarmāyad: man kāna yarjū liqāʾa Allāhi
fa-inna ajala Allāhi la-ātin wa-huwa al-Samīʿu al-ʿAlīm. Ilā qawlihi
taʿālā: wa-alladhīna kafarū bi-āyāti Allāhi wa-liqāʾihi ulāʾika yaʾisū
min raḥmatī wa-ulāʾika lahum ʿadhābun alīm. Va hamchunīn mīfarmāyad:
qālū a-idhā ḍalalnā fī al-arḍi a-innā la-fī khalqin jadīdin bal hum
bi-liqāʾi rabbihim kāfirūn. Va hamchunīn mīfarmāyad: alā innahum fī
miryatin min liqāʾi rabbihim alā innahu bi-kulli shayʾin muḥīṭ. Va
hamchunīn mīfarmāyad: inna alladhīna lā yarjūna liqāʾanā wa-raḍū
bi-al-ḥayāti al-dunyā wa-iṭmaʾannū bihā wa-alladhīna hum ʿan āyātinā
ghāfilūn ulāʾika maʾwāhumu al-nāru bi-mā kānū yaksibūn. Va hamchunīn
mīfarmāyad: wa-idhā tutlā ʿalayhim āyātunā bayyinātin qāla alladhīna lā
yarjūna liqāʾanā iʾti bi-Qurʾānin ghayri hādhā aw baddilhu qul mā yakūnu
lī an ubaddilahu min tilqāʾi nafsī in attabiʿu illā mā yūḥā ilayya innī
akhāfu in ʿaṣaytu rabbī ʿadhāba yawmin ʿaẓīm. Va hamchunīn mīfarmāyad:
thumma ātaynā Mūsā al-kitāba tamāman ʿalā alladhī aḥsana wa-tafṣīlan
li-kulli shayʾin wa-hudan wa-raḥmatan laʿallahum bi-liqāʾi rabbihim
yuʾminūn. Va hamchunīn mīfarmāyad: ulāʾika alladhīna kafarū bi-āyāti
rabbihim wa-liqāʾihi fa-ḥabiṭat aʿmāluhum fa-lā nuqīmu lahum yawma
al-qiyāmati waznan dhālika jazāʾuhum jahannamu bi-mā kafarū
wa-ittakhadhū āyātī wa-rusulī huzūʾan. Va hamchunīn mīfarmāyad: hal
atāka ḥadīthu Mūsā idh raʾā nāran fa-qāla li-ahlihi umkuthū innī ānastu
nāran laʿallī ātīkum minhā bi-qabasin aw ajidu ʿalā al-nāri hudan
fa-lammā atāhā nūdiya yā Mūsā innī anā rabbuka fa-ikhlaʿ naʿlayka innaka
bi-al-wādi al-muqaddasi Ṭuwan wa-anā ikhtartuka fa-istamiʿ li-mā yūḥā
innanī anā Allāhu lā ilāha illā anā fa-uʿbudnī. Va hamchunīn mīfarmāyad:
a-wa-lam yatafakkarū fī anfusihim mā khalaqa Allāhu al-samāwāti
wa-al-arḍa wa-mā baynahumā illā bi-al-ḥaqqi wa-ajalin musamman wa-inna
kathīran mina al-nāsi bi-liqāʾi rabbihim la-kāfirūn. Va hamchunīn
mīfarmāyad: a-lā yaẓunnu ulāʾika annahum mabʿūthūna li-yawmin ʿaẓīmin
yawma yaqūmu al-nāsu li-rabbi al-ʿālamīn. Va hamchunīn mīfarmāyad:
wa-laqad ātaynā Mūsā al-kitāba fa-lā takun fī miryatin min liqāʾihi. Va
mīfarmāyad: kallā idhā dukkati al-arḍu dakkan dakkan wa-jāʾa rabbuka
wa-al-malaku ṣaffan ṣaffan. Va hamchunīn mīfarmāyad: yurīdūna li-yuṭfiʾū
nūra Allāhi bi-afwāhihim wa-Allāhu mutimmu nūrihi wa-law kariha
al-kāfirūn. Va hamchunīn mīfarmāyad: fa-lammā qaḍā Mūsā al-ajala wa-sāra
bi-ahlihi ānasa min jānibi al-Ṭūri nāran qāla li-ahlihi umkuthū innī
ānastu nāran laʿallī ātīkum minhā bi-khabarin aw jadhwatin mina al-nāri
laʿallakum taṣṭalūn fa-lammā atāhā nūdiya min shāṭiʾi al-wādi al-aymani
fī al-buqʿati al-mubārakati mina al-shajarati an yā Mūsā innī anā Allāhu
rabbu al-ʿālamīn. Intahā.

Reflect a little on the verses concerning the meeting with God that have
been revealed in the Qurʾān by the Sovereign of the Kingdom of Names, so
that perhaps you may find the straight path and become the means and
cause of guiding the people. Someone such as you should arise today to
serve the Cause. The abasement of this Wronged One and your honor will
both pass away. Strive, so that perhaps you may accomplish a deed whose
fragrance will never cease in the world. Concerning the meeting with
God, what has been revealed has left, and leaves, those who deny it no
room for rejection or refusal. His words, blessed and exalted is He:
``God is the One who raised the heavens without pillars that you can
see, then established Himself upon the throne and subjected the sun and
the moon, each running for an appointed term. He directs all affairs and
sets forth the signs in detail, so that you may be certain of meeting
your Lord.'' And He says: ``Whoever hopes to meet God---the term
appointed by God will surely come, and He is the All-Hearing, the
All-Knowing.'' Further on, He, exalted is He, says: ``Those who
disbelieve in the signs of God and in meeting Him have despaired of My
mercy, and for them there is a painful punishment.'' He also says:
``They say: `When we have been lost in the earth, will we really be
created anew?' Rather, they deny the meeting with their Lord.'' He also
says: ``Indeed, they are in doubt about meeting their Lord. Indeed, He
encompasses all things.'' He also says: ``Those who do not hope to meet
Us, and are content with the life of this world and feel secure in it,
and those who are heedless of Our signs---their refuge is the Fire
because of what they have earned.'' He also says: ``When Our clear
verses are recited to them, those who do not hope to meet Us say: `Bring
a Qurʾān other than this, or change it.' Say: `It is not for me to
change it of my own accord. I follow only what is revealed to me. I
fear, if I disobey my Lord, the punishment of a mighty day.'\,'' He also
says: ``Then We gave Moses the Book, completing Our favor upon the one
who did good, explaining all things in detail, and providing guidance
and mercy, so that they might believe in meeting their Lord.'' He also
says: ``These are the ones who disbelieve in the signs of their Lord and
in meeting Him, so their deeds have come to nothing, and on the Day of
Resurrection We will give them no weight. That is their
recompense---Hell---because they disbelieved and made a mockery of My
signs and My messengers.'' He also says: ``Has the story of Moses
reached you? When he saw a fire, he said to his family: `Stay here. I
have perceived a fire. Perhaps I can bring you a burning brand from it,
or find guidance at the fire.' When he came to it, he was called: `O
Moses! I am indeed your Lord. Take off your sandals: you are in the
sacred valley of Ṭuwā. I have chosen you, so listen to what is revealed.
I am indeed God; there is no god but Me, so worship Me.'\,'' He also
says: ``Have they not reflected within themselves? God did not create
the heavens and the earth and what is between them except in truth and
for an appointed term. Yet many people indeed deny the meeting with
their Lord.'' He also says: ``Do they not think that they will be raised
up for a mighty day, the day when humanity will stand before the Lord of
the worlds?'' He also says: ``We indeed gave Moses the Book, so do not
be in doubt about meeting Him.'' And He says: ``No indeed! When the
earth is pounded to dust, again and again, and your Lord comes, and the
angels, rank upon rank.'' He also says: ``They wish to extinguish the
light of God with their mouths, but God will perfect His light, however
much the unbelievers may detest it.'' He also says: ``When Moses had
fulfilled the term and was traveling with his family, he perceived a
fire on the side of the Mount. He said to his family: `Stay here. I have
perceived a fire. Perhaps I can bring you news from it, or a burning
ember from the fire, so that you may warm yourselves.' When he came to
it, he was called from the right bank of the valley, in the blessed
spot, from the tree: `O Moses! I am indeed God, the Lord of the
worlds.'\,'' End of quotation.

\subsection{172}\label{section-1458}

\begin{otherlanguage}{persian}

در جمیع کتب الهی وعد‌‌هٔ لقا صریح بوده و هست و مقصود از این لقا لقاء مشرق
آیات و مطلع بیّنات و مظهر اسماء حسنی و مصدر صفات علیای حقّ جلّ جلاله است حقّ
بذاته و بنفسه غیب منیع لایدرک بوده پس مقصود از لقا لقاء نفسی است که
قائم‌مقام اوست مابین عباد و از برای او هم شبه و مثلی نبوده و نیست چه اگر
از برای او شبه و مثلی مشاهده شود کیف یثبت تقدیس ذاته و تنزیه کینونته عن
الأشباه و الأمثال باری در مقامات لقا و تجلّی در کتاب ایقان نازل شده آنچه
که منصفین را کفایت نماید نسأله تعالی ان یؤیّد الکلّ علی الصّدق الخالص و
یقرّبهم الیه انّه هو المقتدر القدیر لا اله الّا هو السّامع النّاطق العزیز
الحمید

\end{otherlanguage}

Dar jamīʿ-i kutub-i Ilāhī vaʿdih-yi liqā ṣarīḥ būdih va hast va maqṣūd
az īn liqā liqāʾ-i Mashriq-i Āyāt va Maṭlaʿ-i Bayyināt va Maẓhar-i
Asmā-yi Ḥusnā va Maṣdar-i Ṣifāt-i ʿUlyā-yi Ḥaqq jalla jalāluhu ast. Ḥaqq
bi-dhātih va bi-nafsih ghayb-i manīʿ-i lā yudrak būdih pas maqṣūd az
liqā liqāʾ-i nafsī ast kih qāʾim-maqām-i Ūst mābayn-i ʿibād va az
barā-yi Ū ham shibh va mithlī nabūdih va nīst chih agar az barā-yi Ū
shibh va mithlī mushāhidih shavad kayfa yathbutu taqdīsu dhātihi
wa-tanzīhu kaynūnatihi ʿani al-ashbāhi wa-al-amthāl. Bārī dar maqāmāt-i
liqā va tajallī dar Kitāb-i Īqān nāzil shudih ānchih kih munṣifīn rā
kifāyat namāyad. Nasʾaluhu taʿālā an yuʾayyida al-kulla ʿalā al-ṣidqi
al-khāliṣi wa-yuqarribahum ilayhi innahu huwa al-Muqtadiru al-Qadīru lā
ilāha illā huwa al-Sāmiʿu al-Nāṭiqu al-ʿAzīzu al-Ḥamīd.

In all the divine Books, the promise of meeting God has been, and
remains, explicit. What is meant by this meeting is meeting the Dawning
Place of the verses, the Dayspring of the clear proofs, the
Manifestation of the most beautiful names, and the Source of the most
exalted attributes of God, glorified be His glory. God in His Essence
and His own Self has ever been an inaccessible, unknowable mystery.
Thus, what is meant by this meeting is meeting the One who stands in His
place among His servants. He, too, has never had, nor does He have, any
peer or likeness, for if any peer or likeness to Him were found, how
could the sanctity of His Essence and the transcendence of His Being
above all peers and likenesses be established? In short, concerning the
stations of meeting God and His self-revelation, enough has been
revealed in the Book of Certitude to satisfy the fair-minded. We ask
Him, exalted is He, to aid everyone to attain pure truthfulness and to
draw them near to Him. He is indeed the All-Powerful, the Almighty.
There is no god but Him, the Hearing, the Speaking, the Invincible, the
Praiseworthy.

\subsection{173}\label{section-1459}

\begin{otherlanguage}{persian}

یا ایّها المعروف بالعلم مُر العباد بالمعروف و لا تکن من المتوقّفین ببصر
حدید نظر نما آفتاب حقیقت بامر مالک ملکوت بیان و سلطان جبروت عرفان از افق
سماء سجن عکّا مشرق و لائح اعراض او را محجوب ننمود و صفوف و الوف منعش نکرد
از برای آن جناب هیچ عذری باقی نمانده یا باید اقرار نمائی و یا نعوذ باللّه
بر انکار کل قیام کنی

\end{otherlanguage}

Yā ayyuhā al-maʿrūfu bi-al-ʿilmi mur al-ʿibāda bi-al-maʿrūfi wa-lā takun
mina al-mutawaqqifīn. Bi-baṣar-i ḥadīd naẓar namā. Āftāb-i Ḥaqīqat
bi-amr-i Mālik-i Malakūt-i Bayān va Sulṭān-i Jabarūt-i ʿIrfān az ufuq-i
samā-yi Sijn-i ʿAkkā mushriq va lāʾiḥ. Iʿrāḍ Ū rā maḥjūb nanamūd va
ṣufūf va ulūf manʿash nakard. Az barā-yi ān janāb hīch ʿudhrī bāqī
namāndih yā bāyad iqrār namāʾī va yā naʿūdhu bi-Allāh bar inkār-i kull
qiyām kunī.

O you who are known for learning! Enjoin what is right upon the servants
and do not be among those who hesitate. Look with keen vision. By the
command of the Sovereign of the Kingdom of Utterance and the Ruler of
the Dominion of Recognition, the Sun of Truth has risen and shines from
the horizon of the heaven of the Prison of ʿAkkā. Rejection has not
veiled Him, nor have ranks and thousands held Him back. No excuse
remains for you: either you must acknowledge Him, or---God forbid---you
must arise to deny them all.

\subsection{174}\label{section-1460}

\begin{otherlanguage}{persian}

یا شیخ در حزب شیعه تفکّر نما چه مقدار عمارتها که بایادی ظنون و اوهام
تعمیر نمودند و چه شهرها بنا نهادند بالأخره آن اوهام برصاص تبدیل شد و بر
سیّد عالم وارد و یک نفس از رؤسای آن حزب در یوم ظهور اقبال ننمود نزد ذکر
اسم مبارک کل به عجّل اللّه فرجه ناطق ولکن در یوم ظهور آن شمس حقیقت کل به
عجّل اللّه فی نقمته متکلّم و ناطق چنانچه دیده شد ساذج وجود و مالک غیب و
شهود را آویختند و عمل نمودند آنچه را که لوح گریست و قلم نوحه نمود و
زفرات مخلصین مرتفع و عبرات مقرّبین نازل

\end{otherlanguage}

Yā Shaykh, dar ḥizb-e Shīʿa tafakkor namā, che meqdār ʿemārathā ke
be-ayādī-ye ẓonūn va awhām taʿmīr namūdand va che shahrhā banā nahādand.
Be-l-ākhera ān awhām be-raṣāṣ tabdīl shod va bar Sayyed-e ʿālam vāred,
va yek nafs az roʾasā-ye ān ḥizb dar yawm-e ẓohūr eqbāl nanamūd. Nazd-e
dhekr-e esm-e mobārak koll be-ʿajjala Allāhu farajahu nāṭeq, va-lāken
dar yawm-e ẓohūr-e ān Shams-e Ḥaqīqat koll be-ʿajjala Allāhu fī
niqmatihi motakallem va nāṭeq. Chenānche dīde shod, Sādhej-e Vojūd va
Mālek-e ghayb va shohūd rā āvīkhtand va ʿamal namūdand ānche rā ke Lawḥ
gerīst va Qalam nawḥe namūd va zafarāt-e mokhleṣīn mortafeʿ va ʿabarāt-e
moqarrabīn nāzel.

O Shaykh! Reflect on the Shiʿi community: how many edifices they raised
with the hands of conjecture and delusion, and how many cities they
built. In the end, those delusions turned into bullets and struck the
Lord of the world, and not one of that community's leaders turned toward
Him on the day of His appearance. At the mention of His blessed name,
all would say, ``May God hasten His advent,'' but on the day that Sun of
Truth appeared, all spoke and said, ``May God hasten His vengeance
against Him.'' As was witnessed, they suspended the Essence of Being,
the Sovereign of the unseen and the seen, and committed deeds that made
the Tablet weep and the Pen lament, raised the sighs of the sincere, and
caused the tears of those brought near to flow.

\subsection{175}\label{section-1461}

\begin{otherlanguage}{persian}

یا شیخ فکر نما و بانصاف تکلّم کن حزب شیخ احسائی باعانت الهی عارف شدند
بآنچه که دون آن حزب از آن محروم و محجوب مشاهده گشتند باری در هر عصری و
هر قرنی ایّام ظهور مشارق وحی و مطالع الهام و مهابط علم الهی اختلاف ظاهر و
سبب و علّت آن نفوس کاذبهٔ ملحده بوده‌اند شرح این مقام جائز نه آن جناب خود
اعرف و اعلمند باوهام متوهّمین و ظنون مریبین

\end{otherlanguage}

Yā Shaykh, fekr namā va be-enṣāf takallom kon. Ḥizb-e Shaykh-e Aḥsāʾī
be-eʿānat-e elāhī ʿāref shodand be-ānche ke dūn-e ān ḥizb az ān maḥrūm
va maḥjūb moshāhede gashtand. Bārī, dar har ʿaṣrī va har qarnī ayyām-e
ẓohūr-e mashāreq-e vaḥy va maṭāleʿ-e elhām va mahābeṭ-e ʿelm-e elāhī
ekhtelāf ẓāher va sabab va ʿellat-e ān nofūs-e kādhebe-ye molḥede
būde-and. Sharḥ-e īn maqām jāʾez na. Ān janāb khod aʿraf va aʿlamand
be-awhām-e motavahhemīn va ẓonūn-e morībīn.

O Shaykh! Reflect and speak fairly. Through divine assistance, the
followers of Shaykh Aḥsāʾī came to recognize that of which those outside
that community were seen to be deprived and from which they were veiled.
In every age and every century, when the Dawning Places of revelation,
the Daysprings of inspiration, and the Repositories of divine knowledge
appeared, dissension arose, and lying, faithless souls were its cause.
It is not appropriate to elaborate on this matter. You yourself are more
fully aware and knowledgeable concerning the delusions of the deluded
and the conjectures of those who sow doubt.

\subsection{176}\label{section-1462}

\begin{otherlanguage}{arabic}

الیوم این مظلوم از آن جناب و سائر علما که از کأس علم الهی نوشیده‌اند و از
کلمات انوار نیّر عدل منوّرند مسئلت مینماید نفسی را معیّن فرمایند من‌دون
اطّلاع احدی و او را باین جهات بفرستند و چندی در جزیرهٔ قبرس توقّف نماید و
با میرزا یحیی معاشر شود شاید بر اصل امر و مصدر اوامر و احکام الهی آگاه
گردد

\end{otherlanguage}

Al-yawm īn Maẓlūm az ān janāb va sāʾer-e ʿolamā ke az kaʾs-e ʿelm-e
elāhī nūshīde-and va az kalemāt-e anvār-e Nayy er-e ʿAdl monavvarand
masʾalat mīnamāyad nafsī rā moʿayyan farmāyand min dūn-e eṭṭelāʿ-e aḥadī
va ū rā be-īn jehāt beferestand va chandī dar jazīre-ye Qebres tavaqqof
namāyad va bā Mīrzā Yaḥyā moʿāsher shavad, shāyad bar aṣl-e amr va
maṣdar-e avāmer va aḥkām-e elāhī āgāh gardad.

Today this Wronged One asks you and the other learned men who have drunk
from the cup of divine knowledge and are illumined by the words of the
Luminary of Justice to appoint someone, without anyone's knowledge, and
send him to these regions. Let him stay for a time on the island of
Cyprus and associate with Mīrzā Yaḥyā, so that he may perhaps become
acquainted with the true nature of the Cause and the source of the
divine commands and laws.

\subsection{177}\label{section-1463}

\begin{otherlanguage}{persian}

اگر قدری تفکّر نمائی شهادت میدهی بر حکمت و قدرت و سلطنت حقّ جلّ جلاله
معدودی که از امر آگاه نه و با ما نبوده‌اند گفته‌اند آنچه را که اشیاء کلّها
و نفوس مطمئنّهٔ راضیهٔ مرضیّه گواهی داده‌اند بر کذب آن نفوس غافله حال اگر آن
جناب همّت فرمایند حقیقت امر بر عالمیان کشف شود و ناس را از این ظلمات مظلم
صیلم نجات بخشد لو لا البهآء من یقدر ان ینطق امام وجوه الأنام و لولاه من
یقدر ان یتکلّم بما امر به من لدی اللّه ربّ الجنود

\end{otherlanguage}

Agar qadrī tafakkor namāʾī shahādat mīdahī bar ḥekmat va qodrat va
salṭanat-e Ḥaqq, jalla jalāluhu. Maʿdūdī ke az amr āgāh na va bā mā
nabūde-and gofte-and ānche rā ke ashyāʾ kolluhā va nofūs-e moṭmaʾenne-ye
rāḍīye-ye marḍīye gavāhī dāde-and bar kedhb-e ān nofūs-e ghāfele. Ḥāl
agar ān janāb hemmat farmāyand ḥaqīqat-e amr bar ʿālamiyān kashf shavad
va nās rā az īn ẓolomāt-e moẓlem-e ṣaylam nejāt bakhshad. Law lā
al-Bahāʾ man yaqdiru an yanṭiqa amāma wujūhi al-anām, wa-lawlāhu man
yaqdiru an yatakallama bi-mā umira bihi min ladā Allāhi Rabbi al-junūd.

If you reflect a little, you will testify to the wisdom, power, and
sovereignty of God, glorified be His glory. A few who know nothing of
the Cause and have not been with Us have made claims whose falsehood all
things, and souls that are assured, content, and pleasing to God, have
attested. Now, if you exert yourself, the truth of the matter will be
disclosed to the world and will deliver people from this dark,
devastating gloom. Were it not for Bahāʾ, who could speak openly before
the faces of humanity? Were it not for Him, who could utter what he has
been commanded by God, the Lord of Hosts?

\subsection{178}\label{section-1464}

\begin{otherlanguage}{persian}

حال نفس غافل بحبل روضه‌خوانی تمسّک جسته لعمر اللّه انّه فی کذب مبین چه که
این حزب را اعتقاد آنکه در ظهور حضرت قائم ائمّه سلام اللّه علیهم از قبور
برخاسته‌اند هذا حقّ لا ریب فیه از حقّ میطلبیم متوهّمین را از کوثر ایقان که
از معین قلم اعلی جاریست قسمت عطا فرماید تا کل فائز شوند بآنچه که سزاوار
ایّام اوست

\end{otherlanguage}

Ḥāl nafs-e ghāfel be-ḥabl-e rawḍe-khvānī tamassok joste. La-ʿamru Allāhi
innahu la-fī kid hbin mubīn. Che ke īn ḥizb rā eʿteqād ānke dar ẓohūr-e
Ḥaḍrat-e Qāʾem aʾemme, salāmu Allāhi ʿalayhim, az qobūr barkhāste-and.
Hādhā ḥaqqun lā rayba fīhi. Az Ḥaqq mīṭalabīm motavahhemīn rā az
kawthar-e īqān ke az maʿīn-e Qalam-e Aʿlā jārīst qesmat ʿaṭā farmāyad tā
koll fāʾez shavand be-ānche ke sazāvār-e ayyām-e Ūst.

Now that heedless soul has clung to the cord of reciting lamentations.
By the life of God, he is plainly lying. For this community believes
that, at the appearance of the Qāʾim, the Imams, God's peace be upon
them, rose from their graves. This is a truth beyond doubt. We ask God
to grant the deluded a share of the Kawthar of certitude flowing from
the wellspring of the Most Exalted Pen, so that all may attain what
befits His days.

\subsection{179}\label{section-1465}

\begin{otherlanguage}{arabic}

یا شیخ در بحبوحهٔ بلایا این مظلوم بتحریر این کلمات مشغول از جمیع جهات نار
ظلم و اعتساف مشهود از یک جهت خبر رسیده اولیا را در ارض طاء اخذ نموده‌اند
مع آنکه آفتاب و ماه و بر و بحر گواه که این حزب بطراز وفا مزیّنند و جز
بارتفاع دولت و نظم مملکت و راحت ملّت بامری تمسّک نجسته و نخواهند جست

\end{otherlanguage}

Yā Shaykh, dar boḥbūḥe-ye balāyā īn Maẓlūm be-taḥrīr-e īn kalemāt
mashghūl. Az jamīʿ-e jehāt nār-e ẓolm va eʿtesāf mashhūd. Az yek jahat
khabar resīde awliyā rā dar arḍ-e Ṭāʾ akhdh namūde-and, maʿ-e ānke āftāb
va māh va barr va baḥr gavāh ke īn ḥizb be-ṭerāz-e vafā mozayyanand va
joz be-ertefāʿ-e dawlat va naẓm-e mamlakat va rāḥat-e mellat be-amrī
tamassok najoste va nakhvāhand jost.

O Shaykh! Amid tribulations, this Wronged One is occupied with writing
these words. On every side the fire of oppression and tyranny is
evident. From one direction, news has come that the loved ones have been
arrested in the land of Ṭāʾ, although the sun and moon, land and sea
bear witness that this community is adorned with loyalty and has sought,
and will seek, nothing but the advancement of the government, the order
of the country, and the well-being of its people.

\subsection{180}\label{section-1466}

\begin{otherlanguage}{persian}

یا شیخ مکرّر گفتیم حضرت پادشاه را در سنین معدودات نصرت نموده‌ایم سالهاست
در ایران امر مغایری ظاهر نشده زمام مفسدین احزاب در قبضهٔ اقتدار مقبوض
احدی از حدّ تجاوز ننمود لعمر اللّه این حزب اهل فساد نبوده و نیستند قلوبشان
بنور تقوی منوّر و بطراز محبّة اللّه مزیّن همّشان اصلاح عالم بوده و هست و
اراده آنکه اختلاف از میان برخیزد و نار ضغینه و بغضا خاموشی پذیرد تا جمیع
ارض قطعهٔ واحده مشاهده گردد

\end{otherlanguage}

Yā Shaykh, mokarrar goftīm Ḥaḍrat-e Pādeshāh rā dar senīn-e maʿdūdāt
noṣrat namūde-īm. Sālhāst dar Īrān amr-e moghāyerī ẓāher nashode.
Zemām-e mofsedīn-e aḥzāb dar qabḍe-ye eqtedār maqbūḍ. Aḥadī az ḥadd
tajāvoz nanamūd. La-ʿamru Allāh, īn ḥizb ahl-e fasād nabūde va nīstand.
Qolūbeshān be-nūr-e taqvā monavvar va be-ṭerāz-e maḥabbat Allāh
mozayyan. Hammes hān eṣlāḥ-e ʿālam būde va hast va erāde ānke ekhtelāf
az miyān barkhīzad va nār-e ḍaghīne va baghḍā khāmūshī padhīrad tā
jamīʿ-e arḍ qeṭʿe-ye vāḥede moshāhede gardad.

O Shaykh! We have repeatedly said that We have aided His Majesty the
Shah over a number of years. For years, nothing contrary to this has
appeared in Iran. The reins of the troublemakers among the factions are
held firmly in the grasp of power. No one has overstepped the bounds. By
the life of God, this community has not been, and is not, a people of
sedition. Their hearts are illumined by the light of piety and adorned
with the love of God. Their concern has been, and remains, the
betterment of the world, and their desire is that dissension may cease
and the fire of malice and hatred be extinguished, so that the whole
earth may be seen as one land.

\subsection{181}\label{section-1467}

\begin{otherlanguage}{persian}

و از جهت دیگر دائرهٔ سفارت ایران در مدینهٔ کبیره بتمام قدرت و قوّت بتضییع
این مظلومان مشغول انّهم ارادوا امراً و اللّه اراد امراً آخر حال تفکّر فرمائید
در آنچه بر امنای حقّ در هر دیار وارد شده گاهی نسبت سرقت و دزدی داده‌اند و
هنگامى بمفتریاتی تکلّم نموده‌اند که در عالم شبه و مثل نداشته حال آن جناب
بانصاف تکلّم فرمایند آیا نسبت سرقتی که از جانب سفارت کبری برعیّت خود داده
ثمر و اثرش در ممالک خارجه چیست از این فقره این مظلوم خجل شد نه از جهت
آنکه سبب و علّت تضییع این عبد بوده بلکه علّت خجلت اطّلاع سفرای اجنبیّه بر
مراتب تدبیر و ادراک جمعی از معتبرین ایران در سفارت کبری بوده متّهم داری
کسانی را که حقّ کرد امین مخزن هفتم طبق باری از مقامی که باید مقامات عالیه
طلب نمایند و اخذ رأی کنند در اطفاء نورش ساعی و جاهدند ولکن در ظهور این
فقره از قرار مذکور جناب سفیر کبیر معین‌الملک میرزا محسن خان ایّده اللّه در
آستانه تشریف نداشته‌اند و این امور نظر بآن است که حضرت پادشاه ایران ایّده
الرّحمن را از محرمان حرم عرفان مکدّر میدانند حقّ شاهد و گواه که این مظلوم
لازال باسبابی که سبب عزّت دولت و ملّت است متمسّک بوده و کفی باللّه شهیداً

\end{otherlanguage}

Va az jahat-e dīgar dāʾere-ye sefārat-e Īrān dar Madīne-ye Kabīre
be-tamām-e qodrat va qovvat be-taḍyīʿ-e īn maẓlūmān mashghūl. Innahum
arādū amran wa-Allāhu arāda amran ākhar. Ḥāl tafakkor farmāʾīd dar ānche
bar omanā-ye Ḥaqq dar har diyār vāred shode. Gāhī nesbat-e serqat va
dozdī dāde-and va hengāmī be-moftarayātī takallom namūde-and ke dar
ʿālam shebh va methl nadāshte. Ḥāl ān janāb be-enṣāf takallom farmāyand,
āyā nesbat-e serqatī ke az jāneb-e sefārat-e kobrā be-raʿīyat-e khod
dāde thamar va atharash dar mamālek-e khāreje chīst? Az īn feqre īn
Maẓlūm khajel shod, na az jahat-e ānke sabab va ʿellat-e taḍyīʿ-e īn
ʿabd būde, balke ʿellat-e khejālat eṭṭelāʿ-e sofarā-ye ajnabīye bar
marāteb-e tadbīr va edrāk-e jamʿī az moʿtabarīn-e Īrān dar sefārat-e
kobrā būde. Mottaham dārī kasānī rā ke Ḥaqq kard amīn-e makhzan-e haftom
ṭabaq. Bārī, az maqāmī ke bāyad maqāmāt-e ʿālīye ṭalab namāyand va
akhdh-e raʾy konand dar eṭfā-ye nūrash sāʿī va jāhedand. Va-lāken dar
ẓohūr-e īn feqre az qarār-e madhkūr Janāb-e Safīr-e Kabīr Moʿīn al-Molk
Mīrzā Moḥsen Khān, ayyadahu Allāh, dar Āstāne tashrīf nadāshte-and. Va
īn omūr naẓar be-ān ast ke Ḥaḍrat-e Pādeshāh-e Īrān, ayyadahu al-Raḥmān,
rā az maḥramān-e ḥaram-e ʿerfān mokaddar mīdānand. Ḥaqq shāhed va gavāh
ke īn Maẓlūm lāzāl be-asbābī ke sabab-e ʿezzat-e dawlat va mellat ast
motamassek būde, wa-kafā bi-Allāhi shahīdan.

From another direction, the Persian embassy in the Great City is
exerting all its power and strength to discredit these wronged ones.
They have intended one thing, and God has intended another. Now reflect
on what has befallen God's trusted ones in every land. At times they
have been accused of theft and robbery, and at others, slanders have
been uttered that have no parallel or likeness in the world. Now speak
fairly: what fruit and effect does the embassy's accusation of theft
against its own subjects have in foreign countries? This affair has made
this Wronged One ashamed---not because it has caused this servant to be
discredited, but because foreign ambassadors have learned the extent of
the judgment and understanding of certain Persian dignitaries at the
embassy. ``You accuse those whom God has made trustees of the treasury
of the seventh heaven.'' In short, they strive and labor to extinguish
the light of the very source from which they should seek exalted
stations and counsel. However, according to what has been reported, His
Excellency the ambassador, Moʿīn al-Molk Mīrzā Moḥsen Khān, may God
assist him, was not in Constantinople when this affair occurred. These
actions arise because they believe His Majesty the Shah of Persia, may
the All-Merciful assist him, to be displeased with the intimates of the
sanctuary of recognition. God bears witness that this Wronged One has
always held fast to whatever promotes the honor of the government and
the people. God is sufficient as a witness.

\subsection{182}\label{section-1468}

\begin{otherlanguage}{arabic}

در وصف اهل بها از قلم اعلی این کلمات نازل انّهم رجال لو یمرّون علی مدائن
الذّهب لا یلتفتون الیها و لو یمرّون علی ملکوت الجمال لا یتوجّهون الیه کذلک
نزّل من القلم الأعلی لأهل البهآء من لدن ناصح علیم و در آخر لوح حضرت
امپراطور پاریس این کلمهٔ علیا نازل هل تفرح بما عندک من الزّخارف بعد اذ
تعلم انّها ستفنی او تسترّ بما تحکم علی شبر من الأرض بعد اذ کلّها لم تکن عند
اهل البهآء الّا کسواد عین نملة میّتة دعها لأهلها ثمّ اقبل الى مقصود
العالمین

\end{otherlanguage}

Dar vaṣf-e ahl-e Bahā az Qalam-e Aʿlā īn kalemāt nāzel: Innahum rijālun
law yamurrūna ʿalā madāʾini al-dhahabi lā yaltafitūna ilayhā, wa-law
yamurrūna ʿalā malakūti al-jamāli lā yatawajjahūna ilayhi. Kadhālika
nuzzila mina al-Qalami al-Aʿlā li-ahli al-Bahāʾ min ladun nāṣiḥin ʿalīm.
Va dar ākhar-e lawḥ-e Ḥaḍrat-e Emperāṭor-e Pārīs īn kaleme-ye ʿolyā
nāzel: Hal tafraḥu bi-mā ʿindaka mina al-zakhārifi baʿda idh taʿlamu
annahā satafnā, aw tastarru bi-mā taḥkumu ʿalā shibrin mina al-arḍi
baʿda idh kulluhā lam takun ʿinda ahli al-Bahāʾ illā ka-sawādi ʿayni
namlatin mayyita. Daʿhā li-ahlihā thumma aqbil ilā Maqṣūdi al-ʿālamīn.

These words were revealed by the Most Exalted Pen in praise of the
people of Bahā: ``They are men who, were they to pass through cities of
gold, would pay them no attention, and were they to pass through the
kingdom of beauty, would not turn toward it. Thus has it been revealed
by the Most Exalted Pen for the people of Bahā, from One who counsels
and knows.'' And at the end of the Tablet to His Majesty the Emperor of
Paris, these exalted words were revealed: ``Do you rejoice in the finery
you possess, when you know that it will perish? Or are you pleased that
you rule over a span of earth, when all of it is, in the eyes of the
people of Bahā, no more than the black of a dead ant's eye? Leave it to
its people, then turn toward the Desire of the worlds.''

\subsection{183}\label{section-1469}

\begin{otherlanguage}{persian}

جز حقّ جلّ جلاله احدی بر آنچه بر این مظلوم وارد شده آگاه نه در هر یوم ذکری
در دائرۀ‌ سفارت کبری در آستانه اصغا میشود سبحان‌اللّه جمیع تدابیر منحصر شده
باسبابی که سبب و علّت تضییع این عبد است غافل از آنکه ذلّت در سبیل الهی نفس
عزّتست در ورق‌ اخبار این کلمات مذکور قوله در تقلّبکاری بعضی از منفیّان عکّا و
تعدّیاتی که از آنان بر بعضی وارد شده الى آخر قوله نزد مظاهر عدل و مطالع
انصاف قصدش معلوم و مقصودش واضح باری بانواع اذیّت و ظلم و اعتساف قیام نمود
لعمر اللّه این مظلوم این منفی را بوطن اعلی تبدیل نمینماید نزد اهل بصر
آنچه در سبیل الهی وارد شود عزّیست مبین و مقامیست کبیر از قبل گفتیم سبحانک
یا الهی لو لا البلایا فی سبیلک من این یظهر مقام عاشقیک و لو لا الرّزایا
فی حبّک بأیّ شیء یثبت شأن مشتاقیک

\end{otherlanguage}

Joz Ḥaqq, jalla jalāluhu, aḥadī bar ānche bar īn Maẓlūm vāred shode āgāh
na. Dar har yawm dhekrī dar dāʾere-ye sefārat-e kobrā dar Āstāne eṣghā
mīshavad. Subḥān Allāh, jamīʿ-e tadābīr monḥaṣer shode be-asbābī ke
sabab va ʿellat-e taḍyīʿ-e īn ʿabd ast, ghāfel az ānke dhellat dar
sabīl-e elāhī nafs-e ʿezzat ast. Dar varaq-e akhbār īn kalemāt madhkūr,
qawluhu: Dar taqallobkārī-ye baʿḍī az manfīyān-e ʿAkkā va taʿaddīyātī ke
az ānān bar baʿḍī vāred shode, ilā ākhir-e qawlihi. Nazd-e maẓāher-e
ʿadl va maṭāleʿ-e enṣāf qaṣdash maʿlūm va maqṣūdash vāḍeḥ. Bārī,
be-anvāʿ-e adhīyat va ẓolm va eʿtesāf qiyām namūd. La-ʿamru Allāh, īn
Maẓlūm īn manfā rā be-vaṭan-e aʿlā tabdīl namīnamāyad. Nazd-e ahl-e
baṣar ānche dar sabīl-e elāhī vāred shavad ʿezzīst mobīn va maqāmīst
kabīr. Az qabl goftīm: Subḥānaka yā Ilāhī, law lā al-balāyā fī sabīlika
min ayna yaẓharu maqāmu ʿāshiqīka, wa-law lā al-razāyā fī ḥubbika
bi-ayyi shayʾin yathbutu shaʾnu mushtāqīk.

No one but God, glorified be His glory, knows what has befallen this
Wronged One. Every day some report is heard in the embassy offices in
Constantinople. Glory be to God! All their planning has become confined
to means of discrediting this servant, heedless that abasement in the
path of God is honor itself. These words appeared in the newspaper:
``Concerning the fraudulent conduct of some of the exiles of ʿAkkā and
the acts of aggression they have committed against certain people,'' and
so on. To the manifestations of justice and the daysprings of fairness,
his intention is known and his purpose clear. In short, he has resorted
to every kind of harassment, oppression, and tyranny. By the life of
God, this Wronged One would not exchange this place of exile for the
most exalted homeland. In the eyes of those endowed with insight,
whatever befalls one in the path of God is manifest honor and a lofty
station. We have previously said: ``Glory be to You, O my God! Were it
not for tribulations in Your path, how would the station of Your lovers
be revealed? Were it not for afflictions endured out of love for You, by
what means would the rank of those who yearn for You be established?''

\subsection{184}\label{section-1470}

\begin{otherlanguage}{persian}

ذلّت بمقامی رسیده که هر یوم بانتشار مفتریات مشغولند ولکن این مظلوم بصبر
جمیل تمسّک جسته ایکاش حضرت پادشاه ممالک ایران آنچه در آستانه وارد شده
صورت آن را طلب فرمایند تا بر حقیقت امر آگاه شوند یا سلطان اقسمک بربّک
الرّحمن در این فقره بنظر عدل ملاحظه فرما هل من ذی عدل یحکم الیوم بما
انزله اللّه فی الکتاب و هل من ذی انصاف ینصف فیما ورد علینا من دون بیّنة و
برهان

\end{otherlanguage}

Dhellat be-maqāmī resīde ke har yawm be-enteshār-e moftarayāt
mashghūland. Va-lāken īn Maẓlūm be-ṣabr-e jamīl tamassok joste. Ey kāsh
Ḥaḍrat-e Pādeshāh-e mamālek-e Īrān ānche dar Āstāne vāred shode ṣūrat-e
ān rā ṭalab farmāyand tā bar ḥaqīqat-e amr āgāh shavand. Yā Sulṭān,
uqsimuka bi-Rabbika al-Raḥmān, dar īn feqre be-naẓar-e ʿadl molāḥeẓe
farmā. Hal min dhī ʿadlin yaḥkumu al-yawma bi-mā anzalahu Allāhu fī
al-Kitāb, wa-hal min dhī inṣāfin yunṣifu fī-mā warada ʿalaynā min dūni
bayyinatin wa-burhān.

Abasement has reached such a point that every day they are occupied with
spreading slanders. Yet this Wronged One has held fast to gracious
patience. If only His Majesty the Shah of Persia would request an
account of what has occurred in Constantinople, so that he might learn
the truth of the matter. O King! I implore you by your Lord, the
All-Merciful: look upon this affair with the eye of justice. Is there
anyone just who will judge today according to what God has revealed in
the Book? Is there anyone fair-minded who will judge fairly concerning
what has befallen Us without evidence or proof?

\subsection{185}\label{section-1471}

\begin{otherlanguage}{persian}

یا شیخ در اطوار نفوس تفکّر نما ساکنین مدائن علم و حکمت متحیّر آیا چه واقع
شده که حزب شیعه که خود را اعلم و ازهد و اتقی از جمیع احزاب عالم میشمردند
در یوم ظهور اعراض نمودند و ظلمی از آن حزب ظاهر شد که شبه نداشته و ندارد
فی‌الجمله تفکّر لازم از اوّل ظهور آن حزب الی حین چه مقدار از علما که آمدند
و یک نفس از ایشان بر کیفیّت ظهور آگاه نه آیا این غفلت را سبب چه بوده لو
نذکره لتنفطر ارکانهم تفکّر لازم بل تفکّر هزارهزار سنه لازم شاید برشحی از
بحر علم فائز شوند و بیابند آنچه را که الیوم از آن غافلند

\end{otherlanguage}

Yā Shaykh, dar aṭvār-e nofūs tafakkor namā. Sākenīn-e madāʾen-e ʿelm va
ḥekmat motaḥayyer. Āyā che vāqeʿ shode ke ḥizb-e Shīʿa ke khod rā aʿlam
va azhad va atqā az jamīʿ-e aḥzāb-e ʿālam mīshomordand dar yawm-e ẓohūr
eʿrāḍ namūdand va ẓolmī az ān ḥizb ẓāher shod ke shebh nadāshte va
nadārad? Fī al-jomle tafakkor lāzem. Az avval-e ẓohūr-e ān ḥizb ilā ḥīn
che meqdār az ʿolamā ke āmadand va yek nafs az īshān bar kayfīyat-e
ẓohūr āgāh na. Āyā īn ghaflat rā sabab che būde? Law nadhkurhu
la-tanfaṭiru arkānuhum. Tafakkor lāzem, bal tafakkor-e hezār hezār sane
lāzem, shāyad be-rashḥī az baḥr-e ʿelm fāʾez shavand va biyāband ānche
rā ke al-yawm az ān ghāfeland.

O Shaykh! Reflect on the conduct of people. The inhabitants of the
cities of knowledge and wisdom are bewildered. What has happened that
the Shiʿi community, who considered themselves more learned, more
ascetic, and more pious than all the world's communities, turned away on
the day of the Manifestation and displayed an oppression without
precedent or parallel? Some reflection is needed. From the beginning of
that community's appearance until now, how many learned men have come,
yet not one of them has understood how the Manifestation would appear.
What has caused this heedlessness? Were We to mention it, their very
foundations would be rent asunder. Reflection is needed---indeed, a
thousand thousand years of reflection---so that they may perhaps attain
a drop from the ocean of knowledge and discover what they are heedless
of today.

\subsection{186}\label{section-1472}

\begin{otherlanguage}{arabic}

قد کنت ماشیاً فی ارض الطّآء مشرق آیات ربّک سمعت حنین المنابر و مناجاتها مع
اللّه تبارک و تعالی نادت و قالت یا اله العالم و سیّد الأمم تری حالنا و ما
ورد علینا من ظلم عبادک قد خلقتنا و اظهرتنا لذکرک و ثنائک اذاً تسمع ما
یقول الغافلون علینا فی ایّامک وعزّتک ذابت اکبادنا و اضطربت ارکاننا آه آه
یا لیت ما خلقتنا و ما اظهرتنا

\end{otherlanguage}

Qad kuntu māshiyan fī arḍi al-Ṭāʾ, mashriqi āyāti Rabbika, samiʿtu
ḥanīna al-manābiri wa-munājātahā maʿa Allāhi tabāraka wa-taʿālā. Nādat
wa-qālat: Yā Ilāha al-ʿālami wa-Sayyida al-umam, tarā ḥālanā wa-mā
warada ʿalaynā min ẓulmi ʿibādika. Qad khalaqtanā wa-aẓhartanā
li-dhikrika wa-thanāʾika, idhan tasmaʿu mā yaqūlu al-ghāfilūna ʿalaynā
fī ayyāmika. Wa-ʿizzatika dhābat akbādunā wa-iḍṭarabat arkānunā. Āh āh,
yā layta mā khalaqtanā wa-mā aẓhartanā.

I was walking in the land of Ṭāʾ, the dawning place of your Lord's
verses, when I heard the lamentation of the pulpits and their intimate
prayers to God, blessed and exalted is He. They called out and said: ``O
God of the world and Lord of the nations! You see our condition and what
has befallen us through the oppression of Your servants. You created us
and brought us into being for Your remembrance and praise; yet now You
hear what the heedless say upon us in Your days. By Your might, our
hearts have melted and our foundations tremble. Alas, alas! If only You
had not created us or brought us into being!''

\subsection{187}\label{section-1473}

\begin{otherlanguage}{persian}

قلوب مقرّبین از این کلمات محترق و زفرات مخلصین از آن متصاعد مکرّر علمای
اعلام را لوجه اللّه نصیحت نمودیم و بافق اعلی دعوت فرمودیم که شاید در ایّام
ظهور از امواج بحر بیان مقصود عالمیان قسمت برند و بالمرّه محروم نمانند

\end{otherlanguage}

Qolūb-e moqarrabīn az īn kalemāt moḥtareq va zafarāt-e mokhleṣīn az ān
motaṣāʿed. Mokarrar ʿolamā-ye aʿlām rā li-wajhi Allāh naṣīḥat namūdīm va
be-Ofoq-e Aʿlā daʿvat farmūdīm ke shāyad dar ayyām-e ẓohūr az amvāj-e
baḥr-e bayān-e Maqṣūd-e ʿālamiyān qesmat barand va be-l-marre maḥrūm
namānand.

The hearts of those brought near are consumed by these words, and the
sighs of the sincere rise because of them. Repeatedly, for the sake of
God, We have counseled the eminent learned men and invited them to the
Most Exalted Horizon, so that in the days of the Manifestation they
might perhaps receive a share of the waves of the ocean of utterance of
the Desire of the worlds and not remain wholly deprived.

\subsection{188}\label{section-1474}

\begin{otherlanguage}{persian}

در اکثری از الواح این نصیحت کبری از سماء رحمت سابقه نازل قلنا یا معشر
الأمرآء و العلمآء اسمعوا النّدآء من افق عکّا انّه یرشدکم و یقرّبکم و یهدیکم
الی مقام جعله اللّه مطلع الوحی و مشرق الأنوار یا اهل العالم قد اتی الاسم
الأعظم من لدن مالک القدم و بشّر العباد بهذا الظّهور الّذی کان مکنوناً فی
العلم و مخزوناً فی کنز العصمة و مرقوماً من القلم الأعلی فی صحف اللّه ربّ
الأرباب یا اهل الشّین أ نسیتم عنایتی و رحمتی الّتی سبقت الأشیآء من لدی
اللّه مالک الرّقاب

\end{otherlanguage}

Dar aktharī az alvāḥ īn naṣīḥat-i kubrā az samā-yi raḥmat-i sābiqih
nāzil. Qulnā yā maʿshar al-umarāʾ wa-l-ʿulamāʾ ismaʿū al-nidāʾ min ufuq
ʿAkkā innahu yurshidukum wa-yuqarribukum wa-yahdīkum ilā maqām jaʿalahu
Allāh maṭlaʿ al-waḥy wa-mashriq al-anwār. Yā ahl al-ʿālam qad atā al-ism
al-aʿẓam min ladun Mālik al-Qidam wa-bashshara al-ʿibād bi-hādhā
al-ẓuhūr alladhī kāna maknūnan fī al-ʿilm wa-makhzūnan fī kanz al-ʿiṣma
wa-marqūman min al-qalam al-aʿlā fī ṣuḥuf Allāh Rabb al-Arbāb. Yā ahl
al-Shīn a-nasītum ʿināyatī wa-raḥmatī allatī sabaqat al-ashyāʾ min ladā
Allāh Mālik al-Riqāb.

In most of the Tablets, this supreme counsel has descended from the
heaven of antecedent mercy. We said: ``O rulers and learned ones! Hear
the call from the horizon of `Akkā. It guides you, draws you near, and
leads you to a station that God has made the dawning place of revelation
and the dayspring of lights. O people of the world! The Most Great Name
has come from the Lord of Eternity and has announced to the servants
this Revelation, which lay concealed in divine knowledge, treasured in
the treasury of infallibility, and inscribed by the Most Exalted Pen in
the scriptures of God, the Lord of lords. O people of Shīn! Have you
forgotten My loving care and My mercy, which preceded all things,
bestowed by God, the Master of all?''

\subsection{189}\label{section-1475}

\begin{otherlanguage}{arabic}

و در کتاب اقدس نازل قل یا معشر العلمآء لا تزنوا کتاب اللّه بما عندکم من
القواعد و العلوم انّه لقسطاس الحقّ بین الخلق قد یوزن ما عند الأمم بهذا
القسطاس الأعظم و انّه بنفسه لو انتم تعلمون تبکی علیکم عین عنایتی لأنّکم ما
عرفتم الّذی دعوتموه فی العشیّ و الاشراق و فی کلّ اصیل و بکور توجّهوا یا قوم
بوجوه بیضآء و قلوب نورآء الی البقعة المبارکة الحمرآء الّتی فیها تنادی
سدرة المنتهی انّه لا اله الّا انا المهیمن القیّوم یا معشر العلمآء فی ایران
هل یقدر احد منکم ان یستنّ معی فی میدان المکاشفة و العرفان او یجول فی
مضمار الحکمة و التّبیان لا وربّی الرّحمن کلّ من علیها فان و هذا وجه ربّکم
العزیز المحبوب یا قوم انّا قدّرنا العلوم لعرفان المعلوم و انتم احتجبتم بها
عن مشرقها الّذى به ظهر کلّ امر مکنون قل هذه سمآء فیها کنز امّ الکتاب لو
انتم تعلمون هذا لهو الّذی به صاحت الصّخرة و نادت السّدرة علی الطّور المرتفع
علی الأرض المبارکة الملک للّه الملک العزیز الودود انّا ما دخلنا المدارس و
ما طالعنا المباحث اسمعوا ما یدعوکم به هذا الأمّیّ الی اللّه الأبدیّ انّه خیر
لکم عمّا کنز فی الأرض لو انتم تفقهون انّ الّذی یأوّل ما نزّل من سمآء الوحی و
یخرجه عن الظّاهر انّه ممّن حرّف کلمة اللّه العلیا و کان من الأخسرین فی کتاب
مبین

\end{otherlanguage}

Va dar Kitāb-i Aqdas nāzil. Qul yā maʿshar al-ʿulamāʾ lā tazinū Kitāb
Allāh bi-mā ʿindakum min al-qawāʿid wa-l-ʿulūm innahu la-qisṭās al-ḥaqq
bayn al-khalq qad yūzanu mā ʿinda al-umam bi-hādhā al-qisṭās al-aʿẓam
wa-innahu bi-nafsihi law antum taʿlamūn. Tabkī ʿalaykum ʿayn ʿināyatī
li-annakum mā ʿaraftum alladhī daʿawtumūhu fī al-ʿashiyy wa-l-ishrāq
wa-fī kull aṣīl wa-bukūr. Tawajjahū yā qawm bi-wujūh bayḍāʾ wa-qulūb
nawrāʾ ilā al-buqʿa al-mubāraka al-ḥamrāʾ allatī fīhā tunādī Sidrat
al-Muntahā innahu lā ilāha illā anā al-Muhaymin al-Qayyūm. Yā maʿshar
al-ʿulamāʾ fī Īrān hal yaqdiru aḥad minkum an yastanna maʿī fī maydān
al-mukāshafa wa-l-ʿirfān aw yajūla fī miḍmār al-ḥikma wa-l-tibyān. Lā
wa-rabbī al-Raḥmān kull man ʿalayhā fān wa-hādhā wajh rabbikum al-ʿAzīz
al-Maḥbūb. Yā qawm innā qaddarnā al-ʿulūm li-ʿirfān al-maʿlūm wa-antum
iḥtajabtum bihā ʿan mashriqihā alladhī bihi ẓahara kull amr maknūn. Qul
hādhihi samāʾ fīhā kanz Umm al-Kitāb law antum taʿlamūn. Hādhā la-huwa
alladhī bihi ṣāḥat al-ṣakhra wa-nādat al-sidra ʿalā al-ṭūr al-murtafiʿ
ʿalā al-arḍ al-mubāraka al-mulk li-llāh al-Malik al-ʿAzīz al-Wadūd. Innā
mā dakhalnā al-madāris wa-mā ṭālaʿnā al-mabāḥith ismaʿū mā yadʿūkum bihi
hādhā al-ummī ilā Allāh al-Abadī innahu khayr lakum ʿammā kuniza fī
al-arḍ law antum tafqahūn. Inna alladhī yuʾawwilu mā nuzzila min samāʾ
al-waḥy wa-yukhrijuhu ʿan al-ẓāhir innahu mimman ḥarrafa kalimat Allāh
al-ʿulyā wa-kāna min al-akhsarīn fī kitāb mubīn.

And in the Most Holy Book it is revealed: ``Say: O learned ones! Do not
weigh the Book of God by the rules and sciences you possess. It is the
true Balance among creation. Whatever the nations possess is weighed in
this Most Great Balance, while it is its own standard, if you only knew.
The eye of My loving care weeps over you, because you have not
recognized the One you invoked at evening and at dawn, and at every dusk
and daybreak. Turn, O people, with radiant faces and luminous hearts,
toward the blessed Crimson Spot, where the Tree of the Utmost Boundary
calls out: `There is no God but Me, the All-Protecting, the
Self-Subsisting.' O learned ones in Iran! Can any of you race with Me in
the arena of unveiling and recognition, or range over the field of
wisdom and exposition? No, by My Lord, the All-Merciful! All who are on
the earth pass away, and this is the face of your Lord, the Almighty,
the Beloved. O people! We ordained the sciences for the recognition of
the Object of knowledge, yet you have made them a veil between
yourselves and their Dayspring, through whom every hidden thing has been
revealed. Say: This is a heaven in which lies the treasure of the Mother
Book, if you only knew. This is the One through whom the Rock cried out
and the Tree proclaimed upon the Mount rising above the blessed land:
`Sovereignty belongs to God, the King, the Almighty, the All-Loving.' We
did not enter the schools or study scholarly treatises. Listen to the
words with which this unlettered One calls you to God, the Eternal. They
are better for you than all that is treasured in the earth, if you only
understood. Whoever interprets what has descended from the heaven of
revelation so as to divert it from its evident meaning is among those
who have distorted God's exalted Word, and is numbered among the
greatest losers in a clear Book.''

\subsection{190}\label{section-1476}

\begin{otherlanguage}{arabic}

عند ذلک سمعنا ضجیج الفطرة قلنا ما لی یا فطرة اسمع فی اللّیالى صریخک و فی
الأیّام ضجیجک و فی الأسحار حنینک قالت یا سیّد العالم الظّاهر بالاسم الأعظم
قد عقر الغافلون ناقتک البیضآء و غرقوا سفینتک الحمرآء و ارادوا ان یطفئوا
نورک و یستروا وجه امرک بذلک ارتفع حنینی و حنین الأشیآء کلّها و النّاس
اکثرهم من الغافلین امروز فطرت باذیال کرم تشبّث نموده و طائف حول گشته

\end{otherlanguage}

ʿInda dhālika samiʿnā ḍajīj al-fiṭra. Qulnā mā lī yā fiṭra asmaʿ fī
al-layālī ṣarīkhaki wa-fī al-ayyām ḍajījaki wa-fī al-asḥār ḥanīnaki.
Qālat yā Sayyid al-ʿĀlam al-ẓāhir bi-l-ism al-aʿẓam qad ʿaqara
al-ghāfilūn nāqataka al-bayḍāʾ wa-gharraqū safīnataka al-ḥamrāʾ wa-arādū
an yuṭfiʾū nūraka wa-yasturū wajh amrika bi-dhālika irtafaʿa ḥanīnī
wa-ḥanīn al-ashyāʾ kullihā wa-l-nās aktharuhum min al-ghāfilīn. Imrūz
fiṭrat bi-adhyāl-i karam tashabbuth namūdih va ṭāʾif-i ḥawl gashtih.

At that moment We heard the lamentation of inborn nature. We said: ``O
inborn nature! Why do I hear your cries at night, your lamentation by
day, and your moaning at dawn?'' She said: ``O Lord of the world,
manifest through the Most Great Name! The heedless have hamstrung Your
white she-camel and sunk Your crimson ark. They have sought to
extinguish Your light and veil the face of Your Cause. Thus my
lamentation and the lamentation of all things have risen, yet most
people remain heedless.'' Today, inborn nature has clung to the hem of
bounty and circles around Him.

\subsection{191}\label{section-1477}

\begin{otherlanguage}{arabic}

یا شیخ احضر لتری ما لا رأت عیون الابداع و تسمع ما لا سمعت آذان الاختراع
لعلّک تخلّص نفسک من طین الأوهام و تتوجّه الی المقام الأعلی الّذى فیه ینادی
المظلوم الملک للّه العزیز الحمید امید آنکه بهمّت آن جناب اجنحهٔ عباد از طین
نفس و هوی مطهّر شود و قابل طیران در هواء محبّت الهی گردد پرهای بطین آلوده
قادر بر طیران نبوده و نیست یشهد بذلک مظاهر العدل و الانصاف ولکنّ القوم فی
ریب مبین

\end{otherlanguage}

Yā shaykh uḥḍur li-tarā mā lā raʾat ʿuyūn al-ibdāʿ wa-tasmaʿ mā lā
samiʿat ādhān al-ikhtirāʿ laʿallaka tukhalliṣ nafsaka min ṭīn al-awhām
wa-tatawajjaha ilā al-maqām al-aʿlā alladhī fīhi yunādī al-Maẓlūm
al-mulk li-llāh al-ʿAzīz al-Ḥamīd. Umīd ānkih bi-himmat-i ān janāb
ajniḥiy-i ʿibād az ṭīn-i nafs va havā muṭahhar shavad va qābil-i ṭayarān
dar havā-yi maḥabbat-i ilāhī gardad. Parhā-yi bi-ṭīn ālūdih qādir bar
ṭayarān nabūdih va nīst. Yashhadu bi-dhālika maẓāhir al-ʿadl wa-l-inṣāf
wa-lākinna al-qawm fī rayb mubīn.

O Shaykh! Come into Our presence so that you may see what the eyes of
creation have never seen and hear what the ears of all that has been
brought into being have never heard. Perhaps you may free yourself from
the clay of vain imaginings and turn toward the Most Exalted Station,
where the Wronged One proclaims: ``Sovereignty belongs to God, the
Almighty, the Praiseworthy.'' It is hoped that, through your efforts,
the wings of the servants may be cleansed of the clay of self and desire
and become capable of flight in the atmosphere of divine love. Wings
soiled with clay have never been, and are not, capable of flight. To
this the manifestations of justice and fairness bear witness, yet the
people remain in evident doubt.

\subsection{192}\label{section-1478}

\begin{otherlanguage}{persian}

یا شیخ از هر جهتی از جهات عباد اعتراضاتی نموده‌اند که قلم از تحریر آن
استغفار مینماید مع‌ذلک نظر برحمت کبری جواب على قدر مراتب ناس داده شد که
شاید از نار نفی و انکار بنور اثبات و اقرار منوّر گردند انصاف کمیاب و عدل
مفقود

\end{otherlanguage}

Yā shaykh az har jihatī az jihāt ʿibād iʿtirāḍātī namūdih-and kih qalam
az taḥrīr-i ān istighfār mīnamāyad. Maʿa-dhālik naẓar bi-raḥmat-i kubrā
javāb ʿalā qadr-i marātib-i nās dādih shud kih shāyad az nār-i nafy va
inkār bi-nūr-i ithbāt va iqrār munavvar gardand. Inṣāf kamyāb va ʿadl
mafqūd.

O Shaykh! From every quarter, people have raised objections that the Pen
seeks forgiveness for recording. Nevertheless, in view of the Most Great
Mercy, answers were given according to people's capacities, so that they
might perhaps leave the fire of negation and denial and be illumined by
the light of affirmation and acknowledgment. Fairness is scarce, and
justice is absent.

\subsection{193}\label{section-1479}

\begin{otherlanguage}{persian}

از جمله این آیات محکمات در جواب بعضی از ملکوت علم الهی ظاهر و نازل یا
ایّها المتوجّه الی انوار الوجه قد احاطت الأوهام سکّان الأرض و منعتهم عن
التّوجّه الی افق الیقین و اشراقه و ظهوراته و انواره بالظّنون منعوا عن
القیّوم یتکلّمون بأهوائهم و لا یشعرون منهم من قال هل الآیات نزّلت قل ای وربّ
السّموات هل اتت السّاعة بل قضت ومظهر البیّنات قد جآءت الحاقّة و اتی الحقّ
بالحجّة و البرهان قد برزت السّاهرة و البریّة فی وجل و اضطراب قد اتت الزّلازل
و ناحت القبائل من خشیة اللّه المقتدر الجبّار قل الصّاخّة صاحت و الیوم للّه
الواحد المختار هل الطّامّة تمّت قل ای وربّ الأرباب هل القیامة قامت بل القیّوم
بملکوت الآیات هل تری النّاس صرعی بلی وربّی العلیّ الأعلی هل انقعرت الأعجاز
بل نسفت الجبال ومالک الصّفات قال این الجنّة و النّار قل الأولی لقائی و
الأخری نفسک یا ایّها المشرک المرتاب قال انّا ما نری المیزان قل ای وربّی
الرّحمن لا یراه الّا اولو الأبصار هل سقطت النّجوم قل ای اذ کان القیّوم فی
ارض السّرّ فاعتبروا یا اولی الأنظار قد ظهرت العلامات کلّها اذ اخرجنا ید
القدرة من جیب العظمة و الاقتدار قد نادی المناد اذ اتی المیعاد و انصعق
الطّوریّون فی تیه الوقوف من سطوة ربّک مالک الایجاد یقول النّاقور هل نفخ فی
الصّور قل بلی وسلطان الظّهور اذ استقرّ علی عرش اسمه الرّحمن قد اضآء الدّیجور
من فجر رحمة ربّک مطلع الأنوار قد سرت نسمة الرّحمن و اهتزّت الأرواح فی قبور
الأبدان کذلک قضی الأمر من لدی اللّه العزیز المنّان قال الّذین غفلوا متی
انفطرت السّمآء قل اذ کنتم فی اجداث الغفلة و الضّلال من الغافلین من یمسح
عینیه و ینظر الیمین و الشّمال قل قد عمیت لیس لک الیوم من ملاذ منهم من قال
هل حشرت النّفوس قل ای وربّی اذ کنت فی مهاد الأوهام و منهم من قال هل نزّل
الکتاب بالفطرة قل انّها فی الحیرة اتّقوا یا اولی الألباب و منهم من قال أ
حشرت اعمی قل بلی وراکب السّحاب قد تزیّنت الجنّة بأوراد المعانی و سعّر السّعیر
من نار الفجّار قل قد اشرق النّور من افق الظّهور و اضآءت الآفاق اذ اتی مالک
یوم المیثاق قد خسر الّذین ارتابوا و ربح من اقبل بنور الیقین الی مطلع
الایقان طوبی لک یا ایّها النّاظر بما نزّل لک هذا اللّوح الّذی منه تطیر
الأرواح احفظه ثمّ اقرأه لعمری انّه باب رحمة ربّک طوبی لمن یقرأه فی العشیّ و
الاشراق انّا نسمع ذکرک فی هذا الأمر الّذی منه اندکّ جبل العلم و زلّت الأقدام
البهآء علیک و علی کلّ مقبل اقبل الی العزیز الوّهاب قد انتهی و ما تمّ اصبر
انّ ربّک هو الصّبّار

\end{otherlanguage}

Az jumlih īn āyāt-i muḥkamāt dar javāb-i baʿḍī az malakūt-i ʿilm-i ilāhī
ẓāhir va nāzil. Yā ayyuhā al-mutawajjih ilā anwār al-wajh qad aḥāṭat
al-awhām sukkān al-arḍ wa-manaʿathum ʿan al-tawajjuh ilā ufuq al-yaqīn
wa-ishrāqihi wa-ẓuhūrātihi wa-anwārihi. Bi-l-ẓunūn muniʿū ʿan al-Qayyūm
yatakallamūn bi-ahwāʾihim wa-lā yashʿurūn. Minhum man qāla hal al-āyāt
nuzzilat qul ay wa-Rabb al-Samāwāt. Hal atat al-sāʿa bal qaḍat wa-Maẓhar
al-Bayyināt. Qad jāʾat al-ḥāqqa wa-atā al-Ḥaqq bi-l-ḥujja wa-l-burhān.
Qad barazat al-sāhira wa-l-bariyya fī wajal wa-iḍṭirāb. Qad atat
al-zalāzil wa-nāḥat al-qabāʾil min khashyat Allāh al-Muqtadir al-Jabbār.
Qul al-ṣākhkha ṣāḥat wa-l-yawm li-llāh al-Wāḥid al-Mukhtār. Hal al-ṭāmma
tammat qul ay wa-Rabb al-Arbāb. Hal al-qiyāma qāmat bal al-Qayyūm
bi-malakūt al-āyāt. Hal tarā al-nās ṣarʿā balā wa-rabbī al-ʿAliyy
al-Aʿlā. Hal inqaʿarat al-aʿjāz bal nusifat al-jibāl wa-Mālik al-Ṣifāt.
Qāla ayna al-janna wa-l-nār qul al-ūlā liqāʾī wa-l-ukhrā nafsuka yā
ayyuhā al-mushrik al-murtāb. Qāla innā mā narā al-mīzān qul ay wa-rabbī
al-Raḥmān lā yarāhu illā ulū al-abṣār. Hal saqaṭat al-nujūm qul ay idh
kāna al-Qayyūm fī arḍ al-sirr fa-ʿtabirū yā ulī al-anẓār. Qad ẓaharat
al-ʿalāmāt kulluhā idh akhrajnā yad al-qudra min jayb al-ʿaẓama
wa-l-iqtidār. Qad nādā al-munād idh atā al-mīʿād wa-inṣaʿaqa al-ṭūriyyūn
fī tīh al-wuqūf min saṭwat rabbika Mālik al-Ījād. Yaqūlu al-nāqūr hal
nufikha fī al-ṣūr qul balā wa-Sulṭān al-Ẓuhūr idh istaqarra ʿalā ʿarsh
ismihi al-Raḥmān. Qad aḍāʾa al-dayjūr min fajr raḥmat rabbika Maṭlaʿ
al-Anwār. Qad sarat nasmat al-Raḥmān wa-ihtazzat al-arwāḥ fī qubūr
al-abdān kadhālika quḍiya al-amr min ladā Allāh al-ʿAzīz al-Mannān. Qāla
alladhīna ghafalū matā infaṭarat al-samāʾ qul idh kuntum fī ajdāth
al-ghafla wa-l-ḍalāl. Min al-ghāfilīn man yamsaḥ ʿaynayhi wa-yanẓur
al-yamīn wa-l-shimāl qul qad ʿamīta laysa laka al-yawm min malādh.
Minhum man qāla hal ḥushirat al-nufūs qul ay wa-rabbī idh kunta fī mihād
al-awhām. Wa-minhum man qāla hal nuzzila al-kitāb bi-l-fiṭra qul innahā
fī al-ḥayra ittaqū yā ulī al-albāb. Wa-minhum man qāla a-ḥushirtu aʿmā
qul balā wa-Rākib al-Saḥāb. Qad tazayyanat al-janna bi-awrād al-maʿānī
wa-suʿʿira al-saʿīr min nār al-fujjār. Qul qad ashraqa al-nūr min ufuq
al-ẓuhūr wa-aḍāʾat al-āfāq idh atā Mālik Yawm al-Mīthāq. Qad khasira
alladhīna irtābū wa-rabiḥa man aqbala bi-nūr al-yaqīn ilā maṭlaʿ
al-īqān. Ṭūbā laka yā ayyuhā al-nāẓir bi-mā nuzzila laka hādhā al-lawḥ
alladhī minhu taṭīr al-arwāḥ. Iḥfaẓhu thumma iqraʾhu la-ʿamrī innahu bāb
raḥmat rabbika ṭūbā li-man yaqraʾuhu fī al-ʿashiyy wa-l-ishrāq. Innā
nasmaʿ dhikraka fī hādhā al-amr alladhī minhu indakka jabal al-ʿilm
wa-zallat al-aqdām. Al-Bahāʾ ʿalayka wa-ʿalā kull muqbil aqbala ilā
al-ʿAzīz al-Wahhāb. Qad intahā wa-mā tamm iṣbir inna rabbaka huwa
al-Ṣabbār.

Among these are the following unequivocal verses, which appeared and
descended from the kingdom of divine knowledge in answer to certain
people: ``O you who turn toward the lights of the Countenance! Vain
imaginings have encompassed the inhabitants of the earth and prevented
them from turning toward the horizon of certainty, its radiance, its
manifestations, and its lights. Through their conjectures they have been
kept from the Self-Subsisting. They speak according to their desires and
do not perceive. Some of them have said: `Have the verses been sent
down?' Say: `Yes, by the Lord of the heavens!' `Has the Hour come?'
`Rather, it has passed, by the Manifestation of clear proofs! The
Inevitable has come, and the True One has appeared with testimony and
proof. The plain of resurrection has emerged, and creation is in fear
and turmoil. The earthquakes have come, and the tribes have lamented in
fear of God, the All-Powerful, the All-Compelling.' Say: `The Deafening
Blast has sounded, and this day belongs to God, the One, the
Unconstrained.' `Has the Overwhelming Calamity been fulfilled?' Say:
`Yes, by the Lord of lords!' `Has the Resurrection arisen?' `Rather, the
Self-Subsisting has appeared with the kingdom of verses.' `Do you see
the people struck down?' `Yes, by My Lord, the Most High, the Most
Exalted!' `Have the trunks been uprooted?' `Rather, the mountains have
been scattered, by the Lord of attributes!' He said: `Where are Paradise
and Hell?' Say: `The first is reunion with Me, and the other is your own
self, O doubting idolater!' He said: `We do not see the Balance.' Say:
`Indeed, by My Lord, the All-Merciful, none sees it but those endowed
with insight.' `Have the stars fallen?' Say: `Yes, when the
Self-Subsisting was in the Land of Mystery. Take heed, then, O people of
discernment!' All the signs appeared when We drew forth the hand of
power from the bosom of grandeur and might. The Crier called out when
the appointed time came, and those upon Sinai were struck senseless in
the wilderness of hesitation by the might of your Lord, the Maker of all
existence. The trumpet says: `Has the horn been sounded?' Say: `Yes, by
the Sovereign of Revelation, when He established Himself upon the throne
of His name, the All-Merciful!' The darkness has been illumined by the
dawn of the mercy of your Lord, the Dawning Place of lights. The breeze
of the All-Merciful has blown, and the spirits have stirred in the tombs
of bodies. Thus has the decree been fulfilled by God, the Almighty, the
All-Bountiful. Those who were heedless said: `When was heaven split
apart?' Say: `When you were in the graves of heedlessness and error.'
Among the heedless is one who rubs his eyes and looks to the right and
left. Say: `You have become blind; today you have no refuge.' Some of
them have said: `Have the souls been gathered?' Say: `Yes, by My Lord,
while you lay in the cradle of vain imaginings!' Some of them have said:
`Has the Book been sent down through inborn nature?' Say: `She is
bewildered. Fear God, O people of understanding!' Some of them have
said: `Have I been gathered blind?' Say: `Yes, by the Rider upon the
clouds!' Paradise has been adorned with the roses of inner meanings, and
the blaze has been kindled by the fire of the wicked. Say: `The light
has shone from the horizon of Revelation, and all horizons have been
illumined, for the Lord of the Day of the Covenant has come.' Those who
doubted have lost, and whoever has turned with the light of certainty
toward the Dawning Place of certitude has gained. Blessed are you, O
beholder, because this Tablet, through which spirits take flight, has
been sent down for you. Preserve it, then read it. By My life! It is the
door to the mercy of your Lord. Blessed is whoever reads it at evening
and at dawn. We hear your mention in this Cause, through which the
mountain of knowledge has been crushed and feet have slipped. Glory rest
upon you and upon everyone who has turned toward the Almighty, the
All-Bestowing. It has ended, yet it is not complete. Be patient; your
Lord is the Ever-Patient.''

\subsection{194}\label{section-1480}

\begin{otherlanguage}{arabic}

هذه آیات انزلناها من قبل فی اوّل ورودنا فی سجن عکّا و ارسلناها الیک لتعرف
ما نطقت به السنتهم الکذبة اذ اتی الأمر بقدرة و سلطان قد تزعزع بنیان
الظّنون و انفطرت سمآء الأوهام و القوم فی مریة و شقاق قد انکروا حجّة اللّه و
برهانه بعد اذ اتی من افق الاقتدار بملکوت الآیات ترکوا ما امروا به و
ارتکبوا ما منعوا عنه فی الکتاب وضعوا الههم اخذوا اهوائهم الا انّهم فی
غفلة و ضلال یقرؤون الآیات و ینکرونها یرون البیّنات و یعرضون عنها الا انّهم
فی ریب عجاب

\end{otherlanguage}

Hādhihi āyāt anzalnāhā min qabl fī awwal wurūdinā fī sijn ʿAkkā
wa-arsalnāhā ilayka li-taʿrifa mā naṭaqat bihi alsinatuhum al-kadhaba
idh atā al-amr bi-qudra wa-sulṭān. Qad tazaʿzaʿa bunyān al-ẓunūn
wa-infaṭarat samāʾ al-awhām wa-l-qawm fī mirya wa-shiqāq. Qad ankarū
ḥujjat Allāh wa-burhānahu baʿda idh atā min ufuq al-iqtidār bi-malakūt
al-āyāt. Tarakū mā umirū bihi wa-irtakabū mā muniʿū ʿanhu fī al-kitāb.
Waḍaʿū ilāhahum akhadhū ahwāʾahum alā innahum fī ghafla wa-ḍalāl.
Yaqraʾūn al-āyāt wa-yunkirūnahā yarawna al-bayyināt wa-yuʿriḍūna ʿanhā
alā innahum fī rayb ʿujāb.

These are verses We sent down previously, upon Our first arrival in the
prison of `Akkā. We have sent them to you so that you may know what
their lying tongues uttered when the Cause appeared with power and
sovereignty. The foundations of conjecture have been shaken, and the
heaven of vain imaginings has been split apart, yet the people remain in
doubt and discord. They have denied God's testimony and His proof after
He came from the horizon of might with the kingdom of verses. They have
abandoned what they were commanded to do and committed what they were
forbidden in the Book. They have set aside their God and embraced their
desires. Indeed, they are in heedlessness and error. They read the
verses and deny them; they see the clear proofs and turn away from them.
Indeed, they are in astonishing doubt.

\subsection{195}\label{section-1481}

\begin{otherlanguage}{arabic}

انّا وصّینا اولیائنا بتقوی اللّه الّذی کان مطلع الأعمال و الأخلاق انّه قائد
جنود العدل فی مدینة البهآء طوبی لمن دخل فی ظلّ رایته النّورآء و تمسّک به
انّه من اصحاب السّفینة الحمرآء الّتی نزّل ذکرها فی قیّوم الأسمآء

\end{otherlanguage}

Innā waṣṣaynā awliyāʾanā bi-taqwā Allāh alladhī kāna maṭlaʿ al-aʿmāl
wa-l-akhlāq innahu qāʾid junūd al-ʿadl fī madīnat al-Bahāʾ. Ṭūbā li-man
dakhala fī ẓill rāyatihi al-nawrāʾ wa-tamassaka bihi innahu min aṣḥāb
al-safīna al-ḥamrāʾ allatī nuzzila dhikruhā fī Qayyūm al-Asmāʾ.

We have enjoined Our loved ones to fear God, which is the source of
deeds and character. It is the commander of the hosts of justice in the
city of Bahā. Blessed is whoever enters beneath its luminous banner and
holds fast to it. Such a person is among the companions of the Crimson
Ark, mention of which was revealed in the Qayyūm al-Asmāʾ.

\subsection{196}\label{section-1482}

\begin{otherlanguage}{arabic}

قل یا حزب اللّه زیّنوا هیاکلکم بطراز الأمانة و الدّیانة ثمّ انصروا ربّکم
بجنود الأعمال و الأخلاق انّا منعناکم عن الفساد و الجدال فی کتبی و صحفی و
زبری و الواحی و ما اردنا بذلک الّا علوّکم و سموّکم تشهد بذلک السّمآء و
انجمها و الشّمس و اشراقها و الأشجار و اوراقها و البحار و امواجها و الأرض
و کنوزها نسأل اللّه ان یمدّ اولیائه و یؤیّدهم علی ما ینبغی لهم فى هذا
المقام المبارک العزیز البدیع

\end{otherlanguage}

Qul yā ḥizb Allāh zayyinū hayākilakum bi-ṭirāz al-amāna wa-l-diyāna
thumma unṣurū rabbakum bi-junūd al-aʿmāl wa-l-akhlāq. Innā manaʿnākum
ʿan al-fasād wa-l-jidāl fī kutubī wa-ṣuḥufī wa-zuburī wa-alwāḥī wa-mā
aradnā bi-dhālika illā ʿuluwwakum wa-sumuwwakum. Tashhadu bi-dhālika
al-samāʾ wa-anjumuhā wa-l-shams wa-ishrāquhā wa-l-ashjār wa-awrāquhā
wa-l-biḥār wa-amwājuhā wa-l-arḍ wa-kunūzuhā. Nasʾalu Allāh an yumidda
awliyāʾahu wa-yuʾayyidahum ʿalā mā yanbaghī lahum fī hādhā al-maqām
al-mubārak al-ʿazīz al-badīʿ.

Say: O people of God! Adorn yourselves with the vesture of
trustworthiness and piety, then aid your Lord with the hosts of deeds
and character. We have forbidden you to engage in corruption and
contention in My books, My scriptures, My writings, and My Tablets. In
this We have desired nothing but your exaltation and elevation. To this
bear witness heaven and its stars, the sun and its radiance, the trees
and their leaves, the seas and their waves, and the earth and its
treasures. We ask God to assist His loved ones and strengthen them to do
what befits them in this blessed, exalted, and wondrous station.

\subsection{197}\label{section-1483}

\begin{otherlanguage}{arabic}

الی ان قلنا فی لوح آخر یا ایّها النّاظر الی الوجه وصّ العباد بتقوی اللّه
تاللّه هو القائد الأوّل فی عساکر ربّک و جنوده الأخلاق المرضیّة و الأعمال
الطّیّبة و بها فتحت فی الأعصار و القرون مدائن الأفئدة و القلوب و نصبت
رایات النّصر و الظّفر علی اعلی الأعلام

\end{otherlanguage}

Ilā an qulnā fī lawḥ ākhar yā ayyuhā al-nāẓir ilā al-wajh waṣṣi al-ʿibād
bi-taqwā Allāh ta-llāh huwa al-qāʾid al-awwal fī ʿasākir rabbika
wa-junūduhu al-akhlāq al-marḍiyya wa-l-aʿmāl al-ṭayyiba wa-bihā futiḥat
fī al-aʿṣār wa-l-qurūn madāʾin al-afʾida wa-l-qulūb wa-nuṣibat rāyāt
al-naṣr wa-l-ẓafar ʿalā aʿlā al-aʿlām.

Further, We said in another Tablet: ``O you who gaze upon the
Countenance! Counsel the servants to fear God. By God! This is the
foremost commander in the armies of your Lord, and its troops are
pleasing character and good deeds. Through these, throughout the ages
and centuries, the cities of hearts and souls have been conquered, and
the banners of victory and triumph have been raised upon the loftiest
heights.''

\subsection{198}\label{section-1484}

\begin{otherlanguage}{arabic}

انّا نذکر لک الأمانة و مقامها عند اللّه ربّک ربّ العرش العظیم انّا قصدنا یوماً
من الأیّام جزیرتنا الخضرآء و لمّا وردنا رأینا انهارها جاریة و اشجارها
ملتفّة و کانت الشّمس تلعب فی خلال الأشجار اذاً توجّهنا الی الیمین رأینا ما
لا یتحرّک القلم علی ذکره و ذکر ما شهدت عین مولی الوری فی ذاک المقام
الألطف الأشرف المبارک الأعلی ثمّ اقبلنا الی الیسار شاهدنا طلعة من طلعات
الفردوس الأعلی قائمة علی عمود من النّور و نادت بأعلی النّدآء یا ملأ الأرض
و السّمآء انظروا جمالی و نوری و ظهوری و اشراقی تاللّه الحقّ انا الأمانة و
ظهورها و حسنها و اجر لمن تمسّک بها و عرف شأنها و مقامها و تشبّث بذیلها انا
الزّینة الکبری لأهل البهآء و طراز العزّ لمن فی ملکوت الانشآء و انا السّبب
الأعظم لثروة العالم و افق الاطمینان لأهل الامکان کذلک انزلنا لک ما یقرّب
العباد الی مالک الایجاد

\end{otherlanguage}

Innā nadhkuru laka al-amānata wa-maqāmahā ʿinda Allāhi rabbika rabbi
al-ʿarshi al-ʿaẓīm. Innā qaṣadnā yawman mina al-ayyāmi jazīratanā
al-khaḍrāʾa wa-lammā waradnā raʾaynā anhārahā jāriyatan wa-ashjārahā
multaffatan wa-kānati al-shamsu talʿabu fī khilāli al-ashjār. Idhan
tawajjahnā ilā al-yamīni raʾaynā mā lā yataḥarraku al-qalamu ʿalā
dhikrihi wa-dhikri mā shahidat ʿaynu mawlā al-warā fī dhāka al-maqāmi
al-alṭafi al-ashrafi al-mubāraki al-aʿlā. Thumma aqbalnā ilā al-yasāri
shāhadnā ṭalʿatan min ṭalaʿāti al-firdawsi al-aʿlā qāʾimatan ʿalā
ʿamūdin mina al-nūri wa-nādat bi-aʿlā al-nidāʾi yā malaʾa al-arḍi
wa-al-samāʾi unẓurū jamālī wa-nūrī wa-ẓuhūrī wa-ishrāqī. Ta-Allāhi
al-ḥaqqi anā al-amānatu wa-ẓuhūruhā wa-ḥusnuhā wa-ajrun li-man tamassaka
bihā wa-ʿarafa shaʾnahā wa-maqāmahā wa-tashabbatha bi-dhaylihā. Anā
al-zīnatu al-kubrā li-ahli al-Bahāʾi wa-ṭirāzu al-ʿizzi li-man fī
malakūti al-inshāʾi wa-anā al-sababu al-aʿẓamu li-tharwati al-ʿālami
wa-ufuqu al-iṭmiʾnāni li-ahli al-imkān. Kadhālika anzalnā laka mā
yuqarribu al-ʿibāda ilā māliki al-ījād.

We speak to you of trustworthiness and its station before God, your
Lord, the Lord of the mighty throne. One day We set out for Our verdant
island. When We arrived, We saw its streams flowing and its trees
intertwined, while the sun played among the trees. Then We turned to the
right and saw what the pen cannot move to describe, nor can it describe
what the eye of the Lord of humanity witnessed in that most delightful,
most noble, blessed, and exalted place. Then We turned to the left and
beheld a figure, one of the figures of the highest Paradise, standing
upon a pillar of light. She called out in the loudest voice: ``O people
of earth and heaven! Behold My beauty, My light, My manifestation, and
My radiance. By God, the True One! I am trustworthiness, its
manifestation and its beauty, and a reward for whoever holds fast to it,
recognizes its standing and station, and clings to its hem. I am the
greatest adornment of the people of Bahāʾ and the ornament of glory for
all in the kingdom of creation. I am the greatest means of the world's
prosperity and the horizon of assurance for all who dwell in the realm
of existence.'' Thus We have sent down for you what brings the servants
nearer to the Lord of creation.

\subsection{199}\label{section-1485}

\begin{otherlanguage}{arabic}

این مظلوم لازال اهل عالم را به ما یرفعهم و یقرّبهم دعوت نموده و از افق
اعلی اشراق نموده آنچه که از برای احدی مجال توقّف و اعراض و اعتراض نمانده
ولکن غافلین را نفعى نبخشیده و لا یزیدهم الّا خسارا

\end{otherlanguage}

Īn maẓlūm lā-zāl ahl-e ʿālam rā be mā yarfaʿuhum wa-yuqarribuhum daʿvat
namūde va az ofoq-e aʿlā eshrāq namūde ānche ke az barā-ye aḥadī majāl-e
tavaqqof va eʿrāż va eʿterāż namānde va-līkan ghāfelīn rā nafʿī
nabakhshīde wa-lā yazīduhum illā khasārā.

This Wronged One has always summoned the people of the world to what
will elevate them and bring them near, and from the most exalted horizon
there has shone forth that which leaves no one any room for hesitation,
rejection, or objection. Yet it has brought no benefit to the heedless,
and only increases their loss.

\subsection{200}\label{section-1486}

\begin{otherlanguage}{persian}

یا شیخ باید علما با حضرت سلطان ایّده اللّه متّحد شوند و در لیالی و ایّام به
ما یرتفع به شأن الدّولة و الملّة تمسّک نمایند این حزب بتمام همّت بتهذیب نفوس
و اصلاح امور متشبّث و مشغولند یشهد بذلک ما نزّل من القلم الأعلی فی هذا
اللّوح المبین چه بسا از امور که سهل و آسان بنظر میآید ولکن اکثرى از آن
غافل و باموری که سبب تضییع ایّام است مشغول

\end{otherlanguage}

Yā Shaykh, bāyad ʿolamā bā Ḥażrat-e Solṭān, ayyadahu Allāh, mottaḥed
shavand va dar layālī va ayyām be mā yartafiʿu bihi shaʾnu al-dawlati
wa-al-millati tamassok namāyand. Īn ḥezb be-tamām-e hemmat be-tahdhīb-e
nofūs va eṣlāḥ-e omūr motashabbeth va mashghūland. Yashhadu bi-dhālika
mā nuzzila mina al-qalami al-aʿlā fī hādhā al-lawḥi al-mubīn. Che basā
az omūr ke sahl va āsān be-naẓar mīāyad va-līkan aktharī az ān ghāfel va
be-omūrī ke sabab-e tażyīʿ-e ayyām ast mashghūl.

O Shaykh! The religious scholars must unite with His Majesty the
sovereign, may God strengthen him, and hold fast, night and day, to what
will elevate the standing of the state and the nation. This community is
devoted with all its energy to refining souls and setting affairs right.
What has been sent down from the Most Exalted Pen in this clear Tablet
bears witness to this. Many things appear simple and easy, yet most
people are heedless of them and occupied with matters that waste their
days.

\subsection{201}\label{section-1487}

\begin{otherlanguage}{persian}

در آستانه یومی از ایّام کمال پاشا نزد مظلوم حاضر و از امور نافعه ذکری
بمیان آمد ذکر نمودند که السن متعدّده آموخته‌اند در جواب ذکر شد عمری را تلف
نموده‌اید باید مثل آن جناب و سائر وکلای دولت مجلسی بیارایند و در آن مجلس
یک لسان از السن مختلفه و همچنین یک خطّ از خطوط موجوده را اختیار نمایند و
یا خطّ و لسانی بدیع ترتیب دهند و در مدارس عالم اطفال را بآن تعلیم فرمایند
در این صورت دارای دو لسان میشوند یکی لسان وطن و دیگری لسانی که عموم اهل
عالم بآن تکلّم نمایند اگر بآنچه ذکر شد تمسّک جویند جمیع ارض قطعهٔ واحده
مشاهده شود و از تعلیم و تعلّم السن مختلفه فارغ و آزاد شوند در حضور قبول
فرمودند و بسیار هم اظهار فرح و مسرّت کردند و بعد بایشان ذکر شد که این
فقره را بوکلا و وزرای دولت برسانند تا حکمش در ممالک جاری گردد و بعد مکرّر
تشریف آوردند و از این فقره ذکرى ننمودند و حال آنکه آنچه ذکر شد سبب اتّحاد
اهل عالم و اتّفاق بوده

\end{otherlanguage}

Dar Āstāne yawmī az ayyām Kamāl Pāshā nazd-e maẓlūm ḥāżer va az omūr-e
nāfeʿe dhekrī be-miyān āmad. Dhekr namūdand ke alson-e motaʿaddede
āmūkhte-and. Dar javāb dhekr shod ʿomrī rā talaf namūde-īd. Bāyad
methl-e ān janāb va sāyer-e vokalā-ye dawlat majlesī biyārāyand va dar
ān majles yek lesān az alson-e mokhtalefe va hamchenīn yek khaṭṭ az
khoṭūṭ-e mawjūde rā ekhtiyār namāyand va yā khaṭṭ va lesānī badīʿ tartīb
dehand va dar madāres-e ʿālam aṭfāl rā be-ān taʿlīm farmāyand. Dar īn
ṣūrat dārā-ye do lesān mīshavand, yekī lesān-e vaṭan va dīgarī lesānī ke
ʿomūm-e ahl-e ʿālam be-ān takallom namāyand. Agar be-ānche dhekr shod
tamassok jūyand jamīʿ-e arż qeṭʿe-ye vāḥede moshāhede shavad va az
taʿlīm va taʿallom-e alson-e mokhtalefe fāregh va āzād shavand. Dar
ḥożūr qabūl farmūdand va besyār ham eẓhār-e faraḥ va masarrat kardand va
baʿd be-īshān dhekr shod ke īn faqare rā be-vokalā va vozarā-ye dawlat
berasānand tā ḥokmash dar mamālek jārī gardad va baʿd mokarrar tashrīf
āvardand va az īn faqare dhekrī nanamūdand va ḥāl ānke ānche dhekr shod
sabab-e etteḥād-e ahl-e ʿālam va ettefāq būde.

One day in Constantinople, Kamāl Pāshā came into the presence of this
Wronged One, and the conversation turned to useful matters. He mentioned
that he had learned several languages. In reply, he was told: ``You have
wasted a lifetime. People such as yourself and the other representatives
of the government should convene an assembly, and in that assembly
choose one language from among the various languages, and likewise one
script from among the existing scripts, or devise a new script and
language, and have children taught these in the schools of the world.
They would then possess two languages: one the language of their
homeland, and the other a language spoken by all the people of the
world. If they hold fast to what has been mentioned, the whole earth
will be seen as a single land, and they will be freed and relieved of
teaching and learning different languages.'' In Our presence he accepted
this and expressed much joy and delight. He was then told to convey this
matter to the representatives and ministers of the government so that
its ruling might be implemented throughout the territories. Afterward he
came repeatedly, but made no mention of this matter, even though what
was proposed would bring about the unity and concord of the people of
the world.

\subsection{202}\label{section-1488}

\begin{otherlanguage}{persian}

امید آنکه دولت ایران بآن تمسّک فرماید و اجرا دارد حال خطّ بدیعی و لسان
جدیدی اختراع شده اگر طالب باشند ارسال شود مقصود آنکه کل تمسّک نمایند
باموری که از زحمت و مشقّت بکاهد و ایّام در آنچه سزاوار است صرف شود و
بانتها رسد انّ اللّه هو المؤیّد العلیم و المدبّر الخبیر

\end{otherlanguage}

Omīd ānke dawlat-e Īrān be-ān tamassok farmāyad va ejrā dārad. Ḥāl
khaṭṭ-e badīʿī va lesān-e jadīdī ekhterāʿ shode agar ṭāleb bāshand ersāl
shavad. Maqṣūd ānke koll tamassok namāyand be-omūrī ke az zaḥmat va
mashaqqat bekāhad va ayyām dar ānche sazāvār ast ṣarf shavad va
be-entehā resad. Inna Allāha huwa al-muʾayyidu al-ʿalīmu wa-al-mudabbiru
al-khabīr.

It is hoped that the government of Iran will adopt and implement it. A
new script and a new language have now been devised; if they wish, these
will be sent. The purpose is that all should hold fast to measures that
lessen toil and hardship, and that their days should be spent and
brought to a close in worthy pursuits. God is indeed the all-knowing
Sustainer and the all-informed Ordainer.

\subsection{203}\label{section-1489}

\begin{otherlanguage}{arabic}

انشآءاللّه ایران مزیّن گردد و فائز شود بآنچه که الی حین از آن محروم بوده
قل یا سلطان همّت فرما تا جمیع اهل عالم بتجلّیات انوار نیّر عدل آن حضرت منوّر
گردند این مظلوم جز بامانت و صدق و صفا و امور نافعه بامری ناظر نبوده و
نیست او را از خائنین مشمرید سبحانک یا الهی و سیّدی و سندی ایّد حضرة
السّلطان علی اجرآء اوامرک و احکامک و اظهار عدلک بین عبادک انّک انت الفضّال
الفیّاض المقتدر القدیر قد اتى امر اللّه بالفضل طوبى للعاملین و طوبى
للعارفین و طوبى لمن تمسّک بالحقّ منقطعاً عمّن فى السّموات و الأرضین

\end{otherlanguage}

In shāʾa Allāh Īrān mozayyan gardad va fāʾez shavad be-ānche ke elā ḥīn
az ān maḥrūm būde. Qul yā Solṭān, hemmat farmā tā jamīʿ-e ahl-e ʿālam
be-tajalliyāt-e anvār-e nayyer-e ʿadl-e ān Ḥażrat monavvar gardand. Īn
maẓlūm joz be-amānat va ṣedq va ṣafā va omūr-e nāfeʿe be-amrī nāẓer
nabūde va nīst. Ū rā az khāʾenīn mashmorīd. Subḥānaka yā ilāhī
wa-sayyidī wa-sanadī, ayyid Ḥaḍrata al-Sulṭāni ʿalā ijrāʾi awāmirika
wa-aḥkāmika wa-iẓhāri ʿadlika bayna ʿibādika. Innaka anta al-faḍḍālu
al-fayyāḍu al-muqtadiru al-qadīr. Qad atā amru Allāhi bi-al-faḍli ṭūbā
li-al-ʿāmilīna wa-ṭūbā li-al-ʿārifīna wa-ṭūbā li-man tamassaka
bi-al-Ḥaqqi munqaṭiʿan ʿamman fī al-samāwāti wa-al-araḍīn.

God willing, Iran will be adorned and attain what it has until now been
deprived of. Say: O sovereign! Make an effort so that all the people of
the world may be illumined by the radiance of the sun of Your Majesty's
justice. This Wronged One has never had, nor does He now have, any aim
other than trustworthiness, truthfulness, sincerity, and useful
endeavors. Do not count Him among the treacherous. Glory be to You, my
God, my Lord, and my Support! Strengthen His Majesty the sovereign to
carry out Your commands and laws and to manifest Your justice among Your
servants. You are indeed the All-Bountiful, the All-Munificent, the
All-Powerful, the Almighty. The Cause of God has come through His grace.
Blessed are those who act, blessed are those who recognize, and blessed
is whoever holds fast to the True One, detached from all who are in the
heavens and on earth.

\subsection{204}\label{section-1490}

\begin{otherlanguage}{arabic}

یا شیخ قصد شاطی بحر اعظم فرما ثمّ ادخل فی السّفینة الحمرآء الّتی قدّرها اللّه
لأهل البهآء فی قیّوم الأسمآء و انّها تمرّ علی البرّ و البحر من دخل فیها نجا
و من اعرض هلک اذا دخلت و فزت ولّ وجهک شطر کعبة اللّه المهیمن القیّوم و قل
اللّهمّ انّی اسألک من بهائک بأبهاه و کلّ بهائک بهیّ اذاً تفتح علی وجهک ابواب
الملکوت و تری ما لا رأت العیون و تسمع ما لا سمعت الآذان انّ المظلوم یعظک
کما وعظک و ما اراد لک الّا الدّخول فی بحر احدیّة اللّه ربّ العالمین هذا یوم
فیه تنادی کلّ الأشیآء و تبشّر النّاس بهذا الظّهور الّذی به ظهر ما کان مکنوناً
مخزوناً فی علم اللّه العزیز الحمید

\end{otherlanguage}

Yā Shaykh, qaṣd-e shāṭī-e baḥr-e aʿẓam farmā. Thumma udkhul fī
al-safīnati al-ḥamrāʾi allatī qaddarahā Allāhu li-ahli al-Bahāʾi fī
Qayyūmi al-Asmāʾi wa-innahā tamurru ʿalā al-barri wa-al-baḥri man
dakhala fīhā najā wa-man aʿraḍa halak. Idhā dakhalta wa-fuzta walli
wajhaka shaṭra kaʿbati Allāhi al-muhaymini al-qayyūmi wa-qul Allāhumma
innī asʾaluka min bahāʾika bi-abhāhu wa-kullu bahāʾika bahiyy. Idhan
tuftaḥu ʿalā wajhika abwābu al-malakūti wa-tarā mā lā raʾati al-ʿuyūnu
wa-tasmaʿu mā lā samiʿati al-ādhān. Inna al-maẓlūma yaʿiẓuka kamā
waʿaẓaka wa-mā arāda laka illā al-dukhūla fī baḥri aḥadiyyati Allāhi
rabbi al-ʿālamīn. Hādhā yawmun fīhi tunādī kullu al-ashyāʾi
wa-tubashshiru al-nāsa bi-hādhā al-ẓuhūri alladhī bihi ẓahara mā kāna
maknūnan makhzūnan fī ʿilmi Allāhi al-ʿazīzi al-ḥamīd.

O Shaykh! Make your way to the shore of the Most Great Ocean. Then enter
the Crimson Ark, which God ordained for the people of Bahāʾ in the
Qayyūm al-Asmāʾ. It passes over land and sea. Whoever enters it is
saved, and whoever turns away perishes. When you have entered and
attained it, turn your face toward the Kaaba of God, the All-Protecting,
the Self-Subsisting, and say: ``O God! I ask You by the most glorious of
Your glory, and all Your glory is glorious.'' Then the gates of the
Kingdom will be opened before you, and you will see what eyes have never
seen and hear what ears have never heard. This Wronged One counsels you
as He has counseled you before, and desires for you nothing but entry
into the ocean of the oneness of God, the Lord of the worlds. This is a
day on which all things call out and announce to humanity the glad
tidings of this Revelation, through which what lay concealed and
treasured in the knowledge of God, the Almighty, the Praiseworthy, has
been made manifest.

\subsection{205}\label{section-1491}

\begin{otherlanguage}{persian}

یا شیخ تغرّدات حمامات بیان را بر افنان سدرهٔ عرفان اصغا نمودی حال تغنّیات
طیور حکمت را که در جنّت علیا مرتفع است بشنو انّه یعرّفک ما کنت غافلاً عنه
اسمع ما نطق به لسان القوّة و القدرة فی کتب اللّه مقصود العارفین فی هذا
الحین ارتفع النّدآء من سدرة المنتهی فی قطب الفردوس الأعلی و امرنی بأن
اذکر لجنابک ما نزل فی الزّبر و الألواح و ما نطق به مبشّری الّذى فدی بنفسه
لهذا النّبإ العظیم و الصّراط المستقیم قال و قوله الحقّ و قد کتبت جوهرة فی
ذکره و هو انّه لا یشار باشارتی و لا بما ذکر فی البیان الی قوله عزّ و جلّ فی
ذکر هذا الظّهور الأعظم و النّبإ العظیم انّه اجلّ و اعلی من ان یکون معروفاً
بدونه او مستشیراً باشارة خلقه و انّنی انا اوّل عبد قد آمنت به و بآیاته و
اخذت من ابکار حدائق جنّة عرفانه حدائق کلماته بلی وعزّته هو الحقّ لا اله الّا
هو کلّ بأمره قائمون انتهی حمامهٔ حقیقی بر اغصان سدرهٔ الهی باین کلمات تغنّی
فرموده طوبی از برای نفسی که باصغای آن فائز شد و از بحور بیان الهی که در
هر کلمه‌ئی مستور اخذ نمود و آشامید و همچنین در مقام دیگر ندای بیان از
اعلی الأغصان مرتفع قوله تبارک و تعالی و فی سنة التّسع انتم کلّ خیر تدرکون
و در مقام دیگر میفرماید و فی سنة التّسع انتم بلقآء اللّه ترزقون انتهی این
نغمات که از طیور مدائن عرفان ظاهر گشته مطابقست به ما انزله الرّحمن فی
الفرقان طوبی للمتبصّرین و طوبی للفائزین

\end{otherlanguage}

Yā Shaykh, tagharrodāt-e ḥamāmāt-e bayān rā bar afnān-e sedre-ye ʿerfān
eṣghā namūdī ḥāl taghanniyāt-e ṭoyūr-e ḥekmat rā ke dar jannat-e ʿolyā
mortafeʿ ast beshnow. Innahu yuʿarrifuka mā kunta ghāfilan ʿanhu. Ismaʿ
mā naṭaqa bihi lisānu al-quwwati wa-al-qudrati fī kutubi Allāhi maqṣūdi
al-ʿārifīn. Fī hādhā al-ḥīni irtafaʿa al-nidāʾu min sidrati al-muntahā
fī quṭbi al-firdawsi al-aʿlā wa-amaranī bi-an adhkura li-janābika mā
nazala fī al-zuburi wa-al-alwāḥi wa-mā naṭaqa bihi mubashshirī alladhī
fadā bi-nafsihi li-hādhā al-nabaʾi al-ʿaẓīmi wa-al-ṣirāṭi al-mustaqīm.
Qāla wa-qawluhu al-ḥaqqu wa-qad katabtu jawharatan fī dhikrihi wa-huwa
innahu lā yushāru bi-ishāratī wa-lā bi-mā dhukira fī al-Bayān. Ilā
qawlihi ʿazza wa-jalla fī dhikri hādhā al-ẓuhūri al-aʿẓami wa-al-nabaʾi
al-ʿaẓīm: Innahu ajallu wa-aʿlā min an yakūna maʿrūfan bi-dūnihi aw
mustashīran bi-ishārati khalqihi wa-innanī anā awwalu ʿabdin qad āmantu
bihi wa-bi-āyātihi wa-akhadhtu min abkāri ḥadāʾiqi jannati ʿirfānihi
ḥadāʾiqa kalimātihi. Balā wa-ʿizzatihi huwa al-Ḥaqqu lā ilāha illā huwa
kullun bi-amrihi qāʾimūn. Intahā. Ḥamāme-ye ḥaqīqī bar aghṣān-e sedre-ye
elāhī be-īn kalemāt taghannī farmūde. Ṭūbā az barā-ye nafsī ke
be-eṣghā-ye ān fāʾez shod va az boḥūr-e bayān-e elāhī ke dar har
kaleme-ī mastūr akhdh namūd va āshāmīd. Va hamchenīn dar maqām-e dīgar
nedā-ye bayān az aʿlā al-aghṣān mortafeʿ, qawluhu tabāraka wa-taʿālā:
Wa-fī sanati al-tisʿi antum kulla khayrin tudrikūn. Va dar maqām-e dīgar
mīfarmāyad: Wa-fī sanati al-tisʿi antum bi-liqāʾi Allāhi turzaqūn.
Intahā. Īn naghamāt ke az ṭoyūr-e madāʾen-e ʿerfān ẓāher gashte moṭābeq
ast be mā anzalahu al-Raḥmānu fī al-Furqān. Ṭūbā li-al-mutabaṣṣirīna
wa-ṭūbā li-al-fāʾizīn.

O Shaykh! You have listened to the warbling of the doves of utterance
upon the boughs of the Tree of recognition; now hear the songs of the
birds of wisdom rising in the highest Paradise. He makes known to you
what you were heedless of. Hear what the tongue of strength and power
has spoken in the books of God, the Desire of those who recognize Him.
At this moment the call has risen from the Tree beyond which there is no
passing, at the heart of the highest Paradise, and has commanded Me to
relate to you what has been sent down in the scriptures and Tablets, and
what My Herald spoke, who sacrificed Himself for this momentous
Announcement and this straight Path. He said---and His word is the
truth---``I have written a gem concerning Him, namely, that He cannot be
indicated by My allusions, nor by what has been mentioned in the
Bayān.'' Further on He said, exalted and glorious is He, concerning this
Most Great Revelation and this momentous Announcement: ``He is too
majestic and exalted to be known through anyone other than Himself, or
to seek counsel through the indications of His creatures. I am indeed
the first servant to have believed in Him and His verses, and from the
virgin gardens of the paradise of His recognition I have gathered the
gardens of His words. Yes, by His might! He is the True One; there is no
god but Him. All stand by His command.'' Here ends the quotation. The
true Dove has sung these words upon the branches of the divine Tree.
Blessed is the soul that has attained to hearing them and has drawn and
drunk from the oceans of divine utterance concealed in every word.
Likewise, in another passage, the call of the Bayān rises from the
highest branches. His words, blessed and exalted is He, are: ``In the
year nine you will attain all good.'' And in another passage He says:
``In the year nine you will be granted the meeting with God.'' Here ends
the quotation. These melodies, which have issued from the birds of the
cities of recognition, accord with what the All-Merciful has sent down
in the Qurʾān. Blessed are those who see clearly, and blessed are those
who attain.

\subsection{206}\label{section-1492}

\begin{otherlanguage}{persian}

یا شیخ لعمر اللّه فرات رحمت جاری و بحر بیان موّاج و شمس ظهور مشرق و منیر
بقلب فارغ و صدر منشرح و لسان صدق مبین این کلمات عالیات که از مبشّر یعنی
نقطهٔ اولی ظاهر شده قرائت نما قوله عزّ بیانه مخاطباً لحضرة العظیم هذا ما قد
وعدناک قبل حین الّذى اجبناک اصبر حتّی یقضی عن البیان تسعة فاذاً قل فتبارک
اللّه احسن المبدعین قل هو نبأ لم یحط بعلمه احد الّا اللّه ولکن انتم یومئذ
لا تعلمون انتهی و در سنهٔ تسع این ظهور اعظم از مشرق ارادهٔ الهی مشرق و
لائح لا ینکره الّا کلّ غافل مریب نسأل اللّه ان یؤیّد عباده علی الرّجوع الیه و
الاستغفار عمّا عملوا به فی الحیوة الباطلة انّه هو التّوّاب الغفور الرّحیم و
در مقام دیگر میفرماید انّنی انا اوّل عبد قد آمنت به و بآیاته انتهی و
همچنین در بیان فارسی میفرماید انّه هو الّذی ینطق فی کلّ شأن انّنی انا اللّه
الی آخر بیانه جلّ و عزّ و مقصود از ذکر ربوبیّت و الوهیّت از قبل ذکر شد قد
خرقنا الأحجاب و اظهرنا ما یقرّب النّاس الی اللّه مالک الرّقاب طوبی لمن فاز
بالعدل و الانصاف فی هذا الفضل الّذی احاط من فی السّموات و الأرضین امراً من
لدی اللّه ربّ العالمین

\end{otherlanguage}

Yā Shaykh, la-ʿamru Allāh, Forāt-e raḥmat jārī va baḥr-e bayān mavvāj va
shams-e ẓohūr moshreq va monīr. Be-qalb-e fāregh va ṣadr-e monshareḥ va
lesān-e ṣedq-e mobīn īn kalemāt-e ʿāliyāt ke az mobashsher yaʿnī
Noqṭe-ye Ūlā ẓāher shode qerāʾat namā. Qawluhu ʿazza bayānuhu mukhāṭiban
li-Ḥaḍrati al-ʿAẓīm: Hādhā mā qad waʿadnāka qabla ḥīna alladhī ajabnāka
iṣbir ḥattā yuqḍā ʿani al-Bayāni tisʿatun fa-idhan qul fa-tabāraka
Allāhu aḥsanu al-mubdiʿīn. Qul huwa nabaʾun lam yuḥiṭ bi-ʿilmihi aḥadun
illā Allāhu wa-lākin antum yawmaʾidhin lā taʿlamūn. Intahā. Va dar
sane-ye tesʿ īn ẓohūr-e aʿẓam az mashreq-e erāde-ye elāhī moshreq va
lāʾeḥ. Lā yunkiruhu illā kullu ghāfilin murīb. Nasʾalu Allāha an
yuʾayyida ʿibādahu ʿalā al-rujūʿi ilayhi wa-al-istighfāri ʿammā ʿamilū
bihi fī al-ḥayāti al-bāṭilati innahu huwa al-tawwābu al-ghafūru
al-raḥīm. Va dar maqām-e dīgar mīfarmāyad: Innanī anā awwalu ʿabdin qad
āmantu bihi wa-bi-āyātihi. Intahā. Va hamchenīn dar Bayān-e Fārsī
mīfarmāyad: Innahu huwa alladhī yanṭiqu fī kulli shaʾnin innanī anā
Allāh. Ilā ākhiri bayānihi jalla wa-ʿazza. Va maqṣūd az dhekr-e
robūbiyyat va olūhiyyat az qabl dhekr shod. Qad kharaqnā al-aḥjāba
wa-aẓharnā mā yuqarribu al-nāsa ilā Allāhi māliki al-riqāb. Ṭūbā li-man
fāza bi-al-ʿadli wa-al-inṣāfi fī hādhā al-faḍli alladhī aḥāṭa man fī
al-samāwāti wa-al-araḍīna amran min ladā Allāhi rabbi al-ʿālamīn.

O Shaykh! By the life of God! The Euphrates of mercy is flowing, the
ocean of utterance is surging, and the sun of Revelation is shining and
radiant. With an unburdened heart, an open breast, and a clearly
truthful tongue, read these exalted words that have issued from the
Herald, namely, the Primal Point. His words, exalted is His utterance,
addressed to His Honor ʿAẓīm, are: ``This is what We promised you
before, when We answered you: Wait until nine have elapsed from the
Bayān; then say: `Blessed be God, the best of originators!' Say: This is
an announcement whose knowledge no one but God has encompassed, yet on
that day you will not know.'' Here ends the quotation. In the year nine
this Most Great Revelation shone and gleamed from the dawning-place of
the divine will. None denies it except the heedless and doubting. We ask
God to help His servants turn back to Him and seek forgiveness for what
they have done in this vain life. He is indeed the Ever-Relenting, the
Ever-Forgiving, the Merciful. And in another passage He says: ``I am
indeed the first servant to have believed in Him and His verses.'' Here
ends the quotation. Likewise, in the Persian Bayān He says: ``He is the
One who declares in every circumstance: `I am indeed God'\,''---and so
continues His utterance, majestic and exalted is He. The intended
meaning of the mention of lordship and divinity has already been
explained. We have rent the veils and revealed what brings people nearer
to God, the Master of all. Blessed is whoever attains justice and
fairness amid this grace, which has encompassed all who are in the
heavens and on earth, by a command from God, the Lord of the worlds.

\subsection{207}\label{section-1493}

\begin{otherlanguage}{persian}

یا شیخ نغمات انجیل را باذن انصاف بشنو قوله عزّ بیانه که از بعد اخبار
میفرماید و امّا ذلک الیوم و تلک السّاعة فلا یعلم بهما احد و لا الملائکة
الّذین فی السّمآء و لا الابن الّا الأب و مقصود از اب در این مقام حقّ جلّ
جلاله است و اوست مربّی حقیقی و معلّم معنوی

\end{otherlanguage}

Yā Shaykh, naghamāt-e Enjīl rā be-odhon-e enṣāf beshnow. Qawluhu ʿazza
bayānuhu ke az baʿd ekh bār mīfarmāyad: Wa-ammā dhālika al-yawmu
wa-tilka al-sāʿatu fa-lā yaʿlamu bihimā aḥadun wa-lā al-malāʾikatu
alladhīna fī al-samāʾi wa-lā al-ibnu illā al-ab. Va maqṣūd az ab dar īn
maqām Ḥaqq jalla jalālohu ast va ūst morabbī-e ḥaqīqī va moʿallem-e
maʿnavī.

O Shaykh! Hear the melodies of the Gospel with the ear of fairness. His
words, exalted is His utterance, in which He speaks of what is to come,
are: ``But of that day and that hour no one knows, neither the angels
who are in heaven nor the Son, but only the Father.'' By ``Father'' in
this passage is meant God, exalted be His glory. He is the true Educator
and spiritual Teacher.

\subsection{208}\label{section-1494}

\begin{otherlanguage}{persian}

یوئیل میفرماید لأنّ یوم الرّبّ عظیم و مخوف جدّا فمن یطیقه اوّل در بیان جلیل
مرقوم در انجیل میفرماید احدی بحین ظهور آگاه نه لم یحط به الّا اللّه العلیم
الخبیر و در ثانی عظمت ظهور را ذکر مینماید و همچنین در فرقان میفرماید عمّ
یتساءلون عن النّبإ العظیم اینست آن نبأئی که عظمتش در اکثری از کتب قبل و
بعد مذکور اینست آن نبأئی که ارتعدت به فرائص العالم الّا من شآء اللّه
الحافظ النّاصر المعین چنانچه بابصار ظاهره مشاهده شد که جمیع عباد و من فى
البلاد منقلب و متحیّر گشتند الّا من شآء اللّه

\end{otherlanguage}

Yūʾīl mīfarmāyad: «Li-anna yawma al-Rabbi ʿaẓīmun wa-makhūfun jiddan
fa-man yuṭīquhu.» Avval dar bayān-i jalīl-i marqūm dar Injīl mīfarmāyad
aḥadī bi-ḥīn-i ẓuhūr āgāh nah: «Lam yuḥiṭ bihi illā Allāhu al-ʿAlīmu
al-Khabīr.» Va dar thānī ʿaẓamat-i ẓuhūr rā dhikr mīnamāyad va hamchunīn
dar Furqān mīfarmāyad: «ʿAmma yatasāʾalūna ʿani al-nabaʾi al-ʿaẓīm.»
Īnast ān nabaʾī kih ʿaẓamatash dar aktharī az kutub-i qabl va baʿd
madhkūr. Īnast ān nabaʾī kih «irtaʿadat bihi farāʾiṣu al-ʿālam illā man
shāʾa Allāhu al-Ḥāfiẓu al-Nāṣiru al-Muʿīn». Chunānchih bi-abṣār-i
ẓāhirih mushāhadih shud kih jamīʿ-i ʿibād va man fī al-bilād munqalib va
mutaḥayyir gashtand illā man shāʾa Allāh.

Joel says: ``For the day of the Lord is great and very terrible; who can
endure it?'' First, in the glorious utterance recorded in the Gospel, He
says that no one knows the time of the Revelation: ``None has
encompassed it except God, the All-Knowing, the All-Aware.'' Second, He
speaks of the greatness of the Revelation. Likewise, in the Qurʾān He
says: ``About what are they questioning one another? About the momentous
announcement.'' This is that announcement whose greatness is mentioned
in most of the Books, both earlier and later. This is that announcement
at which ``the limbs of the world trembled, except those whom God willed
to spare, the Preserver, the Helper, the Sustainer.'' As was seen with
outward eyes, all the servants and inhabitants of the lands became
agitated and bewildered, except those whom God willed to spare.

\subsection{209}\label{section-1495}

\begin{otherlanguage}{persian}

یا شیخ امر عظیم است و نبأ عظیم بصبر و سکون در آیات باهرات و کلمات عالیات
و ما ظهر فی هذه الأیّام تفکّر نما شاید اسرار مکنونهٔ در کتب را بیابی و بر
هدایت عباد همّت نمائی ندای ارمیا را بسمع حقیقی بشنو میفرماید آه لأنّ ذلک
الیوم عظیم و لیس مثله اگر آن جناب بانصاف نظر نمایند بر عظمت یوم آگاه
شوند اسمع ندآء هذا النّاصح العلیم و لا تجعل نفسک محرومة عن الرّحمة الّتی
سبقت الوجود من الغیب و الشّهود نغمهٔ حضرت داود را بشنو میفرماید من یقودنی
الی المدینة المحصّنة مدینهٔ محصّنه عکّاست که سجن اعظم نامیده شده و دارای حصن
و قلعهٔ محکم است

\end{otherlanguage}

Yā Shaykh! Amr ʿaẓīm ast va nabaʾ ʿaẓīm. Bi-ṣabr va sukūn dar āyāt-i
bāhirāt va kalimāt-i ʿāliyāt va mā ẓahara fī hādhihi al-ayyām tafakkur
namā shāyad asrār-i maknūnih-yi dar kutub rā biyābī va bar hidāyat-i
ʿibād himmat namāʾī. Nidā-yi Irmiyā rā bi-samʿ-i ḥaqīqī bishnaw.
Mīfarmāyad: «Āh li-anna dhālika al-yawma ʿaẓīmun wa-laysa mithlahu.»
Agar ān janāb bi-inṣāf naẓar namāyand bar ʿaẓamat-i yawm āgāh shavand.
Ismaʿ nidāʾa hādhā al-Nāṣiḥi al-ʿAlīm wa-lā tajʿal nafsaka maḥrūmatan
ʿani al-raḥmati allatī sabaqati al-wujūda mina al-ghaybi wa-al-shuhūd.
Naghmih-yi Ḥaḍrat-i Dāvūd rā bishnaw. Mīfarmāyad: «Man yaqūdunī ilā
al-madīnati al-muḥaṣṣanah.» Madīnih-yi muḥaṣṣanih ʿAkkāst kih Sijn-i
Aʿẓam nāmīdih shudih va dārā-yi ḥiṣn va qalʿih-yi muḥkam ast.

O Shaykh! Great is the Cause, and momentous the announcement. With
patience and composure, reflect on the resplendent verses, the exalted
words, and what has appeared in these days, so that you may discover the
mysteries hidden in the Books and strive to guide the servants. Hear the
call of Jeremiah with the true ear. He says: ``Alas! For that day is
great, and none is like it.'' If you look with fairness, you will
recognize the greatness of the Day. Hear the call of this all-knowing
Counsellor, and do not deprive yourself of the mercy that has preceded
all existence, both unseen and seen. Hear the melody of David. He says:
``Who will lead me to the fortified city?'' The fortified city is ʿAkkā,
which has been named the Most Great Prison and possesses a strong
fortress and citadel.

\subsection{210}\label{section-1496}

\begin{otherlanguage}{persian}

یا شیخ اقرأ ما نطق به اشعیا فى کتابه قوله علی جبل عال اصعدی یا مبشّرة
صهیون ارفعی صوتک بقوّة یا مبشّرة اورشلیم ارفعی لا تخافی قولی لمدن یهوذا هو
ذا الهک هو ذا السّیّد الرّبّ بقوّة یأتی و ذراعه تحکم له امروز جمیع علامات
نمودار مدینهٔ بزرگی از آسمان نازل شده و صهیون از ظهور حقّ مهتزّ و مسرور چه
که ندآء اللّه را از کلّ جهات اصغا نموده امروز اورشلیم ببشارت جدید فائز چه
که مقام جمّیز سرو آزاد ایستاده اورشلیم محلّ زیارت جمیع احزاب عالم است و به
قدس نامیده شده و آن و صهیون و فلسطین جمیع در این اراضی واقع اینست که
میفرماید طوبی لمن هاجر الی عکّا

\end{otherlanguage}

Yā Shaykh! Iqraʾ mā naṭaqa bihi Ishaʿyā fī kitābihi. Qawluhu: «ʿAlā
jabalin ʿālin iṣʿadī yā mubashshirata Ṣahyūn, irfaʿī ṣawtaki bi-quwwatin
yā mubashshirata Ūrashalīm, irfaʿī lā takhāfī, qūlī li-muduni Yahūdhā:
Huwa dhā ilāhuki, huwa dhā al-Sayyidu al-Rabbu bi-quwwatin yaʾtī
wa-dhirāʿuhu taḥkumu lahu.» Imrūz jamīʿ-i ʿalāmāt namūdār. Madīnih-yi
buzurgī az āsmān nāzil shudih va Ṣahyūn az ẓuhūr-i Ḥaqq muhtazz va
masrūr chih kih nidāʾ Allāh rā az kull-i jihāt iṣghā namūdih. Imrūz
Ūrashalīm bi-bishārat-i jadīd fāʾiz chih kih maqām-i jummayz sarv-i āzād
īstādih. Ūrashalīm maḥall-i ziyārat-i jamīʿ-i aḥzāb-i ʿālam ast va
bi-Quds nāmīdih shudih va ān va Ṣahyūn va Filasṭīn jamīʿ dar īn arāḍī
vāqiʿ. Īnast kih mīfarmāyad: «Ṭūbā li-man hājara ilā ʿAkkā.»

O Shaykh! Read what Isaiah has spoken in his book. His words are:
``Ascend a high mountain, O herald of good news to Zion! Raise your
voice with strength, O herald of good news to Jerusalem! Raise it, do
not fear. Say to the cities of Judah: `Here is your God! Here the Lord
God comes with strength, and His arm rules for Him.'\,'' Today all the
signs are manifest. A great city has descended from heaven, and Zion is
stirred and joyful at the revelation of God, for it has heard the call
of God from every direction. Today Jerusalem has received new glad
tidings, for in place of the sycamore stands the noble cypress.
Jerusalem is a place of pilgrimage for all the peoples of the world and
has been named Quds. It, Zion, and Palestine all lie in these lands.
Thus He says: ``Blessed is whoever migrates to ʿAkkā.''

\subsection{211}\label{section-1497}

\begin{otherlanguage}{persian}

عاموص میفرماید انّ الرّبّ یزمجر من صهیون و یأتی صوته من اورشلیم فتنوح مراعی
الرّعاة و ییبس رأس الکرمل و کرمل در کتاب الهی به کوم اللّه و کرم اللّه
نامیده شده کوم تپه را میگویند و این مقامیست که در این ایّام از فضل دارای
ظهور خباء مجد بر آن مرتفع گشته طوبی للواردین و طوبى للمقبلین و همچنین
میفرماید یأتی الهنا و لا یصمت

\end{otherlanguage}

ʿĀmūṣ mīfarmāyad: «Inna al-Rabba yuzamjiru min Ṣahyūn wa-yaʾtī ṣawtuhu
min Ūrashalīm fa-tanūḥu marāʿī al-ruʿāti wa-yaybasu raʾsu al-Karmal.» Va
Karmal dar Kitāb-i Ilāhī bi-Kawm Allāh va Karm Allāh nāmīdih shudih.
Kawm tappih rā mīgūyand va īn maqāmīst kih dar īn ayyām az faḍl-i
Dārā-yi Ẓuhūr khibāʾ-i majd bar ān murtafiʿ gashtih. Ṭūbā li-al-wāridīn
wa-ṭūbā li-al-muqbilīn. Va hamchunīn mīfarmāyad: «Yaʾtī ilāhunā wa-lā
yaṣmut.»

Amos says: ``The Lord roars from Zion, and His voice comes from
Jerusalem; the shepherds' pastures mourn, and the summit of Carmel
withers.'' Carmel has been named the ``Mound of God'' and the ``Vineyard
of God'' in the Book of God. Kawm means a mound, and this is the place
on which, in these days, through the grace of the Lord of Revelation,
the tabernacle of glory has been raised. Blessed are those who arrive,
and blessed are those who turn toward it. He also says: ``Our God comes
and does not remain silent.''

\subsection{212}\label{section-1498}

\begin{otherlanguage}{persian}

یا شیخ در این بیان مقصود عالمیان که به عاموص فرموده تفکّر نما قوله فاستعدّ
للقآء الهک یا اسرائیل فانّه هو ذا الّذی صنع الجبال و خلق الرّیح و اخبر
الانسان ما هو فکره الّذی یجعل الفجر ظلاماً و یمشی علی مشارف الأرض یهوه اله
الجنود اسمه میفرماید فجر را تاریک مینماید مقصود آنکه اگر در حین ظهور
مکلّم طور نفسی خود را صبح صادق داند بقوّت و قدرت الهی تاریک میشود صبح
کاذبست و خود را صادق میداند ویل له و ویل لمن اتّبعه من دون بیّنة من اللّه
ربّ العالمین

\end{otherlanguage}

Yā Shaykh! Dar īn bayān-i Maqṣūd-i ʿālamiyān kih bi-ʿĀmūṣ farmūdih
tafakkur namā. Qawluhu: «Fa-istaʿidd li-liqāʾi ilāhika yā Isrāʾīl.
Fa-innahu huwa dhā alladhī ṣanaʿa al-jibāla wa-khalaqa al-rīḥa
wa-akhbara al-insāna mā huwa fikruhu, alladhī yajʿalu al-fajra ẓalāman
wa-yamshī ʿalā mashārifi al-arḍ, Yahwah ilāhu al-junūdi ismuhu.»
Mīfarmāyad fajr rā tārīk mīnamāyad. Maqṣūd ānkih agar dar ḥīn-i ẓuhūr-i
Mukallim-i Ṭūr nafsī khud rā ṣubḥ-i ṣādiq dānad bi-quvvat va qudrat-i
ilāhī tārīk mīshavad. Ṣubḥ-i kādhib ast va khud rā ṣādiq mīdānad. Waylun
lahu wa-waylun li-man ittabaʿahu min dūni bayyinatin mina Allāhi Rabbi
al-ʿālamīn.

O Shaykh! Reflect on this utterance of the Desire of all peoples, which
He spoke to Amos. His words are: ``Prepare to meet your God, O Israel.
For here is He who formed the mountains, created the wind, and tells
humanity what His thought is; who turns dawn into darkness and walks
upon the heights of the earth---Yahweh, the God of hosts, is His name.''
He says that He darkens the dawn. The meaning is that if, at the
appearance of Him who spoke on Sinai, a soul considers himself the true
dawn, he will be darkened by divine strength and power. He is a false
dawn, yet considers himself true. Woe to him, and woe to whoever follows
him without a clear proof from God, the Lord of the worlds.

\subsection{213}\label{section-1499}

\begin{otherlanguage}{arabic}

اشعیا میفرماید یسمو الرّبّ وحده فی ذلک الیوم و در عظمت ظهور میفرماید ادخل
الی الصّخرة و اختبئ فی التّراب من امام هیبة الرّبّ و من بهآء عظمته و در مقام
دیگر میفرماید تفرح البرّیّة و الأرض الیابسة و یبتهج القفر و یزهر کالنّرجس
یزهر ازهاراً و یبتهج ابتهاجاً و یرنّم یدفع الیه مجد لبنان بهآء کرمل و شارون
هم یرون مجد الرّبّ بهآء الهنا انتهى

\end{otherlanguage}

Ishaʿyā mīfarmāyad: «Yasmū al-Rabbu waḥdahu fī dhālika al-yawm.» Va dar
ʿaẓamat-i ẓuhūr mīfarmāyad: «Udkhul ilā al-ṣakhrati wa-ikhtabiʾ fī
al-turābi min amāmi haybati al-Rabbi wa-min bahāʾi ʿaẓamatihi.» Va dar
maqām-i dīgar mīfarmāyad: «Tafraḥu al-barriyyatu wa-al-arḍu al-yābisatu
wa-yabtahiju al-qafru wa-yuzhiru ka-al-narjis. Yuzhiru izhāran
wa-yabtahiju ibtihājan wa-yurannimu. Yudfaʿu ilayhi majdu Lubnān, bahāʾu
Karmal wa-Shārūn. Hum yarawna majda al-Rabbi, bahāʾa ilāhinā.» Intahā.

Isaiah says: ``The Lord alone will be exalted on that day.'' Concerning
the greatness of the Revelation, he says: ``Enter the rock and hide in
the dust before the awe of the Lord and the splendor of His majesty.''
Elsewhere he says: ``The wilderness and the dry land will rejoice; the
desert will be glad and blossom like the narcissus. It will blossom
abundantly, rejoice greatly, and sing. The glory of Lebanon will be
given to it, the splendor of Carmel and Sharon. They will see the glory
of the Lord, the splendor of our God.'' Here ends the quotation.

\subsection{214}\label{section-1500}

\begin{otherlanguage}{arabic}

این فقرات احتیاج بتفسیر ندارد بمثابهٔ آفتاب مشرق و ظاهر و بمثابهٔ نور ساطع
و لامع هر منصفی از عرف این بیان بحدیقهٔ عرفان راه یابد و فائز شود بآنچه
اکثر اهل عالم از آن محجوب و ممنوعند قل اتّقوا اللّه یا قوم و لا تتّبعوا
شبهات النّاعقین الّذین نقضوا عهد اللّه و میثاقه و انکروا رحمته الّتی سبقت من
فی السّموات و الأرضین

\end{otherlanguage}

Īn faqarāt iḥtiyāj bi-tafsīr nadārad. Bi-mathābih-yi āftāb mushriq va
ẓāhir va bi-mathābih-yi nūr sāṭiʿ va lāmiʿ. Har munṣifī az ʿarf-i īn
bayān bi-ḥadīqih-yi ʿirfān rāh yābad va fāʾiz shavad bi-ānchih akthar-i
ahl-i ʿālam az ān maḥjūb va mamnūʿand. Qul ittaqū Allāha yā qawmu wa-lā
tattabiʿū shubuhāti al-nāʿiqīna alladhīna naqaḍū ʿahda Allāhi
wa-mīthāqahu wa-ankarū raḥmatahu allatī sabaqat man fī al-samāwāti
wa-al-araḍīn.

These passages need no interpretation. They are as manifest and radiant
as the sun, as brilliant and shining as light. Through the fragrance of
this utterance, every fair-minded person will find the way to the garden
of recognition and attain what most of the world's people are veiled
from and denied. Say: Fear God, O people, and do not follow the doubts
raised by those who clamor, who have broken God's covenant and pledge
and denied His mercy, which has preceded all who are in the heavens and
on earth.

\subsection{215}\label{section-1501}

\begin{otherlanguage}{persian}

و همچنین میفرماید قولوا لخائفی القلوب تشدّدوا لا تخافوا هو ذا الهکم انتهی
این آیهٔ مبارکه دلیل است بر عظمت ظهور و عظمت امر چه که نفخهٔ صور عالم را
مضطرب مینماید زلزله و خوف کل را احاطه میکند طوبی از برای نفسی که بنور
توکّل و انقطاع منوّر گشت شدائد آن یوم او را منع نکند و نترساند کذلک نطق
لسان البیان امراً من لدی الرّحمن انّه هو المقتدر القویّ الغالب القدیر حال از
برای صاحبان آذان و ابصار واجب که در این کلمات عالیات که در هر یک بحور
معانی و بیان مستور تفکّر نمایند لعلّ بیان مالک ادیان عباد خود را بکمال روح
و ریحان بمقصد اقصی و ذروهٔ علیا که مطلع افق این ندا است فائز فرماید

\end{otherlanguage}

Va hamchunīn mīfarmāyad: «Qūlū li-khāʾifī al-qulūbi tashaddadū lā
takhāfū huwa dhā ilāhukum.» Intahā. Īn āyih-yi mubārakih dalīl ast bar
ʿaẓamat-i ẓuhūr va ʿaẓamat-i amr chih kih nafkhih-yi ṣūr ʿālam rā
muḍṭarib mīnamāyad. Zilzilih va khawf kull rā iḥāṭih mīkunad. Ṭūbā az
barā-yi nafsī kih bi-nūr-i tavakkul va inqiṭāʿ munavvar gasht.
Shadāʾid-i ān yawm ū rā manʿ nakunad va natarsānad. Kadhālika naṭaqa
Lisānu al-Bayāni amran min ladā al-Raḥmān. Innahu huwa al-Muqtadiru
al-Qawiyyu al-Ghālibu al-Qadīr. Ḥāl az barā-yi ṣāḥibān-i ādhān va abṣār
vājib kih dar īn kalimāt-i ʿāliyāt kih dar har yak buḥūr-i maʿānī va
bayān mastūr tafakkur namāyand laʿall bayān-i Mālik-i adyān ʿibād-i khud
rā bi-kamāl-i rūḥ va rayḥān bi-maqṣad-i aqṣā va dhurvih-yi ʿulyā kih
maṭlaʿ-i ufuq-i īn nidā ast fāʾiz farmāyad.

Likewise, he says: ``Say to the fearful of heart: `Be strong, do not
fear. Here is your God.'\,'' Here ends the quotation. This blessed verse
testifies to the greatness of the Revelation and the greatness of the
Cause, for the trumpet blast throws the world into turmoil. Trembling
and fear encompass everyone. Blessed is the soul illumined by the light
of trust and detachment. The hardships of that Day will neither hold him
back nor frighten him. Thus has the Tongue of Utterance spoken at the
command of the All-Merciful. He is indeed the All-Powerful, the Strong,
the All-Conquering, the Almighty. Now it is incumbent on those with ears
and eyes to reflect on these exalted words, in each of which oceans of
meaning and utterance lie concealed, so that the utterance of the Lord
of religions may enable His servants, with utmost joy and delight, to
attain the furthest goal and the highest summit, the dawning-place on
the horizon of this call.

\subsection{216}\label{section-1502}

\begin{otherlanguage}{persian}

یا شیخ اگر نفحات بیان را اقلّ از سمّ ابره بیابی عالم و عالمیان را بگذاری و
بانوار وجه حضرت مقصود توجّه نمائی باری در کلمات حضرت روح معانی لاتحصى
مستور ذکر امور کثیره فرموده ولکن نظر بعدم وجود سامعین و ناظرین اکثری را
ستر نموده چنانچه میفرماید ولکن لا تستطیعون ان تحتملوا الآن آن مطلع وحی
میفرماید در آن یوم اخبار بعد را حضرت موعود ذکر میفرماید چنانچه در کتاب
اقدس و الواح حضرات ملوک و لوح رئیس و لوح فؤاد از قلم اعلی اکثر اموری که
در ارض واقع از قبل جاری و نازل گشته

\end{otherlanguage}

Yā Shaykh! Agar nafaḥāt-i bayān rā aqall az samm-i ibrih biyābī ʿālam va
ʿālamiyān rā bigudhārī va bi-anvār-i vajh-i Ḥaḍrat-i Maqṣūd tavajjuh
namāʾī. Bārī dar kalimāt-i Ḥaḍrat-i Rūḥ maʿānī-yi lā tuḥṣā mastūr.
Dhikr-i umūr-i kathīrih farmūdih va-lākin naẓar bi-ʿadam-i vujūd-i
sāmiʿīn va nāẓirīn aktharī rā satr namūdih. Chunānchih mīfarmāyad:
«Wa-lākin lā tastaṭīʿūna an taḥtamilū al-ān.» Ān Maṭlaʿ-i Vaḥy
mīfarmāyad dar ān yawm akhbār-i baʿd rā Ḥaḍrat-i Mawʿūd dhikr
mīfarmāyad. Chunānchih dar Kitāb-i Aqdas va Alvāḥ-i Ḥaḍarāt-i Mulūk va
Lawḥ-i Raʾīs va Lawḥ-i Fuʾād az Qalam-i Aʿlā akthar-i umūrī kih dar arḍ
vāqiʿ az qabl jārī va nāzil gashtih.

O Shaykh! If you were to perceive the fragrances of utterance through an
opening smaller than the eye of a needle, you would abandon the world
and its peoples and turn toward the lights of the face of the Desired
One. Countless meanings lie hidden in the words of the Spirit. He
mentioned many things, but, because there were none to hear or see, He
concealed most of them. Thus He says: ``But you cannot bear them now.''
That Dawning-Place of Revelation says that on that Day the Promised One
will speak of things to come. Thus, in the Most Holy Book, the Tablets
to the sovereigns, the Tablet of the Chief, and the Tablet of Fuʾād,
most of the events occurring on earth flowed from the Most Exalted Pen
and were revealed beforehand.

\subsection{217}\label{section-1503}

\begin{otherlanguage}{arabic}

در کتاب اقدس نازل یا ارض الطّآء لا تحزنی من شیء قد جعلک اللّه مطلع فرح
العالمین لو یشآء یبارک سریرک بالّذى یحکم بالعدل و یجمع اغنام اللّه الّتی
تفرّقت من الذّئاب انّه یواجه اهل البهآء بالفرح و الانبساط الا انّه من جوهر
الخلق لدی الحقّ علیه بهآء اللّه و بهآء من فی ملکوت الأمر فی کلّ حین این
آیات از قبل نازل ولکن در این حین این آیه نازل الهی الهی یدعوک البهآء و
یسألک بأنوار وجهک و امواج بحر امرک و تجلّیات شمس بیانک ان تؤیّد السّلطان
علی العدل و الانصاف و لو ترید بارک به سریر الأمر و الحکم انّک انت المقتدر
علی ما تشآء لا اله الّا انت السّامع المجیب افرحی یا ارض الطّآء بما جعلک
اللّه افق النّور بما ولد فیک مطلع الظّهور و سمّیت بهذا الاسم الّذی به لاح نیّر
الفضل و اشرقت السّموات و الأرضون سوف تنقلب فیک الأمور و یحکم علیک جمهور
النّاس انّ ربّک هو العلیم المحیط اطمئنّی بفضل ربّک انّه لا تنقطع عنک لحظات
الألطاف سوف یأخذک الاطمینان بعد الاضطراب کذلک قضی الأمر فی کتاب بدیع

\end{otherlanguage}

Dar Kitāb-i Aqdas nāzil: «Yā arḍa al-Ṭāʾ lā taḥzanī min shayʾin qad
jaʿalaki Allāhu maṭlaʿa faraḥi al-ʿālamīn. Law yashāʾu yubāriku sarīraki
bi-alladhī yaḥkumu bi-al-ʿadli wa-yajmaʿu aghnāma Allāhi allatī
tafarraqat mina al-dhiʾāb. Innahu yuwājihu ahla al-Bahāʾi bi-al-faraḥi
wa-al-inbisāṭ. Alā innahu min jawhari al-khalqi ladā al-Ḥaqq. ʿAlayhi
bahāʾu Allāhi wa-bahāʾu man fī malakūti al-amri fī kulli ḥīn.» Īn āyāt
az qabl nāzil va-lākin dar īn ḥīn īn āyih nāzil: «Ilāhī Ilāhī yadʿūka
al-Bahāʾu wa-yasʾaluka bi-anwāri wajhika wa-amwāji baḥri amrika
wa-tajalliyāti shamsi bayānika an tuʾayyida al-sulṭāna ʿalā al-ʿadli
wa-al-inṣāf. Wa-law turīdu bārik bihi sarīra al-amri wa-al-ḥukm. Innaka
anta al-Muqtadiru ʿalā mā tashāʾu lā ilāha illā anta al-Sāmiʿu al-Mujīb.
Ifraḥī yā arḍa al-Ṭāʾ bi-mā jaʿalaki Allāhu ufuqa al-nūr bi-mā wulida
fīki Maṭlaʿu al-Ẓuhūr wa-summīti bi-hādhā al-ismi alladhī bihi lāḥa
nayyiru al-faḍli wa-ashraqati al-samāwātu wa-al-araḍūn. Sawfa tanqalibu
fīki al-umūru wa-yaḥkumu ʿalayki jumhūru al-nās. Inna rabbaki huwa
al-ʿAlīmu al-Muḥīṭ. Iṭmaʾinnī bi-faḍli rabbaki innahu lā tanqaṭiʿu ʿanki
laḥaẓātu al-alṭāf. Sawfa yaʾkhudhuki al-iṭmiʾnānu baʿda al-iḍṭirāb.
Kadhālika quḍiya al-amru fī kitābin badīʿ.»

In the Most Holy Book it is revealed: ``O Land of Ṭāʾ! Do not grieve
over anything. God has made you the dawning-place of the joy of all
peoples. If He wills, He will bless your throne with one who rules with
justice and gathers the sheep of God that the wolves have scattered. He
will meet the people of Bahāʾ with joy and gladness. Indeed, in the
sight of the True One, he is among the choicest of creation. Upon him be
the glory of God and the glory of all who are in the kingdom of His
Cause at every moment.'' These verses were revealed before, but at this
time this verse is revealed: ``My God, my God! Bahāʾ calls upon You and
asks You, by the lights of Your face, the waves of the ocean of Your
Cause, and the effulgences of the sun of Your utterance, to assist the
sovereign in justice and fairness. If You wish, bless through him the
throne of authority and government. You are indeed powerful over
whatever You will. There is no God but You, the All-Hearing, the
Responsive. Rejoice, O Land of Ṭāʾ, for God has made you the horizon of
light, since the Dawning-Place of Revelation was born within you, and
you were called by this name through which the luminary of grace shone
and the heavens and earth were illumined. Affairs within you will be
transformed, and the people at large will rule over you. Your Lord is
indeed the All-Knowing, the All-Encompassing. Be assured of your Lord's
grace, for His loving regard will never be withdrawn from you.
Tranquillity will overtake you after turmoil. Thus has the matter been
decreed in a wondrous Book.''

\subsection{218}\label{section-1504}

\begin{otherlanguage}{persian}

و همچنین در لوح فؤاد و لوح پاریس و الواح اخری نازل شده آنچه که هر صاحب
انصافی شهادت میدهد بر قدرت و عظمت و علم حقّ جلّ جلاله اگر بعدل ملاحظه
نمایند بر سرّ این آیهٔ مبارکهٔ لا رطب و لا یابس الّا فی کتاب مبین آگاه
میشوند و ادراک مینمایند ولکنّ الیوم منع النّاس اعراضهم عن ادراک ما نزّل
بالحقّ من لدن منزل قدیم سبحان‌اللّه آیات باهرات از جمیع جهات احاطه نموده
مع‌ذلک اکثری از عباد از مشاهده و آگاهی آن محرومند از حقّ میطلبیم توفیق عطا
فرماید تا جمیع بر لآلی مستورهٔ در اصداف بحر اعظم آگاه شوند و به لک الحمد
یا اله العالم ناطق گردند

\end{otherlanguage}

Va hamchunīn dar Lawḥ-i Fuʾād va Lawḥ-i Pārīs va Alvāḥ-i ukhrā nāzil
shudih ānchih kih har ṣāḥib-i inṣāfī shahādat mīdahad bar qudrat va
ʿaẓamat va ʿilm-i Ḥaqq jalla jalāluhu. Agar bi-ʿadl mulāḥiẓih namāyand
bar sirr-i īn āyih-yi mubārakih-yi «lā raṭbin wa-lā yābisin illā fī
kitābin mubīn» āgāh mīshavand va idrāk mīnamāyand. Wa-lākinna al-yawma
manaʿa al-nāsa iʿrāḍuhum ʿan idrāki mā nuzzila bi-al-ḥaqqi min ladun
Munzilin Qadīm. Subḥān Allāh! Āyāt-i bāhirāt az jamīʿ-i jihāt iḥāṭih
namūdih maʿa-dhālik aktharī az ʿibād az mushāhadih va āgāhī-yi ān
maḥrūmand. Az Ḥaqq mīṭalabīm tawfīq ʿaṭā farmāyad tā jamīʿ bar laʾālī-yi
mastūrih-yi dar aṣdāf-i Baḥr-i Aʿẓam āgāh shavand va bi-«laka al-ḥamdu
yā Ilāha al-ʿālam» nāṭiq gardand.

Likewise, in the Tablet of Fuʾād, the Tablet of Paris, and other
Tablets, things have been revealed that lead every fair-minded person to
testify to the power, greatness, and knowledge of God, exalted be His
glory. If they consider them justly, they will recognize and understand
the mystery of this blessed verse: ``Nothing moist or dry exists but
that it is in a clear Book.'' Yet today their turning away has prevented
people from understanding what has been revealed in truth by the Ancient
Revealer. Glory be to God! Resplendent signs surround them on every
side, yet most of the servants are deprived of seeing and recognizing
them. We ask God to grant His assistance, so that all may become aware
of the pearls concealed within the shells of the Most Great Ocean and
proclaim: ``Praise be to You, O God of the world!''

\subsection{219}\label{section-1505}

\begin{otherlanguage}{persian}

یا معشر المنصفین در امواج بحر بیان و علم الهی نظر نمائید و تفکّر کنید تا
بلسان ظاهر و باطن شهادت دهید بأنّ عنده علم کلّ شیء فی الکتاب لا یعزب عن
علمه من شیء انّه اظهر ما کان مکنوناً اذ استوی علی عرش البیان فی المآب آنچه
نازل شده کلمه بکلمه در ارض ظاهر گشته و میشود از برای احدی مجال اعراض و
اعتراض نه ولکن چون انصاف مخذول و مستور اکثری باوهامات خود تکلّم مینمایند

\end{otherlanguage}

Yā maʿshara al-munṣifīn! Dar amvāj-i baḥr-i bayān va ʿilm-i ilāhī naẓar
namāʾīd va tafakkur kunīd tā bi-lisān-i ẓāhir va bāṭin shahādat dahīd
bi-anna ʿindahu ʿilma kulli shayʾin fī al-kitāb. Lā yaʿzubu ʿan ʿilmihi
min shayʾ. Innahu aẓhara mā kāna maknūnan idh istawā ʿalā ʿarshi
al-bayāni fī al-maʾāb. Ānchih nāzil shudih kalimih bi-kalimih dar arḍ
ẓāhir gashtih va mīshavad. Az barā-yi aḥadī majāl-i iʿrāḍ va iʿtirāḍ nah
va-lākin chūn inṣāf makhdhūl va mastūr aktharī bi-avhāmāt-i khud
takallum mīnamāyand.

O company of the fair-minded! Look at the waves of the ocean of divine
utterance and knowledge, and reflect, so that with both your outward and
inward tongues you may testify that with Him is the knowledge of all
things in the Book. Nothing eludes His knowledge. He revealed what had
been hidden when He ascended the throne of utterance in the place of
return. Whatever has been revealed has appeared, and continues to
appear, word for word on earth. No one has grounds for turning away or
objecting. Yet because fairness is forsaken and concealed, most people
speak according to their own imaginings.

\subsection{220}\label{section-1506}

\begin{otherlanguage}{arabic}

الهی الهی لا تمنع عبادک عن التّوجّه الی نور الایقان الّذی اشرق من افق سمآء
مشیّتک و لا تجعلهم یا الهی محرومین عن بحور آیاتک ای ربّ هم عبادک فی بلادک
و ارقّائک فی دیارک ان لم ترحمهم من یرحمهم خذ یا الهی ایادی الّذین غرقوا فی
بحر الأوهام و خلّصهم بقدرتک و سلطانک ثمّ انقذهم بذراعی اقتدارک انّک انت
المقتدر علی ما تشآء و فی یمینک زمام من فی السّموات و الأرضین

\end{otherlanguage}

Ilāhī Ilāhī lā tamnaʿ ʿibādaka ʿani al-tawajjuhi ilā nūri al-īqāni
alladhī ashraqa min ufuqi samāʾi mashīʾatika wa-lā tajʿalhum yā Ilāhī
maḥrūmīna ʿan buḥūri āyātik. Ay Rabbi hum ʿibāduka fī bilādika
wa-ariqqāʾuka fī diyārik. In lam tarḥamhum man yarḥamhum? Khudh yā Ilāhī
ayādiya alladhīna ghariqū fī baḥri al-awhāmi wa-khalliṣhum bi-qudratika
wa-sulṭānika thumma anqidhhum bi-dhirāʿay iqtidārik. Innaka anta
al-Muqtadiru ʿalā mā tashāʾu wa-fī yamīnika zimāmu man fī al-samāwāti
wa-al-araḍīn.

My God, my God! Do not prevent Your servants from turning toward the
light of certitude that has shone from the horizon of the heaven of Your
will, and do not deprive them, O my God, of the oceans of Your verses. O
Lord! They are Your servants in Your lands and Your bondsmen in Your
dominions. If You do not have mercy on them, who will? Take the hands, O
my God, of those who have drowned in the ocean of vain imaginings,
deliver them through Your power and sovereignty, and rescue them with
the arms of Your might. You are indeed powerful over whatever You will,
and in Your right hand are the reins of all who are in the heavens and
on earth.

\subsection{221}\label{section-1507}

\begin{otherlanguage}{persian}

و همچنین نقطهٔ اولی میفرماید بچشم او او را ببینید و اگر بچشم غیر ملاحظه
کنید هرگز بشناسائی و آگاهی فائز نشوید انتهی و این فقره مخصوصست باین ظهور
اعظم طوبی للمنصفین و همچنین میفرماید نطفهٔ یکسالهٔ ظهور او اقوی است از کلّ
بیان این بشارتهای بیان و کتب قبل مکرّر در کتب عدیده باسامی مختلفه ذکر شده
لعلّ النّاس ینصفون فیما اشرق و لاح من افق ارادة اللّه ربّ العرش العظیم

\end{otherlanguage}

Va hamchunīn Nuqṭiy-i Ūlā mī-farmāyad: bi chashm-i ū ū rā bibīnīd va
agar bi chashm-i ghayr mulāḥiẓih kunīd hargiz bi shināsāʾī va āgāhī
fāʾiz nashavīd. Intahā. Va īn faqarih makhṣūṣ ast bi īn ẓuhūr-i aʿẓam.
Ṭūbā li-l-munṣifīn. Va hamchunīn mī-farmāyad: nuṭfiy-i yik-sāliy-i
ẓuhūr-i ū aqvā ast az kull-i Bayān. Īn bishārathā-yi Bayān va kutub-i
qabl mukarrar dar kutub-i ʿadīdih bi asāmī-yi mukhtalifih dhikr shudih,
laʿalla al-nāsa yunṣifūna fīmā ashraqa va lāḥa min ufuqi irādati Allāhi
Rabbi al-ʿarshi al-ʿaẓīm.

The Primal Point likewise says: ``See Him through His own eyes, for if
you regard Him through the eyes of another, you will never attain
recognition and knowledge.'' End of quotation. This passage applies
specifically to this Most Great Revelation. Blessed are the fair-minded!
He also says: ``The one-year-old embryo of His Revelation is mightier
than the whole Bayán.'' These glad tidings of the Bayán and the earlier
Books have been repeatedly mentioned in numerous books under various
names, so that people may judge fairly what has shone and gleamed from
the horizon of the will of God, the Lord of the mighty throne.

\subsection{222}\label{section-1508}

\begin{otherlanguage}{arabic}

یا شیخ قل لملإ البیان در این کلمهٔ مبارکه تفکّر نمائید میفرماید جمیع بیان
ورقیست از اوراق جنّت او انصفوا یا قوم و لا تکونوا من الأخسرین فی کتاب
اللّه ربّ العالمین امروز سدرهٔ مبارکه با اثمار جنیّهٔ جدیدهٔ بدیعه امام وجهت
موجود انظر الیها منقطعاً عن دونها کذلک نطق لسان القوّة و القدرة فی هذا
المقام الّذی جعله اللّه مزیّناً بقدوم اسمه الأعظم و نبإه العظیم

\end{otherlanguage}

Yā Shaykh, qul li-malaʾi al-Bayān: dar īn kalimiy-i mubārakih tafakkur
namāyīd. Mī-farmāyad: jamīʿ-i Bayān varaqī-st az avrāq-i jannat-i ū.
Anṣifū yā qawmu va lā takūnū mina al-akhsarīna fī kitābi Allāhi Rabbi
al-ʿālamīn. Imrūz Sidriy-i Mubārakih bā athmār-i janiyyiy-i jadīdiy-i
badīʿih imām-i vajhat mawjūd. Unẓur ilayhā munqaṭiʿan ʿan dūnihā.
Kadhālika naṭaqa lisānu al-quwwati va al-qudrati fī hādhā al-maqāmi
alladhī jaʿalahu Allāhu muzayyanan bi-qudūmi ismihi al-Aʿẓami va
nabaʾihi al-ʿaẓīm.

O Shaykh! Say to the people of the Bayán: Reflect on these blessed
words. He says: ``The whole Bayán is a leaf among the leaves of His
Paradise.'' Be fair, O people, and do not be among those accounted the
greatest losers in the Book of God, the Lord of the worlds. Today the
blessed Lote-Tree stands before your face, bearing fresh, wondrous
fruits ready to be gathered. Look upon it, detached from all else. Thus
has the tongue of strength and power spoken in this station, which God
has adorned with the arrival of His Most Great Name and His mighty
Announcement.

\subsection{223}\label{section-1509}

\begin{otherlanguage}{persian}

و همچنین میفرماید من اوّل ذلک الأمر الی قبل ان یکمل تسعة کینونات الخلق لم
تظهر و انّ کلّ ما قد رأیت من النّطفة الی ما کسوناه لحماً ثمّ اصبر حتّی تشهد
خلق الآخر قل فتبارک اللّه احسن الخالقین و همچنین از اقتدار ظهور ذکر
فرموده و قال حلّ لمن یظهره اللّه ان یردّ من لم یکن فوق الأرض اعلی منه اذ
ذلک خلق فی قبضته و کلّ له قانتون فانّ لکم بعد حین امر ستعلمون و میفرماید
فاعرف بالیقین الأقطع و الأمر المثبت الأحتم بأنّه جلّ جلاله و عزّ اعزازه و
قدّس اقداسه و کبّر کبریائه و مجّد شؤوناته یعرّف کلّ شیء نفسه بنفسه فمن یقدر
ان یعرفه بغیره الی قوله عزّ و جلّ ایّاک ایّاک ایّام ظهوره ان تحتجب بالواحد
البیانیّة فانّ ذلک الواحد خلق عنده و ایّاک ایّاک ان تحتجب بکلمات ما نزّل فی
البیان الی قوله تعالی و لا تنظر الیه الّا بعینه فانّ من ینظر الیه بعینه
یدرکه والّا یحتجب ان اردت اللّه و لقائه فأرده و انظر الیه و همچنین
میفرماید اگر یک آیه از آیات من یظهره اللّه را تلاوت کنی اعزّتر خواهد بود
از آنکه کلّ بیان را ثبت کنی زیرا که آن روز آن یک آیه ترا نجات میدهد ولی
کلّ بیان نمیدهد انتهی

\end{otherlanguage}

Va hamchunīn mī-farmāyad: min awwali dhālika al-amri ilā qabli an
yakmula tisʿatan kaynūnātu al-khalqi lam taẓhar, va inna kulla mā qad
raʾayta mina al-nuṭfati ilā mā kasawnāhu laḥman, thumma iṣbir ḥattā
tashhada khalqa al-ākhar. Qul fa-tabāraka Allāhu aḥsanu al-khāliqīn. Va
hamchunīn az iqtidār-i ẓuhūr dhikr farmūdih, va qāla: ḥalla li-man
yuẓhiruhu Allāhu an yarudda man lam yakun fawqa al-arḍi aʿlā minhu, idh
dhālika khalqun fī qabḍatihi va kullun lahu qānitūn. Fa-inna lakum baʿda
ḥīnin amran sa-taʿlamūn. Va mī-farmāyad: fa-iʿraf bi-l-yaqīni al-aqṭaʿi
va al-amri al-muthbati al-aḥtami bi-annahu jalla jalāluhu va ʿazza
iʿzāzuhu va quddisa aqdāsuhu va kubbira kibriyāʾuhu va mujjida
shuʾūnātuhu yuʿarrifu kulla shayʾin nafsahu bi-nafsihi, fa-man yaqdiru
an yaʿrifahu bi-ghayrihi? Ilā qawlihi ʿazza va jalla: iyyāka iyyāka
ayyāma ẓuhūrihi an taḥtajiba bi-l-wāḥidi al-bayāniyyati, fa-inna dhālika
al-wāḥida khalqun ʿindahu, va iyyāka iyyāka an taḥtajiba bi-kalimāti mā
nuzzila fī al-Bayān. Ilā qawlihi taʿālā: va lā tanẓur ilayhi illā
bi-ʿaynihi, fa-inna man yanẓuru ilayhi bi-ʿaynihi yudrikuhu va illā
yaḥtajib. In aradta Allāha va liqāʾahu fa-aridhu va unẓur ilayhi. Va
hamchunīn mī-farmāyad: agar yik āyih az āyāt-i Man Yuẓhiruhu Allāh rā
tilāvat kunī aʿazz-tar khvāhad būd az ānkih kull-i Bayān rā thabt kunī,
zīrā kih ān rūz ān yik āyih turā najāt mī-dahad valī kull-i Bayān
namī-dahad. Intahā.

He likewise says: ``From the beginning of that Cause until before nine
is completed, the realities of creation have not appeared. All that you
have seen is from the stage of the embryo to that which We clothed with
flesh. Then wait until you witness the other creation. Say: `Blessed,
then, is God, the best of creators!'\,'' He has also spoken of the power
of the Revelation, saying: ``It is lawful for Him Whom God will make
manifest to reject one above whom there is no one on earth, for such a
one is a creature in His grasp, and all are obedient to Him. After a
time you will indeed have a Cause that you will come to know.'' He also
says: ``Know, then, with the most decisive certainty and as an
established and most irrevocable matter, that He---glorified be His
glory, exalted be His might, sanctified be His holiness, magnified be
His grandeur, and glorified be His attributes---makes Himself known to
all things through Himself. Who, then, can know Him through anyone
else?'' And further on He, exalted and glorified be He, says: ``Beware,
beware lest, in the days of His appearance, you be veiled by the Váḥid
of the Bayán, for that Váḥid is a creature in His presence. Beware,
beware lest you be veiled by the words of what has been revealed in the
Bayán.'' And further on He, exalted be He, says: ``Look upon Him only
through His own eyes, for whoever looks upon Him through His eyes will
perceive Him; otherwise he will be veiled. If you desire God and His
presence, desire Him and look upon Him.'' He also says: ``If you recite
one of the verses of Him Whom God will make manifest, it will be more
precious than if you transcribe the whole Bayán, for on that day that
one verse will save you, whereas the whole Bayán will not.'' End of
quotation.

\subsection{224}\label{section-1510}

\begin{otherlanguage}{arabic}

قل یا معشر البیان انصفوا انصفوا ثمّ انصفوا انصفوا و لا تکونوا من الّذین
ذکروا مظهر امر اللّه فی اللّیالی و الأیّام و لمّا اتی بالفضل و اشرق افق
الظّهور افتوا علیه بما ناح به سکّان الملکوت و الجبروت و الّذین طافوا حول
ارادة اللّه العلیم الحکیم

\end{otherlanguage}

Qul yā maʿshara al-Bayān, anṣifū anṣifū thumma anṣifū anṣifū, va lā
takūnū mina alladhīna dhakarū maẓhara amri Allāhi fī al-layālī va
al-ayyāmi va lammā atā bi-l-faḍli va ashraqa ufuqu al-ẓuhūri aftaw
ʿalayhi bi-mā nāḥa bihi sukkānu al-malakūti va al-jabarūti va alladhīna
ṭāfū ḥawla irādati Allāhi al-ʿAlīmi al-Ḥakīm.

Say: O people of the Bayán! Be fair, be fair; again, be fair, be fair!
Do not be among those who spoke of the Manifestation of the Cause of God
by night and day, but when He came with grace and the horizon of
Revelation shone forth, pronounced against Him a judgment that caused
the inhabitants of the heavenly kingdom and the realm of power, and
those who circled around the will of God, the All-Knowing, the All-Wise,
to lament.

\subsection{225}\label{section-1511}

\begin{otherlanguage}{arabic}

در این کلمهٔ علیا تفکّر نمائید میفرماید انّی مؤمن به و بدینه و بکتابه و
بأدلّائه و بمناهجه و بما یظهر من عنده فی کلّ ذلک مفتخراً بنسبتی الیه و
متعزّزاً بایمانی به و همچنین میفرماید ان یا کلّ شیء فی البیان فلتعرفنّ حدّ
انفسکم فانّ مثل نقطة البیان یؤمن بمن یظهره اللّه قبل کلّ شیء و انّنی انا
بذلک افتخرنّ علی من فی ملکوت السّموات و الأرض انتهی لعمر اللّه جمیع ذرّات
کائنات بنوحه و ندبه مشغولند از ظلم معرضین بیان آیا صاحبان ابصار و آذان
کجا رفتند نسأل اللّه تبارک و تعالی ان یحضرهم و ینصحهم بما ینفعهم و یمنعهم
عما یضرّهم انّه هو القویّ الغالب القدیر

\end{otherlanguage}

Dar īn kalimiy-i ʿulyā tafakkur namāyīd. Mī-farmāyad: innī muʾminun bihi
va bi-dīnihi va bi-kitābihi va bi-adillāʾihi va bi-manāhijihi va bi-mā
yaẓharu min ʿindihi, fī kulli dhālika muftakhiran bi-nisbatī ilayhi va
mutaʿazzizan bi-īmānī bihi. Va hamchunīn mī-farmāyad: an yā kulla
shayʾin fī al-Bayān, fa-la-taʿrifunna ḥadda anfusikum, fa-inna mithla
Nuqṭati al-Bayāni yuʾminu bi-man yuẓhiruhu Allāhu qabla kulli shayʾin,
va innanī anā bi-dhālika aftakhiranna ʿalā man fī malakūti al-samāwāti
va al-arḍ. Intahā. La-ʿamru Allāh, jamīʿ-i dharrāt-i kāʾināt bi nawḥih
va nudbih mashghūland az ẓulm-i muʿriḍīn-i Bayān. Āyā ṣāḥibān-i abṣār va
ādhān kujā raftand? Nasʾalu Allāha tabāraka va taʿālā an yuḥḍirahum va
yanṣaḥahum bi-mā yanfaʿuhum va yamnaʿuhum ʿammā yaḍurruhum. Innahu huwa
al-Qawiyyu al-Ghālibu al-Qadīr.

Reflect on these exalted words. He says: ``I believe in Him, in His
religion, in His Book, in His guides, in His ways, and in whatever
appears from Him, glorying in all this in My relationship to Him and
finding honor in My faith in Him.'' He likewise says: ``O all things in
the Bayán! You must surely know your own limits, for One such as the
Point of the Bayán believes in Him Whom God will make manifest before
all things, and in this I indeed glory above all who are in the kingdom
of the heavens and the earth.'' End of quotation. By the life of God!
All the atoms of existence are lamenting and mourning because of the
injustice of those among the people of the Bayán who have turned away.
Where have those with eyes and ears gone? We ask God, blessed and
exalted be He, to bring them forth and counsel them with what will
benefit them and keep them from what will harm them. He is indeed the
Strong, the All-Subduing, the All-Powerful.

\subsection{226}\label{section-1512}

\begin{otherlanguage}{arabic}

و همچنین میفرماید لا تحتجبنّ عن اللّه بعد ظهوره فانّ کلّ ما رفع البیان کخاتم
فی یدی و انّنی انا خاتم فی یدی من یظهره اللّه جلّ ذکره یقلّب کیف یشآء لما
یشآء بما یشآء انّه لهو المهیمن المتعال و همچنین میفرماید فانّه لو یجعل ما
علی الأرض نبیّاً لیکوننّ انبیآء عند اللّه و همچنین میفرماید و اذاً یوم ظهور
من یظهره اللّه کلّ من علی الأرض عنده سوآء فمن یجعله نبیّاً کان نبیّاً من اوّل
الّذی لا اوّل له الی آخر الّذی لا آخر له لأنّ ذلک ممّا قد جعله اللّه و من
یجعله ولیّاً فذلک ما کان ولیّاً فی کلّ العوالم فانّ ذلک ممّا قد جعله اللّه لأنّ
مشیّة اللّه لن یظهر الّا بمشیّته و ارادة اللّه لم یظهر الّا بارادته و انّه لهو
القاهر المقتدر المنیع انتهی

\end{otherlanguage}

Va hamchunīn mī-farmāyad: lā taḥtajibanna ʿani Allāhi baʿda ẓuhūrihi,
fa-inna kulla mā rafaʿa al-Bayānu ka-khātamin fī yadī, va innanī anā
khātamun fī yaday man yuẓhiruhu Allāhu jalla dhikruhu, yuqallibu kayfa
yashāʾu li-mā yashāʾu bi-mā yashāʾu. Innahu la-huwa al-Muhayminu
al-Mutaʿāl. Va hamchunīn mī-farmāyad: fa-innahu law yajʿalu mā ʿalā
al-arḍi nabiyyan la-yakūnanna anbiyāʾa ʿinda Allāh. Va hamchunīn
mī-farmāyad: va idhan yawma ẓuhūri man yuẓhiruhu Allāhu kullu man ʿalā
al-arḍi ʿindahu sawāʾun, fa-man yajʿaluhu nabiyyan kāna nabiyyan min
awwali alladhī lā awwala lahu ilā ākhiri alladhī lā ākhira lahu, li-anna
dhālika mimmā qad jaʿalahu Allāhu, va man yajʿaluhu waliyyan fa-dhālika
mā kāna waliyyan fī kulli al-ʿawālimi, fa-inna dhālika mimmā qad
jaʿalahu Allāhu, li-anna mashīʾata Allāhi lan yaẓhara illā
bi-mashīʾatihi va irādata Allāhi lam yaẓhar illā bi-irādatihi, va innahu
la-huwa al-Qāhiru al-Muqtadiru al-Manīʿ. Intahā.

He likewise says: ``Do not be veiled from God after His appearance, for
all that the Bayán has exalted is like a ring in My hand, and I Myself
am a ring in the hands of Him Whom God will make manifest, glorified be
His mention. He turns it however He wills, for whatever He wills, by
whatever He wills. He is indeed the All-Protecting, the Most Exalted.''
He also says: ``If He were to make all that is on earth a prophet, they
would assuredly be prophets in the sight of God.'' He likewise says:
``Thus, on the day of the appearance of Him Whom God will make manifest,
all who are on earth are equal before Him. Whomever He makes a prophet
is a prophet from the beginning that has no beginning to the end that
has no end, for this is what God has ordained. And whomever He makes a
saint is a saint in all the worlds, for this is what God has ordained.
For the will of God will appear only through His will, and the purpose
of God has appeared only through His purpose. He is indeed the
All-Subduing, the All-Powerful, the Inaccessible.'' End of quotation.

\subsection{227}\label{section-1513}

\begin{otherlanguage}{persian}

باری در هر مقامی ذکر فرموده‌اند آنچه را که سبب اقبال و ارتفاع و ارتقا و
هدایت خلق است ولکن معدودی غیر منصفین حجاب شدند و سدّی حائل گشته‌اند و عباد
را از توجّه بانوار وجه منع کرده‌اند نسأل اللّه ان یطردهم بسلطانه و یأخذهم
بأخذه انّه هو المقتدر العزیز الحکیم

\end{otherlanguage}

Bārī dar har maqāmī dhikr farmūdih-and ānchih rā kih sabab-i iqbāl va
irtifāʿ va irtiqā va hidāyat-i khalq ast, va-likin maʿdūdī ghayr-i
munṣifīn ḥijāb shudand va saddī ḥāʾil gashtih-and va ʿibād rā az
tavajjuh bi anvār-i vajh manʿ kardih-and. Nasʾalu Allāha an yaṭrudahum
bi-sulṭānihi va yaʾkhudhahum bi-akhdhihi. Innahu huwa al-Muqtadiru
al-ʿAzīzu al-Ḥakīm.

In short, in every instance He has set forth what will lead people to
turn toward God, to rise, to advance, and to find guidance. Yet a few
unfair people have become a veil and an intervening barrier, and have
prevented His servants from turning toward the light of His face. We ask
God to drive them away through His sovereignty and seize them in His
grasp. He is indeed the All-Powerful, the Almighty, the All-Wise.

\subsection{228}\label{section-1514}

\begin{otherlanguage}{arabic}

و همچنین میفرماید فانّ مثله جلّ ذکره کمثل الشّمس لو یقابلنه الی ما لا نهایة
مرایا کلّهنّ لیستعکسن من تجلّی الشّمس فى حدّهم و ان لم یقابلها من احد فیطلع
الشّمس و یغرب و الحجاب للمرایا و انّی ما قصّرت عن نصحی ذلک الخلق و تدبیری
لاقبالهم الی اللّه ربّهم و ایمانهم باللّه بارئهم و ان یؤمننّ به یوم ظهوره کلّ
ما على الأرض فاذاً یسرّ کینونتی حیث کلّ قد بلغوا الی ذروة وجودهم و وصلوا
الی طلعة محبوبهم و ادرکوا ما یمکن فی الامکان من تجلّی مقصودهم والّا یحزن
فؤادی و انّی قد ربّیت کلّ شیء لذلک فکیف یحتجب احد علی هذا قد دعوت اللّه و
لأدعونّه انّه قریب مجیب

\end{otherlanguage}

Va hamchunīn mī-farmāyad: fa-inna mathalahu jalla dhikruhu ka-mathali
al-shamsi, law yuqābilannahu ilā mā lā nihāyata marāyā kulluhunna
la-yastaʿkisna min tajallī al-shamsi fī ḥaddihim, va in lam yuqābilhā
min aḥadin fa-yaṭluʿu al-shamsu va yaghrubu va al-ḥijābu li-l-marāyā. Va
innī mā qaṣṣartu ʿan nuṣḥī dhālika al-khalqa va tadbīrī li-iqbālihim ilā
Allāhi Rabbihim va īmānihim bi-Allāhi Bāriʾihim. Va in yuʾminanna bihi
yawma ẓuhūrihi kullu mā ʿalā al-arḍi fa-idhan yusarru kaynūnatī ḥaythu
kullun qad balaghū ilā dhurwati wujūdihim va waṣalū ilā ṭalʿati
maḥbūbihim va adrakū mā yumkinu fī al-imkāni min tajallī maqṣūdihim, va
illā yaḥzanu fuʾādī. Va innī qad rabbaytu kulla shayʾin li-dhālika
fa-kayfa yaḥtajibu aḥadun? ʿAlā hādhā qad daʿawtu Allāha va
la-adʿuwannahu. Innahu qarībun mujīb.

He likewise says: ``His likeness, glorified be His mention, is that of
the sun. Were an endless number of mirrors to face it, they would all
reflect the sun's radiance according to their capacity. If none should
face it, the sun would still rise and set, and it is the mirrors that
would be veiled. I have not fallen short in My counsel to these people
or in My arrangements for them to turn toward God, their Lord, and
believe in God, their Maker. If all that is on earth should believe in
Him on the day of His appearance, My being would rejoice, for all would
have attained the summit of their existence, reached the presence of
their Beloved, and perceived as much of the revelation of their heart's
Desire as is possible in the contingent world. Otherwise My heart would
grieve. I have nurtured all things for this purpose; how, then, could
anyone remain veiled? For this I have prayed to God, and I will
assuredly continue to pray to Him. He is near and responsive.''

\subsection{229}\label{section-1515}

\begin{otherlanguage}{arabic}

و همچنین میفرماید بقدر اسم مؤمن هم در حقّ آن شجرهٔ لاشرقیّه و لاغربیّه راضی
نمیشوند چه اگر راضی شوند حزن بر او وارد نمیآورند انتهى یا اذن عالم آیا
شنیدی بچه عجز این بیانات از مشرق ارادهٔ مطلع اسماء ظاهر گشته میفرماید کل
را تربیت نمودم از برای عرفان این ظهور ولکن اهل بیان بقدر اسم مؤمن هم در
حقّ آن سدرهٔ مبارکهٔ لاشرقیّه راضی نمیشوند آه آه عمّا ورد علی نفسى لعمر اللّه
قد ورد علیّ من الّذی ربّیته فی اللّیالی و الأیّام ما ناح به روح القدس و اهل
خبآء عظمة اللّه مالک هذا الیوم البدیع

\end{otherlanguage}

Va hamchunīn mī-farmāyad: bi qadr-i ism-i muʾmin ham dar ḥaqq-i ān
shajariy-i lā-sharqiyyih va lā-gharbiyyih rāḍī namī-shavand, chih agar
rāḍī shavand ḥuzn bar ū vārid namī-āvarand. Intahā. Yā udhun-i ʿālam,
āyā shinīdī bi chih ʿajz īn bayānāt az mashriq-i irādiy-i Maṭlaʿ-i Asmā
ẓāhir gashtih? Mī-farmāyad: kull rā tarbiyat namūdam az barā-yi ʿirfān-i
īn ẓuhūr, va-likin ahl-i Bayān bi qadr-i ism-i muʾmin ham dar ḥaqq-i ān
Sidriy-i Mubārakiy-i lā-sharqiyyih rāḍī namī-shavand. Āh āh ʿammā warada
ʿalā nafsī! La-ʿamru Allāhi qad warada ʿalayya mina alladhī rabbaytuhu
fī al-layālī va al-ayyāmi mā nāḥa bihi Rūḥu al-Qudusi va ahlu khibāʾi
ʿaẓamati Allāhi Māliki hādhā al-yawmi al-badīʿ.

He likewise says: ``They are not willing to accord that Tree, neither of
the East nor of the West, even the name of believer, for if they were
willing, they would not bring sorrow upon Him.'' End of quotation. O ear
of the world! Have you heard with what helplessness these words have
issued from the dawning place of the will of the Dayspring of Names? He
says: ``I nurtured all for the recognition of this Revelation, yet the
people of the Bayán are not willing to accord that blessed Lote-Tree,
which is not of the East, even the name of believer.'' Alas, alas for
what has befallen Me! By the life of God! From the one whom I nurtured
by night and day, there has befallen Me what has caused the Holy Spirit
and the inhabitants of the pavilion of the grandeur of God, the Master
of this wondrous day, to lament.

\subsection{230}\label{section-1516}

\begin{otherlanguage}{persian}

و همچنین ردّاً لبعضی از معرضین میفرماید چه کسی عالم بظهور نیست غیر اللّه هر
وقت شود باید کل تصدیق بنقطهٔ حقیقت نمایند و شکر الهی بجای آورند انتهى
معرضین بمثابهٔ امّت یحیی تکلّم نمودند چه که آن نفوس هم بر حضرت روح اعتراض
کردند که شریعت یحیی تمام نشد تو از برای چه آمدی حال هم معرضین مع آنکه
ابداً با ما نبودند و اطّلاع از اصل امر نداشته و ندارند که از که بوده و چه
بوده گفته‌اند آنچه را که اشیاء کلّها بنوحه و ندبه مشغول لعمری انّ الأخرس لا
یقدر ان یقوم امام ملکوت البیان اتّقوا اللّه یا قوم ثمّ اقرؤوا ما نزّل بالحقّ
فی الباب الثّامن من الواحد السّادس من البیان و لا تکونوا من المعرضین و
همچنین امر فرموده که در هر نوزده روز یک دفعه در این باب نظر کنند لعلّ در
ظهور من یظهره اللّه محتجب نشوند بشئونی دون شأن آیات که اعظم حجج و براهین
بوده و هست انتهی

\end{otherlanguage}

Va hamchunīn raddan li-baʿḍī az muʿriḍīn mī-farmāyad: chih kasī ʿālim bi
ẓuhūr nīst ghayr Allāh; har vaqt shavad bāyad kull taṣdīq bi Nuqṭiy-i
Ḥaqīqat namāyand va shukr-i ilāhī bi jāy āvarand. Intahā. Muʿriḍīn bi
mathābiy-i ummat-i Yaḥyā takallum namūdand, chih kih ān nufūs ham bar
Ḥaḍrat-i Rūḥ iʿtirāḍ kardand kih sharīʿat-i Yaḥyā tamām nashud, tu az
barā-yi chih āmadī? Ḥāl ham muʿriḍīn maʿa ānkih abadan bā mā nabūdand va
iṭṭilāʿ az aṣl-i amr nadāshtih va nadārand kih az kih būdih va chih
būdih, guftih-and ānchih rā kih ashyā kulluhā bi nawḥih va nudbih
mashghūl. La-ʿamrī inna al-akhrasa lā yaqdiru an yaqūma imāma malakūti
al-bayān. Ittaqū Allāha yā qawmu thumma iqraʾū mā nuzzila bi-l-ḥaqqi fī
al-bābi al-thāmini mina al-wāḥidi al-sādisi mina al-Bayāni va lā takūnū
mina al-muʿriḍīn. Va hamchunīn amr farmūdih kih dar har nūzdah rūz yik
dafʿih dar īn bāb naẓar kunand, laʿall dar ẓuhūr-i Man Yuẓhiruhu Allāh
muḥtajib nashavand bi shuʾūnī dūn-i shaʾn-i āyāt kih aʿẓam-i ḥujaj va
barāhīn būdih va hast. Intahā.

He likewise says, in refutation of certain opponents: ``For no one knows
of the Revelation except God. Whenever it occurs, all must acknowledge
the Point of Truth and offer thanks to God.'' End of quotation. Those
who have turned away have spoken like the followers of John, for they
too objected to Him Who is the Spirit, saying: ``The dispensation of
John has not ended; why have you come?'' Now, too, those who have turned
away, although they were never with Us and neither had nor have any
knowledge of the origin of the Cause---of whom it came from or what it
was---have said things that have set all things lamenting and mourning.
By My life! The mute cannot stand before the kingdom of utterance. Fear
God, O people, then read what has been revealed in truth in the eighth
chapter of the sixth Váḥid of the Bayán, and do not be among those who
turn away. He has likewise commanded them to look at this chapter once
every nineteen days, so that, at the appearance of Him Whom God will
make manifest, they may not be veiled by anything other than the
revelation of verses, which has been and remains the greatest of proofs
and testimonies. End of quotation.

\subsection{231}\label{section-1517}

\begin{otherlanguage}{persian}

حضرت یحیی بن زکریّا فرموده آنچه را که مبشّر فرموده قائلاً توبوا لأنّه قد
اقترب ملکوت السّموات انّی اعمّدکم بمآء للتّوبة ولکنّ الّذی یأتی بعدی هو اقوی
منّی الّذی لست اهلاً ان احمل حذائه اینست که مبشّر میفرماید در مقام خضوع و
خشوع که جمیع بیان ورقی است از اوراق جنّت او و همچنین میفرماید من اوّل
عابدینم و افتخار میکنم بنسبتی الیه یا قوم مع‌ذلک اهل بیان عمل نمودند آنچه
را که ذی‌الجوشن و ابن انس و اصبحی از آن پناه بخدا برده و میبرند این مظلوم
در شب و روز بارتفاع امر الهی امام وجوه ادیان مشغول و آن نفوس باسبابی که
سبب ذلّت و علّت اذیّت بوده متمسّک

\end{otherlanguage}

Ḥaḍrat-i Yaḥyā ibn-i Zakariyyā farmūdih ānchih rā kih Mubashshir
farmūdih, qāʾilan: tūbū li-annahu qad iqtaraba malakūtu al-samāwāt. Innī
uʿammidukum bi-māʾin li-l-tawbati va lākinna alladhī yaʾtī baʿdī huwa
aqwā minnī, alladhī lastu ahlan an aḥmila ḥidhāʾahu. Īn ast kih
Mubashshir mī-farmāyad dar maqām-i khuḍūʿ va khushūʿ kih jamīʿ-i Bayān
varaqī ast az avrāq-i jannat-i ū. Va hamchunīn mī-farmāyad: man avval-i
ʿābidīnam va iftikhār mī-kunam bi-nisbatī ilayhi. Yā qawm, maʿa-dhālik
ahl-i Bayān ʿamal namūdand ānchih rā kih Dhī al-Jawshan va Ibn-i Anas va
Aṣbaḥī az ān panāh bi Khudā burdih va mī-barand. Īn Maẓlūm dar shab va
rūz bi irtifāʿ-i Amr-i Ilāhī imām-i vujūh-i adyān mashghūl va ān nufūs
bi asbābī kih sabab-i dhillat va ʿillat-i adhiyyat būdih mutamassik.

John, son of Zechariah, said what the Herald has said, declaring:
``Repent, for the kingdom of heaven has drawn near. I baptize you with
water for repentance, but He who comes after me is mightier than I,
whose sandals I am not worthy to carry.'' Thus the Herald says, in
humility and reverence: ``The whole Bayán is a leaf among the leaves of
His Paradise.'' He likewise says: ``I am the first of His worshippers,
and I glory in My relationship to Him.'' O people! Despite this, the
people of the Bayán have committed deeds from which Dhí'l-Jawshan, Ibn
Anas, and Aṣbaḥí have sought and still seek refuge in God. This Wronged
One is occupied night and day with exalting the Cause of God before the
eyes of the followers of all religions, while those people cling to
means that have brought abasement and caused suffering.

\subsection{232}\label{section-1518}

\begin{otherlanguage}{persian}

و همچنین میفرماید او را شناخته بآیات او و احتیاط در عرفان او نکرده که
بقدر همان در نار محتجب خواهید بود انتهی ای معرضین بیان در این کلمهٔ علیا
که از مطلع بیان نقطهٔ عرفان جاری شده تفکّر نمائید و در این حین آن کلمه را
بشنوید میفرماید در آن روز آن آفتاب حقیقت اهل بیان را خطاب مینماید و این
سورهٔ فرقان را تلاوت میفرماید قل یا ایّها الکافرون لا اعبد ما تعبدون و لا
انتم عابدون ما اعبد و لا انا عابد ما عبدتم و لا انتم عابدون ما اعبد لکم
دینکم و لی دین سبحان‌اللّه مع این بیانات واضحه و آثار لامعهٔ مشرقه کل
باوهام خود مشغولند و از حضرت مقصود غافل و محجوب یا معرضین از نوم غفلت
بیدار شوید و این کلمهٔ مبشّر را بشنوید میفرماید شجرهٔ اثبات باعراضش از او
از نفی محسوب و شجرهٔ نفی باقبالش باو از اثبات محسوب و همچنین میفرماید اگر
نفسی ادّعا نماید و اتیان بحجّت ننماید تعرّض منمائید و حزن وارد میاورید
انتهی باری این مظلوم در لیالی و ایّام به قل یا ایّها الکافرون ناطق که شاید
سبب تنبّه شود و خلق را بطراز انصاف مزیّن دارد

\end{otherlanguage}

Va hamchunīn mīfarmāyad: Ū rā shinākhte bi-āyāt-i ū va iḥtiyāṭ dar
ʿirfān-i ū nakarde ki bi-qadr-i hamān dar nār muḥtajib khvāhīd būd.
Intahā. Ay muʿriḍīn-i Bayān, dar īn kalimiy-i ʿulyā ki az maṭlaʿ-i
Bayān, Nuqṭiy-i ʿirfān, jārī shude tafakkur namāyīd va dar īn ḥīn ān
kalima rā bishnavīd. Mīfarmāyad: Dar ān rūz ān Āftāb-i Ḥaqīqat ahl-i
Bayān rā khiṭāb mīnamāyad va īn sūriy-i Furqān rā tilāvat mīfarmāyad:
Qul yā ayyuhā al-kāfirūn, lā aʿbudu mā taʿbudūn, wa-lā antum ʿābidūna mā
aʿbud, wa-lā anā ʿābidun mā ʿabadtum, wa-lā antum ʿābidūna mā aʿbud,
lakum dīnukum wa-liya dīn. Subḥān Allāh! Maʿ-i īn bayānāt-i vāḍiḥa va
āthār-i lāmiʿiy-i mushriqa kull bi-avhām-i khvud mashghūland va az
Ḥaḍrat-i Maqṣūd ghāfil va maḥjūb. Yā muʿriḍīn, az nawm-i ghaflat bīdār
shavīd va īn kalimiy-i Mubashshir rā bishnavīd. Mīfarmāyad: Shajariy-i
ithbāt bi-iʿrāḍash az ū az nafy maḥsūb va shajariy-i nafy bi-iqbālash
bi-ū az ithbāt maḥsūb. Va hamchunīn mīfarmāyad: Agar nafsī iddiʿā
namāyad va ityān bi-ḥujjat nanamāyad taʿarruḍ manamāyīd va ḥuzn vārid
mayāvarīd. Intahā. Bārī īn Maẓlūm dar layālī va ayyām bi-Qul yā ayyuhā
al-kāfirūn nāṭiq ki shāyad sabab-i tanabbuh shavad va khalq rā
bi-ṭirāz-i inṣāf muzayyan dārad.

He likewise says: ``Recognize Him by His verses, and do not hesitate in
recognizing Him, for to that extent you will remain veiled in the
fire.'' End of quotation. O you who turn away from the Bayán! Reflect on
this exalted utterance that has flowed from the Dayspring of the Bayán,
the Point of recognition, and now listen to that utterance. He says:
``On that day, that Sun of Truth will address the people of the Bayán
and recite this surah of the Qur'an: `Say: O unbelievers! I do not
worship what you worship, nor do you worship what I worship. Nor am I a
worshipper of what you have worshipped, nor are you worshippers of what
I worship. To you your religion, and to Me Mine.'\,'' Glorified be God!
Despite these clear utterances and these radiant, resplendent writings,
all are occupied with their own imaginings, heedless of and veiled from
the Desired One. O you who turn away! Awaken from the sleep of
heedlessness and listen to this utterance of the Herald. He says: ``The
tree of affirmation, by turning away from Him, is counted as negation,
and the tree of negation, by turning toward Him, is counted as
affirmation.'' He likewise says: ``If anyone advances a claim and
produces no proof, do not assail him or cause him sorrow.'' End of
quotation. In short, this Wronged One utters, night and day, ``Say: O
unbelievers!'' so that it may awaken them and adorn humanity with the
ornament of fairness.

\subsection{233}\label{section-1519}

\begin{otherlanguage}{arabic}

حال در این کلمات که عرف یأس از آن متضوّع تفکّر نمائید فی مناجاته مع
الأحزان الی اللّه ربّ العالمین قوله سبحانک اللّهمّ فاشهد بأنّی بذلک الکتاب قد
اخذت عهد ولایة من تظهرنّه عن کلّ شیء قبل عهد ولایتی و کفی بک و بمن آمن
بآیاتک علیّ شهدآء و انّک انت حسبی علیک توکّلت و انّک کنت علی کلّ شیء حسیبا

\end{otherlanguage}

Ḥāl dar īn kalimāt ki ʿarf-i yaʾs az ān mutaḍavviʿ tafakkur namāyīd. Fī
munājātihi maʿa al-aḥzān ilā Allāh Rabbi al-ʿālamīn, qawluhu: Subḥānaka
Allāhumma fa-shhad bi-annī bi-dhālika al-kitāb qad akhadhtu ʿahda
wilāyati man tuẓhirannahu ʿan kulli shayʾ qabla ʿahdi wilāyatī, wa-kafā
bika wa-bi-man āmana bi-āyātika ʿalayya shuhadāʾ, wa-innaka anta ḥasbī,
ʿalayka tawakkaltu, wa-innaka kunta ʿalā kulli shayʾin ḥasībā.

Now reflect on these words, from which the fragrance of despair wafts.
In His sorrowful communion with God, the Lord of the worlds, He says:
``Glory be to You, O God! Bear witness that through that Book I have
taken from all things a covenant of allegiance to Him Whom You will make
manifest before the covenant of allegiance to Me. You and those who have
believed in Your verses suffice as witnesses concerning Me. You are
indeed sufficient for Me. In You I have placed My trust, and You have
ever taken account of all things.''

\subsection{234}\label{section-1520}

\begin{otherlanguage}{persian}

و در مقام دیگر میفرماید ان یا شموس المرایا انتم الی شمس الحقیقة تنظرون
فانّ قیامکم بها لو انتم تتبصّرون کلّکم کحیتان بالمآء فی البحر تتحرّکون و
تحتجبون عن المآء و تسألون عمّا انتم به قائمون و همچنین میفرماید لأشکونّ
الیک ان یا مرآة جودی عن کلّ المرایا کلّ بألوانهم الیّ لینظرون انتهی این
خطاب از مصدر امر حضرت وهّاب بجناب آقا سیّد جواد مشهور به کربلائی نازل حقّ
شاهد و عالم گواه که آن سیّد با این مظلوم بوده و در ردّ معرضین هم شرحی
مبسوط نوشته و دو لوح هم که در آن از ظهور حقّ گواهی داده و از دونش اشارات
اعراض ظاهر و هویدا نزد جناب حیدر قبل علی فرستادیم و خطّ او نزد کل معلوم و
واضح است و مقصود از این عمل آنکه شاید منکرین بفرات اقرار فائز شوند و
معرضین بنور اقبال منوّر گردند حقّ شاهد که این مظلوم مقصودی جز القای کلمهٔ
الهی نداشته طوبی للمنصفین و ویل للمعرضین معرضین تدبیرها نموده‌اند و
بحیلها تمسّک جسته‌اند از جمله صورت این سیّد را گرفته‌اند و همچنین صور بعضی
را بعد هر کدام را در ورقی چسبانده‌اند و فوق اینها صورت میرزا یحیی را باری
بهر اسبابی از برای انکار حقّ تمسّک جسته‌اند قل حقّ عیان چون مهر رخشان آمده
حیف کاندر شهر کوران آمده و سیّد مذکور منکرین را نصیحت نموده و بافق اعلی
دعوت کرده ولکن در احجار ملساء تأثیری ننمود و در باره‌اش گفته‌اند آنچه را
که او از آن بحقّ جلّ جلاله پناه برده و حال عرائضی که بساحت اقدس ارسال
داشته موجود است طوبی للمنصفین

\end{otherlanguage}

Va dar maqām-i dīgar mīfarmāyad: An yā shumūsa al-marāyā, antum ilā
Shamsi al-Ḥaqīqa tanẓurūn, fa-inna qiyāmakum bihā law antum
tatabaṣṣarūn. Kullukum ka-ḥītānin bi-l-māʾi fī al-baḥri tataḥarrakūn,
wa-taḥtajibūna ʿani al-māʾi wa-tasʾalūna ʿammā antum bihi qāʾimūn. Va
hamchunīn mīfarmāyad: La-ashkuwanna ilayka an yā mirʾāta jūdī ʿan kulli
al-marāyā, kullun bi-alwānihim ilayya la-yanẓurūn. Intahā. Īn khiṭāb az
maṣdar-i amr-i Ḥaḍrat-i Vahhāb bi-janāb-i Āqā Siyyid Javād mashhūr
bi-Karbalāʾī nāzil. Ḥaqq shāhid va ʿālam gavāh ki ān Siyyid bā īn Maẓlūm
būde va dar radd-i muʿriḍīn ham sharḥī mabsūṭ nivishte va du lawḥ ham ki
dar ān az ẓuhūr-i Ḥaqq gavāhī dāde va az dūnash ishārāt-i iʿrāḍ ẓāhir va
huvaydā nazd-i janāb-i Ḥaydar qabl-i ʿAlī firistādīm va khaṭṭ-i ū nazd-i
kull maʿlūm va vāḍiḥ ast va maqṣūd az īn ʿamal ānki shāyad munkirīn
bi-Furāt-i iqrār fāʾiz shavand va muʿriḍīn bi-nūr-i iqbāl munavvar
gardand. Ḥaqq shāhid ki īn Maẓlūm maqṣūdī juz ilqāy-i kalimiy-i ilāhī
nadāshte. Ṭūbā li-l-munṣifīn wa-waylun li-l-muʿriḍīn. Muʿriḍīn tadbīrhā
namūde-and va bi-ḥīlhā tamassuk juste-and. Az jumla ṣūrat-i īn Siyyid rā
girifte-and va hamchunīn ṣuvar-i baʿḍī rā, baʿd har kudām rā dar varaqī
chasbānde-and va fawq-i īnhā ṣūrat-i Mīrzā Yaḥyā rā. Bārī bi-har asbābī
az barāy-i inkār-i Ḥaqq tamassuk juste-and. Qul: Ḥaqq ʿayān chūn mihr-i
rakhshān āmade, ḥayf k-andar shahr-i kūrān āmade. Va Siyyid-i madhkūr
munkirīn rā naṣīḥat namūde va bi-Ufuq-i Aʿlā daʿvat karde va-lākin dar
aḥjār-i malsāʾ taʾthīrī nanamūd va dar bāra-ash gufte-and ānchi rā ki ū
az ān bi-Ḥaqq jalla jalāluhu panāh burde va ḥāl ʿarāʾiḍī ki bi-Sāḥat-i
Aqdas irsāl dāshte mawjūd ast. Ṭūbā li-l-munṣifīn.

Elsewhere He says: ``O suns of the mirrors! You look toward the Sun of
Truth, for your existence depends upon it, if only you could perceive.
You are all like fish moving through the sea by virtue of the water, yet
you are veiled from the water and ask about that by which you subsist.''
He likewise says: ``I will surely complain to you, O mirror of My
bounty, about all the mirrors: each looks toward Me through its own
colors.'' End of quotation. This address was revealed from the Source of
the command of the All-Bountiful to Áqá Siyyid Javád, known as
Karbalá'í. God is witness, and the world bears testimony, that this
Siyyid was with this Wronged One and wrote a detailed account refuting
those who turned away. We also sent to Ḥaydar-ʿAlí two letters in which
he testified to the revelation of God, and in which the indications of
his rejection of all others are clear and evident. His handwriting is
known and recognizable to all. The purpose of this action was that those
who deny might perhaps attain the Euphrates of acknowledgment, and those
who turn away might be illumined by the light of acceptance. God is
witness that this Wronged One has had no purpose but to impart the word
of God. Blessed are the fair-minded, and woe to those who turn away!
Those who turn away have devised schemes and resorted to stratagems.
Among other things, they have obtained a portrait of this Siyyid, as
well as portraits of certain others, then pasted each onto a sheet and
placed a portrait of Mírzá Yaḥyá above them. In short, they have
resorted to every means of denying God. Say: ``The Truth has appeared,
manifest as the radiant sun; alas, it has come into a city of the
blind.'' The aforementioned Siyyid admonished those who denied and
invited them to the Most Exalted Horizon, but this had no effect on
those smooth stones. They have said things about him from which he has
sought refuge in God, exalted be His glory. The petitions he sent to the
Most Holy Court are now extant. Blessed are the fair-minded!

\subsection{235}\label{section-1521}

\begin{otherlanguage}{persian}

حال در شکایت نقطهٔ اولی از مرایا تفکّر نمائید شاید سبب انتباه گردد و عباد
از شمال اوهام و ظنون بیمین ایمان و ایقان توجّه نمایند و آگاه گردند بآنچه
که از او محجوبند مع آنکه از عالم نیستی بهستی مخصوص عرفان این امر اعظم
آمده‌اند و همچنین میفرماید و اجعل اللّهمّ تلک الشّجرة کلّها له لیظهرن منها
ثمرات ما قد خلق اللّه فیها لمن قد اراد اللّه ان یظهر به ما اراد فانّنی انا
وعزّتک ما اردت ان یکون على تلک الشّجرة من غصن و لا ورق و لا ثمر لن یسجد له
یوم ظهوره و لا یسبّحک به بما ینبغی لعلوّ علوّ ظهوره و سموّ سموّ بطونه و ان
شهدت یا الهی علیّ من غصن او ورق او ثمر لم یسجد له یوم ظهوره فاقطعه اللّهمّ
عن تلک الشّجرة فانّه لم یکن منّی و لا یرجع الیّ انتهی

\end{otherlanguage}

Ḥāl dar shikāyat-i Nuqṭiy-i Ūlā az marāyā tafakkur namāyīd shāyad
sabab-i intibāh gardad va ʿibād az shimāl-i avhām va ẓunūn bi-yamīn-i
īmān va īqān tavajjuh namāyand va āgāh gardand bi-ānchi ki az ū
maḥjūband, maʿ-i ānki az ʿālam-i nīstī bi-hastī makhṣūṣ-i ʿirfān-i īn
Amr-i Aʿẓam āmade-and. Va hamchunīn mīfarmāyad: Wa-jʿal Allāhumma tilka
al-shajarata kullahā lahu li-yaẓharanna minhā thamarātu mā qad khalaqa
Allāhu fīhā li-man qad arāda Allāhu an yuẓhira bihi mā arād. Fa-innanī
anā wa-ʿizzatika mā aradtu an yakūna ʿalā tilka al-shajarati min ghuṣnin
wa-lā waraqin wa-lā thamarin lan yasjuda lahu yawma ẓuhūrihi wa-lā
yusabbiḥaka bihi bi-mā yanbaghī li-ʿuluwwi ʿuluwwi ẓuhūrihi wa-sumūwwi
sumūwwi buṭūnihi. Wa-in shahidta yā ilāhī ʿalayya min ghuṣnin aw waraqin
aw thamarin lam yasjud lahu yawma ẓuhūrihi fa-qṭaʿhu Allāhumma ʿan tilka
al-shajarati fa-innahu lam yakun minnī wa-lā yarjiʿu ilayya. Intahā.

Now reflect on the Primal Point's complaint against the mirrors, that it
may perhaps awaken the servants and turn them from the left hand of
imaginings and suppositions toward the right hand of faith and
certitude, and make them aware of that from which they are veiled,
although they have come from the realm of nonexistence into existence
expressly to recognize this Most Great Cause. He likewise says: ``O God,
make that tree wholly His, so that from it may appear the fruits of what
God has created within it for Him through Whom God has willed to
manifest whatever He has willed. By Your might! I have not wished that
upon that tree there should be any branch, leaf, or fruit that will not
bow before Him on the day of His appearance and glorify You through Him
as befits the loftiest heights of His manifestation and the sublimest
heights of His concealment. And if, O my God, You witness upon Me any
branch, leaf, or fruit that does not bow before Him on the day of His
appearance, cut it off from that tree, O God, for it is not of Me and
does not return to Me.'' End of quotation.

\subsection{236}\label{section-1522}

\begin{otherlanguage}{persian}

یا اهل بیان لعمر اللّه این مظلوم خیالی نداشته الّا اظهار امری که بآن مأمور
بوده اگر بسمع فطرت توجّه نمائید از جمیع ارکان و اعضا و عروق حتّی از موهای
مظلوم اصغا مینمائید آنچه را که سبب توجّه و انجذاب ملأ اعلى و ناسوت انشاء
است

\end{otherlanguage}

Yā ahl-i Bayān, la-ʿamru Allāh, īn Maẓlūm khiyālī nadāshte illā iẓhār-i
amrī ki bi-ān maʾmūr būde. Agar bi-samʿ-i fiṭrat tavajjuh namāyīd az
jamīʿ-i arkān va aʿḍā va ʿurūq ḥattā az mūhāy-i Maẓlūm iṣghā mīnamāyīd
ānchi rā ki sabab-i tavajjuh va injidhāb-i Malaʾ-i Aʿlā va nāsūt-i inshā
ast.

O people of the Bayán! By the life of God, this Wronged One has had no
thought but to make known the Cause with which He was entrusted. If you
listen with the ear of your innate nature, you will hear from every
part, limb, and vein, even from the hairs of this Wronged One, that
which draws the attention and attracts the hearts of the Concourse on
high and the human world of creation.

\subsection{237}\label{section-1523}

\begin{otherlanguage}{persian}

یا هادی حمیّهٔ جاهلیّهٔ قبل عباد بیچاره را از صراط مستقیم منع نمود در حزب
شیعه تفکّر نما هزار و دویست سنه یا قائم گفتند و بالأخره جمیع بر شهادتش
فتوی دادند و شهیدش نمودند مع آنکه بحقّ جلّ جلاله و حضرت خاتم و اوصیا قائل
و مقرّ و معترف بودند حال قدری فکر لازم که شاید معلوم شود سببی که مابین حقّ
و خلق حائل شد چه بوده و علّت اعتراض و انکار چه عملی شده

\end{otherlanguage}

Yā Hādī, ḥamiyyiy-i jāhiliyyiy-i qabl ʿibād-i bīchāra rā az ṣirāṭ-i
mustaqīm manʿ namūd. Dar ḥizb-i Shīʿa tafakkur namā. Hizār va duvīst
sana Yā Qāʾim guftand va bi-l-ākhara jamīʿ bar shahādatash fatvā dādand
va shahīdash namūdand, maʿ-i ānki bi-Ḥaqq jalla jalāluhu va Ḥaḍrat-i
Khātam va awṣiyā qāʾil va muqirr va muʿtarif būdand. Ḥāl qadrī fikr
lāzim ki shāyad maʿlūm shavad sababī ki mābayn-i Ḥaqq va khalq ḥāʾil
shud chi būde va ʿillat-i iʿtirāḍ va inkār chi ʿamalī shude.

O Hádí! The fanatical zeal of the former age of ignorance barred the
helpless servants from the straight path. Reflect on the Shíʿih
community. For twelve hundred years they cried, ``O Qá'im!'' and in the
end they all issued a verdict for His death and put Him to death,
although they professed belief in, acknowledged, and confessed God,
exalted be His glory, the Seal of the Prophets, and His successors. Now
some reflection is needed, so that it may perhaps become clear what
cause came between God and His creatures, and what action occasioned
their objection and denial.

\subsection{238}\label{section-1524}

\begin{otherlanguage}{arabic}

یا هادی نالهٔ منابر را اصغا نمودیم که علمای عصر ظهور بقول جمیع بر آن
ارتقا جستند و حقّ را سب نموده و بر آن جوهر وجود و اصحابش وارد آوردند آنچه
که چشم عالم ندید و گوش عالم نشنید حال تو باسم وصایت و مرآتیّت مردم را
دعوت نموده و مینمائی مع آنکه اطّلاع از امر نداشته و نداری چه که با ما
نبودی

\end{otherlanguage}

Yā Hādī, nāliy-i manābir rā iṣghā namūdīm ki ʿulamāy-i ʿaṣr-i ẓuhūr
bi-qawl-i jamīʿ bar ān irtiqā justand va Ḥaqq rā sabb namūde va bar ān
Jawhar-i Vujūd va aṣḥābash vārid āvardand ānchi ki chashm-i ʿālam nadīd
va gūsh-i ʿālam nashnīd. Ḥāl tu bi-ism-i viṣāyat va mirʾātiyyat mardum
rā daʿvat namūde va mīnamāʾī maʿ-i ānki iṭṭilāʿ az amr nadāshte va
nadārī chi ki bā mā nabūdī.

O Hádí! We heard the lamentation of the pulpits upon which, as all
attest, the religious scholars at the time of the Revelation ascended,
reviling God and inflicting upon that Essence of Being and His
companions what the eye of the world had never seen and the ear of the
world had never heard. Now you have summoned, and continue to summon,
people in the name of successorship and mirrorhood, although you had and
have no knowledge of the matter, for you were not with Us.

\subsection{239}\label{section-1525}

\begin{otherlanguage}{persian}

جمیع این حزب میدانند که سیّد محمّد یکی از خدّام بوده در ایّامی که حسب
الخواهش دولت علیّهٔ عثمانی بآن شطر توجّه نمودیم همراه بوده و بعد از او ظاهر
شد آنچه که لعمر اللّه قلم اعلی گریست و لوح نوحه نمود لذا طردش نمودیم و به
میرزا یحیی پیوست و عمل نمود آنچه را که هیچ ظالمی ننمود ترکناه و قلنا
اخرج یا غافل بعد از ظهور این کلمه به مولوی‌خانه رفته با آنها بوده الی ان
جآء امر الخروج

\end{otherlanguage}

Jamīʿ-i īn ḥizb mīdānand ki Siyyid Muḥammad yakī az khuddām būde, dar
ayyāmī ki ḥasb al-khvāhish-i Dawlat-i ʿAliyyiy-i ʿUthmānī bi-ān shaṭr
tavajjuh namūdīm hamrāh būde va baʿd az ū ẓāhir shud ānchi ki la-ʿamru
Allāh Qalam-i Aʿlā girīst va Lawḥ nawḥa namūd. Lidhā ṭardash namūdīm va
bi-Mīrzā Yaḥyā payvast va ʿamal namūd ānchi rā ki hīch ẓālimī nanamūd.
Taraknāhu wa-qulnā ukhruj yā ghāfil. Baʿd az ẓuhūr-i īn kalima
bi-Mawlavī-khāna rafte bā ānhā būde ilā an jāʾa amru al-khurūj.

Everyone in this community knows that Siyyid Muḥammad was one of the
attendants. He accompanied Us when, at the request of the exalted
Ottoman government, We set out for that region. Afterward, he displayed
conduct that, by the life of God, made the Most Exalted Pen weep and the
Tablet lament. We therefore expelled him, and he joined Mírzá Yaḥyá and
did what no oppressor had done. We left him and said, ``Depart, O
heedless one!'' After these words were uttered, he went to the Mawlaví
lodge and remained with them until the order to depart arrived.

\subsection{240}\label{section-1526}

\begin{otherlanguage}{persian}

یا هادی سبب اوهام جدیده مشو مرّةً اخری بترتیب حزبی مثل حزب شیعه راضی مباش
تفکّر نما چه مقدار خون ریخته شد از جمله خود تو که مدّعی علم بودی و همچنین
علمای شیعه طرّاً در سنهٔ اولی و اخری حقّ را لعن نمودید و فتوی بر سفک دم
اطهرش دادید اتّق اللّه یا هادی راضی مشو عباد مجدّد باوهام قبل مبتلا شوند
اتّق اللّه و لا تکن من الظّالمین این ایّام اصغا شد در جمع بیان و محو آن همّت
نموده‌ئی این مظلوم از تو میطلبد این اراده را لوجه اللّه ترک نمائی ادراک و
عقل تو از سیّد عالم برتر و بالاتر نبوده و نیست حقّ شاهد و گواهست که این
مظلوم بیان را تلاوت ننموده و مطالبش را ندیده اینقدر معلوم و واضح و
مبرهنست که کتاب بیان را اسّ کتب خود قرار فرموده از خدا بترس در اموری که
شأن تو نیست داخل مشو هزار و دویست سنه شیعهٔ بیچاره را امثال تو ببئر اوهام
و ظنون مبتلا نمودند بالأخره در یوم جزا ظاهر شد آنچه که ظالمهای قبل بحقّ
پناه بردند

\end{otherlanguage}

Yā Hādī, sabab-i avhām-i jadīda mashav. Marratan ukhrā bi-tartīb-i ḥizbī
mithl-i ḥizb-i Shīʿa rāḍī mabāsh. Tafakkur namā chi miqdār khūn rīkhte
shud. Az jumla khvud-i tu ki muddaʿiy-i ʿilm būdī va hamchunīn ʿulamāy-i
Shīʿa ṭurran dar saniy-i ūlā va ukhrā Ḥaqq rā laʿn namūdīd va fatvā bar
safk-i dam-i aṭharash dādīd. Ittaqi Allāh yā Hādī. Rāḍī mashav ʿibād
mujaddad bi-avhām-i qabl mubtalā shavand. Ittaqi Allāh wa-lā takun mina
al-ẓālimīn. Īn ayyām iṣghā shud dar jamʿ-i Bayān va maḥv-i ān himmat
namūde-ī. Īn Maẓlūm az tu mīṭalabad īn irāda rā li-vajh Allāh tark
namāʾī. Idrāk va ʿaql-i tu az Siyyid-i ʿālam bartar va bālātar nabūde va
nīst. Ḥaqq shāhid va gavāh ast ki īn Maẓlūm Bayān rā tilāvat nanamūde va
maṭālibash rā nadīde. Īnqadr maʿlūm va vāḍiḥ va mubarhan ast ki Kitāb-i
Bayān rā uss-i kutub-i khvud qarār farmūde. Az Khudā bitars. Dar umūrī
ki shaʾn-i tu nīst dākhil mashav. Hizār va duvīst sana Shīʿiy-i bīchāra
rā amthāl-i tu bi-biʾr-i avhām va ẓunūn mubtalā namūdand. Bi-l-ākhara
dar Yawm-i Jazā ẓāhir shud ānchi ki ẓālimhāy-i qabl bi-Ḥaqq panāh
burdand.

O Hádí! Do not become the cause of new imaginings. Do not consent to the
formation, once again, of a community like the Shíʿih community. Reflect
on how much blood was shed. You yourself, who claimed knowledge, and
likewise all the Shíʿih scholars, cursed God in the first year and the
last and issued a verdict for the shedding of His most pure blood. Fear
God, O Hádí! Do not consent to the servants being afflicted once again
with the imaginings of the past. Fear God, and do not be among the
oppressors. In these days We have heard that you have set yourself to
collecting the Bayán and obliterating it. This Wronged One asks you to
abandon this intention for the sake of God. Your understanding and
intellect have never been, and are not, superior to those of the Lord of
the world. God bears witness and testimony that this Wronged One has not
read the Bayán or examined its contents. This much is known, clear, and
proven: He made the Book of the Bayán the foundation of His books. Fear
God. Do not involve yourself in matters that are not your concern. For
twelve hundred years, people like you consigned the helpless Shíʿih
community to the pit of imaginings and suppositions. In the end, on the
Day of Recompense, there appeared that from which the oppressors of the
past sought refuge in God.

\subsection{241}\label{section-1527}

\begin{otherlanguage}{persian}

حال صریخ حضرت نقطه را از بیانش ادراک نما عرض میکند الهی اگر از این سدره
که آن وجود مبارک باشد ثمری یا ورقی یا غصنی ظاهر شود که بتو ایمان نیاورد
او را این حین قطع نما و همچنین میفرماید اگر نفسی کلمه‌ئی بگوید من غیر
برهان ردّش ننمائید و حال با صد جلد کتاب ردّش کرده‌ئی و مسروری

\end{otherlanguage}

Ḥāl ṣarīkh-i Ḥaḍrat-i Nuqṭa rā az bayānash idrāk namā. ʿArḍ mīkunad:
Ilāhī agar az īn Sidra ki ān Vujūd-i Mubārak bāshad thamarī yā varaqī yā
ghuṣnī ẓāhir shavad ki bi-tu īmān nayāvarad ū rā īn ḥīn qaṭʿ namā. Va
hamchunīn mīfarmāyad: Agar nafsī kalima-ī bigūyad min ghayr-i burhān
raddash nanamāʾīd. Va ḥāl bā ṣad jild kitāb raddash karde-ī va masrūrī.

Now perceive the cry of the Point in His words. He entreats: ``My God!
If from this Tree''---which is His blessed Being---``there should appear
any fruit, leaf, or branch that does not believe in You, cut it off this
very moment.'' He likewise says: ``If anyone speaks a word without
proof, do not reject him.'' Yet now you have rejected Him despite a
hundred volumes of scripture, and you rejoice.

\subsection{242}\label{section-1528}

\begin{otherlanguage}{persian}

مجدّد میگویم و التماس میکنم ببصر حدید در آنچه نازل شده ملاحظه نما نفحات
بیان این ظهور غیر بیانات قبل است این مظلوم لازال مبتلا بوده و مقرّ امنی
که در کتب حضرت اعلی و یا غیر نظر نماید نداشته بعد از ورود در عراق بامر
حضرت پادشاه ایران ایّده اللّه دو ماه او ازید فاصله میرزا یحیی وارد شد ذکر
نمودیم ما را حسب الأمر باینجا فرستاده‌اند تو خوبست در ایران باشی و میرزا
موسای اخوی را میفرستیم بطرف دیگر چه که اسم شما در دیوان نیامده میتوانید
بخدمتی قیام نمائید بعد این مظلوم دو سنه منقطعاً عن العالم هجرت کرد از
بغداد بعد از رجوع مشاهده شد توقّف کرده و سفرش در عهدهٔ تأخیر مانده این
مظلوم بسیار محزون شد حقّ شاهد و گواه که در جمیع احیان بنشر امر مشغول
بوده‌ایم سلاسل و اغلال منع ننمود کند و حبس از اظهار بازنداشت و در آن ارض
بر منع از فساد و اعمال غیر مرضیّهٔ غیر طیّبه قیام نمودیم و لیلاً و نهاراً
بجمیع اطراف الواح ارسال شد و مقصودی جز تهذیب نفوس و اعلای کلمهٔ مبارکه
نبوده

\end{otherlanguage}

Mujaddad mīgūyam va iltimās mīkunam bi-baṣar-i ḥadīd dar ānchih nāzil
shudih mulāḥiẓih namā. Nafaḥāt-i bayān-i īn ẓuhūr ghayr-i bayānāt-i qabl
ast. Īn Maẓlūm lāzāl mubtalā būdih va maqarr-i amnī kih dar kutub-i
Ḥaḍrat-i Aʿlā va yā ghayr naẓar namāyad nadāshtih. Baʿd az vurūd dar
ʿIrāq bi-amr-i Ḥaḍrat-i Pādishāh-i Īrān, ayyadahu Allāh, du māh aw azyad
fāṣilih Mīrzā Yaḥyā vārid shud. Dhikr namūdīm mā rā ḥasb al-amr bi-īnjā
firistādih-and; tu khūb ast dar Īrān bāshī va Mīrzā Mūsā-yi akhavī rā
mīfiristīm bi-ṭaraf-i dīgar, chih kih ism-i shumā dar dīvān nayāmadih,
mītavānīd bi-khidmatī qiyām namāyīd. Baʿd īn Maẓlūm du sanih munqaṭiʿan
ʿan al-ʿālam hijrat kard az Baghdād. Baʿd az rujūʿ mushāhidih shud
tavaqquf kardih va safarash dar ʿuhdih-yi taʾkhīr māndih. Īn Maẓlūm
bisyār maḥzūn shud. Ḥaqq shāhid va guvāh kih dar jamīʿ-i aḥyān
bi-nashr-i amr mashghūl būdih-īm. Salāsil va aghlāl manʿ nanamūd; kand
va ḥabs az iẓhār bāz nadāsht. Va dar ān arḍ bar manʿ az fasād va aʿmāl-i
ghayr-i marḍiyyih-yi ghayr-i ṭayyibih qiyām namūdīm va laylan va nahāran
bi-jamīʿ-i aṭrāf alvāḥ irsāl shud va maqṣūdī juz tahdhīb-i nufūs va
iʿlā-yi kalimih-yi mubārakih nabūdih.

Once again I say, and implore you: examine with keen sight what has been
revealed. The fragrances of the utterance of this Revelation are unlike
those of earlier utterances. This Wronged One has always been afflicted
and has had no secure place in which to examine the books of His
Holiness the Most Exalted or any others. Two months or more after Our
arrival in Iraq by order of His Majesty the Shah of Iran---may God
sustain him---Mírzá Yaḥyá arrived. We said: ``We have been sent here by
official command. It would be best for you to remain in Iran, and We
will send Our brother Mírzá Músá elsewhere, since your names have not
been entered in the government register and you can arise to render some
service.'' After this, the Wronged One departed from Baghdad for two
years, detached from the world. Upon returning, We found that he had
remained and that his journey had been postponed. This Wronged One was
deeply saddened. God bears witness and testimony that at all times We
have been occupied with spreading the Cause. Chains and shackles did not
prevent Us; fetters and imprisonment did not keep Us from proclaiming
it. In that land We arose to prevent wrongdoing and objectionable,
unwholesome deeds, and day and night Tablets were sent in every
direction. Our purpose has been nothing other than the refinement of
souls and the exaltation of the blessed Word.

\subsection{243}\label{section-1529}

\begin{otherlanguage}{persian}

مخصوص چند نفر معیّن نمودیم بر جمع آثار نقطه و بعد از جمع میرزا یحیى و
میرزا وهّاب خراسانی که به میرزا جواد معروف بود این دو را در محلّی جمع
نمودیم و دو دوره کتب حضرت نقطه را حسب الأمر نوشته و تمام نموده‌اند لعمر
اللّه این مظلوم از کثرت مراوده با ناس کتب را ندیده و از آثار نقطه ببصر
ظاهر مشاهده ننموده و این آثار نزد این دو بوده که هجرت واقع شد و قرار شد
میرزا یحیی این نوشتجات را برداشته بشطر ایران توجّه نماید و در آن اراضی
انتشار دهد و این مظلوم حسب الاستدعای وزرای دولت علیّه بآن شطر توجّه نمود
بعد از ورود در موصل مشاهده شد میرزا یحیی پیش از حرکت مظلوم رفته و منتظر
است باری کتب و آثار در بغداد ماند و او خود بشطر علیّه توجّه نمود و جزء این
عباد شد حال حقّ شاهد است بر این مظلوم چه گذشت چه که بعد از زحمتهای زیاد
آثار را گذاشت و خود بمهاجرین پیوست مدّتها این مظلوم باحزان نامتناهیه
مبتلا تا آنکه بتدبیری که غیر حقّ کسى آگاه نیست آثار را بمقام دیگر و ارض
دیگر فرستادیم چه که در عراق عرب باید اوراق را در هر شهر ملاحظه نمود والّا
از هم میریخت و ضایع میشد ولکنّ اللّه حفظها و ارسلها الی مقام قدّره اللّه من
قبل انّه هو الحافظ المعین

\end{otherlanguage}

Makhṣūṣ chand nafar muʿayyan namūdīm bar jamʿ-i āthār-i Nuqṭih va baʿd
az jamʿ Mīrzā Yaḥyā va Mīrzā Vahhāb-i Khurāsānī kih bi-Mīrzā Javād
maʿrūf būd, īn du rā dar maḥallī jamʿ namūdīm va du dawrih kutub-i
Ḥaḍrat-i Nuqṭih rā ḥasb al-amr nivishtih va tamām namūdih-and. La-ʿamru
Allāh, īn Maẓlūm az kathrat-i murāvadih bā nās kutub rā nadīdih va az
āthār-i Nuqṭih bi-baṣar-i ẓāhir mushāhidih nanamūdih va īn āthār nazd-i
īn du būdih kih hijrat vāqiʿ shud va qarār shud Mīrzā Yaḥyā īn
nivishtijāt rā bardāshtih bi-shaṭr-i Īrān tavajjuh namāyad va dar ān
arāḍī intishār dihad va īn Maẓlūm ḥasb al-istidʿā-yi vuzarā-yi Dawlat-i
ʿAliyyih bi-ān shaṭr tavajjuh namūd. Baʿd az vurūd dar Mawṣil mushāhidih
shud Mīrzā Yaḥyā pīsh az ḥarakat-i Maẓlūm raftih va muntaẓir ast. Bārī
kutub va āthār dar Baghdād mānd va ū khud bi-shaṭr-i ʿAliyyih tavajjuh
namūd va juzʾ-i īn ʿibād shud. Ḥāl Ḥaqq shāhid ast bar īn Maẓlūm chih
gudhasht, chih kih baʿd az zaḥmathā-yi ziyād āthār rā gudhāsht va khud
bi-muhājirīn payvast. Muddat-hā īn Maẓlūm bi-aḥzān-i nāmutanāhiyih
mubtalā tā ānkih bi-tadbīrī kih ghayr-i Ḥaqq kasī āgāh nīst āthār rā
bi-maqām-i dīgar va arḍ-i dīgar firistādīm, chih kih dar ʿIrāq-i ʿArab
bāyad awrāq rā dar har shahr mulāḥiẓih namūd va illā az ham mīrīkht va
ḍāyiʿ mīshud. Wa-lākinna Allāha ḥafiẓahā wa-arsalahā ilā maqāmin
qaddarahu Allāhu min qablu, innahu huwa al-Ḥāfiẓu al-Muʿīn.

We specifically appointed several people to collect the writings of the
Point. After they had been collected, We brought Mírzá Yaḥyá and Mírzá
Vahháb of Khurásán, who was known as Mírzá Javád, together in one place.
At Our direction, these two transcribed and completed two sets of the
books of His Holiness the Point. By the life of God! Because of His
extensive dealings with people, this Wronged One did not see the books
or examine the writings of the Point with His outward eyes. These
writings were in the keeping of those two when the departure took place.
It was arranged that Mírzá Yaḥyá should take these manuscripts, set out
for Iran, and disseminate them in those lands, while this Wronged One,
at the request of the ministers of the Sublime State, set out in that
direction. After arriving in Mosul, We found that Mírzá Yaḥyá had gone
there before the departure of the Wronged One and was waiting. In short,
the books and writings remained in Baghdad, while he himself set out
toward the Sublime Porte and joined these servants. God now bears
witness to what this Wronged One endured, for after so much trouble he
left the writings behind and himself joined those who were departing.
For a long time this Wronged One was afflicted with boundless sorrows,
until, by a measure of which no one but God is aware, We sent the
writings to another place and another land. For in Arab Iraq the papers
had to be inspected in every town; otherwise they would fall apart and
be lost. But God preserved them and sent them to a place God had
previously ordained. He is indeed the Preserver, the Helper.

\subsection{244}\label{section-1530}

\begin{otherlanguage}{persian}

هر جا این مظلوم رفت میرزا یحیی از عقب آمد خود تو گواهی و میدانی که آنچه
ذکر شد صدقست ولکن در سرّ سیّد اصفهانی او را اغوا نمود و عمل نمودند آنچه را
که سبب فزع اکبر شد ایکاش از مأمورین دولت سؤال مینمودید عمل میرزا یحیی را
در آن ارض از همه گذشته اقسمک باللّه الفرد الواحد المقتدر القدیر که در
نوشتجاتی که باسم او نزد نقطهٔ اولی رفته ملاحظه نما تا آثار حقّ را بمثابهٔ
آفتاب ممتاز مشاهده نمائی و همچنین از کلمات نقطهٔ بیان روح ما سواه فداه
ظاهر شده آنچه که هیچ سترى آن را منع ننماید و سبحات جلال و حجبات اهل ضلال
آن را از ظهور بازندارد قد خرقت الأحجاب من اصبع ارادة ربّک القویّ الغالب
القدیر بلی مفترین و مغلّین را چاره نبوده و نیست چندی قبل مذکور شد کتاب
ایقان و بعضی از الواح را نسبت بغیر داده‌ئی لعمر اللّه هذا ظلم عظیم غیر از
ادراک آن عاجز است تا چه رسد بتنزیل آن

\end{otherlanguage}

Har jā īn Maẓlūm raft Mīrzā Yaḥyā az ʿaqab āmad. Khud tu guvāhī va
mīdānī kih ānchih dhikr shud ṣidq ast, valīkin dar sirr Sayyid-i
Iṣfahānī ū rā ighvā namūd va ʿamal namūdand ānchih rā kih sabab-i
fazaʿ-i akbar shud. Ay kāsh az maʾmūrīn-i dawlat suʾāl mīnamūdīd ʿamal-i
Mīrzā Yaḥyā rā dar ān arḍ. Az hamih gudhashtih, uqsimuka bi-Allāhi
al-Fardi al-Wāḥidi al-Muqtadiri al-Qadīr, kih dar nivishtijātī kih
bi-ism-i ū nazd-i Nuqṭih-yi Ūlā raftih mulāḥiẓih namā tā āthār-i Ḥaqq rā
bi-mathābih-yi āftāb mumtāz mushāhidih namāʾī. Va hamchunīn az kalimāt-i
Nuqṭih-yi Bayān, rūḥu mā siwāhu fidāhu, ẓāhir shudih ānchih kih hīch
sitrī ān rā manʿ nanamāyad va subuḥāt-i jalāl va ḥujubāt-i ahl-i ḍalāl
ān rā az ẓuhūr bāz nadārad. Qad khuriqat al-aḥjābu min iṣbaʿi irādati
rabbika al-Qawiyyi al-Ghālibi al-Qadīr. Balī muftarīn va mughillīn rā
chārih nabūdih va nīst. Chandī qabl madhkūr shud Kitāb-i Īqān va baʿḍī
az alvāḥ rā nisbat bi-ghayr dādih-ī. La-ʿamru Allāh, hādhā ẓulmun ʿaẓīm.
Ghayr az idrāk-i ān ʿājiz ast tā chih risad bi-tanzīl-i ān.

Wherever this Wronged One went, Mírzá Yaḥyá followed. You yourself are a
witness and know that what has been stated is true. In secret, however,
the Sayyid of Isfahan led him astray, and they did what caused the
greatest consternation. If only you had asked the government officials
about Mírzá Yaḥyá's conduct in that land! Apart from all this, I adjure
you by God, the Unique, the One, the All-Powerful, the Omnipotent, to
examine the writings sent in his name to the Primal Point, so that you
may see the writings of God as distinctly as the sun. Likewise, from the
words of the Point of the Bayán---may the souls of all besides Him be
sacrificed for Him---there has appeared that which no covering can
conceal, and which neither the veils of glory nor the coverings of the
people of error can prevent from appearing. The veils have been rent by
the finger of the will of your Lord, the Strong, the All-Conquering, the
Omnipotent. Yet the fabricators and the malicious have had, and still
have, no other recourse. Some time ago it was reported that you had
attributed the Book of Certitude and some of the Tablets to someone
else. By the life of God, this is a grievous injustice! Another is
incapable of comprehending it, let alone revealing it.

\subsection{245}\label{section-1531}

\begin{otherlanguage}{persian}

حسن مازندرانی حامل هفتاد لوح بوده و چون فوت شد آن الواح را بصاحبانش
ندادند و بیکى از اختهای این مظلوم که من غیر جهت اعراض نموده سپردند اللّه
یعلم ما ورد علی الواحه و آن اخت ابداً با ما نبوده قسم بآفتاب حقیقت بعد از
ظهور این امور میرزا یحیی را ندیده و از امر مطّلع نبوده چه که آن ایّام
موافق نبوده‌اند ایشان در محلّه‌ئی و این مظلوم در محلّهٔ دیگر ساکن ولکن محض
عنایت و محبّت و شفقت چند یوم قبل از حرکت نزد او و والده‌اش رفته که شاید از
کوثر ایمان بیاشامند و بآنچه الیوم سبب تقرّب الی اللّه است فائز گردند حقّ
میداند و شاهد و خود او گواه که غیر این بهیچوجه خیالی نبوده تا آنکه الحمد
للّه از فضل الهی فائز شد و بطراز محبّت مزیّن گشت ولکن بعد از اسیری و هجرت
ما از عراق به آستانه دیگر از او خبری نرسید و بعد از تفریق در ارض طاء با
جناب اخوی میرزا رضا قلی ملاقات نشد و خبری از او مخصوصاً نرسید در اوّل ایّام
کل در یک بیت ساکن بودیم و بعد آن بیت در خریمه بقیمت نازلی رفت و
فرمانفرما و حسام‌السّلطنه این دو برادر خریدند و قسمت نمودند بعد از این
حادثه مابین ما و اخوی تفریق حاصل ایشان بدر مسجد شاه و ما در دروازهٔ
شمیران ساکن ولکن از اخت از بعد من غیر جهت آثار عناد ظاهر این مظلوم
بهیچوجه سخنی نگفته الّا آنکه بنت اخوی مرحوم میرزا محمّد حسن علیه بهآء اللّه
و سلامه و رحمته که مخطوبهٔ غصن اعظم بوده او را اخت این مظلوم از نور بخانهٔ
خود برده و بمقرّ دیگر فرستاد جمعی از اصحاب و دوستان از اطراف شکایت نمودند
چه که این امر بسیار عظیم بود موافق رأی هیچیک از اولیای حقّ واقع نشد عجب
در اینکه اخت او را بمحلّ خود برده و ترتیبات داده بمقام دیگر فرستاده مع‌ذلک
این مظلوم ساکت و صامت بوده و هست مگر آنکه مخصوص تسکین احباب یک کلمه
اظهار رفت و حقّ شاهد و گواه است که آنچه گفته شد حقیقت بوده و براستی گفته
و احدی از اولیای این اطراف و آن اراضی گمان نمیکردند که از اخت چنین امری
که خلاف حمیّت و محبّت و دوستی است واقع شود بعد از ظهور این امر سبیل را
مقطوع دیدند و عمل نمودند آنچه را که کل میدانید و میدانند دیگر معلومست که
چه مرتبه حزن از این عمل بر مظلوم وارد شد و بعد به میرزا یحیی پیوست و حال
مختلف شنیده میشود معلوم نیست چه میگوید و چه میکند نسأل اللّه تبارک و
تعالی ان یرجعها الیه و یؤیّدها علی الانابة لدی باب فضله انّه هو العزیز
التّوّاب و هو المقتدر الغفّار

\end{otherlanguage}

Ḥasan-i Māzandarānī ḥāmil-i haftād lawḥ būdih va chūn fawt shud ān alvāḥ
rā bi-ṣāḥibānash nadādand va bi-yikī az ukht-hā-yi īn Maẓlūm kih min
ghayr-i jahat iʿrāḍ namūdih sipurdand. Allāhu yaʿlamu mā warada ʿalā
alwāḥihi. Va ān ukht abadan bā mā nabūdih. Qasam bi-Āftāb-i Ḥaqīqat,
baʿd az ẓuhūr-i īn umūr Mīrzā Yaḥyā rā nadīdih va az amr muṭṭaliʿ
nabūdih, chih kih ān ayyām muvāfiq nabūdih-and. Īshān dar maḥallih-ī va
īn Maẓlūm dar maḥallih-yi dīgar sākin, valīkin maḥḍ-i ʿināyat va
maḥabbat va shafaqat chand yawm qabl az ḥarakat nazd-i ū va vālidih-ash
raftih kih shāyad az kawthar-i īmān biyāshāmand va bi-ānchih al-yawm
sabab-i taqarrub ilā Allāh ast fāʾiz gardand. Ḥaqq mīdānad va shāhid va
khud-i ū guvāh kih ghayr-i īn bi-hīch-vajh khayālī nabūdih tā ānkih
al-ḥamdu li-llāh az faḍl-i ilāhī fāʾiz shud va bi-ṭirāz-i maḥabbat
muzayyan gasht. Valīkin baʿd az asīrī va hijrat-i mā az ʿIrāq bi-Āstānih
dīgar az ū khabarī narasīd va baʿd az tafrīq dar arḍ-i Ṭāʾ bā jināb-i
akhavī Mīrzā Riḍā Qulī mulāqāt nashud va khabarī az ū makhṣūṣan narasīd.
Dar avval-i ayyām kull dar yik bayt sākin būdīm va baʿd ān bayt dar
kharīmih bi-qīmat-i nāzilī raft va Farmān-Farmā va Ḥusām al-Salṭanih īn
du barādar kharīdand va qismat namūdand. Baʿd az īn ḥādithih mābayn-i mā
va akhavī tafrīq ḥāṣil, īshān bi-dar-i Masjid-i Shāh va mā dar
Darvāzih-yi Shimīrān sākin. Valīkin az ukht az baʿd min ghayr-i jahat
āthār-i ʿinād ẓāhir. Īn Maẓlūm bi-hīch-vajh sukhanī naguفتih illā ānkih
bint-i akhavī-yi marḥūm Mīrzā Muḥammad Ḥasan, ʿalayhi bahāʾu Allāhi
wa-salāmuhu wa-raḥmatuhu, kih makhṭūbih-yi Ghuṣn-i Aʿẓam būdih, ū rā
ukht-i īn Maẓlūm az Nūr bi-khānih-yi khud burdih va bi-maqarr-i dīgar
firistād. Jamʿī az aṣḥāb va dūstān az aṭrāf shikāyat namūdand, chih kih
īn amr bisyār ʿaẓīm būd. Muvāfiq-i raʾy-i hīch-yik az awliyā-yi Ḥaqq
vāqiʿ nashud. ʿAjab dar īnkih ukht ū rā bi-maḥall-i khud burdih va
tartībāt dādih bi-maqām-i dīgar firistādih, maʿa-dhālik īn Maẓlūm sākit
va ṣāmit būdih va hast, magar ānkih makhṣūṣ-i taskīn-i aḥbāb yik kalimih
iẓhār raft va Ḥaqq shāhid va guvāh ast kih ānchih guftih shud ḥaqīqat
būdih va bi-rāstī guftih va aḥadī az awliyā-yi īn aṭrāf va ān arāḍī
gumān namīkardand kih az ukht chunīn amrī kih khilāf-i ḥamiyyat va
maḥabbat va dūstī ast vāqiʿ shavad. Baʿd az ẓuhūr-i īn amr sabīl rā
maqṭūʿ dīdand va ʿamal namūdand ānchih rā kih kull mīdānīd va mīdānand.
Dīgar maʿlūm ast kih chih martabih ḥuzn az īn ʿamal bar Maẓlūm vārid
shud va baʿd bi-Mīrzā Yaḥyā payvast va ḥāl mukhtalif shinīdih mīshavad.
Maʿlūm nīst chih mīgūyad va chih mīkunad. Nasʾalu Allāha tabāraka
wa-taʿālā an yurjiʿahā ilayhi wa-yuʾayyidahā ʿalā al-inābati ladā bābi
faḍlihi, innahu huwa al-ʿAzīzu al-Tawwābu wa-huwa al-Muqtadiru
al-Ghaffār.

Ḥasan of Mázandarán was carrying seventy Tablets. When he died, those
Tablets were not given to their recipients but were entrusted to one of
this Wronged One's sisters, who had turned away without any reason. God
knows what befell His Tablets. That sister had never been with Us. I
swear by the Sun of Truth that, after these events came to pass, she had
not seen Mírzá Yaḥyá and had not been informed of the Cause, for in
those days they had not been in agreement. She lived in one quarter and
this Wronged One in another. Yet, solely out of kindness, love, and
compassion, a few days before departing We went to see her and her
mother, so that they might drink from the river of faith and attain that
which today brings one near to God. God knows and bears witness, and she
herself can testify, that there was no other intention whatsoever,
until, praise be to God, through divine grace she attained and was
adorned with the ornament of love. But after Our captivity and departure
from Iraq for Constantinople, no further news came from her. After Our
separation in the Land of Ṭá, We did not meet Our brother Mírzá Riḍá
Qulí, nor did any particular news of him arrive. In the early days we
all lived in one house. Later that house was sold off at a low price,
and the two brothers Farmán-Farmá and Ḥusám al-Salṭanih bought and
divided it. After this event, We and Our brother lived apart: he near
the entrance to the Shah Mosque and We at the Shimírán Gate.
Subsequently, however, signs of hostility appeared from Our sister
without any reason. This Wronged One said nothing whatsoever, except
concerning the daughter of Our late brother Mírzá Muḥammad Ḥasan---upon
him be the glory of God, His peace, and His mercy---who was betrothed to
the Most Great Branch. This Wronged One's sister took her from Núr to
her own house and sent her elsewhere. A number of companions and friends
from various places complained, for this was a very grave matter. It met
with the approval of none of the friends of God. It is astonishing that
Our sister took her to her own home, made arrangements, and sent her to
another place. Nevertheless, this Wronged One has remained, and still
remains, silent and uncomplaining, save that a single word was spoken
specifically to soothe the friends. God bears witness and testimony that
what was said was true and was spoken truthfully. None of the friends in
these parts or in those lands imagined that Our sister would do such a
thing, contrary to honor, love, and friendship. After this affair came
to light, they saw the way cut off and did what you all know and what
they know. It is clear, then, how much sorrow this action brought upon
the Wronged One. Afterward she joined Mírzá Yaḥyá, and now conflicting
reports are heard. It is not known what she says or what she does. We
ask God, blessed and exalted, to bring her back to Him and assist her to
turn repentantly toward the door of His grace. He is indeed the
Almighty, the Ever-Relenting, and He is the All-Powerful, the
Ever-Forgiving.

\subsection{246}\label{section-1532}

\begin{otherlanguage}{persian}

و همچنین در مقام دیگر میفرماید اگر در این حین ظاهر شود من اوّل عابدین و
اوّل ساجدینم انتهى یا قوم انصاف دهید مقصود حضرت اعلی آنکه قرب ظهور ناس را
از شریعهٔ باقیهٔ الهیّه منع ننماید چنانچه اصحاب یحیی را از اقرار به روح منع
نمود مکرّر فرموده و میفرماید بیان و آنچه در او نازل شده شما را از آن ساذج
وجود و مالک غیب و شهود منع ننماید مع این حکم محکم اگر کسی به بیان تمسّک
نماید از ظلّ سدرهٔ مبارکهٔ علیا خارج است انصفوا یا قوم و لا تکونوا من
الغافلین

\end{otherlanguage}

Va hamchunīn dar maqām-i dīgar mīfarmāyad agar dar īn ḥīn ẓāhir shavad
man avval-i ʿābidīn va avval-i sājidīnam. Intahā. Yā qawm, inṣāf dihīd.
Maqṣūd-i Ḥaḍrat-i Aʿlā ānkih qurb-i ẓuhūr nās rā az sharīʿih-yi
bāqiyih-yi ilāhiyyih manʿ nanamāyad, chunānchih aṣḥāb-i Yaḥyā rā az
iqrār bi-Rūḥ manʿ namūd. Mukarrar farmūdih va mīfarmāyad Bayān va ānchih
dar ū nāzil shudih shumā rā az ān sādhij-i vujūd va Mālik-i ghayb va
shuhūd manʿ nanamāyad. Maʿ-i īn ḥukm-i muḥkam agar kasī bi-Bayān
tamassuk namāyad az ẓill-i Sidrih-yi Mubārakih-yi ʿUlyā khārij ast.
Anṣifū yā qawmu wa-lā takūnū min al-ghāfilīn.

Likewise, in another passage He says: ``If He should appear at this very
moment, I am the first to worship and the first to bow down.'' End of
quotation. People, be fair! The purpose of His Holiness the Most Exalted
was that the nearness of the Revelation should not keep people from the
everlasting river of God, just as it kept the followers of John from
acknowledging the Spirit. He has repeatedly said, and continues to say:
``Do not let the Bayán and what has been revealed in it keep you from
that pure Essence of existence and Lord of the unseen and the seen.''
Despite this firm decree, if anyone clings to the Bayán, he is outside
the shade of the blessed, exalted Lote Tree. Be fair, people, and do not
be among the heedless.

\subsection{247}\label{section-1533}

\begin{otherlanguage}{persian}

و همچنین میفرماید باسماء از مالک آن محتجب ممانید حتّی اسم النّبیّ فانّ ذلک
الاسم یخلق بقوله و همچنین در باب سابع از واحد ثانی میفرماید ای اهل بیان
نکرده آنچه اهل فرقان کردند که ثمرات لیل خود را باطل کنید الی ان قال عزّ
ذکره میفرماید اگر بظهور او فائز شدی و اطاعت نمودی ثمرهٔ بیان را ظاهر کردی
والّا لائق ذکر نیستی نزد خداوند ترحّم بر خود کرده اگر نصرت نمیکنی مظهر
ربوبیّت را محزون نکرده الی قوله جلّ شأنه اگر بلقآء اللّه فائز نمیگردی آیة
اللّه را هم محزون نکرده باشی از نفع مدیّنین به بیان میگذرد هرگاه شما از
ضرر باو بگذرید اگرچه میدانم نخواهید کرد انتهی

\end{otherlanguage}

Va hamchunīn mīfarmāyad bi-asmāʾ az Mālik-i ān muḥtajib mamānīd, ḥattā
isma al-nabiyyi fa-inna dhālika al-isma yukhlaqu bi-qawlihi. Va
hamchunīn dar bāb-i sābiʿ az vāḥid-i thānī mīfarmāyad ay ahl-i Bayān
nakardih ānchih ahl-i Furqān kardand kih thamarāt-i layl-i khud rā bāṭil
kunīd. Ilā an qāla ʿazza dhikruhu, mīfarmāyad agar bi-ẓuhūr-i ū fāʾiz
shudī va iṭāʿat namūdī thamarِih-yi Bayān rā ẓāhir kardī va illā lāyiq-i
dhikr nīstī nazd-i Khudāvand. Taraḥḥum bar khud kardih agar nuṣrat
namīkunī Maẓhar-i Rubūbiyyat rā maḥzūn nakardih. Ilā qawlihi jalla
shaʾnuhu, agar bi-liqāʾ-i Allāh fāʾiz namīgardī Āyat-i Allāh rā ham
maḥzūn nakardih bāshī. Az nafʿ-i mudayyinīn bi-Bayān mīgudharad hargāh
shumā az ḍarar bi-ū bigudharīd, agarchih mīdānam nakhāhīd kard. Intahā.

He likewise says: ``Do not remain veiled by names from their Lord, even
the name `Prophet,' for that name is created by His word.'' Likewise, in
the seventh chapter of the second Váḥid, He says: ``People of the Bayán,
do not do what the people of the Qur'an did, thereby nullifying the
fruits of your night.'' Further on He says---exalted be His mention:
``If you attain His Revelation and obey Him, you have brought forth the
fruit of the Bayán; otherwise you are not worthy of mention before God.
Have mercy on yourselves: if you do not aid the Manifestation of
Lordship, do not sadden Him.'' And further on He says---glorious is His
station: ``If you do not attain the meeting with God, at least do not
sadden the Sign of God. He will forgo any benefit from those who profess
the Bayán if you will forgo harming Him, although I know that you will
not.'' End of quotation.

\subsection{248}\label{section-1534}

\begin{otherlanguage}{persian}

یا هادی گویا بسبب این بیانات حقّه است اراده نموده‌ئی بیان را محو کنی بشنو
ندای مظلوم را و از این ظلم که ارکان بیان از آن مضطربست بگذر من در چهریق
نبوده‌ام و در ماکو نبوده‌ام حال مابین مریدهای شما حرفها ظاهر شده بعینه
آنچه حزب شیعه میگفتند که این قرآن تمام نیست حضرات هم میگویند این بیان آن
بیان نیست خطّ جناب آقا سیّد حسین موجود خطّ میرزا احمد موجود

\end{otherlanguage}

Yā Hādī, gūyā bi-sabab-i īn bayānāt-i ḥaqqih ast irādih namūdih-ī Bayān
rā maḥv kunī. Bishnaw nidā-yi Maẓlūm rā va az īn ẓulm kih arkān-i Bayān
az ān muḍṭarib ast bigudhar. Man dar Chihrīq nabūdih-am va dar Mākū
nabūdih-am. Ḥāl mābayn-i murīdhā-yi shumā ḥarfhā ẓāhir shudih bi-ʿaynihi
ānchih ḥizb-i Shīʿih mīguftand kih īn Qurʾān tamām nīst. Ḥaḍarāt ham
mīgūyand īn Bayān ān Bayān nīst. Khaṭṭ-i jināb-i Āqā Sayyid Ḥusayn
mawjūd, khaṭṭ-i Mīrzā Aḥmad mawjūd.

O Hádí! It seems that it is because of these true statements that you
have resolved to obliterate the Bayán. Hear the call of the Wronged One
and desist from this injustice, which causes the foundations of the
Bayán to tremble. I was not in Chihríq, nor was I in Mákú. Now
statements have emerged among your followers exactly like those made by
the Shíʿih community, who said that this Qur'an was incomplete. These
gentlemen likewise say that this Bayán is not that Bayán. The
handwriting of Áqá Sayyid Ḥusayn exists; the handwriting of Mírzá Aḥmad
exists.

\subsection{249}\label{section-1535}

\begin{otherlanguage}{arabic}

نفسىی که یک لطمه در دنیا نخورده و همیشه پنج نفس از امآء اللّه نزدش بوده
او را مظلوم میگوئی و حقّ که از اوّل عمر تا بحال دست اعدا بوده و ببدترین
عذابهای عالم معذّب گشته باو نسبت داده‌ئی آنچه را که یهود در حقّ مسیح نگفته
اسمع ندآء المظلوم و لا تکن من الأخسرین

\end{otherlanguage}

Nafsī kih yik laṭmih dar dunyā nakhurdih va hamīshih panj nafs az imāʾ-i
Allāh nazdash būdih ū rā maẓlūm mīgūʾī va Ḥaqq kih az avval-i ʿumr tā
bi-ḥāl dast-i aʿdā būdih va bi-badtarīn ʿadhābhā-yi ʿālam muʿadhdhab
gashtih bi-ū nisbat dādih-ī ānchih rā kih Yahūd dar ḥaqq-i Masīḥ
naguفتih. Ismaʿ nidāʾa al-Maẓlūmi wa-lā takun min al-akhsarīn.

You call a man wronged who has not suffered a single blow in this world
and has always had five of the handmaidens of God with him, while to the
Truth, who from the beginning of His life until now has been in the
hands of enemies and has suffered the worst torments in the world, you
have attributed what the Jews did not say of Christ. Hear the call of
the Wronged One and do not be among those who suffer the greatest loss.

\subsection{250}\label{section-1536}

\begin{otherlanguage}{arabic}

و همچنین میفرماید چه بسا ناری را که خدا نور میکند به من یظهره اللّه و چه
بسا نوری را که نار میفرماید به او و می‌بینم ظهور او را مثل این شمس در وسط
السّمآء و غروب کل را بمثل نجوم لیل در نهار انتهی یا عالم هل لک اذن لتسمع
ندآء الحقّ و تنصف فى هذا الظّهور الّذی اذ ظهر نطق الطّور قد اتی مکلّمی بآیات
واضحات و براهین لائحات رغماً لکلّ غافل بعید و کلّ مفتر کذّاب الّذین ارادوا ان
یطفئوا نور اللّه بمفتریاتهم و یمحوا آثار اللّه بغلّهم الا انّهم من الظّالمین
فی کتاب اللّه ربّ العالمین

\end{otherlanguage}

Va hamchunīn mīfarmāyad: Chih basā nārī rā kih Khudā nūr mīkunad bih Man
Yuẓhiruhu Allāh va chih basā nūrī rā kih nār mīfarmāyad bih ū va mībīnam
ẓuhūr-i ū rā mithl-i īn shams dar vasaṭ al-samāʾ va ghurūb-i kull rā
bi-mithl-i nujūm-i layl dar nahār. Intahā. Yā ʿālim, hal laka udhun
li-tasmaʿa nidāʾ al-Ḥaqq wa tunṣifa fī hādhā al-ẓuhūr alladhī idh ẓahara
naṭaqa al-Ṭūr: Qad atā mukallimī bi-āyāt wāḍiḥāt wa barāhīn lāʾiḥāt
raghman li-kulli ghāfil baʿīd wa kulli muftarin kadhdhāb alladhīna arādū
an yuṭfiʾū nūr Allāh bi-muftarayātihim wa yamḥū āthār Allāh
bi-ghillihim. Alā innahum min al-ẓālimīn fī kitāb Allāh Rabb al-ʿālamīn.

He also says: ``How often God turns fire into light through Him Whom God
will make manifest, and how often He turns light into fire through Him!
I see His appearance like this sun in the midst of heaven, and the
setting of all others like the stars of the night in the daytime.'' End
of quotation. O learned one! Do you have an ear to hear the call of the
Truth and judge fairly concerning this Revelation, at whose appearance
the Mount declared: ``He who speaks to me has come with clear verses and
manifest proofs, despite every heedless one who is far away and every
lying fabricator who sought to extinguish the light of God with their
fabrications and erase the traces of God through their malice.'' Indeed,
they are among the unjust in the Book of God, the Lord of the worlds.

\subsection{251}\label{section-1537}

\begin{otherlanguage}{arabic}

و همچنین میفرماید بیان از اوّل تا آخر مکمن جمیع صفات اوست و خزانهٔ نار و
نور او انتهی سبحان‌اللّه عرف این بیان انسان را اخذ مینماید چه که بکمال حزن
میفرماید آنچه را که مشاهده مینماید و همچنین بجناب حرف حیّ یعنی ملّا باقر
علیه بهآء اللّه و عنایته میفرماید لعلّک فی ثمانیة سنة یوم ظهوره تدرک لقائه
انتهی

\end{otherlanguage}

Va hamchunīn mīfarmāyad: Bayān az avval tā ākhir makman-i jamīʿ-i
ṣifāt-i ūst va khazānih-yi nār va nūr-i ū. Intahā. Subḥān Allāh, ʿarf-i
īn bayān insān rā akhdh mīnamāyad chih kih bi-kamāl-i ḥuzn mīfarmāyad
ānchih rā kih mushāhadih mīnamāyad. Va hamchunīn bih janāb-i Ḥarf-i Ḥayy
yaʿnī Mullā Bāqir, ʿalayhi Bahāʾ Allāh wa ʿināyatuhu, mīfarmāyad:
Laʿallaka fī thamāniyati sanah yawma ẓuhūrihi tudriku liqāʾahu. Intahā.

He also says: ``The Bayán, from beginning to end, is the repository of
all His attributes and the treasury of His fire and light.'' End of
quotation. Glory be to God! The fragrance of this utterance seizes the
soul, for with the deepest sorrow He speaks of what He beholds. He also
says to the Letter of the Living, namely Mullā Bāqir, upon whom be the
glory of God and His loving care: ``Perhaps in the eighth year, on the
day of His appearance, you will attain His presence.'' End of quotation.

\subsection{252}\label{section-1538}

\begin{otherlanguage}{persian}

اعرف یا هادی و کن من السّامعین انصاف ده اکثر اصحاب الهی و ادلّای حقّ شهید
شدند تو هنوز موجودی آیا حفظ تو از چه بوده لعمر اللّه از انکار و شهادت
نفوس مقدّسه از اقرار هر صاحب عدل و انصافی باین فقره گواهی میدهد چه که سبب
و علّت این دو بمثابهٔ آفتاب ظاهر و مشهود

\end{otherlanguage}

Iʿraf yā Hādī wa kun min al-sāmiʿīn. Inṣāf dih, akthar-i aṣḥāb-i ilāhī
va adillā-yi Ḥaqq shahīd shudand, tu hanūz mawjūdī. Āyā ḥifẓ-i tu az
chih būdih? La-ʿamr Allāh, az inkār va shahādat-i nufūs-i muqaddasih az
iqrār. Har ṣāḥib-i ʿadl va inṣāfī bi-īn fiqrih guvāhī mīdihad chih kih
sabab va ʿillat-i īn du bi-mathābih-yi āftāb ẓāhir va mashhūd.

Recognize the truth, O Hādī, and be among those who listen. Judge
fairly: most of the companions of God and the guides to the Truth were
martyred, while you still remain alive. What preserved you? By the life
of God, it was denial, while the martyrdom of those holy souls resulted
from their acknowledgment. Every just and fair-minded person bears
witness to this, for the cause and reason for both are as evident and
visible as the sun.

\subsection{253}\label{section-1539}

\begin{otherlanguage}{persian}

و همچنین خطاب به دیّان مظلوم شهید میفرماید ستعرفنّ قدرک بقول من یظهره اللّه
و همچنین او را حرف ثالث مؤمن به من یظهره اللّه فرموده بقوله و انّک انت یا
حرف الثّالث المؤمن بمن یظهره اللّه و همچنین میفرماید ولکنّ اللّه اذا شآء
لیعرّفنّک بقول من یظهره اللّه حضرت دیّان که بقول نقطه روح ما سواه فداه مخزن
امانت حقّ جلّ جلاله و مکمن لآلی علم اوست او را بظلمی شهید نمودند که ملأ
اعلى گریست و نوحه نمود و اوست نفسی که علم مکنون مخزون را باو تعلیم
فرموده و در او ودیعه گذاشته بقوله ان یا اسم الدّیّان هذا علم مکنون مخزون
قد اودعناک و آتیناک عزّاً من عند اللّه اذ عین فؤادک لطیف تعرف قدره و تعزّ
بهائه و قد منّ اللّه علی نقطة البیان بعلم مکنون مخزون ما نزّل اللّه قبل ذلک
الظّهور و هو اعزّ من کلّ علم عند اللّه سبحانه و قد جعله حجّة من عنده بمثل ما
قد جعل الآیات حجّة من عنده انتهی آن مظلوم که دارای خزینهٔ علم الهی بود مع
جناب میرزا على اکبر از منتسبین نقطه علیه بهآء اللّه و رحمته و جناب آقا
ابوالقاسم کاشی و جمعی دیگر بفتوای میرزا یحیی کل را شهید نمودند

\end{otherlanguage}

Va hamchunīn khiṭāb bih Dayyān-i maẓlūm-i shahīd mīfarmāyad:
Sa-taʿrifanna qadraka bi-qawli Man Yuẓhiruhu Allāh. Va hamchunīn ū rā
ḥarf-i thālith-i muʾmin bih Man Yuẓhiruhu Allāh farmūdih bi-qawlihi: Wa
innaka anta yā ḥarf al-thālith al-muʾmin bi-Man Yuẓhiruhu Allāh. Va
hamchunīn mīfarmāyad: Wa lākinna Allāha idhā shāʾa la-yuʿarrifannaka
bi-qawli Man Yuẓhiruhu Allāh. Ḥaḍrat-i Dayyān kih bi-qawl-i Nuqṭih,
rūḥ-i mā sivāh fidāh, makhzan-i amānat-i Ḥaqq jalla jalāluhu va makman-i
laʾālī-yi ʿilm-i ūst, ū rā bih ẓulmī shahīd namūdand kih malaʾ-i aʿlā
girīst va nawḥih namūd va ūst nafsī kih ʿilm-i maknūn-i makhzūn rā bi-ū
taʿlīm farmūdih va dar ū vadīʿih guẕāshtih bi-qawlihi: An yā ism
al-Dayyān, hādhā ʿilm maknūn makhzūn qad awdaʿnāka wa ātaynāka ʿizzan
min ʿind Allāh idh ʿayn fuʾādika laṭīf taʿrifu qadrahu wa taʿizzu
bahāʾahu wa qad manna Allāhu ʿalā Nuqṭat al-Bayān bi-ʿilm maknūn makhzūn
mā nazzala Allāhu qabla dhālika al-ẓuhūr wa huwa aʿazzu min kulli ʿilm
ʿinda Allāh subḥānahu wa qad jaʿalahu ḥujjatan min ʿindihi bi-mithli mā
qad jaʿala al-āyāti ḥujjatan min ʿindihi. Intahā. Ān maẓlūm kih dārā-yi
khazīnih-yi ʿilm-i ilāhī būd maʿa janāb-i Mīrzā ʿAlī Akbar az
muntasibīn-i Nuqṭih, ʿalayhi Bahāʾ Allāh wa raḥmatuhu, va janāb-i Āqā
Abu'l-Qāsim Kāshī va jamʿī dīgar bi-fatvā-yi Mīrzā Yaḥyā kull rā shahīd
namūdand.

He also addresses Dayyān, the wronged martyr, saying: ``You will
assuredly recognize your worth through the words of Him Whom God will
make manifest.'' He also called him the third Letter to believe in Him
Whom God will make manifest, in His words: ``You indeed, O third Letter,
are the believer in Him Whom God will make manifest.'' He also says:
``But God, if He wills, will assuredly make you known through the words
of Him Whom God will make manifest.'' Dayyān, who, according to the
words of the Point---may the soul of all else be sacrificed for Him---is
the repository of the trust of the Truth, exalted be His glory, and the
treasury of the pearls of His knowledge, was martyred with such
injustice that the celestial Assembly wept and lamented. He is the one
whom He taught the hidden, treasured knowledge, depositing it within
him, in His words: ``O you who bear the name Dayyān! This is a hidden,
treasured knowledge that We have entrusted to you, and We have bestowed
upon you honor from God, for the eye of your heart is discerning: you
recognize its worth and cherish its splendor. God has indeed favored the
Point of the Bayán with a hidden, treasured knowledge that God had not
sent down before that Revelation. It is more precious than all knowledge
in the sight of God, glory be to Him, and He has made it a proof from
Himself, just as He has made the verses a proof from Himself.'' End of
quotation. That wronged one, who possessed the treasury of divine
knowledge, together with Mīrzā ʿAlī Akbar, one of the relatives of the
Point, upon whom be the glory of God and His mercy, and Āqā Abu'l-Qāsim
Kāshī and a number of others---all were martyred by the decree of Mīrzā
Yaḥyā.

\subsection{254}\label{section-1540}

\begin{otherlanguage}{persian}

یا هادی کتابش نزد تو حاضر است آنکه اسمش را مستیقظ گذارده بخوان اگرچه
دیده‌ئی ولکن مکرّر مشاهده نما لعلّک تتّخذ لنفسک فی خبآء الصّدق مقاماً رفیعاً

\end{otherlanguage}

Yā Hādī, kitābash nazd-i tu ḥāḍir ast, ānkih ismash rā Mustayqiẓ
guẕāshtih. Bikhvān agarchih dīdih-ī va lākin mukarrar mushāhadih namā
laʿallaka tattakhidhu li-nafsika fī khibāʾ al-ṣidq maqāman rafīʿan.

O Hādī! His book is there with you, the one he named Mustayqiẓ. Read it;
although you have seen it, examine it again, so that you may attain for
yourself a lofty station within the tent of truthfulness.

\subsection{255}\label{section-1541}

\begin{otherlanguage}{persian}

و همچنین آقا سیّد ابرهیم که در باره‌اش از قلم نقطهٔ اولی این کلمات جاری
قوله تعالی ان یا خلیلی فی الصّحف و ان یا ذکری فی الکتب من بعد الصّحف و ان
یا اسمی فى البیان انتهی او و دیّان را ابوالشّرور و ابوالدّواهی نامیده حال
انصاف ده که بر این مظلومها چه وارد شده مع آنکه یکی در خدمتش مشغول و
دیگری بر او وارد باری لعمر اللّه باعمالی عامل بوده که قلم حیا میکند از
ذکرش

\end{otherlanguage}

Va hamchunīn Āqā Siyyid Ibrāhīm kih dar bārih-ash az qalam-i Nuqṭih-yi
Ūlā īn kalimāt jārī, qawluhu taʿālā: An yā khalīlī fī al-ṣuḥuf wa an yā
dhikrī fī al-kutub min baʿd al-ṣuḥuf wa an yā ismī fī al-Bayān. Intahā.
Ū va Dayyān rā Abu'l-Shurūr va Abu'l-Davāhī nāmīdih. Ḥāl inṣāf dih kih
bar īn maẓlūmhā chih vārid shudih maʿa ānkih yikī dar khidmatash
mashghūl va dīgarī bar ū vārid. Bārī la-ʿamr Allāh bi-aʿmālī ʿāmil būdih
kih qalam ḥayā mīkunad az dhikrash.

Likewise, concerning Āqā Siyyid Ibrāhīm, these words flowed from the Pen
of the Primal Point---exalted are His words: ``O My friend in the
Scrolls! O My remembrance in the Books after the Scrolls! O My name in
the Bayán!'' End of quotation. He has called him and Dayyān the ``Father
of Evils'' and the ``Father of Calamities.'' Now judge fairly what has
befallen these wronged ones, even though one was occupied in his service
and the other had come to visit him. In short, by the life of God, he
has committed deeds that the pen is ashamed to mention.

\subsection{256}\label{section-1542}

\begin{otherlanguage}{persian}

قدری در عصمت نقطهٔ اولی تفکّر کن ملاحظه نما چه ظاهر گشته وقتی که این مظلوم
از هجرت دوساله که در صحاری و جبال سالک بود و بسبب بعضی از نفوس که مدّتها
در بیابانها دویدند رجوع به دارالسّلام نمود میرزا محمّد على نامی رشتی بحضور
آمد و امام جمعی بکلمه‌ئی نطق نمود در بارهٔ عصمت آن حضرت که فی‌الحقیقه حزن
جمیع اقطار را اخذ نمود سبحان‌اللّه چگونه راضی شدند که باین خیانت اعظم تمسّک
جستند باری از حقّ میطلبیم که عامل را توفیق بخشد بر توبه و انابه انّه هو
المؤیّد الحکیم

\end{otherlanguage}

Qadrī dar ʿiṣmat-i Nuqṭih-yi Ūlā tafakkur kun. Mulāḥaẓih namā chih ẓāhir
gashtih. Vaqtī kih īn Maẓlūm az hijrat-i du-sālih kih dar ṣaḥārī va
jibāl sālik būd va bi-sabab-i baʿḍī az nufūs kih muddat-hā dar biyābānhā
davīdand rujūʿ bih Dār al-Salām namūd, Mīrzā Muḥammad ʿAlī nāmī Rashtī
bi-ḥuḍūr āmad va amām-i jamʿī bih kalimih-ī nuṭq namūd dar bārih-yi
ʿiṣmat-i ān Ḥaḍrat kih fī al-ḥaqīqih ḥuzn jamīʿ-i aqṭār rā akhdh namūd.
Subḥān Allāh, chigūnih rāḍī shudand kih bi-īn khiyānat-i aʿẓam tamassuk
justand. Bārī az Ḥaqq mīṭalabīm kih ʿāmil rā tawfīq bakhshad bar tawbih
va inābih. Innahu huwa al-Muʾayyid al-Ḥakīm.

Reflect a little on the infallibility of the Primal Point. Consider what
has come to light. When this Wronged One returned to the Abode of Peace
from His two-year withdrawal, during which He wandered through deserts
and mountains, at the urging of certain souls who had spent a long time
searching the wilderness, a man from Rasht named Mīrzā Muḥammad ʿAlī
came into His presence. Before a gathering he made a statement
concerning the infallibility of that exalted Being which truly caused
sorrow to seize every region. Glory be to God! How could they consent to
resort to this most grievous betrayal? In short, We ask the Truth to
enable its perpetrator to repent and turn back to Him. He is indeed the
One who aids, the All-Wise.

\subsection{257}\label{section-1543}

\begin{otherlanguage}{persian}

و جناب دیّان علیه بهآء اللّه و رحمته بحضور فائز مطابق آنچه از قلم نقطهٔ
اولی ظاهر شد نسأل اللّه ان یؤیّد الغافلین علی التّوجّه الیه و المعرضین علی
الاقبال الی شطره و المنکرین علی التّصدیق علی هذا الأمر الّذی اذ ظهر نطقت
الأشیآء کلّها قد اتی من کان مکنوناً فی کنز العلم و مرقوماً من القلم الأعلی
فی الکتب و الصّحف و الزّبر و الألواح

\end{otherlanguage}

Va janāb-i Dayyān, ʿalayhi Bahāʾ Allāh wa raḥmatuhu, bi-ḥuḍūr fāʾiz
muṭābiq-i ānchih az qalam-i Nuqṭih-yi Ūlā ẓāhir shud. Nasʾalu Allāha an
yuʾayyida al-ghāfilīna ʿalā al-tawajjuh ilayhi wa al-muʿriḍīna ʿalā
al-iqbāl ilā shaṭrihi wa al-munkirīna ʿalā al-taṣdīq ʿalā hādhā al-Amr
alladhī idh ẓahara naṭaqat al-ashyāʾ kulluhā: Qad atā man kāna maknūnan
fī kanz al-ʿilm wa marqūman min al-Qalam al-Aʿlā fī al-kutub wa al-ṣuḥuf
wa al-zubur wa al-alwāḥ.

Dayyān, upon whom be the glory of God and His mercy, attained His
presence in accordance with what had been revealed by the Pen of the
Primal Point. We ask God to help the heedless turn toward Him, those who
have turned away to approach Him, and the deniers to acknowledge this
Cause, at whose appearance all things declared: ``He has come who was
concealed in the treasury of knowledge and recorded by the Most Exalted
Pen in the Books, the Scrolls, the Scriptures, and the Tablets.''

\subsection{258}\label{section-1544}

\begin{otherlanguage}{persian}

در این مقام لازم شد احادیثی که در شأن این مدینهٔ مبارکهٔ مشرّفه یعنی عکّا
وارد شده ذکر شود لعلّک تتّخذ یا هادی الى الصّدق طریقاً و الی اللّه سبیلاً

\end{otherlanguage}

Dar īn maqām lāzim shud aḥādīthī kih dar shaʾn-i īn madīnih-yi
mubārakih-yi musharrafih yaʿnī ʿAkkā vārid shudih dhikr shavad laʿallaka
tattakhidhu yā Hādī ilā al-ṣidq ṭarīqan wa ilā Allāh sabīlan.

At this point it has become necessary to mention the traditions
concerning this blessed and honored city, namely ʿAkkā, so that you, O
Hādī, may take a path toward truthfulness and a way toward God.

\subsection{259}\label{section-1545}

\begin{otherlanguage}{arabic}

بسم اللّه الرّحمن الرّحیم

\end{otherlanguage}

Bism Allāh al-Raḥmān al-Raḥīm.

In the name of God, the All-Merciful, the Ever-Merciful.

\subsection{260}\label{section-1546}

\begin{otherlanguage}{arabic}

ما ورد فی فضل عکّا و البحر و عین البقر الّتی بعکّا حدّثنا

\end{otherlanguage}

Mā warada fī faḍl ʿAkkā wa al-baḥr wa ʿAyn al-Baqar allatī bi-ʿAkkā.
Ḥaddathanā

What has been related concerning the excellence of ʿAkkā, the sea, and
the Spring of the Cow in ʿAkkā. It was related to us by

\subsection{261}\label{section-1547}

\begin{otherlanguage}{arabic}

عبدالعزیز بن عبدالسّلام عن النّبیّ صلّی اللّه علیه و سلّم انّه قال انّ عکّا مدینة
بالشّام قد اختصّها اللّه برحمته

\end{otherlanguage}

ʿAbd al-ʿAzīz ibn ʿAbd al-Salām ʿan al-Nabī ṣallā Allāhu ʿalayhi wa
sallam annahu qāla: Inna ʿAkkā madīnah bi-al-Shām qad ikhtaṣṣahā Allāhu
bi-raḥmatihi.

ʿAbd al-ʿAzīz ibn ʿAbd al-Salām, on the authority of the Prophet, may
God bless him and grant him peace, that he said: ``ʿAkkā is a city in
Syria that God has singled out for His mercy.''

\subsection{262}\label{section-1548}

\begin{otherlanguage}{arabic}

و قال ابن مسعود رضی اللّه عنه انّ النّبیّ صلّی اللّه علیه و سلّم قال الا انّ
افضل السّواحل عسقلان و انّ عکّا افضل من عسقلان و فضل عکّا علی عسقلان و علی
جمیع السّواحل کفضل محمّد علی جمیع الأنبیآء الا اخبرکم بمدینة بین جبلین فی
الشّام فی وسط المرج یقال لها عکّا الا و انّ من دخلها راغباً فیها و فی
زیارتها غفر اللّه له ما تقدّم من ذنبه و ما تأخّر و من خرج منها غیر زائر الا
لم یبارک اللّه له فی خروجه الا و انّ فیها عیناً یقال لها عین البقر من شرب
منها شربةً ملأ اللّه قلبه نوراً و امّنه من العذاب الأکبر یوم القیامة

\end{otherlanguage}

Wa qāla Ibn Masʿūd raḍiya Allāhu ʿanhu inna al-Nabī ṣallā Allāhu ʿalayhi
wa sallam qāla: Alā inna afḍala al-sawāḥil ʿAsqalān wa inna ʿAkkā afḍalu
min ʿAsqalān wa faḍlu ʿAkkā ʿalā ʿAsqalān wa ʿalā jamīʿ al-sawāḥil
ka-faḍli Muḥammad ʿalā jamīʿ al-anbiyāʾ. Alā ukhbirukum bi-madīnah bayna
jabalayn fī al-Shām fī wasaṭ al-marj yuqālu lahā ʿAkkā? Alā wa inna man
dakhalahā rāghiban fīhā wa fī ziyāratihā ghafara Allāhu lahu mā
taqaddama min dhanbihi wa mā taʾakhkhara wa man kharaja minhā ghayra
zāʾir alā lam yubārik Allāhu lahu fī khurūjihi. Alā wa inna fīhā ʿaynan
yuqālu lahā ʿAyn al-Baqar man shariba minhā sharbatan malaʾa Allāhu
qalbahu nūran wa āmmanahu min al-ʿadhāb al-akbar yawma al-qiyāmah.

Ibn Masʿūd, may God be pleased with him, said that the Prophet, may God
bless him and grant him peace, said: ``Indeed, the best of the coastal
towns is Ascalon, and ʿAkkā is better than Ascalon. The excellence of
ʿAkkā over Ascalon and all the coastal towns is like the excellence of
Muḥammad over all the Prophets. Shall I not tell you of a city between
two mountains in Syria, in the midst of a meadow, called ʿAkkā? Indeed,
whoever enters it out of longing for it and to visit it, God will
forgive his past and future sins; and whoever leaves it without having
visited it, God will not bless his departure. Indeed, within it is a
spring called the Spring of the Cow. Whoever drinks a draught from it,
God will fill his heart with light and protect him from the greatest
torment on the Day of Resurrection.''

\subsection{263}\label{section-1549}

\begin{otherlanguage}{arabic}

و عن انس بن مالک رضی اللّه عنه قال قال رسول اللّه صلّی اللّه علیه و سلّم انّ
فی السّواحل مدینة معلّقة تحت ساق العرش یقال لها عکّا من بات فیها مرابطاً
احتساباً للّه تعالی کتب اللّه له ثواب الصّابرین و القائمین و الرّاکعین و
السّاجدین الی یوم القیامة

\end{otherlanguage}

Wa ʿan Anas ibn Mālik raḍiya Allāhu ʿanhu qāla: Qāla Rasūl Allāh ṣallā
Allāhu ʿalayhi wa sallam: Inna fī al-sawāḥil madīnatan muʿallaqatan
taḥta sāq al-ʿarsh yuqālu lahā ʿAkkā man bāta fīhā murābiṭan iḥtisāban
li-llāh taʿālā kataba Allāhu lahu thawāba al-ṣābirīn wa al-qāʾimīn wa
al-rākiʿīn wa al-sājidīn ilā yawmi al-qiyāmah.

Anas ibn Mālik, may God be pleased with him, related: ``The Messenger of
God, may God bless him and grant him peace, said: `Among the coastal
towns is a city suspended beneath the leg of the Throne, called ʿAkkā.
Whoever spends a night there guarding the frontier, seeking the reward
of God, exalted be He, God will record for him the reward of those who
are patient, who stand in prayer, who bow, and who prostrate themselves,
until the Day of Resurrection.'\,''

\subsection{264}\label{section-1550}

\begin{otherlanguage}{arabic}

و قال صلّی اللّه علیه و سلّم الا اخبرکم بمدینة علی شاطئ البحر بیضآء حسن
بیاضها عند اللّه تعالی یقال لها عکّا و انّ من قرصه برغوث من براغیثها کان
عند اللّه افضل من طعنة نافذة فی سبیل اللّه الا و انّ من اذّن فیها کان له مدّ
صوته فی الجنّة و من قعد فیها سبعة ایّام مقابل العدوّ حشره اللّه مع الخضر
علیه السّلام و امّنه اللّه من الفزع الأکبر یوم القیامة و قال صلّی اللّه تعالی
علیه و سلّم الا و انّ فی الجنّة ملوکاً و سادات و فقرآء عکّا ملوک الجنّة و
ساداتها و انّ شهراً فی عکّا افضل من الف سنة فی غیرها

\end{otherlanguage}

Wa qāla ṣallā Allāhu ʿalayhi wa sallam: Alā ukhbirukum bi-madīnah ʿalā
shāṭiʾ al-baḥr bayḍāʾ ḥasan bayāḍuhā ʿinda Allāh taʿālā yuqālu lahā
ʿAkkā wa inna man qaraṣahu burghūth min barāghīthihā kāna ʿinda Allāh
afḍala min ṭaʿnah nāfidhah fī sabīl Allāh. Alā wa inna man adhdhana fīhā
kāna lahu maddu ṣawtihi fī al-jannah wa man qaʿada fīhā sabʿata ayyām
muqābila al-ʿaduww ḥasharahu Allāhu maʿa al-Khiḍr ʿalayhi al-salām wa
āmmanahu Allāhu min al-fazaʿ al-akbar yawma al-qiyāmah. Wa qāla ṣallā
Allāhu taʿālā ʿalayhi wa sallam: Alā wa inna fī al-jannah mulūkan wa
sādāt wa fuqarāʾ ʿAkkā mulūk al-jannah wa sādātuhā wa inna shahran fī
ʿAkkā afḍalu min alfi sanah fī ghayrihā.

He, may God bless him and grant him peace, also said: ``Shall I not tell
you of a white city on the seashore, whose whiteness is beautiful in the
sight of God, exalted be He, called ʿAkkā? If someone is bitten by one
of its fleas, that is better in the sight of God than a piercing wound
received in the path of God. Indeed, whoever gives the call to prayer
there will have a place in Paradise as far-reaching as his voice.
Whoever remains there for seven days facing the enemy, God will raise
him up with al-Khiḍr, peace be upon him, and protect him from the
greatest terror on the Day of Resurrection.'' He, may God, exalted be
He, bless him and grant him peace, also said: ``Indeed, in Paradise
there are kings and lords, and the poor of ʿAkkā are the kings and lords
of Paradise. A month in ʿAkkā is better than a thousand years
elsewhere.''

\subsection{265}\label{section-1551}

\begin{otherlanguage}{arabic}

و عن رسول اللّه صلّی اللّه علیه و سلّم انّه قال طوبی لمن زار عکّا و طوبی لمن
زار زائر عکّا طوبی لمن شرب من عین البقر و اغتسل من مائها فانّ الحور العین
یشربن الکافور الّذی فی الجنّة من عین البقر و عین سلوان و بئر زمزم طوبی لمن
شرب من هؤلآء العیون و اغتسل من مائهنّ فقد حرّم اللّه علیه و علی جسده نار
جهنّم یوم القیامة

\end{otherlanguage}

Wa ʿan Rasūl Allāh ṣallā Allāhu ʿalayhi wa sallam annahu qāla: Ṭūbā
li-man zāra ʿAkkā wa ṭūbā li-man zāra zāʾira ʿAkkā. Ṭūbā li-man shariba
min ʿAyn al-Baqar wa ightasala min māʾihā fa-inna al-ḥūr al-ʿīn
yashrabna al-kāfūr alladhī fī al-jannah min ʿAyn al-Baqar wa ʿAyn Silwān
wa Biʾr Zamzam. Ṭūbā li-man shariba min hāʾulāʾ al-ʿuyūn wa ightasala
min māʾihinna fa-qad ḥarrama Allāhu ʿalayhi wa ʿalā jasadihi nāra
jahannam yawma al-qiyāmah.

It is related of the Messenger of God, may God bless him and grant him
peace, that he said: ``Blessed is the one who visits ʿAkkā, and blessed
is the one who visits a visitor to ʿAkkā. Blessed is the one who drinks
from the Spring of the Cow and bathes in its water, for the wide-eyed
maidens of Paradise drink the camphor of Paradise from the Spring of the
Cow, the Spring of Silwān, and the Well of Zamzam. Blessed is the one
who drinks from these springs and bathes in their waters, for God has
forbidden the fire of hell to touch him or his body on the Day of
Resurrection.''

\subsection{266}\label{section-1552}

\begin{otherlanguage}{persian}

و عن النّبیّ صلّی اللّه علیه و سلّم انّه قال فی عکّا نوافل و فواضل یخصّ اللّه بها
من یشآء من قال فی عکّا سبحان اللّه و الحمد للّه و لا اله الّا اللّه و اللّه
اکبر و لا حول و لا قوّة الّا باللّه العلیّ العظیم کتب اللّه له الف حسنة و محا
عنه الف سیّئة و رفع له الف درجة فی الجنّة و غفر له ذنوبه و من قال فی عکّا
استغفر اللّه غفر اللّه له ذنوبه کلّها و من ذکر اللّه فی عکّا بالغدوّ و الآصال
و العشیّ و الابکار کان عند اللّه افضل من نقل السّیوف و الرّماح و السّلاح فی
سبیل اللّه تعالی

\end{otherlanguage}

Wa ʿan al-Nabī ṣallā Allāhu ʿalayhi wa sallam annahu qāla: Fī ʿAkkā
nawāfil wa fawāḍil yakhuṣṣu Allāhu bihā man yashāʾ. Man qāla fī ʿAkkā
subḥān Allāh wa al-ḥamdu li-llāh wa lā ilāha illā Allāh wa Allāhu akbar
wa lā ḥawla wa lā quwwata illā bi-llāh al-ʿAlī al-ʿAẓīm kataba Allāhu
lahu alfa ḥasanah wa maḥā ʿanhu alfa sayyiʾah wa rafaʿa lahu alfa
darajah fī al-jannah wa ghafara lahu dhunūbahu. Wa man qāla fī ʿAkkā
astaghfiru Allāh ghafara Allāhu lahu dhunūbahu kullahā wa man dhakara
Allāha fī ʿAkkā bi-al-ghuduww wa al-āṣāl wa al-ʿashiyy wa al-ibkār kāna
ʿinda Allāh afḍala min naql al-suyūf wa al-rimāḥ wa al-silāḥ fī sabīl
Allāh taʿālā.

It is related of the Prophet, may God bless him and grant him peace,
that he said: ``In ʿAkkā there are bounties and favors that God bestows
on whomever He wills. Whoever says in ʿAkkā, `Glory be to God, and
praise be to God; there is no god but God, and God is most great; there
is no power and no strength except through God, the Most High, the
Almighty,' God will record for him a thousand good deeds, erase a
thousand evil deeds, raise him a thousand degrees in Paradise, and
forgive his sins. Whoever says in ʿAkkā, `I ask God for forgiveness,'
God will forgive all his sins. Whoever remembers God in ʿAkkā in the
mornings and late afternoons, in the evenings and at daybreak, will be
better in the sight of God than the carrying of swords, spears, and
weapons in the path of God, exalted be He.''

\subsection{267}\label{section-1553}

\begin{otherlanguage}{arabic}

و قال رسول اللّه صلّی اللّه علیه و سلّم من نظر بالبحر عند الزّوال و کبّر اللّه
عند الغروب غفر اللّه له ذنوبه ولو کانت مثل رمل عالج و من عدّ اربعین موجة و
هو یکبّر اللّه تعالی غفر اللّه له ما تقدّم من ذنبه و ما تأخّر

\end{otherlanguage}

Wa-qāla rasūlu Allāh ṣallā Allāhu ʿalayhi wa-sallam: man naẓara
bi-l-baḥr ʿinda al-zawāl wa-kabbara Allāh ʿinda al-ghurūb ghafara Allāhu
lahu dhunūbahu wa-law kānat mithla ramli ʿĀlij, wa-man ʿadda arbaʿīna
mawjatan wa-huwa yukabbiru Allāha taʿālā ghafara Allāhu lahu mā
taqaddama min dhanbihi wa-mā taʾakhkhar.

And the Messenger of God, may God bless him and grant him peace, said:
``Whoever gazes at the sea at noon and proclaims God's greatness at
sunset, God will forgive his sins, even if they are like the sands of
ʿĀlij. Whoever counts forty waves while proclaiming the greatness of
God, exalted be He, God will forgive his past and future sins.''

\subsection{268}\label{section-1554}

\begin{otherlanguage}{arabic}

و قال رسول اللّه صلّی اللّه علیه و سلّم من نظر الی البحر لیلة کاملة کان افضل
من شهرین کاملین بین الرّکن و المقام و من تربّی فی السّواحل خیر ممّن تربّی فی
غیرها و النّائم فی السّواحل کالقائم فی غیرها انتهی

\end{otherlanguage}

Wa-qāla rasūlu Allāh ṣallā Allāhu ʿalayhi wa-sallam: man naẓara ilā
al-baḥr laylatan kāmilatan kāna afḍala min shahrayn kāmilayn bayna
al-rukn wa-l-maqām, wa-man tarabbā fī al-sawāḥil khayrun mimman tarabbā
fī ghayrihā, wa-l-nāʾim fī al-sawāḥil ka-l-qāʾim fī ghayrihā. Intahā.

And the Messenger of God, may God bless him and grant him peace, said:
``For whoever gazes at the sea for an entire night, this is better than
two full months between the Corner and the Station. Whoever is brought
up on the shores is better than one brought up elsewhere, and one who
sleeps on the shores is like one who keeps vigil elsewhere.'' End of
quotation.

\subsection{269}\label{section-1555}

\begin{otherlanguage}{arabic}

صدق رسول اللّه صلّی اللّه تعالی علیه و سلّم

\end{otherlanguage}

Ṣadaqa rasūlu Allāh ṣallā Allāhu taʿālā ʿalayhi wa-sallam.

The Messenger of God, may God, exalted be He, bless him and grant him
peace, has spoken the truth.

\bookmarksetup{startatroot}

\chapter{Lawh-i-Sarraj (Tablet to
Sarraj)}\label{lawh-i-sarraj-tablet-to-sarraj}

\section{BH00006}\label{bh00006}

\emph{Translated by GPT 6.1 Sol on 15 Qudrat 183}

\subsection{1}\label{section-1556}

\begin{otherlanguage}{arabic}

هو المجیب

\end{otherlanguage}

Huwa al-Mujīb

He is the Answerer.

\subsection{2}\label{section-1557}

\begin{otherlanguage}{persian}

نامهٴ آن جناب بین یدی العرش حاضر و بر مقرّ انّه لا یعرف بما سواه واصل و ما
فیه بلحاظ اللّه ملحوظ آمد سؤال شده بود از اینکه چگونه میشود که حرف علّیّین
بسجّین تبدیل شود و یا اثبات بنفی راجع گردد و یا ثمرهٴ جنیّه از لطافت ممنوع
شود و یا مرآت از اشراقات انوار آفتاب معانی محروم ماند فنعم ما سألت و کنت
من السّائلین

\end{otherlanguage}

Nāma-yi ān janāb bayna yaday al-ʿarsh ḥāḍir va bar maqarr-i innahu lā
yuʿrafu bi-mā siwāhu vāṣil va mā fīhi bi-liḥāẓ-i Allāh malḥūẓ āmad.
Suʾāl shuda būd az īnki chigūna mīshavad ki ḥarf-i ʿilliyyīn bi-sijjīn
tabdīl shavad va yā ithbāt bi-nafy rājiʿ gardad va yā thamara-yi janiyya
az liṭāfat mamnūʿ shavad va yā mirʾāt az ishrāqāt-i anvār-i āftāb-i
maʿānī maḥrūm mānad. Fa-niʿma mā saʾalta va kunta mina al-sāʾilīn.

Your letter was presented before the Throne and reached the seat of ``He
is not known through anything other than Himself,'' and its contents
came under the gaze of God. You had asked how a letter of the highest
heaven can be changed into one of the lowest abyss, or affirmation
revert to negation, or a ripe fruit be deprived of its sweetness, or a
mirror remain deprived of the radiance of the lights of the Sun of inner
meanings. Excellent indeed is what you have asked, and you are among
those who inquire.

\subsection{3}\label{section-1558}

\begin{otherlanguage}{persian}

بسیار سؤال شما مقبول افتاد چه که الیوم لازم است هر نفسی که از عرفان
معضلات مسائل الهیّه عاجز شود از شریعهٴ علم ربّانیّه و فرات حکمت صمدانیّه
سائل و آمل گردد که شاید برشحی از آن مشروب شود و بر بساط سکون و ایقان
مستریح گردد و نسأل اللّه بأن یصعدک الی مقام تشهد جمال القدم ببصرک و تنقطع
عن بصر العالمین و تسمع نغماته بأذنک و تنقطع عن اذن الخلایق اجمعین و تعرف
نفسه بقلبک و تنقطع عن افئدة کلّ من فی السّموات و الأرضین و یطهّرک عن دنس
الدّنیا و ما فیها بحیث لن تمرّ علی شیء الّا و قد تسمع منه بأنّه لا اله الّا
هو و انّ طلعة الأعلی لبهائه فی الملإ الأعلی و ضیائه بین الأرض و السّمآء و
کبریائه لمن فی ملکوت الأمر و الخلق و کذلک ینطق کلّ شیء ان انت من السّامعین

\end{otherlanguage}

Bisyār suʾāl-i shumā maqbūl uftād chi ki al-yawm lāzim ast har nafsī ki
az ʿirfān-i muʿḍalāt-i masāʾil-i ilāhiyya ʿājiz shavad az sharīʿa-yi
ʿilm-i rabbāniyya va furāt-i ḥikmat-i ṣamadāniyya sāʾil va āmil gardad
ki shāyad bi-rashḥī az ān mashrūb shavad va bar bisāṭ-i sukūn va īqān
mustarīḥ gardad. Va nasʾalu Allāha bi-an yuṣʿidaka ilā maqāmin tashhadu
jamāla al-qidami bi-baṣarika va tanqaṭiʿu ʿan baṣari al-ʿālamīn va
tasmaʿu naghamātihi bi-udhunika va tanqaṭiʿu ʿan udhuni al-khalāʾiqi
ajmaʿīn va taʿrifu nafsahu bi-qalbika va tanqaṭiʿu ʿan afʾidati kulli
man fī al-samāwāti va al-araḍīn va yuṭahhiraka ʿan danasi al-dunyā va mā
fīhā bi-ḥaythu lan tamurra ʿalā shayʾin illā va qad tasmaʿu minhu
bi-annahu lā ilāha illā huwa va inna ṭalʿata al-Aʿlā la-bahāʾuhu fī
al-malaʾi al-aʿlā va ḍiyāʾuhu bayna al-arḍi va al-samāʾ va kibriyāʾuhu
li-man fī malakūti al-amri va al-khalq. Va kadhālika yanṭiqu kullu
shayʾin in anta mina al-sāmiʿīn.

Your question was most welcome, for today every soul who is unable to
comprehend the difficulties of divine questions must seek and hope to
receive from the river of divine knowledge and the Euphrates of eternal
wisdom, that perhaps they may drink a sprinkling of it and find repose
upon the carpet of tranquillity and certitude. We ask God to raise you
to a station where you behold the Ancient Beauty with your own eyes and
detach yourself from the sight of the peoples of the world, hear His
melodies with your own ears and detach yourself from the hearing of all
created beings, and recognize His Self with your own heart and detach
yourself from the hearts of all who are in the heavens and on earth; and
to purify you from the defilement of the world and all that is in it, so
that you will pass by nothing without hearing it declare, ``There is no
God but Him, and the countenance of the Most Exalted is indeed His glory
in the Concourse on high, His radiance between earth and heaven, and His
majesty to those in the kingdom of command and creation.'' Thus does
everything speak, if you are among those who listen.

\subsection{4}\label{section-1559}

\begin{otherlanguage}{persian}

چه که الیوم بر هر نفسی من عند اللّه فرض شده که بچشم و گوش و فؤاد خود در
امر او ملاحظه نماید و تفکّر کند تا از بدایع مرحمت رحمن و فیوضات حضرت
سبحان باشراقات شمس معانی مستنیر و فایز شود و علّتی که جمیع ناس را از لقآء
اللّه محروم نموده و به ما سواه مشغول داشته اینست که بوهم صرف کفایت
نموده‌اند و بآنچه از امثال خود شنیده قناعت کرده‌اند براههای مهلک تقلید مشی
نموده‌اند و از مناهج تجرید محروم شده‌اند امر الهی لازال چون صبح نورانی
ظاهر و لایح بوده اینکه بعضی از ادراک او محتجب مانده‌اند نظر بآن است که
گوش و قلب را بآلایش کلمات ناس آلوده‌اند والّا اگر ناظر باصل میزان معرفت
الهی باشند هرگز از سبیل هدایت محروم نگردند

\end{otherlanguage}

Chi ki al-yawm bar har nafsī min ʿindi Allāh farḍ shuda ki bi-chashm va
gūsh va fuʾād-i khud dar amr-i ū mulāḥaẓa namāyad va tafakkur kunad tā
az badāyiʿ-i marḥamat-i Raḥmān va fuyūḍāt-i Ḥaḍrat-i Subḥān
bi-ishrāqāt-i shams-i maʿānī mustanīr va fāʾiz shavad. Va ʿillatī ki
jamīʿ-i nās rā az liqāʾ-i Allāh maḥrūm namūda va bi-mā siwāhu mashghūl
dāshta īnast ki bi-vahm-i ṣirf kifāyat namūda-and va bi-ānchi az
amthāl-i khud shanīda qanāʿat karda-and. Bi-rāhhā-yi muhlik-i taqlīd
mashy namūda-and va az manāhij-i tajrīd maḥrūm shuda-and. Amr-i ilāhī
lāzāl chūn ṣubḥ-i nūrānī ẓāhir va lāʾiḥ būda. Īnki baʿḍī az idrāk-i ū
muḥtajib mānda-and naẓar bi-ān ast ki gūsh va qalb rā bi-ālāyish-i
kalimāt-i nās ālūda-and va illā agar nāẓir bi-aṣl-i mīzān-i maʿrifat-i
ilāhī bāshand hargiz az sabīl-i hidāyat maḥrūm nagardand.

For today God has enjoined upon every soul to examine and reflect upon
His Cause with their own eyes, ears, and heart, so that, through the
wondrous mercy of the All-Merciful and the outpourings of the
All-Glorious, they may be illumined by the radiance of the Sun of inner
meanings and attain its grace. The reason all people have been deprived
of meeting God and occupied with anything other than Him is that they
have contented themselves with mere fancy and been satisfied with what
they have heard from others like themselves. They have walked the deadly
paths of imitation and been deprived of the ways of detachment. The
Cause of God has always been as manifest and radiant as the luminous
dawn. That some have remained veiled from perceiving it is because they
have defiled their ears and hearts with the pollution of people's words.
Otherwise, if they looked to the fundamental standard of the knowledge
of God, they would never be deprived of the path of guidance.

\subsection{5}\label{section-1560}

\begin{otherlanguage}{arabic}

حال خود آن جناب ملاحظه نماید امر از دو قسم بیرون نه یا آنکه اهل بیان
مقرّند بقدرت الهیّه یا نه اگر معترف نیستند باین نفوس حرفی نداریم چه که از
ملل قبل محسوبند که ید اللّه را مغلول دانسته‌اند چنانچه ربّ العزّة خبر داده
بقوله تعالی ید اللّه مغلولة و اگر معترفند ببدایع قدرت ربّانی در این صورت
اعتکاف باینگونه مسائل لغو بوده و باطل خواهد بود چه که عجز شأن خلق بوده و
آن ذات قدم لازال بر عرش قدرت و اقتدار مستوی و اگر اراده فرماید بحرفی
جمیع من علی الأرض را بسموات امر متصاعد فرماید و بحرف دیگر بادنی رتبهٴ
خلق راجع نماید و لیس لأحد ان یقول لم و بم و من قال فقد کفر باللّه و اعرض
عن قدرته و حارب بنفسه و نازع بسلطانه و کان من المشرکین فی الواح عزّ حفیظ

\end{otherlanguage}

Ḥāl khud-i ān janāb mulāḥaẓa namāyad amr az du qism bīrūn na: yā ānki
ahl-i Bayān muqirr-and bi-qudrat-i ilāhiyya yā na. Agar muʿtarif nīstand
bi-īn nufūs ḥarfī nadārīm chi ki az milal-i qabl maḥsūb-and ki yad-i
Allāh rā maghlūl dānista-and chunānchi Rabb al-ʿIzza khabar dāda
bi-qawlihi taʿālā: Yadu Allāhi maghlūla. Va agar muʿtarif-and
bi-badāyiʿ-i qudrat-i rabbānī dar īn ṣūrat iʿtikāf bi-īngūna masāʾil
laghv būda va bāṭil khwāhad būd chi ki ʿajz shaʾn-i khalq būda va ān
Dhāt-i Qidam lāzāl bar ʿarsh-i qudrat va iqtidār mustavī va agar irāda
farmāyad bi-ḥarfī jamīʿ-i man ʿalā al-arḍ rā bi-samāvāt-i amr mutaṣāʿid
farmāyad va bi-ḥarf-i dīgar bi-adnā rutba-yi khalq rājiʿ namāyad. Va
laysa li-aḥadin an yaqūla lima va bima va man qāla fa-qad kafara
bi-Allāhi va aʿraḍa ʿan qudratihi va ḥāraba bi-nafsihi va nāzaʿa
bi-sulṭānihi va kāna mina al-mushrikīn fī alwāḥi ʿizzin ḥafīẓ.

Now consider for yourself: the matter admits of only two possibilities.
Either the people of the Bayān acknowledge divine power, or they do not.
If they do not acknowledge it, We have nothing to say to such souls, for
they are reckoned among the religious communities of the past who
considered God's hand to be fettered, as the Lord of Glory has related
in His exalted words: ``God's hand is fettered.'' But if they
acknowledge the wonders of divine power, then dwelling on questions of
this kind is futile and will be of no avail. For powerlessness is a
condition of creation, while that Eternal Essence has always been
established upon the throne of power and might. If He so wills, with a
word He can raise all who are on earth to the heavens of His Cause, and
with another word return them to the lowest rank of creation. No one has
the right to say ``Why?'' or ``Whereby?'' Whoever says this has
disbelieved in God, turned away from His power, made war upon His Self,
and contested His sovereignty, and is numbered among those who associate
partners with God in the tablets of inviolable glory.

\subsection{6}\label{section-1561}

\begin{otherlanguage}{persian}

و هم‌چنین قادر است باینکه هر وقت اراده فرماید مظهر نفس خود را در بین بریّه
مبعوث نماید و در حین ظهور باید از نفس ظهور او تعالی حجّت و دلیل خواست اگر
بآن حجّتی که لازال مابین ناس بوده اتیان فرمود دیگر توقّف باطل است بلکه اگر
کلّ من فی السّموات و الأرض اقلّ من حین توقّف نمایند از اهل نار محسوب اگرچه
جمیع ادّعای ولایت نمایند عزیز میکند هر که را اراده فرماید و ذلیل میفرماید
هر که را بخواهد لا یسأل عمّا یفعل

\end{otherlanguage}

Va hamchunīn qādir ast bi-īnki har vaqt irāda farmāyad maẓhar-i nafs-i
khud rā dar bayn-i bariyya mabʿūth namāyad va dar ḥīn-i ẓuhūr bāyad az
nafs-i ẓuhūr-i ū taʿālā ḥujjat va dalīl khwāst. Agar bi-ān ḥujjatī ki
lāzāl mābayn-i nās būda ityān farmūd dīgar tavaqquf bāṭil ast balki agar
kull-i man fī al-samāvāt va al-arḍ aqall-i min ḥīn tavaqquf namāyand az
ahl-i nār maḥsūb agarchi jamīʿ iddiʿā-yi vilāyat namāyand. ʿAzīz mīkunad
har ki rā irāda farmāyad va dhalīl mīfarmāyad har ki rā bikhwāhad. Lā
yusʾalu ʿammā yafʿal.

Likewise, He has the power to raise up the Manifestation of His own Self
among creation whenever He wills. At the time of His appearance, proof
and evidence must be sought from His exalted Manifestation Himself. If
He produces that proof which has always existed among people, further
hesitation is unwarranted. Indeed, if all who are in the heavens and on
earth hesitate for less than an instant, they are reckoned among the
people of the fire, even if they all claim sainthood. He exalts whomever
He wills and abases whomever He wishes. He is not questioned about what
He does.

\subsection{7}\label{section-1562}

\begin{otherlanguage}{persian}

حال آن جناب بانصاف ملاحظه نمایند اگر نفسی بجمیع شئونات قدرتیّهٴ الهیّه
ظاهر شود و علاوه بر آن اتیان نماید بآن حجّتی که لازال بآن اثبات دین ناس
شده و امر اللّه بین بریّهٴ او ثابت گشته و مع‌ذلک از چنین ظهور که نفس ظهور
اللّه بوده نفسی اعراض نماید و باعراض هم کفایت ننموده بر قتلش قیام نماید
آیا بر چنین نفسی چه حکم جاری قل حکمه عند اللّه یحکم ما یشآء کما حکم بالحقّ
ولکنّ النّاس هم لا یشعرون آیا در هیچ عهدی و ملّتی چنین امری جایز بوده لا
فونفسه الرّحمن الرّحیم

\end{otherlanguage}

Ḥāl ān janāb bi-inṣāf mulāḥaẓa namāyand agar nafsī bi-jamīʿ-i shuʾūnāt-i
qudratiyya-yi ilāhiyya ẓāhir shavad va ʿalāva bar ān ityān namāyad bi-ān
ḥujjatī ki lāzāl bi-ān ithbāt-i dīn-i nās shuda va amr-i Allāh bayn-i
bariyya-yi ū thābit gashta va maʿa dhālik az chunīn ẓuhūr ki nafs-i
ẓuhūr-i Allāh būda nafsī iʿrāḍ namāyad va bi-iʿrāḍ ham kifāyat nanamūda
bar qatl-ash qiyām namāyad āyā bar chunīn nafsī chi ḥukm jārī? Qul
ḥukmuhu ʿinda Allāh yaḥkumu mā yashāʾ kamā ḥakama bi-al-ḥaqq va lākinna
al-nāsa hum lā yashʿurūn. Āyā dar hīch ʿahdī va millatī chunīn amrī
jāʾiz būda? Lā fa-wa-nafsihi al-Raḥmāni al-Raḥīm.

Now consider fairly: if a person appears invested with all the
expressions of divine power and, moreover, produces that proof by which
people's religion has always been substantiated and the Cause of God
established among His creatures, and yet someone turns away from such a
manifestation, which is the very Manifestation of God, and, not content
with turning away, rises up to kill Him, what judgment applies to such a
person? Say: His judgment rests with God. He decrees whatever He wills,
just as He has decreed in truth, but the people do not perceive. Has
such a deed ever been permissible in any age or religious community? No,
by His own Self, the All-Merciful, the Compassionate!

\subsection{8}\label{section-1563}

\begin{otherlanguage}{persian}

اگر بگویند این ظهور نباید ظاهر شود چنانچه مشرکین گفته‌اند در این صورت
قدرت و ارادهٴ حقّ منوط و معلّق بارادهٴ خلق میشود فتعالی عن ذلک علوّاً کبیراً
چنانچه در ظهور ستّین کلّ ناس از عالم و جاهل باین سخنهای مزخرف بی‌معنی از حقّ
محروم شده‌اند و از علّیّین بقا بسجّین فنا راجع گشته‌اند و بگمان خود بر اعلی
مقعد ایمان مستقرّند فبئس ما ظنّوا فی انفسهم و کانوا من المتوهّمین فی امّ
الألواح مذکورا

\end{otherlanguage}

Agar bigūyand īn ẓuhūr nabāyad ẓāhir shavad chunānchi mushrikīn
gufta-and dar īn ṣūrat qudrat va irāda-yi Ḥaqq manūṭ va muʿallaq
bi-irāda-yi khalq mīshavad. Fa-taʿālā ʿan dhālika ʿuluwwan kabīrā.
Chunānchi dar ẓuhūr-i sittīn kull-i nās az ʿālim va jāhil bi-īn
sukhanhā-yi muzakhraf-i bī-maʿnī az Ḥaqq maḥrūm shuda-and va az
ʿilliyyīn-i baqā bi-sijjīn-i fanā rājiʿ gashta-and va bi-gumān-i khud
bar aʿlā maqʿad-i īmān mustaqirr-and. Fa-biʾsa mā ẓannū fī anfusihim va
kānū mina al-mutawahhimīn fī Umm al-Alwāḥ madhkūrā.

If they say that this Manifestation must not appear, as those who
associate partners with God have said, then God's power and will become
dependent upon and subject to the will of His creatures. Exalted,
immeasurably exalted, is He above that! Thus, at the Revelation of the
year sixty, all people, learned and ignorant alike, were deprived of God
through such meaningless, spurious words, and returned from the highest
heaven of eternity to the lowest abyss of extinction, while imagining
themselves established upon the loftiest seat of faith. How evil is what
they have imagined concerning themselves! In the Mother of Tablets they
are recorded among those who are deluded.

\subsection{9}\label{section-1564}

\begin{otherlanguage}{arabic}

باری نظر را از ما سوی اللّه بردار و بحقّ ناظر شو و به ما یظهر من عنده چه
که دون او لاشیء محض بوده و خواهد بود و اگر الیوم کلّ من فی السّموات و
الأرض حروفات بیانیّه شوند که بصدهزار رتبه از حروفات فرقانیّه اعظم و اکبرند
و اقلّ من آن در این امر توقّف نمایند از معرضین عنداللّه محسوبند و از احرف
نفی منسوب

\end{otherlanguage}

Bārī naẓar rā az mā sivā-yi Allāh bardār va bi-Ḥaqq nāẓir shaw va bi-mā
yaẓharu min ʿindihi chi ki dūn-i ū lā-shayʾ-i maḥḍ būda va khwāhad būd
va agar al-yawm kull-i man fī al-samāvāt va al-arḍ ḥurūfāt-i bayāniyya
shavand ki bi-ṣad hazār rutba az ḥurūfāt-i furqāniyya aʿẓam va akbar-and
va aqall-i min ān dar īn amr tavaqquf namāyand az muʿriḍīn ʿinda Allāh
maḥsūb-and va az aḥruf-i nafy mansūb.

In short, turn your gaze away from everything but God and look to Him
and to whatever appears from Him, for all besides Him has been and will
remain sheer nothingness. If today all who are in the heavens and on
earth were to become letters of the Bayān, who are a hundred thousand
degrees greater and mightier than the letters of the Qurʾān, and were to
hesitate in this Cause for less than an instant, they would be reckoned
before God among those who turn away and numbered among the letters of
negation.

\subsection{10}\label{section-1565}

\begin{otherlanguage}{persian}

حقّ جلّ و عزّ را باحدی نسبت و ربط و مشابهت و مشاکلت نه و کلّ به نسبتهم الی
عرفانه مفتخر و معزّز بوده و خواهند بود جمال سبحان بر عرش رحمن مستوی و
پرتو انوار شمس فضلش بر کلّ اشیاء بالسّویّه اشراق و تجلّی فرموده و جمیع من فی
الملک بین یدی الفضل در صقع واحد قائمند و ذرّه را بر ذرّه افتخار و زیادتی
نه الّا بسبقتها الی عرفان اللّه و لقائه فطوبی لمن عرفه بنفسه و انقطع عمّا
سواه

\end{otherlanguage}

Ḥaqq jalla va ʿazza rā bi-aḥadī nisbat va rabṭ va mushābahat va
mushākalat na va kull bi-nisbatihim ilā ʿirfānihi muftakhir va muʿazzaz
būda va khwāhand būd. Jamāl-i Subḥān bar ʿarsh-i Raḥmān mustavī va
partaw-i anvār-i shams-i faḍl-ash bar kull-i ashyāʾ bi-al-saviyya ishrāq
va tajallī farmūda va jamīʿ-i man fī al-mulk bayna yaday al-faḍl dar
ṣuqʿ-i vāḥid qāʾim-and va dharra rā bar dharra iftikhār va ziyādatī na
illā bi-sabqatihā ilā ʿirfāni Allāhi va liqāʾihi. Fa-ṭūbā li-man
ʿarafahu bi-nafsihi va inqaṭaʿa ʿammā siwāhu.

God, glorified and exalted is He, bears no relation, connection,
resemblance, or likeness to anyone, and all have been and will be
honored and exalted according to their relation to the recognition of
Him. The All-Glorious Beauty is established upon the throne of the
All-Merciful, and the rays of the light of the sun of His grace have
shone and appeared equally upon all things. All who dwell in the realm
of existence stand on a single plane before His grace, and no atom has
any distinction or superiority over another except through its
precedence in recognizing God and meeting Him. Blessed, then, is the one
who has recognized Him through Himself and detached himself from
everything but Him.

\subsection{11}\label{section-1566}

\begin{otherlanguage}{persian}

ای علیّ بشنو ندآء اللّه را و بمقرّی وارد شو که لازال مقدّس از اسماء بوده و
خواهد بود تا بهیچ اسمی از جمال مسمّی و سلطانی که بارادهٴ قلمش ملکوت اسماء
خلق شده محروم نگردی فواللّه الّذی لا اله الّا هو که مقصود از این بیان آن
است که شاید آن جناب و معدودی خرق حجاب نموده بسرادق قدس محبوب که مقدّس از
ظنون و اوهام عباد بوده درآیند والّا انّه لمتعالی عن اقبال الخلق و اعراضهم
و مقدّس عن العالمین

\end{otherlanguage}

Ay ʿAlī bishnaw nidā-yi Allāh rā va bi-maqarrī vārid shaw ki lāzāl
muqaddas az asmāʾ būda va khwāhad būd tā bi-hīch ismī az jamāl-i musammā
va sulṭānī ki bi-irāda-yi qalam-ash malakūt-i asmāʾ khalq shuda maḥrūm
nagardī. Fa-wa-Allāhi alladhī lā ilāha illā huwa ki maqṣūd az īn bayān
ān ast ki shāyad ān janāb va maʿdūdī kharq-i ḥijāb namūda bi-surādiq-i
quds-i Maḥbūb ki muqaddas az ẓunūn va awhām-i ʿibād būda darāyand va
illā innahu la-mutaʿālin ʿan iqbāli al-khalqi va iʿrāḍihim va muqaddasun
ʿan al-ʿālamīn.

O ʿAlī, hear the call of God and enter a realm that has always been and
will forever remain sanctified above names, so that no name may deprive
you of the beauty of the One named, the Sovereign through the will of
whose Pen the kingdom of names was created. By God, besides whom there
is no other god, the purpose of these words is that perhaps you and a
few others may rend the veil and enter the tabernacle of the Beloved's
holiness, which is sanctified above the suppositions and imaginings of
His servants. Otherwise, He is indeed exalted above the turning of His
creatures toward Him or away from Him, and sanctified above all the
worlds.

\subsection{12}\label{section-1567}

\begin{otherlanguage}{persian}

آیا در حین اشراق شمس لایق است نفسی سؤال نماید که چگونه میشود نور انجم
اخذ شود و حال آنکه ملاحظه مینماید که نور آفتاب روشنی او را معدوم نموده
بلکه در این مقام نجوم طالب ظلمت لیلند و از نور نهار معرض چه که قدر و
ضیاء نجوم در لیل مشهود است و از تجلّی نیّر یوم معدوم و مفقود میگردند
فسبحانه عن المثل و الأمثال چه که لازال نیّر جمالش مستضیء بوده و احدی با
او نبوده و کلّ ما سواه در امکنهٴ ترابیّه بمشیّت امکانیّه خلق شده‌اند و باو
راجع خواهند شد و انّه جلّ و عزّ در مقعد امتناع و مقرّ ارتفاع خود لم‌یزل و
لایزال مقدّس از کلّ بوده و خواهد بود بسیار عجیب است که از تغییر و تبدیل
اسماء ناس تعجّب مینمایند و متحیّر شده‌اند با اینکه جمیع در کلّ حین تغییر و
تبدیل مظاهر اسماء و مطالع آن را ببصر ظاهر مشاهده مینمایند و مع‌ذلک بحجبات
وهمیّه و کلمات شرکیّه چنان محتجب مانده‌اند که از آنچه ببصر ملاحظه مینمایند
غافل شده‌اند

\end{otherlanguage}

Āyā dar ḥīn-i ishrāq-i shams lāyiq ast nafsī suʾāl namāyad ki chigūna
mīshavad nūr-i anjum akhdh shavad va ḥāl ānki mulāḥaẓa mīnamāyad ki
nūr-i āftāb rawshanī-yi ū rā maʿdūm namūda balki dar īn maqām nujūm
ṭālib-i ẓulmat-i layl-and va az nūr-i nahār muʿriḍ chi ki qadr va
ḍiyā-yi nujūm dar layl mashhūd ast va az tajallī-yi nayyir-i yawm maʿdūm
va mafqūd mīgardand. Fa-subḥānahu ʿan al-mathali va al-amthāl chi ki
lāzāl nayyir-i jamāl-ash mustaḍīʾ būda va aḥadī bā ū nabūda va kull-i mā
sivāhu dar amkina-yi turābiyya bi-mashiyyat-i imkāniyya khalq shuda-and
va bi-ū rājiʿ khwāhand shud va innahu jalla va ʿazza dar maqʿad-i
imtināʿ va maqarr-i irtifāʿ-i khud lam-yazal va lā-yazāl muqaddas az
kull būda va khwāhad būd. Bisyār ʿajīb ast ki az taghyīr va tabdīl-i
asmāʾ-i nās taʿajjub mīnamāyand va mutaḥayyir shuda-and bā īnki jamīʿ
dar kull-i ḥīn taghyīr va tabdīl-i maẓāhir-i asmāʾ va maṭāliʿ-i ān rā
bi-baṣar-i ẓāhir mushāhada mīnamāyand va maʿa dhālik bi-ḥujubāt-i
vahmiyya va kalimāt-i shirkiyya chunān muḥtajib mānda-and ki az ānchi
bi-baṣar mulāḥaẓa mīnamāyand ghāfil shuda-and.

When the sun shines, is it fitting for anyone to ask how the light of
the stars can be taken away, while observing that the sunlight has
effaced their brightness? Indeed, in this condition the stars seek the
darkness of night and turn away from the light of day, for the worth and
radiance of the stars are seen at night, and through the shining of the
orb of day they vanish and disappear. Glorified is He above comparison
and likenesses! For the orb of His beauty has always been radiant, and
no one has been with Him. All besides Him have been created in earthly
abodes through the will pertaining to contingent existence, and to Him
they will return. He, glorified and exalted is He, in His inaccessible
seat and His lofty abode, has ever been and will forever remain
sanctified above all. It is most strange that people are astonished and
bewildered by the change and alteration of names, although at every
moment they all observe with their outward eyes the change and
alteration of the manifestations of names and their dawning-places. Yet
they have remained so veiled by imaginary veils and words that associate
partners with God that they have become heedless of what they see with
their own eyes.

\subsection{13}\label{section-1568}

\begin{otherlanguage}{persian}

ای سائل اسماء و صفات الهی را موهوم مدان بدان که جمیع اشیاء از آنچه مابین
ارض و سماء خلق شده مظاهر اسماء و مطالع صفات حقّ تعالی شأنه بوده و خواهد
بود غایت این است که انسان نسبت بدون خود اعظم رتبةً و اکبر مقاماً خلق شده و
اگر در سماء ما تری فی خلق الرّحمن من تفاوت ارتقا نمائی در خلق رحمانی و
مطالع صنع سبحانی تفاوت و فطور نبینی فطوبی لمن طار فی هذا الهوآء الّذی ما
طارت فیه اجنحة المریبین

\end{otherlanguage}

Ay sāʾil asmāʾ va ṣifāt-i ilāhī rā mawhūm madān. Bidān ki jamīʿ-i ashyāʾ
az ānchi mābayn-i arḍ va samāʾ khalq shuda maẓāhir-i asmāʾ va maṭāliʿ-i
ṣifāt-i Ḥaqq taʿālā shaʾnuhu būda va khwāhad būd. Ghāyat īn ast ki insān
nisbat bi-dūn-i khud aʿẓam rutbatan va akbar maqāman khalq shuda va agar
dar samā-yi mā tarā fī khalqi al-Raḥmāni min tafāwut irtiqā namāʾī dar
khalq-i raḥmānī va maṭāliʿ-i ṣunʿ-i subḥānī tafāvut va fuṭūr nabīnī.
Fa-ṭūbā li-man ṭāra fī hādhā al-hawāʾ alladhī mā ṭārat fīhi ajniḥatu
al-murībīn.

O questioner, do not regard the divine names and attributes as
imaginary. Know that all things created between earth and heaven have
been and will remain manifestations of the names and dawning-places of
the attributes of God, exalted is His station. It is only that the human
being has been created greater in rank and loftier in station than all
else. If you ascend into the heaven of ``You see no disparity in the
creation of the All-Merciful,'' you will see neither disparity nor flaw
in the creation of the All-Merciful and the dawning-places of His
glorious handiwork. Blessed, then, is the one who has soared in this
atmosphere, where the wings of the doubters have never flown.

\subsection{14}\label{section-1569}

\begin{otherlanguage}{persian}

و حال ملاحظه نما در جمیع این مظاهر اسمیّهٴ الهیّه از اشجار و افنان و اغصان
و اثمار و هم‌چنین در اوراد و ازهار و کلّما ینبت علی وجه الأرض که در اوّل
بچه مقدار طراوت و لطافت و نضارت ظاهر میشوند و بعد از مدّتی کلّ از خلع
لطیفه عاری شده بارض راجع شوند چه مقدار از ثمرات جنیّه که تغییر نماید
بشأنی که از رایحهٴ او انسان اجتناب نماید باری کلّ در علوّ و دنوّ و تغییر و
تبدیل مگر مظاهر کلّیّهٴ الهیّه که بنفسه لنفسه قائم و باقیند

\end{otherlanguage}

Va ḥāl mulāḥaẓa namā dar jamīʿ-i īn maẓāhir-i ismiyya-yi ilāhiyya az
ashjār va afnān va aghṣān va athmār va ham-chunīn dar avrād va azhār va
kulli mā yanbutu ʿalā wajhi al-arḍ ki dar avval bi-chi miqdār ṭarāvat va
laṭāfat va niḍārat ẓāhir mī-shavand va baʿd az muddatī kull az khilaʿ-i
laṭīfa ʿārī shuda bi-arḍ rājiʿ shavand. Chi miqdār az thamarāt-i janiyya
ki taghyīr namāyad bi-shaʾnī ki az rāʾiḥa-yi ū insān ijtināb namāyad.
Bārī kull dar ʿuluvv va dunuvv va taghyīr va tabdīl magar maẓāhir-i
kulliyya-yi ilāhiyya ki bi-nafsihi li-nafsihi qāʾim va bāqī-and.

Now consider all these manifestations of the divine names---trees,
boughs, branches, and fruits, as well as roses, blossoms, and everything
that grows upon the face of the earth---with what freshness, delicacy,
and verdure they first appear, and how, after a time, they are all
stripped of their delicate robes and return to the earth. How many ripe
fruits undergo such a change that people recoil from their smell! In
short, all undergo ascent and decline, change and transformation, except
the universal Manifestations of God, who subsist and endure through
themselves and for themselves.

\subsection{15}\label{section-1570}

\begin{otherlanguage}{persian}

قسم بآفتاب صبح معانی که لسان الهی بشأن و اندازهٴ ناس تکلّم میفرماید چه که
اکثری از ناس ببلوغ نرسیده‌اند والّا بابی از علم بر وجه عباد مفتوح میفرمود
که کلّ من فی السّموات و الأرض بافاضهٴ قلمیّهٴ او از علم ماسوی خود را غنیّ
مشاهده نموده بر اعراش سکون مستقرّ میشدند و نظر بعدم استعداد ناس جواهر علم
ربّانی و اسرار حکمت صمدانی در سماء مشیّت الهی محفوظ و مستور مانده و تا حین
حرفی از آن نازل نشده و بعد الأمر بیده یفعل ما یشآء و لا یسأل عمّا شآء و
هو العلیم الخبیر

\end{otherlanguage}

Qasam bi-āftāb-i ṣubḥ-i maʿānī ki lisān-i ilāhī bi-shaʾn va andāza-yi
nās takallum mī-farmāyad, chi ki aktharī az nās bi-bulūgh na-rasīda-and,
va illā bābī az ʿilm bar vajh-i ʿibād maftūḥ mī-farmūd ki kullu man fī
al-samāwāti wa-al-arḍ bi-ifāḍa-yi qalamiyya-yi ū az ʿilm-i mā-sivā khud
rā ghanī mushāhada namūda bar aʿrāsh-i sukūn mustaqarr mī-shudand. Va
naẓar bi-ʿadam-i istiʿdād-i nās javāhir-i ʿilm-i rabbānī va asrār-i
ḥikmat-i ṣamadānī dar samā-yi mashiyyat-i ilāhī maḥfūẓ va mastūr mānda
va tā ḥīn ḥarfī az ān nāzil na-shuda va baʿd al-amru bi-yadihi yafʿalu
mā yashāʾu wa lā yusʾalu ʿammā shāʾa wa huwa al-ʿAlīmu al-Khabīr.

I swear by the Sun of the morning of inner meanings that the divine
Tongue speaks in keeping with the condition and capacity of the people,
for most people have not reached maturity. Otherwise, He would open a
door of knowledge before His servants such that all who are in the
heavens and on earth, through the outpouring of His Pen, would find
themselves independent of the knowledge of all besides Him and would be
established upon thrones of tranquility. In view of the people's lack of
readiness, the gems of divine knowledge and the mysteries of eternal
wisdom have remained guarded and concealed in the heaven of God's will,
and until now not a word of them has been revealed. Hereafter, the
matter rests in His hand: He does whatever He wills and is not
questioned about what He has willed. He is the All-Knowing, the
All-Aware.

\subsection{16}\label{section-1571}

\begin{otherlanguage}{persian}

و اگر از این عبد میشنوی پرهای تحدید و تقلید را بیفکن و بپرهای تجرید در
این هوای قدس توحید پرواز کن تا از شبهات وهمیّه و اشارات ریبیّه خود را مقدّس
یابی و بانوار یقین خود را منوّر بینی بگوش جان کلمات رحمانی را اصغا نما که
شاید قلب از غبار مکدّرهٴ اوهام که در این ایّام کلّ من فی الملک را احاطه
نموده مطهّر شده بمنظر اکبر راجع شود

\end{otherlanguage}

Va agar az īn ʿabd mī-shinavī par-hā-yi taḥdīd va taqlīd rā biyafkan va
bi-par-hā-yi tajrīd dar īn havā-yi quds-i tawḥīd parvāz kun tā az
shubuhāt-i vahmiyya va ishārāt-i raybiyya khud rā muqaddas yābī va
bi-anvār-i yaqīn khud rā munavvar bīnī. Bi-gūsh-i jān kalimāt-i raḥmānī
rā iṣghā namā ki shāyad qalb az ghubār-i mukaddira-yi avhām ki dar īn
ayyām kulla man fī al-mulk rā iḥāṭa namūda muṭahhar shuda bi-manẓar-i
akbar rājiʿ shavad.

If you will listen to this servant, cast away the wings of limitation
and imitation, and soar on the wings of detachment in this holy
atmosphere of divine unity, so that you may find yourself sanctified
from fanciful doubts and dubious allusions, and see yourself illumined
by the lights of certitude. Listen with the ear of your soul to the
words of the All-Merciful, that perhaps the heart may be cleansed of the
obscuring dust of imaginings that has enveloped all who dwell in the
world in these days, and may turn toward the Most Great Vision.

\subsection{17}\label{section-1572}

\begin{otherlanguage}{persian}

و چون باین مقام اقدس امنع اطهر فایز شوی ملاحظه مینمائی که مقصود ما فی
البیان که از سماء مشیّت ظهور قبلم نازل شده این ظهور بوده و خواهد بود
فوالّذی نفس حسین فی قبضة قدرته که اعظم از این امر در بیان نازل نشده

\end{otherlanguage}

Va chūn bi-īn maqām-i aqdas-i amnaʿ-i aṭhar fāʾiz shavī mulāḥaẓa
mī-namāʾī ki maqṣūd-i mā fī al-Bayān ki az samā-yi mashiyyat-i ẓuhūr-i
qablam nāzil shuda īn ẓuhūr būda va khvāhad būd. Fa-wa-alladhī nafsu
Ḥusaynin fī qabḍati qudratihi ki aʿẓam az īn amr dar Bayān nāzil
na-shuda.

When you attain this most holy, most inaccessible, most pure station,
you will perceive that the purpose of all that was revealed in the Bayān
from the heaven of the will of My previous Manifestation has been and
will remain this Revelation. By Him in the grasp of whose power is the
soul of Ḥusayn, nothing greater than this Cause has been revealed in the
Bayān.

\subsection{18}\label{section-1573}

\begin{otherlanguage}{persian}

ببصر منیر حدید در او ملاحظه نمائید تا بر مقصود کلمات قدس ربّانی مطّلع شوید
و در جمیع اوراق و الواح و رقاع و صحف و زبر و کتب کلّ ناس را وصیّت فرموده و
از جمیع اخذ عهد نموده که مباد در حین ظهور بشیئی از آنچه خلق شده متمسّک
شوند و از نفس ظهور محتجب مانند چه که در آن یوم هیچ شیئی نفع نمی‌بخشد الّا
بعد اذنه

\end{otherlanguage}

Bi-baṣar-i munīr-i ḥadīd dar ū mulāḥaẓa namāʾīd tā bar maqṣūd-i
kalimāt-i quds-i rabbānī muṭṭaliʿ shavīd va dar jamīʿ-i avrāq va alvāḥ
va riqāʿ va ṣuḥuf va zubur va kutub kull-i nās rā vaṣiyyat farmūda va az
jamīʿ akhdh-i ʿahd namūda ki mabād dar ḥīn-i ẓuhūr bi-shayʾī az ānchi
khalq shuda mutamassik shavand va az nafs-i ẓuhūr muḥtajib mānand, chi
ki dar ān yawm hīch shayʾī nafʿ na-mī-bakhshad illā baʿda idhnihi.

Examine it with clear and penetrating vision, so that you may understand
the intent of the holy words of the Lord. In all His pages, tablets,
missives, scriptures, writings, and books, He has counseled all people
and taken a covenant from them all, lest at the time of the Revelation
they cling to anything that has been created and remain veiled from the
Manifestation Himself. For on that Day nothing is of any avail except by
His permission.

\subsection{19}\label{section-1574}

\begin{otherlanguage}{arabic}

به وحید اکبر میفرمایند و انتظروا من یذکرکم اللّه وجهه فانّکم ما خلقتم الّا
للقائه و هو الّذی علّق کلّ شیء بأمره ایّاک ایّاک ایّام ظهوره ان تحتجب بالواحد
البیانیّة فانّ ذلک الواحد خلق عنده و ایّاک ایّاک ان تحتجب بکلمات ما نزّلت فی
البیان فانّها کلمات نفسه فی هیکل ظهوره من قبل

\end{otherlanguage}

Bi-Vaḥīd-i Akbar mī-farmāyand: wa-intaẓirū man yudhakkirukumu Allāhu
wajhahu fa-innakum mā khuliqtum illā li-liqāʾihi wa huwa alladhī ʿallaqa
kulla shayʾin bi-amrihi. Iyyāka iyyāka ayyāma ẓuhūrihi an taḥtajiba
bi-al-wāḥidi al-bayāniyya fa-inna dhālika al-wāḥida khuliqa ʿindahu wa
iyyāka iyyāka an taḥtajiba bi-kalimāti mā nuzzilat fī al-Bayān fa-innahā
kalimātu nafsihi fī haykali ẓuhūrihi min qabl.

He says to Vaḥīd-i Akbar: ``Await Him whose countenance God will call to
your remembrance, for you were created only to meet Him. He is the One
who has made all things dependent upon His command. Beware, beware lest
in the days of His appearance you allow the Vāḥid of the Bayān to veil
you, for that Vāḥid is a creation in His sight. Beware, beware lest you
be veiled by the words revealed in the Bayān, for they are His own words
in the temple of His previous Manifestation.''

\subsection{20}\label{section-1575}

\begin{otherlanguage}{arabic}

حال در این کلمات سلطان اسماء و صفات تفکّر نمائید با این آیات محکمه و
کلمات متقنه دیگر مجال اعراض برای نفسی باقی میماند لا فوالّذی انطق الرّوح
فی صدری مگر آنکه بالمرّه از حقّ اعراض نماید و واحد بیان خلق اوّلیّه‌اند و ما
دونهم خلقوا فی ظلّهم

\end{otherlanguage}

Ḥāl dar īn kalimāt-i sulṭān-i asmā va ṣifāt tafakkur namāʾīd. Bā īn
āyāt-i muḥkama va kalimāt-i mutqana dīgar majāl-i iʿrāḍ barā-yi nafsī
bāqī mī-mānad? Lā fa-wa-alladhī anṭaqa al-Rūḥa fī ṣadrī, magar ānki
bi-al-marra az Ḥaqq iʿrāḍ namāyad. Va Vāḥid-i Bayān khalq-i
avvaliyya-and wa mā dūnahum khuliqū fī ẓillihim.

Now reflect on these words of the Sovereign of names and attributes.
With these firm verses and conclusive words, does any room remain for
anyone to turn away? No, by Him who has caused the Spirit to speak
within My breast---unless one should turn away from God altogether. The
Vāḥid of the Bayān are the first creation, and all besides them were
created beneath their shadow.

\subsection{21}\label{section-1576}

\begin{otherlanguage}{arabic}

و هم‌چنین به وحید اکبر میفرمایند و ربّما یأتیک من انت قد سألت عن علوّ ذکره
و ارتفاع امره و انّ من فی البیان یقرؤون تلک الکلمات و هم لا یلتفتون
بظهوره و لا یؤمنون باللّه الّذی خلقهم بظهور قبله و هم راقدون

\end{otherlanguage}

Va ham-chunīn bi-Vaḥīd-i Akbar mī-farmāyand: wa rubbamā yaʾtīka man anta
qad saʾalta ʿan ʿuluwwi dhikrihi wa irtifāʿi amrihi wa inna man fī
al-Bayān yaqraʾūna tilka al-kalimāti wa hum lā yaltafitūna bi-ẓuhūrihi
wa lā yuʾminūna bi-Allāhi alladhī khalaqahum bi-ẓuhūrin qablihi wa hum
rāqidūn.

He likewise says to Vaḥīd-i Akbar: ``Perhaps He will come to you about
the loftiness of whose mention and the exaltation of whose Cause you
have asked. Those who are in the Bayān read these words, yet they pay no
heed to His appearance, nor do they believe in God, who created them
through a Revelation before Him, and they lie asleep.''

\subsection{22}\label{section-1577}

\begin{otherlanguage}{persian}

چنانچه حال ملاحظه میشود که کلّ تلاوت کتاب اللّه مینمایند و در لیل و نهار
مینویسند و مع‌ذلک بحرفی از کتاب مستشعر نشده‌اند بلکه مقصود از ترویج جز
تحقّق ریاست و اثبات آن نبوده و نخواهد بود کذلک یشهد لسان اللّه الملک
العزیز العلیم

\end{otherlanguage}

Chunānchi ḥāl mulāḥaẓa mī-shavad ki kull tilāvat-i Kitāb Allāh
mī-namāyand va dar layl va nahār mī-nivīsand va maʿa-dhālik bi-ḥarfī az
Kitāb mustashʿir na-shuda-and balki maqṣūd az tarvīj juz taḥaqquq-i
riyāsat va ithbāt-i ān na-būda va na-khvāhad būd. Kadhālika yashhadu
lisānu Allāhi al-Maliki al-ʿAzīzi al-ʿAlīm.

Thus it is now seen that all recite the Book of God and write by night
and day, yet they have not understood a single word of the Book. Rather,
their purpose in promoting it has been and will remain nothing but to
secure and establish their own leadership. Thus testifies the Tongue of
God, the King, the Almighty, the All-Knowing.

\subsection{23}\label{section-1578}

\begin{otherlanguage}{arabic}

و میفرمایند من اوّل ذلک الأمر الی قبل ان یکمل تسعة کینونات الخلق لم تظهر
و انّ کلّما قد رأیت من النّطفة الی ما کسوناه لحماً ثمّ اصبر حتّی تشهد خلق
الآخر قل فتبارک اللّه احسن الخالقین

\end{otherlanguage}

Va mī-farmāyand: min awwali dhālika al-amri ilā qabli an yakmula
tisʿatun kaynūnātu al-khalqi lam taẓhar wa inna kulla mā qad raʾayta min
al-nuṭfati ilā mā kasawnāhu laḥman thumma iṣbir ḥattā tashhada khalqa
al-ākhar. Qul fa-tabāraka Allāhu aḥsanu al-khāliqīn.

He also says: ``From the beginning of that Cause until before nine is
completed, the realities of creation have not appeared. All that you
have seen is from the stage of the drop of seed to that which We clothed
with flesh. Then be patient until you witness the other creation. Say:
`Blessed, then, is God, the best of creators!'\,''

\subsection{24}\label{section-1579}

\begin{otherlanguage}{arabic}

و هم‌چنین به عظیم میفرمایند هذا ما وعدناک قبل الّذی اجبناک اصبر حتّی یقضی
عن البیان تسعة فاذاً قل فتبارک اللّه احسن المبدعین

\end{otherlanguage}

Va ham-chunīn bi-ʿAẓīm mī-farmāyand: hādhā mā waʿadnāka qabla alladhī
ajabnāka iṣbir ḥattā yuqḍā ʿani al-Bayāni tisʿatun fa-idhan qul
fa-tabāraka Allāhu aḥsanu al-mubdiʿīn.

He likewise says to ʿAẓīm: ``This is what We promised you before We
answered you. Be patient until nine has elapsed from the Bayān. Then
say: `Blessed, then, is God, the best of originators!'\,''

\subsection{25}\label{section-1580}

\begin{otherlanguage}{arabic}

و میفرمایند فلتراقبنّ فرق القائم و القیّوم ثمّ فی سنة التّسع کلّ خیر تدرکون

\end{otherlanguage}

Va mī-farmāyand: fa-la-turāqibanna farqa al-Qāʾimi wa-al-Qayyūmi thumma
fī sanati al-tisʿi kulla khayrin tudrikūn.

He also says: ``Be sure to observe the difference between the Qāʾim and
the Qayyūm. Then, in the year nine, you will attain all good.''

\subsection{26}\label{section-1581}

\begin{otherlanguage}{arabic}

حال قدری در این کلمات تفکّر فرمائید و هم‌چنین در فرق قائم و قیّوم تفکّر لازم
چه که این عبد از احزان وارده قادر بر تفسیر کلمات اللّه نه فواللّه الّذی لا
اله الّا هو که تبلیغ امر اللّه این عبد را بر تحریر این لوح مضطرّ نموده که
شاید مصباحی چند در مشکاة امکان مستضیء شوند و بقدرة اللّه بر نصر این مظلوم
فرید قائم گردند

\end{otherlanguage}

Ḥāl qadrī dar īn kalimāt tafakkur farmāʾīd va ham-chunīn dar farq-i
Qāʾim va Qayyūm tafakkur lāzim, chi ki īn ʿabd az aḥzān-i vārida qādir
bar tafsīr-i kalimāt Allāh na. Fa-wa-Allāhi alladhī lā ilāha illā huwa
ki tablīgh-i Amr Allāh īn ʿabd rā bar taḥrīr-i īn lawḥ muḍṭarr namūda ki
shāyad miṣbāḥī chand dar mishkāt-i imkān mustaḍīʾ shavand va bi-qudrat
Allāh bar naṣr-i īn maẓlūm-i farīd qāʾim gardand.

Now reflect a little on these words. Reflection on the difference
between the Qāʾim and the Qayyūm is also necessary, for the sorrows that
have befallen this servant leave Him unable to expound the words of God.
By God, besides whom there is no other god, the proclamation of the
Cause of God has compelled this servant to write this Tablet, that
perhaps a few lamps may be kindled in the niche of contingent existence
and may arise, through the power of God, to aid this wronged and
solitary One.

\subsection{27}\label{section-1582}

\begin{otherlanguage}{persian}

اذاً کلّ شیء یبکی علی ضرّی و بما ورد علیّ من الّذین هم خلقوا بقولی و از همه
گذشته امر بمقامی منجرّ شده که با آنکه ظهور آیات قدس ربّانی بمثابهٴ غیث
هاطل از سماء مشیّت من غیر مکث و سکون نازل و من دونها آیات قدرتیّه و ظهورات
الهیّه که عالم را احاطه نموده بشأنی که ملل قبل مذعن و معترف شده‌اند مع‌ذلک
باید بادلّه استدلال نمایم و امری را که لایزال مقدّس از دلیل بوده بدلیل
ثابت نمایم که لعلّ معدودی بسماء شهود صعود نمایند

\end{otherlanguage}

Idhan kullu shayʾin yabkī ʿalā ḍurrī wa bi-mā warada ʿalayya mina
alladhīna hum khuliqū bi-qawlī va az hama gudhashta amr bi-maqāmī
munjarr shuda ki bā ānki ẓuhūr-i āyāt-i quds-i rabbānī bi-mathāba-yi
ghayth-i hāṭil az samā-yi mashiyyat min ghayr-i makth va sukūn nāzil va
min dūnihā āyāt-i qudratiyya va ẓuhūrāt-i ilāhiyya ki ʿālam rā iḥāṭa
namūda bi-shaʾnī ki milal-i qabl mudhʿin va muʿtarif shuda-and
maʿa-dhālik bāyad bi-adilla istidlāl namāyam va amrī rā ki lā-yazāl
muqaddas az dalīl būda bi-dalīl thābit namāyam ki laʿalla maʿdūdī
bi-samā-yi shuhūd ṣuʿūd namāyand.

Now all things weep over My suffering and over what has befallen Me at
the hands of those who were created by My word. Beyond all this, matters
have reached such a pass that, although the holy verses of the Lord pour
down like a torrent of rain from the heaven of His will without pause or
cessation, and besides these there are signs of power and divine
manifestations that have encompassed the world to such a degree that the
peoples of earlier religions have acknowledged and confessed them, I
must nevertheless argue by proofs and establish by evidence a Cause that
has always been sanctified above proof, so that perhaps a few may ascend
to the heaven of direct witness.

\subsection{28}\label{section-1583}

\begin{otherlanguage}{persian}

ظلمی فوق این در علم الهی نه که جمال قدم بدون خود استدلال بر حقّیّت خود
نماید بعد از آنکه چون شمس در قطب زوال سماء لایزال مشهود و لایح است
فسیعلم الّذین ظلموا نفس اللّه ایّ مرجع یرجعون و بلایای محتومه و قضایای
مثبته بشأنی وارد که جز حقّ احدی بر احصای آن قادر نه در کلّ حین رماح بغضا
از شطر اعدا بر هیکل بقا وارد و لقد جآءنی مظاهر نفس اللّه فی ضحی من الیوم
یبکون و ینوحون قالوا یا اسفا علی یوسف اللّه المهیمن القیّوم فقد اودعوه
عبیده فی الجبّ ثمّ فی انفسهم یستصرخون

\end{otherlanguage}

Ẓulmī fawq-i īn dar ʿilm-i ilāhī na ki Jamāl-i Qidam bi-dūn-i khud
istidlāl bar ḥaqqiyyat-i khud namāyad baʿd az ānki chūn shams dar quṭb-i
zavāl-i samā-yi lā-yazāl mashhūd va lāʾiḥ ast. Fa-sa-yaʿlamu alladhīna
ẓalamū nafsa Allāhi ayya marjaʿin yarjiʿūn. Va balāyā-yi maḥtūma va
qaḍāyā-yi muthbata bi-shaʾnī vārid ki juz Ḥaqq aḥadī bar iḥṣā-yi ān
qādir na. Dar kull-i ḥīn rimāḥ-i baghḍā az shaṭr-i aʿdā bar haykal-i
baqā vārid. Wa laqad jāʾanī maẓāhiru nafsi Allāhi fī ḍuḥan mina al-yawmi
yabkūna wa yanūḥūna qālū yā asafā ʿalā Yūsufi Allāhi al-Muhaymini
al-Qayyūmi fa-qad awdaʿūhu ʿabīduhu fī al-jubbi thumma fī anfusihim
yastaṣrikhūn.

In God's knowledge there is no injustice greater than this: that the
Ancient Beauty should prove His own truth by something other than
Himself, when He is manifest and radiant as the sun at the zenith of the
everlasting heaven. Those who have wronged the very Self of God will
soon know to what place they will return. Inescapable afflictions and
decreed ordeals have descended in such measure that no one but God can
number them. At every moment, the spears of hatred strike the Temple of
eternity from the direction of His enemies. The Manifestations of the
Self of God came to Me in the forenoon, weeping and lamenting. They
said: ``Alas for the Joseph of God, the All-Protecting, the
Self-Subsisting! His servants have cast Him into the pit.'' Then within
themselves they cry out for help.

\subsection{29}\label{section-1584}

\begin{otherlanguage}{arabic}

قل یا ملأ البیان أ تقتلون نقطة الأولی و تقرؤون آیاته فی کلّ عشیّ و بکور
تاللّه قد فعلتم ما لا فعلت امم القبل و یشهد بذلک عباد مکرمون ان انتم
تنکرون أ تقتلون اللّه بأسیاف النّفس و الهوی ثمّ علی مقاعدکم بذکره تشتغلون و
بآیاته تستدلّون کذلک فعل کلّ امّة باللّه حین ظهوره بمظهر نفسه فی کلّ عصر و
کذلک فعلوا و کانوا ان یفعلون قل الیوم لن یحرّک علی اسمآء هؤلآء قلم اللّه
المهیمن العزیز القیّوم و لن یرتدّ الیهم طرف اللّه و لن یأخذهم نفحات قدس
محبوب الیوم

\end{otherlanguage}

Qul yā malaʾa al-Bayān a-taqtulūna Nuqṭata al-Ūlā wa taqraʾūna āyātihi
fī kulli ʿashiyyin wa bukūr? Ta-Allāhi qad faʿaltum mā lā faʿalat umamu
al-qabli wa yashhadu bi-dhālika ʿibādun mukramūn in antum tunkirūn.
A-taqtulūna Allāha bi-asyāfi al-nafsi wa-al-hawā thumma ʿalā maqāʿidikum
bi-dhikrihi tashtaghilūna wa bi-āyātihi tastadillūn? Kadhālika faʿala
kullu ummatin bi-Allāhi ḥīna ẓuhūrihi bi-maẓhari nafsihi fī kulli ʿaṣrin
wa kadhālika faʿalū wa kānū in yafʿalūn. Qul al-yawma lan yuḥarraka ʿalā
asmāʾi hāʾulāʾi qalamu Allāhi al-Muhaymini al-ʿAzīzi al-Qayyūm wa lan
yartadda ilayhim ṭarfu Allāhi wa lan yaʾkhudhahum nafaḥātu qudsi maḥbūbi
al-yawm.

Say: O people of the Bayān! Do you slay the Primal Point while reading
His verses every evening and morning? By God, you have done what the
peoples of former times did not do, and honored servants bear witness to
this, if you deny it. Do you slay God with the swords of self and
desire, then sit in your seats occupying yourselves with His remembrance
and citing His verses as proof? Thus has every community treated God
when He appeared in the Manifestation of Himself in every age. Thus they
acted, and such was their practice. Say: Today the Pen of God, the
All-Protecting, the Almighty, the Self-Subsisting, will not move to
record the names of these people. The gaze of God will not turn toward
them, nor will the holy breezes of the Beloved of this day reach them.

\subsection{30}\label{section-1585}

\begin{otherlanguage}{arabic}

اگر عباد از این آیات قدس الهیّه و ظهورات عزّ صمدانیّه اعراض نمایند بچه حجّت
و دلیل اثبات دین خود مینمایند بگو ای اهل بیان ببصر حقّ بمنظر اکبر ناظر
شوید چه که ببصر دون خود مشهود نگردد اینست که ظهور قبلم میفرماید ایّاکم
فانظروا الیه بعینه فمن ینظر الیه بعین سواه لن یعرفه ابداً و بعد از وصایای
لایحصی عباد خود را در این امر میفرماید فیا الهی انت تعلم بأنّی ما قصّرت فی
نصحی ذلک الخلق و تدبیری لاقبالهم الی اللّه ربّهم و ایمانهم باللّه بارئهم
الی آخر قوله روحی فداه

\end{otherlanguage}

Agar ʿibād az īn āyāt-i quds-i ilāhiyya va ẓuhūrāt-i ʿizz-i ṣamadāniyya
iʿrāḍ namāyand bi-chi ḥujjat va dalīl ithbāt-i dīn-i khud mī-namāyand?
Bi-gū ay ahl-i Bayān bi-baṣar-i Ḥaqq bi-manẓar-i akbar nāẓir shavīd chi
ki bi-baṣar-i dūn-i khud mashhūd na-gardad. Īn ast ki ẓuhūr-i qablam
mī-farmāyad: iyyākum fa-unẓurū ilayhi bi-ʿaynihi fa-man yanẓuru ilayhi
bi-ʿayni siwāhu lan yaʿrifahu abadan. Va baʿd az vaṣāyā-yi lā-yuḥṣā
ʿibād-i khud rā dar īn amr mī-farmāyad: fa-yā ilāhī anta taʿlamu bi-annī
mā qaṣṣartu fī nuṣḥī dhālika al-khalqa wa tadbīrī li-iqbālihim ilā
Allāhi rabbihim wa īmānihim bi-Allāhi bāriʾihim, ilā ākhir-i qawlihi
rūḥī fidāh.

If the servants turn away from these holy divine verses and these
manifestations of eternal glory, by what testimony and proof will they
establish their own religion? Say: O people of the Bayān! Look upon the
Most Great Vision with God's own sight, for He cannot be seen through
any sight other than His own. Thus My previous Manifestation says:
``Beware! Look upon Him with His own eye, for whoever looks upon Him
through the eye of anyone besides Him will never know Him.'' After
giving His servants countless counsels concerning this Cause, He says:
``O My God, You know that I have not fallen short in counseling these
creatures or in making provision for them to turn toward God, their
Lord, and to believe in God, their Maker''---to the end of His words,
may my soul be sacrificed for Him.

\subsection{31}\label{section-1586}

\begin{otherlanguage}{persian}

قسم بجمال قدم که از این بیان که از قلم رحمن جاری شده قلب کلّ اشیاء محترق
گشته و هر ذی بصری از این کلمات نوحه و حنین مظهر اسماء و صفات را ادراک
مینماید ولکن لا یزید الظّالمین الّا غروراً و خساراً ای علیّ گوش جان بگشا و
کلمات رحمن که در قیّوم اسماء نازل شده اصغا نما که میفرماید یا قرّة العین
لا تجعل یدک مبسوطة علی الأمر لأنّ النّاس فی سکران من السّرّ و انّ لک الکرّة
بعد هذه الدّورة بالحقّ الأکبر هنالک فأظهر من السّرّ سرّاً علی قدر سمّ الابرة فی
الطّور الأکبر لیموتنّ الطّوریّون فی السّینا عند مطلع رشح من ذلک النّور المهیمن
الحمرآء باذن اللّه الحکیم و هو اللّه قد کان علیک بالحقّ علی الحقّ حفیظا

\end{otherlanguage}

Qasam bi-Jamāl-i Qidam ki az īn bayān ki az qalam-i Raḥmān jārī shuda
qalb-i kull-i ashyā muḥtariq gashta va har dhī-baṣarī az īn kalimāt
nawḥa va ḥanīn-i maẓhar-i asmā va ṣifāt rā idrāk mī-namāyad wa lākin lā
yazīdu al-ẓālimīna illā ghurūran wa khasāran. Ay ʿAlī gūsh-i jān
bi-gushā va kalimāt-i Raḥmān ki dar Qayyūm-i Asmā nāzil shuda iṣghā namā
ki mī-farmāyad: yā Qurrata al-ʿAyn lā tajʿal yadaka mabsūṭatan ʿalā
al-amri li-anna al-nāsa fī sukrānin mina al-sirri wa inna laka
al-karrata baʿda hādhihi al-dawrati bi-al-ḥaqqi al-akbar. Hunālika
fa-aẓhir mina al-sirri sirran ʿalā qadri sammi al-ibrati fī al-Ṭūri
al-akbari li-yamūtanna al-ṭūriyyūna fī al-Sīnā ʿinda maṭlaʿi rashḥin min
dhālika al-nūri al-muhaymini al-ḥamrāʾi bi-idhni Allāhi al-Ḥakīmi wa
huwa Allāhu qad kāna ʿalayka bi-al-ḥaqqi ʿalā al-ḥaqqi ḥafīẓan.

I swear by the Ancient Beauty that this utterance, which has flowed from
the Pen of the All-Merciful, has set the heart of all things ablaze.
Everyone endowed with sight perceives in these words the lamentation and
yearning of the Manifestation of the names and attributes. Yet it only
increases the wrongdoers in delusion and loss. O ʿAlī, open the ear of
your soul and listen to the words of the All-Merciful revealed in the
Qayyūm-i Asmā, where He says: ``O Qurrat al-ʿAyn! Do not stretch out
Your hand over the Cause, for the people are intoxicated by the mystery.
Yours will be the return after this cycle, through the Most Great Truth.
Then disclose, upon the Most Great Mount, a mystery of the mystery to
the measure of a needle's eye, so that the dwellers of the Mount in
Sinai will surely die at the dawning of a sprinkling of that
all-protecting, crimson Light, by the permission of God, the All-Wise.
He is God, who has indeed been Your Guardian in truth and according to
truth.''

\subsection{32}\label{section-1587}

\begin{otherlanguage}{arabic}

و حال کرّة ظاهر و طوریّون معدوم و مفقود چنانچه مشاهده میشود و با اینکه
میفرماید طوریّون میّت و لاشیء و معدوم میشوند مع‌ذلک تعجّب مینمایند از اینکه
ثمرهٴ طوبی حنظل شود چنانچه بعینه همین را سؤال کرده‌اند و یا صور علّیّین
بسجّین تبدیل گردد لازال امر الهی محدود نبوده و نخواهد بود کلّ مرایاء
موجودات اگر در حین بشمس سماء قدم مقابل شوند در کلّ انوار شمس ظاهر و لایح
و منطبع و مرتسم و بمجرّد انحراف از کلّ اخذ میشود فانظر فی الشّمس ثمّ فی
المرایا لکی تجد الی ما یلقیک الرّوح سبیلا

\end{otherlanguage}

Va ḥāl karra ẓāher va ṭūriyyūn maʿdūm va mafqūd chonānche moshāhede
mīshavad va bā īnke mīfarmāyad ṭūriyyūn mayyet va lā-shayʾ va maʿdūm
mīshavand maʿa dhālek taʿajj ob mīnamāyand az īnke thamare-ye Ṭūbā
ḥanẓal shavad chonānche be-ʿayneh hamīn rā soʾāl karde-and va yā ṣovar-e
ʿilliyyīn be-sijjīn tabdīl gardad lāzāl amr-e elāhī maḥdūd nabūde va
nakhvāhad būd koll-e marāyā-ye mawjūdāt agar dar ḥīn be-shams-e samā-ye
qedam moqābel shavand dar koll anvār-e shams ẓāher va lāyeḥ va monṭabeʿ
va mortasem va be-mojarrad-e enḥerāf az koll akhdh mīshavad fa-nẓur fī
al-shams thumma fī al-marāyā li-kay tajida ilā mā yulqīka al-rūḥ
sabīlan.

And now the return is manifest, while the dwellers of the Mount are
nonexistent and lost, as can be seen. Although He says that the dwellers
of the Mount become dead, nothing, and nonexistent, they nevertheless
express surprise that the fruit of the Tree of Bliss should become
colocynth---as indeed they have asked this very question---or that the
forms of the highest heaven should be transformed into those of the
nethermost prison. The Cause of God has never been, and will never be,
confined by limits. If all the mirrors of existing things face the sun
in the heaven of eternity for a moment, the sun's lights appear, shine,
are imprinted, and are reflected in them all; the instant they turn
away, those lights are withdrawn from them all. Look, then, at the sun
and then at the mirrors, so that you may find a way to what the Spirit
imparts to you.

\subsection{33}\label{section-1588}

\begin{otherlanguage}{persian}

ابن نبیل مرفوع در اثبات امر اللّه بما القی اللّه علی فؤاده الواحی نوشته و
در ابتدا باین آیه که از سماء مشیّت ظهور قبلم نازل شده استدلال نموده قوله
عزّ ذکره قل اللّهمّ انّک انت الهان الالهین لتؤتینّ الألوهیّة من تشآء و لتنزعنّ
الألوهیّة عمّن تشآء الی آخر و کذلک قل اللّهمّ انّک انت ربّاب السّموات و الأرض
لتؤتینّ الرّبوبیّة من تشآء و لتنزعنّ الرّبوبیّة عمّن تشآء الی آخر با آنکه سلطان
وجود باین صریحی فرموده که عطا میفرماید الوهیّت و ربوبیّت را بهر نفسی که
اراده فرماید و اخذ میفرماید از هر که بخواهد خداوند قادری که مقام الوهیّت
و ربوبیّت که اعلی مقامات است اخذ فرماید قادر نیست بر اینکه از هیکلی قمیص
اسم خود را نزع نماید یا آنکه حلوی را بمرّ تبدیل فرماید سبحان اللّه عمّا
یتوهّمون العباد فی قدرته فتعالی عمّا یصفون حال مشاهده نمائید که طیر قدرت و
عظمت در چه هوا طیران مینماید و ناس در چه اماکن توقّف نموده‌اند آیا آیهٴ
انّه علی کلّ شیء قدیر را چه معنی نموده‌اند و از یفعل ما یشآء و لا یسأل عمّا
شآء چه ادراک کرده‌اند

\end{otherlanguage}

Ibn-e Nabīl-e marfūʿ dar ethbāt-e amr-e Allāh bi-mā alqā Allāh ʿalā
fuʾādihi alvāḥī neveshte va dar ebt edā be-īn āye ke az samā-ye
mashiyyat-e ẓohūr-e qablam nāzel shode estedlāl namūde qawluhu ʿazza
dhikruhu qul Allāhumma innaka anta ilāhān al-ilāhīn la-tuʾtiyanna
al-ulūhiyya man tashāʾ wa-la-tanziʿanna al-ulūhiyya ʿamman tashāʾ ilā
ākhir va kadhālika qul Allāhumma innaka anta rabbāb al-samāwāt wa-al-arḍ
la-tuʾtiyanna al-rubūbiyya man tashāʾ wa-la-tanziʿanna al-rubūbiyya
ʿamman tashāʾ ilā ākhir bā ānke solṭān-e vojūd be-īn ṣarīḥī farmūde ke
ʿaṭā mīfarmāyad olūhiyyat va robūbiyyat rā be-har nafsī ke erāde
farmāyad va akhdh mīfarmāyad az har ke bekhvāhad Khodāvand-e qāderī ke
maqām-e olūhiyyat va robūbiyyat ke aʿlā-ye maqāmāt ast akhdh farmāyad
qāder nīst bar īnke az haykalī qamīṣ-e esm-e khod rā nazʿ namāyad yā
ānke ḥolvī rā be-morr tabdīl farmāyad subḥāna Allāh ʿammā yatawahhamūna
al-ʿibād fī qudratihi fa-taʿālā ʿammā yaṣifūn ḥāl moshāhede namāyīd ke
ṭayr-e qodrat va ʿaẓamat dar che havā ṭayarān mīnamāyad va nās dar che
amāken tavaqqof namūde-and āyā āye-ye innahu ʿalā kulli shayʾin qadīr rā
che maʿnī namūde-and va az yafʿalu mā yashāʾ wa-lā yusʾalu ʿammā shāʾ
che edrāk karde-and.

The exalted Ibn-i Nabīl has written tablets in proof of the Cause of God
through what God has inspired in his heart. At the outset he has adduced
this verse, revealed from the heaven of the will of My previous
Manifestation. His words---exalted be His remembrance---are: ``Say: O
God, You are indeed the God of gods. You will surely bestow divinity on
whomever You will, and You will surely strip divinity from whomever You
will''---and so forth. Likewise: ``Say: O God, You are indeed the Lord
of the heavens and the earth. You will surely bestow lordship on
whomever You will, and You will surely strip lordship from whomever You
will''---and so forth. Although the Sovereign of existence has stated so
explicitly that He bestows divinity and lordship on any soul He wishes
and takes them from whomever He wills, is the mighty God who takes away
the station of divinity and lordship, the highest of stations, unable to
strip the robe of His name from a person, or to turn something sweet
into something bitter? Glory be to God, far above what His servants
imagine concerning His power! Exalted is He above what they describe!
Now observe in what atmosphere the bird of power and grandeur is
soaring, and in what places the people have come to a halt. What meaning
have they given to the verse, ``He has power over all things''? And what
have they understood from ``He does whatever He wills'' and ``He is not
questioned about what He has willed''?

\subsection{34}\label{section-1589}

\begin{otherlanguage}{persian}

ای عباد از اوهن بیوت برکن شدید متمسّک شوید و از ظلمات جهل و نادانی بفجر
منیر علم ربّانی توجّه نمائید و کاش اهل بیان در آیهٴ مبارکه که ابن نبیل ذکر
نموده تفکّر نمایند که شاید از سبل وهم بصراط یقین درآیند ای علیّ یک قدح از
این ماء عذب حیوان که در ظلمات کلمات سلطان اسماء و صفات مستور شده بیاشام
تا از کدورات ایّام و شبهات انام و اشارات غافلین و دلالات مغلّین پاک و مقدّس
شوی و ابواب علوم نامتناهی ربّانی بر وجه قلبت مفتوح شود تا آنکه موقن شوی
باینکه سلطان قدم قادر است بر آنکه در ساعتی جمیع اشیاء را بخلع اسماء حسنی
مفتخر و معزّز فرماید و در ساعت اخری از جمیع اخذ نماید و انّی اشکو الی اللّه
من هؤلآء العباد لأنّهم ینظرون الیّ بما عندهم لا بما عندی و یقاسون نفس اللّه
بأنفسهم و کلماته بکلماتهم فوالّذی نفسی نفسه لو ینقطعنّ الیوم کلّ من فی
السّموات و الأرض و یقبلنّ الی اللّه لیعلّمهم من بدایع علمه ما یغنیهم عن
العالمین

\end{otherlanguage}

Ey ʿebād az awhan-e boyūt be-rokn-e shadīd motamassek shavīd va az
ẓolomāt-e jahl va nādānī be-fajr-e monīr-e ʿelm-e rabbānī tavajjoh
namāyīd va kāsh ahl-e Bayān dar āye-ye mobārake ke Ibn-e Nabīl dhekr
namūde tafakkor namāyand ke shāyad az sobol-e vahm be-ṣerāṭ-e yaqīn
dar-āyand ey ʿAlī yek qadaḥ az īn māʾ-e ʿadhb-e ḥayavān ke dar ẓolomāt-e
kalemāt-e solṭān-e asmā va ṣefāt mastūr shode biyāshām tā az kodūrāt-e
ayyām va shobohāt-e anām va eshārāt-e ghāfelīn va dalālāt-e moghellīn
pāk va moqaddas shavī va abvāb-e ʿolūm-e nāmotanāhī-ye rabbānī bar
vajh-e qalbat maftūḥ shavad tā ānke mūqen shavī be-īnke solṭān-e qedam
qāder ast bar ānke dar sāʿatī jamīʿ-e ashyā rā be-khelaʿ-e asmā-ye ḥosnā
moftakhar va moʿazzaz farmāyad va dar sāʿat-e okhrā az jamīʿ akhdh
namāyad wa-innī ashkū ilā Allāh min hāʾulāʾ al-ʿibād li-annahum
yanẓurūna ilayya bi-mā ʿindahum lā bi-mā ʿindī wa-yuqāsūna nafs Allāh
bi-anfusihim wa-kalimātihi bi-kalimātihim fa-wa-alladhī nafsī nafsuhu
law yanqaṭiʿanna al-yawm kullu man fī al-samāwāt wa-al-arḍ wa-yuqbilanna
ilā Allāh la-yuʿallimahum min badāʾiʿ ʿilmihi mā yughnīhim ʿan
al-ʿālamīn.

O servants! Turn from the frailest of houses and take hold of a strong
support, and turn from the darkness of ignorance and unknowing toward
the luminous dawn of divine knowledge. Would that the people of the
Bayān would ponder the blessed verse cited by Ibn-i Nabīl, so that they
might leave the paths of illusion and enter the path of certainty. O
ʿAlī! Drink a cup of this sweet water of life concealed in the darkness
of the words of the Sovereign of names and attributes, so that you may
be purified and sanctified from the troubles of the days, the doubts of
humanity, the allusions of the heedless, and the arguments of the
malicious, and so that the gates of infinite divine knowledge may be
opened before your heart. Then you may become certain that the Sovereign
of eternity has the power, in one hour, to honor and exalt all things
with the robes of the most beautiful names, and in another hour to take
them from them all. I complain to God of these servants, for they look
upon Me through what they possess, not through what I possess, and they
measure God's own Self against themselves and His words against their
words. By Him whose Self is My own! If all who are in the heavens and on
earth were today to detach themselves and turn toward God, He would
assuredly teach them such wonders of His knowledge as would make them
independent of all the worlds.

\subsection{35}\label{section-1590}

\begin{otherlanguage}{persian}

و در این سنهٴ شداد بعضی از عباد اینگونه مسائل سؤال نموده‌اند و عبد حاضر
لدی العرش جوابهای محکمهٴ شافیهٴ کافیه نوشته ارسال داشته عجب است که شما
ندیده‌اید و در این ارض هم بعضی از شبهات القا نموده‌اند که شاید نفوس قدسیّه
را بحجبات کلمات قبلیّه محتجب دارند ولکن غافل از اینکه نفسی که در هوای
بافضای قدس الهی طیران نمود و از خرمنهای علوم نامتناهی ربّانی القاط فرمود
بشبهات وهمیّه ممنوع نشود و محروم نگردد قل من ورد علی بحر الأعظم لن یلتفت
الی سراب بقیعة و لن یشرب من مآء الحمیم اگرچه لایق نه که قلم اعلی بذکر
کلمات اولو البغضآء بیالاید و یا حرکت نماید ولکن نظر بتبلیغ رسالات ربّانی
مفرّی نه که شاید جاهلی ببحر علم درآید و یا گم‌گشتهٴ وادی غفلت و نسیان بیمن
رحمن خرامد و انّه لیهدی من یشآء الی صراط البهآء و انّه علی کلّ شیء قدیر

\end{otherlanguage}

Va dar īn sane-ye shedād baʿżī az ʿebād īngūne masāʾel soʾāl namūde-and
va ʿabd-e ḥāżer laday al-ʿarsh javābhā-ye moḥkame-ye shāfiye-ye kāfiye
neveshte ersāl dāshte ʿajab ast ke shomā nadīde-īd va dar īn arż ham
baʿżī az shobohāt elqā namūde-and ke shāyad nofūs-e qodsiyye rā
be-ḥojobāt-e kalemāt-e qabliyye moḥtajeb dārand va-lāken ghāfel az īnke
nafsī ke dar havā-ye bā-faẓā-ye qods-e elāhī ṭayarān namūd va az
kharmanhā-ye ʿolūm-e nāmotanāhī-ye rabbānī elqāṭ farmūd be-shobohāt-e
vahmiyye mamnūʿ nashavad va maḥrūm nagardad qul man warada ʿalā baḥr
al-aʿẓam lan yaltafita ilā sarābin bi-qīʿa wa-lan yashraba min māʾ
al-ḥamīm agarche lāyeq na ke qalam-e aʿlā be-dhekr-e kalemāt-e ulū
al-baghḍā biyālāyad va yā ḥarakat namāyad va-lāken naẓar be-tablīgh-e
resālāt-e rabbānī mafarrī na ke shāyad jāhelī be-baḥr-e ʿelm dar-āyad va
yā gom-gashte-ye vādī-ye ghaflat va nesyān be-yomn-e Raḥmān kherāmad
wa-innahu la-yahdī man yashāʾ ilā ṣirāṭ al-Bahāʾ wa-innahu ʿalā kulli
shayʾin qadīr.

In this year of hardships, some of the servants have asked questions of
this kind, and the servant present before the Throne has written and
sent firm, healing, and sufficient answers. It is strange that you have
not seen them. In this land, too, some have cast doubts, hoping to veil
holy souls behind the veils of earlier words. Yet they are unaware that
a soul who has soared in the spacious atmosphere of divine holiness and
gleaned from the harvests of infinite divine knowledge will neither be
held back nor deprived by fanciful doubts. Say: Whoever has reached the
Most Great Ocean will never turn toward a mirage on a plain, nor drink
of scalding water. Although it is not fitting that the Most Exalted Pen
should be sullied by mentioning the words of those possessed of hatred,
or move to write them, there is nevertheless no escape from doing so, in
order to convey the divine messages, so that perhaps an ignorant person
may enter the ocean of knowledge, or one lost in the valley of
heedlessness and forgetfulness may, through the grace of the
All-Merciful, advance. He indeed guides whomever He wills to the path of
Bahā, and He has power over all things.

\subsection{36}\label{section-1591}

\begin{otherlanguage}{persian}

ولکن این ذکر و بیان برای نفوسی است که فی‌الحقیقه طالب سبیل هدایت باشند
والّا علیل غلّ و بغضا را سلسبیل طاهر ننماید چنانچه الیوم اکثری اهل بیان
جمیع ظهورات قدرتیّه و شئونات الهیّه و آیات منزله را بچشم خود دیده‌اند و
بگوش خود شنیده‌اند مع‌ذلک ببغضی قیام نموده‌اند که ذکر آن ممکن نه الّا من فتح
اللّه بصره و ایّده علی امره و اخرجه عن ظلمات الوهم و هداه الی صراطه العزیز
الحمید و تمسّک باموری جسته‌اند که لم‌یزل عنداللّه مذکور نبوده و بشأنی غافلند
که بحجر ساجد شده‌اند و از منظر اکبر و جمال اظهر انور معرض گشته‌اند وربّ
معبود شعر من عابدیه خیر منه ان یا علیّ قل الیوم انتم و معبودکم عند اللّه
فی حدّ سوآء فمن آمن منکم فقد امن من فزع الأکبر و من اعرض فقد خرج عن صراط
اللّه و انّ هذا لهو الحقّ و ما بعد الحقّ الّا الضّلال

\end{otherlanguage}

Va-lāken īn dhekr va bayān barā-ye nofūsī ast ke fī al-ḥaqīqe ṭāleb-e
sabīl-e hedāyat bāshand va-ellā ʿalīl-e ghell va baghżā rā salsabīl
ṭāher nanamāyad chonānche al-yawm aktharī-ye ahl-e Bayān jamīʿ-e
ẓohūrāt-e qodratiyye va shoʾūnāt-e elāhiyye va āyāt-e monzale rā
be-chashm-e khod dīde-and va be-gūsh-e khod shenīde-and maʿa dhālek
be-boghżī qiyām namūde-and ke dhekr-e ān momken na illā man fataḥa Allāh
baṣarahu wa-ayyadahu ʿalā amrihi wa-akhrajahu ʿan ẓulumāt al-wahm
wa-hadāhu ilā ṣirāṭihi al-ʿazīz al-ḥamīd va tamassok be-omūrī joste-and
ke lam-yazal ʿinda Allāh madhkūr nabūde va be-shaʾnī ghāfeland ke
be-ḥajar sājed shode-and va az manẓar-e akbar va jamāl-e aẓhar-e anvar
moʿreż gashte-and wa-rubba maʿbūdin shaʿrun min ʿābidīhi khayrun minhu
an yā ʿAlī qul al-yawm antum wa-maʿbūdukum ʿinda Allāh fī ḥaddi sawāʾ
fa-man āmana minkum fa-qad amina min fazaʿ al-akbar wa-man aʿraḍa fa-qad
kharaja ʿan ṣirāṭ Allāh wa-inna hādhā la-huwa al-ḥaqq wa-mā baʿda
al-ḥaqq illā al-ḍalāl.

But this discourse and exposition are for souls who truly seek the way
of guidance; otherwise, even Salsabīl will not purify one afflicted with
malice and hatred. Thus today most of the people of the Bayān have seen
with their own eyes and heard with their own ears all the manifestations
of power, the divine signs, and the revealed verses, yet they have
arisen in such hatred that it cannot be described---except those whose
eyes God has opened, whom He has strengthened in His Cause, brought out
of the darkness of illusion, and guided to His mighty and praiseworthy
path. They have clung to things that have never had any standing before
God, and they are so heedless that they have bowed down before a stone
and turned away from the Most Great Vision and the most manifest, most
luminous Beauty. How many an object of worship is worth less than a hair
of its worshipers! O ʿAlī! Say: Today you and the one you worship are on
an equal footing before God. Whoever among you believes is safe from the
Most Great Terror, and whoever turns away has departed from the path of
God. This is indeed the truth, and beyond truth there is nothing but
error.

\subsection{37}\label{section-1592}

\begin{otherlanguage}{arabic}

ایّاکم ایّاکم یا ملأ البیان لا تکفروا باللّه و لا تحاربوا بمظهر نفسه و لا
تجادلوا بالّذی جآءکم عن مشرق الأمر بسلطان مبین آیا گمان مینمائید که امر
اللّه باعراض معرضین ممنوع شود و یا انوار شمس عزّ باقی باکمام انفس ظلمانی
مستور ماند لا فوالّذی نطق فی صدری و بعثنی بالحقّ و ارسلنی علی العالمین

\end{otherlanguage}

Iyyākum iyyākum yā malaʾ al-Bayān lā takfurū bi-Allāh wa-lā tuḥāribū
bi-maẓhar nafsihi wa-lā tujādilū bi-alladhī jāʾakum ʿan mashriq al-amr
bi-sulṭānin mubīn āyā gomān mīnamāyīd ke amr-e Allāh be-eʿrāż-e moʿreżīn
mamnūʿ shavad va yā anvār-e shams-e ʿezz-e bāqī be-akmām-e anfos-e
ẓolmānī mastūr mānad lā fa-wa-alladhī naṭaqa fī ṣadrī wa-baʿathanī
bi-al-ḥaqq wa-arsalanī ʿalā al-ʿālamīn.

Beware, beware, O people of the Bayān! Do not disbelieve in God, do not
wage war against the Manifestation of His own Self, and do not dispute
with Him who has come to you from the Dayspring of the Cause with
manifest authority. Do you imagine that the Cause of God can be
obstructed by those who turn away, or that the light of the sun of
everlasting glory can remain concealed within the coverings of darkened
souls? No, by Him who has spoken within My breast, raised Me up in
truth, and sent Me to all the worlds!

\subsection{38}\label{section-1593}

\begin{otherlanguage}{persian}

و بعضی از مشرکین از جمله شبهات که در این ارض القا نموده‌اند اینست که آیا
میشود ذهب نحاس شود قل ای وربّی ولکن عندنا علمه نعلّم من نشآء بعلم من لدنّا
و من کان فی ریب فلیسأل اللّه ربّه بأن یشهده و یکون من الموقنین و در رسیدن
نحاس برتبهٴ ذهبیّت همان دلیلی است واضح بر عود ذهب بحالت اوّل لو هم یشعرون
جمیع فلزّات بوزن و صورت و مادّه یکدیگر میرسند ولکن علمه عندنا فی کتاب
مکنون

\end{otherlanguage}

Va baʿżī az moshrekīn az jomle-ye shobohāt ke dar īn arż elqā namūde-and
īnast ke āyā mīshavad dhahab noḥās shavad qul ī wa-rabbī wa-lākin
ʿindanā ʿilmuhu nuʿallimu man nashāʾ bi-ʿilmin min ladunnā wa-man kāna
fī raybin fa-l-yasʾal Allāh rabbahu bi-an yushhidahu wa-yakūna min
al-mūqinīn va dar resīdan-e noḥās be-rotbe-ye dhahabiyyat hamān dalīlī
ast vāżeḥ bar ʿawd-e dhahab be-ḥālat-e avval law hum yashʿurūn jamīʿ-e
felezzāt be-vazn va ṣūrat va mādde-ye yekdīgar mīresand va-lākin ʿilmuhu
ʿindanā fī kitābin maknūn.

Among the doubts that some of those who associate partners with God have
raised in this land is this: Can gold become copper? Say: Yes, by My
Lord! But the knowledge of it is with Us. We teach whomever We will
through knowledge from Our presence. Whoever is in doubt should ask God,
his Lord, to enable him to witness it and become one of those who are
certain. The very fact that copper attains the rank of gold is a clear
proof that gold can return to its former state, if only they understood.
All metals attain the weight, form, and substance of one another, but
the knowledge of this is with Us in a hidden Book.

\subsection{39}\label{section-1594}

\begin{otherlanguage}{persian}

میگوئیم علم معرضین باین مقام صعود ننموده که ادراک نمایند ذهب نحاس میشود
آنقدر هم ادراک ننموده که تراب میشود این رتبه که مشهود هر ذی شعوری بوده
که کلّ از تراب ظاهر و بتراب راجع و تراب در قدر و قیمت ارخص از نحاس است چه
که او از اجسام محسوب و نحاس از اجساد و این بسی ظاهر و هویدا است و اگر
ناس لایق و بالغ مشاهده میشدند هرآینه در این مقام ذکر بعضی از علوم
مستورهٴ الهیّه میشد ولکن قضی ما قضی بر هر ذی بصری مشهود است که حقّ تعالی
ذکره بر کلّ شیء قادر بوده و خواهد بود البتّه اگر بخواهد بمجرّد اراده ذهب را
بنحاس تبدیل میفرماید و این عجز در موجودات موجود و مشهود و انّه لهو
المقتدر العزیز القادر المحمود

\end{otherlanguage}

Mīgūʾīm ʿelm-e moʿreżīn be-īn maqām ṣoʿūd nanamūde ke edrāk namāyand
dhahab noḥās mīshavad ānqadr ham edrāk nanamūde ke torāb mīshavad īn
rotbe ke mashhūd-e har dhī shoʿūrī būde ke koll az torāb ẓāher va
be-torāb rājeʿ va torāb dar qadr va qīmat arkhaṣ az noḥās ast che ke ū
az ajsām maḥsūb va noḥās az ajsād va īn basī ẓāher va hoveydā ast va
agar nās lāyeq va bālegh moshāhede mīshodand har-āyene dar īn maqām
dhekr-e baʿżī az ʿolūm-e mastūre-ye elāhiyye mīshod va-lāken quḍiya mā
quḍiya bar har dhī baṣarī mashhūd ast ke Ḥaqq taʿālā dhikruhu bar koll-e
shayʾ qāder būde va khvāhad būd albatte agar bekhvāhad be-mojarrad-e
erāde dhahab rā be-noḥās tabdīl mīfarmāyad va īn ʿajz dar mawjūdāt
mawjūd va mashhūd va-innahu la-huwa al-muqtadir al-ʿazīz al-qādir
al-maḥmūd.

We say: The knowledge of those who turn away has not risen to the point
of understanding that gold can become copper; they have not even
understood that it becomes dust. Yet this much has been evident to
everyone endowed with understanding: all things emerge from dust and
return to dust. Dust is lower in worth and price than copper, for it is
reckoned among bodies, while copper is reckoned among metallic bodies.
This is plain and evident. If people had been found worthy and mature,
some of the hidden divine sciences would certainly have been mentioned
here. But what has been decreed has been decreed. It is evident to
everyone endowed with sight that God---exalted be His remembrance---has
always had, and will always have, power over all things. Certainly, if
He wishes, He transforms gold into copper by the mere act of willing. It
is in created things that this powerlessness exists and is seen. He is
indeed the All-Powerful, the Almighty, the Omnipotent, the Praised.

\subsection{40}\label{section-1595}

\begin{otherlanguage}{arabic}

نظر را مطهّر نموده بمنظر اکبر توجّه نمائید و از اشجار لایغنی و لایثمر
منقطع شوید اینست از بدایع امر الهی فمن شآء فلیقبل و من شآء فلیعرض فمن
اقبل فلنفسه و من اعرض فلها و انّه لهو المقدّس عن الخلایق اجمعین

\end{otherlanguage}

Naẓar rā moṭahhar namūde be-manẓar-e akbar tavajjoh namāyīd va az
ashjār-e lā-yoghnī va lā-yothmer monqaṭeʿ shavīd īnast az badāyeʿ-e
amr-e elāhī fa-man shāʾ fa-l-yuqbil wa-man shāʾ fa-l-yuʿriḍ fa-man
aqbala fa-li-nafsihi wa-man aʿraḍa fa-lahā wa-innahu la-huwa al-muqaddas
ʿan al-khalāʾiq ajmaʿīn.

Purify your sight and turn toward the Most Great Vision, and detach
yourselves from trees that neither profit nor bear fruit. This is among
the wonders of the Cause of God. Whoever wishes, let him turn toward it,
and whoever wishes, let him turn away. Whoever turns toward it does so
for his own sake, and whoever turns away does so to his own detriment.
He indeed is sanctified above all created things.

\subsection{41}\label{section-1596}

\begin{otherlanguage}{persian}

در شیطان تفکّر نما که معلّم ملکوت بوده در ملأ اعلی و در مداین اسماء بأسماء
حسنی معروف و بعد باعراضی از اعلی رفیق اعلی بادنی ارض سفلی مقرّ گرفته کذلک
یفعل ربّک ما یشآء ان انت من الموقنین از اینها گذشته بیت عتیق که کعبهٴ
موجودات بوده و محلّ طواف مظاهر اسماء و صفات چرا از این فضل کبری محروم شد
اذاً تفکّروا یا اولو الألباب کلّ شیء در قبضهٴ قدرت الهی اسیرند و در کلّ حین
بآنچه اراده فرماید قادر و مقتدر است و نفاذ قدرت محیطه‌اش در هیچ اوان از
مظاهر امکان و اکوان سلب نشده و نخواهد شد

\end{otherlanguage}

Dar Shayṭān tafakkor namā ke moʿallem-e malakūt būde dar malaʾ-e aʿlā va
dar madāyen-e asmā be-asmā-ye ḥosnā maʿrūf va baʿd be-eʿrāżī az aʿlā-ye
rafīq-e aʿlā be-adnā-ye arż-e soflā maqarr gerefte kadhālika yafʿalu
rabbuka mā yashāʾ in anta min al-mūqinīn az īnhā goẕashte bayt-e ʿatīq
ke Kaʿbe-ye mawjūdāt būde va maḥall-e ṭavāf-e maẓāher-e asmā va ṣefāt
cherā az īn fażl-e kobrā maḥrūm shod idhan tafakkarū yā ulū al-albāb
koll-e shayʾ dar qabże-ye qodrat-e elāhī asīrand va dar koll-e ḥīn
be-ānche erāde farmāyad qāder va moqtader ast va nafādh-e qodrat-e
moḥīṭe-ash dar hīch avān az maẓāher-e emkān va akvān salb nashode va
nakhvāhad shod.

Reflect on Satan, who was the teacher of the Kingdom in the Concourse on
high and was known by the most beautiful names in the cities of names.
Yet through one act of turning away, he went from the loftiest heights
of the supreme companionship to make his abode in the lowest depths of
the earth below. Thus does your Lord do whatever He wills, if you are
one of those who are certain. Beyond these examples, why was the Ancient
House, which had been the Kaaba of all beings and the place around which
the manifestations of the names and attributes circled, deprived of this
supreme grace? Reflect, then, O people of understanding! All things are
held captive in the grasp of divine power, and at every moment He is
able and powerful to do whatever He wills. The sway of His
all-encompassing power over the manifestations of contingent existence
and created things has never been, and will never be, withdrawn.

\subsection{42}\label{section-1597}

\begin{otherlanguage}{persian}

اسرار مکنونه‌اش را هر گوشی لایق استماع نه و حوریّات معانی مقصورهٴ علمیّه‌اش
را هر چشمی قابل مشاهده نه چه مقدار از هیاکل ظلم که بقمیص عدل بین عباد
معروف شده‌اند و چه مقدار از حقایق عدلیّه که در اثواب ظلم اشتهار یافته‌اند
نظر در اصنام نما که حال نصف من علی الأرض باو عاکف شده‌اند و من دون اللّه
معبود اخذ نموده‌اند و باین مرض مبتلا نشده‌اند مگر آنکه بوهم و تقلید اکتفا
کرده‌اند و از سلطان توحید اعراض نموده‌اند باری الیوم مظاهر کلّ اسماء و صفات
در صقع واحد و موقف واحد مشهودند الّا من صعد الی اللّه کذلک نلقی علیک لعلّ
تطهّر نفسک و صدرک عن کلمات العالمین و تسمع ما غرّد الرّوح علی افنان هذه
السّدرة الّتی احاطت کلّ من فی السّموات و الأرضین

\end{otherlanguage}

Asrār-e maknūne-ash rā har gūshī lāyeq-e estemāʿ na va ḥūriyyāt-e
maʿānī-ye maqṣūre-ye ʿelmiyye-ash rā har chashmī qābel-e moshāhede na
che meqdār az hayākel-e ẓolm ke be-qamīṣ-e ʿadl bayn-e ʿebād maʿrūf
shode-and va che meqdār az ḥaqāyeq-e ʿadliyye ke dar athvāb-e ẓolm
eshtehār yāfte-and naẓar dar aṣnām namā ke ḥāl neṣf-e man ʿalā al-arż
be-ū ʿākef shode-and va min dūn-e Allāh maʿbūd akhdh namūde-and va be-īn
maraż mobtalā nashode-and magar ānke be-vahm va taqlīd ektefā karde-and
va az solṭān-e tawḥīd eʿrāż namūde-and bārī al-yawm maẓāher-e koll-e
asmā va ṣefāt dar ṣoqʿ-e vāḥed va mawqef-e vāḥed mashhūdand illā man
ṣaʿida ilā Allāh kadhālika nulqī ʿalayka laʿalla tuṭahhiru nafsaka
wa-ṣadraka ʿan kalimāt al-ʿālamīn wa-tasmaʿu mā gharrada al-rūḥ ʿalā
afnān hādhihi al-sidra allatī aḥāṭat kulla man fī al-samāwāt
wa-al-araḍīn.

Not every ear is worthy to hear His hidden mysteries, nor is every eye
capable of beholding the maidens of meaning secluded within the chambers
of His knowledge. How many embodiments of oppression have become known
among the servants in the robe of justice, and how many realities of
justice have become renowned in the garments of oppression! Consider the
idols: half of those on earth now devote themselves to them and have
taken them as objects of worship instead of God. They have been
afflicted with this disease only because they have contented themselves
with illusion and imitation and turned away from the Sovereign of divine
unity. In short, today the manifestations of all names and attributes
are seen on one plane and in one station, except those who have ascended
to God. Thus do We impart to you, so that you may purify your soul and
your breast of the words of the peoples of the world and hear what the
Spirit has sung upon the branches of this Lote Tree, which encompasses
all who are in the heavens and the earths.

\subsection{43}\label{section-1598}

\begin{otherlanguage}{arabic}

قل یا ملأ البیان تاللّه الحقّ لم یکن هذا من تلقآء نفسی بل بما نطق اللّه فی
صدری و ما ظهر من سلطانی و جری من قلمی برهانی ثمّ حجّتی ثمّ دلیلی ان انتم من
المنصفین قل انتم بأیّ حجّة آمنتم بعلیّ من قبل حین الّذی ظهر بالحقّ و جآءکم
بسلطان مبین و بأیّ برهان صدّقتم آیاته و اذعنتم برهانه و خضعتم عند ظهورات
امره المهیمن المتعالی العزیز المنیر

\end{otherlanguage}

Qul yā malaʾ al-Bayān ta-llāh al-Ḥaqq lam yakun hādhā min tilqāʾ nafsī
bal bi-mā naṭaqa Allāh fī ṣadrī wa-mā ẓahara min sulṭānī wa-jarā min
qalamī burhānī thumma ḥujjatī thumma dalīlī in antum min al-munṣifīn.
Qul antum bi-ayyi ḥujja āmantum bi-ʿAlī min qabl ḥīna alladhī ẓahara
bi-l-ḥaqq wa-jāʾakum bi-sulṭān mubīn wa-bi-ayyi burhān ṣaddaqtum āyātihi
wa-adhʿantum burhānahu wa-khaḍaʿtum ʿinda ẓuhūrāt amrihi al-Muhaymin
al-Mutaʿālī al-ʿAzīz al-Munīr.

Say: O people of the Bayān! By God, the Truth, this has not come from My
own prompting, but through what God has spoken within My breast.
Whatever has appeared through My sovereignty and flowed from My Pen is
My proof, My testimony, and My evidence, if you are fair-minded. Say: By
what testimony did you formerly believe in ʿAlī, when He appeared in
truth and came to you with manifest authority? By what proof did you
acknowledge His verses, assent to His demonstration, and bow before the
revelations of His Cause, the All-Protecting, the Most Exalted, the
Almighty, the Luminous?

\subsection{44}\label{section-1599}

\begin{otherlanguage}{persian}

و ان تقولوا انّا آمنّا به بنفسه و اکتفینا بحجّیّة نفسه عمّا سواه قل تاللّه هذا
نفسه قد قام بین العباد و ظهر بسلطان اسمه المقتدر المهیمن العلیّ العظیم و
ان تقولوا بأنّا آمنّا بما نزّل علیه من آیات اللّه العزیز الغالب القدیر فتلک
آیاته ملأت شرق الأرض و غربها اذاً فاستمعوا لما یوحی عن شطر القدس الّذی
یسمع من اریاحها ما سمع اذن الحبیب فی معراج القدس ثمّ اذن الرّوح فی سموات
الأمر ثمّ اذن الکلیم علی طور العزّ عن شجرة اللّه النّاطق العزیز الحکیم و من
دونهما قد ظهر هذا الغلام بسلطنة الّتی علت علی الممکنات و یشهد بذلک السن
الکائنات ان انتم من السّامعین

\end{otherlanguage}

Wa-in taqūlū innā āmannā bihi bi-nafsihi wa-ktafaynā bi-ḥujjiyyat
nafsihi ʿammā siwāhu qul ta-llāh hādhā nafsuhu qad qāma bayna al-ʿibād
wa-ẓahara bi-sulṭān ismihi al-Muqtadir al-Muhaymin al-ʿAlī al-ʿAẓīm.
Wa-in taqūlū bi-annā āmannā bi-mā nuzzila ʿalayhi min āyāt Allāh
al-ʿAzīz al-Ghālib al-Qadīr fa-tilka āyātuhu malaʾat sharq al-arḍ
wa-gharbahā. Idhan fa-stamiʿū li-mā yūḥā ʿan shaṭr al-quds alladhī
yusmaʿ min aryāḥihā mā samiʿa udhun al-Ḥabīb fī miʿrāj al-quds thumma
udhun al-Rūḥ fī samāwāt al-amr thumma udhun al-Kalīm ʿalā Ṭūr al-ʿizz
ʿan shajarat Allāh al-Nāṭiq al-ʿAzīz al-Ḥakīm. Wa-min dūnihimā qad
ẓahara hādhā al-ghulām bi-salṭanat allatī ʿalat ʿalā al-mumkināt
wa-yashhad bi-dhālika alsun al-kāʾināt in antum min al-sāmiʿīn.

And if you say, ``We believed in Him through His own Self and were
content with the testimony of His Self, apart from all else,'' say: By
God, this is His very Self, who has arisen among His servants and
appeared with the sovereignty of His name, the All-Powerful, the
All-Protecting, the Most High, the Most Great. And if you say, ``We
believed through the verses of God, the Almighty, the All-Conquering,
the All-Powerful, that were sent down to Him,'' these are His verses,
which have filled the east and the west of the earth. Listen, then, to
what is revealed from the realm of holiness, upon whose breezes can be
heard what the ear of the Beloved heard on His ascent to holiness, then
the ear of the Spirit in the heavens of the Cause, then the ear of Moses
upon the Mount of glory, from the Tree of God, the Speaker, the
Almighty, the All-Wise. Apart from these two, this Youth has appeared
with a sovereignty that transcends all contingent things. To this the
tongues of all beings bear witness, if you are among those who listen.

\subsection{45}\label{section-1600}

\begin{otherlanguage}{persian}

ثمّ قل لرؤسآء البیان این کنتم حین الّذی اضطربت فیه انفس العباد و زلّت فیه
الأقدام و غشی الرّعب قلوب الرّاسخین و قام علیّ کلّ العباد من مذاهب شتّی و ما
استنصرت من احد الّا اللّه الّذی بعثنی و ارسلنی علی العالمین تاللّه هم کانوا
مستوراً خلف قناع النّسآء فلمّا ظهر الأمر بسلطانه اطمأنّوا فی انفسهم و خرجوا
عن الحجاب فأوّل ما فعلوا اعرضوا عن الّذی به ثبت ایمانهم کذلک کان الأمر و
انّک کنت من الشّاهدین و انّک ان لن تصدّقنی بعد الّذی شهدت بعینک یصدّقنی کلّ
الأشیآء و عن ورائها لسان اللّه الصّادق الأمین

\end{otherlanguage}

Thumma qul li-ruʾasāʾ al-Bayān ayna kuntum ḥīna alladhī iḍṭarabat fīhi
anfus al-ʿibād wa-zallat fīhi al-aqdām wa-ghashiya al-ruʿb qulūb
al-rāsikhīn wa-qāma ʿalayya kull al-ʿibād min madhāhib shattā wa-mā
istanṣartu min aḥad illā Allāh alladhī baʿathanī wa-arsalanī ʿalā
al-ʿālamīn. Ta-llāh hum kānū mastūran khalfa qināʿ al-nisāʾ fa-lammā
ẓahara al-amr bi-sulṭānihi iṭmaʾannū fī anfusihim wa-kharajū ʿan
al-ḥijāb fa-awwal mā faʿalū aʿraḍū ʿan alladhī bihi thabata īmānuhum.
Kadhālika kāna al-amr wa-innaka kunta min al-shāhidīn. Wa-innaka in lan
tuṣaddiqanī baʿda alladhī shahidta bi-ʿaynika yuṣaddiqunī kull al-ashyāʾ
wa-ʿan warāʾihā lisān Allāh al-Ṣādiq al-Amīn.

Then say to the leaders of the Bayān: Where were you when the souls of
the servants were shaken, when feet slipped and terror overwhelmed the
hearts of the steadfast, and all the servants, of diverse creeds, rose
against Me? I sought help from no one but God, who raised Me up and sent
Me to the peoples of the world. By God, they were concealed behind the
veil of women! Then, when the Cause appeared in its sovereignty, they
felt secure within themselves and emerged from behind the veil. The
first thing they did was to turn away from Him through whom their faith
had been established. Such was the matter, and you were among its
witnesses. And if you will not acknowledge My truth after what you have
witnessed with your own eyes, all things will acknowledge My truth, and
beyond them, the tongue of God, the Truthful, the Trustworthy.

\subsection{46}\label{section-1601}

\begin{otherlanguage}{persian}

ان یا جمال الکبریآء بین الأرض و السّمآء غیّر اللّحن لأهل الانشآء ثمّ غنّ علی
افنان البقآء علی لحن عجمیّ منیع لیکشف اسرار الأمر فیما رقم من هذا القلم
المحکم المتین چه که جمیع ناس بلغات عربیّه مطّلع نه و ادراک کلمات پارسیّه
اسهل است نزد اهل لسان ای سائل آنچه از اسامی در کتاب الهی از ذکر طوبی و
سدرهٴ منتهی و شجرهٴ قصوی و ورقه و ثمره و امثال آن مشاهده مینمائی موهوم
مدان مقصود از جمیع این اسامی عنداللّه مؤمن باللّه بوده و خواهد بود و مؤمن
تا در ظلّ سدرهٴ الهیّه ساکن از سدرهٴ طوبی و علّیّین عنداللّه محسوب و بعد از
اعراض از سدرهٴ نار و سجّین مذکور و در حین ایمان افنان و اغصان و اوراق و
اثمار او جمیع از اثبات مشهود و بعد از اعراض جمیع از نفی محسوب میشود و
بسا نفسی که در اصیل از ابهی سدرهٴ بقا است و در بکور از ادنی شجرهٴ فنا و
کذلک بالعکس لو انت من العارفین

\end{otherlanguage}

An yā Jamāl al-Kibriyāʾ bayna al-arḍ wa-l-samāʾ ghayyir al-laḥn li-ahl
al-inshāʾ thumma ghanni ʿalā afnān al-baqāʾ ʿalā laḥn ʿajamī manīʿ
li-yukshafa asrār al-amr fīmā ruqima min hādhā al-qalam al-muḥkam
al-matīn. Che ke jamīʿ-e nās be-loghāt-e ʿarabiyye moṭṭaleʿ na va
edrāk-e kalemāt-e pārsiyye ashal ast nazd-e ahl-e lesān. Ey sāʾel ānche
az asāmī dar ketāb-e elāhī az dhekr-e Ṭūbā va Sedre-ye Montahā va
Shajare-ye Qoṣvā va varaqe va thamare va amthāl-e ān moshāhede mīnamāʾī
mowhūm madān. Maqṣūd az jamīʿ-e īn asāmī ʿinda Allāh moʾmen be-Allāh
būde va khāhad būd va moʾmen tā dar ẓell-e Sedre-ye Elāhiyye sāken az
Sedre-ye Ṭūbā va ʿIlliyyīn ʿinda Allāh maḥsūb va baʿd az eʿrāḍ az
Sedre-ye Nār va Sijjīn madhkūr va dar ḥīn-e īmān afnān va aghṣān va
owrāq va athmār-e ū jamīʿ az ethbāt mashhūd va baʿd az eʿrāḍ jamīʿ az
nafy maḥsūb mīshavad va basā nafsī ke dar aṣīl az abhā Sedre-ye Baqā ast
va dar bokūr az adnā Shajare-ye Fanā wa-kadhālika bi-l-ʿaks law anta min
al-ʿārifīn.

O Beauty of majesty between earth and heaven! Change the melody for the
people of creation, then sing upon the branches of eternity in an
exalted Persian melody, so that the mysteries of the Cause may be
disclosed in what is inscribed by this firm and steadfast Pen. For not
everyone is acquainted with the Arabic tongue, and understanding Persian
words is easier for those who speak that language. O questioner! Do not
regard as imaginary the names you encounter in the Book of God, such as
Ṭūbā, the Lote Tree of the Utmost Limit, the Farthest Tree, leaf, fruit,
and the like. In the sight of God, what is meant by all these names has
been and will be the believer in God. As long as a believer dwells
beneath the shade of the divine Lote Tree, that believer is reckoned in
God's sight as belonging to the Lote Tree of Ṭūbā and to ʿIlliyyīn;
after turning away, that believer is described as belonging to the Tree
of Fire and to Sijjīn. While in a state of faith, that believer's
boughs, branches, leaves, and fruits are all seen as belonging to
affirmation; after turning away, they are all reckoned as belonging to
negation. How many a soul belongs in the evening to the most glorious
Lote Tree of eternity, and in the morning to the lowliest Tree of
extinction! And likewise the reverse, if you are among those who
understand.

\subsection{47}\label{section-1602}

\begin{otherlanguage}{persian}

مؤمن را در حین اقبال او الی اللّه جنّتی مشاهده کن با کمال تزیین بشأنی که
جمیع آنچه در جنّت شنیده‌ئی در او مشاهده نما از افنان علمیّه و اثمار معارف
الهیّه و انهار بیانیّه و ازهار حکمتیّه و فوق ذلک الی ان یشآء اللّه در او
موجود و همین نفس بعینه بعد از اعراض نفس هاویه میشود مع آنچه متعلّق باو
بوده کذلک یبدّل اللّه النّور بالظّلمة و الظّلمة بالنّور لو انتم تفقهون آیا
نشنیده‌ئی که ظهور قبلم در ارض همین حکم جاری فرموده چنانچه میفرماید بر هر
ارضی که مؤمن مستقرّ شود از ارض جنّت عنداللّه مذکور است و اگر مشرک بهمان ارض
مستقرّ شود از ارض جحیم محسوب چنانچه الیوم مقرّ عرش ارضی واقع شده که ابداً
معروف نبوده ولکن الیوم مبغضین جمال ربّ العالمین خود را از علّیّین محروم
ساخته‌اند و در قعر سجّین مقرّ گرفته‌اند و بزعم خود در اعلی مقعد جنّت ساکنند
چنانچه ملل قبل هم بهمین اوهام مشغولند اذاً یبشّرهم قلم الأمر بعذاب یوم
عقیم

\end{otherlanguage}

Moʾmen rā dar ḥīn-e eqbāl-e ū ilā Allāh jannatī moshāhede kon bā kamāl-e
tazyīn be-shaʾnī ke jamīʿ-e ānche dar jannat shenīdeʾī dar ū moshāhede
namā az afnān-e ʿelmiyye va athmār-e maʿāref-e elāhiyye va anhār-e
bayāniyye va azhār-e ḥekmatiyye va fowq-e dhālek ilā an yashāʾa Allāh
dar ū mowjūd va hamīn nafs be-ʿayneh baʿd az eʿrāḍ nafs-e Hāviye
mīshavad maʿ-e ānche motaʿalleq be-ū būde. Kadhālika yubaddilu Allāh
al-nūr bi-l-ẓulma wa-l-ẓulma bi-l-nūr law antum tafqahūn. Āyā
nashenīdeʾī ke ẓohūr-e qablam dar arḍ hamīn ḥokm jārī farmūde chenānche
mīfarmāyad bar har arḍī ke moʾmen mostaqarr shavad az arḍ-e jannat ʿinda
Allāh madhkūr ast va agar moshrek be-hamān arḍ mostaqarr shavad az arḍ-e
jaḥīm maḥsūb chenānche al-yowm maqarr-e ʿarsh arḍī vāqeʿ shode ke abadan
maʿrūf nabūde va-lāken al-yowm mobgheḍīn-e Jamāl-e Rabb al-ʿĀlamīn khod
rā az ʿIlliyyīn maḥrūm sākhte-and va dar qaʿr-e Sijjīn maqarr
gerefte-and va be-zaʿm-e khod dar aʿlā maqʿad-e jannat sākenand
chenānche melal-e qabl ham be-hamīn owhām mashghūland. Idhan
yubashshiruhum Qalam al-Amr bi-ʿadhāb yawm ʿaqīm.

When a believer turns toward God, behold that believer as a paradise
adorned to perfection, and see within that soul everything you have
heard of paradise: branches of knowledge, fruits of divine
understanding, rivers of utterance, and flowers of wisdom. All these,
and more beyond them, as God wills, are present within that soul. Yet
this very same soul, after turning away, becomes hell itself, together
with everything that belonged to it. Thus does God exchange light for
darkness and darkness for light, if you understand. Have you not heard
that My former Manifestation on earth pronounced this very judgment,
saying: ``Whatever land a believer settles in is accounted in God's
sight as the land of paradise; and if an idolater settles in that same
land, it is reckoned as the land of hell''? Thus today the seat of the
Throne has come to be in a land that was never renowned. Yet today those
who hate the Beauty of the Lord of the worlds have deprived themselves
of ʿIlliyyīn and made their abode in the depths of Sijjīn, while
imagining that they dwell in the loftiest seat of paradise, just as the
peoples of the past are occupied with these same delusions. The Pen of
the Cause therefore gives them tidings of the punishment of a barren
day.

\subsection{48}\label{section-1603}

\begin{otherlanguage}{persian}

این قوم را لایق آنکه عجلی من دون اللّه اخذ نمایند و باو ساجد و عاکف شوند
چنانچه شده‌اند کجا لایقند بهوای قدس صمدانی طیران نمایند و یا بسماء عزّ
الهی راجع شوند چنانچه مشاهده میشود آنکه طبعش گل میل نموده ابداً بگل ملتفت
نه کذلک یضرب اللّه مثلاً لعلّ النّاس هم یشعرون آیا نشنیده‌اید که میفرماید بسا
شجرهٴ اثبات که در ظهور بعد از شجرهٴ نفی میشود باری الیوم هر نفسی که از
کلّ آنچه مابین عباد مشهود و مذکور است منقطع نشود و جمیع را چون کفّ طین
مشاهده ننماید ابداً قادر نه که باین هوا طیران نماید و یا بمقرّ سلطان عزّ
تقدیس درآید

\end{otherlanguage}

Īn qowm rā lāyeq ānke ʿejlī min dūni Allāh akhdh namāyand va be-ū sājed
va ʿākef shavand chenānche shode-and. Kojā lāyeqand be-havā-ye qods-e
ṣamadānī ṭayarān namāyand va yā be-samā-ye ʿezz-e elāhī rājeʿ shavand.
Chenānche moshāhede mīshavad ānke ṭabʿash gel meyl namūde abadan be-gol
moltafet na. Kadhālika yaḍribu Allāh mathalan laʿalla al-nās hum
yashʿurūn. Āyā nashenīde-īd ke mīfarmāyad basā Shajare-ye Ethbāt ke dar
ẓohūr-e baʿd az Shajare-ye Nafy mīshavad. Bārī al-yowm har nafsī ke az
koll-e ānche mābeyn-e ʿebād mashhūd va madhkūr ast monqaṭeʿ nashavad va
jamīʿ rā chon kaff-e ṭīn moshāhede nanamāyad abadan qāder na ke be-īn
havā ṭayarān namāyad va yā be-maqarr-e Solṭān-e ʿEzz-e Taqdīs darāyad.

It befits these people to take a calf instead of God, bow before it, and
devote themselves to it, as indeed they have done. How could they be
worthy to soar in the atmosphere of divine holiness or ascend to the
heaven of divine glory? As can be seen, one whose nature inclines toward
clay pays no attention to the rose. Thus does God offer a parable, so
that people may understand. Have you not heard that He says: ``Many a
Tree of Affirmation becomes a Tree of Negation in the next
Manifestation''? In short, today any soul who does not become detached
from everything seen and mentioned among the servants, and does not
regard it all as a handful of clay, will never be able to soar in this
atmosphere or enter the dwelling of the Sovereign of sanctified glory.

\subsection{49}\label{section-1604}

\begin{otherlanguage}{persian}

لازال مؤمن اقلّ از کبریت احمر بوده و خواهد بود و علاوه بر این شموس معانی
که از مشرق اصبع رحمانی اشراق نمود این بسی مبرهن و واضح است که ثمره بنفسه
لنفسه موجود نه بلکه باعانت الهیّه از امکنهٴ ترابیّه صعود نموده تا باین
رتبه ظاهر شده که مقام ثمری باشد و آن مقتدری که او را باین مقام فایز نمود
بهمان قدرت قادر است که او و صدهزار امثال و اشباح را در اقلّ من آن از اعلی
مقرّ بقا بادنی مقعد فنا راجع فرماید و هم‌چنین بالعکس و این از سنن او بوده
و خواهد بود چه که قدرت محیطه و قضایای محتومهٴ آن سلطان احدیّه لازال بر کلّ
شیء نافذ بوده و اقلّ من آن از ظهورات قدرت خود ممنوع نبوده فسبحان اللّه عمّا
یظنّون المتوهّمون

\end{otherlanguage}

Lāzāl moʾmen aqall az kebrīt-e aḥmar būde va khāhad būd va ʿelāve bar īn
shomūs-e maʿānī ke az mashreq-e eṣbaʿ-e raḥmānī eshrāq namūd īn basī
mobarhan va vāḍeḥ ast ke thamare be-nafseh le-nafseh mowjūd na balke
be-eʿānat-e elāhiyye az amkene-ye torābiyye ṣoʿūd namūde tā be-īn rotbe
ẓāher shode ke maqām-e thamarī bāshad va ān moqtaderī ke ū rā be-īn
maqām fāʾez namūd be-hamān qodrat qāder ast ke ū va ṣad-hezār amthāl va
ashbāḥ rā dar aqall-e men ān az aʿlā maqarr-e baqā be-adnā maqʿad-e fanā
rājeʿ farmāyad va ham-chenīn be-l-ʿaks va īn az sonan-e ū būde va khāhad
būd che ke qodrat-e moḥīṭe va qaḍāyā-ye maḥtūme-ye ān Solṭān-e Aḥadiyye
lāzāl bar koll-e shayʾ nāfedh būde va aqall-e men ān az ẓohūrāt-e
qodrat-e khod mamnūʿ nabūde. Fa-subḥān Allāh ʿammā yaẓunnūn
al-mutawahhimūn.

A believer has always been, and will always be, rarer than red sulfur.
Beyond these suns of meaning that have shone from the dawning-place of
the finger of the All-Merciful, it is abundantly evident and clear that
a fruit does not exist by itself and for itself. Rather, through divine
assistance, it has risen from earthly places until it has appeared at
this rank, the station of bearing fruit. The Mighty One who enabled it
to attain this station is able, through that same power, to return it
and a hundred thousand others like it, in less than an instant, from the
loftiest abode of eternity to the lowliest seat of extinction---and
likewise the reverse. This has been and will be among His ways, for the
all-encompassing power and irrevocable decrees of that Sovereign of
oneness have always held sway over all things, and never for less than
an instant has He been prevented from manifesting His power. Exalted,
then, is God above what the deluded imagine!

\subsection{50}\label{section-1605}

\begin{otherlanguage}{persian}

بسا لآلی اسرار که در اصداف بحار اسم ربّک السّتّار مستور بوده و خواهد بود که
اظهار آن سبب اعثار اقدام غیر مستقیمه شود باری الیوم فضلی ظاهر شده که اگر
اراده فرماید از کفی از طین کلّ حروفات اوّلیّه و آخریّه را مبعوث فرماید قادر
است مع‌ذلک بسیار حیف است انسان در این ایّام که جمال رحمن بتمام فضل ظاهر
شده خود را بغیر او مشغول نماید دع کلّ من فی السّموات و الأرض لأهلها ثمّ
ادخل فی غمرات هذا البحر الّذی لن یوجد فیه الّا لآلئ ذکر اسم ربّک العلیّ
المقتدر العظیم

\end{otherlanguage}

Basā laʾālī-ye asrār ke dar aṣdāf-e beḥār-e esm-e rabbika al-Sattār
mastūr būde va khāhad būd ke eẓhār-e ān sabab-e eʿthār-e aqdām-e gheyr-e
mostaqīme shavad. Bārī al-yowm faḍlī ẓāher shode ke agar erāde farmāyad
az kaffī az ṭīn koll-e ḥorūfāt-e avvaliyye va ākheriyye rā mabʿūth
farmāyad qāder ast. Maʿa-dhālek besyār ḥeyf ast ensān dar īn ayyām ke
Jamāl-e Raḥmān be-tamām-e faḍl ẓāher shode khod rā be-gheyr-e ū mashghūl
namāyad. Daʿ kulla man fī al-samāwāt wa-l-arḍ li-ahlihā thumma udkhul fī
ghamarāt hādhā al-baḥr alladhī lan yūjada fīhi illā laʾāliʾ dhikr ism
rabbika al-ʿAlī al-Muqtadir al-ʿAẓīm.

Many pearls of mysteries have been and will remain concealed within the
shells of the seas of the name of your Lord, the All-Concealing, whose
disclosure would cause unsteady feet to stumble. In short, today a grace
has appeared such that, should He wish to raise up all the first and
last Letters from a handful of clay, He has the power to do so.
Nevertheless, it is a great pity for anyone, in these days when the
Beauty of the All-Merciful has appeared in the fullness of His grace, to
be occupied with anything other than Him. Leave all who are in the
heavens and the earth to their own people, then plunge into the depths
of this Ocean, in which nothing will be found but the pearls of the
remembrance of the name of your Lord, the Most High, the All-Powerful,
the Most Great.

\subsection{51}\label{section-1606}

\begin{otherlanguage}{persian}

اینست بدایع ظهورات شمس امر رحمن که از افق اصبع ملیک امکان اشراق فرموده
فطوبی للعاملین و من دون ذلک انّ ربّک لغنیّ عن العالمین ای بسا نفوسی که
الیوم در ابحر نار مستغرقند و مستشعر نیستند بلکه خود را از اهل جنّت
میدانند چنانچه امم قبل هم باین ظنون مسرورند قسم بآفتاب عزّ قدس تجرید که
این ظهور اعظم از آن است که بدلیل محتاج باشد و یا ببرهان منوط گردد قل انّ
دلیله ظهوره و حجّته نفسه و وجوده اثباته و برهانه قیامه بین السّموات و
الأرض فی ایّام الّتی فیها اضطربت کلّ من فی ملکوت الأمر و الخلق اجمعین و ان
لن تقدرنّ ان تعرفنّه بما فصّلنا لکم فاعرفوه بما نزّل من عنده و کذلک قدّر لکم
فضلاً من عنده و انّه لهو الفضّال القدیم

\end{otherlanguage}

Īnast badāyeʿ-e ẓohūrāt-e Shams-e Amr-e Raḥmān ke az ofoq-e eṣbaʿ-e
Malīk-e Emkān eshrāq farmūde fa-ṭūbā li-l-ʿāmilīn wa-min dūni dhālika
inna rabbaka la-ghaniyy ʿan al-ʿālamīn. Ey basā nofūsī ke al-yowm dar
abḥor-e nār mostaghraqand va mostashʿer nīstand balke khod rā az ahl-e
jannat mīdānand chenānche omam-e qabl ham be-īn ẓonūn masrūrand. Qasam
be-Āftāb-e ʿEzz-e Qods-e Tajrīd ke īn ẓohūr aʿẓam az ān ast ke be-dalīl
moḥtāj bāshad va yā be-borhān manūṭ gardad. Qul inna dalīlahu ẓuhūruhu
wa-ḥujjatahu nafsuhu wa-wujūduhu ithbātuhu wa-burhānahu qiyāmuhu bayna
al-samāwāt wa-l-arḍ fī ayyām allatī fīhā iḍṭarabat kull man fī malakūt
al-amr wa-l-khalq ajmaʿīn. Wa-in lan taqdirunna an taʿrifunnahu bi-mā
faṣṣalnā lakum fa-ʿrifūhu bi-mā nuzzila min ʿindihi wa-kadhālika quddira
lakum faḍlan min ʿindihi wa-innahu la-huwa al-Faḍḍāl al-Qadīm.

These are the wondrous revelations of the Sun of the Cause of the
All-Merciful, which has shone from the horizon of the finger of the
Sovereign of contingent existence. Blessed are those who act; otherwise,
your Lord is indeed independent of all the worlds. How many souls today
are submerged in seas of fire and are unaware of it! Rather, they
consider themselves among the people of paradise, just as the peoples of
the past rejoice in these same suppositions. I swear by the Sun of the
glory of holy detachment that this Revelation is too great to need
evidence or depend upon proof. Say: Its evidence is its appearance, its
testimony is its own Self, its existence is its confirmation, and its
proof is its arising between the heavens and the earth in days when all
who are in the kingdoms of revelation and creation were shaken. And if
you cannot recognize Him through what We have explained to you,
recognize Him through what has been sent down from Him. Thus has it been
ordained for you as a grace from Him. He is indeed the Most Bountiful,
the Ancient.

\subsection{52}\label{section-1607}

\begin{otherlanguage}{persian}

قلم اعلی میفرماید ای علیّ یکبار بطور تقدیس قدم گذار و بقلب فارغ و لسان
طاهر ربّ ارنی گو تا لازال از مکمن قدس بی‌زوال انظر ترانی بشنوی و بلقاء
جمال بیمثال حضرت ذو الجلال فایز گردی یعنی لقای مظهر نفس او که بیک تجلّی
از تجلّیات انوار فضلش یتغنّی اشجار الوجود من الغیب و الشّهود بما نطقت سدرة
الطّور این چنین احاطه فرموده فضل سلطان یفعل ما یشآء ولکنّ النّاس هم فی وهم
عظیم و حجاب غلیظ و غفلة مبین اینست شأن این ناس که لایزال بقول حقّ افتخار
مینمایند و از نفس او معرض مثلاً حجر را طواف میکنند و از اماکن بعیده طیّ
سبلهای صعبه مینمایند و از جان و مال میگذرند تا بزیارتش فایز شوند ولکن از
سلطان مقتدری که بقول او صدهزار امثال این حجر خلق میشود غافل بلکه معرضند

\end{otherlanguage}

Qalam-e Aʿlā mīfarmāyad ey ʿAlī yek bār be-Ṭūr-e Taqdīs qadam godhār va
be-qalb-e fāregh va lesān-e ṭāher rabbi arinī gū tā lāzāl az makman-e
qods-e bī-zavāl unẓur tarānī beshnavī va be-leqā-ye Jamāl-e Bī-methāl-e
Ḥaḍrat-e Dhu-l-Jalāl fāʾez gardī yaʿnī leqā-ye Maẓhar-e Nafs-e ū ke
be-yek tajallī az tajalliyāt-e anvār-e faḍlash yataghannā ashjār
al-wujūd min al-ghayb wa-l-shuhūd bi-mā naṭaqat Sidrat al-Ṭūr. Īn chenīn
eḥāṭe farmūde faḍl-e Solṭān-e yafʿalu mā yashāʾ wa-lākinna al-nās hum fī
wahm ʿaẓīm wa-ḥijāb ghalīẓ wa-ghafla mubīn. Īnast shaʾn-e īn nās ke
lāyazāl be-qowl-e Ḥaqq eftekhār mīnamāyand va az nafs-e ū moʿreḍ
mathalan ḥajar rā ṭavāf mīkonand va az amāken-e baʿīde ṭayy-e sobolhā-ye
ṣaʿbe mīnamāyand va az jān va māl mīgodharand tā be-zeyāratash fāʾez
shavand va-lāken az Solṭān-e Moqtaderī ke be-qowl-e ū ṣad-hezār amthāl-e
īn ḥajar khalq mīshavad ghāfel balke moʿreḍand.

The Most Exalted Pen says: O ʿAlī! For once, set foot upon the Mount of
sanctity and, with an unburdened heart and a pure tongue, say, ``My
Lord, show me,'' so that you may continually hear from the imperishable
sanctuary of holiness, ``Look, and you will see Me,'' and attain the
presence of the peerless Beauty of the Lord of majesty---that is, the
presence of the Manifestation of His own Self. Through a single
effulgence of the lights of His grace, the trees of existence, unseen
and seen, sing what the Lote Tree of the Mount has spoken. Such is the
all-encompassing grace of the Sovereign who does whatever He wills. Yet
the people are in great delusion, behind a thick veil, and in manifest
heedlessness. Such is the condition of these people: they continually
take pride in the word of the True One, yet turn away from His own Self.
For example, they circle around the stone, traverse difficult paths from
distant places, and sacrifice their lives and possessions to attain the
privilege of visiting it. Yet they are heedless of---indeed, turn away
from---the mighty Sovereign through whose word a hundred thousand stones
like this are created.

\subsection{53}\label{section-1608}

\begin{otherlanguage}{persian}

چنانچه ملاحظه شد در سنهٴ ستّین و هم‌چنین در این ایّام بصر منیر را بصدهزار
حجبات وهمیّه و سبحات نفسیّه مستور مینمایند و بعد فریاد برآرند که آفتاب
جهانتاب عزّ صمدانی طالع نشده و اگر هم ادراک نمایند سؤال نموده که فلان نجم
چگونه میشود نورش محو گردد و یا زایل شود دیگر غافل از اینکه بوجود شمس
یمحو نوره چنانچه مشاهده میشود که در ایّام انوار نجوم محو و لایظهرند باری
عنقریب ید قدرت محیطهٴ الهیّه نفوسی چند خلق فرماید که جمیع احجاب را خرق
نمایند و بی ستر و حجاب بمکمن ربّ الأرباب درآیند و در سبیل محبوب از هیچ
آبی مخمود نشوند و از هیچ ناری جزع ننمایند غیر معبود را مفقود شمرند و
ماسوای مقصود را معدوم مشاهده نمایند و در کلّ اوان اهل امکان را بسلاسل
بیان برضوان قدس رحمن کشند قسم بآفتاب معانی که انوار این نفوس اهل ملأ
اعلی را مستنیر نماید چنانچه شمس اهل ارض را

\end{otherlanguage}

Chunānchih mulāḥiẓih shud dar sanih-yi sittīn va ham-chunīn dar īn ayyām
baṣar-i munīr rā bi-ṣad-hizār ḥujubāt-i vahmiyyih va subuḥāt-i nafsiyyih
mastūr mīnamāyand va baʿd faryād bar-ārand kih āftāb-i jahān-tāb-i
ʿizz-i ṣamadānī ṭāliʿ nashudih va agar ham idrāk namāyand suʾāl namūdih
kih fulān najm chigūnih mīshavad nūrash maḥv gardad va yā zāʾil shavad
dīgar ghāfil az īnkih bi-vujūd-i shams yamḥū nūrahu chunānchih
mushāhidih mīshavad kih dar ayyām anvār-i nujūm maḥv va lā-yaẓhar-and
bārī ʿan-qarīb yad-i qudrat-i muḥīṭih-yi ilāhiyyih nufūsī chand khalq
farmāyad kih jamīʿ-i aḥjāb rā kharq namāyand va bī sitr va ḥijāb
bi-makman-i Rabb al-Arbāb dar-āyand va dar sabīl-i Maḥbūb az hīch ābī
makhmūd nashavand va az hīch nārī jazaʿ nanamāyand ghayr-i Maʿbūd rā
mafqūd shumurand va mā-sivā-yi Maqṣūd rā maʿdūm mushāhidih namāyand va
dar kull-i avān ahl-i imkān rā bi-salāsil-i bayān bi-Riḍvān-i quds-i
Raḥmān kashand qasam bi-āftāb-i maʿānī kih anvār-i īn nufūs ahl-i
malaʾ-i aʿlā rā mustanīr namāyad chunānchih shams ahl-i arḍ rā

As was observed in the year sixty, and likewise in these days, they veil
their radiant vision with a hundred thousand veils of vain imagining and
coverings of self, and then cry out that the world-illumining Sun of
everlasting glory has not risen. Even if they perceive it, they ask how
the light of a particular star could be effaced or disappear, heedless
that the presence of the sun effaces its light, just as one observes
that in the daytime the lights of the stars are effaced and do not
appear. Soon, however, the hand of God's all-encompassing power will
create souls who will rend every veil and enter, without covering or
veil, the sanctuary of the Lord of lords. In the path of the Beloved, no
water will quench them and no fire will cause them to fear. They will
reckon all but the One they worship as absent, and regard all besides
the Object of their desire as nonexistent. At all times they will draw
the people of the world of existence, with the chains of utterance, into
the paradise of the holiness of the All-Merciful. I swear by the Sun of
inner meanings that the lights of these souls will illumine the dwellers
of the heavenly assembly, just as the sun illumines the people of the
earth.

\subsection{54}\label{section-1609}

\begin{otherlanguage}{arabic}

ان یا علیّ دع ما عندک و خذ ما یأمرک به اللّه و انّ هذا لیغنیک عن العالمین و
من دونه لایسمن و لایغنی بوده و خواهد بود تمسّک بالعروة الوثقی و دع ما
یأمرک به الهوی تاللّه الحقّ انّ هذا لحبل المحکم الّذی ظهر بین الأرض و السّمآء
فمن تمسّک به فقد نجا و من اعرض فقد هلک

\end{otherlanguage}

In yā ʿAlī daʿ mā ʿindaka wa-khudh mā yaʾmuruka bihi Allāh wa-inna hādhā
la-yughnīka ʿan al-ʿālamīn va min dūnihi lā yusminu wa-lā yughnī būdih
va khvāhad būd tamassak bi-l-ʿurwati al-wuthqā wa-daʿ mā yaʾmuruka bihi
al-hawā ta-Allāhi al-Ḥaqq inna hādhā la-ḥablu al-muḥkam alladhī ẓahara
bayna al-arḍi wa-l-samāʾ fa-man tamassaka bihi fa-qad najā wa-man aʿraḍa
fa-qad halak

O ʿAlī! Leave what you possess and take what God commands you. This will
indeed make you independent of all the worlds, while anything besides it
has never nourished or satisfied, nor ever will. Hold fast to the surest
handle and abandon what desire commands you. By God, the Truth! This is
indeed the firm cord that has appeared between earth and heaven. Whoever
holds fast to it has been saved, and whoever turns away has perished.

\subsection{55}\label{section-1610}

\begin{otherlanguage}{persian}

باری الیوم هر نفسی حقّ منیع را بدون او تعالی بخواهد عارف شود و یا ادراک
نماید مثل آن است که از اکمه طلب ارائهٴ سبیل نماید و هذا لم یمکن ابداً
دلیله آیاته و سلطانه اثباته

\end{otherlanguage}

Bārī al-yawm har nafsī Ḥaqq-i manīʿ rā bidūn-i Ū taʿālā bikhvāhad ʿārif
shavad va yā idrāk namāyad mithl-i ān ast kih az akmah ṭalab-i irāʾih-yi
sabīl namāyad wa-hādhā lam yumkin abadan dalīluhu āyātuhu wa-sulṭānuhu
ithbātuhu

In short, whoever today wishes to recognize or comprehend the exalted
True One without Him---exalted is He---is like someone asking a person
blind from birth to show the way. This has never been possible. His
proof is His verses, and His sovereignty is His confirmation.

\subsection{56}\label{section-1611}

\begin{otherlanguage}{arabic}

ای علیّ بلسان بدیع پارسی کلمات ربّانی را اصغا نما و اگر تو نشنوی البتّه
قدرت محیطهٴ الهیّه انفسی بدعاً خلق فرماید که بشنوند ندای او را و بر نصر
امرش قیام نمایند لیس هذا علی اللّه بعزیز

\end{otherlanguage}

Ay ʿAlī bi-lisān-i badīʿ-i pārsī kalimāt-i rabbānī rā iṣghā namā va agar
tū nashnavī al-battih qudrat-i muḥīṭih-yi ilāhiyyih anfusī badʿan khalq
farmāyad kih bishnavand nidā-yi Ū rā va bar naṣr-i amrash qiyām namāyand
laysa hādhā ʿalā Allāhi bi-ʿazīz

O ʿAlī! Listen to the divine words in the wondrous Persian tongue. If
you do not listen, God's all-encompassing power will assuredly create
souls anew who will hear His call and arise to aid His Cause. This is
not difficult for God.

\subsection{57}\label{section-1612}

\begin{otherlanguage}{persian}

و اگر از شمال وهم بیمین یقین راجع شدی و از کوثر عرفان جمال رحمن که در
رضوان معانی جاری شده مرزوق گشتی بگو ای اهل بیان بکدام دین متدیّنید و بر
کدام صراط قائم اگر بگویند بنقطهٴ بیان روح ما سواه فداه بگو بچه حجّت و
دلیل بآن سلطان سبیل موقن شده‌اید و معترف گشته‌اید اگر بگویند او را بنفسه
شناخته‌ایم بگو هذا کذب صراح چه که الی حین بنفس خود عارف نشده‌اید تا چه رسد
بنفس اللّه القائمة علی کلّ شیء و اگر هم این قول از این نفوس مسموع آید چرا
بنفس اللّه الظّاهرة که چون شمس مشرق است مذعن نشده‌اند و اگر بگویند بآیات
منزله مؤمن شده‌ایم چرا باین آیات که بمثابۀ غیث هاطل از سماء غیب در کلّ حین
نازل است کافر شده‌اند قل أ تؤمنون ببعض الکتاب و تکفرون ببعض فویل لکم یا
معشر الظّالمین

\end{otherlanguage}

Va agar az shimāl-i vahm bi-yamīn-i yaqīn rājiʿ shudī va az kawthar-i
ʿirfān-i Jamāl-i Raḥmān kih dar Riḍvān-i maʿānī jārī shudih marzūq
gashtī bigū ay ahl-i Bayān bi-kudām dīn mutadayyinīd va bar kudām ṣirāṭ
qāʾim agar bigūyand bi-Nuqṭih-yi Bayān rūḥu mā siwāhu fidāhu bigū
bi-chih ḥujjat va dalīl bi-ān sulṭān-i sabīl mūqin shudih-īd va muʿtarif
gashtih-īd agar bigūyand ū rā bi-nafsih shinākhtih-īm bigū hādhā
kadhibun ṣurāḥ chih kih ilā ḥīn bi-nafs-i khvud ʿārif nashudih-īd tā
chih rasad bi-nafsi Allāhi al-qāʾimati ʿalā kulli shayʾ va agar ham īn
qawl az īn nufūs masmūʿ āyad chirā bi-nafsi Allāhi al-ẓāhirah kih chūn
shams mushriq ast mudhʿin nashudih-and va agar bigūyand bi-āyāt-i
munzalih muʾmin shudih-īm chirā bi-īn āyāt kih bi-mathābih-yi ghayth-i
hāṭil az samāʾ-i ghayb dar kull-i ḥīn nāzil ast kāfir shudih-and qul
a-tuʾminūna bi-baʿḍi al-kitābi wa-takfurūna bi-baʿḍ fa-waylun lakum yā
maʿshara al-ẓālimīn

If you have turned from the left hand of vain imagining to the right
hand of certitude, and have received your portion of the river of the
recognition of the Beauty of the All-Merciful that flows in the paradise
of inner meanings, say: O people of the Bayān! What religion do you
profess, and on what path do you stand? If they say, ``We believe in the
Point of the Bayān---may the spirit of all besides Him be sacrificed for
Him,'' say: By what proof and evidence have you become certain of that
Sovereign of the path and acknowledged Him? If they say, ``We have
recognized Him through His own Self,'' say: This is a blatant lie, for
you have not yet recognized your own selves, much less the Self of God
who sustains all things. Moreover, if such a claim is heard from these
souls, why have they not acknowledged the manifest Self of God, who
shines like the sun? And if they say, ``We have believed through the
revealed verses,'' why have they disbelieved in these verses, which
descend at every moment from the heaven of the unseen like pouring rain?
Say: Do you believe in part of the Book and disbelieve in part? Woe to
you, O company of wrongdoers!

\subsection{58}\label{section-1613}

\begin{otherlanguage}{persian}

و من‌دون این دو مقام ظهورات قدرتیّه و شئونات الوهیّه که عالم را احاطه نموده
بشأنی که برای نفسی مجال اعراض نمانده الّا بأن یکون معرضاً بکلّه عن اللّه و
انبیائه و اصفیائه و اودّائه قریب بیست سنه میشود که این عبد آنی بر بستر
راحت نیاسوده و در کلّ حین در ارتفاع امر اللّه بنفسه کوشیده قسم بسلطان
لایعرف که از اوّل ابداع تا حال چنین قدرتی ظاهر نشده که نفسی وحده اعلام
قدرت بر {[}اقنان{]} عظمت مرتفع نماید و مع‌ذلک این مشرکین بر جمال مبین ربّ
العالمین وارد آورده‌اند آنچه را که الآن روح الأمین و هیاکل علّیّین نوحه و
ندبه مینمایند

\end{otherlanguage}

Va min-dūn-i īn du maqām ẓuhūrāt-i qudratiyyih va shuʾūnāt-i ulūhiyyih
kih ʿālam rā iḥāṭih namūdih bi-shaʾnī kih barā-yi nafsī majāl-i iʿrāḍ
namāndih illā bi-an yakūna muʿriḍan bi-kullihi ʿan Allāhi wa-anbiyāʾihi
wa-aṣfiyāʾihi wa-awiddāʾihi qarīb-i bīst sanih mīshavad kih īn ʿabd ānī
bar bistar-i rāḥat nayāsūdih va dar kull-i ḥīn dar irtifāʿ-i Amr Allāh
bi-nafsih kūshīdih qasam bi-Sulṭān-i lā-yuʿraf kih az avval-i ibdāʿ tā
ḥāl chunīn qudratī ẓāhir nashudih kih nafsī vaḥdahu aʿlām-i qudrat bar
{[}aqnān{]}-i ʿaẓamat murtafiʿ namāyad va maʿa-dhālik īn mushrikīn bar
Jamāl-i mubīn-i Rabb al-ʿālamīn vārid āvardih-and ānchih rā kih al-ān
Rūḥ al-Amīn va hayākil-i ʿilliyyīn nawḥih va nudbih mīnamāyand

Besides these two stations, there are manifestations of power and signs
of divinity that have encompassed the world in such a way that no soul
has any scope to turn away, unless it turns wholly away from God, His
Prophets, His chosen ones, and His loved ones. For nearly twenty years
this servant has not rested for a moment upon a bed of ease, and has at
all times striven in person to exalt the Cause of God. I swear by the
unknowable Sovereign that, from the beginning of creation until now, no
such power has appeared: a single soul, alone, raising the banners of
power upon the summits of grandeur. Yet these idolaters have inflicted
upon the manifest Beauty of the Lord of the worlds what causes the
Faithful Spirit and the embodiments of the highest heaven to lament and
mourn even now.

\subsection{59}\label{section-1614}

\begin{otherlanguage}{persian}

و اگر میگویند این آیات بدیعه از فطرت الهیّه نازل نشده چنانچه مشرکین قبل
در احیان ظهورات شموس حقیقت این سخن را گفته‌اند بگو الواح منزله که از سحاب
عزّ ربّ العالمین در سنهٴ ستّین نازل شده موجود و این آیات بدیعه که از سماء
قدس ابهی نازل شده حاضر و مشهود هر دو را نزد عدّه‌ئی از اهل قلوب صافیه و
ابصار حدیده و انفس زکیّه و آذان واعیه تلاوت مینمائیم تا نغمات اللّه و
روایح قدسش از قمیص کلمات بدیعش استشمام شود تاللّه الحقّ برایحة کلمة من تلک
الکلمات لتهبّ رایحة اللّه المهیمن القیّوم ولکن کلّ ناس بزکام مبتلا گشته‌اند و
هم لا یجدون ابداً الّا من شآء ربّک العزیز المحبوب اگرچه اعراض معرضین و
افترای مفترین بمقامی رسیده که قلم و بیان هر دو از ذکر بازمانده و مع‌ذلک
انّا نطعم العباد من مائدة العلمیّة الطّریّة الأبدیّة القدمیّة الالهیّة لوجه
اللّه و ما نرید منهم جزآء و لا شکورا

\end{otherlanguage}

Va agar mīgūyand īn āyāt-i badīʿih az fiṭrat-i ilāhiyyih nāzil nashudih
chunānchih mushrikīn-i qabl dar aḥyān-i ẓuhūrāt-i shumūs-i ḥaqīqat īn
sukhan rā guftih-and bigū alvāḥ-i munzalih kih az saḥāb-i ʿizz-i Rabb
al-ʿālamīn dar sanih-yi sittīn nāzil shudih mawjūd va īn āyāt-i badīʿih
kih az samāʾ-i quds-i Abhā nāzil shudih ḥāḍir va mashhūd har du rā
nazd-i ʿiddih-ī az ahl-i qulūb-i ṣāfiyih va abṣār-i ḥadīdih va anfus-i
zakiyyih va ādhān-i vāʿiyih tilāvat mīnamāyīm tā naghamāt Allāh va
ravāyiḥ-i qudsash az qamīṣ-i kalimāt-i badīʿash istishmām shavad
ta-Allāhi al-Ḥaqq bi-rāʾiḥati kalimatin min tilka al-kalimāt la-tahubbu
rāʾiḥatu Allāhi al-Muhaymini al-Qayyūm valīkin kull-i nās bi-zukām
mubtalā gashtih-and wa-hum lā yajidūna abadan illā man shāʾa rabbuka
al-ʿAzīzu al-Maḥbūb agarchih iʿrāḍ-i muʿriḍīn va iftirā-yi muftarīn
bi-maqāmī rasīdih kih qalam va bayān har du az dhikr bāz-māndih va
maʿa-dhālik innā nuṭʿimu al-ʿibāda min māʾidati al-ʿilmiyyati
al-ṭariyyati al-abadiyyati al-qidamiyyati al-ilāhiyyati li-wajhi Allāhi
wa-mā nurīdu minhum jazāʾan wa-lā shukūrā

If they say that these wondrous verses have not descended from the
divine nature, just as the idolaters of the past said when the Suns of
Truth appeared, say: The revealed Tablets that descended from the cloud
of the glory of the Lord of the worlds in the year sixty are available,
and these wondrous verses that have descended from the heaven of the
holiness of the Most Glorious are present and visible. We will recite
both before a number of people with pure hearts, keen sight, stainless
souls, and attentive ears, so that the melodies of God and His holy
fragrances may be perceived from the garment of His wondrous words. By
God, the Truth! With the fragrance of a single one of those words, the
fragrance of God, the All-Protecting, the Self-Subsisting, will
assuredly waft. Yet all the people have been afflicted with a cold, and
they will never perceive it, except those whom your Lord, the Almighty,
the Beloved, wills. Although the rejection of those who turn away and
the slander of the slanderers have reached such a point that pen and
utterance alike are unable to recount them, nevertheless We feed the
servants, for the sake of God, from the fresh, eternal, ancient, divine
table of knowledge, and We desire from them neither reward nor thanks.

\subsection{60}\label{section-1615}

\begin{otherlanguage}{persian}

فوالّذی نطق فی صدری که هیچ یک از ملل قبل بحجبات اهل بیان مشاهده نمیشوند
چه که در جمیع کتب قبل ذکر ظهورات احدیّه بتلویح ذکر شده مثلاً در توریة در
بعضی از مواضع ذکر شده که اگر نفسی بیاید و دعوی نبوّت کند کاذب است چه که
الهی جز اله موسی نبوده و رسولی جز من مبعوث نخواهد شد و احکام الهی جز ما
نزّل فی التّوریة نخواهد آمد و در یک موضع بتلویح اشاره بظهور بعد فرموده و
بشأنی بحجب و استار ذکر شده که اکثری از عباد از عرفانش عاجزند حال در این
صورت اگر امّت او از مشارق احدیّه و مظاهر الهیّه محتجب مانند فی‌الجمله عذری
در دست دارند که بآن متعذّر شوند که معضلات کلمات الهی را ادراک ننمودیم لذا
از منبع کوثر عرفان جمال رحمن محروم ماندیم

\end{otherlanguage}

Fa-wa-lladhī naṭaqa fī ṣadrī kih hīch yak az milal-i qabl bi-ḥujubāt-i
ahl-i Bayān mushāhidih namīshavand chih kih dar jamīʿ-i kutub-i qabl
dhikr-i ẓuhūrāt-i aḥadiyyih bi-talvīḥ dhikr shudih mathalan dar Tawrāt
dar baʿḍī az mavāḍiʿ dhikr shudih kih agar nafsī biyāyad va daʿvā-yi
nubuwwat kunad kādhib ast chih kih ilāhī juz Ilāh-i Mūsā nabūdih va
rasūlī juz man mabʿūth nakhvāhad shud va aḥkām-i ilāhī juz mā nuzzila fī
al-Tawrāt nakhvāhad āmad va dar yak mawḍiʿ bi-talvīḥ ishārih bi-ẓuhūr-i
baʿd farmūdih va bi-shaʾnī bi-ḥujub va astār dhikr shudih kih aktharī az
ʿibād az ʿirfānash ʿājizand ḥāl dar īn ṣūrat agar ummat-i ū az
mashāriq-i aḥadiyyih va maẓāhir-i ilāhiyyih muḥtajib mānand fī-l-jumlih
ʿudhrī dar dast dārand kih bi-ān mutaʿadhdhir shavand kih muʿḍalāt-i
kalimāt-i ilāhī rā idrāk nanamūdīm lidhā az manbaʿ-i kawthar-i ʿirfān-i
Jamāl-i Raḥmān maḥrūm māndīm

By Him who has spoken within My breast! None of the peoples of the past
are seen to be veiled as the people of the Bayān are. For in all the
former Books the manifestations of divine unity were mentioned by
allusion. For example, in certain passages of the Torah it is stated
that if anyone comes and claims prophethood, he is a liar, for there has
been no God but the God of Moses, no messenger but Me will be raised up,
and no divine ordinances will come other than what has been sent down in
the Torah. In one passage, He alludes indirectly to the subsequent
Revelation, but it is mentioned beneath such veils and coverings that
most people are unable to recognize it. In these circumstances, if His
people remain veiled from the Daysprings of divine unity and the
Manifestations of God, they have some measure of excuse to offer: ``We
did not comprehend the difficulties in the divine words, and thus
remained deprived of the source of the river of the recognition of the
Beauty of the All-Merciful.''

\subsection{61}\label{section-1616}

\begin{otherlanguage}{persian}

و هم‌چنین در انجیل روح القدس بنغماتی در علایم ظهور تغنّی و تکلّم فرموده که
ادراک آن هر نفسی را ممکن نه الّا المنقطعین چه که برموزات خفیّه و اشارات
دقیقه بیان شده چنانچه همان عبارات از قلم عزّ باقیه در رسایل فارسیّه مسطور
گشت فانظروا الیها لعلّ تجدون در این صورت این طایفه هم اگر متعذّر شوند
بآنکه عقول و ادراک ما قاصر بود از عرفان این کلمات مرموزهٴ معضله شاید که
بعضی از عباد بپذیرند

\end{otherlanguage}

Va ham-chunīn dar Injīl Rūḥ al-Qudus bi-naghamātī dar ʿalāʾim-i ẓuhūr
taghannī va takallum farmūdih kih idrāk-i ān har nafsī rā mumkin nah
illā al-munqaṭiʿīn chih kih bi-rumūzāt-i khafiyyih va ishārāt-i daqīqih
bayān shudih chunānchih hamān ʿibārāt az qalam-i ʿizz-i bāqiyih dar
rasāʾil-i fārsiyyih masṭūr gasht fa-nẓurū ilayhā laʿallakum tajidūn dar
īn ṣūrat īn ṭāʾifih ham agar mutaʿadhdhir shavand bi-ānkih ʿuqūl va
idrāk-i mā qāṣir būd az ʿirfān-i īn kalimāt-i marmūzih-yi muʿḍalih
shāyad kih baʿḍī az ʿibād bipadhīrand

Likewise, in the Gospel the Holy Spirit has sung and spoken of the signs
of the Revelation in melodies that no soul can comprehend except those
who are detached, for they were expressed in hidden symbols and subtle
allusions. Those very passages have been recorded by the Pen of
everlasting glory in Persian treatises. Look at them, that perhaps you
may discover. In these circumstances, if this community too were to
offer the excuse, ``Our minds and understanding were incapable of
recognizing these difficult, cryptic words,'' some people might accept
it.

\subsection{62}\label{section-1617}

\begin{otherlanguage}{persian}

و هم‌چنین در فرقان که همهٴ شما دیده‌اید که در علایم ظهور بعد بچه شأن و
بیان از سماء سبحان نازل شده مثلاً از جملۀ هل ینظرون الّا ان یأتیهم اللّه فی
ظلل من الغمام و هم‌چنین یوم تأتی السّمآء بدخان مبین و هم‌چنین انفطار سماء و
انشقاق ارض و اندکاک جبال و انسجار بحار و اقامهٴ اموات از قبور و دمیدن در
صور و اشراق شمس عن جهة الغرب و ارتفاع صیحة بین السّموات و الأرض و امثال
این کلمات که در کتاب اللّه مذکور است

\end{otherlanguage}

Va ham-chunīn dar Furqān kih hamih-yi shumā dīdih-īd kih dar ʿalāʾim-i
ẓuhūr-i baʿd bi-chih shaʾn va bayān az samāʾ-i Subḥān nāzil shudih
mathalan az jumlih hal yanẓurūna illā an yaʾtiyahumu Allāhu fī ẓulalin
min al-ghamām va ham-chunīn yawma taʾtī al-samāʾu bi-dukhānin mubīn va
ham-chunīn infiṭār-i samāʾ va inshiqāq-i arḍ va indikāk-i jibāl va
insijār-i biḥār va iqāmih-yi amvāt az qubūr va damīdan dar ṣūr va
ishrāq-i shams ʿan jihati al-gharb va irtifāʿ-i ṣayḥih bayna al-samāwāti
wa-l-arḍ va amthāl-i īn kalimāt kih dar Kitāb Allāh madhkūr ast

Likewise, in the Qurʾān, all of you have seen in what manner and
language the signs of the subsequent Revelation have been sent down from
the heaven of the All-Glorious. Among them, for example, are: ``Do they
await anything but that God should come to them in canopies of cloud?''
and likewise, ``The day when heaven brings forth manifest smoke''; and
likewise the rending of heaven, the splitting of the earth, the crushing
of the mountains, the kindling of the seas, the raising of the dead from
their graves, the blowing of the trumpet, the rising of the sun from the
west, the sounding of a cry between the heavens and the earth, and
similar expressions recorded in the Book of God.

\subsection{63}\label{section-1618}

\begin{otherlanguage}{persian}

و هم‌چنین از ذکر خاتم النّبیّین {[}ص{]} که اصرح کلمات فرقانیّه است با این
عبارات صعبهٴ مستصعبه و اشارات دقیقهٴ خفیّه اگر عباد از شریعهٴ ربّ الایجاد
و عرفان نفس او فی المعاد محروم مانند میتوان گفت که ادراک بیانات الهیّه
ننموده‌اند و از معانی کلمات ربّانیّه غافل شده‌اند چه که بفهم عباد نزدیک
نبوده اگرچه جمیع این معاذیر در حین ظهور غیر مقبول بوده و خواهد بود چه که
در هر ظهور بنفس ظهور و بما یظهر من عنده حجّت بر کلّ من فی السّموات و الأرض
بالغ میشود

\end{otherlanguage}

Va ham-chunīn az dhikr-i Khātam al-Nabiyyīn {[}ṣ{]} kih aṣraḥ-i
kalimāt-i furqāniyyih ast bā īn ʿibārāt-i ṣaʿbih-yi mustaṣʿabih va
ishārāt-i daqīqih-yi khafiyyih agar ʿibād az sharīʿih-yi Rabb al-Ījād va
ʿirfān-i nafs-i ū fī al-maʿād maḥrūm mānand mītavān guft kih idrāk-i
bayānāt-i ilāhiyyih nanamūdih-and va az maʿānī-yi kalimāt-i rabbāniyyih
ghāfil shudih-and chih kih bi-fahm-i ʿibād nazdīk nabūdih agarchih
jamīʿ-i īn maʿādhīr dar ḥīn-i ẓuhūr ghayr-i maqbūl būdih va khvāhad būd
chih kih dar har ẓuhūr bi-nafs-i ẓuhūr va bi-mā yaẓharu min ʿindihi
ḥujjat bar kulli man fī al-samāwāti wa-l-arḍ bāligh mīshavad

Likewise, there is the mention of the Seal of the Prophets {[}God bless
him{]}, which is among the most explicit expressions of the Qurʾān. With
these exceedingly difficult passages and subtle, hidden allusions, if
people remain deprived of the life-giving waters of the Lord of creation
and of recognizing His own Self at the Return, it can be said that they
have not comprehended the divine utterances and have remained unaware of
the meanings of the words of the Lord, for these were not readily
accessible to human understanding. Nevertheless, all such excuses have
always been, and will always be, unacceptable at the time of Revelation,
for in every Revelation the proof is made conclusive for all who are in
the heavens and on earth through the Revelation itself and through what
appears from Him.

\subsection{64}\label{section-1619}

\begin{otherlanguage}{persian}

و بر هر نفسی لازم است که مرآت قلب را از کلّ آنچه در دست ناس بوده طاهر
نماید و بعد از تطهیر و اقبال البتّه انوار شمس مجلّی بر قلبش تجلّی فرماید و
اگر حجّت الهی در احیان ظهور بالغ نباشد تکلیف از کلّ ساقط میگردد مثلاً در
اهل فرقان ملاحظه نما که هر نفسی که قلب را از اشارات کلماتیّه مطهّر ننمود
بعرفان نقطهٴ بیان فایز نشد چه که ذکر ختمیّت که در کتاب مذکور است از کلمات
محکمهٴ فرقان بوده با اثبات این کلمه و تحقّق معنی ظاهر آن در قلب هرگز
اقرار ننماید باینکه نبیّ از اوّل لا اوّل بوده و الی آخر لا آخر خواهد آمد
چنانچه در بیان مذکور است طهّروا قلوبکم عمّا شهدتم لتشهدوا ما لا شهد احد من
العالمین

\end{otherlanguage}

Va bar har nafsī lāzim ast kih mirʾāt-i qalb rā az kull-i ānchih dar
dast-i nās būdih ṭāhir namāyad va baʿd az taṭhīr va iqbāl al-battih
anvār-i shams-i mujallī bar qalbash tajallī farmāyad va agar ḥujjat-i
ilāhī dar aḥyān-i ẓuhūr bāligh nabāshad taklīf az kull sāqiṭ mīgardad
mathalan dar ahl-i Furqān mulāḥiẓih namā kih har nafsī kih qalb rā az
ishārāt-i kalimātiyyih muṭahhar nanamūd bi-ʿirfān-i Nuqṭih-yi Bayān
fāʾiz nashud chih kih dhikr-i khatmiyyat kih dar kitāb madhkūr ast az
kalimāt-i muḥkamih-yi Furqān būdih bā ithbāt-i īn kalimih va taḥaqquq-i
maʿnī-yi ẓāhir-i ān dar qalb hargiz iqrār nanamāyad bi-īnkih nabī az
avval-i lā avval būdih va ilā ākhir-i lā ākhir khvāhad āmad chunānchih
dar Bayān madhkūr ast ṭahhirū qulūbakum ʿammā shahidtum li-tashhadū mā
lā shahida aḥadun min al-ʿālamīn

Every soul must purify the mirror of the heart of all that people
possess. After this purification and turning toward Him, the rays of the
effulgent Sun will assuredly shine upon that heart. If God's proof were
not conclusive at the times of Revelation, all would be absolved of
obligation. Consider, for example, the people of the Qurʾān: whoever did
not purify the heart of verbal allusions failed to attain recognition of
the Point of the Bayān. For the mention of finality recorded in the Book
was among the unequivocal statements of the Qurʾān. With this statement
affirmed and its outward meaning established in the heart, a person will
never acknowledge that Prophets have existed from the beginning that has
no beginning and will continue to come until the end that has no end.
Thus it is stated in the Bayān: ``Purify your hearts of all that you
have witnessed, so that you may witness what no one in all the worlds
has witnessed.''

\subsection{65}\label{section-1620}

\begin{otherlanguage}{persian}

در حین ظهور باید چشم از کلّ برداشت و بطرف اللّه ناظر شد که من دون ذلک لن
یفوز احد باللّه العلیّ العظیم ای بندگان هوی بشنوید نغمات قدس بقا را و بمقرّ
اعرفوا اللّه باللّه بشتابید و از دونش منقطع شوید این بیانات انبیا که در
علامات ظهور فرموده‌اند

\end{otherlanguage}

Dar ḥīn-i ẓuhūr bāyad chashm az kull bardāsht va bi-ṭaraf Allāh nāẓir
shud kih min dūni dhālika lan yafūza aḥadun bi-Allāhi al-ʿAliyyi
al-ʿAẓīm ay bandigān-i havā bishnavīd naghamāt-i quds-i baqā rā va
bi-maqarr-i iʿrifū Allāha bi-Allāh bishtābīd va az dūnash munqaṭiʿ
shavīd īn bayānāt-i anbiyā kih dar ʿalāmāt-i ẓuhūr farmūdih-and

At the time of Revelation, one must turn one's eyes away from everything
and look toward God, for otherwise no one will attain God, the Most
High, the Most Great. O servants of desire! Hear the melodies of
everlasting holiness, hasten to the abode of ``Recognize God through
God,'' and detach yourselves from all besides Him. These utterances of
the Prophets concerning the signs of the Revelation

\subsection{66}\label{section-1621}

\begin{otherlanguage}{persian}

و امّا نقطهٴ بیان روح من فی لجج الأسمآء فداه جمیع این بیانات را مرتفع
فرموده و حجبات را بالمرّه خرق نموده و جمیع این کلمات معضله را بنفس ظهور و
ما یظهر من عنده تمام نموده‌اند و معنی فرموده‌اند و بشأنی ذکر این ظهور عزّ
احدیّه را فرموده که برای احدی مجال توقّف نمانده تا چه رسد باعراض و جمیع
بیان و ما نزّل فیه را منوط بعرفان آن شمس عزّ باقیه فرموده و میفرمایند که
مباد در حین ظهور به بیان از منزل آن محتجب مانند و بکلمات آن تمسّک جسته از
سلطان احدیّه محروم مانید و مخصوص میفرمایند ای حروفات و مرایا شما بقول من
حجّت شده‌اید مبادا در حین ظهور بر جمال مختار استکبار نمائید

\end{otherlanguage}

Va ammā Noqṭe-ye Bayān, rūḥu man fī lujaji al-asmāʾ fidāhu, jamīʿ-e īn
bayānāt rā mortafeʿ farmūde va ḥojobāt rā bel-marre kharq namūde va
jamīʿ-e īn kalemāt-e moʿḍele rā be-nafs-e ẓohūr va mā yaẓharu min
ʿindihi tamām namūde-and va maʿnī farmūde-and va be-shaʾnī dhekr-e īn
ẓohūr-e ʿezz-e aḥadiyye rā farmūde ke barā-ye aḥadī majāl-e tavaqqof
namānde tā che resad be-eʿrāḍ, va jamīʿ-e Bayān va mā nuzzila fīhi rā
manūṭ be-ʿerfān-e ān shams-e ʿezz-e bāqiye farmūde va mīfarmāyand ke
mabād dar ḥīn-e ẓohūr be-Bayān az monzel-e ān moḥtajeb mānand va
be-kalemāt-e ān tamassok joste az solṭān-e aḥadiyye maḥrūm mānīd va
makhṣūṣ mīfarmāyand: Ey ḥorūfāt va marāyā, shomā be-qawl-e man ḥojjat
shode-īd, mabādā dar ḥīn-e ẓohūr bar jamāl-e mokhtār estekbār namāyīd.

As for the Point of the Bayān---may the souls of all who dwell in the
surging seas of names be sacrificed for Him---He has set aside all these
statements, utterly torn through the veils, and brought all these
difficult expressions to their fulfillment and explained their meaning
through the Manifestation Himself and whatever proceeds from Him. He has
spoken of this Revelation of the glory of divine unity in such a way
that no one has any room to hesitate, much less to turn away. He has
made the whole Bayān and all that was revealed in it dependent upon
recognition of that Sun of everlasting glory. He says that, at the time
of the Revelation, they must not let the Bayān veil them from its
Revealer, and that you must not cling to its words and remain deprived
of the Sovereign of oneness. He specifically says: ``O Letters and
Mirrors! You have become proofs through My word. Beware lest, at the
time of the Revelation, you act arrogantly toward the sovereign
Beauty.''

\subsection{67}\label{section-1622}

\begin{otherlanguage}{persian}

و با آنکه در مستغاث اشراق شمس حقیقت را باسم من یظهر عزّ اسمه وعده
فرموده‌اند مع‌ذلک میفرمایند که اگر در ساعت دیگر ظاهر شود احدی را نمیرسد که
لم او بم بگوید چه که آن سلطان امکان لم‌یزل مختار بوده و خواهد بود و اگر
احدی بغیر آنچه ذکر شده قائل شود حقّ جلّ ذکره را مختار ندانسته پناه میبریم
بحقّ از چنین توهّمات باطله و جمیع حدود و حجب و اشاره را از این مقام
برداشته‌اند و آن جمال قدم ازلاً و ابداً مقدّس از حدود و اشارهٴ دون خود بوده
و احدی بر کیفیّت ظهور اطّلاع نداشته و ندارد لن یحیط بعلمه احد و انّه بکلّ
شیء علیم

\end{otherlanguage}

Va bā ānke dar Mostaghāth eshrāq-e shams-e ḥaqīqat rā be-esm-e Man
Yaẓhar, ʿazza ismuhu, vaʿde farmūde-and, maʿa-dhālek mīfarmāyand ke agar
dar sāʿat-e dīgar ẓāher shavad aḥadī rā nemīrasad ke lima aw bima
begūyad, che ke ān solṭān-e emkān lam-yazal mokhtār būde va khvāhad būd
va agar aḥadī be-gheyr-e ānche dhekr shode qāʾel shavad Ḥaqq, jalla
dhikruhu, rā mokhtār nadāneste. Panāh mībarīm be-Ḥaqq az chonīn
tavahhomāt-e bāṭele. Va jamīʿ-e ḥodūd va ḥojob va eshāre rā az īn maqām
bardāshte-and va ān Jamāl-e Qedam azalan va abadan moqaddas az ḥodūd va
eshāre-ye dūn-e khod būde va aḥadī bar keyfiyyat-e ẓohūr eṭṭelāʿ
nadāshte va nadārad. Lan yuḥīṭa bi-ʿilmihi aḥadun wa-innahu bi-kulli
shayʾin ʿalīm.

Although He has promised that the Sun of Truth will shine forth in
Mustaghāth under the name of Him Who Will Appear---exalted be His
name---He nevertheless says that, should He appear in the next hour, no
one has the right to ask ``Why?'' or ``By what means?'' For that
Sovereign of existence has always been, and will always be, free to
choose. Anyone who maintains anything other than what has been stated
has not acknowledged the freedom of God---exalted be His mention---to
choose. We seek refuge in God from such vain imaginings. He has removed
all limitations, veils, and allusions from this station. That Ancient
Beauty has, from eternity and forever, been sanctified above the
limitations and allusions of anyone other than Himself. No one has ever
known, or now knows, the manner of His manifestation. No one will
encompass His knowledge, and He knows all things.

\subsection{68}\label{section-1623}

\begin{otherlanguage}{persian}

و میفرمایند من اوّل ما یطلع شمس البهآء الی ان یغرب خیر فی کتاب اللّه عن کلّ
اللّیل ان انتم تدرکون ما خلق اللّه من شیء الّا لیومئذ اذ کلّ للقآء اللّه ثمّ
رضائه یعملون و در این مقام میفرمایند و لقد قرب الزّوال و انّکم انتم ذلک
الیوم لا تعرفون و من یکن لقائه ذات لقائی لا ترضوا له ما لا یرضی نفس لنفس
الی آخر قوله عزّ و جلّ ای سمع امکان صاحب سمعی مشهود نه تو بشنو کلمات
بالغهٴ ربّانیّه را که تصریحاً من غیر تأویل اخبار فرموده در آن ایّام که شمس
طالع است و نزدیک است در وسط زوال اشراق فرماید ولکن شما ای ملأ بیان آن
یوم عارف بآن جمال سبحان نخواهید شد و تصریح باین اسم هم فرموده‌اند بقوله
عزّ و جلّ من اوّل ما تطلع شمس البهآء الی ان یغرب خیر فی کتاب اللّه الی آخر
بیانه جلّ و عزّ که احدی شبهه ننماید و بدون ما قضی اللّه تفسیر ننماید مع‌ذلک
کلّ بحجبات نفس و هوی از ادراک شمس بها محتجب مانده‌اند و به ما امرهم
الشّیطان متمسّک و کاش باحتجاب کفایت مینمودند بلکه بعضی بتحریف این کلمات
مظهر اسماء و صفات مشغول شده‌اند و این ظلمی است که در ابداع فوق آن ظاهر
نشده فویل للظّالمین من عذاب یوم عظیم

\end{otherlanguage}

Va mīfarmāyand: Min awwali mā yaṭluʿu shamsu al-Bahāʾ ilā an yaghruba
khayrun fī kitābi Allāh ʿan kulli al-layl in antum tudrikūn. Mā khalaqa
Allāhu min shayʾin illā li-yawmiʾidh idh kullun li-liqāʾi Allāh thumma
riḍāʾihi yaʿmalūn. Va dar īn maqām mīfarmāyand: Wa-laqad qaruba al-zawāl
wa-innakum antum dhālika al-yawma lā taʿrifūn. Wa-man yakun liqāʾuhu
dhāta liqāʾī lā tarḍaw lahu mā lā yarḍā nafsun li-nafs. Ilā ākhiri
qawlihi ʿazza wa-jalla. Ey samʿ-e emkān, ṣāḥeb-e samʿī mashhūd na, to
beshnow kalemāt-e bāleghe-ye rabbāniyye rā ke taṣrīḥan min ghayri taʾwīl
ekhbār farmūde dar ān ayyām ke shams ṭāleʿ ast va nazdīk ast dar vasaṭ-e
zavāl eshrāq farmāyad va-lāken shomā ey malaʾ-e Bayān ān yowm ʿāref
be-ān jamāl-e sobḥān nakhvāhīd shod va taṣrīḥ be-īn esm ham farmūde-and
bi-qawlihi ʿazza wa-jalla: Min awwali mā taṭluʿu shamsu al-Bahāʾ ilā an
yaghruba khayrun fī kitābi Allāh. Ilā ākhiri bayānihi jalla wa-ʿazza. Ke
aḥadī shobhe nanamāyad va be-dūn-e mā qaḍā Allāh tafsīr nanamāyad.
Maʿa-dhālek koll be-ḥojobāt-e nafs va havā az edrāk-e Shams-e Bahā
moḥtajeb mānde-and va be-mā amarahumu al-shayṭān motamassek va kāsh
be-eḥtejāb kefāyat mīnamūdand balke baʿḍī be-taḥrīf-e īn kalemāt-e
maẓhar-e asmāʾ va ṣefāt mashghūl shode-and va īn ẓolmī ast ke dar ebdāʿ
fowq-e ān ẓāher nashode. Fa-waylun li-l-ẓālimīn min ʿadhābi yawmin
ʿaẓīm.

He also says: ``From the moment the Sun of Bahāʾ rises until it sets is
better, in the Book of God, than all the night, if you understand. God
has created nothing except for that Day, for all labor to attain the
presence of God and then His good pleasure.'' In this connection He
says: ``Noon has indeed drawn near, yet you do not recognize that Day.
As for Him whose presence is My very presence, do not wish for Him what
no soul wishes for itself''---to the end of His words, exalted and
glorious is He. O ear of existence! No one possessed of hearing is to be
seen; you, then, hear these conclusive divine words. In those days He
explicitly announced, without recourse to interpretation, that the Sun
had risen and was about to shine at the height of noon, yet that you, O
people of the Bayān, would not recognize that all-glorious Beauty on
that Day. He also explicitly stated this name in His words, exalted and
glorious is He: ``From the moment the Sun of Bahāʾ rises until it sets
is better, in the Book of God''---to the end of His utterance, glorious
and exalted is He---so that no one should entertain any doubt or
interpret it otherwise than God has decreed. Nevertheless, all have
remained veiled by self and desire from recognizing the Sun of Bahāʾ and
have clung to what Satan has commanded them. If only they had been
content to remain veiled! Rather, some have occupied themselves with
distorting these words of the Manifestation of the names and attributes.
This is an injustice greater than any that has appeared in creation.
Woe, then, to the wrongdoers, because of the punishment of a tremendous
Day!

\subsection{69}\label{section-1624}

\begin{otherlanguage}{persian}

و بعد از ظهور این شمس عزّ صمدانی بعضی از مرایا توهّم نموده‌اند و رتبهٴ شمسی
ادّعا نموده‌اند ولکن غافل از اینکه در بیان فارسی نصّ فرموده‌اند که اگر مرآت
ادّعای شمسی نماید نزد شمس ظاهر است که شبح او است که او میگوید و هم‌چنین
میفرماید مرایا بنفسها شیئیّت ندارند و در مقام دیگر میفرماید قل ان یا شموس
المرایا انتم الی شمس الحقیقة تنظرون و انّ قیامکم بها لو انتم تتبصّرون کلّکم
کحیتان بالمآء فی البحر تتحرّکون و تحتجبون عن المآء و تسألون عمّا انتم به
قائمون حال ملاحظه نمائید که بشموس مرایا که مرایای اوّلیّه‌اند میفرماید که
شما بشمس حقیقت ناظر باشید چه که وجود و ظهور شما بعنایت او بوده و خواهد
بود و میفرمایند شموس مرایا مثل حیتان در آبند که در بحر حرکت میکنند ولکن
از بحر و ماء محتجبند چنانچه الیوم ملاحظه میشود که مرآت قوم در بحر آیات
حرکت مینماید چنانچه بآیات عزّ صمدانی که از ظهور قبلم نازل شده اثبات خود
مینماید و حرکت و اظهار شأن او از آیات اللّه بوده و مع‌ذلک از جوهر آیات و
منزل آن در این ایّام بالمرّه محتجب مانده در بحر حرکت مینماید و از سلطان
بحر غافل و این بیان در رتبهٴ شموس مرایا از ملکوت بقا نازل تا چه رسد
بمرایائی که در تحت این شموس واقعند

\end{otherlanguage}

Va baʿd az ẓohūr-e īn shams-e ʿezz-e ṣamadānī baʿḍī az marāyā tavahhom
namūde-and va rotbe-ye shamsī eddeʿā namūde-and va-lāken ghāfel az īnke
dar Bayān-e Fārsī naṣṣ farmūde-and ke agar merʾāt eddeʿā-ye shamsī
namāyad nazd-e shams ẓāher ast ke shabaḥ-e ū ast ke ū mīgūyad va
hamchenīn mīfarmāyad marāyā be-nafsehā shayʾiyyat nadārand va dar
maqām-e dīgar mīfarmāyad: Qul in yā shumūsa al-marāyā antum ilā shamsi
al-ḥaqīqa tanẓurūn wa-inna qiyāmakum bihā law antum tatabaṣṣarūn.
Kullukum ka-ḥītānin bi-l-māʾi fī al-baḥri tataḥarrakūn wa-taḥtajibūna
ʿani al-māʾi wa-tasʾalūna ʿammā antum bihi qāʾimūn. Ḥāl molāḥeẓe namāyīd
ke be-shomūs-e marāyā ke marāyā-ye avvaliyye-and mīfarmāyad ke shomā
be-shams-e ḥaqīqat nāẓer bāshīd che ke vojūd va ẓohūr-e shomā
be-ʿenāyat-e ū būde va khvāhad būd va mīfarmāyand shomūs-e marāyā mesl-e
ḥītān dar āband ke dar baḥr ḥarekat mīkonand va-lāken az baḥr va māʾ
moḥtajeband chenānche al-yowm molāḥeẓe mīshavad ke merʾāt-e qowm dar
baḥr-e āyāt ḥarekat mīnamāyad chenānche be-āyāt-e ʿezz-e ṣamadānī ke az
ẓohūr-e qablam nāzel shode ethbāt-e khod mīnamāyad va ḥarekat va eẓhār-e
shaʾn-e ū az āyāt-e Allāh būde va maʿa-dhālek az jowhar-e āyāt va
monzel-e ān dar īn ayyām bel-marre moḥtajeb mānde dar baḥr ḥarekat
mīnamāyad va az solṭān-e baḥr ghāfel va īn bayān dar rotbe-ye shomūs-e
marāyā az malakūt-e baqā nāzel tā che resad be-marāyāʾī ke dar taḥt-e īn
shomūs vāqeʿand.

After the appearance of this Sun of eternal glory, some of the Mirrors
have indulged in vain imaginings and claimed the rank of a sun. Yet they
are heedless that He has explicitly stated in the Persian Bayān that, if
a mirror claims to be a sun, it is clear to the sun that it is its own
reflected image that is speaking. He likewise says that mirrors have no
existence in themselves. Elsewhere He says: ``Say: O suns of the
mirrors! You look toward the Sun of Truth, and your existence depends
upon it, if only you perceive. You are all like fish moving through the
water in the sea, yet you are veiled from the water and ask about that
through which you subsist.'' Now consider how He tells the suns of the
mirrors, who are the foremost Mirrors, to look toward the Sun of Truth,
for their existence and manifestation have depended, and will depend,
upon His favor. He says that the suns of the mirrors are like fish in
the water that move through the sea but are veiled from the sea and its
water. Today, likewise, it can be seen that the Mirror of the people
moves in the ocean of verses, for he establishes his own claim through
the verses of eternal glory revealed by My previous Manifestation. His
movement and the display of his station have derived from the verses of
God. Nevertheless, in these days he has remained utterly veiled from the
Essence of the verses and their Revealer. He moves in the sea but is
heedless of the Sovereign of the sea. This utterance concerning the rank
of the suns of the mirrors has descended from the Kingdom of eternity;
how much more does it apply to the mirrors that stand beneath these
suns!

\subsection{70}\label{section-1625}

\begin{otherlanguage}{persian}

و هم‌چنین میفرماید فانّ یوم ظهوره اعلی الخلق مثل ادناه و انّ اقربکم من یؤمن
به و لا انساب بینکم و لا افتخار الّا بایمانکم به اینست که الیوم مشاهده
میشود که اعلای خلق از ادنی بین یدی اللّه مذکور چه ‌قدر از مظاهر علّیّین که
بسجّین راجع شده‌اند و چه مقدار از مطالع لا که بمقرّ مطهّر الّا وارد شده و
هم‌چنین بجناب آقا سیّد جواد میفرماید لأشکونّ الیک ان یا مرآة جودی عن کلّ
المرایا کلّ بألوانهم الیّ لینظرون ملاحظه نمائید که از مرایا کلّها شکایت
میفرمایند ولکن هذا الجمال یشکو حینئذ عن الّذی اشتهر بالمرآة لأنّه ینظر الیّ
بما عنده لا بما عندی

\end{otherlanguage}

Va hamchenīn mīfarmāyad: Fa-inna yawma ẓuhūrihi aʿlā al-khalqi mithlu
adnāhu wa-inna aqrabakum man yuʾminu bihi wa-lā ansāba baynakum wa-lā
iftikhāra illā bi-īmānikum bihi. Īnast ke al-yowm moshāhede mīshavad ke
aʿlā-ye khalq az adnā bayna yaday Allāh madhkūr. Che qadr az maẓāher-e
ʿelliyyīn ke be-sejjīn rājeʿ shode-and va che meqdār az maṭāleʿ-e lā ke
be-maqarr-e moṭahhar-e illā vāred shode va hamchenīn be-janāb-e Āqā
Seyyed Javād mīfarmāyad: La-ashkuwanna ilayka in yā mirʾāta jūdī ʿan
kulli al-marāyā, kullun bi-alwānihim ilayya la-yanẓurūn. Molāḥeẓe
namāyīd ke az marāyā kollahā shekāyat mīfarmāyand va-lāken hādhā
al-jamālu yashkū ḥīnaʾidhin ʿani alladhī ishtahara bi-l-mirʾāti
li-annahu yanẓuru ilayya bi-mā ʿindahu lā bi-mā ʿindī.

He likewise says: ``On the day of His appearance, the highest of
creation is like the lowest, and the nearest among you is whoever
believes in Him. There are no ties of kinship among you, nor any cause
for pride except your faith in Him.'' Thus today it is seen that the
highest of creation is accounted among the lowest before God. How many
manifestations of the highest heaven have returned to the lowest depths,
and how many dawning-places of negation have entered the sanctified
abode of affirmation! He likewise says to Āqā Siyyid Javād: ``I will
surely complain to you, O mirror of My bounty, of all the mirrors: they
all look toward Me through their own colors.'' Consider how He complains
of all the mirrors. This Beauty, however, now complains of the one who
has become known as the Mirror, for he looks toward Me through what he
possesses, not through what I possess.

\subsection{71}\label{section-1626}

\begin{otherlanguage}{persian}

باری الیوم اگر جمیع مرایا باسم الوهیّت و ربوبیّت و فوق آن بنفس اللّه و ظهور
اللّه و یا ذات اللّه در بیان موسوم باشند کفایت نمینماید مادام که باین ظهور
قدس ربّانی و بطون غیب صمدانی موقن نشوند جمیع معدوم صرف و مفقود بحت
عنداللّه محسوبند چنانچه در ابتدای این الواح من قلم اللّه ثبت شده که فرموده
ایّاک ایّاک یوم ظهوره ان تحتجب بالواحد البیانیّة فانّ ذلک الواحد خلق عنده و
ایّاک ایّاک ان تحتجب بکلمات ما نزّلت فی البیان فانّها کلمات نفسه فی هیکل
ظهوره من قبل تحذیر میفرماید که در حین ظهور مباد بواحد بیان از مظهر امر
محتجب مانی و این واحد نفوسی بوده‌اند که بعد از نقطهٴ بیان احدی بر آن نفوس
مقدّم نبوده مع‌ذلک میفرماید که باین نفوس از حقّ محتجب نمانید و از این بیان
مفهوم میشود که در ظهور بعد بعضی از این حروفات باید موجود باشند و
میفرماید که بآنچه در بیان نازل شده از حقّ محتجب نمانید

\end{otherlanguage}

Bārī al-yowm agar jamīʿ-e marāyā be-esm-e olūhiyyat va robūbiyyat va
fowq-e ān be-Nafs-e Allāh va Ẓohūr-e Allāh va yā Dhāt-e Allāh dar Bayān
mowsūm bāshand kefāyat nemīnamāyad mādām ke be-īn ẓohūr-e qods-e rabbānī
va boṭūn-e ghayb-e ṣamadānī mūqen nashavand jamīʿ maʿdūm-e ṣerf va
mafqūd-e baḥt ʿinda Allāh maḥsūband chenānche dar ebtedā-ye īn alvāḥ min
qalami Allāh thabt shode ke farmūde: Iyyāka iyyāka yawma ẓuhūrihi an
taḥtajiba bi-l-wāḥidi al-bayāniyya fa-inna dhālika al-wāḥida khalqun
ʿindahu wa-iyyāka iyyāka an taḥtajiba bi-kalimāti mā nuzzilat fī
al-Bayān fa-innahā kalimātu nafsihi fī haykali ẓuhūrihi min qabl.
Taḥdhīr mīfarmāyad ke dar ḥīn-e ẓohūr mabād be-Vāḥed-e Bayān az maẓhar-e
amr moḥtajeb mānī va īn Vāḥed nofūsī būde-and ke baʿd az Noqṭe-ye Bayān
aḥadī bar ān nofūs moqaddam nabūde maʿa-dhālek mīfarmāyad ke be-īn nofūs
az Ḥaqq moḥtajeb namānīd va az īn bayān mafhūm mīshavad ke dar ẓohūr-e
baʿd baʿḍī az īn ḥorūfāt bāyad mowjūd bāshand va mīfarmāyad ke be-ānche
dar Bayān nāzel shode az Ḥaqq moḥtajeb namānīd.

In short, even if all the Mirrors are designated in the Bayān by the
names of divinity and lordship, and beyond these as the Self of God, the
Manifestation of God, or the Essence of God, this does not suffice
today. Until they attain certainty in this Revelation of divine holiness
and this inward concealment of the eternal Unseen, they are all
accounted as sheer nothingness and utter nonexistence before God. Thus,
at the beginning of these Tablets, it has been recorded from the Pen of
God that He said: ``Beware, beware lest, on the day of His appearance,
you be veiled by the Vāḥid of the Bayān, for that Vāḥid is but a
creature in His sight. Beware, beware lest you be veiled by the words
revealed in the Bayān, for they are His own words in the temple of His
previous Manifestation.'' He warns you not to let the Vāḥid of the Bayān
veil you from the Manifestation of the Cause at the time of His
appearance. This Vāḥid consisted of souls whom no one preceded in rank
after the Point of the Bayān. Nevertheless, He says that you must not
let these souls veil you from God. From this utterance it is understood
that some of these Letters must be present in the subsequent Revelation.
He also says that you must not let what has been revealed in the Bayān
veil you from God.

\subsection{72}\label{section-1627}

\begin{otherlanguage}{persian}

حال انصاف دهید ای اهل بیان با این بیان اصرح اتمّ کسی میتواند الیوم معارضه
نماید با نفس ظهور که فلان در بیان باسم اللّه مذکور است چگونه میشود الیوم
از مقام خود سلب شود لا فوالّذی ایّد الرّوح بنفسی چه که این اسامی و اذکار و
توصیف بر فرض تسلیم کلّ کلماتی است که در بیان مسطور است و از جمله وصایای
سلطان لایزالی اینست که بآنچه در بیان نازل شده از ظهور اللّه محتجب نشوید

\end{otherlanguage}

Ḥāl enṣāf dehīd ey ahl-e Bayān bā īn bayān-e aṣraḥ-e atamm kasī
mītavānad al-yowm moʿāreḍe namāyad bā nafs-e ẓohūr ke folān dar Bayān
be-esm-e Allāh madhkūr ast chegūne mīshavad al-yowm az maqām-e khod salb
shavad? Lā fa-wa-alladhī ayy ada al-Rūḥa bi-nafsī. Che ke īn asāmī va
adhkār va towṣīf bar farḍ-e taslīm koll kalemātī ast ke dar Bayān masṭūr
ast va az jomle vaṣāyā-ye solṭān-e lā-yazālī īnast ke be-ānche dar Bayān
nāzel shode az Ẓohūr-e Allāh moḥtajeb nashavīd.

Now judge fairly, O people of the Bayān: in the face of this most
explicit and complete statement, can anyone today oppose the
Manifestation Himself, saying, ``So-and-so is mentioned in the Bayān by
the name of God; how can he now be stripped of his station?'' No, by Him
who strengthened the Spirit through My Self! For these names, mentions,
and descriptions, even if all are conceded, are words inscribed in the
Bayān. Among the counsels of the everlasting Sovereign is this: do not
let what has been revealed in the Bayān veil you from the Manifestation
of God.

\subsection{73}\label{section-1628}

\begin{otherlanguage}{persian}

و علاوه بر این قَسَم بشمس عزّ صمدانی که الیوم از افق قدس ربّانی اشراق فرموده
که ذکر آن سلطان قدم خلق بیان را باسم ربوبیّت و یا عبودیّت هر دو در آن ساحت
یکسان است و در سنگ که جماد است و حروف ثلاثهٴ آن سه اسم اعظم الهی را ذکر
فرموده‌اند و متنبّه نموده عباد را که در ملاحظهٴ آن باین اسماء ناظر باشند
تا در مصنوع آیات صانع مشاهده شود و مقتدری که این مراتب را بحجر عنایت
فرموده اگر اخذ نماید لایق است نفسی را که اعتراض نماید تاللّه لن یعترض احد
علی امر اللّه الّا کلّ معتد اثیم اگر ناس از بعث اسماء مطّلع میبودند هرگز
بسلطان یفعل ما یشآء در هیچ امری اعتراض نمینمودند

\end{otherlanguage}

Va ʿalāve bar īn qasam be-shams-e ʿezz-e ṣamadānī ke al-yowm az ofoq-e
qods-e rabbānī eshrāq farmūde ke dhekr-e ān solṭān-e qedam khalq-e Bayān
rā be-esm-e robūbiyyat va yā ʿobūdiyyat har do dar ān sāḥat yeksān ast
va dar sang ke jamād ast va ḥorūf-e thalāthe-ye ān se esm-e aʿẓam-e
elāhī rā dhekr farmūde-and va motanabbeh namūde ʿebād rā ke dar
molāḥeẓe-ye ān be-īn asmāʾ nāẓer bāshand tā dar maṣnūʿ āyāt-e ṣāneʿ
moshāhede shavad va moqtaderī ke īn marāteb rā be-ḥajar ʿenāyat farmūde
agar akhdh namāyad lāyeq ast nafsī rā ke eʿterāḍ namāyad? Ta-Allāh lan
yaʿtariḍa aḥadun ʿalā amri Allāh illā kullu muʿtadin athīm. Agar nās az
baʿth-e asmāʾ moṭṭaleʿ mībūdand harg ez be-solṭān-e yafʿalu mā yashāʾ
dar hīch amrī eʿterāḍ nemīnamūdand.

Moreover, I swear by the Sun of eternal glory that has today shone from
the horizon of divine holiness: whether that Ancient Sovereign refers to
the creatures of the Bayān by the name of lordship or of servitude, both
are alike in that court. In the stone, which is inanimate, and in the
three letters of its name, He has mentioned three supreme divine names.
He has instructed His servants to look to these names when contemplating
it, so that the signs of the Maker may be seen in what is made. If the
Almighty who has bestowed these ranks upon a stone should take them
away, is it fitting for any soul to object? By God, none will object to
the Cause of God except a sinful transgressor. Had people known of the
raising up of names, they would never have objected in any matter to the
Sovereign who does whatever He wills.

\subsection{74}\label{section-1629}

\begin{otherlanguage}{persian}

اینست که بعضی از ناس اسماً مبعوث میشوند و در ملکوت اسماء معروف ولکن
کینونتشان ابداً از سجّین نفس و هوی عروج ننموده لحکمة لا یعلمها الّا اللّه و
ربّ شهرة لا اصل لها و بعضی بکینونت مبعوث شده من‌دون اسم چنانچه حکایت خضر
را استماع نموده که احدی در آن زمان بر علوّ مقام او مطّلع نه الّا اللّه
چنانچه امرش از موسی مع علوّ مقامه و سموّ قدره مستور بوده و بعد از حضور در
خدمتش بخطاب انّک لن تستطع معی صبراً مخاطب شده و اگرچه در کتاب باسم خضر
مذکور شده ولکن ما عرفه احد الّا اللّه و عندنا علمه فی کتاب مبین و بعضی
بکینونت و ذات و نفس و روح و اسم مبعوث میشوند اذاً فاشهد و قل سبحان اللّه
احسن المبدعین و اگر مراتب بعث بتمامه ذکر شود ملاحظه میشد که این مفتخرین
باسماء در چه مقام و رتبه مذکورند ولکن امسکنا القلم الی ان یشآء اللّه و
انّه مبیّن کلّ امر فی کتاب مبین

\end{otherlanguage}

Īnast ke baʿḍī az nās esman mabʿūth mīshavand va dar malakūt-e asmāʾ
maʿrūf va-lāken kaynūnatashān abadan az sejjīn-e nafs va havā ʿorūj
nanamūde li-ḥikmatin lā yaʿlamuhā illā Allāh wa-rubba shuhratin lā aṣla
lahā va baʿḍī be-kaynūnat mabʿūth shode min dūni ism chenānche ḥekāyat-e
Kheḍr rā estemāʿ namūde ke aḥadī dar ān zamān bar ʿolovv-e maqām-e ū
moṭṭaleʿ na illā Allāh chenānche amrash az Mūsā maʿa ʿuluwwi maqāmihi
wa-sumūwi qadrihi mastūr būde va baʿd az ḥoḍūr dar khedmatash
be-kheṭāb-e innaka lan tastaṭīʿa maʿiya ṣabran mokhāṭab shode va agarche
dar ketāb be-esm-e Kheḍr madhkūr shode va-lāken mā ʿarafahu aḥadun illā
Allāh wa-ʿindanā ʿilmuhu fī kitābin mubīn va baʿḍī be-kaynūnat va dhāt
va nafs va rūḥ va esm mabʿūth mīshavand idhan fa-shhad wa-qul subḥāna
Allāhi aḥsani al-mubdiʿīn va agar marāteb-e baʿth be-tamāme dhekr shavad
molāḥeẓe mīshod ke īn moftakherīn be-asmāʾ dar che maqām va rotbe
madhkūrand va-lāken amsaknā al-qalama ilā an yashāʾa Allāh wa-innahu
mubayyinu kulli amrin fī kitābin mubīn.

Thus some people are raised up in name and are renowned in the kingdom
of names, yet their being has never ascended from the prison of self and
desire, for a wisdom known to God alone. Many a reputation has no
foundation. Others are raised up in their being, without a name, as you
have heard in the account of Khiḍr: no one at that time knew the
loftiness of his station except God. His reality was hidden even from
Moses, despite His lofty station and exalted rank; and after Moses
entered his presence, He was addressed with the words, ``You will not be
able to bear patiently with me.'' Although he is mentioned in the Book
by the name of Khiḍr, no one has known him except God, and with Us is
the knowledge of him in a clear Book. Others are raised up in being,
essence, self, spirit, and name. Bear witness, then, and say: ``Glory be
to God, the best of originators!'' Were the degrees of this raising up
to be set forth in full, it would be seen what station and rank are
assigned to those who pride themselves on names. But We have stayed the
Pen until God wills, and He makes every matter clear in a lucid Book.

\subsection{75}\label{section-1630}

\begin{otherlanguage}{persian}

و دیگر آنکه کلّ اشیاء مکمن اسماء الهیّه و مخزن اسرار صنعیّه بوده و هستند و
در هر کوری از هر شیء ظاهر میفرماید آنچه را اراده فرماید و اخذ میکند آنچه
را بخواهد العجب کلّ العجب الیوم که جمال قدم و شمس اسم اعظم در قطب زوال
مشرق و مستضیء شده مع‌ذلک ناس باسمی از اسماء از انوار جمال و عرفان طلعت
بی‌مثال ذو الجلال محروم گشته‌اند قل تاللّه الحقّ بارادة من قلمه خلقت بحور
الأسمآء و ملکوتها ان انتم من العارفین و لو یأخذ الیوم کفّاً من التّراب و
یبعث منه اسم ما کان و ما یکون لیقدر و انّه لهو السّلطان الملک المقتدر
العزیز القدیر

\end{otherlanguage}

Va dīgar ānke koll-e ashyāʾ makman-e asmāʾ-e elāhiyye va makhzan-e
asrār-e ṣonʿiyye būde va hastand va dar har kawrī az har shayʾ ẓāher
mīfarmāyad ānche rā erāde farmāyad va akhdh mīkonad ānche rā bekhāhad.
Al-ʿajab koll al-ʿajab, emrūz ke Jamāl-e Qedam va Shams-e Esm-e Aʿẓam
dar qoṭb-e zavāl moshreq va mostaḍīʾ shode maʿa-dhālek nās be-esmī az
asmāʾ az anvār-e jamāl va ʿerfān-e ṭalʿat-e bī-methāl-e Dhu'l-Jalāl
maḥrūm gashte-and. Qul ta'llāhi al-Ḥaqq bi-irādatin min qalamihi
khuliqat biḥāru al-asmāʾ wa malakūtuhā in antum min al-ʿārifīn. Wa law
yaʾkhudh al-yawm kaffan min al-turāb wa yabʿath minhu isma mā kāna wa mā
yakūn la-yaqdir wa innahu la-huwa al-Sulṭān al-Malik al-Muqtadir
al-ʿAzīz al-Qadīr.

Moreover, all things have been and remain the repositories of the divine
names and the treasuries of the mysteries of creation. In every
dispensation He manifests from each thing whatever He wills and takes
away whatever He wishes. How utterly astonishing that today, when the
Ancient Beauty and the Sun of the Most Great Name have risen and shine
at the zenith, people have nevertheless, because of one of the names,
deprived themselves of the light of His beauty and of recognizing the
peerless Countenance of the Lord of Glory. Say: By God, the True One!
Through an act of will from His Pen, the seas of names and their kingdom
were created, if you are among those who recognize the truth. Were He
today to take a handful of dust and raise from it the name of all that
has been and all that will be, He would indeed have the power. He is
indeed the Sovereign, the King, the All-Powerful, the Almighty, the
Omnipotent.

\subsection{76}\label{section-1631}

\begin{otherlanguage}{persian}

تاللّه اذاً یبکی لنفسی و بما ورد علیّ عیون النّبیّین و المرسلین فی الرّفیق
الأعلی و ینوحنّ فی انفسهم و یصحنّ فی کینوناتهم ولکنّ النّاس هم فی غفلة و وهم
عظیم و بشأنی اهل بیان تنزّل نموده‌اند که لسان رحمن بذکر امثال این کلمات
مشغول شده و مع‌ذلک از او نمی‌پذیرند اینست شأن این قوم و الیوم هر نفسی که
از ملکوت اسماء ارتقا نجوید ابداً باین فیض اعظم فایز نشود اینست که نقطهٴ
بیان میفرماید که نطفهٴ ظهور بعد اقوی از جمیع اهل بیان بوده و خواهد بود
اذاً تفکّروا فی ذلک یا اولی الألباب فوالّذی قام بنفسه که اگر عباد در همین
کلمه تأمّل نمایند خود را مستغنی از دون اللّه مشاهده نمایند و در اثبات این
امر بدیع بجواب و سؤال محتاج نخواهند بود و اللّه یقول الحقّ ولکنّ النّاس لا
یسمعون و

\end{otherlanguage}

Ta'llāhi idhan yabkī li-nafsī wa bi-mā warada ʿalayya ʿuyūnu al-nabiyyīn
wa al-mursalīn fī al-rafīq al-aʿlā wa yanūḥanna fī anfusihim wa
yaṣiḥanna fī kaynūnātihim wa lākinna al-nāsa hum fī ghaflatin wa wahmin
ʿaẓīm. Va be-shaʾnī ahl-e Bayān tanazzol namūde-and ke lesān-e Raḥmān
be-dhekr-e amthāl-e īn kalemāt mashghūl shode va maʿa-dhālek az ū
nemīpadhīrand. Īnast shaʾn-e īn qawm va emrūz har nafsī ke az malakūt-e
asmāʾ erteqā najūyad abadan be-īn fayḍ-e aʿẓam fāʾez nashavad. Īnast ke
Noqṭe-ye Bayān mīfarmāyad ke noṭfe-ye ẓohūr-e baʿd aqvā az jamīʿ-e ahl-e
Bayān būde va khāhad būd. Idhan tafakkarū fī dhālika yā ulī al-albāb.
Fa-wa'lladhī qāma bi-nafsihi ke agar ʿebād dar hamīn kaleme taʾammol
namāyand khod rā mostaghnī az dūn-e Allāh moshāhede namāyand va dar
ethbāt-e īn amr-e badīʿ be-javāb va soʾāl moḥtāj nakhāhand būd. Wa
Allāhu yaqūlu al-ḥaqq wa lākinna al-nāsa lā yasmaʿūn va

By God! The eyes of the Prophets and Messengers in the exalted company
weep for Me and for what has befallen Me; they lament within themselves
and cry out in their inmost beings. Yet the people are sunk in
heedlessness and grievous delusion. The people of the Bayān have fallen
so low that the Tongue of the All-Merciful has occupied itself with such
words as these, and still they will not accept them from Him. Such is
the state of this people. Today, whoever does not ascend beyond the
kingdom of names will never attain this most great outpouring of grace.
Thus the Point of the Bayān states that the embryo of the next
Manifestation has been and will remain mightier than all the people of
the Bayān. Reflect on this, then, you who possess understanding. By Him
who stands through His own Self! If the servants were to ponder this one
utterance, they would find themselves independent of all but God and
would have no need of questions and answers to establish this wondrous
Cause. God speaks the truth, but the people do not listen, and

\subsection{77}\label{section-1632}

\begin{otherlanguage}{persian}

هم‌چنین میفرماید قوله عزّ ذکره فانّ مثله جلّ ذکره کمثل الشّمس لو یقابلنّه الی
ما لا نهایة مرایا کلّهنّ لیستعکسنّ عن تجلّی الشّمس فی حدّهم و ان لن یقابلها من
احد فیطلع الشّمس و یغرب و الحجاب للمرایا حال ملاحظه نمائید که از این بیان
قدس ربّانی مستفاد میشود که ممکن است شمس طالع شود و غروب نماید و مرایا در
حجاب باشند مع‌ذلک میتوان گفت بعد از اشراق شمس الشّموس الّتی یطوفنّ فی حولها
شموس لا یعلم عدّتهنّ احد الّا نفس اللّه العالم العلیم که مرآت چگونه میشود از
انوار شمس و تجلّی آن ممنوع و محروم شود و حال آنکه مسلّم است که وجود مرایا
بنفسه لنفسه نبوده بلکه بوجود شمس قایم و منیرند چنانچه الیوم اگر تمام
مرایای ممکنات بشمس عزّ صمدانی مقابل شوند در جمیع انوار شمس ظاهر و مشهود و
بمجرّد انحراف جمیع معدوم و مفقود بوده و خواهند بود مشاهده در مرآت ظاهره
نمائید تا در مقابل شمس قائم است آثار تجلّی در او ظاهر و بعد از انحراف آن
آثار محو بوده و خواهد بود چنانچه مشهود است

\end{otherlanguage}

Hamchenīn mīfarmāyad, qawluhu ʿazza dhikruhu: Fa-inna mathalahu jalla
dhikruhu ka-mathali al-shams law yuqābilannahu ilā mā lā nihāyata marāyā
kulluhunna la-yastaʿkisanna ʿan tajallī al-shams fī ḥaddihim wa in lan
yuqābilhā min aḥadin fa-yaṭluʿu al-shams wa yaghrub wa al-ḥijābu
li'l-marāyā. Ḥāl molāḥeẓe namāʾīd ke az īn bayān-e qods-e rabbānī
mostafād mīshavad ke momken ast shams ṭāleʿ shavad va ghorūb namāyad va
marāyā dar ḥejāb bāshand. Maʿa-dhālek mītavān goft baʿd az eshrāq-e
Shams al-shumūs allatī yaṭufna fī ḥawlihā shumūsun lā yaʿlamu
ʿiddatahunna aḥadun illā nafsu Allāh al-ʿĀlim al-ʿAlīm ke merʾāt chegūne
mīshavad az anvār-e shams va tajallī-ye ān mamnūʿ va maḥrūm shavad va
ḥāl ānke mosallam ast ke vojūd-e marāyā be-nafsihi li-nafsihi nabūde
balke be-vojūd-e shams qāʾem va monīrand. Chenānche emrūz agar tamām-e
marāyā-ye momkenāt be-shams-e ʿezz-e ṣamadānī moqābel shavand dar jamīʿ
anvār-e shams ẓāher va mashhūd va be-mojarrad-e enḥerāf jamīʿ maʿdūm va
mafqūd būde va khāhand būd. Moshāhede dar merʾāt-e ẓāhere namāʾīd tā dar
moqābel-e shams qāʾem ast āthār-e tajallī dar ū ẓāher va baʿd az enḥerāf
ān āthār maḥv būde va khāhad būd chenānche mashhūd ast.

He also states---exalted be His mention: ``His likeness---glorified be
His mention---is that of the sun. Were mirrors without end to face it,
they would all reflect the sun's radiance according to their own limits.
If none should face it, the sun would still rise and set, and it is the
mirrors that would be veiled.'' Now consider how this holy, divine
utterance conveys that the sun may rise and set while the mirrors remain
veiled. Nevertheless, one may ask: after the dawning of the Sun of suns,
around which revolve suns whose number no one knows except God Himself,
the Knowing, the All-Knowing, how can a mirror be kept from the sun's
light and deprived of its radiance? Yet it is certain that the mirrors
do not exist of themselves and for themselves; rather, they subsist and
shine through the existence of the sun. Thus today, if all the mirrors
of contingent existence turn toward the Sun of eternal glory, the sun's
light will be manifest and visible in them all; and the moment they turn
aside, all have been and will be nonexistent and lost. Look at an
outward mirror: so long as it faces the sun, the traces of its radiance
appear within it; after it turns aside, those traces have vanished and
will vanish, as is plainly seen.

\subsection{78}\label{section-1633}

\begin{otherlanguage}{persian}

و هم‌چنین میفرماید او است آن سلطان مقتدری که اگر بحرکت بیاید لسان قدسش
خلق میفرماید بقولی آنچه بخواهد و اراده فرماید از نبیّ و ولیّ و صدّیق حال
لایق است از مقتدری که بقول او نبیّ و ولیّ خلق میشوند اعراض نمایند و باسمی
از اسماء و یا بذکری از اذکار و یا بمرآتی از مرایا تمسّک جویند ان هذا الّا
لظلم عظیم تاللّه الحقّ امری را مرتکب شده‌اند که احدی ارتکاب ننموده چه که
بعد از ظهور شمس در قطب زوال مجال توقّف برای نفسی نه بلی اگر بنفس او عارف
نباشند و از عرفان آن جمال قدم بنفسه لنفسه خود را عاجز مشاهده نمایند حجّت
و دلیل طلب نمایند همان حجّت و دلیلی که ایمان کلّ عباد باو محقّق و ثابت بوده
و اگر از بحور جود و فضل خود ظاهر فرمود دیگر که قادر است که توقّف نماید
مگر آنکه بالمرّه از حقّ اعراض نماید و در سجّین قهر و هاویهٴ نفی و جحیم
اعراض مقرّ یابد چنانچه اکثر از نفوسی که الیوم خود را بر رفرف ایمان و عرش
ایقان متّکی و مستوی میدانند بالمرّة از حقّ اعراض نموده و در بحر کفر مستغرق
و در کلّ حین بعذاب بدیع معذّبند ولکن من غیر شعور

\end{otherlanguage}

Va hamchenīn mīfarmāyad ū ast ān solṭān-e moqtaderī ke agar be-ḥarakat
bīyāyad lesān-e qodsash khalq mīfarmāyad be-qawlī ānche bekhāhad va
erāde farmāyad az nabī va valī va ṣeddīq. Ḥāl lāyeq ast az moqtaderī ke
be-qawl-e ū nabī va valī khalq mīshavand eʿrāḍ namāyand va be-esmī az
asmāʾ va yā be-dhekrī az adhkār va yā be-merʾātī az marāyā tamassok
jūyand? In hādhā illā la-ẓulmun ʿaẓīm. Ta'llāhi al-Ḥaqq amrī rā mortakeb
shode-and ke aḥadī ertekāb nanamūde che ke baʿd az ẓohūr-e shams dar
qoṭb-e zavāl majāl-e tavaqqof barā-ye nafsī na. Balī agar be-nafs-e ū
ʿāref nabāshand va az ʿerfān-e ān Jamāl-e Qedam bi-nafsihi li-nafsihi
khod rā ʿājez moshāhede namāyand ḥojjat va dalīl ṭalab namāyand hamān
ḥojjat va dalīlī ke īmān-e koll-e ʿebād be-ū moḥaqqaq va thābet būde va
agar az boḥūr-e jūd va faḍl-e khod ẓāher farmūd dīgar ke qāder ast ke
tavaqqof namāyad magar ānke bi'l-marre az Ḥaqq eʿrāḍ namāyad va dar
sijjīn-e qahr va hāviye-ye nafy va jaḥīm-e eʿrāḍ maqarr yābad. Chenānche
akthar az nofūsī ke emrūz khod rā bar rafraf-e īmān va ʿarsh-e īqān
mottakī va mostavī mīdānand bi'l-marre az Ḥaqq eʿrāḍ namūde va dar
baḥr-e kofr mostaghraq va dar koll-e ḥīn be-ʿadhāb-e badīʿ moʿadhdhaband
va lākin min ghayri shuʿūr.

He also states: ``He is that mighty Sovereign whose holy tongue, if it
moves, creates with a word whatever He wishes and wills, whether
prophet, holy one, or righteous one.'' Is it fitting, then, to turn away
from the All-Powerful, by whose word prophets and holy ones are created,
and cling to one of the names, or one of the designations, or one of the
mirrors? This is indeed nothing but a grievous wrong. By God, the True
One! They have committed what no one has ever committed, for after the
sun has appeared at the zenith, no soul has any grounds for hesitation.
Yes, if they do not recognize Him through His own Self and find
themselves unable to recognize that Ancient Beauty by Himself and for
Himself, let them ask for proof and evidence---the very proof and
evidence through which the faith of all the servants has been
established and confirmed. And if He brings it forth from the seas of
His bounty and grace, who can then hesitate, unless he turns away
entirely from the Truth and takes his abode in the dungeon of wrath, the
abyss of negation, and the hell of rejection? Thus most of those who
today consider themselves reclining upon the lofty seat of faith and
established upon the throne of certitude have turned away entirely from
the Truth, are submerged in the sea of unbelief, and at every moment
suffer an unprecedented torment, yet remain unaware.

\subsection{79}\label{section-1634}

\begin{otherlanguage}{persian}

و ایکاش باعراض و انکار کفایت مینمودند لا فوالّذی انطقنی بثنآء نفسه که
ابداً کفایت نخواهند نمود چنانچه نفس الّتی ربّیناه فی مذ الأیّام و علّمناه فی
کلّ الأحیان کلمات الرّحمن کما یعلّمون اهل الهند طیورهم قام علیّ و حارب بنفسی
و اعرض عن جمالی و جادل بآیات اللّه المهیمن القیّوم و باین کفایت ننموده بر
قطع سدرهٴ الوهیّة ایستاد فلمّا دمّره اللّه و ظهر ما فی قلبه اذاً قام علی
المکر بشأن لن یقدر احد ان یحصیه الّا اللّه الملک العزیز القدیر و مفتریاتی
جعل نموده و انتشار داده که قلم عاجز است از ذکرش و شقاوتی که از اوّل دنیا
الی حین ظاهر شده بنفس اللّه القائمة علی کلّ من فی السّموات و الأرض نسبت
داده و مع‌ذلک در مقرّ خود ساکن و مستریح است چه که از حمقای ناس مطمئن است
که کلّ در سبیل وهم سالکند و در بحر تقلید سابح فوالّذی نفسی بیده که اگر اقلّ
من ان یحصی ناس را بابصر مشاهده مینمود باین مزخرفات ارتکاب نمینمود ولکن
غافل از اینکه عنقریب از خلف استار عصمت الهی هیاکلی ظاهر شوند که ببصر
الهی حقّ را از باطل و شمس را از ظلّ فرق گذارند و بهیچ بندی از صراط مستقیم
ممنوع نشوند و بهیچ سدّی از نبأ عظیم محروم نگردند فیا لیت من ذی بصر لینظر
الی کلماتهم الّتی یحکی عن انفسهم و ذواتهم لیعرفهم و یکون من العارفین قسم
بسلطان اعظم که اگر نفسی در کلمات معرضین ببصر حدید ملاحظه نماید شأن و
رتبه و مقام این انفس وهمیّه را ادراک مینماید

\end{otherlanguage}

Va ey kāsh be-eʿrāḍ va enkār kefāyat mīnamūdand. Lā fa-wa'lladhī
anṭaqanī bi-thanāʾi nafsihi ke abadan kefāyat nakhāhand namūd. Chenānche
nafs allatī rabbaynāhu fī mudhdh al-ayyām wa ʿallamnāhu fī kull al-aḥyān
kalimāt al-Raḥmān kamā yuʿallimūna ahl al-Hind ṭuyūrahum qāma ʿalayya wa
ḥāraba bi-nafsī wa aʿraḍa ʿan jamālī wa jādala bi-āyāti Allāh
al-Muhaymin al-Qayyūm va be-īn kefāyat nanamūde bar qaṭʿ-e Sedre-ye
Olūhiyyat īstād. Fa-lammā dammarahu Allāh wa ẓahara mā fī qalbihi idhan
qāma ʿalā al-makr bi-shaʾnin lan yaqdir aḥadun an yuḥṣiyahu illā Allāh
al-Malik al-ʿAzīz al-Qadīr va moftarayātī jaʿl namūde va enteshār dāde
ke qalam ʿājez ast az dhekrash va shaqāvatī ke az avval-e donyā ilā ḥīn
ẓāher shode be-nafs Allāh al-qāʾima ʿalā kulli man fī al-samāwāt wa
al-arḍ nesbat dāde va maʿa-dhālek dar maqarr-e khod sāken va mostarīḥ
ast che ke az ḥomaqā-ye nās moṭmaʾenn ast ke koll dar sabīl-e vahm
sālekand va dar baḥr-e taqlīd sābeḥ. Fa-wa'lladhī nafsī bi-yadihi ke
agar aqall min an yuḥṣā nās rā bā-baṣar moshāhede mīnamūd be-īn
mozakhrafāt ertekāb nemīnamūd va lāken ghāfel az īnke ʿan-qarīb az
khalf-e astār-e ʿeṣmat-e elāhī hayākelī ẓāher shavand ke be-baṣar-e
elāhī Ḥaqq rā az bāṭel va shams rā az ẓell farq godhārand va be-hīch
bandī az ṣerāṭ-e mostaqīm mamnūʿ nashavand va be-hīch saddī az nabaʾ-e
ʿaẓīm maḥrūm nagardand. Fa-yā layta min dhī baṣarin li-yanẓura ilā
kalimātihim allatī yaḥkī ʿan anfusihim wa dhawātihim li-yaʿrifahum wa
yakūna min al-ʿārifīn. Qasam be-Solṭān-e Aʿẓam ke agar nafsī dar
kalemāt-e moʿreḍīn be-baṣar-e ḥadīd molāḥeẓe namāyad shaʾn va rotbe va
maqām-e īn anfos-e vahmiyye rā edrāk mīnamāyad.

Would that they had contented themselves with rejection and denial! No,
by Him who has caused Me to speak in praise of His own Self, they will
never be content with that. Thus the soul whom We nurtured over the days
and taught at all times the words of the All-Merciful, as the people of
India teach their birds, rose against Me, waged war upon My own Self,
turned away from My beauty, and disputed the verses of God, the
All-Protecting, the Self-Subsisting. Not content with this, he arose to
cut down the Tree of Divinity. When God brought him to ruin and what was
in his heart became manifest, he then resorted to such plotting as no
one can reckon except God, the King, the Almighty, the Omnipotent. He
has fabricated and circulated slanders that the pen is powerless to
recount, and has ascribed every wickedness manifested from the beginning
of the world until now to the Self of God, who sustains all who are in
the heavens and on earth. Yet he remains settled and at ease in his own
abode, for he is confident of the fools among the people, who all tread
the path of delusion and swim in the sea of imitation. By Him in whose
hand is My soul! Had he perceived discernment in even an uncountably
small number of people, he would not have committed these absurdities.
But he is heedless that before long, from behind the veils of divine
protection, figures will appear who, with divine sight, will distinguish
truth from falsehood and the sun from the shadow. No bond will keep them
from the straight path, and no barrier will deprive them of the
momentous news. Would that someone possessed sight, to examine their
words, which disclose their own selves and essences, so that he might
recognize them and be among those who understand. I swear by the Most
Great Sovereign that if anyone examines the words of those who turn away
with penetrating sight, he will perceive the condition, rank, and
station of these deluded souls.

\subsection{80}\label{section-1635}

\begin{otherlanguage}{persian}

و الیوم رؤسای بیان بهمان ادلّه که پست‌ترین اهل فرقان بر حقّیّت خود استدلال
مینموده‌اند بهمان دلایل استدلال نموده و مینمایند من حیث لا یشعرون مثلاً
حکم وصایت را که ظهور قبلم بالمرّه از کتاب محو نموده چنانچه جمیع مطّلعند که
جز حروفات و مرایا در بیان از قلم رحمن نازل نشده و مرایا را هم محدود
نفرموده‌اند چنانچه در دعوات میفرمایند الهی در کلّ حین بفرست مرایای ممتنعه
و بلّوریّات صافیه لیحکینّ عنک و یدلّنّ منک و در جمیع الواح من قبل اللّه این
کلمات نازل من یشآء فلینظر الیها و یکون من العارفین حال بمثل امّت فرقان که
خاتم النّبیّین درست نموده و از مرسل و مبعث آن محتجب مانده‌اند این اصحاب هم
اراده نموده خاتم الولیّین برای نفس خود ثابت نمایند چنانچه مشهود شد و نفسی
را که بقول او صدهزار ولیّ خلق شده و میشود در بئر بغضا انداخته‌اند و در کلّ
حین احجار ظنون از کلّ جهات بر وجود مبارکش میاندازند و در نفس خود صیحه
میزنند و دعوی مظلومیّت نموده که شاید غلّ اللّه در انفس ضعیفه القا نمایند
اینست شأن مرایا که ظاهر شده قل الیوم لو یحکم اللّه لاسم الّذی یکون ابغض
الأسمآء عند النّاس بأنّه ربّ لکم او بالعکس لیس لأحد ان یقول لم او بم لأنّه
جلّ و عزّ یحکم ما یرید و لا یسأل عمّا اراد و انّه لهو المقتدر القدیر

\end{otherlanguage}

Va emrūz roʾasā-ye Bayān be-hamān adelle ke past-tarīn-e ahl-e Forqān
bar ḥaqqiyyat-e khod estedlāl mīnamūde-and be-hamān dalāyel estedlāl
namūde va mīnamāyand min ḥaythu lā yashʿurūn. Mathalan ḥokm-e veṣāyat rā
ke ẓohūr-e qablam bi'l-marre az ketāb maḥv namūde chenānche jamīʿ
moṭṭaleʿand ke joz ḥorūfāt va marāyā dar Bayān az qalam-e Raḥmān nāzel
nashode va marāyā rā ham maḥdūd nafarmūde-and chenānche dar daʿavāt
mīfarmāyand Elāhī dar koll-e ḥīn beferest marāyā-ye momtaneʿe va
bellūriyyāt-e ṣāfiye li-yaḥkīna ʿanka wa yadululna minka va dar jamīʿ-e
alvāḥ min qibal Allāh īn kalemāt nāzel man yashāʾ fa'l-yanẓur ilayhā wa
yakūn min al-ʿārifīn. Ḥāl be-methl-e ommat-e Forqān ke khātam
al-nabiyyīn dorost namūde va az morsel va mobʿeth-e ān moḥtajeb
mānde-and īn aṣḥāb ham erāde namūde khātam al-valiyyīn barā-ye nafs-e
khod thābet namāyand chenānche mashhūd shod va nafsī rā ke be-qawl-e ū
ṣad-hezār valī khalq shode va mīshavad dar beʾr-e baghḍā andākhte-and va
dar koll-e ḥīn aḥjār-e ẓonūn az koll-e jehāt bar vojūd-e mobārakash
mīandāzand va dar nafs-e khod ṣayḥe mīzanand va daʿvā-ye maẓlūmiyyat
namūde ke shāyad ghell-e Allāh dar anfos-e ḍaʿīfe elqā namāyand. Īnast
shaʾn-e marāyā ke ẓāher shode. Qul al-yawm law yaḥkum Allāh li-ism
alladhī yakūnu abghaḍ al-asmāʾ ʿinda al-nās bi-annahu rabbun lakum aw
bi'l-ʿaks laysa li-aḥadin an yaqūla lima aw bi-ma li-annahu jalla wa
ʿazza yaḥkumu mā yurīd wa lā yusʾalu ʿammā arāda wa innahu la-huwa
al-Muqtadir al-Qadīr.

Today the leaders of the Bayān have advanced and continue to advance,
without realizing it, the very arguments by which the lowliest of the
people of the Qurʾān argued for the truth of their own claims. One
example is the ordinance of appointed succession, which My previous
Manifestation entirely effaced from the Book. As everyone knows, nothing
concerning this was sent down in the Bayān from the Pen of the
All-Merciful apart from letters and mirrors, nor did He set any limit on
the mirrors. Thus He says in His prayers: ``My God, send forth at every
moment inviolable mirrors and pure crystals, that they may reflect You
and point to You.'' These words have been sent down in all the Tablets
from God. Whoever wishes, let him examine them and be among those who
understand. Now, just as the community of the Qurʾān fashioned a ``Seal
of the Prophets'' and remained veiled from Him who sent and raised him
up, these companions too have sought to establish a ``Seal of the Holy
Ones'' for themselves, as has been witnessed. They have cast into the
well of hatred the One by whose word a hundred thousand holy ones have
been and are created. At every moment they hurl the stones of their
suspicions from every direction at His blessed Being, while they
themselves cry out and claim to be wronged, hoping thereby to instill
hatred of God in weak souls. Such is the condition of the mirrors that
have appeared. Say: If today God should decree that the name most
hateful to the people is your Lord, or the reverse, no one has the right
to say ``Why?'' or ``On what grounds?'' For He, glorified and exalted,
decrees whatever He wills and is not questioned about what He has
willed. He is indeed the All-Powerful, the Omnipotent.

\subsection{81}\label{section-1636}

\begin{otherlanguage}{persian}

و در این ایّام رسائلی در ردّ حقّ بانامل شرکیّه نوشته و ارسال داشته‌اند قسم
بآفتاب معانی که مثل اطفال بکلمات مزخرفهٴ بی‌معنی تکلّم نموده‌اند بل احقر لو
انتم تعرفون حال کتب و رسائل این فئه را با آنچه از خدّام این بیت ظاهر شده
میزان نمائید و خود انصاف دهید که شاید الیوم از فیض بحر معانی که جمیع
ممکنات را احاطه نموده محروم نمانید

\end{otherlanguage}

Va dar īn ayyām rasāʾelī dar radd-e Ḥaqq be-anāmel-e sherkiyye neveshte
va ersāl dāshte-and. Qasam be-Āftāb-e Maʿānī ke methl-e aṭfāl
be-kalemāt-e mozakhrafe-ye bī-maʿnī takallom namūde-and bal aḥqar law
antum taʿrifūn. Ḥāl kotob va rasāʾel-e īn feʾe rā bā ānche az khoddām-e
īn bayt ẓāher shode mīzān namāʾīd va khod enṣāf dahīd ke shāyad emrūz az
fayḍ-e baḥr-e maʿānī ke jamīʿ-e momkenāt rā eḥāṭe namūde maḥrūm namānīd.

In these days they have written and dispatched treatises refuting the
Truth, penned by fingers of idolatry. I swear by the Sun of inner
meanings that, like children, they have uttered meaningless
absurdities---indeed, they are even more contemptible, if you only knew.
Now weigh the books and treatises of this faction against what has come
forth from the servants of this House, and judge fairly for yourselves,
so that today you may not remain deprived of the outpouring of the ocean
of meanings that has encompassed all contingent things.

\subsection{82}\label{section-1637}

\begin{otherlanguage}{persian}

و حال حقّ بشأنی تنزّل نموده که از اطراف متتابعاً از مسائل فرقانیّه سؤال
مینمایند و باید جواب مرقوم دارد مثلاً نفسی مرآت را یاقوت فرض گرفته و بعد
سؤال نموده که یاقوت قبل از بلوغ رتبهٴ او باین مقام میشود تغییر نماید و
بعد از بلوغ چگونه تغییر مینماید یعنی مرآت قبل از وصول باین مقام مرآتیّت
ممکن بوده تبدیل شود و یا تغییر نماید ولکن بعد از بلوغ باین مقام چگونه
تنزّل مینماید

\end{otherlanguage}

Va ḥāl Ḥaqq be-shaʾnī tanazzol namūde ke az aṭrāf motatābeʿan az
masāʾel-e Forqāniyye soʾāl mīnamāyand va bāyad javāb marqūm dārad.
Mathalan nafsī merʾāt rā yāqūt farḍ gerefte va baʿd soʾāl namūde ke
yāqūt qabl az bolūgh-e rotbe-ye ū be-īn maqām mīshavad taghyīr namāyad
va baʿd az bolūgh chegūne taghyīr mīnamāyad yaʿnī merʾāt qabl az voṣūl
be-īn maqām-e merʾātiyyat momken būde tabdīl shavad va yā taghyīr
namāyad va lāken baʿd az bolūgh be-īn maqām chegūne tanazzol mīnamāyad.

Now the Truth has so far condescended that questions concerning matters
of the Qurʾān are continually asked from every quarter, and He must
write replies. For example, someone has likened a mirror to a ruby and
then asked: a ruby may change before attaining this rank, but how can it
change after attaining it? That is, before a mirror reached this station
of mirrorhood, it could have been transformed or changed; but after
attaining this station, how can it decline?

\subsection{83}\label{section-1638}

\begin{otherlanguage}{arabic}

چنانچه اکثری باین مطالب بدیهیّه الیوم از سلطان احدیّه محروم مانده‌اند اوّلاً
بگو ای سائل نسأل اللّه بأن یوفّقک و یؤیّدک بشأن تعرف کلّ شیء فی مقامه و
تنقطع عن کلّ الاشارات و لن تشهد فی الملک الّا تجلّی انوار الّتی احاطت
العالمین و یمحو عن قلبک کلّ الأذکار و تقوم علی ثنآء ربّک العزیز المختار

\end{otherlanguage}

Chenānche aktharī be-īn maṭāleb-e badīhiyye emrūz az Solṭān-e Aḥadiyye
maḥrūm mānde-and. Avvalan begū ey sāʾel nasʾalu Allāh bi-an yuwaffiqaka
wa yuʾayyidaka bi-shaʾnin taʿrifu kulla shayʾin fī maqāmihi wa tanqaṭiʿu
ʿan kull al-ishārāt wa lan tashhada fī al-mulk illā tajallī anwār allatī
aḥāṭat al-ʿālamīn wa yamḥū ʿan qalbika kulla al-adhkār wa taqūmu ʿalā
thanāʾi rabbika al-ʿAzīz al-Mukhtār.

Thus many have today remained deprived of the Sovereign of Oneness
because of these self-evident matters. First say: O questioner, we ask
God to grant you success and strengthen you, so that you may recognize
everything in its own station, detach yourself from all allusions, and
behold in the realm of existence nothing but the radiance of the lights
that have encompassed all worlds; that He may efface every mention from
your heart, and that you may arise to praise your Lord, the Almighty,
the Unconstrained.

\subsection{84}\label{section-1639}

\begin{otherlanguage}{arabic}

لأنّ فی مثل ذلک الیوم لا ینبغی لأحد ان یلتفت الی شیء عمّا خلق بین السّموات
و الأرضین هل ینبغی لنفس بعد اشراق الشّمس فی قطب الزّوال ان یشتغل بذکر
النّجوم او یسأل عن السّراج ولو یکون سراجاً منیرا لا فوجمالی المشرق المقدّس
العزیز المنیر دع کلّ الأذکار عن ورائک ثمّ تمسّک بهذا الذّکر الّذی ظهر بالحقّ و
ینطق فی کلّ شیء بأنّه لا اله الّا انا العزیز العلیم

\end{otherlanguage}

Li-anna fī mithli dhālika al-yawm lā yanbaghī li-aḥadin an yaltafita ilā
shayʾin ʿammā khuliqa bayna al-samāwāt wa-al-araḍīn. Hal yanbaghī
li-nafsin baʿda ishrāqi al-shams fī quṭbi al-zawāl an yashtaghila
bi-dhikri al-nujūm aw yasʾala ʿan al-sirāj wa-law yakūnu sirājan
munīran? Lā fa-wa-jamālī al-mushriq al-muqaddas al-ʿazīz al-munīr. Daʿ
kulla al-adhkār ʿan warāʾika thumma tamassak bi-hādhā al-dhikr alladhī
ẓahara bi-al-ḥaqq wa-yanṭiqu fī kulli shayʾin bi-annahu lā ilāha illā
anā al-ʿazīz al-ʿalīm.

For on such a day no one should turn toward anything created between the
heavens and the earths. After the sun has shone at its zenith, should
anyone occupy himself with mentioning the stars or inquire about a lamp,
even if it is a shining lamp? No, by My radiant, sanctified, mighty,
luminous Beauty! Leave every remembrance behind you, then hold fast to
this Remembrance, who has appeared in truth and proclaims within all
things: ``There is no God but Me, the Almighty, the All-Knowing.''

\subsection{85}\label{section-1640}

\begin{otherlanguage}{persian}

و اگر موفّق بآنچه در این آیات منزله از سماء احدیّه نازل شده نشدی و بجواب
مسئلهٴ خود ناظر باشی بشنو نداء ابهی را از افق عزّ اعلی و اوّلاً بدان که کلّ
شیء بکلّ شیء تبدیل شده و میشود و علم ذلک فی کتاب ربّک الّذی لا یضلّ و لا
ینسی

\end{otherlanguage}

Va agar muwaffaq be-ānchih dar īn āyāt-e munazzaleh az samā-ye aḥadiyyeh
nāzil shudeh nashudī va be-javāb-e masʾaleh-ye khud nāẓir bāshī, bishnav
nidā-ye Abhā rā az ufuq-e ʿizz-e aʿlā va avvalan bidān keh kull-e shayʾ
be-kull-e shayʾ tabdīl shudeh va mīshavad va ʿilmu dhālika fī kitābi
rabbika alladhī lā yaḍillu wa-lā yansā.

And if you have not succeeded in attaining what has been sent down in
these verses revealed from the heaven of oneness, and still seek an
answer to your question, hear the call of the All-Glorious from the
horizon of supreme glory. Know first that everything has been and is
being transformed into everything else, and knowledge of this is in the
Book of your Lord, who neither errs nor forgets.

\subsection{86}\label{section-1641}

\begin{otherlanguage}{persian}

و ثانیاً اگر از تو سؤال شود که حقّ قادر است بر تبدیل یاقوت یا نه جواب چه
خواهی گفت باری عجز لازال شأن خلق بوده و حقّ منیع در قطب اقتدار قائم بحرفی
قادر است که جمیع من فی الأرض را قطعه‌ئی از یاقوت رطبهٴ حمرا فرماید و
بحرفی کلّ را بحجر راجع نماید فتعالی عمّا انتم ظننتم فی قدرته و تظنّون

\end{otherlanguage}

Va thāniyan agar az tu suʾāl shavad keh Ḥaqq qādir ast bar tabdīl-e
yāqūt yā nah, javāb chih khāhī guft? Bārī ʿajz lā-zāl shaʾn-e khalq
būdeh va Ḥaqq-e manīʿ dar quṭb-e iqtidār qāʾim. Be-ḥarfī qādir ast keh
jamīʿ-e man fī al-arḍ rā qiṭʿeh-ī az yāqūt-e raṭbeh-ye ḥamrā farmāyad va
be-ḥarfī kull rā be-ḥajar rājiʿ namāyad. Fa-taʿālā ʿammā antum ẓanantum
fī qudratihi wa-taẓunnūn.

And second, if you are asked whether God has the power to transform a
ruby or not, what answer will you give? Powerlessness has always been
the condition of creation, while God, the Inaccessible, stands at the
center of power. With a single word He can turn all who are on earth
into a piece of moist crimson ruby, and with a single word return them
all to stone. Exalted is He above what you have imagined and still
imagine concerning His power.

\subsection{87}\label{section-1642}

\begin{otherlanguage}{persian}

و از این مقام گذشته یاقوت را در نار بگذار و بعد ملاحظه کن که چه میشود تا
بر تبدیل کلّ شیء الّا اللّه موقن شوی و هم‌چنین یاقوت قلب که در نار نفس و هوی
مبتلا شد البتّه از لطافت و لون و صفای خود محروم ماند و الیوم بسی از قطعات
یاقوت که بحجر راجع شده و خود شاعر نیستند و کذلک بالعکس لو انتم من
العارفین

\end{otherlanguage}

Va az īn maqām gudhashteh yāqūt rā dar nār bigudhār va baʿd mulāḥiẓeh
kun keh chih mīshavad tā bar tabdīl-e kull-e shayʾ illā Allāh mūqin
shavī. Va ham-chunīn yāqūt-e qalb keh dar nār-e nafs va havā mubtalā
shud al-batteh az liṭāfat va lawn va ṣafā-ye khud maḥrūm mānad va
al-yawm basī az qiṭaʿāt-e yāqūt keh be-ḥajar rājiʿ shudeh va khud shāʿir
nīstand va-kadhālika bi-al-ʿaks law antum min al-ʿārifīn.

And beyond this, place a ruby in the fire and then observe what happens,
so that you may become certain of the transformation of everything
except God. Likewise, when the ruby of the heart is afflicted by the
fire of self and desire, it is inevitably deprived of its delicacy,
color, and purity. Today many pieces of ruby have reverted to stone and
are themselves unaware of it; and the reverse is also true, if you are
among those who understand.

\subsection{88}\label{section-1643}

\begin{otherlanguage}{persian}

در یاقوت انواع متعدّده و الوان مختلفه مشهود و جمیع الوان آن در نار تغییر
مینماید مگر دو قسم که به بهرمانی و رمّانی مشهورند و این دو هم اگرچه بنار
شدیده در ایّام عدیده تغییر مینماید ولکن از ورای حجاب اگر در نار متوسّطه
بمانند ضرّی مشهود نه و بعضی بر آنند که نار بر لون این دو قسم بیفزاید و
امّا اگر نار الهی که در علم مکتوم ربّانی بکبریت احمر معروفست بر یاقوت القا
گردد فی‌الحین تکلیس شود چون سایر اجساد

\end{otherlanguage}

Dar yāqūt anvāʿ-e mutaʿaddideh va alvān-e mukhtalifeh mashhūd va jamīʿ-e
alvān-e ān dar nār taghyīr mīnamāyad magar du qism keh be-bahramānī va
rummānī mashhūrand va īn du ham agarchih be-nār-e shadīdeh dar ayyām-e
ʿadīdeh taghyīr mīnamāyad va-lākin az varā-ye ḥijāb agar dar nār-e
mutavassiṭeh bimānand ḍararī mashhūd nah va baʿḍī bar ānand keh nār bar
lawn-e īn du qism biyafzāyad. Va ammā agar nār-e ilāhī keh dar ʿilm-e
maktūm-e rabbānī be-kibrīt-e aḥmar maʿrūf ast bar yāqūt ilqā gardad fī
al-ḥīn taklīs shavad chūn sāʾir-e ajsād.

Rubies are found in numerous varieties and different colors, and all
their colors change in fire except for two varieties known as
saffron-colored and pomegranate-colored. Although these two also change
under intense fire over many days, albeit from behind a veil, no damage
is observed if they remain in a moderate fire. Some hold that fire
deepens the color of these two varieties. But if the divine fire known
in the hidden divine science as the crimson sulfur is applied to a ruby,
it is immediately calcined like other bodies.

\subsection{89}\label{section-1644}

\begin{otherlanguage}{persian}

انّما النّار یصلح کلّ مفسد و یفسد کلّ مصلح تزید الصّالح صلاحاً و الفاسد فساداً
هر نفسی که از علم الهی مطّلع شده آنچه ذکر شد ادراک مینماید و من دونه لیس
لأحد نصیب الّا بأمر من لدن مقتدر قدیر

\end{otherlanguage}

Innamā al-nār yuṣliḥu kulla mufsid wa-yufsidu kulla muṣliḥ, tazīdu
al-ṣāliḥa ṣalāḥan wa-al-fāsida fasādan. Har nafsī keh az ʿilm-e ilāhī
muṭṭaliʿ shudeh ānchih dhikr shud idrāk mīnamāyad va min dūnihi laysa
li-aḥadin naṣīb illā bi-amrin min ladun muqtadirin qadīr.

Fire indeed sets right all that corrupts and corrupts all that sets
right; it increases the sound in soundness and the corrupt in
corruption. Whoever has become acquainted with divine knowledge
understands what has been stated. Without it, no one has any share
except through a command from One who is All-Powerful, Almighty.

\subsection{90}\label{section-1645}

\begin{otherlanguage}{persian}

باری ای سائل در جمیع ما سوی اللّه و نفس ظهور تغییر و تبدیل مشهود است
خلقنا الأشیآء من التّراب و نعیدنّها الیه و نبعث منه حیّاً ذلک وعد من لدنّا
فسوف نوفی بما وعدناکم به و انا خیر الموفین

\end{otherlanguage}

Bārī ey sāʾil dar jamīʿ-e mā sivā Allāh va nafs-e ẓuhūr taghyīr va
tabdīl mashhūd ast. Khalaqnā al-ashyāʾa min al-turāb wa-nuʿīdannahā
ilayhi wa-nabʿathu minhu ḥayyan. Dhālika waʿdun min ladunnā fa-sawfa
nūfī bi-mā waʿadnākum bihi wa-anā khayru al-mūfīn.

In short, O questioner, change and transformation are evident in all
that is other than God and in the Manifestation Himself. We created all
things from dust, and We will surely return them to it and raise them
from it alive. This is a promise from Us, and We will fulfill what We
have promised you. I am the best of those who fulfill their promises.

\subsection{91}\label{section-1646}

\begin{otherlanguage}{persian}

ای قوم ندای طیر بقا را از رضوان اعلی بشنوید و رحیق الطف ارقّ اصفی از
انامل قدس بها بنوشید و از کلّ ما سوی اللّه غنیّ و بی‌نیاز شوید تا بحور علم و
حکمت ربّانی از قلب و لسان جاری شود و از امثال این مسائل فارغ و مقدّس گردید
زینهار قلب را که ودیعهٴ جمال مختار است بآلایش کلمات فجّار و شبهات اشرار
میالائید چه که الیوم بکلّ حیل و مکر ظاهر شده‌اند و بهر نحو که ممکن شود
القای شبهه در قلوب مینمایند

\end{otherlanguage}

Ey qawm nidā-ye ṭayr-e baqā rā az Riḍvān-e aʿlā bishnavīd va raḥīq-e
alṭaf-e araqq-e aṣfā az anāmil-e quds-e Bahā binūshīd va az kull-e mā
sivā Allāh ghanī va bī-niyāz shavīd tā buḥūr-e ʿilm va ḥikmat-e rabbānī
az qalb va lisān jārī shavad va az amthāl-e īn masāʾil fārigh va
muqaddas gardīd. Zīnhār qalb rā keh vadīʿeh-ye Jamāl-e Mukhtār ast
be-ālāyish-e kalimāt-e fujjār va shubuhāt-e ashrār mayālāyīd chih keh
al-yawm be-kull-e ḥiyal va makr ẓāhir shudeh-and va be-har naḥv keh
mumkin shavad ilqā-ye shubheh dar qulūb mīnamāyand.

O people, hear the call of the Bird of Eternity from the highest
Paradise, drink the most exquisite, delicate, and pure wine from the
holy fingers of Bahá, and become independent of all but God, needing
nothing from them, so that oceans of divine knowledge and wisdom may
flow from your hearts and tongues and you may be freed and sanctified
from such questions. Take care not to defile the heart, which is
entrusted to you by the sovereign Beauty, with the pollution of the
words of the wicked and the doubts raised by evildoers. For today they
have come forth with every trick and deception, instilling doubts in
hearts by every possible means.

\subsection{92}\label{section-1647}

\begin{otherlanguage}{arabic}

و از همهٴ این بیانات گذشته اگر در عالم شیئی یافت شود که تبدیل نشود و
تغییر نیابد این چه دخلی بمطلب این قوم دارد مثلاً یاقوت بزعم این قوم اگر
تغییر نیابد چه مناسبت دارد بآنکه مقبل نمیشود که معرض شود و یا موحّد مشرک
و مؤمن کافر در بعضی اشیاء ملیک اسماء امکان تغییر گذاشته و در بعضی
نگذاشته مثلاً نحاس امکان دارد ذهب شود ولکن تراب بالفعل امکان ذهبیّت در او
موجود نه و این مطلب را چون اهل علم از قبل ذکر نموده‌اند این عبد دوست
نداشته که مفصّل ذکر نماید

\end{otherlanguage}

Va az hameh-ye īn bayānāt gudhashteh agar dar ʿālam shayʾī yāft shavad
keh tabdīl nashavad va taghyīr nayābad īn chih dakhlī be-maṭlab-e īn
qawm dārad? Mathalan yāqūt be-zaʿm-e īn qawm agar taghyīr nayābad chih
munāsibat dārad be-ānkeh muqbil namīshavad keh muʿriḍ shavad va yā
muvaḥḥid mushrik va muʾmin kāfir? Dar baʿḍī ashyāʾ Malīk-e Asmāʾ imkān-e
taghyīr gudhāshteh va dar baʿḍī nagudhāshteh. Mathalan nuḥās imkān dārad
dhahab shavad va-lākin turāb bi-al-fiʿl imkān-e dhahabiyyat dar ū mawjūd
nah va īn maṭlab rā chūn ahl-e ʿilm az qabl dhikr namūdeh-and īn ʿabd
dūst nadāshteh keh mufaṣṣal dhikr namāyad.

And beyond all these explanations, if something were found in the world
that did not undergo transformation or change, what bearing would this
have on the argument of these people? For example, if, as these people
suppose, a ruby does not change, what connection does this have with the
claim that one who has turned toward God cannot turn away, or that a
believer in His oneness cannot become an idolater, or a believer an
unbeliever? The King of Names has placed the possibility of change in
some things and not in others. For example, copper can become gold, but
the possibility of becoming gold is not actually present in earth. Since
people of learning have discussed this matter before, this Servant has
not wished to discuss it in detail.

\subsection{93}\label{section-1648}

\begin{otherlanguage}{persian}

ولکن در انسان امکان علم و جهل و اقبال و اعراض و ایمان و کفر موجود و
مشهود است و اگر مقصود سائل از ذکر ذهب و یاقوت در رتبهٴ جماد بوده این از
مقصود او بغایت دور است چه اگر هم بالفرض ذهب نحاس نشود و یاقوت رماد از
این نمیتوان در سایر مصنوعات صانع قیاس نمود و این نزد اهل بصر که بمنظر
اکبر ناظرند واضح و لایح است

\end{otherlanguage}

Va-lākin dar insān imkān-e ʿilm va jahl va iqbāl va iʿrāḍ va īmān va
kufr mawjūd va mashhūd ast va agar maqṣūd-e sāʾil az dhikr-e dhahab va
yāqūt dar rutbeh-ye jamād būdeh īn az maqṣūd-e ū be-ghāyat dūr ast chih
agar ham bi-al-farḍ dhahab nuḥās nashavad va yāqūt ramād az īn namītavān
dar sāʾir-e maṣnūʿāt-e Ṣāniʿ qiyās namūd va īn nazd-e ahl-e baṣar keh
be-Manẓar-e Akbar nāẓirand vāḍiḥ va lāyiḥ ast.

In the human being, however, the possibilities of knowledge and
ignorance, turning toward and turning away, faith and unbelief are
present and evident. If the questioner, in mentioning gold and ruby,
meant them at the mineral level, this is very far from establishing his
point. For even supposing that gold does not become copper and ruby does
not become ash, one cannot draw an analogy from this concerning the
other handiworks of the Maker. This is clear and evident to those
endowed with sight who look toward the Most Great Vision.

\subsection{94}\label{section-1649}

\begin{otherlanguage}{persian}

و اگر مقصود در رتبهٴ انسان است معلوم بوده که ذهب و یاقوت در رتبهٴ مؤمن
ایمانه باللّه و عرفانه نفسه بوده و خواهد بود و این بسی واضح است که تبدیل
میشود چنانچه بسا از عباد که در اوّل ظهور جمال رحمن بایمان فایز شده‌اند و
بعد معرض و اگر از عالم وهم متصاعد شوی و قلب را از اشارهٴ قولیّه مقدّس
نمائی تبدیل کلّ شیء را در کلّ حین بکلّ شیء مشاهده نمائی و من فتح له هذا
البصر انّه من اولی الأبصار لدی اللّه المقتدر المختار

\end{otherlanguage}

Va agar maqṣūd dar rutbeh-ye insān ast maʿlūm būdeh keh dhahab va yāqūt
dar rutbeh-ye muʾmin īmānuhu bi-Allāh va ʿirfānuhu nafsahu būdeh va
khāhad būd va īn basī vāḍiḥ ast keh tabdīl mīshavad chunānchih basā az
ʿibād keh dar avval-e ẓuhūr-e Jamāl-e Raḥmān be-īmān fāʾiz shudeh-and va
baʿd muʿriḍ. Va agar az ʿālam-e vahm mutaṣāʿid shavī va qalb rā az
ishāreh-ye qawliyyih muqaddas namāʾī tabdīl-e kull-e shayʾ rā dar kull-e
ḥīn be-kull-e shayʾ mushāhideh namāʾī va man futiḥa lahu hādhā al-baṣar
innahu min ulī al-abṣār ladā Allāh al-muqtadir al-mukhtār.

And if what is intended concerns the human level, it is known that gold
and ruby, at the level of the believer, have been and will remain his
faith in God and his recognition of Him. It is abundantly clear that
this can change, for many servants attained faith at the beginning of
the revelation of the Beauty of the All-Merciful and afterward turned
away. If you rise above the realm of illusion and sanctify your heart
from verbal allusions, you will observe the transformation of everything
into everything else at every moment. Whoever has this vision opened to
him is indeed reckoned among those endowed with insight before God, the
All-Powerful, the Unconstrained.

\subsection{95}\label{section-1650}

\begin{otherlanguage}{persian}

جوهر بیان طلعت رحمن آنکه امر اللّه لازال مقدّس از ذکر و وصف و تحدید بوده و
خواهد بود و هر نفسی باین حجبات اراده نماید ملیک من فی لجج اللّانهایات را
ادراک نماید ابداً موقن نخواهد شد بقوّت سلطان بقا حجبات اسماء و دون آن و
فوق آن را خرق نمائید و بمدینهٴ اعزّ ارفع امنع اقدس ابهی وارد شوید که
الیوم جز ظلّش نار و جز حبّش مردود بوده و خواهد بود

\end{otherlanguage}

Jawhar-e bayān-e Ṭalʿat-e Raḥmān ānkeh Amr Allāh lā-zāl muqaddas az
dhikr va vaṣf va taḥdīd būdeh va khāhad būd va har nafsī be-īn ḥujubāt
irādeh namāyad Malīk-e man fī lujaj-e lā-nihāyāt rā idrāk namāyad abadan
mūqin nakhāhad shud. Be-quvvat-e Sulṭān-e Baqā ḥujubāt-e asmāʾ va dūn-e
ān va fawq-e ān rā kharq namāʾīd va be-madīneh-ye aʿazz-e arfaʿ-e
amnaʿ-e aqdas-e abhā vārid shavīd keh al-yawm juz ẓillash nār va juz
ḥubbash mardūd būdeh va khāhad būd.

The essence of the utterance of the Countenance of the All-Merciful is
this: the Cause of God has always been and will always remain sanctified
above mention, description, and limitation. Whoever seeks through these
veils to comprehend the King of all who dwell in the billowing oceans of
infinity will never attain certainty. Through the power of the Sovereign
of Eternity, tear through the veils of names and all that is below and
above them, and enter the most mighty, most exalted, most inaccessible,
most holy, most glorious City. For today all but its shade is fire, and
all but love of it has been and will remain rejected.

\subsection{96}\label{section-1651}

\begin{otherlanguage}{persian}

و هم‌چنین نفسی از عدد آیهٴ منزلهٴ اقتربت السّاعة و انشقّ القمر سؤال نموده
بگو ان یا ایّها النّاظر الی الکلمات عدد را بگذار و بخود آیه ملاحظه نما و
اگر الیوم کلّ من فی السّموات و الأرض باین آیهٴ مبارکه ناظر شوند و در معانی
مستورهٴ آن تفکّر نمایند جمیع را کفایت نماید چه که امروز شمس باطن از افق
ظاهر آیه اشراق فرموده و سیفی است قاطع در ردّ نفوسی که غیر اللّه را ولیّ اخذ
نموده‌اند و معنی این آیه تا این ظهور منیع ظاهر نشده چه که در حین ظهور
نقطهٴ اولی ولیّی علی زعم ناس نبوده تا منشقّ شود بلکه علما بوده‌اند و تعبیر
از علما بنجوم شده و از ولیّ بقمر و از نبیّ بشمس اینست که در این ظهور شمس
از افق اللّه مشرق و قمر وهم منشقّ شده اذاً قل فتبارک اللّه اقدر الأقدرین

\end{otherlanguage}

Va ham-chunīn nafsī az ʿadad-e āyeh-ye munazzaleh-ye iqtarabati
al-sāʿatu wa-inshaqqa al-qamar suʾāl namūdeh. Bigū an yā ayyuhā al-nāẓir
ilā al-kalimāt ʿadad rā bigudhār va be-khud-e āyeh mulāḥiẓeh namā va
agar al-yawm kull-e man fī al-samāwāt va al-arḍ be-īn āyeh-ye mubārakeh
nāẓir shavand va dar maʿānī-ye mastūreh-ye ān tafakkur namāyand jamīʿ rā
kifāyat namāyad chih keh imrūz shams-e bāṭin az ufuq-e ẓāhir-e āyeh
ishrāq farmūdeh va sayfī ast qāṭiʿ dar radd-e nufūsī keh ghayr Allāh rā
valī akhdh namūdeh-and va maʿnī-ye īn āyeh tā īn ẓuhūr-e manīʿ ẓāhir
nashudeh chih keh dar ḥīn-e ẓuhūr-e Nuqṭeh-ye Ūlā valī-ī ʿalā zaʿm-e nās
nabūdeh tā munshaqq shavad balkih ʿulamā būdeh-and va taʿbīr az ʿulamā
be-nujūm shudeh va az valī be-qamar va az nabī be-shams. Īn ast keh dar
īn ẓuhūr shams az ufuq Allāh mushriq va qamar-e vahm munshaqq shudeh.
Idhan qul fa-tabāraka Allāh aqdar al-aqdarīn.

Someone has also asked about the numerical value of the revealed verse,
``The Hour has drawn near, and the moon has split.'' Say: O you who look
at the words, set aside the numerical value and consider the verse
itself. If today all who are in the heavens and on earth were to turn to
this blessed verse and reflect on its hidden meanings, it would suffice
them all. For today the sun of its inner meaning has shone from the
horizon of the verse's outward form, and it is a sharp sword refuting
those who have taken anyone other than God as their guardian. The
meaning of this verse had not become manifest until this inaccessible
Revelation, for at the time of the appearance of the Primal Point there
was no guardian, as people conceive it, to be split apart. Rather, there
were religious scholars: scholars have been referred to as stars, the
guardian as the moon, and the prophet as the sun. Thus, in this
Revelation the sun has risen from the horizon of God, and the moon of
illusion has split apart. Then say: Blessed be God, the most powerful of
all who wield power!

\subsection{97}\label{section-1652}

\begin{otherlanguage}{persian}

حال در این آیهٴ منزله تفکّر نمائید و انصاف دهید و لا تکونوا من الّذینهم
یعرفون نعمة اللّه ثمّ ینکرون و ایکاش این ناس باین ولیّ و مرآت که برای خود
من دون اللّه اخذ نموده‌اند مطّلع میشدند قسم بآفتاب عزّ صمدانی که اصل امر که
در بارهٴ مرآت شنیده‌اند بآن نحو نبوده و هر نفسی ادّعای علم نماید کذب بربّ
البقآء و آنچه نظر بحکمة اللّه این عبد مذکور داشته بین النّاس اشتهار یافته
و مقصود از کلمات نقطهٴ اولی جلّ جلاله را احدی ادراک ننموده و این عبد اصل
امر را از کلّ مستور داشته لحکمة لا یعلمها الّا نفسی العلیم الحکیم فواللّه
فعل بموجده الّذی خلقه بنفخة من عنده ما لا فعل احد فی العالمین و اگر خلق
لوحی از الواح بدیعه که از سماء عزّ احدیّه در این ایّام نازل شده ببصر اللّه
ملاحظه نمایند از حقّ محتجب نمیمانند و ماسوایش را معدوم مشاهده میکنند و
موقن میشوند که مادونش غیر مذکور بوده و خواهد بود تا چه رسد بمعرضین

\end{otherlanguage}

Ḥāl dar īn āyeh-ye munazzaleh tafakkur namāʾīd va inṣāf dihīd va lā
takūnū min alladhīna hum yaʿrifūna niʿmata Allāh thumma yunkirūn. Va ey
kāsh īn nās be-īn valī va mirʾāt keh barā-ye khud min dūn Allāh akhdh
namūdeh-and muṭṭaliʿ mīshudand. Qasam be-āftāb-e ʿizz-e ṣamadānī keh
aṣl-e amr keh dar bāreh-ye mirʾāt shinīdeh-and be-ān naḥv nabūdeh va har
nafsī iddiʿā-ye ʿilm namāyad kadhaba bi-rabbi al-baqāʾ va ānchih naẓar
be-ḥikmat Allāh īn ʿabd madhkūr dāshteh bayn al-nās ishtihār yāfteh va
maqṣūd az kalimāt-e Nuqṭeh-ye Ūlā jalla jalāluhu rā aḥadī idrāk
nanamūdeh va īn ʿabd aṣl-e amr rā az kull mastūr dāshteh li-ḥikmatin lā
yaʿlamuhā illā nafsī al-ʿalīm al-ḥakīm. Fa-wa-Allāh faʿala bi-mūjidihi
alladhī khalaqahu bi-nafkhatin min ʿindihi mā lā faʿala aḥadun fī
al-ʿālamīn va agar khalq lawḥī az alvāḥ-e badīʿeh keh az samā-ye ʿizz-e
aḥadiyyeh dar īn ayyām nāzil shudeh be-baṣar Allāh mulāḥiẓeh namāyand az
Ḥaqq muḥtajib namīmānand va mā-sivāyash rā maʿdūm mushāhideh mīkunand va
mūqin mīshavand keh mā-dūnash ghayr-e madhkūr būdeh va khāhad būd tā
chih rasad be-muʿriḍīn.

Now reflect on this revealed verse and judge fairly, and do not be among
those who recognize God's favor and then deny it. If only these people
knew about this guardian and mirror whom they have taken for themselves
apart from God! I swear by the Sun of everlasting glory that the true
circumstances concerning the mirror were not as they have heard. Whoever
claims to know has lied, by the Lord of Eternity! What this Servant has
stated in accordance with God's wisdom has become widely known among the
people, but no one has understood the intent of the words of the Primal
Point---exalted be His glory. This Servant has kept the truth of the
matter concealed from everyone, for a wisdom that none knows but My own
Self, the All-Knowing, the All-Wise. By God, he has done to his Creator,
who created him through a breath from Himself, what no one in all the
worlds has done. If people were to look with the eye of God at a single
one of the wondrous Tablets that have descended in these days from the
heaven of the glory of oneness, they would not remain veiled from God.
They would see all besides Him as nonexistent and become certain that
whatever is other than Him has been and will remain unworthy of mention,
much less those who turn away.

\subsection{98}\label{section-1653}

\begin{otherlanguage}{persian}

و بعضی الیوم بعصمت بعضی من دون اللّه قائل شده‌اند چنانچه اهل فرقان من غیر
شعور تکلّم مینمودند و مقصود از عصمت را ابداً ادراک ننموده‌اند حکم عصمت
الیوم محقّق میشود هر نفسی که بعد از استماع کلمة اللّه و ندائه بکلمهٴ بلی
موقن شد از اهل عصمت بوده و من‌دون آن از عصمت خارج چنانچه نقطهٴ بیان روح
ما سواه فداه میفرمایند مخاطباً للعظیم فانّ الأمر قد رقت عن الحدودات انت
تصفی عباداً حین ما تجلّی اللّه لهم بهم قد عرفوا اللّه بارئهم و ما صبروا فیه
و ما شکوا حتّی اجعلنّهم مثل ما جعلت من قبل من الأنبیآء و الأوصیآء و
الشّهدآء و المقرّبین و لعمری لو تحضرنّ بعدد کلّ شیء لأجعلنّه و لا ینقص عن ملک
اللّه قدر شیء و لا یزید قدر شیء ولکن تری ینبغی لتلک الدّرجة العصمة الکبری
و لم یکن العصمة بما تری عند النّاس من احتیاطاتهم فی دینهم لأنّهم حین ما
سمعوا ندآء أ لست بربّکم ما قالوا بلی قد خرجوا و انّما العصمة للّذینهم حین
ما قال اللّه أ لست بربّکم قالوا بلی و انّ اللّه لم یقل لأحد الّا بمظهر نفسه
فی کلّ ظهور انتهی

\end{otherlanguage}

Va baʿżī al-yawm bi-ʿiṣmat-i baʿżī min dūni Allāh qāʾil shudih-and,
chunānchih ahl-i Furqān min ghayri shuʿūr takallum mīnamūdand va maqṣūd
az ʿiṣmat rā abadan idrāk nanamūdih-and. Ḥukm-i ʿiṣmat al-yawm muḥaqqaq
mīshavad: har nafsī kih baʿd az istimāʿ-i Kalimatuʾllāh va nidāʾih
bi-kalimiy-i balā mūqin shud az ahl-i ʿiṣmat būdih va min dūn-i ān az
ʿiṣmat khārij. Chunānchih Nuqṭiy-i Bayān, rūḥ-i mā sivāh fidāh,
mīfarmāyand mukhāṭiban li-l-ʿAẓīm: fa-inna al-amra qad raqqat ʿan
al-ḥudūdāt, anta tuṣaffī ʿibādan ḥīna mā tajallā Allāhu lahum bihim qad
ʿarafū Allāha bāriʾahum wa-mā ṣabarū fīhi wa-mā shakkū ḥattā
ajʿalannahum mithla mā jaʿalta min qablu min al-anbiyāʾ wa-l-awṣiyāʾ
wa-l-shuhadāʾ wa-l-muqarrabīn. Wa-la-ʿamrī law tuḥḍiranna bi-ʿadadi
kulli shayʾ la-ajʿalannahu wa-lā yanquṣu ʿan mulki Allāhi qadru shayʾ
wa-lā yazīdu qadru shayʾ, wa-lākin tarā yanbaghī li-tilka al-darajati
al-ʿiṣmatu al-kubrā. Wa-lam yakun al-ʿiṣmatu bi-mā tarā ʿinda al-nāsi
min iḥtiyāṭātihim fī dīnihim li-annahum ḥīna mā samiʿū nidāʾa a-lastu
bi-rabbikum mā qālū balā qad kharajū. Wa-innamā al-ʿiṣmatu
li-lladhīnahum ḥīna mā qāla Allāhu a-lastu bi-rabbikum qālū balā.
Wa-inna Allāha lam yaqul li-aḥadin illā bi-maẓhari nafsihi fī kulli
ẓuhūr. Intahā.

Some today have attributed infallibility to certain people apart from
God, just as the people of the Qur'an used to speak without
understanding; they have never understood what is meant by
infallibility. Today the criterion of infallibility is established:
whoever, after hearing the Word of God and His call, has become assured
of the word ``Yes'' belongs to the people of infallibility, and anyone
else is outside it. Thus the Point of the Bayán---may the spirit of all
besides Him be sacrificed for Him---says, addressing ʿAẓím: ``The Cause
has indeed risen beyond limitations. You are to select servants who,
when God revealed Himself to them, recognized through them God, their
Maker, and did not hesitate concerning Him or entertain doubt, that I
may make them like those whom You previously made prophets, successors,
martyrs, and those brought near. By My life, were you to bring forward
as many as the number of all things, I would surely do so, and nothing
whatsoever would diminish from God's dominion or increase it. But you
see that supreme infallibility befits that station. Infallibility is not
what you see among people in their precautions concerning their
religion, for when they heard the call `Am I not your Lord?' and did not
say `Yes,' they departed from it. Infallibility belongs only to those
who, when God said `Am I not your Lord?' said `Yes.' God has never
spoken to anyone except through the Manifestation of His own Self in
every Revelation.'' End of quotation.

\subsection{99}\label{section-1654}

\begin{otherlanguage}{persian}

و بعد از این بیان که اصرح کلّ کلمات من قلم اللّه نازل شده دیگر که میتواند
برای خود عصمت ثابت نماید الّا بعد از تصدیق این امر اعظم افخم الیوم کلّ ناس
فی ایّ رتبة کان از عصمت خارجند مگر آن نفوسی که باین ظهور قدس صمدانی موقن
شده‌اند ای عباد بشنوید ندآء اللّه را و آذان فطرت را از استماع کلمات اهل
فرقان طاهر نموده تا بیان ابدع لسان رحمن را ادراک نمائید آخر تفکّر نمائید
این فئه در کدام قول صادق بوده‌اند هزار و دویست و شصت سنه بختم ظهورات قائل
بوده و هم‌چنین بکون قائم علی هیکل مخصوص در ارض معیّنه و هم‌چنین در علائم
ظهور که بهزار روایات تمسّک جسته بشأنی که منکرین این اقوال را کافر
میدانستند و بعد معلوم شد در ظهور نقطهٴ بیان جلّ ذکره که جمیع مخطی بوده و
مظاهر عزّ احدیّه لم‌یزل طالع و لایزال مشرق خواهند بود و هم‌چنین در سایر
اقوال این هیاکل اضلال ملاحظه کن که کلّ توهّم بوده و خواهد بود معبودهم
هواهم و مسجودهم اوهامهم ان انتم تعرفون

\end{otherlanguage}

Va baʿd az īn bayān kih aṣraḥ-i kull-i kalimāt min Qalamiʾllāh nāzil
shudih, dīgar kih mītavānad barāy-i khud ʿiṣmat thābit namāyad illā baʿd
az taṣdīq-i īn Amr-i aʿẓam-i afkham? Al-yawm kull-i nās fī ayyi rutbatin
kān az ʿiṣmat khārijand magar ān nufūsī kih bi-īn ẓuhūr-i quds-i
ṣamadānī mūqin shudih-and. Ay ʿibād, bishnavīd nidā-yi Allāh rā va
ādhān-i fiṭrat rā az istimāʿ-i kalimāt-i ahl-i Furqān ṭāhir namūdih tā
bayān-i abdaʿ-i Lisān-i Raḥmān rā idrāk namāyīd. Ākhir tafakkur namāyīd,
īn fiʾih dar kudām qawl ṣādiq būdih-and? Hizār va duvīst va shaṣt sanih
bi-khatm-i ẓuhūrāt qāʾil būdih va hamchunīn bi-kawn-i Qāʾim ʿalā
haykal-i makhṣūṣ dar arḍ-i muʿayyanih va hamchunīn dar ʿalāʾim-i ẓuhūr
kih bi-hizār rivāyāt tamassuk justih bi-shaʾnī kih munkirīn-i īn aqvāl
rā kāfir mīdānistand. Va baʿd maʿlūm shud dar ẓuhūr-i Nuqṭiy-i Bayān,
jalla dhikruhu, kih jamīʿ mukhṭī būdih va maẓāhir-i ʿizz-i aḥadiyyih
lam-yazal ṭāliʿ va lā-yazāl mushriq khvāhand būd. Va hamchunīn dar
sāʾir-i aqvāl-i īn hayākil-i iḍlāl mulāḥiẓih kun kih kull tavahhum būdih
va khvāhad būd. Maʿbūduhum hawāhum wa-masjūduhum awhāmuhum in antum
taʿrifūn.

After this statement, the most explicit of all utterances, has descended
from the Pen of God, who can any longer claim infallibility for himself
except after acknowledging this greatest, most exalted Cause? Today all
people, whatever their rank, are outside infallibility except those
souls who have become assured of this holy, eternal Revelation. O
servants, hear the call of God and cleanse the ears of your innate
nature from listening to the words of the people of the Qur'an, so that
you may comprehend the wondrous utterance of the Tongue of the
All-Merciful. Reflect, after all: in which of their statements have
these people been truthful? For one thousand two hundred and sixty years
they maintained that Revelations had ended, and likewise that the Qáʾim
existed in a particular bodily form in a specified land. Likewise,
concerning the signs of His appearance, they clung to a thousand
traditions, to such an extent that they considered those who denied
these assertions unbelievers. Then, in the Revelation of the Point of
the Bayán---exalted be His mention---it became clear that all had been
mistaken, and that the Manifestations of the glory of oneness have
always risen and will forever continue to shine. Consider likewise the
other statements of these embodiments of error: all have been and will
remain mere illusion. Their god is their own desire, and the object of
their worship their own imaginings, if you understand.

\subsection{100}\label{section-1655}

\begin{otherlanguage}{persian}

مع آنکه جمیع اهل بیان مشاهده نموده‌اند که آنچه در دست این فئه بود غیر حقّ
بوده مع‌ذلک متّصلاً بکلمات اهل فرقان استدلال مینمایند و متکلّمند چنانچه تازه
این فئه بذکر وصایت قبل میخواهند ناس را از فرات قدس الهی و شریعهٴ عزّ
صمدانی ممنوع نمایند قسم بجمال اللّه العلیّ الأبهی که اگر هزار مظهر ظهور در
ظهور نقطهٴ بیان موجود بود و در اقرار بر امر ملیک مختار اقلّ من حین اصطبار
مینمود از اهل نار عنداللّه محسوب میشد

\end{otherlanguage}

Maʿa ānkih jamīʿ-i ahl-i Bayān mushāhidih namūdih-and kih ānchih dar
dast-i īn fiʾih būd ghayr-i ḥaqq būdih, maʿa-dhālik muttaṣilan
bi-kalimāt-i ahl-i Furqān istidlāl mīnamāyand va mutakallimand.
Chunānchih tāzih īn fiʾih bi-dhikr-i viṣāyat-i qabl mīkhvāhand nās rā az
Furāt-i quds-i ilāhī va sharīʿiy-i ʿizz-i ṣamadānī mamnūʿ namāyand.
Qasam bi-Jamāl-i Allāh al-ʿAliyy al-Abhā kih agar hizār maẓhar-i ẓuhūr
dar ẓuhūr-i Nuqṭiy-i Bayān mawjūd būd va dar iqrār bar Amr-i Malīk-i
Mukhtār aqall-i min ḥīn iṣṭibār mīnamūd az ahl-i nār ʿinda Allāh maḥsūb
mīshud.

Although all the people of the Bayán have seen that what these people
possessed was untrue, they nevertheless continually base their arguments
and their discourse on the words of the people of the Qur'an. Thus these
people now seek, by invoking the earlier appointment of a successor, to
bar people from the Euphrates of divine holiness and the waterway of
eternal glory. I swear by the Beauty of God, the Most High, the
All-Glorious, that if a thousand manifestations of revelation had
existed in the Revelation of the Point of the Bayán, and one had delayed
for less than an instant in acknowledging the Cause of the sovereign
King who chooses freely, he would have been reckoned among the people of
the fire in God's sight.

\subsection{101}\label{section-1656}

\begin{otherlanguage}{persian}

خرق کنید حجبات وهم را و بسماء عزّ ابهی عروج نمائید و اگر میخواهید حقّ بر
جمیع معلوم شود چند نفر از منصفین حاضر شوند در این مدینه و آثار اللّه و
کلماتش را ملاحظه نمایند و هم‌چنین مدّعیان هم در این ارض حاضرند درست
استفسار نمایند تا آثار حقّ از دونش مبرهن گردد و بر جمیع محقّق شود که دونه
فقرآء لدی بابه و عجزآء عند حضرته و فقدآء لدی ظهورات انواره

\end{otherlanguage}

Kharq kunīd ḥujubāt-i vahm rā va bi-samā-yi ʿizz-i Abhā ʿurūj namāyīd.
Va agar mīkhvāhīd Ḥaqq bar jamīʿ maʿlūm shavad chand nafar az munṣifīn
ḥāḍir shavand dar īn madīnih va āthāruʾllāh va kalimātash rā mulāḥiẓih
namāyand va hamchunīn muddaʿiyān ham dar īn arḍ ḥāḍirand. Durust
istifsār namāyand tā āthār-i Ḥaqq az dūnash mubarhan gardad va bar jamīʿ
muḥaqqaq shavad kih dūnahu fuqarāʾ ladā bābihi wa-ʿujazāʾ ʿinda
ḥaḍratihi wa-fuqadāʾ ladā ẓuhūrāti anwārihi.

Tear through the veils of illusion and ascend to the heaven of
all-glorious might. If you wish the Truth to become known to everyone,
let a few fair-minded people come to this city and examine the writings
of God and His words; the claimants are also present in this land. Let
them inquire thoroughly, so that the writings of the True One may be
clearly distinguished from those of others, and everyone may be certain
that all besides Him are paupers at His gate, powerless in His presence,
and destitute before the revelations of His light.

\subsection{102}\label{section-1657}

\begin{otherlanguage}{persian}

در بیان فارسی لسان الهی فرموده قوله تعالی منتهی جدّ و جهد خود را نموده که
در یوم ظهور حقّ از کلمات او بر او احتجاج نگردد که کلّ بیان کلام ظهور قبل
او است و او است عالم‌تر از آنچه نازل فرموده از کلّ خلق زیرا که روح کلّ در
قبضهٴ او است و در نزد کلّ نیست الّا شبحی اگر بر حقّ مستقرّ باشند والّا که
لایق ذکر نیست انتهی

\end{otherlanguage}

Dar Bayān-i Fārsī Lisān-i Ilāhī farmūdih, qawluhu taʿālā: muntahā-yi
jidd va jahd-i khud rā namūdih kih dar yawm-i ẓuhūr-i Ḥaqq az kalimāt-i
ū bar ū iḥtijāj nagardad kih kull-i Bayān kalām-i ẓuhūr-i qabl-i ū ast
va ū ast ʿālimtar az ānchih nāzil farmūdih az kull-i khalq zīrā kih
rūḥ-i kull dar qabḍiy-i ū ast va dar nazd-i kull nīst illā shabaḥī agar
bar Ḥaqq mustaqarr bāshand va-illā kih lāyiq-i dhikr nīst. Intahā.

In the Persian Bayán the Divine Tongue has spoken---exalted are His
words: ``Make your utmost effort and exertion so that, on the day of the
Revelation of the True One, His words may not be used to argue against
Him. For the whole Bayán is the utterance of His previous Revelation,
and He knows what He has revealed better than all creation, for the
spirit of all is in His grasp. All others possess only a semblance, if
they remain steadfast in the Truth; otherwise, it is not worthy of
mention.'' End of quotation.

\subsection{103}\label{section-1658}

\begin{otherlanguage}{persian}

این بیان جمال رحمن حال ملاحظه نمائید با اینکه کل را منع فرموده‌اند از
احتجاج بآن جوهر وجود و ساذج معبود و صریح فرموده‌اند بکلمات بیان بمنزل آن
در ظهور بعد احتجاج ننمائید مع‌ذلک الیوم مشاهده میشود که کلّ باحتجاج
برخاسته‌اند و میگویند فلان را در بیان باسماء حسنی موسوم نموده‌اند مع‌ذلک
چگونه میشود از این مقام تنزّل نماید

\end{otherlanguage}

Īnast bayān-i Jamāl-i Raḥmān. Ḥāl mulāḥiẓih namāyīd bā īnkih kull rā
manʿ farmūdih-and az iḥtijāj bi-ān Jawhar-i Vujūd va Sādhaj-i Maʿbūd va
ṣarīḥ farmūdih-and bi-kalimāt-i Bayān bi-Munzil-i ān dar ẓuhūr-i baʿd
iḥtijāj nanamāyīd, maʿa-dhālik al-yawm mushāhidih mīshavad kih kull
bi-iḥtijāj barkhāstih-and va mīgūyand fulān rā dar Bayān bi-asmā-yi
ḥusnā mawsūm namūdih-and, maʿa-dhālik chigūnih mīshavad az īn maqām
tanazzul namāyad?

Such is the utterance of the Beauty of the All-Merciful. Now consider:
although He has forbidden everyone to argue with that Essence of being,
that pure Being who is worshipped, and has explicitly said, ``Do not use
the words of the Bayán to argue with its Revealer in the next
Revelation,'' nevertheless today everyone can be seen rising to argue,
saying, ``Such a person has been designated in the Bayán by the most
beautiful names; how, then, can he fall from this station?''

\subsection{104}\label{section-1659}

\begin{otherlanguage}{arabic}

و هم‌چنین میفرماید عزّ ذکره او را شناخته بآیات او و احتیاط در عرفان او
نکرده که بقدر همان در نار خواهی بود انتهی و نظر باین کلمات محکمه الیوم
کلّ محتاطین در این امر امنع اقدس در نارند لا شکّ فی ذلک

\end{otherlanguage}

Va hamchunīn mīfarmāyad, ʿazza dhikruhu: ū rā shinākhtih bi-āyāt-i ū va
iḥtiyāṭ dar ʿirfān-i ū nakardih kih bi-qadr-i hamān dar nār khvāhī būd.
Intahā. Va naẓar bi-īn kalimāt-i muḥkamih al-yawm kull-i muḥtāṭīn dar īn
Amr-i amnaʿ-i aqdas dar nārand, lā shakka fī dhālik.

He likewise says---exalted be His mention: ``Recognize Him through His
verses, and exercise no caution in recognizing Him, for to that same
extent you will be in the fire.'' End of quotation. In view of these
decisive words, all who exercise caution today concerning this most
unassailable, most holy Cause are in the fire. There is no doubt of
this.

\subsection{105}\label{section-1660}

\begin{otherlanguage}{persian}

و قوله عزّ ذکره و اگر در مابین خود و خدا توجّه میکنید مثال او است که در
افئدهٴ شما است باو از او محتجب نگردید و بشناسید کسی را که از برای عرفان
او خلق شده‌اید و قوله تعالی در نزد هر ظهوری از ظهور شجرهٴ حقیقت مؤمنین
باو و کتاب او از قبل بظهور او و کتاب او از بعد محتجب میشوند بحیثیّتی که
نمیماند الّا مؤمن خالص و او اعزّ از کبریت احمر است انتهی

\end{otherlanguage}

Va qawluhu, ʿazza dhikruhu: va agar dar mābayn-i khud va Khudā tavajjuh
mīkunīd mithāl-i ū ast kih dar afʾidiy-i shumā ast, bi-ū az ū muḥtajib
nagardīd va bishināsīd kasī rā kih az barāy-i ʿirfān-i ū khalq
shudih-īd. Va qawluhu taʿālā: dar nazd-i har ẓuhūrī az ẓuhūr-i
Shajariy-i Ḥaqīqat muʾminīn bi-ū va kitāb-i ū az qabl bi-ẓuhūr-i ū va
kitāb-i ū az baʿd muḥtajib mīshavand bi-ḥaythiyyatī kih nimīmānad illā
muʾmin-i khāliṣ va ū aʿazz az kibrīt-i aḥmar ast. Intahā.

And His words---exalted be His mention: ``If, in private between
yourselves and God, you turn toward Him, it is His image that is in your
hearts. Do not let it veil you from Him, and recognize the One for whose
recognition you have been created.'' And His words---exalted is He: ``At
every appearance of the Tree of Truth, those who believed in Him and His
Book before become veiled from His subsequent Revelation and His Book,
to such an extent that none remains but the pure believer, who is rarer
than red sulfur.'' End of quotation.

\subsection{106}\label{section-1661}

\begin{otherlanguage}{persian}

چنانچه الیوم کلّ ممتحن شده‌اند و اکثری از بریّه از سلطان احدیّه اعراض
نموده‌اند و میگویند فلان در بیان باسم اللّه نامیده شده و هم‌چنین بامثال این
کلمات از سلطان عزّ لایزال محروم مانده و ملتفت نشده آنچه را عامل شده‌اند و
بر فرض تسلیم غافلند از اینکه کلّ اسماء در آن ساحت اقدس در صقع واحدند عطا
میفرماید کیف یشآء و اخذ میفرماید کیف یرید

\end{otherlanguage}

Chunānchih al-yawm kull mumtaḥan shudih-and va aktharī az bariyyih az
Sulṭān-i Aḥadiyyih iʿrāḍ namūdih-and va mīgūyand fulān dar Bayān
bi-ism-i Allāh nāmīdih shudih va hamchunīn bi-amthāl-i īn kalimāt az
Sulṭān-i ʿizz-i lā-yazāl maḥrūm māndih va multafit nashudih ānchih rā
ʿāmil shudih-and. Va bar farḍ-i taslīm ghāfiland az īnkih kull-i asmā
dar ān sāḥat-i aqdas dar ṣuqʿ-i vāḥidand. ʿAṭā mīfarmāyad kayfa yashāʾ
va akhdh mīfarmāyad kayfa yurīd.

Thus today all have been tested, and most of humanity have turned away
from the Sovereign of oneness, saying, ``Such a person has been called
by the name of God in the Bayán.'' Through such words they have remained
deprived of the Sovereign of everlasting glory, unaware of what they
have done. Even granting their assertion, they are heedless that all
names stand on a single plane in that most holy court. He bestows as He
wills and takes away as He desires.

\subsection{107}\label{section-1662}

\begin{otherlanguage}{persian}

قوله تعالی چه که کسی عالم بظهور نیست غیر اللّه هر وقت شود باید کلّ تصدیق
بنقطهٴ حقیقت نمایند و شکر الهی بجا آورند انتهی مع‌ذلک بعضی معارضه در قرب
ظهور نموده چنانچه ملل قبل کلّ باین معارضات مشغولند قل ویل لکم چه زود
تبدیل نموده‌اید نعمت الهی را و بسبل وهم و خطا سالک شده‌اید با آنکه هنوز از
ظهور چیزی نگذشته که کلّ بقدم اوّل راجع شده

\end{otherlanguage}

Qawluhu taʿālā: chih kih kasī ʿālim bi-ẓuhūr nīst ghayr-i Allāh, har
vaqt shavad bāyad kull taṣdīq bi-Nuqṭiy-i Ḥaqīqat namāyand va shukr-i
ilāhī bi-jā āvarand. Intahā. Maʿa-dhālik baʿżī muʿāriḍih dar qurb-i
ẓuhūr namūdih chunānchih milal-i qabl kull bi-īn muʿāriḍāt mashghūland.
Qul waylun lakum, chih zūd tabdīl namūdih-īd niʿmat-i ilāhī rā va
bi-subul-i vahm va khaṭā sālik shudih-īd, bā ānkih hanūz az ẓuhūr chīzī
nagudhashtih kih kull bi-qadam-i avval rājiʿ shudih.

His words---exalted is He: ``For no one knows of the Revelation except
God. Whenever it occurs, all must acknowledge the Point of Truth and
offer thanks to God.'' End of quotation. Nevertheless, some have
objected to the nearness of the Revelation, just as the peoples of
earlier religions are all occupied with such objections. Say: Woe to
you! How quickly you have exchanged God's blessing and taken the paths
of illusion and error! Although scarcely any time has passed since the
Revelation, all have returned to their starting point.

\subsection{108}\label{section-1663}

\begin{otherlanguage}{arabic}

باری حقّ لم‌یزل مختار بوده و لایزال مختار خواهد بود بهر نحو بخواهد ظاهر
میشود و انّه لهو المختار و ما سواه مقبوض فی قبضة قدرته المهیمن القیّوم
فوالّذی نفسی بیده لم یکن الیوم ظلم اعظم من انّ الّذی ینطق فی کلّ شیء بأنّی
انا اللّه لا اله الّا هو اراد ان یثبت لعباده بأنّه یکون مقتدراً بأن یبدّل
اسماً من اسمائه بعد الّذی یکون ملکوت الأسمآء خاضعة لسلطنته و جبروت الصّفات
مشفقة من خشیته و لاهوت العمآء منقادة لحضرته و منجعلة بارادته و بدعت
بأمره و یرجع الیه

\end{otherlanguage}

Bārī Ḥaqq lam-yazal mukhtār būdih va lā-yazāl mukhtār khvāhad būd,
bi-har naḥv bikhvāhad ẓāhir mīshavad. Wa-innahu la-huwa al-Mukhtār wa-mā
siwāhu maqbūḍun fī qabḍati qudratihi al-Muhaymin al-Qayyūm. Fa-wa-lladhī
nafsī bi-yadihi lam yakun al-yawma ẓulmun aʿẓamu min anna alladhī
yanṭiqu fī kulli shayʾ bi-annī anā Allāhu lā ilāha illā huwa arāda an
yuthbita li-ʿibādihi bi-annahu yakūnu muqtadiran bi-an yubaddila isman
min asmāʾihi baʿda alladhī yakūnu malakūtu al-asmāʾ khāḍiʿatan
li-salṭanatihi wa-jabarūtu al-ṣifāt mushfiqatan min khashyatihi
wa-lāhūtu al-ʿamāʾ munqādatan li-ḥaḍratihi wa-munjaʿilatan bi-irādatihi
wa-budiʿat bi-amrihi wa-yarjiʿu ilayhi.

In short, the True One has always been free to choose and will forever
remain so. He appears in whatever manner He wishes. He is indeed the One
who chooses freely, and all besides Him is held in the grasp of His
power, the All-Protecting, the Self-Subsisting. By Him in whose hand is
My soul! No injustice today is greater than this: that He who proclaims
in all things, ``I am God; there is no god but Him,'' should have to
establish for His servants that He has the power to change one of His
names, when the kingdom of names is subject to His sovereignty, the
dominion of attributes trembles in awe of Him, and the divine realm of
the cloud submits to His presence, is fashioned by His will, was brought
into being by His command, and returns to Him.

\subsection{109}\label{section-1664}

\begin{otherlanguage}{persian}

اینست شأن این خلق نابالغ ناقابل مقتدری که کلّ اسماء خلق شده بارادهٴ او و
کلّ صفات ظاهر شده بسلطان مشیّت او و اهل ملأ اعلی طائفند حول او باید
استدلال نماید بر اینکه قادر است از نفسی قمیص اسمی از اسماء را انتزاع
فرماید مع آنکه اطّلاع بر آن اسم نداشته و ندارند و مریبین در قدرت سلطان
یفعل ما یشآء و این ظهور ابدع امنع اعلی ید اللّه را مغلول دانسته‌اند ولو
بلسان اقرار ننمایند بذلک یشهد لهم لسان المختار ولو هم ینکرون قوله جلّ
ثنائه

\end{otherlanguage}

Īnast shaʾn-i īn khalq-i nābāligh-i nāqābil. Muqtadirī kih kull-i asmā
khalq shudih bi-irādiy-i ū va kull-i ṣifāt ẓāhir shudih bi-sulṭān-i
mashiyyat-i ū va ahl-i Malaʾ-i Aʿlā ṭāʾifand ḥawl-i ū bāyad istidlāl
namāyad bar īnkih qādir ast az nafsī qamīṣ-i ismī az asmā rā intizāʿ
farmāyad, maʿa ānkih iṭṭilāʿ bar ān ism nadāshtih va nadārand. Va
murībīn dar qudrat-i Sulṭān-i yafʿalu mā yashāʾ va īn ẓuhūr-i abdaʿ-i
amnaʿ-i aʿlā Yaduʾllāh rā maghlūl dānistih-and va-law bi-lisān iqrār
nanamāyand. Bi-dhālika yashhadu lahum Lisānu al-Mukhtār wa-law hum
yunkirūn. Qawluhu jalla thanāʾuhu.

Such is the condition of these immature, unworthy people. The
All-Powerful One, by whose will all names have been created, through
whose sovereign purpose all attributes have appeared, and around whom
the Concourse on high circles, must offer proof that He is able to strip
a soul of the robe of one of the names---even though they have never
had, and still have no, knowledge of that name. Those who doubt the
power of the Sovereign who does whatever He wills and this most
wondrous, most unassailable, most exalted Revelation have regarded the
hand of God as shackled, even if they do not acknowledge it with their
tongues. To this the Tongue of the One who chooses freely bears witness
concerning them, though they deny it. His words---glorious be His
praise:

\subsection{110}\label{section-1665}

\begin{otherlanguage}{arabic}

اذ عند الخلق لم یکن الّا کلماتاً ما اطّلعوا بمرادک فیها و لذا قد احتجبوا
یوم ظهورک انتهی اینست که الیوم کلّ بکلمات بیان از منزل آن محتجب مانده‌اند
مع آنکه بمراد اللّه مطّلع نبوده و نخواهند بود

\end{otherlanguage}

Idh ʿinda al-khalqi lam yakun illā kalimātan mā iṭṭalaʿū bi-murādika
fīhā wa-li-dhā qad iḥtajabū yawma ẓuhūrika. Intahā. Īnast kih al-yawm
kull bi-kalimāt-i Bayān az Munzil-i ān muḥtajib māndih-and maʿa ānkih
bi-murādiʾllāh muṭṭaliʿ nabūdih va nakhvāhand būd.

``For people possessed nothing but words whose intended meaning they did
not understand, and therefore they became veiled on the day of Your
Revelation.'' End of quotation. Thus today all have remained veiled by
the words of the Bayán from its Revealer, although they have not known,
and will not know, what God intended.

\subsection{111}\label{section-1666}

\begin{otherlanguage}{persian}

چنانچه آن جمال قدم برؤسای بیان مرقوم فرموده‌اند که مراد اللّه را در کلمات
او از حرف حیّ که تلقاء وجه بوده سؤال نمائید چه که شما عارف بمراد اللّه
نیستید و حرف حیّ مذکور جناب آقا سیّد حسین بوده و هو استشهد فی سبیل مولاه و
فایز نشد نفسی بلقای او تا از معانی کلمات صمدانیّه استفسار نماید و مقصود
از این بیان که لسان رحمن فرموده آنکه رؤسا بدانند که عارف بر معنی کلمات
الهیّه نبوده و نیستند

\end{otherlanguage}

Chunānchih ān Jamāl-i Qidam bi-ruʾasā-yi Bayān marqūm farmūdih-and kih
murādiʾllāh rā dar kalimāt-i ū az Ḥarf-i Ḥayy kih tilqāʾ-i vajh būdih
suʾāl namāyīd, chih kih shumā ʿārif bi-murādiʾllāh nīstīd. Va Ḥarf-i
Ḥayy-i madhkūr Jināb-i Āqā Siyyid Ḥusayn būdih, wa-huwa istushhida fī
sabīli mawlāhu, va fāʾiz nashud nafsī bi-liqā-yi ū tā az maʿānī-yi
kalimāt-i ṣamadāniyyih istifsār namāyad. Va maqṣūd az īn bayān kih
Lisān-i Raḥmān farmūdih ānkih ruʾasā bidānand kih ʿārif bar maʿnī-yi
kalimāt-i ilāhiyyih nabūdih va nīstand.

Thus that Ancient Beauty wrote to the leaders of the Bayán: ``Ask the
Letter of the Living who has been in My presence about God's intention
in His words, for you do not know what God intends.'' The Letter of the
Living mentioned was Áqá Siyyid Ḥusayn, who was martyred in the path of
his Lord, and no one attained his presence to inquire into the meanings
of the eternal words. The purpose of this statement spoken by the Tongue
of the All-Merciful was that the leaders should know that they have
never understood, and do not understand, the meaning of the divine
words.

\subsection{112}\label{section-1667}

\begin{otherlanguage}{arabic}

و چون ظهور مبین قریب بود لذا سیّد مذکور بمقرّ اعلی ارتقا جست تفکّروا فی ذلک
یا اولی الأفکار

\end{otherlanguage}

Va chūn ẓuhūr-i mubīn qarīb būd li-dhā Siyyid-i madhkūr bi-maqarr-i aʿlā
irtiqā just. Tafakkarū fī dhālika yā ulī al-afkār.

And because the manifest Revelation was near, the aforementioned Siyyid
ascended to the highest abode. Reflect on this, O people of
understanding.

\subsection{113}\label{section-1668}

\begin{otherlanguage}{persian}

و قوله جلّ ثنائه و انّ کلّما ذکرت لک فی مقام الاستدلال رشح من طمطام الظّاهر
و ان اردت سرّ الفؤاد بحکم طلعة الباطن لا تشیر الیها الاشارة و لا یواریها
الحجبات اللّانهایة و لا یحتاج احد بذکر دلیلها لأنّها هو نفس الظّهور و تمام
البطون فسبحان اللّه عمّا یصفون انتهی جمالی را که میفرماید

\end{otherlanguage}

Va qawluhu jalla thanāʾuhu: «Wa inna kullamā dhakartu laka fī maqām
al-istidlāl rashḥun min ṭamṭām al-Ẓāhir, wa in aradta sirr al-fuʾād
bi-ḥukm ṭalʿat al-Bāṭin lā tushīru ilayhā al-ishāra wa lā yuwārīhā
al-ḥujubāt al-lānihāya wa lā yaḥtāju aḥad bi-dhikr dalīlihā li-annahā
huwa nafs al-ẓuhūr wa tamām al-buṭūn fa-subḥān Allāh ʿammā yaṣifūn.»
Intahā. Jamālī rā ke mīfarmāyad

And His words---exalted be His praise: ``Everything I have mentioned to
you by way of demonstration is a spray from the ocean of the Manifest.
And if you seek the secret of the heart, by virtue of the countenance of
the Hidden, no allusion can point to it, nor can endless veils conceal
it, nor does anyone need to mention a proof of it, for it is
manifestation itself and the fullness of concealment. Glorified be God
above what they describe.'' Thus ends the quotation. That Beauty of whom
He says

\subsection{114}\label{section-1669}

\begin{otherlanguage}{persian}

اشاره را بآن ساحت عزّ احدیّه سبیل نه اعزازاً لحضرته و اخباراً لصفآء کینونته
و لطافة رتبته مع‌ذلک بعد از آنکه بکلّ ظهورات الهیّه و تمام شئونات احدیّه شمس
جمال خود را از افق قضا ظاهر فرمود از کلّ جهات سهام کین بر جمال مبین ربّ
العالمین انداخته و اوّل من رمانی هو الّذی حفظته تحت جناحین فضلی و ربّیناه
تحت عبّی و علّمناه فی ایّامی و هو الّذی قتلنی بسیف الاعراض و سفک دمی و ضیّع
حرمة اللّه فی نفسی و انکر آیاتی بعد الّذی کان خاضعاً لحرف منها و خلق بنفحة
عنها و جاحد حقّ اللّه فی حقّی و حارب بنفس اللّه فی نفسی و افتی علی قتلی فی
الظّاهر و لمّا امنعه اللّه بسلطانه و اظهر خفیّات قلبه علی الّذینهم هاجروا فی
سبیل اللّه اذاً قام علی الافترآء و رجع فعله و ظلمه الی نفسی المظلوم
اطمیناناً من الّذینهم اتّبعوه

\end{otherlanguage}

Ishāra rā bi ān sāḥat-i ʿizz-i aḥadiyya sabīl na, iʿzāzan li-ḥaḍratihi
wa ikhbāran li-ṣafāʾ kaynūnatihi wa laṭāfat rutbatihi, maʿa dhālik baʿd
az ānke bi kull-i ẓuhūrāt-i ilāhiyya va tamām-i shuʾūnāt-i aḥadiyya
shams-i jamāl-i khud rā az ufuq-i qaḍā ẓāhir farmūd az kull-i jihāt
sihām-i kīn bar jamāl-i mubīn-i Rabb al-ʿālamīn andākhta. Wa awwalu man
ramānī huwa alladhī ḥafiẓtuhu taḥta janāḥayn faḍlī wa rabbaynāhu taḥta
ʿibbī wa ʿallamnāhu fī ayyāmī wa huwa alladhī qatalanī bi-sayf al-iʿrāḍ
wa safaka damī wa ḍayyaʿa ḥurmat Allāh fī nafsī wa ankara āyātī baʿda
alladhī kāna khāḍiʿan li-ḥarfin minhā wa khuliqa bi-nafḥatin ʿanhā wa
jāḥada ḥaqq Allāh fī ḥaqqī wa ḥāraba bi-nafs Allāh fī nafsī wa aftā ʿalā
qatlī fī al-ẓāhir wa lammā amnaʿahu Allāh bi-sulṭānihi wa aẓhara
khafiyyāt qalbihi ʿalā alladhīna hum hājarū fī sabīl Allāh idhan qāma
ʿalā al-iftirāʾ wa rajaʿa fiʿluhu wa ẓulmuhu ilā nafsī al-Maẓlūm
iṭmiʾnānan min alladhīna hum ittabaʿūhu.

that no allusion has access to that court of the glory of oneness, in
honor of His presence and to declare the purity of His being and the
subtlety of His rank---nevertheless, after He revealed the Sun of His
beauty from the horizon of decree with all divine manifestations and
every expression of oneness, they shot the arrows of malice from every
direction at the manifest Beauty of the Lord of the worlds. The first to
shoot at Me was the one whom I sheltered beneath the wings of My grace,
whom We nurtured in My embrace and taught in My days. He is the one who
killed Me with the sword of rejection, shed My blood, violated the
sanctity of God in My person, and denied My verses after having bowed
before a single letter of them and been created by a breath from them.
He repudiated God's right in My right, waged war against God's own Self
in My person, and openly decreed My death. When God restrained him
through His sovereignty and exposed the secrets of his heart to those
who had emigrated in the path of God, he resorted to slander, and his
deeds and oppression were attributed to My wronged Self, to reassure
those who followed him.

\subsection{115}\label{section-1670}

\begin{otherlanguage}{arabic}

تاللّه انفطرت سموات الأمر من فعله و شقّت اراضی القدس من ظلمه و قدّت ستر
حجاب الملکوت من اقواله و بکت سکّان مداین البقآء من اثمه و هو یضحک فی نفسه
کأنّه ما ارتکب فی الملک شیئاً تاللّه اذاً یبکی کلّ شیء ولکن انّه لا یکوننّ من
الشّاعرین فسوف یعلم حین الّذی یأتیه الموت و یدور عینتاه من سطوة القهر و
یقول هل لی من رجوع لأستغفر عمّا فعلت اذاً یضرب علی فمه طین الهاویة و کذلک
قدّر للمستکبرین علی اللّه المهیمن العزیز القدیر قوله عزّ ثنائه فسبحانک
اللّهمّ یا الهی ان تقل قد عفوت عن البیان و من فیه فاذاً کلّ فی غفرانک و
رضائک و ان تقل لا فانّنی انا الآخذ بأذیال جودک و المستشفع الیک بنفسک
انتهی

\end{otherlanguage}

Ta-llāh infaṭarat samāwāt al-amr min fiʿlihi wa shuqqat arāḍī al-quds
min ẓulmihi wa quddat sitr ḥijāb al-malakūt min aqwālihi wa bakat sukkān
madāʾin al-baqāʾ min ithmihi wa huwa yaḍḥaku fī nafsihi ka-annahu mā
irtakaba fī al-mulk shayʾan. Ta-llāh idhan yabkī kull shayʾ wa lākin
innahu lā yakūnanna min al-shāʿirīn. Fa-sawfa yaʿlamu ḥīna alladhī
yaʾtīhi al-mawt wa yadūru ʿaynatāhu min saṭwat al-qahr wa yaqūlu hal lī
min rujūʿ li-astaghfira ʿammā faʿaltu idhan yuḍrabu ʿalā famihi ṭīn
al-hāwiya wa kadhālika quddira li-l-mustakbirīn ʿalā Allāh al-Muhaymin
al-ʿAzīz al-Qadīr. Qawluhu ʿazza thanāʾuhu: «Fa-subḥānaka Allāhumma yā
ilāhī in taqul qad ʿafawtu ʿan al-Bayān wa man fīhi fa-idhan kullun fī
ghufrānika wa riḍāʾika wa in taqul lā fa-innanī anā al-ākhidh bi-adhyāl
jūdika wa al-mustashfiʿ ilayka bi-nafsika.» Intahā.

By God! The heavens of the Cause have been torn apart by his deeds, the
lands of holiness split asunder by his oppression, the covering of the
Kingdom's veil rent by his words, and the inhabitants of the cities of
eternity moved to tears by his sin, while he laughs to himself as though
he had done nothing in this world. By God! Everything weeps, yet he is
not among those who perceive. Soon he will know, when death comes to him
and his eyes roll beneath the assault of overpowering wrath, and he
says, ``Is there any return for me, that I may ask forgiveness for what
I have done?'' Then the clay of the abyss will be struck against his
mouth. Thus has it been ordained for those who arrogantly oppose God,
the All-Protecting, the Almighty, the All-Powerful. His words---exalted
be His praise: ``Glory be to You, O God, my God! If You say, `I have
pardoned the Bayán and those within it,' then all are embraced by Your
forgiveness and Your good pleasure. And if You say, `No,' then I am the
one who clings to the hem of Your generosity and seeks Your own Self as
an intercessor with You.'' Thus ends the quotation.

\subsection{116}\label{section-1671}

\begin{otherlanguage}{persian}

و شکّی نبوده که ذات قدم تکلّم نمینماید مگر بلسان مظهر نفس خود و مشاهده
نمائید که در این بیانات چه مقدار خضوع از محبوب امکان ظاهر شده که جمیع
بیان و ما فیه را بنعم و لای سلطان بقا در ظهور اخری منوط فرموده‌اند و حال
ملیک غیب و شهود بر جمیع مشهود و بشأنی ضرّ بر وجود مبارکش وارد شده که اگر
جمیع بحور غیب و شهود مداد شوند و کلّ من فی الملک اقلام و جمیع من فی
السّموات و الأرض راقم البتّه از ذکرش عاجز شوند و ایکاش که از اهل ایمانش
میشمردند فلعنة اللّه علی الظّالمین و المشرکین و مع‌ذلک بیان میخوانند و
استکتاب مینمایند فربّ تالی بیان و البیان یلعنه و ربّ عامل فیه و العمل یبرأ
منه و ربّ ذاکر و الذّکر یفرّ منه اعاذنا اللّه و ایّاکم یا ملأ الموحّدین من
هؤلآء و شرّهم و مکرهم و خدعهم

\end{otherlanguage}

Va shakkī nabūda ke Dhāt-i Qidam takallum namīnamāyad magar bi lisān-i
maẓhar-i nafs-i khud va mushāhada namāyīd ke dar īn bayānāt chi miqdār
khuḍūʿ az Maḥbūb-i imkān ẓāhir shuda ke jamīʿ-i Bayān va mā fīhi rā bi
naʿam va lā-yi Sulṭān-i Baqā dar ẓuhūr-i ukhrā manūṭ farmūda-and va ḥāl
Malīk-i ghayb va shuhūd bar jamīʿ mashhūd va bi shaʾnī ḍurr bar vujūd-i
mubārakash vārid shuda ke agar jamīʿ-i buḥūr-i ghayb va shuhūd midād
shavand va kull-i man fī al-mulk aqlām va jamīʿ-i man fī al-samāwāt wa
al-arḍ rāqim albatta az dhikrash ʿājiz shavand va ay kāsh ke az ahl-i
īmānash mīshumurdand. Fa-laʿnat Allāh ʿalā al-ẓālimīn wa al-mushrikīn.
Va maʿa dhālik Bayān mīkհwānand va istiktāb mīnamāyand. Fa-rubba tālī
Bayān wa al-Bayān yalʿanuhu wa rubba ʿāmil fīhi wa al-ʿamal yabraʾu
minhu wa rubba dhākir wa al-dhikr yaf irru minhu. Aʿādhanā Allāh wa
iyyākum yā malaʾ al-muwaḥḥidīn min hāʾulāʾ wa sharrihim wa makrihim wa
khidaʿihim.

There has never been any doubt that the Ancient Essence speaks only
through the tongue of the Manifestation of His own Self. Observe how
much humility the Beloved of all existence has shown in these
utterances: He has made the entire Bayán and all it contains dependent
on the ``Yes'' or ``No'' of the Sovereign of eternity in the next
Revelation. Now the King of the unseen and the seen is manifest to all,
yet such suffering has befallen His blessed person that, were all the
seas of the unseen and the seen to become ink, everyone in the world
pens, and all who are in the heavens and on earth scribes, they would
assuredly be powerless to describe it. If only they had counted Him
among His own believers! May God's curse be upon the oppressors and
those who associate partners with Him. Nevertheless, they read the Bayán
and have copies made of it. How many a reader of the Bayán is cursed by
the Bayán itself! How many a person acts according to it, yet the very
deed disowns him! How many a person makes mention, yet the Remembrance
flees from him! May God protect us and you, O company of those who
affirm His oneness, from these people, their evil, their scheming, and
their deceits.

\subsection{117}\label{section-1672}

\begin{otherlanguage}{arabic}

و هم‌چنین در مناجات با ملیک اسماء و صفات ذکر مینمایند فأسألک اللّهمّ بغنائک
عن کلّ شیء ان تقبلنّ البیان و من فیه فانّی قد سمّیت کلّما قد ظهر منّی علی ذلک
الاسم و جعلت کلّ ذلک کتاباً ورقاً من اوراق الحدایق الیک ان تقبلنّه یا محبوبی
فانّ ذلک من ظهور فضلک و رحمتک و ان تردّنّه یا محبوبی فانّ ذلک من ظهور عدلک و
غنآء ربوبیّتک انتهی

\end{otherlanguage}

Va hamchunīn dar munājāt bā Malīk-i asmāʾ va ṣifāt dhikr mīnamāyand:
«Fa-asʾaluka Allāhumma bi-ghināʾika ʿan kull shayʾ an taqbalanna
al-Bayān wa man fīhi fa-innī qad sammaytu kullamā qad ẓahara minnī ʿalā
dhālika al-ism wa jaʿaltu kulla dhālika kitāban waraqan min awrāq
al-ḥadāʾiq ilayka an taqbalannahu yā maḥbūbī fa-inna dhālika min ẓuhūr
faḍlika wa raḥmatika wa in taruddannahu yā maḥbūbī fa-inna dhālika min
ẓuhūr ʿadlika wa ghināʾ rubūbiyyatika.» Intahā.

Likewise, in a prayer to the King of names and attributes, He says: ``I
ask You, O God, by Your independence of all things, to accept the Bayán
and those within it, for I have called everything that has appeared from
Me by that name and made all of it a book, a leaf among the leaves of
the gardens offered to You. If You accept it, O my Beloved, that is a
manifestation of Your grace and mercy; and if You reject it, O my
Beloved, that is a manifestation of Your justice and the
self-sufficiency of Your lordship.'' Thus ends the quotation.

\subsection{118}\label{section-1673}

\begin{otherlanguage}{arabic}

مع آنکه جمیع بیان بقول آن کینونت سبحان و طلعت رحمن منوط بوده حال به
پست‌ترین خلق او مشغول و از جمالش محروم مانده‌اند فویل لهؤلآء ثمّ ویل لهؤلآء
تاللّه ما یمرّنّ هؤلآء علی شیء الّا و انّه یستعیذ باللّه منه و یبرأ عنه و
بلسان سرّه یقول یا ایّها المشرک الغافل بأیّ حجّة آمنت بظهور علیّ قبل نبیل و
بأیّ برهان کفرت بحسین قبل علیّ لا یا ملعون تاللّه ما آمنت من قبل بأحد من
رسل اللّه و لو آمنت فی عهد ذلک الایمان ما کان ایماناً حقیقیّاً بل کان شبحاً و
لو کان علی التّحقیق ما کفرت بالّذی به حقّق رسالة کلّ رسول و بعث کلّ نبیّ و
ایمان من فی السّموات و الأرضین

\end{otherlanguage}

Maʿa ānke jamīʿ-i Bayān bi qawl-i ān Kaynūnat-i Subḥān va Ṭalʿat-i
Raḥmān manūṭ būda ḥāl bi past-tarīn-i khalq-i ū mashghūl va az jamālash
maḥrūm mānda-and. Fa-waylun li-hāʾulāʾ thumma waylun li-hāʾulāʾ. Ta-llāh
mā yamurranna hāʾulāʾ ʿalā shayʾ illā wa innahu yastaʿīdhu bi-Allāh
minhu wa yabraʾu ʿanhu wa bi-lisān sirrihi yaqūlu yā ayyuhā al-mushrik
al-ghāfil bi-ayyi ḥujjatin āmanta bi-ẓuhūr ʿAlī qabla Nabīl wa bi-ayyi
burhānin kafarta bi-Ḥusayn qabla ʿAlī. Lā yā malʿūn ta-llāh mā āmanta
min qablu bi-aḥadin min rusul Allāh wa law āmanta fī ʿahdin dhālika
al-īmān mā kāna īmānan ḥaqīqiyyan bal kāna shabaḥan wa law kāna ʿalā
al-taḥqīq mā kafarta bi-alladhī bihi ḥuqqiqa risālat kull rasūl wa baʿth
kull nabī wa īmān man fī al-samāwāt wa al-araḍīn.

Although the entire Bayán was made dependent on the word of that
glorified Being and all-merciful Countenance, they are now occupied with
the lowliest of His creatures and remain deprived of His beauty. Woe to
these people, and again woe to them! By God! They pass by nothing that
does not seek refuge in God from them and disown them, saying in its
secret tongue, ``O heedless one who associates partners with God! By
what proof did you believe in the Revelation of `Alí before Nabíl, and
by what demonstration did you disbelieve in Ḥusayn before `Alí? No,
accursed one! By God! You never before believed in any of God's
Messengers. Even if you believed at some time, that faith was not true
faith but a mere semblance. Had it been genuine, you would not have
disbelieved in Him through whom the mission of every Messenger, the
raising up of every Prophet, and the faith of all who are in the heavens
and the earth have been established.''

\subsection{119}\label{section-1674}

\begin{otherlanguage}{persian}

ترجمهٴ این کلمات امنع صمدانی بلسان ابدع فارسی خالصاً لوجه اللّه ذکر میشود
که شاید نفوسی از آلایش نفس و هوی پاک شده بسماء عزّ بقا که مقام عرفان نفس
رحمن است عروج نمایند و له الأمر یهدی من یشآء الی جماله المقدّس المنیر

\end{otherlanguage}

Tarjuma-yi īn kalimāt-i amnaʿ-i ṣamadānī bi lisān-i abdaʿ-i fārsī
khāliṣan li-wajh Allāh dhikr mīshavad ke shāyad nufūsī az ālāyish-i nafs
va havā pāk shuda bi samā-yi ʿizz-i baqā ke maqām-i ʿirfān-i Nafs-i
Raḥmān ast ʿurūj namāyand. Wa lahu al-amr yahdī man yashāʾ ilā jamālihi
al-muqaddas al-munīr.

A translation of these supremely exalted, eternal words into the
wondrous Persian tongue is set forth solely for the sake of God, so that
perhaps some souls may be cleansed of the defilement of self and desire
and ascend to the heaven of everlasting glory, the station of
recognizing the Self of the All-Merciful. His is the command; He guides
whom He wills to His holy, radiant Beauty.

\subsection{120}\label{section-1675}

\begin{otherlanguage}{persian}

الیوم هر نفسی که از وجه معبود حقیقی اعراض نماید مرور نمینماید بر هیچ
شیئی مگر آنکه بلسان سرّ میگوید ای مشرک باللّه بکدام حجّت و برهان بجمال رحمن
در ظهور قبل ایمان آورده و بکدام دلیل از ملیک سبیل در این ظهور کافر شده
قسم بآفتاب عزّ تقدیس که لم‌یزل کافر و مشرک بوده و اگر در عهدی اظهار ایمان
نموده از روی حقیقت نبوده بلکه شبحی بوده در تو و در این ظهور اخذ شد دام
الایمان فی الایمان و رجع الشّیطان الی مقرّه فی اسفل الجحیم و اگر ایمانت
علی التّحقیق بود کافر نمیشدی بنفسی که باو محقّق شد رسالت کلّ رسل و بعث کلّ
انبیا و ایمان جمیع آنچه در آسمانها و زمینها است

\end{otherlanguage}

Al-yawm har nafsī ke az Vajh-i Maʿbūd-i Ḥaqīqī iʿrāḍ namāyad murūr
namīnamāyad bar hīch shayʾī magar ānke bi lisān-i sirr mīgūyad ay
mushrik bi-Allāh bi kudām ḥujjat va burhān bi Jamāl-i Raḥmān dar ẓuhūr-i
qabl īmān āvarda va bi kudām dalīl az Malīk-i sabīl dar īn ẓuhūr kāfir
shuda qasam bi āftāb-i ʿizz-i taqdīs ke lam-yazal kāfir va mushrik būda
va agar dar ʿahdī iẓhār-i īmān namūda az rū-yi ḥaqīqat nabūda balke
shabaḥī būda dar tu va dar īn ẓuhūr akhdh shud. Dāma al-īmān fī al-īmān
wa rajaʿa al-shayṭān ilā maqarrihi fī asfal al-jaḥīm. Va agar īmānat
ʿalā al-taḥqīq būd kāfir namīshudī bi nafsī ke bi ū muḥaqqaq shud
risālat-i kull-i rusul va baʿth-i kull-i anbiyā va īmān-i jamīʿ-i ānchi
dar āsimānhā va zamīnhā ast.

Today, any soul who turns away from the countenance of the true Object
of worship passes by nothing that does not say in its secret tongue, ``O
you who associate partners with God! By what proof and demonstration did
you believe in the Beauty of the All-Merciful in the previous
Revelation, and on what grounds have you disbelieved in the King of the
path in this Revelation? I swear by the Sun of the glory of holiness
that you have always been an unbeliever and one who associates partners
with God. If at any time you professed faith, it was not genuine;
rather, it was a semblance within you, and in this Revelation it was
taken away. Faith has remained in faith, and Satan has returned to his
abode in the lowest depths of hell. Had your faith been genuine, you
would not have disbelieved in the One through whom the mission of all
Messengers, the raising up of all Prophets, and the faith of everything
in the heavens and the earth have been established.''

\subsection{121}\label{section-1676}

\begin{otherlanguage}{arabic}

ای عباد قدری بقلب طاهر و نظر انصاف در امر اللّه ملاحظه نمائید که شاید از
حرم ربّانی محروم نشوید و از کعبهٴ قدس صمدانی ممنوع نگردید

\end{otherlanguage}

Ay ʿibād qadrī bi qalb-i ṭāhir va naẓar-i inṣāf dar amr Allāh mulāḥaẓa
namāyīd ke shāyad az ḥaram-i rabbānī maḥrūm nashavīd va az Kaʿba-yi
quds-i ṣamadānī mamnūʿ nagardīd.

O servants! Consider the Cause of God for a while with a pure heart and
an impartial eye, so that perhaps you may not be deprived of the divine
sanctuary or barred from the Kaaba of eternal holiness.

\subsection{122}\label{section-1677}

\begin{otherlanguage}{persian}

ای علیّ اگرچه تو سائلی و جواب باسم تو مرقوم میشود ولکن بسیار مشکل بنظر
میآید که الیوم مذعن بحقّ شوی و از کلّ ما ذکر فی اللّوح متذکّر گردی چه که
قلبت بحجبات غلیظه محتجب شده و بصرت بسبحات لانهایه مستور ولکن نظر اللّه و
لحظاته سرّاً بقلوب دیگر در این کلمات متوجّه است و عنقریب بقدرتی ظاهر شوند
که جمیع ما سوی اللّه را در ظلّ خود مشاهده نمایند ولکن نظر بسبقت رحمت
رحمانی و احاطهٴ فضل سبحانی باسم تو نازل شده لعلّ یأخذک نفحات ایّام اللّه و
ینقطعک عن العالمین

\end{otherlanguage}

Ay ʿAlī agarchi tu sāʾilī va javāb bi ism-i tu marqūm mīshavad va lākin
bisyār mushkil bi naẓar mīāyad ke al-yawm mudhʿin bi Ḥaqq shavī va az
kull-i mā dhukira fī al-lawḥ mutadhakkir gardī chi ke qalbat bi
ḥujubāt-i ghalīẓa muḥtajib shuda va baṣarat bi subuḥāt-i lānihāya mastūr
va lākin naẓar Allāh va laḥaẓātash sirran bi qulūb-i dīgar dar īn
kalimāt mutavajjih ast va ʿan-qarīb bi qudratī ẓāhir shavand ke jamīʿ-i
mā sivā Allāh rā dar ẓill-i khud mushāhada namāyand va lākin naẓar bi
sabqat-i raḥmat-i raḥmānī va iḥāṭa-yi faḍl-i subḥānī bi ism-i tu nāzil
shuda laʿalla yaʾkhudhuka nafaḥāt ayyām Allāh wa yanqaṭiʿuka ʿan
al-ʿālamīn.

O `Alí! Although you are the questioner and the answer is written in
your name, it seems most unlikely that today you will acknowledge the
Truth and take heed of all that has been mentioned in the Tablet, for
your heart is concealed by thick veils and your sight obscured by
endless veils of glory. Yet in these words God's gaze and glances are
secretly directed toward other hearts. Before long they will arise with
such power that they will behold everything other than God beneath their
shadow. Nevertheless, because His all-merciful compassion precedes all
things and His divine grace embraces all things, this has been revealed
in your name, so that perhaps the breezes of the days of God may seize
you and detach you from all the worlds.

\subsection{123}\label{section-1678}

\begin{otherlanguage}{arabic}

بیانه جلّ احسانه لم یزل لتضربنّ المثل بالشّمس و ما دونها بالمرایا کلّ مرآة
علی قدر ما تجلّت لها بها من نفسها الیها فاذاً یستدلّ علیها کلّ بها قائمون و
لو یرفع عنها آیة الّتی قد تجلّت لها بها فیها فاذاً لم یکن فی المرآة من شیء
یدلّ علی اللّه کذلک اذاً یرفع اللّه ما یرفع لن یحکمنّ من بعده بآیات افئدة
النّاس ولو انّ کلّهم یقولون انّا باللّه و بآیاته موقنون ما شهد اللّه علیهم بما
هم فیه ینطقون

\end{otherlanguage}

Bayānuhu jalla iḥsānuhu: «Lam yazal la-taḍribanna al-mathal bi-l-shams
wa mā dūnahā bi-l-marāyā kull mirʾāt ʿalā qadr mā tajallat lahā bihā min
nafsihā ilayhā fa-idhan yustadallu ʿalayhā kullun bihā qāʾimūn wa law
yurfaʿu ʿanhā āyat allatī qad tajallat lahā bihā fīhā fa-idhan lam yakun
fī al-mirʾāt min shayʾ yadullu ʿalā Allāh kadhālika idhan yarfaʿu Allāh
mā yarfaʿu lan yuḥkamanna min baʿdihi bi-āyāt afʾidat al-nās wa law anna
kullahum yaqūlūna innā bi-Allāh wa bi-āyātihi mūqinūn mā shahida Allāh
ʿalayhim bi-mā hum fīhi yanṭiqūn.»

His utterance---exalted be His beneficence: ``You will always draw the
comparison with the sun, and with mirrors for all that is other than it.
Each mirror, according to the measure in which the sun has shone upon it
through it, from itself toward it, then serves as evidence of the sun:
all subsist through it. Were the sign through which it has shone upon
the mirror, through it and within it, to be removed from it, there would
then be nothing in the mirror that points to God. Likewise, when God
removes what He removes, no judgment thereafter shall be based on the
signs in people's hearts. Even if they all say, `We are certain of God
and His signs,' God has not borne witness to them concerning what they
utter.''

\subsection{124}\label{section-1679}

\begin{otherlanguage}{persian}

ای مدهوشان خمر اوهام بشنوید ندای ربّ علیّ اعلی را که در این کلمات بالغه و
آیات محکمهٴ متقنه مرتفع شده که شاید بر خود و انفس عباد رحم نموده بر
سدرهٴ امریّهٴ الهیّه احجار ظنّیّهٴ بغضیّه میندازید و اگر هم خود مهتدی نشوید
عباد را از شریعهٴ ربّ الایجاد در این ایّام شداد محروم نسازید میفرمایند کلّ
مرآتند و حاکیند از شمس باندازه و مقام خود علی قدر تجلّی شمس در آن مرآت
بنفس مرآت در این وقت استدلال میشود باینکه کلّ بشمس قائم و موجودند و اگر
اخذ شود آیهٴ تجلّی شمس که در مرآت بنفس مرآت تجلّی فرموده باقی نمیماند در
مرآت شیئی تا دلالت کند بر مجلّی و خالق خود

\end{otherlanguage}

Ay madhūshān-i khamr-i awhām bishnavīd nidā-yi Rabb-i ʿAliyy-i Aʿlā rā
ke dar īn kalimāt-i bāligha va āyāt-i muḥkama-yi mutqana murtafiʿ shuda
ke shāyad bar khud va anfus-i ʿibād raḥm namūda bar Sidra-yi amriyya-yi
ilāhiyya aḥjār-i ẓanniyya-yi baghḍiyya mayandāzīd va agar ham khud
muhtadī nashavīd ʿibād rā az sharīʿa-yi Rabb al-ījād dar īn ayyām-i
shidād maḥrūm nasāzīd. Mīfarmāyand kull mirʾātand va ḥākī-and az shams
bi andāza va maqām-i khud ʿalā qadr-i tajallī-yi shams dar ān mirʾāt bi
nafs-i mirʾāt dar īn vaqt istidlāl mīshavad bi īnke kull bi shams qāʾim
va mawjūdand va agar akhdh shavad āya-yi tajallī-yi shams ke dar mirʾāt
bi nafs-i mirʾāt tajallī farmūda bāqī namīmānad dar mirʾāt shayʾī tā
dalālat kunad bar mujallī va khāliq-i khud.

O you who are stupefied by the wine of delusions! Hear the call of the
Lord, the Exalted, the Most High, raised in these compelling words and
these firm, unassailable verses, so that perhaps you may have mercy on
yourselves and on the souls of His servants and cease casting stones of
conjecture and hatred at the divine Lote-Tree of the Cause. Even if you
yourselves are not guided, do not deprive His servants of the water's
edge of the Lord of creation in these grievous days. He says that all
are mirrors and reflect the sun according to their own measure and
station. In proportion to the sun's shining in that mirror, the mirror
itself then serves as evidence that all subsist and exist through the
sun. If the sign of the sun's radiance, which has shone in the mirror
through the mirror itself, is taken away, nothing remains in the mirror
to point to the One who illumines and creates it.

\subsection{125}\label{section-1680}

\begin{otherlanguage}{arabic}

و با این کلمات محکمهٴ غیر متشابههٴ بالغه اهل بیان شبهه نموده‌اند که چگونه
میشود مرآت از تجلّی انوار شمس محروم ماند و در این رتبه بشأنی توقّف
نموده‌اند که ذکر آن ممکن نه و بر نفسی که از یک تجلّی از انوار اسم انورش
صدهزار مرایا مستضیء و مستنور میشوند و از یک تجلّی اسم آخذش کلّ بحدود نفس
خود مردود اعتراض مینمایند و از جمالش اعراض و بعد از اشراق شمس جمال در
قطب زوال میگویند انوار و دلیلش چیست

\end{otherlanguage}

Va bā īn kalimāt-i muḥkama-yi ghayr-i mutashābiha-yi bāligha ahl-i Bayān
shubha namūda-and ke chigūna mīshavad mirʾāt az tajallī-yi anvār-i shams
maḥrūm mānad va dar īn rutba bi shaʾnī tavaqquf namūda-and ke dhikr-i ān
mumkin na va bar nafsī ke az yak tajallī az anvār-i ism-i anvarash
ṣad-hazār marāyā mustaḍīʾ va mustanīr mīshavand va az yak tajallī-yi
ism-i Ākhidhash kull bi ḥudūd-i nafs-i khud mardūd iʿtirāḍ mīnamāyand va
az jamālash iʿrāḍ va baʿd az ishrāq-i shams-i jamāl dar quṭb-i zavāl
mīgūyand anvār va dalīlash chīst.

Despite these firm, unambiguous, compelling words, the people of the
Bayán have raised the objection: how can a mirror be deprived of the
radiance of the sun's light? They have become so arrested at this point
that it is impossible to describe. They object to the One through a
single radiance of whose most luminous name a hundred thousand mirrors
are illumined and made radiant, and through a single radiance of whose
name, the Withholder, all are cast back within the limits of their own
selves. They turn away from His beauty, and after the Sun of beauty has
risen, they remain at the point of its setting and ask, ``Where are its
rays, and what is its proof?''

\subsection{126}\label{section-1681}

\begin{otherlanguage}{persian}

ای بی‌خبران سر از نوم غفلت بردارید و انوارش که عالم را احاطه نموده مشاهده
کنید بعضی میگویند زود اشراق نموده ای بی‌بصران یا قریب و یا بعید حال اشراق
فرموده شما ملاحظه نمائید فی‌الحقیقه این اشراق محقّق است یا نه دیگر قرب و
بعد آن در دست شما و این غلام نبوده و نخواهد بود حکمت الهیّه که از انظر
بریّه مستور است اقتضا نموده یا قوم فارضوا بما رضی اللّه لکم و قضی علیکم

\end{otherlanguage}

Ey bī-khabarān, sar az nowm-e ghaflat bardārīd va anvārash ke ʿālam rā
eḥāṭe namūde moshāhede konīd. Baʿżī mīgūyand zūd eshrāq namūde. Ey
bī-baṣarān, yā qarīb va yā baʿīd, ḥāl eshrāq farmūde; shomā molāḥeẓe
namāyīd fī al-ḥaqīqe īn eshrāq moḥaqqaq ast yā na. Dīgar qorb va boʿd-e
ān dar dast-e shomā va īn gholām nabūde va nakhāhad būd. Ḥekmat-e
elāhiyye ke az {[}corrupted text{]} bariyye mastūr ast eqteżā namūde. Yā
qawm, fa-rḍaw bi-mā raḍiya Allāhu lakum wa-qaḍā ʿalaykum.

O heedless ones! Raise your heads from the sleep of negligence and
behold His lights, which have encompassed the world. Some say, ``He has
shone forth too soon.'' O you who lack sight! Whether sooner or later,
He has now shone forth. Consider whether this radiance is indeed real or
not. Its being sooner or later has never been, and will never be, in
your hands or those of this Youth. Divine wisdom, which is hidden from
the {[}corrupted original text{]} of humankind, has required it. O
people, be content with what God has willed for you and decreed for you.

\subsection{127}\label{section-1682}

\begin{otherlanguage}{arabic}

فواللّه لو کان الأمر بیدی ما اظهرت نفسی ابداً

\end{otherlanguage}

Fa-wa-Allāhi law kāna al-amru bi-yadī mā aẓhartu nafsī abadan.

By God, had the matter been in My hands, I would never have revealed
Myself.

\subsection{128}\label{section-1683}

\begin{otherlanguage}{persian}

ای صاحبان بغضا قسم بآفتاب فلک بقا که اگر امر بدست این عبد بود هرگز خود
را معروف نمینمودم چه که اسم مذکورم ننگ دارد از ذکر این السن غیر طاهرهٴ
کاذبه و در هر حین که ساکن شده و صمت اختیار نمودم روح القدس از یمینم ناطق
شده و روح الأعظم قدّام وجهم و روح الأمین فوق رأسم و روح البهآء در صدرم
ندا فرموده و حال اگر بسمع لطیف استماع شود از جمیع اعضا و احشا و عروق و
اظفار ندآء اللّه را استماع نمائید حتّی از شعراتم میشنوید بأنّه لا اله الّا
هو و انّ هذا الجمال لبهائه لمن فی السّموات و الأرضین

\end{otherlanguage}

Ey ṣāḥebān-e baghżā, qasam be-āftāb-e falak-e baqā ke agar amr be-dast-e
īn ʿabd būd hargez khod rā maʿrūf nemīnamūdam, che ke esm-e madhkūram
nang dārad az dhekr-e īn alson-e gheyr-e ṭāhere-ye kādhebe. Va dar har
ḥīn ke sāken shode va ṣamt ekhtiyār namūdam Rūḥ al-Qudus az yamīnam
nāṭeq shode va Rūḥ al-Aʿẓam qoddām-e vajham va Rūḥ al-Amīn fowq-e raʾsam
va Rūḥ al-Bahāʾ dar ṣadram nedā farmūde. Va ḥāl agar be-samʿ-e laṭīf
estemāʿ shavad az jamīʿ-e aʿżā va aḥshā va ʿorūq va aẓfār nedā-ye Allāh
rā estemāʿ namāyīd, ḥattā az shaʿarātam mīshenavīd bi-annahu lā ilāha
illā Huwa wa-inna hādhā al-jamāla la-bahāʾuhu li-man fī al-samāwāti
wa-al-araḍīn.

O bearers of hatred! I swear by the Sun of the heaven of eternity that,
had the matter been in the hands of this servant, I would never have
made Myself known, for My name is ashamed to be mentioned by these
impure, lying tongues. Whenever I have remained still and chosen
silence, the Holy Spirit has spoken at My right hand, the Most Great
Spirit before My face, the Faithful Spirit above My head, and the Spirit
of Glory within My breast. And now, if you listen with a discerning ear,
you will hear the call of God from all My limbs, My inmost parts, My
veins, and My nails. Even from My hairs you will hear that there is no
God but Him, and that this Beauty is indeed His glory for all who are in
the heavens and on the earth.

\subsection{129}\label{section-1684}

\begin{otherlanguage}{persian}

و لو کان هذا ذنبی تاللّه هذا لیس من عندی بل من لدن من ارسلنی و بعثنی
بالحقّ و جعلنی سراجاً للعالمین ای قوم فواللّه رایحهٴ صدق از کذب بسی معلوم و
واضح است مپسندید بر نفس حقّ آنچه بر خود نمی‌پسندید آخر قدری تأمّل نمائید
مقصود این غلام در این آخر ایّام در این سجن کبری چه بوده و چه خواهد بود
فتأمّلوا یا ملأ الغافلین

\end{otherlanguage}

Wa-law kāna hādhā dhanbī, ta-Allāhi hādhā laysa min ʿindī bal min ladun
man arsalanī wa-baʿathanī bi-al-ḥaqqi wa-jaʿalanī sirājan li-al-ʿālamīn.
Ey qawm, fa-wa-Allāh rāyeḥe-ye ṣedq az kedhb basī maʿlūm va vāżeḥ ast.
Mapasandīd bar nafs-e Ḥaqq ānche bar khod nemīpasandīd. Ākher qadrī
taʾammol namāyīd, maqṣūd-e īn gholām dar īn ākher-e ayyām dar īn sejn-e
kobrā che būde va che khāhad būd. Fa-taʾammalū yā malaʾa al-ghāfilīn.

And if this is My sin, by God, it is not of My own doing, but comes from
the One who sent Me, raised Me up in truth, and made Me a lamp for the
worlds. O people! By God, the fragrance of truthfulness is readily and
clearly distinguished from falsehood. Do not wish for the True One
Himself what you do not wish for yourselves. Reflect a little, after
all: what has been, and what will be, the purpose of this Youth in these
latter days in this Most Great Prison? Reflect, O assembly of the
heedless!

\subsection{130}\label{section-1685}

\begin{otherlanguage}{persian}

یا قوم فاستحیوا عن اللّه و لا تشکّوا فی امره و لا تسأل الدّلیل عن الشّمس و
اشراقها لأنّ دلیلها هو ضیائها و انوارها و انّک لا تکن مریباً فی ظهورها
لأنّها اشرقت و لا مردّ لها بل فاسأل اللّه بأن یفتح بصرک لتدرک انوارها الّتی
احاطت الکونین قل تاللّه قد رقم قلم القدس من رحیق المسک علی جبینی البیضآء
بخطّ ابهی

\end{otherlanguage}

Yā qawm, fa-staḥyū ʿan Allāhi wa-lā tashukkū fī amrihi wa-lā tasʾal
al-dalīla ʿan al-shamsi wa-ishrāqihā li-anna dalīlahā huwa ḍiyāʾuhā
wa-anwāruhā wa-innaka lā takun murīban fī ẓuhūrihā li-annahā ashraqat
wa-lā maradda lahā bal fa-sʾal Allāha bi-an yaftaḥa baṣaraka li-tudrika
anwārahā allatī aḥāṭat al-kawnayn. Qul ta-Allāhi qad raqama qalamu
al-qudsi min raḥīqi al-miski ʿalā jabīnī al-bayḍāʾi bi-khaṭṭin abhā.

O people! Feel shame before God and do not doubt His Cause. Do not ask
for proof of the sun and its shining, for its proof is its radiance and
its light. Do not doubt its appearance, for it has risen and cannot be
turned back. Rather, ask God to open your sight so that you may perceive
its lights, which have encompassed both worlds. Say: By God, the Pen of
Holiness has inscribed upon My radiant brow, with the pure essence of
musk, in a most glorious script:

\subsection{131}\label{section-1686}

\begin{otherlanguage}{persian}

ان یا ملأ الأرض و السّمآء انّ هذا لهو المحبوب الّذی ما شهدت عین الابداع
مثله و لا عین الاختراع شبهه و انّه لهو الّذی قرّت بجماله عین اللّه الملک
العزیز الجمیل حدیث قبل است که صفحه‌ئی از توریة در دست یکی از اصحاب رسول
اللّه بود حضرت فرمود چه در دست داری عرض نمود توریة فرمود بگذار آن را اگر
میبود صاحب آن اطاعت مینمود مرا حال ملاحظه کن کتاب الهی که امر کلّ از نبیّ
و حروفات و مرایا و ابواب و ارکان باو محقّق و ثابت شده اخذ آن و قرائت آن
ممنوع شود جایز ولکن تبدیل یکی از مرایا علی زعمکم که بوجود حروفاتی محقّق
شده که آن حروفات بکتاب ثابت و محقّق شده‌اند جایز نه فویل لکم تقولون ما لا
تشعرون و تحسبون انّکم مهتدون کما یظنّون ملل القبل و بذلک یفرحون تاللّه لا
عاصم لأحد من امر اللّه الّا بأن ینقطع عن کلّ من فی السّموات و الأرض و یتوجّه
الی اللّه المهیمن القیّوم نظر را مطهّر ساخته تا جمیع اشیاء را الیوم در صقع
واحد مشاهده کنید و انوار شمس تقدیس را از افق قمیص بچشم و بصر ظاهر و باطن
مشاهده نمائید

\end{otherlanguage}

An yā malaʾa al-arḍi wa-al-samāʾ, inna hādhā la-huwa al-maḥbūbu alladhī
mā shahidat ʿaynu al-ibdāʿi mithlahu wa-lā ʿaynu al-ikhtirāʿi shibhahu
wa-innahu la-huwa alladhī qarrat bi-jamālihi ʿaynu Allāhi al-maliki
al-ʿazīzi al-jamīl. Ḥadīth-e qabl ast ke ṣafḥeʾī az Towrāt dar dast-e
yekī az aṣḥāb-e Rasūl Allāh būd. Ḥażrat farmūd che dar dast dārī? ʿArż
namūd Towrāt. Farmūd begoẕār ān rā; agar mībūd ṣāḥeb-e ān eṭāʿat mīnamūd
marā. Ḥāl molāḥeẓe kon, ketāb-e elāhī ke amr-e koll az nabī va ḥorūfāt
va marāyā va abvāb va arkān be-ū moḥaqqaq va thābet shode, akhdh-e ān va
qerāʾat-e ān mamnūʿ shavad jāyez, va-lāken tabdīl-e yekī az marāyā ʿalā
zaʿmekom ke be-vojūd-e ḥorūfātī moḥaqqaq shode ke ān ḥorūfāt be-ketāb
thābet va moḥaqqaq shode-and jāyez na. Fa-waylun lakum taqūlūna mā lā
tashʿurūna wa-taḥsabūna annakum muhtadūna kamā yaẓunnūna milalu al-qabli
wa-bi-dhālika yafraḥūn. Ta-Allāhi lā ʿāṣima li-aḥadin min amri Allāhi
illā bi-an yanqaṭiʿa ʿan kulli man fī al-samāwāti wa-al-arḍi
wa-yatawajjaha ilā Allāhi al-muhaymini al-qayyūm. Naẓar rā moṭahhar
sākhte tā jamīʿ-e ashyā rā al-yowm dar ṣoqʿ-e vāḥed moshāhede konīd va
anvār-e shams-e taqdīs rā az ofoq-e qamīṣ be-chashm va baṣar-e ẓāher va
bāṭen moshāhede namāyīd.

``O assembly of earth and heaven! This is indeed the Beloved, whose like
the eye of creation has never beheld, nor has the eye of origination
seen His equal. He is indeed the One in whose beauty the eye of God, the
King, the Almighty, the Beautiful, has found delight.'' There is an
earlier tradition that one of the companions of the Messenger of God
held a page of the Torah in his hand. He asked, ``What do you have in
your hand?'' The companion replied, ``The Torah.'' He said, ``Put it
down. If its author were here, he would obey Me.'' Now consider: you
deem it permissible to forbid taking up and reading the Book of God,
through which the authority of all---Prophet, Letters, Mirrors, Gates,
and Pillars---has been established and confirmed. Yet you deem it
impermissible to change one of the Mirrors, whose standing, by your own
account, was established through the existence of Letters whose own
standing was confirmed and established through the Book. Woe to you! You
say what you do not understand and suppose that you are rightly guided,
just as the religious communities of the past suppose, and in this they
rejoice. By God, no one has any protection from the command of God
unless they detach themselves from everyone in the heavens and on earth
and turn toward God, the All-Protecting, the Self-Subsisting. Purify
your vision so that today you may behold all things on a single plane
and, with both your outward and inward eyes and sight, behold the lights
of the Sun of holiness shining from the horizon of the robe.

\subsection{132}\label{section-1687}

\begin{otherlanguage}{persian}

و از جمیع این براهین گذشته در میرزا اسداللّه که نقطهٴ اولی جلّ ذکره باو
خطاب فرموده ان یا مظهر الأحدیّة حال انصاف دهید کسی که محلّ ظهور سلطان
احدیّه باشد بقول اللّه چگونه ابوالشّرور مینامید و اگر بگوئید این شأن را
مرآت از او نفی نمود بر فرض جواز آن حال اگر نفس نقطه قمیص قبول را از
مرآتی از مرایا انتزاع فرماید چه اعتراض بر آن جمال قدم وارد و اگر بگوئید
از کجا محقّق شده که این نقطهٴ مشهوده و کلمهٴ جامعه در ادّعای خود صادق باشد
میگوئیم بهمان حجّت و برهان که از ظهور قبلم کلّما ظهر من عنده و نزّل من لدنه
را مسلّم داشته‌اید حال بهمان حجّت و برهان ناظر شوید و بخواهید ان رأیتم لا
تکفروا باللّه الّذی خلقکم بأمره ان انتم من المنصفین و هم‌چنین به اسداللّه
میفرماید ان یا حرف الثّالث المؤمن بمن یظهره اللّه جلّ اجلاله

\end{otherlanguage}

Va az jamīʿ-e īn barāhīn goẕashte, dar Mīrzā Asad Allāh ke Noqṭe-ye Ūlā
jalla dhikruhu be-ū kheṭāb farmūde, an yā maẓhara al-aḥadiyya, ḥāl enṣāf
dehīd: kasī ke maḥall-e ẓohūr-e solṭān-e aḥadiyye bāshad be-qowl-e Allāh
chegūne Abu al-Shorūr mīnāmīd? Va agar begūyīd īn shaʾn rā merʾāt az ū
nafy namūd, bar farż-e javāz-e ān, ḥāl agar nafs-e Noqṭe qamīṣ-e qabūl
rā az merʾātī az marāyā entezāʿ farmāyad che eʿterāż bar ān Jamāl-e
Qedam vāred? Va agar begūyīd az kojā moḥaqqaq shode ke īn Noqṭe-ye
mashhūde va Kaleme-ye Jāmeʿe dar eddeʿā-ye khod ṣādeq bāshad, mīgūyīm
be-hamān ḥojjat va borhān ke az ẓohūr-e qablam kulla mā ẓahara min
ʿindihi wa-nuzzila min ladunhu rā mosallam dāshte-īd, ḥāl be-hamān
ḥojjat va borhān nāẓer shavīd va bekhāhīd. In raʾaytum lā takfurū
bi-Allāhi alladhī khalaqakum bi-amrihi in antum min al-munṣifīn. Va
ham-chenīn be Asad Allāh mīfarmāyad, an yā ḥarfa al-thālithi al-muʾmina
bi-man yuẓhiruhu Allāh, jalla ijlāluhu.

Apart from all these proofs, consider Mírzá Asadu'lláh, whom the Primal
Point---exalted be His mention---addressed as ``O Manifestation of
Oneness!'' Now judge fairly: how can you call someone who, according to
the word of God, is the place where the Sovereign of Oneness is
manifested, ``the Father of Evils''? If you say that the Mirror deprived
him of this station, then, assuming that this is permissible, what
objection can be raised against that Ancient Beauty if the Point Himself
now strips one of the Mirrors of the robe of acceptance? And if you ask
how it has been established that this visible Point and all-embracing
Word is truthful in His claim, We say: through the very same proof and
evidence by which you have accepted everything that appeared from Me and
was sent down by Me in My previous revelation. Now look for and seek
that same proof and evidence. If you see it, do not disbelieve in God,
who created you by His command, if you are fair-minded. He likewise says
to Asadu'lláh: ``O Letter of the Third, who believes in Him whom God
will make manifest---exalted be His majesty!''

\subsection{133}\label{section-1688}

\begin{otherlanguage}{persian}

و هم‌چنین در رسالهٴ جفریّه که مخصوص او از سماء هویّه نازل شده در صدر رساله
مکتوب و هذه صورته ما نزّل لحضرة الأسد و الفرد الأحد المستشرق بالنّور الصّمد
اسم اللّه الدّیّان للّه الواحد المعتمد و بعد میفرماید جلّ ذکره ان یا اسم
الدّیّان هذا علم مکنون مخزون قد اودعناک و اتیناک عزّاً من عند اللّه اذ عین
فؤادک لطیف تعرف قدره و تعزّ بهائه و بعد میفرماید عزّ قدره ان یا اسم الأجلّ
لا تنظر الی ذلک الخلق الّا بمثل اغنام یذهبه راعیه کیف یشآء و من‌دون آنچه
ذکر شده کلمات قدس صمدانیّه که در توصیف او نازل در کتاب موجود فانظروا الیه
و این بنده نظر باختصار ذکر ننموده ای اکمهٴ وجود بصر بگشا

\end{otherlanguage}

Va ham-chenīn dar Resāle-ye Jafriyye ke makhṣūṣ-e ū az samā-ye hoviyye
nāzel shode, dar ṣadr-e resāle maktūb, wa-hādhihi ṣūratuhu: mā nuzzila
li-ḥaḍrati al-asadi wa-al-fardi al-aḥadi al-mustashriqi bi-al-nūri
al-ṣamadi ismi Allāhi al-Dayyāni li-llāhi al-wāḥidi al-muʿtamad. Va baʿd
mīfarmāyad jalla dhikruhu, an yā isma al-Dayyān, hādhā ʿilmun maknūnun
makhzūnun qad awdaʿnāka wa-ātaynāka ʿizzan min ʿindi Allāhi idh ʿaynu
fuʾādika laṭīfun taʿrifu qadrahu wa-tuʿizzu bahāʾahu. Va baʿd mīfarmāyad
ʿazza qadruhu, an yā isma al-ajall, lā tanẓur ilā dhālika al-khalqi illā
bi-mithli aghnāmin yudhhibuhu rāʿīhi kayfa yashāʾ. Va men-dūn-e ānche
dhekr shode kalemāt-e qods-e ṣamadāniyye ke dar towṣīf-e ū nāzel dar
ketāb mowjūd, fa-nẓurū ilayhi. Va īn bande naẓar be-ekhteṣār dhekr
nanamūde. Ey akmah-e vojūd, baṣar bogshā.

Likewise, in the Treatise of Jafr, which was sent down especially for
him from the heaven of the divine Essence, the following is written at
the beginning of the treatise: ``What has been sent down for His Honor
the Lion, the Single, the One, illumined by the everlasting Light, the
Name of God, the Judge, for God, the One, the Relied-Upon.'' He then
says---exalted be His mention: ``O Name of the Judge! This is a hidden,
treasured knowledge that We have entrusted to you, and We have bestowed
upon you an honor from God, for the eye of your heart is discerning: you
know its worth and cherish its splendor.'' He then says---exalted be His
station: ``O most exalted Name! Regard these people only as sheep that
their shepherd drives wherever he wills.'' Besides what has been
mentioned, words of everlasting holiness revealed in his praise are
found in the Book; consult it. This servant has not quoted them for the
sake of brevity. O you who are blind from birth in the realm of
existence, open your eyes!

\subsection{134}\label{section-1689}

\begin{otherlanguage}{arabic}

و ملاحظه بوهم خود نما که همان فعلی را که خود مرتکبی و بآن عامل حال از آن
متعجّب و متحیّر مانده چنانچه دیّان را دنی میگوئی و فرد الأحد که از اعظم
اسماء الهی بوده ابوالشّرور نامیده و مع‌ذلک اعتراض بر سلطان یفعل ما یشآء
نموده که چگونه میشود حلو مرّ شود و یا مرآت حجر گردد با اینکه اسداللّه
انکار حقّ ننمود و حین ورود مدینة‌اللّه بین یدی العرش حاضر و بآنچه اهل بیان
در آن یوم بآن مقرّ بوده‌اند اقرار نموده و مع‌ذلک ورد علیه ما ورد بلکه اکتفا
باو وحده ننموده بر جمعی حکم قتل جاری نموده کتاب مستیقظ موجود ملاحظه
نمائید

\end{otherlanguage}

Va molāḥeẓe be-vahm-e khod namā ke hamān feʿlī rā ke khod mortakebī va
be-ān ʿāmel ḥāl az ān motaʿajjeb va motaḥayyer mānde, chenānche Dayyān
rā danī mīgūʾī va Fard al-Aḥad ke az aʿẓam-e asmā-ye elāhī būde Abu
al-Shorūr nāmīde va maʿa-dhālek eʿterāż bar Solṭān-e yafʿalu mā yashāʾ
namūde ke chegūne mīshavad ḥolv morr shavad va yā merʾāt ḥajar gardad.
Bā īnke Asad Allāh enkār-e Ḥaqq nanamūd va ḥīn-e vorūd-e Madīnat Allāh
bayn-e yaday-e ʿarsh ḥāżer va be-ānche ahl-e Bayān dar ān yowm be-ān
moqerr būde-and eqrār namūde va maʿa-dhālek warada ʿalayhi mā warad.
Balke ektefā be-ū vaḥdeh nanamūde bar jamʿī ḥokm-e qatl jārī namūde.
Ketāb-e Mostayqeẓ mowjūd, molāḥeẓe namāyīd.

Consider your own delusion: you now stand astonished and bewildered at
the very act that you yourself commit and practice. You call Dayyán
``base'' and have named the Single, the One---one of the greatest names
of God---``the Father of Evils.'' Yet you object to the Sovereign who
does whatever He wills, asking how the sweet can become bitter or a
mirror turn to stone. Asadu'lláh, however, did not deny the True One. On
arriving in the City of God, he entered the presence of the Throne and
acknowledged what the people of the Bayán acknowledged on that day.
Nevertheless, there befell him what befell him. Nor did you content
yourself with him alone: you pronounced a death sentence against a
number of people. The book Mustayqiẓ is available; consult it.

\subsection{135}\label{section-1690}

\begin{otherlanguage}{persian}

قسم بآفتاب معانی که در آن کتاب بر کلّ مظاهر احدیّه فتوای قتل داده شده با
آنکه حکم قتل بالمرّه از بیان و اهل آن مرتفع شده بنصّ منزل بیان یک دو افترا
بآن بیچاره زده‌اند و حکم قتل جاری نموده‌اند من غیر اذن اللّه و ربّ مقتول شعر
منه خیر من الف قاتلیه چنانچه حال بر جمال احدیّه ظلم باین واضحی و مبرهنی
وارد و بقتلش قیام نموده‌اند مع‌ذلک باطراف نوشته و چه ظلمها که بجوهر عدل و
منبع و معدن آن نسبت داده که شاید غلّ و بغضا در صدور بعضی از متوهّمین القا
نماید ولکن غافل از آنکه و للّه تحت ظلال اسمه البصیر عباد یشقّنّ الشّعر
بالشّعر و لن یمنع ابصارهم حجبات العالمین

\end{otherlanguage}

Qasam be-āftāb-e maʿānī ke dar ān ketāb bar koll-e maẓāher-e aḥadiyye
fatvā-ye qatl dāde shode, bā ānke ḥokm-e qatl be-al-marre az Bayān va
ahl-e ān mortafeʿ shode be-naṣṣ-e Monzel-e Bayān. Yek do efterā be-ān
bīchāre zade-and va ḥokm-e qatl jārī namūde-and min ghayri idhni Allāh.
Wa-rubba maqtūlin shaʿrun minhu khayrun min alfi qātilīhi. Chenānche ḥāl
bar Jamāl-e Aḥadiyye ẓolm be-īn vāżeḥī va mobarhanī vāred va be-qatlash
qiyām namūde-and, maʿa-dhālek be-aṭrāf neveshte va che ẓolmhā ke
be-Jowhar-e ʿAdl va manbaʿ va maʿdan-e ān nesbat dāde ke shāyad ghell va
baghżā dar ṣodūr-e baʿżī az motavahhemīn elqā namāyad, va-lāken ghāfel
az ānke wa-li-llāhi taḥta ẓilāli ismihi al-baṣīri ʿibādun yashaqqanna
al-shaʿra bi-al-shaʿri wa-lan yamnaʿa abṣārahum ḥujubātu al-ʿālamīn.

I swear by the Sun of inner meanings that in that book a sentence of
death has been issued against all the manifestations of Oneness, even
though the death penalty has been entirely removed from the Bayán and
its people by the explicit text of the Revealer of the Bayán. They have
leveled one or two false accusations against that hapless man and
pronounced a death sentence without God's permission. How many a slain
person has a single hair worth more than a thousand of his killers!
Thus, even now, they have inflicted such evident and manifest injustice
upon the Beauty of Oneness and have risen to kill Him. Yet they have
written far and wide, attributing so many injustices to the very Essence
of justice, its Source and Wellspring, hoping perhaps to instill malice
and hatred in the hearts of some who are prey to delusions. But they are
unaware that God has servants beneath the shadow of His name, the
All-Seeing, who split a hair with a hair, and whose sight the veils of
all the worlds will never obstruct.

\subsection{136}\label{section-1691}

\begin{otherlanguage}{persian}

و هم‌چنین به سیّد ابرهیم میفرمایند عزّ ذکره ان یا خلیلی فی الصّحف ان یا ذکری
فی الکتب من بعد الصّحف ان یا اسمی فی البیان الی آخر ما نزّل این بیان الهی
که در ذکر او از قلم قدس صمدانی جاری شده و حال به ابوالدّواهی مشهور شده با
اینکه خود افتخار مینمایند که بذکر اللّه فی العالمین در کتب الهی ذکر
شده‌اند و نفسی که بخلیل در صحف الهیّه و بذکر اللّه در کتب منزله و باسم اللّه
در بیان متّسم شده به ابوالدّواهی تعبیر مینمایند و در دونش که اعتراض از نفس
مشیّة اللّه و عظمته نموده تغییر جایز نمیدانند قدرت الهی را مثل قدرت خود
فرض گرفته‌اند و قلم اعلی را معطّل دانسته‌اند

\end{otherlanguage}

Va ham-chenīn be Seyyed Ebrāhīm mīfarmāyand ʿazza dhikruhu, an yā
khalīlī fī al-ṣuḥufi an yā dhikrī fī al-kutubi min baʿdi al-ṣuḥufi an yā
ismī fī al-Bayān, ilā ākhiri mā nuzzil. Īn bayān-e elāhī ke dar dhekr-e
ū az qalam-e qods-e ṣamadānī jārī shode va ḥāl be Abu al-Davāhī mashhūr
shode, bā īnke khod eftekhār mīnamāyand ke be-Dhekr Allāh fī al-ʿālamīn
dar kotob-e elāhī dhekr shode-and. Va nafsī ke be-Khalīl dar ṣoḥof-e
elāhiyye va be-Dhekr Allāh dar kotob-e monzale va be-Esm Allāh dar Bayān
mottasem shode be-Abu al-Davāhī taʿbīr mīnamāyand, va dar dūnash ke
eʿterāż az nafs-e Mashiyyat Allāh va ʿaẓamat-e ū namūde taghyīr jāyez
nemīdānand. Qodrat-e elāhī rā methl-e qodrat-e khod farż gerefte-and va
Qalam-e Aʿlā rā moʿaṭṭal dāneste-and.

Likewise, He says to Siyyid Ibráhím---exalted be His mention: ``O My
Friend in the Scriptures! O My Remembrance in the Books after the
Scriptures! O My Name in the Bayán!''---and so on to the end of what was
revealed. Such is the divine utterance that flowed from the Pen of
everlasting holiness in his mention; yet he has now become known as
``the Father of Calamities,'' even though they themselves take pride in
having been mentioned in the divine Books as ``the Remembrance of God
among the worlds.'' They call someone designated as the Friend in the
divine Scriptures, the Remembrance of God in the revealed Books, and the
Name of God in the Bayán ``the Father of Calamities.'' Yet in the case
of one beneath him, who has objected to the very Will of God and His
greatness, they consider change impermissible. They have supposed divine
power to be like their own power and have regarded the Most Exalted Pen
as idle.

\subsection{137}\label{section-1692}

\begin{otherlanguage}{persian}

قسم بمظلومیّتم که قلم قدرت کلّ حین در محو و اثبات مشغول و لا یعقل ذلک الّا
المنقطعون کذلک یمحو الرّحمن ما یشآء و یثبت و عنده امّ البیان و عجب از این
نفوس که باین اسماء افتخار مینمایند و از حقّ معرض گشته‌اند بلی این اذکار و
اوصاف در هر نفسی موجود مادامی که از حصن امر اللّه خارج نشده‌اند والّا لیس
لهم قدر عند اللّه علی قدر سواد عین نملة میّتة و از این گذشته کلّ اشیاء ذکر
اللّه فی العالمین بوده و خواهند بود چه که کلّ اشیاء بکینونتها اسم اللّه
بوده و اسم او تعالی شأنه ذکر او بوده بین عباد و لم‌یزل چنین بوده و لایزال
چنین خواهد بود

\end{otherlanguage}

Qasam be-maẓlūmiyyatam ke qalam-e qodrat koll-e ḥīn dar maḥv va ethbāt
mashghūl, wa-lā yaʿqilu dhālika illā al-munqaṭiʿūn. Kadhālika yamḥū
al-Raḥmānu mā yashāʾu wa-yuthbitu wa-ʿindahu Ummu al-Bayān. Va ʿajab az
īn nofūs ke be-īn asmā eftekhār mīnamāyand va az Ḥaqq moʿreż gashte-and.
Balī īn adhkār va owṣāf dar har nafsī mowjūd mā-dāmī ke az ḥeṣn-e Amr
Allāh khārej nashode-and, va-illā laysa lahum qadrun ʿinda Allāhi ʿalā
qadri sawādi ʿayni namlatin mayyita. Va az īn goẕashte koll-e ashyā
Dhekr Allāh fī al-ʿālamīn būde va khāhand būd, che ke koll-e ashyā
bi-kaynūnatihā Esm Allāh būde va esm-e ū taʿālā shaʾnuhu dhekr-e ū būde
bayn-e ʿebād, va lam-yazal chenīn būde va lā-yazāl chenīn khāhad būd.

I swear by the wrongs I have suffered that the Pen of power is at every
moment engaged in effacing and confirming, and none understands this
except those who are detached. Thus the All-Merciful effaces whatever He
wills and confirms, and with Him is the Mother of the Bayán. How
astonishing are these people who take pride in these names yet have
turned away from the True One! Yes, these designations and attributes
remain in each soul so long as they have not left the stronghold of the
Cause of God. Otherwise, they have no worth before God, not even as much
as the black of a dead ant's eye. Beyond this, all things have been and
will be the Remembrance of God among the worlds, for all things, by
their very being, are the Name of God, and His name---exalted be His
glory---is His remembrance among His servants. So it has ever been, and
so it will ever remain.

\subsection{138}\label{section-1693}

\begin{otherlanguage}{persian}

فوالّذی تفرّد بالعزّ و البهآء که اگر نمیبود که نظر اللّه خلق را نابالغ
مشاهده نموده هرآینه در هر ذرّه از ذرّات وجود ابواب علمی میگشودم که جمیع
خلق خود را از عرفانش عاجز و قاصر مشاهده نمایند ولکن چون اغیار موجود
اسرار مستور به کذلک قدّر من لدن مقتدر قدیر قل یا ملأ البیان انتم خلقتم
للقائی بعد استغنائی عنکم و وعدتم بنفسی و عرفانها بعد تقدیسی عن عرفانکم و
بشّرتم بجمالی بعد استعلائی عنکم و عمّا خلق بین السّموات و الأرضین مع‌ذلک
باوهن اسماء الیوم تشبّث نموده از سلطان اسماء دور و غافل مانده‌اید اگر
بکلمات ناظرید معادل بیان از سماء رحمن نازل و بشأنی ظاهر که احدی را مجال
اعراض نه مگر مشرکین که مرض قلبی آن نفوس را از عرفان و ادراک منع نموده

\end{otherlanguage}

Fa-wa-lladhī tafarrada bi-l-ʿizzi wa-l-bahāʾ ke agar namībūd ke naẓar-i
Allāh khalq rā nābāligh mushāhade namūde har-āyine dar har dharre az
dharrāt-i vujūd abvāb-i ʿilmī mīgushūdam ke jamīʿ-i khalq khud rā az
ʿirfānash ʿājiz va qāṣir mushāhade namāyand va-līkin chūn aghyār mawjūd
asrār mastūr beh kadhālika quddira min ladun muqtadirin qadīr qul yā
malaʾa l-bayān antum khuliqtum li-liqāʾī baʿda istighnāʾī ʿankum
wa-wuʿidtum bi-nafsī wa-ʿirfānihā baʿda taqdīsī ʿan ʿirfānikum
wa-bushshirtum bi-jamālī baʿda istiʿlāʾī ʿankum wa-ʿammā khuliqa bayna
l-samāwāti wa-l-araḍīn maʿa-dhālik bi-awhan-i asmāʾ al-yawm tashabbuth
namūde az sulṭān-i asmāʾ dūr va ghāfil mānde-īd agar bi-kalimāt nāẓirīd
muʿādil-i bayān az samāʾ-i raḥmān nāzil va bi-shaʾnī ẓāhir ke aḥadī rā
majāl-i iʿrāḍ na magar mushrikīn ke maraḍ-i qalbī ān nufūs rā az ʿirfān
va idrāk manʿ namūde.

I swear by Him who is peerless in might and glory: were it not that the
gaze of God has found humanity immature, I would open within every atom
of existence doors to knowledge before which all people would find
themselves powerless and incapable of comprehension. But since strangers
are present, it is better that the mysteries remain concealed. Thus has
it been decreed by One who is all-powerful and almighty. Say: ``O people
of the Bayán! You were created to meet Me, though I have no need of you;
you were promised My own Self and its recognition, though I am
sanctified above your recognition; and you were given the glad tidings
of My beauty, though I am exalted above you and whatever has been
created between the heavens and the earths.'' Yet today you have clung
to the feeblest of names and remained far from, and heedless of, the
Sovereign of names. If you look to the words, the equivalent of the
Bayán has descended from the heaven of the All-Merciful and has appeared
in such a manner that no one has any grounds for turning away, except
those who associate partners with God, whose sickness of heart has
prevented them from recognition and understanding.

\subsection{139}\label{section-1694}

\begin{otherlanguage}{persian}

از این گذشته کلمات نفوسی که الیوم حول امر اللّه طایفند با کلمات من
اتّخذتموه من دون اللّه لأنفسکم ولیّاً میزان نمائید تا بر عظمت امر مطّلع شوید
که لعلّ بر جمال قدم بمجرّد وهم ضرّ وارد نیاورید قسم بحقّ که اهل هر ملّتی از
ملل قبل اگر در احتجاب خود از مظاهر عزّ باقیه بمعاذیر نالایقه متعذّر شوند
شاید ولکن از برای اهل بیان هیچ حجّت و عذری باقی نه چه که سلطان احدیّه ذکر
ظهور بعد را باصرح بیان و اوضح تبیان فرموده‌اند بشأنی که کلّ بیان و ما نزّل
فیه را بقول او منوط فرموده چنانچه از قبل از قلم قدم تحریر یافت فطوبی
للعارفین

\end{otherlanguage}

Az īn gudashte kalimāt-i nufūsī ke al-yawm ḥawl-i amr-i Allāh ṭāʾifand
bā kalimāt-i man ittakhadhtumūhu min dūni llāhi li-anfusikum waliyyan
mīzān namāyīd tā bar ʿaẓamat-i amr muṭṭaliʿ shavīd ke laʿall bar jamāl-i
qidam bi-mujarrad-i vahm ḍurr vārid nayāvarīd qasam bi-ḥaqq ke ahl-i har
millatī az milal-i qabl agar dar iḥtijāb-i khud az maẓāhir-i ʿizz-i
bāqiye bi-maʿādhīr-i nālāyiqe mutaʿadhdhir shavand shāyad va-līkin az
barā-yi ahl-i bayān hīch ḥujjat va ʿudhrī bāqī na che ke sulṭān-i
aḥadiyye dhikr-i ẓuhūr-i baʿd rā bi-aṣraḥ-i bayān va awḍaḥ-i tibyān
farmūde-and bi-shaʾnī ke kull-i bayān va mā nuzzila fīhi rā bi-qawl-i ū
manūṭ farmūde chunānche az qabl az qalam-i qidam taḥrīr yāft fa-ṭūbā
li-l-ʿārifīn.

Beyond this, weigh the words of those who today circle around the Cause
of God against the words of the one whom you have taken as your guardian
instead of God, so that you may become aware of the greatness of the
Cause and perhaps refrain from harming the Ancient Beauty on the basis
of mere conjecture. I swear by the True One: if the followers of any
former religion should offer unworthy excuses for remaining veiled from
the Manifestations of everlasting glory, that might be possible; but for
the people of the Bayán no argument or excuse remains. For the Sovereign
of oneness has spoken of the next Revelation in the most explicit
language and the clearest exposition, making the entire Bayán and all
that was revealed in it dependent on His word, as was previously written
by the Pen of eternity. Blessed, then, are those who recognize.

\subsection{140}\label{section-1695}

\begin{otherlanguage}{persian}

و از جمله در ابطال حقّ و اثبات باطل معرض باللّه باین کلمات تمسّک جسته و چون
لایق نبود که کلماتش بعینه ذکر شود لذا ما حرّک علیها قلم اللّه و مضمون آن
این است که از قلم عزّ صادر میشود محبوب شما صاحب بصر بوده یا نه اگر نفی
شود حکم بصر از منظر اکبر هذا ظلم علی اللّه و اگر آن سلطان قدم صاحب بصر و
علم بوده چگونه میشود نفسی را وصف فرماید و باسماء حسنی ذکر نماید از مقام
خود تنزّل کند و از معرضین محسوب گردد

\end{otherlanguage}

Va az jumle dar ibṭāl-i ḥaqq va ithbāt-i bāṭil muʿriḍ bi-llāh bi-īn
kalimāt tamassuk juste va chūn lāyiq nabūd ke kalimātash bi-ʿaynih dhikr
shavad lidhā mā ḥurrika ʿalayhā qalamu llāh va maḍmūn-i ān īn ast ke az
qalam-i ʿizz ṣādir mīshavad maḥbūb-i shumā ṣāḥib-i baṣar būde yā na agar
nafy shavad ḥukm-i baṣar az manẓar-i akbar hādhā ẓulmun ʿalā llāh va
agar ān sulṭān-i qidam ṣāḥib-i baṣar va ʿilm būde chigūne mīshavad nafsī
rā vaṣf farmāyad va bi-asmāʾ-i ḥusnā dhikr namāyad az maqām-i khud
tanazzul kunad va az muʿriḍīn maḥsūb gardad.

Among other things, in seeking to invalidate the truth and establish
falsehood, the one who has turned away from God has resorted to these
words. Since it was not fitting to quote his words exactly, the Pen of
God was not moved to record them. Their substance, as it now issues from
the Pen of glory, is this: ``Did your Beloved possess sight or not? If
sight is denied to the Most Great Countenance, this is an injustice
against God. But if that Sovereign of eternity possessed sight and
knowledge, how can someone whom He praises and mentions by the most
beautiful names fall from his station and be counted among those who
turn away?''

\subsection{141}\label{section-1696}

\begin{otherlanguage}{persian}

اوّلاً اینکه قائل این قول حقّ جلّ ذکره را یفعل ما یشآء ندانسته و قدرت محیطهٴ
الهیّه را انکار نموده و چنین نفسی ابعد عباد بوده و خواهد بود بنصّ نقطهٴ
بیان چنانچه از قبل مذکور شد قل اللّهمّ انّک انت الهان الالهین لتؤتینّ
الألوهیّة من تشآء و لتنزعنّ الألوهیّة عمّن تشآء الی انّک علی کلّ شیء قدیرا قل
اللّهمّ انّک انت ربّاب السّموات و الأرض لتؤتینّ الرّبوبیّة من تشآء و لتنزعنّ
الرّبوبیّة عمّن تشآء الی آخر ما نزّل

\end{otherlanguage}

Avvalan īnke qāʾil-i īn qawl Ḥaqq jalla dhikruhu rā yafʿalu mā yashāʾ
nadāniste va qudrat-i muḥīṭe-yi ilāhiyye rā inkār namūde va chunīn nafsī
abʿad-i ʿibād būde va khwāhad būd bi-naṣṣ-i nuqṭe-yi bayān chunānche az
qabl madhkūr shud qul allāhumma innaka anta ilāhānu l-ilāhīn
la-tuʾtiyanna l-ulūhiyyata man tashāʾ wa-la-tanziʿanna l-ulūhiyyata
ʿamman tashāʾ ilā innaka ʿalā kulli shayʾin qadīran qul allāhumma innaka
anta rabbābu l-samāwāti wa-l-arḍ la-tuʾtiyanna l-rubūbiyyata man tashāʾ
wa-la-tanziʿanna l-rubūbiyyata ʿamman tashāʾ ilā ākhiri mā nuzzil.

First, the speaker of these words has failed to recognize that the True
One---exalted be His mention---does whatever He wills, and has denied
the all-encompassing power of God. Such a person has been and will
remain the most remote of servants, according to the explicit text of
the Point of the Bayán, as previously cited: ``Say: O God, You are
indeed the God of gods. You will surely bestow divinity on whomever You
will, and You will surely strip divinity from whomever You will,''
continuing to ``You are powerful over all things.'' ``Say: O God, You
are indeed the supreme Lord of the heavens and the earth. You will
surely bestow lordship on whomever You will, and You will surely strip
lordship from whomever You will,'' and so on to the end of what was
revealed.

\subsection{142}\label{section-1697}

\begin{otherlanguage}{persian}

حال میگوئیم در این کلمهٴ جامعه که از ساذج فطرت سلطان احدیّه جاری شده چه
میگوئی اگر حقّ جلّ و عزّ را صادق دانسته سلطان مقتدری که قادر است باینکه
الوهیّت و ربوبیّت را که اعظم مقامات بوده عطا فرماید بهر نفسی که اراده
فرماید و هم‌چنین اخذ نماید از هر نفسی که بخواهد البتّه قادر بوده بر آنکه
خلعت وصف را از نفسی انتزاع فرماید این ذکری است که کلّیّه فرموده‌اند تخصیص
باحدی نداشته و ندارد

\end{otherlanguage}

Ḥāl mīgūʾīm dar īn kalime-yi jāmiʿe ke az sādhij-i fiṭrat-i sulṭān-i
aḥadiyye jārī shude che mīgūʾī agar Ḥaqq jalla va ʿazza rā ṣādiq dāniste
sulṭān-i muqtadirī ke qādir ast bi-īnke ulūhiyyat va rubūbiyyat rā ke
aʿẓam-i maqāmāt būde ʿaṭā farmāyad bi-har nafsī ke irāde farmāyad va
ham-chunīn akhdh namāyad az har nafsī ke bikhwāhad albatte qādir būde
bar ānke khilʿat-i vaṣf rā az nafsī intizāʿ farmāyad īn dhikrī ast ke
kulliyye farmūde-and takhṣīṣ bi-aḥadī nadāshte va nadārad.

Now We ask: what do you say about this all-embracing utterance that has
flowed from the pure, innate nature of the Sovereign of oneness? If you
acknowledge the truthfulness of the True One---glorified and exalted be
He---then an all-powerful Sovereign who can bestow divinity and
lordship, the greatest of stations, on any soul He wills, and likewise
take them from any soul He wishes, is certainly able to strip a soul of
the robe of praise. This is a statement He has made in universal terms;
it has never been restricted to any particular person.

\subsection{143}\label{section-1698}

\begin{otherlanguage}{persian}

و اگر بتخصیص وصف ناظری و ید اللّه را مغلول دانسته و قدرت محیطه را انکار
مینمائی پس چه میگوئی در میرزا اسداللّه که بنصّ صریح در بارهٴ او میفرمایند
ان یا حرف الثّالث المؤمن بمن یظهره اللّه و این بسی واضح است که حرف ثالث من
یظهره اللّه جلّ اجلاله بصدهزار رتبه از حروفات بیان و مرایای او عنداللّه
مقدّم بوده چه که خود نقطهٴ بیان روح ما سواه فداه میفرمایند نطفهٴ یکسالهٴ
زمان من یظهره اللّه اقوی است از کلّ من فی البیان

\end{otherlanguage}

Va agar bi-takhṣīṣ-i vaṣf nāẓirī va yadu llāh rā maghlūl dāniste va
qudrat-i muḥīṭe rā inkār mīnamāʾī pas che mīgūʾī dar Mīrzā Asaduʾllāh ke
bi-naṣṣ-i ṣarīḥ dar bāre-yi ū mīfarmāyand in yā ḥarfa l-thālithi
l-muʾminu bi-man yuẓhiruhu llāh va īn basī vāḍiḥ ast ke ḥarf-i thālith-i
man yuẓhiruhu llāh jalla ijlāluhu bi-ṣad-hazār rutbe az ḥurūfāt-i bayān
va marāyā-yi ū ʿinda llāh muqaddam būde che ke khud-i nuqṭe-yi bayān
rūḥ-i mā sivāhu fidāh mīfarmāyand nuṭfe-yi yaksāle-yi zamān-i man
yuẓhiruhu llāh aqvā ast az kull-i man fī l-bayān.

But if you focus on the specific bestowal of praise, regard the hand of
God as fettered, and deny His all-encompassing power, what do you say
about Mírzá Asadu'lláh, concerning whom He explicitly says: ``O Letter
of the Third, who believes in Him whom God will make manifest!'' It is
abundantly clear that the Third Letter of Him whom God will make
manifest---exalted be His majesty---outranks the Letters of the Bayán
and its Mirrors by a hundred thousand degrees in the sight of God. For
the Point of the Bayán Himself---may the souls of all besides Him be
sacrificed for Him---says: ``A one-year-old embryo of the time of Him
whom God will make manifest is mightier than all who are in the Bayán.''

\subsection{144}\label{section-1699}

\begin{otherlanguage}{persian}

حال انصاف ده نقطهٴ بیان صاحب بصیرت بوده یا نه اگر او را صادق و صاحب بصر
میدانی چرا نفسی را که باین علوّ و امتناع و سموّ و ارتفاع در اصل کتاب ذکر
فرموده‌اند ابوالشّرور نامیدی و بر قتلش امر نمودی ان یا ذی عین واحدة شهدت
عیوب النّاس و غفلت عمّا فی نفسک

\end{otherlanguage}

Ḥāl inṣāf dih nuqṭe-yi bayān ṣāḥib-i baṣīrat būde yā na agar ū rā ṣādiq
va ṣāḥib-i baṣar mīdānī chirā nafsī rā ke bi-īn ʿuluvv va imtināʿ va
sumuvv va irtifāʿ dar aṣl-i kitāb dhikr farmūde-and Abuʾl-Shurūr nāmīdī
va bar qatlash amr namūdī in yā dhī ʿaynin wāḥidatin shahidta ʿuyūba
l-nāsi wa-ghafalta ʿammā fī nafsik.

Now judge fairly: did the Point of the Bayán possess insight or not? If
you consider Him truthful and possessed of sight, why did you call a
person whom He mentioned in the very Book with such loftiness,
inaccessibility, sublimity, and exaltation ``the Father of Evils,'' and
order his killing? O you with but one eye! You have seen the faults of
others and remained heedless of what is within yourself.

\subsection{145}\label{section-1700}

\begin{otherlanguage}{persian}

اگر بگوئی از حقّ اعراض نمود لذا این حکم بر او جاری شد بر فرض تسلیم
میگوئیم حرف ثالث من یظهره اللّه که ابهی و اعلای از خلق بیان عنداللّه بوده
بنصّ قلم امر اگر جایز است معرض و کافر و مشرک شود پس چرا جایز نمیدانی که
مرآتی از مرایای لایحصی از حقّ اعراض نماید فبهت الّذی کفر

\end{otherlanguage}

Agar bigūʾī az Ḥaqq iʿrāḍ namūd lidhā īn ḥukm bar ū jārī shud bar farḍ-i
taslīm mīgūʾīm ḥarf-i thālith-i man yuẓhiruhu llāh ke abhā va aʿlā-yi az
khalq-i bayān ʿinda llāh būde bi-naṣṣ-i qalam-i amr agar jāʾiz ast
muʿriḍ va kāfir va mushrik shavad pas chirā jāʾiz namīdānī ke mirʾātī az
marāyā-yi lā yuḥṣā az Ḥaqq iʿrāḍ namāyad fa-buhita lladhī kafar.

If you say, ``He turned away from the True One, and therefore this
judgment was passed upon him,'' We reply, granting this for the sake of
argument: if the Third Letter of Him whom God will make manifest, who,
according to the explicit text of the Pen of the Cause, was more
glorious and exalted in the sight of God than the people of the Bayán,
can become one who turns away, an unbeliever, and one who associates
partners with God, why do you not admit that one Mirror among countless
Mirrors can turn away from the True One? Thus was the one who
disbelieved confounded.

\subsection{146}\label{section-1701}

\begin{otherlanguage}{arabic}

فواللّه احدی بر اصل این امر مطّلع نه و لو یطّلعون لیعرفون

\end{otherlanguage}

Fa-wa-llāh aḥadī bar aṣl-i īn amr muṭṭaliʿ na wa-law yaṭṭaliʿūna
la-yaʿrifūn.

By God, no one is aware of the true nature of this Cause; if they became
aware, they would surely recognize.

\subsection{147}\label{section-1702}

\begin{otherlanguage}{persian}

و در توقیع وحید اکبر ورقاء بقا بابدع نغمات تغنّی فرموده فیا روحا لمن یسمع
و یکون فی آیات اللّه لمن المتفکّرین بیانه عزّ بیانه فوالّذی فلق الحبّة و برأ
النّسمة لو ایقنت بأنّک یوم ظهوره لا تؤمن به لأرفعت عنک حکم الایمان فی ذلک
الظّهور لأنّک ما خلقت الّا له و لو علمت انّ احداً من النّصاری یؤمن به لجعلته
قرّة عینای و احکمت علیه بالایمان فی ذلک الظّهور من دون ان اشهد علیه من شیء
انتهی

\end{otherlanguage}

Va dar tawqīʿ-i Vaḥīd-i Akbar varqā-yi baqā bi-abdaʿ-i naghamāt taghannī
farmūde fa-yā rawḥan li-man yasmaʿu wa-yakūnu fī āyāti llāhi la-mina
l-mutafakkirīn bayānuhu ʿazza bayānuhu fa-wa-lladhī falaqa l-ḥabbata
wa-baraʾa l-nasamata law ayqantu bi-annaka yawma ẓuhūrihi lā tuʾminu
bihi la-arfaʿtu ʿanka ḥukma l-īmāni fī dhālika l-ẓuhūr li-annaka mā
khuliqta illā lahu wa-law ʿalimtu anna aḥadan mina l-naṣārā yuʾminu bihi
la-jaʿaltuhu qurrata ʿaynāya wa-aḥkamtu ʿalayhi bi-l-īmāni fī dhālika
l-ẓuhūr min dūni an ashhada ʿalayhi min shayʾ intahā.

In the Tablet to Vaḥíd-i Akbar, the Dove of eternity has sung in the
most wondrous melodies. What joy for one who listens and is among those
who ponder the verses of God! His words---exalted be His
utterance---are: ``I swear by Him who split the seed and created the
soul: if I were certain that on the day of His appearance you would not
believe in Him, I would remove from you the designation of faith in that
Revelation, for you were created for Him alone. And if I knew that one
of the Christians would believe in Him, I would make him the solace of
My eyes and pronounce him a believer in that Revelation, without
witnessing anything from him.'' Here ends the quotation.

\subsection{148}\label{section-1703}

\begin{otherlanguage}{persian}

میفرماید قسم بمقتدری که شکافت حبّهٴ قلوب را و خلق فرمود انسان را که اگر
موقن بودم که تو در ظهور بعد ایمان نمیآوری بآن مشرق آفتاب حقیقت هرآینه
مرتفع مینمودم از تو حکم ایمان را در این ظهور چه که تو خلق نشده مگر از
برای عرفان آن شمس سماء ایقان و هرگاه بدانم یکی از نصاری موفّق میشود
بعرفان آن جمال رحمن و ایمان باو هرآینه میگردانیدم او را قرّۀ چشمهای خود و
حکم مینمودم بر او در این ظهور بایمان از دون اینکه ملحوظ شود از او امری

\end{otherlanguage}

Mīfarmāyad qasam bi-muqtadirī ke shikāft ḥabbe-yi qulūb rā va khalq
farmūd insān rā ke agar mūqin būdam ke tu dar ẓuhūr-i baʿd īmān
namīāvarī bi-ān mashriq-i āftāb-i ḥaqīqat har-āyine murtafiʿ mīnamūdam
az tu ḥukm-i īmān rā dar īn ẓuhūr che ke tu khalq nashude magar az
barā-yi ʿirfān-i ān shams-i samāʾ-i īqān va har-gāh bidānam yakī az
naṣārā muwaffaq mīshavad bi-ʿirfān-i ān jamāl-i raḥmān va īmān bi-ū
har-āyine mīgardānīdam ū rā qurrat-i chashmhā-yi khud va ḥukm mīnamūdam
bar ū dar īn ẓuhūr bi-īmān az dūn-i īnke malḥūẓ shavad az ū amrī.

He says: ``I swear by the All-Powerful One who split the seed of hearts
and created humanity: if I were certain that in the next Revelation you
would not believe in that Dawning-Place of the Sun of Truth, I would
remove from you the designation of faith in this Revelation, for you
were created solely to recognize that Sun of the heaven of certitude.
And if I should know that one of the Christians would attain recognition
of that Beauty of the All-Merciful and faith in Him, I would assuredly
make him the solace of My eyes and pronounce him a believer in this
Revelation, without any deed of his being taken into account.''

\subsection{149}\label{section-1704}

\begin{otherlanguage}{persian}

حال ای اهل بصر ملاحظه نمائید که امر اللّه چه مقدار الطف و ارقّ بوده و رحمت
منبسطهٴ الهیّه چگونه کلّ وجود را احاطه فرموده که در بارهٴ نفسی از نصاری که
الیوم کافر و مشرک است بمجرّد ایمان در ظهور بعد آن سدرهٴ الهی او را بالطف
و ارقّ اعضای مبارک خود نسبت داده و منسوب ساخته‌اند و مع‌ذلک این همج ارض در
چه رتبه واقف و ناظرند فسحقاً لکم یا ملأ الغافلین

\end{otherlanguage}

Ḥāl ay ahl-i baṣar mulāḥiẓe namāyīd ke amr-i Allāh che miqdār alṭaf va
araqq būde va raḥmat-i munbasiṭe-yi ilāhiyye chigūne kull-i vujūd rā
iḥāṭe farmūde ke dar bāre-yi nafsī az naṣārā ke al-yawm kāfir va mushrik
ast bi-mujarrad-i īmān dar ẓuhūr-i baʿd ān sidre-yi ilāhī ū rā bi-alṭaf
va araqq-i aʿḍā-yi mubārak-i khud nisbat dāde va mansūb sākhte-and va
maʿa-dhālik īn hamaj-i arḍ dar che rutbe vāqif va nāẓirand fa-suḥqan
lakum yā malaʾa l-ghāfilīn.

Now, O people of insight, consider how gracious and tender the Cause of
God is, and how God's all-embracing mercy encompasses the whole of
existence: solely on account of faith in the next Revelation, that
divine Lote-Tree has associated a Christian who today is an unbeliever
and one who associates partners with God with the most delicate and
tender of His blessed bodily members, and has linked him to them. Yet at
what a level these rabble of the earth remain, and from what a
standpoint they look! Away with you, O company of the heedless!

\subsection{150}\label{section-1705}

\begin{otherlanguage}{persian}

فوالّذی جعل الیوم دلیله نفسه و سلطانه ما یجری من قلمه المنیر که اگر کلّ من
فی البیان الیوم معرض شوند و یک نفر از نصاری مؤمن هرآینه اسماء کلّ از لوح
محفوظ محو خواهد شد و اسم نصاری ثابت یمحو اللّه ما یشآء و یثبت و عنده امّ
الکتاب

\end{otherlanguage}

Fa-wa-lladhī jaʿala l-yawma dalīlahu nafsahu wa-sulṭānahu mā yajrī min
qalamihi l-munīr ke agar kull-i man fī l-bayān al-yawm muʿriḍ shavand va
yak nafar az naṣārā muʾmin har-āyine asmāʾ-i kull az lawḥ-i maḥfūẓ maḥv
khwāhad shud va ism-i naṣārā thābit yamḥū llāhu mā yashāʾu wa-yuthbitu
wa-ʿindahu ummu l-kitāb.

I swear by Him who today has made His own Self His proof and what flows
from His luminous Pen His authority: if all who are in the Bayán should
turn away today and one Christian become a believer, the names of all
would assuredly be effaced from the Preserved Tablet, and the name of
the Christian confirmed. God effaces whatever He wills and confirms, and
with Him is the Mother of the Book.

\subsection{151}\label{section-1706}

\begin{otherlanguage}{persian}

و از جمیع گذشته آیا حکایت بلعام که در کتب الهی مسطور است نشنیده‌اند که با
آنکه یکی از اوصیای انبیاء قبل بوده و جمیع اهل اقطار عالم خدمتش را فرض
میشمردند و طاعتش را طاعة اللّه میدانستند و اشتهار اسمش جمیع دیار را احاطه
نموده بود و چون جمال کلیم از افق عزّ تسلیم ظاهر شد بمعارضهٴ آن جمال احدیّه
قیام نمود و هم‌چنین یهودای اسخریوطی که از حروف اثنی‌عشر انجیل بود حضرت روح
را بید یهود تسلیم نمود

\end{otherlanguage}

Va az jamīʿ gudashte āyā ḥikāyat-i Balʿām ke dar kutub-i ilāhī masṭūr
ast nashanīde-and ke bā ānke yakī az awṣiyā-yi anbiyā-yi qabl būde va
jamīʿ-i ahl-i aqṭār-i ʿālam khidmatash rā farḍ mīshumurdand va ṭāʿatash
rā ṭāʿata llāh mīdānistand va ishtihār-i ismash jamīʿ-i diyār rā iḥāṭe
namūde būd va chūn jamāl-i Kalīm az ufuq-i ʿizz-i taslīm ẓāhir shud
bi-muʿāriḍe-yi ān jamāl-i aḥadiyye qiyām namūd va ham-chunīn Yahūdā-yi
Iskharīyūṭī ke az ḥurūf-i ithnā-ʿashar-i injīl būd Ḥaḍrat-i Rūḥ rā
bi-yad-i yahūd taslīm namūd.

Beyond all this, have they not heard the account of Balaam recorded in
the divine Books? Although he was one of the appointed successors of the
former Prophets, and all the peoples throughout the world considered
service to him obligatory and obedience to him obedience to God, and the
fame of his name had spread through every land, when the beauty of Moses
appeared from the horizon of glorious submission, he arose to oppose
that Beauty of oneness. Likewise, Judas Iscariot, who was one of the
twelve Letters of the Gospel, delivered the Spirit into the hands of the
Jews.

\subsection{152}\label{section-1707}

\begin{otherlanguage}{arabic}

زینهار از افتتان و امتحان الهی آسوده نشوید و در کلّ حین پناه بخداوند
متعال برده که شاید الیوم از صراط واضح منحرف نشوید و نلغزید فطوبی
للثّابتین چه که امتحان برای کلّ بوده و خواهد بود و احدی از کمندش خارج نه
الّا من شآء ربّک اگر نسیم عدلش مرور نماید صدهزار جوهر نور را در ظلمت نفس و
هوی مشاهده نمائی و اگر نسیم فضلش هبوب فرماید صدهزار هیاکل فانیه را بمقرّ
عرش باقیه ملاحظه کنی

\end{otherlanguage}

Zīnhār az iftitān va imtiḥān-i ilāhī āsūde nashavīd va dar kull-i ḥīn
panāh bi-Khudāvand-i mutaʿāl burde ke shāyad al-yawm az ṣirāṭ-i vāḍiḥ
munḥarif nashavīd va nalaghzīd fa-ṭūbā li-l-thābitīn che ke imtiḥān
barā-yi kull būde va khwāhad būd va aḥadī az kamandash khārij na illā
man shāʾa rabbuk agar nasīm-i ʿadlash murūr namāyad ṣad-hazār jawhar-i
nūr rā dar ẓulmat-i nafs va havā mushāhade namāʾī va agar nasīm-i
faḍlash hubūb farmāyad ṣad-hazār hayākil-i fāniye rā bi-maqarr-i ʿarsh-i
bāqiye mulāḥiẓe kunī.

Beware lest you feel secure from divine trials and tests. At every
moment seek refuge in God, the Most Exalted, so that perhaps today you
may neither stray from the clear path nor stumble. Blessed are the
steadfast! For testing has been and will remain the lot of all, and no
one is beyond its snare except those whom your Lord wills. Should the
breeze of His justice pass by, you would see a hundred thousand essences
of light in the darkness of self and desire; and should the breeze of
His grace blow, you would behold a hundred thousand mortal forms at the
seat of the everlasting throne.

\subsection{153}\label{section-1708}

\begin{otherlanguage}{arabic}

فواللّه یا قوم مطّلع نیستید و بمجرّد وهم حرکت مینمائید بشنوید نصایح بدیعهٴ
الهیّه را و در امر اللّه تجرّی ننمائید و خود را بعذاب لاینتهی معذّب مسازید
از بحر اعظم موّاج بسراب بقیعه خود را مشغول نکنید و از شمس عزّ باقی بظلّ
فانی قانع مشوید بیان حقّ را از لسان رحمن بشنوید

\end{otherlanguage}

Fa-wa-llāh yā qawm muṭṭaliʿ nīstīd va bi-mujarrad-i vahm ḥarakat
mīnamāʾīd bishnavīd naṣāyiḥ-i badīʿe-yi ilāhiyye rā va dar amr-i Allāh
tajarrī nanamāʾīd va khud rā bi-ʿadhāb-i lā yantahī muʿadhdhab masāzīd
az baḥr-i aʿẓam-i mavvāj bi-sarāb-i baqīʿe khud rā mashghūl nakunīd va
az shams-i ʿizz-i bāqī bi-ẓill-i fānī qāniʿ mashavīd bayān-i Ḥaqq rā az
lisān-i Raḥmān bishnavīd.

By God, O people, you are unaware and act on mere conjecture. Listen to
the wondrous counsels of God, do not act presumptuously in the Cause of
God, and do not subject yourselves to unending torment. Do not turn from
the surging Most Great Ocean to occupy yourselves with a mirage on a
barren plain, and do not forsake the Sun of everlasting glory to content
yourselves with a fleeting shadow. Hear the utterance of the True One
from the tongue of the All-Merciful.

\subsection{154}\label{section-1709}

\begin{otherlanguage}{persian}

احدی الیوم بر بیان و ما کنز فیه مطّلع نه و علمه عندنا لو انتم من العارفین
از سلطان مسمّی که ملکوت اسماء بقول او خلق شده بقمیص اسم محتجب نشوید چه که
جمیع اسماء بمنزلهٴ قمایصند یلبسها علی من یشآء فضلاً من عنده و ینزعها عمّن
یشآء عدلاً من لدنه و لا یسأل عمّا یفعل و انّه لهو المقتدر القدیر

\end{otherlanguage}

Aḥadī al-yawm bar Bayān va mā kuniza fīhi moṭṭaleʿ na, va ʿilmuhu
ʿindanā law antum mina al-ʿārifīn. Az Solṭān-e mosammā ke malakūt-e
asmāʾ be-qawl-e ū khalq shode be-qamīṣ-e esm moḥtajeb nashavīd, che ke
jamīʿ-e asmāʾ be-manzale-ye qamāyeṣ-and; yalbasuhā ʿalā man yashāʾ
faḍlan min ʿindihi wa yanziʿuhā ʿamman yashāʾ ʿadlan min ladunhi wa lā
yusʾalu ʿammā yafʿalu wa innahu la-huwa al-Muqtadir al-Qadīr.

No one today is aware of the Bayán and what is treasured within it. Its
knowledge is with Us, if you are among those who understand. Do not let
the garment of a name veil you from the Sovereign who bears the names,
through whose word the kingdom of names was created. For all names are
like garments: He clothes whomever He wills in them as a grace from Him,
and strips them from whomever He wills in His justice. He is not
questioned about what He does, and He is indeed the All-Powerful, the
Omnipotent.

\subsection{155}\label{section-1710}

\begin{otherlanguage}{persian}

و الیوم اکثری از اهل بیان عبدهٴ اسماء بین یدی اللّه مذکور بوده و خواهند
بود چنانچه اگر بنقطهٴ اولی بنفسه لنفسه عارف شده بودند فواللّه اقلّ من آن
در این ظهور قدس صمدانی و تجلّی عزّ رحمانی توقّف نمینمودند

\end{otherlanguage}

Va al-yawm aktharī az ahl-e Bayān ʿabade-ye asmāʾ bayna yaday Allāh
madhkūr būde va khāhand būd, chenānche agar be-Noqṭe-ye Ūlā bi-nafsihi
li-nafsihi ʿāref shode būdand, fa-wa-Allāh aqall-e men ān dar īn ẓohūr-e
qods-e ṣamadānī va tajallī-ye ʿezz-e raḥmānī tavaqqof nemīnamūdand.

Today most of the people of the Bayán have been, and will remain,
designated in the presence of God as worshippers of names. For had they
recognized the Primal Point through His own Self and for His own sake,
then, by God, they would not have hesitated for less than an instant
before this Revelation of eternal holiness and this effulgence of
merciful glory.

\subsection{156}\label{section-1711}

\begin{otherlanguage}{persian}

پس معلوم شد که عارف بنفس او جلّ و عزّ نبوده‌اند بلکه بقمیص اسمی او تعالی
شأنه عارف شده‌اند چنانچه ایّامی نگذشت و بنفسه لنفسه ظاهر شده و او را
نشناخته‌اند و بر او وارد آورده‌اند آنچه را که قلم حیا میکند از ذکرش

\end{otherlanguage}

Pas maʿlūm shod ke ʿāref be-nafs-e ū jalla wa ʿazza nabūde-and, balke
be-qamīṣ-e esmī-ye ū taʿālā shaʾnuhu ʿāref shode-and, chenānche ayyāmī
nagodhasht va bi-nafsihi li-nafsihi ẓāher shode va ū rā nashenākhte-and
va bar ū vāred āvorde-and ānche rā ke qalam ḥayā mīkonad az dhekrash.

It is thus clear that they had not recognized His own Self, glorious and
mighty is He, but had recognized the garment of His name, exalted is His
station. For only a few days passed before He appeared through His own
Self and for His own sake, yet they failed to recognize Him and
inflicted on Him what the Pen is ashamed to mention.

\subsection{157}\label{section-1712}

\begin{otherlanguage}{persian}

و این عدم عرفان نشده مگر آنکه تبدیل قمیص فرموده یعنی اسم اعلی باسم ابهی
تبدیل شد و احدی عارف بآن نشده الّا الّذینهم انقطعوا عن کلّ من فی الملک و
صعدوا الی اللّه بکلّهم الی ان دخلوا مقرّ عزّ کریم اولئک اولیائی تحت ظلال
عنایتی و

\end{otherlanguage}

Va īn ʿadam-e ʿerfān nashode magar ānke tabdīl-e qamīṣ farmūde, yaʿnī
esm-e Aʿlā be-esm-e Abhā tabdīl shod va aḥadī ʿāref be-ān nashode illā
alladhīna hum inqaṭaʿū ʿan kulli man fī al-mulk wa ṣaʿidū ilā Allāh
bi-kullihim ilā an dakhalū maqarra ʿizzin karīm. Ulāʾika awliyāʾī taḥta
ẓilāli ʿināyatī wa

This failure to recognize Him came about only because He changed His
garment: that is, the name Most Exalted was changed to the name
All-Glorious. No one recognized Him except those who detached themselves
from everyone in the realm of creation and ascended to God with their
whole being, until they entered an abode of noble glory. These are My
loved ones beneath the shade of My loving care, and

\subsection{158}\label{section-1713}

\begin{otherlanguage}{persian}

ینزل علیهم سکینة من عندی و هم علی سرر الاستقلال هم مستقرّون اولئک حفظهم
اللّه عن رمی الشّبهات و هم علی وساید القدس متّکئون و الی طرفی ناظرون علیهم
فی کلّ حین بهائی ثمّ رحمتی و عنایتی ثمّ فضلی و افضالی و کذلک اختصّهم اللّه
لنفسی و جعلهم من المؤمنین

\end{otherlanguage}

Yanzilu ʿalayhim sakīnatun min ʿindī wa hum ʿalā sururi al-istiqlāl hum
mustaqirrūn. Ulāʾika ḥafiẓahumu Allāhu ʿan ramyi al-shubuhāt wa hum ʿalā
wasāʾidi al-quds muttakiʾūn wa ilā ṭarfī nāẓirūn. ʿAlayhim fī kulli
ḥīnin bahāʾī thumma raḥmatī wa ʿināyatī thumma faḍlī wa ifḍālī wa
kadhālika ikhtaṣṣahumu Allāhu li-nafsī wa jaʿalahum mina al-muʾminīn.

Tranquility descends upon them from Me, and they rest upon the couches
of independence. These are the ones God has protected from the onslaught
of doubts. They recline upon the cushions of holiness and gaze toward
Me. Upon them at every moment be My glory, My mercy and loving care, My
grace and bounty. Thus has God chosen them for My own Self and numbered
them among the believers.

\subsection{159}\label{section-1714}

\begin{otherlanguage}{arabic}

ای عباد سعی نموده که شاید از بدایع فضل ربّ الأرباب در ظلّ قباب آفتاب قدس
عنایت وارد شوید و از حرارت نفس و هوی آسوده و محفوظ مانید اینست نصایح قلم
امر فمن سمع فلنفسه و من اعرض فلها و انّه لمقدّس عن العالمین

\end{otherlanguage}

Ey ʿebād, saʿy namūde ke shāyad az badāyeʿ-e faḍl-e Rabb al-Arbāb dar
ẓell-e qebāb-e āftāb-e qods-e ʿenāyat vāred shavīd va az ḥarārat-e nafs
va havā āsūde va maḥfūẓ mānīd. Īnast naṣāyeḥ-e qalam-e amr; fa-man
samiʿa fa-li-nafsihi wa man aʿraḍa fa-lahā wa innahu la-muqaddasun ʿani
al-ʿālamīn.

O servants, strive, so that through the wondrous grace of the Lord of
Lords you may enter beneath the shade of the canopies of the holy Sun of
loving care, and remain sheltered and protected from the heat of self
and desire. These are the counsels of the Pen of the Cause. Whoever
listens does so for his own sake, and whoever turns away does so to his
own detriment. He is indeed sanctified above all the worlds.

\subsection{160}\label{section-1715}

\begin{otherlanguage}{persian}

و ایکاش که اهل بیان در بیان فارسی که از قلم رحمانی نازل نظر مینمودند که
شاید در این ظهور قدس صمدانی خود را از سلسبیل عنایت ربّانی و تسنیم مکرمت
سبحانی محروم نمینمودند و بشأنی شبهات بر قلوب غیر طاهره وارد شده که لن
یعرف ایّاً من ایّ و اگر قلبی از کوثر بی‌زوال سلطان لایزال از شبهه مطهّر شود
فی الحین بشبههٴ دیگر بیالاید سبحان‌اللّه با اینکه اهل بیان خود مقرّ و
معترفند که نبوّت برسول اللّه ختم شده و سنهٴ ستّین اوّل ظهور اللّه است مع‌ذلک و
مع اینکه حقّ حال بنفسه و کینونته مشهود و از افق انّنی انا حیّ فی الأفق
الأبهی ظاهر و طالع مجدّداً ذکر وصایت و امثال آن مینمایند و ببقیّهٴ وهم ما
ترک من ملل الفرقان اراده نموده‌اند انوار شمس یقین را که در سماء قدس
بی‌زوال ذو الجلال مشرق شده ستر نمایند نیکو است ذکر این دو فرد در این مقام

\end{otherlanguage}

Va ey kāsh ke ahl-e Bayān dar Bayān-e Fārsī ke az qalam-e raḥmānī nāzel
naẓar mīnamūdand ke shāyad dar īn ẓohūr-e qods-e ṣamadānī khod rā az
salsabīl-e ʿenāyat-e rabbānī va tasnīm-e makramat-e sobḥānī maḥrūm
nemīnamūdand. Va be-shaʾnī shobohāt bar qolūb-e gheyr-e ṭāhere vāred
shode ke lan yuʿrafa ayyan min ayy, va agar qalbī az kawthar-e
bī-zavāl-e Solṭān-e lā-yazāl az shobhe moṭahhar shavad fī al-ḥīn
be-shobhe-ye dīgar bīālāyad. Sobḥān Allāh, bā īnke ahl-e Bayān khod
moqerr va moʿtaref-and ke nobovvat be-Rasūl Allāh khatm shode va sane-ye
settīn avval-e ẓohūr Allāh ast, maʿa dhālek va maʿa īnke Ḥaqq ḥāl
bi-nafsihi va kaynūnatihi mashhūd va az ofoq-e innanī anā ḥayyun fī
al-ufuq al-Abhā ẓāher va ṭāleʿ, mojaddadan dhekr-e veṣāyat va amthāl-e
ān mīnamāyand va be-baqiyye-ye vahm-e mā taraka min milal al-Furqān
erāde namūde-and anvār-e shams-e yaqīn rā ke dar samā-ye qods-e
bī-zavāl-e Dhu al-Jalāl moshreq shode satr namāyand. Nīkū ast dhekr-e īn
do fard dar īn maqām.

If only the people of the Bayán would study the Persian Bayán, revealed
by the Pen of the All-Merciful, so that they might not deprive
themselves, in this Revelation of eternal holiness, of the Salsabíl of
divine loving care and the Tasním of heavenly generosity. Doubts have so
beset impure hearts that one thing cannot be distinguished from another;
and if a heart is cleansed of doubt by the unfailing Kawthar of the
everlasting Sovereign, it is immediately defiled by another doubt. Glory
be to God! Although the people of the Bayán themselves acknowledge and
confess that prophethood ended with the Messenger of God and that the
year sixty marks the beginning of the Revelation of God, and although
the True One is now visible in His own Self and Being, manifest and
radiant from the horizon of ``Indeed, I am living in the All-Glorious
Horizon,'' they nevertheless once again speak of successorship and the
like. With the remnants of delusion left by the communities of the
Qur'an, they have sought to veil the rays of the Sun of certitude that
has risen in the heaven of the unfading holiness of the Lord of Majesty.
It is fitting to cite these two couplets here.

\subsection{161}\label{section-1716}

\begin{otherlanguage}{arabic}

ای ضیآء الحقّ حسام الدّین و دل ای دل و جان از قدوم تو خجل قصد این دارند
این گل‌پاره‌ها که از حسد پوشند خورشید تو را

\end{otherlanguage}

Ey Ḍiyāʾ al-Ḥaqq Ḥosām al-Dīn va del, ey del o jān az qodūm-e to khejel,
qaṣd-e īn dārand īn gel-pāreh-hā ke az ḥasad pūshand khorshīd-e to rā.

O Radiance of Truth, Sword of Religion and of the heart, O you before
whose coming heart and soul are abashed! These clods of earth intend,
out of envy, to conceal your sun.

\subsection{162}\label{section-1717}

\begin{otherlanguage}{persian}

بگذارید این اقوال لایغنیه را و باصل امر و ما حقّق به الأسمآء ناظر شوید
ظهور قبلم میفرماید اگر ناس ضعیف ملاحظه نمیشدند هرآینه ذکر مشیّت نمیفرمودم
حال تقدیس امر و ارتفاع ما اراد اللّه لکم را مشاهده نمائید و قدر و رتبهٴ
خود را ملاحظه کنید که حال بادلّهٴ پست‌ترین عباد از اهل فرقان تمسّک جسته و
مستدلّ شده‌اند فوا حسرتا علیکم یا ملأ الغافلین

\end{otherlanguage}

Begd hārīd īn aqvāl-e lā-yughniye rā va be-aṣl-e amr va mā ḥuqqiqa bihi
al-asmāʾ nāẓer shavīd. Ẓohūr-e qablam mīfarmāyad agar nās ḍaʿīf molāḥeẓe
nemīshodand har-āyene dhekr-e mashiyyat nemīfarmūdam. Ḥāl taqdīs-e amr
va ertefāʿ-e mā arāda Allāhu lakum rā moshāhede namāyīd va qadr va
rotbe-ye khod rā molāḥeẓe konīd ke ḥāl be-adelle-ye past-tarīn-e ʿebād
az ahl-e Forqān tamassok joste va mostadell shode-and. Fa-wā ḥasratā
ʿalaykum yā malaʾ al-ghāfilīn.

Set aside these fruitless sayings and turn your gaze to the source of
the Cause and that through which the names were established. My previous
Manifestation says, ``Had people not been regarded as weak, I would
assuredly not have mentioned the Will.'' Now behold the sanctity of the
Cause and the loftiness of what God has intended for you, and consider
your own worth and rank: they have now resorted to the arguments of the
lowliest servants among the people of the Qur'an and made these their
proofs. Alas for you, O company of the heedless!

\subsection{163}\label{section-1718}

\begin{otherlanguage}{persian}

باری ایّام ظهور است و حقّ چون شمس فی وسط السّمآء ظاهر و زخرفی از دنیا در
ساحت قدسش موجود نه که محتاج بوصیّ باشد آن را قسمت نماید و اگر میگوئید در
امرش محتاج بوصیّ است هذا بغی منکم علی اللّه المهیمن القیّوم چه که امر طایف
حول نفسش بوده و خواهد بود و لن یفارق منه ابداً ایّاکم ان تجعلوا للّه وزیراً
او نظیراً او ضدّاً او ندّاً او وصیّاً او قریناً او مثالاً و در بیان این اسامی و
ما حدّدتم به انفسکم جمیع را محو فرموده‌اند

\end{otherlanguage}

Bārī ayyām-e ẓohūr ast va Ḥaqq chon shams fī wasaṭ al-samāʾ ẓāher va
zokhrofī az donyā dar sāḥat-e qodsash mawjūd na ke moḥtāj be-vaṣī bāshad
ān rā qesmat namāyad. Va agar mīgūyīd dar amrash moḥtāj be-vaṣī ast,
hādhā baghyun minkum ʿalā Allāh al-Muhaymin al-Qayyūm, che ke amr
ṭāyef-e ḥawl-e nafsash būde va khāhad būd va lan yufāriqa minhu abadan.
Iyyākum an tajʿalū li-llāhi wazīran aw naẓīran aw ḍiddan aw niddan aw
waṣiyyan aw qarīnan aw mithālan. Va dar Bayān īn asāmī va mā ḥaddadtum
bihi anfusakum jamīʿ rā maḥv farmūde-and.

In short, these are the days of Revelation, and the True One is manifest
like the sun in the midst of heaven. There are no worldly riches in His
holy court that He should need an executor to divide them. And if you
say that He needs a successor in His Cause, this is an outrage on your
part against God, the All-Protecting, the Self-Subsisting. For the Cause
has always revolved, and will always revolve, around His own Self, and
will never be separated from Him. Beware of assigning to God a minister,
a peer, an adversary, a rival, a successor, a companion, or a likeness.
In the Bayán He has effaced all these names and everything by which you
have confined yourselves.

\subsection{164}\label{section-1719}

\begin{otherlanguage}{persian}

اینست که ذکر مرایا در بیان من قبل الرّحمن نازل شده و آن هم محدود نبوده و
نخواهد بود امیدواریم که از رشحات سماء تقدیس امطار قدسیّه بر قلوب غبرهٴ
مکدّره مبذول شود که شاید از اشارات طاهر شده بمنظر ملیک اسماء و صفات توجّه
نمایند فواللّه یا قوم مقصود از ذکر این براهین متقنه و ادلّهٴ محکمه و کلمات
غیر متشابهه اثبات نفسم نبوده چه که هر نفسی که بأقلّ عمّا یحصی ببصر معنوی
فایز شده البتّه انوار شمس ابهی را که در کلّ شیء علی هیئة انّه لا اله الّا
انا العزیز الحکیم تجلّی فرموده مشاهده مینماید

\end{otherlanguage}

Īnast ke dhekr-e marāyā dar Bayān min qibal al-Raḥmān nāzel shode va ān
ham maḥdūd nabūde va nakhāhad būd. Omīdvārīm ke az rashaḥāt-e samā-ye
taqdīs amṭār-e qodsiyye bar qolūb-e ghabere-ye mokaddare mabdhūl shavad
ke shāyad az eshārāt ṭāher shode be-manẓar-e Malīk-e asmāʾ va ṣefāt
tavajjoh namāyand. Fa-wa-Allāh yā qawm, maqṣūd az dhekr-e īn barāhīn-e
motqane va adelle-ye moḥkame va kalemāt-e gheyr-e motashābehe ethbāt-e
nafsam nabūde, che ke har nafsī ke be-aqall-e ʿammā yuḥṣā be-baṣar-e
maʿnavī fāyez shode albatte anvār-e shams-e Abhā rā ke dar kull-e shayʾ
ʿalā hayʾati innahu lā ilāha illā anā al-ʿAzīz al-Ḥakīm tajallī farmūde
moshāhede mīnamāyad.

Thus it is that mention of mirrors has been revealed in the Bayán by the
All-Merciful, and this too has not been, and will not be, restricted. We
hope that, from the showers of the heaven of sanctity, holy rains may be
bestowed upon dusty, clouded hearts, so that they may be cleansed of
allusions and turn toward the countenance of the Sovereign of names and
attributes. By God, O people, My purpose in setting forth these
conclusive proofs, these firm arguments, and these unambiguous words has
not been to establish My own claim. For anyone endowed with even an
immeasurably small degree of spiritual vision will assuredly behold the
rays of the All-Glorious Sun, which has shone forth in all things in the
form of ``Indeed, there is no God but Me, the Almighty, the All-Wise.''

\subsection{165}\label{section-1720}

\begin{otherlanguage}{arabic}

بلکه مقصودی من یأتی فی المستغاث تاللّه اذاً یبکی قلمی و اناملی ثمّ عینی و
قلبی و کبدی و عروقی و دمی و عظمی و جلدی بما یرد علیه من هؤلآء الّذینهم
کفروا باللّه و حاربوا بنفسه و جادلوا بآیاته و انکروا فضله و افتوا علیه
بعد الّذی ظهر بکلّ الآیات و الحجج و العلامات

\end{otherlanguage}

Balke maqṣūdam man yaʾtī fī al-Mustaghāth. Ta-Allāh idhan yabkī qalamī
wa anāmilī thumma ʿaynī wa qalbī wa kabidī wa ʿurūqī wa damī wa ʿaẓmī wa
jildī bi-mā yaridu ʿalayhi min hāʾulāʾ alladhīna hum kafarū bi-Allāh wa
ḥārabū bi-nafsihi wa jādalū bi-āyātihi wa ankarū faḍlahu wa aftaw
ʿalayhi baʿda alladhī ẓahara bi-kulli al-āyāt wa al-ḥujaj wa al-ʿalāmāt.

Rather, My purpose concerns Him who will come in the Mustagháth. By God,
at this My pen and My fingers weep, as do My eye, My heart, My liver, My
veins, My blood, My bones, and My skin, over what will befall Him at the
hands of those who have disbelieved in God, fought against His own Self,
disputed His verses, denied His grace, and pronounced judgment against
Him after He appeared with every verse, proof, and sign.

\subsection{166}\label{section-1721}

\begin{otherlanguage}{arabic}

فواللّه فعلتم ما لا فعل احد من ملل القبل و ارتکبتم ما شقّ عنه هیکل العظمة
ردآء الصّبر ولکن انتم تمشون علی وجه الأرض کمن لا ذنب له فأفّ لکم یا ملأ
المشرکین

\end{otherlanguage}

Fa-wa-Allāh faʿaltum mā lā faʿala aḥadun min milal al-qabl wa irtakabtum
mā shaqqa ʿanhu Haykal al-ʿAẓama ridāʾ al-ṣabr wa lākin antum tamshūna
ʿalā wajh al-arḍ ka-man lā dhanba lahu fa-uffin lakum yā malaʾ
al-mushrikīn.

By God, you have done what no one among the communities of the past has
done, and committed that for which the Temple of Grandeur has torn the
robe of patience. Yet you walk upon the earth like those who have
committed no sin. Shame on you, O company of those who associate
partners with God!

\subsection{167}\label{section-1722}

\begin{otherlanguage}{persian}

تاللّه یا قوم انّ فی الکذب روایح و کذلک فی الصّدق لو انتم تطهّرون قلوبکم و
اشمامکم و تستنشقوا لتجدنّ رایحة اللّه المهیمن القیّوم قلم متحیّر است که چه
ذکر نماید و حال امر بمقامی رسیده که ساذج روحی که اگر محلّ خطاب ابهی کلمات
عالمین واقع شود خود را محزون مشاهده مینماید باید باستماع ادنی کلمات
مشغول شود

\end{otherlanguage}

Ta-Allāh yā qawm inna fī al-kadhib rawāʾiḥ wa kadhālika fī al-ṣidq law
antum tuṭahhirūna qulūbakum wa ashmāmakum wa tastanshiqū la-tajidunna
rāʾiḥat Allāh al-Muhaymin al-Qayyūm. Qalam motaḥayyer ast ke che dhekr
namāyad va ḥāl amr be-maqāmī resīde ke sādhej-e rūḥī ke agar maḥall-e
kheṭāb-e abhā kalemāt-e ʿālamīn vāqeʿ shavad khod rā maḥzūn moshāhede
mīnamāyad bāyad be-estemāʿ-e adnā kalemāt mashghūl shavad.

By God, O people, falsehood has its odors, and truth likewise. If you
would purify your hearts and your sense of smell and breathe deeply, you
would surely perceive the fragrance of God, the All-Protecting, the
Self-Subsisting. The Pen is bewildered as to what to mention. Matters
have now reached such a point that the pure essence of spirit, which
would find itself saddened even to be addressed by the most glorious
words in all the worlds, must occupy itself with listening to the
lowliest words:

\subsection{168}\label{section-1723}

\begin{otherlanguage}{persian}

که چگونه میشود ذئب راعی اغنام شود بگو ای فقیر بی‌بصر اوّلاً راعی را نشناخته
و ثانیاً سالها ذیاب راعی شما بوده چنانچه بعد از ظهور نیّر الهی از افق سماء
معنوی ادراک نموده‌اید که متابعت قومی مینمودید که صدهزار مرتبه ذئب از آن
نفوس عنداللّه ارجح بوده چه که بفتوای آن ذیاب مظهر ربّ الأرباب بتمام ظلم
شهید شد لعنة اللّه علیهم و علی اتباعهم

\end{otherlanguage}

Ke chegūne mīshavad dheʾb rāʿī-ye aghnām shavad. Begū ey faqīr-e
bī-baṣar, avvalan rāʿī rā nashenākhte va thāniyan sālhā dheʾāb rāʿī-ye
shomā būde, chenānche baʿd az ẓohūr-e nayyer-e elāhī az ofoq-e samā-ye
maʿnavī edrāk namūde-īd ke motābeʿat-e qawmī mīnamūdīd ke ṣad-hezār
martabe dheʾb az ān nofūs ʿinda Allāh arjaḥ būde, che ke be-fatvā-ye ān
dheʾāb Maẓhar-e Rabb al-Arbāb be-tamām-e ẓolm shahīd shod. Laʿnat Allāh
ʿalayhim wa ʿalā atbāʿihim.

``How can a wolf become the shepherd of the sheep?'' Say: O sightless,
impoverished one! First, you have not recognized the Shepherd; and
second, for years wolves have been your shepherds. For after the divine
Luminary rose above the horizon of the spiritual heaven, you realized
that you had been following a people to whom, in God's sight, a wolf was
a hundred thousand times preferable. For it was through the ruling of
those wolves that the Manifestation of the Lord of Lords was martyred
with utter cruelty. God's curse be upon them and upon their followers.

\subsection{169}\label{section-1724}

\begin{otherlanguage}{persian}

و اگر بگوئی قبل از ظهور متابعت آن قوم مجری و عنداللّه مقبول بود و بعد از
ظهور و توقّف این قوم در امر الهی از نار محسوب شده‌اند حال هم در این ظهور
قدس صمدانی همان حکم را جاری کن بشنوید نصح الهی را و قلب را از اشارات
کلمات اهل حجبات مطهّر سازید بگو ای بندگان هوی قدمی بسوی خدا بردارید و از
ارض ظنّ بهوای یقین طیران نمائید

\end{otherlanguage}

Va agar begūʾī qabl az ẓohūr motābeʿat-e ān qawm mojrā va ʿinda Allāh
maqbūl būd va baʿd az ẓohūr va tavaqqof-e īn qawm dar amr-e elāhī az nār
maḥsūb shode-and, ḥāl ham dar īn ẓohūr-e qods-e ṣamadānī hamān ḥokm rā
jārī kon. Beshnavīd noṣḥ-e elāhī rā va qalb rā az eshārāt-e kalemāt-e
ahl-e ḥojobāt moṭahhar sāzīd. Begū ey bandegān-e havā, qadamī be-sū-ye
Khodā bardārīd va az arḍ-e ẓann be-havā-ye yaqīn ṭayarān namāyīd.

And if you say that before the Revelation following those people was the
accepted practice and was acceptable to God, and that after the
Revelation, because they hesitated in the Cause of God, they were
numbered among the people of the fire, then apply the same judgment now
in this Revelation of eternal holiness. Listen to the counsel of God and
purify your hearts of the allusions in the words of those who are
veiled. Say: O servants of desire, take a step toward God and soar from
the land of conjecture into the air of certitude.

\subsection{170}\label{section-1725}

\begin{otherlanguage}{arabic}

ان یا قلم القدم دع حینئذ ما اردت ان تذکره لعباد اللّه ثمّ اذکر اخیک الّذی
حارب بنفسک و جادل بآیات اللّه لیکون ذکری للعالمین لعلّ یفتح بذلک ابصار
الّذین ارادوا الوجه فی العشیّ و الاشراق و ما زلّت اقدامهم فی هذا الیوم الّذی
فیه التفّت السّاق بالسّاق

\end{otherlanguage}

An yā Qalam al-Qidam daʿ ḥīnaʾidhin mā aradta an tadhkurahu li-ʿibād
Allāh thumma udhkur akhīka alladhī ḥāraba bi-nafsika wa jādala bi-āyāt
Allāh li-yakūna dhikrī li-al-ʿālamīn laʿalla yaftaḥu bi-dhālika abṣār
alladhīna arādū al-wajh fī al-ʿashiyy wa al-ishrāq wa mā zallat
aqdāmuhum fī hādhā al-yawm alladhī fīhi iltaffati al-sāqu bi-al-sāq.

O Pen of Eternity, leave aside for now what you intended to mention to
the servants of God. Speak instead of your brother, who fought against
your own Self and disputed the verses of God, that My account may be for
all the worlds. Perhaps through it the eyes of those who have sought the
Face at evening and at dawn may be opened, whose feet have not slipped
on this Day when leg is entwined with leg.

\subsection{171}\label{section-1726}

\begin{otherlanguage}{persian}

فوالّذی نفسی بیده که بیست سنه در تعلیم و تأدیب و تربیت و حفظش سعی بلیغ و
جهد منیع مبذول شد که ذکرش ممکن نه تا آنکه هجرت واقع و باین ارض وارد شدیم
چندی گذشت سیّد محمّد اصفهانی از بیت ‌اللّه و حصن حبّه خارج و بمقرّ دیگر مستقرّ
و اخوی فواللّه الّذی لا اله الّا هو در صدد قطع سدرهٴ ربّانیّه برآمده الی ان
اظهر اللّه خافیة نفسه و اشتهر مکره بین المهاجرین و ارتفع الضّجیج عن کلّ
اناث و ذکور و صغیر و کبیر

\end{otherlanguage}

Fa-wa-alladhī nafsī bi-yadihi ke bīst sane dar taʿlīm va taʾdīb va
tarbiyat va ḥefẓash saʿy-e balīgh va jahd-e manīʿ mabdhūl shod ke
dhekrash momken na, tā ānke hejrat vāqeʿ va be-īn arḍ vāred shodīm.
Chandī godhasht, Seyyed Moḥammad-e Eṣfahānī az Bayt Allāh va ḥeṣn-e
ḥobbash khārej va be-maqarr-e dīgar mostaqarr va akhavī fa-wa-Allāh
alladhī lā ilāha illā huwa dar ṣadad-e qaṭʿ-e sedre-ye rabbāniyye
bar-āmade, ilā an aẓhara Allāhu khāfiyata nafsihi wa ishtahara makruhu
bayna al-muhājirīn wa irtafaʿa al-ḍajīj ʿan kulli ināthin wa dhukūrin wa
ṣaghīrin wa kabīr.

By Him in whose hand is My soul, for twenty years such strenuous efforts
and such immense exertions were devoted to his instruction, discipline,
upbringing, and protection that they cannot be recounted, until the
departure took place and We arrived in this land. Some time passed.
Siyyid Muḥammad-i-Iṣfahání left the House of God and the fortress of His
love and settled elsewhere. And My brother---by God, besides whom there
is no other god---set about cutting down the divine Lote-Tree, until God
laid bare what was hidden within him, his scheming became known among
the emigrants, and an outcry arose from every woman and man, young and
old.

\subsection{172}\label{section-1727}

\begin{otherlanguage}{persian}

و این عبد بالمرّه از جمیع گذشته خارج شده که شاید حرارت حسد بیفسرد و نار
بغضا ساکن شود و در بیت وحده ساکن شده و ابواب خروج و دخول بر کلّ مسدود
نمودم و اهل حرم جمیع از استماع این امر منکر در نوحه و ندبه مشغول فواللّه
لن یقدر احد ان یذکر ما ورد علینا و لن یستطیع ان یسمعه احد من المحبّین

\end{otherlanguage}

Va īn ʿabd bi-al-marre az jamīʿ godhashte khārej shode ke shāyad
ḥarārat-e ḥasad beyafsorad va nār-e baghḍā sāken shavad, va dar bayt
vaḥdahu sāken shode va abvāb-e khorūj va dokhūl bar koll masdūd namūdam.
Va ahl-e ḥaram jamīʿ az estemāʿ-e īn amr-e monkar dar nawḥe va nodbe
mashghūl. Fa-wa-Allāh lan yaqdira aḥadun an yadhkura mā warada ʿalaynā
wa lan yastaṭīʿa an yasmaʿahu aḥadun mina al-muḥibbīn.

This servant relinquished everything and withdrew, so that perhaps the
heat of envy might cool and the fire of hatred subside. I took up
residence alone in a house and closed the doors to everyone coming or
going. All the members of the household, on hearing of this heinous
affair, were occupied with lamentation and weeping. By God, no one will
ever be able to recount what befell Us, nor will any of Our loved ones
be able to bear hearing it.

\subsection{173}\label{section-1728}

\begin{otherlanguage}{persian}

و در حین خروج سه نفر برای خدمت معیّن شد دو نفر مخصوص اخوی مذکور و کلیم و
یک نفس برای خود و هم‌چنین دراهم معدوده که از دولت علیّه میرسید مابین کلّ
قسمت شده و بعد کلّ متفرّق شده در اماکن متعدّده ساکن و این عبد در بیت وحده
جالس مغلق الباب و مرخی الحجاب و ابداً بنفسی از احباب و غیره معاشرت ننموده
بالمرّه باب لقا مسدود و جمال بقا مستور

\end{otherlanguage}

Va dar ḥīn-i khurūj sih nafar barā-yi khidmat muʿayyan shud, du nafar
makhṣūṣ-i akhavī-yi madhkūr va Kalīm va yak nafs barā-yi khud, va
hamchunīn darāhim-i maʿdūdih kih az dawlat-i ʿaliyyih mīrasīd mābayn-i
kull qismat shudih, va baʿd kull mutafarriq shudih dar amākin-i
mutaʿaddidih sākin, va īn ʿabd dar bayt vaḥdahū jālis, mughlaq al-bāb va
murkhā al-ḥijāb, va abadan bi-nafsī az aḥbāb va ghayrih muʿāshirat
nanamūdih, bi'l-marrih bāb-i liqā masdūd va jamāl-i baqā mastūr.

At the time of departure, three people were appointed to serve: two
specifically for the aforementioned brother and Kalīm, and one for
Myself. Likewise, the small sum received from the imperial government
was divided among everyone. Then all dispersed and settled in various
places, while this servant sat alone in the house, with the door closed
and the curtain drawn. I personally associated with none of the friends
or anyone else. The door of access was completely closed, and the Beauty
of Eternity was concealed.

\subsection{174}\label{section-1729}

\begin{otherlanguage}{persian}

و بعد از ستر جمال قدم سیّد محمّد باخوی پیوست و بغض اللّه سبب و علّت اتّحاد
این دو شد و بخدعه و مکر مشغول شدند چنانچه نفسی را بدراهم فریب داده که
ببلاد رود و آنچه ارتکاب نموده‌اند بسدرهٴ امر نسبت دهد دیگر تفصیل بسیار
ولکن اشتغال باین اذکار را جایز نمیدانم البتّه شنیده‌اید بقدر ضرورت اظهار
میشود که مباد ناس از کلمات کذبهٴ فانیه از صراط امریّه بلغزند

\end{otherlanguage}

Va baʿd az satr-i jamāl-i qidam, Siyyid Muḥammad bi-akhavī payvast va
bughḍ-i Allāh sabab va ʿillat-i ittiḥād-i īn du shud va bi-khudʿih va
makr mashghūl shudand, chunānchih nafsī rā bi-darāhim farīb dādih kih
bi-bilād ravad va ānchih irtikāb namūdih-and bi-sidrih-yi amr nisbat
dihad. Dīgar tafṣīl bisyār, valākin ishtighāl bi-īn adhkār rā jāʾiz
namīdānam. Albattih shinīdih-īd. Bi-qadr-i ḍarūrat iẓhār mīshavad kih
mabād nās az kalimāt-i kadhabih-yi fāniyih az ṣirāṭ-i amriyyih
bilaghzand.

After the Ancient Beauty withdrew into concealment, Siyyid Muḥammad
joined My brother, and hatred of God became the cause and reason for
their union. They occupied themselves with deceit and intrigue, even
bribing someone with money to travel to the provinces and attribute to
the Lote-Tree of the Cause what they themselves had committed. There is
much more to recount, but I do not consider it fitting to occupy Myself
with these matters. You have surely heard. Only as much as is necessary
is disclosed, lest people be led by the words of these perishable liars
to stray from the path of the Cause.

\subsection{175}\label{section-1730}

\begin{otherlanguage}{persian}

تا آنکه نفس مأمور نادم شده نزد بعضی رفته تفصیل امور را مذکور نمود و چون
مراسلاتی که با او همراه نموده مشاهده نمودند جمیع را تحیّر دست داده چه که
آنچه عامل شده بحقّ راجع نموده‌اند و هم‌چنین رساله در ردّ اللّه نوشته که
باطراف برده انتشار دهند فواللّه الّذی لا اله الّا هو بمثل صبیان بل ادنی
تکلّم نموده لو انتم تعرفون و آنچه را بزعم خود در ردّ اللّه نوشته‌اند کلّ مثبت
این امر ابدع بوده لو کان النّاس ببصر اللّه فی کلماتهم ینظرون

\end{otherlanguage}

Tā ānkih nafs-i maʾmūr nādim shudih nazd-i baʿḍī raftih tafṣīl-i umūr rā
madhkūr namūd, va chūn murāsalātī kih bā ū hamrāh namūdih mushāhidih
namūdand jamīʿ rā taḥayyur dast dādih, chih kih ānchih ʿāmil shudih
bi-Ḥaqq rājiʿ namūdih-and. Va hamchunīn risālih dar radd-i Allāh
nivishtih kih bi-aṭrāf burdih intishār dihand. Fa-wa-llāhi alladhī lā
ilāha illā huwa, bi-mithl-i ṣibyān bal adnā takallum namūdih, law antum
taʿrifūn. Va ānchih rā bi-zaʿm-i khud dar radd-i Allāh nivishtih-and
kull muthbit-i īn amr-i abdaʿ būdih, law kāna al-nāsu bi-baṣari Allāhi
fī kalimātihim yanẓurūn.

Eventually, the person they had commissioned became remorseful, went to
certain people, and related the details of these affairs. When they saw
the correspondence he had been given to carry, they were all seized with
astonishment, for everything those men had done they had attributed to
the True One. They had also written a treatise refuting God, to be taken
to various regions and circulated. By God, besides whom there is no
other god, they have spoken like children---indeed, worse---if only you
understood. Everything they imagined they had written in refutation of
God has, in fact, affirmed this wondrous Cause, if only people would
examine their words with the eye of God.

\subsection{176}\label{section-1731}

\begin{otherlanguage}{arabic}

و دیگر چه مقدار مفتریات بملیک اسماء و صفات نسبت داده‌اند و آنچه دون خیر
بوده باصل شجرهٴ ربوبیّه راجع نموده انّما اشکو بثّی و حزنی الی اللّه اذاً
یتکلّم لسان اللّه و یقول

\end{otherlanguage}

Va dīgar chih miqdār muftariyāt bi-Malīk-i asmā va ṣifāt nisbat
dādih-and, va ānchih dūn-i khayr būdih bi-aṣl-i shajari h-yi rubūbiyyih
rājiʿ namūdih. Innamā ashkū baththī wa-ḥuznī ilā Allāh. Idhan
yatakallamu lisānu Allāhi wa-yaqūl.

How many other fabrications they have attributed to the Sovereign of
names and attributes! Whatever was not good they have ascribed to the
very root of the Tree of Lordship. ``I complain of my anguish and sorrow
only to God.'' Then the Tongue of God speaks and says:

\subsection{177}\label{section-1732}

\begin{otherlanguage}{persian}

هل من مآء صافی یطهّر حولی عن هؤلآء المعرضین و هل من سیف لسان یفتح بالبیان
و یدفع عن فنائی شرّ المشرکین و هل من ذی استقامة ینصرنی فی هذه الأیّام الّتی
احاطتنی المنکرین و هل من ذی انقطاع یعرفنی بعینی و ینقطع عن العالمین و هل
من ذی بصر لینظر فی امری و ما ورد علیّ من جنود الشّیاطین و هل من ذی قلب
رقیق طاهر لیبکی ببکائی و بما ورد علیّ من هؤلآء المستکبرین و هل من ذی علم
یصدّ سیوف الرّدّ و البغضآء عن وجه اللّه الملک العزیز الجمیل

\end{otherlanguage}

Hal min māʾin ṣāfin yuṭahhiru ḥawlī ʿan hāʾulāʾi al-muʿriḍīn, wa-hal min
sayfi lisānin yaftaḥu bi'l-bayāni wa-yadfaʿu ʿan fināʾī sharra
al-mushrikīn, wa-hal min dhī istiqāmatin yanṣurunī fī hādhihi al-ayyāmi
allatī aḥāṭatnī al-munkirīn, wa-hal min dhī inqiṭāʿin yaʿrifunī bi-ʿaynī
wa-yanqaṭiʿu ʿan al-ʿālamīn, wa-hal min dhī baṣarin li-yanẓura fī amrī
wa-mā warada ʿalayya min junūdi al-shayāṭīn, wa-hal min dhī qalbin
raqīqin ṭāhirin li-yabkiya bi-bukāʾī wa-bi-mā warada ʿalayya min
hāʾulāʾi al-mustakbirīn, wa-hal min dhī ʿilmin yaṣuddu suyūfa al-raddi
wa'l-baghḍāʾi ʿan wajhi Allāhi al-Maliki al-ʿAzīzi al-Jamīl.

``Is there any pure water to cleanse My surroundings of these people who
turn away? Is there any sword of the tongue to conquer through utterance
and ward off from My court the evil of those who associate partners with
God? Is there anyone steadfast to help Me in these days when the deniers
have surrounded Me? Is there anyone detached to know Me through My own
eye and sever himself from all the worlds? Is there anyone with vision
to examine My Cause and what has befallen Me from the hosts of devils?
Is there anyone with a tender, pure heart to weep at My weeping and at
what these arrogant ones have inflicted upon Me? Is there anyone with
knowledge to turn the swords of rejection and hatred away from the face
of God, the King, the Almighty, the Beautiful?''

\subsection{178}\label{section-1733}

\begin{otherlanguage}{persian}

ان یا قلم الأعلی قل تاللّه الحقّ انّی ما استنصرت من احد حین الّذی حارب معی
کلّ الملل و کان اللّه علی ذلک شهید و علیم و کلّما استنصر من احد و نأمر به
العباد هذا لم یکن الّا لبلوغهم الی ذروة الفضل و من دون ذلک انّ اللّه لغنیّ
حمید

\end{otherlanguage}

In yā Qalama al-Aʿlā, qul ta-llāhi al-Ḥaqq, innī mā istanṣartu min
aḥadin ḥīna alladhī ḥāraba maʿī kullu al-milal, wa-kāna Allāhu ʿalā
dhālika shahīd wa-ʿalīm. Wa-kullamā istanṣara min aḥadin wa-naʾmuru bihi
al-ʿibād, hādhā lam yakun illā li-bulūghihim ilā dhirwati al-faḍl,
wa-min dūni dhālika inna Allāha la-ghaniyyun ḥamīd.

O Most Exalted Pen! Say: ``By God, the True One, I sought help from no
one when all the nations made war upon Me, and God is a witness to this
and knows it well. Whenever He seeks help from anyone, and We enjoin
this upon Our servants, it is solely that they may attain the summit of
grace. Apart from this, God is indeed Self-Sufficient, Praiseworthy.''

\subsection{179}\label{section-1734}

\begin{otherlanguage}{persian}

باری رسائل متعدّده در ردّ سلطان احدیّه نوشته‌اند تاللّه بطل بذلک اعمالهم و ضلّ
سعیهم فی الدّنیا و الآخرة و هم لا یشعرون و آنچه را سبب و علّت اعراض ناس
دانسته‌اند باین عبد نسبت داده‌اند از جمله نسبت دنیا و این بر کلّ ذی شعور
واضح بوده که این عبد در مابین معرضین چون شمس واضح و لایح بود و قسم بجمال
قدم که در هیچ صبحی امید نبود که لیل ادراک شود و هم‌چنین هیچ لیلی گمان
مشاهدهٴ فجر نه چه که مشرکین بکمال جدّ در قطع سدرهٴ ربّ العالمین بوده تا
آنکه شمس قدرت ید اللّه فوق ایدیهم از افق جیب عظمت اشراق فرمود و مشرکین
مغلوب و جمال انّه لهو الغالب علی کلّ شیء بر عرش قدرت و اقتدار مستوی و
باعزاز تمام از عراق حرکت نمود

\end{otherlanguage}

Bārī rasāʾil-i mutaʿaddidih dar radd-i Sulṭān-i aḥadiyyih nivishtih-and.
Ta-llāhi baṭala bi-dhālika aʿmāluhum wa-ḍalla saʿyuhum fī al-dunyā
wa'l-ākhirati wa-hum lā yashʿurūn. Va ānchih rā sabab va ʿillat-i
iʿrāḍ-i nās dānistih-and bi-īn ʿabd nisbat dādih-and, az jumlih nisbat-i
dunyā. Va īn bar kull-i dhī shuʿūr vāḍiḥ būdih kih īn ʿabd dar mābayn-i
muʿriḍīn chūn shams vāḍiḥ va lāʾiḥ būd, va qasam bi-jamāl-i qidam kih
dar hīch ṣubḥī umīd nabūd kih layl idrāk shavad va hamchunīn hīch laylī
gumān-i mushāhidih-yi fajr nah, chih kih mushrikīn bi-kamāl-i jidd dar
qaṭʿ-i sidrih-yi Rabb-i ʿālamīn būdih, tā ānkih shams-i qudrat-i Yadu
Allāhi fawqa aydīhim az ufuq-i jayb-i ʿaẓamat ishrāq farmūd va mushrikīn
maghlūb va jamāl-i Innahu la-huwa al-ghālibu ʿalā kulli shayʾ bar
ʿarsh-i qudrat va iqtidār mustavī va bi-iʿzāz-i tamām az ʿIrāq ḥarakat
namūd.

In short, they have written numerous treatises refuting the Sovereign of
Oneness. By God, through this their deeds have come to nothing and their
efforts have gone astray in this world and the next, yet they do not
perceive it. They have attributed to this servant whatever they thought
would cause people to turn away, including worldliness. Yet it was clear
to every discerning person that this servant was as manifest and radiant
as the sun amid those who turned away. I swear by the Ancient Beauty
that no morning held any hope of living to see the night, nor did any
night bring any expectation of witnessing the dawn, for those who
associate partners with God were striving with all their might to cut
down the Lote-Tree of the Lord of the worlds. At last the sun of the
power of ``The hand of God is above their hands'' shone from the horizon
of the bosom of grandeur, and those who associate partners with God were
vanquished. The Beauty of ``He is indeed supreme over all things'' was
established upon the throne of power and might and departed from Iraq
with the utmost honor.

\subsection{180}\label{section-1735}

\begin{otherlanguage}{persian}

و این معلوم بوده که اراده آن بود که وحده عازم اسلامبول شوم چون جمعی جزع
و فزع لایحصی نمودند ببعضی اذن خروج داده شد و بعضی هم من غیر اذن حرکت
نمودند تا آنکه وارد ارض مشهور شدیم و ابداً بدولت علیّه مطلبی اظهار نشد و
هر نفسی که بین یدی حاضر شد ملاقات واقع والّا فلا چنانچه جمیع استماع نموده
و در نفس مدینۀ کبیر هم شنیده‌اند که ارادهٴ الهی چه بوده تا آنکه چند نفر
مستضعفین حفظاً لأنفسهم و نسائهم مضطرب شده و جزع نمودند لذا ترحّماً لأنفسهم
حکم بدا نازل

\end{otherlanguage}

Va īn maʿlūm būdih kih irādih ān būd kih vaḥdahū ʿāzim-i Islāmbūl
shavam. Chūn jamʿī jazaʿ va fazaʿ-i lā yuḥṣā namūdand bi-baʿḍī idhn-i
khurūj dādih shud va baʿḍī ham min ghayr-i idhn ḥarakat namūdand, tā
ānkih vārid-i arḍ-i mashhūr shudīm. Va abadan bi-dawlat-i ʿaliyyih
maṭlabī iẓhār nashud, va har nafsī kih bayna yaday ḥāḍir shud mulāqāt
vāqiʿ va-illā fa-lā, chunānchih jamīʿ istimāʿ namūdih va dar nafs-i
Madīnih-yi Kabīr ham shinīdih-and kih irādih-yi ilāhī chih būdih, tā
ānkih chand nafar mustaḍʿafīn ḥifẓan li-anfusihim wa-nisāʾihim muḍṭarib
shudih va jazaʿ namūdand. Lidhā taraḥḥuman li-anfusihim ḥukm-i badā
nāzil.

It was known that My intention was to set out alone for Istanbul. When a
number of people raised immeasurable cries of grief and distress, some
were given permission to leave, while others set out without permission,
until We arrived in the well-known land. No request whatsoever was made
to the imperial government. Whoever came into My presence was received;
otherwise, no meeting took place. Everyone has heard this, and in the
Great City itself they also heard what the divine intention had been,
until a few weak souls became anxious for their own safety and that of
their wives and began to lament. Thus, out of compassion for them, a
decree changing the former decree was revealed.

\subsection{181}\label{section-1736}

\begin{otherlanguage}{persian}

حال کدام یک از این امور متعلّق بدنیا بوده باری قیاس بنفس نموده‌اند کذلک
صفت لهم السنتهم الکذبة یقولون ما لا یشعرون واللّه الّذی لا اله الّا هو قلم
حیا مینماید از آنچه عامل شده‌اند ذکر نماید و اگر ذی بصری یافت شود از
رسائل حضرات که از قلم جهل و بغضا ظاهر شده مقام و شأن معرضین مشهود گردد

\end{otherlanguage}

Ḥāl kudām yak az īn umūr mutaʿalliq bi-dunyā būdih? Bārī qiyās bi-nafs
namūdih-and. Kadhālika ṣafat lahum alsinatuhum al-kadhibah, yaqūlūna mā
lā yashʿurūn. Wa-llāhi alladhī lā ilāha illā huwa, qalam ḥayā mīnamāyad
az ānchih ʿāmil shudih-and dhikr namāyad, va agar dhī baṣarī yāft shavad
az rasāʾil-i ḥaḍarāt kih az qalam-i jahl va baghḍā ẓāhir shudih maqām va
shaʾn-i muʿriḍīn mashhūd gardad.

Now, which of these matters had anything to do with the world? They
have, in short, judged Me by themselves. Thus have their lying tongues
described matters to them; they say what they do not comprehend. By God,
besides whom there is no other god, the Pen is ashamed to mention what
they have done. If anyone with discernment could be found, the standing
and character of those who turn away would become apparent from these
gentlemen's treatises, which have issued from the pen of ignorance and
hatred.

\subsection{182}\label{section-1737}

\begin{otherlanguage}{persian}

دو سنهٴ قبل شعری جناب آقا محمّد علی و جناب سلمان سؤال نموده اخوی جواب
نوشته بعد سیّد محمّد اصفهانی بین یدی حاضر و مذکور داشت میرزا علی ملتفت
معنی شعر نشده‌اند و مصلحت نیست این مکتوب به سلمان داده شود این بنده سکوت
نمود بعد با حاجی میرزا احمد نزد اخوی در بیت کلیم ذکر نمودند اخوی بعد از
کلمات لایغنیه بالأخره متعذّر شد که شاید سائل اصل شعر را غلط نوشته باشد

\end{otherlanguage}

Du sanih-yi qabl shiʿrī Janāb-i Āqā Muḥammad ʿAlī va Janāb-i Salmān
suʾāl namūdih, akhavī javāb nivishtih. Baʿd Siyyid Muḥammad-i Iṣfahānī
bayna yaday ḥāḍir va madhkūr dāsht Mīrzā ʿAlī multafit-i maʿnī-yi shiʿr
nashudih-and va maṣlaḥat nīst īn maktūb bi-Salmān dādih shavad. Īn
bandih sukūt namūd. Baʿd bā Ḥājī Mīrzā Aḥmad nazd-i akhavī dar bayt-i
Kalīm dhikr namūdand. Akhavī baʿd az kalimāt-i lā yughniyih bi'l-ākharah
mutaʿadhdhir shud kih shāyad sāʾil aṣl-i shiʿr rā ghalaṭ nivishtih
bāshad.

Two years ago, Āqā Muḥammad ʿAlī and Salmān asked about a verse of
poetry, and My brother wrote a reply. Siyyid Muḥammad-i-Iṣfahānī then
came into My presence and said, ``Mīrzā ʿAlī has not understood the
meaning of the verse, and it is not advisable to give this letter to
Salmān.'' This servant remained silent. Later, together with Ḥājī Mīrzā
Aḥmad, he raised the matter with My brother in Kalīm's house. After some
unhelpful remarks, My brother finally offered the excuse that perhaps
the questioner had copied the original verse incorrectly.

\subsection{183}\label{section-1738}

\begin{otherlanguage}{persian}

بعد شعر را ملاحظه نمودند صحیح نوشته بود و آن نوشته نزد سلمان موجود است و
جمیع ملاحظه نموده‌اند ابداً نفسی شاعر نشده چه که کلّ بحجبات وهمیّه محتجبند
حال ملاحظه نمائید تا مقدار معرضین معلوم شود و مع‌ذلک بمنبع و مصدر علم
ربّانیّه اعتراض مینمایند فوالّذی قامت ملکوت التّقدیر علی فنآء بابه که کلّ
آنچه از قلم جاری نموده‌اند بهمین نحو بوده و خواهد بود ولکن عمت الأبصار
احجاب الوهم و الهوی و لذا لا یلتفتون الّا الّذین جعل اللّه ابصارهم مطهّراً عن
الأحجاب و هم ببصر الحدید فی کلّ شیء ینظرون

\end{otherlanguage}

Baʿd shiʿr rā mulāḥiẓih namūdand, ṣaḥīḥ nivishtih būd, va ān nivishtih
nazd-i Salmān mawjūd ast va jamīʿ mulāḥiẓih namūdih-and. Abadan nafsī
shāʿir nashudih, chih kih kull bi-ḥujubāt-i vahmiyyih muḥtajiband. Ḥāl
mulāḥiẓih namāyīd tā miqdār-i muʿriḍīn maʿlūm shavad, va maʿa-dhālik
bi-manbaʿ va maṣdar-i ʿilm-i rabbāniyyih iʿtirāḍ mīnamāyand.
Fa-wa-alladhī qāmat malakūtu al-taqdīri ʿalā fināʾi bābih, kih kull-i
ānchih az qalam jārī namūdih-and bi-hamīn naḥv būdih va khvāhad būd.
Va-lākin ʿamat al-abṣāra aḥjābu al-wahmi wa'l-hawā, va lidhā lā
yaltafitūna illā alladhīna jaʿala Allāhu abṣārahum muṭahharan ʿan
al-aḥjāb, wa-hum bi-baṣari al-ḥadīdi fī kulli shayʾin yanẓurūn.

They then examined the verse: it had been copied correctly. That
document is still in Salmān's possession, and everyone has seen it. Yet
no one has taken any notice, for all are veiled by imaginary notions.
Now consider this, so that the worth of those who turn away may become
known. Despite this, they object to the fountainhead and source of
divine knowledge. By Him at whose threshold the kingdom of destiny
stands, everything they have set down with their pens has been, and will
be, of this same kind. But the veils of delusion and desire have blinded
their eyes, and thus none take heed except those whose sight God has
purified of veils; they look upon all things with keen vision.

\subsection{184}\label{section-1739}

\begin{otherlanguage}{persian}

ای عباد ملاحظه نمائید که ضرّ بمقامی رسیده که قلم اعلی باین کلمات مشغول
شده چه که ناس ضعیف و نابالغ مشاهده میشوند که مباد از کلمات جعلیّهٴ معرضین
از جمال مبین غافل شوند والّا اگر فی‌الجمله شعور و بصر ملاحظه میشد ابداً ذکر
معرضین و ما ظهر منهم تحریر نمیشد چه که الیوم کلّ مکلّفند که در نفس امر و
ما یظهر من عنده ناظر شوند و من صعد الی ذلک المقام یشهد بأنّه هو الظّاهر
الباهر المشرق المضیء فی وسط السّمآء یحکم ما یشآء و لا یسأل فیما اراد و
انّه لهو المقتدر المختار

\end{otherlanguage}

Ay ʿibād, mulāḥiẓih namāyīd kih ḍurr bi-maqāmī rasīdih kih Qalam-i Aʿlā
bi-īn kalimāt mashghūl shudih, chih kih nās ḍaʿīf va nābāligh mushāhidih
mīshavand, kih mabād az kalimāt-i jaʿliyyih-yi muʿriḍīn az Jamāl-i Mubīn
ghāfil shavand. Va-illā agar fī'l-jumlih shuʿūr va baṣar mulāḥiẓih
mīshud abadan dhikr-i muʿriḍīn va mā ẓahara minhum taḥrīr namīshud, chih
kih al-yawm kull mukallafand kih dar nafs-i amr va mā yaẓharu min ʿindih
nāẓir shavand. Wa-man ṣaʿida ilā dhālika al-maqāmi yashhadu bi-annahu
huwa al-Ẓāhiru al-Bāhiru al-Mushriqu al-Muḍīʾu fī wasaṭi al-samāʾ,
yaḥkumu mā yashāʾu wa-lā yusʾalu fīmā arāda wa-innahu la-huwa
al-Muqtadiru al-Mukhtār.

O servants! Consider how grievous the harm has become, that the Most
Exalted Pen has occupied itself with these words. People are seen to be
weak and immature, and it is feared that the fabricated words of those
who turn away may cause them to become heedless of the Manifest Beauty.
Otherwise, had even a little understanding and discernment been
observed, no mention of those who turn away or of what has proceeded
from them would ever have been written. For today everyone is obliged to
look to the Cause itself and to what appears from its Source. Whoever
ascends to that station will testify that He is the Manifest, the
Resplendent, the Radiant, the Luminous in the midst of heaven. He
ordains whatever He wills and is not questioned concerning what He
desires. He is indeed the All-Powerful, the Unconstrained.

\subsection{185}\label{section-1740}

\begin{otherlanguage}{persian}

و از جمله نسبت ظلم باین مظلوم داده‌اند حال در همین فقره تفکّر نمائید و
انصاف دهید جمیع اهل بیان باین عبد ملاقات نموده و بچشم ظاهر دیده هرگز بضرّ
نفسی راضی شده لا فوربّ العرش لو انتم تنصفون و قد کان نفس الرّحمن فی فم
الثّعبان عشرین من السّنین و مع ذلک ابداً اظهار نشد و حال هم اگر این مکاتیب
مجعوله دیده نمیشد ابداً اظهار نمیرفت و یشهد بذلک کلّ من یعرفنی و کان علی
بصیرة منیر

\end{otherlanguage}

Va az jumlih nisbat-i ẓulm bi-īn maẓlūm dādih-and. Ḥāl dar hamīn fiqrih
tafakkur namāyīd va inṣāf di hīd. Jamīʿ-i ahl-i Bayān bi-īn ʿabd mulāqāt
namūdih va bi-chashm-i ẓāhir dīdih, hargiz bi-ḍurr-i nafsī rāḍī shudih?
Lā fa-wa-Rabbi al-ʿarshi law antum tunṣifūn. Wa-qad kāna nafsu
al-Raḥmāni fī fami al-thuʿbāni ʿishrīna mina al-sinīn, va maʿa dhālik
abadan iẓhār nashud, va ḥāl ham agar īn makātīb-i majʿūlih dīdih
namīshud abadan iẓhār namīraft. Wa-yashhadu bi-dhālika kullu man
yaʿrifunī wa-kāna ʿalā baṣīratin munīr.

Among other things, they have charged this wronged One with oppression.
Now reflect on this very point and judge fairly. All the people of the
Bayān have met this servant and seen Him with their own eyes. Has He
ever consented to anyone being harmed? No, by the Lord of the throne, if
you judge fairly! The very person of the All-Merciful was in the
serpent's mouth for twenty years, yet this was never disclosed. Even
now, had these fabricated letters not been seen, it would never have
been disclosed. Everyone who knows Me and possesses clear insight bears
witness to this.

\subsection{186}\label{section-1741}

\begin{otherlanguage}{persian}

بیست سنه اخوی نزد این عبد بود و احدی از او مطّلع نه و مع قدرت بر او و
اطّلاع ما فی سرّه و قلبه کمال شفقت در حقّش معمول شد و یشهد بذلک کلّ الأشیآء
ان انتم تسمعون ولکن چون ناس را مقلّد و بی‌بصر یافته‌اند لذا باینگونه کلمات
اراده نموده‌اند این فئه را از بحر البحور منع نمایند

\end{otherlanguage}

Bīst sanih akhavī nazd-i īn ʿabd būd va aḥadī az ū muṭṭaliʿ nah, va maʿa
qudrat bar ū va iṭṭilāʿ-i mā fī sirrihi wa-qalbih kamāl-i shafaqat dar
ḥaqqash maʿmūl shud. Wa-yashhadu bi-dhālika kullu al-ashyāʾi in antum
tasmaʿūn. Va-lākin chūn nās rā muqallid va bī-baṣar yāftih-and lidhā
bi-īngūnih kalimāt irādih namūdih-and īn fiʾih rā az baḥr-i buḥūr manʿ
namāyand.

For twenty years My brother was with this servant, and no one knew of
him. Although I had power over him and knew what lay concealed within
him and in his heart, he was treated with the utmost compassion. All
things bear witness to this, if you would listen. But since they have
found people to be blind followers without discernment, they have
sought, through words of this kind, to keep this community from the
Ocean of oceans.

\subsection{187}\label{section-1742}

\begin{otherlanguage}{arabic}

و بر جمیع اهل بیان و غیره واضح و مبرهن که این عبد باحدی معاشرت در این
ارض ننموده و هم‌چنین در ورود مدینهٴ کبیر الّا من حضر بین یدینا و ما حرّک
لسانی الّا علی ما یصعدهم الی جبروت الأمر مع‌ذلک نوشته‌اند که نزد اهل سنّت ما
را رافضی نامیده‌اند فواللّه الّذی لا اله الّا هو که خود این نفوس معرضه نزد
اعجام نسبت دون ایمان بسدرهٴ رحمن داده‌اند و عند اهل سنّت و جماعت نسبت رفض
حال نفسی خالصاً للّه بیاید صدق و کذب را معلوم نماید و میزان حقّ و باطل را
بصدق و کذب همین قول قرار میدهیم اذاً ینادی لسان القدم من عرشه الأعظم و
یقول لعن اللّه اوّل ظالم ظلم نقطة الأولی فی ظهوره الأخری و کفر بآیاته و
حارب بنفسه و جادل بسلطانه و اشرک بذاته و الحد کلماته و اعرض عن جماله و
کان من المشرکین فی الواح اللّه من قلم القدس مکتوبا

\end{otherlanguage}

Va bar jamīʿ-i ahl-i Bayān va ghayrih vāḍiḥ va mubarhan kih īn ʿabd
bi-aḥadī muʿāshirat dar īn arḍ nanamūdih, va hamchunīn dar vurūd-i
Madīnih-yi Kabīr, illā man ḥaḍara bayna yadaynā, wa-mā ḥarraka lisānī
illā ʿalā mā yuṣʿiduhum ilā jabarūti al-amr. Maʿa-dhālik nivishtih-and
kih nazd-i ahl-i sunnat mā rā rāfiḍī nāmīdih-and. Fa-wa-llāhi alladhī lā
ilāha illā huwa, kih khud-i īn nufūs-i muʿriḍih nazd-i aʿjām nisbat-i
dūn-i īmān bi-sidrih-yi Raḥmān dādih-and va ʿinda ahl-i sunnat va
jamāʿat nisbat-i rafḍ. Ḥāl nafsī khāliṣan li-llāh biyāyad ṣidq va kidhb
rā maʿlūm namāyad, va mīzān-i ḥaqq va bāṭil rā bi-ṣidq va kidhb-i hamīn
qawl qarār mīdahīm. Idhan yunādī lisānu al-qidami min ʿarshihi al-aʿẓami
wa-yaqūl: laʿana Allāhu awwala ẓālimin ẓalama Nuqṭata al-Ūlā fī ẓuhūrihi
al-ukhrā wa-kafara bi-āyātihi wa-ḥāraba bi-nafsihi wa-jādala
bi-sulṭānihi wa-ashraka bi-dhātihi wa-alḥada kalimātihi wa-aʿraḍa ʿan
jamālihi wa-kāna mina al-mushrikīna fī alwāḥi Allāhi min qalami al-qudsi
maktūbā.

It is clear and evident to all the people of the Bayān and to others
that this servant has associated with no one in this land, nor upon
arriving in the Great City, except those who came into Our presence. My
tongue moved only to speak what would raise them to the dominion of the
Cause. Nevertheless, they have written, ``He has called us Rāfiḍīs
before the Sunnis.'' By God, besides whom there is no other god, these
very people who turn away have themselves accused the Lote-Tree of the
All-Merciful of unbelief before the Persians and of Rāfiḍism before the
Sunni community. Let someone now come, solely for God's sake, and
ascertain what is true and what is false. We make the truth or falsehood
of this very statement the criterion of truth and error. Then the Tongue
of Eternity calls out from His most great throne and says: ``May God
curse the first oppressor who wronged the Primal Point in His latter
manifestation, denied His verses, made war upon His person, disputed His
sovereignty, associated partners with His Essence, impugned His words,
turned away from His beauty, and was numbered among those who associate
partners with God, as inscribed by the Pen of holiness in the Tablets of
God.''

\subsection{188}\label{section-1743}

\begin{otherlanguage}{persian}

فواللّه در هر وقتی از اوقات که در ارض طاء و یا عراق و مواضع دیگر حرفی در
اخذ این طایفه مذکور میشد اخوی خود از عراق باطراف میرفتند چنانچه اکثری
شنیده و میدانند و این عبد وحده مابین اعدا بوده تا آنکه از قدرت ربّانی امر
مرتفع شد و جمیع از نفوس خود مطمئن شده از خلف ستر بیرون آمده اوّل کاری که
بآن مشغول شدند قطع سدرهٴ امر بوده و چون خود را خائب و خاسر مشاهده نمودند
عمّا ارادوا بعد باطراف نوشته و ارادهٴ خود را بارادة اللّه و مشیّته نسبت
داده با اینکه این عبد حکم قتل را که مابین این فئه شهرت تام داشت مرتفع
نمود و در اکثر الواح من قلم اللّه جاری شده و البتّه بنظر بعضی رسیده

\end{otherlanguage}

Fa-wa-llāh dar har vaqtī az awqāt ke dar arḍ-i Ṭāʾ va yā ʿIrāq va
mavāḍiʿ-i dīgar ḥarfī dar akhdh-i īn ṭāyife madhkūr mīshud akhavī khud
az ʿIrāq bi-aṭrāf mīraftand chunānchi aktharī shinīde va mīdānand va īn
ʿabd vaḥdahu mābayn-i aʿdā būde tā ānke az qudrat-i rabbānī amr murtafiʿ
shud va jamīʿ az nufūs-i khud muṭmaʾinn shude az khalf-i sitr bīrūn
āmade avval kārī ke bi-ān mashghūl shudand qaṭʿ-i sidri-yi amr būde va
chun khud rā khāʾib va khāsir mushāhide nimūdand ʿammā arādū baʿd
bi-aṭrāf nivishte va irādi-yi khud rā bi-irādati Allāh wa mashīʾatihi
nisbat dāde bā īnke īn ʿabd ḥukm-i qatl rā ke mābayn-i īn fiʾe shuhrat-i
tāmm dāsht murtafiʿ nimūd va dar akthar-i alvāḥ min qalami Allāh jārī
shude va albatte bi-naẓar-i baʿḍī risīde

By God, whenever there was any talk in the Land of Ṭāʾ, in Iraq, or
elsewhere of arresting this community, My brother himself would leave
Iraq for other places, as most have heard and know. This servant
remained alone among the enemies until, through divine power, the Cause
was exalted. Once all felt secure for their own lives, they emerged from
behind concealment, and the first task they undertook was to cut down
the Lote-Tree of the Cause. When they saw themselves disappointed and
defeated in what they had intended, they then wrote to various places,
attributing their own intention to the will and purpose of God. Yet this
servant abolished the injunction to kill, which had been widely known
among this community. This has flowed from the Pen of God in most of the
Tablets, and has certainly come to the attention of some.

\subsection{189}\label{section-1744}

\begin{otherlanguage}{persian}

از جمله سؤال شده که مقصود از نصری که در آیات بدیعهٴ الهیّه نازل شده چیست
و ما نزّل عن جهة العرش فی جوابه این است که سواد آن بعینه ذکر میشود و هو
هذا

\end{otherlanguage}

Az jumle suʾāl shude ke maqṣūd az naṣrī ke dar āyāt-i badīʿi-yi ilāhiyye
nāzil shude chīst va mā nuzzila ʿan jihati al-ʿarsh fī jawābihi īn ast
ke savād-i ān bi-ʿaynihi dhikr mīshavad wa huwa hādhā

Among other things, a question was asked about what is meant by the
assistance mentioned in the wondrous divine verses. What was sent down
from the Throne in answer to it is reproduced here word for word, as
follows:

\subsection{190}\label{section-1745}

\begin{otherlanguage}{persian}

مقصود از نصرتی که در الواح منیعه ذکر شده معلوم احبّآء اللّه بوده که حقّ جلّ
ذکره مقدّس است از دنیا و آنچه در او است و مقصود از نصر این نبوده که احدی
بنفسی محاربه و یا مجادله نماید سلطان جلال جمیع ارض از برّ و بحر آن را بید
ملوک گذاشته و ایشانند مظاهر قدرت الهی علی قدر مراتبهم و اگر در ظلّ حقّ
وارد شوند از حقّ محسوب و من دون ذلک انّ ربّک لعلیم خبیر

\end{otherlanguage}

Maqṣūd az nuṣratī ke dar alvāḥ-i manīʿe dhikr shude maʿlūm-i aḥibbāʾ-i
Allāh būde ke Ḥaqq jalla dhikruhu muqaddas ast az dunyā va ānche dar ū
ast va maqṣūd az naṣr īn nabūde ke aḥadī bi-nafsī muḥāribe va yā
mujādile namāyad Sulṭān-i Jalāl jamīʿ-i arḍ az barr va baḥr-i ān rā
bi-yad-i mulūk gudhāshte va īshānand maẓāhir-i qudrat-i ilāhī ʿalā qadri
marātibihim va agar dar ẓill-i Ḥaqq vārid shavand az Ḥaqq maḥsūb wa min
dūni dhālika inna rabbaka la-ʿalīmun khabīr

``The friends of God know what is meant by the assistance mentioned in
the exalted Tablets: God, exalted be His mention, is sanctified above
the world and all that is in it. Assistance has not meant that anyone
should wage war or contend with another. The Sovereign of Glory has
placed the whole earth, both land and sea, in the hands of kings, and
they are the manifestations of divine power according to their
respective ranks. If they enter beneath the shadow of God, they are
accounted as belonging to Him; otherwise, your Lord is indeed
all-knowing, all-informed.

\subsection{191}\label{section-1746}

\begin{otherlanguage}{persian}

و آنچه حقّ جلّ ذکره از برای خود خواسته قلوب عباد او است که کنایز وحی
صمدانیّه و خزاین حبّ الهیّه‌اند و لم‌یزل ارادهٴ سلطان لایزال این بوده که قلوب
را از اشارات دنیا و ما فیها و علیها طاهر نماید تا قابل شوند از برای
انوار تجلّیات ملیک اسماء و صفات

\end{otherlanguage}

Va ānche Ḥaqq jalla dhikruhu az barā-yi khud khāste qulūb-i ʿibād-i ū
ast ke kanāyiz-i vaḥy-i ṣamadāniyye va khazāyin-i ḥubb-i ilāhiyye-and va
lam-yazal irādi-yi Sulṭān-i lā-yazāl īn būde ke qulūb rā az ishārāt-i
dunyā va mā fīhā va ʿalayhā ṭāhir namāyad tā qābil shavand az barā-yi
anvār-i tajalliyāt-i Malīk-i asmāʾ va ṣifāt

What God, exalted be His mention, has desired for Himself is the hearts
of His servants, which are the treasuries of divine revelation and the
repositories of divine love. It has always been the will of the
everlasting Sovereign to purify hearts from the allusions of the world
and all that is in it and upon it, so that they may become receptive to
the lights of the revelations of the King of names and attributes.

\subsection{192}\label{section-1747}

\begin{otherlanguage}{persian}

پس باید در مدینهٴ قلب بیگانه راه نیابد تا دوست بمقرّ خود شتابد یعنی تجلّی
جمالش نه ذات و نفس او چه که لازال مقدّس از صعود و نزول بوده و خواهد بود
پس نصرت امر اللّه الیوم اعتراض باحدی و مجادله بنفسی و محاربه مع شیئی
نبوده و نخواهد بود بلکه محبوب آن است که مداین قلوب بسیف لسان و حکمت بیان
مفتوح شود نه بسیف حدید

\end{otherlanguage}

Pas bāyad dar madīni-yi qalb bīgāne rāh nayābad tā Dūst bi-maqarr-i khud
shitābad yaʿnī tajallī-yi jamālash na dhāt va nafs-i ū chi ke lā-zāl
muqaddas az ṣuʿūd va nuzūl būde va khāhad būd pas nuṣrat-i amr-i Allāh
al-yawm iʿtirāḍ bi-aḥadī va mujādile bi-nafsī va muḥāribe maʿa shayʾī
nabūde va nakhāhad būd balke maḥbūb ān ast ke madāyin-i qulūb bi-sayf-i
lisān va ḥikmat-i bayān maftūḥ shavad na bi-sayf-i ḥadīd

Thus no stranger must find a way into the city of the heart, so that the
Friend may hasten to His own abode---that is, the revelation of His
beauty, not His Essence and His very Self, for He has always been and
will forever remain sanctified above ascent and descent. Assisting the
Cause of God today has therefore never consisted, nor will it ever
consist, in opposing anyone, disputing with anyone, or waging war
against anything. Rather, what is desired is that the cities of hearts
be conquered with the sword of the tongue and the wisdom of utterance,
not with a sword of iron.

\subsection{193}\label{section-1748}

\begin{otherlanguage}{persian}

پس هر نفسی که ارادهٴ نصر الهی نماید باید اوّل بسیف معانی و بیان مدینهٴ
قلب خود را تصرّف نماید و از جمیع ما سوی اللّه او را مطهّر سازد و بعد بمداین
قلوب توجّه کند این است نصرت امر اللّه که الیوم از مشرق اصبع ملیک اسماء
اشراق فرموده ابداً فساد محبوب نبوده و نیست و آنچه از قبل شده من غیر اذن
اللّه بوده

\end{otherlanguage}

Pas har nafsī ke irādi-yi naṣr-i ilāhī namāyad bāyad avval bi-sayf-i
maʿānī va bayān madīni-yi qalb-i khud rā taṣarruf namāyad va az jamīʿ-i
mā sivā Allāh ū rā muṭahhar sāzad va baʿd bi-madāyin-i qulūb tavajjuh
kunad īn ast nuṣrat-i amr-i Allāh ke al-yawm az mashriq-i iṣbaʿ-i
Malīk-i asmāʾ ishrāq farmūde abadan fasād maḥbūb nabūde va nīst va ānche
az qabl shude min ghayri idhni Allāh būde

Thus whoever wishes to assist God must first conquer the city of his own
heart with the sword of inner meanings and utterance, purify it from all
besides God, and then turn toward the cities of hearts. This is the
assistance to the Cause of God that has shone forth today from the
dawning-place of the finger of the King of names. Disorder has never
been, and is not, acceptable, and whatever occurred in the past was
without God's permission.

\subsection{194}\label{section-1749}

\begin{otherlanguage}{persian}

باری الیوم باید احبّای الهی بشأنی در مابین عباد ظاهر شوند که جمیع را
بافعال خود برضوان ذو الجلال هدایت نمایند قسم بآفتاب صبح عزّ تقدیس که ابداً
حقّ و احبّای او که منسوب باویند ناظر بارض و اموال فانیه در او نبوده و
نخواهند بود چه اگر مقصود تصرّف در ارض بود البتّه قادر و مقتدر بوده و
بکلمه‌ئی جمیع عالم را تصرّف میفرمود

\end{otherlanguage}

Bārī al-yawm bāyad aḥibbā-yi ilāhī bi-shaʾnī dar mābayn-i ʿibād ẓāhir
shavand ke jamīʿ rā bi-afʿāl-i khud bi-Riḍvān-i Dhu al-Jalāl hidāyat
namāyand qasam bi-āftāb-i ṣubḥ-i ʿizz-i taqdīs ke abadan Ḥaqq va
aḥibbā-yi ū ke mansūb bi-ūyand nāẓir bi-arḍ va amvāl-i fāniye dar ū
nabūde va nakhāhand būd chi agar maqṣūd taṣarruf dar arḍ būd albatte
qādir va muqtadir būde va bi-kalimiʾī jamīʿ-i ʿālam rā taṣarruf mīfarmūd

In short, today the friends of God must appear among His servants in
such a manner that, through their deeds, they guide everyone to the
Paradise of the Lord of Glory. I swear by the sun of the morning of
sanctified glory that neither God nor His friends, who are associated
with Him, have ever set their sights, or will ever set them, on the
earth and its perishable possessions. For had the purpose been to take
possession of the earth, He was assuredly powerful and mighty, and would
have taken possession of the entire world with a single word.

\subsection{195}\label{section-1750}

\begin{otherlanguage}{persian}

ولکن سلطنت بسلاطین عنایت فرموده و حکمت بمتفرّسین و عرفان بعارفین و حبّ
قلوب عباد خود را مخصوص خود مقرّر داشته و بدون آن ناظر نبوده و نخواهد بود
و اینهم نظر بعنایت کبری است که شاید نفوس فانیه از شئونات ترابیّه طاهر و
مقدّس شوند و بمقام باقیه که رضوان عزّ احدیّه است وارد گردند

\end{otherlanguage}

Va-lākin salṭanat bi-salāṭīn ʿināyat farmūde va ḥikmat bi-mutafarrisīn
va ʿirfān bi-ʿārifīn va ḥubb-i qulūb-i ʿibād-i khud rā makhṣūṣ-i khud
muqarrar dāshte va bi-dūn-i ān nāẓir nabūde va nakhāhad būd va īn ham
naẓar bi-ʿināyat-i kubrā ast ke shāyad nufūs-i fāniye az shuʾūnāt-i
turābiyye ṭāhir va muqaddas shavand va bi-maqām-i bāqiye ke Riḍvān-i
ʿizz-i aḥadiyye ast vārid gardand

But He has bestowed sovereignty upon kings, wisdom upon the discerning,
and spiritual knowledge upon those who recognize Him. He has reserved
for Himself the love of His servants' hearts, and has never sought, nor
will He ever seek, anything else. This, too, is out of His supreme
grace, that perhaps mortal souls may be cleansed and sanctified from
earthly concerns and enter the abiding station, which is the Paradise of
the glory of divine unity.

\subsection{196}\label{section-1751}

\begin{otherlanguage}{persian}

والّا آن سلطان قدم بنفسه لنفسه مستغنی از کلّ بوده نه از حبّ ممکنات نفعی باو
راجع و نه از ضرّشان ضرّی باو واقع کلّ از امکنهٴ ترابیّه ظاهر و باو راجع
خواهند شد و آن جمال قدم متوحّداً متفرّداً بر مقرّ خود که مقدّس از مکان و زمان
و ذکر و اشاره و دلاله و وصف و علوّ و دنوّ بوده مستقرّ و لا یعلم ذلک الّا کلّ
ذی فطن بصیر انتهی

\end{otherlanguage}

Va-illā ān Sulṭān-i Qidam bi-nafsihi li-nafsihi mustaghnī az kull būde
na az ḥubb-i mumkināt nafʿī bi-ū rājiʿ va na az ḍurrishān ḍurrī bi-ū
vāqiʿ kull az amkini-yi turābiyye ẓāhir va bi-ū rājiʿ khāhand shud va ān
Jamāl-i Qidam mutavaḥḥidan mutafarridan bar maqarr-i khud ke muqaddas az
makān va zamān va dhikr va ishāre va dilāle va vaṣf va ʿuluvv va dunuvv
būde mustaqarr wa lā yaʿlamu dhālika illā kullu dhī fiṭanin baṣīr intahā

Otherwise, that Sovereign of Eternity is, in Himself and for Himself,
independent of all. No benefit accrues to Him from the love of created
beings, nor does any harm befall Him from their harm. All emerge from
earthly places and will return to Him, while that Beauty of Eternity
remains, alone and unique, established upon His own seat, which is
sanctified above place and time, mention and allusion, indication and
description, height and nearness. None knows this except those endowed
with discernment and insight.'' End.

\subsection{197}\label{section-1752}

\begin{otherlanguage}{persian}

مع‌ذلک آنچه را خود عامل و فاعل بوده‌اند بحقّ تعالی شأنه نسبت داده‌اند مع
آنکه بر کلّ واضح است که حقّ یفعل ما یشآء بوده و بر هر چه حکم فرماید قادر و
لیس لأحد ان یعترض علیه و مع‌ذلک فوالّذی نفسی بیده که اگر راضی بضرّ نفسی
بوده و یشهد بذلک قلم اللّه ان انتم تشهدون و از جمله در مکاتیب خود ذکر
نموده‌اند که این عبد درس خوانده و این کلمات منزله از علم اکتسابی حاصل شده
چنانچه در احیان هر ظهور اینگونه کلمات نالایقه مذکور شده چنانچه نسبت
بظهور قبلم همین نسبت را داده‌اند و من قبله به محمّد رسول اللّه انّما یعلّمه
بشر لسان الّذی یلحدون الیه اعجمیّ و هذا لسان عربیّ مبین حال نفسی در کلّ
ایران و عراق تفحّص نماید تا صدق و کذب معلوم شود فویل لهم و بما هم به
یتکلّمون

\end{otherlanguage}

Maʿa dhālik ānche rā khud ʿāmil va fāʿil būde-and bi-Ḥaqq taʿālā
shaʾnuhu nisbat dāde-and maʿa ānke bar kull vāḍiḥ ast ke Ḥaqq yafʿalu mā
yashāʾ būde va bar har che ḥukm farmāyad qādir wa laysa li-aḥadin an
yaʿtariḍa ʿalayhi va maʿa dhālik fa-wa-alladhī nafsī bi-yadihi ke agar
rāḍī bi-ḍurr-i nafsī būde wa yashhadu bi-dhālika qalamu Allāh in antum
tashhadūn va az jumle dar makātīb-i khud dhikr nimūde-and ke īn ʿabd
dars khānde va īn kalimāt-i munzale az ʿilm-i iktisābī ḥāṣil shude
chunānchi dar aḥyān-i har ẓuhūr īngūne kalimāt-i nālāyiqe madhkūr shude
chunānchi nisbat bi-ẓuhūr-i qablam hamīn nisbat rā dāde-and va min
qablihi bi-Muḥammad Rasūl-i Allāh innamā yuʿallimuhu bashar lisānu
alladhī yulḥidūna ilayhi aʿjamiyyun wa hādhā lisānun ʿarabiyyun mubīn
ḥāl nafsī dar kull-i Īrān va ʿIrāq tafaḥḥuṣ namāyad tā ṣidq va kidhb
maʿlūm shavad fa-waylun lahum wa bi-mā hum bihi yatakallamūn

Nevertheless, they have attributed to God, exalted be His station, deeds
of which they themselves were the agents and perpetrators. Yet it is
clear to all that God does whatever He wills and has power to carry out
whatever He decrees, and no one has the right to object to Him. Even so,
I swear by Him in whose hand is My soul that I have never consented to
harm anyone, and the Pen of God bears witness to this, if you would bear
witness. Among other things, they have stated in their letters that this
servant has studied and that these revealed words are the product of
acquired learning. Such unworthy statements have been made at the time
of every revelation. They made this very allegation concerning My
previous Manifestation, and before Him concerning Muḥammad, the
Messenger of God: ``A human being merely teaches Him.'' ``The tongue of
the one to whom they allude is foreign, while this is a clear Arabic
tongue.'' Now let someone make inquiries throughout Iran and Iraq, so
that truth and falsehood may become known. Woe to them for what they
say!

\subsection{198}\label{section-1753}

\begin{otherlanguage}{persian}

ای بی‌خردان منبع این علوم ذات قدم بوده و معین این حکم نفس اللّه الأعظم لو
انتم من المستشعرین و اگر بسمع فطرت استماع نمائید آنچه را در حین ظهور
نقطه جلّ و عزّ علما و جهّال عصر اعتراضاً علی اللّه بآن تکلّم مینمودند حال هم
همان کلمات را استماع مینمائید فواللّه یا قوم انّی لأکون متحیّراً فی امری و
ما ورد علیّ فیا لیت ما ولدت من امّی و ان ولدت ما رضعت و ما صرت کبیراً ولکن
ظهر ما ظهر و نزّلت جنود الوحی و قضی الأمر من لدی اللّه الواحد القهّار یا
قوم اسمعوا قولی ثمّ اخرقوا الأحجاب ولو یکون مرایا السّموات و الأرضین ایّاکم
ان لا یمنعکم شیء عن نفحات تلک الأیّام تاللّه لن یعادلها کنایز الابداع و لا
خزاین الاختراع و کان اللّه علی ذلک شهید و علیم

\end{otherlanguage}

Ay bī-khiradān manbaʿ-i īn ʿulūm Dhāt-i Qidam būde va maʿīn-i īn ḥikam
Nafs-i Allāh al-Aʿẓam law antum mina al-mustashʿirīn va agar bi-samʿ-i
fiṭrat istimāʿ namāyīd ānche rā dar ḥīn-i ẓuhūr-i Nuqṭe jalla va ʿazza
ʿulamā va juhhāl-i ʿaṣr iʿtirāḍan ʿalā Allāh bi-ān takallum mīnimūdand
ḥāl ham hamān kalimāt rā istimāʿ mīnamāyīd fa-wa-llāhi yā qawm innī
la-akūnu mutaḥayyiran fī amrī wa mā warada ʿalayya fa-yā layta mā
wulidtu min ummī wa in wulidtu mā raḍiʿtu wa mā ṣirtu kabīran wa-lākin
ẓahara mā ẓahara wa nuzzilat junūdu al-waḥy wa quḍiya al-amru min ladā
Allāhi al-Wāḥidi al-Qahhār yā qawmi ismaʿū qawlī thumma ikhriqū al-aḥjāb
wa-law yakūnu marāyā al-samāwāti wa al-araḍīn iyyākum an lā yamnaʿakum
shayʾun ʿan nafaḥāti tilka al-ayyām ta-llāhi lan yuʿādilahā kanāyizu
al-ibdāʿ wa lā khazāyinu al-ikhtirāʿ wa kāna Allāhu ʿalā dhālika shahīd
wa ʿalīm

You who lack understanding! The source of this knowledge is the eternal
Essence, and the wellspring of these words of wisdom is the most great
Self of God, if only you perceived. If you listen with the ear of your
innate nature, you will hear now the very same words that the learned
and the ignorant of the age uttered in opposition to God at the
appearance of the Point, exalted and glorified be He. By God, O people,
I am bewildered by My plight and by what has befallen Me. If only My
mother had not given birth to Me; or, if I had been born, had I never
been nursed or grown to adulthood! But what has appeared has appeared,
the hosts of revelation have been sent down, and the matter has been
decreed by God, the One, the All-Subduing. O people, listen to My words,
then tear through the veils, even if they should be the mirrors of the
heavens and the earths. Beware lest anything keep you from the breezes
of those days. By God, neither the treasures of creation nor the
repositories of origination can equal them. God is witness to this and
knows it well.

\subsection{199}\label{section-1754}

\begin{otherlanguage}{persian}

قسم بآفتاب فلک باقی که اگر امر بدست این عبد بود یک آیه مابین این قوم
تلاوت نمیشد چه که ناس بالمرّه از فطرت اصلیّهٴ الهیّه منحرف شده‌اند ولکنّ اللّه
قضی ما اراد و امضی ما شآء لا رادّ لمشیّته و لا مردّ لقضائه و انّه بکلّ شیء
حکیم و اگر نه این بود که اهل این مملکت مضطرب ملاحظه میشد هرآینه اذن داده
میشد که کلّ بین یدی اللّه حاضر شوند لیشهدوا ما لا شهد احد من العالمین

\end{otherlanguage}

Qasam bi-āftāb-i falak-i bāqī ke agar amr bi-dast-i īn ʿabd būd yak āye
mābayn-i īn qawm tilāvat namīshud chi ke nās bi-al-marre az fiṭrat-i
aṣliyye-yi ilāhiyye munḥarif shude-and wa-lākinna Allāha qaḍā mā arāda
wa amḍā mā shāʾ lā rādda li-mashīʾatihi wa lā maradda li-qaḍāʾihi wa
innahu bi-kulli shayʾin ḥakīm va agar na īn būd ke ahl-i īn mamlakat
muḍṭarib mulāḥiẓe mīshud har-āyine idhn dāde mīshud ke kull bayna yaday
Allāh ḥāḍir shavand li-yashhadū mā lā shahida aḥadun mina al-ʿālamīn

I swear by the sun of the everlasting heaven that, had the matter been
in this servant's hands, not a single verse would have been recited
among this people, for humanity has wholly departed from its original,
divinely bestowed nature. But God has decreed what He desired and
carried out what He willed. None can repel His will or reverse His
decree, and He is wise in all things. Were it not that the people of
this country were seen to be disturbed, permission would certainly have
been given for everyone to appear in the presence of God, so that they
might witness what no one in all the worlds has witnessed.

\subsection{200}\label{section-1755}

\begin{otherlanguage}{persian}

و از جمله نوشته که مقام شمسیّت در ظلّ مرآت واقع میشود و این رتبه فوق
مقامات است و هم‌چنین از این مرآت شموس مشرق سبحان‌اللّه چه مقدار ناس را جاهل
فرض نموده مع آنکه بسی واضح و مشهود است که خلق مرآت و وجود آن بجهت آن
بوده که حکایت از اشباح و امثال نماید و بر هر ذی ‌بصری واضح است که ظهور
عکوس امثال و اشباح هم در او لنفسه بنفسه تحقّق نداشته و ندارد و این فقره
بر هر اکمهی ظاهر است تا چه رسد بذی بصر مع‌ذلک باین عباراتی که فواللّه الّذی
لا اله الّا هو هیچ جاهلی بآن تکلّم نمینماید در رسائل خود نوشته و اراده
نموده باین کلمات ناس را از موجد و محقّق اسماء و صفات محروم نماید و این
نیست مگر آنکه عباد را در بحور وهم و هوی مستغرق دیده و گمان نموده که آنچه
مذکور دارد مسلّماً مقبول است والّا باین کلمات که کذب آن اظهر از شمس در وسط
السّمآء است مشغول نمیشد

\end{otherlanguage}

Va az jumle nivishte ke maqām-i shamsiyyat dar ẓill-i mirʾāt vāqiʿ
mīshavad va īn rutbe fawq-i maqāmāt ast va ham-chunīn az īn mirʾāt
shumūs mushriq subḥāna Allāh che miqdār nās rā jāhil farḍ nimūde maʿa
ānke basī vāḍiḥ va mashhūd ast ke khalq-i mirʾāt va vujūd-i ān
bi-jihat-i ān būde ke ḥikāyat az ashbāḥ va amthāl namāyad va bar har dhī
baṣarī vāḍiḥ ast ke ẓuhūr-i ʿukūs-i amthāl va ashbāḥ ham dar ū
li-nafsihi bi-nafsihi taḥaqquq nadāshte va nadārad va īn fiqre bar har
akmahī ẓāhir ast tā che rasad bi-dhī baṣar maʿa dhālik bi-īn ʿibārātī ke
fa-wa-llāhi alladhī lā ilāha illā huwa hīch jāhilī bi-ān takallum
namīnamāyad dar rasāʾil-i khud nivishte va irāde nimūde bi-īn kalimāt
nās rā az Mūjid va Muḥaqqiq-i asmāʾ va ṣifāt maḥrūm namāyad va īn nīst
magar ānke ʿibād rā dar buḥūr-i vahm va havā mustaghriq dīde va gumān
nimūde ke ānche madhkūr dārad musallaman maqbūl ast va-illā bi-īn
kalimāt ke kidhb-i ān aẓhar az shams dar vasaṭ-i al-samāʾ ast mashghūl
namīshud

Among other things, he has written that the station of being a sun falls
under the shadow of the mirror, that this rank is above all stations,
and that suns likewise shine forth from this mirror. Glorified be God!
How ignorant he has assumed people to be! Yet it is perfectly clear and
evident that the mirror was created and exists in order to reflect forms
and likenesses. It is clear to everyone with sight that even the
appearance within it of the reflections of likenesses and forms has
never come about, nor does it come about, through the mirror itself and
for its own sake. This point is evident to anyone born blind, let alone
to someone with sight. Nevertheless, he has written in his epistles such
statements as, by God, besides whom there is no other god, no ignorant
person would utter. Through these words he has intended to deprive
people of the One who brings names and attributes into existence and
gives them reality. This is only because he has seen God's servants
submerged in the seas of delusion and desire, and has supposed that
whatever he says will assuredly be accepted. Otherwise, he would not
have occupied himself with these words, whose falsehood is more evident
than the sun in the midst of heaven.

\subsection{201}\label{section-1756}

\begin{otherlanguage}{persian}

و متمسّک باین شده که نقطهٴ بیان روح ما سواه فداه خود را در یک مقام مرآت
اللّه نامیده و لذا مرآتیّت فوق رتبهٴ شمسیّت است فوالّذی انطقنی بالحقّ و
اشرقنی عن افق القدس علی العالمین که اگر نفسی در این یک فقره تفکّر نماید
بر جمیع مطالب این همج رعاع مطّلع میشود اوّلاً اینکه آن ذات قدم بکلّ اسماء
خود را موسوم فرموده‌اند و بکلّ صفات موصوف چه از اعلی مراتب اسماء و چه از
ادنی مراتب آن چنانچه در ابتدای احسن‌القصص میفرمایند قوله جلّ کبریائه الحمد
للّه الّذی انزل الکتاب علی عبده لیکون للعالمین سراجاً وهّاجاً حال میتواند کسی
قائل شود که رتبهٴ سراج فوق شمس بوده لا واللّه نمیگوید این سخنان را مگر
نفسی که بر هوی تکلّم نماید و بخواهد در قلوب صافیهٴ منیره شبهه القا نماید
فبئس ما هم یفعلون

\end{otherlanguage}

Va mutamassik bi-īn shude ke Nuqṭi-yi Bayān rūḥ-i mā sivāhu fidāhu khud
rā dar yak maqām Mirʾāt-i Allāh nāmīde va li-dhā mirʾātiyyat fawq-i
rutbi-yi shamsiyyat ast fa-wa-alladhī anṭaqanī bi-al-ḥaqqi wa ashraqanī
ʿan ufuqi al-qudsi ʿalā al-ʿālamīn ke agar nafsī dar īn yak fiqre
tafakkur namāyad bar jamīʿ-i maṭālib-i īn hamaj-i raʿāʿ muṭṭaliʿ
mīshavad avvalan īnke ān Dhāt-i Qidam bi-kull-i asmāʾ khud rā mawsūm
farmūde-and va bi-kull-i ṣifāt mawṣūf che az aʿlā marātib-i asmāʾ va che
az adnā marātib-i ān chunānchi dar ibtidā-yi Aḥsan al-Qaṣaṣ mīfarmāyand
qawluhu jalla kibriyāʾuhu al-ḥamdu li-llāhi alladhī anzala al-kitāba
ʿalā ʿabdihi li-yakūna li-al-ʿālamīna sirājan wahhājan ḥāl mītavānad
kasī qāʾil shavad ke rutbi-yi sirāj fawq-i shams būde lā wa-llāh
namīgūyad īn sukhanān rā magar nafsī ke bar havā takallum namāyad va
bikhāhad dar qulūb-i ṣāfiye-yi munīre shubhe ilqā namāyad fa-biʾsa mā
hum yafʿalūn

He has also seized upon the fact that the Point of the Bayān---may the
souls of all besides Him be sacrificed for Him---called Himself the
Mirror of God in one passage, and has therefore claimed that the station
of being a mirror is above that of being a sun. By Him who has caused Me
to speak the truth and made Me shine from the horizon of holiness upon
all the worlds, if anyone reflects on this one point, he will understand
all the claims of this rabble. First, that eternal Essence has
designated Himself by every name and described Himself by every
attribute, from the highest degrees of names to the lowest. Thus, at the
beginning of the Aḥsan al-Qaṣaṣ, He says---these are His words, exalted
be His majesty---``Praise be to God, who has sent down the Book to His
servant, that He may be a brightly shining lamp to all the worlds.'' Can
anyone now claim that the rank of a lamp is above that of the sun? No,
by God! No one utters such words except someone who speaks from desire
and wishes to cast doubt into pure and radiant hearts. How evil is what
they do!

\subsection{202}\label{section-1757}

\begin{otherlanguage}{arabic}

باری بیقین دانسته که در آن ساحت قدس جمیع اسماء در صقع واحدند چه که آن
کینونت قدم در حینی که میفرماید انا ربّ الأرباب و انا الأرض و التّراب در آن
منظر اکبر این هر دو در یک رتبه بوده و خواهند بود و این بیان ارقّ الطف
اصفی را قلوب مکدّره ادراک نخواهند نمود چنانچه آن جمال قدم میفرماید قوله
عزّ ذکره قل تنزّلت حتّی قلت انّنی انا ذرّ و دون ذرّ مثل ما انّنی انا قلت ربّ و
مربّب کلّ ذی ربّ لأستغفرنّ اللّه عن کلتیهما و انّی الی اللّه ربّی لمن الرّاجعین
انتهی

\end{otherlanguage}

Bārī bi-yaqīn dāniste ke dar ān sāḥat-i quds jamīʿ-i asmāʾ dar ṣuqʿ-i
vāḥidand, chi ke ān kaynūnat-i qidam dar ḥīnī ke mīfarmāyad anā rabbu
al-arbāb va anā al-arḍu wa-al-turāb, dar ān manẓar-i akbar īn har du dar
yik rutbe būde va khvāhand būd. Va īn bayān-i araqq-i alṭaf-i aṣfā rā
qulūb-i mukaddare idrāk nakhvāhand namūd. Chinānchi ān Jamāl-i Qidam
mīfarmāyad, qawluhu ʿazza dhikruhu: qul tanazzaltu ḥattā qultu innanī
anā dharrun wa-dūna dharrin mithla mā innanī anā qultu rabbun
wa-murabbibu kulli dhī rabbin. La-astaghfiranna Allāha ʿan kiltayhimā
wa-innī ilā Allāhi rabbī la-mina al-rājiʿīn. Intahā.

In any case, know with certainty that in that court of holiness all
names occupy a single plane. For when that eternal Being says, ``I am
the Lord of lords'' and ``I am the earth and the dust,'' both have held,
and will continue to hold, the same rank in that most great vista.
Clouded hearts will not comprehend this most subtle, delicate, and pure
utterance. Thus the Ancient Beauty says---exalted be His
mention---``Say: I descended until I said, `I am a mote, and less than a
mote,' just as I said, `I am the Lord and the nurturer of all that has a
lord.' I will surely ask God's forgiveness for both, and I am indeed
among those who return to God, my Lord.'' Thus ends the quotation.

\subsection{203}\label{section-1758}

\begin{otherlanguage}{persian}

حال ای سالکان بوادی جهل و ضلال که کأس زقوم نوشیده‌اید و سلسبیل قدس مختوم
دانسته‌اید در نار هوی محترق شده‌اید و از اهل رفرف اعلی خود را شمرده‌اید زهی
حسرت که یک آن در حکمت بالغهٴ رحمن تفکّر ننموده‌اید و در یک حین در امر مالک
یوم الدّین ببصر منیر توجّه ننموده‌اید کلمات رحمن را که بکمال انبساط و احاطه
من‌دون حدود نازل شده بحجبات نفس خود چنان محجوب نموده‌اید که السن عالمین از
ذکرش عاجزند

\end{otherlanguage}

Ḥāl ay sālikān bi-vādī-yi jahl va ḍalāl ke kaʾs-i zaqqūm nūshīde-īd va
salsabīl-i quds makhtūm dāniste-īd, dar nār-i havā muḥtariq shude-īd va
az ahl-i rafraf-i aʿlā khud rā shumurde-īd. Zihī ḥasrat ke yik ān dar
ḥikmat-i bālighe-yi Raḥmān tafakkur nanamūde-īd va dar yik ḥīn dar amr-i
Mālik-i Yawm al-Dīn bi-baṣar-i munīr tavajjuh nanamūde-īd. Kalimāt-i
Raḥmān rā ke bi-kamāl-i inbisāṭ va iḥāṭe min-dūn-i ḥudūd nāzil shude
bi-ḥujubāt-i nafs-i khud chinān maḥjūb namūde-īd ke alsun-i ʿālamīn az
dhikrash ʿājizand.

Now, you who journey in the valley of ignorance and error, you have
drunk the cup of zaqqūm and deemed the Salsabīl of holiness sealed. You
have burned in the fire of desire and counted yourselves among the
inhabitants of the highest pavilion. What a pity that not for one moment
have you reflected on the consummate wisdom of the All-Merciful, nor for
a single instant turned with illumined sight toward the Cause of the
Master of the Day of Judgment. You have so veiled with the veils of your
own selves the words of the All-Merciful, revealed in perfect breadth
and all-encompassing scope, without limits, that the tongues of all the
worlds are powerless to describe it.

\subsection{204}\label{section-1759}

\begin{otherlanguage}{persian}

حال در این فقره ملاحظه نمائید که از قلم قدس الهی جاری شده میفرماید تنزّل
فرمودم تا بمقامی که میفرمایم منم ذرّه و دون ذرّه مثل آنکه فرمودم منم ربّ و
مربّب کلّ ربوب بلکه استغفار میکنم از این هر دو ذکر حال ملاحظه نمائید که آن
جمال قدم خود را بکلّ اسماء نامیده‌اند مع‌ذلک میتوان گفت که ذرّه فوق شمس است
و یا تراب فوق ربّ الأرباب فأنصفوا ان انتم من المنصفین

\end{otherlanguage}

Ḥāl dar īn fiqre mulāḥiẓe namāyīd ke az Qalam-i Quds-i Ilāhī jārī shude.
Mīfarmāyad: tanazzul farmūdam tā bi-maqāmī ke mīfarmāyam manam dharre va
dūn-i dharre, mithl-i ānke farmūdam manam rabb va murabbib-i kull-i
rubūb, balke istighfār mīkunam az īn har du dhikr. Ḥāl mulāḥiẓe namāyīd
ke ān Jamāl-i Qidam khud rā bi-kull-i asmāʾ nāmīde-and. Maʿa-dhālik
mītavān guft ke dharre fawq-i shams ast va yā turāb fawq-i Rabb
al-Arbāb? Fa-anṣifū in antum mina al-munṣifīn.

Now consider this passage that has flowed from the holy Pen of God. He
says, ``I descended to the station where I say, `I am a mote, and less
than a mote,' just as I said, `I am the Lord and the nurturer of all
lords.' Rather, I ask forgiveness for both these statements.'' Now
consider how that Ancient Beauty has called Himself by every name. Can
it nevertheless be said that the mote is above the sun, or the dust
above the Lord of lords? Be fair, if you are among the fair-minded.

\subsection{205}\label{section-1760}

\begin{otherlanguage}{arabic}

فوالّذی جعلنی مظهر ذاته ثمّ کینونته و عزّه و بهائه که اگر در این کلمات
مبارکهٴ الهیّه تفکّر نمائید ابواب علوم ما لا نهایه بر وجه قلوب مفتوح شود
بشأنی که دیگر مضطرب نشوید از اینکه فلان را باسمآء اللّه نامیده یا بذات
اللّه و مرآت اللّه چه که جمیع اسماء در آن ساحت مساوی و لا فرق بینها و کلّ
این اسماء طائف حول مظهر نفس قدم بوده و خواهد بود مثل اینکه در این حین
اگر جمال قدم شجری را بکلّ اسمآء حسنی و صفات علیا موسوم فرماید انّه لیقدر و
لیس لأحد ان یعترض علیه لأنّه هو المختار و ما سواه فی قبضة الاقتدار لا اله
الّا هو العزیز الجبّار

\end{otherlanguage}

Fa-wa-alladhī jaʿalanī maẓhara dhātihi thumma kaynūnatihi wa-ʿizzihi
wa-bahāʾihi, ke agar dar īn kalimāt-i mubārake-yi ilāhiyye tafakkur
namāyīd abvāb-i ʿulūm-i mā lā nihāye bar vajh-i qulūb maftūḥ shavad,
bi-shaʾnī ke dīgar muḍṭarib nashavīd az īnke fulān rā bi-asmāʾ-i Allāh
nāmīde yā bi-Dhāt-i Allāh va Mirʾāt-i Allāh. Chi ke jamīʿ-i asmāʾ dar ān
sāḥat musāvī va lā farqa baynahā, va kull-i īn asmāʾ ṭāʾif-i ḥawl-i
Maẓhar-i Nafs-i Qidam būde va khvāhad būd. Mithl-i īnke dar īn ḥīn agar
Jamāl-i Qidam shajarī rā bi-kull-i asmāʾ-i ḥusnā va ṣifāt-i ʿulyā mawsūm
farmāyad, innahu la-yaqdiru wa-laysa li-aḥadin an yaʿtariḍa ʿalayhi
li-annahu huwa al-Mukhtāru wa-mā siwāhu fī qabḍati al-iqtidār. Lā ilāha
illā huwa al-ʿAzīzu al-Jabbār.

By Him who has made Me the Manifestation of His Essence, His Being, His
might, and His glory! If you reflect on these blessed divine words, the
gates of limitless knowledge will open before your hearts, so that you
will no longer be troubled because someone has been called by the names
of God, or called the Essence of God and the Mirror of God. For all
names are equal in that court, with no distinction between them, and all
these names have revolved, and will continue to revolve, around the
Manifestation of the eternal Self. For example, if at this moment the
Ancient Beauty should designate a tree by all the most beautiful names
and most exalted attributes, He assuredly has the power to do so, and no
one has the right to object to Him, for He is the One who chooses, and
all besides Him is within the grasp of His power. There is no god but
Him, the Almighty, the All-Compelling.

\subsection{206}\label{section-1761}

\begin{otherlanguage}{arabic}

و دیگر آنکه کلّ خلق را مرایاء لطیفه نامیده‌اند و میفرمایند اگر مسّ کند آن
مرایا را جذبات حبّ الهی جذب میفرماید و اگر مسّ نماید دون حبّ نازل میشود بر
عین آن مرایا غباری که حایل میشود مابین آن مرآت و بین عرفان امر پروردگار
قادر علیم بقوله جلّ و عزّ انّ الخلق مرایا لطیفة ان یمسّها حبّ یجذب و ان یمسّها
دون حبّ ینزل عینها غبار یحولنّ بینه و بین امر ربّه انتهی اذاً تفکّروا یا ملأ
الغافلین

\end{otherlanguage}

Va dīgar ānke kull-i khalq rā marāyā-yi laṭīfe nāmīde-and va mīfarmāyand
agar mass kunad ān marāyā rā jadhabāt-i ḥubb-i ilāhī jadhb mīfarmāyad,
va agar mass namāyad dūn-i ḥubb nāzil mīshavad bar ʿayn-i ān marāyā
ghubārī ke ḥāyil mīshavad mābayn-i ān mirʾāt va bayn-i ʿirfān-i amr-i
Parvardigār-i Qādir-i ʿAlīm. Bi-qawlihi jalla va ʿazza: inna al-khalqa
marāyā laṭīfatun in yamassahā ḥubbun yajdhibu wa-in yamassahā dūnu
ḥubbin yanzilu ʿaynahā ghubārun yaḥūlanna baynahu wa-bayna amri rabbihi.
Intahā. Idhan tafakkarū yā malaʾa al-ghāfilīn.

Furthermore, He has called all creation delicate mirrors, and says that
if the attractions of divine love touch those mirrors, they draw them
onward; but if anything other than love touches them, dust settles upon
the very surface of those mirrors, coming between the mirror and
recognition of the Cause of the almighty, all-knowing Lord. In His
words---glorious and mighty is He---``Creation consists of delicate
mirrors: if love touches them, it attracts; and if anything other than
love touches them, dust settles upon their surface, surely coming
between one and the Cause of his Lord.'' Thus ends the quotation.
Reflect, then, you heedless ones.

\subsection{207}\label{section-1762}

\begin{otherlanguage}{persian}

و این بیان را نقطهٴ بیان جلّت عظمته عموماً فرموده‌اند و اگر بتخصیص ناظری
بجناب آقا سیّد جواد خطاب میفرمایند قوله عزّ و جلّ که توئی آن مرآت اوّلیّه که
لم‌یزل از حقّ حکایت نموده و لایزال حکایت خواهی نمود و هم‌چنین او را علّت
اوّلیّه فرموده‌اند فواللّه که این رتبه فوق کلّ اسماء است چه از ذات اللّه و
کینونة اللّه و ذکر اللّه و مرآت اللّه چه که از قبل هر نفسی که این رتبه را
در حقّ رسول اللّه قائل شده حکم بکفرش نموده‌اند زیرا که علّت اوّلیّه خود حقّ را
دانسته‌اند

\end{otherlanguage}

Va īn bayān rā Nuqṭe-yi Bayān, jallat ʿaẓamatuhu, ʿumūman farmūde-and,
va agar bi-takhṣīṣ nāẓirī bi-janāb-i Āqā Siyyid Javād khiṭāb
mīfarmāyand, qawluhu ʿazza va jalla: ke tuʾī ān mirʾāt-i avvaliyye ke
lam-yazal az Ḥaqq ḥikāyat namūde va lā-yazāl ḥikāyat khvāhī namūd. Va
ham-chinīn ū rā ʿillat-i avvaliyye farmūde-and. Fa-wa-Allāh ke īn rutbe
fawq-i kull-i asmāʾ ast, chi az Dhāt-i Allāh va Kaynūnat-i Allāh va
Dhikr-i Allāh va Mirʾāt-i Allāh. Chi ke az qabl har nafsī ke īn rutbe rā
dar ḥaqq-i Rasūl-i Allāh qāʾil shude ḥukm bi-kufrash namūde-and, zīrā ke
ʿillat-i avvaliyye khud-i Ḥaqq rā dāniste-and.

The Point of the Bayān---exalted be His grandeur---made this statement
in general terms. If, however, you are looking for a particular
instance, He addresses Āqā Siyyid Javād in these words---mighty and
glorious is He---``You are that Primal Mirror who has always reflected
God and will forever reflect Him.'' He has also called him the Primal
Cause. By God, this rank is above all names, whether the Essence of God,
the Being of God, the Remembrance of God, or the Mirror of God. For in
the past, anyone who ascribed this rank to the Messenger of God was
judged an unbeliever, since the Primal Cause was held to be God Himself.

\subsection{208}\label{section-1763}

\begin{otherlanguage}{persian}

و هم‌چنین در بارهٴ او نازل قوله جلّ ذکره انّا جعلناک نبیّاً علی العالمین حال
با جمیع این مراتب و وصف که فوق آن در بیان ذکر نشده اگر العیاذ باللّه از
حقّ اعراض نماید جمیع اقرب من آن از او مرتفع خواهد شد و اگر ببصر منیر
ملاحظه نمائید کلّ اشیاء را در بیان بکلّ اسماء حسنی نامیده‌اند ولکن لا
یعقلها الّا المخلصون و این کلمات ذکر میشود که همچه ندانند که اوصاف منزله
در بیان مخصوص نفسی دون نفسی بوده بلکه در کلّ ثابت مادامی که در ظلّ حقّ
مستقرّ باشند و بعد از خروج کلّ محروم و ممنوع

\end{otherlanguage}

Va ham-chinīn dar bāre-yi ū nāzil, qawluhu jalla dhikruhu: innā
jaʿalnāka nabiyyan ʿalā al-ʿālamīn. Ḥāl bā jamīʿ-i īn marātib va vaṣf ke
fawq-i ān dar Bayān dhikr nashude, agar al-ʿiyādhu bi-Allāh az Ḥaqq
iʿrāḍ namāyad jamīʿ aqrab min ān az ū murtafiʿ khvāhad shud. Va agar
bi-baṣar-i munīr mulāḥiẓe namāyīd kull-i ashyāʾ rā dar Bayān bi-kull-i
asmāʾ-i ḥusnā nāmīde-and, wa-lākin lā yaʿqiluhā illā al-mukhliṣūn. Va īn
kalimāt dhikr mīshavad ke hamche nadānand ke awṣāf-i munzale dar Bayān
makhṣūṣ-i nafsī dūn-i nafsī būde, balke dar kull thābit mādāmī ke dar
ẓill-i Ḥaqq mustaqarr bāshand, va baʿd az khurūj kull maḥrūm va mamnūʿ.

Likewise, these words were revealed concerning him---glorious be His
mention---``We have made you a prophet over all the worlds.'' Now,
despite all these ranks and this description, above which nothing has
been mentioned in the Bayān, should he---God forbid---turn away from
God, all would be stripped from him in less than an instant. If you
consider with illumined sight, He has called all things in the Bayān by
all the most beautiful names, but none comprehend this except the
sincere. These words are mentioned so that no one should suppose that
the descriptions revealed in the Bayān were exclusive to one person
rather than another. Rather, they hold true of all, as long as they
remain beneath the shadow of God; once they leave it, all are deprived
and debarred.

\subsection{209}\label{section-1764}

\begin{otherlanguage}{persian}

و هم‌چنین در رتبهٴ جناب قدّوس روح العالمین فداه میفرمایند عزّ ذکره که بعدد
هشت واحد از مرآت اللّه بر مقعد خود بوده و از شدّت نار محبّت او کسی قدرت بر
قرب بهم نرسانده انتهی حال ملاحظه نمائید که هشت واحد از مرآت اللّه را در
نفس مبارکشان ذکر فرموده‌اند فتفکّروا یا اولی الأفکار و در این بیان کنایز
علمیّه مستور است و این عبد مذکور نداشت خوفاً من نمرود الظّلم و فرعون الجهل

\end{otherlanguage}

Va ham-chinīn dar rutbe-yi janāb-i Quddūs, rūḥu al-ʿālamīna fidāhu,
mīfarmāyand, ʿazza dhikruhu: ke bi-ʿadad-i hasht vāḥid az Mirʾāt-i Allāh
bar maqʿad-i khud būde va az shiddat-i nār-i maḥabbat-i ū kasī qudrat
bar qurb bi-ham narasānde. Intahā. Ḥāl mulāḥiẓe namāyīd ke hasht vāḥid
az Mirʾāt-i Allāh rā dar nafs-i mubārakishān dhikr farmūde-and.
Fa-tafakkarū yā ulī al-afkār. Va dar īn bayān kanāyiz-i ʿilmiyye mastūr
ast va īn ʿabd madhkūr nadāsht khawfan min Nimrūdi al-ẓulmi wa-Firʿawni
al-jahl.

Likewise, concerning the rank of Quddūs---may the soul of all the worlds
be sacrificed for him---He says, exalted be His mention, ``On his seat
he was equal in number to eight vāḥids of the Mirror of God, and such
was the intensity of the fire of his love that no one gained the power
to draw near.'' Thus ends the quotation. Now consider how He has
mentioned eight vāḥids of the Mirror of God in his blessed person.
Reflect, you people of insight. Treasures of knowledge lie hidden in
this utterance, but this servant has not mentioned them, for fear of the
Nimrod of oppression and the Pharaoh of ignorance.

\subsection{210}\label{section-1765}

\begin{otherlanguage}{persian}

و هم‌چنین در جمیع مقامات از قلم عزّ نازل که مرآت بنفسه لنفسه تحقّقی نداشته
الّا بتقابلها بالشّمس و شمس را لم‌یزل و لایزال واحد فرموده‌اند و احدی را
باین اسم موسوم نفرموده جز ذات قدم را در هر ظهوری مع‌ذلک نوشته مقام شمسیّت
تحت رتبهٴ مرآتیّت است و شموس از مرآت مشرق میشوند فواللّه یا قوم ما انطق عن
الهوی که اگر در همین قول تفکّر نمائید بیقین شهادت میدهید که این اقوال از
جهل صرف و بغضای بحت ظاهر شده

\end{otherlanguage}

Va ham-chinīn dar jamīʿ-i maqāmāt az Qalam-i ʿIzz nāzil ke mirʾāt
bi-nafsihi li-nafsihi taḥaqquqī nadāshte illā bi-taqābulihā bi-al-shams,
va shams rā lam-yazal va lā-yazāl vāḥid farmūde-and va aḥadī rā bi-īn
ism mawsūm nafarmūde juz Dhāt-i Qidam rā dar har ẓuhūrī. Maʿa-dhālik
nivishte maqām-i shamsiyyat taḥt-i rutbe-yi mirʾātiyyat ast va shumūs az
mirʾāt mushriq mīshavand. Fa-wa-Allāh yā qawm mā anṭiqu ʿani al-hawā, ke
agar dar hamīn qawl tafakkur namāyīd bi-yaqīn shahādat mīdahīd ke īn
aqvāl az jahl-i ṣirf va baghḍā-yi baḥt ẓāhir shude.

Likewise, in every instance it has been revealed by the Pen of Glory
that the mirror has no reality in and for itself except through facing
the sun. He has declared the sun to have always been, and forever to
remain, one, and has designated no one by this name except the eternal
Essence in every Revelation. Nevertheless, he has written that the
station of being a sun is beneath the rank of being a mirror, and that
suns rise from the mirror. By God, O people, I do not speak from desire!
If you reflect on this statement alone, you will assuredly testify that
these assertions have arisen from sheer ignorance and unalloyed hatred.

\subsection{211}\label{section-1766}

\begin{otherlanguage}{persian}

اقسمکم یا قوم باللّه الّذی خلقکم و سوّاکم که در اوّل بیان فارسی ملاحظه
نمائید که میفرمایند قوله عزّ ذکره اگر مرآت بگوید در من شمس است در نزد شمس
ظاهر است که شبح او است که او میگوید انتهی مع‌ذلک نوشته که مرآت فوق شمس
است و شموس در ظلّ مرآتند اگر مراد از این شموسی که ذکر نموده شموس حقیقت
مقصود است فتعالی عن ذلک چه که بیک تجلّی از تجلّیات شمس حقیقت صدهزار مرآت
اقلّ من حین مبعوث شوند

\end{otherlanguage}

Uqsimukum yā qawm bi-Allāhi alladhī khalaqakum wa-sawwākum ke dar
avval-i Bayān-i Fārsī mulāḥiẓe namāyīd ke mīfarmāyand, qawluhu ʿazza
dhikruhu: agar mirʾāt bigūyad dar man shams ast dar nazd-i shams ẓāhir
ast ke shabaḥ-i ū ast ke ū mīgūyad. Intahā. Maʿa-dhālik nivishte ke
mirʾāt fawq-i shams ast va shumūs dar ẓill-i mirʾātand. Agar murād az īn
shumūsī ke dhikr namūde shumūs-i ḥaqīqat maqṣūd ast, fa-taʿālā ʿan
dhālik, chi ke bi-yik tajallī az tajalliyāt-i Shams-i Ḥaqīqat ṣad-hizār
mirʾāt aqall min ḥīn mabʿūth shavand.

I adjure you, O people, by God who created and fashioned you, to
consider the beginning of the Persian Bayān, where He says---exalted be
His mention---``If the mirror says, `The sun is within me,' it is
evident to the sun that it is its image of which the mirror speaks.''
Thus ends the quotation. Nevertheless, he has written that the mirror is
above the sun and that the suns are beneath the mirror's shadow. If by
these suns he has mentioned he means the Suns of Truth, exalted is He
above that! For through a single effulgence among the effulgences of the
Sun of Truth, a hundred thousand mirrors are raised up in less than an
instant.

\subsection{212}\label{section-1767}

\begin{otherlanguage}{persian}

و من دون ذلک انّ المرآة خلقت لانطباع تجلّی من تجلّیات الشّمس علیها ان کانت
صافیة و من دون ذلک یحکم علیه حکم الحجر چه که بر هر ذی بصری مشهود است که
مرایا از برای اخذ عکوس اشیاء خلق شده‌اند چنانچه ذکر شده و اگر مقصود شموس
اسماء و دون آن بوده این مخصوص بنفسی نه چه که در کلّ این شموس موجود و یظهر
منه فی حین الّذی قدّر اللّه له و من فتح عین فؤاده یشهد فی کلّ ذرّ شمساً ثمّ فی
هذه الشّمس شمساً ثمّ فی شمس الشّمس شمساً بحیث لا نهایة لها ولکن حکمة اللّه ذکر
این مقام ننموده لئلّا یزلّ اقدام العارفین

\end{otherlanguage}

Va min dūni dhālika inna al-mirʾāta khuliqat li-inṭibāʿi tajallin min
tajalliyāti al-shamsi ʿalayhā in kānat ṣāfiyatan wa-min dūni dhālika
yuḥkamu ʿalayhi ḥukmu al-ḥajar, chi ke bar har dhī baṣarī mashhūd ast ke
marāyā az barā-yi akhdh-i ʿukūs-i ashyāʾ khalq shude-and chinānchi dhikr
shude. Va agar maqṣūd shumūs-i asmāʾ va dūn-i ān būde īn makhṣūṣ
bi-nafsī na, chi ke dar kull īn shumūs mawjūd va yaẓharu minhu fī ḥīni
alladhī qaddara Allāhu lahu. Wa-man fataḥa ʿayna fuʾādihi yashhadu fī
kulli dharrin shamsan thumma fī hādhihi al-shamsi shamsan thumma fī
shamsi al-shamsi shamsan bi-ḥaythu lā nihāyata lahā. Va-lākin ḥikmat-i
Allāh dhikr-i īn maqām nanamūde li-allā yazilla aqdāmu al-ʿārifīn.

Besides this, the mirror was created so that an effulgence among the
sun's effulgences might be imprinted upon it, if it is clear; otherwise
it is judged to be a stone. For it is evident to everyone with sight
that mirrors were created to receive the reflections of things, as has
been mentioned. If, however, he meant the suns of names and anything
below them, this is not exclusive to any one person, for these suns
exist in all and appear from each at the time God has ordained for them.
Whoever opens the eye of his heart beholds a sun in every mote, then a
sun within this sun, then a sun within the sun of that sun, without end.
Yet the wisdom of God has left this station unmentioned, lest the feet
of those endowed with spiritual knowledge should slip.

\subsection{213}\label{section-1768}

\begin{otherlanguage}{arabic}

فیا للّه یا قوم من هذا الظّلم الّذی ورد علی جمال القدم من الّذی یفتخر ان
یقوم بین یدیه فلمّا اشتهر اسمه کفر باللّه المهیمن القیّوم و چون ظهور قبلم
عالم بود بر اینکه مرآت در ظهور بعد دعوی شمسی مینماید لذا این حکم از قلم
عزّ صدور یافت

\end{otherlanguage}

Fa-yā li-Allāh yā qawm min hādhā al-ẓulmi alladhī warada ʿalā Jamāli
al-Qidami mina alladhī yaftakhiru an yaqūma bayna yadayhi fa-lammā
ishtahara ismuhu kafara bi-Allāhi al-Muhaymini al-Qayyūm. Va chūn
Ẓuhūr-i Qablam ʿālim būd bar īnke mirʾāt dar ẓuhūr-i baʿd daʿvā-yi
shamsī mīnamāyad lidhā īn ḥukm az Qalam-i ʿIzz ṣudūr yāft.

For God's sake, O people, consider this injustice inflicted upon the
Ancient Beauty by one who prides himself on standing in His presence!
Yet when his name became well known, he disbelieved in God, the
All-Protecting, the Self-Subsisting. And since My previous Manifestation
knew that the mirror would claim to be the sun in the following
Revelation, this pronouncement issued from the Pen of Glory.

\subsection{214}\label{section-1769}

\begin{otherlanguage}{persian}

و بشأنی در بحر کبر و غرور مستغرق شده که بکلمات قبلم در علوّ شأن خود تمسّک
جسته و بر منبع و معدن و مظهر و موجد و منزل آن معرض شده فبعداً للقوم
الظّالمین و هم‌چنین خود را شجرهٴ عظمت نامیده فواللّه جمیع این کلمات نظر بآن
است که ناس را احمق یافته

\end{otherlanguage}

Va bi-shaʾnī dar baḥr-i kibr va ghurūr mustaghraq shude ke bi-kalimāt-i
Qablam dar ʿuluvv-i shaʾn-i khud tamassuk juste va bar manbaʿ va maʿdin
va maẓhar va mūjid va munzil-i ān muʿriḍ shude. Fa-buʿdan li-al-qawmi
al-ẓālimīn. Va ham-chinīn khud rā Shajare-yi ʿAẓamat nāmīde. Fa-wa-Allāh
jamīʿ-i īn kalimāt naẓar bi-ān ast ke nās rā aḥmaq yāfte.

He has become so submerged in the ocean of pride and arrogance that he
has seized upon My former words concerning the loftiness of his station,
yet turned away from their source, their wellspring, their
manifestation, their originator, and the One who revealed them. Away,
then, with the people of wrongdoing! He has likewise called himself the
Tree of Grandeur. By God, all these words are due to his having found
the people foolish.

\subsection{215}\label{section-1770}

\begin{otherlanguage}{persian}

و اگر از شجره عنداللّه محسوب شود من شجرة الّتی ینبغی لها ان تقطع و تلقی فی
النّار الّا بأن یتوب الی اللّه و یستغفر عمّا فعل و یکون من التّائبین و بر
جمیع واضح است چه از مهاجرین این ارض و چه هر فطن بصیری که ببصر اللّه در
امور ناظرند که این کلمات غروریّه از لسان مجعولهٴ کذبهٴ شرکیّه جاری نشده
مگر آنکه چون این شمس بدیع ابدع اعلی را از افق فجر بقا طالع و بنغمهٴ انّنی
انا حیّ فی افق الأبهی مابین ارض و سماء ناطق دیده لذا تعجیل نموده و من حیث
لا یشعر باین کلمات تکلّم کرده که شأنی فوق شأن اللّه از برای خود ثابت نماید

\end{otherlanguage}

Va agar az shajare ʿinda Allāh maḥsūb shavad min shajaratin allatī
yanbaghī lahā an tuqṭaʿa wa-tulqā fī al-nāri illā bi-an yatūba ilā
Allāhi wa-yastaghfira ʿammā faʿala wa-yakūna mina al-tāʾibīn. Va bar
jamīʿ vāḍiḥ ast chi az muhājirīn-i īn arḍ va chi har faṭin-i baṣīrī ke
bi-baṣar-i Allāh dar umūr nāẓirand ke īn kalimāt-i ghurūriyye az lisān-i
majʿūle-yi kadhabe-yi shirkiyye jārī nashude magar ānke chūn īn Shams-i
Badīʿ-i Abdaʿ-i Aʿlā rā az ufuq-i fajr-i baqā ṭāliʿ va bi-naghme-yi
innanī anā ḥayyun fī ufuqi al-Abhā mābayn-i arḍ va samāʾ nāṭiq dīde
lidhā taʿjīl namūde va min ḥaythu lā yashʿur bi-īn kalimāt takallum
karde ke shaʾnī fawq-i shaʾn-i Allāh az barā-yi khud thābit namāyad.

And if he is reckoned a tree in God's sight, he is among the trees that
should be cut down and cast into the fire, unless he repents to God,
seeks forgiveness for what he has done, and becomes one of the
repentant. It is clear to everyone---both those who have migrated to
this land and every discerning person who considers affairs with the eye
of God---that these arrogant words have flowed from that false, lying,
idolatrous tongue only because he has seen this wondrous, most
marvelous, most exalted Sun rising from the horizon of the dawn of
eternity and proclaiming between earth and heaven the melody, ``I am
indeed alive in the all-glorious horizon.'' He has therefore made haste
and, without realizing it, uttered these words to establish for himself
a station above the station of God.

\subsection{216}\label{section-1771}

\begin{otherlanguage}{arabic}

و هذا لن یمکن ابداً لأنّ ما دونه مخلوق بأمره و منجعل بارادته و متحرّک باذنه
و ما بعده الّا عبده و رعیّته و خلقه و بریّته و عباده هل یکون فوق شأن اللّه
شأن لیثبته احد لنفسه و انّه قام بنفسه لنفسه فی علوّ نفسه مقام الّذی لا له
فوق و لا تحت و لا یمین و لا یسار و لا امام و لا خلف و انّه قد خلق الجهات
لا من جهة بمشیّته و الأشیآء لا من مثال بارادته و انّه لخالق کلّ شأن و منزل
کلّ امر و یستحقّ لدونه بأن یفتخر بنفسه و انّه لا یفتخر بأحد من الموجودات و
انّه لهو المقدّس المنزّه المقتدر المطاع

\end{otherlanguage}

Wa-hādhā lan yumkina abadan li-anna mā dūnahu makhlūqun bi-amrihi
wa-munjaʿalun bi-irādatihi wa-mutaḥarrikun bi-idhnihi wa-mā baʿdahu illā
ʿabduhu wa-raʿiyyatuhu wa-khalquhu wa-bariyyatuhu wa-ʿibāduhu. Hal
yakūnu fawqa shaʾni Allāhi shaʾnun li-yuthbitahu aḥadun li-nafsihi?
Wa-innahu qāma bi-nafsihi li-nafsihi fī ʿuluwwi nafsihi maqāma alladhī
lā lahu fawqun wa-lā taḥtun wa-lā yamīnun wa-lā yasārun wa-lā amāmun
wa-lā khalf. Wa-innahu qad khalaqa al-jihāta lā min jihatin
bi-mashīʾatihi wa-al-ashyāʾa lā min mithālin bi-irādatihi. Wa-innahu
la-khāliqu kulli shaʾnin wa-munzilu kulli amrin wa-yastaḥiqqu li-dūnihi
bi-an yaftakhira bi-nafsihi wa-innahu lā yaftakhiru bi-aḥadin mina
al-mawjūdāti wa-innahu la-huwa al-Muqaddasu al-Munazzahu al-Muqtadiru
al-Muṭāʿ.

This can never be possible, for everything other than Him is created by
His command, brought into being by His will, and set in motion by His
permission. All that comes after Him is but His servant, His subject,
His creation, His creature, and His servants. Is there any station above
the station of God that anyone might claim for himself? He stands by
Himself, for Himself, in the loftiness of His own Self, in a station
that has no above or below, no right or left, no front or back. By His
will He created the directions, not from any direction, and by His
purpose He created all things, not from any pattern. He is indeed the
Creator of every station and the Revealer of every command. It befits
those other than Him to glory in His own Self, while He glories in none
of the beings in existence. He is indeed the Holy, the Transcendent, the
All-Powerful, the One who is obeyed.

\subsection{217}\label{section-1772}

\begin{otherlanguage}{persian}

ای قوم بشنوید ندایم را و بر جمالم وارد نیاورید آنچه را که در کلّ الواح
ممنوع شده‌اید تاللّه این است ید قدرت منبسطهٴ الهیّه که احاطه فرموده کلّ من
فی السّموات و الأرض را و این است لسان حقیقت ربّانیّه که بر اعلی سدرهٴ عزّ
احدیّه تغنّی میفرماید و اینست قلم قدم که در ایّام و لیالی بحور اسرار و
معارف الهی از او جاری و ساری است و اینست مخزن لآلی حکم بالغه و معدن علوم
الهیّه و لم‌یزل و لایزال در افاضه بوده و خواهد بود زینهار که خود را از
رشحات فضلش محروم منمائید و متابعت هر نفس بی‌بصری را واجب مشمرید

\end{otherlanguage}

Ey qawm, beshnavīd nedāyam rā va bar jamālam vāred nayāvarīd ānche rā ke
dar koll-e alvāḥ mamnūʿ shode-īd. Ta-llāh, īn ast yad-e qodrat-e
monbaseṭe-ye elāhiyye ke eḥāṭe farmūde koll-e man fī al-samāwāt
wa-al-arḍ rā va īn ast lesān-e ḥaqīqat-e rabbāniyye ke bar aʿlā sedre-ye
ʿezz-e aḥadiyye taghannī mīfarmāyad va īnast qalam-e qedam ke dar ayyām
va layālī boḥūr-e asrār va maʿāref-e elāhī az ū jārī va sārī ast va
īnast makhzan-e laʾālī-ye ḥekam-e bāleghe va maʿdan-e ʿolūm-e elāhiyye
va lam-yazal va lā-yazāl dar efāḍe būde va khwāhad būd. Zīnhār ke khwod
rā az rashaḥāt-e faḍlash maḥrūm manamāyīd va motābaʿat-e har nafs-e
bī-baṣarī rā vājeb mashmorīd.

O people, hear My call, and do not inflict upon My Beauty what you have
been forbidden to do in all the Tablets. By God, this is the
outstretched hand of divine power that has encompassed all who are in
the heavens and on earth. This is the tongue of divine truth that sings
upon the loftiest Tree of the glory of oneness. This is the Pen of
Eternity from which oceans of divine mysteries and knowledge flow and
stream by day and by night. This is the treasury of the pearls of
consummate wisdom and the mine of divine knowledge, which has always
bestowed its bounty and will continue to do so forever. Beware lest you
deprive yourselves of the sprinklings of its grace, and do not consider
it obligatory to follow every sightless soul.

\subsection{218}\label{section-1773}

\begin{otherlanguage}{persian}

و هم‌چنین در رسالهٴ خود مذکور داشته که بمحض دانستن چند کلام و عبارت فهمی
این دعویها نموده و مقصودش از این کلمات مزخرفه این هیکل قدس صمدانی بوده
فوالّذی تفرّد بالقدرة و الاقتدار که در کلّ لیالی و ایّام مثل آنکه معلّمین
اطفال را حروف هجائیّه تعلیم مینمایند بهمان قسم در تعلیمش جهد منیع مبذول
شد فیا لیت قومی یعلمون ما صغر حدّه و حقر شأنه و کثر غروره و کبر قوله و
این بر هر ذی بصری واضح و مبرهن است احتیاج بذکر تفصیل نداشته و ندارد

\end{otherlanguage}

Va hamchenīn dar resāle-ye khwod madhkūr dāshte ke be-maḥḍ-e dānestan-e
chand kalām va ʿebārat-fahmī īn daʿvīhā namūde va maqṣūdash az īn
kalemāt-e mozakhrafe īn haykal-e qods-e ṣamadānī būde. Fa-wa-alladhī
tafarrada bi-al-qudrati wa-al-iqtidār ke dar koll-e layālī va ayyām
mesl-e ānke moʿallemīn aṭfāl rā ḥorūf-e hejāʾiyye taʿlīm mīnamāyand
be-hamān qesm dar taʿlīmash jahd-e manīʿ mabdhūl shod. Fa-yā layta qawmī
yaʿlamūna mā ṣaghura ḥadduhu wa-ḥaqura shaʾnuhu wa-kathura ghurūruhu
wa-kabura qawluhu. Va īn bar har dhī-baṣarī vāḍeḥ va mobarhan ast
eḥtiyāj be-dhekr-e tafṣīl nadāshte va nadārad.

He has likewise stated in his treatise that merely by knowing a few
words and understanding some phrases, He has made these claims. By these
absurd words he meant this sacred, eternal Temple. By Him who is
peerless in power and might, throughout the nights and days, just as
teachers teach children the letters of the alphabet, strenuous efforts
were devoted to teaching him. If only My people knew how limited his
capacity, how lowly his station, how abundant his conceit, and how
grandiose his speech! This is clear and evident to every discerning
person and has never required, nor does it now require, a detailed
account.

\subsection{219}\label{section-1774}

\begin{otherlanguage}{persian}

و از جمله نسبت داده که این عبد نسبت بنقطهٴ اولی بخلاف ادب تکلّم نموده و
حال اینکه این عبد ناطق است در مابین سموات و ارض بأنّنی انا نفسه و ذاته و
روحه و هیکله و بهائه مع‌ذلک چگونه میشود بآن جوهر اعلی و ساذج ابهی روح من
فی ممالک البقآء فداه بغیر رضآء اللّه تکلّم نمایم باری بعینه مثل ملأ فرقان
که در اوّل فجر ظهور نسبت میدادند که این طایفه به رسول اللّه حبّ ندارند که
شاید باین مزخرفات عباد را از عرفان نفس اللّه فی یوم المعاد محروم نمایند
چنانچه نموده و مینمایند الا لعنة اللّه علی المفترین

\end{otherlanguage}

Va az jomle nesbat dāde ke īn ʿabd nesbat be-Noqṭe-ye Ūlā be-khelāf-e
adab takallom namūde va ḥāl īnke īn ʿabd nāṭeq ast dar mābayn-e samāvāt
va arḍ bi-annanī anā nafsuhu wa-dhātuhu wa-rūḥuhu wa-haykaluhu
wa-bahāʾuhu. Maʿa-dhālek chegūne mīshavad be-ān jawhar-e aʿlā va
sādhaj-e abhā rūḥu man fī mamāliki al-baqāʾ fidāhu be-ghayr-e reḍā-ye
Allāh takallom namāyam? Bārī be-ʿayneh mesl-e malaʾ-e Forqān ke dar
avval-e fajr-e ẓohūr nesbat mīdādand ke īn ṭāyefe be-Rasūl Allāh ḥobb
nadārand ke shāyad be-īn mozakhrafāt ʿebād rā az ʿerfān-e nafs-e Allāh
fī yawm al-maʿād maḥrūm namāyand chenānche namūde va mīnamāyand. Alā
laʿnatu Allāhi ʿalā al-muftarīn.

Among other things, he has alleged that this servant has spoken
disrespectfully of the Primal Point, whereas this servant proclaims
between the heavens and the earth, ``I am indeed His very Self, His
Essence, His Spirit, His Temple, and His Glory.'' How, then, could I
speak of that most exalted Essence and most glorious Reality---may the
spirit of everyone in the realms of eternity be sacrificed for Him---in
a manner contrary to God's good pleasure? In short, this is exactly like
the people of the Qur'an who, at the first dawn of the Revelation,
alleged that this community had no love for the Messenger of God, hoping
through such absurdities to deprive God's servants of recognizing His
own Self on the Day of Return, as they have done and continue to do.
Indeed, God's curse be upon those who fabricate falsehoods.

\subsection{220}\label{section-1775}

\begin{otherlanguage}{persian}

با اینکه جمیع عالم را آثار بدیعه احاطه نموده و جمیع اهل بیان بر این شاهد
و گواهند و معادل ما نزّل فی البیان از این ظهور بدیع نازل و از اهل بیان
نفسی نه که اثر اللّه نزدش موجود نباشد و من دون ذلک آنچه در این ارض موجود
کتّاب از تحریرش عاجز مانده‌اند چنانچه اکثری بی‌سواد مانده مع‌ذلک نوشته بمجرّد
عبارت فهمی و چهار کلمه این دعویها نموده‌اند حال شما انصاف دهید قولی که
کذبش مثل شمس در وسط السّمآء مشرق و لایح و واضح است مع‌ذلک بچه جرئت و جسارت
تکلّم نموده

\end{otherlanguage}

Bā īnke jamīʿ-e ʿālam rā āthār-e badīʿe eḥāṭe namūde va jamīʿ-e ahl-e
Bayān bar īn shāhed va govāhand va moʿādel-e mā nuzzila fī al-Bayān az
īn ẓohūr-e badīʿ nāzel va az ahl-e Bayān nafsī na ke athar-e Allāh
nazdash mawjūd nabāshad va min dūn-e dhālek ānche dar īn arḍ mawjūd
kottāb az taḥrīrash ʿājez mānde-and chenānche aktharī bī-savād mānde
maʿa-dhālek neveshte be-mojarrad-e ʿebārat-fahmī va chahār kaleme īn
daʿvīhā namūde-and. Ḥāl shomā enṣāf dehīd qawlī ke kedhbash mesl-e shams
dar vasaṭ-e al-samāʾ moshreq va lāyeḥ va vāḍeḥ ast maʿa-dhālek be-che
jorʾat va jasārat takallom namūde.

Although wondrous writings have encompassed the whole world, and all the
people of the Bayán bear witness and testify to this; although the
equivalent of what was revealed in the Bayán has descended from this
wondrous Revelation, and there is not one among the people of the Bayán
who does not possess a writing of God; and although, beyond this, the
scribes have been unable to transcribe what exists in this land, so that
most of it remains uncopied---nevertheless, he has written that merely
through understanding phrases and knowing a few words, He has made these
claims. Now judge fairly: with what audacity and impudence he has
uttered a statement whose falsehood is as manifest, apparent, and clear
as the sun in the midst of heaven!

\subsection{221}\label{section-1776}

\begin{otherlanguage}{persian}

اگرچه فواللّه الّذی لا اله الّا هو که این عبد ابداً اراده نداشته که آنچه از
سماء مشیّت بدعاً نازل میشود ببلاد اشتهار یابد چه که انظر مطهّره بسیار قلیل
مشاهده می‌شود که قابل ملاحظهٴ آثار اللّه باشند لذا لازال مستور میداشتم و
چه مقدار آیات اللّه که در عراق نازل و جمیع در شطّ محو شد و حال معادل
صدهزار بیت در این ارض موجود و هنوز سواد نشده تا چه رسد بجلد چه که این
امور که معلّق بخودنمائی بوده لازال این عبد از او مقدّس بوده و خواهد بود

\end{otherlanguage}

Agarche fa-wa-Allāhi alladhī lā ilāha illā huwa ke īn ʿabd abadan erāde
nadāshte ke ānche az samā-ye mashiyyat badʿan nāzel mīshavad be-belād
eshtehār yābad che ke anẓor-e moṭahhare besyār qalīl moshāhede mīshavad
ke qābel-e molāḥaẓe-ye āthār-e Allāh bāshand. Ledhā lā-zāl mastūr
mīdāshtam va che meqdār āyāt-e Allāh ke dar ʿErāq nāzel va jamīʿ dar
shaṭṭ maḥv shod va ḥāl moʿādel-e ṣad-hezār bayt dar īn arḍ mawjūd va
hanūz savād nashode tā che rasad be-jeld che ke īn omūr ke moʿallaq
be-khwodnamāʾī būde lā-zāl īn ʿabd az ū moqaddas būde va khwāhad būd.

Yet, by God, besides whom there is no god, this servant has never
intended that what is newly revealed from the heaven of the divine will
should become widely known throughout the lands, for very few purified
eyes are to be found that are worthy of beholding the writings of God.
Therefore I have always kept them concealed. How many verses of God were
revealed in Iraq, all of which were effaced in the river! And now the
equivalent of a hundred thousand verses exists in this land, not yet
transcribed, let alone bound into volumes. For this servant has always
been, and will always remain, sanctified from such matters as pertain to
self-display.

\subsection{222}\label{section-1777}

\begin{otherlanguage}{persian}

و چند مرتبه بعضی خواسته‌اند که آنچه موجود است مجلّدات نموده باطراف ارسال
دارند و هنوز اذن داده نشده چه که هیچیک از مظاهر الهیّه بنفسه مباشر
اینگونه امور نشده‌اند چنانچه نقطهٴ اولی روح ما سواه فداه خود بنفسه اعتنا
باین امور نفرموده و فرموده‌اند که بعد باحسن نظم منظّم شود و میفرمایند عزّ
ذکره فطوبی لمن ینظر الی نظم بهآءاللّه و یشکر ربّه فانّه یظهر و لا مردّ له من
عند اللّه فی البیان انتهی و هم‌چنین از قبل رسول ‌اللّه که بعد از ارتقاء آن
ذات قدم قرآن جلد شد و من قبله انجیل

\end{otherlanguage}

Va chand martabe baʿḍī khwāste-and ke ānche mawjūd ast mojalladāt namūde
be-aṭrāf ersāl dārand va hanūz edhn dāde nashode che ke hīch-yek az
maẓāher-e elāhiyye be-nafseh mobāsher-e īngūne omūr nashode-and
chenānche Noqṭe-ye Ūlā rūḥ-e mā sevāh fedāh khwod be-nafseh eʿtenā be-īn
omūr nafarmūde va farmūde-and ke baʿd be-aḥsan-e naẓm monaẓẓam shavad va
mīfarmāyand ʿazza dhikruhu fa-ṭūbā li-man yanẓuru ilā naẓmi Bahāʾ Allāhi
wa-yashkuru rabbahu fa-innahu yaẓharu wa-lā maradda lahu min ʿindi
Allāhi fī al-Bayān intahā. Va hamchenīn az qabl Rasūl Allāh ke baʿd az
erteqā-ye ān dhāt-e qedam Qorʾān jeld shod va min qablihi Enjīl.

Several times, some have wished to bind what exists into volumes and
send them abroad, but permission has not yet been given, for none of the
Manifestations of God has personally undertaken such matters. Thus the
Primal Point---may the spirit of all besides Him be sacrificed for
Him---did not Himself attend to these matters, but said that they should
later be arranged in the finest order. He says---exalted be His
mention---``Blessed is the one who beholds the order of Bahá'u'lláh and
gives thanks to his Lord, for He will assuredly appear, and there is no
averting Him, as ordained by God in the Bayán.'' Here ends the
quotation. Likewise, in the past, with the Messenger of God: the Qur'an
was bound into a volume after that eternal Being's ascension, and before
Him, the Gospel.

\subsection{223}\label{section-1778}

\begin{otherlanguage}{persian}

و با آنکه آیات عزّ قدس شرق و غرب عالم را احاطه نموده باین کلمات تشبّث
جسته‌اند و بچهار کلمه تعبیر نموده‌اند چنانچه اعمال حجّ در سنهٴ قبل متعدّد از
قلم عزّ صدور یافت و مع‌ذلک ارسال نشد مگر یک سوره که به سورهٴ حجّ معروف است
هرگز این عبد در اشتهار اینگونه امور اوقات صرف ننموده و نخواهد نمود شأن
نزول شأن حقّ است و انتشار شأن خلق و انّه لناشر امره بید النّاشرات من ملائکة
المقرّبین لابدّ از خلف سرادق عصمت ربّانی عبادی روحانی ظاهر شوند و آثار اللّه
را جمع نمایند و باحسن نظم منظّم سازند و هذا حتم لا ریب فیه

\end{otherlanguage}

Va bā ānke āyāt-e ʿezz-e qods sharq va gharb-e ʿālam rā eḥāṭe namūde
be-īn kalemāt tashabboth joste-and va be-chahār kaleme taʿbīr namūde-and
chenānche aʿmāl-e ḥajj dar sane-ye qabl motaʿadded az qalam-e ʿezz ṣodūr
yāft va maʿa-dhālek ersāl nashod magar yek sūre ke be-Sūre-ye Ḥajj
maʿrūf ast. Hargez īn ʿabd dar eshtehār-e īngūne omūr awqāt ṣarf
nanamūde va nakhwāhad namūd. Shaʾn-e nozūl shaʾn-e Ḥaqq ast va enteshār
shaʾn-e khalq wa-innahu la-nāshiru amrihi bi-yadi al-nāshirāti min
malāʾikati al-muqarrabīn. Lā-bodd az khalf-e sorādeq-e ʿeṣmat-e rabbānī
ʿebādī rūḥānī ẓāher shavand va āthār-e Allāh rā jamʿ namāyand va
be-aḥsan-e naẓm monaẓẓam sāzand wa-hādhā ḥatmun lā rayba fīh.

Although the verses of holy glory have encompassed the east and west of
the world, they have seized upon these words and described them as a
mere handful of words. Thus, in the past year, several sets of
pilgrimage rites issued from the Pen of Glory, yet none was sent out
except one surah, known as the Surah of Pilgrimage. This servant has
never spent time publicizing such matters, nor will He ever do so.
Revelation is the province of God, and dissemination is the province of
His creatures. He will indeed spread His Cause through the hands of
those angels brought near to Him who disseminate it. Spiritual servants
will assuredly emerge from behind the tabernacle of divine protection,
gather the writings of God, and arrange them in the finest order. This
is an irrevocable decree, beyond all doubt.

\subsection{224}\label{section-1779}

\begin{otherlanguage}{arabic}

و دیگر آنکه بعضی عبارات نوشته‌اند و نسبت بنقطهٴ بیان روح ما سواه فداه
داده‌اند هذا کذب صراح و حینئذ یبرأ نقطة البیان منهم و من کلماتهم و کان
اللّه یشهد بذلک ان انتم لا تشهدون

\end{otherlanguage}

Va dīgar ānke baʿḍī ʿebārāt neveshte-and va nesbat be-Noqṭe-ye Bayān
rūḥ-e mā sevāh fedāh dāde-and. Hādhā kid hbun ṣurāḥun wa-ḥīnaʾidhin
yabraʾu Nuqṭatu al-Bayāni minhum wa-min kalimātihim wa-kāna Allāhu
yashhadu bi-dhālika in antum lā tashhadūn.

Furthermore, they have written certain statements and attributed them to
the Point of the Bayán---may the spirit of all besides Him be sacrificed
for Him. This is a blatant lie. The Point of the Bayán therefore
disavows them and their words, and God bears witness to this, even if
you do not.

\subsection{225}\label{section-1780}

\begin{otherlanguage}{persian}

و چه مقدار باسماء افتخار مینمایند واللّه الّذی لا اله الّا هو که شبه این
نفوس در ابداع دیده نشده شب و روز اوقات صرف نموده که کلمه‌ئی در کتاب اللّه
بیابند و بخود تفسیر نمایند و این عبد در حین تفریق از اخوی یک جعبه از
خطوط و دوائر و هیاکل که بخطّ نقطهٴ اولی بوده نزد اخوی فرستادم و پیغام
نمودم که چون تو بسیار مایلی که بعد از اعراض از حقّ بآثار آن افتخار نمائی
لذا نزد تو ارسال شد که این هیاکل را بجهت بعضی از اهل دیار مختلفه ارسال
داری و اظهار شأن نمائی و یا آنکه هر نفسی نزدت حاضر میشود انتشار دهی
چنانچه الیوم بهمان عمل مشغولی و بلکه بعضی کلمات مجعولهٴ خود را بآن کلمات
منضم ساخته لعلّ یزلّ بها اقدام العارفین

\end{otherlanguage}

Va che meqdār be-asmāʾ eftekhār mīnamāyand. Wa-Allāhi alladhī lā ilāha
illā huwa ke shebh-e īn nofūs dar ebdāʿ dīde nashode. Shab va rūz awqāt
ṣarf namūde ke kaleme-ī dar ketāb-e Allāh beyāband va be-khwod tafsīr
namāyand va īn ʿabd dar ḥīn-e tafrīq az akhavī yek jaʿbe az khoṭūṭ va
davāʾer va hayākel ke be-khaṭṭ-e Noqṭe-ye Ūlā būde nazd-e akhavī
ferestādam va payghām namūdam ke chūn to besyār māyelī ke baʿd az eʿrāḍ
az Ḥaqq be-āthār-e ān eftekhār namāʾī ledhā nazd-e to ersāl shod ke īn
hayākel rā be-jehat-e baʿḍī az ahl-e deyār-e mokhtalefe ersāl dārī va
eẓhār-e shaʾn namāʾī va yā ānke har nafsī nazdat ḥāḍer mīshavad enteshār
dehī chenānche al-yawm be-hamān ʿamal mashghūlī va balke baʿḍī kalemāt-e
majʿūle-ye khwod rā be-ān kalemāt monḍamm sākhte laʿalla yazilla bihā
aqdāmu al-ʿārifīn.

How greatly they pride themselves on names! By God, besides whom there
is no god, the like of these souls has never been seen in all creation.
They spend their days and nights seeking a word in the Book of God that
they might interpret as referring to themselves. When I, this servant,
parted from My brother, I sent him a box of calligraphic writings,
circles, and temples in the handwriting of the Primal Point, with this
message: ``Since you are so eager, after turning away from the Truth, to
pride yourself on His writings, these have been sent to you so that you
may send these temples to some of the people in different lands and
display your station, or distribute them to whoever comes into your
presence, as indeed you are occupied in doing today. You have even added
some of your own fabricated words to those words, in the hope that
through them the feet of those who have attained recognition might
slip.''

\subsection{226}\label{section-1781}

\begin{otherlanguage}{persian}

باری ای عباد از این مراتب چشم بردارید و دل را مقدّس نمائید چه که الیوم
فضلی ظاهر شده که در یک یوم و لیل اگر کاتب از عهده برآید معادل بیان فارسی
از سماء قدس ربّانی نازل میشود و کذلک بشأن الآیات علی لسان عربیّ بدیع
فاحضروا بین یدی العرش لتشهدوا بعیونکم و لا تقاسوا کلمات اللّه بکلمات دونه
و تکوننّ علی بصیرة منیر

\end{otherlanguage}

Bārī ey ʿebād az īn marāteb chashm bardārīd va del rā moqaddas namāʾīd
che ke al-yawm faḍlī ẓāher shode ke dar yek yawm va layl agar kāteb az
ʿohde bar-āyad moʿādel-e Bayān-e Fārsī az samā-ye qods-e rabbānī nāzel
mīshavad wa-kadhālika bi-shaʾni al-āyāti ʿalā lisānin ʿarabiyyin badīʿ.
Fa-iḥḍurū bayna yaday al-ʿarshi li-tashhadū bi-ʿuyūnikum wa-lā taqāsū
kalimāti Allāhi bi-kalimāti dūnihi wa-takūnanna ʿalā baṣīratin munīr.

In short, O servants, turn your eyes away from these matters and
sanctify your hearts, for today such grace has appeared that in a single
day and night, if a scribe can keep pace, the equivalent of the Persian
Bayán descends from the heaven of divine holiness. The same is true of
verses in a wondrous Arabic tongue. Come, then, before the Throne, that
you may witness with your own eyes. Do not measure the words of God
against the words of anyone besides Him, and possess an illumined
insight.

\subsection{227}\label{section-1782}

\begin{otherlanguage}{arabic}

یا قوم فاشهدوا بما شهد اللّه و لا تلتفتوا الی دونه و لا تکوننّ من المریبین
شهد اللّه انّه لا اله الّا هو و انّ نقطة الأولی عبده و بهائه کذلک نزّل من قبل
من قلم اللّه العلیّ العظیم

\end{otherlanguage}

Yā qawmi fa-ishhadū bi-mā shahida Allāhu wa-lā taltafitū ilā dūnihi
wa-lā takūnanna mina al-murībīn. Shahida Allāhu annahu lā ilāha illā
huwa wa-anna Nuqṭata al-Ūlā ʿabduhu wa-bahāʾuhu kadhālika nuzzila min
qablu min qalami Allāhi al-ʿAliyyi al-ʿAẓīm.

O people, bear witness to that to which God has borne witness, do not
turn toward anyone besides Him, and do not be among the doubters. God
bears witness that there is no god but Him, and that the Primal Point is
His servant and His glory. Thus was it revealed before by the Pen of
God, the Most High, the Most Great.

\subsection{228}\label{section-1783}

\begin{otherlanguage}{persian}

و اگر بمعنی همین یک آیه بتمامه فایز میشدید ابداً از صراط اعظم الهی محتجب
نمیگشتید چنانچه الیوم اکثری از اهل بیان شهادت میدهند و چون باین اسم
مبارک میرسند اذاً تسودّ وجوههم و تضطرب قلوبهم و تزلّ اقدامهم و یکوننّ من
المتوقّفین

\end{otherlanguage}

Va agar be-maʿnī-ye hamīn yek āye be-tamāmeh fāyez mīshodīd abadan az
ṣerāṭ-e aʿẓam-e elāhī moḥtajeb nemīgashtīd chenānche al-yawm aktharī az
ahl-e Bayān shahādat mīdehand va chūn be-īn esm-e mobārak mīresand idhan
taswaddu wujūhuhum wa-taḍṭaribu qulūbuhum wa-tazillu aqdāmuhum
wa-yakūnanna mina al-mutawaqqifīn.

Had you grasped the full meaning of this one verse, you would never have
been veiled from the most great path of God. Yet today most of the
people of the Bayán bear witness, but when they reach this blessed Name,
their faces darken, their hearts are troubled, their feet slip, and they
become those who hesitate.

\subsection{229}\label{section-1784}

\begin{otherlanguage}{arabic}

تاللّه الحقّ هیچ عملی نزد حقّ مقبول نه مگر باقرار باین اسم مبارک مطهّر فطوبی
للذّاکرین و للرّاسخین

\end{otherlanguage}

Ta-Allāhi al-Ḥaqq hīch ʿamalī nazd-e Ḥaqq maqbūl na magar be-eqrār be-īn
esm-e mobārak-e moṭahhar fa-ṭūbā li-al-dhākirīna wa-li-al-rāsikhīn.

By God, the True One, no deed is accepted before God without
acknowledgment of this blessed, sanctified Name. Blessed, then, are
those who remember and those who stand firm.

\subsection{230}\label{section-1785}

\begin{otherlanguage}{persian}

و دیگر آنکه در رسائل افتخار نموده که جناب طاهره و عظیم باو مذعن شده‌اند
اگرچه این عبد افتخار باین گونه امور را مفقود صرف و معدوم بحت میشمرم و
این نه از غرور و استکبار است چه که این عبد در خود شأنی ندیده و نخواهد
دید ولکن نظر باظهار شأن اللّه ذکر شده چه که کینونات حروفات باصلها و اسرها
بقولی مبعوث بوده و خواهند بود

\end{otherlanguage}

Va dīgar ānke dar rasāʾel eftekhār namūde ke Janāb-e Ṭāhere va ʿAẓīm
be-ū modhʿen shode-and agarche īn ʿabd eftekhār be-īn gūne omūr rā
mafqūd-e ṣerf va maʿdūm-e baḥt mīshmoram va īn na az ghorūr va estekbār
ast che ke īn ʿabd dar khwod shaʾnī nadīde va nakhwāhad dīd va-lāken
naẓar be-eẓhār-e shaʾn-e Allāh dhekr shode che ke kaynūnāt-e ḥorūfāt
be-aṣlehā va asrehā be-qawlī mabʿūth būde va khwāhand būd.

Furthermore, he has boasted in his treatises that Ṭáhirih and ʿAẓím
acknowledged him. Yet I, this servant, regard pride in such matters as
sheer nothingness and utter nonexistence. This is not out of conceit or
arrogance, for this servant has never seen, nor will ever see, any
station in Himself. Rather, this has been mentioned to manifest the
station of God, for the very beings of the Letters, in their origin and
entirety, have been, and will be, called into being by a single
utterance.

\subsection{231}\label{section-1786}

\begin{otherlanguage}{persian}

حقّ را لایق نه در اثبات ظهورش بغیر نفسه و ما یظهر من عنده استدلال فرماید
چه که کلّ دلیل بانتسابه الیه محکم بوده و خواهد بود و کلّ سبیل بنسبته الیه
مستقیم ولکن چون ناس ضعیف مشاهده میشوند و غیر بالغ لذا این عبد ذکر
مینمایم که آنچه از کلمات الهی که مخصوص این امر نازل شده خارج از احصای
این نفوس بوده و خواهد بود

\end{otherlanguage}

Ḥaqq rā lāyeq na dar ethbāt-e ẓohūrash be-gheyr-e nafsihi va mā yaẓharu
min ʿindihi estedlāl farmāyad, che ke koll-e dalīl be-entesābash elayh
moḥkam būde va khwāhad būd va koll-e sabīl be-nesbatash elayh mostaqīm.
Va lāken chon nās ḍaʿīf moshāhede mīshavand va gheyr-e bālegh, ledhā īn
ʿabd dhekr mīnamāyam ke ānche az kalemāt-e elāhī ke makhṣūṣ-e īn amr
nāzel shode khārej az eḥṣā-ye īn nofūs būde va khwāhad būd.

It is not fitting for the True One to establish His manifestation by
adducing anything other than Himself and what proceeds from Him, for
every proof has been and will remain sound through its connection to
Him, and every path is straight through its relation to Him. Yet since
people are seen to be weak and immature, I, this servant, mention that
the divine words revealed specifically concerning this Cause have been,
and will remain, beyond the reckoning of these souls.

\subsection{232}\label{section-1787}

\begin{otherlanguage}{persian}

و اگر ناس ببصر اللّه ناظر باشند همین یک بیان نقطهٴ بیان جلّ ثنائه جمیع را
کافی است که میفرمایند انّنی انا حیّ فی افق الأبهی چنانچه الیوم از افق ابهی
ظاهر و لایح و مشرقند و حینئذ یسمع و یری کلّ ما یقول و یرتکب به الخلایق
اجمعین و من دون ذلک

\end{otherlanguage}

Va agar nās be-baṣar-e Allāh nāẓer bāshand, hamīn yek bayān-e Noqṭe-ye
Bayān, jalla thanāʾuhu, jamīʿ rā kāfī ast ke mīfarmāyand: innanī anā
ḥayyun fī ufuqi al-Abhā. Chenānche al-yowm az ofoq-e Abhā ẓāher va lāyeḥ
va moshreqand, va ḥīnaʾidhin yasmaʿu wa yarā kulla mā yaqūlu wa
yartakibu bihi al-khalāʾiqu ajmaʿīn, va men dūn-e dhālek.

If people were to look with the eye of God, this single utterance of the
Point of the Bayān---exalted be His praise---would suffice them all:
``Indeed, I am alive in the horizon of the All-Glorious.'' Thus today He
is manifest, radiant, and shining from the horizon of the All-Glorious,
and He now hears and sees everything that all created beings say and do.
Besides this,

\subsection{233}\label{section-1788}

\begin{otherlanguage}{persian}

شهادتی که در الواح ذکر فرموده‌اند و جمیع را بقرائت آن امر نموده‌اند دلیل
واضح لمن کان ناظراً بعین اللّه ربّه قوله تعالی شهد اللّه انّه لا اله الّا هو و
انّ نقطة الأولی عبده و بهائه و این آیات دلیلی است واضح و برهانی است قاطع
ولکن لا یزید المتوهّمین الّا خساراً چه که متوهّمین جمیع آیات را بوهم تفسیر
نموده‌اند و هم‌چنین میفرماید عزّ ذکره و انّی انا العرش البهآء بالحقّ الأکبر
قد کنت فوق مطلع یاقوته السّیّال فوق الطّور

\end{otherlanguage}

Shahādatī ke dar alvāḥ dhekr farmūde-and va jamīʿ rā be-qerāʾat-e ān amr
namūde-and dalīl-e vāḍeḥ li-man kāna nāẓiran bi-ʿayni Allāhi rabbihi.
Qawluhu taʿālā: shahida Allāhu annahu lā ilāha illā huwa wa anna Nuqṭata
al-Ūlā ʿabduhu wa bahāʾuhu. Va īn āyāt dalīlī ast vāḍeḥ va borhānī ast
qāṭeʿ, va lāken lā yazīdu al-mutawahhimīna illā khasāran, che ke
motavahhemīn jamīʿ-e āyāt rā be-vahm tafsīr namūde-and. Va hamchenīn
mīfarmāyad, ʿazza dhikruhu: wa innī anā al-ʿarshu al-Bahāʾu bi-al-ḥaqqi
al-akbar, qad kuntu fawqa maṭlaʿi yāqūtihi al-sayyāli fawqa al-Ṭūr.

The testimony He has mentioned in the Tablets, which He has commanded
everyone to recite, is a clear proof for anyone who looks with the eye
of God, his Lord. His words---exalted is He---are: ``God bears witness
that there is no god but Him, and that the Primal Point is His servant
and His glory.'' These verses are a clear proof and a decisive
demonstration, yet they ``increase the deluded in nothing but loss,''
for the deluded have interpreted all the verses according to their own
imaginings. He also says---exalted be His remembrance: ``Indeed, I am
the throne of glory in the greatest truth. I have been above the
dawning-place of His flowing ruby, above the Mount.''

\subsection{234}\label{section-1789}

\begin{otherlanguage}{arabic}

مقصود او در این آیهٴ منزلهٴ مبارکه اگر تفکّر رود بر علوّ ظهور ارفع امنع
اقدس مطّلع می‌شوید و هم‌چنین میفرمایند قوله عزّ شأنه و لقد خلق اللّه فی حول
ذلک الباب بحوراً من مآء الاکسیر محمّراً بالدّهن الوجود و حیواناً بالثّمرة
المقصود و قدّر اللّه له سفناً من یاقوته الرّطبة الحمرآء و لا یرکب فیها الّا
اهل البهآء باذن اللّه العلیّ و هو اللّه قد کان عزیزاً و حکیما

\end{otherlanguage}

Maqṣūd-e ū dar īn āye-ye monzale-ye mobārake agar tafakkor ravad, bar
ʿolovv-e ẓohūr-e arfaʿ-e amnaʿ-e aqdas moṭṭaleʿ mīshavīd. Va hamchenīn
mīfarmāyand, qawluhu ʿazza shaʾnuhu: wa laqad khalaqa Allāhu fī ḥawli
dhālika al-bābi buḥūran min māʾi al-iksīri muḥammaran bi-al-duhni
al-wujūdi wa ḥayawānan bi-al-thamarati al-maqṣūdi, wa qaddara Allāhu
lahu sufunan min yāqūtihi al-raṭbati al-ḥamrāʾi wa lā yarkabu fīhā illā
ahlu al-Bahāʾi bi-idhni Allāhi al-ʿAliyyi wa huwa Allāhu qad kāna
ʿazīzan wa ḥakīman.

If you reflect on His intent in this blessed, revealed verse, you will
become aware of the loftiness of this most exalted, most inaccessible,
most holy Revelation. He also says---exalted be His station: ``God has
indeed created around that Gate seas of the water of the elixir,
reddened by the oil of existence and made living by the intended fruit.
God has ordained for it ships of His moist crimson ruby, and none but
the people of Bahāʾ embark in them, by the permission of God, the Most
High. He is God, ever mighty and wise.''

\subsection{235}\label{section-1790}

\begin{otherlanguage}{persian}

حال سفن یاقوتهٴ رطبهٴ حمرا بر بحر کبریا جاری و ساری است ای اهل بها منقطع
از ما سوی اللّه شده در فلک احدیّه و سفینهٴ عزّ باقیه درآئید و راکب شوید که
هر نفسی که تمسّک جست باین فلک عنداللّه از اهل اثبات و نجات و علّیّین و رضوان
مذکور و هر که تخلّف نمود از اهل نار و هلاک و سجّین و هاویه محسوب

\end{otherlanguage}

Ḥāl sofon-e yāqūte-ye raṭbe-ye ḥamrā bar baḥr-e kebrīyā jārī va sārī
ast. Ey ahl-e Bahā, monqaṭeʿ az mā sevā-ye Allāh shode dar folk-e
aḥadīye va safīne-ye ʿezz-e bāqīye darāyīd va rākeb shavīd, ke har nafsī
ke tamassok jost be-īn folk ʿendaʾllāh az ahl-e ethbāt va najāt va
ʿellīyīn va reḍvān madhkūr, va har ke takhallof namūd az ahl-e nār va
halāk va sejjīn va hāvīye maḥsūb.

Now the ships of moist crimson ruby sail and glide upon the ocean of
majesty. O people of Bahā, detach yourselves from everything but God,
enter the ark of oneness and the vessel of everlasting glory, and
embark. For every soul who clings to this ark is reckoned before God
among the people of affirmation and salvation, of the highest heaven and
paradise; and whoever stays behind is counted among the people of fire
and destruction, of the nethermost prison and the abyss.

\subsection{236}\label{section-1791}

\begin{otherlanguage}{arabic}

زینهار دین را بدنیا مدهید و یوسف عزّ احدیّه را بآلآء دنیا و آخرت مبادله
ننمائید عنقریب آنچه مشهود مفقود خواهید دید پس بکمال جهد و اجتهاد بظلّ ربّ
الایجاد بشتابید که مکمن امنی جز ظلّش نبوده و نخواهد بود

\end{otherlanguage}

Zīnhār dīn rā be-donyā madahīd va Yūsof-e ʿezz-e aḥadīye rā be-ālā-ye
donyā va ākherat mobādele nanamāyīd. ʿAnqarīb ānche mashhūd mafqūd
khwāhīd dīd. Pas be-kamāl-e jahd va ejtehād be-ẓell-e Rabb-e ījād
beshetābīd, ke makman-e amnī joz ẓellash nabūde va nakhwāhad būd.

Take care not to exchange your faith for the world, nor trade the Joseph
of the glory of oneness for the bounties of this world and the next.
Before long, you will see all that is visible vanish. Hasten, then, with
the utmost effort and endeavor to the shelter of the Lord of creation,
for there has never been, nor will there ever be, a safe refuge other
than His shelter.

\subsection{237}\label{section-1792}

\begin{otherlanguage}{arabic}

و قوله عزّ کبریائه و امّا الغلام فهو من نار الشّجرة الخضرآء الموقدة من هذا
العین الصّفرآء قد قتلناه فی هیکل المرئی لما قد قدّر اللّه فی الکتاب حظّاً من
ابویه فخشینا ان یرهقهما خلال النّار فی جمال البهآء طغیاناً بلا علماً

\end{otherlanguage}

Wa qawluhu ʿazza kibriyāʾuhu: wa ammā al-ghulāmu fa-huwa min nāri
al-shajarati al-khaḍrāʾi al-mūqadati min hādhā al-ʿayni al-ṣafrāʾi, qad
qatalnāhu fī haykali al-marʾiyyi li-mā qad qaddara Allāhu fī al-kitābi
ḥaẓẓan min abawayhi, fa-khashīnā an yurhiqahumā khilāla al-nāri fī
jamāli al-Bahāʾi ṭughyānan bi-lā ʿilman.

And His words---exalted be His majesty: ``As for the youth, he is of the
fire of the green tree, kindled from this yellow spring. We killed him
in the visible temple because of what God had ordained in the Book as a
portion from his parents, for We feared that, amid the fire, he would
burden them with rebellion against the Beauty of Bahāʾ, without
knowledge.''

\subsection{238}\label{section-1793}

\begin{otherlanguage}{persian}

حال بسی واضح است که کلّ بر این جمال اظهر انور ابهی طغیان نموده بشأنی که
جز علم اللّه احدی محصی آن نه و کاش بطغیان و استکبار کفایت مینمودند بلکه
در کلّ حین در انهدام بیت اللّه و حرمه مشغول بوده و خواهند بود و لن یسکن
قلوبهم الّا بشرب دمی بعد الّذی خلقناهم و حفظناهم لاعلآء امری و نصرة جمالی

\end{otherlanguage}

Ḥāl basī vāḍeḥ ast ke koll bar īn Jamāl-e aẓhar-e anvar-e Abhā ṭoghyān
namūde be-shaʾnī ke joz ʿelm-e Allāh aḥadī moḥṣī-ye ān na. Va kāsh
be-ṭoghyān va estekbār kefāyat mīnamūdand, balke dar koll-e ḥīn dar
enhedām-e bayt-e Allāh va ḥaramash mashghūl būde va khwāhand būd. Wa lan
yaskuna qulūbuhum illā bi-shurbi damī baʿda alladhī khalaqnāhum wa
ḥafiẓnāhum li-iʿlāʾi amrī wa nuṣrati jamālī.

Now it is abundantly clear that all have rebelled against this most
manifest, most luminous, all-glorious Beauty to a degree that no one can
reckon except through the knowledge of God. If only they had been
content with rebellion and arrogance! Rather, at every moment they have
been, and will continue to be, occupied with destroying the house of God
and His sanctuary. Their hearts will never be stilled except by drinking
My blood, after We created and protected them for the exaltation of My
Cause and the aid of My Beauty.

\subsection{239}\label{section-1794}

\begin{otherlanguage}{persian}

فلمّا اشدّ ظهرهم و وجدوا الاطمینان من انفسهم قاموا علی قتلی بطغیان عظیم
فلمّا بلغت الی ذلک المقام نادی المناد عن خلف سرادق البقآء ان یا قلم
الأعلی لا تجعل امرک محدوداً بذکر دون ذکر قل تاللّه قد نزّل البیان کلّه لأمری
البدیع المنیع و لولاک ما نزّل حرف من البیان و لا صحایف عزّ کریم دع الموتی
لأنّ الّذینهم ما آمنوا بک اموات غیر احیآء لا یسمعون و لا یعقلون ولو یلقی
علیهم کلمات الأوّلین و الآخرین و الّذین احیاهم اللّه بک اولئک فی حبّک قرؤوا
کلّ الکتب و لن یحتاجوا بشیء سواک لأنّ ظهورک بنفسه جعله اللّه حجّة علی
العالمین

\end{otherlanguage}

Fa-lammā ashadda ẓahrahum wa wajadū al-iṭmiʾnāna min anfusihim qāmū ʿalā
qatlī bi-ṭughyānin ʿaẓīm. Fa-lammā balaghtu ilā dhālika al-maqāmi nādā
al-munādi ʿan khalfi surādiqi al-baqāʾi an yā Qalama al-Aʿlā lā tajʿal
amraka maḥdūdan bi-dhikrin dūna dhikr. Qul ta-Allāhi qad nuzzila
al-Bayānu kulluhu li-amrī al-badīʿi al-manīʿi wa lawlāka mā nuzzila
ḥarfun mina al-Bayāni wa lā ṣaḥāʾifu ʿizzin karīm. Daʿi al-mawtā li-anna
alladhīnahum mā āmanū bika amwātun ghayru aḥyāʾin lā yasmaʿūna wa lā
yaʿqilūna wa law yulqā ʿalayhim kalimātu al-awwalīna wa al-ākhirīn. Wa
alladhīna aḥyāhumu Allāhu bika ulāʾika fī ḥubbika qaraʾū kulla al-kutubi
wa lan yaḥtājū bi-shayʾin siwāka li-anna ẓuhūraka bi-nafsihi jaʿalahu
Allāhu ḥujjatan ʿalā al-ʿālamīn.

When He had strengthened them and they felt assured of themselves, they
rose up to kill Me in grievous rebellion. When I reached this point, the
Herald called from behind the tabernacle of eternity: ``O Most Exalted
Pen, do not confine Your Cause to mentioning one thing rather than
another. Say: By God, the entire Bayān was revealed for My wondrous,
inaccessible Cause. But for You, not a letter of the Bayān would have
been revealed, nor any scrolls of noble glory. Leave the dead, for those
who have not believed in You are dead, not living. They neither hear nor
understand, even if the words of the first and the last are recited to
them. Those whom God has brought to life through You have, in their love
for You, read all the Books, and they will never need anything but You,
for God has made Your manifestation itself a proof to all the worlds.''

\subsection{240}\label{section-1795}

\begin{otherlanguage}{arabic}

فلمّا سمعت النّدآء اکتفیت بما رقم من اصبع قدس منیر

\end{otherlanguage}

Fa-lammā samiʿtu al-nidāʾa iktafaytu bi-mā ruqima min iṣbaʿi qudsin
munīr.

When I heard the call, I contented Myself with what had been inscribed
by a luminous finger of holiness.

\subsection{241}\label{section-1796}

\begin{otherlanguage}{persian}

باری آیات عزّ احدیّه در این اسم و ظهور زیاده از آن است که احصا توان نمود و
مع‌ذلک این همج رعاع این اسم مبارک را بارض طاء تفسیر نموده‌اند چه که حضرت
اعلی روح ما سواه فداه او را ارض بها نامیده

\end{otherlanguage}

Bārī āyāt-e ʿezz-e aḥadīye dar īn esm va ẓohūr zeyāde az ān ast ke eḥṣā
tavān namūd, va maʿa-dhālek īn hamaj-e raʿāʿ īn esm-e mobārak rā
be-arḍ-e Ṭā tafsīr namūde-and, che ke Ḥaḍrat-e Aʿlā, rūḥ-e mā sevāh
fedāh, ū rā arḍ-e Bahā nāmīde.

In short, the verses of the glory of oneness concerning this Name and
Revelation are beyond counting. Yet these ignorant rabble have
interpreted this blessed Name as referring to the land of Ṭā, because
the Most Exalted One---may the soul of all else be sacrificed for
Him---called it the land of Bahā.

\subsection{242}\label{section-1797}

\begin{otherlanguage}{persian}

و آنقدر شعور نداشته که ادراک نمایند که مقصود حضرت آن بوده که اخبار
فرماید از ظهور جمال مقدّس در آن ارض ای ارض یمشی علیه البهآء و سکن فیها
هذا الاسم المشرق المنیر چنانچه این اسم مبارک را در کلّ مراتب اسبق از کلّ
اسماء ذکر فرموده‌اند و

\end{otherlanguage}

Va ānqadr shoʿūr nadāshte ke edrāk namāyand ke maqṣūd-e Ḥaḍrat ān būde
ke ekhbār farmāyad az ẓohūr-e Jamāl-e moqaddas dar ān arḍ, ay arḍun
yamshī ʿalayhi al-Bahāʾu wa sakana fīhā hādhā al-ismu al-mushriqu
al-munīr. Chenānche īn esm-e mobārak rā dar koll-e marāteb asbaq az
koll-e asmā dhekr farmūde-and va

They have lacked even the understanding to perceive that His purpose was
to foretell the appearance of the holy Beauty in that land---that is, a
land upon which Bahāʾ walks and in which this radiant, luminous Name has
dwelt. Thus He has mentioned this blessed Name as preceding all names in
every degree, and

\subsection{243}\label{section-1798}

\begin{otherlanguage}{persian}

هم‌چنین در دعای قبل ملاحظه نمائید که اقدم اسماء این اسم مبارک ذکر شده
بقوله اللّهمّ انّی اسألک من بهائک بأبهاه و هم‌چنین شهور که باسماء جدیده
نامیده‌اند اوّل آن را باین اسم منسوب فرموده‌اند در هر شیء اگر ملاحظه نمائید
احسن آن را باین اسم مذکور نموده‌اند

\end{otherlanguage}

Hamchenīn dar doʿā-ye qabl molāḥeẓe namāyīd ke aqdam-e asmā īn esm-e
mobārak dhekr shode be-qawlihi: Allāhumma innī asʾaluka min bahāʾika
bi-abhāhu. Va hamchenīn shohūr ke be-asmā-ye jadīde nāmīde-and, avval-e
ān rā be-īn esm mansūb farmūde-and. Dar har shayʾ agar molāḥeẓe namāyīd,
aḥsan-e ān rā be-īn esm madhkūr namūde-and.

likewise, consider the earlier prayer, in which this blessed Name is
mentioned as the foremost of the names, in His words: ``O God, I ask You
by the most glorious of Your glory.'' Likewise, among the months to
which He has given new names, He has associated the first with this
Name. Whatever you examine, He has designated its finest part by this
Name.

\subsection{244}\label{section-1799}

\begin{otherlanguage}{arabic}

و مع‌ذلک ملاحظه نمائید که بغضا بچه رتبه رسیده که این اسم مبارک را بارضی
تفسیر نموده‌اند که ذرّة من ترابها خیر منهم و من ذواتهم و حقایقهم فأفّ لهم و
بما فی صدورهم من نار الحسد و البغضآء و یحترقون بها فی کلّ حین و لا یشعرون

\end{otherlanguage}

Va maʿa-dhālek molāḥeẓe namāyīd ke baghḍā be-che rotbe resīde ke īn
esm-e mobārak rā be-arḍī tafsīr namūde-and ke dharratun min turābihā
khayrun minhum wa min dhawātihim wa ḥaqāʾiqihim. Fa-uffin lahum wa bi-mā
fī ṣudūrihim min nāri al-ḥasadi wa al-baghḍāʾi wa yaḥtariqūna bihā fī
kulli ḥīnin wa lā yashʿurūn.

Yet consider how far their hatred has gone: they have interpreted this
blessed Name as referring to a land, a speck of whose dust is better
than they are, than their very selves and realities. Shame on them and
on the fire of envy and hatred within their breasts! They burn in it at
every moment, yet do not perceive it.

\subsection{245}\label{section-1800}

\begin{otherlanguage}{persian}

و دیگر جناب قدّوس جلّ ثنائه الواحی چند در ارض بدشت مرقوم فرموده و ارسال
داشته و ابداً این عبد اظهار ننموده و از جمله در این امر بدیع اخبار
فرموده‌اند تصریحاً من غیر تلویح قوله جلّ ثنائه و اذا اظهر الرّبّ سرّاً من افق
البهآء فی ارض او ادنی فقد کان ذلک الطّلعة المتلامعة من نقطة البآء طرزیّاً و
اذا قامت السّموات الجذبیّة بالسّطر بالسّرّ السّطریّة فذلک من امرنا لأهل العمآء
قد کان طلیعاً و آن الواح مبارکه الآن موجود است و جمیع در ذکر این امر
باصرح کلمات شاهد و گواه فوربّ العرش و العمآء که حیا مانع است از ذکر این
کلمات چه که این عبد لازال اراده ننموده که بکلمات قبل شأنی از برای خود
ثابت نماید لأنّ شأنی شأنه و انّه مشهود فی وسط سمآء الاستقلال و ما من ذی
بصر الّا و قد یشهد و یری و للعمیآء لیس له من نصیب

\end{otherlanguage}

Va dīgar Janāb-e Qoddūs, jalla thanāʾuhu, alvāḥī chand dar arḍ-e Badasht
marqūm farmūde va ersāl dāshte va abadan īn ʿabd eẓhār nanamūde, va az
jomle dar īn amr-e badīʿ ekhbār farmūde-and taṣrīḥan min ghayri talwīḥ.
Qawluhu jalla thanāʾuhu: wa idhā aẓhara al-rabbu sirran min ufuqi
al-Bahāʾi fī arḍi aw adnā fa-qad kāna dhālika al-ṭalʿatu al-mutalāmiʿatu
min nuqṭati al-bāʾi ṭarziyyan, wa idhā qāmati al-samāwātu al-jadhbiyyatu
bi-al-saṭri bi-al-sirri al-saṭriyyati fa-dhālika min amrinā li-ahli
al-ʿamāʾi qad kāna ṭalīʿan. Va ān alvāḥ-e mobārake al-ān mowjūd ast va
jamīʿ dar dhekr-e īn amr be-aṣraḥ-e kalemāt shāhed va govāh. Fa-wa-rabbi
al-ʿarshi wa al-ʿamāʾi ke ḥayā māneʿ ast az dhekr-e īn kalemāt, che ke
īn ʿabd lāzāl erāde nanamūde ke be-kalemāt-e qabl shaʾnī az barā-ye
khwod thābet namāyad, li-anna shaʾnī shaʾnuhu wa innahu mashhūdun fī
wasaṭi samāʾi al-istiqlāli wa mā min dhī baṣarin illā wa qad yashhadu wa
yarā wa li-al-ʿamyāʾi laysa lahu min naṣīb.

Furthermore, his honor Quddūs---exalted be his praise---wrote and sent
several Tablets in the land of Badasht, which this servant has never
disclosed. Among other things, he foretold this wondrous Cause
explicitly, without allusion. His words---exalted be his praise---are:
``When the Lord manifests a secret from the horizon of Bahāʾ in the land
of `or nearer,' that resplendent Countenance from the point of the
letter Bāʾ has indeed been an exemplar. And when the heavens of
attraction are raised through the line, through the mystery of the line,
that has indeed been a harbinger of Our Cause to the people of the cloud
of unknowing.'' Those blessed Tablets are now extant, and all bear
witness and testify to this Cause in the most explicit words. By the
Lord of the Throne and the cloud of unknowing, shame prevents Me from
mentioning these words, for this servant has never wished to establish a
station for Himself through earlier utterances. For My station is His
station, and it is manifest in the midmost heaven of independence. There
is no one endowed with sight who does not witness and behold it, while
the blind have no share in it.

\subsection{246}\label{section-1801}

\begin{otherlanguage}{arabic}

قل تاللّه شأنی شأنه و قدرتی قدرته و سلطانی سلطانه و ظهوره حجّتی و ما جری
من قلمه برهانی کما انّ ظهوری حجّته و ما جری من قلمی برهانه ولکن نظر بحفظ
عباد نابالغ تحریر میشود

\end{otherlanguage}

Qul ta-Allāhi shaʾnī shaʾnuhu wa qudratī qudratuhu wa sulṭānī sulṭānuhu
wa ẓuhūruhu ḥujjatī wa mā jarā min qalamihi burhānī, kamā anna ẓuhūrī
ḥujjatuhu wa mā jarā min qalamī burhānuhu. Va lāken naẓar be-ḥefẓ-e
ʿebād-e nābālegh taḥrīr mīshavad.

Say: By God, My station is His station, My power is His power, and My
sovereignty is His sovereignty. His manifestation is My proof, and what
has flowed from His Pen is My testimony, just as My manifestation is His
proof, and what has flowed from My Pen is His testimony. Yet this is
being written for the protection of immature servants.

\subsection{247}\label{section-1802}

\begin{otherlanguage}{persian}

و از جمله نقطهٴ جذبیّه جناب طاء علیها بهآء اللّه الأبهی مدّتها با این عبد
بوده و آنی لقاء این غلام را بملک دنیا و آخرت مبادله نمینمود و ما ارادت
ان تفارق عنّی اقلّ من آن ولکن قضی ما قضی و چه مقدار از آیات و اشعار که در
این امر بدیع ذکر فرموده از جمله در وصف طلعت الهی غزلی گفته و یک فرد آن
این است

\end{otherlanguage}

Va az jomle Noqṭe-ye Jadhbīye, Janāb-e Ṭā, ʿalayhā bahāʾu Allāhi
al-Abhā, moddat-hā bā īn ʿabd būde va ānī leqā-ye īn gholām rā be-molk-e
donyā va ākherat mobādele namīnamūd, wa mā arādat an tufāriqa ʿannī
aqalla min ān, va lāken quḍiya mā quḍiya. Va che meqdār az āyāt va
ashʿār ke dar īn amr-e badīʿ dhekr farmūde. Az jomle dar vaṣf-e Ṭalʿat-e
elāhī ghazalī gofte va yek fard-e ān īn ast.

Among them, too, was the Point of Attraction, her honor Ṭā---upon her be
the glory of God, the All-Glorious---who was with this servant for a
long time. She would not have exchanged a moment in the presence of this
Youth for dominion over this world and the next, nor did she wish to
part from Me for less than an instant. Yet what was decreed came to
pass. How many verses and poems she uttered concerning this wondrous
Cause! Among them, she composed a ghazal in praise of the divine
Countenance, one couplet of which is this:

\subsection{248}\label{section-1803}

\begin{otherlanguage}{arabic}

گر براندازد بها از رخ نقاب صدهزار همچون ازل آید پدید

\end{otherlanguage}

Gar barandāzad Bahā az rokh neqāb ṣad-hezār hamchūn Azal āyad padīd.

Should Bahā cast the veil from His face, a hundred thousand like Azal
would appear.

\subsection{249}\label{section-1804}

\begin{otherlanguage}{persian}

و هم‌چنین در دون این عبد هم بسیار وصف نموده و فوق جمیع این بیانات کتاب
مکنون نقطهٴ اعلی روح ما سواه فداه که کلّ در این امر بدیع نازل شده و از
انظار مستور گشته از خدا بخواهید تا بآن فایز شوید

\end{otherlanguage}

Va hamchenīn dar dūn-e īn ʿabd ham besyār vaṣf namūde, va fowq-e jamīʿ-e
īn bayānāt Ketāb-e Maknūn-e Noqṭe-ye Aʿlā, rūḥ-e mā sevāh fedāh, ke koll
dar īn amr-e badīʿ nāzel shode va az anẓār mastūr gashte. Az Khodā
bekhwāhīd tā be-ān fāyez shavīd.

She likewise offered much praise concerning others besides this servant.
Above all these utterances stands the Hidden Book of the Most Exalted
Point---may the soul of all else be sacrificed for Him---which was
revealed entirely concerning this wondrous Cause and has been concealed
from view. Ask God that you may attain it.

\subsection{250}\label{section-1805}

\begin{otherlanguage}{persian}

فوالّذی نفسی بیده که اگر مدّعی جمال الهی الیوم از این گونه دلایل در کتاب
خود مذکور نمیداشت و تمسّک نمیجست که باین ادلّه سبب اضلال خلق شود هرگز این
عبد باین ادلّه استدلال نمینمود باری ای عباد ندای علیّ اعلی را در ظهور
اخرای او بشنوید و جمیع این دلایل را از قلب محو نمائید چه که اگر بکلمات و
اشارات و صحف و کتب در این ظهور بدع ربّانی تمسّک جوئید ابداً بجوهر عرفان
طلعت رحمن و سلطان امکان فایز نخواهید شد

\end{otherlanguage}

Fa-wa-alladhī nafsī bi-yadihi ke agar moddaʿī-ye Jamāl-e elāhī al-yowm
az īn gūne dalāyel dar ketāb-e khwod madhkūr namīdāsht va tamassok
namījost ke be-īn adelle sabab-e eḍlāl-e khalq shavad, hargez īn ʿabd
be-īn adelle estedlāl namīnamūd. Bārī ey ʿebād, nedā-ye ʿAlīy-e Aʿlā rā
dar ẓohūr-e okhrā-ye ū beshnavīd va jamīʿ-e īn dalāyel rā az qalb maḥv
namāyīd, che ke agar be-kalemāt va eshārāt va ṣoḥof va kotob dar īn
ẓohūr-e bedʿ-e rabbānī tamassok jūyīd, abadan be-jowhar-e ʿerfān-e
Ṭalʿat-e Raḥmān va Solṭān-e emkān fāyez nakhwāhīd shod.

By Him in whose hand is My soul, had the one laying claim to the divine
Beauty not cited such proofs in his book today and clung to them,
thereby making these arguments a means of leading people astray, this
servant would never have adduced these proofs. In short, O servants,
hear the call of the Most Exalted ʿAlī in His latter manifestation, and
erase all these proofs from your hearts. For if you cling to words,
allusions, scrolls, and books in this wondrous divine Revelation, you
will never attain the essence of the recognition of the Countenance of
the All-Merciful and the Sovereign of all existence.

\subsection{251}\label{section-1806}

\begin{otherlanguage}{persian}

بلکه ناظر باشید بهمان حجّتی که نفس نقطهٴ اولی جلّ و عزّ بآن اتیان فرمود و
امر خود را ثابت نمود و هم‌چنین از قبل آن جمال قدم محمّداً رسول اللّه و من
قبله ابن مریم و من قبله کلیم فوالّذی نفسی بیده که یک حرف از آیات اللّه اعزّ
است نزد این عبد از کلّ من فی السّموات و الأرض

\end{otherlanguage}

Balke nāẓer bāshīd be-hamān ḥojjatī ke nafs-e Noqṭe-ye Ūlā jalla wa
ʿazza be-ān etyān farmūd va amr-e khod rā thābet namūd va ham-chenīn az
qabl ān Jamāl-e Qedam Muḥammadan Rasūl Allāh wa min qablihi Ibn Maryam
wa min qablihi Kalīm. Fa-wa-lladhī nafsī bi-yadihi ke yek ḥarf az āyāt-e
Allāh aʿazz ast nazd-e īn ʿabd az kull man fī al-samāwāt wa-l-arḍ.

Rather, look to that very proof which the Primal Point
Himself---glorified and exalted be He---brought forth and by which He
established His Cause, as did before Him that Ancient Beauty, Muḥammad,
the Messenger of God, and before Him the Son of Mary, and before Him
Moses. By Him in whose hand is My soul, a single letter of the verses of
God is dearer to this servant than all who are in the heavens and on
earth.

\subsection{252}\label{section-1807}

\begin{otherlanguage}{persian}

و اینکه ذکر شده که در حین ظهور ناظر بکلمات و اشارات و دلالات و اسماء و
صفات نشوید این حکمی است که خود نقطهٴ اولی روح ما سواه فداه فرموده و
مقصود آن جمال الهی آنکه مباد در حین ظهور اهل بیان بمثل اهل فرقان باحادیث
و آیات تمسّک جویند و بمظهر و موجد و منزل آن اعتراض نمایند

\end{otherlanguage}

Va īnke dhekr shode ke dar ḥīn-e ẓohūr nāẓer be-kalamāt va eshārāt va
dalālāt va asmāʾ va ṣefāt nashavīd īn ḥokmī ast ke khod-e Noqṭe-ye Ūlā
rūḥu mā siwāhu fidāhu farmūde va maqṣūd-e ān Jamāl-e Elāhī ānke mabād
dar ḥīn-e ẓohūr ahl-e Bayān be-methl-e ahl-e Forqān be-aḥādīth va āyāt
tamassok jūyand va be-maẓhar va mūjed va monzel-e ān eʿterāż namāyand.

As for the statement that, at the time of the Revelation, you should not
look to words, allusions, indications, names, and attributes: this is a
command given by the Primal Point Himself---may the soul of all else be
sacrificed for Him. The intention of that divine Beauty was that, at the
time of the Revelation, the people of the Bayān should not, like the
people of the Qurʾān, cling to traditions and verses and object to the
One who manifests them, brings them into being, and reveals them.

\subsection{253}\label{section-1808}

\begin{otherlanguage}{persian}

باری اگر قدری در این ظهور ابدع امنع تفکّر نمائید بر اسرار لایحصی مطّلع
شوید و مشاهده نمائید که از اوّل ابداع تا حین چنین ظهوری ظاهر نشده اگرچه
نظر بعدم استعداد عباد اسباب منع هم موجود لیفصل به الجوهر عن العرض و
الرّوح عن الطّین

\end{otherlanguage}

Bārī agar qadrī dar īn ẓohūr-e abdaʿ-e amnaʿ tafakkor namāyīd bar
asrār-e lā yoḥṣā moṭṭaleʿ shavīd va moshāhede namāyīd ke az avval-e
ebdāʿ tā ḥīn chenīn ẓohūrī ẓāher nashode agarche naẓar be-ʿadam-e
esteʿdād-e ʿebād asbāb-e manʿ ham mowjūd li-yafṣila bihi al-jawhar ʿan
al-ʿaraḍ wa-l-rūḥ ʿan al-ṭīn.

In short, if you reflect a little on this most wondrous, most exalted
Revelation, you will become acquainted with countless mysteries and see
that, from the beginning of creation until now, no such Revelation has
appeared. Yet, in view of the servants' lack of readiness, obstacles are
also present, so that through this He may separate the substance from
the accident and the spirit from the clay.

\subsection{254}\label{section-1809}

\begin{otherlanguage}{persian}

تاللّه الحقّ الیوم اگر ذرّه‌ئی از جوهر در صدهزار من سنگ مخلوط باشد و در خلف
سبعة ابحر مستور هرآینه دست قدرت الهی او را ظاهر فرماید و آن ذرّهٴ جوهر را
از او فصل نماید ای قوم نسمة اللّه در هبوب است و روح اللّه در اهتزاز و جذب
اللّه از افق مشرق جمال رحمن طالع

\end{otherlanguage}

Ta-llāh al-Ḥaqq, al-yawm agar dharre-ī az jowhar dar ṣad-hezār man sang
makhlūṭ bāshad va dar khalf-e sabʿat abḥur mastūr har-āyene dast-e
qodrat-e Elāhī ū rā ẓāher farmāyad va ān dharre-ye jowhar rā az ū faṣl
namāyad. Ey qowm nasmat Allāh dar hobūb ast va rūḥ Allāh dar eht ezāz va
jadhb Allāh az ofoq-e mashreq-e Jamāl-e Raḥmān ṭāleʿ.

By God, the True One! Today, if a particle of a gem were mixed with a
hundred thousand maunds of stone and concealed beyond seven seas, the
hand of divine power would assuredly bring it to light and separate that
particle of the gem from the stone. O people, the breeze of God is
blowing, the Spirit of God is stirring, and the attraction of God is
dawning from the horizon of the dayspring of the Beauty of the
All-Merciful.

\subsection{255}\label{section-1810}

\begin{otherlanguage}{persian}

سر از نوم نفس و هوی بردارید و مشاهده کنید که چه ظاهر شده در ابداع زینهار
گوش بمزخرفات قوم مدهید چه که الیوم اگر سمع طاهر یافت شود میشنود از این
طایفه آنچه از اهل فرقان می‌شنید و آنچه از اهل توریة و انجیل استماع مینمود
و من‌دون اقوال بلکه همان افعال را بعینه ببصر ظاهر و باطن مشاهده مینماید

\end{otherlanguage}

Sar az nowm-e nafs va havā bardārīd va moshāhede konīd ke che ẓāher
shode dar ebdāʿ. Zīnhār gūsh be-mozakhrafāt-e qowm madahīd che ke
al-yawm agar samʿ-e ṭāher yāft shavad mīshenavad az īn ṭāyefe ānche az
ahl-e Forqān mīshenīd va ānche az ahl-e Towrāt va Enjīl estemāʿ mīnamūd
va men-dūn-e aqvāl balke hamān afʿāl rā be-ʿayneh be-baṣar-e ẓāher va
bāṭen moshāhede mīnamāyad.

Raise your heads from the sleep of self and desire, and behold what has
appeared in creation. Beware of lending an ear to the people's idle
words, for today, if a pure ear is found, it will hear from this group
what it used to hear from the people of the Qurʾān and what it heard
from the people of the Torah and the Gospel. Indeed, beyond words, it
will witness with both outward and inward sight the very same deeds.

\subsection{256}\label{section-1811}

\begin{otherlanguage}{persian}

چنانچه الیوم کسی که بنفس حقّ محاربه نموده و استکبر علی من خلقه در کلماتش
مذکور داشته که اگر حقّ ظاهر شود چه از اهل بیان و چه غیر آن چه عبد و چه حرّ
چه سیاه و چه سفید این عبد مطیع است حرفی ندارم حال ملاحظه نمائید که این
سخن را در وقتی میگوید که در همان حین بمحاربه با جمال اللّه مشغول است و در
کمال اعراض معرض حال شما تفکّر کنید و ملاحظه نمائید که حجّت حقّ چه بوده و
علامت ظهور حقّ چه

\end{otherlanguage}

Chenānche al-yawm kasī ke be-nafs-e Ḥaqq moḥārebe namūde wa istakbara
ʿalā man khalaqahu dar kalamātash madhkūr dāshte ke agar Ḥaqq ẓāher
shavad che az ahl-e Bayān va che gheyr-e ān che ʿabd va che ḥorr che
sīyāh va che sefīd īn ʿabd moṭīʿ ast ḥarfī nadāram. Ḥāl molāḥeẓe namāyīd
ke īn sokhan rā dar vaqtī mīgūyad ke dar hamān ḥīn be-moḥārebe bā Jamāl
Allāh mashghūl ast va dar kamāl-e eʿrāż moʿreż. Ḥāl shomā tafakkor konīd
va molāḥeẓe namāyīd ke ḥojjat-e Ḥaqq che būde va ʿalāmat-e ẓohūr-e Ḥaqq
che.

Thus today, the one who has waged war against the True One Himself and
arrogantly defied Him who created him has stated in his writings: ``If
the True One appears, whether from among the people of the Bayān or
others, whether slave or free, black or white, this servant is obedient;
I have no objection.'' Now consider that he says this at the very time
he is engaged in warfare against the Beauty of God and has turned away
in utter rejection. Now reflect and consider what the proof of the True
One has been and what the sign of His appearance is.

\subsection{257}\label{section-1812}

\begin{otherlanguage}{persian}

این ظهور باین مشهودی که انوارش عالم را احاطه نموده از آن معرض و بر قتلش
ایستاده و بعد باین کلمات مشغول شده لأجل توهّم ناس مثل اوّل ظهور که علمای
فرقان سبحه‌ها بر دست و در کمال خضوع و خشوع در معابد بعبادت مشغول و بشأنی
اظهار زهد مینمودند که چون اسم حقّ مذکور میشد برمیخاستند و بکمال خضوع
تواضع مینمودند و در همان حین تواضع و خضوع بر قتلش فتوی میدادند و حکم دون
ایمان در حقّش جاری مینمودند و بزبان متکلّم که چگونه میشود ما معرض از حقّ
باشیم و حال آنکه شبها آمل بوده‌ایم و روزها از حقّ سائلیم هر وقت ظاهر شود
مطیعیم و منقاد

\end{otherlanguage}

Īn ẓohūr be-īn mashhūdī ke anvārash ʿālam rā eḥāṭe namūde az ān moʿreż
va bar qatl ash īstāde va baʿd be-īn kalamāt mashghūl shode li-ajl-e
tavahhom-e nās methl-e avval-e ẓohūr ke ʿolamā-ye Forqān sobḥe-hā bar
dast va dar kamāl-e khożūʿ va khoshūʿ dar maʿābed be-ʿebādat mashghūl va
be-shaʾnī eẓhār-e zohd mīnamūdand ke chon esm-e Ḥaqq madhkūr mīshod
barmīkhāstand va be-kamāl-e khożūʿ tavāżoʿ mīnamūdand va dar hamān ḥīn-e
tavāżoʿ va khożūʿ bar qatl ash fatvā mīdādand va ḥokm-e dūn-e īmān dar
ḥaqqash jārī mīnamūdand va be-zabān motakallem ke chegūne mīshavad mā
moʿreż az Ḥaqq bāshīm va ḥāl ānke shab-hā āmel būde-īm va rūz-hā az Ḥaqq
sāʾelīm har vaqt ẓāher shavad moṭīʿīm va monqād.

He turns away from this Revelation, so plainly manifest that its light
has encompassed the world, and has risen to put Him to death; then he
occupies himself with these words to delude the people. It was the same
at the beginning of the Revelation, when the divines of the Qurʾān,
prayer beads in hand, were occupied with worship in their temples in the
utmost humility and reverence. They made such a display of piety that,
whenever the name of the True One was mentioned, they would rise and bow
with the utmost humility. Yet in that very state of humility and
submission, they issued decrees for His death and pronounced Him an
unbeliever, while saying: ``How could we turn away from the True One,
when by night we have awaited Him with hope and by day we beseech Him?
Whenever He appears, we are obedient and submissive.''

\subsection{258}\label{section-1813}

\begin{otherlanguage}{arabic}

حال معرض باللّه فواللّه الّذی لا اله الّا هو رأس مظهر نفس حقّ را بسکّین بغی و
فحشا و سیف غلّ و بغضا قطع مینماید و در آن حین با گردن خاضع باین کلمه ناطق
که حقّ هر وقت ظاهر شود من مطیعم

\end{otherlanguage}

Ḥāl moʿreż bi-llāh fa-wa-llāh alladhī lā ilāha illā huwa raʾs-e maẓhar-e
nafs-e Ḥaqq rā be-sekkīn-e baghy va faḥshā va seyf-e ghell va baghżā
qaṭʿ mīnamāyad va dar ān ḥīn bā gardan-e khāżeʿ be-īn kalame nāṭeq ke
Ḥaqq har vaqt ẓāher shavad man moṭīʿam.

Now the one who turns away from God---by God, besides whom there is no
other god---severs the head of the Manifestation of the True One Himself
with the knife of iniquity and indecency and the sword of malice and
hatred, while at that very moment, with bowed head, he utters these
words: ``Whenever the True One appears, I am obedient.''

\subsection{259}\label{section-1814}

\begin{otherlanguage}{persian}

فواللّه یا قوم شما را نابالغ یافته که باینگونه هذیانات تکلّم مینماید و
اینقدر شما را متوهّم یافته که حقّی که بجمیع ظهورات ظاهر شده فتوی بر قتلش
داده و در کلّ حین بمحاربه با جمالش مشغول و مع‌ذلک باین اقوال تکلّم نموده
حسرةً و اسفاً علیکم بما غفلتم عن الّذی کان اظهر من کلّ ظهور و اشتغلتم بالّذی
فواللّه لو کنتم مطّلعاً علیه لفررتم منه الف منازل

\end{otherlanguage}

Fa-wa-llāh yā qawm shomā rā nābālegh yāfte ke be-īngūne hadhayānāt
takallom mīnamāyad va īnqadr shomā rā motavahhem yāfte ke Ḥaqqī ke
be-jamīʿ-e ẓohūrāt ẓāher shode fatvā bar qatl ash dāde va dar koll-e ḥīn
be-moḥārebe bā jamālash mashghūl va maʿa-dhālek be-īn aqvāl takallom
namūde. Ḥasratan wa asafan ʿalaykum bi-mā ghafaltum ʿan alladhī kāna
aẓhar min kull ẓuhūr wa ishtaghaltum bi-alladhī fa-wa-llāh law kuntum
muṭṭaliʿan ʿalayhi la-farartum minhu alf manāzil.

By God, O people, he has found you immature, and so he speaks such
nonsense. He has found you so deluded that he has decreed the death of
the True One who has appeared with every manifestation, and is at all
times engaged in warfare against His Beauty, yet still speaks these
words. Alas, and woe to you, that you have remained heedless of Him who
is more manifest than every manifestation and have occupied yourselves
with one from whom---by God---if you knew him, you would flee a thousand
stages.

\subsection{260}\label{section-1815}

\begin{otherlanguage}{arabic}

و نسأل اللّه بأن یفصل بینی و بینه و یطهّر ارض الّتی اکون علیها عن رجس هؤلآء
الفاسقین او یسکننی فی دیار اخری و انّه لذو فضل علی العباد و انّه لعلی کلّ
شیء قدیر

\end{otherlanguage}

Wa nasʾal Allāh bi-an yafṣila baynī wa baynahu wa yuṭahhira arḍ allatī
akūnu ʿalayhā ʿan rijs hāʾulāʾ al-fāsiqīn aw yuskinanī fī diyār ukhrā wa
innahu la-dhū faḍl ʿalā al-ʿibād wa innahu la-ʿalā kull shayʾ qadīr.

And We ask God to separate Me from him and to cleanse the land on which
I dwell of the defilement of these transgressors, or to settle Me in
another land. He is indeed bountiful to His servants, and He has power
over all things.

\subsection{261}\label{section-1816}

\begin{otherlanguage}{persian}

بگو ای جاهل نفس محبوب ظاهر است و چون شمس لایح و تو بعرفان او الیوم فایز
نشده بلکه او را با جمیع حجج و براهین که عالم را احاطه نموده از اهل ایمان
نمیشمری مع‌ذلک بکلماتش تمسّک جسته اراده نموده که ریاستی برای نفس خود ثابت
کنی تاللّه الحقّ اذاً یکذّبکم کلّ الأشیآء ولکن انتم لا تشعرون مثل شما عنداللّه
مثل اهل فرقان است که بقول رسول ریاست خود را ثابت نموده‌اند و بر اعراش عزّت
و حکم جالس و بنفس او وارد آورده‌اند آنچه را که هیچ بصری در ابداع ندیده و
ادراک ننموده ای مست خمر نفس و هوی از سلطان معلوم چشم پوشیده و بموهوم خود
تمسّک جسته همین ذلّت تو را کافی است که انکار نموده آیاتی را که بآن دینت
ثابت شده و باو افتخار مینمائی و از برای خود شأن ثابت میکنی

\end{otherlanguage}

Begū ey jāhel nafs-e Maḥbūb ẓāher ast va chon shams lāyeḥ va to
be-ʿerfān-e ū al-yawm fāyez nashode balke ū rā bā jamīʿ-e ḥojaj va
barāhīn ke ʿālam rā eḥāṭe namūde az ahl-e īmān nemīshomorī maʿa-dhālek
be-kalamātash tamassok joste erāde namūde ke rīyāsatī barā-ye nafs-e
khod thābet konī. Ta-llāh al-Ḥaqq idhan yukadhdhibukum kull al-ashyāʾ wa
lākin antum lā tashʿurūn. Methl-e shomā ʿenda-llāh methl-e ahl-e Forqān
ast ke be-qowl-e Rasūl rīyāsat-e khod rā thābet namūde-and va bar
aʿrāsh-e ʿezzat va ḥokm jāles va be-nafs-e ū vāred āvarde-and ānche rā
ke hīch baṣarī dar ebdāʿ nadīde va edrāk nanamūde. Ey mast-e khamr-e
nafs va havā az Solṭān-e maʿlūm chashm pūshīde va be-mowhūm-e khod
tamassok joste hamīn dhellat to rā kāfī ast ke enkār namūde āyātī rā ke
be-ān dīnat thābet shode va be-ū eftekhār mīnamāʾī va az barā-ye khod
shaʾn thābet mīkonī.

Say: O ignorant one, the Beloved Himself is manifest and shines like the
sun, yet today you have not attained recognition of Him. Rather, despite
all the proofs and testimonies that have encompassed the world, you do
not count Him among the faithful. Nevertheless, you cling to His words,
seeking to establish leadership for yourself. By God, the True One, all
things then declare you false, but you do not perceive it. In God's
sight, you are like the people of the Qurʾān, who have established their
leadership through the words of the Messenger and sit upon thrones of
honor and authority, yet have inflicted on His very person what no eye
in creation has ever seen or perceived. O you who are drunk with the
wine of self and desire! You have turned your eyes from the manifest
Sovereign and clung to your own imagined one. This abasement alone is
enough for you: you have denied the verses by which your religion was
established, in which you take pride, and through which you claim a
station for yourself.

\subsection{262}\label{section-1817}

\begin{otherlanguage}{persian}

و مع‌ذلک شاعر نیستی اذاً لم یکن لک شأن عند اللّه الّا کشأن الّذینهم اعرضوا
بعدما آمنوا و انکروا بعد الّذی اعترفوا ان انت من العارفین تاللّه الحقّ عرّ
روح الأمین رأسه عن فعلک و انّک ما استشعرت بذلک و کنت من الغافلین و جلسن
حوریّات الغرفات علی الرّماد من ظلمک و انّک تکون فی نفسک من الفرحین

\end{otherlanguage}

Va maʿa-dhālek shāʿer nīstī. Idhan lam yakun laka shaʾn ʿinda Allāh illā
ka-shaʾn alladhīna-hum aʿraḍū baʿdamā āmanū wa ankarū baʿda alladhī
iʿtarafū in anta min al-ʿārifīn. Ta-llāh al-Ḥaqq ʿarra Rūḥ al-Amīn
raʾsahu ʿan fiʿlika wa innaka mā istashʿarta bi-dhālika wa kunta min
al-ghāfilīn wa jalasna ḥūriyyāt al-ghurafāt ʿalā al-ramād min ẓulmika wa
innaka takūnu fī nafsika min al-fariḥīn.

Yet you remain unaware. Thus your station in God's sight is no different
from that of those who turned away after believing and denied after
acknowledging, if you are among those who understand. By God, the True
One, the Faithful Spirit has bared His head because of your deed, yet
you have not perceived it and remain among the heedless. The maidens of
the heavenly chambers have sat upon ashes because of your injustice,
while within yourself you rejoice.

\subsection{263}\label{section-1818}

\begin{otherlanguage}{arabic}

تاللّه الحقّ ینوح کلّ شیء فی نفسه و یبکی ولکن انّک غفلت و کنت من المعرضین و
انت الّذی تذکر المحبوب بلسانک لتشتبه علی النّاس و تقتل محبوب الأوّلین و
الآخرین لو یفتح اللّه بصرک لتشهد بأنّ بظلمک قد علّق المحبوب فی الهوآء و
ترمیه فی کلّ حین برمی الحسد و البغضآء ثمّ بسهام الغلّ و العناد ثمّ برصاص
الاعراض و کان اللّه علی ذلک لشهید و علیم

\end{otherlanguage}

Ta-llāh al-Ḥaqq yanūḥu kull shayʾ fī nafsihi wa yabkī wa lākin innaka
ghafalta wa kunta min al-muʿriḍīn wa anta alladhī tadhkuru al-Maḥbūb
bi-lisānika li-tashtabiha ʿalā al-nās wa taqtulu Maḥbūb al-awwalīn
wa-l-ākhirīn. Law yaftaḥ Allāh baṣaraka la-tashhada bi-anna bi-ẓulmika
qad ʿulliqa al-Maḥbūb fī al-hawāʾ wa tarmīhi fī kull ḥīn bi-ramy
al-ḥasad wa-l-baghḍāʾ thumma bi-sihām al-ghill wa-l-ʿinād thumma
bi-raṣāṣ al-iʿrāḍ wa kāna Allāh ʿalā dhālika la-shahīd wa ʿalīm.

By God, the True One, everything laments and weeps within itself, but
you have remained heedless and are among those who turn away. You are
the one who mentions the Beloved with your tongue to mislead the people,
while killing the Beloved of the first and the last. If God were to open
your eyes, you would see that through your injustice the Beloved has
been suspended in the air, and that at every moment you assail Him with
the projectiles of envy and hatred, then with the arrows of malice and
obstinacy, then with the bullets of rejection. God is indeed a witness
to this and knows it.

\subsection{264}\label{section-1819}

\begin{otherlanguage}{arabic}

اذاً ینوح محمّد فی الأفق الأعلی و یبکی الرّوح فی الرّفیق الأبهی ثمّ الکلیم
عند سدرة المنتهی ثمّ عیون النّبیّین و المرسلین ای مست بادهٴ غرور اقلّ من حین
بشعور آی در خود و افعال خود تفکّر نما

\end{otherlanguage}

Idhan yanūḥu Muḥammad fī al-ufuq al-aʿlā wa yabkī al-Rūḥ fī al-rafīq
al-abhā thumma al-Kalīm ʿinda Sidrat al-Muntahā thumma ʿuyūn al-nabiyyīn
wa-l-mursalīn. Ey mast-e bāde-ye ghorūr aqall-e men ḥīn be-shoʿūr āy dar
khod va afʿāl-e khod tafakkor namā.

Thus Muḥammad laments in the highest horizon, and the Spirit weeps in
the all-glorious company, as does Moses beside the Lote Tree beyond
which none may pass, and the eyes of the Prophets and Messengers shed
tears. O you who are drunk with the wine of pride! Come to your senses
for less than a moment; reflect on yourself and your deeds.

\subsection{265}\label{section-1820}

\begin{otherlanguage}{persian}

گیرم بر این گروه نابالغ امر اللّه را مشتبه نمودی و جمیع هم تو را من دون
اللّه ساجد شدند چه نفعی برای تو حاصل لا فوربّ العالمین جز خسران دنیا و
آخرت حاصلی نداشته و ندارد و اگر الیوم بر سریر یاقوت جالس شوی عنداللّه بر
ارض هاویه ساکنی

\end{otherlanguage}

Gīram bar īn gorūh-e nābālegh Amr Allāh rā moshtabeh namūdī va jamīʿ ham
to rā min dūn Allāh sājed shodand che nafʿī barā-ye to ḥāṣel? Lā
fa-wa-Rabb al-ʿālamīn joz khosrān-e donyā va ākherat ḥāṣelī nadāshte va
nadārad va agar al-yawm bar sarīr-e yāqūt jāles shavī ʿenda-llāh bar
arż-e Hāvīye sākenī.

Suppose you have obscured the Cause of God for this immature group, and
all of them have bowed down to you instead of God. What benefit do you
gain? No, by the Lord of the worlds! It has yielded and yields nothing
but loss in this world and the next. Even if today you sit upon a throne
of ruby, in God's sight you dwell on the ground of the abyss.

\subsection{266}\label{section-1821}

\begin{otherlanguage}{persian}

فواللّه اگر بسمع فطرت استماع نمائی میشنوی که همان ارضی که بر او جالسی
پناه میبرد بخدا از تو و میگوید ای غافل بچه حجّت و دلیل حقّ خود را ثابت
مینمائی و حقّ من له الحقّ و به حقّق الحقّ را انکار میکنی فأفّ علیک و علی
الّذینهم اتّخذوک لأنفسهم من دون اللّه حبیب بحقّ موقن نشدی دیگر چرا بر قتلش
قیام نمودی

\end{otherlanguage}

Fa-wa-llāh agar be-samʿ-e feṭrat estemāʿ namāʾī mīshenavī ke hamān arżī
ke bar ū jālesī panāh mībarad be-Khodā az to va mīgūyad ey ghāfel be-che
ḥojjat va dalīl ḥaqq-e khod rā thābet mīnamāʾī va ḥaqq-e man lahu
al-ḥaqq wa bihi ḥuqqiqa al-ḥaqq rā enkār mīkonī? Fa-uff ʿalayka wa ʿalā
alladhīna-hum ittakhadhūka li-anfusihim min dūn Allāh ḥabīb. Be-Ḥaqq
mūqen nashodī dīgar cherā bar qatl ash qīyām namūdī?

By God, if you listen with the ear of your inborn nature, you will hear
the very ground on which you sit seeking refuge in God from you and
saying: ``O heedless one, by what proof and evidence do you establish
your own right while denying the right of Him to whom all right belongs
and through whom truth has been established? Shame on you and on those
who have taken you as their beloved instead of God!'' You did not become
certain of the True One; why, then, did you rise to put Him to death?

\subsection{267}\label{section-1822}

\begin{otherlanguage}{persian}

آخر بیست سنه در حفظ تو سعی نمود مع آنکه عالم بوده بما فی قلبک و اگر آنی
اراده میفرمود موجود نبودی مع‌ذلک متنبّه نشدی و باطراف عالم فعل خود را بحقّ
نسبت داده تاللّه نیست در این قلب مگر انوار تجلّیات فجر بقا و بر ضرّ نفسی
راضی نه ان

\end{otherlanguage}

Ākhar bīst sane dar ḥefẓ-e to saʿy namūd maʿ-e ānke ʿālem būde bi-mā fī
qalbika va agar ānī erāde mīfarmūd mowjūd nabūdī maʿa-dhālek motanabbeh
nashodī va be-aṭrāf-e ʿālam feʿl-e khod rā be-Ḥaqq nesbat dāde. Ta-llāh
nīst dar īn qalb magar anvār-e tajallīyāt-e fajr-e baqā va bar żorr-e
nafsī rāżī na {[}corrupted text{]}

After all, for twenty years He strove to protect you, although He knew
what was in your heart. Had He so willed for an instant, you would not
have existed. Yet you did not take heed, and throughout the world you
have attributed your own deed to the True One. By God, nothing is in
this heart but the lights of the effulgences of the dawn of eternity,
and it is not willing that any soul should suffer harm {[}corrupted
original text{]}

\subsection{268}\label{section-1823}

\begin{otherlanguage}{arabic}

یا اخی اسمع ندآء هذا الّذی وقع من ظلمک فی هذا البئر الّذی لم یکن لها من
قعر و کان قمیصه مرشوشاً بدم صادق من غلّک و فی قعر البئر ینادی و یقول یا
اخی لا تفعل بأخیک کما فعل ابن آدم بأخیه یا اخی اتّق اللّه و لا تجادل
بآیاتی و لا تحارب بنفسی و لا تقتل الّذی جعل صدره محلّاً لسیف الأعدآء لئلّا
یرد علیک من ضرّ و کان ان یحفظک فی کلّ اللّیالی و الأیّام و فی کلّ بکور و اصیل

\end{otherlanguage}

Yā akhī ismaʿ nidāʾ hādhā alladhī waqaʿa min ẓulmika fī hādhā al-biʾr
alladhī lam yakun lahā min qaʿr wa kāna qamīṣuhu marshūshan bi-dam ṣādiq
min ghillika wa fī qaʿr al-biʾr yunādī wa yaqūlu yā akhī lā tafʿal
bi-akhīka ka-mā faʿala Ibn Ādam bi-akhīhi. Yā akhī ittaqi Allāh wa lā
tujādil bi-āyātī wa lā tuḥārib bi-nafsī wa lā taqtul alladhī jaʿala
ṣadrahu maḥallan li-sayf al-aʿdāʾ li-allā yarida ʿalayka min ḍurr wa
kāna an yaḥfaẓaka fī kull al-layālī wa-l-ayyām wa fī kull bukūr wa aṣīl.

O My brother, hear the call of this One who, through your injustice, has
fallen into this bottomless well, whose shirt has been spattered with
true blood through your malice. From the depths of the well He calls out
and says: ``O My brother, do not do to your brother what the son of Adam
did to his brother. O My brother, fear God. Do not dispute My verses, do
not wage war against My person, and do not kill the One who made His
breast a target for the swords of His enemies so that no harm might come
to you, and who protected you through every night and day, every morning
and evening.''

\subsection{269}\label{section-1824}

\begin{otherlanguage}{persian}

یا اخی تاللّه الحقّ ما انطق عن الهوی ان هذا الّا وحی یوحی علّمنی شدید الرّوح
عند سدرة المنتهی تاللّه یا اخی ان هی من تلقآء نفسی بل من لدن عزیز حکیم
کلّما اقبلت الیک لئلّا تحدث من فتنة یفتننّ بها العباد حفظاً لأمر ربّک تاللّه
رمیت نحوی رمی الشّقاق و کلّما تقرّبت بک لعلّ تسکن بذلک نفسک قمت علیّ بالنّفاق
و یشهد بذلک ارکانک ان انت من المنکرین

\end{otherlanguage}

Yā akhī, ta-Allāhi al-Ḥaqqi mā anṭiqu ʿani al-hawā. In hādhā illā waḥyun
yūḥā. ʿAllamanī shadīdu al-rūḥi ʿinda sidrati al-muntahā. Ta-Allāhi yā
akhī, in hiya min tilqāʾi nafsī bal min ladun ʿazīzin ḥakīm. Kullamā
aqbaltu ilayka li-allā tuḥditha min fitnatin yuftananna bihā al-ʿibādu
ḥifẓan li-amri rabbika, ta-Allāhi ramayta naḥwī ramya al-shiqāq.
Wa-kullamā taqarrabtu bika laʿalla taskuna bi-dhālika nafsuka qumta
ʿalayya bi-al-nifāq. Wa-yashhadu bi-dhālika arkānuka in anta mina
al-munkirīn.

O My brother, by God, the True One, I do not speak from desire. This is
nothing but revelation revealed. One mighty in spirit taught Me at the
Lote Tree of the Utmost Boundary. By God, O My brother, it is not of My
own devising, but comes from One who is Almighty, All-Wise. Whenever I
turned toward you to prevent you from stirring up a turmoil by which the
servants would be put to the test, and to safeguard the Cause of your
Lord, by God, you hurled the shafts of discord at Me. And whenever I
drew near to you so that your soul might thereby find peace, you rose
against Me in hypocrisy. Your own limbs bear witness to this, if you
should deny it.

\subsection{270}\label{section-1825}

\begin{otherlanguage}{arabic}

تاللّه یا اخی کلّما کنت صامتاً عن بدایع ذکر ربّی روح القدس قامنی علی امره و
روح الأعظم ایقظنی عن رقدی و انطقنی بالحقّ بین السّموات و الأرضین ان کان
هذا ذنبی فلست انا اوّل من اذنب فقد ارتکب ذلک الّذینهم کانوا من قبلی ثمّ
الّذی سمّی بعلیّ فی ملکوت الأسمآء ثمّ بمحمّد فی جبروت القصوی ثمّ بابنی فی
الملإ الأعلی ثمّ بالکلیم فی هذا السّینآء المقدّس المبارک العزیز الرّفیع

\end{otherlanguage}

Ta-Allāhi yā akhī, kullamā kuntu ṣāmitan ʿan badāʾiʿi dhikri rabbī, rūḥu
al-qudusi qāmanī ʿalā amrihi wa-rūḥu al-aʿẓami ayqaẓanī ʿan raqdī
wa-anṭaqanī bi-al-ḥaqqi bayna al-samāwāti wa-al-araḍīn. In kāna hādhā
dhanbī fa-lastu anā awwala man adhnaba, fa-qad irtakaba dhālika
alladhīnahum kānū min qablī, thumma alladhī summiya bi-ʿAliyyin fī
malakūti al-asmāʾi thumma bi-Muḥammadin fī jabarūti al-quṣwā thumma
bi-ibnī fī al-malaʾi al-aʿlā thumma bi-al-Kalīmi fī hādhā al-Sīnāʾi
al-muqaddasi al-mubāraki al-ʿazīzi al-rafīʿ.

By God, O My brother, whenever I remained silent about the wonders of
the remembrance of My Lord, the Holy Spirit raised Me up for His Cause,
and the Most Great Spirit awakened Me from My slumber and made Me speak
the truth between the heavens and the earths. If this is My sin, I am
not the first to have sinned: those who came before Me committed it, as
did the One who was named ʿAlī in the kingdom of names, then Muḥammad in
the uttermost dominion, then My Son in the Concourse on High, then the
One who conversed with God upon this holy, blessed, mighty, exalted
Sinai.

\subsection{271}\label{section-1826}

\begin{otherlanguage}{persian}

یا اخی فانظر فی کتب النّبیّین و المرسلین ثمّ ما فی ایادیهم من حجج اللّه و
برهانه و دلائله و آثاره و ظهوراته و آیاته ثمّ انصف فی اخیک و لا تکن من
الّذینهم ظلموا نفس اللّه ثمّ استظلموا لیدخلوا الشّکّ و الرّیب فی قلوب المؤمنین

\end{otherlanguage}

Yā akhī, fa-unẓur fī kutubi al-nabiyyīna wa-al-mursalīn, thumma mā fī
ayādīhim min ḥujaji Allāhi wa-burhānihi wa-dalāʾilihi wa-āthārihi
wa-ẓuhūrātihi wa-āyātihi, thumma anṣif fī akhīka wa-lā takun mina
alladhīnahum ẓalamū nafsa Allāhi thumma istaẓlamū li-yudkhilū al-shakka
wa-al-rayba fī qulūbi al-muʾminīn.

O My brother, look into the books of the Prophets and Messengers, and
then into the proofs of God, His testimony, His evidences, His traces,
His manifestations, and His signs that are in their hands. Then judge
your brother fairly, and do not be among those who wronged the very Self
of God and then claimed to have been wronged, so as to instill doubt and
misgiving in the hearts of the believers.

\subsection{272}\label{section-1827}

\begin{otherlanguage}{arabic}

تاللّه لمّا عرف اخیک بأنّک قمت علیه و لن تسکن نار الحسد فی صدرک خرج بنفسه و
اهله وحده و دخل بیتاً اخری و غلق علی وجهه باب الدّخول و الخروج و کان فیه
لمن السّاکنین و مع ذلک انّک یا اخی ما سکنت فی نفسک و کتبت ما کتبت و ارسلت
ما ارسلت تاللّه انّ القلم یستحیی لیجری علی ما ذکرت و بما افتریت علی اخیک
تاللّه بذلک ضیّعت حرمتی و حرمة اللّه بین العباد فسوف تشهد و تری و ان تکون
حینئذ لمن الغافلین

\end{otherlanguage}

Ta-Allāhi lammā ʿarafa akhūka bi-annaka qumta ʿalayhi wa-lan taskuna
nāru al-ḥasadi fī ṣadrika kharaja bi-nafsihi wa-ahlihi waḥdahu
wa-dakhala baytan ukhrā wa-ghalaqa ʿalā wajhihi bāba al-dukhūli
wa-al-khurūji wa-kāna fīhi la-mina al-sākinīn. Wa-maʿa dhālika innaka yā
akhī mā sakanta fī nafsika wa-katabta mā katabta wa-arsalta mā arsalta.
Ta-Allāhi inna al-qalama yastaḥyī li-yajriya ʿalā mā dhakarta wa-bi-mā
iftarayta ʿalā akhīka. Ta-Allāhi bi-dhālika ḍayyaʿta ḥurmatī wa-ḥurmata
Allāhi bayna al-ʿibādi fa-sawfa tashhadu wa-tarā wa-in takūna ḥīnaʾidhin
la-mina al-ghāfilīn.

By God, when your brother realized that you had risen against Him and
that the fire of envy in your breast would not subside, He left on His
own with His family, entered another house, shut its door against His
own coming and going, and remained there. Yet even then, O My brother,
you found no peace within yourself: you wrote what you wrote and sent
what you sent. By God, the Pen is ashamed to set down what you have said
and the accusations you have fabricated against your brother. By God,
through this you have violated My honor and the honor of God among His
servants. You will soon witness and see, even though you are now among
the heedless.

\subsection{273}\label{section-1828}

\begin{otherlanguage}{arabic}

اذاً قم عن رقد نفسک و هواک و قصّ شارب الحسد و قلّم اظفار البغضآء ثمّ تطیّب من
طیب الوفآء ثمّ غسّل بمیاه القدس و ضع وجهک علی التّراب بخضوع و خشوع و انابة
و رجوع

\end{otherlanguage}

Idhan qum ʿan raqdi nafsika wa-hawāka wa-quṣṣ shāriba al-ḥasadi
wa-qallim aẓfāra al-baghḍāʾi thumma taṭayyab min ṭībi al-wafāʾi thumma
ghassil bi-miyāhi al-qudsi wa-ḍaʿ wajhaka ʿalā al-turābi bi-khuḍūʿin
wa-khushūʿin wa-inābatin wa-rujūʿ.

Rise, then, from the slumber of self and desire, trim the mustache of
envy, and clip the nails of hatred. Then perfume yourself with the
fragrance of faithfulness, wash with the waters of holiness, and lay
your face upon the dust in submission, humility, repentance, and return.

\subsection{274}\label{section-1829}

\begin{otherlanguage}{arabic}

محبوب و قل ای ربّ انا الّذی فرّطت فی جنب اخی فی مذ اللّیالی و الأیّام و کنت
غافلاً عن بدایع ذکرک الحکیم اذاً یا الهی فارفع الحجاب عن بصری لأعرف نفسک و
اقوم بثنائک و انقطع عمّا سواک و اقبل الی وجهک الکریم ثمّ اجعلنی من عبادک
الّذین جعلت لهم مقعد صدق عندک ثمّ ارزقنی من تسنیم عنایتک و کوثر افضالک

\end{otherlanguage}

Maḥbūb. Wa-qul ay rabbi anā alladhī farraṭtu fī janbi akhī fī mudh
al-layālī wa-al-ayyāmi wa-kuntu ghāfilan ʿan badāʾiʿi dhikrika al-ḥakīm.
Idhan yā ilāhī fa-irfaʿ al-ḥijāba ʿan baṣarī li-aʿrifa nafsaka wa-aqūma
bi-thanāʾika wa-anqaṭiʿa ʿammā siwāka wa-uqbila ilā wajhika al-karīm.
Thumma ijʿalnī min ʿibādika alladhīna jaʿalta lahum maqʿada ṣidqin
ʿindaka thumma urzuqnī min tasnīmi ʿināyatika wa-kawthari afḍālika.

Beloved. And say: ``O Lord, I am the one who failed in my duty toward my
brother throughout the nights and days and was heedless of the wonders
of Your wise remembrance. Now, O my God, lift the veil from my sight, so
that I may recognize Your Self, arise to praise You, detach myself from
everything but You, and turn toward Your generous face. Then make me one
of those servants of Yours for whom You have appointed a seat of truth
in Your presence, and grant me sustenance from the heavenly spring of
Your care and the river of Your bounty.

\subsection{275}\label{section-1830}

\begin{otherlanguage}{arabic}

ثمّ الحقنی بعبادک المخلصین الّذین ما التفتوا بالدّنیا و لا بریاساتها و لا
بما فیها و علیها و انقطعوا بأنفسهم و اموالهم فی سبیل بارئهم و کانوا من
المنقطعین ای ربّ لا تدعنی بنفسی ثمّ خذ یدی بید قدرتک ثمّ انقذنی من غمرات
النّفس و الهوی و لهیبها و اشتعالها ثمّ اجعل تلک النّار علیّ برداً و سلاماً و
روحاً و ریحاناً ثمّ اکتبنی من عبادک المنقطعین

\end{otherlanguage}

Thumma alḥiqnī bi-ʿibādika al-mukhliṣīna alladhīna mā iltafatū
bi-al-dunyā wa-lā bi-riyāsātihā wa-lā bi-mā fīhā wa-ʿalayhā wa-inqaṭaʿū
bi-anfusihim wa-amwālihim fī sabīli bāriʾihim wa-kānū mina
al-munqaṭiʿīn. Ay rabbi lā tadaʿnī bi-nafsī thumma khudh yadī bi-yadi
qudratika thumma anqidhني min ghamarāti al-nafsi wa-al-hawā wa-lahībihā
wa-ishtiʿālihā thumma ijʿal tilka al-nāra ʿalayya bardan wa-salāman
wa-rūḥan wa-rayḥānan thumma uktubnī min ʿibādika al-munqaṭiʿīn.

Then join me to Your sincere servants, who paid no heed to the world,
its positions of authority, or anything within it or upon it, who
detached themselves and their possessions in the path of their Maker and
were among the detached. O Lord, do not leave me to myself. Take my hand
in the hand of Your power, and rescue me from the engulfing depths of
self and desire, their blaze and their burning. Then make that fire
coolness and peace for me, and a spirit of life and sweet fragrance, and
number me among Your detached servants.

\subsection{276}\label{section-1831}

\begin{otherlanguage}{arabic}

ای ربّ وفّقنی لخدمتک و تبلیغ آیاتک ثمّ اجعلنی ناصراً لأمرک و حافظاً لدینک و
ناطقاً بثنائک و معلناً بدایع افضالک و اکرامک و انّک انت المقتدر علی ما تشآء
و انّک انت العزیز الکریم ای ربّ لا تخیّب من تمسّک بحبل عنایتک و لا تطرد من
علّق سبّابة الرّجآء علی حبل جودک و افضالک

\end{otherlanguage}

Ay rabbi waffiqnī li-khidmatika wa-tablīghi āyātika thumma ijʿalnī
nāṣiran li-amrika wa-ḥāfiẓan li-dīnika wa-nāṭiqan bi-thanāʾika
wa-muʿlinan badāʾiʿa afḍālika wa-ikrāmika wa-innaka anta al-muqtadiru
ʿalā mā tashāʾu wa-innaka anta al-ʿazīzu al-karīm. Ay rabbi lā tukhayyib
man tamassaka bi-ḥabli ʿināyatika wa-lā taṭrud man ʿallaqa sabbābata
al-rajāʾi ʿalā ḥabli jūdika wa-afḍālika.

O Lord, enable me to serve You and convey Your verses. Make me a helper
of Your Cause, a guardian of Your religion, one who voices Your praise
and proclaims the wonders of Your bounty and generosity. You indeed have
power over whatever You will, and You are the Almighty, the Generous. O
Lord, do not disappoint anyone who has grasped the cord of Your care,
and do not drive away anyone who has hooked the finger of hope upon the
cord of Your generosity and bounty.

\subsection{277}\label{section-1832}

\begin{otherlanguage}{arabic}

ای ربّ لا تحتجبنی عن منبع رضاک ثمّ ارضنی بالاقرار بما رضیته و نزّلته عن
غمام فضلک و سحاب عزّ مکرمتک و انت المعطی فی کلّ الأحوال و انّک انت الغفور
الرّحیم ای ربّ لا تعرّ جسدی عن قمیص الانصاف و لا قلبی عن برد الاعتراف بنفسک
الرّحمن الرّحیم

\end{otherlanguage}

Ay rabbi lā taḥtajibnī ʿan manbaʿi riḍāka thumma arḍinī bi-al-iqrāri
bi-mā raḍītahu wa-nazzaltahu ʿan ghamāmi faḍlika wa-saḥābi ʿizzi
makrumatika wa-anta al-muʿṭī fī kulli al-aḥwāli wa-innaka anta
al-ghafūru al-raḥīm. Ay rabbi lā tuʿarri jasadī ʿan qamīṣi al-inṣāfi
wa-lā qalbī ʿan burdi al-iʿtirāfi bi-nafsika al-raḥmāni al-raḥīm.

O Lord, do not veil me from the wellspring of Your good pleasure. Make
me content to acknowledge what You have approved and sent down from the
clouds of Your grace and the clouds of Your glorious generosity. You are
the Bestower in every circumstance, and You are the All-Forgiving, the
Merciful. O Lord, do not strip my body of the shirt of fairness, or my
heart of the robe of acknowledgment of Your Self, the All-Merciful, the
Merciful.

\subsection{278}\label{section-1833}

\begin{otherlanguage}{arabic}

ای ربّ فاجعل قدمی ثابتاً علی صراطک بحیث لا انکر ما دعوتنی به فی کلّ آیاتک و
الواحک و زبرک و کتبک و اسفارک و صحایف قدسک المنیع ای ربّ فاجعلنی ناظراً
الی شطر مواهبک و راجعاً الی بحر غفرانک فلا تعرّنی عن جمیل ثنائک و انّک انت
القادر علی ما تشآء و انّک انت المقتدر علی ما ترید

\end{otherlanguage}

Ay rabbi fa-ijʿal qadamī thābitan ʿalā ṣirāṭika bi-ḥaythu lā unkira mā
daʿawtanī bihi fī kulli āyātika wa-alwāḥika wa-zuburika wa-kutubika
wa-asfārika wa-ṣaḥāyifi qudsika al-manīʿ. Ay rabbi fa-ijʿalnī nāẓiran
ilā shaṭri mawāhibika wa-rājiʿan ilā baḥri ghufrānika fa-lā tuʿarrinī
ʿan jamīli thanāʾika wa-innaka anta al-qādiru ʿalā mā tashāʾu wa-innaka
anta al-muqtadiru ʿalā mā turīd.

O Lord, make my foot firm upon Your path, so that I may not deny that to
which You have summoned me in all Your verses, Your tablets, Your
scriptures, Your books, Your volumes, and the pages of Your unassailable
holiness. O Lord, make me look toward Your gifts and return to the ocean
of Your forgiveness. Do not strip me of the beauty of praising You. You
indeed have power over whatever You will, and You are mighty to do
whatever You desire.

\subsection{279}\label{section-1834}

\begin{otherlanguage}{arabic}

و ایّدنی بأن لا انکر ما حقّق به ایمانی و ثبت ذکری و رفع اسمی و بعث کینونتی
و ذوّت حقیقتی و کنت من المؤمنین ای ربّ لا تحرم من وقف تلقآء مدین رحمتک و
تشبّث بذیل احسانک و فضلک ای ربّ فاجعل لی کلمة من عندک لأکون متذکّراً بها فی
العشیّ و الاشراق و فی کلّ حین

\end{otherlanguage}

Wa-ayyidnī bi-an lā unkira mā ḥuqqiqa bihi īmānī wa-thabata dhikrī
wa-rufiʿa ismī wa-buʿithat kaynūnatī wa-dhuwwitat ḥaqīqatī wa-kuntu mina
al-muʾminīn. Ay rabbi lā taḥrim man waqafa tilqāʾa madīni raḥmatika
wa-tashabbatha bi-dhayli iḥsānika wa-faḍlika. Ay rabbi fa-ijʿal lī
kalimatan min ʿindika li-akūna mutadhakkiran bihā fī al-ʿashiyyi
wa-al-ishrāqi wa-fī kulli ḥīn.

Strengthen me so that I may not deny that through which my faith was
made real, my remembrance established, my name exalted, my being brought
to life, and my reality brought into existence, and through which I
became one of the believers. O Lord, do not deprive the one who has
stood before the city of Your mercy and clung to the hem of Your
kindness and grace. O Lord, appoint for me a word from Your presence, so
that through it I may remain mindful in the evening, at daybreak, and at
every moment.''

\subsection{280}\label{section-1835}

\begin{otherlanguage}{arabic}

کذلک تنطقک لسان المظلوم فی قعر الجبّ لعلّ تتّخذ الی شطر الانصاف من سبیل اذاً
یخاطب اللّه نسیم قدسه الّذی یهبّ عن شطر العرش و جعله رسولاً من عنده علی
العالمین لأنّه لن یجد فی تلک الأیّام من رسول لیرسله الی العباد ببشارات
امره و یجعلهم من الذّاکرین و المستبشرین لأنّه وقع فی سجن الّذی انقطعت عنه
ایادی المریدین و ارجل القاصدین و من دون السّجن وقع فی بئر الحسد الّذی ما
اطّلع بقعره الّا نفسه المحصی العلیم الخبیر کذلک قصّ اللّه من قصص الحقّ بلسانه
الصّادق المتکلّم العلیم الأمین

\end{otherlanguage}

Kadhālika tunṭiquka lisānu al-Maẓlūmi fī qaʿri al-jubbi laʿalla
tattakhidha ilā shaṭri al-inṣāfi min sabīl. Idhan yukhāṭibu Allāhu
nasīma qudsihi alladhī yahubbu ʿan shaṭri al-ʿarshi wa-jaʿalahu rasūlan
min ʿindihi ʿalā al-ʿālamīna li-annahu lan yajida fī tilka al-ayyāmi min
rasūlin li-yursilahu ilā al-ʿibādi bi-bishārāti amrihi wa-yajʿalahum
mina al-dhākirīna wa-al-mustabshirīna li-annahu waqaʿa fī sijni alladhī
inqaṭaʿat ʿanhu ayādī al-murīdīna wa-arjulu al-qāṣidīn. Wa-min dūni
al-sijni waqaʿa fī biʾri al-ḥasadi alladhī mā iṭṭalaʿa bi-qaʿrihi illā
nafsuhu al-muḥṣī al-ʿalīmu al-khabīr. Kadhālika qaṣṣa Allāhu min qaṣaṣi
al-ḥaqqi bi-lisānihi al-ṣādiqi al-mutakallimi al-ʿalīmi al-amīn.

Thus does the tongue of the Wronged One in the depths of the pit give
you words to speak, so that you may find a path toward fairness. God now
addresses His holy breeze, which blows from the direction of the Throne,
and has made it a messenger from His presence to the worlds, for in
these days He will find no messenger to send to His servants with the
glad tidings of His Cause and to make them among those who remember and
rejoice. For He has fallen into a prison beyond the reach of the hands
of His devotees and the feet of those who seek Him. Beyond this
imprisonment, He has fallen into the well of envy, whose depths none
knows but His own Self, the All-Reckoning, the All-Knowing, the
All-Informed. Thus has God recounted a true account through His
truthful, speaking, all-knowing, faithful Tongue.

\subsection{281}\label{section-1836}

\begin{otherlanguage}{persian}

ای نسیم صبا چون قاصدی ملاحظه نمیشود تو برایحهٴ قمیص بها از رضوان بقا بر
مریدین از احبّایم مرور نما و ببیّنات روح و آیات ظهور جمیع را آگاه کن که
شاید بعضی از نفوس از جمیع من علی الأرض و تعلّق بآن پاک و مقدّس شده بفردوس
اعظم راجع شوند ولکن ای نسیم بانقطاع تمام مرور نما بشأنی که اگر ضرّ عالمین
بر تو وارد شود صابر باشی و اگر نعمت آن بر تو نازل گردد توجّه ننمائی چه که
اگر از جهات حسد و بغض و ردّ و قبول و سکون و اضطراب جمیع من علی الأرض مقدّس
نشوی قادر بر این تبلیغ منیع و فایز بحمل اسرار ربّانی نگردی کذلک یأمرک
لسان ربّک لتکون من العاملین

\end{otherlanguage}

Ey nasīm-e ṣabā, chūn qāṣedī molāḥeẓe nemī-shavad to be-rāyeḥe-ye
qamīṣ-e Bahā az Reżvān-e baqā bar morīdīn az aḥebbāyam morūr namā va
be-bayyenāt-e rūḥ va āyāt-e ẓohūr jamīʿ rā āgāh kon ke shāyad baʿżī az
nofūs az jamīʿ-e man ʿalā al-arḍ va taʿalloq be-ān pāk va moqaddas shode
be-Ferdows-e Aʿẓam rājeʿ shavand. Valāken ey nasīm be-enqeṭāʿ-e tamām
morūr namā be-shaʾnī ke agar żorr-e ʿālamīn bar to vāred shavad ṣāber
bāshī va agar neʿmat-e ān bar to nāzel gardad tavajjoh nanamāʾī, che ke
agar az jahāt-e ḥasad va boghż va radd va qabūl va sokūn va eżṭerāb-e
jamīʿ-e man ʿalā al-arḍ moqaddas nashavī qāder bar īn tablīgh-e manīʿ va
fāyez be-ḥaml-e asrār-e rabbānī nagardī. Kadhālika yaʾmuruka lisānu
rabbika li-takūna mina al-ʿāmilīn.

O morning breeze, since no messenger is to be found, pass from the
Riḍván of eternity, bearing the fragrance of the shirt of Bahá, among
the devoted ones of My loved ones. Acquaint them all with the clear
evidences of the spirit and the signs of revelation, so that perhaps
some souls may become pure and sanctified from all who dwell on earth
and from attachment to it, and return to the Most Great Paradise. But, O
breeze, pass with complete detachment, so that if the harm of all the
worlds should assail you, you would remain patient, and if their
blessings should descend upon you, you would pay no heed. For unless you
become sanctified from the envy and hatred, rejection and acceptance,
calm and agitation of all who dwell on earth, you will neither be
capable of this exalted proclamation nor attain the privilege of bearing
divine mysteries. Thus does the tongue of your Lord command you, so that
you may be among those who act.

\subsection{282}\label{section-1837}

\begin{otherlanguage}{persian}

بگو ای احباب و ای اولی الألباب آخر قدری نظر را از توجّه بدنیا و شئونات آن
مقدّس نمائید و دیدهٴ بصیرت در امورات ظاهره و شئونات لامعه که از شطر عرش
ابهی ظاهر و هویدا است ملاحظه نموده تفکّر نمائید که شاید در این ایّام که
سکر غفلت جمیع اهل سموات و ارض را احاطه نموده خود را بمدینهٴ احدیّۀ ابدیّهٴ
الهی کشانید و از بدایع رحمت بی‌زوالش محروم نمانید که مباد نعوذ باللّه از
مقصود اوّلیّهٴ رحمانی محجوب گردید و از معرضین در کتاب ربّ الأرباب محسوب
شوید من‌دون آنکه شاعر باشید فنعوذ باللّه عن ذلک یا اولی الألباب

\end{otherlanguage}

Begū ey aḥbāb va ey olī al-albāb, ākhar qadrī naẓar rā az tavajjoh
be-donyā va shoʾūnāt-e ān moqaddas namāʾīd va dīde-ye baṣīrat dar
omūrāt-e ẓāhere va shoʾūnāt-e lāmeʿe ke az shaṭr-e ʿarsh-e Abhā ẓāher va
hoveydā ast molāḥeẓe namūde tafakkor namāʾīd, ke shāyad dar īn ayyām ke
sokr-e ghaflat jamīʿ-e ahl-e samavāt va arż rā eḥāṭe namūde khod rā
be-madīne-ye aḥadīyye-ye abadīyye-ye elāhī keshānīd va az badāyeʿ-e
raḥmat-e bī-zavālash maḥrūm namānīd, ke mabād naʿūdhu bi-Allāh az
maqṣūd-e avvalīyye-ye raḥmānī maḥjūb gardīd va az moʿreżīn dar ketāb-e
Rabb al-Arbāb maḥsūb shavīd men-dūn-e ānke shāʿer bāshīd. Fa-naʿūdhu
bi-Allāhi ʿan dhālika yā ulī al-albāb.

Say: O loved ones, O people of understanding! At last, free your gaze to
some degree from attention to the world and its affairs. With the eye of
insight, consider and reflect upon the manifest events and radiant signs
that have appeared and become evident from the direction of the
All-Glorious Throne. Perhaps, in these days when the intoxication of
heedlessness has overtaken all the inhabitants of heaven and earth, you
may bring yourselves to the city of God's eternal oneness and not remain
deprived of the wonders of His unfailing mercy. Otherwise---God
forbid---you may be veiled from the primary purpose of the All-Merciful
and be reckoned in the Book of the Lord of Lords among those who turn
away, without being aware of it. We seek refuge in God from this, O
people of understanding!

\subsection{283}\label{section-1838}

\begin{otherlanguage}{persian}

جمیع انبیا و رسل ناس را بسبل عرفان جمال رحمن دعوت نموده‌اند چه که این
مقام اعظم مقامات بوده و خواهد بود قدری ملاحظه نموده در ارسال رسل الهی که
بچه سبب و جهت این هیاکل احدیّه از غیب بعرصهٴ شهود آمده‌اند و جمیع این
بلایا و رزایا که شنیده‌اید جمیع را تحمّل فرموده‌اند شکّی نیست که مقصود جز
دعوت عباد بعرفان جمال رحمن نبوده و نخواهد بود و اگر بگوئید مقصود اوامر و
نواهی آن بوده شکّی نیست که این مقصود اوّلیّه نبوده و نخواهد بود چنانچه اگر
کسی بعبادت اهل سموات و ارض قیام نماید و از عرفان الهی محروم باشد هرگز
نفعی بعاملین آن نبخشیده و نخواهد بخشید و در جمیع کتب سماویّه این مطلب
مشهود و واضح است و اگر نفسی عارف بحقّ باشد و جمیع اوامر الهیّه را ترک
نماید امید نجات هست چنانچه نزد اولی البصر واضح و مبرهن است

\end{otherlanguage}

Jamīʿ-e anbiyā va rosol nās rā be-sobol-e ʿerfān-e Jamāl-e Raḥmān daʿvat
namūde-and, che ke īn maqām aʿẓam-e maqāmāt būde va khāhad būd. Qadrī
molāḥeẓe namūde dar ersāl-e rosol-e elāhī ke be-che sabab va jahat īn
hayākel-e aḥadīyye az gheyb be-ʿarṣe-ye shohūd āmade-and va jamīʿ-e īn
balāyā va razāyā ke shenīde-īd jamīʿ rā taḥammol farmūde-and. Shakkī
nīst ke maqṣūd joz daʿvat-e ʿebād be-ʿerfān-e Jamāl-e Raḥmān nabūde va
nakhāhad būd. Va agar begūʾīd maqṣūd avāmer va navāhī-ye ān būde shakkī
nīst ke īn maqṣūd-e avvalīyye nabūde va nakhāhad būd, chenānche agar
kasī be-ʿebādat-e ahl-e samavāt va arż qiyām namāyad va az ʿerfān-e
elāhī maḥrūm bāshad hargez nafʿī be-ʿāmelīn-e ān nabakhshīde va nakhāhad
bakhshīd. Va dar jamīʿ-e kotob-e samāvīyye īn maṭlab mashhūd va vāżeḥ
ast. Va agar nafsī ʿāref be-Ḥaqq bāshad va jamīʿ-e avāmer-e elāhīyye rā
tark namāyad omīd-e najāt hast, chenānche nazd-e olī al-baṣar vāżeḥ va
mobarhan ast.

All the Prophets and Messengers have summoned humanity to the paths of
recognition of the Beauty of the All-Merciful, for this has been and
will remain the greatest of all stations. Reflect a little on the
sending of the divine Messengers: for what cause and purpose have these
Temples of Oneness come from the unseen into the realm of the visible,
and endured all the afflictions and calamities of which you have heard?
There is no doubt that their purpose has been, and will remain, nothing
but to summon the servants to recognition of the Beauty of the
All-Merciful. If you say that the purpose was His commandments and
prohibitions, there is no doubt that this has not been, and will not be,
the primary purpose. For if someone were to perform the worship of all
the inhabitants of heaven and earth yet remain deprived of the
recognition of God, it would never have benefited, nor will it ever
benefit, those who perform it. This matter is evident and clear in all
the heavenly books. And if a soul recognizes the True One but abandons
all the divine commandments, there is hope of salvation, as is clear and
manifest to those endowed with insight.

\subsection{284}\label{section-1839}

\begin{otherlanguage}{persian}

پس مقصود اوّلیّه از خلق ابداع و ظهور اختراع و ارسال رسل و انزال کتب و حمل
رسل مشقّتهای لایحصی را جمیع بعلّت عرفان جمال سبحان بوده پس حال اگر نفسی
بجمیع اعمال مشغول شود و در تمام عمر بقیام و قعود و ذکر و فکر و مادون آن
از اعمال مشغول گردد و از عرفان اللّه محروم ماند ابداً ثمری باو راجع نخواهد
شد و عرفان اللّه هم لازال عرفان مظهر نفس او بوده در میانهٴ خلق او چنانچه
در جمیع کتب خاصّه در بیان که در جمیع الواح آن این مطلب بلند اعلی و این
لطیفهٴ اعزّ قصوی مذکور گشته و مبرهن آمده فطوبی للعارفین

\end{otherlanguage}

Pas maqṣūd-e avvalīyye az khalq-e ebdāʿ va ẓohūr-e ekhterāʿ va ersāl-e
rosol va enzāl-e kotob va ḥaml-e rosol mashaqqat-hā-ye lā yoḥṣā rā jamīʿ
be-ʿellat-e ʿerfān-e Jamāl-e Sobḥān būde. Pas ḥāl agar nafsī be-jamīʿ-e
aʿmāl mashghūl shavad va dar tamām-e ʿomr be-qiyām va qoʿūd va dhekr va
fekr va mādūn-e ān az aʿmāl mashghūl gardad va az ʿerfān-e Allāh maḥrūm
mānad abadan thamarī be-ū rājeʿ nakhāhad shod. Va ʿerfān-e Allāh ham
lāzāl ʿerfān-e maẓhar-e nafs-e Ū būde dar miyāne-ye khalq-e Ū, chenānche
dar jamīʿ-e kotob khāṣṣe dar Bayān ke dar jamīʿ-e alvāḥ-e ān īn maṭlab-e
boland-e aʿlā va īn laṭīfe-ye aʿazz-e qoṣvā madhkūr gashte va mobarhan
āmade. Fa-ṭūbā li-al-ʿārifīn.

Thus the primary purpose of bringing creation into being, of the
appearance of originated things, of sending Messengers, of revealing
books, and of the Messengers' bearing countless hardships has, in every
case, been recognition of the Beauty of the All-Glorious. Now,
therefore, if a soul should engage in every deed and spend an entire
lifetime in standing and sitting, remembrance and contemplation, and
other such acts, yet remain deprived of the recognition of God, no fruit
will ever accrue to that soul. And recognition of God has always meant
recognition of the Manifestation of His own Self among His creatures, as
set forth in all the books, especially in the Bayán, in all of whose
tablets this lofty and most exalted theme, this most precious and
sublime truth, has been stated and demonstrated. Blessed, then, are
those who recognize.

\subsection{285}\label{section-1840}

\begin{otherlanguage}{arabic}

و اگر ببصر اطهر ملاحظه شود مشهود میگردد که جمیع هیاکل احدیّه که جان و مال
و ننگ و نام را در سبیل محبوب انفاق نموده‌اند در رتبهٴ اوّلیّه مقصودی
نداشته‌اند جز آنکه عباد را بشریعهٴ عرفان کشانند حال ملاحظه در انبیا
نمائید که چه قدر بلایای لایحصی حمل نموده‌اند که شاید ناس حجبات وهم را خرق
نمایند و از کوثر یقین مشروب گردند

\end{otherlanguage}

Va agar bi-baṣar-i aṭhar mulāḥiẓa shavad, mashhūd mīgardad ki jamīʿ-i
hayākil-i aḥadiyya ki jān va māl va nang va nām rā dar sabīl-i maḥbūb
infāq namūda-and, dar rutba-yi avvaliyya maqṣūdī nadāshta-and juz ānki
ʿibād rā bi-sharīʿa-yi ʿirfān kishānand. Ḥāl mulāḥiẓa dar anbiyā namāyīd
ki chi qadr balāyā-yi lāyuḥṣā ḥaml namūda-and ki shāyad nās ḥujubāt-i
vahm rā kharq namāyand va az kawthar-i yaqīn mashrūb gardand.

If viewed with the purest sight, it becomes evident that all the Temples
of Oneness who have sacrificed life and possessions, reputation and name
in the path of the Beloved have had no primary purpose except to draw
the servants to the waters of recognition. Now consider the Prophets:
how many countless afflictions they have borne so that people might tear
through the veils of illusion and drink from the river of certainty.

\subsection{286}\label{section-1841}

\begin{otherlanguage}{persian}

و چون حجبات غلیظیّهٴ وهمیّه در انظار بسیار بزرگ و مهیمن بود لذا هر رسول که
از جانب حقّ ظاهر شد جمیع بر اعتراض بآن نفحهٴ ربّانیّه قیام نموده‌اند تا آنکه
ارسال رسل منتهی بهادی سبل در سنهٴ ستّین شد ملاحظه شد که چه قدر ناس
بوهمیّات انفس خود از شاطی قدس یقین دور مانده‌اند بشأنی که از خدا جز هوی و
از یقین جز ظنّ مبین در مابینشان مذکور نه

\end{otherlanguage}

Va chūn ḥujubāt-i ghalīẓiyya-yi vahmiyya dar anẓār bisyār buzurg va
muhaymin būd, lidhā har rasūl ki az jānib-i Ḥaqq ẓāhir shud, jamīʿ bar
iʿtirāḍ bi-ān nafḥa-yi rabbāniyya qiyām namūda-and, tā ānki irsāl-i
rusul muntahī bi-Hādī-yi subul dar sana-yi sittīn shud. Mulāḥiẓa shud ki
chi qadr nās bi-vahmiyyāt-i anfus-i khud az shāṭī-yi quds-i yaqīn dūr
mānda-and, bi-shaʾnī ki az Khudā juz havā va az yaqīn juz ẓann-i mubīn
dar mābaynishān madhkūr na.

Since the thick veils of illusion loomed so large and held such sway in
people's eyes, whenever a Messenger appeared from God, all rose in
opposition to that divine Breeze, until the succession of Messengers
culminated in the Guide of the paths in the year sixty. It became
evident how far people had remained from the sacred shore of certainty
through their own vain imaginings, to such a degree that among them
nothing was spoken of as God but desire, and nothing as certainty but
manifest conjecture.

\subsection{287}\label{section-1842}

\begin{otherlanguage}{persian}

و چون جمال علیّ اعلی امر بخرق احجاب فرمود جمیع بر اعراضش قیام نمودند تا
آنکه معدودی قلیل باصل مقصود عارف شده جمیع حجبات وهمیّه و سبحات ظنونیّه را
بنار سدرهٴ ربّانیّه محترق نموده بعرفان جمال رحمن فایز گشتند تاللّه الحقّ اگر
نفسی در آنچه بر آن جمال مبارک وارد شده از اعراض علما و بلایای لایحصی
تفکّر نماید در تمام عمر بناله و نوحه و ندبه مشغول گردد تا آنکه اهل کین آن
جمال مبین را در هوا معلّق نمودند و برصاص شرک و بغضا آن سدرهٴ عزّ منتهی را
قطع نمودند

\end{otherlanguage}

Va chūn Jamāl-i ʿAlī-yi Aʿlā amr bi-kharq-i aḥjāb farmūd, jamīʿ bar
iʿrāḍash qiyām namūdand, tā ānki maʿdūdī qalīl bi-aṣl-i maqṣūd ʿārif
shuda, jamīʿ-i ḥujubāt-i vahmiyya va subuḥāt-i ẓunūniyya rā bi-nār-i
sidra-yi rabbāniyya muḥtariq namūda, bi-ʿirfān-i Jamāl-i Raḥmān fāʾiz
gashtand. Ta-Allāh al-Ḥaqq, agar nafsī dar ānchi bar ān Jamāl-i mubārak
vārid shuda az iʿrāḍ-i ʿulamā va balāyā-yi lāyuḥṣā tafakkur namāyad, dar
tamām-i ʿumr bi-nāla va nawḥa va nudba mashghūl gardad, tā ānki ahl-i
kīn ān Jamāl-i mubīn rā dar havā muʿallaq namūdand va bi-raṣāṣ-i shirk
va baghḍā ān sidra-yi ʿizz-i muntahā rā qaṭʿ namūdand.

When the most exalted Beauty of ʿAlī commanded that the veils be torn
away, all rose to turn away from Him, until a small number recognized
the essential purpose and, burning away all the veils of illusion and
the splendors of conjecture with the fire of the divine Tree, attained
recognition of the Beauty of the All-Merciful. By God, the True One! If
anyone were to reflect on what befell that blessed Beauty---the
rejection of the religious scholars and countless afflictions---that
person would spend an entire lifetime lamenting, mourning, and weeping.
At last, the people of malice suspended that manifest Beauty in the air
and, with the bullets of idolatry and hatred, cut down that Tree of
utmost glory.

\subsection{288}\label{section-1843}

\begin{otherlanguage}{persian}

و آن جمال مبارک در حینی که معلّق بهوا بود در سرّ سرّ بلسان ناطق مبین باین
کلمات محکم متین تکلّم میفرمودند که ای اهل بیان قدری تفکّر در این خلق
نمائید که جمیع خود را عارف بحقّ میدانند و سالک سبل یقین میشمرند و در کلّ
احیان باذکار و کلمات رحمن مشغول شده‌اند بقسمی که یومی از ایّام اوامر ملیک
علّام را ترک ننموده‌اند و اگر هم از نفسی ترک شده خود را نادم مشاهده نموده

\end{otherlanguage}

Va ān Jamāl-i mubārak dar ḥīnī ki muʿallaq bi-havā būd, dar sirr-i sirr
bi-lisān-i nāṭiq-i mubīn bi-īn kalimāt-i muḥkam-i matīn takallum
mīfarmūdand ki ay ahl-i Bayān, qadrī tafakkur dar īn khalq namāyīd ki
jamīʿ khud rā ʿārif bi-Ḥaqq mīdānand va sālik-i subul-i yaqīn mīshumrand
va dar kull-i aḥyān bi-adhkār va kalimāt-i Raḥmān mashghūl shuda-and,
bi-qismī ki yawmī az ayyām avāmir-i Malīk-i ʿAllām rā tark nanamūda-and
va agar ham az nafsī tark shuda, khud rā nādim mushāhida namūda.

While that blessed Beauty was suspended in the air, in the innermost
mystery He spoke these firm and weighty words with an eloquent and lucid
tongue: ``People of the Bayān, reflect a little on these people. They
all consider themselves to have recognized God and count themselves
among those who walk the paths of certainty. At all times they have
occupied themselves with the remembrance and words of the All-Merciful,
so much so that they have not neglected the commandments of the
all-knowing King for a single day; and if anyone has neglected them,
that person has felt remorse.

\subsection{289}\label{section-1844}

\begin{otherlanguage}{arabic}

و مع‌ذلک جوهر رحمن و ساذج سبحان را که بعرفان او متمسّک بوده‌اند معلّق داشته
و شهید مینمایند و ابداً عارف باو نگشته و خود را از جمیع فیوضات الهیّه و
عنایات عزّ رحمانیّه محروم داشته‌اند در قعر نارند و خود را از اهل جنّت
میشمرند و در بحر عذاب مستغرقند و خود را از احباب ربّ الأرباب میدانند و در
منتهی بعد از حقّ خود را از اهل قرب فرض گرفته‌اند

\end{otherlanguage}

Va maʿa-dhālik Jawhar-i Raḥmān va Sādhij-i Subḥān rā ki bi-ʿirfān-i ū
mutamassik būda-and muʿallaq dāshta va shahīd mīnamāyand va abadan ʿārif
bi-ū nagashta va khud rā az jamīʿ-i fuyūḍāt-i ilāhiyya va ʿināyāt-i
ʿizz-i raḥmāniyya maḥrūm dāshta-and. Dar qaʿr-i nārand va khud rā az
ahl-i jannat mīshumrand va dar baḥr-i ʿadhāb mustaghraqand va khud rā az
aḥbāb-i Rabb al-Arbāb mīdānand va dar muntahā-yi buʿd az Ḥaqq khud rā az
ahl-i qurb farḍ girifta-and.

Yet they have suspended and are putting to death the Essence of the
All-Merciful, the pure Being of the All-Glorious, to whose recognition
they have professed to hold fast. They have never recognized Him and
have deprived themselves of all divine outpourings and the glorious
favors of the All-Merciful. They are in the depths of the fire, yet
count themselves among the people of paradise. They are submerged in the
sea of torment, yet consider themselves among the beloved of the Lord of
Lords. At the utmost distance from God, they imagine themselves among
those who are near Him.

\subsection{290}\label{section-1845}

\begin{otherlanguage}{persian}

پس شما ای اهل بیان در کلّ حین توکّل بجمال ربّ العالمین نمائید و باو پناه
برید که مبادا بر جمالم در ظهور اخرایم وارد آورید آنچه آن گروه در این حین
وارد آورده‌اند

\end{otherlanguage}

Pas shumā ay ahl-i Bayān dar kull-i ḥīn tavakkul bi-Jamāl-i Rabb
al-ʿĀlamīn namāyīd va bi-ū panāh barīd ki mabādā bar Jamālam dar ẓuhūr-i
ukhrāyam vārid āvarīd ānchi ān gurūh dar īn ḥīn vārid āvarda-and.

Therefore, people of the Bayān, place your trust at all times in the
Beauty of the Lord of the worlds and seek refuge in Him, lest in My next
Revelation you inflict upon My Beauty what those people have inflicted
at this time.

\subsection{291}\label{section-1846}

\begin{otherlanguage}{persian}

ای اهل بیان فئهٴ فرقان حجّتی در دست نداشتند مگر فرقان که فارق بین حقّ و
باطل بوده در آن ایّام و چون جمال عزّ رحمانیم از افق قدس سبحانیم طالع و
مشرق شد بهمان حجّت و دلیل بل اعظم نفس خود را ظاهر فرمودم که شاید از حجّت
قبل بحجّت بعد آگاه شده خود را از حرم تصدیق جمال یقین محروم نسازند و مع‌ذلک
در کلّ ایّام آیات فرقان را تلاوت مینمایند و از این آیات بدیعه که حجّیّتش چون
شمس مابین سماء ظاهر و لایح است ممنوع و محروم گشته‌اند

\end{otherlanguage}

Ay ahl-i Bayān, fiʾa-yi Furqān ḥujjatī dar dast nadāshtand magar Furqān
ki fāriq bayn-i ḥaqq va bāṭil būda dar ān ayyām. Va chūn Jamāl-i ʿizz-i
raḥmāniyam az ufuq-i quds-i subḥāniyam ṭāliʿ va mushriq shud, bi-hamān
ḥujjat va dalīl, bal aʿẓam, nafs-i khud rā ẓāhir farmūdam ki shāyad az
ḥujjat-i qabl bi-ḥujjat-i baʿd āgāh shuda, khud rā az ḥaram-i taṣdīq-i
Jamāl-i yaqīn maḥrūm nasāzand. Va maʿa-dhālik dar kull-i ayyām āyāt-i
Furqān rā tilāvat mīnamāyand va az īn āyāt-i badīʿa ki ḥujjiyyatash chūn
shams mābayn-i samā ẓāhir va lāʾiḥ ast mamnūʿ va maḥrūm gashta-and.

People of the Bayān, the people of the Qurʾān had no proof in hand
except the Qurʾān, which distinguished truth from falsehood in those
days. When My Beauty of merciful glory rose and shone from the horizon
of My transcendent holiness, I manifested Myself with that same proof
and testimony---indeed, with a greater one---so that through the earlier
proof they might recognize the later proof and not deprive themselves of
the sanctuary of acknowledgment of the Beauty of certainty. Yet they
recite the verses of the Qurʾān every day while remaining shut out and
deprived of these wondrous verses, whose authority as proof is as
evident and radiant as the sun in the midst of heaven.

\subsection{292}\label{section-1847}

\begin{otherlanguage}{persian}

چنانچه در حین ظهور آیات عزّیّهٴ الهیّه جمعی مذکور نمودند که این آیات از کتب
قبل سرقت شده و بعضی لسان غافلشان باین تکلّم نموده که این آیات از معدن نفس
و هوی ظاهر کذلک اشهدناهم ان انتم من الشّاهدین تا آنکه امر بمقامی رسید در
حینی که بذکرم توسّل میجستید جسدم را مجروح نمودید و در وقتی که بعرفانم
افتخار مینمودید بر جمال معروفم رصاص کین انداختید

\end{otherlanguage}

Chunānchi dar ḥīn-i ẓuhūr-i āyāt-i ʿizziyya-yi ilāhiyya jamʿī madhkūr
namūdand ki īn āyāt az kutub-i qabl sirqat shuda va baʿḍī lisān-i
ghāfilishān bi-īn takallum namūda ki īn āyāt az maʿdan-i nafs va havā
ẓāhir. Kadhālika ashhadnāhum in antum min al-shāhidīn. Tā ānki amr
bi-maqāmī rasīd, dar ḥīnī ki bi-dhikram tavassul mījustīd jasadam rā
majrūḥ namūdīd va dar vaqtī ki bi-ʿirfānam iftikhār mīnamūdīd bar
Jamāl-i maʿrūfam raṣāṣ-i kīn andākhtīd.

Thus, when the glorious divine verses were revealed, some declared that
these verses had been stolen from earlier books, while others said with
heedless tongues that these verses arose from the source of self and
desire. Thus We made them bear witness, if you are among those who
witness. Matters reached such a point that, even as you sought recourse
through My remembrance, you wounded My body, and while you prided
yourselves on recognizing Me, you fired bullets of malice at My
recognized Beauty.

\subsection{293}\label{section-1848}

\begin{otherlanguage}{persian}

اینست شأن دنیا و اهل آن تا آنکه روح لطیفم از آلایش انفس مشرکه فارغ شده
بمقرّ عزّ ارفع اعلی و رفیق اقدس امنع ابهی طیران نمود و بعد از ارتقای روح
مبارکم بافق ابهی بطرف عنایت و لحاظ مکرمت بمدّعیان محبّتم ناظرم که کدام
بوصایایم عمل مینمایند و بامرم مطیعند

\end{otherlanguage}

Īnast shaʾn-i dunyā va ahl-i ān, tā ānki rūḥ-i laṭīfam az ālāyish-i
anfus-i mushrika fārigh shuda, bi-maqarr-i ʿizz-i arfaʿ-i aʿlā va
Rafīq-i aqdas-i amnaʿ-i abhā ṭayarān namūd. Va baʿd az irtiqā-yi rūḥ-i
mubārakam bi-ufuq-i abhā, bi-ṭarf-i ʿināyat va liḥāẓ-i makramat
bi-muddaʿiyān-i maḥabbatam nāẓiram ki kudām bi-vaṣāyāyam ʿamal
mīnamāyand va bi-amram muṭīʿand.

Such is the way of the world and its people. At last My gentle spirit
was freed from the defilement of idolatrous souls and winged its way to
the loftiest, most exalted seat of glory and the most holy, most
unapproachable, all-glorious Companion. After My blessed spirit ascended
to the all-glorious horizon, I look with the eye of favor and the gaze
of loving-kindness upon those who claim to love Me, to see which of them
act according to My counsels and obey My command.

\subsection{294}\label{section-1849}

\begin{otherlanguage}{arabic}

اذاً ینطق لسان القدم عن افق الأبهی و یقول یا ملأ البیان هذا جمالی قد ظهر
بآیاتی ثمّ ظهوراتی لم کفرتم به و اعرضتم عنه اذاً یثبت بأنّکم ما آمنتم بنفسی
کما ثبت فی حین ظهوری بأنّ ملأ الفرقان ما آمنوا بمحمّد مظهر نفسی کما ظهر فی
ظهور محمّد بأنّ ملأ الانجیل ما آمنوا بابنی کما ظهر حین الّذی جآء الرّوح بأنّ
امّة التّوریة ما آمنوا بالکلیم اذاً فارجعوا ثمّ انظروا الی ان ینتهی الأمر
الی ظهور الأوّل و کذلک نلقی علیکم من اسرار ما کان لعلّ تکوننّ فی انفسکم لمن
الشّاعرین

\end{otherlanguage}

Idhan yanṭiqu lisān al-qidam ʿan ufuq al-Abhā wa-yaqūlu yā malaʾ
al-Bayān, hādhā Jamālī qad ẓahara bi-āyātī thumma ẓuhūrātī, lima
kafartum bihi wa-aʿraḍtum ʿanhu? Idhan yathbutu bi-annakum mā āmantum
bi-nafsī kamā thabata fī ḥīn ẓuhūrī bi-anna malaʾ al-Furqān mā āmanū
bi-Muḥammad maẓhar nafsī, kamā ẓahara fī ẓuhūr Muḥammad bi-anna malaʾ
al-Injīl mā āmanū bi-Ibnī, kamā ẓahara ḥīna alladhī jāʾa al-Rūḥ bi-anna
ummat al-Tawrāt mā āmanū bi-l-Kalīm. Idhan fa-rjiʿū thumma unẓurū ilā an
yantahiya al-amru ilā ẓuhūr al-awwal. Wa-kadhālika nulqī ʿalaykum min
asrār mā kāna laʿalla takūnunna fī anfusikum la-min al-shāʿirīn.

Then the Tongue of Eternity speaks from the all-glorious horizon and
says: `People of the Bayān, this is My Beauty, who has appeared with My
verses and My manifestations. Why have you disbelieved in Him and turned
away from Him? It is thereby established that you have not believed in
Me, just as at the time of My Revelation it was established that the
people of the Qurʾān had not believed in Muḥammad, the Manifestation of
My Self; just as in the Revelation of Muḥammad it became evident that
the people of the Gospel had not believed in My Son; and just as, when
the Spirit came, it became evident that the people of the Torah had not
believed in Him who conversed with God. Go back, then, and consider,
until the matter reaches the appearance of the first. Thus We impart to
you some of the mysteries of what has been, so that within yourselves
you may be among those who perceive.'

\subsection{295}\label{section-1850}

\begin{otherlanguage}{persian}

و از شهادتم ایّامی نگذشته که عنایت جدید از شطر قدس ابهی و مکرمت منیع از
افق عزّ اعلی اشراق فرمود و ساذج قدم بجمال اعظم اکرم از رضوان غیب ظاهر شد
بهمان حجّتی که من حجّت قرار دادم و بهمان برهان که عنداللّه مقبول بوده بلکه
بجمیع شئونات احدیّت و ظهورات عزّ صمدیّت و بطونات غیب لایدرک و دلالات عزّ
لایعرف ظاهر شده

\end{otherlanguage}

Va az shahādatam ayyāmī nagudhashta ki ʿināyat-i jadīd az shaṭr-i quds-i
abhā va makramat-i manīʿ az ufuq-i ʿizz-i aʿlā ishrāq farmūd va Sādhij-i
Qidam bi-Jamāl-i aʿẓam-i akram az Riḍvān-i ghayb ẓāhir shud, bi-hamān
ḥujjatī ki man ḥujjat qarār dādam va bi-hamān burhān ki ʿinda Allāh
maqbūl būda, balki bi-jamīʿ-i shuʾūnāt-i aḥadiyyat va ẓuhūrāt-i ʿizz-i
ṣamadiyyat va buṭūnāt-i ghayb-i lāyudrak va dalālāt-i ʿizz-i lāyuʿraf
ẓāhir shuda.

Only a few days had passed since My martyrdom when a fresh favor shone
from the region of all-glorious holiness and an exalted loving-kindness
from the horizon of the highest glory. The pure Being of Eternity
appeared from the paradise of the unseen in a most great and most
bountiful Beauty, with the very proof that I had established as proof
and the very testimony that had been acceptable before God. Indeed, He
appeared with all the aspects of oneness, the manifestations of the
glory of eternal self-sufficiency, the hidden depths of the unfathomable
unseen, and the signs of unknowable glory.

\subsection{296}\label{section-1851}

\begin{otherlanguage}{persian}

مع‌ذلک شما ای ملأ بیان از کلّ جهات باسیاف غلّ و اشارات بر حول عرش اعظم جمع
شده‌اید و در کلّ حین از سهام کین بر این جمال عزّ منیع وارد آورده‌اید قسم
بجمالم اگر حال ملاحظه نمائید و بطرف حقیقت برفیق اعلی متوجّه شوید ملاحظه
مینمائید که از جسدم بحور دم جاری و از ارکانم آثار اسیاف کین ظاهر

\end{otherlanguage}

Maʿa-dhālik shumā ay malaʾ-i Bayān az kull-i jihāt bi-asyāf-i ghill va
ishārāt bar ḥawl-i ʿarsh-i aʿẓam jamʿ shuda-īd va dar kull-i ḥīn az
sihām-i kīn bar īn Jamāl-i ʿizz-i manīʿ vārid āvarda-īd. Qasam
bi-Jamālam, agar ḥāl mulāḥiẓa namāyīd va bi-ṭarf-i ḥaqīqat bi-Rafīq-i
Aʿlā mutavajjih shavīd, mulāḥiẓa mīnamāyīd ki az jasadam buḥūr-i dam
jārī va az arkānam āthār-i asyāf-i kīn ẓāhir.

Yet you, people of the Bayān, have gathered from every direction around
the most great Throne with the swords of rancor and insinuation, and at
every moment you have struck this Beauty of exalted glory with the
arrows of malice. I swear by My Beauty! If you look now and turn toward
the truth, toward the most exalted Companion, you will see seas of blood
flowing from My body and the marks of the swords of malice evident upon
My limbs.

\subsection{297}\label{section-1852}

\begin{otherlanguage}{persian}

آخر تفکّر ننموده‌اید که بچه سبب نفسم را انفاق فرمودم و آن همه سیوف فحشا و
رصاص بغضا بر او وارد آمد و این بسی واضح است که مقصودی جز عرفان مظهر نفسم
نبوده و چون مظهر نفسم بتمام ظهور ظاهر شد اینگونه معمول داشته‌اید که مشهود
شده

\end{otherlanguage}

Ākhir tafakkur nanamūda-īd ki bi-chi sabab nafsam rā infāq farmūdam va
ān hama suyūf-i faḥshā va raṣāṣ-i baghḍā bar ū vārid āmad? Va īn basī
vāḍiḥ ast ki maqṣūdī juz ʿirfān-i Maẓhar-i nafsam nabūda va chūn
Maẓhar-i nafsam bi-tamām-i ẓuhūr ẓāhir shud, īngūna maʿmūl dāshta-īd ki
mashhūd shuda.

Have you not reflected, after all, on why I sacrificed Myself and why
all those swords of wickedness and bullets of hatred struck Me? It is
abundantly clear that My purpose was nothing but recognition of the
Manifestation of My Self. Yet when the Manifestation of My Self appeared
in the fullness of revelation, you treated Him in the manner that has
been witnessed.

\subsection{298}\label{section-1853}

\begin{otherlanguage}{persian}

حال اگر در موقف حشر اکبر از شما سؤال شود و از آنچه بآن عامل شده‌اید
استفسار رود در جواب ربّ الأرباب چه خواهید گفت لا واللّه ابداً قوّت تکلّم
نخواهید داشت چه که ایّامی از غیبتم نگذشته و جمیع حجّیّت ظهور و امور وارده
را بچشم خود مشاهده نموده‌اید و مع‌ذلک از جمالی که بقول او کلّ شئونات قبلیّه
و بعدیّه محقّق شده غافل شده‌اید و بحجبات نفس و هوی از منظر اعلی محجوب
مانده‌اید

\end{otherlanguage}

Ḥāl agar dar mawqif-i ḥashr-i akbar az shumā suʾāl shavad va az ānchi
bi-ān ʿāmil shuda-īd istifsār ravad, dar javāb-i Rabb al-Arbāb chi
khāhīd guft? Lā wa-Allāh, abadan quvvat-i takallum nakhāhīd dāsht, chi
ki ayyāmī az ghaybatam nagudhashta va jamīʿ-i ḥujjiyyat-i ẓuhūr va
umūr-i vārida rā bi-chashm-i khud mushāhida namūda-īd va maʿa-dhālik az
Jamālī ki bi-qawl-i ū kull-i shuʾūnāt-i qabliyya va baʿdiyya muḥaqqaq
shuda ghāfil shuda-īd va bi-ḥujubāt-i nafs va havā az Manẓar-i Aʿlā
maḥjūb mānda-īd.

Now, if you are questioned at the place of the greatest gathering and
asked about what you have done, what will you say in answer to the Lord
of Lords? No, by God! You will have no power to speak, for only a few
days have passed since My absence, and you have witnessed with your own
eyes the full proof of the Revelation and all that befell Me. Yet you
have become heedless of the Beauty through whose word all things, past
and future, have been brought into being, and have remained shut out
from the most exalted Vision by the veils of self and desire.

\subsection{299}\label{section-1854}

\begin{otherlanguage}{persian}

ای قوم در تمام اوراق بیان کلّ را باین ظهور قدس صمدانی بشارت دادم و فرمودم
که مباد در حین طلوع این نیّر عزّ اعظم بشیئی از آنچه خلق شده مابین سموات و
ارض مشغول شوید و از جمال قدم محتجب مانید

\end{otherlanguage}

Ay qawm, dar tamām-i awrāq-i Bayān kull rā bi-īn ẓuhūr-i quds-i ṣamadānī
bishārat dādam va farmūdam ki mabād dar ḥīn-i ṭulūʿ-i īn nayyir-i ʿizz-i
aʿẓam bi-shayʾī az ānchi khalq shuda mābayn-i samāvāt va arḍ mashghūl
shavīd va az Jamāl-i Qidam muḥtajib mānīd.

People, throughout all the pages of the Bayān I gave everyone the glad
tidings of this holy and eternal Revelation, and said that when this
Luminary of supreme glory arose, you must not become occupied with
anything created between the heavens and the earth and remain veiled
from the Beauty of Eternity.

\subsection{300}\label{section-1855}

\begin{otherlanguage}{persian}

و هم‌چنین تصریحاً فرمودم که ایّاکم ای قوم اگر در حین ظهور بواحد بیان محتجب
شوید و این معلوم بوده که واحد بیان اوّل خلق بیان عنداللّه محسوب و بر کلّ
سبقت داده‌ام

\end{otherlanguage}

Va hamchunīn taṣrīḥan farmūdam ki iyyākum ay qawm agar dar ḥīn-i ẓuhūr
bi-vāḥid-i Bayān muḥtajib shavīd va īn maʿlūm būda ki vāḥid-i Bayān
avval-i khalq-i Bayān ʿinda Allāh maḥsūb va bar kull sabqat dāda-am.

I also expressly said: beware, people, lest at the time of the
Revelation you be veiled by the Vāḥid of the Bayān. It was known that
the Vāḥid of the Bayān is reckoned before God as the first of the
creation of the Bayān, and that I have given it precedence over all.

\subsection{301}\label{section-1856}

\begin{otherlanguage}{persian}

و هم‌چنین بنصّ صریح فرموده‌ام که ایّاکم ای قوم اگر حین ظهور محتجب شوید بآنچه
نازل شده در بیان و حال شما ای قوم بوصفی که در بیان نازل شده و نمیدانید
که مقصود چیست و در بارهٴ کیست چه که ابداً نفسی از مقصودم مطّلع نبوده مع‌ذلک
از موجد وصف و منزل و مظهر و محقّق آن که بقولی از او اینگونه هیاکل خلق شده
و میشوند معرض شده‌اید و کاش باین اکتفا مینمودید بلکه بر قتلش فتوی داده‌اید
واللّه در کلّ احیان مثل ثعبان این جوهر رحمن را اذیّت نموده‌اید مؤمن نشدید و
بحقّ عارف نگشتید دیگر بر قتل و ضرّش چرا راضی شدید

\end{otherlanguage}

Va hamchunīn bi-naṣṣ-i ṣarīḥ farmūda-am ki iyyākum ay qawm agar ḥīn-i
ẓuhūr muḥtajib shavīd bi-ānchi nāzil shuda dar Bayān. Va ḥāl shumā ay
qawm bi-vaṣfī ki dar Bayān nāzil shuda va namīdānīd ki maqṣūd chīst va
dar bāra-yi kīst, chi ki abadan nafsī az maqṣūdam muṭṭaliʿ nabūda,
maʿa-dhālik az mūjid-i vaṣf va munzil va muẓhir va muḥaqqiq-i ān ki
bi-qawlī az ū īngūna hayākil khalq shuda va mīshavand muʿriḍ shuda-īd.
Va kāsh bi-īn iktifā mīnamūdīd, balki bar qatlash fatvā dāda-īd.
Va-Allāh dar kull-i aḥyān mithl-i thuʿbān īn Jawhar-i Raḥmān rā adhiyyat
namūda-īd. Muʾmin nashudīd va bi-Ḥaqq ʿārif nagashtīd, dīgar bar qatl va
ḍarrash chirā rāḍī shudīd?

I also stated in an explicit text: beware, people, lest at the time of
the Revelation you be veiled by what has been revealed in the Bayān. Yet
now, people, because of a description revealed in the Bayān---when you
do not know what is intended or whom it concerns, for no one has ever
been aware of My intention---you have turned away from the One who
brought that description into being, revealed it, manifested it, and
fulfilled it, through a single word from whom such temples have been and
continue to be created. If only you had been content with this! Instead,
you have issued a decree for His death. By God, at every moment you have
tormented this Essence of the All-Merciful like a serpent. You did not
believe and did not recognize God; why, then, did you consent to His
death and injury?

\subsection{302}\label{section-1857}

\begin{otherlanguage}{arabic}

قسم بجمالم که از کأس انصاف ننوشیده‌اید و از بحر رضایم نچشیده‌اید بیمین عدل
قدم نگذارده‌اید و در کوی متبصّرین مرور ننموده‌اید ضلالت را نفس هدایت
شمرده‌اید و صرف شرک را جوهر توحید دانسته‌اید

\end{otherlanguage}

Qasam bi-Jamālam ki az kaʾs-i inṣāf nanūshīda-īd va az baḥr-i riḍāyam
nachishīda-īd, bi-yamīn-i ʿadl qadam nagudhāshta-īd va dar kūy-i
mutabaṣṣirīn murūr nanamūda-īd. Ḍalālat rā nafs-i hidāyat shumurda-īd va
ṣirf-i shirk rā jawhar-i tawḥīd dānista-īd.

I swear by My Beauty that you have not drunk from the cup of fairness or
tasted the ocean of My good pleasure. You have not set foot on the right
side of justice or passed through the quarter of those endowed with
insight. You have counted error as guidance itself and regarded sheer
idolatry as the essence of divine unity.''

\subsection{303}\label{section-1858}

\begin{otherlanguage}{persian}

و جمیع شما بعین ظاهر دیده‌اید نفسی را که بیست سنه حفظ فرمود و الآن موجود
است و با قدرت بر او و اطّلاع به ما فی سرّ او در کنف حفظ رحمانیّتش حفظ فرمود
و مع‌ذلک بر قتلش قیام نمود و چون اراده و فعلش بین مهاجرین انتشار یافت لذا
محض ستر اعمال شنیعهٴ خود و القای شبهه در قلوب متوهّمه بمفتریات مجعوله
مشغول گشت و افعال و اعمال خود را بساذج قدم نسبت داده که شاید عباد را از
یمین یقین بشمال وهم کشاند

\end{otherlanguage}

Va jamīʿ-e shomā be-ʿayn-e ẓāher dīde-īd nafsī rā ke bīst sana ḥefẓ
farmūd va al-ān mowjūd ast va bā qodrat bar ū va eṭṭelāʿ be mā fī sirri
ū dar kanaf-e ḥefẓ-e raḥmāniyyatash ḥefẓ farmūd va maʿa-dhālek bar
qatl-ash qiyām namūd va chūn erāda va feʿl-ash bayn-e mohājerīn enteshār
yāft ledhā maḥż-e satr-e aʿmāl-e shanīʿa-ye khod va elqā-ye shobha dar
qolūb-e motavahhema be-moftarayāt-e majʿūla mashghūl gasht va afʿāl va
aʿmāl-e khod rā be-sādhaj-e qedam nesbat dāde ke shāyad ʿebād rā az
yamīn-e yaqīn be-shomāl-e vahm kashānad.

All of you have seen with your own outward eyes the person whom He
protected for twenty years and who is still alive. Though He had power
over him and knew what lay hidden within him, He sheltered him within
His merciful protection. Nevertheless, that person arose to kill Him.
When his intention and deed became known among the emigrants, he
occupied himself with fabricated accusations solely to conceal his own
vile deeds and cast doubt into minds prone to imagining things. He
attributed his own acts and deeds to the Essence of Eternity, hoping
thereby to draw the servants from the right hand of certitude to the
left hand of delusion.

\subsection{304}\label{section-1859}

\begin{otherlanguage}{persian}

و شما آنچه را ببصر دیده‌اید و ادراک نموده‌اید انکار نموده و بمفتریات
مجعوله از طلعت احدیّه محتجب مانده‌اید قسم بجمال عزّ تقدیسم که آنچه در بارهٴ
این جمال مبین نوشته‌اند حجّت است بر کلّ در کذب قائلین و اثبات این امر مبین
چه که نسبتهائی ذکر نموده که کذبش اظهر از شمس است در وسط سماء چه که جمیع
با این جمال عزّ اعلی معاشرت نموده‌اید و بقدر خود ادراک نموده

\end{otherlanguage}

Va shomā ānche rā be-baṣar dīde-īd va edrāk namūde-īd enkār namūde va
be-moftarayāt-e majʿūla az ṭalʿat-e aḥadiyya moḥtajeb mānde-īd. Qasam
be-jamāl-e ʿezz-e taqdīsam ke ānche dar bāra-ye īn jamāl-e mobīn
neveshte-and ḥojjat ast bar koll dar kedhb-e qāʾelīn va ethbāt-e īn
amr-e mobīn che ke nesbat-hāʾī dhekr namūde ke kedhb-ash aẓhar az shams
ast dar vasaṭ-e samāʾ che ke jamīʿ bā īn jamāl-e ʿezz-e aʿlā moʿāsherat
namūde-īd va be-qadr-e khod edrāk namūde.

You have denied what you have seen with your own eyes and perceived, and
through fabricated accusations you have remained veiled from the
Countenance of Oneness. I swear by the beauty of My transcendent
holiness that what they have written about this manifest Beauty is proof
to all of the falsehood of those who utter it and confirmation of this
manifest Cause. For they have made allegations whose falsehood is more
evident than the sun in the midst of heaven, since you have all
associated with this Beauty of most exalted glory and have perceived
according to your own capacity.

\subsection{305}\label{section-1860}

\begin{otherlanguage}{persian}

و مع‌ذلک باین کلمات مجعوله از منزل آیات احدیّه معرض شده و سلّمنا که صادقند
آیا حقّ را یفعل ما یشآء ندانسته‌اید و بأنّه یحکم ما یرید موقن نشده‌اید

\end{otherlanguage}

Va maʿa-dhālek be-īn kalemāt-e majʿūla az monzel-e āyāt-e aḥadiyya
moʿreż shode va sallamnā ke ṣādeq-and āyā Ḥaqq rā yafʿalu mā yashāʾ
nadāneste-īd va bi-annahu yaḥkumu mā yurīd mūqen nashode-īd.

Nevertheless, because of these fabricated words you have turned away
from the Revealer of the verses of Oneness. Even if we concede that they
are true, have you not known that the True One does whatever He wills,
and have you not become certain that He decrees whatever He desires?

\subsection{306}\label{section-1861}

\begin{otherlanguage}{arabic}

از جمیع این مراتب گذشته این جمعی که در این سفر مع اللّه هجرت نموده و
اکثری از امور را ببصر و فؤاد خود ادراک نموده‌اند و شهادت میدهند که حقّ جلّ
شأنه بجمیع شئونات از کلّ ما سواه ممتاز بوده و خواهد بود مع‌ذلک این نفوس را
کاذب دانسته‌اید و کسانی که اصلاً مطّلع بر امر نبوده و نخواهند بود و در
مناهج وهم و ظنّ سالکند صادق دانسته و میدانید فأفّ لکم یا ملأ المتوهّمین

\end{otherlanguage}

Az jamīʿ-e īn marāteb godhashte īn jamʿī ke dar īn safar maʿa Allāh
hejrat namūde va aktharī az omūr rā be-baṣar va foʾād-e khod edrāk
namūde-and va shahādat mīdehand ke Ḥaqq jalla shaʾnuhu be-jamīʿ-e
shoʾūnāt az koll-e mā-sevāh momtāz būde va khvāhad būd maʿa-dhālek īn
nofūs rā kādheb dāneste-īd va kasānī ke aṣlan moṭṭaleʿ bar amr nabūde va
nakhvāhand būd va dar manāhej-e vahm va ẓann sālek-and ṣādeq dāneste va
mīdānīd. Fa-uffin lakum yā malaʾa al-mutawahhimīn.

Apart from all these considerations, this company who have emigrated
with God on this journey have perceived most of these matters with their
own eyes and hearts, and testify that the True One---exalted be His
glory---has been and will remain distinguished in every respect from all
besides Him. Nevertheless, you have deemed these souls liars, while you
have regarded, and still regard, as truthful those who have never been
aware of the matter and never will be, and who walk the paths of
delusion and conjecture. Shame on you, you company of the deluded!

\subsection{307}\label{section-1862}

\begin{otherlanguage}{persian}

قسم بجمالم که حجّت بالغهٴ الهی بر کلّ تمام شده و کلمهٴ تامّه‌اش از افق صدق
مشرق گشته و الیوم پناهی جز پناهش نه و ظلّی جز ظلّش مشهود نه بشنوید ندای
محبوب خود را و آنچه ببصر دیده‌اید متمسّک شوید و از عروهٴ وهم بگسلید

\end{otherlanguage}

Qasam be-jamālam ke ḥojjat-e bālegha-ye elāhī bar koll tamām shode va
kalema-ye tāmma-ash az ofoq-e ṣedq moshreq gashte va al-yowm panāhī joz
panāh-ash na va ẓellī joz ẓell-ash mashhūd na. Beshnavīd nedā-ye
maḥbūb-e khod rā va ānche be-baṣar dīde-īd motamassek shavīd va az
ʿorva-ye vahm bogsalīd.

I swear by My Beauty that God's conclusive proof has been fully
established for all, and His perfect Word has dawned from the horizon of
truth. Today there is no refuge but His refuge, and no sheltering shade
is to be seen but His. Listen to the call of your Beloved, hold fast to
what you have seen with your own eyes, and loose your hold on the handle
of delusion.

\subsection{308}\label{section-1863}

\begin{otherlanguage}{arabic}

و اگر ببصر حقیقت ملاحظه نمائید تاللّه اهل ملأ اعلی در نوحه و ندبه مشغولند
و جمیع حوریّات غرفات در حنین و ناله اوراق سدرهٴ منتهی از ظلم این ظالمان
پژمرده گشته تاللّه الحقّ اریاح رحمت رحمن از شطر امکان مقطوع و اشراقات
انوار وجه سبحان از اهل اکوان ممنوع

\end{otherlanguage}

Va agar be-baṣar-e ḥaqīqat molāḥeẓa namāʾīd ta-Allāh ahl-e malaʾ-e aʿlā
dar nowḥa va nodba mashghūl-and va jamīʿ-e ḥūriyyāt-e ghorofāt dar ḥanīn
va nāla. Owrāq-e sedra-ye montahā az ẓolm-e īn ẓālemān pazhmorda gashte.
Ta-Allāh al-Ḥaqq aryāḥ-e raḥmat-e Raḥmān az shaṭr-e emkān maqṭūʿ va
eshrāqāt-e anvār-e vajh-e Sobḥān az ahl-e akvān mamnūʿ.

If you look with the eye of truth, by God, the inhabitants of the
Concourse on high are mourning and lamenting, and all the maidens of the
heavenly chambers are sighing and wailing. The leaves of the Lote-Tree
of the Utmost Boundary have withered through the oppression of these
oppressors. By God, the True One, the breezes of the All-Merciful's
mercy have been cut off from the realm of contingent being, and the
radiance of the lights of the All-Glorious Countenance has been withheld
from the inhabitants of creation.

\subsection{309}\label{section-1864}

\begin{otherlanguage}{persian}

تاللّه الحقّ ظلمی نموده‌اید که کلّ اشیاء از حیات خود منقطع شده‌اند و الیوم
خلقی باقی نه چه از اهل ملأ اعلی و چه از اهل مداین بقا و چه عاکفین لجّهٴ
اسماء مگر آنکه کلّ لطایف سرور را بحزن تبدیل نموده‌اند و قمیص سود پوشیده‌اند
و جمیع ملأ کرّوبیین و حقایق انبیا و مرسلین در غرفات عزّ تمکین بنوحه مشغول
و شما ای غافلین ارض مسرورید و در ارض هاویه بکمال فرح سیر مینمائید

\end{otherlanguage}

Ta-Allāh al-Ḥaqq ẓolmī namūde-īd ke koll-e ashyāʾ az ḥayāt-e khod
monqaṭeʿ shode-and va al-yowm khalqī bāqī na che az ahl-e malaʾ-e aʿlā
va che az ahl-e madāyen-e baqā va che ʿākefīn-e lojja-ye asmāʾ magar
ānke koll-e laṭāyef-e sorūr rā be-ḥozn tabdīl namūde-and va qamīṣ-e sūd
pūshīde-and va jamīʿ-e malaʾ-e karūbiyyīn va ḥaqāʾeq-e anbiyā va
morsalīn dar ghorofāt-e ʿezz-e tamkīn be-nowḥa mashghūl va shomā ey
ghāfelīn-e arż masrūrīd va dar arż-e hāviya be-kamāl-e faraḥ seyr
mīnamāʾīd.

By God, the True One, you have committed such oppression that all things
have been severed from their own life. Today no creature
remains---whether among the inhabitants of the Concourse on high, the
dwellers in the cities of eternity, or those abiding in the ocean of
names---who has not exchanged all the delights of joy for sorrow and put
on a garment of black. The entire company of the cherubim and the
realities of the Prophets and Messengers are lamenting in the chambers
of established glory, while you, heedless ones of the earth, are joyful
and wander through the land of the abyss in utter delight.

\subsection{310}\label{section-1865}

\begin{otherlanguage}{persian}

جوهر دین را کشته‌اید و بگمان خود بر سرر دین و ایقان جالسید فواللّه یا قوم
شبه این ظهور ظاهر نشده و چشم امکان ندیده بشنوید ندایم را و نباشید از
احتیاط‌کنندگانی که در ملأ فرقان بودند بشأنی که احتیاط از دم بعوضه
مینمودند و بر سفک دم اللّه فتوی میدادند

\end{otherlanguage}

Jowhar-e dīn rā koshte-īd va be-gomān-e khod bar soror-e dīn va īqān
jāles-īd. Fa-wa-Allāh yā qawm shebh-e īn ẓohūr ẓāher nashode va chashm-e
emkān nadīde. Beshnavīd nedāyam rā va nabāshīd az eḥtiyāṭ-konandegānī ke
dar malaʾ-e Forqān būdand be-shaʾnī ke eḥtiyāṭ az dam-e baʿūża
mīnamūdand va bar safk-e dam-e Allāh fatvā mīdādand.

You have slain the essence of religion, yet imagine yourselves seated on
the thrones of religion and certitude. By God, O people, nothing like
this Revelation has ever appeared, nor has the eye of contingent being
ever beheld its like. Hear My call, and do not be among those
punctilious ones who were found among the people of the Qur'an: so
scrupulous were they that they took precautions over the blood of a
gnat, yet issued rulings authorizing the shedding of the blood of God.

\subsection{311}\label{section-1866}

\begin{otherlanguage}{persian}

بذکر مشغول بودند و چون آیات سلطان ذکر بر آن گروه القا میشد صیحه میزدند
که این اذکار را بگذار و ما را از ذکر اللّه غافل مکن این بوده اعمال و
افعال آن گروه که مشهود گشت و شما ای اهل بیان در این ایّام بجوهر آن اعمال
مرتکب و عاملید و خود را از حقّ شمرده و میشمرید

\end{otherlanguage}

Be-dhekr mashghūl būdand va chūn āyāt-e solṭān-e dhekr bar ān gorūh elqā
mīshod ṣayḥa mīzadand ke īn adhkār rā bogdhār va mā rā az dhekr-e Allāh
ghāfel makon. Īn būde aʿmāl va afʿāl-e ān gorūh ke mashhūd gasht va
shomā ey ahl-e Bayān dar īn ayyām be-jowhar-e ān aʿmāl mortakeb va
ʿāmel-īd va khod rā az Ḥaqq shomorde va mīshomorīd.

They were occupied with remembrance, yet when the verses of the
Sovereign of Remembrance were recited to that company, they cried out,
``Leave aside these remembrances, and do not distract us from the
remembrance of God!'' Such were the deeds and actions witnessed from
that company. And you, people of the Bayān, are committing and carrying
out the very essence of those deeds in these days, while you have
reckoned, and still reckon, yourselves to be of the True One.

\subsection{312}\label{section-1867}

\begin{otherlanguage}{arabic}

اذاً اشهد اللّه و ملائکته و انبیائه و رسله و الّذینهم یطوفنّ فی حول عرشه و
کلّ ما خلق فی السّموات و الأرض بأنّی ما قصّرت فی تبلیغی ایّاکم و بلّغتکم
رسالات اللّه حین ظهوری و حین ارتقائی و هذا الحین الّذی اظهرت نفسی عن افق
الأبهی و القیت علیکم الحکمة و البیان و عرّفتکم جمال الرّحمن و اتممت الحجّة
لکم و الدّلیل علیکم و البرهان فیکم و ما بقی من ذکر الّا و قد القیتکم

\end{otherlanguage}

Idhan ush'hidu Allāha wa-malāʾikatahu wa-anbiyāʾahu wa-rusulahu
wa-alladhīna hum yaṭūfunna fī ḥawli ʿarshihi wa-kulla mā khalaqa fī
al-samāwāti wa-al-arḍi bi-annī mā qaṣṣartu fī tablīghī iyyākum
wa-ballaghtukum risālāti Allāhi ḥīna ẓuhūrī wa-ḥīna irtiqāʾī wa-hādhā
al-ḥīna alladhī aẓhartu nafsī ʿan ufuqi al-Abhā wa-alqaytu ʿalaykumu
al-ḥikmata wa-al-bayāna wa-ʿarraftukum jamāla al-Raḥmāni wa-atmamtu
al-ḥujjata lakum wa-al-dalīla ʿalaykum wa-al-burhāna fīkum wa-mā baqiya
min dhikrin illā wa-qad alqaytukum.

Thus I call God to witness, and His angels, His Prophets, His
Messengers, those who circle around His throne, and all that He has
created in the heavens and on earth, that I have not fallen short in
conveying My message to you. I have delivered to you the messages of God
at the time of My manifestation, at the time of My ascent, and at this
time when I have manifested Myself from the All-Glorious Horizon. I have
imparted to you wisdom and utterance, made known to you the beauty of
the All-Merciful, and completed the testimony for you, the evidence
against you, and the proof among you. No admonition remains that I have
not conveyed to you.

\subsection{313}\label{section-1868}

\begin{otherlanguage}{arabic}

اذاً یا الهی انت تعلم بأنّی ما قصّرت فی امرک و بلّغت هؤلآء ما امرتنی به قبل
خلق السّموات و الأرض و بیّنت لهم مناهج عدلک و اظهرت لهم مسالک رضائک

\end{otherlanguage}

Idhan yā ilāhī anta taʿlamu bi-annī mā qaṣṣartu fī amrika wa-ballaghtu
hāʾulāʾi mā amartanī bihi qabla khalqi al-samāwāti wa-al-arḍi
wa-bayyantu lahum manāhija ʿadlika wa-aẓhartu lahum masālika riḍāʾika.

Thus, O my God, You know that I have not fallen short in Your Cause, and
that I have conveyed to these people what You commanded Me before the
creation of the heavens and the earth. I have explained to them the ways
of Your justice and shown them the paths of Your good pleasure.

\subsection{314}\label{section-1869}

\begin{otherlanguage}{arabic}

اذاً یا الهی فارحم علی هؤلآء و لا تجعلهم من الّذین اعرضوا عنّی و انکروا حقّی
و جادلوا بآیاتی الی ان سفکوا دمی و تقطعوا جوارحی اذاً یا الهی ایّدهم علی
امرک ثمّ انصرهم بنصرتک و لا تجعلهم محروماً عن هذه النّفحات الّتی هبّت عن هذا
الرّضوان الّذی خلقته فی قطب الجنان و لا تمنعهم عن فوحات الّتی ارسلته عن افق
اسمک الرّحمن

\end{otherlanguage}

Idhan yā ilāhī fa-irḥam ʿalā hāʾulāʾi wa-lā tajʿalhum mina alladhīna
aʿraḍū ʿannī wa-ankarū ḥaqqī wa-jādalū bi-āyātī ilā an safakū damī
wa-taqaṭṭaʿū jawāriḥī. Idhan yā ilāhī ayyidhum ʿalā amrika thumma
unṣurhum bi-nuṣratika wa-lā tajʿalhum maḥrūman ʿan hādhihi al-nafaḥāti
allatī habbat ʿan hādhā al-riḍwāni alladhī khalaqtahu fī quṭbi al-jināni
wa-lā tamnaʿhum ʿan fawaḥāti allatī arsaltahu ʿan ufuqi ismika
al-Raḥmān.

Thus, O my God, have mercy on these people, and do not make them among
those who turned away from Me, denied My right, and disputed My verses
until they shed My blood and dismembered My limbs. O my God, strengthen
them in Your Cause, then aid them with Your assistance. Do not deprive
them of these breezes that have blown from this paradise which You have
created at the heart of the gardens of heaven, and do not withhold from
them the fragrances You have sent forth from the horizon of Your name,
the All-Merciful.

\subsection{315}\label{section-1870}

\begin{otherlanguage}{arabic}

اذاً فاحدث یا الهی فی صدورهم من نور کلماتک نار انجذابک لیقلبهم من قدرة
المحضة الی یمین عرش رحمانیّتک ثمّ اشتعل یا الهی فی قلوبهم مشاعل عشقک و
اشتیاقک لیحترق بها حجبات الّتی منعتهم عن ساحة قربک و لقائک

\end{otherlanguage}

Idhan fa-aḥdith yā ilāhī fī ṣudūrihim min nūri kalimātika nāra
injidhābika li-yaqlibahum min qudrati al-maḥḍati ilā yamīni ʿarshi
raḥmāniyyatika thumma ishtaʿil yā ilāhī fī qulūbihim mashāʿila ʿishqika
wa-ishtiyāqika li-yaḥtariqa bihā ḥujubāti allatī manaʿathum ʿan sāḥati
qurbika wa-liqāʾika.

Thus, O my God, bring into being within their breasts, from the light of
Your words, the fire of attraction to You, that through sheer power it
may turn them toward the right hand of the throne of Your mercy. Then
kindle within their hearts, O my God, the torches of love and longing
for You, that through them the veils that have barred them from the
court of Your nearness and Your presence may be burned away.

\subsection{316}\label{section-1871}

\begin{otherlanguage}{arabic}

ثمّ خذ یا الهی ایادیهم بأیدی القدرة و الاقتدار ثمّ انقذهم عن غمرات الوهم و
الهوی و بلّغهم الی مقرّ الّذی قدّسته عن اشارات کلّ ما خلق بین الأرض و السّمآء
ثمّ الق علیهم کلمة الّتی بها تجذب افئدة العارفین الی سمآء عزّ الطافک و قلوب
المقرّبین الی هوآء قدس افضالک

\end{otherlanguage}

Thumma khudh yā ilāhī ayādiyahum bi-aydī al-qudrati wa-al-iqtidāri
thumma anqidhhum ʿan ghamarāti al-wahmi wa-al-hawā wa-ballighhum ilā
maqarri alladhī qaddastahu ʿan ishārāti kulli mā khuliqa bayna al-arḍi
wa-al-samāʾi thumma alqi ʿalayhim kalimata allatī bihā tujdhabu afʾidatu
al-ʿārifīna ilā samāʾi ʿizzi alṭāfika wa-qulūbu al-muqarrabīna ilā
hawāʾi qudsi afḍālika.

Then take their hands, O my God, with the hands of power and might.
Rescue them from the depths of delusion and desire, and bring them to
the abode that You have sanctified beyond the allusions of all that has
been created between earth and heaven. Then impart to them the word
through which the hearts of those who know You are drawn to the heaven
of the glory of Your loving-kindness, and the hearts of those brought
near to You are drawn to the atmosphere of the holiness of Your favors.

\subsection{317}\label{section-1872}

\begin{otherlanguage}{persian}

ثمّ اجعلهم یا محبوبی من الّذین ما منعتهم کلّ من فی السّموات و الأرض عن
التّوجّه الی شطر عنایتک و الاستقرار علی امرک و الاعتراف بحضرتک و الایقان
بلقائک و انّک انت الغفور الرّحیم المعطی العزیز الباذل النّاصر الکریم

\end{otherlanguage}

Thumma ijʿalhum yā maḥbūbī mina alladhīna mā manaʿathum kullu man fī
al-samāwāti wa-al-arḍi ʿani al-tawajjuhi ilā shaṭri ʿināyatika
wa-al-istiqrāri ʿalā amrika wa-al-iʿtirāfi bi-ḥaḍratika wa-al-īqāni
bi-liqāʾika wa-innaka anta al-Ghafūru al-Raḥīmu al-Muʿṭī al-ʿAzīzu
al-Bādhilu al-Nāṣiru al-Karīm.

Then make them, O my Beloved, among those whom all who are in the
heavens and on earth have not prevented from turning toward Your loving
care, remaining steadfast in Your Cause, acknowledging Your presence,
and attaining certainty of meeting You. You are indeed the
Ever-Forgiving, the Merciful, the Bestower, the Almighty, the
Munificent, the Helper, the All-Bountiful.

\bookmarksetup{startatroot}

\chapter{Surih-i-Haykal (Surah of the
Temple)}\label{surih-i-haykal-surah-of-the-temple}

\section{BH00007}\label{bh00007}

\emph{Translated on 16 Qudrat 183 using GPT 6.1 Sol}

\subsection{1}\label{section-1873}

\begin{otherlanguage}{arabic}

هذه سورة الهيكل قد جعلها الله مرآت أسمآئه بين السّموات والأرض وجعلها آية
ذكره بين العالمين\hspace{0pt}

\end{otherlanguage}

Hādhihi sūratu al-haykal, qad jaʿalahā Allāhu mirʾāta asmāʾihi bayna
al-samāwāti wa-l-arḍ, wa-jaʿalahā āyata dhikrihi bayna al-ʿālamīn.

This is the Surah of the Temple. God has made it the mirror of His names
between the heavens and the earth, and the sign of His remembrance among
all peoples.

\subsection{2}\label{section-1874}

\begin{otherlanguage}{arabic}

هو الأبدع الأبهی

\end{otherlanguage}

Huwa al-Abdaʿ al-Abhā.

He is the Most Wondrous, the All-Glorious.

\subsection{3}\label{section-1875}

\begin{otherlanguage}{arabic}

سبحان الّذی نزّل الآیات لقوم یفقهون سبحان الّذی ینزّل الآیات لقوم یشعرون
سبحان الّذی یهدی من یشآء الی صراطه قل انّی لصراط اللّه لمن فی السّموات و
الأرض طوبی لقوم یسرعون

\end{otherlanguage}

Subḥāna alladhī nazzala al-āyāti li-qawmin yafqahūn. Subḥāna alladhī
yunazzilu al-āyāti li-qawmin yashʿurūn. Subḥāna alladhī yahdī man
yashāʾu ilā ṣirāṭihi. Qul innī la-ṣirāṭu Allāhi li-man fī al-samāwāti
wa-l-arḍ. Ṭūbā li-qawmin yusriʿūn.

Glory be to Him who has sent down the verses for people who understand.
Glory be to Him who sends down the verses for people who perceive. Glory
be to Him who guides whomever He wills to His path. Say: I am indeed the
path of God for all who are in the heavens and on earth. Blessed are
those who hasten toward it.

\subsection{4}\label{section-1876}

\begin{otherlanguage}{arabic}

سبحان الّذی ینزّل الآیات لقوم یعلمون سبحان الّذی ینطق من جبروت الأمر لا
یعرفه الّا عباد مکرمون سبحان الّذی یحیی من یشآء بقوله کن فیکون سبحان الّذی
یرفع من یشآء الی سمآء الفضل و ینزّل منها ما اراد علی قدر مقدور

\end{otherlanguage}

Subḥāna alladhī yunazzilu al-āyāti li-qawmin yaʿlamūn. Subḥāna alladhī
yanṭiqu min jabarūti al-amr, lā yaʿrifuhu illā ʿibādun mukramūn. Subḥāna
alladhī yuḥyī man yashāʾu bi-qawlihi kun fa-yakūn. Subḥāna alladhī
yarfaʿu man yashāʾu ilā samāʾi al-faḍli wa-yunazzilu minhā mā arāda ʿalā
qadarin maqdūr.

Glory be to Him who sends down the verses for people who know. Glory be
to Him who speaks from the dominion of command; none recognize Him
except His honored servants. Glory be to Him who brings whomever He
wills to life through His word, ``Be,'' and it is. Glory be to Him who
raises whomever He wills to the heaven of grace and sends down from it
whatever He desires in a predetermined measure.

\subsection{5}\label{section-1877}

\begin{otherlanguage}{arabic}

تبارک الّذی یفعل ما یشآء بأمر من عنده انّه لهو الحقّ علّام العلوم تبارک الّذی
یلهم من یشآء ما اراد بأمره المبرم المکنون تبارک الّذی ینصر من یشآء بجنود
الغیب انّه لهو الفاعل لما اراد و انّه لهو العزیز القیّوم تبارک الّذی یعزّ من
یشآء بسلطان عزّه و یؤیّد من یشآء کیف اراد طوبی لقوم یعرفون

\end{otherlanguage}

Tabāraka alladhī yafʿalu mā yashāʾu bi-amrin min ʿindihi, innahu la-huwa
al-Ḥaqq ʿAllāmu al-ʿulūm. Tabāraka alladhī yulhimu man yashāʾu mā arāda
bi-amrihi al-mubrami al-maknūn. Tabāraka alladhī yanṣuru man yashāʾu
bi-junūdi al-ghayb, innahu la-huwa al-fāʿilu li-mā arāda wa-innahu
la-huwa al-ʿAzīzu al-Qayyūm. Tabāraka alladhī yuʿizzu man yashāʾu
bi-sulṭāni ʿizzihi wa-yuʾayyidu man yashāʾu kayfa arād. Ṭūbā li-qawmin
yaʿrifūn.

Blessed is He who does whatever He wills through a command from His
presence. He is indeed the Truth, the Knower of all knowledge. Blessed
is He who inspires whomever He wills with whatever He desires through
His irrevocable, hidden command. Blessed is He who aids whomever He
wills with the hosts of the unseen. He indeed accomplishes whatever He
desires, and He is the Almighty, the Self-Subsisting. Blessed is He who
exalts whomever He wills through the sovereignty of His might and
strengthens whomever He wills however He desires. Blessed are those who
recognize Him.

\subsection{6}\label{section-1878}

\begin{otherlanguage}{persian}

تبارک الّذی قدّر لکلّ شیء مقداراً فی لوح مخزون تبارک الّذی نزّل علی عبده ما
تستضیء به الأفئدة و العقول تبارک الّذی نزّل علی عبده من البلآء ما احترقت
به اکباد الّذین استقرّوا فی سرادق البقآء ثمّ قلوب المقرّبین تبارک الّذی نزّل
علی عبده من سحاب القضآء سهام البلآء اذاً یرانی فی صبر جمیل تبارک الّذی قدّر
لعبده ما لا قدّره لأحد من عباده انّه لهو الفرد العزیز القیّوم

\end{otherlanguage}

Tabāraka alladhī qaddara li-kulli shayʾin miqdāran fī lawḥin makhzūn.
Tabāraka alladhī nazzala ʿalā ʿabdihi mā tastaḍīʾu bihi al-afʾidatu
wa-l-ʿuqūl. Tabāraka alladhī nazzala ʿalā ʿabdihi mina al-balāʾi mā
iḥtaraqat bihi akbādu alladhīna istaqarrū fī surādiqi al-baqāʾi thumma
qulūbu al-muqarrabīn. Tabāraka alladhī nazzala ʿalā ʿabdihi min saḥābi
al-qaḍāʾi sihāma al-balāʾ, idhan yarānī fī ṣabrin jamīl. Tabāraka
alladhī qaddara li-ʿabdihi mā lā qaddarahu li-aḥadin min ʿibādihi,
innahu la-huwa al-Fardu al-ʿAzīzu al-Qayyūm.

Blessed is He who has ordained a measure for everything in a treasured
Tablet. Blessed is He who has sent down to His servant that which
illumines hearts and minds. Blessed is He who has sent down upon His
servant such affliction as has burned the innermost beings of those
established in the pavilion of eternity, and the hearts of those brought
near. Blessed is He who has sent down upon His servant the arrows of
affliction from the clouds of decree; then He sees Me enduring with
gracious patience. Blessed is He who has ordained for His servant what
He has ordained for none of His other servants. He is indeed the Unique,
the Almighty, the Self-Subsisting.

\subsection{7}\label{section-1879}

\begin{otherlanguage}{arabic}

تبارک الّذی نزّل علی عبده من غمام البغضآء من اولی الاغضآء رماح القضآء اذاً
یراه فی شکر عظیم تبارک الّذی نزّل علی عبده ثقل السّموات و الأرض انّا نحمده
فی ذلک و لا یعرفه الّا العارفون سبحان الّذی اوقع جماله تحت مخالیب الغلّ من
اولی الفحشآء انّا نرضی بذلک و لا یدرکه الّا المدرکون سبحان الّذی اودع
الحسین بین الأحزاب من الأعدآء و یرد فی کلّ حین علی جسده رماح القهر و
البغضآء انّا نشکره علی ما قضی علی عبده المنیب المغموم

\end{otherlanguage}

Tabāraka alladhī nazzala ʿalā ʿabdihi min ghamāmi al-baghḍāʾi min ulī
al-ighḍāʾi rimāḥa al-qaḍāʾ, idhan yarāhu fī shukrin ʿaẓīm. Tabāraka
alladhī nazzala ʿalā ʿabdihi thiqla al-samāwāti wa-l-arḍ, innā naḥmaduhu
fī dhālika wa-lā yaʿrifuhu illā al-ʿārifūn. Subḥāna alladhī awqaʿa
jamālahu taḥta makhālībi al-ghilli min ulī al-faḥshāʾ, innā narḍā
bi-dhālika wa-lā yudrikuhu illā al-mudrikūn. Subḥāna alladhī awdaʿa
al-Ḥusayna bayna al-aḥzābi mina al-aʿdāʾ, wa-yaridu fī kulli ḥīnin ʿalā
jasadihi rimāḥu al-qahri wa-l-baghḍāʾ, innā nashkuruhu ʿalā mā qaḍā ʿalā
ʿabdihi al-munībi al-maghmūm.

Blessed is He who has sent down upon His servant the spears of decree
from the clouds of hatred raised by those who turn away; then He sees
Him filled with profound gratitude. Blessed is He who has laid upon His
servant the weight of the heavens and the earth. We praise Him for this,
though none understand it except those endowed with knowledge. Glory be
to Him who has placed His Beauty beneath the claws of malice wielded by
the perpetrators of indecency. We are content with this, though none
comprehend it except those endowed with understanding. Glory be to Him
who has placed Husayn among the hosts of enemies, while at every moment
the spears of oppression and hatred strike His body. We thank Him for
what He has decreed for His sorrow-laden servant who turns to Him.

\subsection{8}\label{section-1880}

\begin{otherlanguage}{persian}

فلمّا رأیت نفسی علی قطب البلآء سمعت الصّوت الأبدع الأحلی من فوق رأسی فلمّا
توجّهت شاهدت حوریّة ذکر اسم ربّی معلّقةً فی الهوآء محاذی الرّأس و رأیت انّها
مستبشرة فی نفسها کأنّ طراز الرّضوان یظهر من وجهها و نضرة الرّحمن تعلن من
خدّها و کانت تنطق بین السّموات و الأرض بندآء تنجذب منه الأفئدة و العقول و
تبشّر کلّ الجوارح من ظاهری و باطنی ببشارة استبشرت بها نفسی و استفرحت منها
عباد مکرمون

\end{otherlanguage}

Fa-lammā raʾaytu nafsī ʿalā quṭbi al-balāʾi samiʿtu al-ṣawta al-abdaʿa
al-aḥlā min fawqi raʾsī. Fa-lammā tawajjahtu shāhadtu ḥūriyyata dhikri
ismi rabbī muʿallaqatan fī al-hawāʾi muḥādhiya al-raʾs, wa-raʾaytu
annahā mustabshiratun fī nafsihā ka-anna ṭirāza al-riḍwāni yaẓharu min
wajhihā wa-naḍrata al-Raḥmāni tuʿlanu min khaddihā. Wa-kānat tanṭiqu
bayna al-samāwāti wa-l-arḍi bi-nidāʾin tanjadhibu minhu al-afʾidatu
wa-l-ʿuqūl, wa-tubashshiru kulla al-jawāriḥi min ẓāhirī wa-bāṭinī
bi-bishāratin istabsharat bihā nafsī wa-stafraḥat minhā ʿibādun
mukramūn.

When I saw Myself at the very center of affliction, I heard the most
wondrous, sweetest voice above My head. When I turned, I saw the Maiden
of the remembrance of My Lord's name suspended in the air, level with My
head. I saw that she was filled with joy, as though the adornment of
divine pleasure shone from her face and the radiance of the All-Merciful
appeared upon her cheek. Between the heavens and the earth she spoke
with a call that enraptured hearts and minds, bringing to every part of
My outward and inward being glad tidings that rejoiced My soul and
filled the honored servants with delight.

\subsection{9}\label{section-1881}

\begin{otherlanguage}{persian}

و اشارت باصبعها الی رأسی و خاطبت من فی السّموات و الأرض تاللّه هذا لمحبوب
العالمین ولکن انتم لا تفقهون هذا لجمال اللّه بینکم و سلطانه فیکم ان انتم
تعرفون و هذا لسرّ اللّه و کنزه و امر اللّه و عزّه لمن فی ملکوت الأمر و الخلق
ان انتم تعقلون انّ هذا لهو الّذی یشتاق لقائه من فی جبروت البقآء ثمّ الّذینهم
استقرّوا خلف سرادق الأبهی ولکن انتم عن جماله معرضون

\end{otherlanguage}

Wa-ashārat bi-iṣbaʿihā ilā raʾsī wa-khāṭabat man fī al-samāwāti
wa-l-arḍ: Ta-Allāhi hādhā la-maḥbūbu al-ʿālamīn, wa-lākin antum lā
tafqahūn. Hādhā la-jamālu Allāhi baynakum wa-sulṭānuhu fīkum in antum
taʿrifūn. Wa-hādhā la-sirru Allāhi wa-kanzuhu wa-amru Allāhi wa-ʿizzuhu
li-man fī malakūti al-amri wa-l-khalqi in antum taʿqilūn. Inna hādhā
la-huwa alladhī yashtāqu liqāʾahu man fī jabarūti al-baqāʾi thumma
alladhīna hum istaqarrū khalfa surādiqi al-Abhā, wa-lākin antum ʿan
jamālihi muʿriḍūn.

She pointed her finger toward My head and addressed all who are in the
heavens and on earth: ``By God! This is the Beloved of the worlds, but
you do not understand. This is the Beauty of God among you and His
sovereignty in your midst, if only you would recognize Him. This is the
mystery of God and His treasure, the Cause of God and His might, for all
who are in the kingdom of command and creation, if only you would
understand. This is indeed the One whom those in the dominion of
eternity, and those established behind the pavilion of the All-Glorious,
long to meet, yet you turn away from His Beauty.''

\subsection{10}\label{section-1882}

\begin{otherlanguage}{arabic}

ان یا ملأ البیان انتم ان لن تنصروه سوف ینصره اللّه بجنود السّموات و الأرض
ثمّ جنود الغیب بأمره کن فیکون و یبعث بارادته خلقاً ما اطّلع بهم احد الّا
نفسه المهیمن القیّوم و یطهّرهم عن دنس الوهم و الهوی و یرفعهم الی مقام
التّقدیس و یظهر منهم آثار عزّ سلطانه فی الأرض کذلک قدّر من لدی اللّه العزیز
الودود

\end{otherlanguage}

In yā malaʾa al-Bayān, antum in lan tanṣurūhu sawfa yanṣuruhu Allāhu
bi-junūdi al-samāwāti wa-l-arḍi thumma junūdi al-ghaybi bi-amrihi kun
fa-yakūn. Wa-yabʿathu bi-irādatihi khalqan mā iṭṭalaʿa bihim aḥadun illā
nafsuhu al-Muhayminu al-Qayyūm. Wa-yuṭahhiruhum ʿan danasi al-wahmi
wa-l-hawā wa-yarfaʿuhum ilā maqāmi al-taqdīsi wa-yuẓhiru minhum āthāra
ʿizzi sulṭānihi fī al-arḍ. Kadhālika quddira min ladā Allāhi al-ʿAzīzi
al-Wadūd.

People of the Bayan! If you will not aid Him, God will aid Him with the
hosts of the heavens and the earth, and the hosts of the unseen, through
His command, ``Be,'' and it is. By His will He will raise up a people
known to none but Himself, the All-Protecting, the Self-Subsisting. He
will purify them of the defilement of delusion and desire, raise them to
the station of holiness, and manifest through them the signs of His
sovereign might on earth. Thus has it been decreed by God, the Almighty,
the All-Loving.

\subsection{11}\label{section-1883}

\begin{otherlanguage}{arabic}

ان یا ملأ البیان أ تکفرون بالّذی خلقتم للقائه ثمّ علی مقاعدکم تفرحون و
تعترضون علی الّذی شعرة منه خیر عند اللّه عمّن فی السّموات و الأرض ثمّ بنا
تستهزئون ان یا ملأ البیان فأتوا بما عندکم لأعرف بأیّ حجّة آمنتم بمظاهر
الأمر من قبل و الیوم بأیّ برهان تستکبرون

\end{otherlanguage}

In yā malaʾa al-Bayān, a-takfurūna bi-alladhī khuliqtum li-liqāʾihi
thumma ʿalā maqāʿidikum tafraḥūn, wa-taʿtariḍūna ʿalā alladhī shaʿratun
minhu khayrun ʿinda Allāhi ʿamman fī al-samāwāti wa-l-arḍi thumma binā
tastahziʾūn? In yā malaʾa al-Bayān, faʾtū bi-mā ʿindakum li-aʿrifa
bi-ayyi ḥujjatin āmantum bi-maẓāhiri al-amri min qablu wa-l-yawma
bi-ayyi burhānin tastakbirūn.

People of the Bayan! Do you deny the One you were created to meet, then
rejoice in your seats? Do you oppose the One a single hair of whom is
better in God's sight than all who are in the heavens and on earth, and
then mock Us? People of the Bayan! Bring forward what you possess, so
that I may know by what proof you believed in the Manifestations of the
Cause in the past, and by what evidence you now act so arrogantly.

\subsection{12}\label{section-1884}

\begin{otherlanguage}{arabic}

فوالّذی خلقنی من نور جماله ما وجدت غافلاً اغفل منکم و عمیاً اعمی عنکم انّکم
تستدلّون لایمانکم باللّه بما عندکم من الألواح لمّا نزّلت الآیات و اضآء
المصباح کفرتم بالّذی من قلمه قضت الأمور فی لوح محفوظ تقرؤون الآیات و
تکفرون بمطلعها و منزلها کذلک اخذ اللّه ابصارکم جزآء اعمالکم ان انتم
تشعرون و تکتبون الآیات فی العشیّ و الاشراق ثمّ عن منزلها انتم محتجبون

\end{otherlanguage}

Fa-wa-alladhī khalaqanī min nūri jamālihi, mā wajadtu ghāfilan aghfala
minkum wa-ʿamiyan aʿmā ʿankum. Innakum tastadillūna li-īmānikum
bi-Allāhi bi-mā ʿindakum mina al-alwāḥ. Lammā nuzzilati al-āyātu
wa-aḍāʾa al-miṣbāḥu kafartum bi-alladhī min qalamihi quḍiyati al-umūru
fī lawḥin maḥfūẓ. Taqraʾūna al-āyāti wa-takfurūna bi-maṭlaʿihā
wa-munzilihā. Kadhālika akhadha Allāhu abṣārakum jazāʾa aʿmālikum in
antum tashʿurūn. Wa-taktubūna al-āyāti fī al-ʿashiyyi wa-l-ishrāqi
thumma ʿan munzilihā antum muḥtajibūn.

By Him who created Me from the light of His beauty, I have found none
more heedless or more blind than you. You offer the Tablets in your
possession as proof of your faith in God. Yet when the verses were sent
down and the lamp shone forth, you denied the One through whose Pen all
affairs were decreed in a preserved Tablet. You read the verses and deny
their source and the One who reveals them. Thus has God taken away your
sight in recompense for your deeds, if only you perceived. You
transcribe the verses at evening and at daybreak, yet remain veiled from
their Revealer.

\subsection{13}\label{section-1885}

\begin{otherlanguage}{arabic}

اذاً یراکم الملأ الأعلی فی سوء اعمالکم و یتبرّؤنّ منکم و انتم لا تسمعون و
یستخبر بعضهم بعضاً ما یقولون هؤلآء الحمرات و فی ایّ واد هم یرتعون أ ینکرون
ما تشهد به ذواتهم أ یغمضون عیونهم و هم ینظرون تاللّه یا قوم بأفعالکم
تحیّرت سکّان مدائن الأسمآء و انتم فی وادی الجرز هائمون و لا تشعرون

\end{otherlanguage}

Idhan yarākumu al-malaʾu al-aʿlā fī sūʾi aʿmālikum wa-yatabarraʾūna
minkum wa-antum lā tasmaʿūn. Wa-yastakhbiru baʿḍuhum baʿḍan mā yaqūlu
hāʾulāʾi al-ḥumurātu wa-fī ayyi wādin hum yartaʿūn. A-yunkirūna mā
tashhadu bihi dhawātuhum? A-yughmiḍūna ʿuyūnahum wa-hum yanẓurūn?
Ta-Allāhi yā qawm, bi-afʿālikum taḥayyarat sukkānu madāʾini al-asmāʾi
wa-antum fī wādī al-juruzi hāʾimūna wa-lā tashʿurūn.

The Concourse on high sees you engaged in your evil deeds and disowns
you, though you do not hear. They ask one another, ``What are these
asses saying, and in what valley are they grazing? Do they deny what
their own beings testify to? Do they shut their eyes while they can
see?'' By God, people! Your actions have bewildered the inhabitants of
the cities of names, while you wander in a barren valley, unaware.

\subsection{14}\label{section-1886}

\begin{otherlanguage}{persian}

ان یا قلم الأعلی ان استمع ندآء ربّک من سدرة المنتهی فی البقعة الأحدیّة
النّورآء لتجد نفسک علی روح و ریحان من نغمات ربّک الرّحمن و تکون مقدّساً عن
الأحزان من هذه النّفحات الّتی تمرّ من شطر اسمی الغفور ثمّ ابتعث فی هذا
الهیکل هیاکل الأحدیّة لیحکینّ فی ملکوت الانشآء عن ربّهم العلیّ الأبهی و
یکوننّ من الّذینهم بأنوار ربّهم یستضیؤون

\end{otherlanguage}

In yā qalama al-Aʿlā, in istamiʿ nidāʾa rabbika min sidrati al-muntahā
fī al-buqʿati al-aḥadiyyati al-nawrāʾ, li-tajida nafsaka ʿalā rawḥin
wa-rayḥānin min naghamāti rabbika al-Raḥmān, wa-takūna muqaddasan ʿani
al-aḥzāni min hādhihi al-nafaḥāti allatī tamurru min shaṭri ismī
al-Ghafūr. Thumma ibʿath fī hādhā al-haykali hayākila al-aḥadiyyati
li-yaḥkiyanna fī malakūti al-inshāʾi ʿan rabbihimu al-ʿAliyyi al-Abhā
wa-yakūnanna mina alladhīna hum bi-anwāri rabbihim yastaḍīʾūn.

Pen of the Most High! Listen to the call of your Lord from the Tree of
the Utmost Boundary in the radiant realm of oneness, so that the
melodies of your Lord, the All-Merciful, may fill you with joy and
fragrance, and these breaths wafting from the direction of My name, the
Ever-Forgiving, may free you from sorrow. Then raise up within this
Temple the temples of oneness, so that in the kingdom of creation they
may tell of their Lord, the Most High, the All-Glorious, and be among
those illumined by the lights of their Lord.

\subsection{15}\label{section-1887}

\begin{otherlanguage}{arabic}

انّا قدّرنا هذا الهیکل مبدأ الوجود فی الخلق البدیع لیوقننّ الکلّ بأنّی اکون
مقتدراً علی ما اشآء بقولی کن فیکون و فی ظلّ کلّ حرف من حروفات هذا الهیکل
نبعث خلقاً لا یعلم عدّتهم الّا اللّه المهیمن القیّوم سوف یخلق اللّه منه خلقاً
لا تحجبهم اشارات الّذینهم بغوا علی اللّه و هم یشربون فی کلّ الأحیان کوثر
الحیوان الا انّهم هم الفائزون

\end{otherlanguage}

Innā qaddarnā hādhā al-haykala mabdaʾa al-wujūdi fī al-khalqi al-badīʿ,
li-yūqinanna al-kullu bi-annī akūnu muqtadiran ʿalā mā ashāʾu bi-qawlī
kun fa-yakūn. Wa-fī ẓilli kulli ḥarfin min ḥurūfāti hādhā al-haykali
nabʿathu khalqan lā yaʿlamu ʿiddatahum illā Allāhu al-Muhayminu
al-Qayyūm. Sawfa yakhluqu Allāhu minhu khalqan lā taḥjubuhum ishārātu
alladhīna hum baghaw ʿalā Allāhi wa-hum yashrabūna fī kulli al-aḥyāni
kawthara al-ḥayawān. Alā innahum humu al-fāʾizūn.

We have ordained this Temple to be the origin of existence in the
wondrous creation, so that all may know with certainty that I have the
power to accomplish whatever I will through My word, ``Be,'' and it is.
Beneath the shadow of every letter of this Temple We raise up a people
whose number none knows except God, the All-Protecting, the
Self-Subsisting. God will create from it a people whom the allusions of
those who have rebelled against God will not veil, and who will drink at
all times from the river of life. They, indeed, are the ones who attain
success.

\subsection{16}\label{section-1888}

\begin{otherlanguage}{persian}

اولئک العباد الّذین استقرّوا فی ظلّ رحمة ربّهم و ما منعهم المانعون یری من
وجوههم نضرة الرّحمن و یسمع من قلوبهم ذکر اسمی العزیز المکنون اولئک لو
تفتح شفواتهم فی تسبیح ربّهم یسبّحنّ معهم من فی السّموات و الأرض و قلیلاً من
النّاس ما هم یسمعون و اذا یذکرون بارئهم یذکرنّ معهم کلّ الأشیآء کذلک فضّلهم
اللّه علی الخلق ولکنّ النّاس لا یعلمون

\end{otherlanguage}

Ulāʾika al-ʿibādu alladhīna istaqarrū fī ẓilli raḥmati rabbihim wa-mā
manaʿahumu al-māniʿūn. Yurā min wujūhihim naḍratu al-Raḥmāni wa-yusmaʿu
min qulūbihim dhikru ismī al-ʿAzīzi al-Maknūn. Ulāʾika law tuftaḥu
shafawātuhum fī tasbīḥi rabbihim yusabbiḥanna maʿahum man fī al-samāwāti
wa-l-arḍi wa-qalīlan mina al-nāsi mā hum yasmaʿūn. Wa-idhā yadhkurūna
bāriʾahum yadhkuranna maʿahum kullu al-ashyāʾ. Kadhālika faḍḍalahumu
Allāhu ʿalā al-khalqi wa-lākinna al-nāsa lā yaʿlamūn.

These are the servants established beneath the shadow of their Lord's
mercy, whom none seeking to hinder them has held back. The radiance of
the All-Merciful is seen upon their faces, and the remembrance of My
name, the Almighty, the Hidden, is heard from their hearts. Were their
lips to open in glorification of their Lord, all who are in the heavens
and on earth would glorify Him with them, yet few among the people hear.
When they remember their Maker, all things remember Him with them. Thus
has God favored them above creation, but the people do not know.

\subsection{17}\label{section-1889}

\begin{otherlanguage}{persian}

و یتحرّکون حول امر اللّه کما یتحرّک الظّلّ حول الشّمس ان افتحوا الأبصار یا ملأ
البیان لعلّ انتم تشهدون و بحرکة هؤلآء یتحرّک کلّ شیء و بسکونهم یسکن کلّ شیء
ان انتم توقنون بهم یقبل الموحّدون الی قبلة الآفاق و ظهرت السّکینة و الوقار
بین الأخیار ان انتم تعلمون و بهم استقرّت الأرض و امطرت السّحاب و نزّلت
مائدة القدس من سمآء الفضل ان انتم تفقهون

\end{otherlanguage}

Wa-yataḥarrakūna ḥawla amri Allāhi kamā yataḥarraku al-ẓillu ḥawla
al-shams. In iftaḥū al-abṣāra yā malaʾa al-Bayān laʿalla antum
tashhadūn. Wa-bi-ḥarakati hāʾulāʾi yataḥarraku kullu shayʾin
wa-bi-sukūnihim yaskunu kullu shayʾin in antum tūqinūn. Bihim yuqbilu
al-muwaḥḥidūna ilā qiblati al-āfāqi wa-ẓaharati al-sakīnatu wa-l-waqāru
bayna al-akhyāri in antum taʿlamūn. Wa-bihim istaqarrati al-arḍu
wa-amṭarati al-saḥābu wa-nuzzilat māʾidatu al-qudsi min samāʾi al-faḍli
in antum tafqahūn.

They move around the Cause of God as the shadow moves around the sun.
Open your eyes, people of the Bayan, so that you may see. Through their
movement all things move, and through their stillness all things are
still, if only you were certain. Through them those who affirm God's
unity turn toward the focal point of all horizons, and tranquillity and
dignity have appeared among the righteous, if only you knew. Through
them the earth has been made firm, the clouds have poured down rain, and
the table of holiness has been sent down from the heaven of grace, if
only you understood.

\subsection{18}\label{section-1890}

\begin{otherlanguage}{persian}

اولئک حفظة امر اللّه فی الأرض یحفظون جمال الأمر من عجاج کلّ مشرک مبغوض و
لا یخافنّ من انفسهم فی سبیل اللّه بل ینفقونها رجآء للقآء المحبوب و
استعلائه بهذا الاسم المقتدر القادر العزیز القدّوس

\end{otherlanguage}

Ulāʾika ḥafaẓatu amri Allāhi fī al-arḍ, yaḥfaẓūna jamāla al-amri min
ʿajāji kulli mushrikin mabghūḍ. Wa-lā yakhāfanna min anfusihim fī sabīli
Allāhi bal yunfiqūnahā rajāʾan li-liqāʾi al-Maḥbūbi wa-stiʿlāʾihi
bi-hādhā al-ismi al-Muqtadiri al-Qādiri al-ʿAzīzi al-Quddūs.

These are the guardians of the Cause of God on earth. They protect the
beauty of the Cause from the dust raised by every detested idolater.
They have no fear for their lives in the path of God, but expend them in
the hope of meeting the Beloved and of His exaltation through this name,
the All-Powerful, the Omnipotent, the Almighty, the Most Holy.

\subsection{19}\label{section-1891}

\begin{otherlanguage}{persian}

ان یا هذا الهیکل قم بنفسک علی شأن یقومنّ بقیامک کلّ الممکنات ثمّ انصر ربّک
بما اعطیناک من القدرة و الاقتدار ایّاک ان تجزع حین الّذی یجزع فیه کلّ
الأشیآء کن مظهر اسمی القیّوم ثمّ انصر ربّک بما استطعت و لا تنظر الکائنات و
ما یخرج من افواههم الّا کندآء بعوضة فی واد ما حدّد بالحدود ان اشرب کوثر
الحیوان باسمی الرّحمن ثمّ اسق المقرّبین من اهل هذا الرّضوان ما ینقطعون به عن
کلّ الأسمآء و یدخلهم فی هذا الظّلّ المبارک الممدود

\end{otherlanguage}

In yā hādhā al-haykal, qum bi-nafsika ʿalā shaʾnin yaqūmanna bi-qiyāmika
kullu al-mumkināt. Thumma unṣur rabbaka bi-mā aʿṭaynāka mina al-qudrati
wa-l-iqtidār. Iyyāka an tajzaʿa ḥīna alladhī yajzaʿu fīhi kullu
al-ashyāʾ. Kun maẓhara ismī al-Qayyūm, thumma unṣur rabbaka bi-mā
istaṭaʿta wa-lā tanẓuri al-kāʾināti wa-mā yakhruju min afwāhihim illā
ka-nidāʾi baʿūḍatin fī wādin mā ḥuddida bi-l-ḥudūd. In ishrab kawthara
al-ḥayawāni bi-ismī al-Raḥmān, thumma isqi al-muqarrabīna min ahli hādhā
al-riḍwāni mā yanqaṭiʿūna bihi ʿan kulli al-asmāʾi wa-yudkhiluhum fī
hādhā al-ẓilli al-mubāraki al-mamdūd.

Temple! Arise in your own person in such a way that all contingent
beings may arise through your arising. Then aid your Lord with the power
and might We have bestowed upon you. Beware of dismay when all things
are dismayed. Be the manifestation of My name, the Self-Subsisting; then
aid your Lord to the utmost of your ability, and regard created beings
and whatever comes from their mouths as no more than the hum of a gnat
in a boundless valley. Drink from the river of life in My name, the
All-Merciful, then give those brought near among the people of this
paradise a draught that will detach them from every name and bring them
beneath this blessed, outstretched shadow.

\subsection{20}\label{section-1892}

\begin{otherlanguage}{arabic}

ان یا هذا الهیکل انّا حشرنا فیک کلّ الأشیآء عمّا خلق بین الأرض و السّمآء و
سألناهم ما اخذنا به عنهم العهد فی ذرّ البقآء اذاً وجدنا اکثرهم کلیل اللّسان
شاخصة الأبصار و قلیلاً ناضر الوجه طلق البیان و بعثنا من هؤلآء خلق ما کان
و ما یکون اولئک کرّم اللّه وجوههم عن التّوجّه الی وجوه المشرکین و اسکنهم فی
ظلّ سدرة نفسه و انزل علیهم سکینة الأمر و ایّدهم بجنود الغیب و الشّهود

\end{otherlanguage}

In yā hādhā al-haykal, innā ḥasharnā fīka kulla al-ashyāʾ ʿammā khuliqa
bayna al-arḍ wa-al-samāʾ, wa-saʾalnāhum mā akhadhnā bihi ʿanhum al-ʿahd
fī dharr al-baqāʾ. Idhan wajadnā aktharahum kalīla al-lisān shākhiṣata
al-abṣār, wa-qalīlan nāḍira al-wajh ṭalqa al-bayān. Wa-baʿathnā min
hāʾulāʾ khalqa mā kāna wa-mā yakūn. Ulāʾika karrama Allāh wujūhahum ʿan
al-tawajjuh ilā wujūh al-mushrikīn, wa-askanahum fī ẓill sidrat nafsihi,
wa-anzala ʿalayhim sakīnat al-amr, wa-ayyadahum bi-junūd al-ghayb
wa-al-shuhūd.

O Temple! We have gathered within you all things created between earth
and heaven, and questioned them about that concerning which We took
their covenant in the primordial realm of eternity. We found most of
them faltering in speech, their eyes staring fixedly, and only a few
radiant in countenance and fluent in utterance. From these We brought
forth the creation of all that has been and will be. God has honored
their faces above turning toward the faces of those who associate
partners with Him, sheltered them beneath the shade of the Lote-Tree of
His own Self, sent down upon them the tranquility of His Cause, and
strengthened them with the hosts of the unseen and the visible.

\subsection{21}\label{section-1893}

\begin{otherlanguage}{arabic}

ان یا عین هذا الهیکل لا تلتفتی الی السّمآء و ما فیها و لا الی الأرض و من
علیها انّا خلقناک لجمالی ها هو هذا فانظری کیف شئت و لا تمنعی لحاظک عن
جمال ربّک العزیز المحبوب سوف نبعث بک اعیناً حدیدة و ابصاراً ناظرة یرون آیات
بارئهم و یحوّلنّ النّظر عن کلّ ما یدرکه المدرکون و بک نعطی قوّة البصر لمن
نشآء و نأخذ الّذین منعوا عن هذا الفضل الا انّهم من کأس الوهم یکرعون و لا
یفقهون

\end{otherlanguage}

In yā ʿayna hādhā al-haykal, lā taltafitī ilā al-samāʾ wa-mā fīhā wa-lā
ilā al-arḍ wa-man ʿalayhā. Innā khalaqnāki li-jamālī. Hā huwa hādhā,
fa-unẓurī kayfa shiʾti wa-lā tamnaʿī liḥāẓaki ʿan jamāl rabbiki al-ʿazīz
al-maḥbūb. Sawfa nabʿathu biki aʿyunan ḥadīdatan wa-abṣāran nāẓiratan
yarawna āyāt bāriʾihim wa-yuḥawwilunna al-naẓar ʿan kull mā yudrikuhu
al-mudrikūn. Wa-biki nuʿṭī quwwat al-baṣar li-man nashāʾ wa-naʾkhudhu
alladhīna muniʿū ʿan hādhā al-faḍl. Alā innahum min kaʾs al-wahm
yakraʿūn wa-lā yafqahūn.

O Eye of this Temple! Do not turn toward heaven and what it contains, or
toward earth and those upon it. We created you for My beauty. Here it is
before you: gaze as you wish, and do not withhold your glances from the
beauty of your Lord, the Almighty, the Beloved. Through you We will
raise up keen eyes and discerning vision that will behold the signs of
their Creator and turn their gaze away from all that those endowed with
perception can perceive. Through you We bestow the power of sight upon
whomever We will, and seize those who have been deprived of this bounty.
Indeed, they drink deeply from the cup of delusion and do not
understand.

\subsection{22}\label{section-1894}

\begin{otherlanguage}{persian}

ان یا سمع هذا الهیکل طهّر نفسک عن نعیق کلّ ناعق مردود ثمّ استمع نغمات ربّک
انّه یوحی الیک من جهة العرش انّه لا اله الّا انا العزیز المقتدر المهیمن
القیّوم سوف نبعث بک آذاناً مطهّرة لاصغآء کلمة اللّه و ما ظهر من مطلع بیان
ربّک الرّحمن الا انّهنّ یجدن ترنّمات الوحی من هذا الشّطر المبارک المحمود

\end{otherlanguage}

In yā samʿa hādhā al-haykal, ṭahhir nafsaka ʿan naʿīq kull nāʿiq mardūd,
thumma istamiʿ naghamāt rabbika. Innahu yūḥī ilayka min jihat al-ʿarsh:
innahu lā ilāha illā anā al-ʿazīz al-muqtadir al-muhaymin al-qayyūm.
Sawfa nabʿathu bika ādhānan muṭahharatan li-iṣghāʾ kalimat Allāh wa-mā
ẓahara min maṭlaʿ bayān rabbika al-raḥmān. Alā innahunna yajidna
tarannumāt al-waḥy min hādhā al-shaṭr al-mubārak al-maḥmūd.

O Hearing of this Temple! Purify yourself from the clamor of every
rejected clamorer, then listen to the melodies of your Lord. From the
direction of the throne He reveals to you: ``There is no God but Me, the
Almighty, the All-Powerful, the All-Protecting, the Self-Subsisting.''
Through you We will raise up purified ears to listen to the Word of God
and to what has appeared from the dawning-place of the utterance of your
Lord, the All-Merciful. Indeed, they will hear the melodies of
revelation from this blessed and praiseworthy quarter.

\subsection{23}\label{section-1895}

\begin{otherlanguage}{arabic}

ان یا لسان هذا الهیکل انّا خلقناک باسمی الرّحمن و علّمناک ما کنز فی البیان
و انطقناک لذکری العظیم فی الأکوان ان انطق بهذا الذّکر البدیع و لا تخف من
مظاهر الشّیطان لأنّک خلقت لذلک بأمری المهیمن القیّوم و بک فتحنا اللّسان
بالبیان فیما کان و نفتح بسلطانی فیما یکون و بک نبعث السناً ناطقة کلّها
تتحرّک بالثّنآء فی ملإ البقآء و بین ملإ الانشآء کذلک نزّلت الآیات و قضی
الأمر من لدن مالک الأسمآء و الصّفات انّ ربّک لهو الحقّ علّام الغیوب اولئک لا
یمنعهم شیء عن ثنآء بارئهم بهم یقومنّ الأشیآء علی ذکر مالک الأسمآء بأنّه لا
اله الّا انا المقتدر العزیز المحبوب لا تنطق السن الذّاکرین الّا و یمدّها هذا
اللّسان من هذا الرّضوان و قلیلاً من النّاس ما هم یعرفون ان من لسان الّا و قد
یسبّح ربّه و ینطق علی ذکره و من النّاس من یفقه و یذکر و منهم من یذکرون و لا
یفقهون

\end{otherlanguage}

In yā lisāna hādhā al-haykal, innā khalaqnāka bi-ismī al-raḥmān,
wa-ʿallamnāka mā kuniza fī al-bayān, wa-anṭaqnāka li-dhikrī al-ʿaẓīm fī
al-akwān. An inṭiq bi-hādhā al-dhikr al-badīʿ wa-lā takhaf min maẓāhir
al-shayṭān, li-annaka khuliqta li-dhālika bi-amrī al-muhaymin al-qayyūm.
Wa-bika fataḥnā al-lisān bi-al-bayān fīmā kāna, wa-naftaḥu bi-sulṭānī
fīmā yakūn. Wa-bika nabʿathu alsunan nāṭiqatan kulluhā tataḥarraku
bi-al-thanāʾ fī malaʾ al-baqāʾ wa-bayna malaʾ al-inshāʾ. Kadhālika
nuzzilat al-āyāt wa-quḍiya al-amr min ladun mālik al-asmāʾ wa-al-ṣifāt.
Inna rabbaka la-huwa al-ḥaqq ʿallām al-ghuyūb. Ulāʾika lā yamnaʿuhum
shayʾ ʿan thanāʾ bāriʾihim. Bihim yaqūmanna al-ashyāʾ ʿalā dhikr mālik
al-asmāʾ bi-annahu lā ilāha illā anā al-muqtadir al-ʿazīz al-maḥbūb. Lā
tanṭiqu alsun al-dhākirīn illā wa-yumidduhā hādhā al-lisān min hādhā
al-riḍwān, wa-qalīlan min al-nās mā hum yaʿrifūn. In min lisān illā
wa-qad yusabbiḥu rabbahu wa-yanṭiqu ʿalā dhikrihi. Wa-min al-nās man
yafqahu wa-yadhkuru, wa-minhum man yadhkurūna wa-lā yafqahūn.

O Tongue of this Temple! We created you through My name, the
All-Merciful, taught you what was treasured in utterance, and caused you
to speak so that I might be greatly remembered throughout creation.
Speak this wondrous remembrance and do not fear the manifestations of
Satan, for you were created for this by My all-protecting,
self-subsisting command. Through you We have opened the tongue to
utterance in all that has been, and through My sovereignty We will do so
in all that is to come. Through you We will raise up speaking tongues,
all moving in praise among the company of eternity and the company of
creation. Thus have the verses been sent down and the decree issued by
the Possessor of names and attributes. Your Lord is indeed the Truth,
the Knower of all unseen things. Nothing will prevent them from praising
their Creator. Through them all things will arise to proclaim the
Possessor of names: ``There is no God but Me, the All-Powerful, the
Almighty, the Beloved.'' No tongue of those who remember speaks without
being sustained by this Tongue from this paradise, yet few among the
people recognize this. There is no tongue that does not glorify its Lord
and speak His remembrance. Some people understand and remember, while
others remember but do not understand.

\subsection{24}\label{section-1896}

\begin{otherlanguage}{persian}

ان یا حوریّة المعانی ان اخرجی من غرفات الکلمات باذن اللّه مالک الأرضین و
السّموات ثمّ اظهری بطراز اللّاهوت ثمّ اسقی خمر الجبروت بأنامل الیاقوت لعلّ
اهل النّاسوت یطّلعنّ بما اشرقت من افق الملکوت شمس البقآء بطراز البهآء و
یقومنّ علی الثّنآء بین الأرض و السّمآء فی ذکر هذا الفتی الّذی استقرّ علی عرش
اسمه المنّان فی قطب الجنان و من وجهه ظهرت نضرة الرّحمن و عن لحظه لحظات
السّبحان و من شؤونه شؤونات اللّه المهیمن القیّوم

\end{otherlanguage}

In yā ḥūriyyat al-maʿānī, an ukhrujī min ghurufāt al-kalimāt bi-idhn
Allāh mālik al-arḍīn wa-al-samāwāt, thumma iẓharī bi-ṭirāz al-lāhūt,
thumma isqī khamr al-jabarūt bi-anāmil al-yāqūt, laʿalla ahl al-nāsūt
yaṭṭaliʿunna bi-mā ashraqat min ufuq al-malakūt shams al-baqāʾ bi-ṭirāz
al-bahāʾ, wa-yaqūmunna ʿalā al-thanāʾ bayna al-arḍ wa-al-samāʾ fī dhikr
hādhā al-fatā alladhī istaqarra ʿalā ʿarsh ismihi al-mannān fī quṭb
al-jinān, wa-min wajhihi ẓaharat naḍrat al-raḥmān, wa-ʿan laḥẓihi
laḥaẓāt al-subḥān, wa-min shuʾūnihi shuʾūnāt Allāh al-muhaymin
al-qayyūm.

O Maiden of inner meanings! Come forth from the chambers of words by the
permission of God, the Possessor of the earths and the heavens. Appear
adorned with the vesture of divinity, then with ruby fingers serve the
wine of the realm of power, so that the inhabitants of the mortal realm
may become aware of the splendor with which the Sun of eternity has
shone from the horizon of the heavenly kingdom, and arise in praise
between earth and heaven, remembering this Youth who has taken His seat
upon the throne of His name, the Bestower, at the heart of the gardens
of paradise. From His face has appeared the radiance of the
All-Merciful, from His glance the glances of the All-Glorious, and from
His ways the ways of God, the All-Protecting, the Self-Subsisting.

\subsection{25}\label{section-1897}

\begin{otherlanguage}{persian}

و ان لن تجدی احداً ان یأخذ من الید البیضآء الخمر الحمرآء باسم ربّک العلیّ
الأعلی الّذی ظهر مرّة بعد اولی باسمه الأبهی لا تحزنی دعی هؤلآء بأنفسهم ثمّ
ارجعی الی خلف سرادق العظمة و الکبریآء اذاً تجدی قوماً تستضیء انوار وجوههم
کالشّمس فی وسط الزّوال و هم یهلّلون و یسبّحون ربّهم بهذا الاسم الّذی قام علی
مقرّ الاستقلال بسلطان العزّ و الاجلال و انّک لن تسمعی منهم الّا ذکری انّ ربّک
شهید علی ما اقول و ما اطّلع بهؤلآء احد من الّذینهم خلقوا بکلمة اللّه فی ازل
الآزال کذلک فصّلنا لک الأمر و صرّفنا الآیات لعلّ النّاس فی آثار ربّهم یتفکّرون

\end{otherlanguage}

Wa-in lan tajidī aḥadan an yaʾkhudha min al-yad al-bayḍāʾ al-khamr
al-ḥamrāʾ bi-ism rabbiki al-ʿaliyy al-aʿlā alladhī ẓahara marratan baʿda
ūlā bi-ismihi al-abhā, lā taḥzanī. Daʿī hāʾulāʾ bi-anfusihim thumma
irjiʿī ilā khalf surādiq al-ʿaẓama wa-al-kibriyāʾ. Idhan tajidī qawman
tastaḍīʾu anwār wujūhihim ka-al-shams fī wasaṭ al-zawāl, wa-hum
yuhallilūna wa-yusabbiḥūna rabbahum bi-hādhā al-ism alladhī qāma ʿalā
maqarr al-istiqlāl bi-sulṭān al-ʿizz wa-al-ijlāl. Wa-innaki lan tasmaʿī
minhum illā dhikrī. Inna rabbaki shahīd ʿalā mā aqūl. Wa-mā iṭṭalaʿa
bi-hāʾulāʾ aḥad min alladhīna hum khuliqū bi-kalimat Allāh fī azal
al-āzāl. Kadhālika faṣṣalnā laki al-amr wa-ṣarrafnā al-āyāt laʿalla
al-nās fī āthār rabbihim yatafakkarūn.

If you find no one to take the crimson wine from the snow-white hand in
the name of your Lord, the Exalted, the Most High, who has appeared once
again in His name, the All-Glorious, do not grieve. Leave these people
to themselves, then return behind the pavilion of grandeur and majesty.
There you will find a people whose faces shine with light like the sun
at its zenith. They proclaim the oneness of their Lord and glorify Him
through this Name, which has taken its place upon the seat of
independence with the sovereignty of might and glory. You will hear from
them nothing but My remembrance. Your Lord bears witness to what I say.
None of those created by the Word of God from the eternity of eternities
has known these people. Thus have We explained the matter to you and
variously set forth the verses, so that people may reflect upon the
signs of their Lord.

\subsection{26}\label{section-1898}

\begin{otherlanguage}{arabic}

انّهم ما امروا بسجدة الآدم و ما حوّلوا وجوههم عن وجه ربّک و هم من نعمة
التّقدیس فی کلّ حین یتنعّمون کذلک رقم قلم الرّحمن اسرار ما کان و ما یکون لعلّ
النّاس هم یعرفون فسوف یظهر اللّه هؤلآء فی الأرض و یرفع بهم ذکره و ینشر
آثاره و یثبت کلماته و یعلن آیاته رغماً للّذینهم کفروا و انکروا و کانوا
بآیاته یجحدون

\end{otherlanguage}

Innahum mā umirū bi-sajdat al-Ādam wa-mā ḥawwalū wujūhahum ʿan wajh
rabbika, wa-hum min niʿmat al-taqdīs fī kull ḥīn yatanaʿʿamūn. Kadhālika
raqama qalam al-raḥmān asrār mā kāna wa-mā yakūn laʿalla al-nās hum
yaʿrifūn. Fa-sawfa yuẓhiru Allāh hāʾulāʾ fī al-arḍ wa-yarfaʿu bihim
dhikrahu wa-yanshuru āthārahu wa-yuthbitu kalimātihi wa-yuʿlinu āyātihi
raghman li-alladhīna hum kafarū wa-ankarū wa-kānū bi-āyātihi yajḥadūn.

They were not commanded to bow down before Adam, nor have they turned
their faces away from the face of your Lord. At every moment they enjoy
the bounty of holiness. Thus has the Pen of the All-Merciful inscribed
the mysteries of all that has been and will be, so that people may
understand. God will bring these people forth upon the earth, and
through them exalt His remembrance, spread His works, establish His
words, and proclaim His signs, despite those who have disbelieved,
denied, and rejected His signs.

\subsection{27}\label{section-1899}

\begin{otherlanguage}{arabic}

ان یا طلعة الأحدیّة ان وجدتهم و ادرکت لقائهم ان اقصصی لهم ما یقصّ لک
الغلام من قصص نفسه و بما ورد علیه لیطّلعنّ علی ما هو المسطور فی لوح محفوظ
و اخبریهم من نبإ الغلام و ما مسّته من البأسآء و الضّرّآء لیتذکّرنّ بمصائبی و
یکوننّ من الّذینهم متذکّرون ثمّ اذکری لهم بأنّا اصطفینا من اخواننا احداً و
رشّحنا علیه من طمطام بحر العلم رشحاً ثمّ البسناه قمیص اسم من الأسمآء و
رفعناه الی المقام الّذی قام الکلّ علی ثنآء نفسه و حفظناه عن ضرّ کلّ ذی ضرّ
علی شأن یعجز عنه القادرون

\end{otherlanguage}

In yā ṭalʿat al-aḥadiyya, in wajadtihim wa-adrakti liqāʾahum, an uqṣuṣī
lahum mā yaquṣṣu laki al-ghulām min qiṣaṣ nafsihi wa-bi-mā warada
ʿalayhi li-yaṭṭaliʿunna ʿalā mā huwa al-masṭūr fī lawḥ maḥfūẓ.
Wa-akhbirīhim min nabaʾ al-ghulām wa-mā massat-hu min al-baʾsāʾ
wa-al-ḍarrāʾ li-yatadhakkarunna bi-maṣāʾibī wa-yakūnunna min alladhīna
hum mutadhakkirūn. Thumma udhkurī lahum bi-annā iṣṭafaynā min ikhwāninā
aḥadan wa-rashshaḥnā ʿalayhi min ṭamṭām baḥr al-ʿilm rashḥan, thumma
albasnāhu qamīṣ ism min al-asmāʾ wa-rafaʿnāhu ilā al-maqām alladhī qāma
al-kull ʿalā thanāʾ nafsihi, wa-ḥafiẓnāhu ʿan ḍarr kull dhī ḍarr ʿalā
shaʾn yaʿjizu ʿanhu al-qādirūn.

O Countenance of oneness! If you find them and attain their presence,
recount to them what the Youth tells you of His own story and of what
has befallen Him, so that they may learn what is inscribed in a
preserved Tablet. Tell them the tidings of the Youth and the adversity
and affliction that have touched Him, so that they may take heed through
My sufferings and be among those who are mindful. Then tell them that We
chose one of Our brothers and sprinkled upon him a spray from the
surging ocean of knowledge. We then clothed him in the garment of one of
the names and raised him to a station at which all arose to praise him.
We protected him from the harm of everyone intent on harm, in a manner
beyond the ability of the powerful.

\subsection{28}\label{section-1900}

\begin{otherlanguage}{arabic}

و کنّا وحدة فی مقابلة اهل السّموات و الأرض فی ایّام کلّ العباد قاموا علی
قتلی و کنّا بینهم ناطقاً بذکر اللّه و ثنائه و قائماً علی امره الی ان حقّقت
کلمة اللّه بین خلقه و اشتهرت آثاره و علت قدرته و لاحت سلطنته و یشهد بذلک
عباد مکرمون انّ اخی لمّا رأی الأمر ارتفع وجد فی نفسه کبراً و غرورا اذاً خرج
عن خلف الأستار و حارب بنفسی و جادل بآیاتی و کذّب برهانی و جحد آثاری و ما
شبع بطن الحریص الی ان اراد اکل لحمی و شرب دمی و یشهد بذلک العباد الّذینهم
هاجروا مع اللّه و عن ورائهم عباد مقرّبون

\end{otherlanguage}

Wa-kunnā waḥdatan fī muqābalat ahl al-samāwāt wa-al-arḍ fī ayyām kull
al-ʿibād qāmū ʿalā qatlī, wa-kunnā baynahum nāṭiqan bi-dhikr Allāh
wa-thanāʾihi wa-qāʾiman ʿalā amrihi ilā an ḥuqqiqat kalimat Allāh bayna
khalqihi wa-ishtaharat āthāruhu wa-ʿalat qudratuhu wa-lāḥat salṭanatuhu,
wa-yashhadu bi-dhālika ʿibād mukramūn. Inna akhī lammā raʾā al-amr
irtafaʿa wajada fī nafsihi kibran wa-ghurūran. Idhan kharaja ʿan khalf
al-astār wa-ḥāraba bi-nafsī wa-jādala bi-āyātī wa-kadhdhaba burhānī
wa-jaḥada āthārī. Wa-mā shabiʿa baṭn al-ḥarīṣ ilā an arāda akla laḥmī
wa-shurba damī. Wa-yashhadu bi-dhālika al-ʿibād alladhīna hum hājarū
maʿa Allāh, wa-ʿan warāʾihim ʿibād muqarrabūn.

We stood alone against the inhabitants of the heavens and the earth in
the days when all people arose to kill Me. Among them We continued to
proclaim the remembrance and praise of God and to uphold His Cause,
until the Word of God was established among His creatures, His signs
became widely known, His power was exalted, and His sovereignty shone
forth. Honored servants bear witness to this. When My brother saw that
the Cause had risen, pride and conceit arose within him. He then emerged
from behind the veils, made war against Me, disputed My verses, denied
My proof, and rejected My signs. The belly of the covetous one was not
satisfied until he desired to eat My flesh and drink My blood. Those
servants who migrated with God bear witness to this, and beyond them,
servants who have been brought near.

\subsection{29}\label{section-1901}

\begin{otherlanguage}{arabic}

و شاور فی ذلک مع احد من خدّامی و اغواه علی ذلک اذاً نصرنی اللّه بجنود الغیب
و الشّهادة و حفظنی بالحقّ و انزل علیّ ما منعه عمّا اراد و بطل مکر الّذینهم
کفروا بآیات الرّحمن الا انّهم قوم منکرون فلمّا شیّع ما سوّلت له نفسه و اطّلع
به الّذینهم هاجروا ارتفع الضّجیج من هؤلآء و بلغ الی مقام کاد ان یشتهر بین
المدینة اذاً منعناهم و القینا علیهم کلمة الصّبر لیکوننّ من الّذینهم یصبرون

\end{otherlanguage}

Wa-shāwara fī dhālika maʿa aḥad min khuddāmī wa-aghwāhu ʿalā dhālika.
Idhan naṣaranī Allāh bi-junūd al-ghayb wa-al-shahāda wa-ḥafiẓanī
bi-al-ḥaqq wa-anzala ʿalayya mā manaʿahu ʿammā arāda, wa-baṭala makr
alladhīna hum kafarū bi-āyāt al-raḥmān. Alā innahum qawm munkirūn.
Fa-lammā shuyyiʿa mā sawwalat lahu nafsuhu wa-iṭṭalaʿa bihi alladhīna
hum hājarū, irtafaʿa al-ḍajīj min hāʾulāʾ wa-balagha ilā maqām kāda an
yashtahira bayna al-madīna. Idhan manaʿnāhum wa-alqaynā ʿalayhim kalimat
al-ṣabr li-yakūnunna min alladhīna hum yaṣbirūn.

He consulted about this with one of My servants and enticed him into it.
Then God aided Me with the hosts of the unseen and the visible,
protected Me through the truth, and sent down upon Me what prevented him
from carrying out his intent. The scheming of those who disbelieved in
the signs of the All-Merciful came to nothing. Indeed, they are a people
given to denial. When news of what his self had prompted him to do was
spread, and those who had migrated learned of it, an outcry arose among
them, reaching such a pitch that it almost became known throughout the
city. We then restrained them and imparted to them the word of patience,
so that they might be among those who endure patiently.

\subsection{30}\label{section-1902}

\begin{otherlanguage}{arabic}

فواللّه الّذی لا اله الّا هو انّا صبرنا فی ذلک و امرنا العباد بالصّبر و
الاصطبار و خرجنا من بین هؤلآء و سکنّا فی بیت آخر لتسکن نار البغضآء فی
صدره و یکون من الّذینهم مهتدون و ما تعرّضنا به و ما رأیناه من بعد و جلسنا
فی البیت وحدةً مرتقباً فضل اللّه المهیمن القیّوم انّه لمّا اطّلع بأنّ الأمر
اشتهر اخذ قلم الکذب و کتب الی العباد و نسب کلّ ما فعل بجمالی الفرید
المظلوم ابتغآء فتنة فی نفسه و ادخال البغضآء فی صدور الّذینهم آمنوا باللّه
العزیز الودود

\end{otherlanguage}

Fa-wa-Allāh alladhī lā ilāha illā huwa, innā ṣabarnā fī dhālika
wa-amarnā al-ʿibād bi-al-ṣabr wa-al-iṣṭibār, wa-kharajnā min bayn
hāʾulāʾ wa-sakannā fī bayt ākhar li-taskuna nār al-baghḍāʾ fī ṣadrihi
wa-yakūna min alladhīna hum muhtadūn. Wa-mā taʿarraḍnā bihi wa-mā
raʾaynāhu min baʿd, wa-jalasnā fī al-bayt waḥdatan murtaqiban faḍl Allāh
al-muhaymin al-qayyūm. Innahu lammā iṭṭalaʿa bi-anna al-amr ishtahara,
akhadha qalam al-kadhib wa-kataba ilā al-ʿibād wa-nasaba kulla mā faʿala
bi-jamālī al-farīd al-maẓlūm, ibtighāʾ fitna fī nafsihi wa-idkhāl
al-baghḍāʾ fī ṣudūr alladhīna hum āmanū bi-Allāh al-ʿazīz al-wadūd.

By God, besides whom there is no other God! We endured this patiently
and commanded the servants to show patience and forbearance. We departed
from among them and took up residence in another house, so that the fire
of hatred in his breast might subside and he might be among the rightly
guided. We did not interfere with him or see him afterward, but remained
alone in the house, awaiting the grace of God, the All-Protecting, the
Self-Subsisting. When he learned that the matter had become known, he
took up the pen of falsehood, wrote to the servants, and attributed
everything he had done to My peerless and wronged Beauty, seeking within
himself to stir up discord and to instill hatred in the hearts of those
who believed in God, the Almighty, the All-Loving.

\subsection{31}\label{section-1903}

\begin{otherlanguage}{arabic}

فوالّذی نفسی بیده تحیّرنا من مکره بل تحیّر منه کلّ الوجود من الغیب و الشّهود
مع ذلک ما سکن فی نفسه الی ان ارتکب ما لا یجری القلم علیه و به ضیّع حرمتی
و حرمة اللّه المقتدر العزیز المحمود لو اذکر ما فعل بی لن تتمّه بحور الأرض
لو یجعلها اللّه مداداً و لن تنفده الأشیآء ولو یقلّبها اللّه اقلاماً کذلک نلقی
ما ورد علی نفسی ان انتم تعلمون

\end{otherlanguage}

Fa-wa-alladhī nafsī bi-yadihi, taḥayyarnā min makrihi, bal taḥayyara
minhu kull al-wujūd min al-ghayb wa-al-shuhūd. Maʿa dhālika mā sakana fī
nafsihi ilā an irtakaba mā lā yajrī al-qalam ʿalayhi, wa-bihi ḍayyaʿa
ḥurmatī wa-ḥurmat Allāh al-muqtadir al-ʿazīz al-maḥmūd. Law adhkuru mā
faʿala bī lan tutimmahu buḥūr al-arḍ law yajʿaluhā Allāh midādan, wa-lan
tunfidahu al-ashyāʾ wa-law yuqallibuhā Allāh aqlāman. Kadhālika nulqī mā
warada ʿalā nafsī, in antum taʿlamūn.

By Him in whose hand is My soul! We were bewildered by his scheming;
indeed, all existence, unseen and visible, was bewildered by it. Even
so, he found no rest within himself until he committed what the pen
cannot recount, violating thereby My honor and the honor of God, the
All-Powerful, the Almighty, the Praiseworthy. Were I to recount what he
did to Me, the oceans of the earth could not complete the account, even
if God made them ink, nor could all things exhaust it, even if God
turned them into pens. Thus do We relate what has befallen Me, if only
you knew.

\subsection{32}\label{section-1904}

\begin{otherlanguage}{arabic}

ان یا قلم البقآء لا تحزن عمّا ورد علیک فسوف یبعث اللّه خلقاً یرون بأبصارهم
و یذکرون ما ورد علیک خذ القلم عن ذکر هؤلآء ثمّ حرّکه علی ذکر مالک القدم دع
الممکنات ثمّ اشرب من رحیق ذکری المختوم ایّاک ان تشتغل بذکر الّذین لن تجد
منهم الّا روائح البغضآء و اخذهم حبّ الرّیاسة علی مقام یهلکون انفسهم لاعلآء
ذکرهم و ابقآء اسمائهم قد کتب اللّه هؤلآء من عبدة الأسمآء فی لوح محفوظ ان
اذکر ما اردته لهذا الهیکل لیظهر فی الأرض آثاره و یملأ الآفاق انوار هذا
الاشراق و یطهّر الأرض من دنس الّذین کفروا باللّه کذلک نزّلنا الآیات و فصّلنا
الأمر لقوم یعرفون

\end{otherlanguage}

In yā qalam al-baqāʾ, lā taḥzan ʿammā warada ʿalayka. Fa-sawfa yabʿathu
Allāh khalqan yarawna bi-abṣārihim wa-yadhkurūna mā warada ʿalayka.
Khudh al-qalam ʿan dhikr hāʾulāʾ thumma ḥarrikhu ʿalā dhikr mālik
al-qidam. Daʿ al-mumkināt thumma ishrab min raḥīq dhikrī al-makhtūm.
Iyyāka an tashtaghila bi-dhikr alladhīna lan tajida minhum illā rawāʾiḥ
al-baghḍāʾ, wa-akhadhahum ḥubb al-riyāsa ʿalā maqām yuhlikūna anfusahum
li-iʿlāʾ dhikrihim wa-ibqāʾ asmāʾihim. Qad kataba Allāh hāʾulāʾ min
ʿabadat al-asmāʾ fī lawḥ maḥfūẓ. An udhkur mā aradtahu li-hādhā
al-haykal li-yaẓhara fī al-arḍ āthāruhu wa-yamlaʾa al-āfāq anwār hādhā
al-ishrāq wa-yuṭahhira al-arḍ min danas alladhīna kafarū bi-Allāh.
Kadhālika nazzalnā al-āyāt wa-faṣṣalnā al-amr li-qawm yaʿrifūn.

O Pen of eternity! Do not grieve over what has befallen you. God will
raise up a people who will see with their own eyes and recount what has
befallen you. Withdraw the pen from mentioning these people, then move
it in remembrance of the Possessor of eternity. Leave contingent things
behind, then drink of the sealed wine of My remembrance. Beware of
occupying yourself with the mention of those from whom you will find
nothing but the odors of hatred. Love of leadership has seized them to
such a degree that they destroy themselves to exalt their own renown and
perpetuate their names. God has inscribed them among the worshippers of
names in a preserved Tablet. Recount what you intended for this Temple,
so that its signs may appear on earth, the lights of this radiance may
fill the horizons, and it may cleanse the earth of the defilement of
those who have disbelieved in God. Thus have We sent down the verses and
explained the matter for a people who understand.

\subsection{33}\label{section-1905}

\begin{otherlanguage}{arabic}

ان یا هذا الهیکل فابسط یدک علی من فی السّموات و الأرض ثمّ خذ زمام الأمر
بقبضة ارادتک انّا جعلنا فی یمینک ملکوت کلّ شیء ان افعل ما شئت و لا تخف من
الّذینهم لا یعرفون ثمّ ارفع یدک الی اللّوح الّذی اشرق من افق اصبع ربّک و خذه
علی شأن بأخذک تأخذه ایادی من فی الابداع کذلک ینبغی لک ان انت من الّذینهم
یفقهون و بارتفاع یدک الی سمآء فضلی ترتفع ایادی کلّ شیء الی اللّه المقتدر
العزیز الودود سوف نبعث من یدک ایادی القوّة و القدرة و الاقتدار و نظهر بها
قدرتی لمن فی ملکوت الأمر و الخلق لیعرفنّ العباد انّه لا اله الّا انا
المهیمن القیّوم و بها نعطی و نأخذ و لا یعرف ذلک الّا الّذینهم ببصر الرّوح
ینظرون

\end{otherlanguage}

In yā hādhā al-haykal, fa-bsuṭ yadaka ʿalā man fī al-samāwāt wa-al-arḍ,
thumma khudh zimām al-amr bi-qabḍat irādatika. Innā jaʿalnā fī yamīnika
malakūt kull shayʾ. In ifʿal mā shiʾta wa-lā takhaf min alladhīna hum lā
yaʿrifūn. Thumma irfaʿ yadaka ilā al-lawḥ alladhī ashraqa min ufuq iṣbaʿ
rabbika wa-khudhhu ʿalā shaʾn bi-akhdhika taʾkhudhuhu ayādī man fī
al-ibdāʿ. Kadhālika yanbaghī laka in anta min alladhīna hum yafqahūn.
Wa-bi-irtifāʿ yadika ilā samāʾ faḍlī tartafiʿ ayādī kull shayʾ ilā Allāh
al-Muqtadir al-ʿAzīz al-Wadūd. Sawfa nabʿathu min yadika ayādī al-quwwa
wa-al-qudra wa-al-iqtidār wa-nuẓhiru bihā qudratī li-man fī malakūt
al-amr wa-al-khalq li-yaʿrifanna al-ʿibād annahu lā ilāha illā anā
al-Muhaymin al-Qayyūm. Wa-bihā nuʿṭī wa-naʾkhudh wa-lā yaʿrifu dhālika
illā alladhīna hum bi-baṣar al-rūḥ yanẓurūn.

O Temple! Extend your hand over all who are in the heavens and on earth,
then grasp the reins of the Cause in the grip of your will. We have
placed the kingdom of all things in your right hand. Do whatever you
wish, and do not fear those who do not know. Then raise your hand to the
Tablet that has shone from the horizon of your Lord's finger, and grasp
it in such a way that, through your grasp, the hands of all in creation
may grasp it. This is what befits you, if you are among those who
understand. As your hand rises toward the heaven of My grace, the hands
of all things rise toward God, the All-Powerful, the Almighty, the
All-Loving. From your hand We will raise up hands of strength, power,
and might, and through them manifest My power to all in the kingdom of
command and creation, so that the servants may know that there is no God
but Me, the All-Protecting, the Self-Subsisting. Through them We give
and take, and none knows this except those who see with the eye of the
spirit.

\subsection{34}\label{section-1906}

\begin{otherlanguage}{arabic}

قل یا قوم أ تفرّون من قدرة اللّه تاللّه لا مهرب لکم الیوم و لا عاصم لأحد
الّا من رحمه اللّه بفضل من عنده انّه لهو الرّحیم الغفور قل یا قوم دعوا ما
عندکم ثمّ ادخلوا فی ظلّ ربّکم الرّحمن هذا خیر لکم عمّا عملتم او تعملون خافوا
عن اللّه و لا تحرموا انفسکم من نفحات ایّام الله مالک الأسمآء و الصّفات و لا
تبدّلوا کلمة اللّه و لا تحرّفوها عن مقرّها اتّقوا اللّه و کونوا من الّذینهم
یتّقون

\end{otherlanguage}

Qul yā qawm, a-tafirrūna min qudrat Allāh? Ta-Allāh, lā mahraba lakum
al-yawm wa-lā ʿāṣima li-aḥad illā man raḥimahu Allāh bi-faḍl min
ʿindihi. Innahu la-huwa al-Raḥīm al-Ghafūr. Qul yā qawm, daʿū mā
ʿindakum thumma udkhulū fī ẓill rabbikum al-Raḥmān. Hādhā khayr lakum
ʿammā ʿamiltum aw taʿmalūn. Khāfū ʿan Allāh wa-lā taḥrimū anfusakum min
nafaḥāt ayyām Allāh Mālik al-asmāʾ wa-al-ṣifāt. Wa-lā tubaddilū kalimat
Allāh wa-lā tuḥarrifūhā ʿan maqarrihā. Ittaqū Allāh wa-kūnū min
alladhīna hum yattaqūn.

Say: People! Do you flee from the power of God? By God, today you have
no escape, nor does anyone have a protector except those to whom God
shows mercy through His grace. He is indeed the Merciful, the
Ever-Forgiving. Say: People! Leave what you possess, then enter the
shadow of your Lord, the All-Merciful. This is better for you than
whatever you have done or will do. Fear God, and do not deprive
yourselves of the breezes of the days of God, the Possessor of names and
attributes. Do not alter the word of God or displace it from its proper
place. Be mindful of God and be among those who fear Him.

\subsection{35}\label{section-1907}

\begin{otherlanguage}{arabic}

قل یا قوم هذه ید اللّه الّتی لم تزل کانت فوق ایدیکم ان انتم تعقلون و فیها
قدّرنا خیر السّموات و الأرض بحیث لا یظهر من خیر الّا و قد یظهر منها کذلک
جعلناها مطلع الخیر و مخزنه فیما کان و ما یکون قل کلّ ما جری فی الألواح من
انهار المعانی و البیان قد اتّصلت بهذا البحر الأعظم ان انتم تشعرون و ما
فصّل فی الکتب قد انتهی الی هذه الکلمة العلیا الّتی اشرقت من افق فم مشیّة
الأبهی فی هذا الظّهور الّذی به افترّ ثغر الغیب و الشّهود

\end{otherlanguage}

Qul yā qawm, hādhihi yad Allāh allatī lam tazal kānat fawqa aydīkum in
antum taʿqilūn. Wa-fīhā qaddarnā khayr al-samāwāt wa-al-arḍ bi-ḥaythu lā
yaẓharu min khayr illā wa-qad yaẓharu minhā. Kadhālika jaʿalnāhā maṭlaʿ
al-khayr wa-makhzanahu fīmā kāna wa-mā yakūn. Qul kull mā jarā fī
al-alwāḥ min anhār al-maʿānī wa-al-bayān qad ittaṣalat bi-hādhā al-baḥr
al-aʿẓam in antum tashʿurūn. Wa-mā fuṣṣila fī al-kutub qad intahā ilā
hādhihi al-kalima al-ʿulyā allatī ashraqat min ufuq fam mashīʾat al-Abhā
fī hādhā al-ẓuhūr alladhī bihi iftarra thaghr al-ghayb wa-al-shuhūd.

Say: People! This is the hand of God that has always been above your
hands, if you understand. Within it We have ordained the good of the
heavens and the earth, so that no good appears that does not issue from
it. Thus We have made it the dawning place and treasury of good in all
that has been and will be. Say: All the rivers of meaning and utterance
that have flowed in the Tablets have joined this Most Great Ocean, if
you perceive. Whatever has been set forth in detail in the Books has
reached its culmination in this most exalted Word, which has shone from
the horizon of the mouth of the will of the Most Glorious in this
Revelation, through which the unseen and the visible have broken into a
smile.

\subsection{36}\label{section-1908}

\begin{otherlanguage}{persian}

سوف یخرج اللّه من اکمام القدرة ایادی القوّة و الغلبة و ینصرنّ الغلام و
یطهّرنّ الأرض من دنس کلّ مشرک مردود و یقومنّ علی الأمر و یفتحنّ البلاد باسمی
المقتدر القیّوم و یدخلنّ خلال الدّیار و یأخذ رعبهم کلّ العباد هذا من بطش
اللّه انّ بطشه شدید بالعدل انّه لمحیط علی من فی السّموات و الأرض ینزّل ما
یشآء علی قدر مقدور

\end{otherlanguage}

Sawfa yukhriju Allāh min akmām al-qudra ayādī al-quwwa wa-al-ghalaba
wa-yanṣuranna al-ghulām wa-yuṭahhiranna al-arḍ min danas kull mushrik
mardūd wa-yaqūmanna ʿalā al-amr wa-yaftaḥanna al-bilād bi-ismī
al-Muqtadir al-Qayyūm wa-yadkhulanna khilāl al-diyār wa-yaʾkhudhu
ruʿbuhum kull al-ʿibād. Hādhā min baṭsh Allāh, inna baṭshahu shadīd
bi-al-ʿadl. Innahu la-muḥīṭ ʿalā man fī al-samāwāt wa-al-arḍ, yunazzilu
mā yashāʾ ʿalā qadar maqdūr.

God will bring forth from the sheaths of power hands of strength and
triumph. They will assuredly aid the Youth, cleanse the earth of the
defilement of every rejected idolater, arise for the Cause, and conquer
the lands through My name, the All-Powerful, the Self-Subsisting. They
will enter the midst of the dwellings, and dread of them will seize all
people. This is an act of God's might; His grasp is indeed severe in its
justice. He encompasses all who are in the heavens and on earth and
sends down whatever He wills in a determined measure.

\subsection{37}\label{section-1909}

\begin{otherlanguage}{arabic}

و لو یقوم احد من هؤلآء فی مقابلة ما خلق فی الابداع لیکون غالباً بغلبة
ارادتی هذا من قدرتی ولکن خلقی لا یعرفون و هذا من سلطنتی ولکن بریّتی لا
یفقهون و هذا من امری ولکن عبادی لا یشعرون و هذا من غلبتی ولکنّ النّاس لا
یشکرون الّا الّذین نوّر اللّه ابصارهم بنور عرفانه و جعل قلوبهم خزائن وحیه و
انفسهم حملة امره اولئک یجدون روائح الرّحمن من قمیص اسمه و هم فی کلّ
الأحیان بآیات اللّه یفرحون و الّذین‌هم کفروا و اشرکوا اولئک غضب اللّه علیهم
و هم الی النّار یسحبون ثمّ فی اطباقها یجزعون کذلک نفصّل الآیات و نبیّن الحقّ
بالبیّنات لعلّ النّاس فی آیات ربّهم یتفکّرون

\end{otherlanguage}

Wa-law yaqūmu aḥad min hāʾulāʾ fī muqābalat mā khuliqa fī al-ibdāʿ
la-yakūnu ghāliban bi-ghalabat irādatī. Hādhā min qudratī wa-lākin
khalqī lā yaʿrifūn. Wa-hādhā min salṭanatī wa-lākin bariyyatī lā
yafqahūn. Wa-hādhā min amrī wa-lākin ʿibādī lā yashʿurūn. Wa-hādhā min
ghalabatī wa-lākinna al-nās lā yashkurūn illā alladhīna nawwara Allāh
abṣārahum bi-nūr ʿirfānihi wa-jaʿala qulūbahum khazāʾin waḥyihi
wa-anfusahum ḥamalat amrihi. Ulāʾika yajidūna rawāʾiḥ al-Raḥmān min
qamīṣ ismihi wa-hum fī kull al-aḥyān bi-āyāt Allāh yafraḥūn.
Wa-alladhīna hum kafarū wa-ashrakū ulāʾika ghaḍiba Allāh ʿalayhim wa-hum
ilā al-nār yusḥabūn, thumma fī aṭbāqihā yajzaʿūn. Kadhālika nufaṣṣilu
al-āyāt wa-nubayyinu al-ḥaqq bi-al-bayyināt laʿalla al-nās fī āyāt
rabbihim yatafakkarūn.

Were one of these to stand against all that has been brought into being
in creation, he would prevail through the ascendancy of My will. This is
of My power, yet My creatures do not know. This is of My sovereignty,
yet My creation does not understand. This is of My command, yet My
servants do not perceive. This is of My ascendancy, yet people give no
thanks, except those whose eyes God has illumined with the light of His
recognition, whose hearts He has made treasuries of His revelation, and
whose selves He has made bearers of His Cause. They discern the
fragrances of the All-Merciful from the garment of His name, and at all
times they rejoice in the verses of God. But those who have disbelieved
and associated partners with God have incurred His wrath. They will be
dragged into the fire and then cry out in anguish in its depths. Thus We
set forth the verses in detail and make the truth clear through evident
proofs, so that people may reflect on the verses of their Lord.

\subsection{38}\label{section-1910}

\begin{otherlanguage}{arabic}

ان یا هذا الهیکل قد جعلناک آیة عزّی بین ما کان و ما یکون و جعلناک آیة
امری بین السّموات و الأرض بقولی کن فیکون

\end{otherlanguage}

In yā hādhā al-haykal, qad jaʿalnāka āyat ʿizzī bayna mā kāna wa-mā
yakūn wa-jaʿalnāka āyat amrī bayna al-samāwāt wa-al-arḍ bi-qawlī kun
fa-yakūn.

O Temple! We have made you the sign of My glory amid all that has been
and will be, and We have made you the sign of My command between the
heavens and the earth through My saying, ``Be,'' and it is.

\subsection{39}\label{section-1911}

\begin{otherlanguage}{arabic}

ان یا هآء الهویّة فی هذا الاسم قد جعلناک مخزن مشیّتی ثمّ مکمن ارادتی لمن فی
ملکوت الأمر و الخلق فضلاً من لدن مهیمن قیّوم

\end{otherlanguage}

In yā hāʾ al-huwiyya fī hādhā al-ism, qad jaʿalnāka makhzan mashīʾatī
thumma makman irādatī li-man fī malakūt al-amr wa-al-khalq faḍlan min
ladun Muhaymin Qayyūm.

O H of Essence in this name! We have made you the treasury of My will
and the hidden repository of My purpose for all in the kingdom of
command and creation, as a grace from One who is All-Protecting,
Self-Subsisting.

\subsection{40}\label{section-1912}

\begin{otherlanguage}{arabic}

ان یا یآء اسمی القدیر قد جعلناک مظهر سلطانی و مطلع اسمائی و انا المقتدر
علی ما اقول

\end{otherlanguage}

In yā yāʾ ismī al-Qadīr, qad jaʿalnāka maẓhar sulṭānī wa-maṭlaʿ asmāʾī
wa-anā al-Muqtadir ʿalā mā aqūl.

O Y of My name, the Omnipotent! We have made you the manifestation of My
sovereignty and the dawning place of My names, and I have the power to
accomplish whatever I say.

\subsection{41}\label{section-1913}

\begin{otherlanguage}{arabic}

ان یا کاف اسمی الکریم قد جعلناک مشرق کرمی بین بریّتی و منبع جودی بین خلقی
انا المقتدر بسلطانی لن یعزب عن علمی شیء عمّا خلق بین السّموات و الأرض و
انا الحقّ علّام الغیوب

\end{otherlanguage}

In yā kāf ismī al-Karīm, qad jaʿalnāka mashriq karamī bayna bariyyatī
wa-manbaʿ jūdī bayna khalqī. Anā al-Muqtadir bi-sulṭānī, lan yaʿzuba ʿan
ʿilmī shayʾ ʿammā khuliqa bayna al-samāwāt wa-al-arḍ wa-anā al-Ḥaqq
ʿAllām al-ghuyūb.

O K of My name, the Generous! We have made you the dawning place of My
generosity among My creatures and the wellspring of My bounty amid My
creation. I am the All-Powerful through My sovereignty. Nothing created
between the heavens and the earth will escape My knowledge, and I am the
Truth, the Knower of all unseen things.

\subsection{42}\label{section-1914}

\begin{otherlanguage}{arabic}

ان انزل من سحاب کرمک ما یغنی الممکنات لا تمنع فضلک عن الوجود انّک انت
الکریم فی جبروت البقآء و ذو الفضل العظیم لمن فی ملکوت الأسمآء لا تنظر
الی النّاس و ما عندهم فانظر الی جمیل احسانک و بدائع مواهبک ان ادخل العباد
فی ظلّک الممدود ان ابسط ید الجود علی الممکنات و اصابع الکرم علی الکائنات
هذا ینبغی لک ولکنّ النّاس لا یعقلون من اقبل الیک هذا من فضلک و من اعرض انّ
ربّک لهو الغنیّ عمّا خلق فی الامکان یشهد بذلک عباد مخلصون

\end{otherlanguage}

In anzil min saḥāb karamika mā yughnī al-mumkināt. Lā tamnaʿ faḍlaka ʿan
al-wujūd. Innaka anta al-Karīm fī jabarūt al-baqāʾ wa-dhū al-faḍl
al-ʿaẓīm li-man fī malakūt al-asmāʾ. Lā tanẓur ilā al-nās wa-mā ʿindahum
fa-nẓur ilā jamīl iḥsānika wa-badāʾiʿ mawāhibika. In adkhil al-ʿibād fī
ẓillika al-mamdūd. In ubṣuṭ yad al-jūd ʿalā al-mumkināt wa-aṣābiʿ
al-karam ʿalā al-kāʾināt. Hādhā yanbaghī laka wa-lākinna al-nās lā
yaʿqilūn. Man aqbala ilayka hādhā min faḍlika wa-man aʿraḍa inna rabbaka
la-huwa al-Ghanī ʿammā khuliqa fī al-imkān. Yashhadu bi-dhālika ʿibād
mukhliṣūn.

Send down from the clouds of your generosity what will enrich all
contingent beings. Do not withhold your grace from existence. You are
the Generous in the dominion of eternity and the possessor of immense
grace for all in the kingdom of names. Do not look at people and what
they possess; look instead at the beauty of your beneficence and the
wonders of your gifts. Bring the servants beneath your outstretched
shadow. Extend the hand of bounty over all contingent beings and the
fingers of generosity over all created things. This is what befits you,
yet people do not understand. If anyone turns toward you, this is
through your grace; and if anyone turns away, your Lord is indeed
independent of all that has been created in the contingent world. His
sincere servants bear witness to this.

\subsection{43}\label{section-1915}

\begin{otherlanguage}{arabic}

سوف یبعث اللّه بک ایادی غالبة و اعضاداً قاهرة یخرجنّ عن خلف الأستار و ینصرنّ
نفس الرّحمن بین الامکان و یصیحنّ بصیحة تتمیّز منها الصّدور کذلک رقم فی لوح
مسطور و یظهرنّ بسطوة یأخذ الخوف سکّان الأرض علی شأن کلّهم یضطربون

\end{otherlanguage}

Sawfa yabʿathu Allāh bika ayādī ghāliba wa-aʿḍādan qāhira yakhrujanna
ʿan khalf al-astār wa-yanṣuranna nafs al-Raḥmān bayna al-imkān
wa-yaṣiḥnna bi-ṣayḥa tatamayyazu minhā al-ṣudūr. Kadhālika ruqima fī
lawḥ masṭūr. Wa-yaẓharanna bi-saṭwa yaʾkhudhu al-khawf sukkān al-arḍ
ʿalā shaʾn kulluhum yaḍṭaribūn.

Through you God will raise up triumphant hands and conquering arms. They
will come forth from behind the veils, aid the All-Merciful Himself amid
the contingent world, and raise a cry through which hearts will be
distinguished from one another. Thus has it been inscribed in a written
Tablet. They will appear with such might that fear will seize the
inhabitants of the earth and all will tremble.

\subsection{44}\label{section-1916}

\begin{otherlanguage}{arabic}

ایّاکم ان تسفکوا الدّمآء ان اخرجوا سیف اللّسان عن غمد البیان لأنّ به تفتح
مدائن القلوب انّا رفعنا حکم القتل عن بینکم انّ رحمتی سبقت الممکنات ان انتم
تعلمون ثمّ انصروا ربّکم الرّحمن بسیف التّبیان انّه احدّ من البیان و اغلی منه
لو انتم فی کلمات ربّکم تنظرون کذلک نزّلت جنود الوحی من شطر اللّه المهیمن
القیّوم و ظهرت جنود الالهام من مشرق الأمر من لدی الله العزیز المحبوب

\end{otherlanguage}

Iyyākum an tasfikū al-dimāʾ. In akhrijū sayf al-lisān ʿan ghimd al-bayān
li-anna bihi tuftaḥu madāʾin al-qulūb. Innā rafaʿnā ḥukm al-qatl ʿan
baynikum. Inna raḥmatī sabaqat al-mumkināt in antum taʿlamūn. Thumma
unṣurū rabbakum al-Raḥmān bi-sayf al-tibyān. Innahu aḥaddu min al-bayān
wa-aghlā minhu law antum fī kalimāt rabbikum tanẓurūn. Kadhālika
nuzzilat junūd al-waḥy min shaṭr Allāh al-Muhaymin al-Qayyūm wa-ẓaharat
junūd al-ilhām min mashriq al-amr min ladā Allāh al-ʿAzīz al-Maḥbūb.

Beware of shedding blood. Draw the sword of the tongue from the scabbard
of utterance, for through it the cities of hearts are conquered. We have
abolished the ordinance of killing among you. My mercy has preceded all
contingent beings, if you know. Then aid your Lord, the All-Merciful,
with the sword of clear exposition. It is sharper than utterance and
more precious than it, if you reflect on the words of your Lord. Thus
have the hosts of revelation been sent down from God, the
All-Protecting, the Self-Subsisting, and the hosts of inspiration
appeared from the dawning place of the Cause, from God, the Almighty,
the Beloved.

\subsection{45}\label{section-1917}

\begin{otherlanguage}{arabic}

قل قد قدّر مقادیر الأشیآء فی هذا الهیکل المخزون المشهود و کنز فیه علم
السّموات و الأرض و علم ما کان و ما یکون و رقم من اصبع صنع ربّک فی هذا
الکتاب ما یعجز عن ادراکه العارفون و خلق فیه الهیاکل الّتی ما اطّلع بها احد
الّا نفس اللّه ان انتم توقنون طوبی لمن یقرأه و یتفکّر فیه و یکون من الّذینهم
یفقهون

\end{otherlanguage}

Qul qad quddira maqādīr al-ashyāʾ fī hādhā al-haykal al-makhzūn
al-mashhūd wa-kuniza fīhi ʿilm al-samāwāt wa-al-arḍ wa-ʿilm mā kāna
wa-mā yakūn. Wa-ruqima min iṣbaʿ ṣunʿ rabbika fī hādhā al-kitāb mā
yaʿjizu ʿan idrākihi al-ʿārifūn. Wa-khuliqa fīhi al-hayākil allatī mā
iṭṭalaʿa bihā aḥad illā nafs Allāh in antum tūqinūn. Ṭūbā li-man
yaqraʾuhu wa-yatafakkaru fīhi wa-yakūnu min alladhīna hum yafqahūn.

Say: The measures of all things have been ordained in this hidden yet
visible Temple, and within it have been treasured the knowledge of the
heavens and the earth and the knowledge of all that has been and will
be. The finger of your Lord's handiwork has inscribed in this Book what
those endowed with knowledge are powerless to comprehend. Within it have
been created temples of which no one but God Himself has knowledge, if
you are certain. Blessed is the one who reads it, reflects on it, and is
among those who understand.

\subsection{46}\label{section-1918}

\begin{otherlanguage}{arabic}

قل لا یری فی هیکلی الّا هیکل اللّه و لا فی جمالی الّا جماله و لا فی کینونتی
الّا کینونته و لا فی ذاتی الّا ذاته و لا فی حرکتی الّا حرکته و لا فی سکونی
الّا سکونه و لا فی قلمی الّا قلمه العزیز المحمود قل لم یکن فی نفسی الّا
الحقّ و لا یری فی ذاتی الّا اللّه

\end{otherlanguage}

Qul lā yurā fī haykalī illā haykal Allāh wa-lā fī jamālī illā jamāluhu
wa-lā fī kaynūnatī illā kaynūnatuhu wa-lā fī dhātī illā dhātuhu wa-lā fī
ḥarakatī illā ḥarakatuhu wa-lā fī sukūnī illā sukūnuhu wa-lā fī qalamī
illā qalamuhu al-ʿAzīz al-Maḥmūd. Qul lam yakun fī nafsī illā al-Ḥaqq
wa-lā yurā fī dhātī illā Allāh.

Say: Nothing is seen in My temple but the Temple of God, nor in My
beauty but His beauty, nor in My being but His being, nor in My essence
but His essence, nor in My movement but His movement, nor in My
stillness but His stillness, nor in My pen but His Pen, the Almighty,
the Praiseworthy. Say: Within My Self there has been none but the Truth,
and in My essence none is seen but God.

\subsection{47}\label{section-1919}

\begin{otherlanguage}{arabic}

ایّاکم ان تذکروا الآیتین فی نفسی تنطق الذّرّات انّه لا اله الّا هو الواحد
الفرد العزیز الودود لم ازل کنت ناطقاً فی جبروت البقآء انّنی انا اللّه لا
اله الّا انا المهیمن القیّوم و لا ازال انطق فی ملکوت الأسمآء انّنی انا اللّه
لا اله الّا انا العزیز المحبوب قل انّ الرّبوبیّة اسمی قد خلقت لها مظاهر فی
الملک انّا کنّا منزّهاً عنها ان انتم تشهدون و الألوهیّة اسمی قد جعلنا لها
مطالع یحیطنّ العباد و یجعلنّهم عباد للّه ان انتم توقنون کذلک فاعرفوا کلّ
الأسمآء ان انتم تعرفون

\end{otherlanguage}

Iyyākum an tadhkurū al-āyatayn fī nafsī. Tanṭiqu al-dharrāt innahu lā
ilāha illā huwa al-Wāḥid al-Fard al-ʿAzīz al-Wadūd. Lam azal kuntu
nāṭiqan fī jabarūt al-baqāʾ innanī anā Allāh lā ilāha illā anā
al-Muhaymin al-Qayyūm. Wa-lā azālu anṭiqu fī malakūt al-asmāʾ innanī anā
Allāh lā ilāha illā anā al-ʿAzīz al-Maḥbūb. Qul inna al-rubūbiyya ismī
qad khalaqtu lahā maẓāhir fī al-mulk. Innā kunnā munazzahan ʿanhā in
antum tashhadūn. Wa-al-ulūhiyya ismī qad jaʿalnā lahā maṭāliʿ yuḥiṭnna
al-ʿibād wa-yajʿalnnahum ʿibādan li-Allāh in antum tūqinūn. Kadhālika
fa-ʿrifū kull al-asmāʾ in antum taʿrifūn.

Beware of speaking of two signs in My Self. The atoms proclaim that
there is no God but Him, the One, the Unique, the Almighty, the
All-Loving. From all eternity I have proclaimed in the dominion of
eternity, ``I am God; there is no God but Me, the All-Protecting, the
Self-Subsisting.'' And I will forever proclaim in the kingdom of names,
``I am God; there is no God but Me, the Almighty, the Beloved.'' Say:
Lordship is My name, for which I have created manifestations in the
world. We have ever been sanctified above it, if you bear witness. And
Divinity is My name, for which We have appointed dawning places that
encompass the servants and make them servants of God, if you are
certain. Understand all the names in this way, if you understand.

\subsection{48}\label{section-1920}

\begin{otherlanguage}{arabic}

ان یا لام الفضل فی هذا الاسم انّا جعلناک مظهر الفضل بین السّموات و الأرض
منک بدأنا بالفضل بین الممکنات و الیک نرجعه ثمّ منک نظهره مرّة اخری امراً من
لدنّا و انا الفاعل لما اشآء بقولی کن فیکون کلّ فضل ظهر فی الملک بدئ منک و
الیک یعود هذا ما قدّر فی لوح حفظناه خلف سرادق العظمة و عصمناه عن مشاهدة
العیون فیا حبّذا لمن لم یحرم نفسه عن هذا الفضل المسلسل المرسول

\end{otherlanguage}

In yā lām al-faḍl fī hādhā al-ism, innā jaʿalnāka maẓhar al-faḍl bayna
al-samāwāt wa-al-arḍ. Minka badaʾnā bi-al-faḍl bayna al-mumkināt
wa-ilayka nurjiʿuhu thumma minka nuẓhiruhu marratan ukhrā amran min
ladunnā. Wa-anā al-fāʿil li-mā ashāʾ bi-qawlī kun fa-yakūn. Kull faḍl
ẓahara fī al-mulk budiʾa minka wa-ilayka yaʿūd. Hādhā mā quddira fī lawḥ
ḥafiẓnāhu khalf surādiq al-ʿaẓama wa-ʿaṣamnāhu ʿan mushāhadat al-ʿuyūn.
Fa-yā ḥabbadhā li-man lam yaḥrim nafsahu ʿan hādhā al-faḍl al-mutasalsil
al-marsūl.

O L of grace in this name! We have made you the manifestation of grace
between the heavens and the earth. From you We began to bestow grace
among contingent beings, and to you We return it; then from you We
manifest it once again by Our command. I do whatever I will through My
saying, ``Be,'' and it is. Every grace that has appeared in the world
originated from you and returns to you. This is what has been ordained
in a Tablet that We have preserved behind the pavilion of grandeur and
protected from the sight of eyes. How blessed is the one who has not
deprived himself of this grace, sent forth in unbroken succession!

\subsection{49}\label{section-1921}

\begin{otherlanguage}{arabic}

قل الیوم قد هبّت لواقح الفضل علی الأشیآء و حمل کلّ شیء علی ما هو علیه ولکنّ
النّاس عنه معرضون قد حملت الأشجار بالأثمار البدیعة و البحور باللّآٓلئ
المنیرة و الانسان بالمعانی و العرفان و الأکوان بتجلّیات الرّحمن و الأرض
بما لا اطّلع به احد الّا الحقّ علّام الغیوب سوف یضعن کلّ حملها تبارک اللّه
مرسل هذا الفضل الّذی احاط الأشیآء کلّها عمّا ظهر و عمّا هو المکنون کذلک
خلقنا الأکوان بدعاً فی هذا الیوم ولکنّ النّاس اکثرهم لا یشعرون قل لن یعرف
فضل اللّه علی ما هو علیه فکیف نفسه المهیمن القیّوم

\end{otherlanguage}

Qul al-yawm qad habbat lawāqiḥ al-faḍl ʿalā al-ashyāʾ wa-ḥamala kull
shayʾ ʿalā mā huwa ʿalayhi wa-lākinna al-nās ʿanhu muʿriḍūn. Qad ḥamalat
al-ashjār bi-al-athmār al-badīʿa wa-al-buḥūr bi-al-laʾāliʾ al-munīra
wa-al-insān bi-al-maʿānī wa-al-ʿirfān wa-al-akwān bi-tajalliyāt
al-Raḥmān wa-al-arḍ bi-mā lā iṭṭalaʿa bihi aḥad illā al-Ḥaqq ʿAllām
al-ghuyūb. Sawfa yaḍaʿna kull ḥamlahā. Tabāraka Allāh mursil hādhā
al-faḍl alladhī aḥāṭa al-ashyāʾ kullahā ʿammā ẓahara wa-ʿammā huwa
al-maknūn. Kadhālika khalaqnā al-akwān bidʿan fī hādhā al-yawm
wa-lākinna al-nās aktharuhum lā yashʿurūn. Qul lan yuʿrafa faḍl Allāh
ʿalā mā huwa ʿalayhi fa-kayfa nafsuhu al-Muhaymin al-Qayyūm.

Say: Today the fertilizing winds of grace have blown over all things,
and everything has conceived according to its nature, yet people turn
away from it. The trees have conceived wondrous fruits, the seas radiant
pearls, the human being meanings and understanding, created things the
effulgences of the All-Merciful, and the earth what no one knows but the
Truth, the Knower of all unseen things. Soon they will each bring forth
their burden. Blessed is God, who sends forth this grace that has
encompassed all things, both manifest and hidden. Thus We have created
all things anew on this day, yet most people do not perceive. Say: The
grace of God can never be known as it truly is; how much less His own
Self, the All-Protecting, the Self-Subsisting!

\subsection{50}\label{section-1922}

\begin{otherlanguage}{arabic}

ان یا هیکل الأمر ان لن تجد مقبلاً الی مواهبک لا تحزن قد خلقت لنفسی ان
اشتغل بذکری بین عبادی هذا ما قدّر لک فی لوح محفوظ انّا لمّا وجدنا الأیادی
غیر طاهرة فی الأرض لذا جعلنا ذیلک مطهّراً عن مسّها و مسّ الّذینهم مشرکون ان
اصبر فی امر ربّک سوف یبعث اللّه افئدة طاهرة و ابصار منیرة یهربنّ من کلّ
الجهات الی جهة فضلک المحیط المبسوط

\end{otherlanguage}

In yā haykala al-amr, in lan tajid muqbilan ilā mawāhibika lā taḥzan.
Qad khuliqta li-nafsī. In ishtag hil bi-dhikrī bayna ʿibādī. Hādhā mā
quddira laka fī lawḥin maḥfūẓ. Innā lammā wajadnā al-ayādiya ghayra
ṭāhiratin fī al-arḍ li-dhā jaʿalnā dhaylaka muṭahharan ʿan massihā
wa-massi alladhīna hum mushrikūn. In iṣbir fī amri rabbika. Sawfa
yabʿathu Allāhu afʾidatan ṭāhiratan wa-abṣāran munīratan yahrubanna min
kulli al-jihāti ilā jihati faḍlika al-muḥīṭi al-mabsūṭ.

O Temple of the Cause! If you find no one turning toward your gifts, do
not grieve. You have been created for My own Self. Occupy yourself with
My remembrance among My servants. This is what has been ordained for you
in a preserved Tablet. Since We found the hands on earth unclean, We
have kept the hem of your robe pure from their touch and from the touch
of those who associate partners with God. Be patient in the Cause of
your Lord. God will raise up pure hearts and radiant eyes that will
hasten from every direction toward your all-encompassing, outspread
grace.

\subsection{51}\label{section-1923}

\begin{otherlanguage}{arabic}

ان یا هیکل اللّه لمّا نزّلت جنود الوحی برایات الآیات من ملیک الأسمآء و
الصّفات انهزموا اولو الاشارات و کفروا ببیّنات اللّه المهیمن القیّوم و قاموا
علی النّفاق منهم من قال لیست هذه الآیات بیّنات من اللّه و ما نزّلت علی
الفطرة کذلک یداوون المشرکون جرح الصّدور و بذلک یلعنهم من فی السّموات و
الأرض و هم فی انفسهم لا یشعرون

\end{otherlanguage}

In yā haykala Allāh, lammā nuzzilat junūdu al-waḥyi bi-rāyāti al-āyāti
min malīki al-asmāʾi wa-al-ṣifāt, inhazamū ulū al-ishārāti wa-kafarū
bi-bayyināti Allāhi al-Muhaymini al-Qayyūm, wa-qāmū ʿalā al-nifāq.
Minhum man qāla laysat hādhihi al-āyātu bayyinātin mina Allāhi wa-mā
nuzzilat ʿalā al-fiṭra. Kadhālika yudāwūna al-mushrikūna jurḥa al-ṣudūr,
wa-bi-dhālika yalʿanuhum man fī al-samāwāti wa-al-arḍ, wa-hum fī
anfusihim lā yashʿurūn.

O Temple of God! When the hosts of revelation were sent down with the
banners of verses from the Sovereign of names and attributes, those
versed in allusions fled, rejected the clear proofs of God, the
All-Protecting, the Self-Subsisting, and rose up in hypocrisy. Among
them was one who said, ``These verses are not clear proofs from God and
were not revealed through innate inspiration.'' In this way those who
associate partners with God treat the wound in their breasts. For this,
all who are in the heavens and on earth curse them, yet they themselves
remain unaware.

\subsection{52}\label{section-1924}

\begin{otherlanguage}{arabic}

قل انّ روح القدس قد خلق بحرف ممّا نزّل من هذا الرّوح الأعظم ان انتم تفقهون و
انّ الفطرة بکینونتها قد خلقت من آیات اللّه المهیمن العزیز المحبوب قل انّها
تفتخر بنسبتها الی نفسنا الحقّ و انّا لا نفتخر بها و بما دونها لأنّ دونی قد
خلق بقولی ان انتم تعقلون

\end{otherlanguage}

Qul inna rūḥa al-qudusi qad khuliqa bi-ḥarfin mimmā nuzzila min hādhā
al-rūḥi al-aʿẓam, in antum tafqahūn. Wa-inna al-fiṭrata bi-kaynūnatihā
qad khuliqat min āyāti Allāhi al-Muhaymini al-ʿAzīzi al-Maḥbūb. Qul
innahā taftakhiru bi-nisbatihā ilā nafsinā al-Ḥaqq, wa-innā lā
naftakhiru bihā wa-bi-mā dūnahā, li-anna dūnī qad khuliqa bi-qawlī, in
antum taʿqilūn.

Say: The Holy Spirit was created through a single letter of what has
been sent down from this Most Great Spirit, if you understand. Innate
inspiration, in its very essence, was created from the verses of God,
the All-Protecting, the Almighty, the Beloved. Say: It takes pride in
its relationship to Our own Self, the Truth, but We take no pride in it
or in anything besides it, for all that is other than Me was created
through My word, if you would understand.

\subsection{53}\label{section-1925}

\begin{otherlanguage}{arabic}

قل انّا نزّلنا الآیات علی تسعة شؤون کلّ شأن منها یدلّ علی سلطنة اللّه المهیمن
القیّوم شأن منها یکفینّ فی الحجّیّة من فی السّموات و الأرض ولکنّ النّاس اکثرهم
غافلون و لو شئنا لنزّلنا علی شؤون اخری الّتی لا یحصی عدّتها المحصون

\end{otherlanguage}

Qul innā nazzalnā al-āyāti ʿalā tisʿati shuʾūn. Kullu shaʾnin minhā
yadullu ʿalā salṭanati Allāhi al-Muhaymini al-Qayyūm. Shaʾnun minhā
yakfiyanna fī al-ḥujjiyyati man fī al-samāwāti wa-al-arḍ, wa-lākinna
al-nāsa aktharuhum ghāfilūn. Wa-law shiʾnā la-nazzalnā ʿalā shuʾūnin
ukhrā allatī lā yuḥṣī ʿiddatahā al-muḥṣūn.

Say: We have revealed the verses in nine modes, each of which attests to
the sovereignty of God, the All-Protecting, the Self-Subsisting. A
single one of these modes is sufficient as proof for all who are in the
heavens and on earth, yet most people are heedless. If We wished, We
would reveal them in other modes whose number no reckoners could count.

\subsection{54}\label{section-1926}

\begin{otherlanguage}{persian}

قل یا قوم خافوا عن اللّه و لا تحرّکوا السنتکم الکذبة علی ما لا یحبّه اللّه
فاستحیوا عن الّذی خلقکم بقطرة من المآء کما انتم تعلمون قل انّا خلقنا من فی
السّموات و الأرض علی فطرة اللّه فمن اقبل الی هذا الوجه یظهر علی ما خلق
علیه و من احتجب یحتجب عن هذا الفضل المحیط المکنون انّا ما منعنا شیئاً عن
فضل قد خلقنا الأشیآء علی حدّ سوآء و عرضنا علیها امانة حبّنا بکلمة من لدنّا
فمن حمل نجا و امن و کان من الّذینهم من فزع الیوم آمنون و من اعرض کفر
باللّه المهیمن القیّوم و بها فرّقنا بین العباد و فصّلنا بینهم انّا نحن فاصلون

\end{otherlanguage}

Qul yā qawmi khāfū ʿani Allāhi wa-lā tuḥarrikū alsinatakumu al-kadhibata
ʿalā mā lā yuḥibbuhu Allāh. Fa-staḥyū ʿani alladhī khalaqakum
bi-qaṭratin mina al-māʾi kamā antum taʿlamūn. Qul innā khalaqnā man fī
al-samāwāti wa-al-arḍi ʿalā fiṭrati Allāh. Fa-man aqbala ilā hādhā
al-wajhi yaẓharu ʿalā mā khuliqa ʿalayhi, wa-man iḥtajaba yaḥtajibu ʿan
hādhā al-faḍli al-muḥīṭi al-maknūn. Innā mā manaʿnā shayʾan ʿan faḍl.
Qad khalaqnā al-ashyāʾa ʿalā ḥaddin sawāʾ, wa-ʿaraḍnā ʿalayhā amānata
ḥubbinā bi-kalimatin min ladunnā. Fa-man ḥamala najā wa-amina wa-kāna
mina alladhīna hum min fazaʿi al-yawmi āminūn. Wa-man aʿraḍa kafara
bi-Allāhi al-Muhaymini al-Qayyūm. Wa-bihā farraqnā bayna al-ʿibādi
wa-faṣṣalnā baynahum. Innā naḥnu fāṣilūn.

Say: O people, fear God, and do not move your lying tongues to utter
what God does not love. Feel shame before the One who created you from a
drop of water, as you know. Say: We created all who are in the heavens
and on earth according to God's pattern. Whoever turns toward this Face
will manifest the nature in which he was created, while whoever veils
himself will be veiled from this all-encompassing, hidden grace. We have
withheld grace from nothing. We created all things on an equal footing
and, through a word from Us, offered them the trust of Our love. Whoever
bore it was saved and secure and numbered among those who are safe from
the terror of this Day. Whoever turned away disbelieved in God, the
All-Protecting, the Self-Subsisting. Through it We have distinguished
between the servants and separated them from one another. We are indeed
the One who separates.

\subsection{55}\label{section-1927}

\begin{otherlanguage}{arabic}

قل کلمة اللّه لن تشتبه بکلمات خلقه انّها سلطان الکلمات کما انّ نفسه سلطان
النّفوس و امره مهیمن علی ما کان و ما یکون ان ادخلوا یا قوم مصر الایقان
مقرّ عرش ربّکم الرّحمن هذا ما یأمرکم به قلم السّبحان فضلاً من عنده علیکم ان
انتم فی امره لا تختلفون

\end{otherlanguage}

Qul kalimatu Allāhi lan tashtabiha bi-kalimāti khalqih. Innahā sulṭānu
al-kalimāti kamā anna nafsahu sulṭānu al-nufūs, wa-amruhu muhayminun
ʿalā mā kāna wa-mā yakūn. In udkhulū yā qawmi miṣra al-īqāni maqarra
ʿarshi rabbikumu al-Raḥmān. Hādhā mā yaʾmurukum bihi qalamu al-Subḥāni
faḍlan min ʿindihi ʿalaykum, in antum fī amrihi lā takhtalifūn.

Say: The Word of God can never be confused with the words of His
creatures. It is the sovereign of words, just as His own Self is the
sovereign of all souls, and His Cause rules over all that has been and
all that will be. Enter, O people, the city of certitude, the seat of
the throne of your Lord, the All-Merciful. This is what the Pen of the
All-Glorious commands you, as a grace bestowed upon you from His
presence, if you will not dispute His Cause.

\subsection{56}\label{section-1928}

\begin{otherlanguage}{arabic}

و من المشرکین من کفر فی نفسه و قام بالمحاربة و قال هذه الآیات مفتریات
کذلک قالوا من قبل العباد الّذین مضوا و اذاً فی النّار هم یستغیثون قل ویل
لکم بما یخرج من افواهکم ان کانت الآیات مفتریات فبأیّ حجّة آمنتم باللّه
فأتوا بها ان انتم تفقهون کلّما نزّلنا علیهم آیات بیّنات کفروا بها و اذا
رأوا ما عجزت عن الاتیان بمثلها کلّ الوری قالوا هذا سحر

\end{otherlanguage}

Wa-mina al-mushrikīna man kafara fī nafsihi wa-qāma bi-al-muḥārabati
wa-qāla hādhihi al-āyātu muftarayāt. Kadhālika qālū min qablu al-ʿibādu
alladhīna maḍaw, wa-idhan fī al-nāri hum yastaghīthūn. Qul waylun lakum
bi-mā yakhruju min afwāhikum. In kānati al-āyātu muftarayātin fa-bi-ayyi
ḥujjatin āmantum bi-Allāh? Faʾtū bihā in antum tafqahūn. Kullamā
nazzalnā ʿalayhim āyātin bayyinātin kafarū bihā, wa-idhā raʾaw mā
ʿajazat ʿani al-ityāni bi-mithlihā kullu al-warā qālū hādhā siḥr.

Among those who associate partners with God is one who disbelieved in
his heart, rose up in opposition, and said, ``These verses are
fabrications.'' So spoke the people who passed away before them, and now
they cry for help in the Fire. Say: Woe to you for what comes from your
mouths! If the verses are fabrications, by what proof have you believed
in God? Bring it forth, if you understand. Whenever We sent down clear
verses to them, they rejected them, and when they saw that the like of
which all humanity was powerless to produce, they said, ``This is
sorcery.''

\subsection{57}\label{section-1929}

\begin{otherlanguage}{arabic}

ما لهؤلآء القوم یقولون ما لا یعلمون کذلک قالت امّة الفرقان حین ‌الّذی اتی
اللّه بأمره الا انّهم قوم منکرون و منعوا النّاس عن الحضور بین یدی جمال
القدم و الأکل مع احبّائه و قال قائل منهم لا تقربوا هؤلآء انّهم یسحرون
النّاس و یضلّونهم عن سبیل اللّه المهیمن القیّوم تاللّه الحقّ انّ الّذی لن یقدر
ان یتکلّم بین یدینا لیقول ما لا قاله الأوّلون و ارتکب ما لا ارتکب نفس من
الّذینهم کفروا بالرّحمن فی کلّ الأعصار

\end{otherlanguage}

Mā li-hāʾulāʾi al-qawmi yaqūlūna mā lā yaʿlamūn? Kadhālika qālat ummatu
al-Furqāni ḥīna alladhī atā Allāhu bi-amrih. Alā innahum qawmun
munkirūn. Wa-manaʿū al-nāsa ʿani al-ḥuḍūri bayna yaday jamāli al-qidami
wa-al-akli maʿa aḥibbāʾih. Wa-qāla qāʾilun minhum lā taqrabū hāʾulāʾ,
innahum yasḥarūna al-nāsa wa-yuḍillūnahum ʿan sabīli Allāhi al-Muhaymini
al-Qayyūm. Ta-Allāhi al-Ḥaqq, inna alladhī lan yaqdira an yatakallama
bayna yadaynā la-yaqūlu mā lā qālahu al-awwalūn, wa-irtakaba mā lā
irtakaba nafsun mina alladhīna hum kafarū bi-al-Raḥmāni fī kulli
al-aʿṣār.

What is wrong with these people that they say what they do not know? So
spoke the people of the Qur'an when God came with His Cause. They are
indeed a people who deny. They prevented people from entering the
presence of the Ancient Beauty and eating with His loved ones. One of
them said, ``Do not approach these people. They bewitch people and lead
them astray from the path of God, the All-Protecting, the
Self-Subsisting.'' By God, the Truth! The one who cannot speak in Our
presence says what those of former times never said and has committed
what no one who disbelieved in the All-Merciful in any age ever
committed.

\subsection{58}\label{section-1930}

\begin{otherlanguage}{arabic}

یشهد بذلک اقوالهم و افعالهم لو انتم تنصفون من نسب آیات اللّه بالسّحر انّه
ما آمن بأحد من رسل اللّه قد ضلّ سعیه فی الحیوة الباطلة و کان من الّذین
یقولون ما لا یعلمون قل یا عبد خف من اللّه الّذی خلقک و سوّاک و لا تفرّط فی
جنب اللّه ثمّ انصف فی نفسک و کن من الّذینهم یعدلون انّ الّذین اوتوا العلم من
اللّه اولئک یجدنّ من اعتراضاتهم دلائل قویّة فی ابطالهم و اثبات هذا النّور
المشهود قل أ تقولون ما لا قاله المشرکون اذ جآءهم ذکر من ربّهم فویل لکم یا
معشر الجهلآء و بئس ما انتم تکسبون

\end{otherlanguage}

Yashhadu bi-dhālika aqwāluhum wa-afʿāluhum law antum tunṣifūn. Man
nasaba āyāti Allāhi bi-al-siḥri innahu mā āmana bi-aḥadin min rusuli
Allāh. Qad ḍalla saʿyuhu fī al-ḥayāti al-bāṭilati wa-kāna mina alladhīna
yaqūlūna mā lā yaʿlamūn. Qul yā ʿabdu khaf mina Allāhi alladhī khalaqaka
wa-sawwāka wa-lā tufarriṭ fī janbi Allāh, thumma anṣif fī nafsika wa-kun
mina alladhīna hum yaʿdilūn. Inna alladhīna ūtū al-ʿilma mina Allāhi
ulāʾika la-yajidanna min iʿtirāḍātihim dalāʾila qawiyyatan fī ibṭālihim
wa-ithbāti hādhā al-nūri al-mashhūd. Qul a-taqūlūna mā lā qālahu
al-mushrikūna idh jāʾahum dhikrun min rabbihim? Fa-waylun lakum yā
maʿshara al-juhalāʾi wa-biʾsa mā antum taksibūn.

Their words and deeds bear witness to this, if you judge fairly. Whoever
attributes the verses of God to sorcery has not believed in any of God's
Messengers. His efforts have gone astray in this vain life, and he is
among those who say what they do not know. Say: O servant, fear God, who
created and fashioned you, and do not neglect your duty to God. Then
judge fairly within yourself and be among those who act justly. Those
who have been given knowledge by God will surely find in their
objections powerful proofs that refute them and establish this manifest
Light. Say: Do you say what even those who associated partners with God
did not say when a Reminder came to them from their Lord? Woe to you,
company of the ignorant! How evil is what you do!

\subsection{59}\label{section-1931}

\begin{otherlanguage}{persian}

ان یا جمال القدم دع المشرکین و ما عندهم ثمّ عطّر الممکنات بذکر محبوبک
العلیّ العظیم بذکره تحیی الموجودات و تجدّد هیاکل العالمین قل انّه استقرّ علی
عرش العظمة و الجلال من اراد ان ینظر جماله هو هذا تبارک اللّه الّذی ظهر
بهذا الجمال المشرق المنیر من اراد ان یسمع نغماته انّها ارتفعت من هذا الفم
الدّرّیّ البدیع و من اراد ان یستضیء بأنواره قل ان احضر تلقآء العرش هذا ما
اذن اللّه لکم فضلاً من عنده علی العالمین

\end{otherlanguage}

In yā jamāla al-qidam, daʿi al-mushrikīna wa-mā ʿindahum, thumma ʿaṭṭiri
al-mumkināti bi-dhikri maḥbūbika al-ʿAliyyi al-ʿAẓīm. Bi-dhikrihi taḥyā
al-mawjūdātu wa-tujaddadu hayākilu al-ʿālamīn. Qul innahu istaqarra ʿalā
ʿarshi al-ʿaẓamati wa-al-jalāl. Man arāda an yanẓura jamālahu huwa
hādhā. Tabāraka Allāhu alladhī ẓahara bi-hādhā al-jamāli al-mushriqi
al-munīr. Man arāda an yasmaʿa naghamātahu innahā irtafaʿat min hādhā
al-fami al-durriyyi al-badīʿ. Wa-man arāda an yastaḍīʾa bi-anwārihi qul
in iḥḍur tilqāʾa al-ʿarsh. Hādhā mā adhina Allāhu lakum faḍlan min
ʿindihi ʿalā al-ʿālamīn.

O Ancient Beauty! Leave those who associate partners with God and all
they possess, and perfume all contingent beings with the remembrance of
your Beloved, the Most High, the Most Great. Through His remembrance all
beings are brought to life and the temples of the worlds are renewed.
Say: He has taken His seat upon the throne of grandeur and majesty.
Whoever wishes to behold His beauty---here He is! Blessed is God, who
has appeared in this radiant, luminous Beauty. Whoever wishes to hear
His melodies---they have risen from this wondrous, pearl-like mouth. To
whoever wishes to be illumined by His lights, say: Come before the
throne. This is what God has permitted you, as a grace from His presence
to all the worlds.

\subsection{60}\label{section-1932}

\begin{otherlanguage}{arabic}

قل یا قوم انّا نسأل منکم کلمةً علی الصّدق الأکبر و نتّخذ اللّه بیننا و بینکم
شهیداً انّه ولیّ المحسنین ان اجعلوا محضرکم بین یدی العرش ثمّ انصفوا فی القول
و کونوا من المنصفین أ کان اللّه مقتدراً علی امره ام انتم من القادرین أ انّه
کان مختاراً فی نفسه کما تقولون انّه یفعل ما یشآء و لا یسأل عمّا شآء ام انتم
المختارون و تقولون هذه الکلمة علی التّقلید کما تکلّم به آبائکم فی زمن
المرسلین

\end{otherlanguage}

Qul yā qawmi innā nasʾalu minkum kalimatan ʿalā al-ṣidqi al-akbari
wa-nattakhidhu Allāha baynanā wa-baynakum shahīdan. Innahu waliyyu
al-muḥsinīn. In ijʿalū maḥḍarakum bayna yaday al-ʿarshi thumma anṣifū fī
al-qawli wa-kūnū mina al-munṣifīn. A-kāna Allāhu muqtadiran ʿalā amrihi
am antum mina al-qādirīn? A-innahu kāna mukhtāran fī nafsihi kamā
taqūlūna innahu yafʿalu mā yashāʾu wa-lā yusʾalu ʿammā shāʾa am antumu
al-mukhtārūn, wa-taqūlūna hādhihi al-kalimata ʿalā al-taqlīdi kamā
takallama bihi ābāʾukum fī zamani al-mursalīn?

Say: O people, We ask you for a word spoken with the utmost
truthfulness, and We take God as a witness between Us and you. He is the
protector of those who do good. Place yourselves before the throne, then
speak fairly and be among the fair-minded. Does God have power over His
Cause, or are you the ones who possess power? Is He free to choose for
Himself, as you say, ``He does whatever He wills and is not questioned
about what He wills,'' or are you the ones who choose, uttering these
words merely in imitation, as your fathers uttered them in the days of
the Messengers?

\subsection{61}\label{section-1933}

\begin{otherlanguage}{arabic}

لو انّه کان مختاراً فی نفسه قد اظهر مظهر امره بالآیات الّتی لا یقوم معها
شیء لا فی السّموات و لا فی الأرضین و ظهر علی شأن ما ظهر فی الابداع شبهه
کما رأیتم و سمعتم اذ اشرق نیّر الآفاق من افق العراق بسلطان مبین کلّ الأمور
تنتهی الی الآیات و تلک آیات اللّه الملک المهیمن العزیز القدیر و من دونها
قد ظهر بأمر اقرّ لسلطانه کلّ الممکنات و لا ینکر ذلک الّا کلّ مشرک اثیم

\end{otherlanguage}

Law innahu kāna mukhtāran fī nafsihi qad aẓhara maẓhara amrihi
bi-al-āyāti allatī lā yaqūmu maʿahā shayʾun lā fī al-samāwāti wa-lā fī
al-araḍīn. Wa-ẓahara ʿalā shaʾnin mā ẓahara fī al-ibdāʿi shibhuhu kamā
raʾaytum wa-samiʿtum, idh ashraqa nayyiru al-āfāqi min ufuqi al-ʿIrāqi
bi-sulṭānin mubīn. Kullu al-umūri tantahī ilā al-āyāti wa-tilka āyātu
Allāhi al-Maliki al-Muhaymini al-ʿAzīzi al-Qadīr. Wa-min dūnihā qad
ẓahara bi-amrin aqarra li-sulṭānihi kullu al-mumkināt, wa-lā yunkiru
dhālika illā kullu mushrikin athīm.

If He is free to choose for Himself, He has indeed made manifest the
Manifestation of His Cause with verses that nothing in the heavens or on
earth can withstand. He has appeared in a manner whose like has never
appeared in creation, as you have seen and heard, when the Luminary of
the horizons shone from the horizon of Iraq with manifest sovereignty.
All matters ultimately rest upon the verses, and these are the verses of
God, the King, the All-Protecting, the Almighty, the All-Powerful. Apart
from these, He has appeared with a Cause to whose sovereignty all
contingent beings have acknowledged submission. None denies this except
every sinful one who associates partners with God.

\subsection{62}\label{section-1934}

\begin{otherlanguage}{arabic}

قل یا قوم أ اردتم ان تستروا جمال الشّمس بأحجاب انفسکم او تمنعوا الرّوح عن
التّغرّد فی هذا الصّدر الممرّد المنیر خافوا من اللّه و لا تحاربوا مع نفسه و
لا تجادلوا مع الّذی بأمره خلقت الکاف و اتّصلت برکنها العظیم آمنوا بسفرآء
اللّه و سلطانه ثمّ بنفس اللّه و عظمته و لا تعقّبوا الّذینهم کفروا بعد ایمانهم
و اتّخذوا لأنفسهم مقاماً فی هواهم الا انّهم من المشرکین ان اشهدوا بما شهد
اللّه لیستضیء بما یخرج من افواهکم ملأ مقرّبون قولوا انّا آمنّا بما نزّل الی
رسل اللّه من قبل و ما نزّل الی علیّ بالحقّ و ما ینزّل من جهة عرش عظیم کذلک
یعلّمکم اللّه جوداً من عنده و فضلاً من لدنه انّ فضله احاط العالمین

\end{otherlanguage}

Qul yā qawmi a-aradtum an tasturū jamāla al-shamsi bi-aḥjābi anfusikum
aw tamnaʿū al-rūḥa ʿani al-tag harrudi fī hādhā al-ṣadri al-mumarradi
al-munīr? Khāfū mina Allāhi wa-lā tuḥāribū maʿa nafsihi wa-lā tujādilū
maʿa alladhī bi-amrihi khuliqati al-kāfu wa-ittaṣalat bi-ruknihā
al-ʿaẓīm. Āminū bi-sufarāʾi Allāhi wa-sulṭānihi thumma bi-nafsi Allāhi
wa-ʿaẓamatih, wa-lā taʿaqqabū alladhīna hum kafarū baʿda īmānihim
wa-ittakhadhū li-anfusihim maqāman fī hawāhum. Alā innahum mina
al-mushrikīn. In is hhadū bi-mā shahida Allāhu li-yastaḍīʾa bi-mā
yakhruju min afwāhikum malaʾun muqarrabūn. Qūlū innā āmannā bi-mā
nuzzila ilā rusuli Allāhi min qablu wa-mā nuzzila ilā ʿAliyyin
bi-al-ḥaqqi wa-mā yunazzalu min jihati ʿarshin ʿaẓīm. Kadhālika
yuʿallimukumu Allāhu jūdan min ʿindihi wa-faḍlan min ladunh. Inna
faḍlahu aḥāṭa al-ʿālamīn.

Say: O people, do you wish to conceal the beauty of the sun with the
veils of your own selves, or prevent the Spirit from singing in this
polished, luminous breast? Fear God, do not wage war against His own
Self, and do not dispute with the One at whose command the letter Kaf
was created and joined to its mighty pillar. Believe in the emissaries
of God and His sovereignty, and then in God's own Self and His grandeur.
Do not follow those who disbelieved after believing and claimed for
themselves a station according to their own desires. They are indeed
among those who associate partners with God. Bear witness to what God
has witnessed, so that the company of those brought near may be
illumined by what comes from your mouths. Say: ``We believe in what was
sent down to the Messengers of God before, in what was sent down in
truth to Ali, and in what is being sent down from the direction of a
mighty throne.'' Thus God teaches you through His bounty and grace. His
grace has encompassed all the worlds.

\subsection{63}\label{section-1935}

\begin{otherlanguage}{persian}

ان یا رجل هذا الهیکل انّا خلقناک من الحدید ان استقم علی امر ربّک علی شأن
تستقیم به ارجل المنقطعین علی صراط ربّک العزیز الحکیم ایّاک ان تتحرّک من
عواصف البغضآء و قواصف هؤلآء الأشقیآء ان اثبت علی الأمر و کن من الثّابتین
انّا بعثناک باسمنا الّذی به استقام کلّ ذی استقامة و بکلّ اسم من اسمائنا
الحسنی بین السّموات و الأرضین سوف نبعث منک ارجلاً مستقیمة یقومنّ علی الصّراط
و لا یزلّنّ عنه ولو یحارب معهم جنود یعادل جنود الأوّلین و الآخرین انّ الفضل
کلّه فی قبضتنا نعطی من نشآء من عبادنا المقرّبین کذلک مننّا علیک مرّةً بعد مرّة
لتشکر ربّک بشکر یفتح به السن الممکنات علی شکر نفسی الرّحمن الرّحیم

\end{otherlanguage}

In yā rijla hādhā al-haykal, innā khalaqnāka mina al-ḥadīd. In istaqim
ʿalā amri rabbika ʿalā shaʾnin tastaqīmu bihi arjulu al-munqaṭiʿīna ʿalā
ṣirāṭi rabbika al-ʿAzīzi al-Ḥakīm. Iyyāka an tataḥarraka min ʿawāṣifi
al-baghḍāʾi wa-qawāṣifi hāʾulāʾi al-ashqiyāʾ. In uthbut ʿalā al-amri
wa-kun mina al-thābitīn. Innā baʿathnāka bi-isminā alladhī bihi istaqāma
kullu dhī istiqāmatin wa-bi-kulli ismin min asmāʾinā al-ḥusnā bayna
al-samāwāti wa-al-araḍīn. Sawfa nabʿathu minka arjulan mustaqīmatan
yaqūmanna ʿalā al-ṣirāṭi wa-lā yazillanna ʿanhu wa-law yuḥāribu maʿahum
junūdun yuʿādilu junūda al-awwalīna wa-al-ākhirīn. Inna al-faḍla kullahu
fī qabḍatinā nuʿṭī man nashāʾu min ʿibādinā al-muqarrabīn. Kadhālika
manannā ʿalayka marratan baʿda marratin li-tashkura rabbaka bi-shukrin
yaftaḥu bihi alsuna al-mumkināti ʿalā shukri nafsī al-Raḥmāni al-Raḥīm.

O Foot of this Temple! We have created you from iron. Stand firm in the
Cause of your Lord in such a way that the feet of the detached may stand
firm on the path of your Lord, the Almighty, the All-Wise. Beware lest
you be shaken by the storms of hatred and the tempests of these wretched
ones. Remain steadfast in the Cause and be among the steadfast. We have
raised you up through Our name, through which everyone endowed with
steadfastness has stood firm, and through every one of Our most
beautiful names between the heavens and the earth. From you We will
raise up steadfast feet that will stand upon the path and never slip
from it, even if armies equal to the hosts of the first and the last
wage war against them. All grace is in Our grasp; We bestow it upon
whomever We will among Our servants who are brought near. Thus We have
favored you time and again, so that you may thank your Lord with a
thanksgiving that will open the tongues of all contingent beings to give
thanks to My own Self, the All-Merciful, the Merciful.

\subsection{64}\label{section-1936}

\begin{otherlanguage}{persian}

قم علی الأمر بقدرة من لدنّا و سلطان من عندنا ثمّ الق العباد ما القاک روح
اللّه الملک الفرد العزیز العلیم قل یا قوم أ تدعون الحقّ عن ورائکم و تدعون
الّذی خلقناه بکفّ من الطّین هذا ظلم منکم علی انفسکم ان انتم فی آیات ربّکم
لمن المتفکّرین قل یا قوم طهّروا قلوبکم ثمّ ابصارکم لعلّ تعرفون بارئکم فی هذا
القمیص المقدّس اللّمیع قل انّ هذا فتی الهیّ قد استقرّ علی عرش الجلال و ظهر
بسلطان القدرة و الاستقلال و یصیح بین الأرض و السّمآء بندائه الأبدع الأحلی
ان یا اهل الأکوان لم کفرتم بربّکم الرّحمن و اعرضتم عن جمال السّبحان تاللّه
هذا لغیب المستور قد طلع من مشرق الامکان و هذا لجمال المحبوب قد اشرق من
افق هذا الرّضوان بسلطنة اللّه المهیمن العزیز الغالب القدیر

\end{otherlanguage}

Qum ʿalā al-amr bi-qudratin min ladunnā wa-sulṭānin min ʿindinā thumma
alqi al-ʿibāda mā alqāka Rūḥu Allāh al-Malik al-Fard al-ʿAzīz al-ʿAlīm.
Qul yā qawm a-tadaʿūna al-Ḥaqq ʿan warāʾikum wa-tadʿūna alladhī
khalaqnāhu bi-kaffin min al-ṭīn? Hādhā ẓulmun minkum ʿalā anfusikum in
antum fī āyāti rabbikum la-min al-mutafakkirīn. Qul yā qawm ṭahhirū
qulūbakum thumma abṣārakum laʿalla taʿrifūna bāriʾakum fī hādhā al-qamīṣ
al-muqaddas al-lamīʿ. Qul inna hādhā fatan ilāhiyyun qad istaqarra ʿalā
ʿarsh al-jalāl wa-ẓahara bi-sulṭān al-qudra wa-al-istiqlāl wa-yaṣīḥu
bayna al-arḍ wa-al-samāʾ bi-nidāʾihi al-abdaʿ al-aḥlā an yā ahl al-akwān
lima kafartum bi-rabbikum al-Raḥmān wa-aʿraḍtum ʿan jamāl al-Subḥān?
Ta-Allāh hādhā la-ghayb al-mastūr qad ṭalaʿa min mashriq al-imkān
wa-hādhā la-jamāl al-Maḥbūb qad ashraqa min ufuq hādhā al-riḍwān
bi-salṭanat Allāh al-Muhaymin al-ʿAzīz al-Ghālib al-Qadīr.

Arise for the Cause with power from Us and authority from Our presence,
then convey to the servants what the Spirit of God, the King, the
Unique, the Almighty, the All-Knowing, has conveyed to you. Say: O
people, do you leave the True One behind you and invoke the one We
created from a handful of clay? This is an injustice you do to
yourselves, if you reflect on the signs of your Lord. Say: O people,
purify your hearts and then your eyes, so that you may recognize your
Creator in this holy, radiant robe. Say: This is indeed a divine Youth
who has established Himself upon the throne of glory and appeared with
the sovereignty of power and independence. He calls out between earth
and heaven with His wondrous, sweetest call: ``O people of all realms of
being, why have you disbelieved in your Lord, the All-Merciful, and
turned away from the beauty of the All-Glorious? By God, this is the
concealed Unseen, risen from the dawning place of contingent being, and
this is the beauty of the Beloved, shining from the horizon of this
Paradise through the sovereignty of God, the All-Protecting, the
Almighty, the All-Conquering, the All-Powerful.''

\subsection{65}\label{section-1937}

\begin{otherlanguage}{arabic}

ان یا هیکل القدس انّا جعلنا صدرک ممرّداً من اشارات الممکنات و مقدّساً من
دلالات الکائنات لینطبع فیه انوار جمالی و تنعکس منه مرایا العالمین بذلک
اخترناک عمّا خلق فی السّموات و الأرض و اصطفیناک عمّا قدّر فی ملکوت الأمر و
الخلق و اختصصناک لنفسی هذا من فضل اللّه علیک من هذا الیوم الی الیوم الّذی
لن ینتهی فی الملک و یبقی ببقآء اللّه الملک المهیمن العزیز العلیم لأنّ یوم
اللّه هو نفسه اذاً ظهر بالحقّ و لن یعقّبه اللّیل و لن یحدّده الذّکر لو انتم من
العارفین

\end{otherlanguage}

An yā Haykal al-Quds innā jaʿalnā ṣadraka mumarradan min ishārāt
al-mumkināt wa-muqaddasan min dalālāt al-kāʾināt li-yanṭabiʿa fīhi
anwāru jamālī wa-tanʿakisa minhu marāyā al-ʿālamīn. Bi-dhālika
ikhtarnāka ʿammā khuliqa fī al-samāwāt wa-al-arḍ wa-iṣṭafaynāka ʿammā
quddira fī malakūt al-amr wa-al-khalq wa-ikhtaṣaṣnāka li-nafsī. Hādhā
min faḍl Allāh ʿalayka min hādhā al-yawm ilā al-yawm alladhī lan
yantahiya fī al-mulk wa-yabqā bi-baqāʾ Allāh al-Malik al-Muhaymin
al-ʿAzīz al-ʿAlīm. Li-anna yawma Allāh huwa nafsuhu idhan ẓahara
bi-al-ḥaqq wa-lan yuʿaqqibahu al-layl wa-lan yuḥaddidahu al-dhikr law
antum min al-ʿārifīn.

O Temple of Holiness! We have made your breast smooth and free of the
allusions of contingent things and sanctified it from the indications of
created beings, so that the lights of My beauty may be imprinted upon it
and the mirrors of the worlds may be reflected from it. Thus We have
chosen you above all that has been created in the heavens and on earth,
preferred you above all that has been ordained in the kingdom of
revelation and creation, and set you apart for Myself. This is God's
grace to you, from this day until the day that will never end in His
dominion and will endure as long as God, the King, the All-Protecting,
the Almighty, the All-Knowing, endures. For the Day of God is His own
Self when He appears in truth; night will never follow it, nor will any
description define it, if you are among those who understand.

\subsection{66}\label{section-1938}

\begin{otherlanguage}{arabic}

ان یا صدر هذا الهیکل انّا جعلنا الأشیآء مرایا نفسک و جعلناک مرآة نفسی
فأشرق علی صدور الممکنات بما تجلّی علیک من انوار ربّک لیطهّرها عن الحدود و
الاشارات کذلک اشرقت شمس الحکم من افق قلم مالک القدم طوبی للمتوسّمین انّا
بدأنا منک صدوراً ممرّدة و نعیدنّها الیک رحمةً من لدنّا علیک و علی المقرّبین
سوف نبعث بک صدوراً صافیة و ترائب منیرة لن یحکینّ الّا عن جمالی و لن یدلّنّ
الّا عن تجلّیات وجهی انّهم مرایا اسمائی بین الخلائق اجمعین

\end{otherlanguage}

An yā ṣadra hādhā al-haykal innā jaʿalnā al-ashyāʾ marāyā nafsika
wa-jaʿalnāka mirʾāta nafsī fa-ashriq ʿalā ṣudūr al-mumkināt bi-mā
tajallā ʿalayka min anwār rabbika li-yuṭahhirahā ʿan al-ḥudūd
wa-al-ishārāt. Kadhālika ashraqat shams al-ḥikam min ufuq qalam Mālik
al-Qidam. Ṭūbā li-al-mutawassimīn. Innā badaʾnā minka ṣudūran
mumarradatan wa-nuʿīdannahā ilayka raḥmatan min ladunnā ʿalayka wa-ʿalā
al-muqarrabīn. Sawfa nabʿathu bika ṣudūran ṣāfiyatan wa-tarāʾiba
munīratan lan yaḥkīna illā ʿan jamālī wa-lan yadululna illā ʿan
tajalliyāt wajhī. Innahum marāyā asmāʾī bayna al-khalāʾiq ajmaʿīn.

O Breast of this Temple! We have made all things mirrors of yourself and
made you the mirror of My own Self. Shine, then, upon the breasts of
contingent beings with the lights of your Lord that have shone upon you,
so that He may purify them of limitations and allusions. Thus has the
sun of wisdom shone from the horizon of the Pen of the Lord of Eternity.
Blessed are those who discern. From you We have brought forth polished
breasts, and to you We will surely return them, as a mercy from Us to
you and to those brought near. Through you We will raise up pure breasts
and radiant bosoms that will reflect nothing but My beauty and indicate
nothing but the revelations of My face. They are the mirrors of My names
among all created beings.

\subsection{67}\label{section-1939}

\begin{otherlanguage}{arabic}

ان یا هیکل القدس انّا قد جعلنا فؤادک مخزن علم ما کان و ما یکون و مطلع
علمنا الّذی قدّرناه لأهل السّموات و الأرض لیستفضنّ منک الموجودات و یبلغنّ
ببدائع علومک الی عرفان اللّه المقتدر العلیّ العظیم و انّ علمی الّذی ینسب الی
ذاتی ما عرفه احد و لا یعرفه نفس و لن یحمله احد من العالمین لو تظهر منه
کلمة لتضطرب النّفوس و تنعدم ارکان کلّ شیء و تزلّ اقدام البالغین

\end{otherlanguage}

An yā Haykal al-Quds innā qad jaʿalnā fuʾādaka makhzana ʿilm mā kāna
wa-mā yakūn wa-maṭlaʿa ʿilminā alladhī qaddarnāhu li-ahl al-samāwāt
wa-al-arḍ. La-yastafīḍanna minka al-mawjūdāt wa-yablughanna bi-badāʾiʿ
ʿulūmika ilā ʿirfān Allāh al-Muqtadir al-ʿAliyy al-ʿAẓīm. Wa-inna ʿilmī
alladhī yunsabu ilā dhātī mā ʿarafahu aḥad wa-lā yaʿrifuhu nafs wa-lan
yaḥmilahu aḥadun min al-ʿālamīn. Law taẓharu minhu kalimatun
la-taḍṭaribu al-nufūs wa-tanʿadimu arkānu kulli shayʾ wa-tazillu aqdāmu
al-bālighīn.

O Temple of Holiness! We have made your heart the treasury of the
knowledge of all that has been and all that will be, and the dawning
place of Our knowledge that We have ordained for the inhabitants of the
heavens and the earth. All beings will surely draw abundantly from you
and, through the wonders of your knowledge, attain recognition of God,
the All-Powerful, the Most High, the Most Great. But My knowledge that
pertains to My Essence has been known by no one; no soul knows it, and
no one in all the worlds can bear it. Were a single word of it to
appear, souls would be thrown into turmoil, the foundations of all
things would vanish, and the feet of those who have attained maturity
would slip.

\subsection{68}\label{section-1940}

\begin{otherlanguage}{arabic}

و عندنا علم لو نلقی علی الکائنات کلمةً منه لیوقننّ کلّ بظهور اللّه و علمه و
یطّلعنّ بأسرار العلوم کلّها و یبلغنّ الی المقام الّذی یرون انفسهم اغنیآء عن
علوم الأوّلین و الآخرین و لنا علوم اخری الّتی لا نقدر ان نذکر حرفاً منها و
لا النّاس یستطیعنّ ان یسمعنّ ذکراً منها کذلک نبّأناکم من علم اللّه العالم
الخبیر و لو نجد اوعیةً لألقیناها کنوز المعانی و علّمناها ما یحیط بحرف منه
العالمین

\end{otherlanguage}

Wa-ʿindanā ʿilmun law nulqī ʿalā al-kāʾināt kalimatan minhu la-yūqinanna
kullun bi-ẓuhūr Allāh wa-ʿilmihi wa-yaṭṭaliʿanna bi-asrār al-ʿulūm
kullihā wa-yablughanna ilā al-maqām alladhī yarawna anfusahum aghniyāʾ
ʿan ʿulūm al-awwalīn wa-al-ākhirīn. Wa-lanā ʿulūmun ukhrā allatī lā na
qdiru an nadhkura ḥarfan minhā wa-lā al-nāsu yastaṭīʿunna an yasmaʿna
dhikran minhā. Kadhālika nabbaʾnākum min ʿilm Allāh al-ʿĀlim al-Khabīr.
Wa-law najidu awʿiyatan la-alqaynāhā kunūza al-maʿānī wa-ʿallamnāhā mā
yuḥīṭu bi-ḥarfin minhu al-ʿālamīn.

We possess a knowledge of which, were We to impart a single word to
created beings, all would become certain of the manifestation of God and
His knowledge, discover the mysteries of all the sciences, and reach a
station at which they would see themselves as having no need of the
knowledge of the first and the last. We possess other branches of
knowledge of which We cannot mention a single letter, nor can people
bear to hear any mention of them. Thus have We informed you of the
knowledge of God, the Knowing, the All-Aware. Were We to find receptive
vessels, We would impart to them the treasures of inner meanings and
teach them that of which a single letter encompasses all the worlds.

\subsection{69}\label{section-1941}

\begin{otherlanguage}{arabic}

ان یا فؤاد هذا الهیکل انّا جعلناک مطلع علمی و مظهر حکمتی بین السّموات و
الأرضین و اظهرنا منک العلوم و نرجعها الیک ثمّ نبعث منک مرّة اخری وعداً من
لدنّا انّا کنّا فاعلین سوف نبعث منک ذا علوم بدیعة و ذا صنائع قویّة و نظهر
منهما ما لا خطر به قلب احد من العباد کذلک نعطی من نشآء ما نشآء و نأخذ
عمّن نشآء ما اعطیناه و نحکم بأمرنا ما نرید

\end{otherlanguage}

An yā fuʾāda hādhā al-haykal innā jaʿalnāka maṭlaʿa ʿilmī wa-maẓhara
ḥikmatī bayna al-samāwāt wa-al-araḍīn wa-aẓharnā minka al-ʿulūm
wa-nurjiʿuhā ilayka thumma nabʿathu minka marratan ukhrā waʿdan min
ladunnā innā kunnā fāʿilīn. Sawfa nabʿathu minka dhā ʿulūmin badīʿatin
wa-dhā ṣanāʾiʿa qawiyyatin wa-nuẓhiru minhumā mā lā khaṭara bihi qalbu
aḥadin min al-ʿibād. Kadhālika nuʿṭī man nashāʾu mā nashāʾu wa-naʾkhudhu
ʿamman nashāʾu mā aʿṭaynāhu wa-naḥkumu bi-amrinā mā nurīd.

O Heart of this Temple! We have made you the dawning place of My
knowledge and the manifestation of My wisdom between the heavens and the
earths. From you We have brought forth the branches of knowledge, and to
you We return them; then We will bring them forth from you once again.
This is a promise from Us; We will assuredly fulfill it. From you We
will raise up one endowed with wondrous sciences and another endowed
with powerful arts, and through them We will reveal what has never
entered the heart of any of Our servants. Thus We give whatever We will
to whomever We will, take back from whomever We will what We have given
him, and decree by Our command whatever We desire.

\subsection{70}\label{section-1942}

\begin{otherlanguage}{arabic}

قل انّا لو نتجلّی علی مرایا الموجودات بشمس عنایتنا فی ساعة و نأخذ عنهم
انوار تجلّیاتنا فی ساعة اخری لنقدر و لیس لأحد ان یقول لم او بم لأنّا نحن
الفاعل لما نشآء و لا نسأل عمّا فعلناه و لا یشکّ فی ذلک الّا کلّ مشرک مریب قل
لن تمنع قدرتنا و لن یعطّل حکمنا نرفع من نشآء الی جبروت العزّة و الاقتدار
ثمّ نرجعه لو نشآء الی اسفل السّافلین

\end{otherlanguage}

Qul innā law natajallā ʿalā marāyā al-mawjūdāt bi-shams ʿināyatinā fī
sāʿatin wa-naʾkhudhu ʿanhum anwāra tajalliyātinā fī sāʿatin ukhrā
la-naqdiru wa-laysa li-aḥadin an yaqūla lima aw bi-ma li-annā naḥnu
al-fāʿilu li-mā nashāʾu wa-lā nusʾalu ʿammā faʿalnāhu wa-lā yashukku fī
dhālika illā kullu mushrikin murīb. Qul lan tumnaʿa qudratunā wa-lan
yuʿaṭṭala ḥukmunā narfaʿu man nashāʾu ilā jabarūt al-ʿizza wa-al-iqtidār
thumma nurjiʿuhu law nashāʾu ilā asfal al-sāfilīn.

Say: Were We to shine upon the mirrors of all beings with the sun of Our
loving care at one moment and withdraw from them the lights of Our
revelations at another, We would have the power to do so. No one has the
right to say ``Why?'' or ``By what means?'' For We do whatever We will,
and We are not to be questioned about what We have done. None doubts
this except those who associate partners with God and are steeped in
doubt. Say: Our power cannot be restrained, nor Our decree suspended. We
raise whomever We will to the dominion of glory and might, then, if We
will, return him to the lowest of the low.

\subsection{71}\label{section-1943}

\begin{otherlanguage}{arabic}

أ تزعمون یا ملأ الأرض بأنّا لو نصعد احداً الی السّدرة المنتهی اذاً تعزل عنه
قدرتی و سلطانی لا ونفسی بل لو نشآء لنرجعنّه الی التّراب فی اقلّ من حین
فانظروا فی الشّجرة انّا نغرسها فی الجنان و نسقیها من مآء عنایتنا فلمّا
ارتفعت فی نفسها و تورّقت بالأوراق الخضرآء و اثمرّت بالأثمار الحسنی اذاً
نرسل علیها قواصف الأمر و ندعها علی وجه الأرض کذلک کنّا فاعلین و کذلک نفعل
بکلّ شیء هذا من بدائع سنننا من قبل و من بعد فی کلّ الأشیآء ان انتم من
النّاظرین و لا یعلم حکمة ذلک الّا اللّه المقتدر العزیز الحکیم

\end{otherlanguage}

A-tazʿumūna yā malaʾ al-arḍ bi-annā law nuṣʿidu aḥadan ilā al-sidra
al-muntahā idhan tuʿzalu ʿanhu qudratī wa-sulṭānī? Lā wa-nafsī bal law
nashāʾu la-nurjiʿannahu ilā al-turāb fī aqalla min ḥīn. Fa-unẓurū fī
al-shajara innā naghrisuhā fī al-jinān wa-nasqīhā min māʾ ʿināyatinā
fa-lammā irtafaʿat fī nafsihā wa-tawarraqat bi-al-awrāq al-khaḍrāʾ
wa-athmarrat bi-al-athmār al-ḥusnā idhan nursilu ʿalayhā qawāṣifa al-amr
wa-nadaʿuhā ʿalā wajh al-arḍ. Kadhālika kunnā fāʿilīn wa-kadhālika
nafʿalu bi-kulli shayʾ. Hādhā min badāʾiʿ sunaninā min qablu wa-min
baʿdu fī kulli al-ashyāʾ in antum min al-nāẓirīn. Wa-lā yaʿlamu ḥikmata
dhālika illā Allāh al-Muqtadir al-ʿAzīz al-Ḥakīm.

Do you imagine, O people of the earth, that if We raise someone to the
Lote Tree beyond which there is no passing, My power and sovereignty
over him are then removed? No, by My own Self! Rather, should We will,
We would return him to dust in less than an instant. Consider the tree:
We plant it in the gardens and water it with the water of Our loving
care. Then, when it has grown tall, put forth green leaves, and borne
the finest fruits, We send against it the shattering winds of Our
command and leave it lying upon the ground. Thus have We acted, and thus
do We deal with all things. This is among Our wondrous ways with all
things, before and hereafter, if you are among those who see. No one
knows the wisdom of this except God, the All-Powerful, the Almighty, the
All-Wise.

\subsection{72}\label{section-1944}

\begin{otherlanguage}{arabic}

أ تنکرون یا قوم ما ترونه ویل لکم یا ملأ المنکرین و الّذی لن یتغیّر هو نفسه
الرّحمن الرّحیم ان انتم من المتبصّرین و دونه یتغیّر بارادة من عنده و هو
المقتدر العزیز الحکیم

\end{otherlanguage}

A-tunkirūna yā qawm mā tarawnahu? Waylun lakum yā malaʾ al-munkirīn.
Wa-alladhī lan yataghayyara huwa nafsuhu al-Raḥmān al-Raḥīm in antum min
al-mutabaṣṣirīn. Wa-dūnahu yataghayyaru bi-irādatin min ʿindihi wa-huwa
al-Muqtadir al-ʿAzīz al-Ḥakīm.

Do you deny, O people, what you see? Woe to you, O company of deniers!
The One who will never change is His own Self, the All-Merciful, the
Merciful, if you are among those endowed with insight. All besides Him
changes through His will, and He is the All-Powerful, the Almighty, the
All-Wise.

\subsection{73}\label{section-1945}

\begin{otherlanguage}{arabic}

یا قوم لا تتکلّموا فی امری لأنّکم لا تبلغون الی حکمة ربّکم و لن تنالوا
بعلمه العزیز المحیط و من ادّعی عرفان ذاته هو من اجهل النّاس یکذّبه کلّ
الذّرّات و یشهد بهذا لسانی الصّادق الأمین ان اذکروا امری ثمّ تکلّموا فیه و
فیما امرتم به من لدنّا و من دون ذلک لا ینبغی لکم و لیس لأحد الیه سبیل ان
انتم من السّامعین

\end{otherlanguage}

Yā qawm lā tatakallamū fī amrī li-annakum lā tab l ughūna ilā ḥikmat
rabbikum wa-lan tanālū bi-ʿilmihi al-ʿAzīz al-Muḥīṭ. Wa-man iddaʿā
ʿirfāna dhātihi huwa min ajhal al-nās yukadhdhibuhu kullu al-dharrāt
wa-yashhadu bi-hādhā lisānī al-ṣādiq al-amīn. An udhkurū amrī thumma
takallamū fīhi wa-fīmā umirtum bihi min ladunnā wa-min dūni dhālika lā
yanbaghī lakum wa-laysa li-aḥadin ilayhi sabīl in antum min al-sāmiʿīn.

O people, do not speak concerning Me, for you cannot attain the wisdom
of your Lord, nor reach His mighty, all-encompassing knowledge. Whoever
claims to know His Essence is among the most ignorant of people. Every
atom gives the lie to him, and My truthful, trustworthy tongue bears
witness to this. Make mention of My Cause, then speak of it and of what
you have been commanded by Us. Anything beyond this does not befit you,
nor does anyone have a path to it, if you are among those who listen.

\subsection{74}\label{section-1946}

\begin{otherlanguage}{arabic}

ان یا هذا الهیکل قد جعلناک مطلع کلّ اسم من اسمائنا الحسنی و مظهر کلّ صفة
من صفاتنا العلیا و منبع کلّ ذکر من اذکارنا لمن فی الأرض و السّمآء ثمّ
بعثناک علی صورتی بین السّموات و الأرض و جعلناک آیة عزّی لمن فی جبروت الأمر
و الخلق لیهتدینّ بک عبادی و یکوننّ من المهتدین و جعلناک سدرة الجود لمن فی
السّموات و الأرض هنیئاً لمن یستظلّ فی ظلّک و یتقرّب الی نفسک المهیمنة علی
العالمین

\end{otherlanguage}

An yā hādhā al-haykal qad jaʿalnāka maṭlaʿa kulli ismin min asmāʾinā
al-ḥusnā wa-maẓhara kulli ṣifatin min ṣifātinā al-ʿulyā wa-manbaʿa kulli
dhikrin min adhkārinā li-man fī al-arḍ wa-al-samāʾ. Thumma baʿathnāka
ʿalā ṣūratī bayna al-samāwāt wa-al-arḍ wa-jaʿalnāka āyata ʿizzī li-man
fī jabarūt al-amr wa-al-khalq li-yahtadiyanna bika ʿibādī wa-yakūnunna
min al-muhtadīn. Wa-jaʿalnāka sidrata al-jūd li-man fī al-samāwāt
wa-al-arḍ. Hanīʾan li-man yastaẓillu fī ẓillika wa-yataqarrabu ilā
nafsika al-muhaymina ʿalā al-ʿālamīn.

O Temple! We have made you the dawning place of each of Our most
beautiful names, the manifestation of each of Our most exalted
attributes, and the source of every mention of Us for those on earth and
in heaven. Then We raised you up in My image between the heavens and the
earth and made you the sign of My glory for those in the dominion of
revelation and creation, so that My servants may be guided by you and be
among the rightly guided. We have made you the Lote Tree of bounty for
those in the heavens and on earth. Blessed is the one who takes shelter
in your shade and draws near to your Self, which keeps watch over all
the worlds.

\subsection{75}\label{section-1947}

\begin{otherlanguage}{arabic}

قل انّا جعلنا کلّ اسم معیناً و اجرینا منها انهار الحکمة و العرفان فی ریاض
الأمر و لا یعلم عدّتها احد الّا ربّک المقدّس المقتدر العلیم الحکیم قل انّا
بدأنا کلّ الحروف من النّقطة و ارجعناها الیها ثمّ بعثناها علی هیکل بشر تعالی
الصّانع الأبدع البدیع سوف نفصّل منها مرّةً اخری باسمی الأبهی فضلاً من عندی و
انا الفضّال القدیم

\end{otherlanguage}

Qul innā jaʿalnā kulla ismin maʿīnan wa-ajraynā minhā anhāra al-ḥikma
wa-al-ʿirfān fī riyāḍ al-amr wa-lā yaʿlamu ʿiddatahā aḥadun illā rabbuka
al-Muqaddas al-Muqtadir al-ʿAlīm al-Ḥakīm. Qul innā badaʾnā kulla
al-ḥurūf min al-nuqṭa wa-arjaʿnāhā ilayhā thumma baʿathnāhā ʿalā haykal
bashar. Taʿālā al-Ṣāniʿ al-Abdaʿ al-Badīʿ. Sawfa nufaṣṣilu minhā
marratan ukhrā bi-ismī al-Abhā faḍlan min ʿindī wa-anā al-Faḍḍāl
al-Qadīm.

Say: We have made every name a wellspring, and from them We have caused
rivers of wisdom and recognition to flow through the gardens of the
Cause. No one knows their number except your Lord, the Most Holy, the
All-Powerful, the All-Knowing, the All-Wise. Say: We originated all the
letters from the Point and returned them to it, then raised them up in a
human temple. Exalted is the Maker, the Most Wondrous, the Incomparable!
Once again We will unfold them through My name, the All-Glorious, as a
grace from Me. I am the Most Bountiful, the Ancient.

\subsection{76}\label{section-1948}

\begin{otherlanguage}{arabic}

و اشرقنا الأنوار من شمس اسمنا الحقّ و ارجعناها الیها و اظهرناها علی هیکل
الانسان تعالی القادر المقتدر القدیر لن یمنعنی احد عن امری و لن تحجبنی
نفس عن سلطانی و قدرتی انا الّذی بعثت الممکنات بقولی و انا المقتدر علی ما
ارید

\end{otherlanguage}

Wa-ashraqnā al-anwār min shams isminā al-Ḥaqq wa-arjaʿnāhā ilayhā
wa-aẓharnāhā ʿalā haykal al-insān. Taʿālā al-Qādir al-Muqtadir al-Qadīr.
Lan yamnaʿanī aḥadun ʿan amrī wa-lan taḥjubanī nafsun ʿan sulṭānī
wa-qudratī. Anā alladhī baʿathtu al-mumkināt bi-qawlī wa-anā al-Muqtadir
ʿalā mā urīd.

We caused the lights to shine from the sun of Our name, the True One,
returned them to it, and manifested them in the human temple. Exalted is
the Powerful, the All-Compelling, the All-Powerful! No one can prevent
Me from accomplishing My Cause, nor can any soul bar Me from My
sovereignty and power. I am the One who brought contingent beings into
existence through My word, and I have power to do whatever I desire.

\subsection{77}\label{section-1949}

\begin{otherlanguage}{arabic}

قل انّا لو نرید ان نقبض الأرواح من کلّ الأشیآء فی نفس و نبعث منها مرّةً اخری
لنقدر لا یعرف علم ذلک الّا اللّه العالم العلیم و لو نرید ان نظهر من ذرّة
شموساً لا لهنّ بدایة و لا نهایة لنقدر و نظهر کلّهنّ بأمری فی اقلّ من حین و لو
نرید ان نبعث من قطرة بحور السّموات و الأرض و نفصّل من حرف علم ما کان و ما
یکون لنقدر انّ هذا لسهل یسیر کذلک کنت مقتدراً من الأوّل الّذی لا اوّل له الی
الآخر الّذی لا آخر له ولکن خلقی غفلوا عن قدرتی و اعرضوا عن سلطانی و
جادلوا بنفسی العلیم الحکیم

\end{otherlanguage}

Qul innā law nurīdu an naqbiḍa al-arwāḥ min kulli al-ashyāʾ fī nafasin
wa-nabʿatha minhā marratan ukhrā la-naqdiru. Lā yaʿrifu ʿilma dhālika
illā Allāh al-ʿĀlim al-ʿAlīm. Wa-law nurīdu an nuẓhira min dharratin
shumūsan lā lahunna bidāyatun wa-lā nihāyatun la-naqdiru wa-nuẓhiru
kullahunna bi-amrī fī aqalla min ḥīn. Wa-law nurīdu an nabʿatha min
qaṭratin buḥūra al-samāwāt wa-al-arḍ wa-nufaṣṣila min ḥarfin ʿilma mā
kāna wa-mā yakūn la-naqdiru. Inna hādhā la-sahlun yasīr. Kadhālika kuntu
muqtadiran min al-awwal alladhī lā awwala lahu ilā al-ākhir alladhī lā
ākhira lahu wa-lākin khalqī ghafalū ʿan qudratī wa-aʿraḍū ʿan sulṭānī
wa-jādalū bi-nafsī al-ʿAlīm al-Ḥakīm.

Say: Were We to desire to withdraw the spirits from all things in a
single breath and bring them forth from them once again, We would have
the power to do so. No one possesses knowledge of this except God, the
Knowing, the All-Knowing. Were We to desire to bring forth from an atom
suns that have neither beginning nor end, We would have the power to do
so, and would manifest them all at My command in less than an instant.
Were We to desire to bring forth from a drop the oceans of the heavens
and the earth, and unfold from a single letter the knowledge of all that
has been and all that will be, We would have the power to do so. This is
indeed easy and simple. Thus have I been powerful from the beginning
that has no beginning to the end that has no end. Yet My creatures have
remained heedless of My power, turned away from My sovereignty, and
disputed with My own Self, the All-Knowing, the All-Wise.

\subsection{78}\label{section-1950}

\begin{otherlanguage}{arabic}

قل لم یتحرّک شیء بین السّموات و الأرض الّا بعد اذنی و لم تصعد نفس الی
الملکوت الّا بعد امری ولکن بریّتی احتجبوا عن قدرتی و سلطانی و کانوا من
الغافلین قل لا یری فی ظهوری الّا ظهور اللّه و لا فی قدرتی الّا قدرة اللّه لو
انتم من العارفین قل مثل خلقی کمثل الأوراق علی الشّجرة انّها قد کانت ظاهرة
بوجودها و قائمة بنفسها ولکن غافلة عن اصلها کذلک مثّلنا لعبادنا العاقلین
لعلّ یصعدنّ عن رتبة النّبات و یبلغنّ الی مقام البلوغ فی هذا الأمر المبرم
المتین قل انّ مثلهم کمثل الحوت فی المآء انّ حیوته به و انّه لم یعرف ممدّ
حیوته من لدن عزیز حکیم و کان محتجباً عنه بحیث لو یسأل عن المآء و صفاته لن
یعرف کذلک نلقی الأمثال لعلّ النّاس یقبلنّ الی قبلة من فی السّموات و الأرضین

\end{otherlanguage}

Qul lam yataḥarrak shayʾun bayna al-samāwāt wa-al-arḍ illā baʿda idhnī
wa-lam taṣʿad nafsun ilā al-malakūt illā baʿda amrī wa-lākin bariyyatī
iḥtajabū ʿan qudratī wa-sulṭānī wa-kānū min al-ghāfilīn. Qul lā yurā fī
ẓuhūrī illā ẓuhūru Allāh wa-lā fī qudratī illā qudratu Allāh law antum
min al-ʿārifīn. Qul mathalu khalqī ka-mathal al-awrāq ʿalā al-shajara
innahā qad kānat ẓāhiratan bi-wujūdihā wa-qāʾimatan bi-nafsihā wa-lākin
ghāfilatan ʿan aṣlihā. Kadhālika maththalnā li-ʿibādinā al-ʿāqilīn
laʿalla yaṣʿadunna ʿan rutbat al-nabāt wa-yablughunna ilā maqām
al-bulūgh fī hādhā al-amr al-mubram al-matīn. Qul inna mathalahum
ka-mathal al-ḥūt fī al-māʾ inna ḥayātahu bihi wa-innahu lam yaʿrif
mumidda ḥayātihi min ladun ʿAzīzin Ḥakīm wa-kāna muḥtajiban ʿanhu
bi-ḥaythu law yusʾalu ʿan al-māʾ wa-ṣifātihi lan yaʿrif. Kadhālika nulqī
al-amthāl laʿalla al-nāsa yuqbilunna ilā qiblat man fī al-samāwāt
wa-al-araḍīn.

Say: Nothing moves between the heavens and the earth except with My
permission, and no soul ascends to the kingdom except at My command. Yet
My creatures have been veiled from My power and sovereignty and have
remained among the heedless. Say: Nothing is seen in My manifestation
but the manifestation of God, and nothing in My power but the power of
God, if you are among those who understand. Say: My creatures are like
the leaves upon a tree: they are manifest through its existence and
sustained by its very being, yet remain heedless of their root. Thus
have We set forth a likeness for Our servants who understand, so that
they may rise above the rank of plants and attain the station of
maturity in this firmly established, steadfast Cause. Say: They are like
a fish in the water: its life depends upon it, yet it does not recognize
what sustains its life from the Almighty, the All-Wise. It is so veiled
from it that, if asked about the water and its qualities, it would not
know. Thus do We set forth likenesses, so that people may turn toward
the Point of Adoration of all who are in the heavens and the earths.

\subsection{79}\label{section-1951}

\begin{otherlanguage}{arabic}

یا قوم خافوا من اللّه و لا تکفروا بالّذی احاطت رحمته الممکنات و سبق فضله
الموجودات و احاط سلطان امره ظاهرکم و باطنکم و اوّلکم و آخرکم اتّقوا اللّه و
کونوا من المتّقین ایّاکم ان تکونوا مثل الّذین تمرّ علیهم آیات اللّه و هم لا
یعرفونها الا انّهم من الغافلین

\end{otherlanguage}

Yā qawm khāfū min Allāh wa-lā takfurū bi-lladhī aḥāṭat raḥmatuhu
al-mumkināt wa-sabaqa faḍluhu al-mawjūdāt wa-aḥāṭa sulṭān amrihi
ẓāhirakum wa-bāṭinakum wa-awwalakum wa-ākhirakum. Ittaqū Allāh wa-kūnū
min al-muttaqīn. Iyyākum an takūnū mithl alladhīna tamurru ʿalayhim āyāt
Allāh wa-hum lā yaʿrifūnahā. Alā innahum min al-ghāfilīn.

People, fear God and do not disbelieve in Him whose mercy has
encompassed all contingent things, whose grace has preceded all created
things, and the sovereignty of whose Cause has embraced your outward and
inward being, your first and your last. Be mindful of God and be among
the God-fearing. Beware of being like those who fail to recognize the
verses of God as they pass before them. They are indeed among the
heedless.

\subsection{80}\label{section-1952}

\begin{otherlanguage}{persian}

قل أ تعبدون من لا یسمع و لا یبصر و کان احقر العباد و اضلّهم فما لکم لا
تتّبعون الّذی اتی من مطلع الأمر بنبإ اللّه العلیّ العظیم یا قوم لا تکونوا
کالّذین دخلوا تلقآء العرش و ما استشعروا الا انّهم من الصّاغرین کنّا نتلو
علیهم الآیات الّتی انجذب منها اهل الجبروت و سکّان الملکوت و هم رجعوا
محتجبین عنها و مترصّدین ندآء احد من العباد الّذی حیّ بارادة من عند اللّه
کذلک نلقی علیکم ما یهدیکم سبیل المقرّبین

\end{otherlanguage}

Qul a-taʿbudūna man lā yasmaʿ wa-lā yubṣir wa-kāna aḥqar al-ʿibād
wa-aḍallahum fa-mā lakum lā tattabiʿūna alladhī atā min maṭlaʿ al-amr
bi-nabaʾ Allāh al-ʿAliyy al-ʿAẓīm. Yā qawm lā takūnū ka-lladhīna dakhalū
tilqāʾ al-ʿarsh wa-mā istashʿarū. Alā innahum min al-ṣāghirīn. Kunnā
natlū ʿalayhim al-āyāt allatī injadhaba minhā ahl al-jabarūt wa-sukkān
al-malakūt wa-hum rajaʿū muḥtajibīn ʿanhā wa-mutaraṣṣidīn nidāʾ aḥad min
al-ʿibād alladhī ḥayya bi-irāda min ʿind Allāh. Kadhālika nulqī ʿalaykum
mā yahdīkum sabīl al-muqarrabīn.

Say: Do you worship one who neither hears nor sees and who is the most
contemptible and most astray of servants? Why, then, do you not follow
Him who has come from the dawning-place of the Cause with tidings of
God, the Most High, the Most Great? People, do not be like those who
entered the presence of the Throne without perceiving it. They are
indeed among the abased. We were reciting to them the verses that
enraptured the people of the realm of divine power and the inhabitants
of the Kingdom, yet they returned veiled from them, awaiting the call of
one of the servants whose life depended on a will proceeding from God.
Thus We impart to you what will guide you to the path of those brought
near.

\subsection{81}\label{section-1953}

\begin{otherlanguage}{persian}

کم من عباد دخلوا بقعة الفردوس مقرّ العرش بین یدی ربّهم العلیّ العظیم و
سألوا عن ابواب اربعة او عن احد من ائمّة الفرقان کذلک کان شأن هؤلآء ان
انتم من العالمین کما ترون فی تلک الأیّام الّذینهم کفروا و اشرکوا تمسّکوا
باسم من الأسمآء و احتجبوا عن موجدها نشهد انّهم من اهل السّعیر یسألون الشّمس
ما قاله الظّلّ و عن الحقّ ما نطق الخلق ان انتم من الشّاهدین قل یا قوم لم یکن
عند الشّمس الّا اشراقها و ما یظهر منها و ما سواها استضآء بنورها اتّقوا اللّه
و لا تکوننّ من الجاهلین منهم من سأل الظّلمة عن النّور قل ان افتح بصرک لتری
الاشراق الّذی احاط الآفاق انّه یری بالعین هذا نور اشرق و لاح من افق فجر
المعانی بضیآء مبین أ تسألون الیهود هل کان الرّوح علی حقّ من اللّه او
الأصنام هل کان محمّد رسولاً او ملأ الفرقان ذکر اللّه العلیّ العظیم

\end{otherlanguage}

Kam min ʿibād dakhalū buqʿat al-firdaws maqarr al-ʿarsh bayna yaday
rabbihim al-ʿAliyy al-ʿAẓīm wa-saʾalū ʿan abwāb arbaʿa aw ʿan aḥad min
aʾimmat al-Furqān. Kadhālika kāna shaʾn hāʾulāʾ in antum min al-ʿālimīn.
Kamā tarawna fī tilka al-ayyām alladhīnahum kafarū wa-ashrakū tamassakū
bi-ism min al-asmāʾ wa-iḥtajabū ʿan mūjidihā. Nashhadu annahum min ahl
al-saʿīr. Yasʾalūna al-shams mā qālahu al-ẓill wa-ʿan al-Ḥaqq mā naṭaqa
al-khalq in antum min al-shāhidīn. Qul yā qawm lam yakun ʿinda al-shams
illā ishrāquhā wa-mā yaẓharu minhā wa-mā siwāhā istaḍāʾa bi-nūrihā.
Ittaqū Allāh wa-lā takūnanna min al-jāhilīn. Minhum man saʾala al-ẓulma
ʿan al-nūr. Qul in iftaḥ baṣarak li-tarā al-ishrāq alladhī aḥāṭa al-āfāq
innahu yurā bi-l-ʿayn. Hādhā nūr ashraqa wa-lāḥa min ufuq fajr al-maʿānī
bi-ḍiyāʾ mubīn. A-tasʾalūna al-Yahūd hal kāna al-Rūḥ ʿalā ḥaqq min Allāh
aw al-aṣnām hal kāna Muḥammad rasūlan aw malaʾ al-Furqān dhikr Allāh
al-ʿAliyy al-ʿAẓīm.

How many servants entered the precinct of Paradise, the seat of the
Throne, into the presence of their Lord, the Most High, the Most Great,
and asked about the four Gates or one of the Imams of the Qur'an! Such
was their conduct, if you are among those who know. Likewise, you see in
these days that those who have disbelieved and associated partners with
God have clung to one of the names and remained veiled from the One who
brought them into being. We testify that they are among the people of
the blazing Fire. They ask the Sun what the shadow has said, and the
Truth what creation has uttered, if you are among those who witness.
Say: People, the Sun has nothing but its radiance and what appears from
it; everything else has been illumined by its light. Be mindful of God
and do not be among the ignorant. Among them are those who have asked
darkness about the light. Say: Open your eyes so that you may see the
radiance that has encompassed the horizons. It can indeed be seen with
the eye. This is a light that has risen and shone from the horizon of
the dawn of inner meanings with manifest splendor. Do you ask the Jews
whether the Spirit was truly from God, or the idols whether Muhammad was
a Messenger, or the people of the Qur'an about the Remembrance of God,
the Most High, the Most Great?

\subsection{82}\label{section-1954}

\begin{otherlanguage}{persian}

قل یا قوم دعوا ما عندکم عند تجلّی هذا الظّهور خذوا ما امرتم به هذا امر
اللّه لکم انّه هو خیر الآمرین وجمالی لم یکن مقصودی فی تلک الکلمات نفسی بل
الّذی یأتی بعدی و الله علی ذلک لشهید و علیم لا تفعلوا به ما فعلتم بنفسی
اذا نزّلت علیکم آیات اللّه من شطر فضلی لا تقولوا انّها ما نزّلت علی الفطرة
انّ الفطرة قد خلقت بقولی و تطوف فی حولی ان انتم من الموقنین ان استنشقوا
نفحات قمیص المعانی من بیان ربّکم الرّحمن انّها تضوّعت فی الأکوان و تعطّر بها
الامکان طوبی لمن وجد عرفها و اقبل الی اللّه بقلب منیر

\end{otherlanguage}

Qul yā qawm daʿū mā ʿindakum ʿinda tajallī hādhā al-ẓuhūr khudhū mā
umirtum bihi. Hādhā amr Allāh lakum innahu huwa khayr al-āmirīn.
Wa-jamālī lam yakun maqṣūdī fī tilka al-kalimāt nafsī bal alladhī yaʾtī
baʿdī wa-Allāh ʿalā dhālika la-shahīd wa-ʿalīm. Lā tafʿalū bihi mā
faʿaltum bi-nafsī. Idhā nuzzilat ʿalaykum āyāt Allāh min shaṭr faḍlī lā
taqūlū innahā mā nuzzilat ʿalā al-fiṭra. Inna al-fiṭra qad khuliqat
bi-qawlī wa-taṭūfu fī ḥawlī in antum min al-mūqinīn. An istanshiqū
nafaḥāt qamīṣ al-maʿānī min bayān rabbikum al-Raḥmān innahā taḍawwaʿat
fī al-akwān wa-taʿaṭṭara bihā al-imkān. Ṭūbā li-man wajada ʿarfahā
wa-aqbala ilā Allāh bi-qalb munīr.

Say: People, set aside what you possess when this Revelation shines
forth, and take hold of what you have been commanded. This is God's
command to you; He is indeed the best of those who command. By My
beauty! In those words I did not mean Myself, but the One who will come
after Me, and God is indeed a knowing witness to this. Do not treat Him
as you have treated Me. When the verses of God are sent down to you from
the direction of My grace, do not say, ``They have not been revealed
spontaneously.'' Innate nature itself was created by My word and circles
around Me, if you are among those who are certain. Inhale from the
utterance of your Lord, the All-Merciful, the fragrances of the garment
of inner meanings. They have spread through all worlds, and the realm of
existence has been perfumed by them. Blessed is the one who has
perceived their fragrance and turned to God with a radiant heart.

\subsection{83}\label{section-1955}

\begin{otherlanguage}{arabic}

ان یا هذا الهیکل انّا قد جعلناک مرآةً لملکوت الأسمآء لتحکی عن سلطانی بین
الخلائق اجمعین و تدعو النّاس الی لقائی و جمالی و تکون هادیاً الی سبیلی
الواضح المستقیم قد رفعنا اسمک بین العباد فضلاً من عندنا و انا الفضّال
القدیم و زیّنّاک بطراز نفسی و القینا علیک کلمتی لتحکم فی الملک کیف تشآء و
تفعل ما ترید و قدّرنا لک خیر السّموات و الأرض بحیث لم یکن لأحد من خیر الّا
بأن یدخل فی ظلّک امراً من لدن ربّک العلیم الخبیر و اعطیناک عصا الأمر و
فرقان الحکم لتفرق بین کلّ امر حکیم و موّجنا فی صدرک ابحر المعانی و البیان
فی ذکر ربّک الرّحمن لتشکر ربّک و تکون من الشّاکرین و اختصصناک بین خلقی و
جعلناک مظهر نفسی بین السّموات و الأرضین

\end{otherlanguage}

In yā hādhā al-haykal innā qad jaʿalnāka mirʾātan li-malakūt al-asmāʾ
li-taḥkiya ʿan sulṭānī bayna al-khalāʾiq ajmaʿīn wa-tadʿuwa al-nās ilā
liqāʾī wa-jamālī wa-takūna hādiyan ilā sabīlī al-wāḍiḥ al-mustaqīm. Qad
rafaʿnā ismaka bayna al-ʿibād faḍlan min ʿindinā wa-anā al-Faḍḍāl
al-Qadīm. Wa-zayyannāka bi-ṭirāz nafsī wa-alqaynā ʿalayka kalimatī
li-taḥkuma fī al-mulk kayfa tashāʾ wa-tafʿala mā turīd. Wa-qaddarnā laka
khayr al-samāwāt wa-l-arḍ bi-ḥaythu lam yakun li-aḥad min khayr illā
bi-an yadkhula fī ẓillik amran min ladun rabbika al-ʿAlīm al-Khabīr.
Wa-aʿṭaynāka ʿaṣā al-amr wa-furqān al-ḥukm li-tafruqa bayna kull amr
ḥakīm. Wa-mawwajnā fī ṣadrika abḥur al-maʿānī wa-l-bayān fī dhikr
rabbika al-Raḥmān li-tashkura rabbaka wa-takūna min al-shākirīn.
Wa-ikhtaṣaṣnāka bayna khalqī wa-jaʿalnāka maẓhar nafsī bayna al-samāwāt
wa-l-araḍīn.

O Temple! We have made You a mirror of the kingdom of names, so that You
may reflect My sovereignty among all created beings, summon people to My
presence and My beauty, and guide them to My clear and straight path. We
have exalted Your name among the servants as a grace from Us, and I am
the Ever-Bountiful, the Ancient. We have adorned You with the ornament
of My own Self and imparted My word to You, so that You may rule
throughout creation as You wish and do whatever You desire. We have
ordained for You the good of the heavens and the earth, so that no one
may have any share of good unless they enter Your shadow---a decree from
Your Lord, the All-Knowing, the All-Aware. We have given You the rod of
authority and the criterion of judgment, so that You may distinguish
every matter of wisdom. We have caused the oceans of inner meanings and
utterance to surge within Your breast in remembrance of Your Lord, the
All-Merciful, so that You may give thanks to Your Lord and be among the
thankful. We have singled You out among My creatures and made You the
Manifestation of My own Self throughout the heavens and the earths.

\subsection{84}\label{section-1956}

\begin{otherlanguage}{arabic}

ان ابتعث باذن من لدنّا مرایا مستحکیات و حروفات عالیات لیحکین عن سلطانک و
قدرتک و یدللن عن اقتدارک و عظمتک و یکنّ مظاهر اسمائک بین العالمین انّا
جعلناک مبدأ المرایا و مبدعهنّ کما بدأناهنّ منک اوّل مرّة و نعیدک الی نفسی
کما بدأناک انّ ربّک لهو الغالب المقتدر القدیر نبّئ المرایا حین ظهورهنّ بأن
لا یستکبرن علی موجدهنّ و خالقهنّ حین ظهوره و لا تغرّنّهنّ الرّیاسة عن الخضوع
بین یدی اللّه العزیز الجمیل

\end{otherlanguage}

An ibʿath bi-idhn min ladunnā marāyā mustaḥkiyāt wa-ḥurūfāt ʿāliyāt
li-yaḥkīna ʿan sulṭānika wa-qudratika wa-yadlulna ʿan iqtidārika
wa-ʿaẓamatika wa-yakunna maẓāhir asmāʾika bayna al-ʿālamīn. Innā
jaʿalnāka mabdaʾ al-marāyā wa-mubdiʿahunna kamā badaʾnāhunna minka awwal
marra wa-nuʿīduka ilā nafsī kamā badaʾnāk. Inna rabbaka la-huwa
al-Ghālib al-Muqtadir al-Qadīr. Nabbiʾ al-marāyā ḥīna ẓuhūrihinna bi-an
lā yastakbirna ʿalā mūjidihinna wa-khāliqihinna ḥīna ẓuhūrihi wa-lā
taghurrannahunna al-riʾāsa ʿan al-khuḍūʿ bayna yaday Allāh al-ʿAzīz
al-Jamīl.

Raise up, by Our permission, reflecting mirrors and exalted letters, so
that they may reflect Your sovereignty and power, bear witness to Your
might and grandeur, and be manifestations of Your names among all
peoples. We have made You the source of the mirrors and their
originator, just as We brought them forth from You the first time; and
We return You to My own Self, just as We brought You forth. Your Lord is
indeed the All-Conquering, the All-Powerful, the Omnipotent. Tell the
mirrors, when they appear, not to exalt themselves over the One who
brought them into being and created them when He appears, and not to let
leadership beguile them into failing to humble themselves before God,
the Almighty, the All-Beautiful.

\subsection{85}\label{section-1957}

\begin{otherlanguage}{arabic}

قل انتنّ یا ایّتها المرایا قد خلقتنّ بأمری و بعثتنّ بارادتی ایّاکنّ ان تکفرن
بآیات ربّی و تکنّ من الّذینهم ظلموا و کانوا من الخاسرین و تتمسّکن بما عندکنّ
و تفتخرن بارتفاع اسمائکنّ ینبغی لکنّ بأن تنقطعن عمّن فی السّموات و الأرض
کذلک قدّر من لدن مقتدر قدیر

\end{otherlanguage}

Qul antunna yā ayyatuhā al-marāyā qad khuliqtunna bi-amrī wa-buʿithtunna
bi-irādatī. Iyyākunna an takfurna bi-āyāt rabbī wa-takunna min
alladhīnahum ẓalamū wa-kānū min al-khāsirīn wa-tatamassakna bi-mā
ʿindakunna wa-taftakhirna bi-irtifāʿ asmāʾikunna. Yanbaghī lakunna bi-an
tanqaṭiʿna ʿamman fī al-samāwāt wa-l-arḍ. Kadhālika quddira min ladun
muqtadir qadīr.

Say: You mirrors have been created by My command and raised up by My
will. Beware of disbelieving in the verses of My Lord and becoming like
those who have done wrong and are among the losers, or of clinging to
what you possess and boasting of the exaltation of your names. You
should detach yourselves from all who are in the heavens and on earth.
Thus has it been ordained by One who is All-Powerful, Omnipotent.

\subsection{86}\label{section-1958}

\begin{otherlanguage}{arabic}

ان یا هیکل امری قل انّی لو ارید ان اجعل الأشیآء مرایا اسمائی فی اقلّ من
حین لأقدر فکیف ربّی الّذی خلقنی بأمره المبرم المتین و لو ارید ان اقلّب
الممکنات اقرب من لمح البصر لأقدر فکیف الارادة الّتی خزنت فی مشیّة اللّه ربّی
و ربّ العالمین

\end{otherlanguage}

In yā haykal amrī qul innī law urīdu an ajʿala al-ashyāʾ marāyā asmāʾī
fī aqall min ḥīn la-aqdir fa-kayfa rabbī alladhī khalaqanī bi-amrihi
al-mubram al-matīn. Wa-law urīdu an uqalliba al-mumkināt aqrab min lamḥ
al-baṣar la-aqdir fa-kayfa al-irāda allatī khuzinat fī mashīʾat Allāh
rabbī wa-rabb al-ʿālamīn.

O Temple of My Cause! Say: If I wished to make all things mirrors of My
names in less than an instant, I would have the power to do so. How much
more, then, My Lord, who created Me through His irrevocable and mighty
command! And if I wished to transform all contingent things in less than
the blink of an eye, I would have the power to do so. How much more,
then, the purpose treasured within the Will of God, My Lord and the Lord
of all worlds!

\subsection{87}\label{section-1959}

\begin{otherlanguage}{arabic}

قل یا مظاهر اسمائی انتم لو تجاهدون فی سبیل اللّه بأموالکم و انفسکم و
تعبدون الله بعدد رمول الأرض و قطرات الأمطار و امواج البحار و تعترضون علی
مظهر الأمر حین الظّهور لا یذکر اعمالکم عند اللّه و ان ترکتم الأعمال و
آمنتم به فی تلک الأیّام عسی اللّه ان یکفّر عنکم سیّئاتکم انّه لهو العزیز
الکریم کذلک یعلّمکم اللّه ما هو المقصود لعلّ لا تستکبرون علی الّذی به ثبت ما
نزّل فی ازل الآزال طوبی لمن تقرّب الی المنظر الأکبر و سحقاً للمعرضین

\end{otherlanguage}

Qul yā maẓāhir asmāʾī antum law tujāhidūna fī sabīl Allāh bi-amwālikum
wa-anfusikum wa-taʿbudūna Allāh bi-ʿadad rumūl al-arḍ wa-qaṭarāt
al-amṭār wa-amwāj al-biḥār wa-taʿtariḍūna ʿalā maẓhar al-amr ḥīna
al-ẓuhūr lā yudhkar aʿmālukum ʿinda Allāh. Wa-in taraktum al-aʿmāl
wa-āmantum bihi fī tilka al-ayyām ʿasā Allāh an yukaffira ʿankum
sayyiʾātikum innahu la-huwa al-ʿAzīz al-Karīm. Kadhālika yuʿallimukum
Allāh mā huwa al-maqṣūd laʿalla lā tastakbirūna ʿalā alladhī bihi
thabata mā nuzzila fī azal al-āzāl. Ṭūbā li-man taqarraba ilā al-manẓar
al-akbar wa-suḥqan li-l-muʿriḍīn.

Say: Manifestations of My names! If you were to strive in the path of
God with your possessions and your lives, and worship God as many times
as there are grains of sand on earth, drops of rain, and waves of the
seas, yet oppose the Manifestation of the Cause when He appears, your
deeds would have no standing before God. But if you were to leave such
deeds undone and believe in Him in those days, God might pardon your
sins. He is indeed the Almighty, the All-Bountiful. Thus God teaches you
what is intended, so that you may not exalt yourselves over Him through
whom all that was sent down in the eternity of eternities has been
confirmed. Blessed is the one who draws near to the Most Great
Countenance, and away with those who turn aside!

\subsection{88}\label{section-1960}

\begin{otherlanguage}{arabic}

کم من عباد ینفقون اموالهم فی سبیل اللّه ولکن فی حین الظّهور نراهم من
المعرضین و کم من عباد یصومون فی الأیّام و یعترضون علی الّذی بأمره حقّق حکم
الصّوم الا انّهم من الجاهلین و کم من عباد یأکلون خبز الشّعیر و یقعدون علی
ما ینبت من الأرض و یحملون الشّدائد حفظاً لریاساتهم کذلک فصّلنا لک اعمالهم
لتکون ذکری للآخرین اولئک یحملون الشّدائد رئآء النّاس لابقآء اسمائهم بعد
الّذی لن یبقی الّا بما یلعنهم به من فی السّموات و الأرضین

\end{otherlanguage}

Kam min ʿibād yunfiqūna amwālahum fī sabīl Allāh wa-lākin fī ḥīn
al-ẓuhūr narāhum min al-muʿriḍīn. Wa-kam min ʿibād yaṣūmūna fī al-ayyām
wa-yaʿtariḍūna ʿalā alladhī bi-amrihi ḥuqqiqa ḥukm al-ṣawm. Alā innahum
min al-jāhilīn. Wa-kam min ʿibād yaʾkulūna khubz al-shaʿīr wa-yaqʿudūna
ʿalā mā yanbutu min al-arḍ wa-yaḥmilūna al-shadāʾid ḥifẓan
li-riʾāsātihim. Kadhālika faṣṣalnā laka aʿmālahum li-takūna dhikrā
li-l-ākharīn. Ulāʾika yaḥmilūna al-shadāʾid riʾāʾ al-nās li-ibqāʾ
asmāʾihim baʿda alladhī lan yabqā illā bi-mā yalʿanuhum bihi man fī
al-samāwāt wa-l-araḍīn.

How many servants spend their wealth in the path of God, yet at the time
of His appearance We find them among those who turn away! How many
servants fast by day, yet oppose Him by whose command the law of fasting
was established! They are indeed among the ignorant. How many servants
eat barley bread, sit upon what grows from the earth, and bear hardships
to preserve their positions of leadership! Thus We have set out their
deeds for you, so that this may be a reminder to others. They bear
hardships to be seen by people, to perpetuate their names, although
these will endure only through that for which those in the heavens and
the earths will curse them.

\subsection{89}\label{section-1961}

\begin{otherlanguage}{arabic}

قل لو تبقی اسمائکم کما زعمتم هل ینفعکم فی شیء لا وربّ العالمین هل عزّ عزّی
بابقآء اسمه بین الّذین یعبدون الأسمآء لا ونفس اللّه العزیز القدیر و ان لا
یذکرکم احد فی الأرض و کان اللّه راضیاً عنکم اذاً انتم فی کنائز اسمه الباطن
کذلک نزّلنا الآیات لتجذبکم الی مطلع الأنوار و تعرفوا ما اراد ربّکم العلیم
الحکیم ان امسکوا انفسکم عمّا نهیتم عنه فی الکتاب و کلوا ممّا رزقکم اللّه
حلالاً و لا تحرموا انفسکم عن نعمائه انّه لهو الکریم ذو الفضل العظیم لا
تحملوا الشّدائد علی انفسکم ان اعملوا ما بیّنّاه لکم ببراهین واضحات و آیات
لائحات و لا تکوننّ من الغافلین

\end{otherlanguage}

Qul law tabqā asmāʾukum kamā zaʿamtum hal yanfaʿukum fī shayʾ. Lā
wa-rabb al-ʿālamīn. Hal ʿazza ʿUzzā bi-ibqāʾ ismihi bayna alladhīna
yaʿbudūna al-asmāʾ. Lā wa-nafs Allāh al-ʿAzīz al-Qadīr. Wa-in lā
yadhkurukum aḥad fī al-arḍ wa-kāna Allāh rāḍiyan ʿankum idhan antum fī
kanāʾiz ismihi al-Bāṭin. Kadhālika nazzalnā al-āyāt li-tajdhibakum ilā
maṭlaʿ al-anwār wa-taʿrifū mā arāda rabbukum al-ʿAlīm al-Ḥakīm. An
amsikū anfusakum ʿammā nuhītum ʿanhu fī al-kitāb wa-kulū mimmā razaqakum
Allāh ḥalālan wa-lā taḥrimū anfusakum ʿan naʿmāʾihi innahu la-huwa
al-Karīm dhū al-faḍl al-ʿAẓīm. Lā taḥmilū al-shadāʾid ʿalā anfusikum an
iʿmalū mā bayyannāhu lakum bi-barāhīn wāḍiḥāt wa-āyāt lāʾiḥāt wa-lā
takūnanna min al-ghāfilīn.

Say: If your names endure as you imagine, will this benefit you in any
way? No, by the Lord of all worlds! Has `Uzza gained honor through the
perpetuation of its name among those who worship names? No, by the very
Self of God, the Almighty, the Omnipotent! And if no one on earth
mentions you, but God is pleased with you, then you are within the
treasuries of His name, the Hidden. Thus We have sent down the verses to
draw you toward the dawning-place of lights and enable you to know what
your Lord, the All-Knowing, the All-Wise, has intended. Restrain
yourselves from what you have been forbidden in the Book, eat of the
lawful things God has provided for you, and do not deprive yourselves of
His blessings. He is indeed the All-Bountiful, the Possessor of immense
grace. Do not impose hardships upon yourselves. Do what We have made
clear to you through evident proofs and shining verses, and do not be
among the heedless.

\subsection{90}\label{section-1962}

\begin{otherlanguage}{persian}

یا معشر العلمآء انتم لو تجتنبون الخمر و امثالها عمّا نهیتم عنه فی الکتاب
هذا لم یکن فخراً لکم لأنّ بارتکابها تضیّع مقاماتکم عند النّاس و تبدّل امورکم
و تهتک استارکم بل الفخر فی اذعانکم کلمة الحقّ و انقطاعکم فی السّرّ و الجهر
عمّا سوی اللّه العزیز القدیر طوبی لعالم ما جعل العلم حجاباً بینه و بین
المعلوم و اذا اتی القیّوم اقبل الیه بوجه منیر انّه من العلمآء یستبرکنّ
بأنفاسه اهل الفردوس و یستضیئنّ بنبراسه من فی السّموات و الأرضین انّه من
ورثة الأنبیآء من رآه قد رأی الحقّ و من اقبل الیه اقبل الی اللّه العزیز
الحکیم

\end{otherlanguage}

Yā maʿshar al-ʿulamāʾ antum law tajtanibūna al-khamr wa-amthālahā ʿammā
nuhītum ʿanhu fī al-kitāb hādhā lam yakun fakhran lakum li-anna
bi-irtikābihā tuḍayyaʿ maqāmātukum ʿinda al-nās wa-tubaddal umūrukum
wa-tuhtak astārukum. Bal al-fakhr fī idhʿānikum kalimat al-Ḥaqq
wa-inqiṭāʿikum fī al-sirr wa-l-jahr ʿammā siwā Allāh al-ʿAzīz al-Qadīr.
Ṭūbā li-ʿālim mā jaʿala al-ʿilm ḥijāban baynahu wa-bayna al-maʿlūm
wa-idhā atā al-Qayyūm aqbala ilayhi bi-wajh munīr. Innahu min al-ʿulamāʾ
yastabrikanna bi-anfāsihi ahl al-firdaws wa-yastaḍīʾanna bi-nibrāsihi
man fī al-samāwāt wa-l-araḍīn. Innahu min warathat al-anbiyāʾ man raʾāhu
qad raʾā al-Ḥaqq wa-man aqbala ilayhi aqbala ilā Allāh al-ʿAzīz
al-Ḥakīm.

Company of learned ones! If you abstain from wine and similar things
forbidden to you in the Book, this is no cause for pride, for indulging
in them would destroy your standing among people, alter your
circumstances, and expose what you conceal. Rather, honor lies in your
submission to the Word of Truth and your detachment, in private and in
public, from everything but God, the Almighty, the Omnipotent. Blessed
is the learned one who has not made knowledge a veil between himself and
the Object of knowledge, and who, when the Self-Subsisting comes, turns
toward Him with a radiant face. He is indeed among the learned; the
people of Paradise will seek blessings from his breath, and those in the
heavens and the earths will seek illumination from his lamp. He is among
the heirs of the Prophets. Whoever has seen him has seen the Truth, and
whoever has turned toward him has turned toward God, the Almighty, the
All-Wise.

\subsection{91}\label{section-1963}

\begin{otherlanguage}{arabic}

ان یا مطالع العلم ایّاکم ان تتغیّروا فی انفسکم لأنّ بتغییرکم یتغیّر اکثر
العباد انّ هذا ظلم منکم علی انفسکم و علی العباد و یشهد بذلک کلّ عارف خبیر
مثلکم کمثل عین اذا تغیّر تتغیّر الأنهار المنشعبة منها اتّقوا اللّه و کونوا
من المتّقین کذلک الانسان اذا فسد قلبه یفسد ارکانه و کذلک الشّجرة ان فسد
اصلها یفسد اغصانها و افنانها و اوراقها و اثمارها کذلک ضربنا لکم الأمثال
لعلّ لا تحتجبون بما عندکم عمّا قدّر لکم من لدن عزیز کریم

\end{otherlanguage}

In yā maṭāliʿ al-ʿilm iyyākum an tataghayyarū fī anfusikum li-anna
bi-taghayyurikum yataghayyar akthar al-ʿibād. Inna hādhā ẓulm minkum
ʿalā anfusikum wa-ʿalā al-ʿibād wa-yashhadu bi-dhālika kull ʿārif
khabīr. Mathalukum ka-mathal ʿayn idhā taghayyara tataghayyar al-anhār
al-munshaʿiba minhā. Ittaqū Allāh wa-kūnū min al-muttaqīn. Kadhālika
al-insān idhā fasada qalbuhu yafsud arkānuhu wa-kadhālika al-shajara in
fasada aṣluhā yafsud aghṣānuhā wa-afnānuhā wa-awrāquhā wa-athmāruhā.
Kadhālika ḍarabnā lakum al-amthāl laʿalla lā taḥtajibūna bi-mā ʿindakum
ʿammā quddira lakum min ladun ʿazīz karīm.

Dawning-places of knowledge! Beware of allowing yourselves to become
corrupted, for through your corruption most people will become
corrupted. This is an injustice you commit against yourselves and
against the servants, as everyone with knowledge and understanding will
testify. You are like a spring: when it is polluted, the rivers
branching from it become polluted. Be mindful of God and be among the
God-fearing. Likewise, when a person's heart becomes corrupt, the limbs
become corrupt; and when a tree's root becomes corrupt, its branches,
twigs, leaves, and fruits become corrupt. Thus We have set forth these
comparisons for you, so that what you possess may not veil you from what
has been ordained for you by One who is Almighty, All-Bountiful.

\subsection{92}\label{section-1964}

\begin{otherlanguage}{persian}

انّا لو نأخذ کفّاً من التّراب و نزیّنه بطراز الأسمآء لنقدر و هذا من فضلی علیه
من دون استحقاقه کذلک نزّل بالحقّ من لدن منزل علیم فانظروا الی الحجر الأسود
الّذی جعله اللّه مقبل العالمین هل یکون هذا الفضل من نفسه لا ونفسی و هل
یکون هذا العزّ من ذاته لا وذاتی الّذی عجز عن عرفانه من فی العالمین

\end{otherlanguage}

Innā law naʾkhudhu kaffan min al-turāb wa-nuzayyinuhu bi-ṭirāz al-asmāʾ
la-naqdir wa-hādhā min faḍlī ʿalayhi min dūn istiḥqāqihi. Kadhālika
nuzzila bi-l-ḥaqq min ladun munzil ʿalīm. Fa-nẓurū ilā al-ḥajar al-aswad
alladhī jaʿalahu Allāh muqbil al-ʿālamīn. Hal yakūnu hādhā al-faḍl min
nafsihi. Lā wa-nafsī. Wa-hal yakūnu hādhā al-ʿizz min dhātihi. Lā
wa-dhātī alladhī ʿajaza ʿan ʿirfānihi man fī al-ʿālamīn.

If We were to take a handful of dust and adorn it with the ornament of
the names, We would have the power to do so. This would be My grace
bestowed upon it, without any merit on its part. Thus has it been sent
down in truth by an All-Knowing Revealer. Consider the Black Stone,
which God has made the point toward which all peoples turn. Does this
distinction come from the stone itself? No, by My own Self! Does this
honor come from its own essence? No, by My Essence, which all who are in
the worlds have been powerless to comprehend!

\subsection{93}\label{section-1965}

\begin{otherlanguage}{arabic}

کذلک فانظر فی المسجد الأقصی و الأماکن الّتی جعلناها مطاف من فی الأطراف و
الأقطار لم یکن شرفها منها بل بما تنسب الی مظاهرنا الّذین جعلناهم مطالع
وحینا بین العباد ان انتم من العالمین و فی کلّ ذلک لحکمة لا یعلمها الّا
اللّه ان اسألوه لیبیّن لکم ما اراد انّه بکلّ شیء علیم ان انقطعوا یا قوم عن
الدّنیا و زخرفها و لا تلتفتوا الی الّذینهم کفروا و اشرکوا ان اطلعوا من افق
البیان لذکر ربّکم الرّحمن هذا ما اراده اللّه لکم طوبی للعارفین

\end{otherlanguage}

Kadhālika fa-nẓur fī al-Masjid al-Aqṣā wa-al-amākin allatī jaʿalnāhā
maṭāfa man fī al-aṭrāf wa-al-aqṭār. Lam yakun sharafuhā minhā bal bi-mā
tunsabu ilā maẓāhirinā alladhīna jaʿalnāhum maṭāliʿa waḥyinā bayna
al-ʿibād, in antum min al-ʿālimīn. Wa-fī kulli dhālika la-ḥikmatun lā
yaʿlamuhā illā Allāh. In isʾalūhu li-yubayyina lakum mā arāda, innahu
bi-kulli shayʾin ʿalīm. In inqaṭiʿū yā qawmu ʿan al-dunyā wa-zukhrufihā
wa-lā taltafitū ilā alladhīnahum kafarū wa-ashrakū. In iṭlaʿū min ufuq
al-bayān li-dhikri rabbikum al-Raḥmān. Hādhā mā arādahu Allāhu lakum.
Ṭūbā li-l-ʿārifīn.

Likewise, consider the Farthest Mosque and the places We have made
centers of pilgrimage for those in every region and land. Their honor
does not come from themselves, but from their association with Our
Manifestations, whom We have made the dawning-places of Our revelation
among the servants, if you are among those who know. In all this there
is a wisdom known only to God. Ask Him to explain to you what He has
intended; He knows all things. Detach yourselves, people, from the world
and its glitter, and pay no attention to those who have disbelieved and
associated partners with God. Rise above the horizon of utterance to
remember your Lord, the All-Merciful. This is what God has desired for
you. Blessed are those who recognize Him.

\subsection{94}\label{section-1966}

\begin{otherlanguage}{arabic}

قل یا قوم انّا امرناکم فی الألواح بأن تقدّسوا انفسکم حین الظّهور عن الأسمآء
و عن کلّ ما خلق بین الأرض و السّمآء لینطبع فیها تجلّی شمس الحقّ من افق مشیّة
ربّکم العزیز العظیم و امرناکم بأن تطهّروا نفوسکم عن حبّ من علی الأرض و
بغضهم لئلّا یمنعکم شیء عن جهة و یضطرّکم الی جهة اخری و کان هذا من اعظم
نصحی لکم فی کتاب مبین من تمسّک بأحد منهما انّه لا یقدر ان یعرف الأمر علی
ما هو علیه و یشهد بذلک کلّ منصف خبیر

\end{otherlanguage}

Qul yā qawmu innā amarnākum fī al-alwāḥ bi-an tuqaddisū anfusakum ḥīna
al-ẓuhūr ʿan al-asmāʾ wa-ʿan kulli mā khuliqa bayna al-arḍ wa-al-samāʾ
li-yanṭabiʿa fīhā tajallī shams al-Ḥaqq min ufuq mashīʾat rabbikum
al-ʿAzīz al-ʿAẓīm. Wa-amarnākum bi-an tuṭahhirū nufūsakum ʿan ḥubbi man
ʿalā al-arḍ wa-bughḍihim li-allā yamnaʿakum shayʾun ʿan jihatin
wa-yaḍṭarrakum ilā jihatin ukhrā. Wa-kāna hādhā min aʿẓam nuṣḥī lakum fī
kitābin mubīn. Man tamassaka bi-aḥadin minhumā innahu lā yaqdiru an
yaʿrifa al-amra ʿalā mā huwa ʿalayhi, wa-yashhadu bi-dhālika kullu
munṣifin khabīr.

Say: People, We commanded you in the Tablets to sanctify yourselves, at
the time of the Revelation, from names and from everything created
between earth and heaven, so that the radiance of the Sun of Truth,
shining from the horizon of the will of your Lord, the Almighty, the
Magnificent, might be imprinted upon your souls. We commanded you to
purify your souls of love for those on earth and hatred of them, lest
anything hold you back from one direction and compel you toward another.
This was among My greatest counsels to you in a clear Book. Whoever
clings to either cannot recognize the Cause as it truly is. Every
fair-minded and discerning person bears witness to this.

\subsection{95}\label{section-1967}

\begin{otherlanguage}{arabic}

انتم نسیتم عهد اللّه و نقضتم میثاقه الی ان اعرضتم عن الّذی بظهوره قرّت عیون
الموحّدین طهّروا الأنظار عن الحجب و الأستار ثمّ انظروا حجج النّبیّین و
المرسلین لتعرفوا امر اللّه فی تلک الأیّام الّتی فیها اتی الموعود بسلطان
عظیم اتّقوا اللّه و لا تحرموا انفسکم عن مطلع الآیات هذا ما تنتفع به انفسکم
انّ ربّکم لغنیّ عن العالمین انّه لم یزل کان و لم یکن معه من شیء قد ارتفعت
باسمه رایة التّوحید علی طور الوجود من الغیب و الشّهود علی انّه لا اله الّا
انا الواحد العزیز الفرید

\end{otherlanguage}

Antum nasītum ʿahda Allāh wa-naqaḍtum mīthāqahu ilā an aʿraḍtum ʿan
alladhī bi-ẓuhūrihi qarrat ʿuyūnu al-muwaḥḥidīn. Ṭahhirū al-anẓāra ʿan
al-ḥujub wa-al-astār thumma nẓurū ḥujaja al-nabiyyīn wa-al-mursalīn
li-taʿrifū amra Allāh fī tilka al-ayyām allatī fīhā atā al-Mawʿūd
bi-sulṭānin ʿaẓīm. Ittaqū Allāha wa-lā taḥrimū anfusakum ʿan maṭlaʿ
al-āyāt. Hādhā mā tanta fiʿu bihi anfusukum. Inna rabbakum la-ghaniyyun
ʿan al-ʿālamīn. Innahu lam yazal kāna wa-lam yakun maʿahu min shayʾ. Qad
irtafaʿat bi-ismihi rāyatu al-tawḥīd ʿalā ṭūr al-wujūd min al-ghayb
wa-al-shuhūd ʿalā annahu lā ilāha illā anā al-Wāḥid al-ʿAzīz al-Farīd.

You forgot God's pledge and broke His covenant, until you turned away
from the One whose appearance brought joy to the eyes of those who
affirm His unity. Purify your vision of veils and coverings, then
consider the proofs of the Prophets and Messengers, so that you may
recognize the Cause of God in these days when the Promised One has come
with mighty authority. Be mindful of God and do not deprive yourselves
of the Dawning-place of the verses. This is what will benefit your
souls. Your Lord has no need of any of the worlds. He has always
existed, and nothing existed with Him. In His name the banner of divine
unity has been raised upon the Sinai of existence, both unseen and
visible, proclaiming: ``There is no God but Me, the One, the Almighty,
the Peerless.''

\subsection{96}\label{section-1968}

\begin{otherlanguage}{arabic}

انّ الّذینهم خلقوا بارادة من عنده و بعثوا بأمره اعرضوا عنه و اتّخذوا
لأنفسهم ربّاً من دون اللّه الا انّهم من المبعدین کانوا ان یذکروا الرّحمن فی
کلّ الأحیان و اذا ظهر بالحقّ حاربوا معه افّ لهم بما نقضوا المیثاق اذ اشرق
نیّر الآفاق من افق مشیّة اللّه المقدّس العلیم الحکیم سلّوا سیوف البغضآء علی
وجه اللّه و لا یستشعرون فی انفسهم کأنّهم اموات فی قبور اهوائهم بعد الّذی
فاحت نسمة اللّه بین الدّیار الا انّهم فی حجاب عظیم اذا تتلی علیهم آیات اللّه
یصرّون مستکبرین کأنّهم ما عرفوا شیئاً و ما سمعوا نغمة اللّه العلیّ العلیم

\end{otherlanguage}

Inna alladhīnahum khuliqū bi-irādatin min ʿindihi wa-buʿithū bi-amrihi
aʿraḍū ʿanhu wa-ittakhadhū li-anfusihim rabban min dūni Allāh. Alā
innahum min al-mubʿadīn. Kānū an yadhkurū al-Raḥmāna fī kulli al-aḥyān
wa-idhā ẓahara bi-al-ḥaqq ḥārabū maʿahu. Uffin lahum bi-mā naqaḍū
al-mīthāqa idh ashraqa nayyiru al-āfāq min ufuq mashīʾat Allāh
al-Muqaddas al-ʿAlīm al-Ḥakīm. Sallū suyūfa al-baghḍāʾ ʿalā wajh Allāh
wa-lā yastashʿirūna fī anfusihim ka-annahum amwātun fī qubūri ahwāʾihim
baʿda alladhī fāḥat nasmatu Allāh bayna al-diyār. Alā innahum fī ḥijābin
ʿaẓīm. Idhā tutlā ʿalayhim āyātu Allāh yuṣirrūna mustakbirīn ka-annahum
mā ʿarafū shayʾan wa-mā samiʿū nagh mata Allāh al-ʿAliyy al-ʿAlīm.

Those who were created through a will proceeding from Him and raised up
at His command have turned away from Him and taken for themselves a lord
other than God. Truly, they are among those who are cast far away. They
used to remember the All-Merciful at all times, yet when He appeared in
truth, they waged war against Him. Shame on them for breaking the
covenant when the Luminary of the horizons shone from the horizon of the
will of God, the Holy, the All-Knowing, the All-Wise. They drew the
swords of hatred against the face of God, yet have no awareness within
themselves, as though they were dead in the graves of their desires,
even after the breeze of God had wafted through the lands. Truly, they
are wrapped in a mighty veil. When the verses of God are recited to
them, they persist in arrogance, as though they knew nothing and had not
heard the melody of God, the Most High, the All-Knowing.

\subsection{97}\label{section-1969}

\begin{otherlanguage}{arabic}

قل وا حسرة علیکم أ تدّعون الایمان فی انفسکم و تکفرون بآیات اللّه العزیز
العلیم قل یا قوم ولّوا وجوهکم شطر ربّکم الرّحمن ایّاکم ان یحجبنّکم ما نزّل فی
البیان انّه ما نزّل الّا لذکری العزیز المنیع و ما کان مقصوده الّا جمالی قد
ملئت الآفاق {[}من{]} برهانی لو انتم من المنصفین

\end{otherlanguage}

Qul wā ḥasratan ʿalaykum! A-taddaʿūna al-īmāna fī anfusikum wa-takfurūna
bi-āyāt Allāh al-ʿAzīz al-ʿAlīm? Qul yā qawmu wallū wujūhakum shaṭra
rabbikum al-Raḥmān. Iyyākum an yaḥjubannakum mā nuzzila fī al-Bayān.
Innahu mā nuzzila illā li-dhikrī al-ʿazīz al-manīʿ, wa-mā kāna maqṣūduhu
illā jamālī. Qad muliʾat al-āfāqu {[}min{]} burhānī law antum min
al-munṣifīn.

Say: Alas for you! Do you claim faith for yourselves while disbelieving
in the verses of God, the Almighty, the All-Knowing? Say: People, turn
your faces toward your Lord, the All-Merciful. Beware lest what was sent
down in the Bayān veil you. It was sent down only for My mighty,
unassailable remembrance, and its sole purpose was My Beauty. The
horizons have been filled with My proof, if only you were fair-minded.

\subsection{98}\label{section-1970}

\begin{otherlanguage}{persian}

لو کان النّقطة الأولی علی زعمکم غیری و یدرک لقائی لن یفارق منّی و یستأنس
بنفسی و استأنست بنفسه فی ایّامی انّه ناح لفراقی قد سبقنی لیبشّر النّاس
بملکوتی کذلک نزّل فی الألواح ان انتم من النّاظرین فیا لیت یکون من ذی سمع
لیسمع ضجیجه فی البیان بما ورد علی نفسی من هؤلآء الغافلین و یعرف حنینه فی
فراقی و شغفه الی لقائی العزیز البدیع اذاً یشاهد محبوبه بین العباد الّذینهم
خلقوا لأیّامه و السّجود بین یدیه بالذّلّة الّتی اعترف القلم بالعجز عن ذکرها
بما ورد علیه من هؤلآء الظّالمین

\end{otherlanguage}

Law kāna al-Nuqṭatu al-Ūlā ʿalā zaʿmikum ghayrī wa-yudriku liqāʾī lan
yufāriqa minnī wa-yastaʾnisu bi-nafsī wa-istaʾnastu bi-nafsihi fī
ayyāmī. Innahu nāḥa li-firāqī. Qad sabaqanī li-yubashshira al-nāsa
bi-malakūtī. Kadhālika nuzzila fī al-alwāḥ in antum min al-nāẓirīn.
Fa-yā layta yakūnu min dhī samʿin li-yasmaʿa ḍajījahu fī al-Bayān bi-mā
warada ʿalā nafsī min hāʾulāʾ al-ghāfilīn wa-yaʿrifa ḥanīnahu fī firāqī
wa-shaghafahu ilā liqāʾī al-ʿazīz al-badīʿ. Idhan yushāhidu maḥbūbahu
bayna al-ʿibād alladhīnahum khuliqū li-ayyāmihi wa-al-sujūd bayna
yadayhi bi-al-dhillati allatī iʿtarafa al-qalamu bi-al-ʿajz ʿan dhikrihā
bi-mā warada ʿalayhi min hāʾulāʾ al-ẓālimīn.

If the Primal Point were, as you suppose, other than Me and were to
attain My presence, He would never part from Me and would find
companionship in My Self, and I would find companionship in His self in
My days. He lamented His separation from Me. He preceded Me to announce
to the people the glad tidings of My Kingdom. Thus it was sent down in
the Tablets, if you are among those who look. If only there were someone
with hearing to hear His lament in the Bayān over what has befallen Me
at the hands of these heedless people, and to recognize His yearning in
separation from Me and His ardent desire to attain My mighty, wondrous
presence. Then he would behold His Beloved among the servants who were
created for His days and to bow down before Him, in such abasement that
the Pen has confessed its inability to describe it, because of what has
befallen Him at the hands of these wrongdoers.

\subsection{99}\label{section-1971}

\begin{otherlanguage}{arabic}

قل یا قوم انّا دعوناکم فی الظّهور الأوّل الی المنظر الأکبر هذا المقام
الأطهر و بشّرناکم بأیّام اللّه فلمّا انشقّ السّتر الأعظم و اتی جمال القدم علی
سحاب القدر کفرتم بالّذی آمنتم فویل لکم یا معشر المشرکین خافوا من اللّه و
لا تدحضوا الحقّ بما عندکم اذا اشرقت علیکم شمس الآیات من افق اصبع ملیک
الأسمآء و الصّفات خرّوا بوجوهکم سجّداً للّه ربّ العالمین انّ سجودکم فنآء بابه
لیکون خیراً من عبادة الثّقلین و خضوعکم عند ظهوره خیر لکم عمّا خلق فی
السّموات و الأرضین

\end{otherlanguage}

Qul yā qawmu innā daʿawnākum fī al-ẓuhūr al-awwal ilā al-manẓar
al-akbar, hādhā al-maqām al-aṭhar, wa-bashsharnākum bi-ayyām Allāh.
Fa-lammā inshaqqa al-sitru al-aʿẓam wa-atā Jamālu al-Qidam ʿalā saḥāb
al-qadar kafartum bi-alladhī āmantum. Fa-waylun lakum yā maʿshara
al-mushrikīn. Khāfū min Allāh wa-lā tudḥiḍū al-ḥaqqa bi-mā ʿindakum.
Idhā ashraqat ʿalaykum shamsu al-āyāt min ufuq iṣbaʿ malīk al-asmāʾ
wa-al-ṣifāt khirrū bi-wujūhikum sujjadan li-llāhi rabbi al-ʿālamīn. Inna
sujūdakum fināʾa bābihi la-yakūnu khayran min ʿibādat al-thaqalayn,
wa-khuḍūʿukum ʿinda ẓuhūrihi khayrun lakum ʿammā khuliqa fī al-samāwāt
wa-al-araḍīn.

Say: People, in the first Revelation We called you to the Most Great
Vision, this most pure station, and announced to you the glad tidings of
the days of God. Yet when the greatest veil was torn apart and the
Ancient Beauty came upon the clouds of destiny, you disbelieved in the
One in whom you had believed. Woe to you, you who associate partners
with God! Fear God and do not refute the truth with what you possess.
When the Sun of the verses shines upon you from the horizon of the
finger of the Sovereign of names and attributes, fall on your faces in
prostration before God, the Lord of the worlds. Your prostration at the
threshold of His gate is indeed better than the worship of both
humankind and jinn, and your humility at His appearance is better for
you than whatever has been created in the heavens and the earths.

\subsection{100}\label{section-1972}

\begin{otherlanguage}{arabic}

قل یا قوم اذکّرکم لوجه اللّه و ما ارید منکم جزآء ان اجری الّا علی الله الّذی
فطرنی و بعثنی بالحقّ و جعلنی ذکراً للخلائق اجمعین ان اسرعوا الی منظر اللّه
و مقرّه و لا تتّبعوا الشّیطان فی انفسکم انّه یأمرکم بالبغی و الفحشآء و
یمنعکم عن الصّراط الّذی نصب فی العالم بهذا الأمر المبرم الحکیم

\end{otherlanguage}

Qul yā qawmu udhakkirukum li-wajh Allāh wa-mā urīdu minkum jazāʾan. In
ajriya illā ʿalā Allāh alladhī faṭaranī wa-baʿathanī bi-al-ḥaqq
wa-jaʿalanī dhikran li-l-khalāʾiq ajmaʿīn. In asriʿū ilā manẓar Allāh
wa-maqarrihi wa-lā tattabiʿū al-shayṭāna fī anfusikum. Innahu yaʾmurukum
bi-al-baghy wa-al-faḥshāʾ wa-yamnaʿukum ʿan al-ṣirāṭ alladhī nuṣiba fī
al-ʿālam bi-hādhā al-amr al-mubram al-ḥakīm.

Say: People, I admonish you for the sake of God and desire no reward
from you. My reward rests only with God, who created Me, raised Me up in
truth, and made Me a remembrance for all creatures. Hasten to the place
of God's manifestation and His dwelling, and do not follow Satan within
yourselves. He commands you to commit injustice and indecency and bars
you from the path established in the world through this irrevocable,
wise Command.

\subsection{101}\label{section-1973}

\begin{otherlanguage}{arabic}

قل قد ظهر الشّیطان بشأن ما ظهر شبهه فی الامکان و کذلک ظهر جمال الرّحمن
بالطّراز الّذی ما ادرکت مثله عیون الأوّلین قد ارتفع ندآء الرّحمن و عن ورائه
ندآء الشّیطان طوبی لمن سمع ندآء اللّه و توجّه الی جهة العرش منظر قدس کریم
من کان فی قلبه اقلّ من خردل حبّ دونی لن یقدر ان یدخل ملکوتی و برهانی ما
طرّز به دیباج کتاب الوجود ان انتم من العارفین قل الیوم یوم فیه ظهر الفضل
الأعظم و لم یکن شیء لا فی السّموات العلی و لا فی الأراضی السّفلی الّا و
ینطقنّ بذکری و یغرّدنّ بثنآء نفسی ان انتم من السّامعین

\end{otherlanguage}

Qul qad ẓahara al-shayṭānu bi-shaʾnin mā ẓahara shibhuhu fī al-imkān,
wa-kadhālika ẓahara Jamālu al-Raḥmān bi-al-ṭirāz alladhī mā adrakat
mithlahu ʿuyūnu al-awwalīn. Qad irtafaʿa nidāʾu al-Raḥmān wa-ʿan
warāʾihi nidāʾu al-shayṭān. Ṭūbā li-man samiʿa nidāʾa Allāh wa-tawajjaha
ilā jihat al-ʿarsh, manẓari qudsin karīm. Man kāna fī qalbihi aqallu min
khardalin ḥubbu dūnī lan yaqdira an yadkhula malakūtī, wa-burhānī mā
ṭurriza bihi dībāju kitāb al-wujūd, in antum min al-ʿārifīn. Qul
al-yawmu yawmun fīhi ẓahara al-faḍlu al-aʿẓam wa-lam yakun shayʾun lā fī
al-samāwāt al-ʿulā wa-lā fī al-arāḍī al-suflā illā wa-yanṭiqanna
bi-dhikrī wa-yugharridanna bi-thanāʾ nafsī, in antum min al-sāmiʿīn.

Say: Satan has appeared in a manner whose like has never appeared in the
realm of existence. Likewise, the Beauty of the All-Merciful has
appeared in an adornment the like of which the eyes of earlier
generations have never beheld. The call of the All-Merciful has been
raised, and behind it the call of Satan. Blessed is whoever has heard
the call of God and turned toward the Throne, a noble seat of holiness.
Anyone whose heart contains even less than a mustard seed of love for
anyone other than Me will never be able to enter My Kingdom. My proof is
that with which the brocade of the Book of Existence has been
embroidered, if you are among those who understand. Say: Today is the
day when the greatest grace has appeared, and there is nothing, either
in the heavens above or in the earths below, that does not speak My
remembrance and sing My praise, if you are among those who hear.

\subsection{102}\label{section-1974}

\begin{otherlanguage}{arabic}

ان یا هیکل الظّهور ان انفخ فی الصّور باسمی ثمّ ان یا هیکل الأسرار تنفّس فی
المزمار بذکر ربّک المختار ثمّ ان یا حوریّة الفردوس ان اخرجی من غرف الجنان
ثمّ اخبری اهل الأکوان تاللّه قد ظهر محبوب العالمین و مقصود العارفین و
معبود من فی السّموات و الأرضین و مسجود الأوّلین و الآخرین

\end{otherlanguage}

In yā haykala al-ẓuhūr, in unfukh fī al-ṣūr bi-ismī. Thumma in yā
haykala al-asrār, tanaffas fī al-mizmār bi-dhikri rabbika al-Mukhtār.
Thumma in yā ḥūriyyata al-Firdaws, in ukhrujī min ghuraf al-jinān thumma
akhbirī ahla al-akwān: ta-Allāhi qad ẓahara maḥbūbu al-ʿālamīn
wa-maqṣūdu al-ʿārifīn wa-maʿbūdu man fī al-samāwāt wa-al-araḍīn
wa-masjūdu al-awwalīn wa-al-ākhirīn.

Temple of Revelation, blow the trumpet in My name. Then, Temple of
Mysteries, breathe into the flute in remembrance of your Lord, the
Chosen One. Then, Maiden of Paradise, come out from the chambers of the
gardens and announce to the inhabitants of all realms: ``By God, the
Beloved of the worlds has appeared, the One sought by those who
recognize Him, the One worshipped by all who are in the heavens and the
earths, before whom the first and the last bow down.''

\subsection{103}\label{section-1975}

\begin{otherlanguage}{persian}

ایّاکم ان تتوقّفوا فی هذا الجمال بعد الّذی ظهر بسلطان القدرة و القوّة و
الاستجلال انّه لهو الحقّ و ما سواه معدوم عند احد من عباده و مفقود لدی ظهور
انواره ان اسرعوا الی کوثر الفضل و لا تکوننّ من الصّابرین و من توقّف اقلّ من
آن لیحبط اللّه عمله و یرجعه الی مقرّ القهر فبئس مثوی المتوقّفین

\end{otherlanguage}

Iyyākum an tatawaqqafū fī hādhā al-jamāl baʿda alladhī ẓahara bi-sulṭān
al-qudra wa-al-quwwa wa-al-istijlāl. Innahu la-huwa al-Ḥaqq wa-mā siwāhu
maʿdūmun ʿinda aḥadin min ʿibādihi wa-mafqūdun ladā ẓuhūr anwārihi. In
asriʿū ilā kawthar al-faḍl wa-lā takūnanna min al-ṣābirīn. Wa-man
tawaqqafa aqalla min ānin la-yuḥbiṭi Allāhu ʿamalahu wa-yurjiʿhu ilā
maqarri al-qahr, fa-biʾsa mathwā al-mutawaqqifīn.

Beware of hesitating over this Beauty after He has appeared with the
sovereignty of might, power, and majesty. He is indeed the Truth, and
everything besides Him is as nothing before even one of His servants and
vanishes at the appearance of His lights. Hasten to the river of grace
and do not be among those who wait. Whoever hesitates for less than a
moment, God will surely nullify his deeds and return him to the abode of
wrath. How wretched is the dwelling of those who hesitate!

\bookmarksetup{startatroot}

\chapter{BH00008 (Unnamed)}\label{bh00008-unnamed}

\emph{Translated on 17 Qudrat 183 using GPT 6.1 Sol}

\subsection{1}\label{section-1976}

\begin{otherlanguage}{arabic}

بسم ربّنا الأقدس الأعظم الأبهی

\end{otherlanguage}

Bismi Rabbinā al-Aqdas al-Aʿẓam al-Abhā.

In the name of our Lord, the Most Holy, the Most Great, the
All-Glorious.

\subsection{2}\label{section-1977}

\begin{otherlanguage}{arabic}

الحمد للّه الّذی وفی بالوعد و انزل الموعود من سمآء مشیّته و اظهر ما کان
مسطوراً من قلمه الأعلی فی کتبه و صحفه و زبره و الواحه قد قام القیام و اتت
السّاعة و ظهر الصّراط و سرت نسمة اللّه و فاح عرف قمیصه رغماً لملأ البیان
الّذین نقضوا میثاق اللّه و عهده و ارادوا ان تعمروا بیوت الظّنون و الأوهام
بأیادی الظّلم و الاعتساف کما عمّروها من قبل تبّاً لهم و سحقاً لهم و بعداً لهم
یا ملأ الأوهام قد انار الأفق الأعلی و ظهر ما تستره حجبات انفسکم و لا
اوهامکم انّه قد ظهر بسلطان من عنده و خرق السّبحات و الحجبات بأصابع القدرة
علی شأن سمع الملأ الأعلی صوت خرقها اتّقوا اللّه یا معشر الغافلین لا تصدّوا
النّاس عن السّبیل و لا تمنعوهم عن الدّلیل هذا یوم فیه ماج بحر العلم و هاج
عرف اللّه المهیمن القیّوم خافوا اللّه و لا تدحضوا الحقّ بأقوالکم و افعالکم
انّ الخادم ینصحکم لوجه اللّه یشهد بذلک کلّ منصف بصیر

\end{otherlanguage}

Al-ḥamdu li-llāh alladhī wafā bi-l-waʿd wa-anzala al-Mawʿūd min samāʾ
mashīʾatihi wa-aẓhara mā kāna masṭūran min qalamihi al-aʿlā fī kutubihi
wa-ṣuḥufihi wa-zuburihi wa-alwāḥihi. Qad qāma al-qiyām wa-atati al-sāʿa
wa-ẓahara al-ṣirāṭ wa-sarat nasmatu Allāh wa-fāḥa ʿarfu qamīṣihi raghman
li-malaʾ al-Bayān alladhīna naqaḍū mīthāqa Allāh wa-ʿahdahu wa-arādū an
taʿmurū buyūta al-ẓunūn wa-l-awhām bi-aydī al-ẓulm wa-l-iʿtisāf kamā
ʿammarūhā min qabl. Tabban lahum wa-suḥqan lahum wa-buʿdan lahum. Yā
malaʾ al-awhām, qad anāra al-ufuq al-aʿlā wa-ẓahara mā tasturuhu
ḥujubātu anfusikum wa-lā awhāmukum. Innahu qad ẓahara bi-sulṭānin min
ʿindihi wa-kharaqa al-subuḥāt wa-l-ḥujubāt bi-aṣābiʿ al-qudra ʿalā
shaʾnin samiʿa al-malaʾ al-aʿlā ṣawta kharqihā. Ittaqū Allāh yā maʿshar
al-ghāfilīn. Lā taṣuddū al-nāsa ʿan al-sabīl wa-lā tamnaʿūhum ʿan
al-dalīl. Hādhā yawmun fīhi māja baḥru al-ʿilm wa-hāja ʿarfu Allāh
al-Muhaymin al-Qayyūm. Khāfū Allāh wa-lā tudḥiḍū al-ḥaqqa bi-aqwālikum
wa-afʿālikum. Inna al-khādima yanṣaḥukum li-wajhi Allāh, yashhadu
bi-dhālika kullu munṣifin baṣīr.

Praise be to God, who has fulfilled the promise, sent down the Promised
One from the heaven of His will, and revealed what was inscribed by His
Most Exalted Pen in His books, scriptures, writings, and tablets. The
Resurrection has arisen, the Hour has come, the Path has appeared, the
breeze of God has blown, and the fragrance of His garment has spread,
despite the people of the Bayán, who have broken God's covenant and
pledge and have desired that you build the houses of conjecture and
delusion with the hands of oppression and tyranny, as they built them
before. Perish they; may they be crushed and banished! People of
delusion! The Most Exalted Horizon has shone forth, and there has
appeared what neither the veils of your own selves nor your delusions
conceal. He has appeared with a sovereignty of His own and has torn
through the veils of glory and concealment with the fingers of power, so
that the Concourse on high has heard the sound of their rending. Fear
God, you who are heedless! Do not turn people away from the path or keep
them from the Guide. This is a day on which the ocean of knowledge has
surged and the fragrance of God, the All-Protecting, the
Self-Subsisting, has spread. Fear God, and do not seek to refute the
truth through your words and deeds. This servant counsels you for the
sake of God; every fair-minded, discerning person bears witness to this.

\subsection{3}\label{section-1978}

\begin{otherlanguage}{persian}

و الصّلوة و السّلام و التّکبیر و البهآء علی الّذین شربوا رحیق الاستقامة و
قاموا علی خدمة الأمر بالحکمة و البیان و شهدوا بما شهد به لسان العظمة قبل
خلق الأشیآء اولئک عباد وصفهم اللّه فی الفرقان و فی البیان طوبی لذی اذن
سمع و لذی بصر رأی و ویل لکلّ خادع مکّار الّذی اتّخذ الوهم لنفسه ربّا من دون
اللّه و یمنع النّاس عن فرات رحمة ربّهم الغفور الکریم

\end{otherlanguage}

Wa-l-ṣalātu wa-l-salāmu wa-l-takbīru wa-l-bahāʾ ʿalā alladhīna sharibū
raḥīqa al-istiqāma wa-qāmū ʿalā khidmati al-amr bi-l-ḥikma wa-l-bayān
wa-shahidū bi-mā shahida bihi lisānu al-ʿaẓama qabla khalqi al-ashyāʾ.
Ulāʾika ʿibādun waṣafahumu Allāh fī al-Furqān wa-fī al-Bayān. Ṭūbā
li-dhī udhunin samiʿa wa-li-dhī baṣarin raʾā, wa-waylun li-kulli
khādiʿin makkār alladhī ittakhadha al-wahma li-nafsihi rabban min dūni
Allāh wa-yamnaʿu al-nāsa ʿan furāti raḥmati rabbihim al-Ghafūr al-Karīm.

And blessings, peace, exaltation, and glory be upon those who have drunk
the choice wine of steadfastness, arisen to serve the Cause with wisdom
and eloquence, and testified to what the Tongue of Grandeur testified
before the creation of all things. These are servants whom God has
described in the Qur'an and in the Bayán. Blessed is the one with an ear
who has heard, and the one with sight who has seen; and woe to every
scheming deceiver who has taken delusion as his lord instead of God and
keeps people from the sweet waters of the mercy of their Lord, the
Ever-Forgiving, the All-Bountiful.

\subsection{4}\label{section-1979}

\begin{otherlanguage}{arabic}

عرض میشود نامهٴ آنجناب باین خادم فانی رسید بعد از قرائت و اطّلاع قصد مقام
مالک ابداع نموده تلقاء وجه حاضر و بعد از اذن آنچه در نامه بود بشرف اصغا
فائز گشت و اینکلمات عالیات از سماء عنایت منزل آیات نازل قوله جلّ جلاله

\end{otherlanguage}

ʿArż mī-shavad nāmih-yi ān-janāb bi-īn khādim-i fānī rasīd. Baʿd az
qirāʾat va iṭṭilāʿ, qaṣd-i maqām-i Mālik-i Ibdāʿ namūdih, tilqāʾ-i vajh
ḥāżir va baʿd az idhn ānchih dar nāmih būd bi-sharaf-i iṣghā fāʾiz
gasht, va īn kalimāt-i ʿāliyāt az samā-yi ʿināyat-i Munzil-i Āyāt nāzil.
Qawluhu jalla jalāluhu:

I respectfully submit that your letter reached this transient servant.
After reading it and learning its contents, I proceeded to the abode of
the Lord of Creation and entered His presence. After permission was
granted, the contents of the letter were honored with His hearing, and
these exalted words descended from the heaven of the grace of the
Revealer of Verses. His words---exalted be His glory:

\subsection{5}\label{section-1980}

\begin{otherlanguage}{arabic}

انا المظلوم الغریب

\end{otherlanguage}

Anā al-Maẓlūm al-Gharīb.

I am the Wronged One, the Stranger.

\subsection{6}\label{section-1981}

\begin{otherlanguage}{arabic}

یا ایّها النّاظر الی الوجه نامه‌ات نزد مظلوم حاضر و عبد حاضر تمام آنرا عرض
نمود للّه الحمد رحیق عرفان را نوشیدی و بافق رحمن توجّه نمودی در ایّامیکه
افئده و قلوب از ظلم مشرکین مضطرب و ارکان بیت از اعتساف علما متزلزل در
اوّل هر ظهور از مظهر الهی و مشرق وحی ربّانی اعراض نمودند و فتوی بر سفک
دماء آن نفوس مقدّسه دادند

\end{otherlanguage}

Yā ayyuhā al-nāẓir ilā al-wajh, nāmih-at nazd-i Maẓlūm ḥāżir va ʿabd-i
ḥāżir tamām-i ān-rā ʿarż namūd. Li-llāh al-ḥamd, raḥīq-i ʿirfān-rā
nūshīdī va bi-ufuq-i Raḥmān tavajjuh namūdī dar ayyāmī kih afʾidih va
qulūb az ẓulm-i mushrikīn mużṭarib va arkān-i bayt az iʿtisāf-i ʿulamā
mutazalzil. Dar avval-i har ẓuhūr az Maẓhar-i Ilāhī va Mashriq-i Vaḥy-i
Rabbānī iʿrāż namūdand va fatvā bar safk-i dimā-yi ān nufūs-i muqaddasih
dādand.

You who gaze toward the Face! Your letter is before the Wronged One, and
the servant in attendance has presented it in full. Praise be to God,
you have drunk the choice wine of recognition and turned toward the
horizon of the All-Merciful in days when hearts and souls are agitated
by the oppression of those who associate partners with God, and the
pillars of the House are shaken by the tyranny of the religious
scholars. At the beginning of every revelation, they turned away from
the Manifestation of God and the Dawning-Place of divine revelation, and
issued decrees for the shedding of the blood of those holy souls.

\subsection{7}\label{section-1982}

\begin{otherlanguage}{persian}

هجرت خاتم انبیا روح ما سواه فداه از مکّه سبب که بود و علّت چه و صعود حضرت
روح بفلک چهارم از دست که بقسمی بر آنحضرت سخت شد که ید قدرت او را صعود
داد و بفلک چهارمش برد به عمامه و ردایات حقّ را بکمال ظلم و اعتساف بستند و
باب عزّت و ثروت باسم آنحضرت از برای خود گشودند اگر نفوس مطمئنّه در امورات
قبل تفکّر کنند لعمراللّه بمثابهٴ فاقدین بر مظلوم عالم نوحه و ندبه نمایند

\end{otherlanguage}

Hijrat-i Khātam-i Anbiyā, rūḥ-i mā sivāh fidāh, az Makkih sabab kih būd
va ʿillat chih, va ṣuʿūd-i Ḥaẓrat-i Rūḥ bi-falak-i chahārum az dast-i
kih? Bi-qismī bar ān Ḥaẓrat sakht shud kih yad-i qudrat ū-rā ṣuʿūd dād
va bi-falak-i chahārumash burd. Bih ʿamāmih va ridāyāt Ḥaqq-rā
bi-kamāl-i ẓulm va iʿtisāf bastand va bāb-i ʿizzat va thirvat bi-ism-i
ān Ḥaẓrat az barā-yi khud gushūdand. Agar nufūs-i muṭmaʾinnih dar
umūrāt-i qabl tafakkur kunand, la-ʿamru Allāh, bi-mathābih-yi fāqidīn
bar Maẓlūm-i ʿālam nawḥih va nudbih namāyand.

Who brought about the departure of the Seal of the Prophets---may the
souls of all besides Him be sacrificed for Him---from Mecca, and what
was its cause? And at whose hands did the Spirit ascend to the fourth
heaven? He was so sorely afflicted that the hand of power raised Him up
and carried Him to the fourth heaven. With turbans and robes they bound
the True One with the utmost oppression and tyranny, and in His name
opened for themselves the door to honor and wealth. If assured souls
reflect on the events of the past, by the life of God, they will lament
and weep over the Wronged One of the world like those bereaved.

\subsection{8}\label{section-1983}

\begin{otherlanguage}{persian}

هنگام حرکت از زوراء کلّ را اخبار نمودیم بارتفاع نعاق در آفاق جمیع از خوف
و اضطراب خلف حجاب ساکن و حقّ وحده بر امر قیام نمود و چون اشراقات انوار
آفتاب توحید پرتو افکند و عالم فی الجمله روشن و منیر گشت از خلف حجاب
بیرون آمده اوّل بر نفس حقّ اعتراض نمودند و بسیوف اعراض بر سفک دم مطهّرش
قیام کردند اینست شأن ناس ان اعتبروا یا اولی الألباب یا ایّها المشتعل بنار
محبّة اللّه حمد کن محبوب عالم را که بنعمت باقیهٴ دائمه فائز شدی و آن
ذکریست که از قلم اعلی در باره تو جاری شد آنچه دیده میشود البتّه او را فنا
اخذ نماید مگر کلمه‌ئی که حقّ جلّ جلاله بآن تکلّم فرماید از حقّ بطلب ترا مؤیّد
نماید بر آنچه که سزاوار یوم اوست معادله نمینماید بکلمهٴ او خزائن ارض و
دفائن آن ان اشکر و قل لک الحمد یا من ذکرتنی و لک الثّنآء یا من بیدک زمام
ملکوت ملک السّموات و الأرضین انتهی

\end{otherlanguage}

Hingām-i ḥarakat az Zawrā kull-rā ikhbār namūdīm bi-irtifāʿ-i nuʿāq dar
āfāq. Jamīʿ az khawf va iẓṭirāb khalf-i ḥijāb sākin va Ḥaqq vaḥdahu bar
amr qiyām namūd, va chūn ishrāqāt-i anvār-i āftāb-i tawḥīd partaw afkand
va ʿālam fī al-jumlih rawshan va munīr gasht, az khalf-i ḥijāb bīrūn
āmadih avval bar nafs-i Ḥaqq iʿtirāż namūdand va bi-suyūf-i iʿrāż bar
safk-i dam-i muṭahharash qiyām kardand. Īnast shaʾn-i nās. An iʿtabirū
yā ulī al-albāb. Yā ayyuhā al-mushtaʿil bi-nāri maḥabbati Allāh, ḥamd
kun Maḥbūb-i ʿālam-rā kih bi-niʿmat-i bāqiyih-yi dāʾimih fāʾiz shudī, va
ān dhikrīst kih az Qalam-i Aʿlā dar bārih-yi tu jārī shud. Ānchih dīdih
mī-shavad al-battih ū-rā fanā akhdh namāyad magar kalimih-ī kih Ḥaqq
jalla jalāluhu bi-ān takallum farmāyad. Az Ḥaqq bi-ṭalab tu-rā muʾayyad
namāyad bar ānchih kih sazāvār-i yawm-i ūst. Muʿādilih na-mī-namāyad
bi-kalimih-yi ū khazāʾin-i arż va dafāʾin-i ān. An ishkur wa-qul: Laka
al-ḥamdu yā man dhakartanī, wa-laka al-thanāʾu yā man bi-yadika zimāmu
malakūti mulki al-samāwāti wa-l-araḍīn. Intahā.

At the time of Our departure from Zawrá, We informed everyone that
clamor would rise throughout the world. All remained behind the veil in
fear and agitation, while the True One alone arose for the Cause. When
the radiance of the light of the sun of divine unity shone forth and the
world became somewhat bright and illumined, they emerged from behind the
veil and first objected to the True One Himself, and arose with the
swords of rejection to shed His pure blood. Such is the way of the
people. Take heed, you who possess understanding! You who are aflame
with the fire of the love of God! Praise the Beloved of the world, for
you have attained an abiding, everlasting blessing: the mention of you
that has flowed from the Most Exalted Pen. Whatever is seen will
assuredly perish, except a word uttered by the True One, exalted be His
glory. Ask God to confirm you in what befits His Day. The treasuries of
the earth and its buried riches cannot equal a single word of His. Give
thanks and say: ``Praise be to You, You who have remembered me; and all
acclaim be to You, in whose hand are the reins of the kingdom of
dominion over the heavens and the earth.'' Here ends the quotation.

\subsection{9}\label{section-1984}

\begin{otherlanguage}{persian}

دیگر خادم چه عرض نماید حقّ شاهد و گواهست که ارکانش متزلزل و قلبش مضطرب چه
که دیده میشود بعضی از عباد جاهل که خود را عالم فرض گرفته‌اند در جمیع قرون
و اعصار بر مطالع نور احدیّه اعتراض نموده‌اند و بکمال عناد در اطفاء نور
الهی ساعی و جاهد بوده‌اند و بعد که امر مرتفع میشد بدو دست اخذ مینمودند و
باسباب زهد زمام عباد را بتصرّف میآوردند چنانچه دیده و شنیده‌اید و لهم
الأشباه و الأمثال فی هذه الأیّام نفوسیکه ابداً اطّلاع از این امر نداشته حال
به ظلم و اعتساف قیام کرده‌اند و دعوی اطّلاع کامل نموده‌اند و نفوس ضعیفه را
از شریعهٴ ربّانیّه منع مینمایند و بغدیر منتنه ترغیب و تحریض میکنند حقّ عباد
را حفظ فرماید و از وساوس آن نفوس محفوظ دارد یعنی منافقین از اهل بیان
اعاذنا اللّه و ایّاکم یا معشر الموحّدین من شرّ هؤلآء

\end{otherlanguage}

Dīgar khādim chih ʿarż namāyad? Ḥaqq shāhid va guvāh ast kih arkānash
mutazalzil va qalbash mużṭarib, chih kih dīdih mī-shavad baʿżī az
ʿibād-i jāhil kih khud-rā ʿālim farż giriftih-and dar jamīʿ-i qurūn va
aʿṣār bar maṭāliʿ-i nūr-i aḥadiyyih iʿtirāż namūdih-and va bi-kamāl-i
ʿinād dar iṭfā-yi nūr-i ilāhī sāʿī va jāhid būdih-and, va baʿd kih amr
murtafiʿ mī-shud bi-du dast akhdh mī-namūdand va bi-asbāb-i zuhd zimām-i
ʿibād-rā bi-taṣarruf mī-āvardand, chinānchih dīdih va shinīdih-īd.
Wa-lahumu al-ashbāhu wa-l-amthālu fī hādhihi al-ayyām. Nufūsī kih abadan
iṭṭilāʿ az īn amr nadāshtih ḥāl bih ẓulm va iʿtisāf qiyām kardih-and va
daʿvā-yi iṭṭilāʿ-i kāmil namūdih-and va nufūs-i żaʿīfih-rā az
sharīʿih-yi rabbāniyyih manʿ mī-namāyand va bi-ghadīr-i muntinih targhīb
va taḥrīṣ mī-kunand. Ḥaqq ʿibād-rā ḥifẓ farmāyad va az vasāvis-i ān
nufūs maḥfūẓ dārad, yaʿnī munāfiqīn az ahl-i Bayān. Aʿādhanā Allāhu
wa-iyyākum yā maʿshar al-muwaḥḥidīn min sharri hāʾulāʾ.

What more can this servant say? God bears witness that his limbs are
trembling and his heart is troubled, for it is evident that in every age
and century certain ignorant servants, imagining themselves learned,
have objected to the Dawning-Places of the light of divine unity and
have striven with the utmost obstinacy to extinguish the divine light.
Then, when the Cause became exalted, they would seize it with both hands
and, by means of asceticism, take hold of the reins of the people, as
you have seen and heard. They have their counterparts and their like in
these days. Souls who had no knowledge whatever of this Cause have now
arisen in oppression and tyranny, have claimed complete knowledge, and
are keeping weak souls from the divine waterside, urging and encouraging
them toward a stagnant, foul-smelling pool. May God protect His servants
and preserve them from the whisperings of those souls, namely, the
hypocrites among the people of the Bayán. May God shelter us and you,
believers in His unity, from their evil.

\subsection{10}\label{section-1985}

\begin{otherlanguage}{persian}

این خادم فانی از حقّ جلّ جلاله سائل و آمل که اولیای خود را مؤیّد فرماید تا
از بحر معانی کلمات الهی بیاشامند و از نار مشتعلهٴ در آن بشأنی مشتعل
گردند که همزات شیاطین ایشانرا از مالک یوم الدّین منع ننماید در کاف و را
نعیق مرتفع و مطلع ظنون و نفاق بکمال مکر و حیله و کلمات جعلیّه نفوس ضعیفه
را از مالک احدیّه منع نموده و بناحیهٴ کذبه دعوت مینماید مع آنکه ابداً از
این امر اطّلاع نداشته و ندارد محض هوی از مالک اسماء اعتراض جسته و در تیه
اوهام سرگردان از حقّ میطلبم او را هدایت نماید و از رحیق معانی و سلسبیل
راستی محروم نسازد و از دریای آگاهی قسمت عطا فرماید اوست بر هر شیء قادر و
توانا

\end{otherlanguage}

Īn khādim-i fānī az Ḥaqq jalla jalāluhu sāʾil va āmil kih awliyā-yi
khud-rā muʾayyad farmāyad tā az baḥr-i maʿānī-yi kalimāt-i ilāhī
biyāshāmand va az nār-i mushtaʿilih-yi dar ān bi-shaʾnī mushtaʿil
gardand kih hamazāt-i shayāṭīn īshān-rā az Mālik-i Yawm-i Dīn manʿ
na-namāyad. Dar Kāf va Rā naʿīq murtafiʿ va maṭlaʿ-i ẓunūn va nifāq
bi-kamāl-i makr va ḥīlih va kalimāt-i jaʿliyyih nufūs-i żaʿīfih-rā az
Mālik-i Aḥadiyyih manʿ namūdih va bi-nāḥiyih-yi kadhabih daʿvat
mī-namāyad, maʿa ānkih abadan az īn amr iṭṭilāʿ nadāshtih va na-dārad.
Maḥż-i havā az Mālik-i Asmā iʿtirāż justih va dar tīh-i awhām sargardān.
Az Ḥaqq mī-ṭalabam ū-rā hidāyat namāyad va az raḥīq-i maʿānī va
salsabīl-i rāstī maḥrūm na-sāzad va az daryā-yi āgāhī qismat ʿaṭā
farmāyad. Ūst bar har shayʾ qādir va tavānā.

This transient servant asks and hopes that God, exalted be His glory,
will confirm His loved ones, so that they may drink from the ocean of
the meanings of the divine words and be so enkindled by the fire blazing
within them that the promptings of devils will not keep them from the
Lord of the Day of Judgment. In Káf and Rá a clamor has risen, and the
source of conjecture and hypocrisy, with the utmost cunning and deceit
and with fabricated words, has kept weak souls from the Lord of Oneness
and is summoning them toward the liars, although he has never had, and
does not now have, any knowledge of this Cause. Out of sheer desire he
has taken exception to the Lord of Names and wanders bewildered in the
wilderness of delusions. I ask God to guide him, not to deprive him of
the choice wine of meanings and the flowing spring of truthfulness, and
to grant him a share of the ocean of awareness. He is powerful and
mighty over all things.

\subsection{11}\label{section-1986}

\begin{otherlanguage}{persian}

در کتاب اقدس که از ملکوت مقدّس الهی در اوّل ایّام سجن نازل ذکر ارض کاف و را
و ما یحدث فیها از قلم اعلی جاری و ثبت شد و تا حال مستور ولکن ناعق آن ارض
در این سنه بکلماتی تکلّم نموده که کادت السّموات ان تتفطّرن و تنشقّ الأرض و
تخرّ الجبال مدّا

\end{otherlanguage}

Dar Kitāb-i Aqdas kih az malakūt-i muqaddas-i ilāhī dar avval-i ayyām-i
sijn nāzil, dhikr-i arż-i Kāf va Rā va mā yaḥduthu fīhā az Qalam-i Aʿlā
jārī va thabt shud va tā ḥāl mastūr, va-lākin nāʿiq-i ān arż dar īn
sanih bi-kalimātī takallum namūdih kih kādati al-samāwātu an
tatafaṭṭarna wa-tanshaqqa al-arḍu wa-takhirra al-jibālu maddan.

In the Most Holy Book, which descended from the holy, divine Kingdom in
the early days of the imprisonment, mention of the land of Káf and Rá
and what would occur there flowed from the Most Exalted Pen and was
recorded, and until now has remained concealed. Yet the clamorer of that
land has this year uttered words at which the heavens have almost split
apart, the earth rent asunder, and the mountains fallen flat.

\subsection{12}\label{section-1987}

\begin{otherlanguage}{persian}

حقّ انصاف عنایت فرماید و عدل و راستی بخشد امثال آن نفوس هزار و دویست سنه
بل ازید ناس را به جابلقا و جابلسا دعوت نمودند و هر نفس آگاهی که اینکلمه
را ردّ نمود و مردم را اخبار کرد که حضرت قائم در اصلاب است و باید متولّد
شود فتوی بر قتلش دادند و خونش ریختند قسم بآفتاب افق بیان الرّحمن که از
جمله آن نفوس همین ناعق ارض کاف و راء است چه که بر قدم آن نفوس مشی
مینماید باری چه مقدار از نفوس مقدّسه بفتوای آن نفوس غافلهٴ جاهله کأس
شهادت نوشیدند آنچه ذکر شد جمیع اولیا و دوستان الهی گواهی میدهند بر صدق
آن بگمان خود نصرت امر حقّ مینمودند و ناس را بسبیل مستقیم هدایت میکردند و
حال آنکه آن سبیل سبیل جحیم بود نه صراط مستقیم

\end{otherlanguage}

Ḥaqq inṣāf ʿināyat farmāyad va ʿadl va rāstī bakhshad. Amthāl-i ān nufūs
hizār va duvīst sanih bal azyad nās-rā bih Jābulqā va Jābulsā daʿvat
namūdand va har nafs-i āgāhī kih īn kalimih-rā radd namūd va mardum-rā
ikhbār kard kih Ḥaẓrat-i Qāʾim dar aṣlāb ast va bāyad mutavallid shavad,
fatvā bar qatl-ash dādand va khūnash rīkhtand. Qasam bi-āftāb-i ufuq-i
bayān-i al-Raḥmān kih az jumlih-yi ān nufūs hamīn nāʿiq-i arż-i Kāf va
Rāʾ ast, chih kih bar qadam-i ān nufūs mashy mī-namāyad. Bārī chih
miqdār az nufūs-i muqaddasih bi-fatvā-yi ān nufūs-i ghāfilih-yi jāhilih
kaʾs-i shahādat nūshīdand. Ānchih dhikr shud jamīʿ-i awliyā va dūstān-i
ilāhī guvāhī mī-dihand bar ṣidq-i ān. Bi-gumān-i khud nuṣrat-i amr-i
Ḥaqq mī-namūdand va nās-rā bi-sabīl-i mustaqīm hidāyat mī-kardand va ḥāl
ānkih ān sabīl sabīl-i jaḥīm būd nah ṣirāṭ-i mustaqīm.

May God grant fairness and bestow justice and truthfulness. People like
those souls summoned the people to Jábulqá and Jábulsá for twelve
hundred years, indeed longer. Whenever an enlightened soul rejected this
assertion and informed the people that the Qá'im was in the loins of His
forebears and must be born, they issued a decree for his death and shed
his blood. I swear by the Sun of the horizon of the All-Merciful's
utterance that this very clamorer of the land of Káf and Rá is among
those souls, for he follows in their footsteps. How many holy souls
drank the cup of martyrdom at the decree of those heedless, ignorant
souls! All the loved ones and friends of God bear witness to the truth
of what has been stated. They imagined that they were aiding the Cause
of God and guiding people to the straight way, whereas that way was the
way of hell, not the straight path.

\subsection{13}\label{section-1988}

\begin{otherlanguage}{persian}

آیا ناس را چه سکری اخذ نموده که تا حین بر حیله و مکر گروه منافقین اطّلاع
نیافته‌اند آنحضرت دوستان الهی را آگاه فرمائید که خود را از مظاهر حیل و
نفاق و مطالع خدعه و فریب باسم حقّ جلّ جلاله حفظ نمایند و از آنقوم مطمئن
مشوند چه که مشاهده میشود چون از نفسی مأیوس شوند یعنی مشاهده نمایند
اوهاماتشان در او تصرّف ننمود و تأثیر نکرد بلباس دوستی درآیند و بکمال محبّت
و مودّت در ظاهر معاشرت نمایند و مقصودشان آنکه فرصت یابند و القا کنند آنچه
را که در قلوبشان مستور است در هر حال اجتناب لازم باسباب دوستی اساس دشمنی
را محکم نمایند اقبالشان اعراض و شفقتشان عداوت وصمتشان خدعه و نطقشان فریب

\end{otherlanguage}

Āyā nās-rā chih sukrī akhdh namūdih kih tā ḥīn bar ḥīlih va makr-i
gurūh-i munāfiqīn iṭṭilāʿ na-yāftih-and? Ān ḥaẓrat dūstān-i ilāhī-rā
āgāh farmāyīd kih khud-rā az maẓāhir-i ḥiyal va nifāq va maṭāliʿ-i
khudʿih va farīb bi-ism-i Ḥaqq jalla jalāluhu ḥifẓ namāyand va az ān
qawm muṭmaʾin na-shavand, chih kih mushāhadih mī-shavad chūn az nafsī
maʾyūs shavand, yaʿnī mushāhadih namāyand awhāmātishān dar ū taṣarruf
na-namūd va taʾthīr na-kard, bi-libās-i dūstī dar-āyand va bi-kamāl-i
maḥabbat va mavaddat dar ẓāhir muʿāsharat namāyand va maqṣūdishān ānkih
furṣat yāband va ilqā kunand ānchih-rā kih dar qulūbishān mastūr ast.
Dar har ḥāl ijtināb lāzim. Bi-asbāb-i dūstī asās-i dushmanī-rā muḥkam
namāyand. Iqbālishān iʿrāż va shafaqatishān ʿadāvat va ṣamtishān khudʿih
va nuṭqishān farīb.

What intoxication has seized the people that, even now, they have not
become aware of the deceit and cunning of the band of hypocrites? Please
alert the friends of God to protect themselves, in the name of God,
exalted be His glory, from the embodiments of guile and hypocrisy and
the sources of deception and deceit, and not to place their trust in
those people. For it is observed that when they lose hope of
someone---that is, when they see that their delusions have neither taken
hold of him nor influenced him---they assume the guise of friendship and
outwardly associate with him with the utmost love and affection. Their
purpose is to find an opportunity to instill what lies concealed in
their hearts. Avoidance is necessary in every circumstance. Through the
means of friendship they strengthen the foundations of enmity. Their
approach is rejection, their kindness hostility, their silence
deception, and their speech deceit.

\subsection{14}\label{section-1989}

\begin{otherlanguage}{persian}

بگوئید نای بیخردان فوارس مضمار انقطاع الیوم جز حقّ نبینند و جز حقّ نشناسند
و جز حقّ نگویند چه که یوم یوم اللّه است لایُذکر فیه الّا هو الیوم اگر نفسی
بغیر حقّ نطق نماید و ناظر باسماء موهومه شود او لایق اینمقام و این ذکر و
این امر نبوده و نیست بگوئید ای غافلان از خدا بترسید و ناس بیچاره را بدام
میندازید بگذارید بقوادم انقطاع در این هوا طیران نمایند و بکمال انبساط در
این بساط وارد شوند یا لیت القوم یعرفون و یعلمون چندی قبل در یکی از الواح
اینکلمهٴ علیا از لسان مالک وری جاری و نازل قوله تبارک و تعالی

\end{otherlanguage}

Bigūʾīd: {[}corrupted text{]} bī-khiradān, favāris-i miżmār-i inqiṭāʿ
al-yawm juz Ḥaqq na-bīnand va juz Ḥaqq na-shināsand va juz Ḥaqq
na-gūyand, chih kih yawm Yawm-i Allāh ast, lā yudhkaru fīhi illā Huwa.
Al-yawm agar nafsī bi-ghayr-i Ḥaqq nuṭq namāyad va nāẓir bi-asmā-yi
mawhūmih shavad, ū lāyiq-i īn maqām va īn dhikr va īn amr na-būdih va
nīst. Bigūʾīd: Ay ghāfilān, az Khudā bitarsīd va nās-i bīchārih-rā
bi-dām ma-yandāzīd. Bigudhārīd bi-qavādim-i inqiṭāʿ dar īn havā ṭayarān
namāyand va bi-kamāl-i inbisāṭ dar īn bisāṭ vārid shavand. Yā layta
al-qawma yaʿrifūna wa-yaʿlamūn. Chandī qabl dar yakī az alvāḥ īn
kalimih-yi ʿulyā az lisān-i Mālik-i Varā jārī va nāzil. Qawluhu tabāraka
wa-taʿālā:

Say: {[}corrupted original text{]} foolish ones, the riders in the arena
of detachment today see none but God, recognize none but God, and speak
of none but God, for this day is the Day of God, in which none is
mentioned but Him. Today, if a soul speaks of anyone other than God and
turns his gaze toward imaginary names, he has never been, and is not,
worthy of this station, this mention, and this Cause. Say: Heedless
ones, fear God and do not ensnare the helpless people. Let them fly in
this atmosphere on the wings of detachment and enter this court with the
utmost joy. If only the people would recognize and understand! Some time
ago, in one of the Tablets, this exalted utterance flowed from the
tongue of the Lord of all creation and was sent down. His
words---blessed and exalted be He:

\subsection{15}\label{section-1990}

\begin{otherlanguage}{persian}

ای عبد حاضر آیا این خلق تفکّر نمینمایند در عظمت این یوم امروز روزیست که
بزرگی آن در جمیع کتب الهی از قلم اعلی مذکور و مسطور است اگرچه ایّام ظهور
مظاهر احدیّه و مطالع نور الهیّه در مقامی بایّام اللّه مذکور ولکن این ایّام در
جمیع کتب تخصیص یافته و مخصوص است بحقّ جلّ جلاله چنانچه بخاتم انبیا
میفرماید یوم یقوم النّاس لربّ العالمین و همچنین الملک یومئذ للّه و جئی ربّک
و الملک صفّاً صفّا و امثال آن و به کلیم میفرماید ان اخرج القوم من الظّلمات
الی النّور و ذکّرهم بایّام اللّه جمیع رسل را باین یوم مبارک بشارت میفرماید
مع‌ذلک گاهی مرآت هنگامی وصیّ و وقتی ثمره میگویند باین اسماء عاکفند و شاعر
نیستند یتکلّمون و لا یشعرون یقولون و لا یفقهون قل موتوا بغیظکم هذا یوم
اللّه لو انتم تعرفون و هذا امر اللّه لو انتم تنصفون قد ظهر من کان مکنوناً
فی علم اللّه و مسطوراً فی کتبه و به ظهر کلّ امر مختوم و کلّ سرّ مخزون انتهی

\end{otherlanguage}

Ey ʿabd-e ḥāżer, āyā īn khalq tafakkor nemīnamāyand dar ʿaẓamat-e īn
yowm? Emrūz rūzīst ke bozorgī-ye ān dar jamīʿ-e kotob-e elāhī az qalam-e
aʿlā madhkūr va masṭūr ast. Agarche ayyām-e ẓohūr-e maẓāher-e aḥadīyye
va maṭāleʿ-e nūr-e elāhīyye dar maqāmī be-ayyām Allāh madhkūr, va lāken
īn ayyām dar jamīʿ-e kotob takhṣīṣ yāfte va makhṣūṣ ast be-Ḥaqq, jalla
jalāluhu. Chenānche be-Khātam-e Anbiyā mīfarmāyad: yawma yaqūmu al-nāsu
li-rabbi al-ʿālamīn; va hamchenīn: al-mulku yawmaʾidhin li-llāh; va jāʾa
rabbuka wa-l-malaku ṣaffan ṣaffā; va amthāl-e ān. Va be-Kalīm
mīfarmāyad: an akhrij al-qawma mina al-ẓulumāti ilā al-nūri
wa-dhakkirhum bi-ayyāmi Allāh. Jamīʿ-e rosol rā be-īn yowm-e mobārak
beshārat mīfarmāyad. Maʿa dhālek gāhī merʾāt, hengāmī vaṣī va vaqtī
thamare mīgūyand; be-īn asmā ʿākefand va shāʿer nīstand. Yatakallamūna
wa-lā yashʿurūn, yaqūlūna wa-lā yafqahūn. Qul mūtū bi-ghayẓikum. Hādhā
yawmu Allāh law antum taʿrifūn, wa-hādhā amru Allāh law antum tunṣifūn.
Qad ẓahara man kāna maknūnan fī ʿilmi Allāh wa-masṭūran fī kutubihi,
wa-bihi ẓahara kullu amrin makhtūmin wa-kullu sirrin makhzūn. Intahā.

O servant in attendance, do these people not reflect on the greatness of
this Day? Today is a day whose greatness has been mentioned and recorded
by the Most Exalted Pen in all the divine Books. Although the days of
the appearance of the Manifestations of divine unity and the
Dawning-Places of divine light are, in one sense, called the Days of
God, these days have been singled out in all the Books and belong
especially to God, exalted be His glory. Thus He says to the Seal of the
Prophets: ``The day when the people will stand before the Lord of the
worlds''; and likewise: ``Sovereignty on that day belongs to God''; and
``Your Lord comes, and the angels, rank upon rank,'' and other such
statements. And to Him who conversed with God He says: ``Bring the
people out of darkness into light, and remind them of the Days of God.''
He gives all the Messengers the glad tidings of this blessed Day.
Nevertheless, at one time they speak of a mirror, at another of a
successor, and at another of a fruit. They cling to these names and are
unaware. They speak without awareness; they utter words without
understanding. Say: ``Die in your rage! This is the Day of God, if only
you recognized it, and this is the Cause of God, if only you judged
fairly. He who was concealed in the knowledge of God and recorded in His
Books has appeared, and through Him every sealed matter and every hidden
secret has been disclosed.'' End of quotation.

\subsection{16}\label{section-1991}

\begin{otherlanguage}{persian}

حال از پستی نفوس ذباب بر ربّ الأرباب اعتراض نموده غراب را از عالم معانی
قسمتی نه و ذباب را از ترنّمات الهی نصیبی نه این میدان را فوارس عرفان باید
تا باسم حقّ جلّ جلاله مرکب همّت رانند سبحان‌اللّه بعضی از نفوس محض هوی تکلّم
نموده و مینمایند چه که از خمر اوهام نوشیده‌اند و غیر اللّه را ربّ اخذ
کرده‌اند در بادهای ظنون مشی نمایند و باوهامات خود ناس را گمراه کنند اگر
خادم در اینمقامات آنچه شنیده عرض نماید البتّه اوراق سخن کفایت ننماید
آیاتیکه در اوّل ورود سجن در کتاب اقدس نازل ذکر میشود که شاید غافلین بشعور
آیند و نائمین متنبّه شوند و قاعدین قیام نمایند قوله عزّ بیانه و جلّ مقامه

\end{otherlanguage}

Ḥāl az pastī-ye nofūs, dhobāb bar Rabb al-Arbāb eʿterāż namūde. Ghorāb
rā az ʿālam-e maʿānī qesmatī na va dhobāb rā az tarannomāt-e elāhī
naṣībī na. Īn meydān rā favāres-e ʿerfān bāyad tā be-esm-e Ḥaqq, jalla
jalāluhu, markab-e hemmat rānand. Sobḥān Allāh! Baʿżī az nofūs maḥż-e
havā takallom namūde va mīnamāyand, che ke az khamr-e owhām nūshīde-and
va gheyr Allāh rā rabb akhdh karde-and. Dar bādhā-ye ẓonūn mashy
namāyand va be-owhāmāt-e khod nās rā gomrāh konand. Agar khādem dar īn
maqāmāt ānche shenīde ʿarż namāyad, al-batte owrāq-e sokhan kefāyat
nanamāyad. Āyātī ke dar avval-e vorūd-e sejn dar Ketāb-e Aqdas nāzel
dhekr mīshavad ke shāyad ghāfelīn be-shoʿūr āyand va nāʾemīn motanabbeh
shavand va qāʿedīn qīyām namāyand. Qawluhu ʿazza bayānuhu wa-jalla
maqāmuhu:

Now, such is the baseness of these souls that a fly has raised
objections against the Lord of lords. The raven has no share in the
realm of inner meanings, and the fly has no portion of the divine
melodies. This arena calls for the horsemen of spiritual knowledge to
spur on the steed of endeavor in the name of God, exalted be His glory.
Glory be to God! Some souls have spoken, and continue to speak, out of
sheer desire, for they have drunk the wine of delusions and taken
someone other than God as their lord. They move amid the winds of
conjecture and lead people astray with their imaginings. If this servant
were to relate what he has heard concerning these matters, pages of
discourse would certainly not suffice. The verses revealed in the Most
Holy Book upon first entering the prison are cited here, so that perhaps
the heedless may become aware, the sleeping may awaken, and those who
sit idle may arise. His words---mighty is His utterance and exalted His
station:

\subsection{17}\label{section-1992}

\begin{otherlanguage}{persian}

ان یا ارض الکاف و الرّاء انّا نراک علی ما لایحبّه اللّه و نری منک ما لا اطّلع
به الّا اللّه العلیم الخبیر و نجد ما تمرّ منک فی سرّ السّرّ عندنا علم کلّشیء فی
لوح مبین لا تحزنی بذلک سوف یظهر اللّه فیک اولی بأس شدید یذکروننی باستقامة
لا تمنعهم اشارات العلمآء و لا تحجبهم شبهات المریبین اولئک ینظرنّ اللّه
بأعینهم و ینصرنّه بأنفسهم الا انّه من الرّاسخین انتهی

\end{otherlanguage}

In yā arḍa al-Kāf wa-l-Rāʾ, innā narāki ʿalā mā lā yuḥibbuhu Allāh,
wa-narā minki mā lā iṭṭalaʿa bihi illā Allāhu al-ʿAlīmu al-Khabīr,
wa-najidu mā tamurru minki fī sirri al-sirr. ʿIndanā ʿilmu kulli shayʾin
fī lawḥin mubīn. Lā taḥzanī bi-dhālika. Sawfa yuẓhiru Allāhu fīki ulī
baʾsin shadīdin yadhkurūnanī bi-istiqāmatin lā tamnaʿuhum ishārātu
al-ʿulamāʾ wa-lā taḥjubuhum shubuhātu al-murībīn. Ulāʾika yanẓurunna
Allāha bi-aʿyunihim wa-yanṣurunnahu bi-anfusihim. Alā innahu mina
al-rāsikhīn. Intahā.

O land of Kaf and Ra! We see you in a state displeasing to God, and We
see from you what none but God, the All-Knowing, the All-Aware, has
perceived. We discern what proceeds from you in the deepest secrecy.
With Us is the knowledge of all things in a clear Tablet. Do not grieve
over this. God will raise up within you people of great might who will
remember Me with a steadfastness that neither the allusions of the
learned can hinder nor the doubts of the doubters veil. They will surely
behold God with their own eyes and aid Him with their own selves. Truly,
he is among the steadfast. End of quotation.

\subsection{18}\label{section-1993}

\begin{otherlanguage}{persian}

وقتی این آیات نازل که ذکر ناعق آن ارض نبوده و کتاب اقدس ده سنه قبل آن
نازل و کلّ مطّلعند لازال نفوس جاهله خود را برداء و عمّامه میآراستند و در
لباس علم باب آنرا سدّ مینمودند آیا کدام یک از روایات و احادیث و اخبارشان
صدق و مطابق بوده عباد بیچاره از جهل آن نفوس غافله قرنها در تیه اوهام و
ظنون سائر باری لوجه اللّه گفته میشود که شاید بنیان وهم را تجدید ننمایند و
اساس ظنون را بر ارض نگذارند

\end{otherlanguage}

Vaqtī īn āyāt nāzel ke dhekr-e nāʿeq-e ān arż nabūde va Ketāb-e Aqdas
dah sane qabl-e ān nāzel va koll moṭṭaleʿand. Lā yazāl nofūs-e jāhele
khod rā be-radā va ʿamāme mīārāstand va dar lebās-e ʿelm bāb-e ān rā
sadd mīnamūdand. Āyā kodām yek az revāyāt va aḥādīth va akhbāreshān ṣedq
va moṭābeq būde? ʿEbād-e bīchāre az jahl-e ān nofūs-e ghāfele qarnhā dar
tīh-e owhām va ẓonūn sāʾer. Bārī, li-wajhi Allāh gofte mīshavad ke
shāyad bonyān-e vahm rā tajdīd nanamāyand va asās-e ẓonūn rā bar arż
nagodhārand.

These verses were revealed when there was no mention of the clamorer of
that land, and the Most Holy Book was revealed ten years before that, as
everyone knows. Ignorant souls have always adorned themselves with
cloaks and turbans and, in the guise of learning, barred its door. Which
of their traditions, sayings, and reports has proved true and accorded
with reality? Because of the ignorance of those heedless souls, the
hapless people have wandered for centuries in the wilderness of
delusions and conjectures. In any event, this is said for the sake of
God, so that perhaps they may not rebuild the structure of delusion or
lay the foundations of conjecture upon the earth.

\subsection{19}\label{section-1994}

\begin{otherlanguage}{persian}

امروز باید بقوّت خلیل آفاق سلاسل تقلید را بشکنند و اصنام اوهام را از هم
بریزند گوش را باید از آنچه شنیده‌اند مقدّس نمایند از گفتگوها بگذرند و از
اقوال بافعال توجّه نمایند بفرمائید ای غافلان من قام علی الأمر و من الّذی
دعا الکلّ الی الأفق الأعلی فی ایّام فیها اضطربت افئدة اولی النّهی سبحان‌اللّه
امروز را چه دانسته‌اند و از این ظهور اعظم چه ادراک نموده‌اند

\end{otherlanguage}

Emrūz bāyad be-qovvat-e Khalīl-e āfāq salāsel-e taqlīd rā beshkanand va
aṣnām-e owhām rā az ham berīzand. Gūsh rā bāyad az ānche shenīde-and
moqaddas namāyand, az goftogūhā bogdharand va az aqvāl be-afʿāl tavajjoh
namāyand. Befarmāyīd: ey ghāfelān, man qāma ʿalā al-amri wa-man alladhī
daʿā al-kulla ilā al-ufuq al-aʿlā fī ayyāmin fīhā iḍṭarabat afʾidatu ulī
al-nuhā? Sobḥān Allāh! Emrūz rā che dāneste-and va az īn ẓohūr-e aʿẓam
che edrāk namūde-and?

Today they must, through the power of the Friend of all the world, break
the chains of imitation and shatter the idols of delusion. They must
purify their ears of what they have heard, leave discussions behind, and
turn from words to deeds. Say: ``O heedless ones, who arose to uphold
the Cause, and who called everyone to the Most Exalted Horizon in days
when the hearts of those endowed with understanding were shaken?'' Glory
be to God! What have they understood of this Day, and what have they
perceived of this Most Great Revelation?

\subsection{20}\label{section-1995}

\begin{otherlanguage}{persian}

هذا یوم لا یذکر فیه الّا هو و لا ینبغی لأحد ان یتمسّک بالأمثال و الأشباه و
الأسمآء هذا یوم نطقت بثنائه کتب اللّه و شهد له النّبیّون و المرسلون انسان
متحیّر است چه ذکر نماید حزب شیعه که خود را اعظم و اعلم و اتقی از جمیع
احزاب عالم میشمردند نزد اشراقات انوار آفتاب ظهور واضح و مبرهن شد که
پستترین احزابند چه که بفتوای آن نفوس محبوب عالمیان را در ارض تا شهید
نمودند الا لعنة اللّه علی القوم الظّالمین باینهم قناعت ننمودند فتوی بر قتل
جمیع نفوس مقدّسه دادند چه مقدار از اطفال و نساء که از ظلم آن نفوس شهید
شدند و چه مقدار که بی محلّ و مأوی و بی مربّی ماندند

\end{otherlanguage}

Hādhā yawmun lā yudhkaru fīhi illā huwa, wa-lā yanbaghī li-aḥadin an
yatamassaka bi-l-amthāli wa-l-ashbāhi wa-l-asmāʾ. Hādhā yawmun naṭaqat
bi-thanāʾihi kutubu Allāh wa-shahida lahu al-nabiyyūna wa-l-mursalūn.
Ensān motaḥayyer ast che dhekr namāyad. Ḥezb-e Shīʿe ke khod rā aʿẓam va
aʿlam va atqā az jamīʿ-e aḥzāb-e ʿālam mīshomordand, nazd-e eshrāqāt-e
anvār-e āftāb-e ẓohūr vāżeḥ va mobarhan shod ke past-tarīn-e aḥzāband,
che ke be-fatvā-ye ān nofūs Maḥbūb-e ʿālamīyān rā dar arż-e Tā shahīd
namūdand. Alā laʿnatu Allāhi ʿalā al-qawmi al-ẓālimīn. Be-īn ham qanāʿat
nanamūdand; fatvā bar qatl-e jamīʿ-e nofūs-e moqaddase dādand. Che
meqdār az aṭfāl va nesā ke az ẓolm-e ān nofūs shahīd shodand va che
meqdār ke bī maḥall va maʾvā va bī morabbī māndand.

This is a Day in which none is mentioned but Him, and no one should
cling to likenesses, counterparts, or names. This is a Day whose praise
the Books of God have voiced and to which the Prophets and Messengers
have borne witness. One is bewildered as to what to say. The Shiite
community considered itself greater, more learned, and more pious than
all the communities of the world; yet in the radiance of the light of
the Sun of Revelation, it became clear and evident that they are the
lowest of communities, for it was by the ruling of those souls that the
Beloved of all the worlds was martyred in the land of Ta. Truly, the
curse of God rests upon the unjust people. Nor were they content with
this: they issued a ruling for the killing of all the holy souls. How
many children and women were martyred through the oppression of those
souls, and how many were left without a home, a refuge, or anyone to
care for them!

\subsection{21}\label{section-1996}

\begin{otherlanguage}{persian}

حال تازه جمعی از کذّابهای عالم جمع شده‌اند و در صدد ترتیب اساس قبل بوده و
هستند ناعق ارض کاف و را و ناعقین اخری بر آنند که انوار آفتاب راستی را
بغمام ظنون جدیده ستر نمایند و دو سه هزار سنهٴ دیگر به ردّ و قبول و اعراض
و اعتراض و سبّ و لعن یکدیگر بر منابر و مجالس مشغول شوند قسم بآفتاب افق
تقدیس اینعبد از این اذکار خجل و منفعل است جمیع علمای عصر که هر یک در
ایران پیشوا بودند و صاحب منبر و محراب عرف حقیقت را علی قدر سمّ ابره
نیافتند چه اگر یافته بودند این ظلم جدید لا مثل له در عالم واقع نمیشد محض
هوی و ریاست تکلّم نمودند بآنچه که سبب و علّت شهادت مقصود عالم نقطهٴ اولی
روح ما سواه فداه گشت چه اگر آن اذکار مجعوله نبود حقّ و اولیائش شهید
نمیشدند جمیع علما در سنهٴ اوّل ذکر نمودند که جمیع علما از قبل و بعد بر
این بوده‌اند که آنحضرت باید از جابلقا و جابلسا ظاهر شود و حال این حزب
میگویند در شیراز متولّد شد لذا جمیع حکم بر کفر و سفک دم مطهّر نمودند

\end{otherlanguage}

Ḥāl tāze jamʿī az kadhdhābhā-ye ʿālam jamʿ shode-and va dar ṣadad-e
tartīb-e asās-e qabl būde va hastand. Nāʿeq-e arż-e Kāf va Rā va
nāʿeqīn-e okhrā bar ānand ke anvār-e āftāb-e rāstī rā be-ghamām-e
ẓonūn-e jadīde satr namāyand va do se hezār sane-ye dīgar be-radd va
qabūl va eʿrāż va eʿterāż va sabb va laʿn-e yekdīgar bar manāber va
majāles mashghūl shavand. Qasam be-āftāb-e ofoq-e taqdīs, īn ʿabd az īn
adhkār khajel va monfaʿel ast. Jamīʿ-e ʿolamā-ye ʿaṣr ke har yek dar
Īrān pīshvā būdand va ṣāḥeb-e menbar va meḥrāb, ʿarf-e ḥaqīqat rā ʿalā
qadr-e samm-e ebre nayāftand, che agar yāfte būdand īn ẓolm-e jadīd-e lā
methla lah dar ʿālam vāqeʿ nemīshod. Maḥż-e havā va rīyāsat takallom
namūdand be-ānche ke sabab va ʿellat-e shahādat-e Maqṣūd-e ʿālam,
Noqṭe-ye Ūlā, rūḥ-e mā sevāh fedāh, gasht, che agar ān adhkār-e majʿūle
nabūd Ḥaqq va owlīyāyash shahīd nemīshodand. Jamīʿ-e ʿolamā dar sane-ye
avval dhekr namūdand ke jamīʿ-e ʿolamā az qabl va baʿd bar īn būde-and
ke ān Ḥażrat bāyad az Jābolqā va Jābolsā ẓāher shavad va ḥāl īn ḥezb
mīgūyand dar Shīrāz motavalled shod. Ledhā jamīʿ ḥokm bar kofr va safk-e
dam-e moṭahhar namūdand.

Now a fresh gathering of the world's liars has assembled, and they have
been, and remain, intent on reconstructing the former foundations. The
clamorer of the land of Kaf and Ra and the other clamorers intend to
veil the light of the Sun of Truth with the clouds of new conjectures,
and to spend another two or three thousand years rejecting and
accepting, turning away and objecting, reviling and cursing one another
from pulpits and in assemblies. I swear by the Sun of the horizon of
holiness, this servant is ashamed and abashed to mention these things.
All the divines of the age, each of whom was a leader in Iran and
occupied a pulpit and a prayer niche, failed to perceive even a needle's
eye's measure of the fragrance of truth. Had they perceived it, this new
and unparalleled oppression would not have occurred in the world. Out of
sheer desire and ambition for leadership, they uttered what became the
cause and occasion of the martyrdom of the Desire of the world, the
Primal Point---may the souls of all besides Him be sacrificed for Him.
Had it not been for those fabricated statements, God and His loved ones
would not have been martyred. In the first year, all the divines
declared that the divines, earlier and later alike, had held that He
must appear from Jabulqa and Jabulsa, whereas this community now said
that He had been born in Shiraz. Thus they all pronounced Him an
unbeliever and decreed the shedding of His pure blood.

\subsection{22}\label{section-1997}

\begin{otherlanguage}{persian}

یا حبیب فؤادی قدری تفکّر لازم است تا بر ظلم این فئه مطّلع شوید و بعضی را
آگاه نمائید اکثری از عباد بتقلید متمسّکند و بهوی مبتلا سیل ظنون و
اوهامشان بنیان ایمان را ویران نمود افّ لهم و للّذین اتّبعوهم غفلت بمقامی
رسیده که باوهامات خود شرایط از برای ظهور ترتیب داده‌اند و به اشارات و
ادراکات خود او را محدود نموده‌اند و چون مخالف ظنون خود یافته‌اند ضوضاء
مرتفع بگوئید

\end{otherlanguage}

Yā ḥabība fuʾādī, qadrī tafakkor lāzem ast tā bar ẓolm-e īn feʾe
moṭṭaleʿ shavīd va baʿżī rā āgāh namāyīd. Aktharī az ʿebād be-taqlīd
motamassekand va be-havā mobtalā. Seyl-e ẓonūn va owhāmeshān bonyān-e
īmān rā vīrān namūd. Uffin lahum wa-li-lladhīna ittabaʿūhum. Ghaflat
be-maqāmī resīde ke be-owhāmāt-e khod sharāyeṭ az barā-ye ẓohūr tartīb
dāde-and va be-eshārāt va edrākāt-e khod ū rā maḥdūd namūde-and va chon
mokhālef-e ẓonūn-e khod yāfte-and ḍowḍāʾ mortafeʿ. Begūyīd:

O beloved of my heart, some reflection is necessary, so that you may
become aware of the oppression of this faction and alert others. Most
people cling to imitation and are afflicted by desire. The flood of
their conjectures and delusions has destroyed the foundation of faith.
Fie upon them and those who have followed them! Their heedlessness has
reached such a degree that, through their own imaginings, they have
devised conditions for the Revelation and have confined Him within their
own allusions and conceptions. When they found Him contrary to their
conjectures, an uproar arose. Say:

\subsection{23}\label{section-1998}

\begin{otherlanguage}{arabic}

این نبأ عظیم در اثبات امرش محتاج بتصدیق احدی نبوده و نیست فانظروا ما نطق
به نقطة البیان روح من فی الامکان فداه قوله تبارک و تعالی و قد کتبت جوهرة
فی ذکره و هو انّه لا یستشار باشارتی و لا بما ذکر فی البیان الی قوله جلّ و
عزّ فانّه اجلّ و اعلی من ان یکون معروفا بدونه او مستشیراً باشارة خلقه و انّنی
انا اوّل عبد قد آمنت به و یآیاته و اخذت من ابکار حدائق جنّة عرفانه حدائق
کلماته بلی وعزّته هو الحقّ لا اله الّا هو کلّ بأمره قائمون انتهی

\end{otherlanguage}

Īn Nabaʾ-e ʿAẓīm dar ethbāt-e amrash moḥtāj be-taṣdīq-e aḥadī nabūde va
nīst. Fa-nẓurū mā naṭaqa bihi Nuqṭatu al-Bayān, rūḥu man fī al-imkāni
fidāhu. Qawluhu tabāraka wa-taʿālā: wa-qad katabtu jawharatan fī
dhikrihi wa-huwa annahu lā yustashāru bi-ishāratī wa-lā bi-mā dhukira fī
al-Bayān. Ilā qawlihi jalla wa-ʿazza: fa-innahu ajallu wa-aʿlā min an
yakūna maʿrūfan bi-dūnihi aw mustashīran bi-ishārati khalqihi, wa-innanī
anā awwalu ʿabdin qad āmantu bihi wa-{[}corrupted text{]} wa-akhadhtu
min abkāri ḥadāʾiqi jannati ʿirfānihi ḥadāʾiqa kalimātihi. Balā
wa-ʿizzatihi, huwa al-Ḥaqqu lā ilāha illā huwa, kullun bi-amrihi
qāʾimūn. Intahā.

This Great Announcement has never needed, nor does He now need, anyone's
confirmation to establish His Cause. Consider, then, what the Point of
the Bayan---may the souls of all within the realm of existence be
sacrificed for Him---has uttered. His words, blessed and exalted be He:
``I have written a jewel concerning His mention, namely that He is not
to be consulted through My allusion, nor through what has been mentioned
in the Bayan.'' And further on, His words, glorified and mighty be He:
``For He is too glorious and exalted to be known through anyone other
than Himself or to seek counsel through the allusion of His creation. I
am indeed the first servant to have believed in Him and {[}corrupted
original text{]}, and from the pristine gardens of the paradise of His
recognition I have taken the gardens of His words. Yes, by His might! He
is the Truth; there is no god but Him. All subsist through His
command.'' End of quotation.

\subsection{24}\label{section-1999}

\begin{otherlanguage}{persian}

اَین الآذان الواعیه و اَین الأبصار الحدیده تا اشراقات انوار آفتاب بیانرا
مشاهده نمایند و ندای حقّ جلّ جلاله را استماع کنند اگر جمیع من فی البیان در
همین یکفقره از روی انصاف و عدل ملاحظه و تفکّر نمایند کلّ بشطر اقدس توجّه
کنند و از ما عندالنّاس غافل و بما عنداللّه متمسّک گردند ولکن چه فائده که
آذان را قطن اغراض از اصغا منع نموده و ابصار را حجابهای ظنون و اوهام از
مشاهده باز داشته

\end{otherlanguage}

Ayna al-ādhānu al-wāʿiya wa-ayna al-abṣāru al-ḥadīda, tā eshrāqāt-e
anvār-e āftāb-e bayān rā moshāhede namāyand va nedā-ye Ḥaqq, jalla
jalāluhu, rā estemāʿ konand? Agar jamīʿ-e man fī al-Bayān dar hamīn yek
feqre az rū-ye enṣāf va ʿadl molāḥeẓe va tafakkor namāyand, koll
be-shaṭr-e aqdas tavajjoh konand va az mā ʿenda al-nās ghāfel va be-mā
ʿenda Allāh motamassek gardand. Va lāken che fāʾede ke ādhān rā qoṭn-e
aghrāż az eṣghā manʿ namūde va abṣār rā ḥejābhā-ye ẓonūn va owhām az
moshāhede bāz dāshte.

Where are the attentive ears, and where are the keen eyes, that they may
behold the radiance of the light of the Sun of Utterance and hear the
call of God, exalted be His glory? If all the people of the Bayan were
to consider and reflect on this one passage with fairness and justice,
they would all turn toward the Most Holy, disregard what is with people,
and cling to what is with God. But what use is it, when the cotton of
self-interest has prevented ears from hearing, and the veils of
conjectures and delusions have kept eyes from seeing?

\subsection{25}\label{section-2000}

\begin{otherlanguage}{arabic}

وقتی از اوقات اینکلمهٴ علیا از مشرق وحی مالک اسماء ظاهر و نازل قوله
تبارک و تعالی ای عبد حاضر ملاحظه کن پلک باین رقیقی که مشاهده میشود اگر
حجاب گردد چشم از مشاهدهٴ آسمان و زمین و آنچه در اوست از اشجار و انهار و
بحار و جبال و اثمار و الوان و ما عندهم محروم و ممنوع ماند حال اگر حجاب
غلیظ ظنون بر قلب وارد شود چه خواهد نمود ای عبد حاضر از حقّ بطلب کلّ را از
حجبات مانعه مقدّس فرماید که شاید بمقصود از خلقت خلق پی برند و بعرفان کنز
مکنون مختوم فائز شوند و از این رحیق مختوم بیاشامند انتهی

\end{otherlanguage}

Vaqtī az owqāt īn kalame-ye ʿolyā az mashreq-e vaḥy-e Mālek-e Asmā ẓāher
va nāzel. Qawluhu tabāraka wa-taʿālā: ey ʿabd-e ḥāżer, molāḥeẓe kon,
pelk be-īn raqīqī ke moshāhede mīshavad agar ḥejāb gardad, chashm az
moshāhede-ye āsmān va zamīn va ānche dar ūst az ashjār va anhār va beḥār
va jebāl va athmār va alvān va mā ʿendahom maḥrūm va mamnūʿ mānad. Ḥāl
agar ḥejāb-e ghalīẓ-e ẓonūn bar qalb vāred shavad che khāhad namūd? Ey
ʿabd-e ḥāżer, az Ḥaqq beṭalab koll rā az ḥojobāt-e māneʿe moqaddas
farmāyad ke shāyad be-maqṣūd az khelqat-e khalq pey barand va
be-ʿerfān-e kanz-e maknūn-e makhtūm fāʾez shavand va az īn raḥīq-e
makhtūm bīyāshāmand. Intahā.

On one occasion, this exalted utterance appeared and was revealed from
the Dawning-Place of the revelation of the Lord of Names. His words,
blessed and exalted be He: ``O servant in attendance, consider: the
eyelid, thin as it is seen to be, if it becomes a veil, leaves the eye
deprived and prevented from beholding the heavens and the earth and all
they contain---trees, rivers, seas, mountains, fruits, colors, and all
that is with them. Now, if the thick veil of conjectures should fall
upon the heart, what will it do? O servant in attendance, ask God to
purify everyone of obstructing veils, so that perhaps they may discern
the purpose of the creation of humankind, attain recognition of the
hidden, sealed Treasure, and drink of this sealed choice wine.'' End of
quotation.

\subsection{26}\label{section-2001}

\begin{otherlanguage}{persian}

باری جوهر بیان بتصدیق محبوب عالمیان آنکه مشرق آیات و مطلع بیّنات باشارهٴ
آنسلطان وجود و بما نزّل فی البیان منوط نه و کلّ مأمورند او را بنفس او
بشناسند در اینمقام بعضی از اسرار از لسان مطلع اسماء شنیده شد که اینعبد
من‌غیر اذن جسارت بر عرض آن ندارد یا لیت القوم یعلمون و یعرفون طوبی از
برای نفسیکه ببصر خود بافق اعلی توجّه نماید و انوار وجه را من‌غیر سبحات و
اشارات ادراک کند

\end{otherlanguage}

Bārī, Jowhar-e Bayān be-taṣdīq-e Maḥbūb-e ʿālamīyān ānke Mashreq-e Āyāt
va Maṭlaʿ-e Bayyenāt be-eshāre-ye ān Solṭān-e Vojūd va be-mā nuzzila fī
al-Bayān manūṭ na va koll maʾmūrand ū rā be-nafs-e ū beshenāsand. Dar īn
maqām baʿżī az asrār az lesān-e Maṭlaʿ-e Asmā shenīde shod ke īn ʿabd
men gheyr-e edhn jasārat bar ʿarż-e ān nadārad. Yā layta al-qawma
yaʿlamūna wa-yaʿrifūn. Ṭūbā az barā-ye nafsī ke be-baṣar-e khod
be-Ofoq-e Aʿlā tavajjoh namāyad va anvār-e vajh rā men gheyr-e soboḥāt
va eshārāt edrāk konad.

In short, the essence of the Bayan, as confirmed by the Beloved of all
the worlds, is that the Dawning-Place of the Verses and the Dayspring of
Clear Proofs is not dependent on the allusion of that Sovereign of
existence or on what was revealed in the Bayan, and that all are
commanded to recognize Him through His own Self. Concerning this
station, certain mysteries were heard from the tongue of the Dayspring
of Names, which this servant does not dare to relate without permission.
If only the people knew and recognized! Blessed is the soul who turns
toward the Most Exalted Horizon with his own eyes and perceives the
light of the Face without veils of glory or allusions.

\subsection{27}\label{section-2002}

\begin{otherlanguage}{persian}

سکر بادهٴ غرور و غفلت اکثری را محجوب نموده ساذج قدمی را که نقطهٴ بیان
روح ما سواه فداه میفرماید انّنی انا اوّل عبد قد آمنت به و بآیاته در باره‌اش
گفته‌اند آنچه را که آسمان از او منفطر و ملکوت اسماء از آن مضطرب فی‌الحقیقه
اگر نفسی در اینکلمات مبارکه که از مشرق بیان حقّ جلّ جلاله اشراق نموده تفکّر
نماید و بر لآلی بیان مطّلع گردد خود را غنی و مستغنی از کلّ مشاهده کند
بپرهای اطمینان در هواء قدس رحمن طیران نماید و از کوثر استقامت بیاشامد
امروز اگر نفسی اقلّ من آن توقّف کند از اهل نار محسوب عند ذکرش جمیع اذکار
مفقود و عند بیانش جمیع بیان خاضع

\end{otherlanguage}

Sokr-e bāde-ye ghorūr va ghaflat aksarī rā maḥjūb namūde. Sādhaj-e
qidamī rā ke Noqṭe-ye Bayān, rūḥ-e mā sevāh fedāh, mīfarmāyad: ``Innanī
anā awwalu ʿabdin qad āmantu bihi wa-bi-āyātihi,'' dar bāre-ash
gofte-and ānche rā ke āsmān az ū monfaṭer va malakūt-e asmā az ān
możṭareb. Fī'l-ḥaqīqe agar nafsī dar īn kalemāt-e mobārake ke az
mashreq-e bayān-e Ḥaqq, jalla jalāluhu, eshrāq namūde tafakkor namāyad
va bar laʾālī-ye bayān moṭṭaleʿ gardad, khod rā ghanī va mostaghnī az
koll moshāhede konad, be-parhā-ye eṭmīnān dar havā-ye qods-e Raḥmān
ṭayarān namāyad va az kawthar-e esteqāmat biyāshāmad. Emrūz agar nafsī
aqall-e men ān tavaqqof konad, az ahl-e nār maḥsūb. ʿEnd-e dhekrash
jamīʿ-e adhkār mafqūd va ʿend-e bayānash jamīʿ-e bayān khāżeʿ.

Intoxication with the wine of pride and heedlessness has veiled the
majority. Of Him who is the pure Essence of eternity, concerning whom
the Point of the Bayan---may the spirit of all else be sacrificed for
Him---says, ``I am indeed the first servant to have believed in Him and
His verses,'' they have said things that cause heaven to split apart and
the kingdom of names to tremble. In truth, if anyone reflects on these
blessed words that have shone from the Dayspring of the utterance of
God, exalted be His glory, and discovers the pearls of His utterance, he
will find himself rich and independent of all, soar on the wings of
assurance in the holy atmosphere of the All-Merciful, and drink from the
river of steadfastness. Today, if anyone hesitates for less than an
instant, he is counted among the people of the fire. At His mention, all
other mentions vanish, and before His utterance, every utterance bows.

\subsection{28}\label{section-2003}

\begin{otherlanguage}{persian}

وقتی از اوقات فرمودند یا عبد حاضر امروز ملکوت بیان لدی الرّحمن خاضع و
خاشع جمیع کتب و الواح باراده معلّق و منوط و جمیع من فی البیان باین ذکر
اعظم و نبأ عظیم و عرفانش مأمور طوبی از برای نفسیکه سلاسل اوهام و ظنون را
باسم حضرت قیّوم شکست از ما عندالخلق گذشت و بما عندالحقّ تمسّک جست انتهی

\end{otherlanguage}

Vaqtī az awqāt farmūdand: ``Yā ʿabd-e ḥāżer, emrūz malakūt-e bayān ladā
al-Raḥmān khāżeʿ va khāsheʿ. Jamīʿ-e kotob va alvāḥ be-erāde moʿallaq va
manūṭ va jamīʿ-e man fī al-Bayān be-īn Dhekr-e Aʿẓam va Nabaʾ-e ʿAẓīm va
ʿerfānash maʾmūr. Ṭūbā az barā-ye nafsī ke salāsel-e awhām va ẓonūn rā
be-esm-e Ḥażrat-e Qayyūm shekast, az mā ʿenda al-khalq godhasht va be-mā
ʿenda al-Ḥaqq tamassok jost.'' Entahā.

On one occasion He said: ``O servant in attendance, today the kingdom of
utterance bows in humble submission before the All-Merciful. All books
and tablets depend upon His will and are subject to it, and all who are
in the Bayan are commanded to turn to this Most Great Remembrance and
momentous Announcement and to recognize Him. Blessed is the soul who,
through the name of the Self-Subsisting, has broken the chains of vain
imaginings and conjectures, passed beyond what is with the people, and
held fast to what is with God.'' End of quotation.

\subsection{29}\label{section-2004}

\begin{otherlanguage}{persian}

امثال آن نفوس معرضه لازال بوده‌اند اگر نفسی بعد از ارتقا حضرت عسکری علیه
سلام اللّه در امورات واقعه ودسایس علمای زنادقه تفکّر نماید شهادت میدهد بر
اینکه جمیع فساد عالم از آن نفوس بوده و هست لعمر ربّنا اگر کذب رؤسا نبود
ربّ اعلی بظلم اعدا شهید نمیشد حضرت موهومی بر سریر موهومی معیّن نمودند و
بروایات موهومه ثابت کردند آنچه را که سبب سفک دم سیّد عالم و جواهر امم شد
و گاهی تواقیع جعلیّه ظاهر مینمودند و نسبتش را بمبدء وجود میدادند الا لعنة
اللّه علی القوم الکاذبین در باطن مقصود اعلاء علوّ و سموّ خودشان بود ولکن در
ظاهر بخرائر اشاره مینمودند گاهی جابلقا و هنگامی جابلسا و وقتی ناحیهٴ
کذبه که بمقدّسه مینامیدند سلطان لایزالی را که نقطهٴ بیان روح ما سواه فداه
در ذکرش میفرماید قوله عزّ ذکره انّه اجلّ و اعلی من ان یکون معروفاً بدونه او
مستشیراً باشارة خلقه

\end{otherlanguage}

Amthāl-e ān nofūs-e moʿreże lā yazāl būde-and. Agar nafsī baʿd az
erteqā-ye Ḥażrat-e ʿAskarī, ʿalayhi salāmu Allāh, dar omūrāt-e vāqeʿe va
dasāyes-e ʿolamā-ye zanādeqe tafakkor namāyad, shahādat mīdahad bar īnke
jamīʿ-e fasād-e ʿālam az ān nofūs būde va hast. La-ʿamri Rabbinā, agar
kedhb-e roʾasā nabūd, Rabb-e Aʿlā be-ẓolm-e aʿdā shahīd nemīshod.
Ḥażrat-e mawhūmī bar sarīr-e mawhūmī moʿayyan namūdand va be-revāyāt-e
mawhūme thābet kardand ānche rā ke sabab-e safk-e dam-e Sayyed-e ʿālam
va javāher-e omam shod, va gāhī tavāqīʿ-e jaʿliyye ẓāher mīnamūdand va
nesbatash rā be-mabdaʾ-e vojūd mīdādand. A-lā laʿnatu Allāhi ʿalā
al-qawmi al-kādhibīn. Dar bāṭen maqṣūd eʿlā-ye ʿolovv va somovv-e
khodeshān būd, va-lāken dar ẓāher be-{[}corrupted text{]} eshāre
mīnamūdand, gāhī Jābolqā va hengāmī Jābolsā va vaqtī nāḥiye-ye kedhbe ke
be-moqaddase mīnāmīdand. Solṭān-e lā yazālī rā ke Noqṭe-ye Bayān, rūḥ-e
mā sevāh fedāh, dar dhekrash mīfarmāyad, qawluhu ʿazza dhikruhu:
``Innahu ajallu wa-aʿlā min an yakūna maʿrūfan bi-dūnihi aw mustashīran
bi-ishārati khalqihi.''

There have always been people like those who turn away. If anyone
reflects on the events following the passing of His Holiness ʿAskari,
upon him be the peace of God, and on the intrigues of the godless
religious scholars, he will testify that all the corruption in the world
has come, and continues to come, from such people. By the life of our
Lord! Were it not for the lies of the leaders, the Most Exalted Lord
would not have been martyred through the oppression of His enemies. They
appointed an imaginary personage to an imaginary throne and, through
imaginary traditions, established what led to the shedding of the blood
of the Lord of the world and the jewels of the nations. At times they
produced forged rescripts and attributed them to the Source of
existence. Indeed, may God's curse rest upon the people who lie!
Inwardly, their aim was to enhance their own eminence and exaltation,
but outwardly they pointed to {[}corrupted original text{]}---sometimes
Jabulqa, sometimes Jabulsa, and sometimes the lying seat that they
called holy. As for the everlasting Sovereign, the Point of the
Bayan---may the spirit of all else be sacrificed for Him---says
concerning Him, exalted be His remembrance: ``He is too glorious and
exalted to be known through anyone other than Himself or to seek counsel
through the indications of His creatures.''

\subsection{30}\label{section-2005}

\begin{otherlanguage}{arabic}

حال نفوس غافله بر اعراض قیام نموده‌اند و بقول این و آن تمسّک جسته‌اند امروز
لسان عظمت میفرماید قوله عزّ اجلاله ان یا قلم هل تری الیوم غیری این
الأشیآء و ظهوراتها و این الأسمآء و ملکوتها و البواطن و اسرارها و الظّواهر
و آثارها قد اخذ الفنآء من فی الانشآء و هذا وجهی الباقی المشرق المنیر
انتهی

\end{otherlanguage}

Ḥāl nofūs-e ghāfele bar eʿrāż qiyām namūde-and va be-qawl-e īn va ān
tamassok joste-and. Emrūz Lesān-e ʿAẓamat mīfarmāyad, qawluhu ʿazza
ijlāluhu: ``An yā qalamu, hal tarā al-yawma ghayrī? Ayna al-ashyāʾu
wa-ẓuhūrātuhā, wa-ayna al-asmāʾu wa-malakūtuhā, wa-al-bawāṭinu
wa-asrāruhā, wa-al-ẓawāhiru wa-āthāruhā? Qad akhadha al-fanāʾu man fī
al-inshāʾi, wa-hādhā wajhī al-bāqī al-mushriqu al-munīr.'' Entahā.

Now heedless souls have arisen in opposition and have clung to the words
of this person and that. Today the Tongue of Grandeur declares, exalted
be His majesty: ``O Pen, do you see anyone but Me today? Where are all
things and their manifestations? Where are the names and their kingdom,
the inward realities and their mysteries, the outward forms and their
effects? Annihilation has overtaken all who dwell in creation, and this
is My everlasting, radiant, luminous Face.'' End of quotation.

\subsection{31}\label{section-2006}

\begin{otherlanguage}{arabic}

حال ملاحظه نمائید بچه متشبّثند و بچه حبال متمسّک اراده دارند یک اسباب
موهومهٴ جدیده‌ئی مثل اوهامات قبل از برای ناس بیچاره برپا نمایند مکرّر عرض
مینمایم اعاذنا اللّه یا معشر الموحّدین من شرّ هؤلآء

\end{otherlanguage}

Ḥāl molāḥeẓe namāyīd be-che motashabbethand va be-che ḥebāl motamassek.
Erāde dārand yek asbāb-e mawhūme-ye jadīdeʾī methl-e awhāmāt-e qabl az
barā-ye nās-e bīchāre barpā namāyand. Mokarrar ʿarż mīnamāyam:
``Aʿādhanā Allāhu, yā maʿshara al-muwaḥḥidīn, min sharri hāʾulāʾ.''

Now consider what they cling to and what cords they grasp. They intend
to set up a new imaginary apparatus for the helpless people, like the
vain imaginings of the past. Again I say: ``May God protect us, O
company of those who affirm His oneness, from the evil of these
people.''

\subsection{32}\label{section-2007}

\begin{otherlanguage}{persian}

باری الیوم هیچ کتابی از کتب و هیچ لوحی از الواح و هیچ صحیفه‌ئی از صحف
کفایت نمینماید سبحان‌اللّه عمّا یقولون الظّالمون الیوم کلّ باید بافق ظهور و
ما یظهر من عنده ناظر باشند اگر به اوهامات و ظنونات و روایات انسان مشغول
شود البتّه عاقبت آن بر هلاکت است چنانچه دیده شد

\end{otherlanguage}

Bārī al-yawm hīch ketābī az kotob va hīch lawḥī az alvāḥ va hīch
ṣaḥīfeʾī az ṣoḥof kefāyat nemīnamāyad. Subḥāna Allāhi ʿammā yaqūlu
al-ẓālimūn. Al-yawm koll bāyad be-ofoq-e ẓohūr va mā yaẓharu min ʿindihi
nāẓer bāshand. Agar be-awhāmāt va ẓonūnāt va revāyāt ensān mashghūl
shavad, albatte ʿāqebat-e ān bar halākat ast, chenānche dīde shod.

In short, today no book among the books, no tablet among the tablets,
and no scripture among the scriptures suffices. Glorified be God above
what the wrongdoers say! Today all must look toward the horizon of the
Revelation and whatever appears from Him. If a person occupies himself
with vain imaginings, conjectures, and traditions, his end will
assuredly be destruction, as has been seen.

\subsection{33}\label{section-2008}

\begin{otherlanguage}{arabic}

سبحانک یا اله البرّ و البحر ترانی قائماً بین یدیک و ایادی رجای مرتفعة الی
سمآء جودک و فضلک اسئلک بسلطانک المهیمن علی الکائنات و اسمک الّذی به سخّرت
الممکنات بأن تؤیّد عبادک علی العدل و الانصاف و وفّقهم علی التّوجّه الیک
ببصرک و النّظر الی شطرک بنورک انت الّذی باسمک ماجت بحور الفضل و العطآء و
بمشیّتک تزیّنت سمآء کرمک یا مالک الأسمآء لا اله الّا انت المقتدر المختار

\end{otherlanguage}

Subḥānaka yā Ilāha al-barri wa-al-baḥr. Tarānī qāʾiman bayna yadayka
wa-ayādī rajāya murtafiʿatan ilā samāʾi jūdika wa-faḍlika. Asʾaluka
bi-sulṭānika al-muhaymini ʿalā al-kāʾināti wa-ismika alladhī bihi
sakhkharta al-mumkināti bi-an tuʾayyida ʿibādaka ʿalā al-ʿadli
wa-al-inṣāfi wa-waffiqhum ʿalā al-tawajjuhi ilayka bi-baṣarika
wa-al-naẓari ilā shaṭrika bi-nūrika. Anta alladhī bi-ismika mājat buḥūru
al-faḍli wa-al-ʿaṭāʾi wa-bi-mashīʾatika tazayyanat samāʾu karamika yā
Mālika al-asmāʾ. Lā ilāha illā Anta al-Muqtadiru al-Mukhtār.

Glory be to You, O God of land and sea! You see me standing before You,
with the hands of my hope raised toward the heaven of Your bounty and
grace. I ask You, by Your sovereignty that overshadows all created
things and by Your name through which You have subdued all contingent
beings, to aid Your servants to act with justice and fairness. Enable
them to turn toward You with Your sight and to look in Your direction
through Your light. You are the One through whose name the seas of grace
and bounty have surged, and through whose will the heaven of Your
generosity has been adorned, O Lord of Names. There is no God but You,
the All-Powerful, who chooses as He wills.

\subsection{34}\label{section-2009}

\begin{otherlanguage}{persian}

اگر الیوم نفسی به غیر حقّ و بیان حقّ ناظر باشد البتّه از صراط منحرف شود و
البتّه بسقر مقرّ گیرد میفرماید جمیع بیان ورقی است از اوراق جنّت او اگر
بقبول فائز شود سبحان‌اللّه امر در چه مقام و خراطین در چه طین این خادم فانی
متحیّر واللّه المهیمن القیّوم از اوّل امر تا حین خیالی جز رضای حقّ نداشته و
ندارم حال پیر شده‌ام و قوی ضعیف گشته مع‌ذلک و مع شغل عظیم خالصاً لوجه اللّه
باین نوشته مشغول شده‌ام حقّ گواه و آگاه که جز رضایش را طالب نبوده و نیستم
لوجهه گفته و لوجهه میگویم من شآء فلیقبل و من شآء فلیعرض انّ ربّنا الرّحمن
لهو الغنیّ الحمید

\end{otherlanguage}

Agar al-yawm nafsī be-ghayr-e Ḥaqq va bayān-e Ḥaqq nāẓer bāshad, albatte
az ṣerāṭ monḥaref shavad va albatte be-Saqar maqarr gīrad. Mīfarmāyad:
``Jamīʿ-e Bayān varaqī ast az awrāq-e jannat-e ū agar be-qabūl fāʾez
shavad.'' Subḥāna Allāh, amr dar che maqām va kharāṭīn dar che ṭīn. Īn
khādem-e fānī motaḥayyer. Wa-Allāhi al-Muhaymini al-Qayyūm, az avval-e
amr tā ḥīn khayālī joz reżā-ye Ḥaqq nadāshte va nadāram. Ḥāl pīr
shode-am va qovā żaʿīf gashte. Maʿa dhālek va maʿa shoghl-e ʿaẓīm,
khāliṣan li-wajhi Allāh be-īn neveshte mashghūl shode-am. Ḥaqq govāh va
āgāh ke joz reżāyash rā ṭāleb nabūde va nīstam. Li-wajhihi gofte va
li-wajhihi mīgūyam. Man shāʾa fa-l-yuqbil wa-man shāʾa fa-l-yuʿriḍ. Inna
Rabbanā al-Raḥmāna la-huwa al-Ghaniyyu al-Ḥamīd.

If today anyone looks to anything other than God and His utterance, he
will assuredly stray from the path and assuredly make his abode in hell.
He says: ``The entire Bayan is but a leaf among the leaves of His
paradise, if it attains His acceptance.'' Glory be to God! How lofty is
the station of the Cause, and in what mud do these earthworms dwell!
This mortal servant is bewildered. By God, the All-Protecting, the
Self-Subsisting! From the beginning of the Cause until now, I have had,
and still have, no thought but God's good pleasure. Now I have grown old
and my strength has waned. Nevertheless, despite my pressing duties, I
have occupied myself with this writing solely for the sake of God. God
is witness and knows that I have sought, and still seek, nothing but His
good pleasure. For His sake I have spoken, and for His sake I speak. Let
whoever wishes turn toward Him, and whoever wishes turn away. Our Lord,
the All-Merciful, is indeed the Self-Sufficient, the Praiseworthy.

\subsection{35}\label{section-2010}

\begin{otherlanguage}{persian}

و همچنین در اینکلمهٴ علیا که از مطلع بیان مالک اسماء یعنی نقطهٴ اولی روح
ما سواه فداه نازل شده تفکّر فرمائید و فی‌الحقیقه اینکلمه از جوهر بیان
محسوبست قوله تعالی و فی سنة التّسع انتم کلّ خیر تدرکون و فی سنة التّسع انتم
بلقآء اللّه ترزقون و مخاطباً لحضرة العظیم علیه من کلّ بهآء ابهاه و من کلّ
علآء اعلاه میفرماید قوله تعالی

\end{otherlanguage}

Va hamchenīn dar īn kaleme-ye ʿolyā ke az maṭlaʿ-e bayān-e Mālek-e Asmā,
yaʿnī Noqṭe-ye Ūlā, rūḥ-e mā sevāh fedāh, nāzel shode tafakkor farmāyīd.
Va fī'l-ḥaqīqe īn kaleme az jawhar-e bayān maḥsūb ast. Qawluhu taʿālā:
``Wa-fī sanati al-tisʿi antum kulla khayrin tudrikūn, wa-fī sanati
al-tisʿi antum bi-liqāʾi Allāhi turzaqūn.'' Va mokhāṭaban le-Ḥażrat-e
ʿAẓīm, ʿalayhi min kulli bahāʾin abhāhu wa-min kulli ʿalāʾin aʿlāhu,
mīfarmāyad, qawluhu taʿālā:

Reflect also on this exalted utterance revealed from the Dayspring of
the utterance of the Lord of Names, namely the Primal Point---may the
spirit of all else be sacrificed for Him. In truth, this utterance is to
be counted among the very essence of His words. He, exalted be He, says:
``In the year nine you will attain all good, and in the year nine you
will be granted the presence of God.'' And, addressing His Honor
ʿAzim---upon him be the most resplendent of all glory and the loftiest
of all exaltation---He, exalted be He, says:

\subsection{36}\label{section-2011}

\begin{otherlanguage}{persian}

هذا ما قد وعدناک قبل حین الّذی اجبناک اصبر حتّی یقضی عن البیان تسعةً فاذا
قل فتبارک اللّه احسن المبدعین حال انصاف دهید در سنهٴ تسع کِه ظاهر شد این
الآذان و این الأبصار باری مفرّی از برای احدی نبوده و نیست الّا ان یتوب و
یرجع الی اللّه المهیمن القیّوم مگر از سماء برهان و شمس ایقان و بحر اعظم و
نبأ عظیم چشم پوشد از برای چنین نفسی آیات الهی و بیّنات ربّانی ثمری نداشته
و ندارد هر صاحب انصافی شهادت میدهد که امر اللّه باین ظهور اعظم ظاهر در
ایّامیکه جمیع بحفظ جان مشغول و در زاویها مستور و محجوب حقّ جلّ جلاله بنفسه
قائم و ناطق چنانچه رشحات بحر بیانش عالم را اخذ نموده و آثارش آفاق را
احاطه فرموده طوبی از برای منصفین و سامعین

\end{otherlanguage}

``Hādhā mā qad waʿadnāka qablu ḥīna alladhī ajabnāka. Iṣbir ḥattā yuqḍā
ʿani al-Bayāni tisʿatan, fa-idhā qul fa-tabāraka Allāhu aḥsanu
al-mubdiʿīn.'' Ḥāl enṣāf dahīd, dar sane-ye tesʿ keh ẓāher shod? Ayna
al-ādhānu wa-ayna al-abṣār? Bārī mafarrī az barā-ye aḥadī nabūde va nīst
illā an yatūba wa-yarjiʿa ilā Allāhi al-Muhaymini al-Qayyūm, magar az
samā-ye borhān va shams-e īqān va baḥr-e aʿẓam va nabaʾ-e ʿaẓīm chashm
pūshad. Az barā-ye chonīn nafsī āyāt-e elāhī va bayyenāt-e rabbānī
thamarī nadāshte va nadārad. Har ṣāḥeb-e enṣāfī shahādat mīdahad ke Amr
Allāh be-īn ẓohūr-e aʿẓam ẓāher. Dar ayyāmī ke jamīʿ be-ḥefẓ-e jān
mashghūl va dar zāviyehā mastūr va maḥjūb, Ḥaqq, jalla jalāluhu,
be-nafsehe qāʾem va nāṭeq, chenānche rashaḥāt-e baḥr-e bayānash ʿālam rā
akhdh namūde va āthārash āfāq rā eḥāṭe farmūde. Ṭūbā az barā-ye monṣefīn
va sāmeʿīn.

``This is what We promised you before, at the time when We answered you.
Be patient until nine have elapsed from the Bayan; then say: `Blessed be
God, the best of originators!'\,'' Now judge fairly: who appeared in the
year nine? Where are the ears, and where are the eyes? In short, there
has never been, nor is there now, any refuge for anyone except to repent
and return to God, the All-Protecting, the Self-Subsisting, unless he
shuts his eyes to the heaven of proof, the sun of certitude, the Most
Great Ocean, and the momentous Announcement. For such a person the
divine verses and the Lord's clear proofs have been, and remain,
fruitless. Every fair-minded person testifies that the Cause of God is
manifest through this Most Great Revelation. In the days when all were
occupied with preserving their lives and were concealed and hidden in
corners, God, exalted be His glory, stood forth and spoke in His own
Person, so that the spray from the ocean of His utterance has
encompassed the world and His writings have embraced every horizon.
Blessed are the fair-minded and those who listen.

\subsection{37}\label{section-2012}

\begin{otherlanguage}{persian}

و همچنین نقطهٴ اولی روح ما سواه فداه میفرماید قل هو نبأ لم یحیط بعلمه
احد الّا اللّه ولکن انتم یومئذ لا تعلمون مع‌ذلک این قوم بی‌انصاف گفته‌اند
آنچه گفته‌اند این آن نبأ عظیمی است که در جمیع کتب مذکور و مسطور است و
احدی جز نفس حقّ بر حین ظهور مطّلع نه و آگاه نه اینست آن ساعتیکه حضرت روح
میفرماید و از او اخبار نموده و امّا ذلک الیوم و تلک السّاعة فلا یعلم بهما
احداً و لا الملائکة الّذین فی السّمآء و لا الابن الّا الأب انتهی

\end{otherlanguage}

Va hamchenīn Noqṭe-ye Ūlā, rūḥ-e mā sevāh fedāh, mīfarmāyad: ``Qul huwa
nabaʾun lam yuḥiṭ bi-ʿilmihi aḥadun illā Allāhu wa-lākin antum
yawmaʾidhin lā taʿlamūn.'' Maʿa dhālek īn qawm-e bī-enṣāf gofte-and
ānche gofte-and. Īn ān nabaʾ-e ʿaẓīmī ast ke dar jamīʿ-e kotob madhkūr
va masṭūr ast va aḥadī joz nafs-e Ḥaqq bar ḥīn-e ẓohūr moṭṭaleʿ na va
āgāh na. Īn ast ān sāʿatī ke Ḥażrat-e Rūḥ mīfarmāyad va az ū ekhbār
namūde: ``Wa-ammā dhālika al-yawmu wa-tilka al-sāʿatu fa-lā yaʿlamu
bihimā aḥadan wa-lā al-malāʾikatu alladhīna fī al-samāʾi wa-lā al-ibnu
illā al-ab.'' Entahā.

Likewise, the Primal Point---may the spirit of all else be sacrificed
for Him---says: ``Say: It is an announcement whose knowledge none but
God has encompassed, yet on that day you will not know.'' Despite this,
these unfair people have said what they have said. This is that
momentous Announcement mentioned and recorded in all the books, and no
one but God Himself knows or is aware of the time of its appearance.
This is the hour of which Jesus speaks and which He foretold: ``But of
that day and that hour no one knows, neither the angels in heaven nor
the Son, but only the Father.'' End of quotation.

\subsection{38}\label{section-2013}

\begin{otherlanguage}{persian}

و اینست آن یومیکه حضرت یوئیل در کتاب خود میفرماید لأنّ یوم الرّبّ عظیم و
مخوف جدّاً فمن یطیقه و همچنین حضرت ارمیا میفرماید آه لأنّ ذلک الیوم عظیم و
لیس مثله جمیع آنچه ظاهر شده در کتب مرسلین مذکور حتّی ذکر حصن که حال مقرّ
عرش است و به سجن اعظم موسوم چنانچه داود میفرماید من یقودنی الی المدینة
المحصّنه و مدینهٴ محصّنه عکّاست و حصنش معروف و مشهور اینست آن حصنیکه
میفرماید انوار آفتاب حقیقت بر او تجلّی فرموده و از افقش اشراق کرده و
همچنین بشارت داده‌اند یکدیگر و عباد این ارض را که حقّ در آخر ایّام در این
اراضی ظاهر میشود و تجلّی میفرماید

\end{otherlanguage}

Va īn ast ān yawmī ke Ḥażrat-e Yūʾīl dar ketāb-e khod mīfarmāyad:
``Li-anna yawma al-Rabbi ʿaẓīmun wa-makhūfun jiddan fa-man yuṭīquhu?''
Va hamchenīn Ḥażrat-e Ermiyā mīfarmāyad: ``Āh, li-anna dhālika al-yawma
ʿaẓīmun wa-laysa mithluhu.'' Jamīʿ-e ānche ẓāher shode dar kotob-e
morsalīn madhkūr, ḥattā dhekr-e ḥeṣn ke ḥāl maqarr-e ʿarsh ast va
be-Sejn-e Aʿẓam mawsūm, chenānche Dāvūd mīfarmāyad: ``Man yaqūdunī ilā
al-madīnati al-muḥaṣṣana?'' Va madīne-ye moḥaṣṣane ʿAkkā-st va ḥeṣnash
maʿrūf va mashhūr. Īn ast ān ḥeṣnī ke mīfarmāyad anvār-e Āftāb-e Ḥaqīqat
bar ū tajallī farmūde va az ofoqash eshrāq karde. Va hamchenīn beshārat
dāde-and yekdīgar va ʿebād-e īn arż rā ke Ḥaqq dar ākhar-e ayyām dar īn
arāżī ẓāher mīshavad va tajallī mīfarmāyad.

And this is the day of which Joel says in his book: ``For the day of the
Lord is great and very terrible; who can endure it?'' Likewise Jeremiah
says: ``Alas, for that day is great, and none is like it.'' Everything
that has appeared is mentioned in the books of the Messengers, even the
fortress that is now the seat of the throne and is called the Most Great
Prison, as David says: ``Who will lead me into the fortified city?'' The
fortified city is Akka, and its fortress is well known and renowned.
This is the fortress of which he says that the lights of the Sun of
Truth have shone upon it and risen from its horizon. Likewise, they gave
one another and the inhabitants of this land the glad tidings that God
would appear and reveal His splendor in these lands in the last days.

\subsection{39}\label{section-2014}

\begin{otherlanguage}{persian}

حضرت اشعیا میفرماید علی جبل عال اصعدی یا مبشّرة صهیون ارفعی صوتک بقوّة یا
مبشّرة اورشلیم ارفعی لا تخافی قولی لمدن یهوذا هو ذا الهک هو ذا السّیّد الرّبّ
بقوّة یأتی و ذراعة تحکم له صهیون و اورشلیم در این اراضی است و میفرماید
باسم جدید نامیده میشود چنانچه در ارض سرّ از لسان عظمت بواد النّبیل موسوم
گشته و اینست آن وادی که نداء جمیع انبیا در او مرتفع و به لبّیک اللّهمّ لبّیک
ناطق حضرت خلیل و کلیم و روح و خاتم انبیا روح من فی الوجود فداهم در این
اراضی ظاهر چه که این ارض بارض بطحا و حجاز متّصل میشود اینست که میفرماید
طوبی لمن هاجر الی عکّا قسم بآفتاب افق بیان که مقصود اینفانی اظهار علم و
فضل و اطّلاع نبوده و نیست بلکه لوجه اللّه عرض میشود که نفوس غافله سبب
احتجاب خلق نشوند چنانچه از قبل شده‌اند باری اینست اراضی مقدّسه که در فرقان
و کتب قبل مذکور و در کتاب عاموس نبی (ع) میفرماید

\end{otherlanguage}

Ḥażrat-e Eshʿiyā mīfarmāyad: ``ʿAlā jabalin ʿālin iṣʿadī yā
mubashshirata Ṣahyūn. Irfaʿī ṣawtaki bi-quwwatin yā mubashshirata
Ūrashalīm. Irfaʿī lā takhāfī. Qūlī li-muduni Yahūdhā: Huwa dhā Ilāhuki.
Huwa dhā al-Sayyidu al-Rabbu bi-quwwatin yaʾtī wa-dhirāʿuhu taḥkumu
lahu.'' Ṣahyūn va Ūrashalīm dar īn arāżī ast va mīfarmāyad be-esm-e
jadīd nāmīde mīshavad, chenānche dar Arż-e Serr az Lesān-e ʿAẓamat
be-Vād-e Nabīl mawsūm gashte. Va īn ast ān vādī ke nedā-ye jamīʿ-e
anbiyā dar ū mortafeʿ va be-``Labbayka Allāhumma labbayk'' nāṭeq.
Ḥażrat-e Khalīl va Kalīm va Rūḥ va Khātam-e Anbiyā, rūḥ-e man fī
al-vojūd fedāhom, dar īn arāżī ẓāher, che ke īn arż be-arż-e Baṭḥā va
Ḥejāz mottaṣel mīshavad. Īn ast ke mīfarmāyad: ``Ṭūbā li-man hājara ilā
ʿAkkā.'' Qasam be-Āftāb-e ofoq-e bayān ke maqṣūd-e īn fānī eẓhār-e ʿelm
va fażl va eṭṭelāʿ nabūde va nīst, balke li-wajhi Allāh ʿarż mīshavad ke
nofūs-e ghāfele sabab-e eḥtejāb-e khalq nashavand, chenānche az qabl
shode-and. Bārī īn ast arāżī-ye moqaddase ke dar Forqān va kotob-e qabl
madhkūr, va dar ketāb-e ʿĀmūs-e Nabī, ʿalayhi al-salām, mīfarmāyad:

Isaiah says: ``Go up onto a high mountain, O herald of Zion. Lift up
your voice with strength, O herald of Jerusalem. Lift it up, do not
fear. Say to the cities of Judah: `Behold your God! Behold, the Lord God
comes with power, and His arm rules for Him.'\,'' Zion and Jerusalem are
in these lands, and he says that it will be called by a new name, just
as in the Land of Mystery it was named the Noble Valley by the Tongue of
Grandeur. This is the valley in which the call of all the Prophets is
raised and which proclaims: ``Here I am, O God, here I am!'' Abraham,
Moses, Jesus, and the Seal of the Prophets---may the spirit of all who
exist be sacrificed for them---appeared in these lands, for this land is
connected to the land of Batha and the Hijaz. Thus He says: ``Blessed is
whoever migrates to Akka.'' I swear by the Sun of the horizon of
utterance that this mortal one's purpose has never been, nor is it now,
to display knowledge, learning, or erudition. Rather, these things are
set forth for the sake of God, so that heedless souls may not cause the
people to be veiled, as they have done before. In short, these are the
holy lands mentioned in the Quran and the earlier books, and in the book
of the Prophet Amos, peace be upon him, he says:

\subsection{40}\label{section-2015}

\begin{otherlanguage}{persian}

انّ الرّبّ یزمجر من صهیون و یعطی صوته من اورشلیم فتنوح مراعی الرّعاة و یبلبس
رأس الکرمل و این کرمل مقابل عکّا واقع و در کتب قبل بکَوْم اللّه موسوم و ذکر
آن از قلم اعلی جاری البتّه بعضی بآن فائز شده‌اند و از جمله امور محتومه‌ای
که باید خباء مجد بر او مرتفع شود و للّه الحمد مرتفع شد و از صهیون فرات
معانی جاری رغماً لمن نبذ الانصاف و اخذ الاعتساف

\end{otherlanguage}

Inna al-Rabba yuzamjiru min Ṣahyūn wa-yuʿṭī ṣawtahu min Ūrushalīm
fa-tanūḥu marāʿī al-ruʿāt wa-{[}corrupted text{]} raʾs al-Karmil va īn
Karmil muqābil-i ʿAkkā vāqiʿ va dar kutub-i qabl be-Kawm-i Allāh mawsūm
va dhikr-i ān az Qalam-i Aʿlā jārī albatteh baʿḍī be-ān fāʾiz shodeh-and
va az jumleh umūr-i maḥtūmeh-ī keh bāyad khibā-yi majd bar ū murtafiʿ
shavad va li-llāh al-ḥamd murtafiʿ shod va az Ṣahyūn Furāt-i maʿānī jārī
raghman li-man nabadha al-inṣāf wa-akhadha al-iʿtisāf.

``The Lord roars from Zion and gives forth His voice from Jerusalem; the
pastures of the shepherds mourn, and {[}corrupted original text{]} the
summit of Carmel.'' This Carmel lies opposite Akka and is called the
Mound of God in the earlier books. Its mention has flowed from the Most
Exalted Pen, and some have surely attained it. Among the things
irrevocably ordained is that the tabernacle of glory should be raised
upon it, and, praise be to God, it has been raised. From Zion flows the
Euphrates of meanings, despite those who have cast aside fairness and
embraced injustice.

\subsection{41}\label{section-2016}

\begin{otherlanguage}{persian}

عظمت ظهور بمقامیست که حضرت داود در ذکرش میفرماید علوّا الرّبّ الهنا و
اسجدوا عند موطیء قدمیه قدّوس هو موسی و هرون بین کهنته و صموئیل بین الّذین
یدّعون باسمه حال باید در این آیهٴ مبارکه بسیار تفکّر نمود امروز جمیع اشجار
منجذب و صخور و احجار مهتزّ و متحرّک و جبال از فرح اکبر در مرور چه که ساعة
موعود آمده و مکنون مخزون ظاهر گشته اگر این کلمهٴ علیا از کتب برداشته شود
دیگر بیانی باقی نه و ذکری مشهود نه لکن این بیان از برای متبصّرین و عارفین
است نه از برای هر همج رعاعی که قائد اوهام و ظنون است

\end{otherlanguage}

ʿAẓamat-i ẓuhūr be-maqāmīst keh Ḥaḍrat-i Dāvūd dar dhikrash mī-farmāyad
ʿAllū al-Rabba ilāhanā wa-usjudū ʿinda mawṭiʾ qadamayhi quddūsun huwa
Mūsā wa-Hārūn bayna kahanatihi wa-Ṣamūʾīl bayna alladhīna yadʿūna
bi-ismihi ḥāl bāyad dar īn āyeh-yi mubārakeh bisyār tafakkur namūd imrūz
jamīʿ-i ashjār munjadh ib va ṣukhūr va aḥjār muhtazz va mutaḥarrik va
jibāl az faraḥ-i akbar dar murūr cheh keh sāʿat-i mawʿūd āmadeh va
maknūn-i makhzūn ẓāhir gashteh agar īn kalameh-yi ʿulyā az kutub
bardāshteh shavad dīgar bayānī bāqī nah va dhikrī mashhūd nah lākin īn
bayān az barā-yi mutabaṣṣirīn va ʿārifīn ast nah az barā-yi har hamaj-i
raʿāʿī keh qāʾid-i awhām va ẓunūn ast.

The greatness of the Revelation is such that David says in speaking of
it: ``Exalt the Lord our God and bow down at His footstool. Holy is He.
Moses and Aaron are among His priests, and Samuel among those who call
upon His name.'' Now one must reflect deeply upon this blessed verse.
Today all trees are enraptured, rocks and stones tremble and move, and
the mountains pass along in the greatest joy, for the promised hour has
come and the hidden treasure has been revealed. If this sublime word
were removed from the books, no utterance would remain and no mention
would be found. But this statement is for those endowed with insight and
understanding, not for every base and ignorant person who leads the way
in vain imaginings and conjectures.

\subsection{42}\label{section-2017}

\begin{otherlanguage}{persian}

امروز آفتاب یقین اشراق نموده روز دیگر است و ایّام دیگر لا تقاس بغیرها
ولکن ابصار مریض و مرمود و آذان مسدود از حقّ میطلبم فرح عطا فرماید و خلق
را از این فرح اعظم محروم نسازد و همچنین میفرماید یأتی الیهنا و لا یصمت

\end{otherlanguage}

Imrūz āftāb-i yaqīn ishrāq namūdeh rūz-i dīgar ast va ayyām-i dīgar lā
tuqāsu bi-ghayrihā va-lākin abṣār marīḍ va marmūd va ādhān masdūd az
Ḥaqq mī-ṭalabam faraḥ ʿaṭā farmāyad va khalq rā az īn faraḥ-i aʿẓam
maḥrūm nasāzad va hamchunīn mī-farmāyad Yaʾtī {[}corrupted text{]} wa-lā
yaṣmut.

Today the sun of certainty has shone forth. It is another day, and these
are other days, incomparable with any others. Yet eyes are diseased and
inflamed, and ears are stopped. I ask the True One to bestow joy and not
to deprive His creatures of this greatest joy. He also says:
``{[}corrupted original text{]} comes and will not remain silent.''

\subsection{43}\label{section-2018}

\begin{otherlanguage}{persian}

اگر اینعبد بخواهد از کتب قبل و بیانات منزلهٴ در آن را ذکر نماید قارئین
شاید از طول کلام افسرده شوند لذا بمنتهای قوّه در اختصار عزم جزم شده و از
نفوس مطمئنّهٴ راضیهٴ مرضیّه طلب ذکر مینمایم شاید به دعای ایشان و ذکر ایشان
از عهدهٴ آنچه بآن مأمورم برآیم و از خدمت اولیای حقّ باز نمانم

\end{otherlanguage}

Agar īn ʿabd bikhvāhad az kutub-i qabl va bayānāt-i munazzaleh dar ān rā
dhikr namāyad qāriʾīn shāyad az ṭūl-i kalām afsordeh shavand lidhā
be-muntahā-yi quvveh dar ikhtiṣār ʿazm jazm shodeh va az nufūs-i
muṭmaʾinneh-yi rāḍiyeh-yi marḍiyyeh ṭalab-i dhikr mī-namāyam shāyad
be-duʿā-yi īshān va dhikr-i īshān az ʿohdeh-yi āncheh be-ān maʾmūram
bar-āyam va az khidmat-i awliyā-yi Ḥaqq bāz namānam.

If this servant wished to quote from the earlier books and the
utterances revealed in them, readers might grow weary of the length of
the discourse. I have therefore firmly resolved to be as brief as
possible. I ask those assured souls who are content and well-pleasing to
remember me, so that through their prayers and remembrance I may fulfill
what I have been charged to do and not fall short in serving the loved
ones of the True One.

\subsection{44}\label{section-2019}

\begin{otherlanguage}{persian}

فآه آه امر باین بزرگی و این قدرت ظاهره و قوّت باهره که عالم را احاطه
نموده خلق از مشاهده‌اش محجوب و از معرفتش غافل اوست مقتدریکه حضرت عاموس
نبی در باره‌اش میفرماید فاستعد للقآء الهک یا اسرائیل ١٣ فانّه هو الّذی صنع
الجبال و خلق الرّیح و اخبر الانسان ما هو فکره الّذی یجعل الفجر ظلاماً و
یمشی علی مشارف الأرض یَهْوَه اله الجنود اسمه لعمر محبوبنا و محبوبکم و
مقصودنا و مقصودکم در این آیهٴ عظیمه اسرار عظیمه مکنون چه که اخبار
مینماید از ظهور ثمانین و همچنین صبح کاذب را که میفرماید او را اله الجنود
تاریک مینماید و بر بلندیهای ارض میخرامد که اشاره بسفر مدینه کبیره و
همچنین این اطراف است

\end{otherlanguage}

Fa-āh āh amr be-īn buzurgī va īn qudrat-i ẓāhireh va quvvat-i bāhireh
keh ʿālam rā iḥāṭeh namūdeh khalq az mushāhadeh-ash maḥjūb va az
maʿrifatash ghāfil ūst muqtadirī keh Ḥaḍrat-i ʿĀmūs-i nabī dar bāreh-ash
mī-farmāyad Fa-istaʿidd li-liqāʾ ilāhika yā Isrāʾīl 13 fa-innahu huwa
alladhī ṣanaʿa al-jibāl wa-khalaqa al-rīḥ wa-akhbara al-insāna mā huwa
fikruhu alladhī yajʿalu al-fajra ẓalāman wa-yamshī ʿalā mashārif al-arḍ
Yahwah ilāh al-junūd ismuhu la-ʿamri maḥbūbinā wa-maḥbūbikum
wa-maqṣūdinā wa-maqṣūdikum dar īn āyeh-yi ʿaẓīmeh asrār-i ʿaẓīmeh maknūn
cheh keh ikhbār mī-namāyad az ẓuhūr-i thamānīn va hamchunīn ṣubḥ-i
kādhib rā keh mī-farmāyad ū rā ilāh al-junūd tārīk mī-namāyad va bar
bulandīhā-yi arḍ mī-khirāmad keh ishāreh be-safar-i Madīneh-yi Kabīreh
va hamchunīn īn aṭrāf ast.

Alas, alas! A Cause of such greatness, with this manifest power and
dazzling might encompassing the world, yet people are veiled from
beholding it and heedless of recognizing it. He is the Almighty of whom
the Prophet Amos says: ``Prepare to meet your God, O Israel. 13 For He
is the One who formed the mountains, created the wind, and told humanity
what His thought is; who makes the dawn darkness and walks upon the
heights of the earth. Yahweh, the God of Hosts, is His name.'' By the
life of our Beloved and your Beloved, our Desire and your Desire! Great
mysteries lie concealed in this mighty verse, for it foretells the
appearance of eighty, and likewise the false dawn, which, He says, the
God of Hosts darkens. He walks upon the heights of the earth, which
alludes to the journey to the Great City and also to these regions.

\subsection{45}\label{section-2020}

\begin{otherlanguage}{persian}

اینعبد متحیّر که ناس از قبل بچه قانع شده‌اند و حال بچه قانع میشوند اگر
آیات میطلبند انّها ملئت الآفاق و اگر بیّنات میخواهند علی العبد اثباتها و
اظهارها قسم بنفس حقّ روح من فی ملکوت الأمر و الخلق فداه که امورات عظیمهٴ
عجیبه در اینظهور ظاهر و آنچه در سنین متوالیه از بعد واقع از قبل از قلم
اعلی جاری و نازل من‌غیر رمز و اشاره بل بکمال تصریح هر منصفی عرف راستی از
این عرایض استشمام مینماید چه که آنچه عرض شده و میشود با بیّنه و برهانست

\end{otherlanguage}

Īn ʿabd mutaḥayyir keh nās az qabl be-cheh qāniʿ shodeh-and va ḥāl
be-cheh qāniʿ mī-shavand agar āyāt mī-ṭalaband innahā malaʾat al-āfāq va
agar bayyināt mī-khvāhand ʿalā al-ʿabd ithbātuhā wa-iẓhāruhā qasam
be-nafs-i Ḥaqq rūḥ-i man fī malakūt al-amr wa-al-khalq fidāh keh
umūrāt-i ʿaẓīmeh-yi ʿajībeh dar īn ẓuhūr ẓāhir va āncheh dar sinīn-i
mutavāliyeh az baʿd vāqiʿ az qabl az Qalam-i Aʿlā jārī va nāzil min
ghayri ramz wa-ishārah bal be-kamāl-i taṣrīḥ har munṣifī ʿarf-i rāstī az
īn ʿarāyiḍ istishmām mī-namāyad cheh keh āncheh ʿarḍ shodeh va mī-shavad
bā bayyineh va burhānast.

This servant is bewildered as to what satisfied people in the past and
what satisfies them now. If they seek verses, these have filled the
horizons; if they want clear proofs, it is for this servant to establish
and display them. I swear by the True One Himself---may the spirits of
all in the kingdoms of command and creation be sacrificed for Him---that
great and wondrous things have appeared in this Revelation. Whatever
occurred in successive years afterward had already flowed down and been
revealed from the Most Exalted Pen, without cryptic language or
allusion, but with complete explicitness. Every fair-minded person
inhales the fragrance of truthfulness from these submissions, for
whatever has been or is being set forth is supported by clear evidence
and proof.

\subsection{46}\label{section-2021}

\begin{otherlanguage}{persian}

و اینعبد اگر بخواهد جمیع را ذکر نماید البتّه بیان بطول انجامد ولکن بعضی
از آنرا مختصر ذکر مینمایم لوجه اللّه اگرچه در رسائل متعدّده ذکر شده و چون
بعضی از آن از قرار مسموع دست اعدا افتاده و بعضی از آن مستور مانده لذا
این خادم مکرّر ذکر مینماید و اجره علی الّذی خلقه و فطره و رزقه و کبّره و
علّمه و ارشده و هداه الی صراطه المستقیم

\end{otherlanguage}

Va īn ʿabd agar bikhvāhad jamīʿ rā dhikr namāyad albatteh bayān be-ṭūl
anjāmad va-lākin baʿḍī az ān rā mukhtaṣar dhikr mī-namāyam li-wajhi
Allāh agarch eh dar rasāʾil-i mutaʿaddideh dhikr shodeh va chūn baʿḍī az
ān az qarār-i masmūʿ dast-i aʿdā uftādeh va baʿḍī az ān mastūr māndeh
lidhā īn khādim mukarrar dhikr mī-namāyad wa-ajruhu ʿalā alladhī
khalaqahu wa-faṭarahu wa-razaqahu wa-kabbarahu wa-ʿallamahu wa-arshadahu
wa-hadāhu ilā ṣirāṭihi al-mustaqīm.

If this servant wished to recount everything, the account would
certainly grow lengthy. I will, however, briefly mention some of it for
the sake of God, although it has been mentioned in several epistles.
Since some of these, according to what has been heard, have fallen into
the hands of enemies, while others have remained concealed, this servant
recounts it again. His reward rests with the One who created him,
brought him into being, sustained him, raised him up, taught him,
directed him, and guided him to His straight path.

\subsection{47}\label{section-2022}

\begin{otherlanguage}{arabic}

هنگامیکه آفتاب حقیقت از افق ارض سرّ مشرق یومی از ایّام بغتةً بیت را احاطه
نمودند و ذکر تفصیل آن حال لازم نه مختصر عرض میشود مدّتی ناس را از خروج و
دخول منع نمودند ولکن بعد احزاب مختلفه که بدول اخری منسوب بودند وارد و
تلقاء وجه حاضر و فیکلّ حین نطق لسان العظمة بما عجز عن ذکره السن العالم
قناسل دول بر حسب ظاهر باعانت قیام نمودند و بین یدی حاضر ولکن ابی اللّه عن
اظهار ما ارادوا

\end{otherlanguage}

Hingāmī keh Āftāb-i Ḥaqīqat az ufuq-i Arḍ-i Sirr mushriq yawmī az ayyām
baghtatan bayt rā iḥāṭeh namūdand va dhikr-i tafṣīl-i ān ḥāl lāzim nah
mukhtaṣar ʿarḍ mī-shavad muddatī nās rā az khurūj va dukhūl manʿ
namūdand va-lākin baʿd aḥzāb-i mukhtalifeh keh be-duval-i ukhrā mansūb
būdand vā rid va tilqā-yi vajh ḥāḍir va-fī kulli ḥīn naṭaqa Lisān
al-ʿAẓamah bi-mā ʿajaza ʿan dhikrihi alsun al-ʿālam qanāṣil-i duval bar
ḥasab-i ẓāhir be-iʿānat qiyām namūdand va bayn-i yaday ḥāḍir va-lākin
abā Allāh ʿan iẓhār mā arādū.

While the Sun of Truth was shining from the horizon of the Land of
Mystery, one day they suddenly surrounded the house. It is unnecessary
to recount the details now; a brief account is given. For a time they
prevented people from leaving or entering. Afterward, however, various
groups affiliated with other states came and entered His presence. At
every moment the Tongue of Grandeur uttered what the tongues of the
world are powerless to recount. The consuls of the states outwardly
arose to offer assistance and came into His presence, but God refused to
bring about what they intended.

\subsection{48}\label{section-2023}

\begin{otherlanguage}{persian}

و بعد حضرت ذبیح علیه بهآء اللّه الأبهی وارد و قلم اعلی بسورهٴ مبارکه رئیس
ابتدا فرموده و در عرض راه تمام شد و در اسکلهٴ بحر بر شخص مأمور قرائت
فرمود اینعبد چند آیه آن را ذکر مینماید که شاید غافلان آگاه شوند و نائمین
بیدار گردند و مردگان بحیات جدید فائز شوند الأمر بید محبوبنا و محبوبکم و
مقصودنا و مقصود من فی السّموات و الأرضین قوله جلّ مقامه و عزّ بیانه و کبر
شأنه

\end{otherlanguage}

Va baʿd Ḥaḍrat-i Dhabīḥ ʿalayhi Bahāʾ Allāh al-Abhā vā rid va Qalam-i
Aʿlā be-Sūreh-yi mubārakeh-yi Raʾīs ibtidā farmūdeh va dar ʿarḍ-i rāh
tamām shod va dar iskeleh-yi baḥr bar shak hṣ-i maʾmūr qirāʾat farmūd īn
ʿabd chand āyeh-yi ān rā dhikr mī-namāyad keh shāyad ghāfilān āgāh
shavand va nāʾimīn bīdār gardand va murdegān be-ḥayāt-i jadīd fāʾiz
shavand al-amru bi-yadi maḥbūbinā wa-maḥbūbikum wa-maqṣūdinā wa-maqṣūd
man fī al-samāwāt wa-al-araḍīn qawluhu jalla maqāmuhu wa-ʿazza bayānuhu
wa-kabura shaʾnuhu.

Afterward, Dhabih---upon him be the glory of God, the
All-Glorious---arrived, and the Most Exalted Pen began the blessed Surih
of the Chief. It was completed during the journey, and He recited it to
the official at the seaport. This servant quotes a few of its verses so
that perhaps the heedless may become aware, the sleeping awaken, and the
dead attain a new life. The matter rests in the hand of our Beloved and
your Beloved, our Desire and the Desire of all in the heavens and on
earth. His words---exalted be His station, mighty His utterance, and
great His glory:

\subsection{49}\label{section-2024}

\begin{otherlanguage}{arabic}

بسمه الأبهی

\end{otherlanguage}

Bi-ismihi al-Abhā.

In His name, the All-Glorious.

\subsection{50}\label{section-2025}

\begin{otherlanguage}{arabic}

ان یا رئیس اسمع ندآء اللّه الملک المهیمن القیّوم انّه ینادی بین الأرض و
السّمآء و یدع الکلّ الی المنظر الأبهی و لا یمنعه قباعک و لا نباح من فی
حولک و لا جنود العالمین قد اشتعل العالم من کلمة ربّک الأبهی و انّها ارقّ من
نسیم الصّبا قد ظهرت علی هیئة الانسان و بها احی اللّه عباده المقبلین و فی
باطنها مآء قد طهّر اللّه به افئدة الّذین اقبلوا الیه و غفلوا عن ذکر ما سواه
و قرّبهم الی منظر اسمه العظیم الی ان قال عزّ شأنه

\end{otherlanguage}

In yā raʾīs ismaʿ nidāʾ Allāh al-Malik al-Muhaymin al-Qayyūm innahu
yunādī bayna al-arḍ wa-al-samāʾ wa-yadʿu al-kulla ilā al-manẓar al-Abhā
wa-lā yamnaʿuhu qubāʿuka wa-lā nubāḥu man fī ḥawlika wa-lā junūd
al-ʿālamīn qad ishtaʿala al-ʿālam min kalimati rabbika al-Abhā wa-innahā
araqqu min nasīm al-ṣabā qad ẓaharat ʿalā hayʾati al-insān wa-bihā aḥyā
Allāh ʿibādahu al-muqbilīn wa-fī bāṭinihā māʾun qad ṭahhara Allāh bihi
afʾidata alladhīna aqbalū ilayhi wa-ghafalū ʿan dhikri mā siwāhu
wa-qarrabahum ilā manẓari ismihi al-ʿAẓīm ilā an qāla ʿazza shaʾnuhu.

``O Chief, hear the call of God, the King, the All-Protecting, the
Self-Subsisting. He calls aloud between earth and heaven and summons all
to the All-Glorious Scene. Neither your grunting, nor the barking of
those around you, nor the armies of the worlds can hinder Him. The world
has been set ablaze by the Word of your Lord, the All-Glorious. More
delicate than the morning breeze, it has appeared in human form. Through
it God has brought to life those of His servants who have turned toward
Him. Within it is water with which God has purified the hearts of those
who have turned toward Him and become oblivious to the mention of
anything besides Him, and He has drawn them near to the scene of His
Great Name.'' Further on He says---exalted be His glory:

\subsection{51}\label{section-2026}

\begin{otherlanguage}{persian}

ان یا رئیس قد ارتکبت ما ینوح به محمّد رسول اللّه فی الجنّة العلیا و غرّتک
الدّنیا علی شأن اعرضت عن الوجه الّذی بنوره استضآء الملأ الأعلی فسوف تجد
نفسک فی خسران مبین و اتّحدت مع رئیس العجم فی ضرّی بعد اذ جئتکم من مطلع
العظمة و الکبریآء بامر به قرّت عیون المقرّبین الی ان قال جلّ ثنائه هل ظننت
انّک تقدر ان تطفیء النّار الّتی اوقدها اللّه فی الآفاق لا ونفسه الحقّ لو انت
من العارفین بل بما فعلت زاد لهیبها و اشتعالها فسوف تحیط الأرض و من علیها
کذلک قضی الأمر و لا یقوم معه حکم من فی السّموات و الأرضین فسوف تبدّل ارض
السّرّ و ما دونها و تخرج من ید الملک و یظهر الزّلازل و یرتفع العویل و یظهر
الفساد فی الأقطار و تختلف الأمور بما ورد علی هؤلآء الأسرآء من جنود
الظّالمین و یتغیّر الحکم و یشتدّ الأمر علی شأن ینوح الکتیب فی الهضاب و تبکی
الأشجار فی الجبال و تجری الدّمّ من کلّ الأشیآء و تری النّاس فی اضطراب عظیم
انتهی

\end{otherlanguage}

In yā raʾīs qad irtakabta mā yanūḥu bihi Muḥammad Rasūl Allāh fī
al-jannati al-ʿulyā wa-gharratka al-dunyā ʿalā shaʾnin aʿraḍta ʿan
al-wajh alladhī bi-nūrihi istaḍāʾa al-malaʾ al-aʿlā fa-sawfa tajidu
nafsaka fī khusrānin mubīn wa-ittaḥadta maʿa raʾīs al-ʿAjam fī ḍurrī
baʿda idh jiʾtukum min maṭlaʿ al-ʿaẓamah wa-al-kibriyāʾ bi-amrin bihi
qarrat ʿuyūn al-muqarrabīn ilā an qāla jalla thanāʾuhu hal ẓananta
annaka taqdiru an tuṭfiʾa al-nāra allatī awqadahā Allāh fī al-āfāq lā
wa-nafsihi al-Ḥaqq law anta min al-ʿārifīn bal bi-mā faʿalta zāda
lahībuhā wa-ishtiʿāluhā fa-sawfa tuḥīṭu al-arḍa wa-man ʿalayhā kadhālika
quḍiya al-amru wa-lā yaqūmu maʿahu ḥukmu man fī al-samāwāt wa-al-araḍīn
fa-sawfa tubaddalu Arḍ al-Sirr wa-mā dūnahā wa-takhruju min yad al-malik
wa-yaẓharu al-zalāzil wa-yartafiʿu al-ʿawīl wa-yaẓharu al-fasādu fī
al-aqṭār wa-takhtalifu al-umūr bi-mā warada ʿalā hāʾulāʾ al-asrāʾ min
junūd al-ẓālimīn wa-yataghayyaru al-ḥukm wa-yashtaddu al-amr ʿalā
shaʾnin yanūḥu {[}corrupted text{]} fī al-hiḍāb wa-tabkī al-ashjār fī
al-jibāl wa-tajrī al-damm min kull al-ashyāʾ wa-tarā al-nāsa fī
iḍṭirābin ʿaẓīm intahā.

``O Chief, you have committed what causes Muhammad, the Messenger of
God, to lament in the highest Paradise. The world has so deceived you
that you have turned away from the Face whose light has illumined the
Concourse on high. You will soon find yourself in manifest loss. You
have joined with the chief of Persia to harm Me, after I came to you
from the dawning-place of grandeur and majesty with a Cause that brought
delight to the eyes of those brought near to God.'' Further on He
says---exalted be His praise: ``Did you imagine that you could
extinguish the fire God has kindled throughout the horizons? No, by His
own true Self, if only you understood! Rather, through what you have
done, its flame and its blaze have increased. It will soon encompass the
earth and all who dwell upon it. Thus has the matter been decreed, and
the judgment of those in the heavens and on earth cannot withstand it.
The Land of Mystery and other lands will soon be changed and pass out of
the king's hands. Earthquakes will occur, wailing will rise, corruption
will appear throughout the regions, and affairs will fall into disorder
because of what has befallen these captives at the hands of the armies
of the oppressors. The government will change, and conditions will grow
so severe that {[}corrupted original text{]} will lament in the hills,
the trees will weep in the mountains, blood will flow from all things,
and you will see the people in great turmoil.'' End of quotation.

\subsection{52}\label{section-2027}

\begin{otherlanguage}{persian}

ای عالم آیا سمع در تو یافت میشود که بشنود آیا بصر یافت میشود که ببیند به
ابدع و اصرح بیان محبوب امکان نطق فرموده و مقصود از ارض سرّ ادرنه است که
قبل از سجن اعظم مقرّ عرش آن ارض بود و عدد آن با عدد سِرْ مطابق است میفرماید
ارض سِرْ تبدیل و همچنین ما حول آن و از ید ملک بیرون میرود و ظاهر میشود
زلازل و مرتفع شود فریاد و ناله و ظاهر میگردد فساد در اقطار و مختلف میشود
امور جمیع این بلایا بسبب این مصائبی است که بر این اسرا وارد شده از جنود
ظالمین و جمیع این امور در ظاهر ظاهر دیده شد و هیچ منصف و صادق و عادلی
انکار ننموده و نمینماید و همچنین میفرماید و یتغیّر الحکم یعنی رئیس اوّل
میرود و رئیس دیگر بر مقامش جالس چنانچه مشاهده شد و همچنین میفرماید یشتدّ
الأمر شدید میشود امر بشأنیکه رمل در تپها و کوهها و اشجار در جبال میگریند
و از جمیع اشیا دم جاری و ساری چه که بر نفس حقّ وارد شد آنچه که بر احدی
وارد نه و بعد میفرماید تری النّاس فی اضطراب عظیم

\end{otherlanguage}

Ay ʿālam āyā samʿ dar tu yāft mī-shavad keh bishnavad āyā baṣar yāft
mī-shavad keh bibīnad be-abdaʿ va aṣraḥ bayān Maḥbūb-i imkān nuṭq
farmūdeh va maqṣūd az Arḍ-i Sirr Adirneh ast keh qabl az Sijn-i Aʿẓam
maqarr-i ʿarsh ān arḍ būd va ʿadad-i ān bā ʿadad-i Sirr muṭābiq ast
mī-farmāyad Arḍ-i Sirr tabdīl va hamchunīn mā ḥawl-i ān va az yad-i
malik bīrūn mī-ravad va ẓāhir mī-shavad zalāzil va murtafiʿ shavad
faryād va nāleh va ẓāhir mī-gardad fasād dar aqṭār va mukhtalif
mī-shavad umūr jamīʿ-i īn balāyā be-sabab-i īn maṣāʾibī ast keh bar īn
usarā vā rid shodeh az junūd-i ẓālimīn va jamīʿ-i īn umūr dar ẓāhir
ẓāhir dīdeh shod va hīch munṣif va ṣādiq va ʿādilī inkār nanamūdeh va
namī-namāyad va hamchunīn mī-farmāyad wa-yataghayyaru al-ḥukm yaʿnī
raʾīs-i avval mī-ravad va raʾīs-i dīgar bar maqāmash jālis chunāncheh
mushāhedeh shod va hamchunīn mī-farmāyad yashtaddu al-amr shadīd
mī-shavad amr be-shaʾnī keh raml dar tappeh-hā va kūh-hā va ashjār dar
jibāl mī-giryand va az jamīʿ-i ashyā dam jārī va sārī cheh keh bar
nafs-i Ḥaqq vā rid shod āncheh keh bar aḥadī vā rid nah va baʿd
mī-farmāyad tarā al-nāsa fī iḍṭirābin ʿaẓīm.

O world, is there any hearing within you that will listen? Is there any
sight that will see? The Beloved of all existence has spoken in the most
wondrous and explicit language. By the Land of Mystery is meant
Adrianople, for before the Most Great Prison that land was the seat of
the throne, and its numerical value corresponds to that of ``Sirr.'' He
says that the Land of Mystery will be changed, as will its surroundings,
and will pass out of the king's hands; earthquakes will occur, cries and
lamentations will rise, corruption will appear throughout the regions,
and affairs will fall into disorder. All these calamities are because of
the afflictions brought upon these captives by the armies of the
oppressors. All these things were seen to occur outwardly, and no
fair-minded, truthful, and just person has denied or denies them. He
also says, ``The government will change,'' meaning that the first chief
departs and another chief takes his place, as was witnessed. He also
says, ``Conditions will grow severe'': matters will become so grave that
the sand in the hills and mountains, and the trees upon the mountains,
will weep, and blood will flow and stream from all things, for there
befell the True One Himself what had befallen no one. Then He says,
``You will see the people in great turmoil.''

\subsection{53}\label{section-2028}

\begin{otherlanguage}{persian}

لعمراللّه اضطراب کلّ را احاطه نموده بود جنود اسلام بر قراء نصاری وارد
میشدند زن و مرد صغیر و کبیر جمیع را از حیات محروم مینمودند و بعد خانه‌ها
و اموال را آتش میزدند و همین معامله را عیناً اسلام با نصاری میکردند وحشت
کلّ را اخذ کرده بود و اضطراب جمیع را احاطه یکی از سیّاحان نوشته بود از
محلّی عبور نمودم مشاهده شد در بیست قدم او سی سیصد نعش بر هم ریخته بود
هرگز رعب آنمنظر از نظر نمیرود باری آنفتنه را کلّ شنیده و میدانید و در هیچ
قرنی مثل آن دیده نشده ارباب جرائد نوشته‌اند آنچه را که عقل از قبول آن ابا
مینماید حال باید اهل بیان تفکّر نمایند و بیقین مبین بدانند که آنچه از بعد
وارد شده حرف بحرف از قبل از قلم اعلی نازل سبحان من تحیی العظام و هی رمیم
و سبحان من اخبر الکلّ من قبل و هو علیم

\end{otherlanguage}

La-ʿamru'llāh eżṭerāb koll rā eḥāṭe nemūde būd. Jonūd-e Eslām bar
qorā-ye Naṣārā vāred mī-shodand, zan o mard, ṣaghīr o kabīr, jamīʿ rā az
ḥayāt maḥrūm mī-nemūdand va baʿd khāne-hā va amvāl rā ātash mī-zadand va
hamīn moʿāmele rā ʿaynan Eslām bā Naṣārā mī-kardand. Vaḥshat koll rā
akhdh karde būd va eżṭerāb jamīʿ rā eḥāṭe. Yekī az sayyāḥān neveshte
būd: Az maḥallī ʿobūr nemūdam, moshāhede shod dar bīst qadam
{[}corrupted text{]} sīṣad naʿsh bar ham rīkhte būd. Hargez roʿb-e ān
manẓar az naẓar nemī-ravad. Bārī ān fetne rā koll shenīde va mī-dānīd va
dar hīch qarnī methl-e ān dīde nashode. Arbāb-e jarāʾed neveshte-and
ānche rā ke ʿaql az qabūl-e ān ebā mī-nemāyad. Ḥāl bāyad ahl-e Bayān
tafakkor nemāyand va be-yaqīn-e mobīn be-dānand ke ānche az baʿd vāred
shode ḥarf be-ḥarf az qabl az Qalam-e Aʿlā nāzel. Subḥāna man tuḥyī
al-ʿiẓāma wa-hiya ramīm, wa-subḥāna man akhbara al-kulla min qablu
wa-huwa ʿalīm.

By the life of God, turmoil had engulfed everyone. Muslim troops entered
Christian villages, depriving everyone of life---women and men, young
and old---and then setting fire to their houses and possessions. This
was precisely how the Muslims treated the Christians. Terror had seized
everyone, and turmoil had engulfed them all. One traveler wrote, ``I
passed through a place where, within twenty paces, {[}corrupted original
text{]} three hundred corpses lay piled upon one another. The horror of
that sight never leaves my eyes.'' In short, you have all heard of that
upheaval and know of it, and nothing like it has been seen in any
century. Journalists have written things that the mind refuses to
accept. Now the people of the Bayán must reflect and know with
unmistakable certainty that everything that subsequently occurred had
previously been revealed, word for word, by the Most Exalted Pen. Glory
to Him who revives bones when they have crumbled to dust, and glory to
Him who foretold everything and is all-knowing.

\subsection{54}\label{section-2029}

\begin{otherlanguage}{persian}

ولکن یا محبوب فؤادی و المذکور فی قلبی آنچه ذکر شد از برای اولیا و اصفیا
و راسخین و منصفین بشارتیست بزرگ ولکن لا یزید الظّالمین الّا خسارا چه که کَر
از هدیر حمامهٴ امر محروم و کور از انوار افق اعلی ممنوع

\end{otherlanguage}

Va-lākin yā maḥbūba fuʾādī wa-l-madhkūra fī qalbī, ānche dhekr shod az
barā-ye awliyā va aṣfiyā va rāsekhīn va monṣefīn beshāratīst bozorg,
va-lākin lā yazīdu al-ẓālimīna illā khasāran, che ke kar az hadīr-e
Ḥamāme-ye Amr maḥrūm va kūr az anvār-e Ofoq-e Aʿlā mamnūʿ.

Yet, O beloved of my heart, whose remembrance is in my heart, what has
been mentioned is a great glad tiding for the faithful, the chosen, the
steadfast, and the fair-minded, but ``it increases the wrongdoers in
nothing but loss,'' for the deaf are deprived of the cooing of the Dove
of the Cause, and the blind are barred from the lights of the Most
Exalted Horizon.

\subsection{55}\label{section-2030}

\begin{otherlanguage}{arabic}

و همچنین عزّ ذکره و عمّ فضله و عظم احسانه در سورهٴ هیکل میفرماید مخاطباً
لملک باریس قوله جلّ جلاله

\end{otherlanguage}

Va hamchenīn, ʿazza dhikruhu wa-ʿamma faḍluhu wa-ʿaẓuma iḥsānuhu, dar
Sūre-ye Haykal mī-farmāyad, mukhāṭiban li-maliki Bārīs, qawluhu jalla
jalāluhu:

He also says---exalted be His remembrance, all-embracing His grace, and
great His bounty---in the Súrih of the Temple, addressing the King of
Paris, these words---glorified be His glory:

\subsection{56}\label{section-2031}

\begin{otherlanguage}{arabic}

قل یا ملک باریس نبّیء القسّیس ان لا یدقّ النّواقیس تاللّه الحقّ قد ظهر النّاقوس
الأفخم علی هیکل الاسم الأعظم و تدقّه اصابع مشیّة ربّک العلیّ الأعلی فی جبروت
البقآء باسمه الأبهی کذلک نزّلت آیات ربّک الکبری بارة اخری لتقوم علی ذکر
اللّه فاطر الأرض و السّمآء فی هذه الأیّام الّتی فیها ناحت قبائل الأرض کلّها و
تزلزلت ارکان البلاد و غشّت العباد غبرة الالحاد الّا من شآء ربّک العزیز
الحکیم

\end{otherlanguage}

Qul yā malika Bārīs, nabbiʾi al-qissīsa an lā yaduqqa al-nawāqīsa.
Ta'llāhi al-Ḥaqqi qad ẓahara al-nāqūsu al-afkhamu ʿalā haykali al-Ismi
al-Aʿẓami wa-taduqquhu aṣābiʿu mashīʾati rabbika al-ʿAliyyi al-Aʿlā fī
jabarūti al-baqāʾi bi-smihi al-Abhā. Kadhālika nuzzilat āyātu rabbika
al-kubrā {[}corrupted text{]} ukhrā li-taqūma ʿalā dhikri Allāhi fāṭiri
al-arḍi wa-l-samāʾi fī hādhihi al-ayyāmi allatī fīhā nāḥat qabāʾilu
al-arḍi kulluhā wa-tazalzalat arkānu al-bilādi wa-ghashshat al-ʿibāda
ghubratu al-ilḥādi illā man shāʾa rabbuka al-ʿAzīzu al-Ḥakīm.

``Say: O King of Paris, tell the priest not to ring the bells. By God,
the True One, the Most Mighty Bell has appeared in the form of the Most
Great Name, and the fingers of the will of your Lord, the Most High, the
Most Exalted, ring it in the dominion of eternity in His name, the
All-Glorious. Thus have the most great verses of your Lord been sent
down {[}corrupted original text{]} another, so that you may arise to
remember God, the Creator of earth and heaven, in these days when all
the peoples of the earth have lamented, the foundations of the countries
have been shaken, and the dust of unbelief has enveloped all people
except those whom your Lord, the Almighty, the All-Wise, has willed to
spare.''

\subsection{57}\label{section-2032}

\begin{otherlanguage}{persian}

الی ان قال جلّت عظمته ان یا ملک ان استمع النّدآء من هذه النّار المشتعلة من
هذه الشّجرة الخضرآء فی هذا الطّور المرتفع علی البقعة المقدّسة البیضآء خلف
قلزم البقآء انّه لا اله الّا انا الغفور الرّحیم

\end{otherlanguage}

Ilā an qāla, jallat ʿaẓamatuhu: An yā malik, an istamiʿi al-nidāʾa min
hādhihi al-nāri al-mushtaʿilati min hādhihi al-shajarati al-khaḍrāʾi fī
hādhā al-Ṭūri al-murtafiʿi ʿalā al-buqʿati al-muqaddasati al-bayḍāʾi
khalfa qulzumi al-baqāʾi: Innahu lā ilāha illā anā al-Ghafūru al-Raḥīm.

Further on He says---exalted be His grandeur: ``O King, listen to the
call from this Fire blazing in this verdant Tree upon this lofty Mount,
in the holy, snow-white Spot beyond the ocean of eternity: `There is no
God but Me, the Ever-Forgiving, the Merciful.'\,''

\subsection{58}\label{section-2033}

\begin{otherlanguage}{arabic}

الی ان قال عزّ بیانه انّ شجرة الطّور تنطق فی صدر العالم و روح القدس ینادی
بین الأمم قد اتی المقصود بسلطان مبین الی

\end{otherlanguage}

Ilā an qāla, ʿazza bayānuhu: Inna shajarata al-Ṭūri tanṭiqu fī ṣadri
al-ʿālami wa-Rūḥu al-Qudusi yunādī bayna al-umami: Qad atā al-Maqṣūdu
bi-sulṭānin mubīn. Ilā

Further on He says---exalted be His utterance: ``The Tree of Sinai
speaks in the heart of the world, and the Holy Spirit calls out among
the nations: `The Desired One has come with manifest sovereignty.'\,''
Until

\subsection{59}\label{section-2034}

\begin{otherlanguage}{arabic}

ان قال عزّ شأنه ان یا ملک انّا سمعنا منک کلمة تکلّمت بها اذ سئلک ملک الرّوس
عمّا قضی من حکم الغزا انّ ربّک لهو العلیم الخبیر قلت کنت راقداً فی المهاد
ایقظنی ندآء العباد الّذین ظلموا الی ان غرقوا فی البحر الأسود کذلک سمعنا و
ربّک علی ما اقول شهید یشهد بانّک ما ایقظک النّدآء بل الهوی لانّا تلوناک
وجدناک فی معزل ان اعرف لحن القول و کن من المتفرّسین انّا لا نحبّ ان یرجع
الیک کلمة سوء حفظاً للمقام الّذی اعطیناک فی الحیوة الظّاهره انّا اخترنا
الأدب و جعلناه سجیّة المقرّبین انّه ثوب توافق النّفوس من کلّ صغیر و کبیر طوبی
لمن جعله طراز هیکله ویل لمن جعل محروماً من هذا الفضل العظیم

\end{otherlanguage}

An qāla, ʿazza shaʾnuhu: An yā malik, innā samiʿnā minka kalimatan
takallamta bihā idh saʾalaka maliku al-Rūsi ʿammā quḍiya min ḥukmi
al-ghazā. Inna rabbaka la-huwa al-ʿAlīmu al-Khabīr. Qulta: Kuntu rāqidan
fī al-mihādi ayqaẓanī nidāʾu al-ʿibādi alladhīna ẓulimū ilā an ghariqū
fī al-Baḥri al-Aswad. Kadhālika samiʿnā wa-rabbuka ʿalā mā aqūlu shahīd.
Yashhadu bi-annaka mā ayqaẓaka al-nidāʾu bal al-hawā li-annā
{[}corrupted text{]} wajadnāka fī maʿzil. An iʿrif laḥna al-qawli wa-kun
mina al-mutafarrisīn. Innā lā nuḥibbu an yarjiʿa ilayka kalimatu sūʾin
ḥifẓan li-l-maqāmi alladhī aʿṭaynāka fī al-ḥayāti al-ẓāhirah. Innā
ikhtarnā al-adaba wa-jaʿalnāhu sajiyyata al-muqarrabīn. Innahu thawbun
tuwāfiqu al-nufūsa min kulli ṣaghīrin wa-kabīr. Ṭūbā li-man jaʿalahu
ṭirāza haykalihi, waylun li-man juʿila maḥrūman min hādhā al-faḍli
al-ʿaẓīm.

He says---exalted be His station: ``O King, We heard a statement you
made when the King of Russia questioned you about the decision to wage
war. Your Lord is indeed the All-Knowing, the All-Informed. You said, `I
was asleep in my bed when I was awakened by the cry of the servants who
were wronged until they drowned in the Black Sea.' Thus We heard, and
your Lord bears witness to what I say. He testifies that it was not the
cry that awakened you, but desire, for We {[}corrupted original text{]}
found you aloof. Recognize the meaning of these words and be among those
who discern. We do not wish any word of reproach to be directed at you,
out of regard for the station We have bestowed on you in this outward
life. We have chosen courtesy and made it the characteristic of those
brought near. It is a garment that suits every soul, young and old.
Blessed is the one who makes it the adornment of his person; woe to the
one who is deprived of this great bounty.''

\subsection{60}\label{section-2035}

\begin{otherlanguage}{arabic}

لو کنت صاحب الکلمة ما نبذت کتاب اللّه ورآء ظهرک اذ ارسل الیک من لدن عزیز
حکیم انّا یلوناک به ما وجدناک علی ماادّعیت قم و تدارک ما فات عنک سوف تفنی
الدّنیا و ما عندک و یبقی الملک للّه ربّک و ربّ آبائک الأوّلین لا ینبغی لک ان
تقتصر الأمور علی ما تهوی به هویک اتّق زفرات المظلوم ان احفظه من سهام
الظّالمین بما فعلن تختلف الأمور فی مملکتک و تخرج الملک من کفّک جزآء عملک
اذاً تجد نفسک فی خسران مبین و تأخذ الزّلازل کلّ القبائل هناک الّا بأن تقوم
علی نصرة هذا الامر و تتّبع الرّوح فی هذا السّبیل المستقیم أ عزّک غرّک لعمری
انّه لا یدوم سوف یزول الّا بأن تتمسّک بهذا الحبل المتین قد نری الذّلّة تسعی
عن ورائک و انت من الرّاقدین ینبغی لک اذا سمعت النّدآء من شطر الکبریآء تدع
ما عندک و تقول لبّیک یا اله من فی السّموات و الأرضین

\end{otherlanguage}

Law kunta ṣāḥiba al-kalimati mā nabadhta kitāba Allāhi warāʾa ẓahrika
idh ursila ilayka min ladun ʿAzīzin Ḥakīm. Innā {[}corrupted text{]}
bihi mā wajadnāka ʿalā mā iddaʿayta. Qum wa-tadārak mā fāta ʿanka. Sawfa
tafnā al-dunyā wa-mā ʿindaka wa-yabqā al-mulku li'llāhi rabbika wa-rabbi
ābāʾika al-awwalīn. Lā yanbaghī laka an taqtaṣira al-umūra ʿalā mā tahwā
bihi hawāka. Ittaqi zafarāti al-maẓlūmi an iḥfaẓhu min sihāmi
al-ẓālimīn. Bi-mā {[}corrupted text{]} takhtalifu al-umūru fī
mamlakatika wa-takhruju al-mulku min kaffika jazāʾa ʿamalika, idhan
tajidu nafsaka fī khusrānin mubīn. Wa-taʾkhudhu al-zalāzilu kulla
al-qabāʾili hunāka illā bi-an taqūma ʿalā nuṣrati hādhā al-amri
wa-tattabiʿa al-Rūḥa fī hādhā al-sabīli al-mustaqīm. A-ʿizzuka gharraka?
La-ʿamrī innahu lā yadūmu, sawfa yazūlu illā bi-an tatamassaka bi-hādhā
al-ḥabli al-matīn. Qad narā al-dhillata tasʿā ʿan warāʾika wa-anta mina
al-rāqidīn. Yanbaghī laka idhā samiʿta al-nidāʾa min shaṭri al-kibriyāʾi
tadaʿa mā ʿindaka wa-taqūla: Labbayka yā ilāha man fī al-samāwāti
wa-l-araḍīn.

``Had you been true to your word, you would not have cast the Book of
God behind your back when it was sent to you by One who is Almighty,
All-Wise. We {[}corrupted original text{]} through it and did not find
you as you claimed to be. Arise and make amends for what you have
missed. The world and all you possess will pass away, and sovereignty
will remain with God, your Lord and the Lord of your forefathers. You
should not confine affairs to whatever your desire inclines toward. Fear
the sighs of the Wronged One; protect Him from the arrows of the
oppressors. Because of what {[}corrupted original text{]}, affairs in
your kingdom will fall into disorder, and sovereignty will leave your
grasp as recompense for your actions. Then you will find yourself in
manifest loss. Earthquakes will seize all the peoples there unless you
arise to aid this Cause and follow the Spirit on this straight path. Has
your glory deceived you? By My life, it will not endure; it will pass
away unless you hold fast to this firm cord. We see humiliation
hastening after you while you remain asleep. When you hear the call from
the realm of grandeur, you should cast aside what you possess and say,
`Here I am, O God of all who are in the heavens and on earth!'\,''

\subsection{61}\label{section-2036}

\begin{otherlanguage}{arabic}

ان یا ملک قد کنّا بأمّ العراق الی ان حمّ الفراق توجّهنا الی ملک الاسلام
بأمره فلمّا اتیناه ورد علینا من اولی النّفاق ما لا یتمّ بالأوراق بذلک ناح
سکّان الفردوس و اهل حظائر القدس ولکنّ القوم فی حجاب غلیظ قل أ تعترضون علی
الّذی جائکم ببیّنات اللّه و برهانه ثمّ حجّته و آیاته ان هی من تلقآء نفسه بل
من لدن من بعثه و ارسله بالحقّ و جعله سراجاً للعالمین قد اشتدّ علینا الأمر
فی کلّ یوم بل فی کلّ ساعة الی ان اخرجونا من السّجن و ادخلونا فی السّجن
الأعظم بظلم مبین اذا قیل بأیّ جرم حبسوا قالوا انّهم ارادوا ان یجددوا الدّین
لو کان القدیم هو المختار عندکم لم ترکتم ما شرع فی التّوراة و الانجیل
تبیّنوا یا قوم لعمری لیس لکم الیوم من محیص ان کان هذا جرمی قد سبقنی فی
ذلک محمّد رسول اللّه و من قبله الرّوح و من قبله الکلیم و ان کان ذنبی اعلآء
کلمة اللّه و اظهار امره فاَنَا اوّل المذنبین لا ابدّل هذا الذّنب بملکوت ملک
السّموات و الأرضین انتهی

\end{otherlanguage}

An yā malik, qad kunnā bi-Ummi al-ʿIrāqi ilā an ḥumma al-firāqu
tawajjahnā ilā maliki al-Islāmi bi-amrihi. Fa-lammā ataynāhu warada
ʿalaynā min ulī al-nifāqi mā lā yatimmu bi-l-awrāq. Bi-dhālika nāḥa
sukkānu al-Firdawsi wa-ahlu ḥaẓāʾiri al-qudsi wa-lākinna al-qawma fī
ḥijābin ghalīẓ. Qul a-taʿtariḍūna ʿalā alladhī jāʾakum bi-bayyināti
Allāhi wa-burhānihi thumma ḥujjatihi wa-āyātihi? In hiya min tilqāʾi
nafsihi bal min ladun man baʿathahu wa-arsalahu bi-l-ḥaqqi wa-jaʿalahu
sirājan li-l-ʿālamīn. Qad ishtadda ʿalaynā al-amru fī kulli yawmin bal
fī kulli sāʿatin ilā an akhrajūnā mina al-sijni wa-adkhalūnā fī al-Sijni
al-Aʿẓami bi-ẓulmin mubīn. Idhā qīla bi-ayyi jurmin ḥubisū qālū: Innahum
arādū an yujaddidū al-dīn. Law kāna al-qadīmu huwa al-mukhtāra ʿindakum
lima taraktum mā shuriʿa fī al-Tawrāti wa-l-Injīl? Tabayyanū yā qawm.
La-ʿamrī laysa lakumu al-yawma min maḥīṣ. In kāna hādhā jurmī qad
sabaqanī fī dhālika Muḥammadun rasūlu Allāhi wa-min qablihi al-Rūḥu
wa-min qablihi al-Kalīm. Wa-in kāna dhanbī iʿlāʾa kalimati Allāhi
wa-iẓhāra amrihi fa-anā awwalu al-mudhnibīn. Lā ubaddilu hādhā al-dhanba
bi-malakūti mulki al-samāwāti wa-l-araḍīn. Intahā.

``O King, We were in the chief city of Iraq until the hour of departure
arrived, when We proceeded to the King of Islam at his command. When We
reached him, We suffered at the hands of the hypocrites what no pages
could fully record. At this the inhabitants of Paradise and the dwellers
in the courts of holiness lamented, yet the people remain behind a thick
veil. Say: Do you oppose the One who has come to you with God's clear
proofs and His testimony, His evidence and His verses? These do not
originate with Him, but with the One who raised Him up, sent Him forth
with the truth, and made Him a lamp for all the worlds. Our plight grew
more severe every day, indeed every hour, until they brought Us out of
prison and, with manifest injustice, cast Us into the Most Great Prison.
When asked, `For what crime have they been imprisoned?' they said, `They
wished to renew religion.' If what is old is what you prefer, why have
you abandoned what was prescribed in the Torah and the Gospel? Consider,
O people! By My life, today you have no escape. If this is My crime,
Muhammad, the Messenger of God, committed it before Me, and before Him
the Spirit, and before Him the One who conversed with God. If My sin is
to exalt the Word of God and manifest His Cause, then I am the foremost
of sinners. I would not exchange this sin for sovereignty over the
kingdom of the heavens and the earth.'' End of quotation.

\subsection{62}\label{section-2037}

\begin{otherlanguage}{persian}

تفصیل این لوح مبارک آنکه در ارض سرّ یک لوح امنع اقدس مخصوص ملک پاریس نازل
و ارسال شد و سبب تنزیل آنکه ملک در مجلس عامّ کلمه‌ئی گفته و آن اینکه ما از
برای تمشیت امور متعلّقه بخود نیستیم بلکه اگر مظلومی در عالم یافت شود و
ندا نماید جواب میگوئیم و حقّ او را از ظالمین اخذ مینمائیم اینمضمون کلمه‌ئی
بود که از او ظاهر شد و کلمهٴ دیگر در جواب ملک روس گفته هنگامیکه از او
سؤال نموده که ما و تو هر دو ملّة واحده هستیم چه ترا بر آن داشت که حمایت
اسلام نمودی و با او در محاربه متّحد شدی جواب گفته من اعانت اسلام ننمودم
بلکه ندای مظلومانیکه در بحر اسود کشتیهای آنها را شکستی و غرق نمودی مرا
از خواب بیدار نمود و بر رفع ظلم قیام نمودم

\end{otherlanguage}

Tafṣīl-e īn Lawḥ-e mobārak ānke dar Arż-e Serr yek Lawḥ-e amnaʿ-e aqdas
makhṣūṣ-e malek-e Pārīs nāzel va ersāl shod va sabab-e tanzīl ānke malek
dar majles-e ʿāmm kalame-ī gofte va ān īnke mā az barā-ye tamshīyat-e
omūr-e motaʿalleqe be-khod nīstīm balke agar maẓlūmī dar ʿālam yāft
shavad va nedā nemāyad javāb mī-gūʾīm va ḥaqq-e ū rā az ẓālemīn akhdh
mī-nemāʾīm. Īn mażmūn-e kalame-ī būd ke az ū ẓāher shod va kalame-ye
dīgar dar javāb-e malek-e Rūs gofte hengāmī ke az ū soʾāl nemūde ke mā
va to har do mellat-e vāḥede hastīm che to rā bar ān dāsht ke ḥemāyat-e
Eslām nemūdī va bā ū dar moḥārebe mottaḥed shodī. Javāb gofte man
eʿānat-e Eslām na-nemūdam balke nedā-ye maẓlūmānī ke dar Baḥr-e Asvad
keshtī-hā-ye ānhā rā shekastī va gharq nemūdī marā az khvāb bīdār nemūd
va bar rafʿ-e ẓolm qīyām nemūdam.

The circumstances of this blessed Tablet are as follows: In the Land of
Mystery, a most exalted and holy Tablet was revealed specifically for
the King of Paris and sent to him. The reason for its revelation was
that the king had made a statement in a public assembly, saying, ``We
are not here to manage affairs concerning ourselves alone. Rather, if
anyone in the world is oppressed and cries out, we answer and recover
his rights from the oppressors.'' This was the substance of a statement
he made. He made another statement in reply to the King of Russia, who
had asked him, ``You and I both belong to the same religious community.
What prompted you to support Islam and ally yourself with it in war?''
He replied, ``I did not aid Islam. Rather, the cry of the oppressed,
whose ships you shattered and sank in the Black Sea, awakened me from
sleep, and I arose to put an end to oppression.''

\subsection{63}\label{section-2038}

\begin{otherlanguage}{persian}

بعد از اصغاء این دو فقره لوح الهی نازل و در آن میفرماید دو کلمه در دو
مقام از تو استماع نمودیم که عرف عدل و انصاف از او متضوّع اگر حضرت ملک خود
را ناصر مظلومان عالم میدانند چرا بنصرت اینمظلومان که چشم عالم شبه آنرا
ندیده قیام ننمودند مع آنکه در سنین متوالیه ندای ایشان از هر ارضی مرتفعست
و در اینمدّت تحت سیاط ظلم و غضب بوده‌اند بقسمیکه هر صاحب بصری گریست و هر
صاحب قلبی نوحه نمود کم من ارض صبغت من دمائهم و کم من بلد ارتفع فیه
حنینهم و ضجیجهم و بکائهم در این لوح اکثر از مصائب وارده از قلم اعلی جاری
و نازل

\end{otherlanguage}

Baʿd az eṣghā-ye īn do faqare Lawḥ-e Elāhī nāzel va dar ān mī-farmāyad
do kalame dar do maqām az to estemāʿ nemūdīm ke ʿarf-e ʿadl va enṣāf az
ū motaḍavveʿ. Agar Ḥażrat-e Malek khod rā nāṣer-e maẓlūmān-e ʿālam
mī-dānand cherā be-noṣrat-e īn maẓlūmān ke chashm-e ʿālam shebh-e ān rā
na-dīde qīyām na-nemūdand maʿa ānke dar senīn-e motavālīye nedā-ye īshān
az har arżī mortafeʿ ast va dar īn moddat taḥt-e seyāṭ-e ẓolm va ghażab
būde-and be-qesmī ke har ṣāḥeb-e baṣarī gerīst va har ṣāḥeb-e qalbī
nawḥe nemūd. Kam min arḍin ṣubighat min dimāʾihim wa-kam min baladin
irtafaʿa fīhi ḥanīnuhum wa-ḍajījuhum wa-bukāʾuhum. Dar īn Lawḥ akthar az
maṣāʾeb-e vārede az Qalam-e Aʿlā jārī va nāzel.

After these two statements were heard, the divine Tablet was revealed.
In it He says, ``We heard two statements from you on two occasions, from
which the fragrance of justice and fairness wafted. If His Majesty the
King considers himself a champion of the oppressed of the world, why has
he not arisen to aid these oppressed ones, whose like the eye of the
world has never seen, even though for successive years their cries have
risen from every land? Throughout this time they have been under the
scourges of oppression and wrath, so that everyone with eyes has wept
and everyone with a heart has lamented. How many lands have been dyed
with their blood, and in how many cities have their lamentations, their
outcries, and their weeping risen!'' In this Tablet, most of the
afflictions that had befallen them flowed from the Most Exalted Pen and
were revealed.

\subsection{64}\label{section-2039}

\begin{otherlanguage}{persian}

و بعد از ارسال مدّتی جواب نیامد بالأخره وزیرش عریضه عرض نمود باینمضمون من
نامهٴ شما را با وزیر خارجه بحضور امپراطور بردیم و تفصیل را ذکر نمودیم
نامه را اخذ نمود و پهلوی خود گذاشت و تا حال جوابی در اینفقره از او ظاهر
نه لذا لوح ثانی از سماء مشیّت الهی نازل و شخصی از ملأ ابن که باقبال فائز
بود عرض نمود بلسان فرنساوی ترجمه نموده ارسال داشتیم العلم عنداللّه

\end{otherlanguage}

Va baʿd az ersāl moddatī javāb nayāmad. Bi-l-ākhare vazīrash ʿarīże ʿarż
nemūd be-īn mażmūn: Man nāme-ye shomā rā bā vazīr-e khāreje be-ḥożūr-e
emperāṭūr bordīm va tafṣīl rā dhekr nemūdīm. Nāme rā akhdh nemūd va
pahlū-ye khod goẕāsht va tā ḥāl javābī dar īn faqare az ū ẓāher na.
Ledhā Lawḥ-e thānī az samā-ye mashīyyat-e Elāhī nāzel va shakhṣī az
malaʾ-e Ebn ke be-eqbāl fāʾez būd ʿarż nemūd be-lesān-e Farānsavī
tarjome nemūde ersāl dāshtīm. Al-ʿilmu ʿinda Allāh.

For some time after it was sent, no reply came. Eventually his minister
submitted a letter to this effect: ``The foreign minister and I brought
your letter into the emperor's presence and explained the circumstances.
He took the letter and laid it beside him, and so far he has given no
reply on this matter.'' Therefore a second Tablet was revealed from the
heaven of the divine will, and a member of the community of the Son who
had attained to recognition submitted, ``We translated it into French
and sent it.'' Knowledge rests with God.

\subsection{65}\label{section-2040}

\begin{otherlanguage}{persian}

اینست که میفرماید ای ملک ما شنیدیم از تو کلمه‌ئی را که تکلّم کردی بآن
هنگامیکه سؤال نمود ترا ملک روس از آنچه واقع شد از حکم حرب جواب گفتی در
مهد خوابیده بودم بیدار کرد مرا ندای عباد مظلومیکه غرق شدند در بحر اسود
شهادت میدهیم ندا ترا بیدار ننمود بلکه هوای نفس چه که امتحان کردیم یافتیم
ترا فی معزل یعنی در کناری از کلمه الی ان قال عزّ بیانه اگر بودی تو صاحب
کلمه ما نبذت کتاب اللّه ورآء ظهرک

\end{otherlanguage}

Īn ast ke mī-farmāyad: Ey malek, mā shenīdīm az to kalame-ī rā ke
takallom kardī be-ān hengāmī ke soʾāl nemūd to rā malek-e Rūs az ānche
vāqeʿ shod az ḥokm-e ḥarb. Javāb goftī dar mahd khvābīde būdam bīdār
kard marā nedā-ye ʿebād-e maẓlūmī ke gharq shodand dar Baḥr-e Asvad.
Shahādat mī-dahīm nedā to rā bīdār na-nemūd balke havā-ye nafs che ke
emteḥān kardīm yāftīm to rā fī maʿzil, yaʿnī dar kenārī az kalame. Ilā
an qāla, ʿazza bayānuhu: Agar būdī to ṣāḥeb-e kalame mā nabadhta kitāba
Allāhi warāʾa ẓahrika.

This is why He says, ``O King, We heard a statement you made when the
King of Russia questioned you about what had occurred through the decree
of war. You replied, `I was asleep in my bed when I was awakened by the
cry of the oppressed servants who drowned in the Black Sea.' We testify
that it was not the cry that awakened you, but selfish desire, for We
tested you and found you aloof''---that is, standing apart from your
word. Further on He says---exalted be His utterance: ``Had you been true
to your word, you would not have cast the Book of God behind your
back.''

\subsection{66}\label{section-2041}

\begin{otherlanguage}{persian}

بعد میفرماید بما فعلت بآنچه عمل نمودی مختلف میشود امور در مملکت تو و
بیرون میرود سلطنت از کفّ تو جزای عمل تو در آنوقت مییابی نفس خود را در
زیانی آشکار و اخذ مینماید زلازل جمیع قبائل را در آنجا و از این امور مفرّی
نیست مگر آنکه بر نصرت امر قیام نمائی و حضرت روح را متابعت کنی در این
سبیل مستقیم و همچنین میفرماید آیا عزّت تو ترا مغرور نمود لعمری این عزّت
پاینده نیست زود است که زایل شود مگر آنکه متمسّک شوی بحبل محکم حقّ جلّ جلاله
بعد میفرماید میبینیم ما ذلّت را که بتعجیل از ورای تو حرکت مینماید و تو از
نائمین هستی یعنی در خوابی و مطّلع نیستی

\end{otherlanguage}

Baʿd mīfarmāyad: bi-mā faʿalta, bi-ānchi ʿamal namūdī, mukhtalif
mīshavad umūr dar mamlakat-i tu va bīrūn mīravad salṭanat az kaff-i tu,
jazā-yi ʿamal-i tu. Dar ān vaqt mīyābī nafs-i khud rā dar ziyānī āshkār
va akhdh mīnamāyad zalāzil jamīʿ-i qabāʾil rā dar ānjā va az īn umūr
mafarrī nīst magar ānki bar nuṣrat-i amr qiyām namāʾī va Ḥaḍrat-i Rūḥ rā
mutābaʿat kunī dar īn sabīl-i mustaqīm. Va hamchunīn mīfarmāyad: āyā
ʿizzat-i tu turā maghrūr namūd? La-ʿamrī, īn ʿizzat pāyandih nīst, zūd
ast ki zāʾil shavad magar ānki mutamassik shavī bi-ḥabl-i muḥkam-i Ḥaqq,
jalla jalāluhu. Baʿd mīfarmāyad: mībīnīm mā dhillat rā ki bi-taʿjīl az
varā-yi tu ḥarakat mīnamāyad va tu az nāʾimīn hastī, yaʿnī dar khwābī va
muṭṭaliʿ nīstī.

He then says, ``Because of what you have done, on account of your
actions, affairs in your kingdom will fall into disorder, and
sovereignty will pass from your grasp as the recompense for your deeds.
At that time you will find yourself in manifest loss, and earthquakes
will seize all the peoples there. There is no escape from these things
unless you arise to support the Cause and follow the Spirit on this
straight path.'' He also says, ``Has your glory made you proud? By My
life, this glory will not endure. Soon it will pass away unless you hold
fast to the firm cord of God, exalted be His glory.'' He then says, ``We
see abasement hastening toward you from behind, while you are among
those who sleep''---that is, you are asleep and unaware.

\subsection{67}\label{section-2042}

\begin{otherlanguage}{persian}

حال ملاحظه نمائید این امورات را بکمال تصریح از قبل اخبار فرموده‌اند و
اکثری ظهور این وقایع را ممتنع میدانستند و محال میشمردند چنانچه وقتی شخصی
از اینعبد سؤال نمود که یعنی امپراطور پاریس عزل میشود چگونه سلطنت از او
خارج میگردد و عزّت او چگونه زایل میشود و حال آنکه رأس ملوک عالم است این
بنده عرض نمود که اینعبد آگاه نیست و مطّلع نه ولکن اینقدر میدانم که آنچه
از لسان عظمت نازل جمیع در ظاهر ظاهر واقع میشود میبینید و میشنوید

\end{otherlanguage}

Ḥāl mulāḥiẓih namāʾīd īn umūrāt rā bi-kamāl-i taṣrīḥ az qabl ikhbār
farmūdih-and va aktharī ẓuhūr-i īn vaqāyiʿ rā mumtaniʿ mīdānistand va
muḥāl mīshumurdand. Chunānchi vaqtī shakhṣī az īn ʿabd suʾāl namūd ki
yaʿnī impirāṭūr-i Pārīs ʿazl mīshavad? Chigūnih salṭanat az ū khārij
mīgardad va ʿizzat-i ū chigūnih zāʾil mīshavad va ḥāl ānki raʾs-i
mulūk-i ʿālam ast? Īn bandih ʿarḍ namūd ki īn ʿabd āgāh nīst va muṭṭaliʿ
nah, va-līkin īn qadr mīdānam ki ānchi az lisān-i ʿaẓamat nāzil, jamīʿ
dar ẓāhir-i ẓāhir vāqiʿ mīshavad. Mībīnīd va mīshinavīd.

Now consider: He foretold these events with the utmost explicitness,
although most people regarded their occurrence as impossible and
inconceivable. For example, someone once asked this servant, ``Does this
mean that the emperor of Paris will be deposed? How will sovereignty
pass from him, and how will his glory fade, when he is the foremost of
the world's kings?'' This servant replied, ``This servant has no
knowledge or information, but I know this much: whatever has been
revealed by the Tongue of Grandeur will all come to pass in the most
outward reality. You will see and hear.''

\subsection{68}\label{section-2043}

\begin{otherlanguage}{persian}

باری مقصود خادم از این اذکار آنکه هر حزبی بر آیات و بیّنات حقّ مطّلع و آگاه
شود و بعلم الیقین بداند که آنچه از بعد ظاهر از قبل از قلم اعلی نازل و
چون الواح ملوک که در سورهٴ هیکل نازل و همچنین لوح رئیس نزد اکثری چند سنه
قبل از وقوع موجود لذا ذکر شد والّا اگر اینعبد بخواهد از آیات ظاهره و
بیّنات باهره ذکر نماید دفاتر عدیده لازم و شاید بعضی در قبول آن توقّف
نمایند

\end{otherlanguage}

Bārī maqṣūd-i khādim az īn adhkār ānki har ḥizbī bar āyāt va bayyināt-i
Ḥaqq muṭṭaliʿ va āgāh shavad va bi-ʿilm al-yaqīn bidānad ki ānchi az
baʿd ẓāhir az qabl az Qalam-i Aʿlā nāzil. Va chūn alvāḥ-i mulūk ki dar
Sūrih-yi Haykal nāzil va hamchunīn Lawḥ-i Raʾīs nazd-i aktharī chand
sanih qabl az vuqūʿ mawjūd, lidhā dhikr shud. Va-illā agar īn ʿabd
bikhwāhad az āyāt-i ẓāhirih va bayyināt-i bāhirih dhikr namāyad
dafātir-i ʿadīdih lāzim va shāyad baʿḍī dar qabūl-i ān tavaqquf
namāyand.

In any case, this servant's purpose in recounting these things is that
every group may become informed and aware of the signs and clear proofs
of God, and know with certainty that whatever subsequently appeared had
already been revealed by the Most Exalted Pen. Since the Tablets to the
Kings revealed in the Sūrih of the Temple, as well as the Tablet to the
Chief, were in the possession of many people several years before these
events occurred, they have been mentioned here. Otherwise, were this
servant to recount the manifest signs and brilliant proofs, many volumes
would be needed, and some might hesitate to accept them.

\subsection{69}\label{section-2044}

\begin{otherlanguage}{persian}

باری از کلّ استدعا مینمایم ببصر حقّ در حقّ نظر نمایند چه که کلّ باینفقره
مأمورند و این مخصوص باینظهور است نقطهٴ بیان روح ما سواه فداه میفرماید
بچشم او او را ببینید اگر بچشم غیر ملاحظه کنید هرگز بشناسائی و آگاهی فائز
نشوید

\end{otherlanguage}

Bārī az kull istidʿā mīnamāyam bi-baṣar-i Ḥaqq dar Ḥaqq naẓar namāyand,
chih ki kull bi-īn fiqrih maʾmūrand va īn makhṣūṣ bi-īn ẓuhūr ast.
Nuqṭih-yi Bayān, rūḥ-i mā sivāh fidāh, mīfarmāyad: bi-chashm-i ū ū rā
bibīnīd; agar bi-chashm-i ghayr mulāḥiẓih kunīd hargiz bi-shināsāʾī va
āgāhī fāʾiz nashavīd.

In any case, I ask everyone to look upon God with God's own vision, for
all are commanded to do this, and this is particular to this Revelation.
The Point of the Bayān---may the souls of all besides Him be sacrificed
for Him---says, ``See Him with His own eyes. If you look with the eyes
of another, you will never attain recognition and understanding.''

\subsection{70}\label{section-2045}

\begin{otherlanguage}{persian}

مشرکین به حیل و مکری ظاهر شده‌اند که شبه آن دیده نشده و جمیع این امورات
را از قبل اخبار فرموده‌اند طوبی از برای مستقیمیکه ضوضاء و اشارات و شبهات
اهل عالم را بمثابهٴ طنین ذباب داند ظهور عظیم و امر عظیم از حقّ میطلبم
عباد را بر آنچه سزاوار است مؤیّد فرماید

\end{otherlanguage}

Mushrikīn bi-ḥiyal va makrī ẓāhir shudih-and ki shibh-i ān dīdih
nashudih va jamīʿ-i īn umūrāt rā az qabl ikhbār farmūdih-and. Ṭūbā az
barā-yi mustaqīmī ki ḍawḍāʾ va ishārāt va shubuhāt-i ahl-i ʿālam rā
bi-mathābih-yi ṭanīn-i dhubāb dānad. Ẓuhūr ʿaẓīm va amr ʿaẓīm. Az Ḥaqq
mīṭalabam ʿibād rā bar ānchi sazāvār ast muʾayyad farmāyad.

Those who associate partners with God have appeared with stratagems and
deceit the like of which has never been seen, and He foretold all these
things. Blessed is the steadfast person who regards the clamor,
insinuations, and doubts of the people of the world as the buzzing of
flies. Great is the Revelation, and great is the Cause. I ask God to
assist His servants to do what is fitting.

\subsection{71}\label{section-2046}

\begin{otherlanguage}{persian}

امروز فرقان بافصح بیان میفرماید ای غافلان امّ الکتاب ناطق و افق ایقان
منوّر برهان مشرق و دلیل واضح بظنون از مطلع نور یقین محروم نمانید و باوهام
از مالک انام ممنوع نشوید امثال این بیانات از منزل آیات بکرّات استماع شده
و همچنین اینکلمات از لسان عظمت استماع گشت قوله تبارک و تعالی

\end{otherlanguage}

Imrūz Furqān bi-afṣaḥ-i bayān mīfarmāyad: ay ghāfilān, Umm al-Kitāb
nāṭiq va ufuq-i īqān munavvar, burhān mushriq va dalīl vāḍiḥ. Bi-ẓunūn
az maṭlaʿ-i nūr-i yaqīn maḥrūm namānīd va bi-awhām az Mālik-i Anām
mamnūʿ nashavīd. Amthāl-i īn bayānāt az Munzil-i Āyāt bi-karrāt istimāʿ
shudih va hamchunīn īn kalimāt az Lisān-i ʿAẓamat istimāʿ gasht, qawluhu
tabāraka va taʿālā:

Today the Qurʾān declares in the most eloquent language, ``O heedless
ones! The Mother Book is speaking, and the horizon of certitude is
illumined. The proof shines forth, and the evidence is clear. Do not let
conjectures deprive you of the Dayspring of the light of certainty, or
let vain imaginings keep you from the Lord of all people.'' Utterances
like these have repeatedly been heard from the Revealer of the verses.
These words were also heard from the Tongue of Grandeur---His words,
blessed and exalted be He:

\subsection{72}\label{section-2047}

\begin{otherlanguage}{arabic}

امروز بیان بکمال شفقت و عنایت ندا مینماید و میفرماید ای اهل من و اصحاب
من امروز مالک غیب و شهود بر عرش ظهور مستوی و امّ البیان ظاهر بظنون خود را
از رحیق مختوم منع ننمائید و از کوثر یقین محروم مسازید

\end{otherlanguage}

Imrūz Bayān bi-kamāl-i shafaqat va ʿināyat nidā mīnamāyad va mīfarmāyad:
ay ahl-i man va aṣḥāb-i man, imrūz Mālik-i ghayb va shuhūd bar ʿarsh-i
ẓuhūr mustavī va Umm al-Bayān ẓāhir. Bi-ẓunūn khud rā az raḥīq-i makhtūm
manʿ nanamāʾīd va az kawthar-i yaqīn maḥrūm masāzīd.

``Today the Bayān calls out with the utmost compassion and loving care,
saying, `O My people and My companions! Today the Lord of the unseen and
the seen is established upon the throne of revelation, and the Mother of
the Bayān is manifest. Do not, through conjectures, deny yourselves the
sealed wine or deprive yourselves of the river of certainty.'\,''

\subsection{73}\label{section-2048}

\begin{otherlanguage}{persian}

و وقتی از اوقات لسان عظمت باینکلمهٴ مبارکه ناطق اگر بصر ببزرگی آسمان
باشد وبخواهد اقلّ عمّا یحصی بشطر دیگر توجّه نماید قابل ورود این بساط نه

\end{otherlanguage}

Va vaqtī az awqāt Lisān-i ʿAẓamat bi-īn kalimih-yi mubārakih nāṭiq: agar
baṣar bi-buzurgī-yi āsmān bāshad va bikhwāhad aqall-i ʿammā yuḥṣā
bi-shaṭr-i dīgar tavajjuh namāyad qābil-i vurūd-i īn bisāṭ nah.

On one occasion the Tongue of Grandeur spoke these blessed words: ``Even
if an eye were as vast as the heavens, should it wish to turn in another
direction to a degree too small to reckon, it would not be worthy to
enter this court.''

\subsection{74}\label{section-2049}

\begin{otherlanguage}{arabic}

سبحانک یا من فی قبضتک زمام المعانی و البیان اسئلک باسمک الّذی به اشرقت
شمس القدم من افق العالم بأن تحفظ اصفیائک و احبّائک من مکر الّذین کفروا بک
و بآیاتک و اعرضوا عمّا امرتهم به فی کتبک و صحفک و زبرک و الواحک ای ربّ تری
عبادک الجهلآء الّذین جعلوا الخدعه سلاحهم و الکذب معاذهم و ملاذهم یفترون
علیک فی اللّیالی و الأیّام بما ناح به سکّان خبآء مجدک و عظمتک ای ربّ اسئلک
بأنوار وجهک و اسرار امرک بأن تنزل علیهم ما یعرّفهم و یعلّمهم سبیلک و ما
عندک انّک انت المقتدر الکافی و المهیمن الباقی العلیم الحکیم

\end{otherlanguage}

Subḥānaka yā man fī qabḍatika zimām al-maʿānī wa-l-bayān. Asʾaluka
bi-smika alladhī bihi ashraqat shams al-qidam min ufuq al-ʿālam bi-an
taḥfaẓa aṣfiyāʾaka wa-aḥibbāʾaka min makr alladhīna kafarū bika
wa-bi-āyātika wa-aʿraḍū ʿammā amartahum bihi fī kutubika wa-ṣuḥufika
wa-zuburika wa-alwāḥika. Ay Rabb, tarā ʿibādaka al-juhalāʾ alladhīna
jaʿalū al-khudʿata silāḥahum wa-l-kadhiba maʿādhahum wa-malādhahum
yaftarūna ʿalayka fī al-layālī wa-l-ayyām bi-mā nāḥa bihi sukkān khibāʾ
majdika wa-ʿaẓamatika. Ay Rabb, asʾaluka bi-anwār wajhika wa-asrār
amrika bi-an tunazzila ʿalayhim mā yuʿarrifuhum wa-yuʿallimuhum sabīlaka
wa-mā ʿindaka. Innaka anta al-Muqtadir al-Kāfī wa-l-Muhaymin al-Bāqī
al-ʿAlīm al-Ḥakīm.

Glory be to You, O You in whose grasp are the reins of meaning and
utterance! I ask You, by Your name through which the Sun of Eternity has
shone from the horizon of the world, to protect Your chosen ones and
Your loved ones from the deceit of those who have disbelieved in You and
Your signs and turned away from what You commanded them in Your books,
scriptures, writings, and tablets. O Lord, You see Your ignorant
servants, who have made deception their weapon and falsehood their
refuge and shelter, fabricating against You by night and day things that
have caused the dwellers in the pavilion of Your glory and grandeur to
lament. O Lord, I ask You, by the lights of Your face and the mysteries
of Your Cause, to send down upon them what will acquaint them with, and
teach them, Your path and what is with You. You are indeed the
All-Powerful, the All-Sufficing, the All-Protecting, the Everlasting,
the All-Knowing, the All-Wise.

\subsection{75}\label{section-2050}

\begin{otherlanguage}{arabic}

یا اولیآء اللّه همّت نمائید که مجدّد نفوس موهومه ناس بیچاره را بمثابهٴ قبل
در بئر ظنون مبتلا ننمایند این نفوسند که در حقّشان میفرماید قوله تعالی
بدّلوا نعمة اللّه کفراً و احلّوا قومهم دار البوار این نفوسند که لازال نعمت
الهی را تبدیل نمودند و خلق را از صراط مستقیم منحرف ساختند معدودی از قبل
در کور فرقان بهواهای نفسانیّه در اضلال نفوس جهدها نمودند در جمیع اعصار
اینگونه نفوس بوده و خواهند بود

\end{otherlanguage}

Yā awliyāʾ Allāh, himmat namāʾīd ki mujaddad nufūs-i mawhūmih nās-i
bīchārih rā bi-mathābih-yi qabl dar biʾr-i ẓunūn mubtalā nanamāyand. Īn
nufūsand ki dar ḥaqqishān mīfarmāyad, qawluhu taʿālā: baddalū niʿmata
Allāhi kufran wa-aḥallū qawmahum dāra al-bawār. Īn nufūsand ki lāzāl
niʿmat-i ilāhī rā tabdīl namūdand va khalq rā az ṣirāṭ-i mustaqīm
munḥarif sākhtand. Maʿdūdī az qabl dar kawr-i Furqān bi-havāhā-yi
nafsāniyyih dar iḍlāl-i nufūs jahdhā namūdand. Dar jamīʿ-i aʿṣār īn
gūnih nufūs būdih va khwāhand būd.

O friends of God! Make every effort to prevent deluded souls from once
again casting the hapless people into the pit of conjectures, as they
did before. These are the souls of whom He says---His words, exalted be
He: ``They exchanged God's favor for unbelief and settled their people
in the abode of ruin.'' These are the souls who have always altered the
divine bounty and turned people away from the straight path. In the
past, during the Qurʾānic dispensation, a few people, driven by selfish
desires, made strenuous efforts to lead souls astray. Such people have
existed in every age and will continue to exist.

\subsection{76}\label{section-2051}

\begin{otherlanguage}{persian}

در این ایّام اوراقی وارد بعد از ملاحظه و مشاهده واضح و مبرهن شد که از
جوهر غل و بغضا ظاهر گشته بیانات الهی که نعمت حقیقی است بعضی را تبدیل
نموده ثعبان بالوان جدید ظاهر و بکمال دوستی خواسته عباد را از افق اعلی
منع نماید و از بحر یقین محروم سازد و بغدیر وهم دعوت کند

\end{otherlanguage}

Dar īn ayyām awrāqī vārid. Baʿd az mulāḥiẓih va mushāhidih vāḍiḥ va
mubarhan shud ki az jawhar-i ghill va baghḍā ẓāhir gashtih. Bayānāt-i
ilāhī ki niʿmat-i ḥaqīqī ast baʿḍī rā tabdīl namūdih. Thuʿbān bi-alvān-i
jadīd ẓāhir va bi-kamāl-i dūstī khwāstih ʿibād rā az ufuq-i aʿlā manʿ
namāyad va az baḥr-i yaqīn maḥrūm sāzad va bi-ghadīr-i vahm daʿvat
kunad.

Some papers arrived in these days. Upon examining and reading them, it
became clear and evident that they had issued from the very essence of
malice and hatred. He has altered some of the divine utterances, which
are the true bounty. The serpent has appeared in new colors and, with
every show of friendship, has sought to keep the servants from the Most
Exalted Horizon, deprive them of the ocean of certainty, and summon them
to the pool of delusion.

\subsection{77}\label{section-2052}

\begin{otherlanguage}{arabic}

بگوئید ای بی‌انصاف اعتساف را بگذار و بحبل عدل تمسّک جو تو کجا بودی وقتیکه
ظلمت ظلم عالم را احاطه نموده بود و از سطوت ظالمین کلّ مستور

\end{otherlanguage}

Bigūʾīd: ay bī-inṣāf, iʿtisāf rā bigudhār va bi-ḥabl-i ʿadl tamassuk jū.
Tu kujā būdī vaqtī ki ẓulmat-i ẓulm ʿālam rā iḥāṭih namūdih būd va az
saṭvat-i ẓālimīn kull mastūr?

Say, ``O unjust one! Abandon injustice and hold fast to the cord of
justice. Where were you when the darkness of oppression enveloped the
world and everyone was in hiding for fear of the oppressors' might?''

\subsection{78}\label{section-2053}

\begin{otherlanguage}{persian}

در کلماتیکه ذکر نموده سمّ مهلک مکنون مثلاً در اوّل نامه‌اش ذکر اسم اعظم را
نموده نوشته بسم اللّه الأبهی الأبهی لا اله الّا هو قل کلّ الیه یرجعون
سبحان‌اللّه اسم حقّ را از برای خود شَرَکی قرار داده و دامی گسترده باینکلمه
خواسته موقنین و موحّدین و مخلصین را ساکن نماید و بعد سمّ ناقع را که در
انیاب کلماتش مستور است بچشاند انسان متحیّر است از حیله و مکر امثال آن
نفوس

\end{otherlanguage}

Dar kalimātī ki dhikr namūdih samm-i muhlik maknūn. Mathalan dar avval-i
nāmih-ash dhikr-i Ism-i Aʿẓam rā namūdih, nivishtih: bi-smi Allāhi
al-Abhā al-Abhā, lā ilāha illā Huwa, qul kullun ilayhi yarjiʿūn. Subḥān
Allāh! Ism-i Ḥaqq rā az barā-yi khud sharakī qarār dādih va dāmī
gustardih. Bi-īn kalimih khwāstih mūqinīn va muvaḥḥidīn va mukhliṣīn rā
sākin namāyad va baʿd samm-i nāqiʿ rā ki dar anyāb-i kalimātash mastūr
ast bichashānad. Insān mutaḥayyir ast az ḥīlih va makr-i amthāl-i ān
nufūs.

A deadly poison lies concealed in the words he has written. For example,
at the beginning of his letter he has mentioned the Most Great Name,
writing, ``In the name of God, the All-Glorious, the All-Glorious. There
is no god but Him. Say: All return to Him.'' Glory be to God! He has
made the name of God a snare for himself and spread a trap. With these
words he has sought to put those who are certain, who affirm God's
oneness, and who are sincere at ease, and then make them taste the
lethal poison hidden in the fangs of his words. One is bewildered by the
trickery and deceit of such people.

\subsection{79}\label{section-2054}

\begin{otherlanguage}{persian}

معترف باینکلمه فارغ است از ذکر عالم و ضوضاء امم و اقوال ناس و اختلافات
قوم چه که الیوم یوم اللّه است لا ینبغی فیه ذکر شیء اینست یومیکه میفرماید
کلّشیء هالک الّا وجهه اهل بها امروز بر سفائن حمرا راکب و این سُفُن بر بحر و
برّ جاری و ساری طوبی لمن تمسّک بها منقطعاً عن الدّنیا و ذکرها و ما فیها

\end{otherlanguage}

Muʿtarif bi-īn kalimih fārigh ast az dhikr-i ʿālam va ḍawḍāʾ-i umam va
aqvāl-i nās va ikhtilāfāt-i qawm, chih ki al-yawm Yawm Allāh ast, lā
yanbaghī fīhi dhikru shayʾ. Īn ast yawmī ki mīfarmāyad: kullu shayʾin
hālikun illā wajhahu. Ahl-i Bahā imrūz bar safāʾin-i ḥamrā rākib va īn
sufun bar baḥr va barr jārī va sārī. Ṭūbā li-man tamassaka bihā
munqaṭiʿan ʿan al-dunyā wa-dhikrihā wa-mā fīhā.

Whoever acknowledges these words is free from concern with the talk of
the world, the clamor of nations, the sayings of people, and the
disputes of the community, for today is the Day of God, in which it is
not fitting to mention anything else. This is the Day of which He says,
``All things perish except His face.'' Today the people of Bahā are
aboard the crimson arks, and these vessels move and travel over sea and
land. Blessed is whoever holds fast to them, detached from the world,
its mention, and all that is in it.

\subsection{80}\label{section-2055}

\begin{otherlanguage}{persian}

باری شخص مذکور بعد از ذکر اسم حقّ جلّ جلاله که بگمان خود دام قرار داده بعض
ذکرها که لایق ذکر نیست نموده چه که ذکر عروج خود را بمعارج علوم نوشته از
نحو و صرف و معانی و بیان و منطق و جواهر و اعراض و فقه و اصول و حکمت مشّاء
و اشراق و حکمت مرحوم شیخ و خدمت حاجی کریم خان) و از این اذکار خواسته
اظهار فضیلت نماید اگرچه در ظاهر نوشته از این علوم چیزی نفهمیدم ایکاش
بقدریکه خود را مشغول تحریر اینکلمات نموده ببصر الهی در آیاتش ملاحظه
مینمود و در اسرار علم و حکمت ربّانی تفکّر میکرد

\end{otherlanguage}

Bārī shakhṣ-i madhkūr baʿd az dhikr-i ism-i Ḥaqq, jalla jalāluhu, ki
bi-gumān-i khud dām qarār dādih, baʿḍ dhikrhā ki lāyiq-i dhikr nīst
namūdih, chih ki dhikr-i ʿurūj-i khud rā bi-maʿārij-i ʿulūm nivishtih,
az naḥv va ṣarf va maʿānī va bayān va manṭiq va javāhir va aʿrāḍ va fiqh
va uṣūl va ḥikmat-i Mashshāʾ va Ishrāq va ḥikmat-i marḥūm-i Shaykh va
khidmat-i Ḥājī Karīm Khān) va az īn adhkār khwāstih iẓhār-i faḍīlat
namāyad, agarchih dar ẓāhir nivishtih az īn ʿulūm chīzī nafahmīdam. Ay
kāsh bi-qadrī ki khud rā mashghūl-i taḥrīr-i īn kalimāt namūdih
bi-baṣar-i ilāhī dar āyātash mulāḥiẓih mīnamūd va dar asrār-i ʿilm va
ḥikmat-i rabbānī tafakkur mīkard.

In any case, after mentioning the name of God, exalted be His glory,
which he imagines he has made into a trap, the person in question has
said things unworthy of mention. He has written of his ascent through
the stages of learning: syntax, morphology, semantics, rhetoric, logic,
substances and accidents, jurisprudence and its principles, Peripatetic
and Illuminationist philosophy, the philosophy of the late Shaykh, and
service to Ḥājī Karīm Khān. By mentioning these things he has sought to
display his learning, although he has outwardly written, ``I understood
nothing of these sciences.'' If only, for as much time as he occupied
himself with writing these words, he had examined God's verses with
divine vision and reflected on the mysteries of divine knowledge and
wisdom.

\subsection{81}\label{section-2056}

\begin{otherlanguage}{arabic}

علمای فرقان هزار و دویست سنه از این عبارات نوشتند و در ردّ و اثبات یکدیگر
آنچه آن نفوس از زرع الفاظ حصاد نمودند و ثمر حاصل کردند این نفوس هم
خواهند نمود

\end{otherlanguage}

ʿUlamā-yi Furqān hizār va duvīst sanih az īn ʿibārāt nivishtand va dar
radd va ithbāt-i yakdīgar. Ānchi ān nufūs az zarʿ-i alfāẓ ḥiṣād namūdand
va thamar ḥāṣil kardand īn nufūs ham khwāhand namūd.

For twelve hundred years the scholars of the Qurʾān wrote expressions of
this kind, refuting and affirming one another. Whatever those people
reaped from sowing words, and whatever fruit they obtained, these people
too will obtain.

\subsection{82}\label{section-2057}

\begin{otherlanguage}{arabic}

باری در کمین اولیا و دوستان الهی نشسته و باسباب محبّت و مهربانی و کلمات
لایسمنهٴ لایغنیه در اضلال خلق بیچاره مشغول تا آنکه اظهار تصدیق بشجرهٴ
بیان نموده و نوشته آنچه از کلمات نقطهٴ بیان دیدم حیرت بر حیرتم افزود و
باب معرفتی برایم نگشود

\end{otherlanguage}

Bārī dar kamīn-i awliyā va dūstān-i ilāhī nishastih va bi-asbāb-i
maḥabbat va mihrabānī va kalimāt-i lā yusminuhu lā yughnīhi dar iḍlāl-i
khalq-i bīchārih mashghūl, tā ānki iẓhār-i taṣdīq bi-Shajari h-yi Bayān
namūdih va nivishtih: ānchi az kalimāt-i Nuqṭih-yi Bayān dīdam ḥayrat
bar ḥayratam afzūd va bāb-i maʿrifatī barāyam nagushūd.

In any case, he lies in wait for God's devoted ones and friends and,
through displays of love and kindness and words that neither nourish nor
avail, is busy leading the hapless people astray. He has even professed
belief in the Tree of the Bayān and written, ``Whatever I read of the
words of the Point of the Bayān added bewilderment to my bewilderment
and opened no door of understanding for me.''

\subsection{83}\label{section-2058}

\begin{otherlanguage}{persian}

نه اینکه العیاذ باللّه آنها محتوی بر مطالب عالیه نباشد بلکه من باب آنستکه
چونکه این امر از قرار بیان خود صاحب بیان کنون در مقام نطفه است خیلی کس
میخواهد که احکام کلّشیء را ظاهراً باهراً مشهوداً مشهوراً از نطفه بفهمد یا
آنکه بحالات واردهٴ بر آن برسد الّا من علّمه اللّه علم کلّشیء این کلمات آنشخص
غافل این خادم فانی عرض مینماید

\end{otherlanguage}

Nah īnki, al-ʿiyādhu bi-Allāh, ānhā muḥtavī bar maṭālib-i ʿāliyih
nabāshad, balkih min bāb-i ān ast ki chūnki īn amr az qarār-i bayān-i
khud-i Ṣāḥib-i Bayān kunūn dar maqām-i nuṭfih ast, khaylī kas mīkhwāhad
ki aḥkām-i kull-i shayʾ rā ẓāhiran bāhiran mashhūdan mashhūran az nuṭfih
bifahmad yā ānki bi-ḥālāt-i vāridih-yi bar ān birasad, illā man
ʿallamahu Allāhu ʿilma kulli shayʾ. Īn kalimāt-i ān shakhṣ-i ghāfil. Īn
khādim-i fānī ʿarḍ mīnamāyad:

``Not that---God forbid---they do not contain exalted matters. Rather,
it is because this Cause, according to the words of the Author of the
Bayān Himself, is now at the stage of the embryo. It takes a very
exceptional person to understand from the embryo the laws governing all
things, outwardly, clearly, perceptibly, and manifestly, or to apprehend
the conditions that befall it---except one whom God has taught the
knowledge of all things.'' These are the words of that heedless person.
This humble servant submits:

\subsection{84}\label{section-2059}

\begin{otherlanguage}{persian}

یا ملأ البیان و یا ملأ الفرقان انوار شمس ظهور عالم را احاطه نموده و
اشراقات و تجلّیاتش ظاهر و باهر نقطهٴ بیان روح ما سواه فداه میفرماید نطفهٴ
یکسالهٴ یوم ظهور او اقوی است از کلّ بیان روز روزیست که جمیع کتب و صحف و
زبر او را بحقّ نسبت داده‌اند و کلّ را بآن بشارت فرموده‌اند باید کلّ بابصار
حدیده و آذان مقدّسهٴ از کلّ اذکار بافق اعلی توجّه نمایند و ندائیکه از یمین
عرش اعظم مرتفع است اصغا کنند

\end{otherlanguage}

Yā malaʾ al-Bayān va yā malaʾ al-Furqān, anvār-i shams-i ẓuhūr ʿālam rā
iḥāṭe namūde va ishrāqāt va tajalliyātash ẓāhir va bāhir. Nuqṭe-yi
Bayān, rūḥu mā siwāhu fidāhu, mīfarmāyad: nuṭfe-yi yeksāle-yi yawm-i
ẓuhūr-i ū aqvā ast az kull-i Bayān. Rūz rūzīst ke jamīʿ-i kutub va ṣuḥuf
va zubur ū rā bi-Ḥaqq nisbat dāde-and va kull rā bi-ān bishārat
farmūde-and. Bāyad kull bi-abṣār-i ḥadīde va ādhān-i muqaddase-yi az
kull-i adhkār bi-ufuq-i aʿlā tavajjuh namāyand va nidāʾī ke az yamīn-i
ʿarsh-i aʿẓam murtafiʿ ast iṣghā kunand.

O people of the Bayān and O people of the Qur'an! The light of the Sun
of Revelation has encompassed the world, and its radiance and splendors
are manifest and resplendent. The Point of the Bayān---may the souls of
all others be sacrificed for Him---says: ``A one-year-old embryo in the
day of His appearance is mightier than the whole Bayān.'' This is the
Day that all the Books, Scriptures, and Tablets have ascribed to God,
and of which they have given everyone glad tidings. All must turn toward
the Most Exalted Horizon with keen eyes and ears sanctified from every
mention, and listen to the call raised from the right hand of the Most
Great Throne.

\subsection{85}\label{section-2060}

\begin{otherlanguage}{persian}

امروز هر نفسی به بیان تمسّک نماید و بر اعتراض قیام کند بشهادت نقطهٴ بیان
از اهل حقّ محسوب نه چه که جمیع امور معلّق بتصدیق میفرماید جوهر بیان آنکه
انّه لا یستشار باشارتی و لا بما ذکر فی البیان بعد از اینکلمهٴ مبارکه منصف
خبیر و بصیر علیم به بیان بر حقّ اعتراض نمینماید آیا نشنیده که میفرماید
جمیع بیان ورقی است از اوراق جنّت

\end{otherlanguage}

Imrūz har nafsī bi-Bayān tamassuk namāyad va bar iʿtirāḍ qiyām kunad,
bi-shahādat-i Nuqṭe-yi Bayān az ahl-i Ḥaqq maḥsūb na, chi ke jamīʿ-i
umūr muʿallaq bi-taṣdīq mīfarmāyad. Jawhar-i Bayān ānke: innahu lā
yustashāru bi-ishāratī va lā bi-mā dhukira fī al-Bayān. Baʿd az īn
kalame-yi mubārake munṣif-i khabīr va baṣīr-i ʿalīm bi-Bayān bar Ḥaqq
iʿtirāḍ namīnamāyad. Āyā nashanīde ke mīfarmāyad jamīʿ-i Bayān varaqī
ast az awrāq-i jannat.

Today, anyone who clings to the Bayān and rises in objection is not,
according to the testimony of the Point of the Bayān, counted among the
people of truth, for He makes all things conditional upon recognition.
The essence of the Bayān is this: ``He is not to be consulted through My
allusion, nor through what has been mentioned in the Bayān.'' After this
blessed utterance, no fair-minded and discerning person, no perceptive
and knowledgeable soul, would use the Bayān to object to God. Has he not
heard Him say that the entire Bayān is a leaf among the leaves of the
Paradise

\subsection{86}\label{section-2061}

\begin{otherlanguage}{persian}

او معلوم نیست آنشخص غافل جاهل از آنچه ذکر نموده مقصودش چیست اگر مقصودش
اینست که حال بیان در مقام نطفه است و تکمیل نشده نقطهٴ اولی روح ما سواه
فداه میفرماید قوله جلّ و عزّ من اوّل ذلک الأمر الی قبل ان یکمل تسعة کینونات
الخلق لم تظهر و ان کلّ ما قد رأیت من النّطفه الی ما کسوناه لحماً ثمّ اصبر
حتّی تشهد خلق الآخر قل فتبارک اللّه احسن الخالقین انتهی

\end{otherlanguage}

Ū. Maʿlūm nīst ān shakhs-i ghāfil-i jāhil az ānchi dhikr namūde
maqṣūdash chīst. Agar maqṣūdash īnast ke ḥāl Bayān dar maqām-i nuṭfe ast
va takmīl nashude, Nuqṭe-yi Ūlā, rūḥu mā siwāhu fidāhu, mīfarmāyad,
qawluhu jalla va ʿazza: min awwali dhālika al-amr ilā qabli an yakmula
tisʿatun kaynūnātu al-khalqi lam taẓhar, va inna kulla mā qad raʾayta
mina al-nuṭfati ilā mā kasawnāhu laḥman, thumma iṣbir ḥattā tashhada
khalqa al-ākhar, qul fa-tabāraka Allāhu aḥsanu al-khāliqīn. Intahā.

of His? It is not clear what that heedless, ignorant person means by
what he has said. If he means that the Bayān is now at the stage of the
embryo and has not been completed, the Primal Point---may the souls of
all others be sacrificed for Him---says, exalted and mighty are His
words: ``From the beginning of that Cause until the completion of nine,
the realities of creation have not appeared. All that you have seen is
from the embryo to that which We clothed in flesh. Then wait until you
witness the other creation. Say: `Blessed, then, is God, the best of
creators!'\,'' End of quotation.

\subsection{87}\label{section-2062}

\begin{otherlanguage}{persian}

بگوئید این بی‌انصاف ملاحظه نما و در آنچه از قلم نقطهٴ بیان روح من فی
ملکوت الأمر و الخلق فداه جاری شده تفکّر کن امروز است روز خلق آخر امروز
است که لسان کلّ اشیاء بکلمهٴ مبارکهٴ فتبارک اللّه احسن الخالقین ناطق بظهور
او میفرماید بیان تکمیل میشود و بکمال تصریح میفرماید قبل اکمال تسع
کینونات خلق ظاهر نشده و کلّ آنچه دیده از مقام نطفه است الی اکتساء لحم صبر
کن تا خلق آخر را مشاهده نمائی قل فتبارک اللّه احسن الخالقین اینست مقام و
نفخ فیه اخری الی آخر قوله تعالی

\end{otherlanguage}

Bigūʾīd: īn bī-inṣāf, mulāḥaẓe namā va dar ānchi az qalam-i Nuqṭe-yi
Bayān, rūḥu man fī malakūti al-amri va al-khalqi fidāhu, jārī shude
tafakkur kun. Imrūz ast rūz-i khalq-i ākhar. Imrūz ast ke lisān-i kull-i
ashyā bi-kalame-yi mubārake-yi fa-tabāraka Allāhu aḥsanu al-khāliqīn
nāṭiq. Bi-ẓuhūr-i ū mīfarmāyad Bayān takmīl mīshavad va bi-kamāl-i
taṣrīḥ mīfarmāyad qabl-i ikmāl-i tisʿ kaynūnāt-i khalq ẓāhir nashude va
kull-i ānchi dīde az maqām-i nuṭfe ast ilā iktisāʾ-i laḥm. Ṣabr kun tā
khalq-i ākhar rā mushāhade namāʾī. Qul fa-tabāraka Allāhu aḥsanu
al-khāliqīn. Īnast maqām-i va nufikha fīhi ukhrā ilā ākhir-i qawlihi
taʿālā.

Say: ``You unfair one! Consider and reflect on what has flowed from the
Pen of the Point of the Bayān---may the souls of all in the kingdoms of
command and creation be sacrificed for Him. Today is the day of the
other creation. Today the tongue of all things speaks the blessed words,
`Blessed, then, is God, the best of creators!' He says that through His
appearance the Bayān is completed, and states with complete explicitness
that before the completion of nine the realities of creation have not
appeared, and that all you have seen belongs to the stage of the embryo
up to its being clothed in flesh. Wait until you behold the other
creation. Say: `Blessed, then, is God, the best of creators!' This is
the station of `Then it was sounded again,' to the end of His exalted
words.''

\subsection{88}\label{section-2063}

\begin{otherlanguage}{persian}

عظمت ظهور بمقامیست که نقطهٴ اولی میفرماید قوله جلّ و عزّ حلّ ان یردّ من لم
یکن فوق الأرض اعلی منه اذ ذلک خلق فی قبضته و کلّ له قانتون احدی بر عظمت
اینظهور عارف نه الّا من شآء اللّه این فلک را ملّاحی دیگر باید و این بحر را
غوّاصی دیگر

\end{otherlanguage}

ʿAẓamat-i ẓuhūr bi-maqāmīst ke Nuqṭe-yi Ūlā mīfarmāyad, qawluhu jalla va
ʿazza: ḥalla an yarudda man lam yakun fawqa al-arḍi aʿlā minhu idh
dhālika khalqun fī qabḍatihi va kullun lahu qānitūn. Aḥadī bar ʿaẓamat-i
īn ẓuhūr ʿārif na illā man shāʾa Allāh. Īn fulk rā mallāḥī dīgar bāyad
va īn baḥr rā ghavvāṣī dīgar.

Such is the greatness of this Revelation that the Primal Point says,
exalted and mighty are His words: ``He may reject one above whom there
is none on earth, for that one is a creature in His grasp, and all are
obedient to Him.'' No one understands the greatness of this Revelation
except those whom God wills. This vessel needs another mariner, and this
sea another diver.

\subsection{89}\label{section-2064}

\begin{otherlanguage}{persian}

اگرچه جمیع ایّام مظاهر حقّ و مهابط وحی در مقامی بیوم اللّه مذکور ولکن این
ایّام برداء تخصیص مزیّن چنانچه تفصیل ذکر شده جمیع اقوال مشرکین قبل در
اینظهور از مشرکین جدید استماع شده و میشود بلکه گفته‌اند آنچه را که هیچ
مشرکی نگفته

\end{otherlanguage}

Agarchi jamīʿ-i ayyām-i maẓāhir-i Ḥaqq va mahābiṭ-i vaḥy dar maqāmī
bi-Yawm Allāh madhkūr, va-lākin īn ayyām bi-ridā-yi takhṣīṣ muzayyan,
chinānchi tafṣīl dhikr shude. Jamīʿ-i aqvāl-i mushrikīn-i qabl dar īn
ẓuhūr az mushrikīn-i jadīd istimāʿ shude va mīshavad, balke gufte-and
ānchi rā ke hīch mushrikī nagufte.

Although all the days of the Manifestations of God and the recipients of
revelation are, in one sense, called the Day of God, these days are
adorned with the robe of particular distinction, as has been explained
in detail. All the words of the idolaters of old have been, and continue
to be, heard from the new idolaters in this Revelation. Indeed, they
have said what no idolater has ever said.

\subsection{90}\label{section-2065}

\begin{otherlanguage}{persian}

یا ملأ البیان اتّقوا اللّه لا تعترضوا علی اللّه بأوهامکم انّه یدعوکم الی
النّور و تدعونه الی النّار مالکم لا تفقهون حدیثاً من اللّه العلیم الحکیم علم
احدی باین ظهور احاطه ننموده جمیع علوم عالم نزد بحر علمش بمثابهٴ قطره بل
استغفر اللّه عن هذا التّحدید

\end{otherlanguage}

Yā malaʾ al-Bayān, ittaqū Allāha, lā taʿtariḍū ʿalā Allāhi bi-awhāmikum.
Innahu yadʿūkum ilā al-nūri va tadʿūnahu ilā al-nār. Mā lakum lā
tafqahūna ḥadīthan mina Allāhi al-ʿAlīmi al-Ḥakīm? ʿIlm-i aḥadī bi-īn
ẓuhūr iḥāṭe nanamūde. Jamīʿ-i ʿulūm-i ʿālam nazd-i baḥr-i ʿilmash
bi-mathābe-yi qaṭre, bal astaghfiru Allāha ʿan hādhā al-taḥdīd.

O people of the Bayān! Fear God, and do not object to God on the basis
of your imaginings. He calls you to the light, while you call Him to the
fire. Why do you fail to understand a word from God, the All-Knowing,
the All-Wise? No one's knowledge has encompassed this Revelation. All
the knowledge of the world is like a drop beside the ocean of His
knowledge. Rather, I ask God's forgiveness for this limitation.

\subsection{91}\label{section-2066}

\begin{otherlanguage}{arabic}

مبیّن بیان روح من فی الوجود فداه میفرماید قوله جلّ و عزّ

\end{otherlanguage}

Mubayyin-i Bayān, rūḥu man fī al-wujūdi fidāhu, mīfarmāyad, qawluhu
jalla va ʿazza:

The Expounder of the Bayān---may the souls of all in existence be
sacrificed for Him---says, exalted and mighty are His words:

\subsection{92}\label{section-2067}

\begin{otherlanguage}{arabic}

فان لکم بعد حین امر ستعلمون و انّ یومئذ مثل حین ولکن انتم لا تبصرون تلک
شجرة الطّور لتنطقنّک ان انتم تسمعون قل هو نبأ لن یحیط بعلمه احد ولکن انتم
یومئذ لا تعلمون انتهی

\end{otherlanguage}

Fa-inna lakum baʿda ḥīnin amran sataʿlamūn, va inna yawmaʾidhin mithlu
ḥīnin va-lākin antum lā tubṣirūn. Tilka shajaratu al-Ṭūri
la-tanṭiqannaka in antum tasmaʿūn. Qul huwa nabaʾun lan yuḥīṭa
bi-ʿilmihi aḥadun va-lākin antum yawmaʾidhin lā taʿlamūn. Intahā.

``For after a time there is a Cause for you that you will come to know,
and that day is like a time, but you do not perceive. That Tree of Sinai
will surely speak to you, if you will listen. Say: `It is an
Announcement whose knowledge no one will encompass, yet on that day you
will not know.'\,'' End of quotation.

\subsection{93}\label{section-2068}

\begin{otherlanguage}{persian}

بشهادت بیان و جمیع کتب الهی علم اینظهور نزد احدی نبوده و نیست در ذکر بعد
حین تفکّر نمائید شاید عرف بیان رحمن استشمام شود و گمراهان قصد صراط مستقیم
و نبأ عظیم نمایند و از مفتریات خود را بآب رجوع مطهّر سازند انّ ربّنا الرّحمن
هو التّوّاب الرّحیم

\end{otherlanguage}

Bi-shahādat-i Bayān va jamīʿ-i kutub-i ilāhī, ʿilm-i īn ẓuhūr nazd-i
aḥadī nabūde va nīst. Dar dhikr-i baʿda ḥīn tafakkur namāʾīd, shāyad
ʿarf-i bayān-i Raḥmān istishmām shavad va gumrāhān qaṣd-i ṣirāṭ-i
mustaqīm va nabaʾ-i ʿaẓīm namāyand va az muftariyāt khud rā bi-āb-i
rujūʿ muṭahhar sāzand. Inna Rabbanā al-Raḥmāna huwa al-Tawwābu al-Raḥīm.

According to the testimony of the Bayān and all the divine Books, no one
has possessed, or possesses, knowledge of this Revelation. Reflect on
the mention of ``after a time,'' so that perhaps the fragrance of the
All-Merciful's utterance may be inhaled, and those who have gone astray
may turn toward the straight path and the Great Announcement, and
cleanse themselves of their fabrications with the water of return.
Truly, our Lord, the All-Merciful, is the Ever-Relenting, the Merciful.

\subsection{94}\label{section-2069}

\begin{otherlanguage}{persian}

یا حبیب فؤادی خادم چه عرض نماید و بچه دلیل تمسّک جوید و چه برهان اظهار
دارد لعمر مقصودی کینونة دلیل کلیل است از ذکر این سبیل سبحان‌اللّه از آنچه
گفته و میگویند در قرون و اعصار این همج رعاع بوهمی تمسّک جسته بجدال قیام
نمودند گاهی ذکر قائم اسباب جدال بود و گاهی ذکر ولیّ و وصیّ و نقیب و امثال
آن و در یوم ظهور ثمر گفتگوها و عقاید و اعمال آن بود که دیدند ما هذه
التّماثیل الّتی هم علیها عاکفون کاش بیک کلمه تمسّک میجستند و از روی صدق و
حقیقت اظهار میداشتند افّ لهم و بما اکتسبت ایادیهم

\end{otherlanguage}

Yā ḥabība fuʾādī, khādim chi ʿarḍ namāyad va bi-chi dalīl tamassuk jūyad
va chi burhān iẓhār dārad? La-ʿamri maqṣūdī, kaynūnat-i dalīl kalīl ast
az dhikr-i īn sabīl. Subḥān Allāh az ānchi gufte va mīgūyand! Dar qurūn
va aʿṣār īn hamaj-i raʿāʿ bi-vahmī tamassuk juste bi-jidāl qiyām
namūdand. Gāhī dhikr-i Qāʾim asbāb-i jidāl būd va gāhī dhikr-i valī va
vaṣī va naqīb va amthāl-i ān, va dar yawm-i ẓuhūr thamar-i guftugūhā va
ʿaqāʾid va aʿmāl ān būd ke dīdand. Mā hādhihi al-tamāthīlu allatī hum
ʿalayhā ʿākifūn? Kāsh bi-yek kalame tamassuk mījustand va az rū-yi ṣidq
va ḥaqīqat iẓhār mīdāshtand. Uffin lahum va bi-mā iktasabat ayādīhim.

O beloved of my heart! What can this servant submit, what evidence can
he invoke, and what proof can he present? By the life of my heart's
Desire! The very essence of evidence is powerless to describe this path.
Glory be to God, beyond what they have said and continue to say!
Throughout the centuries and ages, this ignorant rabble has clung to a
delusion and risen in dispute. At times, mention of the Qāʾim was the
cause of contention; at others, mention of the guardian, the successor,
the chief, and the like. On the day of Revelation, the fruit of their
discussions, beliefs, and deeds was what they witnessed. ``What are
these images to which they are so devoted?'' If only they had held fast
to a single word and professed it sincerely and truthfully! Shame on
them and on what their hands have earned!

\subsection{95}\label{section-2070}

\begin{otherlanguage}{arabic}

امروز هیچ فنّی از فنون و هیچ علمی از علوم و هیچ اسمی از اسماء و هیچ شأنی
از شئون و هیچ کتابی از کتب انسانرا غنی نمینماید و کفایت نمیکند الّا بما
نطق به لسان العظمه قوله جلّ جلاله هذا یوم فیه امّ الکتاب ینطق و امّ البیان
ینادی و القوم لا یعرفون و لا یشعرون انّ النّبیّ من آمن بنبئی و الرّسول من
بلّغ رسالاتی و الوصیّ من وصّی النّاس بما نزّل من ملکوتی من بدایع اوامری و
احکامی و الولیّ من امر النّاس بحبّی و دعاهم الی صراطی المستقیم انتهی

\end{otherlanguage}

Imrūz hīch fannī az funūn va hīch ʿilmī az ʿulūm va hīch ismī az asmā va
hīch shaʾnī az shuʾūn va hīch kitābī az kutub insān rā ghanī namīnamāyad
va kifāyat namīkunad illā bi-mā naṭaqa bihi lisānu al-ʿaẓama, qawluhu
jalla jalāluhu: hādhā yawmun fīhi Ummu al-Kitābi yanṭiqu va Ummu
al-Bayāni yunādī va al-qawmu lā yaʿrifūna va lā yashʿurūn. Inna
al-nabiyya man āmana bi-nabaʾī va al-rasūla man ballagha risālātī va
al-waṣiyya man waṣṣā al-nāsa bi-mā nuzzila min malakūtī min badāʾiʿi
awāmirī va aḥkāmī va al-waliyya man amara al-nāsa bi-ḥubbī va daʿāhum
ilā ṣirāṭī al-mustaqīm. Intahā.

Today, no art, no science, no name, no rank, and no book can make a
person independent or suffice for them, except through what the Tongue
of Grandeur has spoken, exalted be His glory: ``This is a Day in which
the Mother Book speaks and the Mother of the Bayān calls aloud, while
the people neither recognize nor perceive. The prophet is one who
believes in My Announcement; the messenger is one who conveys My
messages; the successor is one who enjoins upon the people the wondrous
commandments and laws sent down from My Kingdom; and the guardian is one
who bids the people love Me and calls them to My straight path.'' End of
quotation.

\subsection{96}\label{section-2071}

\begin{otherlanguage}{arabic}

طوبی از برای نفوسیکه شبهات و اشارات قوم ایشانرا از منزل آیات و مظهر
بیّنات منع ننماید و محروم نسازد جمیع اسامی بقولش خلق میشود نقطهٴ اولی روح
ما سواه فداه میفرماید قوله جلّ و عزّ

\end{otherlanguage}

Ṭūbā az barā-yi nufūsī ke shubuhāt va ishārāt-i qawm īshān rā az
munzil-i āyāt va maẓhar-i bayyināt manʿ nanamāyad va maḥrūm nasāzad.
Jamīʿ-i asāmī bi-qawlash khalq mīshavad. Nuqṭe-yi Ūlā, rūḥu mā siwāhu
fidāhu, mīfarmāyad, qawluhu jalla va ʿazza:

Blessed are those souls whom the people's doubts and allusions do not
keep from, or deprive of, the Revealer of the verses and the
Manifestation of clear proofs. All names are created by His word. The
Primal Point---may the souls of all others be sacrificed for Him---says,
exalted and mighty are His words:

\subsection{97}\label{section-2072}

\begin{otherlanguage}{arabic}

یخلق فی قول ما یشآء من نبیّ او ولیّ او صدّیق او نقیّ اذ کلّ ما قد خلق کلّ
ادلّآء من عنده و سفرآء من لدنه الی قوله عزّ و جلّ ایّاک ایّاک ایّام ظهوره ان
تحتجب بالواحد البیانیّه فانّ ذلک الواحد خلق عنده و ایّاک ایّاک ان تحتجب
بکلمات ما نزّلت فی البیان الی قوله تعالی و لا تنظر الیه الّا بعینه فانّ من
ینظر الیه بعینه یدرکه والّا یحتجب ان اردت اللّه و لقائه فارده و انظر الیه
انتهی

\end{otherlanguage}

Yakhluqu fī qawlin mā yashāʾu min nabiyyin aw waliyyin aw ṣiddīqin aw
naqiyyin, idh kullu mā qad khalaqa kullun adillāʾu min ʿindihi va
sufarāʾu min ladunhi. Ilā qawlihi ʿazza va jalla: iyyāka iyyāka ayyāma
ẓuhūrihi an taḥtajiba bi-al-wāḥidi al-bayāniyya, fa-inna dhālika
al-wāḥida khalqun ʿindahu, va iyyāka iyyāka an taḥtajiba bi-kalimāti mā
nazzaltu fī al-Bayān. Ilā qawlihi taʿālā: va lā tanẓur ilayhi illā
bi-ʿaynihi, fa-inna man yanẓuru ilayhi bi-ʿaynihi yudrikuhu va-illā
yaḥtajib. In aradta Allāha va liqāʾahu fa-aridhu va unẓur ilayhi.
Intahā.

``Through a word He creates whomever He wills---a prophet, a guardian, a
supremely truthful one, or a pure one---for all whom He has created are
guides from Him and emissaries from His presence.'' And further, mighty
and exalted are His words: ``Beware, beware lest, in the days of His
appearance, you be veiled by the Vāḥid of the Bayān, for that Vāḥid is a
creature in His presence. And beware, beware lest you be veiled by the
words I have revealed in the Bayān.'' And further, exalted are His
words: ``Look upon Him only through His own eye, for whoever looks upon
Him through His own eye will recognize Him; otherwise he will be veiled.
If you desire God and His presence, desire Him and look upon Him.'' End
of quotation.

\subsection{98}\label{section-2073}

\begin{otherlanguage}{persian}

ای صاحبان آذان اقلّ من آن بانصاف نطق نمائید بشهادت نقطهٴ بیان روح من فی
عوالم الأمر و الخلق فداه امروز اگر نفسی به بیان استدلال نماید و یا باو
تمسّک جوید نفس بیان و ما نزّل فیه من القلم الأعلی او او تبرّی میجوید از قبل
ذکر شد میفرماید انّه لا یستشار باشارتی و لا بما ذکر فی البیان امروز بیان
معلّق بقبول است طوبی از برای نفوسیکه فی‌الجمله تفکّر نمایند و در ثمرات
اعمال اهل فرقان نظر کنند تا از جمیع آنگفتها و استدلالها و اعتراضها فارغ
و آزاد شوند و با قوادم لطیف مقدّس در این هوا طیران نمایند آیا از مفتریات
مطالع نفس و هوی سیر نشده‌اند که تازه مجدّد میخواهند گفتگوئی برپا نمایند و
احادیثی ذکر کنند و ردّ و قبولی احداث نمایند بگوئید بعضد یقین حجبات مانعه
را شقّ نمائید

\end{otherlanguage}

Ay ṣāḥibān-i ādhān, aqalla min ān bi-inṣāf nuṭq namāʾīd. Bi-shahādat-i
Nuqṭe-yi Bayān, rūḥu man fī ʿawālimi al-amri va al-khalqi fidāhu, imrūz
agar nafsī bi-Bayān istidlāl namāyad va yā bi-ū tamassuk jūyad, nafs-i
Bayān va mā nuzzila fīhi mina al-Qalami al-Aʿlā {[}corrupted text{]}
tabarrī mījūyad. Az qabl dhikr shud, mīfarmāyad: innahu lā yustashāru
bi-ishāratī va lā bi-mā dhukira fī al-Bayān. Imrūz Bayān muʿallaq
bi-qabūl ast. Ṭūbā az barā-yi nufūsī ke fī al-jumle tafakkur namāyand va
dar thamarāt-i aʿmāl-i ahl-i Furqān naẓar kunand tā az jamīʿ-i ān
guftehā va istidlālhā va iʿtirāḍhā fārigh va āzād shavand va bā
qavādim-i laṭīf-i muqaddas dar īn havā ṭayarān namāyand. Āyā az
muftariyāt-i maṭāliʿ-i nafs va havā sīr nashude-and ke tāze mujaddad
mīkhvāhand guftugūʾī barpā namāyand va aḥādīthī dhikr kunand va radd va
qabūlī iḥdāth namāyand? Bigūʾīd: bi-ʿaḍud-i yaqīn ḥujubāt-i māniʿe rā
shaqq namāʾīd.

O you who have ears! Speak fairly for less than a moment. According to
the testimony of the Point of the Bayān---may the souls of all in the
worlds of command and creation be sacrificed for Him---if anyone today
argues from the Bayān or clings to it, the Bayān itself and what has
been revealed in it by the Most Exalted Pen {[}corrupted original
text{]} disavow. As mentioned before, He says: ``He is not to be
consulted through My allusion, nor through what has been mentioned in
the Bayān.'' Today the Bayān is conditional upon acceptance. Blessed are
those souls who reflect a little and consider the fruits of the deeds of
the people of the Qur'an, so that they may be released and freed from
all those sayings, arguments, and objections, and soar in this
atmosphere on delicate, sanctified wings. Have they not had enough of
the fabrications of the sources of self and desire, that they now wish
to set up a fresh controversy, cite traditions, and introduce rejection
and acceptance? Say: ``With the arm of certitude, tear through the
obstructing veils.''

\subsection{99}\label{section-2074}

\begin{otherlanguage}{persian}

لعمر ربّنا امروز مکلّم طور ظاهر و ناطق و سرّ مکنون و غیب مخزون باهر چندی
بعضی از اهل بیان بذکر الوهیّت اعتراض نمودند که چرا حقّ جلّ جلاله باینکلمه
ناطق است مع آنکه در جمیع کتب میفرماید حضرت قیّوم در آن یوم بکلمهٴ انّنی
انا اللّه ناطق انظروا ثمّ اذکروا ما انزله الرّحمن فی الفرقان قوله تبارک و
تعالی یوم یأتی ربّک او بعض آیات ربّک و یوم یقوم النّاس لربّ العالمین و جئی
ربّک و الملک صفّاً صفّاً و امثال آن خاتم انبیا روح ما سواه فداه میفرماید
سترون ربّکم کما ترون البدر فی لیلة اربعة عشر

\end{otherlanguage}

La-ʿamri Rabbinā, imrūz Mukallim-i Ṭūr ẓāhir va nāṭiq va sirr-i maknūn
va ghayb-i makhzūn bāhir. Chandī baʿḍī az ahl-i Bayān bi-dhikr-i
ulūhiyyat iʿtirāḍ namūdand ke chirā Ḥaqq jalla jalāluhu bi-īn kalame
nāṭiq ast, maʿa ānke dar jamīʿ-i kutub mīfarmāyad Ḥaḍrat-i Qayyūm dar ān
yawm bi-kalame-yi innanī anā Allāh nāṭiq. Unẓurū thumma udhkurū mā
anzalahu al-Raḥmānu fī al-Furqān, qawluhu tabāraka va taʿālā: yawma
yaʾtī Rabbuka aw baʿḍu āyāti Rabbika, va yawma yaqūmu al-nāsu li-Rabbi
al-ʿālamīn, va jīʾa Rabbuka va al-malaku ṣaffan ṣaffan, va amthāl-i ān.
Khātam-i anbiyā, rūḥu mā siwāhu fidāhu, mīfarmāyad: satarawna Rabbakum
kamā tarawna al-badra fī laylati arbaʿata ʿashar.

By the life of our Lord! Today He who spoke on Sinai is manifest and
speaking, and the hidden secret and treasured mystery are revealed. For
a time, some of the people of the Bayān objected to the mention of
divinity, asking why God, exalted be His glory, speaks these words, even
though throughout the Books He says that on that Day the Self-Subsisting
One will utter the words, ``Truly, I am God.'' Consider, then recall
what the All-Merciful has sent down in the Qur'an, blessed and exalted
are His words: ``The day your Lord comes, or some of your Lord's
signs''; and ``The day when the people stand before the Lord of the
worlds''; and ``Your Lord is brought, and the angels, rank upon rank'';
and the like. The Seal of the Prophets---may the souls of all others be
sacrificed for Him---says: ``You will see your Lord as you see the full
moon on the fourteenth night.''

\subsection{100}\label{section-2075}

\begin{otherlanguage}{arabic}

اشعیاء نبی میفرماید یسمو الرّبّ وحده فی ذلک الیوم در عظمت یوم میفرماید
ادخل الی الصّخره و اختبئی فی التّراب من امام هیبة الرّبّ و من بهآء عظمته و
در مقام دیگر تفرح البرّیّة و الأرض الیابسة و یبتهج القفر و یزهر کالنّرجس
یزهر ازهاراً و یبتهج ابتهاجاً یدفع الیه مجد لبنان بهآء کرمل و شارون هم
یرون مجد الرّبّ بهآء الهنا

\end{otherlanguage}

Ishʿiyā Nabī mīfarmāyad: yasmū al-Rabbu waḥdahu fī dhālika al-yawm. Dar
ʿaẓamat-i yawm mīfarmāyad: udkhul ilā al-ṣakhrati va ikhtabiʾī fī
al-turābi min amāmi haybati al-Rabbi va min bahāʾi ʿaẓamatihi. Va dar
maqām-i dīgar: tafraḥu al-barriyyatu va al-arḍu al-yābisatu va yabtahiju
al-qafru va yuzhiru ka-al-narjis, yuzhiru izhāran va yabtahiju
ibtihājan, yudfaʿu ilayhi majdu Lubnān, bahāʾu Karmil va Shārūn, hum
yarawna majda al-Rabbi bahāʾa ilāhinā.

The prophet Isaiah says: ``The Lord alone will be exalted on that Day.''
Concerning the greatness of the Day, he says: ``Enter into the rock and
hide in the dust before the terror of the Lord and the splendor of His
majesty.'' And elsewhere: ``The wilderness and the dry land will
rejoice, and the desert will exult and blossom like the narcissus. It
will blossom abundantly and rejoice with exultation. The glory of
Lebanon will be given to it, the splendor of Carmel and Sharon. They
will see the glory of the Lord, the splendor of our God.''

\subsection{101}\label{section-2076}

\begin{otherlanguage}{persian}

امروز صاحبان ابصار جمیع ذرّات کائنات را با فرح و سرور مشاهده مینمایند چه
که بآنچه از اوّل لا اوّل بشارت داده‌اند ظاهر شده لعمر محبوبنا و مقصودنا اگر
در قلبی اقلّ از سمّ ابره محبّة اللّه باشد و یا نور انصاف باین بیان مشتعل شود
اشتعالیکه عالم را بحرارت محبّة اللّه فائز نماید

\end{otherlanguage}

Emrūz ṣāḥebān-e abṣār jamīʿ-e dharrāt-e kāʾenāt rā bā faraḥ va sorūr
moshāhede mīnamāyand, che ke be-ānche az avval-e lā-avval beshārat
dāde-and ẓāher shode. La-ʿamri maḥbūbinā wa-maqṣūdinā, agar dar qalbī
aqall az samm-e ebre maḥabbat-e Allāh bāshad va yā nūr-e enṣāf, be-īn
bayān moshtaʿel shavad, eshteʿālī ke ʿālam rā be-ḥarārat-e maḥabbat-e
Allāh fāʾez namāyad.

Today those with eyes to see behold all the atoms of creation filled
with joy and gladness, for what has been foretold from the beginning
that has no beginning has appeared. By the life of our Beloved and our
heart's Desire! If a heart contains even less than a needle's eye of the
love of God, or the light of fairness, it will be set ablaze by this
utterance with a fire that will enable the world to partake of the
warmth of the love of God.

\subsection{102}\label{section-2077}

\begin{otherlanguage}{persian}

ای اهل بها این سست عنصرهای عالم را بگذارید بحال خودشان و با انفس مشتعله
و ارجل و ایادی قویّه و ابصار حدیده و آذان واعیه و قلوب زکیّه و السن طلقه
بگوئید و ببینید و بشنوید رغماً للّذین یدّعون الایمان بالبیان و یکفرون
بمُنزله و مُظهره و مُرسله این برّیّه که در آیهٴ مبارکه فرموده برّیّهٴ فیحاست و
اینست ارض مقدّسه که در کتب قبل و فرقان مذکور است میفرماید مسرور میشود
برّیّه و اراضی یابسه و بیداء قفر و این سرور و فرح و بهجت از چه جهت ظاهر
لعمر ربّنا از قدوم ربّ جلّ جلاله کرمل از قبل ذکر شد که در کتب مقدّسه بکَوم
اللّه معروف و امام وجه عکّا واقع و از جهت دیگر لبنان و همچنین شارون باری
جمیع وعدهای الهی که در کتب منزله از قبل و بعد مذکور است در این اراضی
بوده

\end{otherlanguage}

Ey ahl-e Bahā, īn sost-ʿonṣorhā-ye ʿālam rā begoẕārīd be-ḥāl-e khodeshān
va bā anfos-e moshtaʿele va arjol va ayādī-ye qaviyye va abṣār-e ḥadīde
va ādhān-e vāʿiye va qolūb-e zakiyye va alson-e ṭaleqe begūʾīd va
bebīnīd va beshnavīd, raghman li-lladhīna yaddaʿūna al-īmāna bi-l-Bayāni
wa-yakfurūna bi-munzilihi wa-muẓhirihi wa-mursilihi. Īn barriyye ke dar
āye-ye mobārake farmūde barriyye-ye fayḥāst va īnast arż-e moqaddase ke
dar kotob-e qabl va Forqān madhkūr ast. Mīfarmāyad: masrūr mīshavad
barriyye va arāżī-ye yābese va baydā-ye qafr. Va īn sorūr va faraḥ va
bahjat az che jahat ẓāher? La-ʿamri rabbinā, az qodūm-e Rabb, jalla
jalāluhu. Karmel az qabl dhekr shod ke dar kotob-e moqaddase be-Kawm-e
Allāh maʿrūf va amām-e vajh-e ʿAkkā vāqeʿ va az jahat-e dīgar Lobnān va
hamchenīn Shārūn. Bārī jamīʿ-e vaʿdehā-ye elāhī ke dar kotob-e monzale
az qabl va baʿd madhkūr ast dar īn arāżī būde.

People of Bahá! Leave these weak-willed people of the world to
themselves, and with souls ablaze, strong feet and hands, keen eyes,
attentive ears, pure hearts, and eloquent tongues, speak, see, and hear,
despite those who profess faith in the Bayān yet disbelieve in the One
who sent it down, revealed it, and dispatched it. This wilderness
mentioned in the blessed verse is the spacious wilderness, and this is
the Holy Land mentioned in the earlier Books and the Qur'an. He says:
``The wilderness, the dry lands, and the barren desert will rejoice.''
Why do this gladness, joy, and delight appear? By the life of our Lord!
Because of the coming of the Lord, exalted be His glory. Carmel was
mentioned earlier; it is known in the Holy Books as the Mound of God and
lies facing Acre, while in another direction are Lebanon and likewise
Sharon. In short, all the divine promises mentioned in the revealed
Books, earlier and later, concern these lands.

\subsection{103}\label{section-2078}

\begin{otherlanguage}{persian}

و همچنین اشعیاء نبی میفرماید شدّدوا الأیادی المسترخیه و الرّکب المرتعشة
ثبّتوها میفرماید ایادی مسترخیه را محکم نمائید و رکب مرتعشه را ثابت و
مستقیم کنید قولوا لخائفی القلوب تشدّدوا لا تخافوا هو ذا الهکم

\end{otherlanguage}

Va hamchenīn Eshʿeyā-ye nabī mīfarmāyad: shaddidū al-ayādiya
al-mustarkhiyata wa-l-rukaba al-murtaʿishata thabbitūhā. Mīfarmāyad:
ayādī-ye mostarkhiye rā moḥkam namāyīd va rokab-e mortaʿeshe rā thābet
va mostaqīm konīd. Qūlū li-khāʾifī al-qulūbi tashaddadū lā takhāfū huwa
dhā ilāhukum.

Likewise, the prophet Isaiah says: ``Strengthen the feeble hands and
make firm the trembling knees.'' He says: ``Make the feeble hands
strong, and the trembling knees firm and steadfast. Say to the fearful
of heart: `Be strong; do not fear. Here is your God.'\,''

\subsection{104}\label{section-2079}

\begin{otherlanguage}{persian}

جوهر این بیان آنکه در یوم ظهور از کلمه مبارکه انّنی انا اللّه محجوب نمانید
و سست نشوید بایادی قویه کتاب را اخذ نمائید و بارجل ثابته مستقیمه بر خدمت
امر قیام کنید ای دوستان خود را بما عند النّاس و بما عندکم از ما عند اللّه
محجوب مسازید امروز روز اقبال و تصدیق و بلی و لبیک است میفرماید قوله جل
جلاله لیس هذا یوم السئوال ینبغی لکل نفس اذا سمع النداء من الافق الاعلی
یقول لبیک یا مالک الاسماء و لبیک یا مقصود من فی السموات و الارضین ''
انتهی

\end{otherlanguage}

Jowhar-e īn bayān ānke dar yawm-e ẓohūr az kaleme-ye mobārake-ye innānī
anā Allāh maḥjūb namānīd va sost nashavīd. Be-ayādī-ye qaviyye ketāb rā
akhdh namāyīd va be-arjol-e thābete-ye mostaqīme bar khedmat-e amr qiyām
konīd. Ey dūstān, khod rā be-mā ʿinda al-nās va be-mā ʿindakum az mā
ʿinda Allāh maḥjūb masāzīd. Emrūz rūz-e eqbāl va taṣdīq va balī va
labbayk ast. Mīfarmāyad, qawluhu jalla jalāluhu: laysa hādhā yawma
al-suʾāl yanbaghī li-kulli nafsin idhā samiʿa al-nidāʾa mina al-ufuqi
al-aʿlā yaqūlu labbayka yā Mālika al-asmāʾ wa-labbayka yā maqṣūda man fī
al-samāwāti wa-l-araḍīn. Intahā.

The essence of this utterance is that, on the Day of Revelation, you
should not remain veiled from the blessed words ``Truly, I am God,'' nor
grow faint. Take hold of the Book with strong hands, and arise to serve
the Cause with firm and steadfast feet. Friends! Do not let what people
possess, or what you yourselves possess, veil you from what is with God.
Today is the day of turning toward Him, of affirmation, of ``Yes'' and
``Here I am.'' He says---exalted be His glory: ``This is not the day for
questioning. It befits every soul, upon hearing the call from the
highest horizon, to say: `Here I am, Lord of Names! Here I am, Desire of
all who are in the heavens and the earths!'\,'' End of quotation.

\subsection{105}\label{section-2080}

\begin{otherlanguage}{persian}

این خادم نمیداند چه عرض نماید و بچه مشغول شود حیا مینماید از این اذکار و
سبب تأخیر این اوراق آنکه یومی تلقاء وجه عرض شد اصبع مبارک را حرکت دادند
و فرمودند انّ ربّک لا یحتاج بدلیل و لا بحجة و لا ببرهان دون نفسه یا عبد
حاضر هر نفسی اقل از سم ابره عرف قمیص را یافته جمیع اشیاء را دلیل و حجت و
برهان مشاهده مینماید طوبی از برای نفسی که بآب محبت الهی خود را از اقوال
و افعال و اشارات و شبهات معرضین پاک و مقدس نمود و بافق اعلی وحده متوجه
شد اوست از اهل بها و الذی منع انه فی ضلال مبین انتهی.

\end{otherlanguage}

Īn khādem nemīdānad che ʿarż namāyad va be-che mashghūl shavad. Ḥayā
mīnamāyad az īn adhkār va sabab-e taʾkhīr-e īn owrāq ānke yawmī telqāʾ-e
vajh ʿarż shod. Oṣboʿ-e mobārak rā ḥarekat dādand va farmūdand: inna
rabbaka lā yaḥtāju bi-dalīlin wa-lā bi-ḥujjatin wa-lā bi-burhānin dūna
nafsihi. Yā ʿAbd-e Ḥāżer, har nafsī aqall az samm-e ebre ʿarf-e qamīṣ rā
yāfte, jamīʿ-e ashyā rā dalīl va ḥojjat va borhān moshāhede mīnamāyad.
Ṭūbā az barā-ye nafsī ke be-āb-e maḥabbat-e elāhī khod rā az aqvāl va
afʿāl va eshārāt va shobohāt-e moʿreżīn pāk va moqaddas namūd va be-of
oq-e aʿlā vaḥdahu motavajjeh shod. Ūst az ahl-e Bahā, wa-lladhī muniʿa
innahu fī ḍalālin mubīn. Intahā.

This servant does not know what to say or with what to occupy himself.
He feels ashamed to mention these things. The reason for the delay of
these pages is that one day the matter was submitted in His presence. He
moved His blessed finger and said: ``Your Lord needs no evidence,
argument, or proof other than Himself. O servant in attendance! Any soul
who has caught even less than a needle's eye of the fragrance of the
robe sees all things as evidence, argument, and proof. Blessed is the
soul who, with the water of divine love, has cleansed and sanctified
himself from the words, deeds, insinuations, and doubts of those who
turn away, and has turned toward the highest horizon alone. He is of the
people of Bahá, while whoever has been kept back is indeed in manifest
error.'' End of quotation.

\subsection{106}\label{section-2081}

\begin{otherlanguage}{persian}

و در مقامی این کلمه علیا از مالک اسماء استماع شد فرمودند یا عبد حاضر قلم
اعلی باعلی النداء میفرماید یا ملأ الفرقان و ملأ البیان و یا اهل الادیان
امروز یوم اللّه است و حق بنفسه ظاهر و در اثبات امرش بغیر محتاج نبوده و
نیست اگر آیات را حجت میدانند معادل جمیع کتب الهی از قبل و بعد بل ازید از
سماء مشیتش نازل و اگر بینات میجویند در سور منزله ملاحظه نمایند تا آگاه
شوند و از بحر علم الهی بیاشامند یا عبد حاضر امروز تا از خلیج اسماء عبور
ننمایند ببحر اعظم وارد نشوند و فائز نگردند هذا حق لاریب فیه '' انتهی.

\end{otherlanguage}

Va dar maqāmī īn kaleme-ye ʿolyā az Mālek-e Asmāʾ estemāʿ shod.
Farmūdand: yā ʿAbd-e Ḥāżer, Qalam-e Aʿlā be-aʿlā al-nidāʾ mīfarmāyad: yā
malaʾa al-Furqān wa-malaʾa al-Bayān wa-yā ahla al-adyān, emrūz Yawm-e
Allāh ast va Ḥaqq be-nafsehe ẓāher va dar ethbāt-e amrash be-ghayr
moḥtāj nabūde va nīst. Agar āyāt rā ḥojjat mīdānand, moʿādel-e jamīʿ-e
kotob-e elāhī az qabl va baʿd, bal azyad, az samā-ye mashiyyatash nāzel,
va agar bayyenāt mījūyand dar sovar-e monzale molāḥeẓe namāyand tā āgāh
shavand va az baḥr-e ʿelm-e elāhī biyāshāmand. Yā ʿAbd-e Ḥāżer, emrūz tā
az khalīj-e asmāʾ ʿobūr nanamāyand be-baḥr-e aʿẓam vāred nashavand va
fāʾez nagardand. Hādhā ḥaqqun lā rayba fīhi. Intahā.

On another occasion, this exalted utterance was heard from the Lord of
Names. He said: ``O servant in attendance! The Most Exalted Pen
proclaims at the loudest call: `People of the Qur'an, people of the
Bayān, and followers of all religions! Today is the Day of God, and the
True One is manifest in His own Person. He has never needed, nor does He
now need, anyone else to establish His Cause. If they consider verses a
proof, the equivalent of all the divine Books, earlier and
later---indeed, more---has descended from the heaven of His will. If
they seek clear evidence, let them examine the revealed surahs, so that
they may become aware and drink from the ocean of divine knowledge.' O
servant in attendance! Today, until they cross the gulf of names, they
will neither enter the Most Great Ocean nor attain its bounty. This is
the truth, beyond all doubt.'' End of quotation.

\subsection{107}\label{section-2082}

\begin{otherlanguage}{persian}

این خادم فانی خدمت کل عرض مینماید ایام را باعتراضات ناشایسته و اقوال
نالایقه و اعمال غیر مرضیه صرف ننمائید امر بسیار عزیز و عند اللّه از کلشئ
اعلی و اغلی بسیار حیف است چنین شئ ثمینی بمشتهیات نفسانیه و اوهامات لا
یسمنه لا یغنیه صرف شود سبحان اللّه شخص کرمانی علومیکه ذکر نموده محض
افتخار بوده غافل از آنکه از هر یک از آنچه نوشته باید سالها قلب را از آن
طاهر نماید و بحق رجوع کند

\end{otherlanguage}

Īn khādem-e fānī khedmat-e koll ʿarż mīnamāyad: ayyām rā be-eʿterāżāt-e
nāshāyeste va aqvāl-e nālāyeqe va aʿmāl-e ghayr-e marżiyye ṣarf
nanamāyīd. Amr besyār ʿazīz va ʿinda Allāh az koll-e shayʾ aʿlā va
aghlā. Besyār ḥayf ast chonīn shayʾ-e thamīnī be-moshtahayāt-e
nafsāniyye va owhāmāt-e lā yusminuhu lā yughnīhi ṣarf shavad. Sobḥān-e
Allāh, shakhṣ-e Kermānī ʿolūmī ke dhekr namūde maḥż-e eftekhār būde,
ghāfel az ānke az har yek az ānche neveshte bāyad sālhā qalb rā az ān
ṭāher namāyad va be-Ḥaqq rojūʿ konad.

This transient servant submits to everyone: Do not spend your days on
unseemly objections, unworthy words, and unacceptable deeds. The Cause
is exceedingly precious and, in God's sight, higher and more valuable
than all things. It is a great pity to expend something so precious on
selfish desires and vain imaginings that neither nourish nor enrich one.
Glory be to God! The man from Kirmán has mentioned these branches of
learning merely to boast, unaware that he must spend years purifying his
heart of each of the things he has written and turn back to the True
One.

\subsection{108}\label{section-2083}

\begin{otherlanguage}{arabic}

ای عالم اگر در تو گوشی هست بشنو نقطه بیان روح ما سواه فداه میفرماید اگر
یک آیه از آیات من یظهره اللّه تلاوت کنی عزیزتر خواهد بود عند اللّه از آنکه
کل بیان را ثبت کنی زیرا که در آن روز آن یک آیه تو را نجات میدهد ولی کل
بیان نمیدهد

\end{otherlanguage}

Ey ʿālam, agar dar to gūshī hast beshnow. Noqṭe-ye Bayān, rūḥu mā siwāhu
fidāhu, mīfarmāyad: agar yek āye az āyāt-e Man Yuẓhiruhu Allāh telāvat
konī ʿazīztar khāhad būd ʿinda Allāh az ānke koll-e Bayān rā thabt konī,
zīrā ke dar ān rūz ān yek āye to rā nejāt mīdahad valī koll-e Bayān
nemīdahad.

O world! If you have an ear, listen. The Point of the Bayān---may the
soul of all else be sacrificed for Him---says: ``If you recite a single
verse of the verses of Him whom God will make manifest, it will be more
precious in God's sight than if you transcribe the whole Bayān, for on
that Day that one verse will save you, whereas the whole Bayān will
not.''

\subsection{109}\label{section-2084}

\begin{otherlanguage}{arabic}

و اگر باز در تو گوشی یافت میشود این کلمه نقطه روح ما سواه فداه را اصغاء
نما قوله عز و جل انّی مؤمن به و بدینه و بکتابه و بادلائه و بمناهجه و بما
یظهر من عنده فی کل ذلک مفتخرا بنسبتی الیه و متعززا بایمانی به

\end{otherlanguage}

Va agar bāz dar to gūshī yāft mīshavad īn kaleme-ye Noqṭe, rūḥu mā
siwāhu fidāhu, rā eṣghā namā. Qawluhu ʿazza wa-jalla: innī muʾminun bihi
wa-bi-dīnihi wa-bi-kitābihi wa-bi-adillāʾihi wa-bi-manāhijihi wa-bi-mā
yaẓharu min ʿindihi fī kulli dhālika muftakhiran bi-nisbatī ilayhi
wa-mutaʿazzizan bi-īmānī bihi.

And if an ear can still be found in you, listen to these words of the
Point---may the soul of all else be sacrificed for Him. He says, mighty
and majestic is He: ``I believe in Him, in His religion, His Book, His
guides, His paths, and whatever appears from Him, in all this taking
pride in my relationship to Him and finding honor in my faith in Him.''

\subsection{110}\label{section-2085}

\begin{otherlanguage}{arabic}

و باز اگر در تو گوشی یافت میشود لله بشنو و بحقیقت بشنو و بانصاف بشنو و
باذن واعیه بشنو قوله عز و جل ان یا کلشئ فی البیان فلتعرفن حد انفسکم فان
مثل نقطة البیان یومن بمن یظهره اللّه قبل کل شئ و انّنی انا بذلک افتخرن علی
من فی ملکوت السموات و الارض و قوله عز و جل لا تحتجبن عن اللّه بعد ظهوره
فان کل ما رفع البیان کخاتمة فی یدی و انّنی انا خاتم فی یدی من یظهره اللّه
و قوله جل ذکره انک انت و اعداد واحد الاول و کل ما یحصی اللّه جل جلاله ما
جعل اللّه الا کخاتم فی یدی من یظهره اللّه جل ذکره یقلب کیف یشاء انه لهو
المهیمن المتعال و قوله جل شأنه فانه یجعل ما علی الارض نبیا لیکونن انبیاء
عند اللّه و قوله تبارک و تعالی و اذا یوم ظهور من یظهره اللّه کل من علی
الارض عنده سواء فمن یجعله نبیا کان نبیا من اول الذی لا اول له الی آخر
الذی لا آخر له لان ذلک مما قد جعله اللّه و من یجعله ولیا فذلک ما کان ولیا
فی کل العوالم فان ذلک مما قد جعله

\end{otherlanguage}

Va bāz agar dar to gūshī yāft mīshavad li-llāh beshnow va be-ḥaqīqat
beshnow va be-enṣāf beshnow va be-odhon-e vāʿiye beshnow. Qawluhu ʿazza
wa-jalla: an yā kulla shayʾin fī al-Bayān fa-la-taʿrifunna ḥadda
anfusikum fa-inna mithla Nuqṭati al-Bayān yuʾminu bi-Man Yuẓhiruhu Allāh
qabla kulli shayʾ wa-innanī anā bi-dhālika aftakhiranna ʿalā man fī
malakūti al-samāwāti wa-l-arḍ. Wa-qawluhu ʿazza wa-jalla: lā
taḥtajibanna ʿani Allāhi baʿda ẓuhūrihi fa-inna kulla mā rafaʿa al-Bayān
ka-khātamatin fī yadayya wa-innanī anā khātamun fī yaday Man Yuẓhiruhu
Allāh. Wa-qawluhu jalla dhikruhu: innaka anta wa-aʿdādu Wāḥidi al-Awwal
wa-kullu mā yuḥṣī Allāhu jalla jalāluhu mā jaʿala Allāhu illā
ka-khātamin fī yaday Man Yuẓhiruhu Allāh jalla dhikruhu yuqallibu kayfa
yashāʾu innahu la-huwa al-Muhayminu al-Mutaʿāl. Wa-qawluhu jalla
shaʾnuhu: fa-innahu yajʿalu mā ʿalā al-arḍi nabiyyan li-yakūnunna
anbiyāʾa ʿinda Allāh. Wa-qawluhu tabāraka wa-taʿālā: wa-idhā yawma
ẓuhūri Man Yuẓhiruhu Allāh kullu man ʿalā al-arḍi ʿindahu sawāʾ fa-man
yajʿalhu nabiyyan kāna nabiyyan min awwali alladhī lā awwala lahu ilā
ākhiri alladhī lā ākhira lahu li-anna dhālika mimmā qad jaʿalahu Allāh
wa-man yajʿalhu waliyyan fa-dhālika mā kāna waliyyan fī kulli al-ʿawālim
fa-inna dhālika mimmā qad jaʿalahu.

And again, if an ear can be found in you, listen for God's sake; listen
truly, listen fairly, and listen with an attentive ear. He says, mighty
and majestic is He: ``O all things in the Bayān! Know your own limits,
for one such as the Point of the Bayān believes in Him whom God will
make manifest before all things, and in this I indeed take pride over
all who are in the kingdom of the heavens and the earth.'' And He says,
mighty and majestic is He: ``Do not veil yourself from God after His
appearance, for all that the Bayān has exalted is like a seal-ring in my
hands, and I myself am a seal-ring in the hands of Him whom God will
make manifest.'' And He says, exalted be His mention: ``You, the numbers
of the First Váḥid, and all that God, exalted be His glory,
enumerates---God has made you only as a seal-ring in the hands of Him
whom God will make manifest, exalted be His mention. He turns it however
He wills. He is indeed the All-Protecting, the Most Exalted.'' And He
says, exalted be His station: ``He makes those on earth prophets, so
that they will indeed be prophets in God's sight.'' And He says, blessed
and exalted is He: ``On the day of the appearance of Him whom God will
make manifest, all who are on earth are equal before Him. Whoever He
makes a prophet has been a prophet from the beginning that has no
beginning to the end that has no end, for this is something God has
ordained. And whoever He makes a saint has been a saint in all the
worlds, for this is something He has ordained.''

\subsection{111}\label{section-2086}

\begin{otherlanguage}{persian}

سبحان اللّه از تدبیرات مدبر حقیقی و از غفلت ما عباد ما من نصح الا و قد
ذکره و ما من تدبیر الا و قد انزله بیانی فرموده‌اند که هر صاحب بصر و هر
صاحب فوأدی ناله مینماید و نوحه میکند چه که از کلماتش عالم عالم حزن رو
مینماید و از بیانش عالم عالم هم و غم قوله جل جلاله فان مثله جل ذکره کمثل
الشمس لو یقابلنه الی ما لا نهایه مرایاء کلهن لیستعکسن عن تجلی الشمس فی
حدهم و ان لم یقابلها من احد فیطلع الشمس و یغرب الحجاب للمرایا و انی ما
قصرت عن نصحی ذلک الخلق و تدبیری لاقبالهم الی اللّه علی الارض فاذا یسر
کینونتی حیث کل قد بلغوا الی ذروة وجودهم و وصلوا الی طلعة محبوبهم و
ادرکوا ما یمکن فی الامکان من تجلی مقصودهم و الا یحزن فوأدی بانی قد ربیت
کلشئ لذلک فکیف تحتجب احد علی هذا قد دعوت اللّه و لادعونه انه قریب مجیب

\end{otherlanguage}

Sobḥān-e Allāh az tadbīrāt-e Modabber-e Ḥaqīqī va az ghaflat-e mā ʿebād.
Mā min nuṣḥin illā wa-qad dhakarahu wa-mā min tadbīrin illā wa-qad
anzalahu. Bayānī farmūde-and ke har ṣāḥeb-e baṣar va har ṣāḥeb-e foʾādī
nāle mīnamāyad va nowḥe mīkonad, che ke az kalemātash ʿālam ʿālam ḥozn
rū mīnamāyad va az bayānash ʿālam ʿālam hamm va ghamm. Qawluhu jalla
jalāluhu: fa-inna mithlahu jalla dhikruhu ka-mithli al-shamsi law
yuqābilannahu ilā mā lā nihāyata marāyāʾ kulluhunna la-yastaʿkisna ʿan
tajallī al-shamsi fī ḥaddihim wa-in lam yuqābilhā min aḥadin fa-yaṭluʿu
al-shamsu wa-yaghrubu al-ḥijābu li-l-marāyā wa-innī mā qaṣṣartu ʿan
nuṣḥī dhālika al-khalqa wa-tadbīrī li-iqbālihim ilā Allāhi ʿalā al-arḍ
fa-idhā yusarru kaynūnatī ḥaythu kullun qad balaghū ilā dhurwati
wujūdihim wa-waṣalū ilā ṭalʿati maḥbūbihim wa-adrakū mā yumkinu fī
al-imkāni min tajallī maqṣūdihim wa-illā yaḥzanu fuʾādī bi-annī qad
rabbaytu kulla shayʾin li-dhālika fa-kayfa taḥtajibu aḥadun ʿalā hādhā
qad daʿawtu Allāha wa-la-adʿūnahu innahu qarībun mujīb.

Glory be to God! Consider the provisions of the true Disposer and the
heedlessness of us servants! There is no counsel He has not given, and
no provision He has not sent down. He has spoken words that cause
everyone with sight and everyone with a heart to lament and mourn, for
from His words appear worlds upon worlds of sorrow, and from His
utterance worlds upon worlds of grief and anguish. He says, exalted be
His glory: ``He, exalted be His mention, is like the sun: if an infinite
number of mirrors were to face Him, they would all reflect the sun's
radiance according to their capacity. And if none faces it, the sun
still rises and sets; the veil belongs to the mirrors. I have spared no
effort in counseling these people and providing for their turning toward
God on earth. Then my being rejoices when all have reached the summit of
their existence, attained the countenance of their Beloved, and
apprehended all that is possible in the contingent world of the
revelation of their heart's Desire. Otherwise my heart grieves, for I
have nurtured all things for this. How, then, can anyone remain veiled?
For this I have called upon God, and I will surely call upon Him. He is
near and answers.''

\subsection{112}\label{section-2087}

\begin{otherlanguage}{arabic}

مع تدبیرات حضرت مدبر حقیقی و بیاناتی که دل کوه را میشکافد این خلق بضلالت
قدیم خود باقی وبرقرارند الا من شاء ربّک.

\end{otherlanguage}

Maʿ-e tadbīrāt-e Ḥażrat-e Modabber-e Ḥaqīqī va bayānātī ke del-e kūh rā
mīshekāfad īn khalq be-żalālat-e qadīm-e khod bāqī va barqarārand, illā
man shāʾa rabbuka.

Despite the provisions of the true Disposer and utterances that split
the heart of a mountain, these people remain firmly in their old error,
except those whom your Lord wills.

\subsection{113}\label{section-2088}

\begin{otherlanguage}{persian}

ای آقایان و ای دوستان الهی ناله‌های حضرت نقطه را بگوش جان بشنوید و همچنین
اسرار کتابش را و در ظهورات تدبیر و امتحان و افتتانش تفکر نمائید روحی
لعنایته الفداء و لذکره الفداء و لتدبیره الفداء و لنصحه الفداء و لامتحانه
الفداء میفرماید بیان و آنچه در او نازل شده او را میزان معرفت آن شمس
حقیقی قرار مدهید و میفرماید انه لا یستشار باشارتی باشاره من معروف نمیشود
ونه بآنچه در بیان نازل شده و در مقام دیگر بکلمات تحذیر نصح فرموده
میفرماید ایّاک ایّاک ان تحتجب بالواحد البیانیه و واحد بیانیه بمنزله ائمه
قبلند چنانچه خود حضرت در بیان فارسی میفرمایند بلکه میفرمایند جسد این
حروف بمقام روح ائمه قبلند مع ذلک میفرماید باین حروف از آن شمس حقیقت
محتجب نمانید و مکرر بکلمه تحذیر میفرماید بآیات منزله در بیان از او محجوب
نمانید

\end{otherlanguage}

Ey āqāyān va ey dūstān-e elāhī, nālehā-ye Ḥażrat-e Noqṭe rā be-gūsh-e
jān beshnavīd va hamchenīn asrār-e ketābash rā va dar ẓohūrāt-e tadbīr
va emteḥān va eftetānash tafakkor namāyīd. Rūḥī li-ʿināyatihi al-fidāʾ
wa-li-dhikrihi al-fidāʾ wa-li-tadbīrihi al-fidāʾ wa-li-nuṣḥihi al-fidāʾ
wa-li-imtiḥānihi al-fidāʾ. Mīfarmāyad: Bayān va ānche dar ū nāzel shode
ū rā mīzān-e maʿrefat-e ān Shams-e Ḥaqīqī qarār madahīd. Va mīfarmāyad:
innahu lā yustashāru bi-ishāratī. Be-eshāre-ye man maʿrūf nemīshavad va
na be-ānche dar Bayān nāzel shode. Va dar maqām-e dīgar be-kalemāt-e
taḥdhīr noṣḥ farmūde mīfarmāyad: iyyāka iyyāka an taḥtajiba bi-l-Wāḥidi
al-Bayāniyya. Va Vāḥed-e Bayāniyye be-manzale-ye aʾemme-ye qabland,
chonānche khod-e Ḥażrat dar Bayān-e Fārsī mīfarmāyand, balke mīfarmāyand
jasad-e īn ḥorūf be-maqām-e rūḥ-e aʾemme-ye qabland. Maʿ-e dhālek
mīfarmāyad: be-īn ḥorūf az ān Shams-e Ḥaqīqat moḥtajeb namānīd. Va
mokarrar be-kaleme-ye taḥdhīr mīfarmāyad: be-āyāt-e monzale dar Bayān az
ū maḥjūb namānīd.

Gentlemen and friends of God! Listen with the ear of your soul to the
lamentations of the Point, and likewise to the mysteries of His Book,
and reflect on the manifestations of His providence, His testing, and
His trials. May my soul be sacrificed for His loving care, His mention,
His providence, His counsel, and His testing! He says: ``Do not make the
Bayān and what has been revealed in it the measure by which to recognize
that true Sun.'' And He says: ``He is not to be consulted through my
indication.'' He is not made known by my indication, nor by what has
been revealed in the Bayān. Elsewhere, counseling with words of warning,
He says: ``Beware, beware lest you veil yourself through the Váḥid of
the Bayān.'' The Váḥid of the Bayān occupies the position of the Imams
of the past, as He Himself states in the Persian Bayān. Indeed, He says
that the bodies of these Letters occupy the station of the spirits of
the Imams of the past. Nevertheless, He says: ``Do not let these Letters
veil you from that Sun of Truth.'' And repeatedly He says in words of
warning: ``Do not let the verses revealed in the Bayān veil you from
Him.''

\subsection{114}\label{section-2089}

\begin{otherlanguage}{persian}

قسم بآفتاب افق معانی در مقامات شتی این بیانات از قلم مبارکش جاری چنانچه
بعضی از آن ذکر شد و بنظر اولیای حق میرسد و بآن فائز میشوند مع جمیع این
وصیتها و ذکرها و بیانها بذکر یک کلمه مستغاث کل را امتحان فرمود و کلمه
مستغاث بیان است و میفرماید بجمیع بیان از او محتجب نشوید باین امتحان جزئی
کل را متوقف مشاهده فرمود در این حین گویا منادی عظمت از افق ملکوت بیان
رحمن باعلی النداء ندا فرمود و فرمود الان قد حصحص الحق و ظهر کذب الذین
یدعون الایمان بالبیان و یکفرون بالذی انزله و ارسله

\end{otherlanguage}

Qasam bi-āftāb-i ufuq-i maʿānī dar maqāmāt-i shattā īn bayānāt az
qalam-i mubārakash jārī, chunānchih baʿżī az ān dhikr shud va bi-naẓar-i
awliyā-yi Ḥaqq mīrasad va bi-ān fāʾiz mīshavand. Maʿ-i jamīʿ-i īn
vaṣiyyathā va dhikrhā va bayānhā bi-dhikr-i yak kalimih Mustaghāth kull
rā imtiḥān farmūd va kalimih-yi Mustaghāth Bayān ast va mīfarmāyad
bi-jamīʿ-i Bayān az ū muḥtajib nashavīd. Bi-īn imtiḥān-i juzʾī kull rā
mutavaqqif mushāhidih farmūd. Dar īn ḥīn gūyā munādī-yi ʿaẓamat az
ufuq-i malakūt-i bayān-i Raḥmān bi-aʿlā al-nidāʾ nidā farmūd va farmūd:
al-āna qad ḥaṣḥaṣa al-ḥaqq wa ẓahara kadhib alladhīna yaddaʿūna al-īmān
bi-l-Bayān wa yakfurūna bi-alladhī anzalahu wa arsalahu.

I swear by the Sun of the horizon of meanings! In various passages these
utterances have flowed from His blessed Pen, as some of them have
already been cited and are coming before the eyes of the friends of God,
who are attaining them. Despite all these counsels, reminders, and
explanations, He tested everyone by mentioning the single word
``Mustaghāth.'' Yet the word ``Mustaghāth'' belongs to the Bayān, and He
says: ``Do not let the whole Bayān veil you from Him.'' Through this
slight test He found everyone brought to a standstill. At this moment it
was as though the Herald of grandeur called out at the top of His voice
from the horizon of the kingdom of the All-Merciful's utterance, saying:
``Now the truth has become manifest, and the falsehood of those who
profess belief in the Bayān yet disbelieve in the One who revealed and
sent it down has been exposed.''

\subsection{115}\label{section-2090}

\begin{otherlanguage}{persian}

ایکاش بتوقف اکتفاء مینمودند بعضی بسهام اعراض و بعضی باسنه انکار و برخی
باسیاف آخته بر قطع سدره مبارکه کوشیده و میکوشند چنانچه منزل بیان بکمال
تصریح اخبار فرموده میفرماید بقدراسم مؤمن هم در حق آن شجره لاشرقیه و
لاغربیه راضی نمیشوند چه اگر راضی شوند حزن بر او وارد نمیآورند باری ذکر
اینمقامات چون مطلع حزن است این فانی باظهارش بیش از این جسارت نمینماید

\end{otherlanguage}

Ay kāsh bi-tavaqquf iktifā mīnamūdand. Baʿżī bi-sihām-i iʿrāż va baʿżī
bi-asinnah-yi inkār va barkhī bi-asyāf-i ākhtih bar qaṭʿ-i Sidrih-yi
Mubārakih kūshīdih va mīkūshand, chunānchih Munzil-i Bayān bi-kamāl-i
taṣrīḥ ikhbār farmūdih, mīfarmāyad: bi-qadr-i ism-i muʾmin ham dar
ḥaqq-i ān shajarih-yi lā-sharqiyyih va lā-gharbiyyih rāżī namīshavand,
chih agar rāżī shavand ḥuzn bar ū vārid namīāvarand. Bārī dhikr-i īn
maqāmāt chūn maṭlaʿ-i ḥuzn ast īn fānī bi-iẓhārash bīsh az īn jisārat
namīnamāyad.

If only they had been content to remain at a standstill! Some, with the
arrows of aversion, some with the spears of denial, and others with
drawn swords, have striven and still strive to cut down the blessed
Tree. The Revealer of the Bayān explicitly foretold this, saying: ``They
will not consent to grant that Tree, neither of the East nor of the
West, even the name of believer; for if they consented, they would not
bring sorrow upon Him.'' Since recounting these matters gives rise to
sorrow, this mortal does not venture to speak further of them.

\subsection{116}\label{section-2091}

\begin{otherlanguage}{persian}

سبحان اللّه کل بیک کلمه مستغاث مغرور و متحیر و ممتحن مع آنکه این فقره هم
ذکر فرموده‌اند و تمسک بآن را هم نهی فرموده‌اند قوله تعالی چه که کسی عالم
بظهور نیست غیر اللّه هر وقت شود باید کل تصدیق بنقطه حقیقت نمایند و شکر
الهی بجای آورند هذا حق لاریب فیه

\end{otherlanguage}

Subḥāna Allāh! Kull bi-yak kalimih Mustaghāth maghrūr va mutaḥayyir va
mumtaḥan, maʿ-i ānkih īn faqarih ham dhikr farmūdih-and va tamassuk
bi-ān rā ham nahy farmūdih-and. Qawluhu taʿālā: chih kih kasī ʿālim
bi-ẓuhūr nīst ghayr-i Allāh, har vaqt shavad bāyad kull taṣdīq
bi-Nuqṭih-yi Ḥaqīqat namāyand va shukr-i ilāhī bi-jāy āvarand. Hādhā
ḥaqq lā rayba fīhi.

Glory be to God! All are deluded, bewildered, and tested by the single
word ``Mustaghāth,'' even though He also stated the following and
forbade them to cling to it. His words, exalted is He: ``For no one
knows of the Revelation except God. Whenever it takes place, all must
acknowledge the Point of Truth and offer thanks to God.'' This is the
truth, beyond all doubt.

\subsection{117}\label{section-2092}

\begin{otherlanguage}{persian}

امروز کتب عالم حجاب نمیشود و منع نمینماید وکل بکلمه حق جل جلاله معلق و
منوط ای صاحبان آذان بشنوید ندای مقصود عالمیان را که امام وجوه عالم از
اول امر تا حین ندا فرموده و میفرماید در یکی از الواح این کلمه علیا از
فالق الاصباح ظاهر قوله عمّ نواله امروز روز بصر است چه که افق اعلی مشرق و
یوم سمع است چه که نداء اللّه در کل حین مرتفع انتهی

\end{otherlanguage}

Imrūz kutub-i ʿālam ḥijāb namīshavad va manʿ namīnamāyad va kull
bi-kalimih-yi Ḥaqq jalla jalāluhu muʿallaq va manūṭ. Ay ṣāḥibān-i ādhān,
bishnavīd nidā-yi Maqṣūd-i ʿālamiyān rā kih amām-i vujūh-i ʿālam az
avval-i amr tā ḥīn nidā farmūdih va mīfarmāyad. Dar yakī az alvāḥ īn
kalimih-yi ʿulyā az Fāliq al-Aṣbāḥ ẓāhir, qawluhu ʿamma nawāluhu: imrūz
rūz-i baṣar ast chih kih ufuq-i aʿlā mushriq va yawm-i samʿ ast chih kih
nidāʾ Allāh dar kull-i ḥīn murtafiʿ. Intahā.

Today the books of the world cannot serve as a veil or a hindrance, and
everything depends upon and is subject to the word of God, exalted be
His glory. You who have ears, hear the call of the Desire of all
peoples, who has called out and continues to call out before the eyes of
the world from the beginning of the Cause until now. In one of the
Tablets this exalted utterance appeared from the Cleaver of the Dawn.
His words, all-embracing is His bounty: ``Today is the day of sight, for
the Most Exalted Horizon is radiant, and the day of hearing, for the
call of God is raised at every moment.'' End of quotation.

\subsection{118}\label{section-2093}

\begin{otherlanguage}{persian}

ای دوستان آیا نفوس غافله چه گمان نموده‌اند اگر از ملکوت آیات الهی و سماء
ظهور ربانی محروم مانند خود را بچه امری تسلی میدهند و مؤمن میدانند امروز
بنص جمیع کتب هر نفسی بغیر حق وحده ناظر باشد او از اهل ضلال بوده و خواهد
بود یشهد بذلک کتب اللّه المهیمن القیوم

\end{otherlanguage}

Ay dūstān, āyā nufūs-i ghāfilih chih gumān namūdih-and? Agar az
malakūt-i āyāt-i ilāhī va samāʾ-i ẓuhūr-i rabbānī maḥrūm mānand khud rā
bi-chih amrī tasallī mīdihand va muʾmin mīdānand? Imrūz bi-naṣṣ-i
jamīʿ-i kutub har nafsī bi-ghayr-i Ḥaqq vaḥdahu nāẓir bāshad ū az ahl-i
ḍalāl būdih va khvāhad būd. Yashhadu bi-dhālika kutub Allāh al-Muhaymin
al-Qayyūm.

Friends! What have these heedless souls imagined? If they remain
deprived of the kingdom of divine verses and the heaven of divine
Revelation, with what will they console themselves and consider
themselves believers? Today, according to the explicit text of all the
Books, whoever looks to anyone other than God alone has been and will be
among those who have gone astray. To this bear witness the Books of God,
the All-Protecting, the Self-Subsisting.

\subsection{119}\label{section-2094}

\begin{otherlanguage}{persian}

بمفتریاتی تمسک جسته‌اند که انسان متحیر است از جمله حضرت اسم اللّه جناب حاج
سید جواد کربلائی علیه بهاء اللّه الابهی را نسبت بتوقف در این امر و یا
نعوذ باللّه انکار داده‌اند عرایض و مکتوباتی که بخط خود نوشته‌اند و بساحت
اقدس ارسال نموده‌اند موجود است بعضی از سمت ارض طا و بعضی از شطر یا و خا و
جمعی مطلعند در یک مکتوب ذکر متوهمین را بتفصیل نموده‌اند از رئیس و غیره حق
شاهد و گواه که از روی راستی عرض میشود حق محتاج باین اذکار نبوده و نیست
اگر باین عوالم توجه نمائیم و باین قسم از اذکار دفاتر عدیده موجود شود و
در هر یک عالمی مبسوط

\end{otherlanguage}

Bi-muftarayātī tamassuk justih-and kih insān mutaḥayyir ast. Az jumlih
Ḥaḍrat-i Ism Allāh Jināb-i Ḥājj Sayyid Javād-i Karbalāʾī ʿalayhi Bahāʾ
Allāh al-Abhā rā nisbat bi-tavaqquf dar īn amr va yā naʿūdhu bi-Allāh
inkār dādih-and. ʿArāyiż va maktūbātī kih bi-khaṭṭ-i khud nivishtih-and
va bi-sāḥat-i aqdas irsāl namūdih-and mawjūd ast, baʿżī az samt-i arż-i
Ṭā va baʿżī az shaṭr-i Yā va Khā, va jamʿī muṭṭaliʿand. Dar yak maktūb
dhikr-i mutavahhimīn rā bi-tafṣīl namūdih-and az raʾīs va ghayrih. Ḥaqq
shāhid va guvāh kih az rūy-i rāstī ʿarż mīshavad. Ḥaqq muḥtāj bi-īn
adhkār nabūdih va nīst. Agar bi-īn ʿavālim tavajjuh namāyīm va bi-īn
qism az adhkār dafātir-i ʿadīdih mawjūd shavad va dar har yak ʿālamī
mabsūṭ.

They have clung to fabrications that leave one bewildered. Among other
things, they have accused His Honor the Name of God, Ḥājj Sayyid Javād
of Karbalā---upon him be the glory of God, the Most Glorious---of
hesitating in this Cause or even, God forbid, denying it. The petitions
and letters that he wrote in his own hand and sent to the Most Holy
Presence are extant, some from the direction of the land of Ṭā and some
from the regions of Yā and Khā, and a number of people know of them. In
one letter he discussed in detail those who harbor vain imaginings,
including their leader and others. God bears witness that this is stated
truthfully. God has never needed, nor does He now need, these accounts.
If we were to turn our attention to these matters and accounts of this
kind, numerous volumes would result, each unfolding a world of its own.

\subsection{120}\label{section-2095}

\begin{otherlanguage}{persian}

وقتی از اوقات این کلمات از مظهر بینات نازل قوله جل جلاله و عظم شأنه یا
عبد حاضر یک ذبیح در عالم یافت شد و از آن یوم تا حین السن خلایق و کتب
موجوده بذکرش ناطق و بر علو مقامش شاهد و ما در هر مقامی ذبیحی فرستادیم و
کل از قربانگاه دوست زنده بر نگشتند هر یک اکلیل حیات را رایگان بکمال شوق
و اشتیاق نثار قدوم محبوب امکان نمودند. انتهی

\end{otherlanguage}

Vaqtī az awqāt īn kalimāt az Maẓhar-i bayyināt nāzil, qawluhu jalla
jalāluhu va ʿaẓuma shaʾnuhu: yā ʿabd-i ḥāżir, yak dhabīḥ dar ʿālam yāft
shud va az ān yawm tā ḥīn alsun-i khalāyiq va kutub-i mawjūdih
bi-dhikrash nāṭiq va bar ʿuluvv-i maqāmash shāhid va mā dar har maqāmī
dhabīḥī firistādīm va kull az qurbāngāh-i Dūst zindih bar nagashtand.
Har yak iklīl-i ḥayāt rā rāygān bi-kamāl-i shawq va ishtiyāq nithār-i
qudūm-i Maḥbūb-i imkān namūdand. Intahā.

On one occasion these words were revealed by the Manifestation of clear
proofs. His words, exalted be His glory and great be His station:
``Servant in attendance! One sacrificial victim appeared in the world,
and from that day until now the tongues of the people and the books in
existence have spoken of him and borne witness to the loftiness of his
station. We, however, sent a sacrificial victim to every place, and none
returned alive from the sacrificial ground of the Friend. Each freely
laid the crown of life at the feet of the Beloved of all existence with
the utmost joy and longing.'' End of quotation.

\subsection{121}\label{section-2096}

\begin{otherlanguage}{persian}

از جمله محبوب شهدا جناب آقا سید اسمعیل زوارهٴ علیه بهاء اللّه الابهی که
بدست خود خود را مقبلا الی البیت فدا نمود اگر چه این عمل در ظاهر منکر
ولکن محبت الهی چنان اخذش نمود که از هر عرقی از عروقش شعله نار ظاهر و
باهر و همچنین حضرت ابا بصیر و آقا سید اشرف علیهما بهاء اللّه و عنایته از
افق بیت باسم حق جل جلاله طالع شدند و من غیر ستر و حجاب بذکر محبوب آفاق
ناطق مشرکین هر دو را اخذ نموده اول بسجن فرستادند امّ اشرف را طلب نموده که
ابن خود را نصیحت کند که شاید اقبالش را تبدیل نماید و یا امرش را مستور
دارد ولکن آن امة ثابته راسخه مستقیمه بعد از ورود در سجن فرمود ای پسر در
امر اللّه مستقیم باش مباد خوف نمائی و یا از سطوت مشرکین مضطرب شوی باری او
و ابا بصیر هر دو روح را بکمال روح و ریحان در ره دوست انفاق نمودند

\end{otherlanguage}

Az jumlih Maḥbūb-i shuhadā Jināb-i Āqā Sayyid Ismaʿīl-i Zavārih ʿalayhi
Bahāʾ Allāh al-Abhā kih bi-dast-i khud khud rā muqbilan ilā al-Bayt fidā
namūd. Agarchih īn ʿamal dar ẓāhir munkar valīkin maḥabbat-i ilāhī
chunān akhdhash namūd kih az har ʿirqī az ʿurūqash shuʿlih-yi nār ẓāhir
va bāhir. Va hamchunīn Ḥaḍrat-i Abā Baṣīr va Āqā Sayyid Ashraf ʿalayhimā
Bahāʾ Allāh va ʿināyatuhu az ufuq-i bayt bi-ism-i Ḥaqq jalla jalāluhu
ṭāliʿ shudand va min ghayr-i sitr va ḥijāb bi-dhikr-i Maḥbūb-i āfāq
nāṭiq. Mushrikīn har du rā akhdh namūdih avval bi-sijn firistādand.
Umm-i Ashraf rā ṭalab namūdih kih ibn-i khud rā naṣīḥat kunad kih shāyad
iqbālash rā tabdīl namāyad va yā amrash rā mastūr dārad. Valīkin ān
amat-i thābitih-yi rāsikhih-yi mustaqīmih baʿd az vurūd dar sijn farmūd:
ay pisar, dar Amr Allāh mustaqīm bāsh, mabād khawf namāʾī va yā az
saṭvat-i mushrikīn muḍṭarib shavī. Bārī ū va Abā Baṣīr har du rūḥ rā
bi-kamāl-i rawḥ va rayḥān dar rah-i Dūst infāq namūdand.

Among them was the Beloved of Martyrs, His Honor Āqā Sayyid Ismaʿīl of
Zavārih---upon him be the glory of God, the Most Glorious---who
sacrificed himself by his own hand while facing the House. Although this
act is outwardly reprehensible, the love of God so seized him that a
flame of fire shone forth from every one of his veins. Likewise, His
Honor Abā Baṣīr and Āqā Sayyid Ashraf---upon them both be the glory of
God and His loving care---rose from the horizon of their home in the
name of God, exalted be His glory, and openly, without concealment,
spoke in praise of the Beloved of all lands. The idolaters seized them
both and first sent them to prison. They summoned Ashraf's mother to
counsel her son, hoping that he might turn away from his allegiance or
conceal his cause. But that steadfast, firmly grounded, and unwavering
handmaid said upon entering the prison: ``My son, be steadfast in the
Cause of God. Never be afraid or dismayed by the might of the
idolaters.'' In short, he and Abā Baṣīr both gave their souls in the
path of the Friend with the utmost joy and delight.

\subsection{122}\label{section-2097}

\begin{otherlanguage}{arabic}

و در جناب نجفعلی علیه بهاء اللّه تفکر نمائید هنگامیکه او را بقربانگاه
دوست میبردند باین فرد ناطق ``ما بها و خون بها را یافتیم'' مکرر باین فرد
ناطق تاآنکه بشهادت فائز گشت

\end{otherlanguage}

Va dar Jināb-i NajafʿAlī ʿalayhi Bahāʾ Allāh tafakkur namāʾīd. Hingāmī
kih ū rā bi-qurbāngāh-i Dūst mīburdand bi-īn fard nāṭiq: ``Mā bahā va
khūn-bahā rā yāftīm.'' Mukarrar bi-īn fard nāṭiq tā ānkih bi-shahādat
fāʾiz gasht.

And reflect on His Honor NajafʿAlī---upon him be the glory of God. As
they were taking him to the sacrificial ground of the Friend, he recited
this verse: ``We have found both the price and the price of blood.'' He
repeated this verse until he attained martyrdom.

\subsection{123}\label{section-2098}

\begin{otherlanguage}{persian}

ای آقایان قصه بدیع را بشنوید فی الحقیقه از هر حرکتی از حرکاتش آثار قدرت
الهی و شوکت صمدانی ظاهر و باهر در اول ایّام سجن اعظم وارد شد هنگامی که
مقر عرش قشله عسکریه واقع دو یوم او را در بیت مخصوص طلب فرمودند باب مسدود
و احدی تلقاء وجه غیر او موجود نه و احدی هم مطلع نبود که مقصود چیست تا
آنکه فرمودند حق اراده خلق جدید نموده و خود بدیع هم آگاه نه در مقامی این
کلمه علیا از قلم اعلی در لوحی از الواح نازل قوله عزّ بیانه انا شرعنا فی
خلق البدیع فلما تم خلقه و طاب خلقه ارسلناه ککرة النّار الی آخر قوله تعالی

\end{otherlanguage}

Ay āqāyān, qiṣṣih-yi Badīʿ rā bishnavīd. Fī al-ḥaqīqih az har ḥarakatī
az ḥarakātash āthār-i qudrat-i ilāhī va shawkat-i ṣamadānī ẓāhir va
bāhir. Dar avval-i ayyām-i Sijn-i Aʿẓam vārid shud, hingāmī kih maqarr-i
ʿarsh qishlih-yi ʿaskariyyih vāqiʿ. Du yawm ū rā dar bayt-i makhṣūṣ
ṭalab farmūdand, bāb masdūd va aḥadī tilqāʾ-i vajh ghayr-i ū mawjūd nah
va aḥadī ham muṭṭaliʿ nabūd kih maqṣūd chīst, tā ānkih farmūdand: Ḥaqq
irādih-yi khalq-i jadīd namūdih, va khud-i Badīʿ ham āgāh nah. Dar
maqāmī īn kalimih-yi ʿulyā az Qalam-i Aʿlā dar lawḥī az alvāḥ nāzil,
qawluhu ʿazza bayānuhu: innā sharaʿnā fī khalq al-Badīʿ fa-lammā tamma
khalquhu wa ṭāba khuluquhu arsalnāhu ka-kurat al-nār, ilā ākhir-i
qawlihi taʿālā.

Gentlemen! Hear the story of Badīʿ. Truly, in his every movement the
signs of divine power and eternal majesty were plainly manifest. He
arrived in the early days of the Most Great Prison, when the throne was
established in the military barracks. For two days He summoned him to a
private room; the door was closed, and no one else was present before
Him. No one knew the purpose until He said: ``God has willed a new
creation,'' and Badīʿ himself was unaware of it. In one passage of a
Tablet this exalted utterance was revealed by the Most Exalted Pen. His
words, mighty is His utterance: ``We began the creation of Badīʿ, and
when his creation was complete and his character refined, We sent him
forth like a ball of fire,'' and so on to the end of His words, exalted
is He.

\subsection{124}\label{section-2099}

\begin{otherlanguage}{persian}

و بعد مرخص شد جناب امین علیه بهاء اللّه با لوح حضرت سلطان حسب الامر بوطن
توجه در اسکله بحر حضرت بدیع از لوح اطلاع یافته استدعا نمود او حامل شود و
بعد کل استماع نمودند که چه واقع شد شخصی من غیر سلاح و من دون آلات جز
قمیصی از کرباس در بر نداشت بقدرت و قوت الهی در مقابل سلطان لوح را بلند
نمود و قال قد جئتک من السبا الاعظم بکتاب عظیم یا کلمه اخری انصاف میطلبم
این یکنفس مقابل عالم ایستاد باری از قدرت الهی عجیب نه که عالم قدرت را در
آدمی مبعوث فرماید هو المقتدر علی ما یشاء و هو المهیمن علی ما یرید

\end{otherlanguage}

Va baʿd murakhkhaṣ shud. Jināb-i Amīn ʿalayhi Bahāʾ Allāh bā lawḥ-i
Ḥaḍrat-i Sulṭān ḥasb al-amr bi-vaṭan tavajjuh. Dar iskalih-yi baḥr
Ḥaḍrat-i Badīʿ az lawḥ iṭṭilāʿ yāftih istidʿā namūd ū ḥāmil shavad va
baʿd kull istimāʿ namūdand kih chih vāqiʿ shud. Shakhṣī min ghayr-i
silāḥ va min dūn-i ālāt juz qamīṣī az karbās dar bar nadāsht, bi-qudrat
va quvvat-i ilāhī dar muqābil-i sulṭān lawḥ rā buland namūd wa qāla: qad
jiʾtuka min al-sabā al-aʿẓam bi-kitābin ʿaẓīm, yā kalimih-yi ukhrā.
Inṣāf mīṭalabam, īn yak nafs muqābil-i ʿālam īstād. Bārī az qudrat-i
ilāhī ʿajīb nah kih ʿālam-i qudrat rā dar ādamī mabʿūth farmāyad. Huwa
al-Muqtadir ʿalā mā yashāʾ wa huwa al-Muhaymin ʿalā mā yurīd.

He was then dismissed. His Honor Amīn---upon him be the glory of
God---set out for his homeland, as instructed, with the Tablet to His
Majesty the King. At the seaport, His Honor Badīʿ learned of the Tablet
and asked to be its bearer. Everyone subsequently heard what happened. A
man without weapons or equipment, wearing nothing but a cotton shirt,
raised the Tablet before the king through divine strength and power and
said: ``I have come to you from the greatest captivity with a mighty
Book,'' or words to that effect. I ask for fairness: this one soul stood
against the world. Yet it is no wonder that divine power should raise up
a world of might within a human being. He has power over whatever He
wills, and He is supreme over whatever He desires.

\subsection{125}\label{section-2100}

\begin{otherlanguage}{persian}

بگوئید ای بی انصافان شما شهادت سید الشهداء را حجت و دلیل بر اعظم حقیت
رسول مختار روح ما سواه فداه میدانستید و همچنین ذبح اسمعیل را از علو مقام
آنحضرت و خلیل میشمردید حال در این فدائیها تفکر نمائید لعمر محبوبی و
محبوب العالم و مقصودی و مقصود الامم جناب بدیع بقدرت و اطمینانی ظاهر که
تزلزل و اضطراب در ارکان ظاهر شد البته کل شنیده‌اند و منصفین انکار
نمینمایند و بعد چند یوم بکمال عذاب و شکنجه آن هیکل مقدس معذب و بکمال
استقامت بذکر دوست ذاکر و ناطق و این فقره سبب حیرت کل شد بشأنی در حینی که
بین ایادی میرغضبها با سلاسل و اغلال مبتلا عکس او را اخذ نمودند و یکی از
آنرا بعضی بساحت اقدس ارسال داشتند و حال موجود است

\end{otherlanguage}

Bigūʾīd ay bī-inṣāfān, shumā shahādat-i Sayyid al-Shuhadā rā ḥujjat va
dalīl bar aʿẓam-i ḥaqqiyyat-i Rasūl-i Mukhtār rūḥ-i mā sivāhu fidāhu
mīdānistīd va hamchunīn dhibḥ-i Ismaʿīl rā az ʿuluvv-i maqām-i ān Ḥaḍrat
va Khalīl mīshumurdīd. Ḥāl dar īn fidāʾīhā tafakkur namāʾīd. La-ʿamr-i
Maḥbūbī va Maḥbūb-i ʿālam va Maqṣūdī va Maqṣūd-i umam, Jināb-i Badīʿ
bi-qudrat va iṭmīnānī ẓāhir kih tazalzul va iḍṭirāb dar arkān ẓāhir
shud. Albattih kull shinīdih-and va munṣifīn inkār namīnamāyand. Va baʿd
chand yawm bi-kamāl-i ʿadhāb va shikanjih ān haykal-i muqaddas
muʿadhdhab va bi-kamāl-i istiqāmat bi-dhikr-i Dūst dhākir va nāṭiq va īn
faqarih sabab-i ḥayrat-i kull shud, bi-shaʾnī dar ḥīnī kih bayn-i
ayādī-yi mīrghaḍabhā bā salāsil va aghlāl mubtalā ʿaks-i ū rā akhdh
namūdand va yakī az ān rā baʿżī bi-sāḥat-i aqdas irsāl dāshtand va ḥāl
mawjūd ast.

Say: You who are without fairness! You regarded the martyrdom of the
Prince of Martyrs as proof and evidence of the supreme truth of the
chosen Messenger---may the souls of all others be sacrificed for
Him---and likewise counted the sacrifice of Ishmael as a sign of the
exalted station of both him and the Friend of God. Now reflect on these
self-sacrificing souls. By the life of my Beloved and the Beloved of the
world, my Desire and the Desire of the nations! His Honor Badīʿ appeared
with such power and assurance that trembling and turmoil seized the
pillars of authority. Surely everyone has heard, and the fair-minded do
not deny it. Then, for several days, that sacred person was subjected to
the severest torment and torture, yet with perfect steadfastness
continued to remember and speak of the Friend. This astonished everyone,
so much so that while he was in the hands of the executioners, burdened
with chains and fetters, they took his photograph. Someone sent one of
these to the Most Holy Presence, and it is still extant.

\subsection{126}\label{section-2101}

\begin{otherlanguage}{persian}

باری در این امورات که شبه و مثل نداشته تفکر نمائید این فانی نمیداند کدام
ذبیح را ذکر نماید اب جناب بدیع علیه بهاء اللّه رادر ارض خا اخذ نمودند
آنچه نواب شاهزاده و سایر ناس خواستند ستر نماید تا اورا از ضوضاء علماء
حفظ کنند قبول نفرمودند و فرمودند آخر ایام من است و شهادت از هر شهدی
خوشتر و شیرین‌تر و بآن فائز شد

\end{otherlanguage}

Bārī dar īn umūrāt kih shibh va mithl nadāshtih tafakkur namāʾīd. Īn
fānī namīdānad kudām dhabīḥ rā dhikr namāyad. Ab-i Jināb-i Badīʿ ʿalayhi
Bahāʾ Allāh rā dar arż-i Khā akhdh namūdand. Ānchih Navvāb-i Shāhzādih
va sāyir-i nās khvāstand satr namāyad tā ū rā az ḍawḍāʾ-i ʿulamā ḥifẓ
kunand qabūl nafarmūdand va farmūdand: ākhir-i ayyām-i man ast va
shahādat az har shahdī khvushtar va shīrīntar, va bi-ān fāʾiz shud.

Reflect, then, on these events, which have had no parallel or equal.
This mortal does not know which sacrificial victim to recount. They
arrested the father of His Honor Badīʿ---upon him be the glory of
God---in the land of Khā. However much His Highness the Prince and
others urged him to conceal his faith so that they could protect him
from the clamor of the clergy, he refused, saying: ``These are my last
days, and martyrdom is more delightful and sweeter than any honey.'' And
he attained it.

\subsection{127}\label{section-2102}

\begin{otherlanguage}{persian}

و همچنین ذبیح ارض ن ر در ارض تا الذی سمی بمیرزا مصطفی علیه بهاء اللّه
الابهی و نفوسیکه با او بودند بکمال استقامت و منتهای رضا بمشهد فداء متوجه
و اکلیل حیات را نثار قدوم مالک اسماء و صفات نمود و این نفوس نفوسی هستند
که رایگان جان را فدا نمودند ان اعتبروا یا اولی الابصار

\end{otherlanguage}

Va hamchunīn dhabīḥ-i arż-i Nūn Rā dar arż-i Tā alladhī summiya bi-Mīrzā
Muṣṭafā ʿalayhi Bahāʾ Allāh al-Abhā va nufūsī kih bā ū būdand bi-kamāl-i
istiqāmat va muntahā-yi riżā bi-mashhad-i fidāʾ mutavajjih va iklīl-i
ḥayāt rā nithār-i qudūm-i Mālik-i asmāʾ va ṣifāt namūd. Va īn nufūs
nufūsī hastand kih rāygān jān rā fidā namūdand. An iʿtabirū yā ulī
al-abṣār.

Likewise, the sacrificial victim of the land of Nūn Rā, in the land of
Tā, who was named Mīrzā Muṣṭafā---upon him be the glory of God, the Most
Glorious---and the souls who were with him, proceeded to the place of
sacrifice with perfect steadfastness and utmost contentment, and he laid
the crown of life at the feet of the Lord of names and attributes. These
are souls who freely sacrificed their lives. Take heed, you who have
insight!

\subsection{128}\label{section-2103}

\begin{otherlanguage}{persian}

حال بهتر آنکه بارض صاد رویم و در حضرت سلطان الشهداء و محبوب الشهداء و
انقطاع آن دو نفس مقدس از ما سوی اللّه تفکر نمائیم و انصاف دهیم آنچه
خواستند ایشان ستر نمایند و یا کلمه بگویند قبول نفرمودند بکمال شوق و
اشتیاق قصد وطن اعلی نمودند و بعد لوح برهان از سماء امر رحمن نازل و
درباره آن دو نفس یعنی ذئب و رقشاء نازل شد آنچه هر منصفی عرف قدرت و علم
الهی را از آن ادراک مینماید قوله تبارک و تعالی هو المقتدر العلیم الحکیم
قد احاطت اریاح البغضاء سفینة البطحاء\ldots\ldots{} '' الخ.

\end{otherlanguage}

Ḥāl behtar ānke be-arḍ-e Ṣād ravim va dar Ḥażrat-e Sulṭān al-Shuhadāʾ va
Maḥbūb al-Shuhadāʾ va enqeṭāʿ-e ān do nafs-e moqaddas az mā sivā Allāh
tafakkor namāyim va enṣāf dehim. Ānche khvāstand ishān satr namāyand va
yā kalame beguyand qabūl nafarmūdand. Be-kamāl-e showq va eshtiyāq
qaṣd-e vaṭan-e aʿlā namūdand va baʿd Lawḥ-e Burhān az samā-ye amr-e
Raḥmān nāzel va darbāre-ye ān do nafs yaʿni Dhiʾb va Raqshāʾ nāzel shod
ānche har monṣefi ʿarf-e qodrat va ʿelm-e elāhi rā az ān edrāk
minamāyad. Qawluhu tabāraka wa-taʿālā: Huwa al-Muqtadir al-ʿAlīm
al-Ḥakīm. Qad aḥāṭat aryāḥ al-baghḍāʾ safīnat al-Baṭḥāʾ\ldots\ldots{}
ilā ākhirihi.

It is better now that we turn to the land of Ṣād and reflect on His
Honor the King of Martyrs and the Beloved of Martyrs, and on the
detachment of those two holy souls from all but God, and judge fairly.
However much they were urged to conceal their faith or utter a word,
they refused. With the utmost longing and eagerness, they set out for
the heavenly homeland. Afterward, the Tablet of the Proof descended from
the heaven of the All-Merciful's Cause, and concerning those two
souls---the Wolf and the She-Serpent---there was revealed that from
which every fair-minded person perceives the fragrance of divine power
and knowledge. These are His words, blessed and exalted is He: ``He is
the All-Powerful, the All-Knowing, the All-Wise. The winds of hatred
have encompassed the Ark of Baṭḥā\ldots\ldots'' and so forth.

\subsection{129}\label{section-2104}

\begin{otherlanguage}{persian}

ای صاحبان بصر قدری تفکر نمائید شاید حلاوت بیان رحمن را بیابید و خود را
از هیماء جهل و نادانی نجات دهید از ذکر حجج و براهین و آیات وبینات مقصود
آنکه شاید ابوابهای بسته بگشاید و عهدهای شکسته بسته شود آنچه از قلم اعلی
درباره رقشا و ذئب نازل واقع و کل مشاهده نمودند این عبد متحیر آیا شعور
بالمره تمام شده یا انصاف عنقا گشته چه شده که سکر غفلت کل را احاطه نموده
آیا ذائقه را چه منع نموده و شامه را چه حادث شده نیست اینها مگر از جزای
اعمال از خدا میطلبم تائید عنایت فرماید

\end{otherlanguage}

Ey ṣāḥebān-e baṣar, qadri tafakkor namāyid shāyad ḥalāvat-e bayān-e
Raḥmān rā biyābid va khod rā az haymā-ye jahl va nādāni nejāt dahid. Az
dhekr-e ḥojaj va barāhin va āyāt va bayyenāt maqṣūd ānke shāyad
abvābhā-ye baste bogshāyad va ʿahdhā-ye shekaste baste shavad. Ānche az
Qalam-e Aʿlā darbāre-ye Raqshā va Dhiʾb nāzel vāqeʿ va koll moshāhede
namūdand. In ʿabd motaḥayyer: āyā shoʿūr be-l-marre tamām shode yā enṣāf
ʿanqā gashte? Che shode ke sokr-e ghaflat koll rā eḥāṭe namūde? Āyā
dhāʾeqe rā che manʿ namūde va shāmme rā che ḥādeth shode? Nist inhā
magar az jazā-ye aʿmāl. Az Khodā miṭalabam taʾyid ʿenāyat farmāyad.

You who possess sight, reflect a little, that you may discover the
sweetness of the All-Merciful's utterance and deliver yourselves from
the desert of ignorance and unawareness. The purpose of recounting
proofs, arguments, verses, and clear evidences is that closed doors may
perhaps be opened and broken covenants renewed. Whatever the Most
Exalted Pen revealed concerning the She-Serpent and the Wolf came to
pass, and everyone witnessed it. This servant is bewildered: has
understanding utterly vanished, or has fairness become as elusive as the
phoenix? What has happened that the intoxication of heedlessness has
encompassed everyone? What has impeded the sense of taste, and what has
befallen the sense of smell? These things are nothing but the recompense
for deeds. I ask God to graciously bestow His assistance.

\subsection{130}\label{section-2105}

\begin{otherlanguage}{persian}

بار الها کریما رحیما چه کفران از عباد ظاهر شده که از امواج بحر رحمت
محروم مانده‌اند و از اشراقات انوار آفتاب ظهور ممنوع گشته‌اند الهی اعمال
جهال را باسم ستارت ستر فرما توئی کریمی که ذنوب مذنبین بخششت را منع ننمود
و باب فیضت را سد نکرد یک حرف از کتاب جودت علت وجود شد و یک قطره از دریای
رحمتت سرمایه غیب و شهود با کمال عجز و نادانی و فقر و ناتوانی بدایع فضلت
را میطلبیم بحبلش متمسکیم و بذیلش متشبث ای خطا پوش خطاها را بپوش و بطراز
عفو اکبر مزین دار توئی آن توانائیکه بیک اراده ناتوانان را مطلع اقتدار و
ضعیفان را مصدر اختیار نمائی لا اله الّا انت العلیم الخبیر و انت العزیز
الحکیم و انت الغفور الرحیم و الفضال الکریم.

\end{otherlanguage}

Bār Elāhā, Karimā, Raḥimā, che kofrān az ʿebād ẓāher shode ke az amvāj-e
baḥr-e raḥmat maḥrūm mānde-and va az eshrāqāt-e anvār-e āftāb-e ẓohūr
mamnūʿ gashte-and? Elāhi, aʿmāl-e johhāl rā be-esm-e Sattārat satr
farmā. Toʾi Karimi ke dhonūb-e modhnebin bakhsheshat rā manʿ nanamūd va
bāb-e fayżat rā sadd nakard. Yek ḥarf az ketāb-e jūdat ʿellat-e vojūd
shod va yek qaṭre az daryā-ye raḥmatat sarmāye-ye ghayb va shohūd. Bā
kamāl-e ʿajz va nādāni va faqr va nātavāni badāyeʿ-e fażlat rā
miṭalabim, be-ḥablash motamassekim va be-dhaylash motashabbeth. Ey
khaṭāpūsh, khaṭāhā rā bepūsh va be-ṭerāz-e ʿafv-e akbar mozayyan dār.
Toʾi ān tavānāʾi ke be-yek erāde nātavānān rā maṭlaʿ-e eqtedār va
żaʿifān rā maṣdar-e ekhtiyār namāʾi. Lā ilāha illā anta al-ʿAlīm
al-Khabīr wa-anta al-ʿAzīz al-Ḥakīm wa-anta al-Ghafūr al-Raḥīm
wa-l-Faḍḍāl al-Karīm.

O God, O Generous One, O Merciful One! What ingratitude have Your
servants shown that they have remained deprived of the waves of the
ocean of mercy and barred from the radiant light of the Sun of
Revelation? My God, conceal the deeds of the ignorant through Your name,
the Concealer. You are the Generous One whose giving the sins of sinners
have not prevented, and whose gate of grace they have not closed. One
letter from the book of Your bounty became the cause of existence, and
one drop from the ocean of Your mercy became the substance of the unseen
and the seen. In utter helplessness, ignorance, poverty, and weakness,
we seek the wonders of Your grace; we hold fast to its cord and cling to
its hem. O Concealer of faults, conceal our faults and adorn us with the
vesture of Your most great forgiveness. You are that Mighty One who,
through a single act of will, can make the powerless the dawning places
of power and the weak the sources of authority. There is no God but You,
the All-Knowing, the All-Aware. You are the Almighty, the All-Wise, and
You are the Ever-Forgiving, the Merciful, the Bountiful, the Generous.

\subsection{131}\label{section-2106}

\begin{otherlanguage}{persian}

در ذبیح دیگر ارض صاد حضرت کاظم تفکر نمائید آن ذبیحی که انفاقش آفاق را
منورنمود و اشتعالش ابداع را حرارت بخشود جان را در سبیل دوست نثار نمود
بشأنیکه ملأ اعلی از استقامت و انقطاع و انفاقش متحیر از او گذشته ذبیح
جدید عالم را متحیر نمود و ذرات کائنات بر اقبال و تسلیم و رضا و وفا و
استقامتش شهادت داده

\end{otherlanguage}

Dar dhabīḥ-e digar-e arḍ-e Ṣād, Ḥażrat-e Kāẓem, tafakkor namāyid. Ān
dhabīḥi ke enfāqash āfāq rā monavvar namūd va eshteʿālash ebdāʿ rā
ḥarārat bakhshūd. Jān rā dar sabil-e Dūst nethār namūd be-shaʾni ke
Malaʾ-e Aʿlā az esteqāmat va enqeṭāʿ va enfāqash motaḥayyer. Az ū
godhashte, dhabīḥ-e jadid ʿālam rā motaḥayyer namūd va dharrāt-e kāʾenāt
bar eqbāl va taslim va reżā va vafā va esteqāmatash shahādat dāde.

Reflect on another sacrificial victim of the land of Ṣād, His Honor
Kāẓim---that sacrificial victim whose self-sacrifice illumined the
horizons and whose ardor imparted warmth to creation. He offered his
life in the path of the Friend in such a manner that the Concourse on
high was astonished at his steadfastness, detachment, and
self-sacrifice. Apart from him, a new sacrificial victim astonished the
world, and the atoms of all created things have borne witness to his
devotion, submission, contentment, fidelity, and steadfastness.

\subsection{132}\label{section-2107}

\begin{otherlanguage}{persian}

آن حضرت که در ارض میم مسکن داشت و بملا علیجان موسوم علم توحید در آن ارض
بر افراخت و بحرارت محبة اللّه مشتعل بشأنیکه جمعی را مشتعل نمود تا آنکه
اهل شقاق و نفاق میثاق الهی شکستند و قصد آنمظلوم نمودند آنچه را مالک بود
بتاراج رفت و بعد آن وجود مقدس را با جمعی از منتسبین با سلاسل و قیود بارض
طا آوردند و نزد علی کندیش فرستادند و بهزار حیله و مکر اراده نمودند عهد
الهی را بشکند و میثاقش را محو نماید ولکن آن فارس مضمار حقیقت و برهان از
ما سوی اللّه چشم پوشیده جز ذکر دوست و وصال دوست و قرب دوست یاری و مطلبی و
ذکری نداشت باری معرضین و مشرکین بمراد نرسیدند یعنی کلمه انکار از آن معدن
اقرار اصغا نشد تا آنکه بمراد خود فائز گشت و

\end{otherlanguage}

Ān Ḥażrat ke dar arḍ-e Mim maskan dāsht va be-Mollā ʿAlijān mawsūm,
ʿalam-e towḥid dar ān arḍ bar afrākht va be-ḥarārat-e maḥabbat Allāh
moshtaʿel be-shaʾni ke jamʿi rā moshtaʿel namūd, tā ānke ahl-e sheqāq va
nefāq mithāq-e elāhi shekastand va qaṣd-e ān maẓlūm namūdand. Ānche rā
mālek būd be-tārāj raft va baʿd ān vojūd-e moqaddas rā bā jamʿi az
montasebin bā salāsel va qoyūd be-arḍ-e Ṭā āvardand va nazd-e ʿAli
Kandish ferestādand va be-hezār ḥile va makr erāde namūdand ʿahd-e elāhi
rā beshkanad va mithāqash rā maḥv namāyad. Va-lāken ān fāres-e meżmār-e
ḥaqiqat va borhān az mā sivā Allāh chashm pūshide, joz dhekr-e Dūst va
veṣāl-e Dūst va qorb-e Dūst yāri va maṭlabi va dhekri nadāsht. Bāri,
moʿreżin va moshrekin be-morād naresidand, yaʿni kalame-ye enkār az ān
maʿdan-e eqrār eṣghā nashod, tā ānke be-morād-e khod fāʾez gasht va

That honored man, who lived in the land of Mīm and was known as Mullā
ʿAlījān, raised the banner of divine unity in that land. He was so
aflame with the heat of the love of God that he set a number of others
ablaze, until the people of discord and hypocrisy broke the divine
Covenant and plotted against that wronged one. Everything he possessed
was plundered. Then they brought that holy being, together with a number
of his associates, to the land of Ṭā in chains and fetters and sent them
to ʿAlī Kandish. Through a thousand tricks and schemes, they sought to
make him break God's Covenant and efface His pledge. Yet that horseman
in the arena of truth and proof had turned his eyes from all but God and
had no companion, desire, or remembrance except the remembrance of the
Friend, reunion with the Friend, and nearness to the Friend. In short,
those who turned away and associated partners with God did not achieve
their aim: no word of denial was heard from that mine of acknowledgment,
until he attained his own desire and

\subsection{133}\label{section-2108}

\begin{otherlanguage}{arabic}

بعد از صعود آن روح مطهر و عروج آن جوهر بشر از قلم اعلی روح من فی ملکوت
الامر و الخلق فداه دو زیارت از برای او از سماء مواهب و الطاف نازل یکی از
آن دو ذکر میشود که شاید عرف عدل و انصاف متضوع گردد و مثل قرون و اعصار
قبل باوهام حرکت ننمایند و بظنون تمسک نجویند انه یفعل ما یشا ویحکم ما
یرید و هو العزیز الحمید.

\end{otherlanguage}

Baʿd az ṣoʿūd-e ān rūḥ-e moṭahhar va ʿorūj-e ān jowhar-e bashar, az
Qalam-e Aʿlā---rūḥu man fī malakūt al-amr wa-l-khalq fidāhu---do ziyārat
az barā-ye ū az samā-ye mavāheb va alṭāf nāzel. Yeki az ān do dhekr
mishavad ke shāyad ʿarf-e ʿadl va enṣāf motaḍavveʿ gardad va methl-e
qorūn va aʿṣār-e qabl be-avhām ḥarakat nanamāyand va be-ẓonūn tamassok
najūyand. Innahu yafʿalu mā yashāʾ wa-yaḥkumu mā yurīd wa-huwa al-ʿAzīz
al-Ḥamīd.

After the ascent of that sanctified spirit and the ascension of that
essence of humanity, two Tablets of Visitation for him descended from
the heaven of bounty and grace through the Most Exalted Pen---may the
souls of all in the kingdoms of command and creation be sacrificed for
Him. One of those two is recounted here, that the fragrance of justice
and fairness may perhaps spread, and that people may not act on
illusions and cling to conjectures as in former centuries and ages. He
does whatever He wills and decrees whatever He desires, and He is the
Almighty, the Praiseworthy.

\subsection{134}\label{section-2109}

\begin{otherlanguage}{arabic}

هو المعزی من افقه الاعلی

\end{otherlanguage}

Huwa al-Muʿazzī min ufuqihi al-aʿlā.

He is the Comforter from His most exalted horizon.

\subsection{135}\label{section-2110}

\begin{otherlanguage}{persian}

اول نور ظهر و اشرق و لاح و برق من فجر المعانی علیک یا حفیف سدرة المنتهی
فی ملکوت الاسماء و آیة الظهور فی ناسوت الانشاء اشهد یا مولی العباد و نور
البلاد بانک کنت مطلع الاتفاق فی الافاق و باستقامتک ارتعدت فرائص اهل
النفاق الذین نقضوا المیثاق و کفروا باللّه مالک التلاق

\end{otherlanguage}

Awwalu nūrin ẓahara wa-ashraqa wa-lāḥa wa-baraqa min fajr al-maʿānī
ʿalayka yā ḥafīfa Sidrat al-Muntahā fī malakūt al-asmāʾ wa-āyata
al-ẓuhūr fī nāsūt al-inshāʾ. Ashhadu yā mawlā al-ʿibād wa-nūra al-bilād
bi-annaka kunta maṭliʿa al-ittifāq fī al-āfāq wa-bi-istiqāmatika
irtaʿadat farāʾiṣu ahl al-nifāq alladhīna naqaḍū al-mīthāq wa-kafarū
bi-llāh Mālik al-Talāq.

May the first light that appeared, shone, gleamed, and flashed from the
dawn of inner meanings rest upon you, O rustling of the Tree beyond
which there is no passing in the kingdom of names, and sign of
Revelation in the mortal realm of creation. I bear witness, O master of
the servants and light of the lands, that you were the dawning place of
concord throughout the horizons, and that through your steadfastness the
limbs of the people of hypocrisy trembled---those who broke the Covenant
and disbelieved in God, the Lord of the Day of Meeting.

\subsection{136}\label{section-2111}

\begin{otherlanguage}{arabic}

بک نصبت رایة الحق بین الخلق و ظهرت آیاته و انتشرت احکامه و برزت بیناته
طوبی لمن اقبل الیک و آنس بک و زار قبرک و وجد منک عرف قمیص ربّک و طاف حولک
و فاز بلقائک و نطق بذکرک

\end{otherlanguage}

Bika nuṣibat rāyat al-ḥaqq bayna al-khalq wa-ẓaharat āyātuhu
wa-intasharat aḥkāmuhu wa-barazat bayyinātuhu. Ṭūbā li-man aqbala ilayka
wa-ānasa bika wa-zāra qabraka wa-wajada minka ʿarfa qamīṣi rabbika
wa-ṭāfa ḥawlaka wa-fāza bi-liqāʾika wa-naṭaqa bi-dhikrika.

Through you the banner of truth was raised among creation, His verses
were manifested, His laws were spread abroad, and His clear proofs
appeared. Blessed is whoever turned toward you, found companionship with
you, visited your grave, perceived from you the fragrance of your Lord's
garment, circled around you, attained your presence, and spoke in
remembrance of you.

\subsection{137}\label{section-2112}

\begin{otherlanguage}{persian}

بک ارتفع علم التوحید علی العلم فی العالم و اشرق نور القدم بین الامم
بخضوعک خضعت الاعناق و بتوجهک توجهت الوجوه الی وجه اللّه رب الارباب انت
الذی کنت مترجما لوحی اللّه و مبینا لایاته و حاملا اماناته و مهبطا لاسراره
و مخزنا للئالی علمه و مشرقا لظهور انواره و منبعا لفرات رحمته و مطلعا
لآثاره و عرشا لاستواء هیکل احکامه

\end{otherlanguage}

Bika irtafaʿa ʿalam al-tawḥīd ʿalā al-ʿalam fī al-ʿālam wa-ashraqa nūr
al-qidam bayna al-umam. Bi-khuḍūʿika khaḍaʿat al-aʿnāq wa-bi-tawajjuhika
tawajjahat al-wujūh ilā wajh Allāh Rabb al-arbāb. Anta alladhī kunta
mutarjiman li-waḥy Allāh wa-mubayyinan li-āyātihi wa-ḥāmilan amānātihi
wa-mahbiṭan li-asrārihi wa-makhzanan li-laʾāliʾi ʿilmihi wa-mashriqan
li-ẓuhūr anwārihi wa-manbaʿan li-furāti raḥmatihi wa-maṭliʿan
li-āthārihi wa-ʿarshan li-istiwāʾ haykal aḥkāmihi.

Through you the banner of divine unity rose above the standard in the
world, and the light of eternity shone among the nations. Through your
humility, necks bowed, and through your turning toward Him, faces turned
toward the face of God, the Lord of lords. You were an interpreter of
God's revelation, an expounder of His verses, a bearer of His trusts, a
place where His mysteries descended, a treasury for the pearls of His
knowledge, a dawning place for the appearance of His lights, a source of
the sweet waters of His mercy, a dayspring of His signs, and a throne on
which the temple of His laws was established.

\subsection{138}\label{section-2113}

\begin{otherlanguage}{arabic}

اشهد ان بقیامک قام المخلصون و بندائک انتبه الراقدون و باقبالک اقبل
المقربون و سرع الموحدون بک ظهر الکتاب و فصل الخطاب و انت المنادی باسم
ربّک فی المآب اشهد فی ایامک فتح باب اللقاء علی من فی الارض و السماء و
هطلت من سحاب الکرم امطارالفضل و العطاء

\end{otherlanguage}

Ashhadu anna bi-qiyāmika qāma al-mukhliṣūn wa-bi-nidāʾika intabaha
al-rāqidūn wa-bi-iqbālika aqbala al-muqarrabūn wa-saraʿa al-muwaḥḥidūn.
Bika ẓahara al-kitāb wa-faṣl al-khiṭāb wa-anta al-munādī bi-ismi rabbika
fī al-maʾāb. Ashhadu fī ayyāmika futiḥa bāb al-liqāʾ ʿalā man fī al-arḍ
wa-l-samāʾ wa-haṭalat min saḥāb al-karam amṭār al-faḍl wa-l-ʿaṭāʾ.

I bear witness that through your arising the sincere arose, through your
call the sleeping awoke, and through your turning toward Him those
brought near turned toward Him and the upholders of His unity hastened.
Through you the Book and the decisive utterance were manifested, and you
are the one who calls out in the name of your Lord at the place of
return. I bear witness that in your days the gate of reunion was opened
to all in earth and heaven, and the rains of grace and bounty poured
down from the cloud of generosity.

\subsection{139}\label{section-2114}

\begin{otherlanguage}{arabic}

طوبی لک یا بهجة الاصفیاء و مهجة الاولیاء انت الذی زین اللّه ظاهر الارض
بدمک و باطنها بهیکلک طوبی لک یا رکن الایمان و مطلع العرفان انت الذی خضعت
الاذکار عند ذکرک و الاسماء عند ظهور اسمک

\end{otherlanguage}

Ṭūbā laka yā bahjat al-aṣfiyāʾ wa-muhjat al-awliyāʾ. Anta alladhī
zayyana Allāh ẓāhir al-arḍ bi-damika wa-bāṭinahā bi-haykalika. Ṭūbā laka
yā rukn al-īmān wa-maṭliʿ al-ʿirfān. Anta alladhī khaḍaʿat al-adhkār
ʿinda dhikrika wa-l-asmāʾ ʿinda ẓuhūr ismika.

Blessed are you, O joy of the chosen ones and heart of the holy ones.
You are the one through whose blood God adorned the surface of the
earth, and through whose body He adorned its depths. Blessed are you, O
pillar of faith and dawning place of recognition. You are the one before
whose remembrance all remembrances bowed, and before the appearance of
whose name all names humbled themselves.

\subsection{140}\label{section-2115}

\begin{otherlanguage}{arabic}

فآه آه یا مقصد الاقصی و اسم اللّه مالک العرش و الثری بما ورد علیک من
الامراء و العلماء تاللّه بمصیبتک ناحت کینونة السرور امام وجه الظهور و
هدرت حمامة الاحزان علی الاغصان و بما ورد علیک صعدت زفرات المقربین و نزلت
عبرات الموحدین

\end{otherlanguage}

Fa-āh āh yā maqṣad al-aqṣā wa-ism Allāh Mālik al-ʿarsh wa-l-tharā bi-mā
warada ʿalayka min al-umarāʾ wa-l-ʿulamāʾ. Ta-llāh bi-muṣībatika nāḥat
kaynūnat al-surūr amāma wajh al-ẓuhūr wa-hadarat ḥamāmat al-aḥzān ʿalā
al-aghṣān wa-bi-mā warada ʿalayka ṣaʿidat zafarāt al-muqarrabīn
wa-nazalat ʿabarāt al-muwaḥḥidīn.

Alas, alas, O most distant goal and name of God, the Lord of the throne
and the earth, for what befell you at the hands of rulers and religious
scholars! By God, at your affliction the very essence of joy lamented
before the face of Revelation, and the dove of sorrows moaned upon the
branches. Because of what befell you, the sighs of those brought near
ascended and the tears of the upholders of His unity flowed.

\subsection{141}\label{section-2116}

\begin{otherlanguage}{arabic}

تری و تعلم یا اله الغیب و سلطان الشهود ان الارض من دماء اصفیائک و
اولیائک و لم تکن لهم من امّ لتنوح علیهم و لا من مونس لیعزی ابنائهم و
نسائهم قد ورد علیهم فی ارض الطاء ما ناح به الملأ الاعلی

\end{otherlanguage}

Tarā wa-taʿlamu yā Ilāh al-ghayb wa-Sulṭān al-shuhūd anna al-arḍ min
dimāʾ aṣfiyāʾika wa-awliyāʾika wa-lam takun lahum min ummin li-tanūḥa
ʿalayhim wa-lā min muʾnisin li-yuʿazziya abnāʾahum wa-nisāʾahum. Qad
warada ʿalayhim fī arḍ al-Ṭāʾ mā nāḥa bihi al-Malaʾ al-Aʿlā.

You see and know, O God of the unseen and Sovereign of the visible, that
the earth is awash with the blood of Your chosen ones and Your holy
ones, and that they had no mother to lament them, nor any companion to
comfort their children and their wives. There befell them in the land of
Ṭā that which caused the Concourse on high to lament.

\subsection{142}\label{section-2117}

\begin{otherlanguage}{arabic}

قد رأیت یا الهی عبدک علیا مطروحا علی التراب الذی کان منجذبا بنفحات وحیک
و متحرکا بارادتک و مشتعلا بنار حبک علی شأن سمع کل ذی سمع من کل عرق من
عروقه آیات محبتک و اسرار ولایتک قد جاهد فی سبیلک حق الجهاد و قام علی
خدمتک حق القیام الی ان صار محمرا بالدم امام وجهه العالم و انفق روحه شوقا
للقائک و وصالک یا مالک القدم

\end{otherlanguage}

Qad raʾayta yā Ilāhī ʿabdaka ʿAliyyan maṭrūḥan ʿalā al-turāb alladhī
kāna munjadhiban bi-nafaḥāt waḥyika wa-mutaḥarrikan bi-irādatika
wa-mushtaʿilan bi-nār ḥubbika ʿalā shaʾnin samiʿa kullu dhī samʿin min
kulli ʿirqin min ʿurūqihi āyāt maḥabbatika wa-asrār wilāyatika. Qad
jāhada fī sabīlika ḥaqqa al-jihād wa-qāma ʿalā khidmatika ḥaqqa al-qiyām
ilā an ṣāra muḥammaran bi-l-dam amāma wajhihi al-ʿālim wa-anfaqa rūḥahu
shawqan li-liqāʾika wa-wiṣālika yā Mālik al-qidam.

You have seen, O my God, Your servant ʿAlī cast upon the dust---he who
was enraptured by the breezes of Your revelation, moved by Your will,
and set aflame by the fire of Your love, so that everyone with hearing
heard from every vein of his body the verses of Your love and the
mysteries of Your guardianship. He truly strove in Your path and arose
in full devotion to Your service, until he became crimson with blood
before His all-knowing face and offered up his soul in longing to meet
You and attain reunion with You, O Lord of eternity.

\subsection{143}\label{section-2118}

\begin{otherlanguage}{arabic}

بذلک ناحت حوریة البهاء فی الفردوس الاعلی و اهل الجنة العلیا و خباء المجد
فی ملکوت الاسماء و لم یکن یا الهی من یحمله الی مقره المقام الذی قدرت له
من قلمک الاعلی فی الصحیفة الحمراء و حملناه امتان من امائک اللّاتی وجدن
حلاوة خدمتک فی ایامک و عرفن ما غفل عنه اکثر رجالک و عبادک

\end{otherlanguage}

Bi-dhālika nāḥat Ḥūriyyat al-Bahāʾ fī al-Firdaws al-aʿlā wa-ahl
al-jannat al-ʿulyā wa-khibāʾ al-majd fī malakūt al-asmāʾ wa-lam yakun yā
Ilāhī man yaḥmiluhu ilā maqarrihi al-maqām alladhī qaddarta lahu min
qalamika al-aʿlā fī al-ṣaḥīfat al-ḥamrāʾ wa-{[}corrupted text{]} amatān
min imāʾika allātī wajadna ḥalāwat khidmatika fī ayyāmika wa-ʿarafna mā
ghafala ʿanhu aktharu rijālika wa-ʿibādika.

At this the Maid of Glory in the highest Paradise lamented, as did the
inhabitants of the loftiest heaven and the tabernacle of glory in the
kingdom of names. There was no one, O my God, to carry him to his
resting place, the station You had ordained for him through Your Most
Exalted Pen in the Crimson Scroll, and {[}corrupted original text{]} two
of Your handmaidens who had found the sweetness of Your service in Your
days and recognized what most of Your men and servants had overlooked.

\subsection{144}\label{section-2119}

\begin{otherlanguage}{persian}

ای رب صل علیهما فی جبروتک ثم اکتب لهما ما یجد منه المخلصون عرف فضلک و
عنایتک یا الهی و اله الاسماء و فاطری و فاطر السماء اسئلک بدمه الشریف و
قیامه علی امرک و استقامته علی خدمتک و انفاق روحه فی سبیلک بان تغفر لی و
لوالدی و لمن آمن بک و اقبل الیک ثم اسئلک بان تجعلنی ثابتا علی امرک و
مستقیما علی حبک و معروفا باسمک و غنیا بغنائک لا اله الا انت المهیمن علی
ما کان و ما یکون

\end{otherlanguage}

Ay Rabb, ṣalli ʿalayhimā fī jabarūtika thumma uktub lahumā mā yajidu
minhu al-mukhliṣūn ʿarfa faḍlika wa-ʿināyatika. Yā Ilāhī wa-Ilāh
al-asmāʾ wa-Fāṭirī wa-Fāṭir al-samāʾ, asʾaluka bi-damihi al-sharīf
wa-qiyāmihi ʿalā amrika wa-istiqāmatihi ʿalā khidmatika wa-infāqi rūḥihi
fī sabīlika bi-an taghfira lī wa-li-wālidayya wa-li-man āmana bika
wa-aqbala ilayka. Thumma asʾaluka bi-an tajʿalanī thābitan ʿalā amrika
wa-mustaqīman ʿalā ḥubbika wa-maʿrūfan bi-ismika wa-ghaniyyan
bi-ghināʾika. Lā ilāha illā anta al-Muhaymin ʿalā mā kāna wa-mā yakūn.

O Lord, bless them both in Your realm of power, then ordain for them
that from which the sincere may perceive the fragrance of Your grace and
loving care. O my God and God of names, my Creator and Creator of
heaven, I ask You by his noble blood, by his arising for Your Cause, by
his steadfastness in Your service, and by his offering up his soul in
Your path, to forgive me, my parents, and whoever has believed in You
and turned toward You. Then I ask You to make me firm in Your Cause,
steadfast in Your love, known by Your name, and enriched by Your riches.
There is no God but You, the Guardian over all that has been and all
that will be.

\subsection{145}\label{section-2120}

\begin{otherlanguage}{persian}

ای آقایان لله بگوئید و لله بشنوید از قبل شهادت حضرت سید الشهداء علیه
سلام اللّه را حجت عظمی و بیّنه کبری از برای امر حضرت خاتم روح ما سواه فداه
ذکر مینمودند چنانچه عرض شد حال این نفوس مشتعله مطمئنه راضیه مرضیه که
باجنحه انقطاع بمقر فدا طیران نمودند و جان رایگان در سبیل محبوب امکان
انفاق کردند بنظر غافلین و مشرکین نیامده قد حجبتهم اهوائهم و غشتهم
اوهامهم لعمر اللّه شهادت هر یک از نفوس مذکوره دلیلی است اعظم و برهانی است
اکمل و اتم ولکن حق جل جلاله جز بما نزل من عنده و ظهر من لدنه در اثبات
امر مبارکش بامری نطق نفرموده

\end{otherlanguage}

Ey āqāyān, li-llāh begūʾīd va li-llāh beshnavīd. Az qabl shahādat-e
Ḥaẓrat-e Sayyid al-Shuhadāʾ, ʿalayhi salāmuʾllāh, rā ḥojjat-e ʿoẓmā va
bayyene-ye kobrā az barā-ye amr-e Ḥaẓrat-e Khātam, rūḥu mā siwāhu fidāh,
dhekr mīnamūdand, chenānche ʿarẓ shod. Ḥāl īn nofūs-e moshtaʿele-ye
moṭmaʾenne-ye rāẓiye-ye marẓiyye ke bā ajneḥe-ye enqeṭāʿ be-maqarr-e
fedā ṭayarān namūdand va jān rāyegān dar sabīl-e Maḥbūb-e emkān enfāq
kardand, be-naẓar-e ghāfelīn va moshrekīn nayāmade. Qad ḥajabatʹhum
ahwāʾuhum wa-ghashiyatʹhum awhāmuhum. La-ʿamruʾllāh, shahādat-e har yek
az nofūs-e madhkūre dalīlī ast aʿẓam va borhānī ast akmal va atamm,
va-lāken Ḥaqq, jalla jalāluhu, joz bi-mā nazala min ʿindihi wa-ẓahara
min ladunhu dar ethbāt-e amr-e mobārakash be-amrī noṭq nafarmūde.

Gentlemen, speak for God's sake and listen for God's sake. Formerly, as
has been stated, they cited the martyrdom of His Holiness the Lord of
Martyrs---may God's peace be upon him---as the greatest proof and
mightiest testimony to the Cause of His Holiness the Seal---may the
souls of all others be sacrificed for Him. Yet these fervent, assured
souls, content with God and pleasing to Him, who soared on the wings of
detachment to the realm of sacrifice and freely gave their lives in the
path of the Beloved of all existence, have escaped the notice of the
heedless and those who associate partners with God. Their desires have
veiled them, and their delusions have enveloped them. By the life of
God, the martyrdom of each of these souls is a greater proof and a more
perfect and complete demonstration. Yet God, exalted be His glory, has
spoken of nothing in proof of His blessed Cause except what has been
sent down from Him and has appeared from His presence.

\subsection{146}\label{section-2121}

\begin{otherlanguage}{persian}

یا اولی النهی در آنچه در باب ثامن از واحد سادس از قلم نقطه اولی روح ما
سواه فداه جاری شده تفکر نمائید قوله تعالی و امر شده که در هر نوزده روز
یک دفعه در این باب نظر کنند لعل در ظهور من یظهره اللّه محتجب نشوند بشئونی
دون شأن آیات که اعظم حجج و براهین بوده و هست

\end{otherlanguage}

Yā ulī al-nuhā, dar ānche dar bāb-e thāmen az vāḥed-e sādes az qalam-e
Noqṭe-ye Ūlā, rūḥu mā siwāhu fidāh, jārī shode tafakkor namāyīd. Qawluhu
taʿālā: Va amr shode ke dar har nūzdah rūz yek dafʿe dar īn bāb naẓar
konand, laʿall dar ẓohūr-e Man Yuẓhiruhuʾllāh moḥtajeb nashavand
be-shoʾūnī dūn-e shaʾn-e āyāt ke aʿẓam-e ḥojaj va barāhīn būde va hast.

O people of understanding, reflect on what flowed from the pen of the
Primal Point---may the souls of all others be sacrificed for Him---in
the eighth chapter of the sixth Váḥid. His words, exalted be He: ``It
has been commanded that once every nineteen days they consider this
chapter, so that, at the appearance of Him Whom God will make manifest,
they may not be veiled by matters other than the verses, which have been
and remain the greatest of proofs and demonstrations.''

\subsection{147}\label{section-2122}

\begin{otherlanguage}{persian}

در این باب میفرماید قوله تبارک و تعالی بعد از غروب شمس حقیقت امتناع دارد
که از غیر او آیه ظاهر شود بر نهج فطرت و قدرت بدون تعلم و شئونیکه نزد اهل
علم متصور است با وجود این امتناع که غیر از من یظهره اللّه کسی نتواند مدعی
شد این امر را فرض شده در بیان که اگر نفسی ادعا کند و آیاتی از او ظاهر
گردد احدی متعرض نگردد او را لعل بر آن شمس حقیقت حزنی وارد نیاید الی قوله
عز و جل

\end{otherlanguage}

Dar īn bāb mīfarmāyad, qawluhu tabāraka wa-taʿālā: Baʿd az ghorūb-e
Shams-e Ḥaqīqat emtenāʿ dārad ke az gheyr-e ū āye ẓāher shavad bar
nahj-e feṭrat va qodrat bedūn-e taʿallom va shoʾūnī ke nazd-e ahl-e ʿelm
motaṣavvar ast. Bā vojūd-e īn emtenāʿ ke gheyr az Man Yuẓhiruhuʾllāh
kasī natavānad moddaʿī shod īn amr rā, farẓ shode dar Bayān ke agar
nafsī eddeʿā konad va āyātī az ū ẓāher gardad, aḥadī motaʿarreẓ nagardad
ū rā, laʿall bar ān Shams-e Ḥaqīqat ḥoznī vāred nayāyad. Ilā qawlihi
ʿazza wa-jall.

In this chapter He says---His words, blessed and exalted be He: ``After
the setting of the Sun of Truth, it is impossible for a verse to appear
from anyone other than Him, through innate capacity and power, without
learning or the means that people of learning can conceive. Although it
is impossible for anyone other than Him Whom God will make manifest to
claim this, it has been enjoined in the Bayán that if anyone makes such
a claim and verses appear from him, no one should oppose him, lest any
sorrow befall that Sun of Truth.'' And further, in His words, mighty and
glorious is He:

\subsection{148}\label{section-2123}

\begin{otherlanguage}{persian}

اگر شنوید چنین امری و یقین نکنید تکسب امری ننموده که سبب حزن او باشد اگر
چه در واقع غیر او باشد اگر چه این تصوریست محال ولی همینقدر که ذکر اسم او
کرد بر صاحبان حب او بعید است که او را محزون کنند احتراما لاسمه زیرا که
امر از دو شق بیرون نیست یا اوست و حال آنکه غیر او ممکن نیست که آیات بر
نهج فطرت نازل فرماید که چرا نفسی تکذیب حق کرده باشد و حال آنکه شب وروز
در انتظار او عمل کرده باشد و اگر بر فرض امتناع کسی خود را نسبت داد
واگذارند حکم او را با خدا بر خلق نیست که حکم او نمایند اجلالا لاسم
محبوبهم و حال آنکه چنین نفسی نیست که تواند چنین مقامی را ادعا کند اگر در
کور قرآن بهمرسید در این کور هم خواهد رسید

\end{otherlanguage}

Agar beshnavīd chenīn amrī va yaqīn nakonīd, takassob-e amrī nanamūde ke
sabab-e ḥozn-e ū bāshad, agar che dar vāqeʿ gheyr-e ū bāshad. Agar che
īn taṣavvorīst moḥāl, valī hamīn-qadr ke dhekr-e esm-e ū kard, bar
ṣāḥebān-e ḥobb-e ū baʿīd ast ke ū rā maḥzūn konand, iḥtirāman li-smihi.
Zīrā ke amr az do sheqq bīrūn nīst: yā ūst, va ḥāl ānke gheyr-e ū momken
nīst ke āyāt bar nahj-e feṭrat nāzel farmāyad, ke cherā nafsī takdhīb-e
Ḥaqq karde bāshad va ḥāl ānke shab va rūz dar enteẓār-e ū ʿamal karde
bāshad; va agar bar farẓ-e emtenāʿ kasī khod rā nesbat dād, vāgodār?
vāgo dhārand ḥokm-e ū rā bā Khodā. Bar khalq nīst ke ḥokm-e ū namāyand,
ijlālan li-smi maḥbūbihim. Va ḥāl ānke chenīn nafsī nīst ke tavānad
chenīn maqāmī rā eddeʿā konad. Agar dar kowr-e Qorʾān be-ham-resīd, dar
īn kowr ham khāhad resīd.

``If you hear of such a claim and are not convinced, do nothing that
would cause him sorrow, even if in reality he is someone other than Him.
Although this is an impossible supposition, the mere fact that he has
mentioned His name makes it unworthy of those who love Him to cause him
sorrow, out of respect for His name. For the matter admits of only two
possibilities: either it is He---and it is impossible for anyone other
than Him to reveal verses through innate capacity---so why should anyone
deny the Truth when night and day he has acted in expectation of Him?
Or, supposing the impossible, someone has attributed himself to Him;
then let them leave his judgment to God. It is not for people to judge
him, out of reverence for the name of their Beloved. Yet there is no
such soul who can claim such a station. If such a person appeared in the
cycle of the Qurʼán, one will appear in this cycle as well.

\subsection{149}\label{section-2124}

\begin{otherlanguage}{persian}

آیات او بنفسه دلیل است بر ضیاء شمس وجود او و عجز کل دلیل است بر فقر و
احتیاج بسوی اوسبب این امر اینست لعل در یوم ظهور حق قدمهای ایشان بر صراط
نلغزد و یا به شبحی که در افئده ایشان است بر مکون وجود خود بآن آیه حکمی
نکنند که یک دفعه کل کینونیت و اعمال آنها باطل گردد و خود خبر نشوند

\end{otherlanguage}

Āyāt-e ū be-nafsehe dalīl ast bar ẓiyā-ye shams-e vojūd-e ū va ʿajz-e
koll dalīl ast bar faqr va eḥtiyāj be-sū-ye ū. Sabab-e īn amr īnast:
laʿall dar yowm-e ẓohūr-e Ḥaqq qadam-hā-ye īshān bar ṣerāṭ nalaghzad va
yā be-shabaḥī ke dar afʾede-ye īshān ast bar Mokavven-e vojūd-e khod
be-ān āye ḥokmī nakonand ke yek-dafʿe koll-e kaynūniyyat va aʿmāl-e ānhā
bāṭel gardad va khod khabar nashavand.

``His verses are in themselves proof of the radiance of the sun of His
being, and the powerlessness of all is proof of their poverty and need
for Him. The reason for this command is that, on the day of the
appearance of the Truth, their feet may not slip on the path, nor may
they, because of a phantom in their hearts, use that verse to pass a
judgment against the Creator of their own existence that would at once
nullify their entire being and all their deeds without their even
knowing it.

\subsection{150}\label{section-2125}

\begin{otherlanguage}{persian}

اگر کل بر این یک امر اطاعت کنند بر خداوند است که حق را بر ایشان ظاهر
فرماید و القای حجت و دلیل در قلوب ایشان فرماید بادله ظاهره که از قبل او
مشرق میگردد تا اینکه توانند در حق او اظهار یقین نمود و دون آنها توانند
اظهار وقوف نمود که وقوف در حق دون حق است و کافی است کل اهل بیان را اگر
در این حکم عمل نمایند در نجات ایشان در یوم قیامت زیرا که آن یومی است که
کل ما علی الارض مشرقه در افئده ایشان از شمس حقیقت انا للّه عاملون گویند
ولی صادقین عبادی هستند که نظر بمکون آیات باللّه عزّ و جل نموده که آن شمس
حقیقت باشد که آیات او که ضیاء اوست فاصل است ما بین کل شئ تلک حجة اللّه قد
تمت علیکم ان یا عباد اللّه فاتقون.

\end{otherlanguage}

Agar koll bar īn yek amr eṭāʿat konand, bar Khodāvand ast ke Ḥaqq rā bar
īshān ẓāher farmāyad va elqā-ye ḥojjat va dalīl dar qolūb-e īshān
farmāyad be-adelle-ye ẓāhere ke az qabl-e ū moshreq mīgardad, tā īnke
tavānand dar ḥaqq-e ū eẓhār-e yaqīn namūd va dūn-e ānhā tavānand eẓhār-e
voqūf namūd, ke voqūf dar ḥaqq-e dūn-e Ḥaqq ast. Va kāfī ast koll-e
ahl-e Bayān rā, agar dar īn ḥokm ʿamal namāyand, dar najāt-e īshān dar
yowm-e qiyāmat, zīrā ke ān yowmī ast ke koll-e mā ʿalā al-arẓ, moshreqe
dar afʾede-ye īshān az Shams-e Ḥaqīqat, innā li-llāhi ʿāmilūn gūyand.
Valī ṣādeqīn ʿebādī hastand ke naẓar be-Mokavven-e āyāt bi-llāhi ʿazza
wa-jall namūde, ke ān Shams-e Ḥaqīqat bāshad ke āyāt-e ū ke ẓiyā-ye ūst
fāṣel ast mā bayn-e koll-e shayʾ. Tilka ḥujjatuʾllāhi qad tammat
ʿalaykum, an yā ʿibādaʾllāhi fa-ttaqūn.

``If all obey this one command, it is for God to make the Truth manifest
to them and to instill proof and evidence in their hearts through the
manifest proofs that shine forth from Him, so that they may express
certainty concerning Him, and concerning others may express suspension
of judgment, for suspension of judgment applies to those other than the
Truth. If all the people of the Bayán act according to this command, it
will suffice for their salvation on the Day of Resurrection. For that is
the day when all on earth, through what shines in their hearts from the
Sun of Truth, will say, `We act for God.' But the truthful are those
servants who look, through God, mighty and glorious is He, to the
Creator of the verses, who is that Sun of Truth whose verses, which are
His radiance, distinguish between all things. Such is God's proof, now
made complete for you. O servants of God, then fear Me.''

\subsection{151}\label{section-2126}

\begin{otherlanguage}{persian}

ای صاحبان ابصار در این بیان رحمن که عرفش عالم را احاطه نموده و مقصود از
او ظاهر و باهر و آشکار نظر نمائید فی الحقیقه این باب سپریست محکم و درعی
است متین از برای حفظ مقصود عالمیان ای اهل بیان حیا کنید و قدری تفکر
نمائید آیاتش عالم را احاطه نموده جمیع ذرات کائنات بر عظمتش شهادت داده و
میدهند از حق جل جلاله بترسید وتدبیرات و بیانات حضرت مبشر را باغراض
نفسانیه مهجور منمائید

\end{otherlanguage}

Ey ṣāḥebān-e abṣār, dar īn bayān-e Raḥmān ke ʿarfash ʿālam rā eḥāṭe
namūde va maqṣūd az ū ẓāher va bāher va āshkār, naẓar namāyīd. Fī
al-ḥaqīqe īn bāb separīst moḥkam va derʿī ast matīn az barā-ye ḥefẓ-e
Maqṣūd-e ʿālamiyān. Ey ahl-e Bayān, ḥayā konīd va qadrī tafakkor
namāyīd. Āyātash ʿālam rā eḥāṭe namūde, jamīʿ-e dharrāt-e kāʾenāt bar
ʿaẓamatash shahādat dāde va mīdahand. Az Ḥaqq, jalla jalāluhu, betarsīd
va tadbīrāt va bayānāt-e Ḥaẓrat-e Mobashsher rā be-aghrāẓ-e nafsāniyye
mahjūr manamāyīd.

O people of insight, consider this utterance of the All-Merciful, whose
fragrance has encompassed the world and whose intent is manifest,
evident, and clear. This chapter is indeed a strong shield and sturdy
armor for the protection of the Desire of all peoples. O people of the
Bayán, have some shame and reflect a little. His verses have encompassed
the world; every atom of creation has testified and continues to testify
to His greatness. Fear God, exalted be His glory, and do not set aside
the provisions and utterances of His Holiness the Herald for selfish
aims.

\subsection{152}\label{section-2127}

\begin{otherlanguage}{persian}

حضرت مقصود روح من فی ملکوت الامر و الخلق فداه از اعانت وتصدیق و ایمان
شما گذشت این را فی الحقیقه ذکر مینمایم چه که از لسان مبارک استماع شده
بهر بابی اراده نمودید توجه کنید و از هر کأسی میخواهید بیاشامید و بهر جهت
و هر سو میل کنید بروید ولکن این مظلوم را با این آیات باهرات و بینات
واضحات و قیامش مقابل اعداء بخود او واگذارید از حیل و مکر و خدعه و فریب
بگذرید

\end{otherlanguage}

Ḥaẓrat-e Maqṣūd, rūḥu man fī malakūti al-amri wa-l-khalqi fidāh, az
eʿānat va taṣdīq va īmān-e shomā godhasht. Īn rā fī al-ḥaqīqe dhekr
mīnamāyam, che ke az lesān-e mobārak estemāʿ shode: Be-har bābī erāde
namūdīd tavajjoh konīd va az har kaʾsī mīkhāhīd biyāshāmīd va be-har
jahat va har sū meyl konīd beravīd, va-lāken īn Maẓlūm rā bā īn āyāt-e
bāherāt va bayyenāt-e vāẓeḥāt va qiyāmash moqābel-e aʿdāʾ be-khod-e ū
vāgodārīd. Az ḥiyal va makr va khodʿe va farīb begodharīd.

The Desired One---may the souls of all in the kingdoms of command and
creation be sacrificed for Him---has relinquished your assistance,
acknowledgment, and faith. I say this in truth, for it has been heard
from His blessed lips: ``Turn to whatever door you wish, drink from
whatever cup you desire, and go in whatever direction and toward
whatever place you choose. But leave this Wronged One to Himself, with
these resplendent verses and clear proofs, standing against His enemies.
Abandon your stratagems, scheming, deceit, and deception.''

\subsection{153}\label{section-2128}

\begin{otherlanguage}{arabic}

این فانی عرض مینماید در اراده‌های خود بکمال اطمینان مشغول شوید ابدا تعرض
نمیفرمایند لو تتخذوا لانفسکم فی کل حین ربّا من دون اللّه ولکن مرتکب نشوید
امری را که کبد نقطه اولی از او محترق شود و عیون عظمت بگرید یک سوره از
سور منزله را با بیان و جمیع کتب الهی قرائت نمایند لیظهر لکم ما غفلتم عنه
و بیقین مبین بدانید و شهادت دهید بر آنچه نقطه بیان روح من فی الامکان
فداه شهادت داده میفرماید جمیع بیان معادله بیک آیه او نمینماید هذا حق
لاریب فیه

\end{otherlanguage}

Īn fānī ʿarẓ mīnamāyad: Dar erāde-hā-ye khod be-kamāl-e eṭmīnān mashghūl
shavīd, abadan taʿarroẓ nemīfarmāyand, law tattakhidhū li-anfusikum fī
kulli ḥīnin rabban min dūniʾllāh. Va-lāken mortakeb nashavīd amrī rā ke
kabed-e Noqṭe-ye Ūlā az ū moḥtareq shavad va ʿoyūn-e ʿaẓamat begeryad.
Yek sūre az sovar-e monazzale rā bā Bayān va jamīʿ-e kotob-e elāhī
qerāʾat namāyand, li-yaẓhara lakum mā ghafaltum ʿanhu, va be-yaqīn-e
mobīn bedānīd va shahādat dahīd bar ānche Noqṭe-ye Bayān, rūḥu man fī
al-imkāni fidāh, shahādat dāde. Mīfarmāyad: Jamīʿ-e Bayān moʿādele
be-yek āye-ye ū nemīnamāyad. Hādhā ḥaqqun lā rayba fīh.

This mortal submits: Pursue your own intentions with complete assurance;
He will never interfere, even if at every moment you take for yourselves
a lord other than God. But do not commit an act that would burn the
heart of the Primal Point and make the eyes of grandeur weep. Let them
read one of the revealed surahs alongside the Bayán and all the divine
scriptures, so that what you have overlooked may become clear to you.
Know with evident certainty and testify to that to which the Point of
the Bayán---may the souls of all in existence be sacrificed for
Him---has testified. He says: ``The whole of the Bayán does not equal a
single one of His verses.'' This is the truth, beyond all doubt.

\subsection{154}\label{section-2129}

\begin{otherlanguage}{persian}

حال شخصی را که باو تمسک جسته‌اید لعمر مقصودی و مقصود من فی الارض و السماء
بر امر او آگاه نبوده و نیستید و این عبد هم ممنوع است از ذکرش عسی ان یأتی
حینه ایکاش شمس اذن از افق اعلی اشراق مینمود و این عبد عرض مینمود آنچه را
که بآن آگاه است قلم متحیر مداد متحیر عالم متحیر آیا چه شد بطنین ذباب
مشغولند و از ربّ الارباب معرض عهد اللّه را شکستند و از میثاقش چشم برداشتند

\end{otherlanguage}

Ḥāl shakhṣī rā ke be-ū tamassok josteʾīd, la-ʿamri Maqṣūdī wa-Maqṣūdi
man fī al-arḍi wa-l-samāʾ, bar amr-e ū āgāh nabūde va nīstīd, va īn ʿabd
ham mamnūʿ ast az dhekrash. ʿAsā an yaʾtiya ḥīnuhu. Ey kāsh shams-e edhn
az ofoq-e aʿlā eshrāq mīnamūd va īn ʿabd ʿarẓ mīnamūd ānche rā ke be-ān
āgāh ast. Qalam motaḥayyer, medād motaḥayyer, ʿālam motaḥayyer. Āyā che
shod? Be-ṭanīn-e dhobāb mashghūland va az Rabb al-Arbāb moʿreẓ. ʿAhd-e
Allāh rā shekastand va az mīthāqash chashm bardāshtand.

As for the person to whom you now cling---by the life of my Desired One
and the Desired One of all in earth and heaven---you have never been,
and are not now, aware of his affair, and this servant too is forbidden
to speak of it. Perhaps its time will come. If only the sun of
permission would shine from the highest horizon, and this servant could
relate what he knows. The pen is bewildered, the ink is bewildered, the
world is bewildered. What has happened? They are engrossed in the
buzzing of a fly and turn away from the Lord of Lords. They have broken
God's covenant and turned their eyes away from His pledge.

\subsection{155}\label{section-2130}

\begin{otherlanguage}{persian}

اگر نفسی تفکر نماید در خضوع و خشوع نقطه بیان عند ظهور این اسم شهادت
میدهد بر آنچه از قلم اعلی جاری شد قسم بآفتاب حقیقت بشأنی خاضع و خاشع
بودند که قلم اعلی از ذکرش عاجز است

\end{otherlanguage}

Agar nafsī tafakkor namāyad dar khoẓūʿ va khoshūʿ-e Noqṭe-ye Bayān ʿinda
ẓuhūri īn Esm, shahādat mīdahad bar ānche az Qalam-e Aʿlā jārī shod.
Qasam be-Āftāb-e Ḥaqīqat, be-shaʾnī khāẓeʿ va khāsheʿ būdand ke Qalam-e
Aʿlā az dhekrash ʿājez ast.

If anyone reflects on the humility and submissiveness of the Point of
the Bayán before the manifestation of this Name, he will testify to what
has flowed from the Most Exalted Pen. By the Sun of Truth, so great were
His humility and submissiveness that the Most Exalted Pen is powerless
to describe them.

\subsection{156}\label{section-2131}

\begin{otherlanguage}{persian}

چه نیکوست ذکر این فقره در این مقام که در این حین بخاطر فانی آمد ایامیکه
حضرت ابن زکریا ظاهر و بغسل تعمید مشغول در بریه سائر و ببشارت ناطق قائلا
توبوا لانّه قد اقترب ملکوت السموات میفرماید توبه و انابه نمائید چه که
ملکوت الهی نزدیک شده انی اعمدکم بماء للتوبه ولکن الذی یأتی بعدی هو اقوی
منّی الذی لست اهلا ان احمل حذائه بشارت میداد ناس را بظهور حضرت روح و
میفرمود او اقوی است و من قابل آنکه کفش او را حمل نمایم نیستم

\end{otherlanguage}

Che nīkūst dhekr-e īn feqre dar īn maqām ke dar īn ḥīn be-khāṭer-e fānī
āmad. Ayyāmī ke Ḥaẓrat-e Ebn-e Zakariyyā ẓāher va be-ghosl-e taʿmīd
mashghūl, dar barriyye sāʾer va be-beshārat nāṭeq, qāʾilan: Tūbū
li-annahu qad iqtaraba malakūtu al-samāwāt. Mīfarmāyad: Towbe va enābe
namāyīd, che ke malakūt-e elāhī nazdīk shode. Innī uʿammidukum bi-māʾin
li-l-tawbati wa-lākin alladhī yaʾtī baʿdī huwa aqwā minnī alladhī lastu
ahlan an aḥmila ḥidhāʾahu. Beshārat mīdād nās rā be-ẓohūr-e Ḥaẓrat-e Rūḥ
va mīfarmūd: Ū aqvā ast va man qābel-e ānke kafsh-e ū rā ḥaml namāyam
nīstam.

How fitting it is to mention here this passage, which has just come to
this mortal's mind. In the days when the son of Zechariah appeared and
was engaged in baptism, traveling in the wilderness and proclaiming the
glad tidings, he said: ``Repent, for the kingdom of heaven has drawn
near.'' He says: ``Repent and turn back to God, for the divine kingdom
has drawn near.'' ``I baptize you with water for repentance, but He who
comes after me is mightier than I, whose shoes I am not worthy to
carry.'' He gave the people glad tidings of the appearance of His
Holiness the Spirit, saying: ``He is mightier, and I am not worthy to
carry His shoes.''

\subsection{157}\label{section-2132}

\begin{otherlanguage}{arabic}

مع آنکه حضرت روح و او در یک یوم بودند و لکن چون امر حضرت روح در آن یوم
مستور بود میفرماید و الذی یأتی بعدی یعنی یظهر نفسه بعدی و بعد میفرماید
جاء یسوع من الجلیل الی الاردن الی یوحنا لیعتمد منه ولکن یوحنا منعه قائلا
انا محتاج ان اعتمد منک و انت تأتی الی فاجاب یسوع و قال له اسمع الان لان
هکذا یلیق بنا ان نکمل کل برء

\end{otherlanguage}

Maʿa ānke Ḥaẓrat-e Rūḥ va ū dar yek yowm būdand, va-lāken chon amr-e
Ḥaẓrat-e Rūḥ dar ān yowm mastūr būd, mīfarmāyad: Wa-alladhī yaʾtī baʿdī,
yaʿnī yuẓhiru nafsahu baʿdī. Va baʿd mīfarmāyad: Jāʾa Yasūʿu min
al-Jalīli ilā al-Urdunni ilā Yūḥannā li-yaʿtamida minhu, wa-lākin
Yūḥannā manaʿahu qāʾilan: Anā muḥtājun an aʿtamida minka wa-anta taʾtī
ilayya. Fa-ajāba Yasūʿu wa-qāla lahu: Ismaʿ al-āna li-anna hākadhā
yalīqu binā an nukammila kulla {[}corrupted text{]}.

Although His Holiness the Spirit and he lived in the same day, because
the Cause of His Holiness the Spirit was concealed at that time, he
says, ``And He who comes after me,'' meaning, ``He will manifest Himself
after me.'' Then he says: ``Jesus came from Galilee to the Jordan, to
John, to be baptized by him. But John prevented Him, saying, `I need to
be baptized by You, and do You come to me?' Jesus answered and said to
him, `Listen now, for thus it is fitting for us to fulfill all
{[}corrupted original text{]}.'\,''

\subsection{158}\label{section-2133}

\begin{otherlanguage}{persian}

یسوع حضرت مسیح است و یوحنا ابن زکریا و او را در کتب یوحنا المعمدانی
میگویند چه که غسل تعمید از احکام اوست و حال ما بین مسیحیین مجری و ممضی
است و بآن عاملند فی الجمله نفحهٴ این ایّام در آن ایام متضوع بوده چنانچه
حضرت نقطه روح ما سواه فداه جمیع بیان را ببشارت این ظهور تمام فرموده لعمر
اللّه اگر نفسی بعدل و انصاف کلمات بیان را یک یک استشمام نماید جز عرف محبت
این ظهور از او نیابد یشهد بذلک کل منصف و کل عادل ولکن القوم اکثرهم من
المنکرین.

\end{otherlanguage}

Yasūʿ Ḥaẓrat-e Masīḥ ast va Yūḥannā Ebn-e Zakariyyā, va ū rā dar kotob
Yūḥannā al-Maʿmadānī mīgūyand, che ke ghosl-e taʿmīd az aḥkām-e ūst va
ḥāl mā bayn-e Masīḥiyyīn mojrā va momẓā ast va be-ān ʿāmeland. Fī
al-jomle nafḥe-ye īn ayyām dar ān ayyām motaẓavveʿ būde, chenānche
Ḥaẓrat-e Noqṭe, rūḥu mā siwāhu fidāh, jamīʿ-e Bayān rā be-beshārat-e īn
ẓohūr tamām farmūde. La-ʿamruʾllāh, agar nafsī be-ʿadl va enṣāf
kalemāt-e Bayān rā yek-yek esteshmām namāyad, joz ʿarf-e maḥabbat-e īn
ẓohūr az ū nayābad. Yashhadu bi-dhālika kullu munṣifin wa-kullu ʿādilin,
wa-lākinna al-qawma aktharuhum mina al-munkirīn.

Yasúʿ is His Holiness Christ, and John is the son of Zechariah. In the
books he is called John the Baptist, for baptism is among his
ordinances; it is now practiced and upheld among Christians, and they
observe it. In short, the fragrance of these days was wafting in those
days, just as His Holiness the Point---may the souls of all others be
sacrificed for Him---filled the entire Bayán with the glad tidings of
this Revelation. By the life of God, if anyone inhales with justice and
fairness the fragrance of the words of the Bayán, one by one, he will
find in them nothing but the fragrance of love for this Revelation.
Every fair-minded and just person bears witness to this, but most of the
people are among those who deny.

\subsection{159}\label{section-2134}

\begin{otherlanguage}{arabic}

ای اهل بیان عرض این خادم را بشنوید و بکمال اطمینان متوکلا علی اللّه بر
صراط وارد شوید احتیاط امروز از اعظم سیئات محسوب و در نار بوده و خواهد
بود قد شهد بذلک قلم نقطه الاولی قوله تبارک و تعالی او را شناخته بآیات و
احتیاط در عرفان او نکرده که بقدر همان در نار محتجب خواهید بود'' انتهی

\end{otherlanguage}

Ey ahl-e Bayān, ʿarẓ-e īn khādem rā beshnavīd va be-kamāl-e eṭmīnān,
mutawakkilan ʿalā Allāh, bar ṣerāṭ vāred shavīd. Eḥtiyāṭ emrūz az
aʿẓam-e sayyeʾāt maḥsūb va dar nār būde va khāhad būd. Qad shahida
bi-dhālika qalamu Nuqṭati al-Ūlā. Qawluhu tabāraka wa-taʿālā: Ū rā
shenākhte be-āyāt va eḥtiyāṭ dar ʿerfān-e ū nakarde ke be-qadr-e hamān
dar nār moḥtajeb khāhīd būd. Intahā.

O people of the Bayán, listen to what this servant submits, and enter
the path with complete assurance, placing your trust in God. Caution
today is counted among the greatest sins and has been and will remain in
the fire. The pen of the Primal Point has testified to this. His words,
blessed and exalted be He: ``Recognize Him by His verses and exercise no
caution in recognizing Him, for to that same extent you will remain
veiled in the fire.'' End of quotation.

\subsection{160}\label{section-2135}

\begin{otherlanguage}{persian}

اصنام را بشکنید حجبات ظنون را خرق نمائید بافق اعلی وحده ناظر باشید که
شاید در یوم الهی بتوحید حقیقی فائز گردید قسم بآفتاب افق سماء فضل آنچه در
این ظهور اعظم ظاهر هر یک بمثابه شمس است در وسط زوال لا ینکرها الا کل
غافل مرتاب این کلمه را لوجه اللّه بشنوید و عمل نمائید حجت و برهانی که بآن
تصدیق نقطه بیان و رسل قبل را نموده‌اید حال اعظم و اکمل آن موجود ملاحظه
نمائید و بعدل و انصاف تکلم کنید

\end{otherlanguage}

Aṣnām rā beshkanīd, ḥojobāt-e ẓonūn rā kharq namāyīd, be ofoq-e aʿlā
vaḥdahu nāẓer bāshīd ke shāyad dar yawm-e elāhī be tawḥīd-e ḥaqīqī fāʾez
gardīd. Qasam be āftāb-e ofoq-e samā-ye fażl, ānche dar īn ẓohūr-e aʿẓam
ẓāher, har yek be mathābe-ye shams ast dar vasaṭ-e zavāl. Lā yunkiruhā
illā kullu ghāfilin murtāb. Īn kalame rā li-wajhi Allāh beshnavīd va
ʿamal namāyīd. Ḥojjat va borhānī ke be ān taṣdīq-e Noqṭe-ye Bayān va
rosol-e qabl rā namūde-īd, ḥāl aʿẓam va akmal-e ān mawjūd. Molāḥeẓe
namāyīd va be ʿadl va enṣāf takallom konīd.

Break the idols, tear apart the veils of conjecture, and look to the
Most Exalted Horizon alone, so that you may attain true recognition of
divine unity on the Day of God. I swear by the sun of the horizon of the
heaven of grace: each thing manifested in this Most Great Revelation is
like the sun at its zenith. None denies them except those who are
heedless and beset by doubt. Hear these words for the sake of God and
act upon them. The proof and testimony by which you acknowledged the
Point of the Bayán and the Messengers of the past now exist in a greater
and more perfect form. Consider this, and speak with justice and
fairness.

\subsection{161}\label{section-2136}

\begin{otherlanguage}{persian}

استغفر اللّه من ذکری و بیانی و نطقی عالم از اهل انصاف و عدل و دانش خالی
نه البته ارباب دانش در آنچه از قلم اعلی در این ظهور از عربی و فارسی نازل
شده ملاحظه نمایند و همچنین آنچه از قبل بر انبیاء و مرسلین نازل و بانصاف
تکلم کنند و بعدل سخن گویند

\end{otherlanguage}

Astaghfiru Allāha min dhikrī wa-bayānī wa-nuṭqī. ʿĀlam az ahl-e enṣāf va
ʿadl va dānesh khālī na. Albatte arbāb-e dānesh dar ānche az Qalam-e
Aʿlā dar īn ẓohūr az ʿarabī va fārsī nāzel shode molāḥeẓe namāyand va
hamchenīn ānche az qabl bar anbiyā va morsalīn nāzel, va be enṣāf
takallom konand va be ʿadl sokhan gūyand.

I ask God's forgiveness for my mention, my exposition, and my speech.
The world is not devoid of people of fairness, justice, and knowledge.
Those possessed of knowledge should certainly examine what has been sent
down in Arabic and Persian from the Most Exalted Pen in this Revelation,
as well as what was previously sent down to the Prophets and Messengers,
and speak with fairness and express themselves with justice.

\subsection{162}\label{section-2137}

\begin{otherlanguage}{persian}

استغفرک یا اله الوجود من ذکر یجد منه المقربین عرف التشبیه و الحدود ای ربّ
طهر ابصار خلقک عن کل ما یمنعها عن النظر الیک و التوکل علیک ای ربّ زینهم
بطراز العدل و الانصاف و نور قلوبهم بانوار معرفتک انک انت علی کلشئ قدیر

\end{otherlanguage}

Astaghfiruka yā Ilāha al-wujūd min dhikrin yajidu minhu al-muqarrabīn
ʿarfa al-tashbīhi wa-l-ḥudūd. Ay Rabbi, ṭahhir abṣāra khalqika ʿan kulli
mā yamnaʿuhā ʿani al-naẓari ilayka wa-l-tawakkuli ʿalayk. Ay Rabbi,
zayyin-hum bi-ṭirāzi al-ʿadli wa-l-inṣāfi wa-nawwir qulūbahum bi-anwāri
maʿrifatik. Innaka anta ʿalā kulli shayʾin qadīr.

I ask Your forgiveness, O God of existence, for any mention from which
those brought near to You perceive the scent of comparison and
limitation. O Lord, purify the sight of Your creatures from everything
that prevents them from looking to You and placing their trust in You. O
Lord, adorn them with the vesture of justice and fairness, and illumine
their hearts with the lights of Your knowledge. You are indeed powerful
over all things.

\subsection{163}\label{section-2138}

\begin{otherlanguage}{persian}

وقتی از اوقات مشرق آیات و مظهر بینات فرمودند یا عبد حاضر بیقین مبین بدان
حق جل جلاله حقش ثابت و امرش ظاهر و نورش باهر جمیع اعتراضات اهل بیان را
بیک کلمه که از قبل بآن نطق فرموده معدوم و مفقود و باطل مینماید قوله عز
بیانه یفعل ما یشاء و یحکم ما یرید

\end{otherlanguage}

Vaqtī az awqāt Mashreq-e Āyāt va Maẓhar-e Bayyenāt farmūdand: Yā ʿabd-e
ḥāżer, be yaqīn-e mobīn bedān Ḥaqq jalla jalāluhu ḥaqqash thābet va
amrash ẓāher va nūrash bāher. Jamīʿ-e eʿterāżāt-e ahl-e Bayān rā be yek
kalame ke az qabl be ān noṭq farmūde maʿdūm va mafqūd va bāṭel
mīnamāyad. Qawluhu ʿazza bayānuhu: Yafʿalu mā yashāʾu wa-yaḥkumu mā
yurīd.

On one occasion, the Dawning-place of the Verses and Manifestation of
Clear Proofs said: ``O servant in attendance, know with manifest
certainty that God---exalted be His glory---has an established right, a
manifest Cause, and a resplendent light. He renders all the objections
of the people of the Bayán nonexistent, null, and void through a single
utterance He previously spoke.'' His words---exalted be His utterance:
``He does whatever He wills and decrees whatever He desires.''

\subsection{164}\label{section-2139}

\begin{otherlanguage}{persian}

اگر نفسی فی الحقیقه باین کلمه علیا متمسک و موقن باشد جز باو توجه ننماید
و غیر او را نبیند و از غیر او نپذیرد و بلسان ظاهر و باطن میگوید

\end{otherlanguage}

Agar nafsī fī al-ḥaqīqe be īn kalame-ye ʿolyā motamassek va mūqen
bāshad, joz be Ū tavajjoh nanamāyad va gheyr-e Ū rā nabīnad va az
gheyr-e Ū napaẕīrad va be lesān-e ẓāher va bāṭen mīgūyad.

If a soul truly holds fast to this exalted utterance and is certain of
it, that soul will turn to none but Him, see none but Him, accept
nothing from anyone but Him, and say with both the outward and the
inward tongue:

\subsection{165}\label{section-2140}

\begin{otherlanguage}{arabic}

یا الهی و مقصودی لک الحمد بما اظهرت نفسک و انزلت آیاتک و ابرزت بیناتک و
ارسلت ما وجد منه کل ذی شم عرف قمیصک و نفحات ایامک و فوحات ظهورک

\end{otherlanguage}

Yā Ilāhī wa-maqṣūdī, laka al-ḥamdu bi-mā aẓharta nafsaka wa-anzalta
āyātika wa-abrazta bayyinātika wa-arsalta mā wajada minhu kullu dhī
shammin ʿarfa qamīṣika wa-nafaḥāti ayyāmika wa-fawaḥāti ẓuhūrik.

``O my God and my heart's desire, praise be to You for revealing
Yourself, sending down Your verses, bringing forth Your clear proofs,
and sending that from which everyone endowed with a sense of smell has
perceived the fragrance of Your garment, the sweet scents of Your days,
and the perfumes of Your Revelation.

\subsection{166}\label{section-2141}

\begin{otherlanguage}{persian}

ای ربّ قوّ قلبی لئلا تمنعه اشارات المعرضین من خلقک و الملحدین من عبادک عن
الاقبال الی افقک قوّ یا الهی یدی لاخذ کتابک بیقین تضطرب به افئدة المشرکین
من بریتک و نور وجهی بانوار وجهک و زین هیکلی بطراز الاستقامة فی حبک

\end{otherlanguage}

Ay Rabbi, qawwi qalbī li-allā tamnaʿahu ishārātu al-muʿriḍīna min
khalqika wa-l-mulḥidīna min ʿibādika ʿani al-iqbāli ilā ufuqik. Qawwi yā
Ilāhī yadī li-akhdhi kitābika bi-yaqīnin taḍṭaribu bihi afʾidatu
al-mushrikīna min bariyyatik, wa-nawwir wajhī bi-anwāri wajhika
wa-zayyin haykalī bi-ṭirāzi al-istiqāmati fī ḥubbik.

O Lord, strengthen my heart so that the allusions of those among Your
creatures who turn away, and of the faithless among Your servants, may
not prevent it from turning toward Your horizon. Strengthen my hand, O
my God, to take hold of Your Book with a certainty that makes the hearts
of those among Your creatures who associate partners with You tremble.
Illumine my face with the lights of Your face, and adorn my being with
the vesture of steadfastness in Your love.

\subsection{167}\label{section-2142}

\begin{otherlanguage}{arabic}

ای ربّ هذا یومک اشهد فیه بظهورک ظهرت الاسرار و اثمرت الاشجار و ثبت حکم
التوحید و محت شئونات التحدید و نصب علم البیان علی اعلی بقعة الامکان و
ارتفع ندائک فی کل الاحیان

\end{otherlanguage}

Ay Rabbi, hādhā yawmuka, ashhadu fīhi bi-ẓuhūrika ẓaharati al-asrāru
wa-athmarati al-ashjāru wa-thabata ḥukmu al-tawḥīdi wa-maḥat shuʾūnātu
al-taḥdīdi wa-nuṣiba ʿalamu al-bayāni ʿalā aʿlā buqʿati al-imkāni
wa-rtafaʿa nidāʾuka fī kulli al-aḥyān.

O Lord, this is Your Day. I testify that in it, through Your Revelation,
mysteries have been disclosed, trees have borne fruit, the rule of
divine unity has been established, all forms of limitation have been
effaced, the banner of utterance has been raised upon the highest spot
in the contingent world, and Your call has risen at all times.

\subsection{168}\label{section-2143}

\begin{otherlanguage}{arabic}

ای ربّ اسئلک باسمک الکریم و باسمک الحافظ و باسمک المعین ان لا تدعنی بنفسی
و انقذنی بذراعی قدرتک ثمّ احفظنی من اهل البیان الذین نقضوا عهدک و کسروا
میثاقک و نبذوا احکامک و جادلوا بآیاتک و کفروا ببرهانک اذ جئتهم من سماء
الایقان برایات الحکمة و البیان انّک انت المقتدر علی ما تشاء و فی قبضتک
ازمة الاشیاء لا اله الّاانت القوی القدیر '' انتهی

\end{otherlanguage}

Ay Rabbi, asʾaluka bi-smika al-Karīmi wa-bi-smika al-Ḥāfiẓi wa-bi-smika
al-Muʿīni an lā tadaʿanī bi-nafsī wa-anqidhni bi-dhirāʿay qudratika
thumma iḥfaẓnī min ahli al-Bayāni alladhīna naqaḍū ʿahdaka wa-kasarū
mīthāqaka wa-nabadhū aḥkāmaka wa-jādalū bi-āyātika wa-kafarū
bi-burhānika idh jiʾtahum min samāʾi al-īqāni bi-rāyāti al-ḥikmati
wa-l-bayān. Innaka anta al-muqtadiru ʿalā mā tashāʾu wa-fī qabḍatika
azimmatu al-ashyāʾ. Lā ilāha illā anta al-Qawiyyu al-Qadīr. Intahā.

O Lord, I ask You by Your name the All-Bountiful, by Your name the
Preserver, and by Your name the Helper, not to leave me to myself.
Rescue me with the arms of Your power, and then protect me from those
people of the Bayán who have broken Your covenant, shattered Your
pledge, cast aside Your laws, disputed Your verses, and rejected Your
proof when You came to them from the heaven of certitude with the
banners of wisdom and utterance. You are indeed powerful to do whatever
You will, and in Your grasp are the reins of all things. There is no God
but You, the Strong, the Almighty.'' End of quotation.

\subsection{169}\label{section-2144}

\begin{otherlanguage}{persian}

سبحان اللّه چه واقع شده که ناس از حق ظاهر باهر مشهود خود را ممنوع و محروم
ساخته‌اند و بگفته این و آن متمسک و مشغولند ظهور این امور تازگی نداشته و
ندارد در قرون و اعصار بحبال اوهام متمسک و بعد از ظهور آفتاب حقیقت و کشف
مفتریات متنبه نشده بر قدم اول راجع گشته و شاعر نیستند گویا مستوجب باشند
مجدد بوعظ واعظها و شرک جهال مبتلا شوند.

\end{otherlanguage}

Subḥāna Allāh, che vāqeʿ shode ke nās az Ḥaqq-e ẓāher-e bāher-e mashhūd
khod rā mamnūʿ va maḥrūm sākhte-and va be gofte-ye īn va ān motamassek
va mashghūland. Ẓohūr-e īn omūr tāzegī nadāshte va nadārad. Dar qorūn va
aʿṣār be ḥebāl-e awhām motamassek va baʿd az ẓohūr-e Āftāb-e Ḥaqīqat va
kashf-e moftariyāt motanabbeh nashode, bar qadam-e avval rājeʿ gashte va
shāʿer nīstand. Gūyā mostawjeb bāshand mojaddad be vaʿẓ-e vāʿeẓ-hā va
sharak-e johhāl mobtalā shavand.

Glory be to God! What has happened that people have barred and deprived
themselves of God, manifest, resplendent, and plainly visible, and cling
to and occupy themselves with the words of this person and that? Such
things have never been new, nor are they now. Throughout the centuries
and ages, they have clung to the cords of delusion, and after the Sun of
Truth has appeared and the fabrications have been exposed, they have
failed to awaken and have returned to their starting point without
realizing it. It is as though they deserve to be afflicted once again by
the preaching of preachers and the snares of the ignorant.

\subsection{170}\label{section-2145}

\begin{otherlanguage}{persian}

ای آقایان در امورات قبل و ثمرات و نتایج آن تفکر نمائید تا از امثال آن
نفوس طاغیه یاغیه محفوظ مانید فی الحقیقه ثمرات اعمال آن نفوس پندی است
بزرگ و نصحی است عظیم و عبرتی است عالی

\end{otherlanguage}

Ey āqāyān, dar omūrāt-e qabl va thamarāt va natāyej-e ān tafakkor
namāyīd tā az amthāl-e ān nofūs-e ṭāghiye-ye yāghiye maḥfūẓ mānīd. Fī
al-ḥaqīqe thamarāt-e aʿmāl-e ān nofūs pandī ast bozorg va noṣḥī ast
ʿaẓīm va ʿebratī ast ʿālī.

Gentlemen, reflect on the events of the past and their fruits and
consequences, so that you may remain protected from the likes of those
rebellious, defiant souls. Truly, the fruits of their deeds are a great
admonition, a mighty counsel, and a profound lesson.

\subsection{171}\label{section-2146}

\begin{otherlanguage}{persian}

باری خادم چه عرض نماید لعمر محبوبنا و محبوب من فی الارض و السماء
اوهاماتی در بعضی از اهل بیان مشاهده میشود که تا حین شبه آن ظاهر نشده شخص
عارفی از حزب اللّه در یکی از مدن با یکی متوهمین گفتگو نموده و آیات منزله
بیان را قرائت کرده و شاهد بر مطلب خود آورده شخص متوهم انکار نمود که آنچه
میگوئی در بیان نیست بعد بیان را آورد و عبارت حضرت را نشان داد آن شخص ذکر
نمود این بیان صحیح نیست گفت بیانی که در نزد شما هست بروید بیاورید گفت
آنهم صحیح نیست گفت پس بیان صحیح کجاست گفت نزد حضرت

\end{otherlanguage}

Bārī khādem che ʿarż namāyad? La-ʿamri maḥbūbinā wa-maḥbūbi man fī
al-arḍi wa-l-samāʾ, awhāmātī dar baʿżī az ahl-e Bayān moshāhede mīshavad
ke tā ḥīn shebh-e ān ẓāher nashode. Shakhṣ-e ʿārefī az Ḥezb Allāh dar
yekī az modon bā yekī motavahhemīn goftogū namūde va āyāt-e monzale-ye
Bayān rā qerāʾat karde va shāhed bar maṭlab-e khod āvarde. Shakhṣ-e
motavahhem enkār namūd ke ānche mīgūʾī dar Bayān nīst. Baʿd Bayān rā
āvard va ʿebārat-e Ḥażrat rā neshān dād. Ān shakhṣ dhekr namūd īn Bayān
ṣaḥīḥ nīst. Goft Bayānī ke dar nazd-e shomā hast beravīd biyāvarīd. Goft
ān ham ṣaḥīḥ nīst. Goft pas Bayān-e ṣaḥīḥ kojāst? Goft nazd-e Ḥażrat.

What, then, can this servant say? By the life of our Beloved and the
Beloved of all who are in earth and heaven, delusions are seen among
some of the people of the Bayán the like of which have never appeared
before. In one of the cities, a knowledgeable member of the people of
God conversed with one of these deluded people, reciting the revealed
verses of the Bayán and citing them as evidence for his point. The
deluded man denied it, saying, ``What you say is not in the Bayán.'' The
other then brought the Bayán and showed him the passage written by His
Holiness. The man said, ``This Bayán is not authentic.'' He replied,
``Go and bring the Bayán that you have.'' He said, ``That is not
authentic either.'' He asked, ``Where, then, is the authentic Bayán?''
He answered, ``With His Holiness.''

\subsection{172}\label{section-2147}

\begin{otherlanguage}{persian}

از رائحه منتنه این کلمه رایحه کلمات قبل مرور نموده چنانچه میگفتند ده
جزوه قرآن را سرقت نموده‌اند و قرآن تمام نزد قائم است و او میآورد فیالله
فیالله انهم جعلوا القرآن عضین یعنی قسمت نمودند یعنی پارچه پارچه نمودند
بعضی را دزد برد و بعضی را بدل نمودند

\end{otherlanguage}

Az rāʾeḥe-ye montene-ye īn kalame rāyeḥe-ye kalamāt-e qabl morūr namūde,
chenānche mīgoftand dah jozve-ye Qorʾān rā serqat namūde-and va Qorʾān-e
tamām nazd-e Qāʾem ast va Ū mīāvarad. Fa-yā Allāh, fa-yā Allāh, innahum
jaʿalū al-Qurʾāna ʿiḍīn, yaʿnī qesmat namūdand, yaʿnī pārche pārche
namūdand, baʿżī rā dozd bord va baʿżī rā badal namūdand.

The foul odor of this statement carries the odor of words spoken in the
past, as when they used to say, ``Ten sections of the Qur'an have been
stolen, and the complete Qur'an is with the Qá'im, who will bring it.''
O God, O God! ``They have made the Qur'an into fragments,'' meaning that
they divided it, that is, tore it into pieces: a thief carried off some
of it, and they altered some of it.

\subsection{173}\label{section-2148}

\begin{otherlanguage}{persian}

این خادم فانی شهادت میدهد که قرآن تمام بوده و بیان هم تمام است و قرآن
همان است که در دست کل بوده و بیان هم همین است که حال موجود است ولکن آن
کذاب مقصودش آنکه طرق را مسدود نماید و سبل را ممنوع و آنچه از ناحیهٴ ظنون
ظاهر شود سند نماید و ناس را مبتلا کند چنانچه از قبل کرده‌اند

\end{otherlanguage}

Īn khādem-e fānī shahādat mīdahad ke Qorʾān tamām būde va Bayān ham
tamām ast va Qorʾān hamān ast ke dar dast-e koll būde va Bayān ham hamīn
ast ke ḥāl mawjūd ast, va lāken ān kadhdhāb maqṣūdash ānke ṭoroq rā
masdūd namāyad va sobol rā mamnūʿ va ānche az nāḥeye-ye ẓonūn ẓāher
shavad sanad namāyad va nās rā mobtalā konad, chenānche az qabl
karde-and.

This evanescent servant testifies that the Qur'an was complete and that
the Bayán is also complete. The Qur'an is the very one that has been in
everyone's hands, and the Bayán is likewise the very one that now
exists. But that liar's purpose is to close the roads and bar the paths,
to make whatever arises from conjecture an authority, and to afflict the
people, just as others have done before.

\subsection{174}\label{section-2149}

\begin{otherlanguage}{persian}

باری امید چنان است ید قدرت حجبات جدیده را منع نماید و ناس بیچاره را
آگاهی بخشد آن قوم بی شرمند حرف آن غافل سبب ظهور اوهاماتست سبحان اللّه چه
شده اهل بیان را که صادق را کاذب میدانند و کاذب را صادق مثل فرقه هالکه که
خود را ناجیه میشمردند هزار و دویست سنه بل ازید جعفر را کذاب گفته و
میگویند مع آنکه آن مظلوم یک کلمه از روی حقیقت و راستی و صدق بیان نمود

\end{otherlanguage}

Bārī omīd chenān ast Yad-e Qodrat ḥojobāt-e jadīde rā manʿ namāyad va
nās-e bīchāre rā āgāhī bakhshad. Ān qawm bī-sharmand. Ḥarf-e ān ghāfel
sabab-e ẓohūr-e awhāmāt ast. Subḥāna Allāh, che shode ahl-e Bayān rā ke
ṣādeq rā kādheb mīdānand va kādheb rā ṣādeq, methl-e ferqe-ye hāleke ke
khod rā nājiye mīshomordand. Hezār va devīst sane bal azyad Jaʿfar rā
kadhdhāb gofte va mīgūyand, maʿa ānke ān maẓlūm yek kalame az rū-ye
ḥaqīqat va rāstī va ṣedq bayān namūd.

In any case, the hope is that the Hand of Power will ward off these new
veils and awaken the hapless people. Those people are shameless. The
words of that heedless man give rise to delusions. Glory be to God! What
has happened to the people of the Bayán that they consider the truthful
a liar and the liar truthful, like the lost sect that regarded itself as
saved? For twelve hundred years---indeed, longer---they have called and
continue to call Ja`far a liar, although that wronged man spoke a single
statement in accordance with truth, honesty, and sincerity.

\subsection{175}\label{section-2150}

\begin{otherlanguage}{arabic}

حال نفوس غافله معرضه متوهمه بسیارند از جمله میرزا هادی دولت آبادی خود را
هادی و قائد قوم میشمرد قسم بنفس حق که در لیالی و ایام بتحریر آیات منزله
تلقاء وجهش مشغول که شخص مستور تلقاء وجه قادر بر تکلم نبود این بیچاره‌ها
نمیدانند بهوی شنیده‌اند و بهوی متمسکند و بهوی متشبث و عاقبت ثمرات اعمال و
اقوالشان بمثابه اهل فرقان هبا بوده و خواهد بود

\end{otherlanguage}

Ḥāl nofūs-e ghāfele-ye moʿreże-ye motavahheme besyārand, az jomle Mīrzā
Hādī Dawlatābādī khod rā hādī va qāʾed-e qawm mīshomorad. Qasam be
Nafs-e Ḥaqq ke dar layālī va ayyām be taḥrīr-e āyāt-e monzale telqā-ye
vajhash mashghūl, ke shakhṣ-e mastūr telqā-ye vajh qāder bar takallom
nabūd. Īn bīchāre-hā nemīdānand, be havā shenīde-and va be havā
motamassekand va be havā motashabbeth, va ʿāqebat thamarāt-e aʿmāl va
aqvāl-eshān be mathābe-ye ahl-e Forqān habā būde va khāhad būd.

There are now many heedless, deluded souls who turn away. Among them is
Mírzá Hádí Dawlatábádí, who considers himself the people's guide and
leader. I swear by God Himself that, by night and day, he was occupied
with transcribing the verses revealed in His presence, while the
concealed person was unable to speak in His presence. These hapless
people do not know. They have listened out of desire, cling to desire,
and hold fast to desire. In the end, the fruits of their deeds and
words, like those of the people of the Qur'an, have been and will be
scattered dust.

\subsection{176}\label{section-2151}

\begin{otherlanguage}{persian}

از حق میطلبم انصاف عطا فرماید و ضغینه و بغضای قبل را از میان خلق بردارد
ولکن بسیار مشکل است چه که نفوسی که در قرون و اعصار بظنون و اوهام تربیت
شده‌اند کجا لایق مشاهده انوار آفتاب یقینند

\end{otherlanguage}

Az Ḥaqq mīṭalabam enṣāf ʿaṭā farmāyad va żaghīne va baghżā-ye qabl rā az
miyān-e khalq bardārad, va lāken besyār moshkel ast, che ke nofūsī ke
dar qorūn va aʿṣār be ẓonūn va awhām tarbiyat shode-and kojā lāyeq-e
moshāhede-ye anvār-e Āftāb-e Yaqīn-and?

I ask God to bestow fairness and remove the old rancor and hatred from
among His creatures. Yet this is very difficult, for how can souls
nurtured in conjectures and delusions throughout the centuries and ages
be worthy to behold the lights of the Sun of Certitude?

\subsection{177}\label{section-2152}

\begin{otherlanguage}{persian}

این عبد از اظهار و ذکر بعضی از امور ممنوع است و الا گفته میشد آنچه را که
تا حین عباد از او غافل و بیخبرند قوله جل جلاله و عم نواله یا عبد حاضر
سخن باندازه میگوئیم لا بما ینبغی لملکوت بیانی و جبروت علمی الا ان نری
البیت خالیا من دونی انتهی.

\end{otherlanguage}

Īn ʿabd az eẓhār va dhekr-e baʿżī az omūr mamnūʿ ast, va ellā gofte
mīshod ānche rā ke tā ḥīn ʿebād az ū ghāfel va bīkhabarand. Qawluhu
jalla jalāluhu wa-ʿamma nawāluhu: Yā ʿabd-e ḥāżer, sokhan be andāze
mīgūʾīm, lā bi-mā yanbaghī li-malakūti bayānī wa-jabarūti ʿilmī illā an
narā al-bayta khāliyan min dūnī. Intahā.

This servant is forbidden to disclose and mention certain matters;
otherwise, things would be said of which the servants have until now
remained heedless and unaware. His words---exalted be His glory and
all-embracing His bounty: ``O servant in attendance, We speak within
measure, not as befits the kingdom of My utterance and the dominion of
My knowledge, unless We see the House empty of all but Me.'' End of
quotation.

\subsection{178}\label{section-2153}

\begin{otherlanguage}{arabic}

ایم اللّه این خادم لوجه اللّه گفته و خواهد گفت خیالی نداشته و ندارد انّ ربّی
الرّحمن هو الغنی عن ایمان اهل البیان الذین نقضوا میثاقه و حاربوا بنفسه

\end{otherlanguage}

Aymu Allāh, īn khādem li-wajhi Allāh gofte va khāhad goft, khiyālī
nadāshte va nadārad. Inna rabbī al-Raḥmāna huwa al-ghaniyyu ʿan īmāni
ahli al-Bayāni alladhīna naqaḍū mīthāqahu wa-ḥārabū bi-nafsih.

By God, this servant has spoken and will speak for the sake of God, and
has had and has no ulterior aim. My Lord, the All-Merciful, is indeed
independent of the faith of those people of the Bayán who have broken
His covenant and made war against Him.

\subsection{179}\label{section-2154}

\begin{otherlanguage}{arabic}

نقطه بیان روح ما سواه فداه میفرماید در آن روز آن آفتاب حقیقت اهل بیان را
خطاب مینماید و این سوره فرقان را تلاوت میفرماید قوله تعالی قل یا ایّها
الکافرون لا اعبد ما تعبدون و لا انتم عابدون ما اعبد و لا انا عابد ما
عبدتم و لا انتم عابدون ما اعبد لکم دینکم ولی دینی.

\end{otherlanguage}

Noqṭe-ye Bayān, rūḥ-e mā sevāh fedāh, mīfarmāyad dar ān rūz ān Āftāb-e
Ḥaqīqat ahl-e Bayān rā kheṭāb mīnamāyad va īn sūre-ye Forqān rā telāvat
mīfarmāyad. Qawluhu taʿālā: Qul yā ayyuhā al-kāfirūn, lā aʿbudu mā
taʿbudūn, wa-lā antum ʿābidūna mā aʿbud, wa-lā anā ʿābidun mā ʿabadtum,
wa-lā antum ʿābidūna mā aʿbud. Lakum dīnukum wa-liya dīnī.

The Point of the Bayán---may the souls of all others be sacrificed for
Him---says that on that Day the Sun of Truth will address the people of
the Bayán and recite this surah of the Qur'an. His words, exalted be He:
``Say: O unbelievers, I do not worship what you worship, nor do you
worship what I worship. Nor am I a worshipper of what you have
worshipped, nor are you worshippers of what I worship. To you your
religion, and to me mine.''

\subsection{180}\label{section-2155}

\begin{otherlanguage}{persian}

این ایام مکرر از لسان قدم این سوره مبارکه شنیده شد طوبی للمنصفین و طوبی
للمتفرسین و طوبی للمتفکرین بر هر نفسی الیوم لازم که بقدرت و قوت حق جل
جلاله افق آفتاب یقین را از سحاب ظنون و اوهام حفظ نماید اگر چه او بنفسه
حافظ کل بوده و هست و آنچه در این مقامات ذکر میشود مقصود آنکه عبادش بخدمت
امرش فائز شوند یک ساعت دو ساعت دنیا قابل ذکر نبوده و نیست امید آنکه باب
فضل را مسدود نفرماید و اولیاء خود را از عرفش محروم نسازد.

\end{otherlanguage}

Īn ayyām mokarrar az Lesān-e Qedam īn sūre-ye mobārake shenīde shod.
Ṭūbā li-l-munṣifīn wa-ṭūbā li-l-mutafarrisīn wa-ṭūbā li-l-mutafakkirīn.
Bar har nafsī al-yawm lāzem ke be qodrat va qovvat-e Ḥaqq jalla jalāluhu
ofoq-e Āftāb-e Yaqīn rā az saḥāb-e ẓonūn va awhām ḥefẓ namāyad, agarche
Ū be nafsihī ḥāfeẓ-e koll būde va hast. Va ānche dar īn maqāmāt dhekr
mīshavad maqṣūd ānke ʿebādash be khedmat-e Amrash fāʾez shavand. Yek
sāʿat do sāʿat-e donyā qābel-e dhekr nabūde va nīst. Omīd ānke bāb-e
fażl rā masdūd nafarmāyad va awliyā-ye khod rā az ʿarfash maḥrūm
nasāzad.

In these days this blessed surah has repeatedly been heard from the
Tongue of Eternity. Blessed are the fair-minded, blessed are the
discerning, and blessed are those who reflect. It is incumbent on every
soul today, through the power and might of God---exalted be His
glory---to protect the horizon of the Sun of Certitude from the clouds
of conjecture and delusion, although He Himself has ever been and
remains the Protector of all. The purpose of what is said in these
passages is that His servants may attain the privilege of serving His
Cause. The world's one or two hours have never been and are not worthy
of mention. The hope is that He will not close the door of grace or
deprive His loved ones of His fragrance.

\subsection{181}\label{section-2156}

\begin{otherlanguage}{persian}

ای آقایان مخصوص یک باب از بیان در ذکر این بیان است میفرماید شجره اثبات
باعراضش از او از نفی محسوب و شجره نفی باقبالش باو از اثبات محسوب باری
بمتوهمین بگوئید امروز روزیست عظیم و بنور توحید منور مجال تحدید شرکأ
نبوده و نیست چنانچه سالها بآن مشغول بودند عصمت کبری مخصوص ذات مقدس بوده
وهست احدی را از آن نصیبی نه و هو ینطق فی قطب الابداع لا اله الّا هو
المهیمن القیوم

\end{otherlanguage}

Ey āqāyān, makhṣūṣ yek bāb az Bayān dar ẕekr-e īn bayān ast. Mīfarmāyad:
shajare-ye ethbāt be-eʿrāżash az ū az nafy maḥsūb va shajare-ye nafy
be-eqbālash be-ū az ethbāt maḥsūb. Bārī, be-motavahhemīn begūʾīd emrūz
rūzīst ʿaẓīm va be-nūr-e towḥīd monavvar. Majāl-e taḥdīd-e shorakā
nabūde va nīst, chenānche sālhā be-ān mashghūl būdand. ʿEṣmat-e kobrā
makhṣūṣ-e ẕāt-e moqaddas būde va hast; aḥadī rā az ān naṣībī na. Wa-huwa
yanṭiqu fī quṭb al-ibdāʿ: lā ilāha illā huwa al-Muhaymin al-Qayyūm.

Gentlemen! An entire chapter of the Bayán is devoted specifically to
this statement. He says: ``The tree of affirmation, by turning away from
Him, is reckoned among those of negation, and the tree of negation, by
turning toward Him, is reckoned among those of affirmation.'' In short,
tell those caught in delusion: today is a mighty day, illumined by the
light of divine unity. There has never been, nor is there now, any scope
for defining partners with God, as they spent years doing. Supreme
infallibility has always belonged, and still belongs, exclusively to the
sanctified Essence; no one else has any share in it. And He proclaims at
the heart of creation: ``There is no God but Him, the All-Protecting,
the Self-Subsisting.''

\subsection{182}\label{section-2157}

\begin{otherlanguage}{arabic}

ای اهل بیان بشنوید عرض این خادم فانی را و بمثل اهل فرقان بظلم قیام
منمائید جمیع علمای آن فرقه کمال ظلم را مینمودند و خود را مظلوم میشمردند
اتقوا اللّه یا قوم و لا تکونوا من الظالمین مقامی را که میفرماید اگر نفسی
ادعا نماید و اتیان به حجتی ننماید تعرض منمائید و حزن وارد نیاورید حال مع
حجج و براهینی که عالم را احاطه نموده و نورش آفاق و انفس را منور کرده
گفته‌اند آنچه را که هیچ ظالمی از قبل و بعد نگفته

\end{otherlanguage}

Ey ahl-e Bayān, beshnavīd ʿarż-e īn khādem-e fānī rā va be-methl-e ahl-e
Forqān be-ẓolm qiyām manamāyīd. Jamīʿ-e ʿolamā-ye ān ferqe kamāl-e ẓolm
rā mīnamūdand va khod rā maẓlūm mīshomordand. Ittaqū Allāha yā qawm
wa-lā takūnū min al-ẓālimīn. Maqāmī rā ke mīfarmāyad agar nafsī eddeʿā
namāyad va etyān be-ḥojjatī nanamāyad taʿarroż manamāyīd va ḥozn vāred
nayāvarīd, ḥāl maʿ-e ḥojaj va barāhīnī ke ʿālam rā eḥāṭe namūde va
nūrash āfāq va anfos rā monavvar karde, gofte-and ānche rā ke hīch
ẓālemī az qabl va baʿd nagofte.

People of the Bayán! Listen to the words of this transitory servant, and
do not rise up in oppression as did the people of the Qur'an. All the
divines of that community practiced the utmost oppression, yet
considered themselves oppressed. Fear God, people, and do not be among
the oppressors. Although He says that if a soul makes a claim without
producing any proof, you must not interfere with him or cause him
sorrow, now, despite proofs and testimonies that have encompassed the
world and whose light has illumined the outer world and the world
within, they have said what no oppressor, earlier or later, has ever
said.

\subsection{183}\label{section-2158}

\begin{otherlanguage}{persian}

لله تفکر نمائید شاید از بحر معانی لئالی بیان را بیابید جمیع بیانات آن
حضرت رحمتی بود از برای عباد لا تجعلوها نقمة لکم و لمن بعدکم برچینید
بساطهای قبل وبعد را و بکمال فرح و انبساط در این بساط مبسوط وارد شوید از
حق جل جلاله میطلبم معرضین و منکرین را تائید فرماید بر رجوع و توبه و
انابه بسیار حیف است خود را مستحق سیاط و غضب دائمی نمائید

\end{otherlanguage}

Li-llāh tafakkor namāyīd shāyad az baḥr-e maʿānī laʾālī-ye bayān rā
beyābīd. Jamīʿ-e bayānāt-e ān Ḥażrat raḥmatī būd az barā-ye ʿebād. Lā
tajʿalūhā niqmatan lakum wa-li-man baʿdakum. Barchīnīd besāṭhā-ye qabl
va baʿd rā va be-kamāl-e faraḥ va enbesāṭ dar īn besāṭ-e mabsūṭ vāred
shavīd. Az Ḥaqq jalla jalāluhu mīṭalabam moʿreżīn va monkerīn rā taʾyīd
farmāyad bar rojūʿ va towbe va enābe. Besyār ḥeyf ast khod rā
mostaḥaqq-e seyāṭ va ghażab-e dāʾemī namāyīd.

For God's sake, reflect, that you may find the pearls of utterance in
the ocean of meanings. All His utterances were a mercy for His servants.
Do not make them a scourge for yourselves and those who come after you.
Roll up the arrangements of past and future, and enter with the utmost
joy and delight into this gathering now spread before you. I ask
God---exalted be His glory---to assist those who turn away and deny in
returning, repenting, and turning humbly to Him. What a great pity it
would be to make yourselves deserving of scourges and unending wrath.

\subsection{184}\label{section-2159}

\begin{otherlanguage}{persian}

آیا میتوان این سدرهٴ مبارکهٴ مرتفعه که فرعش از صد هزار سماءِ مرتفع‌تر است
انکار نمود؟ آیا مجال توقف وگفتگوباقی؟ لا و نفسه الحق مگر قهر الهی بسبب
جزای اعمال اخذ نماید سبحان اللّه نفوس غافله اراده نموده‌اند آفتاب حقیقت را
سترنمایند این فرد در این مقام مناسب است:

\end{otherlanguage}

Āyā mītavān īn Sedre-ye mobārake-ye mortafeʿe ke farʿash az ṣad hezār
samāʾ mortafeʿtar ast enkār namūd? Āyā majāl-e tavaqqof va goftogū bāqī?
Lā wa-nafsihi al-Ḥaqq. Magar qahr-e elāhī be-sabab-e jazā-ye aʿmāl akhẕ
namāyad. Subḥān Allāh! Nofūs-e ghāfele erāde namūde-and Āftāb-e Ḥaqīqat
rā satr namāyand. Īn fard dar īn maqām monāseb ast:

Can this blessed, exalted Tree, whose branch rises higher than a hundred
thousand heavens, be denied? Is there any room left for hesitation or
debate? No, by His own Self, the True One! Unless divine wrath should
seize someone in requital for their deeds. Glory be to God! Heedless
souls have sought to conceal the Sun of Truth. This couplet is fitting
here:

\subsection{185}\label{section-2160}

\begin{otherlanguage}{arabic}

``قصد آن دارند این گل پاره‌ها کز حسد پوشند خورشید ترا''

\end{otherlanguage}

``Qaṣd-e ān dārand īn gel pārehā kaz ḥasad pūshand khorshīd-e to rā.''

``These clods of clay intend, out of envy, to conceal Your sun.''

\subsection{186}\label{section-2161}

\begin{otherlanguage}{persian}

بلایائی بر سدره مبارکه وارد که اگر فی الجمله انسان در او تفکر نماید خود
را هلاک کند چون این فقره منع شده لذا این عبد و امثال این عبد باقی مانده
قسم بآفتاب افق راستی اقلام عالم از ذکرش عاجز و قلوب ملأ اعلی در احتراق

\end{otherlanguage}

Balāyāʾī bar Sedre-ye mobārake vāred ke agar fī al-jomle ensān dar ū
tafakkor namāyad khod rā halāk konad. Chūn īn faqare manʿ shode leẕā īn
ʿabd va amthāl-e īn ʿabd bāqī mānde. Qasam be-Āftāb-e ofoq-e rāstī,
aqlām-e ʿālam az ẕekrash ʿājez va qolūb-e malaʾ-e aʿlā dar eḥterāq.

Such calamities have befallen the blessed Tree that anyone who reflected
on them even briefly would take their own life. Since this has been
forbidden, this servant and others like him have remained alive. I swear
by the Sun of the horizon of truthfulness: the pens of the world are
powerless to recount them, and the hearts of the Concourse on high are
burning.

\subsection{187}\label{section-2162}

\begin{otherlanguage}{persian}

از جمله در این ارض بعد از ورود در سجن اعظم دو نفس غافل محتجب با هر کسی
مراوده مینمودند و در هر مجلسی با هر شخصی فصلی مذکور میداشتند اللّه یعلم
ما ورد منها علیه فی دیار الغربة و بلاد الغربة بعضی از سور را با مجعولات
خود ترکیب نموده و بهر کسی ورقی دادند و اللّه الذی لا اله الّا هو حتی بکناس
بلده تا آنکه امر بمقامی رسید که آتش فتنه مشتعل و لهیب نار بغضا از قلوب و
نفوس ظاهر بالاخره از مالک رقاب بسد باب حکم صادر چند شهر با احدی از داخل
و خارج ملاقات نفرمودند باب مسدود و طالبین ممنوع و محجوب

\end{otherlanguage}

Az jomle dar īn arż baʿd az vorūd dar Sejn-e Aʿẓam do nafs-e ghāfel-e
moḥtajeb bā har kasī morāvede mīnamūdand va dar har majlesī bā har
shakhṣī faṣlī maẕkūr mīdāshtand. Allāhu yaʿlamu mā warada minhā ʿalayhi
fī diyār al-ghurba wa-bilād al-ghurba. Baʿżī az sovar rā bā majʿūlāt-e
khod tarkīb namūde va be-har kasī varaqī dādand, va Allāh allaẕī lā
ilāha illā huwa, ḥattā be-kannās-e balde. Tā ānke amr be-maqāmī resīd ke
ātash-e fetne moshtaʿel va lahīb-e nār-e baghżā az qolūb va nofūs ẓāher.
Bel-ākhare az Mālek-e Reqāb be-sadd-e bāb ḥokm ṣāder. Chand shahr bā
aḥadī az dākhel va khārej molāqāt nafarmūdand. Bāb masdūd va ṭālebīn
mamnūʿ va maḥjūb.

Among these afflictions was this: in this land, after the arrival in the
Most Great Prison, two heedless, veiled souls associated with everyone
and, in every gathering, gave each person an account. God knows what
befell Him through them in the regions and lands of exile. They mixed
some of the surahs with their own fabrications and gave a sheet to
everyone---by God, besides whom there is no other God, even to the
town's scavenger. Matters reached the point where the fire of sedition
was kindled and the blaze of hatred flared up from hearts and souls. At
last, the Lord of all issued an order that the door be closed. For
several months He met with no one, whether from within or without. The
door was closed, and seekers were barred and kept away.

\subsection{188}\label{section-2163}

\begin{otherlanguage}{persian}

تا آنکه شبی از شبها این فانی تلقاء وجه بتحریر آیات مشغول بغتة لحن مبارک
تغییر نمود این آیات منیعه قویه غالبه قاهره نازل بشأنیکه فرائص این عبد
متزلزل بود حق شاهد و گواه که بهیچوجه اطمینان و سکون در نفس باقی نماند
قوله جل جلاله و عم نواله و عظم بیانه قد ماج بحرالبلاء و احاطت الامواج
فلک اللّه المهیمن القیوم از اصغاء این آیه مبارکه روح از برای این عبد باقی
نه بعد این آیه نازل انّک انت یا ملاح لا تضطرب من الاریاح لان فالق الاصباح
معنا فی هذه الظلمة التی احاطت العالمین از این آیه فی الجمله سکون حاصل
ولکن واضح و مبرهن شد که فتنه دهماء از عقب است

\end{otherlanguage}

Tā ānke shabī az shabhā īn fānī telqāʾ-e vajh be-taḥrīr-e āyāt mashghūl,
baghtatan laḥn-e mobārak taghyīr namūd. Īn āyāt-e manīʿe-ye qaviyye-ye
ghālebe-ye qāhere nāzel be-shaʾnī ke farāʾeṣ-e īn ʿabd motazalzel būd.
Ḥaqq shāhed va govāh ke be-hīch vajh eṭmīnān va sokūn dar nafs bāqī
namānd. Qawluhu jalla jalāluhu wa-ʿamma nawāluhu wa-ʿaẓuma bayānuhu: qad
māja baḥr al-balāʾ wa-aḥāṭat al-amwāj fulk Allāh al-Muhaymin al-Qayyūm.
Az eṣghā-ye īn āye-ye mobārake rūḥ az barā-ye īn ʿabd bāqī na. Baʿd īn
āye nāzel: innaka anta yā mallāḥ lā taḍṭarib min al-aryāḥ li-anna Fāliq
al-Iṣbāḥ maʿanā fī hāẕihi al-ẓulma allatī aḥāṭat al-ʿālamīn. Az īn āye
fī al-jomle sokūn ḥāṣel, va-līkan vāżeḥ va mobarhan shod ke fetne-ye
dahmā az ʿaqab ast.

Then, one night, while this transitory one was in His presence, occupied
in transcribing the verses, the blessed tone suddenly changed. These
exalted, mighty, triumphant, all-subduing verses were revealed in such a
manner that this servant's limbs trembled. God is my witness and
testimony that no assurance or tranquility whatsoever remained in my
soul. His words---exalted be His glory, all-embracing His bounty, and
mighty His utterance: ``The ocean of tribulation has surged, and the
waves have encompassed the Ark of God, the All-Protecting, the
Self-Subsisting.'' On hearing this blessed verse, no life remained in
this servant. Then this verse was revealed: ``O Mariner, do not be
troubled by the winds, for the Cleaver of the Dawn is with Us in this
darkness that has encompassed the worlds.'' From this verse some
tranquility was gained, but it became clear and certain that a dark
upheaval lay ahead.

\subsection{189}\label{section-2164}

\begin{otherlanguage}{persian}

باری آنچه از بعد واقع حرف بحرف در آن لوح مبارک نازل از جمله میفرماید
وقتی بایادی امرا و سفرا و احبا نصرت مینمائیم و گاهی بایادی اعدا یوم بعد
هنگام عصر این عبد تلقاء وجه بتحریر آیات مشغول که غوغا بلند شد جمیع اهل
بلد و عسکریه با اسیاف شاهره امام بیت حاضر اعمال شنیعهٴ آن نفوس مجسم شد و
ظاهر شد آنچه شد مدتی این عبد مضطرب و متزلزل که مبادا بر جمال قدم حزنی
وارد آید لعمر ربنا اگر بحق واگذار شده بود خود کفایت میفرمود چنانچه ذئب و
رقشاء وامثال آن نفوس را کفایت نمود

\end{otherlanguage}

Bārī ānche az baʿd vāqeʿ ḥarf be-ḥarf dar ān Lowḥ-e mobārak nāzel. Az
jomle mīfarmāyad vaqtī be-ayādī-ye omarā va sofarā va aḥebbā noṣrat
mīnamāyīm va gāhī be-ayādī-ye aʿdā. Yowm-e baʿd hengām-e ʿaṣr īn ʿabd
telqāʾ-e vajh be-taḥrīr-e āyāt mashghūl ke ghowghā boland shod. Jamīʿ-e
ahl-e balad va ʿaskarīye bā asyāf-e shāhere amām-e bayt ḥāżer. Aʿmāl-e
shanīʿe-ye ān nofūs mojassam shod va ẓāher shod ānche shod. Moddatī īn
ʿabd możṭareb va motazalzel ke mabādā bar Jamāl-e Qedam ḥoznī vāred
āyad. La-ʿamr-e Rabbenā, agar be-Ḥaqq vāgoẕār shode būd khod kefāyat
mīfarmūd, chenānche Ẕeʾb va Raqshā va amthāl-e ān nofūs rā kefāyat
namūd.

In short, everything that happened afterward was revealed, word for
word, in that blessed Tablet. Among other things, He says: ``At times We
render assistance through the hands of rulers, ambassadors, and loved
ones, and at times through the hands of enemies.'' The following
afternoon, this servant was in His presence, occupied in transcribing
the verses, when an uproar arose. All the townspeople and the soldiery
appeared before the house with drawn swords. The shameful deeds of those
souls took tangible form, and what happened, happened. For a time this
servant was anxious and trembling, lest any sorrow befall the Ancient
Beauty. By the life of our Lord, had it been left to God, He Himself
would have sufficed, just as He sufficed against the Wolf, the
She-Serpent, and others like them.

\subsection{190}\label{section-2165}

\begin{otherlanguage}{persian}

بعد اخت رضا قلی یک بسته سور و الواح منزله را برد و بدست پاشای بلد داد
ازجمله سوره رئیس بگمان اینکه پاشای مزبور بعناد قیام نماید ولکن حکم آیه
منزله ظاهر پاشا قدری مطالعه نمود صبح برخاست و خدمت حضرت غصن اللّه الاعظم
روحی و ذاتی لقدومه الفداء رسید و جمیع سور را آورد و تسلیم نمود مذکور
داشت شب گذشته در این الواح ملاحظه کردم عظمت امر مرا اخذ نمود بیقین
دانستم این امر حق است چه که اگر او را انکار نمائیم جمیع رسل و صحف را
باید انکار نمود و استدعای لقا وحضور نمود

\end{otherlanguage}

Baʿd okht-e Reżā-Qolī yek baste sovar va alvāḥ-e monazzale rā bord va
be-dast-e pāshā-ye balad dād, az jomle Sūre-ye Raʾīs, be-gomān-e īnke
pāshā-ye mazbūr be-ʿenād qiyām namāyad. Va-līkan ḥokm-e āye-ye monazzale
ẓāher. Pāshā qadrī moṭāleʿe namūd, ṣobḥ barkhāst va khedmat-e Ḥażrat-e
Ghoṣn Allāh al-Aʿẓam, rūḥī va ẕātī le-qodūmehi al-fedāʾ, resīd va
jamīʿ-e sovar rā āvard va taslīm namūd. Maẕkūr dāsht: shab-e goẕashte
dar īn alvāḥ molāḥeẓe kardam, ʿaẓamat-e Amr marā akhẕ namūd. Be-yaqīn
dānestam īn Amr ḥaqq ast, che ke agar ū rā enkār namāyīm jamīʿ-e rosol
va ṣoḥof rā bāyad enkār namūd. Va estedʿā-ye leqā va ḥożūr namūd.

Afterward, the sister of Riḍā-Qulī took a bundle of revealed surahs and
Tablets, including the Surah of the Chief, and handed it to the governor
of the town, supposing that he would rise up in opposition. Instead, the
decree of the revealed verse became manifest. The governor read some of
them, rose in the morning, and went to His Holiness the Most Great
Branch of God---may my soul and my being be sacrificed at His
feet---bringing all the surahs and handing them over. He said: ``Last
night I examined these Tablets, and the greatness of the Cause seized
me. I knew with certainty that this Cause is true, for if we deny Him,
we must deny all the Messengers and Scriptures.'' And he requested an
audience and admission to His presence.

\subsection{191}\label{section-2166}

\begin{otherlanguage}{persian}

تا آنکه یومی عبود نام که صاحب بیت بود بین یدی حاضر و استدعا نمود که حضرت
پاشا بحضور فائز شود قبول فرمودند و بعد از حضور تلقاء وجه بایمان و ایقان
فائز و آنچه باو امر شد جمیع را عمل نمود. از برای عامه محبوسین که در این
شهر بودند فرج عظیم حاصل

\end{otherlanguage}

Tā ānke yowmī ʿAbbūd nām ke ṣāḥeb-e bayt būd bayn-e yaday ḥāżer va
estedʿā namūd ke Ḥażrat-e Pāshā be-ḥożūr fāʾez shavad. Qabūl farmūdand
va baʿd az ḥożūr telqāʾ-e vajh be-īmān va īqān fāʾez va ānche be-ū amr
shod jamīʿ rā ʿamal namūd. Az barā-ye ʿāmme-ye maḥbūsīn ke dar īn shahr
būdand faraj-e ʿaẓīm ḥāṣel.

Eventually, one day, a man named ʿAbbūd, who owned the house, came
before Him and asked that His Excellency the governor be admitted to His
presence. He consented, and after entering His presence, the governor
attained faith and certitude and carried out everything he was commanded
to do. Great relief came to all the prisoners who were in this city.

\subsection{192}\label{section-2167}

\begin{otherlanguage}{persian}

مقصود آنکه آنچه از قوه نفوس غافله بر میآید مجری نمودند جمیع خلق این ارض
را بشأنی اغوی نمودند که هر یک از این اسرا را میدیدند کلمات نا پسندیده
میگفتند. باری سبب شدند چند شهر باب رحمت بر روی کل مسدود اللّه یعلم ما ورد
علیه من اعدائه هر نفسی بذره از عدل و انصاف فائز باشد شهادت میدهد که اعلی
و ابهی و اظهر از این ظهور در عالم نیامده

\end{otherlanguage}

Maqṣūd ānke ānche az qovve-ye nofūs-e ghāfele bar mīāyad mojrā namūdand.
Jamīʿ-e khalq-e īn arż rā be-shaʾnī eghvā namūdand ke har yek az īn
osarā rā mīdīdand kalemāt-e nāpasandīde mīgoftand. Bārī sabab shodand
chand shahr bāb-e raḥmat bar rū-ye koll masdūd. Allāhu yaʿlamu mā warada
ʿalayhi min aʿdāʾihi. Har nafsī be-ẕarre az ʿadl va enṣāf fāʾez bāshad
shahādat mīdahad ke aʿlā va abhā va aẓhar az īn Ẓohūr dar ʿālam
nayāmade.

The point is that they carried out everything within the power of
heedless souls. They so misled all the people of this land that whenever
they saw any of these captives, they uttered offensive words. In short,
they caused the door of mercy to be closed to everyone for several
months. God knows what befell Him at the hands of His enemies. Anyone
possessing even an atom of justice and fairness will testify that no
Revelation more exalted, more glorious, or more manifest than this
Revelation has appeared in the world.

\subsection{193}\label{section-2168}

\begin{otherlanguage}{persian}

آیا معرضین حکایت مباهله واقع شده در ارض سرّ را نشنیده‌اند لعمر مقصودنا
ظاهر شد آنچه که اهل مدینه را متحیر نمود چون احتمال میرود بعضی نشنیده
باشند لذا مختصر عرض میشود تا منصفین آگاه شوند و شاید معرضین هم راه یابند

\end{otherlanguage}

Āyā moʿreżīn ḥekāyat-e mobāhele-ye vāqeʿ-shode dar Arż-e Serr rā
nashnīde-and? La-ʿamr-e Maqṣūdenā, ẓāher shod ānche ke ahl-e madīne rā
motaḥayyer namūd. Chūn eḥtemāl mīravad baʿżī nashnīde bāshand leẕā
mokhtaṣar ʿarż mīshavad tā monṣefīn āgāh shavand va shāyad moʿreżīn ham
rāh yāband.

Have those who turn away not heard the account of the challenge to
invoke God's judgment that took place in the Land of Mystery? By the
life of our Desired One, what came to pass astonished the people of the
city. Since some may not have heard of it, a brief account is therefore
given, so that the fair-minded may be informed and perhaps those who
turn away may also find the way.

\subsection{194}\label{section-2169}

\begin{otherlanguage}{persian}

بعد از ورود ارض سرّ یومی از ایام سید محمد اصفهانی تلقاء وجه حاضر عرض
مطلبی نمود قبول نشد و چند یوم قبل هم لسان عظمت مخاطبا ایاه باینکلمات
ناطق یا محمد تو از مشی انبیاء و سجیه اصفیاء اطلاع نداری و آگاه نیستی
انتهی. xx

\end{otherlanguage}

Baʿd az vorūd-e Arż-e Serr yowmī az ayyām Seyyed Moḥammad-e Eṣfahānī
telqāʾ-e vajh ḥāżer, ʿarż-e maṭlabī namūd; qabūl nashod. Va chand yowm
qabl ham Lesān-e ʿAẓamat mokhāṭaban iyyāhu be-īn kalemāt nāṭeq: yā
Moḥammad, to az mashy-e anbiyā va sajiyye-ye aṣfiyā eṭṭelāʿ nadārī va
āgāh nīstī. Intahā. {[}corrupted text{]}

After the arrival in the Land of Mystery, one day Siyyid
Muḥammad-i-Iṣfahānī came into His presence and presented a matter, which
was not accepted. A few days earlier, too, the Tongue of Grandeur had
addressed him in these words: ``Muḥammad, you have no knowledge or
awareness of the ways of the Prophets and the character of the chosen
ones.'' End of quotation. {[}corrupted original text{]}

\subsection{195}\label{section-2170}

\begin{otherlanguage}{persian}

از این کلمه مبارکه همچه معلوم شد بعضی از حرفهای عرفای کهنه را نزد بعضی
گفته و همچنین چند یوم بعد هم اذن خواست بعد از حضور عرض نمود استدعا آنکه
بمیرزا یحیی امر فرمائید چیزی ننویسد چه که آقا محمد علی اصفهانی یک شعر
فارسی سئوال نموده معنی آنرا ملتفت نشده جمال قدم فرمودند سیّد تو را چه
باین فضولی‌ها تا آنکه بالاخره طردش فرمودند

\end{otherlanguage}

Az īn kaleme-ye mobārake hamcho maʿlūm shod baʿżī az ḥarfhā-ye ʿorafā-ye
kohne rā nazd-e baʿżī gofte. Va hamchenīn chand yowm baʿd ham eẕn khāst.
Baʿd az ḥożūr ʿarż namūd: estedʿā ānke be-Mīrzā Yaḥyā amr farmāyīd chīzī
nanavīsad, che ke Āqā Moḥammad-ʿAlī-ye Eṣfahānī yek sheʿr-e fārsī soʾāl
namūde maʿnī-ye ān rā moltafet nashode. Jamāl-e Qedam farmūdand: Seyyed,
to rā che be-īn fożūlīhā? Tā ānke bel-ākhare ṭardash farmūdand.

From this blessed utterance it appeared that he had repeated to certain
people some of the sayings of the mystics of old. A few days later he
again requested permission to enter. Having entered His presence, he
said: ``I ask that You instruct Mīrzā Yaḥyā not to write anything, for
Āqā Muḥammad-ʿAlī-i-Iṣfahānī asked him about a Persian verse, and he
failed to understand its meaning.'' The Ancient Beauty said: ``Siyyid,
what business do you have with such meddling?'' At last He dismissed
him.

\subsection{196}\label{section-2171}

\begin{otherlanguage}{persian}

بعد رفته میرزا یحیی را اغوی نموده و در ظاهر نزد بعضی اظهار ایمان باو
میکرد مع آنکه و نفس الحق روح من فی ملکوت الامر و الخلق فداه مکرر در سبیل
مابینشان در بعضی امور گفتگو شد بالاخره بمعارضه انجامید و تلقاء وجه هر دو
حاضر و از یکدیگر شکایت آغاز نمودند این خادم فانی دوست ندارد تفصیل آن
نفوس را عرض نماید.

\end{otherlanguage}

Baʿd rafte Mīrzā Yaḥyā rā eghvā namūde va dar ẓāher nazd-e baʿżī eẓhār-e
īmān be-ū mīkard, maʿ-e ānke, va Nafs-e Ḥaqq, rūḥ-e man fī malakūt
al-amr va al-khalq fedāh, mokarrar dar sabīl mābeynashān dar baʿżī omūr
goftogū shod. Bel-ākhare be-moʿāreże anjāmīd va telqāʾ-e vajh har do
ḥāżer va az yekdīgar shekāyat āghāz namūdand. Īn khādem-e fānī dūst
nadārad tafṣīl-e ān nofūs rā ʿarż namāyad.

He then went and misled Mīrzā Yaḥyā, outwardly professing faith in him
before certain people, although---by the Self of the True One, may the
souls of all in the kingdoms of revelation and creation be sacrificed
for Him---during the journey the two repeatedly argued over certain
matters. This finally led to confrontation, and both came into His
presence and began complaining about one another. This transitory
servant has no wish to give a detailed account of those souls.

\subsection{197}\label{section-2172}

\begin{otherlanguage}{persian}

ای برادران و اللّه الذی لا اله الّا هو بصدق میگویم مقصودی ندارم جز آنکه
مجدد حجبات غلیظهٴ موهومه و سبحات کذبه عباد بیچاره را منع ننماید و محروم
نسازد شخص مذکور با هفتاد نفر در این سبیل لیلا و نهارا معاشر معذلک باطراف
نوشته من با احدی معاشر نبودم و متوهمین قبول نمودند این است شأن ناس از حق
میطلبم اولیای خود را تائید فرماید تا بمثابه قبل ابواب کذب مفتوح نشود و
مدن و دیار موهومه تعمیر نیابد

\end{otherlanguage}

Ey barādarān, va Allāh allaẕī lā ilāha illā huwa, be-ṣedq mīgūyam
maqṣūdī nadāram joz ānke mojaddad ḥojobāt-e ghalīẓe-ye mowhūme va
soboḥāt-e kaẕabe ʿebād-e bīchāre rā manʿ nanamāyad va maḥrūm nasāzad.
Shakhṣ-e maẕkūr bā haftād nafar dar īn sabīl laylan va nahāran moʿāsher,
maʿa-ẕālek be-aṭrāf neveshte man bā aḥadī moʿāsher nabūdam va
motavahhemīn qabūl namūdand. Īn ast shaʾn-e nās. Az Ḥaqq mīṭalabam
owliyā-ye khod rā taʾyīd farmāyad tā be-mathābe-ye qabl abvāb-e keẕb
maftūḥ nashavad va modon va diyār-e mowhūme taʿmīr nayābad.

Brothers! By God, besides whom there is no other God, I speak
truthfully: I have no purpose except that thick veils of illusion and
false splendors should not once again obstruct and deprive these
helpless servants. The person mentioned associated with seventy people
day and night during this journey; nevertheless, he wrote to various
places: ``I associated with no one,'' and those caught in delusion
accepted it. Such is the condition of the people. I ask God to
strengthen His loved ones, so that the doors of falsehood may not be
opened as before, nor imaginary cities and lands be built up.

\subsection{198}\label{section-2173}

\begin{otherlanguage}{persian}

بعد از فصل اکبر سید مذکور باغوای چند نفر از اعجام خارج که در آن مدینه
بودند قیام نمود و لیلا و نهارا با آن نفوس معاشر تا آنکه در حضور آن نفوس
ذکر نمود که قرار بر مباهله گذاشته‌ایم و ذکر این فقره نظر باطمینانی بود که
داشت که جمال قدم بامثال آن نفوس اعتنا نمیفرمایند و البته با آن نفوس در
یک مجلس جمع نمیشوند باینجهة با کمال یقین نزد اعجام باین کلمات وهمیه
مشغول شد

\end{otherlanguage}

Baʿd az faṣl-e akbar Sayyid-e madhkūr be-ighvā-ye chand nafar az aʿjām-e
khārej ke dar ān madīne būdand qiyām namūd va laylan va nahāran bā ān
nufūs muʿāsher tā ānke dar ḥuḍūr-e ān nufūs dhekr namūd ke qarār bar
mubāhale godhāshte-īm va dhekr-e īn feqre naẓar be-eṭmīnānī būd ke dāsht
ke Jamāl-e Qedam be-amthāl-e ān nufūs eʿtenā nemīfarmāyand va albatte bā
ān nufūs dar yek majles jamʿ nemīshavand be-īn jehat bā kamāl-e yaqīn
nazd-e aʿjām be-īn kalamāt-e vahmiyye mashghūl shod.

After the Most Great Separation, the aforementioned Sayyid arose at the
instigation of several Persians outside the community who were in that
city, and associated with them day and night. Eventually, in their
presence, he stated, ``We have agreed to a mutual invocation of God's
judgment.'' He mentioned this because he was confident that the Ancient
Beauty would pay no attention to such people and would certainly not
join them in a gathering. Thus, with complete assurance, he occupied
himself among the Persians with these fanciful statements.

\subsection{199}\label{section-2174}

\begin{otherlanguage}{persian}

در آن بین میر محمد نام که در ارض شین خدمت نقطه اولی روح ما سواه فداه
فائز شده و در سفر مدینه کبیره هم ملتزم رکاب بود رسید بر گفتگوهای سید بین
اعجام مطلع شد گفت وقت را معین کنید من میروم و در ساحت اقدس عرض مینمایم
تشریف بیاورند تا حق از باطل واضح و مبرهن گردد قرار شد در مسجد سلطان سلیم
این امر واقع شود

\end{otherlanguage}

Dar ān bayn Mīr Muḥammad nām ke dar arḍ-e Shīn khedmat-e Noqṭe-ye Ūlā
rūḥu mā siwāhu fidāh fāʾez shode va dar safar-e Madīne-ye Kabīre ham
multazem-e rekāb būd resīd bar goftogūhā-ye Sayyid bayn-e aʿjām moṭṭaleʿ
shod goft vaqt rā moʿayyan konīd man mīravam va dar sāḥat-e aqdas ʿarḍ
mīnamāyam tashrīf beyāvarand tā ḥaqq az bāṭel vāḍeḥ va mobarhan gardad
qarār shod dar masjed-e Sulṭān Salīm īn amr vāqeʿ shavad.

Meanwhile, a man named Mīr Muḥammad arrived. He had attained the
presence of the Primal Point---may the souls of all besides Him be
sacrificed for Him---in the Land of Shīn, and had also accompanied Him
on the journey to the Great City. He learned of the Sayyid's statements
among the Persians and said, ``Set the time. I will go and ask in the
Most Holy Court that He come, so that truth may be clearly and
conclusively distinguished from falsehood.'' It was agreed that this
should take place in the Mosque of Sulṭān Salīm.

\subsection{200}\label{section-2175}

\begin{otherlanguage}{persian}

میر مذکور وارد و تفصیل را نزد بعضی از دوستان ذکر نمود تا آنکه در پیشگاه
حضور عرض شد این فانی آن حین حاضر نبود بمجرد اصغا جمال قدم از مقر بر
خاستند و فردا واحدا متوجها علی مقر المعلوم توجه فرمودند و در عرض راه
بشأنی فرات آیات الهی جاری و نازل که احدی قادر بر وصف آن نه

\end{otherlanguage}

Mīr-e madhkūr vāred va tafṣīl rā nazd-e baʿḍī az dūstān dhekr namūd tā
ānke dar pīshgāh-e ḥuḍūr ʿarḍ shod īn fānī ān ḥīn ḥāḍer nabūd
be-mojarrad-e eṣghā Jamāl-e Qedam az maqarr bar khāstand va fardan
wāḥidan mutawajjihan ʿalā maqarri al-maʿlūm tavajjoh farmūdand va dar
ʿarḍ-e rāh be-shaʾnī Forāt-e āyāt-e elāhī jārī va nāzel ke aḥadī qāder
bar vaṣf-e ān na.

The aforementioned Mīr arrived and related the details to some of the
friends, until the matter was submitted in His presence. This transitory
servant was not present at that moment. As soon as He heard it, the
Ancient Beauty rose from His seat and, alone and unaccompanied, set out
toward the appointed place. Along the way, the Euphrates of divine
verses flowed and descended in such a manner that no one could describe
it.

\subsection{201}\label{section-2176}

\begin{otherlanguage}{persian}

و این عبد در آن حین در سوق بوده بجهة تمشیت بعض امور بیت بعد از استماع
بسرعت تمام توجه نمودم در اثنای راه ملاحظه شد جمعی اهل بلد از طرفین معبر
متحیر ایستاده‌اند

\end{otherlanguage}

Va īn ʿabd dar ān ḥīn dar sūq būde be-jehat-e tamshīyat-e baʿḍ-e omūr-e
bayt baʿd az estemāʿ be-sorʿat-e tamām tavajjoh namūdam dar athnā-ye rāh
molāḥaẓe shod jamʿī ahl-e balad az ṭarafayn-e maʿbar motaḥayyer
īstāde-and.

At that moment this servant was in the market attending to some
household affairs. After hearing the news, I set out with all speed.
Along the way, I saw a group of townspeople standing in astonishment on
both sides of the road.

\subsection{202}\label{section-2177}

\begin{otherlanguage}{persian}

این عبد را که دیدند با دست اشاره نمودند یعنی از این سمت تشریف بردند تا
اینکه این عبد وارد جامع شد مشاهده نمود بحر آیات مواج لعمر مقصودی و مقصود
من فی السموات و الارض بشأنی آیات نازل که کل متحیر که آیا چه واقع شده یکی
از اعجام مشاهده شد که میلرزید و میگریست و در حضور از اصحاب جز میر محمد و
این عبد که از بعد رسید احدی حاضر نه چه که کل را از حضور نهی فرمودند

\end{otherlanguage}

Īn ʿabd rā ke dīdand bā dast eshāre namūdand yaʿnī az īn samt tashrīf
bordand tā īnke īn ʿabd vāred-e jāmeʿ shod moshāhede namūd baḥr-e āyāt
mavvāj la-ʿamri maqṣūdī wa-maqṣūdi man fī al-samāwāti wa-al-arḍ
be-shaʾnī āyāt nāzel ke koll motaḥayyer ke āyā che vāqeʿ shode yekī az
aʿjām moshāhede shod ke mīlarzīd va mīgerīst va dar ḥuḍūr az aṣḥāb joz
Mīr Muḥammad va īn ʿabd ke az baʿd resīd aḥadī ḥāḍer na che ke koll rā
az ḥuḍūr nahy farmūdand.

When they saw this servant, they gestured with their hands to indicate,
``He went this way.'' When this servant entered the congregational
mosque, he saw the ocean of verses surging. By the life of Him who is my
Desire and the Desire of all in the heavens and on earth! The verses
were descending in such a manner that everyone was bewildered, wondering
what had happened. One of the Persians was seen trembling and weeping.
None of the companions was in His presence except Mīr Muḥammad and this
servant, who arrived afterward, for He had forbidden them all to attend.

\subsection{203}\label{section-2178}

\begin{otherlanguage}{persian}

بعد وجه مبارک بمیر محمد متوجه قال یا محمد اذهب الیهم و قل تعالوا بحبالکم
و عصیکم میر بسرعت تمام رفت بعد از مراجعت عرض نمود حضرات از گفته نادم
شدند و عذر آوردند که امروز ممکن نیست فرمودند بگو تا سه یوم مهلت هر یوم
حاضر شوید ما حاضریم بعد ساعتی در جامع بصلوة مشغول یعنی صلاة بدیعی و بعد
از مراجعت بمجرد ورود در بیت یک لوح بخط مبارک نازل و بمهر مبارک مختوم که
تا سه یوم هر یوم بیائید ما حاضریم و آن لوح رانزد اصحاب فرستادند تا کل
مطلع باشند و جناب میر هم رفت و تفصیل لوح را اخبار نمود ولکن از آن نفوس
اثری ظاهر نشد

\end{otherlanguage}

Baʿd vajh-e mobārak be-Mīr Muḥammad motavajjeh qāla yā Muḥammad idhhab
ilayhim wa-qul taʿālaw bi-ḥibālikum wa-ʿiṣiyyikum Mīr be-sorʿat-e tamām
raft baʿd az morājaʿat ʿarḍ namūd ḥaḍarāt az gofte nādem shodand va
ʿodhr āvardand ke emrūz momken nīst farmūdand begū tā se yawm mohlat har
yawm ḥāḍer shavīd mā ḥāḍerīm baʿd sāʿatī dar jāmeʿ be-ṣalāt mashghūl
yaʿnī ṣalāt-e badīʿī va baʿd az morājaʿat be-mojarrad-e vorūd dar bayt
yek lawḥ be-khaṭṭ-e mobārak nāzel va be-mohr-e mobārak makhtūm ke tā se
yawm har yawm beyāyīd mā ḥāḍerīm va ān lawḥ rā nazd-e aṣḥāb ferestādand
tā koll moṭṭaleʿ bāshand va Janāb-e Mīr ham raft va tafṣīl-e lawḥ rā
ekhbār namūd va lāken az ān nufūs atharī ẓāher nashod.

Then He turned His blessed face toward Mīr Muḥammad and said,
``Muḥammad, go to them and say, `Come with your ropes and your
rods.'\,'' Mīr went with all speed. On returning, he reported, ``The
gentlemen have regretted what they said and have offered the excuse that
today is not possible.'' He said, ``Tell them they have three days'
grace. Whichever day you present yourselves, We are ready.'' He then
spent a while in prayer in the mosque---a wondrous prayer. After
returning, as soon as He entered the house, a Tablet was revealed in His
own blessed handwriting and sealed with His blessed seal, stating,
``Come on any day within three days; We are ready.'' He sent that Tablet
to the companions so that all would be informed. Mīr also went and
conveyed the details of the Tablet, but no sign of those people
appeared.

\subsection{204}\label{section-2179}

\begin{otherlanguage}{persian}

ای آقایان تفکر نمائید امر بچه شأن ظاهر و ناس بچه شأن غافل غیر اهل آن
مدینه قریب صد نفر از مهاجرین و مسافرین جمع بودند و کل بر این مطلب شاهد و
گواه و آن لوح مبارک موجود و سوره مباهله در تفصیل آن یوم نازل و حسب الامر
نزد یکی از اولیا ارسال شد که بر احبا و اولیای الهی قرائت نماید و البته
اکثری باصغاء آن فائز شده‌اند و بعضی از آن در این مقام ثبت میشود تا امکان
از آیات رحمن منور گردد و افئده عالم بطراز عدل و انصاف مزین شود الامر بید
اللّه هذا ما نزل من جبروت مشیة ربّنا العزیز العلام قوله جل جلاله ان

\end{otherlanguage}

Ay āqāyān tafakkor namāyīd amr be-che shaʾn ẓāher va nās be-che shaʾn
ghāfel ghayr-e ahl-e ān madīne qarīb-e ṣad nafar az mohājerīn va
mosāferīn jamʿ būdand va koll bar īn maṭlab shāhed va govāh va ān lawḥ-e
mobārak mawjūd va Sūre-ye Mubāhale dar tafṣīl-e ān yawm nāzel va ḥasb
al-amr nazd-e yekī az awliyā ersāl shod ke bar aḥebbā va awliyā-ye elāhī
qerāʾat namāyad va albatte aktharī be-eṣghā-ye ān fāʾez shode-and va
baʿḍī az ān dar īn maqām thabt mīshavad tā emkān az āyāt-e Raḥmān
monavvar gardad va afʾede-ye ʿālam be-ṭerāz-e ʿadl va enṣāf mozayyan
shavad al-amru bi-yadi Allāh hādhā mā nazala min jabarūti mashīʾati
Rabbinā al-ʿAzīzi al-ʿAllām qawluhu jalla jalāluhu in.

Gentlemen, reflect on how manifest the Cause is and how heedless the
people are. Apart from the inhabitants of that city, nearly a hundred
emigrants and travelers were gathered there, and all bear witness and
testify to this matter. That blessed Tablet still exists. The Sūrih of
the Mutual Invocation was revealed, detailing that day, and was sent, as
instructed, to one of the faithful to read to the loved ones and friends
of God. Surely most have had the privilege of hearing it. A portion of
it is recorded here, so that the world of existence may be illumined by
the verses of the All-Merciful and the hearts of the world adorned with
justice and fairness. Authority rests in God's hand. This is what has
descended from the dominion of the will of our Lord, the Almighty, the
All-Knowing. His words---exalted be His glory---are: ``Indeed,

\subsection{205}\label{section-2180}

\begin{otherlanguage}{arabic}

یا قلم الاعلی اذکر لاسمک نبا هذه المدینة لیکون من المطلعین اذ جائک محمد
من لدی المشرکین و قال انهم ارادوا ان یجتمعوا معک فی مقر آخر فلما سمعنا
خرجنا عن البیت بسلطان مبین

\end{otherlanguage}

Yā Qalama al-Aʿlā udhkur li-ismika nabaʾa hādhihi al-madīnati li-yakūna
mina al-muṭṭaliʿīn idh jāʾaka Muḥammadun min ladā al-mushrikīna wa-qāla
innahum arādū an yajtamiʿū maʿaka fī maqarrin ākhar fa-lammā samiʿnā
kharajnā ʿani al-bayti bi-sulṭānin mubīn.

O Most Exalted Pen, recount to your namesake the news of this city, that
he may be informed. When Muḥammad came to you from those who associate
partners with God and said, `They wish to meet with You in another
place,' We left the house with manifest authority as soon as We heard.

\subsection{206}\label{section-2181}

\begin{otherlanguage}{arabic}

و قلنا یا محمد خرج الروح من مقره و خرجت معه ارواح الاصفیاء ثمّ حقایق
المرسلین انظر لتری الملأ الاعلی فوق رأسی و فی قبضتی حجج النبیین افتح
عینک هذا علی ثمّ محمد رسول اللّه قد طلع من افق البیت بسلطان مبین قل انه
اتی برایات الایات من لدی اللّه المقتدر القدیر

\end{otherlanguage}

Wa-qulnā yā Muḥammad kharaja al-rūḥu min maqarrihi wa-kharajat maʿahu
arwāḥu al-aṣfiyāʾ thumma ḥaqāʾiqu al-mursalīn unẓur li-tarā al-malaʾa
al-aʿlā fawqa raʾsī wa-fī qabḍatī ḥujaju al-nabiyyīn iftaḥ ʿaynaka hādhā
ʿAliyyun thumma Muḥammadun Rasūlu Allāh qad ṭalaʿa min ufuqi al-bayti
bi-sulṭānin mubīn qul innahu atā bi-rāyāti al-āyāti min ladā Allāhi
al-Muqtadiri al-Qadīr.

And We said, `Muḥammad, the Spirit has left His dwelling, and with Him
have come forth the souls of the chosen ones and the realities of the
Messengers. Look, that you may see the Concourse on High above My head,
while in My grasp are the proofs of the Prophets. Open your eyes: this
is ʿAlī, and then Muḥammad, the Messenger of God, who has appeared above
the horizon of the house with manifest authority. Say: He has come with
the banners of the verses from God, the All-Powerful, the Omnipotent.'

\subsection{207}\label{section-2182}

\begin{otherlanguage}{arabic}

قل لو اجتمع من علی الارض من العلماء و العرفأ ثم الملوک و السلاطین انّی
احضر امام وجوههم ناطقا آیات اللّه الملک العزیز الحکیم انا الذی لا اخاف من
احد لو یجتمع علی ضری من فی السموات و الارضین قد خضعت الایات لوجهی و خشعت
الاصوات لندائی العزیز البدیع

\end{otherlanguage}

Qul law ijtamaʿa man ʿalā al-arḍi mina al-ʿulamāʾi wa-al-ʿurafāʾ thumma
al-mulūki wa-al-salāṭīn innī aḥḍuru amāma wujūhihim nāṭiqan āyāti Allāhi
al-Maliki al-ʿAzīzi al-Ḥakīm anā alladhī lā akhāfu min aḥadin law
yajtamiʿu ʿalā ḍurrī man fī al-samāwāti wa-al-arḍīn qad khaḍaʿati
al-āyātu li-wajhī wa-khashaʿati al-aṣwātu li-nidāʾī al-ʿAzīzi al-Badīʿ.

Say: Were all the learned and the mystics on earth, and then the kings
and sovereigns, to gather together, I would appear before them,
proclaiming the verses of God, the King, the Almighty, the All-Wise. I
am He who fears no one, even if all in the heavens and the earths should
unite to harm Me. The verses have bowed before My face, and voices have
been hushed before My mighty, wondrous call.

\subsection{208}\label{section-2183}

\begin{otherlanguage}{arabic}

قل هذه یدی قد جعلها اللّه بیضاء للناظرین و هی عصای لو نلقیها لتبلع الارض
و من علیها و انا العلیم الخبیر

\end{otherlanguage}

Qul hādhihi yadī qad jaʿalahā Allāhu bayḍāʾa li-al-nāẓirīn wa-hiya
ʿaṣāya law nulqīhā li-tablaʿa al-arḍa wa-man ʿalayhā wa-anā al-ʿAlīmu
al-Khabīr.

Say: This is My hand, which God has made white for all to behold, and it
is My rod: were We to cast it down, it would swallow the earth and all
who are upon it. I am the All-Knowing, the All-Aware.

\subsection{209}\label{section-2184}

\begin{otherlanguage}{arabic}

یا محمد اذهب الی المشرکین برسالات اللّه رب العالمین قل قد جائکم الغلام و
معه جنود الوحی و الالهام یمشی امامه الملکوت و عن ورائه قبائل الملأ
الاعلی و اهل مدائن الاسماء ثمّ ملائکة المقربون

\end{otherlanguage}

Yā Muḥammad idhhab ilā al-mushrikīna bi-risālāti Allāhi Rabbi al-ʿālamīn
qul qad jāʾakumu al-ghulāmu wa-maʿahu junūdu al-waḥyi wa-al-ilhām yamshī
amāmahu al-malakūtu wa-ʿan warāʾihi qabāʾilu al-malaʾi al-aʿlā wa-ahlu
madāʾini al-asmāʾ thumma malāʾikatu al-muqarrabūn.

Muḥammad, go to those who associate partners with God with the messages
of God, the Lord of the worlds. Say: The Youth has come to you, and with
Him are the hosts of revelation and inspiration. The Kingdom walks
before Him, and behind Him are the tribes of the Concourse on High, the
inhabitants of the cities of names, and the angels who are brought near.

\subsection{210}\label{section-2185}

\begin{otherlanguage}{arabic}

قل خافوا اللّه و لا تعترضوا علی الذی بامره قدر کل امر حکیم قل تعالوا
بحبالکم و عصیکم و ما عندکم لو انتم من القادرین انا نذهب الی بیت اللّه
الذی بناه احد من الملوک وسمی بالسلیم و اتوقف هناک الی ان تغرب الشمس فی
مغربها لیتم علیهم حجة اللّه و برهانه و یکونن من التائبین

\end{otherlanguage}

Qul khāfū Allāha wa-lā taʿtariḍū ʿalā alladhī bi-amrihi quddira kullu
amrin ḥakīm qul taʿālaw bi-ḥibālikum wa-ʿiṣiyyikum wa-mā ʿindakum law
antum mina al-qādirīn anā nadhhabu ilā bayti Allāhi alladhī banāhu
aḥadun mina al-mulūki wa-summiya bi-al-Salīm wa-atawaqqafu hunāka ilā an
taghruba al-shamsu fī maghribihā li-yatimma ʿalayhim ḥujjatu Allāhi
wa-burhānuhu wa-yakūnanna mina al-tāʾibīn.

Say: Fear God, and do not oppose Him by whose command every wise decree
has been ordained. Say: Come with your ropes, your rods, and whatever
you possess, if you are able. We are going to the house of God which one
of the kings built and which is called Salīm, and I shall remain there
until the sun sets in its place of setting, so that God's proof and
demonstration may be made complete for them and they may be among those
who repent.

\subsection{211}\label{section-2186}

\begin{otherlanguage}{arabic}

قل قد انقطعت الیوم نسبة کل ذی نسبة الّا من دخل فی ظل هذا الامر المشرق
المنیر من تمسک باسم من الاسماء معرضا عن موجدها انه لیس منی و کان اللّه
بریئا منه انا خلقنا الاسماء کخلق کل شئ ان انتم من العارفین

\end{otherlanguage}

Qul qad inqaṭaʿati al-yawma nisbatu kulli dhī nisbatin illā man dakhala
fī ẓilli hādhā al-amri al-mushriqi al-munīr man tamassaka bi-ismin mina
al-asmāʾi muʿriḍan ʿan mūjidihā innahu laysa minnī wa-kāna Allāhu
barīʾan minhu anā khalaqnā al-asmāʾa ka-khalqi kulli shayʾin in antum
mina al-ʿārifīn.

Say: Today every bond of affiliation has been severed, except for those
who have entered beneath the shadow of this radiant, luminous Cause.
Whoever clings to any name while turning away from the One who brought
the names into being is not of Me, and God disavows him. We created the
names just as We created all things, if you are among those who
understand.

\subsection{212}\label{section-2187}

\begin{otherlanguage}{arabic}

قل انّی قد اظهرت نفسی لله و خرجت لله و اتکلم امام من علی الارض کما اتکلم
امام وجهک الا اخاف من احد یشهد بذلک عملی لو انت من العارفین

\end{otherlanguage}

Qul innī qad aẓhartu nafsī li-llāhi wa-kharajtu li-llāhi wa-atakallamu
amāma man ʿalā al-arḍi kamā atakallamu amāma wajhika alā akhāfu min
aḥadin yashhadu bi-dhālika ʿamalī law anta mina al-ʿārifīn.

Say: I have manifested Myself for God, and I have come forth for God. I
speak before all who are on earth just as I speak before you. Truly, I
fear no one. My actions bear witness to this, if you are among those who
understand.

\subsection{213}\label{section-2188}

\begin{otherlanguage}{arabic}

قل ان آثار الحق ظاهرة کضیاء الشمس و انوارها تاللّه انها اظهر عند کل ذی
بصر حدید هل تظن دونه یقدر ان یقوم معه لا و نفسه العلی العظیم ان الذی
اعتکفوا علی الاسماء اولئک من عبدة الاصنام لو انت من الشاهدین

\end{otherlanguage}

Qul inna āthāra al-Ḥaqqi ẓāhiratun ka-ḍiyāʾi al-shamsi wa-anwārihā
ta-Allāhi innahā aẓharu ʿinda kulli dhī baṣarin ḥadīd hal taẓunnu dūnahu
yaqdiru an yaqūma maʿahu lā wa-nafsihi al-ʿAliyyi al-ʿAẓīm inna alladhī
iʿtakafū ʿalā al-asmāʾi ulāʾika min ʿabadati al-aṣnāmi law anta mina
al-shāhidīn.

Say: The signs of the True One are as manifest as the sun's radiance and
light. By God, they are even more manifest to everyone with keen sight.
Do you imagine that anyone besides Him can stand alongside Him? No, by
His own exalted and mighty Self! Those who have devoted themselves to
names are worshippers of idols, if you are among those who see.

\subsection{214}\label{section-2189}

\begin{otherlanguage}{persian}

ان النبی من اخبر الناس بهذا النبأ الاقوم الاقدم القدیم و الرسول من بلغ
رسالاتی و هدی الذین ضلوا السبیل و الامام من قام امام وجهی بخضوع مبین و
الولی من استحصن فی حصن ولایتی المحکم المتین و الوصی من وصی النّاس بما
اتاه فی لوح حفیظ

\end{otherlanguage}

Inna al-nabiyya man akhbara al-nāsa bi-hādhā al-nabaʾi al-aqwami
al-aqdami al-qadīm wa-al-rasūlu man ballagha risālātī wa-hadā alladhīna
ḍallū al-sabīl wa-al-imāmu man qāma amāma wajhī bi-khuḍūʿin mubīn
wa-al-waliyyu man istaḥṣana fī ḥiṣni wilāyatī al-muḥkami al-matīn
wa-al-waṣiyyu man waṣṣā al-nāsa bi-mā atāhu fī lawḥin ḥafīẓ.

The Prophet is the one who has informed the people of this most upright,
most ancient, eternal Announcement. The Messenger is the one who has
conveyed My messages and guided those who have lost the way. The Imam is
the one who has stood before My face in manifest humility. The guardian
is the one who has taken refuge in the firm and mighty stronghold of My
guardianship. The appointed successor is the one who has enjoined upon
the people what has come to him in a preserved Tablet.

\subsection{215}\label{section-2190}

\begin{otherlanguage}{arabic}

ان الذی اتخذه المشرکون ربا لانفسهم و اعرضوا به علی اللّه المقتدر العزیز
الجمیل انه کان کاحد من العباد و کان ان یحضر تلقاء الوجه فی کل بکور و
اصیل هبت علیه من شطر نفسه اریاح الهوی بما اغواه احد من الشیاطین انا
شهرنا اسمه بین العباد حکمة من لدنا ان ربک هو الحاکم علی ما یرید ''
انتهی.

\end{otherlanguage}

Inna alladhī ittakhadhahu al-mushrikūna rabban li-anfusihim wa-aʿraḍū
bihi ʿalā Allāhi al-Muqtadiri al-ʿAzīzi al-Jamīl innahu kāna ka-aḥadin
mina al-ʿibādi wa-kāna an yaḥḍura tilqāʾa al-wajhi fī kulli bukūrin
wa-aṣīl habbat ʿalayhi min shaṭri nafsihi aryāḥu al-hawā bi-mā aghwāhu
aḥadun mina al-shayāṭīn anā shahharnā ismahu bayna al-ʿibādi ḥikmatan
min ladunnā inna rabbaka huwa al-ḥākimu ʿalā mā yurīd intahā.

The one whom those who associate partners with God have taken as a lord
for themselves, and on whose account they have turned against God, the
All-Powerful, the Almighty, the All-Beautiful, was but one of the
servants, and used to appear before the Face every morning and evening.
The winds of desire blew upon him from his own self, because one of the
devils led him astray. We made his name renowned among the servants
through a wisdom of Our own. Your Lord is indeed the One who decrees
whatever He wills.'' End of quotation.

\subsection{216}\label{section-2191}

\begin{otherlanguage}{arabic}

ای دوستان صاحب بصر و فؤاد را یک ذکر از روی حقیقت کفایت مینماید در این
حین بحضور فائز لسان عظمت باین کلمه ناطق قوله عز بیانه یا عبد حاضر بگو ای
بی انصافها آیا امواج بحر بیان را مشاهده ننموده‌اید و یا اشراقات انوار
آفتاب ظهور را ندیده‌اید اتقوا اللّه از غدیر ببحر و از وهم بیقین توجه
نمائید امروز جمیع کتب او بلا و نعم معلق و منوط هذا بدیع السموات و الارض
لو انتم تعرفون بدیع من کل جهات این ظهور بوده و هست طوبی للعارفین

\end{otherlanguage}

Ay dūstān ṣāḥeb-e baṣar va fuʾād rā yek dhekr az rū-ye ḥaqīqat kefāyat
mīnamāyad dar īn ḥīn be-ḥuḍūr fāʾez lesān-e ʿaẓamat be-īn kalame nāṭeq
qawluhu ʿazza bayānuhu yā ʿabd-e ḥāḍer begū ay bī-enṣāfhā āyā amvāj-e
baḥr-e bayān rā moshāhede nanamūde-īd va yā eshrāqāt-e anvār-e āftāb-e
ẓuhūr rā nadīde-īd ittaqū Allāh az ghadīr be-baḥr va az vahm be-yaqīn
tavajjoh namāyīd emrūz jamīʿ-e kotob-e ū be-lā va naʿam moʿallaq va
manūṭ hādhā badīʿu al-samāwāti wa-al-arḍi law antum taʿrifūn badīʿ men
koll-e jehāt īn ẓuhūr būde va hast ṭūbā li-al-ʿārifīn.

Friends, a single truthful reminder suffices for one possessed of sight
and understanding. At this moment I attained His presence, and the
Tongue of Grandeur spoke these words---exalted be His utterance:
``Servant who is present, say: You who are unfair! Have you not beheld
the waves of the ocean of utterance, or have you not seen the radiance
of the lights of the sun of Revelation? Fear God. Turn from the pool to
the ocean, and from illusion to certainty. Today all His Books depend
upon and are subject to `No' and `Yes.' This is the Originator of the
heavens and the earth, if you understand. This Revelation has been and
remains wondrous in every respect. Blessed are those who recognize it.''

\subsection{217}\label{section-2192}

\begin{otherlanguage}{arabic}

قل لیت افکاری افکارکم و لا امشی فی طرقکم قد اوضحت سبیلی و اظهرت دلیلی و
انزلت آیاتی و دعوت العباد الی افقی منهم من اقبل و منهم من اعرض ان الذی
اقبل انه من الفائزین فی کتاب مبین '' انتهی.

\end{otherlanguage}

Qul layta afkārī afkārukum wa-lā amshī fī ṭuruqikum qad awḍaḥtu sabīlī
wa-aẓhartu dalīlī wa-anzaltu āyātī wa-daʿawtu al-ʿibāda ilā ufuqī minhum
man aqbala wa-minhum man aʿraḍa inna alladhī aqbala innahu mina
al-fāʾizīna fī kitābin mubīn intahā.

``Say: Would that My thoughts were your thoughts! I do not walk in your
ways. I have made My path clear, manifested My proof, sent down My
verses, and summoned the servants to My horizon. Some have turned toward
Me, and others have turned away. Whoever has turned toward Me is indeed
numbered among those who have attained success in a clear Book.'' End of
quotation.

\subsection{218}\label{section-2193}

\begin{otherlanguage}{persian}

ایکاش معرضین یک ورق از کتاب ایقان را ببصر عدل نظر مینمودند و بانصاف حکم
میکردند معادله نمینماید بآن آنچه باین شأن از قبل نازل شده امروز کل در
صقع واحد مشاهده میشوند هر نفسی بحبل متین تمسک نمود و بذیل منیر تشبث جست
او از اهل بها در صحیفه حمراء از قلم اعلی مسطور امروز جمیع ذرات بر عظمت
امر شاهد و گواهند و معرضین غافل و گمراه

\end{otherlanguage}

Ay kāsh moʿreḍīn yek varaq az Ketāb-e Īqān rā be-baṣar-e ʿadl naẓar
mīnamūdand va be-enṣāf ḥokm mīkardand moʿādele nemīnamāyad be-ān ānche
be-īn shaʾn az qabl nāzel shode emrūz koll dar ṣoqʿ-e vāḥed moshāhede
mīshavand har nafsī be-ḥabl-e matīn tamassok namūd va be-dhayl-e monīr
tashabboth jost ū az ahl-e Bahā dar ṣaḥīfe-ye ḥamrā az Qalam-e Aʿlā
masṭūr emrūz jamīʿ-e dharrāt bar ʿaẓamat-e amr shāhed va govāhand va
moʿreḍīn ghāfel va gomrāh.

Would that those who turn away would look with the eye of justice at a
single page of the Book of Certitude and judge fairly. Nothing
previously revealed in this manner equals it. Today all are seen on a
single plane. Whoever has held fast to the firm cord and clung to the
luminous hem is inscribed by the Most Exalted Pen in the Crimson Scroll
as one of the people of Bahā. Today every atom bears witness and
testifies to the greatness of the Cause, while those who turn away
remain heedless and astray.

\subsection{219}\label{section-2194}

\begin{otherlanguage}{persian}

حضرت اعلی روح ما سواه فداه بر اثبات حقیتشان در آخر تفسیرها بشهادت دو
عالم استدلال فرموده‌اند ملا عبد الخالق و حاج ملا محمد علی برقانی قزوینی
میفرمایند قوله جل و عزّ و کفی بشادتهما فی حقی علی ذلک الامر شهیدا انتهی.

\end{otherlanguage}

Ḥażrat-e Aʿlā, rūḥu mā siwāhu fidāhu, bar ethbāt-e ḥaqqīyatashān dar
ākhar-e tafsīrhā be-shahādat-e do ʿālem estedlāl farmūde-and: Mollā ʿAbd
al-Khāleq va Ḥājj Mollā Moḥammad ʿAlī Barqānī Qazvīnī. Mīfarmāyand,
qawluhu jalla wa ʿazza: «Wa kafā {[}corrupted text{]} fī ḥaqqī ʿalā
dhālika al-amri shahīdan.» Intahā.

The Most Exalted One---may the spirit of all besides Him be sacrificed
for Him---appealed, at the end of His commentaries, to the testimony of
two scholars as proof of His truth: Mullā ʿAbd al-Khāliq and Ḥājj Mullā
Muḥammad ʿAlī Barqānī Qazvīnī. He says---exalted and mighty are His
words: ``And {[}corrupted original text{]} suffices as a witness
concerning Me in this matter.'' End of quotation.

\subsection{220}\label{section-2195}

\begin{otherlanguage}{persian}

مع آنکه بعد از ملاحظه لوح مبارک که میفرماید اننی انا القائم الذی انتم
بظهوره توعدون نفس اول اعراض نمود اعراض شدید و نفس ثانی هم من بعد اقبال
ننمود و با احدی تقرب نجست چه که این دو عریضه در اول امر عرض نمودند حال
ملاحظه نمائید چه مقدار از علما و عرفا و فقها نزد اشراق انوار آفتاب
اینظهور اعظم خاضع و خاشع مشاهده گشتند اگر اسامی ذکر شود بر حسب ظاهر بیم
ضرّ است از برای ایشان

\end{otherlanguage}

Maʿ ānke baʿd az molāḥeẓe-ye Lawḥ-e mobārak ke mīfarmāyad «Innanī anā
al-Qāʾimu alladhī antum bi-ẓuhūrihi tūʿadūn», nafs-e avval eʿrāż namūd,
eʿrāż-e shadīd, va nafs-e thānī ham men baʿd eqbāl nanamūd va bā aḥadī
taqarrub najost, che ke īn do ʿarīże dar avval-e amr ʿarż namūdand. Ḥāl
molāḥeẓe namāyīd che meqdār az ʿolamā va ʿorafā va foqahā nazd-e
eshrāq-e anvār-e āftāb-e īn Ẓohūr-e Aʿẓam khāżeʿ va khāsheʿ moshāhede
gashtand. Agar asāmī dhekr shavad, bar ḥasab-e ẓāher bīm-e żorr ast az
barā-ye īshān.

Yet, after reading the blessed Tablet in which He says, ``I am indeed
the Qāʾim whose appearance you have been promised,'' the first person
turned away, and did so vehemently; the second likewise did not turn
toward Him thereafter, nor did he seek proximity to anyone, for these
two had submitted their petitions at the beginning of the Cause. Now
consider how many scholars, mystics, and jurists have been seen humble
and submissive before the radiance of the Sun of this Most Great
Revelation. If their names were mentioned, there would, outwardly, be
reason to fear harm to them.

\subsection{221}\label{section-2196}

\begin{otherlanguage}{arabic}

هر منصفی گواهی میدهد عالم دانائی ازاعمال معرضین و اقوال منکرین متحیر
گاهی ذکر مرآت مینمایند و هنگامی ذکر وصی این ایام از کرمان الوهیت از برای
مرجع خود ترتیب داده‌اند و نفس دیگر از اهل کاظمین ساکن کرمانشاه مثل رومی و
چینی گفته تاللّه انهم و ما عندهم کبیوت العنکبوت بل اوهن لوهم ینصفون

\end{otherlanguage}

Har monṣefī gavāhī mīdahad. ʿĀlem-e dānāʾī az aʿmāl-e moʿreżīn va
aqvāl-e monkerīn motaḥayyer. Gāhī dhekr-e Merʾāt mīnamāyand va hengāmī
dhekr-e vaṣī. Īn ayyām az Kermān olūhīyat az barā-ye marjaʿ-e khod
tartīb dāde-and va nafs-e dīgar az ahl-e Kāẓemayn, sāken-e Kermānshāh,
mathal-e Rūmī va Chīnī gofte. Ta-llāhi innahum wa mā ʿindahum ka-buyūti
al-ʿankabūt, bal awhanu law hum yunṣifūn.

Every fair-minded person bears witness. A learned and knowledgeable
person is bewildered by the actions of those who turn away and the words
of the deniers. At times they speak of the Mirror, and at others of the
appointed successor. These days, people in Kirmān have contrived
divinity for their own leader, while another person, from Kāẓimayn and
residing in Kirmānshāh, has recounted the parable of the Byzantine and
the Chinese. By God, they and all they possess are like spiders'
webs---indeed, frailer still, if only they would judge fairly.

\subsection{222}\label{section-2197}

\begin{otherlanguage}{persian}

بگوئید ای بیچاره‌ها شما کجا بوده‌اید آیا غیر نفس حق کسی ظاهر بوده آیا از
دون او قدرتی مشاهده شده عجب است از نفوسی که از امر اول آگاهی نداشته‌اند و
فتوی میدهند

\end{otherlanguage}

Begūʾīd: Ey bīchāre-hā, shomā kojā būde-īd? Āyā gheyr-e nafs-e Ḥaqq kasī
ẓāher būde? Āyā az dūn-e Ū qodratī moshāhede shode? ʿAjab ast az nofūsī
ke az amr-e avval āgāhī nadāshte-and va fatvā mīdahand.

Say: You hapless people, where have you been? Has anyone other than the
True One Himself been manifest? Has any power been witnessed from anyone
besides Him? How astonishing that people who have had no knowledge of
the beginning of the Cause should issue rulings!

\subsection{223}\label{section-2198}

\begin{otherlanguage}{persian}

یا محمد لعمر اللّه غافلی خبر نداری این عبد شمارا ندیده و لکن لوجه اللّه
میگوید اتق اللّه بحق راجع شو تا باب مفتوح خود را محروم منما آن نفوسی که
تو را القا نموده‌اند خود آگاه نبوده‌اند و از بحر دانش نصیب برنداشته‌اند
نفسی که من غیر اللّه او را رب اخذ نموده آثارش نزدت موجود دیده و خوانده
حال یک صفحه آیات حق را قرائت نما و خود انصاف ده

\end{otherlanguage}

Yā Moḥammad, la-ʿamru Allāh, ghāfelī, khabar nadārī. Īn ʿabd shomā rā
nadīde va lāken li-wajhi Allāh mīgūyad: Ittaqi Allāh. Be-Ḥaqq rājeʿ
show. Tā bāb maftūḥ, khod rā maḥrūm manamā. Ān nofūsī ke to rā elqā
namūde-and khod āgāh nabūde-and va az baḥr-e dānesh naṣīb
barnadāshte-and. Nafsī ke min ghayri Allāh ū rā rabb akhdh namūde,
āthārash nazdat mowjūd, dīde va khānde. Ḥāl yek ṣafḥe āyāt-e Ḥaqq rā
qerāʾat namā va khod enṣāf deh.

O Muḥammad, by the life of God, you are heedless and unaware. This
servant has not seen you, but for the sake of God he says: Fear God.
Return to the True One. While the door is open, do not deprive yourself.
Those who have prompted you were themselves unaware and had gained no
share of the ocean of knowledge. The writings of the person you have
taken as a lord besides God are in your possession; you have seen and
read them. Now read one page of the verses of the True One and judge
fairly for yourself.

\subsection{224}\label{section-2199}

\begin{otherlanguage}{persian}

نقطه اولی روح ما سواه فداه میفرماید قوله جل و عزّ اگر کل اهل بیان در جوهر
علم مثل او گردند که ثمر نمی بخشد الا بتصدیق او فلتعتبرن ان یا اولی العلم
ثم ایاه تتقون انتهی.

\end{otherlanguage}

Noqṭe-ye Ūlā, rūḥu mā siwāhu fidāhu, mīfarmāyad, qawluhu jalla wa ʿazza:
«Agar koll-e ahl-e Bayān dar jowhar-e ʿelm methl-e ū gardand ke thamar
namībakhshad ellā be-taṣdīq-e Ū. Fa-ltaʿtabirunna in yā ulī al-ʿilm,
thumma iyyāhu tattaqūn.» Intahā.

The Primal Point---may the spirit of all besides Him be sacrificed for
Him---says, exalted and mighty are His words: ``Even if all the people
of the Bayān should become like him in the essence of knowledge, it
bears no fruit except through acknowledgment of Him. Take heed, then, O
possessors of knowledge, and fear Him.'' End of quotation.

\subsection{225}\label{section-2200}

\begin{otherlanguage}{persian}

ای برادر من اگر نفسی این امر را انکار نماید و از این ظهور اعظم غافل شود
قادر بر اثبات هیچ امری نبوده و نیست این فانی را حزن از کل جهات بشأنی
احاطه نموده که از اظهارش عاجز است چه که از برای شمس حقیقت بعد از ظهور و
اشراق دلیل ذکر مینماید و این خطائی است کبیر باعلی النداء میفرماید انّی
اکون مقدسا عن کل ما ذکر و سطر ان دلیلی ما ظهر من عندی و حجتی نفسی لو
انتم من العارفین انتهی.

\end{otherlanguage}

Ey barādar-e man, agar nafsī īn Amr rā enkār namāyad va az īn Ẓohūr-e
Aʿẓam ghāfel shavad, qāder bar ethbāt-e hīch amrī nabūde va nīst. Īn
fānī rā ḥozn az koll-e jehāt be-shaʾnī eḥāṭe namūde ke az eẓhārash ʿājez
ast, che ke az barā-ye Shams-e Ḥaqīqat baʿd az ẓohūr va eshrāq dalīl
dhekr mīnamāyad va īn khaṭāʾī ast kabīr. Be-aʿlā al-nedā mīfarmāyad:
«Innī akūnu muqaddasan ʿan kulli mā dhukira wa suṭira. Inna dalīlī mā
ẓahara min ʿindī wa ḥujjatī nafsī, law antum mina al-ʿārifīn.» Intahā.

O my brother, anyone who denies this Cause and remains heedless of this
Most Great Revelation has never been, and is not now, capable of proving
any cause. Sorrow has surrounded this mortal one on every side to such a
degree that he is powerless to express it, for he is adducing evidence
for the Sun of Truth after its appearance and radiance, and this is a
grave error. He proclaims in the loudest voice: ``I am sanctified above
all that has been mentioned and written. My proof is what has appeared
from Me, and My testimony is My own Self, if you are among those who
understand.'' End of quotation.

\subsection{226}\label{section-2201}

\begin{otherlanguage}{persian}

ولکن این عبد چون بعضی را ضعیف مشاهده نموده بحبل استدلال تمسک جسته که
شاید از معین کلمات الهی کوثر باقی بیاشامند و بزندگی دائمی فائز گردند و
الا ان الدلیل یستحیی ان ینسب الیه و البرهان یخجل ان یسجد بین یدیه نعم ما
قیل :

\end{otherlanguage}

Va lāken īn ʿabd chon baʿżī rā żaʿīf moshāhede namūde, be-ḥabl-e
estedlāl tamassok joste ke shāyad az maʿīn-e kalemāt-e elāhī Kowthar-e
bāqī biyāshāmand va be-zendegī-ye dāʾemī fāʾez gardand. Va ellā inna
al-dalīla yastaḥyī an yunsaba ilayhi wa al-burhāna yakhjalu an yasjuda
bayna yadayhi. Neʿma mā qīl:

But since this servant has seen that some are weak, he has clung to the
cord of reasoning, so that perhaps they may drink the everlasting
Kawthar from the wellspring of divine words and attain eternal life.
Otherwise, evidence itself is ashamed to be attributed to Him, and proof
blushes to bow down before Him. How well it has been said:

\subsection{227}\label{section-2202}

\begin{otherlanguage}{arabic}

حق عیان چون مهر رخشان آمده حیف کاندر شهر کوران آمده

\end{otherlanguage}

Ḥaqq ʿayān chon mehr-e rakhshān āmade, ḥeyf k-andar shahr-e kūrān āmade.

The True One has appeared, plain as the radiant sun; alas that He has
come into a city of the blind.

\subsection{228}\label{section-2203}

\begin{otherlanguage}{persian}

سزاوار انسان آنکه بطراز عدل و انصاف مزین باشد مالک وجودی که بنفسه امام
عالم قیام فرموده و باسم قیومش رحیق مختوم را گشوده و معادل ما نزل من قبل
و من بعد از قلم مبارکش جاری و نازل چگونه میتواند انسان در این مقام توقف
نماید و یا نعوذ باللّه کلمهٴ نا لایق ادا کند بیان را نقطه اولی بید قدرت
بیفشرد زلال آن این کلمه مبارکه است انه لا یستشار باشارتی

\end{otherlanguage}

Sazāvār-e ensān ānke be-ṭerāz-e ʿadl va enṣāf mozayyan bāshad. Mālek-e
vojūdī ke be-nafsehī amām-e ʿālam qīyām farmūde va be-esm-e Qayyūmash
raḥīq-e makhtūm rā goshūde va moʿādel-e mā nazala min qabl va min baʿd
az qalam-e mobārakash jārī va nāzel, chegūne mītavānad ensān dar īn
maqām tavaqqof namāyad va yā naʿūdhu bi-llāh kalame-ye nā-lāyeq adā
konad? Bayān rā Noqṭe-ye Ūlā be-yad-e qodrat bīfeshord; zolāl-e ān īn
kalame-ye mobārake ast: «Innahu lā yustashāru bi-ishāratī.»

It befits a human being to be adorned with justice and fairness. The
Lord of existence has Himself arisen before the world, has unsealed the
choice wine through His name, the Self-Subsisting, and from His blessed
Pen has flowed and descended the equivalent of all that has been sent
down before and after. How can anyone hesitate in this station, or, God
forbid, utter an unworthy word? The Primal Point pressed the Bayān in
the hand of power; its pure essence is this blessed utterance: ``He is
not to be consulted on the basis of My indication.''

\subsection{229}\label{section-2204}

\begin{otherlanguage}{arabic}

یومی از ایام مقصود عالمیان باین کلمه علیا ناطق قوله جل جلاله یا عبد حاضر
امروز اگر نفس نقطه در این امر توقف نماید از کلمهٴ رضا محروم ماند ناس در
نوم غفلت و بصر اللّه از افق اعلی ناظر و شاهد انتهی.

\end{otherlanguage}

Yowmī az ayyām Maqṣūd-e ʿālamīyān be-īn kalame-ye ʿolyā nāṭeq, qawluhu
jalla jalāluhu: «Yā ʿabd-e ḥāżer, emrūz agar nafs-e Noqṭe dar īn Amr
tavaqqof namāyad, az kalame-ye reżā maḥrūm mānad. Nās dar nowm-e ghaflat
va baṣar-e Allāh az ofoq-e aʿlā nāẓer va shāhed.» Intahā.

One day the Desire of the worlds spoke these exalted words---glorified
be His glory: ``O servant who is present, today, if the Point Himself
should hesitate in this Cause, He would be deprived of the word of
approval. The people are in the sleep of heedlessness, while the eye of
God watches and witnesses from the Most Exalted Horizon.'' End of
quotation.

\subsection{230}\label{section-2205}

\begin{otherlanguage}{persian}

هر صاحب شمی عرف بیان رحمن را ادراک مینماید حق بمثل وجودش ظاهر و هویدا
حجباترا خرق فرمود سبحات را شق نمود مجدد سبب نشوید و ناس بیچاره را بمثل
قبل در تیه اوهام مبتلا منمائید اتقوا اللّه یا معشر البیان و لا تکونوا من
الظالمین

\end{otherlanguage}

Har ṣāḥeb-e shammī ʿarf-e bayān-e Raḥmān rā edrāk mīnamāyad. Ḥaqq
be-methl-e vojūdash ẓāher va hoveydā. Ḥojobāt rā kharq farmūd, soboḥāt
rā shaqq namūd. Mojaddad sabab nashavīd va nās-e bīchāre rā be-methl-e
qabl dar tīh-e owhām mobtalā manamāyīd. Ittaqū Allāha yā maʿshara
al-Bayān wa lā takūnū mina al-ẓālimīn.

Anyone with a sense of smell perceives the fragrance of the
All-Merciful's utterance. The True One is as manifest and evident as His
own existence. He has torn through the veils and rent the coverings of
glory. Do not become a cause of these again, and do not, as before,
leave the hapless people afflicted in the wilderness of vain imaginings.
Fear God, O people of the Bayān, and do not be among the unjust.

\subsection{231}\label{section-2206}

\begin{otherlanguage}{persian}

نقطه اولی روح ما سواه فداه از اول ظهور الی آخر بذکر این اسم مبارک مشغول
بهیچوجه در تشویق اهل بیان توقف نفرمودند هر لفظ مبارکی و هر کلمه مقدسهٴ
طیبهٴ که در بیان نازل شده بحق جل جلاله نسبت داده‌اند

\end{otherlanguage}

Noqṭe-ye Ūlā, rūḥu mā siwāhu fidāhu, az avval-e ẓohūr elā ākhar
be-dhekr-e īn Esm-e mobārak mashghūl. Be-hīch-vajh dar tashvīq-e ahl-e
Bayān tavaqqof nafarmūdand. Har lafẓ-e mobārakī va har kalame-ye
moqaddase-ye ṭayyebe-ī ke dar Bayān nāzel shode be-Ḥaqq, jalla jalāluhu,
nesbat dāde-and.

The Primal Point---may the spirit of all besides Him be sacrificed for
Him---was occupied, from the beginning of His revelation to its end,
with the mention of this blessed Name. He never ceased to encourage the
people of the Bayān. Every blessed expression and every holy and goodly
word revealed in the Bayān He ascribed to the True One, glorified be His
glory.

\subsection{232}\label{section-2207}

\begin{otherlanguage}{persian}

سبحان اللّه جهالت و نادانی بمقامی است و غل و بغضاء بشأنی است که اذان را
از اصغاء منع نموده قسم بآفتاب افق معانی اگر بعضی از آیات منزله در ذکر
این ظهور اعظم بر حجر القاء شود لینفجر منه الانهار

\end{otherlanguage}

Subḥāna Allāh, jahālat va nādānī be-maqāmī ast va ghell va baghżā
be-shaʾnī ast ke ādhān rā az eṣghā manʿ namūde. Qasam be-āftāb-e ofoq-e
maʿānī, agar baʿżī az āyāt-e monzale dar dhekr-e īn Ẓohūr-e Aʿẓam bar
ḥajar elqā shavad, la-yanfajiru minhu al-anhār.

Glory be to God! Ignorance and folly have reached such a degree, and
malice and hatred such an intensity, that they have prevented ears from
listening. I swear by the Sun of the horizon of meanings: if some of the
verses revealed concerning this Most Great Revelation were cast upon a
stone, rivers would surely gush from it.

\subsection{233}\label{section-2208}

\begin{otherlanguage}{persian}

وقتی از اوقات این کلمه علیا از لسان مولی الوری استماع شد یا عبد حاضر
بیاناتی که از ملکوت علم الهی در این ظهور ظاهر شده گوش بخشد و القا نماید
محزون مباش از گفته‌های این و آن یک نفس اگر حلاوت بیان رحمن را بیابد او
عند اللّه از من علی الارض اقدم و افضل و اظهر و اعلم و اعلی بوده انتهی.

\end{otherlanguage}

Vaqtī az owqāt īn kalame-ye ʿolyā az lesān-e Mawlā al-Warā estemāʿ shod:
«Yā ʿabd-e ḥāżer, bayānātī ke az malakūt-e ʿelm-e elāhī dar īn Ẓohūr
ẓāher shode gūsh bakhshad va elqā namāyad. Maḥzūn mabāsh az gofte-hā-ye
īn va ān. Yek nafs agar ḥalāvat-e bayān-e Raḥmān rā biyābad, ū ʿinda
Allāh az man ʿalā al-arḍ aqdam va afżal va aẓhar va aʿlam va aʿlā būde.»
Intahā.

On one occasion these exalted words were heard from the tongue of the
Lord of all people: ``O servant who is present, the utterances that have
appeared in this Revelation from the kingdom of divine knowledge bestow
ears and impart their message. Do not be saddened by what this person or
that person says. If a single soul discovers the sweetness of the
All-Merciful's utterance, that soul is, in God's sight, above all who
are on earth in precedence, excellence, manifestness, knowledge, and
loftiness.'' End of quotation.

\subsection{234}\label{section-2209}

\begin{otherlanguage}{persian}

نقطه اولی روح ما سواه فداه میفرماید اگردر یوم ظهور او یک نفس مؤمن باشد
بعث کل نبیین در نفس او میشود این عبد گاهی که در قدرت و عظمت و سلطنت حق
جل جلاله نظر مینماید ابواب رجا مفتوح میشود لذا استدعا مینماید که جمیع من
علی الارض را از نفحات ایّامش محروم نفرماید و چون در غفلت اهل بیان و ضغینه
و بغضای آن نفوس ملاحظه میشود یأس کلی حاصل

\end{otherlanguage}

Noqṭe-ye Ūlā, rūḥu mā siwāhu fidāhu, mīfarmāyad: «Agar dar yowm-e
ẓohūr-e Ū yek nafs moʾmen bāshad, baʿth-e koll-e nabīyīn dar nafs-e ū
mīshavad.» Īn ʿabd gāhī ke dar qodrat va ʿaẓamat va salṭanat-e Ḥaqq,
jalla jalāluhu, naẓar mīnamāyad, abvāb-e rajā maftūḥ mīshavad. Ledhā
estedʿā mīnamāyad ke jamīʿ-e man ʿalā al-arḍ rā az nafaḥāt-e ayyāmash
maḥrūm nafarmāyad. Va chon dar ghaflat-e ahl-e Bayān va żaghīne va
baghżā-ye ān nofūs molāḥeẓe mīshavad, yaʾs-e kollī ḥāṣel.

The Primal Point---may the spirit of all besides Him be sacrificed for
Him---says: ``If on the day of His appearance there is a single
believing soul, the resurrection of all the Prophets takes place in that
soul.'' Whenever this servant considers the power, grandeur, and
sovereignty of the True One, glorified be His glory, the doors of hope
open. He therefore entreats Him not to deprive anyone on earth of the
breezes of His days. But when the heedlessness of the people of the
Bayān and the rancor and hatred of those souls are considered, utter
despair ensues.

\subsection{235}\label{section-2210}

\begin{otherlanguage}{persian}

اگر چه این عبد بگمان خود از این تحریرات انتباه و آگاهی نفوس را خواسته
ولکن مشاهده میشود نقطه اولی روح ما سواه فداه بشأنی در این ظهور اعظم ذکر
فرموده‌اند که ما فوق آن متصوّر نه مع ذلک میفرمایند بقدر اسم مؤمن هم در حق
او راضی نمیشوند چه که اگر راضی شوند حزن بر او وارد نمیآوردند

\end{otherlanguage}

Agar che īn ʿabd be-gomān-e khod az īn taḥrīrāt entebāh va āgāhī-ye
nofūs rā khāste, va lāken moshāhede mīshavad Noqṭe-ye Ūlā, rūḥu mā
siwāhu fidāhu, be-shaʾnī dar īn Ẓohūr-e Aʿẓam dhekr farmūde-and ke mā
fowq-e ān motaṣavvar na. Maʿa dhālek mīfarmāyand: «Be-qadr-e esm-e
moʾmen ham dar ḥaqq-e Ū rāżī namīshavand, che ke agar rāżī shavand ḥozn
bar Ū vāred namīāvardand.»

Although this servant, as he sees it, has sought through these writings
to awaken souls and make them aware, it is evident that the Primal
Point---may the spirit of all besides Him be sacrificed for Him---has
spoken of this Most Great Revelation in such a manner that nothing
beyond it can be conceived. Nevertheless, He says: ``They will not be
willing to accord Him even the name of believer, for if they were
willing, they would not bring sorrow upon Him.''

\subsection{236}\label{section-2211}

\begin{otherlanguage}{persian}

فی الحقیقه این بیان نقطهٴ بیان روح من فی الامکان فداه کتاب وعظ است و
دفتر نصح و صحیفهٴ انتباه ولکن لا یزید الظالمین الا خسارا و لا یزید
الغافلین الا انکارا و لا یزید الفاسقین الّا فرارا

\end{otherlanguage}

Fī al-ḥaqīqe īn bayān-e Noqṭe-ye Bayān, rūḥu man fī al-imkān fidāhu,
ketāb-e vaʿẓ ast va daftar-e noṣḥ va ṣaḥīfe-ye entebāh, va lāken lā
yazīdu al-ẓālimīna illā khasāran wa lā yazīdu al-ghāfilīna illā inkāran
wa lā yazīdu al-fāsiqīna illā firāran.

In truth, this utterance of the Point of the Bayān---may the spirit of
all in the contingent world be sacrificed for Him---is a book of
admonition, a volume of counsel, and a scroll of awakening. Yet it only
increases the loss of the unjust, the denial of the heedless, and the
flight of the transgressors.

\subsection{237}\label{section-2212}

\begin{otherlanguage}{persian}

این فانی امام وجوه کل عرض مینماید در حالتیکه قلبش فارغ است از ظنون و
اوهام و ضغینه وبغضای نفسی از نفوس آیا اگر نقطه بیان یعنی منزل آن بفرماید
یا معشر البیان بچه حجت و برهان از این ظهور اعظم محتجب مانده‌اید بچه جواب
متمسک و متشبث میشوید

\end{otherlanguage}

Īn fānī amām-e vojūh-e koll ʿarż mīnamāyad, dar ḥālatī ke qalbash fāregh
ast az ẓonūn va owhām va żaghīne va baghżā-ye nafsī az nofūs: Āyā agar
Noqṭe-ye Bayān, yaʿnī Monzel-e ān, befarmāyad «Yā maʿshara al-Bayān,
be-che ḥojjat va borhān az īn Ẓohūr-e Aʿẓam moḥtajeb mānde-īd?», be-che
javāb motamassek va motashabbeth mīshavīd?

This mortal one submits before everyone, his heart free of suspicions
and vain imaginings and of rancor and hatred toward any soul: If the
Point of the Bayān, that is, its Revealer, were to say, ``O people of
the Bayān, on what evidence and proof have you remained veiled from this
Most Great Revelation?'', what answer would you seize upon and cling to?

\subsection{238}\label{section-2213}

\begin{otherlanguage}{persian}

اگر میگوئید بمن محتجب شده‌اید جمیع ذرات کائنات شاهد و گواه که من اخذ عهد
او را قبل از اخذ عهد خود نموده‌ام سبحانک اللّهم فاشهد بانی بذلک الکتاب قد
اخذت عهد ولایة من تظهرنه عن کل شئ قبل عهد ولایتی و کفی بک و بمن آمن
بآیاتک علی شهداء و انّک انت حسبی علیک توکلت و انت کنت علی کل شئ حسیبا

\end{otherlanguage}

Agar mīgūʾīd be-man moḥtajeb shode-īd, jamīʿ-e dharrāt-e kāʾenāt shāhed
va gavāh ke man akhdh-e ʿahd-e Ū rā qabl az akhdh-e ʿahd-e khod
namūde-am. Subḥānaka Allāhumma fa-shhad bi-annī bi-dhālika al-kitābi qad
akhadhtu ʿahda wilāyati man tuẓhirannahu ʿan kulli shayʾin qabla ʿahdi
wilāyatī, wa kafā bika wa bi-man āmana bi-āyātika ʿalayya shuhadāʾa, wa
innaka anta ḥasbī, ʿalayka tawakkaltu wa anta kunta ʿalā kulli shayʾin
ḥasīban.

If you say that you have become veiled because of Me, every atom of
creation bears witness and testifies that I took the covenant of
allegiance to Him before taking the covenant of allegiance to Myself.
``Glory be to You, O God! Bear witness that through that Book I have
taken from all things the covenant of allegiance to Him whom You will
make manifest before the covenant of allegiance to Myself. You and those
who have believed in Your verses suffice as witnesses concerning Me. You
are indeed sufficient for Me. In You have I placed My trust, and You
have ever taken account of all things.''

\subsection{239}\label{section-2214}

\begin{otherlanguage}{persian}

و به بیانی که اظهر از شمس است در وسط سماء کل را وصیت نموده‌ام که بمن از
آن مقصود عالمیان محجوب نمانید نفس کتاب شاهد و گواه و اگر به بیان از آن
سلطان ایام محروم مانده‌اید هر سطری و هر ورقی از آن شاهد وگواه که من باعلی
النداء گفته‌ام بجمیع بیان و منزل و مظهر وسلطان آن محروم نمانید

\end{otherlanguage}

Va be-bayānī ke aẓhar az shams ast dar vasaṭ-e samā koll rā vaṣīyat
namūde-am ke be-man az ān Maqṣūd-e ʿālamīyān maḥjūb namānīd. Nafs-e
Ketāb shāhed va gavāh. Va agar be-Bayān az ān Solṭān-e ayyām maḥrūm
mānde-īd, har saṭrī va har varaqī az ān shāhed va gavāh ke man be-aʿlā
al-nedā gofte-am: Be-jamīʿ-e Bayān va Monzel va Maẓhar va Solṭān-e ān
maḥrūm namānīd.

And in words more evident than the sun in the midst of heaven, I
counseled everyone not to remain veiled by Me from that Desire of the
worlds. The Book itself bears witness and testifies. And if, because of
the Bayān, you have remained deprived of that Sovereign of days, every
line and every leaf of it bears witness and testifies that I proclaimed
in the loudest voice: Do not let the whole Bayān, or its Revealer, its
Manifestation, and its Sovereign, deprive you of Him.

\subsection{240}\label{section-2215}

\begin{otherlanguage}{persian}

اگر بوصفی ناظرید و از مالک اوصاف ممنوع گشته‌اید بتصریح تمام گفته‌ام که کل
آنچه ذکر خیر در بیان نازل شده مقصود آن شمس ظهور بوده و اگر بلفظ مرآت از
شمس معانی محتجب شده‌اید قد نزلت فی الکتاب ان یا شموس المرایا انتم الی شمس
الحقیقه تنظرون فان قیامکم بها لو انتم تتبصرون کلکم کحیتان بالماء فی
البحر تتحرکون و تحتجبون عن المأ و تسئلون عما انتم به قائمون الی آخر قولی
لاشکون الیک ان یا مرآت جودی عن کل المرایا کل بالوانهم الی لینظرون

\end{otherlanguage}

Agar be-vaṣfī nāẓerīd va az Mālek-e owṣāf mamnūʿ gashteh-īd, be-taṣrīḥ-e
tamām gofteh-am keh koll-e āncheh dhekr-e kheyr dar Bayān nāzel shodeh
maqṣūd ān Shams-e Ẓohūr būdeh, va agar be-lafẓ-e merʾāt az Shams-e
Maʿānī moḥtajeb shodeh-īd, qad nazalat fī al-Kitāb: an yā shumūs
al-marāyā, antum ilā Shams al-Ḥaqīqa tanẓurūn, fa-inna qiyāmakum bihā
law antum tatabaṣṣarūn. Kullukum ka-ḥītān bi-l-māʾ fī al-baḥr
tataḥarrakūn wa taḥtajibūn ʿan al-māʾ wa tasʾalūn ʿammā antum bihi
qāʾimūn. Ilā ākhir-i qawlī: la-ashkuwanna ilayka, an yā Mirʾāt Jūdī, ʿan
kull al-marāyā; kullun bi-alwānihim ilayya la-yanẓurūn.

If you are looking to an attribute and have been kept from the Possessor
of all attributes, I have stated with complete explicitness that
everything spoken of favorably in the Bayān was intended to refer to
that Sun of Revelation. And if the word ``mirror'' has veiled you from
the Sun of Meanings, this has been revealed in the Book: ``O suns of the
mirrors! You look toward the Sun of Truth, for your existence depends on
it, if only you could perceive. You are all like fish moving through the
water in the sea, yet you are veiled from the water and ask about that
upon which your existence depends.'' And so on, to My words: ``I will
surely complain to You, O Mirror of My bounty, of all the mirrors: each
looks toward Me in its own colors.''

\subsection{241}\label{section-2216}

\begin{otherlanguage}{arabic}

و اگر بذکر ثمره محتجب مانده‌اید هذا عهدی و میثاقی مع ربی من قبل و من بعد
و اجعل اللّهم تلک الشجرة کلها له لیظهرن منها ثمرات ما خلق اللّه فیها لمن
قد اراد اللّه ان یظهر به ما اراد فاننی انا و عزتک ما اردت ان یکون علی تلک
الشجرة من غصن و لا ورق و لا ثمر لن یسجد له یوم ظهوره و لا یسبحک به بما
ینبغی لعلو علو ظهوره و سموّ سموّ بطونه و ان شهدت یا الهی علیّ من غصن او ورق
او ثمر لم یسجد له یوم ظهوره فاقطعه اللّهم عن تلک الشجرة فانه لم یکن منی و
لا یرجع الیّ

\end{otherlanguage}

Va agar be-dhekr-e thamareh moḥtajeb māndeh-īd, hādhā ʿahdī wa-mīthāqī
maʿa Rabbī min qablu wa-min baʿdu. Wa-jʿal, Allāhumma, tilka
al-shajarata kullahā lahu, li-yaẓharanna minhā thamarātu mā khalaqa
Allāhu fīhā li-man qad arāda Allāhu an yuẓhira bihi mā arāda. Fa-innanī
anā wa-ʿizzatika mā aradtu an yakūna ʿalā tilka al-shajarati min ghuṣnin
wa-lā waraqin wa-lā thamarin lan yasjuda lahu yawma ẓuhūrihi wa-lā
yusabbiḥaka bihi bi-mā yanbaghī li-ʿuluwwi ʿuluwwi ẓuhūrihi wa-sumūwwi
sumūwwi buṭūnihi. Wa-in shahidta, yā Ilāhī, ʿalayya min ghuṣnin aw
waraqin aw thamarin lam yasjud lahu yawma ẓuhūrihi, fa-qṭaʿhu,
Allāhumma, ʿan tilka al-shajarati, fa-innahu lam yakun minnī wa-lā
yarjiʿu ilayya.

And if mention of the fruit has left you veiled: ``This is My covenant
and My pledge with My Lord, before and after. Make that Tree, O God,
entirely His, so that from it may appear the fruits of what God has
created within it for Him through whom God has willed to manifest
whatever He has willed. By Your might, I have never desired that there
should be upon that Tree any branch, leaf, or fruit that will not bow
down before Him on the day of His appearance, or will not glorify You
through Him as befits the loftiest heights of His manifestation and the
most exalted sublimity of His concealment. And if You witness in Me, O
my God, any branch, leaf, or fruit that does not bow down before Him on
the day of His appearance, cut it off from that Tree, O God, for it is
not of Me and does not return to Me.''

\subsection{242}\label{section-2217}

\begin{otherlanguage}{arabic}

اگر بذکر مستغاث محتجب مانده‌اید این ندا از جمیع بیان مرتفع اگر در این حین
ظاهر شود من اول عابدین و اول ساجدینم اگر بکلماتی متمسکید بتأکید تمام ذکر
نموده‌ام که جمیع بیان معادله بیک آیه او نمیکند

\end{otherlanguage}

Agar be-dhekr-e Mostaghāth moḥtajeb māndeh-īd, īn nedā az jamīʿ-e Bayān
mortafeʿ: agar dar īn ḥīn ẓāher shavad, man avval-e ʿābedīn va avval-e
sājedīnam. Agar be-kalamātī motamassekīd, be-taʾkīd-e tamām dhekr
namūdeh-am keh jamīʿ-e Bayān moʿādeleh be-yek āyeh-ye Ū nemīkonad.

If mention of the Mustaghāth has left you veiled, this call rises from
the whole Bayān: ``If He were to appear at this very moment, I would be
the first to worship Him and the first to bow down before Him.'' If you
cling to certain words, I have stated with the utmost emphasis that the
whole Bayān does not equal a single verse of His.

\subsection{243}\label{section-2218}

\begin{otherlanguage}{persian}

حال ای قوم آیا آن نفسی که من دون او اخذ نموده‌اید و کلمات او را آیات فرض
گرفته‌اید از من و بیان من اعلی و احبّ است نزد شما افّ لکم و لقلة شعورکم و
عدم وفائکم اعظم اوهامات و ظنون شما آنکه ظهورم را مثل ظهور قبل فرض
نموده‌اید خلیفه و وصی تعیین مینمائید بالمره از یوم اللّه و شأن او محجوبید
و از اشراقات انوار آفتاب توحید حقیقی محروم و بی خبر

\end{otherlanguage}

Ḥāl ey qowm, āyā ān nafsī keh min dūn-e Ū akhdh namūdeh-īd va kalamāt-e
ū rā āyāt farż gerefteh-īd az man va Bayān-e man aʿlā va aḥabb ast
nazd-e shomā? Uffin lakum wa-li-qillati shuʿūrikum wa-ʿadami wafāʾikum.
Aʿẓam-e owhāmāt va ẓonūn-e shomā ānkeh ẓohūram rā methl-e ẓohūr-e qabl
farż namūdeh-īd, khalīfeh va vaṣī taʿyīn mīnamāyīd. Be-l-marreh az
Yowm-e Allāh va shaʾn-e Ū maḥjūbīd va az eshrāqāt-e anvār-e āftāb-e
towḥīd-e ḥaqīqī maḥrūm va bī-khabar.

Now, people, is the one you have chosen instead of Him, whose words you
have assumed to be divine verses, more exalted and dearer to you than I
and My Bayān? Shame on you for your lack of understanding and your
faithlessness! Your greatest delusion and vain imagining is that you
have assumed My Revelation to be like the Revelation before it, and you
appoint a successor and an executor. You are utterly veiled from the Day
of God and its station, and are deprived and unaware of the radiance of
the Sun of true unity.

\subsection{244}\label{section-2219}

\begin{otherlanguage}{persian}

ای ظالمان هنوز از قمیصم عرف دم اطهرم متضوع سیف‌های آن ملحدین که در هزار و
دویست سنه بصیقل غل و بغضا مهیا نمودند و از سموم آبش دادند آیا کفایت
ننموده که مجدد در تیه این اوهامات هائمید بآنچه در طول هزار و دویست سنه
عامل شدید قناعت نمائید آیا در بیان نازل نشده باسماء از مالک آن محتجب
نمانید حتی اسم النبی فان ذلک الاسم یخلق بقوله

\end{otherlanguage}

Ey ẓālemān, hanūz az qamīṣam ʿarf-e dam-e aṭharam motaḍavveʿ. Seyf-hā-ye
ān molḥedīn keh dar hezār va devīst saneh be-ṣeyqal-e ghell va baghżā
mohayyā namūdand va az somūm ābash dādand, āyā kefāyat nanamūdeh keh
mojaddad dar tīh-e īn owhāmāt hāʾemīd? Be-āncheh dar ṭūl-e hezār va
devīst saneh ʿāmel shodīd qanāʿat namāyīd. Āyā dar Bayān nāzel nashodeh:
be-asmāʾ az Mālek-e ān moḥtajeb namānīd, ḥattā ism al-Nabī, fa-inna
dhālika al-isma yukhlaqu bi-qawlihi?

You oppressors! The fragrance of My most pure blood still wafts from My
garment. Have the swords of those unbelievers, which they prepared over
twelve hundred years with the polish of malice and hatred and tempered
with poisons, not sufficed, that you wander once again in the wilderness
of these delusions? Be content with what you have done over the course
of twelve hundred years. Has it not been revealed in the Bayān: ``Do not
let names veil you from their Possessor, even the name `Prophet,' for
that name is created by His word''?

\subsection{245}\label{section-2220}

\begin{otherlanguage}{persian}

سبحان اللّه با امواج بحر بیان حضرت رحمن و اشراقات انوار آفتاب نصح و
ظهورات نیر تدبیر و تعلیم از یمین قلب مبارک نقطه روح ما سواه فداه اریاح
یاسیه در هبوب بشأنیکه حزن عالم را احاطه نموده و سبب یأس افئده مقدسه
منوره شده مخصوص این فانی را بشأنی اخذ نموده که ظهور غیر آن را از محالات
شمرده

\end{otherlanguage}

Subḥān Allāh! Bā amvāj-e baḥr-e bayān-e Ḥażrat-e Raḥmān va eshrāqāt-e
anvār-e āftāb-e noṣḥ va ẓohūrāt-e nayyer-e tadbīr va taʿlīm az yamīn-e
qalb-e mobārak-e Noqṭeh, rūḥ-e mā sevāh fedāh, aryāḥ-e yaʾsīyeh dar
hobūb, be-shaʾnī-keh ḥozn ʿālam rā eḥāṭeh namūdeh va sabab-e yaʾs-e
afʾedeh-ye moqaddaseh-ye monavvareh shodeh. Makhṣūṣ īn fānī rā be-shaʾnī
akhdh namūdeh keh ẓohūr-e gheyr-e ān rā az moḥālāt shomordeh.

Glory be to God! Despite the waves of the ocean of the All-Merciful's
utterance, the radiance of the sun of counsel, and the dawning of the
luminary of guidance and instruction from the right side of the blessed
heart of the Point---may the spirit of all else be sacrificed for
Him---the winds of despair are blowing so fiercely that sorrow has
encompassed the world and brought despair to sanctified, illumined
hearts. It has seized this transient one in particular so completely
that he considers the appearance of anything else impossible.

\subsection{246}\label{section-2221}

\begin{otherlanguage}{persian}

در باب سابع از واحد ثانی میفرماید قوله جل و عزّ ای اهل بیان نکرده آنچه
اهل قرآن کردند که ثمرات لیل خود را باطل کنید اگر آنچه که مؤمن به بیان
هستید در حین ظهور آیات او گفتید اللّه ربنا و لا نشرک به احد او ان هذا ما
وعدنا اللّه من مظهر نفسه لن ندعو معه شیئا و بر آنچه بر او هستید اطاعت او
کردید ثمره بیان را ظاهر کرده‌اید و الا لایق ذکر نیستید نزد خداوند ترحم بر
خود کرده اگر نصرت نمیکنید مظهر ربوبیت را محزون نکرده الی قوله جل شأنه
اگر بلقاء اللّه فائز نمیگردید آیة اللّه را هم محزون نکرده باشید از نفع
مدینین به بیان میگذرد هر گاه شما از ضرر باو بگذارید اگر چه میدانم
نخواهید کرد انتهی.

\end{otherlanguage}

Dar bāb-e sābeʿ az vāḥed-e thānī mīfarmāyad, qawluhu jalla wa-ʿazza: Ey
ahl-e Bayān, nakardeh āncheh ahl-e Qorʾān kardand keh thamarāt-e leyl-e
khod rā bāṭel konīd. Agar āncheh keh moʾmen be-Bayān hastīd, dar ḥīn-e
ẓohūr-e āyāt-e Ū goftīd: Allāhu Rabbunā wa-lā nushriku bihi aḥad, aw
inna hādhā mā waʿadanā Allāhu min maẓhari nafsihi, lan nadʿuwa maʿahu
shayʾan, va bar āncheh bar ū hastīd eṭāʿat-e Ū kardīd, thamareh-ye Bayān
rā ẓāher kardeh-īd, va ellā lāyeq-e dhekr nīstīd nazd-e Khodāvand.
Taraḥḥom bar khod kardeh, agar noṣrat nemīkonīd, Maẓhar-e Robūbīyat rā
maḥzūn nakardeh. Ilā qawlihi jalla shaʾnuhu: Agar be-leqāʾ Allāh fāʾez
nemīgardīd, Āyat Allāh rā ham maḥzūn nakardeh bāshīd. Az nafʿ-e
motadayyenīn be-Bayān mīgodharad har gāh shomā az żarar be-Ū bogdhārīd,
agarcheh mīdānam nakhāhīd kard. Intahā.

In the seventh chapter of the second Vāḥid, He says---exalted and mighty
are His words: ``People of the Bayān! Do not do what the people of the
Qurʾān did, thereby rendering the fruits of your own night futile. If,
to the extent that you believe in the Bayān, you say at the appearance
of His verses, `God is our Lord, and we associate no one with Him,' or
`This is what God promised us through the Manifestation of Himself; we
will invoke nothing alongside Him,' and if, in whatever state you are,
you obey Him, you will have brought forth the fruit of the Bayān.
Otherwise, you are unworthy of mention before God. Have mercy on
yourselves: if you do not aid the Manifestation of Lordship, do not
grieve Him.'' And further, He says---exalted is His station: ``If you do
not attain the presence of God, at least do not grieve the Sign of God.
He will forgo any benefit from the followers of the Bayān if you will
refrain from harming Him, although I know you will not.'' End of
quotation.

\subsection{247}\label{section-2222}

\begin{otherlanguage}{arabic}

سبحان اللّه مشاهده نمائید آن طلعت مبارک را چه مقدار یأس اخذ نموده که مع
وصایای لا تحصی و عهدهای محکمه و میثاقهای متقنه و تدبیرهای معلنه یأس اکبر
از بیانات مبارک ظاهر و مشهود طوبی للباکین

\end{otherlanguage}

Subḥān Allāh! Moshāhedeh namāyīd ān Ṭalʿat-e Mobārak rā cheh meqdār yaʾs
akhdh namūdeh keh maʿ-e vaṣāyā-ye lā toḥṣā va ʿahd-hā-ye moḥkameh va
mīthāq-hā-ye motqaneh va tadbīr-hā-ye moʿlaneh, yaʾs-e akbar az
bayānāt-e mobārak ẓāher va mashhūd. Ṭūbā li-l-bākīn.

Glory be to God! Consider how deeply despair had seized that blessed
Countenance: despite His countless counsels, firm covenants, binding
pledges, and openly declared measures, the greatest despair is evident
and manifest in His blessed utterances. Blessed are those who weep.

\subsection{248}\label{section-2223}

\begin{otherlanguage}{persian}

و در مقام شفقت کبری و عالم آگاهی میفرمایند قوله عزّ و جل چه بسا نوری را
که نار میفرماید باو و همچنین میفرماید قوله عزّ ذکره اگر چه می‌بینم ظهور او
رامثل این شمس در وسط السماء و غروب کل را بمثل نجوم لیل در نهار هر صاحب
سمع و بصری شهادت میدهد بر اینکه آنحضرت روح ما سواه فداه در جمیع احیان
بذکر محبوب امکان مشغول بوده‌اند و بجمیع تدبیرات قلمیه و بیانیه و ذکریه
اخبار فرموده‌اند بآنچه از اهل بیان ظاهر شده و میشود

\end{otherlanguage}

Va dar maqām-e shafaqat-e kobrā va ʿālam-e āgāhī mīfarmāyand, qawluhu
ʿazza wa-jalla: Cheh basā nūrī rā keh nār mīfarmāyad be-ū. Va hamchenīn
mīfarmāyad, qawluhu ʿazza dhikruhu: Agarcheh mībīnam ẓohūr-e Ū rā
methl-e īn shams dar vasaṭ-e samāʾ va ghorūb-e koll rā be-methl-e
nojūm-e leyl dar nahār. Har ṣāḥeb-e samʿ va baṣarī shahādat mīdahad bar
īnkeh ān Ḥażrat, rūḥ-e mā sevāh fedāh, dar jamīʿ-e aḥyān be-dhekr-e
Maḥbūb-e emkān mashghūl būdeh-and va be-jamīʿ-e tadbīrāt-e qalamīyeh va
bayānīyeh va dhekrīyeh ekhbār farmūdeh-and be-āncheh az ahl-e Bayān
ẓāher shodeh va mīshavad.

And from the station of supreme compassion and the realm of awareness,
He says---exalted and glorious are His words: ``How often He calls a
light `fire.'\,'' He also says---exalted be His remembrance: ``Although
I see His appearance as this sun in the midst of heaven, and the setting
of all others as the stars of night in the daytime.'' Everyone with
hearing and sight bears witness that He---may the spirit of all else be
sacrificed for Him---was at all times occupied with mention of the
Beloved of all existence, and that through every means of writing,
exposition, and admonition He foretold what has proceeded and will
proceed from the people of the Bayān.

\subsection{249}\label{section-2224}

\begin{otherlanguage}{persian}

لعمر مقصودی و مقصود من فی الارض و السماء ان اسفه عظیم و حزنه کبیر و یأسه
عمن فی البیان اعظم من کل عظیم ای معشر غافلین بشنوید عرض این خادم فانی را
و بهواهای نفسانیه و مفتریات نالایقه ذیل اطهر اقدس را میالائید اگر اقبال
نمینمائید ساکت باشید میفرماید عزّ ذکره در بیان هیچ ذکری نیست مگر ذکر او
لعل در وقت ظهور مشاهده حزن نفرماید از مؤمنین بخود و همچنین باین کلمه
مبارکه که بمثابه نور است ما بین عالم نطق فرموده قوله جل و عزّ بیان از اول
تا آخر مکمن جمیع صفات اوست و خزانه نار و نور او

\end{otherlanguage}

La-ʿamri Maqṣūdī wa-Maqṣūdi man fī al-arḍi wa-l-samāʾ, inna asafahu
ʿaẓīmun wa-ḥuznahu kabīrun wa-yaʾsahu ʿamman fī al-Bayān aʿẓamu min
kulli ʿaẓīm. Ey maʿshar-e ghāfelīn, beshnavīd ʿarż-e īn khādem-e fānī rā
va be-havā-hā-ye nafsānīyeh va moftarayāt-e nālāyeqeh dheyl-e aṭhar-e
aqdas rā mayālāyīd. Agar eqbāl nemīnamāyīd, sāket bāshīd. Mīfarmāyad,
ʿazza dhikruhu: Dar Bayān hīch dhekrī nīst magar dhekr-e Ū, laʿall dar
vaqt-e ẓohūr moshāhedeh-ye ḥozn nafarmāyad az moʾmenīn be-khod. Va
hamchenīn be-īn kalameh-ye mobārakeh keh be-mathābeh-ye nūr ast mā
beyn-e ʿālam noṭq farmūdeh, qawluhu jalla wa-ʿazza: Bayān az avval tā
ākhar makman-e jamīʿ-e ṣefāt-e Ūst va khazāneh-ye nār va nūr-e Ū.

By the life of my Desire and the Desire of all who are in earth and
heaven! His anguish is immense, His sorrow great, and His despair of
those in the Bayān greater than any greatness. You heedless ones! Listen
to the plea of this transient servant, and do not stain that most pure
and holy hem with selfish desires and unworthy fabrications. If you will
not turn toward Him, remain silent. He says---exalted be His
remembrance: ``There is no mention in the Bayān except mention of Him,
so that at the time of His appearance He may perhaps experience no
sorrow from those who believe in Me.'' He has also spoken these blessed
words, which are like a light amid the world---exalted and mighty are
His words: ``The Bayān, from beginning to end, is the repository of all
His attributes and the treasury of His fire and His light.''

\subsection{250}\label{section-2225}

\begin{otherlanguage}{persian}

امید هست که ناس غافل از آنچه ذکر شد بسماء آگاهی عروج نمایند چنانچه
میفرماید قوله عز ذکره اگر کل بتربیت نقطه بیان مربی شوند و ناظر بمبدء و
حجت و دلیل باشند در حین ظهور محتجب نمیمانند در این مقام میفرماید قوله
تعالی گویا دیده میشود که کل این اصحاب بشرف ایمان بآن شمس حقیقت مستضئی
گشته در تلویحات وتصریحات ملاحظه نمائید میفرمایند امید هست که این اصحاب
یعنی اصحاب موجوده بشرف ایمان فائز شوند چنانچه شدند

\end{otherlanguage}

Omīd hast keh nās-e ghāfel az āncheh dhekr shod be-samā-ye āgāhī ʿorūj
namāyand. Chenāncheh mīfarmāyad, qawluhu ʿazza dhikruhu: Agar koll
be-tarbīyat-e Noqṭeh-ye Bayān morabbā shavand va nāẓer be-mabdaʾ va
ḥojjat va dalīl bāshand, dar ḥīn-e ẓohūr moḥtajeb nemīmānand. Dar īn
maqām mīfarmāyad, qawluhu taʿālā: Gūyā dīdeh mīshavad keh koll-e īn
aṣḥāb be-sharaf-e īmān be-ān Shams-e Ḥaqīqat mostażīʾ gashteh. Dar
talvīḥāt va taṣrīḥāt molāḥeẓeh namāyīd. Mīfarmāyand: Omīd hast keh īn
aṣḥāb, yaʿnī aṣḥāb-e mowjūdeh, be-sharaf-e īmān fāʾez shavand,
chenāncheh shodand.

It is hoped that, through what has been mentioned, the heedless may
ascend to the heaven of awareness. As He says---exalted be His
remembrance: ``If all are educated by the Point of the Bayān and look to
the Source, the proof, and the evidence, they will not remain veiled at
the time of the appearance.'' In this connection He says---exalted are
His words: ``It is as though one sees all these companions illumined by
the honor of faith in that Sun of Truth.'' Consider the allusions and
explicit statements. He says that it is hoped these companions, meaning
the companions then living, will attain the honor of faith, as indeed
they did.

\subsection{251}\label{section-2226}

\begin{otherlanguage}{persian}

خدای واحد شاهد و حق آگاه گواه که این فانی خالصانه لوجهه و لاولیائه و
لاحبائه ذکر نمود آنچه که منصفین را نفع بخشد و نجات دهد از بیانات نقطه
اولی روح ما سواه فداه چه عربی و چه فارسی آنچه در نظر بود بعضی از آن ذکر
شد اگر چه الیوم ام البیان میفرماید جمیع بیان ذکری است از او و برد و قبول
او معلق و منوط

\end{otherlanguage}

Khodā-ye Vāḥed shāhed va Ḥaqq-e Āgāh govāh keh īn fānī khāleṣāneh
li-wajhihi wa-li-awliyāʾihi wa-li-aḥibbāʾihi dhekr namūd āncheh keh
monṣefīn rā nafʿ bakhshad va nejāt dahad. Az bayānāt-e Noqṭeh-ye Ūlā,
rūḥ-e mā sevāh fedāh, cheh ʿarabī va cheh fārsī, āncheh dar naẓar būd
baʿżī az ān dhekr shod, agarcheh al-yowm Omm al-Bayān mīfarmāyad:
Jamīʿ-e Bayān dhekrī ast az Ū va be-radd va qabūl-e Ū moʿallaq va manūṭ.

The one God is my witness, and the all-knowing True One bears testimony,
that this transient one has sincerely set forth, for His sake and for
His chosen ones and His loved ones, what will benefit and save the
fair-minded. Some of the utterances of the Primal Point---may the spirit
of all else be sacrificed for Him---that were in mind, whether Arabic or
Persian, have been cited, although today the Mother of Utterance
declares: ``The whole Bayān is a mention of Him and is dependent upon
and subject to His rejection or acceptance.''

\subsection{252}\label{section-2227}

\begin{otherlanguage}{arabic}

اگر از فضل حق جل جلاله افق عدل و انصاف از سحاب ظنون و اوهام مقدس شود
جمیع شهادت دهند بر آنچه این خادم ذکر نمود چه که کل میدانند نفسی بقدر این
عبد از این امر مطلع نبوده ونیست مع ذلک نفوسی که از حقیقت آگاه نبوده‌اند
گفته‌اند آنچه را که بعضی از بی خبرهای عالم اخذ نموده و بر طلعت قدم وارد
آورده‌اند آنچه را که صخره صیحه زد\ldots.. انتهی.

\end{otherlanguage}

Agar az fażl-e Ḥaqq, jalla jalāluhu, ofoq-e ʿadl va enṣāf az saḥāb-e
ẓonūn va owhām moqaddas shavad, jamīʿ shahādat dahand bar āncheh īn
khādem dhekr namūd, cheh keh koll mīdānand nafsī be-qadr-e īn ʿabd az īn
amr moṭṭaleʿ nabūdeh va nīst. Maʿ-e dhālek nofūsī keh az ḥaqīqat āgāh
nabūdeh-and gofteh-and āncheh rā keh baʿżī az bī-khabar-hā-ye ʿālam
akhdh namūdeh va bar Ṭalʿat-e Qedam vāred āvordeh-and āncheh rā keh
ṣakhreh ṣeyḥeh zad\ldots.. Intahā.

If, through the grace of the True One---exalted be His glory---the
horizon of justice and fairness is cleared of the clouds of conjecture
and delusion, all will bear witness to what this servant has stated, for
everyone knows that no one has been or is as well informed about this
matter as this servant. Nevertheless, people who were unaware of the
truth have said things that some of the uninformed in the world have
taken up, and they have inflicted upon the Ancient Beauty what has made
the rock cry out\ldots.. End of quotation.

\subsection{253}\label{section-2228}

\begin{otherlanguage}{arabic}

و در لوح ملا علی بجستانی نازل قوله تعالی :

\end{otherlanguage}

Va dar lawḥ-e Mollā ʿAlī-ye Bajestānī nāzel, qawluhu taʿālā:

And in the Tablet to Mullā ʿAlī Bajestānī, this has been
revealed---exalted are His words:

\subsection{254}\label{section-2229}

\begin{otherlanguage}{persian}

آنچه از قلم اعلی نازل شد البته ظاهر خواهد گشت سورهٴ رئیس را تلاوت
فرمائید و همچنین لوح ملک پاریس را که از اجزای سوره مبارکه هیکل است و
همچنین لوح فوأد که مخصوص یکی از احباب نازل شده و این لوح در وقتی نازل شد
که پاشا که وزیر خارجه روم بود بمقر خود راجع شده بود و سبب فتنه اخیره و
مهاجرت از ارض سرّ بعکا او شده بود

\end{otherlanguage}

Āncheh az Qalam-e Aʿlā nāzel shod albatteh ẓāher khāhad gasht. Sūreh-ye
Raʾīs rā telāvat farmāyīd va hamchenīn Lawḥ-e Malek-e Pārīs rā keh az
ajzā-ye Sūreh-ye Mobārakeh-ye Heykal ast, va hamchenīn Lawḥ-e Fuʾād keh
makhṣūṣ-e yekī az aḥbāb nāzel shodeh. Va īn lawḥ dar vaqtī nāzel shod
keh pāshā keh vazīr-e khārejeh-ye Rūm būd be-maqarr-e khod rājeʿ shodeh
būd va sabab-e fetneh-ye akhīreh va mohājerat az Arż-e Serr be-ʿAkkā ū
shodeh būd.

Whatever has been revealed by the Most Exalted Pen will assuredly come
to pass. Read the Sūrih of the Chief, and likewise the Tablet to the
King of Paris, which forms part of the blessed Sūrih of the Temple, and
also the Tablet of Fuʾād, which was revealed specifically for one of the
friends. This Tablet was revealed when the Pasha who was the Ottoman
foreign minister had returned to his abode; he had been the cause of the
latest upheaval and of the departure from the Land of Mystery to ʿAkkā.

\subsection{255}\label{section-2230}

\begin{otherlanguage}{arabic}

دو نفر بودند که بعد از سلطان رئیس کل بودند یکی فؤاد پاشا و یکی عالی پاشا
گاهی این صدر اعظم بود و آن وزیر اول خارجه و گاهی بالعکس در آن لوح
میفرمایند قوله عزّ کبریائه '' سوف نعزل الذی کان مثله و نأخذ امیرهم الذی
یحکم علی البلاد و انا العزیز الجبار

\end{otherlanguage}

Do nafar būdand keh baʿd az solṭān raʾīs-e koll būdand, yekī Fuʾād Pāshā
va yekī ʿĀlī Pāshā. Gāhī īn ṣadr-e aʿẓam būd va ān vazīr-e avval-e
khārejeh va gāhī be-l-ʿaks. Dar ān lawḥ mīfarmāyand, qawluhu ʿazza
kibriyāʾuhu: ``Sawfa naʿzilu alladhī kāna mithlahu wa-naʾkhudhu amīrahum
alladhī yaḥkumu ʿalā al-bilādi wa-anā al-ʿAzīzu al-Jabbār.''

There were two men who, after the Sultan, were the chief authorities:
one was Fuʾād Pasha and the other ʿĀlī Pasha. Sometimes the latter was
grand vizier and the former chief minister of foreign affairs, and
sometimes the reverse. In that Tablet He says---exalted be His majesty:
``We will dismiss the one who was like him and seize their ruler who
governs the lands, and I am the Almighty, the All-Compelling.''

\subsection{256}\label{section-2231}

\begin{otherlanguage}{persian}

و همچنین در کتاب افدس در نقطه واقعه بین البحرین ملاحظه نمائید که مقصود
از آن نقطه اسلامبول است چه که از یک جهتش بحر ابیض است و جهة دیگر بحر
اسود باری آنچه از قلم اعلی جاری کل ظاهر شده و اخبارهای دیگر هم که در
الواح هست کل ظاهر خواهد شد نشهد انه هو العالم المقتدر السامع البصیر
الخبیر\ldots.

\end{otherlanguage}

Va hamchenīn dar Ketāb-e {[}corrupted text{]} dar noqṭeh-ye vāqeʿeh beyn
al-baḥrayn molāḥeẓeh namāyīd keh maqṣūd az ān noqṭeh Eslāmbūl ast, cheh
keh az yek jahatash Baḥr-e Abyaż ast va jahat-e dīgar Baḥr-e Asvad. Bārī
āncheh az Qalam-e Aʿlā jārī koll ẓāher shodeh va ekhbār-hā-ye dīgar ham
keh dar alvāḥ hast koll ẓāher khāhad shod. Nashhadu annahu huwa
al-ʿĀlimu al-Muqtadiru al-Sāmiʿu al-Baṣīru al-Khabīr\ldots.

Likewise, consider in the Book of {[}corrupted original text{]} the
place situated between the two seas. That place means Istanbul, for on
one side of it is the White Sea and on the other the Black Sea. In
short, everything that has flowed from the Most Exalted Pen has come to
pass, and all the other predictions contained in the Tablets will also
come to pass. We bear witness that He is the All-Knowing, the
All-Powerful, the All-Hearing, the All-Seeing, the All-Aware\ldots.

\bookmarksetup{startatroot}

\chapter{BH00009 (Unnamed)}\label{bh00009-unnamed}

\emph{Translated on 17 Qudrat 183 using GPT 6.1 Sol}

\subsection{1}\label{section-2232}

\begin{otherlanguage}{arabic}

بسم ربّنا الأقدس الأعظم العلیّ الأبهی

\end{otherlanguage}

Bi-smi rabbinā al-aqdas al-aʿẓam al-ʿaliyy al-abhā.

In the name of our Lord, the Most Holy, the Most Great, the Most High,
the Most Glorious.

\subsection{2}\label{section-2233}

\begin{otherlanguage}{persian}

امروز حفیف سدرهٴ منتهی مرتفع و آفتاب بیان از افق سماء معانی مشرق و لائح
صراط بافصح بیان ناطق و میزان باتمّ کلمات ذاکر مع‌ذلک رایة لیس فی الملک الّا
هو منصوب و مرتفع چه که عالم مخمود و افسرده بل مرده گوشی نه تا اصغا نماید
و چشمی نه تا مشاهده کند الّا معدودیکه سبحات اسماء و حجبات اشارات و کلمات
ایشانرا منع ننمود غیر حقّ را گذاشتند و باو پیوستند ندا را شنیدند و رحیق
را نوشیدند و بر امر قیام نمودند قیامیکه سبب و علّت اضطراب قلوب معتدین و
تزلزل اقدام مشرکین گشت فی‌الحقیقه هر گوشی لایق استماع نه و ابصار مرمودهٴ
غیر حدیده قابل مشاهده نه امروز روز ذکر و ثناست کِه را اولیا آنانکه بکمال
استقامت بر امر قائمند و بذکر مقصود عالم ناطق هر ذکر و هر ثنائی آن نفوس
مطمئنّهٴ راضیهٴ مرضیّه را سزاوار

\end{otherlanguage}

Emrūz ḥafīf-e Sedre-ye Montahā mortafeʿ va āftāb-e bayān az ofoq-e
samā-ye maʿānī moshreq va lāʾeḥ, ṣerāṭ be-afṣaḥ-e bayān nāṭeq va mīzān
be-atamm-e kalemāt dhāker. Maʿa-dhālek rāyat-e «laysa fī al-mulk illā
huwa» manṣūb va mortafeʿ, che ke ʿālam khomūd va afsorde bal morde.
Gūshī na tā eṣghā namāyad va chashmī na tā moshāhede konad, ellā maʿdūdī
ke soboḥāt-e asmā va ḥojobāt-e eshārāt va kalemāt īshān rā manʿ nanomūd.
Gheyr-e Ḥaqq rā godhāshtand va be-ū peyvastand, nedā rā shenīdand va
raḥīq rā nūshīdand va bar amr qīyām namūdand, qīyāmī ke sabab va
ʿellat-e eżṭerāb-e qolūb-e moʿtadīn va tazalzol-e aqdām-e moshrekīn
gasht. Fī al-ḥaqīqe har gūshī lāyeq-e estemāʿ na va abṣār-e marmūde-ye
gheyr-e ḥadīde qābel-e moshāhede na. Emrūz rūz-e dhekr va thanā-st. Ke
rā? Awlīyā, ānān ke be-kamāl-e esteqāmat bar amr qāʾemand va be-dhekr-e
Maqṣūd-e ʿālam nāṭeq. Har dhekr va har thanāʾī ān nofūs-e moṭmaʾenne-ye
rāżīye-ye marżīye rā sazāvār.

Today the rustling of the Tree beyond which there is no passing rises
high, and the sun of utterance shines resplendent from the horizon of
the heaven of meanings. The Path speaks with the most eloquent
expression, and the Balance makes mention in the most perfect words. Yet
the banner of ``There is none in the realm but Him'' is planted and
raised aloft, for the world is lifeless and despondent---indeed, dead.
There is no ear to listen and no eye to behold, except for a few whom
the veils of names and the coverings of allusions and words did not hold
back. They left behind all but the True One and joined themselves to
Him. They heard the call, drank the choice wine, and arose for the Cause
with such resolve that it caused the hearts of the aggressors to tremble
and the feet of those who associate partners with God to falter. In
truth, not every ear is fit to hear, nor are diseased eyes, lacking
clear vision, capable of beholding. Today is the day of remembrance and
praise. For whom? For the friends, those who stand firm in the Cause
with perfect steadfastness and speak in remembrance of the Desire of the
world. Every remembrance and every praise befits those assured souls,
content with God and pleasing to Him.

\subsection{3}\label{section-2234}

\begin{otherlanguage}{arabic}

سبحانک یا من ناح المقرّبون فی فراقک و صاح المخلصون فی هجرک اسألک بانوار
صبح ظهورک الّذی به احییت عبادک و خلقک و بشّرت من فی البلاد بوصالک و لقائک
ان تؤیّد اولیائک علی ما ترتفع به ظهورات قوّتک و شئونات قدرتک ای ربّ زیّن
بأسمائهم دفاتر العالم و کتب الأمم انّک انت المقتدر علی ما تشآء لا اله الّا
انت المؤیّد المعطی المقتدر العلیم الحکیم یا الهی و اله العالم و مقصودی و
مقصود الأمم اسألک بندائک الأحلی الّذی ارتفع بین الأرض و السّمآء و باشراقات
انوار وجهک من الأفق الأعلی بأن تحفظ من قام علی خدمة امرک و نطق بثنائک ثمّ
ایّده بجنود الغیب و الشّهاده ثمّ انصره بقبیل من الملائکة و الرّوح انّک انت
المقتدر الّذی لا تعجزک قدرة الأمرآء و لا سطوتهم و لا شبهات العلمآء و
اشاراتهم لا اله الّا انت المهیمن المقتدر العزیز العظیم

\end{otherlanguage}

Subḥānaka yā man nāḥa al-muqarrabūn fī firāqika wa-ṣāḥa al-mukhliṣūn fī
hajrika. Asʾaluka bi-anwāri ṣubḥi ẓuhūrika alladhī bihi aḥyayta ʿibādaka
wa-khalqaka wa-bashsharta man fī al-bilād bi-wiṣālika wa-liqāʾika an
tuʾayyida awliyāʾaka ʿalā mā tartafiʿu bihi ẓuhūrātu quwwatika
wa-shuʾūnātu qudratika. Ay rabb, zayyin bi-asmāʾihim dafātira al-ʿālam
wa-kutuba al-umam. Innaka anta al-muqtadir ʿalā mā tashāʾ, lā ilāha illā
anta al-muʾayyid al-muʿṭī al-muqtadir al-ʿalīm al-ḥakīm. Yā ilāhī
wa-ilāha al-ʿālam wa-maqṣūdī wa-maqṣūda al-umam, asʾaluka bi-nidāʾika
al-aḥlā alladhī irtafaʿa bayna al-arḍ wa-al-samāʾ wa-bi-ishrāqāti anwāri
wajhika min al-ufuq al-aʿlā bi-an taḥfaẓa man qāma ʿalā khidmati amrika
wa-naṭaqa bi-thanāʾika, thumma ayyidhu bi-junūdi al-ghayb wa-al-shahāda,
thumma unṣurhu bi-qabīlin min al-malāʾika wa-al-rūḥ. Innaka anta
al-muqtadir alladhī lā tuʿjizuka qudratu al-umarāʾ wa-lā saṭwatuhum
wa-lā shubuhātu al-ʿulamāʾ wa-ishārātuhum. Lā ilāha illā anta
al-muhaymin al-muqtadir al-ʿazīz al-ʿaẓīm.

Glory be to You, O You in separation from whom those brought near have
lamented, and at whose absence the sincere have cried out. I ask You, by
the lights of the dawn of Your revelation, through which You have
brought Your servants and Your creation to life and given those in every
land the glad news of union with You and attainment to Your presence, to
assist Your friends in that through which the manifestations of Your
might and the evidences of Your power may be exalted. O Lord, adorn the
records of the world and the books of the nations with their names. You
indeed have power over whatever You will. There is no God but You, the
Helper, the Bestower, the All-Powerful, the All-Knowing, the All-Wise. O
my God and the God of the world, my Desire and the Desire of the
nations! I ask You, by Your sweetest call raised between earth and
heaven, and by the radiance of the lights of Your face from the most
exalted horizon, to protect whoever has arisen to serve Your Cause and
spoken Your praise. Then strengthen him with the hosts of the unseen and
the seen, and aid him with a company of angels and the Spirit. You are
indeed the All-Powerful, whom neither the power of rulers nor their
might, nor the doubts raised by the learned nor their insinuations, can
thwart. There is no God but You, the All-Protecting, the All-Powerful,
the Almighty, the Most Great.

\subsection{4}\label{section-2235}

\begin{otherlanguage}{arabic}

یا محبوب فؤادی کنت ناظراً متفکّراً متحسّراً متفرّساً فی ظهورات قدرة ربّنا و فیما
ورد علیه و ما ظهر من الأحزاب اذاً فتح الباب و دخل من کان حاملاً کتابکم
الّذی وجد منه کلّ ذی شمّ عرف خلوصکم و حبّکم و اقبالکم الی اللّه ربّنا و ربّکم و
قیامکم علی خدمة الأمر فی ایّام فیها اخذت سطوت الغافلین سکّان مدائن الأسمآء
الّا من حفظه اللّه مولی الوری و سلطان الآخرة و الأولی فلمّا قرأت وجدت من کلّ
حرف من حروفاته ما علّمنی ما عندکم و عرّفنی ما انجذب به فؤادکم و عروقکم و
ظاهرکم و باطنکم امر اللّه الّذی تمسّکتم بحبله و تشبّثتم بذیله لاعلآء کلمته و
هدایة عباده و خلقه فلمّا عرفت طرت بأجنحة الاشتیاق الی اللّه مالک المیثاق
حضرت و عرضت الی ان انتهی اذاً ارتفع خفیف سدرة المنتهی و نطق مولی الوری
بما لا یقدر الخادم ان یذکره حقّ الذّکر اسأله بأن یمدّنی بفضله و جوده هل
یقدر کاتب ان یحصی امواج البحر اذ ظهرت و ارتفعت لا ونفسه الحقّ الّا بتأیید
من عنده و فضل من لدنه لا اله الّا هو محصی کلّشیء فی کتاب مبین

\end{otherlanguage}

Yā maḥbūba fuʾādī, kuntu nāẓiran mutafakkiran mutaḥassiran mutafarrisan
fī ẓuhūrāti qudrati rabbinā wa-fīmā warada ʿalayhi wa-mā ẓahara min
al-aḥzāb. Idhan futiḥa al-bāb wa-dakhala man kāna ḥāmilan kitābakum
alladhī wajada minhu kullu dhī shammin ʿarfa khulūṣikum wa-ḥubbikum
wa-iqbālikum ilā Allāh rabbinā wa-rabbikum wa-qiyāmikum ʿalā khidmati
al-amr fī ayyāmin fīhā akhadhat saṭwatu al-ghāfilīn sukkāna madāʾin
al-asmāʾ illā man ḥafiẓahu Allāh mawlā al-warā wa-sulṭān al-ākhira
wa-al-ūlā. Fa-lammā qaraʾtu wajadtu min kulli ḥarfin min ḥurūfātihi mā
ʿallamanī mā ʿindakum wa-ʿarrafanī mā injadhaba bihi fuʾādukum
wa-ʿurūqukum wa-ẓāhirukum wa-bāṭinukum, amru Allāh alladhī tamassaktum
bi-ḥablihi wa-tashabbathtum bi-dhaylihi li-iʿlāʾi kalimatihi wa-hidāyati
ʿibādihi wa-khalqihi. Fa-lammā ʿaraftu ṭirtu bi-ajniḥati al-ishtiyāq ilā
Allāh mālik al-mīthāq. Ḥāżer shod va ʿarż shod? Ḥaḍartu wa-ʿaraḍtu ilā
an intahā. Idhan irtafaʿa khafīfu sidrat al-muntahā wa-naṭaqa mawlā
al-warā bi-mā lā yaqdiru al-khādim an yadhkurahu ḥaqqa al-dhikr.
Asʾaluhu bi-an yumiddanī bi-faḍlihi wa-jūdihi. Hal yaqdiru kātib an
yuḥṣiya amwāja al-baḥr idh ẓaharat wa-irtafaʿat? Lā wa-nafsihi al-ḥaqq,
illā bi-taʾyīdin min ʿindihi wa-faḍlin min ladunhu. Lā ilāha illā huwa
muḥṣī kulli shayʾin fī kitābin mubīn.

O beloved of my heart! I was looking, reflecting, grieving, and
pondering the manifestations of our Lord's power, what had befallen Him,
and what had proceeded from the factions. Then the door opened, and the
bearer of your letter entered---a letter from which everyone endowed
with a sense of smell perceived the fragrance of your sincerity, your
love, your turning toward God, our Lord and your Lord, and your resolve
to serve the Cause in days when the might of the heedless had seized the
inhabitants of the cities of names, except those whom God, the Lord of
all and the Sovereign of the next world and this one, had protected.
When I read it, I found in each of its letters that which acquainted me
with what was in your hearts and made known to me what had drawn your
heart, your very veins, your outward and inward being: the Cause of God,
to whose cord you have held fast and to whose hem you have clung, to
exalt His Word and guide His servants and His creation. When I
understood this, I flew on the wings of longing toward God, the Lord of
the Covenant. I entered His presence and presented it to the end. Then
the gentle rustling of the Tree beyond which there is no passing arose,
and the Lord of all spoke words that this servant cannot fittingly
recount. I ask Him to sustain me through His grace and generosity. Can
any writer number the waves of the sea as they appear and rise? No, by
His own true Self, except through assistance from Him and grace from His
presence. There is no God but Him, who records all things in a clear
Book.

\subsection{5}\label{section-2236}

\begin{otherlanguage}{arabic}

و فی آخر بیاناته روح من فی ملکوت الأمر و الخلق فداه قال عزّ بیانه و جلّ
برهانه

\end{otherlanguage}

Wa-fī ākhiri bayānātihi, rūḥu man fī malakūt al-amr wa-al-khalq fidāhu,
qāla, ʿazza bayānuhu wa-jalla burhānuhu:

And at the close of His utterances---may the souls of all in the
kingdoms of command and creation be sacrificed for Him---He said,
exalted be His utterance and glorious His proof:

\subsection{6}\label{section-2237}

\begin{otherlanguage}{arabic}

هو النّاظر الشّاهد العالم السّمیع البصیر

\end{otherlanguage}

Huwa al-nāẓir al-shāhid al-ʿālim al-samīʿ al-baṣīr.

He is the Watchful, the Witness, the Knowing, the All-Hearing, the
All-Seeing.

\subsection{7}\label{section-2238}

\begin{otherlanguage}{arabic}

سدرة البیان فی قطب الامکان تنادی انّه لا اله الّا انا المهیمن القیّوم سدرة
المنتهی علی الأرض النّورآء خلف قلزم الکبریآء تاللّه قد اتی المکنون بسلطان
مشهود سدرة الایقان المغرسة علی ارض الزّعفران تنادی قد اتی المالک الملک
للّه مالک الوجود شجرة البرهان علی البقعة البیضآء تنطق قد تضوّع عرف القمیص
و اشرق النّور من افق الظّهور اسرعوا یا ملأ الارض و لا تتّبعوا کلّ جاهل محجوب
هذا یوم فیه خضعت الجبال اذ اتی الغنیّ المتعال برایة الاقتدار و علم انّه لا
اله الّا انا مالک الغیب و الشّهود قل یا ملأ الایران اتّقوا الرّحمان و
لاتکونوا کالّذینهم لا یعلمون انّا نفخنا فی الصّور و هو قلمی الأعلی و احیینا
به من کان فی علمی المحفوظ

\end{otherlanguage}

Sidratu al-bayān fī quṭb al-imkān tunādī innahu lā ilāha illā anā
al-muhaymin al-qayyūm. Sidratu al-muntahā ʿalā al-arḍ al-nawrāʾ khalfa
qulzum al-kibriyāʾ: ta-Allāh qad atā al-maknūn bi-sulṭānin mashhūd.
Sidratu al-īqān al-mughrasa ʿalā arḍ al-zaʿfarān tunādī qad atā
al-mālik, al-mulku li-llāh mālik al-wujūd. Shajaratu al-burhān ʿalā
al-buqʿa al-bayḍāʾ tanṭiqu qad taḍawwaʿa ʿarfu al-qamīṣ wa-ashraqa
al-nūr min ufuq al-ẓuhūr. Asriʿū yā malaʾ al-arḍ wa-lā tattabiʿū kulla
jāhilin maḥjūb. Hādhā yawmun fīhi khaḍaʿat al-jibāl idh atā al-ghaniyy
al-mutaʿāl bi-rāyat al-iqtidār wa-ʿalami innahu lā ilāha illā anā mālik
al-ghayb wa-al-shuhūd. Qul yā malaʾ al-Īrān ittaqū al-raḥmān wa-lā
takūnū ka-alladhīna hum lā yaʿlamūn. Innā nafakhnā fī al-ṣūr wa-huwa
qalamī al-aʿlā wa-aḥyaynā bihi man kāna fī ʿilmī al-maḥfūẓ.

The Tree of Utterance at the pivot of the contingent world calls out:
``There is no God but Me, the All-Protecting, the Self-Subsisting.'' The
Tree beyond which there is no passing, upon the luminous earth beyond
the ocean of grandeur: ``By God, the Hidden One has come with manifest
sovereignty!'' The Tree of Certitude, planted in the land of saffron,
calls out: ``The Sovereign has come! Sovereignty belongs to God, the
Lord of existence.'' The Tree of Proof upon the white spot declares:
``The fragrance of the garment has spread, and the light has shone from
the horizon of revelation. Hasten, people of the earth, and do not
follow anyone who is ignorant and veiled.'' This is a day on which the
mountains have bowed down, for the Self-Sufficient, the Most Exalted,
has come with the banner of power and the standard of ``There is no God
but Me, the Lord of the unseen and the seen.'' Say: ``People of Iran,
fear the All-Merciful and do not be among those who do not know. We have
blown the trumpet, which is My Most Exalted Pen, and through it We have
brought to life those who were in My preserved knowledge.''

\subsection{8}\label{section-2239}

\begin{otherlanguage}{arabic}

قل یا ملأ البیان هل سمعتم و انکرتم و هل رأیتم و اعرضتم اتّقوا اللّه و
انصفوا فی هذا الأمر الّذی اذ ظهر اضطربت النّفوس الّا من اخذه اللّه و انقذه
من ظلمات الأوهام و رفعه الی مقام سمع ندآء الملکوت یا ورقآء اسمع النّدآء
کما سمعت من قبل و من بعد

\end{otherlanguage}

Qul yā malaʾ al-bayān hal samiʿtum wa-ankartum wa-hal raʾaytum
wa-aʿraḍtum? Ittaqū Allāh wa-anṣifū fī hādhā al-amr alladhī idh ẓahara
iḍṭarabat al-nufūs illā man akhadhahu Allāh wa-anqadhahu min ẓulumāt
al-awhām wa-rafaʿahu ilā maqāmin samiʿa nidāʾ al-malakūt. Yā Warqāʾ
ismaʿ al-nidāʾ kamā samiʿta min qablu wa-min baʿd.

Say: ``People of the Bayān, have you heard and yet denied? Have you seen
and yet turned away? Fear God and judge fairly concerning this Cause, at
whose appearance souls were shaken, except those whom God took hold of,
rescued from the darkness of vain imaginings, and raised to a station
where they heard the call of the Kingdom.'' O Warqāʾ, hear the call as
you have heard it before and since.

\subsection{9}\label{section-2240}

\begin{otherlanguage}{arabic}

قد حضر العبد الحاضر بکتابک فلمّا قرأه لدی المظلوم تحرّک قلم ربّک القیّوم و
انزل لک ما افتر به ثغر البیان و اخذ الاهتزاز من فی القبور قد انزلنا لک
ما کان نوراً لکتب العالم و صراطاً مستقیماً لأهل الأدیان و رکناً عظیماً لمن
اقرّ بهذا النّبأ العظیم قل یا ملأ المعرضین اسمعوا ثمّ انصفوا لعمر اللّه انّ
النّدآء یجذبکم الی مقام تضعون ما عندکم رجآء ما عند اللّه العزیز الودود قل
هذا مضمار البرهان فأتوا بما عندکم لعمر اللّه لا تجول فیه الّا فوارس
الایقان الّذین نبذوا العالم مقبلین الی اللّه مالک الیوم الموعود کذلک اظهر
البحر اصداف الحکمة و البیان و السّدره اثمارها ولکنّ القوم لا یعلمون

\end{otherlanguage}

Qad ḥaḍara al-ʿabd al-ḥāḍir bi-kitābika, fa-lammā qaraʾahu ladā
al-maẓlūm taḥarraka qalamu rabbika al-qayyūm wa-anzala laka mā iftarra
bihi thaghr al-bayān wa-akhadha al-ihtizāzu man fī al-qubūr. Qad anzalnā
laka mā kāna nūran li-kutub al-ʿālam wa-ṣirāṭan mustaqīman li-ahl
al-adyān wa-ruknan ʿaẓīman li-man aqarra bi-hādhā al-nabaʾ al-ʿaẓīm. Qul
yā malaʾ al-muʿriḍīn ismaʿū thumma anṣifū. La-ʿamru Allāh inna al-nidāʾ
yajdhibukum ilā maqāmin taḍaʿūna mā ʿindakum rajāʾa mā ʿinda Allāh
al-ʿazīz al-wadūd. Qul hādhā miḍmār al-burhān faʾtū bi-mā ʿindakum.
La-ʿamru Allāh lā tajūlu fīhi illā fawāris al-īqān alladhīna nabadhū
al-ʿālam muqbilīna ilā Allāh mālik al-yawm al-mawʿūd. Kadhālika aẓhara
al-baḥru aṣdāfa al-ḥikma wa-al-bayān wa-al-sidratu athmārahā wa-lākinna
al-qawma lā yaʿlamūn.

The servant attending here arrived with your letter. When he read it in
the presence of the Wronged One, the Pen of your Lord, the
Self-Subsisting, began to move and revealed for you that which made the
mouth of utterance smile and stirred those in the graves. We have
revealed for you that which is a light to the books of the world, a
straight path for the followers of religions, and a mighty pillar for
whoever acknowledges this momentous announcement. Say: ``You who turn
away, listen and then judge fairly. By the life of God, the call draws
you to a station where you lay aside what you possess in the hope of
what is with God, the Almighty, the All-Loving.'' Say: ``This is the
arena of proof; bring forward what you have. By the life of God, none
ride in it but the horsemen of certitude, who have cast the world aside
and turned toward God, the Lord of the promised Day.'' Thus has the sea
brought forth the shells of wisdom and utterance, and the Tree its
fruits, but the people do not know.

\subsection{10}\label{section-2241}

\begin{otherlanguage}{persian}

بلسان پارسی ندای مظلوم را بشنو ظاهر شد آنچه مکنون بود و نازل شد آنچه در
سماء مشیّت لازال مخزون و در لوح محفوظ محفوظ ولکن جزاء اعمال کلّ را از
فیوضات یوم مآل محروم ساخت چنانچه در حزب فرقان اعمال و جزاء آن دیده شد
مخصوص حزب شیعه یا ورقا در عدم درایت نفوس موهومهٴ اهل بیان تفکّر نما مع
آنکه حجبات خرق شد و ظنون و اوهام و کفر و نفاق و ظلم و اعتساف حزب قبل را
ببصر ظاهر مشاهده نمودند حال مجدّد در ترتیب حزب شیعه مشغولند میخواهند قومی
را شیعه نمایند باری اوهام و اسماء ناس را از بحر معانی محروم نمود و از
توحید حقیقی بی‌بهره ساخت فی‌الحقیقه در یک مقام اسماء اصنامند چنانچه مشاهده
شده و میشود الی حین حزب قبل متنبّه نشدند و از نوم غفلت برنخاستند هر یوم
نار ظلم مشتعل‌تر و آتش ضغینه و بغضا روشن‌تر امروز بر کلّ لازم و واجب که بر
نصرت امر قیام نمایند و آن حفظ نفوس ضعیفه است از ظنون و اوهام

\end{otherlanguage}

Be-lesān-e pārsī nedā-ye Maẓlūm rā beshnow. Ẓāher shod ānche maknūn būd
va nāzel shod ānche dar samā-ye mashīyyat lā-zāl makhzūn va dar lawḥ-e
maḥfūẓ maḥfūẓ, va-lāken jazā-ye aʿmāl koll rā az foyūżāt-e yawm-e maʾāl
maḥrūm sākht, chenānche dar ḥezb-e Forqān aʿmāl va jazā-ye ān dīde shod,
makhṣūṣ ḥezb-e Shīʿe. Yā Varqā, dar ʿadam-e derāyat-e nofūs-e mowhūme-ye
ahl-e Bayān tafakkor namā. Maʿ-e ānke ḥojobāt kharq shod va ẓonūn va
owhām va kofr va nefāq va ẓolm va eʿtesāf-e ḥezb-e qabl rā be-baṣar-e
ẓāher moshāhede namūdand, ḥāl mojaddad dar tartīb-e ḥezb-e Shīʿe
mashghūland, mīkhāhand qowmī rā Shīʿe namāyand. Bārī owhām va asmā nās
rā az baḥr-e maʿānī maḥrūm namūd va az towḥīd-e ḥaqīqī bī-bahre sākht.
Fī al-ḥaqīqe dar yek maqām asmā aṣnāmand, chenānche moshāhede shode va
mīshavad. Ilā ḥīn ḥezb-e qabl motanabbeh nashodand va az nowm-e ghaflat
barnakhāstand. Har yowm nār-e ẓolm moshtaʿel-tar va ātash-e żaghīne va
baghżā rowshan-tar. Emrūz bar koll lāzem va vājeb ke bar noṣrat-e amr
qīyām namāyand va ān ḥefẓ-e nofūs-e żaʿīfe ast az ẓonūn va owhām.

Hear the call of the Wronged One in the Persian tongue. What was hidden
has become manifest, and what had ever been stored in the heaven of the
divine will and preserved in the Preserved Tablet has been sent down.
Yet the recompense for their deeds deprived all of the outpourings of
the Day of Return, just as the deeds and their recompense were seen
among the people of the Qurʾān, particularly the Shiʿi community. O
Warqāʾ, reflect on the lack of discernment of those deluded souls among
the people of the Bayān. Although the veils were torn apart, and they
saw with their outward eyes the conjectures, vain imaginings, unbelief,
hypocrisy, oppression, and tyranny of the former community, they are now
once again occupied in organizing a Shiʿi community; they wish to make a
people into Shiʿis. In short, vain imaginings and names have deprived
people of the ocean of meanings and denied them a share in true
recognition of divine unity. In truth, in one respect names are idols,
as has been and continues to be witnessed. Even now the former community
has not taken heed or risen from the sleep of heedlessness. Each day the
fire of oppression blazes more fiercely, and the flame of malice and
hatred burns more brightly. Today it is necessary and incumbent upon all
to arise to aid the Cause, which consists in protecting weak souls from
conjectures and vain imaginings.

\subsection{11}\label{section-2242}

\begin{otherlanguage}{persian}

اولیاء آن ارض را از قبل مظلوم تکبیر برسانید بگو جهد نمائید شاید بانوار
آفتاب توحید حقیقی فائز شوید هر نفسی بآن فائز شد او خود را از عالم فارغ و
آزاد مشاهده نماید هنیئاً للعاملین و هنیئاً للعارفین البهآء الظّاهر اللّائح
من افق سمآء ملکوتی علیک و علی من یسمع قولک فی امر اللّه ربّ العالمین انتهی

\end{otherlanguage}

Awlīyā-ye ān arż rā az qebal-e Maẓlūm takbīr berasānīd. Begū jahd
namāʾīd shāyad be-anvār-e āftāb-e towḥīd-e ḥaqīqī fāʾez shavīd. Har
nafsī be-ān fāʾez shod ū khod rā az ʿālam fāregh va āzād moshāhede
namāyad. Hanīʾan li-l-ʿāmilīn wa-hanīʾan li-l-ʿārifīn. Al-bahāʾ al-ẓāhir
al-lāʾiḥ min ufuqi samāʾi malakūtī ʿalayka wa-ʿalā man yasmaʿu qawlaka
fī amri Allāh rabb al-ʿālamīn. Intahā.

Convey to the friends in that land, on behalf of the Wronged One, the
greeting ``God is Most Great.'' Say: ``Strive, that you may attain the
lights of the sun of true recognition of divine unity. Whoever attains
it will find himself free and released from the world.'' Blessed are
those who act, and blessed are those who recognize. May the glory
manifest and shining from the horizon of the heaven of My Kingdom rest
upon you and upon whoever listens to your words concerning the Cause of
God, the Lord of the worlds. End.

\subsection{12}\label{section-2243}

\begin{otherlanguage}{persian}

یا محبوب فؤادی آفتاب بیان رحمن بشأنی مشرق و لایح که بیان عالم از اوّل
ابداع تا حین عند ظهور و نزولش خاضع و خاشع دیگر این مفقود چه عرض نماید و
چه اظهار دارد این خادم فانی متحیّر که ناعقین چه میگویند و بچه متمسّکند حال
قدر و مقام بعضی معلوم میشود مع آنکه کلّ جسارت حزب قبل را ببصر ظاهر مشاهده
نموده‌اند و شرک و غفلت و نفاق و ظلم آن نفوس را دیده مع‌ذلک گاهی مرآت گاهی
وصیّ گاهی وکیل و نجیب و نقیب و امثال آن ذکر مینمایند سبحان‌اللّه هرچه مخالف
اهواء نفوس معرضه از قبل نازل شده او را انکار نموده‌اند بیان فارسی که
فی‌الحقیقه از اعظم کتب نقطهٴ اولی روح ما سواه فداه است بالمرّه انکار
کرده‌اند گفته‌اند تحریف شده باید کلّ بنوحه مشغول شویم و وااسفا گوئیم کذبهای
قبل مبعوث شده و مفتریات قبل بمیان آمده انّا للّه و انّا الیه راجعون

\end{otherlanguage}

Yā maḥbūba fuʾādī, āftāb-e bayān-e Raḥmān be-shaʾnī moshreq va lāyeḥ ke
bayān-e ʿālam az avval-e ebdāʿ tā ḥīn ʿend-e ẓohūr va nozūlash khāżeʿ va
khāsheʿ. Dīgar īn mafqūd che ʿarż namāyad va che eẓhār dārad? Īn
khādem-e fānī motaḥayyer ke nāʿeqīn che mīgūyand va be-che
motamassekand. Ḥāl qadr va maqām-e baʿżī maʿlūm mīshavad. Maʿ-e ānke
koll jasārat-e ḥezb-e qabl rā be-baṣar-e ẓāher moshāhede nomūde-and va
sherk va ghaflat va nefāq va ẓolm-e ān nofūs rā dīde, maʿa-dhālek gāhī
merʾāt, gāhī vaṣī, gāhī vakīl va najīb va naqīb va amthāl-e ān dhekr
mīnamāyand. Sobḥān Allāh, harche mokhālef-e ahvā-ye nofūs-e moʿreże az
qabl nāzel shode ū rā enkār nomūde-and. Bayān-e Fārsī ke fī al-ḥaqīqe az
aʿẓam-e kotob-e Noqṭe-ye Ūlā, rūḥ-e mā sevāh fedāh, ast be-al-marre
enkār karde-and. Gofte-and taḥrīf shode. Bāyad koll be-nowḥe mashghūl
shavīm va vā asafā gūʾīm. Kedhb-hā-ye qabl mabʿūth shode va moftarayāt-e
qabl be-mīyān āmade. Innā li-llāh wa-innā ilayhi rājiʿūn.

O beloved of my heart! The sun of the All-Merciful's utterance shines
with such radiance that all the utterance of the world, from the
beginning of creation until now, bows in humility and reverence at its
appearance and revelation. What, then, can this lost one submit, and
what can he express? This mortal servant is bewildered by what those
clamorous voices are saying and what they cling to. The worth and
station of certain people are now becoming clear. Although all have seen
with their outward eyes the audacity of the former community and
witnessed the idolatry, heedlessness, hypocrisy, and oppression of those
souls, they nevertheless speak now of a mirror, now of a successor, now
of a deputy, a noble one, a chief, and the like. Glory be to God! They
have denied whatever was previously revealed that runs counter to the
desires of those who turn away. They have utterly denied the Persian
Bayān, which is in truth among the greatest books of the Primal
Point---may the soul of all besides Him be sacrificed for Him. They have
said it has been altered. We must all lament and cry, ``Alas!'' The lies
of the past have been resurrected, and the former fabrications have
reappeared. We belong to God, and to Him we return.

\subsection{13}\label{section-2244}

\begin{otherlanguage}{arabic}

یا محبوب فؤادی بر عظمت ظهور و احاطه‌اش و همچنین ما نزّل من سمآء مشیّته
منصفین و متبصّرین عالم و عارفند و بمثابه آفتاب انوار ظهور مالک الرّقاب را
مشاهده مینمایند مع‌ذلک این قوم نسناس مکدّرهٴ مغبرّهٴ کثیفه مثل علماء حزب
شیعه که در زاویه‌های مدارس بکثافت اظهار زهد مینمودند و بخباثت ادّعای علم و
ایمان مشاهده میشوند فی‌الحقیقه لایق این بساط نبوده و نیستند الحمد للّه
الّذی قدّس بساطه عن لقائهم و نوره عن مشاهدتهم و ظهوره عن عرفانهم و ندائه
عن اصغائهم به مکر و تزویر و خدعه تمسّک نموده‌اند بعضی اعمال از آن نفوس
غافلهٴ جاهله ذکر نموده‌اند که هر صاحب عقلی متحیّر اعاذنا اللّه و ایّاکم من
شرّ هؤلآء

\end{otherlanguage}

Yā maḥbūb-i fuʾādī, bar ʿaẓamat-i ẓuhūr va iḥāṭe-ash va hamchunīn mā
nuzzila min samāʾi mashīʾatihi munṣifīn va mutabaṣṣirīn ʿālim va
ʿārifand va bi-mathābe-yi āftāb anvār-i ẓuhūr-i Mālik al-Riqāb rā
mushāhede mīnamāyand. Maʿa dhālik īn qawm-i nasnās-i mukaddare-yi
mughabbare-yi kathīfe mithl-i ʿulamā-yi ḥizb-i Shīʿe ki dar zāviyehā-yi
madāris bi-kithāfat iẓhār-i zuhd mīnamūdand va bi-khibāthat iddiʿā-yi
ʿilm va īmān mushāhede mīshavand. Fi'l-ḥaqīqe lāyiq-i īn bisāṭ nabūde va
nīstand. Al-ḥamdu li-llāhi alladhī qaddasa bisāṭahu ʿan liqāʾihim va
nūrahu ʿan mushāhadatihim va ẓuhūrahu ʿan ʿirfānihim va nidāʾahu ʿan
iṣghāʾihim. Bi-makr va tazvīr va khudʿe tamassuk namūde-and. Baʿḍī aʿmāl
az ān nufūs-i ghāfile-yi jāhile dhikr namūde-and ki har ṣāḥib-i ʿaqlī
mutaḥayyir. Aʿādhanā Allāhu va iyyākum min sharri hāʾulāʾ.

O beloved of my heart! The fair-minded and discerning know and recognize
the greatness of the Revelation, its all-encompassing power, and
likewise what has been sent down from the heaven of His will. They
behold, as plainly as the sun, the lights of the revelation of the Lord
of all. Nevertheless, this tribe of apish creatures, defiled,
dust-covered, and filthy, are seen to resemble the divines of the Shiite
community, who, in the corners of their seminaries, made a filthy show
of asceticism and wickedly claimed knowledge and faith. They have never
been, and are not now, worthy of this court. Praise be to God, who has
sanctified His court from meeting them, His light from their sight, His
Revelation from their recognition, and His call from their hearing. They
have resorted to scheming, deceit, and trickery. Some deeds of those
heedless, ignorant souls have been recounted that leave every rational
person bewildered. May God protect us and you from the evil of these
people.

\subsection{14}\label{section-2245}

\begin{otherlanguage}{persian}

اینکه مرقوم داشتند مراسلات اینعبد کمتر در حرکت است فی‌الحقیقه تصدیق
مینمایم اینکلمه را چه که باید در کلّ حین ارسال نمایم آنچه را که عرف عنایت
و شفقت محبوبنا و محبوبکم و مقصودنا و مقصودکم و مقصود من فی السّموات و
الارض از آن متضوّع باشد ولکن یا محبوب فؤادی بشأنی از اطراف نامه‌ها رسیده و
میرسد که فی‌الحقیقه اینعبد فرصت اینکه ما فی الضّمیر را نزد اولیا اظهار
نماید نبوده و نیست دو سنه گذشته و جواب بعضی از آقایان علیهم بهآء اللّه و
عنایته در عهدهٴ تعویق بوده و هست از حقّ جلّ جلاله سائل و آمل که تأیید
فرماید و توفیق عطا نماید اوست قادر و توانا ولکن جواب دستخطّهای آنحضرت
غالباً اظهار شده خدای واحد شاهد هر هنگام رقیمهٴ آنمحبوب رسیده فرح و
انبساط جدیدی حاصل و از آنگذشته چندیست در این ارض هم بعضی گفتگوها بمیان
آمده حکمت هم گاهی ملاحظه میشود

\end{otherlanguage}

Īnki marqūm dāshtand murāsalāt-i īn ʿabd kamtar dar ḥarakat ast,
fi'l-ḥaqīqe taṣdīq mīnamāyam īn kalame rā, chi ki bāyad dar kull-i ḥīn
irsāl namāyam ānchi rā ki ʿarf-i ʿināyat va shafaqat-i maḥbūbinā va
maḥbūbikum va maqṣūdinā va maqṣūdikum va maqṣūdi man fī al-samāvāti va
al-arḍ az ān mutaḍavviʿ bāshad. Va-lākin yā maḥbūb-i fuʾādī, bi-shaʾnī
az aṭrāf nāmehā rasīde va mīrasad ki fi'l-ḥaqīqe īn ʿabd furṣat-i īnki
mā fī al-ḍamīr rā nazd-i awliyā iẓhār namāyad nabūde va nīst. Du sane
gudhashte va javāb-i baʿḍī az āqāyān, ʿalayhim bahāʾu Allāhi va
ʿināyatuhu, dar ʿuhde-yi taʿvīq būde va hast. Az Ḥaqq, jalla jalāluhu,
sāʾil va āmil ki taʾyīd farmāyad va tawfīq ʿaṭā namāyad. Ūst qādir va
tavānā. Va-lākin javāb-i dastkhaṭṭhā-yi ān ḥaḍrat ghāliban iẓhār shude.
Khudā-yi vāḥid shāhid, har hangām raqīme-yi ān maḥbūb rasīde faraḥ va
inbisāṭ-i jadīdī ḥāṣil, va az ān gudhashte chandīst dar īn arḍ ham baʿḍī
guftugūhā bi-miyān āmade. Ḥikmat ham gāhī mulāḥaẓe mīshavad.

As to your remark that this servant's correspondence has been
infrequent, I fully acknowledge the truth of this statement, for I
should send at every moment something fragrant with the loving care and
compassion of our Beloved and yours, our Desired One and yours, and the
Desired One of all in the heavens and on earth. Yet, O beloved of my
heart, letters have arrived and continue to arrive from all quarters in
such numbers that this servant has truly had, and still has, no
opportunity to express to the friends what is in his heart. Two years
have passed, and replies to some of the gentlemen---may God's glory and
loving care rest upon them---have been and remain overdue. I ask and
hope that God, exalted be His glory, will grant assistance and success.
He is powerful and mighty. Most of your letters, however, have been
answered. The one God is my witness: whenever a letter from that beloved
friend has arrived, it has brought fresh joy and delight. Besides this,
for some time certain discussions have arisen in this land as well, and
prudence must sometimes be observed.

\subsection{15}\label{section-2246}

\begin{otherlanguage}{persian}

و از همه گذشته آنحضرت میدانند که لیلاً و نهاراً در حضور بتحریر آیات الهی
مشغول ولکن از خدمات اینعبد آنمحبوب هم قسمت و نصیب دارند اینکلمه بشهادت
مقصود عالمیان مزیّن هنوا لحضرتک ولی ذکر محبوبی جناب امین علیه بهآء اللّه و
عنایته را فرموده بودند خود ایشانهم مکتوبی باین بنده مرقوم داشته‌اند لازال
در جمیع مکاتب ذکر آنحضرت را نموده‌اند فی‌الحقیقه خیرخواه اولیا بوده و
هستند از حقّ سائل و آمل که ایشانرا در جمیع احوال مؤیّد فرماید انّه هو
المشفق انّه هو الکریم

\end{otherlanguage}

Va az hame gudhashte ān ḥaḍrat mīdānand ki laylan va nahāran dar ḥuḍūr
bi-taḥrīr-i āyāt-i ilāhī mashghūl, va-lākin az khadamāt-i īn ʿabd ān
maḥbūb ham qismat va naṣīb dārand. Īn kalame bi-shahādat-i Maqṣūd-i
ʿālamiyān muzayyan. {[}corrupted text{]} li-ḥaḍratika. Valī dhikr-i
maḥbūbī Jināb-i Amīn, ʿalayhi bahāʾu Allāhi va ʿināyatuhu, rā farmūde
būdand. Khud-i īshān ham maktūbī bi-īn bande marqūm dāshte-and. Lā yazāl
dar jamīʿ-i makātīb dhikr-i ān ḥaḍrat rā namūde-and. Fi'l-ḥaqīqe
khayrkhvāh-i awliyā būde va hastand. Az Ḥaqq sāʾil va āmil ki īshān rā
dar jamīʿ-i aḥvāl muʾayyad farmāyad. Innahu huwa al-mushfiqu, innahu
huwa al-karīm.

Above all, you know that day and night I am occupied in His presence
with transcribing the divine verses. Yet that beloved friend also has a
share and portion in this servant's services. This statement is adorned
with the testimony of the Desired One of the worlds. {[}corrupted
original text{]} to you. You had also mentioned my beloved friend,
Jinab-i Amin---may God's glory and loving care rest upon him. He himself
has written a letter to this servant. He has always mentioned you in all
his letters. He has truly been, and remains, a well-wisher of the
friends. I ask and hope that God will assist him in every circumstance.
He is indeed the Compassionate; He is indeed the Generous.

\subsection{16}\label{section-2247}

\begin{otherlanguage}{persian}

اینکه بشارت ترک عادات مضرّهٴ بعضی از دوستان الهی را فرمودند بعد از عرض
اینفقره در ساحت امنع اقدس فرمودند از حقّ جلّ جلاله تقدیس و تنزیه کلّ را
میطلبیم امروز اعمال طیّبه و اخلاق مرضیّه جند اللّهند در ارض و باین جند حقّ
نصرت فرموده و میفرماید یا ایّها النّاطق باسمی و النّاظر الی افقی و الطّائر
فی هوآء حبّی اولیاء آن ارض را مکرّر از قبل اینمظلوم تکبیر برسان فرداً فرداً
در ساحت حقّ مذکورند و ذکرشان از قلم اعلی جاری

\end{otherlanguage}

Īnki bishārat-i tark-i ʿādāt-i muḍirre-yi baʿḍī az dūstān-i ilāhī rā
farmūdand, baʿd az ʿarḍ-i īn fiqre dar sāḥat-i amnaʿ-i aqdas farmūdand:
Az Ḥaqq, jalla jalāluhu, taqdīs va tanzīh-i kull rā mīṭalabīm. Imrūz
aʿmāl-i ṭayyibe va akhlāq-i marḍiyye jund-i Allāhand dar arḍ va bi-īn
jund Ḥaqq nuṣrat farmūde va mīfarmāyad. Yā ayyuhā al-nāṭiqu bi-ismī va
al-nāẓiru ilā ufuqī va al-ṭāʾiru fī havāʾi ḥubbī, awliyā-yi ān arḍ rā
mukarrar az qibal-i īn Maẓlūm takbīr birasān. Fardan fardan dar sāḥat-i
Ḥaqq madhkūrand va dhikrishān az Qalam-i Aʿlā jārī.

As to the glad news you conveyed that some of the friends of God have
abandoned harmful habits, after this matter was presented at the most
holy, most exalted threshold, He said: ``We ask God, exalted be His
glory, for the sanctification and purification of all. Today, good deeds
and praiseworthy character are the hosts of God on earth, and through
these hosts God has granted and continues to grant victory. O you who
speak in My name, gaze toward My horizon, and soar in the atmosphere of
My love! Repeatedly convey to the friends in that land, on behalf of
this Wronged One, the greeting `God is Most Great.' Each and every one
is remembered in the presence of God, and their remembrance flows from
the Most Exalted Pen.''

\subsection{17}\label{section-2248}

\begin{otherlanguage}{persian}

جناب میم و صاد علیه بهآء اللّه را تکبیر برسان از حقّ میطلبیم او را مؤیّد
فرماید بر آنچه سزاوار یوم او است و همچنین بر عملیکه عرفش منقطع نشود للّه
الحمد باقبال فائز شدند و بذکر قلم اعلی مزیّن انتهی

\end{otherlanguage}

Jināb-i Mīm va Ṣād, ʿalayhi bahāʾu Allāh, rā takbīr birasān. Az Ḥaqq
mīṭalabīm ū rā muʾayyad farmāyad bar ānchi sazāvār-i Yawm-i Ū ast va
hamchunīn bar ʿamalī ki ʿarfash munqaṭiʿ nashavad. Li-llāhi al-ḥamd,
bi-iqbāl fāʾiz shudand va bi-dhikr-i Qalam-i Aʿlā muzayyan. Intahā.

``Convey the greeting `God is Most Great' to Jinab-i Mim and Sad---may
God's glory rest upon him. We ask God to assist him in what befits His
Day, and likewise in a deed whose fragrance will never cease. Praise be
to God, he has attained the blessing of turning toward Him and has been
adorned with the remembrance of the Most Exalted Pen.'' End.

\subsection{18}\label{section-2249}

\begin{otherlanguage}{persian}

اینفانی هم از حقّ جلّ جلاله مسئلت مینماید که ایشانرا از شئونات عرضیّه مقدّس
فرماید و بشئوناتی مزیّن نماید که سبب اعلاء کلمة اللّه گردد ارواح از برای
این خلق شده و اجساد از برای این ظاهر گشته و آفتاب و ماه و اشیاء کلّ از
برای اینخدمت بوجود آمده انشآءاللّه به ذکر و ثنا و تبلیغ امر موفّق شوند
تکبیر و سلام خدمت ایشان میرسانم و از حقّ از برای ایشان تأیید میطلبم

\end{otherlanguage}

Īn fānī ham az Ḥaqq, jalla jalāluhu, masʾalat mīnamāyad ki īshān rā az
shuʾūnāt-i ʿaraḍiyye muqaddas farmāyad va bi-shuʾūnātī muzayyan namāyad
ki sabab-i iʿlā-yi Kalimat Allāh gardad. Arvāḥ az barā-yi īn khalq shude
va ajsād az barā-yi īn ẓāhir gashte va āftāb va māh va ashyā kull az
barā-yi īn khidmat bi-vujūd āmade. In shāʾ Allāh bi-dhikr va thanā va
tablīgh-i Amr muvaffaq shavand. Takbīr va salām khidmat-i īshān
mīrasānam va az Ḥaqq az barā-yi īshān taʾyīd mīṭalabam.

This mortal one also beseeches God, exalted be His glory, to sanctify
him from transitory concerns and adorn him with qualities that will
exalt the Word of God. Spirits were created for this, bodies were
brought into being for this, and the sun, the moon, and all things came
into existence for this service. God willing, may he be enabled to
remember and praise Him and to teach the Cause. I convey to him the
greeting `God is Most Great' and my salutations, and ask God to assist
him.

\subsection{19}\label{section-2250}

\begin{otherlanguage}{persian}

اینکه ذکر جناب سهراب و توجّه ایشان بافق اعلی فرمودند در ساحت امنع اقدس
بشرف اصغا فائز قوله جلّ ثنائه و عزّ برهانه یا سهراب معزّی آمد تا نگفته‌ها را
بگوید و نشنیده‌ها را بشنواند و آنچه گفته شده او را تمام کند حمد کن محبوب
عالمیانرا که بذروهٴ علیا و غایة قصوی که عرفان نفس حقّ جلّ جلاله است توجّه
نمودی خرق کن سبحات موهومه را و بشکن اصنام ظنون را تا بانوار آفتاب حقیقت
و ایقان منوّر شوی و فائز گردی اب شفیق آمد ابنا نشناختند او را و از عرفانش
محروم الّا معدودی چند که شبهات بطارقه و مطارنه و اساقفه و خوارنه ایشانرا
از مطلع نور احدیّت منع ننمود ندا مرتفع و لسانیکه بکلمهٴ بلی فائز شود
نایاب از حقّ بطلب آگاهی عطا فرماید و دانائی بخشد انّه هو المقتدر علی ما
یشآء بقوله کن فیکون انتهی

\end{otherlanguage}

Īnki dhikr-i Jināb-i Suhrāb va tavajjuh-i īshān bi-Ufuq-i Aʿlā
farmūdand, dar sāḥat-i amnaʿ-i aqdas bi-sharaf-i iṣghā fāʾiz. Qawluhu
jalla thanāʾuhu va ʿazza burhānuhu: Yā Suhrāb, Muʿazzī āmad tā naguftehā
rā bigūyad va nashunīdehā rā bishnavānad va ānchi gufte shude ū rā tamām
kunad. Ḥamd kun Maḥbūb-i ʿālamiyān rā ki bi-dhirve-yi ʿulyā va ghāyat-i
quṣvā ki ʿirfān-i nafs-i Ḥaqq, jalla jalāluhu, ast tavajjuh namūdī.
Kharq kun subuḥāt-i mawhūme rā va bishkan aṣnām-i ẓunūn rā tā bi-anvār-i
āftāb-i ḥaqīqat va īqān munavvar shavī va fāʾiz gardī. Ab-i shafīq āmad,
abnā nashinākhtand ū rā va az ʿirfānash maḥrūm, illā maʿdūdī chand ki
shubahāt-i baṭāriqe va maṭārine va asāqife va khavārine īshān rā az
Maṭlaʿ-i nūr-i aḥadiyyat manʿ nanamūd. Nidā murtafiʿ va lisānī ki
bi-kalame-yi balā fāʾiz shavad nāyāb. Az Ḥaqq biṭalab āgāhī ʿaṭā
farmāyad va dānāʾī bakhshad. Innahu huwa al-muqtadiru ʿalā mā yashāʾu
bi-qawlihi kun fa-yakūn. Intahā.

What you mentioned concerning Jinab-i Suhrab and his turning toward the
Most Exalted Horizon attained the honor of being heard at the most holy,
most exalted threshold. His words---exalted be His praise and mighty His
proof: ``O Suhrab! The Comforter has come to speak what has remained
unspoken, to make heard what has remained unheard, and to bring to
completion what has been said. Praise the Beloved of the worlds that you
have turned toward the loftiest summit and the ultimate goal, which is
recognition of God Himself, exalted be His glory. Rend the veils of vain
imaginings and shatter the idols of conjecture, that you may be
illumined by the rays of the Sun of truth and certitude and attain your
goal. The loving Father came, but the children did not recognize Him and
were deprived of recognizing Him, except for a few whom the doubts
raised by patriarchs, metropolitans, bishops, and priests did not keep
from the Dawning-Place of the light of divine unity. The call is raised,
but a tongue that attains the word `Yes' is nowhere to be found. Ask God
to grant awareness and bestow understanding. He indeed has power over
whatever He wills through His word `Be,' and it is.'' End.

\subsection{20}\label{section-2251}

\begin{otherlanguage}{persian}

و اینکه در بارهٴ حبیب روحانی جناب ع و با علیه بهآء اللّه الأبهی مرقوم
داشتید للّه الحمد ایشان بنفحات ایّام فائز گشته‌اند در کتاب الهی مذکور و در
صحیفهٴ حمرا از قلم اعلی مسطور این فضل را هیچ فضلی معادله ننماید و
اینمقام را هر مقامی خاضع از حقّ جلّ جلاله و عمّ نواله سائل و آمل که ایشانرا
تأیید فرماید بر حفظ اینمقام اعلی و رتبهٴ علیا نعیماً له و هنیئاً له و
مریئاً له و طوبی له و اینمراتب در پیشگاه حضور مالک ظهور عرض شد این بیان
از ملکوت عنایت محبوب عالمیان ظاهر قوله جلّت عظمته و عزّ سلطانه

\end{otherlanguage}

Va īnki dar bāre-yi ḥabīb-i rūḥānī Jināb-i ʿAyn va Bā, ʿalayhi bahāʾu
Allāhi al-Abhā, marqūm dāshtīd, li-llāhi al-ḥamd īshān bi-nafaḥāt-i
ayyām fāʾiz gashte-and. Dar Kitāb-i ilāhī madhkūr va dar Ṣaḥīfe-yi Ḥamrā
az Qalam-i Aʿlā masṭūr. Īn faḍl rā hīch faḍlī muʿādile nanamāyad va īn
maqām rā har maqāmī khāḍiʿ. Az Ḥaqq, jalla jalāluhu va ʿamma navāluhu,
sāʾil va āmil ki īshān rā taʾyīd farmāyad bar ḥifẓ-i īn maqām-i aʿlā va
rutbe-yi ʿulyā. Naʿīman lahu va hanīʾan lahu va marīʾan lahu va ṭūbā
lahu. Va īn marātib dar pīshgāh-i ḥuḍūr-i Mālik-i Ẓuhūr ʿarḍ shud. Īn
bayān az malakūt-i ʿināyat-i Maḥbūb-i ʿālamiyān ẓāhir. Qawluhu jallat
ʿaẓamatuhu va ʿazza sulṭānuhu:

As to what you wrote concerning the spiritual friend, Jinab-i Ayn and
Ba---may the glory of God, the All-Glorious, rest upon him---praise be
to God, he has received the fragrant breezes of these days. He is
mentioned in the divine Book and inscribed by the Most Exalted Pen in
the Crimson Scroll. No favor can equal this favor, and every station
bows before this station. I ask and hope that God, exalted be His glory
and all-embracing His bounty, will assist him in preserving this most
exalted station and lofty rank. May joy be his, may it bring him delight
and satisfaction, and blessed be he! These matters were presented in the
presence of the Lord of Revelation. This utterance appeared from the
kingdom of the loving care of the Beloved of the worlds. His
words---exalted be His grandeur and mighty His sovereignty:

\subsection{21}\label{section-2252}

\begin{otherlanguage}{arabic}

هو الشّاهد من افقه الأعلی

\end{otherlanguage}

Huwa al-shāhidu min ufuqihi al-aʿlā.

He is the Witness from His most exalted horizon.

\subsection{22}\label{section-2253}

\begin{otherlanguage}{persian}

انّا ذکرناه بذکر لا تغیّره حوادث الامکان و لا تبدّله القرون و الأعصار یکون
مخلّداً فی کتاب اللّه ربّ العالمین لعمری قد قدّر له ما لا تعادله الأرض و
خزائنها یشهد بذلک امّ الکتاب فی هذا القصر الّذی تنطق ذرّاتها تاللّه قد استوی
مکلّم الطّور علی عرش الظّهور و نطق اللّسان الملک و الملکوت ثمّ العزّة و
الجبروت ثمّ المقام و اللّاهوت لنفسی العزیز القدیم نسأل اللّه ان یفتح علی
وجهه ابواب العنایة و الألطاف و یقرّبه الیه فیکلّ الأحوال انّه هو المقتدر
علی ما یشآء و هو الفضّال الکریم قد ذکرناه فی لوح آخر لیفرح و یکون من
الشّاکرین کذلک اظهرت الرّحمة امطارها و القلم اسراره و شمس المعانی اشراقها
انّه هو المقتدر الغفور الرّحیم البهآء المشرق من افق سمآء فضلی علیه و علی
الّذین آمنوا باللّه العزیز الحمید انتهی بحر جود شاهد جود است و امواجش دلیل
بر او آفتاب خود مبیّن مقام خود است و اشراقاتش برهان او دیگر این خادم فانی
چه ذکر نماید و چه عرض کند که لایق و سزاوار باشد

\end{otherlanguage}

Innā dhakarnāhu bi-dhikrin lā tughayyiruhu ḥavādithu al-imkāni va lā
tubaddiluhu al-qurūnu va al-aʿṣāru yakūnu mukhalladan fī Kitābi Allāhi
Rabbi al-ʿālamīn. La-ʿamrī qad quddira lahu mā lā tuʿādiluhu al-arḍu va
khazāʾinuhā. Yashhadu bi-dhālika Ummu al-Kitābi fī hādhā al-qaṣri
alladhī tanṭiqu dharrātuhu: Ta-llāhi qad istavā Mukallimu al-Ṭūri ʿalā
ʿarshi al-ẓuhūri va naṭaqa al-lisānu: Al-mulku va al-malakūtu thumma
al-ʿizzatu va al-jabarūtu thumma al-maqāmu va al-lāhūtu li-nafsī
al-ʿazīzi al-qadīm. Nasʾalu Allāha an yaftaḥa ʿalā vajhihi abvāba
al-ʿināyati va al-alṭāfi va yuqarribahu ilayhi fī kulli al-aḥvāl. Innahu
huwa al-muqtadiru ʿalā mā yashāʾu va huwa al-faḍḍālu al-karīm. Qad
dhakarnāhu fī lawḥin ākhara li-yafraḥa va yakūna mina al-shākirīn.
Kadhālika aẓharati al-raḥmatu amṭārahā va al-qalamu asrārahu va shamsu
al-maʿānī ishrāqahā. Innahu huwa al-muqtadiru al-ghafūru al-raḥīm.
Al-bahāʾu al-mushriqu min ufuqi samāʾi faḍlī ʿalayhi va ʿalā alladhīna
āmanū bi-llāhi al-ʿazīzi al-ḥamīd. Intahā. Baḥr-i jūd shāhid-i jūd ast
va amvājash dalīl bar ū. Āftāb khud mubayyin-i maqām-i khud ast va
ishrāqātash burhān-i ū. Dīgar īn khādim-i fānī chi dhikr namāyad va chi
ʿarḍ kunad ki lāyiq va sazāvār bāshad?

``We have remembered him with a remembrance that the happenings of the
contingent world will not alter, nor the centuries and ages change; it
will abide forever in the Book of God, the Lord of the worlds. By My
life! There has been ordained for him that which the earth and its
treasures cannot equal. To this the Mother Book bears witness in this
palace, whose very atoms proclaim: `By God! He who spoke on Sinai has
ascended the throne of Revelation,' and the Tongue has declared:
`Sovereignty and the heavenly kingdom, might and dominion, station and
divinity belong to My own Self, the Almighty, the Ancient.' We ask God
to open before his face the doors of loving care and gracious favor, and
to draw him near to Himself in every circumstance. He indeed has power
over whatever He wills, and He is the Most Bountiful, the Generous. We
have remembered him in another Tablet, that he may rejoice and be among
the thankful. Thus has mercy revealed its showers, the Pen its
mysteries, and the Sun of inner meanings its radiance. He is indeed the
All-Powerful, the All-Forgiving, the Merciful. May the glory shining
from the horizon of the heaven of My grace rest upon him and upon those
who have believed in God, the Almighty, the Praiseworthy.'' End. The
ocean of bounty bears witness to bounty, and its waves are evidence of
it. The sun itself makes its own station clear, and its rays are its
proof. What more, then, can this mortal servant mention or submit that
would be fitting and worthy?

\subsection{23}\label{section-2254}

\begin{otherlanguage}{persian}

ذکر جناب خلیل علیه بهآء اللّه را نمودند للّه الحمد موفّق بوده و هستند
ذکرشان در ساحت امنع اقدس بوده اینفانی از حقّ جلّ جلاله میطلبد ایشانرا موفّق
فرماید بر خدمت امرش انّه هو السّامع المجیب یا محبوب فؤادی آثار قلم اعلی
اکثری در مقدّمهٴ طا بدست معاندین افتاد هر که را اخذ نمودند مع کتب و زبر و
الواح بوده باید آنمحبوب حزب اللّه را آگاه فرمایند تا کلّ در حفظ و اعزاز
الواح و آثار جهد بلیغ مبذول دارند و کمال سعی را منظور طوبی از برای
دیاریکه عرف صریر قلم اعلی در او متضوّعست حفظ الواح منیعهٴ ممتنعهٴ مبارکه
بسیار لازم و عنداللّه محبوبست که مباد صفحه و خطوط بطراز خود باقی نمانند
طریق حفظ را بفرمائید و تعلیم نمائید آنچه ذکر شد در اینمقام حسب الأمر
اقدس اعلی بوده طوبی للعاملین

\end{otherlanguage}

Dhikr-i Jināb-i Khalīl, ʿalayhi bahāʾu Allāh, rā namūdand. Li-llāhi
al-ḥamd muvaffaq būde va hastand. Dhikrishān dar sāḥat-i amnaʿ-i aqdas
būde. Īn fānī az Ḥaqq, jalla jalāluhu, mīṭalabad īshān rā muvaffaq
farmāyad bar khidmat-i Amrash. Innahu huwa al-sāmiʿu al-mujīb. Yā
maḥbūb-i fuʾādī, āthār-i Qalam-i Aʿlā aktharī dar muqaddame-yi Ṭā
bi-dast-i muʿānidīn uftād. Har ki rā akhdh namūdand maʿ-i kutub va zubur
va alvāḥ būde. Bāyad ān maḥbūb Ḥizb Allāh rā āgāh farmāyand tā kull dar
ḥifẓ va iʿzāz-i alvāḥ va āthār jahd-i balīgh mabdhūl dārand va kamāl-i
saʿy rā manẓūr. Ṭūbā az barā-yi diyārī ki ʿarf-i ṣarīr-i Qalam-i Aʿlā
dar ū mutaḍavviʿ ast. Ḥifẓ-i alvāḥ-i manīʿe-yi mumtaniʿe-yi mubārake
bisyār lāzim va ʿinda Allāh maḥbūb ast ki mabād ṣafḥe va khuṭūṭ
bi-ṭirāz-i khud bāqī namānand. Ṭarīq-i ḥifẓ rā bifarmāyīd va taʿlīm
namāyīd. Ānchi dhikr shud dar īn maqām ḥasb al-amr-i aqdas-i aʿlā būde.
Ṭūbā li-l-ʿāmilīn.

You mentioned Jinab-i Khalil---may God's glory rest upon him. Praise be
to God, he has been and remains favored with success. He has been
mentioned at the most holy, most exalted threshold. This mortal one asks
God, exalted be His glory, to enable him to serve His Cause. He is
indeed the Hearer, the Answerer. O beloved of my heart! Most of the
writings of the Most Exalted Pen fell into the hands of the adversaries
during the affair of Ta. Everyone who was arrested had books,
scriptures, and Tablets with them. That beloved friend should alert the
people of God, so that all may make strenuous efforts to preserve and
honor the Tablets and writings, exerting themselves to the utmost.
Blessed is the land in which the fragrance of the sound of the Most
Exalted Pen is diffused. Preserving the exalted, unassailable, blessed
Tablets is most necessary and is beloved in God's sight, lest the pages
and their written lines fail to retain their original form. Explain and
teach the means of preserving them. What has been mentioned here is in
accordance with the most holy, most exalted command. Blessed are those
who act.

\subsection{24}\label{section-2255}

\begin{otherlanguage}{persian}

ذکر جناب میرزا محمّد حسین خ ا و جناب عبّاسقلی خ ا علیهما بهآء اللّه فرموده
بودند بعد از عرض امام حضور دو لوح امنع اقدس مخصوص ایشان نازل و ارسال شد
از حقّ جلّ جلاله میطلبم بسماء مرتفعهٴ در لوح فائز شوند و از انوار شموس و
اقمار آن منوّر گردند بشأنیکه حوادث ایّام آنرا تغییر ندهد و ظلم فراعنهٴ عصر
و جبابرهٴ بلاد آنرا محجوب ننماید لیس هذا علی اللّه ربّنا و ربّکم بعزیز

\end{otherlanguage}

Dhikr-i Jināb-i Mīrzā Muḥammad Ḥusayn Khā Alif va Jināb-i ʿAbbāsqulī Khā
Alif, ʿalayhimā bahāʾu Allāh, farmūde būdand. Baʿd az ʿarḍ imām-i ḥuḍūr
du lawḥ-i amnaʿ-i aqdas makhṣūṣ-i īshān nāzil va irsāl shud. Az Ḥaqq,
jalla jalāluhu, mīṭalabam bi-samā-yi murtafiʿe-yi dar lawḥ fāʾiz shavand
va az anvār-i shumūs va aqmār-i ān munavvar gardand, bi-shaʾnī ki
ḥavādith-i ayyām ān rā taghyīr nadahad va ẓulm-i farāʿine-yi ʿaṣr va
jabābire-yi bilād ān rā maḥjūb nanamāyad. Laysa hādhā ʿalā Allāhi
Rabbinā va Rabbikum bi-ʿazīz.

You had mentioned Jinab-i Mirza Muhammad Husayn Kha Alif and Jinab-i
Abbasquli Kha Alif---may God's glory rest upon them both. After the
matter was presented in His presence, two most holy, most exalted
Tablets were revealed especially for them and dispatched. I ask God,
exalted be His glory, that they may attain the lofty heaven within the
Tablet and be illumined by the lights of its suns and moons, in such a
way that the happenings of the days may not alter it, nor the oppression
of the pharaohs of the age and the tyrants of the lands obscure it. This
is not difficult for God, our Lord and your Lord.

\subsection{25}\label{section-2256}

\begin{otherlanguage}{arabic}

قول الرّبّ تعالی و تقدّس هو النّور المشرق من افق العالم

\end{otherlanguage}

Qawlu al-Rabbi taʿālā va taqaddasa: Huwa al-nūru al-mushriqu min ufuqi
al-ʿālam.

The words of the Lord, exalted and sanctified is He: ``He is the Light
shining from the horizon of the world.''

\subsection{26}\label{section-2257}

\begin{otherlanguage}{persian}

هذا یوم فیه ینادی المقصود و یبشّر النّاس بامره المحتوم و رحیقه المختوم
طوبی للّذین اقبلوا بوجوه بیضآء و قلوب نورآء و اخذوا کأس العطآء من ید مالک
الأسمآء و فاطر السّمآء و شربوا باسمه المهیمن القیّوم قد ظهرت النّار من
السّدرة و نطق مکلّم الطّور من افق سمآء الظّهور یا ملأ الأرض تاللّه قد اتی
الوعد و الموعود من مقامه المحمود یدعوکم الی اللّه مالک الغیب و الشّهود قد
تجلّی الرّحمن علی من فی الامکان و ظهر المنظر الأکبر و مالک القدر یقول یا
معشر البشر اقبلوا الی الأفق الأبهی ثمّ استمعوا صریر قلمی الأعلی انّه
یهدیکم الی صراط اللّه ربّ ما کان و ما یکون

\end{otherlanguage}

Hādhā yawm fīhi yunādī al-Maqṣūd wa-yubashshiru al-nās bi-amrihi
al-maḥtūm wa-raḥīqihi al-makhtūm. Ṭūbā li-lladhīna aqbalū bi-wujūh
bayḍāʾ wa-qulūb nawrāʾ wa-akhadhū kaʾs al-ʿaṭāʾ min yad Mālik al-Asmāʾ
wa-Fāṭir al-Samāʾ wa-sharibū bi-ismihi al-Muhaymin al-Qayyūm. Qad
ẓaharat al-nār min al-sidra wa-naṭaqa Mukallim al-Ṭūr min ufuq samāʾ
al-ẓuhūr: yā malaʾ al-arḍ, ta-Allāh qad atā al-waʿd wa-l-Mawʿūd min
maqāmihi al-maḥmūd yadʿūkum ilā Allāh Mālik al-ghayb wa-l-shuhūd. Qad
tajallā al-Raḥmān ʿalā man fī al-imkān wa-ẓahara al-Manẓar al-Akbar
wa-Mālik al-qadar yaqūl: yā maʿshar al-bashar, aqbilū ilā al-ufuq
al-abhā thumma istamiʿū ṣarīr qalamī al-aʿlā, innahu yahdīkum ilā ṣirāṭ
Allāh Rabb mā kāna wa-mā yakūn.

This is a day on which the One sought calls out and brings humanity glad
tidings of His irrevocable decree and His sealed wine. Blessed are those
who have turned toward Him with radiant faces and luminous hearts, taken
the cup of bounty from the hand of the Lord of Names and Creator of
heaven, and drunk in His name, the All-Subduing, the Self-Subsisting.
The fire has appeared from the Lote-Tree, and He who spoke on Sinai has
spoken from the horizon of the heaven of revelation: ``People of the
earth! By God, the promise has come, and the Promised One, from His
praiseworthy station, calls you to God, the Lord of the unseen and the
seen.'' The All-Merciful has revealed Himself to all who dwell in the
realm of existence, and the Most Great Vision has appeared. The Lord of
destiny says: ``Humankind! Turn toward the all-glorious horizon, then
listen to the sound of My Most Exalted Pen. It guides you to the path of
God, the Lord of all that has been and all that will be.''

\subsection{27}\label{section-2258}

\begin{otherlanguage}{arabic}

یا ایّها المذکور لدی المظلوم انّا سمعنا ذکرک ذکرناک بما لا یعادله ما عند
الأمم و لا ما عند الأمرآء و الملوک طوبی لک بما فزت بآثار قلمی الأعلی و
نطق باسمک المظلوم اذ کان مستویاً علی عرش اسمه الودود انّا انزلنا الآیات و
اظهرنا البیّنات و خرقنا الحجبات امام الوجوه بحیث ما منعتنا سطوة العالم
یشهد بذلک من عنده کتاب مرقوم

\end{otherlanguage}

Yā ayyuhā al-madhkūr ladā al-Maẓlūm, innā samiʿnā dhikraka dhakarnāka
bi-mā lā yuʿādiluhu mā ʿinda al-umam wa-lā mā ʿinda al-umarāʾ
wa-l-mulūk. Ṭūbā laka bi-mā fuzta bi-āthār qalamī al-aʿlā wa-naṭaqa
bi-ismika al-Maẓlūm idh kāna mustawiyan ʿalā ʿarsh ismihi al-Wadūd. Innā
anzalnā al-āyāt wa-aẓharnā al-bayyināt wa-kharaqnā al-ḥujubāt amāma
al-wujūh bi-ḥaythu mā manaʿatnā saṭwat al-ʿālam. Yashhadu bi-dhālika man
ʿindahu kitāb marqūm.

You who have been mentioned before the Wronged One! We have heard you
mentioned and have remembered you with that which nothing possessed by
the nations, or by princes and kings, can equal. Blessed are you for
having received the writings of My Most Exalted Pen and for having had
your name spoken by the Wronged One as He sat upon the throne of His
name, the All-Loving. We have sent down the verses, made the clear
proofs manifest, and torn through the veils before all eyes, undeterred
by the might of the world. To this bears witness the one who possesses
an inscribed Book.

\subsection{28}\label{section-2259}

\begin{otherlanguage}{arabic}

قد احاطتنا البلایا من کلّ الأشطار و الرّزایا من الأقطار و قلم المختار کان
ینطق فی العشیّ و الاشراق و فی البکور و الآصال انّه لا اله الّا انا الحقّ
علّام الغیوب طوبی لمن سمع ندائی و اقبل الی افقی و نطق بثنائی و قام علی
خدمة امری و تمسّک بحبلی الممدود انّک اذا اجتذبتک نفحات آیاتی و اخذک کوثر
بیانی

\end{otherlanguage}

Qad aḥāṭatnā al-balāyā min kull al-ashṭār wa-l-razāyā min al-aqṭār
wa-qalam al-Mukhtār kāna yanṭiqu fī al-ʿashiyy wa-l-ishrāq wa-fī
al-bukūr wa-l-āṣāl: innahu lā ilāha illā anā al-Ḥaqq ʿAllām al-ghuyūb.
Ṭūbā li-man samiʿa nidāʾī wa-aqbala ilā ufuqī wa-naṭaqa bi-thanāʾī
wa-qāma ʿalā khidmat amrī wa-tamassaka bi-ḥablī al-mamdūd. Innaka idhā
ijtadhabatka nafaḥāt āyātī wa-akhadhaka kawthar bayānī

Afflictions have surrounded Us on every side, and calamities from every
quarter, while the Pen of the Chosen One has continued to proclaim, at
dusk and at sunrise, in the early morning and in the evening: ``There is
no God but Me, the Truth, the Knower of all hidden things.'' Blessed is
the one who has heard My call, turned toward My horizon, voiced My
praise, arisen to serve My Cause, and held fast to My outstretched cord.
When the breezes of My verses draw you and the river of My utterance
takes hold of you,

\subsection{29}\label{section-2260}

\begin{otherlanguage}{persian}

قل الهی الهی تری البعید اراد قربک و الفقیر بحر غنائک و العطشان کوثر
عطائک اسألک بانوار نیّر بیانک و اسرار کتابک و بافقک الأعلی و ما کان
مخزوناً فی خزائن قلمک و کنائز علمک یا مولی الأسمآء و فاطر السّمآء بأن
تؤیّدنی علی الاستقامة علی امرک و القیام علی خدمتک ای ربّ ترانی مقرّاً بما
نطق به لسان عظمتک فی ملکوت بیانک و معترفاً بما نزل من سمآء مشیّتک قوّ یا
الهی قلبی و جوارحی بحیث لا تضعفها قوّة الأقویآء و لا شبهات العلمآء ثمّ
اجعلنی مشتعلاً بنار سدرتک و منوّراً بأنوار عرشک اسألک یا مسخّر الآیات بأمرک
الّذی به نصبت رایات ظهورک فی الآفاق و اعلام نصرک فی البلاد بان تکتب لی من
قلم فضلک ما یکون معی فیکلّ عالم من عوالمک انّک انت المقتدر علی ما تشآء و
فی قبضتک زمام الأشیآء تفعل ما تشآء و تحکم ما ترید لا اله الّا انت الفرد
الواحد المهیمن العزیز الفضّال انتهی

\end{otherlanguage}

Qul: ilāhī ilāhī, tarā al-baʿīd arāda qurbaka wa-l-faqīr baḥr ghināʾika
wa-l-ʿaṭshān kawthar ʿaṭāʾika. Asʾaluka bi-anwār nayyir bayānika
wa-asrār kitābika wa-bi-ufuqika al-aʿlā wa-mā kāna makhzūnan fī khazāʾin
qalamika wa-kanāʾiz ʿilmika, yā Mawlā al-Asmāʾ wa-Fāṭir al-Samāʾ, bi-an
tuʾayyidanī ʿalā al-istiqāma ʿalā amrika wa-l-qiyām ʿalā khidmatika. Ay
Rabb, tarānī muqirran bi-mā naṭaqa bihi lisān ʿaẓamatika fī malakūt
bayānika wa-muʿtarifan bi-mā nazala min samāʾ mashīʾatika. Qawwi yā
ilāhī qalbī wa-jawāriḥī bi-ḥaythu lā tuḍʿifuhā quwwat al-aqwiyāʾ wa-lā
shubuhāt al-ʿulamāʾ. Thumma ijʿalnī mushtaʿilan bi-nār sidratika
wa-munaww aran bi-anwār ʿarshika. Asʾaluka yā Musakhkhir al-āyāt
bi-amrika alladhī bihi naṣabta rāyāt ẓuhūrika fī al-āfāq wa-aʿlām
naṣrika fī al-bilād bi-an taktuba lī min qalam faḍlika mā yakūnu maʿī fī
kull ʿālam min ʿawālimika. Innaka anta al-Muqtadir ʿalā mā tashāʾ wa-fī
qabḍatika zimām al-ashyāʾ. Tafʿalu mā tashāʾ wa-taḥkumu mā turīd. Lā
ilāha illā anta al-Fard al-Wāḥid al-Muhaymin al-ʿAzīz al-Faḍḍāl. Intahā.

Say: ``My God, my God! You see the distant one seeking nearness to You,
the poor one the ocean of Your riches, and the thirsty one the river of
Your bounty. I ask You, by the lights of the sun of Your utterance, the
mysteries of Your Book, Your most exalted horizon, and all that lies
stored in the treasuries of Your Pen and the repositories of Your
knowledge, O Lord of Names and Creator of heaven, to strengthen me to
remain steadfast in Your Cause and arise to serve You. My Lord! You see
me affirming what the Tongue of Your grandeur has spoken in the kingdom
of Your utterance, and acknowledging what has descended from the heaven
of Your will. Strengthen my heart and my limbs, my God, so that neither
the might of the powerful nor the doubts raised by the learned may
weaken them. Then set me ablaze with the fire of Your Lote-Tree and
illumine me with the lights of Your throne. I ask You, O Ruler of the
signs, by Your command through which You have raised the banners of Your
revelation throughout the horizons and the standards of Your victory
throughout the lands, to inscribe for me with the Pen of Your grace that
which will remain with me in every world of Your worlds. You are indeed
powerful to do whatever You will, and in Your grasp are the reins of all
things. You do what You will and decree what You desire. There is no God
but You, the Unique, the One, the All-Subduing, the Almighty, the
All-Bountiful.'' End.

\subsection{30}\label{section-2261}

\begin{otherlanguage}{arabic}

و هذا ما نزّل لجناب عبّاسقلیخان علیه بهآء اللّه قول الرّبّ تعالی و تقدّس

\end{otherlanguage}

Wa-hādhā mā nuzzila li-Janāb-i ʿAbbāsqulī Khān, ʿalayhi bahāʾ Allāh.
Qawl al-Rabb, taʿālā wa-taqaddasa.

And this is what was revealed for Jinab-i Abbasquli Khan, may the glory
of God rest upon him. The words of the Lord, exalted and sanctified is
He:

\subsection{31}\label{section-2262}

\begin{otherlanguage}{arabic}

بسمی المهیمن علی الأسمآء

\end{otherlanguage}

Bi-ismī al-Muhaymin ʿalā al-asmāʾ.

In My name, the All-Subduing over all names.

\subsection{32}\label{section-2263}

\begin{otherlanguage}{persian}

در این حین سدرهٴ بیان که در فردوس اعلی بید عنایت الهی غرس شده باینکلمه
ناطق یا اهل الأرض دعوا ما عندکم من الأوهام و الظّنون هذا یوم فیه فکّ ختم
الرّحیق المختوم باسم اللّه المهیمن القیّوم اقبلوا ثمّ اشربوا منه رغماً لکلّ
غافل مردود و کلّ عالم محجوب طوبی از برای نفسیکه ندایش را شنید و بحلاوت
بیانش فائز گشت قسم بانوار نیّر معانی که از افق سمآء بیان مشرقست هر نفسی
باصغا فائز شود مشتعل گردد اشتعالیکه اعراض عالم و اعتراض امم او را منع
ننماید و محروم نسازد و آثارش ظاهر و هویدا گردد از اوّل ایّام تا حین من‌غیر
ستر و حجاب امام وجوه امرا و علما و فقها ندا نمودیم ندائیکه آذان از برای
اصغاء ان از عدم بوجود آمده ولکن عباد را زخارف ظاهره و الوان فانیه از
نعمتهای باقیهٴ سرمدیّه منع نمود و محروم ساخت

\end{otherlanguage}

Dar īn ḥīn Sidrih-yi Bayān kih dar Firdaws-i Aʿlā bi-yad-i ʿināyat-i
ilāhī ghars shudih bi-īn kalimih nāṭiq: yā ahl al-arḍ, daʿū mā ʿindakum
min al-awhām wa-l-ẓunūn. Hādhā yawm fīhi fukka khatm al-raḥīq al-makhtūm
bi-ism Allāh al-Muhaymin al-Qayyūm. Aqbilū thumma ishrabū minhu raghman
li-kull ghāfil mardūd wa-kull ʿālim maḥjūb. Ṭūbā az barā-yi nafsī kih
nidāyash rā shinīd va bi-ḥalāvat-i bayānash fāʾiz gasht. Qasam
bi-anvār-i nayyir-i maʿānī kih az ufuq-i samā-yi bayān mushriq ast, har
nafsī bi-iṣghā fāʾiz shavad mushtaʿil gardad ishtiʿālī kih iʿrāḍ-i ʿālam
va iʿtirāḍ-i umam ū rā manʿ nanamāyad va maḥrūm nasāzad va āthārash
ẓāhir va huvaydā gardad. Az avval-i ayyām tā ḥīn min ghayr-i sitr va
ḥijāb imām-i vujūh-i umarā va ʿulamā va fuqahā nidā namūdīm, nidāyī kih
ādhān az barā-yi iṣghā-yi ān az ʿadam bi-vujūd āmadih, va-lākin ʿibād rā
zakhārif-i ẓāhirih va alvān-i fāniyih az niʿmat-hā-yi bāqiyih-yi
sarmadiyyih manʿ namūd va maḥrūm sākht.

At this moment the Lote-Tree of utterance, planted in the most exalted
Paradise by the hand of divine providence, speaks these words: ``People
of the earth! Abandon your vain imaginings and conjectures. This is the
day on which the seal of the sealed wine has been broken in the name of
God, the All-Subduing, the Self-Subsisting. Come forward and drink from
it, despite every heedless outcast and every veiled scholar.'' Blessed
is the soul who has heard its call and tasted the sweetness of its
utterance. By the lights of the sun of meanings shining from the horizon
of the heaven of utterance! Every soul who attains to hearing will be
set ablaze with such a fire that neither the world's rejection nor the
nations' objections will hinder or deprive that soul, and its effects
will become plain and manifest. From the earliest days until now,
without concealment or veil, We have called out before the faces of
rulers, scholars, and jurists---a call to hear which ears have been
brought from nonexistence into existence. Yet outward ornaments and
fleeting colors have barred the servants from everlasting, eternal
blessings and deprived them of these.

\subsection{33}\label{section-2264}

\begin{otherlanguage}{arabic}

انّک اذا سمعت نغمات طیور المعانی علی اغصان البیان طهّر قلبک و صدرک بمآء
التّوجّه و الانقطاع و قل

\end{otherlanguage}

Innaka idhā samiʿta naghamāt ṭuyūr al-maʿānī ʿalā aghṣān al-bayān ṭahhir
qalbaka wa-ṣadraka bi-māʾ al-tawajjuh wa-l-inqiṭāʿ wa-qul:

When you hear the melodies of the birds of meaning upon the branches of
utterance, cleanse your heart and breast with the water of turning
toward God and detachment, and say:

\subsection{34}\label{section-2265}

\begin{otherlanguage}{arabic}

الهی الهی لک الحمد بما اظهرت صراطک بین عبادک و دعوت النّاس الی مشرق وحیک
و مطلع الهامک و مصدر اوامرک و احکامک اشهد انّک اظهرت السّبیل و انزلت
الدّلیل و امرت الکلّ بما یقرّبهم الیک و ینفعهم فی کلّ عالم من عوالمک اسألک
یا مالک الوجود بأمواج بحر جودک و انوار شمس فضلک بأن تؤیّدنی فیکلّ الأحوال
لأکون ناطقاً بذکرک و مشتعلاً بنار حبّک و متذکّراً بآیاتک و طائراً فی هوآئک و
متمسّکاً بحبل عطائک و متشبّثاً بذیل کرمک ای ربّ تری عبدک مقبلاً الی افقک
الأعلی و معترفاً بوحدانیّتک و فردانیّتک و مقرّاً بعظمتک و سلطانک اسألک بالّذی
سرع الی مقرّ الفدآء شوقاً للقائک و اقبل الی سهام البلآء حبّاً لجمالک بأن
ترزقنی نعمتک الّتی انزلتها من سمآء امرک و المآئدة الّتی ارسلتها من ملکوت
بیانک انّک انت الّذی لا یمنعک صفوف العالم و لا جنوده تفعل ما تشآء و تحکم
ما ترید لا اله الّا انت العزیز الحمید ثمّ اسألک یا مالک القدم بالاسم
الأعظم الّذی به نوّرت العالم بأن تقدّر لی خیر الآخرة و الأولی انّک انت مالک
الأسمآء و فاطر السّمآء لا اله الّا انت المقتدر العزیز الوهّاب انتهی

\end{otherlanguage}

Ilāhī ilāhī, laka al-ḥamd bi-mā aẓharta ṣirāṭaka bayna ʿibādika
wa-daʿawta al-nās ilā mashriq waḥyika wa-maṭlaʿ ilhāmika wa-maṣdar
awāmirika wa-aḥkāmika. Ashhadu annaka aẓharta al-sabīl wa-anzalta
al-dalīl wa-amarta al-kull bi-mā yuqarribuhum ilayka wa-yanfaʿuhum fī
kull ʿālam min ʿawālimika. Asʾaluka yā Mālik al-wujūd bi-amwāj baḥr
jūdika wa-anwār shams faḍlika bi-an tuʾayyidanī fī kull al-aḥwāl
li-akūna nāṭiqan bi-dhikrika wa-mushtaʿilan bi-nār ḥubbika
wa-mutadhakkiran bi-āyātika wa-ṭāʾiran fī hawāʾika wa-mutamassikan
bi-ḥabl ʿaṭāʾika wa-mutashabbithan bi-dhayl karamika. Ay Rabb, tarā
ʿabdaka muqbilan ilā ufuqika al-aʿlā wa-muʿtarifan bi-waḥdāniyyatika
wa-fardāniyyatika wa-muqirran bi-ʿaẓamatika wa-sulṭānika. Asʾaluka
bi-lladhī saraʿa ilā maqarr al-fidāʾ shawqan li-liqāʾika wa-aqbala ilā
sihām al-balāʾ ḥubban li-jamālika bi-an tarzuqanī niʿmataka allatī
anzaltahā min samāʾ amrika wa-l-māʾida allatī arsaltahā min malakūt
bayānika. Innaka anta alladhī lā yamnaʿuka ṣufūf al-ʿālam wa-lā
junūduhu. Tafʿalu mā tashāʾ wa-taḥkumu mā turīd. Lā ilāha illā anta
al-ʿAzīz al-Ḥamīd. Thumma asʾaluka yā Mālik al-qidam bi-l-ism al-aʿẓam
alladhī bihi nawwarta al-ʿālam bi-an tuqaddira lī khayr al-ākhira
wa-l-ūlā. Innaka anta Mālik al-Asmāʾ wa-Fāṭir al-Samāʾ. Lā ilāha illā
anta al-Muqtadir al-ʿAzīz al-Wahhāb. Intahā.

``My God, my God! Praise be to You for having made Your path manifest
among Your servants and called humanity to the dawning-place of Your
revelation, the dayspring of Your inspiration, and the source of Your
commandments and laws. I testify that You have shown the way, sent down
the proof, and commanded everyone to do what will draw them near to You
and benefit them in every world of Your worlds. I ask You, O Lord of
existence, by the waves of the ocean of Your generosity and the lights
of the sun of Your grace, to sustain me in all circumstances, that I may
speak Your remembrance, burn with the fire of Your love, be mindful of
Your verses, soar in Your heaven, hold fast to the cord of Your bounty,
and cling to the hem of Your generosity. My Lord! You see Your servant
turning toward Your most exalted horizon, acknowledging Your oneness and
uniqueness, and affirming Your grandeur and sovereignty. I ask You, by
Him who hastened to the place of sacrifice out of longing to meet You
and welcomed the arrows of affliction out of love for Your beauty, to
grant me the blessing You have sent down from the heaven of Your Cause
and the table You have sent from the kingdom of Your utterance. You are
the One whom neither the ranks of the world nor its armies can hinder.
You do what You will and decree what You desire. There is no God but
You, the Almighty, the Praiseworthy. I further ask You, O Lord of
eternity, by the Most Great Name through which You have illumined the
world, to ordain for me the good of the next life and this life. You are
indeed the Lord of Names and Creator of heaven. There is no God but You,
the All-Powerful, the Almighty, the Ever-Bestowing.'' End.

\subsection{35}\label{section-2266}

\begin{otherlanguage}{persian}

العنایة للّه الفضل له و العطآء له هر نفسیکه در ساحت اقدس حاضر و یا اسمش
مذکور نازل شده از برای او آنچه که در هر کلمهٴ آن بحور فضل و عنایت مستور
طوبی از برای بصریکه بآن فائز شد و قلبیکه از آن نوشید اینعبد خدمت هر یک
از جنابان تکبیر و سلام میرسانم و از حقّ باقی میطلبم ایشانرا تأیید فرماید
بر خدمت امرش تا باسمش مدائن افئده را تصرّف نمایند ونفسه الحقّ سیف بیان احدّ
است و نفوذ کلمه مقدّس از احصا در هر حال حسب الأمر باید حکمت ملاحظه شود
این حکم مقدّمست بر اکثری از احکام چه که در الواح شتّی از قلم اعلی نازل و
جاری طوبی للعاملین

\end{otherlanguage}

Al-ʿināya li-llāh, al-faḍl lahu wa-l-ʿaṭāʾ lahu. Har nafsī kih dar
sāḥat-i aqdas ḥāḍir va yā ismash madhkūr, nāzil shudih az barā-yi ū
ānchih kih dar har kalimih-yi ān buḥūr-i faḍl va ʿināyat mastūr. Ṭūbā az
barā-yi baṣarī kih bi-ān fāʾiz shud va qalbī kih az ān nūshīd. Īn ʿabd
khidmat-i har yik az janābān takbīr va salām mīrasānam va az Ḥaqq-i bāqī
mīṭalabam īshān rā taʾyīd farmāyad bar khidmat-i amrash tā bi-ismash
madāʾin-i afʾidih rā taṣarruf namāyand. Wa-nafsihi al-Ḥaqq, sayf-i bayān
aḥadd ast va nufūdh-i kalimih muqaddas az iḥṣā. Dar har ḥāl ḥasab al-amr
bāyad ḥikmat mulāḥiẓih shavad. Īn ḥukm muqaddam ast bar aktharī az
aḥkām, chih kih dar alvāḥ-i shattā az qalam-i aʿlā nāzil va jārī. Ṭūbā
li-l-ʿāmilīn.

Gracious care belongs to God; grace is His, and bounty is His. For every
soul who is present in the most holy court or whose name is mentioned
there, something has been revealed in each word of which oceans of grace
and loving care lie concealed. Blessed is the eye that has attained it
and the heart that has drunk from it. This servant conveys greetings and
salutations to each of these honored friends and asks the everlasting
God to strengthen them in service to His Cause, so that through His name
they may conquer the cities of hearts. By His own Self, the Truth! The
sword of utterance is sharper, and the influence of the Word is beyond
reckoning. In every circumstance, as commanded, wisdom must be observed.
This command takes precedence over most ordinances, for it has been
revealed and issued by the Most Exalted Pen in numerous Tablets. Blessed
are those who act accordingly.

\subsection{36}\label{section-2267}

\begin{otherlanguage}{persian}

اینکه مرقوم داشتند جناب آقا خلیل ح ما علیه بهآء اللّه و عنایته در اوّل با
نفسین مذکورین در القای کلمه اقدام نموده بعد از عرض اینفقره در ساحت امنع
اقدس یک لوح مخصوص ایشان از سماء فضل نازل و ارسال شد انشآءاللّه بآن فائز
شوند و بحور بحور از بیان مقصود عالمیان قسمت برند

\end{otherlanguage}

Īnkih marqūm dāshtand Janāb-i Āqā Khalīl Ḥā Mā, ʿalayhi bahāʾ Allāh
wa-ʿināyatuhu, dar avval bā nafsayn-i madhkūrayn dar ilqā-yi kalimih
iqdām namūdih, baʿd az ʿarḍ-i īn fiqrih dar sāḥat-i amnaʿ-i aqdas yik
lawḥ makhṣūṣ-i īshān az samā-yi faḍl nāzil va irsāl shud. In shāʾ Allāh
bi-ān fāʾiz shavand va buḥūr buḥūr az bayān-i Maqṣūd-i ʿālamiyān qismat
barand.

You wrote that Jinab-i Aqa Khalil Ha Ma---may the glory of God and His
loving care rest upon him---had initially undertaken to convey the Word
together with the two aforementioned souls. After this matter was
presented in the most inaccessible, most holy court, a Tablet especially
for him was revealed from the heaven of grace and dispatched. God
willing, he will receive it and gain oceans upon oceans from the
utterance of the Desire of all worlds.

\subsection{37}\label{section-2268}

\begin{otherlanguage}{arabic}

قول الرّبّ تعالی و تقدّس هو المغرّد علی الأغصان

\end{otherlanguage}

Qawl al-Rabb, taʿālā wa-taqaddasa: huwa al-mugharrid ʿalā al-aghṣān.

The words of the Lord, exalted and sanctified is He: ``He is the One
singing upon the branches.''

\subsection{38}\label{section-2269}

\begin{otherlanguage}{arabic}

قد انزلنا ما کان مکنوناً فی ملکوت البیان و اظهرنا ما کان مخزوناً فی کتب
اللّه ربّ العالمین من النّاس من وجد عرف البیان و اقبل الی الرّحمن و منهم من
اعرض و انکر ما اتی من سمآء الفضل من لدی اللّه العلیم الحکیم نعیماً لمن سمع
النّدآء و انجذب بآیات اللّه و سحقاً لکلّ غافل بعید انّ الّذین وجدوا حلاوة
البیان انّهم من اهل البهآء فی کتاب اللّه العزیز الحمید

\end{otherlanguage}

Qad anzalnā mā kāna maknūnan fī malakūt al-bayān wa-aẓharnā mā kāna
makhzūnan fī kutub Allāh Rabb al-ʿālamīn. Min al-nās man wajada ʿarf
al-bayān wa-aqbala ilā al-Raḥmān wa-minhum man aʿraḍa wa-ankara mā atā
min samāʾ al-faḍl min ladā Allāh al-ʿAlīm al-Ḥakīm. Naʿīman li-man
samiʿa al-nidāʾ wa-injadhaba bi-āyāt Allāh wa-suḥqan li-kull ghāfil
baʿīd. Inna alladhīna wajadū ḥalāwat al-bayān innahum min ahl al-Bahāʾ
fī kitāb Allāh al-ʿAzīz al-Ḥamīd.

We have sent down what lay concealed in the kingdom of utterance and
made manifest what was stored in the Books of God, the Lord of the
worlds. Among the people are those who have perceived the fragrance of
utterance and turned toward the All-Merciful, and among them are those
who have turned away and denied what has come from the heaven of grace,
from God, the All-Knowing, the All-Wise. Bliss be upon the one who has
heard the call and been drawn by the verses of God, and perdition upon
every heedless one who remains far away. Those who have tasted the
sweetness of utterance are indeed numbered among the people of Baha in
the Book of God, the Almighty, the Praiseworthy.

\subsection{39}\label{section-2270}

\begin{otherlanguage}{arabic}

یا ایّها الخلیل یذکرک المظلوم من شطر السّجن و یبشّرک بعنایات ربّک و یذکّرک
بآیاته انّه هو الفضّال الکریم قد انزلنا لک و للّذین حضرت اسمآئهم لدی الوجه
ما یغنیهم عن کتب العالم یشهد بذلک هذا اللّوح المبین الّذی نزّل فی هذا الحین
من قلم اللّه مالک یوم الدّین انّک اذا شربت رحیق الوحی من کأس البیان و اخذک
جذب آیات ربّک الرّحمن

\end{otherlanguage}

Yā ayyuhā al-Khalīl, yadhkuruka al-Maẓlūm min shaṭr al-sijn
wa-yubashshiruka bi-ʿināyāt Rabbika wa-yudhakkiruka bi-āyātihi. Innahu
huwa al-Faḍḍāl al-Karīm. Qad anzalnā laka wa-li-lladhīna ḥaḍarat
asmāʾuhum ladā al-wajh mā yughnīhim ʿan kutub al-ʿālam. Yashhadu
bi-dhālika hādhā al-lawḥ al-mubīn alladhī nuzzila fī hādhā al-ḥīn min
qalam Allāh Mālik yawm al-dīn. Innaka idhā sharibta raḥīq al-waḥy min
kaʾs al-bayān wa-akhadhaka jadhb āyāt Rabbika al-Raḥmān

Khalil! The Wronged One remembers you from the prison and brings you
glad tidings of your Lord's loving favors, reminding you through His
verses. He is indeed the All-Bountiful, the Generous. We have revealed
for you and for those whose names have been brought before Our face that
which makes them independent of the books of the world. To this bears
witness this clear Tablet, revealed at this moment by the Pen of God,
the Lord of the Day of Judgment. When you have drunk the wine of
revelation from the cup of utterance and the attraction of the verses of
your Lord, the All-Merciful, has taken hold of you,

\subsection{40}\label{section-2271}

\begin{otherlanguage}{arabic}

قل الهی الهی اشهد بوحدانیّتک و فردانیّتک و بما جری من قلمک الأبهی و نزل من
افقک الأعلی اسألک یا مقصود العالم و محبوب الأمم باسمک الاعظم الّذی جعلته
نوراً للمقرّبین من عبادک و المخلصین من خلقک و ناراً لمن اعرض عنک و عن آیاتک
و اعترض علیک فی ایّامک بأن تؤیّد عبدک هذا علی خرق حجبات عبادک و سبحات خلقک
و وفّقه فیکلّ الأحوال علی ما یتضوّع به عرف رضائک و اسألک یا مولی الوری و ربّ
العرش و الثّری بأن تجعلنی منقطعاً عن دونک و قائماً علی نصرة امرک و ناطقاً
بما یقرّب عبادک الی ساحة عزّک وجمالک یا مطلع الجمال وجلالک یا مشرق الجلال
لا احبّ الوجود الّا لذکرک و ثنآئک و لا ارید الّا ما قدّرته لی من قلم تقدیرک
ای ربّ ایّدنی علی ما ینبغی لأیّامک و ظهورک و ما انزلته فی کتابک

\end{otherlanguage}

Qul: ilāhī ilāhī, ashhadu bi-waḥdāniyyatika wa-fardāniyyatika wa-bi-mā
jarā min qalamika al-abhā wa-nazala min ufuqika al-aʿlā. Asʾaluka yā
Maqṣūd al-ʿālam wa-Maḥbūb al-umam bi-ismika al-aʿẓam alladhī jaʿaltahu
nūran li-l-muqarrabīn min ʿibādika wa-l-mukhliṣīn min khalqika wa-nāran
li-man aʿraḍa ʿanka wa-ʿan āyātika wa-iʿtaraḍa ʿalayka fī ayyāmika bi-an
tuʾayyida ʿabdaka hādhā ʿalā kharq ḥujubāt ʿibādika wa-subuḥāt khalqika
wa-waffiqhu fī kull al-aḥwāl ʿalā mā yataḍawwaʿu bihi ʿarf riḍāka.
Wa-asʾaluka yā Mawlā al-warā wa-Rabb al-ʿarsh wa-l-tharā bi-an tajʿalanī
munqaṭiʿan ʿan dūnika wa-qāʾiman ʿalā nuṣrat amrika wa-nāṭiqan bi-mā
yuqarribu ʿibādaka ilā sāḥat ʿizzika wa-jamālika. Wa-jamālika yā Maṭlaʿ
al-jamāl wa-jalālika yā Mashriq al-jalāl, lā uḥibbu al-wujūd illā
li-dhikrika wa-thanāʾika wa-lā urīdu illā mā qaddartahu lī min qalam
taqdīrika. Ay Rabb, ayyidnī ʿalā mā yanbaghī li-ayyāmika wa-ẓuhūrika
wa-mā anzaltahu fī kitābika.

Say: ``My God, my God! I testify to Your oneness and uniqueness, and to
what has flowed from Your all-glorious Pen and descended from Your most
exalted horizon. I ask You, O Desire of the world and Beloved of the
nations, by Your Most Great Name, which You have made a light for those
of Your servants who are near to You and those of Your creatures who are
sincere, and a fire for those who have turned away from You and Your
verses and opposed You in Your days, to strengthen this servant of Yours
to tear through the veils of Your servants and the veils of glory of
Your creatures. Enable him in all circumstances to do that through which
the fragrance of Your good pleasure may spread. And I ask You, O Lord of
all creation and Lord of the throne and the earth, to make me detached
from all but You, steadfast in aiding Your Cause, and eloquent in that
which draws Your servants near to the court of Your might and beauty. By
Your beauty, O Dayspring of beauty, and Your majesty, O Dawning-Place of
majesty! I cherish existence only to remember and praise You, and I
desire only what You have ordained for me through the Pen of Your
decree. My Lord! Strengthen me to do what befits Your days and Your
revelation, and what You have sent down in Your Book.''

\subsection{41}\label{section-2272}

\begin{otherlanguage}{arabic}

اسألک بآثار قلمک و نفوذ مشیّتک و اقتدار ارادتک بأن تؤیّد عبادک علی الرّجوع
الیک و الانابة لدی باب فضلک و الخضوع عند ظهورات عظمتک ای ربّ لا تمنعهم عن
امواج بحر بیانک و لا عن مشاهدة آثار قدرتک فی ایّامک انت الّذی لا تمنعک
شئونات العالم و لا اعراض الأمم و انت الّذی یا الهی اظهرت الکلیم من بیت
فرعون رغماً له و لأعدائک و نجّیت الخلیل من نار ذات الوقود اظهاراً لقدرتک و
سلطنتک اسألک بآیات امرک و رایات عزمک بأن تقدّر لأولیائک خیر الآخرة و
الأولی و ما یقرّبهم الیک یا مالک الأشیاء لا اله الّا انت المقتدر القدیر
انتهی

\end{otherlanguage}

Asʾaluka bi-āthāri qalamika wa-nufūdhi mashīʾatika wa-iqtidāri irādatika
bi-an tuʾayyida ʿibādaka ʿalā al-rujūʿi ilayka wa-l-inābati ladā bābi
faḍlika wa-l-khuḍūʿi ʿinda ẓuhūrāti ʿaẓamatika. Ay rabbi lā tamnaʿhum
ʿan amwāji baḥri bayānika wa-lā ʿan mushāhadati āthāri qudratika fī
ayyāmika. Anta alladhī lā tamnaʿuka shuʾūnu al-ʿālami wa-lā iʿrāḍu
al-umami wa-anta alladhī yā ilāhī aẓharta al-Kalīma min bayti Firʿawna
raghman lahu wa-li-aʿdāʾika wa-najjayta al-Khalīla min nārin dhāti
al-waqūdi iẓhāran li-qudratika wa-salṭanatika. Asʾaluka bi-āyāti amrika
wa-rāyāti ʿazmika bi-an tuqaddira li-awliyāʾika khayra al-ākhirati
wa-l-ūlā wa-mā yuqarribuhum ilayka yā Mālika al-ashyāʾi. Lā ilāha illā
anta al-Muqtadiru al-Qadīru. Intahā.

I ask You, by the traces of Your Pen, the pervasive power of Your will,
and the might of Your purpose, to strengthen Your servants to return to
You, to turn in repentance at the door of Your grace, and to bow before
the manifestations of Your grandeur. My Lord! Do not withhold them from
the waves of the ocean of Your utterance or from beholding the signs of
Your power in Your days. You are the One whom neither the affairs of the
world nor the nations' rejection can hinder. You are the One, my God,
who brought forth Moses from Pharaoh's house despite him and Your
enemies, and delivered Abraham from a blazing fire, manifesting Your
power and sovereignty. I ask You, by the verses of Your Cause and the
banners of Your resolve, to ordain for Your loved ones the good of the
next world and this world, and whatever will draw them near to You, O
Possessor of all things. There is no God but You, the All-Powerful, the
Omnipotent.'' End.

\subsection{42}\label{section-2273}

\begin{otherlanguage}{arabic}

از حقّ جلّ جلاله سائل و آمل که جناب مذکور علیه بهآء اللّه را مؤیّد فرماید و
موفّق دارد بر خدمت امر به روح و ریحان و حکمت و بیان انّ ربّنا الرّحمن هو
المقتدر علی ما یشآء و هو المشفق الکریم

\end{otherlanguage}

Az Ḥaqq jalla jalāluhu sāʾil va āmil kih janāb-i madhkūr ʿalayhi
Bahāʾuʾllāh rā muʾayyad farmāyad va muwaffaq dārad bar khidmat-i Amr bih
rūḥ va rayḥān va ḥikmat va bayān. Inna Rabbanā al-Raḥmāna huwa
al-Muqtadiru ʿalā mā yashāʾu wa-huwa al-Mushfiqu al-Karīmu.

I ask and hope that God, glorified be His glory, will strengthen the
gentleman mentioned---may the glory of God rest upon him---and enable
him to serve the Cause with joy and delight, wisdom and utterance. Our
Lord, the All-Merciful, is indeed powerful over whatever He wills, and
He is the Compassionate, the Generous.

\subsection{43}\label{section-2274}

\begin{otherlanguage}{persian}

ذکر جناب آقا میرزا عبداللّه خ ا علیه بهآء اللّه که بهدایت حبیب روحانی جناب
هدی علیه بهآء اللّه و عنایته بافق اعلی توجّه نموده‌اند در پیشگاه حضور مالک
ظهور بشرف اصغا فائز و یک لوح امنع اقدس مخصوص نازل و ارسال شد

\end{otherlanguage}

Dhikr-i janāb-i Āqā Mīrzā ʿAbduʾllāh Kh Ā ʿalayhi Bahāʾuʾllāh kih bih
hidāyat-i ḥabīb-i rūḥānī janāb-i Hudā ʿalayhi Bahāʾuʾllāh wa-ʿināyatuhu
bih ufuq-i aʿlā tavajjuh namūdih-and dar pīshgāh-i ḥuḍūr-i Mālik-i Ẓuhūr
bih sharaf-i iṣghā fāʾiz va yik lawḥ-i amnaʿ-i aqdas-i makhṣūṣ nāzil va
irsāl shud.

The mention of Aqa Mirza ʿAbdu'llah Kh. A.---may the glory of God rest
upon him---who has turned toward the Most Exalted Horizon through the
guidance of the spiritual friend, his honor Huda---may the glory of God
and His loving care rest upon him---was honored with a hearing in the
presence of the Lord of Revelation. A most exalted, most holy Tablet was
revealed especially for him and dispatched.

\subsection{44}\label{section-2275}

\begin{otherlanguage}{arabic}

قول الرّبّ تعالی و تقدّس هو النّاطق فی ملکوت البیان

\end{otherlanguage}

Qawlu al-Rabbi taʿālā wa-taqaddasa: Huwa al-Nāṭiqu fī Malakūti
al-Bayāni.

The words of the Lord, exalted and sanctified be He: ``He is the One who
speaks in the Kingdom of Utterance.''

\subsection{45}\label{section-2276}

\begin{otherlanguage}{arabic}

کتاب انزله المظلوم لمن اراد ان یشرب الرّحیق المختوم من ایادی عطآء ربّه
المهیمن القیّوم لیقرّبه النّدآء الی الأفق الأعلی و یؤیّده علی الذّکر و الثّنآء
طوبی لمن سمع و فاز و ویل لکلّ غافل محجوب

\end{otherlanguage}

Kitābun anzalahu al-Maẓlūmu li-man arāda an yashraba al-raḥīqa
al-makhtūma min ayādī ʿaṭāʾi Rabbihi al-Muhaymini al-Qayyūmi
li-yuqarribahu al-nidāʾu ilā al-Ufuqi al-Aʿlā wa-yuʾayyidahu ʿalā
al-dhikri wa-l-thanāʾi. Ṭūbā li-man samiʿa wa-fāza wa-waylun li-kulli
ghāfilin maḥjūbin.

A Book sent down by the Wronged One for whoever desires to drink the
sealed wine from the hands of the bounty of his Lord, the
All-Protecting, the Self-Subsisting, so that the call may draw him near
to the Most Exalted Horizon and strengthen him in remembrance and
praise. Blessed is the one who has heard and attained, and woe to every
heedless one shut out by veils.

\subsection{46}\label{section-2277}

\begin{otherlanguage}{arabic}

انّا سمعنا ذکرک ذکرناک من قبل و فی هذا الحین فضلاً من لدی اللّه مالک الغیب
و الشّهود قد حضرت العبد الحاضر امام الوجه بکتاب من احبّی و تمسّک بحبلی و
طار فی هواء حبّی و قام علی نصرة امری بین عبادی و کان فیه اسمک انزلنا من
سمآء الفضل آیات لا یعادلها شیء من الأشیاء لتشکر ربّک مالک الوجود

\end{otherlanguage}

Innā samiʿnā dhikraka dhakarnāka min qablu wa-fī hādhā al-ḥīni faḍlan
min ladā Allāhi Māliki al-ghaybi wa-l-shuhūdi. Qad ḥaḍḍarta al-ʿabda
al-ḥāḍira amāma al-wajhi bi-kitābi man aḥabbanī wa-tamassaka bi-ḥablī
wa-ṭāra fī hawāʾi ḥubbī wa-qāma ʿalā nuṣrati amrī bayna ʿibādī wa-kāna
fīhi ismuka. Anzalnā min samāʾi al-faḍli āyātin lā yuʿādiluhā shayʾun
mina al-ashyāʾi li-tashkura Rabbaka Mālika al-wujūdi.

We have heard mention of you. We remembered you before, and do so at
this time, as a bounty from God, the Lord of the unseen and the seen.
You have brought the servant who attends before Our face with a letter
from one who has loved Me, held fast to My cord, soared in the heaven of
My love, and arisen to aid My Cause among My servants; your name was in
it. We have sent down from the heaven of grace verses that nothing
whatsoever can equal, so that you may thank your Lord, the Possessor of
existence.

\subsection{47}\label{section-2278}

\begin{otherlanguage}{arabic}

ایّاک ان یمنعک حوادث العالم عن مالک القدم او یحجبک شبهات الأمم عن هذا
الأمر المختوم قل یا معشر البشر قد انار افق الظّهور بمالک القدر و ظهر ما
کان مکنوناً فی کتب اللّه العزیز الودود ضعوا ما عندکم خذوا ما امرتم به من
لدی الحقّ علّام الغیوب

\end{otherlanguage}

Iyyāka an yamnaʿaka ḥawādithu al-ʿālami ʿan Māliki al-Qidami aw
yaḥjubaka shubuhātu al-umami ʿan hādhā al-Amri al-makhtūmi. Qul yā
maʿshara al-bashari qad anāra ufuqu al-ẓuhūri bi-Māliki al-Qadari
wa-ẓahara mā kāna maknūnan fī kutubi Allāhi al-ʿAzīzi al-Wadūdi. Ḍaʿū mā
ʿindakum khudhū mā umirtum bihi min ladā al-Ḥaqqi ʿAllāmi al-ghuyūbi.

Beware lest the events of the world keep you from the Lord of eternity,
or the doubts raised by the nations veil you from this sealed Cause.
Say: ``People of the world! The horizon of revelation has been illumined
by the Lord of destiny, and what lay hidden in the Books of God, the
Almighty, the All-Loving, has appeared. Set aside what you possess and
take hold of what you have been commanded by the True One, the Knower of
all hidden things.''

\subsection{48}\label{section-2279}

\begin{otherlanguage}{persian}

از لغة فصحی بلغة نوراء توجّه نمودیم امروز عالم بانوار ظهور منوّر و جمیع
اشیاء به ذکر و ثنا و فرح و سرور مشغول در کتب الهی از قبل و بعد بیاد این
یوم مبارک عیش اعظم برپا طوبی از برای نفسیکه فائز شد و بمقام یوم آگاه گشت
امروز صریر قلم الهی مرتفع و افقش مشهود طوبی از برای نفسیکه شبهات و
اشارات اهل عالم او را از مالک قدم منع ننمود و از بحر اعظم محروم نساخت

\end{otherlanguage}

Az lughat-i fuṣḥā bih lughat-i nawrā tavajjuh namūdīm. Imrūz ʿālam bih
anvār-i ẓuhūr munavvar va jamīʿ-i ashyā bih dhikr va thanā va faraḥ va
surūr mashghūl. Dar kutub-i ilāhī az qabl va baʿd bih yād-i īn yawm-i
mubārak ʿaysh-i aʿẓam bar pā. Ṭūbā az barā-yi nafsī kih fāʾiz shud va
bih maqām-i yawm āgāh gasht. Imrūz ṣarīr-i qalam-i ilāhī murtafiʿ va
ufuqash mashhūd. Ṭūbā az barā-yi nafsī kih shubuhāt va ishārāt-i ahl-i
ʿālam ū rā az Mālik-i Qidam manʿ nanamūd va az baḥr-i aʿẓam maḥrūm
nasākht.

We have turned from the most eloquent tongue to the luminous tongue.
Today the world is illumined by the lights of revelation, and all things
are occupied with remembrance and praise, joy and gladness. In the
divine Books, earlier and later, the greatest rejoicing has been raised
in commemoration of this blessed Day. Blessed is the soul who has
attained and become aware of the station of this Day. Today the sound of
the divine Pen is raised, and its horizon is manifest. Blessed is the
soul whom the doubts and allusions of the people of the world have not
kept from the Lord of eternity or deprived of the Most Great Ocean.

\subsection{49}\label{section-2280}

\begin{otherlanguage}{arabic}

طوبی لک و لمن بشّرک و هداک الی صراطی المستقیم و نبأی العظیم کذلک اسمعناک
ندائی الأحلی و ارسلنا الیک ما قرّت به اعین المقرّبین

\end{otherlanguage}

Ṭūbā laka wa-li-man bashsharaka wa-hadāka ilā ṣirāṭī al-mustaqīmi
wa-nabaʾī al-ʿaẓīmi. Kadhālika asmaʿnāka nidāʾī al-aḥlā wa-arsalnā
ilayka mā qarrat bihi aʿyunu al-muqarrabīna.

Blessed are you and the one who brought you glad tidings and guided you
to My straight path and My mighty Announcement. Thus have We enabled you
to hear My sweetest call and sent you that which has delighted the eyes
of those brought near.

\subsection{50}\label{section-2281}

\begin{otherlanguage}{arabic}

و نذکر فی هذا المقام ضلعک و نذکّرها بآیاتی و نبشّرها بعنایتی و رحمتی الّتی
سبقت من فی السّموات و الأرضین البهآء المشرق من افق سمآء عنایتی علیک و
علیها و علی الّذین ما خوّفتهم جنود الفراعنه و صفوف الجبابره عن اللّه ربّ
العالمین الحمد للّه الملک الحقّ العدل المبین

\end{otherlanguage}

Wa-nadhkuru fī hādhā al-maqāmi ḍilʿaka wa-nudhakkiruhā bi-āyātī
wa-nubashshiruhā bi-ʿināyatī wa-raḥmatī allatī sabaqat man fī
al-samāwāti wa-l-araḍīna. Al-Bahāʾu al-mushriqu min ufuqi samāʾi
ʿināyatī ʿalayka wa-ʿalayhā wa-ʿalā alladhīna mā khawwafathum junūdu
al-farāʿinati wa-ṣufūfu al-jabābirati ʿani Allāhi Rabbi al-ʿālamīna.
Al-ḥamdu li-llāhi al-Maliki al-Ḥaqqi al-ʿAdli al-Mubīni.

Here We also remember your wife, remind her through My verses, and give
her glad tidings of My loving care and My mercy, which has preceded all
who are in the heavens and on earth. May the glory shining from the
horizon of the heaven of My loving care rest upon you, upon her, and
upon those whom the armies of the Pharaohs and the ranks of tyrants have
not frightened away from God, the Lord of the worlds. Praise be to God,
the King, the True One, the Just, the Manifest.

\subsection{51}\label{section-2282}

\begin{otherlanguage}{persian}

انتهی ذکر جناب ایشان از قبل هم شده این خادم فانی از حقّ تعالی شأنه سائل و
آمل مدد فرماید ایشانرا تا بتمامه بر خدمت امرش قیام نمایند و در این ایّام
فانیه کسب مقامات باقیه فرمایند اوست مقتدر و توانا

\end{otherlanguage}

Intahā. Dhikr-i janāb-i īshān az qabl ham shudih. Īn khādim-i fānī az
Ḥaqq taʿālā shaʾnuhu sāʾil va āmil madad farmāyad īshān rā tā bih
tamāmih bar khidmat-i Amrash qiyām namāyand va dar īn ayyām-i fāniyih
kasb-i maqāmāt-i bāqiyih farmāyand. Ūst muqtadir va tavānā.

End. He has also been mentioned before. This transitory servant asks and
hopes that God, exalted be His station, will assist him to arise
wholeheartedly in service to His Cause and attain enduring stations in
these fleeting days. He is the Powerful, the Mighty.

\subsection{52}\label{section-2283}

\begin{otherlanguage}{persian}

و همچنین خدمت جناب هادی یعنی هدی علیه بهآء اللّه الأعلی تکبیر و سلام
میرسانم للّه الحمد بخدمت مشغولند و بعنایت فائز نسأل اللّه ربّنا و ربّکم ان
یودع فی بیانه ما یجذب النّاس الی افقه المبین و امره العظیم و همچنین ذکر
جناب حس خ ا و جناب میرزا محمّد خ ا علیهما بهآء اللّه را فرمودند عرض شد و
مخصوص هر یک لوح بدیع از افق سماء فضل اشراق نمود و ارسال گشت لتقرّ بهما
عیناهما و تشکرا ربّهما المشفق الکریم این بنده هم خدمت ایشان تکبیر میرساند
و از حقّ جلّ جلاله از برای ایشان تأیید میطلبد انّه علی کلّشیء قدیر

\end{otherlanguage}

Va hamchunīn khidmat-i janāb-i Hādī yaʿnī Hudā ʿalayhi Bahāʾuʾllāhi
al-Aʿlā takbīr va salām mīrasānam. Li-llāhi al-ḥamd bih khidmat
mashghūland va bih ʿināyat fāʾiz. Nasʾalu Allāha Rabbanā wa-Rabbakum an
yūdiʿa fī bayānihi mā yajdhibu al-nāsa ilā ufuqihi al-mubīni wa-amrihi
al-ʿaẓīmi. Va hamchunīn dhikr-i janāb-i Ḥas Kh Ā va janāb-i Mīrzā
Muḥammad Kh Ā ʿalayhimā Bahāʾuʾllāh rā farmūdand ʿarḍ shud va makhṣūṣ-i
har yik lawḥ-i badīʿ az ufuq-i samāʾ-i faḍl ishrāq namūd va irsāl gasht
li-taqarra bihimā ʿaynāhumā wa-tashkurā Rabbahumā al-Mushfiqa al-Karīma.
Īn bandih ham khidmat-i īshān takbīr mīrasānad va az Ḥaqq jalla jalāluhu
az barā-yi īshān taʾyīd mīṭalabad. Innahu ʿalā kulli shayʾin qadīrun.

I also convey greetings and salutations to his honor Hadi, that is,
Huda---may the glory of God, the Most Exalted, rest upon him. Praise be
to God, he is occupied with service and has received loving favor. We
ask God, our Lord and your Lord, to place in his utterance that which
will draw people toward His manifest horizon and His mighty Cause. Your
mention of his honor Has Kh. A. and his honor Mirza Muhammad Kh.
A.---may the glory of God rest upon them both---was likewise presented.
A wondrous Tablet shone forth from the horizon of the heaven of grace
especially for each of them and was dispatched, so that their eyes might
be gladdened by them and they might thank their Lord, the Compassionate,
the Generous. This servant also conveys greetings to them and asks God,
glorified be His glory, to strengthen them. He is powerful over all
things.

\subsection{53}\label{section-2284}

\begin{otherlanguage}{persian}

ذکر دوستان ارض یا علیهم بهآء اللّه را فرموده بودند در ساحت اقدس عرض شد
واحداً بعد واحد بذکر حقّ جلّ جلاله فائز و این آیات از افق اراده مشرق قوله
جلّ جلاله و عمّ نواله

\end{otherlanguage}

Dhikr-i dūstān-i arḍ-i Yā ʿalayhim Bahāʾuʾllāh rā farmūdih būdand dar
sāḥat-i aqdas ʿarḍ shud. Vāḥidan baʿda vāḥid bih dhikr-i Ḥaqq jalla
jalāluhu fāʾiz va īn āyāt az ufuq-i irādih mushriq. Qawluhu jalla
jalāluhu wa-ʿamma nawāluhu.

Your mention of the friends in the land of Ya---may the glory of God
rest upon them---was presented in the Most Holy Court. One by one, they
were favored with the remembrance of God, glorified be His glory, and
these verses shone from the horizon of His will. His words, glorified be
His glory and all-embracing His bounty:

\subsection{54}\label{section-2285}

\begin{otherlanguage}{arabic}

انا السّامع المجیب

\end{otherlanguage}

Anā al-Sāmiʿu al-Mujību.

``I am the One who hears and answers.''

\subsection{55}\label{section-2286}

\begin{otherlanguage}{arabic}

در این حین که مکلّم طور بر عرش ظهور مستوی پیک بیان از حدیقهٴ معانی
اینکلمه آورد یا اهل البهآء زیّنوا هیاکلکم بطراز الأعمال و الأخلاق لتستضیء
بها الآفاق هذه کلمة انزلناها فی لوح آخر و ذکّرناها لکم فی هذا الکتاب
المبین اسمعوا ندائی من شطر سجنی انّه لا اله الّا انا المقتدر القدیر

\end{otherlanguage}

Dar īn ḥīn kih Mukallim-i Ṭūr bar ʿarsh-i ẓuhūr mustavī payk-i bayān az
ḥadīqiy-i maʿānī īn kalimih āvard: Yā ahla al-Bahāʾi zayyinū hayākilakum
bi-ṭirāzi al-aʿmāli wa-l-akhlāqi li-tastaḍīʾa bihā al-āfāqu. Hādhihi
kalimatun anzalnāhā fī lawḥin ākhara wa-dhakkarnāhā lakum fī hādhā
al-kitābi al-mubīni. Ismaʿū nidāʾī min shaṭri sijnī innahu lā ilāha illā
anā al-Muqtadiru al-Qadīru.

At this moment, when the One who spoke on Sinai is established upon the
throne of revelation, the messenger of utterance has brought these words
from the garden of meanings: ``People of Baha! Adorn yourselves with the
ornament of deeds and character, so that the horizons may be illumined
by them. These are words We revealed in another Tablet and have recalled
for you in this clear Book. Hear My call from the direction of My
prison: there is no God but Me, the All-Powerful, the Omnipotent.''

\subsection{56}\label{section-2287}

\begin{otherlanguage}{arabic}

قد اظهرت نفسی و دعوت العباد الی بحر فضلی منهم من اقبل و منهم من اعرض
طوبی لمقبل ما منعه اعراض المعتدین انّا ذکرناکم من قبل بآیاتی و فی هذا
الحین بما انجذبت به افئدة العارفین انّا قدّرنا النّصرة بتقوی اللّه طوبی لمن
تمسّک به ویل للغافلین انصروا ربّکم الرّحمن بجنود الأخلاق هذا ما امرتم به فی
الواح شتّی و انا الآمر الحکیم خذوا اقداح الفلاح باسم فالق الأصباح ثمّ
اشربوا منها بذکری العزیز البدیع قولوا یا ملأ الامکان قد احاطتکم انوار
البیان اتّقوا الرّحمن و لا تکونوا من المعرضین اقبلوا الی اللّه سیفنی ما
ترونه الیوم و ترون انفسکم من الأخسرین فی محضر اللّه العزیز العظیم این
الأکاسرة و ابنائهم و الفراعنة و آثارهم انظروا ثمّ اذکروا اذ اتی محمّد رسول
اللّه اعترضتم علیه و اردتم اطفاء نوره ولکنّ اللّه حفظه بجنود الغیب و
الشّهاده و رفع امره یا اولیائی فی یا کیف انّه هو المقتدر علی ما یشآء و
الحاکم علی ما یرید لا اله الّا هو العزیز الحمید

\end{otherlanguage}

Qad aẓhartu nafsī wa-daʿawtu al-ʿibāda ilā baḥri faḍlī minhum man aqbala
wa-minhum man aʿraḍa. Ṭūbā li-muqbilin mā manaʿahu iʿrāḍu al-muʿtadīna.
Innā dhakarnākum min qablu bi-āyātī wa-fī hādhā al-ḥīni bi-mā injadhabat
bihi afʾidatu al-ʿārifīna. Innā qaddarnā al-nuṣrata bi-taqwā Allāhi ṭūbā
li-man tamassaka bihi waylun li-l-ghāfilīna. Unṣurū Rabbakumu al-Raḥmāna
bi-junūdi al-akhlāqi hādhā mā umirtum bihi fī alwāḥin shattā wa-anā
al-Āmiru al-Ḥakīmu. Khudhū aqdāḥa al-falāḥi bi-smi Fāliqi al-aṣbāḥi
thumma ishrabū minhā bi-dhikrī al-ʿazīzi al-badīʿi. Qūlū yā malaʾa
al-imkāni qad aḥāṭatkum anwāru al-bayāni ittaqū al-Raḥmāna wa-lā takūnū
mina al-muʿriḍīna. Aqbilū ilā Allāhi sa-yafnā mā tarawnahu al-yawma
wa-tarawna anfusakum mina al-akhsarīna fī maḥḍari Allāhi al-ʿAzīzi
al-ʿAẓīmi. Ayna al-akāsiratu wa-abnāʾuhum wa-l-farāʿinatu wa-āthāruhum.
Unẓurū thumma udhkurū idh atā Muḥammadun Rasūlu Allāhi iʿtaraḍtum
ʿalayhi wa-aradtum iṭfāʾa nūrihi wa-lākinna Allāha ḥafiẓahu bi-junūdi
al-ghaybi wa-l-shahādati wa-rafaʿa amrahu yā awliyāʾī fī Yā Kayf. Innahu
huwa al-Muqtadiru ʿalā mā yashāʾu wa-l-Ḥākimu ʿalā mā yurīdu lā ilāha
illā huwa al-ʿAzīzu al-Ḥamīdu.

``I have manifested Myself and summoned the servants to the ocean of My
grace. Some have turned toward Me and some have turned away. Blessed is
the one who has turned toward Me and whom the rejection of the
transgressors has not hindered. We remembered you before through My
verses, and at this time through that which has enraptured the hearts of
those who recognize Me. We have ordained that aid be rendered through
the fear of God. Blessed is the one who holds fast to it; woe to the
heedless. Aid your Lord, the All-Merciful, with the hosts of good
character. This is what you have been commanded in various Tablets, and
I am the Ordainer, the All-Wise. Take up the cups of salvation in the
name of the One who causes the dawn to break, then drink from them in My
mighty and wondrous remembrance. Say: `People of the world of being! The
lights of utterance have encompassed you. Fear the All-Merciful and do
not be among those who turn away. Turn toward God. What you see today
will perish, and you will see yourselves among those who have suffered
the greatest loss in the presence of God, the Almighty, the Most Great.
Where are the kings of Persia and their sons, and the Pharaohs and their
traces? Consider, then remember how, when Muhammad, the Messenger of
God, came, you opposed Him and sought to extinguish His light. But God
protected Him with the hosts of the unseen and the seen and exalted His
Cause, O My loved ones in Ya Kayf. He is indeed powerful over whatever
He wills and rules over whatever He desires. There is no God but Him,
the Almighty, the Praiseworthy.'\,''

\subsection{57}\label{section-2288}

\begin{otherlanguage}{arabic}

یا احبّائی یا اولیائی اشکروا اللّه بما اسمعکم ندآئه و اراکم آثار قلمه و
هداکم الی صراطه المستقیم طوبی لکم و نعیماً لکم بما فزتم بما لا فاز به
علمآء الأرض و فقهائها احفظوا هذا المقام الأعلی باسم ربّکم النّاطق فی هذا
المقام الرّفیع کذلک نطق القلم اذ یمشی مالک القدم فی سجنه العظیم البهآء
المشرق من افق سمآء عنایتی علیکم و علی من اقبل الی اللّه الفرد الخبیر
انتهی

\end{otherlanguage}

Yā aḥibbāʾī yā awliyāʾī ushkurū Allāha bi-mā asmaʿakum nidāʾahu
wa-arākum āthāra qalamihi wa-hadākum ilā ṣirāṭihi al-mustaqīmi. Ṭūbā
lakum wa-naʿīman lakum bi-mā fuztum bi-mā lā fāza bihi ʿulamāʾu al-arḍi
wa-fuqahāʾuhā. Iḥfaẓū hādhā al-maqāma al-aʿlā bi-smi Rabbikumu al-Nāṭiqi
fī hādhā al-maqāmi al-rafīʿi. Kadhālika naṭaqa al-qalamu idh yamshī
Māliku al-Qidami fī sijnihi al-ʿaẓīmi. Al-Bahāʾu al-mushriqu min ufuqi
samāʾi ʿināyatī ʿalaykum wa-ʿalā man aqbala ilā Allāhi al-Fardi
al-Khabīri. Intahā.

``My beloved ones! My loved ones! Thank God for enabling you to hear His
call, showing you the traces of His Pen, and guiding you to His straight
path. Blessed are you, and joy be yours, for attaining what the learned
and the religious scholars of the earth have not attained. Preserve this
most exalted station in the name of your Lord, who speaks in this lofty
station. Thus has the Pen spoken as the Lord of eternity walks in His
mighty prison. May the glory shining from the horizon of the heaven of
My loving care rest upon you and upon whoever has turned toward God, the
One, the All-Aware.'' End.

\subsection{58}\label{section-2289}

\begin{otherlanguage}{arabic}

للّه الحمد کلّ بآثار فائز و بانوار شمس بیان منوّر از حقّ میطلبم کلّ را توفیق
عطا فرماید و بر تبلیغ امر مؤیّد نماید و اسأله تعالی ان یمدّهم بجنود البیان
انّه هو المقتدر العزیز المستعان

\end{otherlanguage}

Li-llāhi al-ḥamd kull bih āthār fāʾiz va bih anvār-i shams-i bayān
munavvar. Az Ḥaqq mīṭalabam kull rā tawfīq ʿaṭā farmāyad va bar
tablīgh-i Amr muʾayyad namāyad wa-asʾaluhu taʿālā an yumiddahum
bi-junūdi al-bayāni innahu huwa al-Muqtadiru al-ʿAzīzu al-Mustaʿānu.

Praise be to God, all have been favored with the writings and illumined
by the rays of the sun of utterance. I ask God to grant success to all
and strengthen them in teaching the Cause, and I ask Him, exalted be He,
to reinforce them with the hosts of utterance. He is indeed the
All-Powerful, the Almighty, the One whose help is sought.

\subsection{59}\label{section-2290}

\begin{otherlanguage}{arabic}

دستخطّ دیگر آنمحبوب فؤاد که نیّر اتّحاد از افقش طالع رسید و کار خود دید
ذکرش نافذ و کلمه‌اس مسری للّه الحمد نفحه‌اش فرح آورد و عرفش بساط انبساط
گسترد و بعد از عرض در ساحت امنع اقدس این بیان از ملکوت الهی نازل قول
الرّبّ تعالی و تقدّس

\end{otherlanguage}

Dastkhaṭṭ-i dīgar-i ān maḥbūb-i fuʾād kih nayyir-i ittiḥād az ufuqash
ṭāliʿ rasīd va kār-i khud dīd. Dhikrash nāfidh va {[}corrupted text{]}
musrī. Li-llāhi al-ḥamd nafḥih-ash faraḥ āvard va ʿarfash bisāṭ-i
inbisāṭ gustard va baʿd az ʿarḍ dar sāḥat-i amnaʿ-i aqdas īn bayān az
malakūt-i ilāhī nāzil. Qawlu al-Rabbi taʿālā wa-taqaddasa.

Another letter from that beloved of the heart, from whose horizon the
sun of unity shines, arrived and accomplished its work. Its mention was
penetrating and {[}corrupted original text{]} pervasive. Praise be to
God, its breeze brought joy and its fragrance spread a banquet of
delight. After it was presented in the most exalted, most holy Court,
this utterance descended from the divine Kingdom. The words of the Lord,
exalted and sanctified be He:

\subsection{60}\label{section-2291}

\begin{otherlanguage}{arabic}

بسم الّذی به فتح باب الملکوت علی الملوک و المملوک

\end{otherlanguage}

Bi-smi alladhī bihi futiḥa bābu al-malakūti ʿalā al-mulūki wa-l-mamlūki.

``In the name of Him through whom the gate of the Kingdom has been
opened to kings and subjects.''

\subsection{61}\label{section-2292}

\begin{otherlanguage}{arabic}

علیک بهآء اللّه و عنایته و رحمته انّی اذکر نفسی فیهذا الحین من قبلک و نصفه
بما ینبغی لی و یلیق لساحة اقدس جزآء اعمالک فی سبیلی و انفاقک روحک و جسدک
و لسانک و قلمک فی ذکری و بیانی و خدمة امری العزیز المنیع لک کنز عند اللّه
انّه یوفّی اجرک انّه ولیّک فی الآخرة و الأولی و هو المشفق الکریم البهآء
المشرق من افق سمآء رحمتی علیک و علی من معک و علی الّذین ما منعتهم شئونات
العالم عن امر ربّه العلیم انتهی

\end{otherlanguage}

ʿAlayka Bahāʾuʾllāhi wa-ʿināyatuhu wa-raḥmatuhu. Innī adhkuru nafsī fī
hādhā al-ḥīni min qibalika wa-naṣifuhu bi-mā yanbaghī lī wa-yalīqu
li-sāḥati aqdas jazāʾa aʿmālika fī sabīlī wa-infāqika rūḥaka wa-jasadaka
wa-lisānaka wa-qalamaka fī dhikrī wa-bayānī wa-khidmati amrī al-ʿAzīzi
al-Manīʿi. Laka kanzun ʿinda Allāhi innahu yuwaffī ajraka innahu
waliyyuka fī al-ākhirati wa-l-ūlā wa-huwa al-Mushfiqu al-Karīmu.
Al-Bahāʾu al-mushriqu min ufuqi samāʾi raḥmatī ʿalayka wa-ʿalā man
maʿaka wa-ʿalā alladhīna mā manaʿathum shuʾūnu al-ʿālami ʿan amri
Rabbihi al-ʿAlīmi. Intahā.

``May the glory of God, His loving care, and His mercy rest upon you. At
this moment I make mention of My own Self on your behalf, and We
describe Him in terms befitting Me and worthy of a most holy Court, as a
recompense for your deeds in My path and for your offering of your soul,
your body, your tongue, and your pen in My remembrance, My praise, and
service to My mighty and unassailable Cause. You have a treasure with
God. He will pay your reward in full. He is your Guardian in the next
world and this world, and He is the Compassionate, the Generous. May the
glory shining from the horizon of the heaven of My mercy rest upon you,
upon those with you, and upon those whom the affairs of the world have
not kept from the Cause of his Lord, the All-Knowing.'' End.

\subsection{62}\label{section-2293}

\begin{otherlanguage}{persian}

دو آیه از سماء عنایت کبری نازل فی‌الحقیقه دو دنیا باو معادله نمینماید
هنیئاً لحضرتکم مریئاً لجنابکم اینکه ذکر اجتماع اولیای الهی را در بیت جناب
آقا میرزا سلیمان و آقا میرزا مهدی علیهما بهآء اللّه الأبدی فرمودند که از
جانب جناب امین علیه بهآء اللّه الأبهی مدعوّ بضیافت بودند بعد از عرض
اینمراتب و ذکر اسامی مجتمعه در ساحت امنع اقدس اعلی این آیات از افق سماء
عنایت حقّ جلّ جلاله اشراق فرمود

\end{otherlanguage}

Do āye az samā-ye ʿenāyat-e kobrā nāzel, fī'l-ḥaqīqe do donyā bā ū
moʿādele nemīnamāyad. Hanīʾan li-ḥaḍratikum, marīʾan li-janābikum. Īnke
dhekr-e ejtemāʿ-e awlīyā-ye elāhī rā dar bayt-e Janāb-e Āqā Mīrzā
Solaymān va Āqā Mīrzā Mahdī, ʿalayhimā bahāʾu'llāhi al-abadī, farmūdand
ke az jāneb-e Janāb-e Amīn, ʿalayhi bahāʾu'llāhi al-abhā, madʿovv
be-żīyāfat būdand, baʿd az ʿarż-e īn marāteb va dhekr-e asāmī-ye
mojtameʿe dar sāḥat-e amnaʿ-e aqdas-e aʿlā, īn āyāt az ofoq-e samā-ye
ʿenāyat-e Ḥaqq, jalla jalāluhu, eshrāq farmūd.

Two verses have descended from the heaven of supreme loving care; truly,
two worlds cannot equal them. May they bring you joy and delight! You
mentioned the gathering of the friends of God at the home of Áqá Mírzá
Sulaymán and Áqá Mírzá Mihdí---may the eternal glory of God rest upon
them---who had been invited to a feast by Amín---may the all-glorious
glory of God rest upon him. After these matters and the names of those
assembled were presented in the most exalted, most holy, most
inaccessible Court, these verses shone forth from the horizon of the
heaven of the loving care of God, exalted be His glory.

\subsection{63}\label{section-2294}

\begin{otherlanguage}{arabic}

قول الرّبّ تعالی و تقدّس هو الذّاکر اولیآئه بذکر لا تعادله الأذکار

\end{otherlanguage}

Qawlu al-Rabbi, taʿālā wa-taqaddasa: Huwa al-dhākiru awliyāʾahu
bi-dhikrin lā tuʿādiluhu al-adhkār.

The words of the Lord, exalted and sanctified be He: ``He is the One who
remembers His friends with a remembrance that no other remembrance can
equal.''

\subsection{64}\label{section-2295}

\begin{otherlanguage}{persian}

ا ح انّ المظلوم من شطر السّجن توجّه الیک و یقول لک الحمد یا اله الغیب و
الشّهود بما ایّدت عبدک علی الاقبال الیک و التّوجّه الی انوار وجهک و النّظر
الی افقک و التّمسّک بحبلک اسألک بأن تؤیّده علی الاستقامة علی امرک و تجعل
عمله کنزاً له عندک ثمّ اغفر له فضلاً من لدنک انّک انت الفضّال الکریم

\end{otherlanguage}

Alif Ḥāʾ. Inna al-Maẓlūma min shaṭri al-sijni tawajjaha ilayka
wa-yaqūlu: Laka al-ḥamdu yā ilāha al-ghaybi wa-l-shuhūd, bi-mā ayyadta
ʿabdaka ʿalā al-iqbāli ilayka wa-l-tawajjuhi ilā anwāri wajhika
wa-l-naẓari ilā ufuqika wa-l-tamassuki bi-ḥablika. Asʾaluka bi-an
tuʾayyidahu ʿalā al-istiqāmati ʿalā amrika wa-tajʿala ʿamalahu kanzan
lahu ʿindaka, thumma ighfir lahu faḍlan min ladunka. Innaka anta
al-Faḍḍālu al-Karīm.

A. Ḥ. The Wronged One has turned toward you from the prison and says:
``Praise be to You, O God of the unseen and the seen, for enabling Your
servant to turn toward You, to direct himself toward the lights of Your
face, to look toward Your horizon, and to hold fast to Your cord. I ask
You to strengthen him in steadfastness in Your Cause and to make his
deeds a treasure for him in Your presence. Then forgive him through Your
grace. You are indeed the Most Bountiful, the Generous.''

\subsection{65}\label{section-2296}

\begin{otherlanguage}{arabic}

یا نبیل قبل با انّ المظلوم ذکرک فی الواح شتّی و انزل لک فی هذه الایّام ما
لو یتضوّع عرفه لینجذب به افئدة العارفین قل الحمد للّه ربّ العالمین

\end{otherlanguage}

Yā Nabīl qabla Bāʾ, inna al-Maẓlūma dhakaraka fī alwāḥin shattā
wa-anzala laka fī hādhihi al-ayyāmi mā law yataḍawwaʿu ʿarfuhu
la-yanjadhibu bihi afʾidatu al-ʿārifīn. Qul: al-ḥamdu li-llāhi Rabbi
al-ʿālamīn.

O Nabíl before Bá! The Wronged One has remembered you in many Tablets
and has sent down for you in these days that whose fragrance, were it to
spread, would enrapture the hearts of those who recognize Him. Say:
``Praise be to God, the Lord of the worlds.''

\subsection{66}\label{section-2297}

\begin{otherlanguage}{arabic}

ج و قد فزت بذکری و آثار قلمی اعرف قدر هذا المقام الأعلی و قل

\end{otherlanguage}

Jīm Vāv. Qad fuzta bi-dhikrī wa-āthāri qalamī. Iʿrif qadra hādhā
al-maqāmi al-aʿlā wa-qul:

J. V. You have attained My remembrance and the traces of My Pen.
Recognize the worth of this most exalted station and say:

\subsection{67}\label{section-2298}

\begin{otherlanguage}{arabic}

لک الحمد یا مولی الوری بما ایّدتنی علی الاقبال الیک و الاقرار بما نطق به
لسان عظمتک اسألک بان تؤیّدنی علی حفظ ما اعطیتنی بجودک و کرمک انّک انت
الجواد الکریم

\end{otherlanguage}

Laka al-ḥamdu yā Mawlā al-warā bi-mā ayyadtanī ʿalā al-iqbāli ilayka
wa-l-iqrāri bi-mā naṭaqa bihi lisānu ʿaẓamatika. Asʾaluka bi-an
tuʾayyidanī ʿalā ḥifẓi mā aʿṭaytanī bi-jūdika wa-karamika. Innaka anta
al-Jawādu al-Karīm.

``Praise be to You, O Lord of all creation, for enabling me to turn
toward You and acknowledge what the Tongue of Your grandeur has spoken.
I ask You to enable me to preserve what You have given me through Your
bounty and generosity. You are indeed the All-Bountiful, the Generous.''

\subsection{68}\label{section-2299}

\begin{otherlanguage}{arabic}

یا اباطالب انّ المطلوب یذکرک و المظلوم توجّه الیک طوبی لطالب فاز بالمطلوب
و لعارف عرف المعروف و لعالم اقبل الی المعلوم خذ کتاب اللّه بقوّة من عنده
رغماً للّذین انکروا حجّته و برهانه الا انّهم من الأخسرین عند ملأ الأعلی فی
هذا المقام الکریم

\end{otherlanguage}

Yā Abā Ṭālib, inna al-Maṭlūba yadhkuruka wa-l-Maẓlūmu tawajjaha ilayka.
Ṭūbā li-ṭālibin fāza bi-l-Maṭlūbi wa-li-ʿārifin ʿarafa al-Maʿrūfa
wa-li-ʿālimin aqbala ilā al-Maʿlūm. Khudh kitāba Allāhi bi-quwwatin min
ʿindihi raghman li-lladhīna ankarū ḥujjatahu wa-burhānahu. Alā innahum
mina al-akhsarīna ʿinda malaʾi al-aʿlā fī hādhā al-maqāmi al-karīm.

O Abú Ṭálib! The One sought remembers you, and the Wronged One has
turned toward you. Blessed is the seeker who has attained the One
sought, the one endowed with recognition who has recognized its Object,
and the one endowed with knowledge who has turned toward the One to be
known. Take hold of the Book of God with strength from Him, despite
those who have denied His testimony and His proof. Truly, in the eyes of
the Concourse on high in this noble station, they are among those who
have suffered the greatest loss.

\subsection{69}\label{section-2300}

\begin{otherlanguage}{persian}

یا محمّد علی مکلّم طور بر عرش ظهور مستوی و اسمعک ندآئه و انزل لک آیاته و
ارسل الیک ما عجز عن ذکره کلّ عالم علیم و کلّ عارف خبیر یا علی مظلوم در سجن
ترا ذکر مینماید بگو ای اهل عالم امروز نیّر اینکلمه از افق سماء نصح الهی
مشرق و لائح خذوا سبیلی ضعوا سبل سوائی طوبی از برای نفسیکه باصغا فائز شد
و از کأس استقامت آشامید اشکر اللّه بما ذکرک انّه یذکر من احبّه انّه هو
الذّاکر العلیم

\end{otherlanguage}

Yā Muḥammad ʿAlī, Mokallem-e Ṭūr bar ʿarsh-e ẓohūr mostavī, wa-asmaʿaka
nidāʾahu wa-anzala laka āyātihi wa-arsala ilayka mā ʿajaza ʿan dhikrihi
kullu ʿālimin ʿalīmin wa-kullu ʿārifin khabīr. Yā ʿAlī, Maẓlūm dar sejn
torā dhekr mīnamāyad. Begū ey ahl-e ʿālam, emrūz nayyer-e īn kalame az
ofoq-e samā-ye noṣḥ-e elāhī moshreq va lāʾeḥ: Khudhū sabīlī, ḍaʿū subula
siwāʾī. Ṭūbā az barā-ye nafsī ke be-eṣghā fāʾez shod va az kaʾs-e
esteqāmat āshāmīd. Ushkur Allāha bi-mā dhakaraka. Innahu yadhkuru man
aḥabbahu. Innahu huwa al-Dhākiru al-ʿAlīm.

O Muḥammad ʿAlí! He who spoke on Sinai is established upon the throne of
revelation. He has caused you to hear His call, sent down His verses for
you, and sent you that which every learned scholar and every discerning
knower is powerless to describe. O ʿAlí! The Wronged One remembers you
in prison. Say: ``O people of the world! Today the sun of this utterance
shines resplendent from the horizon of the heaven of divine counsel:
`Take My path; leave aside all paths other than Mine.'\,'' Blessed is
the soul who has attained a hearing and drunk from the cup of
steadfastness. Thank God for remembering you. He remembers whoever loves
Him. He is indeed the One who remembers, the All-Knowing.

\subsection{70}\label{section-2301}

\begin{otherlanguage}{persian}

یا محمّد قبل علی بآثار قلم اعلی فائز شدی للّه الحمد حجبات اوهام را باسم
مالک انام شقّ نمودی امروز میزان اعظم امام وجوه عالم منصوب و ینادی بأعلی
النّدآء انّی انا الممیّز العلیم طوبی لمن ما منعته قدرة و ما خوّفته سطوة و
اقبل الی الأفق الأعلی باستقامة کبری انّه من اهل البهآء فی لوح مبین

\end{otherlanguage}

Yā Muḥammad qabl-e ʿAlī, be-āthār-e Qalam-e Aʿlā fāʾez shodī. Li-llāhi
al-ḥamd, ḥojobāt-e awhām rā be-esm-e Mālek-e anām shaqq nemūdī. Emrūz
mīzān-e aʿẓam emām-e vojūh-e ʿālam manṣūb, wa-yunādī bi-aʿlā al-nidāʾ:
Innī anā al-Mumayyizu al-ʿAlīm. Ṭūbā li-man mā manaʿathu qudratun wa-mā
khawwafathu saṭwatun wa-aqbala ilā al-ufuqِ al-aʿlā bi-istiqāmatin kubrā.
Innahu min ahli al-Bahāʾ fī lawḥin mubīn.

O Muḥammad before ʿAlí! You have attained the writings of the Most
Exalted Pen. Praise be to God, you have torn through the veils of vain
imaginings in the name of the Lord of humankind. Today the Most Great
Balance is set before the eyes of the world and proclaims at the highest
pitch: ``I am indeed the All-Knowing Discerner.'' Blessed is the one
whom no power has hindered and no might has frightened, and who has
turned toward the Most Exalted Horizon with supreme steadfastness. He is
numbered among the people of Bahá in a manifest Tablet.

\subsection{71}\label{section-2302}

\begin{otherlanguage}{persian}

یا رضا طوبی از برای نفسیکه امروز بکلمهٴ رضا فائز شد او از اهل فردوس اعلی
مذکور عباد را اوهام آمال از مآل منع نموده بایّامی معدوده از ملکوت باقی
محروم و ممنوع ذرهم فی خوضهم و قل لک الحمد یا الهی بما ایّدتنی علی الاقبال
و انطقتنی بثنائک الجمیل اسألک بان تقدّر لی ما یقرّبنی الیک انّک انت السّامع
المجیب

\end{otherlanguage}

Yā Reżā, ṭūbā az barā-ye nafsī ke emrūz be-kalame-ye reżā fāʾez shod. Ū
az ahl-e Ferdaws-e Aʿlā madhkūr. ʿEbād rā awhām-e āmāl az maʾāl manʿ
nemūde, be-ayyāmī maʿdūde az malakūt-e bāqī maḥrūm va mamnūʿ. Dharhum fī
khawḍihim wa-qul: Laka al-ḥamdu yā ilāhī bi-mā ayyadtanī ʿalā al-iqbāli
wa-anṭaqtanī bi-thanāʾika al-jamīl. Asʾaluka bi-an tuqaddira lī mā
yuqarribunī ilayka. Innaka anta al-Sāmiʿu al-Mujīb.

O Riḍá! Blessed is the soul who today has attained the word of good
pleasure. He is reckoned among the people of the Most Exalted Paradise.
Delusive hopes have kept people from their ultimate end; for the sake of
a few numbered days, they are deprived of and barred from the
everlasting Kingdom. Leave them to their vain pursuits and say: ``Praise
be to You, O my God, for enabling me to turn toward You and causing me
to voice Your beautiful praise. I ask You to ordain for me what will
draw me near to You. You are indeed the All-Hearing, the Responsive.''

\subsection{72}\label{section-2303}

\begin{otherlanguage}{persian}

یا حاجی علی قبل اکبر (م‌ی) در منظر اکبر ذکرت مذکور و بعنایت حقّ فائز امروز
دیباج کتاب عدل باینکلمهٴ علیا مزیّن ای اهل عالم شما قطره‌های یکبحر و اوراق
یکشجرید اختلاف و نفاق را بگذارید به اتّحاد و اتّفاق تمسّک نمائید هذا ما
امرتم به من لدی اللّه مالک یوم التّلاق ولکن آذان اهل عالم از اصغا محروم و
از تقریر ممنوع چه که زمام کلّ در قبضهٴ جهلا بوده و هست الّا من شآء اللّه
طوبی لک بما عرفت المقصود اذ انکره المشرکون

\end{otherlanguage}

Yā Ḥājī ʿAlī qabl-e Akbar (Mīm Yāʾ), dar manẓar-e akbar dhekrat madhkūr
va be-ʿenāyat-e Ḥaqq fāʾez. Emrūz dībā j-e ketāb-e ʿadl be-īn kalame-ye
ʿolyā mozayyan: Ey ahl-e ʿālam, shomā qaṭrehā-ye yek baḥr va awrāq-e yek
shajarīd. Ekhtelāf va nefāq rā begodhārīd, be etteḥād va ettefāq
tamassok namāʾīd. Hādhā mā umirtum bihi min ladā Allāhi Māliki yawmi
al-talāq. Va-lāken ādhān-e ahl-e ʿālam az eṣghā maḥrūm va az taqrīr
mamnūʿ, che ke zemām-e koll dar qabże-ye johalā būde va hast, illā man
shāʾa Allāh. Ṭūbā laka bi-mā ʿarafta al-Maqṣūda idh ankarahu
al-mushrikūn.

O Ḥájí ʿAlí before Akbar (M. Y.)! You are remembered in the Most Great
Scene and have received the loving care of God. Today the preface to the
Book of Justice is adorned with this exalted utterance: ``O people of
the world! You are the drops of one sea and the leaves of one tree. Set
aside dissension and discord; hold fast to unity and concord.'' This is
what you have been commanded by God, the Lord of the Day of Meeting. Yet
the ears of the people of the world are deprived of hearing and barred
from the exposition, for the reins of all have been and remain in the
grasp of the ignorant, except those whom God wills. Blessed are you for
recognizing the Desired One when those who associate partners with God
denied Him.

\subsection{73}\label{section-2304}

\begin{otherlanguage}{persian}

یا علی اکبر امروز نیّر ذکر از افق سماء قلم اعلی مشرق و لائح مقصود آنکه
نفوسی در ظلّ امر تربیت شوند و بر اطفاء نار بغضا که در صدور احزاب مختلفهٴ
عالم مشتعل است قیام نمایند شاید بصر عالم بچنین یومی فائز شود عناد سبب
خراب و اتّحاد علّة عمار بوده و هست از حقّ بطلب نفوسیکه خود را بافق اعلی
نسبت میدهند باصغاء حقیقی فائز شوند و عمل نمایند آنچه را که بآن مأمورند
طوبی لک بما سمعت النّدآء و اجبت مولاک

\end{otherlanguage}

Yā ʿAlī Akbar, emrūz nayyer-e dhekr az ofoq-e samā-ye Qalam-e Aʿlā
moshreq va lāʾeḥ. Maqṣūd ānke nofūsī dar ẓell-e Amr tarbīyat shavand va
bar eṭfā-ye nār-e baghżā ke dar ṣodūr-e aḥzāb-e mokhtalefe-ye ʿālam
moshtaʿel ast qīyām namāyand. Shāyad baṣar-e ʿālam be-chenīn yawmī fāʾez
shavad. ʿEnād sabab-e kharāb va etteḥād ʿellat-e ʿamār būde va hast. Az
Ḥaqq beṭalab nofūsī ke khod rā be-Ofoq-e Aʿlā nesbat mīdehand
be-eṣghā-ye ḥaqīqī fāʾez shavand va ʿamal namāyand ānche rā ke be-ān
maʾmūrand. Ṭūbā laka bi-mā samiʿta al-nidāʾa wa-ajabta Mawlāka.

O ʿAlí-Akbar! Today the sun of remembrance shines resplendent from the
horizon of the heaven of the Most Exalted Pen. The purpose is that souls
should be educated under the shelter of the Cause and arise to
extinguish the fire of hatred burning in the breasts of the diverse
peoples of the world. May the eyes of the world behold such a day.
Enmity has always been the cause of ruin, and unity the source of
prosperity. Ask God that those who claim allegiance to the Most Exalted
Horizon may attain true hearing and carry out what they are commanded to
do. Blessed are you for hearing the call and answering your Lord.

\subsection{74}\label{section-2305}

\begin{otherlanguage}{persian}

یا محمّد ق امروز عالم بانوار ذکر اعظم فائز و منوّر و مظلوم بأعلی النّداء
عباد را آگاه فرمود و راه نمود طوبی از برای نفوسیکه بیوم فائز شدند و
مقصود را شناختند از حقّ میطلبیم کلّ را تأیید فرماید و از نفحات آیاتش محروم
ننماید

\end{otherlanguage}

Yā Muḥammad Qāf, emrūz ʿālam be-anvār-e Dhekr-e Aʿẓam fāʾez va monavvar
va Maẓlūm be-aʿlā al-nedā ʿebād rā āgāh farmūd va rāh nemūd. Ṭūbā az
barā-ye nofūsī ke be-Yawm fāʾez shodand va Maqṣūd rā shenākhtand. Az
Ḥaqq mīṭalabīm koll rā taʾyīd farmāyad va az nafaḥāt-e āyātash maḥrūm
nanamāyad.

O Muḥammad Q.! Today the world has received and been illumined by the
lights of the Most Great Remembrance, and the Wronged One has, with His
loudest call, awakened His servants and shown them the way. Blessed are
the souls who have attained the Day and recognized the Desired One. We
ask God to strengthen everyone and not deprive them of the fragrant
breezes of His verses.

\subsection{75}\label{section-2306}

\begin{otherlanguage}{arabic}

یا کربلائی صادق صدق مبین را معرضین بیان انکار نمودند نفوسیکه از مبدء و
منتهی غافل بر ضرّ امر قیام کردند و بجمیع اسباب متمسّک که شاید بآمال خود که
وبال دنیا و آخرتست برسند از حقّ میطلبیم اولیای خود را مستقیم فرماید
بشأنیکه ضوضای عالم ایشانرا منع ننماید و محروم نسازد

\end{otherlanguage}

Yā Karbalāʾī Ṣādeq, Ṣedq-e Mobīn rā moʿreżīn-e Bayān enkār nemūdand.
Nofūsī ke az mabdaʾ va montahā ghāfel bar żorr-e Amr qīyām kardand va
be-jamīʿ-e asbāb motamassek ke shāyad be-āmāl-e khod ke vabāl-e donyā va
ākherat ast beresand. Az Ḥaqq mīṭalabīm awlīyā-ye khod rā mostaqīm
farmāyad be-shaʾnī ke żawżā-ye ʿālam īshān rā manʿ nanamāyad va maḥrūm
nasāzad.

O Karbaláʾí Ṣádiq! The opponents among the people of the Bayán have
denied the manifest Truth. Souls heedless of the Beginning and the End
have arisen to harm the Cause and have seized upon every means in the
hope of fulfilling their ambitions, which bring ruin in this world and
the next. We ask God to make His friends so steadfast that the clamor of
the world will neither hinder nor deprive them.

\subsection{76}\label{section-2307}

\begin{otherlanguage}{arabic}

یا محمّد مهدی اذا اخذک جذب ندائی و فزت برحیق بیانی قل الهی الهی زد نور
بصری لمشاهدة افقک الأعلی و قوّة سمعی لاصغاء صریر قلمک الأبهی ای ربّ انا
الّذی توکّلت علیک و فوّضت امری الیک قدّر لی خیر الآخرة و الأولی لا اله الّا
انت القویّ القدیر

\end{otherlanguage}

Yā Muḥammad Mahdī, idhā akhadhaka jadhbu nidāʾī wa-fuzta bi-raḥīqi
bayānī, qul: Ilāhī ilāhī, zid nūra baṣarī li-mushāhadati ufuqika al-aʿlā
wa-quwwata samʿī li-iṣghāʾi ṣarīri qalamika al-abhā. Ay Rabbi, anā
alladhī tawakkaltu ʿalayka wa-fawwaḍtu amrī ilayka. Qaddir lī khayra
al-ākhirati wa-l-ūlā. Lā ilāha illā anta al-Qawiyyu al-Qadīr.

O Muḥammad Mihdí! When the attraction of My call seizes you and you
attain the choice wine of My utterance, say: ``My God, my God! Increase
the light of my sight that I may behold Your Most Exalted Horizon, and
the strength of my hearing that I may hear the movement of Your
All-Glorious Pen. O Lord! I am the one who has placed his trust in You
and committed his affairs to You. Ordain for me the good of the next
world and this world. There is no God but You, the Strong, the
All-Powerful.''

\subsection{77}\label{section-2308}

\begin{otherlanguage}{persian}

یا مهدی هادی آمد انکارش نمودند حقّ آمد تکذیبش کردند سمع باطل قابل اصغاء
نداء حقّ نه و بصر عاطل لایق مشاهده نه للّه الحمد معشر اولیا بابصار حدیده
فائز گشته‌اند ایشانند نفوسیکه میفرماید اوتوا بصائر من اللّه شنیدند و دیدند
و فائز شدند بآنچه که مقصود از آفرینش بوده

\end{otherlanguage}

Yā Mahdī, Hādī āmad, enkārash nemūdand. Ḥaqq āmad, takdhībash kardand.
Samʿ-e bāṭel qābel-e eṣghā-ye nedā-ye Ḥaqq na, va baṣar-e ʿāṭel lāyeq-e
moshāhade na. Li-llāhi al-ḥamd, maʿshar-e awlīyā be-abṣār-e ḥadīde fāʾez
gashte-and. Īshānand nofūsī ke mīfarmāyad: Ūtū baṣāʾira min Allāh.
Shenīdand va dīdand va fāʾez shodand be-ānche ke maqṣūd az āfarīnesh
būde.

O Mihdí! The Guide came; they denied Him. The Truth came; they accused
Him of falsehood. The ear given to falsehood cannot hear the call of
Truth, nor is the idle eye fit to behold it. Praise be to God, the
company of His friends has been endowed with keen sight. They are the
souls of whom He says: ``They were given insight by God.'' They heard
and saw and attained that which was the purpose of creation.

\subsection{78}\label{section-2309}

\begin{otherlanguage}{persian}

یا ورقا علیک بهائی و عنایتی منتهی شد اذکار بذکر تو طوبی لک و نعیماً لک
لازال در ساحت اقدس مذکوری و بعنایات لانهایهٴ مقصود عالمیان فائز عزیز و
روح را من قبل اللّه تکبیر برسان و آثار شفقت و عنایت و ظهورات مرحمت از قبل
مظلوم اظهار نما تا بعنایت حقّ نموّ نمایند و ثمر دهند انّه هو المشفق الکریم
و هو الغفور الرّحیم انتهی

\end{otherlanguage}

Yā Varqā, ʿalayka bahāʾī wa-ʿināyatī. Montahī shod adhkār be-dhekr-e to.
Ṭūbā laka wa-naʿīman laka. Lāzāl dar sāḥat-e aqdas madhkūrī va
be-ʿenāyāt-e lānehāye-ye Maqṣūd-e ʿālamīyān fāʾez. ʿAzīz va Rūḥ rā men
qebal-e Allāh takbīr beresān va āthār-e shafaqat va ʿenāyat va ẓohūrāt-e
marḥamat az qebal-e Maẓlūm eẓhār namā tā be-ʿenāyat-e Ḥaqq nomovv
namāyand va thamar dehand. Innahu huwa al-Mushfiqu al-Karīmu wa-huwa
al-Ghafūru al-Raḥīm. Entahā.

O Varqá! My glory and My loving care rest upon you. These remembrances
conclude with your remembrance. Blessed are you, and may bliss be yours!
You are ever remembered in the Most Holy Court and receive the boundless
favors of the Desired One of all peoples. Convey to ʿAzíz and Rúḥ a
greeting of glory on behalf of God, and show them signs of affection and
loving care and expressions of mercy on behalf of the Wronged One, that
through God's loving care they may grow and bear fruit. He is indeed the
Compassionate, the Generous, and He is the Ever-Forgiving, the Merciful.
End.

\subsection{79}\label{section-2310}

\begin{otherlanguage}{persian}

للّه الفضل و العطآء انجم سماء عرفان هر یک کسب ضیاء نموده و از انوار نیّر
اعظم منوّر و روشن این اذکار و اوصاف فی‌الحقیقه حاکی از حدود نفوسست چگونه
قابل این مقام آنچه از خادم ظاهر قابل بساط اللّه نبوده و نیست له الحمد و
الثّنآء و له الشّکر و البهآء که مقامات و مراتب عجز و نیستی را القا فرمود و
قبول نمود اینمقام را هیچ عمل و هیچ ذکر و هیچ فکر چه از عالم ظاهر و چه از
عالم باطن معادله نمینماید

\end{otherlanguage}

Li-llāhi al-faḍlu wa-l-ʿaṭāʾ. Anjom-e samā-ye ʿerfān har yek kasb-e żīyā
nemūde va az anvār-e Nayyer-e Aʿẓam monavvar va rowshan. Īn adhkār va
awṣāf fī'l-ḥaqīqe ḥākī az ḥodūd-e nofūs ast, chegūne qābel-e īn maqām?
Ānche az khādem ẓāher qābel-e besāṭ-e Allāh nabūde va nīst. Lahu
al-ḥamdu wa-l-thanāʾu wa-lahu al-shukru wa-l-bahāʾ, ke maqāmāt va
marāteb-e ʿajz va nīstī rā elqā farmūd va qabūl nemūd. Īn maqām rā hīch
ʿamal va hīch dhekr va hīch fekr, che az ʿālam-e ẓāher va che az ʿālam-e
bāṭen, moʿādele nemīnamāyad.

To God belong grace and bounty. Each star in the heaven of recognition
has acquired radiance and is illumined and brightened by the lights of
the Most Great Luminary. These remembrances and descriptions in truth
reflect the limitations of souls; how could they be worthy of this
station? Whatever has proceeded from this servant has never been, and is
not, worthy of the court of God. To Him belong praise and acclaim,
thanks and glory, for He imparted the stations and degrees of
helplessness and nothingness and accepted them. No deed, no remembrance,
and no thought, whether of the outward or the inward world, can equal
this station.

\subsection{80}\label{section-2311}

\begin{otherlanguage}{persian}

ذکر امآء اللّه علیهنّ بهآء اللّه را نمودند بعد از عرض در حضور هر یک بطراز
ذکر مزیّن و بعنایات ملکوت بیان فائز نسأل ربّنا و ربّ الأرض و السّمآء ان یقدّر
لهنّ خیر الآخرة و الأولی سبحان‌اللّه سرّ امر و حکمت بالغه از نظر مستور
عالمیکه سالها بگمان خود طلب علم نمود از برای معرفت معلوم و چون انوار وجه
معلوم مشرق معرض مشاهده شد و ورقات و امآء اللّه که در دبستان وارد نشده‌اند
و از علوم بر حسب ظاهر بهره نبرده‌اند بنور عرفان منوّر و از کأس استقامت
مرزوق گشتند انّ الفضل بیده و فی قبضة اقتداره یعطی و یمنع و هو العزیز
المختار

\end{otherlanguage}

Dhekr-e emāʾ-e Allāh, ʿalayhinna bahāʾu'llāh, rā nemūdand. Baʿd az ʿarż
dar ḥożūr, har yek be-ṭerāz-e dhekr mozayyan va be-ʿenāyāt-e malakūt-e
bayān fāʾez. Nasʾalu Rabbanā wa-Rabba al-arḍi wa-l-samāʾi an yuqaddira
lahunna khayra al-ākhirati wa-l-ūlā. Sobḥān Allāh, serr-e Amr va
ḥekmat-e bāleghe az naẓar mastūr. ʿĀlemī ke sālhā be-gomān-e khod
ṭalab-e ʿelm nemūd az barā-ye maʿrefat-e Maʿlūm, va chon anvār-e vajh-e
Maʿlūm moshreq, moʿreż-e moshāhade shod, va varaqāt va emāʾ-e Allāh ke
dar dabestān vāred nashode-and va az ʿolūm bar ḥasab-e ẓāher bahre
naborde-and be-nūr-e ʿerfān monavvar va az kaʾs-e esteqāmat marzūq
gashtand. Inna al-faḍla bi-yadihi wa-fī qabḍati iqtidārihi, yuʿṭī
wa-yamnaʿu wa-huwa al-ʿAzīzu al-Mukhtār.

You mentioned the handmaids of God---may the glory of God rest upon
them. After this was presented in His presence, each was adorned with
the ornament of remembrance and received the favors of the kingdom of
utterance. We ask our Lord, the Lord of earth and heaven, to ordain for
them the good of the next world and this world. Glory be to God! The
mystery of the Cause and its consummate wisdom are hidden from sight. A
scholar who for years, as he imagined, sought knowledge in order to
recognize the One to be known turned away from beholding Him when the
lights of His face shone forth; yet the leaves and handmaids of God, who
had not attended school and had outwardly acquired no share of learning,
were illumined by the light of recognition and given to drink from the
cup of steadfastness. Grace is indeed in His hand and in the grasp of
His might. He gives and withholds, and He is the Almighty, the One who
chooses.

\subsection{81}\label{section-2312}

\begin{otherlanguage}{arabic}

و اینکه مرقوم داشتند از انفسیکه موعود بوده و بعد بواسطهٴ تکسّر مزاج نیامد
و رجای ذکر اسم در محضر انور نمود جناب آقا علی میلانی علیه بهآء اللّه
میباشند بعد از عرض اینفقره امام وجه مولی البریّه این آیات از ملکوت بیان
رحمن نازل

\end{otherlanguage}

Va īnke marqūm dāshtand az anfosī ke mowʿūd būde va baʿd be-vāseṭe-ye
takassor-e mezāj nayāmad va rajā-ye dhekr-e esm dar maḥżar-e anvar
nemūd, Janāb-e Āqā ʿAlī Mīlānī, ʿalayhi bahāʾu'llāh, mībāshand. Baʿd az
ʿarż-e īn feqre emām-e vajh-e Mawlā al-barīye, īn āyāt az malakūt-e
bayān-e Raḥmān nāzel.

As for what you wrote concerning the person who was expected but later
did not come because of ill health, and who expressed the hope that his
name might be mentioned in the radiant Presence: he is Áqá ʿAlí
Mílání---may the glory of God rest upon him. After this matter was
presented before the face of the Lord of creation, these verses
descended from the kingdom of the utterance of the All-Merciful.

\subsection{82}\label{section-2313}

\begin{otherlanguage}{persian}

قول الرّبّ تعالی و تقدّس یا علی ذکرت مذکور و کتاب مبین بتو متوجّه و ناطق
قبول نمودی آنچه را که کل از او معرضند و اعتراف نمودی آنچه را که قوم
منکرند رحیق عرفان را در یومی آشامیدی که اهل عالم از او محروم و محجوب حال
فکر نما که چه لسان باید و چه بیان شاید که از عهدهٴ شکر برآید

\end{otherlanguage}

Qawl al-Rabb taʿālā wa-taqaddasa: Yā ʿAlī, dhekrat madhkūr va Kitāb-i
Mubīn be-to motavajjeh va nāṭeq. Qabūl namūdī ānche rā ke koll az ū
moʿreżand va eʿterāf namūdī ānche rā ke qowm monkerand. Raḥīq-i ʿerfān
rā dar yowmī āshāmīdī ke ahl-i ʿālam az ū maḥrūm va maḥjūb. Ḥāl fekr
namā ke che lesān bāyad va che bayān shāyad ke az ʿohde-ye shokr
barāyad.

The words of the Lord, exalted and sanctified is He: O ʿAlī, you are
remembered, and the Perspicuous Book turns toward you and speaks. You
accepted what all have turned away from, and acknowledged what the
people deny. You drank the choice wine of recognition on a day when the
people of the world were deprived of it and veiled from it. Now consider
what tongue is needed, and what utterance would be fitting, to offer
adequate thanks.

\subsection{83}\label{section-2314}

\begin{otherlanguage}{arabic}

قل لک الحمد یا الهی و لک الفضل یا موجدی و رجائی اشهد انّک انزلت النّعمة و
ارسلت المآئدة و اظهرت الحجّة اسألک یا موجد العالم بالاسم الأعظم بان
تؤیّدنی علی ما یقرّبنی الیک علی شأن لا تمنعنی عن ساحة القرب ضوضآء خلقک و
لا تبعدنی عن الحضور شبهات عبادک ای ربّ ترانی متمسّکاً بحبل جودک اسألک
بأنوار نیّر علمک و اسرار کتابک بأن تکتب لی ما یکون باقیاً ببقآء ملکوتک و
جبروتک لا اله الّا انت المعطی الباذل العلیم الحکیم انتهی

\end{otherlanguage}

Qul: Laka al-ḥamdu yā ilāhī wa-laka al-faḍlu yā mūjidī wa-rajāʾī.
Ashhadu annaka anzalta al-niʿmata wa-arsalta al-māʾidata wa-aẓharta
al-ḥujjata. Asʾaluka yā mūjida al-ʿālami bi-al-ismi al-aʿẓami bi-an
tuʾayyidanī ʿalā mā yuqarribunī ilayka ʿalā shaʾnin lā tamnaʿunī ʿan
sāḥati al-qurbi ḍawḍāʾu khalqika wa-lā tubʿidanī ʿan al-ḥuḍūri shubuhātu
ʿibādika. Ay rabbi, tarānī mutamassikan bi-ḥabli jūdika. Asʾaluka
bi-anwāri nayyiri ʿilmika wa-asrāri kitābika bi-an taktuba lī mā yakūnu
bāqiyan bi-baqāʾi malakūtika wa-jabarūtika. Lā ilāha illā anta al-Muʿṭī
al-Bādhil al-ʿAlīm al-Ḥakīm. Intahā.

Say: Praise be to You, my God, and all bounty is Yours, O my Creator and
my Hope. I testify that You have sent down Your blessing, sent forth the
table, and made manifest the proof. I ask You, O Creator of the world,
by the Most Great Name, to aid me in what will draw me near to You, so
that neither the clamor of Your creatures may bar me from the court of
nearness nor the doubts of Your servants keep me far from Your presence.
O Lord, You see me holding fast to the cord of Your generosity. I ask
You, by the lights of the sun of Your knowledge and the mysteries of
Your Book, to ordain for me what will endure as long as Your kingdom and
Your dominion endure. There is no God but You, the Bestower, the
Munificent, the All-Knowing, the All-Wise. End.

\subsection{84}\label{section-2315}

\begin{otherlanguage}{persian}

و همچنین ذکر مجلس آخر بیت جناب خلیل علیه ٩ ٦٦ {[}بهآء الله{]} را نمودند
و ذکر جناب آقا ملّا مصطفی علیه بهآء اللّه و جناب آقا میرزا عباداللّه علیه ٩
٦٦ ٩{[}بهآء الله{]} فرمودند بعد از عرض لسان عظمت باین آیات مشرقات ناطق
قوله تبارک و تعالی

\end{otherlanguage}

Va hamchenīn dhekr-i majles-i ākhar-i bayt-i Janāb-i Khalīl ʿalayhi 9 66
{[}Bahāʾ Allāh{]} rā namūdand va dhekr-i Janāb-i Āqā Mollā Moṣṭafā
ʿalayhi Bahāʾ Allāh va Janāb-i Āqā Mīrzā ʿEbād Allāh ʿalayhi 9 66 9
{[}Bahāʾ Allāh{]} farmūdand. Baʿd az ʿarż, Lesān-i ʿAẓamat be-īn āyāt-i
moshreqāt nāṭeq. Qawluhu tabāraka wa-taʿālā:

You also mentioned the last gathering at the home of Jināb-i
Khalīl---upon him be 9 66 {[}the glory of God{]}---and mentioned Jināb-i
Āqā Mullā Muṣṭafā---may the glory of God rest upon him---and Jināb-i Āqā
Mīrzā ʿIbād Allāh---upon him be 9 66 9 {[}the glory of God{]}. After
this was presented, the Tongue of Grandeur spoke these radiant verses.
His words, blessed and exalted is He:

\subsection{85}\label{section-2316}

\begin{otherlanguage}{arabic}

یا مصطفی علیک بهائی سدرهٴ باقیهٴ ابهی در فردوس اعلی باینکلمهٴ علیا ناطق
طوبی لمن عرف مقام یوم ربّه و عمل بما ینبغی له للّه الحمد ندای الهی را اصغا
نمودی و بافق اعلی اقبال کردی همّت را بلند نما انّه معک و یؤیّدک علی ذکره و
ثنآئه و خدمة امره

\end{otherlanguage}

Yā Muṣṭafā, ʿalayka bahāʾī. Sedre-ye bāqīye-ye abhā dar ferdows-i aʿlā
be-īn kalame-ye ʿolyā nāṭeq: Ṭūbā li-man ʿarafa maqāma yawmi rabbihi
wa-ʿamila bi-mā yanbaghī lahu. Li-llāh al-ḥamd, nedā-ye elāhī rā eṣghā
namūdī va be-of oq-i aʿlā eqbāl kardī. Hemmat rā boland namā. Innahu
maʿaka wa-yuʾayyiduka ʿalā dhikrihi wa-thanāʾihi wa-khidmati amrihi.

O Muṣṭafā, My glory rest upon you. The eternal, all-glorious Tree in the
highest Paradise speaks this exalted word: ``Blessed is the one who has
recognized the station of the Day of his Lord and acted as befits it.''
Praise be to God, you have heeded the divine call and turned toward the
Most Exalted Horizon. Set your sights high. He is with you and aids you
in remembering Him, praising Him, and serving His Cause.

\subsection{86}\label{section-2317}

\begin{otherlanguage}{persian}

قل الهی الهی انا عبدک و ابن عبدک و ابن امتک اشهد انّک خلقت لسانی لذکرک و
بصری لمشاهدة آثارک و اذنی لاصغآء ندائک و قلبی لعرفانک و یدی لأخذ کتابک و
رجلی للاستقامة علی صراطک ای ربّ ایّدنی علی ما یتضوّع به عرف رضائک هذا خیر
لی عمّا فی الآخرة و الأولی انّک انت مولی الوری و ربّ العرش و الثّری لا اله
الّا انت المبیّن الحکیم انتهی

\end{otherlanguage}

Qul: Ilāhī ilāhī, anā ʿabduka wa-ibnu ʿabdika wa-ibnu amatika. Ashhadu
annaka khalaqta lisānī li-dhikrika wa-baṣarī li-mushāhadati āthārika
wa-udhunī li-iṣghāʾi nidāʾika wa-qalbī li-ʿirfānika wa-yadī li-akhdhi
kitābika wa-rijlī li-al-istiqāmati ʿalā ṣirāṭika. Ay rabbi, ayyidnī ʿalā
mā yataḍawwaʿu bihi ʿarfu riḍāʾika. Hādhā khayrun lī ʿammā fī
al-ākhirati wa-al-ūlā. Innaka anta mawlā al-warā wa-rabbu al-ʿarshi
wa-al-tharā. Lā ilāha illā anta al-Mubayyin al-Ḥakīm. Intahā.

Say: My God, my God, I am Your servant, the son of Your servant and Your
handmaid. I testify that You created my tongue to remember You, my sight
to behold Your signs, my ear to heed Your call, my heart to recognize
You, my hand to take hold of Your Book, and my foot to remain steadfast
on Your path. O Lord, aid me in that through which the fragrance of Your
good pleasure is diffused. This is better for me than all that is in the
world to come and in this world. You are the Master of all creation and
the Lord of the throne and the earth. There is no God but You, the
Expounder, the All-Wise. End.

\subsection{87}\label{section-2318}

\begin{otherlanguage}{arabic}

و هذا ما نزّل لجناب آقا میرزا عباداللّه علیه ٩ ٦٦ {[}بهآء الله{]} قوله
تبارک و تعالی

\end{otherlanguage}

Wa-hādhā mā nuzzila li-Janāb-i Āqā Mīrzā ʿIbād Allāh ʿalayhi 9 66
{[}Bahāʾ Allāh{]}. Qawluhu tabāraka wa-taʿālā:

And this is what was revealed for Jināb-i Āqā Mīrzā ʿIbād Allāh---upon
him be 9 66 {[}the glory of God{]}. His words, blessed and exalted is
He:

\subsection{88}\label{section-2319}

\begin{otherlanguage}{arabic}

هو الشّاهد الخبیر

\end{otherlanguage}

Huwa al-Shāhid al-Khabīr.

He is the Witness, the All-Aware.

\subsection{89}\label{section-2320}

\begin{otherlanguage}{persian}

حقّ جلّ جلاله مبشّر فرستاد و مژده داد تا اهل ارض بطراز فرح و سرور مزیّن
گردند و با کمال انبساط قصد بساط حقّ نمایند ولکن چون افئدهٴ عباد باوهام
مشغول از انوار صبح یقین محروم مانده‌اند و این مژدهٴ جانبخش سبب احزان شد
چه که قلب بغیر مشغول و آذان بانکر الأصوات مأنوس سبحان‌اللّه فرات شیرین
امام وجوه جاری و نیّر اعظم از اعلی افق عالم مشرق مع‌ذلک خلق غافل در هر حال
باید از حقّ جلّ جلاله مسئلت نمود عباد خود را محروم نفرماید حمد کن مقصود
عالمیان را که بشارتش جان عطا فرمود و ترا بر توجّه و اقبال و عرفان تأیید
نمود انّه هو الفضّال الکریم لا اله الّا هو المقتدر العلیم الحکیم انتهی

\end{otherlanguage}

Ḥaqq jalla jalāluhu mobashsher ferestād va mozhde dād tā ahl-i arż
be-ṭerāz-i faraḥ va sorūr mozayyan gardand va bā kamāl-i enbesāṭ qaṣd-i
besāṭ-i Ḥaqq namāyand. Va-līkan chon afʾede-ye ʿebād be-awhām mashghūl,
az anvār-i ṣobḥ-i yaqīn maḥrūm mānde-and va īn mozhde-ye jānbakhsh
sabab-i aḥzān shod, che ke qalb be-ghayr mashghūl va ādhān be-ankar
al-aṣwāt maʾnūs. Subḥān Allāh, forāt-i shīrīn amām-i vojūh jārī va
nayyer-i aʿẓam az aʿlā ofoq-i ʿālam moshreq, maʿa dhālek khalq ghāfel.
Dar har ḥāl bāyad az Ḥaqq jalla jalāluhu masʾalat namūd ʿebād-i khod rā
maḥrūm nafarmāyad. Ḥamd kon Maqṣūd-i ʿālamiyān rā ke beshāratash jān
ʿaṭā farmūd va to rā bar tavajjoh va eqbāl va ʿerfān taʾyīd namūd.
Innahu huwa al-Faḍḍāl al-Karīm. Lā ilāha illā huwa al-Muqtadir al-ʿAlīm
al-Ḥakīm. Intahā.

God, glorious is His glory, sent a Herald and gave glad tidings so that
the people of the earth might be adorned with joy and gladness and, with
the utmost delight, approach His court. Yet because the hearts of His
servants are occupied with vain imaginings, they have remained deprived
of the lights of the morning of certainty, and these life-giving tidings
have become a cause of sorrow, for their hearts are occupied with
something else and their ears are accustomed to the harshest of voices.
Glory be to God! Sweet, refreshing water flows before their faces, and
the Most Great Luminary shines from the loftiest horizon of the world,
yet the people are heedless. In every circumstance one must beseech God,
glorious is His glory, not to deprive His servants. Praise the Desire of
all the worlds, whose glad tidings bestowed life and who aided you to
turn toward Him, approach Him, and recognize Him. He is indeed the Most
Bountiful, the Generous. There is no God but Him, the All-Powerful, the
All-Knowing, the All-Wise. End.

\subsection{90}\label{section-2321}

\begin{otherlanguage}{persian}

اینعبد فانیهم خدمت هر یک سلام میرساند و تکبیر معروض میدارد و از حقّ از
برای کلّ مسئلت مینماید آنچه را که ذکرش باقی و دائمست ذکر حبیب روحانی جناب
ع ب علیه بهآء اللّه الأبهی را فرمودند و همچنین ذکر اشیاء مرسلهٴ ایشانرا
للّه الحمد رسید و چون بنیّت خالص ارسال شده بود هر یک بمقام خود فائز و ذکر
حبیب مرسل مکرّر در ساحت اقدس و انجمن طائفین مذکور جناب محبوبی امین علیه
بهآء اللّه مکرّر ذکر فرمودند و بیدمشک که فی‌الحقیقه اکلیل عطریّات عالم بود
در ساحت امنع اقدس و اهل سرادق عفّت و عصمت صرف شد و محبوبی جناب امین علیه
بهآء اللّه هم عمل نمود بآنچه که محبوب بود اینفانی بصدهزار لسان میطلبد که
به شکر و حمد الهی ناطق شود چه که ایشانرا تأیید فرمود و در ساحت اقدس
بقبول فائز گشت حقّ جلّ جلاله اولیای خود را حفظ فرماید و بأعلی ما یمکن فائز
نماید انّه علی کلّشیء قدیر فی‌الحقیقه آنمحبوب فؤاد مع فرداً فرداً آن اشیاء
امام وجه حاضر اینکلمه‌ایست که آنحضرت رجا نموده بودند للّه الحمد بنور اجابت
منوّر گشت

\end{otherlanguage}

Īn ʿabd-i fānī ham khedmat-i har yek salām mīresānad va takbīr maʿrūż
mīdārad va az Ḥaqq az barā-ye koll masʾalat mīnamāyad ānche rā ke
dhekrash bāqī va dāʾem ast. Dhekr-i ḥabīb-i rūḥānī Janāb-i ʿ B ʿalayhi
Bahāʾ Allāh al-Abhā rā farmūdand va hamchenīn dhekr-i ashyā-ye
morsale-ye īshān rā. Li-llāh al-ḥamd, rasīd va chon be-niyyat-i khāleṣ
ersāl shode būd har yek be-maqām-i khod fāʾez va dhekr-i ḥabīb-i morsel
mokarrar dar sāḥat-i aqdas va anjoman-i ṭāʾefīn madhkūr. Janāb-i maḥbūbī
Amīn ʿalayhi Bahāʾ Allāh mokarrar dhekr farmūdand va bīdm eshk ke
fī-al-ḥaqīqe eklīl-i ʿeṭrīyāt-i ʿālam būd dar sāḥat-i amnaʿ-i aqdas va
ahl-i sorādeq-i ʿeffat va ʿeṣmat ṣarf shod va maḥbūbī Janāb-i Amīn
ʿalayhi Bahāʾ Allāh ham ʿamal namūd be-ānche ke maḥbūb būd. Īn fānī
be-ṣad-hezār lesān mīṭalabad ke be-shokr va ḥamd-i elāhī nāṭeq shavad,
che ke īshān rā taʾyīd farmūd va dar sāḥat-i aqdas be-qabūl fāʾez gasht.
Ḥaqq jalla jalāluhu awliyā-ye khod rā ḥefẓ farmāyad va be-aʿlā mā yomken
fāʾez namāyad. Innahu ʿalā kulli shayʾin qadīr. Fī-al-ḥaqīqe ān maḥbūb-i
foʾād maʿa fardan fardan ān ashyā amām-i vajh ḥāżer. Īn kalame-īst ke ān
ḥażrat rajā namūde būdand. Li-llāh al-ḥamd, be-nūr-i ejābat monavvar
gasht.

This transient servant, too, sends greetings to each one, offers an
expression of God's greatness, and asks God to grant everyone that whose
remembrance is lasting and enduring. You mentioned the spiritual friend,
Jināb-i ʿ. B.---may the glory of God, the All-Glorious, rest upon
him---and also the things he sent. Praise be to God, they arrived, and
since they had been sent with pure intent, each attained its own
station, and the friend who sent them has repeatedly been remembered in
the Most Holy Court and in the gathering of those who circle around it.
My beloved Jināb-i Amīn---may the glory of God rest upon him---was
repeatedly mentioned, and the musk-willow water, which was truly the
crown of the world's fragrances, was used in the most exalted, most holy
Court and by the people of the pavilion of chastity and purity. My
beloved Jināb-i Amīn---may the glory of God rest upon him---also did
what was desired. This transient one longs to speak with a hundred
thousand tongues in gratitude and praise to God, for He aided him, and
he attained acceptance in the Most Holy Court. May God, glorious is His
glory, protect His loved ones and enable them to attain the highest that
is possible. He is indeed powerful over all things. Truly, that beloved
of the heart is present before His face together with every one of those
things. This is the word for which he had expressed his hope. Praise be
to God, it was illumined with the light of an answer.

\subsection{91}\label{section-2322}

\begin{otherlanguage}{persian}

ذکر جناب کربلائی عبدالصّمد علیه بهآء اللّه از اهل م ر فرموده بودند در ساحت
امنع اقدس عرض شد و همچنین مناجاتیکه از لسان مقصود جاری و از آنحضرت مسئلت
نمود که از قبل او بساحت اقدس معروض دارید مطابق مسئلت ایشان واقع شد چه که
در شبی از شبها این تفصیل عرض شد و همچنین مناجات قرائت گشت هذا ما نطق به
لسان العظمة فی الجواب

\end{otherlanguage}

Dhekr-i Janāb-i Karbalāʾī ʿAbd al-Ṣamad ʿalayhi Bahāʾ Allāh az ahl-i M R
farmūde būdand, dar sāḥat-i amnaʿ-i aqdas ʿarż shod va hamchenīn
monājātī ke az lesān-i Maqṣūd jārī va az ān ḥażrat masʾalat namūd ke az
qebal-i ū be-sāḥat-i aqdas maʿrūż dārīd. Moṭābeq-i masʾalat-i īshān
vāqeʿ shod, che ke dar shabī az shabhā īn tafṣīl ʿarż shod va hamchenīn
monājāt qerāʾat gasht. Hādhā mā naṭaqa bihi lisānu al-ʿaẓamati fī
al-jawāb:

You had mentioned Jināb-i Karbalāʾī ʿAbd al-Ṣamad---may the glory of God
rest upon him---of M. R. This was presented in the most exalted, most
holy Court, as was the prayer that had flowed from the tongue of the
Desired One and which he had asked you to present on his behalf in the
Most Holy Court. His request was fulfilled, for on one of the nights
these details were presented and the prayer was also recited. This is
what the Tongue of Grandeur spoke in reply:

\subsection{92}\label{section-2323}

\begin{otherlanguage}{arabic}

قوله جلّ اجلاله هو النّاطق فیملکوت البیان

\end{otherlanguage}

Qawluhu jalla ijlāluhu: Huwa al-Nāṭiq fī malakūti al-bayān.

His words, exalted is His majesty: He is the Speaker in the kingdom of
utterance.

\subsection{93}\label{section-2324}

\begin{otherlanguage}{arabic}

یا صمد علیک بهائی قد ذکرک من احبّنی و ارادنی و تمسّک بحبلی و تشبّث بذیلی و
نطق بثنائی بین عبادی و قام علی نصرة امری ذکرناک بآیات لو تلقیها علی
الجبل لتسمع منه لبّیک لبّیک یا مقصود الملأ الأعلی و یا محبوب من فی السّموات
و الأرضین یا صمد اذ کنت حاضراً لدی العرش و سمعت بأذنک ندآء ربّک و رأیت
بعینیک افق ظهور اللّه ربّ العالمین

\end{otherlanguage}

Yā Ṣamad, ʿalayka bahāʾī. Qad dhakaraka man aḥabbanī wa-arādanī
wa-tamassaka bi-ḥablī wa-tashabbatha bi-dhaylī wa-naṭaqa bi-thanāʾī
bayna ʿibādī wa-qāma ʿalā nuṣrati amrī. Dhakarnāka bi-āyātin law tulqīhā
ʿalā al-jabali la-tasmaʿa minhu labbayka labbayka yā maqṣūda al-malaʾi
al-aʿlā wa-yā maḥbūba man fī al-samāwāti wa-al-araḍīn. Yā Ṣamad, idh
kunta ḥāḍiran ladā al-ʿarshi wa-samiʿta bi-udhunika nidāʾa rabbika
wa-raʾayta bi-ʿaynayka ufuqa ẓuhūri Allāhi rabbi al-ʿālamīn:

O Ṣamad, My glory rest upon you. You have been mentioned by one who
loved Me, sought Me, held fast to My cord, clung to the hem of My robe,
voiced My praise among My servants, and arose to aid My Cause. We have
remembered you with verses which, were you to cast them upon a mountain,
would cause you to hear it cry, ``Here I am! Here I am, O Desire of the
Concourse on high and Beloved of all who are in the heavens and the
earths!'' O Ṣamad, since you were present before the throne, heard with
your own ear the call of your Lord, and saw with your own eyes the
horizon of the Revelation of God, the Lord of the worlds:

\subsection{94}\label{section-2325}

\begin{otherlanguage}{arabic}

قل الهی الهی لک الحمد بما اسمعتنی ما خلقتنی له انت الّذی یا الهی ایقظتنی
اذ کنت راقداً و اقمتنی اذ کنت قاعداً و علّمتنی اذ کنت جاهلاً و قرّبتنی اذ کنت
بعیداً و سقیتنی رحیقک المختوم و آویتنی فی جوار اسمک القیّوم اسألک بظهورات
قدرتک و شئونات عظمتک بأن تجعلنی مستقیماً علی امرک و ثابتاً علی حبّک و راسخاً
فی کلمتک انّک انت المقتدر العلیم الحکیم انتهی

\end{otherlanguage}

Qul: Ilāhī ilāhī, laka al-ḥamdu bi-mā asmaʿtanī mā khalaqtanī lahu. Anta
alladhī yā ilāhī ayqaẓtanī idh kuntu rāqidan wa-aqamtanī idh kuntu
qāʿidan wa-ʿallamtanī idh kuntu jāhilan wa-qarrabtanī idh kuntu baʿīdan
wa-saqaytanī raḥīqaka al-makhtūma wa-āwaytanī fī jiwāri ismika
al-Qayyūm. Asʾaluka bi-ẓuhūrāti qudratika wa-shuʾūnāti ʿaẓamatika bi-an
tajʿalanī mustaqīman ʿalā amrika wa-thābitan ʿalā ḥubbika wa-rāsikhan fī
kalimatika. Innaka anta al-Muqtadir al-ʿAlīm al-Ḥakīm. Intahā.

Say: My God, my God, praise be to You for enabling me to hear that for
which You created me. You are the One, my God, who awakened me when I
was asleep, raised me when I was seated, taught me when I was ignorant,
drew me near when I was far away, gave me Your sealed choice wine to
drink, and sheltered me in the refuge of Your name, the Self-Subsisting.
I ask You, by the manifestations of Your power and the expressions of
Your grandeur, to make me steadfast in Your Cause, constant in Your
love, and firmly rooted in Your Word. You are indeed the All-Powerful,
the All-Knowing, the All-Wise. End.

\subsection{95}\label{section-2326}

\begin{otherlanguage}{persian}

فی‌الحقیقه جناب صمد علیه ٩ ٦٦ {[}بهآء الله{]} را فضل تازه رسید و عنایت
احد مدد نمود هنیئاً له و لمن فاز بعنایة ربّه المشفق الکریم دستخطّ دیگر
آنمحبوب باب فرح گشود و بساط انبساط گسترد صدهزار آفرین بر فرح و سروریکه
از شطر اولیا ظاهر میشود آن فرح ملک الأفراحست و آن سرور سرور سرورهای عالم
و نور بهجتهای امم سبحان‌اللّه قلم مکث نکند فکر قطع نشود میترسم کثرت تحریر
آنمحبوب صدرا از امور لازمه باز دارد یعنی از تحریر آیات و اظهار بیّنات و
تذکّر عباد الحمد للّه آنحضرت موفّقند بر ذکر و ثنا و خدمت شاهد اینمقال قول
خود آنمحبوب که مرقوم داشتند آنچه بافتخار اولیای این ارض علیهم بهآء اللّه
بود در سه مقام که با ملاحظهٴ حکمت حضور بهمرسانیدند تلاوت شد و هر یک
بتمام وجود به شکر و ثنای حضرت مقصود جلّ فضله و عزّ عطائه شاکر و ذاکر
گردیدند و آنچه بافتخار اصفیای بلاد اخری بود فرداً فرداً سواد برداشته ارسال
داشت و سوادیهم ناتمام از مکتوب برداشته که بجمیع بلاد بعد از اتمام ارسال
شود چون در این ایّام غیر از این فانی کسی نیست که مشغول تحریر باشد الی آخر
بیانکم صدهزار فرح معادله نمینماید باین فرح که حاصل شد چه که از فضلهای
نامتناهی الهی بصر و سمع و لسان و ید و قلم و مداد آنمحبوب کلّ بخدمت مشغول
و باعلآء کلمه متوجّه از حقّ جلّ جلاله سائل و آمل که از عمل هر عضوی اعضای
عالم بر اعمال مأموره قیام نمایند و عامل شوند

\end{otherlanguage}

Fī-al-ḥaqīqe Janāb-i Ṣamad ʿalayhi 9 66 {[}Bahāʾ Allāh{]} rā fażl-i tāze
rasīd va ʿenāyat-i Aḥad madad namūd. Hanīʾan lahu wa-li-man fāza
bi-ʿināyati rabbihi al-Mushfiq al-Karīm. Dastkhaṭṭ-i dīgar-i ān maḥbūb
bāb-i faraḥ goshūd va besāṭ-i enbesāṭ gostard. Ṣad-hezār āfarīn bar
faraḥ va sorūrī ke az shaṭr-i awliyā ẓāher mīshavad. Ān faraḥ malek
al-afrāḥ ast va ān sorūr sarvar-i sorūrhā-ye ʿālam va nūr-i bahjathā-ye
omam. Subḥān Allāh, qalam makth nakonad, fekr qaṭʿ nashavad. Mītarsam
kathrat-i taḥrīr ān maḥbūb Ṣadr rā az omūr-i lāzeme bāz dārad, yaʿnī az
taḥrīr-i āyāt va eẓhār-i bayyenāt va tadhakkor-i ʿebād. Al-ḥamdu li-llāh
ān ḥażrat movaffaqand bar dhekr va thanā va khedmat. Shāhed-i īn maqāl
qowl-i khod-i ān maḥbūb ke marqūm dāshtand: Ānche be-eftek hār-i
awliyā-ye īn arż ʿalayhim Bahāʾ Allāh būd dar se maqām ke bā molāḥaẓe-ye
ḥekmat ḥożūr be-ham rasānīdand telāvat shod va har yek be-tamām-i vojūd
be-shokr va thanā-ye Ḥażrat-i Maqṣūd jalla faḍluhu wa-ʿazza ʿaṭāʾuhu
shāker va dhāker gardīdand va ānche be-eftek hār-i aṣfiyā-ye belād-i
okhrā būd fardan fardan savād bardāshte ersāl dāsht va savādī ham
nātamām az maktūb bardāshte ke be-jamīʿ-i belād baʿd az etmām ersāl
shavad chon dar īn ayyām ghayr az īn fānī kasī nīst ke mashghūl-i taḥrīr
bāshad. Ilā ākhiri bayānikum. Ṣad-hezār faraḥ moʿādele namīnamāyad be-īn
faraḥ ke ḥāṣel shod, che ke az fażlhā-ye nāmotanāhī-ye elāhī baṣar va
samʿ va lesān va yad va qalam va medād-i ān maḥbūb koll be-khedmat
mashghūl va be-eʿlā-ye kalame motavajjeh. Az Ḥaqq jalla jalāluhu sāʾel
va āmel ke az ʿamal-i har ʿożvī aʿżā-ye ʿālam bar aʿmāl-i maʾmūre qīyām
namāyand va ʿāmel shavand.

Truly, fresh bounty reached Jināb-i Ṣamad---upon him be 9 66 {[}the
glory of God{]}---and the favor of the One came to his aid. Blessed is
he, and whoever has attained the favor of his compassionate, generous
Lord. Your other letter opened the door of joy and spread a banquet of
delight. A hundred thousand blessings upon the joy and gladness that
appear from the direction of the loved ones! That joy is the king of
joys, and that gladness is the chief of the world's gladness and the
light of the nations' delights. Glory be to God! The pen does not pause,
and thought does not cease. I fear that so much writing may keep that
beloved Ṣadr from his necessary tasks: transcribing the verses, making
known the clear proofs, and reminding the servants. Praise be to God,
you are assisted in remembrance, praise, and service. Evidence of this
is your own statement, which you wrote: ``What had been revealed in
honor of the loved ones in this land---may the glory of God rest upon
them---was recited on three occasions when, with due regard for wisdom,
they gathered, and each, with his entire being, offered thanks and
praise to the Desired One---glorious is His bounty and mighty His
bestowal. As for what had been revealed in honor of the chosen ones in
other lands, I copied and sent each item separately. I have also made a
copy of the letter, still unfinished, to be sent to all the lands once
completed, since in these days there is no one but this transient one
engaged in transcribing''---and so on to the end of your statement. A
hundred thousand joys cannot equal the joy this has brought, for through
God's infinite bounties your sight, hearing, tongue, hand, pen, and ink
are all engaged in service and intent upon exalting the Word. I ask and
hope of God, glorious is His glory, that through the work of each of
these members, the members of the world may arise to perform the deeds
enjoined upon them and put them into practice.

\subsection{96}\label{section-2327}

\begin{otherlanguage}{persian}

اینکه در بارهٴ رسالهٴ استدلالیّه که از یراعة اللّه نازل شده ذکر فرمودند
معلومست آنحضرت عرف بیان رحمن را ادراک مینمایند شامّهٴ عالم مسدود است الّا
من شآء ربّنا ربّ الأرض و السّمآء و آن نفوس اربعه را مخصوص احضار فرمودند تا
لدی العرش حاضر شوند ولکن موفّق نشدند و یکی از آن اربعه جناب آقا ملّا محمّد
علیه بهآء اللّه بوده لأجل حکمت با آن نفوس بوده و آن سهٴ دیگر در ظاهر از
اهل فرقان محسوب ولکن لدی العرش بشجرهٴ اعراض نسبت داده شده اگر جمیع عالم
خالصاً عن الأغراض اصغا نمایند همان اوراق مبارکهٴ معدوده کلّ را هدایت نماید
و ببحر اعظم کشاند از حقّ میطلبم کلّ را از اغراض نفسانیّه مطهّر فرماید تا
لایق اصغا و قابل مشاهده گردند

\end{otherlanguage}

Īnke dar bāre-ye resāle-ye estedlālīye ke az Yarāʿat Allāh nāzel shode
dhekr farmūdand, maʿlūm ast ān ḥażrat ʿarf-i bayān-i Raḥmān rā edrāk
mīnamāyand. Shāmme-ye ʿālam masdūd ast illā man shāʾa rabbunā rabbu
al-arḍi wa-al-samāʾ. Va ān nofūs-i arbaʿe rā makhṣūṣ eḥżār farmūdand tā
ladā al-ʿarsh ḥāżer shavand va-līkan movaffaq nashodand va yekī az ān
arbaʿe Janāb-i Āqā Mollā Moḥammad ʿalayhi Bahāʾ Allāh būde, li-ajl-i
ḥekmat bā ān nofūs būde va ān se-ye dīgar dar ẓāher az ahl-i Forqān
maḥsūb va-līkan ladā al-ʿarsh be-shajare-ye eʿrāż nesbat dāde shode.
Agar jamīʿ-i ʿālam khāliṣan ʿan al-aghrāḍ eṣghā namāyand hamān awrāq-i
mobārake-ye maʿdūde koll rā hedāyat namāyad va be-baḥr-i aʿẓam keshānad.
Az Ḥaqq mīṭalabam koll rā az aghrāż-i nafsānīye moṭahhar farmāyad tā
lāyeq-i eṣghā va qābel-i moshāhade gardand.

Regarding your mention of the treatise of proofs revealed by the Pen of
God, it is evident that you perceive the fragrance of the utterance of
the All-Merciful. The world's sense of smell is blocked, except for
those whom our Lord, the Lord of earth and heaven, wills. He
specifically summoned those four individuals to appear before the
throne, but they did not succeed in doing so. One of those four was
Jināb-i Āqā Mullā Muḥammad---may the glory of God rest upon him---who
had been with those individuals for reasons of wisdom. The other three
are outwardly counted among the people of the Qurʾān, but before the
throne they have been associated with the tree of opposition. If the
whole world were to listen free of ulterior motives, those few blessed
pages alone would guide everyone and draw them to the Most Great Ocean.
I ask God to purify all of selfish motives, so that they may become
worthy to listen and capable of beholding.

\subsection{97}\label{section-2328}

\begin{otherlanguage}{arabic}

اللّهمّ یا الهی و سیّدی و ربّی و مقصودی قوّ قلب من اراد ذکرک و خدمة امرک ورد
فی طلاقة لسانه فی ثنائک و حفظه فی قرائة آیاتک و رجله علی القیام علی
اجرآء حدودک انّک انت القویّ القدیر

\end{otherlanguage}

Allāhumma yā ilāhī wa sayyidī wa rabbī wa maqṣūdī, qawwi qalba man arāda
dhikraka wa khidmata amrika {[}corrupted text{]} fī ṭalāqati lisānihi fī
thanāʾika wa ḥifẓihi fī qirāʾati āyātika wa rijlihi ʿalā al-qiyāmi ʿalā
ijrāʾi ḥudūdika. Innaka anta al-qawiyyu al-qadīr.

O God, my God, my Master, my Lord, and my Desire! Strengthen the heart
of whoever desires to remember You and serve Your Cause, {[}corrupted
original text{]} in the fluency of his tongue in praising You, his
memory in reciting Your verses, and his foot in arising to enforce Your
laws. You are indeed the Strong, the All-Powerful.

\subsection{98}\label{section-2329}

\begin{otherlanguage}{persian}

اینکه اینعبد را وکیل فرمودند که از قبل آنمحبوب و جناب علی قبل با علیه
بهآء اللّه الأبهی بزیارت فائز شوم در شبی از شبها قصد مقصد اقصی نموده
تلقاء وجه حاضر و بزیارت فائز و لسان عظمت بالکلمة قبول ناطق هنیئاً لحضرتکم
و له و لهذا العبد دستخطّ دیگر آنحضرت که مطلع صحیح و صریح بود در فراق
محبوب عالمیان و مشعر بر مراتب شوق و اشتیاق و جذب و بهجت و ابتهاج حالت
تحسّر حاصل و از حقّ جلّ جلاله آمل و سائل که در جمیع احوال این حالات را
باینعبد فانی عنایت فرماید له الحمد و العطآء چه که عنایت فرموده آنچه را
که شبه و مثل نداشته هم ذکر فراق آنمحبوب در قلب نار وداد مشتعل نمود و هم
شوق وصال و لقا نور ابتهاج پدید آورد له الحمد فیکلّ الأمور و کلّ الأحوال

\end{otherlanguage}

Īnkih īn ʿabd rā vakīl farmūdand kih az qibal-i ān maḥbūb va Jināb-i
ʿAlī qabl-i Bā, ʿalayhi bahāʾu Allāhi al-Abhā, bih ziyārat fāʾiz shavam,
dar shabī az shab-hā qaṣd-i maqṣad-i aqṣā namūdih, tilqāʾ-i vajh ḥāḍir
va bih ziyārat fāʾiz va lisān-i ʿaẓamat bi'l-kalimat-i qabūl nāṭiq.
Hanīʾan li-ḥaḍratikum wa lahu wa li-hādhā al-ʿabd. Dastkhaṭṭ-i dīgar-i
ān ḥaḍrat kih maṭlaʿ-i ṣaḥīḥ va ṣarīḥ būd dar firāq-i Maḥbūb-i ʿālamiyān
va mushʿir bar marātib-i shawq va ishtiyāq va jadhb va bahjat va
ibtihāj, ḥālat-i taḥassur ḥāṣil va az Ḥaqq jalla jalāluhu āmil va sāʾil
kih dar jamīʿ-i aḥvāl īn ḥālāt rā bih īn ʿabd-i fānī ʿināyat farmāyad.
Lahu al-ḥamdu wa al-ʿaṭāʾ, chih kih ʿināyat farmūdih ānchih rā kih shibh
va mithl nadāshtih. Ham dhikr-i firāq-i ān maḥbūb dar qalb nār-i vidād
mushtaʿil namūd va ham shawq-i viṣāl va liqā nūr-i ibtihāj padīd āvard.
Lahu al-ḥamdu fī kulli al-umūr wa kulli al-aḥwāl.

As you appointed this servant to make the visitation on behalf of that
beloved one and Jināb-i ʿAlī before Bā---may the all-glorious glory of
God rest upon him---one night I set out for the most distant Goal,
entered His presence, and attained the visitation. The Tongue of
Grandeur pronounced the word of acceptance. May this bring joy to you,
to him, and to this servant! Your other letter, whose opening spoke
plainly and explicitly of separation from the Beloved of the worlds and
conveyed the degrees of your longing, yearning, attraction, joy, and
delight, aroused a sense of regret. I hope and ask of God, glorious is
His glory, that He may graciously bestow these states upon this
transient servant at all times. His are praise and bounty, for He has
bestowed what has neither likeness nor equal. The mention of separation
from that beloved one kindled the fire of love in the heart, while the
longing for reunion and attainment to His presence brought forth the
light of delight. Praise be to Him in every matter and every
circumstance.

\subsection{99}\label{section-2330}

\begin{otherlanguage}{persian}

هر حرفی مدلّ بر عنایتش و هر کلمه‌ئی گواه بر فضلش ذکر حضرت اسم‌اللّه الجمال
علیه من کلّ بهآء ابهاه را فرمودند و آنچه از قلمش در ذکر و ثنای حقّ جلّ
جلاله جاری امام بحر اعظم عرض شد و بشرف اصغا فائز و چون منتهی گشت لسان
عظمت در ملکوت بیان باین آیات باهرات ناطق قول الرّبّ تعالی و تقدّس

\end{otherlanguage}

Har ḥarfī mudill bar ʿināyatash va har kalimih-ī guvāh bar faḍlash.
Dhikr-i Ḥaḍrat-i Ismu'llāh al-Jamāl, ʿalayhi min kulli bahāʾin abhāhu,
rā farmūdand va ānchih az qalamash dar dhikr va thanā-yi Ḥaqq jalla
jalāluhu jārī, imām-i Baḥr-i Aʿẓam ʿarḍ shud va bih sharaf-i iṣghā
fāʾiz, va chūn muntahī gasht lisān-i ʿaẓamat dar malakūt-i bayān bih īn
āyāt-i bāhirāt nāṭiq. Qawlu al-Rabb taʿālā wa taqaddas.

Every letter betokens His loving care, and every word bears witness to
His grace. You mentioned His Honor Ismu'llāh al-Jamāl---may the most
glorious of all glory rest upon him---and whatever had flowed from his
pen in remembrance and praise of God, glorious is His glory, was
presented before the Most Great Ocean and received the honor of being
heard. When it was finished, the Tongue of Grandeur spoke these
resplendent verses in the kingdom of utterance. These are the words of
the Lord, exalted and sanctified is He:

\subsection{100}\label{section-2331}

\begin{otherlanguage}{persian}

یا ورقآء یا ایّها الشّارب رحیقی المختوم و المتمسّک باسمی القیّوم و القائم
علی امری المحتوم بحر عنایت در این حین موّاج و عمّان سرور در اعلی مقام
امکان ظاهر چه که عرف محبّت کبری از آنچه ذکر شد در اظهار خضوع و خشوع و
استقامت و تضرّع و ابتهال اسمی جمال علیه بهائی ساطع و متضوّع یا جمال اظهار
عجز نمودی تلقاء قدرت حقّ جلّ جلاله حال قدیر ترا ذکر مینماید و قادر بتو
توجّه نموده و اظهار نیستی نمودی حقّ مطلق از شطر سجن بتو اقبال فرموده عجزت
در اینمقام فخر قدرتهای اهل ارضست و نیستیت اکلیل هستیهای امم قد ذکرک و
ذکر ذکرک من طار فی هوآء قربی و حبّی و عنایتی و ارسله الی سجنی و عرضه عبد
حاضری و اجابک لسان عظمتی رشحی طلب نمودی قد انهدر الیک بحری قطره خواستی
هطلت علیک امطار رحمتی طوبی لمن ذکرک در نیستی و مقام آن تفکّر نما که مطلع
این عنایت کبری شده و در اظهار فقر و ثمراتش ناظر باش چه که ملکوت غنا را
بتو متوجّه نموده هم ذکر حبّ تو من احبّی را و ذکر او ترا از این دو ذکر نوری
در عوالم محبّت و وداد و الفت و اتّحاد ساطع هنیئاً لکما و مریئاً لکما از حقّ
میطلبیم از این نور آفاق قلوب اولیا را منوّر فرماید تا کلّ بحبل اعظم حبّ
الهی تمسّک نمایند اینست کوثر حیوان که از قلم حضرت رحمن جاری و ساریست
البهآء علیکما الرّحمة علیکما العنایة علیکما الفضل علیکما

\end{otherlanguage}

Yā Warqāʾ, yā ayyuhā al-shāribu raḥīqī al-makhtūm wa al-mutamassiku
bi-smī al-Qayyūm wa al-qāʾimu ʿalā amrī al-maḥtūm. Baḥr-i ʿināyat dar īn
ḥīn mavvāj va ʿummān-i surūr dar aʿlā maqām-i imkān ẓāhir, chih kih
ʿarf-i maḥabbat-i kubrā az ānchih dhikr shud dar iẓhār-i khuḍūʿ va
khushūʿ va istiqāmat va taḍarruʿ va ibtihāl-i Ismī Jamāl, ʿalayhi
bahāʾī, sāṭiʿ va mutaḍawwiʿ. Yā Jamāl, iẓhār-i ʿajz namūdī tilqāʾ-i
qudrat-i Ḥaqq jalla jalāluhu; ḥāl Qadīr turā dhikr mīnamāyad va Qādir
bih tu tavajjuh namūdih. Va iẓhār-i nīstī namūdī; Ḥaqq-i Muṭlaq az
shaṭr-i sijn bih tu iqbāl farmūdih. ʿAjzat dar īn maqām fakhr-i
qudrat-hā-yi ahl-i arḍ ast va nīstīyat iklīl-i hastī-hā-yi umam. Qad
dhakaraka wa dhakara dhikraka man ṭāra fī hawāʾi qurbī wa ḥubbī wa
ʿināyatī wa arsalahu ilā sijnī wa ʿaraḍahu ʿabdun ḥāḍirī wa ajābaka
lisānu ʿaẓamatī. Rashḥī ṭalab namūdī; qad inhadara ilayka baḥrī. Qaṭrih
khvāstī; haṭalat ʿalayka amṭāru raḥmatī. Ṭūbā li-man dhakarak. Dar nīstī
va maqām-i ān tafakkur namā kih maṭlaʿ-i īn ʿināyat-i kubrā shudih va
dar iẓhār-i faqr va thamarātash nāẓir bāsh, chih kih malakūt-i ghinā rā
bih tu mutavajjih namūdih. Ham dhikr-i ḥubb-i tu man aḥabba rā va
dhikr-i ū turā, az īn du dhikr nūrī dar ʿavālim-i maḥabbat va vidād va
ulfat va ittiḥād sāṭiʿ. Hanīʾan lakumā wa marīʾan lakumā. Az Ḥaqq
mīṭalabīm az īn nūr āfāq-i qulūb-i awliyā rā munavvar farmāyad tā kull
bih ḥabl-i aʿẓam-i ḥubb-i ilāhī tamassuk namāyand. Īn ast Kawthar-i
ḥayavān kih az qalam-i Ḥaḍrat-i Raḥmān jārī va sārīst. Al-bahāʾu
ʿalaykumā, al-raḥmatu ʿalaykumā, al-ʿināyatu ʿalaykumā, al-faḍlu
ʿalaykumā.

``O Warqāʾ, you who drink My sealed wine, hold fast to My name, the
Self-Subsisting, and uphold My irrevocable Cause! The ocean of loving
care is now surging, and the sea of joy is manifest at the highest
station of the contingent world, for the fragrance of the most great
love rises and spreads from what was related of the submissiveness,
humility, steadfastness, supplication, and entreaty expressed by My name
Jamāl---upon him be My glory. O Jamāl! You expressed your powerlessness
before the might of God, glorious is His glory. Now the All-Powerful
makes mention of you, and the Almighty has turned toward you. You
expressed your nothingness; the absolute Truth has turned toward you
from the direction of the prison. In this station, your powerlessness is
the pride of the powers of the people of the earth, and your nothingness
is the crown of the existence of nations. He who soared in the
atmosphere of My nearness, My love, and My loving care mentioned you and
your words of remembrance and sent them to My prison. A servant
attending Me presented them, and the Tongue of My Grandeur answered you.
You asked for a sprinkling; My ocean has poured down upon you. You
sought a drop; the rains of My mercy have showered upon you. Blessed is
he who mentioned you! Reflect on nothingness and its station, for it has
become the dawning-place of this most great loving care. Consider the
expression of poverty and its fruits, for it has turned the kingdom of
wealth toward you. From both your mention of your love for him who loved
and his mention of you, a light shines in the worlds of love, affection,
fellowship, and unity. May it bring joy and delight to you both! We ask
God to illumine the horizons of the hearts of His loved ones with this
light, so that all may hold fast to the most great cord of divine love.
This is the Kawthar of life that flows and streams from the Pen of the
All-Merciful. Glory be upon you both; mercy be upon you both; loving
care be upon you both; grace be upon you both.''

\subsection{101}\label{section-2332}

\begin{otherlanguage}{persian}

انتهی مقام اتّفاق و اتّحاد مخصوص نفوسیکه از ایادی امر محسوبند لدی اللّه
مرغوب و محبوب بلکه بعد از قرائت نامهٴ حضرت اسم‌اللّه علیه بهآء اللّه الأبهی
که بآنمحبوب نوشته بودند و همچنین ذکر آنمحبوب آن نامه را ابتسم ثغر
المحبوب از آنچه نازل شد در اینمقام گواهیست امین و شاهدیست مشهود بر علوّ
مقام اتّفاق و اتّحاد عند اللّه مالک الایجاد اینعبد فانی خدمت حضرت ایشان
سلام و تکبیر عرض مینماید لازال از غنیّ متعال مسئلت نموده‌ام آنچه را که سبب
علوّ و سموّ اولیای اوست انّ ربّی هو السّامع المجیب

\end{otherlanguage}

Intahā. Maqām-i ittifāq va ittiḥād, makhṣūṣ-i nufūsī kih az Ayādī-yi Amr
maḥsūband, ladā Allāh marghūb va maḥbūb. Balkih baʿd az qirāʾat-i
nāmih-yi Ḥaḍrat-i Ismu'llāh, ʿalayhi bahāʾu Allāhi al-Abhā, kih bih ān
maḥbūb nivishtih būdand va hamchunīn dhikr-i ān maḥbūb ān nāmih rā,
ibtasama thaghru al-Maḥbūb. Az ānchih nāzil shud dar īn maqām guvāhīst
amīn va shāhidīst mashhūd bar ʿuluvv-i maqām-i ittifāq va ittiḥād ʿinda
Allāh Mālik al-ījād. Īn ʿabd-i fānī khidmat-i ḥaḍrat-i īshān salām va
takbīr ʿarḍ mīnamāyad. Lā yazāl az Ghanī-yi Mutaʿāl masʾalat namūdih-am
ānchih rā kih sabab-i ʿuluvv va sumuvv-i awliyā-yi ūst. Inna rabbī huwa
al-sāmiʿu al-mujīb.

Here it ends. The station of concord and unity, especially among those
souls counted as Hands of the Cause, is desired and beloved in the sight
of God. Indeed, after the letter that His Honor Ismu'llāh---may the
all-glorious glory of God rest upon him---had written to that beloved
one was read, together with that beloved one's mention of the letter,
the Beloved's lips smiled. What was revealed on this occasion is a
trustworthy testimony and a manifest witness to the loftiness of the
station of concord and unity before God, the Master of creation. This
transient servant offers His Honor greetings and praise. I have
continually asked the All-Sufficing, the Most Exalted, for whatever will
bring about the elevation and exaltation of His loved ones. My Lord is
indeed the All-Hearing, the Answerer.

\subsection{102}\label{section-2333}

\begin{otherlanguage}{persian}

ذکر سمیّ حضرت مقصود ابن من صعد الی اللّه آقا میرزا محسن طبیب علیهما بهآء
اللّه و اخوی ایشان جناب آقا میرزا جعفر علیه بهآء اللّه را نمودند بعد از
عرض امام عرش فرمودند قول الرّبّ تعالی و تقدّس

\end{otherlanguage}

Dhikr-i samiyy-i Ḥaḍrat-i Maqṣūd, ibn-i man ṣaʿida ilā Allāh, Āqā Mīrzā
Muḥsin-i Ṭabīb, ʿalayhimā bahāʾu Allāh, va akhavī-yi īshān Jināb-i Āqā
Mīrzā Jaʿfar, ʿalayhi bahāʾu Allāh, rā namūdand. Baʿd az ʿarḍ imām-i
ʿarsh farmūdand, qawlu al-Rabb taʿālā wa taqaddas.

You mentioned the namesake of the Desired One, the son of Āqā Mīrzā
Muḥsin the physician, who ascended to God---may the glory of God rest
upon them both---and his brother, Jināb-i Āqā Mīrzā Jaʿfar---may the
glory of God rest upon him. After this was presented before the throne,
He spoke these words, the words of the Lord, exalted and sanctified is
He:

\subsection{103}\label{section-2334}

\begin{otherlanguage}{persian}

علیک بهائی و عنایتی انّا نذکر ابن من فاز بحبّی و عاش بحبّی و صعد بحبّی از
قبل مظلوم باو تکبیر برسان و بگو یا سمیّ مقصود للّه الحمد نیّر فضل از افق
سماء عطا اشراق نموده اشراقیکه عالم را احاطه کرده از حقّ میطلبیم ترا تأیید
فرماید بر اعمالیکه هر ذی شمّی عرف رضای الهی از او استشمام نماید باب فضل
مفتوح و کرسیّ بیان در قطب امکان مشهود و خطیب الهی بر آن مستوی در لیالی و
ایّام بأعلی النّدآء بآیاتیکه سبب عزّ وجود و نصایحیکه علّت علوّ بوده ندا فرمود
و وصیّت نمود یا ورقا غبار اوهام بعضی را اخذ نموده انّ ربّک یعلم خائنة
الأعین و خافیة الصّدور

\end{otherlanguage}

ʿAlayka bahāʾī wa ʿināyatī. Innā nadhkuru ibna man fāza bi-ḥubbī wa
ʿāsha bi-ḥubbī wa ṣaʿida bi-ḥubbī. Az qibal-i Maẓlūm bih ū takbīr
birasān va bigū yā samiyy-i Maqṣūd, li'llāh al-ḥamd, nayyir-i faḍl az
ufuq-i samā-yi ʿaṭā ishrāq namūdih, ishrāqī kih ʿālam rā iḥāṭih kardih.
Az Ḥaqq mīṭalabīm turā taʾyīd farmāyad bar aʿmālī kih har dhī shammī
ʿarf-i riḍā-yi ilāhī az ū istishmām namāyad. Bāb-i faḍl maftūḥ va
kursī-yi bayān dar quṭb-i imkān mashhūd va khaṭīb-i ilāhī bar ān
mustavī. Dar layālī va ayyām bi-aʿlā al-nidāʾ bih āyātī kih sabab-i
ʿizz-i vujūd va naṣāyiḥī kih ʿillat-i ʿuluvv būdih nidā farmūd va
vaṣiyyat namūd. Yā Warqā, ghubār-i awhām baʿḍī rā akhdh namūdih. Inna
rabbaka yaʿlamu khāʾinata al-aʿyun wa khāfiyata al-ṣudūr.

``Upon you be My glory and My loving care. We remember the son of him
who attained My love, lived in My love, and ascended in My love. Convey
greetings to him on behalf of the Wronged One and say: O namesake of the
Desired One! Praise be to God, the luminary of grace has shone from the
horizon of the heaven of bounty with a radiance that has encompassed the
world. We ask God to confirm you in deeds from which everyone endowed
with the sense of smell may inhale the fragrance of divine good
pleasure. The door of grace is open, the throne of utterance is visible
at the center of the contingent world, and the divine Preacher is seated
upon it. By night and day He has called out in the loudest voice with
verses that bring honor to existence and counsels that lead to
exaltation, and has offered His exhortations. O Warqā! The dust of vain
imaginings has enveloped some. Your Lord indeed knows the treachery of
the eyes and what lies hidden in the hearts.''

\subsection{104}\label{section-2335}

\begin{otherlanguage}{arabic}

میرزا جعفر را ذکر مینمائیم فضلاً من عندنا و انا الفضّال الکریم

\end{otherlanguage}

Mīrzā Jaʿfar rā dhikr mīnamāyīm faḍlan min ʿindinā wa anā al-Faḍḍālu
al-Karīm.

``We mention Mīrzā Jaʿfar as a token of grace from Us, and I am the Most
Bountiful, the Generous.''

\subsection{105}\label{section-2336}

\begin{otherlanguage}{persian}

یا جعفر خذ ما نزّل فی الکتاب و ضِع ما عند الأحزاب بعضی از نفوس موهومه
باغوی مشغول و برخی بمفتریات ناطق و حزبی به اوهام و ظنون متمسّک حقّ مقدّس از
ادراک و عرفان و امرش اظهر از آفتاب از حقّ میطلبیم سمع و بصر را محروم
ننماید و کوثر بیان را از من علی الأرض منع نکند باید اولیای الهی به اخلاق
و آداب و اعمال ناس را هدایت کنند و ببحر اعظم راه نمایند انتهی

\end{otherlanguage}

Yā Jaʿfar, khudh mā nuzzila fī al-kitāb wa ḍiʿ mā ʿinda al-aḥzāb. Baʿḍī
az nufūs-i mawhūmih bih ighvā mashghūl va barkhī bih muftarayāt nāṭiq va
ḥizbī bih awhām va ẓunūn mutamassik. Ḥaqq muqaddas az idrāk va ʿirfān va
amrash aẓhar az āftāb. Az Ḥaqq mīṭalabīm samʿ va baṣar rā maḥrūm
nanamāyad va Kawthar-i bayān rā az man ʿalā al-arḍ manʿ nakunad. Bāyad
awliyā-yi ilāhī bih akhlāq va ādāb va aʿmāl nās rā hidāyat kunand va bih
Baḥr-i Aʿẓam rāh namāyand. Intahā.

``O Jaʿfar! Take what has been revealed in the Book and leave aside what
the factions possess. Some deluded souls are occupied with leading
others astray, some voice fabrications, and a group clings to vain
imaginings and conjectures. God is sanctified beyond comprehension and
recognition, and His Cause is more manifest than the sun. We ask God not
to deprive hearing and sight, and not to withhold the Kawthar of
utterance from those on earth. The loved ones of God must guide people
through their character, conduct, and deeds, and show them the way to
the Most Great Ocean.'' Here it ends.

\subsection{106}\label{section-2337}

\begin{otherlanguage}{arabic}

هر ذی بصر و سمع و فؤادی متحیّر که چه سکری عباد را اخذ نموده که مابین
آفتاب و کرم شب‌تاب فرق نمیگذارند و تمیز نمیدهند وقتی از اوقات اینکلمهٴ
علیا از مطلع نور بقا اشراق نمود قول الرّبّ تعالی و تقدّس

\end{otherlanguage}

Har dhī baṣar va samʿ va fuʾādī mutaḥayyir kih chih sukrī ʿibād rā akhdh
namūdih kih mābayn-i āftāb va kirm-i shab-tāb farq namīgudhārand va
tamīz namīdihand. Vaqtī az awqāt īn kalimih-yi ʿulyā az maṭlaʿ-i nūr-i
baqā ishrāq namūd. Qawlu al-Rabb taʿālā wa taqaddas.

Everyone endowed with sight, hearing, and a heart is bewildered as to
what intoxication has seized the servants, that they neither distinguish
nor discern between the sun and a glowworm. On one occasion, this
exalted utterance shone from the dawning-place of the light of eternity.
These are the words of the Lord, exalted and sanctified is He:

\subsection{107}\label{section-2338}

\begin{otherlanguage}{persian}

یا عبد حاضر الی حین نمیتوان حقّ مبین را ذکر نمود مع صریر قلم اعلی و نداء
مالک اسماء سمع و بصر حقیقی مشاهده نشد ناقابل مشاهده و نالایق اصغا گردد
در حزب بیان ملاحظه کن مع آنکه حزب شیعه ضلالت و گمراهی و خسارتشان واضح و
آشکار مجدّد بترتیب چنین حزبی مشغولند لعمراللّه آنحزب رائحهٴ توحید را
نیافته‌اند عبده اصنام و اوهام بوده و هستند حقّ الی حین باندازه نطق فرموده
چه که عباد ضعیفند در قرون و اعصار باوهام تربیت شده‌اند عیون مرموده لایق
مشاهدهٴ انوار آفتاب نه نفوس موهومه قابل تجلّیات نیّر یقین نه مبدئهم الظّنون
و منتهاهم الاوهام کذلک نطق لسان دیک العزیز العلّام انتهی ذکر جناب آقا سیّد
فرج‌اللّه علیه ٩ ٦٦ {[}بهآء الله{]} از قلم آنمحبوب جاری بعد از عرض لدی
العرش نیّر این بیان از افق سماء برهان مشرق

\end{otherlanguage}

Yā ʿabd-i ḥāḍir, ilā ḥīn namītavān Ḥaqq-i Mubīn rā dhikr namūd. Maʿa
ṣarīr-i Qalam-i Aʿlā va nidā-yi Mālik-i Asmā samʿ va baṣar-i ḥaqīqī
mushāhadih nashud. Nāqābil-i mushāhadih va nālāyiq-i iṣghā gardad. Dar
ḥizb-i Bayān mulāḥiẓih kun; maʿ ānkih ḥizb-i Shīʿih ḍalālat va gumrāhī
va khasāratishān vāḍiḥ va āshkār, mujaddad bih tartīb-i chunīn ḥizbī
mashghūland. La-ʿamru'llāh, ān ḥizb rāʾiḥih-yi tawḥīd rā nayāftih-and.
ʿAbadih-yi aṣnām va awhām būdih va hastand. Ḥaqq ilā ḥīn bih andāzih
nuṭq farmūdih, chih kih ʿibād ḍaʿīfand. Dar qurūn va aʿṣār bih awhām
tarbiyat shudih-and. ʿUyūn-i marmūdih lāyiq-i mushāhadih-yi anvār-i
āftāb nah; nufūs-i mawhūmih qābil-i tajalliyāt-i nayyir-i yaqīn nah.
Mabdaʾuhum al-ẓunūn wa muntahāhum al-awhām. Kadhālika naṭaqa lisānu
{[}corrupted text{]} al-ʿAzīz al-ʿAllām. Intahā. Dhikr-i Jināb-i Āqā
Sayyid Faraju'llāh, ʿalayhi 9 66 {[}bahāʾu Allāh{]}, az qalam-i ān
maḥbūb jārī. Baʿd az ʿarḍ ladā al-ʿarsh nayyir-i īn bayān az ufuq-i
samā-yi burhān mushriq.

``O servant in attendance! Even now it is not possible to speak of the
manifest Truth. Despite the sound of the Most Exalted Pen and the call
of the Master of Names, no true hearing or sight has been seen. He is
deemed unworthy to be beheld and undeserving to be heard. Consider the
people of the Bayān: although the error, waywardness, and loss of the
Shīʿih community are plain and evident, they are once again engaged in
forming such a community. By the life of God, that community has not
perceived the fragrance of divine unity. They have been, and remain,
worshippers of idols and vain imaginings. Until now, God has spoken with
restraint, for the servants are weak. Over centuries and ages they have
been brought up on vain imaginings. Inflamed eyes are unfit to behold
the sun's rays; deluded souls are incapable of receiving the effulgences
of the luminary of certitude. Their beginning is conjecture, and their
end is vain imaginings. Thus has spoken the tongue of {[}corrupted
original text{]}, the Almighty, the All-Knowing.'' Here it ends. The
mention of Jināb-i Āqā Sayyid Faraju'llāh---upon him be 9 66 {[}the
glory of God{]}---flowed from that beloved one's pen. After it was
presented before the throne, the luminary of this utterance shone from
the horizon of the heaven of proof:

\subsection{108}\label{section-2339}

\begin{otherlanguage}{persian}

یا فرج علیک بهائی ذکرت مذکور و اسمت مسطور دیدیم و شنیدیم و باین بیان در
ملکوت برهان نطق فرمودیم این ایّام مکرّر ذکرت از قلم اعلی جاری اشکر ربّک از
حقّ بطلب ترا تأیید فرماید بر ذکر و ثنا و ترویج امر انّه علی کلّشیء قدیر اهل
آن ارض را از قبل مظلوم تکبیر برسان و بما ینبغی لأیّام اللّه متذکّر دار
انتهی

\end{otherlanguage}

Yā Faraj, ʿalayka bahāʾī. Dhikrat madhkūr va ismat masṭūr. Dīdīm va
shinīdīm va bih īn bayān dar malakūt-i burhān nuṭq farmūdīm. Īn ayyām
mukarrar dhikrat az Qalam-i Aʿlā jārī. Ushkur rabbak. Az Ḥaqq biṭalab
turā taʾyīd farmāyad bar dhikr va thanā va tarvīj-i Amr. Innahu ʿalā
kulli shayʾin qadīr. Ahl-i ān arḍ rā az qibal-i Maẓlūm takbīr birasān va
bi-mā yanbaghī li-ayyāmi Allāh mutadhakkir dār. Intahā.

``O Faraj! Upon you be My glory. Your mention is remembered and your
name recorded. We saw and heard, and spoke this utterance in the kingdom
of proof. In these days your mention has repeatedly flowed from the Most
Exalted Pen. Thank your Lord. Ask God to confirm you in remembrance,
praise, and the promotion of the Cause. He indeed has power over all
things. Convey greetings to the people of that land on behalf of the
Wronged One, and remind them of what befits the days of God.'' Here it
ends.

\subsection{109}\label{section-2340}

\begin{otherlanguage}{persian}

و همچنین ذکر جناب آقا میرزا آقا ش علیه بهآء اللّه نمودند بعد از عرض این
آیات نازل قوله تبارک و تعالی شماتت اعدا و ظلم اشقیا و بلایای دنیا کلّ را
لأجل خلاص عالم قبول نمودیم و حمل فرمودیم ولکن کلّ غافل الّا من شآء اللّه
انّک احمد ربّک بما ایّدک و علّمک و عرّفک و سقّاک کوثر العرفان من ایادی الفضل و
العطآء نسأله بأن یوفّقک علی الاستقامة علی امره المبرم المتین انتهی

\end{otherlanguage}

Va hamchunīn dhikr-i Jināb-i Āqā Mīrzā Āqā Sh, ʿalayhi bahāʾu Allāh,
namūdand. Baʿd az ʿarḍ īn āyāt nāzil, qawluhu tabāraka wa taʿālā.
Shamātat-i aʿdā va ẓulm-i ashqiyā va balāyā-yi dunyā, kull rā li-ajl-i
khalāṣ-i ʿālam qabūl namūdīm va ḥaml farmūdīm, va lākin kull ghāfil illā
man shāʾa Allāh. Innaka iḥmad rabbaka bi-mā ayyadaka wa ʿallamaka wa
ʿarrafaka wa saqāka Kawthara al-ʿirfān min ayādī al-faḍli wa al-ʿaṭāʾ.
Nasʾaluhu bi-an yuwaffiqaka ʿalā al-istiqāmati ʿalā amrihi al-mubrami
al-matīn. Intahā.

You also mentioned Jināb-i Āqā Mīrzā Āqā Sh---may the glory of God rest
upon him. After this was presented, these verses were revealed. These
are His words, blessed and exalted is He: ``The gloating of enemies, the
oppression of the wicked, and the afflictions of the world: We accepted
and bore them all for the deliverance of the world. Yet all are heedless
except those whom God wills. Praise your Lord for having confirmed you,
taught you, granted you recognition, and given you to drink of the
Kawthar of recognition from the hands of grace and bounty. We ask Him to
enable you to remain steadfast in His irrevocable, firm Cause.'' Here it
ends.

\subsection{110}\label{section-2341}

\begin{otherlanguage}{persian}

از توجّه جناب آقا میرزا مجید علیه بهآء اللّه بآنجهت و همچنین استقامت و
اشتعال ایشان ذکر فرمودند بعد از عرض در ساحت اقدس اینکلمات نازل قوله
تبارک و تعالی

\end{otherlanguage}

Az tavajjuh-i Jināb-i Āqā Mīrzā Majīd, ʿalayhi bahāʾu Allāh, bih ān
jihat va hamchunīn istiqāmat va ishtiʿāl-i īshān dhikr farmūdand. Baʿd
az ʿarḍ dar sāḥat-i aqdas īn kalimāt nāzil, qawluhu tabāraka wa taʿālā.

You mentioned Jināb-i Āqā Mīrzā Majīd's journey in that direction---may
the glory of God rest upon him---as well as his steadfastness and ardor.
After this was presented in the Most Holy Court, these words were
revealed. These are His words, blessed and exalted is He:

\subsection{111}\label{section-2342}

\begin{otherlanguage}{persian}

یا مجید عالم را اوهام اخذ نموده و اسباب نزاع و فساد از هر شطری مشهود در
این قرن عالم جرائد را رونقی پدیدار و صاحبان آن آنرا علّت عمار عالم و
تربیت امم و ترقّی و تهذیب و دانائی و آگاهی عباد میدانند فی‌الحقیقه اسباب
حشر و نشر است یوم تحدّث الأرض اخبارها مشهود و مسموع و الأرض جمیعاً قبضته
یوم القیامة ظاهر و آشکار ارض در سرّ حدیثهای بدیعهٴ منیعه دارد و باخبارهای
جدیده و عتیقه ناطق و در ظاهر هم طیر اخبار بکلّ جهات طائر ولکن اسف در
اینکه بعضی از جرائد اخبارات مجعوله را هم من‌غیر تحقیق نشر مینمایند هر
نفسی آنچه بگوید میگویند و آنچه بنگارد مینگارند چنانچه مطالع اوهام بجرائد
اوهامات خود را نشر مینمایند یا مجید در یوم تأتی السّمآء بدخان مبین تفکّر
نما این ایّام کذب را رونق عظیمی است غبار اوهام انام را از مالک ایّام منع
نموده از حقّ میطلبیم نفوسی برانگیزاند بحقّ نطق نمایند و سد و صدّ عباد
ایشانرا از یوم معاد محروم نسازد طوبی لمن سمع ندائی و قام علی خدمتی و
خدمة اولیائی و نطق بثنائی الجمیل الحمد للّه الفرد الواحد العزیز الجمیل
انتهی

\end{otherlanguage}

Yā Majīd, ʿālam rā awhām akhdh namūdeh va asbāb-i nizāʿ va fasād az har
shaṭrī mashhūd. Dar īn qarn-i ʿālam jarāʾid rā rawnaqī padīdār va
ṣāḥibān-i ān ān rā ʿillat-i ʿimār-i ʿālam va tarbiyat-i umam va taraqqī
va tahdhīb va dānāʾī va āgāhī-yi ʿibād mīdānand. Fī-l-ḥaqīqeh asbāb-i
ḥashr va nashr ast. Yawma tuḥaddithu al-arḍu akhbārahā mashhūd va
masmūʿ, va-l-arḍu jamīʿan qabḍatuhu yawma al-qiyāmah ẓāhir va āshkār.
Arḍ dar sirr ḥadīthhā-yi badīʿeh-yi manīʿeh dārad va bi-akhbārhā-yi
jadīdeh va ʿatīqeh nāṭiq va dar ẓāhir ham ṭayr-i akhbār bi-kull-i jihāt
ṭāʾir. Va-lākin asaf dar īnkih baʿḍī az jarāʾid akhbārāt-i majʿūleh rā
ham min ghayr-i taḥqīq nashr mīnamāyand. Har nafsī ānchih bigūyad
mīgūyand va ānchih bingārad mīngārand, chunānchih maṭāliʿ-i awhām
bi-jarāʾid awhāmāt-i khud rā nashr mīnamāyand. Yā Majīd, dar yawma taʾtī
al-samāʾu bi-dukhānin mubīn tafakkur namā. Īn ayyām kidhb rā rawnaq-i
ʿaẓīmī ast. Ghubār-i awhām anām rā az Mālik-i Ayyām manʿ namūdeh. Az
Ḥaqq mīṭalabīm nufūsī bar-angīzānad bi-ḥaqq nuṭq namāyand va sadd va
ṣadd-i ʿibād īshān rā az Yawm-i Maʿād maḥrūm nasāzad. Ṭūbā li-man samiʿa
nidāʾī wa qāma ʿalā khidmatī wa khidmati awliyāʾī wa naṭaqa bi-thanāʾī
al-jamīl. Al-ḥamdu li-llāhi al-Fardi al-Wāḥidi al-ʿAzīzi al-Jamīl.
Intahā.

O Majīd, delusions have seized the world, and causes of conflict and
disorder are evident on every side. In this century, newspapers have
gained prominence in the world, and their proprietors regard them as a
means of building up the world, educating the nations, and advancing,
refining, informing, and enlightening its people. They are, in truth,
instruments of gathering and dissemination. ``The day when the earth
recounts its tidings'' is seen and heard, and ``the entire earth will be
in His grasp on the Day of Resurrection'' is plain and manifest. In
secret, the earth has wondrous and sublime accounts to tell and speaks
of tidings new and old; outwardly, too, the bird of news flies in every
direction. Yet it is regrettable that some newspapers also publish
fabricated reports without investigation. Whatever anyone says, they
repeat, and whatever anyone writes, they reproduce. Thus, those who are
the sources of delusions disseminate their imaginings through
newspapers. O Majīd, reflect on ``the day when the sky brings forth a
manifest smoke.'' In these days, falsehood flourishes greatly. The dust
of delusions has kept the people from the Lord of Days. We ask God to
raise up souls who will speak the truth and whom the barriers and
opposition of the people will not deprive of the Day of Return. Blessed
is anyone who has heard My call, arisen to serve Me and My loved ones,
and voiced My beautiful praise. Praise be to God, the Unique, the One,
the Almighty, the All-Beauteous. Here it ends.

\subsection{112}\label{section-2343}

\begin{otherlanguage}{persian}

فی‌الحقیقه جراید گاهی طیّ الأرض مینمایند و هنگامی بقوادم اخبار در طیران و
وقتی تغییر اسم مینمایند و سیّالهٴ برقیّه میشود تتطایر منها الأذکار و
الأخبار و الآثار ولکن از حقّ جلّ جلاله مسئلت مینمایم صاحبان آنرا تأیید
فرماید تا کلّ بآنچه سزاوار است عمل نمایند گمان نمیرفت که بعضی از جرائد
اروپا من‌غیر تحقیق ذکری نماید و یا تابع گفتهای مجعولهٴ عباد شود مظاهر
اوهام در لیالی و ایّام بجراید مکتوب میفرستند و من‌غیر تحقیق طبع میشود چون
بعضی طالب اخبار تازه هستند طبع مینمایند از حقّ میطلبم کلّ را تأیید فرماید
علی العمل بقوله تعالی لا تصدّقوا کلّ قائل و لا تطمئنّوا من کلّ وارد مکرّر از
هر محل بعضی از نفوس اراده نموده‌اند شرحی نوشته و ارسال شود تا در کتب و
جرائد نشر نمایند مقبول نیفتاد الی حین کلمه داده نشد الأمر بیده یفعل ما
یشآء و یحکم ما یرید و هو العزیز الحمید

\end{otherlanguage}

Fī-l-ḥaqīqeh jarāyid gāhī ṭayy-i arḍ mīnamāyand va hangāmī bi-qawādim-i
akhbār dar ṭayarān va vaqtī taghyīr-i ism mīnamāyand va sayyāleh-yi
barqiyyeh mīshavad. Tataṭāyaru minhā al-adhkāru wa-l-akhbāru wa-l-āthār.
Va-lākin az Ḥaqq jalla jalāluhu masʾalat mīnamāyam ṣāḥibān-i ān rā
taʾyīd farmāyad tā kull bi-ānchih sazāvār ast ʿamal namāyand. Gumān
namīraft kih baʿḍī az jarāʾid-i Urūpā min ghayr-i taḥqīq dhikrī namāyad
va yā tābiʿ-i guftahā-yi majʿūleh-yi ʿibād shavad. Maẓāhir-i awhām dar
layālī va ayyām bi-jarāyid maktūb mīfiristand va min ghayr-i taḥqīq ṭabʿ
mīshavad. Chūn baʿḍī ṭālib-i akhbār-i tāzeh hastand ṭabʿ mīnamāyand. Az
Ḥaqq mīṭalabam kull rā taʾyīd farmāyad ʿalā al-ʿamali bi-qawlihi taʿālā
lā tuṣaddiqū kulla qāʾilin wa lā taṭmaʾinnū min kulli wārid. Mukarrar az
har maḥall baʿḍī az nufūs irādeh namūdeh-and sharḥī nivishteh va irsāl
shavad tā dar kutub va jarāʾid nashr namāyand. Maqbūl nayuftād. Ilā ḥīn
kalameh dādeh nashud. Al-amru bi-yadihi yafʿalu mā yashāʾu wa yaḥkumu mā
yurīdu wa huwa al-ʿAzīzu al-Ḥamīd.

Indeed, newspapers sometimes traverse the earth, sometimes fly on the
pinions of news, and sometimes change their name and become an electric
current. From it, accounts, reports, and writings fly forth. Yet I
beseech God, exalted be His glory, to assist their proprietors, so that
all may act as is fitting. It was not expected that some European
newspapers would make statements without investigation or follow the
fabricated accounts of individuals. Day and night, the embodiments of
delusion send letters to newspapers, and these are printed without
investigation. Since some are eager for fresh news, they print them. I
ask God to assist everyone to act according to His exalted words: ``Do
not believe every speaker, nor place your trust in every report that
arrives.'' Repeatedly, certain individuals in various places have wished
that an account be written and sent so that they might publish it in
books and newspapers. This was not accepted. To this day, no word has
been given. The matter is in His hands. He does what He wills and
ordains what He desires, and He is the Almighty, the Praiseworthy.

\subsection{113}\label{section-2344}

\begin{otherlanguage}{persian}

در بارهٴ رضائین مرقوم داشتند بعد از عرض یک لوح امنع اقدس مخصوص جناب آقا
محمّد رضای میلانی علیه بهآء اللّه از سماء عنایت الهی نازل و ارسال شد
فرمودند هذا افضل من الاذن و انا الفضّال الکریم مخصوص جناب آقا میرزا حسن
علیه بهآء اللّه که مرقوم داشتند یک لوح امنع اقدس نازل و ارسال شد از حقّ
میطلبم ایشان را موفّق فرماید بر حفظ این عنایت کبری انّه هو المجیب

\end{otherlanguage}

Dar bāreh-yi Riḍāyayn marqūm dāshtand. Baʿd az ʿarḍ, yik lawḥ-i amnaʿ-i
aqdas makhṣūṣ-i Jināb-i Āqā Muḥammad-Riḍā-yi Mīlānī ʿalayhi Bahāʾuʾllāh
az samāʾ-i ʿināyat-i ilāhī nāzil va irsāl shud. Farmūdand hādhā afḍalu
mina al-idhni wa anā al-Faḍḍālu al-Karīm. Makhṣūṣ-i Jināb-i Āqā Mīrzā
Ḥasan ʿalayhi Bahāʾuʾllāh kih marqūm dāshtand yik lawḥ-i amnaʿ-i aqdas
nāzil va irsāl shud. Az Ḥaqq mīṭalabam īshān rā muvaffaq farmāyad bar
ḥifẓ-i īn ʿināyat-i kubrā. Innahu huwa al-Mujīb.

You wrote concerning the two Riḍās. After this was presented, a most
exalted, most holy Tablet was revealed from the heaven of divine favor
specifically for Jināb-i Āqā Muḥammad-Riḍā Mīlānī---may the glory of God
rest upon him---and was sent. He said: ``This is better than permission,
and I am the Most Bountiful, the Generous.'' A most exalted, most holy
Tablet was also revealed and sent specifically for Jināb-i Āqā Mīrzā
Ḥasan---may the glory of God rest upon him---whom you mentioned. I ask
God to enable him to preserve this supreme favor. He is indeed the
Answerer.

\subsection{114}\label{section-2345}

\begin{otherlanguage}{persian}

در بارهٴ ابن جناب مرحوم آقا میرزا ع ب علیه بهآء اللّه آنچه مرقوم داشتند
ملاحظه شد بعد از عرض در ساحت اقدس فرمودند لوحی در بارهٴ او نازل اگر
مقتضی است برسانند والّا فلا ولکن آنمحبوب باید بسیار بحکمت حرکت فرمایند چه
که امثال آن نفوس را اگر انسان پند بگوید و یا نصیحت نماید همان کلمه سبب
بغضا شود و سبب فساد گردد

\end{otherlanguage}

Dar bāreh-yi ibn-i Jināb-i marḥūm Āqā Mīrzā ʿ B ʿalayhi Bahāʾuʾllāh
ānchih marqūm dāshtand mulāḥaẓeh shud. Baʿd az ʿarḍ dar sāḥat-i aqdas
farmūdand lawḥī dar bāreh-yi ū nāzil, agar muqtaḍī ast birasānand va
illā fa-lā. Va-lākin ān maḥbūb bāyad bisyār bi-ḥikmat ḥarakat farmāyand,
chih kih amthāl-i ān nufūs rā agar insān pand bigūyad va yā naṣīḥat
namāyad hamān kalameh sabab-i baghḍā shavad va sabab-i fasād gardad.

What you wrote concerning the son of the late Jināb-i Āqā Mīrzā ʿ.
B.---may the glory of God rest upon him---was read. After it was
presented in the Most Holy Court, He said: ``A Tablet has been revealed
concerning him. If it is appropriate, let them deliver it; otherwise,
no.'' That beloved one, however, must act with great wisdom, for if one
offers counsel or advice to people of that kind, those very words become
a cause of hatred and disorder.

\subsection{115}\label{section-2346}

\begin{otherlanguage}{persian}

آنچه در آخر این نامه از خضوع و خشوع و ابتهال و مناجات با حقّ جلّ جلاله ذکر
شده بود بشرف اصغا فائز قوله تبارک و تعالی انّا سمعنا منک ما ینبغی لساحة
عزّی و عرش کبریائی طوبی للسانک انّه نطق بما تضوّع منه عرف عنایتی و فضلی و
رحمتی و لعینک بما شهدت و رأت و لأذنک بما سمعت ندآء اللّه النّاطق العلیم
الحکیم قد فاز بالقبول توجّهک و حضورک و اقبالک و ابتهالک فی ایّام اللّه ربّ
العالمین نسأل اللّه ان یمدّک و یؤیّدک علی نشر امره و اعلآء کلمته و انتشار
احکامه انّه علی کلّشیء قدیر انتهی

\end{otherlanguage}

Ānchih dar ākhir-i īn nāmeh az khuḍūʿ va khushūʿ va ibtihāl va munājāt
bā Ḥaqq jalla jalāluhu dhikr shudeh būd bi-sharaf-i iṣghā fāʾiz. Qawluhu
tabāraka wa taʿālā innā samiʿnā minka mā yanbaghī li-sāḥati ʿizzī wa
ʿarshi kibriyāʾī. Ṭūbā li-lisānika innahu naṭaqa bi-mā taḍawwaʿa minhu
ʿarfu ʿināyatī wa faḍlī wa raḥmatī wa li-ʿaynika bi-mā shahidat wa raʾat
wa li-udhunika bi-mā samiʿat nidāʾa Allāhi al-Nāṭiqi al-ʿAlīmi al-Ḥakīm.
Qad fāza bi-l-qabūli tawajjuhuka wa ḥuḍūruka wa iqbāluka wa ibtihāluka
fī ayyāmi Allāhi Rabbi al-ʿālamīn. Nasʾalu Allāha an yumiddaka wa
yuʾayyidaka ʿalā nashri amrihi wa iʿlāʾi kalimatihi wa intishāri
aḥkāmihi innahu ʿalā kulli shayʾin qadīr. Intahā.

What was mentioned at the end of this letter concerning humility,
submissiveness, supplication, and prayer to God, exalted be His glory,
attained the honor of being heard. These are His words, blessed and
exalted is He: ``We have heard from you what befits the court of My
might and the throne of My grandeur. Blessed is your tongue, for it
spoke words from which the fragrance of My favor, My bounty, and My
mercy spread; and blessed is your eye for what it witnessed and saw, and
your ear for having heard the call of God, the Speaker, the All-Knowing,
the All-Wise. Your turning toward Me, your presence, your devotion, and
your supplication have been accepted in the days of God, the Lord of the
worlds. We ask God to aid and confirm you in spreading His Cause,
exalting His Word, and disseminating His laws. He has power over all
things.'' Here it ends.

\subsection{116}\label{section-2347}

\begin{otherlanguage}{arabic}

بعد از عرض مناجات و مراتب خضوع و ابتهال آنمحبوب فرات رحمت جاری و آفتاب
عنایت از افق سماء فضل مشرق بشأنی بحر بیان رحمن موّاج که اینعبد فانی از
ذکرش عاجز و قاصر از بدایع فضلش میطلبیم رشحی از این فضل بر عالمیان مبذول
فرماید تا کلّ فائز شوند بآنچه که سبب بقاء ابدی و حیوة سرمدیست

\end{otherlanguage}

Baʿd az ʿarḍ-i munājāt va marātib-i khuḍūʿ va ibtihāl-i ān maḥbūb,
Furāt-i raḥmat jārī va āftāb-i ʿināyat az ufuq-i samāʾ-i faḍl mushriq.
Bi-shaʾnī baḥr-i bayān-i Raḥmān mavvāj kih īn ʿabd-i fānī az dhikrash
ʿājiz va qāṣir. Az badāyiʿ-i faḍlash mīṭalabīm rashḥī az īn faḍl bar
ʿālamiyān mabdhūl farmāyad tā kull fāʾiz shavand bi-ānchih kih sabab-i
baqā-yi abadī va ḥayāt-i sarmadīst.

After that beloved one's prayer and expressions of humility and
supplication were presented, the Euphrates of mercy flowed and the sun
of loving-kindness shone from the horizon of the heaven of bounty. The
ocean of the All-Merciful's utterance surged to such a degree that this
evanescent servant is powerless and unable to describe it. We ask,
through the wonders of His bounty, that He bestow a sprinkling of this
grace upon the people of the world, so that all may attain what brings
everlasting existence and eternal life.

\subsection{117}\label{section-2348}

\begin{otherlanguage}{arabic}

در دستخطّ دیگر آنمحبوب بعد از مراتب ذکر و ثنا و ارتقاء بسماء عرفان مالک
اسماء ذکر حبیب روحانی جناب نوری علیه بهآء اللّه الأبهی را نمودند و همچنین
فقرات مذکورهٴ ایشانرا مرقوم داشتند بعد از عرض حضور لسان عظمت باین آیات
ناطق قوله تبارک و تعالی

\end{otherlanguage}

Dar dastkhaṭṭ-i dīgar-i ān maḥbūb baʿd az marātib-i dhikr va thanā va
irtiqā bi-samāʾ-i ʿirfān-i Mālik-i Asmāʾ dhikr-i ḥabīb-i rūḥānī Jināb-i
Nūrī ʿalayhi Bahāʾuʾllāhi al-Abhā rā namūdand va hamchunīn faqarāt-i
madhkūreh-yi īshān rā marqūm dāshtand. Baʿd az ʿarḍ-i ḥuḍūr lisān-i
ʿaẓamat bi-īn āyāt nāṭiq. Qawluhu tabāraka wa taʿālā.

In that beloved one's other letter, after expressions of remembrance and
praise and ascent to the heaven of recognition of the Lord of Names, you
mentioned the spiritual friend Jināb-i Nūrī---may the glory of God, the
All-Glorious, rest upon him---and also set down the passages he had
mentioned. After this was presented in His presence, the Tongue of
Grandeur spoke these verses. These are His words, blessed and exalted is
He:

\subsection{118}\label{section-2349}

\begin{otherlanguage}{persian}

ع ب لازال مذکور بوده و هستی ذکرش ذکرت تعلیم نمود عنایتش یدت را اخذ فرمود
رحمتش راه نمود و فضلش احاطه فرمود دانش عطا کرد بینش بخشید مدد فرمود
تأیید نمود للّه الحمد فائز شدی بآنچه که شبه از برای او نبوده و نیست در
جمیع نامه‌های ورقا علیه بهائی و عنایتی ذکرت بوده نسأله تعالی بأن یمدّک و
یؤیّدک علی ما یحبّ و یرضی و یقدّر لک ما یکون نوراً امام وجهک فی کلّ عالم من
عوالمه انّه علی کلّشیء قدیر انتهی

\end{otherlanguage}

ʿ B, lāzāl madhkūr būdeh va hastī. Dhikrash dhikrat taʿlīm namūd.
ʿInāyatash yadat rā akhdh farmūd. Raḥmatash rāh numūd va faḍlash iḥāṭeh
farmūd. Dānish ʿaṭā kard, bīnish bakhshīd, madad farmūd, taʾyīd namūd.
Li-llāhi al-ḥamd fāʾiz shudī bi-ānchih kih shibh az barā-yi ū nabūdeh va
nīst. Dar jamīʿ-i nāmehā-yi Varqā ʿalayhi bahāʾī va ʿināyatī dhikrat
būdeh. Nasʾaluhu taʿālā bi-an yumiddaka wa yuʾayyidaka ʿalā mā yuḥibbu
wa yarḍā wa yuqaddira laka mā yakūnu nūran amāma wajhika fī kulli
ʿālamin min ʿawālimihi innahu ʿalā kulli shayʾin qadīr. Intahā.

``ʿ. B., you have always been and remain remembered. His remembrance
taught you to remember Him. His loving-kindness took you by the hand.
His mercy showed you the way, and His bounty encompassed you. He granted
knowledge, bestowed insight, gave assistance, and lent support. Praise
be to God, you have attained what has never had and has no equal. You
have been mentioned in all the letters of Varqā---upon him be My glory
and My loving-kindness. We ask Him, exalted is He, to aid and confirm
you in what He loves and approves, and to ordain for you what will be a
light before your face in every one of His worlds. He has power over all
things.'' Here it ends.

\subsection{119}\label{section-2350}

\begin{otherlanguage}{arabic}

خبر ایشان و صحّت و سلامتی پی در پی از ارض طا رسیده احبّای آن ارض ذکر خیر
ایشانرا مکرّر نموده‌اند فی‌الحقیقه بحقّ ناطقند و بخدمت قائم ذکر جناب اخوی
آقا میرزا حسین علیه بهآء اللّه را نمودند بعد از عرض در ساحت اقدس این آیات
از ملکوت بیان رحمن نازل

\end{otherlanguage}

Khabar-i īshān va ṣiḥḥat va salāmatī pay dar pay az Arḍ-i Ṭā risīdeh.
Aḥibbā-yi ān arḍ dhikr-i khayr-i īshān rā mukarrar namūdeh-and.
Fī-l-ḥaqīqeh bi-ḥaqq nāṭiqand va bi-khidmat qāʾim. Dhikr-i Jināb-i
akhavī Āqā Mīrzā Ḥusayn ʿalayhi Bahāʾuʾllāh rā namūdand. Baʿd az ʿarḍ
dar sāḥat-i aqdas īn āyāt az malakūt-i bayān-i Raḥmān nāzil.

News of him and of his health and well-being has arrived repeatedly from
the Land of Ṭā. The friends in that land have often spoken well of him.
In truth, he speaks the truth and is engaged in service. You mentioned
your honored brother, Āqā Mīrzā Ḥusayn---may the glory of God rest upon
him. After this was presented in the Most Holy Court, these verses were
revealed from the kingdom of the All-Merciful's utterance:

\subsection{120}\label{section-2351}

\begin{otherlanguage}{persian}

قول الرّبّ تعالی و تقدّس یا حسین علیک بهائی انشآءاللّه فائز شوی بآنچه که
سزاوار است و از نهالهای باثمر باشی که ظهوراتش در ایّام الهی ظاهر شود حقّ
بید عنایت نهالهای جدیده کشته و بآب عنایت سقایه فرموده که شاید ثمر دهند و
سایه بخشند طوبی لمن فاز بهذا المقام الأسنی قل یا اله الأسمآء لک الحمد
بما شرّفتنی و اریتنی و اسمعتنی اسألک بأن تؤیّدنی علی ما تحبّ و ترضی انّک انت
مولی الوری و ربّ الآخرة و الأولی انتهی

\end{otherlanguage}

Qawlu al-Rabbi taʿālā wa taqaddas. Yā Ḥusayn, ʿalayka bahāʾī. In shāʾ
Allāh fāʾiz shavī bi-ānchih kih sazāvār ast va az nihālhā-yi bā-thamar
bāshī kih ẓuhūrātash dar ayyām-i ilāhī ẓāhir shavad. Ḥaqq bi-yad-i
ʿināyat nihālhā-yi jadīdeh kishteh va bi-āb-i ʿināyat siqāyeh farmūdeh
kih shāyad thamar dihand va sāyeh bakhshand. Ṭūbā li-man fāza bi-hādhā
al-maqāmi al-asnā. Qul yā Ilāha al-Asmāʾ laka al-ḥamdu bi-mā sharraftanī
wa araytanī wa asmaʿtanī. Asʾaluka bi-an tuʾayyidanī ʿalā mā tuḥibbu wa
tarḍā innaka anta Mawlā al-warā wa Rabbu al-ākhirati wa-l-ūlā. Intahā.

The words of the Lord, exalted and sanctified is He: ``O Ḥusayn, upon
you be My glory. God willing, may you attain what is fitting and be
among the fruitful saplings whose growth becomes manifest in the divine
days. God has planted new saplings with the hand of loving-kindness and
watered them with the water of His favor, that they may bear fruit and
give shade. Blessed is anyone who attains this most exalted station.
Say: `O God of Names, praise be to You for having honored me, enabled me
to see, and enabled me to hear. I ask You to confirm me in what You love
and approve. You are the Lord of all creation and the Lord of the world
to come and of this world.'\,'' Here it ends.

\subsection{121}\label{section-2352}

\begin{otherlanguage}{persian}

اینعبد هم خدمت ایشان سلام و تکبیر میرساند و از حقّ توفیق و تأیید میطلبد
اینکه در بارهٴ جناب استاد حسین علیه ٩ ٦٦ {[}بهآء الله{]} مرقوم داشتند که
تازه بشریعهٴ عرفان وارد شده بعد فوز ذکره باصغاء مولی الوری نیّر فضل از
افق ملکوت بیان مشرق و این آیات نازل قوله جلّ جلاله و عمّ نواله

\end{otherlanguage}

Īn ʿabd ham khidmat-i īshān salām va takbīr mīrasānad va az Ḥaqq tawfīq
va taʾyīd mīṭalabad. Īnkih dar bāreh-yi Jināb-i Ustād Ḥusayn ʿalayhi 9
66 {[}Bahāʾuʾllāh{]} marqūm dāshtand kih tāzeh bi-sharīʿeh-yi ʿirfān
vārid shudeh, baʿda fawzi dhikrihi bi-iṣghāʾi Mawlā al-warā nayyir-i
faḍl az ufuq-i malakūt-i bayān mushriq va īn āyāt nāzil. Qawluhu jalla
jalāluhu wa ʿamma nawāluhu.

This servant also conveys greetings and praise to him and asks God to
grant him success and support. Concerning what you wrote about Jināb-i
Ustād Ḥusayn---upon him be 9 66 {[}the glory of God{]}---that he has
newly reached the water's edge of recognition: after the mention of him
attained the hearing of the Lord of all creation, the luminary of bounty
shone from the horizon of the kingdom of utterance, and these verses
were revealed. These are His words, exalted be His glory and
all-embracing His bounty:

\subsection{122}\label{section-2353}

\begin{otherlanguage}{persian}

احمد اللّه بما ایّدک علی الاقبال فی یوم اخذ فیه الدّخان سکّان الأرض و علّمک
اذ کان العباد علی جهل مبین قل لک البهآء یا مولی الأسماء و لک العنایة یا
مولی البریّه بما ایّدتنی علی عرفان بحر بیانک و سمآء عطائک و نوّرتنی بأنوار
شمس ظهورک ای ربّ ایّدنی علی الاستقامة علی امرک ثمّ وفّقنی علی ذکرک و ثنائک
انّک انت المقتدر العزیز القدیر و بالاجابة جدیر انتهی

\end{otherlanguage}

Iḥmadi Allāha bi-mā ayyadaka ʿalā al-iqbāli fī yawmin akhadha fīhi
al-dukhānu sukkāna al-arḍi wa ʿallamaka idh kāna al-ʿibādu ʿalā jahlin
mubīn. Qul laka al-bahāʾu yā Mawlā al-Asmāʾ wa laka al-ʿināyatu yā Mawlā
al-bariyyah bi-mā ayyadtanī ʿalā ʿirfāni baḥri bayānika wa samāʾi
ʿaṭāʾika wa nawwartanī bi-anwāri shamsi ẓuhūrika. Ay Rabbi ayyidnī ʿalā
al-istiqāmati ʿalā amrika thumma waffiqnī ʿalā dhikrika wa thanāʾika
innaka anta al-Muqtadiru al-ʿAzīzu al-Qadīru wa bi-l-ijābati jadīr.
Intahā.

``Praise God for having enabled you to turn toward Him on a day when
smoke has enveloped the inhabitants of the earth, and for having taught
you when the people were in manifest ignorance. Say: `Yours is the
glory, O Lord of Names, and Yours is the loving-kindness, O Lord of
creation, for having enabled me to recognize the ocean of Your utterance
and the heaven of Your bounty, and for having illumined me with the rays
of the sun of Your revelation. O Lord, confirm me in steadfastness in
Your Cause, and enable me to remember and praise You. You are the
All-Powerful, the Almighty, the Omnipotent, and worthy to answer
prayers.'\,'' Here it ends.

\subsection{123}\label{section-2354}

\begin{otherlanguage}{arabic}

و اینکلمات درّیّات از مشرق ارادهٴ مولی الأسمآء و الصّفات مخصوص دو ابن جناب
ملّا مصطفی و آقا غلامحسین علیهم بهآء اللّه نازل قول الرّب تعالی و تقدّس

\end{otherlanguage}

Va īn kalimāt-i durriyyāt az mashriq-i irādeh-yi Mawlā al-Asmāʾ
wa-l-Ṣifāt makhṣūṣ-i du ibn-i Jināb-i Mullā Muṣṭafā va Āqā Ghulām-Ḥusayn
ʿalayhim Bahāʾuʾllāh nāzil. Qawlu al-Rabbi taʿālā wa taqaddas.

And these pearl-like words were revealed from the dawning-place of the
will of the Lord of Names and Attributes specifically for the two sons
of Jināb-i Mullā Muṣṭafā and Āqā Ghulām-Ḥusayn---may the glory of God
rest upon them. The words of the Lord, exalted and sanctified is He:

\subsection{124}\label{section-2355}

\begin{otherlanguage}{persian}

یا ابنآء الانسان جهد نمائید از شما ظاهر شود آنچه لایق ایّام الهی است
اینمظلوم شما را وصیّت میفرماید بآنچه که سبب علوّ و سموّ است زیّنوا هیاکلکم
بطراز الأعمال و رؤوسکم باکلیل الأخلاق لتستضیء بهما الآفاق نور ظاهر و
مشرق و نیّر امر مشهود و لائح دیوان یعنی نفوس غافله در کمین انسان بوده و
هستند خود را باسم حقّ حفظ نمائید آگاه باشید بیدار باشید اینست وصیّت او
طوبی للعاملین انتهی

\end{otherlanguage}

Yā abnāʾa al-insān, jahd namāyīd az shumā ẓāhir shavad ānchih lāyiq-i
ayyām-i ilāhī ast. Īn Maẓlūm shumā rā vaṣiyyat mīfarmāyad bi-ānchih kih
sabab-i ʿuluvv va sumuvv ast. Zayyinū hayākilakum bi-ṭirāzi al-aʿmāli wa
ruʾūsakum bi-iklīli al-akhlāqi li-tastaḍīʾa bi-himā al-āfāq. Nūr ẓāhir
va mushriq va nayyir-i amr mashhūd va lāʾiḥ. Dīvān yaʿnī nufūs-i
ghāfileh dar kamīn-i insān būdeh va hastand. Khud rā bi-ism-i Ḥaqq ḥifẓ
namāyīd. Āgāh bāshīd, bīdār bāshīd. Īnast vaṣiyyat-i ū. Ṭūbā
li-l-ʿāmilīn. Intahā.

``O children of humanity, strive to manifest what befits the divine
days. This Wronged One counsels you to do what will bring exaltation and
elevation. Adorn your persons with the ornament of deeds and your heads
with the crown of good character, so that the world may be illumined by
them both. The light is manifest and shining, and the luminary of the
Cause is visible and radiant. Demons---that is, heedless souls---have
always lain and still lie in wait for human beings. Protect yourselves
through the name of God. Be aware; be watchful. This is His counsel.
Blessed are those who act upon it.'' Here it ends.

\subsection{125}\label{section-2356}

\begin{otherlanguage}{persian}

شکر مقصود عالم را عنایتش کلّ را شامل نیسان رحمتش لازال بر اراضی قلوب
من‌غیر توقّف هاطل و نازل جودش وجود را احاطه نموده و فضلش عطا فرموده آنچه
را که شبه و مثل نداشته ولکن آنمحبوب فؤاد و اینعبد باید از حقّ جلّ جلاله
بطلبیم تا حجبات مانعه را خرق فرماید شاید ببینند و بشناسند اینکه در بارهٴ
مخدّرهٴ مرفوعه امّ حبیبان جناب آقا میرزا سلیمان و جناب آقا میرزا مهدی
علیها و علیهما بهآء اللّه ذکر نمودند بعد از عرض در پیشگاه حضور اینکلمات
عالیات از مطلع ارادهٴ حقّ جلّ جلاله نازل و ظاهر قوله جلّ بیانه و جلّ برهانه

\end{otherlanguage}

Shukr-i Maqṣūd-i ʿālam rā. ʿInāyatash kull rā shāmil. Naysān-i raḥmatash
lāzāl bar arāḍī-yi qulūb min ghayr-i tavaqquf hāṭil va nāzil. Jūdash
vujūd rā iḥāṭeh namūdeh va faḍlash ʿaṭā farmūdeh ānchih rā kih shibh va
mithl nadāshteh. Va-lākin ān maḥbūb-i fuʾād va īn ʿabd bāyad az Ḥaqq
jalla jalāluhu biṭalabīm tā ḥujubāt-i māniʿeh rā kharq farmāyad, shāyad
bibīnand va bishināsand. Īnkih dar bāreh-yi mukhaddareh-yi marfūʿeh
umm-i ḥabībayn Jināb-i Āqā Mīrzā Sulaymān va Jināb-i Āqā Mīrzā Mahdī
ʿalayhā wa ʿalayhimā Bahāʾuʾllāh dhikr namūdand, baʿd az ʿarḍ dar
pīshgāh-i ḥuḍūr īn kalimāt-i ʿāliyāt az maṭlaʿ-i irādeh-yi Ḥaqq jalla
jalāluhu nāzil va ẓāhir. Qawluhu jalla bayānuhu wa jalla burhānuhu.

Thanks be to the Desire of the world. His loving-kindness embraces
everyone. The spring rain of His mercy has always fallen ceaselessly
upon the soil of hearts. His generosity has encompassed existence, and
His bounty has bestowed what has had neither equal nor likeness. Yet
that beloved of the heart and this servant must ask God, exalted be His
glory, to rend the obstructing veils, that perhaps they may see and
recognize. Concerning your mention of the exalted lady, the mother of
the two dear ones, Jināb-i Āqā Mīrzā Sulaymān and Jināb-i Āqā Mīrzā
Mahdī---may the glory of God rest upon her and upon them both---after
this was presented in His presence, these exalted words were revealed
and made manifest from the dawning-place of the will of God, exalted be
His glory. These are His words, exalted be His utterance and exalted His
proof:

\subsection{126}\label{section-2357}

\begin{otherlanguage}{arabic}

یا قلم اذکر الورقة الّتی خرجت من دیارها مقبلة الی الأفق الأعلی لتحضر امام
الوجه و تقوم لدی الباب و تسمع ندآء المظلوم و تری افق الظّهور اشهد انّها
سمعت ندآء ربّها و آمنت به و اقبلت الیه و شهدت بما شهد به لسان العظمة طوبی
لک یا امتی و یا ورقتی انت الّتی فزت بذکری قبل صعودک و بعده طوبی لک بما
فزت بأیّام اللّه و آثاره فی یوم فیه اعرضن امآء الأرض و ملکاتها انّا وجدنا
منک عرف حبّی و التّوجّه الی انوار وجهی نسأل اللّه تعالی ان یدخلک فی الفردوس
الأعلی و یزیّنک بالاکلیل المزیّنة باسمه الأبهی و ینزل علیک فیکلّ حین رحمة
من عنده و نوراً من لدنه انّه هو المقتدر الغفور الکریم البهآء من لدنّا علیک
و علی الامآء اللّائی یذکرنک بما نطق القلم الأعلی فی هذه اللّیلة الدّلمآء
انتهی

\end{otherlanguage}

Yā qalam, udhkur al-waraqata allatī kharajat min diyārihā muqbilatan ilā
al-ufuq al-aʿlā li-taḥḍura amāma al-wajh wa-taqūma ladā al-bāb
wa-tasmaʿa nidāʾ al-Maẓlūm wa-tarā ufuq al-ẓuhūr. Ashhadu annahā samiʿat
nidāʾ rabbihā wa-āmanat bihi wa-aqbalat ilayhi wa-shahidat bimā shahida
bihi lisān al-ʿaẓama. Ṭūbā laki yā amatī wa-yā waraqatī, anti allatī
fuzti bi-dhikrī qabla ṣuʿūdiki wa-baʿdahu. Ṭūbā laki bimā fuzti bi-ayyām
Allāh wa-āthārihi fī yawmin fīhi aʿraḍna imāʾ al-arḍ wa-malikātuhā. Innā
wajadnā minki ʿarfa ḥubbī wa-l-tawajjuh ilā anwār wajhī. Nasʾalu Allāh
taʿālā an yudkhilaki fī al-firdaws al-aʿlā wa-yuzayyinaki bi-l-iklīl
al-muzayyana bi-smihi al-Abhā wa-yunzila ʿalayki fī kulli ḥīnin raḥmatan
min ʿindihi wa-nūran min ladunhu. Innahu huwa al-Muqtadir al-Ghafūr
al-Karīm. Al-bahāʾ min ladunnā ʿalayki wa-ʿalā al-imāʾ allāʾī
yadhkurnaki bimā naṭaqa al-qalam al-aʿlā fī hādhihi al-layla al-dalmāʾ.
Intahā.

``O Pen, remember the leaf who left her homeland, turning toward the
Most Exalted Horizon, to appear before the Countenance, stand at the
door, hear the call of the Wronged One, and behold the horizon of
Revelation. I testify that she heard the call of her Lord, believed in
Him, turned toward Him, and bore witness to that to which the Tongue of
Grandeur bore witness. Blessed are you, My handmaiden and My leaf! You
attained My remembrance before your ascension and after it. Blessed are
you for having attained the days of God and His writings on a day when
the women of the earth and its queens turned away. We perceived from you
the fragrance of love for Me and of turning toward the lights of My
face. We ask God, exalted be He, to admit you into the highest Paradise,
adorn you with the crown adorned with His All-Glorious name, and send
down upon you at every moment mercy from Himself and light from His
presence. He is indeed the All-Powerful, the Ever-Forgiving, the
Generous. Glory from Us rest upon you and upon the handmaidens who
remember you through what the Most Exalted Pen has uttered on this dark
night.'' Here it ends.

\subsection{127}\label{section-2358}

\begin{otherlanguage}{persian}

فی‌الحقیقه این ایّام ایّامیست انسان از یکجهت تفکّر مینماید موت میطلبد چه که
بعد از ارتقا بآثار قلم اعلی فائز میشود و از جهت آخر بقا میطلبد بامید
آنکه خدمتی از او ظاهر شود ولکن حقّ آنکه هرچه محبوب اراده فرموده آن
محبوبست والسّلام خیر ختام مخصوص فرمودند جناب میرزا سلیمان و میرزا مهدی
علیهما بهائی را تسلّی میدهیم و ذکر مینمائیم بآنچه که ابواب فرح بگشاید و
حزن بسرور اکبر تبدیل شود در حین صعود از قلم اعلی در بارهٴ او نازل شد
آنچه که باقی و دائمست انتهی

\end{otherlanguage}

Fī al-ḥaqīqe īn ayyām ayyāmīst insān az yek jahat tafakkor mīnamāyad
mawt mīṭalabad, che ke baʿd az irtiqā be-āthār-e qalam-e aʿlā fāʾez
mīshavad, va az jahat-e ākhar baqā mīṭalabad be-omīd-e ānke khedmatī az
ū ẓāher shavad. Va-lāken ḥaqq ānke harche Maḥbūb erāde farmūde ān maḥbūb
ast. Va-l-salām khayr khitām. Makhṣūṣ farmūdand: Janāb-e Mīrzā Solaymān
va Mīrzā Mahdī, ʿalayhimā bahāʾī, rā tasallī mīdahīm va dhekr mīnamāʾīm
be-ānche ke abvāb-e faraḥ bogshāyad va ḥozn be-sorūr-e akbar tabdīl
shavad. Dar ḥīn-e ṣoʿūd az qalam-e aʿlā dar bāre-ye ū nāzel shod ānche
ke bāqī va dāʾem ast. Intahā.

Truly, these are days when, reflecting from one perspective, a person
longs for death, since after ascending one attains the writings of the
Most Exalted Pen; yet from another perspective, one longs to remain
alive, in the hope of rendering some service. The truth, however, is
that whatever the Beloved has willed is what is beloved. Peace is the
best conclusion. He said specifically: ``We console Jināb-i Mīrzā
Sulaymān and Mīrzā Mahdī---My glory rest upon them both---and remember
them with words that will open the gates of joy and transform sorrow
into the greatest happiness. At the time of her ascension, there was
revealed concerning her from the Most Exalted Pen that which is abiding
and everlasting.'' Here it ends.

\subsection{128}\label{section-2359}

\begin{otherlanguage}{persian}

از قبل هم در بارهٴ مخدّرهٴ مذکوره از قلم اعلی نازل شد آنچه که عرفش بدوام
ملک و ملکوت باقی و پاینده است هنیئاً لها آنچه در بارهٴ جناب آقا میرزا
محمود مراغه علیه بهآء اللّه ذکر فرمودند باجابت مقرون و اینعبد فانی بنیابت
آنحضرت و ایشان در وقتی مخصوص بحضور فائز و بقبول مزیّن اسأله تعالی ان یجعل
قطرة عملی بحراً من عنده انّه علی کلّشیء قدیر و همچنین ذکر جناب آقا میرزا
مهدی شفافی علیه ٩ ٦٦ {[}بهآء الله{]} را فرمودند و استدعای عنایت مخصوص
نمودند بعد از عرض لسان عظمت باین آیات ناطق قول الرّبّ تعالی و تقدّس

\end{otherlanguage}

Az qabl ham dar bāre-ye mokhaddare-ye madhkūre az qalam-e aʿlā nāzel
shod ānche ke ʿarfash be-davām-e molk va malakūt bāqī va pāyande ast.
Hanīʾan lahā. Ānche dar bāre-ye Janāb-e Āqā Mīrzā Maḥmūd-e Marāghe,
ʿalayhi bahāʾu Allāh, dhekr farmūdand be-ejābat maqrūn, va īn ʿabd-e
fānī be-niyābat-e ān ḥażrat va īshān dar vaqtī makhṣūṣ be-ḥożūr fāʾez va
be-qabūl mozayyan. Asʾaluhu taʿālā an yajʿala qaṭrata ʿamalī baḥran min
ʿindihi, innahu ʿalā kulli shayʾin qadīr. Va hamchenīn dhekr-e Janāb-e
Āqā Mīrzā Mahdī-ye Shaffāfī, ʿalayhi 9 66 {[}bahāʾu Allāh{]}, rā
farmūdand va estedʿā-ye ʿenāyat-e makhṣūṣ namūdand. Baʿd az ʿarż lesān-e
ʿaẓamat be-īn āyāt nāṭeq. Qawl al-Rabb taʿālā wa-taqaddas.

Previously, too, there was revealed from the Most Exalted Pen concerning
the aforementioned lady that whose fragrance will endure as long as the
earthly realm and the Kingdom endure. Blessed is she! What you mentioned
concerning Jināb-i Āqā Mīrzā Maḥmūd of Marāghih---may the glory of God
rest upon him---was granted, and this evanescent servant, acting on your
behalf and his, attained His presence at a specially appointed time and
was adorned with acceptance. I ask Him, exalted be He, to make the drop
of my deed an ocean from Himself. He is powerful over all things. You
also mentioned Jināb-i Āqā Mīrzā Mahdī Shaffāfī---upon him be 9 66
{[}the glory of God{]}---and requested special favor for him. After this
was presented, the Tongue of Grandeur uttered these verses. The words of
the Lord, exalted and sanctified be He:

\subsection{129}\label{section-2360}

\begin{otherlanguage}{arabic}

بسمی المظلوم الغریب

\end{otherlanguage}

Bi-smī al-Maẓlūm al-Gharīb.

``In My name, the Wronged One, the Stranger.''

\subsection{130}\label{section-2361}

\begin{otherlanguage}{persian}

حقّ جلّ جلاله در جمیع احیان بذکر دوستان مشغول لسان عظمت ناطق و قلمش متحرّک
لعمراللّه بآثار قلمش معادله نمینماید آنچه ظاهر و مشهود است یا مهدی اشکر
ربّک بهذه الکلمة العلیا انّک اذا وجدت نفحات بیان الرّحمن قل

\end{otherlanguage}

Ḥaqq jalla jalāluhu dar jamīʿ-e aḥyān be-dhekr-e dūstān mashghūl,
lesān-e ʿaẓamat nāṭeq va qalamash motaḥarrek. La-ʿamru Allāh, be-āthār-e
qalamash moʿādele namīnamāyad ānche ẓāher va mashhūd ast. Yā Mahdī,
ushkur rabbaka bi-hādhihi al-kalima al-ʿulyā. Innaka idhā wajadta
nafaḥāt bayān al-Raḥmān qul.

``God, exalted be His glory, is at all times occupied with remembering
His friends. The Tongue of Grandeur speaks and His Pen moves. By the
life of God, nothing visible or manifest can equal the writings of His
Pen. O Mahdī, thank your Lord for this exalted utterance. When you sense
the fragrances of the All-Merciful's words, say:

\subsection{131}\label{section-2362}

\begin{otherlanguage}{arabic}

الهی الهی اشهد انّک خلقتنی لعرفانک و التّوجّه الی انوار وجهک و الاقبال الی
افق ظهورک و القیام لدی باب عظمتک و الحضور فی ساحة عزّک اسألک بآیاتک
الکبری و ظهورات قدرتک فی ناسوت الانشآء و ببحر فضلک و سمآء علمک بأن
تجعلنی مستقیماً علی امرک و منادیاً باسمک ای ربّ اکتب من قلم جودک لعبدک هذا
ما یقرّبه الیک فی کلّ عالم من عوالمک ثمّ اکتب له ما کتبته لعبادک الّذین ما
نقضوا عهدک و میثاقک و شربوا رحیق الوحی فی ایّامک من ایادی عطائک انّک انت
المقتدر المهیمن العزیز الحکیم انتهی

\end{otherlanguage}

Ilāhī ilāhī, ashhadu annaka khalaqtanī li-ʿirfānika wa-l-tawajjuh ilā
anwār wajhika wa-l-iqbāl ilā ufuq ẓuhūrika wa-l-qiyām ladā bāb
ʿaẓamatika wa-l-ḥuḍūr fī sāḥat ʿizzika. Asʾaluka bi-āyātika al-kubrā
wa-ẓuhūrāt qudratika fī nāsūt al-inshāʾ wa-bi-baḥr faḍlika wa-samāʾ
ʿilmika bi-an tajʿalanī mustaqīman ʿalā amrika wa-munādiyan bi-smika. Ay
rabb, uktub min qalam jūdika li-ʿabdika hādhā mā yuqarribuhu ilayka fī
kulli ʿālamin min ʿawālimika. Thumma uktub lahu mā katabtahu li-ʿibādika
alladhīna mā naqaḍū ʿahdaka wa-mīthāqaka wa-sharibū raḥīq al-waḥy fī
ayyāmika min ayādī ʿaṭāʾika. Innaka anta al-Muqtadir al-Muhaymin
al-ʿAzīz al-Ḥakīm. Intahā.

`My God, my God! I testify that You created me to know You, to turn
toward the lights of Your face, to direct myself toward the horizon of
Your Revelation, to stand at the gate of Your grandeur, and to appear in
the court of Your glory. I ask You, by Your greatest signs, by the
manifestations of Your power in the earthly realm of creation, by the
ocean of Your grace and the heaven of Your knowledge, to make me
steadfast in Your Cause and a herald of Your name. O Lord, ordain with
the Pen of Your generosity for this servant of Yours what will draw him
near to You in every world among Your worlds. Then ordain for him what
You have ordained for Your servants who have not broken Your covenant
and pledge, and who, in Your days, have drunk the choice wine of
revelation from the hands of Your bounty. You are indeed the
All-Powerful, the All-Protecting, the Almighty, the All-Wise.'\,'' Here
it ends.

\subsection{132}\label{section-2363}

\begin{otherlanguage}{persian}

و دستخطّ آنمحبوب فؤاد که باسم حبیب مکرّم جناب الف و حا علیه بهآء اللّه و
عنایته ارسال شد بشرف اصغاء مقصود عالمیان فائز قوله تبارک و تعالی

\end{otherlanguage}

Va dastkhaṭṭ-e ān maḥbūb-e foʾād ke be-esm-e ḥabīb-e mokarram Janāb-e
Alif va Ḥā, ʿalayhi bahāʾu Allāh va ʿenāyatuhu, ersāl shod be-sharaf-e
eṣghā-ye Maqṣūd-e ʿālamiyān fāʾez. Qawluhu tabāraka wa-taʿālā.

And your handwritten letter, beloved of the heart, which was sent in the
name of the honored friend, Jināb-i Alif and Ḥā---may the glory of God
and His loving-kindness rest upon him---attained the honor of being
heard by the Desire of all worlds. His words, blessed and exalted be He:

\subsection{133}\label{section-2364}

\begin{otherlanguage}{arabic}

الحمد للّه الّذی ایّد الألف و الحآء علی الاقبال الیه و التّمسّک بحبل فضله و
التّشبّث بأذیال رحمته و القیام علی خدمة امره انّه هو المشفق المؤیّد الغفور
الکریم یا ورقآء علیک بهائی و عنایتی قلم اعلی شهادت میدهد بر اقبال و توجّه
و خدمت او در ایّامیکه از سطوت و طغیان اهل ارض فرائص وجود مضطرب او بحقّ
مقبل و بعرفانش فائز و بذکرش ناطق و بکلمه‌اش مقرّ و معترف طوبی له و نعیماً
له بشّره من قبلی و ذکّره بما نزّل من قلمی لیفرح و یکون من الشّاکرین انتهی

\end{otherlanguage}

Al-ḥamdu li-llāh alladhī ayy ada al-Alif wa-l-Ḥāʾ ʿalā al-iqbāl ilayhi
wa-l-tamassuk bi-ḥabl faḍlihi wa-l-tashabbuth bi-adhyāl raḥmatihi
wa-l-qiyām ʿalā khidmat amrihi. Innahu huwa al-Mushfiq al-Muʾayyid
al-Ghafūr al-Karīm. Yā Varqāʾ, ʿalayka bahāʾī wa-ʿināyatī. Qalam-e aʿlā
shahādat mīdahad bar eqbāl va tavajjoh va khedmat-e ū. Dar ayyāmī ke az
saṭvat va ṭoghyān-e ahl-e arż farāʾeṣ-e vojūd możṭareb, ū be-Ḥaqq moqbel
va be-ʿerfānash fāʾez va be-dhekrash nāṭeq va be-kalame-ash moqerr va
moʿtaref. Ṭūbā lahu wa-naʿīman lahu. Bashshirhu min qibalī wa-dhakkirhu
bimā nuzzila min qalamī li-yafraḥa wa-yakūna min al-shākirīn. Intahā.

``Praise be to God, who has aided Alif and Ḥā to turn toward Him, hold
fast to the cord of His grace, cling to the hem of His mercy, and arise
to serve His Cause. He is indeed the Compassionate, the Helper, the
Ever-Forgiving, the Generous. O Varqā, My glory and loving-kindness rest
upon you. The Most Exalted Pen testifies to his turning toward God, his
devotion, and his service. In days when the limbs of existence trembled
at the might and tyranny of the people of the earth, he turned toward
God, attained the knowledge of Him, spoke His remembrance, and
acknowledged and professed His Word. Blessed is he, and bliss be his!
Give him glad tidings on My behalf and remind him of what has been sent
down from My Pen, that he may rejoice and be among the thankful.'' Here
it ends.

\subsection{134}\label{section-2365}

\begin{otherlanguage}{arabic}

جناب محبوب روحانی امین علیه بهآء اللّه موافق ارادهٴ ایشان عمل نمودند مکرّر
ذکر ایشان و ذکر ابناء ایشان و ذکر آقا میر علی اکبر و آقا میر علی اصغر
علیهم بهآء اللّه را نموده و همچنین اولیای دیگر را از اهل میم و یا و سیسان
جمیع را ذکر نموده و در اکثر نامه‌های ایشان ذکر آنحضرت هم بوده از حقّ
میطلبم ایشان را تأیید فرماید در جمیع احوال بمحبّت اولیائش انّه هو المؤیّد
المقتدر العلیم الحکیم

\end{otherlanguage}

Janāb-e maḥbūb-e rūḥānī Amīn, ʿalayhi bahāʾu Allāh, movāfeq-e erāde-ye
īshān ʿamal namūdand. Mokarrar dhekr-e īshān va dhekr-e abnā-ye īshān va
dhekr-e Āqā Mīr ʿAlī Akbar va Āqā Mīr ʿAlī Aṣghar, ʿalayhim bahāʾu
Allāh, rā namūde va hamchenīn awliyā-ye dīgar rā az ahl-e Mīm va Yā va
Sīsān jamīʿ rā dhekr namūde, va dar akthar-e nāmehā-ye īshān dhekr-e ān
ḥażrat ham būde. Az Ḥaqq mīṭalabam īshān rā taʾyīd farmāyad dar jamīʿ-e
aḥvāl be-maḥabbat-e awliyāyash. Innahu huwa al-Muʾayyid al-Muqtadir
al-ʿAlīm al-Ḥakīm.

The spiritually beloved Jināb-i Amīn---may the glory of God rest upon
him---acted in accordance with his wishes. He has repeatedly mentioned
him, his sons, Āqā Mīr ʿAlī Akbar, and Āqā Mīr ʿAlī Aṣghar---may the
glory of God rest upon them---and has likewise mentioned all the other
friends among the people of Mīm and Yā and Sīsān. In most of his
letters, he has also mentioned you. I ask God to aid him under all
circumstances in loving His friends. He is indeed the Helper, the
All-Powerful, the All-Knowing, the All-Wise.

\subsection{135}\label{section-2366}

\begin{otherlanguage}{persian}

ورقهٴ اخری که از سدرهٴ محبّت محبوب عالمیان روئید بصورت و هیئة نامه رسید
آفرین بر آن خامه و آمه که آب معنوی از او جاری و از آن آب که از عین محبّت
جاری چون بر ارض اوراق رسید ازهار خلوص مخصوص و خضوع و خشوع و ذکر و ثنا
روئید صدهزار آفرین بر آنخامه و بر آن آمه و آن نامه و بعد از استشمام
ازهار و اطّلاع بر آثار قصد مقام اعزّ اعلی نموده بشرف اصغا فائز

\end{otherlanguage}

Varaqe-ye okhrā ke az sedre-ye maḥabbat-e Maḥbūb-e ʿālamiyān rūʾīd
be-ṣūrat va heyʾat-e nāme resīd. Āfarīn bar ān khāme va āme ke āb-e
maʿnavī az ū jārī, va az ān āb ke az ʿayn-e maḥabbat jārī chon bar arż-e
awrāq resīd azhār-e kholūṣ-e makhṣūṣ va khożūʿ va khoshūʿ va dhekr va
thanā rūʾīd. Ṣad hezār āfarīn bar ān khāme va bar ān āme va ān nāme. Va
baʿd az esteshmām-e azhār va eṭṭelāʿ bar āthār qaṣd-e maqām-e aʿazz-e
aʿlā namūde be-sharaf-e eṣghā fāʾez.

Another leaf, which grew from the tree of love for the Beloved of all
worlds, arrived in the form and shape of a letter. Praise be to that pen
and inkwell from which spiritual water flows! When that water, flowing
from the spring of love, reached the soil of the pages, flowers of
singular sincerity, humility, reverence, remembrance, and praise sprang
forth. A hundred thousand praises to that pen, that inkwell, and that
letter! After inhaling the fragrance of the flowers and acquainting
myself with the writings, I proceeded to the most glorious, most exalted
station, and the letter attained the honor of being heard.

\subsection{136}\label{section-2367}

\begin{otherlanguage}{arabic}

قول الرّبّ تعالی و تقدّس بسم الّذی به حمل اللّوح و تحرّک القلم

\end{otherlanguage}

Qawl al-Rabb taʿālā wa-taqaddas: Bi-sm alladhī bihi ḥumila al-lawḥ
wa-taḥarraka al-qalam.

The words of the Lord, exalted and sanctified be He: ``In the name of
Him through whom the Tablet was borne and the Pen set in motion.''

\subsection{137}\label{section-2368}

\begin{otherlanguage}{persian}

علیک بهآئی و عنایتی و رحمتی عبد حاضر با نامه حاضر ندای لسان فؤاد بشرف
اصغا فائز شنیدیم و دیدیم طوبی لقلبک بما اقبل الی افقی المنیر و للسانه
بما نطق بثنائی الجمیل و لوجهه بما استضآء من انوار الوجه و لذاته بما تذوت
بحبّی العزیز المنیع یا ورقآء علیک عنایة اللّه و رحمته ایّام حضورت لازال در
پیشگاه حضور مذکور قد قمت فی اوّل ایّامک علی خدمة امری و هدایة عبادی الی
صراطی و تهذیب النّفوس بآیاتی و تطهیرهم بمآء عنایتی الّذی جری من سحاب سمآء
مشیّتی طوبی لک و نعیماً لک و لمن احبّک لوجه اللّه ربّک و ربّ العالمین البهآء
علیک و علی من معک انتهی

\end{otherlanguage}

ʿAlayka bahāʾī wa-ʿināyatī wa-raḥmatī. ʿAbd ḥāżer bā nāme ḥāżer, nedā-ye
lesān-e foʾād be-sharaf-e eṣghā fāʾez. Shenīdīm va dīdīm. Ṭūbā
li-qalbika bimā aqbala ilā ufuqī al-munīr, wa-li-lisānihi bimā naṭaqa
bi-thanāʾī al-jamīl, wa-li-wajhihi bimā istaḍāʾa min anwār al-wajh,
wa-li-dhātihi bimā tadhawwata bi-ḥubbī al-ʿAzīz al-Manīʿ. Yā Varqāʾ,
ʿalayka ʿināyatu Allāh wa-raḥmatuhu. Ayyām-e ḥożūrat lā zāl dar
pīshgāh-e ḥożūr madhkūr. Qad qumta fī awwal ayyāmika ʿalā khidmat amrī
wa-hidāyat ʿibādī ilā ṣirāṭī wa-tahdhīb al-nufūs bi-āyātī wa-taṭhīrihim
bi-māʾ ʿināyatī alladhī jarā min saḥāb samāʾ mashīʾatī. Ṭūbā laka
wa-naʿīman laka wa-li-man aḥabbaka li-wajh Allāh rabbika wa-rabb
al-ʿālamīn. Al-bahāʾ ʿalayka wa-ʿalā man maʿaka. Intahā.

``My glory, My loving-kindness, and My mercy rest upon you. The servant
is present with the letter; the call of the tongue of the heart has
attained the honor of being heard. We have heard and seen. Blessed is
your heart, for it has turned toward My luminous horizon; and its
tongue, for it has voiced My beautiful praise; and its face, for it has
been illumined by the lights of the Countenance; and its essence, for it
has gained its being through love for Me, the Almighty, the
Inaccessible. O Varqā, the loving-kindness of God and His mercy rest
upon you. The days of your presence are ever remembered in this court.
In your early days you arose to serve My Cause, guide My servants to My
path, refine souls through My verses, and purify them with the water of
My loving-kindness that flowed from the clouds of the heaven of My will.
Blessed are you, and bliss be yours and his who loves you for the sake
of God, your Lord and the Lord of all worlds. Glory rest upon you and
upon those who are with you.'' Here it ends.

\subsection{138}\label{section-2369}

\begin{otherlanguage}{persian}

له الحمد و الشّکر چه که در این لیلهٴ مبارکه عنایت مخصوصه شامل و آنحضرت
بطراز مخصوص مزیّن انشآءاللّه مؤیّد و موفّق باشید بر تربیت عباد و تهذیب نفوس
و تجذیب افئده و قلوب بعد از عرض مناجات امام وجه بحر اعظم بموجی ظاهر که
اینعبد از ذکرش عاجز در آخر لسان عظمت باینکلمهٴ مبارکه ناطق ورقآء علیه
بهائی و عنایتی بوجه قلب و لسان قلب و قلب قلب و ذات قلب توجّه نمود انّه
اشتعل بنار حبّ ربّک و فاز بکوثر ذکره الجذّاب الفضّال المبعث المحیی المحرّک
العلیم الحکیم و بعد باینکلمهٴ مبارکه ختم شد فرمودند هنیئاً له

\end{otherlanguage}

Lahu al-ḥamd wa-l-shukr, che ke dar īn layle-ye mobārake ʿenāyat-e
makhṣūṣe shāmel va ān ḥażrat be-ṭerāz-e makhṣūṣ mozayyan. In shāʾ Allāh
moʾayyad va movaffaq bāshīd bar tarbiyat-e ʿebād va tahdhīb-e nofūs va
tajdhīb-e afʾede va qolūb. Baʿd az ʿarż-e monājāt amām-e vajh baḥr-e
aʿẓam be-mawjī ẓāher ke īn ʿabd az dhekrash ʿājez. Dar ākhar lesān-e
ʿaẓamat be-īn kalame-ye mobārake nāṭeq: Varqāʾ, ʿalayhi bahāʾī
wa-ʿināyatī, be-vajh-e qalb va lesān-e qalb va qalb-e qalb va dhāt-e
qalb tavajjoh namūd. Innahu ishtaʿala bi-nār ḥubb rabbika wa-fāza
bi-kawthar dhikrihi al-Jadhdhāb al-Faḍḍāl al-Mubʿith al-Muḥyī
al-Muḥarrik al-ʿAlīm al-Ḥakīm. Va baʿd be-īn kalame-ye mobārake khatm
shod. Farmūdand: Hanīʾan lahu.

Praise and thanks be to Him, for on this blessed night special
loving-kindness has encompassed you, and you have been adorned with a
special distinction. God willing, may you be aided and enabled to
educate His servants, refine souls, and attract hearts and inmost
beings. After the prayer was presented before the Countenance, the Most
Great Ocean surged with a wave that this servant is powerless to
describe. At the end, the Tongue of Grandeur uttered these blessed
words: ``Varqā---My glory and loving-kindness rest upon him---turned
toward Us with the face of his heart, the tongue of his heart, the heart
of his heart, and the essence of his heart. He has been set ablaze with
the fire of love for your Lord and has attained the heavenly stream of
His remembrance, the All-Attracting, the Most Bountiful, the Quickener,
the Giver of Life, the Mover, the All-Knowing, the All-Wise.'' Then it
concluded with these blessed words. He said: ``Blessed is he!''

\subsection{139}\label{section-2370}

\begin{otherlanguage}{persian}

در بارهٴ فقرهٴ معلوم آنچه واقع شد صحیح و مقبول بلی امثال این نفوس را
نباید مأیوس نمود و همچنین نباید متابعت اهوی نمود بین بین باید نگاه داشت
عجب در اینکه الی حین جناب حرف سین هم بر مقصود آگاه نه و حال آنکه دو کرّه
در سجن اعظم فائز و هر کرّه مدّتی در جوار ساکن از حقّ جلّ جلاله سائل و آمل که
او را بر مقصود آگاه فرماید وقتی از اوقات فرمودند از حقّ بطلبید او را مؤیّد
فرماید بر آنچه سبب حفظ او و نجات اوست تقرّب در خدمت ملوک چون از حدّ تجاوز
نماید بیم ضرّ است انتهی

\end{otherlanguage}

Dar bāre-ye faqare-ye maʿlūm ānche vāqeʿ shod ṣaḥīḥ va maqbūl. Balī
amthāl-e īn nofūs rā nabāyad maʾyūs namūd va hamchenīn nabāyad
motābaʿat-e ahvā namūd. Bayn bayn bāyad negāh dāsht. ʿAjab dar īnke elā
ḥīn Janāb-e Ḥarf-e Sīn ham bar maqṣūd āgāh na, va ḥāl ānke do karre dar
sejn-e aʿẓam fāʾez va har karre moddatī dar javār sāken. Az Ḥaqq jalla
jalāluhu sāʾel va āmel ke ū rā bar maqṣūd āgāh farmāyad. Vaqtī az awqāt
farmūdand: Az Ḥaqq beṭalabīd ū rā moʾayyad farmāyad bar ānche sabab-e
ḥefẓ-e ū va najāt-e ūst. Taq arrob dar khedmat-e molūk chon az ḥadd
tajāvoz namāyad bīm-e żarr ast. Intahā.

Regarding the matter in question, what took place was correct and
acceptable. Indeed, people such as these should not be discouraged, nor
should their desires be followed. A middle course should be maintained.
It is surprising that even Jināb-i Letter Sīn has not yet understood the
purpose, although he attained the Most Great Prison twice and on each
occasion resided nearby for a time. I ask God, exalted be His glory, and
hope that He will acquaint him with the purpose. On one occasion He
said: ``Ask God to aid him in that which will bring about his protection
and salvation. When seeking closeness in the service of kings exceeds
its bounds, there is a danger of harm.'' Here it ends.

\subsection{140}\label{section-2371}

\begin{otherlanguage}{persian}

این خادم فانی از حقّ میطلبد حضرت معین را مؤیّد فرماید بر اعلاء کلمه و
انتشار امر اگر باینمقام توجّه فرمایند البتّه تأیید الهی و رحمت رحمانی مدد
نماید و آنچه سزاوار است عطا کند ذکر مخدوم مکرّم جناب میرزا ف ض علیه ٩ ٦٦
{[}بهآ الله{]} نموده بودند چندی قبل یک لوح امنع اقدس از سماء فضل مخصوص
ایشان نازل و ارسال شد از هر حرفی از آن لوح مبارک نفحات عنایات لایزالی
متضوّع للّه الحمد مؤیّدند بر استقامت و خدمت و ذکر و ثنا خدمت ایشان سلام و
ثنا و تکبیر و بهاء میرسانم انشآءاللّه آثار فضل ایشان از اشطار و اقطار
بامر مقتدر مختار ظاهر و هویدا گردد

\end{otherlanguage}

Īn khādem-e fānī az Ḥaqq mīṭalabad Ḥażrat-e Moʿīn rā moʾayyad farmāyad
bar eʿlā-ye kalame va enteshār-e amr. Agar be-īn maqām tavajjoh
farmāyand albatte taʾyīd-e elāhī va raḥmat-e raḥmānī madad namāyad va
ānche sazāvār ast ʿaṭā konad. Dhekr-e makhdūm-e mokarram Janāb-e Mīrzā F
Ḍ, ʿalayhi 9 66 {[}bahā Allāh{]}, namūde būdand. Chandī qabl yek lawḥ-e
amnaʿ-e aqdas az samā-ye fażl makhṣūṣ-e īshān nāzel va ersāl shod. Az
har ḥarfī az ān lawḥ-e mobārak nafaḥāt-e ʿenāyāt-e lā yazālī motaḍavveʿ.
Li-llāh al-ḥamd moʾayyadand bar esteqāmat va khedmat va dhekr va thanā.
Khedmat-e īshān salām va thanā va takbīr va bahā mīresānam. In shāʾ
Allāh āthār-e fażl-e īshān az ashṭār va aqṭār be-amr-e Moqtader-e
Mokhtār ẓāher va hovaydā gardad.

This evanescent servant asks God to aid Ḥaḍrat-i Muʿīn in exalting the
Word and spreading the Cause. If he turns his attention to this station,
divine assistance and the mercy of the All-Merciful will assuredly aid
him and bestow what is fitting. You had mentioned the honored gentleman,
Jināb-i Mīrzā F. Ḍ.---upon him be 9 66 {[}the glory of God{]}. Some time
ago, a most exalted, most holy Tablet was revealed from the heaven of
grace especially for him and was sent. From every letter of that blessed
Tablet waft the fragrances of everlasting loving-kindness. Praise be to
God, he is aided in steadfastness, service, remembrance, and praise. I
convey to him greetings, praise, magnification, and glory. God willing,
the effects of his excellence will become manifest and evident
throughout the regions and lands at the command of the All-Powerful, the
Unconstrained.

\subsection{141}\label{section-2372}

\begin{otherlanguage}{persian}

جمیع حاضرین و طائفین تکبیر میرسانند و سلام میگویند و بثنا ذاکرند و در
باب فقرهٴ معهود باید بروایت شیعه بحدیث قبل عمل نمود که فرموده‌اند لا جبر
و لا تفویض بل امر بین الأمرین سبحان‌اللّه الی حین نزد حضرات عنایت و شفقّت و
مرحمت حقّ جلّ جلاله معلوم نه مع آنکه جمیع الواح شاهد و گواهست یسأل الخادم
ربّه بان یؤیّدهم علی عرفان ما اراد انّه هو السّامع المجیب

\end{otherlanguage}

Jamīʿ-e ḥāżerīn va ṭāʾefīn takbīr mīrasānand va salām mīgūyand va
be-thanā dhākerand va dar bāb-e faqare-ye maʿhūd bāyad be-revāyat-e
Shīʿe be-ḥadīth-e qabl ʿamal namūd ke farmūde-and lā jabr wa-lā tafwīḍ
bal amr bayna al-amrayn. Subḥān Allāh, elā ḥīn nazd-e ḥażarāt ʿenāyat va
shafaqat va marḥamat-e Ḥaqq jalla jalāluhu maʿlūm na maʿ-e ānke jamīʿ-e
alvāḥ shāhed va govāh ast. Yasʾalu al-khādim rabbahu bi-an yuʾayyidahum
ʿalā ʿirfān mā arād. Innahu huwa al-Sāmiʿ al-Mujīb.

All those present and those who circle around convey greetings of
glorification, send their salutations, and offer praise. Regarding the
matter in question, one should act according to the previously cited
tradition transmitted by the Shiʿa, which states: ``Neither compulsion
nor delegation, but a matter between the two.'' Glory be to God! To this
day, these gentlemen have not understood the loving-kindness,
compassion, and mercy of God, exalted be His glory, although all the
Tablets bear witness and testify to it. This servant asks his Lord to
help them recognize what He has intended. He is indeed the All-Hearing,
the Responsive.

\subsection{142}\label{section-2373}

\begin{otherlanguage}{arabic}

ذکر قرّهٴ عین الموسوم به ولیّ‌اللّه علیه بهآء اللّه در ساحت اقدس معروض این
آیات باهرات مخصوص آن قرّهٴ عین و سایر قرر عیون علیهم حفظ اللّه و عنایته
نازل قول الرّبّ تعالی و تقدّس

\end{otherlanguage}

Dhekr-e qorrat-e ʿayn al-mawsūm be-Valīyuʾllāh ʿalayhi Bahāʾ Allāh dar
sāḥat-e aqdas maʿrūż. Īn āyāt-e bāherāt makhṣūṣ-e ān qorrat-e ʿayn va
sāyer-e qorar-e ʿoyūn ʿalayhim ḥifẓ Allāh wa-ʿināyatuhu nāzel. Qawl
al-Rabb taʿālā wa-taqaddas:

Mention of that solace of the eyes named Valīyuʾllāh---upon him be the
glory of God---was presented at the Most Holy Court. These resplendent
verses were revealed especially for that solace of the eyes and the
other cherished children---may God's protection and loving-kindness rest
upon them. These are the words of the Lord, exalted and sanctified is
He:

\subsection{143}\label{section-2374}

\begin{otherlanguage}{arabic}

هو النّاطق فیملکوت البیان

\end{otherlanguage}

Huwa al-Nāṭiq fī malakūt al-bayān.

He is the One who speaks in the kingdom of utterance.

\subsection{144}\label{section-2375}

\begin{otherlanguage}{arabic}

طوبی لک بما اقبلت فی اوّل ایّامک الی ربّک ربّ العرش العظیم قل الهی الهی
ترانی مقبلاً الی سمآء ظهورک و ناطقاً بثنائک و آیاتک و معترفاً بما اشرق من
افق ملکوت عرفانک اسألک بأن تجعلنی مؤیّداً علی نصرة امرک و اعلآء کلمتک بحیث
ترتفع رایات امرک فی مدنک و دیارک انّک انت مولی العالم و مربّی الأمم لا اله
الّا انت القویّ القدیر

\end{otherlanguage}

Ṭūbā laka bi-mā aqbalta fī awwal ayyāmika ilā rabbika rabb al-ʿarsh
al-ʿaẓīm. Qul ilāhī ilāhī tarānī muqbilan ilā samāʾ ẓuhūrika wa-nāṭiqan
bi-thanāʾika wa-āyātika wa-muʿtarifan bi-mā ashraqa min ufuq malakūt
ʿirfānika. Asʾaluka bi-an tajʿalanī muʾayyadan ʿalā nuṣrat amrika
wa-iʿlāʾ kalimatika bi-ḥaythu tartafiʿ rāyāt amrika fī mudunika
wa-diyārika. Innaka anta mawlā al-ʿālam wa-murabbī al-umam. Lā ilāha
illā anta al-Qawī al-Qadīr.

Blessed are you for having turned, in your earliest days, toward your
Lord, the Lord of the mighty throne. Say: ``My God, my God! You see me
turning toward the heaven of Your revelation, voicing Your praise and
Your verses, and acknowledging what has shone from the horizon of the
kingdom of Your recognition. I ask You to aid me in championing Your
Cause and exalting Your Word, so that the banners of Your Cause may be
raised in Your cities and Your lands. You are indeed the Lord of the
world and the Educator of the nations. There is no God but You, the
Strong, the All-Powerful.''

\subsection{145}\label{section-2376}

\begin{otherlanguage}{arabic}

و هذا ما نزل لعزیزاللّه علیه بهآء اللّه

\end{otherlanguage}

Wa-hādhā mā nazala li-ʿAzīzuʾllāh ʿalayhi Bahāʾ Allāh.

And this is what was revealed for ʿAzīzuʾllāh---upon him be the glory of
God:

\subsection{146}\label{section-2377}

\begin{otherlanguage}{arabic}

بسمی المشرق من افق سمآء البرهان

\end{otherlanguage}

Bi-ismī al-mushriq min ufuq samāʾ al-burhān.

In My name, shining from the horizon of the heaven of proof.

\subsection{147}\label{section-2378}

\begin{otherlanguage}{arabic}

قد سمع المظلوم ندائک و اجابک بما کان للأراضی الطّیّبة المبارکه رذاذ الرّحمة
و لریاض المعرفة نسمة العنایة و لسمآء البیان شمس الایقان و للعطشان کوثر
الحیوان و لأهل البهآء رائحة الرّحمن اشکر قلمی الأعلی انّه ذکرک بما یکون
باقیاً بدوام ملکوت اللّه ربّ العالمین

\end{otherlanguage}

Qad samiʿa al-Maẓlūm nidāʾaka wa-ajābaka bi-mā kāna li-l-arāḍī
al-ṭayyiba al-mubāraka radhādh al-raḥma wa-li-riyāḍ al-maʿrifa nasmat
al-ʿināya wa-li-samāʾ al-bayān shams al-īqān wa-li-l-ʿaṭshān kawthar
al-ḥayawān wa-li-ahl al-Bahāʾ rāʾiḥat al-Raḥmān. Ushkur qalamī al-aʿlā
innahu dhakaraka bi-mā yakūnu bāqiyan bi-dawām malakūt Allāh rabb
al-ʿālamīn.

The Wronged One has heard your call and answered you with that which is
a gentle rain of mercy for the good and blessed lands, a breeze of
loving-kindness for the gardens of knowledge, a sun of certitude for the
heaven of utterance, a river of life for the thirsty, and the fragrance
of the All-Merciful for the people of Bahāʾ. Give thanks to My Most
Exalted Pen: it has indeed made mention of you in words that will endure
as long as the kingdom of God, the Lord of the worlds.

\subsection{148}\label{section-2379}

\begin{otherlanguage}{arabic}

هذا ما نزّل لروح اللّه علیه بهآء اللّه

\end{otherlanguage}

Hādhā mā nuzzila li-Rūḥuʾllāh ʿalayhi Bahāʾ Allāh.

This is what was revealed for Rūḥuʾllāh---upon him be the glory of God:

\subsection{149}\label{section-2380}

\begin{otherlanguage}{persian}

بنام آنکه پنهانست در حینیکه ظاهر و هویداست

\end{otherlanguage}

Be-nām-e ānke penhān ast dar ḥīnī ke ẓāher va hovaydāst.

In the name of Him who is hidden even while He is manifest and evident.

\subsection{150}\label{section-2381}

\begin{otherlanguage}{persian}

آنچه لسان فؤاد بآن نطق نمود بشرف اصغا فائز قل الهی الهی لک الحمد بما
خلقتنی بکلمتک العلیا و اظهرتنی فی ایّامک یا مولی الوری و ربّ العرش و الثّری
اسألک بالسّفینة الّتی استوی علیها البحر الأعظم و بأمرک الّذی به سخّرت العالم
بأن تجعلنی ساکناً فی ظلّ قباب رحمتک و سمآء فضلک ثمّ قدّر لی ما یقرّبنی الیک
فی کلّ الأحوال و یؤیّدنی علی نصرة امرک انّک انت الغنیّ المتعال لا اله الّا
انت العزیز الفضّال انتهی للّه الحمد

\end{otherlanguage}

Ānche lesān-e foʾād be-ān noṭq namūd be-sharaf-e eṣghā fāʾez. Qul ilāhī
ilāhī laka al-ḥamd bi-mā khalaqtanī bi-kalimatika al-ʿulyā wa-aẓhartanī
fī ayyāmika yā mawlā al-warā wa-rabb al-ʿarsh wa-l-tharā. Asʾaluka
bi-l-safīna allatī istawā ʿalayhā al-baḥr al-aʿẓam wa-bi-amrika alladhī
bihi sakhkharta al-ʿālam bi-an tajʿalanī sākinan fī ẓill qibāb raḥmatika
wa-samāʾ faḍlika. Thumma qaddir lī mā yuqarribunī ilayka fī kull
al-aḥwāl wa-yuʾayyidunī ʿalā nuṣrat amrika. Innaka anta al-Ghanī
al-Mutaʿāl. Lā ilāha illā anta al-ʿAzīz al-Faḍḍāl. Intahā. Li-llāh
al-ḥamd.

What the tongue of your heart uttered attained the honor of being heard.
Say: ``My God, my God! Praise be to You for having created me through
Your most exalted Word and brought me into being in Your days, O Lord of
all creation and Lord of the throne and the earth. I ask You, by the Ark
upon which the Most Great Ocean has established itself, and by Your
Cause through which You have subdued the world, to let me dwell beneath
the shelter of the domes of Your mercy and the heaven of Your grace.
Then ordain for me what will draw me nearer to You in every circumstance
and aid me in championing Your Cause. You are indeed the
Self-Sufficient, the Most Exalted. There is no God but You, the
Almighty, the All-Bountiful.'' Here it ends. Praise be to God.

\subsection{151}\label{section-2382}

\begin{otherlanguage}{persian}

قرر عیون بعنایات مقصود عالمیان فائز گشتند از برای هر یک نعمت باقیه و
مائدهٴ معنویّه از سمآء فضل نازل هنیئاً لهم در حین عرض تبسّم فرمودند و بحر
عنایت به موجی بدیع و اوجی عظیم ظاهر و هویدا له الحمد و الثّناء و له الشّکر
و البهآء و مولود جدید علیه بهآء اللّه العزیز الحمید که بطراز اسم بدیع‌اللّه
مزیّن در آنحین بشرف امضاء مشرّف نعیماً له این خادم از حقّ سائل و آمل که هر
یک را تأیید فرماید تأییدیکه از آیات منزله واضح و مشهود یا ورقا علیک
بهائی و عنایتی للّه الحمد آنجناب خود باین شرافت کبری فائزید و باین موهبت
عظمی مزیّن یشهد بذلک لسان عظمتی و لسان فضلی و لسان عنایتی و لسان بیانی
البدیع انتهی

\end{otherlanguage}

Qorar-e ʿoyūn be-ʿenāyāt-e Maqṣūd-e ʿālamiyān fāʾez gashtand. Az barā-ye
har yek neʿmat-e bāqiye va māʾede-ye maʿnaviyye az samāʾ-e fażl nāzel.
Hanīʾan lahum. Dar ḥīn-e ʿarż tabassom farmūdand va baḥr-e ʿenāyat
be-mowjī badīʿ va owjī ʿaẓīm ẓāher va hovaydā. Lahu al-ḥamd wa-l-thanāʾ
wa-lahu al-shukr wa-l-bahāʾ. Va mowlūd-e jadīd ʿalayhi Bahāʾ Allāh
al-ʿAzīz al-Ḥamīd ke be-ṭerāz-e esm-e Badīʿuʾllāh mozayyan dar ān ḥīn
be-sharaf-e emżā mosharraf. Naʿīman lahu. Īn khādem az Ḥaqq sāʾel va
āmel ke har yek rā taʾyīd farmāyad taʾyīdī ke az āyāt-e monzale vāżeḥ va
mashhūd. Yā Warqāʾ ʿalayka bahāʾī wa-ʿināyatī. Li-llāh al-ḥamd ān janāb
khod be-īn sharāfat-e kobrā fāʾezīd va be-īn mowhebat-e ʿoẓmā mozayyan.
Yashhadu bi-dhālika lisān ʿaẓamatī wa-lisān faḍlī wa-lisān ʿināyatī
wa-lisān bayānī al-badīʿ. Intahā.

The cherished children have received the favors of the Desire of all the
worlds. For each of them, an enduring blessing and a spiritual feast
have descended from the heaven of grace. May it bring them joy! As their
names were presented, He smiled, and the ocean of loving-kindness rose
in a wondrous wave to a mighty height. His are all praise and
glorification, and His are all thanksgiving and glory. At that time, the
newborn---upon him be the glory of God, the Almighty, the
Praiseworthy---who is adorned with the name Badīʿuʾllāh, was honored
with approval. May bliss be his! This servant asks God and hopes that He
will strengthen each one with the aid that is clear and evident in the
revealed verses. ``O Warqāʾ, upon you be My glory and My
loving-kindness.'' Praise be to God, you yourself have attained this
supreme distinction and are adorned with this most great gift. ``To this
testify the tongue of My grandeur, the tongue of My grace, the tongue of
My loving-kindness, and the tongue of My wondrous utterance.'' Here it
ends.

\subsection{152}\label{section-2383}

\begin{otherlanguage}{arabic}

هذا فضل آخر و نعمة اخری و رحمة اخری و موهبة اخری من لدی اللّه ربّنا و ربّ
من فی السّموات و الأرضین

\end{otherlanguage}

Hādhā faḍl ākhar wa-niʿma ukhrā wa-raḥma ukhrā wa-mawhiba ukhrā min ladā
Allāh rabbinā wa-rabb man fī al-samāwāt wa-l-arḍīn.

This is yet another grace, another blessing, another mercy, and another
gift from God, our Lord and the Lord of all who are in the heavens and
the earth.

\subsection{153}\label{section-2384}

\begin{otherlanguage}{arabic}

اینکه در بارهٴ جناب آقا میرزا آقا ش علیه بهآء اللّه و ارادهٴ توجّه ایشان
بانوار وجه محبوب و اشتعالشان بنار فراق مرقوم داشتند اینمراتب بسمع من
خلقت الآذان لاصغآء ندائه فائز هذا ما نطق به لسان العظمة فی الجواب قول
الرّبّ تعالی و تقدّس

\end{otherlanguage}

Īnke dar bāre-ye janāb-e Āqā Mīrzā Āqā Sh. ʿalayhi Bahāʾ Allāh va
erāde-ye tavajjoh-e īshān be-anvār-e vajh-e Maḥbūb va eshteʿāl-e shān
be-nār-e ferāq marqūm dāshtand īn marāteb be-samʿ-e man khuliqat
al-ādhān li-iṣghāʾ nidāʾihi fāʾez. Hādhā mā naṭaqa bihi lisān al-ʿaẓama
fī al-jawāb. Qawl al-Rabb taʿālā wa-taqaddas:

What you wrote concerning Jināb-i Āqā Mīrzā Āqā Sh.---upon him be the
glory of God---his wish to turn toward the lights of the Beloved's
countenance, and his burning in the fire of separation, attained the
hearing of Him for whose call ears were created to listen. This is what
the Tongue of Grandeur uttered in reply. These are the words of the
Lord, exalted and sanctified is He:

\subsection{154}\label{section-2385}

\begin{otherlanguage}{arabic}

هو الشّاهد الخبیر

\end{otherlanguage}

Huwa al-Shāhid al-Khabīr.

He is the Witness, the All-Aware.

\subsection{155}\label{section-2386}

\begin{otherlanguage}{arabic}

قد سمعنا ندآئک و رأینا توجّهک و اقبالک و اشتعالک فی هجر مولیک انّا ذکرناک
من قبل بما فاحت به نفحة الرّحمن فی الامکان نسأل اللّه ان یؤیّدک و یوفّقک علی
ذکره و ثنائه و خدمة امره انّه هو المقتدر المهیمن القیّوم انّک اذا شربت رحیق
الوحی من کأس بیانی و وجدت عرف الرّحمن من کلماتی

\end{otherlanguage}

Qad samiʿnā nidāʾaka wa-raʾaynā tawajjuhaka wa-iqbālaka wa-ishtiʿālaka
fī hajr mawlāka. Innā dhakarnāka min qabl bi-mā fāḥat bihi nafḥat
al-Raḥmān fī al-imkān. Nasʾalu Allāh an yuʾayyidaka wa-yuwaffiqaka ʿalā
dhikrihi wa-thanāʾihi wa-khidmat amrihi. Innahu huwa al-Muqtadir
al-Muhaymin al-Qayyūm. Innaka idhā sharibta raḥīq al-waḥy min kaʾs
bayānī wa-wajadta ʿarf al-Raḥmān min kalimātī

We have heard your call and seen your turning toward Us, your devotion,
and your burning in separation from your Lord. We have previously made
mention of you in words through which the breeze of the All-Merciful has
wafted throughout existence. We ask God to aid you and enable you to
remember and praise Him and to serve His Cause. He is indeed the
All-Powerful, the All-Protecting, the Self-Subsisting. When you have
drunk the choice wine of revelation from the cup of My utterance and
perceived the fragrance of the All-Merciful from My words,

\subsection{156}\label{section-2387}

\begin{otherlanguage}{arabic}

قل الهی الهی هجرک اهلکنی و فراقک احرقنی و ظهورک حیّرنی و آیاتک اشعلتنی و
بیّناتک جذبتنی اسألک بالکلمة الّتی بها سرع المقرّبون الی مقرّ الفدآء بأن
تکتب لی من قلمک الأعلی اجر لقائک و الحضور امام وجهک و القیام لدی باب
عظمتک ای ربّ ترانی منجذباً من نفحات وحیک و طائراً فی هوآء حبّک اسألک بأمطار
فجر ظهورک و انوار وجهک بأن تجعلنی فیکلّ الأحوال متمسّکاً بحبل فضلک و عاملاً
بما امرتنی به فی کتابک انّک انت المقتدر علی ما تشآء لا اله الّا انت القویّ
القدیر

\end{otherlanguage}

Qul ilāhī ilāhī hajruka ahlakanī wa-firāquka aḥraqanī wa-ẓuhūruka ḥay
yaranī wa-āyātuka ashʿalatanī wa-bayyinātuka jadhabatnī. Asʾaluka
bi-l-kalima allatī bihā saraʿa al-muqarrabūn ilā maqarr al-fidāʾ bi-an
taktuba lī min qalamika al-aʿlā ajr liqāʾika wa-l-ḥuḍūr amāma wajhika
wa-l-qiyām ladā bāb ʿaẓamatika. Ay rabb tarānī munjadhiban min nafaḥāt
waḥyika wa-ṭāʾiran fī hawāʾ ḥubbika. Asʾaluka bi-amṭār fajr ẓuhūrika
wa-anwār wajhika bi-an tajʿalanī fī kull al-aḥwāl mutamassikan bi-ḥabl
faḍlika wa-ʿāmilan bi-mā amartanī bihi fī kitābika. Innaka anta
al-Muqtadir ʿalā mā tashāʾ. Lā ilāha illā anta al-Qawī al-Qadīr.

say: ``My God, my God! Separation from You has brought me to ruin, and
distance from You has consumed me. Your revelation has bewildered me,
Your verses have set me aflame, and Your clear proofs have drawn me to
You. I ask You, by the Word through which those who are near to You
hastened to the place of sacrifice, to inscribe for me, with Your Most
Exalted Pen, the reward of meeting You, of being present before Your
face, and of standing at the door of Your grandeur. O Lord, You see me
enraptured by the breezes of Your revelation and soaring in the
atmosphere of Your love. I ask You, by the rains of the dawn of Your
revelation and the lights of Your countenance, to enable me in every
circumstance to hold fast to the cord of Your grace and to act according
to what You have commanded me in Your Book. You indeed have power over
whatever You will. There is no God but You, the Strong, the
All-Powerful.''

\subsection{157}\label{section-2388}

\begin{otherlanguage}{persian}

حقّ اولیای خود را دوست داشته و دارد و از برای لقا و ظهورات عنایت و مرحمت
آمده ولکن ظلم ظالمین و غفلت غافلین حایل شده چندیست ارض سجن مضطرب مشاهده
میشود چه که نفسی از اهل بغضا مدّتی بر سریر حکومت جالس اگرچه بسقر راجع
یعنی عزل شد ولکن الی حین ارض ساکن نشده لذا در اظهار کلمهٴ اذن توقّف
نمودیم انّه هو الفضّال الکریم و هو الغفور الرّحیم از او میطلبیم آنچه را که
خیر دنیا و آخرت در آنست انّه علی کلّشیء قدیر و بالاجابة جدیر انتهی

\end{otherlanguage}

Ḥaqq owliyā-ye khod rā dūst dāshte va dārad va az barā-ye leqā va
ẓohūrāt-e ʿenāyat va marḥamat āmade va-līkan ẓolm-e ẓālemīn va ghaflat-e
ghāfelīn ḥāʾel shode. Chandīst arż-e sejn możṭareb moshāhede mīshavad
che ke nafsī az ahl-e baghżā moddatī bar sarīr-e ḥokūmat jāles. Agarche
be-Saqar rājeʿ yaʿnī ʿazl shod va-līkan elā ḥīn arż sāken nashode. Ledhā
dar eẓhār-e kalame-ye edhn tavaqqof namūdīm. Innahu huwa al-Faḍḍāl
al-Karīm wa-huwa al-Ghafūr al-Raḥīm. Az ū mīṭalabīm ānche rā ke khayr-e
donyā va ākherat dar ān ast. Innahu ʿalā kull shayʾ qadīr wa-bi-l-ijāba
jadīr. Intahā.

God has loved and continues to love His loved ones, and has come so that
they may meet Him and witness the manifestations of His loving-kindness
and mercy. But the oppression of the oppressors and the heedlessness of
the heedless have intervened. For some time, the land of the prison has
been seen to be in turmoil, for a person from among the people of hatred
occupied the seat of government for a while. Although he returned to
hell---that is, he was dismissed---the land has not yet become calm. We
have therefore refrained from uttering the word of permission. He is
indeed the All-Bountiful, the Generous, and He is the Ever-Forgiving,
the Merciful. We ask Him for that in which lies the good of this world
and the next. He indeed has power over all things and is worthy to
answer prayers. Here it ends.

\subsection{158}\label{section-2389}

\begin{otherlanguage}{persian}

للّه الحمد فائز شدند بآنچه که شبه و مثل نداشته و ندارد تفصیل این ارض
البتّه الی حین بسمع آنجناب رسیده اینعبد هم خدمت ایشان تکبیر میرساند و از
حقّ جلّ جلاله خیر دنیا و آخرت میطلبد اینکه مجدّد در بارهٴ جناب حاجی محمّد
علی علیه ٩ ٦٦ {[}بهآ الله{]} مرقوم داشتند یوم ورود دستخطّ آنمحبوب آنچه
اراده نمودند در ساحت اقدس عرض شد و بشرف امضا فائز هنیئاً له و جناب محبوب
فؤاد امین علیه بهآء اللّه هم در ساحت اقدس مطالب ایشانرا عرض نمودند و امضا
را باذن خود شنیدند و همچنین لوحی هم مخصوص ایشان از قبل نازل و ارسال شد
یسأل الخادم ربّه بأن یؤیّد اولیآئه علی ما یحبّ و یرضی و یقدّر لهم خیر ما نزل
فی کتابه المبین

\end{otherlanguage}

Li-llāh al-ḥamd fāʾez shodand be-ānche ke shebh va methl nadāshte va
nadārad. Tafṣīl-e īn arż albatte elā ḥīn be-samʿ-e ān janāb resīde. Īn
ʿabd ham khedmat-e īshān takbīr mīrasānad va az Ḥaqq jalla jalāluhu
khayr-e donyā va ākherat mīṭalabad. Īnke mojaddad dar bāre-ye janāb-e
Ḥājī Moḥammad ʿAlī ʿalayhi 9 66 {[}Bahāʾ Allāh{]} marqūm dāshtand yowm-e
vorūd-e dastkhaṭṭ-e ān maḥbūb ānche erāde namūdand dar sāḥat-e aqdas
ʿarż shod va be-sharaf-e emżā fāʾez. Hanīʾan lahu. Va janāb-e maḥbūb-e
foʾād Amīn ʿalayhi Bahāʾ Allāh ham dar sāḥat-e aqdas maṭāleb-e īshān rā
ʿarż namūdand va emżā rā be-ozon-e khod shenīdand va hamchenīn lowḥī ham
makhṣūṣ-e īshān az qabl nāzel va ersāl shod. Yasʾalu al-khādim rabbahu
bi-an yuʾayyida awliyāʾahu ʿalā mā yuḥibbu wa-yarḍā wa-yuqaddira lahum
khayr mā nazala fī kitābihi al-mubīn.

Praise be to God, he has attained what has never had, and does not have,
any likeness or equal. The details of events in this land have surely
reached your hearing by now. This servant also conveys greetings of
glorification to him and asks God, exalted be His glory, to grant him
the good of this world and the next. Regarding what you wrote again
concerning Jināb-i Ḥājī Muḥammad ʿAlī---upon him be 9 66 {[}the glory of
God{]}---on the day the letter from you, that beloved friend, arrived,
what he desired was presented at the Most Holy Court and received the
honor of approval. May it bring him joy! Jināb-i Amīn, the beloved of
the heart---upon him be the glory of God---also presented his requests
at the Most Holy Court and heard the approval with his own ears.
Likewise, a Tablet had previously been revealed especially for him and
sent. This servant asks his Lord to aid His loved ones in what He loves
and approves, and to ordain for them the best of what has been revealed
in His lucid Book.

\subsection{159}\label{section-2390}

\begin{otherlanguage}{persian}

و اینکه در ذکر الوهیّت مرقوم فرمودند که شخصی را از شنیدن آن پریشانی دست
داد بعد از عرض اینفقر اینکلمات عالیات از مشرق اراده ظاهر و

\end{otherlanguage}

Va īnke dar dhekr-e olūhiyyat marqūm farmūdand ke shakhṣī rā az
shenīdan-e ān parīshānī dast dād baʿd az ʿarż-e īn faqre īn kalemāt-e
ʿāliyāt az mashreq-e erāde ẓāher va

And regarding what you wrote concerning divinity, that someone had
become distressed on hearing it mentioned, after this matter was
presented, these exalted words appeared from the dayspring of the Will
and

\subsection{160}\label{section-2391}

\begin{otherlanguage}{persian}

مشرق قوله تبارک و تعالی یا ورقآء علیک بهائی و عنایتی ناس ضعیفند الی حین
قلم اعلی بما ینبغی ذکر ننموده مکلّم طور مع آنکه آیاتش عالم را احاطه نموده
خود را صامت مشاهده مینماید و الی حین جهل اهل فرقان بتمامه ذکر نشده
لعمراللّه از کأس توحید محرومند و از قدح تفرید ممنوع الی حین مقام توحید
حقیقی مستور و این ستر فضلی است بزرگ از برای بعضی از نفوس که اقبال
نموده‌اند اگرچه نفحهٴ اینمقام از قلم اعلی متضوّع مقصود آنکه آذان مستعدّه
لاصغآء ندآء اللّه و ابصار حدیده لمشاهدة ما ظهر من خزائن القلم قلیل مشاهده
میشود آنجناب باید مدارا نمایند و باذکاری ناطق شوند که سبب توجّه و اقبال و
انجذابست چه که مراتب و مقامات اولیای فرقان و مراتب آنچه سزاوار ظهور است
اگر بعضی مستور ماند اقرب بتقوی است حال محبّت و اقبال کافیست و بعد از
اقبال و توجّه اکثری بحبل بیان ارتقا مینمایند و بقدر استعداد بر مراتب واقف
میشوند قدّرنا لک خیراً کثیراً فی الکتاب انتهی

\end{otherlanguage}

Moshreq. Qawluhu tabāraka wa-taʿālā: Yā Warqāʾ ʿalayka bahāʾī
wa-ʿināyatī. Nās żaʿīfand. Elā ḥīn Qalam-e Aʿlā bi-mā yanbaghī dhekr
nanamūde. Mokallem-e Ṭūr maʿ-e ānke āyātash ʿālam rā eḥāṭe namūde khod
rā ṣāmet moshāhede mīnamāyad va elā ḥīn jahl-e ahl-e Forqān be-tamāme
dhekr nashode. La-ʿamru Allāh az kaʾs-e towḥīd maḥrūmand va az qadaḥ-e
tafrīd mamnūʿ. Elā ḥīn maqām-e towḥīd-e ḥaqīqī mastūr va īn satr fażlī
ast bozorg az barā-ye baʿżī az nofūs ke eqbāl namūde-and agarche
nafḥe-ye īn maqām az Qalam-e Aʿlā motaḍavveʿ. Maqṣūd ānke ādhān-e
mostaʿedde li-iṣghāʾ nidāʾ Allāh va abṣār-e ḥadīde li-mushāhadat mā
ẓahara min khazāʾin al-qalam qalīl moshāhede mīshavad. Ān janāb bāyad
modārā namāyand va be-adhkārī nāṭeq shavand ke sabab-e tavajjoh va eqbāl
va enjedhāb ast che ke marāteb va maqāmāt-e owliyā-ye Forqān va
marāteb-e ānche sazāvār-e ẓohūr ast agar baʿżī mastūr mānad aqrab
be-taqvā ast. Ḥāl maḥabbat va eqbāl kāfīst va baʿd az eqbāl va tavajjoh
aktharī be-ḥabl-e bayān erteqā mīnamāyand va be-qadr-e esteʿdād bar
marāteb vāqef mīshavand. Qaddarnā laka khayran kathīran fī al-kitāb.
Intahā.

shone forth. His words, blessed and exalted is He: ``O Warqāʾ, upon you
be My glory and My loving-kindness. People are weak. To this day, the
Most Exalted Pen has not spoken as is befitting. Although His verses
have encompassed the world, He who spoke on Sinai regards Himself as
silent, and the ignorance of the people of the Qurʾān has not yet been
fully described. By the life of God, they are deprived of the cup of
divine unity and barred from the chalice of singleness. To this day, the
station of true divine unity remains concealed. This concealment is a
great favor to some of the souls who have turned toward Him, although
the fragrance of this station wafts from the Most Exalted Pen. The point
is that few ears are found ready to hear the call of God, and few eyes
keen enough to behold what has appeared from the treasuries of the Pen.
You should show forbearance and speak of things that will cause people
to turn toward Him, embrace His Cause, and be attracted to Him. For if
some aspects of the ranks and stations of the holy ones of the Qurʾān,
and of the degrees of what is fitting for this Revelation, remain
concealed, this is nearer to piety. For now, love and acceptance
suffice. After accepting and turning toward Him, most ascend by the cord
of utterance and become acquainted with these degrees according to their
capacity. We have ordained abundant good for you in the Book.'' Here it
ends.

\subsection{161}\label{section-2392}

\begin{otherlanguage}{arabic}

یا محبوب فؤادی اکثری رضیعند شیر میطلبند قابل فوق آن نبوده و نیستند نسأل
اللّه تعالی ان یؤیّد عباده علی ذکره و ثنآئه و الاقبال الیه و التّوکّل علیه
انّه هو السّامع المجیب و هو الفرد الواحد العزیز الحمید

\end{otherlanguage}

Yā maḥbūb-e foʾādī aktharī rażīʿand shīr mīṭalaband qābel-e fowq-e ān
nabūde va nīstand. Nasʾalu Allāh taʿālā an yuʾayyida ʿibādahu ʿalā
dhikrihi wa-thanāʾihi wa-l-iqbāl ilayhi wa-l-tawakkul ʿalayhi. Innahu
huwa al-Sāmiʿ al-Mujīb wa-huwa al-Fard al-Wāḥid al-ʿAzīz al-Ḥamīd.

O beloved of my heart, most are nursing infants: they seek milk and have
never been, nor are they now, capable of taking anything beyond it. We
ask God, exalted is He, to aid His servants in remembering and praising
Him, turning toward Him, and placing their trust in Him. He is indeed
the All-Hearing, the Responsive, and He is the Unique, the One, the
Almighty, the Praiseworthy.

\subsection{162}\label{section-2393}

\begin{otherlanguage}{persian}

عرایضیکه بساحت اقدس ارسال شد جواب نازل از حقّ میطلبم از اشراقات انوار
آفتاب حقیقی که از افق لوح مشرقست اولیای آن دیار طرّاً منوّر شوند و از نار
بدیعهٴ منیعه مشتعل گردند و به اتّحاد و اتّفاق قیام نمایند و با صدور منیرهٴ
مطهّرهٴ از ضغینه و بغضا باصلاح عالم پردازند تا نور اتّفاق عالم را اخذ
نماید و عداوتهای مکنونهٴ در صدور ساکن شود و زایل گردد

\end{otherlanguage}

ʿArāyeżī ke be sāḥat-e aqdas ersāl shod javāb nāzel. Az Ḥaqq mīṭalabam
az eshrāqāt-e anvār-e āftāb-e ḥaqīqī ke az ofoq-e lawḥ moshreq ast
awliyā-ye ān diyār ṭorran monavvar shavand va az nār-e badīʿe-ye manīʿe
moshtaʿel gardand va be etteḥād va ettefāq qiyām namāyand va bā ṣodūr-e
monīre-ye moṭahhare az żaghīne va baghżā be eṣlāḥ-e ʿālam pardāzand tā
nūr-e ettefāq ʿālam rā akhdh namāyad va ʿadāvat-hā-ye maknūne dar ṣodūr
sāken shavad va zāyel gardad.

Replies have been revealed to the petitions sent to the Most Holy Court.
I ask God that all the friends in that land may be illumined by the
radiance of the light of the Sun of Truth shining from the horizon of
the Tablet, become ablaze with the wondrous, exalted fire, and arise in
unity and concord. May they devote themselves to the betterment of the
world with radiant hearts purified of rancor and hatred, so that the
light of concord may encompass the world and the enmities concealed in
human hearts may subside and disappear.

\subsection{163}\label{section-2394}

\begin{otherlanguage}{arabic}

امید هست که اهل بها روح العالم فداهم کمر همّت را محکم نمایند و فساد و
نزاع و جدال و سفک دماء را از عالم بردارند تا عالم بیچاره نفس راحتی برآرد
من این را میگویم ولکن هیهات هیهات میشنوم چه که بعضی از نفوس امر اللّه را
سهل شمرده‌اند یلعبون و لا یعرفون و یقولون ما لا یشعرون

\end{otherlanguage}

Omīd hast ke ahl-e Bahā rūḥu al-ʿālam fidāhum kamar-e hemmat rā moḥkam
namāyand va fasād va nezāʿ va jedāl va safk-e demā rā az ʿālam bardārand
tā ʿālam-e bīchāre nafas-e rāḥatī barārad. Man īn rā mīgūyam va lāken
hayhāt hayhāt mīshenavam che ke baʿżī az nofūs amr-e Allāh rā sahl
shemorde-and. Yalʿabūna wa-lā yaʿrifūna wa-yaqūlūna mā lā yashʿurūn.

It is hoped that the people of Bahā---may the soul of the world be
sacrificed for them---will firmly brace themselves for action and rid
the world of corruption, conflict, contention, and bloodshed, so that
this hapless world may breathe a sigh of relief. I say this, but I hear,
``Far from it, far from it!'' For some souls have treated the Cause of
God lightly. They play and do not understand, and say things of which
they have no awareness.

\subsection{164}\label{section-2395}

\begin{otherlanguage}{arabic}

در جمیع احوال این خادم فانی تأیید میطلبد و توفیق میخواهد انّ ربّنا الرّحمن
هو المشفق الکریم و هو العلیم الحکیم

\end{otherlanguage}

Dar jamīʿ-e aḥvāl īn khādem-e fānī taʾyīd mīṭalabad va tawfīq mīkhāhad.
Inna rabbanā al-Raḥmāna huwa al-Mushfiqu al-Karīm wa-huwa al-ʿAlīmu
al-Ḥakīm.

In all circumstances, this mortal servant seeks assistance and asks for
success. Our Lord, the All-Merciful, is indeed the Compassionate, the
Generous, and He is the All-Knowing, the All-Wise.

\subsection{165}\label{section-2396}

\begin{otherlanguage}{persian}

و امّا فقرهٴ خطوط دو قسم حضرت غصن‌اللّه الأکبر روحی و ذاتی لتراب قدومه
الفدآء مرقوم فرموده‌اند بسیار خوبست تفصیل را خود حضرتشان مرقوم میفرمایند
اینکه مرقوم فرمودند که یکی از دوستان سؤال از غیبت نموده اینمراتب در
الواح الهی بمثابهٴ آفتاب ظاهر و لائح این ایّام باید کلّ بنار کلمه مشتعل
شوند و بخدمت امر مشغول وقت را غنیمت شمارند و به حکمت و بیان اهل امکانرا
آگاه کنند هر نفسی الیوم باعتراف بما انزله اللّه فی الکتاب موفّق گشت او از
اهل فردوس اعلی مذکور و محسوب

\end{otherlanguage}

Va ammā faqare-ye khoṭūṭ-e do qesm Ḥażrat-e Ghoṣn Allāh al-Akbar rūḥī
wa-dhātī li-turābi qudūmihi al-fidāʾ marqūm farmūde-and besyār khūb ast.
Tafṣīl rā khod-e Ḥażrat-eshān marqūm mīfarmāyand. Īnke marqūm farmūdand
ke yekī az dūstān soʾāl az ghībat namūde īn marāteb dar alvāḥ-e elāhī be
mathābe-ye āftāb ẓāher va lāʾeḥ. Īn ayyām bāyad koll be nār-e kalame
moshtaʿel shavand va be khedmat-e amr mashghūl. Vaqt rā ghanīmat
shomārand va be ḥekmat va bayān ahl-e emkān rā āgāh konand. Har nafsī
al-yawm be eʿterāf be mā anzalahu Allāhu fī al-kitāb movaffaq gasht ū az
ahl-e ferdaws-e aʿlā madhkūr va maḥsūb.

As for the matter of the two kinds of script, His Holiness the Greater
Branch of God---may my soul and my being be sacrificed for the dust
beneath His feet---has written concerning it. It is very good. He
Himself will write out the details. Regarding your statement that one of
the friends has asked about backbiting, these matters are as manifest
and clear as the sun in the divine Tablets. In these days, all should be
ablaze with the fire of the Word and engaged in service to the Cause.
They should seize the opportunity the time affords and, through wisdom
and utterance, awaken the people of the world. Every soul who today
succeeds in acknowledging what God has sent down in the Book is named
and counted among the inhabitants of the highest Paradise.

\subsection{166}\label{section-2397}

\begin{otherlanguage}{persian}

امروز سیّد ایّام است و مالک ظهورات بر عرش بیان مستوی طوبی للفائزین ایّام
مظاهر الهی در مقامی بیوم اللّه مذکور ولکن این ایّام در جمیع کتب مخصوص است
بظهور غیب مکنون و سرّ مخزون الّذی ینطق بأعلی النّدآء امام وجوه العالم در
کتب قبل و فرقان عظمت این ایّام مذکور و مسطور است

\end{otherlanguage}

Emrūz sayyed-e ayyām ast va Mālek-e ẓohūrāt bar ʿarsh-e bayān mostavī.
Ṭūbā li-l-fāʾizīn. Ayyām-e maẓāher-e elāhī dar maqāmī be yawm Allāh
madhkūr va lāken īn ayyām dar jamīʿ-e kotob makhṣūṣ ast be ẓohūr-e
ghayb-e maknūn va serr-e makhzūn alladhī yanṭiqu bi-aʿlā al-nidāʾ amāma
wujūhi al-ʿālam. Dar kotob-e qabl va Forqān ʿaẓamat-e īn ayyām madhkūr
va masṭūr ast.

Today is the lord of days, and the Lord of Revelations is established
upon the throne of utterance. Blessed are those who have attained! In
one sense, the days of the Manifestations of God are called the Day of
God, but in all the Books these days are specifically associated with
the revelation of the hidden Unseen and the treasured Mystery, who
speaks with the loudest call before the faces of the world. The
greatness of these days is mentioned and recorded in the former Books
and in the Qurʾān.

\subsection{167}\label{section-2398}

\begin{otherlanguage}{persian}

حضرت مبشّر یعنی نقطهٴ اولی روح ما سواه فداه در ایّام توجّه به ماکو بر حسب
ظاهر مستوراً عن الکلّ بشرف لقا فائز و جناب قدّوس علیه بهائی و عنایتی مکرّر
فائز در سفر ارض خا مخصوص او را طلبیدیم و اشهر معدودات در حضور بلقا فائز
ایشان و جناب باب اعظم را حضرت نقطه روح ما سواه فداه باسم اوّل و آخر
نامیده و همچنین باوّل وارد و آخر نازل تعبیر فرموده و بشارت اینظهور را بهر
یک داده و بجناب سیّاح علیه بهائی مخصوص فرموده‌اند برو شاید بیابی

\end{otherlanguage}

Ḥażrat-e Mobashsher yaʿnī Noqṭe-ye Ūlā rūḥu mā siwāhu fidāh dar ayyām-e
tavajjoh be Mākū bar ḥasab-e ẓāher mastūran ʿan al-koll be sharaf-e leqā
fāʾez va Janāb-e Qoddūs ʿalayhi bahāʾī wa-ʿināyatī mokarrar fāʾez. Dar
safar-e arż-e Khā makhṣūṣ ū rā ṭalabīdīm va ashhor-e maʿdūdāt dar ḥożūr
be leqā fāʾez. Īshān va Janāb-e Bāb-e Aʿẓam rā Ḥażrat-e Noqṭe rūḥu mā
siwāhu fidāh be esm-e avval va ākhar nāmīde va hamchenīn be avval vāred
va ākhar nāzel taʿbīr farmūde va beshārat-e īn ẓohūr rā be har yek dāde
va be Janāb-e Sayyāḥ ʿalayhi bahāʾī makhṣūṣ farmūde-and boro shāyad
beyābī.

The Herald, namely the Primal Point---may the soul of all besides Him be
sacrificed for Him---attained the honor of meeting, outwardly concealed
from everyone, during the days of His journey toward Mākū. And
Quddūs---upon him be My glory and My loving care---attained repeatedly.
During the journey to the land of Khā, We expressly summoned him, and
for several months he was in Our presence and attained the meeting. The
Point---may the soul of all besides Him be sacrificed for Him---named
him and the Most Great Bāb the First and the Last, and also referred to
them as the first to arrive and the last to alight. He gave each of them
the glad tidings of this Revelation, and specifically told Sayyāḥ---upon
him be My glory---``Go; perhaps you will find.''

\subsection{168}\label{section-2399}

\begin{otherlanguage}{persian}

مظاهر کلّیّه قبل از ظهور آگاه بوده‌اند ایشانند مطالع علم الهی نسأل اللّه ان
یؤیّد اولیائه علی عرفان ما انزله فی کتابه المحکم المبین بعضی اذکار سبب
احتجاب میشود و همچنین علّت اخماد نار سدره در وجود البهآء علی من اقتصر
النّظر علی التّوجّه الی الأفق الأعلی منقطعاً عن الظّنون و الأوهام یا محبوب
فؤادی فی‌الحقیقه بعضی از اذکار سبب غفلت و کسالت میشود مخصوص اذکار ما لم
یرض به اللّه امید هست کلّ از رحیق بیان رحمن بنوشند و در انجمن عالم بطراز
امتیاز ظاهر شوند البتّه حزب اللّه باید از دونش ممتاز باشند خدمت اولیای آن
ارض تکبیر و سلام از قبل و بعد عرض شده و میشود و این خادم فانی از حقّ باقی
سائل و آمل که جمیع را تأیید فرماید بر آنچه سزاوار است

\end{otherlanguage}

Maẓāher-e kolliyye qabl az ẓohūr āgāh būde-and. Īshān-and maṭāleʿ-e
ʿelm-e elāhī. Nasʾalu Allāha an yuʾayyida awliyāʾahu ʿalā ʿirfāni mā
anzalahu fī kitābihi al-muḥkami al-mubīn. Baʿżī adhkār sabab-e eḥtejāb
mīshavad va hamchenīn ʿellat-e ekhmād-e nār-e sedre dar vojūd. Al-bahāʾu
ʿalā man iqtaṣara al-naẓara ʿalā al-tawajjuh ilā al-ufuq al-aʿlā
munqaṭiʿan ʿan al-ẓunūn wa-l-awhām. Yā maḥbūb-e foʾādī fī al-ḥaqīqe
baʿżī az adhkār sabab-e ghaflat va kesālat mīshavad makhṣūṣ adhkār-e mā
lam yarḍa bihi Allāh. Omīd hast koll az raḥīq-e bayān-e Raḥmān benūshand
va dar anjoman-e ʿālam be ṭerāz-e emtiyāz ẓāher shavand. Albatte ḥezb-e
Allāh bāyad az dūn-ash momtāz bāshand. Khedmat-e awliyā-ye ān arż takbīr
va salām az qabl va baʿd ʿarż shode va mīshavad va īn khādem-e fānī az
Ḥaqq-e bāqī sāʾel va āmel ke jamīʿ rā taʾyīd farmāyad bar ānche sazāvār
ast.

The universal Manifestations were aware before their appearance. They
are the dawning places of divine knowledge. We ask God to assist His
friends in recognizing what He has sent down in His firm and lucid Book.
Certain discussions cause veiling and also extinguish the fire of the
divine Lote-Tree within one's being. Glory rest upon whoever confines
his gaze to turning toward the Most Exalted Horizon, detached from
conjectures and vain imaginings. O beloved of my heart, certain
discussions truly cause heedlessness and listlessness, especially the
mention of what God does not approve. It is hoped that all will drink of
the choice wine of the All-Merciful's utterance and appear in the
gathering of the world adorned with distinction. The people of God must
assuredly be distinguished from others. Greetings of praise and peace
have been and continue to be conveyed to the friends in that land, and
this mortal servant asks and hopes that the abiding God will assist
everyone to do what is fitting.

\subsection{169}\label{section-2400}

\begin{otherlanguage}{persian}

و در فقرهٴ معهود بسی واضح و معلومست که عنداللّه امریکه سبب ضرّ کسی و
ابتلای نفسی بلکه سبب احزان و کدورات نفوس شود در این ظهور اعظم مردود و
غیرمقبول و آنچه سبب ارتقاء و علوّ و سموّ و علّت فرح و انبساط قلوب اهل عالم
است مقبول و محبوب ولکن نظر بحکمت باید بقسمی رفتار شود که رائحهٴ یأس مرور
ننماید انّ ربّنا الرّحمن قد منع الکلّ عن الفساد و النّزاع و امرهم بالمحبّة و
الشّفقّة و الوداد ولکن حکمت مقتضی مدارا بوده و هست از حقّ میطلبم حضرت مذکور
را تأیید فرماید بر حفظ اغنام در اینصورت اعانت الهی در بارهٴ او البتّه
میرسد انّه یحفظ من اقبل الیه و یقدّر له خیر الآخرة و الأولی و یرفعه بجوده
و رحمته و سلطانه و ینصره بجنود الغیب و الشّهاده انّه هو المقتدر القدیر

\end{otherlanguage}

Va dar faqare-ye maʿhūd basī vāżeḥ va maʿlūm ast ke ʿinda Allāh amrī ke
sabab-e żorr-e kasī va ebtelā-ye nafsī balke sabab-e aḥzān va kodūrāt-e
nofūs shavad dar īn ẓohūr-e aʿẓam mardūd va ghayr-e maqbūl va ānche
sabab-e erteqā va ʿolovv va somovv va ʿellat-e faraḥ va enbesāṭ-e
qolūb-e ahl-e ʿālam ast maqbūl va maḥbūb. Va lāken naẓar be ḥekmat bāyad
be qesmī raftār shavad ke rāʾeḥe-ye yaʾs morūr nanamāyad. Inna rabbanā
al-Raḥmāna qad manaʿa al-kulla ʿan al-fasād wa-l-nizāʿ wa-amarahum
bi-l-maḥabba wa-l-shafaqa wa-l-widād. Va lāken ḥekmat moqtażī-ye modārā
būde va hast. Az Ḥaqq mīṭalabam Ḥażrat-e madhkūr rā taʾyīd farmāyad bar
ḥefẓ-e aghnām. Dar īn ṣūrat eʿānat-e elāhī dar bāre-ye ū albatte
mīresad. Innahu yaḥfaẓu man aqbala ilayhi wa-yuqaddiru lahu khayra
al-ākhirati wa-l-ūlā wa-yarfaʿuhu bi-jūdihi wa-raḥmatihi wa-sulṭānihi
wa-yanṣuruhu bi-junūdi al-ghaybi wa-l-shahāda. Innahu huwa al-Muqtadiru
al-Qadīr.

Regarding the matter in question, it is very clear and evident that, in
God's sight, anything that harms someone or brings affliction upon a
soul---indeed, anything that causes people sorrow and distress---is
rejected and unacceptable in this Most Great Revelation. Whatever brings
advancement, elevation, and exaltation, and causes joy and gladness in
the hearts of the people of the world, is acceptable and beloved. Yet,
in the interests of wisdom, one should act in such a way that no hint of
despair arises. Our Lord, the All-Merciful, has forbidden everyone
corruption and conflict and commanded them to show love, compassion, and
affection. Wisdom, however, has always required and still requires
forbearance. I ask God to assist the person mentioned in protecting the
flock. In that case, divine assistance will assuredly reach him. He
protects whoever turns toward Him, ordains for him the good of the next
world and of this one, exalts him through His bounty, His mercy, and His
sovereignty, and aids him with the hosts of the unseen and the seen. He
is indeed the All-Powerful, the Omnipotent.

\subsection{170}\label{section-2401}

\begin{otherlanguage}{persian}

و امّا فقرهٴ اذکار حسب الاستدعای محبوب مکرّم جناب میرزا ف ض علیه بهآء اللّه
از سماء عنایت نازل و ارسال شد و همچنین جواب حبیب روحانی جناب آقا میرزا
جلیل علیه ٩ ٦٦ {[}بهآء الله{]} و سایر اولیا علیهم بهآء اللّه آنمحبوب فؤاد
بعد از زیارت و اطّلاع بصاحبانش برسانند اینکه اذکاری مخصوص نوم و یقظه
مسئلت نمودند بعد از عرض امام عرش این اذکار مشرقات از افق بیان الهی اشراق
فرمود قوله جلّت عظمته و عظمت قدرته

\end{otherlanguage}

Va ammā faqare-ye adhkār ḥasab-e estedʿā-ye maḥbūb-e mokarram Janāb-e
Mīrzā F Ḍ ʿalayhi bahāʾu Allāh az samā-ye ʿenāyat nāzel va ersāl shod va
hamchenīn javāb-e ḥabīb-e rūḥānī Janāb-e Āqā Mīrzā Jalīl ʿalayhi 9 66
{[}bahāʾ Allāh{]} va sāʾer-e awliyā ʿalayhim bahāʾu Allāh. Ān maḥbūb-e
foʾād baʿd az zeyārat va eṭṭelāʿ be ṣāḥebān-ash berasānand. Īnke adhkārī
makhṣūṣ-e nawm va yaqẓe masʾalat namūdand baʿd az ʿarż-e amām-e ʿarsh īn
adhkār-e moshreqāt az ofoq-e bayān-e elāhī eshrāq farmūd. Qawluhu jallat
ʿaẓamatuhu wa-ʿaẓumat qudratuh.

As for the prayers, they were sent down from the heaven of loving care
and dispatched in accordance with the request of the honored beloved,
Mīrzā F. Ḍ.---upon him be the glory of God---as were the reply to the
spiritual friend, Āqā Mīrzā Jalīl---upon him be 9 66 {[}the glory of
God{]}---and those to the other friends---upon them be the glory of God.
May that beloved of my heart deliver them to their recipients after
reading and acquainting himself with them. As prayers specifically for
sleeping and waking were requested, after the request was presented
before the Throne, these radiant prayers shone forth from the horizon of
divine utterance. His words---exalted be His grandeur and great be His
power:

\subsection{171}\label{section-2402}

\begin{otherlanguage}{arabic}

انت الذّاکر و انت المذکور

\end{otherlanguage}

Anta al-dhākiru wa-anta al-madhkūr.

You are the One who remembers, and You are the One remembered.

\subsection{172}\label{section-2403}

\begin{otherlanguage}{persian}

یا الهی و سیّدی و مقصودی اراد عبدک ان ینام فی جوار رحمتک و یستریح فی ظلّ
قباب فضلک مستعیناً بحفظک و حراستک ای ربّ اسألک بعینک الّتی لا ینام بأن تحفظ
عینی عن النّظر الی دونک ثمّ زد نورها لمشاهدة آثارک و النّظر الی افق ظهورک
انت الّذی ضعفت کینونة القدرة عند ظهورات قدرتک لا اله الّا انت القویّ الغالب
المختار

\end{otherlanguage}

Yā ilāhī wa-sayyidī wa-maqṣūdī arāda ʿabduka an yanāma fī jiwāri
raḥmatika wa-yastarīḥa fī ẓilli qibābi faḍlika mustaʿīnan bi-ḥifẓika
wa-ḥirāsatika. Ay rabbi asʾaluka bi-ʿaynika allatī lā yanāmu bi-an
taḥfaẓa ʿaynī ʿan al-naẓari ilā dūnika thumma zid nūrahā li-mushāhadati
āthārika wa-l-naẓari ilā ufuqi ẓuhūrika. Anta alladhī ḍaʿufat kaynūnatu
al-qudrati ʿinda ẓuhūrāti qudratika. Lā ilāha illā anta al-Qawiyyu
al-Ghālibu al-Mukhtār.

O my God, my Lord, and my Desire! Your servant wishes to sleep in the
shelter of Your mercy and rest beneath the domes of Your grace, seeking
the aid of Your protection and Your watchful care. O Lord, I ask You by
Your eye that never sleeps to guard my eye from looking toward anything
but You. Then increase its light, that it may behold Your signs and gaze
toward the horizon of Your Revelation. You are the One before the
revelations of whose power the very essence of power has grown weak.
There is no God but You, the Strong, the All-Conquering, who chooses as
He wills.

\subsection{173}\label{section-2404}

\begin{otherlanguage}{arabic}

انت اللّه لا اله الّا انت

\end{otherlanguage}

Anta Allāhu lā ilāha illā anta.

You are God; there is no God but You.

\subsection{174}\label{section-2405}

\begin{otherlanguage}{arabic}

لک الحمد یا الهی بما ایقظتنی بعد نومی و اظهرتنی بعد غیبتی و اقمتنی بعد
رقدی اصبحت متوجّهاً الی انوار فجر ظهورک الّذی به انارت آفاق سموات قدرتک و
عظمتک و معترفاً بآیاتک و موقناً بکتابک و متمسّکاً بحبلک

\end{otherlanguage}

Laka al-ḥamdu yā ilāhī bi-mā ayqaẓtanī baʿda nawmī wa-aẓhartanī baʿda
ghaybatī wa-aqamtanī baʿda raqdī. Aṣbaḥtu mutawajjihan ilā anwāri fajri
ẓuhūrika alladhī bihi anārat āfāqu samawāti qudratika wa-ʿaẓamatika
wa-muʿtarifan bi-āyātika wa-mūqinan bi-kitābika wa-mutamassikan
bi-ḥablika.

Praise be to You, O my God, for awakening me after my sleep, bringing me
forth after my absence, and raising me after my slumber. I have entered
the morning turning toward the lights of the dawn of Your Revelation,
through which the horizons of the heavens of Your power and grandeur
have been illumined, acknowledging Your signs, assured of Your Book, and
holding fast to Your cord.

\subsection{175}\label{section-2406}

\begin{otherlanguage}{arabic}

اسألک باقتدار مشیّتک و نفوذ ارادتک بأن تجعل ما اریتنی فی منامی امتن اساس
لثبوت حبّک فی افئدة اولیائک و احسن اسباب لظهورات فضلک و عنایتک ای ربّ قدّر
لی من قلمک الأعلی خیر الآخرة و الأولی اشهد انّ فی قبضتک زمام الأمور
تبدّلها کیف تشآء لا اله الّا انت القویّ الأمین انت الّذی بأمرک تبدّل الذّلّة
بالعزّة و الضّعف بالقوّة و العجز بالاقتدار و الاضطراب بالاطمینان و الرّیب
بالایقان لا اله الّا انت العزیز المنّان لا تخیّب من سئلک و لا تمنع من ارادک
قدّر ما ینبغی لسمآء جودک و بحر کرمک انّک انت المقتدر القدیر

\end{otherlanguage}

Asʾaluka bi-iqtidāri mashīʾatika wa-nufūdhi irādatika bi-an tajʿala mā
araytanī fī manāmī amtana asāsin li-thubūti ḥubbika fī afʾidati
awliyāʾika wa-aḥsana asbābin li-ẓuhūrāti faḍlika wa-ʿināyatika. Ay rabbi
qaddir lī min qalamika al-aʿlā khayra al-ākhirati wa-l-ūlā. Ashhadu anna
fī qabḍatika zimāma al-umūri tubaddiluhā kayfa tashāʾu. Lā ilāha illā
anta al-Qawiyyu al-Amīn. Anta alladhī bi-amrika tubaddilu al-dhillata
bi-l-ʿizzati wa-l-ḍaʿfa bi-l-quwwati wa-l-ʿajza bi-l-iqtidāri
wa-l-iḍṭirāba bi-l-iṭmiʾnāni wa-l-rayba bi-l-īqān. Lā ilāha illā anta
al-ʿAzīzu al-Mannān. Lā tukhayyib man saʾalaka wa-lā tamnaʿ man arādaka.
Qaddir mā yanbaghī li-samāʾi jūdika wa-baḥri karamika. Innaka anta
al-Muqtadiru al-Qadīr.

I ask You, by the might of Your will and the compelling power of Your
purpose, to make what You showed me in my sleep the firmest foundation
for establishing Your love in the hearts of Your friends and the best
means for the revelation of Your grace and loving care. O Lord, ordain
for me through Your Most Exalted Pen the good of the next world and of
this one. I bear witness that the reins of all affairs are in Your
grasp; You change them as You will. There is no God but You, the Strong,
the Faithful. You are the One who, by Your command, changes abasement
into glory, weakness into strength, helplessness into power, agitation
into tranquillity, and doubt into certainty. There is no God but You,
the Almighty, the Bestower. Do not disappoint whoever has asked of You,
nor withhold from whoever has sought You. Ordain what befits the heaven
of Your bounty and the ocean of Your generosity. You are indeed the
All-Powerful, the Omnipotent.

\subsection{176}\label{section-2407}

\begin{otherlanguage}{arabic}

هو المهیمن علی الأسمآء

\end{otherlanguage}

Huwa al-Muhayminu ʿalā al-asmāʾ.

He is the Guardian over all names.

\subsection{177}\label{section-2408}

\begin{otherlanguage}{persian}

الهی الهی کیف اختار النّوم و عیون مشتاقیک ساهرة فی فراقک و کیف استریح علی
الفراش و افئدة عاشقیک مضطربة من هجرک ای ربّ اودعت روحی و ذاتی فی یمین
اقتدارک و امانک و اوضع رأسی علی الفراش بحولک و ارفع عنه بمشیّتک و ارادتک
انّک انت الحافظ الحارس المقتدر القدیر

\end{otherlanguage}

Ilāhī ilāhī kayfa akhtāru al-nawma wa-ʿuyūnu mushtāqīka sāhiratun fī
firāqika wa-kayfa astarīḥu ʿalā al-firāshi wa-afʾidatu ʿāshiqīka
muḍṭaribatun min hajrika. Ay rabbi awdaʿtu rūḥī wa-dhātī fī yamīni
iqtidārika wa-amānika wa-awḍaʿu raʾsī ʿalā al-firāshi bi-ḥawlika
wa-arfaʿu ʿanhu bi-mashīʾatika wa-irādatika. Innaka anta al-Ḥāfiẓu
al-Ḥārisu al-Muqtadiru al-Qadīr.

My God, my God! How can I choose sleep when the eyes of those who long
for You remain wakeful in separation from You? How can I rest upon my
bed when the hearts of Your lovers are troubled by their estrangement
from You? O Lord, I have entrusted my soul and my being to the right
hand of Your might and Your protection. Through Your power I lay my head
upon my bed, and through Your will and purpose I raise it from there.
You are indeed the Preserver, the Watchful Guardian, the All-Powerful,
the Omnipotent.

\subsection{178}\label{section-2409}

\begin{otherlanguage}{arabic}

وعزّتک لا ارید من النّوم و لا من الیقظه الّا ما انت ترید انا عبدک و فی
قبضتک ایّدنی علی ما یتضوّع به عرف رضائک هذا املی و امل المقرّبین الحمد لک
یا اله العالمین

\end{otherlanguage}

Wa-ʿizzatika lā urīdu min al-nawmi wa-lā min al-yaqẓati illā mā anta
turīdu. Anā ʿabduka wa-fī qabḍatika. Ayyidnī ʿalā mā yataḍawwaʿu bihi
ʿarfu riḍāʾika. Hādhā amalī wa-amalu al-muqarrabīn. Al-ḥamdu laka yā
ilāha al-ʿālamīn.

By Your glory, I desire from sleep or wakefulness only what You desire.
I am Your servant and am in Your grasp. Assist me to do that through
which the fragrance of Your good pleasure may spread. This is my hope
and the hope of those brought near to You. Praise be to You, O God of
the worlds.

\subsection{179}\label{section-2410}

\begin{otherlanguage}{arabic}

هو السّامع المجیب

\end{otherlanguage}

Huwa al-Sāmiʿu al-Mujīb.

He is the Hearer, the Responsive.

\subsection{180}\label{section-2411}

\begin{otherlanguage}{persian}

یا الهی اصبحت فی جوارک و الّذی استجارک ینبغی ان یکون فی کنف حفظک و حصن
حمایتک ای ربّ نوّر باطنی بأنوار فجر ظهورک کما نوّرت ظاهری بنور صباح عطائک

\end{otherlanguage}

Yā ilāhī aṣbaḥtu fī jiwārika wa-alladhī istajāraka yanbaghī an yakūna fī
kanafi ḥifẓika wa-ḥiṣni ḥimāyatika. Ay rabbi nawwir bāṭinī bi-anwāri
fajri ẓuhūrika kamā nawwarta ẓāhirī bi-nūri ṣabāḥi ʿaṭāʾika.

O my God, I have entered the morning in Your shelter, and whoever has
sought Your shelter should be under the care of Your protection and
within the stronghold of Your defense. O Lord, illumine my inner being
with the lights of the dawn of Your Revelation, just as You have
illumined my outer being with the light of the morning of Your bounty.

\subsection{181}\label{section-2412}

\begin{otherlanguage}{arabic}

هو المهیمن القیّوم

\end{otherlanguage}

Huwa al-Muhayminu al-Qayyūm.

He is the All-Protecting, the Self-Subsisting.

\subsection{182}\label{section-2413}

\begin{otherlanguage}{arabic}

اصبحت یا الهی بفضلک و اخرج من البیت متوکّلاً علیک و مفوّضاً امری الیک فانزل
علیّ من سمآء رحمتک برکة من عندک ثمّ ارجعنی الی البیت سالماً کما اخرجنی منه
سالماً مستقیماً

\end{otherlanguage}

Aṣbaḥtu yā ilāhī bi-faḍlika wa-akhruju min al-bayti mutawakkilan ʿalayka
wa-mufawwiḍan amrī ilayka. Fa-anzil ʿalayya min samāʾi raḥmatika
barakatan min ʿindika thumma arjiʿnī ilā al-bayti sāliman kamā akhrijnī
minhu sāliman mustaqīman.

I have entered the morning, O my God, through Your grace, and I leave
the house trusting in You and committing my affairs to You. Send down
upon me from the heaven of Your mercy a blessing from Your presence.
Then return me to the house in safety, just as You bring me out of it
safe and steadfast.

\subsection{183}\label{section-2414}

\begin{otherlanguage}{arabic}

لا اله الّا انت الفرد الواحد العلیم الحکیم انتهی

\end{otherlanguage}

Lā ilāha illā anta al-Fardu al-Wāḥidu al-ʿAlīmu al-Ḥakīm. Intahā.

There is no God but You, the Unique, the One, the All-Knowing, the
All-Wise. Here it ends.

\subsection{184}\label{section-2415}

\begin{otherlanguage}{persian}

آفتاب فضل الهی لازال مشرق و لائح دلیله نفسه و ظهورات عنایته آیاته هو
الدّلیل و هو الظّاهر و هو الشّاهد و هو المبیّن عنایتش شامل و رحمتش محیط
احتیاج باستدلال نداشته و ندارد

\end{otherlanguage}

Āftāb-e faḍl-e elāhī lāzāl moshreq o lāʾeḥ. Dalīluhu nafsuhu wa-ẓuhūrātu
ʿināyatihi āyātuhu. Huwa al-dalīl wa-huwa al-Ẓāhir wa-huwa al-shāhid
wa-huwa al-mubayyin. ʿEnāyatash shāmel o raḥmatash moḥīṭ. Eḥtīāj
be-estedlāl nadāshteh va nadārad.

The sun of divine grace has always been radiant and resplendent. He
Himself is His proof, and the manifestations of His loving care are His
signs. He is the Proof, the Manifest, the Witness, and the Expounder.
His loving care embraces all, and His mercy encompasses everything.
There has never been, nor is there now, any need for argument.

\subsection{185}\label{section-2416}

\begin{otherlanguage}{persian}

در بارهٴ جناب حاجی علی عسکر م ی علیه ٩ ٦٦ {[}بهآء الله{]} مرقوم داشتند
بعد از عرض در ساحت اقدس اذن رجوع نازل با بودن عیال البتّه باید بکتاب اقدس
عمل نمایند صدهزار طوبی از برای عاملین اینکه

\end{otherlanguage}

Dar bāreh-ye Janāb-e Ḥājī ʿAlī ʿAskar M. Y. ʿalayhi 9 66 {[}Bahāʾ
Allāh{]} marqūm dāshtand. Baʿd az ʿarḍ dar sāḥat-e aqdas edhn-e rojūʿ
nāzel. Bā būdan-e ʿiyāl al-batteh bāyad be-Ketāb-e Aqdas ʿamal namāyand.
Ṣad-hezār ṭūbā az barā-ye ʿāmelīn. Īn-keh

You wrote concerning the honorable Ḥājī ʿAlī ʿAskar M. Y.---upon him be
9 66 {[}the glory of God{]}. After the matter was presented at the Most
Holy Court, permission to return was sent down. Since he has a family,
he must certainly act in accordance with the Most Holy Book. A hundred
thousand blessings upon those who act accordingly! As for the fact that

\subsection{186}\label{section-2417}

\begin{otherlanguage}{persian}

در بارهٴ مسجون تفلیس مشهدی اسمعیل علاقه‌بند مرقوم داشتند در حضور عرض شد و
به عنایت و شفقت الهی فائز ذکرش از لسان عظمت جاری و فرمودند بنویس باو
محزون مباش مولی العالم در سجن اعظم قل الهی الهی توکّلت علیک و فوّضت امری
الیک قدّر لی ما ینفعنی فیکلّ عالم من عوالمک انّک انت العلیم الخبیر انتهی

\end{otherlanguage}

Dar bāreh-ye masjūn-e Teflīs, Mashhadī Esmāʿīl-e ʿAlāqeh-band, marqūm
dāshtand. Dar ḥoḍūr ʿarḍ shod va be-ʿenāyat o shafaqat-e elāhī fāʾez.
Dhekrash az lesān-e ʿaẓamat jārī va farmūdand benevīs be-ū maḥzūn
mabāsh, Mawlā al-ʿālam dar sejn-e aʿẓam. Qul ilāhī ilāhī tawakkaltu
ʿalayka wa-fawwaḍtu amrī ilayka. Qaddir lī mā yanfaʿunī fī kulli ʿālamin
min ʿawālimika. Innaka anta al-ʿAlīm al-Khabīr. Intahā.

You wrote concerning the prisoner in Tiflis, Mashhadī Ismāʿīl the
cord-maker. The matter was presented in His presence, and he received
divine loving care and compassion. His name was mentioned by the Tongue
of Grandeur, and He said: ``Write to him: Do not grieve. The Lord of the
world is in the Most Great Prison. Say: My God, my God! I have placed my
trust in You and committed my affairs to You. Ordain for me what will
benefit me in every world among Your worlds. You are indeed the
All-Knowing, the All-Aware.'' Here it ends.

\subsection{187}\label{section-2418}

\begin{otherlanguage}{persian}

ذکر جناب آقا محمّد صادق الّذی هاجر و فاز فرمودند مخصوص ایشان لوح امنع اقدس
از سماء مشیّت نازل و ارسال شد انشاءاللّه بآن فائز شوند و از بحر معانی
بیاشامند اینعبد هم خدمت ایشان سلام و تکبیر میرساند و توفیق میطلبد

\end{otherlanguage}

Dhekr-e Janāb-e Āqā Moḥammad Ṣādeq alladhī hājara wa-fāza farmūdand.
Makhṣūṣ-e īshān lawḥ-e amnaʿ-e aqdas az samā-ye mashīyyat nāzel va ersāl
shod. In shāʾ Allāh be-ān fāʾez shavand va az baḥr-e maʿānī biyāshāmand.
Īn ʿabd ham khedmat-e īshān salām o takbīr mīresānad va tawfīq
mīṭalabad.

You mentioned the honorable Āqā Muḥammad Ṣādiq, who emigrated and
attained His presence. A most exalted, most holy Tablet was sent down
from the heaven of the divine will especially for him and dispatched.
God willing, may he receive it and drink from the ocean of meanings.
This servant also conveys greetings and praise of God to him and prays
for his success.

\subsection{188}\label{section-2419}

\begin{otherlanguage}{persian}

در بارهٴ آنچه که رحمت عالم از او ظاهر و همچنین غضبش مرقوم داشتند یعنی
مرکّب که مراکب عالم از او جاری و ارکان امم از او محکم فی‌الحقیقه اگر اقلام
عالم بخواهند اولاد این ام را احصا نمایند خود را عاجز مشاهده کنند فکر باو
نرسد دفاتر او را کفایت نکند باری اگر بر جمیع اوراق فی‌الحقیقه وصف او ذکر
شود نزد تصرّفات عالیهٴ قاهرهٴ او معدوم و مفقود است باری آنچه باین اوصاف و
فوق آن مذکور رسید و بسیار خوب و مرغوب لذا {[}هی{]} رحمت از او ظاهر

\end{otherlanguage}

Dar bāreh-ye āncheh keh raḥmat-e ʿālam az ū ẓāher va hamchenīn ghaḍabash
marqūm dāshtand, yaʿnī morakkab keh marākeb-e ʿālam az ū jārī va arkān-e
omam az ū moḥkam. Fī al-ḥaqīqeh agar aqlām-e ʿālam bekhāhand awlād-e īn
omm rā eḥṣā namāyand khod rā ʿājez moshāhedeh konand. Fekr be-ū naresad,
dafāter ū rā kefāyat nakonad. Bārī agar bar jamīʿ-e awrāq fī al-ḥaqīqeh
vaṣf-e ū dhekr shavad nazd-e taṣarrofāt-e ʿāliyeh-ye qāhereh-ye ū maʿdūm
o mafqūd ast. Bārī āncheh be-īn awṣāf va fawq-e ān madhkūr rasīd va
besyār khūb o marghūb, ledhā {[}hiya{]} raḥmat az ū ẓāher.

You wrote concerning that through which mercy for the world is
manifested, and likewise His wrath---that is, ink, through which the
vessels of the world sail and the pillars of nations are made firm. In
truth, if all the pens of the world wished to count the offspring of
this mother, they would find themselves powerless. Thought cannot
encompass it, nor can volumes suffice for it. In short, even if its
description were recorded on every sheet of paper, it would truly be as
nothing before its lofty, all-subduing powers. In short, what has been
described in these terms, and in terms beyond them, is very good and
desirable. Thus {[}it is{]} mercy that is manifested through it.

\subsection{189}\label{section-2420}

\begin{otherlanguage}{arabic}

سبحانک یا اله الممکنات و مقصود الکائنات انت الّذی اودعت فی قطره شیء حالک
ما اهتزّنّه اهل القبور به احییت و به اخذت و قبضت اسألک بقدرتک المهیمنة علی
العالم بأن توفّق الأمم علی قبول امرک انّک انت المقتدر علی ما تشآء و انّک
انت المهیمن القیّوم اسألک یا الهی ببحر عنایتک و سمآء فضلک بأن تنزل علی من
احبّک و ذکر خادمک ببهآء من عندک و سلام من لدنک و روح من جانبک انّک انت
المقتدر الّذی لا تمنعک ما تشهد و تری و انّک انت مولی الوری و ربّ العرش و
الثّری لا اله الّا انت العلیّ الأبهی

\end{otherlanguage}

Subḥānaka yā ilāha al-mumkināt wa-maqṣūda al-kāʾināt. Anta alladhī
awdaʿta fī qaṭrati shayʾin ḥālikin mā {[}corrupted text{]} ahlu
al-qubūr. Bihi aḥyayta wa-bihi akhadhta wa-qabaḍta. Asʾaluka
bi-qudratika al-muhayminati ʿalā al-ʿālam bi-an tuwaffiqa al-umama ʿalā
qabūli amrika. Innaka anta al-Muqtadir ʿalā mā tashāʾ wa-innaka anta
al-Muhaymin al-Qayyūm. Asʾaluka yā ilāhī bi-baḥri ʿināyatika wa-samāʾi
faḍlika bi-an tunzila ʿalā man aḥabbaka wa-dhakara khādimaka bi-bahāʾin
min ʿindika wa-salāmin min ladunka wa-rūḥin min jānibika. Innaka anta
al-Muqtadir alladhī lā tamnaʿuka mā tashhadu wa-tarā wa-innaka anta
Mawlā al-warā wa-Rabb al-ʿarsh wa-al-tharā. Lā ilāha illā anta al-ʿAlī
al-Abhā.

Glory be to You, O God of all contingent beings and Desire of all
created things! You are the One who placed within a drop of a dark
substance that which {[}corrupted original text{]} the people of the
graves. Through it You gave life, and through it You seized and took
away. I ask You, by Your power that holds sway over the world, to enable
the nations to accept Your Cause. You are indeed powerful over whatever
You will, and You are the All-Protecting, the Self-Subsisting. I ask
You, O my God, by the ocean of Your loving care and the heaven of Your
grace, to send down upon whoever has loved You and mentioned Your
servant glory from You, peace from Your presence, and spirit from Your
side. You are indeed the All-Powerful, whom nothing You witness and see
can hinder, and You are the Lord of all creation and the Lord of the
throne and the earth. There is no God but You, the Most High, the
All-Glorious.

\bookmarksetup{startatroot}

\chapter{BH00010 (Unnamed)}\label{bh00010-unnamed}

\emph{Translated on 17 Qudrat 183 using GPT 6.1 Sol}

\subsection{1}\label{section-2421}

\begin{otherlanguage}{arabic}

بسم ربّنا الأقدس الأعظم العلیّ الأبهی

\end{otherlanguage}

Bismi rabbinā al-Aqdas al-Aʿẓam al-ʿAliyy al-Abhā.

In the name of our Lord, the Most Holy, the Most Great, the Most
Exalted, the All-Glorious.

\subsection{2}\label{section-2422}

\begin{otherlanguage}{arabic}

الحمد للّه الّذی تفرّد بالمحبّة و توحّد بالوداد انّه لمولی العباد و مالک
الایجاد و النّاطق بما اراد فی المبدء و المعاد و النّفحة الّتی تضوّعت من قمیص
اللّه و النّور المشرق من افق عنایته و السّلام الظّاهر اللّائح من ملکوت فضله
علی الّذین قاموا علی اعلآء کلمته و اظهار امره و عملوا بما امروا به فی
کتابه اولئک عباد بهم جرت الأوامر و الأحکام بین الأنام و ثبت حکم التّقدیس
بین الأحزاب بهم رفعت الرّایات و ظهرت العلامات یصلّینّ علیهم الملأ الأعلی و
عن ورائهم مالک الأسمآء و فاطر السّمآء

\end{otherlanguage}

Al-ḥamdu li-llāhi alladhī tafarrada bi-l-maḥabbati wa-tawaḥḥada
bi-l-widād. Innahu la-mawlā al-ʿibādi wa-māliku al-ījādi wa-l-nāṭiqu
bi-mā arāda fī al-mabdaʾi wa-l-maʿād. Wa-l-nafḥatu allatī taḍawwaʿat min
qamīṣi Allāhi wa-l-nūru al-mushriqu min ufuqi ʿināyatihi wa-l-salāmu
al-ẓāhiru al-lāʾiḥu min malakūti faḍlihi ʿalā alladhīna qāmū ʿalā iʿlāʾi
kalimatihi wa-iẓhāri amrihi wa-ʿamilū bi-mā umirū bihi fī kitābih.
Ulāʾika ʿibādun bihim jarati al-awāmiru wa-l-aḥkāmu bayna al-anāmi
wa-thabata ḥukmu al-taqdīsi bayna al-aḥzāb. Bihim rufiʿati al-rāyātu
wa-ẓaharati al-ʿalāmāt. Yuṣalliyanna ʿalayhimu al-malaʾu al-aʿlā wa-ʿan
warāʾihim Māliku al-Asmāʾi wa-Fāṭiru al-Samāʾ.

Praise be to God, who is peerless in love and unique in affection. He is
indeed the Lord of all servants, the Master of creation, who speaks
whatever He wills in the beginning and the return. May the fragrance
wafting from the garment of God, the light shining from the horizon of
His loving care, and the peace manifest and radiant from the kingdom of
His grace rest upon those who have arisen to exalt His Word and make His
Cause manifest, and have acted as they were commanded in His Book. These
are servants through whom the commandments and laws have gained currency
among humanity and the rule of sanctity has been established among its
peoples. Through them the banners have been raised and the signs have
appeared. The Concourse on high will surely bless them, and beyond them,
the Lord of Names and Creator of heaven.

\subsection{3}\label{section-2423}

\begin{otherlanguage}{persian}

و بعد قد اشرق من افق سمآء ودادکم ما تهیّجت به اریاح المحبّة انّه لکتاب ترشّح
فیه انآء الفطرة فلمّا وجدت العرف و عرفت ما انتم علیه فی امر اللّه ربّی و
ربّکم و ربّ من فی السّموات و الأرض قصدت المقام المحمود الّذی فیه استقرّ عرش
مالک الوجود و سلطان الغیب و الشّهود حضرت و عرضت ما نطق به عندلیب خلوصکم
علی غصن الوفآء اذاً توجّه الیّ وجه مالک الأسمآء تاللّه قد نطق بما طارت به
الجبال ولکن مرآة قلبی ما انطبع علیها الّا علی قدر ضرّی و مسکنتی و فقری و
فاقتی یا محبوب فؤادی کیف تقدر القطرة ان تحصی بحر البیان لعمر محبوبی و
محبوبکم لا اتکلّم الّا بالحقّ براستی عرض مینمایم و آنمحبوب قلب شاهد و
گواهند چه که بفضل اللّه آگاه بوده و هستند

\end{otherlanguage}

Wa-baʿdu qad ashraqa min ufuqi samāʾi widādikum mā tahayyajat bihi
aryāḥu al-maḥabbah. Innahu la-kitābun tarashshaḥa fīhi ināʾu al-fiṭrah.
Fa-lammā wajadtu al-ʿarfa wa-ʿaraftu mā antum ʿalayhi fī amri Allāhi
rabbī wa-rabbikum wa-rabbi man fī al-samāwāti wa-l-arḍi qaṣadtu
al-maqāma al-maḥmūda alladhī fīhi istaqarra ʿarshu Māliki al-Wujūdi
wa-Sulṭāni al-Ghaybi wa-l-Shuhūd. Ḥaḍartu wa-ʿaraḍtu mā naṭaqa bihi
ʿandalību khulūṣikum ʿalā ghuṣni al-wafāʾ. Idhan tawajjaha ilayya wajhu
Māliki al-Asmāʾ. Ta-llāhi qad naṭaqa bi-mā ṭārat bihi al-jibālu wa-lākin
mirʾāta qalbī mā inṭabaʿa ʿalayhā illā ʿalā qadri ḍurrī wa-maskanatī
wa-faqrī wa-fāqatī. Yā maḥbūba fuʾādī kayfa taqdiru al-qaṭratu an
tuḥṣiya baḥra al-bayān. La-ʿamri maḥbūbī wa-maḥbūbikum lā atakallamu
illā bi-l-ḥaqq. Bi-rāstī ʿarḍ mī-namāyam va ān maḥbūb-i qalb shāhid va
guvāhand chih kih bi-faḍl-i Allāh āgāh būdih va hastand.

To continue: from the horizon of the heaven of your affection there has
shone forth that which has stirred the winds of love. It is indeed a
letter in which the vessel of your innate nature has overflowed. When I
perceived its fragrance and recognized your condition in the Cause of
God, my Lord and your Lord and the Lord of all who are in the heavens
and on earth, I set out for the praised station where the throne of the
Lord of existence, the Sovereign of the unseen and the seen, is
established. I entered His presence and presented what the nightingale
of your sincerity had sung upon the branch of faithfulness. Then the
face of the Lord of Names turned toward me. By God, He spoke words that
made the mountains take flight, yet nothing was imprinted upon the
mirror of my heart except in the measure of my affliction, lowliness,
poverty, and need. O beloved of my heart, how can a drop encompass the
ocean of utterance? By the life of my Beloved and yours, I speak only
the truth. I tell you truly, and that Beloved of the heart bears witness
and testimony, for through the grace of God He has been and remains
aware.

\subsection{4}\label{section-2424}

\begin{otherlanguage}{arabic}

قال جلّت عظمته و عظم کبریائه

\end{otherlanguage}

Qāla jallat ʿaẓamatuhu wa-ʿaẓuma kibriyāʾuh.

He said---exalted be His grandeur and great His majesty:

\subsection{5}\label{section-2425}

\begin{otherlanguage}{arabic}

بنام مظلوم آفاق

\end{otherlanguage}

Bi-nām-i Maẓlūm-i āfāq.

In the name of the Wronged One of the world.

\subsection{6}\label{section-2426}

\begin{otherlanguage}{persian}

یا علی الحمد للّه بصر عالم از کحل حکمت و بیان احبّای الهی روشنست اگرچه این
روشنی را سحاب تیرهٴ اوهام خلق از ظهور و بروز منع نموده ولکن البتّه اریاح
اراده این ابر را متفرّق نماید اوست قادر و توانا و قدرتش الیوم مستور ولکن
در سرّ بفتح مدائن قلوب مشغول بیان رحمن مفتاحی است قوی و فاتح و آفتابیست
مشرق و لائح و بحریست محیط بگشاید و نور بخشد و حیات عطا نماید طوبی از
برای بصیریکه شکّ و ریب و اوهام و ظنون و سطوت و غضب اهل عالم او را از
ذروهٴ علیا منع ننمود بیک کلمه عالم مضطرب و بکلمهٴ اخری ساکن و مطمئن و
بیک عمل مشتعل و بعمل آخر برد و سلام اینست که از قلم اعلی در زبر و الواح
اخلاق روحانیّه و اعمال مرضیّه و کلمهٴ طیّبه از جنود قویّهٴ الهیّه ذکر شده از
جمیع مراتب گذشته اگر ناس بانصاف تمسّک نمایند هم خود بانوار یقین فائز
گردند و آثار و حجج و براهین حقّ جلّ جلاله را بمثابهٴ آفتاب مشاهده نمایند و
هم کرسیّ اقتدار حقّ مابین عالم ظاهراً باهراً مشاهده گردد

\end{otherlanguage}

Yā ʿAlī, al-ḥamdu li-llāh, baṣar-i ʿālam az kuḥl-i ḥikmat va bayān-i
aḥibbā-yi ilāhī rawshan ast. Agarchih īn rawshanī rā saḥāb-i tīrih-yi
awhām-i khalq az ẓuhūr va burūz manʿ nimūdih va-lākin al-battih aryāḥ-i
irādih īn abr rā mutafarriq namāyad. Ūst qādir va tavānā va qudratash
al-yawm mastūr va-lākin dar sirr bi-fatḥ-i madāʾin-i qulūb mashghūl.
Bayān-i Raḥmān miftāḥī ast qavī va fātiḥ va āftābī-st mushriq va lāʾiḥ
va baḥrī-st muḥīṭ. Bi-gushāyad va nūr bakhshad va ḥayāt ʿaṭā namāyad.
Ṭūbā az barā-yi baṣīrī kih shakk va rayb va awhām va ẓunūn va saṭvat va
ghaḍab-i ahl-i ʿālam ū rā az dhurvih-yi ʿulyā manʿ nanimūd. Bi-yik
kalimih ʿālam muḍṭarib va bi-kalimih-yi ukhrā sākin va muṭmaʾinn va
bi-yik ʿamal mushtaʿil va bi-ʿamal-i ākhar bard va salām. Īn ast kih az
Qalam-i Aʿlā dar zubur va alvāḥ akhlāq-i rūḥāniyyih va aʿmāl-i marḍiyyih
va kalimih-yi ṭayyibih az junūd-i qaviyyih-yi ilāhiyyih dhikr shudih. Az
jamīʿ-i marātib gudhashtih agar nās bi-inṣāf tamassuk namāyand ham khud
bi-anvār-i yaqīn fāʾiz gardand va āthār va ḥujaj va barāhīn-i Ḥaqq jalla
jalāluhu rā bi-mathābih-yi āftāb mushāhadih namāyand va ham kursī-yi
iqtidār-i Ḥaqq mābayn-i ʿālam ẓāhiran bāhiran mushāhadih gardad.

O ʿAlī, praise be to God, the sight of the world is illumined by the
kohl of the wisdom and utterance of the loved ones of God. Although the
dark clouds of the people's vain imaginings have prevented this light
from appearing and becoming manifest, the winds of His will shall
assuredly disperse these clouds. He is powerful and mighty, and His
power, though concealed today, is secretly engaged in conquering the
cities of hearts. The utterance of the All-Merciful is a mighty key that
unlocks, a sun shining and radiant, and an all-encompassing ocean. It
unlocks, sheds light, and bestows life. Blessed is the discerning one
whom doubt and uncertainty, vain imaginings and conjectures, and the
power and wrath of the people of the world have not kept from the
loftiest summit. Through one word the world is thrown into turmoil, and
through another it becomes still and assured; through one deed it is set
ablaze, and through another it becomes coolness and peace. Thus, in the
scriptures and tablets, the Most Exalted Pen has numbered spiritual
character, pleasing deeds, and a good word among the mighty hosts of
God. Beyond all these considerations, if people hold fast to fairness,
they themselves will attain the lights of certitude and behold the
signs, proofs, and testimonies of God---exalted be His glory---as
clearly as the sun, and the throne of God's power will also be seen,
manifest and resplendent, amidst the world.

\subsection{7}\label{section-2427}

\begin{otherlanguage}{persian}

یا اکبر دوستان را تکبیر برسان و از قبل مظلوم القا نما شما از سایر اهل
عالم باید ممتاز باشید و در جمیع احیان بشطر امتیاز ناظر قدر خود را بدانید
و مقام اولیای حقّ را که از اوّل امر بخدمت قیام نموده‌اند بشناسید قلم اعلی
جمیعرا ذکر فرموده و بخلع عنایت مزیّن داشته ولکن از الواح نفوس مستقیمهٴ
ثابته عرف امتیاز متضوّع طوبی لمن وجد و کان من العاملین

\end{otherlanguage}

Yā Akbar, dūstān rā takbīr birasān va az qibal-i Maẓlūm ilqā namā: shumā
az sāʾir-i ahl-i ʿālam bāyad mumtāz bāshīd va dar jamīʿ-i aḥyān
bi-shaṭr-i imtiyāz nāẓir. Qadr-i khud rā bidānīd va maqām-i awliyā-yi
Ḥaqq rā kih az avval-i amr bi-khidmat qiyām nimūdih-and bishināsīd.
Qalam-i Aʿlā jamīʿ rā dhikr farmūdih va bi-khilaʿ-i ʿināyat muzayyan
dāshtih va-lākin az alvāḥ-i nufūs-i mustaqīmih-yi thābitih ʿarf-i
imtiyāz mutaḍavviʿ. Ṭūbā li-man wajada wa-kāna mina al-ʿāmilīn.

O Akbar, convey greetings of praise to the friends and tell them on
behalf of the Wronged One: You must be distinguished from the other
people of the world and at all times keep distinction in view. Know your
own worth and recognize the station of the loved ones of God who have
arisen to serve from the beginning of the Cause. The Most Exalted Pen
has mentioned everyone and adorned them with robes of loving care, yet
from the tablets of steadfast, unwavering souls the fragrance of
distinction is diffused. Blessed is the one who has perceived it and is
among those who act.

\subsection{8}\label{section-2428}

\begin{otherlanguage}{arabic}

یا ایّها النّاظر الی اللّه اشهد انّک شربت رحیق حبّی فی اوّل ایّامی و اعترفت بما
اعترف به اللّه و تمسّکت به و نشرت آثاره و کنت قائماً علی خدمته فی یوم فیه
زلّت الأقدام کبّر من قبلی علی احبّائی فی الطّآء و من حولها و بشّرهم بعنایة
اللّه ربّ العالمین

\end{otherlanguage}

Yā ayyuhā al-nāẓiru ilā Allāh, ashhadu annaka sharibta raḥīqa ḥubbī fī
awwali ayyāmī wa-iʿtarafta bi-mā iʿtarafa bihi Allāhu wa-tamassakta bihi
wa-nasharta āthārahu wa-kunta qāʾiman ʿalā khidmatihi fī yawmin fīhi
zallati al-aqdām. Kabbir min qibalī ʿalā aḥibbāʾī fī al-Ṭāʾi wa-man
ḥawlahā wa-bashshirhum bi-ʿināyati Allāhi Rabbi al-ʿālamīn.

O you who look toward God, I testify that you drank the choice wine of
My love in the earliest days of My revelation, acknowledged what God
acknowledged, held fast to it, spread His writings, and stood ready to
serve Him on a day when feet faltered. Convey greetings of praise on My
behalf to My loved ones in Ṭāʾ and those around it, and give them the
glad tidings of the loving care of God, the Lord of the worlds.

\subsection{9}\label{section-2429}

\begin{otherlanguage}{arabic}

انّا نوصیهم بالتّقوی الخالص و بما یرتفع به امر اللّه العلیم الخبیر قل ان
ارحموا علی انفسکم و علی سدرة الأمر و علی الدّمآء الّتی سفکت لاعلآء کلمة
اللّه المطاعة اتّقوا اللّه یا قوم و کونوا من المنصفین لا تضیّعوا اعمالکم و
ما ورد علیکم فی هذا النّبأ العظیم لکم ان تفتخروا علی من علی الأرض انّکم
فزتم بما لا فاز به احد من الخلق یشهد بذلک من نطق امّ الکتاب بذکره العزیز
البدیع لا تجعلوا انفسکم ملعب جهلآء الأرض ینبغی لکلّ عمل ینسب الیکم ان
یکون مشرقاً کالشّمس بین العالم کذلک یعظکم اللّه حقّ الوعظ رحمةً من عنده انّه
یحبّکم و یرید لکم ما ینبغی لمقاماتکم عند اللّه العلیّ العظیم ایّاکم ان
تعملوا ما تتوقّف به الرّحمة فی سمآء المشیّة کذلک ینصحکم من کان اشفق بکم
منکم انّ ربّکم الرّحمن لهو العلیم الحکیم البهآء المشرق من افق عنایتی و
النّور اللّائح من سمآء رحمتی علیک و علی من معک و یحبّک فی سبیل اللّه مالک
هذا الیوم المبارک المبین انتهی

\end{otherlanguage}

Innā nūṣīhim bi-l-taqwā al-khāliṣi wa-bi-mā yartafiʿu bihi amru Allāhi
al-ʿAlīmi al-Khabīr. Qul an irḥamū ʿalā anfusikum wa-ʿalā sidrati
al-amri wa-ʿalā al-dimāʾi allatī sufikat li-iʿlāʾi kalimati Allāhi
al-muṭāʿah. Ittaqū Allāha yā qawmu wa-kūnū mina al-munṣifīn. Lā
tuḍayyiʿū aʿmālakum wa-mā warada ʿalaykum fī hādhā al-nabaʾi al-ʿaẓīm.
Lakum an taftakhirū ʿalā man ʿalā al-arḍi annakum fuztum bi-mā lā fāza
bihi aḥadun mina al-khalq. Yashhadu bi-dhālika man naṭaqa Ummu al-Kitābi
bi-dhikrihi al-ʿAzīzu al-Badīʿ. Lā tajʿalū anfusakum malʿaba juhalāʾi
al-arḍ. Yanbaghī li-kulli ʿamalin yunsabu ilaykum an yakūna mushriqan
ka-l-shamsi bayna al-ʿālam. Kadhālika yaʿiẓukumu Allāhu ḥaqqa al-waʿẓi
raḥmatan min ʿindihi innahu yuḥibbukum wa-yurīdu lakum mā yanbaghī
li-maqāmātikum ʿinda Allāhi al-ʿAliyyi al-ʿAẓīm. Iyyākum an taʿmalū mā
tatawaqqafu bihi al-raḥmatu fī samāʾi al-mashīʾah. Kadhālika yanṣaḥukum
man kāna ashfaqa bikum minkum. Inna rabbakumu al-Raḥmāna la-huwa
al-ʿAlīmu al-Ḥakīm. Al-Bahāʾu al-mushriqu min ufuqi ʿināyatī wa-l-nūru
al-lāʾiḥu min samāʾi raḥmatī ʿalayka wa-ʿalā man maʿaka wa-yuḥibbuka fī
sabīli Allāhi Māliki hādhā al-yawmi al-mubāraki al-mubīn. Intahā.

We counsel them to observe pure godliness and to do whatever exalts the
Cause of God, the All-Knowing, the All-Aware. Say: Have mercy on
yourselves, on the Tree of the Cause, and on the blood shed to exalt the
Word of God, which all must obey. Fear God, O people, and be
fair-minded. Do not squander your deeds and what you have endured for
this mighty Announcement. You may well take pride before all who dwell
on earth, for you have attained what no one else in creation has
attained. To this bears witness He whose mention the Mother Book has
voiced, the Almighty, the Wondrous. Do not make yourselves the
playthings of the ignorant of the earth. Every deed attributed to you
should shine like the sun amidst the world. Thus God admonishes you with
true admonition, as a mercy from Him. He loves you and desires for you
what befits your stations before God, the Most Exalted, the Most Great.
Beware of doing anything that would hold back mercy in the heaven of His
will. Thus counsels you One who is more compassionate toward you than
you are toward yourselves. Your Lord, the All-Merciful, is indeed the
All-Knowing, the All-Wise. May the glory shining from the horizon of My
loving care and the light gleaming from the heaven of My mercy rest upon
you and upon those who are with you and love you in the path of God, the
Lord of this blessed and manifest Day. End.

\subsection{10}\label{section-2430}

\begin{otherlanguage}{persian}

انشآءاللّه جمیع دوستان الهی بقمیص تقدیس مزیّن شوند باید بشأنی ظاهر باشند
که اگر در انجمن عالم بوعظ و نصح عباد قیام نمایند و بأعلی البیان تنطّق
کنند نه خود در سرّ خجل باشند و نه من علی الأرض قادر بر اعتراض حال باید در
لیالی و ایّام در تدبیر آن باشند که نفسیرا هدایت نمایند و از رحیق عرفان
بچشانند این ایّام ایّامی نیست که فکر مختلف گردد و اراده متوقّف و قصد مشوب و
اعمال منافی و اخلاق خشن باید جمیع اعمال و افعال و اخلاق مثل نقطهٴ واحده
مشاهده شود یعنی مرجع کلّ معروف باشد تا مدلّ بر معروف حقیقی و حاکی از او
گردد کلّ تحت لحاظ عنایت بوده و هستند ما من احد الّا و قد فاز باکلیل الذّکر
من لدی اللّه المهیمن القیّوم

\end{otherlanguage}

In shāʾ Allāh jamīʿ-i dūstān-i ilāhī bi-qamīṣ-i taqdīs muzayyan shavand.
Bāyad bi-shaʾnī ẓāhir bāshand kih agar dar anjuman-i ʿālam bi-vaʿẓ va
nuṣḥ-i ʿibād qiyām namāyand va bi-aʿlā al-bayān tanaṭṭuq kunand nah khud
dar sirr khajil bāshand va nah man ʿalā al-arḍ qādir bar iʿtirāḍ. Ḥāl
bāyad dar layālī va ayyām dar tadbīr-i ān bāshand kih nafsī rā hidāyat
namāyand va az raḥīq-i ʿirfān bichashānand. Īn ayyām ayyāmī nīst kih
fikr mukhtalif gardad va irādih mutavaqqif va qaṣd mashūb va aʿmāl
munāfī va akhlāq khashin. Bāyad jamīʿ-i aʿmāl va afʿāl va akhlāq mithl-i
nuqṭih-yi vāḥidih mushāhadih shavad yaʿnī marjaʿ-i kull maʿrūf bāshad tā
mudill bar Maʿrūf-i Ḥaqīqī va ḥākī az ū gardad. Kull taḥt-i liḥāẓ-i
ʿināyat būdih va hastand. Mā min aḥadin illā wa-qad fāza bi-iklīli
al-dhikri min ladā Allāhi al-Muhaymini al-Qayyūm.

God willing, all the friends of God will be adorned with the garment of
sanctity. They should conduct themselves in such a way that, were they
to arise in the assembly of the world to admonish and counsel its people
and speak with the utmost eloquence, they would neither be inwardly
ashamed nor give anyone on earth grounds for objection. Now, night and
day, they should devise ways to guide a soul and give it a taste of the
choice wine of recognition. These are not days for divided thought,
faltering will, impure intent, inconsistent deeds, or harsh character.
All deeds, actions, and qualities should be seen as a single point; that
is, the focus of them all should be known, so that they may point to the
true Object of recognition and reflect Him. All have been and remain
beneath the gaze of His loving care. There is no one who has not
attained the crown of remembrance from God, the All-Protecting, the
Self-Subsisting.

\subsection{11}\label{section-2431}

\begin{otherlanguage}{persian}

حال آنچه لازمست باید احبّای الهی بآن تمسّک نمایند ایّام راحت و فرح و عزّت را
هم البتّه مشاهده کنند لازال میفرمایند اگر منادیرا بشناسند و لذّت بیان رحمن
را بیابند عالم را معدوم شمرند و بما یرتفع به الأمر ناظر و متشبّث و متمسّک
گردند

\end{otherlanguage}

Ḥāl ānchih lāzim ast bāyad aḥibbā-yi ilāhī bi-ān tamassuk namāyand.
Ayyām-i rāḥat va faraḥ va ʿizzat rā ham al-battih mushāhadih kunand. Lā
yazāl mī-farmāyand agar munādī rā bishināsand va lidhdhat-i bayān-i
Raḥmān rā biyāband ʿālam rā maʿdūm shumurand va bi-mā yartafiʿu bihi
al-amr nāẓir va mutashabbith va mutamassik gardand.

Now the loved ones of God must hold fast to what is necessary. They will
assuredly also witness days of ease, joy, and honor. He continually says
that, if they recognize the Summoner and taste the sweetness of the
All-Merciful's utterance, they will count the world as nothing and
direct their gaze toward whatever exalts the Cause, cleaving and holding
fast to it.

\subsection{12}\label{section-2432}

\begin{otherlanguage}{persian}

باری این عبد خدمت هر یک عرض سلام و تکبیر و فنا و نیستی میرساند و بکمال
عجز و ابتهال از معادن جواهر محبّت الهی استدعا مینماید که بمقامات عالیهٴ
خود ناظر باشند و آنچه سزاوار است عمل فرمایند العرض و العجز من العبد و
القبول من احبّآء اللّه العزیز المحبوب

\end{otherlanguage}

Bārī īn ʿabd khidmat-i har yik ʿarḍ-i salām va takbīr va fanā va nīstī
mī-rasānad va bi-kamāl-i ʿajz va ibtihāl az maʿādin-i javāhir-i
maḥabbat-i ilāhī istidʿā mī-namāyad kih bi-maqāmāt-i ʿāliyih-yi khud
nāẓir bāshand va ānchih sazāvār ast ʿamal farmāyand. Al-ʿarḍu wa-l-ʿajzu
mina al-ʿabdi wa-l-qabūlu min aḥibbāʾi Allāhi al-ʿAzīzi al-Maḥbūb.

In any case, this servant conveys to each one greetings of peace and
praise, in self-effacement and nothingness, and with the utmost
helplessness and supplication begs these mines of the jewels of divine
love to keep their lofty stations in view and act as befits them. The
plea and the helplessness are this servant's; acceptance rests with the
loved ones of God, the Almighty, the Beloved.

\subsection{13}\label{section-2433}

\begin{otherlanguage}{persian}

اینکه ذکر نفسین مقبلین جناب آقا سیّد محمّد و استاد محمّد قلی علیهما بهآء
اللّه فرموده بودند و هم‌چنین ذکر عمل ایشانرا فی سبیل اللّه در ساحت امنع
اقدس اعلی عرض شد نسبت بهر یک آفتاب عنایت و قمر رأفت و نجم شفقت و کلمه
علیا ظاهر و مشرق قوله جلّ و عزّ

\end{otherlanguage}

Īnkih dhikr-i nafsayn-i muqbilayn Janāb-i Āqā Siyyid Muḥammad va Ustād
Muḥammad-Qulī ʿalayhimā bahāʾu Allāh farmūdih būdand va hamchunīn
dhikr-i ʿamal-i īshān rā fī sabīli Allāh dar sāḥat-i amnaʿ-i aqdas-i
aʿlā ʿarḍ shud. Nisbat bi-har yik āftāb-i ʿināyat va qamar-i raʾfat va
najm-i shafaqat va kalimih-yi ʿulyā ẓāhir va mushriq. Qawluhu jalla
wa-ʿazz.

The mention you had made of the two devoted souls, the honored Āqā
Siyyid Muḥammad and Ustād Muḥammad-Qulī---may the glory of God rest upon
them both---and likewise of their work in the path of God, was presented
in the most exalted, most holy, most inaccessible Court. For each of
them the sun of loving care, the moon of tenderness, the star of
compassion, and the exalted Word appeared and shone forth. His
words---exalted and glorified be He:

\subsection{14}\label{section-2434}

\begin{otherlanguage}{arabic}

بسمی المنادی فی العالم

\end{otherlanguage}

Bismī al-Munādī fī al-ʿālam.

In My name, the Summoner in the world.

\subsection{15}\label{section-2435}

\begin{otherlanguage}{arabic}

یا محمّد قد ذکرت لدی المظلوم مرّة بعد مرّة و ذکرناک فیکلّ مرّة ما یجد منه
المخلصون عرف عنایتی ان افرح و کن من الشّاکرین و نذکرک فیهذا الحین کرّة
اخری فضلاً من لدنّا انّ ربّک لهو العزیز الکریم

\end{otherlanguage}

Yā Muḥammad, qad dhukirta ladā al-Maẓlūmi marratan baʿda marratin
wa-dhakarnāka fī kulli marratin mā yajidu minhu al-mukhliṣūna ʿarfa
ʿināyatī. An ifraḥ wa-kun mina al-shākirīn. Wa-nadhkuruka fī hādhā
al-ḥīni karratan ukhrā faḍlan min ladunnā. Inna rabbaka la-huwa
al-ʿAzīzu al-Karīm.

O Muḥammad, you have been mentioned before the Wronged One time after
time, and each time We have remembered you in words from which the
sincere perceive the fragrance of My loving care. Rejoice and be among
the thankful. At this moment We remember you once again as a bounty from
Our presence. Your Lord is indeed the Almighty, the All-Bountiful.

\subsection{16}\label{section-2436}

\begin{otherlanguage}{arabic}

ان احمد اللّه بما ذکر عملک فی الألواح و انزل لک ما یزیّنک بطراز الخلوص انّه
لهو المقتدر القدیر قد فاز ما عملته فی سبیلی برضائی و طراز قبولی انّ ربّک
لهو العلیم الحکیم

\end{otherlanguage}

An iḥmadi Allāha bi-mā dhakara ʿamalaka fī al-alwāḥi wa-anzala laka mā
yuzayyinuka bi-ṭirāzi al-khulūṣ. Innahu la-huwa al-Muqtadiru al-Qadīr.
Qad fāza mā ʿamiltahu fī sabīlī bi-riḍāʾī wa-ṭirāzi qabūlī. Inna rabbaka
la-huwa al-ʿAlīmu al-Ḥakīm.

Praise God, for He has mentioned your work in the Tablets and sent down
for you that which adorns you with the ornament of sincerity. He is
indeed the All-Powerful, the Omnipotent. What you have done in My path
has attained My good pleasure and the ornament of My acceptance. Your
Lord is indeed the All-Knowing, the All-Wise.

\subsection{17}\label{section-2437}

\begin{otherlanguage}{arabic}

طوبی لک و لعملک و لمن عمل بما امر به فی کتابی المبین بقیامکم علی خدمتی و
عملکم فی حبّی نطق الملأ الأعلی طوبی لمن وفی بمیثاق اللّه فی هذا الیوم
المبارک البدیع البهآء المشرق الظّاهر اللّائح اللّمیع علیک و علی کلّ عبد
مستقیم

\end{otherlanguage}

Ṭūbā laka wa-li-ʿamalika wa-li-man ʿamila bi-mā umira bihi fī kitābī
al-mubīn. Bi-qiyāmikum ʿalā khidmatī wa-ʿamalikum fī ḥubbī naṭaqa
al-malaʾu al-aʿlā: Ṭūbā li-man wafā bi-mīthāqi Allāhi fī hādhā al-yawmi
al-mubāraki al-badīʿ. Al-Bahāʾu al-mushriqu al-ẓāhiru al-lāʾiḥu
al-lamīʿu ʿalayka wa-ʿalā kulli ʿabdin mustaqīm.

Blessed are you and your work, and whoever has acted as commanded in My
clear Book. Through your arising to serve Me and your work out of love
for Me, the Concourse on high has proclaimed: ``Blessed is the one who
has been faithful to the Covenant of God on this blessed and wondrous
Day.'' May the glory, shining, manifest, radiant, and resplendent, rest
upon you and upon every steadfast servant.

\subsection{18}\label{section-2438}

\begin{otherlanguage}{persian}

و هذا ما نزّل من جبروت الفضل لجناب استاد محمّد قلی علیه بهآء اللّه قوله جلّ
کبریائه

\end{otherlanguage}

Wa-hādhā mā nuzzila min jabarūti al-faḍli li-janābi Ustād Muḥammad-Qulī
ʿalayhi bahāʾu Allāh. Qawluhu jalla kibriyāʾuh.

And this is what was sent down from the dominion of grace for the
honored Ustād Muḥammad-Qulī, may the glory of God rest upon him. His
words---exalted be His majesty:

\subsection{19}\label{section-2439}

\begin{otherlanguage}{arabic}

بسمی المهیمن علی الظّاهر و الباطن

\end{otherlanguage}

Bismī al-Muhaymini ʿalā al-ẓāhiri wa-l-bāṭin.

In My name, the Guardian over the outward and the inward.

\subsection{20}\label{section-2440}

\begin{otherlanguage}{arabic}

یا محمّد قبل قلی یذکرک المظلوم من شطره و یذکرک بعهده و ینزل لک ما عملته
متمسّکاً بکتابه طوبی لک بما وفیت و عملت بما امرت به فی حبّه کم من عبد تمسّک
بسبل الأوهام معرضاً عن سبیله و کم من عبد نبذ العالم عن ورائه مقبلاً الی
انوار وجهه ان افرح بما ایّدک علی العمل بما یقرّبک فیکلّ الأحوال الی نفسه
نعیماً لک بما اقبلت الی اللّه و وجدت عرف قمیصه انّا ذکرناک من قبل و من بعد
انّ ربّک یجزی من تمسّک بحبل عنایته البهآء المشرق من افق سمآء الفضل علیک و
علی من اتّبع امره انتهی

\end{otherlanguage}

Yā Muḥammad qabla Qulī, yadhkuruka al-Maẓlūm min shaṭrihi
wa-yudhakkiruka bi-ʿahdihi wa-yunazzilu laka mā ʿamiltahu mutamassikan
bi-kitābihi. Ṭūbā laka bi-mā wafayta wa-ʿamilta bi-mā umirta bihi fī
ḥubbihi. Kam min ʿabdin tamassaka bi-subuli al-awhām muʿriḍan ʿan
sabīlihi wa-kam min ʿabdin nabadha al-ʿālama warāʾihi muqbilan ilā
anwāri wajhihi. An ifraḥ bi-mā ayyadaka ʿalā al-ʿamali bi-mā yuqarribuka
fī kulli al-aḥwāl ilā nafsihi. Naʿīman laka bi-mā aqbalta ilā Allāh
wa-wajadta ʿarfa qamīṣihi. Innā dhakarnāka min qablu wa-min baʿd. Inna
rabbaka yajzī man tamassaka bi-ḥabli ʿināyatihi. Al-Bahāʾ al-mushriq min
ufuqi samāʾi al-faḍli ʿalayka wa-ʿalā man ittabaʿa amrahu. Intahā.

O Muḥammad-Qulī, the Wronged One remembers you from where He dwells,
reminds you of His covenant, and sends down for you what you have done
while holding fast to His Book. Blessed are you for having remained
faithful and done what you were commanded out of love for Him. How many
a servant has clung to the paths of delusion, turning away from His
path, and how many a servant has cast the world behind him, turning
toward the lights of His face! Rejoice that He has enabled you to do
what draws you closer to Him in every circumstance. Happy are you for
having turned to God and perceived the fragrance of His garment. We have
remembered you before and afterward. Your Lord rewards whoever holds
fast to the cord of His loving care. May the glory shining from the
horizon of the heaven of grace rest upon you and upon whoever follows
His command. End of quotation.

\subsection{21}\label{section-2441}

\begin{otherlanguage}{persian}

طوبی لهما و نعیماً لهما چه که عرف خلوص از ایشان بشهادة اللّه متضوّع و بروح
و ریحان عمل نمودند آنچه را بآن امر شدند انشآءاللّه اینمقام اعلی را باسم
حقّ جلّ جلاله حفظ فرمائید از حقّ منیع این خادم فانی سائل و آمل که ایشانرا
باستقامت کبری فائز فرماید انّه علی کلّشیء قدیر تکبیر و سلام خدمت ایشان
منوط باراده و عنایت آنمحبوبست

\end{otherlanguage}

Ṭūbā lahumā va naʿīman lahumā, che ke ʿarf-e khulūṣ az īshān
bi-shahādati Allāh mutaḍavviʿ va bi-rūḥ va rayḥān ʿamal namūdand ānche
rā bi-ān amr shudand. In shāʾ Allāh īn maqām-e aʿlā rā bi-ism-e Ḥaqq
jalla jalāluhu ḥifẓ farmāyīd. Az Ḥaqq-e manīʿ īn khādim-e fānī sāʾil va
āmil ke īshān rā bi-istiqāmat-e kubrā fāʾiz farmāyad. Innahu ʿalā kulli
shayʾin qadīr. Takbīr va salām khidmat-e īshān manūṭ bi-irādih va
ʿināyat-e ān maḥbūb ast.

Blessed and happy are they both, for, as God bears witness, the
fragrance of sincerity wafts from them, and with joy and gladness they
have done what they were commanded. God willing, may you preserve this
most exalted station in the name of the True One, exalted be His glory.
This servant, who is passing away, asks and hopes that the exalted True
One will grant them supreme steadfastness. He has power over all things.
Conveying greetings and salutations to them depends upon the will and
kindness of that beloved friend.

\subsection{22}\label{section-2442}

\begin{otherlanguage}{persian}

اینکه در بارهٴ ورقهٴ مخدّره ح‌و علیها بهآء اللّه مرقوم داشتید عرض شد
فرمودند نعم ما عمل من المعروف فی حقّها انشآءاللّه مؤیّد باشند بر اعمال طیّبه
و ما یکون حاکیاً عن الحکمة المحضة منتسبین انشآءاللّه بعنایت حقّ فائز باشند
و باستقامت کبری مزیّن تا جلاجل داهیه ایشانرا از صراط احدیّه منع ننماید

\end{otherlanguage}

Īnke dar bārih-yi varaqih-yi mukhaddarih Ḥāʾ-Vāv, ʿalayhā bahāʾu Allāh,
marqūm dāshtīd ʿarḍ shud. Farmūdand: niʿma mā ʿumila mina al-maʿrūfi fī
ḥaqqihā. In shāʾ Allāh muʾayyad bāshand bar aʿmāl-e ṭayyibih va mā
yakūnu ḥākiyan ʿani al-ḥikmati al-maḥḍah. Muntasibīn in shāʾ Allāh
bi-ʿināyat-e Ḥaqq fāʾiz bāshand va bi-istiqāmat-e kubrā muzayyan tā
jalājil-e dāhiyih īshān rā az ṣirāṭ-e aḥadiyyih manʿ nanamāyad.

What you wrote concerning the honored leaf Ḥāʾ-Vāv, may the glory of God
rest upon her, was presented. He said: ``Excellent is the kindness shown
to her.'' God willing, may they be assisted in performing good deeds and
whatever reflects pure wisdom. God willing, may the relatives receive
the loving care of the True One and be adorned with supreme
steadfastness, so that the clamors of calamity may not keep them from
the path of divine unity.

\subsection{23}\label{section-2443}

\begin{otherlanguage}{arabic}

یا علی قبل اکبر علیک بهائی و رحمتی و عنایتی ان انظر ثمّ اذکر فی امة رجعت
الی مقرّها انّها بدّلت النّور بالنّار و باعت ملک النّاس بقطعة من الألماس مع
کمال محبّت و اقبال و ادّعای عرفان داهیه بیک گل الماس او را فریب داد از حقّ
معرض و بغیر او مقبل اینست شأن دنیا و اهل آن از اوّل امر معرض بوده و ابداً
اطّلاع نداشته و ندارد چه که با ما نبوده خضوع و خشوع بیمعنی حضرات ق آن
نفوس را بر آن داشت که اظهار ایمان نمایند انّه یعلم ما فی الصّدور و ما فی
القلوب و هو العلیم الخبیر

\end{otherlanguage}

Yā ʿAlī qabla Akbar, ʿalayka bahāʾī wa-raḥmatī wa-ʿināyatī. An unẓur
thumma udhkur fī amatin rajaʿat ilā maqarrihā. Innahā baddalati al-nūra
bi-al-nār wa-bāʿat malika al-nāsi bi-qiṭʿatin mina al-almās. Bā kamāl-e
maḥabbat va iqbāl va iddiʿā-yi ʿirfān, dāhiyih bi-yik gul-e almās ū rā
farīb dād. Az Ḥaqq muʿriḍ va bi-ghayr-e ū muqbil. Īn ast shaʾn-e dunyā
va ahl-e ān. Az avval-e amr muʿriḍ būdih va abadan iṭṭilāʿ nadāshtih va
nadārad, che ke bā mā nabūdih. Khuḍūʿ va khushūʿ-e bī-maʿnī-yi ḥaḍarāt-e
Qāf ān nufūs rā bar ān dāsht ke iẓhār-e īmān namāyand. Innahu yaʿlamu mā
fī al-ṣudūri wa-mā fī al-qulūbi wa-huwa al-ʿAlīmu al-Khabīr.

O ʿAlī-Akbar, upon you be My glory, My mercy, and My loving care.
Consider, then speak of a woman who returned to her abode. She exchanged
light for fire and sold the King of humanity for a piece of diamond.
Despite all her love, devotion, and claims to recognition, the calamity
deceived her with a diamond flower. She has turned away from the True
One and toward another. Such is the condition of the world and its
people. From the beginning of the Cause she has turned away and has
never had, nor does she now have, any understanding, for she has not
been with Us. The meaningless deference and humility of the gentlemen of
Qāf induced those souls to profess belief. He knows what is within the
breasts and what is within the hearts, and He is the All-Knowing, the
All-Informed.

\subsection{24}\label{section-2444}

\begin{otherlanguage}{arabic}

اگر یک ورقه از اوراق در افعال و اعمال و حرکات او تفکّر نماید تنطق و تقول
انّی بریئة منها و عنکلّ ما عندها و تحفظ وجهها عن وجهها آن جناب میدانند قلم
اعلی در ذکر آن نفوس متوقّف بود ولکن نظر بظلم و تعدّی و شرارهٴ نفس و هوی از
قلم اعلی اینکلمات جاری و نازل و اینهم نظر بآنکه او را بشناسند و بر اعمال
و افعال او اطّلاع یابند اعاذنا اللّه و ایّاکم منها و امثالها انتهی

\end{otherlanguage}

Agar yik varaqih az avrāq dar afʿāl va aʿmāl va ḥarakāt-e ū tafakkur
namāyad, tanṭiqu wa-taqūlu innī barīʾatun minhā wa-ʿan kulli mā ʿindahā
wa-taḥfaẓu wajhahā ʿan wajhihā. Ān janāb mīdānand Qalam-e Aʿlā dar
dhikr-e ān nufūs mutavaqqif būd, va-lākin naẓar bi-ẓulm va taʿaddī va
sharārih-yi nafs va havā az Qalam-e Aʿlā īn kalimāt jārī va nāzil, va īn
ham naẓar bi-ānke ū rā bishināsand va bar aʿmāl va afʿāl-e ū iṭṭilāʿ
yāband. Aʿādhanā Allāhu wa-iyyākum minhā wa-amthālihā. Intahā.

If a single leaf among the leaves reflects upon her actions, deeds, and
conduct, she will speak and say, ``I am free of her and of everything
she possesses,'' and will shield her face from hers. You know that the
Most Exalted Pen had refrained from mentioning those souls. Yet, in view
of the oppression, transgression, and flame of selfish desire, these
words have flowed and descended from the Most Exalted Pen, and this too
is so that they may recognize her and become acquainted with her actions
and deeds. May God protect us and you from her and others like her. End
of quotation.

\subsection{25}\label{section-2445}

\begin{otherlanguage}{arabic}

اینفانی از حقّ جلّ جلاله سائل و آمل که دوستان آن ارض را بصراط مستقیم امانت
و دیانت و صدق و صفا هدایت فرماید تا کلّ در ظلّ حقّ و عنایتش مستریح باشند

\end{otherlanguage}

Īn fānī az Ḥaqq jalla jalāluhu sāʾil va āmil ke dūstān-e ān arḍ rā
bi-ṣirāṭ-e mustaqīm-e amānat va diyānat va ṣidq va ṣafā hidāyat farmāyad
tā kull dar ẓill-e Ḥaqq va ʿināyatash mustarīḥ bāshand.

This passing soul asks and hopes that the True One, exalted be His
glory, will guide the friends in that land to the straight path of
trustworthiness, integrity, truthfulness, and purity, so that all may
rest at ease beneath the shelter of the True One and His loving care.

\subsection{26}\label{section-2446}

\begin{otherlanguage}{persian}

وقتی اینکلمهٴ علیا از لسان مالک اسماء استماع شد فرمودند یا عبد حاضر از
برای هر شیء مبدئی بوده و هست اینکه عوام گفته و میگویند بی‌مایه خمیر فطیر
است حرفیست تمام و این قول حکیم عارف است الیوم اصحاب الهی مایهٴ احزاب
عالمند باید کلّ از امانت و صدق و استقامت و اعمال و اخلاق ایشان اقتباس
نمایند از حقّ بطلب عرف آگاهیرا بیابند و بما یرتفع به شؤونهم و مقاماتهم
تمسّک نمایند انتهی

\end{otherlanguage}

Vaqtī īn kalimih-yi ʿulyā az lisān-e Mālik-e Asmāʾ istimāʿ shud.
Farmūdand: yā ʿabd-e ḥāḍir, az barā-yi har shayʾ mabdaʾī būdih va hast.
Īnke ʿavāmm guftih va mīgūyand bī-māyih khamīr faṭīr ast ḥarfīst tamām
va īn qawl-e ḥakīm-e ʿārif ast. Al-yawm aṣḥāb-e ilāhī māyih-yi aḥzāb-e
ʿālamand. Bāyad kull az amānat va ṣidq va istiqāmat va aʿmāl va akhlāq-e
īshān iqtibās namāyand. Az Ḥaqq biṭalab ʿarf-e āgāhī rā biyāband va
bi-mā yartafiʿu bihi shuʾūnuhum wa-maqāmātuhum tamassuk namāyand.
Intahā.

On one occasion this exalted utterance was heard from the tongue of the
Lord of Names. He said: ``O servant who is present, everything has had
and has an originating principle. The saying that ordinary people have
repeated and still repeat, `Without leaven, dough remains unleavened,'
is entirely sound; it is the saying of a wise and discerning person.
Today the companions of God are the leaven of the peoples of the world.
All should draw upon their trustworthiness, truthfulness, steadfastness,
deeds, and character. Ask the True One that they may perceive the
fragrance of awareness and hold fast to what will elevate their standing
and stations.'' End of quotation.

\subsection{27}\label{section-2447}

\begin{otherlanguage}{persian}

اینکه بعضی از معرضین بیان اعتراض نموده‌اند با ادّعای حقیّت اصحابشان بفجور و
فحشا عامل و اصحاب نقطهٴ اولی بتقوی عامل و بطراز امانت و دیانت مزیّن
اینکلمه از صراط حقّ و راستی خارج چه که از ورود جمال قدم به عراق که این
عبد مطّلع است متّصلاً آیات و کلمات در سرّ سرّ نازل و ارسال میشد و جمیع امر
بما اراده اللّه و نهی عمّا نهی عنه فی الکتاب

\end{otherlanguage}

Īnke baʿḍī az muʿriḍīn-e Bayān iʿtirāḍ namūdih-and bā iddiʿā-yi
ḥaqqiyyat aṣḥābashān bi-fujūr va faḥshā ʿāmil va aṣḥāb-e Nuqṭih-yi Ūlā
bi-taqvā ʿāmil va bi-ṭirāz-e amānat va diyānat muzayyan, īn kalimih az
ṣirāṭ-e ḥaqq va rāstī khārij, che ke az vurūd-e Jamāl-e Qidam bi-ʿIrāq
ke īn ʿabd muṭṭaliʿ ast muttaṣilan āyāt va kalimāt dar sirr-e sirr nāzil
va irsāl mīshud va jamīʿ amr bi-mā arāda Allāh va nahy ʿammā nahā ʿanhu
fī al-Kitāb.

Some of the dissenters among the people of the Bayān have objected that,
despite the claim to truth, His companions practice wickedness and
indecency, whereas the companions of the Primal Point practice piety and
are adorned with trustworthiness and integrity. This statement departs
from the path of truth and honesty, for, as this servant knows, ever
since the Ancient Beauty arrived in Iraq, verses and utterances were
continually revealed and dispatched in the utmost secrecy, all enjoining
what God willed and forbidding what He prohibited in the Book.

\subsection{28}\label{section-2448}

\begin{otherlanguage}{persian}

چندی قبل لسان عظمت باینکلمهٴ علیا ناطق از عراق و ارض سرّ و ارض سجن متّصل
اهل بیان را نصیحت نمودیم و بما ینبغی لأمر اللّه امر فرمودیم هنوز بغی و
فساد مرتفع نشده اکثری آلوده‌اند و بتقدیس و تنزیه فائز نه انتهی

\end{otherlanguage}

Chandī qabl Lisān-e ʿAẓamat bi-īn kalimih-yi ʿulyā nāṭiq: az ʿIrāq va
Arḍ-e Sirr va Arḍ-e Sijn muttaṣil ahl-e Bayān rā naṣīḥat namūdīm va
bi-mā yanbaghī li-amri Allāh amr farmūdīm. Hanūz baghy va fasād murtafiʿ
nashudih. Aktharī ālūdih-and va bi-taqdīs va tanzīh fāʾiz nah. Intahā.

Some time ago the Tongue of Grandeur spoke these exalted words: ``From
Iraq, the Land of Mystery, and the Land of the Prison, We continually
counseled the people of the Bayān and commanded what befits the Cause of
God. Still, rebellion and corruption have not ceased. Most remain
defiled and have not attained purity and sanctity.'' End of quotation.

\subsection{29}\label{section-2449}

\begin{otherlanguage}{persian}

آنمحبوب خود مطّلعند که چه اعمال شنیعه مابین اینفرقه ظاهر فی‌الحقیقه انسان
قادر بر تقریر نه احدی اختیار زن و فرزند و مال و اموال خود را نداشت اکثری
باغی و یاغی و طاغی ولکن خود را از اصحاب خلّص الهی میدانستند معدودی مقدّس
از شئونات و شبهات دیده شدند و آنهم در اوّل امر و بعد خود آنجناب شاهد و
گواهند که چه کردند و چه شد بقسمیکه حضرت اعلی روح ما سواه فداه اظهار حیرت
فرمودند و در ظاهر چیزی نفرمودند چه که ناصر و معین موجود نه در ایّام آخر
لوحی مرقوم داشتند و منتهای شکایت از آن نفوس فرموده‌اند و از اکثری عدم
رضایت اظهار نموده‌اند ولکن ستر شد چه که اظهار آن مقتضی حکمت نه

\end{otherlanguage}

Ān maḥbūb khud muṭṭaliʿand ke che aʿmāl-e shanīʿih mābayn-e īn firqih
ẓāhir. Fī al-ḥaqīqih insān qādir bar taqrīr nah. Aḥadī ikhtiyār-e zan va
farzand va māl va amvāl-e khud rā nadāsht. Aktharī bāghī va yāghī va
ṭāghī va-lākin khud rā az aṣḥāb-e khullaṣ-e ilāhī mīdānistand. Maʿdūdī
muqaddas az shuʾūnāt va shubuhāt dīdih shudand va ān ham dar avval-e amr
va baʿd khud-e ān janāb shāhid va guvāhand ke che kardand va che shud,
bi-qismī ke Ḥaḍrat-e Aʿlā rūḥ-e mā sivāhu fidāh iẓhār-e ḥayrat farmūdand
va dar ẓāhir chīzī nafarmūdand che ke nāṣir va muʿīn mawjūd nah. Dar
ayyām-e ākhir lawḥī marqūm dāshtand va muntahā-yi shikāyat az ān nufūs
farmūdih-and va az aktharī ʿadam-e riḍāyat iẓhār namūdih-and va-lākin
satr shud che ke iẓhār-e ān muqtaḍī-yi ḥikmat nah.

That beloved friend is himself aware of the atrocious deeds that
appeared among this group. Truly, one cannot describe them. No one had
control over his own wife, children, possessions, or property. Most were
rebellious, insubordinate, and tyrannical, yet considered themselves
among the sincere companions of God. A few were seen to be free of such
conduct and suspicions, and that only at the beginning of the Cause.
Afterward, you yourself witnessed what they did and what happened, to
such an extent that the Most Exalted One---may the life of all else be
sacrificed for Him---expressed astonishment. Outwardly, however, He said
nothing, for there was no supporter or helper. In His last days He wrote
a Tablet in which He complained most strongly of those souls and
expressed dissatisfaction with the majority. It was, however, kept
concealed, since wisdom did not permit its disclosure.

\subsection{30}\label{section-2450}

\begin{otherlanguage}{persian}

ناس ضعیفند چنانچه مشاهده شد با وجود اشراقات آفتاب حقیقت در وسط زوال و
ظهور سدرهٴ منتهی از یمین بقعهٴ بیضا هر نفسی بهوای خود وهمی را بطراز یقین
جلوه میداد باری اگر در اینمقامات گفتگو نمائیم اوراق و دفاتر کفایت ننماید

\end{otherlanguage}

Nās ḍaʿīfand, chunānchih mushāhidih shud bā vujūd-e ishrāqāt-e Āftāb-e
Ḥaqīqat dar vasaṭ-e zavāl va ẓuhūr-e Sidrih-yi Muntahā az yamīn-e
Buqʿih-yi Bayḍā har nafsī bi-havā-yi khud vahmī rā bi-ṭirāz-e yaqīn
jilvih mīdād. Bārī agar dar īn maqāmāt guftugū namāyīm avrāq va dafātir
kifāyat nanamāyad.

People are weak. As was witnessed, despite the radiance of the Sun of
Truth at high noon and the appearance of the Tree beyond which there is
no passing on the right side of the White Spot, each person, following
his own desire, presented a delusion in the guise of certainty. In any
event, were we to discuss these matters, sheets and volumes would not
suffice.

\subsection{31}\label{section-2451}

\begin{otherlanguage}{persian}

و امّا در اینظهور اعظم حتّی اعدا شهادت داده و میدهند بر تقدیس و تنزیه و
علوّ و سموّ یکی از علمای یزد گفته بود در اوّل هرج و مرج اعمال اینطایفه خارج
از وصف و بیان حال مشاهده میشود در محفلها جمع میشوند و تلاوت آیات
مینمایند و مناجات میخوانند و بکمال امانت و دیانت ظاهر و مشهود معلوم
میشود از امر رئیس جدید است ملاحظه فرمائید چه مقدار از نفوس در اینظهور
رایگان شربت شهادت آشامیدند و چه بدست خود خود را فدا نمودند با آنکه حقّ
راضی نبود حال هم نفوسی هستند از جمیع عیوبات مقدّس و منزّه و مبرّا انّ ربّنا
العلیم لهو الشّاهد الخبیر

\end{otherlanguage}

Va ammā dar īn Ẓuhūr-e Aʿẓam ḥattā aʿdā shahādat dādih va mīdihand bar
taqdīs va tanzīh va ʿuluvv va sumuvv. Yikī az ʿulamā-yi Yazd guftih būd
dar avval-e harj va marj aʿmāl-e īn ṭāʾifih khārij az vaṣf va bayān. Ḥāl
mushāhidih mīshavad dar maḥfilhā jamʿ mīshavand va tilāvat-e āyāt
mīnamāyand va munājāt mīkhvānand va bi-kamāl-e amānat va diyānat ẓāhir
va mashhūd. Maʿlūm mīshavad az amr-e raʾīs-e jadīd ast. Mulāḥiẓih
farmāyīd che miqdār az nufūs dar īn Ẓuhūr rāygān sharbat-e shahādat
āshāmīdand va che bi-dast-e khud khud rā fidā namūdand bā ānke Ḥaqq rāḍī
nabūd. Ḥāl ham nufūsī hastand az jamīʿ-e ʿuyūbāt muqaddas va munazzah va
mubarrā. Inna rabbanā al-ʿAlīma la-huwa al-Shāhidu al-Khabīr.

In this Most Great Revelation, however, even the enemies have testified
and continue to testify to purity, sanctity, loftiness, and exaltation.
One of the divines of Yazd had said: ``At the beginning of the turmoil,
the deeds of this community were beyond description. Now they are seen
gathering in assemblies, reciting verses and prayers, and displaying the
utmost trustworthiness and integrity. Evidently this is at the command
of the new leader.'' Consider how many souls in this Revelation freely
drank the cup of martyrdom, and how many sacrificed themselves by their
own hands, although the True One did not approve. Even now there are
souls who are pure, sanctified, and free from every fault. Our Lord, the
All-Knowing, is indeed the Witness, the All-Informed.

\subsection{32}\label{section-2452}

\begin{otherlanguage}{persian}

این عبد یک کلمه از امانت و وفای بعهد ذکر مینماید که شاید بعضی متنبّه شوند
و بیقین بدانند که للّه تحت قبابه عباد اولی امانة عظیمة و اولی اعمال طیّبة
و اولی اخلاق روحانیّة در اسکندریّه چند نفر از احبّای الهی حسب الأمر بکسب
مشغول و کسب در اینظهور اعظم از عبادت محسوب شده نسئله بأن یوفّق الکلّ علی
ما اراد و

\end{otherlanguage}

Īn ʿabd yik kalimih az amānat va vafā-yi bi-ʿahd dhikr mīnamāyad ke
shāyad baʿḍī mutanabbih shavand va bi-yaqīn bidānand ke li-llāhi taḥta
qibābihi ʿibādun ulū amānatin ʿaẓīmatin wa-ulū aʿmālin ṭayyibatin wa-ulū
akhlāqin rūḥāniyyah. Dar Iskandariyyih chand nafar az aḥibbā-yi ilāhī
ḥasb al-amr bi-kasb mashghūl va kasb dar īn Ẓuhūr-e Aʿẓam az ʿibādat
maḥsūb shudih. Nasʾaluhu bi-an yuwaffiqa al-kulla ʿalā mā arāda va

This servant will relate one account of trustworthiness and fidelity to
a pledge, so that perhaps some may take heed and know with certainty
that beneath His canopies God has servants possessed of great
trustworthiness, good deeds, and spiritual qualities. In Alexandria,
several of the friends of God were engaged in earning a livelihood in
accordance with the command, and in this Most Great Revelation earning a
livelihood has been accounted worship. We ask Him to enable everyone to
fulfill what He has willed, and

\subsection{33}\label{section-2453}

\begin{otherlanguage}{persian}

آن کلمه اینکه جناب آقا شیخ علی ابن جناب حاجی عبدالرّحیم علیهما بهآء اللّه
به خرطوم تشریف بردند بجهت مؤانست با جناب محبوب قلب و فؤاد حاجی میرزا
حیدر علی و جناب حبیب روحانی حا و سین علیهما بهآء اللّه چه که اسمین
مسجونین نوشته بودند و یکنفر از اصحابرا خواسته بودند لأجل مؤانست و اطّلاع
بر امورات حادثهٴ ارض مقصود و ما حولها و هم‌چنین ایران جناب شیخ مذکور مصمّم
شدند و لوجه اللّه حرکت نمودند

\end{otherlanguage}

Ān kalimih īnke janāb-e Āqā Shaykh ʿAlī ibn-e janāb-e Ḥājī ʿAbd
al-Raḥīm, ʿalayhimā bahāʾu Allāh, bi-Kharṭūm tashrīf burdand bi-jihat-e
muʾānisat bā janāb-e maḥbūb-e qalb va fuʾād Ḥājī Mīrzā Ḥaydar ʿAlī va
janāb-e ḥabīb-e rūḥānī Ḥāʾ va Sīn, ʿalayhimā bahāʾu Allāh, che ke
ismayn-e masjūnayn nivishtih būdand va yik nafar az aṣḥāb rā khvāstih
būdand li-ajl-e muʾānisat va iṭṭilāʿ bar umūrāt-e ḥādithih-yi Arḍ-e
Maqṣūd va mā ḥawlahā va hamchunīn Īrān. Janāb-e Shaykh-e madhkūr
muṣammam shudand va li-wajhi Allāh ḥarakat namūdand.

The account is this: the honored Āqā Shaykh ʿAlī, son of the honored
Ḥājī ʿAbd al-Raḥīm, may the glory of God rest upon them both, went to
Khartoum to keep company with the beloved of the heart, the honored Ḥājī
Mīrzā Ḥaydar ʿAlī, and the spiritual friend, the honored Ḥāʾ and Sīn,
may the glory of God rest upon them both. For the two imprisoned names
had written requesting that one of the companions come to keep them
company and inform them of events in the Land of Desire and its
surroundings, as well as in Iran. The aforementioned Shaykh resolved to
go and set out for the sake of God.

\subsection{34}\label{section-2454}

\begin{otherlanguage}{arabic}

مدّتی چون اوراق اشجار محبّت مجتمع و از هواء لطیف و آب نورانی خفیف تر و
تازه بودند ناگاه باد خریف وزیدن گرفت و اسباب تفرقه بمیان آمد اسمین
مذکورین علیهما ثنآء اللّه باین شطر توجّه نمودند و بما هو المقصود فائز
گشتند و جناب شیخ در همان ارض بکسب مشغول

\end{otherlanguage}

Muddatī chun avrāq-e ashjār-e maḥabbat mujtamiʿ va az havā-yi laṭīf va
āb-e nūrānī-yi khafīf tar va tāzih būdand. Nāgāh bād-e kharīf vazīdan
girift va asbāb-e tafriqih bi-miyān āmad. Ismayn-e madhkūrayn, ʿalayhimā
thanāʾu Allāh, bi-īn shaṭr tavajjuh namūdand va bi-mā huwa al-maqṣūd
fāʾiz gashtand va janāb-e Shaykh dar hamān arḍ bi-kasb mashghūl.

For a time they were gathered like leaves on the trees of love, fresh
and flourishing from the gentle air and the light, luminous water.
Suddenly the autumn wind began to blow, and circumstances arose that
brought about separation. The two aforementioned names, may the praise
of God rest upon them both, turned toward this region and attained their
desire, while the honored Shaykh remained in that land earning a
livelihood.

\subsection{35}\label{section-2455}

\begin{otherlanguage}{persian}

بعد از سنین متوالیات حسب الأمر مبرم جناب حاجی عبداللّه علیه ٦٦٩ {[}بهآء
اللّه{]} بآنسمت توجّه نمودند لأجل مصاحبت با جناب شیخ و هم‌چنین اعانت نماید
ایشانرا در اشغال ظاهره و تأیید کند در مراجعت ولکن نظر بتشتّت امور بطول
انجامید تا آنکه حضرت شیخ علیه بهآء اللّه صعود فرمودند و بعد از صعود جناب
حاجی عبداللّه مذکور در تمشیت امور پرداخت

\end{otherlanguage}

Baʿd az sinīn-e mutavāliyāt ḥasb al-amr-e mubram janāb-e Ḥājī ʿAbd
Allāh, ʿalayhi 669 {[}bahāʾu Allāh{]}, bi-ān samt tavajjuh namūdand
li-ajl-e muṣāḥabat bā janāb-e Shaykh va hamchunīn iʿānat namāyad īshān
rā dar ashghāl-e ẓāhirih va taʾyīd kunad dar murājaʿat, va-lākin naẓar
bi-tashattut-e umūr bi-ṭūl anjāmīd tā ānke Ḥaḍrat-e Shaykh, ʿalayhi
bahāʾu Allāh, ṣuʿūd farmūdand va baʿd az ṣuʿūd janāb-e Ḥājī ʿAbd Allāh-e
madhkūr dar tamshiyat-e umūr pardākht.

After a number of successive years, in accordance with the binding
command, the honored Ḥājī ʿAbd Allāh, upon him be 669 {[}the glory of
God{]}, set out for that region to keep the honored Shaykh company,
assist him in his worldly occupations, and help him return. Owing to the
unsettled state of affairs, however, this was prolonged until the
honored Shaykh, may the glory of God rest upon him, ascended. After his
ascension, the aforementioned Ḥājī ʿAbd Allāh attended to the settling
of his affairs.

\subsection{36}\label{section-2456}

\begin{otherlanguage}{persian}

قنسل ایران که در آنمدینه بود تفصیل امورات مرحوم و هم‌چنین امتعه و اموالش
را بجناب کارپرداز اوّل دولت علیّهٴ ایران که مقیم قاهره است نوشت ایشان هم
رسمیکه حسب القانون مجری شده از اموال و ترکه طلب نمودند جناب اسمین اعلیین
محمّد و علی علیهما بهآء اللّه از اسکندریّه نوشتند که آنچه حقّ شماست میرسد
بعضی از تجّار قاهره که اقبال نموده‌اند بجناب کارپرداز مذکور گفتند که این
وجه بشما میرسد چه که قول حضرات صحیحست و مقدّس از شکّ و ریب

\end{otherlanguage}

Qonsol-e Irān ke dar ān madine būd tafṣil-e omūrāt-e marḥūm va hamchenin
amteʿe va amvālash rā be janāb-e kārpardāz-e avval-e Dowlat-e ʿAliyye-ye
Irān ke moqim-e Qāhere ast nevesht. Ishān ham rasmi ke ḥasb al-qānūn
mojrā shode az amvāl va tarake ṭalab namūdand. Janāb-e esmayn-e aʿlayayn
Moḥammad va ʿAli, ʿalayhimā bahāʾu'llāh, az Eskandariyye neveshtand ke
ānche ḥaqq-e shomāst mirasad. Baʿżi az tojjār-e Qāhere ke eqbāl
namūde-and be janāb-e kārpardāz-e madhkūr goftand ke in vajh be shomā
mirasad che ke qowl-e ḥażarāt ṣaḥiḥ ast va moqaddas az shakk va reyb.

The Iranian consul in that city wrote an account of the deceased's
affairs, as well as his goods and property, to the chief agent of the
Imperial Government of Iran, who resides in Cairo. The latter requested
from the property and estate the fee prescribed by law. The two most
exalted names, the honored Muḥammad and ʿAlī, may the glory of God rest
upon them both, wrote from Alexandria that whatever was rightfully his
would reach him. Some merchants of Cairo who had embraced the Cause told
the aforementioned agent, ``This sum will reach you, for these
gentlemen's word is true and beyond doubt or question.''

\subsection{37}\label{section-2457}

\begin{otherlanguage}{persian}

بعد اسکندریّه منقلب و قاهره مضطرب و دمدمهٴ مدافع مرتفع و طنطنهٴ صفوف و
الوف و بنادق ظاهر بشأنیکه جمیع ناس مضطرب و خائف و اضطراب تجاوز نمود تا
باین اراضی رسید لولا فضل اللّه و رحمته لخسفت الأرض بمن علیها چه که نظر
بغفلت و عصیان امّت مستحقّ بودند آنچه را که واقع شد

\end{otherlanguage}

Baʿd Eskandariyye monqaleb va Qāhere moẓṭareb va damdame-ye madāfeʿ
mortafeʿ va ṭanṭane-ye ṣofūf va olūf va banādeq ẓāher be shaʾni ke
jamiʿ-e nās moẓṭareb va khāʾef va eżṭerāb tajāvoz namūd tā be in arāżi
rasid. Lawlā faḍlu'llāhi wa raḥmatuhu la-khasafati al-arḍu bi-man
ʿalayhā, che ke naẓar be ghaflat va ʿeṣyān-e ommat mostaḥaqq būdand
ānche rā ke vāqeʿ shod.

Then Alexandria was thrown into turmoil and Cairo into agitation. The
thunder of cannons rose, and the din of ranks, multitudes, and rifles
broke out, so that everyone was alarmed and afraid. The unrest spread
until it reached these lands. Were it not for God's grace and mercy, the
earth would have swallowed up those upon it, for, in view of the
people's heedlessness and rebellion, they deserved what befell them.

\subsection{38}\label{section-2458}

\begin{otherlanguage}{persian}

باری اسکندریّه در ساعات معدوده عالیها سافلها مشاهده گشت چنانچه از قبل در
بعضی از عرایض این فانی ذکر نموده و هم‌چنین از قلم اعلی نازل احبّای
اسکندریّه و قاهره عریضه بساحت اقدس ارسال نمودند و تکلیف خود را خواستند و
هم‌چنین بحضرت اسم‌اللّه م‌ه علیه منکلّ بهآء ابهاه که در آن اوقات در این ارض
تشریف داشتند عریضه ارسال نمودند فرمودند احبّای الف و سین باقتضا حرکت
نمایند و بشطر اقدس توجّه کنند و احبّای مصر در مقرّ خود ساکن و مستریح باشند

\end{otherlanguage}

Bāri Eskandariyye dar sāʿāt-e maʿdūde ʿāliyahā sāfilahā moshāhede gasht
chenānche az qabl dar baʿżi az ʿarāyeż in fāni dhekr namūde va hamchenin
az Qalam-e Aʿlā nāzel. Aḥebbā-ye Eskandariyye va Qāhere ʿariże be
sāḥat-e aqdas ersāl namūdand va taklif-e khod rā khāstand va hamchenin
be ḥażrat-e Esmo'llāh M-H, ʿalayhi min kulli bahāʾin abhāhu, ke dar ān
owqāt dar in arż tashrif dāshtand ʿariże ersāl namūdand. Farmūdand
aḥebbā-ye Alef va Sin be eqteżā ḥarakat namāyand va be shaṭr-e aqdas
tavajjoh konand va aḥebbā-ye Meṣr dar maqarr-e khod sāken va mostariḥ
bāshand.

In short, within a few hours Alexandria was seen turned upside down, as
this mortal one has previously mentioned in some letters, and as has
also been revealed by the Most Exalted Pen. The friends in Alexandria
and Cairo sent a petition to the Most Holy Court asking what they should
do. They also sent a petition to his honor Ismu'llāh M-H, upon him be
the most resplendent of all glory, who was then staying in this land.
The instruction was given that the friends of Alif and Sīn should act as
circumstances required and turn toward the Most Holy Presence, while the
friends in Cairo should remain where they were and be at ease.

\subsection{39}\label{section-2459}

\begin{otherlanguage}{persian}

و بعد حضرات بارض مقصود وارد و غافلین و خائنین اسکندریّه را آتش زدند و
اموال ناس را از هر قبیل چه سوختند و چه بتاراج بردند و اموال دوستانهم
جمیع بتاراج رفت و مخزن ایشان منهدم و ویران و طرف خرطوم هم نار وغی مشتعل
و شخصی علم عصیان برافراخت نظر باغتشاش مصر آن اطراف جمیع منقلب و مضطرب و
جناب حاجی عبداللّه در آنجا محصور تا آنکه حکم خروج رسید و متوکّلاً علی اللّه
بیرون آمد و از صدهزار بلا حقّ جلّ جلاله حفظش فرمود و الحمد للّه بارض مقصود
وارد و مشرّف

\end{otherlanguage}

Va baʿd ḥażarāt be arż-e maqṣūd vāred va ghāfelin va khāʾenin
Eskandariyye rā ātash zadand va amvāl-e nās rā az har qabil che sūkhtand
va che be tārāj bordand va amvāl-e dūstān ham jamiʿ be tārāj raft va
makhzan-e ishān monhadem va virān va ṭaraf-e Kharṭūm ham nār-e vaghā
moshtaʿel va shakhṣi ʿalam-e ʿeṣyān bar afrākht. Naẓar be eghteshāsh-e
Meṣr ān aṭrāf jamiʿ monqaleb va moẓṭareb va janāb-e Ḥāji ʿAbdo'llāh dar
ānjā maḥṣūr tā ānke ḥokm-e khorūj rasid va motavakkelan ʿala'llāh birūn
āmad va az ṣad hezār balā Ḥaqq jalla jalāluhu ḥefẓash farmūd va al-ḥamdu
li'llāh be arż-e maqṣūd vāred va mosharraf.

After these gentlemen arrived in the land of their desire, the heedless
and treacherous set Alexandria on fire, burning or plundering people's
possessions of every kind. All the friends' property was also plundered,
and their warehouse was demolished and laid waste. In the region of
Khartoum, too, the fire of war was blazing, and someone raised the
banner of rebellion. Owing to the disturbances in Egypt, all those
regions were in turmoil and agitation, and the honored Ḥājī ʿAbd Allāh
was besieged there until the order to leave arrived. He came out,
placing his trust in God, and God, exalted be His glory, preserved him
from a hundred thousand calamities. Praise be to God, he reached the
land of his desire and was admitted to the sacred presence.

\subsection{40}\label{section-2460}

\begin{otherlanguage}{persian}

و امّا حضرات اسکندریّه مع آنکه کمال افلاس موجود و جناب کارپرداز قاهره هم
معزول و خبر ورود کارپرداز جدید در اسلامبل واصل مع‌ذلک نظر بوعده‌ئی که شده
بود جناب آقا سیّد علی علیه بهآء اللّه در صدد آن برآمدند که بوعده وفا
نمایند و از فضل و عنایت حقّ وجهی موجود شد او را اخذ نمود و روانهٴ مصر گشت
و مبلغ پنجاه جنیه انگلیزی برد و تسلیم نمود

\end{otherlanguage}

Va ammā ḥażarāt-e Eskandariyye maʿ-e ānke kamāl-e eflās mowjūd va
janāb-e kārpardāz-e Qāhere ham maʿzūl va khabar-e vorūd-e kārpardāz-e
jadid dar Eslāmbol vāṣel maʿa dhālek naẓar be vaʿde-i ke shode būd
janāb-e Āqā Seyyed ʿAli, ʿalayhi bahāʾu'llāh, dar ṣadad-e ān bar āmadand
ke be vaʿde vafā namāyand va az fażl va ʿenāyat-e Ḥaqq vajhi mowjūd shod
ū rā akhdh namūd va ravāne-ye Meṣr gasht va mablagh-e panjāh joneyh-e
englisi bord va taslim namūd.

As for the gentlemen from Alexandria, although they were utterly
destitute, the agent in Cairo had been dismissed, and news had arrived
of the new agent's arrival in Istanbul, nevertheless, in view of the
promise that had been made, the honored Āqā Siyyid ʿAlī, may the glory
of God rest upon him, sought to fulfill it. Through God's grace and
favor, a sum became available. He took it and set out for Cairo,
carrying fifty British pounds, which he handed over.

\subsection{41}\label{section-2461}

\begin{otherlanguage}{persian}

مؤمنین تجّار و خود جناب کارپرداز بسیار متحیّر و متعجّب و اوّل کارپرداز از
اخذ آن ابا نمود قسم یاد کرد که من حرفی ندارم این صفات شما مرا کافیست و
در مدح و ثنای شما همین قدر بس که همان نفوسیکه شما را دشمن میدارند بمدح
شما ناطقند باری این فقره سبب اعلاء شد در آنجهات و جناب آقا سیّد علی علیه
بهآء اللّه نوشته‌اند شرحی از محبّت دوستان مصر بقسمیکه طیور حبّ در افئدهٴ
اولیا بحرکت و طیران میآید محبّت للّه فی‌الحقیقه حیات حقیقی بخشد و مرده را
زنده نماید

\end{otherlanguage}

Moʾmenin-e tojjār va khod-e janāb-e kārpardāz besyār motaḥayyer va
motaʿajjeb va avval kārpardāz az akhdh-e ān ebā namūd. Qasam yād kard ke
man ḥarfi nadāram in ṣefāt-e shomā marā kāfist va dar madḥ va thanā-ye
shomā hamin qadr bas ke hamān nofūsi ke shomā rā doshman midārand be
madḥ-e shomā nāṭeqand. Bāri in faqare sabab-e eʿlā shod dar ān jahāt va
janāb-e Āqā Seyyed ʿAli, ʿalayhi bahāʾu'llāh, neveshte-and sharḥi az
maḥabbat-e dūstān-e Meṣr be qesmi ke ṭoyūr-e ḥobb dar afʾede-ye owliyā
be ḥarakat va ṭayarān miāyad. Maḥabbat li'llāh fi'l-ḥaqiqe ḥayāt-e
ḥaqiqi bakhshad va morde rā zende namāyad.

The believing merchants and the agent himself were greatly astonished
and amazed. At first the agent refused to accept it, swearing, ``I have
no claim. These qualities of yours are enough for me. It is sufficient
praise of you that the very people who regard you as enemies speak in
your praise.'' In short, this affair brought honor to the Cause in those
regions. The honored Āqā Siyyid ʿAlī, may the glory of God rest upon
him, has written an account of the love of the friends in Cairo such
that the birds of love in the hearts of the faithful stir and take
flight. Love for the sake of God truly bestows real life and brings the
dead to life.

\subsection{42}\label{section-2462}

\begin{otherlanguage}{persian}

این وفای بعهد ایشان از سلطان اعمال عند غنیّ متعال مذکور و مسطور از حقّ جلّ
جلاله میطلبم احبّای آن ارض و اراضی اخری را مؤیّد فرماید بر امانت و راستی

\end{otherlanguage}

In vafā-ye be ʿahd-e ishān az solṭān-e aʿmāl ʿenda Ghaniyy-e Motaʿāl
madhkūr va masṭūr. Az Ḥaqq jalla jalāluhu miṭalabam aḥebbā-ye ān arż va
arāżi-ye okhrā rā moʾayyad farmāyad bar amānat va rāsti.

Their fulfillment of this pledge is mentioned and recorded before the
All-Sufficient, the Most Exalted, as a sovereign among deeds. I ask God,
exalted be His glory, to confirm the friends in that land and in other
lands in trustworthiness and truthfulness.

\subsection{43}\label{section-2463}

\begin{otherlanguage}{arabic}

ارض سجن بمقامی مرتفع که جمیع اطراف باو ناظر بودند و هر یک از اصحاب اگر
چیزی طلب مینمودند لأجل بیع و شری از مدن قریبه و بعیده بأنّه هو حاضر جواب
میشنیدند و بعد بخیانت بعضی از نفوس و ظلمت ما فی قلوبهم از مقام اوّل بمقام
ثانی راجع بل ثالث و مع‌ذلک نظر بآنکه مقرّ کرسی واقع علوّ و سموّ مقامات احباب
علی قدر معلوم مابین عباد باقی و برقرار کار یک عیب دیگر پیدا کرده هر
نفسیکه بسیار خبیث واقع میشود در آخر مقام دعوی حقّیّت هم مینماید ید قدرت
انشآءاللّه عباد خود را حفظ فرماید و بما یحبّ و یرضی فائز نماید

\end{otherlanguage}

Arż-e Sejn be maqāmi mortafeʿ ke jamiʿ-e aṭrāf be ū nāẓer būdand va har
yek az aṣḥāb agar chizi ṭalab minamūdand li-ajl-e beyʿ va sherā az
modon-e qaribe va baʿide bi-annahu huwa ḥāḍir javāb mishenidand va baʿd
be khiyānat-e baʿżi az nofūs va ẓolmat-e mā fi qolūbehem az maqām-e
avval be maqām-e thāni rājeʿ bal thāleth va maʿa dhālek naẓar be ānke
maqarr-e korsi vāqeʿ ʿolovv va somovv-e maqāmāt-e aḥbāb ʿalā qadr-e
maʿlūm mābeyn-e ʿebād bāqi va barqarār. Kār yek ʿeyb-e digar peydā karde
har nafsi ke besyār khabith vāqeʿ mishavad dar ākhar maqām daʿvā-ye
ḥaqqiyyat ham minamāyad. Yad-e qodrat en shāʾa'llāh ʿebād-e khod rā ḥefẓ
farmāyad va be mā yoḥebb va yarżā fāʾez namāyad.

The Prison Land had risen to such a station that all regions looked
toward it. If any of the companions requested something from cities near
or far for purposes of trade, they would hear the answer, ``It is
available.'' Then, through the treachery of certain souls and the
darkness of what was in their hearts, it fell from the first rank to the
second, indeed to the third. Nevertheless, since it is the seat of the
Throne, the loftiness and sublimity of the friends' stations remain
established among the people to a certain degree. Another defect has
arisen in this affair: anyone who proves exceedingly malicious
eventually even lays claim to divinity. God willing, may the hand of
power preserve His servants and enable them to attain what He loves and
approves.

\subsection{44}\label{section-2464}

\begin{otherlanguage}{persian}

در جمیع احوال جمیع جوارح و اعضا باید بکلمهٴ لا حول و لا قوّة الّا باللّه
ناطق باشد در این اثنا لدی العرش حاضر و مقصود غیب و شهود باینکلمه ناطق
بنویس بجناب علی قبل اکبر علیه بهائی و رحمتی و عنایتی یا علی مر احبّائی
بما یرتفع به مقاماتهم بأمری قل یا اولیآء الرّحمن فی البلدان سخّروا مدائن
القلوب بجنود الأعمال و الأخلاق هذا ما حکم به مولی العالم حین استوائه علی
العرش الأعظم طوبی لعبد انجذب بندائی و عمل بما امر به فی کتابی المبین
انتهی

\end{otherlanguage}

Dar jamiʿ-e aḥvāl jamiʿ-e javāreḥ va aʿżā bāyad be kalame-ye lā ḥawla wa
lā quwwata illā bi'llāh nāṭeq bāshad. Dar in athnā ladā al-ʿarsh ḥāżer
va Maqṣūd-e gheyb va shohūd be in kalame nāṭeq: Benevis be janāb-e ʿAli
qabl-e Akbar, ʿalayhi bahāʾī wa raḥmatī wa ʿināyatī. Yā ʿAliyy, mur
aḥibbāʾī bi-mā yartafiʿu bihi maqāmātuhum bi-amrī. Qul yā awliyāʾa
al-Raḥmāni fi'l-buldān, sakhkhirū madāʾina al-qulūbi bi-junūdi al-aʿmāli
wa'l-akhlāq. Hādhā mā ḥakama bihi Mawlā al-ʿālami ḥīna istiwāʾihi ʿalā
al-ʿarshi al-aʿẓam. Ṭūbā li-ʿabdin injadhaba bi-nidāʾī wa ʿamila bi-mā
umira bihi fī kitābiya al-mubīn. Entahā.

In every circumstance, all one's limbs and members should proclaim the
words, ``There is no power or strength except through God.'' At this
moment I was present before the Throne, and the Desire of the unseen and
the seen spoke these words: ``Write to the honored ʿAlī before Akbar,
upon him be My glory, My mercy, and My loving-kindness: O ʿAlī, by My
command instruct My loved ones in what will exalt their stations. Say: O
friends of the All-Merciful throughout the lands! Conquer the cities of
hearts with the armies of deeds and character. This is what the Lord of
the world has decreed upon His establishment on the Most Great Throne.
Blessed is the servant who is enraptured by My call and acts according
to what he has been commanded in My perspicuous Book.'' End of
quotation.

\subsection{45}\label{section-2465}

\begin{otherlanguage}{persian}

از این قبیل بیانات محکمات لاتحصی از لسان مالک وری استماع شد طوبی
للعاملین

\end{otherlanguage}

Az in qabil bayānāt-e moḥkamāt lā toḥṣā az lesān-e Mālek-e varā estemāʿ
shod. Ṭūbā li'l-ʿāmilīn.

Countless firm utterances of this kind were heard from the tongue of the
Lord of creation. Blessed are those who act.

\subsection{46}\label{section-2466}

\begin{otherlanguage}{persian}

اینکه در بارهٴ محبوب فؤاد جناب حاجی میرزا حیدر علی علیه بهآء اللّه مرقوم
داشتید و ذکر خدمات ایشانرا فی سبیل اللّه فرمودند بعد از عرض در ساحت اقدس
اینکلمه مکرّر از لسان عظمت جاری فرمودند علیه بهآء اللّه علیه بهآء اللّه

\end{otherlanguage}

Inke dar bāre-ye maḥbūb-e foʾād janāb-e Ḥāji Mirzā Ḥeydar ʿAli, ʿalayhi
bahāʾu'llāh, marqūm dāshtid va dhekr-e khadamāt-e ishān rā fi sabil-e
Allāh farmūdand baʿd az ʿarż dar sāḥat-e aqdas in kalame mokarrar az
lesān-e ʿaẓamat jāri farmūdand: ʿAlayhi bahāʾu'llāh, ʿalayhi
bahāʾu'llāh.

Concerning what you wrote about the beloved of the heart, the honored
Ḥājī Mīrzā Ḥaydar ʿAlī, may the glory of God rest upon him, and your
mention of his services in the path of God: after this was presented in
the Most Holy Court, these words repeatedly flowed from the tongue of
grandeur: ``Upon him be the glory of God. Upon him be the glory of
God.''

\subsection{47}\label{section-2467}

\begin{otherlanguage}{persian}

خبر مازندران را آنمحبوب بتفصیل مرقوم داشته‌اند و این ایّام از هر شطری
عرایض رسیده و حاوی ذکر خدمت ایشان در ق‌و ولکن بعد از عرض بر حسب ظاهر چند
یوم در اینفقره امری از مطلع ظاهر نه و بعد یومی از ایّام اینکلمهٴ علیا از
مشرق بیان مالک اسماء ظاهر قوله جلّت عظمته و جلّ کبریائه یا عبد حاضر اگر
بگوئیم بحکمت واقع شده مشاهده میشود قباع خنازیر و ضوضای جهلا مرتفع و اگر
تمام هم بجهة ضوضا ناظر باشیم ستر این خدمت باهر را دوست نداشته و نداریم
لذا از مصدر امر الهی این حکم جاری میگوئیم بنطق و بیان ملجأ بوده و حقّ جلّ
جلاله نظر بفضل و عنایت او را محلّ و مصدر اعلای کلمه فرموده و این اعلاء
اگرچه سبب ضوضا شد ولکن نظر بالجاء بأسی نبوده و نیست بلکه لدی العرش مقبول
انّا زیّنّا عمله بطراز القبول و ذکرناه فی الواح شتّی انّ ربک لهو العطوف
الغفور الرّحیم الکریم انتهی

\end{otherlanguage}

Khabar-e Māzandarān rā ān maḥbūb be tafṣil marqūm dāshte-and va in ayyām
az har shaṭri ʿarāyeż reside va ḥāvi-ye dhekr-e khedmat-e ishān dar Q-V
va lāken baʿd az ʿarż bar ḥasb-e ẓāher chand yowm dar in faqare amri az
maṭlaʿ ẓāher na va baʿd yowmi az ayyām in kalame-ye ʿolyā az mashreq-e
bayān-e Mālek-e asmāʾ ẓāher, qawluhu jallat ʿaẓamatuhu wa jalla
kibriyāʾuhu: Yā ʿabd-e ḥāżer, agar begūʾim be ḥekmat vāqeʿ shode
moshāhede mishavad qobāʿ-e khanāzir va żowżā-ye johalā mortafeʿ va agar
tamām ham be jahat-e żowżā nāẓer bāshim satr-e in khedmat-e bāher rā
dūst nadāshte va nadārim. Ledhā az maṣdar-e amr-e elāhi in ḥokm jāri:
Migūʾim be noṭq va bayān moljaʾ būde va Ḥaqq jalla jalāluhu naẓar be
fażl va ʿenāyat ū rā maḥall va maṣdar-e eʿlā-ye kalame farmūde va in
eʿlā agarche sabab-e żowżā shod va lāken naẓar be eljāʾ baʾsi nabūde va
nist balke ladā al-ʿarsh maqbūl. Innā zayyannā ʿamalahu bi-ṭirāzi
al-qabūli wa dhakarnāhu fī alwāḥin shattā. Inna rabbaka la-huwa
al-ʿAṭūfu al-Ghafūru al-Raḥīmu al-Karīm. Entahā.

That beloved one has written the news of Māzandarān in detail. In recent
days letters have arrived from every direction, mentioning his service
in Q-V. After these were presented, however, for several days no command
concerning this matter outwardly appeared from the Dayspring. Then, one
day, these exalted words appeared from the dawning-place of the
utterance of the Lord of Names---His words, exalted be His grandeur and
magnified be His majesty: ``O servant in attendance! If We say that it
was done with wisdom, it is evident that the grunting of swine and the
clamor of the ignorant have arisen. Yet if We were to consider only the
clamor, We have not wished, nor do We wish, to conceal this manifest
service. Therefore, from the Source of the divine command this decree
proceeds: We say that he was compelled to speak and explain, and God,
exalted be His glory, through His grace and favor, made him the locus
and source of the exaltation of the Word. Although this exaltation
caused an uproar, since he was compelled, there was and is no blame;
rather, it is accepted before the Throne. We have adorned his deed with
the ornament of acceptance and have mentioned him in various Tablets.
Your Lord is indeed the Compassionate, the Ever-Forgiving, the Merciful,
the Generous.'' End of quotation.

\subsection{48}\label{section-2468}

\begin{otherlanguage}{persian}

عرض اینفانی آنکه صورت این عریضه اگر نوشته شود و بتصحیح فائز و باطراف
ارسال گردد خالی از نفع نبوده و نیست چه که مزیّن است بآیات و کلمات حقّ جلّ
جلاله که فی‌الحقیقه سبب اشتعال عالم و امم است شاید حلاوت آنرا بیابند و
بما یحبّه اللّه قیام نمایند و شاید عرف اعمال و اخلاق طیّبه در اطراف متضوّع
شود و ناس غافل را آگاهی بخشد الأمر بیده و ما انا الّا عبد منیب

\end{otherlanguage}

ʿArż-e in fāni ānke ṣūrat-e in ʿariże agar neveshte shavad va be taṣḥiḥ
fāʾez va be aṭrāf ersāl gardad khāli az nafʿ nabūde va nist che ke
mozayyan ast be āyāt va kalemāt-e Ḥaqq jalla jalāluhu ke fi'l-ḥaqiqe
sabab-e eshteʿāl-e ʿālam va omam ast. Shāyad ḥalāvat-e ān rā beyāband va
be mā yoḥebbuhu Allāh qiyām namāyand va shāyad ʿarf-e aʿmāl va akhlāq-e
ṭayyebe dar aṭrāf motażavveʿ shavad va nās-e ghāfel rā āgāhi bakhshad.
Al-amru bi-yadihi wa mā anā illā ʿabdun munīb.

This mortal one submits that it would not be without benefit if a copy
of this letter were made, corrected, and sent to various regions, for it
is adorned with the verses and words of God, exalted be His glory, which
truly set the world and its peoples aflame. Perhaps they will taste its
sweetness and arise to do what God loves. Perhaps the fragrance of good
deeds and noble character will spread abroad and awaken the heedless.
The command is in His hand, and I am only a servant who turns to Him.

\subsection{49}\label{section-2469}

\begin{otherlanguage}{persian}

و آنچه از الواح مخصوص احبّای ارض میم خواسته‌اند همان ایّام نازل و نظر
بانقلابات و حکمت در ارسال آن تأخیر رفت این ایّام امر بارسال فرمودند و
انشآءاللّه ارسال میشود آنچه آنمحبوب فؤاد در بارهٴ ایشان مرقوم فرمودند
فی‌الحقیقه سبب مسرّت قلب شد و عالم عالم محبّت ظاهر انشآءاللّه آنمحبوب و
ایشان در جمیع احیان بما اراده اللّه ناطق و بر خدمت امر قائم باشند

\end{otherlanguage}

Va ānche az alvāḥ-e makhṣūṣ-e aḥebbā-ye arż-e Mim khāste-and hamān ayyām
nāzel va naẓar be enqelābāt va ḥekmat dar ersāl-e ān taʾkhir raft. In
ayyām amr be ersāl farmūdand va en shāʾa'llāh ersāl mishavad. Ānche ān
maḥbūb-e foʾād dar bāre-ye ishān marqūm farmūdand fi'l-ḥaqiqe sabab-e
masarrat-e qalb shod va ʿālam ʿālam-e maḥabbat ẓāher. En shāʾa'llāh ān
maḥbūb va ishān dar jamiʿ-e aḥyān be mā arādahu Allāh nāṭeq va bar
khedmat-e amr qāʾem bāshand.

As for the Tablets requested specifically for the friends in the land of
Mīm, they were revealed in those same days, but owing to the upheavals
and considerations of wisdom, their dispatch was delayed. In recent days
the order has been given to send them, and, God willing, they will be
sent. What that beloved of the heart wrote about them truly brought joy
to the heart, and a world of love became manifest. God willing, may that
beloved one and those friends at all times speak what God wills and
arise to serve the Cause.

\subsection{50}\label{section-2470}

\begin{otherlanguage}{persian}

و دستخطّ ثانی آنحضرت که تاریخ آن غرّهٴ شهر رمضان المبارک بود گوش تازه
بخشید و بنغمهٴ جدیده در محبّت مالک احدیّه مترنّم و بعد از استماع و ادراک
قصد مقام مقدّس از ذکر و بیان نموده و تفصیل بعد از اذن عرض شد اذاً جری فرات
الرّحمة و الألطاف و ماج بحر الفضل علی شأن لا یذکر بالبیان قال و قوله الحقّ

\end{otherlanguage}

Va dastkhaṭṭ-e thāni-ye ān ḥażrat ke tārikh-e ān ghorre-ye shahr-e
Ramażān al-mobārak būd gūsh-e tāze bakhshid va be naghme-ye jadide dar
maḥabbat-e Mālek-e aḥadiyye motarannem va baʿd az estemāʿ va edrāk
qaṣd-e maqām-e moqaddas az dhekr va bayān namūde va tafṣil baʿd az edhn
ʿarż shod. Idhan jarā Furātu al-raḥmati wa'l-alṭāfi wa māja baḥru
al-faḍli ʿalā shaʾnin lā yudhkaru bi'l-bayān. Qāla wa qawluhu al-ḥaqq.

Your second letter, dated the first of the blessed month of Ramadan,
bestowed fresh hearing and sang a new melody of love for the Lord of
oneness. After hearing and grasping it, I made my way to the station
sanctified above mention and description, and, permission having been
granted, presented the details. Then the Euphrates of mercy and
loving-kindness flowed, and the ocean of grace surged in a manner that
words cannot describe. He said---and His word is the truth:

\subsection{51}\label{section-2471}

\begin{otherlanguage}{arabic}

بسمی المظلوم الفرید

\end{otherlanguage}

Bi-smī al-Maẓlūmi al-Farīd.

In My name, the Wronged One, the Peerless.

\subsection{52}\label{section-2472}

\begin{otherlanguage}{arabic}

یا ایّها النّاظر الی شطر السّجن و النّاطق بثنآء المقصود و الشّارب من رحیقی
المختوم علیک بهآء اللّه المهیمن القیّوم قد حضر العبد الحاضر بما ارسلته
الیه و عرض تلقآء المظلوم ما غنّت حمامة فؤادک فی ثنآء اللّه و ذکره طوبی
للسانک و لبیانک و لمن یسمع قولک فی امر اللّه العزیز الودود قد انزلنا علی
المقرّبین ما لو یلقی علی التّراب لتطلع منه شمس العلم و البیان کذلک نطقت
یراعة الرّحمن و النّاس اکثرهم لا یفقهون انت الّذی حضرت و رأیت البحر و
امواجه و الشّمس و انوارها و سمعت ما تکلّم به مکلّم الطّور اذ کان مستویاً علی
عرش الظّهور انّ ربّک یذکر من اقبل الیه و نطق بثنائه و قام علی خدمة امره
العزیز الممنوع لا تحزن عمّا یعملون المقبلون لعمری لو انصفوا و عرفوا
لعملوا بما امروا به فی الواح شتّی من لدی اللّه مالک الوری کذلک نطق امّ
الکتاب من قبل و اللّوح المحفوظ فیهذا المقام المحمود

\end{otherlanguage}

Yā ayyuhā al-nāẓir ilā shaṭr al-sijn wa-l-nāṭiq bi-thanāʾ al-Maqṣūd
wa-l-shārib min raḥīqī al-makhtūm, ʿalayka bahāʾ Allāh al-Muhaymin
al-Qayyūm. Qad ḥaḍara al-ʿabd al-ḥāḍir bi-mā arsaltahu ilayhi wa ʿaraḍa
tilqāʾ al-Maẓlūm mā ghannat ḥamāmat fuʾādika fī thanāʾ Allāh wa
dhikrihi. Ṭūbā li-lisānika wa li-bayānika wa li-man yasmaʿu qawlaka fī
amr Allāh al-ʿAzīz al-Wadūd. Qad anzalnā ʿalā al-muqarrabīn mā law yulqā
ʿalā al-turāb la-taṭluʿa minhu shams al-ʿilm wa-l-bayān. Kadhālika
naṭaqat yarāʿat al-Raḥmān wa-l-nāsu aktharuhum lā yafqahūn. Anta alladhī
ḥaḍarta wa raʾayta al-baḥr wa amwājahu wa-l-shams wa anwārahā wa samiʿta
mā takallama bihi Mukallim al-Ṭūr idh kāna mustawiyan ʿalā ʿarsh
al-ẓuhūr. Inna rabbaka yadhkuru man aqbala ilayhi wa naṭaqa bi-thanāʾihi
wa qāma ʿalā khidmat amrihi al-ʿazīz al-mamnūʿ. Lā taḥzan ʿammā
yaʿmalūna al-muqbilūn. La-ʿamrī law anṣafū wa ʿarafū la-ʿamilū bi-mā
umirū bihi fī alwāḥ shattā min ladā Allāh Mālik al-warā. Kadhālika
naṭaqa Umm al-Kitāb min qablu wa-l-Lawḥ al-Maḥfūẓ fī hādhā al-maqām
al-maḥmūd.

O you who turn your gaze toward the prison, speak in praise of the
Desired One, and drink of My sealed wine! Upon you be the glory of God,
the All-Protecting, the Self-Subsisting. The servant in attendance has
arrived with what you sent him and has presented before the Wronged One
what the dove of your heart has sung in praise and remembrance of God.
Blessed are your tongue and your utterance, and whoever hears your words
concerning the Cause of God, the Almighty, the All-Loving. We have sent
down to those brought near what, if cast upon the dust, would cause the
sun of knowledge and utterance to rise from it. Thus has the Pen of the
All-Merciful spoken, yet most people do not understand. You are the one
who came into Our presence and saw the ocean and its waves, and the sun
and its lights, and heard what He who spoke on Sinai uttered while
seated upon the throne of revelation. Your Lord remembers whoever has
turned to Him, spoken His praise, and arisen to serve His mighty,
unassailable Cause. Do not grieve over what those who have turned toward
Us are doing. By My life, had they been fair and attained recognition,
they would have acted as they were commanded in various Tablets from
God, the Sovereign of all creation. Thus spoke the Mother Book before,
and the Preserved Tablet in this praiseworthy station.

\subsection{53}\label{section-2473}

\begin{otherlanguage}{arabic}

کبّر من قبلی احبّائی و ذکّرهم بما ینفعهم و یحفظهم و یرفعهم الی ملکوتی الّذی
جعله اللّه مطاف من فی الغیب و الشّهود قل یا اولیائی انّا نوصیکم بأعمال یجد
منها الملأ الأعلی عرف التّقدیس لعمری بها یرتفع امری و مقاماتکم فیظلّ قباب
عنایتی و خبآء مجدی و فسطاط رحمتی الّتی سبقت الوجود قل انّا نکون معکم فیکلّ
الأحوال نسمع و نری و انا السّامع البصیر انّا نری و نستر نسمع و نصمت لو
نفصّل لک ما ظهر و سیظهر تأخذک الأحزان فاسئل اللّه بأن یبدّله بسلطان من عنده
انّ ربّک لهو القویّ القدیر البهآء و النّور علیک و علی اولیائی و علی الّذین
وفوا بمیثاق اللّه و عهده فی هذا الیوم العزیز البدیع انتهی

\end{otherlanguage}

Kabbir min qibalī aḥibbāʾī wa dhakkirhum bi-mā yanfaʿuhum wa yaḥfaẓuhum
wa yarfaʿuhum ilā malakūtī alladhī jaʿalahu Allāh maṭāf man fī al-ghayb
wa-l-shuhūd. Qul yā awliyāʾī innā nūṣīkum bi-aʿmāl yajidu minhā al-Malaʾ
al-Aʿlā ʿarf al-taqdīs. La-ʿamrī bihā yartafiʿu amrī wa maqāmātukum fī
ẓill qibāb ʿināyatī wa khibāʾ majdī wa fusṭāṭ raḥmatī allatī sabaqat
al-wujūd. Qul innā nakūnu maʿakum fī kull al-aḥwāl nasmaʿu wa narā wa
anā al-Sāmiʿ al-Baṣīr. Innā narā wa nasturu nasmaʿu wa naṣmutu. Law
nufaṣṣil laka mā ẓahara wa sa-yaẓhar taʾkhudhuka al-aḥzān fa-sʾal Allāh
bi-an yubaddilahu bi-sulṭān min ʿindihi. Inna rabbaka la-huwa al-Qawī
al-Qadīr. Al-bahāʾ wa-l-nūr ʿalayka wa ʿalā awliyāʾī wa ʿalā alladhīna
wafaw bi-mīthāq Allāh wa ʿahdihi fī hādhā al-yawm al-ʿazīz al-badīʿ.
Intahā.

Convey to My loved ones, on My behalf, the greeting of magnification,
and remind them of what will benefit and protect them and raise them to
My Kingdom, which God has made the center around which all in the unseen
and visible realms revolve. Say: O My friends, We counsel you to perform
deeds from which the Concourse on high may perceive the fragrance of
holiness. By My life, through such deeds My Cause is exalted, as are
your stations beneath the domes of My care, the canopy of My glory, and
the pavilion of My mercy, which has preceded all existence. Say: We are
with you in every condition; We hear and see, and I am the All-Hearing,
the All-Seeing. We see and conceal; We hear and remain silent. Were We
to detail for you what has appeared and will appear, sorrows would seize
you. Ask God, then, to transform it through a power from Himself. Your
Lord is indeed the Strong, the All-Powerful. Glory and light be upon
you, upon My friends, and upon those who have remained faithful to God's
covenant and pledge in this mighty, wondrous Day. End.

\subsection{54}\label{section-2474}

\begin{otherlanguage}{persian}

این فانی خدمت آنمحبوب عرض مینماید که اگرچه جواب دستخطّها تأخیر میشود ولکن
بعد جمیع مطالب واحداً بعد واحد از سماء مشیّت الهی نازل و ارسال میگردد از
حسن حظّ آنکه سبب تأخیر خدمت آنحضرت مشهود و واضح و مبرهنست در هر حال امید
عفو است و دیگر از آنحضرت این خادم فانی مطمئن است چه که مطّلع و آگاهند
اینکه بشارت وصول الواح الهیرا مرقوم داشتید و هم‌چنین اجتماع دوستان و
القاء آن بر ایشان و ذکر مراتب شوق و ذوق و جذب و انجذاب حزبه تعالی
فی‌الحقیقه مسرّت بی‌اندازه عنایت فرمود انشآءاللّه در جمیع احوال از حلاوت
بیان رحمن قسمت برند و بما ینبغی قیام نمایند چه‌قدر محبوبست از آن نفوس
منجذبهٴ مشتعله ظهور اعمال و اخلاق طیّبهٴ مرضیّه آفاق عالم انسان باین انوار
روشن و منیر عمل مقبل مستقیم بمثابهٴ شجر است اصلش در ارض استقامت و فرعش
الی ما شآء اللّه و اثمارش بدوام ملک و ملکوت فی‌الحقیقه جز مرآت علم الهی
احصای اینمراتب ننموده و نخواهد نمود هنیئاً لحضرتکم و لهم از حقّ میطلبیم
جمیع اهل عالم را از این نعمت طیّبه و مائدهٴ سمائیّه قسمت عطا فرماید و
محروم ننماید انّه علی کلّ شیء قدیر

\end{otherlanguage}

Īn fānī khedmat-e ān maḥbūb ʿarż mīnamāyad ke agarche javāb-e
dastkhaṭṭ-hā taʾkhīr mīshavad valīkan baʿd jamīʿ-e maṭāleb vāḥedan baʿd
vāḥed az samā-ye mashiyyat-e elāhī nāzel va ersāl mīgardad. Az ḥosn-e
ḥaẓẓ ānke sabab-e taʾkhīr khedmat-e ān ḥażrat mashhūd va vāżeḥ va
mobarhan ast. Dar har ḥāl omīd-e ʿafv ast va dīgar az ān ḥażrat īn
khādem-e fānī moṭmaʾenn ast che ke moṭṭaleʿ va āgāhand. Īnke beshārat-e
voṣūl-e alvāḥ-e elāhī rā marqūm dāshtīd va hamchenīn ejtemāʿ-e dūstān va
elqā-ye ān bar īshān va dhekr-e marāteb-e showq va dhowq va jadhb va
enjedhāb-e ḥezb-e taʿālā fī-l-ḥaqīqe masarrat-e bī-andāze ʿenāyat
farmūd. In shāʾ Allāh dar jamīʿ-e aḥvāl az ḥalāvat-e bayān-e Raḥmān
qesmat barand va bi-mā yanbaghī qiyām namāyand. Che-qadr maḥbūb ast az
ān nofūs-e monjadhebe-ye moshtaʿele ẓohūr-e aʿmāl va akhlāq-e ṭayyebe-ye
marżiyye. Āfāq-e ʿālam-e ensān be īn anvār rowshan va monīr. ʿAmal-e
moqbel-e mostaqīm be-mathābe-ye shajar ast aṣlash dar arż-e esteqāmat va
farʿash ilā mā shāʾ Allāh va athmārash be-davām-e molk va malakūt.
Fī-l-ḥaqīqe joz merʾāt-e ʿelm-e elāhī eḥṣā-ye īn marāteb nanamūde va
nakhāhad namūd. Hanīʾan li-ḥaḍratikum wa lahum. Az Ḥaqq mīṭalabīm
jamīʿ-e ahl-e ʿālam rā az īn neʿmat-e ṭayyebe va māʾede-ye samāʾiyye
qesmat ʿaṭā farmāyad va maḥrūm nanamāyad. Innahu ʿalā kull shayʾ qadīr.

This transient one submits to that beloved friend that, although replies
to the letters are delayed, all the matters raised are subsequently
revealed, one by one, from the heaven of the divine will and dispatched.
Fortunately, the reason for the delay is evident, clear, and established
to you. In every circumstance, forgiveness is hoped for; moreover, this
transient servant feels assured concerning you, since you are informed
and aware. What you wrote conveying the glad news of the arrival of the
divine Tablets, the gathering of the friends and the reading of the
Tablets to them, and the degrees of longing, joy, attraction, and
rapture among His exalted company truly bestowed boundless happiness.
God willing, may they partake in every circumstance of the sweetness of
the All-Merciful's utterance and arise to do what is fitting. How
welcome is the appearance of good, pleasing deeds and qualities from
those enraptured, enkindled souls! The horizons of the human world are
bright and radiant with these lights. The deed of a steadfast believer
is like a tree: its root is in the soil of steadfastness, its branches
reach as far as God wills, and its fruits endure as long as the earthly
and heavenly kingdoms. In truth, none but the Mirror of divine knowledge
has reckoned, or will ever reckon, these degrees. May this blessing be
yours and theirs! We ask God to grant all the people of the world a
share of this goodly bounty and heavenly banquet, and not to deprive
them. He is powerful over all things.

\subsection{55}\label{section-2475}

\begin{otherlanguage}{persian}

و اینکه اظهار تشکّر و بهجت و سرور از مائدهٴ ثانیه و شمس جدیدهٴ مشرقه و
نعمت مکرّره فرمودند ارکانرا بهجت بسیار لطیفی بل الطف عمّا یحصی عطا نمود در
ذکر اشراق الواح الهی از افق عنایت مرّهٴ اخری این کلمه را مرقوم داشتید
قولکم دیگر این لاشیء این کرّه چنان جرعه‌نوش شده و از این عطیّهٴ کبری محظوظ
گشته که خود را مالک بر من علی الأرض یافته بعد از عرض این کلمه در ساحت
امنع اقدس لسان الهی باین کلمهٴ مبارکه ناطق قوله جلّ و عزّ این بیان جناب
علی قبل اکبر بمثابهٴ شاخهٴ محبّت است که بعشق الهی پیوند شده انتهی

\end{otherlanguage}

Va īnke eẓhār-e tashakkor va bahjat va sorūr az māʾede-ye thāniye va
shams-e jadīde-ye moshreqe va neʿmat-e mokarrarre farmūdand arkān rā
bahjat-e besyār laṭīfī bal alṭaf ʿammā yuḥṣā ʿaṭā namūd. Dar dhekr-e
eshrāq-e alvāḥ-e elāhī az ofoq-e ʿenāyat marre-ye okhrā īn kalame rā
marqūm dāshtīd qowlokom dīgar īn lā-shayʾ īn karre chenān jorʿe-nūsh
shode va az īn ʿaṭiyye-ye kobrā maḥẓūẓ gashte ke khod rā mālek bar man
ʿalā al-arḍ yāfte. Baʿd az ʿarż-e īn kalame dar sāḥat-e amnaʿ-e aqdas
lesān-e elāhī be īn kalame-ye mobārake nāṭeq qawluhu jalla wa ʿazza īn
bayān-e Jenāb-e ʿAlī qabl-e Akbar be-mathābe-ye shākhe-ye maḥabbat ast
ke be ʿeshq-e elāhī peyvand shode. Intahā.

Your expression of gratitude, delight, and joy for the second banquet,
the newly risen sun, and the renewed bounty bestowed upon every limb a
most exquisite joy---indeed, more exquisite than can be reckoned. In
mentioning the renewed dawning of the divine Tablets from the horizon of
grace, you wrote these words: ``Furthermore, this nothingness has this
time drunk so deeply and so enjoyed this supreme gift that he has found
himself sovereign over all who are on earth.'' After these words were
presented in the most exalted, most holy court, the divine Tongue spoke
these blessed words---exalted and glorified be He: ``This utterance of
Jinab-i `Ali before Akbar is like a branch of love grafted onto divine
love.'' End.

\subsection{56}\label{section-2476}

\begin{otherlanguage}{persian}

صدهزار روح و بحور آن فدای اینکلمهٴ مبارکه که بمثابهٴ آفتاب از افق عنایت
طالع شده سبّح بحمد ربّک و کن من الشّاکرین

\end{otherlanguage}

Ṣad-hezār rūḥ va boḥūr-e ān fadā-ye īn kalame-ye mobārake ke
be-mathābe-ye āftāb az ofoq-e ʿenāyat ṭāleʿ shode. Sabbiḥ bi-ḥamd
rabbika wa kun min al-shākirīn.

May a hundred thousand souls, and oceans of them, be sacrificed for this
blessed utterance, which has risen like the sun from the horizon of
grace. Glorify your Lord with praise and be among the thankful.

\subsection{57}\label{section-2477}

\begin{otherlanguage}{persian}

و اینکه در ذکر دوستان الهی مرقوم داشتید که ایشان هم از کوثر بیان رحمن
منجذب شدند بشأنیکه مجذوب‌وار بثنای الهی ناطق گشتند و هم‌چنین بنفی شعور و
ادراک و انسانیّت از نفوس معرضه از اهل بیان بعد از عرض این فقره این کلمهٴ
تامّهٴ مطاعه از مطلع بیان احدیّه اشراق فرمود فرمودند یا عبد حاضر انشآءاللّه
دوستان الهی همیشه جمع باشند و بذکر و ثنای حقّ مشغول هر نفسی از بحر اعظم
آشامید عالم را در ظلّ خود مشاهده نماید این طیور را پرواز حقّ آموخته و فوق
طیور ادراک و عقول طیران نمایند طوبی لهم طوبی لهم انتهی

\end{otherlanguage}

Va īnke dar dhekr-e dūstān-e elāhī marqūm dāshtīd ke īshān ham az
Kowthar-e bayān-e Raḥmān monjadheb shodand be-shaʾnī ke majdhūb-vār be
thanā-ye elāhī nāṭeq gashtand va hamchenīn be nafy-e shoʿūr va edrāk va
ensāniyyat az nofūs-e moʿreże az ahl-e Bayān baʿd az ʿarż-e īn feqre īn
kalame-ye tāmme-ye moṭāʿe az maṭlaʿ-e bayān-e aḥadiyye eshrāq farmūd.
Farmūdand yā ʿAbd-e Ḥāżer in shāʾ Allāh dūstān-e elāhī hamīshe jamʿ
bāshand va be dhekr va thanā-ye Ḥaqq mashghūl. Har nafsī az baḥr-e aʿẓam
āshāmīd ʿālam rā dar ẓell-e khod moshāhede namāyad. Īn ṭoyūr rā parvāz
Ḥaqq āmūkhte va fowq-e ṭoyūr-e edrāk va ʿoqūl ṭayarān namāyand. Ṭūbā
lahum ṭūbā lahum. Intahā.

You wrote, in speaking of the friends of God, that they too were so
enraptured by the Kawthar of the All-Merciful's utterance that, like
souls transported, they began to voice the praise of God; you also wrote
denying awareness, understanding, and humanity to those among the people
of the Bayan who have turned away. After this passage was presented,
this perfect, authoritative utterance shone forth from the Dayspring of
the utterance of divine unity. He said: ``O servant in attendance! God
willing, may the friends of God always be gathered together and occupied
with the remembrance and praise of God. Any soul who has drunk from the
Most Great Ocean will behold the world beneath his shadow. God has
taught these birds to fly, and they soar above the birds of
understanding and intellect. Blessed are they! Blessed are they!'' End.

\subsection{58}\label{section-2478}

\begin{otherlanguage}{arabic}

فدای وجودیکه حلاوت بیانرا بیابد و شموس شفقت و الطاف را از آفاق سماء
کلمات کتاب الهی مشاهده نماید

\end{otherlanguage}

Fadā-ye vojūdī ke ḥalāvat-e bayān rā beyābad va shomūs-e shafaqat va
alṭāf rā az āfāq-e samā-ye kalemāt-e Ketāb-e elāhī moshāhede namāyad.

May I be sacrificed for the one who tastes the sweetness of the
utterance and beholds the suns of compassion and loving-kindness shining
from the horizons of the heaven of the words of the divine Book.

\subsection{59}\label{section-2479}

\begin{otherlanguage}{arabic}

و اینکه در بارهٴ حرکت و مشورت مرقوم داشتند فی‌الحقیقه کون آنمحبوب در آن
ایّام در آن ارض لازم بود امید هست از عنایت حقّ جلّ جلاله که آن ارض مبارکه
را از شئونات مردودهٴ معرضین از اهل بیان مقدّس فرماید

\end{otherlanguage}

Va īnke dar bāre-ye ḥarekat va mashverat marqūm dāshtand fī-l-ḥaqīqe
kown-e ān maḥbūb dar ān ayyām dar ān arż lāzem būd. Omīd hast az
ʿenāyat-e Ḥaqq jalla jalāluhu ke ān arż-e mobārake rā az shoʾūnāt-e
mardūde-ye moʿreżīn az ahl-e Bayān moqaddas farmāyad.

As to what you wrote concerning departure and consultation, that beloved
friend's presence in that land during those days was indeed necessary.
It is hoped that, through His grace, God---glorified be His glory---will
sanctify that blessed land from the reprehensible conduct of those among
the people of the Bayan who have turned away.

\subsection{60}\label{section-2480}

\begin{otherlanguage}{persian}

سبحان‌اللّه چه مقدار مردم غافلند هر ذیشعور و درایتی متحیّر است از آن نفوس
مجعوله جهل را بر کرسیّ قلوب مقرّ داده‌اند و علم و انصاف را خارج باب
داشته‌اند افّ لهم و لمن ضلّ و اضلّهم و لداهیة تمشی و لا تعرف الی من و لها
هاویة تستعیذ الهاویة منها

\end{otherlanguage}

Sobḥān Allāh che meqdār mardom ghāfeland. Har dhī-shoʿūr va derāyatī
motaḥayyer ast az ān nofūs-e majʿūle. Jahl rā bar korsī-ye qolūb maqarr
dāde-and va ʿelm va enṣāf rā khārej-e bāb dāshte-and. Uffin lahum wa
li-man ḍalla wa aḍallahum wa li-dāhiyatin tamshī wa lā taʿrifu ilā man
wa lahā hāwiyatun tastaʿīdhu al-hāwiyatu minhā.

Glory be to God! How heedless people are! Everyone endowed with
awareness and discernment is astonished at those counterfeit souls. They
have enthroned ignorance in their hearts and kept knowledge and fairness
outside the door. Fie upon them, upon the one who has gone astray and
led them astray, and upon a calamity that walks without knowing toward
whom, and that has an abyss from which the abyss itself seeks refuge!

\subsection{61}\label{section-2481}

\begin{otherlanguage}{persian}

و اینکه در بارهٴ جناب آقا میرزا لطف‌اللّه علیه بهآء اللّه مرقوم داشتید در
پیشگاه حضور مالک قدم عرض شد هذا ما نطق به لسان الرّحمن فی ملکوت البیان
قوله جلّ اجلاله

\end{otherlanguage}

Va īnke dar bāre-ye Jenāb-e Āqā Mīrzā Loṭf-Allāh ʿalayhi bahāʾ Allāh
marqūm dāshtīd dar pīshgāh-e ḥożūr-e Mālek-e Qedam ʿarż shod. Hādhā mā
naṭaqa bihi lisān al-Raḥmān fī malakūt al-bayān qawluhu jalla ijlāluhu.

What you wrote concerning Jinab-i Aqa Mirza Lutfu'llah---upon him be the
glory of God---was presented in the presence of the Lord of Eternity.
This is what the Tongue of the All-Merciful uttered in the kingdom of
utterance---His words, exalted be His majesty:

\subsection{62}\label{section-2482}

\begin{otherlanguage}{arabic}

بسمی النّاصح المشفق الخبیر

\end{otherlanguage}

Bi-smī al-Nāṣiḥ al-Mushfiq al-Khabīr.

In My name, the Counselor, the Compassionate, the All-Aware.

\subsection{63}\label{section-2483}

\begin{otherlanguage}{persian}

یا ایّها المقبل الی الوجه ایّام ایّام اللّه است بیک آن آن قرون و اعصار
معادله نمینماید ذرّهٴ این یوم بمثابهٴ خورشید مشاهده میگردد و قطره‌اش مانند
دریا اگر نفسی یک نفس فی حبّ اللّه و لخدمته برآرد او از سیّد اعمال از قلم
اعلی مسطور اگر فضائل این یوم ذکر شود کلّ منصعق مشاهده شوند الّا من شآء ربّک
لذا محبوب آنکه در این یوم در صدد آن باشی که بانوار شمس رضای حقّ فائز شوی
یک کلمه از صدهزار اولاد بهتر و برتر و پاینده‌تر است انشآءاللّه باینمقام
بلند اعلی فائز گردی تا ذکرت بدوام ملک و ملکوت پاینده ماند در هیچ حال
مأیوس مباش بذیل حقّ متشبّث باش و باو متوکّل قل یا الهی انا الّذی اقبلت الیک
بروحی و ذاتی و نفسی و ما عندی و وضعت املی و رجائی و مشیّتی و ارادتی امام
کرسیّ عظمتک ایّدنی علی ما اردته بجودک و فضلک انت تعلم بأنّی لا علم لی بما
ینفعنی و یضرّنی قدّر لی ما هو خیر لی انّک انت الشّاهد العالم المشفق العلیم
الحکیم انتهی

\end{otherlanguage}

Yā ayyuhā al-muqbil ilā al-Wajh ayyām ayyām-e Allāh ast. Be yek ān-e ān
qorūn va aʿṣār moʿādele namīnamāyad. Dharre-ye īn yowm be-mathābe-ye
khorshīd moshāhede mīgardad va qaṭre-ash mānand-e daryā. Agar nafsī yek
nafas fī ḥubb Allāh wa li-khidmatihi barārad ū az sayyed-e aʿmāl az
Qalam-e Aʿlā masṭūr. Agar fażāʾel-e īn yowm dhekr shavad koll monṣaʿeq
moshāhede shavand illā man shāʾa rabbuka. Ledhā maḥbūb ānke dar īn yowm
dar ṣadad-e ān bāshī ke be anvār-e shams-e reżā-ye Ḥaqq fāʾez shavī. Yek
kalame az ṣad-hezār owlād behtar va bartar va pāyande-tar ast. In shāʾ
Allāh be īn maqām-e boland-e aʿlā fāʾez gardī tā dhekrat be-davām-e molk
va malakūt pāyande mānad. Dar hīch ḥāl maʾyūs mabāsh be dheyl-e Ḥaqq
motashabbeth bāsh va be ū motavakkel. Qul yā ilāhī anā alladhī aqbaltu
ilayka bi-rūḥī wa dhātī wa nafsī wa mā ʿindī wa waḍaʿtu amalī wa rajāʾī
wa mashiyyatī wa irādatī amāma kursī ʿaẓamatika. Ayyidnī ʿalā mā
aradtahu bi-jūdika wa faḍlika. Anta taʿlamu bi-annī lā ʿilma lī bi-mā
yanfaʿunī wa yaḍurrunī. Qaddir lī mā huwa khayrun lī. Innaka anta
al-Shāhid al-ʿĀlim al-Mushfiq al-ʿAlīm al-Ḥakīm. Intahā.

O you who turn toward the Face! These days are the days of God.
Centuries and ages cannot equal a single moment of them. An atom of this
Day is seen as a sun, and its drop as an ocean. If a soul draws a single
breath in the love of God and in His service, it is recorded by the Most
Exalted Pen among the foremost of deeds. Were the virtues of this Day to
be recounted, all would be seen struck senseless, except those whom your
Lord wills. It is therefore desirable that on this Day you strive to
attain the rays of the sun of God's good pleasure. One word is better,
loftier, and more enduring than a hundred thousand children. God
willing, may you attain this lofty, most exalted station, so that your
remembrance may endure as long as the earthly and heavenly kingdoms.
Never despair in any circumstance. Cling to the hem of God's garment and
place your trust in Him. Say: ``O my God, I am the one who has turned to
You with my spirit, my being, my self, and all that I possess, and has
laid my aspiration, my hope, my will, and my desire before the throne of
Your grandeur. Through Your generosity and grace, assist me to
accomplish what You have willed. You know that I have no knowledge of
what benefits me or harms me. Ordain for me what is best for me. You are
indeed the Witness, the Knowing, the Compassionate, the All-Knowing, the
All-Wise.'' End.

\subsection{64}\label{section-2484}

\begin{otherlanguage}{arabic}

و اینکه در بارهٴ جناب حاجی آقا علیه بهآء اللّه مرقوم داشتید بعد از عرض در
ساحت اقدس این لوح مبارک از سماء مشیّت نازل انشآءاللّه از بحر بیان رحمن
بیاشامند و بانوار معرفتش منوّر گردند قوله جلّ کبریائه

\end{otherlanguage}

Va īnke dar bāre-ye Jenāb-e Ḥājī Āqā ʿalayhi bahāʾ Allāh marqūm dāshtīd
baʿd az ʿarż dar sāḥat-e aqdas īn lowḥ-e mobārak az samā-ye mashiyyat
nāzel. In shāʾ Allāh az baḥr-e bayān-e Raḥmān biyāshāmand va be anvār-e
maʿrefatash monavvar gardand. Qawluhu jalla kibriyāʾuhu.

After what you wrote concerning Jinab-i Haji Aqa---upon him be the glory
of God---was presented in the Most Holy Court, this blessed Tablet
descended from the heaven of the divine will. God willing, may he drink
from the ocean of the All-Merciful's utterance and be illumined by the
lights of His knowledge. His words---exalted be His grandeur:

\subsection{65}\label{section-2485}

\begin{otherlanguage}{arabic}

هو الأقدس الأعظم الأبهی

\end{otherlanguage}

Huwa al-Aqdas al-Aʿẓam al-Abhā.

He is the Most Holy, the Most Great, the All-Glorious.

\subsection{66}\label{section-2486}

\begin{otherlanguage}{persian}

انّ النّور ینادی و النّار تنادی طوبی لأذن سمعت و فازت بما اراده اللّه فی
یومه العزیز البدیع انّه اتی لحیوة العالم و الأمم ارادوا قتله فویل لهم انّ
ربّک لهو الشّاهد الخبیر قل تاللّه انّ الطّور یطوف حولی و الأنوار سجدت لهذا
النّور اللّمیع من النّاس من نطق بالهوی و منهم من قام علی الافترآء و منهم من
تمسّک بالّذین نقضوا میثاق اللّه و عهده کذلک قضی الأمر و انا من الشّاهدین یا
ملأ الأرض دعوا ما عندکم رجآء ما عند اللّه کذلک یأمرکم من ینطق بالحقّ انّه
لا اله الّا هو المقتدر القدیر طوبی لک بما ذکرت لدی المظلوم و نزل لک من
عنده ما یغنیک فیکلّ عالم من عوالم ربّک ان انت من الرّاسخین تمسّک بکتاب اللّه
و ما نزّل من عنده انّ ربّک لهو المؤیّد الغفور الرّحیم قد شهدت الکتب لکتابی و
هو هذا الهیکل الّذی استوی علی العرش رغماً لکلّ عالم بعید کذلک نطق قلمی
الأعلی فضلاً من عندی ان افرح بهذا الفضل العظیم البهآء المشرق من افق
ملکوتی علی الّذین نبذوا سوائی و تمسّکوا بعنایتی و تشبّثوا بذیلی المنیر
انتهی

\end{otherlanguage}

Inna al-nūra yunādī wa-al-nāra tunādī. Ṭūbā li-udhunin samiʿat wa-fāzat
bi-mā arādahu Allāhu fī yawmihi al-ʿazīzi al-badīʿ. Innahu atā li-ḥayāti
al-ʿālami wa-al-umamu arādū qatlahu, fa-waylun lahum. Inna rabbaka
la-huwa al-shāhidu al-khabīr. Qul ta-Allāhi inna al-Ṭūra yaṭūfu ḥawlī
wa-al-anwāru sajadat li-hādhā al-nūri al-lamīʿ. Mina al-nāsi man naṭaqa
bi-al-hawā wa-minhum man qāma ʿalā al-iftirāʾ wa-minhum man tamassaka
bi-alladhīna naqaḍū mīthāqa Allāhi wa-ʿahdahu. Kadhālika quḍiya al-amru
wa-anā mina al-shāhidīn. Yā malaʾa al-arḍi daʿū mā ʿindakum rajāʾa mā
ʿinda Allāh. Kadhālika yaʾmurukum man yanṭiqu bi-al-ḥaqq. Innahu lā
ilāha illā huwa al-muqtadiru al-qadīr. Ṭūbā laka bi-mā dhukirta ladā
al-Maẓlūmi wa-nazala laka min ʿindihi mā yughnīka fī kulli ʿālamin min
ʿawālimi rabbika in anta mina al-rāsikhīn. Tamassak bi-kitābi Allāhi
wa-mā nuzzila min ʿindihi. Inna rabbaka la-huwa al-muʾayyidu al-ghafūru
al-raḥīm. Qad shahidati al-kutubu li-kitābī wa-huwa hādhā al-haykalu
alladhī istawā ʿalā al-ʿarshi raghman li-kulli ʿālimin baʿīd. Kadhālika
naṭaqa qalamī al-aʿlā faḍlan min ʿindī. An ifraḥ bi-hādhā al-faḍli
al-ʿaẓīm. Al-bahāʾu al-mushriqu min ufuqi malakūtī ʿalā alladhīna
nabadhū siwāʾī wa-tamassakū bi-ʿināyatī wa-tashabbathū bi-dhaylī
al-munīr. Intahā.

The light calls out, and the fire calls out. Blessed is the ear that has
heard and attained what God has intended in His mighty, wondrous Day. He
has come to bring life to the world, yet the peoples have sought to kill
Him. Woe to them! Your Lord is indeed the Witness, the All-Aware. Say:
``By God, the Mount circles around Me, and the lights have bowed down
before this resplendent Light.'' Among the people are those who speak
out of desire, those who arise to fabricate falsehoods, and those who
cling to the ones who have broken God's covenant and His pledge. Thus
has the matter been decreed, and I am among the witnesses. People of the
earth, relinquish what you possess in the hope of what is with God. Thus
commands you the One who speaks the truth. There is no god but Him, the
All-Powerful, the Almighty. Blessed are you, for you have been mentioned
in the presence of the Wronged One, and He has sent down for you what
will suffice you in every world among the worlds of your Lord, if you
remain steadfast. Hold fast to the Book of God and to what has been sent
down from Him. Your Lord is indeed the One who strengthens, the
Ever-Forgiving, the Merciful. The Scriptures have borne witness to My
Book, which is this Temple that has ascended the throne despite every
wayward scholar. Thus has My Most Exalted Pen spoken, as a favor from
Me. Rejoice in this great favor. May the glory shining from the horizon
of My kingdom rest upon those who have cast aside all but Me, held fast
to My loving care, and clung to My luminous robe. End.

\subsection{67}\label{section-2487}

\begin{otherlanguage}{persian}

انوار فضل احاطه فرمود و فجر عنایت دمید صدهزار طوبی از برای نفوسیکه الیوم
بمقصود فائز شدند و شئونات و ظهورات ایّام معدوده ایشانرا از مقام باقی منع
ننمود تکبیر و سلام این عبد خدمت جناب میرزا و ایشان و هر یک از دوستان
الهی اگر ممکن شود در هر صباح و مساء بلکه در هر ساعت بل در هر آن معلّق
بعنایت آنمحبوبست

\end{otherlanguage}

Anvār-e faḍl eḥāṭeh farmūd va fajr-e ʿināyat damīd. Ṣad-hizār ṭūbā az
barā-ye nofūsī keh al-yawm beh maqṣūd fāʾiz shodand va shuʾūnāt va
ẓuhūrāt-e ayyām-e maʿdūdeh īshān-rā az maqām-e bāqī manʿ nanamūd. Takbīr
va salām-e īn ʿabd khedmat-e Jināb-e Mīrzā va īshān va har yek az
dūstān-e ilāhī, agar momken shavad dar har ṣabāḥ va masāʾ balkeh dar har
sāʿat bal dar har ān, moʿallaq beh ʿināyat-e ān maḥbūb ast.

The lights of grace have encompassed all, and the dawn of loving care
has broken. A hundred thousand blessings upon those souls who have
attained the Desired One today, and whom the circumstances and
appearances of these numbered days have not kept from the everlasting
station. The conveyance of this servant's greetings and salutations to
Jinab-i Mirza, to them, and to each of the friends of God---if possible,
every morning and evening, indeed every hour, even every
moment---depends upon the kindness of that beloved friend.

\subsection{68}\label{section-2488}

\begin{otherlanguage}{persian}

و اینکه ذکر جناب رضا قلیخان علیه ٦٦٩ {[}بهآء اللّه{]} فرموده بودند امام
وجه من لا یعزب عن علمه من شیء عرض شد هذا ما نطق به لسان الشّفقة و الکرم و
الفضل قوله جلّ شأنه و عظم کرمه یا علی قبل اکبر علیک رحمتی الّتی سبقت البشر
حقّ جلّ جلاله جمیع عالمرا از برای معرفت خود خلق فرموده و از برای کلّ راحت
خواسته و اسبابیهم بمقتضیات حکمتش ظاهر فرموده لعمری اگر اهل ارض بانوار
نیّر سماء علم منوّر شوند یعنی بآنچه از افق اعلی و ذروهٴ علیا نازل تمسّک
نمایند و عمل کنند کلّ خود را در کمال سرور و امنیّت و راحت مشاهده نمایند
آنچه در بارهٴ ایشان عمل نمودید لدی العرش محبوب از قبل مظلوم تکبیر برسان
بگو محزون مباش از آنچه بر تو وارد شد بر حقّ اعظم از آن وارد قسم بانوار
وجه الهی این عالم باینعظمت که مشاهده میشود لایق یکنفر از دوستان الهی
نبوده و نیست ان استقم علی الأمر انّه یفعل ما یشآء و یحکم ما یرید
انشآءاللّه در کلّ احیان باو ناظر باشند و بذیلش متشبّث انّه یقدّر له ما اراد و
هو المشفق الکریم البهآء علیه و الصّبر له و الاصطبار لک و الأمر لی و انا
الآمر الحکیم انتهی

\end{otherlanguage}

Va īn-keh dhekr-e Jināb-e Riḍā Qulī Khān, ʿalayhi 669 {[}Bahāʾ Allāh{]},
farmūdeh būdand, imāma wajhi man lā yaʿzubu ʿan ʿilmihi min shayʾ, ʿarḍ
shod. Hādhā mā naṭaqa bihi lisānu al-shafaqati wa-al-karami wa-al-faḍl.
Qawluhu jalla shaʾnuhu wa-ʿaẓuma karamuhu: Yā ʿAlī qabla Akbar, ʿalayka
raḥmatī allatī sabaqati al-bashar. Ḥaqq jalla jalāluhu jamīʿ-e ʿālam-rā
az barā-ye maʿrefat-e khod khalq farmūdeh va az barā-ye koll rāḥat
khvāsteh va asbābī ham beh moqtaḍayāt-e ḥekmatash ẓāher farmūdeh.
La-ʿamrī agar ahl-e arḍ beh anvār-e nayyer-e samā-ye ʿelm monavvar
shavand, yaʿnī beh āncheh az ofoq-e aʿlā va dhorveh-ye ʿulyā nāzel
tamassok namāyand va ʿamal konand, koll khod-rā dar kamāl-e sorūr va
amniyyat va rāḥat moshāhedeh namāyand. Āncheh dar bāreh-ye īshān ʿamal
namūdīd ladā al-ʿarsh maḥbūb. Az qebal-e Maẓlūm takbīr berasān. Begū
maḥzūn mabāsh az āncheh bar to vāred shod. Bar Ḥaqq aʿẓam az ān vāred.
Qasam beh anvār-e vajh-e ilāhī, īn ʿālam beh īn ʿaẓamat keh moshāhedeh
mīshavad lāyeq-e yek nafar az dūstān-e ilāhī nabūdeh va nīst. An istaqim
ʿalā al-amr. Innahu yafʿalu mā yashāʾu wa-yaḥkumu mā yurīd. In shāʾ
Allāh dar koll-e aḥyān beh ū nāẓer bāshand va beh dhaylash motashabbeth.
Innahu yuqaddiru lahu mā arāda wa-huwa al-mushfiqu al-karīm. Al-bahāʾu
ʿalayhi wa-al-ṣabru lahu wa-al-iṣṭibāru laka wa-al-amru lī wa-anā
al-āmiru al-ḥakīm. Intahā.

What you had written concerning Jinab-i Rida Quli Khan---upon him be 669
{[}the glory of God{]}---was presented before the face of Him from whose
knowledge nothing is hidden. This is what the Tongue of compassion,
generosity, and grace spoke. His words---exalted be His station and
great His generosity: ``O Ali before Akbar, upon you be My mercy, which
has preceded humankind. God, exalted be His glory, has created the
entire world for the knowledge of Himself. He has desired comfort for
all and has provided the means in accordance with the requirements of
His wisdom. By My life, if the people of the earth were illumined by the
lights of the sun of the heaven of knowledge---that is, if they held
fast to what has been sent down from the most exalted horizon and the
loftiest summit, and acted upon it---they would all find themselves in
perfect joy, security, and comfort. What you have done for him is
pleasing before the Throne. Convey greetings from the Wronged One. Say:
`Do not grieve over what has befallen you. Something greater has
befallen God. I swear by the lights of the divine countenance, this
world, with all the greatness it appears to possess, has never been and
is not now worthy of a single one of the friends of God. Remain
steadfast in the Cause. He does what He wills and decrees what He
desires.' God willing, may he look to Him at all times and cling to His
robe. He ordains for him whatever He desires, and He is the
Compassionate, the Generous. Glory be upon him, patience be his,
perseverance be yours, and command be Mine. I am the Ordainer, the
All-Wise.'' End.

\subsection{69}\label{section-2489}

\begin{otherlanguage}{persian}

اینکه در بارهٴ جناب آقا سیّد محمّد طبیب علیه بهآء اللّه مرقوم داشتید و
هم‌چنین مراتب ذهن و ذکا و اقبال و کمال و استعداد ایشان اینمراتب بتمامه در
حین زوال یوم یکشنبه شهر صفر مقرّ عرش تلقاء وجه عرض شد هذا ما نطق به لسان
العظمة قوله جلّ کبریائه انشآءاللّه بعنایت حقّ فائز باشند و بانوار معرفتش
منیر و از رحیقش مشروب و از مائده‌اش مرزوق حبّذا هذا الیوم المبارک العزیز
عرفش عالم را احاطه نموده و نسیمش بر کلّ مرور کرده طوبی لمن انتبه من هذا
النّسیم و عرف و قام و قال مقبلاً الی الأفق الأعلی قد کنت نائماً یا الهی
هزّتنی نسمات عنایتک و اقامتنی امام افق فضلک اسئلک بأن تؤیّدنی علی خدمتک و
الاستقامة علی امرک انّک انت المقتدر علی ما تشآء و فی قبضتک زمام الوجود من
الغیب و الشّهود یا علی لحاظ عنایت باو متوجّه بشّره بما نطق به المظلوم فی
سجنه العظیم ترقّی از برای وجود بوده و هست در جمیع عوالم هذا حقّ لا ریب فیه
البهآء علیک و علیه و علی من تمسّک بالمعروف و ما یرتفع به امر اللّه ربّ
العالمین انتهی

\end{otherlanguage}

Īn-keh dar bāreh-ye Jināb-e Āqā Siyyid Muḥammad Ṭabīb, ʿalayhi Bahāʾ
Allāh, marqūm dāshtīd va hamchenīn marāteb-e dhehn va dhakā va eqbāl va
kamāl va esteʿdād-e īshān, īn marāteb beh tamāmeh dar ḥīn-e zavāl-e
yawm-e yekshanbeh shahr-e Ṣafar, maqarr-e ʿarsh, telqāʾ-e vajh ʿarḍ
shod. Hādhā mā naṭaqa bihi lisānu al-ʿaẓama. Qawluhu jalla kibriyāʾuhu:
In shāʾ Allāh beh ʿināyat-e Ḥaqq fāʾiz bāshand va beh anvār-e
maʿrefatash monīr va az raḥīqash mashrūb va az māʾedeh-ash marzūq.
Ḥabbadhā hādhā al-yawmu al-mubāraku al-ʿazīz. ʿArfash ʿālam-rā eḥāṭeh
namūdeh va nasīmash bar koll morūr kardeh. Ṭūbā li-man intabaha min
hādhā al-nasīmi wa-ʿarafa wa-qāma wa-qāla muqbilan ilā al-ufuqi al-aʿlā:
Qad kuntu nāʾiman yā ilāhī. Hazzatnī nasamātu ʿināyatika wa-aqāmatnī
imāma ufuqi faḍlika. Asʾaluka bi-an tuʾayyidanī ʿalā khidmatika
wa-al-istiqāmati ʿalā amrika. Innaka anta al-muqtadiru ʿalā mā tashāʾu
wa-fī qabḍatika zimāmu al-wujūdi mina al-ghaybi wa-al-shuhūd. Yā ʿAlī,
leḥāẓ-e ʿināyat beh ū motavajjeh. Bashshirhu bi-mā naṭaqa bihi
al-Maẓlūmu fī sijnihi al-ʿaẓīm. Taraqqī az barā-ye vojūd būdeh va hast
dar jamīʿ-e ʿavālem. Hādhā ḥaqqun lā rayba fīhi. Al-bahāʾu ʿalayka
wa-ʿalayhi wa-ʿalā man tamassaka bi-al-maʿrūfi wa-mā yartafiʿu bihi amru
Allāhi rabbi al-ʿālamīn. Intahā.

What you wrote concerning Jinab-i Aqa Siyyid Muhammad the
physician---upon him be the glory of God---and likewise the particulars
of his understanding, intelligence, devotion, accomplishments, and
capacity were all presented in His presence at the seat of the Throne,
at noon on Sunday in the month of Safar. This is what the Tongue of
grandeur spoke. His words---exalted be His majesty: ``God willing, may
he receive God's loving care, be illumined by the lights of His
knowledge, drink His choice wine, and be nourished from His table. How
blessed is this mighty, hallowed Day! Its fragrance has encompassed the
world, and its breeze has passed over all. Blessed is the one who has
been awakened by this breeze, recognized, arisen, and said, turning
toward the most exalted horizon: `I was asleep, O my God. The breezes of
Your loving care stirred me and raised me up before the horizon of Your
grace. I ask You to strengthen me to serve You and to remain steadfast
in Your Cause. You are powerful to do whatever You will, and within Your
grasp are the reins of existence, both unseen and seen.' O Ali, the gaze
of loving care is directed toward him. Give him the glad tidings of what
the Wronged One has spoken in His mighty prison. Progress has ever been
and remains possible for existence in all the worlds. This is a truth
beyond doubt. Glory be upon you, upon him, and upon whoever holds fast
to what is right and to that through which the Cause of God, the Lord of
the worlds, is exalted.'' End.

\subsection{70}\label{section-2490}

\begin{otherlanguage}{persian}

اینکه جواب سؤال ایشان مختصر از ملکوت اعلی نازل آنمحبوب میدانند که در
اینمراتب و مقامات نازل شده آنچه که من علی الأرض را کفایت نماید و این
ایّام قلم اعلی یحبّ ان ینطق فیکلّ شأن انّه لا اله الّا هو المهیمن القیّوم از
حقّ جلّ جلاله این خادم فانی سائل و آمل که ایشانرا در جمیع احوال مؤیّد
فرماید بر آنچه که سبب اعلاء کلمة اللّه است

\end{otherlanguage}

Īn-keh javāb-e suʾāl-e īshān mokhtaṣar az malakūt-e aʿlā nāzel, ān
maḥbūb mīdānand keh dar īn marāteb va maqāmāt nāzel shodeh āncheh keh
man ʿalā al-arḍ-rā kefāyat namāyad, va īn ayyām qalam-e aʿlā yuḥibbu an
yanṭiqa fī kulli shaʾnin: Innahu lā ilāha illā huwa al-muhayminu
al-qayyūm. Az Ḥaqq jalla jalāluhu īn khādem-e fānī sāʾel va āmel keh
īshān-rā dar jamīʿ-e aḥvāl moʾayyad farmāyad bar āncheh keh sabab-e
eʿlā-ye kalimat Allāh ast.

As to the brevity of the answer to his question sent down from the
kingdom on high, that beloved friend knows that concerning these matters
and stations, enough has been sent down to suffice everyone on earth. In
these days, the Most Exalted Pen wishes to proclaim under all
circumstances: ``There is no god but Him, the All-Protecting, the
Self-Subsisting.'' This transitory servant asks and hopes of God,
exalted be His glory, that He may strengthen him under all conditions to
do whatever will exalt the Word of God.

\subsection{71}\label{section-2491}

\begin{otherlanguage}{arabic}

و اینکه در بارهٴ جنابان آقا خداداد و آقا شاه‌ویردی علیهما بهآء اللّه مرقوم
داشتید فی‌الحقیقه ایشان بعنایت حقّ فائزند بعد از عرض در ساحت امنع اقدس
اعلی لسان کبریا باین کلمات عالیات ناطق قوله جلّ کبریائه

\end{otherlanguage}

Va īn-keh dar bāreh-ye jinābān Āqā Khodādād va Āqā Shāh-Vīrdī, ʿalayhimā
Bahāʾ Allāh, marqūm dāshtīd, fī al-ḥaqīqeh īshān beh ʿināyat-e Ḥaqq
fāʾizand. Baʿd az ʿarḍ dar sāḥat-e amnaʿ-e aqdas-e aʿlā, lesān-e kebriyā
beh īn kalemāt-e ʿāliyāt nāṭeq. Qawluhu jalla kibriyāʾuhu:

Concerning what you wrote about Aqa Khodadad and Aqa Shah-Virdi---upon
them both be the glory of God---they have indeed received God's loving
care. After it was presented in the most inaccessible, most holy, most
exalted Court, the Tongue of majesty spoke these lofty words. His
words---exalted be His majesty:

\subsection{72}\label{section-2492}

\begin{otherlanguage}{arabic}

هو الشّاهد السّمیع

\end{otherlanguage}

Huwa al-shāhidu al-samīʿ.

He is the Witness, the All-Hearing.

\subsection{73}\label{section-2493}

\begin{otherlanguage}{persian}

یا خداداد قلم اعلی شهادت میدهد بر اینکه هر هنگام ذکرت عرض شد بکلمهٴ الهی
و ذکرش فائز شدی از حقّ سائل شو ترا تأیید فرماید بر عرفان اینمقام کریم
کبیر از اوّل اقبالت بافق اعلی در هر یوم که ذکرت شده لوحی مخصوص نازل و من
غیر آن لحاظ عنایت بتو متوجّه بوده ان اشکر اللّه ربّک و ربّ من فی الأرض و
السّمآء الّذی اخذ یدک و انقذک من ظلمات الأرض و ادخلک فیظلّ سدرة عنایته و
رفع ذکرک بین عباده المقرّبین فضل و رحمتش لازال سبقت داشته و دارد انّه نجّاک
و هداک الی سوآء الصّراط و سقاک خمر الطّهور بید عطائه و جعلک من اهل المقام
کن قائماً علی خدمتی و ناطقاً بثنائی و متمسّکاً بفضلی الّذی احاط الآفاق انتهی

\end{otherlanguage}

Yā Khodādād, qalam-e aʿlā shahādat mīdahad bar īn-keh har hangām dhekrat
ʿarḍ shod beh kalemeh-ye ilāhī va dhekrash fāʾiz shodī. Az Ḥaqq sāʾel
show to-rā taʾyīd farmāyad bar ʿerfān-e īn maqām-e karīm-e kabīr. Az
avval-e eqbālat beh ofoq-e aʿlā dar har yawm keh dhekrat shodeh lowḥī
makhṣūṣ nāzel va men gheyr-e ān leḥāẓ-e ʿināyat beh to motavajjeh būdeh.
An ishkur Allāha rabbaka wa-rabba man fī al-arḍi wa-al-samāʾ, alladhī
akhadha yadaka wa-anqadhaka min ẓulumāti al-arḍi wa-adkhalaka fī ẓilli
sidrati ʿināyatihi wa-rafaʿa dhikraka bayna ʿibādihi al-muqarrabīn. Faḍl
va raḥmatash lā-zāl sabqat dāshteh va dārad. Innahu najjāka wa-hadāka
ilā sawāʾi al-ṣirāṭi wa-saqāka khamra al-ṭuhūri bi-yadi ʿaṭāʾihi
wa-jaʿalaka min ahli al-maqām. Kun qāʾiman ʿalā khidmatī wa-nāṭiqan
bi-thanāʾī wa-mutamassikan bi-faḍlī alladhī aḥāṭa al-āfāq. Intahā.

O Khodadad, the Most Exalted Pen bears witness that whenever mention of
you was presented, you received the divine Word and His remembrance. Ask
God to enable you to recognize this noble and great station. From the
time you first turned toward the most exalted horizon, a special Tablet
has been sent down every day on which you have been mentioned, and
beyond this, the gaze of loving care has been directed toward you. Give
thanks to God, your Lord and the Lord of all who are on earth and in
heaven, who has taken your hand, rescued you from the darkness of the
earth, brought you beneath the shade of the Lote Tree of His loving
care, and exalted your remembrance among His servants who are near to
Him. His grace and mercy have ever taken precedence and continue to do
so. He has delivered you, guided you to the straight path, given you
pure wine to drink from the hand of His bounty, and made you one of the
people of this station. Arise to serve Me, speak My praise, and hold
fast to My grace, which has encompassed the horizons. End.

\subsection{74}\label{section-2494}

\begin{otherlanguage}{arabic}

و هم‌چنین این آیات مشرقات از سماء مشیّت حقّ جلّ جلاله در بارهٴ جناب حاجی
شاه‌ویردی علیه بهآء اللّه نازل قوله جلّت عظمته

\end{otherlanguage}

Va hamchenīn īn āyāt-e moshreqāt az samā-ye mashiyyat-e Ḥaqq jalla
jalāluhu dar bāreh-ye Jināb-e Ḥājī Shāh-Vīrdī, ʿalayhi Bahāʾ Allāh,
nāzel. Qawluhu jallat ʿaẓamatuhu:

Likewise, these radiant verses were sent down from the heaven of the
will of God, exalted be His glory, concerning Jinab-i Haji
Shah-Virdi---upon him be the glory of God. His words---exalted be His
grandeur:

\subsection{75}\label{section-2495}

\begin{otherlanguage}{arabic}

بسمی المشفق الکریم

\end{otherlanguage}

Bi-ismī al-mushfiqi al-karīm.

In My name, the Compassionate, the Generous.

\subsection{76}\label{section-2496}

\begin{otherlanguage}{persian}

یا شاهویردی فائز شدی بآنچه که اهل عالم از آن غافل و محجوبند الّا من شآء
اللّه سلکت و اقبلت و توجّهت الی ان بلغت الی مطلع الهدایة و سمعت ما تکلّم به
لسان اللّه المقتدر العلیم الحکیم کم من ولیّ مات فی حسرة لقائی و کم من حبیب
قصد مقامی و کم من امیر ناح فی فراقی و کم من انیس انفق روحه لوصالی و انّک
قصدت و عزمت متوکّلاً علی اللّه الی ان حضرت فی ساحة طافه الرّوح الأمین و سمعت
من المظلوم ما سمعت اذن المقرّبین ان اشکر اللّه بما ایّدک و قرّبک الی بحر
الحیوان و انزل لک ما لا یعادله شیء من الأشیآء و اسمعک ما سمع الکلیم در
عنایت حقّ تفکّر نمائید تا شما را از جمیع احزان محدودهٴ دوروزه فارغ و آزاد
نماید لایق این فضل توجّه است و لایق توجّه فراغت از احزان البهآء علیک و علی
من معک انتهی

\end{otherlanguage}

Yā Shāh-Vīrdī, fāʾiz shodī beh āncheh keh ahl-e ʿālam az ān ghāfel va
maḥjūband, illā man shāʾa Allāh. Salakta wa-aqbalta wa-tawajjahta ilā an
balaghta ilā maṭlaʿi al-hidāyati wa-samiʿta mā takallama bihi lisānu
Allāhi al-muqtadiri al-ʿalīmi al-ḥakīm. Kam min waliyyin māta fī ḥasrati
liqāʾī wa-kam min ḥabībin qaṣada maqāmī wa-kam min amīrin nāḥa fī firāqī
wa-kam min anīsin anfaqa rūḥahu li-wiṣālī. Wa-innaka qaṣadta wa-ʿazamta
mutawakkilan ʿalā Allāhi ilā an ḥaḍarta fī sāḥatin ṭāfahu al-rūḥu
al-amīn, wa-samiʿta mina al-Maẓlūmi mā samiʿat udhunu al-muqarrabīn. An
ishkur Allāha bi-mā ayyadaka wa-qarrabaka ilā baḥri al-ḥayawāni
wa-anzala laka mā lā yuʿādiluhu shayʾun mina al-ashyāʾi wa-asmaʿaka mā
samiʿa al-Kalīm. Dar ʿināyat-e Ḥaqq tafakkor namāyīd tā shomā-rā az
jamīʿ-e aḥzān-e maḥdūdeh-ye do-rūzeh fāregh va āzād namāyad. Lāyeq-e īn
faḍl tavajjoh ast va lāyeq-e tavajjoh farāghat az aḥzān. Al-bahāʾu
ʿalayka wa-ʿalā man maʿak. Intahā.

O Shah-Virdi, you have attained that of which the people of the world
are unaware and from which they are veiled, except those whom God wills.
You journeyed, turned toward Him, and directed yourself to Him until you
reached the Dawning-place of guidance and heard what the Tongue of God,
the All-Powerful, the All-Knowing, the All-Wise, spoke. How many a
devoted soul died longing to meet Me; how many a loving soul sought My
station; how many a prince lamented separation from Me; and how many a
companion gave his life to attain union with Me! You set out and
resolved, placing your trust in God, until you entered a court around
which the faithful Spirit circled, and heard from the Wronged One what
the ears of those brought near have heard. Give thanks to God for
strengthening you, bringing you near to the ocean of life, sending down
for you that which nothing can equal, and causing you to hear what Moses
heard. Reflect upon God's loving care, that it may free and release you
from all the circumscribed sorrows of these fleeting days. This grace
calls for turning toward Him, and turning toward Him calls for freedom
from sorrows. Glory be upon you and upon those who are with you. End.

\subsection{77}\label{section-2497}

\begin{otherlanguage}{persian}

این خادم فانی لازال از برای جنابان که از قلم اعلی بابناء خلیل و ورّاث
کلیم نامیده شده‌اند خیر و برکت و عنایت و استقامت و فضل و ذکر جاودانی طلب
نموده و مینماید و البتّه وقتی مقرون باجابت گردد چه که بشهادت قلم اعلی و
شهادت آنمحبوب اکثری اوقات را در خدمت امر مصروف داشته و دارند چنین نفوسرا
همهٴ عالم ذاکر شده و خواهند شد این اشیاء متفرّقه و الوان مختلفه عنقریب
معدوم و ما عند اللّه لیبقی طوبی لهم فی‌الحقیقه مؤیّدند از جمله تأیید آنکه
آنمحبوب از ایشان اظهار رضایت نموده‌اند انشآءاللّه لازال موفّق باشند

\end{otherlanguage}

Īn khādem-e fānī lā-zāl az barā-ye jinābān keh az qalam-e aʿlā beh
abnā-ye Khalīl va vorrāth-e Kalīm nāmīdeh shodeh-and kheyr va barakat va
ʿināyat va esteqāmat va faḍl va dhekr-e jāvedānī ṭalab namūdeh va
mīnamāyad va albatteh vaqtī maqrūn beh ejābat gardad, cheh keh beh
shahādat-e qalam-e aʿlā va shahādat-e ān maḥbūb aktharī-ye owqāt-rā dar
khedmat-e Amr maṣrūf dāshteh va dārand. Chonīn nofūs-rā hameh-ye ʿālam
dhāker shodeh va khvāhand shod. Īn ashyā-ye motefarreqeh va alvān-e
mokhtalefeh ʿan-qarīb maʿdūm, wa-mā ʿinda Allāhi la-yabqā. Ṭūbā lahum.
Fī al-ḥaqīqeh moʾayyadand. Az jomleh-ye taʾyīd ān-keh ān maḥbūb az īshān
eẓhār-e riḍāyat namūdeh-and. In shāʾ Allāh lā-zāl movaffaq bāshand.

This transitory servant has always sought, and continues to seek,
goodness, blessing, loving care, steadfastness, grace, and everlasting
remembrance for these two gentlemen, whom the Most Exalted Pen has named
the sons of Abraham and the heirs of Moses. Surely, at some time this
prayer will be answered, for, as the Most Exalted Pen and that beloved
friend testify, they have devoted and continue to devote most of their
time to serving the Cause. The whole world has remembered and will
remember such souls. These scattered things and varied colors will soon
pass away, but what is with God will surely endure. Blessed are they!
They are indeed divinely assisted. Among the signs of this assistance is
that that beloved friend has expressed satisfaction with them. God
willing, may they always be aided to succeed.

\subsection{78}\label{section-2498}

\begin{otherlanguage}{persian}

و اینکه جناب آقا خداداد علیه ٦٦٩ {[}بهآء اللّه{]} فرموده‌اند قریب یک سنه
میشود آنچه ارسال نموده‌اند جواب از این عبد نرفته الأمر کما قال ولکن این
عبد نظر باطمینان از ایشان و هم‌چنین نظر بآنکه در اکثر از الواح ذکرشان از
قلم اعلی جاری این عبد در جواب مراسلات بعضی که تازه بشریعه اقبال نموده‌اند
مشغول و خود ایشان نظر بخدمت و استقامتیکه در امر دارند البتّه راضیند بآنچه
واقعشده حقّ تعالی شأنه شاهد و گواهست که لازال ذکرشان و قیامشان و خدمتشان
نصب عین بوده و هست و امید است از فضل الهی جوابهم ارسال شود

\end{otherlanguage}

Va īn-keh Jināb-e Āqā Khodādād, ʿalayhi 669 {[}Bahāʾ Allāh{]},
farmūdeh-and qarīb-e yek saneh mīshavad āncheh ersāl namūdeh-and javāb
az īn ʿabd narafteh, al-amru kamā qāl. Va-līkin īn ʿabd naẓar beh
eṭmīnān az īshān va hamchenīn naẓar beh ān-keh dar akthar az alvāḥ
dhekrashān az qalam-e aʿlā jārī, īn ʿabd dar javāb-e morāsalāt-e baʿḍī
keh tāzeh beh sharīʿeh eqbāl namūdeh-and mashghūl, va khod-e īshān naẓar
beh khedmat va esteqāmatī keh dar Amr dārand albatteh rāḍī-and beh
āncheh vāqeʿ shodeh. Ḥaqq taʿālā shaʾnuhu shāhed va govāh ast keh lā-zāl
dhekrashān va qiyāmashān va khedmatashān naṣb-e ʿayn būdeh va hast va
omīd ast az faḍl-e ilāhī javāb ham ersāl shavad.

As to Jinab-i Aqa Khodadad's statement---upon him be 669 {[}the glory of
God{]}---that nearly a year has passed without this servant sending a
reply to what he has sent, the matter is as he has said. However,
because of my confidence in him, and because mention of him flows from
the Most Exalted Pen in most of the Tablets, this servant has been
occupied with answering the letters of some who have newly turned toward
the water's edge. In view of his service and steadfastness in the Cause,
he himself is surely content with what has occurred. God, exalted be His
station, bears witness that the remembrance of him, his efforts, and his
service have always been and remain before my eyes. It is hoped that,
through divine grace, a reply will also be sent.

\subsection{79}\label{section-2499}

\begin{otherlanguage}{persian}

و اینکه در بارهٴ جناب آقا محمّد و آقا علی اصغر علیهما بهآء اللّه مرقوم
داشتند الحمد للّه ایشان بذکر الهی از قبل و بعد فائز شده‌اند قد نزّل ذکرهما
من سمآء عنایة ربّنا و ربّکم و محبوبنا و محبوبکم و رقما من القلم الأعلی هذا
من فضل مقصودنا و مقصودکم این کلمات عالیات در اینمقام از لسان مالک وری
جاری و نازل

\end{otherlanguage}

Va īnke dar bāre-ye Janāb-i Āqā Muḥammad va Āqā ʿAlī-Aṣghar, ʿalayhimā
bahāʾuʾllāh, marqūm dāshtand, al-ḥamdu li-llāh, īshān be-dhekr-e ilāhī
az qabl va baʿd fāʾez shode-and. Qad nuzzila dhikruhumā min samāʾi
ʿināyati rabbinā va rabbikum va maḥbūbinā va maḥbūbikum va ruqimā
minaʾl-qalamiʾl-aʿlā. Hādhā min faḍli maqṣūdinā va maqṣūdikum. Īn
kalemāt-e ʿāliyāt dar īn maqām az lesān-e Mālek-e varā jārī va nāzel.

As to what was written concerning Jinab-i Aqa Muhammad and Aqa
ʿAli-Asghar---upon them both be the glory of God---praise be to God,
they have been favored with divine remembrance both before and since.
Their remembrance has indeed been sent down from the heaven of the
loving care of our Lord and your Lord, our Beloved and your Beloved, and
they have been inscribed by the Most Exalted Pen. This is through the
grace of the One we seek and you seek. At this point these exalted words
flowed from the tongue of the Sovereign of all creation and were
revealed:

\subsection{80}\label{section-2500}

\begin{otherlanguage}{arabic}

قوله جلّ و عزّ یا محمّد طوبی لک بما وفیت بعهدک و طوبی لمن عمل فیه ما کان
مقبولاً لدی العرش قد زیّن عملک بطراز القبول فی هذا المقام العزیز المحمود
ان اشکر اللّه بهذا الفضل و قل لک الحمد یا مالک الغیب و الشّهود و نذکر علی
قبل اصغر و نکبّر علیه و نبشّره برحمتی التّی سبقته طوبی له بما فاز بذکری من
قبل و من بعد و شرب رحیقی المختوم انتهی

\end{otherlanguage}

Qawluhu jalla va ʿazza: Yā Muḥammad, ṭūbā laka bimā wafayta bi-ʿahdika
va ṭūbā li-man ʿamila fīhi mā kāna maqbūlan ladāʾl-ʿarsh. Qad zuyyina
ʿamaluka bi-ṭirāziʾl-qabūli fī hādhāʾl-maqāmiʾl-ʿazīziʾl-maḥmūd.
Aniʾshkuriʾllāha bi-hādhāʾl-faḍli va qul: Lakaʾl-ḥamdu yā
Mālikaʾl-ghaybi vaʾl-shuhūd. Va nadhkuru ʿAlī qabla Aṣghar va nukabbiru
ʿalayhi va nubashshiruhu bi-raḥmatiyaʾllatī sabaqat-hu. Ṭūbā lahu bimā
fāza bi-dhikrī min qablu va min baʿdu va shariba raḥīqiyaʾl-makhtūm.
Intahā.

His words, glorified and exalted is He: ``O Muhammad! Blessed are you
for having fulfilled your covenant, and blessed is whoever has done
within it what was acceptable before the Throne. Your deed has been
adorned with the ornament of acceptance in this mighty and praiseworthy
station. Give thanks to God for this favor and say: `Praise be to You, O
Sovereign of the unseen and the seen!' We also remember ʿAli before
Asghar, convey Our greetings to him, and give him the glad tidings of My
mercy, which has preceded him. Blessed is he for having attained My
remembrance before and since, and for having drunk My sealed wine.'' End
of quotation.

\subsection{81}\label{section-2501}

\begin{otherlanguage}{arabic}

عمل آن حضرت بسیار مقبول افتاد و عمل جناب آقا محمّد هم بطراز قبول فائز قد
شهد بذلک آیات اللّه ربّنا و ربّ العالمین

\end{otherlanguage}

ʿAmal-e ān ḥaḍrat besyār maqbūl oftād va ʿamal-e Janāb-i Āqā Muḥammad
ham be-ṭarāz-e qabūl fāʾez. Qad shahidat bi-dhālika āyātuʾllāhi rabbinā
va rabbiʾl-ʿālamīn.

That gentleman's deed was most acceptable, and the deed of Jinab-i Aqa
Muhammad was also favored with the ornament of acceptance. The verses of
God, our Lord and the Lord of the worlds, have indeed borne witness to
this.

\subsection{82}\label{section-2502}

\begin{otherlanguage}{arabic}

اینکه در بارهٴ جناب آقا میرزا محمّد علی و اهل ایشان علیهما بهآء اللّه
مرقوم داشتید بعد از عرض در ساحت اقدس امواج بحر بیان الهی بصورت این کلمات
باهرات ظاهر و لائح قال جلّت عظمته

\end{otherlanguage}

Īnke dar bāre-ye Janāb-i Āqā Mīrzā Muḥammad-ʿAlī va ahl-e īshān,
ʿalayhimā bahāʾuʾllāh, marqūm dāshtīd, baʿd az ʿarḍ dar sāḥat-e aqdas,
amvāj-e baḥr-e bayān-e ilāhī be-ṣūrat-e īn kalemāt-e bāherāt ẓāher va
lāʾeḥ. Qāla jallat ʿaẓamatuhu:

Concerning what you wrote about Jinab-i Aqa Mirza Muhammad-ʿAli and his
wife---upon them both be the glory of God---after it was presented at
the Most Holy Court, the waves of the ocean of divine utterance appeared
and shone forth in the form of these radiant words. He said, exalted is
His grandeur:

\subsection{83}\label{section-2503}

\begin{otherlanguage}{arabic}

بسمی العزیز العظیم

\end{otherlanguage}

Bi-smī al-ʿAzīzi al-ʿAẓīm.

In My name, the Almighty, the Magnificent.

\subsection{84}\label{section-2504}

\begin{otherlanguage}{arabic}

یا محمّد قبل علی قد سمعنا ذکرک فی کتاب من احبّنی و احبّه اللّه و انزلنا لک
آیات لا تعادلها ما فی خزائن الملوک طوبی لمن فاز الیوم بعرفان اللّه الّذی
اذا ظهر نطقت الأشیآء قد اتی الوعد و هذا هو الموعود انّا اظهرنا الأمر و
انزلنا الآیات و دعونا الکلّ الی اللّه المهیمن القیّوم من النّاس من نبذ الهدی
و اخذ الهوی بما اتّبع کلّ غافل محجوب قد ورد علی المظلوم فی سبیل اللّه ما
ناح به الملأ الأعلی و سکّان مدائن الأسمآء ولکنّ القوم اکثرهم لا یفقهون
فیکلّ الأحوال قمنا علی الأمر علی شأن ما منعتنا مدافع العالم و لا صفوف
الجنود و اوّل من اعترض علینا علمآء الأرض قد نقضوا میثاق اللّه و عهده و
قاموا علینا کما قاموا من قبل علی محمّد رسول اللّه و من قبله علی الرّوح
ینادیهم المیزان و هم لا یسمعون و ینذرهم الصّراط و هم لا یشعرون قد تمسّکوا
بالظّنون و الأوهام معرضین عن انوار الیقین کذلک یقصّ لک قلمی الأعلی فی هذا
السّجن الممنوع ان احفظ آیاتی ثمّ اقرأها بربوات اهل البهآء هذا ما یأمرک من
عنده کتاب محفوظ البهآء علیک و علی ضلعک الّتی فازت برحیق محبّة ربّها العزیز
الودود انتهی

\end{otherlanguage}

Yā Muḥammad qabla ʿAlī, qad samiʿnā dhikraka fī kitābi man aḥabbanī va
aḥabbahuʾllāh, va anzalnā laka āyātin lā tuʿādiluhā mā fī
khazāʾiniʾl-mulūk. Ṭūbā li-man fāzaʾl-yawma bi-ʿirfāniʾllāhiʾlladhī idhā
ẓahara naṭaqatiʾl-ashyāʾu: Qad atāʾl-waʿdu va hādhā huwaʾl-mawʿūd. Innā
aẓharnāʾl-amra va anzalnāʾl-āyāti va daʿawnāʾl-kulla
ilāʾllāhiʾl-muhayminiʾl-qayyūm. Minaʾl-nāsi man nabadhaʾl-hudā va
akhadhaʾl-hawā bimāʾttabaʿa kulla ghāfilin maḥjūb. Qad warada
ʿalāʾl-maẓlūmi fī sabīliʾllāhi mā nāḥa bihiʾl-malaʾuʾl-aʿlā va sukkānu
madāʾiniʾl-asmāʾ, va lākinnaʾl-qawma aktharuhum lā yafqahūn. Fī
kulliʾl-aḥwāli qumnā ʿalāʾl-amri ʿalā shaʾnin mā manaʿatnā
madāfiʿuʾl-ʿālami va lā ṣufūfuʾl-junūd. Va awwalu man iʿtaraḍa ʿalaynā
ʿulamāʾuʾl-arḍ. Qad naqaḍū mīthāqaʾllāhi va ʿahdahu va qāmū ʿalaynā kamā
qāmū min qablu ʿalā Muḥammadin rasūliʾllāhi va min qablihi ʿalāʾl-Rūḥ.
Yunādīhimuʾl-mīzānu va hum lā yasmaʿūn, va yundhiruhumuʾl-ṣirāṭu va hum
lā yashʿurūn. Qad tamassakū biʾl-ẓunūni vaʾl-awhāmi muʿriḍīna ʿan
anwāriʾl-yaqīn. Kadhālika yaquṣṣu laka qalamī al-aʿlā fī
hādhāʾl-sijniʾl-mamnūʿ. Aniʾḥfaẓ āyātī thummaʾqraʾhā bi-rabawāti
ahliʾl-bahāʾ. Hādhā mā yaʾmuruka man ʿindahu kitābun maḥfūẓ. Al-bahāʾu
ʿalayka va ʿalā ḍilʿikaʾllatī fāzat bi-raḥīqi maḥabbati
rabbihāʾl-ʿazīziʾl-wadūd. Intahā.

O Muhammad before ʿAli! We have heard mention of you in a letter from
one who has loved Me and whom God has loved, and We have sent down for
you verses that nothing in the treasuries of kings can equal. Blessed is
whoever has attained today the recognition of God, at whose appearance
all things declared: ``The promise has come, and this is the Promised
One!'' We have indeed manifested the Cause, sent down the verses, and
summoned everyone to God, the All-Protecting, the Self-Subsisting. Among
the people are those who have cast aside guidance and embraced desire
because they have followed every heedless and veiled soul. In the path
of God, the Wronged One has suffered what has caused the Concourse on
high and the inhabitants of the cities of names to lament; yet most of
the people do not understand. Under all circumstances We arose to uphold
the Cause in such a way that neither the cannons of the world nor the
ranks of armies could hinder Us. The first to oppose Us were the
religious scholars of the earth. They broke God's covenant and His
pledge and rose against Us, just as they had previously risen against
Muhammad, the Messenger of God, and, before Him, against the Spirit. The
Balance calls to them, but they do not hear; the Path warns them, but
they do not perceive. They have clung to conjectures and illusions,
turning away from the lights of certainty. Thus does My Most Exalted Pen
recount to you in this inaccessible prison. Preserve My verses, then
recite them upon the heights among the people of Baha. This is what He
with whom is a preserved Book commands you. Glory be upon you and upon
your wife, who has attained the choice wine of the love of her Lord, the
Almighty, the All-Loving. End of quotation.

\subsection{85}\label{section-2505}

\begin{otherlanguage}{arabic}

یا محبوب فؤادی حسب الأمر این آیات عظیمهٴ منیعه را بر ایشان القا نمائید و
از قبل حقّ تکبیر برسانید شاید از نفحات بیان سلطان امکان بفرح حقیقی فائز
شوند و از فزع اکبر محفوظ مانند اینعبد هم خدمت ایشان تکبیر و سلام عرض
مینماید

\end{otherlanguage}

Yā maḥbūb-e foʾādī, ḥasb-e amr īn āyāt-e ʿaẓīme-ye manīʿe rā bar īshān
elqā namāʾīd va az qebal-e Ḥaqq takbīr beresānīd. Shāyad az nafaḥāt-e
bayān-e Solṭān-e emkān be-faraḥ-e ḥaqīqī fāʾez shavand va az fazaʿ-e
akbar maḥfūẓ mānand. Īn ʿabd ham khedmat-e īshān takbīr va salām ʿarḍ
mīnamāyad.

O beloved of my heart! In accordance with the command, convey these
mighty and exalted verses to them and deliver greetings on behalf of
God, that through the breezes of the utterance of the Sovereign of all
existence they may attain true joy and remain safe from the most great
terror. This servant also offers them greetings and peace.

\subsection{86}\label{section-2506}

\begin{otherlanguage}{persian}

ذکر جناب آقا سیّد علی علیه بهآء اللّه از اهل خا فرموده بودند و هم‌چنین ذکر
اشتعالشان را بنار محبّت الهی بعد از حضور و اذن و عرض آفتاب عنایت بصورت
اینکلمات عالیات تجلّی فرمود قوله جلّ اجلاله

\end{otherlanguage}

Dhekr-e Janāb-i Āqā Siyyid ʿAlī, ʿalayhi bahāʾuʾllāh, az ahl-e Khā
farmūde būdand va hamchenīn dhekr-e eshteʿāl-e shān rā be-nār-e
maḥabbat-e ilāhī. Baʿd az ḥoḍūr va edhn va ʿarḍ, āftāb-e ʿenāyat
be-ṣūrat-e īn kalemāt-e ʿāliyāt tajallī farmūd. Qawluhu jalla ijlāluhu:

Mention had been made of Jinab-i Aqa Siyyid ʿAli---upon him be the glory
of God---of Kha, and likewise of his being aflame with the fire of
divine love. After entering the Presence, obtaining permission, and
presenting the matter, the sun of loving care shone forth in the form of
these exalted words. His words, exalted is His majesty:

\subsection{87}\label{section-2507}

\begin{otherlanguage}{arabic}

بسمی النّاطق من افق الملکوت

\end{otherlanguage}

Bi-smī al-nāṭiqi min ufuqiʾl-malakūt.

In My name, speaking from the horizon of the Kingdom.

\subsection{88}\label{section-2508}

\begin{otherlanguage}{arabic}

کتاب انزله المظلوم لمن اقبل الی اللّه فی یوم بشّر به رسل اللّه من قبل و کلّ
کتاب منزول یا علی فاذکر ما اقامک علی هذا الأمر و ما ایقظک فی هذا الیوم و
من اسمعک ندآء المظلوم الّذی رأی فی سبیل اللّه ما ناح به المقرّبون و من
اخبرک بهذا المسجون و من حدّثک هذا الأمر الّذی به ارتفع حفیف سدرة المنتهی و
انار الأفق الأعلی و نطقت الأشیآء فیملکوت الانشآء قد اتی من کان مستوراً فی
علم اللّه و مسطوراً من قلم الأمر فی کتابه المختوم طوبی لک بما سمعت و عرفت
و اقبلت الی من اتی من سمآء الأمر بسلطان مشهود دع النّاس بأهوائهم و تشبّث
بذیل عنایة ربّک مالک الوجود انّا نوصیک و الّذین آمنوا بما یفرح به قلب
المظلوم و یرتفع امر اللّه العزیز الودود انّه لو یجد عرف اخلاقه و ما امر به
عباده فی کتابه لیفرح کذلک نطق لسان البیان فی ملکوت العرفان فضلاً من عنده
و هو الحقّ علّام الغیوب البهآء علیک و علی من اشتعل بنار محبّة اللّه فی ایّامه
و علی کلّ قائم نبذ القعود فیهذا الیوم الموعود ثمّ الذّکر و البهآء علی من
انتبه من نسمات الوحی التّی تمرّ من شطر سجنی العزیز المحبوب انتهی

\end{otherlanguage}

Kitābun anzalahuʾl-maẓlūmu li-man aqbala ilāʾllāhi fī yawmin bashshara
bihi rusuluʾllāhi min qablu va kullu kitābin munzūl. Yā ʿAlī, faʾdhkur
mā aqāmaka ʿalā hādhāʾl-amri va mā ayqaẓaka fī hādhāʾl-yawmi va man
asmaʿaka nidāʾaʾl-maẓlūmiʾlladhī raʾā fī sabīliʾllāhi mā nāḥa
bihiʾl-muqarrabūn, va man akhbaraka bi-hādhāʾl-masjūni va man ḥaddathaka
hādhāʾl-amraʾlladhī bihiʾrtafaʿa ḥafīfu sidratiʾl-muntahā va
anāraʾl-ufuquʾl-aʿlā va naṭaqatiʾl-ashyāʾu fī malakūtiʾl-inshāʾ. Qad atā
man kāna mastūran fī ʿilmiʾllāhi va masṭūran min qalamiʾl-amri fī
kitābihiʾl-makhtūm. Ṭūbā laka bimā samiʿta va ʿarafta va aqbalta ilā man
atā min samāʾiʾl-amri bi-sulṭānin mashhūd. Daʿiʾl-nāsa bi-ahwāʾihim va
tashabbath bi-dhayli ʿināyati rabbika mālikiʾl-wujūd. Innā nūṣīka
vaʾlladhīna āmanū bimā yafraḥu bihi qalbuʾl-maẓlūmi va yartafiʿu
amruʾllāhiʾl-ʿazīziʾl-wadūd. Innahu law yajidu ʿarfa akhlāqihi va mā
amara bihi ʿibādahu fī kitābihi la-yafraḥ. Kadhālika naṭaqa
lisānuʾl-bayāni fī malakūtiʾl-ʿirfāni faḍlan min ʿindihi va huwaʾl-ḥaqqu
ʿallāmuʾl-ghuyūb. Al-bahāʾu ʿalayka va ʿalā man ishtaʿala bi-nāri
maḥabbatiʾllāhi fī ayyāmihi va ʿalā kulli qāʾimin nabadhaʾl-quʿūda fī
hādhāʾl-yawmiʾl-mawʿūd. Thummaʾl-dhikru vaʾl-bahāʾu ʿalā man intabaha
min nasamātiʾl-waḥyiʾllatī tamurru min shaṭri sijnī al-ʿazīziʾl-maḥbūb.
Intahā.

A Tablet that the Wronged One has sent down for one who has turned
toward God on a Day whose glad tidings were proclaimed beforehand by the
Messengers of God and every revealed Book. O ʿAli! Call to mind what
raised you up to uphold this Cause, what awakened you on this Day, and
who caused you to hear the call of the Wronged One, who has witnessed in
the path of God what has made those brought near to Him lament. Who
informed you of this Prisoner, and who told you of this Cause, through
which the rustling of the Lote Tree of the Utmost Boundary rose, the
most exalted horizon shone, and all things spoke in the kingdom of
creation? He has come who was concealed in the knowledge of God and
inscribed by the Pen of the Cause in His sealed Book. Blessed are you
for having heard, recognized, and turned toward Him who has come from
the heaven of the Cause with manifest sovereignty. Leave the people to
their desires and cling to the hem of the loving care of your Lord, the
Sovereign of existence. We exhort you and those who have believed to do
what will gladden the heart of the Wronged One and exalt the Cause of
God, the Almighty, the All-Loving. Should He perceive the fragrance of
His own character and of what He has commanded His servants in His Book,
He will indeed rejoice. Thus has the Tongue of Utterance spoken in the
kingdom of recognition, as a grace from Him. He is the Truth, the Knower
of all hidden things. Glory be upon you, upon whoever has been set
ablaze with the fire of the love of God in His days, and upon everyone
who has arisen and cast aside inaction on this promised Day. Remembrance
and glory also be upon whoever has awakened through the breezes of
revelation that pass from the direction of My mighty and beloved prison.
End of quotation.

\subsection{89}\label{section-2509}

\begin{otherlanguage}{persian}

در اینحین مبارک بشأنی آیات نازل و بیّنات باهر که فی‌الحقیقه صدهزار مثل این
عبد از ذکر و وصف و تحریر و تقریرش عاجز و قاصر العظمة للّه ربّنا و ربّکم و
مقصودنا و مقصودکم و مقصود من فی الملک و الجبروت نمیدانم این خلقرا چه
سکری اخذ نموده امّت قبل استدلال مینمودند که چند فقره در کتاب نازل که حقّ
جلّ جلاله از امور بعد خبر داده و این امور واقع گشته و نظر بوقوع آن او را
از حجّت و برهان میشمردند و حال آن جناب و جمیع دوستان شاهد و گواهند که
اموراتیکه هیچ عاقلی تصوّر آن نمینمود از قلم اعلی بکمال تصریح نازل و جمیع
واحداً بعد واحد از قوّه بفعل آمد و از باطن بظاهر جلوه نمود بلکه آنچه از
بعد ظاهر شده و یا بشود واحداً بعد واحد ذکرش در کتاب نازل و در ظاهر ظاهر
مشاهده شد اینفقره هم محدود نبوده چه که از یک و دو و سه و فوق آن و فوق
فوق آن ظاهر و باهر مع‌ذلک کلّ غافل نائم میّت الّا من حفظته ید الفضل آیا احدی
میتواند بگوید در هیچ ظهوری آیاتش اکمل و یا بیّناتش اظهر و یا سدرهٴ حجّت
اثمر از این سدرهٴ مبارکهٴ باسقه بوده لا وربّی مگر بالمرّه از انصاف عاری
باشد این خادم آنچه عرض نموده و یا مینماید در پرده بوده و هست چه که اذن
بر عرایض بعضی از امور نداشته و ندارد ختم اللّه علی فمه الأمر بیده یفعل و
یحکم کیف یشآء و هو المقتدر المختار

\end{otherlanguage}

Dar īn ḥīn-e mobārak be-shaʾnī āyāt nāzel va bayyenāt bāher ke
fīʾl-ḥaqīqe ṣad hezār methl-e īn ʿabd az dhekr va vaṣf va taḥrīr va
taqrīr-ash ʿājez va qāṣer. Al-ʿaẓamatu li-llāhi rabbinā va rabbikum va
maqṣūdinā va maqṣūdikum va maqṣūdi man fīʾl-mulki vaʾl-jabarūt.
Nemīdānam īn khalq rā che sokrī akhdh namūde. Ommat-e qabl estedlāl
mīnamūdand ke chand feqre dar ketāb nāzel ke Ḥaqq jalla jalāluhu az
omūr-e baʿd khabar dāde va īn omūr vāqeʿ gashte va naẓar be-voqūʿ-e ān ū
rā az ḥojjat va borhān mīshomordand. Va ḥāl ān janāb va jamīʿ-e dūstān
shāhed va govāhand ke omūrātī ke hīch ʿāqelī taṣavvor-e ān nemīnamūd az
Qalam-e Aʿlā be-kamāl-e taṣrīḥ nāzel va jamīʿ vāḥedan baʿd-e vāḥed az
qovve be-feʿl āmad va az bāṭen be-ẓāher jelve namūd. Balke ānche az baʿd
ẓāher shode va yā beshavad vāḥedan baʿd-e vāḥed dhekr-ash dar ketāb
nāzel va dar ẓāher ẓāher moshāhede shod. Īn feqre ham maḥdūd nabūde che
ke az yek va do va se va fowq-e ān va fowq-e fowq-e ān ẓāher va bāher.
Maʿa dhālek koll ghāfel, nāʾem, mayyet, illā man ḥafiẓat-hu yaduʾl-faḍl.
Āyā aḥadī mītavānad begūyad dar hīch ẓohūrī āyāt-ash akmal va yā
bayyenāt-ash aẓhar va yā sedre-ye ḥojjat athmar az īn sedre-ye
mobārake-ye bāseqe būde? Lā va rabbī, magar beʾl-marre az enṣāf ʿārī
bāshad. Īn khādem ānche ʿarḍ namūde va yā mīnamāyad dar parde būde va
hast, che ke edhn bar ʿarāyeḍ-e baʿḍī az omūr nadāshte va nadārad.
Khatamaʾllāhu ʿalā famihi. Al-amru bi-yadihi yafʿalu va yaḥkumu kayfa
yashāʾu va huwaʾl-muqtadiruʾl-mukhtār.

At this blessed time, verses are being revealed and clear proofs are
shining forth to such a degree that a hundred thousand like this servant
would truly be powerless and unable to recount, describe, record, or
explain them. Grandeur belongs to God, our Lord and your Lord, the One
we seek and you seek, and the One sought by all who dwell in the earthly
world and the realm of power. I do not know what intoxication has seized
these people. The former community used to argue that a few passages had
been revealed in the Book in which God, exalted be His glory, foretold
future events, and that these events had taken place; because they had
occurred, they counted this as evidence and proof. Now you and all the
friends bear witness that events no rational person could have imagined
were revealed with complete explicitness by the Most Exalted Pen, and
all, one after another, passed from potentiality into actuality and
emerged from concealment into manifestation. Indeed, whatever has since
appeared or may yet appear has, one after another, been mentioned in the
Book and has been visibly witnessed in the outward world. Nor has this
been limited, for not one, two, or three instances, but more than these
and far more still, have become manifest and evident. Nevertheless, all
are heedless, asleep, and dead, except those whom the hand of grace has
protected. Can anyone say that in any revelation its verses have been
more perfect, its clear proofs more evident, or its tree of proof more
fruitful than this blessed, towering Tree? No, by my Lord---not unless
he is utterly devoid of fairness. Whatever this servant has submitted or
submits has been and remains veiled, for he has not had, and does not
have, permission to report certain matters. God has sealed his mouth.
The command is in His hand; He acts and decrees as He wills, and He is
the All-Powerful, the Unconstrained.

\subsection{90}\label{section-2510}

\begin{otherlanguage}{persian}

ذکر مخدوم معظّم حضرت زین ب ر الّذی هاجر الی اللّه و فاز فی ارض السّرّ فرمودید
انشآءاللّه در جمیع احوال بعنایات مخصوصهٴ حقّ و فیوضاتش فائز باشند لازال
امام عین بوده و هستند هیچوقت فراموش نشده و نخواهند شد انشآءاللّه در جمیع
قرون و اعصار باسم حقّ و خدمت حقّ مذکور باشند بعض احیان مخصوص ذکرشان از
لسان مالک قدم جاری یا محبوب فؤادی از جمیع مراتب گذشته شجر وفا شجر
مبارکیست اثمارش محبوب و آثار و اسرارش باقی و دائم اگر خوب ملاحظه شود
اوست قائد جنود حبّ و اوست پیشرو اهل ودّ و راستی از صمیم قلب خدمت ایشان
تکبیر و سلام عرض مینمایم

\end{otherlanguage}

Dhekr-e makhdūm-e moʿaẓẓam Ḥaḍrat-e Zayn bā rā, alladhī hājara ilāʾllāhi
va fāza fī arḍiʾl-sirr, farmūdīd. In shāʾaʾllāh dar jamīʿ-e aḥvāl
be-ʿenāyāt-e makhṣūṣe-ye Ḥaqq va foyūḍāt-ash fāʾez bāshand. Lā zāl
emām-e ʿayn būde va hastand. Hīch vaqt farāmūsh nashode va nakhāhand
shod. In shāʾaʾllāh dar jamīʿ-e qorūn va aʿṣār be-nām-e Ḥaqq va
khedmat-e Ḥaqq madhkūr bāshand. Baʿḍ-e aḥyān makhṣūṣ dhekr-e shān az
lesān-e Mālek-e qedam jārī. Yā maḥbūb-e foʾādī, az jamīʿ-e marāteb
godhashte, shajar-e vafā shajar-e mobārakīst. Athmār-ash maḥbūb va āthār
va asrār-ash bāqī va dāʾem. Agar khūb molāḥeẓe shavad, ūst qāʾed-e
jonūd-e ḥobb va ūst pīshrow-e ahl-e vodd va rāstī. Az ṣamīm-e qalb
khedmat-e īshān takbīr va salām ʿarḍ mīnamāyam.

You mentioned the honored gentleman, his honor Zayn B. R., who migrated
toward God and attained His presence in the Land of Mystery. God
willing, may he under all circumstances receive God's special favors and
outpourings of grace. He has always been and remains before our eyes. He
has never been forgotten and never will be. God willing, may he be
remembered throughout all centuries and ages in association with God's
name and service to Him. At times, mention of him in particular flows
from the tongue of the Sovereign of Eternity. O beloved of my heart!
Beyond all other considerations, the tree of faithfulness is a blessed
tree. Its fruits are cherished, and its effects and mysteries are
lasting and enduring. If one considers carefully, it is the commander of
the hosts of love and the leader of the people of affection and
truthfulness. From the depths of my heart I offer him greetings and
peace.

\subsection{91}\label{section-2511}

\begin{otherlanguage}{persian}

عریضهٴ ایشان خدمت حضرت غصن اللّه الأعظم روحی و ذاتی و کینونتی لتراب قدومه
الفدآء رسید و انشآءاللّه جواب عنایت میفرمایند اینکه در بعض امور توقّف رفته
نظر بحکمت است ایّام ایّامیست که اگر امتیاز ظاهر شود بسا میشود سبب ضرّ گردد
یعنی از برای نفس ممتازه اعدا نظر بآنکه حال بیقین مبین دانسته‌اند که حقّ جلّ
جلاله فتنه و فساد و اذیّت و قتل را منع فرموده‌اند لذا در بعض امور فی‌الجمله
متوقّفند ولکن در سرّ سرّ مترصّد اگر فی‌الحقیقه نفسیرا منشأ اثر کلّی در آنجهات
بدانند بتشبّثات مختلفه تقصیری بر حسب ظاهر اثبات مینمایند مانند تقصیریکه
آن طیر ظالم از برای کبک بیچاره ثابت نمود و او را درید مقصود اینفانی آنکه
باینجهات بعض اوامر الهیّه متوقّف

\end{otherlanguage}

ʿArīḍe-ye īshān khedmat-e Ḥaḍrat-e Ghoṣnuʾllāhiʾl-Aʿẓam, rūḥī va dhātī
va kaynūnatī li-turābi qudūmihiʾl-fidāʾ, resīd va in shāʾaʾllāh javāb
ʿenāyat mīfarmāyand. Īnke dar baʿḍ-e omūr tavaqqof rafte naẓar be-ḥekmat
ast. Ayyām ayyāmīst ke agar emtīāz ẓāher shavad basā mīshavad sabab-e
ḍorr gardad, yaʿnī az barā-ye nafs-e momtāze. Aʿdā naẓar be-ānke ḥāl
be-yaqīn-e mobīn dāneste-and ke Ḥaqq jalla jalāluhu fetne va fasād va
adhiyyat va qatl rā manʿ farmūde-and, ledhā dar baʿḍ-e omūr fīʾl-jomle
motavaqqef-and, va lāken dar serr-e serr motaraṣṣed. Agar fīʾl-ḥaqīqe
nafsī rā manshaʾ-e athar-e kollī dar ān jahāt bedānand,
be-tashabbothāt-e mokhtalefe taqṣīrī bar ḥasab-e ẓāher ethbāt
mīnamāyand, mānand-e taqṣīrī ke ān ṭayr-e ẓālem az barā-ye kabk-e
bīchāre thābet namūd va ū rā darīd. Maqṣūd-e īn fānī ānke be-īn jahāt
baʿḍ-e avāmer-e ilāhiyye motavaqqef.

His letter reached His Holiness the Most Great Branch of God---may my
spirit, my self, and my being be sacrificed for the dust beneath His
feet---and, God willing, He will graciously provide a reply. The delay
in certain matters is out of wisdom. These are days when, if distinction
becomes apparent, it may well cause harm---that is, to the distinguished
person. Since the enemies now know with clear certainty that God,
exalted be His glory, has forbidden sedition, disorder, harassment, and
killing, they are somewhat restrained in certain matters; in the utmost
secrecy, however, they lie in wait. If they truly recognize someone as a
source of substantial influence in those regions, they resort to various
devices to establish some outward fault against that person, like the
fault that the cruel bird established against the helpless partridge
before tearing it apart. What this transitory servant means is that, for
these reasons, certain divine directives are being held in abeyance.

\subsection{92}\label{section-2512}

\begin{otherlanguage}{persian}

اینکه در بارهٴ امة‌اللّه هدهد علیها بهآء اللّه مرقوم داشتید وقتی از اوقات
مخصوص این فانی مجدّد ذکر ایشانرا نموده و بطراز عنایت حقّ جلّ جلاله فائز
شدند و لسان عظمت باینکلمهٴ علیا ناطق قوله تعالی

\end{otherlanguage}

Īnke dar bāre-ye Amatuʾllāh Hudhud, ʿalayhā bahāʾuʾllāh, marqūm dāshtīd,
vaqtī az avqāt makhṣūṣ īn fānī mojaddad dhekr-e īshān rā namūde va
be-ṭarāz-e ʿenāyat-e Ḥaqq jalla jalāluhu fāʾez shodand va lesān-e
ʿaẓamat be-īn kaleme-ye ʿolyā nāṭeq. Qawluhu taʿālā:

Concerning what you wrote about the handmaid of God, Hudhud---upon her
be the glory of God---on one occasion this transitory servant again made
particular mention of her, and she was favored with the ornament of the
loving care of God, exalted be His glory. The Tongue of Grandeur spoke
these exalted words. His words, exalted is He:

\subsection{93}\label{section-2513}

\begin{otherlanguage}{arabic}

یا ایّها النّاظر الی وجهی و القائم علی خدمة امری بشّرها من قبلی و کبّر علیها
امراً من لدن ربّک المشفق الکریم انّا ذکرناها من قبل و بنت اسمی الأصدق الّذی
اقبل و توجّه الی شطر اللّه امراً من عنده الی ان دخل السّجن و قام لدی الباب و
سمع ندآء اللّه المقتدر العزیز الوهّاب انّا نکبّر علیها و علی امّها و علی
امائی اللّائی اقبلن الی اللّه مالک الرّقاب و نوصیهنّ بالعمل الخالص و ما
ینبغی لأیّام اللّه ربّ الأرباب کذلک ماج بحر البیان اذ یمشی الرّحمن فی اعلی
المقام انتهی

\end{otherlanguage}

Yā ayyuhā al-nāẓiru ilā wajhī wa-al-qāʾimu ʿalā khidmati amrī,
bashshirhā min qibalī wa-kabbir ʿalayhā amran min ladun rabbika
al-mushfiqi al-karīm. Innā dhakarnāhā min qablu wa-binta ismī al-Aṣdaq
alladhī aqbala wa-tawajjaha ilā shaṭri Allāhi amran min ʿindihi ilā an
dakhala al-sijna wa-qāma ladā al-bābi wa-samiʿa nidāʾa Allāhi
al-muqtadiri al-ʿazīzi al-wahhāb. Innā nukabbiru ʿalayhā wa-ʿalā ummihā
wa-ʿalā imāʾī allāʾī aqbalna ilā Allāhi māliki al-riqāb, wa-nūṣīhinna
bi-al-ʿamali al-khāliṣi wa-mā yanbaghī li-ayyāmi Allāhi rabbi al-arbāb.
Kadhālika māja baḥru al-bayāni idh yamshī al-Raḥmānu fī aʿlā al-maqām.
Intahā.

O you who look toward My face and arise to serve My Cause! Give her glad
tidings on My behalf and convey to her My greeting of glory, at the
bidding of your compassionate, generous Lord. We previously remembered
her and the daughter of him who bears My name, the Most Truthful, who
turned and directed himself toward God at His command, until he entered
the prison, stood at the gate, and heard the call of God, the
All-Powerful, the Almighty, the All-Bestowing. We bestow Our greeting of
glory upon her, upon her mother, and upon My handmaids who have turned
toward God, the Master of all, and We exhort them to pure deeds and
whatever befits the days of God, the Lord of lords. Thus has the ocean
of utterance surged as the All-Merciful walks in the highest station.
End.

\subsection{94}\label{section-2514}

\begin{otherlanguage}{arabic}

الحمد للّه فازت بما ارادت قد ذکرت و نزّلت لها و لبنتها ما تضوّع به عرف
عنایة ربّنا الکریم آخر دستخطّ آنمحبوب جان بمناجات مع اللّه منتهی در حین عرض
آن چون باینکلمهٴ آنجناب رسید که عرض نموده‌اند الهی و محبوبی معزّز بدار
نفسیرا که امر ترا معزّز بدارد و علّت اعزاز و اعلای امر تو گردد در اینمقام
شمس این بیان از افق ارادهٴ مالک امکان اشراق نمود قوله جلّ اجلاله

\end{otherlanguage}

Al-ḥamdu li-llāh, fāzat bi-mā arādat. Qad dhukirat wa-nuzzilat lahā
wa-li-bintihā mā taḍawwaʿa bihi ʿarfu ʿināyati rabbinā al-karīm. Ākhar-e
dastkhaṭṭ-e ān maḥbūb-e jān be-monājāt maʿa Allāh montahī. Dar ḥīn-e
ʿarż-e ān chūn be-īn kalame-ye ān janāb resīd ke ʿarż namūde-and, Ilāhī
va maḥbūbī, moʿazzaz bedār nafsī-rā ke amr-e to-rā moʿazzaz bedārad va
ʿellat-e eʿzāz va eʿlā-ye amr-e to gardad, dar īn maqām shams-e īn bayān
az ofoq-e erāde-ye Mālek-e emkān eshrāq namūd. Qawluhu jalla ijlāluhu.

Praise be to God, she attained what she desired. She was indeed
remembered, and there was revealed for her and her daughter that through
which the fragrance of our generous Lord's loving care was diffused. The
letter of that beloved of the soul concludes with a supplication to God.
While it was being presented, when the reading reached these words you
had written, ``My God and my Beloved, honor the soul who honors Your
Cause and becomes the means of exalting and elevating Your Cause,'' at
this point the sun of this utterance shone from the horizon of the will
of the Lord of all existence. His words, exalted be His majesty:

\subsection{95}\label{section-2515}

\begin{otherlanguage}{arabic}

یا ایّها الشّارب رحیق بیانی انّا اعززناک و نصرناک و رفعناک و ذکرناک لتذکر
امر ربّک العزیز الکریم

\end{otherlanguage}

Yā ayyuhā al-shāribu raḥīqa bayānī, innā aʿzaznāka wa-naṣarnāka
wa-rafaʿnāka wa-dhakarnāka li-tadhkura amra rabbika al-ʿazīzi al-karīm.

O you who drink the choice wine of My utterance! We have indeed honored
you, aided you, exalted you, and remembered you, so that you may make
mention of the Cause of your Lord, the Almighty, the Generous.

\subsection{96}\label{section-2516}

\begin{otherlanguage}{persian}

این عزّت تغییر نیابد و جنود عالم بر ستر آن قادر نه چه که حقّ عنایت فرموده
و از جانب حقّ است و زود است که آثارش در ارض ظاهر شود تمسّک بذکری و قم علی
خدمة امری بالرّوح و الرّیحان و الحکمة و البیان کذلک یأمرک الرّحمن من هذا
المقام العزیز المنیع آن جناب و بعضی از دوستان که صادق و محرمند و بطراز
محبّت فی‌الحقیقه مزیّن مشورت نمائید و بآنچه سبب اتّحاد و اتّفاق است قیام کنید
لعمری سوف یرفع اللّه من یحبّه و یظهر مقامه بین خلقه انّه ولیّ من والاه و
ناصر من نصره فی یومه العزیز البدیع انتهی

\end{otherlanguage}

Īn ʿezzat taghyīr nayābad va jonūd-e ʿālam bar satr-e ān qāder na, che
ke Ḥaqq ʿenāyat farmūde va az jāneb-e Ḥaqq ast va zūd ast ke āthārash
dar arż ẓāher shavad. Tamassak bi-dhikrī wa-qum ʿalā khidmati amrī
bi-al-rūḥi wa-al-rayḥāni wa-al-ḥikmati wa-al-bayān. Kadhālika yaʾmuruka
al-Raḥmānu min hādhā al-maqāmi al-ʿazīzi al-manīʿ. Ān janāb va baʿżī az
dūstān ke ṣādeq va maḥramand va be-ṭerāz-e maḥabbat fī al-ḥaqīqe
mozayyan, mashverat namāyīd va be-ānche sabab-e etteḥād va ettefāq ast
qiyām konīd. La-ʿamrī sawfa yarfaʿu Allāhu man yuḥibbuhu wa-yuẓhiru
maqāmahu bayna khalqihi. Innahu waliyyu man wālāhu wa-nāṣiru man
naṣarahu fī yawmihi al-ʿazīzi al-badīʿ. Intahā.

This honor will not change, nor can the armies of the world conceal it,
for God has bestowed it and it comes from God; soon its effects will
appear on earth. Hold fast to My remembrance and arise to serve My Cause
with joy and gladness, wisdom and utterance. Thus does the All-Merciful
command you from this mighty, unassailable station. Consult with some of
the friends who are sincere and trustworthy and truly adorned with love,
and undertake whatever will foster unity and concord. By My life! God
will exalt whoever loves Him and reveal his station among His creatures.
He is indeed the protector of whoever befriends Him and the helper of
whoever helps Him in His mighty, wondrous Day. End.

\subsection{97}\label{section-2517}

\begin{otherlanguage}{persian}

در اینمقام بحر بیان بشأنی موّاج که حفظ و خیال و قوّه از این عبد رفت انّا
للّه و انّا الیه راجعون و الأمر بید اللّه دیگر نمیدانم کاد ان یطیر الجسد و
کیف الرّوح بعد از عرض مناجات و تنزیل آیات و اظهار عنایات و ظهورات رحمت و
بروزات فضل ذکر نفوس مقبلهٴ مستقیمه علیهم بهآء اللّه و عنایته که در دستخطّ
آن محبوب بود در ساحت امنع اقدس اعلی عرض شد هذا ما اشرقت به شمس الحقیقة
من افق سمآء البیان قوله جلّت عظمته و عظم احسانه و کبر کبریائه

\end{otherlanguage}

Dar īn maqām baḥr-e bayān be-shaʾnī mavvāj ke ḥefẓ va khayāl va qovve az
īn ʿabd raft. Innā li-llāhi wa-innā ilayhi rājiʿūn, wa-al-amru bi-yadi
Allāh. Dīgar nemīdānam. Kāda an yaṭīra al-jasadu wa-kayfa al-rūḥ. Baʿd
az ʿarż-e monājāt va tanzīl-e āyāt va eẓhār-e ʿenāyāt va ẓohūrāt-e
raḥmat va borūzāt-e fażl, dhekr-e nofūs-e moqbele-ye mostaqīme, ʿalayhim
bahāʾu Allāhi wa-ʿināyatuhu, ke dar dastkhaṭṭ-e ān maḥbūb būd dar
sāḥat-e amnaʿ-e aqdas-e aʿlā ʿarż shod. Hādhā mā ashraqat bihi shamsu
al-ḥaqīqati min ufuqi samāʾi al-bayān. Qawluhu jallat ʿaẓamatuhu
wa-ʿaẓuma iḥsānuhu wa-kabura kibriyāʾuhu.

At this point the ocean of utterance surged so mightily that this
servant lost his memory, thought, and strength. We belong to God and to
Him we return, and the matter is in God's hands. I know no more. The
body almost took flight---how much more the spirit! After the
supplication had been presented, verses revealed, loving care expressed,
mercy manifested, and grace displayed, mention of the steadfast souls
who had turned toward Him---upon them be God's glory and loving
care---whose names appeared in that beloved one's letter was presented
in the most unassailable, most holy, most exalted court. This is what
the Sun of Truth shone forth with from the horizon of the heaven of
utterance. His words, exalted be His grandeur, great His beneficence,
and immense His majesty:

\subsection{98}\label{section-2518}

\begin{otherlanguage}{arabic}

بسمی المذکور فی الصّحف و الزّبر و الألواح

\end{otherlanguage}

Bi-ismī al-madhkūri fī al-ṣuḥufi wa-al-zuburi wa-al-alwāḥ.

In My name, mentioned in the Scriptures, the Books, and the Tablets.

\subsection{99}\label{section-2519}

\begin{otherlanguage}{arabic}

کتاب نزّل من لدی اللّه مالک القدم علی الأحزاب و الأمم طوبی لسمیع سمع و
لبصیر رأی و ویل لکلّ منکر مرتاب ان یا قلمی الأعلی ان اذکر الّذین تشبّثوا
بأذیال ردآء عنایة ربّک مالک الأسمآء بذکر تنجذب به القلوب و الأرواح ان
اذکر عبدی العظیم الّذی اعترف بهذا النّبأ الأعظم و شرب رحیق العرفان من
ایادی عطآء ربّک الرّحمن الّذی ینطق فیکلّ شأن انّه لا اله الّا انا العزیز
الوهّاب

\end{otherlanguage}

Kitābun nuzzila min ladā Allāhi māliki al-qidami ʿalā al-aḥzābi
wa-al-umam. Ṭūbā li-samīʿin samiʿa wa-li-baṣīrin raʾā wa-waylun li-kulli
munkirin murtāb. An yā qalamī al-aʿlā, an udhkur alladhīna tashabbathū
bi-adhyāli ridāʾi ʿināyati rabbika māliki al-asmāʾi bi-dhikrin
tanjadhibu bihi al-qulūbu wa-al-arwāḥ. An udhkur ʿabdī al-ʿAẓīm alladhī
iʿtarafa bi-hādhā al-nabaʾi al-aʿẓami wa-shariba raḥīqa al-ʿirfāni min
ayādī ʿaṭāʾi rabbika al-Raḥmāni alladhī yanṭiqu fī kulli shaʾnin, innahu
lā ilāha illā anā al-ʿazīzu al-wahhāb.

A Book sent down from God, the Lord of eternity, to the peoples and
nations. Blessed is the hearer who has heard and the discerning one who
has seen, and woe to every doubting denier! O My Most Exalted Pen!
Remember those who have clung to the hem of the robe of the loving care
of your Lord, the Lord of Names, with a remembrance that will draw
hearts and spirits to Him. Remember My servant ʿAẓīm, who acknowledged
this Most Great Announcement and drank the choice wine of recognition
from the hands of the bounty of your Lord, the All-Merciful, who
proclaims in every circumstance, ``There is indeed no God but Me, the
Almighty, the All-Bestowing.''

\subsection{100}\label{section-2520}

\begin{otherlanguage}{arabic}

انّا ذکرناک من قبل من شطر البقعة البیضآء من سدرة المنتهی بذکر تضوّع منه
عرف البقآء بین الأحزاب ان اشکر اللّه بهذا الذّکر الأعظم فی سجنی الأعظم و
قل لک الحمد یا مطلع الوحی و منزل الآیات لعمری لا تعادل بحرف من آیاتی
کنوز الأرض کلّها یشهد بذلک من عنده امّ الکتاب

\end{otherlanguage}

Innā dhakarnāka min qablu min shaṭri al-buqʿati al-bayḍāʾi min Sidrati
al-Muntahā bi-dhikrin taḍawwaʿa minhu ʿarfu al-baqāʾi bayna al-aḥzāb. An
ushkur Allāha bi-hādhā al-dhikri al-aʿẓami fī sijnī al-aʿẓami wa-qul
laka al-ḥamdu yā maṭliʿa al-waḥyi wa-munzila al-āyāt. La-ʿamrī lā
tuʿādalu bi-ḥarfin min āyātī kunūzu al-arḍi kulluhā. Yashhadu bi-dhālika
man ʿindahu Ummu al-Kitāb.

We previously remembered you from the direction of the White Spot, from
the Tree beyond which there is no passing, with a remembrance from which
the fragrance of eternity spread among the peoples. Thank God for this
most great remembrance in My Most Great Prison, and say, ``Praise be to
You, O Dawning-Place of revelation and Revealer of verses!'' By My life!
All the treasures of the earth cannot equal a single letter of My
verses. To this bears witness He with whom is the Mother Book.

\subsection{101}\label{section-2521}

\begin{otherlanguage}{arabic}

قل یا ملأ الأحزاب تاللّه انّ الکتاب ینطق فی قطب العالم و یدع الأمم الی
مالک القدم ولکنّ القوم فی نوم عجاب قل اتی من خضعت له کتب الأرض کلّها ان
اعتبروا یا اولی الأبصار اتّقوا الرّحمن یا ملأ الامکان ان انصفوا فیهذا
الأمر الّذی اذا ظهر نفخ فی الصّور و مرّت الجبال قل فاتوا یا ملأ الأرض بما
عندکم و نزن ما عندنا و عندکم بهذا المیزان الّذی وضع بالعدل امراً من لدی
اللّه ربّ الأرباب هذا یوم فیه ینادی المیزان بأعلی النّدآء بأنّی انا الممیّز
العلیم من لدن مالک الأدیان

\end{otherlanguage}

Qul yā malaʾa al-aḥzāb, ta-Allāhi inna al-kitāba yanṭiqu fī quṭbi
al-ʿālami wa-yadʿu al-umama ilā māliki al-qidami wa-lākinna al-qawma fī
nawmin ʿujāb. Qul atā man khaḍaʿat lahu kutubu al-arḍi kulluhā. An
iʿtabirū yā ulī al-abṣār. Ittaqū al-Raḥmāna yā malaʾa al-imkān. An
anṣifū fī hādhā al-amri alladhī idhā ẓahara nufikha fī al-ṣūri
wa-marrati al-jibāl. Qul faʾtū yā malaʾa al-arḍi bi-mā ʿindakum wa-nazin
mā ʿindanā wa-ʿindakum bi-hādhā al-mīzāni alladhī wuḍiʿa bi-al-ʿadli
amran min ladā Allāhi rabbi al-arbāb. Hādhā yawmun fīhi yunādī al-mīzānu
bi-aʿlā al-nidāʾi bi-annī anā al-mumayyizu al-ʿalīmu min ladun māliki
al-adyān.

Say: O peoples! By God, the Book speaks at the pivot of the world and
summons the nations to the Lord of eternity, yet the people are sunk in
an astonishing sleep. Say: He has come before whom all the books of the
earth have bowed. Take heed, you who have insight! Be mindful of the
All-Merciful, O inhabitants of the realm of existence! Judge fairly this
Cause, at whose appearance the trumpet was sounded and the mountains
passed away. Say: O peoples of the earth, bring what you possess, and We
will weigh what We possess and what you possess in this Balance, which
has been established with justice at the command of God, the Lord of
lords. This is a day when the Balance calls out at the top of its voice,
``I am indeed the discerning, the all-knowing, from the Lord of
religions.''

\subsection{102}\label{section-2522}

\begin{otherlanguage}{arabic}

و نذکر من سمّی بحیدر من هذا المنظر الأکبر لیبشّره برحمتی و عنایتی و یقرّبه
الی مقام جعله اللّه مطاف الفردوس الأعلی و اهل ملکوت الأسمآء فی العشیّ و
الاشراق لعمری لو تطّلع بما عند ربّک لیأخذک الفرح و الابتهاج علی شأن تعجز
عن ذکره الأقلام کن قائماً علی خدمة امری و ناطقاً بثنائی و ناظراً الی افقی و
متشبّثاً بأذیال ردآء عنایتی کذلک امرناک من قبل و فی هذا اللّوح الّذی لاحت من
افقه شمس الحکمة و البیان یا حیدر وصّ عبادی بأخلاقی و ما نزّل فی کتابی انّ
ربّک لهو الآمر العزیز العلّام ایّاک ان تمنعک شئونات الخلق عن الحقّ ضع ما
سوائی ثمّ ولّ وجهک الی شطری و قل اشهد انّک کنت مسطوراً من القلم الأعلی و
مکنوناً فی علم اللّه مالک الرّقاب ان احمد اللّه بما ذکرت لدی المظلوم و انزل
لک ما انارت به الآفاق

\end{otherlanguage}

Wa-nadhkuru man summiya bi-Ḥaydar min hādhā al-manẓari al-akbari
li-yubashshirahu bi-raḥmatī wa-ʿināyatī wa-yuqarribahu ilā maqāmin
jaʿalahu Allāhu maṭāfa al-firdawsi al-aʿlā wa-ahli malakūti al-asmāʾi fī
al-ʿashiyyi wa-al-ishrāq. La-ʿamrī law taṭṭaliʿu bi-mā ʿinda rabbika
la-yaʾkhudhuka al-faraḥu wa-al-ibtihāju ʿalā shaʾnin taʿjizu ʿan
dhikrihi al-aqlām. Kun qāʾiman ʿalā khidmati amrī wa-nāṭiqan bi-thanāʾī
wa-nāẓiran ilā ufuqī wa-mutashabbithan bi-adhyāli ridāʾi ʿināyatī.
Kadhālika amarnāka min qablu wa-fī hādhā al-lawḥi alladhī lāḥat min
ufuqihi shamsu al-ḥikmati wa-al-bayān. Yā Ḥaydar, waṣṣi ʿibādī
bi-akhlāqī wa-mā nuzzila fī kitābī. Inna rabbaka la-huwa al-āmiru
al-ʿazīzu al-ʿallām. Iyyāka an tamnaʿaka shuʾūnātu al-khalqi ʿani
al-Ḥaqq. Ḍaʿ mā siwāʾī thumma walli wajhaka ilā shaṭrī wa-qul ashhadu
annaka kunta masṭūran mina al-qalami al-aʿlā wa-maknūnan fī ʿilmi Allāhi
māliki al-riqāb. An iḥmadi Allāha bi-mā dhukirta ladā al-maẓlūmi
wa-unzila laka mā anārat bihi al-āfāq.

We also remember him who is named Ḥaydar from this Most Great Outlook,
that it may gladden him with My mercy and loving care and draw him near
to a station that God has made the place around which the highest
Paradise and the inhabitants of the Kingdom of Names circle at evening
and at dawn. By My life! Were you to learn what is with your Lord, joy
and gladness would seize you in a manner that pens are powerless to
describe. Arise to serve My Cause, speak My praise, look toward My
horizon, and cling to the hem of the robe of My loving care. Thus have
We commanded you before and in this Tablet, from whose horizon the sun
of wisdom and utterance has shone. O Ḥaydar! Exhort My servants to adopt
My character and follow what has been sent down in My Book. Your Lord is
indeed the Ordainer, the Almighty, the All-Knowing. Beware lest the
affairs of creation keep you from the True One. Set aside all but Me,
then turn your face toward Me and say, ``I bear witness that You were
inscribed by the Most Exalted Pen and concealed in the knowledge of God,
the Master of all.'' Praise God that you were mentioned in the presence
of the Wronged One and that there was sent down for you that through
which the horizons were illumined.

\subsection{103}\label{section-2523}

\begin{otherlanguage}{arabic}

یا علی قبل اکبر علیک بهائی و رحمتی ان افرح بما توجّه الیک البحر الأعظم من
هذا المقام الأنور الأقدم و اراد ان یذکر عمّک و ابنه لیفرحا بذکر اللّه فی
یوم فیه اسودّت الوجوه و زلّت الأقدام طوبی لمقبل سمع و اقبل ویل لمعرض اعرض
عن الّذی ذکره فی اللّیالی و الأیّام انّا نبشّرهما و نکبّر من هذا المقام علیهما
لیقوما علی خدمة امری الّذی خضعت له الأعناق و نذکر احبّآء اللّه و اولیائه فی
ارض السّین و النّون و نبشّرهم بالشّمس و اشراقها و البحر و امواجه و السّمآء و
ارتفاعها و الأنجم و انوارها و ما رقم لهم من قلمی الأعلی فی الصّحیفة
الحمرآء کذلک یجزی اللّه عباده فضلاً من عنده و هو العزیز الفضّال

\end{otherlanguage}

Yā ʿAlī qabla Akbar, ʿalayka bahāʾī wa-raḥmatī. An ifraḥ bi-mā tawajjaha
ilayka al-baḥru al-aʿẓamu min hādhā al-maqāmi al-anwari al-aqdami
wa-arāda an yadhkura ʿammaka wa-ibnahu li-yafraḥā bi-dhikri Allāhi fī
yawmin fīhi iswaddati al-wujūhu wa-zallati al-aqdām. Ṭūbā li-muqbilin
samiʿa wa-aqbala, waylun li-muʿriḍin aʿraḍa ʿani alladhī dhakarahu fī
al-layālī wa-al-ayyām. Innā nubashshiruhumā wa-nukabbiru min hādhā
al-maqāmi ʿalayhimā li-yaqūmā ʿalā khidmati amrī alladhī khaḍaʿat lahu
al-aʿnāq. Wa-nadhkuru aḥibbāʾa Allāhi wa-awliyāʾahu fī arḍi al-sīni
wa-al-nūni wa-nubashshiruhum bi-al-shamsi wa-ishrāqihā wa-al-baḥri
wa-amwājihi wa-al-samāʾi wa-irtifāʿihā wa-al-anjumi wa-anwārihā wa-mā
ruqima lahum min qalamī al-aʿlā fī al-ṣaḥīfati al-ḥamrāʾ. Kadhālika
yajzī Allāhu ʿibādahu faḍlan min ʿindihi wa-huwa al-ʿazīzu al-faḍḍāl.

O ʿAlī, whose name precedes Akbar! Upon you be My glory and My mercy.
Rejoice that the Most Great Ocean has turned toward you from this most
luminous, most ancient station and has wished to remember your paternal
uncle and his son, that they may rejoice in the remembrance of God on a
day when faces have darkened and feet have slipped. Blessed is the
receptive one who has heard and turned toward Him; woe to the one who
has turned away from Him whom he remembered by night and day! We give
them both glad tidings and bestow upon them Our greeting of glory from
this station, that they may arise to serve My Cause, before which all
necks have bowed. We also remember God's loved ones and His friends in
the lands of Sīn and Nūn, and gladden them with the sun and its
radiance, the ocean and its waves, the heaven and its loftiness, the
stars and their lights, and what has been inscribed for them by My Most
Exalted Pen in the Crimson Scroll. Thus does God reward His servants as
a grace from Him, and He is the Almighty, the All-Bountiful.

\subsection{104}\label{section-2524}

\begin{otherlanguage}{persian}

یا یحیی ان استمع النّدآء من شطر عکّآء عن یمین البقعة النّورآء من السّدرة انّه
لا اله الّا انا الغنیّ المتعال هل تعرف من ینطقک و هل تعلم من اقبل الیک قل
ای ونفسک انّ مکلّم الطّور ینطق لی و مالک الظّهور قد اقبل الی وجهی فضلاً من
عنده و هو الغنیّ بالحقّ و انا الفقیر المحتاج یا یحیی وصّ عبادی بالأمانة و
الصّدق و بما یرتفع به ذکر اللّه فی المدن و الدّیار ان افتحوا القلوب بجنود
الأخلاق کذلک امرنا العباد من قبل و فی هذا الکتاب و کذلک زیّنّا دیباج کتاب
الوجود بذکر المقصود لتشکر ربّک مالک الأنام

\end{otherlanguage}

Yā Yaḥyā, an istamiʿ al-nidāʾa min shaṭri ʿAkkāʾa ʿan yamīni al-buqʿati
al-nūrāʾi mina al-sidrati, innahu lā ilāha illā anā al-ghaniyyu
al-mutaʿāl. Hal taʿrifu man yunṭiquka wa-hal taʿlamu man aqbala ilayka?
Qul ay wa-nafsika inna mukallima al-Ṭūri yanṭiqu lī wa-mālika al-ẓuhūri
qad aqbala ilā wajhī faḍlan min ʿindihi wa-huwa al-ghaniyyu bi-al-ḥaqqi
wa-anā al-faqīru al-muḥtāj. Yā Yaḥyā, waṣṣi ʿibādī bi-al-amānati
wa-al-ṣidqi wa-bi-mā yartafiʿu bihi dhikru Allāhi fī al-muduni
wa-al-diyār. An iftaḥū al-qulūba bi-junūdi al-akhlāq. Kadhālika amarnā
al-ʿibāda min qablu wa-fī hādhā al-kitāb. Wa-kadhālika zayyannā dībāja
kitābi al-wujūdi bi-dhikri al-maqṣūdi li-tashkura rabbaka mālika
al-anām.

O Yaḥyā! Listen to the call from the direction of ʿAkkā, from the Tree
on the right side of the luminous Spot: ``There is indeed no God but Me,
the Self-Sufficient, the Most Exalted.'' Do you recognize Him who causes
you to speak, and do you know who has turned toward you? Say: ``Yes, by
Your own Self! He who spoke on Sinai speaks to me, and the Lord of
Revelation has turned His face toward me as a grace from Him. He is in
truth the Self-Sufficient, while I am poor and needy.'' O Yaḥyā! Exhort
My servants to trustworthiness and truthfulness and to whatever will
exalt the remembrance of God in cities and lands. Conquer hearts with
the hosts of good character. Thus have We commanded the servants before
and in this Book. Thus have We adorned the brocade of the Book of
Existence with the remembrance of the Desired One, that you may thank
your Lord, the Lord of humankind.

\subsection{105}\label{section-2525}

\begin{otherlanguage}{arabic}

و نذکر من سمّی بمولی فیهذا المقام الأعلی بندائی الأحلی و انا المقتدر علی
ما اراد قم علی الذّکر و الثّنآء و سبّح بحمد ربّک مالک المآب انّا ذکرناکم بذکر
سالت به البطحآء و جری فرات رحمتی من کلّ الجهات قل یا قوم لا تمنعوا انفسکم
عن بحر الحیوان اتّقوا الرّحمن الّذی اتی بجنود الوحی و رایات الآیات قل هل
یبقی لکم ما ترونه الیوم لا ونفسی الحقّ سیفنی من علی الأرض و یبقی الملک
لنفسه وحدها یشهد بذلک کلّ شجر و حجر و مدر و حصاة

\end{otherlanguage}

Wa-nadhkuru man summiya bi-Mawlā fī hādhā al-maqāmi al-aʿlā bi-nidāʾī
al-aḥlā wa-anā al-muqtadiru ʿalā mā arād. Qum ʿalā al-dhikri
wa-al-thanāʾi wa-sabbiḥ bi-ḥamdi rabbika māliki al-maʾāb. Innā
dhakarnākum bi-dhikrin sālat bihi al-baṭḥāʾu wa-jarā furātu raḥmatī min
kulli al-jihāt. Qul yā qawmi lā tamnaʿū anfusakum ʿan baḥri al-ḥayawān.
Ittaqū al-Raḥmāna alladhī atā bi-junūdi al-waḥyi wa-rāyāti al-āyāt. Qul
hal yabqā lakum mā tarawnahu al-yawm? Lā wa-nafsī al-Ḥaqq. Sa-yafnā man
ʿalā al-arḍi wa-yabqā al-mulku li-nafsihi waḥdahā. Yashhadu bi-dhālika
kullu shajarin wa-ḥajarin wa-madarin wa-ḥaṣāt.

We also remember him who is named Mawlā in this most exalted station
with My sweetest call, and I have power over whatever I will. Arise in
remembrance and praise, and glorify your Lord, the Lord of the place of
return. We have remembered you with a remembrance through which the
valley has flowed and the sweet waters of My mercy have streamed from
every direction. Say: O people! Do not deprive yourselves of the ocean
of life. Be mindful of the All-Merciful, who has come with the hosts of
revelation and the banners of verses. Say: Will what you see today
remain yours? No, by My own Self, the True One! Those on earth will
perish, and sovereignty will remain His alone. To this every tree,
stone, clod, and pebble bears witness.

\subsection{106}\label{section-2526}

\begin{otherlanguage}{arabic}

و نذکر الهادی الّذی هدیناه الی سوآء الصّراط یا هادی انّ المظلوم یذکرک و
یأمرک و الّذین آمنوا بالمعروف ان اتّبعوا ما امرتم به من لدن فالق الأصباح
تاللّه یوم اللّه ینطق بالحقّ فی هذا القصر الّذی جعله اللّه المنظر الأکبر ان
اعرفوا یا اولی الألباب قل ضعوا ما عند القوم قد اتی سلطان الوجود بأعلام
الوحی و الالهام یا ملأ الأرض اتّقوا اللّه و لا تحرموا انفسکم عن الّذی به
ظهر ما کان مکنوناً فی الزّبر و الألواح خذوا کتاب اللّه بقوّة و لا تتّبعوا کلّ
جاهل مرتاب کذلک نوّرنا افق سمآء العرفان بنور البیان طوبی لمن سمع و رأی و
ویل لکلّ غافل مکّار

\end{otherlanguage}

Wa-nadhkuru al-Hādī alladhī hadaynāhu ilā sawāʾi al-ṣirāṭ. Yā Hādī, inna
al-maẓlūma yadhkuruka wa-yaʾmuruka wa-alladhīna āmanū bi-al-maʿrūf. An
ittabiʿū mā umirtum bihi min ladun fāliqi al-aṣbāḥ. Ta-Allāhi yawmu
Allāhi yanṭiqu bi-al-ḥaqqi fī hādhā al-qaṣri alladhī jaʿalahu Allāhu
al-manẓara al-akbar. An iʿrifū yā ulī al-albāb. Qul ḍaʿū mā ʿinda
al-qawmi qad atā sulṭānu al-wujūdi bi-aʿlāmi al-waḥyi wa-al-ilhām. Yā
malaʾa al-arḍi ittaqū Allāha wa-lā taḥrimū anfusakum ʿani alladhī bihi
ẓahara mā kāna maknūnan fī al-zuburi wa-al-alwāḥ. Khudhū kitāba Allāhi
bi-quwwatin wa-lā tattabiʿū kulla jāhilin murtāb. Kadhālika nawwarnā
ufuqa samāʾi al-ʿirfāni bi-nūri al-bayān. Ṭūbā li-man samiʿa wa-raʾā
wa-waylun li-kulli ghāfilin makkār.

We also remember Hādī, whom We guided to the straight path. O Hādī! The
Wronged One remembers you and enjoins you and those who have believed to
do what is right. Follow what you have been commanded by the Cleaver of
the dawn. By God! The Day of God speaks the truth in this palace, which
God has made the Most Great Outlook. Recognize this, you who possess
understanding! Say: Set aside what the people possess; the Sovereign of
existence has come with the banners of revelation and inspiration. O
peoples of the earth! Be mindful of God, and do not deprive yourselves
of Him through whom what was concealed in the Books and Tablets has been
revealed. Take hold of the Book of God with strength, and do not follow
any ignorant doubter. Thus have We illumined the horizon of the heaven
of recognition with the light of utterance. Blessed is whoever has heard
and seen, and woe to every heedless schemer!

\subsection{107}\label{section-2527}

\begin{otherlanguage}{arabic}

ان یا قلمی الأعلی ان اذکر من سمع صریرک و اقبل الیک الّذی سمّی بیوسف لیفرح
بهذا الذّکر الّذی اذا ظهر سجدت له الأذکار انّا ظهرنا و اظهرنا الأمر علی شأن
اضطربت القلوب و شاخصت الأبصار الّا الّذین نبذوا الأوهام عن ورائهم مقبلین
الی اللّه مطلع الأنوار طوبی لمن فاز باصغآء ندائی و وجد عرف قمیصی و انتبه
من نسمات عنایتی اذ سرت فی الأسحار هذا یوم فیه انجذبت الأشیآء من ندآء
مالک الأسمآء ولکنّ النّاس فی غفلة و نعاس نعیماً لمن وجد عرف البیان و قام
علی خدمة الأمر فی الغدوّ و الآصال انّا سمعنا ذکرک ذکرناک و رأینا اقبالک
اقبلنا الیک من مشرق الأذکار

\end{otherlanguage}

In yā qalamī al-aʿlā in udhkur man samiʿa ṣarīraka wa-aqbala ilayka
alladhī summiya bi-Yūsuf li-yafraḥa bi-hādhā al-dhikr alladhī idhā
ẓahara sajadat lahu al-adhkār. Innā ẓaharnā wa-aẓharnā al-amr ʿalā shaʾn
iḍṭarabat al-qulūb wa-shākhiṣat al-abṣār illā alladhīna nabadhū al-awhām
ʿan warāʾihim muqbilīna ilā Allāh maṭlaʿ al-anwār. Ṭūbā li-man fāza
bi-iṣghāʾ nidāʾī wa-wajada ʿarf qamīṣī wa-intabaha min nasamāt ʿināyatī
idh sarat fī al-asḥār. Hādhā yawm fīhi injadhabat al-ashyāʾ min nidāʾ
mālik al-asmāʾ wa-lākinna al-nās fī ghafla wa-nuʿās. Naʿīman li-man
wajada ʿarf al-bayān wa-qāma ʿalā khidmat al-amr fī al-ghuduww
wa-al-āṣāl. Innā samiʿnā dhikraka dhakarnāka wa-raʾaynā iqbālaka aqbalnā
ilayka min mashriq al-adhkār.

O My Most Exalted Pen! Remember the one who heard your rustling and
turned toward you, who is named Yūsuf, so that he may rejoice in this
remembrance, before which all remembrances bowed down when it appeared.
We have appeared and manifested the Cause in such a way that hearts were
shaken and eyes stared in astonishment, except those who cast vain
imaginings behind them and turned toward God, the Dawning-Place of
lights. Blessed is whoever has attained to hearing My call, perceived
the fragrance of My robe, and awakened through the breezes of My loving
care as they wafted at dawn. This is a day when all things have been
enraptured by the call of the Lord of Names, yet the people remain
heedless and asleep. Blessed is whoever has perceived the fragrance of
utterance and arisen to serve the Cause morning and evening. We heard
your remembrance and remembered you; We saw your turning toward Us and
turned toward you from the Dawning-Place of remembrances.

\subsection{108}\label{section-2528}

\begin{otherlanguage}{arabic}

و نذکر علی قبل محمّد و نبشّره و الّذین آمنوا فیهناک برحمتی و عنایتی و فضلی
الّذی به شهدت الذّرّات توکّل فی کلّ الأمور علی مکلّم الطّور و تمسّک بما نزّل فی
الکتاب من لدن مرسل الأریاح ایّاک ان تمنعک اوهام العلمآء عن فاطر السّمآء او
تخوّفک سطوة الأمرآء عن مطلع الایقان لا تلتفتوا الی ایّام معدودات انّها
ستمضی و یرون الموحّدون انفسهم فی اعلی الجنان انّا نأمرکم فیکلّ الأحوال
بالحکمة لئلّا یرتفع ضوضآء الّذین کفروا بالمآل کذلک تضوّع عرف البیان فی
الامکان طوبی لمن وجد و ویل لکلّ منکر مکّار

\end{otherlanguage}

Wa-nadhkuru ʿAlī qabla Muḥammad wa-nubashshiruhu wa-alladhīna āmanū fī
hunāka bi-raḥmatī wa-ʿināyatī wa-faḍlī alladhī bihi shahidat al-dharrāt.
Tawakkal fī kull al-umūr ʿalā mukallim al-Ṭūr wa-tamassak bi-mā nuzzila
fī al-kitāb min ladun mursil al-aryāḥ. Iyyāka an tamnaʿaka awhām
al-ʿulamāʾ ʿan fāṭir al-samāʾ aw tukhawwifaka saṭwat al-umarāʾ ʿan
maṭlaʿ al-īqān. Lā taltafitū ilā ayyām maʿdūdāt innahā sa-tamḍī
wa-yarawna al-muwaḥḥidūn anfusahum fī aʿlā al-jinān. Innā naʾmurukum fī
kull al-aḥwāl bi-al-ḥikma li-allā yartafiʿa ḍawḍāʾ alladhīna kafarū
bi-al-maʾāl. Kadhālika taḍawwaʿa ʿarf al-bayān fī al-imkān. Ṭūbā li-man
wajada wa-wayl li-kull munkir makkār.

We also remember ʿAlī before Muḥammad, and give him and those who have
believed in that place the good news of My mercy, My loving care, and My
grace, to which the atoms have borne witness. In all matters, place your
trust in Him who spoke on Sinai, and hold fast to what has been sent
down in the Book by the Sender of the winds. Beware lest the vain
imaginings of the learned keep you from the Creator of heaven, or the
might of rulers frighten you away from the Dawning-Place of certitude.
Pay no heed to these few days: they will pass, and those who affirm
God's unity will see themselves in the loftiest gardens of paradise. We
enjoin you to exercise wisdom under all circumstances, lest the clamor
of those who have denied the ultimate return arise. Thus has the
fragrance of utterance spread throughout existence. Blessed is whoever
has perceived it, and woe to every scheming denier!

\subsection{109}\label{section-2529}

\begin{otherlanguage}{arabic}

و نذکر احبّائی فی الشّین و الهآء و نبشّرهم باقبالی الیهم و نکبّر علی وجوههم
من هذا المقام الّذی تزیّن بأنوار وجه ربّهم العزیز الغفّار یا احبّائی ان
افرحوا بما نزّل لکم من سمآء الأمر ثمّ اشکروا ربّکم الرّحمن فیکلّ الأحیان انّه
غفرکم و هداکم و فضّلکم علی اکثر العباد انّا ذکرنا کلّ من ذکر اسمه تلقآء
الوجه ثمّ الّذین ما ذکرت اسمائهم انّه یفعل کیف یشآء بسلطان من عنده و هو
المقتدر العزیز العلّام

\end{otherlanguage}

Wa-nadhkuru aḥibbāʾī fī al-Shīn wa-al-Hāʾ wa-nubashshiruhum bi-iqbālī
ilayhim wa-nukabbiru ʿalā wujūhihim min hādhā al-maqām alladhī tazayyana
bi-anwār wajh rabbihim al-ʿAzīz al-Ghaffār. Yā aḥibbāʾī in ifraḥū bi-mā
nuzzila lakum min samāʾ al-amr thumma ushkurū rabbakum al-Raḥmān fī kull
al-aḥyān. Innahu ghafarakum wa-hadākum wa-faḍḍalakum ʿalā akthar
al-ʿibād. Innā dhakarnā kulla man dhukira ismuhu tilqāʾ al-wajh thumma
alladhīna mā dhukirat asmāʾuhum. Innahu yafʿalu kayfa yashāʾ bi-sulṭān
min ʿindihi wa-huwa al-Muqtadir al-ʿAzīz al-ʿAllām.

We also remember My loved ones in Shīn and Hāʾ, give them the good news
that I have turned toward them, and greet their faces with words of
glorification from this station, which is adorned with the lights of the
face of their Lord, the Almighty, the Ever-Forgiving. O My loved ones!
Rejoice in what has been sent down for you from the heaven of the Cause,
and thank your Lord, the All-Merciful, at all times. He has forgiven
you, guided you, and favored you above most of His servants. We have
remembered everyone whose name was mentioned before Our face, as well as
those whose names were not mentioned. He does as He wills through a
sovereignty of His own, and He is the All-Powerful, the Almighty, the
All-Knowing.

\subsection{110}\label{section-2530}

\begin{otherlanguage}{arabic}

و نذکر من سمّی بخانلر و نبشّره بعنایاتی الّتی لا یأخذها النّفاد ان افرح بما
ذکرناک و الّذین آمنوا بآیات تفدی لها الأرواح ان اقتصر امورک علی ذکر اللّه
و ثنائه اذا ورد علیک احد من اولیائه قم علی خدمته بالرّوح و الرّیحان ایّاک
ان تحزنک شئونات الدّنیا و ما یرد فیها دعها عن ورائک و تمسّک بحبل اللّه مالک
الایجاد البهآء علیک و علی ابنائک و اهلک و علی کلّ موقن صبّار انتهی

\end{otherlanguage}

Wa-nadhkuru man summiya bi-Khānlar wa-nubashshiruhu bi-ʿināyātī allatī
lā yaʾkhudhuhā al-nafād. In ifraḥ bi-mā dhakarnāka wa-alladhīna āmanū
bi-āyāt tufdā lahā al-arwāḥ. In iqtaṣir umūraka ʿalā dhikr Allāh
wa-thanāʾihi. Idhā warada ʿalayka aḥad min awliyāʾihi qum ʿalā
khidmatihi bi-al-rūḥ wa-al-rayḥān. Iyyāka an tuḥzinaka shuʾūnāt al-dunyā
wa-mā yaridu fīhā. Daʿhā ʿan warāʾika wa-tamassak bi-ḥabl Allāh mālik
al-ījād. Al-bahāʾ ʿalayka wa-ʿalā abnāʾika wa-ahlika wa-ʿalā kull mūqin
ṣabbār. Intahā.

We also remember the one named Khānlar and give him the good news of My
inexhaustible favors. Rejoice that We have remembered you and those who
have believed in verses for which souls would be sacrificed. Confine
your affairs to the remembrance and praise of God. Whenever one of His
loved ones comes to you, arise to serve him with joy and gladness.
Beware lest the affairs of this world and what happens in it grieve you.
Leave it behind you and hold fast to the cord of God, the Lord of
creation. Glory rest upon you, your sons, your family, and everyone who
is assured and steadfastly patient. End.

\subsection{111}\label{section-2531}

\begin{otherlanguage}{arabic}

آیا بعد از امواج بحر معانی و ظهورات انوار بیان الهی اینعبد چه ذکر نماید
و چه عرض کند وقتی از اوقات اینکلمهٴ علیا از لسان مالک اسماء ظاهر
فی‌الحقیقه قلب از آن محترق و زفرات متصاعد و عبرات نازل قوله جلّت عظمته

\end{otherlanguage}

Āyā baʿd az amvāj-i baḥr-i maʿānī va ẓuhūrāt-i anvār-i bayān-i ilāhī īn
ʿabd che dhikr namāyad va che ʿarż konad? Vaqtī az awqāt īn kalimiy-i
ʿulyā az lisān-i Mālik-i Asmāʾ ẓāhir, fī al-ḥaqīqa qalb az ān moḥtareq
va zafarāt motaṣāʿid va ʿabarāt nāzil. Qawluhu jallat ʿaẓamatuhu.

After the waves of the ocean of meanings and the manifestations of the
lights of divine utterance, what can this servant mention, and what can
he present? On one occasion, these exalted words issued from the tongue
of the Lord of Names. Truly, the heart is set aflame by them, sighs
ascend, and tears fall. His words---exalted be His grandeur:

\subsection{112}\label{section-2532}

\begin{otherlanguage}{persian}

یا عبد حاضر مشاهده مینمائی و اصغا میکنی که در لیالی و ایّام صریر قلم فضل
مرتفع و ندای مظلوم بلند و جمیع از برای هدایت عباد و نجات من علی الأرض
بوده ای کاش در جزای اینموهبت کبری و شفقت عظمی بر امر الهی متّحد می‌شدند و
بآنچه سبب و علّت سموّ و علوّ اوست متشبّث که شاید عرف خوشی از اعمال خالصه و
اخلاق روحانیّه ساطع شود و عالمرا اخبار نماید بآنچه ظاهر شده این عرف اگر
مرتفع شود قاصدیست صادق و رسولیست کامل و مبیّن طوبی لمن فاز ببیانی و عرف
ما اراد ملکوت مشیّتی انتهی

\end{otherlanguage}

Yā ʿabd-i ḥāżir, moshāhede mī-namāʾī va iṣghā mī-konī ke dar layālī va
ayyām ṣarīr-i qalam-i fażl mortafeʿ va nidā-yi Maẓlūm boland va jamīʿ az
barā-yi hidāyat-i ʿibād va najāt-i man ʿalā al-arḍ būde. Ey kāsh dar
jazā-yi īn mawhebat-i kubrā va shafaqat-i ʿuẓmā bar amr-i ilāhī mottaḥid
mī-shodand va bi-ānche sabab va ʿellat-i somovv va ʿolovv-i ūst
motashabbith ke shāyad ʿarf-i khoshī az aʿmāl-i khāliṣe va akhlāq-i
rūḥāniyye sāṭeʿ shavad va ʿālam-rā ikhbār namāyad bi-ānche ẓāhir shode.
Īn ʿarf agar mortafeʿ shavad qāṣedīst ṣādiq va rasūlīst kāmil va
mobayyin. Ṭūbā li-man fāza bi-bayānī wa-ʿarafa mā arāda malakūt
mashīʾatī. Intahā.

``O servant in attendance! You see and hear how, night and day, the
rustling of the Pen of grace rises and the call of the Wronged One rings
out, all for the guidance of the servants and the salvation of those on
earth. If only, in return for this supreme gift and this immense
compassion, they would unite in the Cause of God and hold fast to
whatever brings about its elevation and exaltation, so that a sweet
fragrance might spread from pure deeds and spiritual qualities and
acquaint the world with what has appeared. If this fragrance rises, it
is a truthful envoy and a perfect messenger who makes things clear.
Blessed is whoever has attained to My utterance and recognized what the
kingdom of My will has intended.'' End.

\subsection{113}\label{section-2533}

\begin{otherlanguage}{persian}

این فقیر لازال در آنچه ظاهر شده متحیّر و مبهوت بیک طنین عنکبوت اقبال
مینمایند و از صریر قلم اعلی اعراض رایحهٴ کریهه را محبوب میشمرند و از عرف
فردوس اعلی محروم جمیع امور ظاهر و مشهود سوف یراه حضرتک کذلک اخبرنی من
عنده امّ الکتاب

\end{otherlanguage}

Īn faqīr lā-zāl dar ānche ẓāhir shode motaḥayyer va mabhūt. Bi-yek
ṭanīn-i ʿankabūt iqbāl mī-namāyand va az ṣarīr-i Qalam-i Aʿlā iʿrāż.
Rāyiḥiy-i karīhe-rā maḥbūb mī-shomarand va az ʿarf-i Firdaws-i Aʿlā
maḥrūm. Jamīʿ-i omūr ẓāhir va mashhūd. Sawfa yarāhu ḥaḍratuka kadhālika
akhbaranī man ʿindahu umm al-kitāb.

This poor one remains bewildered and astonished at what has appeared.
They turn toward the buzzing of a spider and turn away from the rustling
of the Most Exalted Pen. They cherish a foul odor and are deprived of
the fragrance of the most exalted Paradise. All matters are manifest and
evident. You will see it; thus has He who possesses the Mother Book
informed me.

\subsection{114}\label{section-2534}

\begin{otherlanguage}{persian}

در این حین دستخطّ ثالث آنحضرت بمثابهٴ نسیم ربیع حقیقی بر این جسد یابس
مطروح عبور نمود بخششها فرمود و عنایتها اظهار داشت للّه الحمد بشیر محبّت در
ظهور و مرور و انشآءاللّه این بشیر ممنوع نشود و از حرکت باز نماند

\end{otherlanguage}

Dar īn ḥīn dast-khaṭṭ-i thālith-i ān ḥażrat bi-mathābiy-i nasīm-i
rabīʿ-i ḥaqīqī bar īn jasad-i yābis-i maṭrūḥ ʿobūr namūd, bakhshesh-hā
farmūd va ʿināyat-hā iẓhār dāsht. Li-llāh al-ḥamd, bashīr-i maḥabbat dar
ẓuhūr va morūr va in shāʾ Allāh īn bashīr mamnūʿ nashavad va az ḥarakat
bāz namānad.

At this moment, your third handwritten letter passed over this dry,
prostrate body like the breeze of the true spring, bestowing gifts and
showing kindness. Praise be to God, the herald of love continues to
appear and pass by. God willing, this herald will not be hindered or
cease to move.

\subsection{115}\label{section-2535}

\begin{otherlanguage}{persian}

و بعد از قرائت و اطّلاع بمقام لا عدل له توجّه نموده بعد از حضور و اذن
مناجات آنمحبوب بکمال روح و ریحان تلقاء عرش رحمن عرض شد هذا ما نطق به
لسان القدم و مولی العالم قوله تبارک و تعالی

\end{otherlanguage}

Va baʿd az qirāʾat va iṭṭilāʿ bi-maqām-i lā ʿidla lahu tavajjoh namūde,
baʿd az ḥuḍūr va idhn, monājāt-i ān maḥbūb bi-kamāl-i rūḥ va rayḥān
tilqāʾ-i ʿarsh-i Raḥmān ʿarż shod. Hādhā mā naṭaqa bihi lisān al-qidam
wa-mawlā al-ʿālam. Qawluhu tabāraka wa-taʿālā.

After reading it and acquainting myself with its contents, I turned
toward the peerless station. After entering the presence and receiving
permission, the prayer of that dear friend was presented with the utmost
joy and gladness before the throne of the All-Merciful. This is what the
Tongue of Eternity and Lord of the world spoke. His words---blessed and
exalted is He:

\subsection{116}\label{section-2536}

\begin{otherlanguage}{arabic}

یا ایّها المتوجّه الی وجهی و المنجذب بآیاتی انّا سمعنا ندائک و ما ناجیت به
اللّه ربّک و سمعنا زفراتک و رأینا عبراتک فی حبّ اللّه محبوبک و مقصودک و
شاهدنا نار حبّک و اشتعال قلبک و خضوعک و خشوعک و اهتزاز ارکانک فی عشق اللّه
و ودّه انّا کنّا معک فی المجالس و المحافل اذ کنت ناطقاً بهذا الاسم الأعظم و
هذا النّبأ الّذی بشّرنا به النّبیّین و المرسلین نعیماً لک و للّذین ما منعتهم
شئونات الأوهام قاموا و قالوا اللّه ربّنا و ربّ العرش العظیم کبّر من قبلی علی
احبّائی قل تاللّه ینوح قلبی و تبکی عینی بما ورد علی امری العزیز البدیع
تمسّکوا بحبل الاتّحاد و ما فات عنکم فی ایّام اللّه ربّ العالمین کذلک اوقدنا
سراج النّصح بدهن الفضل و نار الحبّ ان لن تحفظوه عن الأریاح لا تکونوا سبباً
لاطفائه اتّقوا الرّحمن و تفکّروا فی عنایة ربّکم الغفور الرّحیم قوموا علی خدمة
الأمر انّه خیر لکم عمّا خلق فی السّموات و الأرضین انتهی

\end{otherlanguage}

Yā ayyuhā al-mutawajjih ilā wajhī wa-al-munjadhib bi-āyātī, innā samiʿnā
nidāʾaka wa-mā nājayta bihi Allāh rabbaka wa-samiʿnā zafarātaka
wa-raʾaynā ʿabarātaka fī ḥubb Allāh maḥbūbika wa-maqṣūdika wa-shāhadnā
nār ḥubbika wa-ishtiʿāl qalbika wa-khuḍūʿaka wa-khushūʿaka wa-ihtizāz
arkānika fī ʿishq Allāh wa-wuddihi. Innā kunnā maʿaka fī al-majālis
wa-al-maḥāfil idh kunta nāṭiqan bi-hādhā al-ism al-aʿẓam wa-hādhā
al-nabaʾ alladhī bashsharnā bihi al-nabiyyīn wa-al-mursalīn. Naʿīman
laka wa-li-alladhīna mā manaʿathum shuʾūnāt al-awhām qāmū wa-qālū Allāh
rabbunā wa-rabb al-ʿarsh al-ʿaẓīm. Kabbir min qibalī ʿalā aḥibbāʾī. Qul
ta-Allāh yanūḥu qalbī wa-tabkī ʿaynī bi-mā warada ʿalā amrī al-ʿazīz
al-badīʿ. Tamassakū bi-ḥabl al-ittiḥād wa-mā fāta ʿankum fī ayyām Allāh
rabb al-ʿālamīn. Kadhālika awqadnā sirāj al-nuṣḥ bi-duhn al-faḍl wa-nār
al-ḥubb. In lan taḥfaẓūhu ʿan al-aryāḥ lā takūnū sababan li-iṭfāʾihi.
Ittaqū al-Raḥmān wa-tafakkarū fī ʿināyat rabbikum al-Ghafūr al-Raḥīm.
Qūmū ʿalā khidmat al-amr innahu khayr lakum ʿammā khuliqa fī al-samāwāt
wa-al-arḍīn. Intahā.

``O you who turn toward My face and are enraptured by My verses! We have
heard your call and what you whispered in prayer to God, your Lord. We
have heard your sighs and seen your tears in your love for God, your
Beloved and your Goal. We have witnessed the fire of your love, the
burning of your heart, your submissiveness and humility, and the
trembling of your limbs in love and affection for God. We were with you
in gatherings and assemblies when you spoke of this Most Great Name and
this Announcement of which We gave the prophets and messengers glad
tidings. Blessed are you and those whom vain imaginings have not held
back, who arose and said, `God is our Lord and the Lord of the mighty
throne.' Greet My loved ones on My behalf with words of glorification.
Say: By God! My heart laments and My eyes weep at what has befallen My
mighty and wondrous Cause. Hold fast to the cord of unity and to what
you have missed in the days of God, the Lord of the worlds. Thus have We
kindled the lamp of counsel with the oil of grace and the fire of love.
If you will not protect it from the winds, do not be the cause of its
extinction. Be mindful of the All-Merciful, and reflect on the loving
care of your Lord, the Ever-Forgiving, the Merciful. Arise to serve the
Cause: it is better for you than all that has been created in the
heavens and the earths.'' End.

\subsection{117}\label{section-2537}

\begin{otherlanguage}{persian}

بصدهزار لسان این خادم فانی از دوستان الهی مسئلت مینماید که از جان
بخروشند و جمعاً بکوشند که شاید اهل عالم مطّلع و آگاه شوند بآنچه که حقّ جلّ
جلاله اراده فرموده اگر احصای اختلاف و احساس اعمال غیرطیّبه نمایند لعمر
محبوبی و محبوبکم یفرّون و یستهزئون در یکی از الواح میفرماید باید دوستان
بمثابهٴ رگ شریان در جسد عالم متحرّک باشند لذا باید جمیع باعمال و اخلاقی
ظاهر شوند که جمیع عالم بسبب آن اقبال نمایند و بحرکت آیند صدهزار جان فدای
عاملین للّه کجایند نفوسیکه حقّرا بر خود مقدّم دارند و در سبیلش عمل نمایند
آنچه که سبب ترقّی و محبّت اهل عالم است از فضل حقّ امید هست که حبّهای عنایت
مودعهٴ در قلوب انبات نماید و باثمار لطیفه مزیّن گردد الأمر بیده و هو
القویّ القدیر

\end{otherlanguage}

Bi-ṣad-hizār lisān īn khādim-i fānī az dūstān-i ilāhī masʾalat
mī-namāyad ke az jān bi-khorūshand va jamʿan bi-kūshand ke shāyad ahl-i
ʿālam moṭṭaliʿ va āgāh shavand bi-ānche ke Ḥaqq jalla jalāluhu irāde
farmūde. Agar iḥṣā-yi ikhtilāf va iḥsās-i aʿmāl-i ghayr-i ṭayyebe
namāyand, la-ʿamr maḥbūbī wa-maḥbūbikum yafirrrūna wa-yastahziʾūn. Dar
yekī az alvāḥ mī-farmāyad bāyad dūstān bi-mathābiy-i rag-i sharyān dar
jasad-i ʿālam motaḥarrek bāshand. Lidhā bāyad jamīʿ bi-aʿmāl va akhlāqī
ẓāhir shavand ke jamīʿ-i ʿālam bi-sabab-i ān iqbāl namāyand va
bi-ḥarakat āyand. Ṣad-hizār jān fidā-yi ʿāmilīn li-llāh. Kojāyand nofūsī
ke Ḥaqq-rā bar khod moqaddam dārand va dar sabīlash ʿamal namāyand ānche
ke sabab-i taraqqī va maḥabbat-i ahl-i ʿālam ast? Az fażl-i Ḥaqq omīd
hast ke ḥabb-hā-yi ʿināyat-i mūdaʿiy-i dar qolūb inbāt namāyad va
bi-athmār-i laṭīfe mozayyan gardad. Al-amr bi-yadihi wa-huwa al-Qawiyy
al-Qadīr.

With a hundred thousand tongues, this ephemeral servant beseeches the
friends of God to cry out from their souls and strive together, so that
the people of the world may become informed and aware of what
God---exalted be His glory---has intended. If they reckon up the
dissension and perceive unworthy deeds, by the life of my Beloved and
yours, they will flee and scoff. In one of the Tablets He says that the
friends must be like a pulsating artery in the body of the world.
Therefore, all must manifest such deeds and qualities that the whole
world will thereby turn toward Him and be stirred into action. May a
hundred thousand lives be sacrificed for those who act for God! Where
are the souls who put God before themselves and do in His path what
brings about the advancement and love of the world's people? Through
God's grace, there is hope that the seeds of loving care deposited in
hearts will sprout and become adorned with exquisite fruits. The matter
is in His hand, and He is the Strong, the All-Powerful.

\subsection{118}\label{section-2538}

\begin{otherlanguage}{persian}

ذکر جناب آقا شیخ عبّاس علیه بهآء اللّه از قلم آنحضرت مرقوم اینفقره تلقاء
عرش مطلع نور احدیّه عرض شد و این لوح امنع مخصوص ایشان از سماء عنایت نازل
انشآءاللّه از فرات جاریهٴ در لوح الهی بیاشامند و باجر لقا و عنایت ربّنا
مالک الأسمآء فائز گردند قوله عزّ بیانه و جلّ فضله و جلّت عظمته

\end{otherlanguage}

Dhikr-i Janāb-i Āqā Shaykh ʿAbbās ʿalayhi bahāʾ Allāh az qalam-i ān
ḥażrat marqūm. Īn faqare tilqāʾ-i ʿarsh-i maṭlaʿ-i nūr-i aḥadiyye ʿarż
shod va īn lawḥ-i amnaʿ makhṣūṣ-i īshān az samā-yi ʿināyat nāzil. In
shāʾ Allāh az Furāt-i jāriyiy-i dar lawḥ-i ilāhī biyāshāmand va bi-ajr-i
liqā va ʿināyat-i Rabbinā Mālik al-Asmāʾ fāʾiz gardand. Qawluhu ʿazza
bayānuhu wa-jalla faḍluhu wa-jallat ʿaẓamatuhu.

Your pen had made mention of the honored Āqā Shaykh ʿAbbās, upon whom be
the glory of God. This matter was presented before the throne of the
Dawning-Place of the light of divine unity, and this most exalted Tablet
was sent down especially for him from the heaven of loving care. God
willing, he will drink from the Euphrates flowing in the divine Tablet
and attain the reward of the Presence and the favor of our Lord, the
Lord of Names. His words---mighty be His utterance, exalted His grace,
and great His grandeur:

\subsection{119}\label{section-2539}

\begin{otherlanguage}{arabic}

هو الشّاهد السّمیع العلیم

\end{otherlanguage}

Huwa al-Shāhid al-Samīʿ al-ʿAlīm.

He is the Witness, the All-Hearing, the All-Knowing.

\subsection{120}\label{section-2540}

\begin{otherlanguage}{persian}

یا عبّاس یذکرک المظلوم بذکر تنجذب به المقرّبون الی اللّه المهیمن القیّوم قد
اتت السّاعة و ارتفع النّدآء و ظهرت الصّیحة و مرّت الجبال و القوم هم لا
یشعرون قد انقلب العالم و استقرّ العرش و استوی علیه مالک الظّهور و مکلّم
الطّور و النّاس اکثرهم لا یفقهون قد ظهرت امّ الکتاب و تنادی بین الأرض و
السّمآء بأعلی النّدآء و تدع الکلّ الی المقام المحمود من النّاس من انکر و
منهم من اعرض و منهم من جادل بآیات اللّه ربّ ما کان و ما یکون و منهم من نقض
میثاق اللّه و عهده بظلم ناح به اهل مدائن الأسمآء و سکّان الملکوت فلمّا
تفرّسنا وجدنا اکثر اعدائنا العلمآء کذلک یقصّ لک لسان الأبهی من الأفق
الأعلی انّه لهو الحقّ علّام الغیوب

\end{otherlanguage}

Yā ʿAbbās yadhkuruka al-Maẓlūm bi-dhikr tanjadhibu bihi al-muqarrabūn
ilā Allāh al-Muhaymin al-Qayyūm. Qad atat al-sāʿa wa-irtafaʿa al-nidāʾ
wa-ẓaharat al-ṣayḥa wa-marrat al-jibāl wa-al-qawm hum lā yashʿurūn. Qad
inqalaba al-ʿālam wa-istaqarra al-ʿarsh wa-istawā ʿalayhi mālik al-ẓuhūr
wa-mukallim al-Ṭūr wa-al-nās aktharuhum lā yafqahūn. Qad ẓaharat umm
al-kitāb wa-tunādī bayna al-arḍ wa-al-samāʾ bi-aʿlā al-nidāʾ wa-tadʿu
al-kull ilā al-maqām al-maḥmūd. Min al-nās man ankara wa-minhum man
aʿraḍa wa-minhum man jādala bi-āyāt Allāh rabb mā kāna wa-mā yakūn
wa-minhum man naqaḍa mīthāq Allāh wa-ʿahdahu bi-ẓulm nāḥa bihi ahl
madāʾin al-asmāʾ wa-sukkān al-malakūt. Fa-lammā tafarrasnā wajadnā
akthar aʿdāʾinā al-ʿulamāʾ. Kadhālika yaquṣṣu laka lisān al-Abhā min
al-ufuq al-aʿlā innahu la-huwa al-Ḥaqq ʿAllām al-ghuyūb.

O ʿAbbās! The Wronged One remembers you with a remembrance through which
those brought near are drawn toward God, the All-Protecting, the
Self-Subsisting. The Hour has come, the call has risen, the cry has
sounded, and the mountains have passed away, yet the people are unaware.
The world has been overturned, the throne has been established, and the
Lord of Revelation, who spoke on Sinai, has seated Himself upon it, yet
most people do not understand. The Mother Book has appeared and calls
out between earth and heaven with the loudest call, summoning everyone
to the praised station. Among the people are those who have denied,
those who have turned away, those who have disputed the verses of God,
the Lord of all that has been and will be, and those who have broken
God's covenant and His pledge with an injustice that has caused the
inhabitants of the cities of names and the dwellers of the Kingdom to
lament. When We looked discerningly, We found that most of Our enemies
were the learned. Thus does the Tongue of the All-Glorious recount to
you from the most exalted horizon. He is indeed the Truth, the Knower of
things unseen.

\subsection{121}\label{section-2541}

\begin{otherlanguage}{arabic}

انّا اظهرنا العلم لعرفان المعلوم فلمّا الاح افق الظّهور و اتی المعلوم انکره
العلمآء فی اوّل الأمر الّا من شآء اللّه مالک الوجود قل یا معشر العلمآء ضعوا
ما عندکم تاللّه قد ارتفع صریر القلم الأعلی و یبشّر الکلّ بظهور المعلوم الّذی
توجّهتم الیه فیکلّ اصیل و بکور قل ما لی اراکم یا معشر الغافلین تقصدون
البیت و تعرضون عن الّذی رفعه بأمر من عنده اتّقوا اللّه و لا تتّبعوا اهوائکم
ان اتّبعوا من اتاکم بکتاب مشهود الّذی اذا ظهر خضعت له کتب العالم کلّها ان
انصفوا یا قوم و لا تکونوا من الّذین کفروا بالشّاهد و المشهود هذا یوم وعدتم
به فی صحف اللّه و کتبه و هذا امر اخذ عهده من عنده علم ما کان و ما یکون

\end{otherlanguage}

Innā aẓharnā al-ʿilm li-ʿirfān al-maʿlūm, fa-lammā alāḥa ufuq al-ẓuhūr
wa-atā al-maʿlūm ankarahu al-ʿulamāʾ fī awwal al-amr illā man shāʾa
Allāh Mālik al-wujūd. Qul yā maʿshar al-ʿulamāʾ, ḍaʿū mā ʿindakum.
Ta-Allāh qad irtafaʿa ṣarīr al-Qalam al-Aʿlā wa-yubashshiru al-kull
bi-ẓuhūr al-maʿlūm alladhī tawajjahtum ilayhi fī kull aṣīl wa-bukūr. Qul
mā lī arākum yā maʿshar al-ghāfilīn taqṣidūna al-Bayt wa-tuʿriḍūna ʿan
alladhī rafaʿahu bi-amr min ʿindih. Ittaqū Allāh wa-lā tattabiʿū
ahwāʾakum. In ittabiʿū man atākum bi-kitāb mashhūd alladhī idhā ẓahara
khaḍaʿat lahu kutub al-ʿālam kulluhā. In anṣifū yā qawm wa-lā takūnū min
alladhīna kafarū bi-al-shāhid wa-al-mashhūd. Hādhā yawm wuʿidtum bihi fī
ṣuḥuf Allāh wa-kutubih, wa-hādhā amr akhadha ʿahdahu man ʿindahu ʿilm mā
kāna wa-mā yakūn.

We brought forth knowledge so that the One who is its object might be
recognized. Yet when the horizon of revelation shone and that One came,
the learned denied Him at the very outset, except those whom God, the
Master of existence, willed otherwise. Say: O company of the learned!
Set aside what you possess. By God, the sound of the Most Exalted Pen
has risen, announcing to all the appearance of the One toward whom you
turned every evening and morning. Say: Why do I see you, O company of
the heedless, making your way to the House while turning away from Him
who raised it by a command from Himself? Fear God and do not follow your
desires. Rather, follow Him who has come to you with a manifest Book,
before which, when it appeared, all the books of the world bowed down.
Be fair, O people, and do not be among those who have denied the witness
and the witnessed. This is the Day promised to you in the scrolls of God
and His books, and this is the Cause whose covenant was taken by Him who
possesses the knowledge of all that has been and will be.

\subsection{122}\label{section-2542}

\begin{otherlanguage}{arabic}

قل یا ملأ الظّالمین قد عملتم ما صاح به الملأ الأعلی و ناح اهل الفردوس و
صعدت زفرات الجنّة العلیا و نزلت عبرات الملائکة و الرّوح قل ان کان عندکم
اعظم عمّا اتی من جبروت الارادة فأتوا به و لا تتّبعوا کلّ جاهل مردود قل
تعالوا لأریکم ما غفلتم عنه فی هذا الیوم الّذی نسب الی اللّه العزیز الودود
هذا یوم فیه تضوّعت رائحة الرّحمن و فاحت نفحات الوحی و نطق عندلیب الأمر علی
الأغصان الملک للّه مالک الملوک

\end{otherlanguage}

Qul yā malaʾ al-ẓālimīn, qad ʿamiltum mā ṣāḥa bihi al-Malaʾ al-Aʿlā
wa-nāḥa ahl al-Firdaws wa-ṣaʿidat zafarāt al-janna al-ʿulyā wa-nazalat
ʿabarāt al-malāʾika wa-al-Rūḥ. Qul in kāna ʿindakum aʿẓam ʿammā atā min
jabarūt al-irāda faʾtū bihi wa-lā tattabiʿū kulla jāhil mardūd. Qul
taʿālaw li-uriyakum mā ghafaltum ʿanhu fī hādhā al-yawm alladhī nusiba
ilā Allāh al-ʿAzīz al-Wadūd. Hādhā yawm fīhi taḍawwaʿat rāʾiḥat
al-Raḥmān wa-fāḥat nafaḥāt al-waḥy wa-naṭaqa ʿandalīb al-amr ʿalā
al-aghṣān: al-mulk li-Allāh Mālik al-mulūk.

Say: O company of the unjust! You have done what has caused the
Concourse on High to cry out, the people of Paradise to lament, the
sighs of the highest heaven to ascend, and the tears of the angels and
the Spirit to fall. Say: If you possess anything greater than what has
come from the dominion of the divine Will, bring it forth, and do not
follow every ignorant outcast. Say: Come, that I may show you what you
have overlooked in this Day that has been attributed to God, the
Almighty, the All-Loving. This is a Day in which the fragrance of the
All-Merciful has spread abroad, the breezes of revelation have wafted,
and the Nightingale of the Cause has spoken upon the branches:
``Sovereignty belongs to God, the King of kings.''

\subsection{123}\label{section-2543}

\begin{otherlanguage}{arabic}

یا عبّاس قم علی خدمة الأمر باسم ربّک مالک القدم و ذکّر النّاس بهذا النّبأ
الأعظم الّذی اذا ظهر ارتعدت فرائص الأوهام و الظّنون ایّاک ان تمنعک عن مالک
الأسمآء شوکة الأمرآء و شبهات العلمآء کن کالجبل الرّاسخ علی هذا الأمر
المحتوم قل انّه اتی بالحقّ و لا یلتفت الی اقبالکم و اعراضکم ظهر و اظهر ما
اراد رغماً لأنفکم لو انتم تعلمون انّه هو الّذی ینطق فیکلّ شأن انّه لا اله الّا
انا الفرد الواحد العزیز المحبوب

\end{otherlanguage}

Yā ʿAbbās, qum ʿalā khidmat al-amr bi-ism rabbika Mālik al-qidam,
wa-dhakkir al-nās bi-hādhā al-nabaʾ al-aʿẓam alladhī idhā ẓahara
irtaʿadat farāʾiṣ al-awhām wa-al-ẓunūn. Iyyāka an tamnaʿaka ʿan Mālik
al-asmāʾ shawkat al-umarāʾ wa-shubuhāt al-ʿulamāʾ. Kun ka-al-jabal
al-rāsikh ʿalā hādhā al-amr al-maḥtūm. Qul innahu atā bi-al-ḥaqq wa-lā
yaltafitu ilā iqbālikum wa-iʿrāḍikum. Ẓahara wa-aẓhara mā arāda raghman
li-anfikum law antum taʿlamūn. Innahu huwa alladhī yanṭiqu fī kull
shaʾn: innahu lā ilāha illā anā al-Fard al-Wāḥid al-ʿAzīz al-Maḥbūb.

O ʿAbbās! Arise to serve the Cause in the name of your Lord, the Master
of eternity, and remind the people of this Most Great Announcement,
whose appearance has made the limbs of vain imaginings and idle
suppositions tremble. Beware lest the power of rulers and the doubts
raised by the learned keep you from the Master of names. Be as a firmly
rooted mountain in this irrevocable Cause. Say: He has come in truth and
pays no heed to your acceptance or rejection. He has appeared and
brought forth whatever He wished, despite you, if you only knew. He is
the One who declares in every circumstance: ``There is no God but Me,
the Single, the One, the Almighty, the Beloved.''

\subsection{124}\label{section-2544}

\begin{otherlanguage}{arabic}

کذلک فتحنا علی وجهک باب الفضل بمفتاح العدل اذا فزت به قل لک الحمد یا اله
الغیب و الشّهود اشهد انّک نبذت الرّاحة لارتقآء الأمم و قبلت الذّلة لعزّ
العالم و اخترت الغربة لبلوغ العباد الی الوطن الأعلی المقام الّذی جعله
اللّه مقدّساً عمّا تدرکه العقول یا ایّها المتوجّه الی الوجه خذ الرّحیق المختوم
باسم المظلوم ثمّ اشرب منه بهذا الذّکر الّذی به انجذبت الأرواح و القلوب انّا
کتبنا لک اجر اللّقآء من قلمی الأعلی فی الصّحیفة الحمرآء ان افرح ثمّ اشکر ربّ
البیت المعمور انّا قبلنا اقبالک و توجّهک و عملک فی سبیل اللّه طوبی لک بما
زیّنت اعمالک بطراز القبول انتهی

\end{otherlanguage}

Kadhālika fataḥnā ʿalā wajhika bāb al-faḍl bi-miftāḥ al-ʿadl. Idhā fuzta
bihi qul: laka al-ḥamd yā Ilāh al-ghayb wa-al-shuhūd. Ashhadu annaka
nabadhta al-rāḥa li-irtiqāʾ al-umam wa-qabilta al-dhilla li-ʿizz
al-ʿālam wa-ikhtarta al-ghurba li-bulūgh al-ʿibād ilā al-waṭan al-aʿlā,
al-maqām alladhī jaʿalahu Allāh muqaddasan ʿammā tudrikuhu al-ʿuqūl. Yā
ayyuhā al-mutawajjih ilā al-wajh, khudh al-raḥīq al-makhtūm bi-ism
al-Maẓlūm thumma ishrab minhu bi-hādhā al-dhikr alladhī bihi injadhabat
al-arwāḥ wa-al-qulūb. Innā katabnā laka ajr al-liqāʾ min qalamī al-aʿlā
fī al-ṣaḥīfa al-ḥamrāʾ. In ifraḥ thumma ishkur Rabb al-Bayt al-Maʿmūr.
Innā qabilnā iqbālaka wa-tawajjuhaka wa-ʿamalaka fī sabīl Allāh. Ṭūbā
laka bi-mā zayyanat aʿmāluka bi-ṭirāz al-qabūl. Intahā.

Thus We have opened before you the door of grace with the key of
justice. When you attain it, say: ``Praise be to You, O God of the
unseen and the seen! I bear witness that You have cast aside comfort for
the advancement of the nations, accepted abasement for the exaltation of
the world, and chosen exile so that Your servants might reach the most
exalted homeland, the station that God has sanctified beyond the grasp
of minds.'' O you who have turned toward the Face! Take the sealed wine
in the name of the Wronged One, then drink from it with this remembrance
through which souls and hearts have been drawn in rapture. With My Most
Exalted Pen We have inscribed for you the reward of attaining Our
presence upon the Crimson Scroll. Rejoice, then give thanks to the Lord
of the Frequented House. We have accepted your turning toward Us, your
devotion, and your work in the path of God. Blessed are you, for your
deeds have been adorned with the ornament of acceptance. End.

\subsection{125}\label{section-2545}

\begin{otherlanguage}{persian}

یا محبوب فؤادی گویا بر هر کلمه‌ئی از کلمات الهی یک منادی موجود و بأعلی
النّدآء من فی الوجود را ندا مینماید سبحان‌اللّه عالم را چه واقع که از این
فضل اکبر خود را محروم نموده‌اند و بما لا یسمن و لا یغنی مشغول گشته‌اند
باری الحمد للّه جناب مذکور بشرافت کبری فائز و بذکر قلم اعلی مشرّف
انشآءاللّه قدر اینمقام اعلی را بدانند و از سلسبیل و کوثر و تسنیم بیان
الهی بنوشند و بنوشانند و آنمحبوب و این عبد بکلمهٴ هنیئاً ناطق شویم و
مریئاً گوئیم حبّذا هذا العرف نمیدانم از کجا میآید الحمد لمظهره و مرسله و
مبعثه و فاطره و خالقه و محدثه اینفانیهم خدمت ایشان تکبیر و سلام معروض
میدارد و دعای توفیق استدعا مینماید انّ ربّنا الرّحمن علی کلّشیء قدیر

\end{otherlanguage}

Yā maḥbūb-i fuʾādī, gūyā bar har kalamiʾī az kalimāt-i ilāhī yak munādī
mawjūd va bi-aʿlā al-nidāʾ man fī al-wujūd rā nidā mīnamāyad. Subḥān
Allāh, ʿālam rā chih vāqiʿ kih az īn faḍl-i akbar khud rā maḥrūm
namūdih-and va bi-mā lā yusmin wa-lā yughnī mashghūl gashtih-and. Bārī
al-ḥamdu li-Allāh, jināb-i madhkūr bi-sharāfat-i kubrā fāʾiz va
bi-dhikr-i Qalam-i Aʿlā musharraf. In shāʾa Allāh qadr-i īn maqām-i aʿlā
rā bidānand va az Salsabīl va Kawthar va Tasnīm-i bayān-i ilāhī
binūshand va binūshānand va ān maḥbūb va īn ʿabd bi-kalami-yi hanīʾan
nāṭiq shavīm va marīʾan gūʾīm. Ḥabbadhā hādhā al-ʿarf, nimīdānam az kujā
mīāyad. Al-ḥamd li-muẓhirihi wa-mursilihi wa-mabʿathihi wa-fāṭirihi
wa-khāliqihi wa-muḥdithih. Īn fānī ham khidmat-i īshān takbīr va salām
maʿrūḍ mīdārad va duʿā-yi tawfīq istidʿā mīnamāyad. Inna Rabbanā
al-Raḥmān ʿalā kull shayʾ qadīr.

O beloved of my heart! It is as though a herald stands upon every word
of the divine utterances, calling at the top of his voice to all who
exist. Glory be to God! What has happened to the world that its people
have deprived themselves of this greatest grace and occupied themselves
with what neither nourishes nor satisfies? In any event, praise be to
God, the gentleman mentioned has attained the greatest honor and has
been honored by the remembrance of the Most Exalted Pen. God willing,
may he appreciate the value of this most exalted station, drink from the
Salsabīl, Kawthar, and Tasnīm of divine utterance, and give others to
drink, so that you, that beloved one, and this servant may utter the
word ``Delightful!'' and say, ``Wholesome!'' How wonderful is this
fragrance! I do not know where it comes from. Praise be to Him who
manifests it, sends it forth, is its source, fashions it, creates it,
and brings it into being. This mortal also offers him greetings of ``God
is most great'' and peace, and requests his prayers for divine
assistance. Our Lord, the All-Merciful, is indeed powerful over all
things.

\subsection{126}\label{section-2546}

\begin{otherlanguage}{arabic}

و ما ذکر حضرتکم فی جناب محمّد کریم الّذی هاجر فی سنة القبل قد فاز باصغآء
ربّنا الرّحمن و نزّل له ما تفرح به افئدة العارفین قوله تمّت کلمته و کمل
برهانه

\end{otherlanguage}

Wa-mā dhakarat ḥaḍratukum fī janāb Muḥammad Karīm alladhī hājara fī
sanat al-qabl qad fāza bi-iṣghāʾ Rabbinā al-Raḥmān, wa-nuzzila lahu mā
tafraḥu bihi afʾidat al-ʿārifīn. Qawluhu, tammat kalimatuhu wa-kamula
burhānuh:

What you mentioned concerning Muḥammad Karīm, who emigrated last year,
has reached the hearing of our Lord, the All-Merciful, and there has
been revealed for him that which gladdens the hearts of those who
recognize Him. His words---perfect is His word and complete His proof:

\subsection{127}\label{section-2547}

\begin{otherlanguage}{arabic}

بسمی الّذی به نزل کلّ کتاب مبین

\end{otherlanguage}

Bi-ismī alladhī bihi nazala kull kitāb mubīn.

In My name, through which every clear Book has been sent down.

\subsection{128}\label{section-2548}

\begin{otherlanguage}{persian}

شهد اللّه انّه لا اله الّا انا المهیمن القیّوم قد انزلنا الآیات و اظهرنا
البیّنات و ارسلنا الرّسل لیبشّروا العباد بهذا الظّهور الّذی به ظهر ما کان
مذکوراً بلسان الأنبیآء و مسطوراً فی کتب اللّه ربّ العالمین اذا ظهر الأمر و
نزّلت الآیات و استقرّ هیکل الظّهور علی العرش قام النّاس علی الاعراض الّا من
حفظته ید الاقتدار من لدن قویّ قدیر یا کریم ان استمع ما ینادیک به المظلوم
من شطر السّجن و یذکرک بما تجد منه عرف البقآء انّ ربّک لهو العلیم الحکیم
طوبی لوجه توجّه الی الأفق الأعلی و لرجل سلک سبیل اللّه العزیز الحمید

\end{otherlanguage}

Shahida Allāh annahu lā ilāha illā anā al-Muhaymin al-Qayyūm. Qad
anzalnā al-āyāt wa-aẓharnā al-bayyināt wa-arsalnā al-rusul
li-yubashshirū al-ʿibād bi-hādhā al-ẓuhūr alladhī bihi ẓahara mā kāna
madhkūran bi-lisān al-anbiyāʾ wa-masṭūran fī kutub Allāh Rabb
al-ʿālamīn. Idhā ẓahara al-amr wa-nuzzilat al-āyāt wa-istaqarra haykal
al-ẓuhūr ʿalā al-ʿarsh qāma al-nās ʿalā al-iʿrāḍ illā man ḥafiẓathu yad
al-iqtidār min ladun Qawiyy Qadīr. Yā Karīm, in istamiʿ mā yunādīka bihi
al-Maẓlūm min shaṭr al-sijn wa-yadhkuruka bi-mā tajidu minhu ʿarf
al-baqāʾ. Inna rabbaka la-huwa al-ʿAlīm al-Ḥakīm. Ṭūbā li-wajh tawajjaha
ilā al-ufuq al-aʿlā wa-li-rajul salaka sabīl Allāh al-ʿAzīz al-Ḥamīd.

God bears witness that there is no God but Me, the All-Protecting, the
Self-Subsisting. We have sent down the verses, brought forth the clear
proofs, and dispatched the Messengers to announce to the servants this
Revelation, through which there has appeared what was spoken by the
tongues of the Prophets and inscribed in the books of God, the Lord of
the worlds. When the Cause appeared, the verses were sent down, and the
Temple of Revelation was established upon the throne, the people rose up
in rejection, except those whom the hand of might protected on behalf of
One who is Strong and Powerful. O Karīm! Listen to what the Wronged One
calls to you from the direction of the prison, remembering you with
words from which you may perceive the fragrance of everlasting life.
Your Lord is indeed the All-Knowing, the All-Wise. Blessed is the face
that has turned toward the most exalted horizon, and the man who has
walked the path of God, the Almighty, the Praiseworthy.

\subsection{129}\label{section-2549}

\begin{otherlanguage}{arabic}

انّا سمعنا ندائک اذ وجدت عرف قمیصی من شطر الحجاز المقام الّذی نوّرناه
بأنوار وجه محمّد رسول اللّه و خاتم النّبیّین انّا سمعنا ذکرک من قبل و ذکرناک
فی کتاب ما اطّلع به الّا الفرد الخبیر طوبی لمن نبذ الهوی و اتّبع ما امر به
من لدی اللّه مالک الوری انّه من اعلی الخلق فی کتابه المبین نعیماً لک بما
نبذت الأوهام عن ورائک مقبلاً الی مشرق وحی ربّک الرّحمن الرّحیم انّه ینصر من
یشآء بسلطان من عنده و یهدی المقبلین الی صراطه المستقیم هذا یوم فیه ظهر
ما لا ظهر فی ازل الآزال یشهد بذلک من ینطق بالحقّ فیکلّ شأن انّه لا اله الّا
انا العلیّ العظیم انّک اذا فزت بکتابی و وجدت عرف بیانی سبّح بحمد ربّک ثمّ
اشکره بهذا الفضل الّذی ما رأت عین الابداع شبهه یشهد بذلک کلّ عالم منصف و
کلّ عارف بصیر لا تحزن من شیء قد قبل ما عملته فی سبیله انّه لا یضیع اجر
المحسنین کذلک انار افق البیان بنیّر ذکر ربّک المشفق الکریم انتهی

\end{otherlanguage}

Innā samiʿnā nidāʾaka idh wajadta ʿarf qamīṣī min shaṭr al-Ḥijāz,
al-maqām alladhī nawwarnāhu bi-anwār wajh Muḥammad Rasūl Allāh wa-Khātam
al-Nabiyyīn. Innā samiʿnā dhikraka min qabl wa-dhakarnāka fī kitāb mā
iṭṭalaʿa bihi illā al-Fard al-Khabīr. Ṭūbā li-man nabadha al-hawā
wa-ittabaʿa mā umira bihi min ladā Allāh Mālik al-warā. Innahu min aʿlā
al-khalq fī kitābihi al-mubīn. Naʿīman laka bi-mā nabadhta al-awhām ʿan
warāʾika muqbilan ilā mashriq waḥy rabbika al-Raḥmān al-Raḥīm. Innahu
yanṣuru man yashāʾ bi-sulṭān min ʿindihi wa-yahdī al-muqbilīn ilā
ṣirāṭihi al-mustaqīm. Hādhā yawm fīhi ẓahara mā lā ẓahara fī azal
al-āzāl. Yashhadu bi-dhālika man yanṭiqu bi-al-ḥaqq fī kull shaʾn:
innahu lā ilāha illā anā al-ʿAliyy al-ʿAẓīm. Innaka idhā fuzta bi-kitābī
wa-wajadta ʿarf bayānī sabbiḥ bi-ḥamd rabbika thumma ishkurhu bi-hādhā
al-faḍl alladhī mā raʾat ʿayn al-ibdāʿ shibhah. Yashhadu bi-dhālika kull
ʿālim munṣif wa-kull ʿārif baṣīr. Lā taḥzan min shayʾ, qad qubila mā
ʿamiltahu fī sabīlih. Innahu lā yuḍīʿu ajr al-muḥsinīn. Kadhālika anāra
ufuq al-bayān bi-nayyir dhikr rabbika al-Mushfiq al-Karīm. Intahā.

We heard your call when you perceived the fragrance of My robe from the
direction of the Ḥijāz, the place We illumined with the lights of the
face of Muḥammad, the Messenger of God and the Seal of the Prophets. We
heard your remembrance before and remembered you in a Book known to none
but the Unique, the All-Informed. Blessed is the one who has cast aside
desire and followed what he was commanded by God, the Master of
creation. He is among the most exalted of creation in His clear Book.
Joy be yours, for you have cast vain imaginings behind you and turned
toward the Dawning-Place of the revelation of your Lord, the
All-Merciful, the Merciful. He aids whomever He wills with authority
from Himself and guides those who turn toward Him to His straight path.
This is a Day in which there has appeared what never appeared in the
eternity of eternities. To this bears witness the One who declares the
truth in every circumstance: ``There is no God but Me, the Most High,
the Most Great.'' When you receive My letter and perceive the fragrance
of My utterance, glorify your Lord with praise, then thank Him for this
grace, the like of which the eye of creation has never seen. Every
fair-minded scholar and every discerning person of understanding bears
witness to this. Do not grieve over anything. What you have done in His
path has been accepted. He does not let the reward of those who do good
be lost. Thus has the horizon of utterance been illumined by the
luminary of the remembrance of your Lord, the Compassionate, the
Generous. End.

\subsection{130}\label{section-2550}

\begin{otherlanguage}{persian}

هذه نعمة اخری و رحمة اخری و فضل آخر لمن اقبل الیه و تقرّب الی ساحة قدسه و
استنشق الهوآء الّذی کان مجاوراً هوآء هذا السّجن المنیع انشآءاللّه از نار
محبّت الهی که از هر حرفی از حروفاتش ظاهر و باهر است مشتعل شوند و باستقامت
تمام بحکمت و بیان بذکرش ناطق له الفضل و المنّة فیکلّ الأحوال و هو الفضّال
فی ازل الآزال

\end{otherlanguage}

Hādhihi niʿma ukhrā wa-raḥma ukhrā wa-faḍl ākhar li-man aqbala ilayhi
wa-taqarraba ilā sāḥat qudsihi wa-istanshaqa al-hawāʾ alladhī kāna
mujāwiran hawāʾ hādhā al-sijn al-manīʿ. In shāʾa Allāh az nār-i
maḥabbat-i ilāhī kih az har ḥarfī az ḥurūfātash ẓāhir va bāhir ast
mushtaʿil shavand va bi-istiqāmat-i tamām bi-ḥikmat va bayān bi-dhikrih
nāṭiq. Lahu al-faḍl wa-al-minna fī kull al-aḥwāl wa-huwa al-Faḍḍāl fī
azal al-āzāl.

This is another blessing, another mercy, and another grace for one who
has turned toward Him, drawn near to the court of His holiness, and
breathed the air that was adjacent to the air of this inaccessible
prison. God willing, may he be set ablaze by the fire of divine love
that is manifest and radiant in every one of its letters, and with
complete steadfastness, wisdom, and eloquence speak His remembrance. His
are grace and bounty in all circumstances, and He is the All-Bountiful
from the eternity of eternities.

\subsection{131}\label{section-2551}

\begin{otherlanguage}{persian}

و اینکه آنمحبوب ذکر نفوس مقبله علیهم بهآء اللّه فرموده بودند در ساحت اقدس
کان اللّه و لم یکن معه من شیء واحداً بعد واحد عرض شد و مخصوص هر یک از سماء
مشیّت نازل شد آنچه که غریب را بوطن هدایت کند و تشنه را بسلسبیل و عاشقرا
بمعشوق رساند و طالب را بشطر مطلوب کشاند قال عظم احسانه و جلّ کبریائه و
احاط فضله و علا سلطانه

\end{otherlanguage}

Va īnkih ān maḥbūb dhikr-i nufūs-i muqbilih, ʿalayhim Bahāʾ Allāh,
farmūdih būdand, dar sāḥat-i aqdas, kāna Allāh wa-lam yakun maʿahu min
shayʾ, wāḥidan baʿda wāḥid ʿarḍ shud va makhṣūṣ-i har yak az samāʾ-i
mashiyyat nāzil shud ānchih kih gharīb rā bi-waṭan hidāyat kunad va
tishnih rā bi-Salsabīl va ʿāshiq rā bi-maʿshūq rasānad va ṭālib rā
bi-shaṭr-i maṭlūb kishānad. Qāla, ʿaẓuma iḥsānuhu wa-jalla kibriyāʾuhu
wa-aḥāṭa faḍluhu wa-ʿalā sulṭānuh:

As for those souls who have turned toward Him---upon them be the glory
of God---whom you, that beloved one, had mentioned, they were presented
one by one in the Most Holy Court: ``God was, and nothing was with
Him.'' For each one specifically there descended from the heaven of the
divine Will that which guides the stranger to his homeland, brings the
thirsty to Salsabīl and the lover to the beloved, and draws the seeker
toward the One he seeks. He said---great is His beneficence, glorious
His majesty, all-encompassing His grace, and exalted His sovereignty:

\subsection{132}\label{section-2552}

\begin{otherlanguage}{arabic}

بسمی المشرق من افق سمآء ملکوتی العظیم

\end{otherlanguage}

Bi-ismī al-mushriq min ufuq samāʾ malakūtī al-ʿaẓīm.

In My name, shining from the horizon of the heaven of My mighty Kingdom.

\subsection{133}\label{section-2553}

\begin{otherlanguage}{persian}

یا سلیم یذکرک الکریم من شطر السّجن خالصاً لوجهه لتسمع صریر قلمی الأعلی و
هدیر طیر البقآء علی اغصان سدرة المنتهی انّه لا اله الّا انا العلیم الخبیر
طوبی لقلب انجذب بآیات اللّه و لوجه توجّه الیه و للّسان نطق بذکره البدیع لا
تنظر الی العلمآء و کبریائهم ان انظر الی من یذکرک فی هذا السّجن العظیم قل
انّه ظهر بالحقّ و قام علی الأمر بسلطان ما منعته شئونات الأوهام و لا حجّبته
حجبات المعرضین کم من عالم منع عن فضل اللّه و رحمته و کم من امّیّ فاز بهذا
الرّحیق المنیر ان اشکر اللّه بما عرّفک مطلع الأمر و مشرق الوحی و عصمک عن
شبهات الّذین کفروا ببرهانه القویم ان استقم علی حبّ اللّه و امره علی شأن یجد
منک الموحّدون عرف الاستقامة فیهذا النّبأ الّذی استبشر به الملأ الأعلی و اهل
الفردوس علی مقام کریم

\end{otherlanguage}

Yā Salīm, yadhkuruka al-Karīm min shaṭr al-sijn khāliṣan li-wajhihi
li-tasmaʿa ṣarīr qalamī al-aʿlā wa-hadīr ṭayr al-baqāʾ ʿalā aghṣān
Sidrat al-Muntahā. Innahu lā ilāha illā anā al-ʿAlīm al-Khabīr. Ṭūbā
li-qalb injadhaba bi-āyāt Allāh wa-li-wajh tawajjaha ilayhi wa-li-lisān
naṭaqa bi-dhikrihi al-badīʿ. Lā tanẓur ilā al-ʿulamāʾ wa-kibriyāʾihim.
In inẓur ilā man yadhkuruka fī hādhā al-sijn al-ʿaẓīm. Qul innahu ẓahara
bi-al-ḥaqq wa-qāma ʿalā al-amr bi-sulṭān mā manaʿathu shuʾūnāt al-awhām
wa-lā ḥajjabathu ḥujubāt al-muʿriḍīn. Kam min ʿālim muniʿa ʿan faḍl
Allāh wa-raḥmatihi wa-kam min ummiyy fāza bi-hādhā al-raḥīq al-munīr. In
ishkur Allāh bi-mā ʿarrafaka maṭlaʿ al-amr wa-mashriq al-waḥy
wa-ʿaṣamaka ʿan shubuhāt alladhīna kafarū bi-burhānihi al-qawīm. In
istaqim ʿalā ḥubb Allāh wa-amrihi ʿalā shaʾn yajidu minka al-muwaḥḥidūn
ʿarf al-istiqāma fī hādhā al-nabaʾ alladhī istabshara bihi al-Malaʾ
al-Aʿlā wa-ahl al-Firdaws ʿalā maqām karīm.

O Salīm! The Generous One remembers you from the direction of the
prison, solely for His own sake, so that you may hear the sound of My
Most Exalted Pen and the call of the Bird of everlasting life upon the
branches of the Lote-Tree of the Utmost Boundary. There is no God but
Me, the All-Knowing, the All-Informed. Blessed is the heart enraptured
by the verses of God, the face that has turned toward Him, and the
tongue that has spoken His wondrous remembrance. Do not look at the
learned and their arrogance. Rather, look toward Him who remembers you
in this mighty prison. Say: He has appeared in truth and arisen for the
Cause with an authority that the workings of vain imaginings have
neither hindered nor the veils of those who turn away obscured. How many
a learned person has been kept from the grace of God and His mercy, and
how many an unlettered person has attained this luminous wine! Thank God
for enabling you to recognize the Dawning-Place of the Cause and the
Dayspring of revelation, and for protecting you from the doubts raised
by those who have denied His firm proof. Be steadfast in the love of God
and His Cause, so that those who affirm His unity may perceive from you
the fragrance of steadfastness in this Announcement, which has gladdened
the Concourse on High and the people of Paradise in a noble station.

\subsection{134}\label{section-2554}

\begin{otherlanguage}{arabic}

و هذا ما نزل لمن سمّی بمیرزا ابوالحسن

\end{otherlanguage}

Wa-hādhā mā nazala li-man summiya bi-Mīrzā Abū al-Ḥasan.

And this is what was revealed for the one named Mīrzā Abū al-Ḥasan.

\subsection{135}\label{section-2555}

\begin{otherlanguage}{arabic}

بسمی العلیم الحکیم

\end{otherlanguage}

Bi-ismī al-ʿAlīm al-Ḥakīm.

In My name, the All-Knowing, the All-Wise.

\subsection{136}\label{section-2556}

\begin{otherlanguage}{persian}

یا اباالحسن ذکرت نزد مظلوم مذکور و سماء بیان رحمن بتو متوجّه این آسمانیست
که سیّارات آن تا حین احصا نشده و نخواهد شد و شموس و اقمار آن لازال فی قطب
الزّوال لائح و مشرق هر حکیمی از احصایش متحیّر و هر قویّ قادری از ادراک و
تحدیدش عاجز و ترا وصیّت مینماید براستی و انصاف چه اگر جمیع خلق عالم بطراز
انصاف مزیّن می‌شدند از عرفان اسم اعظم و مالک قدم محروم نمیگشتند اینمظلوم
بگمان اهل ظنون بغیر ما حکم به اللّه حکم نمود حضرت رسول و خاتم انبیا روح
ما سواه فداه چه تقصیر بر او ثابت که در سنین متوالیات بر اذیّتش قیام
نمودند برخی کاذبش گفتند و بعضی مجنون و حزبی ساحر و قومی مفتری جمیع آنچه
ذکر شد آنعبد و سایر عباد شنیده‌اند و دانسته‌اند از آن وجود مبارک گذشته در
حضرت روح تفکّر نما که سبب چه بود که من علی الأرض بر انکارش قیام نمودند
وارد شد بر مظاهر امر الهی آنچه که هر صاحب انصافی گریست و نوحه نمود سبب و
علّت اعراض در هر عصر و قرن علمای آن زمان بودند که حبّ ریاست ایشانرا از
اقبال منع مینمود الا انّهم فی ضلال مبین امید هست که آنجناب در مامضی تفکّر
نماید و بصراط مستقیم الهی پی‌برد

\end{otherlanguage}

Yā Abā al-Ḥasan, dhekrat nazd-e Maẓlūm madhkūr va samā-ye bayān-e Raḥmān
be to motavajjeh. Īn āsmānīst ke sayyārāt-e ān tā ḥīn eḥṣā nashodeh va
nakhāhad shod va shomūs va aqmār-e ān lā yazāl fī quṭb al-zawāl lāʾeḥ va
moshreq. Har ḥakīmī az eḥṣāyash motaḥayyer va har qaviyy-e qāderī az
edrāk va taḥdīdash ʿājez. Va to rā vaṣiyyat mīnamāyad be rāstī va enṣāf,
che agar jamīʿ-e khalq-e ʿālam be ṭerāz-e enṣāf mozayyan mīshodand az
ʿerfān-e Esm-e Aʿẓam va Mālek-e Qedam maḥrūm nemīgashtand. Īn Maẓlūm be
gomān-e ahl-e ẓonūn be ghayr-e mā ḥakama bihi Allāh ḥokm namūd. Ḥażrat-e
Rasūl va Khātam-e Anbiyā, rūḥ-e mā sevāh fedāh, che taqṣīr bar ū thābet
ke dar senīn-e motavāleyāt bar adhiyyatash qīyām namūdand? Barkhī
kādhebash goftand va baʿżī majnūn va ḥezbī sāḥer va qowmī moftarī.
Jamīʿ-e ānche dhekr shod ān ʿabd va sāyer-e ʿebād shenīdeh-and va
dānesteh-and. Az ān vojūd-e mobārak godhashteh dar Ḥażrat-e Rūḥ tafakkor
namā ke sabab che būd ke man ʿalā al-arḍ bar enkārash qīyām namūdand.
Vāred shod bar maẓāher-e Amr-e Elāhī ānche ke har ṣāḥeb-e enṣāfī gerīst
va nowḥeh namūd. Sabab va ʿellat-e eʿrāż dar har ʿaṣr va qarn ʿolamā-ye
ān zamān būdand ke ḥobb-e reyāsat īshān rā az eqbāl manʿ mīnamūd. A-lā
innahum fī ḍalālin mubīn. Omīd hast ke ān janāb dar mā mażā tafakkor
namāyad va be ṣerāṭ-e mostaqīm-e Elāhī pey barad.

O Abū al-Ḥasan, you have been mentioned in the presence of the Wronged
One, and the heaven of the All-Merciful's utterance is turned toward
you. This is a heaven whose planets have not yet been counted and never
will be, and whose suns and moons are forever radiant and shining at the
pole of decline. Every sage is bewildered at counting them, and every
mighty and capable person is powerless to comprehend or delimit it. He
counsels you to be truthful and fair, for if all the people of the world
had adorned themselves with fairness, they would not have been deprived
of recognizing the Most Great Name and the Lord of Eternity. In the
estimation of those given to conjecture, this Wronged One has ruled
contrary to what God has decreed. What fault was established against the
Messenger and Seal of the Prophets---may the souls of all others be
sacrificed for Him---that they rose up to torment Him for successive
years? Some called Him a liar, others a madman, one faction a sorcerer,
and another an impostor. That servant and the other servants have heard
and known all that has been mentioned. Passing beyond that blessed
Being, reflect on the Spirit: what caused those on earth to rise up in
denial of Him? What befell the Manifestations of the Cause of God made
every fair-minded person weep and lament. In every age and century, the
cause and source of rejection were the religious scholars of that time,
whose love of leadership prevented them from turning toward Him. Indeed,
they are in manifest error. It is hoped that you will reflect on what
has passed and find the straight path of God.

\subsection{137}\label{section-2557}

\begin{otherlanguage}{arabic}

شک و شبهه‌ئی نبوده و نیست که اینمظلوم للّه میگوید و الی اللّه میخواند لا
یضرّه الاعراض و لا ینفعه الاقبال ظنون و اوهام خلقرا بخلق گذار و بقلب منیر
بحقّ جلّ جلاله توجّه نما که شاید باب معرفت بر وجهت بگشاید و از انوار وجه
بعد فناء اشیاء ترا محروم نفرماید انّه یقول الحقّ و یهدی النّاس الی افقه
المنیر

\end{otherlanguage}

Shakk va shobheʾī nabūdeh va nīst ke īn Maẓlūm li-llāh mīgūyad va ilā
Allāh mīkhānad. Lā yaḍurruhu al-iʿrāḍ wa-lā yanfaʿuhu al-iqbāl. Ẓonūn va
owhām-e khalq rā be khalq godhār va be qalb-e monīr be Ḥaqq, jalla
jalāluhu, tavajjoh namā ke shāyad bāb-e maʿrefat bar vajhat bogshāyad va
az anvār-e vajh baʿd-e fanā-ye ashyā to rā maḥrūm nafarmāyad. Innahu
yaqūlu al-ḥaqq wa-yahdī al-nāsa ilā ufuqihi al-munīr.

There has never been, nor is there now, any doubt that this Wronged One
speaks for God and summons people to God. Rejection does not harm Him,
nor does acceptance benefit Him. Leave the people's conjectures and
imaginings to them, and turn with a radiant heart toward the True One,
exalted be His glory, so that He may open the gate of recognition before
you and not deprive you of the lights of His face after all things have
passed away. He speaks the truth and guides people to His luminous
horizon.

\subsection{138}\label{section-2558}

\begin{otherlanguage}{arabic}

و هذا ما نزل لجناب مشهدی نصراللّه

\end{otherlanguage}

Wa-hādhā mā nazala li-Janāb-i Mashhadī Naṣr Allāh.

And this is what was revealed for Mashhadī Naṣr Allāh.

\subsection{139}\label{section-2559}

\begin{otherlanguage}{arabic}

بسمی المهیمن علی الأسمآء

\end{otherlanguage}

Bi-smī al-Muhaymin ʿalā al-asmāʾ.

In My name, the Guardian over all names.

\subsection{140}\label{section-2560}

\begin{otherlanguage}{arabic}

یا نصراللّه یذکرک مولی الوری اذ کان مسجوناً فی هذا المقام البعید و یبشّرک
بفضل اللّه و رحمته انّه ما من اله الّا هو یفعل ما یشآء و یحکم ما یرید ان
اشهد بما شهد اللّه انّه لا اله الّا انا الفرد الخبیر قد ذکرت لدی المظلوم و
نزل لک هذه الآیات الّتی لا یعادلها ما خلق فی الأرض یشهد بذلک ربّک العلیم
طوبی لک بما نبذت شبهات القوم و شربت من بحر الیقین ان اعرف مقام من هداک
ثمّ اشکر ربّک المؤیّد الکریم ایّاک ان تحزنک شئونات الخلق او تخوّفک الملوک و
السّلاطین ینبغی لکلّ من اقبل الی افقی ان یترک الأرض و ما علیها لهم و یتوجّه
الی فتح القلوب باسمی العزیز الحکیم انّ الأموال لأهل الجلال و القلوب للّه
ربّ العالمین

\end{otherlanguage}

Yā Naṣr Allāh, yadhkuruka Mawlā al-warā idh kāna masjūnan fī hādhā
al-maqām al-baʿīd wa-yubashshiruka bi-faḍli Allāh wa-raḥmatihi. Innahu
mā min ilāhin illā Huwa, yafʿalu mā yashāʾ wa-yaḥkumu mā yurīd. An
ishhad bi-mā shahida Allāh: innahu lā ilāha illā anā al-Fard al-Khabīr.
Qad dhukirta ladā al-Maẓlūm wa-nazala laka hādhihi al-āyāt allatī lā
yuʿādiluhā mā khuliqa fī al-arḍ. Yashhadu bi-dhālika rabbuka al-ʿAlīm.
Ṭūbā laka bi-mā nabadhta shubuhāti al-qawm wa-sharibta min baḥr
al-yaqīn. An iʿraf maqāma man hadāka thumma ushkur rabbaka al-Muʾayyid
al-Karīm. Iyyāka an tuḥzinaka shuʾūnu al-khalq aw tukhawwifaka al-mulūk
wa-al-salāṭīn. Yanbaghī li-kulli man aqbala ilā ufuqī an yatruka al-arḍ
wa-mā ʿalayhā lahum wa-yatawajjaha ilā fatḥ al-qulūb bi-smī al-ʿAzīz
al-Ḥakīm. Inna al-amwāla li-ahl al-jalāl wa-al-qulūba li-llāhi rabb
al-ʿālamīn.

O Naṣr Allāh, the Lord of humanity remembers you while imprisoned in
this remote place, and gives you the glad tidings of God's grace and
mercy. There is no god but Him. He does what He wills and decrees what
He desires. Bear witness to what God has testified: ``There is no god
but Me, the Unique, the All-Aware.'' You have been mentioned in the
presence of the Wronged One, and these verses have been revealed for
you, verses that nothing created on earth can equal. Your Lord, the
All-Knowing, bears witness to this. Blessed are you for having cast
aside the people's doubts and drunk from the ocean of certitude.
Recognize the station of Him who guided you, then thank your Lord, the
Sustainer, the Generous. Beware lest the affairs of the people grieve
you or kings and sovereigns frighten you. Everyone who has turned toward
My horizon should leave the earth and all that is upon it to them, and
set about winning hearts through My name, the Almighty, the All-Wise.
Wealth belongs to the people of worldly grandeur, and hearts belong to
God, the Lord of the worlds.

\subsection{141}\label{section-2561}

\begin{otherlanguage}{persian}

و نذکر فیهذا المقام من سمّی بأبی‌القاسم الّذی هاجر الی ان فاز بلقآء ربّه
العزیز الکریم و نوصیه بالاستقامة علی هذا الأمر الأعظم و ندعو له ما ینبغی
لفضل اللّه العلیّ العظیم البهآء علی اهل البهآء الّذین فازوا بذکری الجمیل
انتهی

\end{otherlanguage}

Wa-nadhkuru fī hādhā al-maqām man summiya bi-Abī al-Qāsim alladhī hājara
ilā an fāza bi-liqāʾi rabbihi al-ʿAzīz al-Karīm, wa-nūṣīhi
bi-al-istiqāmati ʿalā hādhā al-Amr al-Aʿẓam wa-nadʿū lahu mā yanbaghī
li-faḍli Allāh al-ʿAliyy al-ʿAẓīm. Al-Bahāʾ ʿalā ahl al-Bahāʾ alladhīna
fāzū bi-dhikrī al-jamīl. Intahā.

And here We remember the one named Abū al-Qāsim, who emigrated until he
attained the presence of his Lord, the Almighty, the Generous. We
counsel him to remain steadfast in this Most Great Cause, and pray for
him as befits the grace of God, the Most High, the Magnificent. Glory
rest upon the people of Bahāʾ who have attained My gracious remembrance.
End.

\subsection{142}\label{section-2562}

\begin{otherlanguage}{persian}

این خادم فانی از رحمت و عنایت مقصود عالمیان امیدوار است که ید عنایتش کلّ
را اخذ نماید و از جحیم سفلی بجنّت علیا کشاند انّه لا یعجزه من شیء و لا
یعزب عن علمه من امر یعطی بالفضل و یمنع بالحکمة و هو المقتدر الآمر المرید

\end{otherlanguage}

Īn khādem-e fānī az raḥmat va ʿenāyat-e Maqṣūd-e ʿālamiyān omīdvār ast
ke yad-e ʿenāyatash koll rā akhdh namāyad va az jaḥīm-e soflā be
jannat-e ʿolyā kashānad. Innahu lā yuʿjizuhu min shayʾ wa-lā yaʿzubu ʿan
ʿilmihi min amr. Yuʿṭī bi-al-faḍl wa-yamnaʿu bi-al-ḥikmah wa-Huwa
al-Muqtadir al-Āmir al-Murīd.

This ephemeral servant hopes, through the mercy and loving-kindness of
the Desire of all the worlds, that the hand of His loving-kindness will
take hold of everyone and draw them from the lowest hell to the highest
paradise. Nothing can frustrate Him, and no matter escapes His
knowledge. He gives through grace and withholds through wisdom. He is
the All-Powerful, the Ordainer, the One who wills.

\subsection{143}\label{section-2563}

\begin{otherlanguage}{persian}

و اینکه در بارهٴ ملّا محمّد علیه ٦٦٩ {[}بهآء اللّه{]} از اهل یا مرقوم
داشتید و هم‌چنین ارتقاء ایشان بمرقاة البیان الی سماء عرفان ربّنا الرّحمن
عرض شد اینکلمات عالیات از مظهر علم الهی و مصدر امر ربّانی و فی‌الحقیقه
مقام لا یعرف و لا یذکر در بارهٴ ایشان ظاهر و نازل قوله تبارک و تعالی

\end{otherlanguage}

Va īnke dar bāre-ye Mollā Moḥammad, ʿalayhi 669 {[}Bahāʾ Allāh{]}, az
ahl-e Yā marqūm dāshtīd va ham-chenīn erteqā-ye īshān be merqāt al-bayān
ilā samāʾi ʿirfāni Rabbinā al-Raḥmān ʿarż shod. Īn kalemāt-e ʿāliyāt az
Maẓhar-e ʿElm-e Elāhī va Maṣdar-e Amr-e Rabbānī va fī al-ḥaqīqeh Maqām-e
lā yuʿraf wa-lā yudhkar dar bāre-ye īshān ẓāher va nāzel. Qawluhu
tabāraka wa-taʿālā:

What you wrote concerning Mullā Muḥammad, upon him be 669 {[}the glory
of God{]}, a native of Yā, and likewise his ascent by the ladder of
utterance to the heaven of recognition of our Lord, the All-Merciful,
was presented. These exalted words concerning him appeared and were
revealed from the Manifestation of divine knowledge and the Source of
the Lord's command, and, in truth, from the Station that can neither be
known nor mentioned. His words, blessed and exalted is He:

\subsection{144}\label{section-2564}

\begin{otherlanguage}{arabic}

بنام گویندهٴ پاینده

\end{otherlanguage}

Be nām-e Gūyande-ye Pāyandeh.

In the name of the Everlasting Speaker.

\subsection{145}\label{section-2565}

\begin{otherlanguage}{persian}

یا محمّد از حقّ میطلبیم ترا مؤیّد فرماید بر عرفان و ایقان و خدمت امرش
بشأنیکه اگر اهل عالم بر اعراض قیام نمایند قادر بر منع آنجناب نباشند
امروز روح اللّه در برّیّهٴ ارض مقدّسه ندا مینماید و نور اللّه از افق اراده
مشرق و عرف قمیص متضوّع و نسمات وحی در مرور مع‌ذلک عباد در ضلالت خود باقی و
برقرار جهد نما تا از کوثر عرفان رحمن بیاشامی و از مائدهٴ منزله قسمت بری
و بحکمت و بیان اهل امکانرا بافق اعلی هدایت نمائی ناس غافلند مبیّن لازم بر
خدمت امر حقّ قیام نما تا مخدوم عالم شوی و بنار حبّش مشتعل شو تا محبوب آفاق
گردی اگرچه کشف بعضی از اسرار شده ولکن اکثری از آن در علم الهی مستور و
مکنون چه اگر اسرار یوم اللّه و مقامات نفوس موقنهٴ مطمئنّه اقلّ از سمّ ابره
ظاهر شود جمیع ناسرا طائف مطاف ماکان و مایکون مشاهده نمائی باسم حقّ بایست
و بتدارک مافات قیام نما وقت بسیار عزیز است بل اعزّ از کبریت احمر ناسرا
باخلاق روحانیّه و اعمال طیّبه وصیّت نما طوبی از برای نفوسیکه للّه عاملند و
الی اللّه متوجّه

\end{otherlanguage}

Yā Moḥammad, az Ḥaqq mīṭalabīm to rā moʾayyad farmāyad bar ʿerfān va
īqān va khedmat-e Amrash be shaʾnī ke agar ahl-e ʿālam bar eʿrāż qīyām
namāyand qāder bar manʿ-e ān janāb nabāshand. Emrūz Rūḥ Allāh dar
barriyye-ye Arż-e Moqaddaseh nedā mīnamāyad va Nūr Allāh az ofoq-e
erādeh moshreq va ʿarf-e qamīṣ motaḍavveʿ va nasamāt-e vaḥy dar morūr.
Maʿa dhālek ʿebād dar żalālat-e khod bāqī va barqarār. Jahd namā tā az
kowthar-e ʿerfān-e Raḥmān biyāshāmī va az māʾede-ye monazzaleh qesmat
barī va be ḥekmat va bayān ahl-e emkān rā be Ofoq-e Aʿlā hedāyat namāʾī.
Nās ghāfeland, mobayyen lāzem. Bar khedmat-e Amr-e Ḥaqq qīyām namā tā
makhdūm-e ʿālam shavī va be nār-e ḥobbash moshtaʿel show tā maḥbūb-e
āfāq gardī. Agarche kashf-e baʿżī az asrār shodeh va-lāken aktharī az ān
dar ʿelm-e Elāhī mastūr va maknūn, che agar asrār-e Yawm Allāh va
maqāmāt-e nofūs-e mūqene-ye moṭmaʾenneh aqall az samm-e ebreh ẓāher
shavad jamīʿ-e nās rā ṭāʾef-e maṭāf-e mā kān va mā yakūn moshāhedeh
namāʾī. Be nām-e Ḥaqq beʾīst va be tadārok-e mā fāt qīyām namā. Vaqt
besyār ʿazīz ast bal aʿazz az kebrīt-e aḥmar. Nās rā be akhlāq-e
rūḥāniyyeh va aʿmāl-e ṭayyebeh vaṣiyyat namā. Ṭūbā az barā-ye nofūsī ke
li-llāh ʿāmeland va ilā Allāh motavajjeh.

O Muḥammad, We ask the True One to strengthen you in recognition,
certitude, and service to His Cause, so that even if all the people of
the world should rise up in rejection, they would be powerless to hinder
you. Today the Spirit of God calls out in the wilderness of the Holy
Land, the Light of God shines from the horizon of His will, the
fragrance of the garment is diffused, and the breezes of revelation are
passing. Yet the servants remain fixed in their error. Strive to drink
from the Kawthar of recognition of the All-Merciful, to receive a share
of the table sent down from heaven, and to guide the inhabitants of the
contingent world to the Most Exalted Horizon through wisdom and
utterance. People are heedless; an expounder is needed. Arise to serve
the Cause of the True One, so that you may become one whom the world
serves, and be kindled with the fire of His love, so that you may become
beloved throughout the world. Although some mysteries have been
disclosed, most remain concealed and hidden in the knowledge of God. For
if less than a needle's eye of the mysteries of the Day of God and the
stations of assured souls possessed of certitude were revealed, you
would see all people circling the One around whom all that has been and
will be revolves. Stand firm in the name of the True One and arise to
make good what has been lost. Time is exceedingly precious, indeed more
precious than red sulphur. Counsel people to adopt spiritual qualities
and perform good deeds. Blessed are those souls who act for God and turn
toward God.

\subsection{146}\label{section-2566}

\begin{otherlanguage}{arabic}

قل یا ملأ الأرض ایّاکم ان تمنعکم کتب العالم عن کتاب اللّه المهیمن القیّوم
الّذی نزّل من سمآء الفضل من لدی الحقّ علّام الغیوب قل تاللّه لا ینفعکم ما عند
النّاس اتّقوا الرّحمن ثمّ اقبلوا الیه بوجوه بیضآء و قلوب نورآء

\end{otherlanguage}

Qul yā malaʾ al-arḍ, iyyākum an tamnaʿakum kutubu al-ʿālam ʿan Kitābi
Allāh al-Muhaymin al-Qayyūm alladhī nuzzila min samāʾ al-faḍl min ladā
al-Ḥaqq ʿAllām al-ghuyūb. Qul ta-Allāh, lā yanfaʿukum mā ʿinda al-nās.
Ittaqū al-Raḥmān thumma aqbilū ilayhi bi-wujūhin bayḍāʾ wa-qulūbin
nawrāʾ.

Say: O peoples of the earth, beware lest the books of the world keep you
from the Book of God, the All-Protecting, the Self-Subsisting, which has
been sent down from the heaven of grace by the True One, the Knower of
the unseen. Say: By God, what people possess will not benefit you. Be
mindful of the All-Merciful, then turn toward Him with radiant faces and
luminous hearts.

\subsection{147}\label{section-2567}

\begin{otherlanguage}{persian}

کذلک یأمرکم من فی قبضته زمام العلوم قل یا معشر العلمآء طهّروا قلوبکم عمّا
عندکم لتسمعوا صریر القلم الأعلی و حفیف سدرة المنتهی من المقام المحمود
هذا یوم لا تغنیکم کنوز العالم و لا صفوف الأمم دعوها عن ورائکم رجآء ما
عند اللّه ربّ ما کان و ما یکون انّا نوصیک و الّذین آمنوا بالاستقامة الکبری
علی هذا الأمر الّذی به اضطربت ارکان الأسمآء الّا من شآء اللّه مالک الوجود
کذلک ماج بحر بیانی و هاج عرف عنایتی انّک اذا رأیت و وجدت قم و قل اشهد انّک
قبلت اذی العالم لارتفاع امر اللّه و حملت البلایا فی حبّ اللّه فاطر السّمآء و
مالک الأسمآء الی ان قبلت السّجن الأعظم لنجاة الأمم اسئلک بالاسم الأعظم
الّذی به سخّرت ملکوت الأسمآء بأن لا تخیّبنی عمّا عندک و لا تمنعنی عمّا قدّرته
لأمنائک و اصفیائک ایربّ ترانی مقبلاً الیک و متمسّکاً بحبلک اسئلک بأن تؤیّدنی
علی الاستقامة علی هذا الأمر الّذی به زلّت اقدام اکثر خلقک لا اله الّا انت
المتعالی المشفق الغفور الکریم انتهی

\end{otherlanguage}

Kadhālika yaʾmurukum man fī qabḍatihi zimāmu al-ʿulūm. Qul yā maʿshar
al-ʿulamāʾ, ṭahhirū qulūbakum ʿammā ʿindakum li-tasmaʿū ṣarīra al-Qalam
al-Aʿlā wa-ḥafīfa Sidrat al-Muntahā min al-maqām al-maḥmūd. Hādhā yawmun
lā tughnīkum kunūzu al-ʿālam wa-lā ṣufūfu al-umam. Daʿūhā ʿan warāʾikum
rajāʾa mā ʿinda Allāhi rabbi mā kāna wa-mā yakūn. Innā nūṣīka
wa-alladhīna āmanū bi-al-istiqāmati al-kubrā ʿalā hādhā al-Amr alladhī
bihi iḍṭarabat arkānu al-asmāʾ illā man shāʾa Allāhu Mālik al-wujūd.
Kadhālika māja baḥru bayānī wa-hāja ʿarfu ʿināyatī. Innaka idhā raʾayta
wa-wajadta qum wa-qul: Ashhadu annaka qabilta adhā al-ʿālam li-irtifāʿi
Amri Allāh wa-ḥamalta al-balāyā fī ḥubbi Allāhi Fāṭir al-samāʾ wa-Mālik
al-asmāʾ ilā an qabilta al-Sijn al-Aʿẓam li-najāti al-umam. Asʾaluka
bi-al-Ism al-Aʿẓam alladhī bihi sakhkharta malakūt al-asmāʾ bi-an lā
tukhayyibanī ʿammā ʿindaka wa-lā tamnaʿanī ʿammā qaddartahu li-umanāʾika
wa-aṣfiyāʾika. Ay Rabb, tarānī muqbilan ilayka wa-mutamassikan
bi-ḥablika. Asʾaluka bi-an tuʾayyidanī ʿalā al-istiqāmati ʿalā hādhā
al-Amr alladhī bihi zallat aqdāmu akthari khalqika. Lā ilāha illā anta
al-Mutaʿālī al-Mushfiq al-Ghafūr al-Karīm. Intahā.

Thus commands you the One who holds the reins of knowledge in His grasp.
Say: O company of religious scholars, cleanse your hearts of what you
possess, so that you may hear the scratching of the Most Exalted Pen and
the rustling of the Tree of the Utmost Boundary from the praised
station. This is a day when neither the treasures of the world nor the
ranks of the nations will avail you. Leave them behind you, hoping for
what is with God, the Lord of all that has been and will be. We counsel
you and those who have believed to show the greatest steadfastness in
this Cause, through which the pillars of the names have been shaken,
except those whom God, the Lord of existence, has willed. Thus has the
ocean of My utterance surged and the fragrance of My loving-kindness
spread. When you have seen and perceived, arise and say: ``I testify
that You accepted the world's afflictions for the exaltation of the
Cause of God, and bore calamities out of love for God, the Creator of
heaven and Lord of all names, until You accepted the Most Great Prison
for the deliverance of the nations. I ask You by the Most Great Name,
through which You have subdued the kingdom of names, not to disappoint
me in what is with You or deprive me of what You have ordained for Your
trusted and chosen ones. O Lord, You see me turning toward You and
holding fast to Your cord. I ask You to strengthen me in steadfastness
in this Cause, through which the feet of most of Your creatures have
slipped. There is no god but You, the Most Exalted, the Compassionate,
the Ever-Forgiving, the Generous.'' End.

\subsection{148}\label{section-2568}

\begin{otherlanguage}{persian}

انشآء اللّه مؤیّد شوند بر آنچه از لسان عظمت بآن مأمور شدند و از کؤوس عنایت
بیاشامند آشامیدنیکه از منع ظالمین بازنمانند فی‌الحقیقه اگر انسان در این
ایّام قسمت نبرد و از فیوضات فیّاض بی‌نصیب ماند بچه چیز دلخوش و مسرور و
مشغول بسیار حیف است این اوقات عزیزه بگذرد و انسان غافل باشد انشآءاللّه
باید جمیع دوستان همّت بر خدمت گمارند اوّل کلّ بکمال تقدیس و تنزیه و انصاف و
مشورت قیام نمایند آنچه را سبب تفریق مشاهده کنند اوّل باصلاح آن توجّه کنند
و بعد بروح و ریحان و حکمت بخدمت مشغول گردند اصلاح امور ایشان و عالم همه
خواهد شد و امثال این بیان مکرّر از لسان قدم استماع شده صدهزار نعیم و طوبی
از برای نفوسیکه قیام نمایند بر آنچه که سبب ارتفاع امر و راحت اهل عالم
است و فی‌الحقیقه نفس ارتفاع نفس راحت است ولکن ناس جاهل از اینفقره غافل
انشآءاللّه آنمحبوب و اینعبد دعا میکنیم که شاید از شمال جهل بیمین دانائی
فائز شوند

\end{otherlanguage}

En shāʾ Allāh moʾayyad shavand bar ānche az Lesān-e ʿAẓamat be ān maʾmūr
shodand va az koʾūs-e ʿenāyat biyāshāmand, āshāmīdanī ke az manʿ-e
ẓālemīn bāz namānand. Fī al-ḥaqīqeh agar ensān dar īn ayyām qesmat
nabarad va az foyūżāt-e Fayyāż bī-naṣīb mānad be che chīz delkhosh va
masrūr va mashghūl? Besyār ḥeyf ast īn owqāt-e ʿazīzeh bogdharad va
ensān ghāfel bāshad. En shāʾ Allāh bāyad jamīʿ-e dūstān hemmat bar
khedmat gomārand. Avval koll be kamāl-e taqdīs va tanzīh va enṣāf va
mashverat qīyām namāyand. Ānche rā sabab-e tafrīq moshāhedeh konand
avval be eṣlāḥ-e ān tavajjoh konand va baʿd be rūḥ va reyḥān va ḥekmat
be khedmat mashghūl gardand. Eṣlāḥ-e omūr-e īshān va ʿālam hameh khāhad
shod va amthāl-e īn bayān mokarrar az Lesān-e Qedam estemāʿ shodeh. Ṣad
hezār naʿīm va ṭūbā az barā-ye nofūsī ke qīyām namāyand bar ānche ke
sabab-e ertefāʿ-e Amr va rāḥat-e ahl-e ʿālam ast. Va fī al-ḥaqīqeh
nafs-e ertefāʿ nafs-e rāḥat ast va-lāken nās-e jāhel az īn faqareh
ghāfel. En shāʾ Allāh ān maḥbūb va īn ʿabd doʿā mīkonīm ke shāyad az
shemāl-e jahl be yamīn-e dānāʾī fāʾez shavand.

God willing, may they be strengthened to carry out what they were
commanded by the Tongue of Grandeur, and drink from the cups of
loving-kindness so deeply that the opposition of the oppressors cannot
hold them back. In truth, if a person receives no share in these days
and remains deprived of the outpourings of the All-Bountiful, what can
give that person contentment and joy, and what can occupy them? It is a
great pity for these precious times to pass while a person remains
heedless. God willing, all the friends should devote their efforts to
service. First, let all arise with the utmost purity, sanctity,
fairness, and consultation. Whatever they observe to be a cause of
division, let them first attend to remedying it, and then engage in
service with joy, gladness, and wisdom. Their affairs and those of the
whole world will be set right; statements like this have repeatedly been
heard from the Tongue of Eternity. A hundred thousand blessings and
gladness be upon those souls who arise to do what brings about the
exaltation of the Cause and the well-being of the people of the world.
In truth, that exaltation is itself well-being, but ignorant people are
heedless of this matter. God willing, you, beloved, and this servant
pray that they may pass from the left hand of ignorance to the right
hand of knowledge.

\subsection{149}\label{section-2569}

\begin{otherlanguage}{persian}

و اینکه در بارهٴ اخت حضرت سلطان الشّهدآء روح ما سواه فداه مرقوم داشتید
تلقآء عرش عرض شد هذا ما نطق به لسان العظمة یا علی لحاظ عنایت باو و
منتسبین نورین بوده و خواهد بود و به اسم‌اللّه مهدی امر نمودیم که در بارهٴ
حقوق او بجناب صادق علیه بهائی سفارش کامل نماید و هم‌چنین از قلم اعلی
مرقوم اگر اخت حقّی دارد باید ادا نمائید بحقّ و عدل رفتار شود انّا نکبّر من
هذا المقام علیها و نبشّرها بفضلی و رحمتی انتهی

\end{otherlanguage}

Va īnke dar bāre-ye okht-e Ḥażrat-e Solṭān al-Shohadāʾ, rūḥ-e mā sevāh
fedāh, marqūm dāshtīd telqāʾ-e ʿArsh ʿarż shod. Hādhā mā naṭaqa bihi
Lisān al-ʿAẓamah: Yā ʿAlī, leḥāẓ-e ʿenāyat be ū va montasebīn-e Nūrayn
būdeh va khāhad būd va be Esm Allāh Mahdī amr namūdīm ke dar bāre-ye
ḥoqūq-e ū be Janāb-e Ṣādeq, ʿalayhi Bahāʾī, sefāresh-e kāmel namāyad. Va
ham-chenīn az Qalam-e Aʿlā marqūm: agar okht ḥaqqī dārad bāyad adā
namāʾīd. Be ḥaqq va ʿadl raftār shavad. Innā nukabbiru min hādhā
al-maqām ʿalayhā wa-nubashshiruhā bi-faḍlī wa-raḥmatī. Intahā.

What you wrote concerning the sister of the King of Martyrs---may the
souls of all others be sacrificed for him---was presented before the
Throne. This is what the Tongue of Grandeur uttered: ``O ʿAlī, the gaze
of loving-kindness has rested, and will continue to rest, upon her and
the relatives of the Twin Lights. We have instructed the Name of God,
Mahdī, to give full directions to Ṣādiq, upon him be My glory,
concerning her rights. It has likewise been written by the Most Exalted
Pen: `If the sister has any entitlement, you must discharge it.' Let
conduct accord with truth and justice. From this station We send her Our
greeting and give her the glad tidings of My grace and mercy.'' End.

\subsection{150}\label{section-2570}

\begin{otherlanguage}{persian}

و امّا توجّه ایشان بشطر ص و قیام ورقه‌ئی از اوراق علی ما یحبّه اللّه و یرضی
اینفقره بشرف اصغای مالک اسماء فائز فرمودند یا عبد حاضر طوبی لها ثمّ طوبی
لها ینبغی ان ینزل علیها من سمآء عنایتی ما یبقی به ذکرها جزآء عملها و
یکتب لها قلمی الأعلی و یبشّرها بقبول ما ظهر من عندها و نأمر الغصن بأن
یکتب لأمتی و ورقتی ما نزّل من سمآء مشیّة موجدها و ظهر من افق عنایة بارئها
لیکون ذخراً و شرفاً لها فیکلّ عالم من عوالم ربّها انتهی

\end{otherlanguage}

Va ammā tavajjoh-e īshān be shaṭr-e Ṣād va qiyām-e varaqeʾī az avrāq
ʿalā mā yuḥibbuhu Allāhu va yarḍā, īn feqre be sharaf-e iṣghā-ye Mālek-e
asmāʾ fāʾez. Farmūdand: Yā ʿabda ḥāḍir, ṭūbā lahā thumma ṭūbā lahā.
Yanbaghī an yanzila ʿalayhā min samāʾi ʿināyatī mā yabqā bihi dhikruhā
jazāʾa ʿamalihā, va yaktuba lahā qalamī al-aʿlā va yubashshiruhā
bi-qabūli mā ẓahara min ʿindihā. Va naʾmuru al-ghuṣna bi-an yaktuba
li-amatī va varaqatī mā nuzzila min samāʾi mashīʾati mūjidihā va ẓahara
min ufuqi ʿināyati bāriʾihā li-yakūna dhukhran va sharafan lahā fī kulli
ʿālamin min ʿavālimi rabbihā. Intahā.

As for their journey toward Ṣād and the action of one of the leaves in
accordance with what God loves and approves, this matter attained the
honor of being heard by the Lord of Names. He said: ``O servant in
attendance! Blessed is she, and again, blessed is she! It is fitting
that there descend upon her from the heaven of My loving-kindness that
which will perpetuate her remembrance as a recompense for her deed, and
that My Most Exalted Pen write to her and give her the glad tidings that
what she has done has been accepted. We command the Branch to write to
My handmaiden and My leaf what has descended from the heaven of her
Creator's will and appeared from the horizon of her Maker's
loving-kindness, that it may be a treasure and an honor for her in every
world among the worlds of her Lord.'' End.

\subsection{151}\label{section-2571}

\begin{otherlanguage}{arabic}

پرتو آفتاب عنایت عوالم غیب و شهود را منوّر نمود یا محبوب فؤادی مشاهده
نمائید عمل مرغوب چه مقدار محبوبست عنداللّه اگر بحور اکسیر و کنوز عالم را
احدی انفاق کند البتّه باین ذکر فائز نشود عملی ظاهر و بطراز قبول فائز
بشأنیکه از مطلع بیان رحمن ظاهر آنچه که نفحاتش بدوام ملک و ملکوت منقطع
نشود از جانب اینعبد تکبیر و سلام بآنورقه برسانید علیها بهآء ربّها علیها
بهآء ربّها و هذا ما نزّل لها من سمآء عنایة ربّنا الغفور الکریم

\end{otherlanguage}

Partow-e āftāb-e ʿenāyat ʿavālem-e ghayb va shohūd rā monavvar namūd. Yā
maḥbūb-e foʾādī, moshāhede namāʾīd ʿamal-e marghūb che meqdār maḥbūb ast
ʿinda Allāh. Agar boḥūr-e eksīr va konūz-e ʿālam rā aḥadī enfāq konad
albatte be īn dhekr fāʾez nashavad. ʿAmalī ẓāher va be ṭerāz-e qabūl
fāʾez be shaʾnī ke az maṭlaʿ-e bayān-e Raḥmān ẓāher ānche ke nafaḥātesh
be davām-e molk va malakūt monqaṭeʿ nashavad. Az jāneb-e īn ʿabd takbīr
va salām be ān varaqe berasānīd. ʿAlayhā bahāʾu rabbihā, ʿalayhā bahāʾu
rabbihā. Va hādhā mā nuzzila lahā min samāʾi ʿināyati rabbinā al-ghafūri
al-karīm.

The radiance of the sun of loving-kindness has illumined the worlds of
the unseen and the seen. O beloved of my heart, observe how dearly a
worthy deed is loved by God. Were anyone to give away oceans of elixir
and the treasures of the world, they would assuredly not attain such a
remembrance. A deed was performed and adorned with acceptance, in such a
manner that there appeared from the dawning-place of the All-Merciful's
utterance that whose fragrances will never cease for as long as the
earthly and heavenly kingdoms endure. Convey to that leaf greetings of
praise and peace on behalf of this servant. Upon her be the glory of her
Lord; upon her be the glory of her Lord. And this is what has been sent
down for her from the heaven of the loving-kindness of our Lord, the
Ever-Forgiving, the All-Bountiful.

\subsection{152}\label{section-2572}

\begin{otherlanguage}{arabic}

بسمی الظّاهر المشهود

\end{otherlanguage}

Bi-smī al-ẓāhiri al-mashhūd.

In My name, the Manifest, the Witnessed.

\subsection{153}\label{section-2573}

\begin{otherlanguage}{arabic}

یا ورقتی علیک بهائی ان افرحی بما نزّل لک من القلم الأعلی انّه ذکرک و بشّرک
بما ظهر منک فی سبیل اللّه المهیمن القیّوم ان اشکری و تشبّثی بذیل العنایة و
قولی لک الحمد یا اله الوجود و مالک الغیب و الشّهود بما ایّدتنی علی عمل فاز
بطراز قبولک و تزیّن برضائک لعمرک یا محبوب قلبی لا یعادل ما اعطیتنی بما
خلق فی سمائک و ارضک انّک انت المعطی الکریم انتهی

\end{otherlanguage}

Yā varaqatī, ʿalayki bahāʾī. An ifraḥī bi-mā nuzzila laki mina al-qalami
al-aʿlā. Innahu dhakaraki va bashsharaki bi-mā ẓahara minki fī sabīli
Allāhi al-muhaymini al-qayyūm. An ishkurī va tashabbathī bi-dhayli
al-ʿināyati va qūlī: Laka al-ḥamdu yā ilāha al-vujūdi va mālika
al-ghaybi va al-shuhūd, bi-mā ayyadtanī ʿalā ʿamalin fāza bi-ṭirāzi
qabūlika va tazayyana bi-riḍāka. La-ʿamrika yā maḥbūba qalbī, lā
yuʿādilu mā aʿṭaytanī bi-mā khuliqa fī samāʾika va arḍika. Innaka anta
al-muʿṭī al-karīm. Intahā.

O My leaf, upon you be My glory! Rejoice in what has been sent down for
you from the Most Exalted Pen. It has remembered you and given you glad
tidings of what you have done in the path of God, the All-Protecting,
the Self-Subsisting. Give thanks, hold fast to the hem of
loving-kindness, and say: ``Praise be to You, O God of existence and
Lord of the unseen and the seen, for enabling me to perform a deed that
has been adorned with Your acceptance and embellished with Your good
pleasure. By Your life, O Beloved of my heart, nothing created in Your
heaven and on Your earth can equal what You have bestowed upon me. You
are indeed the Bestower, the All-Bountiful.'' End.

\subsection{154}\label{section-2574}

\begin{otherlanguage}{persian}

حقّ جلّ جلاله جمیعرا مؤیّد فرماید بر آنچه نفحات رضایش از او متضوّع یک عمل
مبرور الیوم معادلست با صدهزار عمل بل استغفر اللّه من هذا التّحدید و هم‌چنین
عمل غیر خیر وقتی از اوقات لسان عظمت باینکلمهٴ علیا ناطق قوله جلّ و عزّ یا
عبد حاضر عنایات حقّ خلقرا جسور نموده بشأنیکه ذکرش محبوب نه باید نفوس
مقبله را آگاه نمود الیوم چون آفتاب ظهور مشرقست و پرتو انوار وجه عالم را
احاطه نموده لذا اکثری از اعمال نالایقه ستر شده و میشود انّه هو السّتّار
الحکیم

\end{otherlanguage}

Ḥaqq jalla jalāluhu jamīʿ rā moʾayyad farmāyad bar ānche nafaḥāt-e
reḍāyash az ū motaḍavveʿ. Yek ʿamal-e mabrūr al-yowm moʿādel ast bā
ṣad-hezār ʿamal, bal astaghfiru Allāha min hādhā al-taḥdīd, va
ham-chenīn ʿamal-e ghayr-e khayr. Vaqtī az owqāt lesān-e ʿaẓamat be īn
kalame-ye ʿolyā nāṭeq, qawluhu jalla va ʿazza: Yā ʿabd-e ḥāḍer,
ʿenāyāt-e Ḥaqq khalq rā jasūr namūde be shaʾnī ke dhekrash maḥbūb na.
Bāyad nofūs-e moqbele rā āgāh namūd. Al-yowm chūn āftāb-e ẓohūr moshreq
ast va partow-e anvār-e vajh ʿālam rā eḥāṭe namūde, ledhā aktharī az
aʿmāl-e nālāyeqe satr shode va mīshavad. Innahu huwa al-sattāru
al-ḥakīm.

May God, exalted be His glory, assist everyone to do that from which the
fragrances of His good pleasure spread. One righteous deed today is
equal to a hundred thousand deeds---indeed, I ask God's forgiveness for
setting such a limit---and likewise with an evil deed. On one occasion
the Tongue of Grandeur spoke these exalted words---His words, glorified
and exalted be He: ``O servant in attendance! God's favors have
emboldened people to a degree that it is not pleasing to mention. Those
who have turned toward Him must be made aware. Today, because the sun of
Revelation is shining and the radiance of the lights of His countenance
has encompassed the world, most unseemly deeds have been and are being
concealed. He is indeed the Concealer, the All-Wise.''

\subsection{155}\label{section-2575}

\begin{otherlanguage}{persian}

امروز روز شادیست و روز فرح اکبر است بحر عنایت موّاج و نسیم جود در مرور و
آفتاب فضل مشرق ولکن حال مشاهده میشود که اعمال شنیعه بمقامی رسیده که
نزدیک بآنست غبارش بذیل حقّ جلّ جلاله رسد و حجاب حرمت را بردرد جمیع آنچه از
اوّل امر تا حال وارد شد و قبول نمودیم لأجل اعزاز کلمهٴ الهی و ارتقاء نفوس
مطمئنّهٴ مستقیمهٴ راضیهٴ مرضیّه بود حال از اعمال بعضی بباد رفته و میرود
مگر اینکه حقّ جلّ جلاله بقدرت کامله حفظ نماید چنانچه نمود

\end{otherlanguage}

Emrūz rūz-e shādī ast va rūz-e faraḥ-e akbar ast. Baḥr-e ʿenāyat mavvāj
va nasīm-e jūd dar morūr va āftāb-e faḍl moshreq. Va lāken ḥāl moshāhede
mīshavad ke aʿmāl-e shanīʿe be maqāmī rasīde ke nazdīk be ān ast
ghobārash be dhayl-e Ḥaqq jalla jalāluhu rasad va ḥejāb-e ḥormat rā
bardarad. Jamīʿ-e ānche az avval-e amr tā ḥāl vāred shod va qabūl
namūdīm li-ajl-e eʿzāz-e kalame-ye elāhī va erteqā-ye nofūs-e
moṭmaʾenne-ye mostaqīme-ye rāḍīye-ye marḍīye būd. Ḥāl az aʿmāl-e baʿḍī
be bād rafte va mīravad magar īnke Ḥaqq jalla jalāluhu be qodrat-e
kāmele ḥefẓ namāyad chenānche namūd.

``Today is a day of joy and a day of supreme gladness. The ocean of
loving-kindness is surging, the breeze of bounty is blowing, and the sun
of grace is shining. Yet now it is seen that reprehensible deeds have
reached such a point that their dust is close to reaching the hem of
God, exalted be His glory, and tearing the veil of sanctity. Everything
that has befallen Us from the beginning of the Cause until now, and that
We have accepted, has been for the exaltation of the Word of God and the
advancement of souls who are assured, steadfast, content, and
well-pleasing to Him. Now, through the deeds of some, this has been and
is being cast to the winds, unless God, exalted be His glory, preserves
it through His consummate power, as He has done.''

\subsection{156}\label{section-2576}

\begin{otherlanguage}{persian}

یا عبد حاضر بجناب علی علیه بهائی و عنایتی بنویس باید جمیع دوستان آن ارض
را متذکّر دارد و بما ینفعهم آگاه نماید شاید از مافات عنهم رجوع نمایند و
بحقّ تمسّک جویند و شاید آنچه سبب ضرّ گشته و علّت ذلّت نفوس شده بماء رجوع طاهر
و مقدّس گردد انّه هو التّواب و هو الفضّال الکریم انتهی

\end{otherlanguage}

Yā ʿabd-e ḥāḍer, be Janāb-e ʿAlī ʿalayhi bahāʾī va ʿināyatī benevīs:
Bāyad jamīʿ-e dūstān-e ān arḍ rā motadhakker dārad va bi-mā yanfaʿuhum
āgāh namāyad, shāyad az mā fāta ʿanhum rojūʿ namāyand va be Ḥaqq
tamassok jūyand, va shāyad ānche sabab-e ḍorr gashte va ʿellat-e
dhellat-e nofūs shode be māʾ-e rojūʿ ṭāher va moqaddas gardad. Innahu
huwa al-tavvābu va huwa al-faḍḍālu al-karīm. Intahā.

``O servant in attendance! Write to ʿAlī, upon him be My glory and My
loving-kindness: He must admonish all the friends in that land and
acquaint them with what will benefit them, so that they may perhaps turn
back from what has escaped them and hold fast to God, and that whatever
has caused harm and brought humiliation upon souls may be cleansed and
purified by the water of return. He is indeed the Ever-Relenting, the
Most Bountiful, the All-Generous.'' End.

\subsection{157}\label{section-2577}

\begin{otherlanguage}{persian}

مکرّر این بیانات از لسان منزل آیات جاری و بعضی از اوقات اینعبد وارد و
جمال قدم را در کدورت و حزنی مشاهده مینمود که لرزه اندام اینفانی را اخذ
مینمود و حزن جمال قدم معلوم و واضحست که در چه بوده و هست احتیاج بذکر
ندارد امروز اعمال طیّبهٴ دوستان الهی منادی امرند بین عباد از این قبیل در
الواح الهی مکرّر نازل شاید نفع بخشد و سبب تنبّه گردد و الّا انّه لغنیّ عن
العالمین

\end{otherlanguage}

Mokarrar īn bayānāt az lesān-e Monzel-e āyāt jārī, va baʿḍī az owqāt īn
ʿabd vāred va Jamāl-e Qedam rā dar kodūrat va ḥoznī moshāhede mīnamūd ke
larze andām-e īn fānī rā akhdh mīnamūd. Va ḥozn-e Jamāl-e Qedam maʿlūm
va vāḍeḥ ast ke dar che būde va hast, eḥtīāj be dhekr nadārad. Emrūz
aʿmāl-e ṭayyebe-ye dūstān-e elāhī monādī-ye amr and bayn-e ʿebād. Az īn
qabīl dar alvāḥ-e elāhī mokarrar nāzel, shāyad nafʿ bakhshad va sabab-e
tanabboh gardad. Va ellā innahu la-ghaniyyun ʿani al-ʿālamīn.

These statements have repeatedly flowed from the tongue of the Revealer
of verses. At times this servant would enter and behold the Ancient
Beauty in such distress and sorrow that this mortal frame would tremble.
It is clear and well known what has been and remains the cause of the
Ancient Beauty's sorrow; it needs no mention. Today the pure deeds of
the friends of God proclaim the Cause among His servants. Statements of
this kind have repeatedly been revealed in the divine Tablets, that they
may perhaps prove beneficial and awaken people. Otherwise, He is indeed
independent of all the worlds.

\subsection{158}\label{section-2578}

\begin{otherlanguage}{persian}

چندی قبل بعضی از نفوس بر حسب ظاهر شکایت آغاز نمودند مع‌ذلک از ساحت اقدس
جواب نازل نه فرمودند ستری سبق الکشف و صبری سبق العجل انتهی ولکن بیم آنست
تجاوزات ما حجاب را بردارد و ستر را خرق نماید اتشبّث بذیل عنایته و اسئله
بأن یؤیّد احبائه علی ما یحبّ و یرضی و یمنعهم عمّا یضرّهم و یهدیهم الی ما
ینفعهم انّه هو المقتدر المهیمن العطوف الغفور

\end{otherlanguage}

Chandī qabl baʿḍī az nofūs bar ḥasab-e ẓāher shekāyat āghāz namūdand.
Maʿa dhālek az sāḥat-e aqdas javāb nāzel na. Farmūdand: Satrī sabaqa
al-kashfa va ṣabrī sabaqa al-ʿajala. Intahā. Va lāken bīm ān ast
tajāvozāt-e mā ḥejāb rā bardārad va satr rā kharq namāyad. Atashabbathu
bi-dhayli ʿināyatihi va asʾaluhu bi-an yuʾayyida aḥibbāʾahu ʿalā mā
yuḥibbu va yarḍā va yamnaʿahum ʿammā yaḍurruhum va yahdiyahum ilā mā
yanfaʿuhum. Innahu huwa al-muqtadiru al-muhayminu al-ʿaṭūfu al-ghafūr.

Some time ago, certain souls began to complain outwardly. Nevertheless,
no answer was sent down from the Most Holy Court. He said: ``My
concealment has preceded disclosure, and My patience has preceded
haste.'' End. Yet there is a fear that our transgressions may lift the
veil and rend the covering. I cling to the hem of His loving-kindness
and ask Him to assist His loved ones to do what He loves and approves,
to keep them from what harms them, and to guide them to what benefits
them. He is indeed the All-Powerful, the All-Protecting, the
All-Compassionate, the Ever-Forgiving.

\subsection{159}\label{section-2579}

\begin{otherlanguage}{persian}

حسب الامر مهاجرین ارض صاد علیهم بهآء اللّه را از قبل حقّ تکبیر برسانید
انشآءاللّه بنور اتّفاق منوّر باشند و بر خدمت امر قائم تا اثر نور در آفاق
اشراق نماید و ظاهر شود جمیع فی‌الجمله برضای حقّ پی‌برده‌اند بآن تمسّک نمایند
و بآن عامل شوند اینست سبب رستگاری در دنیا و آخرت هجرت و بلایای شما لازال
تحت لحاظ بوده انشآءاللّه محفوظ ماند انتهی

\end{otherlanguage}

Ḥasab al-amr mohājerīn-e arḍ-e Ṣād ʿalayhim bahāʾu Allāh rā az qebal-e
Ḥaqq takbīr berasānīd. In shāʾ Allāh be nūr-e ettefāq monavvar bāshand
va bar khedmat-e amr qāʾem tā athar-e nūr dar āfāq eshrāq namāyad va
ẓāher shavad. Jamīʿ fī al-jomle be reḍā-ye Ḥaqq pey borde-and, be ān
tamassok namāyand va be ān ʿāmel shavand. Īn ast sabab-e rastegārī dar
donyā va ākherat. Hejrat va balāyā-ye shomā lā yazāl taḥt-e leḥāẓ būde,
in shāʾ Allāh maḥfūẓ mānad. Intahā.

As commanded, convey greetings of praise on behalf of God to the
emigrants of the land of Ṣād, upon them be the glory of God. God
willing, may they be illumined by the light of unity and arise to serve
the Cause, so that the effect of that light may shine forth and become
manifest throughout the horizons. All have come to some understanding of
God's good pleasure; let them hold fast to it and act in accordance with
it. This is the means of salvation in this world and the next. Your
emigration and afflictions have always been under His regard; God
willing, may they remain preserved. End.

\subsection{160}\label{section-2580}

\begin{otherlanguage}{persian}

اینکه در بارهٴ جناب آقا میرزا آقا بابا علیه بهآء اللّه مرقوم داشتید
ذکرشان در ساحت اقدس بوده و بانوار عنایت شمس حقیقت فائز طوبی له در این
ایّام مکتوبی جناب آقا میرزا احمد علیه ٦٦٩ {[}بهآء اللّه{]} از اهل یا ساکن
ش‌اه بمصلحت جناب میرزا ح‌ی علیه بهآء اللّه و عنایته باین عبد ارسال داشته‌اند
و اسامی معدّده در آن مذکور و از جمله اسامی جناب آقا میرزا آقا بابا علیه
٦٦٩ {[}بهآء اللّه{]} بوده مخصوص هر یک لوح امنع نازل و ارسال شد انشآءاللّه
فائز میشوند فی‌الحقیقه امر بقسمی است که آنمحبوب مرقوم داشته‌اند چه که سفر
جناب میرزا بآنسمتها سبب ارتفاع کلمه شد هذا یکفیه وربّ العالمین چون
اینکلمه از لسان محبوب استماع گشت ذکر شد و الّا اینعبد قابل آنکه ذکر خدمات
ایشان را نماید نبوده و نیست

\end{otherlanguage}

Īnke dar bāre-ye Janāb-e Āqā Mīrzā Āqā Bābā ʿalayhi bahāʾu Allāh marqūm
dāshtīd, dhekreshān dar sāḥat-e aqdas būde va be anvār-e ʿenāyat-e
Shams-e Ḥaqīqat fāʾez, ṭūbā lahu. Dar īn ayyām maktūbī Janāb-e Āqā Mīrzā
Aḥmad ʿalayhi 669 {[}bahāʾu Allāh{]} az ahl yā sāken-e Sh-Ah be
maṣlaḥat-e Janāb-e Mīrzā Ḥ-Yā ʿalayhi bahāʾu Allāhi va ʿināyatuhu be īn
ʿabd ersāl dāshte-and va asāmī-ye motaʿaddede dar ān madhkūr va az
jomle-ye asāmī Janāb-e Āqā Mīrzā Āqā Bābā ʿalayhi 669 {[}bahāʾu Allāh{]}
būde. Makhṣūṣ-e har yek lowḥ-e amnaʿ nāzel va ersāl shod, in shāʾ Allāh
fāʾez mīshavand. Fī al-ḥaqīqe amr be qesmī ast ke ān maḥbūb marqūm
dāshte-and, che ke safar-e Janāb-e Mīrzā be ān samthā sabab-e ertefāʿ-e
kalame shod. Hādhā yakfīhi va Rabbi al-ʿālamīn. Chūn īn kalame az
lesān-e Maḥbūb estemāʿ gasht dhekr shod, va ellā īn ʿabd qābel-e ānke
dhekr-e khadamāt-e īshān rā namāyad nabūde va nīst.

Concerning what you wrote about Āqā Mīrzā Āqā Bābā, upon him be the
glory of God: he has been remembered in the Most Holy Court and has
attained the lights of the loving-kindness of the Sun of Truth. Blessed
is he! In these days Āqā Mīrzā Aḥmad, upon him be 669 {[}the glory of
God{]}, a native or resident of Sh-Ah, has sent this servant a letter on
the advice of Mīrzā Ḥ-Yā, upon him be the glory of God and His
loving-kindness. Several names were mentioned in it, including that of
Āqā Mīrzā Āqā Bābā, upon him be 669 {[}the glory of God{]}. A most
exalted Tablet was revealed specifically for each one and sent. God
willing, they will receive them. In truth, the matter is as you,
beloved, have written, for Mīrzā's journey to those regions became the
cause of the exaltation of the Word. ``This suffices him, by the Lord of
the worlds!'' Since these words were heard from the tongue of the
Beloved, they have been mentioned; otherwise, this servant has never
been, nor is he now, worthy to recount his services.

\subsection{161}\label{section-2581}

\begin{otherlanguage}{persian}

اینکه در بارهٴ ورقه اهل جناب میرزا ح و سین علیهما بهآء اللّه و عنایته
مرقوم فرمودند سبب حزن شد و اهل سرادق عصمت جمیع بکمال حزن مشهود و ظاهر و
در ساحت اقدس هم مذکور خدمت جناب میرزا از جانب اینفانی تعزیت بگوئید و
انشآءاللّه خود اینعبد هم عرض مینماید آنچه را که از ضعف و مسکنتش حاکی است

\end{otherlanguage}

Īnke dar bāre-ye varaqe, ahl-e Janāb-e Mīrzā Ḥā va Sīn ʿalayhimā bahāʾu
Allāhi va ʿināyatuhu marqūm farmūdand sabab-e ḥozn shod va ahl-e
sorādeq-e ʿeṣmat jamīʿ be kamāl-e ḥozn mashhūd va ẓāher va dar sāḥat-e
aqdas ham madhkūr. Khedmat-e Janāb-e Mīrzā az jāneb-e īn fānī taʿzīyat
begūʾīd va in shāʾ Allāh khod-e īn ʿabd ham ʿarḍ mīnamāyad ānche rā ke
az ḍaʿf va maskanatash ḥākī ast.

What you wrote concerning the leaf, the wife of Mīrzā Ḥā and Sīn, upon
them both be the glory of God and His loving-kindness, caused sorrow.
All the inhabitants of the tabernacle of purity were visibly and
manifestly filled with the deepest sorrow, and the matter was also
mentioned in the Most Holy Court. Convey this mortal one's condolences
to Mīrzā. God willing, this servant himself will also offer words
expressive of his weakness and lowliness.

\subsection{162}\label{section-2582}

\begin{otherlanguage}{persian}

و اینکه ذکر صعود جناب حاجی میرزا حسن هراتی علیه بهآء اللّه نمودند در ساحت
امنع اقدس عرض شد اینکلمات عالیات محض فضل و عنایت مخصوص ایشان نازل قوله
عزّ جلاله و عزّ ذکره و عزّ ثنائه

\end{otherlanguage}

Va īnke dhekr-e ṣoʿūd-e Janāb-e Ḥājī Mīrzā Ḥasan-e Herātī ʿalayhi bahāʾu
Allāh namūdand dar sāḥat-e amnaʿ-e aqdas ʿarḍ shod. Īn kalemāt-e ʿālīyāt
maḥḍ-e faḍl va ʿenāyat makhṣūṣ-e īshān nāzel, qawluhu ʿazza jalāluhu va
ʿazza dhikruhu va ʿazza thanāʾuhu:

Your mention of the ascension of Ḥājī Mīrzā Ḥasan of Herat, upon him be
the glory of God, was presented in the Most Exalted, Most Holy Court.
These lofty words were revealed especially for him, out of pure grace
and loving-kindness---His words, exalted be His glory, exalted be His
remembrance, and exalted be His praise:

\subsection{163}\label{section-2583}

\begin{otherlanguage}{arabic}

هو الذّاکر المعین

\end{otherlanguage}

Huwa al-dhākiru al-muʿīn.

He is the One Who remembers, the Helper.

\subsection{164}\label{section-2584}

\begin{otherlanguage}{arabic}

یا حسن یزورک المظلوم من مقامه الأعلی الّذی جعله اللّه مطاف اهل ذروة العلیا
و یذکرک بذکر فاح به عرف العنایة بین البریّة طوبی لک و لصعودک و لتوجّهک و
لورودک فی الرّفیق الأعلی و المقام الأبهی اشهد انّک اقبلت الی الغایة القصوی
و الأفق الأعلی اذ ارتفع حفیف سدرة المنتهی بین الأرض و السّمآء و شهدت بما
شهد به لسان الکبریآء و قبلت ما نزّل من عنده و ظهر من لدنه طوبی لک بما فزت
بالغفران من لدی الرّحمن و وجد منک الملأ الأعلی عرف حبّی طوبی لمن یذکرک بما
ذکرک المظلوم فیهذا المقام الممنوع قد کنت للّه و رجعت الی بحر رحمته ان هذا
الّا من فضله العزیز المحبوب انتهی

\end{otherlanguage}

Yā Ḥasan, yazūruka al-maẓlūmu min maqāmihi al-aʿlā alladhī jaʿalahu
Allāhu maṭāfa ahli dhurvati al-ʿulyā, va yadhkuruka bi-dhikrin fāḥa bihi
ʿarfu al-ʿināyati bayna al-bariyya. Ṭūbā laka va li-ṣuʿūdika va
li-tavajjuhika va li-vurūdika fī al-rafīqi al-aʿlā va al-maqāmi al-abhā.
Ashhadu annaka aqbalta ilā al-ghāyati al-quṣvā va al-ufuqi al-aʿlā idh
irtafaʿa ḥafīfu sidrati al-muntahā bayna al-arḍi va al-samāʾ, va
shahidta bi-mā shahida bihi lisānu al-kibriyāʾ, va qabilta mā nuzzila
min ʿindihi va ẓahara min ladunhu. Ṭūbā laka bi-mā fuzta bi-al-ghufrāni
min ladā al-raḥmān, va vajada minka al-malaʾu al-aʿlā ʿarfa ḥubbī. Ṭūbā
li-man yadhkuruka bi-mā dhakaraka al-maẓlūmu fī hādhā al-maqāmi
al-mamnūʿ. Qad kunta li-llāhi va rajaʿta ilā baḥri raḥmatihi. In hādhā
illā min faḍlihi al-ʿazīzi al-maḥbūb. Intahā.

O Ḥasan! The Wronged One visits you from His most exalted station, which
God has made the center around which circle the inhabitants of the
loftiest summit, and remembers you with a remembrance through which the
fragrance of loving-kindness has spread among all people. Blessed are
you, and blessed are your ascension, your turning toward Him, and your
entry into the companionship on high and the all-glorious station. I
bear witness that you turned toward the ultimate Goal and the most
exalted Horizon when the rustling of the Lote-Tree of the Utmost
Boundary rose between earth and heaven, that you bore witness to what
the Tongue of Grandeur bore witness to, and that you accepted what was
sent down from Him and appeared from His presence. Blessed are you for
having attained forgiveness from the All-Merciful, and because the
Concourse on High perceived from you the fragrance of My love. Blessed
is whoever remembers you with that by which the Wronged One has
remembered you in this inaccessible station. You belonged to God and
have returned to the ocean of His mercy. This is nothing but His mighty,
well-beloved grace. End.

\subsection{165}\label{section-2585}

\begin{otherlanguage}{persian}

الحمد للّه باین عنایت کبری مفتخر شدند طوبی از برای نفوسیکه باین فیض اکبر
فائزند و مقامش را میشناسند نفوس مستقیمهٴ طاهرهٴ مطمئنّه از سماء عنایت
فیّاض قسمت داشته و دارند و من غیر تعطیل بر ایشان نازل و جاری چه در حیات و
چه در ممات

\end{otherlanguage}

Al-ḥamdu li-llāh be īn ʿenāyat-e kobrā moftakhar shodand. Ṭūbā az
barā-ye nofūsī ke be īn feyḍ-e akbar fāʾez and va maqāmash rā
mīshenāsand. Nofūs-e mostaqīme-ye ṭāhere-ye moṭmaʾenne az samā-ye
ʿenāyat-e Fayyāḍ qesmat dāshte va dārand va min ghayr-e taʿṭīl bar īshān
nāzel va jārī che dar ḥayāt va che dar mamāt.

Praise be to God, he has been honored with this supreme favor. Blessed
are the souls who attain this greatest bounty and recognize its station.
Steadfast, pure, assured souls have received and continue to receive
their share from the heaven of the loving-kindness of the All-Bountiful,
and without interruption it descends and flows upon them, whether in
life or in death.

\subsection{166}\label{section-2586}

\begin{otherlanguage}{persian}

و اینکه در صعود حضرت مرفوع مبرور آقا سیّد ابوطالب علیه منکلّ بهآء ابهاه
مرقوم داشتید چندی قبل از طرف حدبا حضرت اسم‌اللّه زین علیه منکلّ بهآء ابهاه
ذکر ایشان و همچنین بعض مطالب از لسان اهل ایشان علیها ٦٦٩ {[}بهآء اللّه{]}
بساحت اقدس ارسال داشتند و جواب بتفصیل نازل و ارسال شد و ذکرشانهم در آن
لوح مبارک نازل له الحمد و المنّة فیکلّ الأحوال البهآء المشرق اللّائح من افق
عنایة ربّنا علی حضرتکم و علی من معکم من الذّکور و الاناث

\end{otherlanguage}

Va īnke dar ṣoʿūd-e Ḥaḍrat-e marfūʿ-e mabrūr Āqā Sayyed Abū Ṭāleb
ʿalayhi min kulli bahāʾin abhāhu marqūm dāshtīd, chandī qabl az ṭaraf-e
Ḥadbā Ḥaḍrat-e Ism Allāh Zayn ʿalayhi min kulli bahāʾin abhāhu dhekr-e
īshān va hamchenīn baʿḍ-e maṭāleb az lesān-e ahl-e īshān ʿalayhā 669
{[}bahāʾu Allāh{]} be sāḥat-e aqdas ersāl dāshtand va javāb be tafṣīl
nāzel va ersāl shod va dhekreshān ham dar ān lowḥ-e mobārak nāzel. Lahu
al-ḥamdu va al-minnatu fī kulli al-aḥvāl. Al-bahāʾu al-mushriqu
al-lāʾiḥu min ufuqi ʿināyati rabbinā ʿalā ḥaḍratikum va ʿalā man maʿakum
mina al-dhukūri va al-ināth.

Concerning what you wrote about the ascension of the honored, exalted,
righteous Āqā Sayyid Abū Ṭālib, upon him be the most glorious of all
glory: some time ago, the honored Name of God, Zayn, upon him be the
most glorious of all glory, sent from Ḥadbā to the Most Holy Court a
mention of him, as well as certain matters related by his wife, upon her
be 669 {[}the glory of God{]}. A detailed answer was revealed and sent,
and a remembrance of him was also revealed in that blessed Tablet. To
Him belong praise and gratitude in all circumstances. May the glory
shining and gleaming from the horizon of our Lord's loving-kindness rest
upon you and upon those with you, both men and women.

\subsection{167}\label{section-2587}

\begin{otherlanguage}{persian}

عرض میشود دو لوح امنع اقدس مدّتهاست مخصوص ابناء مرفوع مبرور جناب ذبیح
علیه بهآء اللّه و بحر رحمته از سماء مشیّت نازل و اینمدّت ارسال نشد و بتأخیر
افتاد تا این ایّام ارسال گشت شاید بعنایات الهیّه متذکّر شوند و بما ینبغی و
یحبّه اللّه قیام نمایند خدمت هر نفسی در ساحت اقدس مذکور و منظور ولو بمقدار
سمّ ابره

\end{otherlanguage}

ʿArż mī-shavad du lawḥ-i amnaʿ-i aqdas muddat-hāst makhṣūṣ-i abnā-yi
marfūʿ-i mabrūr Janāb-i Dhabīḥ ʿalayhi bahāʾu Allāh wa-baḥru raḥmatihi
az samā-yi mashiyyat nāzil va īn muddat irsāl nashud va bi-taʾkhīr uftād
tā īn ayyām irsāl gasht shāyad bi-ʿināyāt-i ilāhiyya mutadhakkir shavand
va bi-mā yanbaghī wa-yuḥibbuhu Allāh qiyām namāyand khidmat-i har nafsī
dar sāḥat-i aqdas madhkūr va manẓūr va-law bi-miqdār-i samm-i ibra.

I submit that two most exalted, most holy Tablets were revealed long ago
from the heaven of the divine will specifically for the sons of the
exalted and righteous Dhabīḥ, upon him be the glory of God and the ocean
of His mercy. They were not sent during this time, and their dispatch
was delayed until these days, when they were sent, so that perhaps,
through divine favors, they may become mindful and arise to do what is
fitting and what God loves. Every soul's service is mentioned and
regarded in the Most Holy Court, even if it is no greater than a
needle's eye.

\subsection{168}\label{section-2588}

\begin{otherlanguage}{persian}

مرفوع ذبیح علیه بهآء اللّه در اوّل ایّام بشرافت اقبال و ایمان فائز شدند
لازال در ساحت اقدس مذکور بوده و هستند بعد از استماع ندا اقلّ من آن توقّف
ننمودند یشهد بذلک من عنده کتاب مبین لذا مراعات ایشان در هر حال شده و
انشآءاللّه میشود و اینفانی از حقّ جلّ جلاله تأیید ایشان را طلب مینماید
آنمحبوبهم باید بمهربانی ایشانرا حفظ نمایند اوّل عمر است اریاح شبابرا مهبّی
دیگر است

\end{otherlanguage}

Marfūʿ Dhabīḥ ʿalayhi bahāʾu Allāh dar avval-i ayyām bi-sharāfat-i iqbāl
va īmān fāʾiz shudand lā-zāl dar sāḥat-i aqdas madhkūr būdih va hastand
baʿd az istimāʿ-i nidā aqall min ān tavaqquf nanamūdand yashhadu
bi-dhālika man ʿindahu kitābun mubīn lidhā murāʿāt-i īshān dar har ḥāl
shudih va in shāʾa Allāh mī-shavad va īn fānī az Ḥaqq jalla jalāluhu
taʾyīd-i īshān rā ṭalab mī-namāyad ān maḥbūb ham bāyad bi-mihrabānī
īshān rā ḥifẓ namāyand avval-i ʿumr ast aryāḥ-i shabāb rā mahabbī dīgar
ast.

The exalted Dhabīḥ, upon him be the glory of God, attained the honor of
turning toward God and believing in the earliest days. He has always
been, and continues to be, remembered in the Most Holy Court. After
hearing the call, he did not hesitate for even an instant. He with whom
is a clear Book bears witness to this. Care has therefore been taken of
them in every circumstance and, God willing, will continue to be taken.
This humble servant asks God, exalted be His glory, to sustain them. You
too, dear friend, should watch over them with kindness. They are at the
beginning of life; the winds of youth blow in a different way.

\subsection{169}\label{section-2589}

\begin{otherlanguage}{persian}

عرض دیگر مخدّره خدیجه بیگم علیها بهآء اللّه چند کرّه باین عبد مکتوب ارسال
داشته و این عبد نظر بانقلابات و محاربهٴ اطراف جواب ارسال نداشت تا در این
کرّه مکتوبیکه چند شهر قبل نوشته شده بود ارسال داشت که آنمحبوب برسانند

\end{otherlanguage}

ʿArż-i dīgar mukhaddara Khadīja Bīgum ʿalayhā bahāʾu Allāh chand karra
bi-īn ʿabd maktūb irsāl dāshtih va īn ʿabd naẓar bi-inqilābāt va
muḥāraba-yi aṭrāf javāb irsāl nadāsht tā dar īn karra maktūbī kih chand
shahr qabl nivishtih shudih būd irsāl dāsht kih ān maḥbūb birasānand.

Another matter to submit: the honored lady Khadīja Bīgum, upon her be
the glory of God, has sent letters to this servant several times. In
view of the upheavals and warfare in the surrounding regions, this
servant sent no reply until this occasion, when he sent a letter written
several months earlier for you, dear friend, to deliver.

\subsection{170}\label{section-2590}

\begin{otherlanguage}{persian}

و همچنین جواب مکتوبیکه جناب آقا محمّد کریم ع‌ط علیه بهآء اللّه بحضرت اسم
جود علیه بهآء اللّه الأبهی ارسال داشتند اینفانی نوشته ارسال داشت
انشآءاللّه از اقداح فلاح در جمیع احیان بیاشامند و بما ینبغی لأیّام اللّه
عامل باشند از اطراف خبر رسید که جناب ایشان بما انزله اللّه فی الکتاب عامل
شدند و مخصوص تبلیغ وکیل اخذ نمودند طوبی له و نعیماً له هر نفسی بما اراده
اللّه عمل نماید البتّه در دنیا و آخرت ثمر بخشد و اثر پدیدار نماید و در هر
دو جهان بکار آید

\end{otherlanguage}

Va hamchunīn javāb-i maktūbī kih Janāb-i Āqā Muḥammad Karīm ʿ--Ṭ ʿalayhi
bahāʾu Allāh bi-Ḥażrat-i Ism-i Jūd ʿalayhi bahāʾu Allāh al-Abhā irsāl
dāshtand īn fānī nivishtih irsāl dāsht in shāʾa Allāh az aqdāḥ-i falāḥ
dar jamīʿ-i aḥyān biyāshāmand va bi-mā yanbaghī li-ayyāmi Allāh ʿāmil
bāshand az aṭrāf khabar rasīd kih Janāb-i īshān bi-mā anzalahu Allāh fī
al-kitāb ʿāmil shudand va makhṣūṣ-i tablīgh vakīl akhdh namūdand ṭūbā
lahu wa-naʿīman lahu har nafsī bi-mā arādahu Allāh ʿamal namāyad albatta
dar dunyā va ākhirat thamar bakhshad va athar padīdār namāyad va dar har
du jahān bi-kār āyad.

Likewise, this humble servant wrote and sent a reply to the letter that
Āqā Muḥammad Karīm ʿ--Ṭ, upon him be the glory of God, had sent to His
Honor the Name of Generosity, upon him be the glory of God, the
All-Glorious. God willing, may he drink at all times from the cups of
salvation and act as befits the days of God. News reached us from the
surrounding regions that he has acted according to what God has revealed
in the Book and has appointed an agent specifically for teaching the
Faith. Blessed is he, and may joy be his! Whoever acts according to what
God has willed will assuredly bear fruit and produce an effect in this
world and the next, and will be of service in both worlds.

\subsection{171}\label{section-2591}

\begin{otherlanguage}{persian}

دیگر دوستان آن ارض را مکرّر تکبیر و سلام عرض نموده بجناب علی پاشا خان
علیه بهآء اللّه بشارت دهید چه که فرش ایشان رسید و یومی از ایّام گسترده شد
و جمال قدم تشریف بردند در بستان و بقدوم حضرت مقصود روح من فی اللّاهوت
لقدومه الفدآء مشرّف گشت و فائز شد و جناب حاجی غلامعلی مسافر علیه بهآء
اللّه را احضار فرمودند و فرمودند اگر جناب پاشا علیه بهائی را دیدی او را
بشارت ده بقبول آنچه ارسال داشته نعیماً له و هنیئاً له انشآءاللّه مشتعل
باشند بشأنیکه روشنیشان ظاهر و هویدا باشد این فانی از حقّ جلّ جلاله میطلبد
اولیا و اصفیا را مؤیّد فرماید بر عرفان ما قدّر لهم چه اگر مطّلع شوند سیل
عالم قادر بر اطفاء نار محبّتشان نباشد و روشنی آنرا تاریکی اخذ ننماید

\end{otherlanguage}

Dīgar dūstān-i ān arż rā mukarrar takbīr va salām ʿarż namūdih
bi-Janāb-i ʿAlī Pāshā Khān ʿalayhi bahāʾu Allāh bishārat dahīd chih kih
farsh-i īshān rasīd va yawmī az ayyām gustardih shud va Jamāl-i Qidam
tashrīf burdand dar būstān va bi-qudūm-i Ḥażrat-i Maqṣūd rūḥu man fī
al-lāhūt li-qudūmihi al-fidāʾ musharraf gasht va fāʾiz shud va Janāb-i
Ḥājī Ghulām-ʿAlī Musāfir ʿalayhi bahāʾu Allāh rā iḥżār farmūdand va
farmūdand agar Janāb-i Pāshā ʿalayhi bahāʾī rā dīdī ū rā bishārat dih
bi-qabūl-i ānchih irsāl dāshtih naʿīman lahu wa-hanīʾan lahu in shāʾa
Allāh mushtaʿil bāshand bi-shaʾnī kih rawshanī-shān ẓāhir va huvaydā
bāshad īn fānī az Ḥaqq jalla jalāluhu mī-ṭalabad awliyā va aṣfiyā rā
muʾayyad farmāyad bar ʿirfān-i mā quddira lahum chih agar muṭṭaliʿ
shavand sayl-i ʿālam qādir bar iṭfāʾ-i nār-i maḥabbat-shān nabāshad va
rawshanī-yi ān rā tārīkī akhdh nanamāyad.

Further, convey repeated greetings of praise and peace to the friends in
that land, and give ʿAlī Pāshā Khān, upon him be the glory of God, the
glad news that his carpet arrived. One day it was spread out, and the
Ancient Beauty went into the garden, and it was honored and blessed by
the footsteps of the Desired One---may the soul of whoever dwells in the
realm of divinity be a sacrifice to His footsteps. He summoned Ḥājī
Ghulām-ʿAlī Musāfir, upon him be the glory of God, and said, ``If you
see Pāshā, upon him be My glory, give him the glad news that what he
sent has been accepted. May joy and delight be his!'' God willing, may
they be so enkindled that their light is manifest and evident. This
humble servant asks God, exalted be His glory, to aid His friends and
chosen ones in recognizing what has been ordained for them. For if they
become aware of it, the flood of the world will be powerless to
extinguish the fire of their love, and darkness will not overtake its
light.

\subsection{172}\label{section-2592}

\begin{otherlanguage}{persian}

از جمله عنایات جدیده آنکه مخصوص اسامی مذکوره الواح بدیعهٴ منیعه نازل هم
بخطّ این فانی که در مکتوب ثبت شده و هم بخطّ مبارک حضرت غصن اللّه الأکبر
روحی و ذاتی و کینونتی لتراب قدومه الفدآء فی‌الحقیقه صاحبان الواح ذو
المقامین و ذو اللّوحین و ذو الذّکرین شده‌اند ینبغی ان نقول بکلّ جوارحنا و
ارکاننا لک الحمد یا فضّال القدیم و لک العنایة یا مقصود العالمین النّور و
البهآء من افق عنایة ربّنا مالک الأسمآء علی حضرتکم و علی من معکم و علی
اولیآء اللّه و احبّائه

\end{otherlanguage}

Az jumla-yi ʿināyāt-i jadīda ānkih makhṣūṣ-i asāmī-yi madhkūra alvāḥ-i
badīʿa-yi manīʿa nāzil ham bi-khaṭṭ-i īn fānī kih dar maktūb thabt
shudih va ham bi-khaṭṭ-i mubārak-i Ḥażrat-i Ghuṣn-i Allāh al-Akbar rūḥī
wa-dhātī wa-kaynūnatī li-turābi qudūmihi al-fidāʾ fī al-ḥaqīqa ṣāḥibān-i
alvāḥ dhū al-maqāmayn va dhū al-lawḥayn va dhū al-dhikrayn shudih-and
yanbaghī an naqūla bi-kulli jawāriḥinā wa-arkāninā laka al-ḥamdu yā
faḍḍāl al-qadīm wa-laka al-ʿināyatu yā maqṣūda al-ʿālamīn al-nūru
wa-al-bahāʾu min ufuqi ʿināyati Rabbinā Māliki al-Asmāʾ ʿalā ḥażratikum
wa-ʿalā man maʿakum wa-ʿalā awliyāʾi Allāh wa-aḥibbāʾihi.

Among the new favors is that wondrous, exalted Tablets have been
revealed specifically for the names mentioned, both in the handwriting
of this humble servant, as recorded in the letter, and in the blessed
handwriting of His Holiness the Most Great Branch of God---may my soul,
my essence, and my being be a sacrifice to the dust of His footsteps. In
truth, the recipients of the Tablets have become possessors of two
stations, two Tablets, and two remembrances. With all our limbs and
members we should say, ``Praise be to You, O ancient Bestower, and Yours
is loving-kindness, O Desire of all the worlds.'' May light and glory
from the horizon of the loving-kindness of our Lord, the King of Names,
rest upon you, upon those with you, and upon the friends and loved ones
of God.

\subsection{173}\label{section-2593}

\begin{otherlanguage}{arabic}

خ‌ادم فی ١٣ شهر صفر ١٣٠٠

\end{otherlanguage}

Khādim fī 13 shahr Ṣafar 1300.

The servant, on the 13th of the month of Ṣafar, 1300.

\bookmarksetup{startatroot}

\chapter{Lawḥ Shams al-Ḥaqq (Tablet of the Sun of
Truth)}\label{lawux1e25-shams-al-ux1e25aqq-tablet-of-the-sun-of-truth}

\section{BH00011}\label{bh00011}

\emph{Translated on 17 Qudrat 183 using GPT 6.1 Sol}

\subsection{1}\label{section-2594}

\begin{otherlanguage}{arabic}

باسم الّذی کلّ عنه لغافلون

\end{otherlanguage}

Bi-smi alladhī kullun ʿanhu la-ghāfilūn.

In the name of Him of whom all are heedless.

\subsection{2}\label{section-2595}

\begin{otherlanguage}{persian}

حمد مقدّس از السن ممکنات و منزّه از ذکر موجودات ساحت قدس سلطان قدمی را سزا
است که از رشحی از رشحات قلزم عزّ بی‌زوال خود جمیع ذرّات وجود را بمدینهٴ
حیات باقیهٴ طیّبهٴ خود هدایت فرموده و بنفحه‌ئی از نفحات قدس خود جمیع حقایق
اشیاء را برضوان عزّ احدیّت خود دلالت نموده سحاب فیض سبحانیش چنان رحمتی
مبذول فرموده که هر ذرّه خود را غنیّ و مستغنی از دون خود فرض گرفته و شمس
عنایت ربّانیش از افق عزّ ابهائی چنان مشرق شده که در جمیع مرایای وجود
ممکنات در ذات او بنفس او هیئت تجلّی انوار شمسی منطبع گشته جلّت عظمته و علت
سلطنته و کمل احسانه و ظهر امره و لاح وجهه و انّه لهو السّلطان الباقی
العلیم الحکیم و نظر ناظرین باین مراتب است که بعضی دعاوی نموده و از خمر
مرحمت یزدانی و کوثر مکرمت سبحانی زمام از دست داده و بمقامات فوق خود صعود
نموده‌اند در عالم ذکر و لفظ نه در عالم معنی و حقیقت و جوهر

\end{otherlanguage}

Ḥamd-e muqaddas az alsun-e mumkināt va munazzah az dhikr-e mawjūdāt
sāḥat-e quds-e sulṭān-e qidamī rā sazā ast ke az rashḥī az rashaḥāt-e
qulzum-e ʿizz-e bī-zavāl-e khud jamīʿ-e dharrāt-e vujūd rā bi-madīne-ye
ḥayāt-e bāqiye-ye ṭayyibe-ye khud hidāyat farmūde va bi-nafḥe-ī az
nafaḥāt-e quds-e khud jamīʿ-e ḥaqāyiq-e ashyā rā bi-riḍvān-e ʿizz-e
aḥadiyyat-e khud dilālat namūde. Saḥāb-e fayḍ-e subḥāniyash chinān
raḥmatī mabdhūl farmūde ke har dharre khud rā ghanī va mustaghnī az
dūn-e khud farḍ girifte va shams-e ʿināyat-e rabbāniyash az ufuq-e
ʿizz-e abhāʾī chinān mushriq shude ke dar jamīʿ-e marāyā-ye vujūd-e
mumkināt dar dhāt-e ū bi-nafs-e ū hayʾat-e tajallī-ye anvār-e shamsī
munṭabiʿ gashte. Jallat ʿaẓamatuhu wa-ʿalat salṭanatuhu wa-kamula
iḥsānuhu wa-ẓahara amruhu wa-lāḥa wajhuhu wa-innahu la-huwa al-Sulṭānu
al-Bāqī al-ʿAlīmu al-Ḥakīm. Va naẓar-e nāẓirīn bi-īn marātib ast ke
baʿḍī daʿāvī namūde va az khamr-e marḥamat-e yazdānī va kawthar-e
makramat-e subḥānī zimām az dast dāde va bi-maqāmāt-e fawq-e khud ṣuʿūd
namūde-and dar ʿālam-e dhikr va lafẓ na dar ʿālam-e maʿnī va ḥaqīqat va
jawhar.

Praise sanctified above the tongues of contingent beings and exalted
beyond the remembrance of created things befits the holy court of the
King of Eternity, who, through a sprinkling from the boundless ocean of
His unfading glory, has guided every atom of existence to the city of
His pure, everlasting life, and through a breath of His holiness has
directed the realities of all things to the paradise of the glory of His
oneness. The cloud of His transcendent grace has bestowed such mercy
that every atom considers itself rich and independent of all else. The
sun of His divine providence has shone so brightly from the horizon of
all-glorious majesty that, in all the mirrors of contingent existence,
within each one's own essence and through its own self, the image of the
radiance of the sun's light has been imprinted. Exalted is His grandeur,
supreme His sovereignty, perfect His beneficence, manifest His Cause,
and radiant His face. He is indeed the Sovereign, the Everlasting, the
All-Knowing, the All-Wise. It is because observers have fixed their gaze
on these degrees that some have advanced claims and, intoxicated by the
wine of divine mercy and the Kawthar of transcendent generosity, have
lost their restraint and ascended to stations above their own---in the
realm of mention and words, not in the realm of meaning, reality, and
essence.

\subsection{3}\label{section-2596}

\begin{otherlanguage}{arabic}

صلوة اصلیّه و ساذج رحمت منیعه بر اوّل شجرهٴ مرتفعهٴ الهیّه و اوّل ظهور شمس
احدیّه الّذی به ظهر سرّ الوجود و طلع جمال المقصود و رجعت الممکنات الی مقاعد
قدس محمود و به لاح جمال الغیب و نادت عندلیب العزّ و غنّت ورقآء الأمر و ظهر
کلّ امر مکنون و علی ادلّائه و احبّائه الی یوم الّذی فیه یستشرق شمس القدم عن
افق عزّ محبوب

\end{otherlanguage}

Ṣalāt-e aṣliyye va sādhij-e raḥmat-e manīʿe bar avval shajare-ye
murtafiʿe-ye ilāhiyye va avval ẓuhūr-e shams-e aḥadiyye alladhī bihi
ẓahara sirru al-wujūd wa-ṭalaʿa jamālu al-Maqṣūd wa-rajaʿati
al-mumkinātu ilā maqāʿidi qudsi maḥmūd wa-bihi lāḥa jamālu al-ghayb
wa-nādat ʿandalību al-ʿizz wa-ghannat warqāʾu al-amr wa-ẓahara kullu
amrin maknūn wa-ʿalā adillāʾihi wa-aḥibbāʾihi ilā yawmi alladhī fīhi
yastashriqu shamsu al-qidam ʿan ufuqi ʿizzi maḥbūb.

Primal blessing and the pure essence of exalted mercy rest upon the
first lofty divine Tree and the first appearance of the Sun of Oneness,
through whom the mystery of existence was revealed, the beauty of the
Desired One rose, and contingent beings returned to seats of
praiseworthy holiness; through whom the beauty of the Unseen shone
forth, the nightingale of glory called out, the dove of the Cause sang,
and every hidden thing was revealed; and upon His guides and His loved
ones until the day on which the Sun of Eternity rises from the horizon
of the glory of the Beloved.

\subsection{4}\label{section-2597}

\begin{otherlanguage}{persian}

و بعد سؤال از شمس حقیقت و مرایای مستحکیه از او شده بود نزد آنجناب معلوم
بوده که نظر بآنکه اریاح اختلاف از جمیع اطراف وزیده و نفحات غلّ و بغضا از
تمام اشطار احاطه نموده این عبد متعذّر است از ذکر بدایع علوم مکنونه و
جواهر حکمت مخزونه چه اگر فی‌الحقیقه اینمطلب علی ما اراد اللّه جلّ جلاله ذکر
شود جمیع را مدهوش و منصعق خواهی دید بلکه اگر کشف حجاب نمائی و از مدینهٴ
وهمیّه که حال جمیع ناس حین غفلة عن اللّه در آن مقرّ یافته و معتکف گشته‌اند
بدرآئی هیاکل صافین و انفس عالین را معدوم و فانی و لاشیء محض خواهی یافت

\end{otherlanguage}

Va baʿd suʾāl az shams-e ḥaqīqat va marāyā-ye mustaḥkiye az ū shude būd.
Nazd-e ān-janāb maʿlūm būde ke naẓar bi-ānke aryāḥ-e ikhtilāf az jamīʿ-e
aṭrāf vazīde va nafaḥāt-e ghill va baghḍā az tamām-e ashṭār iḥāṭe namūde
īn ʿabd mutaʿadhdhir ast az dhikr-e badāyiʿ-e ʿulūm-e maknūne va
javāhir-e ḥikmat-e makhzūne. Chi agar fī al-ḥaqīqe īn maṭlab ʿalā mā
arāda Allāh jalla jalāluhu dhikr shavad jamīʿ rā madhūsh va munṣaʿiq
khāhī dīd. Balke agar kashf-e ḥijāb namāʾī va az madīne-ye vahmiyye ke
ḥāl jamīʿ-e nās ḥīna ghaflatin ʿani Allāh dar ān maqarr yāfte va
muʿtakif gashte-and bi-dar āʾī hayākil-e ṣāffīn va anfus-e ʿāllīn rā
maʿdūm va fānī va lā-shayʾ-e maḥḍ khāhī yāft.

Now, a question had been asked concerning the Sun of Truth and the
mirrors that reflect it. It is known to you that, since the winds of
discord have blown from every direction and the breaths of malice and
hatred have encompassed every quarter, this servant is unable to speak
of the wonders of hidden knowledge and the jewels of treasured wisdom.
For if this subject were truly set forth as God---exalted be His
glory---has willed, you would see everyone bewildered and thunderstruck.
Indeed, if you were to lift the veil and leave the city of illusion in
which all people, heedless of God, have now taken up residence and
secluded themselves, you would find the forms of those arrayed in ranks
and the souls of the exalted to be nonexistent, perishing, and sheer
nothingness.

\subsection{5}\label{section-2598}

\begin{otherlanguage}{arabic}

ولکن چون حقّ منیع دوست نداشته که سائلین را از باب رحمت خود محروم نماید
لهذا رشحی از طمطام یمّ جود بر تو و اهل امکان مبذول میشود تا که بنصیب
مقدّرهٴ خود فایز گردد

\end{otherlanguage}

Valīkin chūn Ḥaqq-e manīʿ dūst nadāshte ke sāʾilīn rā az bāb-e raḥmat-e
khud maḥrūm namāyad li-hādhā rashḥī az ṭamṭām-e yamm-e jūd bar tu va
ahl-e imkān mabdhūl mī-shavad tā ke bi-naṣīb-e muqaddare-ye khud fāʾiz
gardad.

But since God, the Exalted, has not wished to turn seekers away deprived
from the door of His mercy, a sprinkling from the vast ocean of
generosity is bestowed upon you and the people of the contingent world,
so that each may attain their allotted share.

\subsection{6}\label{section-2599}

\begin{otherlanguage}{persian}

باری نزد اولی العلم و البصر معلوم بوده که شمس احدیّه و ساذج صمدیّه در مقعد
عزّ تقدیس و مکمن قدس تنزیه خود لم‌یزل و لایزال طالع و مشرق و مضیء بوده و
خواهد بود چنانچه هر ذی بصر دقیق و ذی نظر رقیق و صاحبان فطرت ربّانی و
شاربان بادهٴ لطایف معانی آن شمس الهی و جوهر عزّ سلطانی را در مقرّ خود بعین
بصیرت که مقدّس از این عیون تیرهٴ محتجبه است ملاحظه مینمایند علی مقام لا
یأخذها غروب و لا افول و لا کسوف

\end{otherlanguage}

Bārī nazd-e ulī al-ʿilm va al-baṣar maʿlūm būde ke shams-e aḥadiyye va
sādhij-e ṣamadiyye dar maqʿad-e ʿizz-e taqdīs va makman-e quds-e
tanzīh-e khud lam-yazal va lā-yazāl ṭāliʿ va mushriq va muḍīʾ būde va
khāhad būd. Chinānchi har dhī-baṣar-e daqīq va dhī-naẓar-e raqīq va
ṣāḥibān-e fiṭrat-e rabbānī va shāribān-e bāde-ye laṭāʾif-e maʿānī ān
shams-e ilāhī va jawhar-e ʿizz-e sulṭānī rā dar maqarr-e khud bi-ʿayn-e
baṣīrat ke muqaddas az īn ʿuyūn-e tīre-ye muḥtajibe ast mulāḥaẓe
mī-namāyand ʿalā maqāmin lā yaʾkhudhuhā ghurūbun wa-lā ufūlun wa-lā
kusūf.

It is known to those endowed with knowledge and vision that the Sun of
Oneness and pure Essence of Eternal Self-Sufficiency has always risen,
shone, and shed its light, and will forever do so, upon its seat of
sanctified glory and within its holy retreat of transcendence. Thus, all
who possess keen vision and subtle perception, those endowed with a
divine nature, and those who drink the wine of subtle meanings behold
that divine Sun and Essence of sovereign glory in its own abode, with
the eye of insight, which is sanctified above these darkened, veiled
eyes, in a station where no setting, decline, or eclipse can overtake
it.

\subsection{7}\label{section-2600}

\begin{otherlanguage}{persian}

و آنانکه محتجب بهوای نفس شده و از ملاحظهٴ انوار صبح صادق که لم‌یزل بعنایت
شمس مستنیر بوده محروم شوند بأسی بر تشعشعات انوار شمس نبوده چه که او در
مقرّ خود لایح و ساطع و مستضیء بوده و این عدم ادراک ببصر باصر و نظر ناظر
راجع بوده و خواهد بود و هم‌چنین جمیع این نفی و اثبات چه از مقبلین و چه از
معرضین بمقامات خود این انفس ضعیفه و قویّه راجع است بهیچوجه بآن شمس حقیقت
راجع نخواهد بود چه که آن جوهر وجود در علوّ ارتفاع و سموّ امتناع خود بوده و
مقدّس از تصدیق مقبل و تکذیب معرض

\end{otherlanguage}

Va ānānke muḥtajib bi-havā-ye nafs shude va az mulāḥaẓe-ye anvār-e
ṣubḥ-e ṣādiq ke lam-yazal bi-ʿināyat-e shams mustanīr būde maḥrūm
shavand baʾsī bar tashaʿshuʿāt-e anvār-e shams nabūde chi ke ū dar
maqarr-e khud lāʾiḥ va sāṭiʿ va mustaḍīʾ būde va īn ʿadam-e idrāk
bi-baṣar-e bāṣir va naẓar-e nāẓir rājiʿ būde va khāhad būd. Va
ham-chinīn jamīʿ-e īn nafy va ithbāt chi az muqbilīn va chi az muʿriḍīn
bi-maqāmāt-e khud-e īn anfus-e ḍaʿīfe va qaviyye rājiʿ ast bi-hīch-vajh
bi-ān shams-e ḥaqīqat rājiʿ nakhāhad būd chi ke ān jawhar-e vujūd dar
ʿuluvv-e irtifāʿ va sumuvv-e imtināʿ-e khud būde va muqaddas az taṣdīq-e
muqbil va takdhīb-e muʿriḍ.

As for those who become veiled by selfish desire and are deprived of
beholding the lights of the true dawn, which has always been illumined
through the sun's bounty, this brings no harm to the radiance of the
sun's light. For it shines, gleams, and radiates in its own abode, while
this failure of perception has always pertained, and will always
pertain, to the sight of the beholder and the vision of the observer.
Likewise, all this denial and affirmation, whether from those who turn
toward it or those who turn away, pertains to the stations of these weak
or strong souls themselves; it can in no way pertain to that Sun of
Truth. For that Essence of existence abides in its own lofty exaltation
and inaccessible sublimity, sanctified above the affirmation of those
who turn toward it and the denial of those who turn away.

\subsection{8}\label{section-2601}

\begin{otherlanguage}{persian}

ولکن هر بصری که بفطرت اصلیّه آن ساذج قدم را ادراک نمود از او محسوب و در
ظلّ او لم‌یزل و لایزال محشور خواهد بود و در جنّت ابهی که اعلی الجنان است
وارد و مقرّ یابد و دون آن در ظلمت نفس و نار نفی داخل بلکه الی ما شآء اللّه
خالد خواهد بود

\end{otherlanguage}

Valīkin har baṣarī ke bi-fiṭrat-e aṣliyye ān sādhij-e qidam rā idrāk
namūd az ū maḥsūb va dar ẓill-e ū lam-yazal va lā-yazāl maḥshūr khāhad
būd va dar jannat-e abhā ke aʿlā al-jinān ast vārid va maqarr yābad va
dūn-e ān dar ẓulmat-e nafs va nār-e nafy dākhil balke ilā mā shāʾa Allāh
khālid khāhad būd.

Yet every eye that perceives that pure Essence of Eternity through its
original nature will be reckoned as belonging to Him and will be
gathered beneath His shadow for all eternity. Such a soul will enter and
dwell in the all-glorious Paradise, the highest of all paradises.
Whoever does not will enter the darkness of self and the fire of
negation, and will indeed remain there for as long as God wills.

\subsection{9}\label{section-2602}

\begin{otherlanguage}{persian}

و علاوه بر آن اگر آن شمس مشرق نباشد و این فیض مطلق از مرایای موجودات قطع
فیض نماید ابداً شیئی بر وجه ارض موجود نخواهی یافت و دیگر آنکه وجود مرایای
احدیّه بنفسه لنفسه نخواهد بود چنانچه اینمطلب در جمیع الواح و زبر مندرج و
ثبت شده و بر هیچ عاقل روحانی و بالغ نورانی مستور نیست بلکه وجود مرایا
بوجود شمس است خواه شاعر باین رتبهٴ عالی بوده و خواه غافل چه که عنایت آن
شمس مشرقه از افق احدیّه من دون جهة بوده و خواهد بود و اینمقام رحمت
منبسطهٴ ربّانیّه است که من دون سؤال و طلب و یا عرفان و ایقان عنایت
میفرماید غایت اینست که بعضی از مرایا که این رتبهٴ اعلی و مقام ابهی را
ادراک نموده‌اند در رتبه و مقام اعلای از دون خود خواهند بود

\end{otherlanguage}

Va ʿalāve bar ān agar ān shams mushriq nabāshad va īn fayḍ-e muṭlaq az
marāyā-ye mawjūdāt qaṭʿ-e fayḍ namāyad abadan shayʾī bar vajh-e arḍ
mawjūd nakhāhī yāft. Va dīgar ānke vujūd-e marāyā-ye aḥadiyye bi-nafsihi
li-nafsihi nakhāhad būd chinānchi īn maṭlab dar jamīʿ-e alvāḥ va zubur
mundarij va thabt shude va bar hīch ʿāqil-e rūḥānī va bāligh-e nūrānī
mastūr nīst balke vujūd-e marāyā bi-vujūd-e shams ast khvāh shāʿir bi-īn
rutbe-ye ʿālī būde va khvāh ghāfil chi ke ʿināyat-e ān shams-e mushriqe
az ufuq-e aḥadiyye min dūni jihatin būde va khāhad būd. Va īn maqām-e
raḥmat-e munbasiṭe-ye rabbāniyye ast ke min dūn-e suʾāl va ṭalab va yā
ʿirfān va īqān ʿināyat mī-farmāyad. Ghāyat īnast ke baʿḍī az marāyā ke
īn rutbe-ye aʿlā va maqām-e abhā rā idrāk namūde-and dar rutbe va
maqām-e aʿlāy az dūn-e khud khāhand būd.

Moreover, if that Sun did not shine and this absolute grace were to
withhold its bounty from the mirrors of created things, you would find
nothing whatsoever existing on the face of the earth. Furthermore, the
mirrors of oneness do not exist through themselves or for themselves.
This matter has been recorded and inscribed in all the Tablets and
Scriptures and is hidden from no spiritually discerning or enlightened,
mature soul. Rather, the mirrors exist through the existence of the Sun,
whether they are aware of this lofty degree or heedless of it, for the
bounty of that Sun shining from the horizon of oneness has always been,
and will always be, unconditional. This is the station of God's
all-embracing mercy, which He bestows without entreaty or seeking,
recognition or certitude. The only distinction is that those mirrors
that have perceived this most exalted degree and all-glorious station
will occupy a higher rank and station than the others.

\subsection{10}\label{section-2603}

\begin{otherlanguage}{persian}

و دیگر آنکه ملاحظه نما که اگر شمس موجود نباشد مرایا ابداً وجود نخواهند
داشت و لایزال در حجب عدم خواهند بود و چون شمس قدم از افق قدسی طالع جمیع
مرایا از او حکایت مینمایند بعد از مقابله و قرب و اگر شمس غروب نماید
البتّه صور منطبعه در مرایا که صور شبحیّه‌اند غروب خواهند نمود و اینمطلب
اظهر از آن است که هیچ عاقل کامل ادراک ننماید

\end{otherlanguage}

Va dīgar ānke mulāḥaẓe namā ke agar shams mawjūd nabāshad marāyā abadan
vujūd nakhāhand dāsht va lā-yazāl dar ḥujub-e ʿadam khāhand būd va chūn
shams-e qidam az ufuq-e qudsī ṭāliʿ jamīʿ-e marāyā az ū ḥikāyat
mī-namāyand baʿd az muqābile va qurb. Va agar shams ghurūb namāyad
albatte ṣuvar-e munṭabiʿe dar marāyā ke ṣuvar-e shabaḥiyye-and ghurūb
khāhand namūd va īn maṭlab aẓhar az ān ast ke hīch ʿāqil-e kāmil idrāk
nanamāyad.

Consider further that, if the sun did not exist, the mirrors would have
no existence whatsoever and would remain forever behind the veils of
nonexistence. When the Sun of Eternity rises from the horizon of
holiness, all the mirrors reflect it once they face it and draw near. If
the sun sets, the images imprinted in the mirrors, which are reflected
likenesses, will certainly set as well. This matter is too evident for
anyone of sound understanding to fail to perceive.

\subsection{11}\label{section-2604}

\begin{otherlanguage}{persian}

و هر نفس منکر عنود که این لطیفهٴ مشهود را انکار نماید در شمس ظاهره و
مرایای مستحکیهٴ او ملاحظه نماید که تا مقابل شمس واقع است تجلّی و ظهور در
مرایا ملحوظ و بعد از غروب شمس اثری نخواهد یافت بلکه جمیع را معدوم محض و
فانی مطلق و مفقود صرف مشاهده خواهد نمود باری ای سایل چشم بصیرت معنویّه که
ودیعهٴ ربّانیّه است از غبار کلمات مؤتفکهٴ ظلمانیّه پاک نما تا آنکه شمس
صمدانی را حال در علوّ مقام خود که وسط سماء عدل است ملاحظه نمائی تاللّه قد
قرّت عیناک ان رأیتها و فزت بجمالها

\end{otherlanguage}

Va har nafs-e munkir-e ʿanūd ke īn laṭīfe-ye mashhūd rā inkār namāyad
dar shams-e ẓāhire va marāyā-ye mustaḥkiye-ye ū mulāḥaẓe namāyad ke tā
muqābil-e shams vāqiʿ ast tajallī va ẓuhūr dar marāyā malḥūẓ va baʿd az
ghurūb-e shams atharī nakhāhad yāft balke jamīʿ rā maʿdūm-e maḥḍ va
fānī-ye muṭlaq va mafqūd-e ṣirf mushāhade khāhad namūd. Bārī ay sāʾil
chashm-e baṣīrat-e maʿnaviyye ke vadīʿe-ye rabbāniyye ast az ghubār-e
kalimāt-e muʾtafake-ye ẓulmāniyye pāk namā tā ānke shams-e ṣamadānī rā
ḥāl dar ʿuluvv-e maqām-e khud ke vasaṭ-e samā-ye ʿadl ast mulāḥaẓe
namāʾī. Ta-Allāhi qad qarrat ʿaynāka in raʾaytahā wa-fuzta bi-jamālihā.

Let any obstinate denier who rejects this evident subtle truth consider
the visible sun and the mirrors that reflect it: as long as a mirror
faces the sun, its radiance and appearance can be seen in the mirrors,
but after the sun has set, no trace will be found. Rather, all will be
seen to be utterly nonexistent, wholly vanished, and entirely absent. In
short, O seeker, cleanse the eye of spiritual insight, which is a divine
trust, of the dust of dark, false words, so that you may now behold the
Sun of Eternal Self-Sufficiency in the loftiness of its own station, at
the center of the heaven of justice. By God, your eyes would indeed be
gladdened if you beheld it and attained its beauty.

\subsection{12}\label{section-2605}

\begin{otherlanguage}{persian}

بلی اینکه مسطور در الواح منزلهٴ بیان است حقّ لا ریب فیه و انّا کلّ به
مؤمنون چنانچه آن شمس الهیّه و ساذج قدمیّه روح ما سواه فداه میفرماید که
شموس حقیقت در هر کوری و دوری از افق قدس مشرق و ظاهر میشوند مقصود ظهور
کلّیّه است یعنی در نزول آیات و انزال کلمات و تغییر شریعت و تجدّد امور ملّت
نظر باینمقام است که از آن ساذج فطرت امثال این بیانات ظاهر شده

\end{otherlanguage}

Balī īnke masṭūr dar alvāḥ-e munazzale-ye Bayān ast ḥaqqun lā rayba fīhi
wa-innā kullun bihi muʾminūn. Chinānchi ān shams-e ilāhiyye va sādhij-e
qidamiyye rūḥu mā siwāhu fidāhu mī-farmāyad ke shumūs-e ḥaqīqat dar har
kawrī va dawrī az ufuq-e quds mushriq va ẓāhir mī-shavand maqṣūd ẓuhūr-e
kulliyye ast yaʿnī dar nuzūl-e āyāt va inzāl-e kalimāt va taghyīr-e
sharīʿat va tajaddud-e umūr-e millat. Naẓar bi-īn maqām ast ke az ān
sādhij-e fiṭrat amthāl-e īn bayānāt ẓāhir shude.

Yes, what is recorded in the revealed Tablets of the Bayan is truth in
which there is no doubt, and we all believe in it. When that divine Sun
and pure Essence of Eternity---may the souls of all else be sacrificed
for Him---says that the Suns of Truth rise and appear from the horizon
of holiness in every cycle and age, what is intended is universal
manifestation: namely, the revelation of verses, the sending down of
words, the changing of religious law, and the renewal of the affairs of
the faith. It is with reference to this station that such statements
have issued from that Essence of pristine nature.

\subsection{13}\label{section-2606}

\begin{otherlanguage}{persian}

و قسم بآن جمال معنوی و غیب قدس صمدانی که اگر ناس را بالغ مشاهده
میفرمودند ذکر امثال این بیانات نمیفرمودند چه که آن مقام مقدّس از ذکر کلّی
و جزئی و ظهور و خفا بوده چه که تغییر و تبدیل در عالم حدود واقع است و آن
جوهر مقصود در مقرّ قدس واقف فسبحان اللّه عمّا انتم تظنّون و چه مقدار از
شمسهای عنایت بی‌زوال که مشرق شده و ابداً تغییر نفرموده و تبدیل امر
ننموده‌اند و ناس را بهمان شریعت جاریهٴ قبل دعوت مینمودند چه کلّ من عند
اللّه بوده و جمیع الی اللّه راجع و من اللّه ظاهر

\end{otherlanguage}

Va qasam bi-ān jamāl-e maʿnavī va ghayb-e quds-e ṣamadānī ke agar nās rā
bāligh mushāhade mī-farmūdand dhikr-e amthāl-e īn bayānāt namī-farmūdand
chi ke ān maqām muqaddas az dhikr-e kullī va juzʾī va ẓuhūr va khafā
būde chi ke taghyīr va tabdīl dar ʿālam-e ḥudūd vāqiʿ ast va ān jawhar-e
maqṣūd dar maqarr-e quds vāqif. Fa-subḥāna Allāhi ʿammā antum taẓunnūn.
Va chi miqdār az shams-hā-ye ʿināyat-e bī-zavāl ke mushriq shude va
abadan taghyīr nafarmūde va tabdīl-e amr nanamūde-and va nās rā bi-hamān
sharīʿat-e jāriye-ye qabl daʿvat mī-namūdand chi kull min ʿindi Allāh
būde va jamīʿ ilā Allāh rājiʿ va min Allāh ẓāhir.

I swear by that spiritual Beauty and that unseen holiness of the
Eternally Self-Sufficient that, had He found the people mature, He would
not have uttered such statements. For that station is sanctified above
any mention of universal and particular, manifestation and concealment,
since change and alteration occur in the realm of limitations, while
that sought-after Essence abides in the seat of holiness. Glorified,
then, is God above what you imagine. How many suns of unfading
providence have risen without ever making a change or altering the
religious order, instead summoning people to the same religious law that
had previously been in force! For all are from God, all return to God,
and all appear from God.

\subsection{14}\label{section-2607}

\begin{otherlanguage}{arabic}

باری نظر از ملاحظهٴ ارض و اهل آن بردار و بر منظر اکبر و مقام اظهر ملاحظه
کن تا شمس ربّانی را در قمیص ستر و حجاب مشاهده نمائی اگرچه لم‌یزل چون شمس
سماء ظاهر و درّیّ و منیر بوده ولکن سنین معدوده میگذرد که بحجاب حسد و بغضا
از انظر عارفین مستور مانده تا چه رسد بمحدودین و محتجبین باری در هر حال
بنظر دقیق در بحر عمیق معانی سبّاح شو تا بر لآلی بحر احدیّه فایز گردی و
مستغنی شوی از کلّ ما سوی

\end{otherlanguage}

Bārī naẓar az mulāḥaẓe-ye arḍ va ahl-e ān bar-dār va bar manẓar-e akbar
va maqām-e aẓhar mulāḥaẓe kun tā shams-e rabbānī rā dar qamīṣ-e sitr va
ḥijāb mushāhade namāʾī. Agarchi lam-yazal chūn shams-e samā ẓāhir va
durrī va munīr būde valīkin sinīn-e maʿdūde mī-gudharad ke bi-ḥijāb-e
ḥasad va baghḍā az anẓar-e ʿārifīn mastūr mānde tā chi rasad bi-maḥdūdīn
va muḥtajibīn. Bārī dar har ḥāl bi-naẓar-e daqīq dar baḥr-e ʿamīq-e
maʿānī sabbāḥ shaw tā bar laʾālī-ye baḥr-e aḥadiyye fāʾiz gardī va
mustaghnī shavī az kull-e mā sivā.

Turn your gaze, then, away from the earth and its inhabitants, and look
toward the Most Great Scene and the most manifest Station, so that you
may behold the divine Sun clothed in concealment and veiling. Although
it has always been as manifest, brilliant, and luminous as the sun in
the sky, several years have now passed during which it has remained
hidden by the veil of envy and hatred from the eyes of those endowed
with recognition, let alone those who are confined and veiled. In every
circumstance, then, swim with keen perception in the deep ocean of
meanings, so that you may attain the pearls of the ocean of oneness and
become independent of all else.

\subsection{15}\label{section-2608}

\begin{otherlanguage}{persian}

و دیگر آنکه تمام زمام بصر را بغیر مده و بظاهر بعضی عبارات و اشارات کفایت
مکن چه که جمیع این الفاظ ظاهره و اشارات قدسیّه و دلالات عزّیّه در هر کوری
تغییر مینماید و از همین تغییر و تبدیل و تجدید محقّق میشود که فوق هر بیانی
در علم الهی مکنون بوده و اعلی از هر حکمتی در کنز حکمت ربّانی مخزون خواهد
بود چنانچه در امم سابقه ملاحظه شد که حکم کلمات و طرز بیانات توریة بعد از
ظهور نغمات روح القدس تغییر یافت و تبدیل نمود و هم‌چنین تا بظهور نقطهٴ
بیان و طلعت سبحان روح من فی الملک فداه رسید جمیع جواهر و لآلی اصداف بحور
فرقانیّه را از قعر آن برآورده بهیاکل جدیدهٴ منیعه مزیّن فرمود

\end{otherlanguage}

Va dīgar ānkih tamām-i zimām-i baṣar rā bi-ghayr madih va bi-ẓāhir-i
baʿḍī ʿibārāt va ishārāt kifāyat makun chih kih jamīʿ-i īn alfāẓ-i
ẓāhirih va ishārāt-i qudsīyyih va dalālāt-i ʿizzīyyih dar har kūrī
taghyīr mīnamāyad va az hamīn taghyīr va tabdīl va tajdīd muḥaqqaq
mīshavad kih fawq-i har bayānī dar ʿilm-i ilāhī maknūn būdih va aʿlā az
har ḥikmatī dar kanz-i ḥikmat-i rabbānī makhzūn khvāhad būd chunānchih
dar umam-i sābiqih mulāḥiẓih shud kih ḥikam-i kalimāt va ṭarz-i
bayānāt-i Tawrāt baʿd az ẓuhūr-i naghamāt-i Rūḥu'l-Qudus taghyīr yāft va
tabdīl namūd va hamchunīn tā bi-ẓuhūr-i Nuqṭiy-i Bayān va Ṭalʿat-i
Subḥān rūḥu man fī'l-mulki fidāh rasīd jamīʿ-i javāhir va laʾālī-yi
aṣdāf-i buḥūr-i furqānīyyih rā az qaʿr-i ān bar-āvurdih bi-hayākil-i
jadīdiy-i manīʿih muzayyan farmūd.

Furthermore, do not surrender the entire reins of your vision to others,
or content yourself with the outward meaning of certain expressions and
allusions. For all these outward words, holy allusions, and exalted
indications change in every cycle. This very change, transformation, and
renewal establish that beyond every utterance something has lain
concealed in divine knowledge, and above every wisdom something will
remain stored in the treasury of heavenly wisdom. Thus it was observed
among former peoples that the wisdom of the words and the manner of
expression in the Torah changed and were transformed after the melodies
of the Holy Spirit appeared. So it continued until the advent of the
Point of the Bayán and the Countenance of the All-Glorious---may the
souls of all in the kingdom be sacrificed for Him. He brought forth all
the gems and pearls within the shells of the oceans of the Qur'an from
their depths and adorned them with new and exalted forms.

\subsection{16}\label{section-2609}

\begin{otherlanguage}{persian}

و هم‌چنین در عهد سلطان قدم و ملیک اعظم من یظهره اللّه جمیع این کلمات منیعه
و عبارات بدیعه تغییر خواهد یافت گویا ملاحظه می‌شود که آن صرف احدیّه و جوهر
عزّ قدمیّه چون از مشرق لاهوتیّه من دون کیف و این وجهة و اشارة ظاهر شوند و
بر جنّت اعلی و عرش ابهی که همان ارض ساکنه که مقرّ آن شمس قدم است مقرّ یابند
جمیع اسرار مودعه در بیان را ظاهر فرماید و جز آن جمال قدم احدی بر لآلی
مستوره در این بحر اعظم کما هو حقّه اطّلاع نیابد و او است مظهر و مبیّن اسرار
و علوم مقدّره در بیان

\end{otherlanguage}

Va hamchunīn dar ʿahd-i Sulṭān-i Qidam va Malīk-i Aʿẓam Man Yuẓhiruhu
Allāh jamīʿ-i īn kalimāt-i manīʿih va ʿibārāt-i badīʿih taghyīr khvāhad
yāft gūyā mulāḥiẓih mīshavad kih ān ṣirf-i aḥadīyyih va jawhar-i ʿizz-i
qidamīyyih chūn az mashriq-i lāhūtīyyih min dūni kayf wa-ayn wijha
wa-ishāra ẓāhir shavand va bar jannat-i aʿlā va ʿarsh-i abhā kih hamān
arḍ-i sā kinih kih ma qarr-i ān shams-i qidam ast ma qarr yāband jamīʿ-i
asrār-i mūdiʿih dar Bayān rā ẓāhir farmāyad va juz ān Jamāl-i Qidam
aḥadī bar laʾālī-yi mastūrih dar īn baḥr-i aʿẓam kamā huwa ḥaqquhu
iṭṭilāʿ nayābad va ū ast maẓhar va mubayyin-i asrār va ʿulūm-i
muqaddarih dar Bayān.

Likewise, in the age of the Sovereign of Eternity and the Most Great
King, Him Whom God will make manifest, all these exalted words and
wondrous expressions will change. It is as though one can behold that
pure Essence of Oneness and that Essence of ancient glory appearing from
the dayspring of divinity, without how or where, direction or
indication, and taking His seat in the highest Paradise and upon the
most glorious Throne---that same still earth which is the dwelling place
of that Sun of Eternity. He will reveal all the mysteries deposited in
the Bayán, and no one but that Ancient Beauty will know, as they deserve
to be known, the pearls concealed in this most great Ocean. He is the
Revealer and Interpreter of the mysteries and knowledge ordained in the
Bayán.

\subsection{17}\label{section-2610}

\begin{otherlanguage}{persian}

بلی مرایای منطبعه ادراک نمایند علی قدر مراتبهم و استعداداتهم لا علی ما
هو علیه لهذا احدی قادر بر تقرّب بسدرهٴ معنی در طور بیان نخواهد بود الّا من
شآء ربّک و لهذا ذکر شده که بظاهر بعضی عبارات من عند خود کفایت مکن چه که
عقول قاصر است از ادراک آن والّا حرفی از این کتب منزله نزد این عبد اعزّ است
از کلّ من فی السّموات و الأرض و لا یعرف ذلک الّا اولی الألباب

\end{otherlanguage}

Balī marāyā-yi munṭabiʿih idrāk namāyand ʿalā qadri marātibihim
wa-istiʿdādātihim lā ʿalā mā huwa ʿalayhi li-hādhā aḥadī qādir bar
taqarrub bi-Sidriy-i Maʿnī dar Ṭūr-i Bayān nakhvāhad būd illā man shāʾa
rabbuka va li-hādhā dhikr shudih kih bi-ẓāhir-i baʿḍī ʿibārāt min ʿindi
khud kifāyat makun chih kih ʿuqūl qāṣir ast az idrāk-i ān va-illā ḥarfī
az īn kutub-i munzalih nazd-i īn ʿabd aʿazz ast az kulli man
fī'l-samāwāti wa'l-arḍ wa-lā yaʿrifu dhālika illā ulū'l-albāb.

Yes, the mirrors that receive His image will perceive according to their
ranks and capacities, not as He is in Himself. Therefore no one will be
able to draw near to the Tree of Meaning upon the Sinai of the Bayán
except those whom your Lord wills. This is why it has been said: do not,
of your own accord, content yourself with the outward meaning of certain
expressions, for minds fall short of comprehending it. Otherwise, a
single letter of these revealed Books is dearer to this servant than all
who are in the heavens and on earth. None knows this except those
endowed with understanding.

\subsection{18}\label{section-2611}

\begin{otherlanguage}{persian}

و هم‌چنین آن ذات مقدّس این عبارات مرتفعهٴ الیوم را مرتفع نمایند مثلاً ذکر
شمس که در بیان اعظم ذکر و اکبر بیان است شاید این ذکر و رتبه را بیکی از
عباد خود عنایت فرمایند من دون استحقاق بل بفضل و رحمت و کرم خود این فیض
اکبر و فضل اعظم را مبذول میفرمایند و نفس شریف خود را که محلّ ظهور نار
شجرهٴ اوّلیّه و طلوع نور آخریّه است بسماء موسوم خواهند فرمود چه که شمس از
افق او طالع خواهد شد و بوجود او قائم خواهد بود و چنانچه الیوم در جمیع
بیان ذکر شمس ملاحظه می‌نمائی در الواح آن سماء قدس معانی ذکر سماء ملاحظه
خواهی نمود بهمان ذات غیب مقدّس و ساذج عزّ اقدس

\end{otherlanguage}

Va hamchunīn ān Dhāt-i Muqaddas īn ʿibārāt-i murtafiʿiy-i al-yawm rā
murtafiʿ namāyand mathalan dhikr-i shams kih dar Bayān aʿẓam-i dhikr va
akbar-i bayān ast shāyad īn dhikr va rutbih rā bi-yikī az ʿibād-i khud
ʿināyat farmāyand min dūni istiḥqāq bal bi-faḍl va raḥmat va karam-i
khud īn fayḍ-i akbar va faḍl-i aʿẓam rā mabdhūl mīfarmāyand va nafs-i
sharīf-i khud rā kih maḥall-i ẓuhūr-i nār-i shajariy-i avvalīyyih va
ṭulūʿ-i nūr-i ākhirīyyih ast bi-samāʾ mawsūm khvāhand farmūd chih kih
shams az ufuq-i ū ṭāliʿ khvāhad shud va bi-vujūd-i ū qāʾim khvāhad būd
va chunānchih al-yawm dar jamīʿ-i Bayān dhikr-i shams mulāḥiẓih mīnamāʾī
dar alvāḥ-i ān Samāʾ-i Quds-i Maʿānī dhikr-i samāʾ mulāḥiẓih khvāhī
namūd bi-hamān Dhāt-i Ghayb-i Muqaddas va sādhij-i ʿizz-i aqdas.

Likewise, that holy Being will set aside these expressions that are
exalted today. For example, the designation ``Sun,'' which in the Bayán
is the greatest designation and the most sublime expression: He may
bestow this designation and rank upon one of His servants. He grants
this greatest bounty and most mighty grace without any deserving, but
through His own grace, mercy, and generosity. He will designate His own
noble Self, which is the place where the fire of the Primal Tree appears
and the light of the Last rises, as ``Heaven,'' for the Sun will rise
from His horizon and depend upon His existence. Just as today you find
mention of the Sun throughout the Bayán, so in the Tablets of that
Heaven of holy meanings you will find mention of Heaven. By that same
holy, unseen Essence and that purest, most sacred Glory!

\subsection{19}\label{section-2612}

\begin{otherlanguage}{persian}

اگر بخواهم تا آخر لا آخر بآنچه در اوّل ظهور مظاهر قدسیّه بآن تکلّم نمایند
ذکر نمایم قادرم من فضل ربّی ولکن چون در ظنون افکیّه مسجون بوده و هستم جایز
ندیده که تغنّیات ورقاء غیب صمدانی را ابلاغ نمایم چه که ذکر بلبل معنوی را
جمال گل باقی و رضوان عزّ احدی لایق نه در میان این انفس ضعیفهٴ مشرکه که در
حجبات وهمیّه چنان معدوم و مفقود شده‌اند که اگر ساذج غیب را باظهر بیانات
قدسی و ابهی تجلّیات الهی ملاحظه نمایند اعراض نمایند و اغماض کنند باری ای
سایل خطوهٴ شیطان را بقدم رحمن تبدیل کن و در ارض معارف منبسطه بقوّهٴ
مشاهده و مکاشفه مشی نما تا بسرمنزل یقین که مقام معرفت منزل این کلمات است
فایز شوی و بمصر احدیّه که نفس جنّت ابهائیّه است مستبصر گردی فطوبی للفائزین

\end{otherlanguage}

Agar bikhvāham tā ākhir-i lā ākhir bi-ānchih dar avval-i ẓuhūr-i
maẓāhir-i qudsīyyih bi-ān takallum namāyand dhikr namāyam qādiram min
faḍli rabbī va-lākin chūn dar ẓunūn-i ifkīyyih masjūn būdih va hastam
jāyiz nadīdih kih taghannīyāt-i Varqā-yi Ghayb-i Ṣamadānī rā iblāgh
namāyam chih kih dhikr-i bulbul-i maʿnavī rā jamāl-i gul-i bāqī va
riḍvān-i ʿizz-i aḥadī lāyiq nah dar miyān-i īn anfus-i ḍaʿīfiy-i
mushrikih kih dar ḥujubāt-i vahmīyyih chunān maʿdūm va mafqūd shudih-and
kih agar sādhij-i ghayb rā bi-aẓhar-i bayānāt-i qudsī va abhā
tajallīyāt-i ilāhī mulāḥiẓih namāyand iʿrāḍ namāyand va ighmāḍ kunand
bārī ay sāʾil khuṭviy-i shayṭān rā bi-qadam-i Raḥmān tabdīl kun va dar
arḍ-i maʿārif-i munbasiṭih bi-quvviy-i mushāhadih va mukāshifih mashy
namā tā bi-sar-manzil-i yaqīn kih maqām-i maʿrifat-i Munzil-i īn kalimāt
ast fāʾiz shavī va bi-Miṣr-i Aḥadīyyih kih nafs-i Jannat-i Abhāʾīyyih
ast mustabṣir gardī fa-ṭūbā li'l-fāʾizīn.

If I wished to speak until the end that has no end of what the holy
Manifestations will utter at the beginning of their appearance, I could
do so through the grace of my Lord. Yet because I have been and remain
imprisoned amid false suspicions, I have not deemed it fitting to convey
the songs of the Dove of the unseen, eternal realm. For the beauty of
the everlasting Rose and the paradise of the glory of Oneness befit the
song of the spiritual Nightingale---not the company of these weak,
idolatrous souls, who are so utterly lost and effaced within the veils
of their imaginings that, even if they beheld the pure unseen Essence
with the clearest holy utterances and the most glorious divine
revelations, they would turn away and close their eyes. In brief, O
questioner, exchange the step of Satan for the tread of the
All-Merciful, and walk through the outspread land of understanding by
the power of direct vision and spiritual disclosure, so that you may
attain the destination of certitude, which is the station of recognizing
the Revealer of these words, and gain insight into the City of Oneness,
which is the most glorious Paradise itself. Blessed, then, are those who
attain!

\subsection{20}\label{section-2613}

\begin{otherlanguage}{persian}

و این طلوع و غروب که در مظاهر قدمیّه ذکر شده نه این طلوع و غروب است که
این نفوس ادراک مینمایند زیرا که این طلوع و غروب تحت بصر واقع است و آن
جوهر مقصود فوق نظر در منظر اکبر واقف بلکه مقصود از غروب عدم ادراک انظار
محتجبین و ابصار محدودین است مثلاً شمس ظاهره اگر در خلف غمام واقع شود
ابصار ظاهره از مشاهدهٴ او ممنوع میشوند و در نظر شمس غایب است اگرچه در
مقام خود درّیّ و روشن خواهد بود ولکن بصر صاحبان بصر معنوی را غمام منع
ننماید و بسبحات محتجب نشوند و از ورای صدهزار حجب و ستر مشاهده می‌نمایند

\end{otherlanguage}

Va īn ṭulūʿ va ghurūb kih dar maẓāhir-i qidamīyyih dhikr shudih nah īn
ṭulūʿ va ghurūb ast kih īn nufūs idrāk mīnamāyand zīrā kih īn ṭulūʿ va
ghurūb taḥt-i baṣar vāqiʿ ast va ān jawhar-i maqṣūd fawq-i naẓar dar
manẓar-i akbar vāqif balkih maqṣūd az ghurūb ʿadam-i idrāk-i anẓār-i
muḥtajibīn va abṣār-i maḥdūdīn ast mathalan shams-i ẓāhirih agar dar
khalf-i ghamām vāqiʿ shavad abṣār-i ẓāhirih az mushāhadiy-i ū mamnūʿ
mīshavand va dar naẓar shams ghāʾib ast agarchih dar maqām-i khud durrī
va rawshan khvāhad būd va-lākin baṣar-i ṣāḥibān-i baṣar-i maʿnavī rā
ghamām manʿ nanamāyad va bi-subuḥāt muḥtajib nashavand va az varā-yi
ṣad-hizār ḥujub va sitr mushāhidih mīnamāyand.

The rising and setting mentioned in connection with the eternal
Manifestations are not the rising and setting that these souls
understand. For this rising and setting fall within the compass of
sight, whereas that sought-after Essence stands beyond vision in the
Most Great Scene. Rather, ``setting'' means that the eyes of the veiled
and the vision of those confined by limitations fail to perceive. For
example, if the outward sun passes behind a cloud, outward eyes are
prevented from beholding it, and to their sight the sun is absent,
although in its own station it remains radiant and bright. But clouds do
not obstruct the vision of those endowed with spiritual sight; they are
not shut out by veils of glory, and they behold it from beyond a hundred
thousand veils and coverings.

\subsection{21}\label{section-2614}

\begin{otherlanguage}{persian}

ملاحظه فرما که اکثر اهل ارض با علوّ ارتفاع شمس مشرقه در بیان ابداً ملتفت
نشده و در ظلمت باقی و مستور و محتجب بوده و خواهند بود و هم‌چنین اهل بیان
را ملاحظه نما که با علوّ رفعت و سموّ عظمت این امر مرتفع اعلی که در سماء
قدس ابهی ظاهر و لایح است احدی ادراک ننموده و کلّ بهوای خود سالک بوده و
خواهند بود و حال این عدم التفات دلیل بر عدم وجود آن شمس قدم نخواهد بود
با آنکه آن قلم اعلی و جوهر ابهی و ساذج اسنی در جمیع کتب منزله و صحف
مرسله وصیّت فرموده که اگر نفسی و یا کینونتی و یا نوری و ظهوری از افق مشیّت
الهی طالع شود بهیچ امری محتجب نشده بآن ظلّ مبارک که نفس ظلّ اللّه است وارد
شوند و جمیع بیان و اعمال مقرّره و افعال مقدّره در آن را منتهی بمعرفت آن
نفس مقدّسه فرموده‌اند مع‌ذلک بسبحات ظنون و نفس و وهم چنان محتجب مانده‌اند که
ابداً ارادهٴ شطّ الهی و بحر قدمی و رضوان باقی ربّانی ننموده بلکه ارادهٴ
معرفت آن ذات مقدّس که عین معرفة اللّه بوده ننموده‌اند و جمیع خود را بر مقرّ
ایمان و مقعد ایقان ساکن میدانند و غافل از آنکه الیوم عنداللّه مذکور نبوده
و نخواهند بود لن ینظرهم اللّه و لن یتکلّم معهم و لن یلتفت الیهم ابداً با
آنکه جمیع حجّت بأعلی ظهور خود و منتهی صعود ظاهر شده بقسمیکه احدی را مجال
اعراض و اغماض نمانده مگر آنکه بالمرّه از جمیع کتب منزله و رسل مرسله اعراض
نمایند

\end{otherlanguage}

Mulāḥiẓih farmā kih akthar-i ahl-i arḍ bā ʿuluvv-i irtifāʿ-i shams-i
mushriqih dar Bayān abadan multafit nashudih va dar ẓulmat bāqī va
mastūr va muḥtajib būdih va khvāhand būd va hamchunīn ahl-i Bayān rā
mulāḥiẓih namā kih bā ʿuluvv-i rifʿat va sumuvv-i ʿaẓamat-i īn amr-i
murtafiʿ-i aʿlā kih dar samāʾ-i quds-i abhā ẓāhir va lāyiḥ ast aḥadī
idrāk nanamūdih va kull bi-havā-yi khud sālik būdih va khvāhand būd va
ḥāl īn ʿadam-i iltifāt dalīl bar ʿadam-i vujūd-i ān shams-i qidam
nakhvāhad būd bā ānkih ān Qalam-i Aʿlā va Jawhar-i Abhā va Sādhij-i Asnā
dar jamīʿ-i kutub-i munzalih va ṣuḥuf-i mursalih vaṣīyyat farmūdih kih
agar nafsī va yā kaynūnatī va yā nūrī va ẓuhūrī az ufuq-i mashīyyat-i
ilāhī ṭāliʿ shavad bi-hīch amrī muḥtajib nashudih bi-ān ẓill-i mubārak
kih nafs-i Ẓillu'llāh ast vārid shavand va jamīʿ-i Bayān va aʿmāl-i
muqarrarih va afʿāl-i muqaddarih dar ān rā muntahī bi-maʿrifat-i ān
nafs-i muqaddasih farmūdih-and maʿa-dhālik bi-subuḥāt-i ẓunūn va nafs va
vahm chunān muḥtajib māndih-and kih abadan irādiy-i shaṭṭ-i ilāhī va
baḥr-i qidamī va riḍvān-i bāqī-yi rabbānī nanamūdih balkih irādiy-i
maʿrifat-i ān dhāt-i muqaddas kih ʿayn-i maʿrifatu'llāh būdih
nanamūdih-and va jamīʿ khud rā bar ma qarr-i īmān va maqʿad-i īqān sākin
mīdānand va ghāfil az ānkih al-yawm ʿinda'llāh madhkūr nabūdih va
nakhvāhand būd lan yanẓurahumu Allāhu wa-lan yatakallama maʿahum wa-lan
yaltafita ilayhim abadan bā ānkih jamīʿ-i ḥujjat bi-aʿlā ẓuhūr-i khud va
muntahā ṣuʿūd ẓāhir shudih bi-qismīkih aḥadī rā majāl-i iʿrāḍ va ighmāḍ
namāndih magar ānkih bi'l-marrih az jamīʿ-i kutub-i munzalih va rusul-i
mursalih iʿrāḍ namāyand.

Consider how most of the people of the earth, despite the lofty ascent
of the Sun shining in the Bayán, have paid no attention at all, and have
remained and will remain in darkness, concealed and veiled. Likewise,
consider the people of the Bayán: despite the lofty elevation and
surpassing grandeur of this exalted, most sublime Cause, which is
manifest and resplendent in the heaven of most glorious holiness, no one
has perceived it, and all have followed and will follow their own
desires. Yet this lack of attention is no proof that that Sun of
Eternity does not exist. That Most Exalted Pen, that Most Glorious
Essence, that purest and most resplendent Being, has enjoined in all the
revealed Books and dispatched Scriptures that, should a soul or a being,
a light or a revelation, rise from the horizon of the divine Will, they
should allow nothing to veil them and should enter beneath that blessed
shadow, which is the very shadow of God. He has made the whole Bayán,
and the prescribed deeds and ordained acts within it, culminate in the
recognition of that holy Soul. Nevertheless, they have remained so
thoroughly veiled by their suspicions, selves, and imaginings that they
have never sought the divine River, the Ocean of Eternity, or the
everlasting heavenly Paradise. Indeed, they have not sought to recognize
that holy Being, recognition of whom is the very recognition of God.
They all consider themselves established in the abode of faith and the
seat of certitude, unaware that today they have not been and will not be
remembered before God. God will never look upon them, speak with them,
or turn toward them, although the proof in its entirety has appeared in
its fullest manifestation and at its highest ascent, leaving no one any
room to turn away or close their eyes unless they turn away completely
from all the revealed Books and all the Messengers sent forth.

\subsection{22}\label{section-2615}

\begin{otherlanguage}{persian}

ای عباد ناقابل غیر بالغ اگر قطره‌ئی از سحاب مرحمت بر جمیع ارض مبذول شود
جمیع هیاکل سدرهای وجود را بنار سبحانی و بیان انا اللّهی مشهود بینی و اگر
رشحی از این بحر اعظم بر جمیع مرایای موجودات افاضه شود همهٴ موجودات را
بطراز عزّ قدسی مشاهده نمائید باری پرواز این باز الهی بلندتر از آن است که
طیور ارضی بدو رسند و نغمهٴ این ورقاء عزّ قدسی از آن لطیف‌تر است که بگوش
این بیهوشان ارض فانیه آلوده شود و اللّه یقول الحقّ ولکن النّاس هم لا یسمعون

\end{otherlanguage}

Ay ʿibād-i nā-qābil-i ghayr-i bāligh agar qaṭrih-ī az saḥāb-i marḥamat
bar jamīʿ-i arḍ mabdhūl shavad jamīʿ-i hayākil-i sidrihā-yi vujūd rā
bi-nār-i subḥānī va bayān-i anā Allāhī mashhūd bīnī va agar rashḥī az īn
baḥr-i aʿẓam bar jamīʿ-i marāyā-yi mawjūdāt ifāḍih shavad hamiy-i
mawjūdāt rā bi-ṭirāz-i ʿizz-i qudsī mushāhidih namāʾīd bārī parvāz-i īn
bāz-i ilāhī bulandtar az ān ast kih ṭuyūr-i arḍī bidū rasand va
naghmiy-i īn varqā-yi ʿizz-i qudsī az ān laṭīftar ast kih bi-gūsh-i īn
bī-hūshān-i arḍ-i fāniyih ālūdih shavad wa-Allāhu yaqūlu'l-ḥaqqa
wa-lākinna'l-nāsa hum lā yasmaʿūn.

O unworthy, immature servants! If a single drop from the cloud of mercy
were bestowed upon the whole earth, you would behold all the forms of
the trees of existence bearing the divine fire and the utterance ``I am
God.'' If a sprinkling from this most great Ocean were poured upon all
the mirrors of created things, you would behold all beings adorned with
holy glory. In brief, the flight of this divine Falcon is too high for
earthly birds to reach, and the melody of this Dove of holy glory is too
subtle to be sullied by the ears of these heedless inhabitants of the
perishable earth. God speaks the truth, but the people do not hear.

\subsection{23}\label{section-2616}

\begin{otherlanguage}{persian}

بلی از این فیوضات نازلهٴ غیر مقطوعه جمیع امکان و اکوان مستعدّ شده دیگر تا
چه وقت و زمان بمقام قرب بین یدی الرّحمن که منتهی رتبهٴ عرش ظهور و بطون
است حاضر شوند خواه بقلوب و خواه باجساد و شاید که نسیم قدس لا یعرف از
رضوان عزّ لا یدرک بوزد و انفس بعیده را که در ادنی رتبهٴ حدود ساکن شده‌اند
باین مقرّ اعلی و رفرف قصوی کشاند که شاید از این فیض بدیع و فضل منیع محروم
نگردند و اللّه یرزق من یشآء و ما انا بمرزق الموتی و انّه لهو الرّزّاق المعطی
الکریم

\end{otherlanguage}

Balī az īn fuyūḍāt-i nāziliy-i ghayr-i maqṭūʿih jamīʿ-i imkān va akvān
mustaʿidd shudih dīgar tā chih vaqt va zamān bi-maqām-i qurb bayna yaday
al-Raḥmān kih muntahā rutbiy-i ʿarsh-i ẓuhūr va buṭūn ast ḥāḍir shavand
khvāh bi-qulūb va khvāh bi-ajsād va shāyad kih nasīm-i quds-i lā yuʿraf
az riḍvān-i ʿizz-i lā yudrak bivazad va anfus-i baʿīdih rā kih dar adnā
rutbiy-i ḥudūd sākin shudih-and bi-īn ma qarr-i aʿlā va rafraf-i quṣvā
kishānad kih shāyad az īn fayḍ-i badīʿ va faḍl-i manīʿ maḥrūm nagardand
wa-Allāhu yarzuqu man yashāʾu wa-mā anā bi-murziqi'l-mawtā wa-innahu
la-huwa'l-Razzāqu'l-Muʿṭī al-Karīm.

Yes, through these unceasing outpourings of bounty, all contingent
existence and all beings have been made ready. It remains to be seen at
what time and hour they will present themselves, whether in heart or in
body, in the station of nearness before the All-Merciful, which is the
highest rank of the throne of manifestation and concealment. Perhaps the
holy breeze of the unknowable will blow from the paradise of
inaccessible glory and draw those distant souls, who dwell in the lowest
degree of limitation, to this most exalted abode and this loftiest
retreat, so that they may not be deprived of this wondrous bounty and
exalted grace. God provides for whomever He wills, but I am not one who
provides for the dead. He, indeed, is the Provider, the Bestower, the
Generous.

\subsection{24}\label{section-2617}

\begin{otherlanguage}{persian}

و دیگر آنکه سؤال از تعدّد مطالع احدیّه و مظاهر قدسیّه نموده بودی مشهود
آنجناب بوده که امر الهی و مشیّت سبحانی و ارادهٴ ربّانی هرگز محدود بحدود
نبوده و نیست در بعضی از اعصار نفس واحده بوده و در بعضی از ازمان هیاکل
متعدّده از مشرق امر ظاهرشده‌اند چنانچه در جمیع کتب قبل و بعد مسطور و مبرهن
و مشهود است که اگر این عبد بخواهد تفصیل اظهار نماید کتابی مبسوط باید
تحریر شود

\end{otherlanguage}

Va dīgar ānkih suʾāl az taʿaddud-i maṭāliʿ-i aḥadīyyih va maẓāhir-i
qudsīyyih namūdih būdī mashhūd-i ān-janāb būdih kih amr-i ilāhī va
mashīyyat-i subḥānī va irādiy-i rabbānī hargiz maḥdūd bi-ḥudūd nabūdih
va nīst dar baʿḍī az aʿṣār nafs-i vāḥidih būdih va dar baʿḍī az azmān
hayākil-i mutaʿaddidih az mashriq-i amr ẓāhir shudih-and chunānchih dar
jamīʿ-i kutub-i qabl va baʿd masṭūr va mubarhan va mashhūd ast kih agar
īn ʿabd bikhvāhad tafṣīl iẓhār namāyad kitābī mabsūṭ bāyad taḥrīr
shavad.

Furthermore, you had asked about the plurality of the Daysprings of
Oneness and the holy Manifestations. You have seen that the Cause of
God, the divine Will, and the heavenly Purpose have never been, and are
not, confined within limits. In some ages there has been a single
Person, while at other times multiple Figures have appeared from the
dayspring of the Cause, as is recorded, demonstrated, and evident in all
the Books, earlier and later. Were this servant to wish to set forth the
details, an extensive book would have to be written.

\subsection{25}\label{section-2618}

\begin{otherlanguage}{persian}

و اگر در بعضی از بیانات مظاهر امریّه و مطالع قدسیّه تحدید بواحد ملاحظه
مینمائی بآن محتجب نشده باصل سدرهٴ امر و شجرهٴ مقدّسه ملاحظه نمائید که
لم‌یزل و لایزال امر اللّه من غیر حدّ و عدّ ظاهر شده بلکه این تحدیدات مرایای
مستحکیه مقصود حفظ و ستر امر اللّه است از انفس طاغوتیّه و هیاکل فرعونیّه که
تجرّی ننمایند و حصن عصمت امر را بخیانت نظر و قلب خود نیالایند و مظاهر
شیطانیّه بمقعد قدس سلیمانیّه قدم نگذارند

\end{otherlanguage}

Va agar dar baʿḍī az bayānāt-i maẓāhir-i amrīyyih va maṭāliʿ-i qudsīyyih
taḥdīd bi-vāḥid mulāḥiẓih mīnamāʾī bi-ān muḥtajib nashudih bi-aṣl-i
Sidriy-i Amr va Shajariy-i Muqaddasih mulāḥiẓih namāʾīd kih lam-yazal va
lā-yazāl amru'llāh min ghayri ḥadd va ʿadd ẓāhir shudih balkih īn
taḥdīdāt-i marāyā-yi mustaḥkiyih maqṣūd ḥifẓ va satr-i amru'llāh ast az
anfus-i ṭāghūtīyyih va hayākil-i firʿawnīyyih kih tajarrī nanamāyand va
ḥiṣn-i ʿiṣmat-i amr rā bi-khiyānat-i naẓar va qalb-i khud nayālāyand va
maẓāhir-i shayṭānīyyih bi-maqʿad-i quds-i sulaymānīyyih qadam
nagudhārand.

If in some utterances of the Manifestations of the Cause and the holy
Daysprings you find a restriction to one, do not let it veil you. Look
instead to the root of the Tree of the Cause and the holy Tree, for the
Cause of God has always appeared and will forever appear without limit
or number. Rather, these restrictions upon the reflecting mirrors are
intended to safeguard and conceal the Cause of God from rebellious souls
and pharaonic figures, so that they may not act presumptuously, or
defile the stronghold of the Cause's inviolability with the treachery of
their gaze and hearts, and so that satanic manifestations may not set
foot upon Solomon's seat of holiness.

\subsection{26}\label{section-2619}

\begin{otherlanguage}{persian}

چنانچه الیوم مشاهده میشود که هر نفسی را هوائی در سر و هر دیو نفس امّاره
در ظنون نفس خود بتوهّم افتاده و از هر گوشه ندائی مرتفع شده بسا جسد و
قلبهای مرده که بنعیق هوس خود قیام نموده‌اند و در ارض فرعونی قدم گذارده‌اند
ولکن نزد صاحبان افئده و بالغان مدینهٴ احدیّه حقّ از باطل ممتاز و مفصول
بوده و خواهد بود چنانچه هرگز مشتبه نشده و نخواهد شد

\end{otherlanguage}

Chonānche al-yowm moshāhede mī-shavad ke har nafsī rā havāʾī dar sar va
har dīv-e nafs-e ammāre dar ẓonūn-e nafs-e khod be-tavahhom oftāde va az
har gūshe nedāʾī mortafeʿ shode. Basā jasad va qalbhā-ye morde ke
be-naʿīq-e havas-e khod qiyām namūde-and va dar arḍ-e ferʿownī qadam
goẕārde-and, va-lāken nazd-e ṣāḥebān-e afʾede va bāleghān-e madīne-ye
aḥadīye ḥaqq az bāṭel momtāz va mafṣūl būde va khāhad būd, chonānche
hargez moshtabeh na-shode va na-khāhad shod.

As can be seen today, every soul has a desire in its head, and every
demon of the self that commands evil has fallen into delusion amid its
own imaginings, while a cry has risen from every corner. Many dead
bodies and hearts have risen up with the braying of their own desires
and set foot in the land of Pharaoh. Yet for those endowed with
understanding who have attained maturity in the city of oneness, truth
has always been and will remain distinct and separate from falsehood,
just as the two have never been, and will never be, confused.

\subsection{27}\label{section-2620}

\begin{otherlanguage}{persian}

باری حقّ بجوهر مقدّس خود در هیچ ظهوری از مظاهر و هیچ امری از اوامر محدود
بحدود قلمیّه که از مداد ظاهر است و محتجب بحجبات کلامیّه نشده الأمر بیده
یظهر کیف یشآء و انّه لفعّال لما یشآء و کلّ ما سواه لن یذکر عنده و لن یصل
الیه و انّه لن یعدّ بعدّ و لن یحدّ بحدّ و لن یشیر باشارة و لن یقترن بدلالة و
لن یستر بحجاب و لن یمنع بکلام و انّه لهو الفرد الصّمد الّذی لم‌یزل کان مریداً
بارادته و لا یزال یکون مقتدراً بمشیّته و کلّ بأمره یعملون

\end{otherlanguage}

Bārī, Ḥaqq be-jowhar-e moqaddas-e khod dar hīch ẓohūrī az maẓāher va
hīch amrī az avāmer maḥdūd be-ḥodūd-e qalamīye ke az medād ẓāher ast va
moḥtajeb be-ḥojobāt-e kalāmīye na-shode. Al-amru bi-yadihi yuẓhiru kayfa
yashāʾu wa-innahu la-faʿʿālun li-mā yashāʾu wa-kullu mā siwāhu lan
yudhkara ʿindahu wa-lan yaṣila ilayhi wa-innahu lan yuʿadda bi-ʿaddin
wa-lan yuḥadda bi-ḥaddin wa-lan yushīra bi-ishāratin wa-lan yaqtarina
bi-dalālatin wa-lan yustara bi-ḥijābin wa-lan yumnaʿa bi-kalāmin
wa-innahu la-huwa al-Fardu al-Ṣamadu alladhī lam yazal kāna murīdan
bi-irādatihi wa-lā yazālu yakūnu muqtadiran bi-mashīʾatihi wa-kullun
bi-amrihi yaʿmalūn.

In short, God, in His holy Essence, has not, in any manifestation among
the Manifestations or in any affair among His affairs, been confined
within the limits of writing that appears in ink, or veiled by the veils
of words. The Cause is in His hand; He brings forth whatever He wills,
and He indeed does whatever He wills. Nothing other than Him will ever
be mentioned in His presence or attain to Him. He will never be numbered
by any reckoning or circumscribed by any limit. No gesture will ever
point to Him, nor will He be joined to any indication, concealed by any
veil, or hindered by any words. He is indeed the Unique, the
All-Sufficing, who has always willed through His own will and will
forever exercise power through His own purpose. All act at His command.

\subsection{28}\label{section-2621}

\begin{otherlanguage}{persian}

و از عدم بلوغ اولو الکتاب بآن مقام است که در هر یک از ظهورات مظاهر الهیّه
از مبدء و مبدع امر محتجب مانده و در ظلمات وهمیّه مطروح بل معدوم و لاشیء
مانده‌اند سبحان‌اللّه چگونه سبقت میگیرد بر ارادهٴ حضرتش ارادهٴ دون او و حال
آنکه دون او بکلمهٴ او حیات یافته و بحشر وجود خود در امکنهٴ مقدّره محشور
شده چنانچه در لوح پارسی که چند سنه قبل مرقوم شد اکثر از امور که از انظر
عارفین مستور بوده ذکر یافته و اگر کسی بدیدهٴ قلب خود ملاحظه نماید جمیع
معارف قدسیّه و معانی کلمات بالغه را ادراک نماید بلکه مشهود مشاهده فرماید

\end{otherlanguage}

Va az ʿadam-e bolūgh-e olu'l-ketāb be-ān maqām ast ke dar har yek az
ẓohūrāt-e maẓāher-e elāhīye az mabdaʾ va mobdeʿ-e amr moḥtajeb mānde va
dar ẓolomāt-e vahmīye maṭrūḥ bal maʿdūm va lā-shayʾ mānde-and. Sobḥān
Allāh, chegūne sabqat mī-gīrad bar erāde-ye Ḥaẓratash erāde-ye dūn-e ū,
va ḥāl ānke dūn-e ū be-kaleme-ye ū ḥayāt yāfte va be-ḥashr-e vojūd-e
khod dar amkene-ye moqaddare maḥshūr shode. Chonānche dar Lowḥ-e Pārsī
ke chand sane qabl marqūm shod aksar az omūr ke az anẓar-e ʿārefīn
mastūr būde ẕekr yāfte, va agar kasī be-dīde-ye qalb-e khod molāḥeẓe
namāyad jamīʿ-e maʿāref-e qodsīye va maʿānī-ye kalemāt-e bāleghe rā
edrāk namāyad, balke mashhūd moshāhede farmāyad.

It is because the people of the Book have not attained that station
that, in every appearance of the divine Manifestations, they have
remained veiled from the Source and Originator of the Cause, cast into
the darkness of delusion---indeed, reduced to nonexistence and
nothingness. Glory be to God! How can the will of anyone other than Him
take precedence over His will, when all besides Him have received life
through His word and have been gathered, at the resurrection of their
being, in their appointed places? Thus, in the Persian Tablet written
some years ago, most of the matters that had been hidden from the eyes
of the mystic knowers were mentioned. If anyone considers it with the
eye of the heart, they will grasp all sacred knowledge and the meanings
of the consummate words---indeed, they will behold them directly.

\subsection{29}\label{section-2622}

\begin{otherlanguage}{persian}

باری از اوهن بیوت که ظنون انفس فانیهٴ مردوده است بمدینهٴ مقدّسهٴ طیّبهٴ
مستحکمه وارد شو و اقوال همج ارض را که بظلمت مفطور شده و مجعول شده بنغمات
ورقاء عزّ احدی مبادله مکن و از سلسبیل قدس سبحانی که از یمین عرش رحمانی
جاریست بجرعهٴ فانیهٴ کدره قانع مشو بدل شعلهٴ موسوی برافروز و بقلب نفحهٴ
روح القدس عیسوی بوز و چون خلیل عزّ تسلیم و تسکین در این نار مشتعلهٴ قدمیّه
و این شعلهٴ ملتهبه داخل شو دل و قلب و ارادهٴ خود را در دل و قلب و ارادهٴ
محبوب فانی کن لیجعل اللّه لک نوراً و رحمةً و انّه هو ولیّ المحسنین

\end{otherlanguage}

Bārī, az awhan-e boyūt ke ẓonūn-e anfos-e fānīye-ye mardūde ast
be-madīne-ye moqaddase-ye ṭayyebe-ye mostaḥkame vāred show, va aqvāl-e
hamaj-e arḍ rā ke be-ẓolmat mafṭūr shode va majʿūl shode be-naghamāt-e
Varqā-ye ʿEzz-e Aḥadī mobādele makon, va az salsabīl-e qods-e sobḥānī ke
az yamīn-e ʿarsh-e raḥmānī jārīst be-jorʿe-ye fānīye-ye kadere qāneʿ
mashow. Be-del shoʿle-ye mūsavī barafrūz va be-qalb nafḥe-ye
Rūḥo'l-Qodos-e ʿīsavī bevaz, va chon Khalīl-e ʿezz-e taslīm va taskīn
dar īn nār-e moshtaʿele-ye qedamīye va īn shoʿle-ye moltahebe dākhel
show. Del va qalb va erāde-ye khod rā dar del va qalb va erāde-ye Maḥbūb
fānī kon, li-yajʿala Allāhu laka nūran wa-raḥmatan wa-innahu huwa
waliyyu al-muḥsinīn.

In short, leave the frailest of houses---the imaginings of perishable,
rejected souls---and enter the holy, pure, fortified city. Do not
exchange the melodies of the Dove of the glory of oneness for the words
of the rabble of the earth, who have been formed and fashioned in
darkness. Do not forsake the heavenly Salsabīl of holiness flowing from
the right hand of the throne of the All-Merciful and content yourself
with a fleeting, murky draught. Kindle in your heart the flame of Moses,
and breathe into your inmost heart the breath of the Holy Spirit of
Jesus. Like Abraham, the Friend endowed with the glory of submission and
serenity, enter this blazing, ancient fire and this burning flame.
Efface your heart, your inmost being, and your will in the heart, the
inmost being, and the will of the Beloved, so that God may grant you
light and mercy. He is indeed the guardian of those who do good.

\subsection{30}\label{section-2623}

\begin{otherlanguage}{persian}

و بکمال توجّه مراقب بوده که از اسم اوّل که در دعای لیلة القدر نازل شده
محتجب نمانی و اعراض معرضین و انکار منکرین تو را از وصول و ورود در شاطی
بحر اعظم منع ننماید چه که این اسم مبارک که مقدّم اسماء است لم‌یزل باشراق
انوار شمس قدس مستشرق بوده و لایزال بضیاء بوارق حضرت غیبی مستضیء خواهد
بود دیدهٴ منیر باید تا ظهورات نار سدرهٴ سینا را از مهبط طور قدسیّهٴ طیّبه
ملاحظه نماید و نقطهٴ اعلی و جمال ابهی در ذکر این اسم و تقدّم آن بر جمیع
اسماء مفصّل و مبسوط ذکر فرموده فارجعوا الی البیان ان انتم تریدون ان
تعرفون

\end{otherlanguage}

Va be-kamāl-e tavajjoh morāqeb būde ke az esm-e avval ke dar doʿā-ye
Laylato'l-Qadr nāzel shode moḥtajeb na-mānī va eʿrāḍ-e moʿreẓīn va
enkār-e monkerīn to rā az voṣūl va vorūd dar shāṭī-ye baḥr-e aʿẓam manʿ
na-namāyad, che ke īn esm-e mobārak ke moqaddam-e asmā ast lam yazal
be-eshrāq-e anvār-e shams-e qods mostashreq būde va lā yazāl be-ẓiyā-ye
bavāreq-e Ḥaẓrat-e Ghaybī mostaẓīʾ khāhad būd. Dīde-ye monīr bāyad tā
ẓohūrāt-e nār-e Sedre-ye Sīnā rā az mahbeṭ-e Ṭūr-e qodsīye-ye ṭayyebe
molāḥeẓe namāyad, va Noqṭe-ye Aʿlā va Jamāl-e Abhā dar ẕekr-e īn esm va
taqaddom-e ān bar jamīʿ-e asmā mofaṣṣal va mabsūṭ ẕekr farmūde. Fa-rjiʿū
ilā al-Bayāni in antum turīdūna an taʿrifūn.

Be watchful with the utmost attention, lest you remain veiled from the
first name revealed in the prayer for the Night of Destiny, and lest the
rejection of those who turn away and the denial of the deniers prevent
you from reaching and entering the shore of the Most Great Ocean. For
this blessed name, which precedes all names, has always been illumined
by the radiance of the Sun of holiness and will forever be brightened by
the flashing splendors of the unseen Presence. An illumined eye is
needed to behold the manifestations of the fire of the Lote Tree of
Sinai from the place of revelation on the holy, pure Mount. The Most
Exalted Point and the All-Glorious Beauty have spoken at length and in
detail of this name and its precedence over all names. Turn, then, to
the Bayān, if you wish to know.

\subsection{31}\label{section-2624}

\begin{otherlanguage}{persian}

اگرچه آیات منزله و بیّنات ظاهره اظهر از آن است که استدلال بدون آن شود
ولکن اهل وهم و ظنّ که جز بهوی تکلّم ننمایند باینگونه دلایل ناظر بوده‌اند و
الیوم هم بمثل ما کان خواهند بود و اگر کسی الیوم دعوی نماید که عالم امریّه
بحدّی محدود شده و یا باشاره‌ئی مستشیر چنین نفس ابعد ناس است از صراط سویّ
عدل چه که عالم اطلاق و امر ازلاً قدماً مقدّس از تحدید و تقیید بوده اگر قلوب
صافیهٴ زکیّه اقلّ من آن مقابله بشمس احدیّه که الیوم در ظنون ناس محبوس و در
اوهام عباد مسجون مانده کنند جمیع اسرار ماکان و مایکون علی قدر تقابلهم و
استعدادهم در آن قلوب منطبع خواهد شد

\end{otherlanguage}

Agarche āyāt-e monzale va bayyenāt-e ẓāhere aẓhar az ān ast ke estedlāl
be-dūn-e ān shavad, va-lāken ahl-e vahm va ẓann ke joz be-havā takallom
na-namāyand be-īn-gūne dalāyel nāẓer būde-and va al-yowm ham be-methl-e
mā kān khāhand būd. Va agar kasī al-yowm daʿvā namāyad ke ʿālam-e amrīye
be-ḥaddī maḥdūd shode va yā be-eshāreʾī mostashīr, chonīn nafs abʿad-e
nās ast az ṣerāṭ-e savī-ye ʿadl, che ke ʿālam-e eṭlāq va amr azalan
qedaman moqaddas az taḥdīd va taqyīd būde. Agar qolūb-e ṣāfīye-ye zakīye
aqall-e men ān moqābele be-shams-e aḥadīye ke al-yowm dar ẓonūn-e nās
maḥbūs va dar owhām-e ʿebād masjūn mānde konand, jamīʿ-e asrār-e mā kān
va mā yakūn ʿalā qadri taqābulihim wa-istiʿdādihim dar ān qolūb monṭabeʿ
khāhad shod.

Although the revealed verses and manifest proofs are too evident for
anything else to be used as evidence, the people of delusion and
conjecture, who speak only from desire, have looked to arguments of this
kind, and today they will remain as they were before. If anyone today
claims that the realm of the divine Cause has been confined within a
limit or indicated by an allusion, such a person is the furthest of all
people from the straight path of justice. For the realm of the Absolute
and of the Cause has, from all eternity, been sanctified above
restriction and limitation. If pure, stainless hearts were to turn, for
less than an instant, toward the Sun of oneness, which today remains
imprisoned in people's conjectures and confined within the imaginings of
His servants, all the mysteries of what has been and what is to come
would be imprinted upon those hearts, according to the degree of their
turning and their capacity.

\subsection{32}\label{section-2625}

\begin{otherlanguage}{persian}

و حال علاوه بر آنکه چنین فضل اکبر را مقبل نشده غلّ و بغضا از اکثر نفوس
ظاهر و مع‌ذلک خود را از مؤمنین بیان دانسته و مسرورند و غافل از اینکه جمیع
بیان بجمیع لسان بر معرضین این مقام لعن مینماید باری حال جمیع محتجبین در
شرک باللّه و نار نفی و بعد و سجّین عذاب و هاویهٴ قهر مقرّ گرفته‌اند بقسمیکه
احدی از این غافلان گویا هرگز شاعر نبوده و نیستند و بقسمی میّت لاشیء گشته
که از نسیم عنایت که از مکمن جود ربّ العزّه در کلّ حین بنفحات لایح مبین
میوزد بحرکت نیامده‌اند چه جای آنکه خیمه بر شطر قدم زنند و یا خرگاه بر
شاطی بحر اعظم مرتفع نمایند

\end{otherlanguage}

Va ḥāl ʿalāve bar ānke chonīn faẓl-e akbar rā moqbel na-shode, ghell va
baghẓā az akthar-e nofūs ẓāher va maʿa dhālek khod rā az moʾmenīn-e
Bayān dāneste va masrūrand, va ghāfel az īnke jamīʿ-e Bayān be-jamīʿ-e
lesān bar moʿreẓīn-e īn maqām laʿn mī-namāyad. Bārī, ḥāl jamīʿ-e
moḥtajebīn dar sherk be-Allāh va nār-e nafy va boʿd va sejjīn-e ʿadhāb
va hāviye-ye qahr maqarr gerefte-and, be-qesmī ke aḥadī az īn ghāfelān
gūyā hargez shāʿer na-būde va nīstand, va be-qesmī mayyet-e lā-shayʾ
gashte ke az nasīm-e ʿenāyat ke az makman-e jūd-e Rabbo'l-ʿEzze dar
koll-e ḥīn be-nafaḥāt-e lāyeḥ-e mobīn mī-vazad be-ḥarakat na-yāmade-and,
che jāy-e ānke kheyme bar shaṭr-e qedam zanand va yā khargāh bar
shāṭī-ye baḥr-e aʿẓam mortafeʿ namāyand.

Yet now, besides failing to turn toward so great a bounty, most souls
display malice and hatred. Nevertheless, they count themselves among the
believers in the Bayān and rejoice, unaware that the whole Bayān, in
every tongue, curses those who turn away from this station. In short,
all who are veiled have now taken up their abode in associating partners
with God, in the fire of denial and remoteness, in the dungeon of
torment, and in the abyss of wrath. It is as though not one of these
heedless people has ever possessed, or now possesses, any awareness.
They have become so utterly dead and reduced to nothing that they have
not been stirred by the breeze of loving-kindness, which at every moment
blows in manifest, evident breaths from the hidden sanctuary of the
bounty of the Lord of Glory---much less have they pitched their tent in
the realm of eternity or raised their pavilion on the shore of the Most
Great Ocean.

\subsection{33}\label{section-2626}

\begin{otherlanguage}{persian}

و جمیع این بیانات نظر بتوقّع و طلب عباد در الواح ذکر میشود والّا شمس
مشعشعهٴ مشرقهٴ لایحه مقدّس از آن است که بوصفی موصوف و یا بنعتی منعوت و یا
بذکری مذکور گردد جمیع اذکار و اوصاف باشراقی از او مستشرق و ظاهرند بشنو
ندای روح القدس منطق در این سدرهٴ منیرهٴ مرتفعه را که بنغمه‌ئی از او اجساد
فانیهٴ معدومه بجبروت باقیهٴ دائمه درآیند تا جمیع ارض و خلق آن را از
اعالی و ادانی در ظلّ ظلّ مشاهده نمائی و از مؤتفکات ظنون و وهم بمحکمات این
کلمات درّیّات و بیّنات واضحات بفردوس حقیقت که مقام استواء عرش سبحانی است
درآئی و بثمرهٴ جنیّهٴ جنّت احدیّه فایز گردی

\end{otherlanguage}

Va jamīʿ-e īn bayānāt naẓar be-tavaqqoʿ va ṭalab-e ʿebād dar alvāḥ ẕekr
mī-shavad, va-ellā shams-e moshaʿshaʿe-ye moshreqe-ye lāyeḥe moqaddas az
ān ast ke be-vaṣfī mowṣūf va yā be-naʿtī manʿūt va yā be-ẕekrī maẕkūr
gardad. Jamīʿ-e aẕkār va owṣāf be-eshrāqī az ū mostashreq va ẓāherand.
Beshnow nedā-ye Rūḥo'l-Qodos-e monṭeq dar īn Sedre-ye monīre-ye
mortafeʿe rā ke be-naghmeʾī az ū ajsād-e fānīye-ye maʿdūme be-jabarūt-e
bāqīye-ye dāʾeme darāyand, tā jamīʿ-e arḍ va khalq-e ān rā az aʿālī va
adānī dar ẓell-e ẓell moshāhede namāʾī va az moʾtafekāt-e ẓonūn va vahm
be-moḥkamāt-e īn kalemāt-e dorrīyāt va bayyenāt-e vāẓeḥāt be-ferdows-e
ḥaqīqat ke maqām-e estevā-ye ʿarsh-e sobḥānī ast darāʾī va be-thamare-ye
janīye-ye jannat-e aḥadīye fāyez gardī.

All these utterances are recorded in the Tablets in consideration of the
expectations and requests of His servants. Otherwise, the radiant,
dawning, resplendent Sun is sanctified above being described by any
description, characterized by any attribute, or mentioned in any
remembrance. All remembrances and descriptions are illumined and made
manifest by a single ray from Him. Hear the call of the Holy Spirit
speaking in this luminous, lofty Tree, through a single melody of which
perishable, nonexistent bodies enter the everlasting, enduring realm of
might, so that you may behold the whole earth and its people, high and
low, in the shadow of a shadow. Thus may you leave the overturned cities
of conjecture and delusion and, through the firm foundations of these
pearl-like words and clear proofs, enter the paradise of reality, where
the heavenly throne is established, and partake of the ripe fruit of the
garden of oneness.

\subsection{34}\label{section-2627}

\begin{otherlanguage}{persian}

آیا در ربیع الهی که ایّام ظهور شجرهٴ قصوی و اشراق طلوع طلعت اعلی بود
مشاهده نشد که چه‌قدر از نفوس قدسیّه و هیاکل عزّیّه که از آنشمس قدمیّه مشرق و
مستشرق شده و هر کدام مدّعی مقامات عالیه و مراتب ما لا نهایه شده چنانچه
جمیع این خلق موجود بعین رأس خود دیده‌اند مثل جناب قدّوس که اعلی و ابهی
ثمرهٴ سدره بودند بسلطنتی ظاهر و مشهود گشتند که گویا شجرهٴ نار از یک تجلّی
او ندا نموده و کذلک طلعت طا و من دونهما الّذی یکون موجوداً بینکم

\end{otherlanguage}

Āyā dar rabīʿ-e elāhī ke ayyām-e ẓohūr-e Shajare-ye Qoṣvā va eshrāq-e
ṭolūʿ-e Ṭalʿat-e Aʿlā būd moshāhede na-shod ke che-qadr az nofūs-e
qodsīye va hayākel-e ʿezzīye ke az ān shams-e qedamīye moshreq va
mostashreq shode va har kodām moddaʿī-ye maqāmāt-e ʿālīye va marāteb-e
mā lā nehāye shode, chonānche jamīʿ-e īn khalq-e mowjūd be-ʿayn-e raʾs-e
khod dīde-and. Methl-e Janāb-e Qoddūs ke aʿlā va abhā thamare-ye Sedre
būdand be-salṭanatī ẓāher va mashhūd gashtand ke gūyā shajare-ye nār az
yek tajallī-ye ū nedā namūde, va ka-dhālika Ṭalʿatu Ṭā wa-man dūnahumā
alladhī yakūnu mawjūdan baynakum.

In the divine springtime, the days when the furthermost Tree appeared
and the Most Exalted Countenance rose in splendor, was it not witnessed
how many holy souls and glorious figures shone forth from that ancient
Sun and were illumined by it, each claiming exalted stations and
limitless ranks, as all these people now living have seen with their own
eyes? His honor Quddūs, for example, who was the highest and most
glorious fruit of the Tree, appeared and was witnessed with such
sovereignty that it was as though the tree of fire had called out
through a single revelation of his splendor. So too did the Countenance
of Ṭā, and the one below those two who is present among you.

\subsection{35}\label{section-2628}

\begin{otherlanguage}{persian}

و مع‌ذلک که جمیع این مراتب را جمیع این بریّه بچشم خود و گوش خود شنیده‌اند و
ادراک نموده‌اند مع‌ذلک تمسّک ببعضی از اقوال و کلمات جسته و از منزل و مرسل
رحمت واسعه محروم مانده‌اند با اینکه ایّامی از آنچه بچشم خود دیده‌اند نگذشته
مع‌ذلک مشاهد بصر خود را چنان تشکیک و تردید نموده‌اند که مجدّد سؤال مینمایند
و بقسمی در ظنون مستغرقند که اگر نفسی بهوای خود دو دلیل اوهن از بیت
عنکبوت اقامه نماید و جمیع فیوضات نامتناهی الهی را منتهی و محدود نماید
جمیع بالطّوع من دون کره اخذ نمایند و مسموع و مشهود خود را کأن لم یکن فرض
نمایند صدهزار حیرت و حسرت بر چنین نفوس ضعیفهٴ مضطربهٴ غیر مستقیمه که
باندک ریحی از یمین علم و استقامت بشمال ظنّ و زلّت میل نمایند

\end{otherlanguage}

Va maʿa dhālek ke jamīʿ-e īn marāteb rā jamīʿ-e īn barīye be-chashm-e
khod va gūsh-e khod shenīde-and va edrāk namūde-and, maʿa dhālek
tamassok be-baʿẓī az aqvāl va kalemāt joste va az monzel va morsel-e
raḥmat-e vāseʿe maḥrūm mānde-and. Bā īnke ayyāmī az ānche be-chashm-e
khod dīde-and na-goẕashte, maʿa dhālek mashāhed-e baṣar-e khod rā chonān
tashkīk va tardīd namūde-and ke mojaddad soʾāl mī-namāyand, va be-qesmī
dar ẓonūn mostaghraqand ke agar nafsī be-havā-ye khod do dalīl-e awhan
az bayt-e ʿankabūt eqāme namāyad va jamīʿ-e foyūẓāt-e nā-motanāhī-ye
elāhī rā montahī va maḥdūd namāyad, jamīʿ be'l-ṭowʿ men dūn-e korh akhẕ
namāyand va masmūʿ va mashhūd-e khod rā ka-an lam yakon farẓ namāyand.
Ṣad-hezār ḥayrat va ḥasrat bar chonīn nofūs-e ẓaʿīfe-ye moẓṭarebe-ye
ghayr-e mostaqīme ke be-andak rīḥī az yamīn-e ʿelm va esteqāmat
be-shemāl-e ẓann va zellat meyl namāyand.

Yet although all these people have witnessed, heard, and understood all
these degrees with their own eyes and ears, they have nevertheless clung
to certain sayings and words and remained deprived of the One who sends
down and bestows all-embracing mercy. Although only a few days have
passed since what they saw with their own eyes, they have cast such
doubt and uncertainty upon their own observations that they ask about
them again. They are so immersed in conjecture that, if someone were to
advance, out of personal desire, two arguments frailer than a spider's
house and declare all the infinite outpourings of God to be finite and
limited, they would all accept them willingly, without compulsion, and
regard what they themselves had heard and witnessed as though it had
never been. A hundred thousand expressions of astonishment and regret
for such weak, wavering, inconstant souls, who at the slightest breeze
turn from the right hand of knowledge and steadfastness to the left hand
of conjecture and error!

\subsection{36}\label{section-2629}

\begin{otherlanguage}{arabic}

بمثل امم ماضیه که جمیع فیوضات متتابعه و ظهورات غیر منتهیهٴ الهیّه را
بظنون وهمیّهٴ خود منتهی نموده‌اند و عبدهٴ هوای خود من دون اللّه بوده و
خواهند بود و حال آنکه در جمیع کتب و الواح و زبر و اوراق و صحف و رقاع
تصریحاً من دون تلویح ذکر عدم تعطیل فیض فیّاض و عدم تأخیر و قطع فضل فضّال و
عدم تناهی ظهورات ما لا نهایه مذکور و مبسوطست و مع‌ذلک حال جمیع من علی
الأرض را مشاهده میشود که بتحدید و انقطاع فیض و فضل مقرّ و معترف بلکه راضی
و مسرورند چنانچه اهل فرقان باسمائی چند که مردّف و منظّم ضبط نموده‌اند مشغول
و مسرور گشته و از جوهر اسماء و موجد و مظهر او که ابحر اسمیّه و صفتیّه
بجذبهٴ اسم مبارکش موّاج و متلاطم گشته غافل و محروم مانده‌اند و اللّه یهدی
من یشآء الی صراطه و انّه هو هادی السّبیل و انّه لهو المبدء و المعید و الیه
یصعد ندآء العبید و انّه لهو العزیز الجمیل و

\end{otherlanguage}

Be-methl-e omam-e māẓīye ke jamīʿ-e foyūẓāt-e motatābeʿe va ẓohūrāt-e
ghayr-e montahīye-ye elāhī rā be-ẓonūn-e vahmīye-ye khod montahī
namūde-and va ʿabade-ye havā-ye khod men dūn-e Allāh būde va khāhand
būd. Va ḥāl ānke dar jamīʿ-e kotob va alvāḥ va zobor va ovrāq va ṣoḥof
va reqāʿ taṣrīḥan men dūn-e talvīḥ ẕekr-e ʿadam-e taʿṭīl-e fayẓ-e Fayyāẓ
va ʿadam-e taʾkhīr va qaṭʿ-e faẓl-e Faẓẓāl va ʿadam-e tanāhī-ye
ẓohūrāt-e mā lā nehāye maẕkūr va mabsūṭ ast. Va maʿa dhālek ḥāl-e
jamīʿ-e man ʿalā al-arḍ rā moshāhede mī-shavad ke be-taḥdīd va enqeṭāʿ-e
fayẓ va faẓl moqerr va moʿtaref balke rāẓī va masrūrand. Chonānche ahl-e
Forqān be-asmāʾī chand ke moraddaf va monaẓẓam ẓabṭ namūde-and mashghūl
va masrūr gashte va az jowhar-e asmā va mūjed va moẓher-e ū ke abḥor-e
esmīye va ṣefatīye be-jaẕbe-ye esm-e mobārakash mavvāj va motalāṭem
gashte ghāfel va maḥrūm mānde-and. Wa-Allāhu yahdī man yashāʾu ilā
ṣirāṭihi wa-innahu huwa Hādī al-Sabīli wa-innahu la-huwa al-Mubdiʾu
wa-al-Muʿīdu wa-ilayhi yaṣʿadu nidāʾu al-ʿabīdi wa-innahu la-huwa
al-ʿAzīzu al-Jamīlu wa-

They are like the peoples of the past, who, through their deluded
conjectures, have brought all the successive outpourings and unending
manifestations of God to an end, and have been, and will remain,
worshippers of their own desires rather than God. Yet in all the books,
tablets, scriptures, pages, scrolls, and epistles, it is explicitly
stated and set forth at length, not merely hinted, that the grace of the
All-Bountiful can never be suspended, that the favor of the Bestower can
never be delayed or cut off, and that the infinite manifestations can
never end. Nevertheless, the condition of all who dwell on earth can now
be seen: they acknowledge and confess that grace and favor are limited
and have ceased---indeed, they are content and rejoice in this. Thus the
people of the Qurʾān have become engrossed in, and delighted with, a few
names that they have recorded in ordered sequence, while remaining
heedless of and deprived of the Essence of the names and its Originator
and Manifestation, through the attraction of whose blessed name the
oceans of names and attributes have surged and tossed. God guides
whomever He wills to His path. He is indeed the Guide of the way. He is
indeed the Originator and the One who brings back, and to Him ascends
the cry of His servants. He is indeed the Almighty, the Beautiful, and

\subsection{37}\label{section-2630}

\begin{otherlanguage}{persian}

از همهٴ این بیانات منیرهٴ مستضیئهٴ متقنهٴ محکمه گذشته خود آن طلعت احدیّه
و شمس اوّلیّه در جمیع بیانات و آیات و خطب که از آن سماء مشیّت و مصدر و مطلع
علم و حکمت نازل و ظاهر شده از سلطان غیب و حقیقت لاریب که مقام خفیّهٴ
باطنیّهٴ غیبیّه است طلب فرموده‌اند که در هر سنه مرآتی ظاهر فرمایند بلکه در
هر شهری بلکه در هر ساعتی مع‌ذلک چگونه میشود که هر نفسی بخیال خود هوسی
نماید و امر اللّه را محدود نماید فتعالی عمّا هم یعرفون و یقولون سراب فانی
را کأن لم یکن ملاحظه نما و بکوثر باقی سبحانی بشتاب فرّ بنفسک و روحک و
قلبک و رجلک الی اللّه الّذی الیه منقلبک و مثواک و لا تلتفت الی مثوی
الظّالمین و المشرکین

\end{otherlanguage}

Az hame-ye īn bayānāt-e monīre-ye mostaḍīʾe-ye motqane-ye moḥkame
godhashte, khod-e ān ṭalʿat-e aḥadīye va shams-e avvalīye dar jamīʿ-e
bayānāt va āyāt va khoṭab ke az ān samā-ye mashīyat va maṣdar va
maṭlaʿ-e ʿelm va ḥekmat nāzel va ẓāher shode, az solṭān-e ghayb va
ḥaqīqat-e lā-rayb ke maqām-e khafīye-ye bāṭenīye-ye ghaybīye ast ṭalab
farmūde-and ke dar har sane merʾātī ẓāher farmāyand, balke dar har
shahrī, balke dar har sāʿatī. Maʿa-dhālek chegūne mīshavad ke har nafsī
be-khiyāl-e khod havasī namāyad va amr-e Allāh rā maḥdūd namāyad?
Fa-taʿālā ʿammā hum yaʿrifūna va yaqūlūn. Sarāb-e fānī rā ka-an lam
yakun molāḥeẓe namā va be-kawthar-e bāqī-ye sobḥānī beshetāb. F arra
bi-nafsika va rūḥika va qalbika va rijlika ilā Allāh alladhī ilayhi
munqalabuka va mathwāka va lā taltafit ilā mathwā al-ẓālimīna va
al-mushrikīn.

Beyond all these luminous, radiant, sound, and firmly established
utterances, that Countenance of Oneness and primal Sun has Himself, in
all the utterances, verses, and discourses sent down and manifested from
that heaven of the Will and source and dayspring of knowledge and
wisdom, asked the Sovereign of the Unseen and the indubitable Reality,
whose station is hidden, inward, and unseen, to manifest a Mirror every
year---indeed, every month, indeed, every hour. How, then, can anyone
indulge a desire of his own imagining and set limits to the Cause of
God? Exalted is He above what they understand and say. Regard the
fleeting mirage as though it had never been, and hasten to the
everlasting, divine Kawthar. Flee with your self, your spirit, your
heart, and your feet to God, to whom you will return and with whom you
will dwell, and do not turn toward the abode of the wrongdoers and those
who associate partners with Him.

\subsection{38}\label{section-2631}

\begin{otherlanguage}{persian}

ثمّ اشکره فی کلّ حینک بما انزل علیک الآیات بالحقّ لعلّ النّاس کانوا بآیات
اللّه ان یهتدون و اریاح قدس لقا را که از نفحهٴ مشکیّهٴ این کلمات
{[}منیعهٴ{]} معلنهٴ محکمه در هبوب و مرور است التفات نموده و از جمیع
بیانات شبحیّهٴ ظلّیّه و کلمهٴ غیر مسموعهٴ مردوده تغافل نما و این اریاح
لم‌یزل در هبوب بوده و خواهد بود ولکن احدی شاعر نشده و حیات کلّ من فی
السّموات و الأرض از او ولکن کلّ محجوب از آن جوهر فطرت و ساذج طینت و اگر
این عنایت مشرقهٴ منیرهٴ لایحه از مرایای وجود موجودات منقطع شود کلّ را
لاشیء محض خواهی دید بل معدوم صرف فتعالی از این هبوب روحانی که باقبال و
شعور و التفات و یا دون آن محدود نشده لم‌یزل بر اجساد میّته و عظام رمیمه و
هیاکل باقیه وزیده و خواهد وزید خواه احدی ملتفت بوده و یا غافل و محجوب
مانده

\end{otherlanguage}

Thumma ushkurhu fī kulli ḥīnika bi-mā anzala ʿalayka al-āyāti
bi-al-ḥaqqi laʿalla al-nāsa kānū bi-āyāti Allāhi an yahtadūn. Va aryāḥ-e
qods-e leqā rā ke az nafḥe-ye meshkīye-ye īn kalemāt-e {[}manīʿe-ye{]}
moʿlane-ye moḥkame dar hobūb va morūr ast eltefāt namūde va az jamīʿ-e
bayānāt-e shabaḥīye-ye ẓellīye va kalame-ye gheyr-e masmūʿe-ye mardūde
taghāfol namā. Va īn aryāḥ lam-yazal dar hobūb būde va khwāhad būd,
valāken aḥadī shāʿer nashode, va ḥayāt-e kulli man fī al-samāwāti va
al-arḍ az ū, valāken koll maḥjūb az ān jawhar-e feṭrat va sādhaj-e
ṭīnat. Va agar īn ʿenāyat-e moshreqe-ye monīre-ye lāʾeḥe az marāyā-ye
vojūd-e mawjūdāt monqaṭeʿ shavad, koll rā lā-shayʾ-e maḥḍ khwāhī dīd,
bal maʿdūm-e ṣerf. Fa-taʿālā az īn hobūb-e rūḥānī ke be-eqbāl va shoʿūr
va eltefāt va yā dūn-e ān maḥdūd nashode; lam-yazal bar ajsād-e mayyete
va ʿeẓām-e ramīme va hayākel-e bāqīye vazīde va khwāhad vazīd, khwāh
aḥadī moltafet būde va yā ghāfel va maḥjūb mānde.

Then thank Him at every moment for having sent down the verses to you in
truth, so that people might be guided by the verses of God. Attend to
the holy winds of reunion that blow and pass through the musk-laden
breath of these {[}exalted{]}, proclaimed, firmly established words, and
disregard all shadowy, phantom-like utterances and every unheard,
rejected word. These winds have always blown and will continue to blow,
yet no one has perceived them. The life of all who are in the heavens
and on earth comes from them, yet all are veiled from that essence of
creation and purest substance of nature. If this resplendent, luminous,
manifest favor were cut off from the mirrors of the existence of created
things, you would see them all as sheer nothingness---indeed, as utter
nonexistence. Exalted is this spiritual breeze, which is not limited by
anyone's turning toward it, awareness, or attention, or by their
absence. It has always blown and will continue to blow over dead bodies,
decayed bones, and surviving forms, whether anyone has attended to it or
remained heedless and veiled.

\subsection{39}\label{section-2632}

\begin{otherlanguage}{persian}

مثلاً نسیم ربیع بر جمیع اشیاء میوزد و بر کلّ من علی الأرض مرور مینماید و
انفسیکه در نوم غفلت مدهوشند و یا از رحمت الهی بهوش بر هر دو مرور مینماید
ولکن هنیئاً للشّاعرین چه که غافلین لم‌یزل عنداللّه مردود و غیر مقبول بوده
چنانچه مذکور شد و الیوم جمیع من دان فی البیان و دون آن از ملل مختلفه
مأمورند بدخول در ظلّ این سدرهٴ مرتفعه و شجرهٴ منطقه و کلمهٴ جامعه و حجّت
محکمه و دلیل واضحه و سبیل مشهوده و مفرّ و مقرّ و مستقرّ از برای نفسی نبوده
و نیست و جمیع بنفس خود مکلّفند باین امر بدیع منیع و دیگر آنکه اقرار بر
این کلمه منوط و مشروط بتصدیق نفسی نبوده و نخواهد بود چنانچه بر هر ذی
بصری مشهود و واضح است

\end{otherlanguage}

Mathalan nasīm-e rabīʿ bar jamīʿ-e ashyā mīvazad va bar kulli man ʿalā
al-arḍ morūr mīnamāyad, va anfosī ke dar nawm-e ghaflat madhūshand va yā
az raḥmat-e elāhī be-hūsh, bar har do morūr mīnamāyad. Valāken hanīʾan
li-al-shāʿirīn, che ke ghāfelīn lam-yazal ʿinda Allāh mardūd va gheyr-e
maqbūl būde, chenānche madhkūr shod. Va al-yawm jamīʿ-e man dāna fī
al-Bayān va dūn-e ān az melal-e mokhtalefe maʾmūrand be-dokhūl dar
ẓell-e īn sedre-ye mortafeʿe va shajare-ye monṭeqe va kalame-ye jāmeʿe
va ḥojjat-e moḥkame va dalīl-e vāḍeḥe va sabīl-e mashhūde, va mafarr va
maqarr va mostaqarr az barā-ye nafsī nabūde va nīst. Va jamīʿ be-nafs-e
khod mokallafand be-īn amr-e badīʿ-e manīʿ. Va dīgar ānke eqrār bar īn
kalame manūṭ va mashrūṭ be-taṣdīq-e nafsī nabūde va nakhwāhad būd,
chenānche bar har dhī-baṣarī mashhūd va vāḍeḥ ast.

For example, the spring breeze blows over all things and passes over
everyone on earth. It passes alike over those who lie insensible in the
sleep of heedlessness and those whom divine mercy has awakened. Yet
blessed are those who perceive, for the heedless have always been
rejected and unaccepted in God's sight, as has been mentioned. Today,
all who profess the Bayān, as well as those of other, diverse religions,
are commanded to enter the shade of this exalted Lote Tree, this
speaking Tree, this all-embracing Word, this firmly established proof,
this clear evidence, and this manifest path. There has never been, nor
is there now, any refuge, retreat, or resting place for anyone
elsewhere. All are personally charged with this wondrous, exalted Cause.
Furthermore, acknowledgment of this Word has never been, and will never
be, contingent or conditional upon anyone else's affirmation, as is
evident and clear to everyone with sight.

\subsection{40}\label{section-2633}

\begin{otherlanguage}{persian}

چه میشود که مسئول عنه خود در غفلت باشد چنانچه امم فرقانیّه بهمین جهت ناظر
بوده از محلّ وحی الهی و معدن و منبع فیض مطلق ربّانی محروم مانده‌اند و جمیع
چشم و گوش را بقول و فعل رؤسای خود معلّق و منوط داشته و نقطهٴ وجود و جمال
مقصود روح من فی الملک فداه در مواضع عدیده از بیان در اینمقام وعر خشن
ذکوان وصایای منیعه و تأکیدات بلیغه فرموده‌اند

\end{otherlanguage}

Che mīshavad ke masʾūl ʿanhu khod dar ghaflat bāshad, chenānche omam-e
forqānīye be-hamīn jahat nāẓer būde az maḥall-e vaḥy-e elāhī va maʿdan
va manbaʿ-e fayḍ-e moṭlaq-e rabbānī maḥrūm mānde-and, va jamīʿ chashm va
gūsh rā be-qawl va feʿl-e roʾasā-ye khod moʿallaq va manūṭ dāshte. Va
noqṭe-ye vojūd va jamāl-e maqṣūd, rūḥu man fī al-mulk fidāhu, dar
mavāḍeʿ-e ʿadīde az Bayān dar īn maqām-e vaʿr-e khashen-e dhakavān
vaṣāyā-ye manīʿe va taʾkīdāt-e balīghe farmūde-and.

For the person consulted may himself be heedless, just as the peoples of
the Qurʾān, looking in this very direction, have remained deprived of
the seat of divine revelation and the mine and wellspring of God's
unrestricted grace. They have all made their eyes and ears dependent
upon the words and deeds of their leaders. The Point of existence and
the Beauty of the Desired One---may the souls of all in the realm be
sacrificed for Him---has, in numerous passages of the Bayān, given
exalted counsels and emphatic injunctions concerning this rugged, harsh,
and blazing pass.

\subsection{41}\label{section-2634}

\begin{otherlanguage}{persian}

چنانچه در یک مقام میفرماید که از اهل بیان که ادراک یوم ظهور مشرقه شموس
عزّ احدیّه نمایند باید خود بنفس خود بآن سدرهٴ ربّانیّه تقرّب جوید و مقام
عرفان او که حقیقت جنّت بدیعه و رضوان احدیّه است بعقل و فؤاد و سمع خود حاصل
نماید نه آنکه خود را معلّق بردّ و قبول رؤسای عهد گذارند و از اینمقام ترقّی
داده‌اند و مخصوص ذکر فرموده‌اند که اگر از حروفات من در آن یوم موجود باشند
و انکار نمایند بمنع ایشان ممنوع نشده خود را از فیض سماء عزّ محروم ننمایند
چه میشود که ایشان نظر باسباب فانیه و یا ریاست ظاهره از قرب سدره محروم
مانند

\end{otherlanguage}

Chenānche dar yek maqām mīfarmāyad ke az ahl-e Bayān ke edrāk-e yawm-e
ẓohūr-e moshreqe-ye shomūs-e ʿezz-e aḥadīye namāyand bāyad khod
be-nafs-e khod be-ān sedre-ye rabbānīye taqarrub jūyad va maqām-e
ʿerfān-e ū ke ḥaqīqat-e jannat-e badīʿe va reḍvān-e aḥadīye ast be-ʿaql
va foʾād va samʿ-e khod ḥāṣel namāyad, na ānke khod rā moʿallaq be-radd
va qabūl-e roʾasā-ye ʿahd godhārand. Va az īn maqām taraqqī dāde-and va
makhṣūṣ dhekr farmūde-and ke agar az ḥorūfāt-e man dar ān yawm mawjūd
bāshand va enkār namāyand, be-manʿ-e īshān mamnūʿ nashode khod rā az
fayḍ-e samā-ye ʿezz maḥrūm nanamāyand, che mīshavad ke īshān naẓar
be-asbāb-e fānīye va yā reyāsat-e ẓāhere az qorb-e sedre maḥrūm mānand.

Thus, in one passage He states that whoever among the people of the
Bayān lives to see the day of the radiant manifestation of the Suns of
the glory of Oneness must personally draw near to that divine Lote Tree
and attain recognition of Him---which is the very reality of the
wondrous Paradise and the Riḍvān of Oneness---through his own reason,
heart, and hearing, rather than making himself dependent on the
rejection or acceptance of the leaders of the age. He goes further and
expressly states that if any of My Letters are present on that day and
deny Him, people should not be deterred by their opposition or deprive
themselves of the grace of the heaven of glory, for those Letters may
remain deprived of nearness to the Lote Tree through regard for fleeting
advantages or outward leadership.

\subsection{42}\label{section-2635}

\begin{otherlanguage}{persian}

و در مقام دیگر ذکر میفرمایند که در آن یوم بشیئی محتجب نشده که شاید ادلّای
نور من هم محتجب میمانند فاعتصموا یا قوم بحبل اللّه و عنایته و لا تکوننّ من
السّائلین و لا تتوجّهوا فی عرفانه بغیره لأنّ دلیله آیاته و وجوده اثباته ان
انتم من الموقنین و ایّاکم یا ملأ البیان انتم لا تحتجبوا بکتب ولو کان ملأ
السّموات و الأرض من الکتب المحکمة و لا بالآیات المنزلة لأنّ دونه لن یسبق
علیه فی شیء و کذلک فاعرفه ان انتم من العارفین لأنّ کلّ ما فی الکتب و الصّحف
و الزّبر من الکلمات العالیة و الحکم البالغة کلّها یثبت بقوله ان انتم من
العارفین

\end{otherlanguage}

Va dar maqām-e dīgar dhekr mīfarmāyand ke dar ān yawm be-shayʾī moḥtajeb
nashode, ke shāyad adellā-ye nūr-e man ham moḥtajeb mīmānand.
Fa-iʿtaṣimū yā qawmu bi-ḥabli Allāhi va ʿināyatihi va lā takūnanna mina
al-sāʾilīn. Va lā tatawajjahū fī ʿirfānihi bi-ghayrihi li-anna dalīlahu
āyātuhu va vujūduhu ithbātuhu in antum mina al-mūqinīn. Va iyyākum yā
malaʾa al-Bayān antum lā taḥtajibū bi-kutubin va law kāna milʾa
al-samāwāti va al-arḍi mina al-kutubi al-muḥkamati va lā bi-al-āyāti
al-munzala, li-anna dūnahu lan yasbiqa ʿalayhi fī shayʾin va kadhālika
fa-iʿrifhu in antum mina al-ʿārifīn. Li-anna kulla mā fī al-kutubi va
al-ṣuḥufi va al-zuburi mina al-kalimāti al-ʿāliyati va al-ḥikami
al-bālighati kulluhā yathbutu bi-qawlihi in antum mina al-ʿārifīn.

In another passage He states that on that day they should not be veiled
by anything, for even the guides to My light may remain veiled. Hold
fast, then, O people, to the cord of God and His loving care, and do not
be among those who question. Do not turn to anyone other than Him to
recognize Him, for His proof is His verses, and His existence is His
evidence, if you are certain. Beware, O people of the Bayān, lest you be
veiled by books, even if the heavens and the earth were filled with
firmly established books, or by revealed verses, for nothing other than
Him can take precedence over Him in anything. Recognize Him in this way,
if you are among those who recognize. For all the exalted words and
consummate wisdom contained in the books, scriptures, and scrolls are
established by His utterance, if you are among those who recognize.

\subsection{43}\label{section-2636}

\begin{otherlanguage}{persian}

باری کمال حفظ را مرعی داشته و مراقب بوده که بحر اعظم ملتطمهٴ متموّجه را
بقطره میزان ننمائی چه که دون آن نزد او مذکور نه فسبحان اللّه عمّا انتم
تقولون مقصود از ذکر این کلمات آنکه مبادا تو را غفلت احاطه و اخذ نماید و
در حین اشراق شموس مشرقه تمسّک بردّ و قبول احدی نمائی باری تمسّک باحدی نجسته
و خود بیکنونت خود باشراق آن شمس جود و فضل اقبال نما که همان اشراق اظهر
از کلّ دلائل متقنه و اعظم از جمیع حجج محکمه بوده و خواهد بود و این مضمون
در کتب منزله مذکور است

\end{otherlanguage}

Bārī kamāl-e ḥefẓ rā marʿī dāshte va morāqeb būde ke baḥr-e aʿẓam-e
moltaṭeme-ye motamavveje rā be-qaṭre mīzān nanamāʾī, che ke dūn-e ān
nazd-e ū madhkūr na. Fa-subḥāna Allāhi ʿammā antum taqūlūn. Maqṣūd az
dhekr-e īn kalemāt ānke mabādā to rā ghaflat eḥāṭe va akhdh namāyad va
dar ḥīn-e eshrāq-e shomūs-e moshreqe tamassok be-radd va qabūl-e aḥadī
namāʾī. Bārī tamassok be-aḥadī najoste va khod be-kaynūnat-e khod
be-eshrāq-e ān shams-e jūd va faḍl eqbāl namā ke hamān eshrāq aẓhar az
koll-e dalāʾel-e motqane va aʿẓam az jamīʿ-e ḥojaj-e moḥkame būde va
khwāhad būd, va īn maḍmūn dar kotob-e monzale madhkūr ast.

In short, exercise the utmost care and vigilance, lest you measure the
mighty, surging, billowing Ocean by a drop, for nothing other than that
Ocean has any standing in His presence. Glory be to God, beyond what you
say! The purpose of these words is to prevent heedlessness from
encompassing and seizing you, and to keep you from relying on anyone's
rejection or acceptance when the radiant Suns shine forth. Rely, then,
on no one, but turn with your own being toward the radiance of that Sun
of bounty and grace, for that very radiance has always been, and will
remain, more evident than every sound argument and greater than all
firmly established proofs. This meaning is set forth in the revealed
books.

\subsection{44}\label{section-2637}

\begin{otherlanguage}{persian}

حال ای اهل ظنون و هوی بشعور آمده در علوّ امر اللّه و سموّ کلمة اللّه تفکّر
نمائید که شاید انفس محدوده و حجبات غلیظه مابین شما و عرفان ملیک اسماء
حایل نشوند و از همهٴ این مقام گذشته انسان یا مجنون است و یا عاقل مجنون
حکمی باو راجع نه و بامری مأمور نه و عاقل که صاحب عقل مستقیم و فؤاد
مستضیئه و بصر باصره و سمع سامعه شده البتّه استعداد عرفان آن شمس ایقان در
او بوده چه اگر صاحب اینمقام اعظم نمیبود مأمور باین امر معظم و عرفان
کلمهٴ اتمّ نمیشد کذلک یتمّ اللّه لکم حجّته و یظهر علیکم سبل عرفانه لیکون
الحجّة بالغة علیکم لعلّ انتم فی آیات اللّه تتعقّلون

\end{otherlanguage}

Ḥāl ey ahl-e ẓonūn va havā, be-shoʿūr āmade dar ʿolovv-e amr-e Allāh va
somovv-e kalimat-e Allāh tafakkor namāʾīd, ke shāyad anfos-e maḥdūde va
ḥojobāt-e ghalīẓe mā-bayn-e shomā va ʿerfān-e malīk-e asmā ḥāʾel
nashavand. Va az hame-ye īn maqām godhashte, ensān yā majnūn ast va yā
ʿāqel. Majnūn ḥokmī be-ū rājeʿ na va be-amrī maʾmūr na, va ʿāqel ke
ṣāḥeb-e ʿaql-e mostaqīm va foʾād-e mostaḍīʾe va baṣar-e bāṣere va samʿ-e
sāmeʿe shode, albatte esteʿdād-e ʿerfān-e ān shams-e īqān dar ū būde,
che agar ṣāḥeb-e īn maqām-e aʿẓam nemībūd maʾmūr be-īn amr-e moʿaẓẓam va
ʿerfān-e kalame-ye atamm nemīshod. Kadhālika yutimmu Allāhu lakum
ḥujjatahu va yuẓhiru ʿalaykum subula ʿirfānihi li-yakūna al-ḥujjatu
bālighatan ʿalaykum laʿalla antum fī āyāti Allāhi tataʿaqqalūn.

Now, O people of conjecture and desire, awaken to awareness and reflect
on the loftiness of the Cause of God and the sublimity of the Word of
God, so that limited souls and thick veils may not stand between you and
recognition of the Sovereign of Names. Beyond all this, a person is
either insane or rational. No ruling applies to the insane, nor are they
charged with any command. But the rational person, endowed with sound
reason, an illumined heart, seeing eyes, and hearing ears, certainly has
the capacity to recognize that Sun of certitude. For if he did not
possess this supreme capacity, he would not have been charged with this
mighty Cause and with recognizing the most perfect Word. Thus God
completes His proof for you and reveals to you the paths to His
recognition, so that the proof against you may be conclusive, and that
you may reflect on the verses of God.

\subsection{45}\label{section-2638}

\begin{otherlanguage}{persian}

باری ملخّص کلام و جوهر آن آنکه الیوم بر هر نفسی لازم و واجب است که تمسّک
بفلک هدایت نوح امر جسته که شاید از غرق بحر نفس و هوی و وهم و ظنون اهل
بغی و فحشا نجات یابد فهنیئاً للمستقربین بهذه السّدرة الالهیّة و المستروحین
بهذه الأریاح الأحدیّة و المستقبلین بهذا النّور القدمیّة و هر نفسیکه در مغرب
بعد مستعدّ شود برای فیوضات این سحاب مرتفعه و غمام منیعه البتّه از جمیع
جهات افاضات روح القدس و نفحهٴ آن اقرب من آن او را احاطه فرماید و جمیع
مقامات متعالیه و کلمات جامعه و مراتب عالیه که در او مکنون و مستور بوده
از نفس او باو ظاهر نماید و از قلب او باو القا فرماید

\end{otherlanguage}

Bārī molakhkhaṣ-e kalām va jawhar-e ān ānke al-yawm bar har nafsī lāzem
va vājeb ast ke tamassok be-folk-e hedāyat-e Nūḥ-e amr joste, ke shāyad
az gharq-e baḥr-e nafs va havā va vahm va ẓonūn-e ahl-e baghy va faḥshā
nejāt yābad. Fa-hanīʾan li-al-mustaqribīna bi-hādhihi al-sidrati
al-ilāhīyati va al-mustarwiḥīna bi-hādhihi al-aryāḥi al-aḥadīyati va
al-mustaqbilīna bi-hādhā al-nūri al-qidamīya. Va har nafsī ke dar
maghreb-e boʿd mostaʿedd shavad barā-ye foyūḍāt-e īn saḥāb-e mortafeʿe
va ghamām-e manīʿe, albatte az jamīʿ-e jehāt efāḍāt-e Rūḥ al-Qodos va
nafḥe-ye ān aqrab men ān ū rā eḥāṭe farmāyad, va jamīʿ-e maqāmāt-e
motaʿālīye va kalemāt-e jāmeʿe va marāteb-e ʿālīye ke dar ū maknūn va
mastūr būde az nafs-e ū be-ū ẓāher namāyad va az qalb-e ū be-ū elqā
farmāyad.

In short, the sum and essence of the matter is that today it is
necessary and incumbent upon every soul to hold fast to the ark of
guidance of the Noah of the Cause, so that he may be saved from drowning
in the sea of self and desire and in the fancies and conjectures of the
people of rebellion and depravity. Blessed, then, are those who draw
near to this divine Lote Tree, who find refreshment in these winds of
Oneness, and who turn toward this ancient Light. Any soul who, in the
west of remoteness, becomes receptive to the outpourings of this lofty
cloud and exalted rain cloud will assuredly be encompassed from every
direction, in less than an instant, by the effusions of the Holy Spirit
and its breath. It will reveal to him, from within his own self, all the
transcendent stations, all-embracing words, and lofty degrees that lay
concealed and hidden within him, and will impart them to him from his
own heart.

\subsection{46}\label{section-2639}

\begin{otherlanguage}{arabic}

مثلاً در حبّه مشاهده نما که چه مقدار جثّه و وزن دارد ولکن چون در ارض طیّبه و
مدینهٴ مقدّسه مقرّ یابد بتربیت شمس و هبوب اریاح چگونه جلوه‌های مستوره در او
را از اغصان و افنان و اوراق و اثمار از نفس او بنفس او مشهود و مرئی نماید
حال ای محتجبان حجبات وهمیّه مستعدّ این فضل اکبر و فیض اعظم شوید که برشحی
از آن نعمتهای مکنونه و لآلی مخزونه از صدور منیره ظاهر و مشهود گردد از
رحمت بی‌زوال محروم مشو و از نفحهٴ لایزال مأیوس مباش انّه لا ییأس من روح
اللّه الّا القوم الفاسقون

\end{otherlanguage}

Mathalan dar ḥabbe moshāhede namā ke che meqdār joth the va vazn dārad,
valāken chūn dar arḍ-e ṭayyebe va madīne-ye moqaddase maqarr yābad,
be-tarbīyat-e shams va hobūb-e aryāḥ chegūne jelve-hā-ye mastūre dar ū
rā az aghṣān va afnān va awrāq va athmār az nafs-e ū be-nafs-e ū mashhūd
va marʾī namāyad. Ḥāl ey moḥtajebān-e ḥojobāt-e vahmīye, mostaʿedd-e īn
faḍl-e akbar va fayḍ-e aʿẓam shavīd ke be-rashḥī az ān neʿmat-hā-ye
maknūne va laʾālī-ye makhzūne az ṣodūr-e monīre ẓāher va mashhūd gardad.
Az raḥmat-e bī-zavāl maḥrūm mashow va az nafḥe-ye lā-yazāl maʾyūs
mabāsh. Innahu lā yayʾasu min rawḥi Allāhi illā al-qawmu al-fāsiqūn.

Consider, for example, a seed: how little bulk and weight it has. Yet
when it finds a resting place in good soil and the holy city, see how,
through the sun's nurturing and the blowing of the winds, it makes the
splendors concealed within it---branches, boughs, leaves, and
fruit---manifest and visible to itself from its own being. Now, O you
who are veiled by the veils of fancy, make yourselves receptive to this
greatest favor and this supreme grace, so that through a mere sprinkling
of it, hidden blessings and treasured pearls may appear and become
manifest from illumined hearts. Do not deprive yourself of unfailing
mercy, and do not despair of the everlasting breath. Indeed, none
despairs of God's solace except the transgressors.

\subsection{47}\label{section-2640}

\begin{otherlanguage}{persian}

و دیگر سؤال شده بود که شمس بها حال در چه مقام ظاهر فواللّه تکاد السّموات
ان تتفطّرنّ من قولک و تنشقّ ارض الارادة و تندکّ کلّ جبل شامخ عظیم تب الی اللّه
الّذی خلقک و سوّاک ثمّ ارجع الیه عن قولک ثمّ استغفره عمّا ظننت فی نفسک و کن
من التّائبین ایّاک ان لا تذکر بمثل ما ذکرت لأنّ بذلک یضطرب کلّ ذی علوّ و تزلّ
اقدام العارفین و لما کان العبد لن یحبّ ان یخیّب آملیه او یردّ سائلیه او
یحزن محبّیه او یکدّر مشتاقیه لذا یمطر علیک من هذا السّحاب الطّاهر المنیر علی
ما اراد اللّه لک لا علی ما کان علیه لتکون من الشّاکرین تعهّد فی نفسک بأن لا
تسئل فی مثل ذلک عنّی و لا من غیری و لا یخطر مثله فی قلبک و لا تحرّک لسانک
علیه و کن فی حصن عصیم اتّق اللّه و لا تسئل عن اشیآء ان تبد لک تسؤک فاقنع
بما اظهرناه لکم فی سنة الثّمانین و من قبله فی التّسع ان انتم من الموقنین

\end{otherlanguage}

Va dīgar soʾāl shode būd ke shams-e Bahā ḥāl dar che maqām ẓāher.
Fa-wa-Allāhi takādu al-samāwātu an yatafaṭṭarna min qawlika va tanshaqqu
arḍu al-irādati va tandakku kullu jabalin shāmikhin ʿaẓīm. Tub ilā
Allāhi alladhī khalaqaka va sawwāka thumma irjiʿ ilayhi ʿan qawlika
thumma istaghfirhu ʿammā ẓananta fī nafsika va kun mina al-tāʾibīn.
Iyyāka an lā tadhkura bi-mithli mā dhakarta li-anna bi-dhālika yaḍṭaribu
kullu dhī ʿulūwin va tazillu aqdāmu al-ʿārifīn. Va lammā kāna al-ʿabdu
lan yuḥibba an yukhayyiba āmilīhi aw yarudda sāʾilīhi aw yuḥzina
muḥibbīhi aw yukaddira mushtāqīhi li-dhā yumṭiru ʿalayka min hādhā
al-saḥābi al-ṭāhiri al-munīri ʿalā mā arāda Allāhu laka lā ʿalā mā kāna
ʿalayhi li-takūna mina al-shākirīn. Taʿahhad fī nafsika bi-an lā tasʾala
fī mithli dhālika ʿannī va lā min ghayrī va lā yakhṭura mithluhu fī
qalbika va lā tuḥarrika lisānaka ʿalayhi va kun fī ḥiṣnin ʿaṣīm. Ittaqi
Allāha va lā tasʾal ʿan ashyāʾa in tubda laka tasuʾka fa-iqnaʿ bi-mā
aẓharnāhu lakum fī sanati al-thamānīna va min qablihi fī al-tisʿi in
antum mina al-mūqinīn.

Another question had been asked: in what station is the Sun of Bahā now
manifest? By God, the heavens are almost rent apart by your words, the
earth of the Will splits asunder, and every lofty, mighty mountain is
shattered. Repent to God, who created and fashioned you; then turn back
to Him from what you have said, ask His forgiveness for what you have
imagined within yourself, and be among the repentant. Beware of saying
anything like what you have said, for it causes every exalted one to
tremble and the feet of those who recognize to slip. Since this servant
does not wish to disappoint those who place their hopes in Him, turn
away those who question Him, grieve those who love Him, or distress
those who long for Him, He therefore rains upon you from this pure and
luminous cloud according to what God has intended for you, not according
to the station He Himself occupies, so that you may be thankful. Pledge
within yourself not to ask Me or anyone else anything of this kind, nor
to let its like cross your heart, nor to move your tongue to speak of
it, and remain within an impregnable stronghold. Fear God, and do not
ask about things which, if disclosed to you, would distress you. Be
content, then, with what We revealed to you in the year eighty and,
before that, in the year nine, if you are among those who are certain.

\subsection{48}\label{section-2641}

\begin{otherlanguage}{arabic}

فواللّه لو لم یکن حبّی ایّاک لأمسکت زمام القلم و ما اطلقته الی ابد الآبدین
لئلّا یحرّک علی ما حرّک علیه قلمک فسبحان اللّه عمّا اردت فی نفسک فتعالی عمّا
یتکلّم به السن هؤلآء القائلین ایّاک ایّاک یا عبد لا تذکر ما یرفع به ضجیج
المغلّین ثمّ صریخ الموحّدین ثمّ ضوضآء العارفین و لو کنت قابلاً فی نفسک
لأمرناک بأن تستغفر اللّه ربّک سبعین الف قرن من قرون الأوّلین و ان کنت
مستطیعاً فی الزّخارف لأمرناک بأن تنفق قنطاراً من الذّهب جزآء ما خرج عن فمک و
ظهر من قلمک و انّا کنّا حاکمین و ان کنت ذات جسد قوّیة لجلّدناک الف جلدة جزآء
ما سئلت لتکون من المتنبّهین

\end{otherlanguage}

Fa-wa-Allāhi law lam yakun ḥubbī iyyāka la-amsaktu zimāma al-qalami
wa-mā aṭlaqtuhu ilā abadi al-ābidīn li-allā yuḥarraka ʿalā mā ḥurrika
ʿalayhi qalamuka fa-subḥāna Allāhi ʿammā aradta fī nafsika fa-taʿālā
ʿammā yatakallamu bihi alsunu hāʾulāʾi al-qāʾilīn. Iyyāka iyyāka yā
ʿabdu lā tadhkur mā yurfaʿu bihi ḍajīju al-mughillīn thumma ṣarīkhu
al-muwaḥḥidīn thumma ḍawḍāʾu al-ʿārifīn wa-law kunta qābilan fī nafsika
la-amarnāka bi-an tastaghfira Allāha rabbaka sabʿīna alfa qarnin min
qurūni al-awwalīn wa-in kunta mustaṭīʿan fī al-zakhārifi la-amarnāka
bi-an tunfiqa qinṭāran mina al-dhahabi jazāʾa mā kharaja ʿan famika
wa-ẓahara min qalamika wa-innā kunnā ḥākimīn wa-in kunta dhāta jasadin
qawiyyatin la-jalladnāka alfa jaldatin jazāʾa mā saʾalta li-takūna mina
al-mutanabbihīn.

By God, were it not for My love for you, I would have restrained the Pen
and never released it for all eternity, lest it be moved to write what
your pen has written. Glorified be God above what you have intended
within yourself, and exalted be He above what the tongues of these
speakers utter. Beware, beware, O servant! Do not mention what raises
the clamor of those harboring rancor, the cries of those who affirm
God's unity, and the uproar of those who recognize Him. Were you
inwardly receptive, We would have commanded you to seek the forgiveness
of God, your Lord, for seventy thousand ages of the ages of the
ancients. Were you possessed of worldly wealth, We would have commanded
you to spend a hundredweight of gold in recompense for what has issued
from your mouth and appeared from your pen. We have indeed passed
judgment. And had you a strong body, We would have scourged you with a
thousand lashes in recompense for what you asked, so that you might be
among those who take heed.

\subsection{49}\label{section-2642}

\begin{otherlanguage}{arabic}

هل ینبغی للعلقة بأن یسئل عن الّذی لن یبلغ الی عرفانه کلّ ذی فضل رفیع او
للتّراب بأن یسئل عن ربّ الأرباب فتعالی عن کلّ ما انتم علیه و عن کلّ ما انتم
عنه لمن السّائلین هل للمفقود ذکر عند سلطان الوجود او للعدم تلقآء القدم او
للفانی بأن یرتقی الی جبروت البقآء مقرّ قدس منیع و ان الأعمی لن یقدر ان
یعرف الشّمس بأسرها و لابدّ بأن یکتفی بما یتجلّی علیه من حرارة الّتی کانت من
اثرها کذلک فاعرف شأنک و شأن السّائلین و انّ العلقة حقّ له بأن یسئل ربّه مقام
المضغة و کذلک المضغة مقام العظام ثمّ العظام مقام اکتسآء اللّحم الی ان یبلغ
مقام الّذی قدّر له فی رتبته ان انتم من المدرکین و انّ المضغة لو یسئل عن
اللّه بأن یرفعه فی الحین الی مقام کلّ بالغ علیم هذا الطّلب منه مردود الیه
لأنّه لن یستطیع ذلک من دون ان یتعارج الی المعارج الی ان یصل الی هذا
المقام العالی المنیع فانظر فی البذر یبذر فی الأرض ثمّ یمضی علیه شهور
معلومة و یدبّره ایادی التّدبیر بما قدّر اللّه فی نفس الأرض الی ان یخرج منه
کما انتم تشهدونه و تکوننّ من الشّاهدین و اذا اخرجه الأرض لن یغنی به احد
الّا بأن یبلغ الی مقامه کذلک فاعرفوا حکمة اللّه ان انتم من الطّالبین

\end{otherlanguage}

Hal yanbaghī li-l-ʿalaqati bi-an yasʾala ʿani alladhī lan yablugha ilā
ʿirfānihi kullu dhī faḍlin rafīʿ aw li-l-turābi bi-an yasʾala ʿan rabbi
al-arbāb fa-taʿālā ʿan kulli mā antum ʿalayhi wa-ʿan kulli mā antum
ʿanhu la-mina al-sāʾilīn. Hal li-l-mafqūdi dhikrun ʿinda sulṭāni
al-wujūd aw li-l-ʿadami tilqāʾa al-qidam aw li-l-fānī bi-an yartaqiya
ilā jabarūti al-baqāʾi maqarri qudsin manīʿ wa-inna al-aʿmā lan yaqdira
an yaʿrifa al-shamsa bi-asrihā wa-lā budda bi-an yaktafiya bi-mā yata
jallā ʿalayhi min ḥarārati allatī kānat min atharihā kadhālika fa-iʿrif
shaʾnaka wa-shaʾna al-sāʾilīn wa-inna al-ʿalaqata ḥaqq lahu bi-an
yasʾala rabbahu maqāma al-muḍghati wa-kadhālika al-muḍghatu maqāma
al-ʿiẓāmi thumma al-ʿiẓāmu maqāma iktisāʾi al-laḥmi ilā an yablugha
maqāma alladhī quddira lahu fī rutbatihi in antum mina al-mudrikīn
wa-inna al-muḍghata law yasʾalu ʿani Allāhi bi-an yarfaʿahu fī al-ḥīni
ilā maqāmi kulli bālighin ʿalīm hādhā al-ṭalabu minhu mardūdun ilayhi
li-annahu lan yastaṭīʿa dhālika min dūni an yataʿāraja ilā al-maʿāriji
ilā an yaṣila ilā hādhā al-maqāmi al-ʿālī al-manīʿ fa-nẓur fī al-badhri
yubdharu fī al-arḍi thumma yamḍī ʿalayhi shuhūrun maʿlūmatun
wa-yudabbiruhu ayādī al-tadbīri bi-mā qaddara Allāhu fī nafsi al-arḍi
ilā an yakhruja minhu ka-mā antum tashhadūnahu wa-takūnanna mina
al-shāhidīn wa-idhā akhrajahu al-arḍu lan yaghnā bihi aḥadun illā bi-an
yablugha ilā maqāmihi kadhālika fa-iʿrifū ḥikmata Allāhi in antum mina
al-ṭālibīn.

Is it fitting for a blood clot to ask about Him whom no possessor of
lofty excellence can ever recognize, or for dust to ask about the Lord
of lords? Exalted is He above all that you hold to and all that you
inquire about. Can the nonexistent be mentioned in the presence of the
Sovereign of existence, or nothingness before the Eternal? Can the
perishable ascend to the dominion of everlasting life, an inaccessible
seat of holiness? The blind can never know the sun in its entirety, but
must be content with the warmth that reaches them as an effect of it.
Recognize, then, your own condition and that of those who ask. A blood
clot is entitled to ask its Lord for the stage of a lump of flesh;
likewise, the lump of flesh may ask for the stage of bones, and the
bones for the stage of being clothed with flesh, until it reaches the
station ordained for it in its own rank, if you are among those who
understand. If the lump of flesh should ask God to raise it instantly to
the station of a mature and knowledgeable person, this request would be
rejected, for it cannot attain this without ascending through the stages
until it reaches this lofty and inaccessible station. Consider a seed:
it is sown in the earth, then a certain number of months pass over it,
and the hands of providence tend it through what God has ordained within
the earth itself, until it emerges from it, as you observe. Be among
those who bear witness. Once the earth has brought it forth, no one can
benefit from it until it reaches its proper stage. Recognize, then, the
wisdom of God, if you are among those who seek.

\subsection{50}\label{section-2643}

\begin{otherlanguage}{arabic}

و کذلک فانظر فی الشّجرة اذا غرس لن یثمر فی الحین الّا بأن یعرج الی مراتب
الّتی قدّر اللّه لها اذاً یأتی بفواکه عزّ منیع ولو انّ ربّک یقدر بأن یظهر کلّ ما
کان و ما یکون فی کفّ من الطّین لأنّه هو القادر المقتدر العلیم الحکیم

\end{otherlanguage}

Wa-kadhālika fa-nẓur fī al-shajarati idhā ghurisa lan yuthmira fī
al-ḥīni illā bi-an yaʿruja ilā marātibi allatī qaddara Allāhu lahā idhan
yaʾtī bi-fawākihi ʿizzin manīʿ wa-law anna rabbaka yaqdiru bi-an yuẓhira
kulla mā kāna wa-mā yakūnu fī kaffin mina al-ṭīni li-annahu huwa
al-Qādiru al-Muqtadiru al-ʿAlīmu al-Ḥakīm.

Likewise, consider a tree: when planted, it does not bear fruit at once,
but must ascend through the stages God has ordained for it. Then it
produces fruits of inaccessible glory. Yet your Lord has the power to
bring forth all that has been and all that will be from a handful of
clay, for He is the Powerful, the All-Powerful, the All-Knowing, the
All-Wise.

\subsection{51}\label{section-2644}

\begin{otherlanguage}{arabic}

و لو یشآء ان یبعث من النّطفة خلق الأوّلین و الآخرین فی اقلّ من الحین لیقدر
و انّه لهو المقدّر المدبّر الحکیم العلیم ولکن لما قدّر لکلّ شیء مراتب و درایج
و مقادیر یحبّ بلوغ کلّ شیء بما قدّر له من هذه المراتب لتکوننّ فی حکمة اللّه
لمن المتفرّسین و انت فاجهد فی نفسک بأن تری کلّ شیء فی مقامه علی ما هو علیه
لتکون علی صراط عدل مستقیم

\end{otherlanguage}

Wa-law yashāʾu an yabʿatha mina al-nuṭfati khalqa al-awwalīna
wa-al-ākhirīna fī aqalla mina al-ḥīni la-yaqdiru wa-innahu la-huwa
al-Muqaddiru al-Mudabbiru al-Ḥakīmu al-ʿAlīm wa-lākin lammā qaddara
li-kulli shayʾin marātiba wa-darāyija wa-maqādīra yuḥibbu bulūgha kulli
shayʾin bi-mā qaddara lahu min hādhihi al-marātibi li-takūnanna fī
ḥikmati Allāhi la-mina al-mutafarrisīn wa-anta fa-jhad fī nafsika bi-an
tarā kulla shayʾin fī maqāmihi ʿalā mā huwa ʿalayhi li-takūna ʿalā
ṣirāṭi ʿadlin mustaqīm.

If He wished to bring forth from a drop of seed the creation of the
first and the last in less than an instant, He would have the power to
do so. He is indeed the Ordainer, the Disposer, the All-Wise, the
All-Knowing. Yet, since He has ordained stages, degrees, and measures
for every thing, He desires that each thing attain its fulfillment
through the stages He has ordained for it, so that you may be among
those who discern the wisdom of God. Strive within yourself to see every
thing in its own station, as it truly is, so that you may follow a
straight path of justice.

\subsection{52}\label{section-2645}

\begin{otherlanguage}{arabic}

و انّی لمّا اجد حبّی ایّاک اوصیک بأن تحفظ نفسک فی تلک الأیّام الّتی اخذت کلّ
نفس سکرها و کانوا علی غفلة مبین لأنّ اللّه خلقهم و رزقهم و اماتهم و احیاهم
ثمّ بعثهم لعرفان مظهر امره و اذاً کلّهم محتجبون و تمسّکوا بما عندهم من
الظّنون و الأوهام ثمّ اعترضوا بالّذی خلقوا له و لا یشعرون و منهم من یقرأ
کلمات اللّه فی کلّ حین ثمّ عن جماله هم معرضون و یتکلّمون بما نزل فی البیان
ثمّ فی انفسهم لا یعلمون و یشهدون انفسهم فی اعلی الجنان و هم فی اسفل
الجحیم لو کانوا یشعرون و یأمرون النّاس بالبرّ ثمّ ینسون انفسهم کما انتم فی
اعمالهم تشهدون و ینسبون انفسهم الی اللّه ثمّ علی جماله یستکبرون

\end{otherlanguage}

Wa-innī lammā ajidu ḥubbī iyyāka ūṣīka bi-an taḥfaẓa nafsaka fī tilka
al-ayyāmi allatī akhadhat kulla nafsin sukruhā wa-kānū ʿalā ghaflatin
mubīn li-anna Allāha khalaqahum wa-razaqahum wa-amātahum wa-aḥyāhum
thumma baʿathahum li-ʿirfāni maẓhari amrihi wa-idhan kulluhum muḥtajibūn
wa-tamassakū bi-mā ʿindahum mina al-ẓunūni wa-al-awhāmi thumma iʿtaraḍū
bi-alladhī khuliqū lahu wa-lā yashʿurūn wa-minhum man yaqraʾu kalimāti
Allāhi fī kulli ḥīnin thumma ʿan jamālihi hum muʿriḍūn wa-yatakallamūna
bi-mā nuz zila fī al-Bayāni thumma fī anfusihim lā yaʿlamūn
wa-yashhadūna anfusahum fī aʿlā al-jināni wa-hum fī asfali al-jaḥīmi law
kānū yashʿurūn wa-yaʾmurūna al-nāsa bi-al-birri thumma yansawna
anfusahum ka-mā antum fī aʿmālihim tashhadūn wa-yansibūna anfusahum ilā
Allāhi thumma ʿalā jamālihi yastakbirūn.

Because I feel My love for you, I counsel you to guard yourself in these
days, when every soul has been seized by its own intoxication and they
remain in manifest heedlessness. God created them, sustained them,
caused them to die, brought them to life, and then raised them up to
recognize the Manifestation of His Cause; yet they are all veiled. They
have clung to their own suppositions and imaginings, and then opposed
Him for whose sake they were created, without perceiving it. Among them
are those who read the words of God at every moment, yet turn away from
His Beauty. They speak of what has been sent down in the Bayān, yet
within themselves they do not understand. They regard themselves as
being in the highest gardens of Paradise, while they are in the lowest
depths of hell, if only they perceived it. They enjoin righteousness on
others, then forget themselves, as you observe in their deeds. They
claim a relationship to God, then act arrogantly toward His Beauty.

\subsection{53}\label{section-2646}

\begin{otherlanguage}{persian}

فانظر فی تلک الأیّام کلّما سمعت و رأیت فی ایّام اللّه من قبل بحیث ما بدّل امر
منه ان انتم تعرفون فلمّا جائهم علیّ بالحقّ بکتاب مشهود اعترضوا ثمّ استکبروا
علیه علی ما کانوا علیه ان یستطیعون و اعرضوا عنه ثمّ استدلّوا بما عندهم من
کتاب اللّه المهیمن القیّوم و فی کلّ یوم اتّخذوا من الفرقان ما به یستدلّون
لیدحضوا به الحقّ کما انتم من هؤلآء تشهدون بحیث یحرّکون السنتهم بالباطل و
یقولون هذا من عند اللّه و ما هو من عند اللّه بل من انفسهم لو انتم تفقهون و
لو کان من عند اللّه لم یعترضوا علی اللّه الّذی خلقهم و سوّاهم و کذلک سوّلت
لهم انفسهم ما کانوا یکسبون و فی کلّ یوم یتمسّکون بکلمة من الکتاب و یعترضون
علی الّذی منه ظهرت صحایف اللّه العزیز المحبوب

\end{otherlanguage}

Fa-nẓur fī tilka al-ayyāmi kulla-mā samiʿta wa-raʾayta fī ayyāmi Allāhi
min qablu bi-ḥaythu mā buddila amrun minhu in antum taʿrifūn fa-lammā
jāʾahum ʿAliyyun bi-al-ḥaqqi bi-kitābin mashhūd iʿtaraḍū thumma
istakbarū ʿalayhi ʿalā mā kānū ʿalayhi an yastaṭīʿūn wa-aʿraḍū ʿanhu
thumma istadallū bi-mā ʿindahum min kitābi Allāhi al-Muhaymini al-Qayyūm
wa-fī kulli yawmin ittakhadhū mina al-Furqāni mā bihi yastadillūna
li-yudḥiḍū bihi al-ḥaqqa ka-mā antum min hāʾulāʾi tashhadūn bi-ḥaythu
yuḥarrikūna alsinatahum bi-al-bāṭili wa-yaqūlūna hādhā min ʿindi Allāhi
wa-mā huwa min ʿindi Allāhi bal min anfusihim law antum tafqahūn wa-law
kāna min ʿindi Allāhi lam yaʿtariḍū ʿalā Allāhi alladhī khalaqahum
wa-sawwāhum wa-kadhālika sawwalat lahum anfusuhum mā kānū yaksibūn wa-fī
kulli yawmin yatamassakūna bi-kalimatin mina al-kitābi wa-yaʿtariḍūna
ʿalā alladhī minhu ẓaharat ṣaḥāʾifu Allāhi al-ʿAzīzi al-Maḥbūb.

Observe in these days all that you heard and saw in the days of God in
the past: nothing of it has changed, if you recognize it. When ʿAlī came
to them with the truth and a manifest Book, they objected and acted
arrogantly toward Him to the utmost of their ability. They turned away
from Him and then argued from what they possessed of the Book of God,
the All-Protecting, the Self-Subsisting. Every day they took from the
Qur'an something to use as evidence to refute the truth, just as you
observe these people doing. They move their tongues with falsehood and
say, ``This is from God.'' But it is not from God; rather, it is from
themselves, if only you understood. Had it been from God, they would not
have opposed God, who created and fashioned them. Thus their own selves
made what they were doing seem pleasing to them. Every day they cling to
a word of the Book and oppose Him from whom the scriptures of God, the
Almighty, the Beloved, have appeared.

\subsection{54}\label{section-2647}

\begin{otherlanguage}{arabic}

قل لن یکفیکم بعد ظهوره شیء ان انتم بکلّ شیء تتمسّکون قل کلّ الکتب یشهد بذلک
و کلّ الصّحف دلیل علیه ان انتم تقرؤون بحیث لن تجد فی الکتب حرفاً الّا و قد
شهدت فی نفسها بأنّ هذا لحقّ علی الحقّ من لدن عزیز قیّوم و من دون ذلک اسئلکم
یا ملأ البیان کلمة بالصّدق الخالص ان انتم تنصفون لو کان عندکم الف کتاب و
هذا فی حدّ الاحصآء بل یکون علی قدر الّذی انتم لن تقدروا ان تحصون و فی ظاهر
هذه الکتب انتم تشهدون و تستدلّون بما فیها علی امر عندکم کما انّ کلّ الملل
بمثل ذلک یعملون و یأتیکم احد علی غیر ما عندکم عمّا عرفتم من حکم الکتاب
اذاً ما تقولون أ تجاهدون بالّذی جائکم بآیات اللّه المهیمن القیّوم کما تشهدون
حینئذ بالحقّ و بأبصارکم انتم تنظرون او ترجعون الخطأ علی انفسکم بحیث قرأتم
کتب اللّه و ما عرفتم المقصود فوا حسرة علیکم یا ملأ الغافلون

\end{otherlanguage}

Qul lan yakfiyakum baʿda ẓuhūrihi shayʾun in antum bi-kulli shayʾin
tatamassakūn qul kullu al-kutubi yashhadu bi-dhālika wa-kullu al-ṣuḥufi
dalīlun ʿalayhi in antum taqraʾūn bi-ḥaythu lan tajida fī al-kutubi
ḥarfan illā wa-qad shahidat fī nafsihā bi-anna hādhā la-ḥaqqun ʿalā
al-ḥaqqi min ladun ʿAzīzin Qayyūm wa-min dūni dhālika asʾalukum yā
malaʾa al-Bayāni kalimatan bi-al-ṣidqi al-khāliṣi in antum tunṣifūn law
kāna ʿindakum alfu kitābin wa-hādhā fī ḥaddi al-iḥṣāʾi bal yakūnu ʿalā
qadri alladhī antum lan taqdirū an tuḥṣūn wa-fī ẓāhiri hādhihi al-kutubi
antum tashhadūna wa-tastadillūna bi-mā fīhā ʿalā amrin ʿindakum ka-mā
anna kulla al-milali bi-mithli dhālika yaʿmalūn wa-yaʾtīkum aḥadun ʿalā
ghayri mā ʿindakum ʿammā ʿaraftum min ḥukmi al-kitābi idhan mā taqūlūn
a-tujāhidūna bi-alladhī jāʾakum bi-āyāti Allāhi al-Muhaymini al-Qayyūmi
ka-mā tashhadūna ḥīnaʾidhin bi-al-ḥaqqi wa-bi-abṣārikum antum tanẓurūn
aw tarjiʿūna al-khaṭaʾa ʿalā anfusikum bi-ḥaythu qaraʾtum kutuba Allāhi
wa-mā ʿaraftumu al-maqṣūda fa-wā ḥasratan ʿalaykum yā malaʾa
al-ghāfilūn.

Say: After His appearance, nothing will suffice you, even if you cling
to every thing. Say: All the Books bear witness to this, and all the
scriptures are proof of it, if you read them. You will not find a single
letter in the Books that does not in itself testify that this is indeed
the truth, founded upon truth, from One who is Almighty and
Self-Subsisting. Beyond this, O people of the Bayān, I ask you for a
word in pure truth, if you judge fairly. Suppose you possessed a
thousand books---and this is still a countable number; rather, suppose
their number were beyond your power to count---and you considered the
outward meaning of these books and argued from their contents for some
matter you upheld, as all religious communities do. If someone then came
to you with something other than what you held, contrary to what you
understood of the Book's decree, what would you say? Would you contend
against Him who came to you with the verses of God, the All-Protecting,
the Self-Subsisting, as you now truly witness and behold with your own
eyes? Or would you attribute the error to yourselves, acknowledging that
you had read the Books of God but had not understood their intent? Alas
for you, O heedless people!

\subsection{55}\label{section-2648}

\begin{otherlanguage}{persian}

کما ان بعد محمّد رسول اللّه کلّ من دان فی الفرقان استدلّوا من الکتاب بأنّ
الرّسالة و النّبوّة ختمت به و لن یأتی بعده احد بأمر من اللّه و شرع من عنده و
هذا ما شهدتم بأبصارکم و بسمعهم عنهم تسمعون و انتم یا ملأ البیان کنتم
منهم و بذلک استدللتم فی ایّامکم و کنتم به مستدلّون الی ان جآء علیّ بالحقّ
اذاً رجع ظنونکم و عرفانکم الی انفسکم و اشرق امر اللّه بما اراد بقوله کن
فیکون کذلک فاعرفوا الأمر فی تلک الأیّام الی آخر الّذی لا آخر له ان انتم لا
تشکون فیما رأیتم و تکوننّ من الّذینهم بما یشهدون لموقنون

\end{otherlanguage}

Ka-mā anna baʿda Muḥammadin rasūli Allāhi kullu man dāna fī al-Furqāni
istadallū mina al-kitābi bi-anna al-risālata wa-al-nubuwwata khutimat
bihi wa-lan yaʾtiya baʿdahu aḥadun bi-amrin mina Allāhi wa-sharʿin min
ʿindihi wa-hādhā mā shahidtum bi-abṣārikum wa-bi-samʿihim ʿanhum
tasmaʿūn wa-antum yā malaʾa al-Bayāni kuntum minhum wa-bi-dhālika
istadlaltum fī ayyāmikum wa-kuntum bihi mustadillūna ilā an jāʾa
ʿAliyyun bi-al-ḥaqqi idhan rajaʿa ẓunūnukum wa-ʿirfānukum ilā anfusikum
wa-ashraqa amru Allāhi bi-mā arāda bi-qawlihi kun fa-yakūn kadhālika
fa-iʿrifū al-amra fī tilka al-ayyāmi ilā ākhiri alladhī lā ākhira lahu
in antum lā tashukkūna fī-mā raʾaytum wa-takūnanna mina alladhīna-hum
bi-mā yashhadūna la-mūqinūn.

Thus, after Muḥammad, the Messenger of God, all who professed faith in
the Qur'an argued from the Book that messengership and prophethood had
been sealed with Him, and that no one would come after Him with a Cause
from God and a law from His presence. This you have witnessed with your
own eyes, and you hear from them what they themselves have heard. You, O
people of the Bayān, were among them. In your own days you argued in
this way and continued to do so until ʿAlī came with the truth. Then
your suppositions and your understanding recoiled upon yourselves, and
the Cause of God shone forth as He desired, through His utterance,
``Be,'' and it is. Recognize, then, the Cause in these days and until
the end that has no end, if you do not doubt what you have seen, and be
among those who are certain of what they witness.

\subsection{56}\label{section-2649}

\begin{otherlanguage}{arabic}

قل انّه قد ظهر بالحقّ و لن یحجبه کلّما عندکم و لن یصدّه شیء عمّا فی السّموات و
الأرض و ینطق بالحقّ بأعلی صوته فی جبروت الأمر و الخلق و لن یخاف من احد
کما انتم تشهدون و انتم ان تقبلوا او تعرضوا فی حدّ سوآء و انّه لغنیّ عمّا
انتم تعقلون

\end{otherlanguage}

Qul innahu qad ẓahara bi-al-ḥaqqi wa-lan yaḥjubahu kullu-mā ʿindakum
wa-lan yaṣuddahu shayʾun ʿammā fī al-samāwāti wa-al-arḍi wa-yanṭiqu
bi-al-ḥaqqi bi-aʿlā ṣawtihi fī jabarūti al-amri wa-al-khalqi wa-lan
yakhāfa min aḥadin ka-mā antum tashhadūn wa-antum in tuqbilū aw tuʿriḍū
fī ḥaddi sawāʾin wa-innahu la-ghaniyyun ʿammā antum taʿqilūn.

Say: He has indeed appeared with the truth. Nothing you possess will
veil Him, and nothing in the heavens or on earth will hinder Him. He
proclaims the truth at the top of His voice in the dominion of command
and creation, and He fears no one, as you witness. Whether you turn
toward Him or turn away makes no difference, and He is indeed
independent of all that you comprehend.

\subsection{57}\label{section-2650}

\begin{otherlanguage}{arabic}

و انّا حینئذ نمثل لکم مثلاً بالفضل لعلّ انتم به بما تحبّون تصلون فأنّ البحر
علی ما هو علیه یکون علی مقامه و ان لن یتوجّه الیه احد لیأخذ عنه المآء او
یخرج منه لآلیه هل ینقص عنه شیء ان انتم تعلمون و کذلک الشّمس یضیء علی ما
کان و ان لن ینظر احد الیها هل ینقص عنها نورها و اشراقها اذاً تفکّروا فی
حکمة اللّه ان انتم تتفکّرون بل کلّ من یتوجّه الیه هذا من فضل اللّه علیه لأنّه
ینتفع منها ما لا ینتفع عن دونها و کذلک نضرب لکم الأمثال لعلّ انتم تفقهون

\end{otherlanguage}

Wa-innā ḥīnaʾidhin numaththilu lakum mathalan bi-al-faḍli laʿalla antum
bihi bi-mā tuḥibbūna taṣilūn fa-anna al-baḥra ʿalā mā huwa ʿalayhi
yakūnu ʿalā maqāmihi wa-in lan yatawajjaha ilayhi aḥadun li-yaʾkhudha
ʿanhu al-māʾa aw yukhrija minhu laʾāliʾahu hal yanquṣu ʿanhu shayʾun in
antum taʿlamūn wa-kadhālika al-shamsu yuḍīʾu ʿalā mā kāna wa-in lan
yanẓura aḥadun ilayhā hal yanquṣu ʿanhā nūruhā wa-ishrāquhā idhan
tafakkarū fī ḥikmati Allāhi in antum tatafakkarūn bal kullu man
yatawajjahu ilayhi hādhā min faḍli Allāhi ʿalayhi li-annahu yantafiʿu
minhā mā lā yantafiʿu ʿan dūnihā wa-kadhālika naḍribu lakumu al-amthāla
laʿalla antum tafqahūn.

Now, through Our grace, We set forth a parable for you, so that through
it you may attain what you desire. The sea remains as it is, in its own
station. If no one turns to it to draw water from it or bring forth its
pearls, is anything diminished from it, if you know? Likewise, the sun
shines as it always has. If no one looks toward it, are its light and
radiance diminished? Reflect, then, on the wisdom of God, if you are
among those who reflect. Rather, whoever turns toward it does so through
God's grace toward him, for he derives from it a benefit he cannot
derive from anything else. Thus We set forth parables for you, so that
you may understand.

\subsection{58}\label{section-2651}

\begin{otherlanguage}{persian}

یا قوم دعوا ما عندکم و خذوا ما عند اللّه بقوّة من عنده و ایّاکم ان لا
تحتجبون ثمّ اصتبغوا بصبغ اللّه فی تلک الأیّام ثمّ ادخلوا فی سرادق الّذی ارتفع
بالحقّ و اتّکلوا علیه و کونوا من الّذین لا یمنعهم لومة لائم و لا شماتة فاسق
و لا اعراض کلّ معرض مردود

\end{otherlanguage}

Yā qawmi daʿū mā ʿindakum wa-khudhū mā ʿinda Allāhi bi-quwwatin min
ʿindihi wa-iyyākum an lā taḥtajibūn thumma iṣtabighū bi-ṣibghi Allāhi fī
tilka al-ayyāmi thumma udkhulū fī surādiqi alladhī irtafaʿa bi-al-ḥaqqi
wa-ittakilū ʿalayhi wa-kūnū mina alladhīna lā yamnaʿuhum lawmatu lāʾimin
wa-lā shamātatu fāsiqin wa-lā iʿrāḍu kulli muʿriḍin mardūd.

O people, leave aside what you possess and take hold of what is with God
through strength from Him. Beware lest you become veiled. Take on the
coloring of God in these days, then enter the tabernacle that has been
raised in truth. Put your trust in Him, and be among those whom neither
the blame of a critic, nor the gloating of a wrongdoer, nor the
rejection of anyone who has turned away and been cast out can deter.

\subsection{59}\label{section-2652}

\begin{otherlanguage}{arabic}

قل انّا نشاهد الیوم بطرف القدس عباد الّذینهم یدّعون الایمان فی انفسهم بأنّهم
احتجبوا عمّا عرفوه من قبل و بذلک ظهر کذبهم و هم لا یشعرون اذاً فکّر فی نفسک
ان احتجبوا الیوم کلّ من فی السّموات و الأرض عن بارئهم هل یرجع الیه من نقص
لا فوعزّة اللّه المهیمن القیّوم بل یرجع کلّ دون خیر الی انفسهم لو هم یعرفون
و انّا بذلک نجد نفسنا فی علوّ الّذی لن یقاس بعلوّ و سرور الّذی لو یطفح رشحاً
منه علی الوجود و مظاهره کلّهم ینصعقون

\end{otherlanguage}

Qul innā nushāhidu al-yawma bi-ṭarfi al-qudsi ʿibāda alladhīna-hum
yaddaʿūna al-īmāna fī anfusihim bi-annahum iḥtajabū ʿammā ʿarafūhu min
qablu wa-bi-dhālika ẓahara kadhibuhum wa-hum lā yashʿurūn idhan fakkir
fī nafsika in iḥtajaba al-yawma kullu man fī al-samāwāti wa-al-arḍi ʿan
bāriʾihim hal yarjiʿu ilayhi min naqṣin lā fa-wa-ʿizzati Allāhi
al-Muhaymini al-Qayyūm bal yarjiʿu kullu dūni khayrin ilā anfusihim law
hum yaʿrifūn wa-innā bi-dhālika najidu nafsanā fī ʿuluwwi alladhī lan
yuqāsa bi-ʿuluwwin wa-surūri alladhī law yaṭfaḥu rashḥan minhu ʿalā
al-wujūdi wa-maẓāhirihi kulluhum yanṣaʿiqūn.

Say: Today, with the eye of holiness, We behold those servants who claim
faith for themselves, yet have become veiled from what they formerly
recognized. Thus their falsehood has become manifest, though they do not
perceive it. Reflect within yourself: if all who are in the heavens and
on earth were today to become veiled from their Creator, would any
deficiency fall upon Him? No, by the glory of God, the All-Protecting,
the Self-Subsisting! Rather, every ill would fall upon themselves, if
only they knew. Amid this, We find Ourself in an exaltation that cannot
be compared with any exaltation, and in a joy of which, were a mere drop
to overflow upon existence and its manifestations, they would all be
struck senseless.

\subsection{60}\label{section-2653}

\begin{otherlanguage}{arabic}

قل سبحانک ربّ یا محبوبی ثبّتنی علی امرک ثمّ اجعلنی من الّذین ما نقضوا میثاقک
و ما اتّبعوا اصنام ظنونهم ثمّ اجعل لی مقعد صدق عندک و هبنی من لدنک رحمة و
الحقنی بعبادک الّذین لا خوف علیهم و لا هم یحزنون ای ربّ لا تدعنی بنفسی و
لا تجعلنی محروماً عن عرفان مظهر نفسک و لا تکتبنی من الّذینهم غفلوا عن
لقائک

\end{otherlanguage}

Qul subḥānaka rabbi yā maḥbūbī thabbitnī ʿalā amrika thumma ijʿalnī mina
alladhīna mā naqaḍū mīthāqaka wa-mā ittabaʿū aṣnāma ẓunūnihim thumma
ijʿal lī maqʿada ṣidqin ʿindaka wa-habnī min ladunka raḥmatan wa-alḥiqnī
bi-ʿibādika alladhīna lā khawfun ʿalayhim wa-lā hum yaḥzanūn. Ay rabbi
lā tadaʿnī bi-nafsī wa-lā tajʿalnī maḥrūman ʿan ʿirfāni maẓhari nafsika
wa-lā taktubnī mina alladhīnahum ghafalū ʿan liqāʾika.

Say: Glory be to You, my Lord, my Beloved! Make me steadfast in Your
Cause, and make me one of those who have not broken Your covenant or
followed the idols of their own imaginings. Grant me a seat of truth in
Your presence, bestow mercy upon me from Yourself, and join me with Your
servants who have no fear and do not grieve. My Lord, do not leave me to
myself, do not deprive me of recognizing the Manifestation of Your own
Self, and do not number me among those who have been heedless of meeting
You.

\subsection{61}\label{section-2654}

\begin{otherlanguage}{arabic}

و اجعلنی یا الهی من الّذینهم الی جمالک ینظرون و منه یستلذّون بحیث لن
یبدّلوا آناً منه بملکوت ملک السّموات و الأرض و بکلّ ما کان و ما یکون ای ربّ
فارحمنی فی تلک الأیّام الّتی اخذت الغفلة کل سکّان ارضک ثمّ ارزقنی یا الهی
خیر ما عندک و انّک انت المقتدر العزیز الکریم الغفور

\end{otherlanguage}

Wa-ijʿalnī yā ilāhī mina alladhīnahum ilā jamālika yanẓurūn wa-minhu
yastalidhdhūn bi-ḥaythu lan yubaddilū ānan minhu bi-malakūti mulki
al-samāwāti wa-al-arḍ wa-bi-kulli mā kāna wa-mā yakūn. Ay rabbi
fa-irḥamnī fī tilka al-ayyāmi allatī akhadhati al-ghaflatu kulla sukkāni
arḍika thumma urzuqnī yā ilāhī khayra mā ʿindaka wa-innaka anta
al-Muqtadiru al-ʿAzīzu al-Karīmu al-Ghafūr.

And make me, my God, one of those who gaze upon Your beauty and delight
in it, so that they would not exchange a single moment of it for
sovereignty over the dominion of the heavens and the earth, or for all
that has been and will be. My Lord, have mercy on me in these days when
heedlessness has seized all the inhabitants of Your earth. Grant me, my
God, the best of what is with You. You are indeed the All-Powerful, the
Almighty, the Generous, the Ever-Forgiving.

\subsection{62}\label{section-2655}

\begin{otherlanguage}{arabic}

و لا تجعلنی یا الهی من الّذینهم بالأذن صمّآء و بالعین عمیآء و باللّسان
بکمآء و بالقلب هم لا یفقهون ای ربّ خلّصنی عن نار الجهل و الهوی ثمّ ادخلنی
فی جوار رحمتک الکبری ثمّ انزل علیّ ما قدّرته لأصفیائک و انّک انت المقتدر علی
ما تشآء و انّک انت المهیمن القیّوم

\end{otherlanguage}

Wa-lā tajʿalnī yā ilāhī mina alladhīnahum bi-al-udhuni ṣammāʾ
wa-bi-al-ʿayni ʿamyāʾ wa-bi-al-lisāni bakmāʾ wa-bi-al-qalbi hum lā
yafqahūn. Ay rabbi khalliṣnī ʿan nāri al-jahli wa-al-hawā thumma
adkhilnī fī jiwāri raḥmatika al-kubrā thumma anzil ʿalayya mā qaddartahu
li-aṣfiyāʾika wa-innaka anta al-Muqtadiru ʿalā mā tashāʾ wa-innaka anta
al-Muhayminu al-Qayyūm.

And do not make me, my God, one of those whose ears are deaf, whose eyes
are blind, whose tongues are mute, and whose hearts do not understand.
My Lord, deliver me from the fire of ignorance and desire, bring me into
the shelter of Your most great mercy, and send down upon me what You
have ordained for Your chosen ones. You are indeed powerful over
whatever You will, and You are indeed the All-Protecting, the
Self-Subsisting.

\subsection{63}\label{section-2656}

\begin{otherlanguage}{arabic}

و انتم یا ملأ البیان خافوا عن اللّه و لا تتّبعوا هویکم فاتّبعوا حکم اللّه
فیما نزل فی کلّ الألواح من سمآء عزّ مرفوع قل عنده علم کلّ شیء یعلم غیب
السّموات و الأرض و علم ما کان و ما یکون و یعلم ما فی قلوبکم و نفوسکم ممّا
تخفون او تجهرون لن یخفی عن اللّه ربّکم من شیء و لن یعزب عن علمه من شیء ان
انتم تعرفون

\end{otherlanguage}

Wa-antum yā malaʾa al-Bayān khāfū ʿani Allāhi wa-lā tattabiʿū hawākum
fa-ittabiʿū ḥukma Allāhi fīmā nazala fī kulli al-alwāḥi min samāʾi
ʿizzin marfūʿ. Qul ʿindahu ʿilmu kulli shayʾ yaʿlamu ghayba al-samāwāti
wa-al-arḍ wa-ʿilma mā kāna wa-mā yakūn wa-yaʿlamu mā fī qulūbikum
wa-nufūsikum mimmā tukhfūna aw tajharūn lan yakhfā ʿani Allāhi rabbikum
min shayʾ wa-lan yaʿzuba ʿan ʿilmihi min shayʾ in antum taʿrifūn.

And you, people of the Bayān, fear God and do not follow your desires.
Follow instead the decree of God in what has been sent down in all the
Tablets from a heaven of exalted glory. Say: With Him is the knowledge
of all things. He knows the unseen realities of the heavens and the
earth, and has knowledge of all that has been and will be. He knows what
is in your hearts and souls, whether you conceal it or proclaim it.
Nothing will be hidden from God, your Lord, and nothing will escape His
knowledge, if only you recognize this.

\subsection{64}\label{section-2657}

\begin{otherlanguage}{arabic}

و ان ترید ان یحفظک اللّه من کلّ فتنة و عن هؤلآء اتّق اللّه حقّ تقاتة ثمّ اعرض
عنهم فتوکّل علی اللّه المهیمن القیّوم قل انّا آمنّا باللّه و ملائکته و رسله و
بکلّ ما جائوا به من کلّ صحایف و الواح و زبر مکنون و بما نزل علی محمّد رسول
اللّه و بما نزل علی علیّ و بما کان علیه و بما یکون و بما نزل و ینزل حینئذ
من آیات اللّه المهیمن القیّوم کذلک وصّیناک من قبل و حینئذ لتکون من الّذینهم
موقنون

\end{otherlanguage}

Wa-in turīd an yaḥfaẓaka Allāhu min kulli fitnatin wa-ʿan hāʾulāʾ ittaqi
Allāha ḥaqqa tuqātah thumma aʿriḍ ʿanhum fa-tawakkal ʿalā Allāhi
al-Muhaymini al-Qayyūm. Qul innā āmannā bi-Allāhi wa-malāʾikatihi
wa-rusulihi wa-bi-kulli mā jāʾū bihi min kulli ṣaḥāyifa wa-alwāḥin
wa-zuburin maknūn wa-bi-mā nazala ʿalā Muḥammadin rasūli Allāhi wa-bi-mā
nazala ʿalā ʿAliyyin wa-bi-mā kāna ʿalayhi wa-bi-mā yakūn wa-bi-mā
nazala wa-yanzilu ḥīnaʾidhin min āyāti Allāhi al-Muhaymini al-Qayyūm.
Kadhālika waṣṣaynāka min qablu wa-ḥīnaʾidhin li-takūna mina alladhīnahum
mūqinūn.

And if you wish God to protect you from every trial and from these
people, fear God as He should be feared, then turn away from them and
put your trust in God, the All-Protecting, the Self-Subsisting. Say: We
believe in God, His angels, and His Messengers, and in everything they
have brought---all the scriptures, tablets, and hidden books---and in
what was sent down to Muḥammad, the Messenger of God, and in what was
sent down to ʿAlī, and in what has been with him and what will be, and
in the verses of God, the All-Protecting, the Self-Subsisting, that have
been and are now being sent down. Thus have We counseled you before and
now, so that you may be among those who are certain.

\subsection{65}\label{section-2658}

\begin{otherlanguage}{arabic}

قل هذا سبیلی ان انتم بسبیل اللّه تنظرون و انّ هذا دلیلی ان انتم بدلیل اللّه
تعتنون قل هذا حجّتی ان انتم بحجّة اللّه توقنون قل هذا صراط اللّه فی السّموات
و الأرض ان انتم تریدون ان تمرّون قل ان یعذّب اللّه احد بما آمن بهذه الآیات
المنزل المرسول فبأیّ حجّة یعذّب الّذینهم ما آمنوا بنقطة البیان و من قبله
بمحمّد رسول اللّه و من قبله بابن مریم و من قبله بموسی النّبیّ الی ان یرجع
الأمر الی البدیع الأوّل ان انتم تنصفون

\end{otherlanguage}

Qul hādhā sabīlī in antum bi-sabīli Allāhi tanẓurūn wa-inna hādhā dalīlī
in antum bi-dalīli Allāhi taʿtanūn. Qul hādhā ḥujjatī in antum
bi-ḥujjati Allāhi tūqinūn. Qul hādhā ṣirāṭu Allāhi fī al-samāwāti
wa-al-arḍ in antum turīdūna an tamurrūn. Qul in yuʿadhdhibi Allāhu aḥad
bi-mā āmana bi-hādhihi al-āyāti al-munzali al-mursal fa-bi-ayyi ḥujjatin
yuʿadhdhibu alladhīnahum mā āmanū bi-Nuqṭati al-Bayāni wa-min qablihi
bi-Muḥammadin rasūli Allāhi wa-min qablihi bi-Ibni Maryama wa-min
qablihi bi-Mūsā al-nabiyyi ilā an yarjiʿa al-amru ilā al-Badīʿi al-awwal
in antum tunṣifūn.

Say: This is My way, if you look to the way of God, and this is My
evidence, if you attend to the evidence of God. Say: This is My proof,
if you are certain of the proof of God. Say: This is the path of God in
the heavens and on earth, if you wish to travel it. Say: If God were to
punish anyone for believing in these verses that have been sent down and
dispatched, by what proof would He punish those who did not believe in
the Point of the Bayān, and before Him in Muḥammad, the Messenger of
God, and before Him in the Son of Mary, and before Him in Moses the
Prophet, until the matter returns to the first Originator, if you judge
fairly?

\subsection{66}\label{section-2659}

\begin{otherlanguage}{arabic}

ثمّ اعلم بأنّ الیوم من یتوقّف فی هذا الأمر الطّالع فقد توقّف فی کلّ امر ان
انتم تعرفون و من لن یؤمن بهذا فقد کفر بکلّ رسل من قبل الی ان ینتهی الی
الفطرة الأوّل و هذا لحقّ معلوم فکّر فی نفسک انّک ان لن تؤمن بنقطة البیان هل
یصدق علیک حکم الایمان بأحد من رسل اللّه من قبله و کذلک ان لن تؤمن بأحد من
الرّسل الی آخر الّذی لا آخر له هل یصدق علیکم حکم الایمان بأحد منهم فسبحان
اللّه عمّا انتم تظنّون و لا شکّ بأن لن یصدق الایمان لأحد الّا بأن یؤمن بکلّ ما
نزل من عند اللّه او ینزل و هذا ما رقم علی کلّ لوح مسطور

\end{otherlanguage}

Thumma iʿlam bi-anna al-yawma man yatawaqqaf fī hādhā al-amri al-ṭāliʿi
fa-qad tawaqqafa fī kulli amr in antum taʿrifūn wa-man lan yuʾmin
bi-hādhā fa-qad kafara bi-kulli rusulin min qablu ilā an yantahiya ilā
al-fiṭrati al-awwal wa-hādhā la-ḥaqqun maʿlūm. Fakkir fī nafsika innaka
in lan tuʾmin bi-Nuqṭati al-Bayāni hal yaṣduqu ʿalayka ḥukmu al-īmāni
bi-aḥadin min rusuli Allāhi min qablihi wa-kadhālika in lan tuʾmin
bi-aḥadin mina al-rusuli ilā ākhiri alladhī lā ākhira lahu hal yaṣduqu
ʿalaykum ḥukmu al-īmāni bi-aḥadin minhum fa-subḥāna Allāhi ʿammā antum
taẓunnūn wa-lā shakka bi-an lan yaṣduqa al-īmānu li-aḥadin illā bi-an
yuʾmina bi-kulli mā nazala min ʿindi Allāhi aw yanzilu wa-hādhā mā
ruqima ʿalā kulli lawḥin masṭūr.

Know, then, that today whoever hesitates over this dawning Cause has
hesitated over every Cause, if only you recognize this. Whoever does not
believe in this has rejected all the Messengers of the past, back to the
first creation. This is indeed an established truth. Reflect within
yourself: if you do not believe in the Point of the Bayān, can you truly
be considered a believer in any of God's Messengers before Him?
Likewise, if you do not believe in any one of the Messengers, to the end
that has no end, can you truly be considered believers in any of them?
God is exalted above what you imagine! There is no doubt that faith
cannot be truly ascribed to anyone unless that person believes in
everything that has been or will be sent down from God. This is what has
been inscribed upon every written Tablet.

\subsection{67}\label{section-2660}

\begin{otherlanguage}{arabic}

کذلک فاعرف امر اللّه و سنّته فی کلّ الأعصار لعلّ لن یشتبه علیکم الأمور
فانظروا فی المظاهر نسبتهم الی اللّه فیما ینزل علیهم لا بما هم به مختلفون
لا تشهد فی التّجلی الّا جمال المجلّی و لا فی التّنزیل الّا منزله ان انتم
تبصرون ثمّ اشهد جمال کلّ ظهور فی کلماته لیظهر لک الحقّ و تکون من العارفین
لأنّ اللّه قد جعل کلماتهم مرآة انفسهم و یلقی علی کلّ واحد علی قدر ما قابلوا
الشّمس و کذلک انتم فاشهدون و بهذا المراتب یظهر مراتبهم علیه ان انتم
تفقهون کلّ ذلک مراتب فی مقامات ظهور المرایا عند تجلّی الشّمس ولکن عن الشّمس
فی هذه الأیّام انتم لا تسئلون

\end{otherlanguage}

Kadhālika fa-iʿrif amra Allāhi wa-sunnatahu fī kulli al-aʿṣāri laʿalla
lan yashtabiha ʿalaykumu al-umūr fa-unẓurū fī al-maẓāhiri nisbatahum ilā
Allāhi fīmā yanzilu ʿalayhim lā bi-mā hum bihi mukhtalifūn. Lā tashhad
fī al-tajallī illā jamāla al-mujallī wa-lā fī al-tanzīli illā munzilahu
in antum tubṣirūn. Thumma ishhad jamāla kulli ẓuhūrin fī kalimātihi
li-yaẓhara laka al-ḥaqqu wa-takūna mina al-ʿārifīn li-anna Allāha qad
jaʿala kalimātihim mirʾāta anfusihim wa-yulqī ʿalā kulli wāḥidin ʿalā
qadri mā qābalū al-shamsa wa-kadhālika antum fa-ishhadūn wa-bi-hādhā
al-marātibi yaẓharu marātibuhum ʿalayhi in antum tafqahūn. Kullu dhālika
marātibu fī maqāmāti ẓuhūri al-marāyā ʿinda tajallī al-shamsi wa-lākin
ʿani al-shamsi fī hādhihi al-ayyāmi antum lā tusʾalūn.

Thus recognize the Cause of God and His way in every age, so that
matters may not become confused for you. Consider the Manifestations'
relationship to God in what is sent down to Them, not in the respects in
which They differ. Behold in the revelation of glory only the beauty of
the One who reveals it, and in what is sent down only the One who sends
it down, if you have insight. Then behold the beauty of every
Manifestation in His words, so that the truth may become clear to you
and you may be among those who recognize it. For God has made Their
words the mirror of Their own selves, and casts upon each according to
the degree to which They face the sun. Thus should you behold, and
through these degrees Their ranks in relation to Him become manifest, if
you understand. All these are degrees in the stations of the mirrors'
manifestation when the sun shines upon them; but in these days you are
not being asked about the sun.

\subsection{68}\label{section-2661}

\begin{otherlanguage}{arabic}

ثمّ اشهد بأنّها لم یزل کان واحداً فی ذاته و واحداً فی صفاته و واحداً فی
افعاله و من دون ذلک لن یثبت علیه حکم التّفرید و لن یصدق تنزیه ذاته عن
المثلیّة و التّحدید ان انتم تفقهون ایّاک ان لا تشرک باللّه و لا تدع معه الهاً
آخر و لا تتّخذ له شریکاً و لا شبیهاً و لا وزیراً و لا نظیراً و هذا حقّ الایمان
ان انتم تعلمون و انّا نهیناک عن ذکر ذلک نهیاً عظیماً و امرناک بأن تستغفر
ربّک و ترجع الیه بما اکتسبت اناملک فی ذکر هذا الأمر المستور و هذا الکنز
المخزون

\end{otherlanguage}

Thumma ishhad bi-annahā lam yazal kāna wāḥidan fī dhātihi wa-wāḥidan fī
ṣifātihi wa-wāḥidan fī afʿālihi wa-min dūni dhālika lan yathbuta ʿalayhi
ḥukmu al-tafrīdi wa-lan yaṣduqa tanzīhu dhātihi ʿani al-mithliyyati
wa-al-taḥdīd in antum tafqahūn. Iyyāka an lā tushrik bi-Allāhi wa-lā
tadʿu maʿahu ilāhan ākhara wa-lā tattakhidh lahu sharīkan wa-lā shabīhan
wa-lā wazīran wa-lā naẓīran wa-hādhā ḥaqqu al-īmāni in antum taʿlamūn
wa-innā nahaynāka ʿan dhikri dhālika nahyan ʿaẓīman wa-amarnāka bi-an
tastaghfira rabbaka wa-tarjiʿa ilayhi bi-mā iktasabat anāmiluka fī
dhikri hādhā al-amri al-mastūri wa-hādhā al-kanzi al-makhzūn.

Then bear witness that He has always been one in His essence, one in His
attributes, and one in His acts. Otherwise, absolute singleness could
not be affirmed of Him, nor could His essence truly be held exalted
above likeness and limitation, if you understand. Take care not to
associate anything with God, invoke no other god alongside Him, and
assign Him no partner, likeness, minister, or equal. This is true faith,
if you know. We have most emphatically forbidden you to speak of such
things, and have commanded you to seek your Lord's forgiveness and turn
back to Him for what your fingers have committed in speaking of this
concealed matter and this hidden treasure.

\subsection{69}\label{section-2662}

\begin{otherlanguage}{persian}

ولکن لمّا صعد هذا الطّیر فی قطب هذا السّمآء الی مقام الّذی یشهد کلّ شیء فی
ظلّه و وجدک فی ارض الحیرة لذا یلقی علیک ما یجری اللّه من قلمی علی قدر ما
کان النّاس یقدرون ان یعرفون فاعلم بأنّ اللّه ربّک ما انقطع عن شیء فیضه و
فضله من کلّ ما کان و ما یکون بل خلق کلّ شیء و قدّر فیه کلّ فضل معروف و ما لا
اطّلع به احد ان انتم تفقهون و هذا مقام استعدادات الّتی قدّرت فی کلّ شیء و
یظهر منه علی مرور الأیّام الی ان یصل الی غایة ما قدّر له کما انتم فی کلّ
شیء تشهدون

\end{otherlanguage}

Wa-lākin lammā ṣaʿida hādhā al-ṭayru fī quṭbi hādhā al-samāʾi ilā maqāmi
alladhī yashhadu kulla shayʾin fī ẓillihi wa-wajadaka fī arḍi al-ḥayrati
lidhā yulqī ʿalayka mā yujrī Allāhu min qalamī ʿalā qadri mā kāna
al-nāsu yaqdirūna an yaʿrifūn. Fa-iʿlam bi-anna Allāha rabbaka mā
inqaṭaʿa ʿan shayʾin fayḍuhu wa-faḍluhu min kulli mā kāna wa-mā yakūn
bal khalaqa kulla shayʾin wa-qaddara fīhi kulla faḍlin maʿrūfin wa-mā lā
iṭṭalaʿa bihi aḥad in antum tafqahūn wa-hādhā maqāmu istiʿdādāti allatī
quddirat fī kulli shayʾin wa-yaẓharu minhu ʿalā murūri al-ayyāmi ilā an
yaṣila ilā ghāyati mā quddira lahu kamā antum fī kulli shayʾin
tashhadūn.

But when this Bird ascended at the pole of this heaven to the station of
one who beholds all things beneath His shadow, and found you in the land
of bewilderment, He therefore imparts to you what God causes to flow
from My pen, according to what people are able to comprehend. Know that
the outpouring and grace of God, your Lord, have never been withheld
from anything that has been or will be. Rather, He created all things
and ordained within each every known bounty, as well as what no one has
discovered, if you understand. This is the station of the capacities
ordained within all things, which become manifest from them with the
passage of days until they reach the utmost of what has been ordained
for them, as you witness in all things.

\subsection{70}\label{section-2663}

\begin{otherlanguage}{arabic}

و انّک ان تقدر ان تنقطع عن کلّ ما سمعت و تصعد الی هذا المقام العظیم
المحمود فاشهد کلّ الأسمآء شموس مشرقة کما لا انتهآء للأسمآء و کذلک لا
انتهآء للشّموس مثلاً واحد من اسم اللّه فهو کریم و انّه شمس فی مقامه و یتجلّی
علی مرایا الممکنات و یظهر علی کلّ شیء انوار هذا التّجلّی علی ما کان علیه
کما انتم تشهدون و اثر هذا الشّمس اوّلاً فاشهد فی الانسان بحیث تجد فی کلّ
انسان اثر تجلّی هذا الاسم

\end{otherlanguage}

Wa-innaka in taqdir an tanqaṭiʿa ʿan kulli mā samiʿta wa-taṣʿada ilā
hādhā al-maqāmi al-ʿaẓīmi al-maḥmūd fa-ishhad kulla al-asmāʾi shumūsan
mushriqatan kamā lā intihāʾa li-al-asmāʾi wa-kadhālika lā intihāʾa
li-al-shumūs. Mathalan wāḥidun min ismi Allāhi fa-huwa Karīmun wa-innahu
shamsun fī maqāmihi wa-yatajallā ʿalā marāyā al-mumkināti wa-yaẓharu
ʿalā kulli shayʾin anwāru hādhā al-tajallī ʿalā mā kāna ʿalayhi kamā
antum tashhadūn wa-athara hādhā al-shamsi awwalan fa-ishhad fī al-insāni
bi-ḥaythu tajidu fī kulli insānin athara tajallī hādhā al-ism.

And if you can detach yourself from everything you have heard and ascend
to this great and praiseworthy station, behold all the names as shining
suns. Just as there is no end to the names, so there is no end to the
suns. For example, one of God's names is the Generous. It is a sun in
its own station and shines upon the mirrors of contingent beings. The
lights of this revelation appear upon each thing according to its
condition, as you witness. Behold the effect of this sun first in the
human being, so that in every person you may find a trace of the
revelation of this name.

\subsection{71}\label{section-2664}

\begin{otherlanguage}{arabic}

و الفرق انّ فی بعض یظهر علی حدّ الکمال و یکرم علی نفسه فی کلّ ما رزقه اللّه
و علی اقربائه و احبّائه و علی الّذینهم الیه یتقرّبون و کذلک علی کلّ من دخل
فی الایمان و من دونه علی کلّ عاقل و مجنون و فی البعض لن تجد علی هذا
الکمال بحیث یکرم علی نفسه و الّذینهم کانوا اولی النّسبة و القرابة و یمنع
عن دون هؤلآء و فی البعض تجد الاکرام علی نفسه و المنع عمّا سواه کذلک فاعرف
مراتب ظهور هذه الشّمس فی هذه الهیاکل علی ما هم یعملون

\end{otherlanguage}

Wa-al-farqu anna fī baʿḍin yaẓharu ʿalā ḥaddi al-kamāli wa-yukrimu ʿalā
nafsihi fī kulli mā razaqahu Allāhu wa-ʿalā aqribāʾihi wa-aḥibbāʾihi
wa-ʿalā alladhīnahum ilayhi yataqarrabūn wa-kadhālika ʿalā kulli man
dakhala fī al-īmāni wa-min dūnihi ʿalā kulli ʿāqilin wa-majnūn wa-fī
al-baʿḍi lan tajida ʿalā hādhā al-kamāli bi-ḥaythu yukrimu ʿalā nafsihi
wa-alladhīnahum kānū ulī al-nisbati wa-al-qarābati wa-yamnaʿu ʿan dūni
hāʾulāʾ wa-fī al-baʿḍi tajidu al-ikrāma ʿalā nafsihi wa-al-manʿa ʿammā
siwāhu kadhālika fa-iʿrif marātiba ẓuhūri hādhihi al-shamsi fī hādhihi
al-hayākili ʿalā mā hum yaʿmalūn.

The difference is that in some it appears in its fullness: a person is
generous with everything God has provided, toward himself, his
relatives, his loved ones, and those who draw close to him, and likewise
toward everyone who has entered the faith and, beyond them, toward every
person, sane or insane. In others you will not find this fullness: a
person is generous toward himself and those related to him by family and
kinship, but withholds from everyone else. In still others you find
generosity toward oneself and withholding from all others. Thus
recognize the degrees to which this sun appears in these human forms by
what they do.

\subsection{72}\label{section-2665}

\begin{otherlanguage}{arabic}

فهنیئاً لمن لن ینظر الی نفسه و الی ذی قرابته بل ینظر بطرف الفضل بحیث یکرم
علی نفسه و دونه و لا یلتفت الی جهة الّا جهة اللّه المهیمن القیّوم و یکرم
علی الّذینهم آمنوا ما اعطاه اللّه بجوده من اخلاق الحسنة و الصّفات المنیعة و
الآداب الممدوحة و الأموال المجتمعة علی نفسه و دونه لیحکی عن هذا الشّمس
بتمامها و یکون من الّذینهم انفقوا بکلّ ما عندهم فی سبل اللّه العزیز المقتدر
المحبوب

\end{otherlanguage}

Fa-hanīʾan li-man lan yanẓura ilā nafsihi wa-ilā dhī qarābatihi bal
yanẓuru bi-ṭarfi al-faḍli bi-ḥaythu yukrimu ʿalā nafsihi wa-dūnihi wa-lā
yaltafitu ilā jihatin illā jihata Allāhi al-Muhaymini al-Qayyūm
wa-yukrimu ʿalā alladhīnahum āmanū mā aʿṭāhu Allāhu bi-jūdihi min
akhlāqi al-ḥasanati wa-al-ṣifāti al-manīʿati wa-al-ādābi al-mamdūḥati
wa-al-amwāli al-mujtamiʿati ʿalā nafsihi wa-dūnihi li-yaḥkiya ʿan hādhā
al-shamsi bi-tamāmihā wa-yakūna mina alladhīnahum anfaqū bi-kulli mā
ʿindahum fī subuli Allāhi al-ʿAzīzi al-Muqtadiri al-Maḥbūb.

Blessed, then, is the one who does not look to himself and his kindred,
but looks with the eye of grace, showing generosity toward himself and
others, and turning in no direction except toward God, the
All-Protecting, the Self-Subsisting. He generously bestows upon those
who believe what God in His bounty has given him---good character,
exalted qualities, praiseworthy conduct, and the wealth accumulated for
himself and others---so that he may fully reflect this sun and be among
those who have expended everything they possess in the paths of God, the
Almighty, the All-Powerful, the Beloved.

\subsection{73}\label{section-2666}

\begin{otherlanguage}{arabic}

و انّک لو تنظر الی کلّ من فی السّموات و الأرض لن تجد شیئاً الّا و قد تجد فیه
انوار هذا الشّمس الّا من لن یقابلها و لکن انّا ما وجدنا مثل هذا ان انتم
تجدون و کذلک فاشهد تجلّی هذا الشّمس فی دون الانسان فی کلّ ما خلق ثمّ فیه
تتفکّرون مثلاً فی الشّجر کیف یکرم کلّ نفس عمّا یظهر منه من فواکه ممنوع هذا
اثر هذا التّجلّی فیه ان انتم توقنون و کذلک الی الأرض یکرم کلّ ما اخزنه اللّه
فی بطنها من کلّ شیء و من السّنبلات و انتم منها تأکلون

\end{otherlanguage}

Wa-innaka law tanẓuru ilā kulli man fī al-samāwāti wa-al-arḍi lan tajida
shayʾan illā wa-qad tajidu fīhi anwāra hādhā al-shamsi illā man lan
yuqābilahā wa-lākin innā mā wajadnā mithla hādhā in antum tajidūn
wa-kadhālika fa-ishhad tajallī hādhā al-shamsi fī dūni al-insāni fī
kulli mā khuliqa thumma fīhi tatafakkarūn. Mathalan fī al-shajari kayfa
yukrimu kulla nafsin ʿammā yaẓharu minhu min fawākihin mamnūʿ hādhā
atharu hādhā al-tajallī fīhi in antum tūqinūn wa-kadhālika ilā al-arḍi
yukrimu kulla mā akhzanahu Allāhu fī baṭnihā min kulli shayʾin wa-mina
al-sunbulāti wa-antum minhā taʾkulūn.

If you look at all who are in the heavens and on earth, you will find
nothing in which you do not find the lights of this sun, except one who
does not face it. Yet We have found no such thing, if you can find one.
Likewise, behold the revelation of this sun in everything created other
than the human being, and reflect upon it. Consider, for example, the
tree: how it generously bestows upon every soul the forbidden fruits it
produces. This is the effect of this revelation within it, if you are
certain. Likewise, look to the earth: it generously bestows everything
God has stored within it, including the ears of grain from which you
eat.

\subsection{74}\label{section-2667}

\begin{otherlanguage}{arabic}

و کذلک انت فاعرف کلّ شموس الأسمآء کما علّمناک بالحقّ لتکوننّ من الّذینهم
یعلمون و فی حکمة اللّه هم یتفکّرون ثمّ اشهد ظهور تجلّی الشّمس من اسمه المدبّر
فی کلّ شیء عمّا کان و عمّا یکون و تشهد ظهور هذا الاسم و من دونه من الأسمآء
فی الانسان اتمّها و اکملها کما انتم تعرفون و تنظرون بحیث یظهر منه کلّ
التّدابیر بأمر من لدن عزیز قیّوم

\end{otherlanguage}

Wa-kadhālika anta fa-iʿrif kulla shumūsi al-asmāʾi kamā ʿallamnāka
bi-al-ḥaqqi li-takūnanna mina alladhīnahum yaʿlamūn wa-fī ḥikmati Allāhi
hum yatafakkarūn thumma ishhad ẓuhūra tajallī al-shamsi min ismihi
al-Mudabbiri fī kulli shayʾin ʿammā kāna wa-ʿammā yakūn wa-tashhadu
ẓuhūra hādhā al-ismi wa-min dūnihi mina al-asmāʾi fī al-insāni atammahā
wa-akmalahā kamā antum taʿrifūna wa-tanẓurūn bi-ḥaythu yaẓharu minhu
kullu al-tadābīri bi-amrin min ladun ʿAzīzin Qayyūm.

In this way recognize all the suns of the names, as We have taught you
in truth, so that you may be among those who know and reflect upon the
wisdom of God. Then behold the manifestation of the sun's radiance from
His name, the Ordainer, in everything that has been and will be. You
will behold the fullest and most perfect manifestation of this name, and
of the other names, in the human being, as you recognize and observe, so
that every act of ordering proceeds from him at the command of One who
is Almighty, Self-Subsisting.

\subsection{75}\label{section-2668}

\begin{otherlanguage}{arabic}

و لم یکن هذا التّجلّی من هذا الاسم المبارک المشهود مخصوصاً بأحد دون احد کما
انتم فی غیر الحقّ من الفرق المختلفة تسمعون و تشهدون و قد بلغوا فی التّدبیر
الی مقام الّذی ما سبقهم ایدی احد و هذا من فضل اللّه علیهم و هم لا یشعرون
کذلک سبقت رحمته کلّ الموجودات و احاط فضله من اعلی الخلق الی ان تنتهی الی
جماد محدود فلمّا جهدوا هؤلآء فی ظهور هذا التّجلّی فی انفسهم و قابلوا هذا
الشّمس ما امنعهم ید الفضل من هذا المقام و هذا من فضله المبسوط الّذی احاط
کلّ شیء ولکنّ اکثر النّاس هم لا یعقلون

\end{otherlanguage}

Wa-lam yakun hādhā al-tajallī min hādhā al-ism al-mubārak al-mashhūd
makhṣūṣan bi-aḥadin dūna aḥad kamā antum fī ghayr al-Ḥaqq min al-firaq
al-mukhtalifa tasmaʿūna wa-tashhadūn wa-qad balaghū fī al-tadbīr ilā
maqām alladhī mā sabaqahum aydī aḥad wa-hādhā min faḍl Allāh ʿalayhim
wa-hum lā yashʿurūn kadhālika sabaqat raḥmatuhu kulla al-mawjūdāt
wa-aḥāṭa faḍluhu min aʿlā al-khalq ilā an tantahī ilā jamād maḥdūd
fa-lammā jahadū hāʾulāʾ fī ẓuhūr hādhā al-tajallī fī anfusihim wa-qābalū
hādhā al-shams mā amnaʿahum yad al-faḍl min hādhā al-maqām wa-hādhā min
faḍlihi al-mabsūṭ alladhī aḥāṭa kulla shayʾ wa-lākinna akthar al-nās hum
lā yaʿqilūn

This revelation of this blessed, manifest name has not been confined to
one person to the exclusion of another, as you hear and witness among
the various groups outside the Truth. They have attained a degree of
skill in ordering affairs that no one before them has achieved. This is
a token of God's grace toward them, though they are unaware. Thus His
mercy has preceded all created things, and His grace has encompassed
everything from the highest of creation down to a finite, inanimate
thing. Since these people strove to manifest this revelation within
themselves and turned toward this Sun, the hand of grace did not bar
them from this station. This is a token of His outspread grace, which
has encompassed all things, but most people do not understand.

\subsection{76}\label{section-2669}

\begin{otherlanguage}{arabic}

ولکن هؤلآء اشتغلوا بهذا الاسم فی هذا المقام و غفلوا عن دونه من اسمآء
اللّه الحسنی و مظاهرها کما انتم تشهدون و وقفوا فی ذلک الشّأن ثمّ اعتکفوا
علیه علی ما هم یریدون لا علی ما اراد اللّه لهم و لذا کتب اسمائهم من عبدة
الأسمآء بما کانوا علیها عاکفون لأنّ اللّه اراد من هذا الاسم تدبیرهم فی
دینهم فی‌الحقیقة الأوّلیّة و هم احتجبوا عن ذلک کما انتم تعلمون فلمّا ما
آمنوا بمظاهر العدل یثبت بأنّهم ما بلغوا الی حقّ التّدبیر من هذا الاسم
المبارک المحبوب

\end{otherlanguage}

Wa-lākin hāʾulāʾ ishtaghalū bi-hādhā al-ism fī hādhā al-maqām wa-ghafalū
ʿan dūnihi min asmāʾ Allāh al-ḥusnā wa-maẓāhirihā kamā antum tashhadūn
wa-waqafū fī dhālika al-shaʾn thumma iʿtakafū ʿalayhi ʿalā mā hum
yurīdūn lā ʿalā mā arāda Allāh lahum wa-li-dhā kutiba asmāʾuhum min
ʿabadat al-asmāʾ bi-mā kānū ʿalayhā ʿākifūn li-anna Allāh arāda min
hādhā al-ism tadbīrahum fī dīnihim fī al-ḥaqīqa al-awwaliyya wa-hum
iḥtajabū ʿan dhālika kamā antum taʿlamūn fa-lammā mā āmanū bi-maẓāhir
al-ʿadl yathbutu bi-annahum mā balaghū ilā ḥaqq al-tadbīr min hādhā
al-ism al-mubārak al-maḥbūb

But these people occupied themselves with this name at this station and
neglected the other most beautiful names of God and their
manifestations, as you witness. They stopped at that point and then
devoted themselves to it according to their own desires, not according
to what God desired for them. Therefore their names were recorded among
the worshippers of names, because of their devotion to them. For God's
intention in this name was that they should order their religious
affairs in the primary reality, yet they veiled themselves from this, as
you know. Since they did not believe in the manifestations of justice,
it is established that they did not attain the true ordering of affairs
that proceeds from this blessed and beloved name.

\subsection{77}\label{section-2670}

\begin{otherlanguage}{arabic}

ثمّ انظر تجلّی هذا الاسم فی التّراب بحیث کلّما یودع فی جوّه لن یخان فیه و
یدبّره فی ایّام و شهور الی ان یبلغ مقامه اذاً یخرج عنه بطراز و لون محبوب
کذلک فاشهد تجلّی هذا الشّمس من هذا الاسم فی کلّ شیء لتطّلع بشموس الأسمآء و
تکون من الّذینهم باسرار الأمر هم مطّلعون

\end{otherlanguage}

Thumma unẓur tajallī hādhā al-ism fī al-turāb bi-ḥaythu kullamā yūdaʿu
fī jawwihi lan yukhāna fīhi wa-yudabbiruhu fī ayyām wa-shuhūr ilā an
yablugha maqāmahu idhan yakhruju ʿanhu bi-ṭirāz wa-lawn maḥbūb kadhālika
fa-shhad tajallī hādhā al-shams min hādhā al-ism fī kull shayʾ
li-taṭṭaliʿa bi-shumūs al-asmāʾ wa-takūna min alladhīna hum bi-asrār
al-amr hum muṭṭaliʿūn

Then consider the revelation of this name in the soil: whatever is
entrusted to its depths is not betrayed there. It tends it over days and
months until it reaches its proper stage, whereupon it emerges from the
soil in a pleasing form and color. Likewise, behold the revelation of
this Sun through this name in everything, so that you may become
acquainted with the suns of the names and be among those who are
acquainted with the mysteries of the Cause.

\subsection{78}\label{section-2671}

\begin{otherlanguage}{arabic}

و انّی لو اذکر شموس الأسمآء و تجلّیاتها علی ما قدّر اللّه لها الی آخر الّذی
لا آخر له لن ینتهی ان انتم تعلمون و اکتفینا بما بیّنّاه لک و انت فاکف بما
قدّر لک من هذا القلم المکنون و فتحنا علی وجهک ابواب العلم و الحکمة فیما
القیناک بالحقّ لتخرج من هذا البحر ما ترید من لآلئ علم مکتوم و لو ترید ان
تخرج من کلّ باب مائة الف باب لتقدر ان تقابل هذا الاسم المستور المخزون
الّذی ابتلی بین یدی الّذینهم ما آمنوا باللّه طرفة عین ولو انّهم فی انفسهم
جوهر الایمان یدّعون

\end{otherlanguage}

Wa-innī law adhkuru shumūs al-asmāʾ wa-tajalliyātihā ʿalā mā qaddara
Allāh lahā ilā ākhir alladhī lā ākhira lahu lan yantahī in antum
taʿlamūn wa-ktafaynā bi-mā bayyannāhu laka wa-anta fa-kfi bi-mā quddira
laka min hādhā al-qalam al-maknūn wa-fataḥnā ʿalā wajhika abwāb al-ʿilm
wa-al-ḥikma fī-mā alqaynāka bi-al-ḥaqq li-tukhrija min hādhā al-baḥr mā
turīdu min laʾāliʾ ʿilm maktūm wa-law turīdu an tukhrija min kull bāb
miʾata alf bāb la-taqdiru in tuqābila hādhā al-ism al-mastūr al-makhzūn
alladhī ubtuliya bayna yaday alladhīna hum mā āmanū bi-Allāh ṭarfata
ʿayn wa-law annahum fī anfusihim jawhar al-īmān yaddaʿūn

Were I to recount the suns of the names and their revelations as God has
ordained for them, even to the end that has no end, the account would
never be exhausted, if you but knew. We have contented Ourselves with
what We have explained to you; you too should be content with what has
been apportioned to you by this hidden Pen. We have opened before your
face the doors of knowledge and wisdom through what We have imparted to
you in truth, so that you may draw from this ocean whatever pearls of
concealed knowledge you desire. Should you wish to bring forth a hundred
thousand doors from every door, you will be able to do so if you turn
toward this veiled and treasured Name, who has been afflicted at the
hands of those who have not believed in God for the blink of an eye,
though they claim for themselves the very essence of faith.

\subsection{79}\label{section-2672}

\begin{otherlanguage}{arabic}

قل انّ العلم یظهر من عنده و انّه لهو الحقّ علّام العلوم و لو احد یأتی بعلوم
الأوّلین و الآخرین و لن یدخل فی هذا الظلّ لن یقبل عنه شیء ان انتم توقنون
کما ان علمآء الانجیل حین الّذی ظهر محمّد بالحقّ ما قبل عنهم من شیء و کذلک
فی البیان فانظرون اذا جآء علیّ بالحقّ ما نفع علمآء الفرقان علمهم الّا لمن
دخل فی ظلّه ان انتم تشعرون و کذلک فاعرف من اوّل الّذی لا اوّل له الی آخر
الّذی لا آخر له و کن من الّذینهم ببصر اللّه فی کلّ شیء یتفرّسون و ان یأتی احد
و لم یکن عنده حرف من العلم و یدخل فی هذا الظّلّ انّه اعلم من کلّ عالم ذی
فنون لأنّ العلم الّذی لم یزل کان مذکوراً عند اللّه هو عرفانه و عرفان مظاهر
امره ان انتم الی کتب اللّه تنظرون

\end{otherlanguage}

Qul inna al-ʿilm yaẓharu min ʿindihi wa-innahu la-huwa al-Ḥaqq ʿallām
al-ʿulūm wa-law aḥad yaʾtī bi-ʿulūm al-awwalīn wa-al-ākhirīn wa-lan
yadkhula fī hādhā al-ẓill lan yuqbala ʿanhu shayʾ in antum tūqinūn kamā
anna ʿulamāʾ al-Injīl ḥīna alladhī ẓahara Muḥammad bi-al-ḥaqq mā qubila
ʿanhum min shayʾ wa-kadhālika fī al-Bayān fa-nẓurūn idhā jāʾa ʿAlī
bi-al-ḥaqq mā nafaʿa ʿulamāʾ al-Furqān ʿilmuhum illā li-man dakhala fī
ẓillihi in antum tashʿurūn wa-kadhālika fa-ʿrif min awwal alladhī lā
awwala lahu ilā ākhir alladhī lā ākhira lahu wa-kun min alladhīna hum
bi-baṣar Allāh fī kull shayʾ yatafarrasūn wa-in yaʾti aḥad wa-lam yakun
ʿindahu ḥarf min al-ʿilm wa-yadkhul fī hādhā al-ẓill innahu aʿlam min
kull ʿālim dhī funūn li-anna al-ʿilm alladhī lam yazal kāna madhkūran
ʿinda Allāh huwa ʿirfānuhu wa-ʿirfān maẓāhir amrihi in antum ilā kutub
Allāh tanẓurūn

Say: Knowledge proceeds from Him, and He is indeed the Truth, the Knower
of all knowledge. Even if someone were to bring the knowledge of the
first and the last, yet did not enter this shade, nothing would be
accepted from him, if you are certain. Thus, when Muḥammad appeared in
truth, nothing was accepted from the scholars of the Gospel. Consider
likewise what occurred in the Bayān: when ʿAlī came in truth, the
knowledge of the scholars of the Qur'an did not benefit them, except
those who entered His shade, if you perceive. Understand it thus from
the beginning that has no beginning to the end that has no end, and be
among those who discern all things with the sight of God. If someone
comes who possesses not a single letter of knowledge and enters this
shade, he is more learned than every scholar versed in many disciplines.
For the knowledge that has always been accounted knowledge in God's
sight is recognition of Him and recognition of the Manifestations of His
Cause, if you look into the Books of God.

\subsection{80}\label{section-2673}

\begin{otherlanguage}{arabic}

و تجدون کلّ ذلک فی البیان ان انتم تقرؤون و لن یحرّک قلمی الّا علی ما حرّک
علیه قلم اللّه المهیمن القیّوم قل یا ملأ البیان لا تأخذونی بما نزل علیّ
الآیات فواللّه هذا الذّنب لم یکن منّی ولکن الرّوح یقومنی فی کلّ حین ان انتم
تسمعون و انّکم لو تقدّسون نفوسکم لتجدون اریاح اللّه تسطع من هذا المسک
المهبوب ولکن لمّا منعتم انفسکم عن هذا الرّضوان لن تجدوا من شیء کما لم
یجدوا امم الفرقان رایحة اللّه المهیمن القیّوم

\end{otherlanguage}

Wa-tajidūna kulla dhālika fī al-Bayān in antum taqraʾūn wa-lan yuḥarraka
qalamī illā ʿalā mā ḥurrika ʿalayhi qalam Allāh al-Muhaymin al-Qayyūm
qul yā malaʾ al-Bayān lā taʾkhudhūnī bi-mā nazala ʿalayya al-āyāt
fa-wa-Allāh hādhā al-dhanb lam yakun minnī wa-lākin al-Rūḥ yuqawwimunī
fī kull ḥīn in antum tasmaʿūn wa-innakum law tuqaddisūna nufūsakum
la-tajidūna aryāḥ Allāh tasṭaʿu min hādhā al-misk al-mahbūb wa-lākin
lammā manaʿtum anfusakum ʿan hādhā al-riḍwān lan tajidū min shayʾ kamā
lam yajidū umam al-Furqān rāʾiḥat Allāh al-Muhaymin al-Qayyūm

You will find all this in the Bayān, if you read it. My pen will move
only as the Pen of God, the All-Protecting, the Self-Subsisting, has
moved. Say: O people of the Bayān, do not hold Me to account because
verses have descended upon Me. By God, this fault is not Mine; rather,
the Spirit raises Me up at every moment, if you listen. Were you to
purify yourselves, you would find the fragrances of God wafting from
this diffusing musk. But since you have deprived yourselves of this
paradise, you will find nothing, just as the peoples of the Qur'an did
not find the fragrance of God, the All-Protecting, the Self-Subsisting.

\subsection{81}\label{section-2674}

\begin{otherlanguage}{arabic}

و انّک انت یا ایّها السّائل تفکّر فیما القیناک و خذ لنفسک ما ترید و انّک انت
ان اردت من الشّمس هذه الشّموس المشرقة فقد اظهرنا لک اسمائها و اسرارها و
اریناک تجلّیاتها علی مرایا الموجودات و اشراقها علی ما کان و ما یکون و ان
اردت شمس الشّموس و سلطان الوجود و ملیک المقصود الّذی یطوفنّ فی حوله هذه
الشّموس و حقایقها قل سبحان ذی العظمة الّذی خضع له کلّ ذی عظمة و اقتدار
فسبحان ذی القدرة الّذی بخع له کلّ ذی قدرة و ارتفاع فسبحان ذی العزّة الّذی ذلّ
له کلّ ذی عزّة و امتناع فسبحان ذی القوّة الّذی قد خفض له کلّ ذی قوّة و ذی بهآء

\end{otherlanguage}

Wa-innaka anta yā ayyuhā al-sāʾil tafakkar fī-mā alqaynāka wa-khudh
li-nafsika mā turīdu wa-innaka anta in aradta min al-shams hādhihi
al-shumūs al-mushriqa fa-qad aẓharnā laka asmāʾahā wa-asrārahā
wa-araynāka tajalliyātihā ʿalā marāyā al-mawjūdāt wa-ishrāqahā ʿalā mā
kāna wa-mā yakūn wa-in aradta shams al-shumūs wa-sulṭān al-wujūd
wa-malīk al-maqṣūd alladhī yaṭūfanna fī ḥawlihi hādhihi al-shumūs
wa-ḥaqāʾiquhā qul subḥāna dhī al-ʿaẓama alladhī khaḍaʿa lahu kull dhī
ʿaẓama wa-iqtidār fa-subḥāna dhī al-qudra alladhī bakhaʿa lahu kull dhī
qudra wa-irtifāʿ fa-subḥāna dhī al-ʿizza alladhī dhalla lahu kull dhī
ʿizza wa-imtināʿ fa-subḥāna dhī al-quwwa alladhī qad khafaḍa lahu kull
dhī quwwa wa-dhī bahāʾ

As for you, O questioner, reflect on what We have imparted to you and
take for yourself whatever you desire. If by the sun you meant these
shining suns, We have disclosed their names and mysteries to you and
shown you their revelations in the mirrors of created things and their
radiance upon all that has been and will be. But if you meant the Sun of
suns, the Sovereign of existence and the King who is the object of all
desire, around whom these suns and their realities revolve, say: Glory
be to the Lord of grandeur, before whom all who possess grandeur and
power bow down! Glory be to the Lord of might, before whom all who
possess might and eminence submit! Glory be to the Lord of glory, before
whom all who possess glory and invulnerability humble themselves! Glory
be to the Lord of strength, before whom all who possess strength and
splendor lower themselves!

\subsection{82}\label{section-2675}

\begin{otherlanguage}{arabic}

ثمّ اعلم بأن دونه لن یبلغ الیه و ما سواه معدوم لدیه و انّه لهو المقتدر
العزیز المحبوب و کلّ الأسمآء یبعث بقوله و کلّ الصّفات یظهر بأمره ان انتم
تعلمون ایّاک ان لا تنسب هذا المقام بأحد دونه و لا تذکر عند ذکره احد و لا
تقترن به اشارة و لا عبارة لأنّه جلّ و عزّ لن یشیر باشارة غیره و لن یقترن
بعرفان خلقه ان انتم تعرفون و انّ اعلی عرفان العارفین و بلوغ البالغین یرجع
الی خلقه الّذی خلق بقوله کن فیکون و کلّ ما سواه مخلوق بأمره و منجعل
بارادته و ساجد لطلعته و خاضع لحضرته و کلّ عنده فی لوح محفوظ

\end{otherlanguage}

Thumma iʿlam bi-anna dūnahu lan yablugha ilayhi wa-mā siwāhu maʿdūm
ladayhi wa-innahu la-huwa al-Muqtadir al-ʿAzīz al-Maḥbūb wa-kull
al-asmāʾ yubʿathu bi-qawlihi wa-kull al-ṣifāt yaẓharu bi-amrihi in antum
taʿlamūn iyyāka an lā tansuba hādhā al-maqām bi-aḥad dūnahu wa-lā
tadhkur ʿinda dhikrihi aḥad wa-lā taqtarin bihi ishāra wa-lā ʿibāra
li-annahu jalla wa-ʿazza lan yushīra bi-ishārat ghayrihi wa-lan
yaqtarina bi-ʿirfān khalqihi in antum taʿrifūn wa-inna aʿlā ʿirfān
al-ʿārifīn wa-bulūgh al-bālighīn yarjiʿu ilā khalqihi alladhī khuliqa
bi-qawlihi kun fa-yakūn wa-kull mā siwāhu makhlūq bi-amrihi wa-munjaʿil
bi-irādatihi wa-sājid li-ṭalʿatihi wa-khāḍiʿ li-ḥaḍratihi wa-kull
ʿindahu fī lawḥ maḥfūẓ

Then know that none other than Him can attain to Him, and all besides
Him is nonexistent in His presence. He is indeed the All-Powerful, the
Almighty, the Beloved. All names are called into being by His word, and
all attributes appear at His command, if you know. Beware of ascribing
this station to anyone other than Him. Mention no one when mentioning
Him, and associate with Him neither allusion nor expression. For He,
exalted and mighty is He, cannot be indicated by the allusion of
another, nor can He be encompassed by the understanding of His
creatures, if you understand. The highest understanding of those who
know and the attainment of those who attain refer only to His creation,
brought into being by His saying, ``Be,'' and it is. Everything besides
Him is created at His command, fashioned by His will, prostrate before
His countenance and submissive in His presence. All is with Him in a
preserved Tablet.

\subsection{83}\label{section-2676}

\begin{otherlanguage}{arabic}

و انّه یظهر بالحقّ کما وعدکم اللّه فی الکتاب اتّق اللّه ثمّ فی امره لا تتکلّمون
و انتظر یومه کما ینتظر ایّامه کلّ من فی السّموات و الأرض و ان کانوا فی
انفسهم لا یشعرون ثمّ اعلم بأنّی آمنت به قبل ظهوره و بعد ظهوره و حین ظهوره
و انتظر ایّامه بما رقم علی الواح اللّه المهیمن القیّوم ایّاک ان لا تذکر هذا
الاسم و لا تظنّ ظنّ السّوء فی امره و کن من الّذینهم فی کلّ یوم یستغفرون

\end{otherlanguage}

Wa-innahu yaẓharu bi-al-ḥaqq kamā waʿadakum Allāh fī al-kitāb ittaqi
Allāh thumma fī amrihi lā tatakallamūn wa-ntaẓir yawmahu kamā yantaẓiru
ayyāmahu kull man fī al-samāwāt wa-al-arḍ wa-in kānū fī anfusihim lā
yashʿurūn thumma iʿlam bi-annī āmantu bihi qabla ẓuhūrihi wa-baʿda
ẓuhūrihi wa-ḥīna ẓuhūrihi wa-ntaẓir ayyāmahu bi-mā ruqima ʿalā alwāḥ
Allāh al-Muhaymin al-Qayyūm iyyāka an lā tadhkura hādhā al-ism wa-lā
taẓunna ẓann al-sūʾ fī amrihi wa-kun min alladhīna hum fī kull yawm
yastaghfirūn

He will appear in truth, as God has promised you in the Book. Fear God,
and do not speak concerning His Cause. Await His day, just as all who
are in the heavens and on earth await His days, though they themselves
are unaware. Then know that I have believed in Him before His
appearance, after His appearance, and at the time of His appearance.
Await His days in accordance with what has been inscribed upon the
Tablets of God, the All-Protecting, the Self-Subsisting. Beware of
mentioning this name, and entertain no evil thought concerning His
Cause. Be among those who seek forgiveness every day.

\subsection{84}\label{section-2677}

\begin{otherlanguage}{arabic}

و انّا نهیناک عن ذکر هذا الاسم لئلّا یتوهّم احد فی نفسه التّقرّب الی هذه
الشّجرة المکنون و لئلّا تسئل عن احد مثل ما سئلت و لا تذکر احد بهذا الاسم
لأنّ هذا مخصوص به ان انتم تسمعون ما یعلّمکم هذه الحمامة بما القی الرّوح
علیه من کلمات اللّه العزیز القیّوم یا قوم لا تفرطوا فی الکتاب و لا تحرفوا
کلمة اللّه و لا تکوننّ من الّذینهم یتکلّمون فیما لا یعرفون

\end{otherlanguage}

Wa-innā nahaynāka ʿan dhikr hādhā al-ism li-allā yatawahhama aḥad fī
nafsihi al-taqarrub ilā hādhihi al-shajara al-maknūn wa-li-allā tasʾala
ʿan aḥad mithla mā saʾalta wa-lā tadhkur aḥad bi-hādhā al-ism li-anna
hādhā makhṣūṣ bihi in antum tasmaʿūna mā yuʿallimukum hādhihi al-ḥamāma
bi-mā alqā al-Rūḥ ʿalayhi min kalimāt Allāh al-ʿAzīz al-Qayyūm yā qawm
lā tufarriṭū fī al-kitāb wa-lā tuḥarrifū kalimat Allāh wa-lā takūnanna
min alladhīna hum yatakallamūna fī-mā lā yaʿrifūn

We have forbidden you to mention this name lest anyone imagine himself
to have drawn near to this hidden Tree, and lest you ask about anyone as
you have asked. Do not designate anyone by this name, for it belongs
exclusively to Him, if you listen to what this Dove teaches you through
the words of God, the Almighty, the Self-Subsisting, which the Spirit
has imparted to it. O people, do not be remiss concerning the Book, do
not distort the word of God, and do not be among those who speak of what
they do not know.

\subsection{85}\label{section-2678}

\begin{otherlanguage}{arabic}

فواللّه ترب قدمه لیکون اعلی عن کلّ ما خلق و یخلق و عن کلّ ما کان و ما یکون
و ان عرفت بأنّ النّفوس لن یضطرب لقلت انّه عرش اللّه فی الأرض لأنّ علیه یقع
رجله العزیز المحبوب و انّک فاثبت فیما القیناک بحیث لن یزلّ قدماک و لا
تکونن من الّذینهم لا یعرفون و ینسبون هذا المقام الی الّذینهم یریدون فواللّه
هذا خطآء لم یکن فی الامکان کشبهه اتّقوا اللّه یا قوم و لا تتجاوزوا عمّا حدّد
فی کتاب اللّه المهیمن القیّوم

\end{otherlanguage}

Fa-wa-Allāh turbu qadamihi la-yakūnu aʿlā ʿan kull mā khuliqa
wa-yukhlaqu wa-ʿan kull mā kāna wa-mā yakūn wa-in ʿaraftu bi-anna
al-nufūs lan yaḍṭariba la-qultu innahu ʿarsh Allāh fī al-arḍ li-anna
ʿalayhi yaqaʿu rijluhu al-ʿazīz al-maḥbūb wa-innaka fa-thbut fī-mā
alqaynāka bi-ḥaythu lan yazilla qadamāka wa-lā takūnanna min alladhīna
hum lā yaʿrifūna wa-yansibūna hādhā al-maqām ilā alladhīna hum yurīdūn
fa-wa-Allāh hādhā khaṭaʾ lam yakun fī al-imkān ka-shibhihi ittaqū Allāh
yā qawm wa-lā tatajāwazū ʿammā ḥuddida fī kitāb Allāh al-Muhaymin
al-Qayyūm

By God, the dust beneath His feet is more exalted than all that has been
or will be created, and than all that has been and will be. Were I
assured that souls would not be thrown into turmoil, I would say that it
is the throne of God on earth, for upon it rests His mighty and beloved
foot. Remain so steadfast in what We have imparted to you that your feet
will never slip. Do not be among those who do not understand and ascribe
this station to whomever they wish. By God, this is an error without
parallel in the world of being. Fear God, O people, and do not
transgress the bounds set in the Book of God, the All-Protecting, the
Self-Subsisting.

\subsection{86}\label{section-2679}

\begin{otherlanguage}{arabic}

فسبحانک اللّهمّ یا الهی انت الّذی الهمتنی بدایع ذکرک ثمّ منعتنی عن اظهارها
بین بریّتک و اخذتنی علی مقام الّذی صرت ممنوعاً عن اظهار جوامع اسرار حکمتک
بما ظهرت آثار الغفلة فی خلقک لأنّی یا الهی ان اکتم عن النّاس ما اعطیتنی
بجودک لأکون فی تلک الحالة کافراً بنعمتک و ان احدث لهم ما اکرمتنی بفضلک
اذاً یرفع ضجیج ارقّائک و صریخ عبادک و لم ار کلّ ذلک الّا منک لأنّک بعثتنی و
اظهرتنی فی ایّام الّتی ما بعثت فی مثلها احد من صفوتک و لم ادر بذلک ما اردت
من ظهورات قضائک و تقدیرک و شئونات امضائک و تدبیرک

\end{otherlanguage}

Fa-subḥānaka Allāhumma yā ilāhī anta alladhī alhamtanī badāʾiʿa dhikrika
thumma manaʿtanī ʿan iẓhārihā bayna bariyyatika wa-akhadhtanī ʿalā maqām
alladhī ṣirtu mamnūʿan ʿan iẓhār jawāmiʿ asrār ḥikmatika bi-mā ẓaharat
āthār al-ghafla fī khalqika li-annī yā ilāhī in aktum ʿan al-nās mā
aʿṭaytanī bi-jūdika la-akūna fī tilka al-ḥāla kāfiran bi-niʿmatika wa-in
uḥaddith lahum mā akramtanī bi-faḍlika idhan yurfaʿu ḍajīj ariqqāʾika
wa-ṣarīkh ʿibādika wa-lam ara kulla dhālika illā minka li-annaka
baʿathtanī wa-aẓhartanī fī ayyām allatī mā baʿathta fī mithlihā aḥad min
ṣafwatika wa-lam adri bi-dhālika mā aradta min ẓuhūrāt qaḍāʾika
wa-taqdīrika wa-shuʾūnāt imḍāʾika wa-tadbīrika

Glory be to You, O God, my God! You are the One who inspired Me with the
wondrous expressions of Your remembrance, then prevented Me from
disclosing them among Your creatures. You have brought Me to a station
where I am prevented from disclosing the all-embracing mysteries of Your
wisdom because the signs of heedlessness have appeared among Your
creation. For, O my God, if I conceal from the people what You have
bestowed upon Me through Your generosity, I would then be ungrateful for
Your favor; and if I tell them of what You have graciously conferred
upon Me, the clamor of Your servants and the outcry of Your worshippers
will rise. Yet I see all this as coming only from You, for You have
raised Me up and made Me manifest in days the like of which You have
never chosen for raising up any of Your chosen ones. I do not know what
You have intended thereby through the manifestations of Your decree and
ordainment, and the workings of Your enactment and ordering.

\subsection{87}\label{section-2680}

\begin{otherlanguage}{persian}

فوعزّتک یا محبوبی اشهد بأنّ علمک ما احاط یوماً اعظم و اشدّ من هذه الأیّام لأنّ
فیها یمحّص کلّ العباد من کلّ وضیع و شریف ولکن انّی مع عجزی و فقری و ضرّی و
مسکنتی لأکون مستقیماً علی امرک بحولک و قوّتک و لو لن یقبل احد بهذا الوجه
طلباً للقائک و اشکرک بما اعطیتنی و اکرمتنی من استقامة الّتی لن یقدر احد ان
یردّها عن مقامها ولو یرفع علی رأسی سیوف المنکرین و اعراض المشرکین و اشهد
بأنّ ذرّة من هذه الاستقامة لخیر من عبادة الثّقلین و بذلک یستریح قلبی و تسکن
نفسی و لن اشاهد کلّ من فی السّموات و الأرض الّا ککفّ من الطّین بل احقر من ذلک
الّا من دخل فی ظلّک فی هذا الاسم الأتمّ الأعظم الّذی به انقلبت کلّ الأسمآء و
تغیّرت کلّ الصّفات فی جبروت البدآء لأنّ الّذینهم لن یعرفوک لم یکن لهم شأن
عندک و لا لهم ذکر بین یدیک اذاً اسئلک یا محبوبی بأن لا تزلّ قدمی عن امرک
الّذی لن یقوم معه شیء ثمّ ثبّتنی بجودک و فضلک ثمّ انقطعنی عن دونک و انّک انت
العزیز القادر العلیم الحکیم القیّوم

\end{otherlanguage}

Fa-wa-ʿizzatika yā maḥbūbī ashhadu bi-anna ʿilmaka mā aḥāṭa yawman
aʿẓama wa-ashadda min hādhihi al-ayyām li-anna fīhā yumaḥḥaṣu kullu
al-ʿibād min kulli waḍīʿin wa-sharīf wa-lākin innanī maʿa ʿajzī wa-faqrī
wa-ḍurrī wa-maskanatī la-akūnu mustaqīman ʿalā amrika bi-ḥawlika
wa-quwwatika wa-law lan yuqbila aḥadun bi-hādhā al-wajh ṭalaban
li-liqāʾika wa-ashkuruka bi-mā aʿṭaytanī wa-akramtanī min istiqāmati
allatī lan yaqdira aḥadun an yaruddahā ʿan maqāmihā wa-law yurfaʿu ʿalā
raʾsī suyūfu al-munkirīn wa-iʿrāḍu al-mushrikīn wa-ashhadu bi-anna
dharratan min hādhihi al-istiqāma la-khayrun min ʿibādati al-thaqalayn
wa-bi-dhālika yastarīḥu qalbī wa-taskunu nafsī wa-lan ushāhida kulla man
fī al-samāwāti wa-al-arḍ illā ka-kaffin min al-ṭīn bal aḥqara min
dhālika illā man dakhala fī ẓillika fī hādhā al-ismi al-atammi al-aʿẓam
alladhī bihi inqalabat kullu al-asmāʾ wa-taghayyarat kullu al-ṣifāt fī
jabarūti al-badāʾ li-anna alladhīna hum lan yaʿrifūka lam yakun lahum
shaʾnun ʿindaka wa-lā lahum dhikrun bayna yadayka idhan asʾaluka yā
maḥbūbī bi-an lā tazilla qadamī ʿan amrika alladhī lan yaqūma maʿahu
shayʾ thumma thabbitnī bi-jūdika wa-faḍlika thumma inqaṭiʿnī ʿan dūnika
wa-innaka anta al-ʿAzīz al-Qādir al-ʿAlīm al-Ḥakīm al-Qayyūm.

By Your might, O my Beloved, I testify that Your knowledge has never
encompassed a day greater or more severe than these days, for in them
all Your servants, both lowly and eminent, are put to the test. Yet
despite my helplessness, poverty, affliction, and destitution, I will
remain steadfast in Your Cause through Your power and strength, even if
no one turns toward this Face seeking to meet You. I thank You for the
steadfastness You have bestowed upon me and with which You have honored
me, which no one can dislodge from its station, even if the swords of
the deniers and the rejection of those who associate partners with You
are raised over my head. I testify that a particle of this steadfastness
is better than the worship of both humankind and jinn. Through it my
heart finds rest and my soul becomes calm, and I regard all who are in
the heavens and on earth as no more than a handful of clay---indeed, as
less than that---except those who have entered beneath Your shadow in
this most perfect, most great Name, through which all names have been
transformed and all attributes changed in the dominion of the revision
of divine decree. For those who do not recognize You have no standing
with You, nor are they mentioned in Your presence. I therefore ask You,
O my Beloved, not to let my foot slip from Your Cause, alongside which
nothing can stand. Make me steadfast through Your generosity and grace,
and detach me from all besides You. You are indeed the Almighty, the
All-Powerful, the All-Knowing, the All-Wise, the Self-Subsisting.

\subsection{88}\label{section-2681}

\begin{otherlanguage}{arabic}

و انّا کنّا ان نکتب هذه الألواح بالفارسی و امسک اللّه القلم عنه و ارجعه الی
لسان عربیّ مبین لذا قد القیناک من اسرار العلم و الحکمة بهذا اللّسان ثمّ
نرجع الی لسان الأعجمین لیتمّ علیک حجّة اللّه و برهانه و دللیه و آیاته لعلّ
تکوننّ من المقبلین الی اللّه الّذی خلقنی و ایّاکم و نحن له لمن العابدین و
الحمد للّه ربّ العالمین

\end{otherlanguage}

Wa-innā kunnā an naktuba hādhihi al-alwāḥ bi-al-fārisī wa-amsaka Allāhu
al-qalama ʿanhu wa-arjaʿahu ilā lisānin ʿarabiyyin mubīn li-dhā qad
alqaynāka min asrāri al-ʿilmi wa-al-ḥikma bi-hādhā al-lisān thumma
narjiʿu ilā lisāni al-aʿjamīn li-yatimma ʿalayka ḥujjatu Allāhi
wa-burhānuhu wa-{[}corrupted text{]} wa-āyātuhu laʿalla takūnanna min
al-muqbilīn ilā Allāhi alladhī khalaqanī wa-iyyākum wa-naḥnu lahu la-min
al-ʿābidīn wa-al-ḥamdu li-llāhi rabbi al-ʿālamīn.

We had intended to write these Tablets in Persian, but God held the Pen
back from it and returned it to a clear Arabic tongue. Thus We have
imparted to you some of the mysteries of knowledge and wisdom in this
tongue. Now We return to the tongue of the Persians, so that God's
testimony, His proof, {[}corrupted original text{]}, and His signs may
be complete for you, that you may be among those who turn toward God,
who created Me and you. We are indeed among His worshippers. Praise be
to God, the Lord of all worlds.

\subsection{89}\label{section-2682}

\begin{otherlanguage}{persian}

و از این بیانات مشرقه از افق تبیان معلوم و مبرهن شده که اطلاق اسم شمس بر
هر یک از مظاهر این اسماء مشرقهٴ مذکوره من اولی الحکمة شده و میشود پس اگر
در بیانات مظاهر احدیّه و مطالع قدس صمدیّه ذکر شمس ملحوظ و مذکور شود مقصود
این شموس مشرقه و انوار منیره خواهند بود الّا من شآء ربّک

\end{otherlanguage}

Va az īn bayānāt-e moshreqe az ofoq-e tebyān maʿlūm va mobarhan shode ke
eṭlāq-e esm-e shams bar har yek az maẓāher-e īn asmā-ye moshreqe-ye
madhkūre min ulī al-ḥikma shode va mīshavad pas agar dar bayānāt-e
maẓāher-e aḥadīye va maṭāleʿ-e qods-e ṣamadīye dhekr-e shams malḥūẓ va
madhkūr shavad maqṣūd īn shomūs-e moshreqe va anvār-e monīre khāhand būd
illā man shāʾa rabbuka.

From these utterances shining from the horizon of exposition, it has
become clear and demonstrable that those endowed with wisdom have
applied, and continue to apply, the name ``sun'' to each of the
manifestations of these radiant names already mentioned. Thus, if
mention of the sun is found in the utterances of the Manifestations of
divine unity and the Daysprings of eternal holiness, what is intended
will be these shining suns and radiant lights, except as your Lord may
will.

\subsection{90}\label{section-2683}

\begin{otherlanguage}{persian}

و دیگر آنکه بر صاحبان علم و حکمت و فطرت سلیم و طینت منیر معلوم و مبرهن
بوده که در هر عهد و عصریکه کشف سبحات جلال از وجه جمال شد و طلعت مستورهٴ
مقنّعهٴ مغطئه از حجب غیب ظاهر و مشهود گشت آن شجرهٴ حقیقت ربّانیّه و ثمرهٴ
جنیّهٴ الهیّه را بخود او باید ادراک نمود چه که دون او در ظلّ او بوده و
خواهد بود و مقامات مستورهٴ آن کنز احدیّه از بیانات مشرقهٴ او ادراک میشود

\end{otherlanguage}

Va dīgar ānke bar ṣāḥebān-e ʿelm va ḥekmat va feṭrat-e salīm va ṭīnat-e
monīr maʿlūm va mobarhan būde ke dar har ʿahd va ʿaṣrī ke kashf-e
soboḥāt-e jalāl az vajh-e jamāl shod va ṭalʿat-e mastūre-ye moqannaʿe-ye
moghaṭṭaʾe az ḥojob-e gheyb ẓāher va mashhūd gasht ān shajare-ye
ḥaqīqat-e rabbānīye va thamare-ye janīye-ye elāhīye rā be-khod-e ū bāyad
edrāk namūd che ke dūn-e ū dar ẓell-e ū būde va khāhad būd va maqāmāt-e
mastūre-ye ān kanz-e aḥadīye az bayānāt-e moshreqe-ye ū edrāk mīshavad.

Furthermore, it has been clear and demonstrable to those possessed of
knowledge and wisdom, sound dispositions and illumined natures, that in
every age and era when the veils of glory are lifted from the face of
beauty, and that concealed, veiled, and hidden Countenance emerges
visibly from the veils of the unseen, that Tree of divine Reality and
ripe Fruit of God must be recognized through Himself. For all besides
Him has been, and will remain, beneath His shadow, and the hidden
stations of that Treasure of unity are apprehended through His radiant
utterances.

\subsection{91}\label{section-2684}

\begin{otherlanguage}{persian}

اگرچه بصر حدید مقدّس از مشاهده فطور و تحدید از جمیع ما یظهر منه آثار عظمة
اللّه مشاهده مینماید چه که تجلّیات آن شمس مشرقهٴ محیطه اظهر از آن است که
از انظار عارفین مستور ماند و تموّجات آن بحر اعظم ابین از آن است که از
ابصار مقرّبین محجوب گردد ولکن چون ناس متفاوتند در معارج علم و ایقان لهذا
محض ظهور فضل و جود و بروز فیض و عنایت بر کلّ وجود آیات نازله از سماء مشیّت
و بیّنات مشرقه از افق فضل را حجّت خود قرار فرموده‌اند که تا جمیع ناس از
اعالی و ادانی از این حجّت لایحهٴ بالغه و آیات مشرقهٴ درّیّه آن ذات قدم و
کلمهٴ اتمّ را ادراک نمایند

\end{otherlanguage}

Agarche baṣar-e ḥadīd-e moqaddas az moshāhede-ye foṭūr va taḥdīd az
jamīʿ-e mā yaẓharu minhu āthār-e ʿaẓamat Allāh moshāhede mīnamāyad che
ke tajallīyāt-e ān shams-e moshreqe-ye moḥīṭe aẓhar az ān ast ke az
anẓār-e ʿārefīn mastūr mānad va tamavvojāt-e ān baḥr-e aʿẓam abyan az ān
ast ke az abṣār-e moqarrabīn maḥjūb gardad va-lāken chūn nās
motafāvetand dar maʿārej-e ʿelm va īqān le-hādhā maḥḍ-e ẓohūr-e faḍl va
jūd va borūz-e feyḍ va ʿenāyat bar koll-e vojūd āyāt-e nāzele az samā-ye
mashīyat va bayyenāt-e moshreqe az ofoq-e faḍl rā ḥojjat-e khod qarār
farmūde-and ke tā jamīʿ-e nās az aʿālī va adānī az īn ḥojjat-e lāʾeḥe-ye
bāleghe va āyāt-e moshreqe-ye dorrīye ān dhāt-e qedam va kaleme-ye atamm
rā edrāk namāyand.

A penetrating sight, sanctified from perceiving flaw and limitation,
does indeed behold the traces of God's greatness in all that proceeds
from Him. For the splendors of that radiant, all-encompassing Sun are
too manifest to remain hidden from the eyes of those who possess true
knowledge, and the surging of that Most Great Ocean is too evident to be
concealed from the sight of those brought near. Yet because people
differ in their degrees of knowledge and certitude, solely to manifest
grace and generosity and to extend bounty and favor to all existence, He
has appointed the verses sent down from the heaven of His will and the
clear evidences shining from the horizon of His grace as His testimony,
so that all people, high and low alike, may recognize that eternal
Essence and most perfect Word through this manifest and conclusive
testimony and these radiant, resplendent verses.

\subsection{92}\label{section-2685}

\begin{otherlanguage}{persian}

کلّ ذلک من فضله علی العباد من دون استحقاقهم بذلک بوده و خواهد بود که شاید
از فیوضات ربّانیّه و عنایات قدسیّهٴ الهیّه محروم نمانند والّا اعراض جمیع
موجودات و اقبال جمیع ممکنات در آن ساحت اقدس یکسان خواهد بود و دیگر آنکه
در یوم ظهور آن نیّر اعظم هیچ عذری از هیچ نفسی مسموع نخواهد بود مثلاً اگر
جمیع من فی الأرض از جمیع صحف سماویّه و کتب منزله استدلال نمایند بر امری و
آن جوهر وجود بر خلاف آن امر مسلّمهٴ محکمه ظاهر شوند احدیرا نرسد که اعتراض
نماید چه که جمیع امر و نهی باراده و قبول او مشروط و منوط است و بیانات
ظاهره از آن سمآء احدیّه مشروط و معلّق بامری نبوده و نخواهد بود

\end{otherlanguage}

Kullu dhālika min faḍlihi ʿalā al-ʿibād min dūni istiḥqāqihim bi-dhālika
būde va khāhad būd ke shāyad az foyūḍāt-e rabbānīye va ʿenāyāt-e
qodsīye-ye elāhīye maḥrūm namānand va-illā eʿrāḍ-e jamīʿ-e mowjūdāt va
eqbāl-e jamīʿ-e momkenāt dar ān sāḥat-e aqdas yeksān khāhad būd va dīgar
ānke dar yowm-e ẓohūr-e ān nayyer-e aʿẓam hīch ʿodhrī az hīch nafsī
masmūʿ nakhāhad būd mathalan agar jamīʿ-e man fī al-arḍ az jamīʿ-e
ṣoḥof-e samāvīye va kotob-e monzale estedlāl namāyand bar amrī va ān
jowhar-e vojūd bar khelāf-e ān amr-e mosallame-ye moḥkame ẓāher shavand
aḥadī rā narasad ke eʿterāḍ namāyad che ke jamīʿ-e amr va nahy be-erāde
va qabūl-e ū mashrūṭ va manūṭ ast va bayānāt-e ẓāhere az ān samā-ye
aḥadīye mashrūṭ va moʿallaq be-amrī nabūde va nakhāhad būd.

All this has been, and will remain, a token of His grace to His
servants, without their deserving it, so that they may not be deprived
of divine outpourings and God's holy favors. Otherwise, the rejection of
all beings and the acceptance of all created things would be alike in
that most holy court. Furthermore, on the day of the appearance of that
Most Great Luminary, no excuse from anyone will be accepted. For
example, if everyone on earth were to draw upon all the heavenly
scriptures and revealed books to establish a matter, and that Essence of
being were to appear contrary to that firmly established and
acknowledged matter, no one would have the right to object. For every
command and prohibition is conditional upon and dependent on His will
and acceptance, whereas the utterances issuing from that Heaven of unity
have never been, and will never be, conditional upon or dependent on
anything.

\subsection{93}\label{section-2686}

\begin{otherlanguage}{persian}

ای طالب با سینهٴ منیر باین سینای روح درآی تا تجلّی شمس معانی را در طور دل
مشاهده نمائی و اصنام وهم و ظنّ را بقوّهٴ خلیلی بشکنی و از تقلید و تقیید و
تحدید خالص و مقدّس شده بمعارج یقین و ثبوت و استقامت عروج نمائی

\end{otherlanguage}

Ey ṭāleb bā sīne-ye monīr be-īn Sīnā-ye rūḥ darāy tā tajallī-ye shams-e
maʿānī rā dar Ṭūr-e del moshāhede namāʾī va aṣnām-e vahm va ẓann rā
be-qovve-ye khalīlī beshkanī va az taqlīd va taqyīd va taḥdīd khāleṣ va
moqaddas shode be-maʿārej-e yaqīn va thobūt va esteqāmat ʿorūj namāʾī.

O seeker, enter this Sinai of the spirit with an illumined heart, so
that you may behold the splendor of the Sun of meanings upon the Sinai
of the heart, shatter the idols of fancy and conjecture with the power
of Abraham, and, freed and sanctified from imitation, restriction, and
limitation, ascend to the heights of certainty, constancy, and
steadfastness.

\subsection{94}\label{section-2687}

\begin{otherlanguage}{persian}

و دیگر ملاحظه نما که اگر جمیع من فی البیان بجمیع ما فی البیان استدلال
بامری نمایند در شرایط ظهور بعد و مسلّم دارند که آن جمال موعود و شمس وجود
باین علامت مذکوره باید ظاهر شود و آن جوهر اختیار بر خلاف آن ظاهر شود حال
چه خواهی کرد آیا بآنچه ادراک نموده مستغنی خواهی بود از عرفان آن شمس
معارف لا فوالّذی علّم هذا القلم من اسرار القدم هیچ چیز تو را کفایت
نمینماید مگر تشبّث بذیل عطوفت او

\end{otherlanguage}

Va dīgar molāḥeẓe namā ke agar jamīʿ-e man fī al-Bayān be-jamīʿ-e mā fī
al-Bayān estedlāl be-amrī namāyand dar sharāyeṭ-e ẓohūr-e baʿd va
mosallam dārand ke ān jamāl-e mowʿūd va shams-e vojūd be-īn ʿalāmat-e
madhkūre bāyad ẓāher shavad va ān jowhar-e ekhtīyār bar khelāf-e ān
ẓāher shavad ḥāl che khāhī kard āyā be-ānche edrāk namūde mostaghnī
khāhī būd az ʿerfān-e ān shams-e maʿāref lā fa-wa-alladhī ʿallama hādhā
al-qalama min asrāri al-qidam hīch chīz to rā kefāyat nemīnamāyad magar
tashabboth be-dheyl-e ʿoṭūfat-e ū.

Consider also: if all the people of the Bayán were to draw upon
everything in the Bayán to establish a condition for the next
Revelation, and were to take it as certain that the promised Beauty and
Sun of existence must appear with this stated sign, but that Essence of
free choice were to appear otherwise, what would you do then? Would what
you had understood make you independent of recognizing that Sun of
knowledge? No, by the One who taught this Pen the mysteries of eternity!
Nothing will suffice you except holding fast to the hem of His
loving-kindness.

\subsection{95}\label{section-2688}

\begin{otherlanguage}{persian}

آیا ملاحظه نشد که امم سابقه باعتکاف بر اصنام ظنون و هوی چگونه از فضل یوم
الفصل محروم مانده‌اند باری در جمیع حال و احوال پناه بحقّ برده که بآنچه در
سموات و ارض مشهود بینی چه از بیانات محقّقه و چه از امورات ظاهره از شجرهٴ
مقصود و عرفان او که که عین عرفان مظاهر لا اوّلیّه و لا آخریّهٴ الهیّه است
باز نمانی جهد کنید تا از غیث هاطلهٴ سحاب مرتفعهٴ الهیّه محجوب نمانید

\end{otherlanguage}

Āyā molāḥeẓe nashod ke omam-e sābeqe be-eʿtekāf bar aṣnām-e ẓonūn va
havā chegūne az faḍl-e yowm al-faṣl maḥrūm mānde-and bārī dar jamīʿ-e
ḥāl va aḥvāl panāh be-Ḥaqq borde ke be-ānche dar samāvāt va arḍ mashhūd
bīnī che az bayānāt-e moḥaqqaqe va che az omūrāt-e ẓāhere az shajare-ye
maqṣūd va ʿerfān-e ū ke ke ʿeyn-e ʿerfān-e maẓāher-e lā-avvalīye va
lā-ākherīye-ye elāhīye ast bāz namānī jahd konīd tā az gheyth-e
hāṭele-ye saḥāb-e mortafeʿe-ye elāhīye maḥjūb namānīd.

Have you not observed how the peoples of the past, through their
devotion to the idols of conjecture and desire, have remained deprived
of the grace of the Day of Decision? In every condition and
circumstance, take refuge in God, so that nothing you see in the heavens
and on earth, whether established utterances or outward events, may keep
you from the Tree of desire and from recognizing Him, which---which is
itself the recognition of the divine Manifestations who have neither
beginning nor end. Strive not to remain veiled from the pouring rain of
God's uplifted clouds.

\subsection{96}\label{section-2689}

\begin{otherlanguage}{persian}

اگرچه سحاب فضل در جمیع اوقات مرتفع است و او را انقطاع پدید نه ولکن هر
سحابی مظهر فیض کلّیّه نبوده و نخواهد بود در شمس و صعود و نزول او ملاحظه
نما اگرچه او در ذات خود مقدّس از صعود و نزول است ولکن بتغییر منازل و بروج
بمقادیر مختلفه ظاهر و مجلّی مثل آنکه در برج اسد بکمال ظهور و حرارت ظاهر
است و در دون آن غیر آن مشهود و حال که بفضل اللّه سحاب رحمت جمیع ارض و
ساکنین آن را احاطه نموده و شمس عنایت بکمال تجلّی و حرارت ظاهر و لایح است
خود را بوساوس شیاطین نفسیّه و هیاکل بغضیّهٴ حسدیّه مشغول ننمائید

\end{otherlanguage}

Agarche saḥāb-e faḍl dar jamīʿ-e owqāt mortafeʿ ast va ū rā enqeṭāʿ
padīd na va-lāken har saḥābī maẓhar-e feyḍ-e kollīye nabūde va nakhāhad
būd dar shams va ṣoʿūd va nozūl-e ū molāḥeẓe namā agarche ū dar dhāt-e
khod moqaddas az ṣoʿūd va nozūl ast va-lāken be-taghīr-e manāzel va
borūj be-maqādīr-e mokhtalefe ẓāher va mojallī methl-e ānke dar borj-e
Asad be-kamāl-e ẓohūr va ḥarārat ẓāher ast va dar dūn-e ān gheyr-e ān
mashhūd va ḥāl ke be-faḍl Allāh saḥāb-e raḥmat jamīʿ-e arḍ va sākenīn-e
ān rā eḥāṭe namūde va shams-e ʿenāyat be-kamāl-e tajallī va ḥarārat
ẓāher va lāʾeḥ ast khod rā be-vasāves-e shayāṭīn-e nafsīye va hayākel-e
boghḍīye-ye ḥasadīye mashghūl nanamāʾīd.

Although the cloud of grace is raised at all times and never ceases, not
every cloud has been, or will be, the manifestation of all-encompassing
bounty. Consider the sun and its ascent and descent. Although in its own
essence it is sanctified from ascent and descent, with the changing of
its stations and zodiacal signs it appears and shines in different
degrees. In Leo, for example, it appears with the fullness of its
radiance and heat, while elsewhere it is seen otherwise. Now that, by
God's grace, the cloud of mercy has encompassed the whole earth and its
inhabitants, and the sun of loving favor shines visibly with the
fullness of its splendor and heat, do not occupy yourselves with the
whisperings of the devils of self and the embodiments of hatred and
envy.

\subsection{97}\label{section-2690}

\begin{otherlanguage}{persian}

و در جمیع حال و احوال حقّ را فاعل مختار و فعّال لما یشآء دانسته و باین
کلمهٴ جامعه که اعظم کلمات منزله است ثابت و راسخ بوده لعلّ در دین اللّه و
امر او ثابت مانی و این است حقّ صرف و صدق خالص و ما بعد الحقّ الّا الضّلال و
هم‌چنانکه ذات مقدّس حقّ دلیل بر او بوده و دون او قاصر و عاجز از عرفان او
کما هو حقّه و مستحقّه هم‌چنین آیات منزله از سماء امر او مدّل بر ظهور او است
و دون او مفقود صرف و لاشیء محض و غیر مذکور در آنساحت بوده و کفی باللّه
شهیدا

\end{otherlanguage}

Va dar jamīʿ-e ḥāl va aḥvāl Ḥaqq rā fāʿel-e mokhtār va faʿʿālun li-mā
yashāʾ dāneste va be-īn kaleme-ye jāmeʿe ke aʿẓam-e kalemāt-e monzale
ast thābet va rāsekh būde laʿall dar dīn Allāh va amr-e ū thābet mānī va
īn ast ḥaqq-e ṣerf va ṣedq-e khāleṣ wa-mā baʿda al-ḥaqqi illā al-ḍalāl
va ham-chenānke dhāt-e moqaddas-e Ḥaqq dalīl bar ū būde va dūn-e ū qāṣer
va ʿājez az ʿerfān-e ū kamā huwa ḥaqquhu wa-mustaḥaqquhu ham-chenīn
āyāt-e monzale az samā-ye amr-e ū modall bar ẓohūr-e ū ast va dūn-e ū
mafqūd-e ṣerf va lā-shayʾ-e maḥḍ va gheyr-e madhkūr dar ān sāḥat būde
wa-kafā bi-llāhi shahīdan.

In every condition and circumstance, know God to be a freely choosing
Agent who does whatever He wills, and remain firm and unwavering in this
all-embracing statement, the greatest of revealed words, so that you may
remain steadfast in God's religion and His Cause. This is pure truth and
unalloyed veracity, and what is there beyond truth but error? Just as
God's sanctified Essence is His own proof, and all besides Him falls
short and is powerless to recognize Him as He truly deserves, so the
verses sent down from the heaven of His Cause bear witness to His
Revelation. All besides Him is sheer nonexistence, absolute nothingness,
and unmentioned in that court. God is sufficient as a witness.

\subsection{98}\label{section-2691}

\begin{otherlanguage}{persian}

باید در این کلمات جامعهٴ منزله کمال تعقّل و تفکّر مبذول داشت تا از هر بابی
از آن ابواب علوم لا نهایه بما لا نهایه بر وجه قلوب طالبین و قاصدین مفتوح
شود این مقامی است که احدی جز بجوهر طینت و ساذج فطرت نتواند قدم گذاشت و
اگر مرایای قلوب اقلّ من الشّعر بغبار ما سوی بیالاید ابداً تجلّی این نور منیع
و امر بدیع در آن مرایا منطبع و مرتسم و متجلّی نشود چه‌قدر ارفع و امنع است
اینمقام که جز تقدیس صرف نپذیرد و چه‌قدر عالی و متعالی است که دون تنزیه
بحت قبول نفرماید

\end{otherlanguage}

Bāyad dar īn kalemāt-e jāmeʿe-ye monzale kamāl-e taʿaq qol va tafakkor
mabdhūl dāsht tā az har bābī az ān abvāb-e ʿolūm-e lā-nehāye bi-mā
lā-nehāye bar vajh-e qolūb-e ṭālebīn va qāṣedīn maftūḥ shavad īn maqāmī
ast ke aḥadī joz be-jowhar-e ṭīnat va sādhaj-e feṭrat natavānad qadam
godhāsht va agar marāyā-ye qolūb aqall min al-shaʿr be-ghobār-e mā-sevā
bīālāyad abadan tajallī-ye īn nūr-e manīʿ va amr-e badīʿ dar ān marāyā
monṭabeʿ va mortasem va motajallī nashavad che-qadr arfaʿ va amnaʿ ast
īn maqām ke joz taqdīs-e ṣerf napa dhīrad va che-qadr ʿālī va motaʿālī
ast ke dūn-e tanzīh-e baḥt qabūl nafarmāyad.

One must devote the fullest measure of understanding and reflection to
these all-embracing revealed words, so that from each of them gates of
endlessly infinite knowledge may open before the hearts of seekers and
aspirants. This is a station into which no one can set foot except with
the essence of a pure nature and an unalloyed disposition. If the
mirrors of the heart are stained, even by less than a hair's breadth,
with the dust of anything besides Him, the splendor of this exalted
Light and wondrous Cause will never be imprinted, reflected, or
manifested in those mirrors. How lofty and inaccessible is this station,
which admits nothing but pure sanctity, and how high and exalted,
accepting nothing short of utter purity!

\subsection{99}\label{section-2692}

\begin{otherlanguage}{arabic}

فآه فوق کلّ آه فآه تحت کلّ آه فآه یملأ السّموات من حنین العاشقین فآه یملأ
الأرض من ضجیج المشتاقین فآه یجری عن العیون مدامع حمر منیر فآه ثمّ آه ثمّ
آه من اوّل الّذی لا اوّل له الی آخر الّذی لن یحصیه کلّ عالم علیم فآه الّذی
یظهر فی الوجود ناره و یحدث فی الظّهور حسرته و یحزن کلّ الممکنات و عن ورائه
اهل ملأ الصّافین فآه احترقت به الأکباد و تزلزلت عنه البلاد فآه لن یعرفه
الّا اللّه العزیز الجمیل

\end{otherlanguage}

Fa-āhin fawqa kulli āh fa-āhin taḥta kulli āh fa-āhin yamlaʾu
al-samāwāti min ḥanīni al-ʿāshiqīn fa-āhin yamlaʾu al-arḍa min ḍajīji
al-mushtāqīn fa-āhin yujrī ʿan al-ʿuyūni madāmiʿa ḥumrin munīr fa-āhin
thumma āhin thumma āhin min awwali alladhī lā awwala lahu ilā ākhiri
alladhī lan yuḥṣiyahu kullu ʿālimin ʿalīm fa-āhi alladhī yaẓharu fī
al-wujūdi nāruhu wa-yaḥduthu fī al-ẓuhūri ḥasratuhu wa-yuḥzinu kulla
al-mumkināt wa-ʿan warāʾihi ahla malaʾi al-ṣāffīn fa-āhin iḥtaraqat bihi
al-akbād wa-tazalzalat ʿanhu al-bilād fa-āhin lan yaʿrifahu illā Allāhu
al-ʿAzīzu al-Jamīl.

Ah, a sigh above every sigh! Ah, a sigh beneath every sigh! Ah, a sigh
that fills the heavens with the lovers' lamentation! Ah, a sigh that
fills the earth with the cries of those who yearn! Ah, a sigh that makes
radiant crimson tears flow from the eyes! Ah, and again ah, and yet
again ah, from the beginning that has no beginning to the end that no
learned scholar can reckon! Ah, a sigh whose fire appears in existence,
whose anguish arises in manifestation, and which grieves all created
things and, beyond them, the members of the assembly of those arrayed in
ranks! Ah, a sigh by which hearts are burned and lands are shaken! Ah, a
sigh that none can know but God, the Almighty, the All-Beautiful!

\subsection{100}\label{section-2693}

\begin{otherlanguage}{arabic}

عمّا احتجبوا النّاس عن جمال اللّه و ظهوره و طلعة اللّه و طلوعه و وقع بین یدی
هؤلآء بالذّلّة الّتی لن تقاس بذلّة و بالحزن الّذی لن یقابله حزن لا فی الأوّلین
و لا فی الآخرین و بذلک بدّل عیش کلّ ذی شعور و تزلزلت ارکان عرش عظیم ولکن
النّاس هم فی غفلة و سرور و فی فرح مبین

\end{otherlanguage}

ʿAmmā iḥtajabū al-nāsu ʿan jamāli Allāhi wa-ẓuhūrihi wa-ṭalʿati Allāhi
wa-ṭulūʿihi wa-waqaʿa bayna yaday hāʾulāʾi bi-al-dhillati allatī lan
tuqāsa bi-dhillatin wa-bi-al-ḥuzni alladhī lan yuqābilahu ḥuznun lā fī
al-awwalīna wa-lā fī al-ākhirīn wa-bi-dhālika buddila ʿayshu kulli dhī
shuʿūr wa-tazalzalat arkānu ʿarshin ʿaẓīm wa-lākin al-nāsu hum fī
ghaflatin wa-surūrin wa-fī faraḥin mubīn.

For the people have veiled themselves from the Beauty of God and His
manifestation, from the Countenance of God and His dawning, and He has
fallen into the hands of these people in a humiliation beyond comparison
and a sorrow unequaled among either former or later generations. Through
this, the life of every conscious being has been altered, and the
pillars of a mighty throne have been shaken. Yet the people remain
heedless and happy, in manifest delight.

\subsection{101}\label{section-2694}

\begin{otherlanguage}{arabic}

کذلک یبطل اللّه عمل الّذینهم غفلوا عن ذکره و کفروا بما عندهم من آیات اللّه
المقتدر العزیز القدیر و الرّوح علی اهل الرّوح من هذا البهآء المشرق القدیم

\end{otherlanguage}

Kadhālika yubṭilu Allāhu ʿamala alladhīna hum ghafalū ʿan dhikrihi
wa-kafarū bi-mā ʿindahum min āyāti Allāhi al-Muqtadiri al-ʿAzīzi
al-Qadīr. Wa-al-rūḥu ʿalā ahli al-rūḥi min hādhā al-Bahāʾi al-mushriqi
al-qadīm.

Thus God brings to nothing the deeds of those who have neglected His
remembrance and disbelieved in the verses of God, the All-Powerful, the
Almighty, the Omnipotent, that are in their possession. And may the
Spirit rest upon the people of the Spirit from this dawning, ancient
Glory.

\subsection{102}\label{section-2695}

\begin{otherlanguage}{persian}

در این ایّام تمحیص کبری خود را حفظ نموده که مبادا از صراط مستقیم بلغزی چه
که عنقریب جمیع ناس را مضطرب و متزلزل بل معرض مشاهده خواهی نمود الّا من
شآء ربّک و دو سنه قبل در مدینه ورقه نازل شد و جمیع امورات محدثهٴ بعد در
آن ثبت شده اگر اقبال بمعرفت الیوم نمائی بآن ورقهٴ مبارکه رجوع کن که شاید
نفس خود را از غرقاب فنا نجات داده بمقرّ بقا وارد شوی و از مدینهٴ اسماء و
ملکوت آن گذشته بمقام قدس کریم که مقام تقدیس صرف و تنزیه بحت باتّ است
درآئی و بغنای حقّ از ذلّت سؤال فارغ شوی چه که سؤال بنفسه دلیل بر تزلزل و
عدم استقامت سائل بوده و خواهد بود اگرچه حقّ منیع دوست داشته و میدارد که
مسئول واقع شود لأجل اعطای فضل و اعطای فیض ولکن حقّ عرفان را مقام دیگر است
و او راجع ببصر و قلب و فؤاد است نه بذکر قلم و مداد کذلک یلقیک قلم السّداد
فی هذا الیوم الّذی الیه یرجع حکم المبدأ و المعاد

\end{otherlanguage}

Dar īn ayyām-e tamḥīṣ-e kubrā khod rā ḥefẓ namūde ke mabādā az ṣerāṭ-e
mostaqīm belaghzī, che ke ʿan-qarīb jamīʿ-e nās rā moẓṭareb va
motazalzel bal moʿreż moshāhede khwāhī namūd, illā man shāʾa rabbuka. Va
do sane qabl dar Madīne varaqe nāzel shod va jamīʿ-e omūrāt-e
moḥdathe-ye baʿd dar ān thabt shode. Agar eqbāl be-maʿrefat-e al-yawm
namāʾī, be-ān varaqe-ye mobārake rojūʿ kon ke shāyad nafs-e khod rā az
gharqāb-e fanā nejāt dāde be-maqarr-e baqā vāred shavī va az madīne-ye
asmā va malakūt-e ān godhashte be-maqām-e qods-e karīm ke maqām-e
taqdīs-e ṣerf va tanzīh-e baḥt-e bātt ast darāʾī va be-ghenā-ye Ḥaqq az
dhellat-e soʾāl fāregh shavī. Che ke soʾāl be-nafse dalīl bar tazalzol
va ʿadam-e esteqāmat-e sāʾel būde va khwāhad būd, agarche Ḥaqq-e manīʿ
dūst dāshte va mīdārad ke masʾūl vāqeʿ shavad li-ajl-e eʿṭā-ye fażl va
eʿṭā-ye feyż. Va-lāken ḥaqq-e ʿerfān rā maqām-e dīgar ast va ū rājeʿ
be-baṣar va qalb va foʾād ast na be-dhekr-e qalam va medād. Kadhālika
yulqīka qalamu al-sadādi fī hādhā al-yawmi alladhī ilayhi yarjiʿu ḥukmu
al-mabdaʾi wa-al-maʿād.

In these days of the supreme testing, guard yourself lest you slip from
the straight path, for before long you will see all people agitated and
shaken---indeed, turning away---except those whom your Lord wills. Two
years ago, a sheet was revealed in the City, and all the events that
have since occurred were recorded in it. If you would turn toward
recognition today, consult that blessed sheet, so that you may perhaps
rescue yourself from the abyss of extinction and enter the abode of
everlasting life, pass beyond the city of names and its kingdom, and
enter a noble station of holiness, the station of pure sanctity and
utter, absolute transcendence. Through the wealth of God, you may then
be freed from the abasement of asking. For asking has always been, and
will remain, in itself evidence of the questioner's wavering and lack of
steadfastness, although the exalted Truth has loved, and still loves, to
be asked, so that He may bestow grace and bounty. Yet true recognition
has another station: it pertains to sight, heart, and inmost being, not
to what pen and ink set down. Thus does the Pen of rectitude convey to
you on this Day, to which the decree of the beginning and the return
reverts.

\subsection{103}\label{section-2696}

\begin{otherlanguage}{arabic}

و دیگر آنکه سؤال نموده بودی که این عبد چه رتبه و مقام را ادّعا نموده‌اند و
ما انا الّا عبد آمنت باللّه و رسله و کتبه و بکلّ ما نزل من عنده و بما نزل
علی علیّ بالحقّ و بما نزل حینئذ و بما ینزل فی زمن المستغاث حین الّذی یجدّد
فیه کلّ من فی السّموات و الأرضین هنالک تقوم السّاعة و ینصب المیزان و یرفع
الصّراط و یأتی اللّه بأمره ان انتم من العارفین و یشهد بذلک کلّ ما قدّره اللّه
فی نفسی ان انتم من الشّاهدین

\end{otherlanguage}

Va dīgar ānke soʾāl namūde būdī ke īn ʿabd che rotbe va maqām rā eddeʿā
namūde-and. Wa-mā anā illā ʿabdun āmantu bi-Allāhi wa-rusulihi
wa-kutubihi wa-bi-kulli mā nazala min ʿindihi wa-bi-mā nazala ʿalā
ʿAliyyin bi-al-ḥaqqi wa-bi-mā nazala ḥīnaʾidhin wa-bi-mā yanzilu fī
zamani al-Mustaghāthi ḥīna alladhī yujaddadu fīhi kullu man fī
al-samāwāti wa-al-araḍīn. Hunālika taqūmu al-sāʿatu wa-yunṣabu al-mīzānu
wa-yurfaʿu al-ṣirāṭu wa-yaʾtī Allāhu bi-amrihi in antum mina al-ʿārifīn.
Wa-yashhadu bi-dhālika kullu mā qaddarahu Allāhu fī nafsī in antum mina
al-shāhidīn.

You also asked what rank and station this servant has claimed. I am only
a servant who has believed in God, His Messengers, His Books, and
everything sent down from Him; in what was sent down to ʿAlī in truth;
in what has been sent down at this time; and in what will be sent down
in the time of the Mustaghāth, when everyone in the heavens and the
earths will be made anew. Then the Hour will come, the Balance will be
set up, the Bridge will be raised, and God will come with His Cause, if
you are among those who recognize. Everything God has ordained within Me
bears witness to this, if you are among those who bear witness.

\subsection{104}\label{section-2697}

\begin{otherlanguage}{persian}

و سنین معدوده است که هر آن آن خمسین الف سنه میگذرد از تتابع بلایا و
ترادف قضایا که آنی تعطیل نشده و علاوه بر این بلایای ظاهریّه فوالّذی ظهر
نفس الممکنات بسلطانه و قدرته بلایای باطنیّه که از جمیع عیون و عقول و نفوس
مستور بوده در کلّ حین و حان بوده که غیر اللّه بر آن مطّلع نشده و من دون
آنچه مذکور شد البتّه شنیده‌اید که اهل بیان چه مقدار مفتریات ظنون افکیّهٴ
خود را بر این شجرهٴ مبارکهٴ طیّبهٴ احدیّه نسبت داده و راجع نموده‌اند چنانچه
اکتفا بآنچه که اعدا وارد آورده نشده این فئه هم بآنچه قادر و مقتدر
بوده‌اند کوتاهی ننموده با اینکه وجود این عبد در جمیع اوقات سبب ارتفاع
کلمهٴ طیّبه بوده چنانچه جمیع دانسته و دیده‌اید

\end{otherlanguage}

Va senīn-e maʿdūde ast ke har ān-e ān khamsīn alf sane mīgodharad az
tatāboʿ-e balāyā va tarādof-e qażāyā ke ānī taʿṭīl nashode. Va ʿelāve
bar īn balāyā-ye ẓāheriyye, fa-wa-alladhī ẓahara nafsu al-mumkināti
bi-sulṭānihi wa-qudratihi, balāyā-ye bāṭeniyye ke az jamīʿ-e ʿoyūn va
ʿoqūl va nofūs mastūr būde dar koll-e ḥīn va ḥān būde ke gheyr-e Allāh
bar ān moṭṭaleʿ nashode. Va men dūn-e ānche madhkūr shod albatte
shenīde-īd ke ahl-e Bayān che meqdār moftariyāt-e ẓonūn-e efkiyye-ye
khod rā bar īn shajare-ye mobārake-ye ṭayyebe-ye aḥadiyye nesbat dāde va
rājeʿ namūde-and. Chenānche ektefā be-ānche ke aʿdā vāred āvarde
nashode, īn feʾe ham be-ānche qāder va moqtader būde-and kūtāhī
nanamūde, bā īnke vojūd-e īn ʿabd dar jamīʿ-e awqāt sabab-e ertefāʿ-e
kaleme-ye ṭayyebe būde, chenānche jamīʿ dāneste va dīde-īd.

For a number of years, every moment has passed like fifty thousand
years, owing to successive afflictions and relentless adversities that
have not ceased for an instant. Besides these outward afflictions---by
Him through whose sovereignty and power the very being of all contingent
things has appeared!---there have at every moment been inward
afflictions, hidden from every eye, mind, and soul, of which no one but
God has been aware. Apart from what has been mentioned, you have surely
heard how many fabrications born of their lying suppositions the people
of the Bayān have ascribed and attributed to this blessed, pure Tree of
oneness. As though what the enemies had inflicted were not enough, this
faction too has held back from nothing within its power and ability,
although the existence of this servant has at all times been the cause
of the exaltation of the good Word, as all of you have known and seen.

\subsection{105}\label{section-2698}

\begin{otherlanguage}{persian}

دیگر اگر بعضی را افکیّات ظنونیّه حایل شود بین قلوبهم و صدقهم و انکار این
فضل منیع نمایند آن امر دیگر است و الی اللّه یرجع الأمور و باطن را هم مثل
ظاهر فرض نمائید لو انتم ببصر الرّوح تنظرون و در جمیع این بلایا و محن و
رزایای سرّ و علن جز عبودیّت محضه امری و ذکری از این عبد مسجون ظهور ننموده
و چنانچه ملاحظه شد که این عبد در مدّت توقّف در مدینة السّلام با جمیع اهل آن
ارض و اطراف آن از عالم و دون آن معاشرت نموده و بقدر وسع بل فوق آن در
نصرت امر و ارتفاع آن دقیقه‌ئی تغافل نشده چنانچه ارتفاع کلمه بمقامی رسید
که جمیع اهل ارض را احاطه نمود

\end{otherlanguage}

Dīgar agar baʿżī rā efkiyyāt-e ẓonūniyye ḥāʾel shavad bayna qulūbihim
wa-ṣidqihim va enkār-e īn fażl-e manīʿ namāyand, ān amr-e dīgar ast,
wa-ilā Allāhi yarjiʿu al-umūr. Va bāṭen rā ham methl-e ẓāher farż
namāʾīd law antum bi-baṣari al-rūḥi tanẓurūn. Va dar jamīʿ-e īn balāyā
va meḥan va razāyā-ye serr va ʿalan joz ʿobūdiyyat-e maḥże amrī va
dhekrī az īn ʿabd-e masjūn ẓohūr nanamūde. Va chenānche molāḥeẓe shod ke
īn ʿabd dar moddat-e tavaqqof dar Madīnat al-Salām bā jamīʿ-e ahl-e ān
arż va aṭrāf-e ān az ʿālem va dūn-e ān moʿāsherat namūde va be-qadr-e
vosʿ bal fawq-e ān dar noṣrat-e Amr va ertefāʿ-e ān daqīqe-ī taghāfol
nashode, chenānche ertefāʿ-e kaleme be-maqāmī resīd ke jamīʿ-e ahl-e arż
rā eḥāṭe namūd.

If, however, conjectural falsehoods should come between some people's
hearts and their sincerity, and they should deny this exalted grace,
that is another matter; to God all affairs return. Consider the inward
to be like the outward as well, if you look with the eye of the spirit.
Amid all these afflictions, trials, and calamities, hidden and manifest,
nothing but pure servitude has appeared in any act or utterance of this
imprisoned servant. As has been observed, throughout His stay in the
City of Peace this servant associated with all the people of that land
and its surroundings, learned and otherwise. To the extent of His
ability---indeed, beyond it---not a moment was neglected in aiding and
exalting the Cause, until the exaltation of the Word reached such a
degree that it encompassed all the people of the earth.

\subsection{106}\label{section-2699}

\begin{otherlanguage}{persian}

و ابداً با اعدای الهی مداهنه ننموده بکمال جدّ و اجتهاد در ذکر کلمهٴ ذکریّه
سعی بلیغ و جهد منیع مبذول شده و بالأخره امر منجرّ بحرکت از مدینه شد آن هم
بر هر ذی بصری معلوم شده که بچه نحو در این خروج و هجرت اعتزاز این امر و
علوّ آن ظاهر شده تا آنکه باین مقرّ مسجون و محبوس گشته‌ایم و نشکر اللّه فی کلّ
ذلک و نحمده و نذکره و نصبر علی قضایاه و انّه موفی اجور الصّابرین و الحمد
للّه علی کلّ ذلک و لا ملجأ الّا به و لا مستعان الّا هو و انّه لهو المقتدر
العزیز الرّحیم

\end{otherlanguage}

Va abadan bā aʿdā-ye elāhī modāhene nanamūde, be-kamāl-e jedd va ejtehād
dar dhekr-e kaleme-ye dhekriyye saʿy-e balīgh va jahd-e manīʿ mabdhūl
shode. Va be-al-ākhere amr monjarr be-ḥarakat az Madīne shod. Ān ham bar
har dhī-baṣarī maʿlūm shode ke be-che naḥv dar īn khorūj va hejrat
eʿtezāz-e īn Amr va ʿolovv-e ān ẓāher shode, tā ānke be-īn maqarr masjūn
va maḥbūs gashte-īm. Wa-nashkuru Allāha fī kulli dhālika wa-naḥmaduhu
wa-nadhkuruhu wa-naṣbiru ʿalā qaḍāyāhu wa-innahu muwaffī ujūri
al-ṣābirīn. Wa-al-ḥamdu li-Allāhi ʿalā kulli dhālika wa-lā maljaʾa illā
bihi wa-lā mustaʿāna illā huwa wa-innahu la-huwa al-Muqtadiru al-ʿAzīzu
al-Raḥīm.

Never was any compromise made with the enemies of God. With the utmost
earnestness and diligence, strenuous effort and resolute exertion were
devoted to proclaiming the Word of remembrance. In the end, matters led
to departure from the City. Every discerning person has also come to
know how the might and exaltation of this Cause were manifested in that
departure and migration, until We became imprisoned and confined in this
place. In all this We thank God, praise Him, remember Him, and bear
patiently with His decrees. He indeed pays in full the rewards of the
patient. Praise be to God for all this. There is no refuge except in Him
and no one whose help can be sought but Him. He is indeed the
All-Powerful, the Almighty, the Merciful.

\subsection{107}\label{section-2700}

\begin{otherlanguage}{persian}

بلی در جمیع این سنین ماضیه بحمد اللّه ابحر معانی و انوار ربّانی در دل
متموّج و مستضیء بوده چنانچه من غیر اراده رشحات آن بر ممکنات مبذول میشد و
بقسمی عنایت ربّ العزّة احاطه فرموده که احصا و تحدید آن از امکان و اهل آن
خارج و اگر از این عبد در بیان این مراتب در الواح منزله ذکری رفته لأجل
اتّباع این کلمهٴ منیعه بوده که میفرماید و امّا بنعمة ربّک فحدّث ولکن آنجناب
نظر باصل امر داشته که بدون خود هرگز مشتبه نشده و نخواهد شد

\end{otherlanguage}

Balī dar jamīʿ-e īn senīn-e māżiye be-ḥamd-e Allāh abḥor-e maʿānī va
anvār-e rabbānī dar del motamavvej va mostażīʾ būde, chenānche men
gheyr-e erāde rashaḥāt-e ān bar momkenāt mabdhūl mīshod. Va be-qesmī
ʿenāyat-e Rabb al-ʿEzzat eḥāṭe farmūde ke eḥṣā va taḥdīd-e ān az emkān
va ahl-e ān khārej. Va agar az īn ʿabd dar bayān-e īn marāteb dar
alvāḥ-e monzale dhekrī rafte, li-ajl-e ettebāʿ-e īn kaleme-ye manīʿe
būde ke mīfarmāyad: wa-ammā bi-niʿmati rabbika fa-ḥaddith. Va-lāken
ān-jenāb naẓar be-aṣl-e amr dāshte ke be-dūn-e khod harg ez moshtabeh
nashode va nakhwāhad shod.

Indeed, throughout all these past years, praise be to God, oceans of
meaning have surged and divine lights have shone within the heart, so
that, without any act of will, their sprinklings were bestowed upon
contingent beings. The favor of the Lord of might has so encompassed Us
that to number or delimit it is beyond the contingent world and its
inhabitants. If this servant has made any mention of these degrees in
the revealed Tablets, it has been in obedience to this exalted
utterance: ``As for the bounty of your Lord, proclaim it.'' You,
however, should look to the essence of the matter, which has never been,
and will never be, confused with anything other than itself.

\subsection{108}\label{section-2701}

\begin{otherlanguage}{arabic}

و اگر حقّ منیع را نتوانی بنفس او ادراک نمائی و مستبصر شوی و دون آن بکلمات
منزله از سماء قدس احدیّه ناظر باشی رجوع بکتب منزله از سماء مشیّت نما که در
اکثر آن ذکر این کلمهٴ جامعهٴ مستوره تصریحاً و تلویحاً مشاهده خواهی نمود در
توقیع عظیم و دون آن ملاحظه نما و هم‌چنین اکثر از آیات قیّوم اسماء مشعر بر
قیّومیّت این سدرهٴ مرتفعهٴ منیعه بوده قسم بخدا که شرم دارم از ذکر این
بیانات چه که ایّامی از ظهور آن شمس مشرقهٴ احدیّه نگذشته و کلّ کیفیّت ظهور را
بعین خود مشاهده نموده‌اند و مع‌ذلک جمیع محتجب گشته‌اند

\end{otherlanguage}

Va agar Ḥaqq-e manīʿ rā natavānī be-nafs-e ū edrāk namāʾī va mostabṣer
shavī va dūn-e ān be-kalemāt-e monzale az samā-ye qods-e aḥadiyye nāẓer
bāshī, rojūʿ be-kotob-e monzale az samā-ye mashiyyat namā ke dar
akthar-e ān dhekr-e īn kaleme-ye jāmeʿe-ye mastūre taṣrīḥan va talvīḥan
moshāhede khwāhī namūd. Dar Tawqīʿ-e ʿAẓīm va dūn-e ān molāḥeẓe namā. Va
ham-chenīn akthar az āyāt-e Qayyūm-e Asmā moshʿer bar qayyūmiyyat-e īn
sedre-ye mortafeʿe-ye manīʿe būde. Qasam be-Khodā ke sharm dāram az
dhekr-e īn bayānāt, che ke ayyāmī az ẓohūr-e ān shams-e moshreqe-ye
aḥadiyye nagodhashte va koll keyfiyyat-e ẓohūr rā be-ʿayn-e khod
moshāhede namūde-and va maʿa-dhālek jamīʿ moḥtajeb gashte-and.

If you cannot apprehend the exalted Truth through His own Self and gain
insight, but look instead to the words sent down from the heaven of the
holiness of oneness, turn to the Books sent down from the heaven of the
divine Will. In most of them you will find mention, both explicit and
implicit, of this all-embracing, concealed Word. Examine the Great
Epistle and other writings. Likewise, most of the verses of the Qayyūm
al-Asmāʾ indicate the all-sustaining sovereignty of this lofty, exalted
Tree. I swear by God that I am ashamed to mention these statements, for
only a few days have passed since the appearance of that dawning Sun of
oneness, and all have witnessed with their own eyes the manner of that
appearance. Yet all have become veiled.

\subsection{109}\label{section-2702}

\begin{otherlanguage}{persian}

و کدام ذکر و بیان اظهر از نفس خود ظهور است تا بآن استدلال شود باری حال
ملاحظه نما که چه مقدار ستر و تعویق این امر مشعر و مدلّ بر علوّ این ظهور
است که گویا تعویق آن را سببی جز امتحان و افتتان کلّ من فی السّموات و الأرض
نبوده و هم‌چنین میفرماید انتم فی التّسع بلقآء اللّه ترزقون اگرچه بعضی این
ظهور قدس صمدانی و طلوع نور سبحانی را اراده نموده‌اند که بدون سدرهٴ خود
راجع نمایند و بجدّ تمام در تحریف این کلمهٴ جامعه سعی نمایند ولکن الأمر
بید اللّه یهدی من یشآء الی صراط مستقیم

\end{otherlanguage}

Va kodām dhekr va bayān aẓhar az nafs-e khod-e ẓohūr ast tā be-ān
estedlāl shavad? Bārī ḥāl molāḥeẓe namā ke che meqdār satr va taʿvīq-e
īn Amr moshʿer va modell bar ʿolovv-e īn ẓohūr ast, ke gūyā taʿvīq-e ān
rā sababī joz emteḥān va eftetān-e kulli man fī al-samāwāti wa-al-arḍ
nabūde. Va ham-chenīn mīfarmāyad: antum fī al-tisʿi bi-liqāʾi Allāhi
turzaqūn. Agarche baʿżī īn ẓohūr-e qods-e ṣamadānī va ṭolūʿ-e nūr-e
sobḥānī rā erāde namūde-and ke be-dūn-e sedre-ye khod rājeʿ namāyand va
be-jedd-e tamām dar taḥrīf-e īn kaleme-ye jāmeʿe saʿy namāyand, va-lāken
al-amru bi-yadi Allāhi yahdī man yashāʾu ilā ṣirāṭin mustaqīm.

What mention or statement is more manifest than the Revelation itself,
that it might be used as proof of it? Consider now how greatly the
concealment and delay of this Cause indicate and attest the exaltation
of this Revelation, as though its delay had no purpose other than to
test and try everyone in the heavens and on earth. He likewise says:
``In the year nine you will be granted the meeting with God.'' Although
some have sought to refer this manifestation of eternal holiness and
this dawning of divine light to something other than His own Tree, and
strive with all earnestness to distort this all-embracing utterance, the
matter is in God's hands: He guides whomever He wills to a straight
path.

\subsection{110}\label{section-2703}

\begin{otherlanguage}{arabic}

ای اهل بیان کمال سعی مبذول داشته که مثل اهل فرقان نباشید که جمیع صفات و
اسماء را از منتهی مقام عزّ جبروت الی ادنی رتبهٴ ناسوت بطلعات غیب راجع
نموده‌اند و از مطالع احدیّه و مظاهر قدمیّه که بالفعل موجودند غافل بوده و
خواهند بود

\end{otherlanguage}

Ey ahl-e Bayān, kamāl-e saʿy mabdhūl dāshte ke methl-e ahl-e Forqān
nabāshīd ke jamīʿ-e ṣefāt va asmā rā az montahā maqām-e ʿezz-e jabarūt
ilā adnā rotbe-ye nāsūt be-ṭalaʿāt-e gheyb rājeʿ namūde-and va az
maṭāleʿ-e aḥadiyye va maẓāher-e qadamiyye ke be-al-feʿl mawjūdand ghāfel
būde va khwāhand būd.

People of the Bayān! Make every effort not to be like the people of the
Qurʾān, who have referred all attributes and names, from the highest
station of the glory of dominion to the lowest degree of the human
realm, to unseen countenances, while remaining heedless of the
Dawning-Places of oneness and Manifestations of eternity who actually
exist.

\subsection{111}\label{section-2704}

\begin{otherlanguage}{persian}

حال ای عباد ببصر عدل و انصاف مشاهده نمائید که این عبد قسم بظهور من یظهره
اللّه که اگر در مدّت اقامت در دارالسّلام اگر یومی و یا لیلی بر حسب ظاهر
گمان حیات مینمودم چه که اعدا از یمین و یسار و جنوب و شمال بکمال جدّ در
صدد بوده‌اند چنانچه جمیع عباد ساکنهٴ آن محلّ مطّلعند تا آنکه تقدیر ربّانی بر
تزویر شیطانی فایق آمد و سلطان علّیّین بر صور سجّین غلبه نمود و با این علوّ و
اشتهار با جمعی از مهاجران هجرت نموده و ارتفاع این اسم و شهرت و ظهور آن
من غیر اراده واقع شده

\end{otherlanguage}

Ḥāl ey ʿebād, be-baṣar-e ʿadl va enṣāf moshāhede namāʾīd ke īn ʿabd,
qasam be-ẓohūr-e Man Yuẓhiruhu Allāh, ke agar dar moddat-e eqāmat dar
Dār al-Salām agar yawmī va yā laylī bar ḥasab-e ẓāher gomān-e ḥayāt
mīnamūdam, che ke aʿdā az yamīn va yasār va jonūb va shamāl be-kamāl-e
jedd dar ṣadad būde-and, chenānche jamīʿ-e ʿebād-e sākene-ye ān maḥall
moṭṭaleʿand. Tā ānke taqdīr-e rabbānī bar tazvīr-e sheyṭānī fāʾeq āmad
va solṭān-e ʿelliyyīn bar ṣovar-e sejjīn ghalabe namūd va bā īn ʿolovv
va eshtehār bā jamʿī az mohājerān hejrat namūde va ertefāʿ-e īn esm va
shohrat va ẓohūr-e ān men gheyr-e erāde vāqeʿ shode.

Now, servants, look with the eye of justice and fairness: I swear by the
Revelation of Him whom God will make manifest that, throughout this
servant's residence in the Abode of Peace, not for a day or a night did
outward circumstances give Me reason to expect to remain alive. For
enemies on the right and left, to the south and north, were intent upon
their purpose with the utmost determination, as all the inhabitants of
that place know. At last divine destiny prevailed over satanic
deception, and the sovereignty of the heights overcame the forms of the
depths. With such exaltation and renown, this servant departed with a
company of emigrants. The exaltation of this Name, its fame, and its
manifestation came about without any act of will.

\subsection{112}\label{section-2705}

\begin{otherlanguage}{persian}

دیگر چه ذکر نمایم که چه بر این عبد گذشته و وارد شده فواللّه لن یقدر ان
یسمعه ذی حبّ و لا ذی قلب و لا ذی نظر و لا ذی شعور

\end{otherlanguage}

Dīgar che dhekr namāyam ke che bar īn ʿabd godhashte va vāred shode?
Fa-wa-Allāhi lan yaqdira an yasmaʿahu dhī ḥubbin wa-lā dhī qalbin wa-lā
dhī naẓarin wa-lā dhī shuʿūr.

What more can I say of what this servant has endured and what has
befallen Him? By God, no one possessed of love, a heart, vision, or
awareness could bear to hear it.

\subsection{113}\label{section-2706}

\begin{otherlanguage}{persian}

باری لم‌یزل اسرار الهیّه در کنوز عصمتیّه مستور بوده و لایزال در حجبات حفظیّه
مصون و محفوظ خواهد بود و اگر جبل انفس عباد مندکّ نمیشد رمزی از اسرار
محجّبهٴ مقنّعهٴ مستره ذکر مینمودم ولکن اللّه مظهر ما یشآء و ما انا بمظهر ما
هو المستور و بر همهٴ شما و اهل ارض معلوم بوده که چه مقدار عبودیّت و خضوع
از این عبد بجمیع ادلّا و احبّا از اوّل ظهور شمس حقیقت الی الیوم ظاهر شده و
مع‌ذلک چه نسبتهای کذبه باین عبد راجع و احدی در صدد رفع آن برنیامده الی
اللّه فوّضت امری و علیه توکّلت و علیه فلیتوکّلنّ عباده المنقطعین

\end{otherlanguage}

Bārī lam-yazal asrār-e elāhiyye dar konūz-e ʿeṣmatiyye mastūr būde va
lā-yazāl dar ḥojobāt-e ḥefẓiyye maṣūn va maḥfūẓ khwāhad būd. Va agar
jabal-e anfos-e ʿebād mondakk namīshod ramzī az asrār-e moḥajjabe-ye
moqannaʿe-ye mostatar dhekr mīnamūdam, wa-lākin Allāhu muẓhiru mā
yashāʾu wa-mā anā bi-muẓhiri mā huwa al-mastūr. Va bar hame-ye shomā va
ahl-e arż maʿlūm būde ke che meqdār ʿobūdiyyat va khożūʿ az īn ʿabd
be-jamīʿ-e adellā va aḥebbā az avval-e ẓohūr-e Shams-e Ḥaqīqat ilā
al-yawm ẓāher shode. Va maʿa-dhālek che nesbat-hā-ye kadhebe be-īn ʿabd
rājeʿ va aḥadī dar ṣadad-e rafʿ-e ān barnayāmade. Ilā Allāhi fawwaḍtu
amrī wa-ʿalayhi tawakkaltu wa-ʿalayhi fa-l-yatawakkalanna ʿibāduhu
al-munqaṭiʿīn.

In any case, the divine mysteries have always been concealed in the
treasuries of protection and will forever remain secure and guarded
within sheltering veils. Were the mountain of the servants' souls not to
be crushed, I would mention a hint of the veiled, shrouded, hidden
mysteries. But God reveals whatever He wills, and I will not reveal what
is concealed. All of you, and the people of the earth, have known how
much servitude and humility this servant has shown toward all the guides
and loved ones, from the first appearance of the Sun of Truth until
today. Yet how many false accusations have been directed against this
servant, and no one has sought to dispel them. To God I have committed
My affair, and in Him I have placed My trust. In Him let His detached
servants place their trust.

\subsection{114}\label{section-2707}

\begin{otherlanguage}{persian}

و دیگر سؤال شده بود که چگونه میشود کلمات عالین بسجّین تبدیل شوند و حروف
اثبات از مظاهر نفی محسوب شود این معلوم آنجناب بوده که آن شمس فیض ربّانی و
بحر فضل سبحانی اشراق فرمود بر جمیع مرایای موجودات در بعضی که مقابله
نموده‌اند بقرصه منطبع میشود و بر دون آن همان روشنی شعاع ظاهر اینست که آن
شمس قدم بعد از انطباع او در صور مرایا آنچه وصف نماید مرایا را باوصاف غیر
متناهیه خود را وصف نموده چه که انوار وجه خود را در مرایا مشاهده مینماید
و جمیع این اوصاف منیعه فی‌الحقیقه بخود او راجع است و این فضل در مرایا
موجود است بالتّقابل الیها و بعد از انحراف فانی و معدوم خواهند بود

\end{otherlanguage}

Va dīgar suʾāl shode būd ke chegūne mīshavad kalimāt-i ʿālīn bi-Sijjīn
tabdīl shavand va ḥurūf-i ithbāt az maẓāhir-i nafy maḥsūb shavad. Īn
maʿlūm-i ān janāb būde ke ān shams-i fayḍ-i rabbānī va baḥr-i faḍl-i
subḥānī ishrāq farmūd bar jamīʿ-i marāyā-yi mawjūdāt. Dar baʿḍī ke
muqābele namūde-and bi-qurṣihi munṭabiʿ mīshavad va bar dūn-i ān hamān
rawshanī-yi shuʿāʿ ẓāhir. Īnast ke ān shams-i qidam baʿd az inṭibāʿ-i ū
dar ṣuvar-i marāyā ānche vaṣf namāyad marāyā rā bi-awṣāf-i ghayr-i
mutanāhiye khwad rā vaṣf namūde, che ke anvār-i vajh-i khwad rā dar
marāyā mushāhede mīnamāyad va jamīʿ-i īn awṣāf-i manīʿe fī'l-ḥaqīqe
bi-khwad-i ū rājiʿ ast va īn faḍl dar marāyā mawjūd ast bi'l-taqābul
ilayhā va baʿd az inḥirāf fānī va maʿdūm khwāhand būd.

Another question was asked: how can the words of the exalted be changed
into Sijjin, and the letters of affirmation be counted among the
manifestations of negation? You know that the Sun of divine grace and
the Ocean of heavenly bounty shone upon all the mirrors of existence. In
some that have turned to face it, its disk is imprinted, while upon the
others only the brightness of its rays appears. Thus, after that Sun of
eternity has been imprinted in the forms of the mirrors, whatever
description it gives of the mirrors, with attributes beyond number, is a
description of itself, for it beholds the lights of its own face in the
mirrors. All these exalted attributes in reality refer to it alone. This
bounty exists in the mirrors through their facing it; once they turn
away, these attributes will vanish and cease to exist.

\subsection{115}\label{section-2708}

\begin{otherlanguage}{persian}

حال نظر را پاک نما از غبار ما سوی تا عارف بر شمسی شوی که اگر جمیع مرایای
موجودات اقلّ من آن باو مقابل شوند در جمیع انوار حقیقت و آثار عزّ احدیّت
لم‌یزل و لایزال منطبع و مرتسم شود و من دون ذلک یمحو اللّه ما یشآء و یثبت و
عنده امّ الکتاب فاعل مختار مختار است محو میفرماید بعد از اثبات و کذلک
بالعکس و در این مقام است که میفرماید لا تفرحوا بما اتاکم و لا تیأسوا عمّا
فاتکم و اینمقام نقطهٴ عدل است که جمیع ما سوی خائف و راجی بوده و خواهند
بود هر ذی علوّی خائف است از محو و هر ذی دنوّی راجی بمقام اثبات فوالّذی نفس
حسین بیده که اگر شرح اینمقام را بما قدّر فیه ذکر نمایم اکثر از ناس را
متزلزل بینی قوله عزّ ذکره را متذکّر شو لو شئنا لنذهبنّ ما اوحینا الیک
فوعمری لو القی علیک ازید من ذلک لیرتاب المبطلون

\end{otherlanguage}

Ḥāl naẓar rā pāk namā az ghubār-i mā sivā tā ʿārif bar shamsī shavī ke
agar jamīʿ-i marāyā-yi mawjūdāt aqall-i min ān bi-ū muqābil shavand dar
jamīʿ anvār-i ḥaqīqat va āthār-i ʿizz-i aḥadiyyat-i lam-yazal va
lā-yazāl munṭabiʿ va murtasim shavad va min dūni dhālika yamḥū Allāhu mā
yashāʾu wa-yuthbitu wa-ʿindahu ummu'l-kitāb. Fāʿil-i mukhtār mukhtār
ast. Maḥv mīfarmāyad baʿd az ithbāt va kadhālika bi'l-ʿaks va dar īn
maqām ast ke mīfarmāyad lā tafraḥū bi-mā ātākum wa-lā tayʾasū ʿammā
fātakum va īn maqām nuqṭe-yi ʿadl ast ke jamīʿ-i mā sivā khāʾif va rājī
būde va khwāhand būd. Har dhī ʿuluvvī khāʾif ast az maḥv va har dhī
dunuvvī rājī bi-maqām-i ithbāt. Fa-wa'lladhī nafsu Ḥusaynin bi-yadihi ke
agar sharḥ-i īn maqām rā bi-mā quddira fīhi dhikr namāyam akthar az nās
rā mutazalzil bīnī. Qawluhu ʿazza dhikruhu rā mutadhakkir shaw: law
shiʾnā la-nadhhabanna mā awḥaynā ilayka. Fa-wa-ʿamrī law ulqiya ʿalayka
azyadu min dhālika la-yartābu'l-mubṭilūn.

Now cleanse your vision of the dust of all else, so that you may
recognize a Sun such that, were all the mirrors of existence to face it
for the briefest instant, the lights of truth and the signs of the glory
of eternal and everlasting unity would be imprinted and reflected in
them all. Otherwise, ``God effaces whatever He wills and establishes
whatever He wills, and with Him is the Mother Book.'' The freely
choosing Agent is free to choose. He effaces after establishing, and
likewise the reverse. It is in this connection that He says, ``Do not
rejoice in what has come to you, nor despair over what has escaped
you.'' This station is the point of justice, at which all else has ever
stood, and will ever stand, in fear and hope. Everyone who occupies a
lofty position fears effacement, and everyone in a lowly position hopes
for the station of affirmation. By Him in whose hand is the soul of
Husayn! Were I to explain this station according to what has been
ordained within it, you would see most people shaken. Recall His
words---exalted be His mention---``If We willed, We would surely take
away what We have revealed to you.'' By My life! Were more than this
imparted to you, those who follow falsehood would surely be filled with
doubt.

\subsection{116}\label{section-2709}

\begin{otherlanguage}{persian}

دیگر در اینمقام بیش از این دوست نمیدارم که ذکر نمایم چه که ناس در وهم
صرف مستغرقند و دیگر از بعضی امور واقعه که سؤال نمودی در غیبت این عبد
واقع شده و حسنه و سیّئهٴ آن بعاملین و فاعلین راجع است باری جمیع امور واضح
و لایح است اگر حجبات وهمیّهٴ ناس خرق شود و ببصر خود در امر اللّه تفرّس
نمایند جمیع اسرار را مشهود مشاهده نمایند انشآءاللّه امیدواریم که فضلی از
منبع جود و افضال ربّ العزّة و العظمة ظاهر شود و جمیع محتجبین را از حجبات
ظلمت بسرادقات وحدت کشاند واللّه معطی ما یشآء و انّه لهو الفضّال القدیم و
الحمد للّه ربّ العالمین

\end{otherlanguage}

Dīgar dar īn maqām bīsh az īn dūst namīdāram ke dhikr namāyam che ke nās
dar vahm-i ṣirf mustaghriqand va dīgar az baʿḍī umūr-i vāqiʿe ke suʾāl
namūdī dar ghaybat-i īn ʿabd vāqiʿ shode va ḥasane va sayyiʾe-yi ān
bi-ʿāmilīn va fāʿilīn rājiʿ ast. Bārī jamīʿ-i umūr vāḍiḥ va lāyiḥ ast
agar ḥujubāt-i vahmiyye-yi nās kharq shavad va bi-baṣar-i khwad dar
amr-i Allāh tafarrus namāyand jamīʿ-i asrār rā mashhūd mushāhede
namāyand. In shāʾ Allāh umīdvārīm ke faḍlī az manbaʿ-i jūd va ifḍāl-i
Rabb-i ʿIzzat va ʿAẓamat ẓāhir shavad va jamīʿ-i muḥtajibīn rā az
ḥujubāt-i ẓulmat bi-surādiqāt-i vaḥdat kashānad. Wa-Allāhu muʿṭī mā
yashāʾu wa-innahu la-huwa'l-Faḍḍālu'l-Qadīm. Wa'l-ḥamdu li'llāhi
Rabbi'l-ʿālamīn.

I do not wish to say anything further on this subject, for people are
immersed in sheer illusion. As for some of the events about which you
asked, they occurred in this servant's absence, and their good and evil
rest with those who carried them out. In any case, everything is clear
and manifest. If the veils of people's illusions were torn apart and
they examined the Cause of God with their own eyes, they would plainly
behold all its mysteries. God willing, We hope that a bounty may appear
from the source of the generosity and grace of the Lord of might and
grandeur, drawing all who are veiled from the veils of darkness into the
pavilions of unity. God bestows whatever He wills, and He is indeed the
Most Bountiful, the Ancient. Praise be to God, the Lord of all worlds.

\subsection{117}\label{section-2710}

\begin{otherlanguage}{persian}

و دیگر وصیّت این عبد جمیع را آنکه آنی غفلت ننموده مراقب امر اللّه باشند و
نهایت جدّ و اجتهاد را در حفظ و صیانت آن مرعی داشته که این است کلّ خیر و
اصل آن لأجل آنکه اوّل امر معرفة اللّه بوده و خواهد بود و کلّ آنچه مسموع و
مشهود شده طایف حول این امر است چه که جمیع از او ظاهر و باو راجع و بعد از
معرفت ذات احدیّه اتّباع امر و اجرای حدود او از اعظم امور بوده و لم‌یزل و
لایزال نفحات ذکریّه از مدینهٴ امریّه در هبوب بوده و آثار الهی از این نفحهٴ
قدس صمدانی بین بریّه قائم خواهد بود باری انشآءاللّه بقلب صافی خالص توجّه
بمبدأ خود در کلّ اوقات و احیان نموده و از ما سوی اللّه منقطع شده که شاید
در این ایّام مظلمه از معرفت آن شمس احدیّه محجوب نمانند

\end{otherlanguage}

Va dīgar vaṣiyyat-i īn ʿabd jamīʿ rā ānke ānī ghaflat nanamūde murāqib-i
amr-i Allāh bāshand va nihāyat-i jidd va ijtihād rā dar ḥifẓ va
ṣiyānat-i ān marʿī dāshte ke īn ast kull-i khayr va aṣl-i ān li-ajl-i
ānke avval-i amr maʿrifat-i Allāh būde va khwāhad būd va kull-i ānche
masmūʿ va mashhūd shode ṭāyif-i ḥawl-i īn amr ast che ke jamīʿ az ū
ẓāhir va bi-ū rājiʿ va baʿd az maʿrifat-i dhāt-i aḥadiyye ittibāʿ-i amr
va ijrā-yi ḥudūd-i ū az aʿẓam-i umūr būde va lam-yazal va lā-yazāl
nafaḥāt-i dhikriyye az madīne-yi amriyye dar hubūb būde va āthār-i ilāhī
az īn nafḥe-yi quds-i ṣamadānī bayn-i bariyye qāʾim khwāhad būd. Bārī in
shāʾ Allāh bi-qalb-i ṣāfī-yi khāliṣ tavajjuh bi-mabdaʾ-i khwad dar
kull-i awqāt va aḥyān namūde va az mā siva'llāh munqaṭiʿ shode ke shāyad
dar īn ayyām-i muẓlime az maʿrifat-i ān shams-i aḥadiyye maḥjūb
namānand.

Further, this servant counsels everyone not to be heedless for a moment,
but to watch over the Cause of God and exercise the utmost diligence and
effort in preserving and protecting it, for this is the whole of
goodness and its root. The first concern has been, and will ever be, the
recognition of God, and everything that has been heard and witnessed
revolves around this matter, for everything proceeds from Him and
returns to Him. After recognition of the one divine Essence, following
His command and enforcing His laws are among the greatest concerns. The
breezes of remembrance have always blown, and will forever blow, from
the city of His Cause, and through this holy, eternal breath the signs
of God will remain established among His creatures. God willing, may all
turn toward their Source with pure and sincere hearts at every moment
and at all times, detaching themselves from everything but God, so that
in these darkened days they may not remain veiled from recognizing that
Sun of oneness.

\subsection{118}\label{section-2711}

\begin{otherlanguage}{arabic}

و این عبد از ربّ العزّة مسئلت مینماید که بصر و سمعی مقدّس از تحدید و تقلید
بجمیع عنایت فرماید و اللّه واهب ما یشآء و انّه لهو الوهّاب العلیم الکریم و
اختم نغمات الرّوح بذکر اللّه الأعظم الأکبر العظیم لیکون ختامها مسک الا انّ
بذلک تستجذب ارواح المخلصین فی کلّ حین و الحمد للّه ربّ العالمین

\end{otherlanguage}

Va īn ʿabd az Rabb-i ʿIzzat masʾalat mīnamāyad ke baṣar va samʿī
muqaddas az taḥdīd va taqlīd bi-jamīʿ ʿināyat farmāyad. Wa-Allāhu wāhibu
mā yashāʾu wa-innahu la-huwa'l-Wahhābu'l-ʿAlīmu'l-Karīm. Wa-akhtimu
naghamāti'l-rūḥi bi-dhikri Allāhi'l-Aʿẓami'l-Akbari'l-ʿAẓīm li-yakūna
khitāmuhā misk. Alā inna bi-dhālika tustajdhabu arwāḥu'l-mukhliṣīna fī
kulli ḥīn. Wa'l-ḥamdu li'llāhi Rabbi'l-ʿālamīn.

This servant asks the Lord of might to bestow on everyone sight and
hearing sanctified from limitation and imitation. God bestows whatever
He wills, and He is indeed the Ever-Bestowing, the All-Knowing, the
Generous. I conclude the melodies of the Spirit with the remembrance of
God, the Most Great, the Supreme, the Mighty, so that their close may be
musk. Truly, through this the souls of the sincere are drawn at every
moment. Praise be to God, the Lord of all worlds.

\bookmarksetup{startatroot}

\chapter{Javahiru'l-Asrar (Gems of Divine
Mysteries)}\label{javahirul-asrar-gems-of-divine-mysteries}

\section{BH00012}\label{bh00012}

\emph{Translated on 17 Qudrat 183 using GPT 6.1 Sol}

\subsection{1}\label{section-2712}

\begin{otherlanguage}{arabic}

جواهر الأسرار فی معارج الأسفار لمن اراد ان یتقرّب باللّه المقتدر الغفّار
فهنیاً للأبرار الّذین یشربون من هذه الأنهار

\end{otherlanguage}

Jawāhir al-asrār fī maʿārij al-asfār li-man arāda an yataqarraba
bi-Allāh al-Muqtadir al-Ghaffār, fa-hanīʾan li-l-abrār alladhīna
yashrabūna min hādhihi al-anhār.

Gems of mysteries in the ascending stages of journeys, for whoever
desires to draw near to God, the All-Powerful, the Ever-Forgiving.
Blessed are the righteous who drink from these rivers.

\subsection{2}\label{section-2713}

\begin{otherlanguage}{arabic}

هو العلیّ الأعلی

\end{otherlanguage}

Huwa al-ʿAliyy al-Aʿlā.

He is the Exalted, the Most High.

\subsection{3}\label{section-2714}

\begin{otherlanguage}{arabic}

یا ایّها السّالک فی سبل العدل و النّاظر الی طلعة الفضل قد بلغ کتابک و عرفت
سؤالک و سمعت لحنات قلبک فی سرادق فؤادک اذاً قد رفعت سحاب الارادة لتمطر
علیک من امطار الحکمة لتأخذ عنک کلّ ما اخذت من قبل و تقلّبک عن جهات الضّدیّة
الی مکمن الأحدیّة و تصلک الی شریعة القدسیّة لتشرب عنها و تستریح نفسک فیها
و یسکن عطشک و یبرد فؤادک و تکون من الّذینهم کانوا الیوم بنور اللّه لمهتدین

\end{otherlanguage}

Yā ayyuhā al-sālik fī subul al-ʿadl wa-l-nāẓir ilā ṭalʿat al-faḍl, qad
balagha kitābuka wa-ʿaraftu suʾālaka wa-samiʿtu laḥanāt qalbika fī
surādiq fuʾādika. Idhan qad rafaʿtu saḥāb al-irāda li-tumṭira ʿalayka
min amṭār al-ḥikma li-taʾkhudha ʿanka kulla mā akhadhta min qabl
wa-tuqallibaka ʿan jihāt al-ḍiddiyya ilā makman al-aḥadiyya wa-tuṣilaka
ilā sharīʿat al-qudsiyya li-tashraba ʿanhā wa-tastarīḥa nafsuka fīhā
wa-yaskuna ʿaṭashuka wa-yabruda fuʾāduka wa-takūna min alladhīnahum kānū
al-yawm bi-nūr Allāh la-muhtadīn.

O you who travel the paths of justice and gaze upon the countenance of
grace! Your letter has arrived, I have understood your question, and I
have heard the melodies of your heart within the pavilion of your inmost
being. I have therefore raised the clouds of the will to shower upon you
the rains of wisdom, that they may take from you all you have acquired
before, turn you from the realms of opposition toward the sanctuary of
oneness, and bring you to the shore of holiness, so that you may drink
from it, your soul may find rest there, your thirst may subside, your
heart may be refreshed, and you may be among those who are guided today
by the light of God.

\subsection{4}\label{section-2715}

\begin{otherlanguage}{arabic}

ولو انّی فی تلک الأیّام الّتی احاطتنی کلاب الأرض و سبع البلاد خفیت فی وکر
سرّی و اکون ممنوعاً عن اظهار ما اعطانی اللّه من بدائع علمه و جواهر حکمته و
شؤونات قدرته ولکن مع کلّ ذلک ما احبّ ان اخیّب من قام لدی حرم الکبریآء و
یرید ان یدخل فی رفرف البقآء و یحبّ ان یطیر فی سمآء هذا البدآء فی فجر
القضآء لذا اذکر لک بعض ما اکرمنی اللّه عمّا تطیقه النّفوس و تحمله العقول
لئلّا یرفع ضوضآء المبغضین و اعلام المنافقین و اسأل اللّه بأن یؤیّدنی بذلک
اذ هو ارحم الرّاحمین و معطی السّائلین

\end{otherlanguage}

Wa-law annī fī tilka al-ayyām allatī aḥāṭatnī kilāb al-arḍ wa-sabuʿ
al-bilād khafītu fī wakr sirrī wa-akūnu mamnūʿan ʿan iẓhār mā aʿṭānī
Allāh min badāʾiʿ ʿilmihi wa-jawāhir ḥikmatihi wa-shuʾūnāt qudratihi,
wa-lākin maʿa kull dhālika mā uḥibbu an ukhayyiba man qāma ladā ḥaram
al-kibriyāʾ wa-yurīdu an yadkhula fī rafraf al-baqāʾ wa-yuḥibbu an
yaṭīra fī samāʾ hādhā al-badāʾ fī fajr al-qaḍāʾ. Li-dhā adhkuru laka
baʿḍ mā akramanī Allāh ʿammā tuṭīquhu al-nufūs wa-taḥmiluhu al-ʿuqūl
li-allā yurfaʿa ḍawḍāʾ al-mubghiḍīn wa-aʿlām al-munāfiqīn. Wa-asʾalu
Allāh bi-an yuʾayyidanī bi-dhālika idh huwa arḥam al-rāḥimīn wa-muʿṭī
al-sāʾilīn.

Although in these days, when the dogs of the earth and the beasts of the
lands have surrounded Me, I have concealed Myself in My secret retreat
and am prevented from revealing the wonders of His knowledge, the gems
of His wisdom, and the manifestations of His power that God has bestowed
upon Me, nevertheless I do not wish to disappoint one who has stood
before the sanctuary of grandeur, desires to enter the pavilion of
eternity, and longs to soar in the heaven of this new beginning at the
dawn of decree. I will therefore recount to you something of what God
has graciously bestowed upon Me, within what souls can bear and minds
can sustain, lest the clamor of the malicious and the banners of the
hypocrites be raised. I ask God to assist Me in this, for He is the most
merciful of the merciful and the Giver to those who ask.

\subsection{5}\label{section-2716}

\begin{otherlanguage}{arabic}

فاعلم بأنّ لجنابک ینبغی بأن تفکّر فی اوّل الأمر بأنّ امم المختلفة الّذینهم
کانوا الیوم فی الأرض لم ما آمنوا برسل اللّه الّذین ارسلهم اللّه بقدرته و
اقامهم علی امره و جعلهم سراج ازلیّته فی مشکوة احدیّته و بم اعرضوا عنهم و
اختلفوا فیهم و خالفوا بهم و نازعوا معهم و حاربوا بهم و بأیّ جهة ما اقرّوا
برسالتهم و لا بولایتهم بل کفّروهم و سبّوهم حتّی قتلوهم و اخرجوهم

\end{otherlanguage}

Fa-iʿlam bi-anna li-janābika yanbaghī bi-an tufakkira fī awwal al-amr
bi-anna umam al-mukhtalifa alladhīnahum kānū al-yawm fī al-arḍ li-ma mā
āmanū bi-rusul Allāh alladhīna arsalahum Allāh bi-qudratihi wa-aqāmahum
ʿalā amrihi wa-jaʿalahum sirāj azaliyyatihi fī mishkāt aḥadiyyatihi,
wa-bi-ma aʿraḍū ʿanhum wa-khtalafū fīhim wa-khālafū bihim wa-nāzaʿū
maʿahum wa-ḥārabū bihim, wa-bi-ayyi jiha mā aqarrū bi-risālatihim wa-lā
bi-walāyatihim bal kaffarūhum wa-sabbūhum ḥattā qatalūhum wa-akhrajūhum.

Know that you should first consider why the various peoples now on earth
did not believe in the Messengers of God, whom God sent through His
power, raised up to establish His Cause, and made the lamp of His
eternity within the niche of His oneness. Why did they turn away from
them, disagree about them, oppose them, contend with them, and wage war
against them? For what reason did they acknowledge neither their mission
nor their authority, but instead denounced them as unbelievers and
reviled them, even killing and banishing them?

\subsection{6}\label{section-2717}

\begin{otherlanguage}{arabic}

و انّک یا ایّها الماشی فی بیدآء المعرفة و السّاکن فی سفینة الحکمة لو لا
تعرف سرّ ما ذکرناه لک ما تصل الی مراتب الایمان و لست بموقن فی امر اللّه و
مظاهر امره و مطالع حکمه و مخازن وحیه و معادن علمه و تکون من الّذین ما
جاهدوا فی امر اللّه و ما وجدوا رائحة الایمان من قمص الایقان و ما بلغوا
الی معارج التّوحید و ما وصلوا الی مدارج التّفرید فی هیاکل التّحمید و جواهر
التّجرید

\end{otherlanguage}

Wa-innaka yā ayyuhā al-māshī fī bayḍāʾ al-maʿrifa wa-l-sākin fī safīnat
al-ḥikma, law lā taʿrifu sirra mā dhakarnāhu laka mā taṣilu ilā marātib
al-īmān wa-lasta bi-mūqin fī amr Allāh wa-maẓāhir amrihi wa-maṭāliʿ
ḥikamihi wa-makhāzin waḥyihi wa-maʿādin ʿilmihi, wa-takūnu min alladhīna
mā jāhadū fī amr Allāh wa-mā wajadū rāʾiḥat al-īmān min qumuṣ al-īqān
wa-mā balaghū ilā maʿārij al-tawḥīd wa-mā waṣalū ilā madārij al-tafrīd
fī hayākil al-taḥmīd wa-jawāhir al-tajrīd.

O you who walk in the wilderness of knowledge and dwell in the ark of
wisdom! Unless you understand the mystery of what We have recounted to
you, you will not attain the degrees of faith or be assured of the Cause
of God, the Manifestations of His Cause, the Dawning-Places of His
wisdom, the Treasuries of His revelation, and the Mines of His
knowledge. You will be among those who have not striven in the Cause of
God, have not caught the fragrance of faith from the garments of
certitude, have not attained the ascending heights of unity, and have
not reached the stages of singleness within the embodiments of praise
and the essences of detachment.

\subsection{7}\label{section-2718}

\begin{otherlanguage}{arabic}

فاجهد یا اخی فی معرفة هذا المقام لیکشف الغطآء عن وجه قلبک و تکون من
الّذین جعل اللّه بصرهم حدیداً لتشهد جراثیم الجبروت و تطّلع بأسرار الملکوت و
رموزات الهویّة فی اراضی النّاسوت و تصل الی مقام الّذی ما تری فی خلق الرّحمن
من تفاوت و لا فی خلق السّموات و الأرض من فطور

\end{otherlanguage}

Fa-jhad yā akhī fī maʿrifat hādhā al-maqām li-yukshafa al-ghiṭāʾ ʿan
wajh qalbika wa-takūna min alladhīna jaʿala Allāh baṣarahum ḥadīdan
li-tashhada jarāthīm al-jabarūt wa-taṭṭaliʿa bi-asrār al-malakūt
wa-rumūzāt al-huwiyya fī arāḍī al-nāsūt wa-taṣila ilā maqām alladhī mā
tarā fī khalq al-Raḥmān min tafāwut wa-lā fī khalq al-samāwāt wa-l-arḍ
min fuṭūr.

Strive, then, My brother, to understand this station, so that the veil
may be lifted from the face of your heart and you may be among those
whose sight God has made keen, that you may behold the primal realities
of the realm of might, discover the mysteries of the kingdom and the
symbols of the divine Essence in the lands of humanity, and attain the
station in which you see no disparity in the creation of the
All-Merciful and no flaw in the creation of the heavens and the earth.

\subsection{8}\label{section-2719}

\begin{otherlanguage}{arabic}

فلمّا بلغ الأمر الی هذا المقام الأوعر الأعلی و هذا الرّمز الخشن الأسنی
فاعرف بأنّ هؤلآء الأمم من الیهود و النّصاری لمّا ما عرفوا لحن القول و ما
بلغوا الی ما وعدهم اللّه فی کتابه انکروا امر اللّه و اعرضوا عن رسل اللّه و
انکروا حجج اللّه و انّهم لو کانوا ناظرین الی الحجّة بنفسها و ما اتّبعوا کلّ
همج رعاع من علمائهم و رؤسائهم لبلغوا الی مخزن الهدی و مکمن التّقی و شربوا
من مآء الحیّ الحیوان فی مدینة الرّحمن و حدیقة السّبحان و حقیقة الرّضوان و
انّهم لمّا ما شهدوا الحجّة بعیونهم الّتی خلق اللّه لهم بهم و ارادوا بغیر ما
اراد اللّه لهم من فضله بعدوا عن رفرف القرب و منعوا عن کوثر الوصل و منبع
الفضل و کانوا فی حجبات انفسهم میّتین

\end{otherlanguage}

Fa-lammā balagha al-amru ilā hādhā al-maqām al-awʿar al-aʿlā wa-hādhā
al-ramz al-khashin al-asnā fa-iʿraf bi-anna hāʾulāʾ al-umam min al-yahūd
wa-l-naṣārā lammā mā ʿarafū laḥn al-qawl wa-mā balaghū ilā mā waʿadahum
Allāh fī kitābihi ankarū amr Allāh wa-aʿraḍū ʿan rusul Allāh wa-ankarū
ḥujaj Allāh. Wa-innahum law kānū nāẓirīn ilā al-ḥujja bi-nafsihā wa-mā
ittabaʿū kulla hamaj raʿāʿ min ʿulamāʾihim wa-ruʾasāʾihim la-balaghū ilā
makhzan al-hudā wa-makman al-tuqā wa-sharibū min māʾ al-ḥayy al-ḥayawān
fī madīnat al-Raḥmān wa-ḥadīqat al-Subḥān wa-ḥaqīqat al-riḍwān.
Wa-innahum lammā mā shahidū al-ḥujja bi-ʿuyūnihim allatī khalaqa Allāh
lahum bihim wa-arādū bi-ghayr mā arāda Allāh lahum min faḍlihi baʿudū
ʿan rafraf al-qurb wa-muniʿū ʿan kawthar al-waṣl wa-manbaʿ al-faḍl
wa-kānū fī ḥujubāt anfusihim mayyitīn.

Now that the discussion has reached this most lofty and difficult
station and this most radiant yet formidable mystery, know that these
peoples, the Jews and the Christians, denied the Cause of God, turned
away from His Messengers, and rejected His proofs because they did not
understand the purport of the words or attain what God had promised them
in His Book. Had they looked at the proof itself and not followed every
ignorant rabble-rouser among their scholars and leaders, they would have
reached the treasury of guidance and the sanctuary of piety, and drunk
of the living, life-giving water in the city of the All-Merciful, the
garden of the All-Glorious, and the reality of paradise. But because
they did not behold the proof with their own eyes, which God had created
for them through those Messengers, and desired something other than what
God, through His grace, desired for them, they became far removed from
the pavilion of nearness, were barred from the river of reunion and the
wellspring of grace, and lay dead within the veils of their own selves.

\subsection{9}\label{section-2720}

\begin{otherlanguage}{arabic}

و انّی بحول اللّه و قوّته حینئذ اذکر بعض ما ذکره اللّه فی کتب القبل و علائم
ظهورات الأحدیّة فی هیاکل الأنزعیّة لتعرف مقام الفجر فی هذا الصّبح الأزلیّة و
تشاهد هذه النّار المشتعلة فی سدرة لا شرقیّة و لا غربیّة و تفتح عیناک فی
وصولک الی مولاک و یمذق قلبک من نعمآء المکنونة فی هذه الأوعیة المخزونة و
تشکر اللّه ربّک فیما اختصّک بذلک و جعلک من الّذینهم کانوا بلقآء ربّهم موقنون

\end{otherlanguage}

Wa-innī bi-ḥawl Allāh wa-quwwatihi ḥīnaʾidh adhkuru baʿḍ mā dhakarahu
Allāh fī kutub al-qabl wa-ʿalāʾim ẓuhūrāt al-aḥadiyya fī hayākil
al-anzaʿiyya li-taʿrifa maqām al-fajr fī hādhā al-ṣubḥ al-azaliyya
wa-tushāhida hādhihi al-nār al-mushtaʿila fī sidra lā sharqiyya wa-lā
gharbiyya wa-taftaḥa ʿaynāka fī wuṣūlika ilā mawlāka wa-yamdhaqa qalbuka
min naʿmāʾ al-maknūna fī hādhihi al-awʿiya al-makhzūna wa-tashkura
Allāha rabbaka fīmā ikhtaṣṣaka bi-dhālika wa-jaʿalaka min alladhīnahum
kānū bi-liqāʾ rabbihim mūqinūn.

Now, through the power and strength of God, I will recount some of what
God mentioned in the books of old, and the signs of the revelations of
oneness in the embodiments of detachment, so that you may recognize the
station of daybreak in this eternal morning, behold this fire blazing in
a tree neither of the East nor of the West, open your eyes upon
attaining your Lord, and let your heart partake of the blessings
concealed in these hidden vessels. Thus may you thank God, your Lord,
for singling you out for this and making you one of those who are
certain of meeting their Lord.

\subsection{10}\label{section-2721}

\begin{otherlanguage}{arabic}

هذا صورة ما نزل من قبل فی انجیل المتّی فی سفر الأوّل فیه یذکر علائم ظهور
الّذی یأتی بعده و یقول الویل للحبالی و المرضعات فی تلک الأیّام الی ان تغنّ
الورقآء فی قطب البقآء و یدلع دیک العرش فی شجرة القصوی و سدرة المنتهی و
یقول و للوقت من بعد ضیق تلک الأیّام تظلم الشّمس و القمر لا یعطی ضوئه و
الکواکب تتساقط من السّمآء و قوّات السّمآء ترتجّ حینئذ یظهر علامة ابن الانسان
فی السّمآء و ینوح حینئذ کلّ قبائل الأرض و یرون ابن الانسان آتیاً علی سحاب
السّمآء مع قوّات و مجد کبیر و یرسل ملائکته مع صوت السّافور العظیم انتهی

\end{otherlanguage}

Hādhā ṣūrat mā nazala min qabl fī Injīl al-Mattā fī sifr al-awwal, fīhi
yadhkuru ʿalāʾim ẓuhūr alladhī yaʾtī baʿdahu wa-yaqūlu al-wayl
li-l-ḥabālā wa-l-murḍiʿāt fī tilka al-ayyām, ilā an tughanna al-Warqāʾ
fī quṭb al-baqāʾ wa-yadlaʿa dīk al-ʿarsh fī shajarat al-quṣwā wa-sidrat
al-muntahā wa-yaqūlu wa-li-l-waqt min baʿd ḍīq tilka al-ayyām tuẓlamu
al-shams wa-l-qamar lā yuʿṭī ḍawʾahu wa-l-kawākib tatasāqaṭu min
al-samāʾ wa-quwwāt al-samāʾ tartajju, ḥīnaʾidh yaẓharu ʿalāmat ibn
al-insān fī al-samāʾ wa-yanūḥu ḥīnaʾidh kull qabāʾil al-arḍ wa-yarawna
ibn al-insān ātiyan ʿalā saḥāb al-samāʾ maʿa quwwāt wa-majd kabīr
wa-yursilu malāʾikatahu maʿa ṣawt al-sāfūr al-ʿaẓīm. Intahā.

This is the text of what was revealed before in the Gospel of Matthew,
in the first book, where He mentions the signs of the appearance of the
One who will come after Him and says: ``Woe to pregnant women and
nursing mothers in those days,'' until the Dove sings at the center of
eternity and the Cockerel of the Throne raises its call in the farthest
tree, the Lote-Tree of the Utmost Boundary, and says: ``Immediately
after the distress of those days, the sun will be darkened, the moon
will not give its light, the stars will fall from heaven, and the powers
of heaven will be shaken. Then the sign of the Son of Man will appear in
heaven, and then all the tribes of the earth will mourn, and they will
see the Son of Man coming upon the clouds of heaven with power and great
glory. And He will send His angels with the sound of the mighty
trumpet.'' Here ends the quotation.

\subsection{11}\label{section-2722}

\begin{otherlanguage}{arabic}

و فی سفر الثّانی فی انجیل المرقس فیما یتکلّم حمامة القدس فیقول بأنّ فی تلک
الأیّام ضیق لم یکن مثله من البدء الّذی خلق اللّه الی الآن و لا یکون انتهی و
بعد ترنّ بمثل ما رنّت من قبل من دون تغییر و لا تبدیل و کان اللّه علی ما
اقول وکیل

\end{otherlanguage}

Wa-fī sifr al-thānī fī Injīl al-Marqus fīmā yatakallamu ḥamāmat al-quds
fa-yaqūlu bi-anna fī tilka al-ayyām ḍīq lam yakun mithluhu min al-badʾ
alladhī khalaqa Allāh ilā al-ān wa-lā yakūn. Intahā. Wa-baʿd tarinnu
bi-mithl mā rannat min qabl min dūn taghyīr wa-lā tabdīl wa-kāna Allāhu
ʿalā mā aqūlu wakīl.

In the second book, in the Gospel of Mark, where the Dove of holiness
speaks, it says: ``In those days there will be distress such as has not
occurred from the beginning of what God created until now, and never
will.'' Here ends the quotation. After this, it sings as it sang before,
without change or alteration. God is the witness to what I say.

\subsection{12}\label{section-2723}

\begin{otherlanguage}{arabic}

و فی سفر الثّالث فی انجیل اللّوقا یقول علامات فی الشّمس و القمر و النّجوم و
تحدث علی الأرض ضیق الأمم من هول صوت البحر و الزّلازل و قوّات السّمآء یضطرب
و ینظرون ابن الانسان آتیاً فی السّحاب مع قوّات و مجد عظیم و اذا رأیتم هذا
کلّه کائناً اعلموا انّ ملکوت اللّه قد اقتربت انتهی

\end{otherlanguage}

Wa-fī sifr al-thālith fī Injīl al-Lūqā yaqūlu ʿalāmāt fī al-shams
wa-l-qamar wa-l-nujūm wa-taḥduthu ʿalā al-arḍ ḍīq al-umam min hawl ṣawt
al-baḥr wa-l-zalāzil wa-quwwāt al-samāʾ yaḍṭaribu wa-yanẓurūna ibn
al-insān ātiyan fī al-saḥāb maʿa quwwāt wa-majd ʿaẓīm wa-idhā raʾaytum
hādhā kullahu kāʾinan iʿlamū anna malakūt Allāh qad iqtarabat. Intahā.

In the third book, in the Gospel of Luke, He says: ``There will be signs
in the sun, the moon, and the stars, and on earth distress among the
nations at the terror of the sea's roar and the earthquakes, and the
powers of heaven will be shaken. And they will see the Son of Man coming
in the clouds with power and great glory. When you see all this taking
place, know that the kingdom of God has drawn near.'' Here ends the
quotation.

\subsection{13}\label{section-2724}

\begin{otherlanguage}{persian}

و فی سفر الرّابع فی انجیل الیوحنّا یقول اذا جآء المعزّی الّذی ارسله الیکم
روح الحقّ الآتی من الحقّ فهو یشهد لی و انتم تشهدون و فی مقام آخر یقول و
اذا جآء روح القدس المعزّی الّذی یرسله ربّی باسمی فهو یعلّمکم کلّ شیء و یذکرکم
کلّ ما قلت لکم و الآن فانّی منطلق الی من ارسلنی و لیس احد منکم یسألنی الی
این اذهب لأنّی قلت لکم هذا و فی مقام آخر یقول انّی اقول لکم الحقّ انّه خیر
لکم ان انطلق لأنّی ان لم انطلق لم یأتکم المعزّی فاذا انطلقت ارسلنه الیکم
فاذا جآء روح الحقّ ذاک فهو یرشدکم الی جمیع الحقّ لأنّه لیس ینطق من عنده بل
یتکلّم بما یسمع و یخبرکم بما یأتی

\end{otherlanguage}

Wa-fī sifr al-rābiʿ fī Injīl al-Yūḥannā yaqūlu idhā jāʾa al-Muʿazzī
alladhī ursiluhu ilaykum Rūḥ al-Ḥaqq al-ātī min al-Ḥaqq fa-huwa yashhadu
lī wa-antum tashhadūn. Wa-fī maqām ākhar yaqūlu wa-idhā jāʾa Rūḥ al-Quds
al-Muʿazzī alladhī yursiluhu rabbī bi-ismī fa-huwa yuʿallimukum kulla
shayʾ wa-yudhakkirukum kulla mā qultu lakum wa-l-āna fa-innī munṭaliq
ilā man arsalanī wa-laysa aḥad minkum yasʾalunī ilā ayna adhhabu li-annī
qultu lakum hādhā. Wa-fī maqām ākhar yaqūlu innī aqūlu lakum al-ḥaqq
innahu khayr lakum an anṭaliqa li-annī in lam anṭaliq lam yaʾtikum
al-Muʿazzī fa-idhā inṭalaqtu ursilannahu ilaykum fa-idhā jāʾa Rūḥ
al-Ḥaqq dhāka fa-huwa yurshidukum ilā jamīʿ al-ḥaqq li-annahu laysa
yanṭiqu min ʿindihi bal yatakallamu bi-mā yasmaʿu wa-yukhbirukum bi-mā
yaʾtī.

In the fourth book, in the Gospel of John, He says: ``When the Comforter
comes, whom I will send to you, the Spirit of Truth who comes from the
Truth, He will testify about Me, and you too will testify.'' In another
passage He says: ``When the Holy Spirit, the Comforter whom My Lord will
send in My name, comes, He will teach you everything and remind you of
everything I have told you. Now I am going to the One who sent Me, and
none of you asks Me, `Where are You going?' because I have told you
this.'' In another passage He says: ``I tell you the truth: it is better
for you that I go away, for if I do not go away, the Comforter will not
come to you. But if I go, I will surely send Him to you. When that
Spirit of Truth comes, He will guide you into all truth, for He will not
speak on His own authority, but will speak what He hears and tell you
what is to come.''

\subsection{14}\label{section-2725}

\begin{otherlanguage}{arabic}

هذا صورة ما نزل من قبل و انّی فواللّه الّذی لا اله الّا هو لاختصرت و لو ارید
ان اذکر کلمات الأنبیآء فیما نزل من جبروت العظمة و ملکوت السّلطنة علیهم
لتملأ الأوراق و الألواح من قبل ان اصل الی آخرها و فی کلّ الزّبرات و
المزامیر و الصّحائف لموجود و مذکور بمثل ما ذکرت لک و القیت علیک بل اعلی و
اعظم عن کلّ ما ذکرت و فصّلت و انّی لو ارید ان اذکر کلّ ما نزّل من قبل لأقدر
بما اعطانی اللّه من بدائع علمه و قدرته ولکن اکتفیت بما بیّنت لک لئلّا تکسل
فی سفرک و لا تنقلب علی عقبیک و لئلّا یأخذک من حزن و لا کدورة و لا من نصب
و لا من ذلّ و لا من لغوب

\end{otherlanguage}

Hādhā ṣūrat mā nazala min qabl wa-innī fa-wa-Allāh alladhī lā ilāha illā
huwa la-khtaṣartu. Wa-law urīdu an adhkura kalimāt al-anbiyāʾ fīmā
nazala min jabarūt al-ʿaẓama wa-malakūt al-salṭana ʿalayhim la-tumlaʾu
al-awrāq wa-l-alwāḥ min qabl an aṣila ilā ākhirihā. Wa-fī kull
al-zuburāt wa-l-mazāmīr wa-l-ṣaḥāʾif la-mawjūd wa-madhkūr bi-mithl mā
dhakartu laka wa-alqaytu ʿalayka bal aʿlā wa-aʿẓam ʿan kull mā dhakartu
wa-faṣṣaltu. Wa-innī law urīdu an adhkura kulla mā nuzzila min qabl
la-aqdiru bi-mā aʿṭānī Allāh min badāʾiʿ ʿilmihi wa-qudratihi wa-lākin
iktafaytu bi-mā bayyantu laka li-allā taksala fī safarika wa-lā
tanqaliba ʿalā ʿaqibayka wa-li-allā yaʾkhudhaka min ḥuzn wa-lā kudūra
wa-lā min naṣab wa-lā min dhull wa-lā min lughūb.

This is the text of what was revealed before, and, by God, besides whom
there is no other god, I have indeed been brief. Were I to recount the
words of the Prophets among what was revealed to them from the dominion
of grandeur and the kingdom of sovereignty, pages and tablets would be
filled before I reached their end. In all the scriptures, psalms, and
sacred writings there is to be found and mentioned the like of what I
have recounted and imparted to you---indeed, what is loftier and greater
than everything I have mentioned and explained. Were I to wish to
recount all that was revealed before, I could indeed do so through the
wonders of His knowledge and power that God has bestowed upon Me. But I
have contented Myself with what I have explained to you, lest you grow
weary on your journey and turn back on your heels, and lest sorrow,
gloom, hardship, humiliation, or exhaustion overtake you.

\subsection{15}\label{section-2726}

\begin{otherlanguage}{arabic}

اذاً فأنصف ثمّ فکّر فی تلک العبارات المتعالیات ثمّ اسأل عن الّذین یدّعون العلم
من دون بیّنة من عند اللّه و لا حجّة من لدنه و غفلوا عن تلک الأیّام الّتی
اشرقت شمس العلم و الحکمة عن افق الألوهیّة و تعطی کلّ ذی حقّ حقّه و کلّ ذی قدر
مقداره و مقامه ما یقولون فی هذه الاشارات الّتی ذهلت العقول عن ادراکها و
حارت النّفوس المقدّسة عن عرفان ما ستر فیها من حکمة اللّه البالغة و علم اللّه
المودعة

\end{otherlanguage}

Idhan fa-anṣif thumma fakkir fī tilka al-ʿibārāt al-mutaʿāliyāt thumma
isʾal ʿan alladhīna yaddaʿūna al-ʿilm min dūn bayyina min ʿind Allāh
wa-lā ḥujja min ladunhi wa-ghafalū ʿan tilka al-ayyām allatī ashraqat
shams al-ʿilm wa-l-ḥikma ʿan ufuq al-ulūhiyya wa-tuʿṭī kulla dhī ḥaqq
ḥaqqahu wa-kulla dhī qadr miqdārahu wa-maqāmahu mā yaqūlūna fī hādhihi
al-ishārāt allatī dhahilat al-ʿuqūl ʿan idrākihā wa-ḥārat al-nufūs
al-muqaddasa ʿan ʿirfān mā sutira fīhā min ḥikmat Allāh al-bāligha
wa-ʿilm Allāh al-mūdaʿa.

Be fair, then, and reflect on these exalted utterances. Then ask those
who claim knowledge without any clear evidence from God or any proof
from Him, and who have remained heedless of these days in which the sun
of knowledge and wisdom has risen from the horizon of divinity and gives
everyone their due and everyone of rank their proper measure and
station: what do they say about these allusions, which minds are
bewildered to comprehend and holy souls are perplexed to fathom, with
the consummate wisdom of God and the knowledge of God concealed within
them?

\subsection{16}\label{section-2727}

\begin{otherlanguage}{persian}

ان یقولون هذه الکلمات من عند اللّه و لم یکن لها من تأویل و تکون علی ظاهر
القول فی ظاهر الظّاهر فکیف یعترضون علی هؤلآء الکفرة من اهل الکتاب لأنّهم
لمّا شهدوا فی کتابهم ما ذکرناه لک و فسّروا لهم علمائهم علی ظاهر القول لذا
ما اقرّوا باللّه فی مظاهر التّوحید و مطالع التّفرید و هیاکل التّجرید و ما
آمنوا بهم و ما اطاعوهم لأنّهم ما شهدوا بأن تظلم الشّمس و تساقط الکواکب من
السّمآء علی وجه الأرض و تنزلنّ الملائکة علی ظاهر الهیکل علی الأرض لذا
اعترضوا علی النّبیّین و المرسلین بل لمّا وجدوهم مخالفاً لدینهم و شرائعهم
وردوا علیهم ما استحیی ان اذکر لک من الکذب و الجنون و الکفر و الضّلال
فأرجع البصر فی القرآن لتجد کلّ ذلک و تکون فیه من العارفین و من یومئذ الی
حینئذ ینتظرون هذه الفئة ظهورات ما عرفوا من علمائهم و ایقنوا من فقهائهم و
یقولون متی تظهر هذه العلامات انّا حینئذ لآمنون و لو کان الأمر کذلک کیف
انتم تدحضون حجّتهم و تبطلون برهانهم و تحتجّون بهم فی امر دینهم و ما عرفوا
من کتبهم و سمعوا من صنادیدهم

\end{otherlanguage}

In yaqūlūn hādhihi al-kalimāt min ʿind Allāh wa-lam yakun lahā min
taʾwīl wa-takūn ʿalā ẓāhir al-qawl fī ẓāhir al-ẓāhir fa-kayfa yaʿtariḍūn
ʿalā hāʾulāʾ al-kafara min ahl al-kitāb li-annahum lammā shahidū fī
kitābihim mā dhakarnāhu laka wa-fassara lahum ʿulamāʾuhum ʿalā ẓāhir
al-qawl li-dhā mā aqarrū bi-Allāh fī maẓāhir al-tawḥīd wa-maṭāliʿ
al-tafrīd wa-hayākil al-tajrīd wa-mā āmanū bihim wa-mā aṭāʿūhum
li-annahum mā shahidū bi-an tuẓlim al-shams wa-tasāqaṭ al-kawākib min
al-samāʾ ʿalā wajh al-arḍ wa-tanzilanna al-malāʾika ʿalā ẓāhir al-haykal
ʿalā al-arḍ li-dhā iʿtaraḍū ʿalā al-nabiyyīn wa-l-mursalīn bal lammā
wajadūhum mukhālifan li-dīnihim wa-sharāʾiʿihim waradū ʿalayhim mā
astaḥyī an adhkura laka min al-kadhib wa-l-junūn wa-l-kufr wa-l-ḍalāl
fa-arjiʿ al-baṣar fī al-Qurʾān li-tajida kull dhālika wa-takūn fīhi min
al-ʿārifīn wa-min yawmiʾidh ilā ḥīniʾidh yantaẓirūn hādhihi al-fiʾa
ẓuhūrāt mā ʿarafū min ʿulamāʾihim wa-ayqanū min fuqahāʾihim wa-yaqūlūn
matā taẓhar hādhihi al-ʿalāmāt innā ḥīnaʾidh la-āminūn wa-law kāna
al-amr ka-dhālika kayfa antum tudḥiḍūn ḥujjatahum wa-tubṭilūn burhānahum
wa-taḥtajjūn bihim fī amr dīnihim wa-mā ʿarafū min kutubihim wa-samiʿū
min ṣanādīdihim

If they say that these words are from God, admit of no interpretation,
and are to be taken according to the outward wording in its most literal
sense, how can they object to these unbelievers among the people of the
Book? For when those people saw in their Book what We have recounted to
you, and their learned men explained it to them according to the outward
wording, they therefore did not acknowledge God in the Manifestations of
oneness, the Dawning-places of singleness, and the Temples of
detachment. They neither believed in them nor obeyed them, because they
had not seen the sun darkened, the stars falling from heaven onto the
face of the earth, or the angels descending to earth in outward bodily
form. Thus they opposed the Prophets and Messengers. Indeed, when they
found them at variance with their religion and their laws, they assailed
them with charges of falsehood, madness, unbelief, and error that I am
ashamed to recount to you. Look again in the Qurʾan, and you will find
all this and understand it. From that day to this, this people have
awaited the appearance of what they have learned from their learned men
and accepted with certainty from their jurists, saying, ``When these
signs appear, we will surely believe.'' If the matter were so, how could
you refute their argument, invalidate their proof, and contend with them
concerning their religion, what they have learned from their books, and
what they have heard from their leaders?

\subsection{17}\label{section-2728}

\begin{otherlanguage}{persian}

و ان یقولون هذه الأسفار الّتی تکون بین یدی هذه الفئة و یسمّونها بالانجیل و
ینسبونها بعیسی بن مریم ما نزلت من عند اللّه و مظاهر نفسه یلزم تعطیل الفیض
عن مبدإ الفیّاض و لم تکن الحجّة من عند اللّه بالغة علی عباده و لم تکن
النّعمة کاملة و لا العنایة مشرقة و لا الرّحمة واسعة لأنّه لمّا رفع عیسی الی
السّمآء و رفع کتابه فبأیّ شیء یحتجّ اللّه بهم یوم القیامة و یعذّبهم کما هو
المکتوب من ائمّة الدّین و المنصوص من علمآء الرّاشدین

\end{otherlanguage}

Wa-in yaqūlūn hādhihi al-asfār allatī takūn bayna yaday hādhihi al-fiʾa
wa-yusammūnahā bi-l-Injīl wa-yansibūnahā bi-ʿĪsā ibn Maryam mā nazalat
min ʿind Allāh wa-maẓāhir nafsihi yalzam taʿṭīl al-fayḍ ʿan mabdaʾ
al-fayyāḍ wa-lam takun al-ḥujja min ʿind Allāh bāligha ʿalā ʿibādihi
wa-lam takun al-niʿma kāmila wa-lā al-ʿināya mushriqa wa-lā al-raḥma
wāsiʿa li-annahu lammā rufiʿa ʿĪsā ilā al-samāʾ wa-rufiʿa kitābuhu
fa-bi-ayy shayʾ yaḥtajj Allāh bihim yawm al-qiyāma wa-yuʿadhdhibuhum
ka-mā huwa al-maktūb min aʾimmat al-dīn wa-l-manṣūṣ min ʿulamāʾ
al-rāshidīn

And if they say that these scriptures, which this people possess, which
they call the Gospel and attribute to Jesus son of Mary, did not descend
from God and the Manifestations of His own Self, it follows that the
flow of grace from its all-bountiful Source was interrupted, that God's
proof to His servants was not conclusive, that His blessing was not
complete, His providence not radiant, and His mercy not all-embracing.
For if Jesus was raised to heaven and His Book was taken up, by what
could God establish His case against them on the Day of Resurrection and
punish them, as the leaders of religion have written and rightly guided
scholars have expressly stated?

\subsection{18}\label{section-2729}

\begin{otherlanguage}{persian}

اذاً فکّر فی نفسک لمّا تشهد الأمر کذلک و نشهد کذلک من این تفرّ و الی من ترکض
و الی من تتوجّه و بأیّ ارض تسکن و بأیّ فراش تجلس و بأیّ صراط تستقیم و بأیّ
ساعة تنوم و بأیّ امر تنتهی امرک و بأیّ شیء تشدّ عروة دینک و حبل طاعتک لا
فوالّذی تجلّی بالوحدانیّة و تشهد لنفسه بالفردانیّة لو یحدث فی قلبک قبساً من
نار محبّة اللّه ما تنوم و ما تسکن و ما تضحک و ما تستریح بل تفرّ الی قلل
الجبال فی ساحة القرب و القدس و الجمال و تنوح کنوح الفاقدین و تبکی کبکآء
المشتاقین و لا ترجع الی بیتک و محلّک الّا بأن یکشف اللّه لک امره

\end{otherlanguage}

Idhan fakkir fī nafsika lammā tashhad al-amr ka-dhālika wa-nashhad
ka-dhālika min ayna tafirr wa-ilā man tarkuḍ wa-ilā man tatawajjah
wa-bi-ayy arḍ taskun wa-bi-ayy firāsh tajlis wa-bi-ayy ṣirāṭ tastaqīm
wa-bi-ayy sāʿa tanūm wa-bi-ayy amr tantahī amrak wa-bi-ayy shayʾ tashudd
ʿurwat dīnika wa-ḥabl ṭāʿatika lā fa-wa-alladhī tajallā bi-l-waḥdāniyya
wa-tashhad li-nafsihi bi-l-fardāniyya law yaḥduth fī qalbika qabasan min
nār maḥabbat Allāh mā tanūm wa-mā taskun wa-mā taḍḥak wa-mā tastarīḥ bal
tafirr ilā qulal al-jibāl fī sāḥat al-qurb wa-l-quds wa-l-jamāl wa-tanūḥ
ka-nawḥ al-fāqidīn wa-tabkī ka-bukāʾ al-mushtāqīn wa-lā tarjiʿ ilā
baytika wa-maḥallika illā bi-an yakshif Allāh laka amrahu

Reflect within yourself, then: when you see the matter to be so, as We
also see it, where will you flee, to whom will you hasten, and toward
whom will you turn? In what land will you dwell, on what resting place
will you sit, and on what path will you stand firm? At what hour will
you sleep, by what means will you bring your affair to its conclusion,
and with what will you secure the handle of your religion and the cord
of your obedience? No, by Him who has revealed Himself in oneness and to
whose singleness you bear witness! If a spark of the fire of God's love
were to arise in your heart, you would neither sleep nor remain still,
neither laugh nor rest. Rather, you would flee to the mountain peaks in
the court of nearness, holiness, and beauty. You would lament as the
bereaved lament and weep as those who yearn weep, and you would not
return to your home and dwelling until God disclosed His Cause to you.

\subsection{19}\label{section-2730}

\begin{otherlanguage}{persian}

و انّک انت یا ایّها المتعارج الی جبروت الهدی و المتصاعد الی ملکوت التّقی لو
ترید ان تعرف هذه الاشارات القدسیّة و تشهد اسرار العلمیّة و تطّلع علی کلمة
الجامعة لا بدّ لجنابک ان تسأل کلّ ذلک و کلّ ما یرد علیک فی امر مبدئک و
معادک عن الّذین جعلهم اللّه منبع علمه و سمآء حکمته و سفینة سرّه لأنّ من دون
هذه الأنوار المشرقة عن افق الهویّة ما یعرفون النّاس یمینهم عن شمائلهم و
کیف یقدرنّ ان یتعارجنّ الی افق الحقائق او یصلنّ الی مخزن الدّقائق اذاً نسأل
اللّه بأن یدخلنا فی هذه البحور المتموّجة و یشرّفنا الی هذه الأرواح المرشّحة
و ینزلنا فی هذه المعارج الالهیّة لننزع عن هیاکلنا کلّ ما اخذنا من عند
انفسنا و نخلع عن اجسادنا کلّ الأثواب العاریة الّتی سرقنا عن امثالنا
لیلبسنا اللّه من قمص عنایته و اثواب هدایته

\end{otherlanguage}

Wa-innaka anta yā ayyuhā al-mutaʿārij ilā jabarūt al-hudā wa-l-mutaṣāʿid
ilā malakūt al-tuqā law turīd an taʿrif hādhihi al-ishārāt al-qudsiyya
wa-tashhad asrār al-ʿilmiyya wa-taṭṭaliʿ ʿalā kalimat al-jāmiʿa lā budda
li-janābika an tasʾal kull dhālika wa-kull mā yarid ʿalayka fī amr
mabdaʾika wa-maʿādika ʿan alladhīna jaʿalahum Allāh manbaʿ ʿilmihi
wa-samāʾ ḥikmatihi wa-safīnat sirrihi li-anna min dūn hādhihi al-anwār
al-mushriqa ʿan ufuq al-huwiyya mā yaʿrifūn al-nās yamīnahum ʿan
shamāʾilihim wa-kayfa yaqdiranna an yataʿārajna ilā ufuq al-ḥaqāʾiq aw
yaṣilna ilā makhzan al-daqāʾiq idhan nasʾal Allāh bi-an yudkhilanā fī
hādhihi al-biḥār al-mutamawwija wa-yusharrifanā ilā hādhihi al-arwāḥ
al-murashshaḥa wa-yunzilanā fī hādhihi al-maʿārij al-ilāhiyya li-nanzaʿ
ʿan hayākilinā kull mā akhadhnā min ʿind anfusinā wa-nakhlaʿ ʿan
ajsādinā kull al-athwāb al-ʿāriya allatī saraqnā ʿan amthālinā
li-yulbisanā Allāh min qumuṣ ʿināyatihi wa-athwāb hidāyatihi

As for you, who ascend to the dominion of guidance and rise to the
kingdom of piety: if you wish to understand these sacred allusions,
behold the mysteries of knowledge, and become acquainted with the
all-embracing Word, you must ask about all this, and about everything
that comes before you concerning your origin and your return, of those
whom God has made the wellspring of His knowledge, the heaven of His
wisdom, and the ark of His mystery. For without these lights shining
from the horizon of the divine Essence, people cannot tell their right
from their left. How, then, could they ascend to the horizon of
realities or reach the treasury of subtleties? We therefore ask God to
admit us into these surging seas, to honor us with access to these
outpouring spirits, and to establish us on these divine ascents, so that
we may strip from our persons everything we have acquired from our own
selves and remove from our bodies all the borrowed garments we have
stolen from others like ourselves, that God may clothe us in the robes
of His providence and the garments of His guidance.

\subsection{20}\label{section-2731}

\begin{otherlanguage}{arabic}

و یدخلنا فی مدینة العلم الّذی من دخل فیها لیعرف کلّ العلوم قبل ان یلتفت
الی اسرارها و یعرف کلّ العلم و الحکمة من اسرار الرّبوبیّة المودعة فی کنائز
الخلیقة من اوراقها الّتی تورّقت من اشجارها فسبحان اللّه موجدها و مبدعها عمّا
خلق فیها و قدّر لها و انّی فواللّه المهیمن المقتدر القیّوم لو ارینّک ابواب
هذه المدینة الّتی خلقت عن یمین القدرة و القوّة لتری ما لا رأی احد من قبلک
و تشهد ما لا شهدت نفس دونک و تعرف غوامض الدّلالات و معضلات الاشارات و
تبرهن لک اسرار البدئیّة فی نقطة الختمیّة و تسهل علیک الأمور و تجعل النّار
لک نوراً و علماً و رحمةً و تکون فی بساط القدس لمن المستریحین

\end{otherlanguage}

Wa-yudkhilanā fī madīnat al-ʿilm alladhī man dakhala fīhā la-yaʿrif kull
al-ʿulūm qabl an yaltafit ilā asrārihā wa-yaʿrif kull al-ʿilm wa-l-ḥikma
min asrār al-rubūbiyya al-mūdaʿa fī kanāʾiz al-khalīqa min awrāqihā
allatī tawarraqat min ashjārihā fa-subḥān Allāh mūjidihā wa-mubdiʿihā
ʿammā khalaqa fīhā wa-qaddara lahā wa-innī fa-wa-Allāh al-Muhaymin
al-Muqtadir al-Qayyūm law uriyannaka abwāb hādhihi al-madīna allatī
khuliqat ʿan yamīn al-qudra wa-l-quwwa la-tarā mā lā raʾā aḥad min
qablika wa-tashhad mā lā shahidat nafs dūnaka wa-taʿrif ghawāmiḍ
al-dalālāt wa-muʿḍilāt al-ishārāt wa-tubarhan laka asrār al-badʾiyya fī
nuqṭat al-khatmiyya wa-tashul ʿalayka al-umūr wa-tujʿal al-nār laka
nūran wa-ʿilman wa-raḥmatan wa-takūn fī bisāṭ al-quds la-min
al-mustarīḥīn

And may He admit us into the City of Knowledge, where whoever enters
will know all sciences before turning toward its mysteries, and will
learn all knowledge and wisdom concerning the mysteries of lordship
deposited in the treasuries of creation from the leaves that have
sprouted from its trees. Exalted is God, who brought it into being and
originated it, above all He has created within it and ordained for it!
By God, the All-Protecting, the All-Powerful, the Self-Subsisting! If I
were to show you the gates of this City, which was created by the right
hand of power and might, you would see what no one before you has seen,
behold what no soul other than you has beheld, and understand the
obscure meanings of its indications and the perplexing subtleties of its
allusions. The mysteries of the beginning would be demonstrated to you
in the point of the end, all things would become easy for you, fire
would be made light, knowledge, and mercy for you, and you would be
among those at rest upon the expanse of holiness.

\subsection{21}\label{section-2732}

\begin{otherlanguage}{arabic}

و من دون ذلک و کلّ ما القیناک من جواهر اسرار الحکمة فی غیاهب هذه الکلمات
المبارکة الرّوحیّة ما تقدر ان تعرف رشحاً من طمطام ابحر العلم و قمقام انهر
العزّ و تکون من اصبع الهویّة علی قلم الأحدیّة فی امّ الکتاب بالجهل مکتوباً و
لن تحلّ لک حرفاً من الکتاب و لا کلمات آل اللّه فی اسرار المبدإ و المآب

\end{otherlanguage}

Wa-min dūn dhālika wa-kull mā alqaynāka min jawāhir asrār al-ḥikma fī
ghayāhib hādhihi al-kalimāt al-mubāraka al-rūḥiyya mā taqdir an taʿrif
rashḥan min ṭamṭām abḥur al-ʿilm wa-qamqām anhur al-ʿizz wa-takūn min
iṣbaʿ al-huwiyya ʿalā qalam al-aḥadiyya fī umm al-kitāb bi-l-jahl
maktūban wa-lan tuḥall laka ḥarfan min al-kitāb wa-lā kalimāt āl Allāh
fī asrār al-mabdaʾ wa-l-maʾāb

Without this, and without all the gems of the mysteries of wisdom We
have imparted to you in the depths of these blessed, spiritual words,
you cannot comprehend even a drop from the vast oceans of knowledge and
the deep rivers of glory. By the finger of the divine Essence, through
the Pen of Oneness, you would be inscribed as ignorant in the Mother
Book, and neither a single letter of the Book nor the words of the
family of God concerning the mysteries of the beginning and the return
would be unraveled for you.

\subsection{22}\label{section-2733}

\begin{otherlanguage}{arabic}

اذاً فأنصف یا ایّها العبد الّذی ما رأیناک فی الظّاهر ولکن وجدنا حبّک فی
الباطن ثمّ اجعل محضرک بین یدی الّذی انّک ان لن تراه انّه هو یراک و انّک ان لن
تعرفه انّه هو یعرفک هل یقدر احد ان یفسّر تلک الکلمات بدلائل متقنة و براهین
واضحة و اشارات لائحة علی قدر الّذی یستریح قلب السّائل و یسکن فؤاد المخاطب
لا فوالّذی نفسی بیده لن یقدر احد ان یشرب رشحاً منها الّا من یدخل فی ظلّ هذه
المدینة الّتی بنیت ارکانها علی جبال الیاقوت المحمّرة و جدارها من زبرجد
الأحدیّة و ابوابها من الماس الصّمدیّة و ترابها من طیب المکرمة

\end{otherlanguage}

Idhan fa-anṣif yā ayyuhā al-ʿabd alladhī mā raʾaynāka fī al-ẓāhir
wa-lākin wajadnā ḥubbaka fī al-bāṭin thumma ijʿal maḥḍaraka bayna yaday
alladhī innaka in lan tarāhu innahu huwa yarāka wa-innaka in lan
taʿrifhu innahu huwa yaʿrifuka hal yaqdir aḥad an yufassir tilka
al-kalimāt bi-dalāʾil mutqana wa-barāhīn wāḍiḥa wa-ishārāt lāʾiḥa ʿalā
qadr alladhī yastarīḥ qalb al-sāʾil wa-yaskun fuʾād al-mukhāṭab lā
fa-wa-alladhī nafsī bi-yadihi lan yaqdir aḥad an yashrab rashḥan minhā
illā man yadkhul fī ẓill hādhihi al-madīna allatī buniyat arkānuhā ʿalā
jibāl al-yāqūt al-muḥammara wa-jidāruhā min zabarjad al-aḥadiyya
wa-abwābuhā min almās al-ṣamadiyya wa-turābuhā min ṭīb al-makruma

Be fair, then, you servant whom We have not seen outwardly, but whose
love We have found inwardly. Place yourself in the presence of Him who
sees you even if you do not see Him, and who knows you even if you do
not know Him. Can anyone explain those words with sound evidence, clear
proofs, and manifest allusions to such a degree that the questioner's
heart finds rest and the one addressed is reassured? No, by Him in whose
hand is My soul! No one can drink even a drop from them except one who
enters the shade of this City, whose foundations are built upon
mountains of crimson ruby, whose wall is of the chrysolite of oneness,
whose gates are of the diamond of eternal self-sufficiency, and whose
dust is of the fragrance of generosity.

\subsection{23}\label{section-2734}

\begin{otherlanguage}{arabic}

و لمّا ذکرنا و القینا علیک من بعض الأسرار مع الحجب و الأستار نرجع الی ما
کنّا فیه فیما عرفنا من کتب القبل لئلّا یزلّ قدمک فی شیء و تکون موقناً فی کلّ
ما رشّحنا علیک من تموّجات ابحر الحیاة فی لاهوت الأسمآء و الصّفات

\end{otherlanguage}

Wa-lammā dhakarnā wa-alqaynā ʿalayka min baʿḍ al-asrār maʿa al-ḥujub
wa-l-astār narjiʿ ilā mā kunnā fīhi fīmā ʿarafnā min kutub al-qabl
li-allā yazill qadamuka fī shayʾ wa-takūn mūqinan fī kull mā rashshaḥnā
ʿalayka min tamawwujāt abḥur al-ḥayāt fī lāhūt al-asmāʾ wa-l-ṣifāt

Now that We have mentioned and imparted to you some of the mysteries,
though veiled and concealed, We return to Our earlier discussion of what
We have learned from the books of the past, so that your foot may not
slip in anything and you may be certain of all that We have sprinkled
upon you from the surging oceans of life in the divine realm of names
and attributes.

\subsection{24}\label{section-2735}

\begin{otherlanguage}{arabic}

و هو مکتوب فی جمیع اسفار الانجیل و هو هذا حین الّذی تکلّمت الرّوح بالنّور و
قال لتلامیذه فاعلموا بأنّ السّموات و الأرض یمکن ان تزولان ولکن کلامی لن
یزول ابداً و کان معلوم عند جنابکم بأنّ المعنی فی هذا الکلام علی ظاهر
العبارة لن یدلّ الّا بأنّ هذه الأسفار من الانجیل تکون باقیة بین العباد الی
ابد الدّهر و لا تنفد احکامها و لا یبید برهانها و کلّ ما شرع فیها و حدّد لها
و قدّر بها یبقی و لا یفنی ابداً

\end{otherlanguage}

Wa-huwa maktūb fī jamīʿ asfār al-Injīl wa-huwa hādhā ḥīn alladhī
takallamat al-Rūḥ bi-l-nūr wa-qāla li-talāmīdhihi fa-iʿlamū bi-anna
al-samāwāt wa-l-arḍ yumkin an tazūlān wa-lākin kalāmī lan yazūl abadan
wa-kāna maʿlūm ʿinda janābikum bi-anna al-maʿnā fī hādhā al-kalām ʿalā
ẓāhir al-ʿibāra lan yadull illā bi-anna hādhihi al-asfār min al-Injīl
takūn bāqiya bayna al-ʿibād ilā abad al-dahr wa-lā tanfad aḥkāmuhā wa-lā
yabīd burhānuhā wa-kull mā sharaʿa fīhā wa-ḥaddada lahā wa-qaddara bihā
yabqā wa-lā yafnā abadan

It is written in all the books of the Gospel as follows, when the Spirit
spoke with light and said to His disciples, ``Know, then, that the
heavens and the earth may pass away, but My words will never pass
away.'' You are aware that the meaning of these words, according to
their outward expression, can only be that these books of the Gospel
will remain among the servants forever, that their ordinances will not
expire nor their proof perish, and that everything He prescribed in
them, defined for them, and ordained through them will endure and never
pass away.

\subsection{25}\label{section-2736}

\begin{otherlanguage}{arabic}

اذاً یا اخی طهّر قلبک و نوّر فؤادک و حدّ بصرک لتعرف الحان طیور الهویّة و
نغمات حمامات القدسیّة فی ملکوت البقائیّة لتعرف تأویل الکلمات و اسرارها
والّا لو تفسّر علی ظاهر العبارة لن تقدر ان تثبت امر من جآء بعد عیسی و لا
تستطیع ان تلزم الخصم و تفوق علی المعاندین من هؤلآء المشرکین لأنّ بهذه
الآیة تستدلّون علمآء الانجیل بأنّ الانجیل ما ینسخ ابداً و لو تظهر تلک
العلامات الّتی کانت مکتوباً فی کتبنا و یظهر هیکل المعهود لا بدّ له بأن یحکم
بین العباد بأحکام الانجیل و لو تظهر کلّ العلامات المکتوبة فی الکتب و یحکم
بغیر ما حکم به عیسی ما نقرّ به و ما نتّبعه لأنّ هذا المطلب من مسلّمات
مطالبهم بمثل ما انتم تشهدون الیوم من علمآء القوم و جهلائهم فیما یعترضون

\end{otherlanguage}

Idhan yā akhī ṭahhir qalbaka wa-nawwir fuʾādaka wa-ḥudd baṣaraka
li-taʿrif alḥān ṭuyūr al-huwiyya wa-naghamāt ḥamāmāt al-qudsiyya fī
malakūt al-baqāʾiyya li-taʿrif taʾwīl al-kalimāt wa-asrārahā wa-illā law
tufassir ʿalā ẓāhir al-ʿibāra lan taqdir an tuthbit amr man jāʾa baʿda
ʿĪsā wa-lā tastaṭīʿ an tulzim al-khaṣm wa-tafūq ʿalā al-muʿānidīn min
hāʾulāʾ al-mushrikīn li-anna bi-hādhihi al-āya tastadillūn ʿulamāʾ
al-Injīl bi-anna al-Injīl mā yunsakh abadan wa-law taẓhar tilka
al-ʿalāmāt allatī kānat maktūban fī kutubinā wa-yaẓhar haykal al-maʿhūd
lā budda lahu bi-an yaḥkum bayna al-ʿibād bi-aḥkām al-Injīl wa-law
taẓhar kull al-ʿalāmāt al-maktūba fī al-kutub wa-yaḥkum bi-ghayr mā
ḥakama bihi ʿĪsā mā nuqirr bihi wa-mā nattabiʿuhu li-anna hādhā
al-maṭlab min musallamāt maṭālibihim bi-mithl mā antum tashhadūn al-yawm
min ʿulamāʾ al-qawm wa-juhalāʾihim fīmā yaʿtariḍūn

Therefore, My brother, purify your heart, illumine your inmost being,
and sharpen your sight, so that you may recognize the melodies of the
birds of the divine Essence and the songs of the doves of holiness in
the kingdom of everlasting life, and understand the interpretation of
the words and their mysteries. Otherwise, if you interpret them
according to their outward expression, you will never be able to
establish the Cause of Him who came after Jesus, nor will you be able to
confute your opponent and prevail over the obstinate among these people
who associate partners with God. For the scholars of the Gospel invoke
this verse as evidence that the Gospel can never be abrogated: ``Even if
those signs written in our books should appear and the promised Person
should be manifested, He must judge among the servants according to the
ordinances of the Gospel. Even if all the signs written in the books
should appear, if He rules otherwise than Jesus ruled, we will neither
acknowledge nor follow Him.'' This is one of the tenets they accept as
certain, just like what you witness today in the objections raised by
the learned and the ignorant among the people.

\subsection{26}\label{section-2737}

\begin{otherlanguage}{arabic}

و یقولون بأنّ الشّمس ما اشرقت من المغرب و ما صاح الصّائح بین السّمآء و الأرض
و ما غرق بعض البلاد و ما ظهر الدّجّال و ما قام السّفیانی و ما ظهر الهیکل فی
الشّمس و انّی بسمعی سمعت عن واحد من علمائهم یقول لو یظهر کلّ تلک العلامات و
یظهر قائم المأمول و یحکم بغیر ما نزّل فی القرآن فیما یکون بین ایدینا من
الفروع لنکذّبه و نقتله و ما نقرّ به ابداً و امثال ذلک عمّا یقولون هؤلآء
المکذّبون بعد الّذی قام القیامة و نفخ فی الصّور و حشر کلّ من فی السّموات و
الأرض و المیزان نصبت و الصّراط وضعت و الآیات نزلت و الشّمس اشرقت و النّجوم
طمست و النّفوس بعثت و الرّوح نفخت و الملائکة صفّت و الجنّة ازلفت و النّار
سعّرت و قضی کلّ ذلک و الی حینئذ ما عرف احد منهم کأنّهم فی غشواتهم میّتون الّا
الّذینهم آمنوا و رجعوا الی اللّه و کانوا الیوم فی رضوان القدس یحبرون و فی
رضی اللّه یسلکون

\end{otherlanguage}

Wa-yaqūlūn bi-anna al-shams mā ashraqat min al-maghrib wa-mā ṣāḥa
al-ṣāʾiḥ bayna al-samāʾ wa-l-arḍ wa-mā ghariqa baʿḍ al-bilād wa-mā
ẓahara al-Dajjāl wa-mā qāma al-Sufyānī wa-mā ẓahara al-haykal fī
al-shams wa-innī bi-samʿī samiʿtu ʿan wāḥid min ʿulamāʾihim yaqūl law
yaẓhar kull tilka al-ʿalāmāt wa-yaẓhar qāʾim al-maʾmūl wa-yaḥkum
bi-ghayr mā nuzzila fī al-Qurʾān fīmā yakūn bayna aydīnā min al-furūʿ
la-nukadhdhibannahu wa-naqtuluhu wa-mā nuqirr bihi abadan wa-amthāl
dhālika ʿammā yaqūlūn hāʾulāʾ al-mukadhdhibūn baʿda alladhī qāma
al-qiyāma wa-nufikha fī al-ṣūr wa-ḥushira kull man fī al-samāwāt
wa-l-arḍ wa-l-mīzān nuṣibat wa-l-ṣirāṭ wuḍiʿat wa-l-āyāt nazalat
wa-l-shams ashraqat wa-l-nujūm ṭumisat wa-l-nufūs buʿithat wa-l-rūḥ
nufikhat wa-l-malāʾika ṣuffat wa-l-janna uzlifat wa-l-nār suʿʿirat
wa-quḍiya kull dhālika wa-ilā ḥīniʾidh mā ʿarafa aḥad minhum ka-annahum
fī ghashawātihim mayyitūn illā alladhīnahum āmanū wa-rajaʿū ilā Allāh
wa-kānū al-yawm fī riḍwān al-quds yuḥbarūn wa-fī riḍā Allāh yaslukūn

They say, ``The sun has not risen in the west, the crier has not cried
out between heaven and earth, certain lands have not been submerged, the
Dajjal has not appeared, the Sufyani has not arisen, and the Figure has
not appeared in the sun.'' With My own ears I heard one of their learned
men say, ``Even if all those signs should appear and the awaited Qaʾim
should be manifested, if He rules otherwise than what was revealed in
the Qurʾan concerning the subsidiary laws we possess, we will surely
denounce Him as a liar and kill Him, and will never acknowledge Him.''
Such things these deniers say, even after the Resurrection has come to
pass, the trumpet has been sounded, all who are in the heavens and the
earth have been gathered, the Balance has been set up, the Path has been
laid down, the verses have descended, the sun has shone forth, the stars
have been blotted out, souls have been raised, the spirit has been
breathed forth, the angels have been arrayed, Paradise has been brought
near, and the fire has been kindled. All this has been accomplished, yet
to this very moment none of them has understood, as though they were
dead beneath their veils, except those who have believed and returned to
God, who today rejoice in the Paradise of holiness and walk in God's
good pleasure.

\subsection{27}\label{section-2738}

\begin{otherlanguage}{arabic}

و کلّ النّاس لمّا احتجبوا بغشوات انفسهم ما عرفوا الحان القدس و ما شمّوا
روائح الفضل و ما سألوا عن اهل الذّکر بعد الّذی امرهم اللّه بذلک قال و قوله
الحقّ فاسألوا اهل الذّکر ان کنتم لا تعلمون بل اعرضوا عن اهل الذّکر و اتّبعوا
السّامری بأهوائهم و بذلک بعدوا عن رحمة اللّه و ما فازوا بجماله یوم لقائه
بعد الّذی کلّ انتظروا یوم ظهوره و دعوا اللّه فی اللّیالی و الانهار بأن
یحشرهم بین یدیه لیستشهدوا فی سبیله و یستهدوا بهدایته و یستنوروا بنوره
فلمّا جآءهم بآیة من عند اللّه و حجّة من لدنه کفّروه و سبّوه و فعلوا به ما
فعلوا علی مقام لا انا اقدر ان اذکر و لا انت تقدر ان تسمع و القلم حینئذ
یضجّ و المداد یبکی و یصرخ و انّک لو تتوجّه بسمع الفطرة فواللّه لتسمع ضجیج
اهل السّموات و لو تکشف الحجاب عن عینیک لتشهد بأنّ الحوریّات مغشیّات و
الأرواح منصعقات و یضربن علی وجوههنّ و جلسن علی وجه التّراب

\end{otherlanguage}

Wa-kullu al-nāsi lammā iḥtajabū bi-ghishawāti anfusihim mā ʿarafū alḥāna
al-qudsi wa-mā shammū rawāʾiḥa al-faḍli wa-mā saʾalū ʿan ahli al-dhikri
baʿda alladhī amarahumu Allāhu bi-dhālika. Qāla wa-qawluhu al-ḥaqq:
fa-sʾalū ahla al-dhikri in kuntum lā taʿlamūn. Bal aʿraḍū ʿan ahli
al-dhikri wa-ittabaʿū al-Sāmiriyya bi-ahwāʾihim wa-bi-dhālika baʿudū ʿan
raḥmati Allāhi wa-mā fāzū bi-jamālihi yawma liqāʾihi baʿda alladhī
kullun intaẓarū yawma ẓuhūrihi wa-daʿaw Allāha fī al-layālī wa-al-anhāri
bi-an yaḥshurahum bayna yadayhi li-yustashhadū fī sabīlihi wa-yastahdū
bi-hidāyatihi wa-yastanīrū bi-nūrihi. Fa-lammā jāʾahum bi-āyatin min
ʿindi Allāhi wa-ḥujjatin min ladunhu kaffarūhu wa-sabbūhu wa-faʿalū bihi
mā faʿalū ʿalā maqāmin lā anā aqdiru an adhkura wa-lā anta taqdiru an
tasmaʿa. Wa-al-qalamu ḥīnaʾidhin yaḍijju wa-al-midādu yabkī wa-yaṣrukhu
wa-innaka law tatawajjahu bi-samʿi al-fiṭrati fa-wa-Allāhi la-tasmaʿu
ḍajīja ahli al-samāwāti wa-law takshifu al-ḥijāba ʿan ʿaynayka
la-tashhadu bi-anna al-ḥūriyyāti maghshiyyātun wa-al-arwāḥa munṣaʿiqātun
wa-yaḍribna ʿalā wujūhihinna wa-jalasna ʿalā wajhi al-turāb.

Because all people were veiled by the coverings of their own selves,
they did not recognize the melodies of holiness, did not inhale the
fragrances of grace, and did not inquire of the people of remembrance,
even though God had commanded them to do so. He said---and His word is
true---``Ask the people of remembrance if you do not know.'' Instead,
they turned away from the people of remembrance and followed the Samiri
according to their own desires. Thus they became far removed from God's
mercy and failed to attain His beauty on the day of meeting Him, even
though all had awaited the day of His appearance and had prayed to God
by night and day to gather them into His presence, so that they might
suffer martyrdom in His path, be guided by His guidance, and be
illumined by His light. Yet when He came to them with a sign from God
and a proof from His presence, they declared Him an unbeliever, reviled
Him, and did to Him what they did, to such a degree that neither can I
recount it nor can you bear to hear it. At this the Pen cries out, and
the ink weeps and screams. If you were to listen with the hearing of
your innate nature, by God, you would hear the lamentation of the
inhabitants of the heavens. If you were to lift the veil from your eyes,
you would see that the maidens of heaven have fainted and the spirits
have fallen senseless, striking their faces and sitting upon the dust.

\subsection{28}\label{section-2739}

\begin{otherlanguage}{arabic}

فآه آه عمّا ورد علی مظهر نفس اللّه و ما فعلوا به و بأحبّائه بحیث ما فعل احد
علی احد و لا نفس الی نفس و لا کافر الی مؤمن و لا مؤمن الی کافر فآه آه قد
جلس هیکل البقآء فی التّراب السّودآء و ناحت روح القدس فی رفارف الأعلی و
تهدّمت ارکان العرش فی لاهوت الأسنی و تبدّلت عیش الوجود فی ارض الحمرآء و
خرست لسان الورقآء فی جبروت الصّفرآء افّ لهم و بما اکتسبت ایدیهم و عن کلّ ما
هم کانوا ان یعملون

\end{otherlanguage}

Fa-āh āh ʿammā warada ʿalā maẓhari nafsi Allāhi wa-mā faʿalū bihi
wa-bi-aḥibbāʾihi bi-ḥaythu mā faʿala aḥadun ʿalā aḥadin wa-lā nafsun ilā
nafsin wa-lā kāfirun ilā muʾminin wa-lā muʾminun ilā kāfirin. Fa-āh āh
qad jalasa haykalu al-baqāʾi fī al-turābi al-sawdāʾi wa-nāḥat rūḥu
al-qudusi fī rafārifi al-aʿlā wa-tahaddamat arkānu al-ʿarshi fī lāhūti
al-asnā wa-tabaddalat ʿayshu al-wujūdi fī arḍi al-ḥamrāʾi wa-kharisat
lisānu al-warqāʾi fī jabarūti al-ṣafrāʾi. Uffin lahum wa-bi-mā iktasabat
aydīhim wa-ʿan kulli mā hum kānū an yaʿmalūn.

Alas, alas for what befell the Manifestation of God's own Self, and for
what they did to Him and His loved ones---such things as no one has done
to another, no soul to another soul, no unbeliever to a believer, and no
believer to an unbeliever. Alas, alas! The Temple of eternity has sat
upon the black dust, the Holy Spirit has lamented in the retreats of the
Most Exalted, the pillars of the Throne have collapsed in the most
resplendent realm of divinity, the life of existence has been
transformed in the crimson land, and the tongue of the Dove has fallen
silent in the golden dominion. Shame on them, on what their hands have
earned, and on all that they have done!

\subsection{29}\label{section-2740}

\begin{otherlanguage}{arabic}

فاسمع ما غنّت الورقآء فی شأنهم بأحسن نغمات بدیع و اکمل تغرّدات منیع لیکون
حسرة علیهم من یومئذ الی یوم الّذی یقوم النّاس لربّ العالمین و کانوا من قبل
یستفتحون علی الّذین کفروا فلمّا جآءهم ما عرفوا کفروا به فلعنة اللّه علی
الکاذبین هذا شأنهم و مبلغهم فی حیاة الباطلة و سیردّون الی عذاب السّعیر و
لن یجدوا لأنفسهم لا من ولیّ و لا من نصیر

\end{otherlanguage}

Fa-ismaʿ mā ghannat al-warqāʾu fī shaʾnihim bi-aḥsani naghamāti badīʿin
wa-akmali tagharrudāti manīʿin li-yakūna ḥasratan ʿalayhim min
yawmiʾidhin ilā yawmi alladhī yaqūmu al-nāsu li-rabbi al-ʿālamīn.
Wa-kānū min qablu yastaftiḥūna ʿalā alladhīna kafarū fa-lammā jāʾahum mā
ʿarafū kafarū bihi fa-laʿnatu Allāhi ʿalā al-kādhibīn. Hādhā shaʾnuhum
wa-mablaghuhum fī ḥayāti al-bāṭilati wa-sa-yuraddūna ilā ʿadhābi
al-saʿīri wa-lan yajidū li-anfusihim lā min waliyyin wa-lā min naṣīr.

Listen, then, to what the Dove has sung concerning them in the most
beautiful, wondrous melodies and the most perfect, exalted warblings,
that it may be a cause of remorse to them from that day until the day
when people stand before the Lord of the worlds: ``Previously they had
prayed for victory over those who disbelieved, but when there came to
them what they recognized, they disbelieved in it. So God's curse is
upon the liars.'' Such is their condition and the extent of their
attainment in this vain life. They will be returned to the torment of
the blazing fire, and they will find for themselves neither protector
nor helper.

\subsection{30}\label{section-2741}

\begin{otherlanguage}{arabic}

و لا یحجبک کلّ ما نزل فی الفرقان و ما سمعت عن آثار شموس العصمة و بدور
العظمة فی تحریف الغالین و تبدیل المتحرّفین ما کان مقصودهم من تلک الکلمات
الّا فی بعض الموارد المخصوصة المنصوصة و انّی مع عجزی و فقری لو ارید ان
اذکر لجنابک ما هو المذکور لأقدر ولکن یعزب عنّا المقصود و نبعد عن هذا
الصّراط الممدود و نغرق فی اشارات المحدود و نخرج عمّا هو المحبوب فی ساحة
المحمود

\end{otherlanguage}

Wa-lā yaḥjubka kullu mā nazala fī al-Furqāni wa-mā samiʿta ʿan āthāri
shumūsi al-ʿiṣmati wa-budūri al-ʿaẓamati fī taḥrīfi al-ghālīna
wa-tabdīli al-mutaḥarrifīn. Mā kāna maqṣūduhum min tilka al-kalimāti
illā fī baʿḍi al-mawāridi al-makhṣūṣati al-manṣūṣati wa-innī maʿa ʿajzī
wa-faqrī law urīdu an adhkura li-janābika mā huwa al-madhkūru la-aqdiru
wa-lākin yaʿzubu ʿannā al-maqṣūdu wa-nabʿudu ʿan hādhā al-ṣirāṭi
al-mamdūdi wa-naghraqu fī ishārāti al-maḥdūdi wa-nakhruju ʿammā huwa
al-maḥbūbu fī sāḥati al-maḥmūd.

Do not let all that was revealed in the Qurʾan, or what you have heard
from the writings of the Suns of infallibility and the Moons of grandeur
concerning the distortions of the extremists and the alterations of
those who pervert the text, become a veil for you. By those words they
intended only certain specific, expressly designated instances. Despite
My helplessness and poverty, if I wished to set forth for you what is
mentioned, I could do so. But then the intended aim would elude Us, We
would stray from this outstretched path, become submerged in allusions
to what is limited, and depart from what is beloved in the court of the
All-Praised.

\subsection{31}\label{section-2742}

\begin{otherlanguage}{arabic}

و انّک انت یا ایّها المذکور فی هذا الرّقّ المنشور و المستنور فی هذه الظّلمات
الدّیجور فیما تجلّی علیک من انوار الطّور فی سینآء الظّهور نزّه نفسک عن کلّ ما
عرفت من قبل من اشارات السّوئیّة و الدّلالات الشّرکیّة لتجد رائحة البقآء عن
یوسف الوفآء و تکون داخلاً فی مصر العمآء و تجد روائح طیب السّنآء عن هذا
اللّوح الدّرّیّ البیضآء فیما رقم فیه القلم من اسرار القدم فی اسمآء ربّه العلیّ
الأعلی لتکون من الموقنین فی الواح القدس مکتوباً

\end{otherlanguage}

Wa-innaka anta yā ayyuhā al-madhkūru fī hādhā al-raqqi al-manshūri
wa-al-mustanīru fī hādhihi al-ẓulumāti al-dayjūri fīmā tajallā ʿalayka
min anwāri al-Ṭūri fī Sīnāʾi al-ẓuhūri nazzih nafsaka ʿan kulli mā
ʿarafta min qablu min ishārāti al-sūʾiyyati wa-al-dalālāti
al-shirkiyyati li-tajida rāʾiḥata al-baqāʾi ʿan Yūsufi al-wafāʾi
wa-takūna dākhilan fī Miṣri al-ʿamāʾi wa-tajida rawāʾiḥa ṭībi al-sanāʾi
ʿan hādhā al-lawḥi al-durriyyi al-bayḍāʾi fīmā raqama fīhi al-qalamu min
asrāri al-qidami fī asmāʾi rabbihi al-ʿaliyyi al-aʿlā li-takūna mina
al-mūqinīna fī alwāḥi al-qudsi maktūban.

And you, who are mentioned in this outspread parchment and illumined
amid this pitch-black darkness by the lights of the Mount that have
shone upon you on the Sinai of revelation, purify yourself of all the
evil allusions and idolatrous indications you have previously known, so
that you may inhale the fragrance of eternity from the Joseph of
faithfulness, enter the Egypt of the cloud of unknowing, and perceive
the sweet fragrances of splendor from this radiant, snow-white Tablet,
through the mysteries of eternity that the Pen has inscribed within it
concerning the names of its Lord, the Exalted, the Most High, that you
may be recorded among those who possess certitude in the Tablets of
holiness.

\subsection{32}\label{section-2743}

\begin{otherlanguage}{arabic}

ثمّ اعلم یا ایّها الحاضر بین یدی العبد حین غفلتک عن ذلک لا بدّ لمن یرید ان
یقطع الأسفار فی معارج الأسرار بأن یجاهد فی الدّین علی قدر طاقته و قدرته
لیظهر له السّبیل فی مناهج الدّلیل و ان یجد نفساً یدّعی امراً من اللّه و کان فی
یده حجّة من مولاه الّتی تعجز عنها العالمین لا مفرّ له الّا بأن یتّبعه فی کلّ
ما یأمر و یقول و یحکم ولو یجری علی المآء حکم الأرض او علی الأرض حکم
السّمآء او فوق ذلک او تحت ذلک ولو یحکم بالتّغییر او بالتّبدیل لأنّه اطّلع
بأسرار الهویّة و رموزات الغیبیّة و احکام الالهیّة

\end{otherlanguage}

Thumma iʿlam yā ayyuhā al-ḥāḍiru bayna yaday al-ʿabdi ḥīna ghaflatika
ʿan dhālika lā budda li-man yurīdu an yaqṭaʿa al-asfāra fī maʿāriji
al-asrāri bi-an yujāhida fī al-dīni ʿalā qadri ṭāqatihi wa-qudratihi
li-yaẓhara lahu al-sabīlu fī manāhiji al-dalīli wa-an yajida nafsan
yaddaʿī amran mina Allāhi wa-kāna fī yadihi ḥujjatun min mawlāhu allatī
taʿjazu ʿanhā al-ʿālamīna lā mafarra lahu illā bi-an yattabiʿahu fī
kulli mā yaʾmuru wa-yaqūlu wa-yaḥkumu wa-law yujrī ʿalā al-māʾi ḥukma
al-arḍi aw ʿalā al-arḍi ḥukma al-samāʾi aw fawqa dhālika aw taḥta
dhālika wa-law yaḥkumu bi-al-taghyīri aw bi-al-tabdīli li-annahu
iṭṭalaʿa bi-asrāri al-huwiyyati wa-rumūzāti al-ghaybiyyati wa-aḥkāmi
al-ilāhiyyati.

Then know, you who are present before this Servant while unaware of it,
that whoever wishes to undertake the journeys through the ascending
paths of the mysteries must strive in religion to the full extent of his
strength and ability, so that the way may become manifest to him along
the paths of proof. If he finds a person who lays claim to a Cause from
God and holds in his hand a proof from his Lord that the peoples of the
world are powerless to match, he has no recourse but to follow him in
everything he commands, says, and decrees---even if he applies to water
the law of earth, or to earth the law of heaven, or what is above or
below that, and even if he decrees change or alteration. For he is
acquainted with the mysteries of the divine Essence, the symbols of the
unseen, and the divine ordinances.

\subsection{33}\label{section-2744}

\begin{otherlanguage}{arabic}

و لو انّ کلّ العباد من امم المختلفة یعملون بما ذکرنا حینئذ لیسهل علیهم
امرهم و ما یمنعهم تلک العبارات و الاشارات عن الورود فی غمرات الأسمآء و
الصّفات و لو عرفوا ذلک ما کفروا بأنعم اللّه و ما حاربوا مع النّبیّین و ما
جاحدوهم و ما انکروهم و بمثل تلک العبارات تجدون فی القرآن لو انتم فیه
تتفکّرون

\end{otherlanguage}

Wa-law anna kulla al-ʿibādi min umami al-mukhtalifati yaʿmalūna bi-mā
dhakarnā ḥīnaʾidhin la-yashulu ʿalayhim amruhum wa-mā yamnaʿuhum tilka
al-ʿibārātu wa-al-ishārātu ʿani al-wurūdi fī ghamarāti al-asmāʾi
wa-al-ṣifāti wa-law ʿarafū dhālika mā kafarū bi-anʿumi Allāhi wa-mā
ḥārabū maʿa al-nabiyyīna wa-mā jāḥadūhum wa-mā ankarūhum wa-bi-mithli
tilka al-ʿibārāti tajidūna fī al-Qurʾāni law antum fīhi tatafakkarūn.

If all the servants among the various nations were to act according to
what We have mentioned, their task would then become easy for them, and
those expressions and allusions would not prevent them from entering the
depths of the names and attributes. Had they understood this, they would
not have denied God's blessings, fought against the Prophets, repudiated
them, or rejected them. You will find expressions like these in the
Qurʾan, if you reflect upon it.

\subsection{34}\label{section-2745}

\begin{otherlanguage}{arabic}

ثمّ اعلم بأنّ بمثل تلک الکلمات یمحّص اللّه عباده و یغربلنّهم و یفصّل بین
المؤمن و الکافر و المنقطع و المتمسّک و المحسن و المجرم و التّقیّ و الشّقیّ و
امثال ذلک کما نطق بذلک ورقآء الهویّة الم أ حسب النّاس ان یترکوا ان یقولوا
آمنّا و هم لا یفتنون

\end{otherlanguage}

Thumma iʿlam bi-anna bi-mithli tilka al-kalimāti yumaḥḥiṣu Allāhu
ʿibādahu wa-yugharbilannahum wa-yufaṣṣilu bayna al-muʾmini wa-al-kāfiri
wa-al-munqaṭiʿi wa-al-mutamassiki wa-al-muḥsini wa-al-mujrimi
wa-al-taqiyyi wa-al-shaqiyyi wa-amthāli dhālika kamā naṭaqa bi-dhālika
warqāʾu al-huwiyyati: Alif Lām Mīm. A-ḥasiba al-nāsu an yutrakū an
yaqūlū āmannā wa-hum lā yuftanūn.

Then know that through words such as these God tests His servants, sifts
them, and distinguishes the believer from the unbeliever, the detached
from the attached, the doer of good from the wrongdoer, the devout from
the wretched, and others like them. Thus has the Dove of the divine
Essence proclaimed: ``Alif Lam Mim. Do people suppose that they will be
left to say, `We believe,' without being put to the test?''

\subsection{35}\label{section-2746}

\begin{otherlanguage}{arabic}

لا بدّ للمسافر الی اللّه و المهاجر فی سبیله بأن ینقطع عن کلّ من فی السّموات
و الأرض و یکفّ نفسه عن کلّ ما سواه لیفتح علی وجهه ابواب العنایة و تهبّ علیه
نسمات العطوفة و اذا کتب علی نفسه ما القیناه من جواهر المعانی و البیان
لیعرف کلّ الاشارات من تلک الدّلالات و ینزل اللّه علی قلبه سکینة من عنده و
یجعله من السّاکنین و بمثل هذه الکلمات المتشابهات المنزلة فاعرف ما سألت عن
هذا العبد الّذی جلس علی نقطة الذّلّة و ما یمشی فی الأرض الّا کمثل غریب الّذی
لن یجد لنفسه لا من معین و لا من مؤنس و لا من حبیب و لا من نصیراً و یکون
متوکّلاً علی اللّه و یقول فی کلّ حین انّا للّه و انّا الیه راجعون

\end{otherlanguage}

Lā budda li-al-musāfiri ilā Allāhi wa-al-muhājiri fī sabīlihi bi-an
yanqaṭiʿa ʿan kulli man fī al-samāwāti wa-al-arḍi wa-yakuffa nafsahu ʿan
kulli mā siwāhu li-yuftaḥa ʿalā wajhihi abwābu al-ʿināyati wa-tahubba
ʿalayhi nasamātu al-ʿaṭūfati wa-idhā kataba ʿalā nafsihi mā alqaynāhu
min jawāhiri al-maʿānī wa-al-bayāni li-yaʿrifa kulla al-ishārāti min
tilka al-dalālāti wa-yunazzilu Allāhu ʿalā qalbihi sakīnatan min ʿindihi
wa-yajʿaluhu mina al-sākinīn. Wa-bi-mithli hādhihi al-kalimāti
al-mutashābihāti al-munazzalati fa-iʿraf mā saʾalta ʿan hādhā al-ʿabdi
alladhī jalasa ʿalā nuqṭati al-dhillati wa-mā yamshī fī al-arḍi illā
ka-mithli gharībin alladhī lan yajida li-nafsihi lā min muʿīnin wa-lā
min muʾnisin wa-lā min ḥabībin wa-lā min naṣīran wa-yakūnu mutawakkilan
ʿalā Allāhi wa-yaqūlu fī kulli ḥīnin innā li-llāhi wa-innā ilayhi
rājiʿūn.

The traveler toward God, who leaves his home in His path, must detach
himself from everyone in the heavens and on earth and restrain himself
from everything other than Him, so that the doors of providence may be
opened before his face and the breezes of loving-kindness may blow upon
him. When he has made binding upon himself the gems of meaning and
utterance that We have imparted, he will recognize all the allusions in
those indications, and God will send down upon his heart a tranquility
from His presence and make him one of those who dwell in peace. Through
such revealed words of ambiguous meaning, then, understand what you have
asked of this Servant, who sits at the lowest point of abasement and
walks the earth only as a stranger who finds for himself neither helper,
companion, loved one, nor supporter. He places His trust in God and says
at every moment, ``We belong to God, and to Him we return.''

\subsection{36}\label{section-2747}

\begin{otherlanguage}{arabic}

و انّ ما ذکرنا الکلمات بالمتشابهات هذا لم یکن الّا عند الّذین لن یتعارجوا
الی افق الهدایة و ما وصلوا الی مراتب العرفان فی مکامن العنایة والّا عند
الّذین‌هم عرفوا مواقع الأمر و شهدوا اسرار الولایة فیما القی اللّه علی
انفسهم کلّ الآیات محکمات عندهم و کلّ الاشارات متقنات لدیهم و انّهم یعرفون
اسرار المودعة فی قمص الکلمات بمثل انتم تعرفون من الشّمس الحرارة و من
المآء الرّطوبة بل اظهر من ذلک فتعالی اللّه عمّا کنّا فی ذکر احبّائه فتعالی
عمّا هم یذکرون

\end{otherlanguage}

Wa-inna mā dhakarnā al-kalimāti bi-al-mutashābihāti hādhā lam yakun illā
ʿinda alladhīna lan yataʿārajū ilā ufuqi al-hidāyati wa-mā waṣalū ilā
marātibi al-ʿirfāni fī makāmini al-ʿināyati wa-illā ʿinda alladhīna hum
ʿarafū mawāqiʿa al-amri wa-shahidū asrāra al-wilāyati fīmā alqā Allāhu
ʿalā anfusihim kullu al-āyāti muḥkamātun ʿindahum wa-kullu al-ishārāti
mutqanātun ladayhim wa-innahum yaʿrifūna asrāra al-mūdaʿati fī qumuṣi
al-kalimāti bi-mithli antum taʿrifūna mina al-shamsi al-ḥarārata wa-mina
al-māʾi al-ruṭūbata bal aẓhara min dhālika fa-taʿālā Allāhu ʿammā kunnā
fī dhikri aḥibbāʾihi fa-taʿālā ʿammā hum yadhkurūn.

Our description of these words as ambiguous applies only to those who
have not ascended to the horizon of guidance or attained the degrees of
recognition in the retreats of providence. But for those who have
recognized the stations of the Cause and witnessed the mysteries of
divine guardianship in what God has imparted to their souls, all the
verses are clear and all the allusions firmly established. They
recognize the mysteries deposited within the garments of words just as
you recognize heat from the sun and moisture from water---indeed, more
clearly than that. Exalted is God above what We have said in speaking of
His loved ones! Exalted is He above what they recount!

\subsection{37}\label{section-2748}

\begin{otherlanguage}{arabic}

اذاً لمّا وصلنا الی ذلک المقام الأسنی و بلغنا الی ذروة الأحلی فیما یجری من
هذا القلم من عنایة الکبری من لدی اللّه العلیّ الأعلی اردنا بأن نذکر لک
بعضاً من مقامات سلوک العبد فی اسفاره الی مبدئه لیکشف علی جنابک کلّ ما اردت
و ترید لتکون الحجّة بالغة و النّعمة سابغة

\end{otherlanguage}

Idhan lammā waṣalnā ilā dhālika al-maqāmi al-asnā wa-balaghnā ilā
dhirwati al-aḥlā fīmā yaj\-rī min hādhā al-qalami min ʿināyati al-kubrā
min ladā Allāhi al-ʿaliyyi al-aʿlā aradnā bi-an nadhkura laka baʿḍan min
maqāmāti sulūki al-ʿabdi fī asfārihi ilā mabdaʾihi li-yukshafa ʿalā
janābika kullu mā aradta wa-turīdu li-takūna al-ḥujjatu bālighatan
wa-al-niʿmatu sābighatan.

Now that We have reached this most resplendent station and attained this
sweetest summit through what flows from this Pen by the supreme
providence of God, the Exalted, the Most High, We wish to mention to you
some of the stations of the servant's progress in his journeys toward
his Origin, so that all you have desired and now desire may be disclosed
to you, that the proof may be conclusive and the blessing complete.

\subsection{38}\label{section-2749}

\begin{otherlanguage}{arabic}

فاعلم ثمّ اعرف بأنّ السّالک فی اوّل سلوکه الی اللّه لا بدّ له بأن یدخل فی
حدیقة الطّلب (طلب مدینة قدس اوّلیّة فی مسلک عزّ طلبیّة) و فی هذا السّفر ینبغی
للسّالک بأن ینقطع عن کلّ ما سوی اللّه و یغمض عیناه عن کلّ من فی السّموات و
الأرض و لم یکن فی قلبه بغض احد من العباد و لا حبّ احد علی قدر الّذی یمنعه
عن الوصول الی مکمن الجمال و یقدّس نفسه عن سبحات الجلال و له حقّ بأن لا
یفتخر علی احد فی کلّ ما اعطاه اللّه من زخارف الدّنیا او من علوم الظّاهرة او
غیرها و یطلب الحقّ بکمال جدّه و سعیه لیعلّمه اللّه سبل عنایته و مناهج مکرمته
لأنّه خیر معین بعباده و احسن ناصر لأرقّائه قال و قوله الحقّ الّذین جاهدوا
فینا لنهدینّهم سبلنا و فی مقام آخر اتّقوا اللّه یعلّمکم اللّه

\end{otherlanguage}

Fa-iʿlam thumma iʿraf bi-anna al-sālika fī awwali sulūkihi ilā Allāhi lā
budda lahu bi-an yadkhula fī ḥadīqati al-ṭalabi (ṭalab madīnatu qudsin
awwaliyyatun fī maslaki ʿizzi ṭalabiyyatin) wa-fī hādhā al-safari
yanbaghī li-al-sāliki bi-an yanqaṭiʿa ʿan kulli mā siwā Allāhi
wa-yughmiḍa ʿaynāhu ʿan kulli man fī al-samāwāti wa-al-arḍi wa-lam yakun
fī qalbihi bughḍu aḥadin mina al-ʿibādi wa-lā ḥubbu aḥadin ʿalā qadri
alladhī yamnaʿuhu ʿani al-wuṣūli ilā makmani al-jamāli wa-yuqaddisa
nafsahu ʿan subuḥāti al-jalāli wa-lahu ḥaqqun bi-an lā yaftakhira ʿalā
aḥadin fī kulli mā aʿṭāhu Allāhu min zakhārifi al-dunyā aw min ʿulūmi
al-ẓāhirati aw ghayrihā wa-yaṭluba al-ḥaqqa bi-kamāli jiddihi wa-saʿyihi
li-yuʿallimahu Allāhu subula ʿināyatihi wa-manāhija makramatihi
li-annahu khayru muʿīnin bi-ʿibādihi wa-aḥsanu nāṣirin li-ariqqāʾihi.
Qāla wa-qawluhu al-ḥaqqu alladhīna jāhadū fīnā la-nahdiyannahum subulanā
wa-fī maqāmin ākhara ittaqū Allāha yuʿallimukumu Allāh.

Know, then, and recognize that at the outset of his journey toward God
the wayfarer must enter the garden of search (Search: the first city of
holiness on the path of the glory of seeking). On this journey the
wayfarer should detach himself from everything other than God and close
his eyes to everyone in the heavens and on earth. His heart must harbor
neither hatred for any servant nor love for anyone to a degree that
would prevent him from reaching the retreat of Beauty. He must sanctify
himself from the veils of glory. It is incumbent upon him not to exalt
himself over anyone on account of anything God has given him, whether
the adornments of this world, outward learning, or anything else. He
must seek the Truth with the utmost earnestness and endeavor, so that
God may teach him the ways of His providence and the paths of His
bounty, for He is the best helper of His servants and the finest
supporter of His bondsmen. He said---and His word is true---``Those who
strive for Us, We will surely guide to Our ways,'' and elsewhere, ``Be
mindful of God; God will teach you.''

\subsection{39}\label{section-2750}

\begin{otherlanguage}{arabic}

و فی هذا السّفر یشهد السّالک التّبدیلات و التّغییرات و المختلفات و
المتقارنات و یشهد عجائب الرّبوبیّة فی اسرار الخلیقة و یطّلع علی سبل الهدایة
و طرق الالهیّة هذا مقام الطّالبین و معارج القاصدین

\end{otherlanguage}

Wa-fī hādhā al-safari yashhadu al-sāliku al-tabdīlāti wa-al-taghyīrāti
wa-al-mukhtalifāti wa-al-mutaqārināti wa-yashhadu ʿajāʾiba
al-rubūbiyyati fī asrāri al-khalīqati wa-yaṭṭaliʿu ʿalā subuli
al-hidāyati wa-ṭuruqi al-ilāhiyyati hādhā maqāmu al-ṭālibīna wa-maʿāriju
al-qāṣidīn.

On this journey the wayfarer witnesses alterations and changes,
differences and correspondences. He beholds the wonders of lordship
within the mysteries of creation and becomes acquainted with the ways of
guidance and the divine paths. This is the station of the seekers and
the ascending paths of those who pursue their goal.

\subsection{40}\label{section-2751}

\begin{otherlanguage}{persian}

و اذا استرقی عن ذلک المقام یدخل فی مدینة العشق و الجذب (عشق تموّج بحر
ناریّة فی مدینة جذب عشقیّة) حینئذ تهبّ اریاح المحبّة و تهیج نسمات الرّوحیّة و
یأخذ السّالک فی هذا المقام جذبات الشّوق و نفحات الذّوق بحیث لن یعرف الیمین
عن الشّمال و لا البرّ من البحر و لا الصّحاری عن الجبال و فی کلّ حین یحترق
بنار الاشتیاق و یوقد من سطوة الفراق فی الآفاق و یرکض فی فاران العشق و
حوریب الجذب مرّةً یضحک و مرّةً یبکی و مرّةً یسکن و مرّةً یضطرب و لا یبالی من شیء
و لا یمنعه من امر و لا یسدّه من حکم و ینتظر امر مولاه فی مبدئه و منتهاه و
ینفق روحه فی کلّ حین و یفدی نفسه فی کلّ آن و یقابل صدره فی مقابلة رماح
الأعدآء و یرفع رأسه لسیف القضآء بل یقبّل ایدی من یقتله و ینفق کلّ ما له و
علیه لیفدی روحه و نفسه و جسده فی سبیل مولاه ولکن باذن من محبوبه لا بهوآء
من نفسه و تجده بارداً فی النّار و یابساً فی المآء و یسکن علی کلّ ارض و یمشی
فی کلّ طریق و من یمسّه فی تلک الحالة لیجد حرارة المحبّة منه و انّه یمشی فی
رفرف الانقطاع و یرکض فی وادی الامتناع و لم یزل کانت عیناه منتظراً لبدائع
رحمة اللّه و مشاهدة انوار جماله فهنیئاً للواصلین و هذا مقام العاشقین و شأن
المجتذبین

\end{otherlanguage}

Wa-idhā istarqā ʿan dhālik al-maqām yadkhul fī madīnat al-ʿishq
wa-l-jadhb (ʿishq tamawwuj baḥr nāriyya fī madīnat jadhb ʿishqiyya).
Ḥīnaʾidh tahubb aryāḥ al-maḥabba wa-tahīj nasamāt al-rūḥiyya wa-yaʾkhudh
al-sālik fī hādhā al-maqām jadhabāt al-shawq wa-nafaḥāt al-dhawq
bi-ḥayth lan yaʿrif al-yamīn ʿan al-shamāl wa-lā al-barr min al-baḥr
wa-lā al-ṣaḥārī ʿan al-jibāl. Wa-fī kull ḥīn yaḥtariq bi-nār al-ishtiyāq
wa-yūqad min saṭwat al-firāq fī al-āfāq wa-yarkuḍ fī Fārān al-ʿishq
wa-Ḥūrīb al-jadhb. Marratan yaḍḥak wa-marratan yabkī wa-marratan yaskun
wa-marratan yaḍṭarib wa-lā yubālī min shayʾ wa-lā yamnaʿuhu min amr
wa-lā yasudduhu min ḥukm wa-yantaẓir amr mawlāhu fī mabdaʾihi
wa-muntahāhu wa-yunfiq rūḥahu fī kull ḥīn wa-yafdī nafsahu fī kull ān
wa-yuqābil ṣadrahu fī muqābalat rimāḥ al-aʿdāʾ wa-yarfaʿ raʾsahu li-sayf
al-qaḍāʾ bal yuqabbil aydī man yaqtuluhu wa-yunfiq kull mā lahu
wa-ʿalayhi li-yafdī rūḥahu wa-nafsahu wa-jasadahu fī sabīl mawlāhu
wa-lākin bi-idhn min maḥbūbihi lā bi-hawāʾ min nafsihi. Wa-tajiduhu
bāridan fī al-nār wa-yābisan fī al-māʾ wa-yaskun ʿalā kull arḍ wa-yamshī
fī kull ṭarīq wa-man yamassuhu fī tilka al-ḥāla la-yajid ḥarārat
al-maḥabba minhu wa-innahu yamshī fī rafraf al-inqiṭāʿ wa-yarkuḍ fī wādī
al-imtināʿ wa-lam yazal kānat ʿaynāhu muntaẓiran li-badāʾiʿ raḥmat Allāh
wa-mushāhadat anwār jamālihi. Fa-hanīʾan li-l-wāṣilīn wa-hādhā maqām
al-ʿāshiqīn wa-shaʾn al-mujtadhabīn.

When he ascends beyond that station, he enters the city of love and
attraction (Love: the surging of a fiery sea in the city of love's
attraction). Then the winds of love blow and spiritual breezes stir. In
this station, transports of longing and breaths of spiritual delight
seize the wayfarer, so that he no longer knows right from left, land
from sea, or deserts from mountains. At every moment he burns with the
fire of yearning and is set ablaze throughout the horizons by the
overpowering force of separation. He races through the Paran of love and
the Horeb of attraction. At times he laughs, at times he weeps; at times
he is still, at times he is restless. He cares for nothing; no matter
restrains him and no decree bars his way. He awaits his Lord's command
at his beginning and his end, expends his spirit at every moment, and
sacrifices himself at every instant. He offers his breast to the spears
of his enemies and raises his head to the sword of destiny. Indeed, he
kisses the hands of whoever kills him and gives up everything that is
his and everything incumbent upon him, that he may sacrifice his spirit,
his self, and his body in the path of his Lord---but with his Beloved's
permission, not from his own desire. You find him cool in fire and dry
in water. He dwells in every land and walks every road. Whoever touches
him in that state will feel the heat of love coming from him. He walks
upon the lofty plane of detachment and races through the valley of
inaccessibility. His eyes remain ever expectant of the wonders of God's
mercy and the sight of the lights of His beauty. Blessed are those who
attain! This is the station of the lovers and the condition of those
drawn to Him.

\subsection{41}\label{section-2752}

\begin{otherlanguage}{persian}

و اذا قطع هذا السّفر و استرقی عن هذا المقام الأکبر یدخل فی مدینة التّوحید
(توحید اشراق بدع نزهیّة فی مدینة عزّ احدیّة) و حدیقة التّفرید و بساط التّجرید
و فی هذا المقام یلقی السّالک کلّ الاشارات و الدّلالات و الحجبات و العبارات
و یشهد الأشیآء بعین الّتی تجلّی اللّه له به بنفسه و یشاهد فی هذا السّفر بأنّ
المختلفات کلّها ترجع الی کلمة واحدة و الاشارات تنتهی الی نقطة واحدة کما
شهد بذلک قول من رکب علی فلک النّار و مشی فی قطب الأسفار حتّی وصل الی ذروة
الأعلی فی جبروت البقآء بأنّ العلم نقطة کثرّها الجاهلون و هذا مقام الّذی ذکر
فی الحدیث بأنّی انا هو و هو انا الّا انّه هو هو و انا انا

\end{otherlanguage}

Wa-idhā qaṭaʿa hādhā al-safar wa-istarqā ʿan hādhā al-maqām al-akbar
yadkhul fī madīnat al-tawḥīd (tawḥīd ishrāq bidaʿ nuzhiyya fī madīnat
ʿizz aḥadiyya) wa-ḥadīqat al-tafrīd wa-bisāṭ al-tajrīd. Wa-fī hādhā
al-maqām yulqī al-sālik kull al-ishārāt wa-l-dalālāt wa-l-ḥujubāt
wa-l-ʿibārāt wa-yashhad al-ashyāʾ bi-ʿayn allatī tajallā Allāh lahu bihi
bi-nafsihi wa-yushāhid fī hādhā al-safar bi-anna al-mukhtalifāt kullahā
tarjiʿ ilā kalima wāḥida wa-l-ishārāt tantahī ilā nuqṭa wāḥida kamā
shahida bi-dhālik qawl man rakiba ʿalā fulk al-nār wa-mashā fī quṭb
al-asfār ḥattā waṣala ilā dhirwat al-aʿlā fī jabarūt al-baqāʾ bi-anna
al-ʿilm nuqṭa kaththarahā al-jāhilūn. Wa-hādhā maqām alladhī dhukira fī
al-ḥadīth bi-annī anā huwa wa-huwa anā illā annahu huwa huwa wa-anā anā.

When he completes this journey and ascends beyond this most great
station, he enters the city of unity (Unity: the radiance of wonders of
sanctity in the city of the glory of oneness), the garden of singleness,
and the court of detachment. In this station the wayfarer casts aside
all allusions, indications, veils, and expressions, and beholds things
with the very eye through which God has revealed Himself to him. On this
journey he witnesses how all differences return to a single word and all
allusions end in a single point, as attested by the saying of the one
who boarded the ark of fire and walked at the pivot of journeys until he
reached the loftiest summit in the dominion of everlasting life:
``Knowledge is a point that the ignorant have multiplied.'' This is the
station mentioned in the tradition: ``I am He and He is I, except that
He is He and I am I.''

\subsection{42}\label{section-2753}

\begin{otherlanguage}{arabic}

فی ذلک المقام لو یقول هیکل الختم بأنّی انا نقطة البدء لیصدق و لو یقول
بأنّی انا غیرها لحقّ و لو یقول بأنّی صاحب الملک و الملکوت او ملک الملوک او
سلطان الجبروت او محمّد او علیّ او ابنائهم او غیر ذلک لیکون صادقاً من عند
اللّه و حاکماً علی الممکنات و علی کلّ ما سواه اما سمعت ما ورد من قبل بأنّ
اوّلنا محمّد و آخرنا محمّد و اوسطنا محمّد و فی مقام آخر بأنّ کلّهم من نور واحد

\end{otherlanguage}

Fī dhālik al-maqām law yaqūl haykal al-khatm bi-annī anā nuqṭat al-badʾ
la-yaṣduq wa-law yaqūl bi-annī anā ghayruhā la-ḥaqq wa-law yaqūl bi-annī
ṣāḥib al-mulk wa-l-malakūt aw malik al-mulūk aw sulṭān al-jabarūt aw
Muḥammad aw ʿAlī aw abnāʾihim aw ghayr dhālik la-yakūn ṣādiqan min ʿind
Allāh wa-ḥākiman ʿalā al-mumkināt wa-ʿalā kull mā siwāhu. A-mā samiʿta
mā warada min qabl bi-anna awwalanā Muḥammad wa-ākhiranā Muḥammad
wa-awsaṭanā Muḥammad wa-fī maqām ākhar bi-anna kullahum min nūr wāḥid.

In that station, if the embodiment of the Seal were to say, ``I am the
Point of the beginning,'' he would speak the truth; and if he were to
say, ``I am other than it,'' that too would be true. If he were to say,
``I am the Lord of the earthly and heavenly kingdoms,'' or ``the King of
kings,'' or ``the Sovereign of the dominion of might,'' or ``Muḥammad,''
or ``ʿAlī,'' or ``their descendants,'' or anything else, he would be
speaking truth from God and exercising authority over all contingent
beings and everything besides Him. Have you not heard what was handed
down before: ``Our first is Muḥammad, our last is Muḥammad, and our
middle is Muḥammad,'' and elsewhere, ``They are all of one light''?

\subsection{43}\label{section-2754}

\begin{otherlanguage}{arabic}

و فی ذلک المقام یثبت حکم التّوحید و آیات التّجرید و تجد بأنّ کلّهم رفعوا
رؤوسهم عن جیب قدرة اللّه و یدخلون فی اکمام رحمة اللّه من غیر ان تشاهد
الفرق بین الأکمام و الجیب و التّغییر و التّبدیل فی هذا المقام شرک صرف و
کفر محض لأنّ هذا مقام تجلّی الوحدانیّة و تحکّی الفردانیّة و اشراق انوار فجر
الأزلیّة فی مرایا الرّفیعة المنطبعة و انّی فواللّه لو اذکر هذا المقام علی
قدر الّذی قدّر اللّه فیه لتنقطع الأرواح عن اجسادها و تنزّلت الجوهریّات من
اماکنها و تنصعق کلّ من فی لجج الممکنات و تنعدم کلّ ما یتحرّک فی اراضی
الاشارات

\end{otherlanguage}

Wa-fī dhālik al-maqām yathbut ḥukm al-tawḥīd wa-āyāt al-tajrīd wa-tajid
bi-anna kullahum rafaʿū ruʾūsahum ʿan jayb qudrat Allāh wa-yadkhulūn fī
akmām raḥmat Allāh min ghayr an tushāhid al-farq bayn al-akmām
wa-l-jayb. Wa-l-taghyīr wa-l-tabdīl fī hādhā al-maqām shirk ṣirf wa-kufr
maḥḍ li-anna hādhā maqām tajallī al-waḥdāniyya wa-taḥakkī al-fardāniyya
wa-ishrāq anwār fajr al-azaliyya fī marāyā al-rafīʿa al-munṭabiʿa.
Wa-innī fa-wa-Allāh law adhkur hādhā al-maqām ʿalā qadr alladhī qaddara
Allāh fīhi la-tanqaṭiʿ al-arwāḥ ʿan ajsādihā wa-tanazzalat
al-jawhariyyāt min amākinihā wa-tanṣaʿiq kull man fī lujaj al-mumkināt
wa-tanʿadim kull mā yataḥarrak fī arāḍī al-ishārāt.

In that station the truth of unity and the signs of detachment are
established. You find that all of them have raised their heads from the
bosom of God's power and enter the sleeves of God's mercy, without your
perceiving any distinction between the sleeves and the bosom. To
conceive of change or alteration in this station is sheer association of
partners with God and pure unbelief, for this is the station of the
revelation of oneness, the expression of singleness, and the shining of
the lights of the dawn of eternity in the lofty mirrors that receive
their imprint. By God, if I were to describe this station to the extent
God has ordained for it, spirits would be severed from their bodies,
essences would descend from their places, everyone in the depths of
contingent existence would be struck senseless, and everything moving in
the lands of allusion would cease to exist.

\subsection{44}\label{section-2755}

\begin{otherlanguage}{arabic}

اما سمعت لا تبدیل لخلق اللّه اما قرأت و لن تجد لسنّته من تبدیل و اما شهدت
ما تری فی خلق الرّحمن من تفاوت بلی وربّی من کان من اهل هذه اللّجّة و رکب فی
هذه السّفینة لم یشهد التّبدیل فی خلق اللّه و لا یری التّفاوت فی ارض اللّه و
لمّا لم یکن التّبدیل و التّغییر فی خلق اللّه فکیف یجری علی مظاهر نفس اللّه
فسبحان اللّه عمّا کنّا فی وصف مظاهر امره و تعالی عمّا هم یذکرون

\end{otherlanguage}

A-mā samiʿta lā tabdīl li-khalq Allāh a-mā qaraʾta wa-lan tajid
li-sunnatihi min tabdīl wa-a-mā shahidta mā tarā fī khalq al-Raḥmān min
tafāwut. Balā wa-rabbī man kāna min ahl hādhihi al-lujja wa-rakiba fī
hādhihi al-safīna lam yashhad al-tabdīl fī khalq Allāh wa-lā yarā
al-tafāwut fī arḍ Allāh wa-lammā lam yakun al-tabdīl wa-l-taghyīr fī
khalq Allāh fa-kayfa yajrī ʿalā maẓāhir nafs Allāh. Fa-subḥān Allāh
ʿammā kunnā fī waṣf maẓāhir amrihi wa-taʿālā ʿammā hum yadhkurūn.

Have you not heard, ``There is no alteration in God's creation''? Have
you not read, ``You will find no alteration in His way''? Have you not
witnessed, ``You see no disparity in the creation of the All-Merciful''?
Indeed, by my Lord, whoever belongs to this ocean's depths and has
boarded this ship witnesses no alteration in God's creation and sees no
disparity in God's earth. Since there is no alteration or change in
God's creation, how could it affect the Manifestations of God's own
Self? Glorified, then, is God beyond what We have said in describing the
Manifestations of His Cause, and exalted is He above what they recount.

\subsection{45}\label{section-2756}

\begin{otherlanguage}{arabic}

اللّه اکبر هذا البحر قد ذخرا و هیّج الرّیح موجاً یقذف الدّررا فاخلع ثیابک و
اغرق فیه و دع عنک السّباحة لیس السّبیح مفتخرا

\end{otherlanguage}

Allāh akbar hādhā al-baḥr qad dhakharā wa-hayyaja al-rīḥ mawjan yaqdhif
al-durarā fa-khlaʿ thiyābaka wa-ghraq fīhi wa-daʿ ʿanka al-sibāḥa laysa
al-sabīḥ muftakhirā.

God is most great! This sea has laid up its stores, and the wind has
stirred a wave that casts forth pearls. Strip off your garments and
drown in it; abandon swimming, for the swimmer has nothing to boast of.

\subsection{46}\label{section-2757}

\begin{otherlanguage}{arabic}

و انّک انت لو تکون من اهل هذه المدینة فی هذه اللّجّة الأحدیّة لتری کلّ
النّبیّین و المرسلین کهیکل واحد و نفس واحدة و نور واحد و روح واحدة بحیث
یکون اوّلهم آخرهم و آخرهم اوّلهم و کلّهم قاموا علی امر اللّه و شرعوا شرائع
حکمة اللّه و کانوا مظاهر نفس اللّه و معادن قدرة اللّه و مخازن وحی اللّه و
مشارق شمس اللّه و مطالع نور اللّه و بهم ظهرت آیات التّجرید فی حقائق
الممکنات و علامات التّفرید فی جوهریّات الموجودات و عناصر التّمجید فی ذاتیّات
الأحدیّات و مواقع التّحمید فی ساذجیّات الصّمدیّات و بهم یبدأ الخلق و الیهم
یعید کلّ المذکورات کما انّهم فی حقائقهم کانوا انواراً واحدةً و اسراراً واحدةً
و کذلک فاشهد فی ظواهرهم لتعرف کلّهم علی هیکل واحد بل تجدهم علی لفظ واحد و
کلام واحد و بیان واحد

\end{otherlanguage}

Wa-innaka anta law takūn min ahl hādhihi al-madīna fī hādhihi al-lujja
al-aḥadiyya la-tarā kull al-nabiyyīn wa-l-mursalīn ka-haykal wāḥid
wa-nafs wāḥida wa-nūr wāḥid wa-rūḥ wāḥida bi-ḥayth yakūn awwaluhum
ākhirahum wa-ākhiruhum awwalahum wa-kulluhum qāmū ʿalā amr Allāh
wa-sharaʿū sharāʾiʿ ḥikmat Allāh wa-kānū maẓāhir nafs Allāh wa-maʿādin
qudrat Allāh wa-makhāzin waḥy Allāh wa-mashāriq shams Allāh wa-maṭāliʿ
nūr Allāh. Wa-bihim ẓaharat āyāt al-tajrīd fī ḥaqāʾiq al-mumkināt
wa-ʿalāmāt al-tafrīd fī jawhariyyāt al-mawjūdāt wa-ʿanāṣir al-tamjīd fī
dhātiyyāt al-aḥadiyyāt wa-mawāqiʿ al-taḥmīd fī sādhajiyyāt al-ṣamadiyyāt
wa-bihim yabdaʾ al-khalq wa-ilayhim yuʿīd kull al-madhkūrāt. Kamā
annahum fī ḥaqāʾiqihim kānū anwāran wāḥida wa-asrāran wāḥida
wa-kadhālika fa-shhad fī ẓawāhirihim li-taʿrif kullahum ʿalā haykal
wāḥid bal tajidahum ʿalā lafẓ wāḥid wa-kalām wāḥid wa-bayān wāḥid.

If you belong to this city within these depths of oneness, you will see
all the Prophets and Messengers as one temple, one self, one light, and
one spirit, so that their first is their last and their last is their
first. All arose for the Cause of God, ordained the laws of God's
wisdom, and were the Manifestations of God's Self, the mines of God's
power, the treasuries of God's revelation, the dawning places of God's
sun, and the daysprings of God's light. Through them the signs of
detachment appeared in the realities of contingent beings, the tokens of
singleness in the essences of existing things, the elements of
glorification in the essential realities of oneness, and the seats of
praise in the pure realities of eternal self-sufficiency. Through them
He originates creation, and to them He returns all things that can be
mentioned. Just as in their inner realities they are one in light and
one in mystery, so behold them in their outward appearances, that you
may recognize them all as one temple---indeed, find them to be one
utterance, one discourse, and one exposition.

\subsection{47}\label{section-2758}

\begin{otherlanguage}{arabic}

و انّک فی ذلک المقام لو تطلق اوّلهم باسم آخرهم او بالعکس لحقّ کما نزل حکم
ذلک عن مصدر الألوهیّة و منبع الرّبوبیّة قل ادعوا اللّه او ادعوا الرّحمن ایّا
ما تدعوا فله الأسمآء الحسنی لأنّهم مظاهر اسم اللّه و مطالع صفاته و مواقع
قدرته و مجامع سلطنته و انّه جلّ و عزّ بذاته مقدّس عن کلّ الأسمآء و منزّه عن
معارج الصّفات و کذلک فانظر آثار قدرة اللّه فی آفاق ارواحهم و انفس هیاکلهم
لیطمئنّ قلبک و تکون من الّذینهم کانوا فی آفاق القرب لسائرین

\end{otherlanguage}

Wa-innaka fī dhālik al-maqām law tuṭliq awwalahum bi-ism ākhirihim aw
bi-l-ʿaks la-ḥaqq kamā nazala ḥukm dhālik ʿan maṣdar al-ulūhiyya
wa-manbaʿ al-rubūbiyya qul udʿū Allāh aw udʿū al-Raḥmān ayyan mā tadʿū
fa-lahu al-asmāʾ al-ḥusnā li-annahum maẓāhir ism Allāh wa-maṭāliʿ
ṣifātihi wa-mawāqiʿ qudratihi wa-majāmiʿ salṭanatihi wa-innahu jalla
wa-ʿazza bi-dhātihi muqaddas ʿan kull al-asmāʾ wa-munazzah ʿan maʿārij
al-ṣifāt. Wa-kadhālika fa-nẓur āthār qudrat Allāh fī āfāq arwāḥihim
wa-anfus hayākilihim li-yaṭmaʾinn qalbuka wa-takūn min alladhīnahum kānū
fī āfāq al-qurb la-sāʾirīn.

In that station, if you were to call their first by the name of their
last, or the reverse, it would be true, as the decree concerning this
has descended from the Source of Divinity and the Wellspring of
Lordship: ``Say: Call upon God or call upon the All-Merciful; whichever
you call upon, His are the most beautiful names.'' For they are the
manifestations of God's name, the daysprings of His attributes, the
seats of His power, and the focal points of His sovereignty. He, exalted
and mighty, is in His essence sanctified above all names and transcends
the ascending heights of attributes. Behold, then, the traces of God's
power in the horizons of their spirits and in their very temples, that
your heart may be assured and you may be among those who journey through
the horizons of nearness.

\subsection{48}\label{section-2759}

\begin{otherlanguage}{arabic}

ثمّ اجدّد لک الکلام فی هذا المقام لیکون لک معیناً فی عرفانک بارئک فاعلم بأنّ
اللّه تبارک و تعالی لن یظهر بکینونیّته و لا بذاتیّته لم یزل کان مکنوناً فی
قدم ذاته و مخزوناً فی سرمدیّة کینونیّته فلمّا اراد اظهار جماله فی جبروت
الأسمآء و ابراز جلاله فی ملکوت الصّفات ظهر الأنبیآء من الغیب الی الشّهود
لیمتاز اسمه الظّاهر من اسمه الباطن و یظهر اسمه الأوّل عن اسمه الآخر لیکمل
القول بأنّه هو الأوّل و الآخر و الظّاهر و الباطن و هو بکلّ شیء محیط و جعل
مظاهر تلک الأسمآء الکبری و هذه الکلمات العلیا فی مظاهر نفسه و مرایا
کینونته

\end{otherlanguage}

Thumma ujaddid laka al-kalām fī hādhā al-maqām li-yakūn laka muʿīnan fī
ʿirfānika bāriʾaka fa-ʿlam bi-anna Allāh tabāraka wa-taʿālā lan yaẓhar
bi-kaynūniyyatihi wa-lā bi-dhātiyyatihi. Lam yazal kāna maknūnan fī
qidam dhātihi wa-makhzūnan fī sarmadiyyat kaynūniyyatihi. Fa-lammā arāda
iẓhār jamālihi fī jabarūt al-asmāʾ wa-ibrāz jalālihi fī malakūt al-ṣifāt
ẓahara al-anbiyāʾ min al-ghayb ilā al-shuhūd li-yamtāz ismuhu al-Ẓāhir
min ismihi al-Bāṭin wa-yaẓhar ismuhu al-Awwal ʿan ismihi al-Ākhir
li-yakmul al-qawl bi-annahu huwa al-Awwal wa-l-Ākhir wa-l-Ẓāhir
wa-l-Bāṭin wa-huwa bi-kull shayʾ muḥīṭ. Wa-jaʿala maẓāhir tilka al-asmāʾ
al-kubrā wa-hādhihi al-kalimāt al-ʿulyā fī maẓāhir nafsihi wa-marāyā
kaynūnatihi.

I will now renew My discourse to you on this station, that it may assist
you in recognizing your Creator. Know that God, blessed and exalted,
will never appear in His own being or His own essence. He has always
remained concealed in the eternity of His essence and hidden in the
everlastingness of His being. When He wished to reveal His beauty in the
dominion of names and display His majesty in the kingdom of attributes,
the Prophets appeared from the unseen into the visible world, that His
name the Manifest might be distinguished from His name the Hidden, and
His name the First might become apparent apart from His name the Last,
so that the saying might be fulfilled: ``He is the First and the Last,
the Manifest and the Hidden, and He encompasses all things.'' He placed
the manifestations of those greatest names and these most exalted words
within the Manifestations of His Self and the mirrors of His being.

\subsection{49}\label{section-2760}

\begin{otherlanguage}{arabic}

اذاً ثبت بأنّ کلّ الأسمآء و الصّفات ترجع الی هذه الأنوار المقدّسة المتعالیة و
تجد کلّ الأسمآء فی اسمائهم و کلّ الصّفات فی صفاتهم و فی ذلک المقام لو
تدعوهم بکلّ الأسمآء لحقّ بمثل وجودهم اذاً فاعرف ما هو المقصود فی هذا البیان
ثمّ اکتمها فی سرادق قلبک لتعرف حکم ما سألت و تصل الیه علی قدر ما قدّر اللّه
لک لعلّ تکون من الّذینهم کانوا بمراد اللّه لمن الفائزین

\end{otherlanguage}

Idhan thabata bi-anna kull al-asmāʾ wa-l-ṣifāt tarjiʿ ilā hādhihi
al-anwār al-muqaddasa al-mutaʿāliya wa-tajid kull al-asmāʾ fī asmāʾihim
wa-kull al-ṣifāt fī ṣifātihim wa-fī dhālik al-maqām law tadʿūhum bi-kull
al-asmāʾ la-ḥaqq bi-mithl wujūdihim. Idhan fa-ʿrif mā huwa al-maqṣūd fī
hādhā al-bayān thumma uktumhā fī surādiq qalbika li-taʿrif ḥukm mā
saʾalta wa-taṣil ilayhi ʿalā qadr mā qaddara Allāh laka laʿalla takūn
min alladhīnahum kānū bi-murād Allāh la-min al-fāʾizīn.

It is thus established that all names and attributes refer to these
sanctified, transcendent Lights. You find all names in their names and
all attributes in their attributes. In that station, if you were to call
them by every name, it would be true, just as their existence is true.
Recognize, then, what is intended in this exposition, and conceal it
within the pavilion of your heart, that you may understand the matter
about which you asked and attain it to the extent God has ordained for
you. Perhaps you may be among those who attain what God intends.

\subsection{50}\label{section-2761}

\begin{otherlanguage}{arabic}

و کلّ ما سمعت فی ذکر محمّد بن الحسن روح من فی لجج الأرواح فداه حقّ لا ریب
فیه و انّا کلّ به مؤمنون ولکن ذکروا ائمّة الدّین بأنّه کان فی مدینة جابلقا و
وصفوا هذه المدینة بآثار غریبة و علامات عجیبة و انّک لو ترید ان تفسّر هذه
المدینة علی ظاهر الحدیث لن تقدر و لن تجدها ابداً لأنّک لو تفحّص فی اقطار
العالم و اطراف البلاد لن تجدها بأوصاف الّتی وصفوها من قبل ولو تسیر فی
الأرض بدوام ازلیّة اللّه و بقآء سلطنته لأنّ الأرض بتمامها لن تسعها و لن
تحملها و انّک لو تدلّنی الی هذه المدینة انا ادلّک الی هذه النّفس القدسیّة
الّتی عرفوه النّاس بما عندهم لا بما عنده و لمّا انت لن تقدر علی ذلک لا بدّ
لک التّأویل فی هذه الأحادیث و الأخبار المرویّة عن هؤلآء الأنوار و لمّا
تحتاج الی التّأویل فی هذه الحدیث المرویّة فی ذکر هذه المدینة المذکورة و
کذلک تحتاج الی التّفسیر فی هذه النّفس القدمیّة و لمّا عرفت هذا التّأویل لن
تحتاج الی التّبدیل و لا غیره

\end{otherlanguage}

Wa-kull mā samiʿta fī dhikr Muḥammad ibn al-Ḥasan rūḥ man fī lujaj
al-arwāḥ fidāhu ḥaqq lā rayb fīhi wa-innā kull bihi muʾminūn wa-lākin
dhakarū aʾimmat al-dīn bi-annahu kāna fī madīnat Jābulqā wa-waṣafū
hādhihi al-madīna bi-āthār gharība wa-ʿalāmāt ʿajība wa-innaka law turīd
an tufassir hādhihi al-madīna ʿalā ẓāhir al-ḥadīth lan taqdir wa-lan
tajidahā abadan li-annaka law tafaḥḥaṣ fī aqṭār al-ʿālam wa-aṭrāf
al-bilād lan tajidahā bi-awṣāf allatī waṣafūhā min qabl wa-law tasīr fī
al-arḍ bi-dawām azaliyyat Allāh wa-baqāʾ salṭanatihi li-anna al-arḍ
bi-tamāmihā lan tasaʿahā wa-lan taḥmilahā. Wa-innaka law tadullunī ilā
hādhihi al-madīna anā adulluka ilā hādhihi al-nafs al-qudsiyya allatī
ʿarafūhu al-nās bi-mā ʿindahum lā bi-mā ʿindahu wa-lammā anta lan taqdir
ʿalā dhālik lā budda laka al-taʾwīl fī hādhihi al-aḥādīth wa-l-akhbār
al-marwiyya ʿan hāʾulāʾ al-anwār. Wa-lammā taḥtāj ilā al-taʾwīl fī
hādhihi al-ḥadīth al-marwiyya fī dhikr hādhihi al-madīna al-madhkūra
wa-kadhālika taḥtāj ilā al-tafsīr fī hādhihi al-nafs al-qidamiyya
wa-lammā ʿarafta hādhā al-taʾwīl lan taḥtāj ilā al-tabdīl wa-lā
ghayrihi.

Everything you have heard concerning Muḥammad ibn al-Ḥasan---may the
spirit of everyone in the depths of the world of spirits be sacrificed
for him---is true beyond doubt, and we all believe in it. Yet the Imams
of the Faith stated that he was in the city of Jābulqā, and they
described this city as having strange features and wondrous signs. If
you wish to explain this city according to the outward meaning of the
tradition, you will be unable to do so and will never find it. For even
if you search the regions of the world and the farthest reaches of its
lands, you will not find it with the characteristics they described
before, even if you travel the earth for as long as God's eternity and
the permanence of His sovereignty endure. The whole earth could neither
contain nor support it. If you show Me the way to this city, I will show
you the way to this holy Soul, whom people have known according to what
they themselves possess, not according to what he possesses. Since you
cannot do that, you must interpret these traditions and reports handed
down from these Lights. Just as you need interpretation for the
tradition related concerning this city, so too you need an
interpretation concerning this preexistent Soul. Once you understand
this interpretation, you will have no need for alteration or anything
else.

\subsection{51}\label{section-2762}

\begin{otherlanguage}{arabic}

ثمّ اعلم بأنّه لمّا کان الأنبیآء کلّهم روح و نفس و اسم و رسم واحد و انّک بهذا
العین لتری کلّ الظّهورات اسمهم محمّد و آبائهم حسن و ظهروا من جابلقا قدرة
اللّه و یظهروا من جابلصا رحمة اللّه و جابلقا لم یکن الّا خزائن البقآء فی
جبروت العمآء و مدائن الغیب فی لاهوت العلآء و نشهد بأنّ محمّد بن الحسن کان
فی جابلقا و ظهر منها و من یظهره اللّه یکون فیها الی ان یظهره اللّه علی
مقام سلطنته و انّا بذلک مقرّون و بکلّهم مؤمنون و انّا اختصرنا فی معانی
جابلقا فی هذا المقام ولکن تعرف کلّ المعانی فی اسرار هذه الألواح لو تکون
من الموقنین

\end{otherlanguage}

Thumma iʿlam bi-annahu lammā kāna al-anbiyāʾ kulluhum rūḥ wa-nafs wa-ism
wa-rasm wāḥid wa-innaka bi-hādhā al-ʿayn la-tarā kulla al-ẓuhūrāt
ismuhum Muḥammad wa-ābāʾihim Ḥasan wa-ẓaharū min Jābulqā qudrat Allāh
wa-yaẓharū min Jābulṣā raḥmat Allāh wa-Jābulqā lam yakun illā khazāʾin
al-baqāʾ fī jabarūt al-ʿamāʾ wa-madāʾin al-ghayb fī lāhūt al-ʿalāʾ
wa-nashhad bi-anna Muḥammad ibn al-Ḥasan kāna fī Jābulqā wa-ẓahara minhā
wa-man yuẓhiruhu Allāh yakūnu fīhā ilā an yuẓhirahu Allāh ʿalā maqām
salṭanatihi wa-innā bi-dhālika muqirrūn wa-bi-kullihim muʾminūn wa-innā
ikhtaṣarnā fī maʿānī Jābulqā fī hādhā al-maqām wa-lākin taʿrifu kulla
al-maʿānī fī asrār hādhihi al-alwāḥ law takūnu min al-mūqinīn.

Know, then, that since the Prophets are all one spirit, one soul, one
name, and one designation, you will, with this eye, see that all the
Manifestations bear the name Muḥammad, that their fathers are named
Ḥasan, that they have appeared from the Jābulqā of God's power, and that
they appear from the Jābulṣā of God's mercy. Jābulqā is none other than
the treasuries of everlasting life in the dominion of the cloud of
unknowing and the cities of the unseen in the divine realm on high. We
testify that Muḥammad ibn al-Ḥasan was in Jābulqā and appeared from it,
and that He whom God will make manifest is within it until God manifests
Him in the station of His sovereignty. This we acknowledge, and in them
all we believe. We have been brief concerning the meanings of Jābulqā in
this passage, but you will recognize all these meanings in the mysteries
of these Tablets, if you are among those who are certain.

\subsection{52}\label{section-2763}

\begin{otherlanguage}{arabic}

ولکنّ الّذی ظهر فی السّتّین لا تحتاج فی حقّه لا التّبدیل و لا التّأویل لأنّه کان
اسمه محمّد و کان من ابنآء ائمّة الدّین اذاً یصدق فی حقّه بأنّه ابن الحسن و هذا
معلوم عند جنابک و مشهود لدی حضرتک بل انّه خالق الاسم و مبدعه لنفسه لو
انتم بطرف اللّه تنظرون

\end{otherlanguage}

Wa-lākinna alladhī ẓahara fī al-sittīn lā taḥtāju fī ḥaqqihi lā
al-tabdīl wa-lā al-taʾwīl li-annahu kāna ismuhu Muḥammad wa-kāna min
abnāʾ aʾimmat al-dīn idhan yaṣduqu fī ḥaqqihi bi-annahu ibn al-Ḥasan
wa-hādhā maʿlūm ʿinda janābika wa-mashhūd ladā ḥaḍratika bal innahu
khāliq al-ism wa-mubdiʿuhu li-nafsihi law antum bi-ṭarf Allāh tanẓurūn.

But concerning Him who appeared in the year sixty, you have no need of
either alteration or interpretation, for His name was Muḥammad and He
was a descendant of the Imams of the Faith. Thus it is true of Him that
He is the son of Ḥasan. This is known to you and evident to your
presence. Indeed, He is the creator of the name and its originator for
Himself, if you look with the eye of God.

\subsection{53}\label{section-2764}

\begin{otherlanguage}{arabic}

حینئذ اردنا ان نترک ما کنّا فی ذکره و اذکر ما جری علی نقطة الفرقان و نکون
فیه من الذّاکرین و لتکون علی بصیرة فی کلّ الأمور من لدن عزیز جمیل

\end{otherlanguage}

Ḥīnaʾidh aradnā an natruka mā kunnā fī dhikrihi wa-adhkura mā jarā ʿalā
Nuqṭat al-Furqān wa-nakūna fīhi min al-dhākirīn wa-li-takūna ʿalā baṣīra
fī kull al-umūr min ladun ʿazīz jamīl.

At this point We wish to leave what We have been discussing and recount
what befell the Point of the Qur'an, becoming one of those who tell of
it, so that you may gain insight into all matters from One who is mighty
and beautiful.

\subsection{54}\label{section-2765}

\begin{otherlanguage}{persian}

فاعلم ثمّ فکّر ایّامه حین الّذی اقامه اللّه علی امره و اظهره علی مقام نفسه
کیف هجموا علیه العباد و اعترضوا به و حاججوا معه و کلّما مشی قدّامهم فی
المعابر و الأسواق استهزؤوا به و حرّکوا علیه رؤوسهم و سخروا به و فی کلّ حین
ارادوا قتله بحیث ضاقت علیه الأرض بأوسعها و حارت فی امره سکّان ملإ الأعلی
و تبدّلت ارکان البقآء بالفنآء و بکت علیه عیون اهل العمآء و اصابه من هؤلآء
الکفرة الفجرة ما لا یقدر ان یسمعه اولو الوفآء

\end{otherlanguage}

Fa-iʿlam thumma fakkir ayyāmahu ḥīna alladhī aqāmahu Allāh ʿalā amrihi
wa-aẓharahu ʿalā maqām nafsihi kayfa hajamū ʿalayhi al-ʿibād wa-iʿtaraḍū
bihi wa-ḥājajū maʿahu wa-kullamā mashā quddāmahum fī al-maʿābir
wa-al-aswāq istahzaʾū bihi wa-ḥarrakū ʿalayhi ruʾūsahum wa-sakhirū bihi
wa-fī kull ḥīn arādū qatlahu bi-ḥaythu ḍāqat ʿalayhi al-arḍ bi-awsaʿihā
wa-ḥārat fī amrihi sukkān malaʾ al-aʿlā wa-tabaddalat arkān al-baqāʾ
bi-al-fanāʾ wa-bakat ʿalayhi ʿuyūn ahl al-ʿamāʾ wa-aṣābahu min hāʾulāʾ
al-kafara al-fajara mā lā yaqdiru an yasmaʿahu ulū al-wafāʾ.

Know, then, and reflect on His days, when God raised Him up for His
Cause and manifested Him in the station of His own Self: how the people
attacked Him, opposed Him, and disputed with Him. Whenever He walked
before them in the streets and markets, they mocked Him, shook their
heads at Him, and ridiculed Him. At every moment they sought to kill
Him, until the earth, for all its vastness, became narrow for Him. The
inhabitants of the Concourse on high were bewildered at His plight, the
pillars of everlasting life gave way to extinction, and the eyes of the
people of the cloud of unknowing wept over Him. These wicked unbelievers
inflicted on Him what those who are faithful cannot bear to hear.

\subsection{55}\label{section-2766}

\begin{otherlanguage}{arabic}

و لو انّ هؤلآء الفسقة کانوا ان یفکّروا فی امرهم و یعرفوا نغمات تلک الورقآء
علی افنان هذه الشّجرة البیضآء و یرضوا بما نزّل اللّه علیهم فیما انعمهم به و
یجدوا اثمار الشّجرة علی اغصانها لم اعترضوا علیه و انکروه بعد الّذی کلّهم
کانوا ان یرفعوا اعناقهم لبلوغهم الیه و یسألوا اللّه فی کلّ حین بأن یشرّفهم
جماله و یرزقهم لقائه

\end{otherlanguage}

Wa-law anna hāʾulāʾ al-fasaqa kānū an yufakkirū fī amrihim wa-yaʿrifū
naghamāt tilka al-warqāʾ ʿalā afnān hādhihi al-shajara al-bayḍāʾ
wa-yarḍaw bi-mā nazzala Allāh ʿalayhim fīmā anʿamahum bihi wa-yajidū
athmār al-shajara ʿalā aghṣānihā lam iʿtaraḍū ʿalayhi wa-ankarūhu baʿda
alladhī kulluhum kānū an yarfaʿū aʿnāqahum li-bulūghihim ilayhi
wa-yasʾalū Allāh fī kull ḥīn bi-an yusharrifahum jamāluhu wa-yarzuqahum
liqāʾahu.

Had these evildoers reflected on their own condition, recognized the
melodies of that Dove upon the boughs of this snow-white Tree, been
content with what God had sent down to them among His blessings, and
found the fruits of the Tree upon its branches, they would not have
opposed and denied Him, after they had all stretched their necks in
longing to reach Him and had asked God at every moment that His beauty
might honor them and that He might grant them the bounty of meeting Him.

\subsection{56}\label{section-2767}

\begin{otherlanguage}{arabic}

بلی لمّا ما عرفوا لحن الأحدیّة و اسرار الهویّة و اشارات القدسیّة عمّا ظهر عن
لسان الأحمدیّة و ما تفکّروا فی انفسهم و اتّبعوا علمآء الباطل الّذین صدّوا
عباد اللّه عن ادوار القبل و یصدّون النّاس فی اکوار البعد لذا احتجبوا عن
مراد اللّه و ما شربوا عن کوثر الهویّة و صاروا محرومین عن لقآء اللّه و مظهر
کینونته و مطلع ازلیّته و بذلک سلکوا فی مناهج الضّلالة و سبل الغفلة و رجعوا
الی مقرّهم فی نار الّتی کانت وقودها انفسهم و کانوا فی کتاب القدس من قلم
اللّه بالکفر مکتوبا و ما وجدوا و لن یجدوا الی حینئذ لأنفسهم لا من حبیب و
لا من معینا

\end{otherlanguage}

Balā lammā mā ʿarafū laḥn al-aḥadiyya wa-asrār al-huwiyya wa-ishārāt
al-qudsiyya ʿammā ẓahara ʿan lisān al-Aḥmadiyya wa-mā tafakkarū fī
anfusihim wa-ittabaʿū ʿulamāʾ al-bāṭil alladhīna ṣaddū ʿibād Allāh ʿan
adwār al-qabl wa-yaṣuddūna al-nās fī akwār al-baʿd li-dhā iḥtajabū ʿan
murād Allāh wa-mā sharibū ʿan kawthar al-huwiyya wa-ṣārū maḥrūmīn ʿan
liqāʾ Allāh wa-maẓhar kaynūnatihi wa-maṭlaʿ azaliyyatihi wa-bi-dhālika
salakū fī manāhij al-ḍalāla wa-subul al-ghafla wa-rajaʿū ilā maqarrihim
fī nār allatī kānat waqūduhā anfusuhum wa-kānū fī kitāb al-quds min
qalam Allāh bi-al-kufr maktūban wa-mā wajadū wa-lan yajidū ilā ḥīnaʾidh
li-anfusihim lā min ḥabīb wa-lā min muʿīnan.

Indeed, because they did not recognize the melody of oneness, the
mysteries of the divine Essence, or the intimations of holiness in what
issued from the tongue of Aḥmad, and did not reflect within themselves,
but followed the false divines who turned God's servants away in former
cycles and will turn people away in cycles to come, they were veiled
from God's purpose. They did not drink from the river of the divine
Essence and were deprived of meeting God, the Manifestation of His being
and the Dayspring of His eternity. Thus they walked the paths of error
and the ways of heedlessness and returned to their abode in the fire
whose fuel was their own selves. They were recorded as unbelievers by
the Pen of God in the Book of Holiness, and they found no friend or
helper for themselves, nor will they find one until that time.

\subsection{57}\label{section-2768}

\begin{otherlanguage}{persian}

و لو انّ هؤلآء یتمسّکون بنفس عروة اللّه فی قمیص المحمّدیّة و یقبلون الی اللّه
بتمامهم و یلقون کلّ ما فی ایدیهم من علمائهم لیهدیهم اللّه بفضله و یعرّفهم
معانی القدسیّة فی کلماته الأزلیّة لأنّ اللّه اجلّ و اعظم من ان یردّ السّائل عن
بابه او یخیّب الآمل عن فنائه او یطرد من استجار فی ظلّه او یحرم من تشبّث
بذیل رحمته او یبعّد فقیر الّذی نزل فی شریعة غنائه فلمّا هؤلآء ما اقبلوا الی
اللّه بکلّهم و ما تشبّثوا بذیل رحمته المنبسطة فی ظهور شمس الأحدیّة خرجوا عن
ظلّ الهدایة و وردوا فی مدینة الضّلالة و بذلک فسدوا و افسدوا العباد و ضلّوا
و اضلّوا کلّ من فی البلاد و کانوا من الظّالمین فی کتب السّمآء مسطورا

\end{otherlanguage}

Wa-law anna hāʾulāʾ yatamassakūna bi-nafs ʿurwat Allāh fī qamīṣ
al-Muḥammadiyya wa-yuqbilūna ilā Allāh bi-tamāmihim wa-yulqūna kulla mā
fī aydīhim min ʿulamāʾihim la-yahdīhim Allāh bi-faḍlihi wa-yuʿarrifuhum
maʿānī al-qudsiyya fī kalimātihi al-azaliyya li-anna Allāh ajall
wa-aʿẓam min an yarudda al-sāʾil ʿan bābihi aw yukhayyiba al-āmil ʿan
fināʾihi aw yaṭruda man istajāra fī ẓillihi aw yaḥrima man tashabbatha
bi-dhayl raḥmatihi aw yubaʿʿida faqīr alladhī nazala fī sharīʿat
ghināʾihi fa-lammā hāʾulāʾ mā aqbalū ilā Allāh bi-kullihim wa-mā
tashabbathū bi-dhayl raḥmatihi al-munbasiṭa fī ẓuhūr shams al-aḥadiyya
kharajū ʿan ẓill al-hidāya wa-waradū fī madīnat al-ḍalāla wa-bi-dhālika
fasadū wa-afsadū al-ʿibād wa-ḍallū wa-aḍallū kulla man fī al-bilād
wa-kānū min al-ẓālimīn fī kutub al-samāʾ masṭūran.

Had these people held fast to the very Handle of God in the garment of
Muḥammad, turned to God with their whole being, and cast aside
everything they had received from their divines, God would have guided
them by His grace and made known to them the holy meanings within His
eternal words. For God is too exalted and too great to turn a petitioner
away from His door, disappoint one who hopes at His threshold, drive
away one who seeks refuge in His shade, deprive one who clings to the
hem of His mercy, or banish a poor soul who has alighted at the watering
place of His wealth. But because these people did not turn to God with
their whole being or cling to the hem of His mercy spread abroad through
the appearance of the Sun of oneness, they left the shade of guidance
and entered the city of error. Thus they became corrupt and corrupted
His servants, went astray and led astray everyone in the lands, and were
inscribed among the wrongdoers in the books of heaven.

\subsection{58}\label{section-2769}

\begin{otherlanguage}{arabic}

و حینئذ لمّا بلغ هذا الخادم الفانی الی هذا المقام العالی فی بیان رموز
المعانی اذکر لک علّة اعتراض هؤلآء الغلاظ علی غایة الایجاز لیکون دلیلاً
لأولی الألباب من اولی الأبصار و لیکون موهبة من هذا العبد علی المؤمنین
جمیعاً

\end{otherlanguage}

Wa-ḥīnaʾidh lammā balagha hādhā al-khādim al-fānī ilā hādhā al-maqām
al-ʿālī fī bayān rumūz al-maʿānī adhkuru laka ʿillat iʿtirāḍ hāʾulāʾ
al-ghilāẓ ʿalā ghāyat al-ījāz li-yakūna dalīlan li-ulī al-albāb min ulī
al-abṣār wa-li-yakūna mawhiba min hādhā al-ʿabd ʿalā al-muʾminīn
jamīʿan.

Now that this fleeting servant has reached this lofty station in
explaining the allusions to these meanings, I will tell you, with the
utmost brevity, the reason for the opposition of these hardhearted
people, so that it may be a guide for those of understanding among those
endowed with insight, and a gift from this servant to all believers.

\subsection{59}\label{section-2770}

\begin{otherlanguage}{arabic}

فاعلم بأنّ نقطة الفرقان و نور السّبحان لمّا جآء بآیات محکمات و براهین
ساطعات من الآیات الّتی تعجز عنها کلّ من فی جبروت الموجودات امر الکلّ علی
القیام علی هذا الصّراط المرتفعة الممدودة فی کلّ ما جآء به من عند اللّه و من
اقرّ علیه و اعترف بآیات الوحدانیّة فی فؤاده و جمال الأزلیّة فی جماله حکم
علیه حکم البعث و الحشر و الحیاة و الجنّة لأنّه بعد ایمانه باللّه و مظهر
جماله بعث من مرقد غفلته و حشر فی ارض فؤاده و حیّ بحیاة الایمان و الایقان
و دخل فی جنّة اللّقآء هل یکن الجنّة اعلی من ذلک او الحشر اعظم من هذا او
البعث اکبر من هذا البعث لو یطّلع احد بأسراره لیعرف ما لا عرف احد من
العالمین

\end{otherlanguage}

Fa-iʿlam bi-anna Nuqṭat al-Furqān wa-Nūr al-Subḥān lammā jāʾa bi-āyāt
muḥkamāt wa-barāhīn sāṭiʿāt min al-āyāt allatī taʿjizu ʿanhā kull man fī
jabarūt al-mawjūdāt amara al-kull ʿalā al-qiyām ʿalā hādhā al-ṣirāṭ
al-murtafiʿa al-mamdūda fī kull mā jāʾa bihi min ʿind Allāh wa-man
aqarra ʿalayhi wa-iʿtarafa bi-āyāt al-waḥdāniyya fī fuʾādihi wa-jamāl
al-azaliyya fī jamālihi ḥakama ʿalayhi ḥukm al-baʿth wa-al-ḥashr
wa-al-ḥayāt wa-al-janna li-annahu baʿda īmānihi bi-Allāh wa-maẓhar
jamālihi buʿitha min marqad ghaflatihi wa-ḥushira fī arḍ fuʾādihi
wa-ḥayya bi-ḥayāt al-īmān wa-al-īqān wa-dakhala fī jannat al-liqāʾ hal
yakun al-janna aʿlā min dhālika aw al-ḥashr aʿẓam min hādhā aw al-baʿth
akbar min hādhā al-baʿth law yaṭṭaliʿu aḥad bi-asrārihi la-yaʿrifu mā lā
ʿarafa aḥad min al-ʿālamīn.

Know that when the Point of the Qur'an and the Light of the All-Glorious
came with decisive verses and radiant proofs, verses that everyone in
the dominion of existence is powerless to produce, He commanded all to
stand firm upon this exalted path, stretched forth in everything He had
brought from God. Whoever acknowledged Him and recognized the signs of
oneness in His heart and the beauty of eternity in His beauty, He
declared to have attained resurrection, ingathering, life, and paradise.
For after believing in God and the Manifestation of His beauty, such a
person was raised from the grave of heedlessness, gathered into the land
of His heart, quickened with the life of faith and certitude, and
admitted to the paradise of His presence. Is there any paradise higher
than this, any ingathering greater than this, or any resurrection more
momentous than this resurrection? Were anyone to discover its mysteries,
he would know what no one among the peoples of the world has known.

\subsection{60}\label{section-2771}

\begin{otherlanguage}{arabic}

ثمّ اعلم بأنّ هذه الجنّة فی یوم اللّه اعظم من کلّ الجنان و الطف من حقائق
الرّضوان لأنّ اللّه تبارک و تعالی بعد الّذی ختم مقام النّبوّة فی شأن حبیبه و
صفیّه و خیرته من خلقه کما نزل من ملکوت العزّة ولکنّه رسول اللّه و خاتم
النّبیّین وعد العباد بلقائه یوم القیامة لعظمة ظهور البعد کما ظهر بالحقّ و
لم یکن جنّة اعظم من ذلک و لا رتبة اکبر من هذا ان انتم فی آیات القرآن
تتفکّرون فهنیئاً لمن ایقن بلقائه یوم ظهور جماله

\end{otherlanguage}

Thumma iʿlam bi-anna hādhihi al-janna fī yawm Allāh aʿẓam min kull
al-jinān wa-alṭaf min ḥaqāʾiq al-riḍwān li-anna Allāh tabāraka wa-taʿālā
baʿda alladhī khatama maqām al-nubuwwa fī shaʾn ḥabībihi wa-ṣafiyyihi
wa-khīratihi min khalqihi kamā nazala min malakūt al-ʿizza wa-lākinnahu
rasūl Allāh wa-khātam al-nabiyyīn waʿada al-ʿibād bi-liqāʾihi yawm
al-qiyāma li-ʿaẓamat ẓuhūr al-baʿd kamā ẓahara bi-al-ḥaqq wa-lam yakun
janna aʿẓam min dhālika wa-lā rutba akbar min hādhā in antum fī āyāt
al-Qurʾān tatafakkarūn fa-hanīʾan li-man ayqana bi-liqāʾihi yawm ẓuhūr
jamālihi.

Know, moreover, that this paradise in the Day of God is greater than all
paradises and more sublime than the realities of Riḍwān. For God,
blessed and exalted is He, after sealing the station of prophethood in
His Beloved, His chosen One, the One He selected from among His
creation---as was sent down from the kingdom of glory, ``But He is the
Messenger of God and the Seal of the Prophets''---promised His servants
that they would meet Him on the Day of Resurrection, because of the
greatness of the subsequent Revelation, as has indeed come to pass.
There is no paradise greater than this and no rank higher than this, if
you reflect on the verses of the Qur'an. Blessed, then, is the one who
is certain of meeting Him on the day His beauty is revealed.

\subsection{61}\label{section-2772}

\begin{otherlanguage}{persian}

و انّی لو اذکر لک آیات النّازلة فی هذه الرّتبة العالیة لیطول الکلام و تبعد
عن المرام ولکن اذکر هذه الآیة و نکتفی بها لتقرّ عیناک و تصل الی ما کنز
فیها و خزن بها و هی هذه اللّه الّذی رفع السّموات بغیر عمد ترونها ثمّ استوی
علی العرش و سخّر الشّمس و القمر کلّ یجری لأجل مسمّی یدبّر الأمر یفصّل الآیات
لعلّکم بلقآء ربّکم توقنون

\end{otherlanguage}

Wa-innī law adhkuru laka āyāt al-nāzila fī hādhihi al-rutba al-ʿāliya
la-yaṭūlu al-kalām wa-tabʿudu ʿan al-marām wa-lākin adhkuru hādhihi
al-āya wa-naktafī bihā li-taqarra ʿaynāka wa-taṣila ilā mā kuniza fīhā
wa-khuzina bihā wa-hiya hādhihi Allāh alladhī rafaʿa al-samāwāt bi-ghayr
ʿamad tarawnahā thumma istawā ʿalā al-ʿarsh wa-sakhkhara al-shams
wa-al-qamar kull yajrī li-ajal musamman yudabbiru al-amr yufaṣṣilu
al-āyāt laʿallakum bi-liqāʾ rabbikum tūqinūn.

Were I to recount for you the verses sent down concerning this lofty
rank, the discourse would grow long and you would be drawn away from the
purpose. But I will mention this verse, and We will content Ourselves
with it, so that your eyes may be gladdened and you may attain what is
treasured and stored within it. It is this: ``God is He who raised the
heavens without pillars that you can see, then established Himself upon
the throne and subjected the sun and the moon, each running its course
for an appointed term. He directs all affairs and sets out the signs, so
that you may be certain of meeting your Lord.''

\subsection{62}\label{section-2773}

\begin{otherlanguage}{persian}

اذاً فالتفت یا حبیبی فی ذکر الایقان فی هذه الآیة کأنّ السّموات و الأرض و
العرش و الشّمس و القمر کلّهنّ خلقن لایقان العباد لقائه فی ایّامه فواللّه یا
اخی فانظر عظمة هذا المقام و شأن هؤلآء العباد فی هذه الأیّام کأنّهم حمر
مستنفرة فرّت عن طلعة الالهیّة و جمال الهویّة لو تفکّر فیما نزّلناک لتجد ما
اردنا فی ذکر هذا البیان و تعرف ما احببنا ان نعلّمک فی هذا الرّضوان لتقرّ
عیناک عن النّظر فیها و تلذّ سمعک عن استماع ما قرئ فیها و تحظّ نفسک عن
ادراکها و ینوّر قلبک عن عرفانها و تستبشر روحک عن عطر الّذی نفح منها و تصل
الی غایة فیض اللّه و تکون فی رضوان القدس لمن الخالدین

\end{otherlanguage}

Idhan fa-iltafit yā ḥabībī fī dhikr al-īqān fī hādhihi al-āya ka-anna
al-samāwāt wa-al-arḍ wa-al-ʿarsh wa-al-shams wa-al-qamar kullahunna
khuliqna li-īqān al-ʿibād liqāʾahu fī ayyāmihi fa-wa-Allāh yā akhī
fa-unẓur ʿaẓamat hādhā al-maqām wa-shaʾn hāʾulāʾ al-ʿibād fī hādhihi
al-ayyām ka-annahum ḥumur mustanfira farrat ʿan ṭalʿat al-ilāhiyya
wa-jamāl al-huwiyya law tufakkir fīmā nazzalnāka la-tajidu mā aradnā fī
dhikr hādhā al-bayān wa-taʿrifu mā aḥbabnā an nuʿallimaka fī hādhā
al-riḍwān li-taqarra ʿaynāka ʿan al-naẓar fīhā wa-taladhdha samʿuka ʿan
istimāʿ mā quriʾa fīhā wa-taḥuẓẓa nafsuka ʿan idrākihā wa-yunawwara
qalbuka ʿan ʿirfānihā wa-tastabshira rūḥuka ʿan ʿiṭr alladhī nafaḥa
minhā wa-taṣila ilā ghāyat fayḍ Allāh wa-takūna fī riḍwān al-quds la-min
al-khālidīn.

Consider, then, My beloved, the mention of certitude in this verse. It
is as though the heavens, the earth, the throne, the sun, and the moon
were all created so that His servants might be certain of meeting Him in
His days. By God, My brother! Look at the greatness of this station and
the condition of these servants in these days: they are like startled
asses fleeing from the divine Countenance and the beauty of the divine
Essence. If you reflect on what We have sent down to you, you will
discover Our purpose in this utterance and recognize what We have
desired to teach you in this garden, so that your eyes may be gladdened
by gazing upon it, your hearing delighted by listening to what is
recited within it, your soul enriched by comprehending it, your heart
illumined by recognizing it, and your spirit rejoiced by the fragrance
wafting from it; that you may attain the utmost outpouring of God's
grace and be among those who dwell forever in the garden of holiness.

\subsection{63}\label{section-2774}

\begin{otherlanguage}{arabic}

و من اعرض عن اللّه فی حقّه و ادبر و طغی ثمّ کفر و شقی حکم علیه حکم الشّرک و
الکفر و الموت و النّار و ایّ شرک اعظم من اقباله الی مظاهر الشّیطان و اتّباعه
علمآء النّسیان و اصحاب الطّغیان و ایّ کفر اعلی عن اعراضه عن اللّه فی یوم
الّذی یجدّد فیه الایمان من اللّه المقتدر المنّان و ایّ موت اذلّ عن فراره عن
منبع الحیّ الحیوان و ایّ نار احرّ عن بعده عن جمال الهویّة و جلال الأحدیّة فی
یوم التّغابن و الاحسان

\end{otherlanguage}

Wa-man aʿraḍa ʿan Allāh fī ḥaqqihi wa-adbara wa-ṭaghā thumma kafara
wa-shaqiya ḥakama ʿalayhi ḥukm al-shirk wa-al-kufr wa-al-mawt wa-al-nār
wa-ayy shirk aʿẓam min iqbālihi ilā maẓāhir al-shayṭān wa-ittibāʿihi
ʿulamāʾ al-nisyān wa-aṣḥāb al-ṭughyān wa-ayy kufr aʿlā ʿan iʿrāḍihi ʿan
Allāh fī yawm alladhī yujaddadu fīhi al-īmān min Allāh al-muqtadir
al-mannān wa-ayy mawt adhall ʿan firārihi ʿan manbaʿ al-ḥayy al-ḥayawān
wa-ayy nār aḥarr ʿan buʿdihi ʿan jamāl al-huwiyya wa-jalāl al-aḥadiyya
fī yawm al-taghābun wa-al-iḥsān.

But whoever turned away from God in His person, withdrew and rebelled,
then disbelieved and fell into misery, He pronounced upon him the
judgment of idolatry, unbelief, death, and fire. What idolatry is
greater than turning to the manifestations of Satan and following the
divines of forgetfulness and the companions of rebellion? What unbelief
is greater than turning away from God on the day when faith is renewed
by God, the All-Powerful, the All-Bountiful? What death is more abject
than fleeing from the Source of the Ever-Living, the One who is Life
itself? What fire is hotter than remoteness from the beauty of the
divine Essence and the majesty of oneness on the day of loss and
beneficence?

\subsection{64}\label{section-2775}

\begin{otherlanguage}{persian}

و انّ اعراب الجاهلیّة بهذه العبارات و الکلمات اعترضوا علیه و حکموا علیه ما
حکموا و قالوا هؤلآء الّذین آمنوا بمحمّد هم کانوا معنا و راودونا فی کلّ لیل
و نهار متی ماتوا و بأیّ یوم رجعوا فاسمع ما نزل فیما قالوا ان تعجب فعجب
قولهم ائذا کنّا تراباً و عظاماً ائنّا لمبعوثون و فی مقام اخری و لئن قلت انّکم
مبعوثون من بعد الموت لیقولنّ الّذین کفروا ان هذا الّا سحر مبین و بذلک
استهزؤوا به و سخروا علیه لأنّهم شهدوا فی کتبهم و سمعوا من علمائهم لفظ
الموت و الحیاة و فسّروهما بالموت الظّاهریّة و الحیاة العنصریّة فلمّا ما وجدوا
ما عرفوا من ظنونهم المجتثّة و عقولهم الافکیّة الخبیثة رفعوا اعلام الاختلاف
و رایات الفساد و اشتعلوا نار الحرب ولو اطفأها اللّه بقدرته کما تشهد الیوم
من هؤلآء المشرکین و هؤلآء الفاسقین

\end{otherlanguage}

Wa-inna aʿrāb al-jāhiliyya bi-hādhihi al-ʿibārāt wa-al-kalimāt iʿtaraḍū
ʿalayhi wa-ḥakamū ʿalayhi mā ḥakamū wa-qālū hāʾulāʾ alladhīna āmanū
bi-Muḥammad hum kānū maʿanā wa-rāwadūnā fī kull layl wa-nahār matā mātū
wa-bi-ayy yawm rajaʿū fa-ismaʿ mā nazala fīmā qālū in taʿjab fa-ʿajab
qawluhum a-idhā kunnā turāban wa-ʿiẓāman a-innā la-mabʿūthūn wa-fī maqām
ukhrā wa-la-in qulta innakum mabʿūthūn min baʿd al-mawt la-yaqūlanna
alladhīna kafarū in hādhā illā siḥr mubīn wa-bi-dhālika istahzaʾū bihi
wa-sakhirū ʿalayhi li-annahum shahidū fī kutubihim wa-samiʿū min
ʿulamāʾihim lafẓ al-mawt wa-al-ḥayāt wa-fassarūhumā bi-al-mawt
al-ẓāhiriyya wa-al-ḥayāt al-ʿunṣuriyya fa-lammā mā wajadū mā ʿarafū min
ẓunūnihim al-mujtaththa wa-ʿuqūlihim al-ifkiyya al-khabītha rafaʿū aʿlām
al-ikhtilāf wa-rāyāt al-fasād wa-ishtaʿalū nār al-ḥarb wa-law aṭfaʾahā
Allāh bi-qudratihi kamā tashhadu al-yawm min hāʾulāʾ al-mushrikīn
wa-hāʾulāʾ al-fāsiqīn.

The Arabs of the age of ignorance objected to Him on account of these
expressions and words, passed the judgments against Him that they did,
and said, ``These people who have believed in Muḥammad were with us and
associated with us every night and day. When did they die, and on what
day did they return?'' Listen, then, to what was sent down concerning
what they said: ``If you marvel, marvelous indeed is their saying, `When
we have become dust and bones, will we truly be raised up?'\,'' And in
another passage: ``If you say, `You will surely be raised up after
death,' those who disbelieve will say, `This is nothing but manifest
sorcery.'\,'' Thus they mocked and ridiculed Him, for they had seen the
words ``death'' and ``life'' in their books and heard them from their
divines, and had interpreted them as outward death and bodily life. When
they did not find what they had understood through their groundless
suppositions and their lying, malicious minds, they raised the standards
of discord and the banners of corruption and kindled the fire of war,
even though God extinguished it through His power, just as you witness
today from these idolaters and these evildoers.

\subsection{65}\label{section-2776}

\begin{otherlanguage}{arabic}

و انّی حینئذ لمّا هبّت علیّ رائحة الجذب عن مدینة البقآء و احاطتنی غلبات
الشّوق من شطر الاشراق فیما لاحت شمس الآفاق من رکن العراق و اسمعنی نغمات
الحجاز فی اسرار الفراق ارید ان اذکر لجنابک بعض ما غنّت الورقآء فی قطب
العمآء فی معنی الحیاة و الموت ولو انّ هذا ممتنع لأنّی لو ارید ان افسّر لک
کما هو المکتوب فی الواح المحفوظ لن تحمله الأرواح و لن تسعه الأوراق و لن
تطیقه الأرواح ولکن اذکر علی ما ینبغی لهذا الزّمان و هذه الأوان لیکون
دلیلاً لمن اراد ان یدخل فی رفرف المعانی و یسمع نغمات الرّوحانی من هذا
الطّیر المعنویّ الالهی و یکون من الّذینهم انقطعوا الی اللّه و کانوا الیوم
بلقآء اللّه یستبشرون

\end{otherlanguage}

Wa-innī ḥīnaʾidhin lammā habbat ʿalayya rāʾiḥat al-jadhb ʿan madīnat
al-baqāʾ wa-aḥāṭatnī ghalabāt al-shawq min shaṭr al-ishrāq fīmā lāḥat
shams al-āfāq min rukn al-ʿIrāq wa-asmaʿanī naghamāt al-Ḥijāz fī asrār
al-firāq urīdu an adhkura li-janābika baʿḍa mā ghannat al-Warqāʾ fī quṭb
al-ʿamāʾ fī maʿnā al-ḥayāt wa-al-mawt wa-law anna hādhā mumtaniʿ li-annī
law urīdu an ufassira laka kamā huwa al-maktūb fī alwāḥ al-maḥfūẓ lan
taḥmilahu al-arwāḥ wa-lan tasaʿahu al-awrāq wa-lan tuṭīqahu al-arwāḥ
wa-lākin adhkuru ʿalā mā yanbaghī li-hādhā al-zamān wa-hādhihi al-awān
li-yakūna dalīlan li-man arāda an yadkhula fī rafraf al-maʿānī
wa-yasmaʿa naghamāt al-rūḥānī min hādhā al-ṭayr al-maʿnawiyy al-ilāhī
wa-yakūna min alladhīnahum inqaṭaʿū ilā Allāh wa-kānū al-yawm bi-liqāʾ
Allāh yastabshirūn.

Now that the fragrance of rapture has wafted over Me from the city of
eternity, and the surging forces of longing have surrounded Me from the
direction of the dawning light, as the Sun of the horizons has shone
from the region of Iraq and caused Me to hear the melodies of the Hijaz
concerning the mysteries of separation, I wish to mention to you
something of what the Dove has sung at the heart of the cloud of
unknowing concerning the meaning of life and death. Yet this is
impossible, for if I wished to explain it to you as it is written in the
Preserved Tablets, spirits could not bear it, pages could not contain
it, and souls could not endure it. Nevertheless, I will speak as befits
this time and this hour, that it may be a guide for anyone who wishes to
enter the pavilion of meanings, hear spiritual melodies from this
divine, spiritual Bird, and be among those who have detached themselves
from all but God and rejoice today in meeting Him.

\subsection{66}\label{section-2777}

\begin{otherlanguage}{arabic}

فاعرف بأنّ للحیاة مقامین مقام یتعلّق بظاهر البشریّة فی جسد العنصریّة و هذا
معلوم عند جنابک و عند کلّ من علی الأرض بمثل الشّمس فی وسط السّمآء و هذه
الحیاة تفنی من موت الظّاهریّة و هذا حقّ من عند اللّه و لا مفرّ لأحد و امّا
الحیاة الّتی هی المذکور فی کتب الأنبیآء و الأولیآء لم {[}یکن{]} الّا
الحیاة العرفانیّة ای عرفان العبد آیة تجلّی مجلّیه بما تجلّی له به بنفسه و
ایقانه بلقآء اللّه فی مظاهر امره و هذه هی الحیاة الطّیّبة الباقیة الدّائمة
الّتی من یحیی به لن یموت ابداً و یکون باقیاً ببقآء ربّه و دائماً بدوام بارئه

\end{otherlanguage}

Fa-iʿraf bi-anna li-al-ḥayāt maqāmayn maqām yataʿallaqu bi-ẓāhir
al-bashariyya fī jasad al-ʿunṣuriyya wa-hādhā maʿlūm ʿinda janābika
wa-ʿinda kull man ʿalā al-arḍ bi-mithl al-shams fī wasaṭ al-samāʾ
wa-hādhihi al-ḥayāt tafnā min mawt al-ẓāhiriyya wa-hādhā ḥaqq min ʿind
Allāh wa-lā mafarra li-aḥad wa-ammā al-ḥayāt allatī hiya al-madhkūr fī
kutub al-anbiyāʾ wa-al-awliyāʾ lam {[}yakun{]} illā al-ḥayāt
al-ʿirfāniyya ay ʿirfān al-ʿabd āyat tajallī mujallīhi bi-mā tajallā
lahu bihi bi-nafsihi wa-īqānihi bi-liqāʾ Allāh fī maẓāhir amrihi
wa-hādhihi hiya al-ḥayāt al-ṭayyiba al-bāqiya al-dāʾima allatī man yaḥyā
bihi lan yamūta abadan wa-yakūnu bāqiyan bi-baqāʾ rabbihi wa-dāʾiman
bi-dawām bāriʾihi.

Know that life has two stations. One concerns the outward human
condition in the elemental body. This is as well known to you and to
everyone on earth as the sun in the middle of the sky. This life ends
with outward death. This is a truth from God, and no one can escape it.
But the life mentioned in the books of the Prophets and the holy ones is
nothing other than the life of spiritual recognition: the servant's
recognition of the sign of the revelation of the One who reveals Himself
to him, through that by which He has revealed Himself to him in His own
person, and his certainty of meeting God in the Manifestations of His
Cause. This is the good, abiding, everlasting life. Whoever lives
through it will never die, but will abide through the abiding of his
Lord and endure through the endurance of his Creator.

\subsection{67}\label{section-2778}

\begin{otherlanguage}{arabic}

و الحیاة الأوّلیّة الّتی کانت متعلّقة بالجسد العنصریّة ینفد بما نزل من عند
اللّه کلّ نفس ذائقة الموت و الحیاة الثّانویّة الّتی کانت من المعرفة ما تنفد
کما نزل من قبل فلنحیینّه حیاة طیّبة و فی مقام اخری فی ذکر الشّهدآء بل احیآء
عند ربّهم یرزقون و ما ورد فی الأخبار المؤمن حیّ فی الدّارین و بمثل تلک
الکلمات کثیر فی کتب اللّه و مظاهر عدله و انّا ما اردنا ذکرها للاختصار و
اکتفینا بذلک فیما اردنا لک

\end{otherlanguage}

Wa-al-ḥayāt al-awwaliyya allatī kānat mutaʿalliqa bi-al-jasad
al-ʿunṣuriyya yanfadu bi-mā nazala min ʿind Allāh kull nafs dhāʾiqat
al-mawt wa-al-ḥayāt al-thānawiyya allatī kānat min al-maʿrifa mā tanfadu
kamā nazala min qabl fa-la-nuḥyiyannahu ḥayātan ṭayyiba wa-fī maqām
ukhrā fī dhikr al-shuhadāʾ bal aḥyāʾ ʿinda rabbihim yurzaqūn wa-mā
warada fī al-akhbār al-muʾmin ḥayy fī al-dārayn wa-bi-mithl tilka
al-kalimāt kathīr fī kutub Allāh wa-maẓāhir ʿadlihi wa-innā mā aradnā
dhikrahā li-al-ikhtiṣār wa-iktafaynā bi-dhālika fīmā aradnā laka.

The first life, which is connected with the elemental body, comes to an
end, in accordance with what has been sent down from God: ``Every soul
will taste death.'' But the second life, which comes from spiritual
knowledge, does not come to an end, as was sent down before: ``We will
surely grant him a good life,'' and, in another passage concerning the
martyrs: ``Rather, they are alive, sustained in the presence of their
Lord.'' It is also related in the traditions: ``The believer is alive in
both worlds.'' Many such words are found in the books of God and the
Manifestations of His justice. For brevity, We have not wished to
mention them, and have contented Ourselves with this for what We
intended for you.

\subsection{68}\label{section-2779}

\begin{otherlanguage}{arabic}

اذاً یا اخی فأعرض عن هواک ثمّ اقبل الی مولاک و لا تتّبع الّذین کان الههم
هواهم لتدخل فی قطب الحیاة فی ظلّ النّجاة من مربّی الأسمآء و الصّفات لأنّ
الّذینهم الیوم اعرضوا عن ربّهم اموات ولو یمشون علی الأرض و صمّآء ولو یسمعون
و عمیآء ولو یشهدون کما صرّح بذلک مالک یوم الدّین و لهم قلوب لا یفقهون بها
و لهم اعین لا یبصرون بها الی آخر القول بل انّهم یمشون علی شفا جرف هار او
فی شفا حفرة من النّار لم یکن لهم نصیب من هذا البحر المتموّج الذّخّار و کانوا
فی زخارف اقوالهم یلعبون

\end{otherlanguage}

Idhan yā akhī fa-aʿriḍ ʿan hawāka thumma aqbil ilā mawlāka wa-lā
tattabiʿ alladhīna kāna ilāhuhum hawāhum li-tadkhula fī quṭb al-ḥayāt fī
ẓill al-najāt min murabbī al-asmāʾ wa-al-ṣifāt li-anna alladhīnahum
al-yawm aʿraḍū ʿan rabbihim amwāt wa-law yamshūna ʿalā al-arḍ wa-ṣummāʾ
wa-law yasmaʿūn wa-ʿumyāʾ wa-law yashhadūn kamā ṣarraḥa bi-dhālika mālik
yawm al-dīn wa-lahum qulūb lā yafqahūna bihā wa-lahum aʿyun lā yubṣirūna
bihā ilā ākhir al-qawl bal innahum yamshūna ʿalā shafā juruf hār aw fī
shafā ḥufra min al-nār lam yakun lahum naṣīb min hādhā al-baḥr
al-mutamawwij al-dhakhkhār wa-kānū fī zakhārif aqwālihim yalʿabūn.

Then, My brother, turn away from your desires and turn toward your
Master. Do not follow those who have made their desires their god, so
that you may enter the heart of life beneath the shelter of salvation
granted by the Nurturer of names and attributes. For those who have
turned away from their Lord today are dead, even though they walk the
earth; deaf, even though they hear; and blind, even though they see. The
Master of the Day of Judgment has explicitly declared this: ``They have
hearts with which they do not understand, and eyes with which they do
not see,'' and so on to the end of the passage. Indeed, they walk along
the edge of a crumbling bank, or on the brink of a pit of fire. They
have no share of this surging, overflowing ocean, and amuse themselves
with the embellishments of their own words.

\subsection{69}\label{section-2780}

\begin{otherlanguage}{arabic}

و حینئذ نلقی علیک فی هذا المقام فی ذکر الحیاة ما نزّل من قبل لیقلّبک عن
اشارات النّفس و یخلّصک عن ضیق القفس فی هذا الجوار الخنّس و تکون فی ظلمات
الأرض لمن المهتدین

\end{otherlanguage}

Wa-ḥīnaʾidhin nulqī ʿalayka fī hādhā al-maqām fī dhikr al-ḥayāt mā
nuzzila min qabl li-yuqallibaka ʿan ishārāt al-nafs wa-yukhalliṣaka ʿan
ḍīq al-qafaṣ fī hādhā al-jawār al-khunnas wa-takūna fī ẓulumāt al-arḍ
la-min al-muhtadīn.

Now, in this connection, We convey to you what was sent down before
concerning life, so that it may turn you away from the promptings of the
self, free you from the confines of the cage amid these coursing,
receding stars, and make you one of those who are guided through the
darkness of the earth.

\subsection{70}\label{section-2781}

\begin{otherlanguage}{persian}

قال و قوله الحقّ او من کان میتاً فأحییناه و جعلنا له نوراً یمشی به فی النّاس
کمن مثله فی الظّلمات لیس بخارج منها هذه الآیة نزلت فی شأن الحمزة و ابوجهل
لمّا آمن الأوّل و کفر الثّانی و بذلک استهزؤوا اکثر العلمآء من علمآء
الجاهلیّة و تبلبلوا و تهزّلوا و تصاحوا و قالوا کیف مات الحمزة و کیف رجع
الی حیاة الأولی و بمثل ذلک کثیر فی الکتاب لو انتم فی آیات اللّه تتفرّسون

\end{otherlanguage}

Qāla wa-qawluhu al-ḥaqq a-wa-man kāna maytan fa-aḥyaynāhu wa-jaʿalnā
lahu nūran yamshī bihi fī al-nās ka-man mathaluhu fī al-ẓulumāt laysa
bi-khārij minhā hādhihi al-āya nazalat fī shaʾn al-Ḥamza wa-Abū Jahl
lammā āmana al-awwal wa-kafara al-thānī wa-bi-dhālika istahzaʾū akthar
al-ʿulamāʾ min ʿulamāʾ al-jāhiliyya wa-tabalbalū wa-tahazzalū wa-taṣāḥū
wa-qālū kayfa māta al-Ḥamza wa-kayfa rajaʿa ilā ḥayāt al-ūlā wa-bi-mithl
dhālika kathīr fī al-kitāb law antum fī āyāt Allāh tatafarrasūn.

He said---and His word is the truth---``Is someone who was dead, whom We
brought to life and for whom We made a light by which he walks among
people, like someone who is in darkness and cannot emerge from it?''
This verse was sent down concerning Ḥamza and Abū Jahl, when the former
believed and the latter disbelieved. Most of the divines of the age of
ignorance mocked at this, became confused, made jests, and cried out,
saying, ``How did Ḥamza die, and how did he return to his former life?''
There is much like this in the Book, if you examine the verses of God
discerningly.

\subsection{71}\label{section-2782}

\begin{otherlanguage}{arabic}

فیا لیت وجدت قلوباً صافیةً لألقی علیهم رشحاً من ابحر العلم الّذی علّمنی ربّی
لیطیرنّ فی الهوآء کما یمشون علی الأرض و یرکضنّ علی المآء کما یرکضون علی
التّراب و یأخذوا ارواحهم بأیدیهم و یفدوها فی سبیل بارئهم ولکن ما جآء
الاذن علی القضآء فی هذا الرّمز العظمی و لم یزل کان هذا السّرّ مخزوناً فی
کنوز القدرة و هذا الرّمز مکنوناً فی خزائن القوّة لئلّا یهلکون العباد انفسهم
رجآء لهذا المقام الأعظم فی ممالک القدم و لن یصله الّذین یمشون فی ظلمات
الصّیلم المظلم

\end{otherlanguage}

Fa-yā layta wajadtu qulūban ṣāfiya li-ulqiya ʿalayhim rashḥan min abḥur
al-ʿilm alladhī ʿallamanī rabbī li-yaṭīranna fī al-hawāʾ kamā yamshūna
ʿalā al-arḍ wa-yarkuḍanna ʿalā al-māʾ kamā yarkuḍūna ʿalā al-turāb
wa-yaʾkhudhū arwāḥahum bi-aydīhim wa-yafdūhā fī sabīl bāriʾihim wa-lākin
mā jāʾa al-idhn ʿalā al-qaḍāʾ fī hādhā al-ramz al-ʿuẓmā wa-lam yazal
kāna hādhā al-sirr makhzūnan fī kunūz al-qudra wa-hādhā al-ramz maknūnan
fī khazāʾin al-quwwa li-allā yuhlikūna al-ʿibād anfusahum rajāʾan
li-hādhā al-maqām al-aʿẓam fī mamālik al-qidam wa-lan yaṣilahu alladhīna
yamshūna fī ẓulumāt al-ṣaylam al-muẓlim.

If only I could find pure hearts, that I might cast upon them a
sprinkling from the oceans of the knowledge My Lord has taught Me, so
that they might fly through the air as they walk upon the earth, run
upon water as they run upon the dust, take their souls in their hands,
and sacrifice them in the path of their Creator. But permission has not
been given to pronounce upon this supreme mystery. This secret has
always been stored in the treasuries of power, and this mystery
concealed in the repositories of might, lest the servants destroy
themselves in the hope of attaining this most great station in the
realms of eternity. Those who walk in the darkness of blackest night
will never attain it.

\subsection{72}\label{section-2783}

\begin{otherlanguage}{arabic}

و لقد کرّرنا القول یا اخی فی کلّ المقام لیوضح لک باذن اللّه کلّ الأمور عمّا
سطر فی السّطور و لیغنیک عن الّذینهم یخوضون فی انفس الدّیجور و یمشون فی وادی
الکبر و الغرور و لتکون فی فردوس الحیّ الحیوان لمن السّائرین

\end{otherlanguage}

Wa-la-qad karrarnā al-qawl yā akhī fī kull al-maqām li-yūḍiḥa laka
bi-idhn Allāh kull al-umūr ʿammā suṭira fī al-suṭūr wa-li-yughniyaka ʿan
alladhīnahum yakhūḍūna fī anfas al-dayjūr wa-yamshūna fī wādī al-kibr
wa-al-ghurūr wa-li-takūna fī firdaws al-Ḥayy al-Ḥayawān la-min
al-sāʾirīn.

My brother, We have repeated Our words at every point so that, by God's
permission, everything set down in these lines may become clear to you,
that you may have no need of those who plunge into the depths of
darkness and walk in the valley of pride and delusion, and that you may
be among those who journey in the paradise of the Ever-Living, the One
who is Life itself.

\subsection{73}\label{section-2784}

\begin{otherlanguage}{arabic}

قل یا اهل الملإ {[}الظّنّ{]} انّ شجرة الحیاة قد غرست فی وسط فردوس اللّه و
یعطی الحیاة عن کلّ الجهات کیف انتم لا تشعرون و لا تعرفون و یؤیّدک فی کلّ ما
القیناک من جواهر اسرار الهویّة من هذا النّفس المطمئنّة تغنّی حمامة القدس فی
فردوس البقآء و اذکر لک لتلبس قمیص الجدید من زبر الحدید لیحفظک عن رمی
الشّبهات فی تلک الاشارات و هی هذه انّ من لم یلد من المآء و الرّوح لن یقدر
ان یدخل ملکوت اللّه لأنّ المولود من الجسد جسد هو و المولود من الرّوح فهو
روح فلا تتعجّبنّ من قولی انّه ینبغی لکم ان تولدوا مرّة اخری

\end{otherlanguage}

Qul yā ahl al-malaʾ {[}al-ẓann{]} inna shajarat al-ḥayāt qad ghurisat fī
wasaṭ firdaws Allāh wa-yuʿṭī al-ḥayāt ʿan kull al-jihāt kayfa antum lā
tashʿurūn wa-lā taʿrifūn wa-yuʾayyiduka fī kull mā alqaynāka min jawāhir
asrār al-huwiyya min hādhā al-nafs al-muṭmaʾinna taghannī Ḥamāmat
al-Quds fī firdaws al-baqāʾ wa-adhkuru laka li-talbasa qamīṣ al-jadīd
min zubar al-ḥadīd li-yaḥfaẓaka ʿan ramy al-shubuhāt fī tilka al-ishārāt
wa-hiya hādhihi inna man lam yalid min al-māʾ wa-al-rūḥ lan yaqdira an
yadkhula malakūt Allāh li-anna al-mawlūd min al-jasad jasad huwa
wa-al-mawlūd min al-rūḥ fa-huwa rūḥ fa-lā tataʿajjabanna min qawlī
innahu yanbaghī lakum an tūladū marratan ukhrā.

Say: People of the assembly of {[}conjecture{]}! The Tree of Life has
been planted in the midst of God's paradise and bestows life in every
direction. How is it that you neither perceive nor recognize it? The
song of the Dove of Holiness in the paradise of eternity confirms for
you all the jewels of the mysteries of the divine Essence that We have
conveyed to you from this assured Soul. I will recount it to you so that
you may put on a new garment made of blocks of iron, to protect you
against the doubts hurled at those allusions. It is this: ``Whoever is
not born of water and the Spirit cannot enter the Kingdom of God, for
what is born of the flesh is flesh, and what is born of the Spirit is
spirit. So do not marvel at My saying that you must be born again.''

\subsection{74}\label{section-2785}

\begin{otherlanguage}{arabic}

اذاً طیّر الی شجر الالهی و خذ من ثمراتها ثمّ القط عمّا سقط عنها و کن لها
حافظاً امین و فکّر فیما ذکر واحد من الأنبیآء حین الّذی یبشّر الأرواح بمن
یأتی بعده باشارات مقنّعة و رموزات مغطّئة من دون الجهر من القول لتوقن بأن
لا یعرف کلماتهم الّا اولو الألباب الی ان قال کانت عینتاه کلهیب النّار و
کانت رجلاه کالنّحاس و کان یخرج من فمه سیف ذا فمین حینئذ کیف یفسّر هذه
الکلمات و فی الظّاهر لو یجیء احد بتلک العلامات لم یکن بانسان و کیف یستأنس
به احد بل لمّا یظهر فی مدینة تفرّون منه اهل مدینة اخری و لا یقربوا به احد
ابداً مع انّک لو تفکّر فی هذه العبارات لتجدها علی غایة الفصاحة و نهایة
البلاغة بحیث عرجت الی غایة البیان و وصلت الی منتهی مقام التّبیان کأنّ شموس
البلاغة منها ظهرت و انجم الفصاحة عنها بزغت و لاحت

\end{otherlanguage}

Idhan ṭayyir ilā shajar al-ilāhī wa-khudh min thamarātihā thumma ilquṭ
ʿammā saqaṭa ʿanhā wa-kun lahā ḥāfiẓan amīn wa-fakkir fīmā dhakara wāḥid
min al-anbiyāʾ ḥīna alladhī yubashshiru al-arwāḥ bi-man yaʾtī baʿdahu
bi-ishārāt muqannaʿa wa-rumūzāt mughaṭṭaʾa min dūn al-jahr min al-qawl
li-tūqina bi-an lā yaʿrif kalimātihim illā ulū al-albāb ilā an qāla
kānat ʿaynatāhu ka-lahīb al-nār wa-kānat rijlāhu ka-al-nuḥās wa-kāna
yakhruju min famihi sayf dhā famayn ḥīnaʾidhin kayfa yufassiru hādhihi
al-kalimāt wa-fī al-ẓāhir law yajīʾu aḥad bi-tilka al-ʿalāmāt lam yakun
bi-insān wa-kayfa yastaʾnisu bihi aḥad bal lammā yaẓharu fī madīna
tafirrūna minhu ahl madīna ukhrā wa-lā yaqrabū bihi aḥad abadan maʿa
annaka law tufakkir fī hādhihi al-ʿibārāt la-tajiduhā ʿalā ghāyat
al-faṣāḥa wa-nihāyat al-balāgha bi-ḥayth ʿarajat ilā ghāyat al-bayān
wa-waṣalat ilā muntahā maqām al-tibyān ka-anna shumūs al-balāgha minhā
ẓaharat wa-anjum al-faṣāḥa ʿanhā bazaghat wa-lāḥat.

Then fly to the divine Tree, take of its fruits, gather what has fallen
from it, and be its faithful guardian. Reflect on what one of the
Prophets said when he gave souls the glad tidings of the One who would
come after him through veiled allusions and concealed symbols, rather
than speaking openly, so that you may be certain that none but those
endowed with understanding recognize their words. He went on to say:
``His eyes were like a flame of fire, His feet were like brass, and from
His mouth came a two-edged sword.'' How, then, are these words to be
interpreted? If someone were to appear with these signs in their outward
sense, he would not be a human being, and how could anyone feel at ease
with him? Indeed, when he appeared in one city, you, the people of
another city, would flee from him, and no one would ever approach him.
Yet if you reflect on these expressions, you will find them supremely
eloquent and consummately expressive, having ascended to the summit of
utterance and reached the utmost station of exposition. It is as though
the suns of eloquence have appeared from them and the stars of lucid
speech have risen and shone from them.

\subsection{75}\label{section-2786}

\begin{otherlanguage}{arabic}

اذاً فاعرف هؤلآء الحمرآء من امم الماضیة و الّذین یکونون فی تلک الأیّام
ینتظرون مجیء تلک الانسان و لو لا تجیء هذه النّفس علی هذه الصّورة المذکورة
لن یؤمنوا به ابداً و لمّا ما یجیء هذه ابداً انّهم لن یؤمنوا ابداً هذا مبلغ
هؤلآء الکفرة من انفس المشرکة و انّ الّذین ما یعرفون ما هو ابده البدیهیّات و
اظهر الظّاهریّات فکیف یعرفون غوامض اصول الالهیّة و جواهر اسرار حکمة
الصّمدانیّة

\end{otherlanguage}

Idhan fa-iʿraf hāʾulāʾ al-ḥamrāʾ min umam al-māḍiya wa-alladhīna
yakūnūna fī tilka al-ayyām yantaẓirūna majīʾ tilka al-insān wa-law lā
tajīʾu hādhihi al-nafs ʿalā hādhihi al-ṣūra al-madhkūra lan yuʾminū bihi
abadan wa-lammā mā yajīʾu hādhihi abadan innahum lan yuʾminū abadan
hādhā mablagh hāʾulāʾ al-kafara min anfus al-mushrika wa-inna alladhīna
mā yaʿrifūna mā huwa abdah al-badīhiyyāt wa-aẓhar al-ẓāhiriyyāt fa-kayfa
yaʿrifūna ghawāmiḍ uṣūl al-ilāhiyya wa-jawāhir asrār ḥikmat
al-ṣamadāniyya.

Then recognize these foolish ones among the peoples of the past, and
those who in these days await the coming of that Man. Unless this Soul
comes in the form described, they will never believe in Him; and since
He will never come in this way, they will never believe. Such is the
extent of the understanding of these unbelieving, idolatrous souls. How
can those who do not recognize the most self-evident of truths and the
most manifest of things recognize the subtleties of the divine
principles and the jewels of the mysteries of the wisdom of the Eternal?

\subsection{76}\label{section-2787}

\begin{otherlanguage}{arabic}

و انّی حینئذ افسّر لک هذا الکلام علی سبیل الاختصار لتعرف الأسرار و تکون
فیها من العارفین فاعلم ثمّ انصف فیما نلقی الیک لتکون من اهل الانصاف فی
هذا المصاف بین یدی اللّه مذکورا

\end{otherlanguage}

Wa-innī ḥīnaʾidhin ufassiru laka hādhā al-kalām ʿalā sabīl al-ikhtiṣār
li-taʿrifa al-asrār wa-takūna fīhā min al-ʿārifīn fa-iʿlam thumma anṣif
fīmā nulqī ilayka li-takūna min ahl al-inṣāf fī hādhā al-maṣāff bayna
yaday Allāh madhkūran.

I will now briefly explain these words to you, so that you may recognize
the mysteries and be among those who understand them. Learn, then, and
judge fairly what We convey to you, so that in this arena you may be
remembered before God as one of the fair-minded.

\subsection{77}\label{section-2788}

\begin{otherlanguage}{arabic}

فاعلم بأنّ من تکلّم بهذا المقال فی میادین الجلال اراد ان یذکر اوصاف من
یأتی باضمار و الغاز لئلّا یطّلع علیه اهل المجاز فأمّا قوله کانت عینتاه
کلهیب النّار ما اراد الّا حدّة بصر من یأتی و قوّة بصیرته بحیث بعینتاه یحرق
کلّ الحجبات و السّبحات و بها یعرف اسرار القدمیّة فی عوالم الملکیّة و یمیّز
الّذین ترهق وجوههم قترة من الجحیم عن الّذین تعرف فی وجوههم نضرة النّعیم و
لو لم یکن عینتاه من نار اللّه الموقدة کیف یحرق الحجبات و کلّ ما کان بین
ایدی النّاس و یلاحظ آیات اللّه فی جبروت الأسمآء و ملکوت الأشیآء و یشهد
الأشیآء بعین اللّه النّاظرة و کذلک جعلنا الیوم بصره حدیداً ان انتم بآیات
اللّه موقنا و ایّ نار احرّ من هذه النّار الّتی تجلّی فی طور عینتیه و حرق بها
کلّ ما احتجبوا به العباد فی اراضی الایجاد فسبحان اللّه عمّا ظهر فی الواح
السّداد من اسرار المبدإ و المعاد الی یوم الّذی فیه ینادی المناد اذاً انّا کلّ
الی اللّه منقلبون

\end{otherlanguage}

Fa-iʿlam bi-anna man takallama bi-hādhā al-maqāl fī mayādīn al-jalāl
arāda an yadhkura awṣāf man yaʾtī bi-iḍmār wa-alghāz li-allā yaṭṭaliʿa
ʿalayhi ahl al-majāz fa-ammā qawluhu kānat ʿaynatāhu ka-lahīb al-nār mā
arāda illā ḥiddat baṣar man yaʾtī wa-quwwat baṣīratihi bi-ḥayth
bi-ʿaynatāhu yuḥriqu kull al-ḥujubāt wa-al-subuḥāt wa-bihā yaʿrifu asrār
al-qidamiyya fī ʿawālim al-mulkiyya wa-yumayyizu alladhīna tarhaqu
wujūhahum qatara min al-jaḥīm ʿan alladhīna taʿrifu fī wujūhihim naḍrat
al-naʿīm wa-law lam yakun ʿaynatāhu min nār Allāh al-mūqada kayfa
yuḥriqu al-ḥujubāt wa-kull mā kāna bayna aydī al-nās wa-yulāḥiẓu āyāt
Allāh fī jabarūt al-asmāʾ wa-malakūt al-ashyāʾ wa-yashhadu al-ashyāʾ
bi-ʿayn Allāh al-nāẓira wa-ka-dhālika jaʿalnā al-yawm baṣarahu ḥadīdan
in antum bi-āyāt Allāh mūqinan wa-ayyu nār aḥarr min hādhihi al-nār
allatī tajallā fī ṭūr ʿaynatayhi wa-ḥaraqa bihā kull mā iḥtajabū bihi
al-ʿibād fī arāḍī al-ījād fa-subḥān Allāh ʿammā ẓahara fī alwāḥ al-sadād
min asrār al-mabdaʾ wa-al-maʿād ilā yawm alladhī fīhi yunādī al-munād
idhan innā kull ilā Allāh munqalibūn.

Know that the one who spoke these words in the arenas of glory intended
to describe the attributes of the One who would come through concealed
meanings and riddles, so that those attached to outward appearances
would not discover Him. As for his saying, ``His eyes were like a flame
of fire,'' he meant only the keenness of the sight of the One who would
come and the power of His insight, such that with His eyes He burns away
every veil and covering. Through them He recognizes the mysteries of
eternity in the worlds of earthly dominion and distinguishes those whose
faces are overspread with the gloom of hell from those in whose faces
you recognize the radiance of bliss. If His eyes were not of God's
kindled fire, how could He burn away the veils and everything that is in
people's hands, behold the signs of God in the dominion of names and the
kingdom of things, and witness all things through the all-seeing eye of
God? Thus have We made His sight keen today, if you are certain of the
signs of God. What fire is hotter than this fire that has shone upon the
Sinai of His eyes, through which He has burned away everything by which
the servants have veiled themselves in the lands of existence? Glory be
to God for the mysteries of the beginning and the return that have
appeared in the Tablets of rectitude, until the day when the caller
proclaims: ``Then, indeed, we are all returning to God.''

\subsection{78}\label{section-2789}

\begin{otherlanguage}{persian}

و قوله کانت رجلاه کالنّحاس ما اراد بذلک الّا الاستقامة حین الّذی یسمع ندآء
اللّه فاستقم کما امرت لیستقیم علی امر اللّه و یقیم علی صراط قدرة اللّه بحیث
لو ینکروه کلّ من فی السّموات و الأرض ما تزّل قدماه عن التّبلیغ و ما یفرّ عمّا
امره اللّه فی التّشریع و یکون رجلاه کالجبال الباذخة و القلل الشّامخة و یکون
مستحکماً فی طاعة اللّه و قیّوماً فی اظهار امره و ابراز کلمته و لا یردّه منع
مانع و لا یصدّه نهی معرض و لا یندمه انکار کافر و کلّما یشهد من الانکار و
البغضآء و الکفر و الفحشآء یزداد فی محبّة اللّه و یزید الشّوق فی قلبه و یکثر
الوله فی فؤاده و یبوح العشق فی صدره هل شهدت فی الأرض نحاساً احکم من ذلک
او حدیداً اشدّ من ذلک او جبل اسکن من هذا لأنّه یقوم برجلاه فی مقابلة کلّ من
علی الأرض و لا یخاف من احد مع ما انت تعرف فعل العباد فسبحان اللّه مسکنه و
مبعثه و انّه هو المقتدر علی ما یشآء و انّه هو المهیمن القیّوم

\end{otherlanguage}

Wa-qawluhu kānat rijlāhu ka-l-nuḥās mā arāda bi-dhālika illā al-istiqāma
ḥīna alladhī yasmaʿ nidāʾ Allāh fa-staqim kamā umirta li-yastaqīm ʿalā
amr Allāh wa-yuqīm ʿalā ṣirāṭ qudrat Allāh bi-ḥaythu law yunkirūhu kull
man fī al-samāwāt wa-l-arḍ mā tazill qadamāhu ʿan al-tablīgh wa-mā
yafirr ʿammā amarahu Allāh fī al-tashrīʿ wa-yakūn rijlāhu ka-l-jibāl
al-bādhikha wa-l-qulal al-shāmikha wa-yakūn mustaḥkiman fī ṭāʿat Allāh
wa-qayyūman fī iẓhār amrihi wa-ibrāz kalimatihi wa-lā yarudduhu manʿ
māniʿ wa-lā yaṣudduhu nahy muʿriḍ wa-lā yundimuhu inkār kāfir wa-kullamā
yashhad min al-inkār wa-l-baghḍāʾ wa-l-kufr wa-l-faḥshāʾ yazdād fī
maḥabbat Allāh wa-yazīd al-shawq fī qalbihi wa-yakthur al-walah fī
fuʾādihi wa-yabūḥ al-ʿishq fī ṣadrihi hal shahidta fī al-arḍ nuḥāsan
aḥkam min dhālika aw ḥadīdan ashadd min dhālika aw jabal askan min hādhā
li-annahu yaqūm bi-rijlāhi fī muqābalat kull man ʿalā al-arḍ wa-lā
yakhāf min aḥad maʿa mā anta taʿrif fiʿl al-ʿibād fa-subḥān Allāh
musakkinuhu wa-mubʿithuhu wa-innahu huwa al-Muqtadir ʿalā mā yashāʾ
wa-innahu huwa al-Muhaymin al-Qayyūm.

As for his saying, ``His feet were like brass,'' he meant only
steadfastness when He hears the call of God: ``Stand firm, as you have
been commanded,'' so that He may remain steadfast in the Cause of God
and abide on the path of God's power. Even if everyone in the heavens
and on earth were to deny Him, His feet would not falter in proclaiming
the message, nor would He flee from what God has commanded Him to
ordain. His feet would be like towering mountains and lofty peaks. He
would be firm in obedience to God and unwavering in manifesting His
Cause and setting forth His Word. No one's obstruction would turn Him
back, no opponent's prohibition would deter Him, and no unbeliever's
denial would cause Him regret. Whatever denial, hatred, unbelief, and
wickedness He witnesses, His love for God increases, the longing in His
heart grows, the ardor in His inmost being intensifies, and love
declares itself within His breast. Have you seen on earth brass firmer
than this, iron stronger than this, or a mountain more immovable than
this? For He stands on His feet against everyone on earth and fears no
one, despite what you know of people's deeds. Glory be to God, who gives
Him repose and raises Him up! He indeed has power over whatever He
wills, and He is the All-Protecting, the Self-Subsisting.

\subsection{79}\label{section-2790}

\begin{otherlanguage}{arabic}

و کان یخرج من فمه سیف ذا فمین فاعلم بأنّ السّیف لمّا کان آلة القطع و الفصل
و من فم الأنبیآء و الأولیآء یخرج ما یفصل بین المؤمن و الکافر و یقطع بین
المحبّ و المحبوب لذا سمّی بهذا و انّه ما اراد بذلک الّا القطع و الفصل مثلاً
نقطة الأوّلیّة و الشّمس الأزلیّة فی حین الّذی یرید ان یحشر الخلائق باذن اللّه
و یبعثهم من مراقد نفوسهم و یفصل بینهم لینطق بآیة من عند اللّه و هذه الآیة
تفصل بین الحقّ و الباطل من یومئذ الی یوم القیامة و ایّ سیف احدّ من هذا
السّیف الأحدیّة و ایّ صمصام اشحذ من هذا الصّمصام الصّمدیّة الّذی یقطع کلّ النّسبة
و بذلک یفصل بین المقبل و المعرض و بین الأب و الابن و الأخ و الأخت و
العاشق و المعشوق لأنّ من آمن بما نزّل علیه فهو مؤمن و من اعرض فهو کافر و
یظهر الفصل بین هذا المؤمن و هذا الکافر بحیث لا یعاشرا و لا یجتمعا فی
الملک ابداً و کذلک فی الأب و الابن و انّ الابن لو یؤمن و الأب ینکر یفصل
بینهما و لا یجانسا ابداً بل تشهد بأنّ الابن یقتل الأب و بالعکس و کذلک
فاعرف کلّ ما ذکرنا و بیّنّا و فصّلنا

\end{otherlanguage}

Wa-kāna yakhruj min famihi sayf dhā famayn fa-ʿlam bi-anna al-sayf lammā
kāna ālat al-qaṭʿ wa-l-faṣl wa-min fam al-anbiyāʾ wa-l-awliyāʾ yakhruj
mā yafṣil bayna al-muʾmin wa-l-kāfir wa-yaqṭaʿ bayna al-muḥibb
wa-l-maḥbūb li-dhā summiya bi-hādhā wa-innahu mā arāda bi-dhālika illā
al-qaṭʿ wa-l-faṣl mathalan nuqṭat al-awwaliyya wa-l-shams al-azaliyya fī
ḥīn alladhī yurīd an yaḥshur al-khalāʾiq bi-idhn Allāh wa-yabʿathahum
min marāqid nufūsihim wa-yafṣil baynahum la-yanṭiq bi-āya min ʿind Allāh
wa-hādhihi al-āya tafṣil bayna al-ḥaqq wa-l-bāṭil min yawmiʾidh ilā yawm
al-qiyāma wa-ayy sayf aḥadd min hādhā al-sayf al-aḥadiyya wa-ayy ṣamṣām
ashḥadh min hādhā al-ṣamṣām al-ṣamadiyya alladhī yaqṭaʿ kull al-nisba
wa-bi-dhālika yafṣil bayna al-muqbil wa-l-muʿriḍ wa-bayna al-ab wa-l-ibn
wa-l-akh wa-l-ukht wa-l-ʿāshiq wa-l-maʿshūq li-anna man āmana bi-mā
nuzzila ʿalayhi fa-huwa muʾmin wa-man aʿraḍ fa-huwa kāfir wa-yaẓhar
al-faṣl bayna hādhā al-muʾmin wa-hādhā al-kāfir bi-ḥaythu lā yuʿāshirā
wa-lā yajtamiʿā fī al-mulk abadan wa-ka-dhālika fī al-ab wa-l-ibn
wa-inna al-ibn law yuʾmin wa-l-ab yunkir yufṣal baynahumā wa-lā
yatajānasā abadan bal tashhad bi-anna al-ibn yaqtul al-ab wa-bi-l-ʿaks
wa-ka-dhālika fa-ʿrif kull mā dhakarnā wa-bayyannā wa-faṣṣalnā.

``And from His mouth came a two-edged sword.'' Know that, since the
sword is an instrument of cutting and separation, and from the mouths of
the Prophets and the holy ones comes what separates the believer from
the unbeliever and severs the lover from the beloved, it was given this
name. By this he meant only cutting and separation. For example, when
the Primal Point and the eternal Sun wishes to gather together all
creatures by God's leave, raise them from the resting places of their
own selves, and separate them from one another, He utters a verse from
God. This verse separates truth from falsehood from that day until the
Day of Resurrection. What sword is sharper than this sword of oneness,
and what blade keener than this blade of eternal self-sufficiency, which
severs every tie? Through it He separates those who turn toward Him from
those who turn away, father from son, brother from sister, and lover
from beloved. For whoever believes in what has been sent down to Him is
a believer, and whoever turns away is an unbeliever. The separation
between this believer and this unbeliever becomes such that they will
never associate or come together in the world. So too with father and
son: if the son believes and the father denies, they are separated and
will never be of one kind. Indeed, you see the son killing the father,
and the reverse. In this way, understand everything We have mentioned,
explained, and set forth in detail.

\subsection{80}\label{section-2791}

\begin{otherlanguage}{arabic}

و انّک لو تشهد بعین الیقین لتشهد بأنّ هذا السّیف الالهی لیفصل بین الأصلاب
لو انتم تعلمون و هذه من کلمة الفصل الّتی تظهر فی یوم الفصل و الطّلاق لو
کانوا النّاس فی ایّام ربّهم یتذکّرون بل لو تدقّ بصرک و ترقّ قلبک لتشهد بأنّ کلّ
السّیوف الظّاهریّة الّتی تقتل الکفّار و تجاهد مع الفجّار فی کلّ دهر و زمان یظهر
من هذا السّیف الباطنیّة الالهیّة اذاً فافتح عیناک لتجد کلّ ما اریناک و تبلغ
الی ما لا یبلغ الیه احد من العالمین و نقول الحمد لله اذ هو مالک یوم
الدّین

\end{otherlanguage}

Wa-innaka law tashhad bi-ʿayn al-yaqīn la-tashhad bi-anna hādhā al-sayf
al-ilāhī la-yafṣil bayna al-aṣlāb law antum taʿlamūn wa-hādhihi min
kalimat al-faṣl allatī taẓhar fī yawm al-faṣl wa-l-ṭalāq law kānū al-nās
fī ayyām rabbihim yatadhakkarūn bal law tudiqq baṣaraka wa-turiqq
qalbaka la-tashhad bi-anna kull al-suyūf al-ẓāhiriyya allatī taqtul
al-kuffār wa-tujāhid maʿa al-fujjār fī kull dahr wa-zamān yaẓhar min
hādhā al-sayf al-bāṭiniyya al-ilāhiyya idhan fa-ftaḥ ʿaynāka li-tajid
kull mā araynāka wa-tablugh ilā mā lā yablugh ilayhi aḥad min al-ʿālamīn
wa-naqūl al-ḥamdu li-llāh idh huwa mālik yawm al-dīn.

If you were to behold with the eye of certainty, you would see that this
divine sword divides even the loins, if only you knew. This proceeds
from the word of separation that appears on the day of separation and
sundering, if only people would take heed in the days of their Lord.
Indeed, if you sharpened your sight and softened your heart, you would
see that all the outward swords that slay unbelievers and wage war
against the wicked in every age and time proceed from this inward,
divine sword. Open your eyes, then, to discover everything We have shown
you and attain what no one among all peoples has attained. And We say:
``Praise be to God, for He is the Master of the Day of Judgment.''

\subsection{81}\label{section-2792}

\begin{otherlanguage}{arabic}

و هؤلآء العباد لمّا ما اخذوا العلم من معدنه و محلّه و عن بحر العذب الفرات
السّائغ الّذی یجری باذن اللّه فی قلوب الصّافیة السّاذجیّة لذا احتجبوا عن مراد
اللّه فی کلماته و اشاراته و کانوا فی سجن انفسهم لساکنین

\end{otherlanguage}

Wa-hāʾulāʾ al-ʿibād lammā mā akhadhū al-ʿilm min maʿdinihi wa-maḥallihi
wa-ʿan baḥr al-ʿadhb al-furāt al-sāʾigh alladhī yajrī bi-idhn Allāh fī
qulūb al-ṣāfiya al-sādhijiyya li-dhā iḥtajabū ʿan murād Allāh fī
kalimātihi wa-ishārātihi wa-kānū fī sijn anfusihim la-sākinīn.

Because these servants did not draw knowledge from its source and proper
place, or from the sweet, fresh, palatable sea that flows by God's leave
within pure and unalloyed hearts, they were veiled from God's intention
in His words and allusions, and remained dwelling in the prison of their
own selves.

\subsection{82}\label{section-2793}

\begin{otherlanguage}{arabic}

و انّا نشکر اللّه بما اتانا من فضله و جعلنا موقناً بأمره الّذی لا یقوم معه
السّموات و الأرض و مقرّاً به یوم لقائه و بمن یظهره اللّه فی قیامة الأخری و
جعلنا من الموقنین به قبل ظهوره لتکون النّعمة من عنده بالغة علینا و علی
العالمین

\end{otherlanguage}

Wa-innā nashkur Allāh bi-mā ātānā min faḍlihi wa-jaʿalanā mūqinan
bi-amrihi alladhī lā yaqūm maʿahu al-samāwāt wa-l-arḍ wa-muqirran bihi
yawm liqāʾihi wa-bi-man yuẓhiruhu Allāh fī qiyāmat al-ukhrā wa-jaʿalanā
min al-mūqinīn bihi qabla ẓuhūrihi li-takūn al-niʿma min ʿindihi bāligha
ʿalaynā wa-ʿalā al-ʿālamīn.

We thank God for what He has bestowed upon Us through His grace, and for
making Us certain of His Cause, before which the heavens and the earth
cannot stand, acknowledging Him on the day of meeting Him and
acknowledging Him whom God will make manifest in the next Resurrection.
He has made Us among those certain of Him before His manifestation, so
that His favor to Us and to all the worlds may be complete.

\subsection{83}\label{section-2794}

\begin{otherlanguage}{persian}

ولکن اشکو الیک یا اخی عن الّذین ینسبون انفسهم الی اللّه و مظاهر علمه و
یرتکبون الفواحش و یأکلون اموال النّاس و یشربون الخمر و یقتلون الأنفس و
یسرقون الأموال بینهم و یغتبون بعضهم بعضاً و یفترون علی اللّه و یکذبون فی
اکثر اقوالهم و یرجع النّاس کلّ ذلک الینا و انّهم ما {[}استحیون{]} عن اللّه و
یترکون ما امرهم اللّه و یرتکبون ما نهوا عنه بعد الّذی ینبغی لأهل الحقّ بأن
یظهر آثار الخضوع عن وجوههم و انوار القدس من طلعاتهم و یمشوا فی الأرض
بمثل من یمشی بین یدی اللّه و یکون ممتازاً عن کلّ من علی الأرض بجمیع الحرکات
و السّکنات بحیث یشاهدوا آثار القدرة بعیونهم و یذکروا اللّه بألسنهم و
قلوبهم و یمشوا الی اوطان القرب بأرجلهم و یأخذوا احکام اللّه بأیادیهم و لو
یمضون علی وادی الذّهب و معادن الفضّة ما یعتنون بهما و لا یلتفتون الیهما

\end{otherlanguage}

Wa-lākin ashkū ilayka yā akhī ʿan alladhīna yansibūn anfusahum ilā Allāh
wa-maẓāhir ʿilmihi wa-yartakibūn al-fawāḥish wa-yaʾkulūn amwāl al-nās
wa-yashrabūn al-khamr wa-yaqtulūn al-anfus wa-yasriqūn al-amwāl baynahum
wa-yaghtābūn baʿḍuhum baʿḍan wa-yaftarūn ʿalā Allāh wa-yakdhibūn fī
akthar aqwālihim wa-yurjiʿ al-nās kull dhālika ilaynā wa-innahum mā
{[}istaḥyūn{]} ʿan Allāh wa-yatrukūn mā amarahum Allāh wa-yartakibūn mā
nuhū ʿanhu baʿda alladhī yanbaghī li-ahl al-ḥaqq bi-an yaẓhar āthār
al-khuḍūʿ ʿan wujūhihim wa-anwār al-quds min ṭalaʿātihim wa-yamshū fī
al-arḍ bi-mithl man yamshī bayna yaday Allāh wa-yakūn mumtāzan ʿan kull
man ʿalā al-arḍ bi-jamīʿ al-ḥarakāt wa-l-sakanāt bi-ḥaythu yushāhidū
āthār al-qudra bi-ʿuyūnihim wa-yadhkurū Allāh bi-alsunihim wa-qulūbihim
wa-yamshū ilā awṭān al-qurb bi-arjulihim wa-yaʾkhudhū aḥkām Allāh
bi-ayādīhim wa-law yamḍūn ʿalā wādī al-dhahab wa-maʿādin al-fiḍḍa mā
yaʿtanūn bi-himā wa-lā yaltafitūn ilayhimā.

But I complain to you, My brother, of those who claim affiliation with
God and the Manifestations of His knowledge, yet commit indecencies,
devour people's wealth, drink wine, take lives, steal property from one
another, backbite one another, invent falsehoods about God, and lie in
most of what they say. People attribute all this to Us. They feel no
shame before God, abandon what God has commanded them, and commit what
they have been forbidden, whereas the people of truth should show signs
of humility in their faces and the lights of holiness in their
countenances. They should walk the earth like one walking in God's
presence and be distinguished from everyone on earth in every movement
and moment of stillness, so that with their eyes they witness the signs
of His power, with their tongues and hearts they remember God, with
their feet they walk toward the homelands of nearness, and with their
hands they take hold of God's ordinances. Were they to pass through a
valley of gold and mines of silver, they would pay them no heed and
would not turn toward them.

\subsection{84}\label{section-2795}

\begin{otherlanguage}{arabic}

و انّ هؤلآء اعرضوا عن کلّ ذلک و اقبلوا الی ما تهوی به هواهم و انّهم فی وادی
الکبر و الغرور لیهیمون و اشهد حینئذ بأنّ اللّه کان بریء عنهم و نحن برآء و
نسأل اللّه بأن لا یجمعنا و ایّاهم لا فی الدّنیا و لا فی الآخرة اذ انّه هو
الحقّ لا اله الّا هو و انّه کان علی کلّ شیء قدیر

\end{otherlanguage}

Wa-inna hāʾulāʾ aʿraḍū ʿan kull dhālika wa-aqbalū ilā mā tahwā bihi
hawāhum wa-innahum fī wādī al-kibr wa-l-ghurūr la-yahīmūn wa-ashhad
ḥīnaʾidh bi-anna Allāh kāna barīʾ ʿanhum wa-naḥnu barāʾ wa-nasʾal Allāh
bi-an lā yajmaʿanā wa-iyyāhum lā fī al-dunyā wa-lā fī al-ākhira idh
innahu huwa al-Ḥaqq lā ilāha illā huwa wa-innahu kāna ʿalā kull shayʾ
qadīr.

Yet these people have turned away from all this and turned toward
whatever their desires crave. They wander in the valley of pride and
delusion. I testify, therefore, that God is clear of them, and We are
clear of them. We ask God not to bring Us together with them, either in
this world or in the next. He indeed is the Truth; there is no god but
Him, and He has power over all things.

\subsection{85}\label{section-2796}

\begin{otherlanguage}{arabic}

اذاً فاشرب یا اخی من هذا المآء الّذی اجریناه فی ابحر تلک الکلمات کأنّ بحور
العظمة متموّجات فیها و جواهر الأحدیّة مشعشعات لها و بها و علیها فانّک فاخلع
ثیابک عمّا یحجبک عن الدّخول فی هذا البحر اللّجّیّ الحمرآء فقل بسم اللّه و
باللّه ثمّ ادخل فیها و لا تخف من احد و توکّل علی اللّه ربّک و من یتوکّل علی
اللّه فهو حسبه فانّه هو یحفظک و تکون فیه من الآمنین

\end{otherlanguage}

Idhan fa-shrab yā akhī min hādhā al-māʾ alladhī ajraynāhu fī abḥur tilka
al-kalimāt ka-anna buḥūr al-ʿaẓama mutamawwijāt fīhā wa-jawāhir
al-aḥadiyya mushaʿshiʿāt lahā wa-bihā wa-ʿalayhā fa-innaka fa-khlaʿ
thiyābaka ʿammā yaḥjubuka ʿan al-dukhūl fī hādhā al-baḥr al-lujjī
al-ḥamrāʾ fa-qul bi-smi Allāh wa-bi-llāh thumma udkhul fīhā wa-lā takhaf
min aḥad wa-tawakkal ʿalā Allāh rabbika wa-man yatawakkal ʿalā Allāh
fa-huwa ḥasbuhu fa-innahu huwa yaḥfaẓuka wa-takūn fīhi min al-āminīn.

Drink, then, My brother, of this water that We have caused to flow in
the seas of these words. It is as though the oceans of grandeur surge
within them and the jewels of oneness shine for them, through them, and
upon them. Strip off the garments of whatever prevents you from entering
this deep crimson sea. Say, ``In the name of God and through God,'' then
enter it. Fear no one, and place your trust in God, your Lord. ``Whoever
places his trust in God, He is sufficient for him.'' He will indeed
protect you, and within it you will be among those who are safe.

\subsection{86}\label{section-2797}

\begin{otherlanguage}{persian}

ثمّ اعلم بأنّ فی هذه المدینة الألطف الأبهی تجد السّالک خاضعاً لکلّ الوجوه و
خاشعاً لکلّ الأشیآء لأنّه لا یشهد شیئاً الّا و قد یری اللّه فیه و یشهد نوره
فیما احاطت انوار الظّهور علی طور الممکنات و فی ذلک المقام حقّ علیه بأن لا
یجلس علی صدور المجالس لافتخار نفسه و لا یتقدّم علی نفس لاستکبار نفسه و
یشهد نفسه فی کلّ حین بین یدی مولاه و لا یرضی لوجه ما لا یرضی لوجهه و لا
یقول لأحد ما لا یقدر ان یسمعه من غیره و لا یحبّ لأحد ما لا یحبّه لنفسه و
یحرّک فی الأرض علی خیط الاستوآء فی ملکوت البدآء

\end{otherlanguage}

Thumma iʿlam bi-anna fī hādhihi al-madīna al-alṭaf al-abhā tajid
al-sālik khāḍiʿan li-kull al-wujūh wa-khāshiʿan li-kull al-ashyāʾ
li-annahu lā yashhad shayʾan illā wa-qad yarā Allāh fīhi wa-yashhad
nūrahu fī-mā aḥāṭat anwār al-ẓuhūr ʿalā ṭūr al-mumkināt wa-fī dhālika
al-maqām ḥaqq ʿalayhi bi-an lā yajlis ʿalā ṣudūr al-majālis li-iftikhār
nafsihi wa-lā yataqaddam ʿalā nafs li-istikbār nafsihi wa-yashhad
nafsahu fī kull ḥīn bayna yaday mawlāhu wa-lā yarḍā li-wajh mā lā yarḍā
li-wajhihi wa-lā yaqūl li-aḥad mā lā yaqdir an yasmaʿahu min ghayrihi
wa-lā yuḥibb li-aḥad mā lā yuḥibbuhu li-nafsihi wa-yuḥarrik fī al-arḍ
ʿalā khayṭ al-istiwāʾ fī malakūt al-badāʾ.

Know, moreover, that in this most subtle and most glorious city you will
find the wayfarer humble before every face and reverent before all
things, for he beholds nothing without seeing God within it, and
witnesses His light in all that the lights of revelation have
encompassed upon the Sinai of contingent existence. In this station he
must not occupy the seats of honor in gatherings to glorify himself, nor
take precedence over anyone out of self-importance. He must see himself
at every moment in his Lord's presence, not approve for anyone what he
would not approve for himself, not say to anyone what he could not bear
to hear from another, and not wish for anyone what he would not wish for
himself. He must move upon the earth along the line of rectitude in the
kingdom of unfolding decree.

\subsection{87}\label{section-2798}

\begin{otherlanguage}{arabic}

ولکن اعلم بأنّ السّالک فی اوائل سلوکه کما ذکرنا من قبل لیری التّبدیل و
التّغییر و هذا حقّ لا ریب فیه کما نزل فی وصف تلک الأیّام یوم تبدّل الأرض غیر
الأرض و هذا من ایّام الّذی ما شهدت العیون بمثلها فطوبی لمن ادرکها و عرف
قدرها و لقد ارسلنا موسی بآیاتنا ان اخرج القوم من الظّلمات الی النّور و
ذکّرهم بأیّام اللّه و هذا من ایّام اللّه لو انتم تعرفون

\end{otherlanguage}

Wa-lākin iʿlam bi-anna al-sālik fī awāʾil sulūkihi kamā dhakarnā min
qabl la-yarā al-tabdīl wa-l-taghyīr wa-hādhā ḥaqq lā rayb fīhi kamā
nazala fī waṣf tilka al-ayyām yawm tubaddal al-arḍ ghayr al-arḍ wa-hādhā
min ayyām alladhī mā shahidat al-ʿuyūn bi-mithlihā fa-ṭūbā li-man
adrakahā wa-ʿarafa qadrahā wa-la-qad arsalnā Mūsā bi-āyātinā an akhrij
al-qawm min al-ẓulumāt ilā al-nūr wa-dhakkirhum bi-ayyām Allāh wa-hādhā
min ayyām Allāh law antum taʿrifūn.

But know that, at the beginning of his journey, the wayfarer indeed sees
transformation and change, as We mentioned before. This is a truth
beyond doubt, as was revealed in describing those days: ``The day when
the earth is changed into another earth.'' This is one of those days
whose like no eyes have witnessed. Blessed is the one who attains them
and recognizes their worth. ``We indeed sent Moses with Our signs:
`Bring the people out of darkness into light, and remind them of the
days of God.'\,'' This is one of the days of God, if only you knew.

\subsection{88}\label{section-2799}

\begin{otherlanguage}{arabic}

و فی هذا المقام کلّ المتغایرات و المتبدّلات لموجود بین یدیک و من اقرّ بغیر
ذلک فقد الحد فی امر اللّه و نازعه فی سلطانه و حاربه فی حکومته و من یبدّل
الأرض و یجعلها غیر الأرض لیقدر ان یبدّل کلّ ما علیها و ما یحرّک علی ظهرها و
لا تستعجب عن ذلک کما بدّل الظّلمة بالنّور و النّور بالظّلمة و الجهل بالعلم و
الضّلالة بالهدایة و الموت بالحیاة و الحیاة بالموت و فی ذلک المقام یثبت
حکم التّبدیل ان تکون من اهل هذا السّبیل فکّر فیه لیظهر لک ما طلبت عن هذا
الذّلیل من سرادق هذا الدّلیل لتکون فیه من السّاکنین لأنّه یفعل ما یشآء و
یحکم ما یرید و لا یسأل عمّا یفعل و کلّ عن کلّ یسألون

\end{otherlanguage}

Wa-fī hādhā al-maqām kull al-mutaghāyirāt wa-l-mutabaddilāt la-mawjūd
bayna yadayka wa-man aqarra bi-ghayr dhālika fa-qad alḥada fī amr Allāh
wa-nāzaʿahu fī sulṭānihi wa-ḥārabahu fī ḥukūmatihi wa-man yubaddil
al-arḍ wa-yajʿaluhā ghayr al-arḍ la-yaqdir an yubaddil kull mā ʿalayhā
wa-mā yuḥarrik ʿalā ẓahrihā wa-lā tastaʿjib ʿan dhālika kamā baddala
al-ẓulma bi-l-nūr wa-l-nūr bi-l-ẓulma wa-l-jahl bi-l-ʿilm wa-l-ḍalāla
bi-l-hidāya wa-l-mawt bi-l-ḥayā wa-l-ḥayā bi-l-mawt wa-fī dhālika
al-maqām yathbut ḥukm al-tabdīl in takūn min ahl hādhā al-sabīl fakkir
fīhi li-yaẓhar laka mā ṭalabta ʿan hādhā al-dhalīl min surādiq hādhā
al-dalīl li-takūn fīhi min al-sākinīn li-annahu yafʿal mā yashāʾ
wa-yaḥkum mā yurīd wa-lā yusʾal ʿammā yafʿal wa-kull ʿan kull yusʾalūn.

In this station, all changing and transformed things are present before
you. Whoever affirms otherwise has strayed in the Cause of God, disputed
His sovereignty, and waged war against His rule. He who changes the
earth and makes it another earth is indeed able to change everything
upon it and everything moving across its surface. Do not be astonished
at this, for He has changed darkness into light and light into darkness,
ignorance into knowledge, error into guidance, death into life, and life
into death. In this station the rule of transformation is established,
if you are among those who follow this path. Reflect upon it, so that
what you have sought from this lowly One may appear to you from the
pavilion of this proof, and you may be among those dwelling within it.
For He does whatever He wills and ordains whatever He desires. He is not
questioned about what He does, while all are questioned about
everything.

\subsection{89}\label{section-2800}

\begin{otherlanguage}{arabic}

ولکن یا اخی لتری فی هذه الرّتبة ای فی اوّل السّلوک کما ذکرنا فی مدینة الطّلب
مقامات مختلفة و علامات متفاوتة و کلّها حقّ فی مواقعها و مقاماتها و ینبغی
لجنابک فی هذا المقام بأن تشهد کلّ الأشیآء فی اماکنها من دون ان تنزل شیئاً
عن صعودها و علوّها او ترفع شیئاً من مقامها و دنوّها مثلاً انّک لو تحلّ اللّاهوت
فی النّاسوت هذا شرک محض و لو تصعد النّاسوت الی هوآء اللّاهوت هذا کفر صرف
ولکن لو تذکر اللّاهوت فی اللّاهوت و النّاسوت فی النّاسوت لحقّ لا ریب فیه ای
انّ جنابک لو تشهد التّبدیل فی عوالم التّوحید هذا ذنب لم یکن فی الملک اکبر
من ذلک و ان تشهد التّبدیل فی مقامه و تعرفه علی ما ینبغی لا بأس علیک

\end{otherlanguage}

Wa-lākin yā akhī la-tarā fī hādhihi al-rutba ay fī awwal al-sulūk kamā
dhakarnā fī madīnat al-ṭalab maqāmāt mukhtalifa wa-ʿalāmāt mutafāwita
wa-kulluhā ḥaqq fī mawāqiʿihā wa-maqāmātihā wa-yanbaghī li-janābika fī
hādhā al-maqām bi-an tashhad kull al-ashyāʾ fī amākinihā min dūn an
tunazzil shayʾan ʿan ṣuʿūdihā wa-ʿuluwwihā aw tarfaʿ shayʾan min
maqāmihā wa-dunuwwihā mathalan innaka law tuḥill al-lāhūt fī al-nāsūt
hādhā shirk maḥḍ wa-law tuṣʿid al-nāsūt ilā hawāʾ al-lāhūt hādhā kufr
ṣirf wa-lākin law tadhkur al-lāhūt fī al-lāhūt wa-l-nāsūt fī al-nāsūt
la-ḥaqq lā rayb fīhi ay inna janābaka law tashhad al-tabdīl fī ʿawālim
al-tawḥīd hādhā dhanb lam yakun fī al-mulk akbar min dhālika wa-in
tashhad al-tabdīl fī maqāmihi wa-taʿrifahu ʿalā mā yanbaghī lā baʾs
ʿalayka.

But, My brother, in this degree---that is, at the beginning of the
journey, as We mentioned in the city of search---you will see different
stations and varying signs. All are true in their proper places and
stations. In this station you should behold all things in their proper
places, without bringing anything down from its elevation and loftiness
or raising anything from its station and lowliness. For example, if you
were to make divinity dwell within humanity, this would be sheer
association of partners with God; and if you were to raise humanity into
the atmosphere of divinity, this would be pure unbelief. But if you
speak of divinity in divinity and humanity in humanity, this is a truth
beyond doubt. That is, if you were to behold transformation in the
worlds of divine unity, this would be a sin greater than any in
existence. But if you behold transformation in its proper station and
understand it as you should, no blame rests upon you.

\subsection{90}\label{section-2801}

\begin{otherlanguage}{arabic}

و انّی فوربّی کلّ ما القیناک من اسرار البیان و مقامات التّبیان فی العیان
کأنّی ما ذکرت حرفاً من بحر علم اللّه المکنونة و جوهر حکمة اللّه المخزونة و
سنذکر فی حینها اذا شآء اللّه و اراد و انّه هو ذاکر کلّ شیء فی مقامها و انّا
کلّ له ذاکرون

\end{otherlanguage}

Wa-innī fa-wa-rabbī kullu mā alqaynāka min asrār al-bayān wa-maqāmāt
al-tibyān fī al-ʿiyān ka-annī mā dhakartu ḥarfan min baḥr ʿilm Allāh
al-maknūna wa-jawhar ḥikmat Allāh al-makhzūna wa-sanadhkuru fī ḥīnihā
idhā shāʾa Allāh wa-arāda wa-innahu huwa dhākiru kull shayʾ fī maqāmihā
wa-innā kullun lahu dhākirūn.

By My Lord, despite all We have imparted to you of the mysteries of
expression and the stations of explanation made manifest, it is as
though I had not mentioned a single letter from the ocean of God's
hidden knowledge and the essence of His treasured wisdom. We shall speak
of these in their time, if God so wills and desires. He indeed makes
mention of everything in its proper station, and We all make mention of
Him.

\subsection{91}\label{section-2802}

\begin{otherlanguage}{arabic}

ثمّ اعلم بأنّ طیر الّتی تطیر فی هوآء الجبروت لن تقدر ان تطیر فی سمآء قدس
اللّاهوت و لن تقدر ان تمذق فواکه الّتی خلق اللّه فیها و لن تقدر ان تشرب
انهار الّتی جرت فیها و لو تشرب قطرة منها لتموت فی الحین کما تشهد فی تلک
الأیّام عن الّذین ینسبون انفسهم الینا و یفعلون ما یفعلون و یقولون ما
یقولون و یدّعون ما یدّعون و کأنّهم فی حجباتهم میّتون

\end{otherlanguage}

Thumma iʿlam bi-anna ṭayr allatī taṭīru fī hawāʾ al-jabarūt lan taqdira
an taṭīra fī samāʾ quds al-lāhūt wa-lan taqdira an tamdhuqa fawākih
allatī khalaqa Allāh fīhā wa-lan taqdira an tashraba anhār allatī jarat
fīhā wa-law tashrabu qaṭratan minhā la-tamūtu fī al-ḥīn kamā tashhadu fī
tilka al-ayyām ʿan alladhīna yansibūna anfusahum ilaynā wa-yafʿalūna mā
yafʿalūna wa-yaqūlūna mā yaqūlūna wa-yaddaʿūna mā yaddaʿūna
wa-ka-annahum fī ḥujubātihim mayyitūn.

Know, then, that the bird that flies in the atmosphere of dominion
cannot fly in the heaven of the holiness of divinity, nor can it mingle
the fruits God has created there or drink from the rivers that flow
there. Were it to drink a drop from them, it would die at once, as you
witness in these days among those who ascribe themselves to Us, yet do
what they do, say what they say, and claim what they claim. It is as
though they were dead within their veils.

\subsection{92}\label{section-2803}

\begin{otherlanguage}{arabic}

کذلک فاعرف کلّ المقامات و الاشارات و الدّلالات و تعرف کلّ شیء فی مکانه و
تجد کلّ امر فی مقامه و لهذا المقام ای مقام مدینة الأحدیّة رجال قد رکبوا
علی فلک الهدایة و سافروا فی معارج الأحدیّة {[}و تشهد انوار الجمال عن
وجوههم و اسرار الجلال من هیاکلهم{]} و تجد روائح المسک من کلماتهم و تلاحظ
آیات السّلطنة فی مشیهم و حرکاتهم و سکونهم و لا یحجبک اعمال الّذینهم ما
شربوا من عیون الصّافیة و ما وصلوا الی مدائن القدسیّة و یتّبعون اهوآء انفسهم
و یفسدون فی الأرض و یحسبون بأنّهم مهتدون هم الّذین ورد فی شأنهم همج رعاع
اتباع کلّ ناعق یمیلون بکلّ ریح و مراتب هذا السّفر و هذا المقام و هذا الوطن
معلوم عند جنابک و مشهود عند حضرتک لا یحتاج الی تطویل الکلام

\end{otherlanguage}

Kadhālika fa-iʿraf kulla al-maqāmāt wa-al-ishārāt wa-al-dalālāt
wa-taʿrifu kulla shayʾ fī makānihi wa-tajidu kulla amr fī maqāmihi
wa-li-hādhā al-maqām ay maqām madīnat al-aḥadiyya rijāl qad rakibū ʿalā
fulk al-hidāya wa-sāfarū fī maʿārij al-aḥadiyya {[}wa-tashhadu anwār
al-jamāl ʿan wujūhihim wa-asrār al-jalāl min hayākilihim{]} wa-tajidu
rawāʾiḥ al-misk min kalimātihim wa-tulāḥiẓu āyāt al-salṭana fī mashyihim
wa-ḥarakātihim wa-sukūnihim wa-lā yaḥjubka aʿmāl alladhīnahum mā sharibū
min ʿuyūn al-ṣāfiya wa-mā waṣalū ilā madāʾin al-qudsiyya wa-yattabiʿūna
ahwāʾ anfusihim wa-yufsidūna fī al-arḍ wa-yaḥsabūna bi-annahum muhtadūn
hum alladhīna warada fī shaʾnihim hamajun raʿāʿ atbāʿu kulli nāʿiq
yamīlūna bi-kulli rīḥ wa-marātib hādhā al-safar wa-hādhā al-maqām
wa-hādhā al-waṭan maʿlūm ʿinda janābika wa-mashhūd ʿinda ḥaḍratika lā
yaḥtāju ilā taṭwīl al-kalām.

Understand all the stations, allusions, and indications in this way, and
you will recognize everything in its proper place and find every matter
in its proper station. To this station---that is, the station of the
city of oneness---belong men who have embarked upon the ark of guidance
and journeyed along the ascending paths of oneness. {[}You behold the
lights of beauty shining from their faces and the mysteries of majesty
from their persons.{]} You perceive the fragrance of musk in their words
and observe the signs of sovereignty in their gait, their movements, and
their stillness. Do not let yourself be veiled by the deeds of those who
have not drunk from the pure springs or reached the cities of holiness,
who follow their own desires, spread corruption on earth, and imagine
themselves rightly guided. They are those of whom it has been said: ``A
mindless rabble, followers of every clamorous voice, swaying with every
wind.'' The degrees of this journey, this station, and this homeland are
known to you and evident to your eyes; there is no need to dwell on them
at length.

\subsection{93}\label{section-2804}

\begin{otherlanguage}{arabic}

ثمّ اعلم بأنّ کلّ ما شهدت و سمعت بأنّ شمس الحقیقة و النّقطة الأوّلیّة نسب الی
نفسه من اسمآء القبل لم یکن ذلک الّا من ضعف العباد و هندسة عوالم الایجاد
والّا کلّ الأسمآء و الصّفات یطوفنّ حول ذاته و یدورنّ فی فنآء حرمه بل هو مربّی
الأسمآء و مظهر الصّفات و مذوّت الذّوات و معلن الآیات و مطرّز العلامات بل انّ
جنابک لو تشهد بعین سرّک لتجد ما دونه مفقود عنده و معدوم فی ساحته کان اللّه
و لم یکن معه من شیء و الآن کان بمثل ما قد کان و لمّا ثبت بأنّه جلّ و عزّ کان
و لم یکن معه من شیء کیف یجری حکم التّبدیل و التّغییر و انّک اذاً تفکّر فیما
القیناک لتظهر لک شمس الهدایة فی هذا الصّبح الأزلیّة و تکون فیه من الزّاهدین

\end{otherlanguage}

Thumma iʿlam bi-anna kulla mā shahidta wa-samiʿta bi-anna shams
al-ḥaqīqa wa-al-nuqṭa al-awwaliyya nasaba ilā nafsihi min asmāʾ al-qabl
lam yakun dhālika illā min ḍaʿf al-ʿibād wa-handasat ʿawālim al-ījād
wa-illā kullu al-asmāʾ wa-al-ṣifāt yaṭufna ḥawla dhātihi wa-yadurna fī
fināʾ ḥaramihi bal huwa murabbī al-asmāʾ wa-maẓhar al-ṣifāt wa-mudhawwit
al-dhawāt wa-muʿlin al-āyāt wa-muṭarriz al-ʿalāmāt bal inna janābaka law
tashhadu bi-ʿayn sirrika la-tajidu mā dūnahu mafqūdan ʿindahu
wa-maʿdūman fī sāḥatihi kāna Allāh wa-lam yakun maʿahu min shayʾ
wa-al-āna kāna bi-mithli mā qad kāna wa-lammā thabata bi-annahu jalla
wa-ʿazza kāna wa-lam yakun maʿahu min shayʾ kayfa yajrī ḥukm al-tabdīl
wa-al-taghyīr wa-innaka idhan tafakkar fīmā alqaynāka li-taẓhara laka
shams al-hidāya fī hādhā al-ṣubḥ al-azaliyya wa-takūna fīhi min
al-zāhidīn.

Know, then, that whenever you have seen or heard that the Sun of Truth
and the Primal Point ascribed to Himself any of the names of the past,
this was solely because of the weakness of His servants and the ordering
of the worlds of creation. Otherwise, all names and attributes circle
around His Essence and revolve within the precincts of His sanctuary.
Indeed, He is the nurturer of names, the manifestation of attributes,
the giver of essence to all essences, the revealer of signs, and the
embellisher of tokens. Were you to behold with your inner eye, you would
find all besides Him absent in His presence and nonexistent in His
court. ``God was, and nothing was with Him,'' and now He is as He has
ever been. Since it is established that He---glorified and exalted be
He---was, and nothing was with Him, how can the law of alteration and
change apply? Reflect, then, upon what We have imparted to you, so that
the sun of guidance may appear to you in this eternal dawn and you may
be among those who are detached within it.

\subsection{94}\label{section-2805}

\begin{otherlanguage}{arabic}

ثمّ اعلم بأن کلّ ما ذکرنا فی ذکر الأسفار لم یکن الّا للأحبار من الأخیار و
انّک لو ترکب علی براق المعنوی و تسیر فی حدائق الالهی لتقطع کلّ الأسفار و
تطّلع علی الأسرار من قبل ان ترتدّ الیک الأبصار

\end{otherlanguage}

Thumma iʿlam bi-anna kulla mā dhakarnā fī dhikr al-asfār lam yakun illā
li-al-aḥbār min al-akhyār wa-innaka law tarkabu ʿalā burāq al-maʿnawī
wa-tasīru fī ḥadāʾiq al-ilāhī la-taqṭaʿu kulla al-asfār wa-taṭṭaliʿu
ʿalā al-asrār min qabli an tartadda ilayka al-abṣār.

Know, then, that all We have said in describing these journeys was
intended only for the learned among the righteous. Were you to mount the
steed of the spirit and travel through the divine gardens, you would
traverse all these journeys and discover the mysteries before your eyes
could turn back to you.

\subsection{95}\label{section-2806}

\begin{otherlanguage}{arabic}

اذاً یا اخی ان تکون من فارس هذا المیدان فارکض فی ممالک الایقان لتخلّص نفسک
عن سجن الشّرک فی هذا الزّمان و تجد رائحة المسکیّة من نفحات هذه الحدیقة و من
عطر هذه المدینة تفرّقت نسمات العطریّة فی اقطار العالم و انّک لا تحرم نصیبک
و لا تکن من الغافلین

\end{otherlanguage}

Idhan yā akhī in takūna min fāris hādhā al-maydān fa-urkuḍ fī mamālik
al-īqān li-tukhalliṣa nafsaka ʿan sijn al-shirk fī hādhā al-zamān
wa-tajida rāʾiḥat al-miskiyya min nafaḥāt hādhihi al-ḥadīqa wa-min ʿiṭr
hādhihi al-madīna tafarraqat nasamāt al-ʿiṭriyya fī aqṭār al-ʿālam
wa-innaka lā taḥrim naṣībaka wa-lā takun min al-ghāfilīn.

Then, My brother, if you are a horseman of this arena, gallop through
the realms of certitude, so that you may free yourself from the prison
of associating partners with God in this day and catch the scent of musk
from the breezes of this garden. From the perfume of this city, fragrant
breezes have spread throughout the regions of the world. Do not deprive
yourself of your share, and do not be among the heedless.

\subsection{96}\label{section-2807}

\begin{otherlanguage}{arabic}

فنعم ما قال فلو عبقت فی الشّرق انفاس طیبها و فی الغرب مزکوم لعاد له الشّمّ

\end{otherlanguage}

Fa-niʿma mā qāla fa-law ʿabaqat fī al-sharq anfāsu ṭībihā wa-fī al-gharb
mazkūmun la-ʿāda lahu al-shamm.

How well he said: ``If the breaths of her fragrance spread in the East,
someone in the West whose nose was blocked would regain his sense of
smell.''

\subsection{97}\label{section-2808}

\begin{otherlanguage}{arabic}

و بعد هذا السّفر الالهی و هذا العروج المعنوی یدخل السّالک فی حدیقة الحیرة
(حیرة تعارج روح قدسیّة فی وادی نور حیرتیّة) و هذا مقام الّذی لو القی علیک
لتبکی و تنوح علی هذا العبد الّذی بقی بین یدی هؤلآء المشرکین و صار متحیّراً
فی امره و یکون فی هذه اللّجّة لمن المتحیّرین بحیث فی کلّ یوم یشاورون فی قتلی
و فی کلّ ساعة یریدون خروجی عن هذه البلد کما اخرجونی عن البلاد و هذا العبد
اکون حاضراً بین یدیهم و انتظر ما قضی اللّه علینا و حکم بنا و قدّر لأنفسنا و
ما اخاف من احد و ما احذر من نفس مع ما احاطتنا البأسآء و الضّرّآء من اهل
البغی و البغضآء و اغشت الأحزان فی تلک الأزمان فطوفان نوح عند نوحی کأدمعی
و ایقاد نیران الخلیل کلوعتی و حزنی ما یعقوب بثّ اقلّه و کلّ بلا ایّوب بعض
بلیّتی

\end{otherlanguage}

Wa-baʿda hādhā al-safar al-ilāhī wa-hādhā al-ʿurūj al-maʿnawī yadkhulu
al-sālik fī ḥadīqat al-ḥayra (ḥayrat taʿāruj rūḥ qudsiyya fī wādī nūr
ḥayriyya) wa-hādhā maqām alladhī law ulqiya ʿalayka la-tabkī wa-tanūḥu
ʿalā hādhā al-ʿabd alladhī baqiya bayna yaday hāʾulāʾ al-mushrikīn
wa-ṣāra mutaḥayyiran fī amrihi wa-yakūnu fī hādhihi al-lujja la-min
al-mutaḥayyirīn bi-ḥaythu fī kulli yawm yushāwirūna fī qatlī wa-fī kulli
sāʿa yurīdūna khurūjī ʿan hādhihi al-balad kamā akhrajūnī ʿan al-bilād
wa-hādhā al-ʿabd akūnu ḥāḍiran bayna yadayhim wa-antaẓiru mā qaḍā Allāh
ʿalaynā wa-ḥakama binā wa-qaddara li-anfusinā wa-mā akhāfu min aḥad
wa-mā aḥdharu min nafs maʿa mā aḥāṭatnā al-baʾsāʾ wa-al-ḍarrāʾ min ahl
al-baghy wa-al-baghḍāʾ wa-aghshat al-aḥzān fī tilka al-azmān fa-ṭūfānu
Nūḥ ʿinda nawḥī ka-admuʿī wa-īqādu nīrān al-Khalīl ka-lawʿatī wa-ḥuznī
mā Yaʿqūbu baththa aqallahu wa-kullu balā Ayyūb baʿḍu baliyyatī.

After this divine journey and this spiritual ascent, the wayfarer enters
the garden of bewilderment (the bewilderment of a holy spirit's ascent
in the valley of the light of bewilderment). This is a station which, if
it were disclosed to you, would make you weep and lament for this
Servant, who remains in the hands of these associators of partners with
God, bewildered in His plight and among those lost in wonder in this
ocean. Every day they consult about killing Me, and every hour they seek
My expulsion from this city, just as they expelled Me from other lands.
I, this Servant, stand before them, awaiting what God has decreed for
Us, ruled concerning Us, and ordained for Our souls. I fear no one and
am wary of no soul, although hardship and affliction have encompassed Us
at the hands of the people of oppression and hatred, and sorrows have
enveloped Us in these times. ``Beside My lamentation, Noah's flood is
but as My tears, and the kindling of the fires of Abraham but as My
burning anguish. Jacob did not pour forth even the least of My sorrow,
and all the afflictions of Job are but a portion of My tribulation.''

\subsection{98}\label{section-2809}

\begin{otherlanguage}{arabic}

و لو اذکر لجنابک بلایا النّازلة و القضایا الواردة لتحزن علی شأن ینقطع عنک
کلّ الأذکار و تغفل عن وجودک و عن کلّ ما خلق اللّه فی الملک و انّا لمّا ما
اردنا لجنابک ذلک لذا غطیت اظهار القضآء فی کبد البهآء و احتجبته عمّا یتحرّک
فی ارض الانشآء لیکون مکنوناً فی سرادق الغیب الی ان یظهر اللّه سرّه اذ لا
یعزب عن علمه من شیء لا فی السّموات و لا فی الأرض و انّه کان بکلّ شیء رقیب

\end{otherlanguage}

Wa-law adhkuru li-janābika balāyā al-nāzila wa-al-qaḍāyā al-wārida
la-taḥzanu ʿalā shaʾn yanqaṭiʿu ʿanka kullu al-adhkār wa-taghfulu ʿan
wujūdika wa-ʿan kulli mā khalaqa Allāh fī al-mulk wa-innā lammā mā
aradnā li-janābika dhālika li-dhā ghaṭṭaytu iẓhār al-qaḍāʾ fī kabid
al-Bahāʾ wa-iḥtajabtuhu ʿammā yataḥarraku fī arḍ al-inshāʾ li-yakūna
maknūnan fī surādiq al-ghayb ilā an yuẓhira Allāh sirrahu idh lā yaʿzubu
ʿan ʿilmihi min shayʾ lā fī al-samāwāt wa-lā fī al-arḍ wa-innahu kāna
bi-kulli shayʾ raqīb.

Were I to recount to you the afflictions that have descended and the
events that have befallen Us, you would grieve so deeply that all
remembrance would cease for you, and you would become oblivious to your
own existence and to everything God has created in His dominion. Since
We did not desire this for you, I have concealed the disclosure of this
decree within the heart of Bahāʾ and veiled it from all that moves in
the land of creation, so that it may remain hidden within the pavilion
of the unseen until God reveals its mystery. For nothing in the heavens
or on earth escapes His knowledge, and He is watchful over all things.

\subsection{99}\label{section-2810}

\begin{otherlanguage}{arabic}

و انّا لمّا بعدنا عن ذکر المقصود ترکنا الاشارات و رجعنا الی ما کنّا فیه فی
ذکر هذه المدینة الّتی من دخل فیها نجا و من اعرض عنها هلک

\end{otherlanguage}

Wa-innā lammā baʿudnā ʿan dhikr al-maqṣūd taraknā al-ishārāt wa-rajaʿnā
ilā mā kunnā fīhi fī dhikr hādhihi al-madīna allatī man dakhala fīhā
najā wa-man aʿraḍa ʿanhā halak.

Since We have strayed from the intended subject, We leave these
allusions and return to Our account of this city: whoever enters it is
saved, and whoever turns away from it perishes.

\subsection{100}\label{section-2811}

\begin{otherlanguage}{arabic}

فاعرف یا ایّها المذکور فی هذه الألواح بأنّ من دخل فی هذا السّفر یکون متحیّراً
فی آثار قدرة اللّه و بدائع آیات صنع اللّه و یأخذه الحیرة من کلّ الجهات و من
جمیع الأطراف کما شهد بذلک جوهر البقآء فی ملإ الأعلی فی قوله ربّ زدنی فیک
تحیّراً فنعم ما قال

\end{otherlanguage}

Fa-iʿraf yā ayyuhā al-madhkūr fī hādhihi al-alwāḥ bi-anna man dakhala fī
hādhā al-safar yakūnu mutaḥayyiran fī āthār qudrat Allāh wa-badāʾiʿ āyāt
ṣunʿ Allāh wa-yaʾkhudhuhu al-ḥayra min kull al-jihāt wa-min jamīʿ
al-aṭrāf kamā shahida bi-dhālika jawhar al-baqāʾ fī malaʾ al-aʿlā fī
qawlihi rabb zidnī fīka taḥayyuran fa-niʿma mā qāl.

Know, you who are mentioned in these Tablets, that whoever embarks upon
this journey becomes bewildered by the evidences of God's power and the
wondrous signs of His handiwork. Bewilderment seizes him from every
direction and on every side, as the Essence of Eternity testified in the
Concourse on High with His words: ``My Lord, increase my bewilderment in
You.'' And how well it has been said:

\subsection{101}\label{section-2812}

\begin{otherlanguage}{arabic}

و ما احترت حتّی اخترت حبّک مذهباً فوا حیرتی لو لم تکن فیک حیرتی

\end{otherlanguage}

Wa-mā iḥtartu ḥattā ikhtartu ḥubbaka madhhaban fa-wā ḥayratī law lam
takun fīka ḥayratī.

``I was not bewildered until I chose love for You as my way. How
bewildered I would be if my bewilderment were not in You!''

\subsection{102}\label{section-2813}

\begin{otherlanguage}{arabic}

و فی ذلک الوادی تضلّون السّالکون و تهلکون و لن تقدروا ان تصلوا الی مثواهم
اللّه اکبر من عظمة هذا الواد و من وسعة هذه المدینة فی جبروت الایجاد کأنّک
لن تجد له من اوّل و لا من آخر فبشری ثمّ بشری لمن کمل فیها سفره و ایّده اللّه
علی طیّ هذه الأرض الطّیّبة فی هذه المدینة الالهیّة الّتی تتحیّر فیها کلّ
المقرّبین و المخلصین و نقول الحمد للّه ربّ العالمین

\end{otherlanguage}

Wa-fī dhālika al-wādī taḍillūna al-sālikūn wa-tahlikūna wa-lan taqdirū
an taṣilū ilā mathwāhum Allāhu akbar min ʿaẓamat hādhā al-wād wa-min
wusʿat hādhihi al-madīna fī jabarūt al-ījād ka-annaka lan tajida lahu
min awwal wa-lā min ākhir fa-bushrā thumma bushrā li-man kamula fīhā
safaruhu wa-ayyadahu Allāh ʿalā ṭayy hādhihi al-arḍ al-ṭayyiba fī
hādhihi al-madīna al-ilāhiyya allatī tataḥayyaru fīhā kullu
al-muqarrabīn wa-al-mukhliṣīn wa-naqūlu al-ḥamdu li-llāh rabb
al-ʿālamīn.

In that valley you wayfarers lose your way and perish, and you cannot
reach their abode. God is greater than the grandeur of this valley and
the vastness of this city in the dominion of creation! It is as though
you could find neither its beginning nor its end. Glad tidings, then,
glad tidings to the one whose journey within it is complete, whom God
has aided to traverse this goodly land in this divine city, where all
who are brought near and all who are sincere are bewildered. And We say:
``Praise be to God, the Lord of the worlds.''

\subsection{103}\label{section-2814}

\begin{otherlanguage}{arabic}

و لو یتعارج العبد و یسافر عن هذا الوطن التّرابی و یرید ان یتعارج الی وطن
الالهی لیدخل من هذه المدینة الی مدینة الفنآء (فنآء تغرّد محو {[}الهیّة{]}
فی حدیقة صعق فنائیّة) لفنائه عن نفسه و بقائه باللّه و السّالک فی هذا المقام
و هذا الوطن البحت الأعلی و هذا السّفر المحو الکبری لینسی نفسه و روحه و
جسده و ذاته و یسبح فی قلزم الفنآء و یکون فی الأرض کمن لم یکن شیئاً مذکورا
و لن یشهد احد منه آثار الوجود لاضمحلاله عن ممالک الشّهود و لبلوغه الی
مقامات المحو

\end{otherlanguage}

Wa-law yataʿāraju al-ʿabd wa-yusāfiru ʿan hādhā al-waṭan al-turābī
wa-yurīdu an yataʿāraja ilā waṭan al-ilāhī li-yadkhula min hādhihi
al-madīna ilā madīnat al-fanāʾ (fanāʾ tagharrud maḥw {[}ilāhiyya{]} fī
ḥadīqat ṣaʿq fanāʾiyya) li-fanāʾihi ʿan nafsihi wa-baqāʾihi bi-llāh
wa-al-sālik fī hādhā al-maqām wa-hādhā al-waṭan al-baḥt al-aʿlā wa-hādhā
al-safar al-maḥw al-kubrā la-yansā nafsahu wa-rūḥahu wa-jasadahu
wa-dhātahu wa-yasbaḥu fī qulzum al-fanāʾ wa-yakūnu fī al-arḍ ka-man lam
yakun shayʾan madhkūran wa-lan yashhada aḥadun minhu āthār al-wujūd
li-iḍmiḥlālihi ʿan mamālik al-shuhūd wa-li-bulūghihi ilā maqāmāt
al-maḥw.

Should the servant ascend and journey away from this earthly homeland,
desiring to ascend to the divine homeland, he will enter from this city
into the city of annihilation (the annihilation of the singing of
{[}divine{]} effacement in the garden of the swoon of annihilation),
through his annihilation of self and his abiding in God. In this
station, this pure and most exalted homeland, this journey of supreme
effacement, the wayfarer forgets his self, his spirit, his body, and his
being. He swims in the ocean of annihilation and is on earth as one who
was not a thing worthy of mention. No one will perceive in him any
traces of existence, for he has vanished from the realms of the visible
and attained the stations of effacement.

\subsection{104}\label{section-2815}

\begin{otherlanguage}{arabic}

و لو انّا نذکر اسرار هذه المدینة لتفنی ممالک الفؤاد لکثرة شوق اهلها الی
هذا المقام السّداد لأنّ هذا المقام مقام تجلّی المعشوق للعاشق الصّادق و ظهور
اشراق انوار المحبوب للحبیب الفارغ

\end{otherlanguage}

Wa-law innā nadhkuru asrār hādhihi al-madīna la-tafnā mamālik al-fuʾād
li-kathrat shawq ahlihā ilā hādhā al-maqām al-sadād li-anna hādhā
al-maqām maqām tajallī al-maʿshūq li-al-ʿāshiq al-ṣādiq wa-ẓuhūr ishrāq
anwār al-maḥbūb li-al-ḥabīb al-fārigh.

Were We to recount the mysteries of this city, the realms of the heart
would perish from their inhabitants' intense longing for this station of
rectitude. For this is the station where the Beloved reveals Himself to
the true lover, and where the radiance of the Loved One's lights appears
to the lover who is free of all attachment.

\subsection{105}\label{section-2816}

\begin{otherlanguage}{persian}

و هل یمکن للعاشق وجود حین تجلّی المعشوق او للظّلّ بقآء عند ظهور الشّمس او
للحبیب دوام عند وجود المحبوب لا فوالّذی نفسی بیده بل السّالک فی هذا المقام
لو تفحّص فی شرق الأرض و غربها و برّها و بحرها و سهلها و جبلها ما یجد نفسه
و لا نفس غیره لشدّة فنائه فی موجده و لطافة محوه فی بارئه

\end{otherlanguage}

Wa-hal yumkinu li-al-ʿāshiq wujūd ḥīna tajallī al-maʿshūq aw li-al-ẓill
baqāʾ ʿinda ẓuhūr al-shams aw li-al-ḥabīb dawām ʿinda wujūd al-maḥbūb lā
fa-wa-alladhī nafsī bi-yadihi bal al-sālik fī hādhā al-maqām law
tafaḥḥaṣa fī sharq al-arḍ wa-gharbihā wa-barrihā wa-baḥrihā wa-sahlihā
wa-jabalihā mā yajidu nafsahu wa-lā nafsa ghayrihi li-shiddat fanāʾihi
fī mūjidihi wa-laṭāfat maḥwihi fī bāriʾihi.

Can the lover have any existence when the Beloved reveals Himself, or
the shadow remain when the sun appears, or the lover endure in the
presence of the Loved One? No, by Him in whose hand is My soul! Indeed,
were the wayfarer in this station to search the east and west of the
earth, its land and sea, its plains and mountains, he would find neither
himself nor anyone else, so complete is his annihilation in the One who
brought him into being, and so subtle his effacement in his Creator.

\subsection{106}\label{section-2817}

\begin{otherlanguage}{arabic}

فسبحان اللّه لو لا خوفی من نمرود الظّلم و حفظی لخلیل العدل لألقی علیک ما
یغنیک عن دونک و لأقرأ لک ما یقرّبک الی هذه المدینة حین غفلة عن نفسک و
هواک ولکن اصبر حتّی یأتی اللّه بأمره و انّه هو یجزی الصّابرین بغیر حساب اذاً
فانشق رائحة الرّوحانی من قمص المعانی و قل یا اهل لجّة الفنآء ان اسرعوا
للدّخول فی مدینة البقآء ان انتم الی معارج البقآء تتعارجون و نقول انّا للّه
و انّا الیه راجعون

\end{otherlanguage}

Fa-subḥāna Allāh law lā khawfī min Nimrūd al-ẓulm wa-ḥifẓī li-Khalīl
al-ʿadl la-ulqī ʿalayka mā yughnīka ʿan dūnika wa-la-aqraʾu laka mā
yuqarribuka ilā hādhihi al-madīna ḥīna ghaflatin ʿan nafsika wa-hawāka
wa-lākin iṣbir ḥattā yaʾtiya Allāh bi-amrihi wa-innahu huwa yajzī
al-ṣābirīn bi-ghayri ḥisāb idhan fa-inshaq rāʾiḥat al-rūḥānī min qumuṣ
al-maʿānī wa-qul yā ahl lujjat al-fanāʾ an asriʿū li-al-dukhūl fī
madīnat al-baqāʾ in antum ilā maʿārij al-baqāʾ tataʿārajūn wa-naqūlu
innā li-llāh wa-innā ilayhi rājiʿūn.

Glory be to God! Were it not for My fear of the Nimrod of oppression and
My protection of the Abraham of justice, I would impart to you what
would make you independent of all besides yourself, and recite to you
what would bring you near this city when you are oblivious to your self
and your desires. But be patient until God brings forth His command. He
indeed rewards the patient without measure. Breathe, then, the spiritual
fragrance from the garments of inner meanings, and say: ``People of the
ocean of annihilation, hasten to enter the city of abiding, if you would
ascend the ascending paths of abiding.'' And We say: ``We belong to God,
and to Him we return.''

\subsection{107}\label{section-2818}

\begin{otherlanguage}{persian}

و من ذلک المقام الأعلی الأعلی و الرّتبة الأعظم الأسنی یدخل فی مدینة
البقآء (بقآء تغنّی ورقآء دائمیّة فی وادی قدس بقائیّة) علی البقآء و فی ذلک
المقام یشهد السّالک نفسه علی عرش الاستغنآء و کرسیّ الاستعلآء اذاً یظهر له
حکم ما ذکر من قبل یوم یغنی اللّه کلّاً من سعته فهنیئاً لمن وصل الی هذا
المقام و شرب من هذا الکأس البیضآء فی هذا الرّکن الحمرآء

\end{otherlanguage}

Wa-min dhālika al-maqām al-aʿlā al-aʿlā wa-al-rutba al-aʿẓam al-asnā
yadkhulu fī madīnat al-baqāʾ (baqāʾ taghannī warqāʾ dāʾimiyya fī wādī
quds baqāʾiyya) ʿalā al-baqāʾ wa-fī dhālika al-maqām yashhadu al-sālik
nafsahu ʿalā ʿarsh al-istighnāʾ wa-kursī al-istiʿlāʾ idhan yaẓharu lahu
ḥukm mā dhukira min qabl yawma yughnī Allāhu kullan min saʿatihi
fa-hanīʾan li-man waṣala ilā hādhā al-maqām wa-shariba min hādhā al-kaʾs
al-bayḍāʾ fī hādhā al-rukn al-ḥamrāʾ.

From that most exalted, most exalted station and that greatest, most
resplendent rank, he enters the city of abiding (the abiding of the
singing of an everlasting dove in the valley of the holiness of
abiding), to abide there. In that station the wayfarer beholds himself
upon the throne of independence and the seat of exaltation. Then the
import of what was said before becomes apparent to him: ``The day when
God enriches everyone out of His abundance.'' Blessed is the one who has
reached this station and drunk from this white cup in this crimson
sanctuary.

\subsection{108}\label{section-2819}

\begin{otherlanguage}{persian}

فانّ السّالک فی هذا السّفر لمّا استغرق فی ابحر البقآء و استفرغ فؤاده عن کلّ
ما سواه و استبلغ الی معارج الحیاة لا یری الفنآء لنفسه و لا لغیره ابداً و
یشرب عن کأس البقآء و یمشی فی ارض البقآء و یطیر فی هوآء البقآء و یجالس مع
هیاکل البقآء و یأکل من نعمة الباقیة الدّائمة من شجرة الدّائمة الأزلیّة و
یکون من اهل البقآء فی علی البقآء بالبقآء مذکورا

\end{otherlanguage}

Fa-inna al-sālik fī hādhā al-safar lammā istaghraqa fī abḥur al-baqāʾ
wa-istafragha fuʾādahu ʿan kull mā siwāhu wa-istablagha ilā maʿārij
al-ḥayāh lā yarā al-fanāʾ li-nafsihi wa-lā li-ghayrihi abadan wa-yashrab
ʿan kaʾs al-baqāʾ wa-yamshī fī arḍ al-baqāʾ wa-yaṭīr fī hawāʾ al-baqāʾ
wa-yujālis maʿa hayākil al-baqāʾ wa-yaʾkul min niʿmat al-bāqiyah
al-dāʾimah min shajarat al-dāʾimah al-azaliyyah wa-yakūn min ahl
al-baqāʾ fī ʿulā al-baqāʾ bi-l-baqāʾ madhkūran.

For when the wayfarer on this journey becomes immersed in the oceans of
abiding, empties his heart of all else, and reaches the ascending
heights of life, he never sees annihilation for himself or for anyone
else. He drinks from the cup of abiding, walks in the land of abiding,
flies in the air of abiding, sits with the embodiments of abiding, and
partakes of enduring, everlasting bounty from the everlasting, eternal
tree. He is numbered among the people of abiding, remembered through
abiding in the heights of abiding.

\subsection{109}\label{section-2820}

\begin{otherlanguage}{arabic}

و کلّ ما یکون فی هذه المدینة لباقیة دائمة لا یفنی و انت لو تدخل باذن اللّه
فی هذه الحدیقة العالیة المتعالیة لتجد شمسها فی قطب الزّوال بحیث لا تکسف و
لا تغرب ابداً و کذلک قمرها و افلاکها و انجمها و اشجرها و ابحرها و کلّ ما
فیها و بها و انّی فواللّه الّذی لا اله الّا هو لو اذکر لک بدائع اوصاف هذه
المدینة من یومئذ الی آخر الّذی لا آخر له ما یفرغ حبّ فؤادی لهذه المدینة
الطّیّبة الدّائمة ولکن اختم القول لضیق الوقت و تعجیل الطّالب و لئلّا تظهر
الأسرار فی الاجهار من دون اذن من اللّه المقتدر القهّار

\end{otherlanguage}

Wa-kull mā yakūn fī hādhihi al-madīnah la-bāqiyah dāʾimah lā yafnā
wa-anta law tadkhul bi-idhn Allāh fī hādhihi al-ḥadīqah al-ʿāliyah
al-mutaʿāliyah la-tajid shamsahā fī quṭb al-zawāl bi-ḥaythu lā tuksaf
wa-lā taghrub abadan wa-kadhālika qamaruhā wa-aflākuhā wa-anjumuhā
wa-ashjaruhā wa-abḥuruhā wa-kull mā fīhā wa-bihā wa-innī fa-wa-Allāh
alladhī lā ilāha illā huwa law adhkur laka badāʾiʿ awṣāf hādhihi
al-madīnah min yawmiʾidh ilā ākhir alladhī lā ākhir lahu mā yafragh ḥubb
fuʾādī li-hādhihi al-madīnah al-ṭayyibah al-dāʾimah wa-lākin akhtim
al-qawl li-ḍīq al-waqt wa-taʿjīl al-ṭālib wa-li-allā taẓhar al-asrār fī
al-ijhār min dūn idhn min Allāh al-Muqtadir al-Qahhār.

Everything in this city is enduring and everlasting, never passing away.
Were you, by God's permission, to enter this lofty and exalted garden,
you would find its sun at the zenith, never eclipsed and never setting.
So too are its moon, its heavens, its stars, its trees, its seas, and
everything within it and belonging to it. By God, besides whom there is
no other god, were I to recount to you the wondrous attributes of this
city from this day until the end that has no end, My heart's love for
this good and everlasting city would not be exhausted. But I bring these
words to a close because time is short and the seeker is pressing, and
so that the mysteries may not be openly disclosed without permission
from God, the All-Powerful, the All-Subduing.

\subsection{110}\label{section-2821}

\begin{otherlanguage}{arabic}

و سینظر الموحّدون فی قیامة الأخری بأنّ من یظهره اللّه مع هذه المدینة ینزل
من سمآء الغیب مع ملائکة المقرّبین العالین فطوبی لمن یحضر بین یدیه و یفوز
بلقائه و انّا کلّ بذلک آملون و نقول الحمد لله اذ هو الحقّ و انّا کلّ الیه
منقلبون

\end{otherlanguage}

Wa-sa-yanẓur al-muwaḥḥidūn fī qiyāmat al-ukhrā bi-anna man yuẓhiruhu
Allāh maʿa hādhihi al-madīnah yanzil min samāʾ al-ghayb maʿa malāʾikat
al-muqarrabīn al-ʿālīn fa-ṭūbā li-man yaḥḍur bayna yadayhi wa-yafūz
bi-liqāʾihi wa-innā kull bi-dhālika āmilūn wa-naqūl al-ḥamdu li-llāh idh
huwa al-Ḥaqq wa-innā kull ilayhi munqalibūn.

In the resurrection to come, those who affirm God's unity will behold
Him whom God will make manifest descending with this city from the
heaven of the unseen, accompanied by the angels who are exalted and
brought near. Blessed is the one who enters His presence and attains
reunion with Him. We all hope for this, and We say: ``Praise be to God,
for He is the Truth, and to Him we shall all return.''

\subsection{111}\label{section-2822}

\begin{otherlanguage}{persian}

ثمّ اعرف بأنّ الواصل الی هذه المقامات و المسافر فی هذه الأسفار لو یناله فی
السّبیل من کبر او غرور لیهلک فی الحین و یرجع الی قدم الأوّل من دون ان یعرف
ذلک و علامة الواصلین و المشتاقین فی هذه الأسفار ان یخفضوا جناحهم للّذین
آمنوا باللّه و آیاته و یبخعوا انفسهم للّذین استقربوا الی اللّه و مظاهر
جماله و یخضعوا ذواتهم للّذین استقرّوا علی رفرف امر اللّه و عظمته

\end{otherlanguage}

Thumma iʿraf bi-anna al-wāṣil ilā hādhihi al-maqāmāt wa-al-musāfir fī
hādhihi al-asfār law yanāluhu fī al-sabīl min kibr aw ghurūr la-yahlak
fī al-ḥīn wa-yarjiʿ ilā qadam al-awwal min dūn an yaʿrif dhālika
wa-ʿalāmat al-wāṣilīn wa-al-mushtāqīn fī hādhihi al-asfār an yakhfiḍū
janāḥahum li-lladhīna āmanū bi-Allāh wa-āyātihi wa-yabkhaʿū anfusahum
li-lladhīna istaqrabū ilā Allāh wa-maẓāhir jamālihi wa-yakhḍaʿū
dhawātahum li-lladhīna istaqarrū ʿalā rafraf amr Allāh wa-ʿaẓamatihi.

Know, then, that if anyone who reaches these stations and travels on
these journeys is touched by pride or conceit along the way, he will
perish at once and return to the first step without realizing it. The
sign of those who have arrived and those who yearn on these journeys is
that they lower their wings before those who believe in God and His
verses, humble themselves before those who have drawn near to God and
the Manifestations of His beauty, and bow before those established upon
the lofty seat of God's Cause and grandeur.

\subsection{112}\label{section-2823}

\begin{otherlanguage}{arabic}

لأنّهم لو یتعارجون الی غایة القصوی فی سلوکهم الی اللّه و وصولهم الیه لن
یصلوا الّا الی مقرّ الّذی خلقت فی افئدتهنّ فکیف یقدرنّ ان یتعارجنّ الی مقامات
الّتی ما قدّرت لهم و ما خلقت لشأنهم و لو یسافرون من الأزل الی الأبد لن
یصلوا الی قطب الوجود و مرکز الموجود الّذی جری عن یمینه بحور العظمة و عن
یساره شطوط القدرة و لن یقدر احد ان ینزل بفنائه و کیف الی مقامه و هو کان
ساکناً فی فلک النّار و یسری علی بحر النّار فی کرة النّار و یمشی فی هوآء
النّار فکیف یقدر من خلق بالأضداد ان یدخل فی النّار او یقرب بها و ان یقربها
لیحترق فی الحین

\end{otherlanguage}

Li-annahum law yataʿārajūn ilā ghāyat al-quṣwā fī sulūkihim ilā Allāh
wa-wuṣūlihim ilayhi lan yaṣilū illā ilā maqarr alladhī khuliqat fī
afʾidatihinna fa-kayfa yaqdiranna an yataʿārajna ilā maqāmāt allatī mā
quddirat lahum wa-mā khuliqat li-shaʾnihim wa-law yusāfirūn min al-azal
ilā al-abad lan yaṣilū ilā quṭb al-wujūd wa-markaz al-mawjūd alladhī
jarā ʿan yamīnihi buḥūr al-ʿaẓamah wa-ʿan yasārihi shuṭūṭ al-qudrah
wa-lan yaqdir aḥad an yanzil bi-fināʾihi wa-kayfa ilā maqāmihi wa-huwa
kāna sākinan fī fulk al-nār wa-yasrī ʿalā baḥr al-nār fī kurat al-nār
wa-yamshī fī hawāʾ al-nār fa-kayfa yaqdir man khuliqa bi-l-aḍdād an
yadkhul fī al-nār aw yaqrub bihā wa-in yaqrabhā la-yaḥtariq fī al-ḥīn.

For even if they ascend to the furthest limit in their journey toward
God and their attainment of Him, they will reach only the dwelling
created within their own hearts. How, then, could they ascend to
stations that have not been ordained for them or created for their
condition? Even if they journey from eternity to eternity, they will
never reach the Pole of existence and the Center of being, to whose
right flow the seas of grandeur and to whose left flow the rivers of
power. No one can alight in His courtyard, much less attain His station.
He dwells in the ark of fire, sails upon the sea of fire within the
sphere of fire, and walks in the air of fire. How, then, can one created
of opposites enter the fire or draw near it? If he approaches it, he
will burn at once.

\subsection{113}\label{section-2824}

\begin{otherlanguage}{arabic}

ثمّ اعلم بأنّ هذا القطب الأعظم لو یقطع خیط مدده عن کلّ من فی السّموات و
الأرض لتنعدم کلّهنّ فسبحان اللّه کیف یصل التّراب الی ربّ الأرباب فسبحان اللّه
عمّا یظنّون فی انفسهم و تعالی عمّا هم یذکرون

\end{otherlanguage}

Thumma iʿlam bi-anna hādhā al-quṭb al-aʿẓam law yaqṭaʿ khayṭ madadihi
ʿan kull man fī al-samāwāt wa-al-arḍ la-tanʿadim kulluhunna fa-subḥān
Allāh kayfa yaṣil al-turāb ilā Rabb al-arbāb fa-subḥān Allāh ʿammā
yaẓunnūn fī anfusihim wa-taʿālā ʿammā hum yadhkurūn.

Know, then, that if this most great Pole were to sever the thread of His
sustaining grace from all who are in the heavens and on earth, they
would all cease to exist. Glory be to God! How can dust reach the Lord
of lords? God is far above what they imagine within themselves, and
exalted beyond what they say.

\subsection{114}\label{section-2825}

\begin{otherlanguage}{arabic}

بلی انّ السّالک یتعارج الی مقام الّذی لا غایة له فیما قدّر له و یجد فی قلبه
نار الحبّ بحیث یأخذ زمام الاختیار عن هؤلآء الأخیار و فی کلّ حین یزداد فی
حبّه مولاه و اقباله الی بارئه بحیث لو کان مولاه فی مشرق القربیّة و هو فی
مغرب البعدیّة و کان له ملء السّموات و الأرض من اللّؤلؤ الحمرآء و الذّهب
الصّفرآء لینفق و یرکض بعینیه لیصل الی ارض الّتی کان المقصود فیها و لو تجد
السّالک بغیر ذلک فاعلم بأنّه کذّاب مفتر انّا لمن یظهره اللّه فی قیامة الأخری
و انّا به لمبعثون

\end{otherlanguage}

Balā inna al-sālik yataʿāraj ilā maqām alladhī lā ghāyah lahu fīmā
quddira lahu wa-yajid fī qalbihi nār al-ḥubb bi-ḥaythu yaʾkhudh zimām
al-ikhtiyār ʿan hāʾulāʾ al-akhyār wa-fī kull ḥīn yazdād fī ḥubbihi
mawlāhu wa-iqbālihi ilā bāriʾihi bi-ḥaythu law kāna mawlāhu fī mashriq
al-qurbiyyah wa-huwa fī maghrib al-buʿdiyyah wa-kāna lahu milʾ
al-samāwāt wa-al-arḍ min al-luʾluʾ al-ḥamrāʾ wa-al-dhahab al-ṣafrāʾ
la-yunfiq wa-yarkuḍ bi-ʿaynayhi li-yaṣil ilā arḍ allatī kāna al-Maqṣūd
fīhā wa-law tajid al-sālik bi-ghayr dhālika fa-iʿlam bi-annahu kadhdhāb
muftarin innā li-man yuẓhiruhu Allāh fī qiyāmat al-ukhrā wa-innā bihi
la-mabʿūthūn.

Indeed, the wayfarer ascends to a station without limit within what has
been ordained for him, and finds in his heart the fire of love, so
intense that it wrests the reins of choice from these righteous souls.
At every moment he grows in love for his Lord and in his turning toward
his Creator, to such a degree that, were his Lord in the east of
nearness and he in the west of remoteness, and were he to possess red
pearls and yellow gold enough to fill the heavens and the earth, he
would spend it all and hasten on his very eyes to reach the land where
the Desired One dwells. If you find the wayfarer otherwise, know that he
is a lying impostor. We belong to Him whom God will make manifest in the
resurrection to come, and through Him we shall be raised.

\subsection{115}\label{section-2826}

\begin{otherlanguage}{arabic}

و فی تلک الأیّام لمّا ما کشفنا الغطآء عن وجه الأمر و ما {[}ظهرنا{]} للعباد
ثمرات هذه المقامات الّتی منعنا عن اظهارها لذا تجدهم فی سکران الغفلة والّا
لو کشف لکلّ من علی الأرض اقلّ من سمّ الابرة من هذا المقام لتشهد کیف یجتمعون
فی فنآء رحمة اللّه و یرکضون من کلّ الأطراف للبلوغ الی ساحة القرب فی رفرف
عزّة اللّه ولکن اخفینا لما ذکرنا من قبل و لیمتاز المؤمنون عن المنکرین و
المقبلون عن المعرضین و اقول لا حول و لا قوّة الّا باللّه المهیمن القیّوم

\end{otherlanguage}

Wa-fī tilka al-ayyām lammā mā kashafnā al-ghiṭāʾ ʿan wajh al-amr wa-mā
{[}ẓaharnā{]} li-l-ʿibād thamarāt hādhihi al-maqāmāt allatī muniʿnā ʿan
iẓhārihā li-dhā tajiduhum fī sukrān al-ghaflah wa-illā law kushifa
li-kull man ʿalā al-arḍ aqall min samm al-ibrah min hādhā al-maqām
la-tashhad kayfa yajtamiʿūn fī fināʾ raḥmat Allāh wa-yarkuḍūn min kull
al-aṭrāf li-l-bulūgh ilā sāḥat al-qurb fī rafraf ʿizzat Allāh wa-lākin
akhfaynā li-mā dhakarnā min qabl wa-li-yamtāz al-muʾminūn ʿan
al-munkirīn wa-al-muqbilūn ʿan al-muʿriḍīn wa-aqūl lā ḥawl wa-lā quwwah
illā bi-Allāh al-Muhaymin al-Qayyūm.

In these days, since We have not lifted the veil from the face of the
Cause or {[}shown{]} the servants the fruits of these stations, which We
have been prevented from revealing, you find them drunk with
heedlessness. Otherwise, were less than a needle's eye of this station
disclosed to everyone on earth, you would see how they would gather in
the courtyard of God's mercy and hasten from every direction to reach
the precinct of nearness upon the lofty seat of God's might. But We have
concealed it for the reasons We mentioned before, and so that the
believers may be distinguished from the deniers, and those who turn
toward Him from those who turn away. And I say: ``There is no strength
or power except through God, the All-Protecting, the Self-Subsisting.''

\subsection{116}\label{section-2827}

\begin{otherlanguage}{arabic}

و یسترقی السّالک من هذا المقام الی مدینة (غیب تصاعد قدس ازلیّة فی طلعة غیب
سرّیّة) الّتی لم یکن لها من اسم و لا رسم و لا ذکر و لا صوت تجری فیها بحور
القدم و تدور فی حول القدم و تشرق فیها شمس الغیب عن افق الغیب و لها افلاک
من نفسها و اقمار من نورها کلّهنّ یطلعن من بحر الغیب و یدخلن فی بحر الغیب و
انّی ما اقدر ان اذکر رشحاً عمّا قدّر فیها و لا یطّلع علی اسرارها احد الّا اللّه
و مظاهر نفسه اذ هو خالقها و مبدعها

\end{otherlanguage}

Wa-yastarqī al-sālik min hādhā al-maqām ilā madīnat (ghayb taṣāʿud quds
azaliyyah fī ṭalʿat ghayb sirriyyah) allatī lam yakun lahā min ism wa-lā
rasm wa-lā dhikr wa-lā ṣawt tajrī fīhā buḥūr al-qidam wa-tadūr fī ḥawl
al-qidam wa-tushriq fīhā shams al-ghayb ʿan ufuq al-ghayb wa-lahā aflāk
min nafsihā wa-aqmār min nūrihā kulluhunna yaṭluʿna min baḥr al-ghayb
wa-yadkhulna fī baḥr al-ghayb wa-innī mā aqdir an adhkur rashḥan ʿammā
quddira fīhā wa-lā yaṭṭaliʿ ʿalā asrārihā aḥad illā Allāh wa-maẓāhir
nafsihi idh huwa khāliquhā wa-mubdiʿuhā.

From this station the wayfarer ascends to the city (the unseen ascent of
eternal holiness in the countenance of a hidden mystery), which has no
name, no trace, no mention, and no sound. Within it flow the seas of
eternity, and they circle around eternity. Within it the sun of the
unseen shines from the horizon of the unseen. It has heavens of its own
and moons of its own light, all rising from the sea of the unseen and
entering the sea of the unseen. I cannot mention even a sprinkling of
what has been ordained within it. No one knows its mysteries except God
and the Manifestations of His own Self, for He is its Creator and
Originator.

\subsection{117}\label{section-2828}

\begin{otherlanguage}{arabic}

ثمّ اعلم بأنّا حین الّذی اردنا ان نتعرّض بتلک الکلمات و کتبنا بعضها اردنا
بأن نفسّر لجنابک کلّ ما ذکرنا من قبل من کلمات النّبیّین و عبارات المرسلین
بنغمات المقرّبین و ربوات المقدّسین ولکن ما وجدنا الفرصة و ما شهدنا المهلة
من هذا المسافر الّذی جآء من عندکم و کان عجولاً فی الأمر و راکضاً فی الحکم
لذا قد اختصرنا و اکتفینا و ما اتممنا ذکر الأسفار بتمامها و ما ینبغی لها
و یلیق بها بل ترکنا ذکر مدائن الکبری و اسفار العظمی و بلغ تعجیل الرّافع
الی مقام الّذی ترکنا ذکر السّفرین الأعلیین فی التّسلیم و الرّضا

\end{otherlanguage}

Thumma iʿlam bi-annā ḥīn alladhī aradnā an nataʿarraḍ bi-tilka
al-kalimāt wa-katabnā baʿḍahā aradnā bi-an nufassir li-janābika kull mā
dhakarnā min qabl min kalimāt al-nabiyyīn wa-ʿibārāt al-mursalīn
bi-naghamāt al-muqarrabīn wa-rabawāt al-muqaddasīn wa-lākin mā wajadnā
al-furṣah wa-mā shahidnā al-muhlah min hādhā al-musāfir alladhī jāʾa min
ʿindikum wa-kāna ʿajūlan fī al-amr wa-rākiḍan fī al-ḥukm li-dhā qad
ikhtaṣarnā wa-iktafaynā wa-mā atmamnā dhikr al-asfār bi-tamāmihā wa-mā
yanbaghī lahā wa-yalīq bihā bal taraknā dhikr madāʾin al-kubrā wa-asfār
al-ʿuẓmā wa-balagha taʿjīl al-rāfiʿ ilā maqām alladhī taraknā dhikr
al-safrayn al-aʿlayayn fī al-taslīm wa-al-riḍā.

Know, then, that when We wished to address these words and had written
some of them, We intended to explain to you everything We had previously
mentioned of the words of the Prophets and the utterances of the
Messengers, in the melodies of those brought near and upon the heights
of the sanctified. But We found no opportunity and were granted no
respite by this traveler who came from you, for he was hasty in the
matter and hurried in his decision. We have therefore been brief and
contented Ourselves with this, and have not completed the account of all
the journeys as would be fitting and appropriate. Rather, We have left
unmentioned the greater cities and the mightier journeys. The bearer's
haste reached such a point that We left unmentioned the two most exalted
journeys, those of surrender and contentment.

\subsection{118}\label{section-2829}

\begin{otherlanguage}{arabic}

ولو انّ جنابک لو تفکّر فی هذه الکلمات المختصرات لتعرف کلّ العلوم و تصل الی
ذروة المعلوم و تقول یکفی کلّ الوجود من المشهود و المفقود

\end{otherlanguage}

Wa-law anna janābaka law tafakkar fī hādhihi al-kalimāt al-mukhtaṣarāt
la-taʿrif kull al-ʿulūm wa-taṣil ilā dhirwat al-Maʿlūm wa-taqūl yakfī
kull al-wujūd min al-mashhūd wa-al-mafqūd.

Yet, were you to reflect on these brief words, you would come to know
all knowledge, attain the summit of the Known, and say: ``This suffices
for all existence, both what is seen and what is absent.''

\subsection{119}\label{section-2830}

\begin{otherlanguage}{arabic}

ولکن لو تجد فی نفسک حرارة المحبّة لتقول هل من مزید و نقول الحمد للّه ربّ
العالمین

\end{otherlanguage}

Wa-lākin law tajid fī nafsika ḥarārat al-maḥabbah la-taqūl hal min mazīd
wa-naqūl al-ḥamdu li-llāh Rabb al-ʿālamīn.

But if you find within yourself the heat of love, you will say: ``Is
there more?'' And We say: ``Praise be to God, the Lord of all worlds.''

\bookmarksetup{startatroot}

\chapter{BH00013 (Unnamed)}\label{bh00013-unnamed}

\emph{Translated on 17 Qudrat 183 using GPT 6.1 Sol}

\subsection{1}\label{section-2831}

\begin{otherlanguage}{arabic}

بسم ربّنا الأقدس الأعظم العلیّ الأبهی

\end{otherlanguage}

Bismi Rabbinā al-Aqdas al-Aʿẓam al-ʿAlī al-Abhā.

In the name of our Lord, the Most Holy, the Most Great, the Most
Exalted, the All-Glorious.

\subsection{2}\label{section-2832}

\begin{otherlanguage}{arabic}

البهآء لمن خرق الأحجاب فی یوم المآب و کسّر اصنام الأوهام باسم مالک الأیّام
و قام علی خدمة الأمر علی شأن ما خوّفته سطوة جبابرة الأرض الّذین زیّنوا
رؤوسهم بالعمائم و اخذوا زمام العباد باسم ربّهم مالک الایجاد کانوا ان
یدخلوا المساجد و المعابد لذکر اللّه و تلاوة ما انزله علی حبیبه و صفیّه
فلمّا انار افق العالم و اشرقت منه شمس اسمه الأعظم الّذی به ثبت امر کلّ نبیّ
و تحقّق رسالة کلّ رسول اعرضوا عنه و انکروه و نبذوا امره و نقضوا عهده و
ارتکبوا ما ناحت به الذّرّات و انصعقت الحوریّات فی الغرفات

\end{otherlanguage}

Al-bahāʾ li-man kharaqa al-aḥjāb fī yawm al-maʾāb wa-kassara aṣnām
al-awhām bi-ismi Mālik al-ayyām wa-qāma ʿalā khidmat al-amr ʿalā shaʾn
mā khawwafat-hu saṭwat jabābirat al-arḍ alladhīna zayyanū ruʾūsahum
bi-al-ʿamāʾim wa-akhadhū zimām al-ʿibād bi-ismi Rabbihim Mālik al-ījād.
Kānū an yadkhulū al-masājid wa-al-maʿābid li-dhikr Allāh wa-tilāwat mā
anzalahu ʿalā Ḥabībihi wa-Ṣafiyyihi. Fa-lammā anāra ufuq al-ʿālam
wa-ashraqat minhu shams ismihi al-Aʿẓam alladhī bihi thabata amr kull
nabī wa-taḥaqqaqa risālat kull rasūl aʿraḍū ʿanhu wa-ankarūhu wa-nabadhū
amrahu wa-naqaḍū ʿahdahu wa-irtakabū mā nāḥat bihi al-dharrāt
wa-inṣaʿaqat al-ḥūriyyāt fī al-ghurafāt.

Glory be upon whoever has torn through the veils on the Day of Return,
shattered the idols of vain imaginings in the name of the Lord of Days,
and arisen to serve the Cause in such a way that the might of the
tyrants of the earth has not frightened him---those who have adorned
their heads with turbans and seized the reins of the people in the name
of their Lord, the Lord of creation. They used to enter mosques and
places of worship to remember God and recite what He had sent down to
His Beloved and Chosen One. Yet when the horizon of the world was
illumined and the Sun of His Most Great Name rose above it---the Name
through which the cause of every prophet was established and the mission
of every messenger fulfilled---they turned away from Him, denied Him,
cast aside His Cause, broke His covenant, and committed deeds that made
the atoms lament and the heavenly maidens swoon in their chambers.

\subsection{3}\label{section-2833}

\begin{otherlanguage}{arabic}

لعمر اللّه باعمالهم ذرفت عیون اهل الفردوس الأعلی و الجنّة العلیا باسمه
افتخروا بین النّاس و وردوا علیه بعد ظهوره ما عجز عن ذکره ملکوت البیان قد
انکروا یوم اللّه الّذی تزیّن بذکره الکتب و الزّبر و الصّحف و الألواح یشهد
الخادم و کلّ الأشیآء بغفلتهم و جهلهم و ضلالتهم اولئک الّذین حقّت علیهم کلمة
العذاب فی الکتاب یکفیهم ما انزله اللّه فی الفرقان قوله تعالی و جعلنا لهم
سمعاً و ابصاراً و افئدةً فما اغنی عنهم سمعهم و لا ابصارهم و لا افئدتهم من
شیء اذ کانوا یجحدون بآیات اللّه و حاق بهم ما کانوا به یستهزئون قد اخذوا
کتاب انفسهم و نبذوا امّ الکتاب عن ورائهم لعمر اللّه هم الأصنام و مطالع
الظّنون و مشارق الأوهام و هم جبابرة الأرض و فراعنها قد تمسّکوا بالنّفس و
الهوی معرضین عن مولی الوری متّکئین علی اقلامهم معرضین علی ما نزل من القلم
الأعلی

\end{otherlanguage}

La-ʿamr Allāh, bi-aʿmālihim dharafat ʿuyūn ahl al-Firdaws al-Aʿlā
wa-al-jannat al-ʿulyā. Bi-ismihi iftakharū bayna al-nās wa-waradū
ʿalayhi baʿda ẓuhūrihi mā ʿajaza ʿan dhikrihi malakūt al-bayān. Qad
ankarū yawm Allāh alladhī tazayyana bi-dhikrihi al-kutub wa-al-zubur
wa-al-ṣuḥuf wa-al-alwāḥ. Yashhadu al-khādim wa-kull al-ashyāʾ
bi-ghaflatihim wa-jahlihim wa-ḍalālatihim. Ulāʾika alladhīna ḥaqqat
ʿalayhim kalimat al-ʿadhāb fī al-kitāb. Yakfīhim mā anzalahu Allāh fī
al-Furqān, qawluhu taʿālā: Wa-jaʿalnā lahum samʿan wa-abṣāran
wa-afʾidatan fa-mā aghnā ʿanhum samʿuhum wa-lā abṣāruhum wa-lā
afʾidatuhum min shayʾ idh kānū yajḥadūna bi-āyāt Allāh wa-ḥāqa bihim mā
kānū bihi yastahziʾūn. Qad akhadhū kitāb anfusihim wa-nabadhū Umm
al-Kitāb ʿan warāʾihim. La-ʿamr Allāh, hum al-aṣnām wa-maṭāliʿ al-ẓunūn
wa-mashāriq al-awhām wa-hum jabābirat al-arḍ wa-farāʿinatuhā. Qad
tamassakū bi-al-nafs wa-al-hawā muʿriḍīna ʿan Mawlā al-warā muttakiʾīna
ʿalā aqlāmihim muʿriḍīna ʿalā mā nazala min al-Qalam al-Aʿlā.

By the life of God! Their deeds have brought tears to the eyes of the
inhabitants of the most exalted Paradise and the highest Garden. In His
name they vaunted themselves among the people, yet after His appearance
they inflicted on Him what the kingdom of utterance is powerless to
describe. They have denied the Day of God, whose mention adorns the
books, scriptures, scrolls, and tablets. This servant and all things
bear witness to their heedlessness, ignorance, and error. These are the
ones against whom the sentence of punishment pronounced in the Book has
become binding. What God sent down in the Qur'an suffices them. Exalted
is His utterance: ``We gave them hearing, sight, and hearts, but neither
their hearing, nor their sight, nor their hearts availed them in the
least, since they denied the signs of God; and what they used to
ridicule engulfed them.'' They have taken up the book of their own
selves and cast the Mother Book behind their backs. By the life of God!
They are the idols, the sources of conjecture and the dawning places of
vain imaginings; they are the tyrants of the earth and its pharaohs.
They have clung to self and desire, turning away from the Lord of
humanity, relying on their own pens and turning away from what has
descended from the Most Exalted Pen.

\subsection{4}\label{section-2834}

\begin{otherlanguage}{arabic}

سبحانک یا من بآیاتک اهتزّ من فی القبور و بنفحات قمیصک طار الطّور عند اشراق
النّور من افق الظّهور اسئلک بسراج امرک و امواج بحر بیانک بأن تؤیّد اولیائک
علی ذکرک و ثنائک و علی القیام علی خدمتک ثمّ احفظهم یا الهی بقدرتک و قوّتک
عن الّذین انکروا حقّک و جاحدوا برهانک و جادلوا بآیاتک ثمّ اجعلهم یا الهی
اعلام نصرک فی بلادک و رایات عزّ ک فی دیارک ثمّ اسئلک یا محبوب العالم و
الظّاهر بالاسم الأعظم بأن توفّق من اقبل الی افقک الأعلی و قام علی خدمة
امرک بین الوری ای ربّ تراه ناصراً لدینک و قائماً علی صراطک و ناطقاً باسمک
بین عبادک ایّده بتأیید لا توقّفه شئونات العالم و لا زماجیر الأمم انت الّذی
فی قبضتک زمام القوّة و الاقتدار و یطوف حول مشیّتک مظاهر العظمة و الاختیار
لا اله الّا انت المقتدر المتعالی العزیز البصّار

\end{otherlanguage}

Subḥānaka yā man bi-āyātika ihtazza man fī al-qubūr wa-bi-nafaḥāt
qamīṣika ṭāra al-Ṭūr ʿinda ishrāq al-nūr min ufuq al-ẓuhūr. Asʾaluka
bi-sirāj amrika wa-amwāj baḥr bayānika bi-an tuʾayyida awliyāʾaka ʿalā
dhikrika wa-thanāʾika wa-ʿalā al-qiyām ʿalā khidmatika. Thumma iḥfaẓhum
yā Ilāhī bi-qudratika wa-quwwatika ʿan alladhīna ankarū ḥaqqaka
wa-jāḥadū burhānaka wa-jādalū bi-āyātika. Thumma ijʿalhum yā Ilāhī aʿlām
naṣrika fī bilādika wa-rāyāt ʿizzika fī diyārika. Thumma asʾaluka yā
Maḥbūb al-ʿālam wa-al-Ẓāhir bi-al-Ism al-Aʿẓam bi-an tuwaffiqa man
aqbala ilā ufuqika al-aʿlā wa-qāma ʿalā khidmat amrika bayna al-warā. Ay
Rabb, tarāhu nāṣiran li-dīnika wa-qāʾiman ʿalā ṣirāṭika wa-nāṭiqan
bi-ismika bayna ʿibādika. Ayyidhu bi-taʾyīd lā tūqifuhu shuʾūnāt
al-ʿālam wa-lā zamājīr al-umam. Anta alladhī fī qabḍatika zimām al-quwwa
wa-al-iqtidār wa-yaṭūfu ḥawla mashīʾatika maẓāhir al-ʿaẓama
wa-al-ikhtiyār. Lā ilāha illā Anta al-Muqtadir al-Mutaʿālī al-ʿAzīz
al-Baṣṣār.

Glory be to You, through whose verses those in the graves have stirred,
and through the fragrances of whose shirt the Mount took flight when the
light shone from the horizon of revelation. I ask You, by the lamp of
Your Cause and the waves of the ocean of Your utterance, to strengthen
Your loved ones in remembering and praising You and in arising to serve
You. Then protect them, my God, by Your power and might, from those who
have denied Your truth, rejected Your proof, and disputed Your verses.
Make them, my God, the standards of Your victory in Your lands and the
banners of Your glory in Your dominions. I further ask You, Beloved of
the world, who are manifest through the Most Great Name, to grant
success to him who has turned toward Your most exalted horizon and
arisen to serve Your Cause among humanity. Lord, You see him championing
Your religion, standing firm on Your path, and speaking Your name among
Your servants. Strengthen him with a strength that neither the affairs
of the world nor the clamors of the nations can arrest. You are the One
who holds in His grasp the reins of strength and power, and around whose
will revolve the manifestations of majesty and sovereignty. There is no
God but You, the All-Powerful, the Most Exalted, the Almighty, the
All-Seeing.

\subsection{5}\label{section-2835}

\begin{otherlanguage}{arabic}

و بعد یا حبیب قلبی و محبوب فؤادی قد خبرنی صفیر طیر بیانکم و اجتذبنی تغرّد
عندلیب ثنائکم و اخذنی سکر رحیق ذکرکم علی شأن اردت ان اطیر بأجنحة
الاشتیاق الی مقام قدّسه اللّه عن ذکر من فی الآفاق فلمّا حضرت و فزت بأنوار
شمس العزّ عرضت ما فی کتابکم عند من نزّل امّ الکتاب بحرکة من قلمه الأعلی فی
ملکوت الانشآء و بعدما فاز ما فیه باصغآء مولی الوری اذاً نطق لسان الکبریآء
قوله جلّ کبریائه

\end{otherlanguage}

Wa-baʿd, yā ḥabīb qalbī wa-maḥbūb fuʾādī, qad khabbaranī ṣafīr ṭayr
bayānikum wa-ijtadhabanī tagharrud ʿandalīb thanāʾikum wa-akhadhanī sukr
raḥīq dhikrikum ʿalā shaʾn aradtu an aṭīra bi-ajniḥat al-ishtiyāq ilā
maqām qaddasahu Allāh ʿan dhikr man fī al-āfāq. Fa-lammā ḥaḍartu
wa-fuztu bi-anwār shams al-ʿizz ʿaraḍtu mā fī kitābikum ʿinda man
nazzala Umm al-Kitāb bi-ḥaraka min qalamihi al-aʿlā fī malakūt
al-inshāʾ. Wa-baʿdamā fāza mā fīhi bi-iṣghāʾ Mawlā al-warā idhan naṭaqa
Lisān al-Kibriyāʾ, qawluhu jalla kibriyāʾuhu:

To continue, beloved of my heart and darling of my soul: the call of the
bird of your utterance has brought me tidings, the song of the
nightingale of your praise has drawn me, and the intoxication of the
choice wine of your remembrance has so seized me that I wished to fly on
wings of longing to a station God has sanctified beyond the mention of
all who dwell in the world. When I entered the presence and attained the
radiance of the Sun of Glory, I presented the contents of your letter
before Him who sent down the Mother Book through a movement of His Most
Exalted Pen in the kingdom of creation. After its contents had received
the attention of the Lord of humanity, the Tongue of Grandeur spoke
these words---glorified be His grandeur:

\subsection{6}\label{section-2836}

\begin{otherlanguage}{arabic}

یا ورقآء یا ایّها النّاظر الی الأفق الأعلی قل

\end{otherlanguage}

Yā Warqāʾ, yā ayyuhā al-nāẓir ilā al-ufuq al-aʿlā, qul:

``Varqá, you who gaze toward the most exalted horizon, say:

\subsection{7}\label{section-2837}

\begin{otherlanguage}{arabic}

لک الحمد یا الهی و المقدّس عن ذکری و بیانی و عمّا تفوّه به خلقک فی مملکتک و
نطق به عبادک فی ارضک بما جعلتنی سائراً فی دیار المعانی الی ان بلغت قطبها
المقام الّذی سمّی بالبقعة النّورآء و الأرض البیضآء و سمیّتها فی مقام آخر
بالأرض المحشر و اریتنی فیها جمالک و سقیتنی رحیق بیانک و عرّفتنی ما اخذنی
عنّی و انسانی ما دونک اسئلک بهذه الموهبة الکبری و العطیّة العظمی بأن تحفظ
لی هذا المقام الأعلی انّک انت مولی الآخرة و الأولی یشهد هذا الفانی یا
الهی ببقائک و هذا الضّعیف بقدرتک و قوّتک و هذا العاشق المسکین باشراقات
انوار شمس جمالک انت الّذی یا الهی و سیّدی و سندی و منای تجلّیت بحرف من
کلمتک العلیا علی طور السّینآء اذاً نطقت الأغصان و الأوراق و عن ورائها
الذّرّات بما نطقت به السّدرة فی طور العرفان من لدی الرّحمن و تجلّیت بحرف اخری
علی ملکوت البیان اذاً نزّلت التّوریة و الانجیل و الفرقان و بقطرة من فرات
رحمتک الجاریة من قلمک الأعلی قد ظهر کوثر الحیوان فی الامکان تعالت قدرتک
و عظمت سلطنتک

\end{otherlanguage}

Laka al-ḥamd yā Ilāhī wa-al-Muqaddas ʿan dhikrī wa-bayānī wa-ʿammā
tafawwaha bihi khalquka fī mamlakatika wa-naṭaqa bihi ʿibāduka fī arḍika
bi-mā jaʿaltanī sāʾiran fī diyār al-maʿānī ilā an balaghtu quṭbahā
al-maqām alladhī summiya bi-al-buqʿat al-nawrāʾ wa-al-arḍ al-bayḍāʾ
wa-sammaytahā fī maqām ākhar bi-al-arḍ al-maḥshar wa-araytanī fīhā
jamālaka wa-saqaytanī raḥīq bayānika wa-ʿarraftanī mā akhadhanī ʿannī
wa-ansānī mā dūnaka. Asʾaluka bi-hādhihi al-mawhibat al-kubrā
wa-al-ʿaṭiyyat al-ʿuẓmā bi-an taḥfaẓa lī hādhā al-maqām al-aʿlā. Innaka
Anta Mawlā al-ākhira wa-al-ūlā. Yashhadu hādhā al-fānī yā Ilāhī
bi-baqāʾika wa-hādhā al-ḍaʿīf bi-qudratika wa-quwwatika wa-hādhā
al-ʿāshiq al-miskīn bi-ishrāqāt anwār shams jamālika. Anta alladhī yā
Ilāhī wa-Sayyidī wa-Sanadī wa-Munāya tajallayta bi-ḥarf min kalimatika
al-ʿulyā ʿalā Ṭūr al-Sīnāʾ, idhan naṭaqat al-aghṣān wa-al-awrāq wa-ʿan
warāʾihā al-dharrāt bi-mā naṭaqat bihi al-Sidra fī Ṭūr al-ʿirfān min
ladā al-Raḥmān. Wa-tajallayta bi-ḥarf ukhrā ʿalā malakūt al-bayān, idhan
nazzalta al-Tawrāh wa-al-Injīl wa-al-Furqān. Wa-bi-qaṭra min furāt
raḥmatika al-jāriya min qalamika al-aʿlā qad ẓahara Kawthar al-ḥayawān
fī al-imkān. Taʿālat qudratuka wa-ʿaẓumat salṭanatuka.

Praise be to You, my God, who are sanctified beyond my remembrance and
expression and beyond whatever Your creatures have spoken in Your
kingdom and Your servants have uttered on Your earth, for having enabled
me to journey through the realms of inner meaning until I reached their
center---that station called the Luminous Spot and the White Land, which
in another place You named the Land of Resurrection. There You showed me
Your beauty, gave me the choice wine of Your utterance to drink, and
made known to me what carried me beyond myself and caused me to forget
all but You. I ask You, by this greatest favor and this supreme gift, to
preserve for me this most exalted station. You are indeed the Lord of
the next world and of this world. This mortal being, my God, bears
witness to Your eternity; this weak one, to Your power and might; and
this destitute lover, to the streaming radiance of the sun of Your
beauty. You, my God, my Master, my Support, and my Desire, are the One
who revealed Yourself through a letter of Your most exalted Word upon
Mount Sinai. Then the branches and leaves, and beyond them the atoms,
uttered what the Lote Tree uttered on the Mount of recognition from the
All-Merciful. Through another letter You revealed Yourself to the
kingdom of utterance, and then You sent down the Torah, the Gospel, and
the Qur'an. Through a drop from the sweet stream of Your mercy flowing
from Your Most Exalted Pen, the river of life appeared in the world of
existence. Exalted is Your power, and great is Your sovereignty.

\subsection{8}\label{section-2838}

\begin{otherlanguage}{arabic}

اسئلک یا ایّها الضّارب النّابض فی جسد العالم باصبعک الأعظم بأن تقرّبنی فی کلّ
الأحوال الیک و تذکّرنی الأیّام الّتی مصصت کوثر وصالک من عرق فضلک و عطائک
انت الّذی یا الهی دعوتنی الیک و ادخلتنی فی قباب عظمتک و اقمتنی تلقآء
الوجه علی خدمتک و ذکرک و ثنائک ای ربّ تری و تعلم انّ فرحی فی ذکرک و اقبالی
الیک و بهجتی فی فضلک و عنایتک فأنزل علیّ من سحاب رحمتک ما یجعلنی متبسّماً
فی ایّامک انّک انت مالک الجود و العطآء و فی قبضتک ملکوت الأسمآء لا اله الّا
انت المقتدر العلیم الحکیم انتهی

\end{otherlanguage}

Asʾaluka yā ayyuhā al-ḍārib al-nābiḍ fī jasad al-ʿālam bi-iṣbaʿika
al-aʿẓam bi-an tuqarribanī fī kull al-aḥwāl ilayka wa-tudhakkiranī
al-ayyām allatī maṣaṣtu Kawthar wiṣālika min ʿirq faḍlika wa-ʿaṭāʾika.
Anta alladhī yā Ilāhī daʿawtanī ilayka wa-adkhaltanī fī qibāb ʿaẓamatika
wa-aqamtanī tilqāʾ al-wajh ʿalā khidmatika wa-dhikrika wa-thanāʾika. Ay
Rabb, tarā wa-taʿlamu anna faraḥī fī dhikrika wa-iqbālī ilayka
wa-bahjatī fī faḍlika wa-ʿināyatika. Fa-anzil ʿalayya min saḥāb
raḥmatika mā yajʿalunī mutabassiman fī ayyāmika. Innaka Anta Mālik
al-jūd wa-al-ʿaṭāʾ wa-fī qabḍatika malakūt al-asmāʾ. Lā ilāha illā Anta
al-Muqtadir al-ʿAlīm al-Ḥakīm. Intahā.

I ask You, who set the pulse beating in the body of the world with Your
Most Great Finger, to draw me near to You in every condition and remind
me of the days when I drank the river of reunion with You from the vein
of Your grace and bounty. You, my God, are the One who called me to
Yourself, admitted me beneath the canopies of Your majesty, and
established me before Your face to serve, remember, and praise You.
Lord, You see and know that my joy lies in remembering You and turning
toward You, and my delight in Your grace and loving care. Send down upon
me from the clouds of Your mercy what will keep me smiling in Your days.
You are indeed the Lord of generosity and bounty, and in Your grasp is
the kingdom of names. There is no God but You, the All-Powerful, the
All-Knowing, the All-Wise.'' End.

\subsection{9}\label{section-2839}

\begin{otherlanguage}{arabic}

هنیئاً لحضرتک یا حبیب فؤادی قد شربت مآء العنایة و الحیوة من کؤوس الطاف
محبوبی و محبوبک و محبوب من فی السّموات و الأرض علی شأن ینبغی ان تقول
الذّرّات هنیئاً لک و مریئاً لک یا من فزت بهذا الفضل العظیم

\end{otherlanguage}

Hanīʾan li-ḥaḍratika yā ḥabīb fuʾādī. Qad sharibta māʾ al-ʿināya
wa-al-ḥayāh min kuʾūs alṭāf Maḥbūbī wa-Maḥbūbika wa-Maḥbūb man fī
al-samāwāt wa-al-arḍ ʿalā shaʾn yanbaghī an taqūla al-dharrāt hanīʾan
laka wa-marīʾan laka yā man fuzta bi-hādhā al-faḍl al-ʿaẓīm.

May it bring you joy, beloved of my heart! You have drunk the water of
loving care and life from the cups of the favors of my Beloved, your
Beloved, and the Beloved of all who are in the heavens and on earth, in
such a measure that the atoms should say, ``May it bring you joy and
refreshment, you who have attained this great bounty!''

\subsection{10}\label{section-2840}

\begin{otherlanguage}{persian}

اینکه مرقوم داشته بودید که جناب آقا میرزا عبداللّه ٩ ٦٦ {[}بهآء اللّه{]}
زحمات آن محبوب را حمل نموده‌اند و طلب مکافات نمودید از حقّ جلّ جلاله در
بارهٴ ایشان این مراتب در ساحت امنع اقدس اعلی عرض شد هذا ما نطق به لسان
اللّه قوله عزّ اجلاله یا ورقآء مالک اسمآء میفرماید عنایتی در بارهٴ او از
قلم اعلی جاری شده که ما فی العالم معادله بآن نمینماید اینست اصل مکافات
که نزد حقّ جلّ جلاله محبوب بوده و هست ولکن عنایات ظاهره هم انشآءاللّه میرسد
بشّره من قبلی بما نطق به لسانی و انا الآمر الحکیم انتهی

\end{otherlanguage}

Īnkih marqūm dāshtih būdīd kih Jināb-i Āqā Mīrzā ʿAbduʾllāh 9 66
{[}Bahāʾ Allāh{]} zaḥamāt-i ān maḥbūb rā ḥaml nimūdih-and va ṭalab-i
mukāfāt nimūdīd az Ḥaqq jalla jalāluhu dar bārih-yi īshān, īn marātib
dar sāḥat-i amnaʿ-i aqdas-i aʿlā ʿarḍ shud. Hādhā mā naṭaqa bihi Lisān
Allāh, qawluhu ʿazza ijlāluhu: Yā Warqāʾ, Mālik-i asmāʾ mīfarmāyad
ʿināyatī dar bārih-yi ū az Qalam-i Aʿlā jārī shudih kih mā fī al-ʿālam
muʿādilih bi-ān namīnamāyad. Īnast aṣl-i mukāfāt kih nazd-i Ḥaqq jalla
jalāluhu maḥbūb būdih va hast, va-lākin ʿināyāt-i ẓāhirih ham in shāʾ
Allāh mīrasad. Bashshirhu min qibalī bi-mā naṭaqa bihi lisānī wa-anā
al-Āmir al-Ḥakīm. Intahā.

You had written that Áqá Mírzá `Abdu'lláh, 9 66 {[}the glory of God{]},
had borne the burdens of that beloved one, and you asked God---glorified
be His glory---to reward him. These matters were presented at the most
inaccessible, most holy, most exalted Court. This is what the Tongue of
God uttered---exalted be His majesty: ``Varqá, the Lord of Names
declares: A favor has flowed from the Most Exalted Pen concerning him
that nothing in the world can equal. This is the essence of recompense,
which has ever been and remains beloved in the sight of God---glorified
be His glory. Nevertheless, outward favors will also reach him, God
willing. Give him, on My behalf, the glad tidings of what My tongue has
uttered. I am the Ordainer, the All-Wise.'' End.

\subsection{11}\label{section-2841}

\begin{otherlanguage}{persian}

اینکه در بارهٴ قرّهٴ عین میرزا روح اللّه مرقوم داشتید عرض شد هذا ما نزل من
ملکوت المشیّة و البیان قوله جلّ کبریائه یا ورقآء لک ان تقرء فی اذنه الیمنی
ثلاث مرّات قد جئت بأمر اللّه و ظهرت لذکره و خلقت لخدمته العزیز المحبوب انّا
ذکرناه من قبل اذ استغاثت {[}امّه{]} و فی هذا الحین و انا الفضّال المبین ان
اذکر من قبلی امّها و بشّرها بما نطق به لسان المظلوم فی سجنه العظیم انتهی

\end{otherlanguage}

Īnkih dar bārih-yi qurrat-i ʿayn Mīrzā Rūḥuʾllāh marqūm dāshtīd ʿarḍ
shud. Hādhā mā nazala min malakūt al-mashīʾa wa-al-bayān, qawluhu jalla
kibriyāʾuhu: Yā Warqāʾ, laka an taqraʾa fī udhunihi al-yumnā thalātha
marrāt: Qad jiʾta bi-amr Allāh wa-ẓaharta li-dhikrihi wa-khuliqta
li-khidmatihi al-ʿAzīz al-Maḥbūb. Innā dhakarnāhu min qabl idh
istaghāthat {[}ummuhu{]} wa-fī hādhā al-ḥīn wa-anā al-Faḍḍāl al-Mubīn.
An udhkur min qibalī ummahā wa-bashshirhā bi-mā naṭaqa bihi lisān
al-Maẓlūm fī sijnihi al-ʿaẓīm. Intahā.

What you wrote concerning the delight of the eye, Mírzá Rúḥu'lláh, was
presented. This is what descended from the kingdom of will and
utterance---glorified be His grandeur: ``Varqá, you are to recite three
times in his right ear: `You have come by the command of God, appeared
to remember Him, and been created to serve Him, the Almighty, the
Beloved.' We made mention of him before, when {[}his mother{]} called
for help, and at this time as well. I am the Most Bountiful, the
All-Manifest. Remember her mother on My behalf, and give her the glad
tidings of what the tongue of the Wronged One has uttered in His mighty
prison.'' End.

\subsection{12}\label{section-2842}

\begin{otherlanguage}{persian}

و اینکه مرقوم فرموده بودید در بعض ادعیه و مناجات منزله نازل و اجعله اگر
در مقام رجا و اجعلنی که اضافه بقائل است تلاوت شود از تحریف محسوب یا نه
عرض میشود این فقره تحریف نیست و هر دو قسم جایز است و هم‌چنین در کلمهٴ
مبارکهٴ قل اللّهمّ و قل لک الحمد حذف کلمهٴ قل و ذکر آن هر دو صحیح است امّا
ذکر نظر بشرافت کلمهٴ مبارکه که از فم مالک احدیّه ظاهر شده و حذف آن جایز
چه که ظهور عبودیّت و اختصاص در اینمقام اظهر است

\end{otherlanguage}

Va īnkih marqūm farmūdih būdīd dar baʿḍ-i adʿiyih va munājāt-i
munazzalih nāzil va-ijʿalhu agar dar maqām-i rajā va-ijʿalnī kih iḍāfih
bi-qāʾil ast tilāvat shavad az taḥrīf maḥsūb yā nah, ʿarḍ mīshavad īn
fiqrih taḥrīf nīst va har du qism jāʾiz ast. Va hamchunīn dar kalimih-yi
mubārakih-yi qul Allāhumma va qul laka al-ḥamd ḥadhf-i kalimih-yi qul va
dhikr-i ān har du ṣaḥīḥ ast. Ammā dhikr naẓar bi-sharāfat-i kalimih-yi
mubārakih kih az fam-i Mālik-i Aḥadiyyih ẓāhir shudih va ḥadhf-i ān
jāʾiz chih kih ẓuhūr-i ʿubūdiyyat va ikhtiṣāṣ dar īn maqām aẓhar ast.

You had also written that ``and make him'' occurs in some revealed
prayers and supplications, asking whether it would be considered an
alteration if, when supplicating, one recited ``and make me,'' which
refers to the speaker. The reply is that this is not an alteration, and
both forms are permissible. Likewise, in the blessed expressions ``Say:
O God'' and ``Say: Praise be to You,'' both omitting the word ``Say''
and retaining it are correct. Retaining it honors the blessed word that
has issued from the mouth of the Lord of Oneness, while omitting it is
permissible because the expression of servitude and personal application
is more evident in this case.

\subsection{13}\label{section-2843}

\begin{otherlanguage}{persian}

اینکه در بارهٴ جناب آقا میرزا غلامعلی علیه ٩ ٦٦ {[}بهآء اللّه{]} مرقوم
داشتید و مراتب اقبال و استقامت و خلوص و خدمتشان را این تفصیل لدی العرش
معروض قال جلّت عظمته الحمد للّه بحرارت محبّت مشتعلند و بنور معرفت مزیّن کوثر
وصال که اثمن جواهر عالم است و اغلی از مآء حیوان الحمد للّه بحفظ آن فائز
شده و ثمرهٴ آن را ظاهر نموده انشآءاللّه همیشهٴ ایّام این آثار از او ظاهر
باشد طوبی از برای نفسیکه بلقا فائز شد و به ما یلیق لهذا المقام قیام نمود
انا نکبّر من هذا المقام علیه و علی احبّائی فی هناک الّذین نبذوا العالم و
اخذوا ما امروا به من لدی اللّه مالک القدم انتهی

\end{otherlanguage}

Īnkih dar bārih-yi Jināb-i Āqā Mīrzā Ghulām-ʿAlī ʿalayhi 9 66 {[}Bahāʾ
Allāh{]} marqūm dāshtīd va marātib-i iqbāl va istiqāmat va khulūṣ va
khidmatishān rā, īn tafṣīl ladā al-ʿarsh maʿrūḍ. Qāla jallat ʿaẓamatuhu:
Al-ḥamdu li-llāh bi-ḥarārat-i maḥabbat mushtaʿiland va bi-nūr-i maʿrifat
muzayyan. Kawthar-i viṣāl kih athman-i javāhir-i ʿālam ast va aghlā az
māʾ-i ḥayavān, al-ḥamdu li-llāh bi-ḥifẓ-i ān fāʾiz shudih va th
amarih-yi ān rā ẓāhir nimūdih. In shāʾ Allāh hamīshih-yi ayyām īn āthār
az ū ẓāhir bāshad. Ṭūbā az barā-yi nafsīkih bi-liqā fāʾiz shud va bi-mā
yalīqu li-hādhā al-maqām qiyām nimūd. Innā nukabbiru min hādhā al-maqām
ʿalayhi wa-ʿalā aḥibbāʾī fī hunāk alladhīna nabadhū al-ʿālam wa-akhadhū
mā umirū bihi min ladā Allāh Mālik al-qidam. Intahā.

What you wrote concerning Áqá Mírzá Ghulám-`Alí---upon him be 9 66
{[}the glory of God{]}---and the particulars of his devotion,
steadfastness, sincerity, and service was presented before the Throne.
He said---glorified be His majesty: ``Praise be to God, he is aflame
with the heat of love and adorned with the light of recognition. The
river of reunion is the most precious of the world's jewels and more
costly than the water of life. Praise be to God, he has succeeded in
preserving it and has manifested its fruit. God willing, these signs
will always appear through him. Blessed is the soul who has attained the
meeting and arisen to do what befits this station. From this station We
send Our greeting of glorification upon him and upon My loved ones
there, who have cast aside the world and embraced what they were
commanded by God, the Lord of eternity.'' End.

\subsection{14}\label{section-2844}

\begin{otherlanguage}{persian}

و اینکه در بارهٴ جناب آقا یوسف (من اهل نراق) علیه ٩ ٦٦ {[}بهآء اللّه{]}
مرقوم داشتید بساحت امنع اقدس اعلی عرض شد هذا ما نزل من سمآء مشیّة ربّنا
الحکیم

\end{otherlanguage}

Va īnkih dar bārih-yi Jināb-i Āqā Yūsuf (min ahl-i Narāq) ʿalayhi 9 66
{[}Bahāʾ Allāh{]} marqūm dāshtīd bi-sāḥat-i amnaʿ-i aqdas-i aʿlā ʿarḍ
shud. Hādhā mā nazala min samāʾ mashīʾat Rabbinā al-Ḥakīm.

What you wrote concerning Áqá Yúsuf (of Naráq)---upon him be 9 66 {[}the
glory of God{]}---was presented at the most inaccessible, most holy,
most exalted Court. This is what descended from the heaven of the will
of our Lord, the All-Wise.

\subsection{15}\label{section-2845}

\begin{otherlanguage}{arabic}

هو الأعظم الأقدس الأبهی

\end{otherlanguage}

Huwa al-Aʿẓam al-Aqdas al-Abhā.

He is the Most Great, the Most Holy, the All-Glorious.

\subsection{16}\label{section-2846}

\begin{otherlanguage}{persian}

یا یوسف یذکرک المظلوم بلحن لو یسمعه کلّ الوجود من الغیب و الشّهود لیدعنّ
الحیوة و البقآء مسرعین الی مقرّ الفدآء فی هذا السّبیل الّذی اذا ظهر نادی
الصّراط المستقیم قد اتی الموعود بصراطه الظّاهر الواضح المبین تاللّه قد خلقت
لهذا الیوم لأطوف حول هذا الصّراط و ابشّر النّاس امامه بهذا الظّهور الّذی به
اخذ اللّه عهده عن النّبیّین و المرسلین یا یوسف قد اهتزّ القمیص من نسمة ارادة
ربّک و به سرت نسمة اللّه علی البلدان و تضوّع عرف الرّحمن بین السّموات و
الأرضین انّا ذکرناک من قبل و فی هذا الحین الّذی استوی المظلوم علی عرش اسمه
العظیم انّا نوصیک و الّذین آمنوا بما تطلع به شمس الأخلاق من افق قلب سلیم
یا اولیائی ان انصروا الرّحمن بالأخلاق المرضیّة هذه کلمة انزلناها فی الواح
شتّی یشهد بذلک من نطق فی کلّ شأن انّه لا اله الّا انا الآمر العلیم البهآء
علی اهل البهآء الّذین شربوا رحیق البقآء من کأس عطآء ربّهم الکریم انتهی

\end{otherlanguage}

Yā Yūsuf, yadhkuruka al-Maẓlūm bi-laḥnin law yasmaʿuhu kullu al-wujūd
min al-ghayb wa-l-shuhūd la-yadaʿunna al-ḥayāta wa-l-baqāʾ musriʿīna ilā
maqarri al-fidāʾ fī hādhā al-sabīl alladhī idhā ẓahara nādā al-ṣirāṭ
al-mustaqīm qad atā al-Mawʿūd bi-ṣirāṭihi al-ẓāhir al-wāḍiḥ al-mubīn.
Ta-Allāh qad khuliqtu li-hādhā al-yawm li-aṭūfa ḥawla hādhā al-ṣirāṭ
wa-ubashshira al-nās amāmahu bi-hādhā al-ẓuhūr alladhī bihi akhadha
Allāh ʿahdahu ʿan al-nabiyyīn wa-l-mursalīn. Yā Yūsuf qad ihtazza
al-qamīṣ min nasamati irādati rabbika wa-bihi sarat nasamatu Allāh ʿalā
al-buldān wa-taḍawwaʿa ʿarfu al-Raḥmān bayna al-samāwāt wa-l-araḍīn.
Innā dhakarnāka min qablu wa-fī hādhā al-ḥīn alladhī istawā al-Maẓlūm
ʿalā ʿarshi ismihi al-ʿAẓīm. Innā nūṣīka wa-alladhīna āmanū bi-mā
taṭluʿu bihi shamsu al-akhlāq min ufuqi qalbin salīm. Yā awliyāʾī an
inṣurū al-Raḥmān bi-l-akhlāq al-marḍiyya. Hādhihi kalimatun anzalnāhā fī
alwāḥin shattā yashhadu bi-dhālika man naṭaqa fī kulli shaʾn innahu lā
ilāha illā anā al-Āmir al-ʿAlīm. Al-Bahāʾ ʿalā ahli al-Bahāʾ alladhīna
sharibū raḥīqa al-baqāʾ min kaʾsi ʿaṭāʾi rabbihim al-Karīm. Intahā.

O Yūsuf, the Wronged One remembers you in a melody which, if all
existence, unseen and visible, were to hear it, would cause all to
forsake life and continued existence, hastening to the abode of
sacrifice in this Path. When it appeared, the Straight Path called out:
``The Promised One has come with His manifest, clear, and unmistakable
Path. By God, I was created for this Day, to circle around this Path
and, in its presence, announce to humanity this Revelation, through
which God took His covenant from the Prophets and Messengers.'' O Yūsuf,
the shirt has stirred in the breeze of your Lord's will, and through it
the breeze of God has wafted over the lands and the fragrance of the
All-Merciful has spread between the heavens and the earths. We have
remembered you before, and at this time when the Wronged One has
ascended the throne of His mighty Name. We counsel you and those who
believe to observe that through which the sun of character rises from
the horizon of a sound heart. O My loved ones, aid the All-Merciful
through praiseworthy character. This is a word We have sent down in
various Tablets. To this bears witness the One who has declared in every
circumstance: ``There is no God but Me, the Ordainer, the All-Knowing.''
Glory be upon the people of Bahāʾ who have drunk the choice wine of
eternity from the cup of the bounty of their generous Lord. End.

\subsection{17}\label{section-2847}

\begin{otherlanguage}{persian}

انشآءاللّه کلّ بما ینبغی فائز شوند کینونت عالم یک لسان مشاهده میشود و
بأعلی النّدآء بکلمهٴ قد اتی مالک السّمآء ناطق ولکن احدی باصغاء فائز نه الّا
من شآء اللّه عجب در اینکه خود او گویا و خود او غافل سبحان من امرنی بالصّمت
و ستر امره العزیز العظیم مع آنکه کشف بمقامی رسیده که السن و اقلام از ذکر
آن عاجز است ایضاً سبحان من ظهر و اظهر ما اراد و نطق و انطق کلّشیء بأنّه هو
المقتدر العلیم الخبیر آفتاب حقیقت بکمال ظهور ظاهر و مشرق و سراج الهی
بکمال نور روشن و مجلّی مع‌ذلک کلّ محجوب و محروم از حقّ این خادم فانی میطلبد
و هم بطلبید که امراض مزمنهٴ ناس را که از جهل و نادانی عارض شده بدریاق
عنایت خود معالجه فرماید انّه هو المبرء الحکیم الحمد للّه اذ هو المقتدر
الغفور الرّحیم

\end{otherlanguage}

In shāʾ Allāh koll be-mā yanbaghī fāʾez shavand. Keynūnat-e ʿālam yek
lesān moshāhedeh mīshavad va bi-aʿlā al-nidāʾ be-kalemeh-ye qad atā
Mālik al-samāʾ nāṭeq, va-līkan aḥadī be-eṣghā fāʾez nah illā man shāʾa
Allāh. ʿAjab dar īn-keh khod-e ū gūyā va khod-e ū ghāfel. Subḥāna man
amaranī bi-l-ṣamt wa-satri amrihi al-ʿAzīz al-ʿAẓīm. Maʿ ān-keh kashf
be-maqāmī resīdeh keh alson va aqlām az dhekr-e ān ʿājez ast. Ayḍan
subḥāna man ẓahara wa-aẓhara mā arāda wa-naṭaqa wa-anṭaqa kulla shayʾ
bi-annahu huwa al-Muqtadir al-ʿAlīm al-Khabīr. Āftāb-e ḥaqīqat
be-kamāl-e ẓohūr ẓāher va moshreq va serāj-e elāhī be-kamāl-e nūr
rowshan va mojallā. Maʿa-dhālek koll maḥjūb va maḥrūm. Az Ḥaqq īn
khādem-e fānī mīṭalabad va ham beṭalabīd keh amrāḍ-e mozmenah-ye nās rā
keh az jahl va nādānī ʿāreż shodeh be-deryāq-e ʿenāyat-e khod moʿālejeh
farmāyad. Innahu huwa al-Mubriʾ al-Ḥakīm. Al-ḥamdu li-llāh idh huwa
al-Muqtadir al-Ghafūr al-Raḥīm.

God willing, may all attain what is fitting. The very being of the world
is seen as a single tongue, proclaiming at the highest pitch the words,
``The Lord of heaven has come,'' yet no one has attained a hearing ear
except those whom God wills. The wonder is that it is itself speaking
and itself heedless. Glorified be the One who has commanded me to remain
silent and conceal His mighty and exalted Cause, even though its
disclosure has reached a degree that tongues and pens are powerless to
describe. Again, glorified be the One who has appeared and revealed
whatever He desired, who has spoken and caused all things to declare
that He is the All-Powerful, the All-Knowing, the All-Aware. The Sun of
Truth is manifest and shining in the fullness of its revelation, and the
divine lamp is bright and radiant in the fullness of its light.
Nevertheless, all are veiled and deprived. This lowly servant asks
God---and you, too, should ask---that He heal, with the antidote of His
loving-kindness, the chronic illnesses that have afflicted people
through ignorance and lack of understanding. He is indeed the Healer,
the All-Wise. Praise be to God, for He is the All-Powerful, the
Ever-Forgiving, the Merciful.

\subsection{18}\label{section-2848}

\begin{otherlanguage}{persian}

و اینکه در بارهٴ جناب آقا علی محمّد ابن جناب حاجی احمد میلانی علیهما بهآء
اللّه مرقوم داشتید آنچه خواسته بودید مخصوص ایشان عنایت شد و بتوسّط اخویشان
ارسال گشت و حال در این حین که این عبد بتحریر این مکتوب مشغول است حضرت
غصن اللّه الأکبر روحی و ذاتی و کینونتی لتراب قدومه الأطهر فداء تشریف
ندارند باطراف این مدینه تشریف برده‌اند انشآءاللّه بعد از حضور و شرفیابی
این خادم فانی خطّ جدید که ارسال فرموده بودند بنظر مبارکشان میرسانم و آنچه
مرقوم داشته بودید معروض میدارم

\end{otherlanguage}

Va īn-keh dar bāreh-ye Janāb-e Āqā ʿAlī-Moḥammad ibn-e Janāb-e Ḥājī
Aḥmad-e Mīlānī, ʿalayhimā Bahāʾ Allāh, marqūm dāshtīd, āncheh khwāsteh
būdīd makhṣūṣ-e īshān ʿenāyat shod va be-tavassoṭ-e akhavī-shān ersāl
gasht. Va ḥāl dar īn ḥīn keh īn ʿabd be-taḥrīr-e īn maktūb mashghūl ast
Ḥażrat-e Ghoṣn Allāh al-Akbar, rūḥī wa-dhātī wa-kaynūnatī li-turābi
qudūmihi al-aṭhar fidāʾ, tashrīf nadārand; be-aṭrāf-e īn madīneh tashrīf
bordeh-and. In shāʾ Allāh baʿd az ḥożūr va sharafyābī īn khādem-e fānī
khaṭṭ-e jadīd keh ersāl farmūdeh būdand be-naẓar-e mobārak-shān
mīresānam va āncheh marqūm dāshteh būdīd maʿrūż mīdāram.

As to what you wrote concerning Āqā ʿAlī-Muḥammad, son of Ḥājī Aḥmad
Mīlānī---upon them both be the glory of God---what you requested was
graciously granted especially for him and was sent through his brother.
Now, at this moment when this servant is engaged in writing this letter,
His Holiness the Greater Branch of God---may my spirit, my being, and my
very existence be sacrificed for the dust beneath His most pure
feet---is not here; He has gone to the outskirts of this city. God
willing, after His return, when this lowly servant is honored to enter
His presence, I will bring the new letter that had been sent to His
blessed attention and present what you had written.

\subsection{19}\label{section-2849}

\begin{otherlanguage}{persian}

و اینکه ذکر جناب آقا فرج آقا و مراتب خلوص او را ذکر فرمودید انشآءاللّه در
کلّ احیان بخدمت محبوب امکان موفّق باشند و بذکر و ثنایش مؤیّد و اینکه در
بارهٴ امانت استفسار رفته بود تمام و کمال رسید انّ الخادم یشهد بذلک و عن
ورائه من خلقه و فطره و انطقه بثنائه الجمیل و چون ذکر ایشانرا نمودید و
این عبد هم دوست داشت که ذکری از ایشان در ساحت امنع اقدس اعلی عرض شود لذا
این تفصیل در لیل هفتم شهر صفر تلقاء وجه عرض شد و این آیات باهرات و کلمات
تامّات از سماء عنایت منزل نازل

\end{otherlanguage}

Va īn-keh dhekr-e Janāb-e Āqā Faraj Āqā va marāteb-e kholūṣ-e ū rā dhekr
farmūdīd, in shāʾ Allāh dar koll-e aḥyān be-khedmat-e Maḥbūb-e emkān
movaffaq bāshand va be-dhekr va thanāyash moʾayyad. Va īn-keh dar
bāreh-ye amānat estefsār rafteh būd, tamām va kamāl resīd. Inna
al-khādima yashhadu bi-dhālika wa-ʿan warāʾihi man khalaqahu wa-faṭarahu
wa-anṭaqahu bi-thanāʾihi al-jamīl. Va chūn dhekr-e īshān rā namūdīd va
īn ʿabd ham dūst dāsht keh dhekrī az īshān dar sāḥat-e amnaʿ-e aqdas-e
aʿlā ʿarż shavad, ledhā īn tafṣīl dar layl-e haftom-e shahr-e Ṣafar
telqāʾ-e vajh ʿarż shod va īn āyāt-e bāherāt va kalemāt-e tāmmāt az
samāʾ-e ʿenāyat-e Monzel nāzel.

As to your mention of Āqā Faraj Āqā and the degree of his sincerity, God
willing, may he be enabled at all times to serve the Beloved of all
existence and be assisted in remembering and praising Him. As to the
inquiry concerning the entrusted goods, they arrived in full and intact.
This servant bears witness to this, and behind him, the One who created
him, fashioned him, and caused him to voice His beautiful praise. Since
you mentioned him, and this servant also wished that he should be
mentioned at the most inaccessible, most holy, most exalted Court, these
particulars were presented before His face on the seventh night of the
month of Ṣafar, and these resplendent verses and perfect words descended
from the heaven of the loving-kindness of the Revealer.

\subsection{20}\label{section-2850}

\begin{otherlanguage}{arabic}

قوله جلّت عظمته هو الظّاهر المستور

\end{otherlanguage}

Qawluhu jallat ʿaẓamatuhu: Huwa al-Ẓāhir al-Mastūr.

His words---exalted be His grandeur: He is the Manifest, the Hidden.

\subsection{21}\label{section-2851}

\begin{otherlanguage}{arabic}

یا فرج قد ذکرناک من قبل بذکر لا تعادله کتب العالم انّ ربّک لهو المحصی
العلیم ان افرح بما ذکرت مرّة بعد مرّة و توجّه الیک وجه المظلوم من هذا الحصن
المتین لو تطّلع بما قدّر لک من قلمی الأعلی لتطیر بأجنحة الفرح و الابتهاج
فی هذا الهوآء الّذی یسمع من هزیز اریاحه ما سمعه طلعة الظّهور فی الطّور انّ
ربّک لهو الصّادق الأمین طوبی لک بما اقبلت و سمعت نداء اللّه اذ ارتفع بین
السّموات و الأرضین لعمری سیفنی ما تراه العیون و تدرکه العقول و یبقی لک ما
نزل من هذا القلم المبین قل لک الحمد یا من ایّدتنی علی الاقبال الی افقک و
اسمعتنی ندائک و زیّنتنی بطراز الیقین اذ کان علمآء العصر فی وهم عظیم کن
راضیاً بما قدّر لک و مستقیماً علی هذا الأمر الّذی اذا ظهر زلّت اقدام البالغین

\end{otherlanguage}

Yā Faraj qad dhakarnāka min qablu bi-dhikrin lā tuʿādiluhu kutubu
al-ʿālam. Inna rabbaka la-huwa al-Muḥṣī al-ʿAlīm. An ifraḥ bi-mā
dhukirta marratan baʿda marra wa-tawajjaha ilayka wajhu al-Maẓlūm min
hādhā al-ḥiṣn al-matīn. Law taṭṭaliʿ bi-mā quddira laka min qalamī
al-Aʿlā la-taṭīru bi-ajniḥati al-faraḥ wa-l-ibtihāj fī hādhā al-hawāʾ
alladhī yusmaʿu min hazīzi aryāḥihi mā samiʿahu Ṭalʿatu al-Ẓuhūr fī
al-Ṭūr. Inna rabbaka la-huwa al-Ṣādiq al-Amīn. Ṭūbā laka bi-mā aqbalta
wa-samiʿta nidāʾa Allāh idh irtafaʿa bayna al-samāwāt wa-l-araḍīn.
La-ʿamrī sa-yafnā mā tarāhu al-ʿuyūn wa-tudrikuhu al-ʿuqūl wa-yabqā laka
mā nazala min hādhā al-qalam al-mubīn. Qul laka al-ḥamdu yā man
ayyadtanī ʿalā al-iqbāl ilā ufuqika wa-asmaʿtanī nidāʾaka wa-zayyantanī
bi-ṭirāzi al-yaqīn idh kāna ʿulamāʾu al-ʿaṣr fī wahmin ʿaẓīm. Kun
rāḍiyan bi-mā quddira laka wa-mustaqīman ʿalā hādhā al-amr alladhī idhā
ẓahara zallat aqdāmu al-bālighīn.

O Faraj, We have remembered you before with a remembrance that the books
of the world cannot equal. Your Lord is indeed the Reckoner, the
All-Knowing. Rejoice that you have been remembered time after time and
that the face of the Wronged One has turned toward you from this mighty
fortress. Were you to become aware of what has been ordained for you by
My Most Exalted Pen, you would soar on the wings of joy and gladness in
this air, from the stirring of whose winds can be heard what the
Countenance of Revelation heard on Sinai. Your Lord is indeed the
Truthful, the Trustworthy. Blessed are you for having turned toward Him
and heard the call of God as it rose between the heavens and the earths.
By My life, whatever eyes see and minds comprehend will pass away, while
what has descended from this perspicuous Pen will remain for you. Say:
``Praise be to You, who have enabled me to turn toward Your horizon,
caused me to hear Your call, and adorned me with the vesture of
certitude when the learned of the age were lost in profound delusion.''
Be content with what has been ordained for you, and remain steadfast in
this Cause, at whose appearance the feet of those of mature
understanding slipped.

\subsection{22}\label{section-2852}

\begin{otherlanguage}{arabic}

و نذکر من سمّی بمحمّد الّذی قطع البرّ و البحر الی ان فاز بالبحر الأعظم فی
هذا البرّ الّذی ارتفع فیه ندآء النّبیّین و المرسلین قد حضر و قام لدی الباب و
سمع ندآء المظلوم و شرب من ید عطائه رحیقه المختوم کذلک ذکرناه فضلاً من
عندنا و انا الفضّال الکریم البهآء علیک و علیه و علی من وجد هذا العرف
العزیز البدیع انتهی

\end{otherlanguage}

Wa-nadhkuru man summiya bi-Muḥammad alladhī qaṭaʿa al-barra wa-l-baḥra
ilā an fāza bi-l-baḥr al-aʿẓam fī hādhā al-barr alladhī irtafaʿa fīhi
nidāʾu al-nabiyyīn wa-l-mursalīn. Qad ḥaḍara wa-qāma ladā al-bāb
wa-samiʿa nidāʾa al-Maẓlūm wa-shariba min yadi ʿaṭāʾihi raḥīqahu
al-makhtūm. Kadhālika dhakarnāhu faḍlan min ʿindinā wa-anā al-Faḍḍāl
al-Karīm. Al-Bahāʾ ʿalayka wa-ʿalayhi wa-ʿalā man wajada hādhā al-ʿarfa
al-ʿazīz al-badīʿ. Intahā.

And We remember the one named Muḥammad, who crossed land and sea until
he attained the Most Great Ocean in this land where the call of the
Prophets and Messengers has been raised. He came, stood at the door,
heard the call of the Wronged One, and drank His sealed wine from the
hand of His bounty. Thus have We remembered him as a grace from Our
presence, and I am the Most Bountiful, the Generous. Glory be upon you,
upon him, and upon whoever has perceived this precious and wondrous
fragrance. End.

\subsection{23}\label{section-2853}

\begin{otherlanguage}{persian}

بلی جناب حاجی آقا محمّد علیه ٩ ٦٦ {[}بهآء اللّه{]} وارد شدند و بشرف لقا
فائز و مفتخر و مرّة بعد مرّة در خارج ارض سجن به ما هو المقصود یعنی بزیارت
محبوب عالمیان مرزوق گشتند

\end{otherlanguage}

Balī Janāb-e Ḥājī Āqā Moḥammad ʿalayhi 9 66 {[}Bahāʾ Allāh{]} vāred
shodand va be-sharaf-e leqā fāʾez va moftakher va marratan baʿda marra
dar khārej-e arż-e sejn be-mā huwa al-maqṣūd yaʿnī be-ziyārat-e Maḥbūb-e
ʿālamiyān marzūq gashtand.

Yes, Ḥājī Āqā Muḥammad---upon him be 9 66 {[}the glory of
God{]}---arrived and attained the honor and distinction of His presence.
Time after time, outside the land of the Prison, he was granted what was
desired, namely, the privilege of visiting the Beloved of all the
worlds.

\subsection{24}\label{section-2854}

\begin{otherlanguage}{arabic}

و اینکه در بارهٴ جناب حاجی علی اصغر و شاگرد ایشان علیهما بهآء اللّه مرقوم
داشتند بعد از عرض آفتاب عنایت مشرق و لوح اقدس اعظم نازل انشآءاللّه
بشرفیابی آن مرزوق گردند قوله عزّ کبریائه

\end{otherlanguage}

Va īn-keh dar bāreh-ye Janāb-e Ḥājī ʿAlī-Aṣghar va shāgerd-e īshān
ʿalayhimā Bahāʾ Allāh marqūm dāshtand, baʿd az ʿarż āftāb-e ʿenāyat
moshreq va lawḥ-e aqdas-e aʿẓam nāzel. In shāʾ Allāh be-sharafyābī-ye ān
marzūq gardand. Qawluhu ʿazza kibriyāʾuhu.

As to what was written concerning Ḥājī ʿAlī-Aṣghar and his
apprentice---upon them both be the glory of God---after it was
presented, the sun of loving-kindness shone forth and a most holy, most
great Tablet descended. God willing, may they be granted the honor of
receiving it. His words---exalted be His majesty:

\subsection{25}\label{section-2855}

\begin{otherlanguage}{arabic}

بسمی السّامع البصیر

\end{otherlanguage}

Bi-smī al-Sāmiʿ al-Baṣīr.

In My name, the All-Hearing, the All-Seeing.

\subsection{26}\label{section-2856}

\begin{otherlanguage}{persian}

یا علی قبل اصغر مالک قدم ظاهر و امم از او محجوب و آفتاب حقیقت مشرق ولکن
ابصار از او محروم فرات رحمت جاری و ناس از آن ممنوع قلم اعلی متحرّک و ناطق
و کل از صریرش غافل الّا من شآء اللّه هر نفسی الیوم بعرفان مالک وجود فائز
شد او از نفوسی است که ذکرشان از قلم الهی در کتب قبل ثبت شد شکر کن محبوب
عالمیان را که تو را مؤیّد فرمود بر قبول فضل و عطایش جمیع اهل دنیا از علما
و عرفا و امرا بجان مشتاق این یوم اعزّ انور بودند و چون آفتاب ظهور اشراق
نمود و ندا مرتفع شد باقبال و اصغا فائز نگشتند و علّت اعظم که امم را منع
نمود علمای عصر بوده و هستند رسول از ظلم آن نفوس غافله بنوحه مشغول و آن
همج رعاع باوهام خود معتکف و مسرور غفلت این نفوس سبب غفلت اهل عالم شده و
ظلمشان علّت ظلم امم ولکن اصبع قدرت احجاب را خرق فرماید چنانچه نمود

\end{otherlanguage}

Yā ʿAlī qabl-e Aṣghar, Mālek-e qedam ẓāher va omam az ū maḥjūb va
āftāb-e ḥaqīqat moshreq va-līkan abṣār az ū maḥrūm. Forāt-e raḥmat jārī
va nās az ān mamnūʿ. Qalam-e aʿlā motaḥarrek va nāṭeq va koll az
ṣarīrash ghāfel illā man shāʾa Allāh. Har nafsī al-yawm be-ʿerfān-e
Mālek-e vojūd fāʾez shod ū az nofūsī ast keh dhekr-shān az qalam-e elāhī
dar kotob-e qabl thabt shod. Shokr kon Maḥbūb-e ʿālamiyān rā keh to rā
moʾayyad farmūd bar qabūl-e fażl va ʿaṭāyash. Jamīʿ-e ahl-e donyā az
ʿolamā va ʿorafā va omarā be-jān moshtāq-e īn yawm-e aʿazz-e anvar
būdand va chūn āftāb-e ẓohūr eshrāq namūd va nedā mortafeʿ shod be-eqbāl
va eṣghā fāʾez nagashtand. Va ʿellat-e aʿẓam keh omam rā manʿ namūd
ʿolamā-ye ʿaṣr būdeh va hastand. Rasūl az ẓolm-e ān nofūs-e ghāfeleh
be-nowḥeh mashghūl va ān hamaj-e raʿāʿ be-awhām-e khod moʿtakef va
masrūr. Ghaflat-e īn nofūs sabab-e ghaflat-e ahl-e ʿālam shodeh va
ẓolm-shān ʿellat-e ẓolm-e omam va-līkan eṣbaʿ-e qodrat aḥjāb rā kharq
farmāyad chenāncheh namūd.

O ʿAlī-Aṣghar, the Lord of eternity is manifest, yet the nations are
veiled from Him; the Sun of Truth is shining, yet eyes are deprived of
Him. The Euphrates of mercy is flowing, yet people are kept from it. The
Most Exalted Pen moves and speaks, yet all are heedless of its voice
except those whom God wills. Whoever has attained recognition of the
Lord of existence today is among those whose mention was inscribed by
the divine Pen in the Books of old. Give thanks to the Beloved of all
the worlds, who has assisted you to accept His grace and bounty. All the
people of the world---scholars, mystics, and rulers---longed with all
their souls for this most glorious, most luminous Day. Yet when the Sun
of Revelation shone and the call was raised, they failed to turn toward
Him and listen. The chief cause that has held the nations back has been,
and remains, the learned of the age. The Messenger laments the
oppression of those heedless souls, while that ignorant rabble clings to
its own delusions and rejoices. The heedlessness of these souls has
caused the heedlessness of the people of the world, and their oppression
has caused the oppression of the nations. Yet the finger of power will
rend the veils, even as it has already done.

\subsection{27}\label{section-2857}

\begin{otherlanguage}{persian}

و نذکر محمّد قبل حسین لیفرح بذکر اللّه المهیمن القیّوم انشآءاللّه بطراز
امانت مزیّن باشی و باخلاق الهی ظاهر صفات مرضیّه بمنزلهٴ انوار است از برای
افق سماء انسان هر نفسی بآن فائز شد نورش بتابد و نسبتش بحقّ جلّ جلاله ظاهر
هویدا شود انشآءاللّه کلّ فائز شوند بآنچه سبب ارتفاع امر است انّه یحکم بما
یرتفع به شأن الانسان فی الامکان یشهد بذلک کلّ منصف بصیر انتهی

\end{otherlanguage}

Wa-nadhkuru Muḥammad qabl Ḥusayn li-yafraḥa bi-dhikri Allāh al-Muhaymin
al-Qayyūm. In shāʾ Allāh be-ṭerāz-e amānat mozayyan bāshī va be-akhlāq-e
elāhī ẓāher. Ṣefāt-e marżiyyeh be-manzaleh-ye anvār ast az barā-ye
ofoq-e samāʾ-e ensān. Har nafsī be-ān fāʾez shod nūrash betābad va
nesbatash be-Ḥaqq jalla jalāluhu ẓāher-e hoveydā shavad. In shāʾ Allāh
koll fāʾez shavand be-āncheh sabab-e ertefāʿ-e amr ast. Innahu yaḥkumu
bi-mā yartafiʿu bihi shaʾnu al-insān fī al-imkān yashhadu bi-dhālika
kullu munṣifin baṣīr. Intahā.

And We remember Muḥammad-Ḥusayn, that he may rejoice in the remembrance
of God, the All-Protecting, the Self-Subsisting. God willing, may you be
adorned with the vesture of trustworthiness and manifest divine
character. Praiseworthy attributes are like lights for the horizon of
the heaven of humanity. Whoever attains them will shine with their
light, and his relationship to God---exalted be His glory---will become
clear and manifest. God willing, may all attain what will exalt the
Cause. He ordains that which will elevate the station of the human being
in the world of existence. To this every fair-minded and discerning
person bears witness. End.

\subsection{28}\label{section-2858}

\begin{otherlanguage}{persian}

الحمد للّه مائدهٴ طریّهٴ الهیّه و نعمت باقیهٴ دائمه از سماء احدیّه من غیر
تعطیل جاری و نازل هر نفسی فی‌الجمله تفکّر نماید از دریای شفقت حقّ جلّ جلاله
آگاه شود در لیالی و ایّام لسان عظمت بذکر دوستانش مشغول و بنصح اعدایش
متکلّم عموم اهل عالم را بکمال شفقت و مرحمت بافق اعلی دعوت فرموده‌اند ولکن
طوبی لمن سمع و اقبل و ویل للغافلین و الحمد للّه ربّ العالمین

\end{otherlanguage}

Al-ḥamdu li-llāh māʾedeh-ye ṭariyyeh-ye elāhiyyeh va neʿmat-e bāqiyeh-ye
dāʾemeh az samāʾ-e aḥadiyyeh min ghayr-e taʿṭīl jārī va nāzel. Har nafsī
fī-l-jomleh tafakkor namāyad az daryā-ye shafaqat-e Ḥaqq jalla jalāluhu
āgāh shavad. Dar layālī va ayyām lesān-e ʿaẓamat be-dhekr-e dūstānash
mashghūl va be-noṣḥ-e aʿdāyash motakallem. ʿOmūm-e ahl-e ʿālam rā
be-kamāl-e shafaqat va marḥamat be-ofoq-e aʿlā daʿvat farmūdeh-and
va-līkan ṭūbā li-man samiʿa wa-aqbala wa-waylun li-l-ghāfilīn wa-l-ḥamdu
li-llāhi rabbi al-ʿālamīn.

Praise be to God, the fresh divine banquet and the enduring, everlasting
bounty flow and descend without interruption from the heaven of oneness.
Anyone who reflects even a little will become aware of the ocean of
God's compassion---exalted be His glory. By night and day, the Tongue of
Grandeur is occupied with remembering His friends and speaking counsel
to His enemies. With the utmost compassion and mercy, He has summoned
all the people of the world to the Most Exalted Horizon. Yet blessed is
the one who has heard and turned toward Him, and woe to the heedless.
Praise be to God, the Lord of all the worlds.

\subsection{29}\label{section-2859}

\begin{otherlanguage}{persian}

و اینکه در بارهٴ ابن جناب من صعد الی اللّه الّذی سمّی برضا الرّوح فی کتاب
الأسمآء ذکر فرمودید عرض شد و این لوح امنع اقدس مخصوص ایشان از سماء مشیّت
محبوب امکان نازل قوله عزّ کبریائه

\end{otherlanguage}

Va īn-keh dar bāreh-ye ibn-e Janāb-e man ṣaʿida ilā Allāh alladhī
summiya bi-Riḍā al-Rūḥ fī Kitāb al-Asmāʾ dhekr farmūdīd, ʿarż shod va īn
lawḥ-e amnaʿ-e aqdas makhṣūṣ-e īshān az samāʾ-e mashiyyat-e Maḥbūb-e
emkān nāzel. Qawluhu ʿazza kibriyāʾuhu.

What you mentioned concerning the son of the honored one who has
ascended to God, who was named Riḍā al-Rūḥ in the Book of Names, was
presented, and this most inaccessible, most holy Tablet descended
especially for him from the heaven of the will of the Beloved of all
existence. His words---exalted be His majesty:

\subsection{30}\label{section-2860}

\begin{otherlanguage}{arabic}

بسمی المشرق من افق البیان

\end{otherlanguage}

Bi-smī al-Mushriq min ufuqi al-bayān.

In My name, shining from the horizon of utterance.

\subsection{31}\label{section-2861}

\begin{otherlanguage}{arabic}

ذکر انزله الرّحمن لمن اقبل الی الوجه فی یوم فیه تشتّت شمل العالم و ارتفع
العویل بین السّموات و الأرضین لیفرح بذکری الّذی جعله اللّه رحیق الحیوان
لأهل الامکان و روح الأعظم بین الأمم و العطیّة الکبری للوری تعالی مالک
الأدیان الّذی اتی من سمآء الفضل بعدل مبین

\end{otherlanguage}

Dhikrun anzalahu al-Raḥmān li-man aqbala ilā al-wajh fī yawmin fīhi
tashtattata shamlu al-ʿālam wa-irtafaʿa al-ʿawīl bayna al-samāwāt
wa-l-araḍīn li-yafraḥa bi-dhikrī alladhī jaʿalahu Allāh raḥīqa
al-ḥayawān li-ahli al-imkān wa-rūḥa al-aʿẓam bayna al-umam
wa-l-ʿaṭiyyata al-kubrā li-l-warā. Taʿālā Mālik al-adyān alladhī atā min
samāʾi al-faḍl bi-ʿadlin mubīn.

A remembrance sent down by the All-Merciful for the one who has turned
toward the Face on a day when the unity of the world has been shattered
and lamentation has risen between the heavens and the earths, that he
may rejoice in My remembrance, which God has made the choice wine of
life for the people of existence, the Most Great Spirit among the
nations, and the supreme gift to humanity. Exalted is the Lord of
religions, who has come from the heaven of grace with manifest justice.

\subsection{32}\label{section-2862}

\begin{otherlanguage}{arabic}

یا هادی قد ارسلنا القائم انکره علمآء العصر کلّها الّا من شآء اللّه ربّ
العالمین فلمّا مضت سنین متوالیات اظهرنا نفسنا باسمنا القیّوم ولکنّ القوم
اعرضوا الی ان افتوا علینا بظلم ناح به الرّوح الأمین قد اخذونا مرّة بعد مرّة
و حبسونا کرّة بعد کرّة تارة فی ارض الطّآء و اخری فی المیم و مرّة اخری فی
الطّآء و فی ارض السّرّ و هذا السّجن العظیم انّا نشکر اللّه فی کلّ ذلک و قبلنا
فی سبیله ما ورد علینا من الّذین غفلوا عن اللّه المقتدر القدیر

\end{otherlanguage}

Yā Hādī, qad arsalnā al-Qāʾim, ankarahu ʿulamāʾ al-ʿaṣr kulluhā illā man
shāʾa Allāh rabb al-ʿālamīn. Fa-lammā maḍat sinīn mutawāliyāt aẓharnā
nafsanā bi-sminā al-Qayyūm, wa-lākinna al-qawm aʿraḍū ilā an aftaw
ʿalaynā bi-ẓulm nāḥa bihi al-Rūḥ al-Amīn. Qad akhadhūnā marratan baʿda
marra wa-ḥabasūnā karratan baʿda karra, tāratan fī arḍ al-Ṭāʾ wa-ukhrā
fī al-Mīm wa-marratan ukhrā fī al-Ṭāʾ wa-fī arḍ al-Sirr wa-hādhā al-sijn
al-ʿaẓīm. Innā nashkur Allāh fī kull dhālik wa-qabilnā fī sabīlihi mā
warada ʿalaynā min alladhīna ghafalū ʿan Allāh al-Muqtadir al-Qadīr.

O Hādī! We sent forth the Qāʾim, but all the scholars of the age denied
Him, except those whom God, the Lord of the worlds, willed otherwise.
When successive years had passed, We manifested Our own Self in Our
name, the Self-Subsisting. Yet the people turned away, until they issued
rulings against Us with such injustice that the Faithful Spirit lamented
over it. They seized Us time after time and imprisoned Us again and
again: once in the land of Ṭāʾ, another time in Mīm, once again in Ṭāʾ,
and in the Land of Mystery and this mighty prison. Through all this We
thank God, and in His path We have accepted whatever befell Us at the
hands of those who were heedless of God, the All-Powerful, the
Omnipotent.

\subsection{33}\label{section-2863}

\begin{otherlanguage}{arabic}

انّا نذکر اباک الّذی به ابتسم ثغر العرفان فی الامکان و فاز بلقآء اللّه
العلیم الحکیم انّا ذکرناه بما لا یعادله شیء فی الأرض و نذکره بدوام الملک
و المکوت و انا المشفق الکریم

\end{otherlanguage}

Innā nadhkur abāka alladhī bihi ibtasama thaghr al-ʿirfān fī al-imkān
wa-fāza bi-liqāʾ Allāh al-ʿAlīm al-Ḥakīm. Innā dhakarnāhu bi-mā lā
yuʿādiluhu shayʾ fī al-arḍ wa-nadhkuruhu bi-dawām al-mulk
wa-{[}corrupted text{]} wa-anā al-Mushfiq al-Karīm.

We remember your father, through whom the mouth of recognition smiled in
the world of existence, and who attained the presence of God, the
All-Knowing, the All-Wise. We have remembered him with that which
nothing on earth can equal, and We will remember him as long as the
dominion and {[}corrupted original text{]} endure. I am the
Compassionate, the All-Bountiful.

\subsection{34}\label{section-2864}

\begin{otherlanguage}{arabic}

یا رضا ان افرح بما یذکرک مالک الأسمآء من شطر السّجن بذکر ینجذب منه الملأ
الأعلی و سکّان الفردوس و الّذین طافوا العرش فی البکور و الأصیل طوبی لمن
اقبل الیک و زارک بما نزل من قلمی الأعلی و حضر تلقآء قبرک الشّریف انت الّذی
اقبلت الی افق اللّه و توجّهت الیه فی یوم فیه اعرض عنه کلّ غافل بعید

\end{otherlanguage}

Yā Riḍā, an ifraḥ bi-mā yadhkuruka Mālik al-Asmāʾ min shaṭr al-sijn
bi-dhikr yanjadhibu minhu al-malaʾ al-aʿlā wa-sukkān al-firdaws
wa-alladhīna ṭāfū al-ʿarsh fī al-bukūr wa-al-aṣīl. Ṭūbā li-man aqbala
ilayka wa-zāraka bi-mā nazala min qalamī al-aʿlā wa-ḥaḍara tilqāʾ
qabrika al-sharīf. Anta alladhī aqbalta ilā ufuq Allāh wa-tawajjahta
ilayhi fī yawm fīhi aʿraḍa ʿanhu kull ghāfil baʿīd.

O Riḍā! Rejoice that the Lord of Names remembers you from the direction
of the prison with a remembrance that enraptures the heavenly Concourse,
the inhabitants of Paradise, and those who circle the Throne at morning
and evening. Blessed is whoever turns toward you, visits you with what
has been sent down from My Most Exalted Pen, and stands before your
noble grave. You are the one who turned toward the horizon of God and
directed yourself toward Him on a day when everyone heedless and far
removed turned away from Him.

\subsection{35}\label{section-2865}

\begin{otherlanguage}{arabic}

اوّل ذکر نطق به لسان مالک القدم بین الأمم علیک یا ایّها النّاطق بذکر اللّه و
المتوجّه الی الوجه بعد فنآء الأشیآء اشهد انّک شربت الرّحیق المختوم من ید
اسمه القیّوم و نبذت الأوهام عن ورائک و اخذت الایقان بقوّة من لدی اللّه
القویّ الغالب العزیز الحمید طوبی لک بما تزیّن رأسک باکلیل القبول و هیکلک
بطراز عنایة اللّه فی یومه البدیع البهآء و الذّکر و الرّوح و السّلام و
النّفحات المتضوّعة من قمیص رحمة ربّک علیک یا من فزت بعنایة ربّک اذ اتی
بسلطان مبین انتهی

\end{otherlanguage}

Awwal dhikr naṭaqa bihi lisān Mālik al-Qidam bayna al-umam ʿalayka yā
ayyuhā al-nāṭiq bi-dhikr Allāh wa-al-mutawajjih ilā al-wajh baʿda fanāʾ
al-ashyāʾ. Ashhadu annaka sharibta al-raḥīq al-makhtūm min yad ismihi
al-Qayyūm wa-nabadhta al-awhām ʿan warāʾika wa-akhadhta al-īqān bi-quwwa
min ladā Allāh al-Qawī al-Ghālib al-ʿAzīz al-Ḥamīd. Ṭūbā laka bi-mā
tazayyana raʾsuka bi-iklīl al-qabūl wa-haykaluka bi-ṭirāz ʿināyat Allāh
fī yawmihi al-badīʿ. Al-bahāʾ wa-al-dhikr wa-al-rūḥ wa-al-salām
wa-al-nafaḥāt al-mutaḍawwiʿa min qamīṣ raḥmat rabbika ʿalayka yā man
fuzta bi-ʿināyat rabbika idh atā bi-sulṭān mubīn. Intahā.

May the first remembrance uttered by the tongue of the Lord of eternity
among the nations rest upon you, O you who speak in remembrance of God
and turn toward the Face after all things have passed away. I bear
witness that you drank the sealed wine from the hand of His name, the
Self-Subsisting, cast vain imaginings behind you, and held fast to
certitude through a strength from God, the Strong, the All-Subduing, the
Almighty, the Praiseworthy. Blessed are you, for your head has been
adorned with the crown of acceptance and your person with the ornament
of God's loving care in His wondrous Day. Glory, remembrance, spirit,
peace, and the fragrant breezes wafting from the garment of your Lord's
mercy rest upon you, O you who attained your Lord's loving care when He
came with manifest authority. End.

\subsection{36}\label{section-2866}

\begin{otherlanguage}{persian}

اینکه ذکر اهل سیسان و نفوس مشتعله نمودید انشآءاللّه کلّ بعنایت حقّ فائز
شوند و این دعا این حین باستجابت فائز گشت یک لوح امنع اقدس مخصوص آن نفوس
مقبلهٴ موقنه نازل و از قبل هم بالواح الهی فائز گشته‌اند قوله تعالی شأنه و
جلّ کبریائه

\end{otherlanguage}

Īnke dhekr-e ahl-e Sīsān va nofūs-e moshtaʿele namūdīd, in shāʾ Allāh
koll be-ʿenāyat-e Ḥaqq fāʾez shavand, va īn doʿā īn ḥīn be-estejābat
fāʾez gasht. Yek lawḥ-e amnaʿ-e aqdas makhṣūṣ-e ān nofūs-e moqbele-ye
mūqene nāzel va az qabl ham be-alvāḥ-e elāhī fāʾez gashte-and. Qawluhu
taʿālā shaʾnuhu wa-jalla kibriyāʾuhu.

You mentioned the people of Sīsān and those enkindled souls. God
willing, may they all attain the loving care of the True One. This
prayer has at this very moment been answered. A most exalted, most holy
Tablet has been sent down especially for those souls who have turned
toward Him with certitude, and they have previously also received divine
Tablets. His words---exalted is His station and glorious His majesty:

\subsection{37}\label{section-2867}

\begin{otherlanguage}{arabic}

بنام بینای یکتا

\end{otherlanguage}

Be-nām-e Bīnā-ye Yektā.

In the name of the One, the All-Seeing.

\subsection{38}\label{section-2868}

\begin{otherlanguage}{persian}

اولیای سیسان را بعنایت رحمن بشارت میدهم جمیع تحت لحاظ مرحمت بوده و هستند
و در ظلّ سدرهٴ باسقهٴ ربّانیّه مقرّ گزیده‌اند در فضل الهی تفکّر نمائید علمای
دنیا که خود را اعلم و اعرف و افقه و اعلای اهل عالم میشمردند بقطره‌ئی از
بحر بیان فائز نشدند و بذکری لدی الرّحمن مفتخر نگشتند ولکن آن نفوس از
عنایت حقّ جلّ جلاله بمقامی فائز گشتند که قلم اعلی در سجن اعظم بذکرشان
مشغول تعالی هذا المقام الأسمی الأسنی و تعالی هذا الذّکر الأعظم الأعلی که
تواند از عهدهٴ شکر الهی برآید قسم بآفتاب حقیقت که از افق سجن طالع است
اگر نفسی بعدد ذرّات دارای لسان باشد و از اوّل لا اوّل له الی آخر لا آخر له
شکر و ثنا گوید هرآینه بیک یا عبدی که از مطلع بیان ظاهر شود معادله ننماید
و حال لسان رحمن به یا احبّائی فی سیسان ناطق است انشآءاللّه کلّ اینمقام را
عظیم شمرند و جهد نمایند تا بعظمتش آگاه شوند و اگر فی‌الحقیقه باینمقام
فائز گردند خود را از عالم و عالمیان فارغ مشاهده نمایند

\end{otherlanguage}

Awliyā-ye Sīsān rā be-ʿenāyat-e Raḥmān beshārat mīdaham. Jamīʿ taḥt-e
leḥāẓ-e marḥamat būde va hastand va dar ẓell-e sedre-ye bāseqe-ye
rabbānīye maqar gozīde-and. Dar faḍl-e elāhī tafakkor namāyīd. ʿOlamā-ye
donyā ke khod rā aʿlam va aʿraf va afqah va aʿlā-ye ahl-e ʿālam
mīshomordand be-qaṭreʾī az baḥr-e bayān fāʾez nashodand va be-dhekrī
ladā al-Raḥmān moftakher nagashtand, va-lāken ān nofūs az ʿenāyat-e Ḥaqq
jalla jalāluhu be-maqāmī fāʾez gashtand ke Qalam-e Aʿlā dar Sejn-e Aʿẓam
be-dhekrashān mashghūl. Taʿālā hādhā al-maqām al-asmā al-asnā wa-taʿālā
hādhā al-dhikr al-aʿẓam al-aʿlā. Ke tavānad az ʿohde-ye shokr-e elāhī
barāyad? Qasam be-Āftāb-e Ḥaqīqat ke az ofoq-e sejn ṭāleʿ ast, agar
nafsī be-ʿadad-e dharrāt dārā-ye lesān bāshad va az avval-e lā awwala
lahu ilā ākhar-e lā ākhira lahu shokr va thanā gūyad, har-āyene be-yek
yā ʿabdī ke az maṭlaʿ-e bayān ẓāher shavad moʿādele nanamāyad. Va ḥāl
Lesān-e Raḥmān be-yā aḥibbāʾī fī Sīsān nāṭeq ast. In shāʾ Allāh koll īn
maqām rā ʿaẓīm shomorand va jahd namāyand tā be-ʿaẓamatash āgāh shavand,
va agar fī al-ḥaqīqe be-īn maqām fāʾez gardand khod rā az ʿālam va
ʿālamiyān fāregh moshāhede namāyand.

I give the friends in Sīsān the glad tidings of the loving care of the
All-Merciful. All have been, and remain, under His merciful gaze, and
have made their home beneath the shade of the lofty, divine Lote Tree.
Reflect on the grace of God. The scholars of the world, who considered
themselves the most learned, the most discerning, the most versed in
religious law, and the most exalted of the world's people, did not
attain a single drop from the ocean of utterance, nor were they honored
with a mention in the presence of the All-Merciful. Yet through the
loving care of the True One---glorious is His glory---those souls have
attained a station such that the Most Exalted Pen is occupied with their
remembrance in the Most Great Prison. Exalted is this loftiest, most
radiant station, and exalted is this greatest, most sublime remembrance!
Who can adequately thank God? I swear by the Sun of Truth shining from
the horizon of the prison: if a soul possessed as many tongues as there
are atoms and offered thanks and praise from the beginning that has no
beginning to the end that has no end, it would assuredly not equal a
single ``O My servant'' issuing from the Dawning-Place of utterance. And
now the tongue of the All-Merciful is proclaiming, ``O My loved ones in
Sīsān!'' God willing, may all regard this station as great and strive to
become aware of its greatness. If they truly attain this station, they
will see themselves freed from the world and all its peoples.

\subsection{39}\label{section-2869}

\begin{otherlanguage}{persian}

ایدوستان محبوب امکان شما را ذکر مینماید و مخصوص هر یک در ملکوت فضل مقامی
معیّن فرموده ان افرحوا بما یبشّرکم اللّه به انّه لهو الغفور الرّحیم قد
ذکرناکم من قبل بآیات ما اطّلع بها الّا اللّه المحصی العلیم انّا نأمرکم
بالحکمة و نهدیکم الی صراطه المستقیم تمسّکوا بحبل الحکمة و داروا مع العباد
لئلّا یرتفع ضوضآء الّذین غفلوا عن امری العظیم البهآء علیکم و علی امائی فی
هناک اللّائی سمعن و اجبن ما ارتفع من افقی المنیر انتهی

\end{otherlanguage}

Ey dūstān, Maḥbūb-e emkān shomā rā dhekr mīnamāyad va makhṣūṣ-e har yek
dar malakūt-e faḍl maqāmī moʿayyan farmūde. An ifraḥū bi-mā
yubashshirukum Allāh bihi, innahu la-huwa al-Ghafūr al-Raḥīm. Qad
dhakarnākum min qabl bi-āyāt mā iṭṭalaʿa bihā illā Allāh al-Muḥṣī
al-ʿAlīm. Innā naʾmurukum bi-al-ḥikma wa-nahdīkum ilā ṣirāṭihi
al-mustaqīm. Tamassakū bi-ḥabl al-ḥikma wa-dārū maʿa al-ʿibād li-allā
yartafiʿa ḍawḍāʾ alladhīna ghafalū ʿan amrī al-ʿaẓīm. Al-bahāʾ ʿalaykum
wa-ʿalā imāʾī fī hunāka allāʾī samiʿna wa-ajabna mā irtafaʿa min ufuqī
al-munīr. Intahā.

O friends! The Beloved of all existence remembers you and has assigned a
particular station to each of you in the kingdom of grace. Rejoice in
the glad tidings God gives you. He is indeed the Ever-Forgiving, the
Merciful. We have remembered you before in verses of which none but God,
the All-Reckoning, the All-Knowing, is aware. We enjoin wisdom upon you
and guide you to His straight path. Hold fast to the cord of wisdom and
deal tactfully with people, so that no clamor may arise from those who
are heedless of My mighty Cause. Glory rest upon you and upon My
handmaidens there who have heard and answered what has risen from My
luminous horizon. End.

\subsection{40}\label{section-2870}

\begin{otherlanguage}{persian}

مخصوص دوستان الهی در سیسان از قبل نازل شده آنچه عرفش از عالم منقطع نشود
الحمد للّه غمام فضل بشأنی مرتفع است که هر ذرّه از ذرّات را احاطه نموده تا
چه رسد بنفوس مقبلهٴ مطمئنّه که از کأس ایقان نوشیده‌اند و بمقام اطمینان
فائز گشته‌اند لهم مقام عند ربّهم العلیم الحکیم اینخادم بکمال عجز و انکسار
خدمت هر یک از ایشان تکبیر میرساند اظهاراً لخلوصی لهم و اعلاناً لحبّی ایّاهم

\end{otherlanguage}

Makhṣūṣ-e dūstān-e elāhī dar Sīsān az qabl nāzel shode ānche ʿarfash az
ʿālam monqaṭeʿ nashavad. Al-ḥamdu li-llāh ghamām-e faḍl be-shaʾnī
mortafeʿ ast ke har dharre az dharrāt rā eḥāṭe namūde, tā che resad
be-nofūs-e moqbele-ye moṭmaʾenne ke az kaʾs-e īqān nūshīde-and va
be-maqām-e eṭmīnān fāʾez gashte-and. Lahum maqām ʿinda rabbihim al-ʿAlīm
al-Ḥakīm. Īn khādem be-kamāl-e ʿajz va enkesār khedmat-e har yek az
īshān takbīr mīresānad, iẓhāran li-khulūṣī lahum wa-iʿlānan li-ḥubbī
iyyāhum.

There has previously been sent down especially for the friends of God in
Sīsān that whose fragrance will never cease to pervade the world. Praise
be to God, the cloud of grace has risen in such a way that it has
encompassed every atom---how much more those assured souls who have
turned toward Him, drunk from the cup of certitude, and attained the
station of assurance. They have a station with their Lord, the
All-Knowing, the All-Wise. With utter helplessness and humility, this
servant conveys to each of them the greeting ``God is most great,'' as
an expression of my devotion to them and a declaration of my love for
them.

\subsection{41}\label{section-2871}

\begin{otherlanguage}{persian}

اینکه در ذکر احبّای الهی و اجتماع ایشان در بیت جناب آقا محمّد حسن علیهم و
علیه بهآء اللّه مرقوم داشتند این مراتب در پیشگاه حضور محبوب عالمیان عرض
شد این لوح امنع اقدس مخصوص آن نفوس و صاحب بیت از افق ارادهٴ مالک احدیّه
ظاهر و مشرق انشآءاللّه از انوارش منوّر شوند و از دریای معانیش بیآشامند انّه
لهو الفضّال الکریم قوله عزّ کبریائه

\end{otherlanguage}

Īnke dar dhekr-e aḥebbā-ye elāhī va ejtemāʿ-e īshān dar bayt-e Janāb-e
Āqā Moḥammad Ḥasan, ʿalayhim wa-ʿalayhi Bahāʾ Allāh, marqūm dāshtand, īn
marāteb dar pīshgāh-e ḥoḍūr-e Maḥbūb-e ʿālamiyān ʿarḍ shod. Īn lawḥ-e
amnaʿ-e aqdas makhṣūṣ-e ān nofūs va ṣāḥeb-e bayt az ofoq-e erāde-ye
Mālek-e Aḥadīye ẓāher va moshreq. In shāʾ Allāh az anvārash monavvar
shavand va az daryā-ye maʿānīsh biyāshāmand. Innahu la-huwa al-Faḍḍāl
al-Karīm. Qawluhu ʿazza kibriyāʾuhu.

What you wrote concerning the friends of God and their gathering in the
home of Āqā Muḥammad Ḥasan---may the glory of God rest upon them and
upon him---was presented in the presence of the Beloved of all peoples.
This most exalted, most holy Tablet, intended especially for those souls
and the owner of the house, has appeared and shone forth from the
horizon of the will of the Lord of Oneness. God willing, may they be
illumined by its lights and drink from the ocean of its meanings. He is
indeed the All-Bountiful, the Generous. His words---exalted is His
majesty:

\subsection{42}\label{section-2872}

\begin{otherlanguage}{arabic}

بسمی المبیّن العلیم

\end{otherlanguage}

Bi-smī al-Mubayyin al-ʿAlīm.

In My name, the All-Expounding, the All-Knowing.

\subsection{43}\label{section-2873}

\begin{otherlanguage}{arabic}

لسان العظمة یبشّر الّذین اجتمعوا لذکر اللّه و شربوا باسمه کوثر الثّنآء اولئک
عباد مخلصون یصلّین علیهم الملأ الفردوس الأعلی و عن ورائهم ملائکة مقرّبون
طوبی لأذن فازت باصغآء آیاتی و لبیت ارتفع فیه ذکری الممنوع قد خلقنا
العباد لذکری و ثنائی و خدمة امری فلمّا اظهرت نفسی اعرضوا عنّی الّا من شآء
اللّه المهیمن القیّوم یا احبّآء الرّحمن قوموا علی الذّکر و الثّنآء علی شأن لا
تمنعکم الدّنیا کذلک یذکرکم مالک الأسمآء انّه لهو المشفق العزیز الودود ان
افرحوا بما ذکرتم لدی المظلوم و نزّل لکم هذا الکتاب المحتوم البهآء علیکم و
علی الّذین ما خوّفتهم سطوة العلمآء عن اللّه ربّ ما کان و ما یکون انتهی

\end{otherlanguage}

Lisān al-ʿaẓama yubashshir alladhīna ijtamaʿū li-dhikr Allāh wa-sharibū
bi-smihi kawthar al-thanāʾ. Ulāʾika ʿibād mukhliṣūn, yuṣallīna ʿalayhim
al-malaʾ al-firdaws al-aʿlā wa-ʿan warāʾihim malāʾika muqarrabūn. Ṭūbā
li-udhun fāzat bi-iṣghāʾ āyātī wa-li-bayt irtafaʿa fīhi dhikrī
al-mamnūʿ. Qad khalaqnā al-ʿibād li-dhikrī wa-thanāʾī wa-khidmat amrī,
fa-lammā aẓhartu nafsī aʿraḍū ʿannī illā man shāʾa Allāh al-Muhaymin
al-Qayyūm. Yā aḥibbāʾ al-Raḥmān, qūmū ʿalā al-dhikr wa-al-thanāʾ ʿalā
shaʾn lā tamnaʿukum al-dunyā. Kadhālika yadhkurukum Mālik al-Asmāʾ,
innahu la-huwa al-Mushfiq al-ʿAzīz al-Wadūd. An ifraḥū bi-mā dhukirtum
ladā al-Maẓlūm wa-nuzzila lakum hādhā al-kitāb al-maḥtūm. Al-bahāʾ
ʿalaykum wa-ʿalā alladhīna mā khawwafathum saṭwat al-ʿulamāʾ ʿan Allāh
rabb mā kāna wa-mā yakūn. Intahā.

The tongue of grandeur gives glad tidings to those who have gathered to
remember God and have drunk, in His name, from the Kawthar of praise.
These are sincere servants, upon whom the concourse of the highest
Paradise bestows blessings, and beyond them the angels brought near.
Blessed is the ear that has attained to hearing My verses, and the house
in which My forbidden remembrance has been raised. We created humanity
to remember Me, praise Me, and serve My Cause. Yet when I manifested
Myself, they turned away from Me, except those whom God, the
All-Protecting, the Self-Subsisting, willed otherwise. O loved ones of
the All-Merciful! Arise to remembrance and praise in such a way that the
world cannot hinder you. Thus does the Lord of Names remember you. He is
indeed the Compassionate, the Almighty, the All-Loving. Rejoice that you
have been remembered in the presence of the Wronged One and that this
irrevocable Tablet has been sent down for you. Glory rest upon you and
upon those whom the might of the scholars has not frightened away from
God, the Lord of all that has been and all that is to come. End.

\subsection{44}\label{section-2874}

\begin{otherlanguage}{persian}

اسامی آن نفوس مطمئنّه که آنجناب ذکر نمودند و در آن بیت لأجل ذکر اللّه حاضر
شدند عرض شد و هر یک بفضل اللّه و منّه بفیوضات حقیقی فائز گشتند یعطی و یمنع
و هو المقتدر القدیر ثمّ الحمد للّه و المنّة که ذکر جناب ع‌ب علیه بهآء اللّه
شده و میشود لحاظ عنایت بایشان متوجّه در همین ایّام این فانی مکرّر ذکرشانرا
از لسان عظمت استماع نمود انشآءاللّه مؤیّد شوند بحکمت منزله یشهد الخادم
بأنّه فاز بالاقبال الی اللّه و حبّه و خدمته و ذکره و ثنائه و البتّه بمکافات
این اعمال فائز خواهند شد انّ ربّنا الرّحمن لهو الغفور الکریم و هو العطوف
الرّحیم و هو المعطی المشفق السّمیع البصیر

\end{otherlanguage}

Asāmī-ye ān nofūs-e moṭmaʾenne ke ān janāb dhekr namūdand va dar ān bayt
li-ajl dhikr Allāh ḥāḍer shodand ʿarḍ shod, va har yek be-faḍl-e Allāh
va manneh be-foyūḍāt-e ḥaqīqī fāʾez gashtand. Yuʿṭī wa-yamnaʿ wa-huwa
al-Muqtadir al-Qadīr. Thumma al-ḥamdu li-llāh wa-al-minna ke dhekr-e
Janāb-e ʿ. B., ʿalayhi Bahāʾ Allāh, shode va mīshavad. Leḥāẓ-e ʿenāyat
be-īshān motavajjeh. Dar hamīn ayyām īn fānī mokarrar dhekrashān rā az
Lesān-e ʿAẓamat estemāʿ namūd. In shāʾ Allāh moʾayyad shavand
be-ḥekmat-e monzale. Yashhad al-khādim bi-annahu fāza bi-al-iqbāl ilā
Allāh wa-ḥubbihi wa-khidmatihi wa-dhikrihi wa-thanāʾihi. Va al-batte
be-mokāfāt-e īn aʿmāl fāʾez khvāhand shod. Inna rabbanā al-Raḥmān
la-huwa al-Ghafūr al-Karīm wa-huwa al-ʿAṭūf al-Raḥīm wa-huwa al-Muʿṭī
al-Mushfiq al-Samīʿ al-Baṣīr.

The names of those assured souls whom you mentioned, and who attended
that house for the remembrance of God, were presented, and each, through
God's grace and bounty, received true outpourings of grace. He gives and
withholds, and He is the All-Powerful, the Omnipotent. Praise and
gratitude be to God, moreover, that ʿ. B.---may the glory of God rest
upon him---has been and continues to be remembered. The gaze of loving
care is directed toward him. In these very days, this lowly one has
repeatedly heard him mentioned by the Tongue of Grandeur. God willing,
may he be aided through the wisdom that has been sent down. This servant
bears witness that he has attained to turning toward God, loving Him,
serving Him, remembering Him, and praising Him. He will assuredly
receive the recompense for these deeds. Our Lord, the All-Merciful, is
indeed the Ever-Forgiving, the All-Bountiful. He is the Tenderhearted,
the Merciful, and He is the Bestower, the Compassionate, the
All-Hearing, the All-Seeing.

\subsection{45}\label{section-2875}

\begin{otherlanguage}{arabic}

و در مراسلهٴ دیگر آنجناب ذکر مرفوع جناب ذبیح علیه بهآء اللّه و رحمته بود
لوحی مخصوص جناب آقا میرزا غلامعلی علیه ٩ ٦٦ {[}بهآء اللّه{]} نازل و ذکر
من صعد در او از قلم اعلی جاری انشآءاللّه لم‌یزل و لایزال بعنایت حقّ فائز
باشند قوله جلّ کبریائه

\end{otherlanguage}

Va dar morāsele-ye dīgar-e ān janāb dhekr-e marfūʿ-e Janāb-e Dhabīḥ,
ʿalayhi Bahāʾ Allāh wa-raḥmatuhu, būd. Lawḥī makhṣūṣ-e Janāb-e Āqā Mīrzā
GholāmʿAlī, ʿalayhi 9 66 {[}Bahāʾ Allāh{]}, nāzel va dhekr-e man ṣaʿida
dar ū az Qalam-e Aʿlā jārī. In shāʾ Allāh lam-yazal va lā-yazāl
be-ʿenāyat-e Ḥaqq fāʾez bāshand. Qawluhu jalla kibriyāʾuhu.

In your other letter there was an exalted mention of Dhabīḥ---may the
glory of God and His mercy rest upon him. A Tablet has been sent down
especially for Āqā Mīrzā Ghulām-ʿAlī---upon him be 9 66 {[}the glory of
God{]}---and in it the remembrance of the one who ascended has flowed
from the Most Exalted Pen. God willing, may they always and forever
receive the loving care of the True One. His words---glorious is His
majesty:

\subsection{46}\label{section-2876}

\begin{otherlanguage}{arabic}

هو الأقدم الأعظم الأبهی

\end{otherlanguage}

Huwa al-Aqdam al-Aʿẓam al-Abhā.

He is the Most Ancient, the Most Great, the Most Glorious.

\subsection{47}\label{section-2877}

\begin{otherlanguage}{arabic}

قد طار المعانی الف سنة لیصل الی مقام هذا اللّفظ الّذی تفوّه به لسان اللّه
المهیمن القیّوم قد طار الرّوح فی قرون لا یحصیها الّا اللّه لیزور هذا الجسم
الّذی یمشی فی الفردوس الأعلی الّذی سمّی بالسّجن الأعظم بلسان اللّه مالک الغیب
و الشّهود شهد اللّه لنفسه انّه ظهر بالحقّ و ینطق فی هذا القصر الّذی اقتصرت
فیه الأمور بذکر هذا الذّکر الأعظم الّذی به ظهر الفزع الأکبر فی العالم و
ناح کلّ غافل محجوب

\end{otherlanguage}

Qad ṭāra al-maʿānī alf sana li-yaṣila ilā maqām hādhā al-lafẓ alladhī
tafawwaha bihi Lisān Allāh al-Muhaymin al-Qayyūm. Qad ṭāra al-rūḥ fī
qurūn lā yuḥṣīhā illā Allāh li-yazūra hādhā al-jism alladhī yamshī fī
al-firdaws al-aʿlā alladhī summiya bi-al-Sijn al-Aʿẓam bi-Lisān Allāh
Mālik al-ghayb wa-al-shuhūd. Shahida Allāh li-nafsihi annahu ẓahara
bi-al-ḥaqq wa-yanṭiqu fī hādhā al-qaṣr alladhī iqtaṣarat fīhi al-umūr
bi-dhikr hādhā al-Dhikr al-Aʿẓam alladhī bihi ẓahara al-fazaʿ al-akbar
fī al-ʿālam wa-nāḥa kull ghāfil maḥjūb.

Meanings have soared for a thousand years to reach the station of this
word uttered by the Tongue of God, the All-Protecting, the
Self-Subsisting. The spirit has soared through centuries that none but
God can number to visit this body that walks in the highest Paradise,
which the Tongue of God, the Lord of the unseen and the seen, has named
the Most Great Prison. God bears witness to His own Self that He has
appeared in truth and speaks, in this palace where affairs have been
curtailed, in remembrance of this Most Great Remembrance, through whom
the greatest terror has appeared in the world and every heedless, veiled
soul has lamented.

\subsection{48}\label{section-2878}

\begin{otherlanguage}{arabic}

یا غلام یذکرک مولی الأنام و یذکّرک اذ کنت لدی الوجه قائماً لدی باب فتح
باصبع القدرة علی من فی السّموات و الأرض کذلک نطق جمال القدم فی هذا المقام
المرفوع ان انظر ثمّ اذکر اذ سمعت ندآء ربّک و تکلّم معک بالشّفقة و العنایة انّ
ربّک لهو العزیز الودود و نذکر اباک الّذی فاز بطراز الایقان و شهد بما شهد
اللّه و نطق بما نزل فی لوح المحفوظ نشهد انّه ما توقّف فی هذا الأمر و اذا
سمع النّدآء من شطر ربّه مالک الأسمآء قال لبیّک لبیّک یا اله الغیب و الشّهود
قد سمع و اجاب و خضع اذ اتاه الأمر من مشرق وحی ربّه مالک الوجود انّا زیّنّا
سمآء قلبه بنیّر حبّی العزیز المحبوب قد فاز بالرّحیق الأطهر اذ کان العراق
منظری الأکبر اخذ و شرب باسمی المبارک العزیز الممنوع مرّة فاز فی العراق و
اخری اذ اخرجنا الغافلون من ارض السّرّ کذلک ینطق جمال القدم امام وجه العالم
ولکنّ الأمم اکثرهم لا یفقهون

\end{otherlanguage}

Yā Ghulām, yadhkuruka Mawlā al-anām wa-yudhakkiruka idh kunta ladā
al-wajh qāʾiman ladā bāb futiḥa bi-iṣbaʿ al-qudra ʿalā man fī al-samāwāt
wa-l-arḍ. Kadhālika naṭaqa Jamāl al-Qidam fī hādhā al-maqām al-marfūʿ.
An unẓur thumma udhkur idh samiʿta nidāʾ rabbika wa-takallama maʿaka
bi-l-shafaqa wa-l-ʿināya. Inna rabbaka la-huwa al-ʿAzīz al-Wadūd.
Wa-nadhkuru abāka alladhī fāza bi-ṭirāz al-īqān wa-shahida bi-mā shahida
Allāh wa-naṭaqa bi-mā nazala fī lawḥ al-maḥfūẓ. Nashhadu annahu mā
tawaqqafa fī hādhā al-amr wa-idhā samiʿa al-nidāʾ min shaṭr rabbihi
Mālik al-asmāʾ qāla labbayka labbayka yā Ilāh al-ghayb wa-l-shuhūd. Qad
samiʿa wa-ajāba wa-khaḍaʿa idh atāhu al-amr min mashriq waḥy rabbihi
Mālik al-wujūd. Innā zayyannā samāʾ qalbihi bi-nayyir ḥubbī al-ʿazīz
al-maḥbūb. Qad fāza bi-l-raḥīq al-aṭhar idh kāna al-ʿIrāq manẓarī
al-akbar. Akhadha wa-shariba bi-ismī al-mubārak al-ʿazīz al-mamnūʿ.
Marratan fāza fī al-ʿIrāq wa-ukhrā idh akhrajanā al-ghāfilūn min Arḍ
al-Sirr. Kadhālika yanṭiqu Jamāl al-Qidam amāma wajh al-ʿālam wa-lākinna
al-umam aktharuhum lā yafqahūn.

O Ghulām! The Lord of humanity remembers you and reminds you of when you
were in His presence, standing at a door opened by the finger of power
to all who are in the heavens and on earth. Thus has the Ancient Beauty
spoken in this exalted station. Look, then recall when you heard the
call of your Lord and He spoke with you with compassion and loving care.
Your Lord is indeed the Almighty, the All-Loving. We also remember your
father, who attained the adornment of certitude, testified to what God
testified, and spoke what was sent down in the Preserved Tablet. We
testify that he did not hesitate in this Cause, and when he heard the
call from the direction of his Lord, the Sovereign of Names, he said,
``Here I am! Here I am, O God of the unseen and the seen!'' He heard,
answered, and bowed in submission when the command came to him from the
Dayspring of the revelation of his Lord, the Sovereign of existence. We
adorned the heaven of his heart with the radiant orb of My mighty and
beloved love. He attained the purest wine when Iraq was My Most Great
Scene. He took it and drank in My blessed, mighty, inviolable name. Once
he attained in Iraq, and another time when the heedless expelled Us from
the Land of Mystery. Thus speaks the Ancient Beauty before the face of
the world, yet most of its peoples do not understand.

\subsection{49}\label{section-2879}

\begin{otherlanguage}{arabic}

یا ذبیح یذکرک مولاک بما ورد علیک فی هذا السّبیل الّذی یطوفه الدّلیل بما نسب
الی اللّه العلیم الخبیر اشهد انّک آمنت باللّه اذ ظهر بالحقّ و ما منعتک حجاب
العلوم و لا ظنون الّذین کفروا باللّه ربّ العالمین قد نبذت ما عند النّاس و
اخذت ما عند اللّه ربّک و ربّ العرش العظیم و اعرضت عن القوم مقبلاً الی اسمی
القیّوم و ترکت الشّئونات شوقاً للقآء اللّه العزیز الحمید انت الّذی ما خوّفک
عوی الذّیاب و لا قباع الّذین اعرضوا عن الحقّ متمسّکین بأوهامهم الا انّهم فی
خسران مبین

\end{otherlanguage}

Yā Dhabīḥ, yadhkuruka mawlāka bi-mā warada ʿalayka fī hādhā al-sabīl
alladhī yaṭūfuhu al-dalīl bi-mā nusiba ilā Allāh al-ʿAlīm al-Khabīr.
Ashhadu annaka āmanta bi-Allāh idh ẓahara bi-l-ḥaqq wa-mā manaʿatka
ḥijāb al-ʿulūm wa-lā ẓunūn alladhīna kafarū bi-Allāh Rabb al-ʿālamīn.
Qad nabadhta mā ʿinda al-nās wa-akhadhta mā ʿinda Allāh rabbika wa-Rabb
al-ʿarsh al-ʿaẓīm wa-aʿraḍta ʿan al-qawm muqbilan ilā ismī al-Qayyūm
wa-tarakta al-shuʾūnāt shawqan li-liqāʾ Allāh al-ʿAzīz al-Ḥamīd. Anta
alladhī mā khawwafaka ʿawī al-dhiʾāb wa-lā qubāʿ alladhīna aʿraḍū ʿan
al-Ḥaqq mutamassikīn bi-awhāmihim. Alā innahum fī khusrān mubīn.

O Dhabīḥ! Your Lord remembers you for what befell you in this path,
around which the guide circles because of its association with God, the
All-Knowing, the All-Aware. I testify that you believed in God when He
appeared in truth, and neither the veil of learning nor the conjectures
of those who disbelieved in God, the Lord of the worlds, kept you back.
You cast aside what people possessed and took hold of what is with God,
your Lord and the Lord of the mighty throne. You turned away from the
people, turning toward My name, the Self-Subsisting, and abandoned
worldly affairs out of longing to meet God, the Almighty, the
Praiseworthy. You are the one whom neither the howling of wolves nor the
grunting of those who turned away from the Truth, clinging to their vain
imaginings, could frighten. Indeed, they are in manifest loss.

\subsection{50}\label{section-2880}

\begin{otherlanguage}{persian}

یا ذبیح قد اقبل الیک مقصود العالم و یذکرک بما حیّ به کلّ عظم رمیم یذکرک و
یذکر ایّامک و ما نطقت فی هذا الصّراط العزیز البدیع انت الّذی اخذت الکتاب
بقوّة ربّک و ترکت الکتاب السّجّین فی اسفل الجحیم انت الّذی اعترفت بما اعترف
به اللّه و شربت رحیق الاستقامة فی یوم فیه ارتعدت فرائص الأسمآء و اضطرب من
فی ناسوت الانشآء انّ ربّک لهو الذّاکر المشفق الکریم قد سمعت فی اللّه شماتة
کلّ ظالم و حملت من المعتدین ما ناح به کلّ منصف بصیر طوبی لک و لمن یتقرّب
الیک و یذکرک بما تفوّه به الذّکر الأعظم الّذی باسمه فتح باب السّمآء و ماج
بحر الأسمآء و هاج عرف اللّه الغفور الرّحیم السّلام الّذی کان مشرقاً من افق
دارالسّلام و البهآء الّذی طلع و لاح من افق البقآء علیک و علی من نطق بذکرک
و ثنائک و زار قبرک امراً من لدن آمر قدیر انتهی

\end{otherlanguage}

Yā Dhabīḥ, qad aqbala ilayka Maqṣūd al-ʿālam wa-yadhkuruka bi-mā ḥayya
bihi kull ʿaẓm ramīm. Yadhkuruka wa-yadhkuru ayyāmaka wa-mā naṭaqta fī
hādhā al-ṣirāṭ al-ʿazīz al-badīʿ. Anta alladhī akhadhta al-kitāb
bi-quwwat rabbika wa-tarakta al-kitāb al-sijjīn fī asfal al-jaḥīm. Anta
alladhī iʿtarafta bi-mā iʿtarafa bihi Allāh wa-sharibta raḥīq
al-istiqāma fī yawm fīhi irtaʿadat farāʾiṣ al-asmāʾ wa-iḍṭaraba man fī
nāsūt al-inshāʾ. Inna rabbaka la-huwa al-Dhākir al-Mushfiq al-Karīm. Qad
samiʿta fī Allāh shamātat kull ẓālim wa-ḥamalta min al-muʿtadīn mā nāḥa
bihi kull munṣif baṣīr. Ṭūbā laka wa-li-man yataqarrabu ilayka
wa-yadhkuruka bi-mā tafawwaha bihi al-Dhikr al-Aʿẓam alladhī bi-ismihi
futiḥa bāb al-samāʾ wa-māja baḥr al-asmāʾ wa-hāja ʿarf Allāh al-Ghafūr
al-Raḥīm. Al-salām alladhī kāna mushriqan min ufuq Dār al-Salām
wa-l-bahāʾ alladhī ṭalaʿa wa-lāḥa min ufuq al-baqāʾ ʿalayka wa-ʿalā man
naṭaqa bi-dhikrika wa-thanāʾika wa-zāra qabraka amran min ladun Āmir
Qadīr. Intahā.

O Dhabīḥ! The Desire of the world has turned toward you and remembers
you with that through which every decayed bone has come to life. He
remembers you, your days, and what you uttered on this mighty and
wondrous path. You are the one who took hold of the Book through the
power of your Lord and abandoned the book of Sijjīn in the lowest depths
of hell. You are the one who acknowledged what God acknowledged and
drank the wine of steadfastness on a day when the names trembled in
terror and all who inhabit the earthly realm of creation were shaken.
Your Lord is indeed the One who remembers, the Compassionate, the
Generous. For the sake of God, you heard the gloating of every oppressor
and endured from the aggressors what caused every fair-minded and
discerning soul to lament. Blessed are you and whoever draws near to you
and remembers you with what the Most Great Remembrance has uttered, He
through whose name the gate of heaven was opened, the ocean of names
surged, and the fragrance of God, the Ever-Forgiving, the Merciful,
spread abroad. May the peace that shone from the horizon of the Abode of
Peace and the glory that rose and gleamed from the horizon of eternity
rest upon you and upon whoever speaks in remembrance and praise of you
and visits your grave, as a command from an all-powerful Ordainer. End.

\subsection{51}\label{section-2881}

\begin{otherlanguage}{persian}

آیا لسانی در عالم موجود هست که قابل شکر باشد و یا لایق ثنا اگر جمیع بحور
بامواج و تمام دنیا باریاح و اشجار باغصان و اوراق و اثمار و ارض بکلأ و ما
فیها و علیها من الأتلال و الجبال و الغیاض و السّمآء و ما فیها من الشّموس و
الکواکب و الأقمار بذکر محبوب آفاق ناطق شوند هرآینه بقصور معترفند و بعدم
بلوغ مقرّ ذرّه‌ئی از اعمال در ساحت اقدس ضایع نشده و نخواهد شد

\end{otherlanguage}

Āyā lisānī dar ʿālam mawjūd hast keh qābil-i shukr bāshad va yā lāyiq-i
thanā? Agar jamīʿ-i buḥūr bi-amvāj va tamām-i dunyā bi-aryāḥ va ashjār
bi-aghṣān va awrāq va athmār va arḍ bi-kalāʾ va mā fīhā wa-ʿalayhā min
al-atlāl wa-l-jibāl wa-l-ghiyāḍ wa-l-samāʾ wa-mā fīhā min al-shumūs
wa-l-kawākib wa-l-aqmār bi-dhikr-i Maḥbūb-i āfāq nāṭiq shavand,
har-āyinih bi-quṣūr muʿtarifand va bi-ʿadam-i bulūgh muqirr. Dharrih-ī
az aʿmāl dar sāḥat-i aqdas ḍāyiʿ nashudih va nakhwāhad shud.

Does any tongue exist in the world that is capable of giving thanks or
worthy of offering praise? If all the seas with their waves, the whole
world with its winds, the trees with their branches, leaves, and fruits,
the earth with its herbage and all the hills, mountains, and forests
within and upon it, and heaven with its suns, stars, and moons were to
speak in remembrance of the Beloved of all regions, they would assuredly
confess their inadequacy and acknowledge their failure to attain. Not an
atom of deeds has been or will be lost in the Most Holy Court.

\subsection{52}\label{section-2882}

\begin{otherlanguage}{persian}

یا محبوب فؤادی در دریاهای رحمت و شفقت و عنایت حقّ جلّ جلاله تفکّر نمائید
اگر نفسی یک کلمه فی‌اللّه تکلّم نموده و یا نماید مکافات آن از قلم اعلی
ببهترین وجهی شده و میشود بسیار امور است که این عبد نظر بعدم اذن ستر
نموده یکمرآت از مرآیای علم الهی مظهر اسم سریع الحساب است اللّه اکبر از
ذکر اینمقام بسیار عظیم و بزرگ است و آن مرآت آناً فآناً از حرکات و سکنات و
بیانات و اشارات و دلالات و ظهورات عالم حکایت مینماید

\end{otherlanguage}

Yā Maḥbūb-i fuʾādī, dar daryāhā-yi raḥmat va shafaqat va ʿināyat-i Ḥaqq
jalla jalāluhu tafakkur namāyīd. Agar nafsī yik kalimih fī Allāh
takallum namūdih va yā namāyad, mukāfāt-i ān az Qalam-i Aʿlā bi-bihtarīn
vajhī shudih va mīshavad. Bisyār umūr ast keh īn ʿabd naẓar bi-ʿadam-i
idhn satr namūdih. Yik mirʾāt az marāyā-yi ʿilm-i ilāhī maẓhar-i ism-i
Sarīʿ al-Ḥisāb ast. Allāhu akbar az dhikr-i īn maqām; bisyār ʿaẓīm va
buzurg ast. Va ān mirʾāt ānan fa-ānan az ḥarakāt va sakanāt va bayānāt
va ishārāt va dalālāt va ẓuhūrāt-i ʿālam ḥikāyat mīnamāyad.

O beloved of my heart! Reflect on the oceans of the mercy, compassion,
and loving care of God, exalted be His glory. If anyone has spoken or
should speak a single word for the sake of God, its recompense has been
and will be granted in the finest manner by the Most Exalted Pen. There
are many matters that this servant has concealed for lack of permission.
One of the mirrors of divine knowledge is the manifestation of the name
``Swift in Reckoning.'' God is greater than any mention of this station;
it is exceedingly mighty and great. Moment by moment, that mirror
reflects the movements and stillnesses, utterances, gestures,
indications, and manifestations of the world.

\subsection{53}\label{section-2883}

\begin{otherlanguage}{persian}

وقتی از اوقات این عبد فانی را محض عنایت احضار فرمودند و حسب الأمر اقلّ من
آن از برای این عبد کشف شد کاد ان اکون منصعقاً این کلمه از ظنون و اوهام
نبوده و نیست بلکه مظهر طلعت یقین است لا حول و لا قوّة الّا به دیگر از
مرایای فوق آن چه ذکر شود کلّ امر فی قبضة قدرته و معلّق باذنه و ارادته و
مشیّته از حقّ جلّ جلاله اینخادم فانی سائل است که او را در کلّ احوال باو
نگذارد و بعنایتش مستقیم دارد انّه ولیّی فی الآخرة و الأولی

\end{otherlanguage}

Vaqtī az awqāt īn ʿabd-i fānī rā maḥḍ-i ʿināyat iḥḍār farmūdand va ḥasab
al-amr aqall-i min ān az barā-yi īn ʿabd kashf shud. Kāda an akūna
munṣaʿiqan. Īn kalimih az ẓunūn va awhām nabūdih va nīst, balkih
maẓhar-i ṭalʿat-i yaqīn ast. Lā ḥawla wa-lā quwwata illā bihi. Dīgar az
marāyā-yi fawq-i ān chih dhikr shavad? Kullu amr fī qabḍat qudratihi
wa-muʿallaq bi-idhnihi wa-irādatihi wa-mashīʾatihi. Az Ḥaqq jalla
jalāluhu īn khādim-i fānī sāʾil ast keh ū rā dar kull-i aḥvāl bi-ū
nagudhārad va bi-ʿināyatash mustaqīm dārad. Innahu waliyyī fī al-ākhira
wa-l-ūlā.

On one occasion, out of sheer loving kindness, He summoned this
evanescent servant into His presence, and, in accordance with His
command, it was unveiled to this servant for less than an instant. I
nearly fell thunderstruck. These words have never arisen from conjecture
or imagination; rather, they manifest the countenance of certitude.
There is no power or strength except through Him. What, then, can be
said of the mirrors beyond it? Every matter is within the grasp of His
power and depends on His permission, His will, and His purpose. This
evanescent servant asks God, exalted be His glory, never to leave him to
himself under any circumstances, but to keep him steadfast through His
loving care. He is indeed my Guardian in the next world and in this one.

\subsection{54}\label{section-2884}

\begin{otherlanguage}{persian}

الحمد للّه آنچه آنجناب از مقصود عالمیان خواستند واحداً بعد واحد باجابت
فائز شد یشهد بذلک ما نزل فیمن اردت ذکره من اللّه الفرد الخبیر آن محبوب
باید آنچه در بارهٴ جناب آقا میرزا غلامعلی علیه بهآء ٦٦ {[}اللّه{]} و من
صعد نازل شده باو برسانند ذکر تکبیر هم از قبل فانی بایشان بسته بمرحمت آن
محبوب است و اینکه در بارهٴ جناب آقا محمّد رضا علیه ٩ ٦٦ {[}بهآء اللّه{]}
ابن من طاف العرش و فاز باللّقآء کوثر الوصال فی اشهر معدودات مرقوم داشتید
هنیئاً له و مریئاً له قد شهد لسان العظمة لخضوعه و خشوعه للّه الغنیّ المتعال
این تفصیل در ساحت من لا یعزب عن علمه من شیء عرض شد و این آیات بدیعهٴ
منیعهٴ محکمهٴ جامعهٴ از افق فضل مالک احدیّه نازل قوله جلّ کبریائه

\end{otherlanguage}

Al-ḥamdu li-llāh, ānchih ān janāb az Maqṣūd-i ʿālamiyān khwāstand
vāḥidan baʿda vāḥid bi-ijābat fāʾiz shud. Yashhadu bi-dhālika mā nazala
fīman aradta dhikrahu min Allāh al-Fard al-Khabīr. Ān maḥbūb bāyad
ānchih dar bārih-yi Jināb-i Āqā Mīrzā Ghulām-ʿAlī ʿalayhi bahāʾ 66
{[}Allāh{]} va man ṣaʿida nāzil shudih bi-ū birasānand. Dhikr-i takbīr
ham az qibal-i fānī bi-īshān bastih bi-marḥamat-i ān maḥbūb ast. Va
īnkih dar bārih-yi Jināb-i Āqā Muḥammad-Riḍā ʿalayhi 9 66 {[}bahāʾ
Allāh{]} ibn man ṭāfa al-ʿarsh wa-fāza bi-l-liqāʾ kawthar al-wiṣāl fī
ashhur maʿdūdāt marqūm dāshtīd, hanīʾan lahu wa-marīʾan lahu. Qad
shahida Lisān al-ʿAẓama li-khuḍūʿihi wa-khushūʿihi li-llāh al-Ghanī
al-Mutaʿāl. Īn tafṣīl dar sāḥat-i man lā yaʿzubu ʿan ʿilmihi min shayʾ
ʿarḍ shud va īn āyāt-i badīʿih-yi manīʿih-yi muḥkamih-yi jāmiʿih az
ufuq-i faḍl-i Mālik-i aḥadiyyih nāzil. Qawluhu jalla kibriyāʾuhu.

Praise be to God, whatever you requested of the Desire of all peoples
has been answered, one request after another. What has been sent down by
God, the One, the All-Aware, concerning the person whose remembrance you
desired bears witness to this. That beloved friend should convey to him
whatever has been revealed concerning Jināb-i Āqā Mīrzā Ghulām-ʿAlī,
upon him be the glory of 66 {[}God{]}, and the one who has ascended.
Conveying a greeting of ``God is Most Great'' to him on behalf of this
evanescent one likewise depends on that beloved friend's kindness. As
for what you wrote concerning Jināb-i Āqā Muḥammad-Riḍā, upon him be 9
66 {[}the glory of God{]}, son of the one who circled the Throne and
attained the meeting---the Kawthar of reunion within a few months---may
it be wholesome and delightful to him. The Tongue of Grandeur has borne
witness to his submission and humility before God, the Self-Sufficient,
the Most Exalted. These details were presented in the court of Him from
whose knowledge nothing escapes, and these wondrous, exalted, firm, and
all-embracing verses were revealed from the horizon of the grace of the
Sovereign of oneness. His words, exalted be His majesty:

\subsection{55}\label{section-2885}

\begin{otherlanguage}{arabic}

بنام یکتا خداوند بینا

\end{otherlanguage}

Bi-nām-i yiktā Khudāvand-i bīnā.

In the name of the one, all-seeing God.

\subsection{56}\label{section-2886}

\begin{otherlanguage}{persian}

عالم را انوار ظهور احاطه نموده بشأنیکه حصاة طور بأعلی النّدآء بما نطقت به
السّدرة ندا مینماید مع‌ذلک عباد ارض باصغای آن فائز نشده‌اند مگر نفوس معدوده
که باسم بها از رحیق بقا آشامیده‌اند و از غبار دنیا طاهر و پاکیزه گشته‌اند
هر صاحب بصری که از قبل بوده بذکر این یوم امنع اعزّ اقدس مشغول و از حقّ جلّ
جلاله لقایش را سائل و آمل و چون ادراک عباد و قوّهٴ عرفانشان باوهام و ظنون
آلوده بود از مطلع غیب مکنون و سرّ مخزون محروم ماندند و از انوار آفتاب
یقین ممنوع گشتند

\end{otherlanguage}

ʿĀlam rā anvār-i ẓuhūr iḥāṭih namūdih bi-shaʾnī keh ḥaṣāt-i Ṭūr bi-aʿlā
al-nidāʾ bi-mā naṭaqat bihi al-Sidra nidā mīnamāyad. Maʿa dhālik ʿibād-i
arḍ bi-iṣghā-yi ān fāʾiz nashudih-and magar nufūs-i maʿdūdih keh
bi-ism-i Bahā az raḥīq-i baqā āshāmīdih-and va az ghubār-i dunyā ṭāhir
va pākīzih gashtih-and. Har ṣāḥib-i baṣarī keh az qabl būdih bi-dhikr-i
īn yawm-i amnaʿ-i aʿazz-i aqdas mashghūl va az Ḥaqq jalla jalāluhu
liqāyash rā sāʾil va āmil. Va chūn idrāk-i ʿibād va quvvih-yi
ʿirfānishān bi-awhām va ẓunūn ālūdih būd, az maṭlaʿ-i ghayb-i maknūn va
sirr-i makhzūn maḥrūm māndand va az anvār-i āftāb-i yaqīn mamnūʿ
gashtand.

The lights of revelation have encompassed the world to such a degree
that the pebbles of Sinai cry out at the top of their voices what the
Lote Tree has uttered. Nevertheless, the inhabitants of the earth have
failed to hear it, except for a few souls who have drunk the wine of
eternity in the name of Bahā and have been cleansed and purified of the
dust of the world. Every person of insight in former times was occupied
with the remembrance of this most exalted, most mighty, most holy Day,
asking God, exalted be His glory, and hoping to attain it. Yet, because
the understanding of His servants and their power of recognition were
tainted with vain imaginings and conjectures, they remained deprived of
the Dayspring of the concealed Unseen and the hidden Mystery and were
denied the lights of the Sun of certitude.

\subsection{57}\label{section-2887}

\begin{otherlanguage}{persian}

حمد کن محبوب عالم را که تو را مؤید فرمود باقبال بر امریکه علمای عصر از
آن اعراض نمودند طوبی از برای نفوسیکه سطوت فراعنه که ارباب عمائم باشند
ایشانرا از رحیق اطهر بازنداشت و از منظر اکبر محروم نساخت ذکر کن ایّامی را
که اسم مرفوع طائف حول بود و بخدمت مشغول طوبی له و لأیّامه انّه ما نقض
میثاق اللّه و عهده قد کان من اهل الوفآء من قلم الأمر مسطورا از دنیا و
اختلافات و شئونات آن محزون مباش آنچه نپاید لایق ذکر و توجّه نبوده و نیست
انشآءاللّه در صدد آن باش که قابل و لایق است امروز بزرگ است و امرش عظیم هر
کوری قابل مشاهده نه و هر بی سمعی لایق اصغاء نه قدر مقامیکه بآن فائز شدی
بدان و باسم حافظ حفظش نما طوبی لمن تمسّک بالاستقامة فی هذا الأمر العظیم
انّا ذکرناک من قبل و اباک انّ ربّک لهو الفضّال الکریم البهآء علیک و علی
الّذین سبقوا العباد فی عرفان اللّه مالک هذا الیوم المنیر انتهی

\end{otherlanguage}

Ḥamd kun Maḥbūb-i ʿālam rā keh tū rā muʾayyad farmūd bi-iqbāl bar amrī
keh ʿulamā-yi ʿaṣr az ān iʿrāḍ namūdand. Ṭūbā az barā-yi nufūsī keh
saṭvat-i farāʿinih keh arbāb-i ʿamāʾim bāshand īshān rā az raḥīq-i aṭhar
bāz nadāsht va az manẓar-i akbar maḥrūm nasākht. Dhikr kun ayyāmī rā keh
ism-i marfūʿ ṭāʾif-i ḥawl būd va bi-khidmat mashghūl. Ṭūbā lahu
wa-li-ayyāmihi. Innahu mā naqaḍa mīthāq Allāh wa-ʿahdahu. Qad kāna min
ahl al-wafāʾ min Qalam al-Amr masṭūran. Az dunyā va ikhtilāfāt va
shuʾūnāt-i ān maḥzūn mabāsh. Ānchih napāyad lāyiq-i dhikr va tavajjuh
nabūdih va nīst. In shāʾ Allāh dar ṣadad-i ān bāsh keh qābil va lāyiq
ast. Imrūz buzurg ast va amrash ʿaẓīm. Har kūrī qābil-i mushāhadih nah
va har bī-samʿī lāyiq-i iṣghā nah. Qadr-i maqāmī keh bi-ān fāʾiz shudī
bidān va bi-ism-i Ḥāfiẓ ḥifẓash namā. Ṭūbā li-man tamassaka
bi-l-istiqāma fī hādhā al-amr al-ʿaẓīm. Innā dhakarnāka min qabl
wa-abāka. Inna rabbaka la-huwa al-Faḍḍāl al-Karīm. Al-bahāʾ ʿalayka
wa-ʿalā alladhīna sabaqū al-ʿibād fī ʿirfān Allāh Mālik hādhā al-yawm
al-munīr. Intahā.

Praise the Beloved of the world, who confirmed you in turning toward a
Cause from which the learned of the age turned away. Blessed are the
souls whom the might of the Pharaohs---the wearers of turbans---did not
keep back from the purest wine or deprive of the Most Great Scene.
Recall the days when the exalted name circled around and was engaged in
service. Blessed are he and his days! He did not break God's covenant
and pledge. He has been inscribed by the Pen of the Cause among the
faithful. Do not grieve over the world, its dissensions, and its
affairs. What does not endure has never been and is not worthy of
mention or attention. God willing, pursue what is fitting and worthy.
This Day is great, and its Cause is mighty. Not every blind person is
capable of seeing, nor every deaf person fit to hear. Appreciate the
station you have attained and preserve it through the name of the
Preserver. Blessed is whoever holds fast to steadfastness in this mighty
Cause. We have remembered you before, and your father. Your Lord is
indeed the All-Bountiful, the Generous. Glory be upon you and upon those
who preceded His servants in recognizing God, the Sovereign of this
luminous Day. End.

\subsection{58}\label{section-2888}

\begin{otherlanguage}{persian}

الحمد للّه آفتاب عنایت مشرق و دوستان و آقایان که بافق اعلی متوجّهند هر یک
را از انوار خود قسمت و نصیب عطا فرموده انّه لهو الفیّاض الباذل المعطی
الکریم

\end{otherlanguage}

Al-ḥamdu li-llāh, āftāb-i ʿināyat mushriq va dūstān va āqāyān keh
bi-ufuq-i aʿlā mutavajjihand har yik rā az anvār-i khud qismat va naṣīb
ʿaṭā farmūdih. Innahu la-huwa al-Fayyāḍ al-Bādhil al-Muʿṭī al-Karīm.

Praise be to God, the sun of loving care is shining, and to each of the
friends and honored gentlemen who turn toward the Most Exalted Horizon
it has granted a portion and share of its light. He is indeed the
All-Bountiful, the Munificent, the Bestower, the Generous.

\subsection{59}\label{section-2889}

\begin{otherlanguage}{persian}

و اینکه طلب اذن نموده بودند این ایّام این ارض فی‌الجمله مضطرب است و اضطراب
باطنی از قبل از قلم اعلی جاری شده و خبر داده انّه لهو العلیم الخبیر لذا
توجّه جایز نه الأمر بیده یکتب لمن اراد اجر لقائه و هو الفضّال العزیز
الحکیم این فانی خدمت جناب مذکور علیه بهآء اللّه تکبیر و سلام میرساند و از
حقّ میطلبد که ایشان لازال در انجمن اهل بها بوجه منیر و قلب منیر بثنای
دوست یکتا گویا باشند

\end{otherlanguage}

Va īnkih ṭalab-i idhn namūdih būdand, īn ayyām īn arḍ fī al-jumlih
muḍṭarib ast va iḍṭirāb-i bāṭinī az qabl az Qalam-i Aʿlā jārī shudih va
khabar dādih. Innahu la-huwa al-ʿAlīm al-Khabīr. Lidhā tavajjuh jāyiz
nah. Al-amr bi-yadihi yaktubu li-man arāda ajr liqāʾihi wa-huwa
al-Faḍḍāl al-ʿAzīz al-Ḥakīm. Īn fānī khidmat-i Jināb-i madhkūr ʿalayhi
bahāʾ Allāh takbīr va salām mīrasānad va az Ḥaqq mīṭalabad keh īshān lā
yazāl dar anjuman-i ahl-i Bahā bi-vajh-i munīr va qalb-i munīr
bi-thanā-yi Dūst-i yiktā gūyā bāshand.

As for his request for permission, this land is somewhat unsettled these
days, and its inward unrest was previously foretold by the Most Exalted
Pen. He is indeed the All-Knowing, the All-Aware. Therefore, coming here
is not permissible. The matter is in His hand; He ordains for whomever
He wills the reward of meeting Him, and He is the All-Bountiful, the
Almighty, the All-Wise. This evanescent one conveys greetings of ``God
is Most Great'' and peace to the aforementioned gentleman, upon him be
the glory of God, and asks God that he may always, with a radiant face
and an illumined heart, speak the praise of the one Beloved in the
gathering of the people of Bahā.

\subsection{60}\label{section-2890}

\begin{otherlanguage}{persian}

و اینکه در بارهٴ جناب آقا علی میلانی الّذی فاز بالاذن و رجع بالاذن مرقوم
داشتید فی‌الحقیقه جناب ایشان بنار محبّة اللّه مشتعلند چند یوم قبل از لسان
مبارک حقّ جلّ جلاله ذکرش را این فانی استماع نمود یشهد الخادم انّه فاز
باللّقآء و شرب رحیق البقآء هنیئاً له بما شرب و لی بما اذکره حبّاً لله ربّنا و
ربّکم و محبوبنا و محبوبکم و محبوب من فی السّموات و الأرض لوح امنع اقدس هم
چندی قبل مخصوص ایشان نازل و ارسال شد لتقرّ به عیون اهل مدائن الأسمآء
الّذین تمسّکوا بحبل عنایة ربّهم العلیّ الأبهی

\end{otherlanguage}

Va īnkih dar bārih-yi Jināb-i Āqā ʿAlī Mīlānī alladhī fāza bi-l-idhn
wa-rajaʿa bi-l-idhn marqūm dāshtīd, fī al-ḥaqīqih Jināb-i īshān bi-nār-i
maḥabbat Allāh mushtaʿiland. Chand yawm qabl az lisān-i mubārak-i Ḥaqq
jalla jalāluhu dhikrash rā īn fānī istimāʿ namūd. Yashhadu al-khādim
annahu fāza bi-l-liqāʾ wa-shariba raḥīq al-baqāʾ. Hanīʾan lahu bi-mā
shariba wa-lī bi-mā adhkuruhu ḥubban li-llāh Rabbinā wa-Rabbikum
wa-Maḥbūbinā wa-Maḥbūbikum wa-Maḥbūb man fī al-samāwāt wa-l-arḍ. Lawḥ-i
amnaʿ-i aqdas ham chandī qabl makhṣūṣ-i īshān nāzil va irsāl shud
li-taqarra bihi ʿuyūn ahl madāʾin al-asmāʾ alladhīna tamassakū bi-ḥabl
ʿināyat rabbihim al-ʿAlī al-Abhā.

As for what you wrote concerning Jināb-i Āqā ʿAlī Mīlānī, who obtained
permission and returned with permission, he is truly aflame with the
fire of the love of God. A few days ago, this evanescent one heard him
mentioned by the blessed tongue of God, exalted be His glory. This
servant testifies that he attained the meeting and drank the wine of
eternity. May joy be his for what he drank, and mine for remembering him
out of love for God, our Lord and your Lord, our Beloved and your
Beloved, and the Beloved of all who are in the heavens and on earth. A
most exalted, most holy Tablet was also revealed especially for him some
time ago and dispatched, that through it the eyes of the inhabitants of
the cities of names who have held fast to the cord of the loving care of
their Lord, the Most High, the Most Glorious, might be gladdened.

\subsection{61}\label{section-2891}

\begin{otherlanguage}{persian}

و اینکه جناب مذکور استدعای ذکر و عنایت از حقّ جلّ جلاله در بارهٴ جناب آقا
محمّد جعفر و آقا محمّد حسین و اخت خود نموده بودند این طلب را مطلوب حقیقی
روح من فی الملکوت الأمر و الخلق فداه امضا فرمود انّه مجیب کلّ سائل و معطی
کلّ آمل و معین کلّ ضعیف و وطن کلّ غریب و مقصد کلّ قاصد امین

\end{otherlanguage}

Va īnke janāb-i madhkūr istidʿā-yi dhikr va ʿināyat az Ḥaqq jalla
jalāluhu dar bāra-yi janāb-i Āqā Muḥammad Jaʿfar va Āqā Muḥammad Ḥusayn
va ukht-i khud namūda būdand, īn ṭalab rā Maṭlūb-i Ḥaqīqī, rūḥu man fī
al-malakūt al-amr wa-al-khalq fidāhu, imḍā farmūd. Innahu mujību kulli
sāʾil wa-muʿṭī kulli āmil wa-muʿīnu kulli ḍaʿīf wa-waṭanu kulli gharīb
wa-maqṣidu kulli qāṣidin amīn.

As for the aforementioned gentleman's request that God, exalted be His
glory, remember and bestow His favor upon Āqā Muḥammad Jaʿfar, Āqā
Muḥammad Ḥusayn, and his own sister, the True Object of Desire---may the
souls of all in the kingdoms of command and creation be sacrificed for
Him---granted this request. He indeed answers every petitioner, gives to
everyone who hopes, aids everyone who is weak, is a homeland to every
stranger, and is the destination of every faithful seeker.

\subsection{62}\label{section-2892}

\begin{otherlanguage}{arabic}

هذا ما نزل من ملکوت البیان لمن سمّی بمحمّد جعفر قوله جلّ کبریائه

\end{otherlanguage}

Hādhā mā nazala min malakūt al-bayān li-man summiya bi-Muḥammad Jaʿfar,
qawluhu jalla kibriyāʾuhu.

This is what was revealed from the kingdom of utterance for the one
named Muḥammad Jaʿfar. These are His words, exalted be His majesty:

\subsection{63}\label{section-2893}

\begin{otherlanguage}{arabic}

بسمی الأقدس الأبهی

\end{otherlanguage}

Bi-smī al-Aqdas al-Abhā.

In My name, the Most Holy, the Most Glorious.

\subsection{64}\label{section-2894}

\begin{otherlanguage}{arabic}

کتاب انزله الوهّاب لمن اجاب مولاه اذ اتی بالحقّ لیهدی الکلّ الی صراطه
المستقیم انّه ظهر بالفضل و اظهر بالعدل ما اراد انّه لهو العلیم الحکیم انّ
الّذین سمعوا و اجابوا انّهم من اهل هذا المنظر الکریم یصلّینّ علیهم الملأ
الأعلی و الّذین طافوا حول ارادة ربّهم الخبیر

\end{otherlanguage}

Kitābun anzalahu al-Wahhāb li-man ajāba mawlāhu idh atā bi-al-ḥaqq
li-yahdiya al-kulla ilā ṣirāṭihi al-mustaqīm. Innahu ẓahara bi-al-faḍl
wa-aẓhara bi-al-ʿadl mā arād. Innahu la-huwa al-ʿAlīm al-Ḥakīm. Inna
alladhīna samiʿū wa-ajābū innahum min ahli hādhā al-manẓar al-karīm.
Yuṣalliyanna ʿalayhim al-malaʾ al-aʿlā wa-alladhīna ṭāfū ḥawla irādati
rabbihim al-Khabīr.

A Book sent down by the All-Bountiful for one who answered his Lord when
He came with the truth to guide everyone to His straight path. He has
appeared through grace and manifested through justice whatever He
willed. He is indeed the All-Knowing, the All-Wise. Those who have heard
and answered are indeed the people of this noble Scene. The Concourse on
high and those who circle around the will of their Lord, the All-Aware,
will surely invoke blessings upon them.

\subsection{65}\label{section-2895}

\begin{otherlanguage}{arabic}

کم من عبد فاز بهذا المقام الأعلی و کم من عبد انکر و اعرض عن اللّه الفرد
الواحد العزیز الحمید طوبی لک یا جعفر بما اردت ان تشرب الرّحیق الأطهر من
ید مالک القدر الّذی ینطق فی منظره الأکبر انّه لا اله الّا انا الظّاهر النّاطق
السّمیع البصیر

\end{otherlanguage}

Kam min ʿabdin fāza bi-hādhā al-maqām al-aʿlā wa-kam min ʿabdin ankara
wa-aʿraḍa ʿan Allāh al-Fard al-Wāḥid al-ʿAzīz al-Ḥamīd. Ṭūbā laka yā
Jaʿfar bi-mā aradta an tashraba al-raḥīq al-aṭhar min yadi Mālik
al-Qadar alladhī yanṭiqu fī manẓarihi al-akbar: innahu lā ilāha illā anā
al-Ẓāhir al-Nāṭiq al-Samīʿ al-Baṣīr.

How many a servant has attained this most exalted station, and how many
a servant has denied and turned away from God, the Peerless, the One,
the Almighty, the All-Praised. Blessed are you, O Jaʿfar, for desiring
to drink the purest wine from the hand of the Lord of destiny, who
proclaims from His most great Scene: ``There is indeed no God but Me,
the Manifest, the Speaker, the All-Hearing, the All-Seeing.''

\subsection{66}\label{section-2896}

\begin{otherlanguage}{arabic}

انّک اذا وجدت عرفی من بیانی و اخذک سکر رحیق آیاتی ولّ وجهک شطر سجنی و قل

\end{otherlanguage}

Innaka idhā wajadta ʿarfī min bayānī wa-akhadhaka sukru raḥīqi āyātī
walli wajhaka shaṭra sijnī wa-qul.

When you perceive My fragrance in My utterance and are seized by
intoxication from the wine of My verses, turn your face toward My prison
and say:

\subsection{67}\label{section-2897}

\begin{otherlanguage}{persian}

اشهد انّک اظهرت ما کان مکنوناً فی کنز الغیب و بیّنت ما کان مستوراً فی العلم
و مسطوراً فی کتب اللّه ربّ العالمین روحی لسجنک الفدآء یا ایّها المظلوم
الغریب نفسی لحزنک الفدآء یا من حملت الشّدائد فی سبیلک المستقیم اسئلک بأن
تؤیّدنی علی الاستقامة علی هذا الأمر الّذی به زلّت اقدام العارفین الحمد للّه
ربّ العالمین انتهی

\end{otherlanguage}

Ashhadu annaka aẓharta mā kāna maknūnan fī kanz al-ghayb wa-bayyanta mā
kāna mastūran fī al-ʿilm wa-masṭūran fī kutub Allāh Rabb al-ʿālamīn.
Rūḥī li-sijnika al-fidāʾ yā ayyuhā al-maẓlūm al-gharīb. Nafsī li-ḥuznika
al-fidāʾ yā man ḥamalta al-shadāʾida fī sabīlika al-mustaqīm. Asʾaluka
bi-an tuʾayyidanī ʿalā al-istiqāma ʿalā hādhā al-amr alladhī bihi zallat
aqdāmu al-ʿārifīn. Al-ḥamdu li-llāhi Rabb al-ʿālamīn. Intahā.

``I testify that You have manifested what lay hidden in the treasury of
the unseen and made clear what was concealed in knowledge and inscribed
in the Books of God, the Lord of the worlds. May my spirit be sacrificed
for Your imprisonment, O wronged and exiled One. May my soul be
sacrificed for Your sorrow, O You who have borne hardships in Your
straight path. I ask You to strengthen me in steadfastness in this
Cause, through which the feet of those possessed of knowledge have
slipped. Praise be to God, the Lord of the worlds.'' End.

\subsection{68}\label{section-2898}

\begin{otherlanguage}{arabic}

و هذا ما نزّل من جبروت البرهان لمن سمّی بمحمّد حسین قوله عزّ اجلاله

\end{otherlanguage}

Wa-hādhā mā nuzzila min jabarūt al-burhān li-man summiya bi-Muḥammad
Ḥusayn, qawluhu ʿazza ijlāluhu.

And this is what was sent down from the dominion of proof for the one
named Muḥammad Ḥusayn. These are His words, exalted be His glory:

\subsection{69}\label{section-2899}

\begin{otherlanguage}{arabic}

بسمی المهیمن علی الأسمآء

\end{otherlanguage}

Bi-smī al-Muhaymin ʿalā al-asmāʾ.

In My name, the Guardian over all names.

\subsection{70}\label{section-2900}

\begin{otherlanguage}{persian}

کتاب انزل مالک العرش و الثّری لمن اقبل الی مشرق الوحی لیفرح و یکون من
الشّاکرین یا محمّد قبل حسین ان استمع ندآء المظلوم الّذی اعرض عنه عباده و
اعترض علیه کلّ عالم بعید انّه ینادی فی کلّ الأحیان من فی الامکان الی الرّحمن
ولکنّ القوم اکثرهم فی ضلال مبین قد نبذوا الههم عن ورائهم مقبلین الی الّذین
اتّبعوا من دون بیّنة من اللّه ربّ العالمین لا تحزن فی شیء توکّل فی کلّ الأمور
علی القویّ القدیر انّه مع الّذین آمنوا یشهد و یری و هو السّمیع البصیر لا
یعزب عن علمه من شیء و لا تضعفه قدرة الأقویآء انّه لهو المقتدر القدیر

\end{otherlanguage}

Kitābun anzala Mālik al-ʿarsh wa-al-tharā li-man aqbala ilā mashriq
al-waḥy li-yafraḥa wa-yakūna min al-shākirīn. Yā Muḥammad qabla Ḥusayn,
an istamiʿ nidāʾ al-maẓlūm alladhī aʿraḍa ʿanhu ʿibāduhu wa-iʿtaraḍa
ʿalayhi kullu ʿālimin baʿīd. Innahu yunādī fī kulli al-aḥyān man fī
al-imkān ilā al-Raḥmān wa-lākinna al-qawma aktharuhum fī ḍalālin mubīn.
Qad nabadhū ilāhahum ʿan warāʾihim muqbilīna ilā alladhīna ittabaʿū min
dūni bayyinatin min Allāh Rabb al-ʿālamīn. Lā taḥzan fī shayʾ, tawakkal
fī kulli al-umūr ʿalā al-Qawiyy al-Qadīr. Innahu maʿa alladhīna āmanū
yashhadu wa-yarā wa-huwa al-Samīʿ al-Baṣīr. Lā yaʿzubu ʿan ʿilmihi min
shayʾ wa-lā tuḍʿifuhu qudratu al-aqwiyāʾ. Innahu la-huwa al-Muqtadir
al-Qadīr.

A Book sent down by the Lord of the throne and the earth for one who has
turned toward the Dawning-Place of revelation, that he may rejoice and
be among the thankful. O Muḥammad before Ḥusayn, listen to the call of
the Wronged One, from whom His servants have turned away and against
whom every wayward scholar has raised objections. At all times He calls
all who dwell in the world of existence to the All-Merciful, yet most of
the people are in manifest error. They have cast their God behind them,
turning toward those whom they have followed without any proof from God,
the Lord of the worlds. Grieve over nothing; in all matters put your
trust in the Strong, the All-Powerful. He is with those who believe; He
witnesses and sees, and He is the All-Hearing, the All-Seeing. Nothing
escapes His knowledge, nor does the power of the mighty weaken Him. He
is indeed the All-Powerful, the Omnipotent.

\subsection{71}\label{section-2901}

\begin{otherlanguage}{arabic}

یا اهل البهآء دعوا ما یدعوکم الی الفنآء خذوا کأس البقآء باسمی الکریم
کذلک نوّرنا افق الکتاب بذکر اللّه مالک الرّقاب و انزلنا ما قرّت به عیون
المقرّبین البهآء علی اهل البهآء الّذین ما نقضوا میثاق اللّه و عهده و اخذهم
سکر رحیق بیان ربّهم علی شأن نبذوا العالم و الأمم عن ورائهم مقبلین الی بحر
النّور فی هذا الظّهور الّذی اذا ظهر سجدت الوجوه للّه ربّ العالمین انتهی

\end{otherlanguage}

Yā ahl al-Bahāʾ, daʿū mā yadʿūkum ilā al-fanāʾ, khudhū kaʾs al-baqāʾ
bi-smī al-Karīm. Kadhālika nawwarnā ufuq al-kitāb bi-dhikr Allāh Mālik
al-riqāb wa-anzalnā mā qarrat bihi ʿuyūnu al-muqarrabīn. Al-Bahāʾ ʿalā
ahl al-Bahāʾ alladhīna mā naqaḍū mīthāq Allāh wa-ʿahdahu wa-akhadhahum
sukru raḥīqi bayāni rabbihim ʿalā shaʾnin nabadhū al-ʿālama wa-al-umama
ʿan warāʾihim muqbilīna ilā baḥr al-nūr fī hādhā al-ẓuhūr alladhī idhā
ẓahara sajadat al-wujūhu li-llāhi Rabb al-ʿālamīn. Intahā.

O people of Bahā, abandon whatever calls you to extinction; take up the
cup of everlasting life in My name, the All-Generous. Thus have We
illumined the horizon of the Book with the remembrance of God, the Lord
of all, and sent down what has gladdened the eyes of those brought near.
Glory rest upon the people of Bahā who have not broken God's covenant
and pledge, and whom intoxication from the wine of their Lord's
utterance has seized so completely that they have cast the world and its
peoples behind them, turning toward the ocean of light in this
Revelation, at whose appearance all faces bowed before God, the Lord of
the worlds. End.

\subsection{72}\label{section-2902}

\begin{otherlanguage}{arabic}

و هذا ما نزل من ملکوت الفضل لأخته

\end{otherlanguage}

Wa-hādhā mā nazala min malakūt al-faḍl li-ukhtihi.

And this is what was revealed from the kingdom of grace for his sister:

\subsection{73}\label{section-2903}

\begin{otherlanguage}{arabic}

بنام محبوب عالمیان

\end{otherlanguage}

Bi-nām-i Maḥbūb-i ʿālamiyān.

In the name of the Beloved of all the worlds.

\subsection{74}\label{section-2904}

\begin{otherlanguage}{arabic}

ای کنیز الهی اگر قدر این کلمهٴ مبارکه را بدانی از شوق ترک عالم گوئی و
بذکر عالم یکتا مشغول گردی چه مقدار از ملوک و سلاطین و از ملکه‌های عالم که
در لیالی و ایّام از حقّ جلّ جلاله لقای مطلع آیاترا مسئلت مینمودند و در یوم
ظهور از او محروم ماندند نه عرف قمیصش را یافتند و نه ندایش را اصغاء
نمودند با غفلت تمام بمقرّ خود راجع گشتند

\end{otherlanguage}

Ay kanīz-i ilāhī, agar qadr-i īn kalima-yi mubāraka rā bidānī az shawq
tark-i ʿālam gūʾī va bi-dhikr-i ʿĀlim-i Yaktā mashghūl gardī. Chi miqdār
az mulūk va salāṭīn va az malika-hā-yi ʿālam ki dar layālī va ayyām az
Ḥaqq jalla jalāluhu liqā-yi Maṭlaʿ-i āyāt rā masʾalat mīnamūdand va dar
yawm-i ẓuhūr az ū maḥrūm māndand. Na ʿarf-i qamīṣash rā yāftand va na
nidāyash rā iṣghā namūdand. Bā ghaflat-i tamām bi-maqarr-i khud rājiʿ
gashtand.

O handmaid of God, if you knew the worth of this blessed word, you would
forsake the world out of longing and occupy yourself with remembering
the One All-Knowing. How many kings and sovereigns and queens of the
world implored God, exalted be His glory, night and day to meet the
Dawning-Place of His verses, yet on the day of His appearance remained
deprived of Him. They neither perceived the fragrance of His garment nor
listened to His call. In utter heedlessness they returned to their
abode.

\subsection{75}\label{section-2905}

\begin{otherlanguage}{persian}

ای امة‌اللّه ذکرت در ساحت اقدس مذکور و مظلوم آفاق باین کلمات درّیّات تو را
ذکر نموده قدر خود را بدان و مقام این امر را بشناس و از حقّ جلّ جلاله در
جمیع احیان بخواه تا تو را بطراز استقامت مزیّن فرماید و از شبهات و اشارات
قوم مقدّس و مطهّر دارد انّ ربّک لهو المقتدر علی ما یشآء و فی قبضته زمام من
فی ملکوت الانشآء انّه لهو المهیمن المقتدر القدیر انتهی

\end{otherlanguage}

Ay Amatuʾllāh, dhikrat dar sāḥat-i aqdas madhkūr va Maẓlūm-i āfāq bi-īn
kalimāt-i durriyyāt tu rā dhikr namūda. Qadr-i khud rā bidān va maqām-i
īn amr rā bishinās va az Ḥaqq jalla jalāluhu dar jamīʿ-i aḥyān bikhvāh
tā tu rā bi-ṭirāz-i istiqāmat muzayyan farmāyad va az shubahāt va
ishārāt-i qawm muqaddas va muṭahhar dārad. Inna rabbaki la-huwa
al-muqtadiru ʿalā mā yashāʾ wa-fī qabḍatihi zimāmu man fī malakūt
al-inshāʾ. Innahu la-huwa al-Muhaymin al-Muqtadir al-Qadīr. Intahā.

O handmaid of God, you have been mentioned in the Most Holy Court, and
the Wronged One of the world has remembered you in these pearl-like
words. Know your own worth and recognize the station of this Cause. At
all times ask God, exalted be His glory, to adorn you with steadfastness
and keep you pure and unsullied by the doubts and insinuations of the
people. Your Lord indeed has power over whatever He wills, and in His
grasp are the reins of all who dwell in the kingdom of creation. He is
indeed the All-Protecting, the All-Powerful, the Omnipotent. End.

\subsection{76}\label{section-2906}

\begin{otherlanguage}{persian}

دیگر آن محبوب هر قسم مصلحت میدانند این الواح را برسانند تا سبب و علّت فرح
و سکون و استقامت گردد انّه یهدی من یشآء الی لقائه العزیز البدیع و ینزل
لمن یشآء ما یقرّبه الیه و هو الغفور الکریم جمیع دستخطّهای آن محبوب که مطرّز
بمناجات با حقّ سبحانه و تعالی بود در ساحت امنع اقدس فائز شد و جمیع آنچه
استدعا شده بود در پیشگاه حضور مقصود عالمیان معروض گشت هذا ما نطق به لسان
القدم قوله عزّ اجلاله

\end{otherlanguage}

Dīgar ān maḥbūb har qism maṣlaḥat mīdānand īn alvāḥ rā birasānand tā
sabab va ʿillat-i faraḥ va sukūn va istiqāmat gardad. Innahu yahdī man
yashāʾ ilā liqāʾihi al-ʿazīz al-badīʿ wa-yunazzilu li-man yashāʾ mā
yuqarribuhu ilayhi wa-huwa al-Ghafūr al-Karīm. Jamīʿ-i dastkhaṭṭ-hā-yi
ān maḥbūb ki muṭarraz bi-munājāt bā Ḥaqq subḥānahu va taʿālā būd dar
sāḥat-i amnaʿ-i aqdas fāʾiz shud va jamīʿ-i ānchi istidʿā shuda būd dar
pīshgāh-i ḥuḍūr-i Maqṣūd-i ʿālamiyān maʿrūḍ gasht. Hādhā mā naṭaqa bihi
Lisān al-Qidam, qawluhu ʿazza ijlāluhu.

Furthermore, let that beloved friend deliver these Tablets in whatever
manner he deems advisable, so that they may bring joy, tranquility, and
steadfastness. He guides whomever He wills to His mighty and wondrous
presence, and sends down for whomever He wills what will draw him near
to Him. He is the All-Forgiving, the All-Generous. All that beloved
friend's handwritten letters, adorned with prayers to God, glorified and
exalted be He, reached the most inaccessible, most holy Court, and all
the requests they contained were presented before the presence of the
Desire of all the worlds. This is what the Tongue of Eternity
proclaimed. These are His words, exalted be His glory:

\subsection{77}\label{section-2907}

\begin{otherlanguage}{persian}

یا ورقآء یا ایّها النّاطق بثنآء مالک الأسمآء حمد کن محبوب عالمیان را که
فائز شدی بآنچه که مقصود از وجود عالم بود قد سمعت و اقبلت و فزت بلقآء
اللّه ربّ العالمین جمیع از برای آن است که بکلمهٴ رضا از قلم اعلی فائز شوند
بصدهزار لسان باید تدارک این شکر را نمود چه که فائز شدی باینمقام ان احفظ
و کن من الشّاکرین لحاظ عنایت حقّ بتو بوده و خواهد بود

\end{otherlanguage}

Yā Varqāʾ, yā ayyuhā al-nāṭiqu bi-thanāʾi Mālik al-asmāʾ, ḥamd kun
Maḥbūb-i ʿālamiyān rā ki fāʾiz shudī bi-ānchi ki maqṣūd az vujūd-i ʿālam
būd. Qad samiʿta wa-aqbalta wa-fuzta bi-liqāʾ Allāh Rabb al-ʿālamīn.
Jamīʿ az barā-yi ān ast ki bi-kalima-yi riḍā az Qalam-i Aʿlā fāʾiz
shavand. Bi-ṣad-hizār lisān bāyad tadāruk-i īn shukr rā namūd, chi ki
fāʾiz shudī bi-īn maqām. An iḥfaẓ wa-kun min al-shākirīn. Liḥāẓ-i
ʿināyat-i Ḥaqq bi-tu būda va khvāhad būd.

O Varqā, O you who speak the praise of the Lord of names, praise the
Beloved of all the worlds, for you have attained that which was the
purpose of the world's existence. You have heard, turned toward Him, and
attained the presence of God, the Lord of the worlds. Everything is for
this: that they may receive a word of good pleasure from the Most
Exalted Pen. With a hundred thousand tongues you must strive to offer
thanks for this, for you have attained this station. Guard it and be
among the thankful. The gaze of God's loving care has rested upon you
and will continue to do so.

\subsection{78}\label{section-2908}

\begin{otherlanguage}{persian}

انشآءاللّه در جمیع احیان به ما یحبّ و یرضی فائز باشی دوستان الهی را در هر
دیار و مقام که ملاقات نمودی تکبیر برسان و به ما ینبغی لهذا الیوم الأقدس
الأنور آگاه نما که شاید شبهات و ظنونات علما و الوان دنیا اولیای حقّ را از
رحیق اطهر منع ننماید و جمیع را باعمال طیّبه و اخلاق روحانیّه امر مینمائیم
تا عرف اعمال عالم را آگاهی بخشد و مقصود از ظهور حقّ جلّ جلاله را بیابد و
هم‌چنین عرف اخلاق ملأ اعلی را مسرور نماید و بهجت عطا کند بگو امروز سدرهٴ
نار در قطب بحار مشتعل و ناطق و مجلّی جهد نمائید تا از انوار ظهور محروم
نمانید انسان نباید از حطب یابس پست‌تر مشاهده شود یک بیدا حطب را یکجذوه
مشتعل نماید و حال یکعالم نار امام وجوه ظاهر و دریای نور از نور جبین لائح
و مشهود بشنوید ندای دوست یکتا را و بکمال استقامت و اطمینان بذکر و ثنا و
تبلیغ امرش که اعظم اعمال است متمسّک شوید و در جمیع احوال باید بحکمت ناظر
باشید یا ورقآء یا ایّها الطّائف حولی و النّاظر الی وجهی و النّاطق بثنائی انّا
نکبّر و نسلّم من هذا المقام علیک و علی الّذین یسمعون قولک فی امر اللّه
المهیمن القیّوم انتهی

\end{otherlanguage}

In shāʾ Allāh dar jamīʿ-i aḥyān bi-mā yuḥibbu wa-yarḍā fāʾiz bāshī.
Dūstān-i ilāhī rā dar har diyār va maqām ki mulāqāt namūdī takbīr
birasān va bi-mā yanbaghī li-hādhā al-yawm al-aqdas al-anwar āgāh namā
ki shāyad shubahāt va ẓunūnāt-i ʿulamā va alvān-i dunyā awliyā-yi Ḥaqq
rā az raḥīq-i aṭhar manʿ nanamāyad. Va jamīʿ rā bi-aʿmāl-i ṭayyiba va
akhlāq-i rūḥāniyya amr mīnamāʾīm tā ʿarf-i aʿmāl ʿālam rā āgāhī bakhshad
va maqṣūd az ẓuhūr-i Ḥaqq jalla jalāluhu rā biyābad va ham-chunīn ʿarf-i
akhlāq malaʾ-i aʿlā rā masrūr namāyad va bahjat ʿaṭā kunad. Bigū imrūz
sidra-yi nār dar quṭb-i biḥār mushtaʿil va nāṭiq va mujallī. Jahd
namāʾīd tā az anvār-i ẓuhūr maḥrūm namānīd. Insān nabāyad az ḥaṭab-i
yābis past-tar mushāhada shavad. Yik bayḍā ḥaṭab rā yik jadhva mushtaʿil
namāyad va ḥāl yik ʿālam nār imām-i vujūh ẓāhir va daryā-yi nūr az nūr-i
jabīn lāʾiḥ va mashhūd. Bishnavīd nidā-yi Dūst-i Yaktā rā va bi-kamāl-i
istiqāmat va iṭmīnān bi-dhikr va thanā va tablīgh-i amrash ki aʿẓam-i
aʿmāl ast mutamassik shavīd va dar jamīʿ-i aḥvāl bāyad bi-ḥikmat nāẓir
bāshīd. Yā Varqāʾ, yā ayyuhā al-ṭāʾifu ḥawlī wa-al-nāẓiru ilā wajhī
wa-al-nāṭiqu bi-thanāʾī, innā nukabbiru wa-nusallimu min hādhā al-maqām
ʿalayka wa-ʿalā alladhīna yasmaʿūna qawlaka fī amr Allāh al-Muhaymin
al-Qayyūm. Intahā.

God willing, may you at all times attain what He loves and approves.
Wherever you meet the friends of God, convey to them the greeting ``God
is Most Great'' and acquaint them with what befits this most holy, most
luminous Day, that perhaps the doubts and conjectures of the scholars
and the world's many attractions may not keep the loved ones of God from
the purest wine. We enjoin everyone to perform good deeds and display
spiritual qualities, so that the fragrance of their deeds may awaken the
world and it may discover the purpose of the revelation of God, exalted
be His glory, and so that the fragrance of their character may likewise
gladden the Concourse on high and impart joy. Say: Today the Tree of
fire is ablaze at the center of the seas, speaking and shining forth.
Strive not to remain deprived of the lights of revelation. A human being
should not be found inferior to dry wood. A single flame sets a desert
of wood ablaze, yet now a world of fire is manifest before your faces,
and an ocean of light shines visibly from the radiance of His brow. Hear
the call of the one Friend, and with perfect steadfastness and
confidence hold fast to His remembrance and praise and to teaching His
Cause, which is the greatest of deeds. In all circumstances you must
have regard for wisdom. O Varqā, O you who circle around Me, gaze upon
My face, and speak My praise, from this station We convey the greeting
``God is Most Great'' and peace to you and to those who listen to your
words concerning the Cause of God, the All-Protecting, the
Self-Subsisting. End.

\subsection{79}\label{section-2909}

\begin{otherlanguage}{arabic}

و اینکه در بارهٴ جناب آقا میرزا مهدی و آقا میرزا سلیمان علیهما بهآء اللّه
ذکر نموده بودید هذا ما انزله اللّه لمن سمّی بمهدی علیه ٩ ٦٦ {[}بهآء
اللّه{]}

\end{otherlanguage}

Va īnke dar bāra-yi janāb-i Āqā Mīrzā Mahdī va Āqā Mīrzā Sulaymān
ʿalayhimā Bahāʾ Allāh dhikr namūda būdīd, hādhā mā anzalahu Allāh li-man
summiya bi-Mahdī ʿalayhi 9 66 {[}Bahāʾ Allāh{]}.

As for what you wrote concerning Āqā Mīrzā Mahdī and Āqā Mīrzā Sulaymān,
upon them both be the glory of God, this is what God has sent down for
the one named Mahdī, upon him be 9 66 {[}the glory of God{]}:

\subsection{80}\label{section-2910}

\begin{otherlanguage}{arabic}

هو الرّاکب علی السّحاب

\end{otherlanguage}

Huwa al-Rākib ʿalā al-saḥāb.

He is the One who rides upon the clouds.

\subsection{81}\label{section-2911}

\begin{otherlanguage}{arabic}

سبحان من اظهر ما اراد و عرّف الکلّ سبیله الواضح المستقیم و هدی العباد الی
مشرق العلم و انزل لهم ما فاح به عرف العرفان فی الامکان انّه لهو المنزل
القدیم مکلّم طور ینادی و یدعو النّاس الی ما اراد و انا الحاکم علی ما ارید

\end{otherlanguage}

Subḥāna man aẓhara mā arāda wa-ʿarrafa al-kulla sabīlahu al-wāḍiḥ
al-mustaqīm wa-hadā al-ʿibāda ilā mashriq al-ʿilm wa-anzala lahum mā
fāḥa bihi ʿarf al-ʿirfān fī al-imkān. Innahu la-huwa al-Munzil al-Qadīm.
Mukallimu Ṭūr yunādī wa-yadʿū al-nāsa ilā mā arāda wa-anā al-Ḥākim ʿalā
mā urīd.

Glory be to Him who has manifested what He willed, made known to
everyone His clear and straight path, guided His servants to the
Dawning-Place of knowledge, and sent down for them that through which
the fragrance of recognition has spread throughout existence. He is
indeed the Revealer, the Ancient of Days. He who spoke on Sinai calls
out and summons the people to what He has willed, and I am the Ruler
over whatever I will.

\subsection{82}\label{section-2912}

\begin{otherlanguage}{arabic}

یا مهدی سبّح باسم ربّک ثمّ اذکره بربوات المخلصین طوبی للّذین خضعوا لآیات
ربّهم و هنیئاً لک بما شربت کوثر حبّی العزیز البدیع انّا ذکرناک من قبل و فی
هذا اللّوح المنیر کن مستقیماً علی الأمر و ناطقاً بهذا الاسم العظیم طوبی لک
و لأبیک الّذی قصد المقصود و صعد الی اللّه ربّ العالمین انّا کنّا حاضراً حین
صعوده و انا السّامع المجیب ایّاک ان تحزنک شئونات الخلق توکّل فی کلّ الأمور
علی الحقّ مالک هذا الیوم المنیع انّا نؤیّدک علی خدمة الأمر و نوفّقک علی ما
قصدته فی ایّام اللّه العلیم الخبیر البهآء علیک و علی کلّ مقبل مستقیم انتهی

\end{otherlanguage}

Yā Mahdī, sabbiḥ bi-smi rabbika thumma udhkurhu bi-rabawāt al-mukhliṣīn.
Ṭūbā li-alladhīna khaḍaʿū li-āyāti rabbihim wa-hanīʾan laka bi-mā
sharibta kawthara ḥubbī al-ʿazīz al-badīʿ. Innā dhakarnāka min qablu
wa-fī hādhā al-lawḥ al-munīr. Kun mustaqīman ʿalā al-amr wa-nāṭiqan
bi-hādhā al-ism al-ʿaẓīm. Ṭūbā laka wa-li-abīka alladhī qaṣada al-Maqṣūd
wa-ṣaʿida ilā Allāh Rabb al-ʿālamīn. Innā kunnā ḥāḍiran ḥīna ṣuʿūdihi
wa-anā al-Sāmiʿ al-Mujīb. Iyyāka an tuḥzinaka shuʾūnu al-khalq. Tawakkal
fī kulli al-umūr ʿalā al-Ḥaqq Mālik hādhā al-yawm al-manīʿ. Innā
nuʾayyiduka ʿalā khidmat al-amr wa-nuwaffiquka ʿalā mā qaṣadtahu fī
ayyām Allāh al-ʿAlīm al-Khabīr. Al-Bahāʾ ʿalayka wa-ʿalā kulli muqbilin
mustaqīm. Intahā.

O Mahdī, glorify the name of your Lord, then remember Him upon the
heights of the sincere. Blessed are those who have humbled themselves
before the verses of their Lord, and joy be yours for having drunk from
the river of My mighty and wondrous love. We have remembered you before
and in this luminous Tablet. Be steadfast in the Cause and proclaim this
great Name. Blessed are you and your father, who sought the Desired One
and ascended to God, the Lord of the worlds. We were present at the time
of his ascension, and I am the Hearer, the Answerer. Beware lest the
affairs of the people grieve you. In all matters put your trust in the
True One, the Lord of this unassailable Day. We strengthen you in
serving the Cause and grant you success in what you have intended in the
days of God, the All-Knowing, the All-Aware. Glory rest upon you and
upon everyone who turns toward Him and remains steadfast. End.

\subsection{83}\label{section-2913}

\begin{otherlanguage}{arabic}

و هذا ما انزله الرّحمن لمن سمّی بسلیمان قوله عزّ کبریائه

\end{otherlanguage}

Wa-hādhā mā anzalahu al-Raḥmān li-man summiya bi-Sulaymān. Qawluhu ʿazza
kibriyāʾuhu:

And this is what the All-Merciful revealed for the one named Sulaymān.
His words---exalted be His majesty:

\subsection{84}\label{section-2914}

\begin{otherlanguage}{arabic}

هو الأقدس الأعظم الأبهی

\end{otherlanguage}

Huwa al-Aqdas al-Aʿẓam al-Abhā.

He is the Most Holy, the Most Great, the All-Glorious.

\subsection{85}\label{section-2915}

\begin{otherlanguage}{persian}

یا سلیمان ایّام نزدیک بآخر است و اهل ارض در غفلت مبین و ضلالت آشکار
مشاهده میشوند لسان صبح به قد اتی الحقّ ناطق و یوم به انا یوم اللّه ذاکر و
اشیاء به قد نزّل امّ الکتاب قائل گوشیکه فائز شود بندا و قلبیکه مزیّن گردد
بحبّ اللّه مالک الأسمآء قلیل بوده و هست در صبح و شام اهل ارض از احزاب
مختلفه بمساجد و معابد و هیاکل و کنائس لذکر اللّه توجّه مینمایند و از مذکور
غافل در حین ذکر مذکور امام عیون قائم ولکن از او محتجب می‌بینند و
نمیشناسند میشنوند و ادراک نمینمایند جمیع این غفلت و ضلالت و نادانی که
عالم را فراگرفته سبب و علّت اوّلیّهٴ آن علمای ارض بوده و هستند چه که
اوهاماتی ترتیب داده‌اند و آن را در اذهان جهّال ارض بصدهزار حیل و وساوس
داخل نموده‌اند تا آن اوهام باقی از اشراقات انوار آفتاب یقین محرومند لذا
خود و عباد بیچاره از فیض فیّاض ممنوع و از ماینفعهم محروم و امرای ارض اگر
فی‌الحقیقه مطیع آن نفوس غافله بودند حقّ و آثارش در عالم یافت نمیشد مگر نزد
نفوسیکه از شبهات و حجبات و اوهام عباد مقدّس بوده و هستند باسم حقّ معزّزند و
در ذلّت حقّ باسباب مختلفه متمسّک و متشبّث کاش باین اکتفا مینمودند لا ونفسی
الحقّ هر صاحب بصری آثار ظلم آن نفوس غافله را مشاهده مینماید

\end{otherlanguage}

Yā Soleymān, ayyām nazdīk be-ākhar ast va ahl-e arḍ dar ghaflat-e mobīn
va ḍalālat-e āshkār moshāhede mīshavand. Lesān-e ṣobḥ be-«qad atā
al-Ḥaqq» nāṭeq va yowm be-«anā yowmu Allāh» dhāker va ashyā be-«qad
nuzzila Ummu al-Kitāb» qāʾel. Gūshī ke fāʾez shavad be-nedā va qalbī ke
mozayyan gardad be-ḥobb-e Allāh Mālik al-Asmāʾ qalīl būde va hast. Dar
ṣobḥ va shām ahl-e arḍ az aḥzāb-e mokhtalefe be-masājed va maʿābed va
hayākel va kanāʾes li-dhikri Allāh tavajjoh mīnamāyand va az madhkūr
ghāfel. Dar ḥīn-e dhekr madhkūr amām-e ʿoyūn qāʾem va lāken az ū
moḥtajeb. Mībīnand va nemīshenāsand, mīshenavand va edrāk nemīnamāyand.
Jamīʿ-e īn ghaflat va ḍalālat va nādānī ke ʿālam rā farā gerefte sabab
va ʿellat-e avvaliyye-ye ān ʿolamā-ye arḍ būde va hastand, che ke
owhāmātī tartīb dāde-and va ān rā dar adhhān-e johhāl-e arḍ be-ṣad hezār
ḥiyal va vasāves dākhel nemūde-and. Tā ān owhām bāqī, az eshrāqāt-e
anvār-e āftāb-e yaqīn maḥrūmand. Ledhā khod va ʿebād-e bīchāre az feyḍ-e
Fayyāḍ mamnūʿ va az mā yanfaʿuhum maḥrūm. Va omarā-ye arḍ agar
fi-l-ḥaqīqe moṭīʿ-e ān nofūs-e ghāfele būdand, Ḥaqq va āthārash dar
ʿālam yāft nemīshod magar nazd-e nofūsī ke az shobohāt va ḥojobāt va
owhām-e ʿebād moqaddas būde va hastand. Be-esm-e Ḥaqq moʿazzazand va dar
dhellat-e Ḥaqq be-asbāb-e mokhtalefe motamassek va motashabbeth. Kāsh
be-īn ektefā mīnamūdand. Lā wa-nafsī al-Ḥaqq, har ṣāḥeb-e baṣarī āthār-e
ẓolm-e ān nofūs-e ghāfele rā moshāhede mīnamāyad.

O Sulaymān, the days are nearing their end, and the people of the earth
are seen to be in manifest heedlessness and evident error. The tongue of
dawn proclaims, ``The Truth has come''; the day declares, ``I am the Day
of God''; and all things say, ``The Mother Book has been sent down.''
Few have been, and remain, the ears that attain to the Call and the
hearts adorned with love for God, the Lord of Names. Morning and
evening, the people of the earth, of diverse denominations, turn toward
mosques, places of worship, temples, and churches to remember God, yet
remain heedless of the One they remember. While they remember Him, the
One remembered stands before their eyes, yet they are veiled from Him.
They see but do not recognize; they hear but do not understand. The
primary cause and source of all this heedlessness, error, and ignorance
that have engulfed the world have been, and remain, the religious
scholars of the earth. For they have devised vain imaginings and,
through a hundred thousand stratagems and insinuations, instilled them
in the minds of the ignorant on earth. As long as those imaginings
endure, they remain deprived of the radiant light of the sun of
certainty. Thus both they and the hapless servants are barred from the
grace of the All-Bountiful and deprived of what would benefit them. Had
the rulers of the earth truly obeyed those heedless souls, neither the
True One nor His traces would have been found in the world, except among
souls who have been, and remain, sanctified from the doubts, veils, and
vain imaginings of the people. They are honored through the name of the
True One, yet cling to and employ various means to abase Him. If only
they had been content with this! No, by My own Self, the Truth! Everyone
with sight beholds the traces of the injustice of those heedless souls.

\subsection{86}\label{section-2916}

\begin{otherlanguage}{persian}

طوبی لعالم عرف المعلوم عجب در اینکه مع این اعمال علمای این اراضی و اطراف
را قبول ندارند و خود را ناجی و هادی و افقه و اتقی میشمرند تبّا لهم و
للّذین اتّبعوهم من دون حجّة و بیّنة و برهان و دلیل یا سلیمان حمد کن محبوب
امکان را که تو را تأیید فرمود بر عرفان این فضلیکه آسمان شبه آن را ندیده
و ارض نظیر آن را حمل ننموده ان افرح و کن من الشّاکرین الحمد للّه ربّ
العالمین انتهی

\end{otherlanguage}

Ṭūbā li-ʿālimin ʿarafa al-Maʿlūm. ʿAjab dar īnke maʿ-e īn aʿmāl
ʿolamā-ye īn arāḍī va aṭrāf rā qabūl nadārand va khod rā nājī va hādī va
afqah va atqā mīshomorand. Tabban lahum wa-li-lladhīna ittabaʿūhum min
dūni ḥujjatin wa-bayyinatin wa-burhānin wa-dalīl. Yā Soleymān, ḥamd kon
Maḥbūb-e emkān rā ke to rā taʾyīd farmūd bar ʿerfān-e īn faḍlī ke āsemān
shebh-e ān rā nadīde va arḍ naẓīr-e ān rā ḥaml nanemūde. An ifraḥ wa-kun
mina al-shākirīn. Al-ḥamdu li-llāhi Rabbi al-ʿālamīn. Intahā.

Blessed is the learned one who has recognized the One who is the object
of knowledge. It is astonishing that, despite these deeds, they do not
accept the religious scholars of these lands and the surrounding
regions, and regard themselves as saved, as guides, as more learned in
religious law and more devout. Woe to them and to those who follow them
without testimony, clear evidence, proof, or reason! O Sulaymān, praise
the Beloved of all existence, who has confirmed you in recognizing this
grace, the like of which heaven has never seen and the earth has never
borne. Rejoice and be among the thankful. Praise be to God, the Lord of
the worlds. End.

\subsection{87}\label{section-2917}

\begin{otherlanguage}{arabic}

این دو لوح امنع اقدس مخصوص دو برادر از سماء عنایت نازل انشآءاللّه بلقای
آن فائز شوند و عرف آن را بیابند و به ما اراده اللّه قیام نمایند

\end{otherlanguage}

Īn do lowḥ-e amnaʿ-e aqdas makhṣūṣ-e do barādar az samā-ye ʿenāyat
nāzel. In shāʾ Allāh be-leqā-ye ān fāʾez shavand va ʿarf-e ān rā
biyāband va be-mā arāda Allāh qiyām namāyand.

These two most unassailable, most holy Tablets were sent down from the
heaven of loving care especially for the two brothers. God willing, they
will attain to them, perceive their fragrance, and arise to carry out
what God has willed.

\subsection{88}\label{section-2918}

\begin{otherlanguage}{arabic}

وقتی از اوقات این کلمهٴ علیا از مشرق ملکوت بیان ظاهر قوله جلّ جلاله

\end{otherlanguage}

Vaqtī az owqāt īn kalame-ye ʿolyā az mashreq-e malakūt-e bayān ẓāher.
Qawluhu jalla jalāluhu:

On one occasion, this exalted utterance appeared from the dawning-place
of the kingdom of utterance. His words---glorified be His glory:

\subsection{89}\label{section-2919}

\begin{otherlanguage}{arabic}

کینونت عرفان بمثابهٴ طیر ظاهر و بابدع الحان بین ارض و سماء باینکلمهٴ
علیا ناطق یا اهل الامکان ضعوا مشیّاتکم و ارادتکم و خذوا ما اراده اللّه
محبوبکم و مقصودکم انتهی

\end{otherlanguage}

Keynūnat-e ʿerfān be-mathābe-ye ṭeyr ẓāher va be-abdaʿ-e alḥān beyn-e
arḍ va samā be-īn kalame-ye ʿolyā nāṭeq: Yā ahla al-imkān, ḍaʿū
mashiyyātikum wa-irādātikum wa-khudhū mā arādahu Allāhu Maḥbūbukum
wa-Maqṣūdukum. Intahā.

The very essence of recognition has appeared in the form of a bird and,
with the most wondrous melodies, voices this exalted utterance between
earth and heaven: ``O people of the world of being! Lay aside your
wishes and desires, and take hold of what God, your Beloved and your
Heart's Desire, has willed.'' End.

\subsection{90}\label{section-2920}

\begin{otherlanguage}{arabic}

انشآءاللّه جمیع دوستان و اولیا و اصفیای حقّ باینمقام اعلی فائز شوند

\end{otherlanguage}

In shāʾ Allāh jamīʿ-e dūstān va owliyā va aṣfiyā-ye Ḥaqq be-īn maqām-e
aʿlā fāʾez shavand.

God willing, all the friends, loved ones, and chosen ones of the True
One will attain this most exalted station.

\subsection{91}\label{section-2921}

\begin{otherlanguage}{persian}

و اینکه در بارهٴ جناب آقا میرزا محمّد تقی علیه بهآء اللّه مرقوم داشتید ذکر
ایشان و تفصیل ایشان در ساحت امنع اقدس اعزّ اعلی عرض شد یک لوح از افق فضل
اشراق نمود انشآءاللّه بزیارتش فائز شوند و بکوثر بیان رحمن مرزوق گردند
قوله عزّ کبریائه

\end{otherlanguage}

Va īnke dar bāre-ye janāb-e Āqā Mīrzā Moḥammad Taqī ʿalayhi Bahāʾu Allāh
marqūm dāshtīd, dhekr-e īshān va tafṣīl-e īshān dar sāḥat-e amnaʿ-e
aqdas-e aʿazz-e aʿlā ʿarḍ shod. Yek lowḥ az ofoq-e faḍl eshrāq nemūd. In
shāʾ Allāh be-ziyāratash fāʾez shavand va be-kowthar-e bayān-e Raḥmān
marzūq gardand. Qawluhu ʿazza kibriyāʾuhu:

As for what you wrote concerning Āqā Mīrzā Muḥammad Taqī, upon him be
the glory of God, mention of him and the details of his circumstances
were presented at the most unassailable, most holy, most mighty, most
exalted Court. A Tablet shone forth from the horizon of grace. God
willing, he will attain to reading it and be nourished by the Kawthar of
the All-Merciful's utterance. His words---exalted be His majesty:

\subsection{92}\label{section-2922}

\begin{otherlanguage}{arabic}

بنام دوست یکتا

\end{otherlanguage}

Be-nām-e Dūst-e Yektā.

In the name of the one incomparable Friend.

\subsection{93}\label{section-2923}

\begin{otherlanguage}{persian}

امروز تبلیغ امر الهی از اعظم اعمال بوده و هست ولکن بحکمتیکه حکم آن از
قلم اعلی جاری حقّ جلّ جلاله از برای اصلاح عالم آمده تا جمیع من علی الأرض
بمثابهٴ یکنفس مشاهده شوند جمیع را تکلّم بیک لسان امر فرموده و هم‌چنین
بنوشتن یک خطّ تا جمیع ارض یکوطن کل محسوب شود وقتی باینکلمهٴ علیا نطق
نمودیم تا انوار آفتاب عدل جمیع احزاب مختلفه را منوّر فرماید و بمقام اتّحاد
کشاند لیس الفخر لمن یحبّ الوطن بل لمن یحبّ العالم اگر نفسی صفیر این طیر
معانی را که در این هوای بافضای نورانی طیران مینماید اصغا کند بطراز فراغت
کبری مزیّن شود اینست معنی حرّیّت حقیقی ولکنّ النّاس اکثرهم لا یشعرون

\end{otherlanguage}

Emrūz tablīgh-e Amr-e Elāhī az aʿẓam-e aʿmāl būde va hast va lāken
be-ḥekmatī ke ḥokm-e ān az Qalam-e Aʿlā jārī. Ḥaqq jalla jalāluhu az
barā-ye eṣlāḥ-e ʿālam āmade tā jamīʿ-e man ʿalā al-arḍ be-mathābe-ye yek
nafs moshāhede shavand. Jamīʿ rā takallom be-yek lesān amr farmūde va
hamchenīn be-neveshtan-e yek khaṭṭ tā jamīʿ-e arḍ yek vaṭan-e koll
maḥsūb shavad. Vaqtī be-īn kalame-ye ʿolyā noṭq nemūdīm tā anvār-e
āftāb-e ʿadl jamīʿ-e aḥzāb-e mokhtalefe rā monavvar farmāyad va
be-maqām-e etteḥād keshānad: Laysa al-fakhru li-man yuḥibbu al-waṭan bal
li-man yuḥibbu al-ʿālam. Agar nafsī ṣafīr-e īn ṭeyr-e maʿānī rā ke dar
īn havā-ye bā-faḍā-ye nūrānī ṭayarān mīnamāyad eṣghā konad, be-ṭerāz-e
farāghat-e kobrā mozayyan shavad. Īnast maʿnī-ye ḥorriyyat-e ḥaqīqī,
wa-lākinna al-nāsa aktharuhum lā yashʿurūn.

Today, teaching the Cause of God has been, and remains, among the
greatest of deeds, but it must be done with the wisdom enjoined by the
Most Exalted Pen. The True One---glorified be His glory---has come to
reform the world, so that all who dwell on earth may be seen as one
soul. He has commanded everyone to speak one language and likewise to
write in one script, so that the whole earth may be regarded as one
homeland for all. On one occasion We voiced this exalted utterance, that
the rays of the sun of justice might illumine all the diverse peoples
and draw them into unity: ``Glory belongs not to one who loves his
homeland, but to one who loves the world.'' Whoever listens to the song
of this bird of inner meanings, soaring in this spacious, luminous
atmosphere, will be adorned with the ornament of supreme detachment.
This is the meaning of true freedom, but most people do not understand.

\subsection{94}\label{section-2924}

\begin{otherlanguage}{persian}

و از برای تبلیغ امر هم شرایطی در کتاب الهی نازل اهل ارض بمثابهٴ طفل
مشاهده میشوند باید ببیانات خفیفهٴ لطیفه که بمثابه لبن است در اوّل تکلّم
نمود تا بایادی تربیت بحدّ بلوغ که مقام عرفان اللّه است فائز شوند از قبل
فرمودیم یک کلمه بمثابهٴ شیر است و اخری مانند شمشیر طوبی از برای حکیمی که
بحکمت عمل نماید و کلمة اللّه را در ارض طیّبهٴ مبارکه ودیعه گذارد

\end{otherlanguage}

Va az barā-ye tablīgh-e Amr ham sharāyeṭī dar Ketāb-e Elāhī nāzel. Ahl-e
arḍ be-mathābe-ye ṭefl moshāhede mīshavand. Bāyad be-bayānāt-e
khafīfe-ye laṭīfe ke be-mathābe-ye laban ast dar avval takallom nemūd tā
be-ayādī-ye tarbiyat be-ḥadd-e bolūgh ke maqām-e ʿerfān-e Allāh ast
fāʾez shavand. Az qabl farmūdīm yek kalame be-mathābe-ye shīr ast va
okhrā mānand-e shamshīr. Ṭūbā az barā-ye ḥakīmī ke be-ḥekmat ʿamal
namāyad va Kalimat Allāh rā dar arḍ-e ṭayyebe-ye mobārake vadīʿe
godhārad.

Conditions for teaching the Cause have also been revealed in the Book of
God. The people of the earth are regarded as children. At first one must
speak in mild, gentle words, which are like milk, so that through the
hands of education they may attain maturity, which is the station of
recognizing God. We have said before that one word is like milk and
another like a sword. Blessed is the wise person who acts with wisdom
and deposits the Word of God in pure and blessed soil.

\subsection{95}\label{section-2925}

\begin{otherlanguage}{persian}

ذکرت از قبل و بعد نزد مظلوم مذکور و لحاظ عنایت بتو متوجّه انشآءاللّه از
دریای حکمت و دانائی بیاشامی و بآنچه الیوم مقبول است مشغول گردی چون این
ارض فی‌الجمله در حرکت است لذا در اذن توقّف نمودیم انّه یقدّر لمن اراد و یکتب
لمن یرید اجر لقائه العزیز البدیع ان احمد اللّه بما جعلک فائزاً بهذا اللّوح
العظیم المنیع و انزل لک ما لا یتغیّر بدوام اسمآء اللّه ربّ العالمین انتهی

\end{otherlanguage}

Dhekrat az qabl va baʿd nazd-e Maẓlūm madhkūr va leḥāẓ-e ʿenāyat be-to
motavajjeh. In shāʾ Allāh az daryā-ye ḥekmat va dānāʾī biyāshāmī va
be-ānche elyowm maqbūl ast mashghūl gardī. Chūn īn arḍ fi-l-jomle dar
ḥarakat ast ledhā dar edhn tavaqqof nemūdīm. Innahu yuqaddiru li-man
arāda wa-yaktubu li-man yurīdu ajra liqāʾihi al-ʿAzīzi al-Badīʿ. An
iḥmadi Allāha bi-mā jaʿalaka fāʾizan bi-hādhā al-Lawḥi al-ʿaẓīmi
al-manīʿ wa-anzala laka mā lā yataghayyaru bi-dawāmi asmāʾi Allāhi Rabbi
al-ʿālamīn. Intahā.

You have been remembered before and since in the presence of this
Wronged One, and the gaze of loving care is turned toward you. God
willing, you will drink from the ocean of wisdom and knowledge and
occupy yourself with what is acceptable today. As this land is somewhat
unsettled, We have withheld permission. He ordains for whomever He
chooses, and records for whomever He wills the reward of attaining His
presence, the Almighty, the Wondrous. Praise God for having enabled you
to receive this mighty, unassailable Tablet and for having sent down for
you what will remain unchanged as long as the names of God, the Lord of
the worlds, endure. End.

\subsection{96}\label{section-2926}

\begin{otherlanguage}{persian}

و اینکه در بارهٴ ایشان مرقوم داشتید که بلقای حضرت سلطان الشّهدآء روحی
لتربته الفدآء و مؤانست با ایشان فائز بوده‌اند هنیئاً له و مریئاً له الحمد
للّه بما اراده اللّه موفّق شده‌اند و بلقای من احبّ اللّه و یحبّه اللّه رسیده‌اند
این خادم فانی خدمتشان عرض تکبیر میرساند و بشارت میدهد ایشانرا بعنایت حقّ
جلّ جلاله

\end{otherlanguage}

Va īnke dar bāre-ye īshān marqūm dāshtīd ke be-leqā-ye Ḥaḍrat-e Solṭān
al-Shohadā rūḥī li-turbatihi al-fidāʾ va moʾānesat bā īshān fāʾez
būde-and, hanīʾan lahu wa-marīʾan lahu. Al-ḥamdu li-llāh be-mā arāda
Allāh movaffaq shode-and va be-leqā-ye man aḥabba Allāha wa-yuḥibbuhu
Allāh resīde-and. Īn khādem-e fānī khedmatashān ʿarḍ-e takbīr mīresānad
va beshārat mīdahad īshān rā be-ʿenāyat-e Ḥaqq jalla jalāluhu.

As for what you wrote concerning him, that he had attained the presence
and enjoyed the companionship of His Holiness the King of Martyrs---may
my soul be a sacrifice for the dust of his resting place---may it bring
him joy and delight! Praise be to God, he has been enabled to accomplish
what God has willed and has met one who loved God and whom God loves.
This transient servant conveys to him a greeting of praise to God and
gives him the glad tidings of the loving care of the True
One---glorified be His glory.

\subsection{97}\label{section-2927}

\begin{otherlanguage}{persian}

و اینکه در بارهٴ جناب آقا میرزا جلیل علیه ٩ ٦٦ {[}بهآء اللّه{]} الملک
العزیز الجمیل مرقوم داشتید چندی قبل لوح امنع اقدس مخصوص ایشان نازل و
بهمراهی ابن جناب حاجی احاء علیه ٩ ٦٦ {[}بهآء اللّه{]} ارسال شد و هذه مرّة
اخری فضلاً علی جنابکم و علیه انّ ربّنا الرّحمن لهو المشفق المجیب

\end{otherlanguage}

Va īnke dar bāre-ye janāb-e Āqā Mīrzā Jalīl ʿalayhi 9 66 {[}Bahāʾu
Allāh{]} al-Malik al-ʿAzīz al-Jamīl marqūm dāshtīd, chandī qabl lowḥ-e
amnaʿ-e aqdas makhṣūṣ-e īshān nāzel va be-hamrāhī-ye ebn-e janāb-e Ḥājī
{[}corrupted text{]} ʿalayhi 9 66 {[}Bahāʾu Allāh{]} ersāl shod.
Wa-hādhihi marratun ukhrā faḍlan ʿalā janābikum wa-ʿalayhi. Inna Rabbanā
al-Raḥmāna la-huwa al-Mushfiqu al-Mujīb.

As for what you wrote concerning Āqā Mīrzā Jalīl, upon him be 9 66
{[}the glory of God{]}, the King, the Almighty, the All-Beauteous, some
time ago a most unassailable, most holy Tablet was revealed especially
for him and sent with the son of Ḥājī {[}corrupted original text{]},
upon him be 9 66 {[}the glory of God{]}. And here is another, as a favor
to you and to him. Indeed, our Lord, the All-Merciful, is the
Compassionate, the Answerer.

\subsection{98}\label{section-2928}

\begin{otherlanguage}{arabic}

قوله عزّ کبریائه هو النّاطق فی ملکوت البیان

\end{otherlanguage}

Qawluhu ʿazza kibriyāʾuhu: Huwa al-Nāṭiqu fī Malakūti al-Bayān.

His words---exalted be His majesty: He is the Speaker in the Kingdom of
Utterance.

\subsection{99}\label{section-2929}

\begin{otherlanguage}{persian}

امروز روز ذکر و بیان است و امروز روزی است که برهان و کینونت آن مابین
عالم ندا مینماید و میگوید ای اهل ارض اگر تشنه‌اید رحیق مختوم از اصبع قیّوم
گشوده و جاری خود را از او محروم مسازید اگر طالب ظهورید حضرت غیب مکنون
ظاهر و مشهود و اگر طالب آیاتید قد ملئت منها الآفاق و اگر بیّنات میطلبید
انّها ظهرت علی شأن لا ینکرها الّا کلّ معتد اثیم بگو ایدوستان آفتاب حقیقت از
افق عالم مشرق و بحر علم امام وجوه ظاهر خود را محروم ننمائید و از فیوضات
نامتناهی الهی ممنوع مسازید امروز روز بصر است چه که وجه اللّه از افق اعلی
مشرق و ظاهر و یوم سمع است چه که ندای حقّ جلّ جلاله مرتفع ای اهل بها سفینهٴ
حمرا در خلیج بحر اعظم منتظر شما است توجّهوا الیها ثم ارکبوا علیها باسم
اللّه مالک الأسمآء و فاطر السّمآء ای جلیل انشآءاللّه از نار خلیل نفوس
افسرده را مشتعل نمائی و عباد تشنه را باسم حقّ کوثر زندگانی عطا کنی
باستقامت تمام بر خدمت امر مالک انام قیام نما و بذکرش افئده و قلوبرا مسخّر
کن تا کلّ فائز شوند بآنچه سبب و علّت خلقت عالمیان است

\end{otherlanguage}

Emrūz rūz-e dhekr va bayān ast va emrūz rūzī ast ke borhān va keynūnat-e
ān mābeyn-e ʿālam nedā mīnamāyad va mīgūyad: Ey ahl-e arḍ, agar
teshne-īd raḥīq-e makhtūm az eṣbaʿ-e Qayyūm goshūde va jārī, khod rā az
ū maḥrūm masāzīd. Agar ṭāleb-e ẓohūrīd, Ḥaḍrat-e Gheyb-e Maknūn ẓāher va
mashhūd, va agar ṭāleb-e āyātīd, qad muliʾat minhā al-āfāq, va agar
bayyenāt mīṭalabīd, innahā ẓaharat ʿalā shaʾnin lā yunkiruhā illā kullu
muʿtadin athīm. Begū ey dūstān, āftāb-e ḥaqīqat az ofoq-e ʿālam moshreq
va baḥr-e ʿelm amām-e vojūh ẓāher. Khod rā maḥrūm nanamāyīd va az
foyūḍāt-e nāmotanāhī-ye Elāhī mamnūʿ masāzīd. Emrūz rūz-e baṣar ast che
ke Wajh Allāh az ofoq-e aʿlā moshreq va ẓāher, va yowm-e samʿ ast che ke
nedā-ye Ḥaqq jalla jalāluhu mortafeʿ. Ey ahl-e Bahā, safīne-ye ḥamrā dar
khalīj-e baḥr-e aʿẓam montaẓer-e shomā ast. Tavajjahū ilayhā thumma
irkabū ʿalayhā bi-smi Allāhi Māliki al-Asmāʾ wa-Fāṭiri al-Samāʾ. Ey
Jalīl, in shāʾ Allāh az nār-e Khalīl nofūs-e afsorde rā moshtaʿel namāʾī
va ʿebād-e teshne rā be-esm-e Ḥaqq kowthar-e zendegānī ʿaṭā konī.
Be-esteqāmat-e tamām bar khedmat-e Amr-e Mālek-e anām qiyām namā va
be-dhekrash afʾede va qolūb rā mosakhkhar kon tā koll fāʾez shavand
be-ānche sabab va ʿellat-e khelqat-e ʿālamiyān ast.

Today is the day of remembrance and utterance, and today is the day when
proof and its very essence call out amid the world, saying: ``O people
of the earth! If you are thirsty, the sealed wine has been opened by the
finger of the Self-Subsisting and is flowing; do not deprive yourselves
of it. If you seek revelation, the hidden, unseen Presence is manifest
and visible. If you seek verses, the horizons have been filled with
them. And if you seek clear evidences, they have appeared in such a
manner that none denies them except every sinful transgressor.'' Say: O
friends! The sun of truth shines from the horizon of the world, and the
ocean of knowledge lies revealed before your faces. Do not deprive
yourselves or bar yourselves from the infinite outpourings of divine
grace. Today is the day of sight, for the Face of God shines forth and
is manifest from the most exalted horizon; and it is the day of hearing,
for the Call of the True One---glorified be His glory---has been raised.
O people of Bahā! The Crimson Ark awaits you in the bay of the Most
Great Ocean. Turn toward it, then board it in the name of God, the Lord
of Names and Creator of heaven. O Jalīl! God willing, you will kindle
despondent souls with the fire of Abraham and, in the name of the True
One, bestow the Kawthar of life upon His thirsty servants. Arise with
complete steadfastness to serve the Cause of the Lord of humanity, and
conquer hearts and souls through His remembrance, so that all may attain
that which is the cause and purpose of the creation of the peoples of
the world.

\subsection{100}\label{section-2930}

\begin{otherlanguage}{arabic}

قل سبحانک یا الهی و محبوبی و مقصودی و مالک ناصیتی و المهیمن علی روحی و
نفسی و ذاتی و هیکلی اسئلک بقلزم بیانک و انوار شمس وجهک بأن {[}٠٠٠{]}
ذکری و ندائی مؤثّراً فی افئدة عبادک و کوثراً حیواناً لأموات اهل مملکتک حرّکهم
یا الهی بأیادی فضلک ثمّ ارزقهم سلسبیل عرفانک ای ربّ انت الّذی سبق فضلک و
احاط کرمک و وسعت رحمتک خذ ایادی عبادک بعنایتک ثمّ انقذهم من غمرات النّفس و
الهوی انّک انت مولی الوری و مالک الآخرة و الأولی لا اله الّا انت الغفور
الکریم انتهی

\end{otherlanguage}

Qul subḥānaka yā Ilāhī wa-Maḥbūbī wa-Maqṣūdī wa-Mālika nāṣiyatī
wa-l-Muhaymina ʿalā rūḥī wa-nafsī wa-dhātī wa-haykalī. Asʾaluka
bi-qulzumi bayānika wa-anwāri shamsi wajhika bi-an {[}corrupted text{]}
dhikrī wa-nidāʾī muʾaththiran fī afʾidati ʿibādika wa-kawtharan
ḥayawānan li-amwāti ahli mamlakatik. Ḥarrikhum yā Ilāhī bi-ayādī faḍlika
thumma urzuqhum salsabīla ʿirfānik. Ay Rabb, anta alladhī sabaqa faḍluka
wa-aḥāṭa karamuka wa-wasiʿat raḥmatuk. Khudh ayādiya ʿibādika
bi-ʿināyatika thumma anqidhhum min ghamarāti al-nafsi wa-l-hawā. Innaka
anta Mawlā al-warā wa-Māliku al-ākhirati wa-l-ūlā. Lā ilāha illā anta
al-Ghafūru al-Karīm. Intahā.

Say: Glory be to You, O my God, my Beloved, my Heart's Desire, the One
who holds me by the forelock and watches over my spirit, my soul, my
being, and my body. I ask You, by the ocean of Your utterance and the
radiance of the sun of Your face, to {[}corrupted original text{]} my
remembrance and my call effective in the hearts of Your servants and a
life-giving Kawthar for the dead among the people of Your kingdom. Stir
them, O my God, with the hands of Your grace, then grant them the
Salsabīl of Your recognition. O Lord! You are the One whose grace has
preceded all, whose bounty has encompassed all, and whose mercy has
embraced all. Take the hands of Your servants in Your loving care, then
rescue them from the depths of self and desire. You are indeed the Lord
of all creation and the Ruler of the world to come and this world. There
is no God but You, the Ever-Forgiving, the All-Bountiful. End.

\subsection{101}\label{section-2931}

\begin{otherlanguage}{arabic}

انشآءاللّه باین لوح امنع اقدس مردگان وادی غفلت زنده گردند و بحیوة تازه
فائز شوند هر نفسی الیوم للّه باعتدال تکلّم نماید البتّه بیانش مؤثر و ذکرش
مهیمن و اراده‌اش غالب شود کلّ امر عزیز علینا و لیس علی اللّه بعزیز خدمت
جناب ایشان تکبیر میرسانم و کمال خلوص اظهار میدارم و از حقّ منیع میطلبم تا
از ایشان ظاهر فرماید آنچه را که ذکرش بطراز خلود فائز شود و در دفتر عالم
مخلّد بماند انّه لهو المقتدر القدیر لا اله الّا هو العلیم الحکیم

\end{otherlanguage}

In shāʾ Allāh be īn lawḥ-e amnaʿ-e aqdas mordegān-e vādī-ye ghaflat
zende gardand va be ḥayāt-e tāze fāʾez shavand. Har nafsī al-yawm
li-llāh be eʿtedāl takallom namāyad al-batte bayānash moʾathther va
dhekrash mohaymen va erāde-ash ghāleb shavad. Kullu amrin ʿazīzun
ʿalaynā wa-laysa ʿalā Allāhi bi-ʿazīz. Khedmat-e janāb-e īshān takbīr
mīresānam va kamāl-e kholūṣ eẓhār mīdāram va az Ḥaqq-e manīʿ mīṭalabam
tā az īshān ẓāher farmāyad ānche rā ke dhekrash be ṭerāz-e kholūd fāʾez
shavad va dar daftar-e ʿālam mokhallad bemānad. Innahu la-huwa
al-Muqtadir al-Qadīr lā ilāha illā huwa al-ʿAlīm al-Ḥakīm.

God willing, through this most inaccessible, most holy Tablet, the dead
in the valley of heedlessness will come to life and attain a new life.
Whoever speaks today with moderation for the sake of God will assuredly
find that their utterance is effective, their remembrance commanding,
and their purpose triumphant. Every undertaking is difficult for us, but
nothing is difficult for God. I convey to him the greeting ``God is most
great'' and express my utmost sincerity. I ask the exalted True One to
bring forth from him that which will adorn his remembrance with
immortality and preserve it forever in the record of the world. He is
indeed the All-Powerful, the Omnipotent. There is no God but Him, the
All-Knowing, the All-Wise.

\subsection{102}\label{section-2932}

\begin{otherlanguage}{persian}

اینکه مرقوم داشته بودند در یوم جمعه بیت جناب آقا میرزا ع‌ب علیه بهآء اللّه
بحضور دوستان فائز شد نعیماً لأهل البهآء الّذین اخذتهم حرارة محبّة اللّه علی
شأن یرون المقرّبون من وجوههم نضرة الرّحمن و یسمعون من السنهم ثنائه الجمیل
این فقره در ساحت امنع اقدس عرض شد و این کلمهٴ محکمهٴ مبارکه از سماء
عنایت نازل قوله عزّ کبریائه

\end{otherlanguage}

Īnke marqūm dāshte būdand dar yawm-e jomʿe bayt-e janāb-e Āqā Mīrzā ʿ-B
ʿalayhi bahāʾ Allāh be ḥożūr-e dūstān fāʾez shod. Naʿīman li-ahl
al-Bahāʾ alladhīna akhadhat-hum ḥarārat maḥabbat Allāh ʿalā shaʾn
yarawna al-muqarrabūn min wujūhihim naḍrat al-Raḥmān wa-yasmaʿūna min
alsinatihim thanāʾahu al-jamīl. Īn feqre dar sāḥat-e amnaʿ-e aqdas ʿarż
shod va īn kalame-ye moḥkame-ye mobārake az samā-ye ʿenāyat nāzel,
qawluhu ʿazza kibriyāʾuhu:

Concerning what he had written, that on Friday the house of the honored
Āqā Mīrzā ʿ-B, upon him be the glory of God, was blessed by the presence
of the friends: bliss be upon the people of Bahā, whom the warmth of the
love of God has so seized that those who are near to Him see the
radiance of the All-Merciful in their faces and hear His beautiful
praise from their tongues. This matter was presented in the most
inaccessible, most holy presence, and this firm and blessed utterance
descended from the heaven of grace. These are His words, exalted be His
majesty:

\subsection{103}\label{section-2933}

\begin{otherlanguage}{persian}

نفوسیکه الیوم بذکر اللّه ناطقند و بر امر الهی ثابت و مستقیم یصلّین علیهم
الملأ الأعلی فی العشیّ و الاشراق طوبی لبیت تنوّر بأنوار ذکری و ارتفع فیه
ندآء اولیائی الّذین نبذوا سوائی و اتّبعوا ما امروا به فی کتابی المنیر جناب
ع‌ب علیه بهائی الحمد للّه بذکر حقّ و عنایت حقّ و خدمت حقّ فائزند اینمقام
عنداللّه افضل است از آنچه دیده شده انشآءاللّه همیشه بآن فائز باشند ولکن
بحکمت

\end{otherlanguage}

Nofūsī ke al-yawm be dhekr-e Allāh nāṭeqand va bar amr-e elāhī thābet va
mostaqīm, yuṣallīna ʿalayhim al-malaʾ al-aʿlā fī al-ʿashiyy wa-l-ishrāq.
Ṭūbā li-baytin tanawwara bi-anwār dhikrī wa-rtafaʿa fīhi nidāʾ awliyāʾī
alladhīna nabadhū siwāʾī wa-ttabaʿū mā umirū bihi fī kitābī al-munīr.
Janāb-e ʿ-B ʿalayhi bahāʾī al-ḥamdu li-llāh be dhekr-e Ḥaqq va ʿenāyat-e
Ḥaqq va khedmat-e Ḥaqq fāʾezand. Īn maqām ʿinda Allāh afżal ast az ānche
dīde shode. In shāʾ Allāh hamīshe be ān fāʾez bāshand va-līkan be
ḥekmat.

Upon those who today speak in remembrance of God and remain firm and
steadfast in the divine Cause, the Concourse on high invokes blessings
at evening and dawn. Blessed is a house illumined by the lights of My
remembrance, in which the call of My loved ones has been raised---those
who have cast aside all but Me and followed what they were commanded in
My luminous Book. The honored ʿ-B, upon him be My glory, has, praise be
to God, attained remembrance of the True One, His favor, and service to
Him. In God's sight, this station surpasses whatever has been seen. God
willing, may he always attain it, but with wisdom.

\subsection{104}\label{section-2934}

\begin{otherlanguage}{persian}

یا ورقآء یا ایّها النّاظر الی الأفق الأعلی در جمیع احوال بحکمت تمسّک نما و
متشبّث شو اجماع دوستان که از حدّ اعتدال تجاوز نماید سبب اضطراب نفوس خواهد
شد چه که اهل ارض و علما و امرای آن تا این حین از ارادهٴ حقّ جلّ جلاله
اطّلاع نیافته‌اند و باغوای بعضی از جهلا که خود را علمای ارض میدانند مطلع
وحی و مشرق الهام را سبب و علّت فتنه و مصدر فساد دانسته و میدانند باین جهت
اجماع دوستان الهی را سبب تفریق نفوس و تشتّت ما عندهم شمرده‌اند لذا باید در
جمیع امور باعتدال ناظر بود لضعف نفوس غافله قسم بحقیقت سدرهٴ منتهی که اگر
اقلّ از سمّ ابره اطّلاع یابند و عارف شوند کلّ بأجنحهٴ شوق و اشتیاق طیران
نمایند و قصد مقصد اقصی و ذروهٴ علیا و غایت قصوی کنند

\end{otherlanguage}

Yā Warqāʾ yā ayyuhā al-nāẓir ilā al-ufuq al-aʿlā, dar jamīʿ-e aḥvāl be
ḥekmat tamassok namā va motashabbeth show. Ejtemāʿ-e dūstān ke az ḥadd-e
eʿtedāl tajāvoz namāyad sabab-e eżṭerāb-e nofūs khāhad shod, che ke
ahl-e arż va ʿolamā va omarā-ye ān tā īn ḥīn az erāde-ye Ḥaqq jalla
jalāluhu eṭṭelāʿ nayāfte-and va be eghvā-ye baʿżī az johalā ke khod rā
ʿolamā-ye arż mīdānand maṭlaʿ-e vaḥy va mashreq-e elhām rā sabab va
ʿellat-e fetne va maṣdar-e fasād dāneste va mīdānand. Be īn jahat
ejtemāʿ-e dūstān-e elāhī rā sabab-e tafrīq-e nofūs va tashattot-e mā
ʿindahum shemorde-and. Ledhā bāyad dar jamīʿ-e omūr be eʿtedāl nāẓer būd
li-ḍaʿf-e nofūs-e ghāfele. Qasam be ḥaqīqat-e Sedre-ye Montahā ke agar
aqall az samm-e ebre eṭṭelāʿ yāband va ʿāref shavand koll be ajneḥe-ye
showq va eshtiyāq ṭayarān namāyand va qaṣd-e maqṣad-e aqṣā va dhorve-ye
ʿolyā va ghāyat-e qoṣvā konand.

O Varqā! O you who gaze toward the most exalted horizon! Hold fast to
wisdom in every circumstance and cling to it. A gathering of the friends
that exceeds the bounds of moderation will cause people anxiety, for the
inhabitants of the earth, its religious scholars, and its rulers have
not yet become aware of the will of the True One, glorified be His
glory. Misled by certain ignorant people who consider themselves the
religious scholars of the earth, they have regarded, and still regard,
the Dawning-Place of revelation and the Dayspring of inspiration as the
cause of sedition and the source of disorder. For this reason, they have
reckoned gatherings of the friends of God to be a cause of division
among people and the disruption of what they possess. Thus, in all
matters, moderation must be kept in view because of the weakness of
heedless souls. I swear by the reality of the Tree beyond which there is
no passing: if they were to gain awareness and recognition to a degree
smaller than the eye of a needle, they would all soar on the wings of
yearning and longing and set out for the farthest destination, the
loftiest summit, and the ultimate goal.

\subsection{105}\label{section-2935}

\begin{otherlanguage}{persian}

آز اهل عالم را از حضرت بی‌نیاز محروم نموده نمیدانند چه شده و چه خبر است
چه که طالب آگاهی نبوده و نیستند در یکی از الواح فرمودیم چشم سر را پلک
بآن نازکی از دیدن جهان و آنچه در او است بی‌بهره نماید دیگر اگر پردهٴ آز
بر چشم دل فرود آید چه خواهد نمود عالم را آتش پندار تمام نموده ظنون و
اوهام عالم شعور انسانی را تیره و تار کرده امرای ارض الیوم بمذهب و ملّت
کسی فی‌الحقیقه متعرّض نبوده و نیستند ولکن از محلّی که او را سبب اجماع و
مراودهٴ ناس مشاهده مینمایند کمال احتراز داشته و دارند چه که از امثال این
امور در ارض از قبل انواع فتنه و فساد ظاهر شده و حال چون در ظاهر توجّه ناس
را باین امر مشاهده مینمایند قیاس بما قبل نموده بسدّ و منع قیام میکنند

\end{otherlanguage}

Āz ahl-e ʿālam rā az Ḥażrat-e Bī-Niyāz maḥrūm namūde. Nemīdānand che
shode va che khabar ast, che ke ṭāleb-e āgāhī nabūde va nīstand. Dar
yekī az alvāḥ farmūdīm chashm-e sar rā pelk be ān nāzokī az dīdan-e
jahān va ānche dar ū ast bī-bahre namāyad. Dīgar agar parde-ye āz bar
chashm-e del forūd āyad che khāhad namūd. ʿĀlam rā ātash-e pendār tamām
namūde, ẓonūn va owhām ʿālam-e shoʿūr-e ensānī rā tīre va tār karde.
Omarā-ye arż al-yawm be madhhab va mellat-e kasī fī al-ḥaqīqe motaʿarreż
nabūde va nīstand va-līkan az maḥallī ke ū rā sabab-e ejtemāʿ va
morāvede-ye nās moshāhede mīnamāyand kamāl-e eḥterāz dāshte va dārand,
che ke az amthāl-e īn omūr dar arż az qabl anvāʿ-e fetne va fasād ẓāher
shode va ḥāl chon dar ẓāher tavajjoh-e nās rā be īn amr moshāhede
mīnamāyand qiyās be mā qabl namūde be sadd va manʿ qiyām mīkonand.

Greed has deprived the people of the world of the Self-Sufficient One.
They do not know what has happened or what is taking place, for they
have not sought, and do not seek, awareness. In one of the Tablets We
said: ``An eyelid, thin as it is, deprives the physical eye of seeing
the world and all that is in it. What, then, will happen if the veil of
greed descends upon the eye of the heart?'' The fire of vain imagining
has consumed the world; conjectures and delusions have darkened and
obscured the realm of human consciousness. The rulers of the earth today
have not, in truth, interfered with anyone's religion or creed, nor do
they do so now. Yet they have been, and remain, extremely wary of any
place they see as a cause of people gathering and associating, for in
the past all manner of sedition and disorder have arisen on earth from
such things. Now, seeing people outwardly turn toward this Cause, they
draw an analogy with the past and arise to obstruct and prevent it.

\subsection{106}\label{section-2936}

\begin{otherlanguage}{persian}

از حقّ بخواه آن نفوس را بقطره‌ئی از دریای معرفت فائز فرماید و بنوری از
انوار آفتاب حکمت منوّر دارد تا کلّ از ارادات خود مقدّس شده بارادهٴ الهی
قیام نمایند و این امر را که مطلع زندگانی عالم و مشرق عمار ان است قیاس
بامورات قبل ننمایند و فساد و صلاح از هم امتیاز دهند حقّ جلّ جلاله قادر است
بر اینکه نفسی از امرا را مبعوث فرماید تا مطلّع شود و سائرین را آگاه نماید
لسان حکمت میفرماید هر یومی را بما یوافقه عطا نمائید مقتضیات هر حین
مشاهده کنید انّه یفعل الحقّ و یهدی السّبیل و الأمر للّه العزیز الجمیل

\end{otherlanguage}

Az Ḥaqq bekhāh ān nofūs rā be qaṭre-ī az daryā-ye maʿrefat fāʾez
farmāyad va be nūrī az anvār-e āftāb-e ḥekmat monavvar dārad tā koll az
erādāt-e khod moqaddas shode be erāde-ye elāhī qiyām namāyand va īn amr
rā ke maṭlaʿ-e zendegānī-ye ʿālam va mashreq-e {[}corrupted text{]} ast
qiyās be omūrāt-e qabl nanamāyand va fasād va ṣalāḥ az ham emtiyāz
dehand. Ḥaqq jalla jalāluhu qāder ast bar īnke nafsī az omarā rā mabʿūth
farmāyad tā moṭṭaleʿ shavad va sāʾerīn rā āgāh namāyad. Lesān-e ḥekmat
mīfarmāyad har yawmī rā be mā yovāfeqoh ʿaṭā namāyīd moqtażayāt-e har
ḥīn moshāhede konīd. Innahu yafʿalu al-ḥaqq wa-yahdī al-sabīl wa-l-amru
li-llāh al-ʿAzīz al-Jamīl.

Ask the True One to grant those souls a drop from the ocean of knowledge
and illumine them with a ray from the sun of wisdom, so that all may
become sanctified from their own desires and act according to the divine
will, neither comparing this Cause---which is the dawning-place of the
world's life and the dayspring of {[}corrupted original text{]}---with
the affairs of the past, nor failing to distinguish corruption from
rectitude. The True One, glorified be His glory, is able to raise up one
of the rulers to become informed and make the others aware. The Tongue
of Wisdom says: ``Give each day what accords with it; consider the
requirements of every moment.'' He indeed does what is right and guides
the way. Authority belongs to God, the Almighty, the All-Beautiful.

\subsection{107}\label{section-2937}

\begin{otherlanguage}{persian}

یا ورقآء از قبل گفته‌ام لا تصدّقوا کلّ قائل و لا تطمئنّوا من کلّ وارد و امثال
آن مع‌ذلک بعضی از دوستان بمجرّد ادّعاء نفسی قبول مینمایند و او را محرم
میشمرند و این خطا بوده و خواهد بود اگرچه الحمد للّه محافل و مجالس و اجماع
احبّاء مقدّس و منزّه است از فتنه و فساد ولکن اطّلاع بر جمعیّت و آمد و شد
دوستان الهی امری است مهیّج افئده و افکار غافلین

\end{otherlanguage}

Yā Warqāʾ az qabl gofte-am lā tuṣaddiqū kulla qāʾil wa-lā taṭmaʾinnū min
kulli wārid va amthāl-e ān. Maʿa dhālek baʿżī az dūstān be mojarrad-e
eddeʿā-ye nafsī qabūl mīnamāyand va ū rā maḥram mīshemorand va īn khaṭā
būde va khāhad būd. Agarche al-ḥamdu li-llāh maḥāfel va majāles va
ejtemāʿ-e aḥebbā moqaddas va monazzah ast az fetne va fasād va-līkan
eṭṭelāʿ bar jamʿiyyat va āmad va shod-e dūstān-e elāhī amrī ast
mohayyej-e afʾede va afkār-e ghāfelīn.

O Varqā! I have said before: ``Do not believe everyone who speaks, nor
place your trust in everyone who arrives,'' and other such admonitions.
Nevertheless, some of the friends accept a person merely on the strength
of a claim and regard that person as a confidant. This has been, and
will remain, a mistake. Although, praise be to God, the assemblies,
meetings, and gatherings of the loved ones are sanctified and free from
sedition and disorder, knowledge of the gatherings and the coming and
going of the friends of God stirs the hearts and minds of the heedless.

\subsection{108}\label{section-2938}

\begin{otherlanguage}{arabic}

چون از اصل امر مطّلع نیستند لذا از این امور مضطرب میشوند و در اطفای نور
الهی بقدر قوّه کوشش مینمایند اگرچه قادر نبوده و نیستند کذلک القینا علیک
ما یحکی عن فضلی و رحمتی و عنایتی و شفقتی و مواهبی لتشکر ربّک المشفق
الکریم انتهی

\end{otherlanguage}

Chon az aṣl-e amr moṭṭaleʿ nīstand ledhā az īn omūr możṭareb mīshavand
va dar eṭfā-ye nūr-e elāhī be qadr-e qovve kūshesh mīnamāyand agarche
qāder nabūde va nīstand. Kadhālika alqaynā ʿalayka mā yaḥkī ʿan faḍlī
wa-raḥmatī wa-ʿināyatī wa-shafaqatī wa-mawāhibī li-tashkura rabbaka
al-Mushfiq al-Karīm. Intahā.

Since they are unaware of the true nature of the Cause, these things
make them anxious, and they strive with all their might to extinguish
the divine light, although they have never been able to do so and remain
unable. Thus have We imparted to you that which tells of My grace, My
mercy, My loving care, My compassion, and My gifts, so that you may
thank your Lord, the Compassionate, the All-Bountiful. End.

\subsection{109}\label{section-2939}

\begin{otherlanguage}{persian}

انشآءاللّه این امور بتوجّه آن محبوب در آن اراضی بطراز اصلاح مزیّن شود در
حفظ الواح الهی سعی بلیغ لازم و جهد منیع واجب و هم‌چنین در ستر رسائل و
الواح که از ساحت اقدس باطراف میرود باید بقدر قوّه اهتمام نمود چه اگر
ضعیفی و یا منافقی اطّلاع یابد و ذکر نماید شاید ضرر آن بساحت اقدس راجع شود
و یا باولیای حقّ در اطراف در هر حال ستر و کتمان این امور لازم و از حزم
شمرده میشود انشآءاللّه آن محبوب از عنایت محبوب عالمیان موفّق و مؤیّدند بر
خدمت امر انّه ولیّ المخلصین

\end{otherlanguage}

In shāʾ Allāh īn omūr be tavajjoh-e ān maḥbūb dar ān arāżī be ṭerāz-e
eṣlāḥ mozayyan shavad. Dar ḥefẓ-e alvāḥ-e elāhī saʿy-e balīgh lāzem va
jahd-e manīʿ vājeb va hamchenīn dar satr-e rasāʾel va alvāḥ ke az
sāḥat-e aqdas be aṭrāf mīravad bāyad be qadr-e qovve ehtemām namūd, che
agar żaʿīfī va yā monāfeqī eṭṭelāʿ yābad va dhekr namāyad shāyad żarar-e
ān be sāḥat-e aqdas rājeʿ shavad va yā be awliyā-ye Ḥaqq dar aṭrāf. Dar
har ḥāl satr va ketmān-e īn omūr lāzem va az ḥazm shemorde mīshavad. In
shāʾ Allāh ān maḥbūb az ʿenāyat-e Maḥbūb-e ʿālamiyān movaffaq va
moʾayyadand bar khedmat-e amr. Innahu waliyy al-mukhliṣīn.

God willing, through the attention of that loved one, these matters will
be set right in those regions. Great care and earnest effort are
required to safeguard the divine Tablets. Likewise, every possible
effort must be made to keep concealed the letters and Tablets sent from
the Most Holy Court to various places, for if a weak or hypocritical
person learns of them and speaks of them, the resulting harm may reach
the Most Holy Court or the loved ones of the True One elsewhere. In all
circumstances, keeping these matters hidden and confidential is
necessary and is considered prudent. God willing, through the grace of
the Beloved of all the worlds, that loved one will be assisted and
confirmed in service to the Cause. He is indeed the Guardian of the
sincere.

\subsection{110}\label{section-2940}

\begin{otherlanguage}{persian}

و اینکه در فقرهٴ تأویل که در کتاب اقدس از قلم اعلی نازل شده سؤال
{[}نمودید{]} در ساحت امنع اقدس عرض شد هذا ما نطق به لسان العظمة

\end{otherlanguage}

Va īnke dar feqre-ye taʾvīl ke dar Ketāb-e Aqdas az Qalam-e Aʿlā nāzel
shode soʾāl {[}namūdīd{]} dar sāḥat-e amnaʿ-e aqdas ʿarż shod. Hādhā mā
naṭaqa bihi lisān al-ʿaẓama:

Your question concerning the passage on interpretation revealed by the
Most Exalted Pen in the Most Holy Book was presented in the most
inaccessible, most holy presence. This is what the Tongue of Grandeur
uttered:

\subsection{111}\label{section-2941}

\begin{otherlanguage}{persian}

مقصود از تأویل اینکه از ظاهر خود را محروم ننمایند و از مقصود محتجب
نمانند مثلاً اگر از سماء مشیّت فاغسلوا وجوهکم نازل شود تأویل ننمایند که
مقصود از غسل غسل وجه باطن است و باید بآب عرفان او را غسل داد و طاهر نمود
و امثال آن بسا میشود نفسی باین تأویلات وجهش با کمال ذفر و وسخ آلوده
میماند و بخیال خود باصل امر اللّه عمل نموده و حال اینکه در این مقام واضح
و معلوم است که شستن روست بآب ظاهر

\end{otherlanguage}

Maqṣūd az taʾvīl īnke az ẓāher khod rā maḥrūm nanamāyand va az maqṣūd
moḥtajeb namānand. Mathalan agar az samā-ye mashiyyat fa-ghsilū
wujūhakum nāzel shavad taʾvīl nanamāyand ke maqṣūd az ghosl ghosl-e
vajh-e bāṭen ast va bāyad be āb-e ʿerfān ū rā ghosl dād va ṭāher namūd
va amthāl-e ān. Basā mīshavad nafsī be īn taʾvīlāt vajhash bā kamāl-e
dhafar va vasakh ālūde mīmānad va be khiyāl-e khod be aṣl-e amr Allāh
ʿamal namūde va ḥāl ānke dar īn maqām vāżeḥ va maʿlūm ast ke shostan-e
rūst be āb-e ẓāher.

What is meant concerning interpretation is that people should not
deprive themselves of the outward meaning or remain veiled from the
intended meaning. For example, if ``Wash your faces'' descends from the
heaven of the divine will, they should not interpret it to mean that
washing signifies washing the inner face, which must be washed and
purified with the water of spiritual knowledge, and so forth. It often
happens that, through such interpretations, a person's face remains
thoroughly foul-smelling and dirty, while that person imagines they have
carried out the very essence of God's command. Yet it is clear and
evident that what is meant here is washing the face with physical water.

\subsection{112}\label{section-2942}

\begin{otherlanguage}{persian}

بعض کلمات الهی را میتوان تأویل نمود یعنی تأویلاتی که سبب و علّت ظنون و
اوهام نشود و از مقصود الهی محروم نماند در ما انزله الرّحمن فی الفرقان
تفکّر نما قوله تعالی و من یؤت الحکمة فقد اوتی خیراً کثیراً بعضی برآنند که
مقصود از این حکمت احکام الهیّه است که در کتاب نازل و برخی برآنند که این
حکمت علم طبّ است و هر نفسی بآن فائز شد بخیر کثیر فائز است چه که این متعلّق
بانسان است و علم ابدان و این علم اشرف از سایر علوم است چنانچه از قبل
لسان حکمت باینکلمهٴ علیا نطق نموده العلم علمان علم الأبدان و علم الأدیان
علم ابدان را در کلمهٴ مبارکه مقدّم داشته و فی الحقیقه مقدّم است چه که ظهور
حقّ و احکام الهی جمیع از برای تربیت انسان و ترقّی او و حفظ اهل عالم و
امثال آن بوده و خواهد بود لذا آنچه سبب و علّت حفظ و صحّت و سلامتی وجود
انسان است مقدّم بوده و خواهد بود

\end{otherlanguage}

Baʿż-e kalemāt-e elāhī rā mītavān taʾvīl namūd yaʿnī taʾvīlātī ke sabab
va ʿellat-e ẓonūn va owhām nashavad va az maqṣūd-e elāhī maḥrūm namānad.
Dar mā anzalahu al-Raḥmān fī al-Furqān tafakkor namā, qawluhu taʿālā:
wa-man yuʾta al-ḥikmata fa-qad ūtiya khayran kathīran. Baʿżī bar ānand
ke maqṣūd az īn ḥekmat aḥkām-e elāhiyye ast ke dar ketāb nāzel va barkhī
bar ānand ke īn ḥekmat ʿelm-e ṭebb ast va har nafsī be ān fāʾez shod be
khayr-e kathīr fāʾez ast, che ke īn motaʿalleq be ensān ast va ʿelm-e
abdān va īn ʿelm ashraf az sāyer-e ʿolūm ast. Chenānche az qabl lesān-e
ḥekmat be īn kalame-ye ʿolyā noṭq namūde: al-ʿilmu ʿilmān ʿilm al-abdān
wa-ʿilm al-adyān. ʿElm-e abdān rā dar kalame-ye mobārake moqaddam dāshte
va fī al-ḥaqīqe moqaddam ast che ke ẓohūr-e Ḥaqq va aḥkām-e elāhī jamīʿ
az barā-ye tarbiyat-e ensān va taraqqī-ye ū va ḥefẓ-e ahl-e ʿālam va
amthāl-e ān būde va khāhad būd. Ledhā ānche sabab va ʿellat-e ḥefẓ va
ṣeḥḥat va salāmatī-ye vojūd-e ensān ast moqaddam būde va khāhad būd.

Some divine words may be interpreted---that is, through interpretations
that do not give rise to conjectures and delusions or deprive one of the
divine intent. Reflect on what the All-Merciful revealed in the Qur'an,
exalted be His words: ``Whoever is given wisdom has indeed been given
abundant good.'' Some hold that this wisdom means the divine laws
revealed in the Book, while others hold that it is the science of
medicine, and that whoever attains it has attained abundant good, since
it concerns the human being and is the science of bodies, and this
science is nobler than the other sciences. Thus, in the past, the tongue
of wisdom uttered this exalted saying: ``Knowledge is of two kinds:
knowledge of bodies and knowledge of religions.'' In this blessed
saying, knowledge of bodies was given precedence, and in truth it has
precedence, for the revelation of the True One and the divine laws have
all been, and will remain, for the education and advancement of
humanity, the protection of the people of the world, and similar
purposes. Thus, whatever brings about the preservation, health, and
well-being of the human body has had, and will continue to have,
precedence.

\subsection{113}\label{section-2943}

\begin{otherlanguage}{persian}

و این فقره واضح و مشهود است و بعضی برآنند که حکمت معرفت حقایق اشیاء است
که هر نفسی باینمقام فائز شود دارای مقامات عالیه خواهد بود و حزبی برآنند
حکمت عمل به ما ینتفع به الانسان است هر که باینمقام موفّق شود او دارای خیر
کثیر است و جمهوری برآنند که حکمت آن است که انسانرا از ما یذلّه حفظ
مینماید و به ما یعزّه هدایت میکند و شرذمه‌ئی برآنند که حکمت علم معاشرت با
خلق است و مدارای با عباد علی شأن لا یتجاوز حدّ الاعتدال و بعضی گفته‌اند که
حکمت علم الهیّاتست که از قواعد حکمای قبل است و هم‌چنین گفته‌اند حکمت علم
جواهر و اعراض است و هم‌چنین علم هیئت و امثال آن و برخی گفته‌اند حکمت علم
بمعاش است در دنیا و شرذمه‌ئی قائلند بر اینکه حکمت علم باصول است و نفوسی
برآنند که حکمت در شئونات عدل است و آن اعطاء کلّ ذی حقّ حقّه بوده و خواهد
بود و برخی برآنند که علم اکسیر است هر نفسی بآن فائز شد بخیر کثیر فائز
است و بعضی برآنند که علم هندسه و امثال آن است و هر حزبی هم بقدر ادراک
خود برهان و دلیل از برای هر یک از آنچه ذکر شد گفته‌اند که اگر بتفصیل این
مقامات ذکر شود بر کسالت بیفزاید سبحان من نطق بهذه الکلمة العلیا رأس
الحکمة مخافة اللّه چه مخافة اللّه و خشیته انسانرا منع مینماید از آنچه سبب
ذلّت و پستی مقام انسان است و تأیید مینماید او را بر آنچه سبب علوّ و سموّ
است انسان عاقل از اعمال شنیعه اجتناب مینماید چه که مجازات را از پی
مشاهده میکند خیمهٴ نظم عالم بدو ستون قائم و برپا مکافات و مجازات چندی
قبل این آیهٴ مبارکه از مشرق فم الهی اشراق نمود للعدل جند و هی مجازات
الأعمال و مکافاتها و بهما ارتفع خبآء النّظم فی العالم و اخذ کلّ طاغ زمام
نفسه من خشیة الجزآء کذلک نطق مالک الأسمآء انّه لهو النّاطق العلیم

\end{otherlanguage}

Va īn faqara vāżeḥ va mashhūd ast va baʿżī bar ānand ke ḥekmat
maʿrefat-e ḥaqāyeq-e ashyāʾ ast ke har nafsī be īn maqām fāʾez shavad
dārā-ye maqāmāt-e ʿāliya khvāhad būd va ḥezbī bar ānand ḥekmat ʿamal be
mā yantafiʿu bihi al-insān ast har ke be īn maqām movaffaq shavad ū
dārā-ye khayr-e kathīr ast va jomhūrī bar ānand ke ḥekmat ān ast ke
ensān rā az mā yudhilluhu ḥefẓ mīnamāyad va be mā yuʿizzuhu hedāyat
mīkonad va sherdhemaʾī bar ānand ke ḥekmat ʿelm-e moʿāsherat bā khalq
ast va modārā-ye bā ʿebād ʿalā shaʾnin lā yata jāwazu ḥadda al-iʿtidāl
va baʿżī gofteʾand ke ḥekmat ʿelm-e elāhīyāt ast ke az qavāʿed-e
ḥokamā-ye qabl ast va hamchenīn gofteʾand ḥekmat ʿelm-e javāher va aʿrāż
ast va hamchenīn ʿelm-e heyʾat va amthāl-e ān va barkhī gofteʾand ḥekmat
ʿelm be maʿāsh ast dar donyā va sherdhemaʾī qāʾeland bar īnke ḥekmat
ʿelm be oṣūl ast va nofūsī bar ānand ke ḥekmat dar shoʾūnāt-e ʿadl ast
va ān iʿṭāʾu kulli dhī ḥaqqin ḥaqqahu būde va khvāhad būd va barkhī bar
ānand ke ʿelm-e eksīr ast har nafsī be ān fāʾez shod be khayr-e kathīr
fāʾez ast va baʿżī bar ānand ke ʿelm-e hendesa va amthāl-e ān ast va har
ḥezbī ham be qadr-e edrāk-e khvod borhān va dalīl az barā-ye har yek az
ānche dhekr shod gofteʾand ke agar be tafṣīl īn maqāmāt dhekr shavad bar
kesālat bīafzāyad subḥāna man naṭaqa bi-hādhihi al-kalimati al-ʿulyā
raʾsu al-ḥikmati makhāfatu Allāh che makhāfat Allāh va khashyatahu ensān
rā manʿ mīnamāyad az ānche sabab-e dhellat va pastī-ye maqām-e ensān ast
va taʾyīd mīnamāyad ū rā bar ānche sabab-e ʿolovv va somovv ast ensān-e
ʿāqel az aʿmāl-e shanīʿa ejtenāb mīnamāyad che ke mojāzāt rā az pey
moshāheda mīkonad kheyma-ye naẓm-e ʿālam be do sotūn qāʾem va barpā
mokāfāt va mojāzāt chandī qabl īn āya-ye mobāraka az mashreq-e fam-e
elāhī eshrāq namūd li-l-ʿadli jundun wa-hiya mujāzātu al-aʿmāli
wa-mukāfaʾātuhā wa-bihimā irtafaʿa khibāʾu al-naẓmi fī al-ʿālami
wa-akhadha kullu ṭāghin zimāma nafsihi min khashyati al-jazāʾ kadhālika
naṭaqa māliku al-asmāʾ innahu la-huwa al-nāṭiqu al-ʿalīm.

This matter is clear and evident. Some hold that wisdom is knowledge of
the realities of things, and that whoever attains this station will
possess exalted stations. Another group holds that wisdom is acting upon
what benefits humanity: whoever succeeds in attaining this station
possesses abundant good. A great many hold that wisdom is what protects
a person from whatever abases them and guides them to whatever exalts
them. A small group holds that wisdom is the knowledge of how to
associate with people and treat God's servants with forbearance, without
exceeding the bounds of moderation. Some have said that wisdom is the
science of theology, which is among the established disciplines of
earlier philosophers. They have also said that wisdom is the knowledge
of substances and accidents, as well as astronomy and similar sciences.
Others have said that wisdom is the knowledge of how to earn a
livelihood in this world. A small group maintains that wisdom is
knowledge of first principles. Some hold that wisdom consists in the
manifestations of justice, which has always meant, and will continue to
mean, giving everyone their due. Others hold that it is the science of
the elixir: whoever attains it has attained abundant good. Some hold
that it is geometry and similar sciences. Each group, according to its
own understanding, has also offered proofs and arguments for each of the
views mentioned. Were these matters to be set out in detail, it would
only increase weariness. Glorified be He who uttered this exalted
saying: ``The beginning of wisdom is the fear of God.'' For the fear and
awe of God restrain a person from whatever brings abasement and
degradation to the human station, and strengthen them in whatever brings
exaltation and elevation. A wise person avoids shameful deeds, for they
see the punishment that follows them. The tent of order in the world
stands firmly upon two pillars: reward and punishment. Some time ago,
this blessed verse shone forth from the dawning-place of the divine
mouth: ``Justice has a host, namely the punishment of deeds and their
reward. Through these two, the tent of order has been raised in the
world, and every transgressor has taken hold of the reins of his own
self for fear of retribution. Thus has the Lord of Names spoken. He is
indeed the Speaker, the All-Knowing.''

\subsection{114}\label{section-2944}

\begin{otherlanguage}{persian}

از قبل بعضی نفوس خود را اهل طریقت نامیده‌اند چنانچه آثار آن نفوس حال در
ارض موجود است طوائف مختلفه هر یک طریقی اخذ کرده‌اند اسلام از آن نفوس
متفرّق شد و از بحر اعظم خلیجها خارج نمودند تا آنکه بالأخره بر دین اللّه
وارد شد آنچه وارد شد اسلام با آن قوّت عظیمه از اعمال و افعال منتسبین او
بکمال ضعف مبدّل شد چنانچه مشاهده نموده و مینمایند بعضی از نفوس که خود را
دراویش مینامند جمیع احکام و اوامر الهی را تأویل نمودند اگر گفته شود صلوة
از احکام محکمهٴ الهی است میگویند صلوة بمعنی دعا است و ما در حین تولّد
بدعا آمده‌ایم و صلوة حقیقی را عمل نموده‌ایم و این بیچاره از ظاهر محروم است
تا چه رسد بباطن اوهامات نفوس غافله زیاده از حدّ احصا بوده و هست

\end{otherlanguage}

Az qabl baʿżī nofūs khvod rā ahl-e ṭarīqat nāmīdeʾand chenānche āthār-e
ān nofūs ḥāl dar arż mowjūd ast ṭavāʾef-e mokhtalefa har yek ṭarīqī
akhdh kardeʾand eslām az ān nofūs motafarreq shod va az baḥr-e aʿẓam
khalījhā khārej namūdand tā ānke bi-l-ākhara bar dīn Allāh vāred shod
ānche vāred shod eslām bā ān qovvat-e ʿaẓīma az aʿmāl va afʿāl-e
montasebīn-e ū be kamāl-e żaʿf mobaddal shod chenānche moshāheda namūde
va mīnamāyand baʿżī az nofūs ke khvod rā darāvīsh mīnāmand jamīʿ-e aḥkām
va avāmer-e elāhī rā taʾvīl namūdand agar gofte shavad ṣalāt az aḥkām-e
moḥkama-ye elāhī ast mīgūyand ṣalāt be maʿnī-ye doʿā ast va mā dar ḥīn-e
tavallod be doʿā āmadeʾīm va ṣalāt-e ḥaqīqī rā ʿamal namūdeʾīm va īn
bīchāra az ẓāher maḥrūm ast tā che rasad be bāṭen owhāmāt-e nofūs-e
ghāfela ziyāda az ḥadd-e eḥṣā būde va hast.

In the past, some people called themselves followers of the mystical
path, and their works still exist in the world today. Various sects each
adopted a different path. Through these people, Islam became divided,
and they drew channels away from the Most Great Ocean, until at last
there befell the religion of God what befell it. Islam, with all its
immense power, was reduced to utter weakness through the deeds and
actions of those who claimed allegiance to it, as people have witnessed
and continue to witness. Some who call themselves dervishes have
interpreted away all the divine laws and commandments. If they are told
that obligatory prayer is one of God's firmly established laws, they
say, ``Obligatory prayer means supplication, and at birth we came into
the world supplicating, and have thus performed the true obligatory
prayer.'' Yet this hapless person is deprived of the outward observance,
let alone the inward reality. The delusions of heedless souls have
always been, and remain, beyond reckoning.

\subsection{115}\label{section-2945}

\begin{otherlanguage}{persian}

باری از تنبلی و کسالت جمیع اوامر الهی را که بمثابهٴ سدّ محکم است از برای
حفظ عالم و امنیّت آن تأویل نمودند و در تکایا انزوا جسته جز خورد و خواب
شغلی اختیار ننموده‌اند و در معارف آنچه گفته شود تصدیق مینمایند ولکن اثر
حرارت محبّة اللّه تا حال از ایشان دیده نشد مگر معدود قلیلی که حلاوت بیانرا
یافتند و فی‌الجمله بر مقصود از ظهور آگاه گشتند بعضی از این تأویلات که در
بیان حکمت ذکر شد هر یک در مقام خود صحیح است چه با اصول احکام الهی مخالف
نیست آیاتیکه در اوامر و نواهی الهی است مثل عبادات و دیات و جنایات و
امثال آن مقصود عمل بظاهر آیات بوده و خواهد بود

\end{otherlanguage}

Bārī az tanbalī va kesālat jamīʿ-e avāmer-e elāhī rā ke be mathāba-ye
sadd-e moḥkam ast az barā-ye ḥefẓ-e ʿālam va amnīyat-e ān taʾvīl
namūdand va dar takāyā enzevā joste joz khvord va khvāb shoghlī ekhtiyār
nanamūdeʾand va dar maʿāref ānche gofte shavad taṣdīq mīnamāyand va
lāken athar-e ḥarārat-e maḥabbat Allāh tā ḥāl az īshān dīde nashod magar
maʿdūd-e qalīlī ke ḥalāvat-e bayān rā yāftand va fī al-jomla bar maqṣūd
az ẓohūr āgāh gashtand baʿżī az īn taʾvīlāt ke dar bayān-e ḥekmat dhekr
shod har yek dar maqām-e khvod ṣaḥīḥ ast che bā oṣūl-e aḥkām-e elāhī
mokhālef nīst āyātī ke dar avāmer va navāhī-ye elāhī ast methl-e ʿebādāt
va diyāt va jenāyāt va amthāl-e ān maqṣūd ʿamal be ẓāher-e āyāt būde va
khvāhad būd.

In short, out of laziness and indolence, they have interpreted away all
the divine commandments, which are like a mighty barrier for the
protection and security of the world. Secluding themselves in dervish
lodges, they have adopted no occupation but eating and sleeping. In
matters of spiritual knowledge, they assent to whatever is said, yet no
trace of the warmth of the love of God has so far been seen in them,
except for a very few who have tasted the sweetness of the utterance and
have gained some awareness of the purpose of the Revelation. Some of the
interpretations mentioned in the discussion of wisdom are each valid in
their own context, for they do not conflict with the principles of the
divine laws. As for verses concerning divine commandments and
prohibitions, such as acts of worship, blood money, criminal offenses,
and the like, the intention has always been, and will remain, that the
outward meaning of the verses be put into practice.

\subsection{116}\label{section-2946}

\begin{otherlanguage}{persian}

ولکن آیات الهی که در ذکر قیامت و ساعت چه در کتب قبل چه در فرقان نازل شده
اکثر مأوّل است و لا یعلم تأویله الّا اللّه این مراتب در کتاب ایقان واضح و
مبرهن است هر نفسی در آن تفکّر نماید آگاه شود بر آنچه از نظر کلّ مستور بوده
انظر فیما انزله الوهّاب فی الکتاب قوله تعالی لا الشّمس ینبغی لها ان تدرک
القمر و لا اللّیل سابق النّهار این آیهٴ مبارکه را علما از قبل و بعد تفسیر
و تأویل نموده‌اند و هر یک باعتقاد خود معنی آن را یافته‌اند و لکن این مقام
العلم و این مراتب الظّنون و الأوهام نفسیکه او را سلطان المفسّرین مینامند
یعنی قاضی بیضاوی گفته این آیه رداً لعبدة الشّمس نازل گشته مشاهده نمائید چه
مقدار بعید است از منبع علم مع تبحّره علی زعم النّاس زمام علوم در قبضهٴ
قدرت الهی است عطا میفرماید بهر که اراده نماید

\end{otherlanguage}

Va lāken āyāt-e elāhī ke dar dhekr-e qiyāmat va sāʿat che dar kotob-e
qabl che dar Forqān nāzel shode akthar moʾavval ast wa-lā yaʿlamu
taʾwīlahu illā Allāh īn marāteb dar Ketāb-e Īqān vāżeḥ va mobarhan ast
har nafsī dar ān tafakkor namāyad āgāh shavad bar ānche az naẓar-e koll
mastūr būde unẓur fīmā anzalahu al-Wahhābu fī al-kitāb qawluhu taʿālā lā
al-shamsu yanbaghī lahā an tudrika al-qamara wa-lā al-laylu sābiqu
al-nahār īn āya-ye mobāraka rā ʿolamā az qabl va baʿd tafsīr va taʾvīl
namūdeʾand va har yek be eʿteqād-e khvod maʿnī-ye ān rā yāfteʾand va
lāken ayna maqāmu al-ʿilm wa-ayna marātibu al-ẓunūni wa-al-awhām nafsī
ke ū rā solṭān al-mofasserīn mīnāmand yaʿnī Qāżī Beyżāvī gofte īn āya
raddan li-ʿabadati al-shams nāzel gashte moshāheda namāʾīd che meqdār
baʿīd ast az manbaʿ-e ʿelm maʿa tabaḥḥurihi ʿalā zaʿmi al-nās zemām-e
ʿolūm dar qabża-ye qodrat-e elāhī ast ʿaṭā mīfarmāyad be har ke erāda
namāyad.

However, most of the divine verses revealed concerning the Resurrection
and the Hour, whether in the earlier scriptures or in the Qur'an, are
subject to interpretation: ``None knows its interpretation but God.''
These matters are clear and demonstrated in the Book of Certitude.
Whoever reflects upon it will become aware of what has been hidden from
everyone's sight. Consider what the All-Bountiful has revealed in the
Book, exalted be His words: ``It is not for the sun to overtake the
moon, nor does the night outstrip the day.'' Scholars, both earlier and
later, have explained and interpreted this blessed verse, and each
believes that he has grasped its meaning. But how far the station of
knowledge is from the levels of conjecture and delusion! The man they
call the sovereign of commentators, namely Qāḍī Bayḍāwī, has said that
this verse was revealed as a refutation of sun worshippers. Observe how
far he is from the source of knowledge, despite his vast learning in the
estimation of the people. The reins of the sciences are held in the
grasp of divine power. He bestows them upon whomever He wills.

\subsection{117}\label{section-2947}

\begin{otherlanguage}{persian}

در کتب قبل هم آنچه در ذکر این ظهور از قلم اعلی جاری شد اکثر مأوّل است
مثلاً میفرماید تذوب الودیان کالشّمع قدّام الرّبّ و هم‌چنین میفرماید خیمهٴ او
بر جبال مرتفع میشود چنانچه شده و در مقامی دیگر میفرماید خیمه‌ایست از برای
او که طنابهای او بریده نشده و نمیشود میخهایش کنده نشده و نمیشود و از
جائی بجائی نقل نمیشود این فقرات مأوّل است انسان منصف ادراک مینماید که چه
مقام مقام تأویل است و چه مقام حکم بر ظاهر تنزیل

\end{otherlanguage}

Dar kotob-e qabl ham ānche dar dhekr-e īn ẓohūr az qalam-e aʿlā jārī
shod akthar moʾavval ast mathalan mīfarmāyad tadhūbu al-widyānu
ka-al-shamʿi quddāma al-rabb va hamchenīn mīfarmāyad kheyma-ye ū bar
jebāl mortafeʿ mīshavad chenānche shode va dar maqāmī dīgar mīfarmāyad
kheymaʾīst az barā-ye ū ke ṭanābhā-ye ū borīde nashode va namīshavad
mīkhhāyash kande nashode va namīshavad va az jāʾī be jāʾī naql
namīshavad īn faqarāt moʾavval ast ensān-e monṣef edrāk mīnamāyad ke che
maqām maqām-e taʾvīl ast va che maqām ḥokm bar ẓāher-e tanzīl.

Likewise, in the earlier scriptures, most of what flowed from the Most
Exalted Pen concerning this Revelation is subject to interpretation. For
example, He says, ``The valleys melt like wax before the Lord.'' He also
says, ``His tent will be raised upon the mountains,'' as has indeed
happened. Elsewhere He says, ``He has a tent whose cords have not been,
and will not be, cut, whose pegs have not been, and will not be,
uprooted, and which will not be moved from place to place.'' These
passages are subject to interpretation. A fair-minded person understands
where interpretation is called for and where the outward meaning of the
revealed text is binding.

\subsection{118}\label{section-2948}

\begin{otherlanguage}{arabic}

باری الیوم آنچه ناس را از آلایش پاک نماید و بآسایش حقیقی رساند آن مذهب
اللّه و دین اللّه و امر اللّه بوده کذلک انهمر غیث البیان من سمآء العرفان
فضلاً علیک اشکر و قل لک الحمد یا محبوب العارفین و لک الثّنآء یا مقصود
العالمین و لک الشّکر یا من فی قبضتک ملکوت ملک السّموات و الأرضین انتهی

\end{otherlanguage}

Bārī al-yowm ānche nās rā az ālāyesh pāk namāyad va be āsāyesh-e ḥaqīqī
rasānad ān madhhab Allāh va dīn Allāh va amr Allāh būde kadhālika
inhamara ghaythu al-bayāni min samāʾi al-ʿirfāni faḍlan ʿalayka ushkur
wa-qul laka al-ḥamdu yā maḥbūba al-ʿārifīn wa-laka al-thanāʾu yā maqṣūda
al-ʿālamīn wa-laka al-shukru yā man fī qabḍatika malakūtu mulki
al-samāwāti wa-al-araḍīn intahā.

In short, today, whatever cleanses people of defilement and brings them
true peace is the way of God, the religion of God, and the Cause of God.
Thus has the rain of utterance poured down from the heaven of divine
knowledge as a bounty to you. Give thanks and say, ``Praise be to You, O
Beloved of those who know You! To You be all acclaim, O Desire of the
worlds! To You be thanks, O You in whose grasp is the dominion of the
kingdom of the heavens and the earth!'' End.

\subsection{119}\label{section-2949}

\begin{otherlanguage}{persian}

از آنچه از ملکوت بیان رحمن نازل معلوم و واضح است که بعضی از آیات الهیّه
که در معاشرت و آداب و سلوک و معاملات و حدودات و امثال آن نازل اگر حکم از
ظاهر آن خارج شود امر اللّه معطّل و معوّق ماند یشهد بذلک کلّ عالم بصیر از حقّ
بطلبید اهل عالم را بطراز انصاف مزیّن فرماید بامورات جزئیه و خارق عادات
محدوده و آیات معدوده اظهار ایمان مینمایند و از دریاهای آن معرض مشاهده
میشوند انسان متحیّر است جلّ من ستر عنهم و کشف لنا

\end{otherlanguage}

Az ānche az malakūt-e bayān-e Raḥmān nāzel maʿlūm va vāżeḥ ast ke baʿżī
az āyāt-e elāhīya ke dar moʿāsherat va ādāb va solūk va moʿāmelāt va
ḥodūdāt va amthāl-e ān nāzel agar ḥokm az ẓāher-e ān khārej shavad amr
Allāh moʿaṭṭal va moʿavvaq mānad yashhadu bi-dhālika kullu ʿālimin baṣīr
az Ḥaqq beṭalabīd ahl-e ʿālam rā be ṭerāz-e enṣāf mozayyan farmāyad be
omūrāt-e jozʾīya va khāreq-e ʿādāt-e maḥdūda va āyāt-e maʿdūda eẓhār-e
īmān mīnamāyand va az daryāhā-ye ān moʿreż moshāheda mīshavand ensān
motaḥayyer ast jalla man satara ʿanhum wa-kashafa lanā.

From what has been revealed from the kingdom of the All-Merciful's
utterance, it is clear and evident that, in the case of certain divine
verses revealed concerning social relations, manners, conduct,
transactions, legal penalties, and the like, if their rulings are
removed from their outward meaning, the Cause of God will be brought to
a standstill and hindered. Every discerning scholar bears witness to
this. Beseech the True One to adorn the people of the world with the
ornament of fairness. They profess faith on the basis of trifling
matters, a limited number of extraordinary occurrences, and a handful of
signs, yet are seen to turn away from oceans of them. One is bewildered.
Exalted be He who has concealed from them and disclosed to us.

\subsection{120}\label{section-2950}

\begin{otherlanguage}{persian}

و اینکه از ماء بکر و حدّ مستعمل آن را سؤال نمودید تلقاء عرش عرض شد هذا ما
نطق به اللّه آب قلیل مثل یک کأس یا دو مقابل یا سه مقابل آن دست و رو در آن
شستن از مستعمل مذکور ولکن اگر بحدّ کر برسد از تغسیل یک وجه یا دو وجه
تغییر نمینماید و در استعمال آن بأسی نبوده و نیست و اگر یکی از اوصاف
ثلاثه در او ظاهر شود یعنی فی‌الجمله لون آب تغییر نماید از مستعمل محسوب
است یا ایّها الشّارب عن کأس عنایتی مقصود لطافت است از حوضهای {[}عجم{]}
مطّلعید هر قدر همّت بر لطافت مصروف شود عنداللّه محبوب است انتهی

\end{otherlanguage}

Va īnke az māʾ-e bekr va ḥadd-e mostaʿmal-e ān rā soʾāl namūdīd telqāʾ-e
ʿarsh ʿarż shod hādhā mā naṭaqa bihi Allāh āb-e qalīl methl-e yek kās yā
do moqābel yā se moqābel-e ān dast va rū dar ān shostan az mostaʿmal
madhkūr va lāken agar be ḥadd-e korr beresad az taghsīl-e yek vajh yā do
vajh taghyīr namīnamāyad va dar esteʿmāl-e ān baʾsī nabūde va nīst va
agar yekī az owṣāf-e thalātha dar ū ẓāher shavad yaʿnī fī al-jomla
lown-e āb taghyīr namāyad az mostaʿmal maḥsūb ast yā ayyuhā al-shāribu
ʿan kaʾsi ʿināyatī maqṣūd leṭāfat ast az ḥowżhā-ye {[}ʿAjam{]}
moṭṭaleʿīd har qadr hemmat bar leṭāfat maṣrūf shavad ʿinda Allāh maḥbūb
ast intahā.

Your question concerning unused water and the point at which it is
considered used was submitted before the Throne. This is what God has
spoken: ``With a small quantity of water, such as a cupful or twice or
three times that amount, washing the hands and face in it renders it
used. But if it reaches the volume of a kurr, washing one or two faces
does not alter it, and there has been and is no objection to using it.
If one of the three properties appears in it---that is, if the color of
the water changes to some degree---it is considered used. O you who
drink from the cup of My loving care! The purpose is cleanliness. You
are familiar with the pools of {[}Persia{]}. However much effort is
devoted to cleanliness is pleasing in the sight of God.'' End.

\subsection{121}\label{section-2951}

\begin{otherlanguage}{persian}

و اینکه در بارهٴ ذکر مرقوم داشتید که بعضی از دوستان الهی قبل از طلوع
برمیخیزند و بذکر حقّ جلّ جلاله مشغول میگردند و حال آنکه این ذکر مال یوم
است و تحدید یوم در ایّام صیام من الطّلوع الی الأفول شده این فقره هم در
ساحت اقدس من لا یعزب عن علمه من شیء عرض شد هذا ما نطق به الحقّ جلّ جلاله
یوم فی‌الحقیقه در بیان من الطّلوع الی الأفول است ولکن نفوسیکه قبل از طلوع
حبّاً للّه برمیخیزند و بذکر مشغول میشوند این لدی العرش محبوب است و مقبول
دیگر اعاده لازم نه و همچنین در مشرق الأذکار دوستان الهی متمسّکین بما نزل
فی کتاب الأقدس در اسحار توجّه بمشرق الأذکار مینمایند و بذکر محبوب عالمیان
مشغول میشوند ایشانهم اگر در طلوع و بعد از طلوع بساعتی بذکر مشغول شوند
مقبولست حقّ در جمیع احوال راحت عباد را خواسته و ذکرش اگر بتمام روح و
ریحان از ذاکری ظاهر شود محبوب و مقبول بوده و خواهد بود طوبی لنفس تنفّست
بحبّ مولاه و طوبی لعبد ایقظته نفحات محبّة اللّه انتهی

\end{otherlanguage}

Va īnke dar bāra-ye dhekr marqūm dāshtīd ke baʿżī az dūstān-e elāhī qabl
az ṭolūʿ barmīkhīzand va be dhekr-e Ḥaqq jalla jalāluhu mashghūl
mīgardand va ḥāl ānke īn dhekr māl-e yowm ast va taḥdīd-e yowm dar
ayyām-e ṣiyām min al-ṭulūʿ ilā al-ufūl shode īn faqara ham dar sāḥat-e
aqdas-e man lā yaʿzubu ʿan ʿilmihi min shayʾ ʿarż shod hādhā mā naṭaqa
bihi al-Ḥaqq jalla jalāluhu yowm fī al-ḥaqīqa dar Bayān min al-ṭulūʿ ilā
al-ufūl ast va lāken nofūsī ke qabl az ṭolūʿ ḥubban li-llāh barmīkhīzand
va be dhekr mashghūl mīshavand īn ladā al-ʿarsh maḥbūb ast va maqbūl
dīgar eʿāda lāzem na va hamchenīn dar Mashreq al-Adhkār dūstān-e elāhī
motamassekīn bi-mā nazala fī Kitāb al-Aqdas dar asḥār tavajjoh be
Mashreq al-Adhkār mīnamāyand va be dhekr-e Maḥbūb-e ʿālamiyān mashghūl
mīshavand īshān ham agar dar ṭolūʿ va baʿd az ṭolūʿ be sāʿatī be dhekr
mashghūl shavand maqbūl ast Ḥaqq dar jamīʿ-e aḥvāl rāḥat-e ʿebād rā
khvāste va dhekrash agar be tamām-e rūḥ va reyḥān az dhākerī ẓāher
shavad maḥbūb va maqbūl būde va khvāhad būd ṭūbā li-nafsin tanaffasat
bi-ḥubbi mawlāhā wa-ṭūbā li-ʿabdin ayqaẓat-hu nafaḥātu maḥabbati Allāh
intahā.

Concerning remembrance, you wrote that some of the friends of God rise
before sunrise and engage in remembrance of the True One, glorious is
His majesty, whereas this remembrance belongs to the day, and during the
days of fasting the day has been defined as extending from sunrise to
sunset. This matter too was submitted in the most holy presence of Him
from whose knowledge nothing is hidden. This is what the True One,
glorious is His majesty, has spoken: ``In the Bayān, the day is indeed
from sunrise to sunset. However, when souls rise before sunrise out of
love for God and engage in remembrance, this is pleasing and acceptable
before the Throne; there is no need to repeat it. Likewise, in
connection with the Mashriqu'l-Adhkār, the friends of God, holding fast
to what has been revealed in the Most Holy Book, make their way to the
Mashriqu'l-Adhkār in the hours before dawn and engage in remembrance of
the Beloved of all the worlds. If they too engage in remembrance at
sunrise or at an hour after sunrise, it is acceptable. In all
circumstances, the True One has desired the ease of His servants, and if
remembrance of Him proceeds from anyone with complete joy and radiance,
it has been and will remain pleasing and acceptable. Blessed is the soul
that has breathed with love for its Lord, and blessed is the servant
whom the breezes of the love of God have awakened.'' End.

\subsection{122}\label{section-2952}

\begin{otherlanguage}{arabic}

اینکه در ذکر جناب آقا خلیل عطّار علیه بهآء اللّه و صعود اخیه مرقوم داشتید
هذا ما انزله اللّه له قوله جلّ کبریائه

\end{otherlanguage}

Īnke dar dhekr-e janāb-e Āqā Khalīl-e ʿAṭṭār ʿalayhi bahāʾu Allāh va
ṣoʿūd-e akhīh marqūm dāshtīd hādhā mā anzalahu Allāhu lahu qawluhu jalla
kibriyāʾuhu.

Concerning what you wrote about Āqā Khalīl ʿAṭṭār, may the glory of God
rest upon him, and the ascension of his brother, this is what God has
revealed for him, glorious is His grandeur in His words:

\subsection{123}\label{section-2953}

\begin{otherlanguage}{arabic}

بسمی الأقدم الأعظم

\end{otherlanguage}

Bi-smī al-aqdami al-aʿẓam.

In My name, the Most Ancient, the Most Great.

\subsection{124}\label{section-2954}

\begin{otherlanguage}{arabic}

کتاب نزّل بالحقّ لمن اخذه جذب بیان الرّحمن و اجاب مولاه اذ اتی بأمر عظیم و
ینادی یوم اللّه فیما سواه طوبی لذی سمع سمع و لذی قلب اقبل الیه بفرح مبین
انّ الشّمس مع اشراقها تقبّس من نوره ولکنّ القوم اکثرهم من الغافلین قد ظهر کلّ
ما کان مستوراً فی حجب الغیب و برز ما کان مسطوراً فی کتب اللّه ربّ العالمین
من النّاس من وجد عرف قمیصی و سرع الی بحر عنایتی و منهم من اعرض و اعترض
علی ما نزل من ملکوت بیانی البدیع کم من عالم نقض میثاق اللّه و عهده و کم
من غافل سمع النّدآء من الأفق الأعلی و اقبل بقلبه الی اللّه العلیم الخبیر

\end{otherlanguage}

Kitābun nuzzila bi-l-ḥaqqi li-man akhadhahu jadhbu bayāni al-Raḥmān
wa-ajāba mawlāhu idh atā bi-amrin ʿaẓīm, wa-yunādī yawmu Allāhi fīmā
siwāhu. Ṭūbā li-dhī samʿin samiʿa wa-li-dhī qalbin aqbala ilayhi
bi-faraḥin mubīn. Inna al-shamsa maʿa ishrāqihā taqtabisu min nūrihi
wa-lākinna al-qawma aktharuhum mina al-ghāfilīn. Qad ẓahara kullu mā
kāna mastūran fī ḥujubi al-ghaybi wa-baraza mā kāna masṭūran fī kutubi
Allāhi rabbi al-ʿālamīn. Mina al-nāsi man wajada ʿarfa qamīṣī wa-saraʿa
ilā baḥri ʿināyatī wa-minhum man aʿraḍa wa-iʿtaraḍa ʿalā mā nazala min
malakūti bayānī al-badīʿ. Kam min ʿālimin naqaḍa mīthāqa Allāhi
wa-ʿahdahu wa-kam min ghāfilin samiʿa al-nidāʾa mina al-ufuqِ al-aʿlā
wa-aqbala bi-qalbihi ilā Allāhi al-ʿalīmi al-khabīr.

A Book sent down in truth for one whom the rapture of the All-Merciful's
utterance has seized, and who has answered his Lord when He came with a
mighty Cause, while the Day of God calls out to all besides Him. Blessed
is the one with ears who has heard, and the one with a heart who has
turned to Him with manifest joy. The sun, despite its radiance, draws
light from His light, yet most of the people are heedless. Everything
concealed within the veils of the unseen has appeared, and what was
inscribed in the Books of God, the Lord of the worlds, has come forth.
Among the people are those who have caught the fragrance of My robe and
hastened to the ocean of My loving care; others have turned away and
objected to what has descended from the kingdom of My wondrous
utterance. How many a learned man has broken God's covenant and pledge,
and how many a heedless one has heard the call from the Most Exalted
Horizon and turned his heart toward God, the All-Knowing, the All-Aware.

\subsection{125}\label{section-2955}

\begin{otherlanguage}{arabic}

یا خلیل انّا ذکرناک من قبل و ارسلنا الیک ما حکی عن ظهوری و اقتداری و
رحمتی و شفقتی و فضلی الّذی احاط من فی السّموات و الأرضین لا تحزن عمّا ورد
علیک قد فاز برفیق الأعلی من آمن باللّه مالک هذا الیوم المنیر انّا ذکرناه
فضلاً من عندنا علیک و علیه و انا الفضّال الکریم طوبی له بما فاز بالاقبال
اذ اعرض عنّی اکثر عبادی الّذین خلقناهم لخدمة امری المحکم المتین انّ ربّک غفر
اخاک حین صعوده و طهّره عن الآثام و ادخله فی مقام عجزت من اوصافه السن
العارفین انّا نکبّر من هذا المقام علی احبّائی فی هناک الّذین شربوا رحیق
بیانی و فازوا بما جری من قلمی الأعلی انّ ربّک لهو الغفور الرّحیم البهآء من
لدی اللّه مالک الأسمآء علیک و علی الّذین نزّلت اسمائهم فی لوحی الحفیظ انتهی

\end{otherlanguage}

Yā Khalīl, innā dhakarnāka min qablu wa-arsalnā ilayka mā ḥakā ʿan
ẓuhūrī wa-iqtidārī wa-raḥmatī wa-shafaqatī wa-faḍlī alladhī aḥāṭa man fī
al-samāwāti wa-l-araḍīn. Lā taḥzan ʿammā warada ʿalayka. Qad fāza
bi-Rafīqi al-Aʿlā man āmana bi-Allāhi māliki hādhā al-yawmi al-munīr.
Innā dhakarnāhu faḍlan min ʿindinā ʿalayka wa-ʿalayhi wa-anā al-faḍḍālu
al-karīm. Ṭūbā lahu bi-mā fāza bi-l-iqbāli idh aʿraḍa ʿannī aktharu
ʿibādī alladhīna khalaqnāhum li-khidmati amrī al-muḥkami al-matīn. Inna
rabbaka ghafara akhāka ḥīna ṣuʿūdihi wa-ṭahharahu ʿani al-āthāmi
wa-adkhalahu fī maqāmin ʿajazat min awṣāfihi alsunu al-ʿārifīn. Innā
nukabbiru min hādhā al-maqāmi ʿalā aḥibbāʾī fī hunāka alladhīna sharibū
raḥīqa bayānī wa-fāzū bi-mā jarā min qalamī al-aʿlā. Inna rabbaka
la-huwa al-ghafūru al-raḥīm. Al-bahāʾu min ladā Allāhi māliki al-asmāʾi
ʿalayka wa-ʿalā alladhīna nuzzilat asmāʾuhum fī lawḥī al-ḥafīẓ. Intahā.

O Khalīl, We previously remembered you and sent you what told of My
manifestation, My power, My mercy, My compassion, and My grace, which
has encompassed all who are in the heavens and on earth. Do not grieve
over what has befallen you. He who believed in God, the Master of this
luminous Day, has attained the companionship of the Most Exalted. We
have remembered him as a grace from Us to you and to him, and I am the
Most Bountiful, the Generous. Blessed is he for having turned toward Me
when most of My servants, whom We created to serve My firm and steadfast
Cause, turned away from Me. Your Lord forgave your brother at the time
of his ascension, purified him of sins, and admitted him to a station
that the tongues of those who know are powerless to describe. From this
station We send greetings to My loved ones there who have drunk the
choice wine of My utterance and attained what has flowed from My Most
Exalted Pen. Your Lord is indeed the Ever-Forgiving, the Merciful. May
glory from God, the Lord of Names, rest upon you and upon those whose
names have been sent down in My preserving Tablet. End.

\subsection{126}\label{section-2956}

\begin{otherlanguage}{persian}

الحمد للّه آفتاب کرم مشرق و بحر عنایت ظاهر و فرات بیان من غیر تعطیل جاری
و ساری این فانی از حقّ سائل و آمل است که جمیع من علی الأرض را از نفحات
این یوم امنع اقدس محروم نفرماید و به ما یحبّ و یرضی موفّق نماید انّه لهو
المقتدر القدیر

\end{otherlanguage}

Al-ḥamdu li-llāh, āftāb-i karam mushriq va baḥr-i ʿināyat ẓāhir va
Furāt-i bayān min ghayr-i taʿṭīl jārī va sārī. Īn fānī az Ḥaqq sāʾil va
āmil ast ki jamīʿ-i man ʿalā al-arḍ rā az nafaḥāt-i īn yawm-i amnaʿ-i
aqdas maḥrūm nafarmāyad va bi-mā yuḥibbu wa-yarḍā muwaffaq namāyad.
Innahu la-huwa al-muqtadiru al-qadīr.

Praise be to God, the sun of generosity is shining, the ocean of loving
care is manifest, and the Euphrates of utterance flows and spreads
without interruption. This mortal asks and hopes that God will deprive
no one on earth of the fragrant breezes of this most exalted and most
holy Day, and will enable everyone to do what He loves and approves. He
is indeed the All-Powerful, the Omnipotent.

\subsection{127}\label{section-2957}

\begin{otherlanguage}{persian}

خدمت جمیع دوستان از جانب این فانی تکبیر برسانید و بگوئید این ایّام غیر
ایّام است و این زمان غیر ازمنهٴ قبل و بعد جهد نمائید تا از انوار فجر الهی
قسمت برید انّه لهو المقتدر علی ما یشآء لا اله الّا هو العزیز العظیم

\end{otherlanguage}

Khidmat-i jamīʿ-i dūstān az jānib-i īn fānī takbīr birasānīd va bigūʾīd
īn ayyām ghayr-i ayyām ast va īn zamān ghayr-i azmina-yi qabl va baʿd.
Jahd namāʾīd tā az anvār-i fajr-i ilāhī qismat barīd. Innahu la-huwa
al-muqtadiru ʿalā mā yashāʾu lā ilāha illā huwa al-ʿazīzu al-ʿaẓīm.

Convey greetings from this mortal to all the friends, and say: ``These
days are unlike other days, and this time is unlike the times before and
after it. Strive to gain a share of the light of the divine dawn.'' He
indeed has power over whatever He wills. There is no God but Him, the
Almighty, the Most Great.

\subsection{128}\label{section-2958}

\begin{otherlanguage}{arabic}

البهآء علی حضرتکم و علی من معکم و علی الّذین اقبلوا بقلوب نورآء الی اللّه
ربّنا و ربّکم و مقصودنا و مقصودکم و مقصود من فی السّموات و الأرضین

\end{otherlanguage}

Al-bahāʾu ʿalā ḥaḍratikum wa-ʿalā man maʿakum wa-ʿalā alladhīna aqbalū
bi-qulūbin nawrāʾa ilā Allāhi rabbinā wa-rabbikum wa-maqṣūdinā
wa-maqṣūdikum wa-maqṣūdi man fī al-samāwāti wa-l-araḍīn.

May glory rest upon you, upon those with you, and upon those who have
turned with radiant hearts toward God, our Lord and your Lord, our
heart's desire and your heart's desire, and the heart's desire of all
who are in the heavens and on earth.

\subsection{129}\label{section-2959}

\begin{otherlanguage}{persian}

و بعد از اتمام مکتوب و حضور در ساحت اقدس لسان محبوب باین کلمهٴ مبارکهٴ
علیا ناطق یا ورقآء علیک بهائی و رحمتی و عنایتی ذکر حکمت و معانی آن
فی‌الجمله اظهار شد ولکن حکمت امروز آداب و اعمال و اقوالی است که ناس را
بافق اعلی هدایت نماید من دون آنکه ضوضائی مرتفع شود و یا فتنه و فسادی
ظاهر گردد انّ النّاس فی نوم عجیب و ریب مبین یعبدون الظّنون و لا یفقهون و
یسرعون الی الأوهام و لا یشعرون باید بمدارا و رأفت غافلین را آگاه نمود و
بدریای عنایت هدایت کرد بعضی از جهّال که بأجنحهٴ نفس و هوی طائرند و بنار
بغضا مشتعل بر مراصد مترصّدند که دست آویزی بیابند و بر حقّ جلّ جلاله و
اولیای او اعتراض نمایند انتهی

\end{otherlanguage}

Va baʿd az itmām-i maktūb va ḥuḍūr dar sāḥat-i aqdas lisān-i Maḥbūb
bi-īn kalima-yi mubāraka-yi ʿulyā nāṭiq: Yā Warqāʾ, ʿalayka bahāʾī
wa-raḥmatī wa-ʿināyatī. Dhikr-i ḥikmat va maʿānī-yi ān fī-l-jumla iẓhār
shud va-līkin ḥikmat-i imrūz ādāb va aʿmāl va aqvālī ast ki nās rā
bi-ufuq-i aʿlā hidāyat namāyad min dūn-i ānki ḍawḍāʾī murtafiʿ shavad va
yā fitna va fasādī ẓāhir gardad. Inna al-nāsa fī nawmin ʿajībin
wa-raybin mubīn. Yaʿbudūna al-ẓunūna wa-lā yafqahūna wa-yusriʿūna ilā
al-awhāmi wa-lā yashʿurūn. Bāyad bi-mudārā va raʾfat ghāfilīn rā āgāh
namūd va bi-daryā-yi ʿināyat hidāyat kard. Baʿḍī az juhhāl ki
bi-ajniḥa-yi nafs va havā ṭāʾirand va bi-nār-i baghḍā mushtaʿil bar
marāṣid mutaraṣṣidand ki dast-āvīzī biyāband va bar Ḥaqq jalla jalāluhu
va awliyā-yi ū iʿtirāḍ namāyand. Intahā.

After the letter was completed and I entered the Most Holy Presence, the
tongue of the Beloved spoke these exalted and blessed words: ``O Warqāʾ,
upon you be My glory, My mercy, and My loving care. Wisdom and its
meanings have been explained to some extent. Today, however, wisdom
consists of conduct, deeds, and words that guide people to the Most
Exalted Horizon without raising an uproar or giving rise to turmoil and
discord. The people are indeed sunk in a strange sleep and manifest
doubt. They worship conjectures and do not understand; they hasten
toward delusions and do not perceive. Through forbearance and
compassion, the heedless must be awakened and guided to the ocean of
loving care. Some of the ignorant, who fly on the wings of self and
desire and burn with the fire of hatred, lie in wait at their
watchposts, seeking a pretext to raise objections against God, exalted
be His glory, and His loved ones.'' End.

\subsection{130}\label{section-2960}

\begin{otherlanguage}{persian}

انشآءاللّه جمیع دوستان الهی بطراز حکمت مزیّن شوند و از هر یک ظاهر شود آنچه
مشعر و مدلّ بر ارتفاع امر است حقّ علمای عصر را انصاف عطا فرماید و بصر
عنایت نماید شاید انوار آفتاب حقیقت را که در قطب زوال منیر و روشن است
مشاهده نمایند و از او محروم نمانند لازال این نفوس سبب ضلالت و گمراهی
بوده و هستند الّا من شآءاللّه یسئل الخادم ربّه بأن یعرّفهم و یعلّمهم ما غفلوا
عنه و یقرّبهم الیه فضلاً من عنده انّه لهو الغفور الرّحیم

\end{otherlanguage}

In shāʾ Allāh jamīʿ-i dūstān-i ilāhī bi-ṭirāz-i ḥikmat muzayyan shavand
va az har yak ẓāhir shavad ānchi mushʿir va mudill bar irtifāʿ-i amr
ast. Ḥaqq ʿulamā-yi ʿaṣr rā inṣāf ʿaṭā farmāyad va baṣar ʿināyat namāyad
shāyad anvār-i āftāb-i ḥaqīqat rā ki dar quṭb-i zavāl munīr va rawshan
ast mushāhada namāyand va az ū maḥrūm namānand. Lā-zāl īn nufūs sabab-i
ḍalālat va gumrāhī būdi va hastand illā man shāʾa Allāh. Yasʾalu
al-khādimu rabbahu bi-an yuʿarrifahum wa-yuʿallimahum mā ghafalū ʿanhu
wa-yuqarribahum ilayhi faḍlan min ʿindihi. Innahu la-huwa al-ghafūru
al-raḥīm.

God willing, may all the friends of God be adorned with wisdom, and may
each manifest what indicates and attests to the exaltation of the Cause.
May God grant the learned men of this age fairness and bestow sight upon
them, so that they may behold the rays of the Sun of Truth, shining
brightly at its zenith, and not remain deprived of it. These souls have
always been, and remain, a cause of error and misguidance, except those
whom God wills. This servant asks his Lord to grant them understanding,
teach them what they have overlooked, and draw them near to Him as a
grace from Himself. He is indeed the Ever-Forgiving, the Merciful.

\subsection{131}\label{section-2961}

\begin{otherlanguage}{persian}

عرض دیگر اینکه در بارهٴ مرفوع مرحوم حضرت والد الّذی فاز باللّقآء حین
الارتقآء مرقوم {[}داشتید{]} در ایّامیکه جناب ابوالحسن علیه ٩ ٦٦ {[}بهآء
اللّه{]} در ارض اقدس حضور داشتند یومی از ایّام این عبد و ایشان هر دو در
ساحت اقدس تلقاء وجه حاضر بغتةً لسان مقصود عالمیان متکلّم و یک لوح امنع
اقدس در ذکر ایشان نازل بجناب مذکور داده شد که برسانند آن محبوب در بارهٴ
حضرت مرفوع علیه من کلّ بهآء ابهاه مطمئن باشند

\end{otherlanguage}

ʿArḍ-i dīgar īnki dar bāra-yi marfūʿ-i marḥūm ḥaḍrat-i vālid alladhī
fāza bi-l-liqāʾi ḥīna al-irtiqāʾ marqūm {[}dāshtīd{]}. Dar ayyāmīki
janāb-i Abu'l-Ḥasan ʿalayhi 9 66 {[}Bahāʾ Allāh{]} dar arḍ-i aqdas ḥuḍūr
dāshtand yawmī az ayyām īn ʿabd va īshān har du dar sāḥat-i aqdas
tilqāʾ-i vajh ḥāḍir, baghtatan lisān-i Maqṣūd-i ʿālamiyān mutakallim va
yak lawḥ-i amnaʿ-i aqdas dar dhikr-i īshān nāzil, bi-janāb-i madhkūr
dādi shud ki birasānand. Ān maḥbūb dar bāra-yi ḥaḍrat-i marfūʿ ʿalayhi
min kulli bahāʾin abhāhu muṭmaʾinn bāshand.

Another matter: you {[}wrote{]} concerning your exalted, departed
father, who attained the divine presence at the time of his ascension.
During the days when Abū'l-Ḥasan, upon him be 9 66 {[}the glory of
God{]}, was in the Holy Land, one day he and this servant were both in
the Most Holy Presence, before His face. Suddenly the tongue of the
Desire of all the worlds spoke, and a most exalted and most holy Tablet
was revealed in your father's remembrance and given to the
aforementioned gentleman to deliver. May you, dear friend, be assured
concerning your exalted father---upon him be the most glorious of all
glory.

\subsection{132}\label{section-2962}

\begin{otherlanguage}{persian}

در مواقع متعدّده ذکر ایشان نازل شده علی شأن تطیر به افئدة العارفین مخصوص
در چند مقام خدمات ایشان و استقامت ایشان در امر اللّه از لسان محبوب
عالمیان شنیده شد قد ذکره بما لا یأخذه الفنآء و یذکره بما یتوجّه الیه
الأفئدة و القلوب

\end{otherlanguage}

Dar mavāqiʿ-i mutaʿaddida dhikr-i īshān nāzil shudi ʿalā shaʾnin taṭīru
bihi afʾidatu al-ʿārifīn. Makhṣūṣ dar chand maqām khadamāt-i īshān va
istiqāmat-i īshān dar Amr-i Allāh az lisān-i Maḥbūb-i ʿālamiyān shinīdi
shud. Qad dhakarahu bi-mā lā yaʾkhudhuhu al-fanāʾu wa-yadhkuruhu bi-mā
yatawajjahu ilayhi al-afʾidatu wa-l-qulūb.

On numerous occasions, words in his remembrance have been revealed in
such a manner that the hearts of those who know soar through them. In
particular, on several occasions the Beloved of all the worlds was heard
speaking of his services and his steadfastness in the Cause of God. He
has remembered him in words that will never perish, and will remember
him in words toward which hearts and souls will turn.

\subsection{133}\label{section-2963}

\begin{otherlanguage}{arabic}

این فانی در آخر کتاب اینقدر عرض مینماید بعنایت محبوب فائزید و هذه بشارة
منّی الیکم

\end{otherlanguage}

Īn fānī dar ākhir-i kitāb īnqadr ʿarḍ mīnamāyad: bi-ʿināyat-i Maḥbūb
fāʾizīd wa-hādhihi bishāratun minnī ilaykum.

At the close of this letter, this mortal offers just this: you have
attained the loving care of the Beloved, and this is glad tidings from
me to you.

\bookmarksetup{startatroot}

\chapter{BH00014 (Unnamed)}\label{bh00014-unnamed}

\emph{Translated on 17 Qudrat 183 using GPT 6.1 Sol}

\subsection{1}\label{section-2964}

\begin{otherlanguage}{arabic}

بسم ربّنا الأقدس الأعظم العلیّ الأبهی

\end{otherlanguage}

Bi-smi Rabbinā al-Aqdas al-Aʿẓam al-ʿAliyy al-Abhā.

In the name of our Lord, the Most Holy, the Most Great, the Most
Exalted, the All-Glorious.

\subsection{2}\label{section-2965}

\begin{otherlanguage}{arabic}

للّه الحمد حمداً یمطر به السّحاب و یجد المقرّبون عرف قمیص ربّنا الوهّاب حمداً
به قصد العشّاق منظر اللّه مالک المیثاق حمداً سرع به المخلصون الی مقام
المکاشفة و اللّقآء و الموحّدون الی الذّورة العلیا حمداً به جرت الأنهار و
اثمرّت الأشجار اشهد انّه لا اله الّا هو وحده لا شریک له التّکبیر و البهآء و
الذّکر و الثّنآء علی اولیائه و اصفیائه الّذین ما اخذتهم فی اللّه لومة لائم و
ما منعتهم سطوة ظالم و ما خوّفتهم ضوضآء جاهل اولئک عباد مکرمون لا خوف
علیهم و لا هم یحزنون

\end{otherlanguage}

Li-llāh al-ḥamd, ḥamdan yumṭiru bihi al-saḥāb wa-yajidu al-muqarrabūn
ʿarfa qamīṣi Rabbinā al-Wahhāb; ḥamdan bihi qaṣada al-ʿushshāq manẓara
Allāh Mālik al-Mīthāq; ḥamdan saraʿa bihi al-mukhliṣūn ilā maqām
al-mukāshafa wa-al-liqāʾ wa-al-muwaḥḥidūn ilā al-dhirwa al-ʿulyā; ḥamdan
bihi jarat al-anhār wa-athmarat al-ashjār. Ashhadu annahu lā ilāha illā
Huwa waḥdahu lā sharīka lahu. Al-takbīr wa-al-bahāʾ wa-al-dhikr
wa-al-thanāʾ ʿalā awliyāʾihi wa-aṣfiyāʾihi alladhīna mā akhadhathum fī
Allāh lawmatu lāʾim wa-mā manaʿathum saṭwatu ẓālim wa-mā khawwafathum
ḍawḍāʾu jāhil. Ulāʾika ʿibādun mukramūn, lā khawfun ʿalayhim wa-lā hum
yaḥzanūn.

Praise be to God---a praise through which the clouds rain down and those
brought near perceive the fragrance of the robe of our Lord, the
All-Bountiful; a praise through which the lovers seek the countenance of
God, the Lord of the Covenant; a praise through which the sincere hasten
to the station of unveiling and reunion, and those who affirm His unity
to the loftiest summit; a praise through which the rivers flow and the
trees bear fruit. I testify that there is no God but Him, alone, without
partner. Exaltation, glory, remembrance, and praise be upon His loved
ones and His chosen ones, whom no blame from any blamer has deterred in
the path of God, whom no tyrant's might has held back, and whom no
ignorant person's clamor has frightened. They are honored servants; no
fear will come upon them, nor will they grieve.

\subsection{3}\label{section-2966}

\begin{otherlanguage}{arabic}

و بعد نامه‌های آنمحبوب فی‌الحقیقه بمثابهٴ امطار از سحاب محبّت مترادفاً
متتابعاً هاطل انّ الخادم یشکر الّذی ایّدک و وفّقک و عرّفک و یسئله بأن یرفعک فی
الملک و یعزّک بین عباده انّه علی کلّشیء قدیر

\end{otherlanguage}

Va baʿd, nāmih-hā-yi ān maḥbūb fī al-ḥaqīqih bi-mathābiy-i amṭār az
saḥāb-i maḥabbat mutarādifan mutatābiʿan hāṭil. Inna al-khādima yashkuru
alladhī ayyadaka wa-waffaqaka wa-ʿarrafaka wa-yasʾaluhu bi-an yarfaʿaka
fī al-mulk wa-yuʿizzaka bayna ʿibādihi, innahu ʿalā kulli shayʾin qadīr.

To proceed: your letters, beloved friend, have indeed been pouring down
in succession, like rain from the clouds of love. This servant thanks
Him who has strengthened you, granted you success, and bestowed
recognition upon you, and asks Him to raise you up in the world and
honor you among His servants. He is powerful over all things.

\subsection{4}\label{section-2967}

\begin{otherlanguage}{persian}

یک نامهٴ آنمحبوب که تاریخ آن ١٢ شهر صفر بود قلب را بنفحات سرور و چشم را
بانوار محبّت مزیّن و منوّر فرمود صدهزار بار شکر مقصود عالمیانرا چه که توفیق
عطا فرموده و مددش در هر آن میرسد والّا یرجع الوجود الی العدم و بعد از
مشاهده و ملاحظه و قرائت قصد مقام اعلی و مقرّ اسمی الأسنی نموده بعد از اذن
تلقاء وجه مالک اسماء عرض شد هذا ما نزّل فی الجواب قوله جلّ جلاله و عمّ
نواله

\end{otherlanguage}

Yik nāmiy-i ān maḥbūb kih tārīkh-i ān 12 shahr-i Ṣafar būd qalb rā
bi-nafaḥāt-i surūr va chashm rā bi-anvār-i maḥabbat muzayyan va munavvar
farmūd. Ṣad-hizār bār shukr-i Maqṣūd-i ʿālamiyān rā, chih kih tawfīq
ʿaṭā farmūdih va madadash dar har ān mīrasad, va-illā yarjiʿu al-wujūdu
ilā al-ʿadam. Va baʿd az mushāhidih va mulāḥiẓih va qirāʾat qaṣd-i
maqām-i aʿlā va maqarr-i asmā al-asnā nimūdih, baʿd az idhn tilqāʾ-i
vajh-i Mālik-i asmāʾ ʿarḍ shud. Hādhā mā nuzzila fī al-jawāb, qawluhu
jalla jalāluhu wa-ʿamma nawāluhu:

One of your letters, beloved friend, dated the twelfth of the month of
Ṣafar, adorned and illumined the heart with the breezes of joy and the
eyes with the lights of love. A hundred thousand thanks to the Desire of
all peoples, for He has granted success, and His assistance arrives at
every moment; otherwise, existence would return to nothingness. After
examining, considering, and reading it, I proceeded to the supreme
station and the most exalted, most sublime abode. Once permission was
given, it was presented before the face of the Lord of Names. This is
what was revealed in reply---His words, exalted be His glory and
all-embracing His bounty:

\subsection{5}\label{section-2968}

\begin{otherlanguage}{arabic}

هو المشرق من افق البیان

\end{otherlanguage}

Huwa al-Mushriq min ufuq al-bayān.

He is the One who shines from the horizon of utterance.

\subsection{6}\label{section-2969}

\begin{otherlanguage}{persian}

یا امین علیک بهائی و عنایتی و رحمتی الّتی سبقت من فی السّموات و الأرضین
حمد کن مالک قدم را که ترا مؤیّد فرمود بر اقبال و ذکر و ثنا و خدمت امرش
انّه عرّفک صراطه و هداک الی سبیله و وفّقک علی خدمة امره و رفعک الی مقام
یذکرک القلم الأعلی من الأفق الأبهی المقام الّذی اقرّ الوری بعجزه عن ذکره
ان افرح بهذا الفضل الأعظم و قل لک الحمد یا مقصود العالمین

\end{otherlanguage}

Yā Amīn, ʿalayka bahāʾī wa-ʿināyatī wa-raḥmatī allatī sabaqat man fī
al-samāwāt wa-al-araḍīn. Ḥamd kun Mālik-i qidam rā kih turā muʾayyad
farmūd bar iqbāl va dhikr va thanā va khidmat-i amrash. Innahu ʿarrafaka
ṣirāṭahu wa-hadāka ilā sabīlihi wa-waffaqaka ʿalā khidmati amrihi
wa-rafaʿaka ilā maqāmin yadhkuruka al-Qalam al-Aʿlā min al-Ufuq al-Abhā,
al-maqām alladhī aqarra al-warā bi-ʿajzihi ʿan dhikrihi. An ifraḥ
bi-hādhā al-faḍl al-aʿẓam wa-qul laka al-ḥamdu yā Maqṣūd al-ʿālamīn.

O Amīn, upon you be My glory, My loving care, and My mercy, which has
preceded all who are in the heavens and the earths. Praise the Lord of
Eternity, who has strengthened you to turn toward Him, to remember and
praise Him, and to serve His Cause. He has made His path known to you,
guided you to His way, enabled you to serve His Cause, and raised you to
a station in which the Most Exalted Pen remembers you from the
All-Glorious Horizon---a station that all humanity has acknowledged its
inability to describe. Rejoice in this most great bounty and say:
``Praise be to You, O Desire of all peoples.''

\subsection{7}\label{section-2970}

\begin{otherlanguage}{persian}

یا امین در اکثری از الواح اولیا ذکر تو از قلم اعلی جاری و نازل طوبی لک و
لمن هداک و اراک و اسمعک و ارشدک و الحمد لمن ایّدک و وفّقک علی عرفان من منع
عن عرفانه اکثر العباد یشهد بذلک مالک الایجاد فی هذا المقام الرّفیع اولیای
حقّ جلّ جلاله را تکبیر برسان و از قبل مظلوم ذکر نما تا بنار ذکر مشتعل شوند
و بحکمت و بیان بتبلیغ امر رحمن مشغول گردند ناس بسیار غافلند اینمظلوم در
جمیع احیان و اوان بأعلی البیان ایشان را بافق رحمن دعوت مینماید ولکن کل
از سمع و بصر محرومند الّا من شآء اللّه چشمها را حجبات از مشاهده منع نموده
و گوشها را قطن اوهام از اصغا بازداشته رحم اللّه العدل و الانصاف از این دو
جز اسم نمانده نام موجود و رسم مفقود الأمر بید اللّه ربّ الغیب و الشّهود

\end{otherlanguage}

Yā Amīn, dar aktharī az alvāḥ-i awliyā dhikr-i tu az Qalam-i Aʿlā jārī
va nāzil. Ṭūbā laka wa-li-man hadāka wa-arāka wa-asmaʿaka wa-arshadaka,
wa-al-ḥamdu li-man ayyadaka wa-waffaqaka ʿalā ʿirfāni man muniʿa ʿan
ʿirfānihi aktharu al-ʿibād. Yashhadu bi-dhālika Mālik al-ījād fī hādhā
al-maqām al-rafīʿ. Awliyā-yi Ḥaqq jalla jalāluhu rā takbīr birasān va az
qibal-i Maẓlūm dhikr namā tā bi-nār-i dhikr mushtaʿil shavand va
bi-ḥikmat va bayān bi-tablīgh-i Amr-i Raḥmān mashghūl gardand. Nās
bisyār ghāfiland. Īn Maẓlūm dar jamīʿ-i aḥyān va avān bi-aʿlā al-bayān
īshān rā bi-Ufuq-i Raḥmān daʿvat mīnamāyad, va-lākin kull az samʿ va
baṣar maḥrūmand, illā man shāʾa Allāh. Chashm-hā rā ḥujubāt az
mushāhidih manʿ nimūdih va gūsh-hā rā quṭn-i awhām az iṣghā bāz dāshtih.
Raḥima Allāhu al-ʿadla wa-al-inṣāf. Az īn du juz ism namāndih, nām
mawjūd va rasm mafqūd. Al-amru bi-yadi Allāh Rabb al-ghayb wa-al-shuhūd.

O Amīn, in most of the Tablets to the loved ones, mention of you has
flowed and been revealed from the Most Exalted Pen. Blessed are you and
He who has guided you, enabled you to see and hear, and directed you
aright. Praise be to Him who has strengthened you and enabled you to
recognize the One whom most people have been kept from recognizing. To
this the Lord of creation bears witness in this exalted station. Convey
greetings of exaltation to the loved ones of God, exalted be His glory,
and remember them on behalf of the Wronged One, that they may be set
ablaze with the fire of remembrance and engage, with wisdom and
utterance, in teaching the Cause of the All-Merciful. The people are
deeply heedless. At every time and moment, this Wronged One calls them,
in the most exalted utterance, to the horizon of the All-Merciful; yet
all are deprived of hearing and sight, except those whom God wills.
Veils have kept their eyes from seeing, and the cotton of vain
imaginings has stopped their ears from listening. May God have mercy on
justice and fairness! Nothing but their names remains: the names exist,
but their traces are lost. The matter rests in the hands of God, the
Lord of the unseen and the seen.

\subsection{8}\label{section-2971}

\begin{otherlanguage}{persian}

یا امین بحر بیان امکان را احاطه نموده و ندآء اللّه از اعلی افق عالم مرتفع
معذلک احدی آگاه نشده مگر معدودی ولکن بأسی نبوده و نیست عنقریب ظاهر شود
آنچه در الواح از قلم اللّه جاری گشته انتهی

\end{otherlanguage}

Yā Amīn, baḥr-i bayān imkān rā iḥāṭih nimūdih va nidāʾ-i Allāh az aʿlā
ufuq-i ʿālam murtafiʿ. Maʿa-dhālik aḥadī āgāh nashudih magar maʿdūdī,
va-lākin baʾsī nabūdih va nīst. ʿAn-qarīb ẓāhir shavad ānchih dar alvāḥ
az Qalam-i Allāh jārī gashtih. Intahā.

O Amīn, the ocean of utterance has encompassed the world of being, and
the call of God has risen from the highest horizon of the world.
Nevertheless, no one has become aware of it except a few. Yet this has
not been, and is not, a cause for concern. Soon what has flowed from the
Pen of God in the Tablets will become manifest. End.

\subsection{9}\label{section-2972}

\begin{otherlanguage}{persian}

قد هطلت امطار رحمة ربّنا علیکم و علینا و اشرق نیّر کرمه امام وجوهنا من
یقدر ان یذکره حقّ الذّکر او یصفه حقّ الوصف او یثنیه حقّ الثّنآء اتوب الیه و
استغفره عمّا خطر ببالی و نطق به لسانی و حرّر به قلمی

\end{otherlanguage}

Qad haṭalat amṭāru raḥmati Rabbinā ʿalaykum wa-ʿalaynā wa-ashraqa
nayyiru karamihi amāma wujūhinā. Man yaqdiru an yadhkurahu ḥaqqa
al-dhikr aw yaṣifahu ḥaqqa al-waṣf aw yuthniyahu ḥaqqa al-thanāʾ? Atūbu
ilayhi wa-astaghfiruhu ʿammā khaṭara bi-bālī wa-naṭaqa bihi lisānī
wa-ḥarrara bihi qalamī.

The rains of our Lord's mercy have poured down upon you and upon us, and
the sun of His generosity has shone before our faces. Who can remember
Him as He deserves to be remembered, describe Him as He deserves to be
described, or praise Him as He deserves to be praised? I turn to Him in
repentance and ask His forgiveness for whatever has crossed my mind,
been spoken by my tongue, or been written by my pen.

\subsection{10}\label{section-2973}

\begin{otherlanguage}{persian}

ذکر آنمحبوب لازال بوده از آیات مشرقات عنایات حقّ جلّ جلاله واضح و معلوم
است صدهزار بار حمد و شکر که این ایّام نتیجهٴ کلمهٴ مبارکهٴ علیا که حین
ورود خبر سجن حضرت اسم‌اللّه و من مع حضرته علیهم بهآء اللّه الأبهی بآن تکلّم
فرمودند ظاهر و آشکار قوله تبارک و تعالی انّهم یستنصرون علینا بقوّة الجنود
و المدافع و نحن نستنصر علیهم بالقویّ الغالب القدیر انتهی

\end{otherlanguage}

Dhikr-i ān maḥbūb lā-zāl būdih, az āyāt-i mushriqāt-i ʿināyāt-i Ḥaqq
jalla jalāluhu vāḍiḥ va maʿlūm ast. Ṣad-hizār bār ḥamd va shukr kih īn
ayyām natījiy-i kalimiy-i mubārakiy-i ʿulyā kih ḥīn-i vurūd-i khabar-i
sijn-i Ḥaḍrat-i Ismu'llāh va man maʿa ḥaḍratihi ʿalayhim Bahāʾ Allāh
al-Abhā bi-ān takallum farmūdand ẓāhir va āshkār. Qawluhu tabāraka
wa-taʿālā: Innahum yastanṣirūna ʿalaynā bi-quwwati al-junūd
wa-al-madāfiʿ wa-naḥnu nastanṣiru ʿalayhim bi-al-Qawiyy al-Ghālib
al-Qadīr. Intahā.

You have always been remembered, beloved friend, as is clear and evident
from the shining verses of God's loving care, exalted be His glory. A
hundred thousand praises and thanks that, in these days, the outcome of
the exalted, blessed words He spoke when news arrived of the
imprisonment of the honored Ismu'llāh and those with him---upon them be
the glory of God, the All-Glorious---is plain and manifest. His words,
blessed and exalted is He: ``They seek victory over Us through the
strength of troops and cannons, while We seek victory over them through
the Strong, the All-Subduing, the All-Powerful.'' End.

\subsection{11}\label{section-2974}

\begin{otherlanguage}{persian}

آفاق عالم انسانی امروز بطراز جدید مزیّن و فائز دعای موحّدین و مخلصین و
آملین و سائلین و قاصدین باستجابت مقرون کل خلاصی خواسته بودند از برای
مسجونین و محبوسین مخصوص حضرت اسم‌اللّه علیه من کلّ بهآء ابهاه و من معه
علیهم بهآء اللّه الأبهی الحمد للّه ظاهر شد آنچه سبب و علّت فرح و سرور است

\end{otherlanguage}

Āfāq-i ʿālam-i insānī imrūz bi-ṭirāz-i jadīd muzayyan va fāʾiz. Duʿā-yi
muvaḥḥidīn va mukhliṣīn va āmilīn va sāʾilīn va qāṣidīn bi-istijābat
maqrūn. Kull khalāṣī khvāstih būdand az barā-yi masjūnīn va maḥbūsīn,
makhṣūṣ Ḥaḍrat-i Ismu'llāh ʿalayhi min kulli bahāʾin abhāhu va man
maʿahu ʿalayhim Bahāʾ Allāh al-Abhā. Al-ḥamdu li-llāh ẓāhir shud ānchih
sabab va ʿillat-i faraḥ va surūr ast.

Today the horizons of the human world have been graced and adorned with
a new splendor. The prayers of those who affirm God's unity, the
sincere, the hopeful, the suppliants, and the seekers have been
answered. All had asked for the release of the imprisoned and the
confined, especially the honored Ismu'llāh---upon him be the most
glorious of all glory---and those with him---upon them be the glory of
God, the All-Glorious. Praise be to God, what brings joy and gladness
has come to pass.

\subsection{12}\label{section-2975}

\begin{otherlanguage}{persian}

یا حبیب فؤادی ناس غافلند اگر یک آن تفکّر نمایند مشاهده میکنند عدم بقای
اعدا را و همچنین ظهورات قوّت و قدرت حقّ جلّ جلاله را مع سجن و حبس و صدهزار
بلایای ظاهره مفسدین ارض صاد را اخذ فرمود اخذ عزیز مقتدر و نفسی متنبّه نشد
و همچنین نفوس قاسیهٴ ارض طاء را از کاذب موسوم به صادق ثمّ الّذی قصد الرّزیّة
الکبری لأهل البهآء حال هر دو بسقر مقرّ دارند معذلک بقدر سمّ ابره خلق بشعور
نیامدند اینعبد اگر بخواهد از اینگونه امور ذکر نماید باید کتابی تألیف کند
حیرت در سکر و غفلت خلق است آیا شعور چه شد و ادراک کجا رفت حیرت اندر حیرت
است

\end{otherlanguage}

Yā ḥabīb-i fuʾādī, nās ghāfiland. Agar yik ān tafakkur namāyand
mushāhidih mīkunand ʿadam-i baqā-yi aʿdā rā va hamchunīn ẓuhūrāt-i
quvvat va qudrat-i Ḥaqq jalla jalāluhu rā. Maʿa sijn va ḥabs va
ṣad-hizār balāyā-yi ẓāhirih, mufsidīn-i arḍ-i Ṣād rā akhdh farmūd,
akhdha ʿazīzin muqtadir, va nafsī mutanabbih nashud. Va hamchunīn
nufūs-i qāsiyiy-i arḍ-i Ṭāʾ rā, az kādhib-i mawsūm bi-Ṣādiq, thumma
alladhī qaṣada al-raziyyata al-kubrā li-ahl al-Bahāʾ. Ḥāl har du
bi-Saqar maqarr dārand. Maʿa-dhālik bi-qadr-i samm-i ibrih khalq
bi-shuʿūr nayāmadand. Īn ʿabd agar bikhvāhad az īn-gūnih umūr dhikr
namāyad bāyad kitābī taʾlīf kunad. Ḥayrat dar sukr va ghaflat-i khalq
ast. Āyā shuʿūr chih shud va idrāk kujā raft? Ḥayrat andar ḥayrat ast.

O beloved of my heart, the people are heedless. If they reflected for a
single moment, they would see that the enemies do not endure, and would
likewise see the manifestations of the strength and power of God,
exalted be His glory. Despite imprisonment, confinement, and a hundred
thousand outward afflictions, He seized the corruptors of the land of
Ṣād with the grip of One mighty and all-powerful, yet not a soul took
heed. Likewise He seized the hard-hearted souls of the land of Ṭāʾ, from
the liar named Ṣādiq to the one who intended the greatest calamity for
the people of Bahāʾ. Now both dwell in hellfire. Nevertheless, the
people have not gained awareness even to the extent of a needle's eye.
If this servant wished to recount such matters, he would have to compose
a book. What is astonishing is the people's stupor and heedlessness.
What has become of awareness, and where has understanding gone? It is
bewilderment upon bewilderment.

\subsection{13}\label{section-2976}

\begin{otherlanguage}{persian}

فی‌الحقیقه در فقرهٴ ارض طاء باید تشکّر نمود علما که خود را پیشوا میدانند و
خود را هادی و قائد میشمرند هر حین بنار عناد مشتعلند هر یوم شعله‌اش بیشتر
و لهیبش واضحتر و ضرّش محکمتر هزار بار بر اهل بها هجوم نموده‌اند و هنوز
بقدر ذرّه ضغینهٴ ایشان ساکن نشده و بغضا آرام نگرفته اگر یکی از آن نفوس
خبیثه قدرت ظاهره مییافت و یا از امرا میشد آیا چه مینمود حال در پیشوایان
خلق و حضرت سلطان تفکّر نمائید مع آنکه از این طائفه رضایتی نداشتند و اعدا
هم نظر باظهار خدمت در لیالی و ایّام بمفتریات مشغول معذلک حضرت سلطان
بصرافت طبع مبارک من غیر واسطه مسجونین را آزادی بخشید و از حبس نجات عنایت
فرمود اگرچه مخلصین آنی از سجن و بلایای واردهٴ فی سبیل اللّه را بعالمی
مبادله نمینمایند قسم بمربّی کونین هر آنی از سجن فی سبیل اللّه نعمتی است که
از برای او شبه و مثل نبوده و نخواهد بود در هر حال باید انسان بعدل ناظر
باشد و بانصاف تکلّم نماید از حقّ جلّ جلاله این خادم فانی سائل و آمل که در
هر آن حضرت سلطان را مؤیّد دارد و موفّق فرماید انّه علی کلّشیء قدیر

\end{otherlanguage}

Fī al-ḥaqīqih dar fiqriy-i arḍ-i Ṭāʾ bāyad tashakkur nimūd. ʿUlamā kih
khud rā pīshvā mīdānand va khud rā hādī va qāʾid mīshumār and har ḥīn
bi-nār-i ʿinād mushtaʿiland, har yawm shuʿlih-ash bīshtar va lahībash
vāḍiḥtar va ḍarrash muḥkamtar. Hizār bār bar ahl-i Bahā hujūm
nimūdih-and va hanūz bi-qadr-i dharrih ḍaghīniy-i īshān sākin nashudih
va baghḍā ārām nagiriftih. Agar yikī az ān nufūs-i khabīthih qudrat-i
ẓāhirih mīyāft va yā az umarā mīshud āyā chih mīnamūd? Ḥāl dar
pīshvāyān-i khalq va Ḥaḍrat-i Sulṭān tafakkur namāyīd. Maʿa ānkih az īn
ṭāʾifih riḍāyatī nadāshtand va aʿdā ham naẓar bi-iẓhār-i khidmat dar
layālī va ayyām bi-muftariyāt mashghūl, maʿa-dhālik Ḥaḍrat-i Sulṭān
bi-ṣirāfat-i ṭabʿ-i mubārak min ghayr-i vāsiṭih masjūnīn rā āzādī
bakhshīd va az ḥabs najāt ʿināyat farmūd. Agarchih mukhliṣīn ānī az sijn
va balāyā-yi vāridiy-i fī sabīl-i Allāh rā bi-ʿālamī mubādilih
namīnamāyand. Qasam bi-Murabbī-yi kawnayn, har ānī az sijn fī sabīl-i
Allāh niʿmatī ast kih az barā-yi ū shibh va mithl nabūdih va nakhvāhad
būd. Dar har ḥāl bāyad insān bi-ʿadl nāẓir bāshad va bi-inṣāf takallum
namāyad. Az Ḥaqq jalla jalāluhu īn khādim-i fānī sāʾil va āmil kih dar
har ān Ḥaḍrat-i Sulṭān rā muʾayyad dārad va muvaffaq farmāyad. Innahu
ʿalā kulli shayʾin qadīr.

In truth, one must offer thanks concerning the affair of the land of
Ṭāʾ. The religious scholars, who regard themselves as leaders and count
themselves as guides and commanders, are at every moment ablaze with the
fire of hostility. Each day its flame grows greater, its blaze more
apparent, and its harm more intense. They have attacked the people of
Bahāʾ a thousand times, yet their rancor has not subsided in the
slightest, nor has their hatred abated. If one of those wicked souls
obtained temporal power or became a ruler, what would he do? Now reflect
on the leaders of the people and His Majesty the Sovereign. Although he
was not pleased with this community, and the enemies, seeking to display
their service, were occupied night and day with fabrications, His
Majesty the Sovereign nevertheless, of his own noble inclination and
without an intermediary, granted the prisoners freedom and bestowed
release from confinement. The sincere, however, would not exchange a
single moment of imprisonment and affliction suffered in the path of God
for an entire world. I swear by the Educator of both worlds: every
moment of imprisonment in the path of God is a blessing that has never
had, and will never have, a likeness or an equal. In every circumstance,
one must regard matters with justice and speak with fairness. This
passing servant asks and hopes of God, exalted be His glory, that He may
strengthen His Majesty the Sovereign at every moment and grant him
success. He is powerful over all things.

\subsection{14}\label{section-2977}

\begin{otherlanguage}{persian}

اینکه در بارهٴ الواح ابناء ناعق ارض کاف و را مرقوم داشتند در منظر اکبر
تلقاء وجه مالک قدر عرض شد هذا ما نطق به لسان العظمة قوله جلّ و عزّ از قبل
و بعد قرار بر این بود هر نفسی الواح بدیعهٴ منیعه طلب نماید از سماء فضل
الهی نازل شاید غافلین آگاه شوند و نائمین بیدار و مصلحت آنرا از قبل
بایادی امر در هر محلّی گذاشتیم و اذن دادیم بمقتضیات حکمت و مصلحت عمل
نمایند تا حکم مشورت در عالم ثابت و راسخ شود انتهی

\end{otherlanguage}

Īnkih dar bāriy-i alvāḥ-i abnā-yi nāʿiq-i arḍ-i Kāf va Rā marqūm
dāshtand dar Manẓar-i Akbar tilqāʾ-i vajh-i Mālik-i qadar ʿarḍ shud.
Hādhā mā naṭaqa bihi Lisān al-ʿAẓama, qawluhu jalla wa-ʿazza: Az qabl va
baʿd qarār bar īn būd har nafsī alvāḥ-i badīʿiy-i manīʿih ṭalab namāyad
az samāʾ-i faḍl-i ilāhī nāzil, shāyad ghāfilīn āgāh shavand va nāʾimīn
bīdār. Va maṣlaḥat-i ān rā az qabl bi-Ayādī-yi Amr dar har maḥallī
guzāshtīm va idhn dādīm bi-muqtaḍayāt-i ḥikmat va maṣlaḥat ʿamal
namāyand tā ḥukm-i mashvarat dar ʿālam thābit va rāsikh shavad. Intahā.

What you wrote concerning Tablets for the sons of the clamorous one of
the land of Kāf and Rā was presented in the Most Great Scene before the
face of the Lord of destiny. This is what the Tongue of Grandeur
spoke---His words, exalted and mighty is He: ``The arrangement, before
and since, has been that whenever anyone requests wondrous, exalted
Tablets, they are revealed from the heaven of divine grace, so that
perhaps the heedless may become aware and the sleeping may awaken. We
previously left the decision as to its advisability to the Hands of the
Cause in each locality, and gave them permission to act as wisdom and
expediency require, so that the rule of consultation may become
established and firmly rooted in the world.'' End.

\subsection{15}\label{section-2978}

\begin{otherlanguage}{persian}

و فرمودند نعم ما عملا یعنی حضرت افنان و آنمحبوب لوحیکه مخصوص احبّای بهرام
(اب) خواسته بودند در ساحت امنع اقدس اعلی عرض شد و این آیات باهرات از افق
ارادهٴ مالک اسماء و صفات اشراق نمود قوله عزّ بیانه و جلّ برهانه

\end{otherlanguage}

Va farmūdand: Niʿma mā ʿamilā. Yaʿnī Ḥaḍrat-i Afnān va ān maḥbūb. Lawḥī
kih makhṣūṣ-i aḥibbā-yi Bahrām (ab) khvāstih būdand dar sāḥat-i amnaʿ-i
aqdas-i aʿlā ʿarḍ shud va īn āyāt-i bāhirāt az ufuq-i irādiy-i Mālik-i
asmāʾ va ṣifāt ishrāq nimūd. Qawluhu ʿazza bayānuhu wa-jalla burhānuhu:

And He said: ``How well they have both acted!''---meaning the honored
Afnān and you, beloved friend. The request for a Tablet specifically for
the friends of Bahrām (AB) was presented in the most inaccessible, most
holy, most exalted court, and these resplendent verses shone from the
horizon of the will of the Lord of names and attributes. His words,
mighty be His utterance and exalted His proof:

\subsection{16}\label{section-2979}

\begin{otherlanguage}{arabic}

هو النّاطق العلیم الحکیم

\end{otherlanguage}

Huwa al-Nāṭiq al-ʿAlīm al-Ḥakīm.

He is the Speaker, the All-Knowing, the All-Wise.

\subsection{17}\label{section-2980}

\begin{otherlanguage}{persian}

یا اهل بهرام یذکرکم مالک الأنام من شطر سجنه الأعظم و یدعوکم الیه انّه هو
المهیمن علی ما یشآء و هو الفرد الواحد المقتدر العلیم الخبیر انّه اتی
بملکوت الآیات و رایات البیّنات و النّاس اکثرهم من الغافلین قد ارتفع النّدآء
و انجذبت به الأشیآء و القوم اکثرهم من الرّاقدین هذا یوم فیه قصد کلّ قاصد
مقصود الأرض و السّمآء و کلّ منصف الأفق الأعلی و کلّ مخلص شطر محبوبه الأبهی
طوبی لهذا الیوم و لمن عرفه و لمن قام علی خدمة امر اللّه العزیز الحمید انّا
نوصیکم بالاستقامة علی هذا الأمر علی شأن لا یمنعکم نعاق الّذین کفروا باللّه
و اعرضوا عن صراطه المستقیم و نوصیکم بذکر اللّه و ثنائه و التّوجّه الیه
بقلوب نورآء و بما یرتفع به امره المحکم المتین انّ الأمین علیه بهائی ذکرکم
ذکرناکم بالحقّ انّه هو الفضّال الغفور الکریم

\end{otherlanguage}

Yā ahla Bahrām, yadhkurukum Mālik al-anām min shaṭri sijnihi al-aʿẓam
wa-yadʿūkum ilayhi. Innahu Huwa al-Muhaymin ʿalā mā yashāʾ wa-Huwa
al-Fard al-Wāḥid al-Muqtadir al-ʿAlīm al-Khabīr. Innahu atā bi-malakūt
al-āyāt wa-rāyāt al-bayyināt wa-al-nāsu aktharuhum min al-ghāfilīn. Qad
irtafaʿa al-nidāʾ wa-injadhabat bihi al-ashyāʾ wa-al-qawmu aktharuhum
min al-rāqidīn. Hādhā yawmun fīhi qaṣada kullu qāṣidin Maqṣūda al-arḍ
wa-al-samāʾ wa-kullu munṣifin al-Ufuqa al-Aʿlā wa-kullu mukhliṣin shaṭra
Maḥbūbihi al-Abhā. Ṭūbā li-hādhā al-yawm wa-li-man ʿarafahu wa-li-man
qāma ʿalā khidmati Amri Allāh al-ʿAzīz al-Ḥamīd. Innā nūṣīkum
bi-al-istiqāmati ʿalā hādhā al-amr ʿalā shaʾnin lā yamnaʿukum nuʿāqu
alladhīna kafarū bi-Allāh wa-aʿraḍū ʿan ṣirāṭihi al-mustaqīm. Wa-nūṣīkum
bi-dhikri Allāh wa-thanāʾihi wa-al-tawajjuhi ilayhi bi-qulūbin nawrāʾ
wa-bi-mā yartafiʿu bihi amruhu al-muḥkam al-matīn. Inna al-Amīn ʿalayhi
bahāʾī dhakarakum dhakarnākum bi-al-ḥaqq. Innahu Huwa al-Faḍḍāl
al-Ghafūr al-Karīm.

O people of Bahrām, the Lord of humanity remembers you from His Most
Great Prison and calls you to Himself. He is the One who holds sway over
whatever He wills, and He is the Unique, the One, the All-Powerful, the
All-Knowing, the All-Informed. He has come with the kingdom of verses
and the banners of clear proofs, yet most people are heedless. The call
has been raised, and all things have been drawn to it, yet most people
are asleep. This is a day on which every seeker has sought the Desire of
earth and heaven, every fair-minded soul the Most Exalted Horizon, and
every sincere soul the presence of his All-Glorious Beloved. Blessed is
this day, and blessed are those who have recognized it and arisen to
serve the Cause of God, the Almighty, the Praiseworthy. We exhort you to
remain steadfast in this Cause, so that the clamor of those who have
disbelieved in God and turned away from His straight path may not deter
you. We exhort you to remember and praise God, to turn toward Him with
radiant hearts, and to do whatever will exalt His firm and steadfast
Cause. Amīn---upon him be My glory---has mentioned you, and We have
remembered you in truth. He is indeed the Most Bountiful, the
Ever-Forgiving, the Generous.

\subsection{18}\label{section-2981}

\begin{otherlanguage}{arabic}

ای دوستان از قبل و بعد اهل مدن الهی را اخبار نمودیم بظهور ناعقین و انکار
منکرین و مکرهای معتدین حال در بعضی از مدن آنچه را قلم اعلی بآن اخبار
نموده ظاهر و باهر و بصدهزار حیله و مکر در اضلال نفوس مشغولند

\end{otherlanguage}

Ey dūstān, az qabl va baʿd ahl-e modon-e elāhī rā ekh bār namūdīm be
ẓohūr-e nāʿeqīn va enkār-e monkerīn va makr-hā-ye moʿtadīn. Ḥāl dar
baʿḍī az modon ānche rā Qalam-e Aʿlā be ān ekh bār namūde ẓāher va bāher
va be ṣad-hezār ḥīle va makr dar eḍlāl-e nofūs mashghūland.

O friends, time and again We informed the inhabitants of the divine
cities of the emergence of clamorous voices, the denial of the deniers,
and the schemes of the transgressors. Now, in some cities, what the Most
Exalted Pen foretold is manifest and evident, and they are busy leading
souls astray through a hundred thousand tricks and schemes.

\subsection{19}\label{section-2982}

\begin{otherlanguage}{persian}

نفوسیکه از اشراقات انوار آفتاب حقیقت منوّرند و از کوثر معانی و سلسبیل
بیان نوشیده‌اند باستقامتی فائزند که جمیع اهل عالم قادر بر منع نبوده و
نیستند طوبی از برای نفسیکه بحبل عنایت تمسّک نمود و بذیل فضل الهی تشبّث

\end{otherlanguage}

Nofūsī ke az eshrāqāt-e anvār-e Āftāb-e Ḥaqīqat monavvarand va az
Kawthar-e maʿānī va Salsabīl-e bayān nūshīde-and be esteqāmatī fāʾezand
ke jamīʿ-e ahl-e ʿālam qāder bar manʿ nabūde va nīstand. Ṭūbā az barā-ye
nafsī ke be ḥabl-e ʿenāyat tamassok namūd va be dheyl-e faḍl-e elāhī
tashabboth.

Those souls who are illumined by the radiance of the Sun of Truth and
have drunk from the Kawthar of meanings and the Salsabīl of utterance
have attained a steadfastness that all the peoples of the world have
never been, and are not now, able to overcome. Blessed is the soul who
has held fast to the cord of loving-kindness and clung to the hem of
divine grace.

\subsection{20}\label{section-2983}

\begin{otherlanguage}{arabic}

هزار و دویست سال علمای آنجهات و فقهای اطراف ناس را بحضرت قائم دعوت
مینمودند و بکمال اشتیاق از مالک آفاق لقائش را میطلبیدند و بعد از جهد و
طلب و اشتیاق و انتظار ظهورش را انکار کردند و بظلم تمام آن سدرهٴ مبارکهٴ
لاشرقیّه و لاغربیّه را قطع نمودند الا لعنة اللّه علی القوم الظّالمین

\end{otherlanguage}

Hezār va devīst sāl ʿolamā-ye ān jahāt va foqahā-ye aṭrāf nās rā be
Ḥaḍrat-e Qāʾem daʿvat mīnamūdand va be kamāl-e eshtiyāq az Mālek-e Āfāq
leqāyash rā mīṭalabīdand va baʿd az jahd va ṭalab va eshtiyāq va enteẓār
ẓohūrash rā enkār kardand va be ẓolm-e tamām ān Sedre-ye mobārake-ye
lā-sharqiyye va lā-gharbiyye rā qaṭʿ namūdand. A-lā laʿnatu Allāhi ʿalā
al-qawmi al-ẓālimīn.

For twelve hundred years, the learned of those regions and the jurists
of the surrounding lands called the people to the Qāʾim and, with the
utmost longing, implored the Lord of all horizons to grant them reunion
with Him. Yet after all their striving, seeking, longing, and waiting,
they denied His appearance and, with utter cruelty, cut down that
blessed Tree, neither of the East nor of the West. Surely, the curse of
God rests upon the people who do wrong.

\subsection{21}\label{section-2984}

\begin{otherlanguage}{persian}

حال هم از همان نفوس در اضلال خلق کمر شقاوت را محکم بسته‌اند یا حزب اللّه
بحکمت و بیان بر خدمت امر قیام نمائید قیامیکه از اریاح نفاق حرکت ننماید و
نلغزد قدر ایّام را بدانید و بما ینبغی مشغول گردید امروز را مثل و شبهی
نبوده و نیست در کتب الهی بیوم اللّه معروف است کرسیّ ربّ جلّ جلاله بر بقعهٴ
بیضا مقرّ یافته

\end{otherlanguage}

Ḥāl ham az hamān nofūs dar eḍlāl-e khalq kamar-e shaqāvat rā moḥkam
baste-and. Yā ḥizba Allāh, be ḥekmat va bayān bar khedmat-e Amr qiyām
namāyīd, qiyāmī ke az aryāḥ-e nefāq ḥarakat nanamāyad va nalaghzad.
Qadr-e ayyām rā bedānīd va be mā yanbaghī mashghūl gardīd. Emrūz rā
methl va shabahī nabūde va nīst. Dar kotob-e elāhī be Yawmu Allāh maʿrūf
ast. Korsī-ye Rabb, jalla jalāluhu, bar boqʿe-ye beyḍā ma qarr yāfte.

Even now, some of those same souls have firmly girded themselves for
wickedness in leading people astray. O people of God, arise to serve the
Cause with wisdom and utterance, with such resolve that the winds of
dissension can neither shake it nor cause it to falter. Appreciate the
worth of these days and occupy yourselves with what is fitting. This day
has never had, and does not have, an equal or a likeness. In the divine
Books it is known as the Day of God. The throne of the Lord, exalted be
His glory, has been established upon the white Spot.

\subsection{22}\label{section-2985}

\begin{otherlanguage}{persian}

بقلوب نورا و وجوه بیضا بذکر مولی الوری مشغول باشید و از عالم و ما فیه
فارغ و آزاد از حقّ جلّ جلاله میطلبیم شما را از بحر ظهور محروم نسازد و از
آفتاب عنایت منع ننماید و جمیع را مؤیّد فرماید بر امریکه بدوام ملک و ملکوت
پاینده و باقی باشد اوست بر هر شیء قادر و توانا انتهی

\end{otherlanguage}

Be qolūb-e nowrā va vojūh-e beyḍā be dhekr-e Mawlā-ye Varā mashghūl
bāshīd va az ʿālam va mā fīh fāregh va āzād. Az Ḥaqq, jalla jalāluhu,
mīṭalabīm shomā rā az baḥr-e ẓohūr maḥrūm nasāzad va az āftāb-e ʿenāyat
manʿ nanamāyad va jamīʿ rā moʾayyad farmāyad bar amrī ke be davām-e molk
va malakūt pāyande va bāqī bāshad. Ūst bar har sheyʾ qāder va tavānā.
Entahā.

With luminous hearts and radiant faces, occupy yourselves with the
remembrance of the Lord of humanity, detached and free from the world
and all that is in it. We ask God, exalted be His glory, not to deprive
you of the ocean of Revelation or keep you from the sun of
loving-kindness, and to assist everyone in a deed that will endure and
remain as long as the earthly and heavenly kingdoms endure. He has power
and might over all things. End.

\subsection{23}\label{section-2986}

\begin{otherlanguage}{persian}

اینفانی از حقّ جلّ جلاله میطلبد ابصار را بینائی جدید عطا فرماید و آذان را
قوّت اصغاء تازه بخشد تا بچشم جدید و گوش تازه آیات الهی را بشنوند و بیّناتش
را ببینند الیوم اعدای حقّ یعنی نفوسیکه خود را از اهل بیان میشمرند و سبب و
علّت اضلال عبادند آن نفوس بهر لباسی درآیند و بهر صورتی ظاهر بعضی بعمّامه و
عصا و ردا و زهد و تقوی و برخی بکمال محبّت و مودّت و گروهی باسبابهای دیگر
در صدد عباد اللّه بیچاره افتاده‌اند

\end{otherlanguage}

Īn fānī az Ḥaqq, jalla jalāluhu, mīṭalabad abṣār rā bīnāʾī-ye jadīd ʿaṭā
farmāyad va ādhān rā qovvat-e eṣghā-ye tāze bakhshad tā be chashm-e
jadīd va gūsh-e tāze āyāt-e elāhī rā beshnavand va bayyenā tash rā
bebīnand. Al-yawm aʿdā-ye Ḥaqq, yaʿnī nofūsī ke khod rā az ahl-e Bayān
mīshomorand va sabab va ʿellat-e eḍlāl-e ʿebādand, ān nofūs be har
lebāsī darāyand va be har ṣūratī ẓāher, baʿḍī be ʿammāme va ʿaṣā va radā
va zohd va taqvā va barkhī be kamāl-e maḥabbat va mavaddat va gorūhī be
asbāb-hā-ye dīgar dar ṣadad-e ʿebād Allāh-e bīchāre oftāde-and.

This evanescent one asks God, exalted be His glory, to grant eyes a new
vision and bestow upon ears a fresh power of listening, so that, with
new eyes and fresh ears, they may hear the divine verses and behold His
clear proofs. Today the enemies of God---that is, those who count
themselves among the people of the Bayān and are the cause of His
servants being led astray---assume every guise and appear in every form.
Some, with turban, staff, robe, asceticism, and piety; others, with the
utmost love and affection; and still others, by other means, have set
themselves upon the helpless servants of God.

\subsection{24}\label{section-2987}

\begin{otherlanguage}{persian}

جمیع علمای ایران مگر قلیلی با آن زهد و تقواها و موعظه‌ها و ریاضتها فتوی
بر قتل سیّد عالم روح من فی ملکوت الأمر و الخلق فداه دادند میدان آذربایجان
شاهد و گواه و از ورای آن هر ذی بصر و سمعی و از ورای کل حقّ جلّ جلاله باری
بفتوای آن نفوس ظالمه آن طلعة احدیّه را آویختند و شهید نمودند حالهم از
همان قبیل بعضی بوساوس خنّاسیّه و خدعه‌های شیطانیّه پیدا شده‌اند قاتلهم اللّه و
اعاذنا اللّه و ایّاکم من شرّ هؤلآء

\end{otherlanguage}

Jamīʿ-e ʿolamā-ye Īrān magar qalīlī bā ān zohd va taqvā-hā va mowʿeẓe-hā
va riyāḍat-hā fatvā bar qatl-e Seyyed-e ʿĀlam, rūḥu man fī malakūti
al-amri wa-al-khalqi fidāh, dādand. Meydān-e Ādharbāyjān shāhed va govāh
va az varā-ye ān har dhī baṣar va samʿī va az varā-ye koll Ḥaqq, jalla
jalāluhu. Bārī be fatvā-ye ān nofūs-e ẓāleme ān Ṭalʿat-e Aḥadiyye rā
āvīkhtand va shahīd namūdand. Ḥāl ham az hamān qabīl baʿḍī be vasāves-e
khannāsiyye va khodʿe-hā-ye sheyṭāniyye peydā shode-and. Qātalahumu
Allāhu wa aʿādhanā Allāhu wa iyyākum min sharri hāʾulāʾ.

All the learned of Iran, except a few, despite all their asceticism,
piety, sermons, and austerities, issued a ruling for the execution of
the Lord of the world---may the souls of all who dwell in the kingdoms
of revelation and creation be sacrificed for Him. The square in
Ādharbāyjān bears witness, and beyond it everyone endowed with sight and
hearing, and beyond them all, God, exalted be His glory. In short, on
the ruling of those unjust souls, they suspended that Countenance of
divine unity and put Him to death. Now, too, some of the same kind have
appeared with insidious whisperings and satanic deceptions. May God
strike them down, and may God protect us and you from their evil.

\subsection{25}\label{section-2988}

\begin{otherlanguage}{arabic}

لوحی هم مخصوص اماء آن ارض و آن دیار نازل قوله جلّ کبریائه و عزّ بیانه

\end{otherlanguage}

Lowḥī ham makhṣūṣ-e emāʾ-e ān arḍ va ān diyār nāzel. Qawluhu jalla
kibriyāʾuhu wa ʿazza bayānuhu:

A Tablet has also been revealed specifically for the handmaidens of that
land and region. His words---exalted be His majesty and glorious His
utterance:

\subsection{26}\label{section-2989}

\begin{otherlanguage}{arabic}

بنام خداوند دانا و توانا

\end{otherlanguage}

Be nām-e Khodāvand-e dānā va tavānā.

In the name of God, the All-Knowing, the All-Powerful.

\subsection{27}\label{section-2990}

\begin{otherlanguage}{persian}

یا امآء اللّه فی ارض البآء و الهآء و ضواحیها در امواج بحر فضل حقّ تفکّر
نمائید در سجن ساکن و شما را لوجه اللّه بافق اعلی دعوت مینماید و بصراط
مستقیم دلالت میفرماید از عالم غیب بشهود آمده تا کل را بفیوضات فیّاض حقیقی
فائز فرماید بحبل متین الهی تمسّک نمائید و بافق ظهورش توجّه کنید امروز هر
نفسی بعرفان حضرت دوست فائز شد او بکلّ خیر فائز است از بساط عالم بگذرید و
بکمال نشاط و انبساط در بساط حقّ درآئید هر نفسی الیوم بافق اعلی ناظر و
بذکر مالک اسماء ذاکر او از اهل بساط لدی اللّه محسوب آنچه مشاهده میشود
عنقریب معدوم جهد نمائید تا بعملی فائز شوید که عرف بقا از او استشمام شود
شکر کنید مقصود عالم را که شما را مؤیّد فرمود بر عرفان امریکه عالم از او
محجوب انّه هو الفضّال الکریم الحمد له اذ هو ربّ العالمین و مقصود العارفین
انتهی

\end{otherlanguage}

Yā imāʾa Allāhi fī arḍi al-Bāʾi wa-al-Hāʾi wa ḍawāḥīhā, dar amvāj-e
baḥr-e faḍl-e Ḥaqq tafakkor namāyīd. Dar sejn sāken va shomā rā li-wajhi
Allāh be Ofoq-e Aʿlā daʿvat mīnamāyad va be ṣerāṭ-e mostaqīm delālat
mīfarmāyad. Az ʿālam-e gheyb be shohūd āmade tā koll rā be foyūḍāt-e
Fayyāḍ-e Ḥaqīqī fāʾez farmāyad. Be ḥabl-e matīn-e elāhī tamassok namāyīd
va be ofoq-e ẓohūrash tavajjoh konīd. Emrūz har nafsī be ʿerfān-e
Ḥaḍrat-e Dūst fāʾez shod ū be koll-e kheyr fāʾez ast. Az besāṭ-e ʿālam
bogdharīd va be kamāl-e nashāṭ va enbesāṭ dar besāṭ-e Ḥaqq darāʾīd. Har
nafsī al-yawm be Ofoq-e Aʿlā nāẓer va be dhekr-e Mālek-e Asmāʾ dhāker ū
az ahl-e besāṭ ladā Allāh maḥsūb. Ānche moshāhede mīshavad ʿan-qarīb
maʿdūm. Jahd namāyīd tā be ʿamalī fāʾez shavīd ke ʿarf-e baqā az ū
esteshmām shavad. Shokr konīd Maqṣūd-e ʿĀlam rā ke shomā rā moʾayyad
farmūd bar ʿerfān-e amrī ke ʿālam az ū maḥjūb. Innahu huwa al-Faḍḍālu
al-Karīm. Al-ḥamdu lahu idh huwa Rabbu al-ʿālamīn wa Maqṣūdu al-ʿārifīn.
Entahā.

O handmaidens of God in the land of Bāʾ and Hāʾ and its surroundings,
reflect upon the waves of the ocean of God's grace. He dwells in prison
and, for the sake of God, calls you to the Most Exalted Horizon and
guides you to the straight path. He has come from the unseen world into
the visible so that He may enable all to partake of the outpourings of
the true Bestower. Hold fast to the strong cord of God and turn toward
the horizon of His Revelation. Whoever today attains recognition of the
Beloved has attained all good. Pass beyond the court of the world and,
with the utmost joy and gladness, enter the court of God. Whoever today
looks toward the Most Exalted Horizon and remembers the Lord of Names is
accounted by God among those admitted to His court. Whatever is seen
will soon pass away. Strive to perform a deed from which the fragrance
of eternity may be perceived. Give thanks to the Desire of the world,
who has enabled you to recognize a Cause from which the world is veiled.
He is indeed the Most Bountiful, the Generous. Praise be to Him, for He
is the Lord of the worlds and the Desire of those who know Him. End.

\subsection{28}\label{section-2991}

\begin{otherlanguage}{persian}

للّه الحمد انوار آفتاب فضل الهی شامل حال امائش شد جمیع عالم مشتاق استماع
کلمهٴ یا عبدی و حال امآء اللّه بکلمهٴ مبارکهٴ یا امتی فائزند این فضل
عظیمست و بزرگ طوبی از برای نفسیکه عارف شد و شناخت و ویل للغافلین

\end{otherlanguage}

Li-llāhi al-ḥamd, anvār-e āftāb-e faḍl-e elāhī shāmel-e ḥāl-e emāʾash
shod. Jamīʿ-e ʿālam moshtāq-e estemāʿ-e kaleme-ye yā ʿabdī va ḥāl emāʾ
Allāh be kaleme-ye mobārake-ye yā amatī fāʾezand. Īn faḍl ʿaẓīm ast va
bozorg. Ṭūbā az barā-ye nafsī ke ʿāref shod va shenākht, wa waylun
li-l-ghāfilīn.

Praise be to God, the rays of the sun of divine grace have embraced His
handmaidens. The whole world longs to hear the words ``O My servant,''
and now the handmaidens of God have been favored with the blessed words
``O My handmaiden.'' This is a mighty and great bounty. Blessed is the
soul who has understood and recognized it, and woe to the heedless.

\subsection{29}\label{section-2992}

\begin{otherlanguage}{persian}

اینکه ذکر جناب استاد احمد (ک‌ف) علیه بهآء اللّه را نمودند در ساحت امنع
اقدس مالک الرّقاب عرض شد هذا ما نزّل فی الجواب قوله جلّ جلاله و عمّ نواله

\end{otherlanguage}

Īnke dhekr-e Janāb-e Ostād Aḥmad (K-F), ʿalayhi Bahāʾu Allāh, rā
namūdand dar sāḥat-e amnaʿ-e aqdas-e Mālek-e Reqāb ʿarḍ shod. Hādhā mā
nuzzila fī al-jawāb. Qawluhu jalla jalāluhu wa ʿamma nawāluhu:

The mention that was made of Ustād Aḥmad (K-F), upon him be the glory of
God, was submitted at the most inviolate, most holy court of the Lord of
all. This is what was revealed in reply. His words---exalted be His
glory and all-embracing His bounty:

\subsection{30}\label{section-2993}

\begin{otherlanguage}{arabic}

هو السّمیع البصیر

\end{otherlanguage}

Huwa al-Samīʿu al-Baṣīr.

He is the All-Hearing, the All-Seeing.

\subsection{31}\label{section-2994}

\begin{otherlanguage}{persian}

یا احمد یذکرک مقصود العالم و یبشّرک بعنایة اللّه المهیمن القیّوم طوبی لوجه
توجّه الیه و لقلب اقبل الی شطره و لبصر رأی الأفق الأعلی و للسان نطق امام
الوجوه انّه لا اله الّا هو الفرد الواحد العزیز الودود انّا جدّدنا العالم من
نفحات الوحی و اظهرنا ربیع المعانی و القوم اکثرهم لا یفقهون قد ماج بحر
البیان و الرّحمن استوی علی عرش البرهان و النّاس هم لا یعرفون قد نبذوا
الههم و اخذوا اهوائهم الا انّهم لا یعلمون طوبی لمن کسّر اصنام الظّنون و آمن
باللّه ربّ ما کان و ما یکون یا احمد قد حضر کتاب من الأمین و کان فیه ذکرک و
الّذین آمنوا هناک ذکرناکم من هذا المقام المحمود ان اعرف هذا المقام الأعلی
لعمر اللّه لا یعادل بذکری ما یشهد و یری یشهد بذلک من عنده الرّحیق المختوم

\end{otherlanguage}

Yā Aḥmad, yadhkuruka Maqṣūdu al-ʿālam wa yubashshiruka bi-ʿināyati
Allāhi al-Muhaymini al-Qayyūm. Ṭūbā li-wajhin tawajjaha ilayhi wa
li-qalbin aqbala ilā shaṭrihi wa li-baṣarin raʾā al-ufuqa al-aʿlā wa
li-lisānin naṭaqa amāma al-wujūhi innahu lā ilāha illā huwa al-Fardu
al-Wāḥidu al-ʿAzīzu al-Wadūd. Innā jaddadnā al-ʿālama min nafaḥāti
al-waḥyi wa aẓharnā rabīʿa al-maʿānī wa-al-qawmu aktharuhum lā yafqahūn.
Qad māja baḥru al-bayāni wa-al-Raḥmānu istawā ʿalā ʿarshi al-burhāni
wa-al-nāsu hum lā yaʿrifūn. Qad nabadhū ilāhahum wa akhadhū ahwāʾahum
a-lā innahum lā yaʿlamūn. Ṭūbā li-man kassara aṣnāma al-ẓunūni wa āmana
bi-Allāhi Rabbi mā kāna wa mā yakūn. Yā Aḥmad, qad ḥaḍara kitābun mina
al-Amīn wa kāna fīhi dhikruka wa-alladhīna āmanū hunāk. Dhakarnākum min
hādhā al-maqāmi al-maḥmūd. An iʿraf hādhā al-maqāma al-aʿlā. La-ʿamru
Allāh, lā yuʿādilu bi-dhikrī mā yushhadu wa yurā. Yashhadu bi-dhālika
man ʿindahu al-raḥīqu al-makhtūm.

O Aḥmad, the Desire of the world remembers you and gives you the glad
tidings of the loving-kindness of God, the All-Protecting, the
Self-Subsisting. Blessed is the face that has turned toward Him, the
heart that has inclined in His direction, the eye that has beheld the
Most Exalted Horizon, and the tongue that has proclaimed before all
faces, ``There is no god but Him, the Unique, the One, the Almighty, the
All-Loving.'' We have indeed renewed the world through the breaths of
revelation and brought forth the springtime of meanings, yet most of the
people do not understand. The ocean of utterance has surged, and the
All-Merciful has ascended the throne of proof, yet the people do not
recognize Him. They have cast aside their God and taken hold of their
own desires. Surely, they do not know. Blessed is the one who has
shattered the idols of idle suppositions and believed in God, the Lord
of all that has been and all that will be. O Aḥmad, a letter has arrived
from Amīn, in which you and those who have believed in that land were
mentioned. We have remembered you from this praiseworthy station.
Recognize this most exalted station. By the life of God, nothing
witnessed or seen can equal My remembrance. To this bears witness the
One with whom is the sealed wine.

\subsection{32}\label{section-2995}

\begin{otherlanguage}{persian}

یا سیّد قبل هاشم ان استمع ندآء اللّه انّه ارتفع مرّة فی الطّآء و اخری فی
الزّورآء و تارةً فی ارض السّرّ و فی هذا المقام الّذی سمّی بالأسمآء الحسنی من
لدی اللّه مولی الوری یشهد بذلک من شهد انّه لا اله الّا هو العزیز المحبوب ان
افرح بذکری ثمّ انطق بثنآء ربّک فی العشیّ و البکور

\end{otherlanguage}

Yā Sayyid qabla Hāshim, an istamiʿ nidāʾa Allāh. Innahu irtafaʿa
marratan fī al-Ṭāʾi wa ukhrā fī al-Zawrāʾi wa tāratan fī Arḍi al-Sirri
wa fī hādhā al-maqāmi alladhī summiya bi-al-asmāʾi al-ḥusnā min ladā
Allāhi Mawlā al-warā. Yashhadu bi-dhālika man shahida annahu lā ilāha
illā huwa al-ʿAzīzu al-Maḥbūb. An ifraḥ bi-dhikrī thumma inṭiq
bi-thanāʾi Rabbika fī al-ʿashiyyi wa-al-bukūr.

O Siyyid, whose name is followed by Hāshim, listen to the call of God.
It rose once in Ṭāʾ, another time in Zawrāʾ, and again in the Land of
Mystery and in this station, which has been designated by the most
beautiful names by God, the Lord of humanity. To this bears witness the
One who has testified that there is no god but Him, the Almighty, the
Beloved. Rejoice in My remembrance, then speak the praise of your Lord
at evening and at dawn.

\subsection{33}\label{section-2996}

\begin{otherlanguage}{arabic}

و نذکر من سمّی بسیّد قبل محمّد لیفرح و یکون من الشّاکرین انّا ذکرنا الّذین
اقبلوا الی الأفق الأعلی و نذکرهم رحمة من عندنا لتبقی به اذکارهم بدوام
اسمآء اللّه الحسنی و صفاته العلیا لا اله الّا هو المقتدر العلیم الخبیر انّ
اللّه جعل ذکره کنزاً لأحبّائه الّذین اقبلوا الیه بوجوه بیضآء و عملوا بما
امروا به فی کتابه المبین یا اولیآء الرّحمن انّا نوصیکم بحفظ کنوزکم کذلک
یعظکم من عنده علم کلّشیء و یطوف حوله الملأ الأعلی فی مقامه الکریم و نذکر
اولیائی هناک الّذین ما نقضوا المیثاق اذ ارتفع النّعاق الا انّهم من الرّاسخین
طوبی لعبد ما زلّته شبهات الفجّار و ما منعته اشارات المغلّین الّذین یغلی فی
صدورهم قدر البغضآء و یری من وجوههم غبرة الجحیم البهآء المشرق من افق سمآء
عنایتی علیکم یا حزب اللّه هناک و علی الّذین یطوفون العرش فی مقامی الرّفیع
انتهی

\end{otherlanguage}

Wa nadhkuru man summiya bi-Sayyid qabla Muḥammad li-yafraḥa wa yakūna
mina al-shākirīn. Innā dhakarnā alladhīna aqbalū ilā al-ufuqi al-aʿlā wa
nadhkuruhum raḥmatan min ʿindinā li-tabqā bihi adhkāruhum bi-dawāmi
asmāʾi Allāhi al-ḥusnā wa ṣifātihi al-ʿulyā. Lā ilāha illā huwa
al-Muqtadiru al-ʿAlīmu al-Khabīr. Inna Allāha jaʿala dhikrahu kanzan
li-aḥibbāʾihi alladhīna aqbalū ilayhi bi-wujūhin bayḍāʾa wa ʿamilū bi-mā
umirū bihi fī kitābihi al-mubīn. Yā awliyāʾa al-Raḥmān, innā nūṣīkum
bi-ḥifẓi kunūzikum. Kadhālika yaʿiẓukum man ʿindahu ʿilmu kulli shayʾin
wa yaṭūfu ḥawlahu al-malaʾu al-aʿlā fī maqāmihi al-karīm. Wa nadhkuru
awliyāʾī hunāka alladhīna mā naqaḍū al-mīthāqa idh irtafaʿa al-nuʿāq.
A-lā innahum mina al-rāsikhīn. Ṭūbā li-ʿabdin mā zallat-hu shubuhātu
al-fujjāri wa mā manaʿat-hu ishārātu al-mughillīna alladhīna yaghlī fī
ṣudūrihim qidru al-baghḍāʾi wa yurā min wujūhihim ghubratu al-jaḥīm.
Al-Bahāʾu al-mushriqu min ufuqi samāʾi ʿināyatī ʿalaykum yā ḥizba Allāhi
hunāka wa ʿalā alladhīna yaṭūfūna al-ʿarsha fī maqāmī al-rafīʿ. Intahā.

And We remember the one named Siyyid followed by Muḥammad, that he may
rejoice and be among the thankful. We have remembered those who have
turned toward the Most Exalted Horizon, and We remember them as a mercy
from Us, so that through it their remembrance may endure as long as
God's most beautiful names and most exalted attributes endure. There is
no god but Him, the All-Powerful, the All-Knowing, the All-Informed. God
has made His remembrance a treasure for His loved ones who have turned
toward Him with radiant faces and acted according to what they were
commanded in His perspicuous Book. O loved ones of the All-Merciful, We
exhort you to safeguard your treasures. Thus admonishes you the One who
possesses the knowledge of all things and around whom the Concourse on
high circles in His noble station. And We remember My loved ones in that
land who did not break the Covenant when the clamor arose. Surely, they
are among the steadfast. Blessed is the servant whom the doubts raised
by the wicked have not caused to slip, nor the insinuations of the
spiteful deterred---those in whose breasts the cauldron of hatred boils
and upon whose faces the dust of hell can be seen. The glory shining
from the horizon of the heaven of My loving-kindness rest upon you, O
people of God in that land, and upon those who circle the throne in My
exalted station. End.

\subsection{34}\label{section-2997}

\begin{otherlanguage}{persian}

و هذا ما نزّل لجناب استاد باقر علیه بهآء اللّه قوله جلّ جلاله و عزّ کبریائه
و عظم شأنه

\end{otherlanguage}

Wa hādhā mā nuzzila li-Janāb-e Ostād Bāqer, ʿalayhi Bahāʾu Allāh.
Qawluhu jalla jalāluhu wa ʿazza kibriyāʾuhu wa ʿaẓuma shaʾnuhu:

And this is what was revealed for Ustād Bāqir, upon him be the glory of
God. His words---exalted be His glory, glorious His majesty, and great
His station:

\subsection{35}\label{section-2998}

\begin{otherlanguage}{arabic}

بسمی المشفق الکریم

\end{otherlanguage}

Bi-smī al-Mushfiqi al-Karīm.

In My name, the Compassionate, the Generous.

\subsection{36}\label{section-2999}

\begin{otherlanguage}{persian}

سبحان الّذی سخّر الملک و الملکوت امراً من عنده و هو القویّ الغالب القدیر
یفعل ما یشآء و یحکم ما یرید و هو الفرد الواحد العلیم الحکیم یا باقر انّا
ذکرناک من قبل و فی هذا الحین لعمری لو تعرف مقام ذکری لیأخذک الفرح علی
شأن لا تحزنک شئونات البشر و لا سطوة الغافلین کن مستقیماً علی امر ربّک و
ناطقاً بثنائه الجمیل هذا ما ینفعک و اولیائی فی الآخرة و الأولی یشهد بذلک
من استوی علی عرشه العظیم کن قائماً علی خدمة امری و مقبلاً بقلبک الی افقی
انّه یؤیّد من اقبل الیه بقلب منیر قد شهد لک قلمی الأعلی ان اعرف و کن من
الشّاکرین ایّاک ان تحجبک شئونات الفجّار الّذین ینطقون بأهوائهم و یمنعون
النّاس عن صراطی المستقیم البهآء علیک و علی من سمع النّدآء و اجاب ربّ الکرسیّ
الرّفیع انتهی

\end{otherlanguage}

Subḥāna alladhī sakhkhara al-mulka wa-al-malakūta amran min ʿindihi wa
huwa al-Qawiyyu al-Ghālibu al-Qadīr. Yafʿalu mā yashāʾu wa yaḥkumu mā
yurīdu wa huwa al-Fardu al-Wāḥidu al-ʿAlīmu al-Ḥakīm. Yā Bāqir, innā
dhakarnāka min qablu wa fī hādhā al-ḥīn. La-ʿamrī, law taʿrifu maqāma
dhikrī la-yaʾkhudhuka al-faraḥu ʿalā shaʾnin lā tuḥzinuka shuʾūnātu
al-bashari wa lā saṭwatu al-ghāfilīn. Kun mustaqīman ʿalā amri Rabbika
wa nāṭiqan bi-thanāʾihi al-jamīl. Hādhā mā yanfaʿuka wa awliyāʾī fī
al-ākhirati wa-al-ūlā. Yashhadu bi-dhālika man istawā ʿalā ʿarshihi
al-ʿaẓīm. Kun qāʾiman ʿalā khidmati amrī wa muqbilan bi-qalbika ilā
ufuqī. Innahu yuʾayyidu man aqbala ilayhi bi-qalbin munīr. Qad shahida
laka qalamī al-aʿlā. An iʿraf wa kun mina al-shākirīn. Iyyāka an
taḥjubaka shuʾūnātu al-fujjāri alladhīna yanṭiqūna bi-ahwāʾihim wa
yamnaʿūna al-nāsa ʿan ṣirāṭī al-mustaqīm. Al-Bahāʾu ʿalayka wa ʿalā man
samiʿa al-nidāʾa wa ajāba Rabba al-kursiyyi al-rafīʿ. Intahā.

Glory be to Him who has subjected the earthly and heavenly kingdoms to a
command from Himself. He is the Strong, the All-Subduing, the
All-Powerful. He does whatever He wills and ordains whatever He desires,
and He is the Unique, the One, the All-Knowing, the All-Wise. O Bāqir,
We have remembered you before and at this time. By My life, if you knew
the station of My remembrance, such joy would seize you that neither
human affairs nor the onslaught of the heedless could grieve you. Be
steadfast in the Cause of your Lord and speak His glorious praise. This
is what will benefit you and My loved ones in the world to come and in
this world. To this bears witness the One who has ascended His mighty
throne. Arise to serve My Cause and turn your heart toward My horizon.
He indeed assists whoever turns toward Him with a luminous heart. My
Most Exalted Pen has borne witness in your favor. Recognize this and be
among the thankful. Beware lest the doings of the wicked veil
you---those who speak according to their own desires and keep people
from My straight path. Glory rest upon you and upon whoever has heard
the call and answered the Lord of the lofty throne. End.

\subsection{37}\label{section-3000}

\begin{otherlanguage}{arabic}

و هذا ما نزّل لجناب ابوالقاسم علیه بهآء اللّه قوله جلّ جلاله و عزّ شأنه

\end{otherlanguage}

Wa-hādhā mā nuzzila li-janāb Abī al-Qāsim ʿalayhi bahāʾ Allāh, qawluhu
jalla jalāluhu wa-ʿazza shaʾnuhu.

And this is what was revealed for Abū al-Qāsim, upon him be the glory of
God. His words---glorious is His majesty and exalted His station:

\subsection{38}\label{section-3001}

\begin{otherlanguage}{arabic}

بسمی النّاطق الأمین

\end{otherlanguage}

Bi-smī al-Nāṭiq al-Amīn.

In My name, the Speaker, the Trustworthy.

\subsection{39}\label{section-3002}

\begin{otherlanguage}{arabic}

یا اباالقاسم یذکرک المظلوم فی مقام یطوفه الملأ الأعلی و یوصیک بتقوی اللّه
انّ ربّک لهو النّاصح العلیم انّه یذکر من اقبل الیه و یرید من اراده انّه خلق
الوفآء و جعله من اسمائه الحسنی فی کتابه المبین طوبی لمن تمسّک به انّه من
فوارس مضمار البیان لدی الرّحمن یشهد بذلک کلّ عارف بصیر ان افرح بذکری ایّاک
ثمّ اشکر ربّک الّذی خلقک و ایّدک و هداک الی افقه المنیر انّا نذکر اولیائی
هناک و نبشّرهم برحمتی الّتی سبقت من فی السّموات و الأرضین قد ذکرک الأمین
ذکرناک ان اشکر من انقذک و نجّاک و انزل لک ما یجد منه المقرّبون عرف رحمة
ربّک الغفور الرّحیم انتهی

\end{otherlanguage}

Yā Abā al-Qāsim, yadhkuruka al-Maẓlūm fī maqām yaṭūfuhu al-malaʾ al-aʿlā
wa-yūṣīka bi-taqwā Allāh. Inna rabbaka la-huwa al-Nāṣiḥ al-ʿAlīm. Innahu
yadhkuru man aqbala ilayhi wa-yurīdu man arādahu. Innahu khalaqa
al-wafāʾ wa-jaʿalahu min asmāʾihi al-ḥusnā fī kitābihi al-mubīn. Ṭūbā
li-man tamassaka bihi, innahu min fawāris miḍmār al-bayān ladā
al-Raḥmān. Yashhadu bi-dhālika kullu ʿārif baṣīr. An ifraḥ bi-dhikrī
iyyāka thumma ushkur rabbaka alladhī khalaqaka wa-ayyadaka wa-hadāka ilā
ufuqihi al-munīr. Innā nadhkuru awliyāʾī hunāka wa-nubashshiruhum
bi-raḥmatī allatī sabaqat man fī al-samāwāt wa-al-araḍīn. Qad dhakaraka
al-Amīn, dhakarnāka. An ishkur man anqadhaka wa-najjāka wa-anzala laka
mā yajidu minhu al-muqarrabūn ʿarfa raḥmati rabbika al-Ghafūr al-Raḥīm.
Intahā.

O Abū al-Qāsim, the Wronged One remembers you in a station around which
the heavenly Concourse circles, and counsels you to be mindful of God.
Your Lord is indeed the Counselor, the All-Knowing. He remembers whoever
turns toward Him and desires whoever desires Him. He created
faithfulness and made it one of His most beautiful names in His clear
Book. Blessed is whoever holds fast to it; in the presence of the
All-Merciful, he is among the horsemen in the arena of utterance. Every
discerning knower bears witness to this. Rejoice in My remembrance of
you, then thank your Lord, who created you, strengthened you, and guided
you to His luminous horizon. We remember My loved ones there and give
them the glad tidings of My mercy, which has preceded all who are in the
heavens and the earths. The Trustworthy One has mentioned you, and We
have remembered you. Give thanks to Him who rescued and saved you and
revealed for you that from which those brought near perceive the
fragrance of the mercy of your Lord, the Ever-Forgiving, the Merciful.
End.

\subsection{40}\label{section-3003}

\begin{otherlanguage}{arabic}

و هذا ما نزّل لجناب ملّا حسن علیه بهآء اللّه قوله عظمت آیاته و جلّ سلطانه

\end{otherlanguage}

Wa-hādhā mā nuzzila li-janāb Mullā Ḥasan ʿalayhi bahāʾ Allāh, qawluhu
ʿaẓumat āyātuhu wa-jalla sulṭānuhu.

And this is what was revealed for Mullā Ḥasan, upon him be the glory of
God. His words---great are His signs and glorious His sovereignty:

\subsection{41}\label{section-3004}

\begin{otherlanguage}{arabic}

هو الصّراط المستقیم

\end{otherlanguage}

Huwa al-ṣirāṭ al-mustaqīm.

He is the straight path.

\subsection{42}\label{section-3005}

\begin{otherlanguage}{persian}

کتاب انزله المظلوم و فیه یذکر من اقبل الی اللّه مولی العالم و اجاب اذ
ارتفع ندائه من هذا المقام الّذی جعله من اعلی المقام قد خلقت الأبصار
لمشاهدة افقی و الآذان لاصغآء ندائی و اللّسان لذکری الّذی به اشرقت الآفاق
یا حسن ان استمع ندآء مکلّم الطّور فی هذا الظّهور ثمّ اشکر ربّک بما ایّدک علی
عرفان هذا الأمر الّذی به زلّت الأقدام انّا سمّینا هذا الیوم فی کتب القبل
بیوم اللّه ولکنّ القوم اکثرهم فی غفلة و حجاب قد منعتهم السّبحات عن النّظر
الی الأفق الأعلی و الشّبهات عن الّذی اتی من سمآء القوّة برایات الآیات کن
ناطقاً بذکر ربّک و عاملاً بما امرت به فی الزّبر و الألواح طوبی لذی ظمأ قصد
الفرات و لذی وجه توجّه الی اللّه مالک الأدیان قل ضعوا الظّنون و الأوهام
مسرعین الی ملکوت الایقان بالرّوح و الرّیحان کذلک هطلت من سمآء عنایتی امطار
رحمتی نعیماً لمن فاز و ویل لکلّ غافل مرتاب البهآء من لدنّا علی الّذین ما
حجّبتهم حجبات العلمآء و ما خوّفتهم سطوة الّذین اعرضوا عن المبدء و المآب

\end{otherlanguage}

Kitāb anzalahu al-Maẓlūm wa-fīhi yadhkuru man aqbala ilā Allāh mawlā
al-ʿālam wa-ajāba idh irtafaʿa nidāʾuhu min hādhā al-maqām alladhī
jaʿalahu min aʿlā al-maqām. Qad khuliqat al-abṣār li-mushāhadati ufuqī
wa-al-ādhān li-iṣghāʾi nidāʾī wa-al-lisān li-dhikrī alladhī bihi
ashraqat al-āfāq. Yā Ḥasan, an istamiʿ nidāʾa Mukallim al-Ṭūr fī hādhā
al-ẓuhūr thumma ushkur rabbaka bi-mā ayyadaka ʿalā ʿirfān hādhā al-amr
alladhī bihi zallat al-aqdām. Innā sammaynā hādhā al-yawm fī kutub
al-qabl bi-yawm Allāh wa-lākinna al-qawm aktharuhum fī ghafla wa-ḥijāb.
Qad manaʿathumu al-subuḥāt ʿan al-naẓar ilā al-ufuq al-aʿlā
wa-al-shubuhāt ʿan alladhī atā min samāʾ al-quwwa bi-rāyāt al-āyāt. Kun
nāṭiqan bi-dhikr rabbika wa-ʿāmilan bi-mā umirta bihi fī al-zubur
wa-al-alwāḥ. Ṭūbā li-dhī ẓamaʾ qaṣada al-Furāt wa-li-dhī wajh tawajjaha
ilā Allāh Mālik al-Adyān. Qul ḍaʿū al-ẓunūn wa-al-awhām musriʿīn ilā
malakūt al-īqān bi-al-rūḥ wa-al-rayḥān. Kadhālika haṭalat min samāʾ
ʿināyatī amṭār raḥmatī. Naʿīman li-man fāz wa-wayl li-kulli ghāfil
murtāb. Al-bahāʾ min ladunnā ʿalā alladhīna mā ḥajjabathum ḥujubāt
al-ʿulamāʾ wa-mā khawwafathum saṭwat alladhīna aʿraḍū ʿan al-Mabdaʾ
wa-al-Maʾāb.

A Book revealed by the Wronged One, in which He remembers whoever has
turned toward God, the Lord of the world, and answered when His call
arose from this station, which He has made one of the loftiest stations.
Eyes were created to behold My horizon, ears to hear My call, and the
tongue to remember Me, through which the horizons have been illumined. O
Ḥasan, listen to the call of Him who spoke on Sinai in this Revelation,
then thank your Lord for enabling you to recognize this Cause, through
which feet have slipped. We named this day the Day of God in the Books
of the past, yet most people remain heedless and veiled. The veils of
glory have prevented them from looking toward the Most Exalted Horizon,
and doubts have kept them from Him who has come from the heaven of power
with the banners of the verses. Speak in remembrance of your Lord and
act according to what you have been commanded in the Scriptures and
Tablets. Blessed is the thirsty one who has sought the Euphrates, and
the one who has turned his face toward God, the Sovereign of all
religions. Say: Set aside conjectures and vain imaginings, hastening
toward the kingdom of certitude with joy and gladness. Thus have the
rains of My mercy poured from the heaven of My loving care. Bliss to
whoever has attained, and woe to every heedless doubter. Glory from Our
presence rest upon those whom the veils of the learned have not veiled
and whom the might of those who have turned away from the Origin and the
Return has not frightened.

\subsection{43}\label{section-3006}

\begin{otherlanguage}{arabic}

و هذا ما نزّل لجناب غلام رضا من لدی اللّه مولی الأسمآء قوله جلّ جلاله و عزّ
بیانه

\end{otherlanguage}

Wa-hādhā mā nuzzila li-janāb Ghulām Riḍā min ladā Allāh Mawlā al-Asmāʾ,
qawluhu jalla jalāluhu wa-ʿazza bayānuhu.

And this is what was revealed for Ghulām Riḍā from God, the Lord of
Names. His words---glorious is His majesty and exalted His utterance:

\subsection{44}\label{section-3007}

\begin{otherlanguage}{arabic}

بسمی المهیمن علی الأسمآء

\end{otherlanguage}

Bi-smī al-Muhaymin ʿalā al-asmāʾ.

In My name, the Guardian over all names.

\subsection{45}\label{section-3008}

\begin{otherlanguage}{persian}

یا غلام رضا عالم را دخان ریب و غبار اوهام از مالک انام منع نموده حجبات
نفس و هوی بشأنی قوم را اخذ کرده که خیر و شر و معروف و منکر را از هم تمیز
نمیدهند اهل شرکند و خود را اهل توحید میدانند و عبدهٴ اصنامند و خود را
اهل مقام میشمرند غفلت بمقامی رسیده که نعمت ابدی را بآلاء فانیه تبدیل
نموده‌اند آنچه حقّ جلّ جلاله از برای کل خواسته از آن گذشته‌اند و بمشتهیات
نفس و هوی مشغولند افّ لهم و بما اکتسبت ایدیهم حمد کن مقصود عالمیانرا که
ترا مؤیّد فرمود بر اقبال قدر این عنایت بزرگ را بدان و لآلی محبّت الهی را
از ابصار خائنه و نفوس سارقه باسمش محفوظ دار اینست نصیحت قلم اعلی در این
لیلهٴ ظلما اسمع و کن من الشّاکرین و الحمد لله ربّ العالمین

\end{otherlanguage}

Yā Gholām Reḍā, ʿālam rā dokhān-e rayb va ghobār-e awhām az Mālek-e anām
manʿ namūde. Ḥojobāt-e nafs va havā be-shaʾnī qawm rā akhdh karde ke
khayr va sharr va maʿrūf va monkar rā az ham tamīz nemīdahand. Ahl-e
sherkand va khod rā ahl-e tawḥīd mīdānand va ʿabade-ye aṣnāmand va khod
rā ahl-e maqām mīshomorand. Ghaflat be-maqāmī resīde ke neʿmat-e abadī
rā be-ālā-ye fāniye tabdīl namūde-and. Ānche Ḥaqq jalla jalāluhu az
barā-ye koll khāste az ān godhashte-and va be-moshtahayāt-e nafs va havā
mashghūland. Uffin lahum wa-bi-mā iktasabat aydīhim. Ḥamd kon Maqṣūd-e
ʿālamiyān rā ke to rā moʾayyad farmūd bar eqbāl. Qadr-e īn ʿenāyat-e
bozorg rā bedān va laʾālī-ye maḥabbat-e elāhī rā az abṣār-e khāʾene va
nofūs-e sāreqe be-esmash maḥfūẓ dār. Īnast naṣīḥat-e Qalam-e Aʿlā dar īn
layle-ye ẓalmā. Ismaʿ wa-kun min al-shākirīn wa-al-ḥamdu li-llāh Rabb
al-ʿālamīn.

O Ghulām Riḍā, the smoke of doubt and the dust of vain imaginings have
kept the world from the Sovereign of humankind. The veils of self and
desire have so seized the people that they cannot distinguish good from
evil or right from wrong. They associate partners with God yet consider
themselves believers in His unity; they worship idols yet count
themselves among those of spiritual rank. Heedlessness has reached such
a degree that they have exchanged the everlasting blessing for fleeting
favors. They have forsaken whatever God---glorious is His majesty---has
desired for all, and are occupied with the cravings of self and desire.
Shame on them and on what their hands have earned! Praise the Desire of
all peoples, who has enabled you to turn toward Him. Appreciate this
great favor, and in His name guard the pearls of divine love from
treacherous eyes and thieving souls. This is the counsel of the Most
Exalted Pen in this dark night. Listen and be among the thankful. Praise
be to God, the Lord of all worlds.

\subsection{46}\label{section-3009}

\begin{otherlanguage}{arabic}

و هذا ما نزّل لجناب علی اکبر (س‌م‌س) من لدی اللّه مالک القدر قوله تبارک و
تعالی

\end{otherlanguage}

Wa-hādhā mā nuzzila li-janāb ʿAlī Akbar (s-m-s) min ladā Allāh Mālik
al-Qadar, qawluhu tabāraka wa-taʿālā.

And this is what was revealed for ʿAlī Akbar (s-m-s) from God, the
Sovereign of destiny. His words---blessed and exalted is He:

\subsection{47}\label{section-3010}

\begin{otherlanguage}{arabic}

هو السّامع و هو المجیب

\end{otherlanguage}

Huwa al-Sāmiʿ wa-huwa al-Mujīb.

He is the Hearer, and He is the Answerer.

\subsection{48}\label{section-3011}

\begin{otherlanguage}{persian}

یا حزب اللّه امروز بر کل ذکر حقّ جلّ جلاله و خدمت امر لازم در قرون و اعصار
قوم منتظر این یوم مبارک بوده‌اند و چون باب عنایت رحمن بمفتاح بیان گشوده
شد و طلعت موعود و مکلّم طور خود را لأجل نجات عالم و هدایت امم ظاهر فرمود
کل محجوب بل منصعق بل میّت مشاهده شدند مگر معدودی قلیل که از احرف وجه در
کتاب الهی مذکورند ای دوستان بقول اکتفا ننمائید امروز یوم عمل است قدر
ایّام را بدانید و بآنچه سزاوار است مشغول گردید بگو الوان مختلفهٴ فانیهٴ
این دو یوم شما را از نعماء باقیه و آلاء دائمهٴ الهی منع ننماید عنقریب کل
مشاهده نمایند آنچه را که الیوم از او غافلند آنحین ندامت ثمری ندارد و
تأسّف را اثری نه امروز روز خدمت است جهد نمائید شاید بآنچه لایق است فائز
گردید انّ ربّکم الرّحمن هو المقتدر علی ما یشآء و فی قبضته زمام الأشیآء و هو
المهیمن علی من فی الغیب و الشّهود

\end{otherlanguage}

Yā Ḥezb Allāh, emrūz bar koll dhekr-e Ḥaqq jalla jalāluhu va khedmat-e
amr lāzem. Dar qorūn va aʿṣār qawm montaẓer-e īn yawm-e mobārak būde-and
va chūn bāb-e ʿenāyat-e Raḥmān be-meftāḥ-e bayān goshūde shod va
Ṭalʿat-e Mawʿūd va Mokallem-e Ṭūr khod rā li-ajli najāt-e ʿālam va
hedāyat-e omam ẓāher farmūd, koll maḥjūb bal monṣaʿeq bal mayyet
moshāhade shodand magar maʿdūdī qalīl ke az aḥrof-e vajh dar ketāb-e
elāhī madhkūrand. Ey dūstān, be-qawl ektefā nanamāyīd. Emrūz yawm-e
ʿamal ast. Qadr-e ayyām rā bedānīd va be-ānche sazāvār ast mashghūl
gardīd. Begū alvān-e mokhtalefe-ye fāniye-ye īn do yawm shomā rā az
naʿmā-ye bāqiye va ālā-ye dāʾeme-ye elāhī manʿ nanamāyad. ʿAn-qarīb koll
moshāhade namāyand ānche rā ke al-yawm az ū ghāfeland. Ānḥīn nadāmat
thamarī nadārad va taʾassof rā atharī na. Emrūz rūz-e khedmat ast. Jahd
namāyīd shāyad be-ānche lāyeq ast fāʾez gardīd. Inna rabbakum al-Raḥmān
huwa al-muqtadir ʿalā mā yashāʾ wa-fī qabḍatihi zimām al-ashyāʾ wa-huwa
al-muhaymin ʿalā man fī al-ghayb wa-al-shuhūd.

O people of God, today it is incumbent on everyone to remember
God---glorious is His majesty---and to serve the Cause. Through
centuries and ages the people awaited this blessed Day, yet when the
door of the All-Merciful's loving care was opened with the key of
utterance, and the promised Beauty, who spoke on Sinai, revealed Himself
to save the world and guide the nations, all were seen to be
veiled---indeed, thunderstruck, indeed, dead---except for a small number
who are mentioned among the letters of the Face in the divine Book. O
friends, do not content yourselves with words. Today is the day for
action. Appreciate these days and occupy yourselves with what is
fitting. Say: Do not let the varied, fleeting colors of these two days
keep you from the abiding blessings and enduring favors of God. Soon all
will behold that of which they are heedless today. Then remorse will
bear no fruit, and regret will have no effect. Today is the day of
service. Strive, that you may attain what is worthy. Your Lord, the
All-Merciful, has power over whatever He wills. In His grasp are the
reins of all things, and He is the Guardian over all who are in the
unseen and the visible realms.

\subsection{49}\label{section-3012}

\begin{otherlanguage}{arabic}

و هذا ما نزّل لجناب میرزا مهدی ابن من صعد علیهما بهآء اللّه قوله جلّت عظمته
و عزّ کبریائه

\end{otherlanguage}

Wa-hādhā mā nuzzila li-janāb Mīrzā Mahdī ibn man ṣaʿida ʿalayhimā bahāʾ
Allāh, qawluhu jallat ʿaẓamatuhu wa-ʿazza kibriyāʾuhu.

And this is what was revealed for Mīrzā Mahdī, son of the one who has
ascended, upon them both be the glory of God. His words---glorious is
His grandeur and exalted His majesty:

\subsection{50}\label{section-3013}

\begin{otherlanguage}{arabic}

هو المشرق من افق الملکوت

\end{otherlanguage}

Huwa al-mushriq min ufuq al-malakūt.

He is the One shining from the horizon of the Kingdom.

\subsection{51}\label{section-3014}

\begin{otherlanguage}{persian}

لوح انزله الرّحمن و منه تمرّ عرف البیان علی الامکان تعالی من انزله و تعالی
من اظهره و تعالی من نطق و انطق کلّشیء علی انّه لا اله الّا هو الفرد الواحد
العلیم الحکیم قد ارتفع النّدآء و نطق الکلیل و القوم اکثرهم من الرّاقدین قد
منعتهم النّفس عن الفردوس الأعلی و الهوی عن مالک الوری الا انّهم من
الأخسرین فی کتاب اللّه ربّ العالمین قد اتی من کان مذکوراً بألسن الأصفیآء و
مسطوراً فی کتب اللّه ربّ العرش العظیم من النّاس من سمع و انکر و منهم من سمع
و سرع الی ان فاز بکوثر البیان و شرب منه باسمه الغفور الکریم یا اهل
البهآء ان استمعوا النّدآء من شطر البقعة النّورآء و الأرض البیضآء انّه
یدعوکم بالحقّ و یبشّرکم بعنایة اللّه مالک یوم الدّین ایّاکم ان یمنعکم حبّ
الدّنیا عن مولی الوری او تحزنکم اشارات الغافلین کونوا ثابتین علی الأمر
علی شأن لا تزلّکم شبهات اهل البیان الّذین اعرضوا عن الحقّ اذ اتیهم بسلطان
مبین لعمر اللّه انّهم ما فازوا بالبیان و ما شربوا رحیق العرفان من ایادی
عنایة الرّحمن و انکروا الّذی خلقوا لخدمة امره المبرم المتین یا معشر البیان
تاللّه لا تنفعکم الیوم کتب العالم و لا ملکوت الأسمآء الّا بهذا الکتاب الّذی
ینطق فی اللّیالی و الأیّام و یبشّر الکلّ بما عند اللّه العلیّ العظیم کذلک نطق
قلمی الأعلی و القوم اکثرهم من الغافلین البهآء من لدنّا علی الّذین نبذوا ما
عند النّاس رجآء ما عند اللّه الا انّهم من الرّاسخین الحمد للّه ربّ العالمین

\end{otherlanguage}

Lawḥ anzalahu al-Raḥmān wa-minhu tamurru ʿarf al-bayān ʿalā al-imkān.
Taʿālā man anzalahu wa-taʿālā man aẓharahu wa-taʿālā man naṭaqa
wa-anṭaqa kulla shayʾ ʿalā annahu lā ilāha illā huwa al-Fard al-Wāḥid
al-ʿAlīm al-Ḥakīm. Qad irtafaʿa al-nidāʾ wa-naṭaqa al-kalīl wa-al-qawm
aktharuhum min al-rāqidīn. Qad manaʿathumu al-nafs ʿan al-Firdaws
al-aʿlā wa-al-hawā ʿan Mālik al-warā. Alā innahum min al-akhsarīn fī
kitāb Allāh Rabb al-ʿālamīn. Qad atā man kāna madhkūran bi-alsun
al-aṣfiyāʾ wa-masṭūran fī kutub Allāh Rabb al-ʿarsh al-ʿaẓīm. Min al-nās
man samiʿa wa-ankara wa-minhum man samiʿa wa-saraʿa ilā an fāza
bi-kawthar al-bayān wa-shariba minhu bi-smihi al-Ghafūr al-Karīm. Yā ahl
al-Bahāʾ, an istamiʿū al-nidāʾ min shaṭr al-buqʿa al-nawrāʾ wa-al-arḍ
al-bayḍāʾ. Innahu yadʿūkum bi-al-ḥaqq wa-yubashshirukum bi-ʿināyat Allāh
Mālik yawm al-dīn. Iyyākum an yamnaʿakum ḥubb al-dunyā ʿan Mawlā al-warā
aw tuḥzinakum ishārāt al-ghāfilīn. Kūnū thābitīn ʿalā al-amr ʿalā shaʾn
lā tuzillakum shubuhāt ahl al-Bayān alladhīna aʿraḍū ʿan al-Ḥaqq idh
atāhum bi-sulṭān mubīn. La-ʿamr Allāh innahum mā fāzū bi-al-Bayān wa-mā
sharibū raḥīq al-ʿirfān min ayādī ʿināyat al-Raḥmān wa-ankarū alladhī
khuliqū li-khidmati amrihi al-mubram al-matīn. Yā maʿshar al-Bayān,
ta-Allāh lā tanfaʿukumu al-yawm kutub al-ʿālam wa-lā malakūt al-asmāʾ
illā bi-hādhā al-kitāb alladhī yanṭiqu fī al-layālī wa-al-ayyām
wa-yubashshiru al-kull bi-mā ʿinda Allāh al-ʿAliyy al-ʿAẓīm. Kadhālika
naṭaqa qalamī al-aʿlā wa-al-qawm aktharuhum min al-ghāfilīn. Al-bahāʾ
min ladunnā ʿalā alladhīna nabadhū mā ʿinda al-nās rajāʾa mā ʿinda
Allāh. Alā innahum min al-rāsikhīn. Al-ḥamdu li-llāh Rabb al-ʿālamīn.

A Tablet revealed by the All-Merciful, from which the fragrance of
utterance passes over the world of being. Exalted is He who revealed it,
exalted is He who made it manifest, and exalted is He who has spoken and
caused all things to proclaim that there is no god but Him, the Unique,
the One, the All-Knowing, the All-Wise. The call has risen and the
faltering tongue has spoken, yet most people are asleep. Self has kept
them from the most exalted Paradise, and desire from the Sovereign of
humankind. Indeed, they are among the greatest losers in the Book of
God, the Lord of all worlds. He has come who was mentioned by the
tongues of the chosen ones and recorded in the Books of God, the Lord of
the mighty throne. Some people heard and denied; others heard and
hastened until they attained the Kawthar of utterance and drank from it
in His name, the Ever-Forgiving, the Generous. O people of Bahāʾ, listen
to the call from the direction of the luminous Spot and the white Land.
He truly summons you and gives you glad tidings of the loving care of
God, the Sovereign of the Day of Judgment. Beware lest love of the world
keep you from the Lord of humankind, or the insinuations of the heedless
grieve you. Be so steadfast in the Cause that the doubts of the people
of the Bayān, who turned away from the Truth when He came to them with
manifest authority, cannot cause you to slip. By the life of God, they
have not attained the Bayān or drunk the choice wine of recognition from
the hands of the All-Merciful's loving care, and they have denied Him
for whose irrevocable and firmly established Cause they were created to
serve. O people of the Bayān, by God, neither the books of the world nor
the kingdom of names will benefit you today except through this Book,
which speaks by night and day and gives everyone glad tidings of what is
with God, the Most High, the Most Great. Thus has My Most Exalted Pen
spoken, yet most people are heedless. Glory from Our presence rest upon
those who have cast aside what people possess in the hope of what is
with God. Indeed, they are among the firmly grounded. Praise be to God,
the Lord of all worlds.

\subsection{52}\label{section-3015}

\begin{otherlanguage}{arabic}

و هذا ما نزّل لجناب باقر ابن (ملّا) محمّد قبل علی قوله عزّ شأنه و جلّ بیانه

\end{otherlanguage}

Wa-hādhā mā nuzzila li-janāb Bāqir ibn (Mullā) Muḥammad qabl ʿAlī,
qawluhu ʿazza shaʾnuhu wa-jalla bayānuhu.

And this is what was revealed for Bāqir, son of (Mullā) Muḥammad before
ʿAlī. His words---exalted is His station and glorious His utterance:

\subsection{53}\label{section-3016}

\begin{otherlanguage}{arabic}

هو المشرق من افق العالم باسمه الأعظم

\end{otherlanguage}

Huwa al-mushriq min ufuq al-ʿālam bi-smihi al-aʿẓam.

He is the One shining from the horizon of the world in His Most Great
Name.

\subsection{54}\label{section-3017}

\begin{otherlanguage}{arabic}

ذکر من لدنّا لمن اراد الوجه و انا الذّاکر العلیم قد ذکرنا الّذین آمنوا
باللّه و اریناهم افقی الأعلی و اسمعناهم ندائی الأحلی و هدیناهم الی صراطی
المستقیم یا اهل البهآء تاللّه الحقّ یذکرکم مولی الوری فی مقامه الأعلی الّذی
سمّی بالسّجن مرّة و اخری بسمآء هذه السّمآء الّتی ارتفعت بأمری المهیمن علی من
فی السّموات و الأرضین ان انصروا ربّکم الرّحمن بالأعمال و الأخلاق انّا
امرناکم بها فی الواح نزّلت من سمآء مشیّة ربّکم العالم الخبیر یا باقر قد
ذکرک عبد من عبادی الّذی طاف حولی و فاز بلقائی و شرب رحیق حبّی و اقتصر
الأمور علی خدمة امری ذکرناک فضلاً من عندی و انا الفضّال الکریم ایّاک ان
یمنعک شیء عن اللّه او تحجبک اشارات الغافلین الّذین یدّعون الایمان و جعلوا
البیان شرکاً لأهل الایقان الا انّهم من الملحدین فی کتاب مبین قل خافوا اللّه
و لا تدحضوا الحقّ بما عندکم ستفنی الدّنیا اذاً تجدون انفسکم فی خسران عظیم
کذلک نطق لسان العظمة اذ کان القوم اکثرهم من الرّاقدین الأمواج الظّاهرة من
بحر عنایتی علی الّذین ما نقضوا میثاق اللّه و عهده و طافوا حول امره البدیع

\end{otherlanguage}

Dhikr min ladunnā li-man arāda al-wajh wa-anā al-Dhākir al-ʿAlīm. Qad
dhakarnā alladhīna āmanū bi-Allāh wa-araynāhum ufuqī al-aʿlā
wa-asmaʿnāhum nidāʾī al-aḥlā wa-hadaynāhum ilā ṣirāṭī al-mustaqīm. Yā
ahl al-Bahāʾ, ta-Allāh al-Ḥaqq yadhkurukum Mawlā al-warā fī maqāmihi
al-aʿlā alladhī summiya bi-al-sijn marratan wa-ukhrā bi-samāʾ hādhihi
al-samāʾ allatī irtafaʿat bi-amrī al-muhaymin ʿalā man fī al-samāwāt
wa-al-araḍīn. An inṣurū rabbakum al-Raḥmān bi-al-aʿmāl wa-al-akhlāq.
Innā amarnākum bihā fī alwāḥ nuzzilat min samāʾ mashīʾat rabbikum
al-ʿĀlim al-Khabīr. Yā Bāqir, qad dhakaraka ʿabd min ʿibādī alladhī ṭāfa
ḥawlī wa-fāza bi-liqāʾī wa-shariba raḥīq ḥubbī wa-iqtaṣara al-umūr ʿalā
khidmati amrī. Dhakarnāka faḍlan min ʿindī wa-anā al-Faḍḍāl al-Karīm.
Iyyāka an yamnaʿaka shayʾ ʿan Allāh aw taḥjubaka ishārāt al-ghāfilīn
alladhīna yaddaʿūna al-īmān wa-jaʿalū al-Bayān sharakan li-ahl al-īqān.
Alā innahum min al-mulḥidīn fī kitāb mubīn. Qul khāfū Allāh wa-lā
tudḥiḍū al-Ḥaqq bi-mā ʿindakum. Sa-tafnā al-dunyā idhan tajidūna
anfusakum fī khusrān ʿaẓīm. Kadhālika naṭaqa lisān al-ʿaẓama idh kāna
al-qawm aktharuhum min al-rāqidīn. Al-amwāj al-ẓāhira min baḥr ʿināyatī
ʿalā alladhīna mā naqaḍū mīthāq Allāh wa-ʿahdahu wa-ṭāfū ḥawla amrihi
al-badīʿ.

A remembrance from Our presence for whoever has desired the Face, and I
am the Rememberer, the All-Knowing. We have remembered those who believe
in God, shown them My most exalted horizon, caused them to hear My
sweetest call, and guided them to My straight path. O people of Bahāʾ,
by God, the Truth, the Lord of humankind remembers you in His most
exalted station, which has been called at one time the Prison and at
another the heaven of this heaven that has been raised by My command,
which holds sway over all who are in the heavens and the earths. Aid
your Lord, the All-Merciful, through deeds and character. We have
enjoined these upon you in Tablets sent down from the heaven of the will
of your Lord, the All-Knowing, the All-Informed. O Bāqir, a servant
among My servants has mentioned you---one who has circled around Me,
attained My presence, drunk the choice wine of My love, and confined his
affairs to the service of My Cause. We have remembered you as a grace
from Me, and I am the Most Bountiful, the Generous. Beware lest anything
keep you from God, or the insinuations of the heedless veil you---those
who claim belief and have made the Bayān a snare for the people of
certitude. Indeed, they are among the unbelievers in a clear Book. Say:
Fear God, and do not seek to refute the Truth with what you possess. The
world will pass away; then you will find yourselves in grievous loss.
Thus has the Tongue of Grandeur spoken while most people were asleep.
May the waves rising from the ocean of My loving care be upon those who
have not broken God's covenant and pledge and have circled around His
wondrous Cause.

\subsection{55}\label{section-3018}

\begin{otherlanguage}{arabic}

و هذا ما نزّل لمن سمّی بحسین (عطّار) علیه ٩ ٦٦ {[}بهآء اللّه{]} قوله تبارک و
تعالی

\end{otherlanguage}

Wa-hādhā mā nuzzila li-man summiya bi-Ḥusayn (ʿAṭṭār) ʿalayhi 9 66
{[}bahāʾ Allāh{]}, qawluhu tabāraka wa-taʿālā.

And this is what was revealed for the one named Ḥusayn (ʿAṭṭār), upon
him be 9 66 {[}the glory of God{]}. His words---blessed and exalted is
He:

\subsection{56}\label{section-3019}

\begin{otherlanguage}{arabic}

بنام یکتا خداوند بیمانند

\end{otherlanguage}

Be-nām-e yaktā Khodāvand-e bīmānand.

In the name of the one, incomparable God.

\subsection{57}\label{section-3020}

\begin{otherlanguage}{persian}

یا حسین بمحبّت الهی و شعلهٴ نار مودّت رحمانی بر خدمت امر قیام نما امروز هر
نفسی ارادهٴ نصرت نماید باید از ما عنده بگذرد و به ما عند اللّه ناظر باشد
یا حسین امروز دریای کرم موّاج و اشراقات انوار آفتاب جود کلّ وجود را احاطه
نموده هر نفسی للّه برخاست قعود او را نبیند و توقّف او را اخذ نکند بجنود
بیان بروح و ریحان گمراهانرا هدایت نماید و ضعیفانرا قوّت بخشد از اینکلمهٴ
علیا نفسی تعجّب ننماید ابن مریم علیه سلام اللّه و سلام انبیائه صیّادی را
ملاحظه فرمود که بصید ماهی مشغول فرمود دام را بگذار و بیا تا ترا صیّاد
انام نمایم بعد از اینکلمه کلیل بود نطق یافت جاهل بود ببحر علم درآمد فقیر
بود بملکوت غنا راه یافت ذلیل بود قصد ذروهٴ عزّت نمود گمراه بود بانوار فجر
هدایت فائز گشت چه که از خود گذشت و بحقّ پیوست آمال را بنار حبّ سوخت و
حجبات را باصبع یقین شقّ نمود تا در دنیا بود بیاد روح اللّه ناطق و چون میل
عقبی نمود بمقامی فائز که اقلام عالم از ذکرش عاجز طوبی از برای نفوسیکه
الیوم شبهات معرضین و اشارات منکرین و ضوضاء اهل بیان ایشانرا از مقصود
امکان منع ننمود امروز روز استقامت است و روز خدمت جهد نمائید که شاید فائز
شوید بآنچه که لدی اللّه مذکور و در کتاب مسطور البهآء علیک و علی کلّ ثابت
مستقیم انتهی

\end{otherlanguage}

Yā Ḥusayn, be-maḥabbat-e ilāhī va shoʿle-ye nār-e mavaddat-e raḥmānī bar
khedmat-e amr qiyām namā. Emrūz har nafsī erāde-ye noṣrat namāyad bāyad
az mā ʿendahu bog dharad va be mā ʿinda Allāh nāẓer bāshad. Yā Ḥusayn,
emrūz daryā-ye karam mavvāj va eshrāqāt-e anvār-e āftāb-e jūd koll-e vo
jūd rā eḥāṭe namūde. Har nafsī li-llāh barkhāst qoʿūd ū rā nabīnad va
tavaqqof ū rā akhdh nakonad. Be-jonūd-e bayān be-rūḥ o rayḥān gomrāhān
rā hedāyat namāyad va żaʿīfān rā qovvat bakhshad. Az īn kaleme-ye ʿolyā
nafsī taʿajj ob nanamāyad. Ibn-e Maryam, ʿalayhi salāmu Allāh wa salāmu
anbiyāʾihi, ṣayyādī rā molāḥeẓe farmūd ke be-ṣayd-e māhī mashghūl.
Farmūd dām rā bog dhār va biyā tā to rā ṣayyād-e anām namāyam. Baʿd az
īn kaleme kalīl būd noṭq yāft, jāhel būd be-baḥr-e ʿelm darāmad, faqīr
būd be-malakūt-e ghenā rāh yāft, dhalīl būd qaṣd-e dh erve-ye ʿezzat
namūd, gomrāh būd be-anvār-e fajr-e hedāyat fāʾez gasht, che ke az khod
go dhasht va be-Ḥaqq peyvast. Āmāl rā be-nār-e ḥobb sūkht va ḥojobāt rā
be-eṣbaʿ-e yaqīn shaqq namūd. Tā dar donyā būd be-yād-e Rūḥuʾllāh nāṭeq
va chon meyl-e ʿoqbā namūd be-maqāmī fāʾez ke aqlām-e ʿālam az dhekrash
ʿājez. Ṭūbā az barā-ye nofūsī ke al-yawm shobohāt-e moʿreżīn va
eshārāt-e monkerīn va żawżā-ye ahl-e Bayān īshān rā az maqṣūd-e emkān
manʿ nanamūd. Emrūz rūz-e esteqāmat ast va rūz-e khedmat. Jahd namāyīd
ke shāyad fāʾez shavīd be-ānche ke ladā Allāh madhkūr va dar ketāb
masṭūr. Al-Bahāʾu ʿalayka wa ʿalā kulli thābitin mustaqīm. Intahā.

O Ḥusayn, arise to serve the Cause with divine love and the flame of the
fire of the All-Merciful's affection. Today, whoever wishes to render
assistance must relinquish what he possesses and look to what is with
God. O Ḥusayn, today the ocean of generosity is surging, and the
radiance of the sun of bounty has encompassed all existence. Whoever
arises for God will never be seen sitting idle, nor will hesitation
seize him. With the hosts of utterance, let him guide the wayward with
joy and delight and strengthen the weak. Let no one be astonished at
this exalted word. The Son of Mary, upon him be the peace of God and the
peace of His Prophets, saw a fisherman occupied in catching fish. He
said, ``Put down the net and come, so that I may make you a fisher of
humanity.'' After these words, though tongue-tied, he gained speech;
though ignorant, he entered the ocean of knowledge; though poor, he
found his way to the kingdom of wealth; though abased, he set his sights
on the summit of glory; though astray, he attained the light of the dawn
of guidance. For he renounced himself and joined himself to God. He
burned his desires in the fire of love and tore the veils apart with the
finger of certitude. As long as he was in this world, he spoke in
remembrance of the Spirit of God, and when he turned toward the world
beyond, he attained a station that the pens of the world are powerless
to describe. Blessed are those souls whom neither the doubts of those
who turn away, nor the insinuations of the deniers, nor the clamor of
the people of the Bayān have prevented today from reaching the Desire of
all existence. Today is the day of steadfastness and the day of service.
Strive, so that you may attain what is mentioned in God's presence and
inscribed in the Book. Glory be upon you and upon everyone who is firm
and steadfast. End.

\subsection{58}\label{section-3021}

\begin{otherlanguage}{persian}

این فضل اعظم را نه ببحور میتوان تشبیه نمود و نه بامطار و نه باوراق اشجار
و نه برمول صحرا و نه بذرّات کائنات آیا عباد مقبلین و موحّدین و مخلصین
چگونه از عهدهٴ شکر برآیند نعمتیکه احصاء آن مشکل است شکر آنهم البتّه فوق
عرفان بشر است چون فضلش و رحمتش سبقت یافت اذن عطا فرمود و این سبب بزرگ و
علّة اعظم شد از برای ذاکرین و مخلصین و موحّدین والّا مقام خاک کجا و روح پاک
کجا چه مناسبت چه مشابهت چه مشاکلت معدوم صرفیم و باراده‌اش موجود هر نفسی
را آنمحبوب ذکر نمود بفیوضات نامتناهی فائز گشت مخصوص هر یک بحری از آیات
جاری این خادم فانی از حقّ میطلبد مذکورین را تأیید فرماید بر آنچه امروز
شایسته و پسندیده است

\end{otherlanguage}

Īn fażl-e aʿẓam rā na be-boḥūr mītavān tashbīh namūd va na be-amṭār va
na be-avrāq-e ashjār va na be-romūl-e ṣaḥrā va na be-dharrāt-e kāʾenāt.
Āyā ʿebād-e moqbelīn va movaḥḥedīn va mokhleṣīn chegūne az ʿohde-ye
shokr barāyand? Neʿmatī ke eḥṣā-ye ān moshkel ast shokr-e ān ham albatte
fawq-e ʿerfān-e bashar ast. Chon fażlash va raḥmatash sebqat yāft edhn
ʿaṭā farmūd va īn sabab-e bozorg va ʿellat-e aʿẓam shod az barā-ye
dhākerīn va mokhleṣīn va movaḥḥedīn. Va-ellā maqām-e khāk kojā va rūḥ-e
pāk kojā? Che monāsebat, che moshābahat, che moshākelat? Maʿdūm-e ṣerfīm
va be-erāde-ash mawjūd. Har nafsī rā ān Maḥbūb dhekr namūd be-foyūżāt-e
nāmotanāhī fāʾez gasht. Makhṣūṣ-e har yek baḥrī az āyāt jārī. Īn
khādem-e fānī az Ḥaqq mīṭalabad madhkūrīn rā taʾyīd farmāyad bar ānche
emrūz shāyeste va pasandīde ast.

This supreme grace cannot be likened to oceans, nor to rains, nor to the
leaves of trees, nor to the sands of the desert, nor to the atoms of
creation. How can His servants who turn toward Him, acknowledge His
oneness, and are sincere discharge their duty of gratitude? A bounty
that is difficult to enumerate is certainly also beyond human
understanding to thank Him for. Since His grace and mercy took
precedence, He granted permission, and this became a great means and a
supreme cause for those who remember Him, are sincere, and acknowledge
His oneness. Otherwise, how far apart are the station of dust and the
pure Spirit! What relationship, what resemblance, what likeness is
there? We are sheer nonexistence, brought into being by His will. Every
soul whom that Beloved mentioned attained boundless outpourings of
grace. For each one in particular, an ocean of verses flows. This
transient servant asks God to assist those mentioned in doing what is
fitting and acceptable today.

\subsection{59}\label{section-3022}

\begin{otherlanguage}{persian}

و اینکه در بارهٴ اولیای ن‌ی علیهم بهآء اللّه از ذکور و اناث ذکر نمودند و
همچنین ذکر نفوسیکه بادای حقوق فائز گشتند کل در ساحت امنع اقدس اعلی عرض
شد و هر یک بذکر مقصود عالمیان فائز و باشراقات انوار آفتاب فضل مزیّن قوله
جلّ جلاله و عمّ نواله

\end{otherlanguage}

Va īnke dar bāre-ye awliyā-ye Nūn-Yāʾ, ʿalayhim bahāʾu Allāh, az dhokūr
va enāth dhekr namūdand va hamchenīn dhekr-e nofūsī ke be-adā-ye Ḥoqūq
fāʾez gashtand, koll dar sāḥat-e amnaʿ-e aqdas-e aʿlā ʿarż shod va har
yek be-dhekr-e Maqṣūd-e ʿālamiyān fāʾez va be-eshrāqāt-e anvār-e āftāb-e
fażl mozayyan. Qawluhu jalla jalāluhu wa ʿamma nawāluhu.

What you mentioned concerning the friends in Nūn-Yāʾ, upon them be the
glory of God, both men and women, and likewise your mention of those
souls who attained the privilege of paying the Ḥuqūq, was all presented
at the most exalted, most holy, most inaccessible Court. Each one was
favored with mention by the Desire of all worlds and adorned with the
radiance of the sun of grace. His words---exalted be His glory and
all-embracing His bounty:

\subsection{60}\label{section-3023}

\begin{otherlanguage}{arabic}

هو النّاظر من افقه الأعلی

\end{otherlanguage}

Huwa al-nāẓiru min ufuqihi al-aʿlā.

He is the One who watches from His most exalted horizon.

\subsection{61}\label{section-3024}

\begin{otherlanguage}{persian}

دوستان الهی را تکبیر میرسانیم و ببحر عنایت حقّ جلّ جلاله بشارت میدهیم ای
اهل نون و یا طوبی لکم و نعیماً لکم چه که از اوّل امر فائز شدید بآنچه که من
علی الأرض از آن غافل بودند در سبیل حقّ وارد شد بر شما آنچه که سبب حزن ملأ
اعلی گشت آنچه فی سبیل اللّه وارد شود و آنچه شده کل در کتاب الهی از قلم
اعلی مذکور و مسطور این بشارتیست بزرگ از برای شما هیچ عملی ضایع نشده و
اجرش عنداللّه مخزون زود است مکافات و جزای کل ظاهر و هویدا شود جناب امین
ذکر شما را نموده از قبل و بعد بذکر الهی فائز بوده و هستید وصیّت مینمائیم
شما را بآنچه که سبب اعلاء کلمة اللّه است با کمال اتّحاد و اتّفاق بذکر حقّ
ناطق باشید و بخدمتش مشغول خدمت حقّ جلّ جلاله تبلیغ امر و هدایت ناس بوده و
خواهد بود و این حاصل نشود مگر باخلاق راضیه و اعمال طیّبه و شفقت و محبّت از
حقّ میطلبیم کل را تأیید فرماید بر آنچه رضای او در اوست

\end{otherlanguage}

Dūstān-e ilāhī rā takbīr mīresānīm va be-baḥr-e ʿenāyat-e Ḥaqq, jalla
jalāluhu, beshārat mīdahīm. Ey ahl-e Nūn va Yā, ṭūbā lakum wa naʿīman
lakum, che ke az avval-e amr fāʾez shodīd be-ānche ke man ʿalā al-arḍ az
ān ghāfel būdand. Dar sabīl-e Ḥaqq vāred shod bar shomā ānche ke sabab-e
ḥozn-e malaʾ-e aʿlā gasht. Ānche fī sabīli Allāh vāred shavad va ānche
shode koll dar ketāb-e ilāhī az qalam-e aʿlā madhkūr va masṭūr. Īn
beshāratīst bozorg az barā-ye shomā. Hīch ʿamalī żāyeʿ nashode va ajrash
ʿinda Allāh makhzūn. Zūd ast mokāfāt va jazā-ye koll ẓāher va hoveydā
shavad. Janāb-e Amīn dhekr-e shomā rā namūde. Az qabl va baʿd be-dhekr-e
ilāhī fāʾez būde va hastīd. Vaṣiyyat mīnamāyīm shomā rā be-ānche ke
sabab-e eʿlā-ye kalematuʾllāh ast. Bā kamāl-e etteḥād va ettefāq
be-dhekr-e Ḥaqq nāṭeq bāshīd va be-khedmatash mashghūl. Khedmat-e Ḥaqq,
jalla jalāluhu, tablīgh-e amr va hedāyat-e nās būde va khvāhad būd va īn
ḥāṣel nashavad magar be-akhlāq-e rāżīye va aʿmāl-e ṭayyebe va shafaqat
va maḥabbat. Az Ḥaqq mīṭalabīm koll rā taʾyīd farmāyad bar ānche reżā-ye
ū dar ūst.

We convey greetings of praise to the friends of God and give them glad
tidings of the ocean of the favor of God, exalted be His glory. O people
of Nūn and Yā, blessed are you, and bliss be yours, for from the
beginning of the Cause you attained that of which those on earth were
heedless. In the path of God, you suffered what brought sorrow to the
Concourse on high. Whatever may befall you in the path of God, and
whatever has befallen you, is all mentioned and inscribed in the divine
Book by the Most Exalted Pen. This is great glad tidings for you. No
deed has been lost, and its reward is stored up with God. Soon the
recompense and reward of all will become manifest and evident. His honor
Amīn has mentioned you. Before and since, you have been, and remain,
favored with divine remembrance. We exhort you to that which exalts the
Word of God. In perfect unity and concord, speak in remembrance of God
and occupy yourselves with His service. Service to God, exalted be His
glory, has been and will remain teaching the Cause and guiding the
people. This can be accomplished only through pleasing character, good
deeds, compassion, and love. We ask God to assist everyone in doing what
meets with His good pleasure.

\subsection{62}\label{section-3025}

\begin{otherlanguage}{persian}

چندی قبل مخصوص جناب ملّا محمّد شفیع لوح نازل و ارسال شد و همچنین زیارت
حضرت وحید و اولیای حقّ جلّ جلاله امروز روز فرح و شادیست چه که حقّ ظاهر و
ندایش مرتفع و افق عالم از نورش منوّر و مزیّن از حقّ طلب استقامت نمائید زود
است نعیق ناعق مرتفع شود این خبریست که قلم اعلی از قبل و بعد داده بعضی از
اهل بیان بسرقت آیات الهی مشغول و بعضی بدام خضوع و خشوع مترصّد در هر لباسی
درآیند که شاید اولیا را بدام کشند از حقّ میطلبیم کل را بطراز استقامت مزیّن
دارد بشأنیکه اهل عالم قادر بر منع نباشند خذوا کتاب اللّه بقوّة من عنده انّه
یکفیکم بالحقّ ایّاکم ان یمنعکم شیء عن هذا الصّراط المستقیم و هذا النّبأ
العظیم الحمد للّه المقتدر الباذل الکریم انتهی

\end{otherlanguage}

Chandī qabl makhṣūṣ-e Janāb-e Mullā Muḥammad-Shafīʿ lawḥ nāzel va ersāl
shod va hamchenīn ziyārat-e Ḥażrat-e Vaḥīd va awliyā-ye Ḥaqq, jalla
jalāluhu. Emrūz rūz-e faraḥ va shādīst che ke Ḥaqq ẓāher va nedāyash
mortafeʿ va ofoq-e ʿālam az nūrash monavvar va mozayyan. Az Ḥaqq ṭalab-e
esteqāmat namāyīd. Zūd ast naʿīq-e nāʿeq mortafeʿ shavad. Īn khabarīst
ke qalam-e aʿlā az qabl va baʿd dāde. Baʿżī az ahl-e Bayān be-serqat-e
āyāt-e ilāhī mashghūl va baʿżī be-dām-e khożūʿ va khoshūʿ motaraṣṣed.
Dar har lebāsī darāyand ke shāyad awliyā rā be-dām kashand. Az Ḥaqq
mīṭalabīm koll rā be-ṭerāz-e esteqāmat mozayyan dārad be-shaʾnī ke ahl-e
ʿālam qāder bar manʿ nabāshand. Khudhū kitāba Allāhi bi-quwwatin min
ʿindihi innahu yakfīkum bi-l-ḥaqqi iyyākum an yamnaʿakum shayʾun ʿan
hādhā al-ṣirāṭi al-mustaqīmi wa hādhā al-nabaʾi al-ʿaẓīm. Al-ḥamdu
li-llāhi al-muqtadiri al-bādhili al-karīm. Intahā.

Some time ago, a Tablet was revealed and sent especially for his honor
Mullā Muḥammad-Shafīʿ, as was a Tablet of visitation for his holiness
Vaḥīd and the friends of God, exalted be His glory. Today is a day of
joy and gladness, for God is manifest, His call is raised, and the
horizon of the world is illumined and adorned with His light. Ask God
for steadfastness. Soon the cawing of the crier will rise. This is news
that the Most Exalted Pen has given before and since. Some of the people
of the Bayān are occupied in stealing the divine verses, while others
lie in wait with the snare of submissiveness and humility. They assume
every guise in the hope of ensnaring the friends. We ask God to adorn
everyone with steadfastness in such a manner that the people of the
world will be powerless to hinder them. Take hold of the Book of God
with strength from Him. He will indeed suffice you in truth. Beware lest
anything bar you from this straight path and this momentous
announcement. Praise be to God, the All-Powerful, the Munificent, the
Generous. End.

\subsection{63}\label{section-3026}

\begin{otherlanguage}{arabic}

اهل آن ارض لازال بعنایت حقّ فائز بوده و هستند این ایّام مکتوبی از جناب آقا
میرزا جلال و اخوی ایشان علیهما بهآء اللّه رسید انشآءاللّه این خادم فانی
مؤیّد شود بر جواب بلسان ظاهر و باطن میطلبم حقّ مددش را قطع نفرماید اوست بر
هر شیء مهیمن و محیط

\end{otherlanguage}

Ahl-e ān arż lā yazāl be-ʿenāyat-e Ḥaqq fāʾez būde va hastand. Īn ayyām
maktūbī az Janāb-e Āqā Mīrzā Jalāl va akhavī-ye īshān, ʿalayhimā bahāʾu
Allāh, resīd. In shāʾ Allāh īn khādem-e fānī moʾayyad shavad bar javāb.
Be-lesān-e ẓāher va bāṭen mīṭalabam Ḥaqq madadash rā qaṭʿ nafarmāyad.
Ūst bar har shayʾ mohaymen va moḥīṭ.

The people of that land have always enjoyed, and continue to enjoy,
God's favor. In recent days, a letter arrived from his honor Āqā Mīrzā
Jalāl and his brother, upon them both be the glory of God. God willing,
this transient servant will be assisted in replying. With both my
outward and inward tongue, I ask that God not withhold His assistance.
He is the One who watches over and encompasses all things.

\subsection{64}\label{section-3027}

\begin{otherlanguage}{persian}

و آنچه در ذکر ف‌ت و ابا مرقوم داشتند هر یک از نفوس موقنهٴ ثابته بذکر حقّ
فائز شدند طوبی لهم و نعمیاً لهم و چون اینمراتب در ساحت امنع اقدس اعلی عرض
شد اینکلمهٴ علیا از لسان ارادهٴ مالک وری ظاهر قوله جلّت عظمته

\end{otherlanguage}

Va ānche dar dhekr-e Fāʾ-Tāʾ va Abā marqūm dāshtand, har yek az nofūs-e
mūqene-ye thābete be-dhekr-e Ḥaqq fāʾez shodand. Ṭūbā lahum wa naʿīman
lahum. Va chon īn marāteb dar sāḥat-e amnaʿ-e aqdas-e aʿlā ʿarż shod īn
kaleme-ye ʿolyā az lesān-e erāde-ye Mālek-e varā ẓāher. Qawluhu jallat
ʿaẓamatuhu.

As to what you wrote in mention of Fāʾ-Tāʾ and Abā, each of the assured
and steadfast souls was favored with God's remembrance. Blessed are
they, and bliss be theirs. When these matters were presented at the most
exalted, most holy, most inaccessible Court, this exalted word issued
from the tongue of the will of the Lord of humanity. His words---exalted
be His grandeur:

\subsection{65}\label{section-3028}

\begin{otherlanguage}{persian}

یا اباالحسن نامه‌ات که بعبد حاضر نوشتی از این آیهٴ مبارکهٴ قرآن حکایت
مینماید لا تغادر فیه صغیرة و لا کبیرة اولیاء الهی در هر دیار که ساکنند
ذکرشانرا نمودی و از برای کل اظهار عنایت طلبیدی طوبی لک چه که خیرت احاطه
نموده هر یک فائز شدند بآنچه که سزاوار است انتهی

\end{otherlanguage}

Yā Abā al-Ḥasan, nāme-at ke be-ʿabd-e ḥāżer neveshtī az īn āye-ye
mobārake-ye Qurʾān ḥekāyat mīnamāyad: lā tughādiru fīhi ṣaghīratan wa lā
kabīratan. Awliyā-ye ilāhī dar har diyār ke sākenand dhekrashān rā
namūdī va az barā-ye koll eẓhār-e ʿenāyat ṭalabīdī. Ṭūbā laka, che ke
kheyrat eḥāṭe namūde. Har yek fāʾez shodand be-ānche ke sazāvār ast.
Intahā.

O Abā al-Ḥasan, the letter you wrote to the servant in attendance
reflects this blessed verse of the Qurʾān: ``It leaves out therein
neither small nor great.'' You mentioned the friends of God in every
land where they dwell and requested an expression of favor for them all.
Blessed are you, for your goodness has encompassed them. Each one has
attained what is fitting. End.

\subsection{66}\label{section-3029}

\begin{otherlanguage}{persian}

و اینکه ذکر جناب ملّا مهدی و احبّای خ‌وی علیهم بهآء اللّه را نمودند در ساحت
امنع اقدس عرض شد هذا ما نزّل لهم من سمآء مشیّة ربّنا و ربّکم و ربّ العرش
العظیم قوله جلّ جلاله

\end{otherlanguage}

Va īnke dhekr-e Janāb-e Mullā Mahdī va aḥebbā-ye Khūy, ʿalayhim bahāʾu
Allāh, rā namūdand dar sāḥat-e amnaʿ-e aqdas ʿarż shod. Hādhā mā nuzzila
lahum min samāʾi mashīʾati rabbinā wa rabbikum wa rabbi al-ʿarshi
al-ʿaẓīm. Qawluhu jalla jalāluhu.

Your mention of his honor Mullā Mahdī and the friends of Khūy, upon them
be the glory of God, was presented at the most holy, most inaccessible
Court. This is what was revealed for them from the heaven of the will of
our Lord and your Lord, the Lord of the mighty throne. His
words---exalted be His glory:

\subsection{67}\label{section-3030}

\begin{otherlanguage}{arabic}

هو المشفق الکریم

\end{otherlanguage}

Huwa al-mushfiqu al-karīm.

He is the Compassionate, the Generous.

\subsection{68}\label{section-3031}

\begin{otherlanguage}{arabic}

یا مهدی قرّت عین العالم بظهور مکلّم الطّور فی هذا الظّهور و الأمم فی حجاب
غلیظ قد ظهر الیوم و اتی من کان موعوداً فی کتب اللّه ربّ العالمین انّه اتی
بالحقّ و دعا النّاس الی الأفق الأعلی من النّاس من انکر و اعرض و منهم من سمع
و قال لبّیک یا مقصود العارفین و منهم من منعته سطوة الأمرآء و منهم من
حجّبته اشارات العلمآء الّذین نبذوا الحقّ ورائهم و کانوا من المعتدین فی کتاب
مبین اولئک اتّبعوا اهوائهم و افتوا علی الّذی کانوا ان ینطقوا حین ذکر اسمه
عجّل اللّه فرجه کذلک سوّلت لهم انفسهم الا انّهم من الظّالمین بظلمهم ناح اهل
الفردوس الأعلی و الجنّة العلیا و اخذ الدّخان سکّان الأرض و السّمآء یشهد بذلک
من ینطق انّه لا اله الّا انا الفرد الخبیر

\end{otherlanguage}

Yā Mahdī, qarrat ʿaynu al-ʿālami bi-ẓuhūri Mukallimi al-Ṭūri fī hādhā
al-ẓuhūri wa al-umamu fī ḥijābin ghalīẓ. Qad ẓahara al-yawmu wa atā man
kāna mawʿūdan fī kutubi Allāhi rabbi al-ʿālamīn. Innahu atā bi-l-ḥaqqi
wa daʿā al-nāsa ilā al-ufuqi al-aʿlā. Mina al-nāsi man ankara wa aʿraḍa
wa minhum man samiʿa wa qāla labbayka yā maqṣūda al-ʿārifīn. Wa minhum
man manaʿathu saṭwatu al-umarāʾi wa minhum man ḥajjabathu ishārātu
al-ʿulamāʾi alladhīna nabadhū al-ḥaqqa warāʾahum wa kānū mina
al-muʿtadīna fī kitābin mubīn. Ulāʾika ittabaʿū ahwāʾahum wa aftaw ʿalā
alladhī kānū an yanṭiqū ḥīna dhikri ismihi ʿajjala Allāhu farajahu.
Kadhālika sawwalat lahum anfusuhum alā innahum mina al-ẓālimīn.
Bi-ẓulmihim nāḥa ahlu al-firdawsi al-aʿlā wa al-jannati al-ʿulyā wa
akhadha al-dukhānu sukkāna al-arḍi wa al-samāʾ. Yashhadu bi-dhālika man
yanṭiqu innahu lā ilāha illā anā al-fardu al-khabīr.

O Mahdī, the eye of the world has been gladdened by the appearance, in
this Revelation, of the One who spoke on Sinai, yet the nations remain
behind a thick veil. The Day has appeared, and He who was promised in
the Books of God, the Lord of all worlds, has come. He has indeed come
with the truth and summoned the people to the most exalted horizon.
Among the people are those who denied and turned away, and among them
are those who heard and said, ``Here I am, O Desire of those who know!''
Among them are those whom the might of rulers held back, and those whom
the insinuations of religious scholars veiled---scholars who cast the
Truth behind them and were numbered among the transgressors in a clear
Book. They followed their desires and pronounced judgment against the
One at the mention of whose name they used to say, ``May God hasten His
deliverance.'' Thus have their own selves enticed them. Truly, they are
among the wrongdoers. Through their oppression, the inhabitants of the
most exalted Paradise and the highest Garden have lamented, and smoke
has engulfed the dwellers of earth and heaven. To this bears witness the
One who declares, ``There is no God but Me, the Unique, the All-Aware.''

\subsection{69}\label{section-3032}

\begin{otherlanguage}{persian}

قل یا اهل الخآء و الواو ان استمعوا ندآء مالک الایجاد فی المعاد ثمّ اذکروه
فی اللّیالی و الأیّام انّ ربّکم الرّحمن هو السّامع المجیب ایّاکم ان تمنعکم
اشارات الظّالمین عن اللّه العزیز الحمید دعوا القوم و ما عندهم متوکّلین علی
اللّه المقتدر القدیر سوف یرون المخلصون ما وعدوا به فی کتاب اللّه و
المشرکون مقاماتهم فی الجحیم ان افرحوا بذکر قلمی الأعلی انّه یقرّبکم الی
اللّه مالک یوم الدّین البهآء المشرق من افق سمآء فضلی علیکم و علی الّذین ما
منعهم نعاق النّاعقین الّذین یدّعون الایمان بالبیان و یکفرون بالّذی اتی
بسلطان عظیم انتهی

\end{otherlanguage}

Qul yā ahla al-khāʾi wa al-wāwi an istamiʿū nidāʾa Māliki al-ījādi fī
al-maʿādi thumma udhkurūhu fī al-layālī wa al-ayyām. Inna rabbakumu
al-Raḥmāna huwa al-sāmiʿu al-mujīb. Iyyākum an tamnaʿakum ishārātu
al-ẓālimīna ʿani Allāhi al-ʿazīzi al-ḥamīd. Daʿū al-qawma wa mā ʿindahum
mutawakkilīna ʿalā Allāhi al-muqtadiri al-qadīr. Sawfa yarawna
al-mukhliṣūna mā wuʿidū bihi fī kitābi Allāhi wa al-mushrikūna
maqāmātihim fī al-jaḥīm. An ifraḥū bi-dhikri qalamī al-aʿlā innahu
yuqarribukum ilā Allāhi māliki yawmi al-dīn. Al-bahāʾu al-mushriqu min
ufuqi samāʾi faḍlī ʿalaykum wa ʿalā alladhīna mā manaʿahum nuʿāqu
al-nāʿiqīna alladhīna yaddaʿūna al-īmāna bi-l-Bayāni wa yakfurūna
bi-alladhī atā bi-sulṭānin ʿaẓīm. Intahā.

Say: O people of Khāʾ and Wāw, listen to the call of the Lord of
creation on the Day of Return, then remember Him by night and day. Your
Lord, the All-Merciful, is indeed the All-Hearing, the Responsive.
Beware lest the insinuations of the wrongdoers bar you from God, the
Almighty, the Praiseworthy. Leave the people and what they possess,
placing your trust in God, the All-Powerful, the Omnipotent. Soon the
sincere will behold what they were promised in the Book of God, and
those who associate partners with Him will behold their stations in
hell. Rejoice in the remembrance of My Most Exalted Pen, for it draws
you near to God, the Lord of the Day of Judgment. The glory shining from
the horizon of the heaven of My grace be upon you and upon those whom
the cawing of the criers has not held back---those criers who profess
faith in the Bayān and disbelieve in the One who has come with mighty
authority. End.

\subsection{70}\label{section-3033}

\begin{otherlanguage}{persian}

للّه الحمد انوار ملکوت بیان الهی اهل آن ارض را احاطه فرمود انشآءاللّه مؤیّد
شوند بر اتّحاد و اتّفاق و نصرت امر یعنی ذکر و بیان اینفانی خدمت جناب آقا
ملّا مهدی علیه بهآء اللّه و سایر دوستان آن ارض تکبیر و سلام میرساند و از
برای هر یک توفیق و تأیید میطلبد

\end{otherlanguage}

Li-llāh al-ḥamd, anvār-e malakūt-e bayān-e ilāhī ahl-e ān arż rā eḥāṭe
farmūd. In shāʾ Allāh moʾayyad shavand bar etteḥād va ettefāq va
noṣrat-e amr, yaʿnī dhekr va bayān. Īn fānī khedmat-e Janāb-e Āqā Mullā
Mahdī, ʿalayhi bahāʾu Allāh, va sāyer-e dūstān-e ān arż takbīr va salām
mīresānad va az barā-ye har yek tawfīq va taʾyīd mīṭalabad.

Praise be to God, the lights of the kingdom of divine utterance have
encompassed the people of that land. God willing, they will be assisted
in unity, concord, and the promotion of the Cause---that is, remembrance
and utterance. This transient one conveys greetings of praise and peace
to his honor Āqā Mullā Mahdī, upon him be the glory of God, and the
other friends of that land, and asks for divine enablement and
assistance for each one.

\subsection{71}\label{section-3034}

\begin{otherlanguage}{arabic}

و اینکه در بارهٴ جناب آقا میرزا محمّد (ط‌ب) ابراهیم علیه بهآء اللّه مرقوم
داشتید یک لوح امنع اقدس مخصوص ایشان از سماء مشیّت ربّانی نازل قوله جلّ
جلاله و عزّ کبریائه

\end{otherlanguage}

Va īnke dar bāre-ye Janāb-e Āqā Mīrzā Muḥammad (Ṭāʾ-Bāʾ) Ibrāhīm,
ʿalayhi bahāʾu Allāh, marqūm dāshtīd, yek lawḥ-e amnaʿ-e aqdas makhṣūṣ-e
īshān az samā-ye mashiyyat-e rabbānī nāzel. Qawluhu jalla jalāluhu wa
ʿazza kibriyāʾuhu.

Concerning what you wrote about his honor Āqā Mīrzā Muḥammad (Ṭāʾ-Bāʾ)
Ibrāhīm, upon him be the glory of God, a most holy, most inaccessible
Tablet was revealed especially for him from the heaven of the divine
will. His words---exalted be His glory and mighty His majesty:

\subsection{72}\label{section-3035}

\begin{otherlanguage}{arabic}

بسمه المهیمن علی الأسمآء

\end{otherlanguage}

Bi-smihi al-muhaymini ʿalā al-asmāʾ.

In His name, the One who watches over all names.

\subsection{73}\label{section-3036}

\begin{otherlanguage}{arabic}

ذکر من لدنّا الی الّذی فاز بأنوار الایمان اذ صعد الفرقان و اتی البیان و
بشّر بظهور الرّحمن فی ایّام فیها اخذ الاضطراب سکّان الأرض الّا من شآء اللّه ربّ
الأرباب طوبی لنفس ما زلّتها شبهات القوم و لا اشارات کلّ فاجر مرتاب الّذین
نبذوا الکتاب و اتّخذوا الأوهام لأنفسهم ارباباً من دون اللّه الا انّهم من
اصحاب النّار طوبی للسان اقرّ باللّه و اعترف بما انزله بالحقّ انّه من الأخیار
فی کتاب اللّه مولی الأنام قد حضر کتاب من الأمین و فیه ذکرک ذکرناک بذکر اذ
ظهر خضعت له الأذکار لعمر اللّه لا یعادل بذکری شیء من الأشیآء یشهد بذلک من
عنده امّ الکتاب طوبی لک و للّذین ما نقضوا عهد بارئهم و اقبلوا بوجوه بیضآء
الی افق جعله اللّه مشرق الأنوار خذ کتابی بقوّة من عندی و قل لک الحمد یا من
باسمک ظهرت الأسرار و نزّلت الآیات و برزت البیّنات و ارتفعت الأعلام و نصبت
الرّایات بأن تکتب من قلمک الأعلی لهذا العبد الّذی اقبل الیک یا مولی الوری
ما یکون سراجاً له فی کلّ عالم من عوالمک انّک انت المقتدر علی ما تشآء لا اله
الّا انت القویّ القدیر انتهی

\end{otherlanguage}

Dhikrun min ladunnā ilā alladhī fāza bi-anwār al-īmān idh ṣaʿida
al-Furqān wa atā al-Bayān wa bashshara bi-ẓuhūr al-Raḥmān fī ayyāmin
fīhā akhadha al-iḍṭirābu sukkān al-arḍ illā man shāʾa Allāhu rabbu
al-arbāb. Ṭūbā li-nafsin mā zallat-hā shubuhātu al-qawm wa lā ishārātu
kulli fājirin murtāb, alladhīna nabadhū al-kitāb wa ittakhadhū al-awhāma
li-anfusihim arbāban min dūni Allāh. Alā innahum min aṣḥāb al-nār. Ṭūbā
li-lisānin aqarra bi-Allāh wa iʿtarafa bi-mā anzalahu bi-l-ḥaqq, innahu
min al-akhyār fī kitābi Allāhi mawlā al-anām. Qad ḥaḍara kitābun min
al-Amīn wa fīhi dhikruk. Dhakarnāka bi-dhikrin idh ẓahara khaḍaʿat lahu
al-adhkār. La-ʿamru Allāh lā yuʿādalu bi-dhikrī shayʾun min al-ashyāʾ.
Yashhadu bi-dhālika man ʿindahu ummu al-kitāb. Ṭūbā laka wa li-lladhīna
mā naqaḍū ʿahda bāriʾihim wa aqbalū bi-wujūhin bayḍāʾ ilā ufuqin
jaʿalahu Allāhu mashriqa al-anwār. Khudh kitābī bi-quwwatin min ʿindī wa
qul laka al-ḥamdu yā man bi-ismika ẓaharat al-asrār wa nuzzilat al-āyāt
wa barazat al-bayyināt wa irtafaʿat al-aʿlām wa nuṣibat al-rāyāt bi-an
taktuba min qalamika al-aʿlā li-hādhā al-ʿabdi alladhī aqbala ilayka yā
mawlā al-warā mā yakūnu sirājan lahu fī kulli ʿālamin min ʿawālimik.
Innaka anta al-muqtadiru ʿalā mā tashāʾ, lā ilāha illā anta al-qawiyyu
al-qadīr. Intahā.

A remembrance from Us to the one who attained the lights of faith when
the Qur'an ascended and the Bayān came and announced the appearance of
the All-Merciful, in days when turmoil seized the inhabitants of the
earth, except those whom God, the Lord of lords, willed to spare.
Blessed is the soul whom neither the people's doubts nor the
insinuations of every wicked doubter caused to stumble---those who cast
aside the Book and adopted vain imaginings as lords for themselves
instead of God. Indeed, they are among the people of the fire. Blessed
is the tongue that acknowledged God and confessed the truth of what He
sent down. It is numbered among the righteous in the Book of God, the
Lord of humankind. A letter from Amīn has arrived, and in it you are
mentioned. We have remembered you with a remembrance before which, when
it appeared, all other remembrances bowed. By the life of God, nothing
whatsoever can equal My remembrance. He who possesses the Mother Book
bears witness to this. Blessed are you and those who have not broken
their Creator's covenant and have turned with radiant faces toward a
horizon that God has made the dawning-place of lights. Take hold of My
Book with strength from Me and say, ``Praise be to You, through whose
name mysteries have appeared, verses have been sent down, clear proofs
have come forth, standards have been raised, and banners have been
unfurled, for writing with Your Most Exalted Pen for this servant who
has turned toward You, O Lord of creation, that which will be a lamp for
him in every world of Your worlds. You are indeed powerful over whatever
You will. There is no god but You, the Strong, the Almighty.'' End.

\subsection{74}\label{section-3037}

\begin{otherlanguage}{persian}

للّه الحمد عنایت مقدّس از احصا و ذکر است بر کل لازم در لیالی و ایّام بحمد و
ثنای مولی الأسمآء مشغول باشند و بکمال عجز و ابتهال از مالک مآل استدعای
قبول نمایند چه که انسان مقامات خود و شأن و ثنای خود را ملاحظه مینماید و
همچنین عنایات لاتحصای حقّ را مشاهده میکند متحیّر میشود ولکن از حضرت کریم
مأیوس نبوده و نیستیم از حقّ جلّ جلاله میطلبم جناب میرزای مذکور را مؤیّد
فرماید و بما ینبغی موفّق دارد

\end{otherlanguage}

Li-llāh al-ḥamd, ʿināyat muqaddas az iḥṣā va dhikr ast. Bar kull lāzim
dar layālī va ayyām bi-ḥamd va thanā-yi Mawlā al-Asmāʾ mashghūl bāshand
va bi-kamāl-i ʿajz va ibtihāl az Mālik-i Maʾāl istidʿā-yi qabūl
namāyand, chih kih insān maqāmāt-i khud va shaʾn va thanā-yi khud rā
mulāḥiẓih mīnamāyad va hamchunīn ʿināyāt-i lā-tuḥṣā-yi Ḥaqq rā
mushāhidih mīkunad, mutaḥayyir mīshavad. Va-lākin az Ḥaḍrat-i Karīm
maʾyūs nabūdih va nīstīm. Az Ḥaqq jalla jalāluhu mīṭalabam Jināb-i
Mīrzā-yi madhkūr rā muʾayyad farmāyad va bi-mā yanbaghī muwaffaq dārad.

Praise be to God, His favor is beyond reckoning and description. All
must occupy themselves, by night and day, with praising and extolling
the Lord of Names, and with utter humility and supplication beseech the
Lord of the final outcome for acceptance. For when a person considers
his own stations, standing, and praise, and also beholds the innumerable
favors of God, he is bewildered. Yet we have not despaired, nor do we
despair, of the Most Generous. I ask God, exalted be His glory, to
sustain the aforementioned Mīrzā and enable him to do what is fitting.

\subsection{75}\label{section-3038}

\begin{otherlanguage}{persian}

در این حین نظر بعرض مطلبی قصد مقام اعلی نموده تا آنکه بساحت اقدس فائز
شدم اینکلمات عالیات از لسان مالک اسماء و صفات شنیده شد قوله جلّ و عزّ

\end{otherlanguage}

Dar īn ḥīn naẓar bi-ʿarḍ-i maṭlabī qaṣd-i maqām-i aʿlā namūdih tā ānkih
bi-sāḥat-i aqdas fāʾiz shudam. Īn kalimāt-i ʿāliyāt az lisān-i Mālik-i
asmāʾ va ṣifāt shinīdih shud. Qawluhu jalla wa ʿazza.

At this time, intending to submit a matter, I set out for the exalted
abode and attained the Most Holy Presence. These exalted words were
heard from the tongue of the Lord of names and attributes. His words,
exalted and mighty is He:

\subsection{76}\label{section-3039}

\begin{otherlanguage}{persian}

یا عبد حاضر مدّتهاست عرف محبّت عبدالوهّاب توقّف نموده سبحان‌اللّه او از نفوس
مذکورهٴ لدی العرش بوده و عظمت حقّ و ظهورش را برایات آیات و اعلام بیّنات
مشاهده نموده چه شده آیا سطوت ظالمین حجاب شده یا شبهات علمای ملحدین بگو
یا عبدالوهّاب ربّ وهّاب میفرماید ذرهم فی خوضهم یلعبون اگر صدهزار سنه عمر
بتو عطا شود معادله نمینماید بقدمیکه امروز فی سبیل اللّه برداری و یا
کلمه‌ئی که للّه بآن تکلّم نمائی بشنو ندای مظلوم آفاق را و بکمال همّت بذکر و
ثنای حقّ مشغول باش آنچه ذکر شده و میشود نظر بعنایات حقّ جلّ جلاله است که
نسبت بآنجناب بوده قم و قل لک الحمد بما اسمعتنی ندائک و عرّفتنی جمالک
اسئلک بأن تجعلنی مشتعلاً بنار حبّک علی شأن ینبغی لظهورک و ایّامک

\end{otherlanguage}

Yā ʿabd-i ḥāḍir, muddat-hāst ʿarf-i maḥabbat-i ʿAbdu'l-Vahhāb tavaqquf
namūdih. Subḥān Allāh, ū az nufūs-i madhkūrih-yi ladā al-ʿarsh būdih va
ʿaẓamat-i Ḥaqq va ẓuhūrash rā bi-rāyāt-i āyāt va aʿlām-i bayyināt
mushāhidih namūdih. Chih shudih? Āyā saṭvat-i ẓālimīn ḥijāb shudih yā
shubuhāt-i ʿulamā-yi mulḥidīn? Bigū yā ʿAbdu'l-Vahhāb, Rabb-i Vahhāb
mīfarmāyad: dharhum fī khawḍihim yalʿabūn. Agar ṣad-hizār sanih ʿumr
bi-tu ʿaṭā shavad muʿādilih namīnamāyad bi-qadamī kih imrūz fī sabīl
Allāh bardārī va yā kalimih-ʾī kih li-llāh bi-ān takallum namāʾī.
Bishnaw nidā-yi Maẓlūm-i āfāq rā va bi-kamāl-i himmat bi-dhikr va
thanā-yi Ḥaqq mashghūl bāsh. Ānchih dhikr shudih va mīshavad naẓar
bi-ʿināyāt-i Ḥaqq jalla jalāluhu ast kih nisbat bi-ān jināb būdih. Qum
wa qul laka al-ḥamdu bi-mā asmaʿtanī nidāʾaka wa ʿarraftanī jamālaka.
Asʾaluka bi-an tajʿalanī mushtaʿilan bi-nāri ḥubbika ʿalā shaʾnin
yanbaghī li-ẓuhūrika wa ayyāmik.

O servant in attendance! For some time the fragrance of ʿAbdu'l-Vahhāb's
love has ceased to waft. Glory be to God! He was among the souls
remembered before the Throne and has beheld the greatness of God and His
Revelation through the banners of verses and the standards of clear
proofs. What has happened? Has the tyranny of the oppressors become a
veil, or the doubts of the impious divines? Say: O ʿAbdu'l-Vahhāb! The
All-Bestowing Lord says, ``Leave them to play in their idle
disputation.'' If a hundred thousand years of life were granted to you,
they would not equal a step you take today in the path of God, or a word
you utter for His sake. Hear the call of the Wronged One of the world,
and devote yourself with the utmost resolve to remembering and praising
God. Whatever has been and is being said is in view of the favors that
God, exalted be His glory, has bestowed upon you. Arise and say,
``Praise be to You for causing me to hear Your call and recognize Your
beauty. I ask You to set me ablaze with the fire of Your love in a
manner befitting Your Revelation and Your days.''

\subsection{77}\label{section-3040}

\begin{otherlanguage}{persian}

بیقین بدان احتیاط و حزم بر حیات نیفزاید و استقامت و قیام از حیات نکاهد
للّه بایست و للّه بگو ولکن بحکمت و بیان در کلّ احوال ناظر باش قدر مقام را
بدان و چون جان حفظش نما انّه هو الحافظ النّاصر المعین الحمد للّه ربّ
العالمین انتهی

\end{otherlanguage}

Bi-yaqīn bidān iḥtiyāṭ va ḥazm bar ḥayāt nayafzāyad va istiqāmat va
qiyām az ḥayāt nakāhad. Li-llāh bi-īst va li-llāh bigū, va-lākin
bi-ḥikmat va bayān dar kull-i aḥvāl nāẓir bāsh. Qadr-i maqām rā bidān va
chūn jān ḥifẓash namā. Innahu huwa al-ḥāfiẓu al-nāṣiru al-muʿīn.
Al-ḥamdu li-llāhi rabbi al-ʿālamīn. Intahā.

Know with certainty that caution and prudence do not prolong life, nor
do steadfastness and rising to act shorten it. Stand up for God and
speak for God, but in all circumstances give heed to wisdom and
utterance. Appreciate the worth of this station and guard it as your own
life. He is indeed the Preserver, the Helper, the Sustainer. Praise be
to God, the Lord of the worlds. End.

\subsection{78}\label{section-3041}

\begin{otherlanguage}{persian}

فی‌الحقیقه چندی است خبری از ایشان بساحت اقدس نیامده نه رایحه‌ئی مرور نموده
و نه عرفی متضوّع چون از قبل در سنین متوالیه در سجن مذکور بودند از حقّ جلّ
جلاله سائل که ایشان را مؤیّد فرماید تأییدیکه عرف استقامت از او استشمام
گردد انّه هو الغفور الکریم

\end{otherlanguage}

Fī al-ḥaqīqih chandī ast khabarī az īshān bi-sāḥat-i aqdas nayāmadih,
nah rāyiḥih-ʾī murūr namūdih va nah ʿarfī mutaḍavviʿ. Chūn az qabl dar
sinīn-i mutavāliyih dar sijn madhkūr būdand, az Ḥaqq jalla jalāluhu
sāʾil kih īshān rā muʾayyad farmāyad, taʾyīdī kih ʿarf-i istiqāmat az ū
istishmām gardad. Innahu huwa al-ghafūru al-karīm.

Indeed, for some time no news of him has reached the Most Holy Presence;
no fragrance has passed by, nor has any sweet scent spread abroad. Since
in previous years he was mentioned in the Prison year after year, I
beseech God, exalted be His glory, to sustain him with such assistance
that the fragrance of steadfastness may be perceived from him. He is
indeed the Ever-Forgiving, the All-Generous.

\subsection{79}\label{section-3042}

\begin{otherlanguage}{persian}

حسب الأمر اگر اقتضا نمود آنچه از سماء مشیّت نازل القا شود فرمودند انّا نحبّ
ان نجد منه ما وجدناه من قبل انتهی

\end{otherlanguage}

Ḥasb al-amr agar iqtiḍā namūd ānchih az samāʾ-i mashiyyat nāzil ilqā
shavad. Farmūdand: innā nuḥibbu an najida minhu mā wajadnāhu min qabl.
Intahā.

In accordance with the command, if circumstances warrant, what has
descended from the heaven of the divine will is to be conveyed. He said,
``We wish to find in him what We found before.'' End.

\subsection{80}\label{section-3043}

\begin{otherlanguage}{persian}

ذکر جناب آقا میرزا عبداللّه علیه ٩ ٦٦ {[}بهآء اللّه{]} {[}ی{]} نموده بودند
در ساحت امنع اقدس عرض شد هذا ما نزّل له من سمآء عنایة ربّنا المجیب قوله
بهر بحر بیانه و احاطت انوار شمس فضله

\end{otherlanguage}

Dhikr-i Jināb-i Āqā Mīrzā ʿAbdu'llāh ʿalayhi 9 66 {[}Bahāʾ Allāh{]}
{[}y{]} namūdih būdand. Dar sāḥat-i amnaʿ-i aqdas ʿarḍ shud. Hādhā mā
nuzzila lahu min samāʾi ʿināyati rabbinā al-mujīb. Qawluhu bahara baḥru
bayānihi wa aḥāṭat anwāru shamsi faḍlih.

You had mentioned his honor Āqā Mīrzā ʿAbdu'llāh, upon him be 9 66
{[}the glory of God{]} {[}Y{]}. This was submitted in the Most
Inaccessible, Most Holy Presence. This is what was sent down for him
from the heaven of the favor of our Lord, the Answerer. His
words---dazzling is the ocean of His utterance, and all-encompassing are
the lights of the sun of His grace:

\subsection{81}\label{section-3044}

\begin{otherlanguage}{persian}

هو المقتدر علی من فی الأرض و السّمآء یا عبداللّه ذکرت لدی المظلوم بوده و
عریضه‌ات چندی قبل بساحت اقدس حاضر و باصغا فائز للّه الحمد بافق اعلی ناظری
و بمحبّت حقّ جلّ جلاله مزیّن آنچه مذکور نمودی لدی العرش معلوم از حقّ جلّ جلاله
میطلبیم ترا مؤیّد فرماید بر استقامت قسم بآفتاب افق بیان امروز استقامت از
اعظم صفاتست لدی الوجه طوبی لنفس فازت بها نفوس مذکورهٴ در عریضه هر یک
بذکر حقّ جلّ جلاله فائز و نسئل اللّه بأن یفتح علی وجهک باب البرکة انّه هو
المقتدر علی ما یشآء و فی قبضته زمام من فی السّموات و الأرضین البهآء علیک
من لدنّا و علی الّذین فازوا بهذا الذّکر الحکیم انتهی

\end{otherlanguage}

Huwa al-muqtadiru ʿalā man fī al-arḍi wa al-samāʾ. Yā ʿAbdu'llāh,
dhikrat ladā al-Maẓlūm būdih va ʿarīḍih-at chandī qabl bi-sāḥat-i aqdas
ḥāḍir va bi-iṣghā fāʾiz. Li-llāh al-ḥamd bi-ufuq-i aʿlā nāẓirī va
bi-maḥabbat-i Ḥaqq jalla jalāluhu muzayyan. Ānchih madhkūr namūdī ladā
al-ʿarsh maʿlūm. Az Ḥaqq jalla jalāluhu mīṭalabīm tu-rā muʾayyad
farmāyad bar istiqāmat. Qasam bi-āftāb-i ufuq-i bayān, imrūz istiqāmat
az aʿẓam-i ṣifāt ast ladā al-wajh. Ṭūbā li-nafsin fāzat bihā. Nufūs-i
madhkūrih-yi dar ʿarīḍih har yik bi-dhikr-i Ḥaqq jalla jalāluhu fāʾiz.
Wa nasʾalu Allāha bi-an yaftaḥa ʿalā wajhika bāba al-barakah. Innahu
huwa al-muqtadiru ʿalā mā yashāʾ wa fī qabḍatihi zimāmu man fī
al-samāwāti wa al-araḍīn. Al-bahāʾu ʿalayka min ladunnā wa ʿalā
alladhīna fāzū bi-hādhā al-dhikri al-ḥakīm. Intahā.

He is the All-Powerful over all who are on earth and in heaven. O
ʿAbdu'llāh! You have been remembered before the Wronged One, and your
letter reached the Most Holy Presence some time ago and was heard.
Praise be to God, you look toward the Most Exalted Horizon and are
adorned with the love of God, exalted be His glory. What you mentioned
is known before the Throne. We ask God, exalted be His glory, to sustain
you in steadfastness. By the sun of the horizon of utterance, today
steadfastness is among the greatest attributes in His presence. Blessed
is the soul that has attained it. Each of the souls mentioned in your
letter has been favored with the remembrance of God, exalted be His
glory. We ask God to open the door of blessing before your face. He
indeed has power over whatever He wills, and in His grasp are the reins
of all who are in the heavens and the earths. Glory rest upon you from
Us, and upon those who have attained this wise remembrance. End.

\subsection{82}\label{section-3045}

\begin{otherlanguage}{persian}

ایشان و اخوی ایشان که در شاهرود بودند و همچنین نفوس اخری مکرّر بعنایت حقّ
فائز شدند اینفانی خدمت هر یک تکبیر و سلام میرساند و از برای کل توفیق
میطلبد اینکه در بارهٴ حقوق مرقوم داشتند در ساحت امنع اقدس عرض شد هذا ما
نطق به لسان الفضل فی الجواب قوله جلّ جلاله

\end{otherlanguage}

Īshān va akhavī-yi īshān kih dar Shāhrūd būdand va hamchunīn nufūs-i
ukhrā mukarrar bi-ʿināyat-i Ḥaqq fāʾiz shudand. Īn fānī khidmat-i har
yik takbīr va salām mīrasānad va az barā-yi kull tawfīq mīṭalabad. Īnkih
dar bārih-yi Ḥuqūq marqūm dāshtand dar sāḥat-i amnaʿ-i aqdas ʿarḍ shud.
Hādhā mā naṭaqa bihi lisānu al-faḍli fī al-jawāb. Qawluhu jalla
jalāluhu.

He and his brother, who were in Shāhrūd, and likewise other souls, have
repeatedly been favored with God's grace. This mortal conveys greetings
of ``God is Most Great'' and peace to each of them and prays for divine
assistance for all. What you wrote concerning the Right of God was
submitted in the Most Inaccessible, Most Holy Presence. This is what the
Tongue of Grace spoke in reply. His words, exalted be His glory:

\subsection{83}\label{section-3046}

\begin{otherlanguage}{persian}

یا امین سالها ما ذکر حقوق ننمودیم مع آنکه اینفقره سبب حفظ و علّت برکت
الهی و آثار دیگر بوده چه مقدار از نفوس که خود از {[}مال{]} خود بحضور
فرستادند و چه آوردند اخذ ننمودیم و واگذاشتیم چه که لازم نبود حال آنجناب
بر امورات این ارض مطّلعند و همچنین اسباب لازم است از برای اعلاء کلمه ولکن
نباید نفسی از کسی طلب نماید اگر نفسی فائز شد بادای حقوق الهی اخذ شود
والّا فلا در جمیع احوال باید اولیای حقّ مراقب اعزاز امر باشند و تا حال امر
نشد احدی مطالبه نماید انّه غنیّ عمّا کان و ما یکون انتهی

\end{otherlanguage}

Yā Amīn, sāl-hā mā dhikr-i Ḥuqūq nanamūdīm, maʿa ānkih īn fiqrih sabab-i
ḥifẓ va ʿillat-i barakat-i ilāhī va āthār-i dīgar būdih. Chih miqdār az
nufūs kih khud az {[}māl{]}-i khud bi-ḥuḍūr firistādand va chih āvardand
akhdh nanamūdīm va vāgudhāshtīm, chih kih lāzim nabūd. Ḥāl ān jināb bar
umūrāt-i īn arḍ muṭṭaliʿand va hamchunīn asbāb lāzim ast az barā-yi
iʿlā-yi kalimih, va-lākin nabāyad nafsī az kasī ṭalab namāyad. Agar
nafsī fāʾiz shud bi-adā-yi Ḥuqūq-i ilāhī akhdh shavad, va-illā fa-lā.
Dar jamīʿ-i aḥvāl bāyad awliyā-yi Ḥaqq murāqib-i iʿzāz-i Amr bāshand va
tā ḥāl amr nashud aḥadī muṭālibih namāyad. Innahu ghaniyyun ʿammā kāna
wa mā yakūn. Intahā.

O Amīn! For years We made no mention of the Right of God, although it
has been a means of protection, a cause of divine blessing, and a source
of other effects. How many souls sent offerings from their own
{[}possessions{]} into Our presence or brought them themselves, yet We
did not accept them and left them untouched, for there was no need. Now
you are aware of affairs in this land, and means are also needed for
exalting the Word. Nevertheless, no one should ask anyone for payment.
If a soul attains the privilege of paying the Right of God, it should be
accepted; otherwise, not. In all circumstances the friends of God must
be mindful of upholding the honor of the Cause. Until now, no command
has been given for anyone to demand payment. He is independent of all
that has been and all that will be. End.

\subsection{84}\label{section-3047}

\begin{otherlanguage}{persian}

و اینکه مرقوم داشتند بعضی اخذ مینمایند و صرف میکنند بادّعای اذن ابداً چنین
اذنی صادر نشده از برای احدی باید اولیا مخصوص آنمحبوب در جمیع احوال
باسبابیکه سبب ارتفاع کلمه است ناظر باشند ایّام فانیست باید دوستان بعدل و
انصاف و صدق و صفا عامل باشند و بآنچه سبب اعزاز امر است متمسّک مکرّر ذکر شد
تا آنمحبوب بدانند که چنین امری صادر نشده

\end{otherlanguage}

Va īnkih marqūm dāshtand baʿḍī akhdh mīnamāyand va ṣarf mīkunand
bi-iddiʿā-yi idhn, abadan chunīn idhnī ṣādir nashudih az barā-yi aḥadī.
Bāyad awliyā makhṣūṣ ān maḥbūb dar jamīʿ-i aḥvāl bi-asbābī kih sabab-i
irtifāʿ-i kalimih ast nāẓir bāshand. Ayyām fānī-st. Bāyad dūstān bi-ʿadl
va inṣāf va ṣidq va ṣafā ʿāmil bāshand va bi-ānchih sabab-i iʿzāz-i Amr
ast mutamassik. Mukarrar dhikr shud tā ān maḥbūb bidānand kih chunīn
amrī ṣādir nashudih.

As for what you wrote, that some receive and spend it claiming
permission, no such permission has ever been issued to anyone. The
friends, and especially you, dear friend, must in all circumstances look
to the means that exalt the Word. These days are fleeting. The friends
must act with justice, fairness, truthfulness, and purity, and hold fast
to whatever upholds the honor of the Cause. This has been stated
repeatedly so that you, dear friend, may know that no such command has
been issued.

\subsection{85}\label{section-3048}

\begin{otherlanguage}{persian}

حینیکه محبوب فؤاد جناب حاجی میرزا ح‌ی ع‌ل و جناب آقا میرزا حسین علیهما
بهآء اللّه و عنایته بحضور فائز بودند لسان عظمت ذکر شما را فرمود للّه الحمد
بعنایت فائزید و بکلمهٴ رضا مزیّن هنیئاً لجنابکم و مریئاً لجنابکم این حین که
بتحریر این ورقه مشغول بودم بحضور فائز مقصود عالمیان باینکلمهٴ مبارکه
ناطق قوله جلّ جلاله و احاط فضله یا عبد حاضر انّا نحبّ ان نذکر الّذین شربوا
رحیق الاستقامة هناک فضلاً من عندنا و انا الذّاکر المشفق الکریم نذکرهم و
نستر اسمائهم حکمة من لدنّا و انا الآمر الحکیم انّا ذکرنا الصّاد بذکر تضوّع
منه عرف عنایتی فی ملکوتی الممتنع المنیع و نذکره فی هذا الحین بذکر جمیل
البهآء اللّائح من افق عنایتی علی الّذین ما ذکرنا اسمائهم و علی الّذین شربوا
رحیقی المختوم انتهی

\end{otherlanguage}

Ḥīnī kih maḥbūb-i fuʾād Jināb-i Ḥājī Mīrzā Ḥ-Y ʿ-L va Jināb-i Āqā Mīrzā
Ḥusayn ʿalayhimā bahāʾu Allāhi wa ʿināyatuhu bi-ḥuḍūr fāʾiz būdand
lisān-i ʿaẓamat dhikr-i shumā rā farmūd. Li-llāh al-ḥamd bi-ʿināyat
fāʾizīd va bi-kalimih-yi riḍā muzayyan. Hanīʾan li-jinābikum wa marīʾan
li-jinābikum. Īn ḥīn kih bi-taḥrīr-i īn varaqih mashghūl būdam bi-ḥuḍūr
fāʾiz, Maqṣūd-i ʿālamiyān bi-īn kalimih-yi mubārakih nāṭiq. Qawluhu
jalla jalāluhu wa aḥāṭa faḍluhu: yā ʿabd-i ḥāḍir, innā nuḥibbu an
nadhkura alladhīna sharibū raḥīqa al-istiqāmati hunāka faḍlan min
ʿindinā wa anā al-dhākiru al-mushfiqu al-karīm. Nadhkuruhum wa nasturu
asmāʾahum ḥikmatan min ladunnā wa anā al-āmiru al-ḥakīm. Innā dhakarnā
al-Ṣāda bi-dhikrin taḍawwaʿa minhu ʿarfu ʿināyatī fī malakūtī
al-mumtaniʿi al-manīʿ, wa nadhkuruhu fī hādhā al-ḥīni bi-dhikrin jamīl.
Al-bahāʾu al-lāʾiḥu min ufuqi ʿināyatī ʿalā alladhīna mā dhakarnā
asmāʾahum wa ʿalā alladhīna sharibū raḥīqī al-makhtūm. Intahā.

When the beloved of my heart, his honor Ḥājī Mīrzā Ḥ-Y ʿ-L, and his
honor Āqā Mīrzā Ḥusayn, upon them both be the glory of God and His
favor, were in His presence, the Tongue of Grandeur mentioned you.
Praise be to God, you have attained His favor and are adorned with the
word of His good pleasure. May it bring you joy and delight! Just now,
while I was engaged in writing this page, I attained His presence, and
the Desire of all the worlds spoke these blessed words. His words,
exalted be His glory and all-encompassing His grace: ``O servant in
attendance! We wish to remember those who have drunk the choice wine of
steadfastness in that place, as a grace from Us. I am the One who
remembers, the Compassionate, the All-Generous. We remember them and
conceal their names out of wisdom on Our part. I am the Ordainer, the
All-Wise. We have remembered Ṣād with a remembrance from which the
fragrance of My favor has spread throughout My inaccessible, impregnable
Kingdom, and We remember him at this time with a beautiful remembrance.
May the glory shining from the horizon of My favor rest upon those whose
names We have not mentioned, and upon those who have drunk My sealed
wine.'' End.

\subsection{86}\label{section-3049}

\begin{otherlanguage}{persian}

هر ذرّه‌ئی از ذرّات کائنات شهادت بر فضل و رحمت و شفقت و عنایت حقّ جلّ جلاله
داده و میدهد از او میطلبیم اینفانی و اولیای خود را برضایش مزیّن و بر
استقامت بر امرش موفّق دارد اینست مقامیکه اولیا و اصفیاء قبل و بعد از حقّ
جلّ جلاله مسئلت نموده‌اند اینفانی بکمال عجز میطلبد از برای کل آنچه را که
از برای خود میطلبد و از برای دوستان مسئلت مینماید آنچه که خود را قابل آن
نمیداند امواج بحر رحمت الهی بشأنی ظاهر که گویا شجرهٴ یأس مقطوع و نهالهای
رجا از جمیع جهات مرتفع تعالی تعالی کرمه و فضله و سلطانه و عزّه و ذکره
التّکبیر و البهآء و الثّنآء علی جنابکم و علی الّذین سمعوا و اجابوا ربّهم
النّاطق العلیم الحکیم

\end{otherlanguage}

Har dharrih-ʾī az dharrāt-i kāʾināt shahādat bar faḍl va raḥmat va
shafaqat va ʿināyat-i Ḥaqq jalla jalāluhu dādih va mīdahad. Az ū
mīṭalabīm īn fānī va awliyā-yi khud rā bi-riḍāyash muzayyan va bar
istiqāmat bar Amrash muwaffaq dārad. Īn ast maqāmī kih awliyā va
aṣfiyā-yi qabl va baʿd az Ḥaqq jalla jalāluhu masʾalat namūdih-and. Īn
fānī bi-kamāl-i ʿajz mīṭalabad az barā-yi kull ānchih rā kih az barā-yi
khud mīṭalabad va az barā-yi dūstān masʾalat mīnamāyad ānchih kih khud
rā qābil-i ān namīdānad. Amvāj-i baḥr-i raḥmat-i ilāhī bi-shaʾnī ẓāhir
kih gūyā shajarih-yi yaʾs maqṭūʿ va nahāl-hā-yi rajā az jamīʿ-i jihāt
murtafiʿ. Taʿālā taʿālā karamuhu wa faḍluhu wa sulṭānuhu wa ʿizzuhu wa
dhikruhu. Al-takbīru wa al-bahāʾu wa al-thanāʾu ʿalā jinābikum wa ʿalā
alladhīna samiʿū wa ajābū rabbahum al-nāṭiqa al-ʿalīma al-ḥakīm.

Every atom of creation has borne and continues to bear witness to the
grace, mercy, compassion, and favor of God, exalted be His glory. We ask
Him to adorn this mortal and His friends with His good pleasure and
enable them to remain steadfast in His Cause. This is the station that
His friends and chosen ones, both former and latter, have sought from
God, exalted be His glory. This mortal, in utter humility, asks for
everyone what he asks for himself, and beseeches for the friends what he
does not consider himself worthy of receiving. The waves of the ocean of
divine mercy are so manifest that it is as though the tree of despair
has been cut down and the saplings of hope have risen on every side.
Exalted, exalted are His generosity, His grace, His sovereignty, His
might, and His remembrance. Greetings of ``God is Most Great,'' glory,
and praise rest upon you and upon those who have heard and answered
their Lord, the Speaker, the All-Knowing, the All-Wise.

\subsection{87}\label{section-3050}

\begin{otherlanguage}{arabic}

بسم ربّنا الأقدس الأعظم العلیّ الأبهی

\end{otherlanguage}

Bi-smi rabbinā al-aqdasi al-aʿẓami al-ʿaliyyi al-abhā.

In the name of our Lord, the Most Holy, the Most Great, the Most High,
the All-Glorious.

\subsection{88}\label{section-3051}

\begin{otherlanguage}{arabic}

شهد اللّه انّه لا اله الّا هو له العزّة و الاقتدار و له القوّة و الاختیار
یرفع من یشآء و یضع الآخرین و هو المقتدر العلیم الحکیم

\end{otherlanguage}

Shahida Allāhu annahu lā ilāha illā huwa, lahu al-ʿizzatu wa al-iqtidār
wa lahu al-quwwatu wa al-ikhtiyār. Yarfaʿu man yashāʾu wa yaḍaʿu
al-ākharīn wa huwa al-muqtadiru al-ʿalīmu al-ḥakīm.

God bears witness that there is no god but Him. His are might and
dominion, and His are strength and the power to choose. He exalts
whomever He wills and abases others. He is the All-Powerful, the
All-Knowing, the All-Wise.

\subsection{89}\label{section-3052}

\begin{otherlanguage}{persian}

و بعد نامهٴ دیگر آنمحبوب که از قبل ارسال داشته بودند نظر بکثرت اشغال و
امور جواب تأخیر افتاد از حقّ جلّ جلاله میطلبم خادم را مؤیّد فرماید بر تدارک
آنچه از او فوت شده نامه رسید و باب وصل گشود و نفحهٴ قرب مرور نمود للّه
الحمد امطار ذکر و بیان از قلم و لسان اولیائش جاریست فی‌الحقیقه ذکر و بیان
مفتاحند از برای باب قرب و وصال

\end{otherlanguage}

Va baʿd, nāme-ye dīgar-e ān maḥbūb ke az qabl ersāl dāshte būdand, naẓar
be-kethrat-e ashghāl va omūr javāb taʾkhīr oftād. Az Ḥaqq jalla jalāluhu
mī-ṭalabam khādem rā moʾayyad farmāyad bar tadārok-e ānche az ū fawt
shode. Nāme rasīd va bāb-e vaṣl goshūd va nafḥe-ye qorb morūr namūd.
Li-llāhi al-ḥamd, amṭār-e dhekr va bayān az qalam va lesān-e awliyāyash
jārīst. Fī al-ḥaqīqe dhekr va bayān meftāḥand az barā-ye bāb-e qorb va
veṣāl.

Furthermore, the reply to the other letter that beloved friend had sent
earlier was delayed because of the press of occupations and affairs. I
ask God, exalted be His glory, to assist this servant in making up for
what he has missed. The letter arrived, opened the door of reunion, and
brought a passing breeze of nearness. Praise be to God, showers of
remembrance and utterance flow from the pens and tongues of His loved
ones. Indeed, remembrance and utterance are keys to the door of nearness
and reunion.

\subsection{90}\label{section-3053}

\begin{otherlanguage}{persian}

ذکر جناب آقا محمّد رضا (محمّدآباد) علیه بهآء اللّه را نموده بودند و همچنین
مراتب شوق و اشتیاق و اشتعال ایشانرا در ساحت امنع اقدس اعلی عرض شد و این
آیات باهرات مخصوص ایشان از سماء عنایت رحمن نازل قوله عزّ بیانه و قدس
مقامه

\end{otherlanguage}

Dhekr-e Janāb-e Āqā Moḥammad Reḍā (Moḥammadābād), ʿalayhi bahāʾu Allāh,
rā namūde būdand va hamchenīn marāteb-e shawq va eshtiyāq va eshteʿāl-e
īshān rā. Dar sāḥat-e amnaʿ-e aqdas-e aʿlā ʿarḍ shod va īn āyāt-e
bāherāt makhṣūṣ-e īshān az samā-ye ʿenāyat-e Raḥmān nāzel. Qawluhu ʿazza
bayānuhu wa quddisa maqāmuhu:

You had mentioned Áqá Muḥammad Riḍá (Muḥammadábád), upon him be the
glory of God, as well as the degree of his longing, ardor, and fervor.
This was presented in the most exalted, most holy, most inaccessible
court, and these resplendent verses were revealed especially for him
from the heaven of the All-Merciful's favor. His words, exalted be His
utterance and sanctified His station:

\subsection{91}\label{section-3054}

\begin{otherlanguage}{arabic}

بسمه المهیمن علی من فی الأرض و السّمآء

\end{otherlanguage}

Bismihi al-Muhaymin ʿalā man fī al-arḍ wa al-samāʾ.

In His name, the Guardian over all who are on earth and in heaven.

\subsection{92}\label{section-3055}

\begin{otherlanguage}{arabic}

یا محمّد قبل رضا قل لک الحمد یا الهی بما فتحت علی وجوه اولیائک باب اللّقآء
و لک الثّنآء یا مولی الأسمآء بما فتحت بمقالید البیان ابواب العرفان و لک
البهآء یا فاطر السّمآء بما نوّرت العالم بأنوار اسمک الأعظم و لک الجلال یا
مطلع الجمال بما هدیت احبّائک الی سوآء الصّراط و لک الفضل یا مشرق العدل بما
دعوت العباد الی افقک الأعلی فی المبدء و المعاد و لک النّور یا مکلّم الطّور
بما عرّفت اهل المعانی افق ظهورک و لک الظّهور یا معطی الأرواح بما جعلتنا
منجذبین بآیاتک و مقبلین الی افقک و مشتعلین بنار حبّک و طائرین فی هوآء
قربک و ناطقین بثنائک الجمیل

\end{otherlanguage}

Yā Muḥammad qabla Riḍā, qul: Laka al-ḥamd yā Ilāhī bi-mā fataḥta ʿalā
wujūhi awliyāʾika bāba al-liqāʾ, wa laka al-thanāʾ yā Mawlā al-asmāʾ
bi-mā fataḥta bi-maqālīdi al-bayān abwāba al-ʿirfān, wa laka al-bahāʾ yā
Fāṭira al-samāʾ bi-mā nawwarta al-ʿālama bi-anwāri ismika al-Aʿẓam, wa
laka al-jalāl yā Maṭliʿa al-jamāl bi-mā hadayta aḥibbāʾaka ilā sawāʾi
al-ṣirāṭ, wa laka al-faḍl yā Mashriqa al-ʿadl bi-mā daʿawta al-ʿibāda
ilā ufuqika al-aʿlā fī al-mabdaʾ wa al-maʿād, wa laka al-nūr yā
Mukallima al-Ṭūr bi-mā ʿarrafta ahla al-maʿānī ufuqa ẓuhūrika, wa laka
al-ẓuhūr yā Muʿṭiya al-arwāḥ bi-mā jaʿaltanā munjadhibīna bi-āyātika wa
muqbilīna ilā ufuqika wa mushtaʿilīna bi-nāri ḥubbika wa ṭāʾirīna fī
hawāʾi qurbika wa nāṭiqīna bi-thanāʾika al-jamīl.

O Muḥammad before Riḍá, say: Praise be to You, my God, for opening
before the faces of Your loved ones the door of meeting You. Adoration
be to You, Lord of Names, for opening the doors of recognition with the
keys of utterance. Glory be to You, Creator of heaven, for illuminating
the world with the lights of Your Most Great Name. Majesty be to You,
Dawning-Place of Beauty, for guiding Your loved ones to the straight
path. Bounty be to You, Dayspring of Justice, for summoning Your
servants to Your most exalted horizon in the beginning and the return.
Light be to You, the One who spoke on Sinai, for making known to the
people of inner meaning the horizon of Your revelation. Manifestation be
to You, Giver of spirits, for making us enraptured by Your verses,
turned toward Your horizon, aflame with the fire of Your love, soaring
in the air of Your nearness, and voicing Your beautiful praise.

\subsection{93}\label{section-3056}

\begin{otherlanguage}{persian}

قل یا الهی و سیّدی و معبودی و مقصودی تسمع ما تنطق به السن الکائنات و تری
ما تعمل به الممکنات اسئلک بلآلئ عمّان علمک و برحمتک الّتی سبقت خلقک و
عبادک و بسرّک المکنون و اسمک المخزون الّذی اظهر نفسه فی ناسوت الانشآء و
باستوائه علی العرش بین الوری بأن تنزل من سمآء جودک امطار رحمتک ثمّ احفظ
اولیائک من شرّ اعدائک اسئلک یا مولی العالم و مقصود الأمم باسمک الأعظم بأن
تنصر احبّائک بقوّتک و قدرتک و جنود غیبک کما نصرتهم من قبل انّک انت المقتدر
علی ما تشآء و فی قبضتک ازمّة من فی الأرض و السّمآء

\end{otherlanguage}

Qul: Yā Ilāhī wa Sayyidī wa Maʿbūdī wa Maqṣūdī, tasmaʿu mā tanṭiqu bihi
alsunu al-kāʾināt wa tarā mā taʿmalu bihi al-mumkināt. Asʾaluka
bi-laʾāliʾi ʿummāni ʿilmika wa bi-raḥmatika allatī sabaqat khalqaka wa
ʿibādaka wa bi-sirrika al-maknūn wa ismika al-makhzūn alladhī aẓhara
nafsahu fī nāsūti al-inshāʾ wa bi-istiwāʾihi ʿalā al-ʿarsh bayna al-warā
bi-an tunzila min samāʾi jūdika amṭāra raḥmatika thumma iḥfaẓ awliyāʾaka
min sharri aʿdāʾika. Asʾaluka yā Mawlā al-ʿālam wa Maqṣūda al-umam
bi-ismika al-Aʿẓam bi-an tanṣura aḥibbāʾaka bi-quwwatika wa qudratika wa
junūdi ghaybika ka-mā naṣartahum min qabl. Innaka anta al-Muqtadir ʿalā
mā tashāʾ wa fī qabḍatika azimmatu man fī al-arḍ wa al-samāʾ.

Say: My God, my Master, the One I worship and seek! You hear what the
tongues of all beings utter and see what all created things do. I ask
You, by the pearls of the ocean of Your knowledge, by Your mercy that
preceded Your creation and Your servants, by Your hidden mystery and
Your treasured Name, who manifested Himself in the human realm of
creation, and by His establishment upon the throne among humanity, to
send down the rains of Your mercy from the heaven of Your generosity.
Then protect Your loved ones from the malice of Your enemies. I ask You,
Lord of the world and Desire of the nations, by Your Most Great Name, to
aid Your loved ones with Your strength, Your power, and Your unseen
hosts, as You aided them before. You indeed have power over whatever You
will, and in Your grasp are the reins of all who are on earth and in
heaven.

\subsection{94}\label{section-3057}

\begin{otherlanguage}{arabic}

ای ربّ تری عبدک هذا شهد بوحدانیّتک و فردانیّتک و بما جری من قلمک الأعلی فی
ملکوت الانشآء و یشهد بما نطق به لسانک یا مالک الأسمآء و فاطر السّمآء ای
ربّ ترانی متمسّکاً بحبلک و متشبّثاً بذیل عنایتک و راجیاً ما عندک اسئلک ان لا
تخیّبنی بجودک و فضلک و الطافک لا اله الّا انت القویّ القادر السّامع المجیب
انتهی

\end{otherlanguage}

Ay Rabb, tarā ʿabdaka hādhā shahida bi-waḥdāniyyatika wa fardāniyyatika
wa bi-mā jarā min qalamika al-aʿlā fī malakūti al-inshāʾ, wa yashhadu
bi-mā naṭaqa bihi lisānuka yā Mālika al-asmāʾ wa Fāṭira al-samāʾ. Ay
Rabb, tarānī mutamassikan bi-ḥablika wa mutashabbithan bi-dhayli
ʿināyatika wa rājiyan mā ʿindaka. Asʾaluka an lā tukhayyibanī bi-jūdika
wa faḍlika wa alṭāfika. Lā ilāha illā anta al-Qawiyy al-Qādir al-Sāmiʿ
al-Mujīb. Intahā.

O Lord, You see this servant of Yours bearing witness to Your oneness
and uniqueness and to what has flowed from Your Most Exalted Pen in the
kingdom of creation. He also bears witness to what Your tongue has
spoken, Lord of Names and Creator of heaven. O Lord, You see me holding
fast to Your cord, clinging to the hem of Your loving care, and hoping
for what is with You. I ask You, by Your generosity, Your bounty, and
Your gracious favors, not to disappoint me. There is no god but You, the
Strong, the Powerful, the Hearer, the Answerer. End of quotation.

\subsection{95}\label{section-3058}

\begin{otherlanguage}{persian}

این آیات که هر حرفی از آن باب رحمتی است از برای خلق مخصوص جناب مذکور
نازل انشآءاللّه بآن فائز شوند عرفش عالم بیان را تجدید نماید و اشجار وجود
را تازه کند از قبل هم ایشان بذکر مقصودنا و مقصودکم و مقصود من فی الأرض و
السّمآء فائز شده‌اند از حقّ جلّ جلاله سائل و آمل که ایشانرا در جمیع احوال
مؤیّد فرماید انّه علی کلّ شیء قدیر

\end{otherlanguage}

Īn āyāt ke har ḥarfī az ān bāb-e raḥmatī ast az barā-ye khalq, makhṣūṣ-e
Janāb-e madhkūr nāzel. In shāʾ Allāh be-ān fāʾez shavand. ʿArfash
ʿālam-e bayān rā tajdīd namāyad va ashjār-e vojūd rā tāze konad. Az qabl
ham īshān be-dhekr-e Maqṣūdinā wa Maqṣūdikum wa Maqṣūdi man fī al-arḍ wa
al-samāʾ fāʾez shode-and. Az Ḥaqq jalla jalāluhu sāʾel va āmel ke īshān
rā dar jamīʿ-e aḥvāl moʾayyad farmāyad. Innahu ʿalā kulli shayʾin qadīr.

These verses, every letter of which is a door of mercy for creation,
were revealed especially for the gentleman mentioned. God willing, he
will receive them. Their fragrance renews the world of utterance and
refreshes the trees of existence. Previously, too, he was favored with
mention by the One who is our Desire, your Desire, and the Desire of all
who are on earth and in heaven. I ask and hope that God, exalted be His
glory, will assist him in every circumstance. He indeed has power over
all things.

\subsection{96}\label{section-3059}

\begin{otherlanguage}{persian}

و همچنین ذکر جناب حاجی محمّد طاهر علیه بهآء اللّه را فرمودند عرض فانی آنکه
هر نفسی الیوم اقلّ من آن بوجه مقصود عالمیان توجّه نمود ذکرش از لسان عظمت
ظاهر فضل بمقام و مرتبه‌ایست که عالم از احصائش عاجز و از ذکرش قاصر صدهزار
طوبی از برای نفسیکه بذکر الهی در یومش فائز شد اینست آن کنزیکه در ظلّ عرش
مکنون و محفوظ است از برای صاحبش اهل ارض غافلند قسم بمحبوب عالم اگر شاعر
شوند از برای یکحرف که از سماء مشیّت مقصود نازل شود جان نثار نمایند هذا حقّ
لا ریب فیه

\end{otherlanguage}

Va hamchenīn dhekr-e Janāb-e Ḥājī Moḥammad Ṭāher, ʿalayhi bahāʾu Allāh,
rā farmūdand. ʿArḍ-e fānī ānke har nafsī al-yawm aqalla min ān be-vajh-e
Maqṣūd-e ʿālamiyān tavajjoh namūd, dhekrash az lesān-e ʿaẓamat ẓāher.
Faḍl be-maqām va martabe-īst ke ʿālam az eḥṣāyash ʿājez va az dhekrash
qāṣer. Ṣad hezār ṭūbā az barā-ye nafsī ke be-dhekr-e elāhī dar yawmash
fāʾez shod. Īnast ān kanzī ke dar ẓell-e ʿarsh maknūn va maḥfūẓ ast az
barā-ye ṣāḥebash. Ahl-e arḍ ghāfeland. Qasam be-Maḥbūb-e ʿālam, agar
shāʿer shavand az barā-ye yek ḥarf ke az samā-ye mashiyyat-e Maqṣūd
nāzel shavad jān nethār namāyand. Hādhā ḥaqqun lā rayba fīh.

You also mentioned Ḥájí Muḥammad Ṭáhir, upon him be the glory of God.
This humble one submits that whoever has turned today, for less than an
instant, toward the face of the Desire of all the worlds has been
mentioned by the Tongue of Grandeur. Such is the measure and degree of
His bounty that the world is powerless to reckon it and unable to
describe it. A hundred thousand blessings upon the soul who has been
favored with God's remembrance in His Day! This is the treasure
concealed and preserved beneath the shadow of the Throne for its owner.
The people of the earth are heedless. I swear by the Beloved of the
world, if they were to perceive it, they would offer their lives for a
single letter sent down from the heaven of the will of the Desired One.
This is a truth beyond all doubt.

\subsection{97}\label{section-3060}

\begin{otherlanguage}{persian}

اینفقره هم وقتی از اوقات تلقاء وجه عرض شد و باصغا فائز گشت هذا ما نزّل له
من ملکوت ربّنا الرّحمن قوله تبارک و تعالی

\end{otherlanguage}

Īn feqre ham vaqtī az awqāt telqā-ye vajh ʿarḍ shod va be-eṣghā fāʾez
gasht. Hādhā mā nuzzila lahu min malakūti Rabbinā al-Raḥmān. Qawluhu
tabāraka wa taʿālā:

This matter, too, was presented in His presence on one occasion and was
graciously heard. This is what was sent down for him from the kingdom of
our Lord, the All-Merciful. His words, blessed and exalted is He:

\subsection{98}\label{section-3061}

\begin{otherlanguage}{arabic}

بسمی السّمیع المجیب

\end{otherlanguage}

Bismī al-Samīʿ al-Mujīb.

In My name, the All-Hearing, the Answerer.

\subsection{99}\label{section-3062}

\begin{otherlanguage}{arabic}

یا محمّد قبل طاهر یذکرک فی هذا اللّیل من طاف حوله الملأ الأعلی و اهل
الفردوس و الملائکة المقرّبون ان اشکر اللّه ثمّ اذکره فی اللّیالی و الأیّام
انّه یسمع و یری و هو الحقّ علّام الغیوب نوصیک و الّذین آمنوا بالأمانة و
الوفآء و بما یرتفع به کلمة اللّه المهیمن القیّوم

\end{otherlanguage}

Yā Muḥammad qabla Ṭāhir, yadhkuruka fī hādhā al-layl man ṭāfa ḥawlahu
al-malaʾu al-aʿlā wa ahlu al-Firdaws wa al-malāʾikatu al-muqarrabūn. An
ishkur Allāh thumma udhkurhu fī al-layālī wa al-ayyām. Innahu yasmaʿu wa
yarā wa huwa al-Ḥaqq ʿAllāmu al-ghuyūb. Nūṣīka wa alladhīna āmanū
bi-al-amānati wa al-wafāʾ wa bi-mā yartafiʿu bihi kalimatu Allāh
al-Muhaymin al-Qayyūm.

O Muḥammad before Ṭáhir, tonight you are remembered by Him around whom
circle the Concourse on high, the people of Paradise, and the angels
brought near. Give thanks to God, then remember Him by night and day. He
indeed hears and sees, and He is the Truth, the Knower of all unseen
things. We exhort you and those who have believed to be trustworthy and
faithful, and to do whatever exalts the Word of God, the All-Protecting,
the Self-Subsisting.

\subsection{100}\label{section-3063}

\begin{otherlanguage}{arabic}

اذا فزت بآیاتی و وجدت عرف بیانی قل سبحانک یا ایّها المظلوم المسجون اشهد
انّک اخترت البلآء لنجاة من فی ملکوت الانشآء و سکنت فی السّجن لخلاص من علی
الأرض لا اله الّا انت المقتدر علی ما تشآء و فی قبضتک زمام الأسمآء انّک انت
العلیم الخبیر

\end{otherlanguage}

Idhā fuzta bi-āyātī wa wajadta ʿarfa bayānī qul: Subḥānaka yā ayyuhā
al-Maẓlūm al-Masjūn. Ashhadu annaka ikhtarta al-balāʾ li-najāti man fī
malakūti al-inshāʾ wa sakanta fī al-sijn li-khalāṣi man ʿalā al-arḍ. Lā
ilāha illā anta al-Muqtadir ʿalā mā tashāʾ wa fī qabḍatika zimāmu
al-asmāʾ. Innaka anta al-ʿAlīm al-Khabīr.

When you receive My verses and perceive the fragrance of My utterance,
say: Glory be to You, the Wronged One, the Imprisoned! I bear witness
that You chose affliction for the salvation of all in the kingdom of
creation and dwelt in prison for the deliverance of all on earth. There
is no god but You, who have power over whatever You will, and in Your
grasp are the reins of the names. You indeed are the All-Knowing, the
All-Informed.

\subsection{101}\label{section-3064}

\begin{otherlanguage}{arabic}

اسئلک یا سرّ الوجود و یا من فی قبضتک زمام الغیب و الشّهود بأن تؤیّدنی علی
ما یلیق لأمرک البدیع و ینبغی لیومک العزیز ثمّ اکتب لی بجودک و الطافک ما
یحفظنی عن النّعیق الّذی اخبرت به و بطغیانه ای ربّ انا العاجز البائس القائم
لدی باب رحمتک اسئلک ان لا تحرمنی عمّا عندک لا اله الّا انت المقتدر
المتعالی العزیز العلیم الحمد لک یا مقصود العارفین و محبوب من فی السّموات
و الأرضین انتهی

\end{otherlanguage}

Asʾaluka yā Sirra al-wujūd wa yā man fī qabḍatika zimāmu al-ghayb wa
al-shuhūd bi-an tuʾayyidanī ʿalā mā yalīqu li-amrika al-badīʿ wa
yanbaghī li-yawmika al-ʿazīz. Thumma uktub lī bi-jūdika wa alṭāfika mā
yaḥfaẓunī ʿan al-naʿīq alladhī akhbarta bihi wa bi-ṭughyānihi. Ay Rabb,
anā al-ʿājiz al-bāʾis al-qāʾim ladā bābi raḥmatika. Asʾaluka an lā
taḥrimanī ʿammā ʿindaka. Lā ilāha illā anta al-Muqtadir al-Mutaʿālī
al-ʿAzīz al-ʿAlīm. Al-ḥamdu laka yā Maqṣūda al-ʿārifīn wa Maḥbūba man fī
al-samāwāt wa al-arḍīn. Intahā.

I ask You, Mystery of existence, in whose grasp are the reins of the
unseen and the seen, to assist me in doing what befits Your wondrous
Cause and is worthy of Your mighty Day. Then, through Your generosity
and gracious favors, decree for me what will protect me from the
croaking one whose coming and rebellion You foretold. O Lord, I am the
helpless, wretched one standing at the door of Your mercy. I ask You not
to deprive me of what is with You. There is no god but You, the
All-Powerful, the Most Exalted, the Almighty, the All-Knowing. Praise be
to You, Desire of those who know You and Beloved of all in the heavens
and the earths. End of quotation.

\subsection{102}\label{section-3065}

\begin{otherlanguage}{persian}

اینکه ذکر آثار حضرت افنان آقای م علیه من کلّ بهآء ابهاه فرموده بودند
مخصوص ایشان الواح بدیعهٴ منیعهٴ مقدّسهٴ مبارکه نازل و از قبل ارسال شد و
اینکه خدمات ایشانرا نوشته بودید معلوم و واضح است در آن ارض نفوسی موجودند
که لدی اللّه از ایادی امر محسوبند نظر بحکمت بر حسب ظاهر ستر فرمودند ولکن
عند اهل ملأ اعلی کمال معروفیّت را دارند لأنّ المقصود نسبهم الیه فی الواح
شتّی آنچه در ذکر ایشان مرقوم داشتند جمیع عرض شد لسان عظمت باینکلمهٴ علیا
ناطق قوله جلّ جلاله لدی الوجه مذکورند و از قلم اعلی در صحیفهٴ حمرا مسطور
انتهی

\end{otherlanguage}

Īnke dhekr-e āthār-e Ḥaḍrat-e Afnān Āqā-ye M, ʿalayhi min kulli bahāʾin
abhāhu, farmūde būdand, makhṣūṣ-e īshān alvāḥ-e badīʿe-ye manīʿe-ye
moqaddase-ye mobārake nāzel va az qabl ersāl shod. Va īnke khadamāt-e
īshān rā neveshte būdīd, maʿlūm va vāḍeḥ ast. Dar ān arḍ nofūsī
mawjūdand ke ladā Allāh az Ayādī-ye Amr maḥsūband. Naẓar be-ḥekmat bar
ḥasab-e ẓāher satr farmūdand, va lāken ʿenda ahl-e malaʾ-e aʿlā kamāl-e
maʿrūfiyyat rā dārand, li-anna al-Maqṣūd nasabahum ilayhi fī alwāḥin
shattā. Ānche dar dhekr-e īshān marqūm dāshtand jamīʿ ʿarḍ shod. Lesān-e
ʿaẓamat be-īn kaleme-ye ʿolyā nāṭeq. Qawluhu jalla jalāluhu: Ladā
al-vajh madhkūrand va az qalam-e aʿlā dar ṣaḥīfe-ye ḥamrā masṭūr.
Intahā.

As to your mention of the writings of his honor the Afnán, Áqá M., upon
him be the most glorious of all glory, wondrous, exalted, holy, and
blessed Tablets were revealed especially for him and sent previously. As
to what you wrote of his services, these are known and evident. In that
land there are souls who are counted before God among the Hands of the
Cause. For wisdom's sake, He has outwardly kept them concealed, but they
are fully known among the Concourse on high, for the Desired One has
associated them with Himself in various Tablets. Everything you wrote
concerning him was presented. The Tongue of Grandeur spoke these exalted
words. His words, exalted be His glory: ``He is remembered in Our
presence and inscribed by the Most Exalted Pen in the Crimson Scroll.''
End of quotation.

\subsection{103}\label{section-3066}

\begin{otherlanguage}{arabic}

و ذکر مخدّره امة‌اللّه بی‌بی صاحب علیها بهآء اللّه نموده بودند هذا ما نزّل لها
من ملکوت بیان ربّنا الرّحمن قوله جلّ جلاله و عمّ نواله

\end{otherlanguage}

Va dhekr-e mokhaddare Amat Allāh Bībī Ṣāḥeb, ʿalayhā bahāʾu Allāh,
namūde būdand. Hādhā mā nuzzila lahā min malakūti bayāni Rabbinā
al-Raḥmān. Qawluhu jalla jalāluhu wa ʿamma nawāluhu:

You had also mentioned the honored lady, the handmaid of God, Bíbí
Ṣáḥib, upon her be the glory of God. This is what was sent down for her
from the kingdom of the utterance of our Lord, the All-Merciful. His
words, exalted be His glory and all-embracing His bounty:

\subsection{104}\label{section-3067}

\begin{otherlanguage}{arabic}

بنام دانای یکتا

\end{otherlanguage}

Be-nām-e Dānā-ye Yektā.

In the name of the One, the All-Knowing.

\subsection{105}\label{section-3068}

\begin{otherlanguage}{persian}

یا امتی حقّ جلّ جلاله مهربانست فضلش کل را احاطه نموده و همچنین علمش صدائی
مرتفع نمیشود مگر آنکه جواب از ملکوت بیان رحمن نازل اناث را بکلمهٴ
مبارکهٴ یا امتی مزیّن فرموده و رجال را بکلمهٴ بدیعهٴ منیعهٴ یا عبدی مطرّز
داشته چه بسیار از خلق اوّلین و آخرین بسبب اصغاء اینکلمه جان انفاق نمودند
اینکلمه بمثابهٴ نهالی است که ید قدرت آنرا غرس نموده باید بآب محبّت و
استقامت سقایه شود و باسم حقّ از اریاح سمومیّه محفوظ ماند قدر اینکلمه را
بدان لعمر اللّه اگر جمیع عالم ادراک نمایند کل مشتعل شوند و بافق اعلی توجّه
نمایند جمیع اماء آن ارض را تکبیر میرسانیم و کل را بتقدیس و تنزیه و عمل
به ما انزلناه فی الکتاب وصیّت مینمائیم طوبی للعاملین و العاملات انتهی

\end{otherlanguage}

Yā amatī, Ḥaqq jalla jalāluhu mehrabānast. Faḍlash koll rā eḥāṭe namūde
va hamchenīn ʿelmash. Ṣedāʾī mortafeʿ nemī-shavad magar ānke javāb az
malakūt-e bayān-e Raḥmān nāzel. Enāth rā be-kaleme-ye mobārake-ye yā
amatī mozayyan farmūde va rejāl rā be-kaleme-ye badīʿe-ye manīʿe-ye yā
ʿabdī moṭarraz dāshte. Che besyār az khalq-e avvalīn va ākherīn
be-sabab-e eṣghā-ye īn kaleme jān enfāq namūdand. Īn kaleme
be-mathābe-ye nahālī ast ke yad-e qodrat ān rā ghars namūde. Bāyad
be-āb-e maḥabbat va esteqāmat seqāye shavad va be-esm-e Ḥaqq az aryāḥ-e
samūmiyye maḥfūẓ mānad. Qadr-e īn kaleme rā bedān. La-ʿamru Allāh, agar
jamīʿ-e ʿālam edrāk namāyand koll moshtaʿel shavand va be-ofoq-e aʿlā
tavajjoh namāyand. Jamīʿ-e emā-ye ān arḍ rā takbīr mī-rasānīm va koll rā
be-taqdīs va tanzīh va ʿamal be-mā anzalnāhu fī al-kitāb vaṣiyyat
mī-namāyīm. Ṭūbā li-al-ʿāmilīn wa al-ʿāmilāt. Intahā.

O My handmaid! God, exalted be His glory, is kind. His bounty
encompasses all, as does His knowledge. No voice is raised without an
answer descending from the kingdom of the All-Merciful's utterance. He
has adorned women with the blessed words ``O My handmaid'' and clothed
men in the wondrous and exalted words ``O My servant.'' How many among
the peoples of earlier and later times have laid down their lives to
hear these words! These words are like a sapling planted by the hand of
power. It must be watered with the water of love and steadfastness and
kept safe, in the name of God, from scorching winds. Appreciate the
worth of these words. By the life of God, if the whole world were to
understand, all would be set aflame and turn toward the most exalted
horizon. We convey Our greetings to all the handmaids in that land and
exhort them all to holiness, purity, and action in accordance with what
We have sent down in the Book. Blessed are the men and women who act
accordingly. End of quotation.

\subsection{106}\label{section-3069}

\begin{otherlanguage}{arabic}

ذکر جناب حاجی محمّد حسین علیه بهآء اللّه را نموده بودند ذکر ایشان در ساحت
امنع اقدس بشرف اصغا فائز هذا ما نزّل له من لدن منزل حکیم قوله عزّ ذکره و
جلّ بیانه

\end{otherlanguage}

Dhekr-e Janāb-e Ḥājī Moḥammad Ḥosayn, ʿalayhi bahāʾu Allāh, rā namūde
būdand. Dhekr-e īshān dar sāḥat-e amnaʿ-e aqdas be-sharaf-e eṣghā fāʾez.
Hādhā mā nuzzila lahu min ladun Munzilin Ḥakīm. Qawluhu ʿazza dhikruhu
wa jalla bayānuhu:

You had mentioned Ḥájí Muḥammad Ḥusayn, upon him be the glory of God.
The mention of him was honored with a hearing in the most holy, most
inaccessible court. This is what was sent down for him by the All-Wise
Revealer. His words, exalted be His remembrance and glorious His
utterance:

\subsection{107}\label{section-3070}

\begin{otherlanguage}{arabic}

بسمی المشرق من الأفق الأعلی

\end{otherlanguage}

Bismī al-Mushriq min al-ufuq al-aʿlā.

In My name, shining from the most exalted horizon.

\subsection{108}\label{section-3071}

\begin{otherlanguage}{persian}

هر نفسی الیوم بعرفان فائز شد او در ملکوت بیان لدی الرّحمن مذکور است جمیع
عالم از برای یوم اللّه خلق شده و چون اسم اعظم بقوّهٴ غالبه خرق حجاب نمود
کل غافل و محجوب مشاهده گشتند الّا من شآء اللّه امروز روزیست که در هر نفسی
آنچه مستور است ظاهر و هویدا گردد یوم تبلی السّرائر امروز است و یوم فیه
تمرّ الجبال امروز است هر مکنونی آشکار شود و هر مستوری هویدا طوبی از برای
نفسیکه بنار محبّت الهی مشتعل شد اشتعالیکه سطوت و قدرت و شوکت اهل انشاء و
شبهات و اشارات علما او را از مالک اسماء منع ننمود و از بحر معانی محروم
نساخت هر عملی الیوم فی سبیل اللّه ظاهر شود اثرش ببقاء ملک و ملکوت باقی و
دائم است یا حزب اللّه جهد نمائید تا فائز شوید بآنچه سزاوار یوم حقّ جلّ
جلاله است قدر بیان رحمن را بدان و چون جان حفظش نما البهآء من لدنّا علیک و
علی احبّائی هناک الّذین ما خوّفتهم شوکة المعتدین و ما اضعفتهم قوّة المعرضین
انتهی

\end{otherlanguage}

Har nafsī al-yawm bi-ʿirfān fāʾiz shud ū dar malakūt-i bayān ladā
al-Raḥmān madhkūr ast. Jamīʿ-i ʿālam az barā-yi yawm Allāh khalq shudih
va chūn ism-i aʿẓam bi-quvvih-yi ghālibih kharq-i ḥijāb namūd kull
ghāfil va maḥjūb mushāhadih gashtand illā man shāʾ Allāh. Imrūz rūzīst
kih dar har nafsī ānchih mastūr ast ẓāhir va huvaydā gardad. Yawm tublā
al-sarāʾir imrūz ast va yawm fīhi tamurr al-jibāl imrūz ast. Har maknūnī
āshkār shavad va har mastūrī huvaydā. Ṭūbā az barā-yi nafsī kih bi-nār-i
maḥabbat-i ilāhī mushtaʿil shud, ishtiʿālī kih saṭvat va qudrat va
shawkat-i ahl-i inshā va shubahāt va ishārāt-i ʿulamā ū rā az Mālik-i
asmā manʿ nanamūd va az baḥr-i maʿānī maḥrūm nasākht. Har ʿamalī al-yawm
fī sabīl Allāh ẓāhir shavad atharash bi-baqā-yi mulk va malakūt bāqī va
dāʾim ast. Yā ḥizb Allāh jahd namāyīd tā fāʾiz shavīd bi-ānchih
sazāvār-i yawm-i Ḥaqq jalla jalāluhu ast. Qadr-i bayān-i Raḥmān rā bidān
va chūn jān ḥifẓash namā. Al-Bahāʾ min ladunnā ʿalayka wa-ʿalā aḥibbāʾī
hunāka alladhīna mā khawwafathum shawkat al-muʿtadīn wa-mā aḍʿafathum
quwwat al-muʿriḍīn. Intahā.

Every soul who has attained recognition today is remembered before the
All-Merciful in the kingdom of utterance. The entire world was created
for the Day of God, yet when the Most Great Name rent the veil through
His all-subduing power, all were found heedless and veiled, except those
whom God willed. Today is the day when whatever is concealed within
every soul becomes manifest and plain. The day when secrets are laid
bare is today, and the day when the mountains pass away is today. Every
hidden thing is revealed, and everything concealed is made plain.
Blessed is the soul who has been set ablaze with the fire of divine
love, with such a flame that neither the might, power, and majesty of
the people of creation nor the doubts and insinuations of the learned
have kept that soul from the Lord of Names or deprived it of the ocean
of meanings. Every deed performed today in the path of God has an
influence that endures as long as the earthly and heavenly kingdoms
endure. O people of God, strive to attain what befits the Day of the
True One, exalted be His glory. Appreciate the utterance of the
All-Merciful and guard it as you would your own life. Glory from Us be
upon you and upon My loved ones there, whom the might of the aggressors
has not frightened and the power of those who turn away has not
weakened. End.

\subsection{109}\label{section-3072}

\begin{otherlanguage}{persian}

ذکر جناب آقا سیّد صادق و آقا سیّد یوسف و آقا ابوالقاسم و آقا سیّد علیرضا
علیهم بهآء اللّه را نموده بودند تلقاء وجه عرض شد و مخصوص هر یک نازل شد
آنچه که بدوام اسماء حسنی باقی و پاینده است هذا ما نزّل لجناب سیّد صادق
قوله عزّ کبریائه

\end{otherlanguage}

Dhikr-i Jināb-i Āqā Siyyid Ṣādiq va Āqā Siyyid Yūsuf va Āqā Abu'l-Qāsim
va Āqā Siyyid ʿAlī-Riḍā ʿalayhim Bahāʾ Allāh rā namūdih būdand. Tilqāʾ-i
vajh ʿarḍ shud va makhṣūṣ-i har yak nāzil shud ānchih kih bi-davām-i
asmā-yi ḥusnā bāqī va pāyandih ast. Hādhā mā nuzzila li-Jināb-i Siyyid
Ṣādiq, qawluhu ʿazza kibriyāʾuhu.

You had mentioned Áqá Siyyid Ṣádiq, Áqá Siyyid Yúsuf, Áqá Abu'l-Qásim,
and Áqá Siyyid ʿAlí-Riḍá, upon them be the glory of God. This was
presented in His presence, and for each one in particular there was
revealed that which will endure as long as the most beautiful names
endure. This is what was sent down for Siyyid Ṣádiq. His words, exalted
be His majesty:

\subsection{110}\label{section-3073}

\begin{otherlanguage}{arabic}

هو الباقی الدّائم الأبدیّ الأحدیّ

\end{otherlanguage}

Huwa al-Bāqī al-Dāʾim al-Abadī al-Aḥadī.

He is the Everlasting, the Enduring, the Eternal, the One.

\subsection{111}\label{section-3074}

\begin{otherlanguage}{arabic}

ذکر من لدنّا لمن ذکره الأمین لیفرح بعنایة ربّه و یکون من الشّاکرین انّ الّذی
فاز الیوم بآیاتی انّه فاز بکلّ الخیر یشهد بذلک من عنده کتاب عظیم یا صادق
انّا اقبلنا الیک من شطر السّجن و انزلنا لک ما یهدی العباد الی صراطی
المستقیم قل یا قوم ضعوا ما عندکم قد ظهر ما کان مکنوناً فی العلم و مسطوراً
فی کتاب اللّه ربّ العالمین کذلک انزلنا الآیات و صرّفناها بالحقّ و ارسلناها
الیک هدیّة من محبوبک الجمیل

\end{otherlanguage}

Dhikr min ladunnā li-man dhakarahu al-Amīn li-yafraḥa bi-ʿināyat rabbihi
wa-yakūna min al-shākirīn. Inna alladhī fāza al-yawm bi-āyātī innahu
fāza bi-kull al-khayr. Yashhadu bi-dhālika man ʿindahu kitāb ʿaẓīm. Yā
Ṣādiq innā aqbalnā ilayka min shaṭr al-sijn wa-anzalnā laka mā yahdī
al-ʿibād ilā ṣirāṭī al-mustaqīm. Qul yā qawm ḍaʿū mā ʿindakum qad ẓahara
mā kāna maknūnan fī al-ʿilm wa-masṭūran fī kitāb Allāh Rabb al-ʿālamīn.
Kadhālika anzalnā al-āyāt wa-ṣarrafnāhā bi-al-ḥaqq wa-arsalnāhā ilayka
hadiyyatan min maḥbūbika al-Jamīl.

A remembrance from Us for the one whom Amín has mentioned, that he may
rejoice in the loving care of his Lord and be among the thankful.
Whoever has attained My verses today has indeed attained all good. To
this bears witness the One with whom is a mighty Book. O Ṣádiq, We have
turned toward you from the direction of the prison and have sent down
for you what guides the servants to My straight path. Say: O people, set
aside what you possess. What was concealed in knowledge and inscribed in
the Book of God, the Lord of the worlds, has now appeared. Thus have We
sent down the verses, expounded them in truth, and sent them to you as a
gift from your Beloved, the All-Beauteous.

\subsection{112}\label{section-3075}

\begin{otherlanguage}{arabic}

و نذکر یوسف فضلاً من لدنّا و انا الفضّال الکریم یا یوسف دع ما عند القوم و
خذ ما اوتیت من لدن ربّک الفضّال القدیم لا تعادل بحرف من آیاتی کنوز الأرض
یشهد بذلک کلّ عارف بصیر البهآء من لدنّا {[}علی الّذین{]} ما نقضوا میثاق
اللّه و عهده الا انّهم من الرّاسخین فی کتاب مبین

\end{otherlanguage}

Wa-nadhkuru Yūsuf faḍlan min ladunnā wa-anā al-Faḍḍāl al-Karīm. Yā Yūsuf
daʿ mā ʿinda al-qawm wa-khudh mā ūtīta min ladun rabbika al-Faḍḍāl
al-Qadīm. Lā tuʿādil bi-ḥarf min āyātī kunūz al-arḍ. Yashhadu bi-dhālika
kull ʿārif baṣīr. Al-Bahāʾ min ladunnā {[}ʿalā alladhīna{]} mā naqaḍū
mīthāq Allāh wa-ʿahdahu alā innahum min al-rāsikhīn fī kitāb mubīn.

And We remember Yúsuf as a token of Our grace, and I am the Most
Bountiful, the All-Generous. O Yúsuf, leave aside what the people
possess and take what you have been given by your Lord, the Most
Bountiful, the Ancient. Do not reckon the treasures of the earth equal
to a single letter of My verses. Every discerning soul endowed with
insight bears witness to this. Glory from Us be {[}upon those who{]}
have not broken God's covenant and His pledge. Truly, they are numbered
among the steadfast in a clear Book.

\subsection{113}\label{section-3076}

\begin{otherlanguage}{arabic}

و هذا ما نزّل لجناب ابوالقاسم قوله جلّ و عزّ

\end{otherlanguage}

Va hādhā mā nuzzila li-Jināb-i Abu'l-Qāsim, qawluhu jalla wa-ʿazza.

And this is what was sent down for Abu'l-Qásim. His words, glorified and
exalted be He:

\subsection{114}\label{section-3077}

\begin{otherlanguage}{arabic}

هو المشفق الرّحیم

\end{otherlanguage}

Huwa al-Mushfiq al-Raḥīm.

He is the Compassionate, the Merciful.

\subsection{115}\label{section-3078}

\begin{otherlanguage}{persian}

یا اباالقاسم قدر ایّام را بدان و بما ینبغی قیام نما امروز خلق نیکو و عمل
پاک سبب اعلاء کلمهٴ الهی است جمیع دوستان باید حقّ مظلوم را بآن نصرت
نمایند و در جمیع احوال در صدد ارتفاع امر باشند اینست وصیّت ما از قبل و
بعد از حقّ میطلبیم کل را مشتعل فرماید اشتعالیکه ظلم ظالمان و ضوضاء جاهلان
آنرا منع ننماید هر نفسی بیک قطره از این بحر فائز شد جنود عالم و مدافع
امم قادر بر منع او نبوده و نیست حقّ علیم شهادت میدهد و باینکلمهٴ علیا نطق
میفرماید طوبی لکلّ سامع مجیب

\end{otherlanguage}

Yā Abā al-Qāsim qadr-i ayyām rā bidān va bi-mā yanbaghī qiyām namā.
Imrūz khulq-i nīkū va ʿamal-i pāk sabab-i iʿlā-yi kalimih-yi ilāhī ast.
Jamīʿ-i dūstān bāyad Ḥaqq-i maẓlūm rā bi-ān nuṣrat namāyand va dar
jamīʿ-i aḥvāl dar ṣadad-i irtifāʿ-i amr bāshand. Īnast vaṣiyyat-i mā az
qabl va baʿd. Az Ḥaqq mīṭalabīm kull rā mushtaʿil farmāyad, ishtiʿālī
kih ẓulm-i ẓālimān va ḍawḍā-yi jāhilān ānrā manʿ nanamāyad. Har nafsī
bi-yak qaṭrih az īn baḥr fāʾiz shud junūd-i ʿālam va madāfiʿ-i umam
qādir bar manʿ-i ū nabūdih va nīst. Ḥaqq-i ʿalīm shahādat mīdihad va
bi-īn kalimih-yi ʿulyā nuṭq mīfarmāyad: Ṭūbā li-kull sāmiʿ mujīb.

O Abu'l-Qásim, appreciate the value of these days and arise to do what
is fitting. Today, good character and pure deeds are the means of
exalting the Word of God. All the friends must assist the wronged True
One through these and, in every circumstance, seek to exalt the Cause.
This is Our counsel, before and hereafter. We ask God to set everyone
ablaze with a fire that neither the oppression of the oppressors nor the
clamor of the ignorant can restrain. Whoever has attained a single drop
of this ocean, the armies of the world and the cannons of the nations
have never been, nor are they now, capable of holding back. The
All-Knowing God bears witness and speaks this exalted word: ``Blessed is
everyone who hears and responds.''

\subsection{116}\label{section-3079}

\begin{otherlanguage}{arabic}

و هذا ما نزّل لجناب سیّد علیرضا قوله تقدّس و تعالی

\end{otherlanguage}

Va hādhā mā nuzzila li-Jināb-i Siyyid ʿAlī-Riḍā, qawluhu taqaddasa
wa-taʿālā.

And this is what was sent down for Siyyid ʿAlí-Riḍá. His words,
sanctified and exalted be He:

\subsection{117}\label{section-3080}

\begin{otherlanguage}{arabic}

بسمی الظّاهر من الأفق الأعلی

\end{otherlanguage}

Bi-ismī al-Ẓāhir min al-Ufuq al-Aʿlā.

In My name, manifest from the Most Exalted Horizon.

\subsection{118}\label{section-3081}

\begin{otherlanguage}{persian}

کتاب انزله المظلوم فی السّجن الأعظم لمن ذکره احد احبّائی الّذی قام علی خدمة
امری العظیم یا عباد الرّحمن ان استمعوا ندآء من یدعوکم الی الفرد الخبیر قد
اتی یوم المآب و النّاس اکثرهم من الغافلین و اتی المعاد و القوم اکثرهم من
الرّاقدین قد نبذوا اوامر اللّه ورائهم الا انّهم من الأخسرین ان اشکر اللّه
بما ایّدک علی الاقرار بما نطق به قلمی الأعلی ثمّ اذکره فی اللّیالی و الأیّام
انّ ربّک هو السّمیع العلیم قد اتی الموعود و ظهر المستور و نطق لسان العظمة
الملک للّه ربّ العالمین انّک اذا فزت بآیاتی و شربت رحیق بیانی من ایادی
عطائی قل لک الحمد یا مقصود العالم و لک البهآء یا محبوب العارفین انتهی

\end{otherlanguage}

Kitāb anzalahu al-Maẓlūm fī al-sijn al-aʿẓam li-man dhakarahu aḥad
aḥibbāʾī alladhī qāma ʿalā khidmat amrī al-ʿaẓīm. Yā ʿibād al-Raḥmān an
istamiʿū nidāʾ man yadʿūkum ilā al-Fard al-Khabīr. Qad atā yawm al-maʾāb
wa-al-nās aktharuhum min al-ghāfilīn wa-atā al-maʿād wa-al-qawm
aktharuhum min al-rāqidīn. Qad nabadhū awāmir Allāh warāʾahum alā
innahum min al-akhsarīn. An ishkur Allāh bi-mā ayyadaka ʿalā al-iqrār
bi-mā naṭaqa bihi qalamī al-Aʿlā thumma udhkurhu fī al-layālī
wa-al-ayyām inna rabbaka huwa al-Samīʿ al-ʿAlīm. Qad atā al-Mawʿūd
wa-ẓahara al-mastūr wa-naṭaqa lisān al-ʿaẓamah: al-mulk li-llāh Rabb
al-ʿālamīn. Innaka idhā fuzta bi-āyātī wa-sharibta raḥīq bayānī min
ayādī ʿaṭāʾī qul laka al-ḥamd yā maqṣūd al-ʿālam wa-laka al-Bahāʾ yā
maḥbūb al-ʿārifīn. Intahā.

A Tablet sent down by the Wronged One in the Most Great Prison for the
one mentioned by one of My loved ones who has arisen to serve My mighty
Cause. O servants of the All-Merciful, listen to the call of Him who
summons you to the Unique, the All-Aware. The Day of Return has come,
yet most people are heedless; the promised return has come, yet most of
the people are asleep. They have cast the commandments of God behind
them. Truly, they are among those who have suffered the greatest loss.
Thank God for having enabled you to acknowledge what My Most Exalted Pen
has spoken, and remember Him by night and day. Your Lord is indeed the
All-Hearing, the All-Knowing. The Promised One has come, the Hidden One
has appeared, and the Tongue of Grandeur has proclaimed: ``Sovereignty
belongs to God, the Lord of the worlds.'' Now that you have attained My
verses and drunk the choice wine of My utterance from the hands of My
bounty, say: ``Praise be to You, O Desire of the world, and glory be to
You, O Beloved of those who know You.'' End.

\subsection{119}\label{section-3082}

\begin{otherlanguage}{persian}

رشحات بحر معانی و بیان عالم را اخذ نموده اگر جمیع ذرّات بلسان اوّلین و
آخرین حمد نمایند البتّه لایق این فضل و قابل این جود نبوده و نخواهد بود
للّه الحمد دوستان این عصر فائز شدند بآنچه که خلق کثیر از اوّلین و آخرین در
حسرت لقایش محزون و مهموم بودند این فرح اکبر نصیب اهل منظر اکبر بوده
اینفانی از حقّ باقی سائل و آمل که جمیع من علی الأرض را از این نعمت باقیه
و مائدهٴ سمائیّه محروم نفرماید انّه هو السّامع المجیب و هو الغفور الرّحیم

\end{otherlanguage}

Rashaḥāt-i baḥr-i maʿānī va bayān ʿālam rā akhdh namūdih. Agar jamīʿ-i
dharrāt bi-lisān-i avvalīn va ākhirīn ḥamd namāyand al-battih lāyiq-i īn
faḍl va qābil-i īn jūd nabūdih va nakhvāhad būd. Li-llāh al-ḥamd
dūstān-i īn ʿaṣr fāʾiz shudand bi-ānchih kih khalq-i kathīr az avvalīn
va ākhirīn dar ḥasrat-i liqāyash maḥzūn va mahmūm būdand. Īn faraḥ-i
akbar naṣīb-i ahl-i manẓar-i akbar būdih. Īn fānī az Ḥaqq-i bāqī sāʾil
va āmil kih jamīʿ-i man ʿalā al-arḍ rā az īn niʿmat-i bāqiyih va
māʾidih-yi samāʾiyyih maḥrūm nafarmāyad. Innahu huwa al-Sāmiʿ al-Mujīb
wa-huwa al-Ghafūr al-Raḥīm.

The sprays of the ocean of meanings and utterance have encompassed the
world. Were every atom to offer praise with the tongues of all who came
before and all who are to come, it would assuredly never have been, nor
ever be, adequate to this grace or worthy of this bounty. Praise be to
God, the friends of this age have attained that for which multitudes
among earlier and later generations were sorrowful and grief-stricken in
their longing to attain it. This greatest joy has been the portion of
the people of the Most Great Scene. This transient one asks and hopes of
the everlasting God that He will deprive no one on earth of this
enduring bounty and heavenly table. He is indeed the One who hears and
answers, and He is the Ever-Forgiving, the Merciful.

\subsection{120}\label{section-3083}

\begin{otherlanguage}{persian}

و اینکه ذکر جناب آقا محمّد علی سمسار علیه بهآء اللّه نموده بودند تلقاء وجه
مالک قدم عرض شد هذا ما نزّل له من سمآء عنایة ربّنا الغفور الکریم قوله عزّ
بیانه و عظم برهانه

\end{otherlanguage}

Va īnkih dhikr-i Jināb-i Āqā Muḥammad ʿAlī Simsār ʿalayhi Bahāʾ Allāh
namūdih būdand tilqāʾ-i vajh-i Mālik-i qidam ʿarḍ shud. Hādhā mā nuzzila
lahu min samāʾ ʿināyat Rabbinā al-Ghafūr al-Karīm, qawluhu ʿazza
bayānuhu wa-ʿaẓuma burhānuhu.

And your mention of Áqá Muḥammad ʿAlí Simsár, upon him be the glory of
God, was presented before the face of the Lord of Eternity. This is what
was sent down for him from the heaven of the loving care of our Lord,
the Ever-Forgiving, the All-Generous. His words, exalted be His
utterance and mighty His proof:

\subsection{121}\label{section-3084}

\begin{otherlanguage}{arabic}

هو السّمیع البصیر

\end{otherlanguage}

Huwa al-Samīʿ al-Baṣīr.

He is the All-Hearing, the All-Seeing.

\subsection{122}\label{section-3085}

\begin{otherlanguage}{persian}

یا محمّد قبل علی حقّ جلّ جلاله ظاهر و جمیع را بصراط مستقیم راه نمود و بنبأ
اعظم که در کتب و صحف از قبل و بعد موعود بود دعوت فرمود بعضی از عباد را
هواهای نفسانیّه از مطلع نور احدیّه منع نمود و برخی را متابعت جهلا از عرفان
مالک اسماء محروم ساخت سبب و علّت کل در مقام اوّل علمای عصر بوده‌اند ایشانند
مطالع و مشارق نفس و هوی و قلیلی از عالم و مافیه چشم برداشتند و بافق اعلی
توجّه نمودند ایشانند کنوز ارض و خزائن آن هر نفسی الیوم حجبات اهل شبهات او
را منع ننمود او از اهل جنّت علیا در کتاب اسماء مذکور و مسطور است جناب
امین جمعی از مقبلین را ذکر نمود و هر یک بعنایت مظلوم فائز {[}انّه هو
الغفور{]} الکریم انّه هو العطوف الرّحیم انتهی

\end{otherlanguage}

Yā Muḥammad qabl-i ʿAlī Ḥaqq jalla jalāluhu ẓāhir va jamīʿ rā bi-ṣirāṭ-i
mustaqīm rāh namūd va bi-nabaʾ-i aʿẓam kih dar kutub va ṣuḥuf az qabl va
baʿd mawʿūd būd daʿvat farmūd. Baʿḍī az ʿibād rā havāhā-yi nafsāniyyih
az maṭlaʿ-i nūr-i aḥadiyyih manʿ namūd va barkhī rā mutābaʿat-i juhalā
az ʿirfān-i Mālik-i asmā maḥrūm sākht. Sabab va ʿillat-i kull dar
maqām-i avval ʿulamā-yi ʿaṣr būdih-and. Īshānand maṭāliʿ va mashāriq-i
nafs va havā. Va qalīlī az ʿālam va mā fīh chashm bardāshtand va
bi-ufuq-i aʿlā tavajjuh namūdand. Īshānand kunūz-i arḍ va khazāʾin-i ān.
Har nafsī al-yawm ḥujubāt-i ahl-i shubahāt ū rā manʿ nanamūd ū az ahl-i
jannat-i ʿulyā dar kitāb-i asmā madhkūr va masṭūr ast. Jināb-i Amīn
jamʿī az muqbilīn rā dhikr namūd va har yak bi-ʿināyat-i Maẓlūm fāʾiz.
{[}Innahu huwa al-Ghafūr{]} al-Karīm. Innahu huwa al-ʿAṭūf al-Raḥīm.
Intahā.

O Muḥammad before ʿAlí, the True One, exalted be His glory, has
appeared, guided everyone to the straight path, and summoned them to the
Most Great Announcement promised in the books and scriptures, earlier
and later. Selfish desires kept some of the servants from the
Dawning-Place of the light of oneness, while following the ignorant
deprived others of recognizing the Lord of Names. The primary cause and
reason for all this have been the learned of the age. They are the
dawning-places and daysprings of self and desire. A few, however, turned
their eyes away from the world and all it contains and directed
themselves toward the Most Exalted Horizon. They are the treasures of
the earth and its treasuries. Every soul whom the veils of the people of
doubt have not held back today is mentioned and recorded in the Book of
Names as one of the inhabitants of the highest Paradise. Amín mentioned
a number of those who have turned toward Us, and each has received the
loving care of the Wronged One. {[}He is indeed the Ever-Forgiving,{]}
the All-Generous. He is indeed the All-Loving, the Merciful. End.

\subsection{123}\label{section-3086}

\begin{otherlanguage}{persian}

و این خادم فانی از حقّ جلّ جلاله مسئلت مینماید که ایشان و سایر دوستانرا
تأیید فرماید بر حفظ آنچه عطا فرموده اوست بر هر شیء قادر و توانا عرض خادم
آنکه هر یک از مقبلین که از رحیق مختوم آشامیده‌اند و بحیوة جدید فائز
گشته‌اند کلمهٴ استقامت را امام وجه خود مشاهده نمایند چه که شیاطین و نفوس
خادعه بسیارند امید چنانکه هر یک مؤیّد شوند بر آنچه سزاوار یوم اللّه‌ است

\end{otherlanguage}

Va īn khādim-i fānī az Ḥaqq jalla jalāluhu masʾalat mīnamāyad kih īshān
va sāʾir-i dūstān rā taʾyīd farmāyad bar ḥifẓ-i ānchih ʿaṭā farmūdih.
Ūst bar har shayʾ qādir va tavānā. ʿArḍ-i khādim ānkih har yak az
muqbilīn kih az raḥīq-i makhtūm āshāmīdih-and va bi-ḥayāt-i jadīd fāʾiz
gashtih-and kalimih-yi istiqāmat rā imām-i vajh-i khud mushāhadih
namāyand chih kih shayāṭīn va nufūs-i khādiʿih bisyārand. Umīd chinānkih
har yak muʾayyad shavand bar ānchih sazāvār-i yawm Allāh ast.

And this transient servant beseeches the True One, exalted be His glory,
to assist them and the other friends to preserve what He has bestowed.
He is powerful and mighty over all things. This servant's plea is that
every one of those who have turned toward Him, who have drunk the sealed
wine and attained a new life, should keep the word of steadfastness
before their eyes, for the devils and deceitful souls are many. The hope
is that each may be aided to do what befits the Day of God.

\subsection{124}\label{section-3087}

\begin{otherlanguage}{arabic}

و دیگر ذکر جناب آقا سیّد جواد پشت‌باغی علیه ٩ ٦٦ {[}بهآء اللّه{]} را نموده
بودند بعد از عرض اینفقره تلقاء وجه این آیات از سماء عنایت نازل قوله سبقت
رحمته و سبغت نعمته

\end{otherlanguage}

Va dīgar dhikr-i Jināb-i Āqā Siyyid Javād-i Pusht-Bāghī ʿalayhi 9 66
{[}Bahāʾ Allāh{]} rā namūdih būdand. Baʿd az ʿarḍ-i īn fiqrih tilqāʾ-i
vajh īn āyāt az samā-yi ʿināyat nāzil, qawluhu sabaqat raḥmatuhu
wa-sabaghat niʿmatuhu.

You had also mentioned Áqá Siyyid Javád Pusht-Bághí, upon him be 9 66
{[}the glory of God{]}. After this matter was presented in His presence,
these verses descended from the heaven of loving care. His words, whose
mercy has preceded all things and whose bounty is all-embracing:

\subsection{125}\label{section-3088}

\begin{otherlanguage}{arabic}

هو الشّاهد الخبیر

\end{otherlanguage}

Huwa al-Shāhid al-Khabīr.

He is the Witness, the All-Aware.

\subsection{126}\label{section-3089}

\begin{otherlanguage}{persian}

یا جواد مظلوم میفرماید جمیع بلایای ارض را قبول نمودیم لارتقآء النّاس الی
الذّروة العلیا و الأفق الأعلی و چون ندا مرتفع شد و نار محبّت الهی مشتعل
نفوس غافله بر اخماد و منع قیام نمودند و بظلم تمام و جهد کامل در اطفاء
نور الهی کوشیدند ولکن ید قدرت امر خود را حفظ نمود و نفوس مقدّسه را آگاه
فرمود بگو ای دوستان شیاطین بر مراصد قائم و منتظر باید هر یک از شما
بمثابهٴ شهاب ثاقب شیاطین را از صعود بسموات قلوب منع نمائید اینست مقام
اوّل و رتبهٴ اولی که از قلم اعلی نازل گشته طوبی لمن فاز و ویل للغافلین در
جمیع احوال بحقّ ناظر باش و باو متمسّک انّه مع عباده الفائزین الحمد للّه ربّ
العالمین انتهی

\end{otherlanguage}

Yā Javād Maẓlūm mīfarmāyad jamīʿ-i balāyā-yi arḍ rā qabūl namūdīm
li-irtiqāʾ al-nās ilā al-dhirwah al-ʿulyā wa-al-ufuq al-aʿlā va chūn
nidā murtafiʿ shud va nār-i maḥabbat-i ilāhī mushtaʿil, nufūs-i ghāfilih
bar ikhmād va manʿ qiyām namūdand va bi-ẓulm-i tamām va jahd-i kāmil dar
iṭfā-yi nūr-i ilāhī kūshīdand. Valākin yad-i qudrat amr-i khud rā ḥifẓ
namūd va nufūs-i muqaddasih rā āgāh farmūd. Bigū ay dūstān shayāṭīn bar
marāṣid qāʾim va muntaẓir. Bāyad har yak az shumā bi-mathābih-yi
shihāb-i thāqib shayāṭīn rā az ṣuʿūd bi-samavāt-i qulūb manʿ namāyīd.
Īnast maqām-i avval va rutbih-yi ūlā kih az qalam-i aʿlā nāzil gashtih.
Ṭūbā li-man fāza wa-wayl li-al-ghāfilīn. Dar jamīʿ-i aḥvāl bi-Ḥaqq nāẓir
bāsh va bi-ū mutamassik. Innahu maʿa ʿibādihi al-fāʾizīn. Al-ḥamd
li-llāh Rabb al-ʿālamīn. Intahā.

O Javád, the Wronged One says: We accepted all the afflictions of the
earth so that people might ascend to the highest summit and the Most
Exalted Horizon. Yet when the call was raised and the fire of divine
love was kindled, heedless souls arose to quench it and obstruct it, and
with utter oppression and every effort they strove to extinguish the
light of God. But the hand of power protected His Cause and awakened
holy souls. Say: O friends, the devils are lying in wait at their
watchposts. Each of you must, like a piercing meteor, prevent the devils
from ascending into the heavens of the hearts. This is the first station
and foremost rank that has been revealed by the Most Exalted Pen.
Blessed is the one who has attained, and woe to the heedless. In every
circumstance look to God and hold fast to Him. He is with those of His
servants who have attained. Praise be to God, the Lord of the worlds.
End.

\subsection{127}\label{section-3090}

\begin{otherlanguage}{arabic}

هذا ما نزّل لابن من صعد الی اللّه (آقا میرزا محسن ط‌ب) علیهما بهآء اللّه

\end{otherlanguage}

Hādhā mā nuzzila li-ibn man ṣaʿida ilā Allāh (Āqā Mīrzā Muḥsin Ṭ-B)
ʿalayhimā Bahāʾ Allāh.

This is what was sent down for the son of the one who ascended to God
(Áqá Mírzá Muḥsin, Ṭ-B), upon them both be the glory of God.

\subsection{128}\label{section-3091}

\begin{otherlanguage}{arabic}

بسمی المهیمن القیّوم

\end{otherlanguage}

Bi-ismī al-Muhaymin al-Qayyūm.

In My name, the All-Protecting, the Self-Subsisting.

\subsection{129}\label{section-3092}

\begin{otherlanguage}{persian}

یا حسین قبل علی ان استمع النّدآء انّه لا اله الّا انا المهیمن القیّوم قد
خلقت الخلق لعرفان نفسی و ادراک یومی فلمّا جئت برایات آلایات اعرضوا بما
اتّبعوا کلّ غافل مردود قل یا اهل البهآء ضعوا ما نهیتم عنه فی الکتاب خذوا
ما امرتم به من لدی اللّه مالک الیوم الموعود قدّسوا انفسکم عن کلّ ما یأمرکم
به اهوائکم کذلک امرناکم به من قبل فی لوحی المسطور طوبی لمن نبذ الأوهام
متمسّکاً بجبل اللّه العزیز الودود طوبی لمن تشبّث بذیلی و عمل بما امر به من
لدی الحقّ علّام الغیوب ذکّر من قبلی من معک قل اتّقوا اللّه و لا تکونوا من
الّذینهم لا یعلمون قد اتی الدّلیل و ظهر السّبیل و النّاس اکثرهم لا یشعرون قل
تاللّه قد ظهر الیوم و مکلّم الطّور ینطق علی عرش الظّهور طوبی لمن سمع و فاز و
ویل لکلّ غافل محجوب البهآء من لدنّا علی ابیک الّذی صعد الی افقی و فاز بفضلی
و رحمتی و عنایتی الّتی احاطت الغیب و الشّهود و علیک و علی من معک من الّذین
ما نقضوا میثاق اللّه و عهده و شهدوا بما شهد اللّه اذ اتی بسلطان مشهود
انتهی

\end{otherlanguage}

Yā Ḥusayn qabl ʿAlī, an istamiʿ al-nidāʾ, innahu lā ilāha illā anā
al-Muhaymin al-Qayyūm. Qad khalaqtu al-khalq li-ʿirfān nafsī wa-idrāk
yawmī, fa-lammā jiʾtu bi-rāyāt al-āyāt aʿraḍū bi-mā ittabaʿū kulla
ghāfil mardūd. Qul yā ahl al-Bahāʾ, ḍaʿū mā nuhītum ʿanhu fī al-Kitāb,
khudhū mā umirtum bihi min ladā Allāh Mālik al-yawm al-mawʿūd. Qaddisū
anfusakum ʿan kulli mā yaʾmurukum bihi ahwāʾukum. Kadhālika amarnākum
bihi min qabl fī lawḥī al-masṭūr. Ṭūbā li-man nabadha al-awhām
mutamassikan bi-jabal Allāh al-ʿAzīz al-Wadūd. Ṭūbā li-man tashabbatha
bi-dhaylī wa-ʿamila bi-mā umira bihi min ladā al-Ḥaqq ʿAllām al-ghuyūb.
Dhakkir min qibalī man maʿak. Qul ittaqū Allāh wa-lā takūnū min
alladhīna hum lā yaʿlamūn. Qad atā al-dalīl wa-ẓahara al-sabīl
wa-al-nāsu aktharuhum lā yashʿurūn. Qul ta-Allāh qad ẓahara al-yawm
wa-Mukallim al-Ṭūr yanṭiqu ʿalā ʿarsh al-ẓuhūr. Ṭūbā li-man samiʿa
wa-fāza wa-waylun li-kulli ghāfil maḥjūb. Al-Bahāʾ min ladunnā ʿalā
abīka alladhī ṣaʿida ilā ufuqī wa-fāza bi-faḍlī wa-raḥmatī wa-ʿināyatī
allatī aḥāṭat al-ghayb wa-al-shuhūd, wa-ʿalayka wa-ʿalā man maʿaka min
alladhīna mā naqaḍū mīthāq Allāh wa-ʿahdahu wa-shahidū bi-mā shahida
Allāh idh atā bi-sulṭān mashhūd. Intahā.

O Ḥusayn-ʿAlī, listen to the call: there is no God but Me, the
All-Protecting, the Self-Subsisting. I created all creation to recognize
My own Self and attain My Day, yet when I came with the banners of the
signs, they turned away because they followed every heedless outcast.
Say: O people of Bahā, put aside what you have been forbidden in the
Book, and embrace what you have been commanded by God, the Master of the
promised Day. Sanctify yourselves from all that your desires command you
to do. Thus We commanded you before in My inscribed Tablet. Blessed is
the one who has cast aside vain imaginings, holding fast to the mountain
of God, the Almighty, the All-Loving. Blessed is the one who has clung
to the hem of My robe and acted on what was commanded by the True One,
the Knower of all that is hidden. Remind those with you on My behalf.
Say: Fear God and do not be among those who do not know. The Guide has
come and the path has appeared, yet most people are unaware. Say: By
God, the Day has appeared, and the One who spoke on Sinai speaks upon
the throne of revelation. Blessed is the one who has heard and attained,
and woe to every heedless one who is veiled. Glory from Us be upon your
father, who ascended to My horizon and attained My grace, My mercy, and
My loving care, which have encompassed the unseen and the seen; and upon
you and those with you who have not broken God's covenant and pledge,
and who have testified to what God testified when He came with manifest
sovereignty. End.

\subsection{130}\label{section-3093}

\begin{otherlanguage}{persian}

اینکه ذکر اخوان از اهل الف و ر جناب نصراللّه و حسین و محمّد علی علیهم بهآء
اللّه نموده بودند در پیشگاه حضور مالک ظهور عرض شد هذا ما نزّل لهم من سمآء
فضل ربّنا العزیز الوهّاب قوله جلّ بیانه و عزّ مقامه

\end{otherlanguage}

Īnke dhikr-e ikhvān az ahl-e Alif va Rā, jināb-e Naṣru'llāh va Ḥusayn va
Muḥammad ʿAlī, ʿalayhim bahāʾu'llāh, namūde būdand, dar pīshgāh-e
ḥuḍūr-e Mālik-e Ẓuhūr ʿarḍ shod. Hādhā mā nuzzila lahum min samāʾ faḍl
Rabbinā al-ʿAzīz al-Wahhāb. Qawluhu jalla bayānuhu wa-ʿazza maqāmuhu.

The mention you made of the brothers from Alif and Rā---Naṣru'llāh,
Ḥusayn, and Muḥammad ʿAlī, upon them be the glory of God---was submitted
in the presence of the Lord of Revelation. This is what was sent down
for them from the heaven of the grace of our Lord, the Almighty, the
All-Bestowing. His words---glorious is His utterance and exalted His
station:

\subsection{131}\label{section-3094}

\begin{otherlanguage}{arabic}

هو الحاکم بالحقّ و هو الآمر الحکیم

\end{otherlanguage}

Huwa al-Ḥākim bi-al-ḥaqq wa-Huwa al-Āmir al-Ḥakīm.

He is the Ruler in truth, and He is the All-Wise Ordainer.

\subsection{132}\label{section-3095}

\begin{otherlanguage}{persian}

یا عباد الرّحمن مقصود امکان در این سجن اعظم شما را ذکر مینماید و وصیّت
میفرماید بآنچه که سبب و علّت بزرگی و بزرگواریست دنیا را سطوت فنا اخذ
نموده و در هر حین آثارش ظاهر و باهر دوستانیکه بحقّ منسوبند و باسمش مذکور
باید بحول و قوّهٴ مقصود یکتا بر اعمالی قیام نمایند که هر صاحب شمّی عرف
تقدیس و تنزیه از آن بیابد امروز روز خدمت و عمل پاک است از حقّ میطلبیم کل
را از آلایش اعمال مردوده مقدّس فرماید و بآسایش حقیقی فائز نماید حمد کنید
حضرت مقصود را که شما را بانوار آفتاب بیان منوّر فرمود قد ذکرکم الأمین
ذکرناکم فضلاً من عندنا و انا الفضّال الکریم انتهی

\end{otherlanguage}

Yā ʿibād al-Raḥmān, Maqṣūd-e imkān dar īn sijn-e aʿẓam shomā rā dhikr
mīnamāyad va vaṣiyyat mīfarmāyad bi-ānche ke sabab va ʿillat-e bozorgī
va bozorgvārīst. Dunyā rā saṭvat-e fanā akhdh namūde va dar har ḥīn
āthārash ẓāhir va bāhir. Dūstānī ke bi-Ḥaqq mansūband va bi-ismash
madhkūr bāyad bi-ḥawl va qovve-ye Maqṣūd-e yiktā bar aʿmālī qiyām
namāyand ke har ṣāḥib-e shammī ʿarf-e taqdīs va tanzīh az ān biyābad.
Imrūz rūz-e khidmat va ʿamal-e pāk ast. Az Ḥaqq mīṭalabīm kull rā az
ālāyish-e aʿmāl-e mardūde muqaddas farmāyad va bi-āsāyish-e ḥaqīqī fāʾiz
namāyad. Ḥamd konīd Ḥaḍrat-e Maqṣūd rā ke shomā rā bi-anvār-e āftāb-e
bayān munavvar farmūd. Qad dhakarakum al-Amīn, dhakarnākum faḍlan min
ʿindinā wa-anā al-Faḍḍāl al-Karīm. Intahā.

O servants of the All-Merciful, the Desire of all existence remembers
you in this Most Great Prison and counsels you to do what brings
greatness and nobility. The power of extinction has seized the world,
and its effects are evident and manifest at every moment. The friends
who are associated with God and known by His name must, through the
power and strength of the one Beloved, undertake deeds from which
everyone endowed with a sense of smell may perceive the fragrance of
holiness and purity. Today is the day for service and pure deeds. We ask
God to sanctify everyone from the stain of unacceptable deeds and enable
them to attain true repose. Praise the Beloved, who has illumined you
with the rays of the sun of utterance. Amīn has mentioned you, and We
have remembered you as a grace from Us. I am the Most Bountiful, the
All-Generous. End.

\subsection{133}\label{section-3096}

\begin{otherlanguage}{arabic}

و هذا ما نزّل لأحبّآء اللّه فی الدّال و الهآء قوله جلّ جلاله و عزّ بیانه

\end{otherlanguage}

Wa-hādhā mā nuzzila li-aḥibbāʾ Allāh fī al-Dāl wa-al-Hāʾ. Qawluhu jalla
jalāluhu wa-ʿazza bayānuhu.

And this is what was sent down for the loved ones of God in Dāl and Hā.
His words---exalted is His glory and mighty His utterance:

\subsection{134}\label{section-3097}

\begin{otherlanguage}{arabic}

هو السّامع البصیر

\end{otherlanguage}

Huwa al-Sāmiʿ al-Baṣīr.

He is the All-Hearing, the All-Seeing.

\subsection{135}\label{section-3098}

\begin{otherlanguage}{arabic}

یا حزب اللّه یذکّرکم المظلوم من شطر السّجن بما یقرّبکم الی اللّه ربّ العالمین
قد ماج بحر البیان و هاج عرف الرّحمن و النّاس اکثرهم من الغافلین قد جری
فرات الرّحمة بین البریّة ولکنّ القوم اکثرهم من المنکرین قد اعرضوا عن الأفق
الأعلی مقبلین الی کلّ جاهل بعید طوبی لمن کسر اصنام الهوی و آمن باللّه
الفرد الخبیر

\end{otherlanguage}

Yā ḥizb Allāh, yudhakkirukum al-Maẓlūm min shaṭr al-sijn bi-mā
yuqarribukum ilā Allāh Rabb al-ʿālamīn. Qad māja baḥr al-bayān wa-hāja
ʿarf al-Raḥmān wa-al-nāsu aktharuhum min al-ghāfilīn. Qad jarā furāt
al-raḥma bayna al-bariyya wa-lākinna al-qawma aktharuhum min
al-munkirīn. Qad aʿraḍū ʿan al-Ufuq al-Aʿlā muqbilīn ilā kulli jāhil
baʿīd. Ṭūbā li-man kasara aṣnām al-hawā wa-āmana bi-Allāh al-Fard
al-Khabīr.

O people of God, from the direction of the Prison the Wronged One
reminds you of what will draw you near to God, the Lord of the worlds.
The ocean of utterance has surged and the fragrance of the All-Merciful
has spread, yet most people are heedless. The sweet waters of mercy have
flowed among creation, yet most people deny them. They have turned away
from the Most Exalted Horizon, turning instead toward every ignorant one
who is far removed. Blessed is the one who has shattered the idols of
desire and believed in God, the One, the All-Informed.

\subsection{136}\label{section-3099}

\begin{otherlanguage}{arabic}

انّا نذکر الّذی اقبل الی الأفق الأعلی و قطع البرّ و البحر الی ان ورد فی واد
فیه ارتفع ندآء اللّه العزیز الحمید انّه حضر و قام لدی الباب و فاز باصغآء
بیان ربّه العزیز الجمیل نشهد انّه بلغ و فاز باللّقآء و شرب رحیق عطآء ربّه
المشفق الکریم

\end{otherlanguage}

Innā nadhkuru alladhī aqbala ilā al-Ufuq al-Aʿlā wa-qaṭaʿa al-barr
wa-al-baḥr ilā an warada fī wādin fīhi irtafaʿa nidāʾ Allāh al-ʿAzīz
al-Ḥamīd. Innahu ḥaḍara wa-qāma ladā al-bāb wa-fāza bi-iṣghāʾ bayān
Rabbihi al-ʿAzīz al-Jamīl. Nashhadu innahu balagha wa-fāza bi-al-liqāʾ
wa-shariba raḥīq ʿaṭāʾ Rabbihi al-Mushfiq al-Karīm.

We remember the one who turned toward the Most Exalted Horizon and
crossed land and sea until he arrived in a valley where the call of God,
the Almighty, the All-Praised, was raised. He came, stood at the door,
and attained the privilege of listening to the utterance of his Lord,
the Almighty, the All-Beautiful. We testify that he reached his
destination, attained the meeting, and drank the choice wine of the
bounty of his Lord, the Compassionate, the All-Generous.

\subsection{137}\label{section-3100}

\begin{otherlanguage}{arabic}

و نذکر ابنه الّذی سمّی بالبدیع و نوصیه بما ینبغی للانسان و انا النّاصح
العلیم یا بدیع یذکرک مولاک بما یبقی به ذکرک بدوام الملک انّ ربّک لهو
الغفور الرّحیم البهآء علیک و علی من معک و علی الّذین فازوا بهذا المقام
الرّفیع

\end{otherlanguage}

Wa-nadhkuru ibnahu alladhī summiya bi-al-Badīʿ wa-nūṣīhi bi-mā yanbaghī
li-al-insān wa-anā al-Nāṣiḥ al-ʿAlīm. Yā Badīʿ, yadhkuruka Mawlāka bi-mā
yabqā bihi dhikruka bi-dawām al-mulk. Inna Rabbaka la-Huwa al-Ghafūr
al-Raḥīm. Al-Bahāʾ ʿalayka wa-ʿalā man maʿaka wa-ʿalā alladhīna fāzū
bi-hādhā al-maqām al-rafīʿ.

And We remember his son, who is named Badīʿ, and counsel him to do what
befits a human being. I am the All-Knowing Counsellor. O Badīʿ, your
Lord remembers you in such a way that your remembrance will endure as
long as creation endures. Your Lord is indeed the Ever-Forgiving, the
All-Merciful. Glory be upon you, upon those with you, and upon those who
have attained this lofty station.

\subsection{138}\label{section-3101}

\begin{otherlanguage}{arabic}

و نذکر علی قبل اکبر الّذی وجد عرف الظّهور و اقبل الی الّذی اعرض عنه کلّ جاهل
مریب و نذکر اخاه الّذی سمّی بموسی و نبشّره برحمتی الّتی سبقت من فی السّموات و
الأرضین و نوصی الکلّ بما یرتفع به امر اللّه بین العباد و بما ینبغی لهذا
الیوم العزیز البدیع

\end{otherlanguage}

Wa-nadhkuru ʿAlī qabl Akbar alladhī wajada ʿarf al-ẓuhūr wa-aqbala ilā
alladhī aʿraḍa ʿanhu kullu jāhil murīb. Wa-nadhkuru akhāhu alladhī
summiya bi-Mūsā wa-nubashshiruhu bi-raḥmatī allatī sabaqat man fī
al-samāwāt wa-al-araḍīn. Wa-nūṣī al-kull bi-mā yartafiʿu bihi Amr Allāh
bayna al-ʿibād wa-bi-mā yanbaghī li-hādhā al-yawm al-ʿAzīz al-Badīʿ.

And We remember ʿAlī-Akbar, who perceived the fragrance of revelation
and turned toward the One from whom every ignorant doubter turned away.
And We remember his brother, who is named Mūsā, and give him the glad
news of My mercy, which has preceded all who are in the heavens and the
earths. We counsel everyone to do what will exalt the Cause of God among
His servants and what befits this mighty and wondrous Day.

\subsection{139}\label{section-3102}

\begin{otherlanguage}{arabic}

قد ذکرکم الأمین ذکرناکم لتشکروا ربّکم المقتدر القدیر یا اهل الدّال و الهآء
ان افرحوا بما اظهر اللّه منکم اسمه الّذی سجن فی سبیله و حمل الشّدائد حبّاً
لجماله و نطق بثنائه اذ کان النّاس فی حجاب مبین و نذکر عبادی و امائی هناک
و نبشّرهم بعنایتی و فضلی و نوصیهم بالعمل الخالص و بما انزله الرّحمن فی
کتابه المبین البهآء علیهم و علیهنّ و علی کلّ ثابت مستقیم انتهی

\end{otherlanguage}

Qad dhakarakum al-Amīn, dhakarnākum li-tashkurū Rabbakum al-Muqtadir
al-Qadīr. Yā ahl al-Dāl wa-al-Hāʾ, an ifraḥū bi-mā aẓhara Allāh minkum
ismahu alladhī sujina fī sabīlihi wa-ḥamala al-shadāʾid ḥubban
li-jamālihi wa-naṭaqa bi-thanāʾihi idh kāna al-nāsu fī ḥijāb mubīn.
Wa-nadhkuru ʿibādī wa-imāʾī hunāka wa-nubashshiruhum bi-ʿināyatī
wa-faḍlī wa-nūṣīhim bi-al-ʿamal al-khāliṣ wa-bi-mā anzalahu al-Raḥmān fī
kitābihi al-mubīn. Al-Bahāʾ ʿalayhim wa-ʿalayhinna wa-ʿalā kulli thābit
mustaqīm. Intahā.

Amīn has mentioned you, and We have remembered you so that you may thank
your Lord, the All-Powerful, the Omnipotent. O people of Dāl and Hā,
rejoice that God has raised up from among you His name, who was
imprisoned in His path, bore hardships out of love for His beauty, and
spoke His praise when the people were manifestly veiled. And We remember
My servants and My handmaidens there, give them the glad news of My
loving care and grace, and counsel them to act with pure intent and to
observe what the All-Merciful has sent down in His clear Book. Glory be
upon them, men and women, and upon everyone who is steadfast and
upright. End.

\subsection{140}\label{section-3103}

\begin{otherlanguage}{arabic}

للّه الحمد هر مقبلی باجر خود فائز بل فوق آن و هر مستقیمی بمکافات خود مزیّن
اکثری از دوستان در منازل و محافل خود معتکف و حقّ جلّ جلاله در سجن مع احزان
وارده و ظهور امورات غیر مرضیّه در لیالی و ایّام بذکر احبّای خود مشغول
ذکریکه فنا او را نبیند و تغییر او را اخذ نکند بقول جناب شیخ معروف علیه
رحمة اللّه و غفرانه از دست و زبان که برآید کز عهدهٴ شکرش بدرآید

\end{otherlanguage}

Li'llāh al-ḥamd, har muqbilī bi-ajr-e khod fāʾiz, bal fawq-e ān, va har
mustaqīmī bi-mukāfāt-e khod muzayyan. Aktharī az dūstān dar manāzil va
maḥāfil-e khod muʿtakif va Ḥaqq jalla jalāluhu dar sijn maʿa aḥzān-e
vārede va ẓuhūr-e umūrāt-e ghayr-e marḍiyye dar layālī va ayyām
bi-dhikr-e aḥibbā-ye khod mashghūl, dhikrī ke fanā ū rā nabīnad va
taghyīr ū rā akhdh nakonad. Bi-qawl-e jināb-e Shaykh-e maʿrūf, ʿalayhi
raḥmatu'llāh va ghufrānuhu, az dast va zabān ke barāyad k-az ʿohde-ye
shokrash badarāyad.

Praise be to God, everyone who has turned toward Him has attained his
reward---indeed, more than that---and everyone who is steadfast is
adorned with his recompense. Most of the friends remain in their own
homes and gatherings, while God---exalted is His glory---in prison,
despite the sorrows that have befallen Him and the emergence of
displeasing affairs, is occupied night and day in remembering His loved
ones: a remembrance that extinction will never touch and change will
never overtake. In the words of the renowned Shaykh, upon him be God's
mercy and forgiveness: ``Whose hand and tongue can ever offer the thanks
He is due?''

\subsection{141}\label{section-3104}

\begin{otherlanguage}{arabic}

قسم بآفتاب حقیقت اگر جمیع عالم بدوام ملک و ملکوت بشکر و حمد ناطق شوند
قابل یکحرف که از قلم اعلی جاری میشود نبوده و نیستند اگر مقام ذکر الهی
ظاهر شود کل ناطق شوند بآنچه اینفانی ذکر نموده ندا مرتفع لکن آذان قلیل
فارس موجود ولکن عرصه تنگ یومی حجابها خرق شود و مقامات مستوره ظاهر و
هویدا گردد

\end{otherlanguage}

Qasam bi-Āftāb-e Ḥaqīqat, agar jamīʿ-e ʿālam bi-davām-e mulk va malakūt
bi-shokr va ḥamd nāṭiq shavand, qābil-e yik ḥarf ke az Qalam-e Aʿlā jārī
mīshavad nabūde va nīstand. Agar maqām-e dhikr-e ilāhī ẓāhir shavad,
kull nāṭiq shavand bi-ānche īn fānī dhikr namūde. Nidā murtafiʿ lākin
ādhān qalīl. Fāris mawjūd va-lākin ʿarṣe tang. Yawmī ḥijābhā kharq
shavad va maqāmāt-e mastūre ẓāhir va huvaydā gardad.

I swear by the Sun of Truth, if the whole world were to voice thanks and
praise for as long as the earthly and heavenly kingdoms endure, they
would not be, nor have they ever been, worthy of a single letter flowing
from the Most Exalted Pen. If the station of divine remembrance were
revealed, everyone would voice what this mortal has said. The call is
raised, but ears are few. The rider is present, but the field is narrow.
One day the veils will be torn away and the hidden stations will become
manifest and plain.

\subsection{142}\label{section-3105}

\begin{otherlanguage}{persian}

و اینکه در بارهٴ جناب خان سرکار س‌ر علیه بهآء اللّه و عنایته مرقوم داشتند
و ذکر آنچه در خورجین آنجناب گذاردند فرمودند اینمراتب تلقاء وجه مالک قدم
عرض شد قوله جلّ جلاله

\end{otherlanguage}

Va īnke dar bāre-ye jināb-e Khān Sarkār Sīn Rā, ʿalayhi bahāʾu'llāh va
ʿināyatuhu, marqūm dāshtand va dhikr-e ānche dar khūrjīn-e ān jināb
gudhāshtand farmūdand, īn marātib tilqāʾ-e vajh-e Mālik-e Qidam ʿarḍ
shod. Qawluhu jalla jalāluhu.

What you wrote concerning Khān Sarkār Sīn Rā, upon him be the glory of
God and His loving care, and your mention of what was placed in his
saddlebag were submitted before the face of the Lord of Eternity. His
words---exalted is His glory:

\subsection{143}\label{section-3106}

\begin{otherlanguage}{arabic}

بسمی الشّاهد العلیم

\end{otherlanguage}

Bi-ismī al-Shāhid al-ʿAlīm.

In My name, the Witness, the All-Knowing.

\subsection{144}\label{section-3107}

\begin{otherlanguage}{persian}

یا امین انّا کنّا معه انّ ربّک یسمع و یری و هو العلیم الخبیر لازال بر خدمت
امر قیام نموده‌اند و فائز شدند بذکریکه بدوام ملک و ملکوت باقی و پاینده
است امروز از خریر ماء حیوان و حفیف سدرهٴ بیان اینکلمهٴ مبارکهٴ علیا اصغا
شد طوبی لمن فاز و وجد و عرف و قام علی خدمة امر اللّه ربّ العالمین و هر
شیئی از اشیاء شهادت داده بر وفای عهد و میثاق او اینمقام بسیار عظیم است
انّا نکبّر من هذا المقام علیه و علی من معه انّک لو رأیته ذکّره من قبلی و
بشّره بعنایتی و ما جری من قلمی و ما شهد به لسانی فی هذا المقام المنیع
البهآء علیه و علی کلّ صابر ثابت مستقیم انتهی

\end{otherlanguage}

Yā Amīn, innā kunnā maʿahu. Inna Rabbaka yasmaʿu wa-yarā wa-Huwa
al-ʿAlīm al-Khabīr. Lāzāl bar khidmat-e Amr qiyām namūde-and va fāʾiz
shodand bi-dhikrī ke bi-davām-e mulk va malakūt bāqī va pāyande ast.
Imrūz az kharīr-e māʾ-e ḥayavān va ḥafīf-e Sidre-ye Bayān īn kalime-ye
mubārake-ye ʿulyā iṣghā shod: Ṭūbā li-man fāza wa-wajada wa-ʿarafa
wa-qāma ʿalā khidmat Amr Allāh Rabb al-ʿālamīn. Va har shayʾī az ashyāʾ
shahādat dāde bar vafā-ye ʿahd va mīthāq-e ū. Īn maqām bisyār ʿaẓīm ast.
Innā nukabbiru min hādhā al-maqām ʿalayhi wa-ʿalā man maʿahu. Innaka law
raʾaytahu dhakkirhu min qibalī wa-bashshirhu bi-ʿināyatī wa-mā jarā min
qalamī wa-mā shahida bihi lisānī fī hādhā al-maqām al-manīʿ. Al-Bahāʾ
ʿalayhi wa-ʿalā kulli ṣābir thābit mustaqīm. Intahā.

O Amīn, We were with him. Your Lord hears and sees, and He is the
All-Knowing, the All-Informed. He has always arisen to serve the Cause
and has attained a remembrance that will remain and endure as long as
the earthly and heavenly kingdoms endure. Today, from the murmur of the
water of life and the rustling of the Tree of Utterance, these exalted,
blessed words were heard: ``Blessed is the one who has attained,
discovered, recognized, and arisen to serve the Cause of God, the Lord
of the worlds.'' Every created thing has testified to his faithfulness
to the pledge and covenant. This station is very great. From this
station We send him and those with him the greeting of God's greatness.
If you see him, remember him on My behalf and give him the glad news of
My loving care, of what has flowed from My Pen, and of what My tongue
has testified in this inaccessible station. Glory be upon him and upon
everyone who is patient, steadfast, and upright. End.

\subsection{145}\label{section-3108}

\begin{otherlanguage}{persian}

فی‌الحقیقه اینعبد را نفحات آیات در اینمقام از ذکر و فکر و خیال منقطع نمود
بشأنی آیات اللّه در این حین نازل و جاری که صد مثل اینفانی از عهدهٴ تقریر
و تحریر آن برنیامده و نخواهد آمد براستی عرض مینمایم و کفی باللّه شهیدا و
چه مقدار فرح حاصل شد از اینکه مشاهده شد لحاظ عنایت متوجّه بجناب خان علیه
بهآء اللّه است اینعبد خدمت ایشان سلام و تکبیر و ذکر میرساند و از حقّ جلّ
جلاله در کلّ حین تأیید میطلبد ایشان بمحبّت کل سزاوارند لو شآء اللّه یظهر ما
قدّر له انّه لهو المقتدر القدیر

\end{otherlanguage}

Fī al-ḥaqīqe īn ʿabd rā nafaḥāt-e āyāt dar īn maqām az dhikr va fikr va
khayāl munqaṭiʿ namūd. Bi-shaʾnī āyāt-e Allāh dar īn ḥīn nāzil va jārī
ke ṣad mithl-e īn fānī az ʿohde-ye taqrīr va taḥrīr-e ān barnayāmade va
nakhvāhad āmad. Bi-rāstī ʿarḍ mīnamāyam, wa-kafā bi-Allāh shahīdan. Va
che miqdār faraḥ ḥāṣil shod az īnke mushāhade shod liḥāẓ-e ʿināyat
mutavajjih bi-jināb-e Khān, ʿalayhi bahāʾu'llāh, ast. Īn ʿabd khidmat-e
īshān salām va takbīr va dhikr mīrasānad va az Ḥaqq jalla jalāluhu dar
kull-e ḥīn taʾyīd mīṭalabad. Īshān bi-maḥabbat-e kull sazāvārand. Law
shāʾa Allāh yuẓhiru mā quddira lahu. Innahu la-Huwa al-Muqtadir
al-Qadīr.

Truly, at this point the breezes of the verses carried this servant
beyond mention, thought, and imagination. The verses of God are being
sent down and flowing forth at this moment in such a manner that a
hundred like this mortal have not been, and never will be, capable of
recounting or recording them. I speak truthfully, and God is sufficient
as a witness. And how much joy arose on seeing that the gaze of loving
care was directed toward Khān, upon him be the glory of God. This
servant conveys to him greetings, the salutation of God's greatness, and
remembrance, and at every moment asks God---exalted is His glory---to
assist him. He is worthy of everyone's love. Should God will, He will
reveal what has been ordained for him. He is indeed the All-Powerful,
the Omnipotent.

\subsection{146}\label{section-3109}

\begin{otherlanguage}{persian}

و اینکه ذکر ورود در سیرجان و ملاقات با جناب حاجی سیّد ح‌س علیه بهآء اللّه
وحده نمودند اینمراتب تلقاء وجه مالک اسماء بشرف اصغا فائز قوله جلّ جلاله
باید در قری بعضی از دوستان مرور نمایند و بکمال حکمت بتهذیب نفوس غافله
مشغول شوند که شاید در ایّام اللّه محروم نمانند امروز فضل اعظم ظاهر بر هر
نفسی لازم که بهدایت نفوس مشغول شود

\end{otherlanguage}

Va īnke dhikr-e vurūd dar Sīrjān va mulāqāt bā jināb-e Ḥājī Siyyid Ḥā
Sīn, ʿalayhi bahāʾu'llāh vaḥdahu, namūdand, īn marātib tilqāʾ-e vajh-e
Mālik-e Asmā bi-sharaf-e iṣghā fāʾiz. Qawluhu jalla jalāluhu: Bāyad dar
qurā baʿḍī az dūstān murūr namāyand va bi-kamāl-e ḥikmat bi-tahdhīb-e
nufūs-e ghāfile mashghūl shavand ke shāyad dar ayyām-e Allāh maḥrūm
namānand. Imrūz faḍl-e aʿẓam ẓāhir. Bar har nafsī lāzim ke bi-hidāyat-e
nufūs mashghūl shavad.

Your mention of arriving in Sīrjān and meeting Ḥājī Siyyid Ḥā Sīn, upon
him be the glory of God alone, attained the honor of being heard before
the face of the Lord of Names. His words---exalted is His glory: Some of
the friends should travel through the villages and, with the utmost
wisdom, engage in refining heedless souls, so that perhaps they may not
remain deprived in the days of God. Today the greatest grace is
manifest. Everyone must engage in guiding souls.

\subsection{147}\label{section-3110}

\begin{otherlanguage}{arabic}

طوبی لأفنانی الّذی کان مذکوراً فی ساحتی و فائزاً بلقائی اذ کان النّور مشرقاً
من افق الزّورآء چه که جناب حیدر قبل علی علیه بهائی و عنایتی ذکر تبلیغ
وکالتی خود را نمودند و همچنین زیارت اینفقره فی‌الحقیقه منصور و بقبول حقّ
فائز انتهی

\end{otherlanguage}

Ṭūbā li-Afnānī alladhī kāna madhkūran fī sāḥatī wa-fāʾizan bi-liqāʾī idh
kāna al-nūr mushriqan min ufuq al-Zawrāʾ. Che ke jināb-e Ḥaydar qabl
ʿAlī, ʿalayhi bahāʾī va ʿināyatī, dhikr-e tablīgh-e vikālatī-ye khod rā
namūdand va hamchunīn ziyārat. Īn fiqre fī al-ḥaqīqe manṣūr va
bi-qabūl-e Ḥaqq fāʾiz. Intahā.

Blessed is My Afnān, who was remembered in My court and attained My
presence when the light was shining from the horizon of Zawrāʾ. For
Ḥaydar-ʿAlī, upon him be My glory and loving care, mentioned his
teaching by proxy and likewise the visit. This undertaking is truly
assisted and has attained God's acceptance. End.

\subsection{148}\label{section-3111}

\begin{otherlanguage}{persian}

جناب حاجی مذکور بذکر حقّ جلّ جلاله فائز گشتند از قبل فانی بایشان تکبیر
برسانید باید در صدد آن باشند که کلمة اللّه در آن ارض انتشار یابد چه بسیار
از نفوس که حرفی از علوم ظاهره را نشنیده‌اند و باشتعال نار محبّت لسان گشوده
و بترویج امر و تهذیب نفوس مؤیّد گشتند امروز روح القدس تأیید میفرماید باید
جمیع دوستان با جمیع عالم بمحبّت و دوستی رفتار نمایند در اینصورت محبّت
احداث شود و قلوب مستعدّ آنحین حین القاء کلمه است ولکن بروح و ریحان الأمر
بید اللّه ربّنا الرّحمن

\end{otherlanguage}

Janāb-i Ḥājī-yi madhkūr bi-dhikr-i Ḥaqq jalla jalāluhu fāʾiz gashtand.
Az qibal-i fānī bi-īshān takbīr birasānīd. Bāyad dar ṣadad-i ān bāshand
kih Kalimat Allāh dar ān arḍ intishār yābad. Chih bisyār az nufūs kih
ḥarfī az ʿulūm-i ẓāhirih rā nashunīdih-and va bi-ishtiʿāl-i nār-i
maḥabbat lisān gushūdih va bi-tarvīj-i Amr va tahdhīb-i nufūs muʾayyad
gashtand. Imrūz Rūḥ al-Qudus taʾyīd mīfarmāyad. Bāyad jamīʿ-i dūstān bā
jamīʿ-i ʿālam bi-maḥabbat va dūstī raftār namāyand. Dar īn ṣūrat
maḥabbat iḥdāth shavad va qulūb mustaʿidd. Ānḥīn ḥīn-i ilqā-yi kalimih
ast, valākin bi-rūḥ va rayḥān. Al-amru bi-yadi Allāhi Rabbinā al-Raḥmān.

The aforementioned Ḥājī was favored with mention by God, exalted be His
glory. Convey greetings of praise to him on behalf of this lowly one. He
should endeavor to spread the Word of God in that land. How many souls
have never heard a word of outward learning, yet, set ablaze by the fire
of love, have opened their lips and been aided in promoting the Cause
and refining souls. Today the Holy Spirit grants assistance. All the
friends must treat all the world with love and friendship. In this way
love will arise and hearts will become receptive. That is the time to
impart the Word, but with joy and gladness. The matter rests in the
hands of God, our Lord, the All-Merciful.

\subsection{149}\label{section-3112}

\begin{otherlanguage}{persian}

ذکر ارض ش و اولیای حقّ علیهم بهآء اللّه را نموده بودند هر یک بعنایت الهی
فائز گشتند و مخصوص جناب حاجی شکراللّه علیه بهآء اللّه لوح امنع اقدس نازل و
بجناب محبوبی آقا میرزا حسین علیه بهآء اللّه و عنایته داده شد که برسانند
این حین لسان قدم باینکلمهٴ مبارکه ناطق قوله جلّ جلاله

\end{otherlanguage}

Dhikr-i arḍ-i Shīn va awliyā-yi Ḥaqq ʿalayhim Bahāʾ Allāh rā namūdih
būdand. Har yik bi-ʿināyat-i ilāhī fāʾiz gashtand va makhṣūṣ-i Janāb-i
Ḥājī Shukrullāh ʿalayhi Bahāʾ Allāh lawḥ-i amnaʿ-i aqdas nāzil va
bi-Janāb-i maḥbūbī Āqā Mīrzā Ḥusayn ʿalayhi Bahāʾ Allāh va ʿināyatuhu
dādih shud kih birasānand. Īn ḥīn Lisān-i Qidam bi-īn kalimih-yi
mubārakih nāṭiq, qawluhu jalla jalāluhu:

You had mentioned the land of Shīn and the friends of God, upon them be
the glory of God. Each was favored with divine loving care, and a most
exalted, most holy Tablet was revealed especially for Ḥājī Shukrullāh,
upon him be the glory of God, and given to my beloved Āqā Mīrzā Ḥusayn,
upon him be the glory of God and His loving care, to deliver. At this
moment the Tongue of Eternity speaks these blessed words---His words,
exalted be His glory:

\subsection{150}\label{section-3113}

\begin{otherlanguage}{persian}

یا اباالحسن یا امین ذکر دوستان الهی را نمودی للّه الحمد مؤیّد شدی بفضل و
عنایت حقّ از بحر وصال مرّة بعد مرّه آشامیدی از کوثر بیان رحمن نوشیدی و افق
اعلی را مشاهده نمودی و بذکر خیر اولیا و اصفیا و دوستان حقّ مؤیّد گشتی از
حقّ میطلبیم لازال ترا مدد فرماید و بر آنچه سبب اعزاز امر او جلّ جلاله است
تأیید نماید لعمر اللّه اگر نفسی کنوز عالم را ایثار نماید و یا جمیع ارواح
فدا شود و ذرّه‌ئی از عزّت امر اللّه بکاهد جایز نبوده و نیست در جمیع احوال
باعزاز امر اللّه و اعلاء کلمة اللّه و ارتفاع مقام ناظر باش هذا ما وصّیناک
به من قبل و فی الحضور و فی هذا اللّیل الّذی ینطق فیه لسان العظمة الملک للّه
ربّ العالمین انتهی

\end{otherlanguage}

Yā Abā al-Ḥasan, yā Amīn! Dhikr-i dūstān-i ilāhī rā namūdī. Li-llāhi
al-ḥamd, muʾayyad shudī. Bi-faḍl va ʿināyat-i Ḥaqq az baḥr-i viṣāl
marratan baʿda marrah āshāmīdī, az Kawthar-i bayān-i Raḥmān nūshīdī va
ufuq-i aʿlā rā mushāhidih namūdī va bi-dhikr-i khayr-i awliyā va aṣfiyā
va dūstān-i Ḥaqq muʾayyad gashtī. Az Ḥaqq mīṭalabīm lā yazāl turā madad
farmāyad va bar ānchih sabab-i iʿzāz-i Amr-i ū jalla jalāluhu ast taʾyīd
namāyad. La-ʿamru Allāh, agar nafsī kunūz-i ʿālam rā īthār namāyad va yā
jamīʿ-i arvāḥ fidā shavad va dharrih-ī az ʿizzat-i Amr Allāh bikāhad
jāʾiz nabūdih va nīst. Dar jamīʿ-i aḥvāl bi-iʿzāz-i Amr Allāh va iʿlā-yi
Kalimat Allāh va irtifāʿ-i maqām nāẓir bāsh. Hādhā mā waṣṣaynāka bihi
min qablu wa-fī al-ḥuḍūr wa-fī hādhā al-layli alladhī yanṭiqu fīhi
Lisānu al-ʿAẓamah: al-mulku li-llāhi Rabbi al-ʿālamīn. Intahā.

``O Abū al-Ḥasan! O Amīn! You have mentioned the friends of God. Praise
be to God, you have been aided. Through the grace and loving care of
God, you have drunk time and again from the ocean of reunion, tasted the
Kawthar of the All-Merciful's utterance, beheld the Most Exalted
Horizon, and been aided in speaking well of the faithful, the chosen
ones, and the friends of God. We ask God to assist you always and to
confirm you in whatever serves to honor His Cause, exalted be His glory.
By the life of God! If someone were to give away the treasures of the
world, or all souls were to be sacrificed, and this diminished the honor
of the Cause of God by even an atom, it would not be permissible, nor
has it ever been. In all circumstances, keep your gaze fixed on honoring
the Cause of God, exalting the Word of God, and elevating its station.
This is what We have enjoined upon you before, in Our presence, and on
this night in which the Tongue of Grandeur proclaims: Sovereignty
belongs to God, the Lord of the worlds.'' End.

\subsection{151}\label{section-3114}

\begin{otherlanguage}{arabic}

اینفانی خدمت آنمحبوب و دوستان الهی صدهزار سلام و تکبیر میرساند و از حقّ
جلّ جلاله میطلبد کل را بر نصرت تأیید فرماید و توفیق عطا نماید انّه علی
کلّشیء قدیر البهآء و التّکبیر و الثّنآء علی جنابکم و علی من معکم و علی من
یسمع قولکم فی امر اللّه ربّ العالمین

\end{otherlanguage}

Īn fānī khidmat-i ān maḥbūb va dūstān-i ilāhī ṣad hizār salām va takbīr
mīrasānad va az Ḥaqq jalla jalāluhu mīṭalabad kull rā bar nuṣrat taʾyīd
farmāyad va tawfīq ʿaṭā namāyad. Innahu ʿalā kulli shayʾin qadīr.
Al-Bahāʾ wa-al-takbīr wa-al-thanāʾ ʿalā janābikum wa-ʿalā man maʿakum
wa-ʿalā man yasmaʿu qawlakum fī Amri Allāhi Rabbi al-ʿālamīn.

This lowly one sends a hundred thousand greetings of peace and praise to
you, my beloved, and to the friends of God, and asks God, exalted be His
glory, to aid everyone in supporting His Cause and to grant them
success. He is indeed powerful over all things. Glory, exaltation, and
praise rest upon you, upon those who are with you, and upon those who
listen to your words concerning the Cause of God, the Lord of the
worlds.

\subsection{152}\label{section-3115}

\begin{otherlanguage}{persian}

فقرهٴ دیگر مکرّر از حبیب روحانی جناب آقا محمّد طاهر علیه بهآء اللّه و بعضی
دوستان رسید آنچه که سبب و علّت سرور فانی بود چه که بذکر حقّ جلّ جلاله ناطق
و باقبال و خضوع و خشوع و قیام بر خدمت لوجه اللّه ربّنا و ربّ من فی السّموات
و الأرض مزیّن خدمت کل تکبیر و سلام میرسانم و امید است که انشآءاللّه کل
بعنایت مخصوصه فائز شوند

\end{otherlanguage}

Faqarih-yi dīgar, mukarrar az ḥabīb-i rūḥānī Janāb-i Āqā Muḥammad Ṭāhir
ʿalayhi Bahāʾ Allāh va baʿḍī dūstān rasīd ānchih kih sabab va ʿillat-i
surūr-i fānī būd, chih kih bi-dhikr-i Ḥaqq jalla jalāluhu nāṭiq va
bi-iqbāl va khuḍūʿ va khushūʿ va qiyām bar khidmat li-wajhi Allāhi
Rabbinā wa-Rabbi man fī al-samāwāti wa-al-arḍ muzayyan. Khidmat-i kull
takbīr va salām mīrasānam va umīd ast kih in shāʾ Allāh kull
bi-ʿināyat-i makhṣūṣih fāʾiz shavand.

On another matter, this lowly one has repeatedly received from the
spiritual friend Āqā Muḥammad Ṭāhir, upon him be the glory of God, and
some of the friends, that which has brought him joy, for they speak in
remembrance of God, exalted be His glory, and are adorned with devotion,
humility, reverence, and readiness to serve for the sake of God, our
Lord and the Lord of all who are in the heavens and on earth. I convey
greetings of praise and peace to everyone, and hope that, God willing,
all may receive His special loving care.

\subsection{153}\label{section-3116}

\begin{otherlanguage}{persian}

والدهٴ جناب آقا محمّد طاهر علیهما بهآء اللّه در این ایّام که قصر مقرّ عرش
واقع بشرف حضور فائزند یفعل ما یشآء و یحکم ما یرید و هو المقتدر القدیر
ملکه‌های عالم بعد از طلب و انتظار فائز نشدند و این مخدّره از بدایع فضل
الهی فائز شد بآنچه که اکثر اهل عالم در حسرتش مردند

\end{otherlanguage}

Vālidih-yi Janāb-i Āqā Muḥammad Ṭāhir ʿalayhimā Bahāʾ Allāh dar īn ayyām
kih qaṣr maqarr-i ʿarsh vāqiʿ bi-sharaf-i ḥuḍūr fāʾizand. Yafʿalu mā
yashāʾu wa-yaḥkumu mā yurīdu wa-huwa al-Muqtadiru al-Qadīr.
Malikih-hā-yi ʿālam baʿd az ṭalab va intiẓār fāʾiz nashudand va īn
mukhaddarih az badāyiʿ-i faḍl-i ilāhī fāʾiz shud bi-ānchih kih akthar-i
ahl-i ʿālam dar ḥasratash murdand.

The mother of Āqā Muḥammad Ṭāhir, upon them both be the glory of God,
has attained the honor of His presence in these days when the palace is
the seat of the Throne. He does what He wills and ordains what He
desires, and He is the All-Powerful, the Almighty. The queens of the
world, despite seeking and waiting, did not attain it, while this
honored lady, through the wondrous gifts of divine grace, has attained
what most of the people of the world died longing for.

\subsection{154}\label{section-3117}

\begin{otherlanguage}{persian}

و همچنین این ایّام مکتوبی از جناب حاجی آقا محمّد (علاقه) علیه بهآء اللّه
رسید در ساحت امنع اقدس عرض شد بذکر حقّ جلّ جلاله فائز گشتند هنیئاً له و
امید هست که از بعد هم ارسال شود آنچه که بمثابهٴ روح معانی است از برای
هیاکل مخلصین الحمد للّه ربّ العالمین و البهآء للمقرّبین و المخلصین

\end{otherlanguage}

Va hamchunīn īn ayyām maktūbī az Janāb-i Ḥājī Āqā Muḥammad (ʿAlāqih)
ʿalayhi Bahāʾ Allāh rasīd. Dar sāḥat-i amnaʿ-i aqdas ʿarḍ shud.
Bi-dhikr-i Ḥaqq jalla jalāluhu fāʾiz gashtand. Hanīʾan lahu. Va umīd
hast kih az baʿd ham irsāl shavad ānchih kih bi-mathābih-yi rūḥ-i maʿānī
ast az barā-yi hayākil-i mukhliṣīn. Al-ḥamdu li-llāhi Rabbi al-ʿālamīn
wa-al-Bahāʾ li-l-muqarrabīn wa-al-mukhliṣīn.

Likewise, in these days a letter arrived from Ḥājī Āqā Muḥammad
(ʿAlāqih), upon him be the glory of God. It was presented in the most
exalted, most holy court, and he was favored with mention by God,
exalted be His glory. Blessed is he! It is hoped that in the future,
too, he will send that which is like the spirit of inner meanings to the
bodies of the sincere. Praise be to God, the Lord of the worlds, and
glory rest upon those who are near to Him and the sincere.

\subsection{155}\label{section-3118}

\begin{otherlanguage}{persian}

الواح منزله حسب الأمر علیحده هم بخطّ اینفانی جز یک لوح امنع اقدس مخصوص
جناب طبیب نوشته ارسال شد اگر میشد در آنجا شخص صاحب خطّی الواح مرسوله را
بنویسد و در مطابقه و مقابلهٴ آن منتهی جهد را مبذول دارد بسیار محبوبست و
همچنین در محلّ امنی حفظ نمایند چه که آیات الهی در اطراف متفرّق است لذا آن
ارض و هر ارضی آنچه در او وارد شود جمع نمایند محبوب و مقبول است

\end{otherlanguage}

Alvāḥ-i munazzalih ḥasb al-amr ʿalāḥidih ham bi-khaṭṭ-i īn fānī juz yik
lawḥ-i amnaʿ-i aqdas makhṣūṣ-i janāb-i ṭabīb nivishtih irsāl shud. Agar
mīshud dar ānjā shakhṣ-i ṣāḥib-khaṭṭī alvāḥ-i marsūlih rā binivīsad va
dar muṭābiqih va muqābilih-yi ān muntahā-yi jahd rā mabdhūl dārad bisyār
maḥbūb ast. Va hamchunīn dar maḥall-i amnī ḥifẓ namāyand, chih kih
āyāt-i ilāhī dar aṭrāf mutafarriq ast. Lidhā ān arḍ va har arḍī ānchih
dar ū vārid shavad jamʿ namāyand maḥbūb va maqbūl ast.

As commanded, the revealed Tablets were also copied out separately and
sent, in this lowly one's handwriting except for one most exalted, most
holy Tablet intended for the physician. It would be most welcome if
someone with a good hand there could copy the Tablets sent and exercise
the utmost care in checking and collating them. They should also be kept
in a safe place, for the divine verses are scattered in various places.
It would therefore be welcome and acceptable for whatever arrives in
that land, and in every land, to be collected.

\subsection{156}\label{section-3119}

\begin{otherlanguage}{arabic}

خ‌ادم فی٢ جمادی‌الأولی سنة {[}١٣٠٢{]}

\end{otherlanguage}

Khādim, fī 2 Jumādā al-Ūlā sanat {[}1302{]}.

The servant, on 2 Jumādā al-Ūlā, in the year {[}1302{]}.

\bookmarksetup{startatroot}

\chapter{Lawh-i-Istintaq (Tablet of the
Interrogation)}\label{lawh-i-istintaq-tablet-of-the-interrogation}

\section{BH00015}\label{bh00015}

\emph{Translated on 17 Qudrat 183 using GPT 6.1 Sol}

\subsection{1}\label{section-3120}

\begin{otherlanguage}{arabic}

بسم ربّنا الأمنع الأقدس الأعزّ الأبهی

\end{otherlanguage}

Bismi rabbinā al-Amnaʿ al-Aqdas al-Aʿazz al-Abhā.

In the name of our Lord, the Most Inaccessible, the Most Holy, the Most
Mighty, the All-Glorious.

\subsection{2}\label{section-3121}

\begin{otherlanguage}{persian}

سبحانک یا من بک حقّق الذّات فی الذّات و تردّی کینونة القدم بردآء الأسمآء و
الصّفات من اسمک الظّاهر ظهرت الظّهورات من اوّل الّذی لا اوّل له و باسمک الباطن
استبطنت البطونات فی غیب الغیب بحیث ما ظهر من قدم القدم و لن یظهر الی آخر
الّذی لا آخر له بطلوع شمس جمالک الأبهی طلع نیّر العلآء من افق البدآء مبشّراً
لمن فی لاهوت العمآء و ملإ البقآء و اهل عوالم الّذی لا یذکر بقلم الانشآء و
لا یدرکها حقایق اولی النّهی ثمّ من فی ناسوت الأدنی بأنّ هذا لجمال الأبهی و
الطّلعة الأسنی قد ظهر بما هو مقدّس من ذکری و اشارتی و عرفانی و دلالتی و
تقدیسی و تسبیحی و تنزیهی و تمجیدی و توحیدی و تفریدی ایّاه و ما هو فی جوهر
جوهری و ساذج ساذجی لدی ذکره و بهائه

\end{otherlanguage}

Subḥānaka yā man bika ḥuqqiqa al-dhāt fī al-dhāt wa-taraddā kaynūnat
al-qidam bi-ridāʾ al-asmāʾ wa-al-ṣifāt. Min ismika al-Ẓāhir ẓaharat
al-ẓuhūrāt min awwal alladhī lā awwala lahu wa-bi-ismika al-Bāṭin
istabṭanat al-buṭūnāt fī ghayb al-ghayb bi-ḥaythu mā ẓahara min qidam
al-qidam wa-lan yaẓhara ilā ākhir alladhī lā ākhira lahu. Bi-ṭulūʿ shams
jamālika al-abhā ṭalaʿa nayyir al-ʿalāʾ min ufuq al-badāʾ mubashshiran
li-man fī lāhūt al-ʿamāʾ wa-malaʾ al-baqāʾ wa-ahl ʿawālim alladhī lā
yudhkar bi-qalam al-inshāʾ wa-lā tudrikuhā ḥaqāʾiq ulī al-nuhā thumma
man fī nāsūt al-adnā bi-anna hādhā la-jamāl al-Abhā wa-al-ṭalʿa al-asnā
qad ẓahara bi-mā huwa muqaddas min dhikrī wa-ishāratī wa-ʿirfānī
wa-dalālatī wa-taqdīsī wa-tasbīḥī wa-tanzīhī wa-tamjīdī wa-tawḥīdī
wa-tafrīdī iyyāhu wa-mā huwa fī jawhar jawharī wa-sādhij sādhijī ladā
dhikrihi wa-bahāʾihi.

Glory be to You, O You through whom the Essence was established in the
Essence, and the being of eternity clothed itself in the robe of names
and attributes. Through Your name, the Manifest, manifestations have
appeared from the beginning that has no beginning; and through Your
name, the Hidden, hidden realities have lain concealed in the unseen of
the unseen, so that they have not appeared from the eternity of
eternities and will never appear until the end that has no end. With the
rising of the sun of Your all-glorious beauty, the Luminary of
exaltation rose from the horizon of origination, bringing glad tidings
to those in the divine realm of the cloud of unknowing, to the Concourse
of eternity, to the inhabitants of worlds that the pen of creation
cannot describe and the realities of those endowed with understanding
cannot comprehend, and then to those in the lowly human realm: ``This is
indeed the All-Glorious Beauty and the most resplendent Countenance. He
has appeared in a manner sanctified beyond my mention and allusion, my
recognition and indication, my hallowing and praise, my declaration of
His transcendence and glorification, my affirmation of His unity and
uniqueness, and whatever lies in the essence of my essence and the
purest core of my being before His remembrance and splendor.''

\subsection{3}\label{section-3122}

\begin{otherlanguage}{arabic}

اذاً یا من خلقتم فی البیان بآیاتی لعرفانه لا تحتجبوا عن الّذی کان البیان
کلّه کحلقة خاتم فی اصبع ارادته و ورقة لرضوان امره فسبحانک یا من لن یقدر
ان یشیر بذکرک ذکر الأعظم و لن یقدر ان یجری علی ثنائک قلم القدم فکیف یقدر
و یستطیع ان یجترح بذکرک هذا الفانی الّذی یکون احقر الخدّام لدی باب فضلک
الّذی فتح علی الأمم الّا بترشّح من ابحر جودک الّذی تموّج علی کلّ الموجودات و
طفح من طمطام عنایتک الّتی تلجلجت علی من فی ملکوت الأوّلیّات و الآخریّات حیث
فتحت السن کلّ الأشیآء بأبدع الأذکار بثنآء نفسک المختار و ارتفع النّدآء من
قلب صخرة الصّمّآء فی قعر البحار و ما فی علی الأطوار بأنّک انت الله ربّ ما
یری و ما لا یری

\end{otherlanguage}

Idhan yā man khuliqtum fī al-Bayān bi-āyātī li-ʿirfānihi lā taḥtajibū
ʿan alladhī kāna al-Bayān kulluhu ka-ḥalqat khātam fī iṣbaʿ irādatihi
wa-waraqa li-riḍwān amrihi. Fa-subḥānaka yā man lan yaqdir an yushīra
bi-dhikrika dhikr al-aʿẓam wa-lan yaqdir an yajriya ʿalā thanāʾika qalam
al-qidam. Fa-kayfa yaqdir wa-yastaṭīʿ an yajtaraḥa bi-dhikrika hādhā
al-fānī alladhī yakūn aḥqar al-khuddām ladā bāb faḍlika alladhī futiḥa
ʿalā al-umam illā bi-tarashshuḥ min abḥur jūdika alladhī tamawwaja ʿalā
kull al-mawjūdāt wa-ṭafḥ min ṭamṭām ʿināyatika allatī talajlajat ʿalā
man fī malakūt al-awwaliyyāt wa-al-ākhiriyyāt ḥaythu futiḥat alsun kull
al-ashyāʾ bi-abdaʿ al-adhkār bi-thanāʾ nafsika al-mukhtār wa-irtafaʿa
al-nidāʾ min qalb ṣakhrat al-ṣammāʾ fī qaʿr al-biḥār wa-mā fī ʿulā
al-aṭwār bi-annaka anta Allāh rabb mā yurā wa-mā lā yurā.

Then, O you who were created in the Bayān through My verses to recognize
Him, do not be veiled from the One in whose finger of will the whole
Bayān is but a ring, and but a leaf in the paradise of His Cause. Glory
be to You, O You whom the Most Great Remembrance can never indicate by
mentioning You, and in whose praise the Pen of eternity can never move.
How, then, could this perishable one, the lowliest of servants at the
door of Your grace that has been opened to the nations, dare to mention
You, except through a sprinkling from the oceans of Your generosity that
have surged over all created things, and an overflow from the vast ocean
of Your providence that has surged over those in the kingdom of the
first and the last? Thus the tongues of all things have been opened with
the most wondrous words of remembrance in praise of Your sovereign Self,
and the call has risen from the heart of the solid rock in the depths of
the seas and from all that is on the heights of the mountains: ``You are
indeed God, the Lord of all that is seen and unseen.''

\subsection{4}\label{section-3123}

\begin{otherlanguage}{arabic}

اذاً یا الهی اسألک بهذه العنایة الّتی احاطت کلّ الأشیآء و هذه الرّحمة الّتی
سبقت من فی الانشآء بأن ترسل حینئذ علی هذا الفانی نسمة من نسمات جود
وهّابیّتک و نفحة من نفحات بدعک الّتی بها اهتزّت عظام الرّمیمة و قامت اجساد
المیّتة بأنّک انت الله الباقی الدّائم القدیم لأقوم بها علی ثنآء احبّتک و
اذکر الّذین جعلتهم انوار الهدایة بین خلقک و مشاعل التّوحید بین عبادک و
بریّتک و انطق بما تؤیّدنی و تلهمنی من عندک لأنّی لا اعلم شیئاً و انّک انت
العلیم الخبیر

\end{otherlanguage}

Idhan yā ilāhī asʾaluka bi-hādhihi al-ʿināya allatī aḥāṭat kull
al-ashyāʾ wa-hādhihi al-raḥma allatī sabaqat man fī al-inshāʾ bi-an
tursila ḥīnaʾidhin ʿalā hādhā al-fānī nasma min nasamāt jūd
wahhābiyyatika wa-nafḥa min nafaḥāt bidʿika allatī bihā ihtazzat ʿiẓām
al-ramīma wa-qāmat ajsād al-mayyita bi-annaka anta Allāh al-Bāqī
al-Dāʾim al-Qadīm li-aqūma bihā ʿalā thanāʾ aḥibbatika wa-adhkura
alladhīna jaʿaltahum anwār al-hidāya bayna khalqika wa-mashāʿil
al-tawḥīd bayna ʿibādika wa-bariyyatika wa-anṭiqa bi-mā tuʾayyidunī
wa-tulhimunī min ʿindika li-annī lā aʿlam shayʾan wa-innaka anta
al-ʿAlīm al-Khabīr.

Then, my God, I ask You by this providence that has encompassed all
things and this mercy that has preceded all creation to send down now
upon this perishable one a breeze from the breezes of Your bountiful
generosity and a breath from the breaths of Your wondrous creation,
through which decayed bones have stirred and dead bodies have arisen,
for You are indeed God, the Everlasting, the Abiding, the Ancient. May I
thereby arise to praise Your loved ones, remember those whom You have
made lights of guidance among Your creatures and torches of unity among
Your servants and Your creation, and speak as You strengthen and inspire
me from Your presence. For I know nothing, and You are indeed the
All-Knowing, the All-Aware.

\subsection{5}\label{section-3124}

\begin{otherlanguage}{arabic}

و انت تعلم یا الهی انّ حبّی ایّاهم و ذکری لدیهم یکون خالصاً لوجهک و مطهّراً عن
مشاهدة غیرک انّی احبّ یا الهی انّهم یطّلعون بمواقع الأمور فی ایّامک حقّ
الاطّلاع و یعرفون ما هو المستور عنهم حقّ العرفان لیبقی بذلک کینوناتهم و
حقایقهم و ارواحهم و امثالهم فی کلّ عالم من عوالمک عارفاً بحقّک و مطّلعاً بما
ظهر من عندک و ما اردت لهم فی سرادق عنایتک و فضلک لأنّهم یا الهی لو یطّلعون
بکلّ الأمور و یحتجبون عن شیء احتجبوا بقدر ذلک عن منظرک الأکبر و ما توجّهوا
الیک بتمام البصر و هذا لم یکن الّا بتقدیسه عن کلّ ما خلقت و تخلق و ظهرت او
تظهر حیث نطق بذلک لسانک الحقّ المبین لو یکون البصر علی کبر السّموات و
الأرض و یتوجّه بطرف طرفه اقلّ عمّا یحصی الی جهة اخری لم یقدر أن ینظر الی
هذا المنظر الأعلی علی ما ینبغی له من حقّ النّظر

\end{otherlanguage}

Wa-anta taʿlam yā ilāhī anna ḥubbī iyyāhum wa-dhikrī ladayhim yakūn
khāliṣan li-wajhika wa-muṭahharan ʿan mushāhadat ghayrika. Innī uḥibb yā
ilāhī annahum yaṭṭaliʿūn bi-mawāqiʿ al-umūr fī ayyāmika ḥaqq al-iṭṭilāʿ
wa-yaʿrifūn mā huwa al-mastūr ʿanhum ḥaqq al-ʿirfān li-yabqā bi-dhālika
kaynūnātuhum wa-ḥaqāʾiquhum wa-arwāḥuhum wa-amthāluhum fī kull ʿālam min
ʿawālimika ʿārifan bi-ḥaqqika wa-muṭṭaliʿan bi-mā ẓahara min ʿindika
wa-mā aradta lahum fī surādiq ʿināyatika wa-faḍlika li-annahum yā ilāhī
law yaṭṭaliʿūn bi-kull al-umūr wa-yaḥtajibūn ʿan shayʾ iḥtajabū bi-qadr
dhālika ʿan manẓarika al-akbar wa-mā tawajjahū ilayka bi-tamām al-baṣar.
Wa-hādhā lam yakun illā bi-taqdīsihi ʿan kull mā khalaqta wa-takhluq
wa-ẓaharta aw tuẓhir ḥaythu naṭaqa bi-dhālika lisānuka al-ḥaqq al-mubīn:
law yakūn al-baṣar ʿalā kibar al-samāwāt wa-al-arḍ wa-yatawajjahu
bi-ṭarf ṭarfihi aqall ʿammā yuḥṣā ilā jiha ukhrā lam yaqdir an yanẓura
ilā hādhā al-manẓar al-aʿlā ʿalā mā yanbaghī lahu min ḥaqq al-naẓar.

You know, my God, that my love for them and my mention in their presence
are solely for Your sake and purified of any regard for anyone but You.
I desire, my God, that they become fully aware of the true circumstances
of affairs in Your days and truly recognize what is concealed from them,
so that their beings, their realities, their spirits, and their
likenesses may thereby remain, in every world of Your worlds, aware of
Your truth and acquainted with what has appeared from You and what You
have intended for them within the tabernacle of Your providence and
grace. For, my God, if they were aware of all things but remained veiled
from one thing, to that extent they would be veiled from Your Most Great
Countenance and would not have turned toward You with the fullness of
their sight. This can come about only by sanctifying that sight from all
that You have created and will create, and have manifested or will
manifest. Thus Your true and lucid Tongue has declared: ``Were the eye
as vast as the heavens and the earth, and were it to turn even the tip
of its gaze toward any other direction by an amount too small to be
reckoned, it would be unable to behold this Most Exalted Countenance
with the true vision that befits it.''

\subsection{6}\label{section-3125}

\begin{otherlanguage}{arabic}

تعالی تعالی امرک یا الهی و تباهی تباهی ظهورک یا محبوبی قد صعب علی هذا
العبد فوق کلّ صعوبة بأن یحتجب احد فی مثل تلک الأیّام الّتی ما رأت عیون
الابداع شبهها و لا ابصر من فی الاختراع مثلها حیث ظهرت فیها بصرف جمالک و
کافور طلعتک و ساذج محبوبیّتک و احاطت السّموات و الأرض ظهورات عنایتک و
بروزات الطافک بحیث طلع کلّ امر مکنون و کنز مخزون و رمز مصون و اشرقت شمس
المعانی عن افق سمآء التّبیان و تجلّت بظهورات اسمائک الحسنی و صفاتک العلیا
علی من فی الامکان و الأکوان

\end{otherlanguage}

Taʿālā taʿālā amruka yā ilāhī wa-tabāhā tabāhā ẓuhūruka yā maḥbūbī. Qad
ṣaʿuba ʿalā hādhā al-ʿabd fawqa kull ṣuʿūba bi-an yaḥtajiba aḥad fī
mithl tilka al-ayyām allatī mā raʾat ʿuyūn al-ibdāʿ shibhahā wa-lā
abṣara man fī al-ikhtirāʿ mithlahā ḥaythu ẓaharta fīhā bi-ṣirf jamālika
wa-kāfūr ṭalʿatika wa-sādhij maḥbūbiyyatika wa-aḥāṭat al-samāwāt
wa-al-arḍ ẓuhūrāt ʿināyatika wa-burūzāt alṭāfika bi-ḥaythu ṭalaʿa kull
amr maknūn wa-kanz makhzūn wa-ramz maṣūn wa-ashraqat shams al-maʿānī ʿan
ufuq samāʾ al-tibyān wa-tajallat bi-ẓuhūrāt asmāʾika al-ḥusnā
wa-ṣifātika al-ʿulyā ʿalā man fī al-imkān wa-al-akwān.

Exalted, exalted is Your Cause, my God, and glorious, glorious is Your
Revelation, my Beloved! Beyond all hardship, it has been grievous to
this servant that anyone should remain veiled in days such as these,
whose like the eyes of creation have never seen, nor have those brought
into being ever beheld. In these days You have appeared in the purity of
Your beauty, the camphor-white radiance of Your countenance, and the
unalloyed essence of Your belovedness. The manifestations of Your
providence and the expressions of Your loving-kindness have encompassed
the heavens and the earth, so that every hidden matter, every stored-up
treasure, and every guarded mystery has come to light. The Sun of inner
meanings has shone from the horizon of the heaven of exposition and shed
its splendor, through the manifestations of Your most beautiful names
and most exalted attributes, upon all who dwell in the realms of
possibility and existence.

\subsection{7}\label{section-3126}

\begin{otherlanguage}{arabic}

و الّذین یریدون ان یدخلوا هذا البساط الممتنع المنیع و هذا المقرّ المتعالی
الرّفیع ینبغی لهم ان تکون قلوبهم مقدّسة عن ذکر اشارات القبل و کلمات الّتی
اسّسوها اولی الجهل الّذین اشتهرت انفسهم بالعلم و الفضل لک الحمد یا الهی
بما جعلت للواردین فی هذا البساط الأرفع الأعلی مقاماً لا تناله طیور افئدة
اهل البقآء و لا حقایق من فی الانشآء کما نطق به مظهر وحیک و مطلع الهامک و
نقطة مشیّتک نطفة سنة ذلک الظّهور تکون اقوی عن کلّ البیان و کذلک نزّل فی قیّوم
الکتب و الواح اخری فی ذکر هذا المقام الأعلی علی شأن تتحیّر منه عقول اولی
الألباب فما اعلی علوّ هذا الخلق الّذی خلقتهم من ساذج امرک و انشأتهم من
جوهر فطرتک اذاً یا الهی کیف یلیق ان یکون هذا الخلق البدیع و هذا الصّنع
الأرفع الرّفیع

\end{otherlanguage}

Wa-alladhīna yurīdūn an yadkhulū hādhā al-bisāṭ al-mumtaniʿ al-manīʿ
wa-hādhā al-maqarr al-mutaʿālī al-rafīʿ yanbaghī lahum an takūna
qulūbuhum muqaddasa ʿan dhikr ishārāt al-qabl wa-kalimāt allatī assasūhā
ulū al-jahl alladhīna ishtaharat anfusuhum bi-al-ʿilm wa-al-faḍl. Laka
al-ḥamd yā ilāhī bi-mā jaʿalta li-al-wāridīn fī hādhā al-bisāṭ al-arfaʿ
al-aʿlā maqāman lā tanāluhu ṭuyūr afʾidat ahl al-baqāʾ wa-lā ḥaqāʾiq man
fī al-inshāʾ kamā naṭaqa bihi maẓhar waḥyika wa-maṭlaʿ ilhāmika
wa-nuqṭat mashīʾatika: nuṭfat sanat dhālika al-ẓuhūr takūn aqwā ʿan kull
al-Bayān. Wa-kadhālika nuzzila fī Qayyūm al-Kutub wa-alwāḥ ukhrā fī
dhikr hādhā al-maqām al-aʿlā ʿalā shaʾn tataḥayyar minhu ʿuqūl ulī
al-albāb. Fa-mā aʿlā ʿuluww hādhā al-khalq alladhī khalaqtahum min
sādhij amrika wa-anshaʾtahum min jawhar fiṭratika. Idhan yā ilāhī kayfa
yalīq an yakūna hādhā al-khalq al-badīʿ wa-hādhā al-ṣunʿ al-arfaʿ
al-rafīʿ.

Those who wish to enter this inaccessible, unassailable court and this
exalted, lofty abode should have hearts sanctified from the mention of
the allusions of the past and the words established by the ignorant who
have become renowned for learning and excellence. Praise be to You, my
God, for having appointed for those who enter this most exalted and
lofty court a station that neither the birds of the hearts of the people
of eternity nor the realities of those in creation can attain. Thus has
the Manifestation of Your revelation, the Dawning-Place of Your
inspiration, and the Point of Your will declared: ``A drop of seed in
the year of that Manifestation is more potent than the whole Bayān.''
Likewise, in the Qayyūm of the Books and other Tablets, words have been
revealed concerning this most exalted station in a manner that leaves
the minds of those endowed with understanding bewildered. How exalted is
the loftiness of these creatures whom You have created from the pure
essence of Your Cause and brought into being from the very substance of
Your creation! Then, my God, how is it fitting that this wondrous
creation and this most exalted and lofty handiwork should be

\subsection{8}\label{section-3127}

\begin{otherlanguage}{arabic}

انّی اکون متحیّراً یا الهی من وقر بعض الآذان و غطآء الأبصار و غشاوة بعض
القلوب بحیث من ندائک اهتزّت الأعراش و انفطرت السّموات و خرقت حجبات
اللّانهایات و اندکّت جبال حقایق الممکنات بتجلّی من اشراقات انوار وجهک و مع
کلّ ذلک ما انتبه بعض عبادک و بریّتک مع انّک تنادی فی کلّ الأحیان بأعلی
النّدآء فوق رؤوسهم و بکلّ الأنوار تکون قائماً تلقآء عیونهم قد اخرجت لهم من
جیب الارادة ید المشیّة و تقول بسلطان القوّة و العظمة فها هذه ید اقتداری
الّتی بسطتها علی لاهوت عظمتی و جبروت قدرتی و ملکوت سلطنتی و ممالک قدمی و
اقالیم امری و اطویت کلّها و ما قدّرت باحاطة علمی فیها فی یدی و انّها بیضآء
من انوار وجهی و شعشاع من ضیآء طلعتی و متلألئ من انوار جبینی قد اعطی
بفضلی و کرمی لمن یدخل فی ملکوتی و جبروتی و هم یذکرون عند تموّجات ابحر
النّور روات اولی الظّلمة و الغرور

\end{otherlanguage}

Innī akūn mutaḥayyiran yā ilāhī min waqr baʿḍ al-ādhān wa-ghiṭāʾ
al-abṣār wa-ghishāwat baʿḍ al-qulūb bi-ḥaythu min nidāʾika ihtazzat
al-aʿrāsh wa-infaṭarat al-samāwāt wa-khuriqat ḥujubāt al-lānihāyāt
wa-indakkat jibāl ḥaqāʾiq al-mumkināt bi-tajallin min ishrāqāt anwār
wajhika wa-maʿa kull dhālika mā intabaha baʿḍ ʿibādika wa-bariyyatika
maʿa annaka tunādī fī kull al-aḥyān bi-aʿlā al-nidāʾ fawqa ruʾūsihim
wa-bi-kull al-anwār takūn qāʾiman tilqāʾ ʿuyūnihim. Qad akhrajta lahum
min jayb al-irāda yad al-mashīʾa wa-taqūl bi-sulṭān al-quwwa
wa-al-ʿaẓama: fa-hā hādhihi yad iqtidārī allatī basaṭtuhā ʿalā lāhūt
ʿaẓamatī wa-jabarūt qudratī wa-malakūt salṭanatī wa-mamālik qidamī
wa-aqālīm amrī wa-aṭwaytu kullahā wa-mā qaddartu bi-iḥāṭat ʿilmī fīhā fī
yadī wa-innahā bayḍāʾ min anwār wajhī wa-shaʿshāʿ min ḍiyāʾ ṭalʿatī
wa-mutalaʾliʾ min anwār jabīnī qad uʿṭī bi-faḍlī wa-karamī li-man
yadkhulu fī malakūtī wa-jabarūtī. Wa-hum yadhkurūn ʿinda tamawwujāt
abḥur al-nūr ruwāt ulī al-ẓulma wa-al-ghurūr.

I am bewildered, my God, by the deafness of some ears, the covering over
eyes, and the veil upon some hearts. At Your call the thrones have
trembled, the heavens have been rent asunder, the veils of endlessness
have been torn apart, and the mountains of the realities of contingent
beings have been crushed by a single effulgence of the radiance of the
lights of Your face. Yet, despite all this, some of Your servants and
creatures have not awakened, although at every moment You call out above
their heads with the loudest call and stand before their eyes in the
fullness of light. You have drawn forth for them the hand of will from
the bosom of purpose and say with the sovereignty of strength and
grandeur: ``Here is the hand of My might, which I have stretched over
the divine realm of My grandeur, the dominion of My power, the kingdom
of My sovereignty, the realms of My eternity, and the regions of My
Cause. I have folded them all in My hand, together with all that I have
ordained within them through My all-encompassing knowledge. It is white
with the lights of My face, radiant with the splendor of My countenance,
and gleaming with the lights of My brow. Through My grace and
generosity, I bestow it upon whoever enters My kingdom and My
dominion.'' Yet, while the oceans of light are surging, they make
mention of the narrators of the people of darkness and delusion.

\subsection{9}\label{section-3128}

\begin{otherlanguage}{arabic}

فیا الهی قد خجل کلّ الوجود من اذکارهم و هم لا یخجلون یستحیی کلّ الممکنات
من اقوالهم و هم لا یستحیون هل الّذی انتبه من ندائک و قام بروح امرک و حیّ
من نفس رحمانیّتک و شرب عن کأس عنایتک یجری لسانه بذکر الموهومات الّتی کانت
فی السن الغافلین من بریّتک و الجاهلین من خلقک لا وجمالک المختار الّا الّذی
جعلت له السّمع آیة العذاب من عندک و البصر حفرة السّقر فی وجهه بأمرک

\end{otherlanguage}

Fa-yā ilāhī qad khajila kull al-wujūd min adhkārihim wa-hum lā
yakhjalūn. Yastaḥyī kull al-mumkināt min aqwālihim wa-hum lā yastaḥyūn.
Hal alladhī intabaha min nidāʾika wa-qāma bi-rūḥ amrika wa-ḥayya min
nafas raḥmāniyyatika wa-shariba ʿan kaʾs ʿināyatika yajrī lisānuhu
bi-dhikr al-mawhūmāt allatī kānat fī alsun al-ghāfilīn min bariyyatika
wa-al-jāhilīn min khalqika? Lā wa-jamālika al-mukhtār illā alladhī
jaʿalta lahu al-samʿ āyat al-ʿadhāb min ʿindika wa-al-baṣar ḥufrat
al-saqar fī wajhihi bi-amrika.

My God, all existence has been ashamed of their words of remembrance,
yet they are not ashamed. All contingent beings feel shame at their
sayings, yet they feel none. Could one who has awakened at Your call,
arisen through the spirit of Your Cause, been quickened by the breath of
Your mercy, and drunk from the cup of Your providence allow his tongue
to utter the vain imaginings current on the tongues of the heedless
among Your creatures and the ignorant among Your creation? No, by Your
sovereign Beauty! None could do so except one for whom You have made
hearing a sign of punishment from You, and sight, by Your command, a pit
of hell in his face.

\subsection{10}\label{section-3129}

\begin{otherlanguage}{arabic}

فیا الهی قد اجترحت بتلک الأذکار بین یدی سلطان عظمتک ولکن انّک تعلم بأنّ
النّار تلتهب من شعراتی بما اسمع اقوال الّذین ما اطّلعوا بأسرار امرک و خفیّات
ما فی علمک و ترکوا مقامات العالیة و اقتنعوا بمراتب الدّانیة کأنّهم نسوا ما
نزّل من جبروت امرک قلت و قولک الحقّ یا بن منظر الأعلی قدّرت لک من شجر
الأبهی فواکه الأصفی کیف اعرضت عنها و رضیت بالّذی هو ادنی فارجع الی ما هو
خیر لک فی الأفق الأعلی

\end{otherlanguage}

Fa-yā ilāhī qad ijtaraḥtu bi-tilka al-adhkār bayna yaday sulṭān
ʿaẓamatika wa-lākin innaka taʿlam bi-anna al-nār taltahib min shaʿarātī
bi-mā asmaʿ aqwāl alladhīna mā iṭṭalaʿū bi-asrār amrika wa-khafiyyāt mā
fī ʿilmika wa-tarakū maqāmāt al-ʿāliya wa-iqtanaʿū bi-marātib al-dāniya
ka-annahum nasū mā nuzzila min jabarūt amrika. Qulta wa-qawluka al-ḥaqq:
yā ibn manẓar al-aʿlā qaddartu laka min shajar al-abhā fawākih al-aṣfā
kayfa aʿraḍta ʿanhā wa-raḍīta bi-alladhī huwa adnā fa-irjaʿ ilā mā huwa
khayr laka fī al-ufuq al-aʿlā.

My God, I have indeed ventured to speak these words before the
sovereignty of Your grandeur. Yet You know that fire blazes from every
hair of mine when I hear the sayings of those who have not become
acquainted with the mysteries of Your Cause and the hidden things within
Your knowledge, who have forsaken exalted stations and contented
themselves with lowly ranks. It is as though they have forgotten what
was sent down from the dominion of Your Cause. You said, and Your word
is truth: ``O son of the most exalted vision! I ordained for you the
purest fruits from the all-glorious tree. How have you turned away from
them and been content with what is inferior? Return, then, to what is
better for you in the most exalted horizon.''

\subsection{11}\label{section-3130}

\begin{otherlanguage}{arabic}

و اری بعض منهم یا الهی احتجبوا بقول من الکتاب بعد الّذی انّه نزّل بالتّصریح
من غیر التّلویح بأنّ کلّ البیان لا یعادل بکلمة من عنده کلّ ما ذکر من الحجّة و
البرهان ثمّ الدّلیل و الآیات و البیّنات و العلامات لم یکن الّا لهذا الظّهور
الأعظم الأعظم بحیث ما ذکرت فیه کلمة و ما وجد فیه شیء و ما نطق بحرف الّا و
یکون آیةً لأمرک و برهاناً لسلطنتک کما تکلّم به لسان المشیّة من عندک کلّ ما
خلق و یخلق یکون آیة له خلق بقوله من عنده لیکون حجّته لنفسه یوم ظهوره بل
جعل یا الهی کلّ البیان معلّقاً باذنک و قبولک قل و قوله الحقّ فوعزّتک لو تغفرنّ
البیان و من فیه لا ینقص عن ملکک من شیء و ان لا تقبلنّ البیان و من فیه لا
یزید فی ملکک من شیء بل ان تقبل ذلک تدخل تلک الأفئدة المتیّمة فی ظلّ ظلال
محبّتک و ان تردّه یفنی کأنّه لم یکن له ذکر من قبل

\end{otherlanguage}

Wa-arā baʿḍ minhum yā ilāhī iḥtajabū bi-qawl min al-kitāb baʿd alladhī
innahu nuzzila bi-al-taṣrīḥ min ghayr al-talwīḥ bi-anna kull al-Bayān lā
yuʿādil bi-kalima min ʿindihi. Kull mā dhukira min al-ḥujja wa-al-burhān
thumma al-dalīl wa-al-āyāt wa-al-bayyināt wa-al-ʿalāmāt lam yakun illā
li-hādhā al-ẓuhūr al-aʿẓam al-aʿẓam bi-ḥaythu mā dhukirat fīhi kalima
wa-mā wujida fīhi shayʾ wa-mā nuṭiqa bi-ḥarf illā wa-yakūn āyatan
li-amrika wa-burhānan li-salṭanatika kamā takallama bihi lisān
al-mashīʾa min ʿindika: kull mā khuliqa wa-yukhlaq yakūn āyatan lahu
khuliqa bi-qawlihi min ʿindihi li-yakūna ḥujjatahu li-nafsihi yawm
ẓuhūrihi. Bal jaʿala yā ilāhī kull al-Bayān muʿallaqan bi-idhnika
wa-qabūlika. Qul wa-qawluhu al-ḥaqq: fa-wa-ʿizzatika law taghfiranna
al-Bayān wa-man fīhi lā yanquṣ ʿan mulkika min shayʾ wa-in lā taqbalanna
al-Bayān wa-man fīhi lā yazīd fī mulkika min shayʾ bal in taqbal dhālika
tudkhil tilka al-afʾida al-mutayyama fī ẓill ẓilāl maḥabbatika wa-in
taruddahu yafnā ka-annahu lam yakun lahu dhikr min qabl.

I see some of them, my God, veiled by a statement in the Book, even
though it was revealed explicitly, not by allusion, that the whole Bayān
does not equal a single word from Him. Whatever was mentioned of proof
and testimony, of evidence, verses, clear signs, and tokens, was solely
for this most great, most great Revelation. No word was mentioned in it,
no thing was found in it, and no letter was uttered except as a sign of
Your Cause and a proof of Your sovereignty. Thus the Tongue of Will has
spoken from Your presence: ``Everything that has been or will be created
is a sign of Him, created by His utterance from Him, to be His proof for
Himself on the day of His manifestation.'' Indeed, my God, He made the
whole Bayān dependent upon Your permission and acceptance. Say---and His
word is truth: ``By Your might! If You forgive the Bayān and all who are
in it, nothing is diminished from Your dominion; and if You do not
accept the Bayān and all who are in it, nothing is added to Your
dominion. Rather, if You accept it, You bring those enraptured hearts
beneath the sheltering shade of Your love; and if You reject it, it
perishes as though it had never been mentioned before.''

\subsection{12}\label{section-3131}

\begin{otherlanguage}{arabic}

اذاً اسألک یا مالک الوجود و ملیک الغیب و الشّهود بأن تؤیّد الکلّ علی التّوجّه
الی افق فضلک و افضالک مقدّساً من اشارات القبل و البعد و تعرّفهم ما نزّلته فی
کتابک لیتمسّکوا بجوهر ما فیه لأنّ هذا مقصودک فیما نزّل بالحقّ و محبوبک فیما
انزلته فی کتبک و الواحک انّک انت المقتدر المهیمن المتعالی العزیز الحکیم

\end{otherlanguage}

Idhan asʾaluka yā Mālik al-wujūd wa-Malīk al-ghayb wa-al-shuhūd bi-an
tuʾayyida al-kull ʿalā al-tawajjuh ilā ufuq faḍlika wa-ifḍālika
muqaddasan min ishārāt al-qabl wa-al-baʿd wa-tuʿarrifahum mā nazzaltahu
fī kitābika li-yatamassakū bi-jawhar mā fīhi li-anna hādhā maqṣūduka
fī-mā nuzzila bi-al-ḥaqq wa-maḥbūbuka fī-mā anzaltahu fī kutubika
wa-alwāḥika. Innaka anta al-Muqtadir al-Muhaymin al-Mutaʿālī al-ʿAzīz
al-Ḥakīm.

Then I ask You, O Possessor of existence and King of the unseen and the
seen, to strengthen everyone to turn toward the horizon of Your grace
and bounty, sanctified from the allusions of past and future, and to
acquaint them with what You have sent down in Your Book, so that they
may hold fast to the essence of its contents. For this is Your purpose
in what has been sent down in truth, and Your desire in what You have
revealed in Your Books and Tablets. You are indeed the All-Powerful, the
All-Protecting, the Most Exalted, the Almighty, the All-Wise.

\subsection{13}\label{section-3132}

\begin{otherlanguage}{persian}

ای اله ابدی تو شاهد و گواهی که این عبد در جمیع احیان و اوان جز طالب رضای
تو نبوده و انشآءالله بعنایت تو نخواهد بود همیشۀ اوقات چشم این عبد بمنظر
اکبر بوده و قلب بمصدر امر متوجّه و آنچه دیده و عارف بآن شده خالصاً لوجهک
خدمت دوستان تو معروض داشته که شاید نفوس قابله از کوثر حیوان که در ایّام
ظهور از اصبع فضل و احسان جاری شده بیاشامند و بمطلع امر و مشرق وحی توجّه
نمایند

\end{otherlanguage}

Ay Ilāh-i abadī, tu shāhid va guvāhī ki īn ʿabd dar jamīʿ-i aḥyān va
avān juz ṭālib-i riḍā-yi tu nabūda va in shāʾ Allāh bi-ʿināyat-i tu
nakhvāhad būd. Hamīsha-yi awqāt chashm-i īn ʿabd bi-manẓar-i akbar būda
va qalb bi-maṣdar-i amr mutavajjih, va ānchi dīda va ʿārif bi-ān shuda,
khāliṣan li-wajhika, khidmat-i dūstān-i tu maʿrūḍ dāshta ki shāyad
nufūs-i qābila az kawthar-i ḥayavān ki dar ayyām-i ẓuhūr az iṣbaʿ-i faḍl
va iḥsān jārī shuda biyāshāmand va bi-maṭlaʿ-i amr va mashriq-i vaḥy
tavajjuh namāyand.

O eternal God, You bear witness and testify that this servant has at all
times sought nothing but Your good pleasure and, God willing, through
Your grace will never do otherwise. This servant's eyes have always been
fixed on the Most Great Scene, and his heart turned toward the Source of
the Cause. Whatever he has seen and come to know he has laid before Your
friends solely for Your sake, so that receptive souls may drink from the
Kawthar of life that has flowed from the finger of grace and bounty in
the days of the Revelation, and may turn toward the Dawning-place of the
Cause and the Dayspring of revelation.

\subsection{14}\label{section-3133}

\begin{otherlanguage}{persian}

و چندی بود که این عبد صمت را بر نطق اختیار نموده تا آنکه در این ایّام
مشاهده شد که غبرۀ تیرۀ ظلمانیّه از شطر نفوس غافله در حرکت آمده و بیم آن
رفت که بعضی از ابصار ضعیفه را تیره نماید لذا این عبد متوکّلاً علیک و
منقطعاً عمّا سواک بر خود لازم شمرد که بعضی از امور واقعه را خدمت احبّای تو
که بر بساط استقامت آرمیده‌اند و از رحیق عنایت نوشیده‌اند معروض دارد که
شاید ضعفا را از ظلمت هواهای نفسانیّه و شبهات انفس شیطانیّه حفظ فرمایند قسم
بعظمت تو ای پروردگار که در این وقت جمیع اعضا مرتعش و ارکان متزلزل این
عبد کجا قابل آنست که کلماتش مقبول آید و یا بیانش مذکور شود نیست بحت
چگونه علم هستی برافرازد و عدم صرف کجا تلقاء ظهور قدم اظهار وجود نماید چه
که عرض همین مطالب اظهار وجود است و آن از اعظم خطیئات محسوب

\end{otherlanguage}

Va chandī būd ki īn ʿabd ṣamt rā bar nuṭq ikhtiyār namūda, tā ānki dar
īn ayyām mushāhada shud ki ghubra-yi tīra-yi ẓulmāniyya az shaṭr-i
nufūs-i ghāfila dar ḥarakat āmada va bīm ān raft ki baʿḍī az abṣār-i
ḍaʿīfa rā tīra namāyad. Lidhā īn ʿabd, mutawakkilan ʿalayka va
munqaṭiʿan ʿammā siwāka, bar khud lāzim shumurd ki baʿḍī az umūr-i
vāqiʿa rā khidmat-i aḥibbā-yi tu ki bar bisāṭ-i istiqāmat ārāmīda-and va
az raḥīq-i ʿināyat nūshīda-and maʿrūḍ dārad, ki shāyad ḍuʿafā rā az
ẓulmat-i havāhā-yi nafsāniyya va shubahāt-i anfus-i shayṭāniyya ḥifẓ
farmāyand. Qasam bi-ʿaẓamat-i tu, ay Parvardigār, ki dar īn vaqt jamīʿ-i
aʿḍā murtaʿish va arkān mutazalzil. Īn ʿabd kujā qābil-i ānast ki
kalimātash maqbūl āyad va yā bayānash madhkūr shavad? Nīst-i baḥt
chigūna ʿalam-i hastī bar-afrāzad va ʿadam-i ṣirf kujā tilqāʾ-i ẓuhūr-i
qidam iẓhār-i vujūd namāyad? Chi ki ʿarḍ-i hamīn maṭālib iẓhār-i vujūd
ast va ān az aʿẓam-i khaṭīʾāt maḥsūb.

For some time this servant had chosen silence over speech, until in
these days he observed a dark, gloomy dust rising from the direction of
heedless souls, and feared that it might cloud some weak eyes.
Therefore, trusting in You and detached from all besides You, this
servant considered it his duty to lay some of the events that have
occurred before Your loved ones, who rest upon the carpet of
steadfastness and have drunk the choice wine of Your grace, so that they
might protect the weak from the darkness of selfish desires and the
doubts cast by satanic souls. I swear by Your grandeur, O Lord, that at
this moment every limb trembles and every part of my frame is shaken.
How could this servant be worthy of having his words accepted or his
utterance mentioned? How can sheer nothingness raise the banner of
existence, and how can pure nonexistence assert its being before the
Revelation of Him who is eternal? For presenting these very matters is
an assertion of existence, and that is reckoned among the greatest
transgressions.

\subsection{15}\label{section-3134}

\begin{otherlanguage}{arabic}

ولکن چون لأجل تقرّب عباد بشاطی قرب و لقا و اقبال نفوس بقبلۀ من فی الأرض و
السّمآء عرض میشود لذا از بحر غفران طلب عفو مینمایم مع آنکه این عبد بیقین
میداند که تو لم‌یزل از اذکار و عرفان و اقبال ناس غنی بوده و هستی لا ینفعک
اقبالهم و لا یضرّک اعراضهم هیچ وصفی بساحت اقدست نرسد و هیچ ذکری به بساط
احدیّتت درنیاید بلکه ذکر احدیّت در ساحت اقدست ذنب صرفست و غفلت محض چه که
شرافت او به نسبتها الیک بوده لا بنسبتک الیها لو تطردها بأسرها من یقدر ان
یمنعک فی ذلک و لو ترفعها کما رفعتها من یقدر ان یعترض علی فعلک انّک انت
المحمود فی کلّ ما اردت و ترید و فی کلّ ما قلت و تقول و فی کلّ ما اظهرت و
تظهر

\end{otherlanguage}

Valākin chūn li-ajl-i taqarrub-i ʿibād bi-shāṭī-yi qurb va liqā va
iqbāl-i nufūs bi-qibla-yi man fī al-arḍ wa-l-samāʾ ʿarḍ mīshavad, lidhā
az baḥr-i ghufrān ṭalab-i ʿafv mīnamāyam. Maʿa ānki īn ʿabd bi-yaqīn
mīdānad ki tu lam-yazal az adhkār va ʿirfān va iqbāl-i nās ghanī būda va
hastī. Lā yanfaʿuka iqbāluhum wa-lā yaḍurruka iʿrāḍuhum. Hīch vaṣfī
bi-sāḥat-i aqdasat narasad va hīch dhikrī bi bisāṭ-i aḥadiyyatat
dar-nayāyad. Balki dhikr-i aḥadiyyat dar sāḥat-i aqdasat dhanb-i ṣirfast
va ghaflat-i maḥḍ, chi ki sharāfat-i ū bi-nisbatihā ilayka būda lā
bi-nisbatika ilayhā. Law taṭrudhā bi-asrihā man yaqdiru an yamnaʿaka fī
dhālika, wa-law tarfaʿhā kamā rafaʿtahā man yaqdiru an yaʿtariḍa ʿalā
fiʿlika? Innaka anta al-maḥmūdu fī kulli mā aradta wa-turīdu wa-fī kulli
mā qulta wa-taqūlu wa-fī kulli mā aẓharta wa-tuẓhiru.

But since this is presented so that Your servants may draw near to the
shore of nearness and reunion, and souls may turn toward the Point of
Adoration of all in earth and heaven, I therefore seek pardon from the
ocean of forgiveness. Yet this servant knows with certainty that You
have always been, and remain, independent of people's remembrance,
recognition, and turning toward You. Their turning toward You does not
benefit You, nor does their turning away harm You. No description
reaches Your most holy court, and no mention enters the precincts of
Your oneness. Indeed, to mention oneness in Your most holy court is
sheer sin and utter heedlessness, for its nobility lies in its relation
to You, not in Your relation to it. If You banish it entirely, who can
prevent You from doing so? And if You exalt it, as You have exalted it,
who can object to Your action? You are indeed the One praised in all
that You have willed and will, in all that You have said and say, and in
all that You have manifested and manifest.

\subsection{16}\label{section-3135}

\begin{otherlanguage}{arabic}

فوعزّتک یا الهی لا احبّ ان اختار لنفسی الّا ما اخترته لی و لا ارید ان اتکلّم
الّا بما امرتنی به یکون طرف کینونتی ناظراً الی افق اوامرک و مشرق احکامک لو
تطوی بساط الأذکار المنبسطة المنتشرة فی الأرض لأقول انت الحاکم فی امرک و
لو تأمرنی بما ینکره من فی الملک لأقول انت المطاع فی حکمک اسألک اللّهمّ بأن
تجعلنی من الّذین لم یتکلّموا الّا باذنک و لم یتحرّکوا الّا بارادتک انّک انت
المقتدر المتعالی المهیمن القیّوم

\end{otherlanguage}

Fa-wa-ʿizzatika yā Ilāhī lā uḥibbu an akhtāra li-nafsī illā mā
ikhtartahu lī wa-lā urīdu an atakallama illā bi-mā amartanī bihi. Yakūnu
ṭarfu kaynūnatī nāẓiran ilā ufuqi awāmirika wa-mashriqi aḥkāmika. Law
taṭwī bisāṭa al-adhkāri al-munbasiṭati al-muntashirati fī al-arḍi
la-aqūlu anta al-Ḥākimu fī amrika, wa-law taʾmuranī bi-mā yunkiruhu man
fī al-mulki la-aqūlu anta al-muṭāʿu fī ḥukmika. Asʾaluka Allāhumma bi-an
tajʿalanī mina alladhīna lam yatakallamū illā bi-idhnika wa-lam
yataḥarrakū illā bi-irādatika. Innaka anta al-Muqtadiru al-Mutaʿālī
al-Muhayminu al-Qayyūm.

By Your might, my God, I wish to choose for myself only what You have
chosen for me, and I desire to speak only what You have commanded me to
say. The eye of my being looks toward the horizon of Your commands and
the dayspring of Your laws. If You roll up the carpet of remembrance
spread far and wide across the earth, I will say: You are sovereign in
Your command. And if You command me to do what all in the world would
reject, I will say: You are to be obeyed in Your decree. I ask You, O
God, to make me one of those who speak only with Your permission and
move only by Your will. You are indeed the All-Powerful, the Most
Exalted, the All-Protecting, the Self-Subsisting.

\subsection{17}\label{section-3136}

\begin{otherlanguage}{persian}

عرض میشود حین ورود جمال قدم در سجن عکّا چندی امر بسیار شدید بود و تفصیل
آن در همان ایّام خدمت بعضی از دوستان عرض شد تکرار آن لازم نیست جمیع در
کمال عسر و مشقّت بودند تا آنکه یومی از ایّام لوحی از مصدر امر نازل و این
عبد تلقاء وجه تحریر مینمود آیه‌ئی از سماء مشیّت نازل و این سجن را سجن اعظم
موسوم فرمودند این عبد گمان نمود که نظر به این شداید وارده است که باین
اسم موسوم شده تا آنکه یومی از ایّام آیاتی نازل مضمون آن اینست که
میفرمایند تفکّر نمائید بچه سبب این سجن را باعظم نامیدیم این بیان که از
مطلع وحی رحمن ظاهر شد این عبد خود را بسیار مضطرب و متزلزل مشاهده نمود
ولکن متوکّلاً علی الله ساکن بودم و بهیچوجه بر سؤال جسارت ننمودم تا آنکه
روزی ذکر محبّت و اقبال یکی از اهل این مدینه تلقاء عرش عرض شد لوحی مخصوص
او از سماء احدیّت نازل در آن لوح مفصّلاً امورات حادثۀ این ارض را ذکر
فرمودند و آن لوح حال موجود است بعد از تنزیل آن لوح این عبد دانست که سبب
اعظمیّت این سجن بلایای کلّیّه است که بر جمال احدیّه وارد میشود تا آنکه یومی
از ایّام یکی از همراهان با بعضی از اصحاب بنزاع و جدال برخاست مع آنکه خود
او خاطی و ظالم بود عریضه بساحت اقدس معروض داشت و در آن عریضه اظهار
مظلومیّت خود نمود در جواب او لوحی بلسان پارسی از سحاب حزن سبحانی نازل و
صورت آن این است

\end{otherlanguage}

ʿArḍ mīshavad ḥīn-i vurūd-i Jamāl-i Qidam dar sijn-i ʿAkkā chandī amr
bisyār shadīd būd va tafṣīl-i ān dar hamān ayyām khidmat-i baʿḍī az
dūstān ʿarḍ shud; takrār-i ān lāzim nīst. Jamīʿ dar kamāl-i ʿusr va
mashaqqat būdand tā ānki yawmī az ayyām lawḥī az maṣdar-i amr nāzil va
īn ʿabd tilqāʾ-i vajh taḥrīr mīnamūd. Āya-ī az samā-yi mashiyyat nāzil
va īn sijn rā Sijn-i Aʿẓam mawsūm farmūdand. Īn ʿabd gumān namūd ki
naẓar bi īn shadāyid-i vārida ast ki bi-īn ism mawsūm shuda, tā ānki
yawmī az ayyām āyātī nāzil, maḍmūn-i ān īnast ki mīfarmāyand: tafakkur
namāyīd bi-chi sabab īn sijn rā bi-aʿẓam nāmīdīm. Īn bayān ki az
maṭlaʿ-i vaḥy-i Raḥmān ẓāhir shud īn ʿabd khud rā bisyār muḍṭarib va
mutazalzil mushāhada namūd, valākin mutawakkilan ʿalā Allāh sākin būdam
va bi-hīch-vajh bar suʾāl jisārat nanamūdam. Tā ānki rūzī dhikr-i
maḥabbat va iqbāl-i yakī az ahl-i īn madīna tilqāʾ-i ʿarsh ʿarḍ shud.
Lawḥī makhṣūṣ-i ū az samā-yi aḥadiyyat nāzil; dar ān lawḥ mufaṣṣalan
umūrāt-i ḥāditha-yi īn arḍ rā dhikr farmūdand va ān lawḥ ḥāl mawjūd ast.
Baʿd az tanzīl-i ān lawḥ īn ʿabd dānist ki sabab-i aʿẓamiyyat-i īn sijn
balāyā-yi kulliyya ast ki bar Jamāl-i Aḥadiyya vārid mīshavad. Tā ānki
yawmī az ayyām yakī az hamrāhān bā baʿḍī az aṣḥāb bi-nizāʿ va jidāl
bar-khāst. Maʿa ānki khud-i ū khāṭī va ẓālim būd ʿarīḍa bi-sāḥat-i aqdas
maʿrūḍ dāsht va dar ān ʿarīḍa iẓhār-i maẓlūmiyyat-i khud namūd. Dar
javāb-i ū lawḥī bi-lisān-i pārsī az saḥāb-i ḥuzn-i subḥānī nāzil va
ṣūrat-i ān īn ast.

I submit that, when the Ancient Beauty entered the prison of `Akkā,
conditions were exceedingly harsh for a time. Their details were
reported to some of the friends in those very days, and need not be
repeated. All were in extreme hardship and distress, until one day a
Tablet was revealed from the Source of the Cause while this servant was
transcribing it in His presence. A verse descended from the heaven of
the divine Will in which He named this prison the Most Great Prison.
This servant supposed that it had been given this name on account of
these severe afflictions, until one day verses were revealed whose
purport was this: ``Reflect on why We have named this prison the Most
Great.'' When this utterance appeared from the Dawning-place of the
revelation of the All-Merciful, this servant found himself deeply
troubled and shaken. Yet, trusting in God, I remained still and did not
venture to ask any question. Then one day the love and devotion of one
of the inhabitants of this city were mentioned before the throne. A
Tablet addressed specifically to him descended from the heaven of
oneness. In that Tablet He spoke in detail of the events occurring in
this land, and that Tablet is still extant. After its revelation this
servant understood that the reason for this prison's surpassing
greatness was the all-encompassing afflictions befalling the Beauty of
Oneness. Then one day one of the companions began to quarrel and dispute
with some of the others. Although he himself was at fault and was the
aggressor, he submitted a petition to the Most Holy Court in which he
claimed to have been wronged. In response to him a Tablet in Persian
descended from the cloud of divine sorrow. Its text is as follows.

\subsection{18}\label{section-3137}

\begin{otherlanguage}{arabic}

هوالله

\end{otherlanguage}

Huwa Allāh.

He is God.

\subsection{19}\label{section-3138}

\begin{otherlanguage}{arabic}

عاشق را نزد معشوق اظهار هستی و خودبینی جایز نه اگر خطوه‌ئی از این سبیل
تجاوز نماید از عشّاق محسوب نه نکند عشق نفس زنده قبول نکند باز موش مرده
شکار

\end{otherlanguage}

ʿĀshiq rā nazd-i maʿshūq iẓhār-i hastī va khudbīnī jāyiz na. Agar
khaṭva-ī az īn sabīl tajāvuz namāyad az ʿushshāq maḥsūb na. Nakunad
ʿishq nafs-i zinda qabūl; nakunad bāz mūsh-i murda shikār.

A lover may not assert his own existence or regard himself in the
presence of the beloved. If he strays even one step from this path, he
is not counted among the lovers. Love does not accept a living self; a
falcon does not hunt a dead mouse.

\subsection{20}\label{section-3139}

\begin{otherlanguage}{persian}

بلی بعضی از عاشقان اظهار صدمات شداید خو را در پیش‌گاه محبوب امکان
نموده‌اند و مقصود از آن اشتغال با محبوب و اصغای حضرت مقصود بوده نه ذکر
نفس و هوی حال قاصدی از اعلی مقاصد عزّ مقصود نازل و بکلماتی ناطق است
بفهمید که که میگوید و چه میگوید تالله لو عرفتم و علمتم ما ورآء ستر
الکبریآء من اسرار ربّکم العلیّ الأعلی لفدیتم انفسکم حبّاً لله مالک الأسمآء

\end{otherlanguage}

Balī baʿḍī az ʿāshiqān iẓhār-i ṣadamāt-i shadāyid-i {[}corrupted text{]}
rā dar pīshgāh-i Maḥbūb-i imkān namūda-and va maqṣūd az ān ishtighāl bā
maḥbūb va iṣghā-yi Ḥaḍrat-i Maqṣūd būda na dhikr-i nafs va havā. Ḥāl
qāṣidī az aʿlā maqāṣid-i ʿizz-i maqṣūd nāzil va bi-kalimātī nāṭiq ast.
Bifahmīd ki ki mīgūyad va chi mīgūyad. Ta-Allāhi law ʿaraftum
wa-ʿalimtum mā warāʾa sitri al-kibriyāʾi min asrāri rabbikumu al-ʿAliyyi
al-Aʿlā la-fadaytum anfusakum ḥubban li-llāhi Māliki al-asmāʾ.

Indeed, some lovers have recounted the blows of the afflictions
{[}corrupted original text{]} before the Beloved of all existence. Their
purpose was to commune with the Beloved and gain the attention of Him
who is the Object of all desire, not to speak of self and desire. Now a
messenger has descended from the loftiest retreats of the glory of the
Desired One and is speaking these words. Understand who is speaking and
what he says. By God, if you recognized and knew the mysteries of your
Lord, the Exalted, the Most High, that lie beyond the veil of grandeur,
you would sacrifice your lives out of love for God, the Lord of Names.

\subsection{21}\label{section-3140}

\begin{otherlanguage}{persian}

باری قاصد معهود حکایتی ذکر نموده که وقتی در طور اشراق بودیم و محبوب آفاق
بقصد جبل‌ها بیرون تشریف بردند و بعجز تمام از مدّعیان محبّت رجا فرمودند که
این سفر و هجرت اگرچه بظاهر سهل و آسانست ولکن در باطن شدید و باب امتحان و
اگر بصورت بسلطنت و اقتدار مشهود ولکن در معنی بمحنت و اضطرار لایحصی مکنون
عرض خود مبرید و زحمت بر خود مدهید و بگذارید تا بنفس خود هجرت نمایم آنچه
کلمات محبّت‌آمیز و شفقت‌انگیز بود تلویحاً و تصریحاً فرمودند مفید نیفتاد این
بنده و جمعی بادّعای آنکه زادی بجز رضای دوست نخواهیم و مقصودی جز وجه محبوب
نداریم بگمان خود معتکف و از نصح و یقین دوست غافل عزم سفر نمودیم و با
طلعت محبوب همسفر و همسیر گشتیم قدری که بادیه پیمودیم نار حبّ مخمود و جمال
شوق محجوب تا آنکه از آن مقام تجاوز نمودیم حسنات را سیّئات مشاهده نمودیم و
سیّئات را عین حسنات شمردیم

\end{otherlanguage}

Bārī qāṣid-i maʿhūd ḥikāyatī dhikr namūda ki vaqtī dar Ṭūr-i ishrāq
būdīm va Maḥbūb-i āfāq bi-qaṣd-i jabal-hā bīrūn tashrīf burdand va
bi-ʿajz-i tamām az muddaʿiyān-i maḥabbat rajā farmūdand ki īn safar va
hijrat agarchi bi-ẓāhir sahl va āsānast valākin dar bāṭin shadīd va
bāb-i imtiḥān, va agar bi-ṣūrat bi-salṭanat va iqtidār mashhūd valākin
dar maʿnī bi-miḥnat va iḍṭirār-i lā yuḥṣā maknūn. ʿIrḍ-i khud mabarīd va
zaḥmat bar khud madahīd va bigudhārīd tā bi-nafs-i khud hijrat namāyam.
Ānchi kalimāt-i maḥabbat-āmīz va shafaqat-angīz būd talvīḥan va taṣrīḥan
farmūdand mufīd nayuftād. Īn banda va jamʿī bi-iddiʿā-yi ānki zādī
bi-juz riḍā-yi dūst nakhvāhīm va maqṣūdī juz vajh-i maḥbūb nadārīm,
bi-gumān-i khud muʿtakif va az nuṣḥ va yaqīn-i dūst ghāfil, ʿazm-i safar
namūdīm va bā Ṭalʿat-i Maḥbūb hamsafar va hamsayr gashtīm. Qadrī ki
bādiya paymūdīm nār-i ḥubb makhmūd va jamāl-i shawq maḥjūb, tā ānki az
ān maqām tajāvuz namūdīm, ḥasanāt rā sayyiʾāt mushāhada namūdīm va
sayyiʾāt rā ʿayn-i ḥasanāt shumurdīm.

In brief, the messenger in question related this account: ``Once we were
on the Sinai of illumination, and the Beloved of the world set out for
the mountains. With the utmost humility He entreated those who professed
love: `Although this journey and migration seem outwardly easy and
simple, inwardly they are arduous and a doorway to testing. Though
outwardly they appear as sovereignty and power, inwardly they conceal
untold suffering and distress. Do not forfeit your honor or bring
hardship upon yourselves. Let Me migrate alone.' Whatever words of love
and compassion He spoke, whether indirectly or explicitly, proved of no
avail. This servant and a number of others claimed that we sought no
provision but the Friend's good pleasure and had no aim but the face of
the Beloved. Devoted, as we imagined, yet heedless of the Friend's
counsel and certainty, we resolved to travel and became the companions
and fellow travelers of the Beloved Countenance. We had traversed only a
little of the wilderness when the fire of love died down and the beauty
of longing was veiled. Then we went beyond that stage: we saw good deeds
as evil and counted evil deeds as goodness itself.

\subsection{22}\label{section-3141}

\begin{otherlanguage}{persian}

تا آنکه وارد جزیزۀ خضرا شدیم فلک الهی در آن ارض روحانی بر جودیّ امر مستوی
گشت و بعد بمراکب هوی در بیدای ظنون و اوهام در صبح و شام سایر بودیم گاهی
مجتمع و گاهی متفرّق و گاهی بحبّ و گاهی بغفلت ایّام و لیالی بسر میبردیم
سلسبیل بیان در کلّ احیان از کوثر فم رحمن جاری ولکن عطش مفقود و انوار وجه
از افق اجلال مشرق ولکن اقبال غیر موجود هر روز بر وهم و گمان افزودیم و از
توجّه بحقّ کاستیم و با این احوال غیر مرضیّه و شئونات غیر لایقه چنان در
غمرات غفلت و هوی غرق شدیم که از احوال خود هم غفلت نمودیم

\end{otherlanguage}

Tā ānki vārid-i {[}corrupted text{]} khaḍrā shudīm. Fulk-i ilāhī dar ān
arḍ-i rūḥānī bar Jūdī-yi amr mustavī gasht va baʿd bi-marākib-i havā dar
baydā-yi ẓunūn va awhām dar ṣubḥ va shām sāyir būdīm. Gāhī mujtamiʿ va
gāhī mutafarriq va gāhī bi-ḥubb va gāhī bi-ghaflat ayyām va layālī
bi-sar mīburdīm. Salsabīl-i bayān dar kull-i aḥyān az kawthar-i fam-i
Raḥmān jārī valākin ʿaṭash mafqūd, va anvār-i vajh az ufuq-i ijlāl
mushriq valākin iqbāl ghayr-i mawjūd. Har rūz bar vahm va gumān afzūdīm
va az tavajjuh bi-Ḥaqq kāstīm va bā īn aḥvāl-i ghayr-i marḍiyya va
shuʾūnāt-i ghayr-i lāyiqa chinān dar ghamarāt-i ghaflat va havā gharq
shudīm ki az aḥvāl-i khud ham ghaflat namūdīm.

Then we entered {[}corrupted original text{]} the verdant. In that
spiritual land the divine Ark came to rest upon the Jūdī of the Cause.
Thereafter, mounted on the steeds of desire, we traveled morning and
evening in the wilderness of conjecture and delusion. Sometimes together
and sometimes apart, sometimes in love and sometimes in heedlessness, we
passed our days and nights. At all times the Salsabīl of utterance
flowed from the Kawthar of the mouth of the All-Merciful, yet thirst was
absent; the lights of His face shone from the horizon of majesty, yet no
one turned toward them. Each day we added to our delusions and
conjectures and diminished our attention to God. With these unworthy
conditions and unbecoming ways, we sank so deeply into the floods of
heedlessness and desire that we became heedless even of our own
condition.

\subsection{23}\label{section-3142}

\begin{otherlanguage}{persian}

و در جمیع احوال طلعت محبوب را با کمال شفقت و ملاطفت ملاحظه مینمودیم بعضی
از ما متحیّر که اگر احاطهٴ علمیّهٴ الهیّه موجود چگونه میشود با این افعال
ردّیّه مقبول شویم بالأخره حقّ را غافل و خود را عاقل و عالم شمردیم غافل از
آنکه رحمت کبری مانعست از هتک استار و اگر خدمتی نمودیم اتبعناه بالمنّ و
الأذی

\end{otherlanguage}

Va dar jamīʿ-i aḥvāl Ṭalʿat-i Maḥbūb rā bā kamāl-i shafaqat va mulāṭafat
mulāḥaẓa mīnamūdīm. Baʿḍī az mā mutaḥayyir ki agar iḥāṭa-yi ʿilmiyya-yi
ilāhiyya mawjūd chigūna mīshavad bā īn afʿāl-i radiyya maqbūl shavīm.
Bi-l-ākhira Ḥaqq rā ghāfil va khud rā ʿāqil va ʿālim shumurdīm, ghāfil
az ānki raḥmat-i kubrā māniʿast az hatk-i astār, va agar khidmatī
namūdīm atbaʿnāhu bi-l-manni wa-l-adhā.

In every circumstance we saw the Beloved Countenance showing the utmost
compassion and tenderness. Some of us were bewildered: if God's
all-encompassing knowledge existed, how could we be accepted despite
these vile deeds? In the end we considered God heedless and ourselves
wise and knowing, unaware that His supreme mercy prevented the tearing
away of veils. And if we rendered any service, we followed it with
reminders of our favor and with hurt.

\subsection{24}\label{section-3143}

\begin{otherlanguage}{persian}

و بعد از آن ارض ارادهٴ هجرت فرمودند و مجدّداً کلّ را از حضور منع نمودند
مستشعر نشدیم و متنبّه نگشتیم که علّت منع چیست و سبب چه مرّة اخری هجرت
نمودیم و با حضرت مقصود بادیه‌ها پیمودیم تا آنکه وارد بارض اخری {[}شدیم{]}
و با دوست در یک محلّ آرمیدیم و سرّاً بهوای نفس مشغول گشتیم تا آنکه آتش هوی
غلبه نمود و از منظر ابهی ممنوع شدیم و از کثرت لقا قدر وصال از نظر افتاد
و پردهٴ حیا از هم درید

\end{otherlanguage}

Va baʿd az ān arḍ irāda-yi hijrat farmūdand va mujaddadan kull rā az
ḥuḍūr manʿ namūdand. Mustashʿir nashudīm va mutanabbih nagashtīm ki
ʿillat-i manʿ chīst va sabab chi. Marratan ukhrā hijrat namūdīm va bā
Ḥaḍrat-i Maqṣūd bādiya-hā paymūdīm tā ānki vārid bi-arḍ-i ukhrā
{[}shudīm{]} va bā dūst dar yak maḥall ārāmīdīm va sirran bi-havā-yi
nafs mashghūl gashtīm. Tā ānki ātash-i havā ghalaba namūd va az manẓar-i
abhā mamnūʿ shudīm va az kathrat-i liqā qadr-i viṣāl az naẓar uftād va
parda-yi ḥayā az ham darīd.

Afterward He resolved to migrate from that land and again forbade
everyone to accompany Him. We neither perceived nor took heed of the
reason for this prohibition or its cause. Once more we migrated and
traversed wildernesses with Him who is the Object of all desire, until
{[}we{]} entered another land. We dwelt in one place with the Friend and
secretly occupied ourselves with selfish desires, until the fire of
desire gained the upper hand and we were barred from the Most Glorious
Presence. Through the frequency of our meetings, we lost sight of the
value of reunion, and the veil of shame was torn apart.

\subsection{25}\label{section-3144}

\begin{otherlanguage}{persian}

و حال محبوب را در محلّی حبس نموده‌ایم و در کلّ حین از سهام ظنون و رماح
اوهام بقدر وسع و قوّه بر او دریغ نمیداریم چه که محبوس و مسجون و فریدش
یافته‌ایم غفلت بمقامی رسیده محلّی را که جمیع اهل ملأ اعلی باو ناظرند و از
او مستمدّ در آن محلّ مبارک جهرة باقبح کلمات ناطقیم و باضلّ اعمال و اخسر آن
عامل فتبّاً لنا ما اثّرت فینا کلمات اللّه و مع‌ذلک متنبّه نشدیم و اقلّاً اگر وفا
ننمودیم جفا ننمائیم نه شبی بذکری ذاکر و نه به توجّهی مشغول

\end{otherlanguage}

Va ḥāl Maḥbūb rā dar maḥallī ḥabs namūda-īm va dar kull-i ḥīn az sihām-i
ẓunūn va rimāḥ-i awhām bi-qadr-i vusʿ va quvva bar ū darīgh namīdārīm,
chi ki maḥbūs va masjūn va farīdash yāfta-īm. Ghaflat bi-maqāmī rasīda,
maḥallī rā ki jamīʿ-i ahl-i malaʾ-i aʿlā bi-ū nāẓirand va az ū
mustamidd, dar ān maḥall-i mubārak jahratan bi-aqbaḥ-i kalimāt nāṭiqīm
va bi-aḍall-i aʿmāl va akhsar-i ān ʿāmil. Fa-tabban lanā, mā aththarat
fīnā kalimātu Allāh. Va maʿa-dhālik mutanabbih nashudīm va aqallan agar
vafā nanamūdīm jafā nanamāyīm. Na shabī bi-dhikrī dhākir va na bi
tavajjuhī mashghūl.

Now we have confined the Beloved in a place and at every moment spare
Him none of the arrows of conjecture and spears of delusion that our
capacity and strength allow, for we have found Him confined, imprisoned,
and alone. Heedlessness has reached such a point that, in that blessed
place toward which all the inhabitants of the heavenly Concourse look
and from which they seek assistance, we openly utter the foulest words
and commit the most misguided and ruinous deeds. Woe to us! The words of
God have had no effect upon us. Yet we have not awakened. At the very
least, if we have shown no faithfulness, let us not inflict cruelty. Not
one night do we engage in remembrance or occupy ourselves with turning
toward Him.

\subsection{26}\label{section-3145}

\begin{otherlanguage}{persian}

حکایت کنند که فضیل خراسانی کان من اشقی العباد و یقطع الطّریق انّه عشق
جاریة و اتاها لیلة فصعد الجدار اذاً سمع احداً یقرأ هذه الآیة أ لم یأن
للّذین آمنوا ان تخشع قلوبهم لذکر اللّه و اثّر فی قلبه فقال بلی یا ربّی آن و
حان فرجع و تاب و قصد بیت اللّه الحرام و اقام فیه ثلاثین سنة الی ان صعد
روحه الی الأفق الأعلی عجب است که کلمة اللّه را از لسان یکی از عباد شنید و
چنان مؤثّر افتاد که در یک آن از حضیض امکان بافق رحمن راجع شد و این عباد
در لیالی و ایّام متتابعاً متوالیاً نغمات نفس رحمانی را از لسان قدرت و عظمت
استماع نمودیم و این قدر تأثیر ننموده که اقلّاً بقبح افعال و اعمال و ظنون و
اوهام خود مطّلع شویم

\end{otherlanguage}

Ḥekāyat konand keh Fuḍayl-e Khurāsānī kāna min ashqā al-ʿibād wa-yaqṭaʿu
al-ṭarīq, innahu ʿashiqa jāriyatan wa-atāhā laylatan fa-ṣaʿida al-jidār,
idhan samiʿa aḥadan yaqraʾu hādhihi al-āya: a-lam yaʾni li-lladhīna
āmanū an takhshaʿa qulūbuhum li-dhikri Allāh. Wa-aththara fī qalbihi
fa-qāla balā yā rabbī āna wa-ḥāna. Fa-rajaʿa wa-tāba wa-qaṣada bayta
Allāh al-ḥarām wa-aqāma fīhi thalāthīna sanatan ilā an ṣaʿida rūḥuhu ilā
al-ufuq al-aʿlā. ʿAjab ast keh kalameh-ye Allāh rā az lesān-e yekī az
ʿebād shenīd va chenān moʾathther oftād keh dar yek ān az ḥażīż-e emkān
be-ofoq-e Raḥmān rājeʿ shod va īn ʿebād dar layālī va ayyām motatābeʿan
motavāliyan naghamāt-e nafs-e Raḥmānī rā az lesān-e qodrat va ʿaẓamat
estemāʿ nemūdīm va īn qadr taʾthīr nanemūdeh keh aqallan be-qobḥ-e afʿāl
va aʿmāl va ẓonūn va awhām-e khod moṭṭaleʿ shavīm.

It is related that Fuḍayl of Khurāsān was among the most wretched of
servants and practiced highway robbery. He fell in love with a young
woman and came to her one night. He climbed the wall and then heard
someone reciting this verse: ``Has the time not come for the hearts of
those who believe to be humbled at the remembrance of God?'' It moved
his heart, and he said, ``Yes, my Lord, the time has come; the hour has
arrived.'' He turned back, repented, and set out for the Sacred House of
God, where he remained for thirty years, until his spirit ascended to
the most exalted horizon. It is astonishing that he heard the Word of
God from the lips of one of His servants, and it affected him so deeply
that in a single instant he turned from the depths of contingent
existence toward the horizon of the All-Merciful. Yet we servants have
heard, day after day and night after night, the melodies of the
All-Merciful Himself from the Tongue of power and grandeur, and they
have not affected us enough for us even to recognize the ugliness of our
own deeds and actions, suppositions and delusions.

\subsection{27}\label{section-3146}

\begin{otherlanguage}{arabic}

حکایت که باین مقام رسید فلک بیان بر جودیّ لسان منصعق و مدهوش و قلم از
بنان منقطع و بیهوش لن یصیبنا الّا ما کتب اللّه لنا نسأل اللّه بأن یفتح
ابصارنا و یعرّفنا انفسنا و اعمالنا ان لم نوفّق علی معرفة نفسه سبحانه نوفّق
علی معرفة انفسنا الغافلة و یکشف عن وجوهنا حجبات المانعة لنراه مشرقاً عن
افق الفضل و العنایة و ننقطع عمّن فی الامکان و الأکوان و نتوجّه الیه بکلّنا
انّه هو ربّنا الرّحمن و نسأله بأن یوفّقنا علی التّوبة و الانابة فی کلّ صباح و
مسآء و یحفظ مشرق اشراق انوار وجهه عن رماح احبّائه لأنّه ما اراد معیناً سواه
و انّه مبدأه و مثواه و یقول لا اله الّا اللّه انتهی

\end{otherlanguage}

Ḥekāyat keh be-īn maqām rasīd, folk-e bayān bar Jūdiyy-e lesān monṣaʿeq
va madhūsh va qalam az banān monqaṭeʿ va bīhūsh. Lan yuṣībanā illā mā
kataba Allāhu lanā. Nasʾalu Allāha bi-an yaftaḥa abṣāranā wa-yuʿarrifanā
anfusanā wa-aʿmālanā. In lam nuwaffaq ʿalā maʿrifati nafsihi subḥānahu
nuwaffaq ʿalā maʿrifati anfusinā al-ghāfila, wa-yakshifa ʿan wujūhinā
ḥujubāt al-māniʿa li-narāhu mushriqan ʿan ufuq al-faḍl wa-l-ʿināya
wa-nanqaṭiʿa ʿamman fī al-imkān wa-l-akwān wa-natawajjaha ilayhi
bi-kullinā. Innahu huwa rabbunā al-Raḥmān. Wa-nasʾaluhu bi-an
yuwaffiqanā ʿalā al-tawba wa-l-ināba fī kulli ṣabāḥin wa-masāʾ
wa-yaḥfaẓa mashriqa ishrāqi anwāri wajhihi ʿan rimāḥi aḥibbāʾihi
li-annahu mā arāda muʿīnan siwāhu wa-innahu mabdaʾuhu wa-mathwāhu
wa-yaqūlu lā ilāha illā Allāh. Intahā.

When the account reached this point, the ark of utterance lay stunned
and bewildered upon the Judi of the tongue, and the pen fell senseless
from the fingers. ``Nothing will befall us except what God has ordained
for us.'' We ask God to open our eyes and acquaint us with ourselves and
our deeds. If we are not enabled to know His own Self---glory be to
Him---may we be enabled to know our heedless selves. May He lift the
obstructing veils from our faces, that we may see Him shining from the
horizon of grace and loving-kindness, detach ourselves from all who
dwell in the realm of possibility and the worlds of being, and turn
toward Him with our whole being. He is indeed our Lord, the
All-Merciful. We ask Him to enable us to repent and turn back to Him
every morning and evening, and to protect the Dawning-place of the
radiance of His countenance from the spears of His loved ones, for He
has desired no helper but Him, and He is His origin and His abode; and
He says, ``There is no God but God.'' End.

\subsection{28}\label{section-3147}

\begin{otherlanguage}{persian}

اگرچه بعموم نازل شده ولکن فرمودند مخصوص است بنفوس معرضه که حال در ظاهر
ادّعای حبّ مینمایند و بعد ما فی سرّهم ظاهر خواهد شد و از برای عدّه‌ئی از احّبا
قرائت شد و از خود او مستور بود جمیع ناس دانسته و میدانند که بعضی از
مدّعیان محبّت رحمانیّه در عراق و مدن اخری باعمال ناشایسته مشغول بودند ولکن
رحمت رحمانیّه مانع از خرق حجاب و کشف اعمال بود چنانچه آن نفس ظالم در ایّام
توقّف عراق سرّاً بشرب خمر و بعضی از اعمال دیگر مشغول بود و در اواخر ایّام
عراق گاهی از سرّ بجهر تجاوز می‌نمود و مع‌ذلک طائفین حول ستر مینمودند مع
آنکه متّصلاً از مطلع آیات الهیّه کلمات نصحیّه نازل بود چنانچه در اکثری از
الواح منزله مذکور است

\end{otherlanguage}

Agarcheh be-ʿomūm nāzel shodeh va-līkan farmūdand makhṣūṣ ast be-nofūs-e
moʿreżeh keh ḥāl dar ẓāher eddeʿā-ye ḥobb mīnamāyand va baʿd mā fī
serr-e-shān ẓāher khwāhad shod va az barā-ye ʿeddeh-ī az aḥebbā qerāʾat
shod va az khod-e ū mastūr būd. Jamīʿ-e nās dānesteh va mīdānand keh
baʿżī az moddaʿiyān-e maḥabbat-e Raḥmāniyyeh dar ʿIrāq va modon-e okhrā
be-aʿmāl-e nāshāyesteh mashghūl būdand va-līkan raḥmat-e Raḥmāniyyeh
māneʿ az kharq-e ḥejāb va kashf-e aʿmāl būd chenāncheh ān nafs-e ẓālem
dar ayyām-e tavaqqof-e ʿIrāq serran be-shorb-e khamr va baʿżī az aʿmāl-e
dīgar mashghūl būd va dar avākher-e ayyām-e ʿIrāq gāhī az serr be-jahr
tajāvoz mīnemūd va maʿa-dhālek ṭāʾefīn-e ḥawl satr mīnemūdand maʿ-e
ānkeh mottaṣelan az maṭlaʿ-e āyāt-e elāhiyyeh kalemāt-e noṣḥiyyeh nāzel
būd chenāncheh dar aktharī az alvāḥ-e monazzaleh madhkūr ast.

Although it was revealed in general terms, He said that it applied
specifically to those who had turned away, who now outwardly profess
love, but whose inner secrets would later be revealed. It was read to a
number of the loved ones, but was kept from the man himself. Everyone
has known and knows that some who professed love for the All-Merciful
were engaged in unworthy deeds in Iraq and other cities. Yet the mercy
of the All-Merciful prevented the veil from being torn and their deeds
from being exposed. Thus, during the sojourn in Iraq, that unjust soul
secretly drank wine and engaged in certain other deeds, and toward the
end of the days in Iraq he sometimes passed from secrecy to open
conduct. Nevertheless, those who circled around Him concealed it, even
though words of counsel were continually being revealed from the
Dawning-place of divine verses, as recorded in most of the revealed
Tablets.

\subsection{29}\label{section-3148}

\begin{otherlanguage}{persian}

و اعمال و افعال بعضی از اهل بیان در اوّل امر بر احدی پوشیده نبوده و نیست
ولکن همچو میدانستند که این امور مقبول است و از جانب حقّ تعالی شأنه
مأذونند مع آنکه بساط اوامر حقّ مقدّس از اعمال ردّیّۀ غیر مرضیّه بوده و ساحت
اقدسش منزّه از شئونات مکرهۀ نالایقه ونفسه الحقّ که در سنین توقّف عراق و
ایّام هجرت لیلاً و نهاراً از سماء اراده آیات واضحۀ صریحه در نهی عباد از غیر
ما حکم به الله نازل و باطراف ارسال میشد تا آنکه الحمد لله بعضی از عباد
به ما اراد به الله عارف شدند و عامل گشتند و نفحۀ تقدیس و تنزیه امریّۀ
الهیّه مابین بریّه مرور نمود نظر بعدم اطّلاع ناس از اوامر الهیّه تفضّلاً لهم
از لسان احدیّه این کلمه استماع شد عفا الله عمّا سلف نسأل الله ان یوفّقهم
علی التّقدیس و التّنزیه و العمل بما امروا به من لدن علیم خبیر

\end{otherlanguage}

Va aʿmāl va afʿāl-e baʿżī az ahl-e Bayān dar avval-e amr bar aḥadī
pūshīdeh nabūdeh va nīst va-līkan hamcho mīdānestand keh īn omūr maqbūl
ast va az jāneb-e Ḥaqq taʿālā shaʾnuhu maʾdhūnand maʿ-e ānkeh besāṭ-e
avāmer-e Ḥaqq moqaddas az aʿmāl-e radiyyeh-ye gheyr-e marżiyyeh būdeh va
sāḥat-e aqdas-ash monazzah az shoʾūnāt-e mokraheh-ye nālāyeqeh.
Wa-nafsihi al-Ḥaqq keh dar senīn-e tavaqqof-e ʿIrāq va ayyām-e hejrat
laylan va nahāran az samā-ye erādeh āyāt-e vāżeḥeh-ye ṣarīḥeh dar nahy-e
ʿebād az ghayri mā ḥakama bihi Allāh nāzel va be-aṭrāf ersāl mīshod tā
ānkeh al-ḥamdu li-llāh baʿżī az ʿebād be-mā arāda bihi Allāh ʿāref
shodand va ʿāmel gashtand va nafḥeh-ye taqdīs va tanzīh-e amriyeh-ye
elāhiyyeh mā-bayn-e bariyyeh morūr nemūd. Naẓar be-ʿadam-e eṭṭelāʿ-e nās
az avāmer-e elāhiyyeh tafaḍḍulan lahum az lesān-e aḥadiyyeh īn kalameh
estemāʿ shod: ʿafā Allāhu ʿammā salaf. Nasʾalu Allāha an yuwaffiqahum
ʿalā al-taqdīs wa-l-tanzīh wa-l-ʿamali bi-mā umirū bihi min ladun
ʿAlīmin Khabīr.

The deeds and actions of some of the people of the Bayān at the
beginning of the Cause were not, and are not, hidden from anyone. Yet
they imagined that such things were acceptable and that they were
authorized by God, exalted be His glory, even though the domain of God's
commandments has been sanctified from vile and unacceptable deeds, and
His most holy court is free of repugnant and unworthy conduct. By His
own Self, the True One! During the years of the sojourn in Iraq and the
days of migration, clear and explicit verses forbidding the servants to
do anything other than what God had decreed were revealed day and night
from the heaven of His will and sent to every region. At last, praise be
to God, some of the servants came to recognize what God intended and
acted upon it, and the breeze of sanctity and purity belonging to God's
Cause passed among His creatures. In view of the people's lack of
knowledge of the divine commandments, this utterance was heard from the
Tongue of Oneness as a gracious favor to them: ``God has forgiven what
is past.'' We ask God to enable them to attain sanctity and purity and
to act according to what they have been commanded by One Who is
All-Knowing, All-Informed.

\subsection{30}\label{section-3149}

\begin{otherlanguage}{persian}

باری نفس ظالم و رضا قلی لازال سرّاً باعمال شنیعه مشغول بودند و این عبد
بکرّات هر دو را نصیحت مینمود چه در عراق و چه در ارض سرّ و چه در سجن اعظم
که شاید متنبّه شوند و از خواب غفلت منتبه گردند ابداً تأثیری ننمود ما تأثّرت
فیها کلمات النّصحیّة و بیانات الشّافیة المعلنة تا آنکه رضا قلی با بعضی از
نصاری جهرة بشرب و اعمال شنیعه مشغول گشت لذا اطرده الله بسلطانه و با
خبیثین هم یعنی سیّد محمّد و آقا جان متّصل شد

\end{otherlanguage}

Bārī nafs-e ẓālem va Riḍā-Qulī lā yazāl serran be-aʿmāl-e shanīʿeh
mashghūl būdand va īn ʿabd be-karrāt har do rā naṣīḥat mīnemūd cheh dar
ʿIrāq va cheh dar Arḍ-e Serr va cheh dar Sejn-e Aʿẓam keh shāyad
motanabbeh shavand va az khwāb-e ghaflat montabeh gardand. Abadan
taʾthīrī nanemūd. Mā taʾaththarat fīhā kalimāt al-nuṣḥiyya wa-bayānāt
al-shāfiya al-muʿlana tā ānkeh Riḍā-Qulī bā baʿżī az naṣārā jahratan
be-shorb va aʿmāl-e shanīʿeh mashghūl gasht. Ledhā aṭradahu Allāhu
bi-sulṭānihi va bā khabīthayn ham yaʿnī Sayyed Muḥammad va Āqā Jān
mottaṣel shod.

In short, that unjust soul and Riḍā-Qulī were continually engaged in
shameful deeds in secret. This servant repeatedly counseled them both,
whether in Iraq, in the Land of Mystery, or in the Most Great Prison, in
the hope that they might take heed and awaken from the sleep of
heedlessness. It had no effect whatsoever. Neither the words of counsel
nor the openly proclaimed healing utterances had any effect upon them,
until Riḍā-Qulī openly engaged in drinking and shameful deeds with some
Christians. Therefore God expelled him by His authority, and he joined
the two wicked men, namely Sayyid Muḥammad and Āqā Jān.

\subsection{31}\label{section-3150}

\begin{otherlanguage}{arabic}

چندی از میان گذشت یومی عریضه‌ بساحت اقدس فرستاد مشعر به توبه از خطیئات
ماقبل و همچنین چند یوم بعد عریضۀ اخری ولکن چون بکرّات توبه نموده و نقض
عهد و میثاق الله از او ظاهر لذا عرایض مقبول نیفتاد و بعضی اعمال از او
ظاهر که شبه و مثل نداشته و سبب تضییع امر الله بین عباد گشته این عبد فانی
دوست نداشته که ذکر نماید شاید که ملائکۀ ناشرات اعمال آن خبیث مردود را در
ارض انتشار دهند الا لعنته الله علی القوم الظّالمین و در اتّصال آن خبیث و
خبیثین اشتعلت نار الفتنة و البغضآء علی شأن لا یحصیه القلم و البیان ولکنّ
الرّحمن اخمدها و اطفأها انّه لهو المقتدر القدیر

\end{otherlanguage}

Chandī az miyān gozasht, yawmī ʿarīżeh be-sāḥat-e aqdas ferestād moshʿer
be-tawbeh az khaṭīʾāt-e mā-qabl va hamchenīn chand yawm baʿd ʿarīżeh-ye
okhrā va-līkan chon be-karrāt tawbeh nemūdeh va naqż-e ʿahd va mīthāq-e
Allāh az ū ẓāher ledhā ʿarāyeż maqbūl nayoftād va baʿżī aʿmāl az ū ẓāher
keh shebh va methl nadāshteh va sabab-e tażyīʿ-e Amr-e Allāh bayn-e
ʿebād gashteh. Īn ʿabd-e fānī dūst nadāshteh keh dhekr namāyad shāyad
keh malāʾekeh-ye nāsherāt aʿmāl-e ān khabīth-e mardūd rā dar arż
enteshār dehand. A-lā laʿnatuhu Allāh ʿalā al-qawm al-ẓālimīn. Va dar
etteṣāl-e ān khabīth va khabīthayn ishtaʿalat nār al-fitna wa-l-baghḍāʾ
ʿalā shaʾnin lā yuḥṣīhi al-qalam wa-l-bayān wa-lākinna al-Raḥmāna
akhmadahā wa-aṭfaʾahā innahu la-huwa al-Muqtadir al-Qadīr.

Some time passed. One day he sent a petition to the most holy court
expressing repentance for his past sins, and a few days later another
petition of the same kind. But since he had repeatedly repented and then
broken God's Covenant and Testament, the petitions were not accepted.
Certain deeds also came to light that were without parallel or likeness
and brought the Cause of God into disrepute among His servants. This
mortal servant has not wished to mention them; perhaps the angels who
spread abroad will disseminate the deeds of that wicked outcast
throughout the earth. Truly, God's curse rests upon the unjust! When
that wicked man joined the other two, the fire of discord and hatred
blazed to a degree that pen and utterance cannot recount. But the
All-Merciful subdued and extinguished it. He is indeed the All-Powerful,
the Almighty.

\subsection{32}\label{section-3151}

\begin{otherlanguage}{persian}

هر ذی بصری از طرد جمال قدم آن نفوس را عرف قمیص تقدیس و تنزیه استنشاق
مینماید چه که اعمال رضا قلی و آقا جان خبیث و ذنبه نزد کلّ واضح و مشهود
است نسأل الله بأن یفتح عیون النّاس و یعرّفهم من اعماله و احکامه جلّ و عزّ ما
یدّل علی تقدیس ذاته و تنزیه نفسه و ساذج امره اگرچه اکثری از ناس نائمند
ولکن بعضی از انفس زکیّه و ابصر حدیده بعنایت رحمانیّه موجود انّهم یمیّزون
الحقّ عن الباطل و یعرفون من عرف کلّ امر ما تطمئنّ به قلوبهم و نفوسهم الا
انّهم عباد مکرمون

\end{otherlanguage}

Har dhī-baṣarī az ṭard-e Jamāl-e Qedam ān nofūs rā ʿarf-e qamīṣ-e taqdīs
va tanzīh estenshāq mīnamāyad cheh keh aʿmāl-e Riḍā-Qulī va Āqā Jān-e
khabīth va dhanab-ash nazd-e koll vāżeḥ va mashhūd ast. Nasʾalu Allāha
bi-an yaftaḥa ʿuyūna al-nās wa-yuʿarrifahum min aʿmālihi wa-aḥkāmihi
jalla wa-ʿazza mā yadullu ʿalā taqdīsi dhātihi wa-tanzīhi nafsihi
wa-sādhaji amrihi. Agarcheh aktharī az nās nāʾemand va-līkan baʿżī az
anfos-e zakiyyeh va abṣor-e ḥadīdeh be-ʿenāyat-e Raḥmāniyyeh mawjūd.
Innahum yumayyizūna al-ḥaqqa ʿan al-bāṭil wa-yaʿrifūna min ʿarfi kulli
amrin mā taṭmaʾinnu bihi qulūbuhum wa-nufūsuhum. A-lā innahum ʿibādun
mukramūn.

Everyone endowed with sight inhales the fragrance of the garment of
sanctity and purity from the Ancient Beauty's expulsion of those souls,
for the deeds of Riḍā-Qulī, the wicked Āqā Jān, and his follower are
clear and manifest to all. We ask God to open the eyes of the people and
acquaint them, through His actions and ordinances---glorified and
exalted is He---with that which testifies to the holiness of His
Essence, the purity of His Self, and the unalloyed nature of His Cause.
Although most people are asleep, through the loving-kindness of the
All-Merciful there are some pure souls and keen eyes. They distinguish
truth from falsehood and discern in the fragrance of every matter that
which brings assurance to their hearts and souls. Truly, they are
honored servants.

\subsection{33}\label{section-3152}

\begin{otherlanguage}{persian}

باری امر بمقامی رسید که جمال قدم در بیت عاکف و امر بسدّ باب از کلّ فرمودند
و بهیچوجه ملاقات از برای احدی ممکن نه و بر حسب ظاهر آنچه از مفتریات
مشرکین و حیل مبغضین استماع میشد ابداً از مطلع اوامر الهیّه امری ظاهر نه

\end{otherlanguage}

Bārī amr be-maqāmī rasīd keh Jamāl-e Qedam dar bayt ʿākef va amr
be-sadd-e bāb az koll farmūdand va be-hīch-vajh molāqāt az barā-ye aḥadī
momken nah va bar ḥasab-e ẓāher āncheh az moftarayāt-e moshrekīn va
ḥiyal-e mobgheżīn estemāʿ mīshod abadan az maṭlaʿ-e avāmer-e elāhiyyeh
amrī ẓāher nah.

In short, matters reached the point where the Ancient Beauty remained
secluded in the house and ordered that the door be closed to everyone.
It was in no way possible for anyone to meet Him, and, outwardly,
despite all that was heard of the fabrications of those who associated
partners with God and the schemes of those who bore Him hatred, no
directive whatsoever issued from the Dawning-place of divine commands.

\subsection{34}\label{section-3153}

\begin{otherlanguage}{persian}

تا آنکه یومی از ایّام در شهر رجب او شعبان مکتوبی از جناب آقا سیّد ١١٠
{[}علی{]} قبل اکبر ابن اخ جناب آقائی اسم ‌الله م‌ه علیه من کلّ بهآء ابهاه
لدی العرش حاضر بعد از عرض ما فی المکتوب لدی الوجه جواب لوح امنع اقدس
الّذی جعله الله رحمة للمخلصین و نقمة للمشرکین از مطلع بیان رحمن نازل بعد
از تنزیل آن لوح مبارک افق این ارض تغییر نمود و بحمرة تمام ظاهر و در هر
یوم در ازدیاد بود تا آنکه یومی از ایّام آیات عنایت از مطلع رحمت نازل بعد
از استماع آن این عبد مطمئن شد که الحمد لله که غضب الهی مخصوص نفوس خبیثه
بوده بعضی از آیات آن لوح مبارک در شوق و اشتیاق عشّاق بوده و بعضی در قهر و
اقتهار اهل نفاق و صورت آن لوح مبارک این است

\end{otherlanguage}

Tā ānkeh yawmī az ayyām dar shahr-e Rajab aw Shaʿbān maktūbī az Janāb-e
Āqā Sayyed 110 {[}ʿAlī{]} qabl-e Akbar ebn-e akh-e Janāb-e Āqāyī Esm-e
Allāh M-H ʿalayhi min kulli bahāʾin abhāhu ladā al-ʿarsh ḥāżer. Baʿd az
ʿarż-e mā fī al-maktūb ladā al-wajh javāb lawḥ-e amnaʿ-e aqdas alladhī
jaʿalahu Allāhu raḥmatan li-l-mukhliṣīn wa-niqmatan li-l-mushrikīn az
maṭlaʿ-e bayān-e Raḥmān nāzel. Baʿd az tanzīl-e ān lawḥ-e mobārak ofoq-e
īn arż taghyīr nemūd va be-ḥomrat-e tamām ẓāher va dar har yawm dar
ezdiyād būd tā ānkeh yawmī az ayyām āyāt-e ʿenāyat az maṭlaʿ-e raḥmat
nāzel. Baʿd az estemāʿ-e ān īn ʿabd moṭmaʾenn shod keh al-ḥamdu li-llāh
keh ghażab-e elāhī makhṣūṣ-e nofūs-e khabītheh būdeh. Baʿżī az āyāt-e ān
lawḥ-e mobārak dar shawq va eshtiyāq-e ʿoshshāq būdeh va baʿżī dar qahr
va eqtehār-e ahl-e nefāq va ṣūrat-e ān lawḥ-e mobārak īn ast.

Then one day, in the month of Rajab or Shaʿbān, a letter from Āqā Sayyid
110 {[}ʿAlī{]} before Akbar, nephew of Āqā, the Name of God, M-H---upon
him be the most glorious of all glory---arrived before the Throne. After
the contents of the letter were presented before His countenance, a most
inaccessible and holy Tablet, which God made a mercy to the sincere and
a punishment to those who associate partners with Him, was revealed in
reply from the Dawning-place of the All-Merciful's utterance. After that
blessed Tablet was revealed, the horizon of this land changed and
appeared intensely red, and this increased each day, until one day
verses of loving-kindness were revealed from the Dawning-place of mercy.
Upon hearing them, this servant was reassured that, praise be to God,
the divine wrath was directed specifically at the wicked souls. Some
verses of that blessed Tablet concerned the longing and yearning of the
lovers, and others the subjugation and abasement of the people of
hypocrisy. The text of that blessed Tablet is as follows.

\subsection{35}\label{section-3154}

\begin{otherlanguage}{arabic}

قوله جلّ کبریائه ان یا اسمی مهدی قد حضر لدی الوجه ما انشأه ابن اخیک فی
ثنآء مولاه و عرفنا منه الشّوق و الاشتیاق و نزّلنا له ما یحدث به الشّغف و
الاحتراق فی حبّ اللّه مالک یوم الطّلاق طوبی لمن یقرأ و یتفکّر فیما نزّل من
لدی اللّه المقتدر القدیر

\end{otherlanguage}

Qawluhu jalla kibriyāʾuhu: an yā ismī Mahdī qad ḥaḍara ladā al-wajh mā
anshaʾahu ibnu akhīka fī thanāʾi mawlāhu wa-ʿarafnā minhu al-shawqa
wa-l-ishtiyāq wa-nazzalnā lahu mā yaḥduthu bihi al-shaghaf wa-l-iḥtirāq
fī ḥubbi Allāh Māliki yawm al-ṭalāq. Ṭūbā li-man yaqraʾu wa-yatafakkaru
fī-mā nuzzila min ladā Allāh al-Muqtadir al-Qadīr.

His words---exalted be His majesty: ``O My name, Mahdī! What your nephew
composed in praise of his Lord has come before Our countenance. We
perceived in it his longing and yearning, and We have revealed for him
that which will kindle ardor and burning love for God, the Lord of the
Day of Separation. Blessed is the one who reads and reflects upon what
has been sent down from God, the All-Powerful, the Almighty.''

\subsection{36}\label{section-3155}

\begin{otherlanguage}{arabic}

بسم الله الأقدم الأعظم

\end{otherlanguage}

Bi-smi Allāh al-Aqdam al-Aʿẓam.

In the name of God, the Most Ancient, the Most Great.

\subsection{37}\label{section-3156}

\begin{otherlanguage}{arabic}

قد احترق المخلصون من نار الفراق این تشعشع انوار لقائک یا محبوب العالمین
قد ترک المقرّبون فی ظلمآء الهجران این اشراق صبح وصالک یا مقصود العالمین
قد تبلبل اجساد الأصفیآء علی ارض البعد این بحر قربک یا جذّاب العالمین قد
ارتفعت ایادی الرّجآء الی سمآء الفضل و العطآء این امطار کرمک یا مجیب
العالمین قد قام المشرکون بالاعتساف فی کلّ الأطراف این تسخیر قلم تقدیرک یا
مسخّر العالمین قد ارتفع نباح الکلاب من کلّ الجهات این غضنفر غیاض سطوتک یا
قهّار العالمین قد اخذت البرودة کلّ البریّة این حرارة محبّتک یا نار العالمین
قد بلغت البلیّة الی الغایة این ظهورات فرجک یا فرج العالمین قد احاطت
الظّلمة اکثر الخلیقة این انوار ضیائک یا ضیآء العالمین قد طالت الأعناق
بالنّفاق این اسیاف انتقامک یا مهلک العالمین قد بلغت الذّلّة الی النّهایة این
آیات عزّتک یا عزّ العالمین قد اخذت الأحزان مطلع اسمک الرّحمن این سرور مظهر
ظهورک یا فرح العالمین قد اخذ الهمّ کلّ الأمم این اعلام ابتهاجک یا بهجة
العالمین

\end{otherlanguage}

Qad iḥtaraqa al-mukhliṣūna min nār al-firāq. Ayna tashaʿshuʿu anwāri
liqāʾika yā Maḥbūb al-ʿālamīn? Qad turika al-muqarrabūna fī ẓulmāʾ
al-hijrān. Ayna ishrāqu ṣubḥi wiṣālika yā Maqṣūd al-ʿālamīn? Qad
tabalbal ajsādu al-aṣfiyāʾ ʿalā arḍ al-buʿd. Ayna baḥru qurbika yā
Jadhdhāb al-ʿālamīn? Qad irtafaʿat ayādī al-rajāʾ ilā samāʾ al-faḍl
wa-l-ʿaṭāʾ. Ayna amṭāru karamika yā Mujīb al-ʿālamīn? Qad qāma
al-mushrikūna bi-l-iʿtisāf fī kull al-aṭrāf. Ayna taskhīru qalami
taqdīrika yā Musakhkhir al-ʿālamīn? Qad irtafaʿa nubāḥu al-kilāb min
kull al-jihāt. Ayna ghaḍanfaru ghiyāḍi saṭwatika yā Qahhār al-ʿālamīn?
Qad akhadhat al-burūdatu kulla al-bariyya. Ayna ḥarāratu maḥabbatika yā
Nār al-ʿālamīn? Qad balaghat al-baliyyatu ilā al-ghāya. Ayna ẓuhūrātu
farajika yā Faraj al-ʿālamīn? Qad aḥāṭat al-ẓulmatu akthara al-khalīqa.
Ayna anwāru ḍiyāʾika yā Ḍiyāʾ al-ʿālamīn? Qad ṭālat al-aʿnāqu
bi-l-nifāq. Ayna asyāfu intiqāmika yā Muhlik al-ʿālamīn? Qad balaghat
al-dhillatu ilā al-nihāya. Ayna āyātu ʿizzatika yā ʿIzz al-ʿālamīn? Qad
akhadhat al-aḥzānu maṭlaʿa ismika al-Raḥmān. Ayna surūru maẓhari
ẓuhūrika yā Faraḥ al-ʿālamīn? Qad akhadha al-hammu kulla al-umam. Ayna
aʿlāmu ibtihājika yā Bahjat al-ʿālamīn?

The sincere have burned in the fire of separation. Where is the radiance
of the lights of Your presence, O Beloved of the worlds? Those brought
near have been left in the darkness of estrangement. Where is the
dawning of the morning of reunion with You, O Desire of the worlds? The
bodies of the chosen ones have been thrown into turmoil in the land of
remoteness. Where is the ocean of Your nearness, O Magnet of the worlds?
The hands of hope have been raised toward the heaven of grace and
bounty. Where are the showers of Your generosity, O Answerer of the
worlds? Those who associate partners with You have risen in tyranny on
every side. Where is the subduing power of Your pen of destiny, O
Subduer of the worlds? The barking of dogs has risen from every
direction. Where is the lion of the thickets of Your might, O
All-Conquering One of the worlds? Coldness has seized all creation.
Where is the heat of Your love, O Fire of the worlds? Affliction has
reached its utmost limit. Where are the manifestations of Your
deliverance, O Deliverance of the worlds? Darkness has encompassed most
of creation. Where are the lights of Your radiance, O Radiance of the
worlds? Necks have stretched high in hypocrisy. Where are the swords of
Your vengeance, O Destroyer of the worlds? Humiliation has reached its
furthest extreme. Where are the signs of Your glory, O Glory of the
worlds? Sorrows have seized the Dawning-place of Your name, the
All-Merciful. Where is the joy of the Manifestation of Your revelation,
O Joy of the worlds? Grief has seized all nations. Where are the banners
of Your gladness, O Delight of the worlds?

\subsection{38}\label{section-3157}

\begin{otherlanguage}{arabic}

تری مشرق الآیات فی سبحات الاشارات این اصبع قدرتک یا اقتدار العالمین قد
اخذت رعدة الظّمأ من فی الانشآء این فرات عنایتک یا رحمة العالمین قد اخذ
الحرص من فی الابداع این مطالع الانقطاع یا مولی العالمین تری المظلوم
فریداً فی الغربة این جند سمآء امرک یا سلطان العالمین قد ترکت وحده فی دیار
الغربة این مشارق وفائک یا وفآء العالمین قد اخذت سکرات الموت کلّ الآفاق
این رشحات بحر حیوانک یا حیوة العالمین قد احاطت وساوس الشّیطان من فی
الامکان این شهاب نارک یا نور العالمین قد تغیّر اکثر الوری من سکر الهوی
این مطالع التّقوی یا مقصود العالمین تری المظلوم فی حجاب الظّلام بین اهل
الشّام این اشراق انوار صباحک یا مصباح العالمین ترانی ممنوعاً عن البیان من
این تظهر نغماتک یا ورقآء العالمین قد غشّت الظّنون و الأوهام اکثر الأنام
این مطالع ایقانک یا سکینة العالمین قد غرق البهآء فی بحر البلآء این فلک
نجاتک یا منجی العالمین تری مطلع آیاتک فی ظلمات الامکان این شمس افق
عنایتک یا نوّار العالمین قد خبت مصابیح الصّدق و الصّفآء و الغیرة و الوفآء
این شؤونات غیرتک یا محرّک العالمین هل تری من ینصر نفسک او یتفکّر فیما ورد
علیها فی حبّک اذاً توقّف القلم یا محبوب العالمین

\end{otherlanguage}

Tarā mashriqa al-āyāti fī subuḥāti al-ishārāt. Ayna iṣbaʿu qudratika yā
iqtidāra al-ʿālamīn? Qad akhadhat raʿdatu al-ẓamaʾi man fī al-inshāʾ.
Ayna furātu ʿināyatika yā raḥmata al-ʿālamīn? Qad akhadha al-ḥirṣu man
fī al-ibdāʿ. Ayna maṭāliʿu al-inqiṭāʿi yā mawlā al-ʿālamīn? Tarā
al-Maẓlūma farīdan fī al-ghurba. Ayna jundu samāʾi amrika yā sulṭāna
al-ʿālamīn? Qad turiktu waḥdahu fī diyāri al-ghurba. Ayna mashāriqu
wafāʾika yā wafāʾa al-ʿālamīn? Qad akhadhat sakarātu al-mawti kulla
al-āfāq. Ayna rashaḥātu baḥri ḥayawānika yā ḥayāta al-ʿālamīn? Qad
aḥāṭat wasāwisu al-shayṭāni man fī al-imkān. Ayna shihābu nārika yā nūra
al-ʿālamīn? Qad taghayyara aktharu al-warā min sukri al-hawā. Ayna
maṭāliʿu al-taqwā yā maqṣūda al-ʿālamīn? Tarā al-Maẓlūma fī ḥijābi
al-ẓalāmi bayna ahli al-Shām. Ayna ishrāqu anwāri ṣabāḥika yā miṣbāḥa
al-ʿālamīn? Tarānī mamnūʿan ʿani al-bayān. Min ayna taẓharu naghamātuka
yā warqāʾa al-ʿālamīn? Qad ghashshat al-ẓunūnu wa-al-awhāmu akthara
al-anām. Ayna maṭāliʿu īqānika yā sakīnata al-ʿālamīn? Qad ghariqa
al-Bahāʾu fī baḥri al-balāʾ. Ayna fulku najātika yā munjiya al-ʿālamīn?
Tarā maṭlaʿa āyātika fī ẓulumāti al-imkān. Ayna shamsu ufuqi ʿināyatika
yā nawwāra al-ʿālamīn? Qad khabat maṣābīḥu al-ṣidqi wa-al-ṣafāʾi
wa-al-ghayrati wa-al-wafāʾ. Ayna shuʾūnātu ghayratika yā muḥarrika
al-ʿālamīn? Hal tarā man yanṣuru nafsaka aw yatafakkaru fīmā warada
ʿalayhā fī ḥubbika? Idhan tawaqqafa al-Qalamu yā maḥbūba al-ʿālamīn.

You see the Dawning-place of the verses amid the veils of allusions.
Where is the finger of Your power, O Might of the worlds? The trembling
of thirst has seized all in creation. Where is the river of Your loving
care, O Mercy of the worlds? Greed has seized all created beings. Where
are the dawning-places of detachment, O Lord of the worlds? You see the
Wronged One alone in exile. Where are the hosts of the heaven of Your
command, O Sovereign of the worlds? I have been left alone in the lands
of exile. Where are the dawning-places of Your faithfulness, O
Faithfulness of the worlds? The agonies of death have seized every
region. Where are the spray-drops of Your ocean of life, O Life of the
worlds? Satan's whisperings have encompassed all contingent beings.
Where is the meteor of Your fire, O Light of the worlds? Most people
have been changed by the intoxication of desire. Where are the
dawning-places of piety, O Desire of the worlds? You see the Wronged One
veiled in darkness among the people of Syria. Where is the radiance of
the lights of Your morning, O Lamp of the worlds? You see Me prevented
from speaking. From where will Your melodies arise, O Dove of the
worlds? Suppositions and vain imaginings have enveloped most of
humanity. Where are the dawning-places of Your certitude, O Tranquillity
of the worlds? Bahāʾ has drowned in the ocean of affliction. Where is
the ark of Your salvation, O Savior of the worlds? You see the
Dawning-place of Your verses in the darkness of contingent existence.
Where is the sun of the horizon of Your loving care, O Illuminator of
the worlds? The lamps of truthfulness, purity, ardor, and faithfulness
have gone out. Where are the manifestations of Your ardor, O Mover of
the worlds? Do You see anyone who will help You or reflect on what has
befallen You out of love for You? Here the Pen has stopped, O Beloved of
the worlds.

\subsection{39}\label{section-3158}

\begin{otherlanguage}{arabic}

قد کسرت اغصان سدرة المنتهی من هبوب اریاح القضآء این رایات نصرتک یا منصور
العالمین قد بقی الوجه فی غبار الافترآء این اریاح رحمتک یا رحمن العالمین
قد تکدّر ذیل التّقدیس من اولی التّدلیس این طراز تنزیهک یا مزیّن العالمین قد
رکد بحر العنایة بما اکتسبت ایدی البریّة این امواج فضلک یا مراد العالمین
قد غلق باب اللّقآء من ظلم الأعدآء این مفتاح جودک یا فتّاح العالمین قد
اصفرّت الأوراق من سموم اریاح النّفاق این جود سحاب جودک یا جوّاد العالمین قد
تغیّر الأکوان من غبار العصیان این نفحات غفرانک یا غفّار العالمین قد بقی
الغلام فی ارض جدبآء این غیث سمآء فضلک یا غیاث العالمین

\end{otherlanguage}

Qad kusirat aghṣānu sidrati al-muntahā min hubūbi aryāḥi al-qaḍāʾ. Ayna
rāyātu nuṣratika yā manṣūra al-ʿālamīn? Qad baqiya al-wajhu fī ghubāri
al-iftirāʾ. Ayna aryāḥu raḥmatika yā raḥmāna al-ʿālamīn? Qad takaddara
dhaylu al-taqdīsi min ulī al-tadlīs. Ayna ṭirāzu tanzīhika yā muzayyina
al-ʿālamīn? Qad rakada baḥru al-ʿināyati bimā iktasabat aydī al-bariyya.
Ayna amwāju faḍlika yā murāda al-ʿālamīn? Qad ghuliqa bābu al-liqāʾi min
ẓulmi al-aʿdāʾ. Ayna miftāḥu jūdika yā fattāḥa al-ʿālamīn? Qad iṣfarrat
al-awrāqu min samūmi aryāḥi al-nifāq. Ayna jūdu saḥābi jūdika yā jawwāda
al-ʿālamīn? Qad taghayyara al-akwānu min ghubāri al-ʿiṣyān. Ayna
nafaḥātu ghufrānika yā ghaffāra al-ʿālamīn? Qad baqiya al-Ghulāmu fī
arḍin jadbāʾ. Ayna ghaythu samāʾi faḍlika yā ghiyātha al-ʿālamīn?

The branches of the Lote Tree of the Utmost Boundary have been broken by
the blowing winds of fate. Where are the banners of Your help, O
Victorious One of the worlds? The Face remains in the dust of calumny.
Where are the winds of Your mercy, O All-Merciful One of the worlds? The
hem of holiness has been sullied by the deceitful. Where is the
adornment of Your sanctity, O Adorner of the worlds? The ocean of loving
care has grown still because of what human hands have done. Where are
the waves of Your grace, O Desire of the worlds? The gate of reunion has
been shut by the oppression of enemies. Where is the key of Your
generosity, O Opener of the worlds? The leaves have yellowed in the
scorching winds of hypocrisy. Where is the bounty of the cloud of Your
generosity, O Generous One of the worlds? All beings have been changed
by the dust of disobedience. Where are the breezes of Your forgiveness,
O Forgiver of the worlds? The Youth remains in a barren land. Where is
the rain of the heaven of Your grace, O Succor of the worlds?

\subsection{40}\label{section-3159}

\begin{otherlanguage}{persian}

ان یا قلم الأعلی قد سمعنا ندائک الأحلی من جبروت البقآء ان استمع ما ینطق
به لسان الکبریآء یا مظلوم العالمین لو لا البرودة کیف تظهر حرارة بیانک یا
مبیّن العالمین و لو لا البلیّة کیف اشرقت شمس اصطبارک یا شعاع العالمین لا
تجزع من الأشرار قد خلقت للاصطبار یا صبر العالمین ما احلی اشراقک من افق
المیثاق بین اهل النّفاق و اشتیاقک بالله یا عشق العالمین بک ارتفع علم
الاستقلال علی اعلی الجبال و تموّج بحر الافضال یا وله العالمین بوحدتک
اشرقت شمس التّوحید و بغربتک زیّن وطن التّجرید ان اصطبر یا غریب العالمین قد
جعلنا الذّلّة قمیص العزّة و البلیّة طراز هیکلک یا فخر العالمین تری القلوب
ملئت من البغضآء و لک الاغضآء یا ستّار العالمین اذا رأیت سیفاً ان اقبل اذا
طار سهم ان استقبل یا فدآء العالمین أ تنوح او انوح بل اصیح من قلّة ناصریک
یا من بک ارتفع نوح العالمین

\end{otherlanguage}

An yā Qalama al-Aʿlā, qad samiʿnā nidāʾaka al-aḥlā min jabarūti
al-baqāʾ. An istamiʿ mā yanṭiqu bihi lisānu al-kibriyāʾi yā Maẓlūma
al-ʿālamīn. Law lā al-burūdatu kayfa taẓharu ḥarāratu bayānika yā
mubayyina al-ʿālamīn? Wa-law lā al-baliyyatu kayfa ashraqat shamsu
iṣṭibārika yā shuʿāʿa al-ʿālamīn? Lā tajzaʿ min al-ashrār. Qad khuliqta
li-al-iṣṭibāri yā ṣabra al-ʿālamīn. Mā aḥlā ishrāqaka min ufuqi
al-mīthāqi bayna ahli al-nifāqi wa-ishtiyāqaka bi-Allāhi yā ʿishqa
al-ʿālamīn! Bika irtafaʿa ʿalamu al-istiqlāli ʿalā aʿlā al-jibāli
wa-tamawwaja baḥru al-ifḍāli yā walaha al-ʿālamīn. Bi-waḥdatika ashraqat
shamsu al-tawḥīdi wa-bi-ghurbatika zuyyina waṭanu al-tajrīd. An iṣṭabir
yā gharība al-ʿālamīn. Qad jaʿalnā al-dhillata qamīṣa al-ʿizzati
wa-al-baliyyata ṭirāza haykalika yā fakhra al-ʿālamīn. Tarā al-qulūba
muliʾat min al-baghḍāʾi wa-laka al-ighḍāʾu yā sattāra al-ʿālamīn. Idhā
raʾayta sayfan an aqbil. Idhā ṭāra sahmun an istaqbil yā fidāʾa
al-ʿālamīn. A-tanūḥu aw anūḥ? Bal aṣīḥu min qillati nāṣirīka yā man bika
irtafaʿa nawḥu al-ʿālamīn.

O Most Exalted Pen, We have heard Your sweetest call from the dominion
of eternity. Listen to what the Tongue of Grandeur proclaims, O Wronged
One of the worlds. Were it not for coldness, how would the heat of Your
utterance appear, O Expounder of the worlds? Were it not for affliction,
how would the sun of Your patient endurance have shone, O Radiance of
the worlds? Do not grieve over the wicked. You were created for patient
endurance, O Patience of the worlds. How sweet is Your dawning from the
horizon of the Covenant among the people of hypocrisy, and Your longing
for God, O Love of the worlds! Through You the banner of independence
has been raised upon the highest mountains, and the ocean of bounty has
surged, O Rapture of the worlds. Through Your solitude the sun of divine
unity has shone, and through Your exile the homeland of detachment has
been adorned. Be patient, O Exile of the worlds. We have made abasement
the garment of glory, and affliction the adornment of Your temple, O
Pride of the worlds. You see hearts filled with hatred, yet it is Yours
to overlook, O Concealer of the worlds. When You see a sword, advance;
when an arrow flies, go forth to meet it, O Sacrifice of the worlds.
Will You lament, or shall I? Rather, I cry aloud at the scarcity of Your
helpers, O You through whom the lamentation of the worlds has risen.

\subsection{41}\label{section-3160}

\begin{otherlanguage}{arabic}

قد سمعت ندائک یا محبوب الأبهی اذاً انار وجه البهآء من حرارة البلآء و
انوار کلمتک النّورآء و قام بالوفآء فی مشهد الفدآء ناظراً رضائک یا مقدّر
العالمین

\end{otherlanguage}

Qad samiʿtu nidāʾaka yā maḥbūba al-Abhā. Idhan anāra wajhu al-Bahāʾi min
ḥarārati al-balāʾi wa-anwāri kalimatika al-nawrāʾi wa-qāma bi-al-wafāʾi
fī mashhadi al-fidāʾi nāẓiran riḍāka yā muqaddira al-ʿālamīn.

I have heard Your call, O Beloved of the All-Glorious. Thereupon the
face of Bahāʾ shone with the heat of affliction and the lights of Your
radiant word, and He arose in faithfulness in the arena of sacrifice,
looking toward Your good pleasure, O Ordainer of the worlds.

\subsection{42}\label{section-3161}

\begin{otherlanguage}{arabic}

ان یا علیّ قبل اکبر ان اشکر الله بهذا اللّوح الّذی تجد منه رایحة مظلومیّتی و
ما انا فیه فی سبیل الله معبود العالمین لو یقرأه العباد طرّاً و یتفکّرون فیه
لیضرم فی کلّ عرق من عروقهم ناراً یشتعل منها العالمین

\end{otherlanguage}

An yā ʿAlī qabla Akbar, an ushkur Allāha bi-hādhā al-Lawḥi alladhī
tajidu minhu rāʾiḥata maẓlūmiyyatī wa-mā anā fīhi fī sabīli Allāhi
maʿbūdi al-ʿālamīn. Law yaqraʾuhu al-ʿibādu ṭurran wa-yatafakkarūna fīhi
la-yuḍrima fī kulli ʿirqin min ʿurūqihim nāran yashtaʿilu minhā
al-ʿālamīn.

O ʿAlī before Akbar, give thanks to God for this Tablet, from which you
perceive the fragrance of the wrongs I have suffered and of what I
endure in the path of God, the One worshiped by the worlds. If all the
servants were to read it and reflect upon it, it would kindle in every
vein of their bodies a fire that would set the worlds ablaze.

\subsection{43}\label{section-3162}

\begin{otherlanguage}{persian}

این عبد تا آن یوم از لسان عظمت آیات قهریّه باین شأن استماع ننموده لذا
بسیار متفکّر که چه واقع خواهد شد و چه امری از مطلع غیب ظاهر شود باری در
هر یوم فساد و اعراض اشقیا در تزاید بوده تا آنکه رضا قلی یک بستۀ کبیر
سواد نوشتجات که نزدش بود بعض فقرات آن را بفقرات کفرآمیز مخلوط نموده و
بدست اکثری از اهل این بلد داده و چنان مذکور نمود که من از حضرات بودم و
حال تائب شدم و مسلم گشتم او و آقا جان و سیّد محمّد هر سه نزد ناس اظهار
ندامت نموده و خود را از اهل اسلام قلم دادند و بقسمی این بلد منقلب شد که
اکثر ناس جهرة بشتم و بغضا قیام نمودند و سیّد محمّد خبیث چون مشاهده نمود که
جمال قدم در بیت ساکن و باب بر حسب ظاهر مسدود و حضرت غصن اعظم هم از مابین
اصحاب خارج فرصت غنیمت نموده با بعضی از اصحاب بنای مراوده و دوستی گذاشتند

\end{otherlanguage}

Īn ʿabd tā ān yawm az lisān-i ʿaẓamat āyāt-i qahriyya bi-īn shaʾn
istimāʿ nanamūda, lidhā bisyār mutafakkir ki chi vāqiʿ khwāhad shud va
chi amrī az maṭlaʿ-i ghayb ẓāhir shavad. Bārī dar har yawm fasād va
iʿrāḍ-i ashqiyā dar tazāyud būda tā ānki Riḍā Qulī yik basta-yi kabīr
savād-i nivishtajāt ki nazdash būd, baʿḍ-i faqarāt-i ān rā bi-faqarāt-i
kufr-āmīz makhlūṭ namūda va bi-dast-i aktharī az ahl-i īn balad dāda va
chunān madhkūr namūd ki man az ḥaḍarāt būdam va ḥāl tāʾib shudam va
muslim gashtam. Ū va Āqā Jān va Sayyid Muḥammad har si nazd-i nās
iẓhār-i nadāmat namūda va khwud rā az ahl-i Islām qalam dādand va
bi-qismī īn balad munqalib shud ki akthar-i nās jahratan bi-shatm va
baghḍā qiyām namūdand. Va Sayyid Muḥammad-i khabīth chūn mushāhada namūd
ki Jamāl-i Qidam dar bayt sākin va bāb bar ḥasab-i ẓāhir masdūd va
Ḥaḍrat-i Ghuṣn-i Aʿẓam ham az mābayn-i aṣḥāb khārij, furṣat ghanīmat
namūda bā baʿḍī az aṣḥāb binā-yi murāvada va dūstī gudhāshtand.

Until that day, this servant had never heard verses of wrath of such a
character from the Tongue of Grandeur. I therefore wondered greatly what
would happen and what event would emerge from the dawning-place of the
unseen. In short, the mischief and opposition of the wicked increased
daily, until Riḍā Qulī took a large bundle of copies of writings in his
possession, mixed some of their passages with blasphemous passages, and
distributed them to many of the people of this city. He declared, ``I
was one of those people, but now I have repented and become a Muslim.''
He, Āqā Jān, and Sayyid Muḥammad all publicly expressed remorse and
represented themselves as Muslims. The city was so stirred up that most
people openly resorted to abuse and hatred. When the wicked Sayyid
Muḥammad saw that the Ancient Beauty remained in the house, that to
outward appearances the door was closed, and that His Holiness the Most
Great Branch was also absent from among the companions, he seized the
opportunity and began associating and forming friendships with some of
the companions.

\subsection{44}\label{section-3163}

\begin{otherlanguage}{persian}

و این عباد ساکن و صابر ابداً در این امور امری از مصدر امر صادر نه متوکّلاً
علی الله در محلّ نشسته تا چه ظاهر شود و از سرادق غیب بعرصۀ ظهور چه جلوه
نماید تا آنکه در یوم اثنین ثانی‌عشر شهر ذی‌القعده ساعت یازده از روز غوغا
برخاست مشاهده شد که پاشای بلد با جمیع عسکریّه مع اسیاف مسلوله بیت را
احاطه نمودند و همچنین کلّ اصحاب را اخذ نمودند و این عبد بتحریر آیات منزله
مشغول بود و بقسمی مدینه در حرکت و اضطراب مشاهده شد که فوق آن ممکن نه
بغتةً کلّ اهل بلد مع عساکر و ضبّاط بهیجان آمده و بقسمی نعره و ضوضا مرتفع که
قلم از ذکر آن عاجز است

\end{otherlanguage}

Va īn ʿibād sākin va ṣābir, abadan dar īn umūr amrī az maṣdar-i amr
ṣādir na, mutavakkilan ʿalā Allāh dar maḥall nishasta tā chi ẓāhir
shavad va az surādiq-i ghayb bi-ʿarṣa-yi ẓuhūr chi jilva namāyad. Tā
ānki dar yawm-i ithnayn thānī-ʿashar-i shahr-i Dhī-al-Qaʿda sāʿat-i
yāzdah az rūz ghawghā barkhāst. Mushāhada shud ki pāshā-yi balad bā
jamīʿ-i ʿaskariyya maʿa asyāf-i maslūla bayt rā iḥāṭa namūdand va
hamchunīn kull-i aṣḥāb rā akhdh namūdand. Va īn ʿabd bi-taḥrīr-i āyāt-i
munzala mashghūl būd va bi-qismī madīna dar ḥarakat va iḍṭirāb mushāhada
shud ki fawq-i ān mumkin na. Baghtatan kull-i ahl-i balad maʿa ʿasākir
va ḍubbāṭ bi-hayajān āmada va bi-qismī naʿra va ḍawḍā murtafiʿ ki qalam
az dhikr-i ān ʿājiz ast.

These servants remained quiet and patient. No command whatever
concerning these matters issued from the Source of command. Trusting in
God, we stayed where we were, awaiting what would appear and what would
emerge from the pavilion of the unseen into the realm of manifestation.
Then, on Monday, the twelfth of the month of Dhū al-Qaʿda, at the
eleventh hour of the day, an uproar arose. We saw that the city's pasha,
together with all the soldiery bearing drawn swords, had surrounded the
house. They also arrested all the companions. This servant was occupied
with transcribing the revealed verses. The city was seen to be in such
turmoil and agitation that nothing greater was possible. Suddenly all
the townspeople, together with the soldiers and officers, were thrown
into a frenzy, and such shouting and clamor arose that the pen is
powerless to describe it.

\subsection{45}\label{section-3164}

\begin{otherlanguage}{persian}

در آن اثنا تلقاء وجه حاضر فرمودند لا تلتفت الی ضوضائهم بعد لسان مبارک
بآیات ناطق و این عبد به تحریر آن مشغول که غصن اعظم بین یدی حاضر و معروض
داشتند که بباب حکومت خواسته‌اند جمال قدم مع غصن اعظم تشریف بردند غصن اکبر
و آقائی آقا میرزا محمّد قلی و این عبد خواستیم همراه برویم منع فرمودند بعد
از ساعتی آمدند غصن اکبر و این عبد را هم بردند بعد از ورود سرایه اهل
حکومت مذکور نمودند که هفت نفر از اعجام رفته‌اند آقا جان و سیّد محمّد و رضا
قلی را کشته‌اند لله درّ من قال فی فتیة من جنود القهر

\end{otherlanguage}

Dar ān athnā tilqāʾ-i vajh ḥāḍir farmūdand: Lā taltafit ilā ḍawḍāʾihim.
Baʿd lisān-i mubārak bi-āyāt nāṭiq va īn ʿabd bi-taḥrīr-i ān mashghūl ki
Ghuṣn-i Aʿẓam bayn-i yaday ḥāḍir va maʿrūḍ dāshtand ki bi-bāb-i ḥukūmat
khwāsta-and. Jamāl-i Qidam maʿa Ghuṣn-i Aʿẓam tashrīf burdand. Ghuṣn-i
Akbar va āqāyī Āqā Mīrzā Muḥammad Qulī va īn ʿabd khwāstīm hamrāh
biravīm, manʿ farmūdand. Baʿd az sāʿatī āmadand, Ghuṣn-i Akbar va īn
ʿabd rā ham burdand. Baʿd az vurūd-i sarāya ahl-i ḥukūmat madhkūr
namūdand ki haft nafar az aʿjām rafta-and, Āqā Jān va Sayyid Muḥammad va
Riḍā Qulī rā kushta-and. Li-Allāhi darru man qāla fī fityatin min junūdi
al-qahr.

Meanwhile, He summoned me into His presence and said, ``Pay no attention
to their clamor.'' Then His blessed tongue uttered verses, and this
servant was occupied with transcribing them, when the Most Great Branch
entered His presence and submitted that He had been summoned to the seat
of government. The Ancient Beauty went with the Most Great Branch. The
Greater Branch, my master Āqā Mīrzā Muḥammad Qulī, and this servant
wished to accompany them, but He forbade us. After an hour they came and
took the Greater Branch and this servant as well. Upon our arrival at
the government building, the officials stated that seven Persians had
gone and killed Āqā Jān, Sayyid Muḥammad, and Riḍā Qulī. God bless the
one who said of a band of youths from the hosts of wrath:

\subsection{46}\label{section-3165}

\begin{otherlanguage}{arabic}

ما ترکت للرّعد کرّاتهم صوتاً و لا صیتا قوم اذا قوبلوا کانوا ملائکة حسنا و
ان قوتلوا کانوا عفاریتا

\end{otherlanguage}

Mā tarakat li-al-raʿdi karrātuhum ṣawtan wa-lā ṣītā. Qawmun idhā qūbilū
kānū malāʾikatan ḥusnan wa-in qūtilū kānū ʿafārītā.

Their onslaughts have left thunder neither voice nor renown. A people
who, when met, are angels in beauty, but when fought, are demons.

\subsection{47}\label{section-3166}

\begin{otherlanguage}{persian}

و جمیع ناس مضطرب و خائف مشاهده شدند و بشأنی آثار قهر احاطه نمود که اکثری
از اهل بلد و اهل حکومت در آن لیل از اکل و نوم ممنوع ماندند عجب در آنست
مع آنکه جمال قدم در اشهر معدودات مع غصنین اعظمین سدّ باب نموده و ابداً با
احدی از خارج و داخل معاشرت نمیفرمودند و خود سیّد محمّد باطراف نوشته که
جمیع اصحاب اعراض نمودند مع‌ذلک بعد از وقوع این امر خبیثۀ کاذبه علیها لعنة
الله و لعنة اولیائه در مجلس حکومت رفته و ذکر نموده که جمال قدم امر
فرمودند که این نفوس را بقتل رسانیدند مع آنکه والله المهیمن القیّوم که
ابداً بر حسب ظاهر از مقرّ امری صادر نشد و لازال کلّ را از ارتکاب این امورات
نهی میفرمودند و میفرمودند انّ الکلب ینبح و الذّئب یعوی ان اترکوهما و لا
تتعرّضوا بهما و کونوا من الصّابرین

\end{otherlanguage}

Va jamīʿ-i nās muḍṭarib va khāʾif mushāhada shudand va bi-shaʾnī āthār-i
qahr iḥāṭa namūd ki aktharī az ahl-i balad va ahl-i ḥukūmat dar ān layl
az akl va nawm mamnūʿ māndand. ʿAjab dar ānast maʿa ānki Jamāl-i Qidam
dar ashhur-i maʿdūdāt maʿa ghuṣnayn-i aʿẓamayn sadd-i bāb namūda va
abadan bā aḥadī az khārij va dākhil muʿāsharat namīfarmūdand va khwud-i
Sayyid Muḥammad bi-aṭrāf nivishta ki jamīʿ-i aṣḥāb iʿrāḍ namūdand,
maʿa-dhālik baʿd az vuqūʿ-i īn amr khabītha-yi kādhiba, ʿalayhā laʿnatu
Allāhi wa-laʿnatu awliyāʾih, dar majlis-i ḥukūmat rafta va dhikr namūda
ki Jamāl-i Qidam amr farmūdand ki īn nufūs rā bi-qatl rasānīdand. Maʿa
ānki wa-Allāhi al-Muhaymini al-Qayyūm ki abadan bar ḥasab-i ẓāhir az
maqarr amrī ṣādir nashud va lāzāl kull rā az irtikāb-i īn umūrāt nahy
mīfarmūdand va mīfarmūdand: Inna al-kalba yanbaḥu wa-al-dhiʾba yaʿwī. An
utrukūhumā wa-lā tataʿarraḍū bihimā wa-kūnū min al-ṣābirīn.

Everyone was seen to be agitated and afraid. The signs of wrath so
encompassed them that many of the townspeople and officials were unable
to eat or sleep that night. The astonishing thing is that, although the
Ancient Beauty had closed His door for several months together with the
two Most Great Branches, and had associated with no one, whether from
outside or within, and Sayyid Muḥammad himself had written to various
places that all the companions had turned away, nevertheless, after this
event that wicked liar---may the curse of God and of His chosen ones be
upon her---went to the government council and stated that the Ancient
Beauty had given the order by which these people were killed. Yet, by
God, the All-Protecting, the Self-Subsisting, no outward command
whatever had issued from His abode. He had always forbidden everyone to
commit such acts, saying, ``The dog barks and the wolf howls. Leave them
alone, do not interfere with them, and be patient.''

\subsection{48}\label{section-3167}

\begin{otherlanguage}{persian}

باری یکی از احبّای الهی تفصیل مجلس حکومت را از این عبد از دیار بعیده
استفسار نمود این عبد آنچه در نظر بود خدمت ایشان معروض داشت و حال در این
ورقه مکرّراً ذکر میشود تا جمیع بریّه از بیانات الهیّه در آن لیله مستفیض شوند
لیلۀ اولی حین صلوة عشا بود که در مجلس پاشا تشریف بردند و {[}در{]} آن
مجلس جمع کثیری مجتمع بوده و جمیع در کمال وحشت و اضطراب جمال قدم در صدر
مجلس متّکأً جالس و ابداً تکلّم نفرمودند چه که اهل مجلس بشورا و مکالمه بین
خود مشغول بودند و در آن مجلس جز حضرت غصن اعظم احدی را همراه نبردند بعد
پاشا برخاسته معروض داشت که در محلّ دیگر تشریف ببرید چه که حال میخواهیم
حضرات آخذین را استنطاق نمائیم لذا در محلّ دیگر که بمحبس اداره موسوم است
تشریف بردند و در آن محبس غصنین اعظمین و آقائی آقا میرزا محمّد قلی و آقا
محمّد علی اصفهانی و آقا محمّد جواد و آقا محمّد حسین ابن حاجی علی عسکر تلقاء
وجه حاضر بودیم و بعضی از اهل سرایه از ضبّاط و غیره و الی ساعة سابعه در آن
محلّ تشریف داشتند و متّصلاً بآیات عظمت ناطق از جمله باین عبد فرمودند در
نظرت هست آیاتی که در لیلۀ قبل نازل شده و بعد بأعلی البیان تلاوت فرمودند
و آن آیات این است

\end{otherlanguage}

Bārī yikī az aḥibbā-yi ilāhī tafṣīl-i majlis-i ḥukūmat rā az īn ʿabd az
diyār-i baʿīda istifsār namūd. Īn ʿabd ānchi dar naẓar būd khidmat-i
īshān maʿrūḍ dāsht va ḥāl dar īn varaqa mukarraran dhikr mīshavad tā
jamīʿ-i bariyya az bayānāt-i ilāhiyya dar ān layla mustafīḍ shavand.
Layla-yi ūlā ḥīn-i ṣalāt-i ʿishā būd ki dar majlis-i pāshā tashrīf
burdand va {[}dar{]} ān majlis jamʿ-i kathīrī mujtamiʿ būda va jamīʿ dar
kamāl-i vaḥshat va iḍṭirāb. Jamāl-i Qidam dar ṣadr-i majlis muttakiʾan
jālis va abadan takallum nafarmūdand, chi ki ahl-i majlis bi-shūrā va
mukālama bayn-i khwud mashghūl būdand. Va dar ān majlis juz Ḥaḍrat-i
Ghuṣn-i Aʿẓam aḥadī rā hamrāh naburdand. Baʿd pāshā barkhāsta maʿrūḍ
dāsht ki dar maḥall-i dīgar tashrīf bibarīd, chi ki ḥāl mīkhwāhīm
ḥaḍarāt-i ākhidhīn rā istinṭāq namāʾīm. Lidhā dar maḥall-i dīgar ki
bi-maḥbas-i idāra mawsūm ast tashrīf burdand va dar ān maḥbas ghuṣnayn-i
aʿẓamayn va āqāyī Āqā Mīrzā Muḥammad Qulī va Āqā Muḥammad ʿAlī Iṣfahānī
va Āqā Muḥammad Javād va Āqā Muḥammad Ḥusayn ibn-i Ḥājī ʿAlī ʿAskar
tilqāʾ-i vajh ḥāḍir būdīm va baʿḍī az ahl-i sarāya az ḍubbāṭ va ghayrih.
Va ilā sāʿatin sābiʿa dar ān maḥall tashrīf dāshtand va muttaṣilan
bi-āyāt-i ʿaẓamat nāṭiq. Az jumla bi-īn ʿabd farmūdand: Dar naẓarat hast
āyātī ki dar layla-yi qabl nāzil shuda? Va baʿd bi-aʿlā al-bayān tilāvat
farmūdand va ān āyāt īn ast.

In brief, one of the loved ones of God in a distant land asked this
servant for the details of the government council. This servant
submitted to him what he recalled, and it is now recounted again on this
page so that all people may benefit from the divine utterances of that
night. On the first night, it was at the time of the evening prayer that
He entered the pasha's council. A large crowd had assembled {[}in{]}
that council, and everyone was in extreme fear and agitation. The
Ancient Beauty sat reclining at the head of the assembly and did not
speak at all, for those present were occupied with consultation and
conversation among themselves. He had taken no one with Him to that
council except His Holiness the Most Great Branch. Then the pasha rose
and submitted, ``Please proceed to another place, for we now wish to
question the men who made the arrests.'' He therefore went to another
place known as the administrative prison. In that prison, the two Most
Great Branches, my master Āqā Mīrzā Muḥammad Qulī, Āqā Muḥammad ʿAlī
Iṣfahānī, Āqā Muḥammad Javād, Āqā Muḥammad Ḥusayn, son of Ḥājī ʿAlī
ʿAskar, and I were in His presence, along with some of the people of the
government building, officers and others. He remained there until the
seventh hour, continually uttering verses of majesty. Among other
things, He said to this servant, ``Do you remember the verses revealed
last night?'' Then He recited them in the loftiest utterance. These are
the verses:

\subsection{49}\label{section-3168}

\begin{otherlanguage}{arabic}

بسمه الباقی الدّائم العزیز العظیم

\end{otherlanguage}

Bismihi al-Bāqī al-Dāʾim al-ʿAzīz al-ʿAẓīm.

In His name, the Everlasting, the Eternal, the Almighty, the Most Great.

\subsection{50}\label{section-3169}

\begin{otherlanguage}{persian}

قد ماج بحر البلآء و احاطت الأمواج فلک اللّه المهیمن القیّوم ان یا ملّاح لا
تضطرب من الأریاح انّ فالق الأصباح معک فی هذه الظّلمة الّتی احاطت العالمین
توکّل علی اللّه فی کلّ الأحوال و لا تخف من هبوب عواصف البغضآء ان استعذ
باللّه ربّک المقتدر العلیم انّه یحفظ من یشآء بسلطان من عنده انّه لهو العلیم
الحکیم

\end{otherlanguage}

Qad māja baḥru al-balāʾ wa-aḥāṭati al-amwāju fulka Allāhi al-Muhaymini
al-Qayyūm. An yā mallāḥ lā taḍṭarib mina al-aryāḥ. Inna fāliqa al-aṣbāḥ
maʿaka fī hādhihi al-ẓulmati allatī aḥāṭati al-ʿālamīn. Tawakkal ʿalā
Allāhi fī kulli al-aḥwāl wa-lā takhaf min hubūbi ʿawāṣifi al-baghḍāʾ. An
istaʿidh bi-Allāhi rabbika al-Muqtadiri al-ʿAlīm. Innahu yaḥfaẓu man
yashāʾu bi-sulṭānin min ʿindihi. Innahu la-huwa al-ʿAlīmu al-Ḥakīm.

The sea of tribulation has surged, and the waves have surrounded the Ark
of God, the All-Protecting, the Self-Subsisting. O mariner, do not be
shaken by the winds. He who causes the dawn to break is with you in this
darkness that has encompassed the worlds. Trust in God under all
circumstances, and do not fear the raging storms of hatred. Seek refuge
in God, your Lord, the All-Powerful, the All-Knowing. He protects
whomever He wills through a sovereignty of His own. He is indeed the
All-Knowing, the All-Wise.

\subsection{51}\label{section-3170}

\begin{otherlanguage}{arabic}

فی بحبوحة الظّلمة کان الوجه مشرقاً بضیآء احاط من فی السّموات و الأرضین انّا
فی تلک الحالة ندع البریّة الی اللّه و لا یخوّفنا اجتماع الّذین‌هم کفروا باللّه
اذ اتی بأمر بدیع

\end{otherlanguage}

Fī buḥbūḥati al-ẓulmati kāna al-wajhu mushriqan bi-ḍiyāʾin aḥāṭa man fī
al-samāwāti wa-al-araḍīn. Innā fī tilka al-ḥālati nadʿu al-bariyyata ilā
Allāhi wa-lā yukhawwifunā ijtimāʿu alladhīna hum kafarū bi-Allāhi idh
atā bi-amrin badīʿ.

In the depths of darkness, the Face shone with a radiance that
encompassed all who are in the heavens and the earth. In that condition,
We call humanity to God, and We are not frightened by the gathering of
those who disbelieved in God when He came with a wondrous Cause.

\subsection{52}\label{section-3171}

\begin{otherlanguage}{arabic}

قد سرق السّارق ما نزّل من لدی العرش و احضره لدی الّذین یحکمون علی العباد
کذلک فعل ذاک المشرک البعید قل مت بغیضک یا ایّها الجاهل هل تظنّ انّک تسبقنا
لا واسمی الّذی به فاحت نفحات الرّوح علی کلّ صغیر و کبیر انّا بلّغنا الأمر
بأیادی الرّسل من لدنّا انّا کنّا قادرین ثمّ نشرناه بأیدی الّذین اعرضوا انّه لا
یضرّه مکر الماکرین

\end{otherlanguage}

Qad saraqa al-sāriqu mā nuzzila min ladā al-ʿarshi wa-aḥḍarahu ladā
alladhīna yaḥkumūna ʿalā al-ʿibād. Kadhālika faʿala dhāka al-mushriku
al-baʿīd. Qul mut bi-ghayẓika yā ayyuhā al-jāhil. Hal taẓunnu annaka
tasbiqunā? Lā wa-ismī alladhī bihi fāḥat nafaḥātu al-rūḥi ʿalā kulli
ṣaghīrin wa-kabīr. Innā ballaghnā al-amra bi-ayādī al-rusuli min ladunnā
innā kunnā qādirīn. Thumma nasharnāhu bi-aydī alladhīna aʿraḍū. Innahu
lā yaḍurruhu makru al-mākirīn.

The thief has stolen what was sent down from the Throne and brought it
before those who rule over the people. Thus has that wayward idolater
acted. Say: ``Die in your rage, you ignorant one! Do you imagine that
you can outstrip Us? No, by My name, through which the breezes of the
Spirit have wafted over everyone, small and great! We conveyed the Cause
through the hands of messengers sent by Us. We have indeed been
powerful. Then We spread it through the hands of those who turned away.
The scheming of the schemers cannot harm it.''

\subsection{53}\label{section-3172}

\begin{otherlanguage}{arabic}

انّ الخنّاس حضر بکتابی بین النّاس و ظنّ بذلک یهجمون علی مطلع الأمر و یحتقر
به شأن اللّه فیما سواه کذلک سوّلت له نفسه و حقّت علیه کلمة العذاب من لدن
غالب محیط قل بذلک یرفع امره و ینتشر آیاته و یعلو هذا الذّکر الّذی به قدّر
کلّ امر حکیم

\end{otherlanguage}

Inna al-khannāsa ḥaḍara bi-kitābī bayna al-nāsi wa-ẓanna bi-dhālika
yahjumūna ʿalā maṭlaʿi al-amri wa-yuḥtaqaru bihi shaʾnu Allāhi fīmā
siwāhu. Kadhālika sawwalat lahu nafsuhu wa-ḥaqqat ʿalayhi kalimatu
al-ʿadhābi min ladun ghālibin muḥīṭ. Qul bi-dhālika yurfaʿu amruhu
wa-yantashiru āyātuhu wa-yaʿlū hādhā al-dhikru alladhī bihi quddira
kullu amrin ḥakīm.

The slinking whisperer appeared among the people with My Book, imagining
that they would thereby assail the Dawning-Place of the Cause and that
God's standing would be diminished among all besides Him. Thus did his
own self beguile him, and the sentence of punishment from One who is
all-conquering and all-encompassing became binding upon him. Say:
``Through this His Cause is exalted, His verses are spread abroad, and
this Remembrance, through which every wise decree has been ordained, is
raised on high.''

\subsection{54}\label{section-3173}

\begin{otherlanguage}{arabic}

یا لیت یجتمعون علینا العباد و یسفکون دمائنا فی هذا السّبیل المستقیم انّا
فدینا ما عندنا فی سبیل اللّه یشهد بذلک ما انا فیه من البلایا و عن ورائها
قلم الّذی به ثبت امر اللّه العزیز الحکیم قل انّی اکون منتظراً یوم الّذی فیه
اری نفسی بین الأحزاب من جنود الظّالمین تاللّه اذاً یسمعون ما لا سمعوا من
قبل و یرون ما لا رأت عیون الّذین‌هم سبقوا انّه لهو الحاکم علی ما اراد و انّه
لهو القویّ القدیر

\end{otherlanguage}

Yā layta yajtamiʿūna ʿalaynā al-ʿibādu wa-yasfikūna dimāʾanā fī hādhā
al-sabīli al-mustaqīm. Innā fadaynā mā ʿindanā fī sabīli Allāh. Yashhadu
bi-dhālika mā anā fīhi mina al-balāyā wa-ʿan warāʾihā qalamu alladhī
bihi thabata amru Allāhi al-ʿAzīzi al-Ḥakīm. Qul innī akūnu muntaẓiran
yawma alladhī fīhi arā nafsī bayna al-aḥzābi min junūdi al-ẓālimīn.
Ta-Allāhi idhan yasmaʿūna mā lā samiʿū min qablu wa-yarawna mā lā raʾat
ʿuyūnu alladhīna hum sabaqū. Innahu la-huwa al-Ḥākimu ʿalā mā arāda
wa-innahu la-huwa al-Qawiyyu al-Qadīr.

If only the people would gather against Us and shed Our blood on this
straight path! We have sacrificed all that We possess in the path of
God. The afflictions I endure bear witness to this, and beyond them, the
Pen of Him through whom the Cause of God, the Almighty, the All-Wise,
has been established. Say: ``I await the day when I shall see Myself
amid the ranks of the armies of the oppressors. By God, they will then
hear what they have never heard before and see what the eyes of those
who preceded them have never seen. He indeed rules over whatever He
wills, and He is the Strong, the All-Powerful.''

\subsection{55}\label{section-3174}

\begin{otherlanguage}{arabic}

قل یا ایّها الذّباب هل تقدر ان تطیر مع الورقآء فی هذا الهوآء الّذی ما طارت
فیه طیور العالمین کلّ ما یرد علینا انّه رحمة لنا یشهد بذلک کلّ موقن بصیر
تنوح الذّرّات لضرّنا و نحن فی فرح مبین قد اظهرنا السّرور من افق الأحزان انّه
یفعل ما یشآء و یحکم ما یرید

\end{otherlanguage}

Qul yā ayyuhā al-dhubābu hal taqdiru an taṭīra maʿa al-warqāʾi fī hādhā
al-hawāʾi alladhī mā ṭārat fīhi ṭuyūru al-ʿālamīn. Kullu mā yaridu
ʿalaynā innahu raḥmatun lanā yashhadu bi-dhālika kullu mūqinin baṣīr.
Tanūḥu al-dharrātu li-ḍurrinā wa-naḥnu fī faraḥin mubīn. Qad aẓharnā
al-surūra min ufuqi al-aḥzān. Innahu yafʿalu mā yashāʾu wa-yaḥkumu mā
yurīd.

Say: ``O fly, can you soar with the Dove in this atmosphere where the
birds of all the worlds have never flown?'' Whatever befalls Us is a
mercy to Us; everyone possessed of certitude and insight bears witness
to this. The atoms lament Our suffering, while We are in manifest joy.
We have made gladness appear from the horizon of sorrows. He does
whatever He wills and ordains whatever He desires.

\subsection{56}\label{section-3175}

\begin{otherlanguage}{persian}

و در آن احوال ابداً بمأمورین اعتنا نفرمودند وجه مبارک متوجّه بحضرت غصن
اعظم و حضرت غصن اکبر روحی لهما الفدآء و بعضی احبّا که در خدمت حاضر بودند
بوده بعد بحر بیان رحمن موّاج از جمله فرمودند

\end{otherlanguage}

Va dar ān aḥvāl abadan bi-maʾmūrīn iʿtinā nafarmūdand. Vajh-i mubārak
mutavajjih bi-Ḥaḍrat-i Ghuṣn-i Aʿẓam va Ḥaḍrat-i Ghuṣn-i Akbar, rūḥī
lahumā al-fidāʾ, va baʿḍī aḥibbā kih dar khidmat ḥāḍir būdand būdih.
Baʿd baḥr-i bayān-i Raḥmān mavvāj, az jumlih farmūdand.

Throughout those circumstances He paid no attention whatever to the
officials. His blessed face was turned toward His Holiness the Most
Great Branch and His Holiness the Greater Branch---may my soul be
sacrificed for them both---and some of the loved ones who were present
in His service. Then the ocean of the All-Merciful's utterance surged,
and among other things He said:

\subsection{57}\label{section-3176}

\begin{otherlanguage}{arabic}

قد اوحی الله تبارک و تعالی الی بعض انبیائه فی بعض وحیه وعزّتی و جلالی
لأقطعنّ رجآء کلّ مؤمّل یؤمّل غیری و لأکسونّه ثوب المذلّة بین النّاس و لأبعدنّه
عن فضلی و فرجی أ یؤمّل عبدی فی الشّدائد غیری و مفتاحها بیدی ما لی اری عبدی
معرضاً عنّی و قد اعطیته ما لم یسألنی أ فأسأل فلا اجود کلّا أ لیس الجود و
الکرم لی أ لیس الدّنیا و الآخرة بیدی و لو انّ کلّ واحد من اهل السّموات و
الأرض سألنی مثل السّموات و الأرض و اعطیته ما نقص ذلک من ملکی علی قدر جناح
بعوضة

\end{otherlanguage}

Qad awḥā Allāhu tabāraka wa-taʿālā ilā baʿḍi anbiyāʾihi fī baʿḍi
waḥyihi: Wa-ʿizzatī wa-jalālī la-aqṭaʿanna rajāʾa kulli muʾammilin
yuʾammilu ghayrī wa-la-aksuwannahu thawba al-madhallati bayna al-nāsi
wa-la-ubʿidannahu ʿan faḍlī wa-farajī. A-yuʾammilu ʿabdī fī al-shadāʾidi
ghayrī wa-miftāḥuhā bi-yadī? Mā lī arā ʿabdī muʿriḍan ʿannī wa-qad
aʿṭaytuhu mā lam yasʾalnī? A-fa-usʾalu fa-lā ajūdu? Kallā. A-laysa
al-jūdu wa-al-karamu lī? A-laysati al-dunyā wa-al-ākhiratu bi-yadī?
Wa-law anna kulla wāḥidin min ahli al-samāwāti wa-al-arḍi saʾalanī
mithla al-samāwāti wa-al-arḍi wa-aʿṭaytuhu mā naqaṣa dhālika min mulkī
ʿalā qadri janāḥi baʿūḍah.

God, blessed and exalted be He, revealed to one of His prophets in one
of His revelations: ``By My might and My glory, I shall cut off the hope
of everyone who places hope in anyone but Me, clothe him in the garment
of abasement among the people, and remove him far from My bounty and My
relief. Does My servant place his hope in anyone but Me amid hardships,
when their key is in My hand? Why do I see My servant turning away from
Me, when I have given him what he did not ask of Me? Would I then be
asked and not give generously? Certainly not! Are not generosity and
bounty Mine? Are not this world and the next in My hand? Were every one
of the inhabitants of the heavens and the earth to ask Me for the
equivalent of the heavens and the earth, and were I to give it to him,
this would not diminish My dominion by so much as a mosquito's wing.''

\subsection{58}\label{section-3177}

\begin{otherlanguage}{arabic}

و امثال این بیانات که اکثری در نظر نیست متّصلاً از کوثر فم رحمن جاری بوده
قد تعاطت اقداح البیان من خمر ذکر ربّک الرّحمن آنچه بخاطر این عبد مانده
معروض داشته چه عرض کنم که در آن لیلۀ مبارکه آیات قدرت و عظمت و سلطنت چه
قسم جاری و سایل و نازل بوده کأنّ سلسبیل البیان بذل علی من فی الامکان و
بحر الکرم ترشّح علی الأمم و سرّ المجلّل بالسّرّ تجلّی علی من فی الغیب و الشّهود
بسیار از بیانات ابداً در نظر نمانده

\end{otherlanguage}

Va amthāl-i īn bayānāt kih aktharī dar naẓar nīst muttaṣilan az
kawthar-i fam-i Raḥmān jārī būdih. Qad taʿāṭat aqdāḥu al-bayāni min
khamri dhikri rabbika al-Raḥmān. Ānchih bi-khāṭir-i īn ʿabd māndih
maʿrūḍ dāshtih. Chih ʿarḍ kunam kih dar ān layliy-i mubārakih āyāt-i
qudrat va ʿaẓamat va salṭanat chih qism jārī va sāyil va nāzil būdih.
Ka-anna salsabīla al-bayāni budhila ʿalā man fī al-imkāni wa-baḥra
al-karami tarashshaḥa ʿalā al-umami wa-sirra al-mujallali bi-al-sirri
tajallā ʿalā man fī al-ghaybi wa-al-shuhūd. Bisyār az bayānāt abadan dar
naẓar namāndih.

Utterances such as these, most of which I cannot recall, flowed
unceasingly from the Kawthar of the All-Merciful's mouth. The cups of
utterance passed from hand to hand, filled with the wine of the
remembrance of your Lord, the All-Merciful. I have submitted what
remains in this servant's memory. What can I say of how the verses of
power, majesty, and sovereignty flowed, streamed, and descended on that
blessed night? It was as though the Salsabil of utterance had been
bestowed upon all in the contingent world, the ocean of bounty had
sprinkled its waters upon the nations, and the mystery of the One veiled
in mystery had shone forth upon all in the unseen and visible realms.
Many of the utterances have entirely escaped my memory.

\subsection{59}\label{section-3178}

\begin{otherlanguage}{persian}

بعد در ساعت سابعه یکی از ضبّاط آمده جمال قدم و غصنین اعظمین و آقائی آقا
میرزا محمّد قلی را بردند در عریضۀ قبل عرض شده که جمال قدم و حضرت غصن اکبر
را در محلّی که جنب لیمانست و آقائی را در محلّ دیگر و حضرت غصن اعظم را در
لیمان تحت حدید این عبد و سایر احبّا را در نفس سرایه محلّ بسیار مظلم تحت
حدید حبس نمودند

\end{otherlanguage}

Baʿd dar sāʿat-i sābiʿih yakī az ḍubbāṭ āmadih Jamāl-i Qidam va
Ghuṣnayn-i Aʿẓamayn va āqāʾī Āqā Mīrzā Muḥammad Qulī rā burdand. Dar
ʿarīḍiy-i qabl ʿarḍ shudih kih Jamāl-i Qidam va Ḥaḍrat-i Ghuṣn-i Akbar
rā dar maḥallī kih janb-i līmān ast va āqāʾī rā dar maḥall-i dīgar va
Ḥaḍrat-i Ghuṣn-i Aʿẓam rā dar līmān taḥt-i ḥadīd, īn ʿabd va sāyir-i
aḥibbā rā dar nafs-i sarāyih maḥall-i bisyār muẓlim taḥt-i ḥadīd ḥabs
nimūdand.

Then, at the seventh hour, one of the officers came, and they took away
the Ancient Beauty, the two Most Great Branches, and my master Āqā Mīrzā
Muḥammad Qulī. As was submitted in the previous letter, they confined
the Ancient Beauty and His Holiness the Greater Branch in a place next
to the prison, my master in another place, and His Holiness the Most
Great Branch in the prison in irons. This servant and the other loved
ones were confined in irons in a very dark place within the government
building itself.

\subsection{60}\label{section-3179}

\begin{otherlanguage}{persian}

در لیلۀ ثانیه تلغراف از والی آمده محلّ جمال قدم را تغییر دادند در قبّۀ
مرتفعه که فوق لیمان است تشریف بردند غصنین اعظمین و آقائی در حضور امنع
اقدس و توقّف در قبّه سی و هشت ساعت و نیم بعد مجدّد یوم خمیس ساعت رابع از
نهار بسرایه طلب نمودند لأجل سؤال و جواب

\end{otherlanguage}

Dar layliy-i thāniyih tilighrāf az vālī āmadih maḥall-i Jamāl-i Qidam rā
taghyīr dādand. Dar qubbiy-i murtafiʿih kih fawq-i līmān ast tashrīf
burdand. Ghuṣnayn-i Aʿẓamayn va āqāʾī dar ḥuḍūr-i amnaʿ-i aqdas va
tavaqquf dar qubbih sī va hasht sāʿat va nīm. Baʿd mujaddad yawm-i
khamīs sāʿat-i rābiʿ az nahār bi-sarāyih ṭalab nimūdand li-ajli suʾālin
wa-jawāb.

On the second night a telegram arrived from the governor, and they
changed the Ancient Beauty's quarters. He went to an elevated domed
chamber above the prison. The two Most Great Branches and my master were
in His most inaccessible, most holy presence, and the stay in the domed
chamber lasted thirty-eight and a half hours. Then, on Thursday, at the
fourth hour of the day, they again summoned Him to the government
building for questioning.

\subsection{61}\label{section-3180}

\begin{otherlanguage}{persian}

حین توجّه بمجلس لسان عظمت باین کلمه ناطق اخذتم اصول انفسکم و نبذتم اصول
الله عن ورائکم ما لکم لا تفقهون این آیه مکرّر از لسان عظمت جاری و بعد از
ورود جمال قدم پاشا و اهل مجلس معذرت خواستند که ما مکلّف به تفحّصیم چه که
سه ٣ نفس کشته شده‌اند و جمیع اهل بلد خائف و هراسان لذا باید بموجب قانون
عمل شود گویا بیچاره از قانون هم اطّلاع نداشته

\end{otherlanguage}

Ḥīn-i tavajjuh bi-majlis lisān-i ʿaẓamat bi-īn kalimih nāṭiq: Akhadhtum
uṣūla anfusikum wa-nabadhtum uṣūla Allāhi ʿan warāʾikum mā lakum lā
tafqahūn. Īn āyih mukarrar az lisān-i ʿaẓamat jārī va baʿd az vurūd-i
Jamāl-i Qidam pāshā va ahl-i majlis maʿdhirat khvāstand kih mā mukallaf
bi-tafaḥḥuṣīm chih kih sih 3 nafs kushtih shudih-and va jamīʿ-i ahl-i
balad khāʾif va harāsān lidhā bāyad bi-mūjib-i qānūn ʿamal shavad. Gūyā
bīchārih az qānūn ham iṭṭilāʿ nadāshtih.

As He proceeded to the council, the Tongue of Grandeur uttered these
words: ``You have adopted your own principles and cast the principles of
God behind your backs. Why do you not understand?'' This verse flowed
repeatedly from the Tongue of Grandeur. After the Ancient Beauty
entered, the pasha and the members of the council apologized, saying,
``We are obliged to investigate, since three---3---people have been
killed and all the townspeople are frightened and alarmed. We must
therefore act in accordance with the law.'' Apparently the poor man knew
nothing of the law either.

\subsection{62}\label{section-3181}

\begin{otherlanguage}{persian}

باری سؤال از اسم مبارک و وطن نمودند و معروض داشتند که موجب قانون این است
که این سؤالات بشود و نوشته گردد فرمودند انّه اظهر من الشّمس مجدّد سؤال شد
فرمودند لا ینبغی ذکر الاسم فانظروا فی فرمان الدّولة الّذی عندکم بعد بکمال
ملایمت عرض نمودند خود شما بفرمائید

\end{otherlanguage}

Bārī suʾāl az ism-i mubārak va vaṭan nimūdand va maʿrūḍ dāshtand kih
mūjib-i qānūn īn ast kih īn suʾālāt bishavad va nivishtih gardad.
Farmūdand: Innahu aẓharu mina al-shams. Mujaddad suʾāl shud. Farmūdand:
Lā yanbaghī dhikru al-ismi fa-unẓurū fī farmāni al-dawlati alladhī
ʿindakum. Baʿd bi-kamāl-i mulāyimat ʿarḍ nimūdand khud-i shumā
bifarmāyīd.

They then asked His blessed name and His homeland, submitting that the
law required these questions to be asked and the answers recorded. He
said, ``It is more evident than the sun.'' The question was repeated. He
said, ``There is no need to mention the name. Consult the government
decree in your possession.'' Then, with the utmost gentleness, they
submitted, ``Please state it Yourself.''

\subsection{63}\label{section-3182}

\begin{otherlanguage}{arabic}

اذاً جلس هیکل العظمة مستویاً علی السّریر و نطق بلسان القدرة و القوّة اسمی
بهآءالله و مسکنی نور اذاً فاعرفوا ثمّ توجّه وجه القدم الی المفتی و قال عزّ
کبریائه لو عرفتم لأسمعناکم هدیر ورقآء العظمة علی غصن سدرة الرّبّانیّة لیظهر
لکم ما قال العبودیّة جوهرة کنهها الرّبوبیّة

\end{otherlanguage}

Idhan jalasa Haykalu al-ʿAẓamati mustawiyan ʿalā al-sarīri wa-naṭaqa
bi-lisāni al-qudrati wa-al-quwwah: Ismī Bahāʾuʾllāh wa-maskanī Nūr idhan
fa-iʿrifū. Thumma tawajjaha Wajhu al-Qidami ilā al-muftī wa-qāla ʿazza
kibriyāʾuhu: Law ʿaraftum la-asmaʿnākum hadīra Warqāʾi al-ʿAẓamati ʿalā
ghuṣni Sidrati al-Rabbāniyyati li-yaẓhara lakum mā qāla al-ʿubūdiyyatu
jawharatun kunhuhā al-rubūbiyyah.

Then the Temple of Grandeur sat upright upon the seat and spoke with the
tongue of power and might: ``My name is Bahá'u'lláh, and My dwelling is
Nūr. Know this, then.'' The Ancient Face then turned toward the mufti
and said---exalted be His majesty---``Had you understood, We would have
let you hear the cooing of the Dove of Grandeur upon the branch of the
Lote Tree of Divinity, so that the meaning of the saying, `Servitude is
a gem whose inmost essence is lordship,' might become clear to you.''

\subsection{64}\label{section-3183}

\begin{otherlanguage}{persian}

فبهت من حضر من کلمة الله الأبدی و بعد مخاطباً الی الکلّ فرمودند ما المقصود
من هذا الاجتماع و السّؤال و الجواب لو تریدون اعترف بما هو مقصودکم و رجآء
قلوبکم لأنّی ما احبّ ان اکون فی الدّنیا علی قدر ساعة و از آیات سور ملوک
تلاوت فرمودند از جمله این فقرۀ مبارکه ما مررت علی شجر الّا و خاطبه فؤادی
یا لیت قطعت لاسمی و صلب علیک جسدی فی سبیل ربّی

\end{otherlanguage}

Fa-buhita man ḥaḍara min kalimati Allāhi al-Abadī. Va baʿd mukhāṭiban
ilā al-kull farmūdand: Mā al-maqṣūdu min hādhā al-ijtimāʿi wa-al-suʾāli
wa-al-jawāb? Law turīdūna aʿtarifu bi-mā huwa maqṣūdukum wa-rajāʾu
qulūbikum li-annī mā uḥibbu an akūna fī al-dunyā ʿalā qadri sāʿah. Va az
āyāt-i Suvar-i Mulūk tilāvat farmūdand az jumlih īn fiqriy-i mubārakih:
Mā marartu ʿalā shajarin illā wa-khāṭabahu fuʾādī yā layta quṭiʿta
li-ismī wa-ṣuliba ʿalayka jasadī fī sabīli rabbī.

Those present were stunned by the word of God, the Eternal. Then,
addressing everyone, He said, ``What is the purpose of this gathering
and this questioning? If you wish, I will confess to what you intend and
what your hearts desire, for I have no wish to remain in this world for
even an hour.'' He also recited verses from the Tablets to the Kings,
including this blessed passage: ``I have never passed a tree without My
heart addressing it: `If only you were cut down in My name and My body
crucified upon you in the path of My Lord!'\,''

\subsection{65}\label{section-3184}

\begin{otherlanguage}{arabic}

و فی اثنآء البیان قد اخذ الاهتزاز ارکان ا لرّحمن ثمّ بعد ذلک قام القیام و
توجّه جمال العلّام الی محلّ آخر خارجاً عن المجمع

\end{otherlanguage}

Wa-fī athnāʾi al-bayāni qad akhadha al-ihtizāzu arkāna al-Raḥmān. Thumma
baʿda dhālika qāma al-qiyāmu wa-tawajjaha Jamālu al-ʿAllāmi ilā maḥallin
ākhara khārijan ʿani al-majmaʿ.

During His utterance, a trembling seized the limbs of the All-Merciful.
Then He stood upright, and the Beauty of the All-Knowing proceeded to
another place outside the assembly.

\subsection{66}\label{section-3185}

\begin{otherlanguage}{persian}

باری عجایب روزی و عجایب وقتی بوده در بعضی از ناس شمس کلمة الله اشراق
نموده اظهار محبّت نمودند فی‌الحقیقه این عبد از ذکر بیانات رحمانی عاجر است
صاحب مثنوی ذکری نموده مناسب این مقام است

\end{otherlanguage}

Bārī ʿajāyib rūzī va ʿajāyib vaqtī būdih. Dar baʿḍī az nās shams-i
Kalimatuʾllāh ishrāq nimūdih iẓhār-i maḥabbat nimūdand. Fī al-ḥaqīqih īn
ʿabd az dhikr-i bayānāt-i Raḥmānī {[}corrupted text{]} ast. Ṣāḥib-i
Mathnavī dhikrī nimūdih munāsib-i īn maqām ast.

In short, it was an extraordinary day and an extraordinary hour. The sun
of the Word of God shone upon some of the people, and they expressed
their love. In truth, this servant is {[}corrupted original text{]} to
recount the utterances of the All-Merciful. The author of the Mathnavī
has said something fitting for this occasion:

\subsection{67}\label{section-3186}

\begin{otherlanguage}{arabic}

آن خطاباتی که گفت آندم نبی گر زند بر شب برآید از شبی روز روشن گردد

\end{otherlanguage}

Ān khiṭābātī kih guft ān dam nabī gar zanad bar shab barāyad az shabī
rūz-i rawshan gardad.

Those words the Prophet spoke at that moment---if he addressed them to
the night, it would emerge from darkness and become bright day.

\subsection{68}\label{section-3187}

\begin{otherlanguage}{arabic}

آن شب چون صباح من نتانم گفت باز آن اصطلاح

\end{otherlanguage}

Ān shab chūn ṣabāḥ man natānam guft bāz ān iṣṭilāḥ.

That night, like morning---I cannot repeat that manner of expression.

\subsection{69}\label{section-3188}

\begin{otherlanguage}{arabic}

خود تو دانی کآفتاب اندر حمل می چه گوید با ریاحین با دقل

\end{otherlanguage}

Khud tu dānī k-āftāb andar Ḥamal mī chih gūyad bā rayāḥīn bā daqal.

You yourself know what the sun in Aries says to the fragrant herbs and
the date palm.

\subsection{70}\label{section-3189}

\begin{otherlanguage}{arabic}

هم تو میدانی که آن آب زلال می چه گوید با ریاحین با نهال

\end{otherlanguage}

Ham tu mīdānī kih ān āb-i zulāl mī chih gūyad bā rayāḥīn bā nihāl.

You also know what that limpid water says to the fragrant herbs and the
sapling.

\subsection{71}\label{section-3190}

\begin{otherlanguage}{arabic}

باری این عبد از ذکر بیانات الهیّه و اصطلاحات ربّانیّه عاجر بوده و خواهد بود

\end{otherlanguage}

Bārī īn ʿabd az dhikr-i bayānāt-i ilāhiyyih va iṣṭilāḥāt-i rabbāniyyih
{[}corrupted text{]} būdih va khvāhad būd.

In short, this servant has been and will remain {[}corrupted original
text{]} to recount the divine utterances and heavenly expressions.

\subsection{72}\label{section-3191}

\begin{otherlanguage}{persian}

بعد پاشا فرستاد که ببیت مراجعت فرمائید و معذرت خواست از آنچه واقع شده
فرمودند رجوع بحبس اولی است چه که جمعی محبوسند بعضی تلقاء وجه حاضر شده
معروض داشتند که حال اهل بلد در حرکتند نمیتوان محبوسین را بیرون آورد چند
یومی اگر توقّف شود بهتر است بعد مراجعب ببیت فرمودند و آن حین ساعت احدی‌عشر
یوم خمیس بود

\end{otherlanguage}

Baʿd pāshā firistād kih bi-bayt murājaʿat farmāyīd va maʿdhirat khvāst
az ānchih vāqiʿ shudih. Farmūdand rujūʿ bi-ḥabs awlā ast chih kih jamʿī
maḥbūsand. Baʿḍī tilqāʾ-i vajh ḥāḍir shudih maʿrūḍ dāshtand kih ḥāl
ahl-i balad dar ḥarakatand namītavān maḥbūsīn rā bīrūn āvard chand yawmī
agar tavaqquf shavad bihtar ast. Baʿd {[}corrupted text{]} bi-bayt
farmūdand va ān ḥīn sāʿat-i iḥdā-ʿashar-i yawm-i khamīs būd.

Then the pasha sent word asking Him to return home and apologized for
what had happened. He said, ``It is preferable to return to prison,
since a number of people are imprisoned.'' Some came into His presence
and submitted, ``The townspeople are now in a state of agitation, and
the prisoners cannot be brought out. It would be better to wait a few
days.'' Then He {[}corrupted original text{]} home, and it was then the
eleventh hour on Thursday.

\subsection{73}\label{section-3192}

\begin{otherlanguage}{persian}

مجمل آنکه ابتدای خروج جمال قدم از حرم الی حین ورود شصت و نه ساعت و نیم
منقضی شد توقّف در سرایه لیلۀ اولی ٥ ساعت و نیم محبس جنب لیمان هیجده ساعت
و نیم حبس قبّه فوق لیمان سی و هشت ساعت و نیم ورود سرایه و توقّف در آن مرّۀ
ثانیه هفت ساعت کذلک قضی الأمر من لدی الله المقتدر العلیم الحکیم

\end{otherlanguage}

Mujmal ānkih ibtidā-yi khurūj-i Jamāl-i Qidam az ḥaram ilā ḥīn-i vurūd
shaṣt u nuh sāʿat u nīm munqaḍī shud. Tavaqquf dar sarāyih layliy-i ūlā
panj sāʿat u nīm, maḥbas-i janb-i līmān hijdah sāʿat u nīm, ḥabs-i
qubbih-yi fawq-i līmān sī u hasht sāʿat u nīm, vurūd-i sarāyih va
tavaqquf dar ān marriy-i thāniyih haft sāʿat. Kadhālika quḍiya al-amru
min ladā Allāhi al-Muqtadiri al-ʿAlīmi al-Ḥakīm.

In short, from the time the Ancient Beauty left the household until His
return, sixty-nine and a half hours elapsed. The stay in the government
house on the first night lasted five and a half hours; in the prison
beside the harbor, eighteen and a half hours; confinement in the domed
chamber above the harbor, thirty-eight and a half hours; and the return
to the government house and the second stay there, seven hours. Thus was
the matter decreed by God, the All-Powerful, the All-Knowing, the
All-Wise.

\subsection{74}\label{section-3193}

\begin{otherlanguage}{persian}

و دیگر آنکه صحبت نامیست از اهل زاء علیه لعنة الله اذا حدث کذب و اذا وعد
اخلف و اذ ائتمن خان این صفات منافقین است و صدهزار فوق آن در آن خبیث
موجود قبل از وقوع این مقدّمۀ مذکوره وارد این بلد شد و در مقرّ مشرکین منزل
داشت و چندی با نفوس خبیثه معاشر و بعد با روایات کاذبۀ سقیمه مراجعت نمود
و بهر محلّ که رسید مفتریاتی ذکر نمود او و آقا جان علیهما غضب الله راوی
مجعولات مفتریه از ناحیۀ کذبه شده‌اند چنانچه بعد از انقلابات مدینۀ کبیره و
عزل مشیر و موت صدر اعظم نسبت دادند این امور لأجل ظلم بر ما واقع شده و
این امور را رأس المشرکین خبر داده کلّ هذا کذب صراح و هر عاقل منصفی که
فی‌الجمله بعرف انصاف فائز شده یشهد بکذبهم چه که جمیع نفوس از عرب و عجم و
ترک در صدد جمال قدم بودند و بعداوت تمام قیام نموده بودند و سایرین را
داخل وجود نمیدانستند و بر سایر آنچه واقع شده بالتّبع بوده

\end{otherlanguage}

Va dīgar ānkih Ṣuḥbat nāmīst az ahl-i Zāʾ, ʿalayhi laʿnatu Allāh. Idhā
ḥaddatha kadhaba wa-idhā waʿada akhlafa wa-idh uʾtumina khāna. Īn
ṣifāt-i munāfiqīn ast va ṣad-hizār fawq-i ān dar ān khabīth mawjūd. Qabl
az vuqūʿ-i īn muqaddamih-yi madhkūrih vārid-i īn balad shud va dar
maqarr-i mushrikīn manzil dāsht va chandī bā nufūs-i khabīthih muʿāshir
va baʿd bā rivāyāt-i kādhibih-yi saqīmih murājaʿat namūd va bih har
maḥall kih rasīd muftariyātī dhikr namūd. Ū va Āqā Jān, ʿalayhimā
ghaḍabu Allāh, rāvī-yi majʿūlāt-i muftariyih az nāḥiyih-yi kadhabih
shudih-and, chunānchih baʿd az inqilābāt-i Madīnih-yi Kabīrih va ʿazl-i
mushīr va mawt-i ṣadr-i aʿẓam nisbat dādand īn umūr li-ajl-i ẓulm bar mā
vāqiʿ shudih va īn umūr rā Raʾs-i Mushrikīn khabar dādih. Kullu hādhā
kadhibun ṣurāḥun. Va har ʿāqil-i munṣifī kih fī al-jumlih bih ʿarf-i
inṣāf fāʾiz shudih yashhadu bi-kadhibihim, chih kih jamīʿ-i nufūs az
ʿArab va ʿAjam va Turk dar ṣadad-i Jamāl-i Qidam būdand va bih ʿadāvat-i
tamām qiyām namūdih būdand va sāyirīn rā dākhil-i vujūd nimīdānistand va
bar sāyir ānchih vāqiʿ shudih bi-al-tabaʿ būdih.

Furthermore, there is a man named Suhbat from the people of Za---may
God's curse be upon him. When he speaks, he lies; when he promises, he
breaks his promise; and when he is trusted, he betrays that trust. These
are the attributes of hypocrites, and a hundred thousand times more than
these are found in that wicked man. Before the aforementioned incident
occurred, he came to this city and stayed at the headquarters of the
idolaters. For a time he associated with wicked souls, then returned
with false and unsound accounts, uttering fabrications wherever he went.
He and Aqa Jan---may God's wrath be upon them both---have become
narrators of falsehoods fabricated by liars. Thus, after the upheavals
in the Great City, the dismissal of the marshal, and the death of the
grand vizier, they alleged, ``These things occurred because of the
injustice done to us, and the Chief of the Idolaters foretold them.''
All this is an outright lie. Every fair-minded, rational person who has
gained even a little of the fragrance of fairness bears witness to their
falsehood, for everyone, Arab, Persian, and Turk alike, was intent on
opposing the Ancient Beauty and had risen in utter hostility. They
regarded the others as nonexistent, and whatever befell those others was
incidental.

\subsection{75}\label{section-3194}

\begin{otherlanguage}{persian}

باری اخذ ظالمین و ما ورد علیهم و ما یرد کلّ در لوح رئیس و الواح اخری
بکمال تصریح نازل و حال این عبد بعضی را ذکر مینماید تا کلّ بدانند که آن
نفوس کاذب بوده و خواهند بود و مطلع آیات الهیّه و منبع فیوضات رحمانیّه
بآنچه وارد شده خبر داده‌اند از جمله لوح جناب ابن نبیل علیه بهآء الله است
که در اوّل ورود سجن اعظم نازل و اصل لوح نزد جناب مذکور موجود و سواد آن
نزد جناب زین‌المقرّبین علیه بهآء الله و احبّا موجود است جویا شوند و تحقیق
نمایند تا بیقین بدانند که زمام علوم بید مقتدر علیم بوده و خواهد بود لا
یعلم الغیب الّا هو یظهر لمن یشآء ما یشآء و یستر لمن اراد ما اراد انّه لهو
العلیم الخبیر بعضی آیات منزله در لوح ابن نبیل اینست

\end{otherlanguage}

Bārī akhdh-i ẓālimīn va mā warada ʿalayhim wa-mā yaridu kull dar Lawḥ-i
Raʾīs va alvāḥ-i ukhrā bih kamāl-i taṣrīḥ nāzil va ḥāl īn ʿabd baʿḍī rā
dhikr mīnamāyad tā kull bidānand kih ān nufūs kādhib būdih va khvāhand
būd va Maṭlaʿ-i Āyāt-i Ilāhiyyih va Manbaʿ-i Fuyūḍāt-i Raḥmāniyyih
bi-ānchih vārid shudih khabar dādih-and. Az jumlih Lawḥ-i Jināb-i Ibn-i
Nabīl, ʿalayhi Bahāʾu Allāh, ast kih dar avval-i vurūd-i Sijn-i Aʿẓam
nāzil va aṣl-i lawḥ nazd-i jināb-i madhkūr mawjūd va savād-i ān nazd-i
Jināb-i Zayn al-Muqarrabīn, ʿalayhi Bahāʾu Allāh, va aḥibbā mawjūd ast.
Jūyā shavand va taḥqīq namāyand tā bih yaqīn bidānand kih zimām-i ʿulūm
bih yad-i Muqtadir-i ʿAlīm būdih va khvāhad būd. Lā yaʿlamu al-ghayba
illā Huwa, yuẓhiru li-man yashāʾu mā yashāʾu wa-yasturu li-man arāda mā
arāda, innahu la-Huwa al-ʿAlīmu al-Khabīr. Baʿḍī āyāt-i munazzalih dar
Lawḥ-i Ibn-i Nabīl īnast.

In any event, the seizure of the oppressors, what has befallen them, and
what will befall them have all been revealed with complete explicitness
in the Tablet of the Chief and other Tablets. This servant will now
mention some of these things so that all may know that those souls have
been and will remain liars, and that the Dayspring of the divine verses
and the Source of the outpourings of the All-Merciful foretold what has
occurred. Among these is the Tablet to Ibn-i-Nabil---upon him be the
glory of God---revealed upon arrival in the Most Great Prison. The
original Tablet is with the aforementioned gentleman, and copies are
with Zayn al-Muqarrabin---upon him be the glory of God---and the
friends. Let them seek it out and investigate, so that they may know
with certainty that the reins of knowledge have been and will remain in
the hand of the All-Powerful, the All-Knowing. No one knows the unseen
except Him. He reveals to whomever He wills whatever He wills, and
conceals from whomever He desires whatever He desires. He is indeed the
All-Knowing, the All-Informed. Some of the verses revealed in the Tablet
to Ibn-i-Nabil are as follows:

\subsection{76}\label{section-3195}

\begin{otherlanguage}{arabic}

هو الأقدس الأبهی

\end{otherlanguage}

Huwa al-Aqdasu al-Abhā.

He is the Most Holy, the All-Glorious.

\subsection{77}\label{section-3196}

\begin{otherlanguage}{persian}

ک‌ظ نادیناک عن ورآء قلزم الکبریآء علی ارض الحمرآء من افق البلآء انّه لا
اله الّا هو العزیز الوهّاب ان استقم علی امری و لا تکن من الّذین اذا اوتوا
ما ارادوا کفروا باللّه ربّ الأرباب سوف یأخذهم اللّه بقهر من عنده انّه لهو
المقتدر القهّار فاعلم انّ الّذین حکموا علینا قد اخذ اللّه کبیرهم بقدرة و
سلطان و بعد آیاتی چند نازل تا به این آیه میرسد قوله جلّ کبریائه سوف نعزل
الّذی کان مثله و نأخذ امیرهم الّذی یحکم علی البلاد و انا العزیز الجبّار

\end{otherlanguage}

Kāf-Ẓāʾ. Nādaynāka ʿan warāʾi Qulzumi al-Kibriyāʾi ʿalā Arḍi al-Ḥamrāʾi
min ufuqi al-balāʾi, innahu lā ilāha illā Huwa al-ʿAzīzu al-Wahhāb. An
istaqim ʿalā amrī wa-lā takun mina alladhīna idhā ūtū mā arādū kafarū
bi-Allāhi Rabbi al-arbāb. Sawfa yaʾkhudhuhumu Allāhu bi-qahrin min
ʿindihi, innahu la-Huwa al-Muqtadiru al-Qahhār. Fa-iʿlam anna alladhīna
ḥakamū ʿalaynā qad akhadha Allāhu kabīrahum bi-qudratin wa-sulṭān. Va
baʿd āyātī chand nāzil tā bih īn āyih mīrasad, qawluhu jalla
kibriyāʾuhu: Sawfa naʿzilu alladhī kāna mithlahu wa-naʾkhudhu amīrahum
alladhī yaḥkumu ʿalā al-bilādi wa-anā al-ʿAzīzu al-Jabbār.

``Kaf-Za. We have called to you from beyond the Ocean of Grandeur, upon
the Crimson Land, from the horizon of tribulation: there is no God but
Him, the Almighty, the All-Bountiful. Be steadfast in My Cause, and do
not be among those who, when given what they desired, disbelieved in
God, the Lord of lords. God will seize them with an overpowering might
from Him. He is indeed the All-Powerful, the All-Subduing. Know, then,
that God has seized the greatest of those who passed judgment against
Us, with power and sovereignty.'' Several more verses are then revealed,
until this verse is reached---His words, exalted be His grandeur: ``We
will dismiss the one who was like him and seize their ruler who governs
the lands. I am the Almighty, the All-Compelling.''

\subsection{78}\label{section-3197}

\begin{otherlanguage}{arabic}

چنانچه چندی نگذشت نفسی که سبب مباشر نفی حقّ و آل الله بود از مدینۀ کبیره
نفی نمودند و کبیرشان که جمیع حکم از او صادر بدرک رفت حال ملاحظه کنید مع
این آیات منزله که از قبل در کتاب الهی نازل و در بلاد منتشر صحبت زنجانی
بعد از وقوع گفته رئیس المغلّین اخبار داده الا لعنة الله علی الکاذبین

\end{otherlanguage}

Chunānchih chandī nagudhasht nafsī kih sabab-i mubāshir-i nafy-i Ḥaqq va
Āl-i Allāh būd az Madīnih-yi Kabīrih nafy namūdand va kabīrashān kih
jamīʿ-i ḥukm az ū ṣādir bih darak raft. Ḥāl mulāḥiẓih kunīd maʿa īn
āyāt-i munazzalih kih az qabl dar Kitāb-i Ilāhī nāzil va dar bilād
muntashir, Ṣuḥbat-i Zanjānī baʿd az vuqūʿ guftih Raʾīs-i Mughillīn
ikhbār dādih. A-lā laʿnatu Allāhi ʿalā al-kādhibīn.

Thus, before long, the person directly responsible for the exile of the
True One and the family of God was banished from the Great City, and
their chief, from whom all the orders proceeded, went to hell. Now
consider: despite these revealed verses, which had previously been
revealed in the Book of God and circulated throughout the lands, Suhbat
of Zanjan said after the events that the Chief of the Malevolent had
foretold them. Indeed, may God's curse be upon the liars.

\subsection{79}\label{section-3198}

\begin{otherlanguage}{persian}

و همچنین سورۀ رئیس را ملاحظه نمائید که بعد از هجرت ارض سرّ و ورود شاطی
بحر نازل شده و همچنین آیات هیکل که مخصوص ملک پاریس و اهل آن دیار نازل و
اخبار فرموده‌اند از آنچه بعد بر او واقع شده قوله جلّ کبریائه

\end{otherlanguage}

Va hamchunīn Sūriy-i Raʾīs rā mulāḥiẓih namāʾīd kih baʿd az hijrat-i
Arḍ-i Sirr va vurūd-i shāṭī-yi baḥr nāzil shudih va hamchunīn āyāt-i
Haykal kih makhṣūṣ-i malik-i Pārīs va ahl-i ān diyār nāzil va ikhbār
farmūdih-and az ānchih baʿd bar ū vāqiʿ shudih, qawluhu jalla
kibriyāʾuhu.

Likewise, consider the Surih of the Chief, revealed after the departure
from the Land of Mystery and the arrival at the seashore, and also the
verses of the Temple, revealed specifically for the king of Paris and
the people of that land, which foretold what later befell him. His
words, exalted be His grandeur:

\subsection{80}\label{section-3199}

\begin{otherlanguage}{persian}

ان یا ملک ان استمع النّدآء من هذه النّار المشتعلة من الشّجرة الخضرآء فی هذا
الطّور المرتفع علی البقعة المقدّسة البیضآء خلف قلزم البقآء انّه لا اله الّا
انا الغفور الرّحیم الی ان قال عزّ کبریائه ان یا ملک انّا سمعنا منک کلمة
تکلّمت بها اذ سألک ملک الرّوس عمّا قضی من حکم الغزا انّ ربّک لهو العلیم
الخبیر قلت کنت راقداً فی المهاد ایقظنی ندآء العباد الّذین ظلموا الی ان
غرقوا فی البحر الأسود کذلک سمعنا و ربّک علی ما اقول شهید نشهد بأنّک ما
ایقظک النّدآء بل الهوی لأنّا بلوناک وجدناک فی معزل ان اعرف لحن القول و کن
من المتفرّسین

\end{otherlanguage}

An yā maliku an istamiʿ al-nidāʾa min hādhihi al-nāri al-mushtaʿilati
mina al-shajarati al-khaḍrāʾi fī hādhā al-Ṭūri al-murtafiʿi ʿalā
al-buqʿati al-muqaddasati al-bayḍāʾi khalfa Qulzumi al-Baqāʾi, innahu lā
ilāha illā anā al-Ghafūru al-Raḥīm. Ilā an qāla ʿazza kibriyāʾuhu: An yā
maliku innā samiʿnā minka kalimatan takallamta bihā idh saʾalaka maliku
al-Rūsi ʿammā quḍiya min ḥukmi al-ghazā, inna rabbaka la-Huwa al-ʿAlīmu
al-Khabīr. Qulta kuntu rāqidan fī al-mihādi ayqaẓanī nidāʾu al-ʿibādi
alladhīna ẓulimū ilā an ghariqū fī al-Baḥri al-Aswad. Kadhālika samiʿnā
wa-rabbuka ʿalā mā aqūlu shahīd. Nashhadu bi-annaka mā ayqaẓaka
al-nidāʾu bal al-hawā li-annā balawnāka wajadnāka fī maʿzil. An iʿrif
laḥna al-qawli wa-kun mina al-mutafarrisīn.

``O king, listen to the call from this fire blazing in the green Tree
upon this lofty Mount, in the holy, white Spot beyond the Ocean of
Eternity: there is no God but Me, the Ever-Forgiving, the Merciful.''
Until He says---exalted be His grandeur: ``O king, We heard a statement
you made when the king of Russia asked you about the decision to wage
war. Your Lord is indeed the All-Knowing, the All-Informed. You said, `I
was asleep in my bed when I was awakened by the cries of the people who
were wronged until they drowned in the Black Sea.' Thus did We hear, and
your Lord is witness to what I say. We testify that it was not their
cries that awakened you, but your own desire, for We tested you and
found you aloof. Discern the meaning of these words, and be among those
who perceive.''

\subsection{81}\label{section-3200}

\begin{otherlanguage}{arabic}

انّا ما نحبّ ان نرجع الیک کلمة سوء حفظاً لمقام الّذی اعطیناک فی الحیوة
الظّاهرة انّا اخترنا الأدب و جعلناه سجیّة المقرّبین انّه ثوب یوافق النّفوس من
کلّ صغیر و کبیر طوبی لمن جعله طراز هیکله ویل لمن جعل محروماً من هذا الفضل
العظیم لو کنت صاحب الکلمة ما نبذت کتاب اللّه ورآء ظهرک اذ ارسل الیک من
لدن عزیز حکیم

\end{otherlanguage}

Innā mā nuḥibbu an nurjiʿa ilayka kalimata sūʾin ḥifẓan li-maqāmi
alladhī aʿṭaynāka fī al-ḥayāti al-ẓāhirah. Innā ikhtarnā al-adaba
wa-jaʿalnāhu sajiyyata al-muqarrabīn. Innahu thawbun yuwāfiqu al-nufūsa
min kulli ṣaghīrin wa-kabīr. Ṭūbā li-man jaʿalahu ṭirāza haykalihi,
waylun li-man juʿila maḥrūman min hādhā al-faḍli al-ʿaẓīm. Law kunta
ṣāḥiba al-kalimati mā nabadhta Kitāba Allāhi warāʾa ẓahrika idh ursila
ilayka min ladun ʿAzīzin Ḥakīm.

``We do not wish to return an evil word to you, out of regard for the
station We have bestowed on you in this outward life. We have chosen
courtesy and made it the disposition of those brought near to Us. It is
a garment that befits every soul, great and small. Blessed is the one
who makes it the adornment of his person; woe to the one deprived of
this great bounty. Had you been true to your word, you would not have
cast the Book of God behind your back when it was sent to you by One who
is Almighty, All-Wise.''

\subsection{82}\label{section-3201}

\begin{otherlanguage}{arabic}

انّا بلوناک به ما وجدناک علی ما ادّعیت قم و تدارک ما فات عنک سوف تفنی
الدّنیا و ما عندک و یبقی الملک للّه ربّک و ربّ آبائک الأوّلین لا ینبغی لک ان
تقتصر الأمور علی ما تهوی به هواک اتّق زفرات المظلوم ان احفظه من سهام
الظّالمین بما فعلت تختلف الأمور فی مملکتک و یخرج الملک من کفّک جزآء عملک
اذاً تجد نفسک فی خسران مبین و یأخذ الزّلازل کلّ القبائل فی هناک الّا بأن
تقوم علی نصرة هذا الأمر و تتّبع الرّوح فی هذا السّبیل المستقیم أ عزّک غرّک
لعمری انّه لا یدوم و سوف یزول الّا بأن تتمسّک بهذا الحبل المتین قد نری
الذّلّة تسعی عن ورائک و انّک من الرّاقدین انتهی

\end{otherlanguage}

Innā balawnāka bihi mā wajadnāka ʿalā mā iddaʿayta. Qum wa-tadārak mā
fāta ʿanka. Sawfa tafnā al-dunyā wa-mā ʿindaka wa-yabqā al-mulku
li-Allāhi rabbika wa-rabbi ābāʾika al-awwalīn. Lā yanbaghī laka an
taqtaṣira al-umūra ʿalā mā tahwā bihi hawāka. Ittaqi zafarāti al-Maẓlūmi
an iḥfaẓhu min sihāmi al-ẓālimīn. Bi-mā faʿalta takhtalifu al-umūru fī
mamlakatika wa-yakhruju al-mulku min kaffika jazāʾa ʿamalika, idhan
tajidu nafsaka fī khusrānin mubīn. Wa-yaʾkhudhu al-zalāzilu kulla
al-qabāʾili fī hunāka illā bi-an taqūma ʿalā nuṣrati hādhā al-amri
wa-tattabiʿa al-Rūḥa fī hādhā al-sabīli al-mustaqīm. A-ʿizzuka gharraka?
La-ʿamrī innahu lā yadūmu wa-sawfa yazūlu illā bi-an tatamassaka
bi-hādhā al-ḥabli al-matīn. Qad narā al-dhillata tasʿā ʿan warāʾika
wa-innaka mina al-rāqidīn. Intahā.

``We tested you through it and did not find you as you had claimed to
be. Rise and make amends for what you have missed. The world and all you
possess will pass away, and sovereignty will remain with God, your Lord
and the Lord of your forefathers. You should not confine affairs to what
your own desire craves. Fear the sighs of the Wronged One; protect Him
from the arrows of the oppressors. Because of what you have done,
affairs in your kingdom will be thrown into confusion, and sovereignty
will slip from your grasp as recompense for your deeds. Then you will
find yourself in manifest loss. Tremors will seize all the peoples
there, unless you arise to support this Cause and follow the Spirit on
this straight path. Has your glory deceived you? By My life, it will not
endure and will soon vanish, unless you hold fast to this strong cord.
We see abasement hastening after you, while you are among those
asleep.'' End of quotation.

\subsection{83}\label{section-3202}

\begin{otherlanguage}{persian}

و معلوم احبّای الهی بوده در ایّام توقّف در ارض سرّ الواح منیعه مخصوص بعضی از
ملوک لاتمام حجّة الله نازل و ارسال شد از جمله بملک پاریس که در آن ایّام
رأس ملوک بود لوحی نازل و سبب آن آنکه روزی تلقاء عرش حاضر بودم فرمودند
بعد از دعوای روس و عثمانی ملک پاریس باعانت عثمانی برخاست و بعد از نزاع و
جدال و قتل و غارت جمعی از ملوک بمیان آمدند و حکم مصالحه محقّق شد بعد ملک
روس از ملک پاریس سؤال نمود که من و تو هر دو اهل یک ملّت بودیم سبب چه بود
که باهانت اهل مذهب خود و اعانت غیر مذهب قیام نمودی جواب نوشت که سببی
نداشت مگر آنکه نفوسی از رعیّت عثمانی را شما در بحر اسود بغتةً بر ایشان
هجوم نمودید و جمیع را غرق کردید ندای آن مظلومان مرا از خواب بیدار نمود و
باعانت برخاستم بعد از اتمام این فقره فرمودند حال ما لوحی باو میفرستیم و
او را امتحان مینمائیم اگر باعانت مظلومین اهل بیان برخاست تصدیق مینمائیم
او را در آنچه گفته والّا یظهر کذبه فیما ادّعی و قال لذا لوحی باو نازل و
ارسال شد ابداً خبری از او نرسید مع آنکه وزیری از وزرای او که سرّاً اظهار حبّ
مینمود بساحت اقدس معروض داشته که مخصوصاً لوح را بملک رساندم و تفصیل را هم
لساناً معروض داشتم مع‌ذلک جوابی نرسید این بود که بعد از عدم وصول جواب آن
لوح این لوح ثانی نازل و بخطّ فرنساوی شخصی نوشته و ارسال داشت حال ملاحظه
در لوح نمائید آنچه بعد بر او وارد شده از قبل تصریحاً من غیر تلویح در آن
لوح نازل طوبی للقارئین طوبی للمتفکّرین طوبی للمنصفین

\end{otherlanguage}

Va maʿlūm-i aḥibbā-yi Ilāhī būdih dar ayyām-i tavaqquf dar Arḍ-i Sirr
alvāḥ-i manīʿih makhṣūṣ-i baʿḍī az mulūk li-itmām-i ḥujjatu Allāh nāzil
va irsāl shud. Az jumlih bih malik-i Pārīs kih dar ān ayyām raʾs-i mulūk
būd lawḥī nāzil va sabab-i ān ānkih rūzī tilqāʾ-i ʿarsh ḥāḍir būdam
farmūdand baʿd az daʿvā-yi Rūs va ʿUthmānī malik-i Pārīs bih iʿānat-i
ʿUthmānī barkhāst va baʿd az nizāʿ va jidāl va qatl va ghārat jamʿī az
mulūk bih miyān āmadand va ḥukm-i muṣāliḥih muḥaqqaq shud. Baʿd malik-i
Rūs az malik-i Pārīs suʾāl namūd kih man va tu har du ahl-i yik millat
būdīm sabab chih būd kih bih ihānat-i ahl-i madhhab-i khud va iʿānat-i
ghayr-i madhhab qiyām namūdī. Javāb nivisht kih sababī nadāsht magar
ānkih nufūsī az raʿiyyat-i ʿUthmānī rā shumā dar Baḥr-i Asvad baghtatan
bar īshān hujūm namūdīd va jamīʿ rā gharq kardīd nidā-yi ān maẓlūmān
marā az khvāb bīdār namūd va bih iʿānat barkhāstam. Baʿd az itmām-i īn
faqarih farmūdand ḥāl mā lawḥī bih ū mīfiristīm va ū rā imtiḥān
mīnamāyīm agar bih iʿānat-i maẓlūmīn-i ahl-i Bayān barkhāst taṣdīq
mīnamāyīm ū rā dar ānchih guftih va-illā yaẓharu kadhibuhu fīmā iddaʿā
wa-qāl. Lidhā lawḥī bih ū nāzil va irsāl shud abadan khabarī az ū
narasīd maʿa ānkih vazīrī az vuzarā-yi ū kih sirran iẓhār-i ḥubb mīnamūd
bih sāḥat-i aqdas maʿrūḍ dāshtih kih makhṣūṣan lawḥ rā bih malik
rasāndam va tafṣīl rā ham lisānan maʿrūḍ dāshtam maʿa dhālik javābī
narasīd. Īn būd kih baʿd az ʿadam-i vuṣūl-i javāb-i ān lawḥ īn Lawḥ-i
Thānī nāzil va bih khaṭṭ-i faransāvī shakhṣī nivishtih va irsāl dāsht.
Ḥāl mulāḥiẓih dar lawḥ namāʾīd ānchih baʿd bar ū vārid shudih az qabl
taṣrīḥan min ghayri talwīḥ dar ān lawḥ nāzil. Ṭūbā li-l-qāriʾīn, ṭūbā
li-l-mutafakkirīn, ṭūbā li-l-munṣifīn.

The friends of God know that during the stay in the Land of Mystery,
exalted Tablets addressed specifically to certain kings were revealed
and dispatched to make God's proof complete. Among these was a Tablet
revealed for the king of Paris, who in those days stood foremost among
the kings. The reason for it was this: one day, when I was present
before the throne, He said, ``After the conflict between Russia and the
Ottoman Empire, the king of Paris arose to aid the Ottomans. Following
strife, warfare, killing, and pillage, a number of kings intervened, and
peace was concluded. Then the king of Russia asked the king of Paris,
`You and I both belonged to the same religion. Why did you rise to
humiliate the people of your own faith and aid those of another?' He
wrote in reply, `There was no reason other than that you suddenly
attacked some Ottoman subjects in the Black Sea and drowned them all.
The cries of those wronged people awakened me from sleep, and I arose to
help.'\,'' After completing this account, He said, ``Now We will send
him a Tablet and test him. If he arises to help the wronged people of
the Bayan, We will confirm the truth of what he has said; otherwise, his
falsehood in what he has claimed and said will become evident.''
Accordingly, a Tablet was revealed for him and dispatched, but no word
ever came from him. Although one of his ministers, who secretly
professed love, had submitted to the Most Holy Presence, ``I personally
delivered the Tablet to the king and also explained the details
orally,'' no reply came. Thus, when no answer to that Tablet arrived,
this second Tablet was revealed, and someone wrote it out in French and
sent it. Now examine the Tablet: what later befell him had already been
explicitly revealed in it, without allusion. Blessed are those who read;
blessed are those who reflect; blessed are the fair-minded.

\subsection{84}\label{section-3203}

\begin{otherlanguage}{persian}

و همچنین در خلق بدیع و ارسال آن بقدرت و عظمت تفکّر نمائید و ملاحظه کنید
بعد از اتمام حجّت و اظهار قدرت چگونه بلایا از قحط و غلا و خوف عجم را
احاطه نمود اخبار این امور کلّ در الواح متعدّده نازل شده فوالله الّذی لا اله
الّا هو اگر نفوس در خلق بدیع و ارسال او و لوح او و آنچه بعد وارد شده تفکّر
نمایند جمیع را کافیست و برهان الهی اظهر از شمس بر کلّ مبرهن و واضح میشود
ولکن کینونات جعلیّه بعالم خود مشغول‌اند از انفاس سبحانیّه و عرف آیات
رحمانیّه قسمتی نبرده و نخواهند برد الّا من شآء الله

\end{otherlanguage}

Va hamchunīn dar khalq-i Badīʿ va irsāl-i ān bih qudrat va ʿaẓamat
tafakkur namāʾīd va mulāḥiẓih kunīd baʿd az itmām-i ḥujjat va iẓhār-i
qudrat chigūnih balāyā az qaḥṭ va ghalā va khawf ʿAjam rā iḥāṭih namūd.
Ikhbār-i īn umūr kull dar alvāḥ-i mutaʿaddidih nāzil shudih.
Fa-wa-Allāhi alladhī lā ilāha illā Huwa, agar nufūs dar khalq-i Badīʿ va
irsāl-i ū va lawḥ-i ū va ānchih baʿd vārid shudih tafakkur namāyand
jamīʿ rā kāfīst va burhān-i Ilāhī aẓhar az shams bar kull mubarhan va
vāḍiḥ mīshavad. Va lākin kaynūnāt-i jaʿliyyih bih ʿālam-i khud
mashghūl-and az anfās-i subḥāniyyih va ʿarf-i āyāt-i raḥmāniyyih qismatī
naburdih va nakhvāhand burd illā man shāʾa Allāh.

Likewise, reflect on the creation of Badi and his dispatch with power
and grandeur, and consider how, after the proof had been made complete
and power had been manifested, afflictions---famine, high prices, and
fear---encompassed Persia. All these things were foretold in various
Tablets. By God, besides whom there is no other God, if people reflect
on the creation of Badi, his dispatch, his Tablet, and what subsequently
occurred, this will suffice for everyone, and the divine proof will
become demonstrably clear to all, more manifest than the sun. Yet these
counterfeit beings are absorbed in their own world. They have had and
will have no share in the holy breaths and the fragrance of the verses
of the All-Merciful, except those whom God wills.

\subsection{85}\label{section-3204}

\begin{otherlanguage}{persian}

چنانچه مشاهده میشود مع این اعلاء کبری و ظهور عظمی و آیات واضحات بعضی از
عباد متابعت نفوسی را اختیار نموده‌اند که همیشه خلف قناع بوده و خواهند بود
و ابداً امری از آن نفوس ظاهر نشده چه فایده که این عبد فرصت ندارد والّا
جمیع آنچه ظاهر شده و میشود کلّ را از آیات الله که بالتّصریح من غیر تأویل
نازل شده معروض میداشت نسأل الله بأن یؤیّد العباد علی الانصاف و محو ما
عندهم من الأوهام و ما ذکر فی الأسلاف

\end{otherlanguage}

Chenānche moshāhede mishavad maʿ-e in eʿlā-ye kobrā va ẓohūr-e ʿoẓmā va
āyāt-e vāżeḥāt baʿżi az ʿebād motābaʿat-e nofūsi rā ekhtiyār namūde-and
ke hamishe khalf-e qenāʿ būde va khvāhand būd va abadan amri az ān nofūs
ẓāher nashode che fāyede ke in ʿabd forṣat nadārad va-ellā jamiʿ-e ānche
ẓāher shode va mishavad koll rā az āyāt-e Allāh ke be'l-taṣriḥ men
gheyr-e taʾvil nāzel shode maʿrūż midāsht nasʾalu Allāha bi-an yuʾayyida
al-ʿibāda ʿalā al-inṣāfi wa-maḥwi mā ʿindahum mina al-awhāmi wa-mā
dhukira fī al-aslāf.

As can be seen, despite this supreme exaltation, this most mighty
revelation, and these clear verses, some people have chosen to follow
souls who have always remained, and will remain, behind a veil, and from
whom nothing has ever appeared. What a pity that this servant has no
time; otherwise, he would present everything that has happened and is
happening from the verses of God in which it was explicitly revealed,
without interpretation. We ask God to aid His servants in fairness and
in effacing the vain imaginings they hold and what was related among
former generations.

\subsection{86}\label{section-3205}

\begin{otherlanguage}{persian}

انّ العجب فی الّذین اتّبعوا الدّفرآء بعد الّذی تضوّع رائحة المحبوب بین
العالمین مثل آنکه بسیّد محمّد اصفهانی و اقوال او که ابداً از اصل امر مطّلع
نبوده و لازال بفساد مشغول از حقّ امنع اقدس محتجب مانده‌اند اگرچه قول عوام
است ولکن مناسب این مقام کوری نگر عصاکش کوری دگر شود فوالله فوالله هر ذی
بصر و ذی شمّی که یک مرتبه او را دیده از وجهش اثر جحیم و از نفسش رایحۀ اهل
سجّین یافته هر روز بتدلیسی مشغول چون وارد سجن اعظم شدیم اسمش را قدّوس
افندی گذاشت بگمان آنکه اسم بی‌معنی سبب اعلای مسمّی خواهد شد ویل له و لمن
اتّبعه و همچنین آقا جان را رئیس المشرکین سیف‌ الحقّ نامید و عراق را باو
وعده داد چنانچه مکتوب خود آقا جان الآن موجود است که برئیس المشرکین نوشته
و استغاثه نموده آن وعده که دادید چند وقت دیگر ظاهر میشود باری در این ارض
نزد هر شخصی ذکر نمود که من سیف‌ الحقّم و اکثر بلاد را عنقریب فتح میکنم الا
لعنة الله علی الکاذبین

\end{otherlanguage}

Inna al-ʿajaba fī alladhīna ittabaʿū al-dufarāʾa baʿda alladhī taḍawwaʿa
rāʾiḥatu al-Maḥbūbi bayna al-ʿālamīn methl-e ānke be-Seyyed Moḥammad-e
Eṣfahāni va aqvāl-e ū ke abadan az aṣl-e amr moṭṭaleʿ nabūde va lāzāl
be-fasād mashghūl az Ḥaqq-e amnaʿ-e aqdas moḥtajeb mānde-and agarche
qowl-e ʿavāmm ast va lāken monāseb-e in maqām kūri negar ʿaṣākash-e
kūri-ye degar shavad fa-wa-Allāhi fa-wa-Allāhi har dhi baṣar va dhi
shammi ke yek martabe ū rā dide az vajhash athar-e jaḥim va az nafasash
rāyeḥe-ye ahl-e sejjin yāfte har rūz be-tadlisi mashghūl chūn vāred-e
Sejn-e Aʿẓam shodim esmash rā Qoddūs Efendi gozāsht be-gomān-e ānke
esm-e bi-maʿni sabab-e eʿlā-ye mosammā khvāhad shod waylun lahu
wa-li-man ittabaʿahu va hamchenin Āqā Jān rā Raʾis-e Moshrekin Seyf
al-Ḥaqq nāmid va ʿErāq rā be-ū vaʿde dād chenānche maktūb-e khvod-e Āqā
Jān al-ān mowjūd ast ke be-Raʾis-e Moshrekin neveshte va esteghāthe
namūde ān vaʿde ke dādid chand vaqt-e digar ẓāher mishavad bāri dar in
arż nazd-e har shakhṣi dhekr namūd ke man Seyf al-Ḥaqqam va akthar-e
belād rā ʿanqarib fatḥ mikonam a-lā laʿnatu Allāhi ʿalā al-kādhibīn.

The wonder lies in those who followed the foul-smelling ones after the
fragrance of the Beloved had spread among the peoples of the world.
Thus, through Sayyid Muhammad of Isfahan and his words, they have
remained veiled from the most unapproachable, most holy Truth, although
he has never known the essence of the Cause and has always been occupied
with corruption. Though it is a common saying, it suits this occasion:
``See how one blind man becomes the guide of another.'' By God, by God,
anyone with sight and a sense of smell who has seen him even once has
found in his face a trace of hellfire and in his breath the odor of the
people of Sijjin. Every day he engages in some deception. When we
entered the Most Great Prison, he named himself Quddus Effendi,
imagining that a meaningless name would exalt its bearer. Woe to him and
to those who follow him! Likewise, the Chief of the Associators named
Aqa Jan ``Sword of Truth'' and promised him Iraq. A letter from Aqa Jan
himself still exists, written to the Chief of the Associators, in which
he pleaded, ``How much longer will it be before the promise you made is
fulfilled?'' In short, in this land he told everyone, ``I am the Sword
of Truth, and soon I will conquer most of the lands.'' Indeed, God's
curse rests upon the liars!

\subsection{87}\label{section-3206}

\begin{otherlanguage}{arabic}

در این مقام آیاتی از سماء مشیّت الهیّه نازل قوله عزّ کبریائه

\end{otherlanguage}

Dar in maqām āyāti az samā-ye mashiyyat-e elāhiyye nāzel qawluhu ʿazza
kibriyāʾuhu.

At this point, verses were revealed from the heaven of the divine will.
His words---exalted be His majesty:

\subsection{88}\label{section-3207}

\begin{otherlanguage}{arabic}

الأقدس الأعظم الأبهی

\end{otherlanguage}

Al-Aqdas al-Aʿẓam al-Abhā.

The Most Holy, the Most Great, the All-Glorious.

\subsection{89}\label{section-3208}

\begin{otherlanguage}{arabic}

انّ فی ابتلآء مالک الامکان فی کلّ الأحیان لآیات لمن فی الأکوان قد قبل
الشّدّة لرخآء البریّة و المشقّة لراحة من فی الامکان نفسی لفضله الفدآء و
کینونتی لرحمته الفدآء و روحی لعنایته الّتی احاطت الآفاق ما اصبح الّا و
احاطته ظلمات الاشارات من الّذین کفروا بالله منزل الآیات و انّه لا یمنعه
شیء عمّا اراد فی امر الله مالک یوم التّناد مرّةً ینادی بلسانه المبین و طوراً
یشیر باصبع الیقین و یدع الکلّ الی الله مالک الرّقاب لو نذکر ما ورد علینا
لتنفطر السّمآء و تخرّ الجبال

\end{otherlanguage}

Inna fī ibtilāʾi Māliki al-imkāni fī kulli al-aḥyāni la-āyātin li-man fī
al-akwān. Qad qabila al-shiddata li-rakhāʾi al-bariyyati
wa-al-mashaqqata li-rāḥati man fī al-imkān. Nafsī li-faḍlihi al-fidāʾu
wa-kaynūnatī li-raḥmatihi al-fidāʾu wa-rūḥī li-ʿināyatihi allatī aḥāṭati
al-āfāq. Mā aṣbaḥa illā wa-aḥāṭathu ẓulumātu al-ishārāti mina alladhīna
kafarū bi-Allāhi Munzili al-āyāt. Wa-innahu lā yamnaʿuhu shayʾun ʿammā
arāda fī amri Allāhi Māliki yawmi al-tanād. Marratan yunādī bi-lisānihi
al-mubīni wa-ṭawran yushīru bi-iṣbaʿi al-yaqīni wa-yadʿu al-kulla ilā
Allāhi Māliki al-riqāb. Law nadhkuru mā warada ʿalaynā la-tanfaṭiru
al-samāʾu wa-takhirru al-jibāl.

In the affliction of the Lord of all existence at every moment there are
signs for all created beings. He has accepted hardship for the
prosperity of humanity and suffering for the comfort of all who dwell in
the world. May my soul be sacrificed for His grace, my being for His
mercy, and my spirit for His loving care, which has encompassed the
horizons. No morning has dawned upon Him without His being surrounded by
the darkness of insinuations from those who disbelieve in God, the
Revealer of verses. Yet nothing prevents Him from accomplishing what He
has willed in the Cause of God, the Master of the Day of Summoning. At
one time He calls out with His lucid tongue, and at another He points
with the finger of certainty, summoning everyone to God, the Lord of
all. Were We to recount what has befallen Us, heaven would be rent
asunder and the mountains would collapse.

\subsection{90}\label{section-3209}

\begin{otherlanguage}{arabic}

انّ الّذین کفروا افتخروا بما عندهم من الألقاب انّ الأخرس سمّی نفسه بالقدّوس و
ادّعی فی نفسه ما ادّعی الخنّاس و الآخر سمّی نفسه بسیف الحقّ و قال انّی انا
فاتح البلاد قد بعث الله من ضرب علی فمه لیوقننّ الکلّ بأنّه ذنب الشّیطان قطع
من سیف الرّحمن قد کان ان ینتظر ایّام عزّه و ظهوره بما وعده من کفر بالله
فالق الأصباح کذلک یأخذ الله من اعرض عنه و قام علی تضییع امره بین العباد

\end{otherlanguage}

Inna alladhīna kafarū iftakharū bi-mā ʿindahum mina al-alqāb. Inna
al-akhrasa sammā nafsahu bi-al-Quddūsi wa-iddaʿā fī nafsihi mā iddaʿā
al-khannās. Wa-al-ākharu sammā nafsahu bi-Sayfi al-Ḥaqqi wa-qāla innī
anā fātiḥu al-bilād. Qad baʿatha Allāhu man ḍaraba ʿalā famihi
li-yūqinanna al-kullu bi-annahu dhanabu al-shayṭāni quṭiʿa min Sayfi
al-Raḥmān. Qad kāna an yantaẓira ayyāma ʿizzihi wa-ẓuhūrihi bi-mā
waʿadahu man kafara bi-Allāhi Fāliqi al-aṣbāḥ. Kadhālika yaʾkhudhu
Allāhu man aʿraḍa ʿanhu wa-qāma ʿalā taḍyīʿi amrihi bayna al-ʿibād.

Those who disbelieved took pride in the titles they bore. The mute one
named himself ``the Most Holy'' and made the same claims for himself as
did the stealthy whisperer. The other named himself ``Sword of Truth''
and said, ``I am indeed the conqueror of lands.'' God raised up someone
who struck him on the mouth, so that all might know with certainty that
he was the tail of Satan, severed by the Sword of the All-Merciful. He
had been awaiting the days of his glory and ascendancy, in accordance
with the promises made to him by one who disbelieved in God, the Cleaver
of the dawn. Thus does God seize whoever turns away from Him and arises
to undermine His Cause among His servants.

\subsection{91}\label{section-3210}

\begin{otherlanguage}{persian}

فلمّا هلکوا سرت اریاح الرّبیع و فتحت ابواب السّمآء و امطر السّحاب طوبی لمن
فاز بعرفان الله فی ایّامه و انقطع بکلّه عن کلّ الجهات قل أ و لم یکفکم ربّ
السّموات و الأرض انّه قد اتی بالحقّ باسمه المهیمن علی الابداع انّک نوّر قلبک
بمصباح الأعظم الّذی اوقده مالک القدم ثمّ استقم علی الأمر بسلطان ربّک
المقتدر المختار

\end{otherlanguage}

Fa-lammā halakū sarat aryāḥu al-rabīʿi wa-futiḥat abwābu al-samāʾi
wa-amṭara al-saḥāb. Ṭūbā li-man fāza bi-ʿirfāni Allāhi fī ayyāmihi
wa-inqaṭaʿa bi-kullihi ʿan kulli al-jihāt. Qul a-wa-lam yakfikum Rabbu
al-samāwāti wa-al-arḍi innahu qad atā bi-al-ḥaqqi bi-ismihi al-Muhaymini
ʿalā al-ibdāʿ. Innaka nawwir qalbaka bi-Miṣbāḥi al-Aʿẓami alladhī
awqadahu Māliku al-qidam. Thumma istaqim ʿalā al-amri bi-sulṭāni rabbika
al-Muqtadiri al-Mukhtār.

When they perished, the winds of spring began to blow, the gates of
heaven were opened, and the clouds poured down rain. Blessed is whoever
has attained the recognition of God in His days and detached himself
wholly from every direction. Say: Has the Lord of the heavens and the
earth not sufficed you? He has indeed come with the truth in His name,
the All-Protecting over all creation. Illuminate your heart with the
Most Great Lamp, which the Lord of eternity has kindled. Then stand firm
in the Cause through the sovereignty of your Lord, the All-Powerful, the
Unconstrained.

\subsection{92}\label{section-3211}

\begin{otherlanguage}{persian}

انتهی شخصی در این مدینه از علمای مدینۀ کبیره بوده آقا جان نزد او مراوده
داشت نزد او هم تفصیل را ذکر نموده که من سیف الحقّم عنقریب فتوحات اکثر
بلاد بدست من خواهد شد آن شخص عالم جاهل اگرچه از سبیل حقّ بعید بوده کلمۀ
خوبی ذکر نموده کأنّ ر وح ‌القدس نطق علی لسانه چنانچه بعد از قتل مشرکین
فی‌‌الفور مع پاشای بلد بر سر نعش آقا جان حاضر مذکور نمود این ملعون میگفت
من سیف الحقّم حال معلوم شد که سیف الحقّ آن بود که بر کمرش خورد و بأسفل
الجحیم مقرّش داد

\end{otherlanguage}

Entahā shakhṣi dar in madine az ʿolamā-ye Madine-ye Kabire būde Āqā Jān
nazd-e ū morāvede dāsht nazd-e ū ham tafṣil rā dhekr namūde ke man Seyf
al-Ḥaqqam ʿanqarib fotūḥāt-e akthar-e belād be-dast-e man khvāhad shod
ān shakhṣ-e ʿālem-e jāhel agarche az sabil-e Ḥaqq baʿid būde kaleme-ye
khūbi dhekr namūde ka-anna Rūḥa al-Qudusi naṭaqa ʿalā lisānihi chenānche
baʿd az qatl-e moshrekin fi'l-fowr maʿ-e pāshā-ye balad bar sar-e
naʿsh-e Āqā Jān ḥāżer madhkūr namūd in malʿūn migoft man Seyf al-Ḥaqqam
ḥāl maʿlūm shod ke Seyf al-Ḥaqq ān būd ke bar kamarash khvord va
be-asfal-e jaḥim maqarresh dād.

End of quotation. There was a man in this city who had been one of the
religious scholars of the Great City. Aqa Jan used to associate with him
and had told him the details as well: ``I am the Sword of Truth. Soon
most of the lands will be conquered at my hands.'' Though far removed
from the path of Truth, that ignorant scholar made a good remark, as
though the Holy Spirit had spoken through his tongue. Immediately after
the killing of the associators, he came with the governor of the town to
Aqa Jan's body and said, ``This accursed man used to say, `I am the
Sword of Truth.' Now it is clear that the Sword of Truth was the one
that struck his back and consigned him to the lowest depths of hell.''

\subsection{93}\label{section-3212}

\begin{otherlanguage}{arabic}

باری الحمد لله بعد از وقوع این فقره و رجوع مشرکین بأسفل الجحیم امطار
رحمت لیلاً و نهاراً باریده بعد از آنکه در چند سنه رحمت ممنوع بود و ناس
بقحط و غلا معذّب و مبتلا

\end{otherlanguage}

Bāri al-ḥamdu li-Allāh baʿd az voqūʿ-e in faqare va rojūʿ-e moshrekin
be-asfal-e jaḥim amṭār-e raḥmat laylan va nahāran bāride baʿd az ānke
dar chand sane raḥmat mamnūʿ būd va nās be-qaḥṭ va ghalā moʿadhdhab va
mobtalā.

In short, praise be to God, after this event and the return of the
associators to the lowest depths of hell, the rains of mercy have fallen
night and day, after mercy had been withheld for several years and the
people had suffered the torment of famine and high prices.

\subsection{94}\label{section-3213}

\begin{otherlanguage}{persian}

باری وقتی که در ادرنه بودیم کاغذی آن خبیث بشخصی نوشته بود از جمله ملاحظه
شد یک فقره از فقرات لوحی که در عراق نازل شده بود سرقت نموده و باسم خود
نوشته و آن فقره این است چون شمس مشرقیم و چون قمر لائح و آن غافل پلید این
قدر ادراک ننموده که عرصۀ سیمرغ جولانگه ذباب نشده و نخواهد شد و ظلمت را
نمی‌رسد که دعوی شمسی نماید فوالله چون بعض ناس را ضعیف و احمق دید لذا
بمفتریات نفسیّه قیام نمود چنانچه بعضی مجعولات قلمیّه او را اخذ نموده و
میخوانند قد خسر الکاتب و القاری و بوساوس و دسائس مختلفه ناس را از ربّ
النّاس منع مینمود و باوهام سابقه بعضی را گمراه نمود

\end{otherlanguage}

Bāri vaqti ke dar Edrene būdim kāghadhi ān khabith be-shakhṣi neveshte
būd az jomle molāḥeẓe shod yek faqare az faqarāt-e lowḥi ke dar ʿErāq
nāzel shode būd serqat namūde va be-esm-e khvod neveshte va ān faqare in
ast chūn shams moshreqim va chūn qamar lāʾeḥ va ān ghāfel-e palid in
qadr edrāk nanamūde ke ʿarṣe-ye Simorgh jowlāngah-e dhobāb nashode va
nakhvāhad shod va ẓolmat rā nemirasad ke daʿvā-ye shamsi namāyad
fa-wa-Allāhi chūn baʿżi nās rā żaʿif va aḥmaq did ledhā be-moftarayāt-e
nafsiyye qiyām namūd chenānche baʿżi majʿūlāt-e qalamiyye-ye ū rā akhdh
namūde va mikhvānand qad khasira al-kātibu wa-al-qāriʾu va be-vasāves va
dasāʾes-e mokhtalefe nās rā az Rabb-e nās manʿ minamūd va be-avhām-e
sābeqe baʿżi rā gomrāh namūd.

In short, when we were in Edirne, that vile man had written a letter to
someone. Among other things, it was noticed that he had stolen a passage
from a Tablet revealed in Iraq and written it in his own name. The
passage is this: ``We shine like the sun and gleam like the moon.'' That
heedless, foul man had not understood even this much: the realm of the
Simurgh has never been, nor will it ever be, a playground for flies, and
darkness has no right to claim to be the sun. By God, because he saw
that some people were weak and foolish, he set about fabricating claims
of his own. Thus, some have taken up his forged writings and read them.
Both the writer and the reader have indeed lost. Through various
whisperings and intrigues he kept people from the Lord of humanity, and
through former delusions he led some astray.

\subsection{95}\label{section-3214}

\begin{otherlanguage}{persian}

چنانچه این بنده را شخصی در سوق ملاقات نمود مذکور داشت که ساعتی میخواهم
تو را ملاقات کنم بشرط آنکه احدی جز من و تو نباشد و این عبد از قبل او را
ندیده بودم گفتم بسیار خوب وقتی معیّن شد و آمد بعضی ذکرها بمیان آمد و
معلوم شد که با مشرک بالله مراوده دارد و ذکر نمود بمن گفته‌اند تو از
شیعیانی در این اثنا این عبد را بشأنی ضحک غلبه نمود که زمام صبر از دست
رفت آن شخص تعجّب نمود گفت سبب ضحک چیست گفتم ای بیچارۀ فقیر شیعیان شما که
در دیار ایران مثل حصاة ریخته‌اند بچه مقامی رسیدند و یا چه شأنی عندالله
داشته‌اند که تازه تو میخواهی بر اثر آن متوهّمین مردوده مشی نمائی آیا ندیدی
که کلّ باطل و در ضلالت صرف بوده‌اند چنانچه حقّ منیع را بایادی خود شهید
کردند الا لعنة الله علی القوم الظّالمین

\end{otherlanguage}

Chenānche in bande rā shakhṣi dar sūq molāqāt namūd madhkūr dāsht ke
sāʿati mikhvāham to rā molāqāt konam be-sharṭ-e ānke aḥadi joz man va to
nabāshad va in ʿabd az qabl ū rā nadide būdam goftam besyār khūb vaqti
moʿayyan shod va āmad baʿżi dhekrhā be-miyān āmad va maʿlūm shod ke bā
moshrek be-Allāh morāvede dārad va dhekr namūd be-man gofte-and to az
shiʿiyāni dar in athnā in ʿabd rā be-shaʾni żeḥk ghalabe namūd ke
zemām-e ṣabr az dast raft ān shakhṣ taʿajj ob namūd goft sabab-e żeḥk
chist goftam ey bichāre-ye faqir shiʿiyān-e shomā ke dar diyār-e Irān
methl-e ḥaṣāt rikhte-and be-che maqāmi residand va yā che shaʾni ʿinda
Allāh dāshte-and ke tāze to mikhvāhi bar athar-e ān motavahhemin-e
mardūde mashy namāʾi āyā nadidi ke koll bāṭel va dar żalālat-e ṣerf
būde-and chenānche Ḥaqq-e maniʿ rā be-ayādi-ye khvod shahid kardand a-lā
laʿnatu Allāhi ʿalā al-qawmi al-ẓālimīn.

For example, a man met this servant in the marketplace and said, ``I
would like to meet with you for an hour, provided that no one is there
but you and me.'' I had never seen him before. I said, ``Very well.'' A
time was arranged, and he came. Various matters came up in conversation,
and it became clear that he associated with the one who joined partners
with God. He said, ``I have been told that you are a Shia.'' At this,
such laughter overcame this servant that I lost all self-control. The
man was astonished and asked, ``Why are you laughing?'' I said, ``You
poor, helpless man! Your Shias, who are scattered like pebbles
throughout Persia---what station have they attained, or what standing
have they had before God, that you should now wish to follow in the
footsteps of those rejected, deluded people? Have you not seen that they
were all in error and utter misguidance, so much so that they martyred
the exalted Truth with their own hands? Indeed, God's curse rests upon
the people who do wrong!''

\subsection{96}\label{section-3215}

\begin{otherlanguage}{persian}

بعد از القاء این کلمه بسیار متفکّر شد بعد لوجه الله بعضی از کلمات منیعه
که از مصدر امریّه استماع نموده بودم باو القا نمودم مشاهده شد که خالی از
استعداد نیست از جمله مذکور داشتم که تو فکر کن از آنچه نزد شیعیان بود و
از اساس مذهب و ملّت خود میشمردند کدام یک صدق بود که حال تو میخواهی بر اثر
آن اوهام و کلمات کذبه به بحر صدق و حکمت ربّانیّه وارد شوی آیا این شیعیان
که میگوئی معنی قیامت را ادراک نموده بودند قال لا گفتم آیا میزان را عارف
شدند قال لا گفتم آیا حشر و نشر را فهمیدند قال لا گفتم آیا آنچه در ذکر
قائم نزدشان مذکور است بقسمی که ادراک نموده‌اند حقّ بوده قال لا بعد مذکور
داشتم که خود شاهدی کلّ کذب بود حال این فقره که بتو گفته‌اند از کجا دانستی
که صدق است

\end{otherlanguage}

Baʿd az elqā-ye in kaleme besyār motafakker shod baʿd li-vajh-e Allāh
baʿżi az kalemāt-e maniʿe ke az maṣdar-e amriyye estemāʿ namūde būdam
be-ū elqā namūdam moshāhede shod ke khāli az esteʿdād nist az jomle
madhkūr dāshtam ke to fekr kon az ānche nazd-e shiʿiyān būd va az asās-e
madhhab va mellat-e khvod mishomordand kodām yek ṣedq būd ke ḥāl to
mikhvāhi bar athar-e ān avhām va kalemāt-e kadhibe be-baḥr-e ṣedq va
ḥekmat-e rabbāniyye vāred shavi āyā in shiʿiyān ke migūʾi maʿni-ye
qiyāmat rā edrāk namūde būdand qāla lā goftam āyā mizān rā ʿāref shodand
qāla lā goftam āyā ḥashr va nashr rā fahmidand qāla lā goftam āyā ānche
dar dhekr-e Qāʾem nazdashān madhkūr ast be-qesmi ke edrāk namūde-and
ḥaqq būde qāla lā baʿd madhkūr dāshtam ke khvod shāhedi koll kedhb būd
ḥāl in faqare ke be-to gofte-and az kojā dānesti ke ṣedq ast.

After these words were spoken, he became very thoughtful. Then, for
God's sake, I conveyed to him some of the exalted words I had heard from
the Source of the Cause. It became apparent that he was not without
capacity. Among other things, I said, ``Consider: of all that the Shias
held and counted among the foundations of their religion and faith,
which was true, that you should now wish to enter the ocean of truth and
divine wisdom by following those delusions and lying words? Had these
Shias you speak of understood the meaning of the Resurrection?'' He
said, ``No.'' I asked, ``Had they recognized the Balance?'' He said,
``No.'' I asked, ``Had they understood the gathering and rising of the
dead?'' He said, ``No.'' I asked, ``Was what they related concerning the
Qaim true in the way they understood it?'' He said, ``No.'' Then I said,
``You yourself bear witness that it was all false. Now, how have you
ascertained that this thing they have told you is true?''

\subsection{97}\label{section-3216}

\begin{otherlanguage}{persian}

بعد جلست مستویاً مقابلاً الیه و نطقت بما حفظت من آیات الله الملک المهیمن
القیّوم و القیت علیه ما امرت به من لدی الله العزیز المحبوب قلت یا عبد اما
سمعت ما نطق به لسان العظمة اذ استوی علی عرشه المهیمن علی کلّ شاهد و مشهود
قال و قوله الحقّ دعوا ما عندکم من الأوهام ثمّ استمعوا ما ینطق به لسان ربّکم
العزیز العلّام الی متی تتّبعون الهوی قد اشرقت شمس الهدی ان اقبلوا الیها
مقدّسین عمّا ذکر من قبل من علمائکم تالله هذا لظهور الله و آیة بطونه لما
سواه لا یقاس بما ذکر فی ازل الآزال و لا یعرف بما عندکم من کلمات اهل
الجدال تقرّبوا الیه بعیون نورآء و وجوه بیضآء کذلک امرتم فی الألواح من لدی
الله المقتدر العلیم الخبیر و قوله جلّ کبریائه مخاطباً لأحد قد نراک متغمّساً
فی غدیر المحتجبین و میاه اوهام الغافلین ان اخرج منه باسمی و سلطانی ثمّ
تغمّس فی هذا البحر الأعظم الّذی ینطق کلّ قطرة منه لا اله الّا انا المقتدر
المتعالی المهیمن العزیز الکریم

\end{otherlanguage}

Baʿd jalastu mustawiyan muqābilan ilayhi wa-naṭaqtu bi-mā ḥafiẓtu min
āyāti Allāhi al-Maliki al-Muhaymini al-Qayyūm. Wa-alqaytu ʿalayhi mā
umirtu bihi min ladā Allāhi al-ʿAzīzi al-Maḥbūb. Qultu yā ʿabdu a-mā
samiʿta mā naṭaqa bihi Lisānu al-ʿAẓamati idh istawā ʿalā ʿarshihi
al-muhaymini ʿalā kulli shāhidin wa-mashhūd. Qāla wa-qawluhu al-ḥaqqu
daʿū mā ʿindakum mina al-awhāmi thumma istamiʿū mā yanṭiqu bihi lisānu
rabbikumu al-ʿAzīzi al-ʿAllām. Ilā matā tattabiʿūna al-hawā qad ashraqat
shamsu al-hudā an aqbilū ilayhā muqaddasīna ʿammā dhukira min qablu min
ʿulamāʾikum. Ta-Allāhi hādhā la-ẓuhūru Allāhi wa-āyatu buṭūnihi li-mā
siwāhu lā yuqāsu bi-mā dhukira fī azali al-āzāli wa-lā yuʿrafu bi-mā
ʿindakum min kalimāti ahli al-jidāl. Taqarrabū ilayhi bi-ʿuyūnin nawrāʾa
wa-wujūhin bayḍāʾa kadhālika umirtum fī al-alwāḥi min ladā Allāhi
al-Muqtadiri al-ʿAlīmi al-Khabīr. Wa-qawluhu jalla kibriyāʾuhu
mukhāṭiban li-aḥadin qad narāka mutaghammisan fī ghadīri al-muḥtajibīna
wa-miyāhi awhāmi al-ghāfilīn. An ukhruj minhu bi-ismī wa-sulṭānī thumma
taghammas fī hādhā al-baḥri al-aʿẓami alladhī yanṭiqu kullu qaṭratin
minhu lā ilāha illā anā al-Muqtadiru al-Mutaʿālī al-Muhayminu al-ʿAzīzu
al-Karīm.

Then I sat upright facing him and recited what I had memorized of the
verses of God, the King, the All-Protecting, the Self-Subsisting. I
conveyed to him what I had been commanded by God, the Almighty, the
Beloved. I said, ``O servant! Have you not heard what the Tongue of
Grandeur uttered when He took His seat upon His throne, which holds sway
over every witness and all that is witnessed? He said---and His word is
truth---`Abandon the delusions you hold, then listen to what the tongue
of your Lord, the Almighty, the All-Knowing, proclaims. How long will
you follow desire? The Sun of guidance has risen. Turn toward it,
sanctified from what your religious scholars have said in the past. By
God, this is indeed the manifestation of God and the sign of His
concealment to all besides Him. It cannot be compared with anything
mentioned from all eternity, nor can it be known through the words of
the disputatious that you possess. Draw near to Him with radiant eyes
and shining faces. Thus have you been commanded in the Tablets by God,
the All-Powerful, the All-Knowing, the All-Informed.' And His
words---glorious be His majesty---addressed to someone: `We see you
immersed in the pool of those who are veiled and in the waters of the
delusions of the heedless. Come out of it through My name and My
sovereignty, then immerse yourself in this Most Great Ocean, every drop
of which proclaims: There is no god but Me, the All-Powerful, the Most
Exalted, the All-Protecting, the Almighty, the All-Bountiful.'\,''

\subsection{98}\label{section-3217}

\begin{otherlanguage}{persian}

گفتم ای مرد لوجه الله میگویم و خالصاً لله ذکر مینمایم بریز این اوهام را و
از این مقامات کثیفۀ محدودۀ متوهّمه صعود نما تا به پرهای رحمت رحمانی
بسموات حکم ربّانی عروج نمائی و فائز شوی الیوم یوم قل الله ثمّ ذرهم فی
خوضهم یلعبون است و این آیۀ مبارکه از قبل مخصوص این یوم نازل شده چه که در
این ظهور واضح و مبرهن گشت که نفس حقّ وحده سلطان است بر کلّ و شریک و شبیه
اخذ ننموده و نخواهد نمود اگرچه از قبل اهل ملل به این بیان مقرّ و معترف
بودند ولکن اکثر کاذب بودند چه که هر روز شریکی از برای حقّ قرار میدادند

\end{otherlanguage}

Goftam ey mard, li-wajhi Allāh mīgūyam va khāliṣan li-llāh dhekr
mīnamāyam: berīz īn awhām rā va az īn maqāmāt-e kathīfe-ye maḥdūde-ye
motavahhame ṣoʿūd namā tā be parhā-ye raḥmat-e raḥmānī be samāvāt-e
ḥekam-e rabbānī ʿorūj namāʾī va fāʾez shavī. Al-yawm yawm-e ``qul Allāh
thumma dharhum fī khawḍihim yalʿabūn'' ast va īn āye-ye mobārake az qabl
makhṣūṣ-e īn yawm nāzel shode, che ke dar īn ẓohūr vāżeḥ va mobarhan
gasht ke nafs-e Ḥaqq vaḥdahu solṭān ast bar koll va sharīk va shabīh
akhdh nanamūde va nakhāhad namūd. Agarche az qabl ahl-e melal be īn
bayān moqerr va moʿtaref būdand, valīkan akthar kādheb būdand, che ke
har rūz sharīkī az barā-ye Ḥaqq qarār mīdādand.

I said, ``Man, I speak for the sake of God and make mention of this
purely for God: cast away these delusions and rise above these foul,
limited, imaginary stations, so that on the wings of the mercy of the
All-Merciful you may ascend to the heavens of divine wisdom and attain
your goal. Today is the day of `Say: God! Then leave them to play in
their vain discourse.' This blessed verse was revealed in the past
specifically for this day, for in this Revelation it has become clear
and manifest that God Himself alone is sovereign over all and has never
taken, nor will ever take, a partner or a likeness. Although the
followers of religions previously acknowledged and professed this
statement, most were lying, for every day they assigned a partner to
God.''

\subsection{99}\label{section-3218}

\begin{otherlanguage}{persian}

و بعد در این مقام آنچه سزاوار بود این عبد معروض داشت ذکر نمود الیوم
تکلیف چیست گفتم عمل بهمین آیه که عرض کردم بگذار این انفس موهومه و
اقوالشان را و در فضای خوش بدیع وارد شو تا خمر باقی از آیۀ رحمت رحمانی که
میفرماید بظهوری ثبت حکم البدع و انّ هذا لبدیع السّموات و الأرضین بیاشامی
از قبل و بعد بگذر بنفس ظهور بعین ظهور ناظر شو چه که کلّ در این ظهور اعظم
به این مأمورند و این مختصّ باین ظهور است و بعد ذکر نمودم که بعضی از نفوس
ضعیفه را آن نفس خبیثه بامثال این اذکار از مختار منع نمود چنانچه یحیی هم
به بعضی مینویسد انت من شیعتی و مقصود از این عبارت معلوم قد ضلّ کلّ کتّاب
کذّاب حال مشاهده کن کجاست مقام اذکار آن نفس خبیث و مقامی که حضرت ربّ
الأرباب از برای عباد خواسته ای برادر جمیع الفاظ محدودۀ نالایقۀ قبل را
بریز و بپرهای بدیع در فضای خوش بدیع طایر شو تا بعنایت الهی از خمر بدیع
بیاشامی و بسرّ این امر بدیع پی‌بری و دیگر آنکه بعضی از عباد از بعضی اعمال
نفوس ضعیفه که طائف حولند شکایت نموده‌اند و آن را بحقّ جلّ و عزّ نسبت داده‌اند
چنانچه شخصی ذکر نموده که چگونه میشود مع اظهار حقّیّت و اعلاء این امر که
عالم را احاطه نموده بعضی نفوس که خود را نسبت بحقّ میدهند باعمال ناشایسته
مشغول باشند باری این فقره از غفلت آن غافل است گویا این شعر که مابین ناس
مشهور است نشنیده که میگویند گر جمله کاینات کافر گردد بر دامن کبریاش
ننشیند گرد

\end{otherlanguage}

Va baʿd dar īn maqām ānche sazāvār būd īn ʿabd maʿrūż dāsht. Dhekr
namūd: al-yawm taklīf chīst? Goftam: ʿamal be hamīn āye ke ʿarż kardam.
Begodhār īn anfos-e mawhūme va aqvāl-eshān rā va dar fażā-ye khosh-e
badīʿ vāred show tā khamr-e bāqī az āye-ye raḥmat-e raḥmānī ke
mīfarmāyad ``bi-ẓuhūrī thabata ḥukmu al-badʿ wa-inna hādhā la-badīʿu
al-samāwāti wa-al-arḍīn'' biyāshāmī. Az qabl va baʿd begodhar, be nafs-e
ẓohūr be ʿeyn-e ẓohūr nāẓer show, che ke koll dar īn ẓohūr-e aʿẓam be īn
maʾmūrand va īn mokhtaṣṣ be īn ẓohūr ast. Va baʿd dhekr namūdam ke baʿżī
az nofūs-e żaʿīfe rā ān nafs-e khabīthe be amthāl-e īn adhkār az Mokhtār
manʿ namūd, chenānche Yaḥyā ham be baʿżī mīnevīsad ``anta min shīʿatī''
va maqṣūd az īn ʿebārat maʿlūm. Qad ḍalla kullu kattābin kadhdhāb. Ḥāl
moshāhede kon kojāst maqām-e adhkār-e ān nafs-e khabīth va maqāmī ke
Ḥażrat-e Rabb al-Arbāb az barā-ye ʿebād khāste. Ey barādar, jamīʿ-e
alfāẓ-e maḥdūde-ye nālāyeqe-ye qabl rā berīz va be parhā-ye badīʿ dar
fażā-ye khosh-e badīʿ ṭāyer show tā be ʿenāyat-e elāhī az khamr-e badīʿ
biyāshāmī va be serr-e īn amr-e badīʿ pey barī. Va dīgar ānke baʿżī az
ʿebād az baʿżī aʿmāl-e nofūs-e żaʿīfe ke ṭāʾef-e ḥowland shekāyat
namūde-and va ān rā be Ḥaqq jalla va ʿazza nesbat dāde-and, chenānche
shakhṣī dhekr namūde ke chegūne mīshavad maʿ-e eẓhār-e ḥaqqīyat va
eʿlā-ye īn amr ke ʿālam rā eḥāṭe namūde baʿżī nofūs ke khod rā nesbat be
Ḥaqq mīdahand be aʿmāl-e nāshāyeste mashghūl bāshand. Bārī īn faqare az
ghaflat-e ān ghāfel ast. Gūyā īn sheʿr ke mābeyn-e nās mashhūr ast
nashenīde ke mīgūyand: gar jomle kāʾenāt kāfer gardad, bar dāman-e
kebriyāsh naneshīnad gard.

Then, on this subject, this servant submitted what was fitting. He
asked, ``What is our duty today?'' I said, ``To act according to this
very verse that I have cited. Leave these deluded souls and their
sayings behind, and enter the delightful, wondrous expanse, so that you
may drink the everlasting wine from the verse of the mercy of the
All-Merciful, which says: `Through My manifestation the rule of new
creation has been established, and this is indeed the Originator of the
heavens and the earths.' Pass beyond before and after, and look upon the
Manifestation Himself with the eye of the Manifestation, for all are
commanded to do this in this Most Great Revelation, and this is peculiar
to this Revelation.'' Then I mentioned that that wicked soul had used
such statements to keep some weak souls away from Him who acts as He
chooses, just as Yaḥyā also writes to some, ``You are among my
followers,'' and the purpose of this expression is clear. Every lying
scribbler has indeed gone astray. Now consider how far apart are the
station of the statements of that wicked soul and the station that the
Lord of Lords has desired for His servants. Brother, cast away all the
limited, unworthy expressions of the past, and soar on wondrous wings in
the delightful, wondrous expanse, so that through divine favor you may
drink the wondrous wine and discover the secret of this wondrous Cause.
Furthermore, some servants have complained about certain deeds of the
weak souls who circle around Him, and have attributed them to God,
exalted and glorified be He. For example, someone has said, ``How is it
possible that, despite the declaration of the truth and the exaltation
of this Cause, which has encompassed the world, some souls who claim a
connection to God should occupy themselves with unworthy deeds?'' This
remark, however, arises from that heedless person's heedlessness. It
seems he has not heard this verse, which is well known among the people:
``Were all created things to become unbelievers, no dust would settle on
the hem of His grandeur.''

\subsection{100}\label{section-3219}

\begin{otherlanguage}{persian}

روزی در ساحت عرش حاضر بودم قد توجّه الیّ وجه الله قال أ سمعت انّ المخلصین
فی خطر عظیم ای عبد حاضر لدی العرش علمای ظاهره در اصحاب رسول اختلاف
نموده‌اند که آیا اسم صحابه بر چه نفسی صادقست بعضی گفته‌اند صحابی نفوسی
هستند که در یک سنه او اکثر در خدمت آن حضرت بودند و بجهاد فی سبیل الله
قیام نمودند و بعضی گفته‌‌اند که هر کس اقرار بر کلمۀ توحید نمود و رسول را
ملاقات کرد او از صحابه محسوب است اگرچه مرّة واحده بوده و بعضی گفته‌اند این
اسم در بارۀ نفوسی صادقست که مخصوص حضرت رسول او را باین اسم خطاب فرموده‌ و
بعضی گفته‌اند نفوسی هستند که عند حضرت موثّق بوده‌اند و در سفر و حضر حاضر
ولکن اکثری از علما گفته‌اند کلّ من اسلم ورآء النّبیّ صلّی الله علیه و صحبه
ولو اقلّ زمان انّه من الصّحابة از این قرار در حجّة الوداع چهل‌هزار نفس با
حضرت بودند و یوم وفات حضرت در مدینه صد و بیست و چهار هزار نفر جمع شدند
بر کلّ این نفوس بقول اخیر اسم صحابه صادق مع‌ذلک معدودی بودند که از زلال
خمر ایقان نوشیده‌اند و بمبدء فیوضات رحماینّه وارد شده‌اند باری امثال این
سخنان از غفلت شده و میشود نسأل الله ان یؤیّد الکلّ علی ما یحبّ و یرضی باری
ای دوستان حقّ بسمع قناعت نکنید و باوهام قبل دل مبندید فوالله الیوم نفوس
خبیثه ضعفا را باوهام قبلیّه از شریعۀ الهیّه منع نموده‌اند

\end{otherlanguage}

Rūzī dar sāḥat-e ʿarsh ḥāżer būdam. Qad tawajjaha ilayya wajhu Allāh,
qāla: ``A-samiʿta inna al-mukhliṣīna fī khaṭarin ʿaẓīm?'' Ey ʿabd-e
ḥāżer laday al-ʿarsh, ʿolamā-ye ẓāhere dar aṣḥāb-e Rasūl ekhtelāf
namūde-and ke āyā esm-e ṣaḥābe bar che nafsī ṣādeq ast. Baʿżī gofte-and
ṣaḥābī nofūsī hastand ke dar yek sane aw akthar dar khedmat-e ān Ḥażrat
būdand va be jihād fī sabīli Allāh qiyām namūdand, va baʿżī gofte-and ke
har kas eqrār bar kalame-ye towḥīd namūd va Rasūl rā molāqāt kard ū az
ṣaḥābe maḥsūb ast agarche marratan wāḥidatan būde, va baʿżī gofte-and īn
esm dar bāre-ye nofūsī ṣādeq ast ke makhṣūṣ Ḥażrat-e Rasūl ū rā be īn
esm kheṭāb farmūde, va baʿżī gofte-and nofūsī hastand ke ʿend-e Ḥażrat
movaththaq būde-and va dar safar va ḥażar ḥāżer. Valīkan aktharī az
ʿolamā gofte-and: ``Kullu man aslama {[}corrupted text{]} al-nabiyya
ṣallā Allāhu ʿalayhi wa-ṣaḥibahu wa-law aqalla zamānin innahu mina
al-ṣaḥāba.'' Az īn qarār dar Ḥijjat al-Vadāʿ chehel-hezār nafs bā Ḥażrat
būdand va yawm-e vafāt-e Ḥażrat dar Madīne ṣad va bīst va chahār hezār
nafar jamʿ shodand. Bar koll-e īn nofūs be qowl-e akhīr esm-e ṣaḥābe
ṣādeq. Maʿa-dhālek maʿdūdī būdand ke az zolāl-e khamr-e īqān nūshīde-and
va be mabdaʾ-e foyūżāt-e {[}corrupted text{]} vāred shode-and. Bārī
amthāl-e īn sokhanān az ghaflat shode va mīshavad. Nasʾalu Allāha an
yuʾayyida al-kulla ʿalā mā yuḥibbu wa-yarḍā. Bārī ey dūstān-e Ḥaqq, be
samʿ qanāʿat nakonīd va be awhām-e qabl del mabandīd. Fa-wa-llāh al-yawm
nofūs-e khabīthe żoʿafā rā be awhām-e qablīye az sharīʿe-ye elāhīye manʿ
namūde-and.

One day I was present in the court of the throne. The Face of God turned
toward me and said, ``Have you heard that the sincere are in great
danger? O servant present before the throne! The outwardly learned have
disagreed about the companions of the Messenger, as to whom the
designation `companion' applies. Some have said that companions are
those souls who were in His presence for a year or more and arose to
wage holy war in the path of God. Others have said that whoever
professed the declaration of God's oneness and met the Messenger is
counted among the companions, even if this happened only once. Others
have said that this designation applies to those souls whom the
Messenger Himself specifically addressed by this designation. Still
others have said that they are those souls who were trusted by Him and
were present both on journeys and at home. Most of the learned, however,
have said, `Whoever embraced Islam {[}corrupted original text{]} the
Prophet, may God bless Him, and accompanied Him, even for the shortest
time, is one of the companions.' Accordingly, at the Farewell Pilgrimage
forty thousand souls were with Him, and on the day of His passing one
hundred and twenty-four thousand people gathered in Medina. According to
the last view, the designation `companion' applies to all these souls.
Nevertheless, only a few had drunk of the limpid wine of certitude and
attained the Source of the outpourings of {[}corrupted original text{]}.
In short, such words have arisen, and continue to arise, from
heedlessness. We ask God to assist everyone to do what He loves and
approves. O friends of God, do not content yourselves with hearsay, and
do not set your hearts on the delusions of the past. By God, today
wicked souls have kept the weak away from the divine watering place
through the delusions of former times.''

\subsection{101}\label{section-3220}

\begin{otherlanguage}{persian}

بشنوید عرض این عبد خادم لدی العرش را و با رجل مستقیمه بر صراط احدیّه قائم
شوید حقّ بشأنی ظاهر که احدی را مجال اعراض نمانده جمیع کتب الهیّه مشعر و
مدلّ بر آن چنانچه چندی قبل یکی از دوستان که از ملل مختلفه بوده عریضه‌ئی
عرض نموده و در آن عریضه دو روایت معروض داشته که در کتب قدیمه از لسان
یونانی بلسان عربی ترجمه شده

\end{otherlanguage}

Beshnavīd ʿarż-e īn ʿabd-e khādem-e laday al-ʿarsh rā va bā rejl-e
mostaqīme bar ṣerāṭ-e aḥadīye qāʾem shavīd. Ḥaqq be shaʾnī ẓāher ke
aḥadī rā majāl-e eʿrāż namānde. Jamīʿ-e kotob-e elāhīye moshʿer va
modell bar ān, chenānche chandī qabl yekī az dūstān ke az melal-e
mokhtalefe būde ʿarīżeʾī ʿarż namūde va dar ān ʿarīże do revāyat maʿrūż
dāshte ke dar kotob-e qadīme az lesān-e yūnānī be lesān-e ʿarabī tarjome
shode.

Listen to the appeal of this servant who serves before the throne, and
stand with steadfast feet upon the path of oneness. God has appeared in
such a manner that no one has any room left to turn away. All the divine
books point to this and attest to it. Some time ago, for example, one of
the friends who had belonged to the various religious communities
submitted a letter, and in that letter presented two traditions
translated from Greek into Arabic in the ancient books.

\subsection{102}\label{section-3221}

\begin{otherlanguage}{persian}

فقرۀ اولی سیظهر الشّیطان فی جزیرة قاف و یمنع النّاس عن الرّحمن اذا حان ذاک
الحین توجّهوا الی الأرض المقدّسة منها تمرّ نسمة الله انتهی و قاف قبرص است و
این مشهور است چنانچه جمیع ترک قبرص را شیطان جزیره‌سی میگویند و ارض مقدّسه
هم معلوم که حال مقرّ عرش واقع شده

\end{otherlanguage}

Faqare-ye ūlā: ``Sa-yaẓharu al-shayṭānu fī jazīrati Qāf wa-yamnaʿu
al-nāsa ʿan al-Raḥmān. Idhā ḥāna dhāka al-ḥīnu tawajjahū ilā al-arḍi
al-muqaddasa, minhā tamurru nasmatu Allāh.'' Intahā. Va Qāf Qebres ast
va īn mashhūr ast, chenānche jamīʿ-e Tork Qebres rā ``şeyṭān cezīresi''
mīgūyand. Va arż-e moqaddase ham maʿlūm ke ḥāl maqarr-e ʿarsh vāqeʿ
shode.

The first passage: ``Satan will appear on the island of Qāf and keep
people away from the All-Merciful. When that time comes, turn toward the
Holy Land; from it the breeze of God passes.'' End of quotation. Qāf is
Cyprus, and this is well known, for all the Turks call Cyprus ``Satan's
island.'' The Holy Land is also known: it is where the seat of the
throne is now established.

\subsection{103}\label{section-3222}

\begin{otherlanguage}{persian}

و فقرۀ ثانی یظهر الحباب فی جزیرة المنسوبة الیه انّه قصیر القامة کثیر
اللّحیة ضیّق الجبهة و الصّدر اصفر العین و الشّعر لظهره وبر کالابل و لصدره
شعر کالمعز اذا اتی ذلک الوقت تقرّبوا الی الکرمل ولو بالکلکل ثمّ اقبلوا الی
الواد المقدّس ارض المحشر بقعة البیضآء انتهی معلوم بوده که حباب اسم شیطان
وحیّه است میفرماید ظاهر میشود شیطان در جزیره‌ئی که منسوب باوست که قبرص
باشد چنانچه بجزیرۀ شیطان معروفست و میفرماید اذا اتی ذلک الوقت تقرّبوا الی
الکرمل و کرمل جبلیست مقابل عکّا ولو بالکلکل اگرچه بسینه باشد ثمّ اقبلوا
الی الواد المقدّس ارض المحشر بقعة البیضآء این سه لقب ارض عکّاست چنانچه بین
کلّ مشهور و در کتب مذکور و کاش نفسی به قبرص میرفت و جمیع این صفات که
مذکور است بعین ظاهر در آن شخص مشاهده مینمود

\end{otherlanguage}

Va faqare-ye thānī: ``Yaẓharu al-Ḥubābu fī jazīrati al-mansūbati ilayhi.
Innahu qaṣīru al-qāma, kathīru al-liḥya, ḍayyiqu al-jabhati wa-al-ṣadr,
aṣfaru al-ʿayni wa-al-shaʿr. Li-ẓahrihi wabarun ka-al-ibil wa-li-ṣadrihi
shaʿrun ka-al-maʿz. Idhā atā dhālika al-waqtu taqarrabū ilā al-Karmal
wa-law bi-al-kalkal, thumma aqbilū ilā al-wādi al-muqaddas, arḍi
al-maḥshar, buqʿati al-bayḍāʾ.'' Intahā. Maʿlūm būde ke Ḥobāb esm-e
shayṭān va ḥayye ast. Mīfarmāyad ẓāher mīshavad shayṭān dar jazīreʾī ke
mansūb be ūst ke Qebres bāshad, chenānche be jazīre-ye shayṭān maʿrūf
ast. Va mīfarmāyad ``idhā atā dhālika al-waqtu taqarrabū ilā al-Karmal''
va Karmel jabalīst moqābel-e ʿAkkā, ``wa-law bi-al-kalkal,'' agarche be
sīne bāshad, ``thumma aqbilū ilā al-wādi al-muqaddas, arḍi al-maḥshar,
buqʿati al-bayḍāʾ.'' Īn se laqab-e arż-e ʿAkkāst, chenānche beyn-e koll
mashhūr va dar kotob madhkūr. Va kāsh nafsī be Qebres mīraft va jamīʿ-e
īn ṣefāt ke madhkūr ast be ʿeyn-e ẓāher dar ān shakhṣ moshāhede mīnamūd.

And the second passage: ``Ḥubāb will appear on the island associated
with him. He is short in stature, thick-bearded, narrow of forehead and
chest, with yellow eyes and hair. On his back is down like a camel's,
and on his chest hair like a goat's. When that time comes, draw near to
Carmel, even on your breast; then turn toward the holy valley, the land
of the gathering, the white spot.'' End of quotation. It has been known
that Ḥubāb is a name for Satan and a serpent. It says that Satan will
appear on the island associated with him, which is Cyprus, since it is
known as Satan's island. And it says, ``When that time comes, draw near
to Carmel''---and Carmel is a mountain facing Acre---``even on your
breast,'' that is, even if it be on your chest; ``then turn toward the
holy valley, the land of the gathering, the white spot.'' These are
three designations of the land of Acre, as is well known to all and
recorded in the books. Would that someone would go to Cyprus and observe
with his outward eyes all these stated characteristics in that person.

\subsection{104}\label{section-3223}

\begin{otherlanguage}{persian}

باری ای عباد حقّ محتاج باین اذکار نیست و امثال این اذکار لأجل تفضّل صرفه
است که این بندۀ ذلیل گمان نموده که بامثال آن بعضی از خواب غفلت بیدار
شوند والّا حقّ مقدّس از ذکر این و آن لا یعرف الّا بنفسه و لا ینعت الّا بما
نعت به ذاته لذاته ما سواه مخلوق بأمره و راجع الی اماکنه فی عوالم ابداعه
و اختراعه

\end{otherlanguage}

Bārī ey ʿebād, Ḥaqq moḥtāj be īn adhkār nīst va amthāl-e īn adhkār
li-ajl-e tafażżol-e ṣerfe ast ke īn bande-ye dhalīl gomān namūde ke be
amthāl-e ān baʿżī az khāb-e ghaflat bīdār shavand. Va-ellā Ḥaqq moqaddas
az dhekr-e īn va ān. Lā yuʿrafu illā bi-nafsihi wa-lā yunʿatu illā bi-mā
naʿata bihi dhātahu li-dhātihi. Mā siwāhu makhlūqun bi-amrihi wa-rājiʿun
ilā amākinihi fī ʿawālimi ibdāʿihi wa-ikhtirāʿihi.

In short, O servants, God has no need of these statements. Such
statements are offered solely as a favor, for this lowly servant has
thought that through such things some might awaken from the sleep of
heedlessness. Otherwise, God is sanctified above the mention of this
person or that. He is known only through Himself and described only by
that with which He has described His own Self to His own Self. All
besides Him is created by His command and returns to its places in the
worlds of His origination and creation.

\subsection{105}\label{section-3224}

\begin{otherlanguage}{persian}

امّیدواریم که عرایض این عبد چون خالصاً لوجه الله معروض شده شفای قلوب واقع
شود و ضیاء صدور تا کلّ بحبّ الله بر شأنی قیام نمایند که احدی را مجال اعراض
و اعتراض نماند و جمیع ما قاله المشرکون او یقولون را لاشیء محض انگارند
اذاً اقول تبت الیک یا الهی بما اجترحت فی ساحتک فاغفر لی بسلطانک و فضلک
انّک انت اکرم الأکرمین و الحمد لک یا اله العالمین

\end{otherlanguage}

Ommīdvārīm ke ʿarāyeż-e īn ʿabd chon khāliṣan li-wajhi Allāh maʿrūż
shode shefā-ye qolūb vāqeʿ shavad va żīyā-ye ṣodūr tā koll be ḥobb-e
Allāh bar shaʾnī qiyām namāyand ke aḥadī rā majāl-e eʿrāż va eʿterāż
namānad va jamīʿ-e ``mā qālahu al-mushrikūna aw yaqūlūn'' rā lā-shayʾ-e
maḥż engārand. Idhan aqūlu tubtu ilayka yā ilāhī bi-mā ijtaraḥtu fī
sāḥatika fa-ighfir lī bi-sulṭānika wa-faḍlika innaka anta akramu
al-akramīn wa-al-ḥamdu laka yā ilāha al-ʿālamīn.

We hope that, since this servant's appeals have been submitted purely
for the sake of God, they may become a healing for hearts and a light
for breasts, so that all may arise in the love of God in such a manner
that no one has any room left to turn away or object, and may regard all
that the idolaters have said or may say as sheer nothingness. Then I
say: I turn to You in repentance, O my God, for what I have committed in
Your court. Forgive me, then, through Your sovereignty and Your grace.
You are indeed the Most Generous of the generous, and praise be to You,
O God of the worlds.

\subsection{106}\label{section-3225}

\begin{otherlanguage}{arabic}

خادم‌الله

\end{otherlanguage}

Khādem-e Allāh

The servant of God

\bookmarksetup{startatroot}

\chapter{BH11924 (Unnamed)}\label{bh11924-unnamed}

\emph{Translated on 2 Qudrat, 183 BE using GPT 6.1 Sol}

\subsection{1}\label{section-3226}

\begin{otherlanguage}{arabic}

يا اسمى لعمر الله احزننى حزنك انت تعلم والله يعلم انّ المظلوم احبّك وانزل
لك ما لا ينقطع عَرفه بدوام ملكوت الله وجبروته واسمآئه وصفاته

\end{otherlanguage}

Yā ismī, la-ʿamru Allāh, aḥzananī ḥuznuka. Anta taʿlamu wa-Allāhu
yaʿlamu anna al-Maẓlūma aḥabbaka wa-anzala laka mā lā yanqaṭiʿu ʿarfuhu
bi-dawāmi malakūti Allāhi wa-jabarūtihi wa-asmāʾihi wa-ṣifātihi.

O My Name, by the life of God, your sorrow has saddened Me. You know,
and God knows, that the Wronged One loves you and has sent down for you
that whose fragrance will endure as long as God's Kingdom and Dominion,
His names and His attributes.

\subsection{2}\label{section-3227}

\begin{otherlanguage}{arabic}

انظر ثم اذكر العنايات الّتى قدّرناها لك وكن من الشّاكرين

\end{otherlanguage}

Unẓur thumma udhkur al-ʿināyāti allatī qaddarnāhā laka wa-kun mina
al-shākirīn.

Consider, then remember the favors We have ordained for you, and be
among the thankful.

\subsection{3}\label{section-3228}

\begin{otherlanguage}{arabic}

اوّل الغاية من هدايتك الى طراط الله ونبأه العظيم والثّانى ايّدناك على
الأقبال والثّالث امددناك بجنود الغيب والشّهادة الى ان حضرت وفزت بما كان
مرقوما من القلم الأعلى فى كتب المرسلين والرّابع اسمعناك ندآئى فى الحضور
والخامس اريناك افقى المنير والسّادس انزلنا لك مآئدة البيان من سمآء رحمة
اسمى الرّحمن وجعلناك من المقرّبين والسّابع فتحنا على وجهك باب الرّزق والغناء
واطعمناك من طعام البرهان واسقيناك كوثر العرفان بعد ما اظهرناك وربّيناك
بايادى عنايتى ورحمتى

\end{otherlanguage}

Awwalu al-ghāyati min hidāyatika ilā {[}corrupted text{]} Allāhi
wa-nabaʾihi al-ʿaẓīm, wa-al-thānī ayyadnāka ʿalā al-iqbāl, wa-al-thālith
amdadnāka bi-junūdi al-ghaybi wa-al-shahādati ilā an ḥaḍarta wa-fuzta
bi-mā kāna marqūman mina al-qalami al-aʿlā fī kutubi al-mursalīn,
wa-al-rābiʿ asmaʿnāka nidāʾī fī al-ḥuḍūr, wa-al-khāmis araynāka ufuqī
al-munīr, wa-al-sādis anzalnā laka māʾidata al-bayāni min samāʾi raḥmati
ismī al-Raḥmān wa-jaʿalnāka mina al-muqarrabīn, wa-al-sābiʿ fataḥnā ʿalā
wajhika bāba al-rizqi wa-al-ghināʾ wa-aṭʿamnāka min ṭaʿāmi al-burhān
wa-asqaynāka kawthara al-ʿirfān baʿda mā aẓharnāka wa-rabbaynāka
bi-ayādī ʿināyatī wa-raḥmatī.

The first was to guide you to {[}corrupted original text{]} of God and
His momentous announcement; the second, We strengthened you to turn
toward Him; the third, We aided you with the hosts of the unseen and the
visible until you came into Our presence and attained what had been
inscribed by the Most Exalted Pen in the books of the messengers; the
fourth, We let you hear My call in Our presence; the fifth, We showed
you My luminous horizon; the sixth, We sent down for you the feast of
utterance from the heaven of the mercy of My name, the All-Merciful, and
made you one of those brought near; and the seventh, We opened before
you the door of sustenance and abundance, fed you the food of proof, and
gave you the choice wine of recognition to drink, after bringing you
forth and nurturing you with the hands of My care and mercy.

\subsection{4}\label{section-3229}

\begin{otherlanguage}{arabic}

انّك ان شكر ما رقم لك فى هذا الحين يشهد لك بجوارحك واركانك حتّى شطراتك قم
وقل اعترفت بعطاياك يا مقصود من فى السّموات والأرضين واقررت بنعمائك وألائك
يا محبوب افئدة المخلصين اشهد انّك دعوتنى الى مقامك الأعلى وأويتنى فى جوار
رحمتك الكبرى وسقيتنى سلسبيل الحبّ والّلقآء وعزّتك يا اله الاسمآء وفاطر
السّمآء انّ عبدك هذا لا يقدر ان يحصى نعمائك وكيف شكرك يا مالك ملكوت الحجّة
والبرهان

\end{otherlanguage}

Innaka in shakara mā ruqima laka fī hādhā al-ḥīn yashhadu laka
bi-jawāriḥika wa-arkānika ḥattā shaṭarātika. Qum wa-qul: Iʿtaraftu
bi-ʿaṭāyāka yā Maqṣūda man fī al-samāwāti wa-al-arḍīn, wa-aqrartu
bi-naʿmāʾika wa-ālāʾika yā Maḥbūba afʾidati al-mukhliṣīn. Ashhadu annaka
daʿawtanī ilā maqāmika al-aʿlā wa-āwaytanī fī jiwāri raḥmatika al-kubrā
wa-saqaytanī salsabīla al-ḥubbi wa-al-liqāʾ. Wa-ʿizzatika yā Ilāha
al-asmāʾi wa-Fāṭira al-samāʾ, inna ʿabdaka hādhā lā yaqdiru an yuḥṣiya
naʿmāʾaka wa-kayfa shukruka yā Mālika malakūti al-ḥujjati wa-al-burhān.

If gratitude for what has been inscribed for you at this moment could
testify on your behalf, your limbs and members, even your every part,
would bear witness. Rise and say: ``I acknowledge Your gifts, O Desire
of all who are in the heavens and on earth, and I confess Your blessings
and favors, O Beloved of the hearts of the sincere. I testify that You
called me to Your most exalted station, sheltered me within the bounds
of Your most great mercy, and gave me the purest draught of love and
reunion to drink. By Your might, O God of names and Creator of heaven,
this servant of Yours cannot count Your blessings---how, then, can he
thank You, O Lord of the Kingdom of testimony and proof?''

\subsection{5}\label{section-3230}

\begin{otherlanguage}{arabic}

يا اسمى ونفسى لو اذكر عناياتى ومقاماتها الّتى اظهرناها لك لينتهى الأقلام
والألواح قبل ان ينتهى فضآئلى وذكر نعمآئى التى قدرناها لك بامرك المحكم
المتين يا اسمى لمليك عناية التى وسلامه كمان اكمه بل يقين الكمه

\end{otherlanguage}

Yā ismī wa-nafsī, law adhkuru ʿināyātī wa-maqāmātihā allatī aẓharnāhā
laka la-yantahī al-aqlāmu wa-al-alwāḥu qabla an yantahiya faḍāʾilī
wa-dhikru naʿmāʾī allatī qaddarnāhā laka bi-amrika al-muḥkami al-matīn.
Yā ismī, {[}corrupted text{]}

O My Name and My Self, were I to recount My favors and the stations of
those favors that We have made manifest to you, the pens and tablets
would be exhausted before My bounties and the account of the blessings
We have ordained for you through your firm and steadfast command came to
an end. O My Name, {[}corrupted original text{]}

\subsection{6}\label{section-3231}

\begin{otherlanguage}{persian}

شما عارفيد بر عنايات حق ونعمتا ئيكه عطا فرموند بيجين فعلتها ئيكه مخصوص
شما از بيت العطا ظاهر هرنمود از اين نائمه مرسله لسيار محزون شدم واكريم
برفوض اقرار كنيم برانكه عباد بكلمه نا لايقه ويا بعبارات غير مرضيّه نطق
نمود انجه اين مسلّم است كه من غير ار انجه ومشيّت حق نطق نمواند ور اينصورت
جار شكر است ومقام فرح است نه شكايت از آبق كه شتد

\end{otherlanguage}

Shumā ʿārifīd bar ʿināyāt-i Ḥaqq va {[}corrupted text{]} ʿaṭā farmūnd
{[}corrupted text{]} makhṣūṣ-i shumā az Bayt al-ʿAṭā ẓāhir {[}corrupted
text{]}. Az īn {[}corrupted text{]}-yi mursala {[}corrupted text{]}
maḥzūn shudam va {[}corrupted text{]} iqrār kunīm bar ānkih ʿibād
bi-kalima-yi nālāyiqa va yā bi-ʿibārāt-i ghayr-i marḍiyya nuṭq namūd.
{[}corrupted text{]} īn musallam ast kih {[}corrupted text{]} va
mashiyyat-i Ḥaqq nuṭq {[}corrupted text{]} va dar īn ṣūrat {[}corrupted
text{]}-i shukr ast va maqām-i faraḥ ast, nah shikāyat az {[}corrupted
text{]}.

You are aware of the favors of the True One and {[}corrupted original
text{]} bestowed {[}corrupted original text{]} especially for you,
manifested from the House of Bounty {[}corrupted original text{]}. I was
deeply saddened by this {[}corrupted original text{]} sent {[}corrupted
original text{]}, and {[}corrupted original text{]} acknowledge that
people spoke an unseemly word or used expressions that were not
pleasing. {[}corrupted original text{]} this is certain: {[}corrupted
original text{]} and the will of the True One spoke {[}corrupted
original text{]}, and in this case there is {[}corrupted original
text{]} for gratitude and an occasion for joy, not complaint about
{[}corrupted original text{]}.

\subsection{7}\label{section-3232}

\begin{otherlanguage}{persian}

آيا نشنيد ركه در آستانه ذيل تقدير رابمفتريات خط آلحصراند در حودم سرقتر
باولياى حق لشبت ولطر وسيدميند وسر جعشر راعكه قرار ولصراند جناب وقت ان
است كه انجناب بر نصرت قيام نمايند قدورد على بشهادتك ما ناح به قلبى وصاح
به قلمى جه كه انجه واردشده سبب تضييع امر الله بحصر ناس ضعيفند براسل امور
رآكاه شوند قد اتى از مان

\end{otherlanguage}

Āyā nashnīd {[}corrupted text{]} dar āstāna {[}corrupted text{]}
bi-awliyā-yi Ḥaqq {[}corrupted text{]}. Janāb, vaqt ān ast kih ān janāb
bar nuṣrat qiyām namāyand. {[}corrupted text{]} bi-shahādatika mā nāḥa
bihi qalbī wa-ṣāḥa bihi qalamī. {[}corrupted text{]} {[}corrupted
text{]} vārid shuda sabab-i taḍyīʿ-i Amr Allāh {[}corrupted text{]} nās
ḍaʿīfand {[}corrupted text{]} qad atā {[}corrupted text{]}.

Have you not heard {[}corrupted original text{]} at the threshold
{[}corrupted original text{]} with the friends of God {[}corrupted
original text{]}? Sir, the time has come for you to arise in support.
{[}corrupted original text{]} By your testimony, what my heart lamented
and my pen cried out. {[}corrupted original text{]} {[}corrupted
original text{]} what has occurred has caused the Cause of God to be
lost {[}corrupted original text{]} the people are weak {[}corrupted
original text{]} has come {[}corrupted original text{]}.

\subsection{8}\label{section-3233}

\begin{otherlanguage}{persian}

بايد نسئل الله ان يؤيّد الكلّ على ما يحبّ ويرضى

\end{otherlanguage}

Bāyad nasʾalu Allāha an yuʾayyida al-kulla ʿalā mā yuḥibbu wa-yarḍā.

We must ask God to strengthen everyone to do what He loves and approves.

\subsection{9}\label{section-3234}

\begin{otherlanguage}{arabic}

البهآء والذّكر والثّنآء عليك وعلى من يحبّك لوجه الله العزيز الحميد

\end{otherlanguage}

Al-bahāʾu wa-al-dhikru wa-al-thanāʾu ʿalayka wa-ʿalā man yuḥibbuka
li-wajhi Allāhi al-ʿAzīzi al-Ḥamīd.

Glory, remembrance, and praise be upon you and upon whoever loves you
for the sake of God, the Almighty, the All-Praised.


\backmatter


\end{document}
